\documentclass{book}
%\usepackage[letterpaper, portrait, margin=1cm]{geometry}
%\usepackage[letterpaper, bindingoffset=0.2in, left=1in,right=1in,top=.5in,bottom=.5in,footskip=.25in,marginparwidth=5em]{geometry}
\usepackage[letterpaper, left=1in,right=1in,top=.5in,bottom=.5in,footskip=.25in,marginparwidth=1cm]{geometry}
% ---------------------
% mini-table-of-content
% ---------------------
\usepackage{minitoc}
\setcounter{minitocdepth}{1}
\setlength{\mtcindent}{24pt}
\setcounter{secnumdepth}{-2}
%\renewcommand{\mtcfont}{\small\rm}
%\renewcommand{\mtcSfont}{\small\bf}
%\usepackage{setspace}
%\usepackage{tocloft}
%\setlength\cftparskip{-1.2pt}
%\setlength\cftbeforesecskip{1.3pt}
%\setlength\cftaftertoctitleskip{2pt}
%\renewcommand{\cftsecafterpnum}{\hspace*{02.0em}}
%\renewcommand{\cftsubsecafterpnum}{\hspace*{02.0em}}

% ---------------------------
% Chinese Characters Packages
% ---------------------------
\usepackage{fontspec} 
\usepackage{xeCJK}
\setmainfont{Times}
\setCJKmainfont{BiauKai}
\newfontfamily\sblgoodhebrew{SBL BibLit}[Script=Hebrew,Contextuals=Alternate]
\newfontfamily\sblgoodgreek{SBL BibLit}[Script=Greek,Contextuals=Alternate]

\usepackage{ifpdf,cite,algorithmic,url,tikz}
\usepackage[cmex10]{amsmath}

% ---------------------------
% Hebrew Characters Packages
% ---------------------------
\usepackage{polyglossia}
\setmainfont{Times New Roman}

% -------
% General
% -------
\usepackage{multicol}
\usepackage{multirow}
\usepackage{color,colortbl}
\usepackage{xparse}
\usepackage{pbox}
\usepackage{stackengine}
\usepackage{titlesec}% http://ctan.org/pkg/titlesec
\usepackage{tabularx}
\usepackage{xltabular}
\usepackage{titlesec}
\usepackage{makecell}
\newcommand{\sectionbreak}{\clearpage}

\author{
  Editor, Michael Chan\\
  \texttt{michaelchan\_wahyan@yahoo.com.hk}
}
\usepackage{tocloft}

\usepackage{hyperref}
\hypersetup{
    colorlinks=true, % set true if you want colored links
    linktoc   =all , % set to all if you want both sections and subsections linked
    linkcolor =blue, % choose some color if you want links to stand out
}

% ----------
% Afterword
% ----------
\usepackage{marginnote}
\usepackage{sectsty}
\usepackage{ragged2e}
\usepackage{lineno}
\usepackage{xcolor}
\usepackage{paracol}

\begin{document}

\clearpage
%% temporary titles
% command to provide stretchy vertical space in proportion
\newcommand\nbvspace[1][3]{\vspace*{\stretch{#1}}}
% allow some slack to avoid under/overfull boxes
\newcommand\nbstretchyspace{\spaceskip0.5em plus 0.25em minus 0.25em}
% To improve spacing on titlepages
\newcommand{\nbtitlestretch}{\spaceskip0.6em}
\pagestyle{empty}
\begin{center}
\bfseries
\nbvspace[1]
\Huge
{%\nbtitlestretch
\Large
\textbf{宣道會錦繡堂 粵語講道逐字稿 2023-28 \\
       Youtube Channel: Fairview Church
       }}

\nbvspace[1]

{\large
Editor: Michael\\
\texttt{michaelchan\_wahyan@yahoo.com.hk}
}

\nbvspace[1]

{\large
Revision: \texttt{v1.3}\\
Last Update: \today
}


\vfill
\begin{tikzpicture}
    %remove comment for OT cover%\node (0,0) [opacity=0.03]{\includegraphics[width=15cm]{../bible_out/ot_frontcover.png}} ;
    %remove comment for NT cover%\node (0,0) [opacity=0.03]{\includegraphics[width=15cm]{../bible_out/christ_on_cross.png}} ;
    %remove comment for Bible cover%\node (0,0) [xshift=0.8cm, yshift=+2cm, opacity=0.03]{\includegraphics[width=10cm]{./christ_on_cross.png}} ;
    %remove comment for Bible cover%\node (0,0) [              yshift=-2cm, opacity=0.03]{\includegraphics[width=14cm]{./ot_frontcover.png}} ;
\end{tikzpicture}
\vfill

\end{center}

\newpage

\setcounter{tocdepth}{0}
\dominitoc
\begin{multicols}{3}
\addtocontents{toc}{\protect\hypertarget{toc}{}}
\tableofcontents
\end{multicols}

\large
%\twocolumn

% the color definition syntax is as follow:
% \definecolor{name}{system}{definition}
% example: a mono-channel color can be defined as
%          \definecolor{Gray}{gray}{0.9}
% example: an rgb-3-channel color can be defined as
%          \definecolor{LightCyan}{rgb}{0.88,1,1}
%          \definecolor{pink}{rgb}{0.68,0,0.68}

\definecolor{CUV1LightRed}{rgb}{1,0.75,0.75}     % for CUV1
\definecolor{LZZVLightGray}{rgb}{0.9,0.9,0.9}    % for LZZ
\definecolor{KJVVLightGreen}{rgb}{0.75,1,0.85}   % for KJV
\definecolor{CUV2LightYellow}{rgb}{1,1,0.75}     % for CUV2
\definecolor{CNVVLightBrown}{rgb}{1,0.85,0.7}    % for CNV
\definecolor{NRSVLightBlue}{rgb}{0.75,1,1}       % for NRSV
\definecolor{WENLLightPurple}{rgb}{0.95,0.85,0.9}% for WENL
\definecolor{TCV19PaleGreen}{rgb}{0.85,1,0.95}   % for TCV19
\definecolor{MSGVLightWhite}{rgb}{0.98,0.98,0.98}% for MSGV
\definecolor{NETSLightRed}{rgb}{1,0.75,0.75}     % for NETS
\definecolor{JPS1917LightYellow}{rgb}{1,1,0.75}  % for JPS1917
\definecolor{SBLGNTPaleRed}{rgb}{1,0.85,0.80}    % for SBLGNT

\section{目錄\small{(順時)}}
\label{sec:index_chronic}
{ \scriptsize


\begin{xltabular}{\textwidth}{|p{0.15\textwidth} p{0.6\textwidth}|p{0.07\textwidth} p{0.1\textwidth}|}
\hline
啟示錄 21:1-7 & \hyperref[sec:IlLnn7WOgiE]{2023-01-01 加恩樓《上帝給我們的盼望及出路》啟示錄 21\_1-7 講員:鄧泰康傳道} & 2023-01-03 & \href{https://youtube.com/watch?v=IlLnn7WOgiE}{\texttt{IlLnn7WOgiE}} \\
路加福音 6:43-49 & \hyperref[sec:d2MEojJVUd4]{2023-01-01 午堂《趁新一年,遵主話而行》路加福音 6\_43-49 講員:曾裔貴牧師} & 2023-01-03 & \href{https://youtube.com/watch?v=d2MEojJVUd4}{\texttt{d2MEojJVUd4}} \\
  & \hyperref[sec:BQHcR2q2KwA]{2023-01-05 周四晚查經班《被擄歸回歷史研讀(2)》 講員:曾裔貴牧師} & 2023-01-06 & \href{https://youtube.com/watch?v=BQHcR2q2KwA}{\texttt{BQHcR2q2KwA}} \\
路加福音 8:1-15 & \hyperref[sec:gAWlXGsKQ90]{2023-01-07 晚堂《存忍耐的心結出果實》路加福音 8\_1-15 講員:曾裔貴牧師} & 2023-01-10 & \href{https://youtube.com/watch?v=gAWlXGsKQ90}{\texttt{gAWlXGsKQ90}} \\
撒迦利亞書 2:6-8 & \hyperref[sec:ejP2s680TsQ]{2023-01-08 加恩樓《上帝眼中的瞳人》撒迦利亞書 2\_6-8 講員:陳韋安牧師} & 2023-01-10 & \href{https://youtube.com/watch?v=ejP2s680TsQ}{\texttt{ejP2s680TsQ}} \\
提摩太後書 4:6-22 & \hyperref[sec:uJ8mvXZ5iqw]{2023-01-08 午堂《心友靈情》提摩太後書 4\_6-22 講員:曹偉彤院長} & 2023-01-10 & \href{https://youtube.com/watch?v=uJ8mvXZ5iqw}{\texttt{uJ8mvXZ5iqw}} \\
羅馬書 12:1-2 & \hyperref[sec:dkl9UnCmWfQ]{2019-01-20 加恩樓崇拜《得勝秘笈》羅馬書12\_1-2 講員:鄧泰康傳道} & 2023-01-13 & \href{https://youtube.com/watch?v=dkl9UnCmWfQ}{\texttt{dkl9UnCmWfQ}} \\
  & \hyperref[sec:ZwtOgh6Qmm4]{2023-01-12 周四晚查經班《被擄歸回歷史研讀(3)》 講員:曾裔貴牧師} & 2023-01-13 & \href{https://youtube.com/watch?v=ZwtOgh6Qmm4}{\texttt{ZwtOgh6Qmm4}} \\
馬太福音 13:24-30-36-43 & \hyperref[sec:c4O5F1DLbjI]{2023-01-15 加恩樓《成為麥子被收割》馬太福音 13\_24-30,36-43 講員:陳麗蟬牧師} & 2023-01-17 & \href{https://youtube.com/watch?v=c4O5F1DLbjI}{\texttt{c4O5F1DLbjI}} \\
以賽亞書 40:1-11-27-31 & \hyperref[sec:ezP3uupXjgk]{2023-01-15 午堂《如鷹展翅上騰》以賽亞書 40\_1-11,27-31 講員:郭鴻標牧師} & 2023-01-26 & \href{https://youtube.com/watch?v=ezP3uupXjgk}{\texttt{ezP3uupXjgk}} \\
羅馬書 12:13-21 & \hyperref[sec:ApJnSQORNII]{2023-01-22 加恩樓《只要祝福,不要咒詛》羅馬書 12\_13-21 講員:鄧泰康傳道} & 2023-01-26 & \href{https://youtube.com/watch?v=ApJnSQORNII}{\texttt{ApJnSQORNII}} \\
列王記上  & \hyperref[sec:9IcRN7FThks]{2023-01-22 新春崇拜《蒙福秘訣》列王記上 3: 3-9 講員:陳一華牧師} & 2023-01-26 & \href{https://youtube.com/watch?v=9IcRN7FThks}{\texttt{9IcRN7FThks}} \\
馬可福音 7:1-23 & \hyperref[sec:S5k9W4TGOeU]{2020-10-03 晚堂《用嘴來尊敬神的偽善者》馬可福音 7\_1-23 講員:高銘謙牧師} & 2023-01-27 & \href{https://youtube.com/watch?v=S5k9W4TGOeU}{\texttt{S5k9W4TGOeU}} \\
  & \hyperref[sec:mXSWAPPoVUc]{2023-01-26 周四晚查經班《被擄歸回歷史研讀(4)》 講員:曾裔貴牧師} & 2023-01-27 & \href{https://youtube.com/watch?v=mXSWAPPoVUc}{\texttt{mXSWAPPoVUc}} \\
哥林多後書 8:9 & \hyperref[sec:4sQReCkRl1k]{2023-01-28 晚堂《無比恩典》哥林多後書 8\_9 講員:陳一華牧師} & 2023-01-31 & \href{https://youtube.com/watch?v=4sQReCkRl1k}{\texttt{4sQReCkRl1k}} \\
馬太福音 18:21-35 & \hyperref[sec:65Ibdpx_Vbg]{2023-01-29 加恩樓《接納和饒恕》馬太福音 18\_21-35 講員:謝靄薇牧師} & 2023-02-01 & \href{https://youtube.com/watch?v=65Ibdpx-Vbg}{\texttt{65Ibdpx-Vbg}} \\
約翰福音 12:20-26 & \hyperref[sec:fpQjtet_7Iw]{2023-01-29 午堂《藉門徒見耶穌》約翰福音 12\_20-26 講員:林誠信牧師} & 2023-02-02 & \href{https://youtube.com/watch?v=fpQjtet-7Iw}{\texttt{fpQjtet-7Iw}} \\
  & \hyperref[sec:i9iPPyuXEWw]{2023-02-02 周四晚查經班《被擄歸回歷史研讀(5)》 講員:曾裔貴牧師} & 2023-02-04 & \href{https://youtube.com/watch?v=i9iPPyuXEWw}{\texttt{i9iPPyuXEWw}} \\
馬可福音  & \hyperref[sec:2Us7kVIsCO8]{2023-02-05 加恩樓《是賊窩?還是神的殿!!》馬可福音 11: 12-25 講員:鄧泰康傳道} & 2023-02-07 & \href{https://youtube.com/watch?v=2Us7kVIsCO8}{\texttt{2Us7kVIsCO8}} \\
馬可福音  & \hyperref[sec:k6XACT4Tfao]{2023-02-05 午堂《有果效的門徒生活》馬可福音 11: 12-26 講員:曾裔貴牧師} & 2023-02-07 & \href{https://youtube.com/watch?v=k6XACT4Tfao}{\texttt{k6XACT4Tfao}} \\
  & \hyperref[sec:lrudeppALpc]{2023-02-09 周四晚查經班《被擄歸回歷史研讀(6)》 講員:曾裔貴牧師} & 2023-02-10 & \href{https://youtube.com/watch?v=lrudeppALpc}{\texttt{lrudeppALpc}} \\
馬可福音 10:13-16 & \hyperref[sec:P577AmrjvF0]{2023-02-12 加恩樓《讓小孩子「到我這裡來」》馬可福音 10\_13-16 講員:吳景霖傳道} & 2023-02-14 & \href{https://youtube.com/watch?v=P577AmrjvF0}{\texttt{P577AmrjvF0}} \\
詩篇 139:1-24 & \hyperref[sec:raKQe4zsFjc]{2023-02-12 早堂《心知.神愛》詩篇 139\_1-24 講員:楊明麗姑娘} & 2023-02-14 & \href{https://youtube.com/watch?v=raKQe4zsFjc}{\texttt{raKQe4zsFjc}} \\
詩篇 139:1-24 & \hyperref[sec:la_oxDk2vYY]{2023-02-12 午堂《不分離不拋棄的愛》詩篇 139\_1-24 講員:陳麗蟬牧師} & 2023-02-17 & \href{https://youtube.com/watch?v=la_oxDk2vYY}{\texttt{la\_oxDk2vYY}} \\
詩篇 104:31-35 & \hyperref[sec:oRiBGTYbBfo]{2023-02-18 晚堂《新年願望》詩篇 104\_31-35 講員:張美薇牧師} & 2023-02-21 & \href{https://youtube.com/watch?v=oRiBGTYbBfo}{\texttt{oRiBGTYbBfo}} \\
詩篇 11:1-7 & \hyperref[sec:sY5rw9BFGjc]{2023-02-19 午堂《義路》詩篇 11\_1-7 講員:吳景霖傳道} & 2023-02-21 & \href{https://youtube.com/watch?v=sY5rw9BFGjc}{\texttt{sY5rw9BFGjc}} \\
加拉太書 3:21-29 & \hyperref[sec:jaVgwb2h0Gg]{2023-02-19 加恩樓《律法之上》加拉太書 3\_21-29 講員:陳韋安牧師} & 2023-02-22 & \href{https://youtube.com/watch?v=jaVgwb2h0Gg}{\texttt{jaVgwb2h0Gg}} \\
  & \hyperref[sec:v6FulXmsPK8]{2023-02-23 周四晚查經班《被擄歸回歷史研讀(7)》 講員:曾裔貴牧師} & 2023-02-25 & \href{https://youtube.com/watch?v=v6FulXmsPK8}{\texttt{v6FulXmsPK8}} \\
使徒行傳 8:14-24 & \hyperref[sec:jYaZ2w15SEQ]{2023-02-25 晚堂《屬靈沒有捷徑》使徒行傳 8\_14-24 講員:陳一華牧師} & 2023-02-28 & \href{https://youtube.com/watch?v=jYaZ2w15SEQ}{\texttt{jYaZ2w15SEQ}} \\
帖撒羅尼迦前書 5:16-18 & \hyperref[sec:Lu3g7hRuGdE]{2023-02-26 加恩樓《時刻經歷主的同在》帖撒羅尼迦前書 5\_16-18 講員:李適清博士} & 2023-02-28 & \href{https://youtube.com/watch?v=Lu3g7hRuGdE}{\texttt{Lu3g7hRuGdE}} \\
詩篇 73:1-28 & \hyperref[sec:d9OO9rhlq8k]{2023-02-26 午堂《當信念被動搖》詩篇 73\_1-28 講員:鄧泰康傳道} & 2023-03-01 & \href{https://youtube.com/watch?v=d9OO9rhlq8k}{\texttt{d9OO9rhlq8k}} \\
  & \hyperref[sec:E8bPEPFjp7c]{2023-03-02 周四晚查經班《被擄歸回歷史研讀(8)》 講員:曾裔貴牧師} & 2023-03-03 & \href{https://youtube.com/watch?v=E8bPEPFjp7c}{\texttt{E8bPEPFjp7c}} \\
歌羅西書 1:21-29 & \hyperref[sec:M8De3gUONVA]{2023-03-05 午堂《全力以赴的福音執事》歌羅西書 1\_21-29 講員:曾裔貴牧師} & 2023-03-07 & \href{https://youtube.com/watch?v=M8De3gUONVA}{\texttt{M8De3gUONVA}} \\
詩篇 46:1-11 & \hyperref[sec:BXGw_xefZws]{2023-03-05 加恩樓《上帝是我們患難中隨時的幫助》詩篇 46\_1-11 講員:鄧泰康傳道} & 2023-03-08 & \href{https://youtube.com/watch?v=BXGw_xefZws}{\texttt{BXGw\_xefZws}} \\
創世記 18:18-19 & \hyperref[sec:GHx9ZnTIMjk]{2023-03-11 晚堂《上帝子民要秉公行義》創世記18\_18-19 講員:黃頌言姑娘} & 2023-03-14 & \href{https://youtube.com/watch?v=GHx9ZnTIMjk}{\texttt{GHx9ZnTIMjk}} \\
使徒行傳 9:26-30-11:22-26 & \hyperref[sec:0LbO5gq5ujg]{2023-03-12 午堂《生命師傅的典範》使徒行傳9\_26-30,11\_22-26 講員:陳一華牧師} & 2023-03-14 & \href{https://youtube.com/watch?v=0LbO5gq5ujg}{\texttt{0LbO5gq5ujg}} \\
創世記 18:16-21 & \hyperref[sec:1tknD6q3Urw]{2023-03-12 加恩樓《守主道‧行公義》創世記18\_16-21 講員:陳麗蟬牧師} & 2023-03-15 & \href{https://youtube.com/watch?v=1tknD6q3Urw}{\texttt{1tknD6q3Urw}} \\
歌羅西書 3:1-17 & \hyperref[sec:BtpgDcAKeBQ]{2023-03-19 午堂《天地之間》歌羅西書3\_1-17 講員:吳景霖傳道} & 2023-03-21 & \href{https://youtube.com/watch?v=BtpgDcAKeBQ}{\texttt{BtpgDcAKeBQ}} \\
歌羅西書 3:7-14 & \hyperref[sec:W3jetg3vhiA]{2023-03-19 早堂《信主之後》歌羅西書3\_7-14 講員:楊明麗姑娘} & 2023-03-21 & \href{https://youtube.com/watch?v=W3jetg3vhiA}{\texttt{W3jetg3vhiA}} \\
馬太福音 25:1-13 & \hyperref[sec:wPsCGHG1pLA]{2023-03-19 加恩樓《心理預備,宏觀資訊的裝備,行動準備》馬太福音25\_1-13 講員:郭鴻標牧師} & 2023-03-22 & \href{https://youtube.com/watch?v=wPsCGHG1pLA}{\texttt{wPsCGHG1pLA}} \\
以弗所書 3:14-21 & \hyperref[sec:rNmdtGHU_n0]{2023-03-26 加恩樓《按神的心意成就一切》以弗所書3\_14-21 講員:鄧泰康傳道} & 2023-03-28 & \href{https://youtube.com/watch?v=rNmdtGHU-n0}{\texttt{rNmdtGHU-n0}} \\
提摩太前書 6:3-10 & \hyperref[sec:pzC3e8fjAp8]{2023-03-26 早堂《敬虔與知足》提摩太前書6\_3-10 講員:陳靜娟姑娘} & 2023-03-28 & \href{https://youtube.com/watch?v=pzC3e8fjAp8}{\texttt{pzC3e8fjAp8}} \\
  & \hyperref[sec:yk5Oh8RBHLY]{2023-03-30 周四晚查經班《被擄歸回歷史研讀(9)》 講員:曾裔貴牧師} & 2023-04-01 & \href{https://youtube.com/watch?v=yk5Oh8RBHLY}{\texttt{yk5Oh8RBHLY}} \\
馬太福音 21:1-11 & \hyperref[sec:x06IlO4FHos]{2023-03-26 午堂《看你的王》馬太福音21\_1-11 講員:李思敬牧師} & 2023-04-03 & \href{https://youtube.com/watch?v=x06IlO4FHos}{\texttt{x06IlO4FHos}} \\
約翰一書 5:20-21 & \hyperref[sec:ythaF42qv_4]{2023-04-02 午堂《信仰求真》約翰一書5\_20-21 講員:曾裔貴牧師} & 2023-04-04 & \href{https://youtube.com/watch?v=ythaF42qv_4}{\texttt{ythaF42qv\_4}} \\
提摩太前書 6:6-10 & \hyperref[sec:_BOzzljBDkA]{2023-03-25 晚堂《知足人生》提摩太前書6\_6-10 講員:盧思源姑娘} & 2023-04-05 & \href{https://youtube.com/watch?v=-BOzzljBDkA}{\texttt{-BOzzljBDkA}} \\
使徒行傳 9:26-30-11:22-26 & \hyperref[sec:lYwkS21XVXg]{2023-04-02 加恩樓《生命師傅的典範》使徒行傳9\_26-30,11\_22-26 講員:陳一華牧師} & 2023-04-05 & \href{https://youtube.com/watch?v=lYwkS21XVXg}{\texttt{lYwkS21XVXg}} \\
  & \hyperref[sec:Am9T43vQMKQ]{2023-04-06 周四晚查經班《被擄歸回歷史研讀(10)》 講員:曾裔貴牧師} & 2023-04-07 & \href{https://youtube.com/watch?v=Am9T43vQMKQ}{\texttt{Am9T43vQMKQ}} \\
  & \hyperref[sec:Pndq2oLIHnI]{2023-04-13 周四晚查經班《被擄歸回歷史研讀(11)》 講員:曾裔貴牧師} & 2023-04-15 & \href{https://youtube.com/watch?v=Pndq2oLIHnI}{\texttt{Pndq2oLIHnI}} \\
路加福音 10:38-42 & \hyperref[sec:7t8JFmETTzs]{2023-04-16 加恩樓《世界不能奪去的福份》路加福音10\_38-42 講員:蔣文忠傳道} & 2023-04-18 & \href{https://youtube.com/watch?v=7t8JFmETTzs}{\texttt{7t8JFmETTzs}} \\
 5:1-14 & \hyperref[sec:cuPMYAewAe8]{2023-04-16 午堂《敬拜那又能又配的羔羊》啟示錄5\_1-14 講員:陳麗蟬牧師} & 2023-04-18 & \href{https://youtube.com/watch?v=cuPMYAewAe8}{\texttt{cuPMYAewAe8}} \\
 5:1-14 & \hyperref[sec:Yd_an_Ish1s]{2023-04-16 早堂《異象中的羔羊是世人惟一的希望》啟示錄5\_1-14 講員:謝靄薇牧師} & 2023-04-19 & \href{https://youtube.com/watch?v=Yd-an_Ish1s}{\texttt{Yd-an\_Ish1s}} \\
  & \hyperref[sec:SLUSAYk3lm4]{2023-04-20 周四晚查經班《被擄歸回歷史研讀(12)》 講員:曾裔貴牧師} & 2023-04-21 & \href{https://youtube.com/watch?v=SLUSAYk3lm4}{\texttt{SLUSAYk3lm4}} \\
哥林多前書 15:12-34 & \hyperref[sec:qxsLxZod7kk]{2023-04-22 晚堂《基督 ─ 我們的初熟果子》哥林多前書 15\_12-34 講員:黃頌言姑娘} & 2023-04-25 & \href{https://youtube.com/watch?v=qxsLxZod7kk}{\texttt{qxsLxZod7kk}} \\
希伯來書 12:1-3 & \hyperref[sec:rnbJeDm_hDc]{2023-04-23 加恩樓《奔跑吧!》希伯來書 12\_1-3 講員:周弘毅牧師} & 2023-04-25 & \href{https://youtube.com/watch?v=rnbJeDm_hDc}{\texttt{rnbJeDm\_hDc}} \\
路加福音 16:1-13 & \hyperref[sec:l5hkKLKHhA0]{2023-04-23 午堂《不義與忠心》路加福音 16\_1-13 講員:楊慶球牧師} & 2023-04-25 & \href{https://youtube.com/watch?v=l5hkKLKHhA0}{\texttt{l5hkKLKHhA0}} \\
創世記 22:1-14 & \hyperref[sec:J_KPU0uZysw]{2023-04-29 晚堂《我在這裡》創世記 22\_1-14 講員:褚永華牧師} & 2023-05-02 & \href{https://youtube.com/watch?v=J_KPU0uZysw}{\texttt{J\_KPU0uZysw}} \\
加拉太書 6:12-16 & \hyperref[sec:bNffFxPN8ik]{2023-04-30 加恩樓《我們所誇的是什麼?》加拉太書 6\_12-16 講員:鄧泰康傳道} & 2023-05-02 & \href{https://youtube.com/watch?v=bNffFxPN8ik}{\texttt{bNffFxPN8ik}} \\
加拉太書 3:21-29 & \hyperref[sec:sHtcf0yp8MI]{2023-04-30 午堂《律法之上》加拉太書 3\_21-29 講員:陳韋安牧師} & 2023-05-02 & \href{https://youtube.com/watch?v=sHtcf0yp8MI}{\texttt{sHtcf0yp8MI}} \\
加拉太書 6:11-16 & \hyperref[sec:FRK_krDh8hE]{2023-04-30 早堂《怎麼誇口?》加拉太書 6\_11-16 講員:楊明麗姑娘} & 2023-05-02 & \href{https://youtube.com/watch?v=FRK_krDh8hE}{\texttt{FRK\_krDh8hE}} \\
  & \hyperref[sec:AGfClw0jrYY]{2023-05-04 周四晚查經班《被擄歸回歷史研讀(13)》 講員:曾裔貴牧師} & 2023-05-06 & \href{https://youtube.com/watch?v=AGfClw0jrYY}{\texttt{AGfClw0jrYY}} \\
約翰一書 3:11-18 & \hyperref[sec:eLIEUCC43fA]{2023-05-07 加恩樓崇拜《共融》約翰一書 3\_11-18 講員:吳景霖傳道} & 2023-05-09 & \href{https://youtube.com/watch?v=eLIEUCC43fA}{\texttt{eLIEUCC43fA}} \\
  & \hyperref[sec:_3DxJVmmiXg]{2023-05-11 周四晚查經班《被擄歸回歷史研讀(14)》 講員:曾裔貴牧師} & 2023-05-12 & \href{https://youtube.com/watch?v=-3DxJVmmiXg}{\texttt{-3DxJVmmiXg}} \\
申命記 6:1-9 & \hyperref[sec:A6GWJTLfDEo]{2023-05-14 加恩樓《一家之主》申命記 6\_1-9 講員:謝靄薇牧師} & 2023-05-16 & \href{https://youtube.com/watch?v=A6GWJTLfDEo}{\texttt{A6GWJTLfDEo}} \\
出埃及記 2:1-10 & \hyperref[sec:noei9nAUNuE]{2023-05-14 早堂《我們這一家》出埃及記 2\_1-10 講員:黃頌言姑娘} & 2023-05-16 & \href{https://youtube.com/watch?v=noei9nAUNuE}{\texttt{noei9nAUNuE}} \\
腓立比書 4:4 & \hyperref[sec:PwPuQ_V9PRU]{2023-05-13 晚堂《喜樂再思》腓立比書 4\_4 講員:陳一華牧師} & 2023-05-17 & \href{https://youtube.com/watch?v=PwPuQ-V9PRU}{\texttt{PwPuQ-V9PRU}} \\
使徒行傳 8:1-8 & \hyperref[sec:_W_gd9p_W_Y]{2023-05-14 午堂《福音帶來的生命改變》使徒行傳 8\_1-8 講員:曾裔貴牧師} & 2023-05-17 & \href{https://youtube.com/watch?v=-W_gd9p_W_Y}{\texttt{-W\_gd9p\_W\_Y}} \\
  & \hyperref[sec:MtaDL5MWzm0]{2023-05-18 周四晚查經班《被擄歸回歷史研讀(15)》 講員:曾裔貴牧師} & 2023-05-19 & \href{https://youtube.com/watch?v=MtaDL5MWzm0}{\texttt{MtaDL5MWzm0}} \\
腓立比書 3:7-16 & \hyperref[sec:We85QzJyMNQ]{2023-05-20 晚堂《丟棄與得著》腓立比書 3\_7-16 講員:盧思源姑娘} & 2023-05-23 & \href{https://youtube.com/watch?v=We85QzJyMNQ}{\texttt{We85QzJyMNQ}} \\
腓立比書 3:7-16 & \hyperref[sec:9tgCzMr4KBk]{2023-05-21 午堂《你得著珍寶了嗎?》腓立比書 3\_7-16 講員:鄧泰康傳道} & 2023-05-23 & \href{https://youtube.com/watch?v=9tgCzMr4KBk}{\texttt{9tgCzMr4KBk}} \\
腓立比書 3:13-15 & \hyperref[sec:vWqPqvo2gRk]{2023-05-21 早堂《與基督同跑》腓立比書 3\_13-15 講員:楊明麗姑娘} & 2023-05-23 & \href{https://youtube.com/watch?v=vWqPqvo2gRk}{\texttt{vWqPqvo2gRk}} \\
路加福音 1:5-25-57-66 & \hyperref[sec:wCpZ46fk4yY]{2023-05-21 加恩樓《不相稱的禱告》路加福音 1\_5-25,57-66 講員:陳韋安牧師} & 2023-05-24 & \href{https://youtube.com/watch?v=wCpZ46fk4yY}{\texttt{wCpZ46fk4yY}} \\
羅馬書 15:22-29 & \hyperref[sec:aPVZvLNc78I]{2023-05-27 晚堂《五餅二魚的團契》羅馬書 15\_22-29 講員:陳麗蟬牧師} & 2023-05-30 & \href{https://youtube.com/watch?v=aPVZvLNc78I}{\texttt{aPVZvLNc78I}} \\
使徒行傳 2:1-21 & \hyperref[sec:hSecMeCWH4s]{2023-05-28 加恩樓《聖靈充滿的教會》使徒行傳 2\_1-21 講員:曾浩賢院長} & 2023-05-30 & \href{https://youtube.com/watch?v=hSecMeCWH4s}{\texttt{hSecMeCWH4s}} \\
傳道書 11:1-6 & \hyperref[sec:H85tdsHmqKI]{2023-05-28 午堂《不知道中的信仰》傳道書 11\_1-6 講員:蔣文忠傳道} & 2023-05-30 & \href{https://youtube.com/watch?v=H85tdsHmqKI}{\texttt{H85tdsHmqKI}} \\
  & \hyperref[sec:C1vy6YO3h0g]{2023-06-01 周四晚查經班《被擄歸回歷史研讀(16)》 講員:曾裔貴牧師} & 2023-06-03 & \href{https://youtube.com/watch?v=C1vy6YO3h0g}{\texttt{C1vy6YO3h0g}} \\
馬太福音 28:11-20 & \hyperref[sec:WLN153QoHOE]{2023-06-04 加恩樓《後真相年代下的信仰反思》馬太福音 28\_11-20 講員:吳景霖傳道} & 2023-06-06 & \href{https://youtube.com/watch?v=WLN153QoHOE}{\texttt{WLN153QoHOE}} \\
路加福音 5:27-32 & \hyperref[sec:jJiE2EQ4f48]{2023-06-04 午堂《悔改的筵席》路加福音 5\_27-32 講員:楊明麗姑娘} & 2023-06-06 & \href{https://youtube.com/watch?v=jJiE2EQ4f48}{\texttt{jJiE2EQ4f48}} \\
路加福音 5:27-39 & \hyperref[sec:KzM_XrVrUvg]{2023-06-04 早堂《悔改的筵席》路加福音 5\_27-39 講員:曾裔貴牧師} & 2023-06-07 & \href{https://youtube.com/watch?v=KzM_XrVrUvg}{\texttt{KzM\_XrVrUvg}} \\
  & \hyperref[sec:qmt5oGwLkF4]{2023-06-08 周四晚查經班《被擄歸回歷史研讀(17)》 講員:曾裔貴牧師} & 2023-06-09 & \href{https://youtube.com/watch?v=qmt5oGwLkF4}{\texttt{qmt5oGwLkF4}} \\
路加福音 7:36-50 & \hyperref[sec:Fl2IkOLpxSI]{2023-06-10 晚堂《罪人的筵席》路加福音 7\_36-50 講員:吳景霖傳道} & 2023-06-13 & \href{https://youtube.com/watch?v=Fl2IkOLpxSI}{\texttt{Fl2IkOLpxSI}} \\
馬太福音 9:35-38 & \hyperref[sec:A_K04hyrSP8]{2023-06-11 加恩樓《建立無牆的教會》馬太福音 9\_35-38 講員:鄧泰康傳道} & 2023-06-13 & \href{https://youtube.com/watch?v=A-K04hyrSP8}{\texttt{A-K04hyrSP8}} \\
路加福音 7:36-50 & \hyperref[sec:lzufA9tLqUc]{2023-06-11 午堂《愛和謙卑的事奉》路加福音 7\_36-50 講員:陳麗蟬牧師} & 2023-06-13 & \href{https://youtube.com/watch?v=lzufA9tLqUc}{\texttt{lzufA9tLqUc}} \\
提摩太後書 3:16-17 & \hyperref[sec:eaFKBNr6R3g]{2023-06-18 加恩樓《有故事的人》提摩太後書 3\_16-17 講員:吳景霖傳道} & 2023-06-20 & \href{https://youtube.com/watch?v=eaFKBNr6R3g}{\texttt{eaFKBNr6R3g}} \\
提摩太後書 3:16-17 & \hyperref[sec:cNC4ab19lyE]{2023-06-18 早堂《說故事的爸媽》提摩太後書 3\_16-17 講員:吳景霖傳道} & 2023-06-20 & \href{https://youtube.com/watch?v=cNC4ab19lyE}{\texttt{cNC4ab19lyE}} \\
傳道書 11:1-6 & \hyperref[sec:I3SP1geCVE8]{2023-06-17 晚堂《不知道中的信仰》傳道書 11\_1-6 講員:蔣文忠傳道} & 2023-06-21 & \href{https://youtube.com/watch?v=I3SP1geCVE8}{\texttt{I3SP1geCVE8}} \\
路加福音 12:1-24 & \hyperref[sec:9x4ci5x2GZ0]{2023-06-18 午堂《基督徒的日常》路加福音 12\_1-24 講員:曾裔貴牧師} & 2023-06-21 & \href{https://youtube.com/watch?v=9x4ci5x2GZ0}{\texttt{9x4ci5x2GZ0}} \\
路加福音 12:35-48 & \hyperref[sec:keP18GcFGPs]{2023-06-24 晚堂《儆醒僕人筵席》路加福音 12\_35-48 講員:黃頌言姑娘} & 2023-06-27 & \href{https://youtube.com/watch?v=keP18GcFGPs}{\texttt{keP18GcFGPs}} \\
列王記上 18:1-6 & \hyperref[sec:i86dI_9QW94]{2023-06-25 加恩樓《掌管萬有的神》列王記上 18\_1-6 講員:陳一華牧師} & 2023-06-27 & \href{https://youtube.com/watch?v=i86dI_9QW94}{\texttt{i86dI\_9QW94}} \\
路加福音 19:12-23 & \hyperref[sec:tu4SyHPIS_U]{2023-06-25 早堂《未來的風險》路加福音 19\_12-23 講員:陳韋安牧師} & 2023-06-27 & \href{https://youtube.com/watch?v=tu4SyHPIS_U}{\texttt{tu4SyHPIS\_U}} \\
羅馬書 6:1-11 & \hyperref[sec:WZaLHtVToyI]{2023-06-25 午堂《克勝罪惡》羅馬書 6\_1-11 講員:褚永華牧師} & 2023-06-28 & \href{https://youtube.com/watch?v=WZaLHtVToyI}{\texttt{WZaLHtVToyI}} \\
  & \hyperref[sec:G4jQV76in2M]{2023-06-29 周四晚查經班《被擄歸回歷史研讀(18)》 講員:曾裔貴牧師} & 2023-06-30 & \href{https://youtube.com/watch?v=G4jQV76in2M}{\texttt{G4jQV76in2M}} \\
列王記上 19:1-18 & \hyperref[sec:W6aFxBk_o0s]{2023-07-02 加恩樓《當以為走到絕境時》列王記上 19\_1-18 講員:鄧泰康傳道} & 2023-07-05 & \href{https://youtube.com/watch?v=W6aFxBk_o0s}{\texttt{W6aFxBk\_o0s}} \\
路加福音 14:1-24 & \hyperref[sec:Vwm387KuCtc]{2023-07-02 早堂《筵席的再思》路加福音 14\_1-24 講員:曾裔貴牧師} & 2023-07-05 & \href{https://youtube.com/watch?v=Vwm387KuCtc}{\texttt{Vwm387KuCtc}} \\
  & \hyperref[sec:LQfxIAo88QM]{2023-07-06 周四晚查經班《被擄歸回歷史研讀(19)》 講員:曾裔貴牧師} & 2023-07-08 & \href{https://youtube.com/watch?v=LQfxIAo88QM}{\texttt{LQfxIAo88QM}} \\
路加福音 19:1-10 & \hyperref[sec:oxDrdn_rwTg]{2023-07-08 晚堂《奇妙的改變》路加福音 19\_1-10 講員:陳麗蟬牧師} & 2023-07-11 & \href{https://youtube.com/watch?v=oxDrdn-rwTg}{\texttt{oxDrdn-rwTg}} \\
提摩太前書詩篇羅馬書 1:19-20 & \hyperref[sec:87awTrKJFrY]{2023-07-09 加恩樓《平衡空間的人生》羅馬書 1\_19-20,提摩太前書 2\_5-6,詩篇 27\_13-14 講員:李盛林牧師} & 2023-07-11 & \href{https://youtube.com/watch?v=87awTrKJFrY}{\texttt{87awTrKJFrY}} \\
路加福音 19:1-10 & \hyperref[sec:QuJcmG3w_7k]{2023-07-09 午堂《愛在生命改變前》路加福音 19\_1-10 講員:楊明麗姑娘} & 2023-07-11 & \href{https://youtube.com/watch?v=QuJcmG3w-7k}{\texttt{QuJcmG3w-7k}} \\
尼希米記 11:1-3 & \hyperref[sec:IDYmD1EXhB4]{2023-07-16 加恩樓《別叫我英雄》尼希米記 11\_1-3 講員:曾裔貴牧師} & 2023-07-22 & \href{https://youtube.com/watch?v=IDYmD1EXhB4}{\texttt{IDYmD1EXhB4}} \\
使徒行傳 1:8-28:31 & \hyperref[sec:vtGjR9atL3U]{2023-07-16 午堂《放膽傳道,無人能禁止》使徒行傳 1\_8,28\_31 講員:郭鴻標牧師} & 2023-07-23 & \href{https://youtube.com/watch?v=vtGjR9atL3U}{\texttt{vtGjR9atL3U}} \\
路加福音 10:1-12 & \hyperref[sec:PIfMPVh5oUA]{2023-07-16 早堂《耶穌訓練門徒》路加福音 10\_1-12 講員:曾裔貴牧師} & 2023-07-23 & \href{https://youtube.com/watch?v=PIfMPVh5oUA}{\texttt{PIfMPVh5oUA}} \\
  & \hyperref[sec:zajnPwak_dU]{2023-07-13 周四晚查經班《被擄歸回歷史研讀(20)》 講員:曾裔貴牧師} & 2023-07-25 & \href{https://youtube.com/watch?v=zajnPwak-dU}{\texttt{zajnPwak-dU}} \\
歌羅西書 2:1-19 & \hyperref[sec:epe4ZHzeWUY]{2023-07-22 晚堂《打穩根基》歌羅西書 2\_1-19 講員:謝靄薇牧師} & 2023-07-25 & \href{https://youtube.com/watch?v=epe4ZHzeWUY}{\texttt{epe4ZHzeWUY}} \\
約翰福音 21:15-22 & \hyperref[sec:dFeC0BqGntE]{2023-07-23 加恩樓《跟隨我》約翰福音 21\_15-22 講員:陳靜娟姑娘} & 2023-07-25 & \href{https://youtube.com/watch?v=dFeC0BqGntE}{\texttt{dFeC0BqGntE}} \\
約翰福音 17:6-19 & \hyperref[sec:_FPkb_FFsjA]{2023-07-23 早堂《寶貴的禱告》約翰福音 17\_6-19 講員:曾敏姑娘} & 2023-07-25 & \href{https://youtube.com/watch?v=_FPkb_FFsjA}{\texttt{\_FPkb\_FFsjA}} \\
以斯帖記  & \hyperref[sec:iuYWy9ExSgQ]{2023-07-27 周四晚查經班《被擄歸回歷史研讀(21):以斯帖記》 講員:曾裔貴牧師} & 2023-07-29 & \href{https://youtube.com/watch?v=iuYWy9ExSgQ}{\texttt{iuYWy9ExSgQ}} \\
以弗所書 4:11-16 & \hyperref[sec:c_zS1WaX1q0]{2023-07-29 晚堂《建立基督的身體》以弗所書 4\_11-16 講員:盧思源姑娘} & 2023-08-08 & \href{https://youtube.com/watch?v=c-zS1WaX1q0}{\texttt{c-zS1WaX1q0}} \\
以弗所書 4:11-16 & \hyperref[sec:qI_KwviwgzU]{2023-07-30 加恩樓《建立基督的身體》以弗所書 4\_11-16 講員:楊明麗姑娘} & 2023-08-08 & \href{https://youtube.com/watch?v=qI_KwviwgzU}{\texttt{qI\_KwviwgzU}} \\
馬太福音 18:1-9 & \hyperref[sec:i3m_mXZnkC4]{2023-08-06 加恩樓《天國的小孩》馬太福音 18\_1-9 講員:鄧泰康傳道} & 2023-08-08 & \href{https://youtube.com/watch?v=i3m-mXZnkC4}{\texttt{i3m-mXZnkC4}} \\
歌羅西書 4:5-6 & \hyperref[sec:1pEPdzdQ8_4]{2023-08-06 午堂《著地氣的傳福音》歌羅西書 4\_5-6 講員:曾裔貴牧師} & 2023-08-08 & \href{https://youtube.com/watch?v=1pEPdzdQ8_4}{\texttt{1pEPdzdQ8\_4}} \\
希伯來書 3:13-14 & \hyperref[sec:FlUdBHop01k]{2023-07-30 午堂《信仰的初心》希伯來書 3\_13-14 講員:陳韋安牧師} & 2023-08-09 & \href{https://youtube.com/watch?v=FlUdBHop01k}{\texttt{FlUdBHop01k}} \\
使徒行傳 17:16-31 & \hyperref[sec:RptdKZSozLI]{2023-08-12 晚堂《進入人群的福音》使徒行傳 17\_16-31 講員:陳麗蟬牧師} & 2023-08-15 & \href{https://youtube.com/watch?v=RptdKZSozLI}{\texttt{RptdKZSozLI}} \\
約翰一書彼得前書 1:1-2 & \hyperref[sec:kFsvDzgp4nI]{2023-08-13 加恩樓《使人蒙福的「道」》約翰一書 1\_1-2,彼得前書 1\_23-25 講員:樓恩德牧師} & 2023-08-15 & \href{https://youtube.com/watch?v=kFsvDzgp4nI}{\texttt{kFsvDzgp4nI}} \\
哥林多後書 4:1-18 & \hyperref[sec:FIbxv36OMMk]{2023-08-20 午堂《傳道者的心》哥林多後書 4\_1-18 講員:鄧泰康傳道} & 2023-08-22 & \href{https://youtube.com/watch?v=FIbxv36OMMk}{\texttt{FIbxv36OMMk}} \\
以賽亞書 58:3-7 & \hyperref[sec:cMmv_XzlkAc]{2023-08-20 早堂《禁食的意義》以賽亞書 58\_3-7 講員:陳韋安牧師} & 2023-08-22 & \href{https://youtube.com/watch?v=cMmv_XzlkAc}{\texttt{cMmv\_XzlkAc}} \\
路加福音 10:25-37 & \hyperref[sec:TEOKFaxZL4g]{2023-08-20 加恩樓《愛鄰如己．照樣做吧》路加福音 10\_25-37 講員:陳麗蟬牧師} & 2023-08-23 & \href{https://youtube.com/watch?v=TEOKFaxZL4g}{\texttt{TEOKFaxZL4g}} \\
約伯記 28:1-28 & \hyperref[sec:zrsNFCgx1bQ]{2022-05-22 午堂《明智藏哪裡》約伯記 28\_1-28 講員:周曉暉牧師} & 2023-08-25 & \href{https://youtube.com/watch?v=zrsNFCgx1bQ}{\texttt{zrsNFCgx1bQ}} \\
出埃及記 23:1-13 & \hyperref[sec:lKgHJObwVL4]{2022-12-18 加恩樓《爭議中持守公正與良善》出埃及記 23\_1-13 講員:周曉暉牧師} & 2023-08-25 & \href{https://youtube.com/watch?v=lKgHJObwVL4}{\texttt{lKgHJObwVL4}} \\
以斯帖記  & \hyperref[sec:iRhu93YiYcU]{2023-08-24 周四晚查經班《被擄歸回歷史研讀(22)以斯帖記》 講員:曾裔貴牧師} & 2023-08-26 & \href{https://youtube.com/watch?v=iRhu93YiYcU}{\texttt{iRhu93YiYcU}} \\
彼得前書 5:1-11 & \hyperref[sec:fSxNb3R9eIk]{2023-08-27 加恩樓《父慈子孝的神家》彼得前書 5\_1-11 講員:鄧泰康傳道} & 2023-08-29 & \href{https://youtube.com/watch?v=fSxNb3R9eIk}{\texttt{fSxNb3R9eIk}} \\
哥林多前書 16:15-16 & \hyperref[sec:xZ5COJOEvOY]{2023-08-27 午堂《為主同工,同勞的人》哥林多前書 16\_15-16 講員:曾裔貴牧師} & 2023-08-29 & \href{https://youtube.com/watch?v=xZ5COJOEvOY}{\texttt{xZ5COJOEvOY}} \\
彼得前書 5:1-5 & \hyperref[sec:fjBMO0wHlN4]{2023-08-27 早堂《彼此順服的教會》彼得前書 5\_1-5 講員:盧思源姑娘} & 2023-08-29 & \href{https://youtube.com/watch?v=fjBMO0wHlN4}{\texttt{fjBMO0wHlN4}} \\
以斯帖記  & \hyperref[sec:9XvZPOyWZdE]{2023-08-31 周四晚查經班《被擄歸回歷史研讀(23)以斯帖記》 講員:曾裔貴牧師} & 2023-09-05 & \href{https://youtube.com/watch?v=9XvZPOyWZdE}{\texttt{9XvZPOyWZdE}} \\
提摩太後書 4:1-8 & \hyperref[sec:ay_HsjJTAHo]{2023-09-03 加恩樓《保羅最後的勸勉》提摩太後書 4\_1-8 講員:鄧泰康傳道} & 2023-09-05 & \href{https://youtube.com/watch?v=ay-HsjJTAHo}{\texttt{ay-HsjJTAHo}} \\
彼得後書 1:3-8 & \hyperref[sec:Bres3qEZ6Yk]{2023-09-03 午堂《在信仰上加倍努力》彼得後書 1\_3-8 講員:曾裔貴牧師} & 2023-09-05 & \href{https://youtube.com/watch?v=Bres3qEZ6Yk}{\texttt{Bres3qEZ6Yk}} \\
撒母耳記下 23:13-19 & \hyperref[sec:naPXdHer1J8]{2023-09-10 加恩樓《大衛與三個沒有名字的勇士》撒母耳記下 23\_13-19 講員:陳韋安牧師} & 2023-09-12 & \href{https://youtube.com/watch?v=naPXdHer1J8}{\texttt{naPXdHer1J8}} \\
提摩太後書 4:1-8 & \hyperref[sec:TXT2__xkqM8]{2023-09-10 早堂《保羅最後的勸勉》提摩太後書 4\_1-8 講員:鄧泰康傳道} & 2023-09-12 & \href{https://youtube.com/watch?v=TXT2_-xkqM8}{\texttt{TXT2\_-xkqM8}} \\
詩篇 8:1-9 & \hyperref[sec:LQDQ7CliZ7U]{2023-09-09 晚堂《神你是誰?而我又算甚麼?》詩篇 8\_1-9 講員:張美薇牧師} & 2023-09-13 & \href{https://youtube.com/watch?v=LQDQ7CliZ7U}{\texttt{LQDQ7CliZ7U}} \\
提摩太後書路加福音 9:57-62 & \hyperref[sec:eRch8741F34]{2023-09-10 午堂《若要跟從主》路加福音 9\_57-62,提摩太後書 1\_1-12 講員:郭熹瑜宣教士} & 2023-09-13 & \href{https://youtube.com/watch?v=eRch8741F34}{\texttt{eRch8741F34}} \\
撒迦利亞書  & \hyperref[sec:Kpbv61GAjrM]{2023-09-14 周四晚查經班《歸回先知書研讀(1)撒迦利亞書》 講員:曾裔貴牧師} & 2023-09-16 & \href{https://youtube.com/watch?v=Kpbv61GAjrM}{\texttt{Kpbv61GAjrM}} \\
詩篇 128:1-6 & \hyperref[sec:BaO8MW3rJk0]{2023-09-16 晚堂《福氣:敬畏主+遵主路》詩篇 128\_1-6 講員:陳麗蟬牧師} & 2023-09-19 & \href{https://youtube.com/watch?v=BaO8MW3rJk0}{\texttt{BaO8MW3rJk0}} \\
約翰福音 14:1-14 & \hyperref[sec:Ah5GT26Za2E]{2023-09-17 早堂《你使我堅壯》約翰福音 14\_1-14 講員:蕭建輝牧師} & 2023-09-19 & \href{https://youtube.com/watch?v=Ah5GT26Za2E}{\texttt{Ah5GT26Za2E}} \\
路加福音 11:5-13 & \hyperref[sec:Hdjnt297x1A]{2023-09-17 加恩樓《好吧!天父!我信靠你!》路加福音 11\_5-13 講員:蕭建輝牧師} & 2023-09-20 & \href{https://youtube.com/watch?v=Hdjnt297x1A}{\texttt{Hdjnt297x1A}} \\
以賽亞書 42:1-9 & \hyperref[sec:tDVwCjoCqSI]{2023-09-17 午堂《以賽亞先知的生命質素與使命》以賽亞書 42\_1-9 講員:郭鴻標牧師} & 2023-09-20 & \href{https://youtube.com/watch?v=tDVwCjoCqSI}{\texttt{tDVwCjoCqSI}} \\
撒迦利亞書  & \hyperref[sec:_MeYYEhxYZE]{2023-09-21 周四晚查經班《歸回先知書研讀(2)撒迦利亞書》 講員:曾裔貴牧師} & 2023-09-23 & \href{https://youtube.com/watch?v=-MeYYEhxYZE}{\texttt{-MeYYEhxYZE}} \\
使徒行傳 21:17-26 & \hyperref[sec:B_Q_tzD4YvY]{2023-09-24 加恩樓《入會四則》使徒行傳 21\_17-26 講員:陳麗蟬牧師} & 2023-09-26 & \href{https://youtube.com/watch?v=B-Q-tzD4YvY}{\texttt{B-Q-tzD4YvY}} \\
哥林多後書 5:17-18 & \hyperref[sec:CTfZJ96vxzI]{2023-09-24 早堂《重生》哥林多後書 5\_17-18 講員:黃頌言姑娘} & 2023-09-26 & \href{https://youtube.com/watch?v=CTfZJ96vxzI}{\texttt{CTfZJ96vxzI}} \\
希伯來書 4:1-11 & \hyperref[sec:QkQu6jc1ZPg]{2023-09-24 午堂《兩個安息》希伯來書 4\_1-11 講員:楊慶球牧師} & 2023-09-27 & \href{https://youtube.com/watch?v=QkQu6jc1ZPg}{\texttt{QkQu6jc1ZPg}} \\
撒迦利亞書  & \hyperref[sec:t6pbod4Qek8]{2023-09-28 周四晚查經班《歸回先知書研讀(3)撒迦利亞書》 講員:曾裔貴牧師} & 2023-09-29 & \href{https://youtube.com/watch?v=t6pbod4Qek8}{\texttt{t6pbod4Qek8}} \\
彌迦書 6:6-8 & \hyperref[sec:ZbYxDMl2Zbg]{2023-10-01 加恩樓《選擇站在上帝那邊》彌迦書 6\_6-8 講員:鄧泰康傳道} & 2023-10-03 & \href{https://youtube.com/watch?v=ZbYxDMl2Zbg}{\texttt{ZbYxDMl2Zbg}} \\
以弗所書 4:11-16 & \hyperref[sec:EvC3MODORZI]{2023-10-01 早堂《天降恩賜齊共享》以弗所書 4\_11-16 講員:陳麗蟬牧師} & 2023-10-03 & \href{https://youtube.com/watch?v=EvC3MODORZI}{\texttt{EvC3MODORZI}} \\
馬太福音 9:35-10:4 & \hyperref[sec:x5xyTkGVbQ0]{2023-09-30 晚堂《流散的羊群》馬太福音 9\_35-10\_4 講員:褚永華牧師} & 2023-10-04 & \href{https://youtube.com/watch?v=x5xyTkGVbQ0}{\texttt{x5xyTkGVbQ0}} \\
撒母耳記下 23:13-19 & \hyperref[sec:JHYqHgNM3GQ]{2023-10-01 午堂《大衛與三個沒有名字的勇士》撒母耳記下 23\_13-19 講員:陳韋安牧師} & 2023-10-04 & \href{https://youtube.com/watch?v=JHYqHgNM3GQ}{\texttt{JHYqHgNM3GQ}} \\
撒迦利亞書  & \hyperref[sec:CXJyuYm4_pE]{2023-10-05 周四晚查經班《歸回先知書研讀(4)撒迦利亞書》 講員:曾裔貴牧師} & 2023-10-06 & \href{https://youtube.com/watch?v=CXJyuYm4-pE}{\texttt{CXJyuYm4-pE}} \\
馬可福音 1:40-44 & \hyperref[sec:G0FSq209Kac]{2023-10-07 晚堂《醫治為要得潔淨》馬可福音 1\_40-44 講員:黃國維院長} & 2023-10-11 & \href{https://youtube.com/watch?v=G0FSq209Kac}{\texttt{G0FSq209Kac}} \\
約翰福音 4:1-42 & \hyperref[sec:iO1qZX8Cmq4]{2023-10-08 早堂《生命轉化之旅》約翰福音 4\_1-42 講員:陳淑娟牧師} & 2023-10-12 & \href{https://youtube.com/watch?v=iO1qZX8Cmq4}{\texttt{iO1qZX8Cmq4}} \\
撒迦利亞書  & \hyperref[sec:wwGm5m9DNT8]{2023-10-12 周四晚查經班《歸回先知書研讀(5)撒迦利亞書》 講員:曾裔貴牧師} & 2023-10-14 & \href{https://youtube.com/watch?v=wwGm5m9DNT8}{\texttt{wwGm5m9DNT8}} \\
馬可福音 5:21-43 & \hyperref[sec:HWMu4dld7_M]{2023-10-15 加恩樓《這是一個怎樣的信心》馬可福音 5\_21-43 講員:老耀雄牧師} & 2023-10-18 & \href{https://youtube.com/watch?v=HWMu4dld7_M}{\texttt{HWMu4dld7\_M}} \\
以弗所書 1:17-23-4:7-16 & \hyperref[sec:cM9RUvSdlG8]{2023-10-15 早堂《尊基督為元首的門徒相交》以弗所書 1\_17-23,4\_7-16 講員:曾裔貴牧師} & 2023-10-19 & \href{https://youtube.com/watch?v=cM9RUvSdlG8}{\texttt{cM9RUvSdlG8}} \\
撒迦利亞書  & \hyperref[sec:c_6waYrNYvI]{2023-10-19 周四晚查經班《歸回先知書研讀(6)撒迦利亞書》 講員:曾裔貴牧師} & 2023-10-21 & \href{https://youtube.com/watch?v=c_6waYrNYvI}{\texttt{c\_6waYrNYvI}} \\
馬太福音 10:34-42 & \hyperref[sec:qZ9s1IR331E]{2023-10-21 晚堂《作主門徒》馬太福音 10\_34-42 講員:陳麗蟬牧師} & 2023-10-24 & \href{https://youtube.com/watch?v=qZ9s1IR331E}{\texttt{qZ9s1IR331E}} \\
加拉太書 6:9-10 & \hyperref[sec:ZmEDOfaqivc]{2023-10-22 加恩樓《及時行善》加拉太書 6\_9-10 講員:何智雄牧師} & 2023-10-24 & \href{https://youtube.com/watch?v=ZmEDOfaqivc}{\texttt{ZmEDOfaqivc}} \\
約拿書 1:1-17-4:1-3 & \hyperref[sec:LM01sAVsGrw]{2023-10-22 午堂《經歷有恩典憐憫的神》約拿書 1\_1-17,4\_1-3 講員:陳燕雯姑娘} & 2023-10-25 & \href{https://youtube.com/watch?v=LM01sAVsGrw}{\texttt{LM01sAVsGrw}} \\
以弗所書馬可福音 4:14-16 & \hyperref[sec:xr_m3_qAqFs]{2023-10-28 晚堂《在愛中成長》以弗所書 4\_14-16,馬可福音 12\_29-30 講員:楊明麗姑娘} & 2023-10-31 & \href{https://youtube.com/watch?v=xr_m3-qAqFs}{\texttt{xr\_m3-qAqFs}} \\
腓利門書 1:4-20 & \hyperref[sec:ROGX9cAg8Cg]{2023-10-29  42周年堂慶聯合崇拜《友伴．前行》腓利門書 1\_4-20 講員:周曉暉牧師} & 2023-10-31 & \href{https://youtube.com/watch?v=ROGX9cAg8Cg}{\texttt{ROGX9cAg8Cg}} \\
撒迦利亞書  & \hyperref[sec:j9CuOv5EhDA]{2023-11-02 周四晚查經班《歸回先知書研讀(7)撒迦利亞書》 講員:曾裔貴牧師} & 2023-11-03 & \href{https://youtube.com/watch?v=j9CuOv5EhDA}{\texttt{j9CuOv5EhDA}} \\
詩篇 116:1-19 & \hyperref[sec:URxw02niNWY]{2023-11-04 晚堂《以感恩為祭》詩篇 116\_1-19 講員:曾裔貴牧師} & 2023-11-07 & \href{https://youtube.com/watch?v=URxw02niNWY}{\texttt{URxw02niNWY}} \\
哥林多後書使徒行傳 9:30-13:2-27:22-24 & \hyperref[sec:EpQdyi4IyoI]{2023-11-05 午堂《不同階段有不同學習:神永遠有奇妙恩典》使徒行傳 9\_30,13\_2,27\_22-24,哥林多後書 12\_9 講員:蔡少琪院長} & 2023-11-08 & \href{https://youtube.com/watch?v=EpQdyi4IyoI}{\texttt{EpQdyi4IyoI}} \\
撒迦利亞書  & \hyperref[sec:z09lg9RcA3o]{2023-11-09 周四晚查經班《歸回先知書研讀(8)撒迦利亞書》 講員:曾裔貴牧師} & 2023-11-11 & \href{https://youtube.com/watch?v=z09lg9RcA3o}{\texttt{z09lg9RcA3o}} \\
馬可福音 16:9-20 & \hyperref[sec:Jspoigh_Nw0]{2023-11-12 午堂《傳揚復活的基督》馬可福音 16\_9-20 講員:曾裔貴牧師} & 2023-11-14 & \href{https://youtube.com/watch?v=Jspoigh-Nw0}{\texttt{Jspoigh-Nw0}} \\
哥林多前書 9:24-27 & \hyperref[sec:lMH9_3_EQe8]{2023-11-12 加恩樓《為主奔跑》哥林多前書 9\_24-27 講員:陳一華牧師} & 2023-11-15 & \href{https://youtube.com/watch?v=lMH9_3-EQe8}{\texttt{lMH9\_3-EQe8}} \\
詩篇 150:1-6 & \hyperref[sec:zCZXBv4_hT8]{2023-11-18 晚堂《當來讚美神》詩篇 150\_1-6 講員:張美薇牧師} & 2023-11-21 & \href{https://youtube.com/watch?v=zCZXBv4_hT8}{\texttt{zCZXBv4\_hT8}} \\
約書亞記 14:6-15 & \hyperref[sec:uKxjJHU8_7c]{2023-11-19 加恩樓《生命的投資》約書亞記 14\_6-15 講員:蔣文忠傳道} & 2023-11-21 & \href{https://youtube.com/watch?v=uKxjJHU8-7c}{\texttt{uKxjJHU8-7c}} \\
以賽亞書 49:1-7 & \hyperref[sec:z_EDsyn25oM]{2023-11-19 午堂《成為萬邦之光》以賽亞書 49\_1-7 講員:陳麗蟬牧師} & 2023-11-22 & \href{https://youtube.com/watch?v=z-EDsyn25oM}{\texttt{z-EDsyn25oM}} \\
歌羅西書 3:15-17 & \hyperref[sec:Th6Drt1kZ1s]{2023-11-25 晚堂《心存感恩》歌羅西書 3\_15-17 講員:盧思源姑娘} & 2023-11-28 & \href{https://youtube.com/watch?v=Th6Drt1kZ1s}{\texttt{Th6Drt1kZ1s}} \\
提摩太後書 1:1-5 & \hyperref[sec:Q0yYJrC6Zoc]{2023-11-26 加恩樓《代代相傳的信仰》提摩太後書 1\_1-5 講員:黃頌言姑娘} & 2023-11-28 & \href{https://youtube.com/watch?v=Q0yYJrC6Zoc}{\texttt{Q0yYJrC6Zoc}} \\
使徒行傳 11:19-30 & \hyperref[sec:nQU0VbnDk9Y]{2023-11-26 午堂《要作榮耀主的教會》使徒行傳 11\_19-30 講員:趙家碩牧師} & 2023-11-28 & \href{https://youtube.com/watch?v=nQU0VbnDk9Y}{\texttt{nQU0VbnDk9Y}} \\
馬太福音 2:1-12 & \hyperref[sec:gy4cK2fsHdc]{2023-12-02 晚堂《聖誕禮物》馬太福音 2\_1-12 講員:楊明麗姑娘} & 2023-12-06 & \href{https://youtube.com/watch?v=gy4cK2fsHdc}{\texttt{gy4cK2fsHdc}} \\
使徒行傳 5:12-32 & \hyperref[sec:7XxklgVuie4]{2023-12-03 午堂《傳講生命之道》使徒行傳 5\_12-32 講員:曾裔貴牧師} & 2023-12-06 & \href{https://youtube.com/watch?v=7XxklgVuie4}{\texttt{7XxklgVuie4}} \\
撒迦利亞書  & \hyperref[sec:Qx_35UZ1A44]{2023-12-07 周四晚查經班《歸回先知書研讀(9)撒迦利亞書》 講員:曾裔貴牧師} & 2023-12-08 & \href{https://youtube.com/watch?v=Qx-35UZ1A44}{\texttt{Qx-35UZ1A44}} \\
路加福音 3:1-9 & \hyperref[sec:_n0nRPZl6sE]{2023-12-09 晚堂《君王來臨的平安喜樂》路加福音 3\_1-9 講員:曾裔貴牧師} & 2023-12-12 & \href{https://youtube.com/watch?v=-n0nRPZl6sE}{\texttt{-n0nRPZl6sE}} \\
馬可福音 1:1-8 & \hyperref[sec:KJG_idsYK_w]{2023-12-10 加恩樓《修直祂的道》馬可福音 1\_1-8 講員:楊明麗姑娘} & 2023-12-12 & \href{https://youtube.com/watch?v=KJG_idsYK-w}{\texttt{KJG\_idsYK-w}} \\
以賽亞書 11:1-9 & \hyperref[sec:m2hnqzcprEM]{2023-12-10 午堂《共融的新天新地》以賽亞書 11\_1-9 講員:陳麗蟬牧師} & 2023-12-12 & \href{https://youtube.com/watch?v=m2hnqzcprEM}{\texttt{m2hnqzcprEM}} \\
約書亞記 14:6-15 & \hyperref[sec:jzt3Ohpru70]{2023-12-16 晚堂《生命的投資》約書亞記 14\_6-15 講員:蔣文忠傳道} & 2023-12-19 & \href{https://youtube.com/watch?v=jzt3Ohpru70}{\texttt{jzt3Ohpru70}} \\
路加福音 2:25-35 & \hyperref[sec:oC6cD_CKYi8]{2023-12-17 加恩樓《必看見所立的基督》路加福音 2\_25-35 講員:張漢強牧師} & 2023-12-19 & \href{https://youtube.com/watch?v=oC6cD-CKYi8}{\texttt{oC6cD-CKYi8}} \\
約伯記 42:1-6 & \hyperref[sec:EJ51POSz_gM]{2023-12-17 午堂《從前風聞有你,現在親眼看見你》約伯記 42\_1-6 講員:鄧泰康傳道} & 2023-12-19 & \href{https://youtube.com/watch?v=EJ51POSz_gM}{\texttt{EJ51POSz\_gM}} \\
詩篇 90:1-17 & \hyperref[sec:JmKumENSOqs]{2023-12-17 早堂《赤子之心》詩篇 90\_1-17 講員:陳韋安牧師} & 2023-12-19 & \href{https://youtube.com/watch?v=JmKumENSOqs}{\texttt{JmKumENSOqs}} \\
撒迦利亞書  & \hyperref[sec:I2kkhTmf8Lk]{2023-12-21 周四晚查經班《歸回先知書研讀(10)撒迦利亞書》 講員:曾裔貴牧師} & 2023-12-23 & \href{https://youtube.com/watch?v=I2kkhTmf8Lk}{\texttt{I2kkhTmf8Lk}} \\
路加福音 2:8-20 & \hyperref[sec:pZkchuW18Ro]{2023-12-24 加恩樓三語聯合崇拜《Good News of Great Joy》路加福音 2\_8-20 講員:Rev. Carol Ilidan} & 2023-12-27 & \href{https://youtube.com/watch?v=pZkchuW18Ro}{\texttt{pZkchuW18Ro}} \\
  & \hyperref[sec:7Vdee0ksMqc]{2023-12-24 平安夜佈道會《不一樣的聖誕》 講員:李炳光牧師} & 2023-12-27 & \href{https://youtube.com/watch?v=7Vdee0ksMqc}{\texttt{7Vdee0ksMqc}} \\
以賽亞書 7:10-17 & \hyperref[sec:nEkocKnhlmk]{2023-12-24 聯合崇拜《以馬內利的再思》以賽亞書 7\_10-17 講員:盧思源姑娘} & 2023-12-27 & \href{https://youtube.com/watch?v=nEkocKnhlmk}{\texttt{nEkocKnhlmk}} \\
  & \hyperref[sec:bpii0IGPrEA]{2023-12-25 聖誕節崇拜《諸天萬物同歡唱》 講員:陳麗蟬牧師} & 2023-12-27 & \href{https://youtube.com/watch?v=bpii0IGPrEA}{\texttt{bpii0IGPrEA}} \\
撒迦利亞書  & \hyperref[sec:MkLoNpPCJk4]{2023-12-28 周四晚查經班《歸回先知書研讀(11)撒迦利亞書》 講員:曾裔貴牧師} & 2023-12-29 & \href{https://youtube.com/watch?v=MkLoNpPCJk4}{\texttt{MkLoNpPCJk4}} \\
哥林多前書 9:24-27 & \hyperref[sec:lIj_mrEVPGg]{2023-12-30 晚堂《為主奔跑》哥林多前書 9\_24-27 講員:陳一華牧師} & 2024-01-03 & \href{https://youtube.com/watch?v=lIj-mrEVPGg}{\texttt{lIj-mrEVPGg}} \\
路加福音 2:1-7 & \hyperref[sec:c3xZAeZuwD0]{2023-12-31 加恩樓《在你心裡有空處給主嗎?》路加福音 2\_1-7 講員:謝靄薇牧師} & 2024-01-03 & \href{https://youtube.com/watch?v=c3xZAeZuwD0}{\texttt{c3xZAeZuwD0}} \\
以賽亞書 61:1 & \hyperref[sec:p_xp8W8_F_Q]{2023-12-31 午堂《禧年執行者》以賽亞書 61\_1 講員:楊明麗姑娘} & 2024-01-03 & \href{https://youtube.com/watch?v=p_xp8W8-F_Q}{\texttt{p\_xp8W8-F\_Q}} \\
撒母耳記上 15:17-23 & \hyperref[sec:R1nryPpcggU]{2023-12-31 早堂《作主聽命的兒女》撒母耳記上 15\_17-23 講員:曾敏姑娘} & 2024-01-03 & \href{https://youtube.com/watch?v=R1nryPpcggU}{\texttt{R1nryPpcggU}} \\
瑪拉基書  & \hyperref[sec:qkHV9uoLu6Y]{2024-01-04 周四晚查經班《瑪拉基書(1)》 講員:曾裔貴牧師} & 2024-01-05 & \href{https://youtube.com/watch?v=qkHV9uoLu6Y}{\texttt{qkHV9uoLu6Y}} \\
箴言 29:18 & \hyperref[sec:SWoA7YUZj_4]{2024-01-07 午堂《2024錦宣家．齊躍動》箴言 29\_18 講員:陳麗蟬牧師} & 2024-01-09 & \href{https://youtube.com/watch?v=SWoA7YUZj_4}{\texttt{SWoA7YUZj\_4}} \\
馬太福音 2:1-12 & \hyperref[sec:Y_5AXHxN_yY]{2024-01-07 早堂《星光引導》馬太福音 2\_1-12 講員:曾浩賢院長} & 2024-01-09 & \href{https://youtube.com/watch?v=Y-5AXHxN_yY}{\texttt{Y-5AXHxN\_yY}} \\
希伯來書 11:8-9 & \hyperref[sec:2um86Ip2zrs]{2024-01-14 午堂《亞伯拉罕信心的成長》希伯來書 11\_8-9 講員:楊明麗姑娘} & 2024-01-16 & \href{https://youtube.com/watch?v=2um86Ip2zrs}{\texttt{2um86Ip2zrs}} \\
使徒行傳 26:1-23 & \hyperref[sec:Rlc0zB_br4Y]{2024-01-14 早堂《沒有違背那從天上來的異象》使徒行傳 26\_1-23 講員:李富成牧師} & 2024-01-16 & \href{https://youtube.com/watch?v=Rlc0zB-br4Y}{\texttt{Rlc0zB-br4Y}} \\
瑪拉基書  & \hyperref[sec:dR8L9GrF4hQ]{2024-01-18 周四晚查經班《瑪拉基書(2)》 講員:曾裔貴牧師} & 2024-01-19 & \href{https://youtube.com/watch?v=dR8L9GrF4hQ}{\texttt{dR8L9GrF4hQ}} \\
創世記 16:1-13 & \hyperref[sec:_CjhJKqUcA8]{2024-01-20 晚堂《沒有被遺忘的女兒》創世記 16\_1-13 講員:曾裔貴牧師} & 2024-01-23 & \href{https://youtube.com/watch?v=-CjhJKqUcA8}{\texttt{-CjhJKqUcA8}} \\
創世記 11:1-9 & \hyperref[sec:DTr4fwkGAak]{2024-01-21 午堂《從神的創造心意再思教會的大使命》創世記 11\_1-9 講員:楊穎兒姑娘} & 2024-01-23 & \href{https://youtube.com/watch?v=DTr4fwkGAak}{\texttt{DTr4fwkGAak}} \\
創世記 16:8-14 & \hyperref[sec:pNQz6hqXFi8]{2024-01-21 早堂《憐憫人,讓人看見的神》創世記 16\_8-14 講員:鄧泰康傳道} & 2024-01-23 & \href{https://youtube.com/watch?v=pNQz6hqXFi8}{\texttt{pNQz6hqXFi8}} \\
瑪拉基書  & \hyperref[sec:jLSZtQ5_sq4]{2024-01-25 周四晚查經班《瑪拉基書(3)》 講員:曾裔貴牧師} & 2024-01-27 & \href{https://youtube.com/watch?v=jLSZtQ5_sq4}{\texttt{jLSZtQ5\_sq4}} \\
約翰福音 12:26 & \hyperref[sec:c0kCL7cehvg]{2024-01-28 午堂《蒙福的事奉》約翰福音 12\_26 講員:陳一華牧師} & 2024-01-31 & \href{https://youtube.com/watch?v=c0kCL7cehvg}{\texttt{c0kCL7cehvg}} \\
路加福音 23:40 & \hyperref[sec:FhL_HPKlT9E]{2024-01-28 早堂《談敬畏》路加福音 23\_40 講員:陳韋安牧師} & 2024-01-31 & \href{https://youtube.com/watch?v=FhL-HPKlT9E}{\texttt{FhL-HPKlT9E}} \\
瑪拉基書  & \hyperref[sec:Srg_1OhzOG8]{2024-02-01 周四晚查經班《瑪拉基書(4)》 講員:曾裔貴牧師} & 2024-02-02 & \href{https://youtube.com/watch?v=Srg_1OhzOG8}{\texttt{Srg\_1OhzOG8}} \\
創世記 26:12-33 & \hyperref[sec:vJ5BK_GiIWM]{2024-02-04 早堂《以撒的神》創世記 26\_12-33 講員:黃頌言姑娘} & 2024-02-06 & \href{https://youtube.com/watch?v=vJ5BK_GiIWM}{\texttt{vJ5BK\_GiIWM}} \\
創世記 26:1-23 & \hyperref[sec:HWk9pevFs74]{2024-02-04 午堂《豐盛生命:依靠順服》創世記 26\_1-23 講員:陳麗蟬牧師} & 2024-02-07 & \href{https://youtube.com/watch?v=HWk9pevFs74}{\texttt{HWk9pevFs74}} \\
瑪拉基書  & \hyperref[sec:ny6z0RWTL8g]{2024-02-08 周四晚查經班《瑪拉基書(5)》 講員:曾裔貴牧師} & 2024-02-09 & \href{https://youtube.com/watch?v=ny6z0RWTL8g}{\texttt{ny6z0RWTL8g}} \\
創世記 26:1-23 & \hyperref[sec:OWpYTLtxpsA]{2024-02-03 晚堂《豐盛生命:依靠順服》創世記 26\_1-23 講員:陳麗蟬牧師} & 2024-02-14 & \href{https://youtube.com/watch?v=OWpYTLtxpsA}{\texttt{OWpYTLtxpsA}} \\
但以理書 6:1-3 & \hyperref[sec:45mO3Udlc8E]{2024-02-11 聯合崇拜《龍年生肖道》但以理書 6\_1-3 講員:陳韋安牧師} & 2024-02-14 & \href{https://youtube.com/watch?v=45mO3Udlc8E}{\texttt{45mO3Udlc8E}} \\
馬太福音 5:3-10 & \hyperref[sec:QYUt9NIhVQs]{2024-02-10 新春崇拜《福從天降》馬太福音 5\_3-10 講員:陳一華牧師} & 2024-02-15 & \href{https://youtube.com/watch?v=QYUt9NIhVQs}{\texttt{QYUt9NIhVQs}} \\
瑪拉基書  & \hyperref[sec:oCSQzasymjk]{2024-02-15 周四晚查經班《瑪拉基書(6)》 講員:曾裔貴牧師} & 2024-02-16 & \href{https://youtube.com/watch?v=oCSQzasymjk}{\texttt{oCSQzasymjk}} \\
列王記上 19:1-18 & \hyperref[sec:C6eEbdGd3Ao]{2024-02-18 午堂《當你以為走到絕境時》列王記上 19\_1-18 講員:鄧泰康傳道} & 2024-02-19 & \href{https://youtube.com/watch?v=C6eEbdGd3Ao}{\texttt{C6eEbdGd3Ao}} \\
以弗所書 4:20-23 & \hyperref[sec:myBJVEoimI0]{2024-02-18 早堂《如果...可以重來》以弗所書 4\_20-23 講員:黃仲安傳道} & 2024-02-19 & \href{https://youtube.com/watch?v=myBJVEoimI0}{\texttt{myBJVEoimI0}} \\
瑪拉基書  & \hyperref[sec:4SEW84zBGCU]{2024-02-22 周四晚查經班《瑪拉基書(7)》 講員:曾裔貴牧師} & 2024-02-23 & \href{https://youtube.com/watch?v=4SEW84zBGCU}{\texttt{4SEW84zBGCU}} \\
馬可福音 8:31-38 & \hyperref[sec:dYh7F8XFvZo]{2024-02-24 晚堂《作門徒的代價》馬可福音 8\_31-38 講員:褚永華牧師} & 2024-02-27 & \href{https://youtube.com/watch?v=dYh7F8XFvZo}{\texttt{dYh7F8XFvZo}} \\
馬太福音 3:1-12 & \hyperref[sec:A8lXH1cPVBk]{2024-02-25 午堂《國度,權柄,榮耀》馬太福音 3\_1-12 講員:李思敬牧師} & 2024-02-27 & \href{https://youtube.com/watch?v=A8lXH1cPVBk}{\texttt{A8lXH1cPVBk}} \\
創世記 25:27-34 & \hyperref[sec:__QTchIrjDM]{2024-02-25 早堂《當你以為是理所當然》創世記 25\_27-34 講員:鄧泰康傳道} & 2024-02-27 & \href{https://youtube.com/watch?v=_-QTchIrjDM}{\texttt{\_-QTchIrjDM}} \\
瑪拉基書  & \hyperref[sec:exEZqq1nt80]{2024-02-29 周四晚查經班《瑪拉基書(8)》 講員:曾裔貴牧師} & 2024-03-02 & \href{https://youtube.com/watch?v=exEZqq1nt80}{\texttt{exEZqq1nt80}} \\
創世記 32:1-12-22-32 & \hyperref[sec:oyj7tRcaGj0]{2024-03-03 早堂《雅各的祈禱》創世記 32\_1-12,22-32,33\_1-4 講員:謝靄薇牧師} & 2024-03-04 & \href{https://youtube.com/watch?v=oyj7tRcaGj0}{\texttt{oyj7tRcaGj0}} \\
創世記 32:22-32 & \hyperref[sec:a1w0x_rlHHg]{2024-03-03 午堂《得著與失去》創世記 32\_22-32 講員:陳麗蟬牧師} & 2024-03-06 & \href{https://youtube.com/watch?v=a1w0x_rlHHg}{\texttt{a1w0x\_rlHHg}} \\
瑪拉基書  & \hyperref[sec:2A_RiSlG6Bg]{2024-03-07 周四晚查經班《瑪拉基書(9)》 講員:曾裔貴牧師} & 2024-03-08 & \href{https://youtube.com/watch?v=2A-RiSlG6Bg}{\texttt{2A-RiSlG6Bg}} \\
加拉太書 5:13 & \hyperref[sec:AndPB0XPtoQ]{2024-03-10 午堂《不被捆綁的自由》加拉太書 5\_13 講員:陳韋安牧師} & 2024-03-12 & \href{https://youtube.com/watch?v=AndPB0XPtoQ}{\texttt{AndPB0XPtoQ}} \\
馬可福音 1:14-20 & \hyperref[sec:U06UnMRzdJ4]{2024-03-10 早堂《作主門徒》馬可福音 1\_14-20 講員:董家驊牧師} & 2024-03-12 & \href{https://youtube.com/watch?v=U06UnMRzdJ4}{\texttt{U06UnMRzdJ4}} \\
詩篇 139:1-13-18-23-24 & \hyperref[sec:pdca3AXE1WI]{2024-03-09 晚堂《在靈修和生活中經歷神》詩篇 139\_1,13-18,23-24 講員:廖炳堂牧師} & 2024-03-13 & \href{https://youtube.com/watch?v=pdca3AXE1WI}{\texttt{pdca3AXE1WI}} \\
創世記 48:15-16 & \hyperref[sec:RZv3JKPRqaw]{2024-03-16 晚堂《一生牧養我的神》創世記 48\_15-16 講員:盧思源姑娘} & 2024-03-19 & \href{https://youtube.com/watch?v=RZv3JKPRqaw}{\texttt{RZv3JKPRqaw}} \\
創世記 47:27-31 & \hyperref[sec:82RgV2iRB78]{2024-03-17 早堂《雅各的敬拜人生》創世記 47\_27-31 講員:曾裔貴牧師} & 2024-03-21 & \href{https://youtube.com/watch?v=82RgV2iRB78}{\texttt{82RgV2iRB78}} \\
帖撒羅尼迦前書  & \hyperref[sec:ih_tESQH0P8]{2024-03-21 周四晚查經班《帖撒羅尼迦前書(1)》 講員:曾裔貴牧師} & 2024-03-22 & \href{https://youtube.com/watch?v=ih_tESQH0P8}{\texttt{ih\_tESQH0P8}} \\
約翰福音 12:12-19 & \hyperref[sec:u5__eddKCBg]{2024-03-23 晚堂《耶穌榮耀進城》約翰福音 12\_12-19 講員:黃頌言姑娘} & 2024-03-26 & \href{https://youtube.com/watch?v=u5--eddKCBg}{\texttt{u5--eddKCBg}} \\
約翰福音 12:1-26 & \hyperref[sec:94pJmtJo_JE]{2024-03-24 早堂《先生,我們想見耶穌》約翰福音 12\_1-26 講員:李炳光牧師} & 2024-03-26 & \href{https://youtube.com/watch?v=94pJmtJo-JE}{\texttt{94pJmtJo-JE}} \\
士師記 16:4-31 & \hyperref[sec:NKxHUXgOsyY]{2024-03-24 午堂《勝過致命誘惑的秘訣》士師記 16\_4-31 講員:林誠信牧師} & 2024-03-28 & \href{https://youtube.com/watch?v=NKxHUXgOsyY}{\texttt{NKxHUXgOsyY}} \\
帖撒羅尼迦前書  & \hyperref[sec:xRsKG7u5A0E]{2024-03-28 周四晚查經班《帖撒羅尼迦前書(2)》 講員:曾裔貴牧師} & 2024-03-30 & \href{https://youtube.com/watch?v=xRsKG7u5A0E}{\texttt{xRsKG7u5A0E}} \\
羅馬書 6:1-14 & \hyperref[sec:Su045rsIjTY]{2024-03-29 受苦節崇拜《耶穌為你死了,你呢?!》羅馬書 6\_1-14 講員:鄧泰康傳道} & 2024-04-02 & \href{https://youtube.com/watch?v=Su045rsIjTY}{\texttt{Su045rsIjTY}} \\
馬太福音 28:1-10 & \hyperref[sec:ez2giRGwvx8]{2024-03-31 復活節崇拜《復活:衝破生命的枷鎖》馬太福音 28\_1-10 講員:陳麗蟬牧師} & 2024-04-02 & \href{https://youtube.com/watch?v=ez2giRGwvx8}{\texttt{ez2giRGwvx8}} \\
創世記 30:14-16 & \hyperref[sec:_PqGg3XbEA0]{2024-04-06 晚堂《缺乏了愛的家庭–雅各,利亞,拉結》創世記 30\_14-16 講員:鄧泰康傳道} & 2024-04-09 & \href{https://youtube.com/watch?v=-PqGg3XbEA0}{\texttt{-PqGg3XbEA0}} \\
路加福音 24:36-43 & \hyperref[sec:TfOT_yj_ii8]{2024-04-07 早堂《轉化我們情緒困擾的耶穌》路加福音 24\_36-43 講員:區祥江博士} & 2024-04-10 & \href{https://youtube.com/watch?v=TfOT_yj-ii8}{\texttt{TfOT\_yj-ii8}} \\
帖撒羅尼迦前書  & \hyperref[sec:CWxLFOO62pc]{2024-04-11 周四晚查經班《帖撒羅尼迦前書(3)》 講員:曾裔貴牧師} & 2024-04-13 & \href{https://youtube.com/watch?v=CWxLFOO62pc}{\texttt{CWxLFOO62pc}} \\
路加福音 10:25-37 & \hyperref[sec:_2QFVuN84qA]{2024-04-13 晚堂《關愛鄰舍再思》路加福音 10\_25-37 講員:陳一華牧師} & 2024-04-16 & \href{https://youtube.com/watch?v=-2QFVuN84qA}{\texttt{-2QFVuN84qA}} \\
申命記 1:1-18 & \hyperref[sec:iqMY9wa3DY0]{2024-04-14 午堂《同行共拓主國度》申命記 1\_1-18 講員:戴浩輝牧師} & 2024-04-16 & \href{https://youtube.com/watch?v=iqMY9wa3DY0}{\texttt{iqMY9wa3DY0}} \\
創世記 37:1-36 & \hyperref[sec:y77Sfudracc]{2024-04-20 晚堂《主必追討》創世記 37\_1-36 講員:陳麗蟬牧師} & 2024-04-23 & \href{https://youtube.com/watch?v=y77Sfudracc}{\texttt{y77Sfudracc}} \\
創世記 37:2-4 & \hyperref[sec:kK4Sp5yFq94]{2024-04-21 午堂《約瑟的華服》創世記 37\_2-4;39\_1-6,19-23;41\_14,38-43 講員:楊明麗姑娘} & 2024-04-23 & \href{https://youtube.com/watch?v=kK4Sp5yFq94}{\texttt{kK4Sp5yFq94}} \\
帖撒羅尼迦前書  & \hyperref[sec:ZVwFpbgglfY]{2024-04-25 周四晚查經班《帖撒羅尼迦前書(4)》 講員:曾裔貴牧師} & 2024-04-28 & \href{https://youtube.com/watch?v=ZVwFpbgglfY}{\texttt{ZVwFpbgglfY}} \\
羅馬書 8:26-30 & \hyperref[sec:a_rRhGfuxUw]{2024-04-27 晚堂《再思天路歷程》羅馬書 8\_26-30  講員:曾裔貴牧師} & 2024-04-30 & \href{https://youtube.com/watch?v=a-rRhGfuxUw}{\texttt{a-rRhGfuxUw}} \\
腓立比書 2:17-30 & \hyperref[sec:w0vSwuU8gCk]{2024-04-28 早堂《同心事主的知己》腓立比書 2\_17-30  講員:馮耀榮牧師} & 2024-04-30 & \href{https://youtube.com/watch?v=w0vSwuU8gCk}{\texttt{w0vSwuU8gCk}} \\
撒迦利亞書 2:6-8 & \hyperref[sec:jRV45ettAd0]{2024-04-28 午堂《上帝眼中的瞳人》撒迦利亞書 2\_6-8  講員:陳韋安牧師} & 2024-05-01 & \href{https://youtube.com/watch?v=jRV45ettAd0}{\texttt{jRV45ettAd0}} \\
帖撒羅尼迦前書  & \hyperref[sec:sCNFnlmgxC0]{2024-05-02 周四晚查經班《帖撒羅尼迦前書(5)》 講員:曾裔貴牧師} & 2024-05-03 & \href{https://youtube.com/watch?v=sCNFnlmgxC0}{\texttt{sCNFnlmgxC0}} \\
創世記 39:1-6-19-23 & \hyperref[sec:U5CngR8W_hQ]{2024-05-04 晚堂《變幻人生主同行》創世記 39\_1-6,19-23  講員:陳麗蟬牧師} & 2024-05-07 & \href{https://youtube.com/watch?v=U5CngR8W-hQ}{\texttt{U5CngR8W-hQ}} \\
加拉太書 2:20 & \hyperref[sec:ltEd3gYtiPs]{2024-05-05 早堂《基督在那裡?》加拉太書 2\_20  講員:朱浩廉博士} & 2024-05-07 & \href{https://youtube.com/watch?v=ltEd3gYtiPs}{\texttt{ltEd3gYtiPs}} \\
馬太福音 11:2-19 & \hyperref[sec:1QRoYKXVPvQ]{2024-05-05 午堂《沒有最大,只有更大》馬太福音 11\_2-19  講員:張雲開老師} & 2024-05-08 & \href{https://youtube.com/watch?v=1QRoYKXVPvQ}{\texttt{1QRoYKXVPvQ}} \\
帖撒羅尼迦前書  & \hyperref[sec:VvxSAXf_Lww]{2024-05-09 周四晚查經班《帖撒羅尼迦前書(6)》 講員:曾裔貴牧師} & 2024-05-10 & \href{https://youtube.com/watch?v=VvxSAXf-Lww}{\texttt{VvxSAXf-Lww}} \\
創世記 3:20 & \hyperref[sec:BHhp4UIkEH0]{2024-05-11 晚堂《一個蒙恩的母親》創世記 3\_20  講員:楊穎兒姑娘} & 2024-05-14 & \href{https://youtube.com/watch?v=BHhp4UIkEH0}{\texttt{BHhp4UIkEH0}} \\
撒母耳記上 1:10-18-20 & \hyperref[sec:rgd8tUJeioY]{2024-05-12 午堂《母親的禱告》撒母耳記上 1\_10-18,20  講員:黃頌言姑娘} & 2024-05-14 & \href{https://youtube.com/watch?v=rgd8tUJeioY}{\texttt{rgd8tUJeioY}} \\
約翰福音 16:5-8 & \hyperref[sec:D8TmPX_lJJQ]{2024-05-18 晚堂《看啊!聖靈駕到!》約翰福音16\_5-8 講員:楊明麗姑娘} & 2024-05-21 & \href{https://youtube.com/watch?v=D8TmPX-lJJQ}{\texttt{D8TmPX-lJJQ}} \\
撒迦利亞書 4:1-10 & \hyperref[sec:AlakgRywpkk]{2024-05-19 午堂《靠主的靈方成事》撒迦利亞書4\_1-10 講員:陳麗蟬牧師} & 2024-05-22 & \href{https://youtube.com/watch?v=AlakgRywpkk}{\texttt{AlakgRywpkk}} \\
約翰福音 11:35 & \hyperref[sec:0Pem3SkLOJI]{2024-05-19 早堂《耶穌的眼淚》約翰福音11\_35 講員:陳韋安牧師} & 2024-05-23 & \href{https://youtube.com/watch?v=0Pem3SkLOJI}{\texttt{0Pem3SkLOJI}} \\
申命記 34:10-12 & \hyperref[sec:Z6ZZRqsI0K0]{2024-05-25 晚堂《面對面的認識》申命記 34\_10-12  講員:盧思源姑娘} & 2024-05-28 & \href{https://youtube.com/watch?v=Z6ZZRqsI0K0}{\texttt{Z6ZZRqsI0K0}} \\
  & \hyperref[sec:luI8_7RFevk]{2024-05-26 早堂 樓恩德牧師《耶穌復活的真諦》} & 2024-05-28 & \href{https://youtube.com/watch?v=luI8_7RFevk}{\texttt{luI8\_7RFevk}} \\
約珥書 2:28-32 & \hyperref[sec:QGag1Zr1Cdk]{2024-05-26 午堂《預備耶和華大而可畏的日子》約珥書 2\_28-32  講員:郭鴻標牧師} & 2024-05-29 & \href{https://youtube.com/watch?v=QGag1Zr1Cdk}{\texttt{QGag1Zr1Cdk}} \\
帖撒羅尼迦前書  & \hyperref[sec:eejVMLNyY8I]{2024-05-30 周四晚查經班《帖撒羅尼迦前書(7)》 講員:曾裔貴牧師} & 2024-05-31 & \href{https://youtube.com/watch?v=eejVMLNyY8I}{\texttt{eejVMLNyY8I}} \\
以弗所書 4:17-32 & \hyperref[sec:7PHC2DLBPCo]{2024-06-01 晚堂《黑暗與光明》以弗所書4\_17-32 講員:曾裔貴牧師} & 2024-06-04 & \href{https://youtube.com/watch?v=7PHC2DLBPCo}{\texttt{7PHC2DLBPCo}} \\
列王記下 5:1-19 & \hyperref[sec:Il1UURetEws]{2024-06-02 午堂《無名的小使女》列王記下5\_1-19 講員:楊明麗姑娘} & 2024-06-04 & \href{https://youtube.com/watch?v=Il1UURetEws}{\texttt{Il1UURetEws}} \\
歷代志上 13:9-14 & \hyperref[sec:ui5VIns3kj4]{2024-06-02 早堂《前車可鑑》歷代志上13\_9-14 講員:陳一華牧師} & 2024-06-05 & \href{https://youtube.com/watch?v=ui5VIns3kj4}{\texttt{ui5VIns3kj4}} \\
帖撒羅尼迦前書  & \hyperref[sec:DjV4K1Dg8GU]{2024-06-06 周四晚查經班《帖撒羅尼迦前書(8)》 講員:曾裔貴牧師} & 2024-06-07 & \href{https://youtube.com/watch?v=DjV4K1Dg8GU}{\texttt{DjV4K1Dg8GU}} \\
詩篇 84:1-12 & \hyperref[sec:CutYNgBop3c]{2024-06-09 午堂《永恆的居所》詩篇 84\_1-12  講員:馮耀榮牧師} & 2024-06-11 & \href{https://youtube.com/watch?v=CutYNgBop3c}{\texttt{CutYNgBop3c}} \\
路加福音 9:46-48 & \hyperref[sec:AT0lX1Ebpa4]{2024-06-09 早堂《接待小子的服侍》路加福音 9\_46-48  講員:蔣文忠傳道} & 2024-06-11 & \href{https://youtube.com/watch?v=AT0lX1Ebpa4}{\texttt{AT0lX1Ebpa4}} \\
帖撒羅尼迦前書  & \hyperref[sec:w32Ot_b7FcE]{2024-06-13 周四晚查經班《帖撒羅尼迦前書(9)》 講員:曾裔貴牧師} & 2024-06-14 & \href{https://youtube.com/watch?v=w32Ot_b7FcE}{\texttt{w32Ot\_b7FcE}} \\
民數記 14:24 & \hyperref[sec:1kzKFPo0unA]{2024-06-15 晚堂《專心跟從神》民數記 14\_24  講員:陳迦勒導師} & 2024-06-18 & \href{https://youtube.com/watch?v=1kzKFPo0unA}{\texttt{1kzKFPo0unA}} \\
詩篇 23:1-6 & \hyperref[sec:pibg95moZiM]{2024-06-16 午堂《有爸爸的愛,你不孤單》詩篇 23\_1-6  講員:鄧泰康傳道} & 2024-06-18 & \href{https://youtube.com/watch?v=pibg95moZiM}{\texttt{pibg95moZiM}} \\
路加福音 15:11-32 & \hyperref[sec:Yzl3vxypdqk]{2024-06-16 早堂《父親的心》路加福音 15\_11-32  講員:許亞光牧師} & 2024-06-19 & \href{https://youtube.com/watch?v=Yzl3vxypdqk}{\texttt{Yzl3vxypdqk}} \\
使徒行傳 1:1-15 & \hyperref[sec:aAKtjwj8FbQ]{2024-06-23 午堂《復興的等候》使徒行傳 1\_1-15  講員:陳麗蟬牧師} & 2024-06-25 & \href{https://youtube.com/watch?v=aAKtjwj8FbQ}{\texttt{aAKtjwj8FbQ}} \\
馬太福音 5:13-15 & \hyperref[sec:vr_pYc_qVuc]{2024-06-23 早堂《好ye 介紹返》馬太福音 5\_13-15  講員:陳韋安牧師} & 2024-06-25 & \href{https://youtube.com/watch?v=vr_pYc-qVuc}{\texttt{vr\_pYc-qVuc}} \\
帖撒羅尼迦前書  & \hyperref[sec:4L9GLJKp4e4]{2024-06-27 周四晚查經班《帖撒羅尼迦前書(10)》 講員:曾裔貴牧師} & 2024-06-29 & \href{https://youtube.com/watch?v=4L9GLJKp4e4}{\texttt{4L9GLJKp4e4}} \\
帖撒羅尼迦後書  & \hyperref[sec:uJ_qz3eVEf0]{2024-07-11 周四晚查經班《帖撒羅尼迦後書(1)》 講員:曾裔貴牧師} & 2024-07-12 & \href{https://youtube.com/watch?v=uJ_qz3eVEf0}{\texttt{uJ\_qz3eVEf0}} \\
帖撒羅尼迦後書  & \hyperref[sec:ntrA1aTVgRk]{2024-07-18 周四晚查經班《帖撒羅尼迦後書(2)》 講員:曾裔貴牧師} & 2024-07-19 & \href{https://youtube.com/watch?v=ntrA1aTVgRk}{\texttt{ntrA1aTVgRk}} \\
帖撒羅尼迦後書  & \hyperref[sec:ewzMm4Y0R5I]{2024-08-01 周四晚查經班《帖撒羅尼迦後書(3)》 講員:曾裔貴牧師} & 2024-08-06 & \href{https://youtube.com/watch?v=ewzMm4Y0R5I}{\texttt{ewzMm4Y0R5I}} \\
帖撒羅尼迦後書  & \hyperref[sec:bkamj1iGLMk]{2024-08-15 周四晚查經班《帖撒羅尼迦後書(4)》 講員:曾裔貴牧師} & 2024-08-16 & \href{https://youtube.com/watch?v=bkamj1iGLMk}{\texttt{bkamj1iGLMk}} \\
帖撒羅尼迦後書  & \hyperref[sec:IJNXTQniUCA]{2024-08-22 周四晚查經班《帖撒羅尼迦後書(5)》 講員:曾裔貴牧師} & 2024-08-23 & \href{https://youtube.com/watch?v=IJNXTQniUCA}{\texttt{IJNXTQniUCA}} \\
  & \hyperref[sec:I4c89OnZC6M]{2024-08-29 周四晚查經班《淺談末世論(1)》 講員:曾裔貴牧師} & 2024-08-30 & \href{https://youtube.com/watch?v=I4c89OnZC6M}{\texttt{I4c89OnZC6M}} \\
  & \hyperref[sec:a8KnWbij7Uk]{2024-09-19 周四晚查經班《淺談末世論(2)》 講員:曾裔貴牧師} & 2024-09-20 & \href{https://youtube.com/watch?v=a8KnWbij7Uk}{\texttt{a8KnWbij7Uk}} \\
  & \hyperref[sec:tl2Gc71LinQ]{2024-09-26 周四晚查經班《淺談末世論(3)》 講員:曾裔貴牧師} & 2024-09-27 & \href{https://youtube.com/watch?v=tl2Gc71LinQ}{\texttt{tl2Gc71LinQ}} \\
  & \hyperref[sec:H4W4IN0LZRY]{2024-10-03 周四晚查經班《淺談末世論(4)》 講員:曾裔貴牧師} & 2024-10-04 & \href{https://youtube.com/watch?v=H4W4IN0LZRY}{\texttt{H4W4IN0LZRY}} \\
  & \hyperref[sec:aTaCCC_F45U]{2024-10-10 周四晚查經班《改革之路:過去與未來.(1)》 講員:曾裔貴牧師} & 2024-10-12 & \href{https://youtube.com/watch?v=aTaCCC-F45U}{\texttt{aTaCCC-F45U}} \\
  & \hyperref[sec:Sb_mbdX9Sa4]{2024-10-24 周四晚查經班《改革之路(2):馬丁路德的十架神學》 講員:曾裔貴牧師} & 2024-10-25 & \href{https://youtube.com/watch?v=Sb-mbdX9Sa4}{\texttt{Sb-mbdX9Sa4}} \\
  & \hyperref[sec:ziJGxjL2HKY]{2024-11-07 周四晚查經班《改革之路(3)》 講員:曾裔貴牧師} & 2024-11-09 & \href{https://youtube.com/watch?v=ziJGxjL2HKY}{\texttt{ziJGxjL2HKY}} \\
  & \hyperref[sec:jvEmZbSGJ_A]{2024-11-14 周四晚查經班《改革之路(4):伊拉斯姆》 講員:曾裔貴牧師} & 2024-11-15 & \href{https://youtube.com/watch?v=jvEmZbSGJ_A}{\texttt{jvEmZbSGJ\_A}} \\
  & \hyperref[sec:DzI9N0JOL0s]{2024-12-12 周四晚查經班《改革之路(5):加爾文(一)》 講員:曾裔貴牧師} & 2024-12-13 & \href{https://youtube.com/watch?v=DzI9N0JOL0s}{\texttt{DzI9N0JOL0s}} \\
  & \hyperref[sec:4EM8nFC2gtc]{2025-01-02 周四晚查經班《改革之路(6):加爾文(二)》 講員:曾裔貴牧師} & 2025-01-03 & \href{https://youtube.com/watch?v=4EM8nFC2gtc}{\texttt{4EM8nFC2gtc}} \\
  & \hyperref[sec:3Sqlgs98HDM]{2025-01-09 周四晚查經班《改革之路(7):加爾文(三)》 講員:曾裔貴牧師} & 2025-01-10 & \href{https://youtube.com/watch?v=3Sqlgs98HDM}{\texttt{3Sqlgs98HDM}} \\
  & \hyperref[sec:Dn07Lw4F37U]{2025-01-16 周四晚查經班《改革之路(8):加爾文(四)》 講員:曾裔貴牧師} & 2025-01-17 & \href{https://youtube.com/watch?v=Dn07Lw4F37U}{\texttt{Dn07Lw4F37U}} \\
  & \hyperref[sec:5Lw6EDNBsNI]{2025-02-06 周四晚查經班《改革之路(9):諾克斯(一)》 講員:曾裔貴牧師} & 2025-02-07 & \href{https://youtube.com/watch?v=5Lw6EDNBsNI}{\texttt{5Lw6EDNBsNI}} \\
  & \hyperref[sec:q7QZXuHzcds]{2025-02-13 周四晚查經班《改革之路(10)完結篇:未來》 講員:曾裔貴牧師} & 2025-02-14 & \href{https://youtube.com/watch?v=q7QZXuHzcds}{\texttt{q7QZXuHzcds}} \\
羅馬書  & \hyperref[sec:mCrcbeplxo4]{2025-02-20 周四晚查經班《羅馬書(1)引言:保羅其人其事》 講員:曾裔貴牧師} & 2025-02-22 & \href{https://youtube.com/watch?v=mCrcbeplxo4}{\texttt{mCrcbeplxo4}} \\
羅馬書  & \hyperref[sec:TdP7rHPK3XI]{2025-03-06 周四晚查經班《羅馬書(2):分題:神的義(一)》 講員:曾裔貴牧師} & 2025-03-08 & \href{https://youtube.com/watch?v=TdP7rHPK3XI}{\texttt{TdP7rHPK3XI}} \\
羅馬書  & \hyperref[sec:WjqHgH3CzNw]{2025-03-13 周四晚查經班《羅馬書(3):分題:神的義(二)》 講員:曾裔貴牧師} & 2025-03-14 & \href{https://youtube.com/watch?v=WjqHgH3CzNw}{\texttt{WjqHgH3CzNw}} \\
羅馬書  & \hyperref[sec:8LIvofDzP5M]{2025-03-20 周四晚查經班《羅馬書(4):分題:神的義(三)》 講員:曾裔貴牧師} & 2025-03-21 & \href{https://youtube.com/watch?v=8LIvofDzP5M}{\texttt{8LIvofDzP5M}} \\
羅馬書  & \hyperref[sec:nWNWvRiUB3k]{2025-03-27 周四晚查經班《羅馬書(5):分題:神的義(四)》 講員:曾裔貴牧師} & 2025-03-29 & \href{https://youtube.com/watch?v=nWNWvRiUB3k}{\texttt{nWNWvRiUB3k}} \\
羅馬書  & \hyperref[sec:7C1Uqd4_AYQ]{2025-04-03 周四晚查經班《羅馬書(6):分題:神的義(五)》 講員:曾裔貴牧師} & 2025-04-05 & \href{https://youtube.com/watch?v=7C1Uqd4-AYQ}{\texttt{7C1Uqd4-AYQ}} \\
羅馬書  & \hyperref[sec:JU6MXStfi9g]{2025-04-17 周四晚查經班《羅馬書(7):分題:神的義(六)》 講員:曾裔貴牧師} & 2025-04-19 & \href{https://youtube.com/watch?v=JU6MXStfi9g}{\texttt{JU6MXStfi9g}} \\
羅馬書  & \hyperref[sec:b2w9fc63m3c]{2025-04-24 周四晚查經班《羅馬書(8):分題:神的義(七)》 講員:曾裔貴牧師} & 2025-04-25 & \href{https://youtube.com/watch?v=b2w9fc63m3c}{\texttt{b2w9fc63m3c}} \\
羅馬書  & \hyperref[sec:DvUFUpNSEFw]{2025-05-08 周四晚查經班《羅馬書(9):分題:神的義(八)》 講員:曾裔貴牧師} & 2025-05-09 & \href{https://youtube.com/watch?v=DvUFUpNSEFw}{\texttt{DvUFUpNSEFw}} \\
羅馬書  & \hyperref[sec:BtALIL64hKM]{2025-05-15 周四晚查經班《羅馬書(10):分題:舊約中的因信稱義(一)》 講員:曾裔貴牧師} & 2025-05-16 & \href{https://youtube.com/watch?v=BtALIL64hKM}{\texttt{BtALIL64hKM}} \\
羅馬書  & \hyperref[sec:QTpL9O5GR3E]{2025-05-22 周四晚查經班《羅馬書(11):分題:稱義帶來與神復和》 講員:曾裔貴牧師} & 2025-05-23 & \href{https://youtube.com/watch?v=QTpL9O5GR3E}{\texttt{QTpL9O5GR3E}} \\
羅馬書  & \hyperref[sec:23gUJZhHg2o]{2025-05-22 周四晚查經班《羅馬書(11):分題:稱義帶來與神復和》 講員:曾裔貴牧師 NEW} & 2025-05-24 & \href{https://youtube.com/watch?v=23gUJZhHg2o}{\texttt{23gUJZhHg2o}} \\
羅馬書  & \hyperref[sec:r2uz84DdreY]{2025-06-05 周四晚查經班《羅馬書(12):分題:因信稱義的福》 講員:曾裔貴牧師} & 2025-06-06 & \href{https://youtube.com/watch?v=r2uz84DdreY}{\texttt{r2uz84DdreY}} \\
羅馬書  & \hyperref[sec:G8aDRsOGytw]{2025-06-12 周四晚查經班《羅馬書(13):分題:克勝罪惡的新生命》 講員:曾裔貴牧師} & 2025-06-13 & \href{https://youtube.com/watch?v=G8aDRsOGytw}{\texttt{G8aDRsOGytw}} \\
羅馬書  & \hyperref[sec:7LZOyo5Zt3c]{2025-06-19 周四晚查經班《羅馬書(14):分題:不可犯罪?可不犯罪》 講員:曾裔貴牧師} & 2025-06-20 & \href{https://youtube.com/watch?v=7LZOyo5Zt3c}{\texttt{7LZOyo5Zt3c}} \\
羅馬書  & \hyperref[sec:wh3ZqIs7b68]{2025-06-26 周四晚查經班《羅馬書(15):分題:不可犯罪?可不犯罪(二)》 講員:曾裔貴牧師} & 2025-06-27 & \href{https://youtube.com/watch?v=wh3ZqIs7b68}{\texttt{wh3ZqIs7b68}} \\
羅馬書  & \hyperref[sec:PTM8SSdjvVE]{2025-07-03 周四晚查經班《羅馬書(16):分題:恩㤟與律法》 講員:曾裔貴牧師} & 2025-07-04 & \href{https://youtube.com/watch?v=PTM8SSdjvVE}{\texttt{PTM8SSdjvVE}} \\
羅馬書  & \hyperref[sec:HoQGEbdeqSk]{2025-07-24 周四晚查經班《羅馬書(17):分題:恩㤟與律法》 講員:曾裔貴牧師} & 2025-07-26 & \href{https://youtube.com/watch?v=HoQGEbdeqSk}{\texttt{HoQGEbdeqSk}} \\
羅馬書  & \hyperref[sec:mDmug5anuis]{2025-07-31 周四晚查經班《羅馬書(18):分題:聖靈帶動萬物復興》 講員:曾裔貴牧師} & 2025-08-02 & \href{https://youtube.com/watch?v=mDmug5anuis}{\texttt{mDmug5anuis}} \\
羅馬書  & \hyperref[sec:HwCOo8EiKp0]{2025-08-14 周四晚查經班《羅馬書(19):分題:聖靈帶動萬物復興二》 講員:曾裔貴牧師} & 2025-08-18 & \href{https://youtube.com/watch?v=HwCOo8EiKp0}{\texttt{HwCOo8EiKp0}} \\
羅馬書  & \hyperref[sec:3_mKpJn_jtA]{2025-08-21 周四晚查經班《羅馬書(20):分題:神的選民猶太人可得救嗎?》 講員:曾裔貴牧師} & 2025-08-24 & \href{https://youtube.com/watch?v=3-mKpJn-jtA}{\texttt{3-mKpJn-jtA}} \\
  & \hyperref[sec:MB59D1BA0Dc]{宣道會錦繡堂2025年8月30日晚堂崇拜直播} & 2025-08-30 & \href{https://youtube.com/watch?v=MB59D1BA0Dc}{\texttt{MB59D1BA0Dc}} \\
羅馬書  & \hyperref[sec:3TbpwI2Yugs]{2025-09-05 周五晚查經班《羅馬書(21):分題:信心的義與福音的普世性》 講員:曾裔貴牧師} & 2025-09-07 & \href{https://youtube.com/watch?v=3TbpwI2Yugs}{\texttt{3TbpwI2Yugs}} \\
羅馬書  & \hyperref[sec:1GDpbDYTYf0]{2025-09-12 周五晚查經班《羅馬書(22):福音:人的選擇與後果》 講員:曾裔貴牧師} & 2025-09-13 & \href{https://youtube.com/watch?v=1GDpbDYTYf0}{\texttt{1GDpbDYTYf0}} \\
羅馬書  & \hyperref[sec:ZdF__6Sl2No]{2025-09-26 周五晚查經班《羅馬書(23):分題:神選民與恩典—神的不變救恩》 講員:曾裔貴牧師} & 2025-09-27 & \href{https://youtube.com/watch?v=ZdF-_6Sl2No}{\texttt{ZdF-\_6Sl2No}} \\
羅馬書  & \hyperref[sec:bw6FVf5f_n0]{2025-10-10 周五晚查經班《羅馬書(24):分題:神國救恩歷史的智慧》 講員:曾裔貴牧師} & 2025-10-11 & \href{https://youtube.com/watch?v=bw6FVf5f_n0}{\texttt{bw6FVf5f\_n0}} \\
羅馬書  & \hyperref[sec:FAt1alHLvdg]{2025-10-17 周五晚查經班《羅馬書(25):分題:從奧秘到信徒的生活》 講員:曾裔貴牧師} & 2025-10-19 & \href{https://youtube.com/watch?v=FAt1alHLvdg}{\texttt{FAt1alHLvdg}} \\
羅馬書  & \hyperref[sec:yNx82XHOkbk]{2025-10-31 周五晚查經班《羅馬書(26):分題:教會是發掘,發展及發揮恩賜的場景》 講員:曾裔貴牧師 1} & 2025-11-03 & \href{https://youtube.com/watch?v=yNx82XHOkbk}{\texttt{yNx82XHOkbk}} \\
羅馬書  & \hyperref[sec:B5vFf1eCUBA]{2025-11-14 周五晚查經班《羅馬書(27):分題:合乎中道的自我認識與恩賜的正確運用》 講員:曾裔貴牧師} & 2025-11-15 & \href{https://youtube.com/watch?v=B5vFf1eCUBA}{\texttt{B5vFf1eCUBA}} \\
羅馬書  & \hyperref[sec:0TRElWqGJ4E]{2025-11-21 周五晚查經班《羅馬書(28):分題:恩賜是要來用的,在教會的合一場景中運用》 講員:曾裔貴牧師} & 2025-11-29 & \href{https://youtube.com/watch?v=0TRElWqGJ4E}{\texttt{0TRElWqGJ4E}} \\
羅馬書  & \hyperref[sec:WdwoyEOQVUg]{2025-11-28 周五晚查經班《羅馬書(29):分題:在群體生活中的自我認知》 講員:曾裔貴牧師} & 2025-11-29 & \href{https://youtube.com/watch?v=WdwoyEOQVUg}{\texttt{WdwoyEOQVUg}} \\
羅馬書  & \hyperref[sec:2B_XWriXVXo]{2025-12-12 周五晚查經班《羅馬書(30):分題:愛與順服》 講員:曾裔貴牧師} & 2025-12-13 & \href{https://youtube.com/watch?v=2B_XWriXVXo}{\texttt{2B\_XWriXVXo}} \\
羅馬書  & \hyperref[sec:ego6cJukrIE]{2025-12-19 周五晚查經班《羅馬書(31):分題:愛與順服(二)》 講員:曾裔貴牧師} & 2025-12-21 & \href{https://youtube.com/watch?v=ego6cJukrIE}{\texttt{ego6cJukrIE}} \\
羅馬書  & \hyperref[sec:aePwbTZJwec]{2026-01-02 周五晚查經班《羅馬書(32):分題:如何處理教會內「非核心信仰」的差異》 講員:曾裔貴牧師} & 2026-01-03 & \href{https://youtube.com/watch?v=aePwbTZJwec}{\texttt{aePwbTZJwec}} \\
羅馬書  & \hyperref[sec:nB_6d0PSa6M]{2026-01-09 周五晚查經班《羅馬書(33):分題:如何處理教會內「非核心信仰」的差異(二)》 講員:曾裔貴牧師} & 2026-01-10 & \href{https://youtube.com/watch?v=nB_6d0PSa6M}{\texttt{nB\_6d0PSa6M}} \\
羅馬書  & \hyperref[sec:pe8biigoai0]{2026-01-16 周五晚查經班《羅馬書(34):分題:擔得起,才配得進神的國》 講員:曾裔貴牧師} & 2026-01-17 & \href{https://youtube.com/watch?v=pe8biigoai0}{\texttt{pe8biigoai0}} \\
羅馬書  & \hyperref[sec:UU4dvSZbobM]{2026-01-23 周五晚查經班《羅馬書(三十五):分題:有盼望,有異象,有擔當的人》 講員:曾裔貴牧師} & 2026-01-24 & \href{https://youtube.com/watch?v=UU4dvSZbobM}{\texttt{UU4dvSZbobM}} \\
羅馬書  & \hyperref[sec:LvE6qlAUC04]{2026-02-06 周五晚查經班《羅馬書(三十六):分題:保羅 有擔當的使徒》 講員:曾裔貴牧師} & 2026-02-08 & \href{https://youtube.com/watch?v=LvE6qlAUC04}{\texttt{LvE6qlAUC04}} \\
羅馬書  & \hyperref[sec:3rvmIgiASuQ]{2026-02-13 周五晚查經班《羅馬書(三十七):分題:福音的實踐:從教義到生活》 講員:曾裔貴牧師} & 2026-02-14 & \href{https://youtube.com/watch?v=3rvmIgiASuQ}{\texttt{3rvmIgiASuQ}} \\
羅馬書  & \hyperref[sec:4pLVUVJIfgU]{2026-03-06 周五晚查經班《羅馬書(三十八):分題:使徒保羅的深情問安》 講員:曾裔貴牧師} & 2026-03-09 & \href{https://youtube.com/watch?v=4pLVUVJIfgU}{\texttt{4pLVUVJIfgU}} \\
羅馬書  & \hyperref[sec:oQdnBIyZabY]{2026-03-13 周五晚查經班《羅馬書(三十九):分題:來自保羅的問安(二)》 講員:曾裔貴牧師} & 2026-03-16 & \href{https://youtube.com/watch?v=oQdnBIyZabY}{\texttt{oQdnBIyZabY}} \\
羅馬書  & \hyperref[sec:HW8W_v51HY8]{2026-03-20 周五晚查經班《羅馬書(四十):分題:來自保羅的問安(三)》 講員:曾裔貴牧師} & 2026-03-24 & \href{https://youtube.com/watch?v=HW8W-v51HY8}{\texttt{HW8W-v51HY8}} \\
羅馬書  & \hyperref[sec:Ex_WASW50ZA]{2026-04-10 周五晚查經班《羅馬書(四十一):分題:來自保羅的問安(四)》 講員:曾裔貴牧師} & 2026-04-11 & \href{https://youtube.com/watch?v=Ex-WASW50ZA}{\texttt{Ex-WASW50ZA}} \\
羅馬書  & \hyperref[sec:WmefrJ6J_Lg]{2026-05-08 周五晚查經班《羅馬書(四十二):總論:福音帶來與神和好,所以教會必須向萬民宣講福音.》 講員:曾裔貴牧師} & 2026-05-10 & \href{https://youtube.com/watch?v=WmefrJ6J-Lg}{\texttt{WmefrJ6J-Lg}} \\
出埃及記  & \hyperref[sec:GROoST1kr4I]{2026-05-15 周五晚查經班《出埃及記研讀 進入應許之地(一)》 講員:曾裔貴牧師} & 2026-05-18 & \href{https://youtube.com/watch?v=GROoST1kr4I}{\texttt{GROoST1kr4I}} \\
出埃及記  & \hyperref[sec:kZF3v6CUrpM]{2026-05-22 周五晚查經班《出埃及記研讀(二)應許在黑暗中增長》 講員:曾裔貴牧師} & 2026-05-23 & \href{https://youtube.com/watch?v=kZF3v6CUrpM}{\texttt{kZF3v6CUrpM}} \\
出埃及記  & \hyperref[sec:4cUSeyCs18Q]{2026-05-29 周五晚查經班《出埃及記研讀(三)神所用的小人物:收生婆,母親,公主如何瓦解暴政》 講員:曾裔貴牧師} & 2026-05-30 & \href{https://youtube.com/watch?v=4cUSeyCs18Q}{\texttt{4cUSeyCs18Q}} \\
出埃及記  & \hyperref[sec:HMEU7PxsF0Y]{2026-06-05 周五晚查經班《出埃及記研讀(四)神國領袖的再造工程》 講員:曾裔貴牧師} & 2026-06-07 & \href{https://youtube.com/watch?v=HMEU7PxsF0Y}{\texttt{HMEU7PxsF0Y}} \\
出埃及記  & \hyperref[sec:cWahrBO910c]{2026-06-26 周五晚查經班《出埃及記研讀(五)五個「我不能」與一個「我是」》 講員:曾裔貴牧師} & 2026-06-27 & \href{https://youtube.com/watch?v=cWahrBO910c}{\texttt{cWahrBO910c}} \\
出埃及記  & \hyperref[sec:dj6fIzaDTMQ]{2026-07-17 周五晚查經班《出埃及記研讀(六)舊約先知的蒙召》 講員:曾裔貴牧師} & 2026-07-20 & \href{https://youtube.com/watch?v=dj6fIzaDTMQ}{\texttt{dj6fIzaDTMQ}} \\
出埃及記  & \hyperref[sec:FKTWi97mjmY]{2026-07-31 周五晚查經班《出埃及記研讀(七)舊約先知的蒙召(二)》 講員:曾裔貴牧師} & 2026-08-03 & \href{https://youtube.com/watch?v=FKTWi97mjmY}{\texttt{FKTWi97mjmY}} \\
\end{xltabular}
}
\newpage

\section{目錄\small{(順卷)}}
\label{sec:index_scriptual}
{ \scriptsize


\begin{xltabular}{\textwidth}{|p{0.15\textwidth} p{0.6\textwidth}|p{0.07\textwidth} p{0.1\textwidth}|}
\hline
啟示錄 21:1-7 & \hyperref[sec:IlLnn7WOgiE]{2023-01-01 加恩樓《上帝給我們的盼望及出路》啟示錄 21\_1-7 講員:鄧泰康傳道} & 2023-01-03 & \href{https://youtube.com/watch?v=IlLnn7WOgiE}{\texttt{IlLnn7WOgiE}} \\
路加福音 6:43-49 & \hyperref[sec:d2MEojJVUd4]{2023-01-01 午堂《趁新一年,遵主話而行》路加福音 6\_43-49 講員:曾裔貴牧師} & 2023-01-03 & \href{https://youtube.com/watch?v=d2MEojJVUd4}{\texttt{d2MEojJVUd4}} \\
  & \hyperref[sec:BQHcR2q2KwA]{2023-01-05 周四晚查經班《被擄歸回歷史研讀(2)》 講員:曾裔貴牧師} & 2023-01-06 & \href{https://youtube.com/watch?v=BQHcR2q2KwA}{\texttt{BQHcR2q2KwA}} \\
路加福音 8:1-15 & \hyperref[sec:gAWlXGsKQ90]{2023-01-07 晚堂《存忍耐的心結出果實》路加福音 8\_1-15 講員:曾裔貴牧師} & 2023-01-10 & \href{https://youtube.com/watch?v=gAWlXGsKQ90}{\texttt{gAWlXGsKQ90}} \\
撒迦利亞書 2:6-8 & \hyperref[sec:ejP2s680TsQ]{2023-01-08 加恩樓《上帝眼中的瞳人》撒迦利亞書 2\_6-8 講員:陳韋安牧師} & 2023-01-10 & \href{https://youtube.com/watch?v=ejP2s680TsQ}{\texttt{ejP2s680TsQ}} \\
提摩太後書 4:6-22 & \hyperref[sec:uJ8mvXZ5iqw]{2023-01-08 午堂《心友靈情》提摩太後書 4\_6-22 講員:曹偉彤院長} & 2023-01-10 & \href{https://youtube.com/watch?v=uJ8mvXZ5iqw}{\texttt{uJ8mvXZ5iqw}} \\
羅馬書 12:1-2 & \hyperref[sec:dkl9UnCmWfQ]{2019-01-20 加恩樓崇拜《得勝秘笈》羅馬書12\_1-2 講員:鄧泰康傳道} & 2023-01-13 & \href{https://youtube.com/watch?v=dkl9UnCmWfQ}{\texttt{dkl9UnCmWfQ}} \\
  & \hyperref[sec:ZwtOgh6Qmm4]{2023-01-12 周四晚查經班《被擄歸回歷史研讀(3)》 講員:曾裔貴牧師} & 2023-01-13 & \href{https://youtube.com/watch?v=ZwtOgh6Qmm4}{\texttt{ZwtOgh6Qmm4}} \\
馬太福音 13:24-30-36-43 & \hyperref[sec:c4O5F1DLbjI]{2023-01-15 加恩樓《成為麥子被收割》馬太福音 13\_24-30,36-43 講員:陳麗蟬牧師} & 2023-01-17 & \href{https://youtube.com/watch?v=c4O5F1DLbjI}{\texttt{c4O5F1DLbjI}} \\
以賽亞書 40:1-11-27-31 & \hyperref[sec:ezP3uupXjgk]{2023-01-15 午堂《如鷹展翅上騰》以賽亞書 40\_1-11,27-31 講員:郭鴻標牧師} & 2023-01-26 & \href{https://youtube.com/watch?v=ezP3uupXjgk}{\texttt{ezP3uupXjgk}} \\
羅馬書 12:13-21 & \hyperref[sec:ApJnSQORNII]{2023-01-22 加恩樓《只要祝福,不要咒詛》羅馬書 12\_13-21 講員:鄧泰康傳道} & 2023-01-26 & \href{https://youtube.com/watch?v=ApJnSQORNII}{\texttt{ApJnSQORNII}} \\
列王記上  & \hyperref[sec:9IcRN7FThks]{2023-01-22 新春崇拜《蒙福秘訣》列王記上 3: 3-9 講員:陳一華牧師} & 2023-01-26 & \href{https://youtube.com/watch?v=9IcRN7FThks}{\texttt{9IcRN7FThks}} \\
馬可福音 7:1-23 & \hyperref[sec:S5k9W4TGOeU]{2020-10-03 晚堂《用嘴來尊敬神的偽善者》馬可福音 7\_1-23 講員:高銘謙牧師} & 2023-01-27 & \href{https://youtube.com/watch?v=S5k9W4TGOeU}{\texttt{S5k9W4TGOeU}} \\
  & \hyperref[sec:mXSWAPPoVUc]{2023-01-26 周四晚查經班《被擄歸回歷史研讀(4)》 講員:曾裔貴牧師} & 2023-01-27 & \href{https://youtube.com/watch?v=mXSWAPPoVUc}{\texttt{mXSWAPPoVUc}} \\
哥林多後書 8:9 & \hyperref[sec:4sQReCkRl1k]{2023-01-28 晚堂《無比恩典》哥林多後書 8\_9 講員:陳一華牧師} & 2023-01-31 & \href{https://youtube.com/watch?v=4sQReCkRl1k}{\texttt{4sQReCkRl1k}} \\
馬太福音 18:21-35 & \hyperref[sec:65Ibdpx_Vbg]{2023-01-29 加恩樓《接納和饒恕》馬太福音 18\_21-35 講員:謝靄薇牧師} & 2023-02-01 & \href{https://youtube.com/watch?v=65Ibdpx-Vbg}{\texttt{65Ibdpx-Vbg}} \\
約翰福音 12:20-26 & \hyperref[sec:fpQjtet_7Iw]{2023-01-29 午堂《藉門徒見耶穌》約翰福音 12\_20-26 講員:林誠信牧師} & 2023-02-02 & \href{https://youtube.com/watch?v=fpQjtet-7Iw}{\texttt{fpQjtet-7Iw}} \\
  & \hyperref[sec:i9iPPyuXEWw]{2023-02-02 周四晚查經班《被擄歸回歷史研讀(5)》 講員:曾裔貴牧師} & 2023-02-04 & \href{https://youtube.com/watch?v=i9iPPyuXEWw}{\texttt{i9iPPyuXEWw}} \\
馬可福音  & \hyperref[sec:2Us7kVIsCO8]{2023-02-05 加恩樓《是賊窩?還是神的殿!!》馬可福音 11: 12-25 講員:鄧泰康傳道} & 2023-02-07 & \href{https://youtube.com/watch?v=2Us7kVIsCO8}{\texttt{2Us7kVIsCO8}} \\
馬可福音  & \hyperref[sec:k6XACT4Tfao]{2023-02-05 午堂《有果效的門徒生活》馬可福音 11: 12-26 講員:曾裔貴牧師} & 2023-02-07 & \href{https://youtube.com/watch?v=k6XACT4Tfao}{\texttt{k6XACT4Tfao}} \\
  & \hyperref[sec:lrudeppALpc]{2023-02-09 周四晚查經班《被擄歸回歷史研讀(6)》 講員:曾裔貴牧師} & 2023-02-10 & \href{https://youtube.com/watch?v=lrudeppALpc}{\texttt{lrudeppALpc}} \\
馬可福音 10:13-16 & \hyperref[sec:P577AmrjvF0]{2023-02-12 加恩樓《讓小孩子「到我這裡來」》馬可福音 10\_13-16 講員:吳景霖傳道} & 2023-02-14 & \href{https://youtube.com/watch?v=P577AmrjvF0}{\texttt{P577AmrjvF0}} \\
詩篇 139:1-24 & \hyperref[sec:raKQe4zsFjc]{2023-02-12 早堂《心知.神愛》詩篇 139\_1-24 講員:楊明麗姑娘} & 2023-02-14 & \href{https://youtube.com/watch?v=raKQe4zsFjc}{\texttt{raKQe4zsFjc}} \\
詩篇 139:1-24 & \hyperref[sec:la_oxDk2vYY]{2023-02-12 午堂《不分離不拋棄的愛》詩篇 139\_1-24 講員:陳麗蟬牧師} & 2023-02-17 & \href{https://youtube.com/watch?v=la_oxDk2vYY}{\texttt{la\_oxDk2vYY}} \\
詩篇 104:31-35 & \hyperref[sec:oRiBGTYbBfo]{2023-02-18 晚堂《新年願望》詩篇 104\_31-35 講員:張美薇牧師} & 2023-02-21 & \href{https://youtube.com/watch?v=oRiBGTYbBfo}{\texttt{oRiBGTYbBfo}} \\
詩篇 11:1-7 & \hyperref[sec:sY5rw9BFGjc]{2023-02-19 午堂《義路》詩篇 11\_1-7 講員:吳景霖傳道} & 2023-02-21 & \href{https://youtube.com/watch?v=sY5rw9BFGjc}{\texttt{sY5rw9BFGjc}} \\
加拉太書 3:21-29 & \hyperref[sec:jaVgwb2h0Gg]{2023-02-19 加恩樓《律法之上》加拉太書 3\_21-29 講員:陳韋安牧師} & 2023-02-22 & \href{https://youtube.com/watch?v=jaVgwb2h0Gg}{\texttt{jaVgwb2h0Gg}} \\
  & \hyperref[sec:v6FulXmsPK8]{2023-02-23 周四晚查經班《被擄歸回歷史研讀(7)》 講員:曾裔貴牧師} & 2023-02-25 & \href{https://youtube.com/watch?v=v6FulXmsPK8}{\texttt{v6FulXmsPK8}} \\
使徒行傳 8:14-24 & \hyperref[sec:jYaZ2w15SEQ]{2023-02-25 晚堂《屬靈沒有捷徑》使徒行傳 8\_14-24 講員:陳一華牧師} & 2023-02-28 & \href{https://youtube.com/watch?v=jYaZ2w15SEQ}{\texttt{jYaZ2w15SEQ}} \\
帖撒羅尼迦前書 5:16-18 & \hyperref[sec:Lu3g7hRuGdE]{2023-02-26 加恩樓《時刻經歷主的同在》帖撒羅尼迦前書 5\_16-18 講員:李適清博士} & 2023-02-28 & \href{https://youtube.com/watch?v=Lu3g7hRuGdE}{\texttt{Lu3g7hRuGdE}} \\
詩篇 73:1-28 & \hyperref[sec:d9OO9rhlq8k]{2023-02-26 午堂《當信念被動搖》詩篇 73\_1-28 講員:鄧泰康傳道} & 2023-03-01 & \href{https://youtube.com/watch?v=d9OO9rhlq8k}{\texttt{d9OO9rhlq8k}} \\
  & \hyperref[sec:E8bPEPFjp7c]{2023-03-02 周四晚查經班《被擄歸回歷史研讀(8)》 講員:曾裔貴牧師} & 2023-03-03 & \href{https://youtube.com/watch?v=E8bPEPFjp7c}{\texttt{E8bPEPFjp7c}} \\
歌羅西書 1:21-29 & \hyperref[sec:M8De3gUONVA]{2023-03-05 午堂《全力以赴的福音執事》歌羅西書 1\_21-29 講員:曾裔貴牧師} & 2023-03-07 & \href{https://youtube.com/watch?v=M8De3gUONVA}{\texttt{M8De3gUONVA}} \\
詩篇 46:1-11 & \hyperref[sec:BXGw_xefZws]{2023-03-05 加恩樓《上帝是我們患難中隨時的幫助》詩篇 46\_1-11 講員:鄧泰康傳道} & 2023-03-08 & \href{https://youtube.com/watch?v=BXGw_xefZws}{\texttt{BXGw\_xefZws}} \\
創世記 18:18-19 & \hyperref[sec:GHx9ZnTIMjk]{2023-03-11 晚堂《上帝子民要秉公行義》創世記18\_18-19 講員:黃頌言姑娘} & 2023-03-14 & \href{https://youtube.com/watch?v=GHx9ZnTIMjk}{\texttt{GHx9ZnTIMjk}} \\
使徒行傳 9:26-30-11:22-26 & \hyperref[sec:0LbO5gq5ujg]{2023-03-12 午堂《生命師傅的典範》使徒行傳9\_26-30,11\_22-26 講員:陳一華牧師} & 2023-03-14 & \href{https://youtube.com/watch?v=0LbO5gq5ujg}{\texttt{0LbO5gq5ujg}} \\
創世記 18:16-21 & \hyperref[sec:1tknD6q3Urw]{2023-03-12 加恩樓《守主道‧行公義》創世記18\_16-21 講員:陳麗蟬牧師} & 2023-03-15 & \href{https://youtube.com/watch?v=1tknD6q3Urw}{\texttt{1tknD6q3Urw}} \\
歌羅西書 3:1-17 & \hyperref[sec:BtpgDcAKeBQ]{2023-03-19 午堂《天地之間》歌羅西書3\_1-17 講員:吳景霖傳道} & 2023-03-21 & \href{https://youtube.com/watch?v=BtpgDcAKeBQ}{\texttt{BtpgDcAKeBQ}} \\
歌羅西書 3:7-14 & \hyperref[sec:W3jetg3vhiA]{2023-03-19 早堂《信主之後》歌羅西書3\_7-14 講員:楊明麗姑娘} & 2023-03-21 & \href{https://youtube.com/watch?v=W3jetg3vhiA}{\texttt{W3jetg3vhiA}} \\
馬太福音 25:1-13 & \hyperref[sec:wPsCGHG1pLA]{2023-03-19 加恩樓《心理預備,宏觀資訊的裝備,行動準備》馬太福音25\_1-13 講員:郭鴻標牧師} & 2023-03-22 & \href{https://youtube.com/watch?v=wPsCGHG1pLA}{\texttt{wPsCGHG1pLA}} \\
以弗所書 3:14-21 & \hyperref[sec:rNmdtGHU_n0]{2023-03-26 加恩樓《按神的心意成就一切》以弗所書3\_14-21 講員:鄧泰康傳道} & 2023-03-28 & \href{https://youtube.com/watch?v=rNmdtGHU-n0}{\texttt{rNmdtGHU-n0}} \\
提摩太前書 6:3-10 & \hyperref[sec:pzC3e8fjAp8]{2023-03-26 早堂《敬虔與知足》提摩太前書6\_3-10 講員:陳靜娟姑娘} & 2023-03-28 & \href{https://youtube.com/watch?v=pzC3e8fjAp8}{\texttt{pzC3e8fjAp8}} \\
  & \hyperref[sec:yk5Oh8RBHLY]{2023-03-30 周四晚查經班《被擄歸回歷史研讀(9)》 講員:曾裔貴牧師} & 2023-04-01 & \href{https://youtube.com/watch?v=yk5Oh8RBHLY}{\texttt{yk5Oh8RBHLY}} \\
馬太福音 21:1-11 & \hyperref[sec:x06IlO4FHos]{2023-03-26 午堂《看你的王》馬太福音21\_1-11 講員:李思敬牧師} & 2023-04-03 & \href{https://youtube.com/watch?v=x06IlO4FHos}{\texttt{x06IlO4FHos}} \\
約翰一書 5:20-21 & \hyperref[sec:ythaF42qv_4]{2023-04-02 午堂《信仰求真》約翰一書5\_20-21 講員:曾裔貴牧師} & 2023-04-04 & \href{https://youtube.com/watch?v=ythaF42qv_4}{\texttt{ythaF42qv\_4}} \\
提摩太前書 6:6-10 & \hyperref[sec:_BOzzljBDkA]{2023-03-25 晚堂《知足人生》提摩太前書6\_6-10 講員:盧思源姑娘} & 2023-04-05 & \href{https://youtube.com/watch?v=-BOzzljBDkA}{\texttt{-BOzzljBDkA}} \\
使徒行傳 9:26-30-11:22-26 & \hyperref[sec:lYwkS21XVXg]{2023-04-02 加恩樓《生命師傅的典範》使徒行傳9\_26-30,11\_22-26 講員:陳一華牧師} & 2023-04-05 & \href{https://youtube.com/watch?v=lYwkS21XVXg}{\texttt{lYwkS21XVXg}} \\
  & \hyperref[sec:Am9T43vQMKQ]{2023-04-06 周四晚查經班《被擄歸回歷史研讀(10)》 講員:曾裔貴牧師} & 2023-04-07 & \href{https://youtube.com/watch?v=Am9T43vQMKQ}{\texttt{Am9T43vQMKQ}} \\
  & \hyperref[sec:Pndq2oLIHnI]{2023-04-13 周四晚查經班《被擄歸回歷史研讀(11)》 講員:曾裔貴牧師} & 2023-04-15 & \href{https://youtube.com/watch?v=Pndq2oLIHnI}{\texttt{Pndq2oLIHnI}} \\
路加福音 10:38-42 & \hyperref[sec:7t8JFmETTzs]{2023-04-16 加恩樓《世界不能奪去的福份》路加福音10\_38-42 講員:蔣文忠傳道} & 2023-04-18 & \href{https://youtube.com/watch?v=7t8JFmETTzs}{\texttt{7t8JFmETTzs}} \\
 5:1-14 & \hyperref[sec:cuPMYAewAe8]{2023-04-16 午堂《敬拜那又能又配的羔羊》啟示錄5\_1-14 講員:陳麗蟬牧師} & 2023-04-18 & \href{https://youtube.com/watch?v=cuPMYAewAe8}{\texttt{cuPMYAewAe8}} \\
 5:1-14 & \hyperref[sec:Yd_an_Ish1s]{2023-04-16 早堂《異象中的羔羊是世人惟一的希望》啟示錄5\_1-14 講員:謝靄薇牧師} & 2023-04-19 & \href{https://youtube.com/watch?v=Yd-an_Ish1s}{\texttt{Yd-an\_Ish1s}} \\
  & \hyperref[sec:SLUSAYk3lm4]{2023-04-20 周四晚查經班《被擄歸回歷史研讀(12)》 講員:曾裔貴牧師} & 2023-04-21 & \href{https://youtube.com/watch?v=SLUSAYk3lm4}{\texttt{SLUSAYk3lm4}} \\
哥林多前書 15:12-34 & \hyperref[sec:qxsLxZod7kk]{2023-04-22 晚堂《基督 ─ 我們的初熟果子》哥林多前書 15\_12-34 講員:黃頌言姑娘} & 2023-04-25 & \href{https://youtube.com/watch?v=qxsLxZod7kk}{\texttt{qxsLxZod7kk}} \\
希伯來書 12:1-3 & \hyperref[sec:rnbJeDm_hDc]{2023-04-23 加恩樓《奔跑吧!》希伯來書 12\_1-3 講員:周弘毅牧師} & 2023-04-25 & \href{https://youtube.com/watch?v=rnbJeDm_hDc}{\texttt{rnbJeDm\_hDc}} \\
路加福音 16:1-13 & \hyperref[sec:l5hkKLKHhA0]{2023-04-23 午堂《不義與忠心》路加福音 16\_1-13 講員:楊慶球牧師} & 2023-04-25 & \href{https://youtube.com/watch?v=l5hkKLKHhA0}{\texttt{l5hkKLKHhA0}} \\
創世記 22:1-14 & \hyperref[sec:J_KPU0uZysw]{2023-04-29 晚堂《我在這裡》創世記 22\_1-14 講員:褚永華牧師} & 2023-05-02 & \href{https://youtube.com/watch?v=J_KPU0uZysw}{\texttt{J\_KPU0uZysw}} \\
加拉太書 6:12-16 & \hyperref[sec:bNffFxPN8ik]{2023-04-30 加恩樓《我們所誇的是什麼?》加拉太書 6\_12-16 講員:鄧泰康傳道} & 2023-05-02 & \href{https://youtube.com/watch?v=bNffFxPN8ik}{\texttt{bNffFxPN8ik}} \\
加拉太書 3:21-29 & \hyperref[sec:sHtcf0yp8MI]{2023-04-30 午堂《律法之上》加拉太書 3\_21-29 講員:陳韋安牧師} & 2023-05-02 & \href{https://youtube.com/watch?v=sHtcf0yp8MI}{\texttt{sHtcf0yp8MI}} \\
加拉太書 6:11-16 & \hyperref[sec:FRK_krDh8hE]{2023-04-30 早堂《怎麼誇口?》加拉太書 6\_11-16 講員:楊明麗姑娘} & 2023-05-02 & \href{https://youtube.com/watch?v=FRK_krDh8hE}{\texttt{FRK\_krDh8hE}} \\
  & \hyperref[sec:AGfClw0jrYY]{2023-05-04 周四晚查經班《被擄歸回歷史研讀(13)》 講員:曾裔貴牧師} & 2023-05-06 & \href{https://youtube.com/watch?v=AGfClw0jrYY}{\texttt{AGfClw0jrYY}} \\
約翰一書 3:11-18 & \hyperref[sec:eLIEUCC43fA]{2023-05-07 加恩樓崇拜《共融》約翰一書 3\_11-18 講員:吳景霖傳道} & 2023-05-09 & \href{https://youtube.com/watch?v=eLIEUCC43fA}{\texttt{eLIEUCC43fA}} \\
  & \hyperref[sec:_3DxJVmmiXg]{2023-05-11 周四晚查經班《被擄歸回歷史研讀(14)》 講員:曾裔貴牧師} & 2023-05-12 & \href{https://youtube.com/watch?v=-3DxJVmmiXg}{\texttt{-3DxJVmmiXg}} \\
申命記 6:1-9 & \hyperref[sec:A6GWJTLfDEo]{2023-05-14 加恩樓《一家之主》申命記 6\_1-9 講員:謝靄薇牧師} & 2023-05-16 & \href{https://youtube.com/watch?v=A6GWJTLfDEo}{\texttt{A6GWJTLfDEo}} \\
出埃及記 2:1-10 & \hyperref[sec:noei9nAUNuE]{2023-05-14 早堂《我們這一家》出埃及記 2\_1-10 講員:黃頌言姑娘} & 2023-05-16 & \href{https://youtube.com/watch?v=noei9nAUNuE}{\texttt{noei9nAUNuE}} \\
腓立比書 4:4 & \hyperref[sec:PwPuQ_V9PRU]{2023-05-13 晚堂《喜樂再思》腓立比書 4\_4 講員:陳一華牧師} & 2023-05-17 & \href{https://youtube.com/watch?v=PwPuQ-V9PRU}{\texttt{PwPuQ-V9PRU}} \\
使徒行傳 8:1-8 & \hyperref[sec:_W_gd9p_W_Y]{2023-05-14 午堂《福音帶來的生命改變》使徒行傳 8\_1-8 講員:曾裔貴牧師} & 2023-05-17 & \href{https://youtube.com/watch?v=-W_gd9p_W_Y}{\texttt{-W\_gd9p\_W\_Y}} \\
  & \hyperref[sec:MtaDL5MWzm0]{2023-05-18 周四晚查經班《被擄歸回歷史研讀(15)》 講員:曾裔貴牧師} & 2023-05-19 & \href{https://youtube.com/watch?v=MtaDL5MWzm0}{\texttt{MtaDL5MWzm0}} \\
腓立比書 3:7-16 & \hyperref[sec:We85QzJyMNQ]{2023-05-20 晚堂《丟棄與得著》腓立比書 3\_7-16 講員:盧思源姑娘} & 2023-05-23 & \href{https://youtube.com/watch?v=We85QzJyMNQ}{\texttt{We85QzJyMNQ}} \\
腓立比書 3:7-16 & \hyperref[sec:9tgCzMr4KBk]{2023-05-21 午堂《你得著珍寶了嗎?》腓立比書 3\_7-16 講員:鄧泰康傳道} & 2023-05-23 & \href{https://youtube.com/watch?v=9tgCzMr4KBk}{\texttt{9tgCzMr4KBk}} \\
腓立比書 3:13-15 & \hyperref[sec:vWqPqvo2gRk]{2023-05-21 早堂《與基督同跑》腓立比書 3\_13-15 講員:楊明麗姑娘} & 2023-05-23 & \href{https://youtube.com/watch?v=vWqPqvo2gRk}{\texttt{vWqPqvo2gRk}} \\
路加福音 1:5-25-57-66 & \hyperref[sec:wCpZ46fk4yY]{2023-05-21 加恩樓《不相稱的禱告》路加福音 1\_5-25,57-66 講員:陳韋安牧師} & 2023-05-24 & \href{https://youtube.com/watch?v=wCpZ46fk4yY}{\texttt{wCpZ46fk4yY}} \\
羅馬書 15:22-29 & \hyperref[sec:aPVZvLNc78I]{2023-05-27 晚堂《五餅二魚的團契》羅馬書 15\_22-29 講員:陳麗蟬牧師} & 2023-05-30 & \href{https://youtube.com/watch?v=aPVZvLNc78I}{\texttt{aPVZvLNc78I}} \\
使徒行傳 2:1-21 & \hyperref[sec:hSecMeCWH4s]{2023-05-28 加恩樓《聖靈充滿的教會》使徒行傳 2\_1-21 講員:曾浩賢院長} & 2023-05-30 & \href{https://youtube.com/watch?v=hSecMeCWH4s}{\texttt{hSecMeCWH4s}} \\
傳道書 11:1-6 & \hyperref[sec:H85tdsHmqKI]{2023-05-28 午堂《不知道中的信仰》傳道書 11\_1-6 講員:蔣文忠傳道} & 2023-05-30 & \href{https://youtube.com/watch?v=H85tdsHmqKI}{\texttt{H85tdsHmqKI}} \\
  & \hyperref[sec:C1vy6YO3h0g]{2023-06-01 周四晚查經班《被擄歸回歷史研讀(16)》 講員:曾裔貴牧師} & 2023-06-03 & \href{https://youtube.com/watch?v=C1vy6YO3h0g}{\texttt{C1vy6YO3h0g}} \\
馬太福音 28:11-20 & \hyperref[sec:WLN153QoHOE]{2023-06-04 加恩樓《後真相年代下的信仰反思》馬太福音 28\_11-20 講員:吳景霖傳道} & 2023-06-06 & \href{https://youtube.com/watch?v=WLN153QoHOE}{\texttt{WLN153QoHOE}} \\
路加福音 5:27-32 & \hyperref[sec:jJiE2EQ4f48]{2023-06-04 午堂《悔改的筵席》路加福音 5\_27-32 講員:楊明麗姑娘} & 2023-06-06 & \href{https://youtube.com/watch?v=jJiE2EQ4f48}{\texttt{jJiE2EQ4f48}} \\
路加福音 5:27-39 & \hyperref[sec:KzM_XrVrUvg]{2023-06-04 早堂《悔改的筵席》路加福音 5\_27-39 講員:曾裔貴牧師} & 2023-06-07 & \href{https://youtube.com/watch?v=KzM_XrVrUvg}{\texttt{KzM\_XrVrUvg}} \\
  & \hyperref[sec:qmt5oGwLkF4]{2023-06-08 周四晚查經班《被擄歸回歷史研讀(17)》 講員:曾裔貴牧師} & 2023-06-09 & \href{https://youtube.com/watch?v=qmt5oGwLkF4}{\texttt{qmt5oGwLkF4}} \\
路加福音 7:36-50 & \hyperref[sec:Fl2IkOLpxSI]{2023-06-10 晚堂《罪人的筵席》路加福音 7\_36-50 講員:吳景霖傳道} & 2023-06-13 & \href{https://youtube.com/watch?v=Fl2IkOLpxSI}{\texttt{Fl2IkOLpxSI}} \\
馬太福音 9:35-38 & \hyperref[sec:A_K04hyrSP8]{2023-06-11 加恩樓《建立無牆的教會》馬太福音 9\_35-38 講員:鄧泰康傳道} & 2023-06-13 & \href{https://youtube.com/watch?v=A-K04hyrSP8}{\texttt{A-K04hyrSP8}} \\
路加福音 7:36-50 & \hyperref[sec:lzufA9tLqUc]{2023-06-11 午堂《愛和謙卑的事奉》路加福音 7\_36-50 講員:陳麗蟬牧師} & 2023-06-13 & \href{https://youtube.com/watch?v=lzufA9tLqUc}{\texttt{lzufA9tLqUc}} \\
提摩太後書 3:16-17 & \hyperref[sec:eaFKBNr6R3g]{2023-06-18 加恩樓《有故事的人》提摩太後書 3\_16-17 講員:吳景霖傳道} & 2023-06-20 & \href{https://youtube.com/watch?v=eaFKBNr6R3g}{\texttt{eaFKBNr6R3g}} \\
提摩太後書 3:16-17 & \hyperref[sec:cNC4ab19lyE]{2023-06-18 早堂《說故事的爸媽》提摩太後書 3\_16-17 講員:吳景霖傳道} & 2023-06-20 & \href{https://youtube.com/watch?v=cNC4ab19lyE}{\texttt{cNC4ab19lyE}} \\
傳道書 11:1-6 & \hyperref[sec:I3SP1geCVE8]{2023-06-17 晚堂《不知道中的信仰》傳道書 11\_1-6 講員:蔣文忠傳道} & 2023-06-21 & \href{https://youtube.com/watch?v=I3SP1geCVE8}{\texttt{I3SP1geCVE8}} \\
路加福音 12:1-24 & \hyperref[sec:9x4ci5x2GZ0]{2023-06-18 午堂《基督徒的日常》路加福音 12\_1-24 講員:曾裔貴牧師} & 2023-06-21 & \href{https://youtube.com/watch?v=9x4ci5x2GZ0}{\texttt{9x4ci5x2GZ0}} \\
路加福音 12:35-48 & \hyperref[sec:keP18GcFGPs]{2023-06-24 晚堂《儆醒僕人筵席》路加福音 12\_35-48 講員:黃頌言姑娘} & 2023-06-27 & \href{https://youtube.com/watch?v=keP18GcFGPs}{\texttt{keP18GcFGPs}} \\
列王記上 18:1-6 & \hyperref[sec:i86dI_9QW94]{2023-06-25 加恩樓《掌管萬有的神》列王記上 18\_1-6 講員:陳一華牧師} & 2023-06-27 & \href{https://youtube.com/watch?v=i86dI_9QW94}{\texttt{i86dI\_9QW94}} \\
路加福音 19:12-23 & \hyperref[sec:tu4SyHPIS_U]{2023-06-25 早堂《未來的風險》路加福音 19\_12-23 講員:陳韋安牧師} & 2023-06-27 & \href{https://youtube.com/watch?v=tu4SyHPIS_U}{\texttt{tu4SyHPIS\_U}} \\
羅馬書 6:1-11 & \hyperref[sec:WZaLHtVToyI]{2023-06-25 午堂《克勝罪惡》羅馬書 6\_1-11 講員:褚永華牧師} & 2023-06-28 & \href{https://youtube.com/watch?v=WZaLHtVToyI}{\texttt{WZaLHtVToyI}} \\
  & \hyperref[sec:G4jQV76in2M]{2023-06-29 周四晚查經班《被擄歸回歷史研讀(18)》 講員:曾裔貴牧師} & 2023-06-30 & \href{https://youtube.com/watch?v=G4jQV76in2M}{\texttt{G4jQV76in2M}} \\
列王記上 19:1-18 & \hyperref[sec:W6aFxBk_o0s]{2023-07-02 加恩樓《當以為走到絕境時》列王記上 19\_1-18 講員:鄧泰康傳道} & 2023-07-05 & \href{https://youtube.com/watch?v=W6aFxBk_o0s}{\texttt{W6aFxBk\_o0s}} \\
路加福音 14:1-24 & \hyperref[sec:Vwm387KuCtc]{2023-07-02 早堂《筵席的再思》路加福音 14\_1-24 講員:曾裔貴牧師} & 2023-07-05 & \href{https://youtube.com/watch?v=Vwm387KuCtc}{\texttt{Vwm387KuCtc}} \\
  & \hyperref[sec:LQfxIAo88QM]{2023-07-06 周四晚查經班《被擄歸回歷史研讀(19)》 講員:曾裔貴牧師} & 2023-07-08 & \href{https://youtube.com/watch?v=LQfxIAo88QM}{\texttt{LQfxIAo88QM}} \\
路加福音 19:1-10 & \hyperref[sec:oxDrdn_rwTg]{2023-07-08 晚堂《奇妙的改變》路加福音 19\_1-10 講員:陳麗蟬牧師} & 2023-07-11 & \href{https://youtube.com/watch?v=oxDrdn-rwTg}{\texttt{oxDrdn-rwTg}} \\
提摩太前書詩篇羅馬書 1:19-20 & \hyperref[sec:87awTrKJFrY]{2023-07-09 加恩樓《平衡空間的人生》羅馬書 1\_19-20,提摩太前書 2\_5-6,詩篇 27\_13-14 講員:李盛林牧師} & 2023-07-11 & \href{https://youtube.com/watch?v=87awTrKJFrY}{\texttt{87awTrKJFrY}} \\
路加福音 19:1-10 & \hyperref[sec:QuJcmG3w_7k]{2023-07-09 午堂《愛在生命改變前》路加福音 19\_1-10 講員:楊明麗姑娘} & 2023-07-11 & \href{https://youtube.com/watch?v=QuJcmG3w-7k}{\texttt{QuJcmG3w-7k}} \\
尼希米記 11:1-3 & \hyperref[sec:IDYmD1EXhB4]{2023-07-16 加恩樓《別叫我英雄》尼希米記 11\_1-3 講員:曾裔貴牧師} & 2023-07-22 & \href{https://youtube.com/watch?v=IDYmD1EXhB4}{\texttt{IDYmD1EXhB4}} \\
使徒行傳 1:8-28:31 & \hyperref[sec:vtGjR9atL3U]{2023-07-16 午堂《放膽傳道,無人能禁止》使徒行傳 1\_8,28\_31 講員:郭鴻標牧師} & 2023-07-23 & \href{https://youtube.com/watch?v=vtGjR9atL3U}{\texttt{vtGjR9atL3U}} \\
路加福音 10:1-12 & \hyperref[sec:PIfMPVh5oUA]{2023-07-16 早堂《耶穌訓練門徒》路加福音 10\_1-12 講員:曾裔貴牧師} & 2023-07-23 & \href{https://youtube.com/watch?v=PIfMPVh5oUA}{\texttt{PIfMPVh5oUA}} \\
  & \hyperref[sec:zajnPwak_dU]{2023-07-13 周四晚查經班《被擄歸回歷史研讀(20)》 講員:曾裔貴牧師} & 2023-07-25 & \href{https://youtube.com/watch?v=zajnPwak-dU}{\texttt{zajnPwak-dU}} \\
歌羅西書 2:1-19 & \hyperref[sec:epe4ZHzeWUY]{2023-07-22 晚堂《打穩根基》歌羅西書 2\_1-19 講員:謝靄薇牧師} & 2023-07-25 & \href{https://youtube.com/watch?v=epe4ZHzeWUY}{\texttt{epe4ZHzeWUY}} \\
約翰福音 21:15-22 & \hyperref[sec:dFeC0BqGntE]{2023-07-23 加恩樓《跟隨我》約翰福音 21\_15-22 講員:陳靜娟姑娘} & 2023-07-25 & \href{https://youtube.com/watch?v=dFeC0BqGntE}{\texttt{dFeC0BqGntE}} \\
約翰福音 17:6-19 & \hyperref[sec:_FPkb_FFsjA]{2023-07-23 早堂《寶貴的禱告》約翰福音 17\_6-19 講員:曾敏姑娘} & 2023-07-25 & \href{https://youtube.com/watch?v=_FPkb_FFsjA}{\texttt{\_FPkb\_FFsjA}} \\
以斯帖記  & \hyperref[sec:iuYWy9ExSgQ]{2023-07-27 周四晚查經班《被擄歸回歷史研讀(21):以斯帖記》 講員:曾裔貴牧師} & 2023-07-29 & \href{https://youtube.com/watch?v=iuYWy9ExSgQ}{\texttt{iuYWy9ExSgQ}} \\
以弗所書 4:11-16 & \hyperref[sec:c_zS1WaX1q0]{2023-07-29 晚堂《建立基督的身體》以弗所書 4\_11-16 講員:盧思源姑娘} & 2023-08-08 & \href{https://youtube.com/watch?v=c-zS1WaX1q0}{\texttt{c-zS1WaX1q0}} \\
以弗所書 4:11-16 & \hyperref[sec:qI_KwviwgzU]{2023-07-30 加恩樓《建立基督的身體》以弗所書 4\_11-16 講員:楊明麗姑娘} & 2023-08-08 & \href{https://youtube.com/watch?v=qI_KwviwgzU}{\texttt{qI\_KwviwgzU}} \\
馬太福音 18:1-9 & \hyperref[sec:i3m_mXZnkC4]{2023-08-06 加恩樓《天國的小孩》馬太福音 18\_1-9 講員:鄧泰康傳道} & 2023-08-08 & \href{https://youtube.com/watch?v=i3m-mXZnkC4}{\texttt{i3m-mXZnkC4}} \\
歌羅西書 4:5-6 & \hyperref[sec:1pEPdzdQ8_4]{2023-08-06 午堂《著地氣的傳福音》歌羅西書 4\_5-6 講員:曾裔貴牧師} & 2023-08-08 & \href{https://youtube.com/watch?v=1pEPdzdQ8_4}{\texttt{1pEPdzdQ8\_4}} \\
希伯來書 3:13-14 & \hyperref[sec:FlUdBHop01k]{2023-07-30 午堂《信仰的初心》希伯來書 3\_13-14 講員:陳韋安牧師} & 2023-08-09 & \href{https://youtube.com/watch?v=FlUdBHop01k}{\texttt{FlUdBHop01k}} \\
使徒行傳 17:16-31 & \hyperref[sec:RptdKZSozLI]{2023-08-12 晚堂《進入人群的福音》使徒行傳 17\_16-31 講員:陳麗蟬牧師} & 2023-08-15 & \href{https://youtube.com/watch?v=RptdKZSozLI}{\texttt{RptdKZSozLI}} \\
約翰一書彼得前書 1:1-2 & \hyperref[sec:kFsvDzgp4nI]{2023-08-13 加恩樓《使人蒙福的「道」》約翰一書 1\_1-2,彼得前書 1\_23-25 講員:樓恩德牧師} & 2023-08-15 & \href{https://youtube.com/watch?v=kFsvDzgp4nI}{\texttt{kFsvDzgp4nI}} \\
哥林多後書 4:1-18 & \hyperref[sec:FIbxv36OMMk]{2023-08-20 午堂《傳道者的心》哥林多後書 4\_1-18 講員:鄧泰康傳道} & 2023-08-22 & \href{https://youtube.com/watch?v=FIbxv36OMMk}{\texttt{FIbxv36OMMk}} \\
以賽亞書 58:3-7 & \hyperref[sec:cMmv_XzlkAc]{2023-08-20 早堂《禁食的意義》以賽亞書 58\_3-7 講員:陳韋安牧師} & 2023-08-22 & \href{https://youtube.com/watch?v=cMmv_XzlkAc}{\texttt{cMmv\_XzlkAc}} \\
路加福音 10:25-37 & \hyperref[sec:TEOKFaxZL4g]{2023-08-20 加恩樓《愛鄰如己．照樣做吧》路加福音 10\_25-37 講員:陳麗蟬牧師} & 2023-08-23 & \href{https://youtube.com/watch?v=TEOKFaxZL4g}{\texttt{TEOKFaxZL4g}} \\
約伯記 28:1-28 & \hyperref[sec:zrsNFCgx1bQ]{2022-05-22 午堂《明智藏哪裡》約伯記 28\_1-28 講員:周曉暉牧師} & 2023-08-25 & \href{https://youtube.com/watch?v=zrsNFCgx1bQ}{\texttt{zrsNFCgx1bQ}} \\
出埃及記 23:1-13 & \hyperref[sec:lKgHJObwVL4]{2022-12-18 加恩樓《爭議中持守公正與良善》出埃及記 23\_1-13 講員:周曉暉牧師} & 2023-08-25 & \href{https://youtube.com/watch?v=lKgHJObwVL4}{\texttt{lKgHJObwVL4}} \\
以斯帖記  & \hyperref[sec:iRhu93YiYcU]{2023-08-24 周四晚查經班《被擄歸回歷史研讀(22)以斯帖記》 講員:曾裔貴牧師} & 2023-08-26 & \href{https://youtube.com/watch?v=iRhu93YiYcU}{\texttt{iRhu93YiYcU}} \\
彼得前書 5:1-11 & \hyperref[sec:fSxNb3R9eIk]{2023-08-27 加恩樓《父慈子孝的神家》彼得前書 5\_1-11 講員:鄧泰康傳道} & 2023-08-29 & \href{https://youtube.com/watch?v=fSxNb3R9eIk}{\texttt{fSxNb3R9eIk}} \\
哥林多前書 16:15-16 & \hyperref[sec:xZ5COJOEvOY]{2023-08-27 午堂《為主同工,同勞的人》哥林多前書 16\_15-16 講員:曾裔貴牧師} & 2023-08-29 & \href{https://youtube.com/watch?v=xZ5COJOEvOY}{\texttt{xZ5COJOEvOY}} \\
彼得前書 5:1-5 & \hyperref[sec:fjBMO0wHlN4]{2023-08-27 早堂《彼此順服的教會》彼得前書 5\_1-5 講員:盧思源姑娘} & 2023-08-29 & \href{https://youtube.com/watch?v=fjBMO0wHlN4}{\texttt{fjBMO0wHlN4}} \\
以斯帖記  & \hyperref[sec:9XvZPOyWZdE]{2023-08-31 周四晚查經班《被擄歸回歷史研讀(23)以斯帖記》 講員:曾裔貴牧師} & 2023-09-05 & \href{https://youtube.com/watch?v=9XvZPOyWZdE}{\texttt{9XvZPOyWZdE}} \\
提摩太後書 4:1-8 & \hyperref[sec:ay_HsjJTAHo]{2023-09-03 加恩樓《保羅最後的勸勉》提摩太後書 4\_1-8 講員:鄧泰康傳道} & 2023-09-05 & \href{https://youtube.com/watch?v=ay-HsjJTAHo}{\texttt{ay-HsjJTAHo}} \\
彼得後書 1:3-8 & \hyperref[sec:Bres3qEZ6Yk]{2023-09-03 午堂《在信仰上加倍努力》彼得後書 1\_3-8 講員:曾裔貴牧師} & 2023-09-05 & \href{https://youtube.com/watch?v=Bres3qEZ6Yk}{\texttt{Bres3qEZ6Yk}} \\
撒母耳記下 23:13-19 & \hyperref[sec:naPXdHer1J8]{2023-09-10 加恩樓《大衛與三個沒有名字的勇士》撒母耳記下 23\_13-19 講員:陳韋安牧師} & 2023-09-12 & \href{https://youtube.com/watch?v=naPXdHer1J8}{\texttt{naPXdHer1J8}} \\
提摩太後書 4:1-8 & \hyperref[sec:TXT2__xkqM8]{2023-09-10 早堂《保羅最後的勸勉》提摩太後書 4\_1-8 講員:鄧泰康傳道} & 2023-09-12 & \href{https://youtube.com/watch?v=TXT2_-xkqM8}{\texttt{TXT2\_-xkqM8}} \\
詩篇 8:1-9 & \hyperref[sec:LQDQ7CliZ7U]{2023-09-09 晚堂《神你是誰?而我又算甚麼?》詩篇 8\_1-9 講員:張美薇牧師} & 2023-09-13 & \href{https://youtube.com/watch?v=LQDQ7CliZ7U}{\texttt{LQDQ7CliZ7U}} \\
提摩太後書路加福音 9:57-62 & \hyperref[sec:eRch8741F34]{2023-09-10 午堂《若要跟從主》路加福音 9\_57-62,提摩太後書 1\_1-12 講員:郭熹瑜宣教士} & 2023-09-13 & \href{https://youtube.com/watch?v=eRch8741F34}{\texttt{eRch8741F34}} \\
撒迦利亞書  & \hyperref[sec:Kpbv61GAjrM]{2023-09-14 周四晚查經班《歸回先知書研讀(1)撒迦利亞書》 講員:曾裔貴牧師} & 2023-09-16 & \href{https://youtube.com/watch?v=Kpbv61GAjrM}{\texttt{Kpbv61GAjrM}} \\
詩篇 128:1-6 & \hyperref[sec:BaO8MW3rJk0]{2023-09-16 晚堂《福氣:敬畏主+遵主路》詩篇 128\_1-6 講員:陳麗蟬牧師} & 2023-09-19 & \href{https://youtube.com/watch?v=BaO8MW3rJk0}{\texttt{BaO8MW3rJk0}} \\
約翰福音 14:1-14 & \hyperref[sec:Ah5GT26Za2E]{2023-09-17 早堂《你使我堅壯》約翰福音 14\_1-14 講員:蕭建輝牧師} & 2023-09-19 & \href{https://youtube.com/watch?v=Ah5GT26Za2E}{\texttt{Ah5GT26Za2E}} \\
路加福音 11:5-13 & \hyperref[sec:Hdjnt297x1A]{2023-09-17 加恩樓《好吧!天父!我信靠你!》路加福音 11\_5-13 講員:蕭建輝牧師} & 2023-09-20 & \href{https://youtube.com/watch?v=Hdjnt297x1A}{\texttt{Hdjnt297x1A}} \\
以賽亞書 42:1-9 & \hyperref[sec:tDVwCjoCqSI]{2023-09-17 午堂《以賽亞先知的生命質素與使命》以賽亞書 42\_1-9 講員:郭鴻標牧師} & 2023-09-20 & \href{https://youtube.com/watch?v=tDVwCjoCqSI}{\texttt{tDVwCjoCqSI}} \\
撒迦利亞書  & \hyperref[sec:_MeYYEhxYZE]{2023-09-21 周四晚查經班《歸回先知書研讀(2)撒迦利亞書》 講員:曾裔貴牧師} & 2023-09-23 & \href{https://youtube.com/watch?v=-MeYYEhxYZE}{\texttt{-MeYYEhxYZE}} \\
使徒行傳 21:17-26 & \hyperref[sec:B_Q_tzD4YvY]{2023-09-24 加恩樓《入會四則》使徒行傳 21\_17-26 講員:陳麗蟬牧師} & 2023-09-26 & \href{https://youtube.com/watch?v=B-Q-tzD4YvY}{\texttt{B-Q-tzD4YvY}} \\
哥林多後書 5:17-18 & \hyperref[sec:CTfZJ96vxzI]{2023-09-24 早堂《重生》哥林多後書 5\_17-18 講員:黃頌言姑娘} & 2023-09-26 & \href{https://youtube.com/watch?v=CTfZJ96vxzI}{\texttt{CTfZJ96vxzI}} \\
希伯來書 4:1-11 & \hyperref[sec:QkQu6jc1ZPg]{2023-09-24 午堂《兩個安息》希伯來書 4\_1-11 講員:楊慶球牧師} & 2023-09-27 & \href{https://youtube.com/watch?v=QkQu6jc1ZPg}{\texttt{QkQu6jc1ZPg}} \\
撒迦利亞書  & \hyperref[sec:t6pbod4Qek8]{2023-09-28 周四晚查經班《歸回先知書研讀(3)撒迦利亞書》 講員:曾裔貴牧師} & 2023-09-29 & \href{https://youtube.com/watch?v=t6pbod4Qek8}{\texttt{t6pbod4Qek8}} \\
彌迦書 6:6-8 & \hyperref[sec:ZbYxDMl2Zbg]{2023-10-01 加恩樓《選擇站在上帝那邊》彌迦書 6\_6-8 講員:鄧泰康傳道} & 2023-10-03 & \href{https://youtube.com/watch?v=ZbYxDMl2Zbg}{\texttt{ZbYxDMl2Zbg}} \\
以弗所書 4:11-16 & \hyperref[sec:EvC3MODORZI]{2023-10-01 早堂《天降恩賜齊共享》以弗所書 4\_11-16 講員:陳麗蟬牧師} & 2023-10-03 & \href{https://youtube.com/watch?v=EvC3MODORZI}{\texttt{EvC3MODORZI}} \\
馬太福音 9:35-10:4 & \hyperref[sec:x5xyTkGVbQ0]{2023-09-30 晚堂《流散的羊群》馬太福音 9\_35-10\_4 講員:褚永華牧師} & 2023-10-04 & \href{https://youtube.com/watch?v=x5xyTkGVbQ0}{\texttt{x5xyTkGVbQ0}} \\
撒母耳記下 23:13-19 & \hyperref[sec:JHYqHgNM3GQ]{2023-10-01 午堂《大衛與三個沒有名字的勇士》撒母耳記下 23\_13-19 講員:陳韋安牧師} & 2023-10-04 & \href{https://youtube.com/watch?v=JHYqHgNM3GQ}{\texttt{JHYqHgNM3GQ}} \\
撒迦利亞書  & \hyperref[sec:CXJyuYm4_pE]{2023-10-05 周四晚查經班《歸回先知書研讀(4)撒迦利亞書》 講員:曾裔貴牧師} & 2023-10-06 & \href{https://youtube.com/watch?v=CXJyuYm4-pE}{\texttt{CXJyuYm4-pE}} \\
馬可福音 1:40-44 & \hyperref[sec:G0FSq209Kac]{2023-10-07 晚堂《醫治為要得潔淨》馬可福音 1\_40-44 講員:黃國維院長} & 2023-10-11 & \href{https://youtube.com/watch?v=G0FSq209Kac}{\texttt{G0FSq209Kac}} \\
約翰福音 4:1-42 & \hyperref[sec:iO1qZX8Cmq4]{2023-10-08 早堂《生命轉化之旅》約翰福音 4\_1-42 講員:陳淑娟牧師} & 2023-10-12 & \href{https://youtube.com/watch?v=iO1qZX8Cmq4}{\texttt{iO1qZX8Cmq4}} \\
撒迦利亞書  & \hyperref[sec:wwGm5m9DNT8]{2023-10-12 周四晚查經班《歸回先知書研讀(5)撒迦利亞書》 講員:曾裔貴牧師} & 2023-10-14 & \href{https://youtube.com/watch?v=wwGm5m9DNT8}{\texttt{wwGm5m9DNT8}} \\
馬可福音 5:21-43 & \hyperref[sec:HWMu4dld7_M]{2023-10-15 加恩樓《這是一個怎樣的信心》馬可福音 5\_21-43 講員:老耀雄牧師} & 2023-10-18 & \href{https://youtube.com/watch?v=HWMu4dld7_M}{\texttt{HWMu4dld7\_M}} \\
以弗所書 1:17-23-4:7-16 & \hyperref[sec:cM9RUvSdlG8]{2023-10-15 早堂《尊基督為元首的門徒相交》以弗所書 1\_17-23,4\_7-16 講員:曾裔貴牧師} & 2023-10-19 & \href{https://youtube.com/watch?v=cM9RUvSdlG8}{\texttt{cM9RUvSdlG8}} \\
撒迦利亞書  & \hyperref[sec:c_6waYrNYvI]{2023-10-19 周四晚查經班《歸回先知書研讀(6)撒迦利亞書》 講員:曾裔貴牧師} & 2023-10-21 & \href{https://youtube.com/watch?v=c_6waYrNYvI}{\texttt{c\_6waYrNYvI}} \\
馬太福音 10:34-42 & \hyperref[sec:qZ9s1IR331E]{2023-10-21 晚堂《作主門徒》馬太福音 10\_34-42 講員:陳麗蟬牧師} & 2023-10-24 & \href{https://youtube.com/watch?v=qZ9s1IR331E}{\texttt{qZ9s1IR331E}} \\
加拉太書 6:9-10 & \hyperref[sec:ZmEDOfaqivc]{2023-10-22 加恩樓《及時行善》加拉太書 6\_9-10 講員:何智雄牧師} & 2023-10-24 & \href{https://youtube.com/watch?v=ZmEDOfaqivc}{\texttt{ZmEDOfaqivc}} \\
約拿書 1:1-17-4:1-3 & \hyperref[sec:LM01sAVsGrw]{2023-10-22 午堂《經歷有恩典憐憫的神》約拿書 1\_1-17,4\_1-3 講員:陳燕雯姑娘} & 2023-10-25 & \href{https://youtube.com/watch?v=LM01sAVsGrw}{\texttt{LM01sAVsGrw}} \\
以弗所書馬可福音 4:14-16 & \hyperref[sec:xr_m3_qAqFs]{2023-10-28 晚堂《在愛中成長》以弗所書 4\_14-16,馬可福音 12\_29-30 講員:楊明麗姑娘} & 2023-10-31 & \href{https://youtube.com/watch?v=xr_m3-qAqFs}{\texttt{xr\_m3-qAqFs}} \\
腓利門書 1:4-20 & \hyperref[sec:ROGX9cAg8Cg]{2023-10-29  42周年堂慶聯合崇拜《友伴．前行》腓利門書 1\_4-20 講員:周曉暉牧師} & 2023-10-31 & \href{https://youtube.com/watch?v=ROGX9cAg8Cg}{\texttt{ROGX9cAg8Cg}} \\
撒迦利亞書  & \hyperref[sec:j9CuOv5EhDA]{2023-11-02 周四晚查經班《歸回先知書研讀(7)撒迦利亞書》 講員:曾裔貴牧師} & 2023-11-03 & \href{https://youtube.com/watch?v=j9CuOv5EhDA}{\texttt{j9CuOv5EhDA}} \\
詩篇 116:1-19 & \hyperref[sec:URxw02niNWY]{2023-11-04 晚堂《以感恩為祭》詩篇 116\_1-19 講員:曾裔貴牧師} & 2023-11-07 & \href{https://youtube.com/watch?v=URxw02niNWY}{\texttt{URxw02niNWY}} \\
哥林多後書使徒行傳 9:30-13:2-27:22-24 & \hyperref[sec:EpQdyi4IyoI]{2023-11-05 午堂《不同階段有不同學習:神永遠有奇妙恩典》使徒行傳 9\_30,13\_2,27\_22-24,哥林多後書 12\_9 講員:蔡少琪院長} & 2023-11-08 & \href{https://youtube.com/watch?v=EpQdyi4IyoI}{\texttt{EpQdyi4IyoI}} \\
撒迦利亞書  & \hyperref[sec:z09lg9RcA3o]{2023-11-09 周四晚查經班《歸回先知書研讀(8)撒迦利亞書》 講員:曾裔貴牧師} & 2023-11-11 & \href{https://youtube.com/watch?v=z09lg9RcA3o}{\texttt{z09lg9RcA3o}} \\
馬可福音 16:9-20 & \hyperref[sec:Jspoigh_Nw0]{2023-11-12 午堂《傳揚復活的基督》馬可福音 16\_9-20 講員:曾裔貴牧師} & 2023-11-14 & \href{https://youtube.com/watch?v=Jspoigh-Nw0}{\texttt{Jspoigh-Nw0}} \\
哥林多前書 9:24-27 & \hyperref[sec:lMH9_3_EQe8]{2023-11-12 加恩樓《為主奔跑》哥林多前書 9\_24-27 講員:陳一華牧師} & 2023-11-15 & \href{https://youtube.com/watch?v=lMH9_3-EQe8}{\texttt{lMH9\_3-EQe8}} \\
詩篇 150:1-6 & \hyperref[sec:zCZXBv4_hT8]{2023-11-18 晚堂《當來讚美神》詩篇 150\_1-6 講員:張美薇牧師} & 2023-11-21 & \href{https://youtube.com/watch?v=zCZXBv4_hT8}{\texttt{zCZXBv4\_hT8}} \\
約書亞記 14:6-15 & \hyperref[sec:uKxjJHU8_7c]{2023-11-19 加恩樓《生命的投資》約書亞記 14\_6-15 講員:蔣文忠傳道} & 2023-11-21 & \href{https://youtube.com/watch?v=uKxjJHU8-7c}{\texttt{uKxjJHU8-7c}} \\
以賽亞書 49:1-7 & \hyperref[sec:z_EDsyn25oM]{2023-11-19 午堂《成為萬邦之光》以賽亞書 49\_1-7 講員:陳麗蟬牧師} & 2023-11-22 & \href{https://youtube.com/watch?v=z-EDsyn25oM}{\texttt{z-EDsyn25oM}} \\
歌羅西書 3:15-17 & \hyperref[sec:Th6Drt1kZ1s]{2023-11-25 晚堂《心存感恩》歌羅西書 3\_15-17 講員:盧思源姑娘} & 2023-11-28 & \href{https://youtube.com/watch?v=Th6Drt1kZ1s}{\texttt{Th6Drt1kZ1s}} \\
提摩太後書 1:1-5 & \hyperref[sec:Q0yYJrC6Zoc]{2023-11-26 加恩樓《代代相傳的信仰》提摩太後書 1\_1-5 講員:黃頌言姑娘} & 2023-11-28 & \href{https://youtube.com/watch?v=Q0yYJrC6Zoc}{\texttt{Q0yYJrC6Zoc}} \\
使徒行傳 11:19-30 & \hyperref[sec:nQU0VbnDk9Y]{2023-11-26 午堂《要作榮耀主的教會》使徒行傳 11\_19-30 講員:趙家碩牧師} & 2023-11-28 & \href{https://youtube.com/watch?v=nQU0VbnDk9Y}{\texttt{nQU0VbnDk9Y}} \\
馬太福音 2:1-12 & \hyperref[sec:gy4cK2fsHdc]{2023-12-02 晚堂《聖誕禮物》馬太福音 2\_1-12 講員:楊明麗姑娘} & 2023-12-06 & \href{https://youtube.com/watch?v=gy4cK2fsHdc}{\texttt{gy4cK2fsHdc}} \\
使徒行傳 5:12-32 & \hyperref[sec:7XxklgVuie4]{2023-12-03 午堂《傳講生命之道》使徒行傳 5\_12-32 講員:曾裔貴牧師} & 2023-12-06 & \href{https://youtube.com/watch?v=7XxklgVuie4}{\texttt{7XxklgVuie4}} \\
撒迦利亞書  & \hyperref[sec:Qx_35UZ1A44]{2023-12-07 周四晚查經班《歸回先知書研讀(9)撒迦利亞書》 講員:曾裔貴牧師} & 2023-12-08 & \href{https://youtube.com/watch?v=Qx-35UZ1A44}{\texttt{Qx-35UZ1A44}} \\
路加福音 3:1-9 & \hyperref[sec:_n0nRPZl6sE]{2023-12-09 晚堂《君王來臨的平安喜樂》路加福音 3\_1-9 講員:曾裔貴牧師} & 2023-12-12 & \href{https://youtube.com/watch?v=-n0nRPZl6sE}{\texttt{-n0nRPZl6sE}} \\
馬可福音 1:1-8 & \hyperref[sec:KJG_idsYK_w]{2023-12-10 加恩樓《修直祂的道》馬可福音 1\_1-8 講員:楊明麗姑娘} & 2023-12-12 & \href{https://youtube.com/watch?v=KJG_idsYK-w}{\texttt{KJG\_idsYK-w}} \\
以賽亞書 11:1-9 & \hyperref[sec:m2hnqzcprEM]{2023-12-10 午堂《共融的新天新地》以賽亞書 11\_1-9 講員:陳麗蟬牧師} & 2023-12-12 & \href{https://youtube.com/watch?v=m2hnqzcprEM}{\texttt{m2hnqzcprEM}} \\
約書亞記 14:6-15 & \hyperref[sec:jzt3Ohpru70]{2023-12-16 晚堂《生命的投資》約書亞記 14\_6-15 講員:蔣文忠傳道} & 2023-12-19 & \href{https://youtube.com/watch?v=jzt3Ohpru70}{\texttt{jzt3Ohpru70}} \\
路加福音 2:25-35 & \hyperref[sec:oC6cD_CKYi8]{2023-12-17 加恩樓《必看見所立的基督》路加福音 2\_25-35 講員:張漢強牧師} & 2023-12-19 & \href{https://youtube.com/watch?v=oC6cD-CKYi8}{\texttt{oC6cD-CKYi8}} \\
約伯記 42:1-6 & \hyperref[sec:EJ51POSz_gM]{2023-12-17 午堂《從前風聞有你,現在親眼看見你》約伯記 42\_1-6 講員:鄧泰康傳道} & 2023-12-19 & \href{https://youtube.com/watch?v=EJ51POSz_gM}{\texttt{EJ51POSz\_gM}} \\
詩篇 90:1-17 & \hyperref[sec:JmKumENSOqs]{2023-12-17 早堂《赤子之心》詩篇 90\_1-17 講員:陳韋安牧師} & 2023-12-19 & \href{https://youtube.com/watch?v=JmKumENSOqs}{\texttt{JmKumENSOqs}} \\
撒迦利亞書  & \hyperref[sec:I2kkhTmf8Lk]{2023-12-21 周四晚查經班《歸回先知書研讀(10)撒迦利亞書》 講員:曾裔貴牧師} & 2023-12-23 & \href{https://youtube.com/watch?v=I2kkhTmf8Lk}{\texttt{I2kkhTmf8Lk}} \\
路加福音 2:8-20 & \hyperref[sec:pZkchuW18Ro]{2023-12-24 加恩樓三語聯合崇拜《Good News of Great Joy》路加福音 2\_8-20 講員:Rev. Carol Ilidan} & 2023-12-27 & \href{https://youtube.com/watch?v=pZkchuW18Ro}{\texttt{pZkchuW18Ro}} \\
  & \hyperref[sec:7Vdee0ksMqc]{2023-12-24 平安夜佈道會《不一樣的聖誕》 講員:李炳光牧師} & 2023-12-27 & \href{https://youtube.com/watch?v=7Vdee0ksMqc}{\texttt{7Vdee0ksMqc}} \\
以賽亞書 7:10-17 & \hyperref[sec:nEkocKnhlmk]{2023-12-24 聯合崇拜《以馬內利的再思》以賽亞書 7\_10-17 講員:盧思源姑娘} & 2023-12-27 & \href{https://youtube.com/watch?v=nEkocKnhlmk}{\texttt{nEkocKnhlmk}} \\
  & \hyperref[sec:bpii0IGPrEA]{2023-12-25 聖誕節崇拜《諸天萬物同歡唱》 講員:陳麗蟬牧師} & 2023-12-27 & \href{https://youtube.com/watch?v=bpii0IGPrEA}{\texttt{bpii0IGPrEA}} \\
撒迦利亞書  & \hyperref[sec:MkLoNpPCJk4]{2023-12-28 周四晚查經班《歸回先知書研讀(11)撒迦利亞書》 講員:曾裔貴牧師} & 2023-12-29 & \href{https://youtube.com/watch?v=MkLoNpPCJk4}{\texttt{MkLoNpPCJk4}} \\
哥林多前書 9:24-27 & \hyperref[sec:lIj_mrEVPGg]{2023-12-30 晚堂《為主奔跑》哥林多前書 9\_24-27 講員:陳一華牧師} & 2024-01-03 & \href{https://youtube.com/watch?v=lIj-mrEVPGg}{\texttt{lIj-mrEVPGg}} \\
路加福音 2:1-7 & \hyperref[sec:c3xZAeZuwD0]{2023-12-31 加恩樓《在你心裡有空處給主嗎?》路加福音 2\_1-7 講員:謝靄薇牧師} & 2024-01-03 & \href{https://youtube.com/watch?v=c3xZAeZuwD0}{\texttt{c3xZAeZuwD0}} \\
以賽亞書 61:1 & \hyperref[sec:p_xp8W8_F_Q]{2023-12-31 午堂《禧年執行者》以賽亞書 61\_1 講員:楊明麗姑娘} & 2024-01-03 & \href{https://youtube.com/watch?v=p_xp8W8-F_Q}{\texttt{p\_xp8W8-F\_Q}} \\
撒母耳記上 15:17-23 & \hyperref[sec:R1nryPpcggU]{2023-12-31 早堂《作主聽命的兒女》撒母耳記上 15\_17-23 講員:曾敏姑娘} & 2024-01-03 & \href{https://youtube.com/watch?v=R1nryPpcggU}{\texttt{R1nryPpcggU}} \\
瑪拉基書  & \hyperref[sec:qkHV9uoLu6Y]{2024-01-04 周四晚查經班《瑪拉基書(1)》 講員:曾裔貴牧師} & 2024-01-05 & \href{https://youtube.com/watch?v=qkHV9uoLu6Y}{\texttt{qkHV9uoLu6Y}} \\
箴言 29:18 & \hyperref[sec:SWoA7YUZj_4]{2024-01-07 午堂《2024錦宣家．齊躍動》箴言 29\_18 講員:陳麗蟬牧師} & 2024-01-09 & \href{https://youtube.com/watch?v=SWoA7YUZj_4}{\texttt{SWoA7YUZj\_4}} \\
馬太福音 2:1-12 & \hyperref[sec:Y_5AXHxN_yY]{2024-01-07 早堂《星光引導》馬太福音 2\_1-12 講員:曾浩賢院長} & 2024-01-09 & \href{https://youtube.com/watch?v=Y-5AXHxN_yY}{\texttt{Y-5AXHxN\_yY}} \\
希伯來書 11:8-9 & \hyperref[sec:2um86Ip2zrs]{2024-01-14 午堂《亞伯拉罕信心的成長》希伯來書 11\_8-9 講員:楊明麗姑娘} & 2024-01-16 & \href{https://youtube.com/watch?v=2um86Ip2zrs}{\texttt{2um86Ip2zrs}} \\
使徒行傳 26:1-23 & \hyperref[sec:Rlc0zB_br4Y]{2024-01-14 早堂《沒有違背那從天上來的異象》使徒行傳 26\_1-23 講員:李富成牧師} & 2024-01-16 & \href{https://youtube.com/watch?v=Rlc0zB-br4Y}{\texttt{Rlc0zB-br4Y}} \\
瑪拉基書  & \hyperref[sec:dR8L9GrF4hQ]{2024-01-18 周四晚查經班《瑪拉基書(2)》 講員:曾裔貴牧師} & 2024-01-19 & \href{https://youtube.com/watch?v=dR8L9GrF4hQ}{\texttt{dR8L9GrF4hQ}} \\
創世記 16:1-13 & \hyperref[sec:_CjhJKqUcA8]{2024-01-20 晚堂《沒有被遺忘的女兒》創世記 16\_1-13 講員:曾裔貴牧師} & 2024-01-23 & \href{https://youtube.com/watch?v=-CjhJKqUcA8}{\texttt{-CjhJKqUcA8}} \\
創世記 11:1-9 & \hyperref[sec:DTr4fwkGAak]{2024-01-21 午堂《從神的創造心意再思教會的大使命》創世記 11\_1-9 講員:楊穎兒姑娘} & 2024-01-23 & \href{https://youtube.com/watch?v=DTr4fwkGAak}{\texttt{DTr4fwkGAak}} \\
創世記 16:8-14 & \hyperref[sec:pNQz6hqXFi8]{2024-01-21 早堂《憐憫人,讓人看見的神》創世記 16\_8-14 講員:鄧泰康傳道} & 2024-01-23 & \href{https://youtube.com/watch?v=pNQz6hqXFi8}{\texttt{pNQz6hqXFi8}} \\
瑪拉基書  & \hyperref[sec:jLSZtQ5_sq4]{2024-01-25 周四晚查經班《瑪拉基書(3)》 講員:曾裔貴牧師} & 2024-01-27 & \href{https://youtube.com/watch?v=jLSZtQ5_sq4}{\texttt{jLSZtQ5\_sq4}} \\
約翰福音 12:26 & \hyperref[sec:c0kCL7cehvg]{2024-01-28 午堂《蒙福的事奉》約翰福音 12\_26 講員:陳一華牧師} & 2024-01-31 & \href{https://youtube.com/watch?v=c0kCL7cehvg}{\texttt{c0kCL7cehvg}} \\
路加福音 23:40 & \hyperref[sec:FhL_HPKlT9E]{2024-01-28 早堂《談敬畏》路加福音 23\_40 講員:陳韋安牧師} & 2024-01-31 & \href{https://youtube.com/watch?v=FhL-HPKlT9E}{\texttt{FhL-HPKlT9E}} \\
瑪拉基書  & \hyperref[sec:Srg_1OhzOG8]{2024-02-01 周四晚查經班《瑪拉基書(4)》 講員:曾裔貴牧師} & 2024-02-02 & \href{https://youtube.com/watch?v=Srg_1OhzOG8}{\texttt{Srg\_1OhzOG8}} \\
創世記 26:12-33 & \hyperref[sec:vJ5BK_GiIWM]{2024-02-04 早堂《以撒的神》創世記 26\_12-33 講員:黃頌言姑娘} & 2024-02-06 & \href{https://youtube.com/watch?v=vJ5BK_GiIWM}{\texttt{vJ5BK\_GiIWM}} \\
創世記 26:1-23 & \hyperref[sec:HWk9pevFs74]{2024-02-04 午堂《豐盛生命:依靠順服》創世記 26\_1-23 講員:陳麗蟬牧師} & 2024-02-07 & \href{https://youtube.com/watch?v=HWk9pevFs74}{\texttt{HWk9pevFs74}} \\
瑪拉基書  & \hyperref[sec:ny6z0RWTL8g]{2024-02-08 周四晚查經班《瑪拉基書(5)》 講員:曾裔貴牧師} & 2024-02-09 & \href{https://youtube.com/watch?v=ny6z0RWTL8g}{\texttt{ny6z0RWTL8g}} \\
創世記 26:1-23 & \hyperref[sec:OWpYTLtxpsA]{2024-02-03 晚堂《豐盛生命:依靠順服》創世記 26\_1-23 講員:陳麗蟬牧師} & 2024-02-14 & \href{https://youtube.com/watch?v=OWpYTLtxpsA}{\texttt{OWpYTLtxpsA}} \\
但以理書 6:1-3 & \hyperref[sec:45mO3Udlc8E]{2024-02-11 聯合崇拜《龍年生肖道》但以理書 6\_1-3 講員:陳韋安牧師} & 2024-02-14 & \href{https://youtube.com/watch?v=45mO3Udlc8E}{\texttt{45mO3Udlc8E}} \\
馬太福音 5:3-10 & \hyperref[sec:QYUt9NIhVQs]{2024-02-10 新春崇拜《福從天降》馬太福音 5\_3-10 講員:陳一華牧師} & 2024-02-15 & \href{https://youtube.com/watch?v=QYUt9NIhVQs}{\texttt{QYUt9NIhVQs}} \\
瑪拉基書  & \hyperref[sec:oCSQzasymjk]{2024-02-15 周四晚查經班《瑪拉基書(6)》 講員:曾裔貴牧師} & 2024-02-16 & \href{https://youtube.com/watch?v=oCSQzasymjk}{\texttt{oCSQzasymjk}} \\
列王記上 19:1-18 & \hyperref[sec:C6eEbdGd3Ao]{2024-02-18 午堂《當你以為走到絕境時》列王記上 19\_1-18 講員:鄧泰康傳道} & 2024-02-19 & \href{https://youtube.com/watch?v=C6eEbdGd3Ao}{\texttt{C6eEbdGd3Ao}} \\
以弗所書 4:20-23 & \hyperref[sec:myBJVEoimI0]{2024-02-18 早堂《如果...可以重來》以弗所書 4\_20-23 講員:黃仲安傳道} & 2024-02-19 & \href{https://youtube.com/watch?v=myBJVEoimI0}{\texttt{myBJVEoimI0}} \\
瑪拉基書  & \hyperref[sec:4SEW84zBGCU]{2024-02-22 周四晚查經班《瑪拉基書(7)》 講員:曾裔貴牧師} & 2024-02-23 & \href{https://youtube.com/watch?v=4SEW84zBGCU}{\texttt{4SEW84zBGCU}} \\
馬可福音 8:31-38 & \hyperref[sec:dYh7F8XFvZo]{2024-02-24 晚堂《作門徒的代價》馬可福音 8\_31-38 講員:褚永華牧師} & 2024-02-27 & \href{https://youtube.com/watch?v=dYh7F8XFvZo}{\texttt{dYh7F8XFvZo}} \\
馬太福音 3:1-12 & \hyperref[sec:A8lXH1cPVBk]{2024-02-25 午堂《國度,權柄,榮耀》馬太福音 3\_1-12 講員:李思敬牧師} & 2024-02-27 & \href{https://youtube.com/watch?v=A8lXH1cPVBk}{\texttt{A8lXH1cPVBk}} \\
創世記 25:27-34 & \hyperref[sec:__QTchIrjDM]{2024-02-25 早堂《當你以為是理所當然》創世記 25\_27-34 講員:鄧泰康傳道} & 2024-02-27 & \href{https://youtube.com/watch?v=_-QTchIrjDM}{\texttt{\_-QTchIrjDM}} \\
瑪拉基書  & \hyperref[sec:exEZqq1nt80]{2024-02-29 周四晚查經班《瑪拉基書(8)》 講員:曾裔貴牧師} & 2024-03-02 & \href{https://youtube.com/watch?v=exEZqq1nt80}{\texttt{exEZqq1nt80}} \\
創世記 32:1-12-22-32 & \hyperref[sec:oyj7tRcaGj0]{2024-03-03 早堂《雅各的祈禱》創世記 32\_1-12,22-32,33\_1-4 講員:謝靄薇牧師} & 2024-03-04 & \href{https://youtube.com/watch?v=oyj7tRcaGj0}{\texttt{oyj7tRcaGj0}} \\
創世記 32:22-32 & \hyperref[sec:a1w0x_rlHHg]{2024-03-03 午堂《得著與失去》創世記 32\_22-32 講員:陳麗蟬牧師} & 2024-03-06 & \href{https://youtube.com/watch?v=a1w0x_rlHHg}{\texttt{a1w0x\_rlHHg}} \\
瑪拉基書  & \hyperref[sec:2A_RiSlG6Bg]{2024-03-07 周四晚查經班《瑪拉基書(9)》 講員:曾裔貴牧師} & 2024-03-08 & \href{https://youtube.com/watch?v=2A-RiSlG6Bg}{\texttt{2A-RiSlG6Bg}} \\
加拉太書 5:13 & \hyperref[sec:AndPB0XPtoQ]{2024-03-10 午堂《不被捆綁的自由》加拉太書 5\_13 講員:陳韋安牧師} & 2024-03-12 & \href{https://youtube.com/watch?v=AndPB0XPtoQ}{\texttt{AndPB0XPtoQ}} \\
馬可福音 1:14-20 & \hyperref[sec:U06UnMRzdJ4]{2024-03-10 早堂《作主門徒》馬可福音 1\_14-20 講員:董家驊牧師} & 2024-03-12 & \href{https://youtube.com/watch?v=U06UnMRzdJ4}{\texttt{U06UnMRzdJ4}} \\
詩篇 139:1-13-18-23-24 & \hyperref[sec:pdca3AXE1WI]{2024-03-09 晚堂《在靈修和生活中經歷神》詩篇 139\_1,13-18,23-24 講員:廖炳堂牧師} & 2024-03-13 & \href{https://youtube.com/watch?v=pdca3AXE1WI}{\texttt{pdca3AXE1WI}} \\
創世記 48:15-16 & \hyperref[sec:RZv3JKPRqaw]{2024-03-16 晚堂《一生牧養我的神》創世記 48\_15-16 講員:盧思源姑娘} & 2024-03-19 & \href{https://youtube.com/watch?v=RZv3JKPRqaw}{\texttt{RZv3JKPRqaw}} \\
創世記 47:27-31 & \hyperref[sec:82RgV2iRB78]{2024-03-17 早堂《雅各的敬拜人生》創世記 47\_27-31 講員:曾裔貴牧師} & 2024-03-21 & \href{https://youtube.com/watch?v=82RgV2iRB78}{\texttt{82RgV2iRB78}} \\
帖撒羅尼迦前書  & \hyperref[sec:ih_tESQH0P8]{2024-03-21 周四晚查經班《帖撒羅尼迦前書(1)》 講員:曾裔貴牧師} & 2024-03-22 & \href{https://youtube.com/watch?v=ih_tESQH0P8}{\texttt{ih\_tESQH0P8}} \\
約翰福音 12:12-19 & \hyperref[sec:u5__eddKCBg]{2024-03-23 晚堂《耶穌榮耀進城》約翰福音 12\_12-19 講員:黃頌言姑娘} & 2024-03-26 & \href{https://youtube.com/watch?v=u5--eddKCBg}{\texttt{u5--eddKCBg}} \\
約翰福音 12:1-26 & \hyperref[sec:94pJmtJo_JE]{2024-03-24 早堂《先生,我們想見耶穌》約翰福音 12\_1-26 講員:李炳光牧師} & 2024-03-26 & \href{https://youtube.com/watch?v=94pJmtJo-JE}{\texttt{94pJmtJo-JE}} \\
士師記 16:4-31 & \hyperref[sec:NKxHUXgOsyY]{2024-03-24 午堂《勝過致命誘惑的秘訣》士師記 16\_4-31 講員:林誠信牧師} & 2024-03-28 & \href{https://youtube.com/watch?v=NKxHUXgOsyY}{\texttt{NKxHUXgOsyY}} \\
帖撒羅尼迦前書  & \hyperref[sec:xRsKG7u5A0E]{2024-03-28 周四晚查經班《帖撒羅尼迦前書(2)》 講員:曾裔貴牧師} & 2024-03-30 & \href{https://youtube.com/watch?v=xRsKG7u5A0E}{\texttt{xRsKG7u5A0E}} \\
羅馬書 6:1-14 & \hyperref[sec:Su045rsIjTY]{2024-03-29 受苦節崇拜《耶穌為你死了,你呢?!》羅馬書 6\_1-14 講員:鄧泰康傳道} & 2024-04-02 & \href{https://youtube.com/watch?v=Su045rsIjTY}{\texttt{Su045rsIjTY}} \\
馬太福音 28:1-10 & \hyperref[sec:ez2giRGwvx8]{2024-03-31 復活節崇拜《復活:衝破生命的枷鎖》馬太福音 28\_1-10 講員:陳麗蟬牧師} & 2024-04-02 & \href{https://youtube.com/watch?v=ez2giRGwvx8}{\texttt{ez2giRGwvx8}} \\
創世記 30:14-16 & \hyperref[sec:_PqGg3XbEA0]{2024-04-06 晚堂《缺乏了愛的家庭–雅各,利亞,拉結》創世記 30\_14-16 講員:鄧泰康傳道} & 2024-04-09 & \href{https://youtube.com/watch?v=-PqGg3XbEA0}{\texttt{-PqGg3XbEA0}} \\
路加福音 24:36-43 & \hyperref[sec:TfOT_yj_ii8]{2024-04-07 早堂《轉化我們情緒困擾的耶穌》路加福音 24\_36-43 講員:區祥江博士} & 2024-04-10 & \href{https://youtube.com/watch?v=TfOT_yj-ii8}{\texttt{TfOT\_yj-ii8}} \\
帖撒羅尼迦前書  & \hyperref[sec:CWxLFOO62pc]{2024-04-11 周四晚查經班《帖撒羅尼迦前書(3)》 講員:曾裔貴牧師} & 2024-04-13 & \href{https://youtube.com/watch?v=CWxLFOO62pc}{\texttt{CWxLFOO62pc}} \\
路加福音 10:25-37 & \hyperref[sec:_2QFVuN84qA]{2024-04-13 晚堂《關愛鄰舍再思》路加福音 10\_25-37 講員:陳一華牧師} & 2024-04-16 & \href{https://youtube.com/watch?v=-2QFVuN84qA}{\texttt{-2QFVuN84qA}} \\
申命記 1:1-18 & \hyperref[sec:iqMY9wa3DY0]{2024-04-14 午堂《同行共拓主國度》申命記 1\_1-18 講員:戴浩輝牧師} & 2024-04-16 & \href{https://youtube.com/watch?v=iqMY9wa3DY0}{\texttt{iqMY9wa3DY0}} \\
創世記 37:1-36 & \hyperref[sec:y77Sfudracc]{2024-04-20 晚堂《主必追討》創世記 37\_1-36 講員:陳麗蟬牧師} & 2024-04-23 & \href{https://youtube.com/watch?v=y77Sfudracc}{\texttt{y77Sfudracc}} \\
創世記 37:2-4 & \hyperref[sec:kK4Sp5yFq94]{2024-04-21 午堂《約瑟的華服》創世記 37\_2-4;39\_1-6,19-23;41\_14,38-43 講員:楊明麗姑娘} & 2024-04-23 & \href{https://youtube.com/watch?v=kK4Sp5yFq94}{\texttt{kK4Sp5yFq94}} \\
帖撒羅尼迦前書  & \hyperref[sec:ZVwFpbgglfY]{2024-04-25 周四晚查經班《帖撒羅尼迦前書(4)》 講員:曾裔貴牧師} & 2024-04-28 & \href{https://youtube.com/watch?v=ZVwFpbgglfY}{\texttt{ZVwFpbgglfY}} \\
羅馬書 8:26-30 & \hyperref[sec:a_rRhGfuxUw]{2024-04-27 晚堂《再思天路歷程》羅馬書 8\_26-30  講員:曾裔貴牧師} & 2024-04-30 & \href{https://youtube.com/watch?v=a-rRhGfuxUw}{\texttt{a-rRhGfuxUw}} \\
腓立比書 2:17-30 & \hyperref[sec:w0vSwuU8gCk]{2024-04-28 早堂《同心事主的知己》腓立比書 2\_17-30  講員:馮耀榮牧師} & 2024-04-30 & \href{https://youtube.com/watch?v=w0vSwuU8gCk}{\texttt{w0vSwuU8gCk}} \\
撒迦利亞書 2:6-8 & \hyperref[sec:jRV45ettAd0]{2024-04-28 午堂《上帝眼中的瞳人》撒迦利亞書 2\_6-8  講員:陳韋安牧師} & 2024-05-01 & \href{https://youtube.com/watch?v=jRV45ettAd0}{\texttt{jRV45ettAd0}} \\
帖撒羅尼迦前書  & \hyperref[sec:sCNFnlmgxC0]{2024-05-02 周四晚查經班《帖撒羅尼迦前書(5)》 講員:曾裔貴牧師} & 2024-05-03 & \href{https://youtube.com/watch?v=sCNFnlmgxC0}{\texttt{sCNFnlmgxC0}} \\
創世記 39:1-6-19-23 & \hyperref[sec:U5CngR8W_hQ]{2024-05-04 晚堂《變幻人生主同行》創世記 39\_1-6,19-23  講員:陳麗蟬牧師} & 2024-05-07 & \href{https://youtube.com/watch?v=U5CngR8W-hQ}{\texttt{U5CngR8W-hQ}} \\
加拉太書 2:20 & \hyperref[sec:ltEd3gYtiPs]{2024-05-05 早堂《基督在那裡?》加拉太書 2\_20  講員:朱浩廉博士} & 2024-05-07 & \href{https://youtube.com/watch?v=ltEd3gYtiPs}{\texttt{ltEd3gYtiPs}} \\
馬太福音 11:2-19 & \hyperref[sec:1QRoYKXVPvQ]{2024-05-05 午堂《沒有最大,只有更大》馬太福音 11\_2-19  講員:張雲開老師} & 2024-05-08 & \href{https://youtube.com/watch?v=1QRoYKXVPvQ}{\texttt{1QRoYKXVPvQ}} \\
帖撒羅尼迦前書  & \hyperref[sec:VvxSAXf_Lww]{2024-05-09 周四晚查經班《帖撒羅尼迦前書(6)》 講員:曾裔貴牧師} & 2024-05-10 & \href{https://youtube.com/watch?v=VvxSAXf-Lww}{\texttt{VvxSAXf-Lww}} \\
創世記 3:20 & \hyperref[sec:BHhp4UIkEH0]{2024-05-11 晚堂《一個蒙恩的母親》創世記 3\_20  講員:楊穎兒姑娘} & 2024-05-14 & \href{https://youtube.com/watch?v=BHhp4UIkEH0}{\texttt{BHhp4UIkEH0}} \\
撒母耳記上 1:10-18-20 & \hyperref[sec:rgd8tUJeioY]{2024-05-12 午堂《母親的禱告》撒母耳記上 1\_10-18,20  講員:黃頌言姑娘} & 2024-05-14 & \href{https://youtube.com/watch?v=rgd8tUJeioY}{\texttt{rgd8tUJeioY}} \\
約翰福音 16:5-8 & \hyperref[sec:D8TmPX_lJJQ]{2024-05-18 晚堂《看啊!聖靈駕到!》約翰福音16\_5-8 講員:楊明麗姑娘} & 2024-05-21 & \href{https://youtube.com/watch?v=D8TmPX-lJJQ}{\texttt{D8TmPX-lJJQ}} \\
撒迦利亞書 4:1-10 & \hyperref[sec:AlakgRywpkk]{2024-05-19 午堂《靠主的靈方成事》撒迦利亞書4\_1-10 講員:陳麗蟬牧師} & 2024-05-22 & \href{https://youtube.com/watch?v=AlakgRywpkk}{\texttt{AlakgRywpkk}} \\
約翰福音 11:35 & \hyperref[sec:0Pem3SkLOJI]{2024-05-19 早堂《耶穌的眼淚》約翰福音11\_35 講員:陳韋安牧師} & 2024-05-23 & \href{https://youtube.com/watch?v=0Pem3SkLOJI}{\texttt{0Pem3SkLOJI}} \\
申命記 34:10-12 & \hyperref[sec:Z6ZZRqsI0K0]{2024-05-25 晚堂《面對面的認識》申命記 34\_10-12  講員:盧思源姑娘} & 2024-05-28 & \href{https://youtube.com/watch?v=Z6ZZRqsI0K0}{\texttt{Z6ZZRqsI0K0}} \\
  & \hyperref[sec:luI8_7RFevk]{2024-05-26 早堂 樓恩德牧師《耶穌復活的真諦》} & 2024-05-28 & \href{https://youtube.com/watch?v=luI8_7RFevk}{\texttt{luI8\_7RFevk}} \\
約珥書 2:28-32 & \hyperref[sec:QGag1Zr1Cdk]{2024-05-26 午堂《預備耶和華大而可畏的日子》約珥書 2\_28-32  講員:郭鴻標牧師} & 2024-05-29 & \href{https://youtube.com/watch?v=QGag1Zr1Cdk}{\texttt{QGag1Zr1Cdk}} \\
帖撒羅尼迦前書  & \hyperref[sec:eejVMLNyY8I]{2024-05-30 周四晚查經班《帖撒羅尼迦前書(7)》 講員:曾裔貴牧師} & 2024-05-31 & \href{https://youtube.com/watch?v=eejVMLNyY8I}{\texttt{eejVMLNyY8I}} \\
以弗所書 4:17-32 & \hyperref[sec:7PHC2DLBPCo]{2024-06-01 晚堂《黑暗與光明》以弗所書4\_17-32 講員:曾裔貴牧師} & 2024-06-04 & \href{https://youtube.com/watch?v=7PHC2DLBPCo}{\texttt{7PHC2DLBPCo}} \\
列王記下 5:1-19 & \hyperref[sec:Il1UURetEws]{2024-06-02 午堂《無名的小使女》列王記下5\_1-19 講員:楊明麗姑娘} & 2024-06-04 & \href{https://youtube.com/watch?v=Il1UURetEws}{\texttt{Il1UURetEws}} \\
歷代志上 13:9-14 & \hyperref[sec:ui5VIns3kj4]{2024-06-02 早堂《前車可鑑》歷代志上13\_9-14 講員:陳一華牧師} & 2024-06-05 & \href{https://youtube.com/watch?v=ui5VIns3kj4}{\texttt{ui5VIns3kj4}} \\
帖撒羅尼迦前書  & \hyperref[sec:DjV4K1Dg8GU]{2024-06-06 周四晚查經班《帖撒羅尼迦前書(8)》 講員:曾裔貴牧師} & 2024-06-07 & \href{https://youtube.com/watch?v=DjV4K1Dg8GU}{\texttt{DjV4K1Dg8GU}} \\
詩篇 84:1-12 & \hyperref[sec:CutYNgBop3c]{2024-06-09 午堂《永恆的居所》詩篇 84\_1-12  講員:馮耀榮牧師} & 2024-06-11 & \href{https://youtube.com/watch?v=CutYNgBop3c}{\texttt{CutYNgBop3c}} \\
路加福音 9:46-48 & \hyperref[sec:AT0lX1Ebpa4]{2024-06-09 早堂《接待小子的服侍》路加福音 9\_46-48  講員:蔣文忠傳道} & 2024-06-11 & \href{https://youtube.com/watch?v=AT0lX1Ebpa4}{\texttt{AT0lX1Ebpa4}} \\
帖撒羅尼迦前書  & \hyperref[sec:w32Ot_b7FcE]{2024-06-13 周四晚查經班《帖撒羅尼迦前書(9)》 講員:曾裔貴牧師} & 2024-06-14 & \href{https://youtube.com/watch?v=w32Ot_b7FcE}{\texttt{w32Ot\_b7FcE}} \\
民數記 14:24 & \hyperref[sec:1kzKFPo0unA]{2024-06-15 晚堂《專心跟從神》民數記 14\_24  講員:陳迦勒導師} & 2024-06-18 & \href{https://youtube.com/watch?v=1kzKFPo0unA}{\texttt{1kzKFPo0unA}} \\
詩篇 23:1-6 & \hyperref[sec:pibg95moZiM]{2024-06-16 午堂《有爸爸的愛,你不孤單》詩篇 23\_1-6  講員:鄧泰康傳道} & 2024-06-18 & \href{https://youtube.com/watch?v=pibg95moZiM}{\texttt{pibg95moZiM}} \\
路加福音 15:11-32 & \hyperref[sec:Yzl3vxypdqk]{2024-06-16 早堂《父親的心》路加福音 15\_11-32  講員:許亞光牧師} & 2024-06-19 & \href{https://youtube.com/watch?v=Yzl3vxypdqk}{\texttt{Yzl3vxypdqk}} \\
使徒行傳 1:1-15 & \hyperref[sec:aAKtjwj8FbQ]{2024-06-23 午堂《復興的等候》使徒行傳 1\_1-15  講員:陳麗蟬牧師} & 2024-06-25 & \href{https://youtube.com/watch?v=aAKtjwj8FbQ}{\texttt{aAKtjwj8FbQ}} \\
馬太福音 5:13-15 & \hyperref[sec:vr_pYc_qVuc]{2024-06-23 早堂《好ye 介紹返》馬太福音 5\_13-15  講員:陳韋安牧師} & 2024-06-25 & \href{https://youtube.com/watch?v=vr_pYc-qVuc}{\texttt{vr\_pYc-qVuc}} \\
帖撒羅尼迦前書  & \hyperref[sec:4L9GLJKp4e4]{2024-06-27 周四晚查經班《帖撒羅尼迦前書(10)》 講員:曾裔貴牧師} & 2024-06-29 & \href{https://youtube.com/watch?v=4L9GLJKp4e4}{\texttt{4L9GLJKp4e4}} \\
帖撒羅尼迦後書  & \hyperref[sec:uJ_qz3eVEf0]{2024-07-11 周四晚查經班《帖撒羅尼迦後書(1)》 講員:曾裔貴牧師} & 2024-07-12 & \href{https://youtube.com/watch?v=uJ_qz3eVEf0}{\texttt{uJ\_qz3eVEf0}} \\
帖撒羅尼迦後書  & \hyperref[sec:ntrA1aTVgRk]{2024-07-18 周四晚查經班《帖撒羅尼迦後書(2)》 講員:曾裔貴牧師} & 2024-07-19 & \href{https://youtube.com/watch?v=ntrA1aTVgRk}{\texttt{ntrA1aTVgRk}} \\
帖撒羅尼迦後書  & \hyperref[sec:ewzMm4Y0R5I]{2024-08-01 周四晚查經班《帖撒羅尼迦後書(3)》 講員:曾裔貴牧師} & 2024-08-06 & \href{https://youtube.com/watch?v=ewzMm4Y0R5I}{\texttt{ewzMm4Y0R5I}} \\
帖撒羅尼迦後書  & \hyperref[sec:bkamj1iGLMk]{2024-08-15 周四晚查經班《帖撒羅尼迦後書(4)》 講員:曾裔貴牧師} & 2024-08-16 & \href{https://youtube.com/watch?v=bkamj1iGLMk}{\texttt{bkamj1iGLMk}} \\
帖撒羅尼迦後書  & \hyperref[sec:IJNXTQniUCA]{2024-08-22 周四晚查經班《帖撒羅尼迦後書(5)》 講員:曾裔貴牧師} & 2024-08-23 & \href{https://youtube.com/watch?v=IJNXTQniUCA}{\texttt{IJNXTQniUCA}} \\
  & \hyperref[sec:I4c89OnZC6M]{2024-08-29 周四晚查經班《淺談末世論(1)》 講員:曾裔貴牧師} & 2024-08-30 & \href{https://youtube.com/watch?v=I4c89OnZC6M}{\texttt{I4c89OnZC6M}} \\
  & \hyperref[sec:a8KnWbij7Uk]{2024-09-19 周四晚查經班《淺談末世論(2)》 講員:曾裔貴牧師} & 2024-09-20 & \href{https://youtube.com/watch?v=a8KnWbij7Uk}{\texttt{a8KnWbij7Uk}} \\
  & \hyperref[sec:tl2Gc71LinQ]{2024-09-26 周四晚查經班《淺談末世論(3)》 講員:曾裔貴牧師} & 2024-09-27 & \href{https://youtube.com/watch?v=tl2Gc71LinQ}{\texttt{tl2Gc71LinQ}} \\
  & \hyperref[sec:H4W4IN0LZRY]{2024-10-03 周四晚查經班《淺談末世論(4)》 講員:曾裔貴牧師} & 2024-10-04 & \href{https://youtube.com/watch?v=H4W4IN0LZRY}{\texttt{H4W4IN0LZRY}} \\
  & \hyperref[sec:aTaCCC_F45U]{2024-10-10 周四晚查經班《改革之路:過去與未來.(1)》 講員:曾裔貴牧師} & 2024-10-12 & \href{https://youtube.com/watch?v=aTaCCC-F45U}{\texttt{aTaCCC-F45U}} \\
  & \hyperref[sec:Sb_mbdX9Sa4]{2024-10-24 周四晚查經班《改革之路(2):馬丁路德的十架神學》 講員:曾裔貴牧師} & 2024-10-25 & \href{https://youtube.com/watch?v=Sb-mbdX9Sa4}{\texttt{Sb-mbdX9Sa4}} \\
  & \hyperref[sec:ziJGxjL2HKY]{2024-11-07 周四晚查經班《改革之路(3)》 講員:曾裔貴牧師} & 2024-11-09 & \href{https://youtube.com/watch?v=ziJGxjL2HKY}{\texttt{ziJGxjL2HKY}} \\
  & \hyperref[sec:jvEmZbSGJ_A]{2024-11-14 周四晚查經班《改革之路(4):伊拉斯姆》 講員:曾裔貴牧師} & 2024-11-15 & \href{https://youtube.com/watch?v=jvEmZbSGJ_A}{\texttt{jvEmZbSGJ\_A}} \\
  & \hyperref[sec:DzI9N0JOL0s]{2024-12-12 周四晚查經班《改革之路(5):加爾文(一)》 講員:曾裔貴牧師} & 2024-12-13 & \href{https://youtube.com/watch?v=DzI9N0JOL0s}{\texttt{DzI9N0JOL0s}} \\
  & \hyperref[sec:4EM8nFC2gtc]{2025-01-02 周四晚查經班《改革之路(6):加爾文(二)》 講員:曾裔貴牧師} & 2025-01-03 & \href{https://youtube.com/watch?v=4EM8nFC2gtc}{\texttt{4EM8nFC2gtc}} \\
  & \hyperref[sec:3Sqlgs98HDM]{2025-01-09 周四晚查經班《改革之路(7):加爾文(三)》 講員:曾裔貴牧師} & 2025-01-10 & \href{https://youtube.com/watch?v=3Sqlgs98HDM}{\texttt{3Sqlgs98HDM}} \\
  & \hyperref[sec:Dn07Lw4F37U]{2025-01-16 周四晚查經班《改革之路(8):加爾文(四)》 講員:曾裔貴牧師} & 2025-01-17 & \href{https://youtube.com/watch?v=Dn07Lw4F37U}{\texttt{Dn07Lw4F37U}} \\
  & \hyperref[sec:5Lw6EDNBsNI]{2025-02-06 周四晚查經班《改革之路(9):諾克斯(一)》 講員:曾裔貴牧師} & 2025-02-07 & \href{https://youtube.com/watch?v=5Lw6EDNBsNI}{\texttt{5Lw6EDNBsNI}} \\
  & \hyperref[sec:q7QZXuHzcds]{2025-02-13 周四晚查經班《改革之路(10)完結篇:未來》 講員:曾裔貴牧師} & 2025-02-14 & \href{https://youtube.com/watch?v=q7QZXuHzcds}{\texttt{q7QZXuHzcds}} \\
羅馬書  & \hyperref[sec:mCrcbeplxo4]{2025-02-20 周四晚查經班《羅馬書(1)引言:保羅其人其事》 講員:曾裔貴牧師} & 2025-02-22 & \href{https://youtube.com/watch?v=mCrcbeplxo4}{\texttt{mCrcbeplxo4}} \\
羅馬書  & \hyperref[sec:TdP7rHPK3XI]{2025-03-06 周四晚查經班《羅馬書(2):分題:神的義(一)》 講員:曾裔貴牧師} & 2025-03-08 & \href{https://youtube.com/watch?v=TdP7rHPK3XI}{\texttt{TdP7rHPK3XI}} \\
羅馬書  & \hyperref[sec:WjqHgH3CzNw]{2025-03-13 周四晚查經班《羅馬書(3):分題:神的義(二)》 講員:曾裔貴牧師} & 2025-03-14 & \href{https://youtube.com/watch?v=WjqHgH3CzNw}{\texttt{WjqHgH3CzNw}} \\
羅馬書  & \hyperref[sec:8LIvofDzP5M]{2025-03-20 周四晚查經班《羅馬書(4):分題:神的義(三)》 講員:曾裔貴牧師} & 2025-03-21 & \href{https://youtube.com/watch?v=8LIvofDzP5M}{\texttt{8LIvofDzP5M}} \\
羅馬書  & \hyperref[sec:nWNWvRiUB3k]{2025-03-27 周四晚查經班《羅馬書(5):分題:神的義(四)》 講員:曾裔貴牧師} & 2025-03-29 & \href{https://youtube.com/watch?v=nWNWvRiUB3k}{\texttt{nWNWvRiUB3k}} \\
羅馬書  & \hyperref[sec:7C1Uqd4_AYQ]{2025-04-03 周四晚查經班《羅馬書(6):分題:神的義(五)》 講員:曾裔貴牧師} & 2025-04-05 & \href{https://youtube.com/watch?v=7C1Uqd4-AYQ}{\texttt{7C1Uqd4-AYQ}} \\
羅馬書  & \hyperref[sec:JU6MXStfi9g]{2025-04-17 周四晚查經班《羅馬書(7):分題:神的義(六)》 講員:曾裔貴牧師} & 2025-04-19 & \href{https://youtube.com/watch?v=JU6MXStfi9g}{\texttt{JU6MXStfi9g}} \\
羅馬書  & \hyperref[sec:b2w9fc63m3c]{2025-04-24 周四晚查經班《羅馬書(8):分題:神的義(七)》 講員:曾裔貴牧師} & 2025-04-25 & \href{https://youtube.com/watch?v=b2w9fc63m3c}{\texttt{b2w9fc63m3c}} \\
羅馬書  & \hyperref[sec:DvUFUpNSEFw]{2025-05-08 周四晚查經班《羅馬書(9):分題:神的義(八)》 講員:曾裔貴牧師} & 2025-05-09 & \href{https://youtube.com/watch?v=DvUFUpNSEFw}{\texttt{DvUFUpNSEFw}} \\
羅馬書  & \hyperref[sec:BtALIL64hKM]{2025-05-15 周四晚查經班《羅馬書(10):分題:舊約中的因信稱義(一)》 講員:曾裔貴牧師} & 2025-05-16 & \href{https://youtube.com/watch?v=BtALIL64hKM}{\texttt{BtALIL64hKM}} \\
羅馬書  & \hyperref[sec:QTpL9O5GR3E]{2025-05-22 周四晚查經班《羅馬書(11):分題:稱義帶來與神復和》 講員:曾裔貴牧師} & 2025-05-23 & \href{https://youtube.com/watch?v=QTpL9O5GR3E}{\texttt{QTpL9O5GR3E}} \\
羅馬書  & \hyperref[sec:23gUJZhHg2o]{2025-05-22 周四晚查經班《羅馬書(11):分題:稱義帶來與神復和》 講員:曾裔貴牧師 NEW} & 2025-05-24 & \href{https://youtube.com/watch?v=23gUJZhHg2o}{\texttt{23gUJZhHg2o}} \\
羅馬書  & \hyperref[sec:r2uz84DdreY]{2025-06-05 周四晚查經班《羅馬書(12):分題:因信稱義的福》 講員:曾裔貴牧師} & 2025-06-06 & \href{https://youtube.com/watch?v=r2uz84DdreY}{\texttt{r2uz84DdreY}} \\
羅馬書  & \hyperref[sec:G8aDRsOGytw]{2025-06-12 周四晚查經班《羅馬書(13):分題:克勝罪惡的新生命》 講員:曾裔貴牧師} & 2025-06-13 & \href{https://youtube.com/watch?v=G8aDRsOGytw}{\texttt{G8aDRsOGytw}} \\
羅馬書  & \hyperref[sec:7LZOyo5Zt3c]{2025-06-19 周四晚查經班《羅馬書(14):分題:不可犯罪?可不犯罪》 講員:曾裔貴牧師} & 2025-06-20 & \href{https://youtube.com/watch?v=7LZOyo5Zt3c}{\texttt{7LZOyo5Zt3c}} \\
羅馬書  & \hyperref[sec:wh3ZqIs7b68]{2025-06-26 周四晚查經班《羅馬書(15):分題:不可犯罪?可不犯罪(二)》 講員:曾裔貴牧師} & 2025-06-27 & \href{https://youtube.com/watch?v=wh3ZqIs7b68}{\texttt{wh3ZqIs7b68}} \\
羅馬書  & \hyperref[sec:PTM8SSdjvVE]{2025-07-03 周四晚查經班《羅馬書(16):分題:恩㤟與律法》 講員:曾裔貴牧師} & 2025-07-04 & \href{https://youtube.com/watch?v=PTM8SSdjvVE}{\texttt{PTM8SSdjvVE}} \\
羅馬書  & \hyperref[sec:HoQGEbdeqSk]{2025-07-24 周四晚查經班《羅馬書(17):分題:恩㤟與律法》 講員:曾裔貴牧師} & 2025-07-26 & \href{https://youtube.com/watch?v=HoQGEbdeqSk}{\texttt{HoQGEbdeqSk}} \\
羅馬書  & \hyperref[sec:mDmug5anuis]{2025-07-31 周四晚查經班《羅馬書(18):分題:聖靈帶動萬物復興》 講員:曾裔貴牧師} & 2025-08-02 & \href{https://youtube.com/watch?v=mDmug5anuis}{\texttt{mDmug5anuis}} \\
羅馬書  & \hyperref[sec:HwCOo8EiKp0]{2025-08-14 周四晚查經班《羅馬書(19):分題:聖靈帶動萬物復興二》 講員:曾裔貴牧師} & 2025-08-18 & \href{https://youtube.com/watch?v=HwCOo8EiKp0}{\texttt{HwCOo8EiKp0}} \\
羅馬書  & \hyperref[sec:3_mKpJn_jtA]{2025-08-21 周四晚查經班《羅馬書(20):分題:神的選民猶太人可得救嗎?》 講員:曾裔貴牧師} & 2025-08-24 & \href{https://youtube.com/watch?v=3-mKpJn-jtA}{\texttt{3-mKpJn-jtA}} \\
  & \hyperref[sec:MB59D1BA0Dc]{宣道會錦繡堂2025年8月30日晚堂崇拜直播} & 2025-08-30 & \href{https://youtube.com/watch?v=MB59D1BA0Dc}{\texttt{MB59D1BA0Dc}} \\
羅馬書  & \hyperref[sec:3TbpwI2Yugs]{2025-09-05 周五晚查經班《羅馬書(21):分題:信心的義與福音的普世性》 講員:曾裔貴牧師} & 2025-09-07 & \href{https://youtube.com/watch?v=3TbpwI2Yugs}{\texttt{3TbpwI2Yugs}} \\
羅馬書  & \hyperref[sec:1GDpbDYTYf0]{2025-09-12 周五晚查經班《羅馬書(22):福音:人的選擇與後果》 講員:曾裔貴牧師} & 2025-09-13 & \href{https://youtube.com/watch?v=1GDpbDYTYf0}{\texttt{1GDpbDYTYf0}} \\
羅馬書  & \hyperref[sec:ZdF__6Sl2No]{2025-09-26 周五晚查經班《羅馬書(23):分題:神選民與恩典—神的不變救恩》 講員:曾裔貴牧師} & 2025-09-27 & \href{https://youtube.com/watch?v=ZdF-_6Sl2No}{\texttt{ZdF-\_6Sl2No}} \\
羅馬書  & \hyperref[sec:bw6FVf5f_n0]{2025-10-10 周五晚查經班《羅馬書(24):分題:神國救恩歷史的智慧》 講員:曾裔貴牧師} & 2025-10-11 & \href{https://youtube.com/watch?v=bw6FVf5f_n0}{\texttt{bw6FVf5f\_n0}} \\
羅馬書  & \hyperref[sec:FAt1alHLvdg]{2025-10-17 周五晚查經班《羅馬書(25):分題:從奧秘到信徒的生活》 講員:曾裔貴牧師} & 2025-10-19 & \href{https://youtube.com/watch?v=FAt1alHLvdg}{\texttt{FAt1alHLvdg}} \\
羅馬書  & \hyperref[sec:yNx82XHOkbk]{2025-10-31 周五晚查經班《羅馬書(26):分題:教會是發掘,發展及發揮恩賜的場景》 講員:曾裔貴牧師 1} & 2025-11-03 & \href{https://youtube.com/watch?v=yNx82XHOkbk}{\texttt{yNx82XHOkbk}} \\
羅馬書  & \hyperref[sec:B5vFf1eCUBA]{2025-11-14 周五晚查經班《羅馬書(27):分題:合乎中道的自我認識與恩賜的正確運用》 講員:曾裔貴牧師} & 2025-11-15 & \href{https://youtube.com/watch?v=B5vFf1eCUBA}{\texttt{B5vFf1eCUBA}} \\
羅馬書  & \hyperref[sec:0TRElWqGJ4E]{2025-11-21 周五晚查經班《羅馬書(28):分題:恩賜是要來用的,在教會的合一場景中運用》 講員:曾裔貴牧師} & 2025-11-29 & \href{https://youtube.com/watch?v=0TRElWqGJ4E}{\texttt{0TRElWqGJ4E}} \\
羅馬書  & \hyperref[sec:WdwoyEOQVUg]{2025-11-28 周五晚查經班《羅馬書(29):分題:在群體生活中的自我認知》 講員:曾裔貴牧師} & 2025-11-29 & \href{https://youtube.com/watch?v=WdwoyEOQVUg}{\texttt{WdwoyEOQVUg}} \\
羅馬書  & \hyperref[sec:2B_XWriXVXo]{2025-12-12 周五晚查經班《羅馬書(30):分題:愛與順服》 講員:曾裔貴牧師} & 2025-12-13 & \href{https://youtube.com/watch?v=2B_XWriXVXo}{\texttt{2B\_XWriXVXo}} \\
羅馬書  & \hyperref[sec:ego6cJukrIE]{2025-12-19 周五晚查經班《羅馬書(31):分題:愛與順服(二)》 講員:曾裔貴牧師} & 2025-12-21 & \href{https://youtube.com/watch?v=ego6cJukrIE}{\texttt{ego6cJukrIE}} \\
羅馬書  & \hyperref[sec:aePwbTZJwec]{2026-01-02 周五晚查經班《羅馬書(32):分題:如何處理教會內「非核心信仰」的差異》 講員:曾裔貴牧師} & 2026-01-03 & \href{https://youtube.com/watch?v=aePwbTZJwec}{\texttt{aePwbTZJwec}} \\
羅馬書  & \hyperref[sec:nB_6d0PSa6M]{2026-01-09 周五晚查經班《羅馬書(33):分題:如何處理教會內「非核心信仰」的差異(二)》 講員:曾裔貴牧師} & 2026-01-10 & \href{https://youtube.com/watch?v=nB_6d0PSa6M}{\texttt{nB\_6d0PSa6M}} \\
羅馬書  & \hyperref[sec:pe8biigoai0]{2026-01-16 周五晚查經班《羅馬書(34):分題:擔得起,才配得進神的國》 講員:曾裔貴牧師} & 2026-01-17 & \href{https://youtube.com/watch?v=pe8biigoai0}{\texttt{pe8biigoai0}} \\
羅馬書  & \hyperref[sec:UU4dvSZbobM]{2026-01-23 周五晚查經班《羅馬書(三十五):分題:有盼望,有異象,有擔當的人》 講員:曾裔貴牧師} & 2026-01-24 & \href{https://youtube.com/watch?v=UU4dvSZbobM}{\texttt{UU4dvSZbobM}} \\
羅馬書  & \hyperref[sec:LvE6qlAUC04]{2026-02-06 周五晚查經班《羅馬書(三十六):分題:保羅 有擔當的使徒》 講員:曾裔貴牧師} & 2026-02-08 & \href{https://youtube.com/watch?v=LvE6qlAUC04}{\texttt{LvE6qlAUC04}} \\
羅馬書  & \hyperref[sec:3rvmIgiASuQ]{2026-02-13 周五晚查經班《羅馬書(三十七):分題:福音的實踐:從教義到生活》 講員:曾裔貴牧師} & 2026-02-14 & \href{https://youtube.com/watch?v=3rvmIgiASuQ}{\texttt{3rvmIgiASuQ}} \\
羅馬書  & \hyperref[sec:4pLVUVJIfgU]{2026-03-06 周五晚查經班《羅馬書(三十八):分題:使徒保羅的深情問安》 講員:曾裔貴牧師} & 2026-03-09 & \href{https://youtube.com/watch?v=4pLVUVJIfgU}{\texttt{4pLVUVJIfgU}} \\
羅馬書  & \hyperref[sec:oQdnBIyZabY]{2026-03-13 周五晚查經班《羅馬書(三十九):分題:來自保羅的問安(二)》 講員:曾裔貴牧師} & 2026-03-16 & \href{https://youtube.com/watch?v=oQdnBIyZabY}{\texttt{oQdnBIyZabY}} \\
羅馬書  & \hyperref[sec:HW8W_v51HY8]{2026-03-20 周五晚查經班《羅馬書(四十):分題:來自保羅的問安(三)》 講員:曾裔貴牧師} & 2026-03-24 & \href{https://youtube.com/watch?v=HW8W-v51HY8}{\texttt{HW8W-v51HY8}} \\
羅馬書  & \hyperref[sec:Ex_WASW50ZA]{2026-04-10 周五晚查經班《羅馬書(四十一):分題:來自保羅的問安(四)》 講員:曾裔貴牧師} & 2026-04-11 & \href{https://youtube.com/watch?v=Ex-WASW50ZA}{\texttt{Ex-WASW50ZA}} \\
羅馬書  & \hyperref[sec:WmefrJ6J_Lg]{2026-05-08 周五晚查經班《羅馬書(四十二):總論:福音帶來與神和好,所以教會必須向萬民宣講福音.》 講員:曾裔貴牧師} & 2026-05-10 & \href{https://youtube.com/watch?v=WmefrJ6J-Lg}{\texttt{WmefrJ6J-Lg}} \\
出埃及記  & \hyperref[sec:GROoST1kr4I]{2026-05-15 周五晚查經班《出埃及記研讀 進入應許之地(一)》 講員:曾裔貴牧師} & 2026-05-18 & \href{https://youtube.com/watch?v=GROoST1kr4I}{\texttt{GROoST1kr4I}} \\
出埃及記  & \hyperref[sec:kZF3v6CUrpM]{2026-05-22 周五晚查經班《出埃及記研讀(二)應許在黑暗中增長》 講員:曾裔貴牧師} & 2026-05-23 & \href{https://youtube.com/watch?v=kZF3v6CUrpM}{\texttt{kZF3v6CUrpM}} \\
出埃及記  & \hyperref[sec:4cUSeyCs18Q]{2026-05-29 周五晚查經班《出埃及記研讀(三)神所用的小人物:收生婆,母親,公主如何瓦解暴政》 講員:曾裔貴牧師} & 2026-05-30 & \href{https://youtube.com/watch?v=4cUSeyCs18Q}{\texttt{4cUSeyCs18Q}} \\
出埃及記  & \hyperref[sec:HMEU7PxsF0Y]{2026-06-05 周五晚查經班《出埃及記研讀(四)神國領袖的再造工程》 講員:曾裔貴牧師} & 2026-06-07 & \href{https://youtube.com/watch?v=HMEU7PxsF0Y}{\texttt{HMEU7PxsF0Y}} \\
出埃及記  & \hyperref[sec:cWahrBO910c]{2026-06-26 周五晚查經班《出埃及記研讀(五)五個「我不能」與一個「我是」》 講員:曾裔貴牧師} & 2026-06-27 & \href{https://youtube.com/watch?v=cWahrBO910c}{\texttt{cWahrBO910c}} \\
出埃及記  & \hyperref[sec:dj6fIzaDTMQ]{2026-07-17 周五晚查經班《出埃及記研讀(六)舊約先知的蒙召》 講員:曾裔貴牧師} & 2026-07-20 & \href{https://youtube.com/watch?v=dj6fIzaDTMQ}{\texttt{dj6fIzaDTMQ}} \\
出埃及記  & \hyperref[sec:FKTWi97mjmY]{2026-07-31 周五晚查經班《出埃及記研讀(七)舊約先知的蒙召(二)》 講員:曾裔貴牧師} & 2026-08-03 & \href{https://youtube.com/watch?v=FKTWi97mjmY}{\texttt{FKTWi97mjmY}} \\
\end{xltabular}
}
\newpage



\section{啟示錄 21:1-7}
\label{sec:IlLnn7WOgiE}
\textbf{2023-01-01 加恩樓《上帝給我們的盼望及出路》啟示錄 21\_1-7 講員:鄧泰康傳道}
\newline
\newline
連結: \href{https://youtube.com/watch?v=IlLnn7WOgiE}{\texttt{https://youtube.com/watch?v=IlLnn7WOgiE}} ~~~~ 語音日期: 2023-01-03
\newline
\newline
\hyperref[sec:FKTWi97mjmY]{\small{< < < PREV SERMON < < <}}
~
\hyperref[sec:index_chronic]{\small{[返順時目]}}
~
\hyperref[sec:index_scriptual]{\small{[返順卷目]}}
~
\hyperref[sec:d2MEojJVUd4]{\small{> > > NEXT SERMON > > >}}
\newline
\newline
啟示錄 21:1-7
\newline
\begin{longtable}{cl}
\hline
\hline
章節 & 經文 (和合本修訂版)\\
\hline
21:1 & \begin{tabularx}{0.7\textwidth}{X} 我又看見一個新天新地，因為先前的天和先前的地已經過去了，海也不再有了。 \end{tabularx} \\ \\ \relax
21:2 & \begin{tabularx}{0.7\textwidth}{X} 我又看見聖城，新耶路撒冷由神那裡，從天而降，預備好了，就如新娘打扮整齊，等候丈夫。 \end{tabularx} \\ \\ \relax
21:3 & \begin{tabularx}{0.7\textwidth}{X} 我聽見有大聲音從寶座出來，說：「看哪，神的帳幕在人間！他要和他們同住，他們要作他的子民。神要親自與他們同在。 \end{tabularx} \\ \\ \relax
21:4 & \begin{tabularx}{0.7\textwidth}{X} 神要擦去他們一切的眼淚；不再有死亡，也不再有悲哀、哭號、痛苦，因為先前的事都過去了。」 \end{tabularx} \\ \\ \relax
21:5 & \begin{tabularx}{0.7\textwidth}{X} 那位坐在寶座上的說：「看哪，我把一切都更新了！」他又說：「你要寫下來，因為這些話是可信靠的，是真實的。」 \end{tabularx} \\ \\ \relax
21:6 & \begin{tabularx}{0.7\textwidth}{X} 他又對我說：「成了！我是阿拉法，我是俄梅戛；我是開始，我是終結。我要把生命的泉水白白賜給那口渴的人喝。 \end{tabularx} \\ \\ \relax
21:7 & \begin{tabularx}{0.7\textwidth}{X} 得勝的要承受這些為業；我要作他的神，他要作我的兒子。 \end{tabularx} \\ \\
[1ex]
\hline
\hline
\end{longtable}
$^{1}$好,弟兄姊妹,新年快樂!.
新年真是快樂嘛.
其實我今天第一個要問大家的問題就是.
2022年大家過得怎樣?.
OK嗎?.
過得好嗎?.
因為在過去的時間,我的確都挺相信.
2022年其實對於很多人來說都不容易.
例如疾病,工作的影響,社交的影響,上學的影響.
各樣東西都非常不容易.
但問題是當我們作為基督徒的時候.
面對這些不容易,困難.
我們應該是怎樣呢?.
在上星期,我跟富哥有幸.
富哥邀請我一起參與了一個話劇表演.
這個話劇表演的名稱是《絕望研究所》.
大概這個表演是透過他們很專業的話劇.
帶出了六個場景.
其實我要這樣形容的.
六個場景所帶給大家的感覺和氣氛.
是很灰的.
甚至有些人的回應是覺得很沉重的.
不容易的.
無論在社會上看到的.
在教會當中看到的各樣的場景.
但這個話劇的設計很有趣.
讓我們這些觀眾一起參與.
入場的時候每人派一張卡.
卡裡的時間是一邊寫著盼望.
一邊寫著絕望.
做完每個場景之後.
他們就說現在我們一起投票.
包括台上台下的.
我們一起投票.
如果大家覺得是盼望的.
你就展示盼望.
如果覺得是絕望的.
你就展示絕望.
當然我們在裡面一起.
不是每個環節我們都有一份投票.

$^{41}$不過每個環節當我看完之後.
我心裡都很認真.
我不是只看一個人的表演.
而是我很認真在面對這樣的處境.
這樣的環境當中.
究竟我是盼望還是絕望.
看完整個話劇之後.
我心裡我實際的表態.
我都盼望.
我也問的.
我說這麼沉重這麼不容易.
為什麼我都會舉盼望.
而是我問自己的時候.
我很清楚.
因為每一個處境裡面.
我都想到上帝.
在這個話劇裡面的時候.
當差不多結束的時候.
我們牧者說了一句話.
我覺得很震撼我.
他說我們的信仰.
給我們最大的寶貝是什麼呢.
就是有盼望.
而且是一個終極的盼望.
當今天跟大家分享這段經文的時候.
我不知道大家知不知道這段經文.
最經常出現在一個什麼場合.
大家知不知道.
是在安息禮拜.
很多時候就是當人離開了.
然後我們透過這段經文.
告訴你不用擔心.
因為你所愛的人.
現在去到一個地方.
是沒有眼淚的.
那個地方是沒有疼痛的.
那個地方是沒有疾病的.
我們這樣告訴他.
因為世人去看死亡就是絕望.
因為死亡一生都完結了.

$^{81}$我沒有回頭路.
但是接著我們的信仰就告訴他.
不是.
我們的信仰是很有盼望的.
因為人離開這個世界之後.
他所進入的並不是絕望.
他所要去的是一個很有盼望的地方.
這是我們一直在跟他們說的.
所以弟兄姊妹.
我想告訴大家.
當今天是1月1號.
昨晚是31號.
我不知道大家昨晚是怎樣慶祝.
在我年輕的時候.
我真的會去商場那些地方大叫.
1987年6月.
我都做過這些的.
但是昨晚我約了周公一起倒數.
你會發覺一件事.
我們都很慶祝.
慶祝什麼.
我都問慶祝什麼.
我覺得2022年最慶祝的就是.
我們還在生.
我們仍然有生命氣息.
但是在裡面的時候.
我覺得更加重要的.
不要停留在這裡.
而是我們要在神的當中看到盼望.
這段經文裡面就是.
如果大家知道.
背景是非常重要的.
是使徒約翰因為當時守住他的信仰.
所以被流放到一個地方叫做拔摩島.
大家去看流放就不要以為今天是旅行.
不要以為只是送到偏僻的地方.
我不知道大家喜不喜歡看韓劇.
如果喜歡看韓劇的時候.
你要看古代的韓劇.
他們很流放.

$^{121}$將大官員流放.
其實是很慘的一件事.
對於當時的以色列人來說.
對於他們那群十二門徒來說.
其實當時是一個很絕望的環境.
羅馬帝國那麼強大.
我們這群連國家都沒有的人.
我們還說什麼盼望.
現在我們被迫去到一個情況.
是我們已經被人流放了.
有人要死就死了.
有些人刀釘十字架就死了.
有些人就被出賣了.
整個環境是很絕望很灰暗的.
但是上帝竟然就在這個時間.
將啟示錄的異象給了約翰.
等於上帝的盼望往往就是.
在我們感受到最絕望的時候.
你是更加看到.
我們都很清楚一個例子.
就是一點光在黑暗的時間.
你就見到 是最光亮的.
你就會發覺原來.
我們在神裡面的時候.
上帝是直接的將盼望賜給我們.
在患難的當中賜給我們安慰.
你發覺整個啟示錄.
最奇妙的地方是什麼呢.
就是其實裡面說了很多災難.
但是基本上分了兩班人.
一班是屬於上帝的.
一班不屬於上帝的.
不屬於上帝的人.
其實他就會跟隨著魔鬼.
一起面對著一波接一波.
有史以來最可怕的大災難.
直至滅亡.
但是一班屬於神的人.
屬於耶穌的人.
卻是在這些情況之下.

$^{161}$竟然進入了一個新天新地.
上帝為他們預備的地方.
是完完全全充滿幸福的地方.
一聽到的時候.
我不知道大家有沒有這種感覺.
完完全全是一套電影.
所以有時候我們經常說.
這些是小朋友跟小朋友說的聖經.
只是哄騙人而已.
弟兄姊妹.
我不知道大家還記不記得.
有一套電影是2004年.
相隔今天差不多已經是18年的時間.
叫做明日之後.
那套電影大家還記得說什麼嗎.
其實是說有一個突然間的冷風.
非常強勁的冷風突然吹到美國的時間.
在很短的時間大降溫.
接著漫天風雪.
鋪滿了整個紐約市的自由神像.
只剩下頭和手.
接著有一幕有直升機去救援.
飛到半空的時間.
整架直升機突然結冰.
裡面的人結冰了.
然後墮機.
之後很多人死亡.
一聽就知道是電影.
這些是娛樂性的電影.
弟兄姊妹.
看到上星期的一件新聞嗎.
美國水牛城風暴.
大家有沒有看過裡面的內容.
知道是怎麼發生的嗎.
原來它的成因.
就是跟明日之後所說的成因.
是一模一樣的.
北極壓不住寒冷的風.
突然間吹襲.
突然間的天氣由10度一跌.

$^{201}$跌到零下50度.
弟兄姊妹有很多時候.
你當作是一件童話故事.
你當作只不過是一件.
我們跟小朋友說聖經故事.
但其實原來我想告訴大家.
有很多事情就在你身邊.
在你不在意的時候就發生了.
所以今天跟大家說這段經文.
其實內容一點都不複雜.
但問題在於你怎樣去相信上帝的話.
你是不是真的會相信.
上帝所說的這些話是重要的.
如果重要的話.
我們在2023年.
將會需要一些更新和改變.
這裡的經文一開始.
就是我見到一個新天新地.
因為先前的天地已經過去.
海也不再有了.
當這些大災難過去之後.
因為這裡已經去到21章.
大災難都已經過去了.
他們就見到有一個新天新地.
什麼叫新天新地呢.
就是他很強調的說.
舊的天地已經過去了.
原來什麼叫做盼望.
盼望就是原來舊的可以過去.
今天你所面對的苦難.
無論多重都好.
原來會過去的.
我們一起戴口罩都戴了三年了.
但大家心裡都很清楚.
大家都知道這個口罩.
遲早有一天會過去的.
我們今天所面對的很多困難.
很多苦難其實是會過去的.
其實這個就是給我們盼望很重要的出路.
上帝不是要你困在今天的困難裡面.

$^{241}$你不是由早到晚只想著.
今天我們所面對的困難.
你這樣想的時候.
只會逼到你去跳樓.
但你要想的就是.
不是的.
原來在上帝的掌管下.
一切的事情都會過去的.
新天新地就是一切舊的都會過去.
而第二件事就是.
你會發覺舊的標準同樣都會過去.
今天在地上我們追求的價值觀.
各樣的東西.
其實是不再適用於這個新天新地.
我們經常都會問.
今天我們追求的是什麼.
什麼是今天值得去紀念.
值得去珍惜的.
我這樣說.
我並不是在說地上所有的東西都是沒有意義.
我不是這個意思.
因為你還沒死.
到你死的那天的時候.
的確所有的東西都沒有意思.
是不是.
但我的意思就是.
當你在追求的過程當中.
請你不要將這些成為了你的上帝.
你不要為了追求祂的時候.
將你一切所有的資源.
將你的心思全部放在裡面.
我不管你是賺錢也好.
你是賺自己的理想也好.
你是賺你追求的任何東西也好.
請你記得.
你不要將所有的東西完全放在當中.
因為這些東西會過去的.
將來的新天新地裡面.
不再需要這些東西.
你能夠早點去醒覺的時候.

$^{281}$你就會懂得從神的角度去重新再看.
什麼將來帶入新天新地是有意義的.
我們今天在地上是一個過客.
我們最重要的是為未來去儲錢.
我們經常都有這種觀念.
弟兄姊妹我們儲了多少錢.
你在天上的帳戶裡有多少東西.
我們今天繼續的享受.
繼續的追求買到其他的東西.
是不是代表著我們能夠在未來.
我們進入神的角度裡面就會很開心.
我經常和大家說這個玩笑.
聖經裡面所形容的就是未來.
就是在天上和上帝天天唱歌.
你練好你的歌聲來.
你都不喜歡唱歌的時候.
你上去會不會悶死你.
你都不享受在地上和上帝敬拜的時間.
將來有一天叫你由早到晚對著上帝.
你還不死嗎.
所以弟兄姊妹我們準備好了沒有.
我們有沒有準備好.
我們在未來去迎接這個新天新地.
另外一個很特別的海也不再有年.
雖然這裡有很多爭論.
但起碼在我自己所學的裡面.
海的其中一個解釋就是邪惡.
他說新天新地裡面是不會再有邪惡的了.
邪惡不再存在了.
是一個完完全全潔淨的地方.
弟兄姊妹什麼叫做盼望.
無論你今天面對著一個怎樣的邪惡.
無論是環境無論是人無論是政權.
無論是任何東西都好.
如果你覺得那些是邪惡的話.
你放心吧.
有一天上帝的角度回來的時候.
這些東西不再存在了.
為什麼要和大家分享這些東西.
因為這些東西很多時候都牢籠著我們的心.

$^{321}$我們被這些東西綁得很緊.
我們把很多在世界上的價值觀.
以為是一個不可破滅的價值觀.
我們把它抓得很緊.
放不下手.
弟兄姊妹請你記得.
你今天你手不手抓住所有的東西.
有一天都會過去的.
是不是值得我們要花這麼多時間去抓住它.
其實每一天過去了就是過去了.
每一天明天都是新的.
所以我不知道大家有沒有這個習慣.
除了早上很趕急去起床.
梳洗上班之外.
當你每一天睜開眼睛的第一步.
我還有生命可以繼續下去.
大家有沒有想過.
睡完之後第二天就起不了床.
我這樣說不是灰心.
我是想說我們的人生本就是如此.
第二件很有趣的事情就是.
我又看見聖神耶路撒冷.
由神那裡從天而降準備好了.
就如辛苦裝飾整齊等候丈夫.
很有趣的地方就是.
當時其實對於一群亡國的人.
耶路撒冷的重生再有.
對他們來說是很重要的.
因為這是他們的盼望.
而這時候上帝也是按照他們的需要.
告訴他們新的耶路撒冷來了.
那裡代表著和上帝的關係.
代表著舊恩.
這是他們亡國的人.
一群人很希望回復當年的光輝歲月.
但是特別的地方.
他用了一個例子.
就好像一個新娘預備好了迎接丈夫.
這才是特別的地方.
弟兄姊妹.

$^{361}$當今天我們說等候上帝的盼望.
等候祂的末後這個時間.
你用什麼態度去等待.
你興不興奮.
你想不想上帝的國早日回來.
其實在這個世代當中我有點擔心.
當我和弟兄姊妹說的時候.
還是上帝不要住了.
我還沒玩夠.
你知道在這個時代.
我們要儲一筆首期也不容易.
如果沒有父母幫忙.
自己要儲一筆首期也不容易.
很辛苦.
每天餐餐捱麵包.
到你儲一筆首期.
明天就去交訂.
然後主耶穌說我今天就回來了.
那我之前在做什麼.
究竟我們心裡哪一樣是我們渴望的.
什麼叫上帝比一切都重要.
就在這裡看到.
我們今天有沒有一個新娘.
你記住這個新娘是什麼.
我渴望期盼著一份幸福.
能夠來到我們當中.
今天這個世代裡面同樣.
基督教永遠都會和你說一些最老土的價值觀.
到今天我仍然都會和年輕人斬釘截鐵.
聖經裡面就告訴你.
不要有分散性行為.
不要同居.
為什麼.
因為上帝就是想告訴我們.
有很多東西是很珍貴的.
很值得我們期待的.
你想想有些人用今天的價值觀.
我們同居.
然後時間很好.
試婚.

$^{401}$試試行不行.
不行就散了.
我想告訴大家.
所以今天的人在外面說這些話是正確的.
我們同居之後.
我們所有事情都做了.
只差那張紙.
我拿這張紙來做什麼.
綁手綁腳.
是不是.
不正確的到時.
一分就搞定了.
又不用離婚.
他們說是對的.
因為他們對於婚姻沒有期待.
他們對於人沒有信任.
他們純粹是覺得我悶.
我找個人在旁邊陪著我.
就是這麼多.
完全不是婚姻的本意.
婚姻裡面就是.
我們都很珍惜對方.
跟著時間.
我們很願意將我們最寶貴留給對方.
我們很期盼我們結婚的那一天.
我們兩個的關係是唯我獨有的.
不可以跟別人分享的.
我們的關係是建立在一個很恩典的裡面.
究竟我們今天對於在上帝的國裡面.
我們有沒有同樣的期盼.
我們有沒有一份這樣的渴望呢.
所以你就會發覺.
上帝在這裡的時候告訴我們.
原來這些是我們應該有的.
作為基督徒我們應該有的.
無論環境如何.
你都應該有的東西.
而到最後的時候就是什麼.
最寶貴最寶貴的地方就是.
我們永遠跟上帝都不再分開.

$^{441}$這個是我們終極的安慰.
他說我聽見大聲音從寶座出來.
說看神的障目在人間.
祂要與人同住.
他們要作祂的子民.
神要親自和我們同在.
作為他們的神.
神要擦去他們一切的眼淚.
不再有死亡.
也不再有悲哀 酷懼 疼痛.
因為以前的事都過去了.
你就會發覺.
原來上帝這個保證是.
上帝直接從祂告訴你.
從祂的寶座告訴你.
是祂的保證.
我說有就有.
我說會發生就會發生.
屬我的人有一天我會回來你面前.
我的國度是會充滿全地的.
祂說與我們同住.
大家即時可以聯想到的就是一家人.
我非常羨慕創世記.
上帝與亞當夏娃的關係.
是可以在旁邊的.
可以與你接觸 與你聊天.
沒有隔膜的.
我們是一家人.
祂要告訴你 你要做子民.
祂要做你的王.
做王的裡面就是祂要保護你.
和為你安排好.
說真的一個好的皇帝是很重要.
如果有讀歷史 我很喜歡讀歷史.
好的皇帝就是他真的會愛人民.
他會為人民的福祉去想.
他不會只是為了自己的利益.
但是當人做得不好的時候.
亦都可以變暴君.
所以問題就不是在於他是不是要做王.

$^{481}$而是他要做皇帝.
問題就不是在於他是不是要做王.
而是這是上帝 祂是愛你的.
祂要用愛來與你相處.
祂要保護你 祂要與你同在.
隨時隨地與你相伴.
你隨時隨地不用打電話.
你都可以找到祂.
最後我要做你們的神.
祂一向都不是我們的神嗎.
你要記住祂要成為你獨一的神.
如果你問我的話.
其實說真的 今天魔鬼最厲害的就是搞混亂.
你以為自己是獨一的神.
其實不是 我們還有很多很多的神.
我有金錢神 我有地產神.
我有名譽神 我有很多很多其他的神.
但是祂真的要成為你獨一的神.
其實這個原來主耶穌一早已經跟你說了.
在主禱文裡面大家記不記得第一句.
我們在天上的父 願人都尊你的名為聖.
願你的國降臨 願你的子兒行在地上.
如同行在天上.
這個豈不是主禱文裡面的圖畫.
所以你會發覺這個愛我們的神要與我們做甚麼.
祂要擦去我們的眼淚.
因為祂知道我們在地上所面對的不容易.
祂知道我們所受的苦.
所以在祂這個地方不再有死亡.
即是說不再絕望.
祂知道我們會受委屈.
祂知道我們會悲哀 祂知道我們會哭耗.
祂知道我們會有疾病 祂知道我們會有疼痛.
祂告訴你 你放心吧.
在這個地方裡面 以前的事一切都過去了.
兄姐妹 又回到剛才的主題.
原來盼望是告訴你一切的事都會過去.
所以為甚麼我們今天在一個最困難的地方裡面.
我們仍然可以有力量.
因為我們知道那些事情是短暫的.

$^{521}$那些事情是上帝看著的.
祂答應過你 你放心吧.
但倒過來說一個時間.
即是說你今天過的生活 以前的事.
帶來給你是甚麼呢.
是眼淚 絕望 悲哀 哭耗 疼痛.
你會發覺原來一個沒有上帝的世界.
原來就是這樣.
人很多時候做的是甚麼呢.
好心做壞事 或者壞心做壞事.
我以為我是愛你嗎.
我以為我是對你好嗎.
但原來不是 我一直在傷害你而不自知.
有些人甚至是為了心腸而害你.
所以我們要受盡很多的委屈.
我們經常都會有一個防人之心不可無.
我們都會有這些情況.
我們沒有辦法可以安心安定的下來.
去過我們的生活.
但上帝告訴你.
你把你的盼望放在我身上.
為甚麼呢 因為祂最後有一句.
因為神將一切都更新了.
祂所有的東西都更新了.
原來真正我們得到的盼望.
是因為一切在上帝裡面被更新.
弟兄姊妹 今天你相信耶穌.
你的生命被上帝更新了沒有.
你有新的價值觀嗎.
你有從神而來真正的平安嗎.
你在祂裡面.
我仍然過著我以前那種很懼怕.
那種認同的世界生活.
還是我過著上帝認可認為的生活.
所以你會發覺.
伊田經文到最後的時候.
祂說我是初 我是中.
由我們歲首到年中.
上帝都在.
在你出生到你離開那一天.

$^{561}$上帝都同樣在.
上帝從來都不會離開你的.
但祂一直在跟你說.
究竟你是否得性.
你是否過得性的生活.
如果你過得性的生活的話.
我就會將生命的泉水給你喝.
即是說你不再面對死亡.
不面對死亡.
除非你的生命被更新改變.
你能夠有從神而來的生活.
你的生命才會有盼望和出望.
所以弟兄姊妹你會發覺.
2022年過去了.
我們不會再追到什麼.
我們今天在座的時間.
起碼網上的弟兄姊妹都是.
我們能夠平平安安的來到2023.
你又有一年了.
你想怎樣過.
你想過一個更新的生命.
還是你想過以前的生命.
回到一個很實在的.
我們都是嘉欣樓的弟兄姊妹.
坦白說我們自己都很清楚.
嘉欣樓今年絕對不能過去年的生活.
因為我們要在神面前更新.
嘉欣樓才有它存在的價值.
所以弟兄姊妹.
有好幾次我上來分享.
我都跟大家說.
我們不是在乎究竟留不留.
不是這些問題.
如果我們這裡是一個被更新.
是一個有盼望.
是一個可以讓人得著恩典的地方.
上帝一定會讓我們留.
但如果我們是得過且過.
我們都是隨隨便便.
我要承認.

$^{601}$大家不可以不承認.
不要以為我經常說.
弟兄姊妹剛剛信完.
今早才教完初信班.
我就提醒弟兄姊妹.
請你不要回到教會.
將每一個基督徒的人.
幫他加一個光環在頭上.
不是的.
我們一不小心的時候.
我們就會過回以前的生活.
在地上的時候.
我都要承認.
有很多的教會.
我們稱為教會.
但我們其實在教會裡.
我們做的事.
是不是真的上帝所喜悅的呢.
不一定的.
所以反過來.
我們不需要批評別人.
但我們要問我們自己.
因為我們都是家恩流裡的一份子.
我們是奸家的一份子.
今日我們所做的事.
是不是被上帝更新改變.
我們都盼望.
我們每一個弟兄姊妹.
由我們的人.
變成我們的家庭.
變成我們的內心.
我們都能夠在上帝裡被更新改變.
我們能夠有一個完全.
讓人得到盼望的群體.
讓他們知道.
原來成為基督徒是真的幸福.
讓他們在這裡找到愛.
這才是上帝給我們.
最大的恩典和祝福.
最後和大家分享.

$^{641}$我們教會的傳統.
就是在最後的星期日.
在感受那邊.
就會有一個感恩分享會.
連續感恩分享會.
我們同宮有一個特別身份.
很多時候我們會做後備.
我們不知道有多少弟兄姊妹.
他們會分享或不分享.
到時肥沙通常幫我們做主持.
還有很多時間.
弟兄姊妹你上來分享.
我不可以吃這些雞.
當然要想好.
萬一他叫我上去分享.
分享什麼.
我會想好.
我那天沒機會分享.
很多弟兄姊妹.
特別是這次很多奸流的弟兄姊妹.
去分享.
我覺得整個童話很奇妙.
接下來的時間.
我去想.
我去分享我會說什麼.
我告訴大家.
如果我真的上去說.
在2022年.
對於我來說.
我真的不覺得是特別困難的一年.
因為整個過程裡.
我很清楚.
我真的.
在聖經裡面的教導.
我真的沒有看著.
這些困難去過我的日子.
我只知道前面有很多事.
我需要去做.
有很多人.
我需要去幫.

$^{681}$有很多事工.
我需要一起努力去處理.
所以上帝你交給我的.
我就去做.
所以後來我在現場.
我又回到原來我今年.
我的行事曆滿了.
但你問我我累不累.
我不累.
你問我.
我心裡會不會很擔憂.
我會很擔憂.
因為我覺得一切.
所有都是上帝的.
原來聖經裡面的教導.
是真實和活潑的.
你跟著去做.
你就能經歷到聖經裡面所說的這份力量.
求神幫助.
願意大家都能夠同樣.
帶著這份盼望.
這條出路去進入2023.
上帝在前面.
交給我們有什麼.
我們就做什麼.
也都求上帝.
令開我們的心.
讓我們看見.
在這個社區.
在我們富有的地方.
有什麼需要.
我們去回應你的需要.
我們同心祈禱.
我們去讚美你.
我們多謝你的保守和看顧.
在你面前的時候.
我們獻上我們的禱告.
因為你平安的帶領我們.
能夠經歷了.
渡過了這個恩典的時間.

$^{721}$你讓我們能夠帶著你的恩典.
一起走我們這條路.
將我們獻上感恩.
我們求你幫助.
無論在過去日子裡面.
我們經歷了什麼都好.
我們求主你給我們眼目.
我們仍然的定睛.
我們看見你.
因為你才是我們真正的盼望和出路.
將一切給我們交託.
我們獻上感恩.
祈禱奉靠主耶穌.
得成名字祈求.
我忘了.
\newpage



\section{路加福音 6:43-49}
\label{sec:d2MEojJVUd4}
\textbf{2023-01-01 午堂《趁新一年,遵主話而行》路加福音 6\_43-49 講員:曾裔貴牧師}
\newline
\newline
連結: \href{https://youtube.com/watch?v=d2MEojJVUd4}{\texttt{https://youtube.com/watch?v=d2MEojJVUd4}} ~~~~ 語音日期: 2023-01-03
\newline
\newline
\hyperref[sec:IlLnn7WOgiE]{\small{< < < PREV SERMON < < <}}
~
\hyperref[sec:index_chronic]{\small{[返順時目]}}
~
\hyperref[sec:index_scriptual]{\small{[返順卷目]}}
~
\hyperref[sec:BQHcR2q2KwA]{\small{> > > NEXT SERMON > > >}}
\newline
\newline
路加福音 6:43-49
\newline
\begin{longtable}{cl}
\hline
\hline
章節 & 經文 (和合本修訂版)\\
\hline
6:43 & \begin{tabularx}{0.7\textwidth}{X} 「沒有好樹結壞果子，也沒有壞樹結好果子。 \end{tabularx} \\ \\ \relax
6:44 & \begin{tabularx}{0.7\textwidth}{X} 每一種樹木可以從其果子看出來。人不是從荊棘上摘無花果的，也不是從蒺藜裡摘葡萄的。 \end{tabularx} \\ \\ \relax
6:45 & \begin{tabularx}{0.7\textwidth}{X} 善人從他心裡所存的善發出善來，惡人從他所存的惡發出惡來；因為心裡所充滿的，口裡就說出來。」 \end{tabularx} \\ \\ \relax
6:46 & \begin{tabularx}{0.7\textwidth}{X} 「你們為甚麼稱呼我『主啊，主啊』，卻不照我的話做呢？ \end{tabularx} \\ \\ \relax
6:47 & \begin{tabularx}{0.7\textwidth}{X} 凡到我這裡來，聽了我的話又去做的，我要告訴你們他像甚麼人： \end{tabularx} \\ \\ \relax
6:48 & \begin{tabularx}{0.7\textwidth}{X} 他像一個人蓋房子，把地挖深，將根基立在磐石上，到發大水的時候，水沖那房子，房子總不動搖，因為蓋造得好。 \end{tabularx} \\ \\ \relax
6:49 & \begin{tabularx}{0.7\textwidth}{X} 但聽了不去做的，就像一個人在土地上蓋房子，沒有根基，水一沖，立刻倒塌了，並且那房子損壞得很厲害。」 \end{tabularx} \\ \\
[1ex]
\hline
\hline
\end{longtable}
$^{1}$請姐妹主女們平安.
新年快樂.
更重要的是福杯滿日.
不過跟今天講的訊息.
這個福杯.
在神,在聖經,在耶穌真正看的福.
是什麼福杯呢?.
我們有沒有錯誤的理解.
表面地上的福是唯一的福呢?.
我們就要思想這個課題了.
跟大家有一個簡單禱告.
多謝主.
在新的一年.
1月1日.
主日.
我們可以朝見你.
敬拜你.
我們可以聆聽你的教誨.
尤其是主耶穌講了一個門徒.
聽話,聽神的道.
聽耶穌的教導.
將來的結局.
以及不聽的後果.
是相當清楚,簡單.
但是要實踐起來的時候.
就需要一份毅力,決心.
求主幫助我們.
我們這樣仰望,禱告.
奉耶穌基督的名.
阿們.
43至49節.
也很清楚是兩個的比喻.
43節提到善的人.
和惡人的分別.
46節就講.
如何可以做一個善的人.
繼續地去內心.
積蓄那個善.
以致去到永恆.
你都不會動搖.

$^{41}$第二個比喻就是講到一個聰明的公頭.
他做.
不單只是做外表.
人見得到的那個房屋.
更重要是先打好那個根基.
在那個盆石上面.
因為重點是會發大水.
是會有雨林,水涌.
來衝擊那個房子.
那個房子如果不在根基上面.
就會倒塌了.
主耶穌緊張聽道的門徒.
你會最終.
你的生命會倒塌.
主耶穌是很清楚.
在這個講道裡面.
苦口婆心的.
來勉勵聽道的人.
整個段落其實是由第17節開始的.
17節就說.
耶穌和他們下了山.
這個他們.
就是當時他在晚上.
在山上.
在眾多的門徒當中.
揀選了十二使徒.
揀選完了.
就和他們一起下山.
但是一下山之後.
聖經告訴我們.
主耶穌站在一塊平地上.
同站的有許多的門徒.
又有許多的百姓.
從猶太傳遞和耶路撒冷.
並推羅西頓的海邊來到去.
要聽耶穌講道.
又指望可以醫治他們的病.
還有一些是被污鬼纏磨的.
又纏又磨.
這樣.

$^{81}$都很慘.
另外.
又指望可以醫治他們的疾病.
接著第19節.
眾人都想要摸他.
因為有能力.
從他身上.
從主耶穌的身上.
發出來.
就可以醫好他們.
原來有許多的群眾.
大家可以想像一下.
昨晚特首出巡.
33萬800人.
只是在尖沙咀的文化中心一帶.
聚集33萬人.
對面岸都有3萬多人.
如果每個人都見到特首來摸摸他.
他有能力出來.
一摸就今年沒辦法了.
家庭和睦.
身體健康.
生意又順利.
每個人都想.
幸好大家都不是這樣瘋狂.
他和小朋友握握手.
就沒有其他了.
但這裡不是.
耶穌下山之後.
那些群眾是蜂擁來找耶穌.
即是說這是一個很好的契機.
現在有眾多的人.
是渴望要聽你的講道.
要得到你的祝福.
現在耶穌在這裡開始.
來教導了.
第20節.
耶穌舉目.
因為人多.
看著那些門徒.

$^{121}$就快要提到.
這裡有一個很特別的牌板.
好像一個牌比的對立.
開始講就先講和.
好像有兩個消息.
你想先聽壞消息還是好消息呢.
每個人都真的會.
心理作用.
你先講壞的.
也不是很壞.
接著有一個好消息.
耶穌怎樣呢.
耶穌在這裡說.
耶穌舉目看著門徒的時候.
你們貧窮的人有福了.
因為神的國是你們的.
你們飢餓的人有福了.
因為你們將要飽足.
你們哀哭的人有福了.
因為你們將要喜笑.
然後你馬上跳過去第24節.
但是你們富足的人.
剛剛相反.
剛才貧窮.
現在就富足.
貧窮就有福.
富足就有和.
這樣做一個牌比.
飽足的有和.
因為你們將要飢餓.
不是現在.
將來你就會飢餓了.
喜笑相對於剛才的哀哭.
喜笑的有禍了.
因為你們將要哀動哭泣.
說到三個的將要.
似乎在描述兩種的態度.
兩種的光景.
是眾多的門徒.
主耶穌是不是站在台上呢.

$^{161}$一早就設定好的呢.
似乎不是了.
可能站在小山丘上面.
高一點.
祂能夠看到眾人.
由推羅西頓.
耶路撒冷.
猶太傳遞.
這樣來蜂擁找耶穌.
相信他們真的有所求的.
耶穌就.
你喜歡聽道.
好啊 我就講道.
相信這裡主耶穌.
一定把握最重要的時機.
講述的一定是.
天國最重要的心靈特質.
所以叫做平原的白福.
就不是山上的補份.
這是平原的補份.
因為下山了.
在一個平地上.
這樣的一個落差.
你試想一下.
如果你在現場的聽眾.
又是觀眾.
你會怎樣理解.
耶穌講的這一種.
這麼奇妙的對立呢.
我期望的.
就是你能夠令我富足.
我期望的.
就是你能夠令到我生活.
我是平安.
而我能夠喜笑.
我能夠在我的飽足裡面.
即是我勤力做事.
我樂田耕種.
我能夠有好的收成.
不會有天災.

$^{201}$不會有旱災.
那我就會有飽足.
我的家人都有飽足.
為什麼會這樣講呢.
然後第26節.
和第22節的對比.
就會看到那個端倪了.
先講那個咒詛.
人都說你們好的時候.
豈不是每個人都想嗎.
誰不想別人說你好呢.
你做得好.
別人欣賞你讚賞你.
怎會是一個禍呢.
人都說你們好的時候.
你們就有禍.
因為他們的祖宗.
代那些假先知.
都是這樣.
原來假先知.
可以騙到一世的人.
騙到他那個時代的人.
講的是一般人喜歡聽的一般的事情.
所以假先知是有市場的.
假先知是不會在那代人.
被人揭穿的.
不過主耶穌說.
人都說你們好的時候.
你們就有禍了.
所以相信不是我們一般做人的好.
是我們有真理不來傳遞.
神的心意我們不講明.
是這一種的好和壞.
我們回頭看回22節.
人為人子恨惡你們.
拒絕你們.
辱罵你們.
棄掉你們的名.
以為是惡.
記住.

$^{241}$人為人子的緣故.
極盡的痛苦.
因為你們想想.
恨惡.
拒絕.
辱罵.
棄掉.
四個的形容.
你今年.
對不起.
去年.
有多少人.
你曾經這樣被人對待過呢.
辱罵.
恨惡.
拒絕.
棄掉.
是為人子的緣故.
這樣.
主耶穌說.
你們就有福了.
這樣是有福了.
接著有恩惠了.
當那日.
你們就要歡喜跳躍.
剛才提到的就是.
張耀.
四次提.
張耀張耀的和.
就會出現發生.
這裡提到就是.
到那日.
即是又告訴我們.
原來今生.
不是那麼明顯可以立刻看到的.
原來是要當那日.
就是主耶穌.
神國到降臨的時候.
就會明顯有這個真正的分野.
假的先知.

$^{281}$就有禍了.
這裡下面就是這樣提了.
因為你們在天上的賞賜.
是大的.
留意.
他們的祖宗代先知.
就是這樣了.
真的先知.
就是會被恨惡拒絕.
辱罵棄掉.
各位.
這一切.
主耶穌說的.
這一種福和的對比.
是一般人.
不可能認為這樣是和.
不可能認為這樣是福.
但耶穌就是這樣說.
就帶下去今天我們要說的.
最重要.
主耶穌要勸勉當時聽道的門徒.
要來到身體力行的.
積蓄我們自己的善.
在我們的心靈裡面.
因為第27節.
立即再補充.
現場人的反應.
只是我告訴你們這些聽道的人.
你們的仇敵要愛.
恨你們的要待他好.
咒詛你們的要為他祝福.
凌辱你們的要為他禱告.
似乎那個程度一路升級.
恨你們的你待他好.
仇敵已經要愛.
咒詛可能你們未必聽到.
在你背後才咒你.
不會在你面前.
曾銳貴我咒你.
沒有人這樣實聊交打.

$^{321}$通常在背後.
和一些你聽不到的人說.
去咒那個人.
一路升級.
升級到去到凌辱.
主耶穌怎樣說呢.
仇敵要愛.
恨你們的要待他好.
咒詛的要祝福.
凌辱的要為他禱告.
似乎那個人已經改變不了.
但你改變了.
你在主耶穌裡面有力量.
還要對準神去為他祈禱.
好像你在這個地方最強WIFI.
你直到天上的天庭.
你去為他祈禱.
不是為自己很慘很慘.
是為他很慘的祈禱.
這樣想的.
主耶穌教的那些人.
聽到的人.
你要這樣想.
兄弟姐妹.
如果我們常常都有這種學習.
不要看自己現在的狀況.
反而就是在主耶穌的眼裡.
看甚麼是真正的福.
真正的和.
和主耶穌怎樣去勸勉.
到底我們的心靈裡面.
我們內心的那個寶庫.
承載著甚麼呢.
感受堂在座.
在五堂的弟兄姊妹.
都是成年人了.
都是應該工作的能力.
都強過天崇還沒出來工作.
就算出來工作.
都是剛剛出來工作.

$^{361}$應該那個積蓄不會那麼多的.
相信在座都不用舉手投票了.
隨便可能三五十萬.
在你的儲蓄戶口一定有.
隨時有現金流可以運用.
這些是恩典.
是勤奮.
這些沒有對錯的.
不過在整大段聖經.
主耶穌對現實的財富.
不是很感興趣.
他感興趣的就是.
指示你們這些聽道的人.
你的內心的寶庫.
是存了多少呢.
你有沒有打保仔呢.
你的內心.
看看你的心靈裡面.
有多少的積蓄呢.
這樣的光景.
你能夠從另外一個角度去看.
你自己面對的遭遇.
這個超出一般人.
超出一般人想的.
你做初一 我做十五.
超出了我們這種的.
以牙還牙.
這樣的心態.
已經比一般以牙還牙.
超級了很多.
這樣的情操.
這樣的心靈特質.
是主耶穌指示我告訴你們.
這些聽道的人.
你們要這樣.
因為接著下去第43節.
就開始說.
最終你的內心的寶庫.
要盤點.
43節開始告訴我們.

$^{401}$因為沒有好樹結壞果子.
也沒有壞樹結好果子.
凡是樹木看果子.
就可以認出他來.
人不是從荊棘上摘無花果.
也不是從摘泥裡摘葡萄.
留意45節.
他說善的人.
在他的心裡傳一個善.
這個傳其實中文沒有翻譯.
英文其實是.
Treasure of the heart.
裡面有個寶庫.
內心裡面原來主耶穌說.
每個人都有一個戶口.
你傳的是善還是惡呢.
接著他告訴我們.
善人從他心裡所傳的善.
就發出善來.
惡人從他心裡所傳的惡.
一樣那個字.
又是Treasure of the heart.
不過是惡的存款.
在座應該不會有人滾爆卡.
欠幾十萬數.
這樣不會吧.
應該不會有這些負面的負債.
反而是正面的積蓄.
三幾十萬隨便都可以拿出來.
的儲蓄.
勤奮的積蓄.
主耶穌在這裡告訴我們.
心靈裡面都是一個存款戶口.
但是這個戶口.
我們在傳的是什麼呢.
其實是很寶貴的真理.
我們能夠學習.
我們的內心原來是有一個戶口.
這個戶口你在傳的是善.
因為在傳惡.

$^{441}$神給人類一個很大的.
我想只有人類才有的.
就是記憶.
可能動物都有的.
不過就很殘存的.
可能是一種條件反射的本能反應.
但是我們人類.
不單純是一種本能反應.
我們有思想.
我們有記憶.
我們有不同的期待的幻想.
動物界似乎不多的.
那就要留心了.
我們在傳很多.
別人對我們不好的那些的恨意.
那些的矛盾.
那些的一些的衝突.
我們繼續傳在裡面.
隨時就能夠拿出來的了.
還是沒有了.
完結了.
一筆勾銷了.
我們傳的就是善.
相信主耶穌在這裡很清楚的告訴我們.
我們必須要在內心.
時時看我們是不是一個善的人.
因為第46節就是第二個比喻.
43-45節就是第一個比喻說的善人.
和惡人的分野.
沒有說過結果的.
只說的是.
你是怎樣就是怎樣的.
你不可以欺瞞其他人的.
你是怎樣你就會怎樣了.
所以你不會在那個好處上面找那些擇泥的.
找不到的.
46節就說.
怎麼可能能夠將那個善繼續的積蓄.
去到有困難的時候.
我們都能夠仍然站得穩呢.

$^{481}$第二個比喻就是一個公投了.
不過相信相當簡單.
是聰明和愚笨的公投.
在現實世界.
沒有人會這樣打樁的.
當然了.
要偷工減料的那些令計.
天水圍的短樁事件都發生過.
所以在邏輯上.
是不會將這樣的房屋.
建造在普通的沙土上面.
而不打樁的.
以色列人他們已經離開了埃及.
不再在曠野搭營幕了.
進入耶路撒冷以色列地很多年了.
耶穌的年代已經很久了.
所以他們絕對不會建築.
是這樣隨便不打樁的.
所以這樣說.
是人盡皆知的.
小朋友都知道的.
不會這樣做的.
不過.
耶穌在這裡特別提到.
第47節.
「凡是到我者裡來.
聽見我的話就去行的.
我要告訴你們.
他到底像什麼人?」.
接著第49節.
「唯有聽見不去行的.
又像什麼人?」.
這樣做一個對比.
所以強調的是.
去到耶穌面前.
想聽耶穌說道.
其實剛才已經很多人了.
已經列舉了.
路加醫生.
告訴我們.

$^{521}$猶太傳遞,耶路撒冷,推羅西頓.
很多人.
眾多的需求.
很想摸耶穌的.
真的以他為一個很神聖的聖者.
但是現在耶穌在這裡.
說出一些他內心最真實.
期望他們會轉變的.
不是有求必應的.
摸一下你.
你摸一下我.
每個人都有求必應.
就快點回家.
開開心心.
過一個新年.
不是的.
耶穌在這裡.
說的就是真正的福杯滿日.
是你怎樣看待主耶穌的教導.
說了.
一個聰明的人.
聽了主耶穌的話就去走的.
他就好像一個人蓋房子.
不過他留意的是.
深深地去挖地.
對應剛才說的.
就是要不斷探索自己的內心.
是不是真的善.
如果有惡的.
盡快求主來清除.
尤其是上文說的是人際關係.
因為我沒有讀的一段聖經.
很大一段的聖經.
就是不可以論斷人.
這個我們懂了.
我們就很容易有我們的判斷.
甚至以為我們的判斷是最精準的.
主耶穌就說.
不是不可以判斷.
但是有一個條件.

$^{561}$就是你要將你自己的良目先拿走.
你才能夠看得清楚.
其實就是看得不清楚.
眼睛有良目.
其實就是誇大.
怎麼會眼睛有良目呢.
眼睛有刺都不行.
都要馬上去看醫生急症.
眼睛有刺.
一定是在說.
那個心靈的成見.
自己的前設.
自己的喜好.
喜歡的人你就很喜歡他.
不喜歡的人呢.
我還找不到你的錯.
挑骨頭的那個茶澡.
所以主耶穌說論斷就由這樣開始.
很容易就會出現了一種.
將自己扮演一種.
好像是名冊秋毫的審判官.
甚至提升到自己就是一個救世主.
這個是相當嚴重的問題.
所以這一大段聖經所說的人際關係.
說那個作門徒的心靈特質.
就是原來我們內心很需要有一個.
很重要的價值.
就是我們有主耶穌的話語.
來教導我們.
我們真的要實踐.
似乎是人盡皆知的.
要來打好根基.
殊不知主耶穌的教導就是說.
原來真的有一些人.
第49節說唯有聽見不去行的.
就像一個人在土地上蓋房子.
沒有根基.
水一沖.
隨即就倒塌.
並且那房子壞的很大.

$^{601}$整個大段落的講道.
一個現場的講道.
被路加醫生完全記錄下來.
就是來告訴我們當日的人.
和今天的信徒聽.
我們要遵主的話而行.
新的一年.
求頂只能夠我們一起.
來學習主的教導.
因為是有結果的.
將來我們要面對.
是否我們可以歡喜跳躍的日子.
還是將要哀哭有禍.
相信我們的心靈需要強健.
很多時候我們需要有一種另外的眼睛.
從永恆的角度去看.
早期的時候我說過.
有一個在2020年.
當選十大傑出青年的.
這位蕭海恩.
海船的恩典.
三歲證實眼睛患了癌症.
眼睛患了癌症.
要手術要拿開兩隻眼.
從此是永久的失明.
不過神給他一個恩賜.
就是對音樂的敏感.
盲的人要學音樂.
怎樣看譜.
那個艱難你都很難想像.
但是父母是很好的基督徒.
他在這樣的氛圍底下.
慢慢慢慢來建立他自己.
對別人的關愛.
就是這樣來祝福很多人.
當然不是今天我們每個人都要這樣坎坷.
才能夠成就一個很大的祝福.
但是主耶穌是對每一個聽祂道的人說.
就是我們要怎樣積蓄我們內心的善.
原來這個積蓄內心的善.

$^{641}$就是打根基.
就是聽耶穌的話.
越聽就越願意打.
就是說我們很難避免別人對我們的誤會.
但是你不要計較.
不要慢慢積蓄在你自己的寶庫.
將不好的不斷收起來.
好的就一直不要.
這樣就困難了.
最終每年的結尾.
你內心都裝著一些很不好的東西.
這樣就很難能夠成為一個真正善的人.
求主幫助我們能夠一起去學習.
還有一本書剛剛留意到.
那個女作家.
講述她人生的一些經歷際遇.
又是通常很困難.
又是在呼天搶地.
為什麼會淋到我身上.
其實那個機會率很低的.
突然間有一年.
丈夫遇到一個很嚴重的意外.
幾乎要死.
接著他的大兒子在越戰死亡.
他的第二個兒子.
有一天開車被一個醉駕的司機撞死.
第三個兒子.
因為性取向的緣故和父母決裂.
從此離家出走.
這個媽媽就.
你想一下 三個兒子.
自己的丈夫.
為什麼會這樣呢.
有一段時間非常徘徊在那種痛苦.
抑鬱裡面.
突然間有一天.
為什麼我要這樣去斷定自己的命運就是這樣呢.
接著完全改變.
開始學習去幫助其他人.
首先就是去幫助那些越戰陣亡的媽媽.

$^{681}$去跟他們聊天.
一起套心事.
那個苦水.
再幫助一些同樣有這樣困難的人.
慢慢慢慢他的內心越來越有陽光.
越來越有力量.
更加可以走出自己的困局.
這本書是美國暢銷書.
他說痛苦是很難避免的.
但命運的坎坷是可以選擇的.
有些事情是發生了沒辦法改變.
但你可以選擇不是這樣去活的.
求主幫我做我們能夠在新的一年.
我們不單純只是面對生活的困難.
而是我們的內心.
都能夠有一個很大的力量.
堅毅.
因為主的道成為我們大家內心的.
很重要的力量.
記得內心的寶庫.
在傳什麼是重要的.
我們一起禱告.
多謝主.
我們能夠將你的道.
將你的教訓.
再一次給錢免禮.
我們每一個都需要來聆聽.
我們將根基.
其實是我們打在耶穌裡面的根基.
是穩固.
以致當困難.
很多的事情.
是不如意的時候.
我們都有力量去面對.
幫助我們能夠學習.
真的要學習.
轉眼仰望耶穌.
而不只是定精在困難裡面.
我們這樣禱告.
奉耶穌基督的名.

$^{721}$阿們.
祝福各位.
\newpage



\section{}
\label{sec:BQHcR2q2KwA}
\textbf{2023-01-05 周四晚查經班《被擄歸回歷史研讀(2)》 講員:曾裔貴牧師}
\newline
\newline
連結: \href{https://youtube.com/watch?v=BQHcR2q2KwA}{\texttt{https://youtube.com/watch?v=BQHcR2q2KwA}} ~~~~ 語音日期: 2023-01-06
\newline
\newline
\hyperref[sec:d2MEojJVUd4]{\small{< < < PREV SERMON < < <}}
~
\hyperref[sec:index_chronic]{\small{[返順時目]}}
~
\hyperref[sec:index_scriptual]{\small{[返順卷目]}}
~
\hyperref[sec:gAWlXGsKQ90]{\small{> > > NEXT SERMON > > >}}
\newline
\newline
$^{1}$我們一起禱告.
多謝主給我們今天晚上開始第二課的歸回歷史的研讀.
我們會從以色列記開始看記載和背後的一個屬靈的含義.
都幫助我們明白裡面的經文和對這個時期是一個神感動猶太人回歸的心靈和心意.
都讓我們後世的人雖然真的和我們不是有直接關係.
但這也是醞釀將來耶路撒冷成為大軍的京城.
成為主耶穌降生的地方.
都讓我們看到和舊恩歷史的進行其實有息息相關.
幫助我們能夠在我們自己現在這個時代.
都為神的國度在我們一點點人生的時間去效力去準備.
求主帶領我們這樣禱告奉耶穌基督的名.
阿們.
我們開始第二課.
上個星期很簡單將歷史和大家交代了.
特別是以色列記一直到以斯帖和尼希米三卷書.
另外也有三卷的先知書.
哈該,撒迦利亞和馬拉基書.
就是被擄歸回三位的先知.
他們在那個時代就是在他們重建面對很大困難的時候.
先知就有神秘的說話來勸勉他們.
所以你又可以想像得到.
有70年的時間的真空期.
先知再一次出現來勸導重建的那些百姓.
特別是那些領袖主要是向所羅巴伯來作勸告.
勉勵他要剛強不要膽怯.
要來帶領自己的國民來建造神的殿.
所以你看到是相輔相成的.
有一些領袖響應回歸的呼召.
也有一些先知在他們面對困難阻力的時候.
成為他們的同路人.
也有神的話語來安慰勸勉.
尤其是撒迦利亞有相當多的寓言.
到時我們有機會可以看看這幾卷的先知書.
平時沒什麼機會看的.
看到歸回歷史就一定要看三卷的先知書.
我們由第一章開始看.
70年就是由606第一次的被擄.
主前的606年.
距離主耶穌還有600多年.
到70年有些人是536或者535.

$^{41}$主前的536或者535.
就是那個被擄的第一次的歸回.
就是以色列記頭六章的經文交代的時間.
也是古列王的元年.
波斯第一位的皇帝.
其實他類似打敗了馬代.
馬代波斯是有些相同的統治.
一個好像合宗的王.
國家有些合作.
擊敗了巴比倫.
馬代就取締了.
馬代又被波斯最終毀滅了.
所以現在來到波斯帝國的時候.
上個星期就和大家看了.
但以理書九章.
但以理的祈禱.
他知道那個應許.
即是說他很早期很年輕的時候.
他就已經有耶利米書.
他已經知道耶利米的寓言.
所以他就自己一路計算.
留在異鄉都一把年紀了.
去到都相當大年紀了.
他能夠看到神的應許.
很快就會兌現了.
所以他就去禱告.
似乎不是等到70年才祈禱.
而是他一直都來到去.
等候神的應許的兌現.
他一路來到去默默的.
忠心的在朝廷裡面辦公.
一個忠心耿耿的.
一個亡國的以色列人.
歷經四個王朝.
他都能夠去到古烈的時候.
他仍然都是健在.
相信都是已經是老人家一把年紀.
都知道自己離開世界.
他們不是回去主的懷裡.
因為他不知道誰是主的.

$^{81}$他就是回到亞伯拉罕的懷裡.
得到永恆的安息.
這個是他們所知道.
舊約對於永生的概念.
所以來到這裡的時候.
作者當然以以斯拉作為編輯者.
以斯拉其實是再過一點點的年日.
不是在這個時候出現的.
第七章開始就介紹以斯拉他自己.
不如我們先看看第七章.
是一個字述.
一到第六節.
這是以後波斯王阿達是西連間.
所以已經是一兩個朝代.
有個以斯拉.
他是西來亞的兒子.
西來亞是阿薩尼亞的兒子.
為什麼要歸宗認祖呢.
因為和他的背景有關.
因為他一定是和憲制有關.
有可能和祭司家族有關.
所以他要歸宗認祖.
自己是哪個支派.
是不是屬於祭司利微之派.
所以這樣來看自己的身份.
因為現在反過來是重建聖殿.
重建聖文.
重建聖城.
所以如果沒有身份是做不到的.
他是西來亞的兒子.
西來亞是阿薩尼亞的兒子.
阿薩尼亞是希臘家的兒子.
希臘家是沙隆的兒子.
沙隆是薩都的兒子.
薩都是阿希特的兒子.
阿希特是阿瑪利亞的兒子.
阿瑪利亞是阿薩利亞的兒子.
阿薩利亞是米拉約的兒子.
米拉約是西拉希亞的兒子.
西拉希亞是烏西的兒子.

$^{121}$烏西是布基的兒子.
布基是阿比舒的兒子.
阿比舒是菲尼哈的兒子.
我們比較熟悉.
菲尼哈我們知道是誰了.
民數記25章.
來將兩個人殺死的菲尼哈.
其實就是以利亞薩的兒子.
他是利微支派的.
而阿薩是大祭司亞倫的兒子.
我們明白了.
他為什麼要這樣一路追他的祖宗是誰呢.
原來就是亞倫.
就是說以斯拉是祭司家族.
他是可以一出生就是祭司的了.
等於新界原居民.
一南丁一到十八歲.
他就會有丁權了.
這個是歷史或者是法例.
以斯拉就是一個祭司.
他可以在聖殿裡面事奉.
現在他在波斯王朝.
留意第六節.
「這以斯拉從巴比倫上來.
他是敏捷的文士.
通達耶和華以色列神所賜摩西的律法書.
王允准他一切所求的.
是因耶和華他神的手幫助他」.
這樣描述自己.
是一個很敏捷的文士.
就是說是抄寫聖經.
教導聖經.
包括他是祭司.
他更加有責任去講解.
因為歷代志告訴我們.
大衛臨死之前設立了.
因為聖殿在大衛王朝.
當然是所羅門.
已經建造好了.
所以不需要再搬搬抬抬了.

$^{161}$所以大衛很聰明.
在他臨死之前.
雖然未見到聖殿.
但他已經設立了不同的分班制.
有師班有利未人去教聖經.
諸如此類.
給了男徒他的兒子所羅門.
所以不只是硬件建造聖殿.
而是已經幫助利未人.
因為已經將來建好聖殿之後.
因為不需要朝行晚策.
所以根本上沒有搬搬抬抬.
這麼麻煩的工作.
不需要又要安營又要賊營.
不需要了.
所以他們反而就是.
聖殿最重要就是.
那個中心就是.
神的兒女來獻祭.
和可能平日需要增重.
需要一些類似是.
擺平一些爭執的.
一些掌控的審判.
需要去到聖殿的範圍.
來處理的時候.
需要祭司的介入.
來評判.
或者有些人他剛剛關了戒倉.
有麻風.
又要隔離七天之後.
就要給大祭司去驗.
看他是不是已經發了.
已經是康復者.
就可以回家了.
這個就是古代的一個隔離政策.
所以都是全部在聖殿的範圍.
和利尾人去經手去發生的.
所以大衛都是相當聰明智慧.
已經設立了24班的利尾人.
所以是有輪班制.

$^{201}$所以也是讓所羅門.
可以當他建好聖殿之後.
就有一個規模.
有獻唱的.
有詩班的.
有守衛的.
有做其他的事務的.
也有可能是教導和抄寫聖經的.
因為王就職典禮.
就要抄一本聖經.
舊約的摩西五經.
來讓他自己平生的中毒.
所以按木團送聖經都是合聖經真理的.
一般就是按木送一本聖經.
這樣也對的.
張牧師送一本聖經給我.
我都保留到今天.
又寫了一些金句.
是重要的.
王是要自己抄寫一本摩西五經.
要平生的中毒的.
免得他心高氣傲.
有了實權之後.
就會忘記自己是一個普通的以色列人.
一定會的.
所以有一本聖經給他提醒和規範.
所以已經有這些崗位的.
所以當他們去到巴比倫之後.
沒有了聖殿.
沒有了憲制.
但是他們有民事的制度.
他們有這些抄寫聖經.
他們有自己的一個社區.
有他們的會堂.
這樣來中毒聖經.
和教導他們的兒女讀聖經.
猶太人就真的沒有文盲.
因為他們小時候必定要上主一學.
第一卷書就是要讀利未記.
所以他們從小就懂聖經.

$^{241}$還有很小的時候.
不知道幾歲就要站出來大庭廣眾讀經.
所以猶太人他們很尊重和體重神的話語.
這些都成為我們小朋友回家背誦聖經的傳統.
不要將人生最好的記憶背誦無謂的東西.
所以以色列是一個很敏捷的民事.
巴比倫可能已經是第二三代了.
他也是一個很敏捷的民事.
他是祭司的家族.
他通達耶和華以色列神所賜摩西的律法書.
就是說對摩西律法非常通透豁達.
不單止只是死背.
而是他能夠運籌維亞這樣去運用.
能夠足類旁通.
我不知道哪裡聽過.
有一個.
我在哪裡聽到呢?是真的.
就是一些穆斯林.
他是會背可蘭經.
背了很多次.
背得很熟.
但是我有印象了.
是誰跟我說.
背完之後.
我記得有一個校長.
那次來探我.
他說他的學校有一個學生.
他放學後.
是去穆斯林的清真寺.
做那些類似拉比.
教可蘭經.
他背誦如流整本可蘭經.
很難想像.
但原來他們是一知半解.
不敢去解釋.
某程度是不准去解釋.
要再高級一些的教長.
才可以有解釋權.
所以他們背很熟.
但原來不懂.

$^{281}$但以色列通達耶和華所賜摩西的律法書.
又是黃允准他一切所求.
因耶和華他神的手幫助他.
第一章我們會看.
是古列王頒布天下.
現在去到第七章.
亞達斯西王.
都是准許以色列.
回去聖城耶路撒冷.
來教導國民.
對神律法的遵守.
給大家一點印象.
以色列這位文士.
所以一至六章肯定不是他那個時代.
因為古列王的時候.
相差幾代.
幾個王朝.
但他可以做編輯者.
記錄這些仔細的資料.
有他們的歷史觀.
以色列就是編輯.
第七到第十章.
以色列記.
極多是以色列個人的字述.
這些類似自轉式的記載.
利希米充斥著自轉式的記載.
每一次說到某些東西.
就祈禱.
神啊求你施恩給我.
敗壞那些敵黨的人的計謀.
類似這樣的.
神啊求你紀念我在你面前的需要等等.
利希米充斥著自我向神呼求.
因為利希米記是以自轉的形式.
來寫日記.
所以第七章的以色列記.
就有這些近似日記的手法.
一至六章就不是了.
是一個歷史的技術.
我們回來看第一章.

$^{321}$我再讀多一次第一章.
留意這個叫切巴薩這個人.
就是第五章提到的那個所羅巴伯.
就是那個領袖所羅巴伯.
也是哈該書一開始的時候.
神感動哈該.
來勸勉這個所羅巴伯.
這個領袖.
叫他不要灰心.
因為聖殿被強迫停止了15年的重建.
一開始立了根基之後.
那些敵人就來了.
第四章會說.
來逼令他們停止.
怎樣逼令呢.
就是尋求王的釋法.
現在有這個人.
這一班人由波斯回來重建.
但是他們居心不良.
這樣來去透過這些上訪.
或者這些告狀.
就被王.
阿達什西就不知道什麼事.
就逼令他們停工.
一停就停到很後期.
這是第二章.
近期的經文告訴我們的事.
應該是第三章.
我們從第一章開始看起.
波斯王古烈元年.
耶和華為要應驗.
直耶利米口所說的話.
就激動波斯王古烈的心.
時得下召通告全國.
所以其實是所有被臣服的國家.
都可以回到他們的原居地.
為他們的神設立奠禦.
因為已經被巴比倫.
用強權全部毀滅了.
巴比倫也是將耶路撒冷聖殿拆毀.

$^{361}$他們是用這種高壓政策.
但是波斯就吸取了教訓.
國度這麼大.
你每次都完全高壓.
你就要派很多的駐軍去駐守.
這樣的做法會勞民傷財.
令到自己的帝國都會弱.
所以反而給他們有自己的自治.
給他們有一定程度的宗教上的自由.
這樣來就容易管理.
所以波斯的帝國就用這樣的懷柔政策.
所以當時猶太人都受惠於這個政策.
不過記得勝經告訴我們.
不是作者不是編輯的人.
不知道當時世界局勢的轉變.
波斯王朝採取柔和政策.
不是作者不知道.
而是作者要告訴讀者.
看到一個神的國度.
那個應許那個旨意的成就.
要從這個角度去看.
這一班被老歸回的神文.
這些選民他們很重要的身份.
和回去重建的神的旨意.
因為如果你做一些這麼重要的聖功.
你沒有一個這麼高的旨意.
慢慢就會遇到一些困難.
就會一盤散沙.
果然第四章就真的已經會有散沙.
但是更加重要就是.
要看到大家共同.
什麼原因感動你回來呢.
因為是神激動波斯帝國的君王古烈的心.
所以第二節.
波斯王古烈如此說.
耶和華天上的神已將天下萬國賜給我.
又祝福我在猶大的耶路撒冷.
為他建造電雨.
所以可能這個就是其中一個詔書.
是寫給猶太人的詔書.

$^{401}$其他的國家.
其他的地方.
譬如摩亞.
譬如其他地方的人.
他又會有另外的詔書.
所以這一份詔書是寫給猶太人的.
在你們中間.
凡作他子民的.
可以上猶大的耶路撒冷.
在耶路撒冷重建耶和華.
以色列神的電.
只有他是神.
願神與這人同在.
凡剩下的人無論寄居何處.
那地的人要用金銀財物.
生畜幫助他.
另外也要為耶路撒冷神的電.
用金心獻上禮物.
於是猶太和汴雅文的族長.
祭司利未人.
首先是族長.
祭司利未人.
就是一切被神激動他心的人.
有激動.
神激動他們的心.
終於等到回歸.
現在要來上耶路撒冷.
有心人就被神激動.
都起來要上耶路撒冷去.
建造耶和華的電.
他們四處的人就拿銀器.
金紙財物生畜.
珍寶幫助他們.
這個幫助可能就是.
出錢不出力.
自己有這麼大的生意.
這麼大的網絡.
已經在這個波斯帝國.
已經做著很大的事業.
就不上了.

$^{441}$可以出錢.
可以來出資源.
但是就不會出人.
自己不會上去了.
所以有相當多的人.
其實沒有回到耶路撒冷.
第七節.
古列王也將.
耶和華電的器皿拿出來.
這器皿是尼泊爾加尼撒從耶路撒冷掠來.
放在自己神之廟中的.
波斯王古列派副官米提利達.
將這器皿拿出來.
按數交給猶大的首領切巴薩.
這個切巴薩就是五章十四節的所羅巴伯.
器皿的數目記在下面.
金盤三十個.
銀盤一千個.
刀二十九把.
金碗三十個.
在聖殿裡面用的器皿.
銀碗之次的四百一十個.
這些器皿.
不可以第二個用途.
只能夠用在聖殿裡面.
用來獻祭.
用來有其他的用途.
別樣的器皿一千件.
金銀器皿共有五千四百件.
被擄的人從巴比倫上耶路撒冷的時候.
切巴薩將這一切都帶上來.
相當多的物件.
這些已經專門用於神用.
聖殿用的.
現在巴比倫他們拿了這些器皿.
放在他們自己的神.
就是大群的廟裡面.
但是不敢將它們融化.
因為古代人可能會很害怕這些神聖的物件.
如果融化成金或銀.

$^{481}$再重新鑄造第二件物件.
他們覺得是很大件事.
所以古代人對於神聖的物件.
他們會看得比較重視.
現代人就真的不多了.
還有就是我們看聖禮.
譬如浸禮.
我們不會自己跳進去當游泳.
我們是不敢的.
從來沒有人敢這樣想的.
浸禮叫小朋友跳進去.
沒有人敢這樣做.
因為知道那一刻是很神聖的.
那潭水充滿了之後.
真的有它的特殊用途.
就是下去浸泡洗禮.
歸入神明下.
聖餐一樣.
沒有人敢隨便拿個餅吃.
因為那一刻那個餅和杯已經神聖了.
所以暫時現代人.
還有這兩個聖禮.
是看得很重.
當然是教徒.
如果非教徒就更加不會有.
但他不會有.
他仍然尊重.
如果不是政府都不會說除了聖餐之外.
其他什麼都不准吃.
他都有這麼大的疫情.
都容許基督教一起脫口罩.
其實三百人脫口罩就很高風險.
但都給你進行.
可想而知.
連未信的人都很緊張.
你們神聖的一刻的聖禮.
所以如果反過來.
我們基督徒不尊重.
其實真的要打手板了.
人家不信都尊重你們.

$^{521}$所以巴比倫去到馬代.
現在來到波斯帝國.
這些器皿還原封不動.
可想而知古代人對於聖器皿.
是看得很緊要的.
第二章.
我們先停一停講第一章.
剛才都講過.
這是其中一封詔書.
是古列王特別寫給猶太人的.
這些肯定有他的招安.
懷柔政策.
你可以想像到.
古列王未必真的知道.
有耶利米蘇記載的預言.
而自己就成為了神的僕人.
去頒佈這個預令.
因為他們的內閣會議.
他們的行政會議.
只不過是想像.
現在我們新的王朝開始.
吸取了以前馬代和巴比倫的高壓.
令到很多被臣服的國民.
深深很大的報復.
或者是要反抗.
既然反抗.
不如准你們回去做你們自己的事.
你只需要進貢.
交差餉就已經解決了.
對於管治相對比較簡單很多.
給你有限度的宗教自由.
是最大的益處.
任何的國家.
你改變他的人民的身份.
那個文化.
或者那個宗教.
他就會淹沒了.
但是你不讓他這樣做.
他就會知道這是關乎滅族的.
所以他一定會反抗的.

$^{561}$但是如果你容許他.
大家就可以有一個.
有限度的空間.
一個換取的和平.
所以古列是一個很聰明的人.
但是聖經第一章已經告訴我們.
是神激動古列.
是神激動這些被老歸回的.
是49,697人.
加上瀑彼.
四萬多人是很少.
很少歸回的人.
所以第二章是值得他們每一個家族.
什麼的支派.
都記錄他們的名字.
當然是以族長做單位.
他們每個家族有多少人丁.
一起回去.
是說真的一起上去.
連根拔起地走.
那是很不容易的.
是不是60年代.
還是50年代.
50年代新中國成立.
呼籲全球的華人.
回去中國重建興建.
真的很多人這樣回去.
你可想而知.
那個熱誠那個心意.
是因為經過了二次大戰.
很多的欺凌.
由日俄戰爭一直開始.
到中日的抗戰.
已經受盡的欺凌.
中國人.
所以在外國的僑胞.
真的連根拔起地回去.
很多的印尼華僑.
很多在外國的人.
都是這樣響應回去.

$^{601}$去建立新中國.
所以猶太人是完全一樣.
所以他們被激動.
就起來回到耶路撒冷.
去建造神的殿.
也有些人.
四處的人就拿銀器.
金紙財物幫助他們.
還有一些金心獻的禮物.
就是說.
用錢去祝福這些.
這麼有心上去的人.
但自己就不走了.
自己就留在波斯帝國.
繼續他們的生活.
因為已經習慣了.
落地生根.
已經很難他們離開.
第二章.
這章相當悶.
因為全章只記載了.
這四萬九千多人.
巴比倫王尼布格尼薩.
從前老道巴比倫之猶大省的人.
現在他們是一個猶大省.
巴比倫裡面的.
類似中國城市.
中國人的地方.
猶大省的人.
現在他們的子孫.
從被老道之地.
回耶路撒冷和猶大.
各歸本城.
為什麼是猶大省呢?.
又講一點歷史.
和大家聊聊.
這是獵王幾下的十七章.
有些背景.
因為北面的以色列.
其實已經亡國了.

$^{641}$所以只有猶大和.
變雅敏茲派.
南面的猶大國.
還沒滅的.
所以最終被巴比倫毀滅.
就是南面的猶大.
猶大國.
歷代志下的十七章.
我不知道有多頁.
大家自己揭了.
因為我這本舊的和學本.
十七章的獵王幾下.
猶大王阿哈斯十二年.
以拉的兒子何世亞.
是北面的以色列最後一位王.
在撒馬利亞登基作以色列王九年.
他行耶和華眼中看為惡的事.
只是不像在他以前的以色列諸王.
亞述王撒曼爾息.
亞述已經是來到冰嵐城下.
要將北面的以色列毀滅.
他又不打猶大.
不知道為什麼.
來到這裡的時候.
攻擊何世亞.
何世亞就服侍他.
給他進貢.
何世亞背叛.
差人去見埃及王梭.
不照往年所行的.
與亞述王進貢.
亞述王知道了.
就把他鎖禁.
囚在監獄.
亞述王上來攻擊以色列遍地.
上到撒馬利亞.
圍困三年.
即是十二年.
何世亞第九年.
亞述王攻取了撒馬利亞.

$^{681}$將以色列人擄到亞述.
把他們安置在哈納與哥薩的哈博河邊.
並馬代人的城邑.
原來以色列人嚴格來說.
第一次的被擄.
就是北面的以色列王.
何世亞的年間.
其實如果是南國.
大概是希西加的年間.
因為猶大國有好的君王.
忠卿.
敬畏神.
所以還沒有真正的被毀滅.
但是北面的以色列.
就在這個時候.
亞述完全的毀滅了.
所以將北面的以色列人.
遷徙到其他地方.
第七節.
這是因以色列人得罪.
那令他們出埃及地.
脫離埃及王法老首的耶和華他們的神.
去敬畏別神.
隨從耶和華在他們面前.
所趕出外邦人的風俗.
和以色列諸王所立的條規.
以色列人暗中行不正的事.
違背耶和華他們的神.
在他們所有的城邑.
從瞭望樓直到堅固城.
建築休壇.
就是拜偶像.
燒香.
拜那些柱像.
就是什麼都拜.
激怒了天上的神.
所以神就透過亞述.
去罰這個以色列.
輾轉亞述王.
就是第24節.

$^{721}$亞述王從巴比倫古他亞瓦.
哈瑪和西法瓦欽而來.
安置在撒瑪利亞的城邑.
代替以色列人.
他們就得了撒瑪利亞住在其中.
留意這些人.
一直住在北面的以色列.
就是撒瑪利亞一帶地方.
你就可以想而知.
他們慢慢就變成.
北面以色列的原居民.
又和以色列當地的人通婚.
慢慢就成為了.
表面住在.
因為再過.
可能大概.
我要計算年份.
巴比倫滅猶太之後.
又是將真正的猶太人.
全部擄去巴比倫.
所以在那個時候.
原居民就變成這些.
亞述王遷徙來的外邦人.
所以當他們回來重建的時候.
這些原居民就提出.
不如我們一起去建造神的殿.
但是猶太人就說.
你們和我們是沒有份的.
因為你們不是猶太人.
你們只不過是被人遷徙來的.
有時那些原居民.
外姓那些.
都可以住但不可以投票.
都有類似這樣的情況.
大概大家知道這個背景.
和我們理解.
以色列記和尼希米記.
是有很大的關係的.
你就知道原來.
回去重建.

$^{761}$原來真正的阻撓.
就不是君王.
是在原居地.
以色列自己的國境的.
這些非猶太人的原居民.
因為他們已經是住了很久.
你就所想而知.
出現了這個問題.
就會引發第四章的反抗.
他們國歸本城.
這是第二章的以色列記.
第一節.
第二節.
他們是同著索羅巴伯.
就是這個領袖.
就是剛才說的那個叫切巴薩.
切巴薩有解經家說.
其實這個名字.
其實就是巴比倫的名字.
或者說波斯帝國.
可能那些官名.
就好像但以理.
但以理知道他的名字.
巴比倫的名字就叫做.
伯提沙薩.
但他就稱呼自己是但以理.
他和自己猶太人聊天.
他就叫自己是但以理.
他不會叫做伯提沙薩.
所以切巴薩其實就是索羅巴伯.
第二節.
他們是同著索羅巴伯.
耶穌亞.
尼希米.
記住這些名字是相同.
不是真正代表尼希米.
因為是一百二十多年之後才出現.
尼希米記的作者.
所以幾代之後才出現.
所以這個名字是很普通的名字.

$^{801}$尼希米西萊瓦利萊瓦.
沒底改.
沒底改又不要搞錯了.
是以斯帖的叔叔.
因為還沒有出現.
所以不要想著叫這個名字就以為是.
其實這個時候是古列王的元年.
以斯帖的時代是阿哈徐魯的時代.
是古列的.
我忘記了是曾孫還是孫.
沒底改.
必山米斯拔.
比格瓦伊利弘巴拿回來的.
以色列人民的數目記在下面.
這個似乎真的是一個文字記載的手筆.
很具體的哪個家族.
多少人回去了.
即是做了人口普查.
才准許他們回到自己的城市.
自己的地方.
巴祿的子孫2172名.
斯伐提瓦的子孫372名.
阿拉的子孫775名.
巴哈摩亞的後裔.
就是耶穌亞和約亞的子孫2812名.
以蘭的子孫1254名.
薩圖的子孫945名.
薩改的子孫760名.
你試一下將這些名套用現在錦宣家的家族.
你就已經有共鳴感了.
第三代第四代.
王坤勝的家族28名.
曾睿桂的家族後裔215名.
你可以想像到過了幾代之後.
這些現在是全家人.
由根拔起回去重建.
所以這麼巨細無遺的記載.
是包含了很具體的.
有血有肉的那份對神的愛.
對自己的故土的尊重.

$^{841}$渴慕日於言表的感情.
如果有幾十萬人沒有回去.
只有49600多人回去.
你可想而知這五萬人.
不是他的重要性.
可能是很普通的人.
但是被神的愛激動.
他們回去建造聖殿.
是將來耶穌基督在那個地方出現教導的地方.
所以他們當下回去.
其實對整個歷史.
他們可能一點塵都沒有.
那個影響.
但是當歷史一直進行的時候.
他們曾經回去用他們的金錢.
用他們的時間.
用他們有限的餘生.
來在耶路撒冷發揮一個貢獻.
被神透過以色列.
或者那些編輯的文士.
記錄在案.
被我們後世的人.
至少今天在錦宣家.
我們在這個晚上就讀他們的名字.
被神來紀念.
他們留名在那個生命冊.
在這個歷史的舊因的歷史裡面.
他們是有一份.
真的曾經有一份他們紀念過的.
他們回去曾經摸過那個聖殿.
曾經幫手來作那個根基.
都是有份於那個建造.
好像這個聖殿.
這個禮堂.
建的時候.
外面有很多那些板.
有些人在那裡畫畫.
跟著大家就來到去為新堂奉獻.
有些是捐獻很多的.
有些物件或者其他的東西.

$^{881}$認捐等等.
我當然不存在.
因為我未出現.
但是看回一些歷史.
看回一些文獻.
有些圖片.
就會見到這些事情.
所以是被紀念.
所以看家譜.
其實我們要看到.
即是來它的那個嚴謹性.
哪個支派回去是被記錄.
被記錄即是被神紀念.
亦都是給當時的人知道.
真是有一個很仔細的記錄.
真是有過一件這樣的歷史事件.
當然還有一個最重要的.
最重要的.
就是去到第二章尾.
有些人找不到他的家譜.
他又聲稱自己是利未人.
還未能夠斷定他的身份.
只能夠暫時不准他公職.
這樣.
即是說原來家譜就有這麼重要的含義.
不好意思.
你轉會嗎.
麻煩你給我一個洗禮證來看.
你都沒有洗禮證.
我怎麼幫你.
除非你可以找到那個教會.
你以前幫你洗禮的牧者.
你告訴我們.
原來遺失了.
施過的會的.
他證明了就沒有問題了.
這個就是.
真是需要尋找那個古籍的.
所以這裡也有它的作用.
所以大家千萬不要扔掉你的洗禮證.

$^{921}$當然我們有記錄的.
扔掉都不怕.
我們都可以找得到.
所以這裡就告訴我們.
有這些人名被紀念.
我們由第八節再讀起.
殺土的子孫.
945名.
殺改的子孫760名.
所以不齊齊齊的.
齊齊齊才可能有問題.
就是弄虛作假.
但是不齊齊齊.
證明他們的家族的繁衍.
有些多有些少.
是這樣的.
第十節.
巴黎的子孫642.
比拜的子孫623.
厄甲的子孫1222.
摩尼干的子孫666.
這些名字我們全部都沒聽過.
不是什麼大的家族.
名門望族.
比格瓦伊的子孫2056名.
阿丁的子孫454名.
阿特的後裔.
就是希西加的子孫98名.
相當少.
如果是皇族的.
喜塞的子孫323名.
約拉的子孫112名.
哈遜的子孫223名.
吉羅爾人95名.
伯利恆人.
接著說哪個地區.
什麼人.
例如東莞人.
按支派.
他是哪個地方人.

$^{961}$汴雅文的地方的人.
以法連山地的人.
或者是士劍的人.
這裡就是伯利恆.
即是猶大的.
123名.
尼陀法人56名.
阿拿特人128名.
河北的.
斯段的.
D段的.
類似這樣.
阿斯馬忽人42名.
基列耶林人.
希拉人.
比祿人.
共743名.
拉瑪人.
加巴人.
共621名.
麥瑪人.
122名.
伯特利人.
艾人.
那個艾城.
伯特利那個地方.
共223名.
尼泊山那個尼泊人.
52名.
莫必人.
156名.
別的爾蘭子孫.
1254名.
哈林的子孫.
320名.
羅德人.
哈第人.
阿羅人.
共725名.
耶利哥人.

$^{1001}$即是他們佔據的第一個地方.
345名.
西拿人.
3630名.
接著開始說那個職位了.
祭司.
耶穌亞加.
耶大亞的子孫.
973名.
欽默的孫.
1052名.
巴斯胡爾的子孫.
1247名.
哈林的子孫.
1017名.
利未人.
何達威亞的後裔.
就是耶穌亞和革滅的子孫.
74名.
歌唱的.
因為這個就是大衛設立的.
那個歌唱的師班長.
其中耶道頓.
希曼.
這兩個師班長.
他的後裔.
歌唱的.
阿薩.
第三個師班長.
阿薩的子孫.
128名.
接著有些守門.
那些叫尼提靈.
即是殿役.
守門的.
因為大家有印象.
讀約書亞記.
有一個民族.
叫基片.
他們很怕.

$^{1041}$以色列人.
怕他們會毀滅.
因為他們知道耶利哥.
和埃及都被他們滅了.
他們很怕.
其實住在以色列人附近.
就扮作很遠的來.
很遠地方住的人.
這樣來.
來和約書亞講和.
說他們很慘.
他們走了很遠的路.
很想來求和.
就騙了約書亞.
就被約書亞沒有防備.
和他們立了一個約.
就不可以毀滅他們.
但後來約書亞發現.
原來基片是一個很近的地方.
所以因為起了誓.
就不可以將他們毀滅.
所以怎麼辦呢.
就要他們一生一世.
做以色列人的.
伏役.
就在回望和將來的聖殿.
要做那個.
即是斬柴.
淡水.
這些粗重的工作.
他們就願意.
所以這些的基片人.
都跟隨以色列人.
離開埃及.
回去進入迦南.
這些守門的.
其實.
類似是已經是一些奴隸.
你如果在巴比倫.
在馬代波斯王朝.

$^{1081}$你已經脫離了猶太人.
其實你不需要回去.
但守門的沙龍的子孫.
阿達的子孫.
阿們的子孫.
阿谷的子孫.
哈底大的子孫.
素拜的子孫.
共139名.
接著就有尼提靈.
就是我剛才說的那些殿役.
其實真的不需要回去.
這些是奴隸.
尼提靈.
西哈的子孫.
哈蘇巴的子孫.
達巴俄的子孫.
基祿的子孫.
西亞的子孫.
巴頓的子孫.
尼巴拿的子孫.
哈加巴的子孫.
阿谷的子孫.
哈甲的子孫.
薩邁的子孫.
哈蘭的子孫.
吉特的子孫.
加哈的子孫.
尼亞的子孫.
尼遜的子孫.
尼哥大的子孫.
加薩的子孫.
烏薩的子孫.
巴西亞的子孫.
比賽的子孫.
厄拿的子孫.
米烏陵的子孫.
尼普森的子孫.
巴卜的子孫.
哈古巴的子孫.

$^{1121}$哈佛的子孫.
巴西律的子孫.
米希大的子孫.
哈沙的子孫.
巴赫的子孫.
西西拉的子孫.
達瑪的子孫.
尼西亞的子孫.
哈提法的子孫.
這一班人.
所羅門僕人的後裔.
就是所泰的子孫.
所肥烈的子孫.
比魯大的子孫.
瓦拉的子孫.
達昆的子孫.
吉德的子孫.
斯法提瓦的子孫.
哈提的子孫.
波克烈.
哈斯巴音的子孫.
阿米的子孫.
李緹玲和所羅門僕人的後裔.
嘩,只有三百九十多人.
但分散了很多的子孫.
即是說.
可能一個兩個這樣回歸.
你可以想像到.
都被紀念.
小妹都被紀念.
做奴隸的都被紀念.
回歸後.
你是做什麼的.
以前是做什麼的.
你的太公是做什麼的.
在聖殿門口看門的.
殿役都回歸了.
都重建了.
即是說都已經對神有這種熱切.
所以你可以想像到.

$^{1161}$每一個崗位.
只要你看那個崗位是尊貴.
神的事奉.
你看那個崗位是尊貴.
你的身份就會尊貴.
如果不是.
被官做了都沒有用.
如果你的態度不正確.
所以這些李緹玲.
感覺自己身份很尊貴.
我回到去能夠看門.
重建好聖殿之後.
我的子孫能夠看門.
我們都覺得是一件很榮耀的事.
這個其實很重要的邏輯.
和心靈的態度.
不要放大自己的內心.
慢慢好像沒有你不行.
其實是可以的.
聖殿都沒有了七十年敬拜.
完全沒有了.
都可以的.
所以最終就是.
聖殿在人的心裡.
聖靈.
主耶穌說.
不是在這個地方.
基里森山.
又不是在耶路撒冷.
神是個靈.
敬拜的要用心靈和誠實去敬拜.
在神的裡面.
在聖靈.
在真理裡面去敬拜神.
所以這個是很重要的一個.
一路的進程.
一路的演進.
最終就是神.
藉著耶穌來到這個世界.
聖靈就可以臨到每一個人身上.

$^{1201}$但當時他們真的需要回去重建這個聖殿.
要有神的同在.
在這個聖殿.
第58節.
尼提靈和所羅門卜人的後裔共392名.
接著第59節.
從特米拉.
特哈薩.
基祿.
厄丹.
欽墨上來的.
留意一下這裡.
不能夠指明他們宗族抱系.
是不是以色列人不是.
因為已經散居七十年了.
已經通婚了.
已經落地生根了.
已經可能說巴比倫語.
比說亞蘭文還要流利.
已經不認識了.
那你怎麼辦呢.
所以都不知道他們是不是以色列人.
但這些人都回去.
丹門士弟萊雅的子孫.
多比亞的子孫.
尼哥大的子孫.
就有652名.
這652名就被抽出來了.
因為他又說自己是.
祭司的家族.
可以事奉.
祭司中哈巴雅的子孫.
哈哥斯的子孫.
巴西萊的子孫.
因為他們的先祖.
娶了基列人巴西萊的女兒為妻.
所以起名叫巴西萊.
就是說當年通婚.
所以現在這三個家族.
第62節.

$^{1241}$這三家的人在族譜之中.
尋查自己的譜系.
卻尋不著.
因此就算為不潔.
不准供祭司的職任.
這個所謂的不潔.
就是祭司他們有那個身份.
是亞倫的後人.
是利美茲派.
和他們公職的時候.
是神規定了他們可以這樣公職.
和大衛王也都設立了一些師班長.
一些不同的崗位.
但都是要利美茲派去做的.
所以這個所謂的不潔.
就是他們的身份上.
未搞清楚那個族譜.
不可以亂來.
就是說他們都很認真看.
就是說原來之前那個家譜的重要.
就在這裡出現了.
不單止是鉅細無遺.
真的有一件歷史的事實.
再加上就是.
回到去的人是被神紀念.
擺在聖經裡面.
再加上原來他們都很嚴謹.
看一下你那個族譜.
你那個支派.
你那個身份.
是不是適合當一些的重建.
特別是聖殿的重建.
就需要有那個真正的利未人身份.
所以其實你知道教會.
有些比較重要的崗位.
我們都是要洗了禮.
是要教會的會友.
才可以做的.
譬如執事.
某程度是一些重要的主學老師的教導.

$^{1281}$堂董就更加要是執事.
做過執事.
為什麼呢.
因為有些崗位都是重要.
代表著教會.
有些決策等等.
你真的要對這個教會有個承擔.
是要有那個commitment.
所以你可以想像得到.
聖經都是這樣教導.
不是不知道他是不是以色列人.
不是說上去就算了.
不要理那麼多了.
剛才不是說了嗎.
四萬九千多而已.
不要那麼麻煩.
今天都沒有人事奉了.
隨便吧.
誰都做吧.
讓他做吧.
現在主席都不夠人手.
隨便吧.
他有心就讓他做吧.
不是不好的.
這個是好的.
不過更好的就是.
要大家都真的裝備好.
那個心靈是對準神.
和態度也是經過很多的試驗.
他真是很真誠的服事.
這樣就不會有時帶來很多問題.
聖經經常都說.
初入教的就不要做監督.
免得他自高自大.
很快就突然間很熱心.
可能先過一過冷河.
看看他是不是真的遇到一些事情.
仍然那麼熱心.
仍然都是那麼commit.
這樣去馴服.

$^{1321}$對神的家.
對神的僕人的順從.
對神的話語的服從.
有時需要觀察的.
太快被他吃重的位置.
其實有連累他.
最終連累教會.
他的眼光如果很細軸.
就帶來很多問題.
這些真是不是說得少的.
所以有沒有搞錯.
652個人你是不是歧視他們.
肯上來你都要這樣去瀏覽人家.
找不到那個家譜.
那怎麼決定呢.
可能真的漏寫.
或者是失傳了.
或者不知道什麼原因.
大火燒了.
那就沒有了.
第62節.
尋查不到他們的家譜.
63節.
省長就是索羅巴伯.
對他們說.
猶大的省長.
記住是在巴比倫.
在波斯帝國的猶大省.
對他們說.
不可以吃至聖丁物.
即是不可以守聖餐.
直到有用污名和土名.
決疑的祭司興起來.
這個就是古代的抽籤.
就是神會興起一些祭司.
這個叫決疑的胸牌.
不知道是什麼.
因為已經失傳了.
這個污名土名.
其實是什麼呢.

$^{1361}$不知道.
舊約有說的.
民數記祭司是有這件事的.
原來就是決疑.
現在有一個疑問.
好像你去政務處宣誓.
你不見了什麼重要文件.
你不就在政務官面前宣誓.
現在就簽發了一份真正的複製給你.
就是類似這樣.
所以現在不是不信任你.
這652人.
不過等到祭司就職.
他們有那個決疑的權柄.
就是說現在我代表神.
我們開始做這個小小的儀式.
你是什麼人.
就是我爸我媽的年代.
逃難來香港.
你說多少歲就多少歲.
信你的.
就立刻登記.
有時候報小一點報大一點.
就是上班工作.
個個都是.
因為根本沒有資料.
逃難來的.
人都沒有.
哪有身份證.
由那個時候開始劃線.
你說多少歲就當是.
類似這樣.
所以有祭司的興起之後.
那叫決疑的祭司.
就是說由祭司去做一個判斷.
是喲你真的說得出.
誰是你的家族的人.
就變成容易了.
現在譬如你在美國.
在紐約或者你在加拿大.

$^{1401}$譬如你在多倫多.
你是誰誰.
你說得出.
當地的僑領.
他們都會知道.
因為每年都有一齊聚會.
又或者他們會有些家譜等等.
所以不是不信.
都准你上來.
不過你暫時不可以吃潔淨的食物.
所謂致勝的食物.
就是他們憲制之後.
例如逾越節.
駐彭節.
第三章就說他們守住彭節.
所以你就可以想像得到.
他們是相當嚴謹的去看這件事.
雖然人少.
貴精不貴多.
但是仍然都是來去看重身份.
免得做一些事情是神所不起悅.
第64節.
會眾共有四萬二千三百六十.
此外還有他們的瀑被.
七千三百三十七.
又有歌唱的男女.
二百名.
所以總共是四九六九七.
五萬人都不夠.
當中有會眾.
有做奴隸做瀑被的.
做小妹的.
有些是歌唱的.
即是師班.
歌唱的男女二百名.
接著馬匹都計算了.
馬匹有七百三十六.
驢二百四十五.
駱駝四百三十五.
驢六千七百二十.

$^{1441}$驢和驢是不同的.
有些族長到了耶路撒冷.
耶和華殿的地方.
便為神的殿金心獻上禮物.
要重新建造.
他們量力捐入宮情婦的金紙.
六萬一千達尼克.
即是大家奉獻在神的裡面.
有一個宮情婦.
收集這些獻禮.
金心獻的禮物.
來重建聖殿.
銀就有五千尼拿.
似乎是波斯那些的.
度量行單位.
祭司的禮服一百件.
因為祭司要公職.
即是說祭司都要有輪班制.
於是祭司利未人.
民眾的一些人.
歌唱的守門的利提靈.
並以色列眾人.
一級一級.
有祭司.
有利未人.
有一些可能領袖.
接著有師班.
接著有守門.
那些守衛.
利提靈.
就是那些做搬搬抬抬的.
開始有了那個軟件.
就可以開始建造.
所以先有人.
有心的人.
你想想.
由波斯帝國.
聽到那個詔書.
省長索羅巴伯.
開始來招兵買馬.

$^{1481}$回歸了.
講全這個消息.
接著很多的家族.
收到這個消息.
就開始去思考了.
終於現在回到上去.
其實還沒有.
第三章才講他們.
回到上去之後的事.
這個是一個宏觀一點的描述.
現在回到上去.
就有接近五萬人.
現在他們回來重建聖殿.
所以是關乎到聖殿的職位.
他們就有齊人了.
就可以開始運作了.
當然還沒有建造聖殿.
就是先有人.
就祭司利未人.
民眾的一些人.
歌唱的守門的尼提靈.
所以每一個崗位.
神都看重的.
都不會說有些是特別重要.
有些是沒有那麼重要.
以色列的眾人.
接著就各在自己的城裡面居住了.
就是回到他們自己的地方.
就意味著先回到安頓好之後.
去到七月.
七月是猶太人最多節期的.
哪幾個節期呢.
七月初一就叫做崔國節.
七月初十就是他們的贖罪日.
接著七月十四.
那裡是守七天的.
就是他們那個叫做鑄棚節.
對不起.
七月初十就是贖罪日.
七月十四.

$^{1521}$其實應該是十五.
七月十五至二十一就是鑄棚節.
所以七月是三個大節期.
就在那個月份.
猶太人那個叫做.
糟了忘記了.
七月好像是彌撒月.
還是.
我要想一想.
我就不是很熟這個月份.
其實我沒有背.
彌撒月應該是四月.
就是他們出埃及的時間.
我忘記了是幾次流月還是.
還是什麼.
我要回去找一找.
七月有他們的日子的名字.
就是這樣.
兩張聖經交代了那個來龍去脈.
很鉅細無遺的.
用了一張聖經.
記載這四萬九千六百五十七人.
他們是什麼的支派.
是什麼的人.
是什麼的家族.
還有他們是做什麼的.
他們那個身份.
譬如他們是不是祭司家族.
是不是利微支派.
是不是其他的一些的離題靈人.
其實是外邦人.
是欺騙人.
諸如此類.
還有帶了多少的財物.
生畜上去.
真的讓我們知道是連根拔起的.
他們的家族.
一來就不上.
一上就決心上去.
就是和整個的聖殿.

$^{1561}$來到去.
同舟共濟.
就是生死與共.
這個就是一段時間.
這個就是主前的.
五百.
類似五百三十五.
我就看為五百三十六.
那個時間.
有這個這樣的舉措.
猶太人響應神的激動.
就回歸上去耶路撒冷.
停一停.
有沒有什麼想問一問.
八點八.
我們還有一點時間.
(記者:七月是第四).
謝謝.
有人幫我找了.
現在這些是上網已經找到的.
不過有時背一下會好些.
好.
謝謝.
第三章到了七月.
以色列人住在國城.
其實意味著.
如果從以色列記第七章說.
以色列用了五個月時間.
由巴比倫上去.
這個以色列.
就是耶路撒冷.
用了五個月時間.
當然有些土匪.
或者有些其他東西.
但是意味著.
不是七月才回到去.
是他們安頓好之後.
但是第一章和第二章.
就沒有說到他們怎樣困難上去.
父老攜幼的.

$^{1601}$不是一個短暫的守節.
而是移民回到耶路撒冷.
你可以想像這四萬多人.
是移居回到耶路撒冷的人.
到了七月.
以色列人住在國城.
那時他們如同一人.
聚集在耶路撒冷.
約撒達的兒子耶穌亞.
和他的弟兄仲濟斯.
這個就是大濟斯約書亞.
就是和所羅巴伯.
一起在這個時候.
宗教的民間領袖.
兩者結合.
來呼喚百姓.
來一起重建.
大濟斯約書亞.
並撒拉帖的兒子所羅巴伯.
與他的弟兄都起來建築.
以色列神的壇.
要照神人摩西律法書上所寫的.
在壇上獻梵祭.
獻梵祭每天早上都要.
其實黃昏要獻一次.
早上就要再獻一次.
獻梵祭.
代表著我們這個民族.
不斷來歸神.
梵祭就是英文叫做.
Hole Offering.
全生的祭.
生就是生口的生.
全生祭就是完全焚燒.
整隻獻上給神.
這叫做梵祭.
有些祭甚至贖罪祭.
是一撿.
如果普通的贖罪祭.
只獻內臟.

$^{1641}$Kidney.
肝臟.
和一些其他的脂油.
就可以了.
還有灑血.
整隻祭司可以吃.
贖罪祭可以吃.
平安祭就更加.
獻祭的都可以吃.
但梵祭.
誰都不可以吃.
當然神都不會吃.
因為已經燒了.
但梵祭就真的要這樣獻.
整隻獻上.
唯一祭司可以得到的.
下蘭的Bonus.
就是梵祭.
祭生的皮是歸祭司的.
歸大祭司的.
譬如你獻一隻牛.
那塊牛皮.
羊皮就歸祭司的了.
就不需要燒的.
其他贖罪祭.
平安祭.
甚至是掃祭.
贖獻祭.
都可以是祭司可以吃的.
因為每次都是獻一汁.
就算是那些麵粉.
都是祭司拿一汁這樣燒.
其實大部份.
一盤的麵粉.
都是歸祭司的.
所以祭司.
很多時候都可以吃到.
那些獻祭的.
譬如那些麵.
或者那些初甲.

$^{1681}$班溝.
牛羊.
這些全部都是祭司的下蘭.
所以他們要獻梵祭.
他們在原有的根基上.
祝壇.
如果真的沒有搞錯.
猶太人是這樣理解.
去過聖地的就知道.
現在全世界.
原來不是第五大.
剛剛收到的消息.
昨天看新聞.
因為以色列有一個總理.
去參拜.
去祈禱.
上祈禱山祈禱.
又是很嚴重的.
引起衝突.
現在在聖殿山.
那個伊斯蘭的清真寺.
是全世界第三大.
猶太人這樣理解.
又似乎是.
見到那些城牆.
所羅門的城牆.
又似乎真的.
在他的下面.
就是聖殿的位置.
所以.
為什麼他們這麼緊張.
要在聖殿山祈禱.
因為猶太人的觀念.
那裡是直接對著.
天上的聖殿.
Wi-Fi是最強的.
5G最強.
直上天庭.
所以一定要在聖殿山祈禱.
所以你不要怪猶太人.

$^{1721}$真的會這樣想.
我們今天.
耶穌說我們.
不需要基里森山.
又不需要去耶路撒冷.
但不代表我們.
不回來禮拜堂祈禱.
我們又不要變成.
完全沒有這種觀念.
有時真的要分別為聖.
這樣又不同.
沒錯.
Anytime Anywhere都可以.
即是你在病床.
你一路來到去要休養.
你都可以祈禱.
各方面.
甚至乎現在科技.
可以在現場一起參與祈禱.
是完全沒有問題.
但不代表我們有一個神聖的地方.
我們一起來約定一個時間.
大家一起公禱.
那個重要性.
今天基督教就失落了.
這個神聖的觀念.
又是奸教會.
你怎樣呼籲.
他都不回來.
他都覺得不關他的事.
是不是真的每間公司.
都是要八點才下班.
八點才准你離開.
又不是.
如果你六點鐘下班.
你那星期的一天.
趕回家.
都可能可以.
我又不評論.
因為每個人都不同.

$^{1761}$免得一高打一旋.
人就好像不太好意思.
但問題就是.
他們真的在這裡.
來到去很緊張要獻祭.
現在就在這個原址.
重建.
所以他們猶太人.
很想將這個第三大的清真寺.
來到去拆了和炸了.
當然就是世界大戰.
但他們就知道.
下面真的有那個聖殿.
甚至他們居民說.
那個藥櫃都好像找到了.
所以.
真是一個很大的conflict.
但是現在就告訴我們.
這個聖殿就是耶穌的聖殿.
不過耶穌的時代.
是希律王再重修.
再修葺和擴大.
用了46年.
但是這個聖殿.
第二個聖殿.
就是這個猶太人.
要回去重建的這個聖殿.
就是耶穌時代的.
這個重建的聖殿.
當然去到主後70年.
又拆了.
又沒有了.
到現在都沒有了.
所以他們現在就要在這裡.
在原有的根基上.
來去建這個聖殿.
第三節.
他們在原有的根基上.
祝壇因懼怕鄰國的民.
就是剛才我所說的.

$^{1801}$列王記下的十七章.
那裡很重要的.
那張聖經.
就是這個亞述王.
這個撒曼爾式.
來將那些外族人.
遷徙到北面的以色列.
所以這些變成鄰國的民.
猶太人這樣看.
這些其實住在他們附近.
就看他們是鄰國的人.
不是猶太人.
又在其上向耶和華.
早晚獻梵濟.
沒錯.
早晚梵濟.
就是說其實無時無刻.
我們的生命都是屬神的.
我們要歸神的.
我們第一時間要想神的事.
又照律法書上.
所有弟兄姊妹其實很好.
包括我們同工.
有時沒辦法.
當值我們.
通常我們去旅行.
都是想星期一.
最遲去到星期六.
一般星期五回來.
星期六都可以休息.
星期就要事奉.
有時都真的會這樣想.
這樣安排.
其他弟兄姊妹都是這樣想.
你想自己的計劃.
都想想我回崇拜.
事奉.
或者可以早點回來.
癱透一下.
星期日回崇拜.

$^{1841}$這個不是一個律法主義.
雖然不是律法主義.
但有些事是你自己安排的.
你可不可以選擇呢.
你是不是要去到深夜.
星期六晚機回來.
三點凌晨才下機.
星期就睡到黃絲白眼.
又起不了床.
沒問題.
有時日積月累.
看你的心對神的那件事.
有時我們都.
如果不是很緊張.
我的小朋友.
星期一一定要上學.
星期不要太晚上學.
找機位都要.
情願貴一點.
都回來要讓他睡.
為什麼.
因為第二天他要上學.
有時你想想.
上學.
回崇拜就不要緊.
有時這樣的比較.
有時日積月累.
你看到那個人.
他背後內心的價值觀.
其實這個不好.
重要是我們每個人都有.
我們自己對神的心意.
記得.
去崇拜不是為誰.
去崇拜是為神.
我敬拜神.
所以我去不了崇拜.
都是因為神暫時讓我不能去.
要搞清楚這些價值觀.
你就會放心.

$^{1881}$我們不是律法主義.
不是你少上一次.
就扣你五分.
不是這樣扣分制.
其後獻常獻的凡濟.
常獻的凡濟.
即是要常常想起神.
不是早晚的凡濟這麼簡單.
常獻的凡濟.
安息日會特別獻多些.
然後月數和耶和華.
一切的聖節期都要獻濟.
即是逾越節.
然後又開始到五旬節.
然後一直這樣數下去.
又向耶和華獻個人的金心濟.
金心濟即是你求恩.
感恩.
你就可以想像到.
基本的就規定了.
金心的就是你自己.
發自內心的.
兩種都重要.
所以不只是基本的一個.
你的責任.
而是你發自內心的.
第六節.
從七月初一.
他們就向耶和華獻凡濟.
但耶和華奠的根基尚未立定.
他們有張銀紙.
吸石像,木像.
這個是現在重建.
就是需要有這些工人.
把糧食走油.
吸西頓人.
推羅人.
因為香柏木.
這個推羅西頓.
這個叫做尼巴嫩.

$^{1921}$的香柏木是出名的.
即是不會住的.
給錢他們買這些香柏木.
當年所羅門也是.
他們就在山上斬了之後.
就浮海沿著這個敘利亞.
這樣就浮下來.
來到這個.
沒錯了.
從尼巴嫩運河.
運到海裡.
浮海運到約柏.
一直這樣流下來.
是按波斯王古列所允許的.
就是這個開始興工建造了.
所以剛才用了兩張聖經.
交代這個回歸的一班人.
他們的心意.
他們的決心.
他們可以被紀念的.
全部是因為我們不是猶太人.
我們不是太知他們的家族.
和在史記列王記上下.
歷代治下記錄的家譜.
是記那些重要的族.
就不是記這些.
這些是普通的平民百姓.
不是什麼名門望族.
所以你可以想像得到.
原來這次的回歸.
都是有不同的組成部分.
當然有祭司.
但是更多是一些.
不是太喜眼的家族.
他們想去重建.
就是說回去定居耶路撒冷.
他們這些是先行者.
其實那個決心很大.
我們又不是太知道.
可能我們要參考.

$^{1961}$坊間的歷史書.
到底在兩約中間.
猶太人他們的生活情況是怎樣.
如果我們看以司帖記.
有時就是那個王會忠義.
比如以司帖皇后得寵.
就會被抬舉那個民族.
但是有時會有些蠶神.
有些哈曼就很敵視.
沒抵該.
就想嫁禍猶太民族.
想將他們毀滅.
就會徘徊在這樣的情況.
就是說他們亡國之後.
有時會高有時會低的身份.
中國人一樣.
你去外國.
有些地方是很高的國民身份.
很尊重中國人.
有些地方就很歧視很排華.
都有.
比如非洲現在就對中國人.
就會很有善.
幫他們興建橋.
整路.
很多基建.
對中國人很有善.
但有些就不是.
所以猶太人一樣.
除非我們有很具體的資料.
知道在巴比倫馬代的年間.
到底猶太人他們的處境是怎樣.
他們是很高的國民身份.
很高的社會地位.
很高的商界發展.
回到去的人可能就是很大決心.
將所有東西全部拿上去的一班人.
所以這個其實就是.
聖經留給我們一個空白.
我們去想.

$^{2001}$我們也有一個共鳴感.
神的聖經有時.
每一個時代都有你那個部分的重大參與.
你不要看其他人.
你自己可以參與多少就參與多少.
神看的是那個心靈.
我們來盡我們的力.
人力幫手也是一個力.
也是一個奉獻.
金錢固然是.
人力也是.
心意也是.
所以這個也是值得我們去紀念.
我們的名字有沒有被紀念呢.
我說不是啊.
我們整個師段就是一起服侍的.
現在我們大團小組.
上次那個年終感恩會.
差不多整個團體來分享.
好像類似哪個人.
以蘭人就多少人.
其實是很好的.
E段整段就來了.
K段就是隔離.
K段整段會來幫手.
譬如我們快開茶水檔了.
如果疫情一直減緩的時候.
有整個團體來幫手.
你就會知道.
聖經就是這樣.
原來神都會被紀念.
這些人都會在神的眼中被紀念.
這些就是很重要.
我們都需要求主讓我們能夠.
有我們自己的名字.
甚至我們有一班人.
一起去參與聖功.
今天就說到這裡.
我們先不開第三章.
我們下一次就會由第八節開始提.

$^{2041}$開始就進入那個conflict.
你一做聖功.
鄰國的人.
就是住在那個地方.
70年的原居民.
就會來攔阻你興建聖殿.
因為你一興建聖殿.
就會涉及到他們的宗教.
我們祈禱.
我們有第二堂的.
這個被擄歸回的研讀.
我們很盼望第一章第二章.
這些決心重建的人.
我們根本從來不認識他們.
這四萬九千多人.
我們一個都不認識.
但是在神的眼中.
在歷史裡面記下一筆.
紀念他們的回歸.
他們的參與.
他們的綿力.
這樣來興建好神的殿.
成為日後.
無論是第二批的以色列.
第三批的尼欽米.
他們在這個根基上.
來建造的.
無論是聖城.
城牆.
或者是其他的國民的認同.
在宗教上.
直至到兩約完結.
主耶穌的降臨.
施洗約翰的出現.
和耶路撒冷的宗教氛圍.
都是因為有聖殿.
而來有這個敬拜.
也有瑪加比時代.
那個重光節的典故.
我們就知道.

$^{2081}$聖殿的重建是很重要的.
也是神激動那個時代的人.
我們很盼望主.
在今天21世紀.
求你激動基督教的弟兄姊妹.
無論是香港.
或者是不同的地方.
譬如澳門.
或者是一些比較難有福音的.
譬如國內.
都還不是全部自由的.
盼望神自己都來感動.
讓更多人可以起來.
為神的福音來禱告效力.
使到神的工作來開展.
我們仰望.
我們禱告.
奉主耶穌基督的名.
阿們.
\newpage



\section{路加福音 8:1-15}
\label{sec:gAWlXGsKQ90}
\textbf{2023-01-07 晚堂《存忍耐的心結出果實》路加福音 8\_1-15 講員:曾裔貴牧師}
\newline
\newline
連結: \href{https://youtube.com/watch?v=gAWlXGsKQ90}{\texttt{https://youtube.com/watch?v=gAWlXGsKQ90}} ~~~~ 語音日期: 2023-01-10
\newline
\newline
\hyperref[sec:BQHcR2q2KwA]{\small{< < < PREV SERMON < < <}}
~
\hyperref[sec:index_chronic]{\small{[返順時目]}}
~
\hyperref[sec:index_scriptual]{\small{[返順卷目]}}
~
\hyperref[sec:ejP2s680TsQ]{\small{> > > NEXT SERMON > > >}}
\newline
\newline
路加福音 8:1-15
\newline
\begin{longtable}{cl}
\hline
\hline
章節 & 經文 (和合本修訂版)\\
\hline
8:1 & \begin{tabularx}{0.7\textwidth}{X} 過了不久，耶穌周遊各城各鄉傳道，宣講神國的福音。和他同去的有十二個使徒， \end{tabularx} \\ \\ \relax
8:2 & \begin{tabularx}{0.7\textwidth}{X} 還有曾被邪靈所附，被疾病所纏，而已經治好的幾個婦女，其中有稱為抹大拉的馬利亞，曾有七個鬼從她身上趕出來， \end{tabularx} \\ \\ \relax
8:3 & \begin{tabularx}{0.7\textwidth}{X} 又有希律的管家苦撒的妻子約亞拿，和蘇撒拿以及好些別的婦女，她們都是用自己的財物供給耶穌和使徒。 \end{tabularx} \\ \\ \relax
8:4 & \begin{tabularx}{0.7\textwidth}{X} 當一大群人聚集，又有人從各城裡出來見耶穌的時候，耶穌用比喻說： \end{tabularx} \\ \\ \relax
8:5 & \begin{tabularx}{0.7\textwidth}{X} 「有一個撒種的出去撒種。他撒的時候，有的落在路旁，被人踐踏，天上的飛鳥又來把它吃掉了。 \end{tabularx} \\ \\ \relax
8:6 & \begin{tabularx}{0.7\textwidth}{X} 有的落在磐石上，一出來就枯乾了，因為得不著滋潤。 \end{tabularx} \\ \\ \relax
8:7 & \begin{tabularx}{0.7\textwidth}{X} 有的落在荊棘裡，荊棘跟它一同生長，把它擠住了。 \end{tabularx} \\ \\ \relax
8:8 & \begin{tabularx}{0.7\textwidth}{X} 又有的落在好土裡，生長起來，結實百倍。」耶穌說完這些話，大聲說：「有耳可聽的，就應當聽！」 \end{tabularx} \\ \\ \relax
8:9 & \begin{tabularx}{0.7\textwidth}{X} 門徒問耶穌這比喻是甚麼意思。 \end{tabularx} \\ \\ \relax
8:10 & \begin{tabularx}{0.7\textwidth}{X} 他說：「神國的奧祕只讓你們知道，至於別人，就用比喻，要他們看也看不見，聽也不明白。」 \end{tabularx} \\ \\ \relax
8:11 & \begin{tabularx}{0.7\textwidth}{X} 「這比喻是這樣的：種子就是神的道。 \end{tabularx} \\ \\ \relax
8:12 & \begin{tabularx}{0.7\textwidth}{X} 那些在路旁的，就是人聽了道，隨後魔鬼來，從他們心裡把道奪去，以免他們信了得救。 \end{tabularx} \\ \\ \relax
8:13 & \begin{tabularx}{0.7\textwidth}{X} 那些在磐石上的，就是人聽道，歡喜領受，但沒有根，不過暫時相信，等到碰上試煉就退後了。 \end{tabularx} \\ \\ \relax
8:14 & \begin{tabularx}{0.7\textwidth}{X} 那落在荊棘裡的，就是人聽了道，走開以後，被今生的憂慮、錢財、宴樂擠住了，結不出成熟的子粒來。 \end{tabularx} \\ \\ \relax
8:15 & \begin{tabularx}{0.7\textwidth}{X} 那落在好土裡的，就是人聽了道，並用純真善良的心持守它，耐心等候結果實。」 \end{tabularx} \\ \\
[1ex]
\hline
\hline
\end{longtable}
$^{1}$請主平安.
我們再一次祈禱.
多謝主.
讓我們能夠.
打開聖經.
來聆聽你自己的教導.
這是千古不變的.
一個很重要的比喻.
來開啟我們的人生.
尤其是我們的內心.
應該怎樣來.
去持守真道.
讓真道在我們心靈裡面.
很重要的主耶穌的教導.
是要生命能夠結實.
怎樣結實.
是一個很重要的人生課題.
求主.
藉著在年初的時候.
當我們去講解這些比喻.
主的聖靈在我們大家.
每一位弟兄姊妹的心靈裡面.
來到正正就好像殺種那樣.
能夠我們大家都是一個好的心田.
去盛載神的道.
我們這樣祈禱.
奉耶穌基督的名.
阿們.
我們在開始的時候.
我們都選取了耶穌很重要的幾個比喻.
作為我們這一季的課題.
今天我們要講的.
相當簡單.
就是有四種的心田.
唯有好的心田.
怎樣好的呢.
耶穌有解釋的.
就是持守在誠實善良的心靈裡面.
那個道種就能夠忍耐著結實.
有好幾個形容詞.

$^{41}$將人的內心去剖析出來.
其他三種.
就是道都是這樣落在他的心靈裡面.
但是不可以結到果.
有些是一開始就看到結不到.
有些是以為結到.
但後來都看到是不行的.
唯有第四種的心.
來龍去脈.
作者路加都給我們很清楚背景.
就是耶穌開始公開的去傳道.
周遊各城的傳道.
就是說在比較早期的時候.
公開向猶太人去呼籲神的國.
他們要進入.
就來去大肆的帶著門徒.
因為已經揀選了十二使徒.
就來去開始傳道.
作者就很有興趣.
除了使徒以外.
原來有些女徒.
女性的門徒.
跟隨耶穌.
不過用不同的方式去跟隨.
還有說他們的背景.
為什麼作者要這樣用意呢.
有沒有什麼特別的原因呢.
當然如果熟路加福音.
你就會知道.
路加福音很多時候將那些小人物.
社會上可能被邊緣化的人物.
會來台舉.
凸顯出來.
告訴我們.
原來這位路加醫生.
很有興趣那些真心跟隨耶穌的人.
他們的背景.
他們的情況.
這裡也不例外.
在這張聖經裡.

$^{81}$當然我們今天只能看到第十五節.
但是往後還有一個段落.
也是凸顯了聽道的重要.
就是第十九節和二十節.
耶穌的母親和祂的弟兄來找耶穌.
因為人很多.
不得去到耶穌的庚前.
有人就告訴耶穌.
指望耶穌的媽媽.
耶穌的弟弟來找祂.
有人就告訴耶穌.
第二十節.
你母親和你的弟兄站在外邊要見你.
真是骨肉的媽媽.
肉身的弟弟.
第二十一節.
耶穌回答說.
聽了神的道.
而遵行的人.
就是我的家人.
似乎說的家是神的家.
似乎說的是更加永恆關係的親屬.
相信我們不會理解錯.
耶穌你這麼大逆不道.
亡還亡.
沒理由在公開大庭講眾.
這樣奚落自己的媽媽.
自己的弟弟.
相信我們明白.
不是因為他現在有公務在身.
所以就要這樣來.
特意對自己的媽媽不敬.
相信主耶穌借這個機會.
又或者在兩個場合.
路加醫生很巧妙地這樣來引述.
要凸顯.
神的道在一個人的心裡面.
如果是重要.
如果是結實.
他可以成為主耶穌神家裡面的人.

$^{121}$有這樣的屬靈親屬關係.
即是意味著如果不聽神的道.
不來進行.
不來結實.
那就可能不是真正在神眼中的.
那個親屬.
那個生命.
所以頭尾都讓我們看到.
中間這個比喻的重要性.
比喻就是很多人聚集.
第四節.
當時有許多人聚集.
又有人從各城裡面出來見耶穌的時候.
所以很肯定這個時候.
是主耶穌初期出來傳道的影響力.
因為他曾經小試牛刀.
行神蹟.
給得的外母.
完全馬上康復.
不用吃Panadol.
不用去搶購.
馬上一叫他就起來.
完全沒有珀爾色痛.
主耶穌的按手就可以起來.
接著一過了安息日.
那些人就蜂擁來.
因為消息傳開了.
這樣在以色列突然間出了一個這麼大的先知.
和可以行神蹟的一個神人.
所以就引發很大的震撼.
很多人聽到風聲就來找耶穌.
所以路加也將這個場景告訴我們.
在這樣的背景下.
主耶穌說一個比喻.
聽道的比喻.
這樣就很特別了.
耶穌就用比喻說.
有一個殺種的出去殺種.
接著我們不詳細解釋了.
因為很快第九節.

$^{161}$門徒就已經來問這個問題.
這個問題也關係到.
不只是怎樣聽道的問題.
而是他們知道這個比喻背後是有些玄機的.
這個比喻背後.
承載著一個更加大的課題.
主耶穌就在這裡解開了.
第九節門徒問耶穌說.
這比喻是什麼意思呢.
耶穌就說了.
神國的奧秘是只叫你們知道.
你們就是這十二使徒.
至於別人就用比喻.
叫他們看也看不見.
聽也聽不明.
表面的字面意思.
相信大家都可能會誤會了.
以為為什麼這麼高竇的.
只是十二使徒內圍.
即是內圍的人就可以得到.
外圍的可能是門徒.
可能是跟隨者.
可能是蜂擁來到去想看熱鬧.
又或者可能想得到神蹟的.
就要拒諸於門外.
叫他們看不到的.
聽不明的.
那你又說.
為什麼會這麼特別呢.
意思當然不是這樣.
所以正正就是用比喻.
讓當時大家不會有任何.
聽不明的奧秘深奧.
但是背後又承載著一些.
好像很重要的哲理.
要怎樣來疏解呢.
怎樣來對號入座呢.
各有各說法呢.
原來不需要了.
因為耶穌在這裡.

$^{201}$直接就解給十二使徒知道了.
第十一節.
亦都是我們今天在這一堂崇拜裡面.
很想跟大家在年頭.
真的需要彼此的來提醒.
我們怎樣看待我們今年.
聽道的工程呢.
怎樣看待道在你的心靈裡面.
產生的作用呢.
第十一節.
這比喻乃是這樣.
種子就是神的道.
那些在路旁的.
就是人聽了道.
隨後魔鬼來.
從他們心裡把道奪去.
恐怕他們信了得救.
所以真是在說神國度的奧秘.
主耶穌很清楚.
一個人來到禮拜堂.
或者不同的地方.
聽到神的道.
而來到不遵守的.
就會原來是背後.
有魔鬼在他心裡面工作.
將他的道.
本來屬於他的.
都搶走了.
奪去了.
因為擔心這個人.
被道改變.
就會信了得救.
第十三節.
那些在盤石上的.
就是人聽道.
歡喜的領受.
是真的在說人.
不是在說教你怎樣殺種.
一個就快要弄水仙花了.
就快要過年.

$^{241}$我們要開班.
教怎樣插水仙.
怎樣弄頭.
不是說這些.
是說每一個人.
當時的猶太人.
當時的群眾.
當時不知道耶穌是誰的群眾.
很想親近聆聽.
和希望得到好處的群眾.
耶穌一直都在說這一班人.
是人聽道.
這裡就說歡喜的領受.
但是心中沒有根.
不過是暫時的相信.
及至遇見試煉就退後了.
試煉令他退後.
試煉令他不能夠.
心中繼續有根來扎.
讓他的道.
可以在心靈裡成長結實.
第十四節.
那樂在荊棘裡的.
就是人聽了道.
走開以後.
被今生的思慮錢財.
現樂積住了.
所以不只是不好的東西.
思慮很多時候就是比較負面.
很擔心.
很憂慮.
我們小時候.
看著爸爸.
一到年關.
很煩的.
很煩惱的.
那些店鋪打電話來追過年的找數.
通常就是.
過到就過吧.
有些就是去到年三十晚才找.

$^{281}$往往就找我們的子女.
來做那個磨心.
告訴他爸爸不在家.
很多時候都要這樣說的.
因為真的去到那兩個禮拜.
是沒有錢周轉.
一定要去到收到那些賬.
再去找數.
總是過到的.
年三十晚.
總是去到年三十晚.
我們就很開心的.
一起出街吃飯.
然後就買新衣服.
但是那個憂慮就是年關.
所以這裡說第三種人.
聽了道之後.
和他個人的那個際遇.
生活處境.
他面對的問題.
當然有錢財.
當然有艷樂.
艷樂就是很開心的日子.
開心就更加容易忘記.
會失去神的道.
不開心過度憂慮.
神的道只是這樣說.
實際是不行的.
年輕人你太過天真了.
我們這些見慣了.
不行的.
很多這些哲理.
都會影響我們對於神的道.
在我們心靈裡面的影響力.
使道結不出成熟的子粒來.
十五節.
如果大家不介意.
和我一起讀這一節聖經.
希望成為我們大家今年.
很重要.

$^{321}$我們每一個人的一個.
擺台頭的座右銘.
因為是關乎到你和我的一生.
那樂在好土裡的.
就是人聽了道.
並用純真善良的心.
持守他.
內心等候結果失.
好幾個比喻.
將人的心靈.
主耶穌.
比任何的物質的.
可以反照內心的.
身體的器官的.
無論Test scan.
無論這個.
磁力共振.
或者X-ray.
都要先進.
那些不是單純可以看到人的心靈的.
但是主耶穌的真理.
可以是歷久常新.
二千年前說的話.
到今天.
在每一個聚會.
每一個人的心裡.
都是到今天都有效.
就是我們怎樣預備.
我們的心靈.
成為一個在耶穌的講道裡的好土樣.
以致神的道.
可以在你的心靈裡.
來到去持守.
然後要內心的忍耐.
直至他能夠結實.
結到果實.
那個重點就是.
我們的人生.
是能夠結到.
記住是神眼中.

$^{361}$可以清許的果實.
而且道會將我們的心靈.
慢慢慢慢調整.
可以變成純真善良.
你期不期待.
我們的內心.
是這樣來到去承載神的道.
不單止道歸道.
心靈歸心靈.
而是那個道會改變自己的心靈.
再結出真正是神.
期許的果實.
所以你看看這節聖經是多麼美麗.
多麼值得我們看為是金句.
在新的年頭.
我們去追趕.
因為我們未見主一日.
都是可以生命改變的.
而且這個比喻.
其中一個重點就是.
沒有人可以說自己永遠是硬土的.
更加不會有任何一個人說.
我由信主那一刻開始.
到今天都是很土的.
不需要再有其他的東西.
不會的.
每一個弟兄姊妹.
總之恩典的福音.
何時你看為是你自己的.
你繼續這樣來到去學習.
一定可以在神的裡面.
來到去有那個結果的可能性.
因為聖靈已經在我們大家的心靈裡面.
來到進駐了.
我們已經有聖靈在我們的心靈裡面.
當然了.
我們都不是不知道.
我們想有一個很土的心靈.
沒有一個弟兄姊妹.
他說你不要再跟我說了牧師.

$^{401}$我只是認定自己是硬土.
我沒有得救了.
沒有一個弟兄姊妹這樣想的.
絕對不會的.
不過問題就在這裡.
那個忍耐.
在我們一生的功夫.
這個忍耐是最重要的.
不忍耐就不能夠結出那個果實.
不忍耐我們的生命.
就不能夠慢慢趨向純真和善良.
正正就是這個點.
是主耶穌要凸顯這方面.
即是主耶穌其實很知道大家的心靈.
當時的群眾.
當有一些的危難,代價.
很多人就立即否認.
很多人就會轉頭就走了.
所以難免難怪.
主耶穌說祂的肉可以吃.
祂的血是可以喝.
吃我肉喝我血的人就有永生.
這是約翰福音第六章.
在加伯隆的會堂.
之前那一晚.
行了五餅二魚的神蹟之後.
那些群眾就蜂擁地去找耶穌.
知道耶穌已經過了對面岸.
他們就第二天趕去.
知道耶穌在加伯隆會堂教導人的時候.
他們就來找耶穌.
耶穌第一句話就說.
你們找我是因為昨晚吃餅.
神蹟的餅.
不是加頓的餅.
是神蹟的餅.
五餅二魚.
你們因為吃餅得飽.
主耶穌怎麼說.
不如為那些必扭壞的食物勞力.

$^{441}$要為永生的食物勞力.
就是人子要賜給你們.
他們馬上就說.
我們怎麼可以做神的功.
去得到這些餅呢.
主耶穌說.
我就是生命的糧.
當時他們就完全不明白主耶穌所說的.
有一個相當偉大的籃球明星.
他曾經說過一句話.
最偉大最大一個人的成長.
就是他經歷最大挫敗的時候.
就是他最大的成長動力.
最大的挫敗.
不知道楊明先生昨天出獄.
那個是不是他很深刻的教訓.
坐了18天.
扣了一天有假期.
17天的監獄.
經過了三個節期.
都不可以和家人一起.
不知道他是不是一個最大的挫敗.
原來不是.
15年前都有前科.
都來到他曾經也是醉酒駕駛.
超速各方面.
最大的挫敗.
可以帶來我們生命的成長.
這個很需要我們來到他在神的面前.
去想一下我們整個基督徒人生.
我們那個忍耐的態度是怎樣呢.
面對危難困難.
我們怎樣去用什麼態度去面對.
而我們知道.
原來背後只會使我們的生命.
我們的生命是更加的成長.
我們能不能夠這樣去看人生呢.
的確有很多事情.
很多人會告訴我們.
我們很難在人生.

$^{481}$你整個人生.
每一手的牌都是好牌.
很多人會告訴我們.
就算是一手很爛的牌.
你怎樣打好它.
就是你的人生了.
沒錯.
知而行難.
我們知道一手爛的牌.
怎樣打好它.
但是那個動力是從哪裡來呢.
相信人生不是只有一場的牌.
不是這麼簡單.
而是我們的內心裡面.
有沒有想過.
主耶穌這樣教導我們的時候.
這個好土的生命.
怎樣能夠獲得呢.
怎樣我們能夠在我們面對.
我們自己困難的人生的時候.
都仍然忍耐地去結實呢.
有少許剛才看的上下文.
可能給我們有少許的線索.
雖然這是其中一個很短的場景.
是耶穌剛剛出來傳道的時候.
但是都類似反映了一些少許的資料.
讓我們看到初期教會的弟兄姊妹.
這些信徒.
他們是怎樣去理解主耶穌的呢.
如果大家回到第一節.
去到第三節.
盧家醫生給了我們.
可能很平鋪直敘的描述.
有不同的婦女跟隨耶穌.
用她們自己能夠服侍耶穌的.
去做她們的工作.
大家應該不會有誤解.
為甚麼聖經福音書.
沒有甚麼機會介紹一些婦女的呢.
如果大家知道那些背景.

$^{521}$你就會明白.
無論是希羅的時代.
或者猶太人本身.
婦女的地位是極低極低的.
是不可能拋頭露面.
公開地跟隨耶穌傳道的.
沒可能他們可以大庭廣眾去宣講教導.
連新約的提摩太前書.
婦女在會中是不准說話的.
不准講道的.
當時的二不所教會.
是這樣安排的.
但是這裡就有些很特別了.
為甚麼以一個男性主導的社會.
勞加醫生他好像不會對婦女有一種歧視呢.
因為或多或少當時猶太人.
對婦女的地位是相當低的.
他們不會覺得婦女有自己的權利.
是男性的附屬品.
不會有一個很高的地位.
他們可以有發言權.
不過這裡告訴我們.
作者真的很巧妙地告訴我們.
原來真正愛主耶穌.
跟隨主耶穌.
服侍主耶穌的婦女.
可能這幾個就是很明顯的代表.
他們的心就是這樣.
值得我們去想一想.
這是一個怎樣的心靈呢.
會不會正正就是這個比喻裡說的.
誠實善良的心靈呢.
很純真的心靈呢.
我再將這一段一到第三節.
稍為讀一讀給大家聽.
過了不多日.
耶穌周遊各城各鄉傳道.
宣講神國的福音.
和他同去的有十二個門徒.
相當清楚了.

$^{561}$十二使徒已經被呼召了.
好像五湖四海一樣跟隨耶穌.
但是這裡作者又有一些強調.
第二節.
還有被惡鬼所附.
被疾病所累.
已經治好的幾個婦女.
當然就是其中這個.
其中有稱為「茉大拉」的瑪利亞.
曾經有七個鬼.
在她身上趕出來.
鬼就不會有數量這樣計算.
一隻兩隻.
可能是吧.
但是靈界就不是這樣.
用物質那個量的.
但是主耶穌能夠在那裡.
用到那個數量.
就是很多的意思了.
曾經有七個鬼從她身上趕出來.
然後第三節.
又有希律的家宰.
撫撒的妻子.
希律的家宰.
就是她的管家.
希律王的管家撫撒.
撫撒的太太約亞娜.
還有可能是一些好的姐妹.
一些內室的閨密.
蘇薩娜.
還有好些的別的婦女.
似乎這些婦女是有錢有地位.
有當時的學識.
不過他們怎樣信主呢.
希律的管家.
信了主.
太太和其他有名的婦女.
用她們的財物.
來供給耶穌和門徒.
弟兄姊妹.

$^{601}$兩類的婦女.
都是耶穌基督的女徒.
婦女在當時.
根本上沒有發言權.
到處去跟耶穌大庭廣眾.
是不可以的.
一定會被人歧視的.
但是作者很巧妙地告訴我們.
這兩類的婦女.
有些是鬼婦.
鬼婦是不潔的.
是在社會上被邊緣的.
即是有鬼在你的身上.
其實是一些污穢.
是甚麼原因導致你鬼婦呢.
他們會調查的.
但是現在完全清除了.
這個莫特拉的瑪利亞.
當然不知道她的實際背景.
她的家底.
她是否已經結婚.
就沒有記載了.
但是有些是在皇宮裡面的有錢人.
愛主 跟隨主.
用自己的財物.
他們能夠做到.
又不會太張揚的.
這樣來供給主耶穌的傳道的即是.
即是他們有寫奉獻.
有放進奉獻袋.
這樣來支持主耶穌的工作.
初期教會有這樣的情況.
途徑是怎樣呢.
慢慢慢慢是否收集獻金的是猶大呢.
有可能是.
因為他是管錢郎的.
一路一路這樣來見到.
那就給我們有很多的聯想了.
弟兄姊妹.
在這十二使徒裡面.

$^{641}$他們因為是男性.
有些是憤外.
喜歡說話的彼得.
或者有些是很沉靜的.
譬如馬太 瑞麗.
有些是有其他的跟隨.
沒甚麼興趣記載他們的細節.
其實每個人都有他們的故事.
有他們的細節.
不過整個重點都是偏向耶穌.
和他們的關係.
和後來在被公開拒絕之後.
耶穌在後半段.
主要集中對十二使徒的教訓.
那個培養.
就不再對外的公開傳道了.
但這個時候是公開性的傳道.
耶穌可以吸引到不同的婦女.
讓她們來到真心樂意的跟隨.
她們未必認識很多的真理.
她們未必在宣講上.
或者甚至行神蹟上.
有參與的權.
因為當時的社會風氣.
是不容許女性公開拋頭露面的.
其實我們媽媽的年代都是這樣.
讀書根本是不用想的.
我對上的姐姐小學.
都不讓她讀書.
你可以想像得到.
中國人都是這樣.
何況二千年前的那個時候.
希羅的文化都是一樣的.
所以你想像得到.
主耶穌的道.
如何可以進入到這些婦人的內心.
而她們是很真心地愛主耶穌.
服侍主耶穌呢.
我們曾經在復活節也說過.
這個茉達拉的瑪利亞.

$^{681}$是第一個去到墳墓見耶穌的復活.
耶穌是第一個向這個婦人顯現的.
即是說沒有記載她.
但其實在那幾年的傳道裡面.
茉達拉的瑪利亞是形影不離的.
可能是心靈的形影不離.
對主耶穌的愛和服侍.
因為她知道自己是一個.
可能是污穢的人.
因為鬼婦.
但主耶穌已經將她完全拯救.
她就是用她可以服侍的方式.
去服侍跟隨耶穌.
弟兄姊妹.
這個是道的結實.
有些貴族人家.
他們會將他們自己的生命.
他們的婦孺.
在神的國度裡面參與.
有份.
有錢人有有錢人的那個場景.
去服侍耶穌.
盼望我們能夠在這個小小的講道裡面.
引發大家不同的反思.
你有你的場景.
我有我自己的那個場景.
神怎樣來使用我們每一位.
讓神的道在我們的內心.
慢慢慢慢的結實.
以致我們能夠真的可以在神的裡面.
將來是被神來清去.
有一個最後的預道故事.
小時候就聽過了.
都不是很明白的其實.
其實很簡單的.
這個醫學預言的故事.
太陽和北風.
在爭誰最厲害.
北風說.
我的力.

$^{721}$吹什麼都吹倒的.
太陽說.
未必.
剛剛有個人走過.
北風就和太陽說.
你信不信.
我一吹.
他的衣服全部掉了.
全都掉了出來.
太陽說好啊.
大家試試吧.
用盡他的力.
不斷地吹.
很強烈的北風.
那個人又冷又大風.
把他的衣服和帽子又弄得那麼緊.
接著北風用盡他的力.
都不行了.
放棄了.
太陽說不如我試試吧.
太陽就用他很溫暖的光線.
來照射這個人.
那些風停了.
風和日麗.
又有溫暖的太陽.
那個人就很開心地.
脫下他的帽子和外套.
享受冬日的溫暖.
弟兄姊妹.
保羅說在臨前的十三章.
神聖的愛是恆久的忍耐又有恩慈.
愛是不嫉妒.
不自誇.
不張狂.
不作害羞的事.
凡是相信.
凡是盼望.
愛是永遠不止息.
說的不只是地上的.
那個不是說婚姻的.

$^{761}$那段聖經是說弟兄姊妹的相愛.
因為哥林多教會有爭拗.
爭誰的恩慈大.
所以保羅就說真正的愛.
神聖的永恆的.
但第一樣就是愛.
神的愛是恆久的忍耐.
所以我們對其他人都要有這份恆久忍耐.
其實二十多年前已經讀這些經文.
但去到現在這個年紀.
越去看的時候.
神的愛如果在弟兄姊妹的心靈裡面.
就會讓你慢慢感受到那種力量.
神愛的那種溫柔.
忍耐的力量.
會使我們大家願意闖開我們大家的心靈.
來愛主.
跟隨他.
和在神的道裡面去結實.
很盼望神的愛.
不只是冬日太陽那麼簡單.
而是真的在我們生命裡面有份量的.
基督耶穌降世.
為要拯救罪人.
這說是可信的.
是可佩服的.
這是史圖保羅和提摩太林.
他生命最結束的時候所說的.
可信可佩服的主.
這樣來將保羅拯救.
其實你知不知道提摩太.
在罪人當中我是一個罪魁.
然而我都蒙了憐憫.
我們一起禱告.
讓我們今天晚上.
能夠有第十五節.
那麼寶貴的神話語的一個應許.
和一個呼喚.
很盼望我們大家一起.
在今年為神的國度我們去追求.

$^{801}$我們能夠很心甘情願的.
在神的道裡面去裝備.
去成長.
去結實.
多謝主聽我們的禱告.
奉耶穌基督的名.
阿們.
\newpage



\section{撒迦利亞書 2:6-8}
\label{sec:ejP2s680TsQ}
\textbf{2023-01-08 加恩樓《上帝眼中的瞳人》撒迦利亞書 2\_6-8 講員:陳韋安牧師}
\newline
\newline
連結: \href{https://youtube.com/watch?v=ejP2s680TsQ}{\texttt{https://youtube.com/watch?v=ejP2s680TsQ}} ~~~~ 語音日期: 2023-01-10
\newline
\newline
\hyperref[sec:gAWlXGsKQ90]{\small{< < < PREV SERMON < < <}}
~
\hyperref[sec:index_chronic]{\small{[返順時目]}}
~
\hyperref[sec:index_scriptual]{\small{[返順卷目]}}
~
\hyperref[sec:uJ8mvXZ5iqw]{\small{> > > NEXT SERMON > > >}}
\newline
\newline
撒迦利亞書 2:6-8
\newline
\begin{longtable}{cl}
\hline
\hline
章節 & 經文 (和合本修訂版)\\
\hline
2:6 & \begin{tabularx}{0.7\textwidth}{X} 耶和華說：「來，來！你們要從北方之地逃回；因我曾把你們分散到天的四方。這是耶和華說的。」 \end{tabularx} \\ \\ \relax
2:7 & \begin{tabularx}{0.7\textwidth}{X} 來！住巴比倫的錫安百姓啊，逃吧！ \end{tabularx} \\ \\ \relax
2:8 & \begin{tabularx}{0.7\textwidth}{X} 萬軍之耶和華在顯出榮耀之後，差遣我到擄掠你們的列國那裡，他如此說：「碰你們的就是碰他自己眼中的瞳人。 \end{tabularx} \\ \\
[1ex]
\hline
\hline
\end{longtable}
$^{1}$評論節目平安,早安.
剛才是我第一次領聖餐,剛才很緊張.
大家一起經歷了新一年的聖餐.
真的,我覺得2023年充滿盼望和期待的一年.
希望節目能夠在新年同樣懷著盼望重新出發.
在當中有一個平安快樂的一年.
這些不是純粹客套的說話.
我們過去幾年都不太容易.
經歷了很多不容易的時間.
好像還沒完.
但當我們踏入2023年,我覺得好像有少許改變.
世界社會來到2023年的時候.
情況似乎有一點點的變化.
說的未必是客觀環境的突破.
可能社會依然都是差不多的景象.
但世界的人似乎都改變了一些東西.
有些人似乎離開了困局.
有些人領悟到一些道理.
有些人可能看到一些新的事物.
能夠成長的可能都能夠成長.
尋找一些出路的都找到一點點出路.
雖然這條路未必光明一片.
但起碼去嘗試踏入2023年.
所以2023年我自己是樂觀的.
改變的未必是這個世界.
而是我們當中的人能夠有少許的不同.
可以成長改變.
所以今天在2023年的第一個崇拜裡.
都可以和大家看看有關一個視野的題目.
不知道大家有沒有聽過一句話.
我們的視野決定我們的高度.
不知道大家是否明白這句話的意思.
視野決定了我們的高度.
當然聽到這句話是相反的.
高度是決定視野.
但高度決定視野大概是一個物理問題.
你站得越高就能夠望得越遠.
是庸窮千里目的朋友.
不過視野決定高度是一個人生的哲理.
我們人生裡面 尤其是我們作為教徒.

$^{41}$我們是天上父上帝的兒女.
我們的眼光正在定義我們的高度.
你看得有多遠 你的人生的生命有多高.
你看得有多闊 你所做的就有多闊.
所以我們的視野和看見.
似乎是我們任何事情的第一步.
所以面對2023年.
我們一齊嘗試思考信仰裡面和眼界視的課題.
所以我們今天的講題叫做.
上帝眼中的同人.
其實這個關乎於.
我在Folk Church裡面有一套劇.
上年聖誕節的時候.
其中一個紀念品.
一個Toe bag.
印了一句I am the apple of his eye.
這個slogan.
其實這個slogan是在聖經裡面.
上帝眼中的同人這句經文.
我上了這麼久都沒有認真研究過這段經文.
所以我嘗試今天一齊去思考一下.
這個經文的意思.
等我按好嗎.
我怕我一會兒會劇透.
我們一齊去想一想.
Apple of his eye.
這句說話的典故.
Apple of his eye.
這句中文的版本叫眼中的同人.
如果你翻舊翻來我也知道.
在聖經裡面經常出現了幾次的經文.
聖經裡面出現了五次Apple of his eye的說話.
第一段是在申命記第32章.
耶和華說.
遇見他在曠野荒涼野獸哭叫之地.
就環繞他,看顧他,保護他.
而同保護眼中的同人.
第二段是在詩篇第17篇.
禱告的.
詩人這樣禱告.

$^{81}$求你保護我,而同保護眼中的同人.
將我隱藏在你翅膀的陰下.
第三段是箴言的經文.
箴言這樣說.
遵守我的命令,就得存活.
保守我的法則,好像保守眼中的同人.
第四段是耶利米埃哥的哀求.
詩人這樣禱告.
石安的城牆,願你流淚如何.
晝夜不息,願你眼中的同人流淚不止.
很慘的經文.
最後就是今天提的經文.
撒迦利亞書第二章第八節.
萬君子耶和華說.
在獻出榮耀之後.
差矣我去懲罰那擄掠你們的列國.
摸你們的就是摸他眼中同人.
其實這個他字就是我字.
就是上帝自己的眼目的意思.
五段經文裡面.
英文的聖經全部都翻譯成.
Apple of the eye.
是這樣的意思.
不知道大家有沒有問問題.
為什麼會是Apple.
為什麼不是Orange.
不是Banana,不是Samsung.
為什麼會是Apple.
Apple of the eye其實是一句英文的成語.
意思是什麼.
就是眼中極度喜愛捍衛愛慕.
最重要最寶貴的那位.
所以Apple of the eye這個英文成語.
最早文獻出現就在第九世紀.
一些很古老的英文文獻裡面有的.
莎士比亞都曾經在眾夏日之夢裡面.
都用過Apple of the eye來寫.
後來自此之後.
全部的英文聖經.
都將剛才的五段經文翻譯成.

$^{121}$Apple of the eye這個字.
雖然說了這麼久都不知道為什麼叫Apple.
不過我們知道什麼呢.
就是知道希伯來文聖經裡面.
原文裡面其實不是Apple這個字.
是什麼字呢.
就是解作瞳孔.
就是我們眼球裡面.
眼球裡面中間黑色點點的位置.
這個部分.
所以無論是希伯來文.
或者後來翻譯的希臘文也好.
這個字直譯是解作女兒的意思.
Daughter的意思.
或者是一個小人兒的意思.
所以不知道有沒有關係.
因為我讀文科的.
以前我讀拜安的時候.
我們瞳孔裡面叫什麼.
記得嗎.
記得嗎 土土.
是什麼.
是叫Pupil的.
Pupil就是瞳孔.
是解作學生的意思.
是有關係的.
聖經和字都是這樣.
這個瞳孔其實也是解作小朋友.
或者是學生的意思.
所以可以這樣理解.
我們眼睛中間的瞳孔.
可以解作蘋果.
可以解作瞳孔.
可以解作學生.
或者是女兒的意思.
不過我覺得最漂亮的.
其實是中文聖經.
解作什麼呢.
是翻譯成同人這個字.
是一個很有意思的字.

$^{161}$中文就沒有蘋果這個字眼.
翻譯成眼中的同人.
你翻譯文言文的話.
你會發現.
一些很古舊的中文聖經.
翻譯成母子.
眼貌的貌.
兒子的子.
論中日本就翻譯成眼中之女.
所以這個字原來.
同人這個字其實是一個.
中文有的字.
原來古代的中國人.
都是會有同人這個字.
出現了很久.
譬如唐朝.
唐朝有一個叫杜家唐二哥.
他就寫.
骨寵神寒天妙氣.
一箱同人輾抽水.
宋朝的一首詩就說.
同人輾水腰如簇.
一幅烏紗果寒玉.
所以原來中國就不是.
只有十八個同人的數字.
是有一個同人這個字.
有一個這樣的稱呼.
是解作同孔的意思.
所以無論是聖經.
或者是唐詩都好.
古人都會將這個同孔.
理解為一個同人.
這樣的一個意思.
蘋果的意思.
這個就是Apple of the eye.
有這樣的典故的資料.
明白了Apple of the eye.
的典故之後就開始去看.
就是聖經裡面的一個蘋果.
就是今天所說的撒迦利亞書.

$^{201}$其中一個Apple of the eye的經文.
先說說一些簡單的撒迦利亞書的背景.
撒迦利亞書是一本比較後期的先知書.
一本是說以色列人被擄流亡之後.
將要歸回的一個小先知書.
所以先知撒迦利亞對著一些流亡海外的以色列人說話.
都對著一些仍然存留在一個廢墟的耶路撒冷的以色列人說話.
公元前586年耶路撒冷淪陷之後.
以色列人被擄.
其實有好幾批不同的人.
有一批是流亡海外的以色列人.
他們是被擄的.
有一批是因為這樣的情況而移居到其他地方的以色列人.
但仍然有些是什麼呢?.
是一些仍然住在滅亡了的耶路撒冷的以色列人.
就好像今天香港一樣.
有些人移民,有些人仍然在香港.
因此撒迦利亞書和哈該書都有一個很重要的目的.
就是要代表耶路撒冷上帝.
去吩咐一群身處異地的以色列人.
以及一群殘留在耶路撒冷的廢墟裡的人.
吩咐他們要重返耶路撒冷.
去重建耶路撒冷.
並且恢復他們的聖殿.
所以可能大家不知道.
其實哈該書和撒迦利亞書頭的八章是一個完整的文獻.
一起都是在說一個很重要的主題.
撒迦利亞書是強調一個所謂漢見.
Vision的詩之書.
它是強調耶和華上帝的漢見.
藉著撒迦利亞先知的漢見.
去表達耶和華上帝漢見耶路撒冷的願景.
雖然是一片廢墟.
但是耶和華上帝漢見什麼呢?.
所以頭八章的撒迦利亞書裡面.
都有八個不同的意象.
八個的Vision.
其實漢見Vision就是意象的意思.
不過,本次我想說的是什麼呢?.
所謂的漢見.

$^{241}$或者是Vision或者是意象都好.
其實是一些不是那麼容易消化的東西.
一些很難以置信.
一些尚未發生的事情.
撒迦利亞先知要向一些移居了海外的斯里文講說話.
要說服他們要歸回.
其實是一個純粹將來的漢見.
一些尚未發生的東西.
所以這個現實景象和將來漢見的差異.
其實是甚大.
你望著一個地盤廢墟.
你要幻想是一個將來的豪宅.
一個很貴的樓.
是一個很不容易的事情.
所以這是一個非常困難的事情.
今天有很多朋友.
大家朋友親戚移民.
去到英國已經很不容易找到工作.
找到屋住.
開始穩定.
甚至乎要分.
你很難再突然間去說服他們回到香港.
開始已經順境了.
開始已經很穩定了.
甚至乎每個星期都已經超納.
但就發覺很不容易回來.
所以發覺整件事就是撒迦利亞書的處境.
仍然是一個頹垣敗瓦.
會演像一個廢墟的耶路撒冷.
撒迦利亞要去說服這些人.
回來重建這個城市.
所以正因為這個現實景象和意象漢見的相差甚遠的情況下.
上帝就說出了這句話.
希望呼籲這些流亡海外的人.
能夠安心不用怕.
可以回來這個山裡.
這個說話就是今天看的說話.
撒迦利亞書第二章一到六節經文.
我從前分散你們在天的四方.
現在你們要從北方之地逃回.

$^{281}$與巴比倫人同住的錫安文南.
應當逃脫.
萬君子說再顯出榮耀之後.
差遣我去懲罰那擄掠你們的列國.
摸你們的就是摸我眼中的銅人.
上帝要求這些流亡海外的以色列人歸回.
叫他們不要怕.
因為上帝會保護他們.
因為他們就是眼中的銅人.
Apple of his eye.
甚至經文裡面有一個很有趣的描述.
摸你們的就是摸我眼中的銅人.
一個非常生動特別的一句說法.
經文裡面說你們回來吧.
你們不用怕.
你們可以回流回來.
因為我會幫你懲罰一切的惡人.
誰擾你們的就是擾我眼中的.
一句很奇怪但很有說服力的說話.
誰擾你們的就是擾我眼中的.
感受到這句說話的氣場.
誰擾你們的就是擾我眼中的.
上帝說你們不要害怕將來.
雖然你們看不到將來.
但你們就是我眼中的銅人.
誰去擾你們的就等於擾我眼中的.
一個切膚之痛的愛護.
上星期1月1日.
星期日我和我女兒去屯門友愛村玩滑板.
星期日很忙所以整個星期都沒有陪她.
兩父女在星期日去友愛村玩Xgame場.
但不知道有沒有去過.
那裡有個政府場是非常潮的.
三層一層是玩滑板一層是玩BMX.
在一個卡拉OK上的場.
去到沒有人.
兩父女就包了整個場.
玩了很久很開心.
本來打算離開.
離開前我女兒說不如玩十分鐘.

$^{321}$怎知道就真的樂極生悲.
離開前突然有事關我事.
我老實說.
滑板失腳.
滑板失平衡就滾到她那邊.
她突然有事衝過來.
就跌倒在地上.
我知道就糟糕了.
我馬上跑過去.
我看到她很神色凝重.
我就心知不妙.
她沒有哭.
不過當她打開口罩的時候.
我發現口罩裡面是她兩隻門牙的三分之一.
就在口罩那裡.
整個門牙沒有了前面的部分.
崩了.
血在口裡.
我當時很後悔.
非常傷心.
非常擔心.
因為兩隻是行刺.
她就覺得怎麼辦.
馬上就走.
去友愛村買點水.
馬上被她擋住.
因為是公隔.
找不到牙醫去看.
不過幸好我有個朋友是牙醫.
就問她怎麼辦.
她說也不是很嚴重.
沒有傷很嚴重.
也可以不急的.
打算星期二才帶她去看.
那晚吃飯.
吃完飯她突然脫牙.
幸好不是門牙.
只是脫牙.
嚇死我.
原來是另一隻牙.

$^{361}$搞了很久.
星期二就帶她去看醫生.
看我朋友那裡.
幾經波折.
搞了很久.
終於補回牙齒.
這張照片是我想這樣看的.
後來這張是補完牙的照片.
馬上在牙醫門口拍給媽媽看.
沒事沒事.
跟她媽媽說話.
那時候我馬上發現.
其實她正正就是我眼中的同人.
她就是我眼中的女兒.
那時候我就明白.
為什麼聖經會把眼中的同人.
直譯成The Daughter of My Eye.
那種心肝寶貝.
Michael也知道.
心肝寶貝如珠如寶.
萬千寵愛的那種感覺.
我的女兒就是我眼中的同人.
當你眼中的同人受到傷害的時候.
你寧願受傷的是你自己.
你寧願自己崩了門牙.
也不相信你的女兒崩了門牙.
我正想說什麼呢.
就是如果作為一個普普通通的人.
都會有我們心目中眼中的同人的時候.
何況我們天上的父上帝呢.
上帝去看我們的時候.
當你有任何的傷害.
你有任何的苦況的時候.
那種心痛.
那種苦況.
正正就是上帝那種心痛感受.
面對著2023年.
我們要認定的.
我們是上帝眼中的同人.
我們這樣來看見這件事情.

$^{401}$不過今天說到不只是停在這裡.
當然我們可以停在這裡.
作為一個很安慰的訊息.
我們知道我們是上帝眼中的同人.
來面對著2023年.
但我想說的不只是這一點.
沒錯,我們是認定我們是上帝眼中的同人.
不過我們更加要知道.
同一雙上帝的眼睛.
塞加利斯所說的不單單是上帝的一種視野.
我們看一段經文.
同樣是撒嘉利亞書的經文.
在第八章裡面.
叫做耶和華上帝的眼睛的另一段經文.
我們一起讀吧.
預備,一二三.
撒嘉利亞書第八章第六節.
先知撒嘉利亞描述了耶和華上帝看見的意象.
經文記載了同一雙上帝的眼睛.
耶和華上帝看見斯蒂盟約要歸回重建的城市.
這個城市當中的景象.
耶和華看到這個城市年老的男女在街上.
看到街上滿有男孩女孩在那裡玩耍.
這個是耶路撒冷和平的景象.
這個是上帝在廢墟裡面看到的和平景象.
所謂的vision 所謂的意象 所謂的看見.
正正是這樣的視野.
一些尚未呈現 尚未發生.
卻是一切行動以先 首先出現的必然要素.
值得提的是第六節.
對於這個視野 這個眼睛的問題.
上帝提出了一個值得我們深思的問題.
上帝說要去反問以色列人一個問題.
他說 這是在愚民眼中看為稀奇.
在我眼中有看為稀奇嗎.
一個很有趣的反問.
這個說話翻譯成什麼意思.
耶和華說 這些圖畫 這些意象 這些事情.
你沒見過嗎 你覺得很出奇嗎.
你以為我和你一樣覺得很出奇嗎.

$^{441}$你以為我都沒見過嗎.
一個很特別的反問.
這個說話正正是說耶和華上帝的視野.
和以色列人視野的差異.
以色列人看到的是一個被摧毀了的.
一個已經是非常慘烈的廢墟.
但是上帝在這個頹垣敗瓦的地方.
卻看到一個和平的國度.
很有趣的是 以色列人作為耶和華上帝眼中的同人.
自己卻看不到上帝所見到的東西.
這幫以色列人不是看不見.
他們只不過是看不到.
什麼意思呢 這幫人不是瞎的.
他們看到很多東西.
但他們看不到他們應該要看到的東西.
經文告訴他們 這幫人不是看不見.
只不過是看不到.
當我閱讀經文的時候.
我發現了第四世紀有一個教父.
叫做Dedimus de Blin.
一個叫盲人的的莫斯的教父.
一個盲的基督徒 寫了很多書.
剛好他也在寫撒迦利亞書的《食經》.
一個盲的人在形容撒迦利亞書的經文.
他在形容撒迦利亞書中說.
"With the eye of his mind enlightened and being divinized by the word that came to him,.
he had wonderful vision and proclaimed the prophetic word.".
原來重要的不是你看到多少東西.
而是你內心的眼睛有沒有被擴闊.
從而能夠看到他應該要看到的異象.
一個盲的人可能比我們雙目見全的人看得更加清楚.
不知道大家有沒有同感.
今天大家活在2023年的年代裡.
大家一打開手機.
發覺現在YouTube, IG, Facebook 全部都很TikTok化.
十五秒的短片全部都是垃圾.
全部都是花你時間吸引你的眼球 是沒有內容的東西.
怎樣能夠讓那個人看不見.
你可以去蒙掉他的雙眼.
你可以去遮住他的眼.

$^{481}$你可以把眼睛畫出來.
但最簡單的方法是什麼.
就是不斷去看很多很多不同的東西.
從而讓他看不到他應該要看的東西.
所以看見的相反詞其實不是看不見.
看見的相反詞是忽略.
明明擺在你眼前.
你卻是有眼無珠.
你卻是眼大看那個籠.
你卻是視而不見.
所謂視而不見.
就是明明存在.
你卻看不到.
你只見到很多不應該看.
阻礙你生命的事情.
所以2023年.
我們用上帝的眼光去看事物.
保持我們的vision.
保持我們的意象.
去看一些尚未出現的事情.
真人 世人看東西是看過去往績.
看過去的歷史 看過去的趨勢.
但是基督徒卻被教導.
要看一些我們未必看得到.
但相要出現的事情.
我想佳樓也是一樣.
我們看到的不是純粹過去的樣子.
而是這個地方可以.
上帝如何興起他的座位.
看見是一個倫理問題.
你應該要看些什麼.
多於你看得多.
看得多快.
所以讓我們用這個眼光去審視我們的世界.
我們如何看2023年的世界.
如何閱讀這個世界.
你看到的是什麼.
看到的是一些偏見 短見 還是成見.
還是我們可以真正看到上帝要我們看的事情.
求主開闊我們的眼睛.

$^{521}$讓我們能夠用新的眼光去看見這個世界.
我想這是我們迎接2023年裡.
一起行動以先最重要的事情.
看得通 看得透.
我自己很好奇.
2023年不單叫人對你刮目相看.
更重要的是你對這個世界刮目相看.
能夠看得見這個世界上帝要我們看見的東西.
我們一起祈禱.
祈求你幫助我們.
讓我們能夠被你的靈去擴闊我們的眼光.
當中我們不同等級的生命裡.
都遇到很多不同要看得見的東西.
求主讓我們能夠在當中去分辨.
能夠知道哪些是應該要看的.
哪些是有需要的.
無論是我們這個社區裡.
我們要看見的人的需要.
或者是一些新的事情的發展.
求主讓我們有信心.
來開闊我們的心眼.
來看見你的作為.
看見你所看見的事情.
求主這樣帶領我們.
讓我們更投入成為一個有意象的教會.
一個有視野的教會.
能夠看見你的作為的教會.
求主這樣幫助.
共主命求 阿們.
\newpage



\section{提摩太後書 4:6-22}
\label{sec:uJ8mvXZ5iqw}
\textbf{2023-01-08 午堂《心友靈情》提摩太後書 4\_6-22 講員:曹偉彤院長}
\newline
\newline
連結: \href{https://youtube.com/watch?v=uJ8mvXZ5iqw}{\texttt{https://youtube.com/watch?v=uJ8mvXZ5iqw}} ~~~~ 語音日期: 2023-01-10
\newline
\newline
\hyperref[sec:ejP2s680TsQ]{\small{< < < PREV SERMON < < <}}
~
\hyperref[sec:index_chronic]{\small{[返順時目]}}
~
\hyperref[sec:index_scriptual]{\small{[返順卷目]}}
~
\hyperref[sec:dkl9UnCmWfQ]{\small{> > > NEXT SERMON > > >}}
\newline
\newline
提摩太後書 4:6-22
\newline
\begin{longtable}{cl}
\hline
\hline
章節 & 經文 (和合本修訂版)\\
\hline
4:6 & \begin{tabularx}{0.7\textwidth}{X} 至於我，我已經被澆獻，離世的時候到了。 \end{tabularx} \\ \\ \relax
4:7 & \begin{tabularx}{0.7\textwidth}{X} 那美好的仗我已經打過了，當跑的路我已經跑盡了，該信的道我已經守住了。 \end{tabularx} \\ \\ \relax
4:8 & \begin{tabularx}{0.7\textwidth}{X} 從此以後，有公義的冠冕為我存留，就是按著公義審判的主到了那日要賜給我的；不但賜給我，也賜給凡愛慕他顯現的人。 \end{tabularx} \\ \\ \relax
4:9 & \begin{tabularx}{0.7\textwidth}{X} 你要趕緊到我這裡來。 \end{tabularx} \\ \\ \relax
4:10 & \begin{tabularx}{0.7\textwidth}{X} 因為底馬貪愛現今的世界，已經離棄我，往帖撒羅尼迦去了；革勒士往加拉太去；提多往撻馬太去； \end{tabularx} \\ \\ \relax
4:11 & \begin{tabularx}{0.7\textwidth}{X} 只有路加在我這裡。你來的時候把馬可帶來，因為他在服事上於我有益。 \end{tabularx} \\ \\ \relax
4:12 & \begin{tabularx}{0.7\textwidth}{X} 我已經打發推基古往以弗所去。 \end{tabularx} \\ \\ \relax
4:13 & \begin{tabularx}{0.7\textwidth}{X} 我在特羅亞留給加布的那件外衣，你來的時候要帶來，那些書也帶來，特別是那幾卷羊皮的書。 \end{tabularx} \\ \\ \relax
4:14 & \begin{tabularx}{0.7\textwidth}{X} 銅匠亞歷山大多方害我；主必照他所行的報應他。 \end{tabularx} \\ \\ \relax
4:15 & \begin{tabularx}{0.7\textwidth}{X} 你也要防備他，因為他極力抗拒我們的話。 \end{tabularx} \\ \\ \relax
4:16 & \begin{tabularx}{0.7\textwidth}{X} 我初次上訴時，沒有人前來幫助，竟都離棄了我，但願這罪不歸在他們身上。 \end{tabularx} \\ \\ \relax
4:17 & \begin{tabularx}{0.7\textwidth}{X} 惟有主站在我身邊，加給我力量，使我能把福音完整地傳開，讓所有的外邦人都聽見；我也從獅子口裡被救出來。 \end{tabularx} \\ \\ \relax
4:18 & \begin{tabularx}{0.7\textwidth}{X} 主必救我脫離一切的兇惡，也必救我進他的天國。願榮耀歸給他，直到永永遠遠。阿們！ \end{tabularx} \\ \\ \relax
4:19 & \begin{tabularx}{0.7\textwidth}{X} 請向百基拉、亞居拉和阿尼色弗一家的人問安。 \end{tabularx} \\ \\ \relax
4:20 & \begin{tabularx}{0.7\textwidth}{X} 以拉都在哥林多住下了。特羅非摩病了，我把他留在米利都。 \end{tabularx} \\ \\ \relax
4:21 & \begin{tabularx}{0.7\textwidth}{X} 你要趕緊在冬天以前到我這裡來。友布羅、布田、利奴、革老底亞和眾弟兄都向你問安。 \end{tabularx} \\ \\ \relax
4:22 & \begin{tabularx}{0.7\textwidth}{X} 願主與你的靈同在！願恩惠與你們同在！ \end{tabularx} \\ \\
[1ex]
\hline
\hline
\end{longtable}
$^{1}$很開心能夠再次和你們一起去敬拜.
分享一起敬拜的喜樂.
我們同心禱告.
親愛的主多謝你讓我們能夠有榮幸.
我們能夠歡喜快樂的一起相聚.
在你的聖所當中我們歌頌讚美你.
我們在今年頭的開始.
我們繼續在面前立願.
我們走人生的道路.
如今我們重弟兄姊妹們都要聽你的話語.
求你讓我們有聰明有智慧.
讓我們的耳朵聽得見.
我們的眼睛看得見.
我們的心會相信.
我們將你的真道走出來.
走在你的真道上面.
我們的祈禱是奉主耶穌基督的名字求.
阿們.
歷代以來基督徒都有不同形式的.
屬靈的操練.
為的是要更加的去親近神.
專注在神的身上.
其中一個方式就是朝聖.
在天主教在基督教的實踐當中都有朝聖的.
早期的一位女朝聖者叫做Igiria.
她來自於今天的西班牙.
就在384個AD.
她曾經去到地中海各處.
最後去到聖地去朝聖.
她朝聖的經歷就寫下來.
成為重要的文獻.
在中世紀朝聖的活動是很頻繁的.
不過變成了商業的活動.
所以宗教改革者嘉爾文.
他就不開心.
他說為什麼人們會花錢去朝聖呢?.
為什麼要跪拜那些死人的骨頭呢?.
這些是迷信.
只會傷害我們的信仰.
如果新教的基督徒.

$^{41}$他不是像嘉爾文那樣去朝聖的話.
我們是不是還要朝聖呢?.
其實要朝聖.
我們可以回到聖經那裡.
我們看看怎樣去朝聖.
在聖經當中.
其實聖經是講到一個朝聖的遊記.
在我們從創世記到希伯來書11章.
我們看到亞伯拉罕和撒拉.
開展了一個很重要的信心的旅程.
在希伯來書的11章第8節就這樣說.
「恩爵信 亞伯拉罕蒙召的時候.
就遵命出去往將來要承受為基業的地方去.
他出去的時候都不知道去到哪裡.
從亞伯拉罕的蒙召.
我們知道後來他們去到迦南這個地方作客旅.
在這個過程當中.
其實他們是面對很多挑戰的.
很多的困難.
人生的道路有很多的困難.
但是他們的信心是很穩定的.
亞伯拉罕的後裔同樣都是要面對很多的挑戰.
但是他們同樣都是有一種信心.
那種的忍耐能力.
因為他們都相信神的應許.
就是他們在等候神所設計的.
所建造的那個新的城市.
是天上的耶路撒冷.
所以在一點上其實就是.
亞伯拉罕的朝聖的故事成為了.
歷代以來.
信徒的很好的朝聖故事.
在新約當中.
我們見到保羅是在朝聖的道路當中.
在朝聖的道路當中.
他是有一些的同伴的.
有些朋友的 有些學生的.
其中一位就是提摩泰.
他就是和提摩泰同行.
他教導提摩泰 鼓勵他.

$^{81}$鼓勵他要成長.
鼓勵他要壯大.
在提摩泰後書一章.
保羅就這樣說.
他說 我記得你無偽的信心.
為者緣故我提醒你要把神.
藉著我按手給你的恩賜.
在如火挑旺起來.
因為神賜給我們的.
不是膽怯的心.
而是剛強人愛自制的心.
保羅就是鼓勵提摩泰要強大.
要強大 要成長.
當保羅這樣鼓勵提摩泰的時候.
其實保羅在艱難的情況當中.
他坐牢 他被遺棄.
他的身體也不好.
以前他曾經坐牢.
在《肥納比書》記載.
他曾經坐牢 但那次坐牢.
他坐得很開心.
因為他坐牢的時候.
他更加剛強壯膽.
繼續努力傳福音.
他很開心的.
但這次坐牢 在提摩泰後書看到.
他的經歷真的很不同.
這次是很慘淡.
這次真的是他生命嚴寒的冬天.
他法庭的申辯沒有成功.
他的同伴 他的同工.
沒有因為他坐牢而剛強壯膽.
每個人都是雞飛狗走.
甚至是敵對他.
保羅知道自己的身體不好.
在生命的終結當中.
所以保羅在四章六節這樣說.
他說至於我.
我已經被驕獻離世的時候到了.
驕獻是什麼呢?.

$^{121}$驕獻的意思是說.
電酒祭就是將酒倒出來.
這個離世的祭司就是說.
是指他的身體離散.
他的身體在離散的季節.
就是這樣的情況.
保羅就是和提摩泰分享他生命的故事.
鼓勵提摩泰.
在四章七節這樣說.
那美好的象我已經打過了.
當跑的路我已經跑盡了.
該信的道我已經守住了.
保羅的故事就是說.
這場仗我打完了.
這場賽事我跑完了.
我是守住了我的信心.
守住了我的信仰.
我已經跑到終點了.
是有這個冠冕留住給我.
所以在第八節保羅這樣說.
從此之後.
有公義的冠冕為我存留.
當作公義審判的主.
到了那一天要賜給我的.
不但要賜給我也.
賜給凡愛無他顯現的人.
保羅的意思就是說.
現在是現在.
冠冕在等著我.
因為冠冕已經為保羅存留.
保羅得到這個冠冕.
不是因為保羅有什麼成就.
或者是神給他任何的補償.
不是.
他是能夠得到冠冕.
因為冠冕是給義人的.
正義的法官是耶穌基督.
那些愛慕耶穌基督顯現的人.
愛慕耶穌基督的人.
他們就會得到生命的冠冕.

$^{161}$所以冠冕是給哪些人的.
給那些保持信心的人.
是那些因信心而稱義的人.
所以保羅說我已經跑完賽事了.
打完這場仗.
我保持了信心.
當我們這樣看的時候.
你說保羅是不是有盼望呢?.
其實保羅是有很多的原因.
是可以絕望的.
因為我們在今天所讀的經文當中.
如果細心去讀的時候.
你會看到他是很孤苦,很淒涼.
因為保羅這樣說.
他說你要講到我這女兒來.
因為底馬貪我現金的世界已經離棄我.
其他的人去了其他地方.
保羅就說我受審的時候.
沒有人在這裡陪我.
沒有人為我辯護.
底馬愛上這個世界.
他離棄保羅.
在其他文獻我們見到.
其實底馬這個人被描述為虛偽的人.
他以前和保羅同工.
其實好像和保羅很好.
其實他在這裡討好保羅.
他不是真的愛保羅的.
現在底馬成為保羅的敵人.
還有保羅接著說.
你要盡快在冬天來到之前.
幫我帶來外衣.
因為監牢很冷.
我想我還沒被行刑處死.
我會冷死.
保羅又說.
同仗亞歷山大害我.
但保羅又說神會報應他的.
四章十四節這樣說.
同仗亞歷山大多方害我.

$^{201}$主必照他所行的報應他.
你也要防備他.
因為他極力抗拒我們的話.
我們不知道亞歷山大做了什麼.
他是不是說了一些說話傷害保羅.
或者傷害了教會.
我們不知道.
不過我們知道他真的傷害了保羅.
而我們細心聽保羅這番說話.
你聽不聽到一些意思.
你聽到他苦澀的.
不開心的.
保羅說他傷害了我.
但神會懲罰他的.
在這個時刻.
當我們再聽多一點的時候.
好像是苦澀.
好像很bitter.
但其實你聽不聽見.
保羅的背後有更加深的意思.
他的意思是提摩泰.
現在不是復仇的時間.
現在是盼望的時間.
保羅知道提摩泰是非常勇敢的.
提摩泰知道保羅是受傷害.
做學生的一定是辯護老師.
氣亞歷山大甚至要報復.
所以保羅叫提摩泰不要理這件事.
不要自己去解決.
不要自己試圖解決這件事.
讓神去解決這件事.
相信神會解決這件事.
不知道你聽不聽到.
你會不會有保羅的想法呢.
當然我們信仰告訴我們.
有些事情是不對的.
就要快點做好.
我們信仰要求我們這樣做.
不對的事情就快點去解決.
但是神的智慧也讓我們知道.

$^{241}$哪些事情我們要快點去修復.
有些事情是我們修復不到的.
是用仰神,盼望神去修復的.
在我們人生道路的經驗.
不知道有沒有經驗.
有些事情我們想把它做好.
不過越做越差.
這個就是保羅的意思.
保羅說有些事情真的要留給神.
有些事情是要相信神能夠去改正.
我們不能夠糾正的事情.
所以我們看到保羅這麼艱難.
他沒有絕望.
他也沒有怨恨.
然後聖經這樣說.
17席他說.
唯有主站在我的身旁.
加給我的力量.
使我能把福音完整的傳出去.
讓所有的外邦人都聽見.
我有一重獅子口裡被救出來.
主必救我脫離一切的凶惡.
也必救我進他的天國.
願榮耀歸給他.
直到永永遠遠.
阿們.
所以保羅的臨別陣言是什麼?.
他說榮耀神.
永遠永遠的榮耀神.
他知道他優先的任務.
是繼續傳揚耶穌基督的福音.
他相信神的應許.
他相信天上的耶路撒冷.
所以他對生命是沒有盼望的.
所以我們看到原來保羅在艱難當中.
有盼望的.
問題是為什麼他這麼有盼望?.
為什麼這麼有盼望?.
他的盼望是來自於他的記憶.
這是其中一個解釋.

$^{281}$保羅不是天真 不是傻.
他記得的.
他記得過往受過種種的痛苦.
他記得.
他記得亞歷山大 底嗎哪些人害他.
他記得.
但是他也記得生命當中.
有一群很好的弟兄姊妹.
一群朝聖的弟兄姊妹.
一群的同伴.
他們是神賜給他的同伴.
在當中有很多的恩惠.
有神的愛 有人間的愛.
保羅說 我記得.
所以保羅列出這群弟兄姊妹的名字.
在11節他就這樣說.
他說 只有路加載我者女.
你來的時候 要將馬可帶來.
因為他在服事上 於我有益.
請向伯記拉 阿居拉和阿尼釋忽一家人問安.
伯記拉都在哥林多住了.
有保羅 布田尼 路加 土底亞和眾弟兄都向你們問安.
保羅說 只有路加在這裡.
路加沒有因為保羅坐牢而離開他.
還有馬可.
馬可就是說.
保羅說 要帶馬可來.
提摩塔要帶馬可來.
因為他對我的工作是有益 有份的.
其實我們都知道.
在傳道開始的時候.
保羅就說 馬可沒有用的.
這個年輕人是不行的.
不過後來他們又和好.
所以我們見到這個.
其實是一個很叫人著藥的時刻.
保羅提起馬可的時候.
因為人間那種和好.
是很難得的.
很珍貴.

$^{321}$保羅接著又說 阿尼釋忽一家人.
首先又說伯記拉這對夫婦.
他說 他們是和我一起做宣教的工作.
是很好的同胞.
接著他又說 阿尼釋忽一家人.
他說 我每次在家裡 我都很清新.
很開心見到他們 精神抖擻的.
他說 還有 還有就是.
另外一位弟兄就是這個.
爾拉多 他可能是哥林多的財務官.
以前保羅在特羅亞都是住在家裡的.
接著保羅列了一幫 一堆人的名字.
我們都不知道他們是什麼人.
由保羅 布田 尼盧 格魯底亞.
我們真的不知道他們是誰.
不過他知道 保羅在離世之前.
他記得這幫的朋友.
所以保羅去回顧他的一生.
他是記得一幫人和他朝聖.
因為這幫弟兄姊妹.
他對生命 對神 對人仍然有很大的盼望.
當然更重要的是他相信神.
他真的保持信心.
他對神有盼望.
他知道有公義的觀念為他存留.
他知道生命在世上的結束.
就是另一個的開始.
所以保羅是滿有盼望.
毫不絕望.
當我們見到保羅.
是和我們去數算他的名單的時候.
你想一下 有哪些人在你的名單.
在宣道會感受堂當中.
有哪個人在你的名單裡.
或者其實你可能會忘記.
那就寫下.
曾牧師 我的名單底下.
一見到那些壞人的時候.
很不高興的那些.
那就很不開心.

$^{361}$看一下名單.
曾牧師在這裡.
還有弟兄 有誰在這裡.
寫下名字 記住.
能夠記得是我們生命有盼望的土壤.
在我的朝聖的道路當中.
我也有我的名單.
我同時也在分享一些名單 一些名字.
其中一位就是曾牧師.
在我剛剛募會的時候.
我就遇上他.
在這間浸信的教會.
很有趣的.
這間浸信教會有一點點像宣道會.
為什麼呢.
因為他是執事會的主席.
是主任牧師.
人事會的主席也是主任牧師.
會務會的主席也是主任牧師.
可能還比感受更加 權責更大.
這位陳牧師在這間浸信會那裡.
他是義務的.
因為他在一個語文中心.
開了一個語文中心.
是教外籍人士學廣東話的.
他在這間教會一個月出現一次.
負責主餐,盛餐和開會,講道.
我畢業後就去這間教會服事.
在半年後.
他就叫我做他的助理會務幫他忙.
在過往的日子當中是怎樣的呢.
每個星期他都叫我要去油麻地.
在他的語文中心見他.
跟我聊天 聊個多小時.
每個星期都是這樣.
聊天的時候很鼓勵的.
比如說 曹先生.
我知道你很喜歡和弟兄姊妹們.
等他們下班跟他們聊天.
你這樣比做輔導好.

$^{401}$比輔導好.
那時候我很年輕.
牧師說我做到比輔導好.
哇 這麼好嗎.
很開心.
因為後來他比較忙.
他就少一點出現.
他請我一位年輕人做他的工作.
叫我帶領執事會.
去人事會裡面做主席.
在會務會做主席.
那時候我很年輕.
就在我開始事奉的時候.
就遇到一位這樣的牧師.
所以我是很平衡的.
我不是一般的年輕傳道人.
開始做的時候.
有很多的難處慢慢學.
我有位很好的牧師.
在這裡帶著我 教我 鼓勵我.
在我成長的另一個名單.
就是林道亮牧師.
2009年是我做院長的第一年.
那時候的神學院是有些問題的.
經過半年 神打掂一切.
我就可以很安定的去到洛杉磯.
去三藩市探望校友和探望教會.
第一站就是洛杉磯.
最後一天我就要離開.
有一位何牧師.
他說曹院長.
你必須一定要探望我們.
在洛杉磯教會的一個老前輩.
就是林道亮牧師.
他說他做過台灣中華福音神學院的院長.
做了十年.
你可以和他聊聊天.
我說好啊 我和他聊天.
我和他聊天.
我當然想聽聽和他說一下.

$^{441}$做院長的故事.
我去到的時候.
我想著他當然要和我說故事.
但是這位老牧師一開始.
他不和我說故事.
他一開始和我說聖經.
一開始說三段的聖經.
何牧師很尊重這位林牧師.
所以見到林牧師在這裡說聖經.
立即找一張紙巾寫著經文.
幫我寫下來.
有三段的經文.
創世記一章二十六節.
啟示錄三章二十一節.
馬太福音七章二十一節.
林牧師在這裡說聖經.
但是那時候我不是很喜歡.
因為我來這裡要聽故事.
不是說聖經.
因為這些聖經我都懂得.
我是要來聽故事.
不是聽聖經.
但是當我聽著聽著的時候.
我的心慢慢地軟化.
我心有少許的震撼.
所以我就聽話.
聽他留心說什麼話.
為什麼要和我說聖經.
說完聖經之後.
他就主動說.
我為你祈禱.
於是為我祈禱的時候.
何牧師也很殷勤.
幫我拍這張照片.
我自己也拍了這張照片.
否則我沒有這張照片.
我的心在這裡在這裡震.
接著林牧師說.
我要為你去幫忙禱告.
那時候我在這裡在這裡震.

$^{481}$我心裡整個人很震驚.
因為這位是一位陌生人.
是我素未謀面的人.
他是成為神差來的使者.
成為我屬命的阿爸.
教導我.
教導我要更加親近神.
更多地被神去感動.
是很奇妙的.
經過這一周.
那時候他99歲了.
我打算100歲做大壽的時候.
我再拿橙機見他.
不過他不在了.
不過他已經安息主懷了.
神就是這麼奇妙.
就是讓我做院長的時候.
就是透過林牧師.
他講聖經為我祈禱.
叫我更加全面地對著神.
用聖經的說話.
用禱告的說話.
是讓我見到我的驕傲.
讓我見到我更加多的親近神.
就是這位前輩.
決定了我怎樣做一個院長.
一個很何等奇妙的經歷.
2022年就是去年了.
去年可以說是我做院長的尾聲.
我就專心培養下一代.
過往是我處理不少行政的東西.
我說不處理了.
交給下一代處理了.
我就是做院長了.
做院長籌款.
是和教會更加深的交往.
不過我就加了一個任務.
就是培養下一代.
在2022年1月開始.
我就招聚年輕的牧者.

$^{521}$一批一批見他們.
聽一下他們事奉的經歷.
吸收一下他們的經驗.
究竟在香港事奉有多艱難.
年輕一代有多艱難.
他們怎樣想.
我聽一下.
從這班年輕的牧者聽一下.
我讀書比他們多.
讀得快的他們.
我講一下書本的東西.
讓他們參考一下.
書本的東西未必對.
不過知道也好.
我讀一下書告訴他們.
他們又講一下牧養的經驗給我們聽.
為什麼要這樣做.
因為我希望他們是更加的強大.
好像保羅提醒提摩塔一樣.
是更加強大.
好像保羅這樣講.
我記得你無畏的信心.
為者懸固.
我提醒你要把神.
藉著我按手給你的恩賜.
在如火添翁起來.
因為神賜給我們的.
是膽怯的心.
而是剛強仁愛自制的心.
我希望年輕的下一代牧者.
他留在香港的話.
他們一起打一場仗.
一起繼續跑.
最重要的是要持守這個信心.
接著在八月份.
我又開始一個叫院長茶室.
因為我事奉的教會主要是浸信教會.
我就容易招聚浸信教會的一些執事.
我和他們有院長的茶座.
我找一些年輕的執事.

$^{561}$我的做法是這樣的.
就是找七間教會.
七間不同的教會.
七個不同的執事.
他們素未謀面的.
配對一下他們.
一起喝茶,聊天.
他們說他們在教會做執事的艱難.
我吸收他們的經驗.
我又說一下我怎樣看法.
因為我去的教會多過他們.
我又會有一些神學的訓練.
我又可以和他們分享聖經的話語.
那裡有團契.
想法是這樣的.
喝一次茶.
喝得一般般.
不太合適.
下次不喝了.
喝得開心.
有第二輪,第三輪.
我和曾牧師喝茶喝得很開心.
今天又有第三輪喝茶了.
喝得不開心.
下次不用喝了.
不要那麼辛苦.
不要那麼艱難.
但是如果能夠遇到好的朋友.
素未謀面.
竟然都可以做朋友.
肯定見多幾次面.
喝多幾次茶.
為什麼會這樣做.
因為我們都在朝聖的路上.
我們要有一些朋友.
而這些朋友又不是我們自己去強求.
是神會給我們.
在人生的道路當中.
一定有很多失望的事情.
一定有些人會害你的.

$^{601}$有些人會傷害你的.
你都很難不記得的.
但是要記得.
不要沒有雄心.
神擺了很多的.
同伴,朝聖者.
神透過這些朋友.
這些弟兄姊妹們愛你.
神愛你.
因此你應該大有盼望.
我仍然有一個很深的信念.
縱然很多人說教會是沒什麼用.
我仍然相信這個世界.
是需要盼望.
而教會是我所知道.
最有盼望的地方.
若我們堅守盼望.
有伴同行.
阿們.
\newpage



\section{羅馬書 12:1-2}
\label{sec:dkl9UnCmWfQ}
\textbf{2019-01-20 加恩樓崇拜《得勝秘笈》羅馬書12\_1-2 講員:鄧泰康傳道}
\newline
\newline
連結: \href{https://youtube.com/watch?v=dkl9UnCmWfQ}{\texttt{https://youtube.com/watch?v=dkl9UnCmWfQ}} ~~~~ 語音日期: 2023-01-13
\newline
\newline
\hyperref[sec:uJ8mvXZ5iqw]{\small{< < < PREV SERMON < < <}}
~
\hyperref[sec:index_chronic]{\small{[返順時目]}}
~
\hyperref[sec:index_scriptual]{\small{[返順卷目]}}
~
\hyperref[sec:ZwtOgh6Qmm4]{\small{> > > NEXT SERMON > > >}}
\newline
\newline
羅馬書 12:1-2
\newline
\begin{longtable}{cl}
\hline
\hline
章節 & 經文 (和合本修訂版)\\
\hline
12:1 & \begin{tabularx}{0.7\textwidth}{X} 所以，弟兄們，我以神的慈悲勸你們，將身體獻上當作活祭，是聖潔的，是神所喜悅的，你們如此事奉乃是理所當然的。 \end{tabularx} \\ \\ \relax
12:2 & \begin{tabularx}{0.7\textwidth}{X} 不要效法這個世界，只要心意更新而變化，叫你們察驗何為神的善良、純全、可喜悅的旨意。 \end{tabularx} \\ \\
[1ex]
\hline
\hline
\end{longtable}
$^{1}$世界,只要心意更新而變化.
叫你們測驗何謂神的善良純全可喜悅的旨意.
我們現在將時間交給鄧泉去分享.
神的善良純全可喜悅的旨意.
鄧英姐妹早安.
和你身邊的人說一聲早安.
也說一聲主耶穌我是你.
又教大家第三招了.
平時我們見到人一定會說早安.
這是禮貌.
回來崇拜我們要和主耶穌說早安.
我們一起說主耶穌早安.
預備一二三.
主耶穌早安.
精神一點.
主耶穌早安.
記住我們每一天的裡面.
第一個要和祂說早安的不是你身邊的那個.
是主耶穌.
因為祂又賜了一天新的生命給大家.
今天和大家分享這段經文.
其實大家一點也不陌生.
甚至乎我相信有很多鄧英姐妹.
你應該懂得的金句.
你應該懂得背.
這段金句我們很鼓勵鄧英姐妹.
你真的要背了它.
所以弟兄們.
我以神的慈悲勸你們.
將身體獻上當作活祭.
是聖潔的.
是神所喜悅的.
你們如此事奉乃是理所當然的.
不要效法這個世界.
只要心意更新而變化.
叫你們測驗何謂神的善良.
純全可喜悅的旨意.
為什麼這麼重要.
因為羅馬書是一本非常非常重要的書.
如果你想認識你的罪.

$^{41}$你想認識人的真實情況.
你必須看羅馬書.
裡面很清楚和你說了.
罪如何影響我們的生命.
其實羅馬書裡面第一章至第十一章的時候.
保羅是和我們說到很深層的.
關於罪帶來給人類的嚴重性.
也都說到猶太人和外邦人之間.
其實在這些信仰裡面.
在上帝的裡面我們的關係如何.
因為為什麼.
當日的猶太人很看不起外邦人.
總覺得你們不行.
但他想告訴他們.
在上帝的救恩裡面.
從來上帝都沒有這樣分過.
在罪裡面大家一樣.
沒有一個人比另外一個人聖潔.
沒有一個人比另外一個人高級.
這個在羅馬書裡面說得很清楚.
當我們去面對這個問題的時候.
其實保羅亦都很坦然的.
和大家分享他生命裡面的實況.
保羅並不是一個高賢大智.
之後就和你說.
你看我今天多厲害.
你看你們多差.
這些是律法師的人.
但他不是.
他反而和大家在羅馬書裡面.
他和你很實在.
躺開心去告訴你.
我生命裡面的真實面.
有一段經文是很經典的.
你真的猜不到竟然是保羅說的.
他曾經說.
因為立志為善由得我.
只是行出來由不得我.
故此我所願意的善.
我反不做.

$^{81}$我所不願意的惡.
我倒去做.
若我去做我所不願意做的.
就不是我做的.
乃是住在我裡頭的罪做的.
我覺得有個律.
就是我願意為善的時候.
便有惡與我同在.
因為按著我裡面的意思.
我是喜歡神的律.
但我覺得肢體中另有一個律.
和我心中的律交戰.
把我擄去.
叫我服從那肢體中犯罪的律.
我真是苦.
誰能救我脫離者取死的身體呢.
我很敢說.
這是每一個基督徒.
今天你信了耶穌之後的掙扎.
是不是.
如果你問我.
我們有沒有辦法在這一天裡.
完完全全過.
完完全全聖潔的生活.
你猜有沒有可能.
其實是沒有可能的.
你會發覺原來罪很有趣.
罪是怎樣.
小小的地方便能進來.
但我就是想告訴你.
弟兄姊妹.
這才是我們今天基督徒的實況.
我經常說.
不要因為我們基督徒.
要在頭上交個光環.
我們好像比外面的人.
好一點.
不是.
不是,我想弄清楚.
我們應該要好一點.

$^{121}$你明白我的意思嗎.
因為靠著上帝我們應該要好的.
但實際我們是不是好一點.
不一定.
我想說.
今天我們每一個人的生命裡.
本身就是這樣.
保羅也說.
連保羅.
我們看他和上帝這麼親近.
他也說我就是這樣.
我們就坦然去承認.
所以我經常說我們基督教裡面.
有一個很寶貴的地方是什麼.
我們是蒙恩的罪人.
我們今天能夠站出來.
是因為有主耶穌.
祂給的恩典.
祂十字架的寶血.
洗淨了我們.
所以我們不需要再為罪的緣故.
我們經常不開心.
我們不需要再去責備自己.
因為上帝已經赦免我們.
頂子梅.
Hallelujah.
這才是上帝給我們的恩典.
所以你會發覺.
當保羅說這些的時候.
他要先和你從第一章到第十一章.
很清楚地說原來根本罪就是這樣釋放.
但是.
我想告訴大家.
當他一去到十一章尾.
33節開始.
他說了一番說話.
就進入第十二章.
請你記得這番說話.
在他剛才這麼軟弱的情況下.
他問了一個問題.

$^{161}$誰能夠救我脫離這個取死的身體.
但是他說.
心哉.
神豐富的智慧和知識.
他的判斷何其難測.
他的蹤跡何其難尋.
誰知道主的心.
誰作過他的謀事.
誰是先給了他.
使他後來償還.
因為萬有都本於他.
依靠他.
為他.
願榮耀歸給他.
直到永遠.
阿們.
保羅他用一個態度.
全然謙卑.
來到上帝的面前.
告訴所有人.
告訴自己.
今天一切都是上帝.
不是你做了什麼好.
不是你做了什麼.
所以我去拯救你.
是上帝白白拯救你.
然後他就告訴你.
真正的出路在哪裡.
在上帝那裡.
所以我經常說.
基督教裡面這個很重要.
弟兄姊妹.
當你今天不是去擁抱上帝.
你不是和上帝建立一個很緊密的關係的話.
我想告訴你.
你搞不定的.
然後第十二章第三節開始.
他就跟你說.
其實基督徒的生命應該怎樣.
所以第一節和第二節是很重要的.

$^{201}$弟兄姊妹.
這個是由你舊生命裡面.
進入新生命裡面的一個很重要的關卡.
你能夠按著第一節和第二節裡面.
你能夠做到.
你能夠向著這個方向做.
我想告訴你.
你就能夠順利進入基督徒的生命裡面.
你不會再虛假.
剛才那裡說.
其實我什麼都不是.
根本我在神的面前.
我就不需要去裝假.
我經常都說.
基督徒最可怕的地方.
你千萬不要裝假.
我在你面前的時候.
哇,我是一個怎樣好的人.
哇,原來背對背.
我是多麼的污穢.
我現在這麼骯髒.
我裡面的時間其實有多少的問題.
這個是最大的問題.
當我們能夠明白第一節和第二節之後.
你就能夠很順利進入第三節開始.
怎樣去學做一個基督徒.
所以我今天給你這段經文.
叫做得勝必及.
原因是因為.
這個真的會令到我們得勝.
但是你記住.
是上帝的話語.
當你行出來的時候.
你就得勝.
為什麼你不敢說.
一開始我以神的慈悲勸你們.
你記住.
是在剛才的一至十一章的時候.
因為上帝對大家做了這麼多的救恩.
給了這麼多的恩典給大家.

$^{241}$所以記住.
我們的能力永遠都是從哪裡來.
十字架.
我經常說我們信的是一個很古舊的十字架.
但是一個萬世都能夠用的十字架.
因為我們全部的力量就是憑十字架.
所以.
我以神的慈悲勸你們.
我們今天每一個人做一件事.
我們都需要一個動機.
我想問問大家.
今天大家為什麼要上班.
為什麼你們要上班.
你不是張活柔吧.
有些人不認識張活柔了.
我小時候看張活柔的節目.
我為人民服務.
是不是.
是不是這麼崇高.
有些是的.
但大部分不是.
為什麼你們要上班.
賺錢.
你試試你老闆不給你薪水.
你上不上班.
是不是.
你明白我的意思嗎.
每一個人做一件事有一個動機.
你今天為什麼我們上教會.
我們有時沒有力量.
我們上來教會都沒有動機.
為什麼上.
敬拜上帝.
敬拜上帝是什麼.
你試試一直問下去.
你再解釋一下.
有時我們原來自己做什麼都不知道.
我過個星期習慣了.
我上來教會.
我敬拜上帝.

$^{281}$敬拜上帝.
用你的心靈和誠實去敬拜他.
你記得嗎.
你人坐在這裡是不是代表你正在敬拜他.
不是的.
你弄清楚.
所以你會發覺這個動機很重要.
原來這個是動機.
因為是神慈悲的緣故.
所以我們要做以下的事.
什麼是上帝的慈悲.
就是今天上帝.
將他的救恩.
神的憐憫.
將主耶穌都已經賜給我們.
這個就是我們的動機.
他用一個無限的愛去愛你.
他用一個你本身不配的去愛你.
這個就是我們的動機.
你就會發覺一件事.
我們今天有動機.
才知道我們怎樣走.
而很實際.
我不想說我們有耶穌的愛.
怎樣愛呢.
我經常鼓勵大家學講見證.
每一個裡面你都找一下上帝怎樣愛你.
你才有動力.
而接著最實在的就是.
請你去回望.
去回望.
你在昔日裡面.
上帝怎樣去愛你.
其實這個是非常重要.
第一件事.
在舊約他們有個習慣.
叫做納石遺跡.
去到每一個很重要的地方.
他們就會納石.
就會在那個地方堆一堆石頭.

$^{321}$很清楚地告訴你.
我們要叫我們的後代.
見到這堆石頭.
就記得上帝的恩典.
是不是呀.
靜姐妹.
大家有沒有這個習慣.
其實我想教大家一樣東西.
就是今天現代很多東西.
你不要讓它去玩你.
你去利用它.
很簡單.
今天有Facebook.
我不知道大家喜歡放什麼在Facebook.
最近在玩一個十年挑戰.
十年前什麼樣子.
十年後什麼樣子.
現在在玩這個.
但反過來其實這個Facebook.
我覺得我想反過來.
你可以用的.
用來做什麼.
你有沒有留意.
你放在上面全部都是.
如果你想放在上面.
全部都只不過是.
我在幾年幾月.
去到哪裡吃什麼呢.
其實是沒有用的.
它只是會令到你更加想放假.
你有沒有發覺.
但如果你在裡面.
全部都是承載著.
上帝所給你的恩典.
Facebook有個功能.
很好的.
是不是.
在多少年之前的今天.
你發生了什麼事.
它會彈出來.

$^{361}$你想像一下多漂亮.
我已經忘記了.
人的記性很奇怪.
我們專門記住一些不開心的事.
有沒有留意.
我們都開心忘記了.
所以叫它數算恩典.
大家就啞了.
一說這些不開心的事.
可以說一個小時.
你就會發覺.
原來這個功能.
很好的幫到大家.
就是什麼.
彈回十年前.
是喔 我都已經忘記了.
十年前原來我曾經去過那個地方.
我曾經去過退休.
原來在那裡我曾經.
和上帝去納了一個約.
是不是.
原來上帝曾經.
在多少年前的時候.
我一個病的裡面.
祂曾經怎樣去醫治過我的.
頂展梅.
當你的生命越儲這些東西的時候.
你就越見到.
上帝在你生命當中.
實際的慈悲.
是不是.
你就能夠有這種動力.
是喔 上帝.
因為你這麼愛我的緣故.
所以什麼.
我要去愛你.
我要去回報你.
頂展梅.
我們今天有沒有將上帝.
這一點一滴的東西.

$^{401}$去放下去.
成為納石遺跡.
有沒有記住上帝的恩典.
當你越是記得的時候.
你就越願意去回應上帝.
很實在.
主耶穌今天為你死.
今天我們是一群得了救恩的人.
你願意為祂做什麼.
我經常都說這個不是一個等價的交換.
有沒有發覺.
就算你傾盡了你一生所有的東西.
你所有的錢財.
你都不夠回報給上帝.
但上帝每次當你能夠回報祂.
祂就已經很快樂.
祂就已經很開心.
因為祂就見到你生命裡面的時候.
你能夠怎樣去跟隨祂.
是不是.
這個是非常重要的.
第二個.
當我們知道如果上帝.
原來我們跟隨祂是有一個動機之後.
第二個是要求.
好了.
我們要跟隨上帝.
上帝要什麼要求.
將身體獻上當作活祭.
是聖潔的.
是神所喜悅的.
你們如此事奉乃是理所當然的.
獻上.
它的意思是誠心和專注.
今天我們回教會.
我們去跟隨上帝.
頂智妹.
你夠不夠誠心和專注.
我要很實在的說.
有時候我不知道大家的祈禱內容.

$^{441}$記住.
鄧先生說的這些.
永遠都不是一糟糕打一算人.
不是的.
我們想說自己問自己.
我不知道你的祈禱內容是怎樣的.
你自己想想.
多數是問上帝要東西.
還是你願意幫上帝做事.
你會發覺一件事.
其實今天我們在神的面前.
我們夠不夠誠心.
有時候我去到一些地方參觀.
特別以前.
我也會的.
特別是外國.
有些人的廟.
我也會進去.
看看.
我真的覺得.
有些其他宗教人士比我們誠心.
大家知不知道拉薩.
他們的人是朝聖的.
大家有沒有看過那些.
戴上那些被子.
膝蓋和手.
他們是怎樣的.
我做不到給你看.
我也想示範一次給你看.
走兩步.
跪下.
趴在那裡.
跪拜一次.
然後上.
他們說的就是上.
然後再起身.
他們從他們的家走到他們所說的聖地.
是說了多久.
一年.
兩年.

$^{481}$三年.
我們在很附近的地方.
大家有沒有聽過一句說話.
中國人很好的一句說話.
叫品神.
有沒有聽過.
什麼叫品神.
就是你真的去到你覺得的神面前.
你的菩薩面前.
然後你就合十.
然後有時間你就在面前說什麼.
一講一輪.
保佑我.
又要什麼.
我不懂的四字詞.
什麼身體健康.
財源廣進.
合格平安.
但他起碼.
大家有沒有見到每一年過年的時候.
投注香的那一刻.
大家有沒有留意.
香真的會對到人的那些.
是吧.
他們真的.
我們說人家是假神.
我們說人家什麼都好.
但是他那種誠心.
無論怎樣都是那種誠心.
他真的去到.
是真的晚上很晚睡.
去到排隊擠.
我都要去裝我的投注香.
弟兄姊妹.
今天我們回來崇拜.
我們有沒有這份心.
我想說是每一間教會都有同一個問題.
都有同一個問題.
就是永遠在講台上面.
我們都不斷呼籲.

$^{521}$是什麼.
弟兄姊妹早點到.
是吧.
所以永遠.
任何一間教會.
你記住.
不只是我們教會.
基本上所有教會差不多都是這樣.
就是開始的時候.
開始崇拜之後的五分鐘.
大概裡面都是三分一人.
然後慢慢慢慢就會來.
差不多我們去到十分鐘.
就差不多坐齊了.
弟兄姊妹.
我們今天在我們的信仰裡面.
我們有沒有這份這樣的誠心.
專注.
上帝要求我們將身體獻上.
第二.
這個祭是全人的.
弟兄姊妹.
是活的.
是聖潔的.
全人的意思就是說.
身心靈.
整個獻給他.
是吧.
其實我們今天最怕的那件事就是.
我們今天.
我經常都講人在心不在.
我們人回來了.
但是我們的心.
我們是不是敬拜上帝.
我們在想其他東西.
甚至我以前聽牧師說.
牧師說.
其實弟兄姊妹聽我一堂道的時候.
已經在想待會去那裡吃飯.
是吧.

$^{561}$我們今天有沒有將我們整個人.
我們全部的獻給上帝.
還有一件更加重要.
我也很欣賞那位牧者提醒我.
他說其實我們今天的信仰.
先是全職.
我們上班是兼職.
什麼意思.
就是說我們今天從起床那一刻.
你已經是活著這個祭.
由你起床的時候.
你用什麼態度去煮你的早餐.
你用什麼態度出門.
你用什麼態度回到公司.
你用什麼態度完成一天的工作.
你用什麼態度下班.
你用什麼態度回到家.
去和你的家人相處.
全日都是上帝的.
是不是弟兄姊妹.
不是分割的.
我們今天很可怕的事.
就是有很多信徒.
他星期天回來的時候.
我是一個信徒.
一出去的時候.
對不起我不是.
我是外面的什麼什麼.
所以我記得上次我講這裡.
我都說有一個人可以隱藏了幾十年.
在他身邊.
一個上師.
到他有事的時候.
差不多要走了.
要離職的時候.
你放心我為你祈禱.
你是驕徒嗎.
我認識你三十多年.
你是驕徒嗎.
是不是很可怕的一件事.

$^{601}$弟兄姊妹.
全日.
我們要將我們整個人獻給上帝.
另外是活著的.
為什麼.
因為從前舊約的時候.
獻的祭是死祭.
是用生畜的.
我們用牛的.
一次過的.
但主耶穌是活著的.
我想大家留意.
主耶穌的活祭並不是自免.
而是他真是復活.
其實他在地上就告訴你.
當你信他的時候.
你就有這個復活的盼望.
而復活的盼望並不是你死了之後.
復活的盼望是你信耶穌的時候開始.
你就可以死了.
重新成為一個新做的人.
是活的.
是活到出來的.
弟兄姊妹.
所以我在之前教會.
不同的情況下.
我們牧師幫助受浸.
但我們這些傳道.
就會幫手下詞.
講得好聽.
就幫手扶.
扶一扶人.
有時候浸的時候.
怕他起不來.
以前我在我的教會.
會教浸禮班.
自己負責自己.
因為我當時是堂主.
我跟他們說.
你這次有福了.

$^{641}$這次是我幫手扶.
但我說你放心.
我不會這麼快放你上來.
我說到時你吸多些氣.
浸多兩秒.
為什麼.
我跟他說.
你感受一下.
死了的感覺.
我認真的.
因為這個在水裡面.
你躺下去的時候.
是什麼.
我死了.
我上來就應該是一個全新的人.
所以我跟他說.
你感受一下.
所以我不是扶那個人.
我是拉那個人.
我拉他下去.
這樣的情況下.
我就會發覺原來.
整個生命當他今天活出來.
一個活祭才有意義.
如果你硬是跟人說.
不是啊我死了之後上天堂.
我們都相信的.
我們很肯定.
這個聖經裡面的教導.
但突然之間我想問.
你又誰見過.
死了之後.
他回來告訴你.
我上了天堂.
是很難信服的.
你明白我的意思嗎.
但實實在在的.
就是今天我們這個活著的人.
我們活出一個新的生命.
讓別人見到不同.

$^{681}$所以我真是.
honest我們以前的教會裡面.
我們去服侍一班吸毒的.
以前的朋友.
到後來他戒了毒之後.
我們見到他生命裡面.
是180度的轉變.
他真是由以前那個.
嘩 粗口濫竊.
之後他後來真是變回好的.
我以前服侍一個基層.
所以我有很多機會.
接觸黑社會人物.
有一次因為我們教會.
有一個小朋友是跟黑社會的.
有一次他叫我.
鄧先生.
鄧先生我回到教會這麼久.
其實我不想再回黑社會.
我想做一個好人.
我說好啊 沒問題.
不過你要見我大哥.
真的 我認真的.
他說你要見我大哥.
好啊 去吧.
我就和他一起去見他.
黑社會大哥.
之後他一見到我.
他說你去哪裡.
我跟他說我去耶穌.
什麼耶穌.
我說我教會的.
我是傳道.
很有趣的.
你不要這樣看那些黑道人.
他知道你是宗教人士.
他會恭恭敬敬的.
是啊 很有趣的.
之後他就和我說.
什麼事.

$^{721}$我說因為這樣.
這個弟兄.
這個弟兄現在回到教會.
他不想再參與黑社會活動.
這樣的情況.
我說他現在想叫我和你講一聲.
好了 這樣的情況之後.
他望一望我.
望一望那個弟弟.
我想說真的猜不到.
其實我見到他的時候.
我有一點心跳.
心跳是啵啵啵啵的.
之後他見到我之後.
望一望他.
望一望我.
接著他說.
這樣啊 好啊.
他和那個弟弟說.
你又乖乖的.
等我知道你再出來玩的話.
我打你.
你現在跟著傳道.
好好跟著傳道.
以後你乖乖的做個好人.
我是你黑社會大哥.
之後很有趣的.
後面的時間他還說了一件事.
很有趣的.
本來其他人.
他講的術語.
過檔又不知道要封多少利是.
這樣的情況之後.
他說你是信耶穌的.
不用了.
阿免了.
我就是這樣帶著弟弟.
走了.
你會發現一件事.
但我後來跟這個黑社會大哥.

$^{761}$還有聯絡.
我有時候約出來.
喝下茶.
他後來.
真的後來回到我們教會的報道會.
他後來真的缺字.
但發生了一件事.
他說鄧先生.
我有一點問題.
我說什麼問題.
我也很想跟耶穌祈禱.
不過我習慣了說髒話.
我又怕我祈禱的時候.
我沒理由說髒話跟耶穌說.
那怎麼辦呢.
我跟他說你放心吧.
耶穌是按你的本相去愛你的.
你說髒話算什麼呢.
但你要學一件事.
不要胡亂說.
你祈禱的時候.
你講的時候停一停.
講著講著又停一停.
慢慢說.
我試給你看.
他真的有很多髒話.
耶穌求你X我.
類似很多.
我跟他說不要緊.
試試吧.
他很好.
有時候.
黑社會的人很有趣.
當你跟他接觸.
他認真的時候.
他真的很認真的.
很小的.
他真的試.
然後過了一個月.
黑社會過了一個月之後.

$^{801}$有一個弟兄回來找我.
我很開心.
我說什麼事.
我戒了髒話.
我現在一個字都沒有.
之後他就說.
我現在很開心.
很釋放.
之後我就跟他說.
你現在連髒話都戒了.
你是不是應該跟黑社會都戒了呢.
當然他就說他們的事.
我們什麼小朋友跟著我.
我說這些是藉口.
我說你放下就放下.
當然有他們的方法.
又不知道怎樣將小朋友過當.
又不知道給了誰.
總之他就走出來.
後來他也成為我們報道團裡.
其中一個很重要的人.
我們每次去做報道會.
去跟學生報道就找他.
所以你就會發覺.
原來上帝要的是活的生命.
原來一個活的生命.
他真的活出來告訴你.
活給你看到.
原來耶穌那種勁是什麼.
耶穌那種生命是什麼.
而另外更加重要的是什麼呢.
上帝要我們這樣事奉.
是理所當然.
所以剛才我說.
今天上帝什麼叫理所當然.
不是我們中國人所說.
應該不是.
原來原文裡面.
原來中文譯本是這樣譯.
他說這就是你們應該獻上的真實敬拜.

$^{841}$女真宗譯本是這樣譯.
是你們的心的事奉.
原來是你整個人.
原來這種這樣的事奉.
這種生命是你整個人.
真實每天每夜的敬拜.
原來你每天就是站在神的面前.
用你的生命去敬拜他.
弟兄姊妹.
用你的生命去敬拜他.
並不是你唱詩歌.
並不是你祈禱的時候.
吃飯的時候也可以.
所以曾經有些年青人問我一個很有趣的問題.
鄧先生如果我在廁所裡祈禱.
是不是不太好.
接著我就跟他說.
有問題嗎.
是你覺得有問題.
上帝不會介意的.
所以你會發覺原來很多時候.
在我們心裡面綁著我們自己.
弟兄姊妹.
今天上帝對你的要求.
就是將整個人專心獻上.
做活祭.
是聖潔的.
你要將你整個人獻給他.
跟著這個是應該的.
是整個的新神的事奉.
交給上帝.
有了這些之後.
接著第二節.
就是得勝必救.
不要效法這個世界.
只要心意更新而變化.
叫你們察現何謂神的善良.
純全可喜悅的旨意.
其實裡面說得白一點是什麼.
不要再按這個世界的價值觀.

$^{881}$去過生活.
用屬天的價值觀去過生活.
其實這個是最挑戰我們基督徒的.
說真的我們每一個人.
每天都有很多的追求.
為什麼.
人都是追求優秀的.
人都是追求好的.
我從來沒有叫過大家不追求優秀.
我從來沒有叫過大家不賺錢.
我從來沒有叫過大家.
你不去追求好.
從來沒有.
但是問題就是.
這些東西是不是就等於全部呢.
不是.
在天國價值觀很清楚.
你若愛你身邊的父母兄弟姐妹.
物質.
而勝過我的話.
你就不配做我的門徒.
我經常說的一件事很簡單.
你不要用這個世界的價值觀.
去成為你的全部就夠了.
所以不要效法這個世界.
它的意思就是.
在李鎮忠的譯本裡面.
別和今世同形狀.
其實很實際.
不要和這個世界上同一個形狀.
如果這個世界上是正方形.
你就不要正方形.
你要三角形.
就是這個意思.
其實說得再實際一點.
落地一點.
其實我們今天說的價值觀.
走不開那幾樣東西.
衣食住行.
穿衣服.

$^{921}$頂紙媒.
我想問一下這個世界.
穿衣服的價值觀是什麼.
是什麼.
有沒有這個世界的價值觀穿衣服.
當然有.
不然姐妹怎麼每一季都要換衣服.
因為設計師告訴你.
你的價值觀是什麼.
這一季流行什麼.
我們就要穿那個.
我們就不落伍了.
我們這樣穿就老套了.
你記住.
是世界給你的.
並不是你自己喜不喜歡.
原本上帝創造這個世界的時候.
我想問一下大家.
當下有沒有衣服穿.
沒有.
我想告訴你.
當然我不是鼓勵大家不穿衣服.
你搞清楚.
我是想說的是.
原來上帝很清楚告訴你.
什麼才是最重要.
是你這個人.
當你將所有外表的東西都脫光.
那個就是你.
那個能不能見人.
很多時候我們今天的標準是不能見人.
所以我昨天在Allen的婚禮也說過.
今天的人很厲害.
為了減少兩條愚昧紋.
死也可以.
欠錢也可以.
什麼都敢做.
但我經常說一件事.
年老是上帝給我們的.
我現在的字眼改一改.

$^{961}$現在叫金陵.
現在叫銀髮一族.
有沒有留意.
其實我覺得他用這個是很好的.
為什麼.
拿出價值.
你有你的價值.
上帝的價值.
你才是最重要.
不是衣服.
不是你穿什麼出來.
不是你穿什麼名牌.
不是你今天穿了什麼.
整個世界都覺得你好看.
不是的,弟兄姊妹.
是你的人.
你才是最重要.
我和你相處不是因為你穿得好看.
是因為你的人.
吃.
今天我想問大家.
吃東西為了什麼.
這些是幼稚園問題.
大家不要覺得我在玩什麼.
什麼老緊急轉彎.
我不喜歡這些.
吃東西為了什麼.
吃麵.
會不會有了飽.
弟兄姊妹.
有錢.
如果上帝給你有錢.
吃什麼.
吃豐富一點不要緊.
沒錢.
吃什麼.
麵包.
煮公仔麵.
都飽了.
但是今天人將吃東西變成了什麼.

$^{1001}$一種等次.
我這樣形容.
大家都吃過甜品.
一個個小的甜品.
用一碟普通的餐廳裝.
多少錢.
幾十元.
我當你幾十元.
你在半島酒店.
那幾件甜品.
用個假裝出來.
擺在你前面.
多少錢.
幾百元.
這樣是高檔一點.
這樣是人眼中.
覺得你高級一點.
並不是上帝.
吃東西是上帝給你享受.
我經常問.
今天大家.
你不需要去比較吃什麼.
有些人喜歡吃素.
就給他吃素.
為什麼說人家不吃肉.
有些人喜歡吃肉.
為什麼要人家硬吃素.
你明白我的意思嗎.
我享受就行.
那樣東西我享受.
所以其實很簡單.
我也說.
我很清楚.
如果將來大家有機會.
跟我吃飯.
我真的不吃魚.
我不吃那些.
特別是蒸到滑溜溜的.
什麼魚骨.
我是最怕的.

$^{1041}$我不知道為什麼.
我不知道怎麼解釋給你聽.
但問題是.
總而言之.
我不吃.
那我就不吃.
你不用不斷勸我.
不是啊.
魚有多好吃.
你講到怎樣好吃.
我都不吃.
但問題是.
我想說最重要是享受.
上帝今天給你.
每一個人都有你享受的範圍.
享受這件事就行了.
住.
住的時間.
我記得我剛結婚.
我和太太當時沒有什麼錢.
當時沒有這個term.
今天用這個相對.
我們租了一間雄厚的tong 房.
其實是挺骨子的.
有地毯.
然後有一個area.
然後有雪櫃.
一個放電視機的架.
然後一張床.
一個地方.
旁邊有廚房 廁所.
我記得大概二百多呎.
我剛結婚的時候住的地方.
但我想告訴你.
很開心的.
而且有新婚.
每次一進屋.
一見到我老婆.
說什麼.
老婆我回來了.

$^{1081}$她說老公.
大招呼.
我們坐在地下.
坐在地上.
地毯.
抱著兩個cushion.
兩個人在.
電視這麼近.
兩個人在播電視.
然後在講.
一坐下就講.
這次公司發生什麼事.
我們有很多溝通.
很開心的.
過了八個月.
我們覺得不行.
要搬去一個地方.
因為我太太的身體問題.
然後我們搬去屯門.
由二百多呎搬去五百多呎.
的兩房一廳.
兩個人住兩房.
你怎麼安排.
我想你和我差不多.
一個睡房.
另外一個做什麼.
書房.
其實書房書什麼.
即是男人打機房.
OK.
是不是.
放電腦在裡面.
好了.
下班.
你去哪.
打機.
換完衣服之後.
進去書房.
躲在裡面.
老婆就在外面.

$^{1121}$在外面弄東西.
回到家兩個小時.
都不說一句話.
後來搬到三房.
一廳兩廳.
重大件事.
每人都有間房.
鬥打機.
真的.
因為看你玩得這麼開心.
就去玩.
是不是.
每人一部電腦.
一回去之後.
他玩他的我玩我的.
吃飯的時候.
我趕快吃完.
然後兩個人又回去打機.
弟兄姊妹.
我們今天追求什麼.
我們今天是一個很舒適的環境.
但那個究竟是不是上帝所說的家.
住的地方.
上帝說的是一個家庭.
家庭裡面是說關係.
地方大了.
反而讓我們的關係疏離了.
當然我同樣都是.
我不會叫大家搬小屋.
OK.
但反過來.
你需要知道.
原來一個家裡面.
需要的是什麼.
是兩個人.
一家人當中.
那種關係的融合.
當我們能夠反正之後.
後來第一.
我們真的搬小屋.

$^{1161}$不要那麼大了.
因為不太需要.
第二.
我們真的每一晚.
真的有那個家庭崇拜.
我們有一個一起祈禱.
還有那個時間真的聊天的時候.
很官方.
因為每人都要分享.
你個人的代表事項.
每人大家分享.
讓我們能夠重新建立這個家.
行.
是不是.
賣車.
我帶你買的第一輛車.
是一輛二手車.
是一輛日本的前七.
但之後的時間.
我那時候很開心.
買第一輛車.
是很興奮的.
嘩.
那輛車很好打理.
是吧.
好了.
用了一輪了.
然後怎樣.
是不是就這樣就OK.
當然不是.
是不是.
然後就要上賣新車了.
上賣新車.
你考慮.
你買日本車.
還是買.
賓士呢.
我不知道.
弟兄姊妹.
你怎樣去考慮.

$^{1201}$你買日本車.
還是買賓士呢.
為什麼你要買賓士呢.
你記住.
我不斷地說.
都不是跟大家說.
買賓士是不是錯.
不是.
你搞清楚.
我要回到.
傳道書說.
上帝給你什麼份.
你用什麼份.
OK.
上帝給今天大家有能力.
你住大一點.
就住大一點.
給你賓士就賓士.
甚至你勞銷.
勞銷都可以.
是不是.
但問題.
重點是.
你買賓士.
我為了要跟別人攀比.
人家現在那輛是C Class.
不行.
我要買一輛S Class.
是不是.
人家那輛.
現在都.
現在賓士都不夠威風.
我要買一輛Tesla.
是不是.
一出來之後.
人家說整間屋都在.
這樣.
弟兄姊妹.
最重要的就是.
究竟我們想為什麼.

$^{1241}$有一次我和一個弟兄姊妹.
我記得很多年.
十幾年前.
我真的問她問題.
我很爽問她問題.
為什麼你買賓士.
我弟兄很直接跟我說.
我不是因為要炫耀.
他講一個例子給我聽.
他以前在美國.
他有時參加日本車.
有一次車禍.
他和別人撞車.
一撞之後.
人家的車.
即是人家的賓士.
就凹了一個泵把.
他那輛日本車摺起來.
之後他很害怕.
自此之後.
他說我真的要買這些車.
硬正.
因為不會出意外.
當然我理解了.
我不是跟他買廣告.
夠清楚.
但我想說原來動機很重要.
在神的裡面.
他不是不給你用.
是你心裡面藏著的是什麼.
我所做的每一樣事.
我不斷和別人比較.
我不斷和別人.
要顯我自己那種富裕.
還是我只是為了我的生活.
所以你會發覺.
原來世界價值觀.
和我們價值觀很不同.
上帝要的是回到你的身上.
你才最重要.

$^{1281}$另外關於戀愛和婚姻.
你會發覺.
這個世界裡面的戀愛和婚姻.
所說的是佔有.
所以有很多今天不認識的人.
即是不認識信仰的人.
很簡單.
你愛不愛我.
你陪我就是愛我.
你經常聽我說就是愛我.
但聖經說的不是.
聖經是在說付出.
是在說犧牲.
聖經的要求是什麼.
我經常說.
女人要馴服丈夫.
這是上帝所創造的原則.
但請你記得背後是什麼.
丈夫可以為老婆死都可以.
你要基督愛教會.
你記住這個圖畫.
老公可以為老婆死都可以.
如果有一個老公可以為老婆死都可以.
老婆怎會不聽話.
你明白我的意思嗎.
原來是在說這種圖畫.
另外生存.
我們每一個人究竟是為自我實現滿足.
還是為了滿足上帝.
我們每一個人都有我們的追求.
但我經常提到一件事.
每一樣東西我們都會失去.
你就會發覺原來做一切的事.
我們今天改變我們的價值觀.
當你能夠去改變.
你能夠回到上帝價值觀裡面的時候.
最後一句說話很重要.
你就能夠測驗什麼是神的旨意.
什麼是美善.
什麼是喜歡.

$^{1321}$什麼是純全.
但鄧英姐妹.
原來很有趣.
有時候你經常問我.
希望我來到嘉欣樓.
大家不要當我是廟祝.
有什麼時間就問一些宗教問題.
就找鄧泉.
我現在那裡有兩個人.
我應該選A還是B.
誰才是神的心意.
你這樣問我.
我真的想說我會滾開地下.
難道我在你面前摘聖杯嗎.
我不會的.
其實我就回到聖經那裡告訴你.
很簡單.
如果一個信耶穌一個不信耶穌.
那上帝選誰呢.
當然選信耶穌.
因為聖經寫的.
不是的.
我意思是.
我不是叫你一定只選那個.
所以我經常提到鄧英姐妹.
有很多是好人.
不過不信耶穌.
但沒有人叫你這麼快拍拖.
你不用急於拍拖.
你先和我做個普通朋友.
你把信仰介紹給她.
你帶她信耶穌.
讓她慢慢和上帝建立關係.
到時再談拍拖.
其實是一個過程.
你不用急於這樣.
你會發覺原來我們要知道上帝的心意.
就是如果你心裡已經有上帝的價值觀.
你一看就知道.
你就知道上帝想怎樣.

$^{1361}$你就很容易分辨到上帝喜歡還是不喜歡.
是不是鄧英姐妹.
所以你就會發覺一個很有趣的圖畫.
當你那段時間和上帝的關係越近.
你會發覺你的生命越有那種動力.
因為你會發覺你和上帝在一起.
這個是大家都經歷過.
但很多時我們也會失落的地方.
所以鄧英姐妹.
我真的很盼望在神的裡面.
大家能夠回到上帝的裡面.
好像剛才那兩段經文所說的一樣.
你有動機.
你知道上帝要甚麼.
之後你回到上帝的價值觀裡面.
用上帝的價值觀去過你的生活.
我可以擔保.
你一定會有一個更豐盛.
更喜樂的生命在你裡面.
你一定會遇到任何困難.
你都不會擔心.
因為你會發覺當你去所有的一切.
我舉個例子.
譬如當你今天有你自己.
或者你的親人在疾病裡面.
我們有時很無望.
我們也做不到甚麼.
就算我想聽你是醫生都一樣.
但反過來你會發覺.
如果你回到神的價值觀.
原來生命是上帝掌管的.
上帝要他留,他怎樣都不會走.
上帝要他走,你怎樣都留不到.
你能夠回到這裡的時候.
你就將他全東西交給上帝.
你能夠管好的是甚麼.
不是擔心,弟兄姊妹.
還有時間,上帝還未接他走.
是不是.
每天怎樣好好地陪伴他.

$^{1401}$一起面對這個疾病.
對不對.
好好地和他一起經歷上帝的同在.
對不對,弟兄姊妹.
當你能夠有上帝這種價值觀.
你就面對一切一切的困難.
所以我想最後補一句.
我從來沒有說過.
當你想通就能做到.
從來都不是.
但弟兄姊妹.
每一樣東西都是操練出來的.
是看你是不是習慣.
這樣去用上帝價值觀.
用他的心思意念去想.
當你這樣想的時候.
你每一次就會越快.
有這個回到上帝裡面的反應.
你越快就會回到上帝面前.
你找到那種盼望.
你就會越是能夠有力量.
所以求神幫助大家.
真的能夠在這個得心比給裡面.
能夠去想到.
最後的時候我想說.
我在09年第一次聽到一首歌.
叫做《恩作信》.
那首歌裡面有機會大家可以回去聽.
詩歌裡面說了很多歷史人物.
說了亞伯拉罕.
說了摩西.
說了達爾尼.
說了約書亞.
說了很多很多很多.
其實裡面一直都說一件事.
這班人其實全部都是普通人.
但因為他們信上帝的緣故.
經歷了很多很多很多的神蹟.
如果你問我.
究竟這個世界上有沒有神蹟.

$^{1441}$我想告訴你.
我是相信.
我相信.
但我不會高舉神蹟.
有一次.
因為我未來這裡之前.
其實我當時經常會回國內.
去做一些教導和宣教工作.
有一次我去到一個很窮的家庭.
如果你想知道什麼叫家圖四壁.
那裡就可以讓你看到.
那是一間農村的屋.
牆身沒有任何的粉色.
全部是原本的磚頭和泥.
家裡沒有任何一件家具.
在牆身只放了一張.
由兩個木板椅.
一張木板鋪了一個床.
睡了一個病人.
其他全部家裡煮食的東西都在地下.
沒有多一張椅子給我坐.
我去探望這個家庭的時候.
為什麼去探望他.
因為那裡的木匠跟我說.
鄧先生,這個弟兄病了很久.
他說很想你去探望他.
他說他祈禱.
當然這些伺服的事我當然會做.
我去到的時候.
我見到的時候真的很感動.
在心裡我都問上帝.
其實他什麼都沒有.
為什麼上帝還要讓他經歷這件事呢.
接著我去到的時候就為他祈禱.
我跟上帝做了一個這樣的祈禱.
我問上帝.
他今天就好像在聖經裡面所說的那個黑蓋.
好像那個大馬蜂.
他今天什麼都沒有了.
上帝你可不可以可憐他.

$^{1481}$你可不可以幫他.
你可不可以醫治他.
我想告訴你.
我真的經歷過有一次.
聖經裡面所說彼得.
記不記得彼得的岳母耶穌.
一講的時候.
他立刻起來煮飯給他吃.
是真的出現了.
我告訴你祈禱之後.
他真的消退了.
他真的康復了.
他真的立刻起來的時候.
把他家裡僅有的幾隻雞蛋煮給我們吃.
是真的發生的.
但是我想告訴你.
我不會高舉原因.
這個是上帝的工作.
關我什麼事.
我只不過是一個渠道.
而上帝就是透過這個人.
是有原因的.
之後才震撼.
因為這個人的緣故.
在一條村.
後來他康復了.
他到處去跟人說這件事.
之後他的村長找了一個人來跟我聊天.
信了耶穌.
整條村都信了耶穌.
在短短的半年裡.
整條村全部信了耶穌.
上帝透過這件事.
只是想透過他去將福音傳給其他人.
跟我一點關係都沒有.
但如果你問我信不信有神蹟.
信.
神蹟是什麼.
不只是這個醫治的事.
而是因為我們願意用愛心.

$^{1521}$去到他的家裡為他禱告.
我們願意做最少的事.
而上帝就使用.
這個就是神蹟.
讓人看到盼望.
這個就是神蹟.
讓人從沒有愛.
沒有人理的裡面.
有人理有人愛.
這個就是神蹟.
所以求神幫助我們好嗎.
我們每天都有神蹟.
求神保守我們.
我們是重心極度.
差不多我們求主你去帶領幫助我們.
我們很願意上帝去教我們.
我們怎樣在神的面前.
立志去成為上帝你自己.
跟隨你的基督徒.
我們真是願意將自己全然的獻給你.
我們真是全然的.
我們願意將自己全部的.
交給神的面前.
我們求你真是叫我們.
能夠將整個人獻給你.
亦都去用你的價值觀.
去過我們的每一天.
亦都求你去叫我們.
去將你每一個每一天.
所給我們的恩典.
我們記下來儲下來.
到年尾到年終.
我們就去數算原來上帝.
你陪我們走過多少多少路.
所以我就求主你親自帶領我們.
幫助我們.
求你讓我們真是能夠.
拿著你這本德性秘笈.
我們能夠可以去過一個.
德性豐盛的生命.

$^{1561}$感恩讚美.
基督教主耶穌德性明智祈求.
多謝鄧全導今天提醒我們.
我們究竟信得主.
到今天我們的生命有沒有更新呢.
我們唱一首詩歌一起起來.
唱一首詩歌.
我們唱一首詩歌一起起來.
唱一首詩歌一起起來.
\newpage



\section{}
\label{sec:ZwtOgh6Qmm4}
\textbf{2023-01-12 周四晚查經班《被擄歸回歷史研讀(3)》 講員:曾裔貴牧師}
\newline
\newline
連結: \href{https://youtube.com/watch?v=ZwtOgh6Qmm4}{\texttt{https://youtube.com/watch?v=ZwtOgh6Qmm4}} ~~~~ 語音日期: 2023-01-13
\newline
\newline
\hyperref[sec:dkl9UnCmWfQ]{\small{< < < PREV SERMON < < <}}
~
\hyperref[sec:index_chronic]{\small{[返順時目]}}
~
\hyperref[sec:index_scriptual]{\small{[返順卷目]}}
~
\hyperref[sec:c4O5F1DLbjI]{\small{> > > NEXT SERMON > > >}}
\newline
\newline
$^{1}$猶太人他們的歸回歷史和當時的先知時代的渲染和預言.
我們都盼望這個其實是朝向彌賽亞的降臨.
主前的大概是四百年.
幫助我們能夠對於上主看歷史.
或者是當時的列強,巴比倫帝國,馬代,波斯帝國,希臘的興起等等.
這些都是我們在歷史裡面看到的.
當時猶太人他們身處在列強這些大國的夾縫裡面.
他們要面對的其實是一個很細小的以色列地.
但是這個是至關重要因為是關乎到神的救恩.
日後藉著猶太人,藉著律法,藉著耶穌基督的降生.
所以這一段歷史都是擺在聖經裡面.
是都算相當重要.
幫助我們能夠在梳理那些歷史的時候都能夠掌握.
雖然我們未必有很大的共鳴感.
但是盼望這個都是我們大家在一些大的方向原則.
特別是教會的發展方向都能夠找到一些重要的線索.
和那些大的真理.
求主帶領我們.
我們一起在你面前等候禱告.
奉耶穌基督的名.
阿們.
要報告一下.
因為下個禮拜四晚.
19號因為有個校長的弟兄.
他媽媽安息禮拜.
我要出席.
所以我下禮拜就停一次.
因為都接近過年.
19號應該都很忙了.
不如我們就停一次.
過了一年我們才恢復這個查經班.
我們回來看這個以色列的第三章.
看看大家是否對於這本書能夠掌握.
一至第六章其實就是記載第一次的被擄歸回.
到第七章已經是另外一個時間.
相隔都已經是二十多年.
就有另外一個叫以色拉.
其實就是這個文士.
這一位的祭司.
來到他回歸.

$^{41}$引發的就是由第七章.
一直到第十章的以色拉記.
就完結了.
其實是一個記載.
一個你可以說是復興.
就是他們要來到去一起做一些事情.
就是有些事情他們已經做錯了.
那要怎樣做呢.
不過這裡又很特別的.
我們將來有機會看到.
你就會知道.
那件事未必是最好的一個手法.
就是那個處理的問題.
未必是一個聖經所讚許的一個手法.
以色拉就是這樣做.
做了之後可能帶來一個很大的震撼.
但是聖經又沒有講到.
到底有沒有結果.
做完之後會怎樣呢.
如果你有興趣想知道.
以色拉和當時的領袖做什麼呢.
你看第十章.
今晚回家.
但是沒有結果的.
講完那件事.
有些人做.
有些人就反對這樣做.
但是最後是沒有結果的.
那就完了以色拉記.
所以讀以色拉記.
其實就相當簡單.
一到第六章.
就是他們建這個聖殿.
那個根基.
建好之後.
就開始招致那個碧巴反對.
就停工.
停工之後.
當然就是建了根基.
停了十幾年之後.

$^{81}$又再來得到那個詔書.
其實已經過了兩個皇帝.
波斯的皇帝.
去到這個應該是.
類似是甘拜西二世的時候.
可以讓他們再重新建造.
一直到第六章就完結了.
第七章到第十章.
其實是一個復興的奮興會.
第三章.
第二章我們講過.
很長很長的名字.
每一個都值得紀念.
雖然不是每一個人名.
其實有的.
因為他們是四萬九千八百多人.
其實我想每一個人都會記錄在案.
登記在冊.
不過就以那個族長.
以那個領袖.
以那個哪一個支派的人.
或者哪一個地方的人.
譬如他是台山.
譬如他是番禺.
用這些來做一個單位.
但是其實.
當時的那些民事.
應該記錄每一個人.
他們的人名.
這樣來追溯他們那個支派.
其中就凸顯了.
有一些他自稱是祭司.
是利未人.
但是找不到那個族譜.
找不到他們那個.
可能是那些太公.
找不到.
所以要等到.
由祭司開始公職.
他們就用污名土名.

$^{121}$來解決那些疑難.
不知道是抽籤.
還是有些事情來做.
已經失傳了.
所以後來就決定了.
原來這些都是可以的.
類似在政務處宣誓.
在政務官面前宣誓.
你遺失了什麼.
現在你要來拿回.
就給你.
類似這樣.
第三章我們讀過一段.
就是頭一段.
就是他們開始都挺同心.
來做這個聖殿.
就是七月一路準備.
七月就有很多節期.
催國節.
贖罪日.
還有鑄平節.
七月裡面有三個節期.
其中要上去聖殿.
即將來.
其實不是將來.
猶太人由摩西開始.
一年有三個重大的節期.
他們一定要上耶路撒冷去守的.
南丁.
逾越節.
五旬節.
還有鑄平節.
三個.
一定要守的.
所以七月他們現在就是被擄歸回.
他們一路去準備.
無論是他們住的地方.
還有預備那些材料.
然後因為他們要起神的殿.
所以需要那些很長.

$^{161}$很穩固.
不會住的香柏木.
在這個藥柏.
在敘利亞高原的地方.
類似是戈蘭高地.
伐木之後.
就除海.
敘利亞的海一直流下來.
流到下面藥柏這個地方.
就可以拿到這些木.
然後起聖殿.
我們由第八節開始看.
百姓到了耶路撒冷神殿的地方.
第二年二月.
過了五個多月.
去到第二年.
一定是以猶太人的曆法.
不是以波斯的曆法.
因為是不同的.
所以正月就是他們的逾月節.
所以他們守完逾月節.
之後就來到第二年的第二月.
開始建造這個聖殿.
你可以說是他們守完那些節期.
預備好心靈.
預備好材料.
來準備.
即是說經過了很多節期.
他們都完了.
記得在這個時候.
以色列的時代.
還沒有那個縱光節.
縱光節或者叫修殿節.
就是去到主前的第二世紀.
就是因為他們繼續.
即是他們一直住了之後.
其實已經過了尼希米的時代.
尼希米是四百多年的主前.
所以連尼希米那個時代都過去之後.
就去到差不多主前的第二世紀.

$^{201}$當時的南方王.
敘利亞王.
去攻打耶路撒冷.
令到耶路撒冷受到圍困.
還有敘利亞王這個埃提瓦古四世.
非常褻瀆神.
非常憎恨猶太人的神.
所以他就想到一個辦法.
去到聖殿親自獻祭.
不過獻豬不是獻牛羊.
即是徹底羞辱你們的神耶和華.
所以就激發了一個很大的.
一個民間的一個部.
即是一個家族叫馬提亞家族.
應該叫馬加比家族.
他的兒子叫馬提亞.
來去揭竿起義.
因為差不多沒有辦法.
你要來去摧毀聖殿.
所以猶太人就更加.
來去有這個誓死來對抗.
所以終於打敗了這個.
埃提瓦古四世.
他們就去重光那個聖殿.
據聞當時聖殿的燈光.
是一直都不熄.
一直都點著.
不知怎樣自己不斷的燒著.
是不熄聖殿的光.
所以當他們一直去潔淨聖殿之後.
就來到去在聖殿裡面.
無論是修葺,點燈.
後來就叫做重光節.
每年都守的,在冬天的.
耶穌都有守那個重光節.
但是以色列的時候.
他們根本都沒有聖殿.
他們現在是回來重建聖殿.
過了大概200多年.
當然你說由他那個536開始回去重建.

$^{241}$去到主前的第二世紀.
好像165年還是166年.
差不多了,是那個時間.
就是耶穌還沒出世之前的主前165年.
就是那個事件.
所以重光節後來是加多的節日.
每年猶太人都守的.
他們都守這個叫做Halucka.
是很開心的一個點燈的節日.
第八節就是第二年的二月.
撒拉提的兒子所羅巴伯.
約撒達的兒子耶穌亞和其餘的弟兄.
就是祭司利未人.
並一切被擄歸回耶路撒冷的人.
都興功建造.
又派利未人從二十歲以外的.
二十歲以外就是他們可以公職.
其實嚴格來說是35歲.
才可以來到去在聖殿裡面服事的.
30歲可以開始來到去服事.
35歲正式的可以來到去做.
所以原來在利未人的年代.
利未人他們都有五年的實習.
實習神學生讀五年.
所以他學怎樣事奉.
就跟那些祭司.
跟那些老的祭司或者利未人.
所以都幾特別的.
這個利未人有講的.
我不記得第幾章了.
就是規定了的.
在這個時候.
凡是二十歲以外的利未人.
因為少.
所以要將年齡降低.
如果不降低.
其實未必夠資格.
所以都感恩.
有些二十歲以外的.
都跟著上去重建.

$^{281}$跟著自己的父母.
或者自己的族兄弟.
上去重建.
是利未人才可以在聖殿裡面公職.
所以利未人.
現在因為未建好聖殿.
所以他們幫手做那個.
可能是領袖.
公頭.
督你建造耶和華殿的工作.
第九節.
於是猶大的後裔.
就是耶穌亞和他的子孫與弟兄.
革滅和他的子孫.
這個第二章讀過了.
利未人希拿達的子孫與弟兄.
都一同起來.
祝你那在神殿作公的人.
即是說這些公頭.
來做一個領導.
猶大之派似乎.
可以出到一些比較強的領袖.
可能回去的人多.
有心的人又有能力的人.
可以來做這些領導的工作.
所以這樣分工都很重要.
第十節.
將人.
即是那些石匠或者鐵匠.
納耶和華殿根基的時候.
祭司街村禮服.
那些禮服的器皿.
聖殿裡面的器皿.
是當時古列王頒佈了.
要全數給猶太人.
他們被擄歸回.
所以這些制服.
這些禮服.
當然就是利未記摩西.
其實應該是出埃及記的時候.

$^{321}$摩西領受了神的寓言.
那個密使.
要他去做.
用什麼材料去做.
用什麼針線.
用什麼質料.
是關乎大祭司的.
有講究的.
進去會幕裡面.
和在會幕外面穿的又有不同.
譬如你見到警務處長.
在不同的場合.
他穿的制服都有不同.
所以祭司都有他不同的.
穿的制服.
有時他要去倒垃圾的灰.
處理了那些憲制完的灰.
他都要換衣服.
倒完回來又要洗身.
所以其實都很講究.
這些其實是利未記很仔細的教導.
留意第十節.
穿禮服吹號.
阿薩的子孫利未人敲拔.
阿薩的子孫就是當時.
由大衛王所頒佈的.
有阿薩歌唱的阿薩.
又有那個.
我又忘記了名字.
我想起了.
阿薩的子孫利未人敲拔.
有不同的樂器.
照以色列王大衛所定的例.
這個在歷代志上又頒佈了.
這個耶道頓.
希曼.
歌唱的希曼的子孫.
都站著去讚美耶和華.
原來那些師班都有細集.
某程度上都是.

$^{361}$因為他們一直去訓練.
越來越在各方面去精於訓練.
所以一個靚的崇拜.
所謂靚就是能夠各方面都和諧.
就有不同的部.
譬如我們有影音部.
有聲樂參與.
有崇拜的師事.
來管理一些秩序.
有童工.
有其他的.
譬如主席.
讓整個崇拜.
在不同的部份一起去參與.
這樣來突顯和諧.
和領導會眾一起去敬拜神.
這個都是很重要的配搭.
任何一個崗位都有它的重要性.
所以現在其實.
他們雖然被擄了70年.
原來這些崗位.
其實一直都沒有失傳.
一直都有準備.
即是說如果亞薩的子孫.
他們在巴比倫.
在馬來.
在波斯.
他們都一直沉潛.
可能白天就是做他們自己的生意工作.
去到星期六.
因為他們是安死人.
猶太人.
安死人.
我不知道當時被擄的時間.
准不准許他們有自由敬拜.
准不准許他們有自己的會堂.
因為當時是未衍生會堂.
是後來主前的400年左右.
舊約和新約中間的400年.
就是由會堂的文化開始.

$^{401}$因為他們已經分散去到不同地方.
尤其是土耳其一帶.
歐洲的地方.
所以他們就要將它化整為零.
所以我們基督徒都要學.
遇到壁壩都不怕.
只要我們大概大家有5至10個.
猶太人就15個左右.
他們就會有自己的會堂.
他們就會有.
重要的就是傳承摩西律法.
他們來到去自己中毒.
其實令到下一代知道自己是有神.
和讀利未記.
神期望猶太人是怎樣.
所以他們是會有這些傳承.
如果你有興趣.
安息日你可以去半山.
中半山上的猶太廟.
我不知道安息日是否讓你進去.
你可以去看人家怎樣敬拜.
他們真的會有些弟兄.
其實好像中間是一個打擂台的擂台.
在正中間.
其他特別低.
椅子都排了.
中間就是人們上去大唱詩歌或者讀經.
就是猶太人安息日的聚會.
婦女坐超等最高的.
圍著兩邊.
應該是三邊.
所以婦女在上面.
居高臨下看到.
原因是.
麻煩你們有這麼遠走這麼遠.
正中間是弟兄.
是猶太人.
所以婦女要戴帽子.
蒙頭.
其實居高臨下更清楚.

$^{441}$所以無所謂誰歧視誰.
真是的.
男性才可以坐下面.
正中間圍著.
男孩子都可以讀經.
這就是他們猶太人的做法.
所以現在避老歸回.
他們一開始就已經有這些師班長.
已經有這些人會唱.
而且他們唱的詩歌.
似乎是一三六篇的詩篇.
耶和華本為善.
他的慈愛永遠常存.
就是說.
一路以來.
70年沒有令到他們滅俗.
70年沒有令到他們的信仰完全瓦解.
被同化.
這個是很重要的.
你想想現在我們很多.
香港的基督徒.
傳道人都移民去英國.
或者是加拿大.
溫哥華.
台灣.
其實沒有其他.
其他都沒有人去.
歐洲講的德文.
荷蘭文就沒有人去.
甚至乎沒有.
你可以想像到.
變了很聚在一起.
還要不是聚那麼簡單.
而是繼續去傳承那個舊恩.
和要令到下一代更加深化.
不要受那個異族.
異地的文化的影響.
譬如你去到外國.
人家現在崇尚的食大麻.
同性戀.

$^{481}$那是不是你們那些下一代都明白.
而抗衡呢.
接受呢.
抗拒呢.
這些是很重要的.
那個傳承.
猶太人就做得很好了.
他們一上來就馬上可以來去敬拜.
所以他們不是唱那些外邦的歌.
他們一直在.
照以色列皇大位所定的例.
都站著讚美耶和華.
接著第十一節.
他們彼此唱和.
讚美.
稱謝耶和華說.
他本為善.
其實是在唱.
唱這些詩歌.
揭一揭一三六篇好嗎.
給大家知道一下.
這些猶太人他們在.
傳承時期的敬拜.
對造物主的敬拜.
一直都是這個.
很長的.
你們要稱謝耶和華.
因他本為善.
他的慈愛永遠長存.
神是本為善.
這些只是一生.
我們去想一下.
其實在你的人生.
造物主是不是真的在你的生命裡.
他本為善呢.
不是要做任何事.
他本身就是善.
是善良.
比我們想的善良.
還要超出很多很多倍.

$^{521}$他的慈愛是永遠長存的.
其實不容易.
假設你的兒子突然遇到意外.
斷了一隻腳.
已經是很輕的.
不是變植物人.
只是斷了一隻腳可以搏.
你想想.
你都不可能會當下.
你能夠馬上就說.
你在十字車上.
耶和華本為善.
他的慈愛永遠長存.
很難.
可能那個救護員都以為你有事.
要找心理輔導.
因為是真的.
你的兒子當時就是這樣.
叫生叫死的時候.
你居然在這裡這樣唱.
其實是很難的.
所以這些是要很多的經歷.
整個民族這樣唱.
是不容易的.
還有.
你們要稱謝萬神之神.
因他的慈愛永遠長存.
主題曲的主調就是.
慈愛永存.
你們要稱謝萬主之主.
因他的慈愛永遠長存.
稱謝那獨行大其事的.
因他的慈愛永遠長存.
稱謝那用智慧做天的.
因他的慈愛永遠長存.
稱謝那鋪地在水以上的.
因他的慈愛永遠長存.
開始有些內容了.
開始說到創世記的創造主.
他的榮耀.

$^{561}$他的權柄.
他的獨行其事.
要有光就有光.
是那一位.
由無創造萬有的主宰.
猶太人信這一位神.
信這樣的一位神.
所以我們基督徒.
由主耶穌的死而復活.
我們偶然間面對突發性.
我們一下子受不了.
未能消化.
不要緊的.
當下人人都會有情緒.
人人都會有悲哀.
但是沉澱之後.
這些詩篇就要浮出來.
否則你的信仰會崩潰.
很多事情解釋不了的人世間.
你無法解釋的.
但是你要記住.
造物主在你的生命裡面.
他不單有主權.
更加有慈愛.
這是我們需要去深化.
這些才是重要.
不是花香常蔓.
天色常藍.
你才說他的愛是長存.
不可以輸打贏要.
基督徒我們真的要有不同的學習.
後面我不詳細讀了.
因為這篇詩篇很長.
猶太人一回來.
第一個建造.
在第二年二月.
他們就開始同心.
其實是四萬九千多人.
來去重建神的殿.
準備要立根基的時候.

$^{601}$他們就開始來去唱歌.
似乎敬拜多過真的做事.
敬拜比要怎樣做更重要.
他們讚美耶和華的時候.
眾民就大聲呼喊.
因耶和華殿的根基已經立定了.
然而有許多祭司利未人族長.
就是見過舊殿的老年人.
這裡有些考究.
和大家揭曉列王記一下.
揭上少少.
列王記一下的廿四章.
應該是很厚的.
因為列王記一下只有25章.
24章.
列王記一下.
約瓦敬年間.
巴比倫王尼布格尼撒上到猶大.
約瓦敬服侍他三年.
猶太人猶大國最後的四位王.
約西亞是最後一位好王.
約西亞的兒子.
約哈斯,約瓦敬,約瓦根.
接著就到西底加.
西底加其實就是約瓦根的叔叔.
即是約哈斯的弟弟.
因為已經沒有兒子.
所以就去約西亞的第二個兒子.
做最後的王.
很好記的.
因為約哈斯,約瓦敬,約瓦根和西底加四位王.
約哈斯是三個月零十一日.
接著就到約瓦敬.
做十一年.
接著就到約瓦根.
也是三個月.
三個月做王.
即是說被巴比倫王廢了.
老了是巴比倫.
去到最後的西底加.

$^{641}$就是十一年了.
所以是三個月十一年.
三個月十一年.
就是最後的四位王了.
所以這裡告訴我們.
約瓦敬年間.
巴比倫王尼布格尼薩上到猶大.
約瓦敬服侍他三年.
然後背叛他.
拆開說.
大家現在已經過了那些記憶最好的日子.
記不到了.
近的記憶.
這裡說完那裡就忘記了.
我試過就是在辦公室走去辦公室.
走到忘記了拿什麼問什麼.
完全忘記了.
那你就可以想像到.
經過禮堂失憶了.
因為這些近的記憶.
大家都應該會開始感受到.
所以剛才我背的那些.
全部是初信的時候.
因為讀聖經多和有興趣去讀.
去背.
刻意去背.
所以是記得.
變成一個比較深層的記憶.
大家其實真的要把握機會.
讀多一點.
坐牢的時候就可以背一下聖經.
不會那麼悶.
那時候什麼都沒有了.
收了你的手機.
什麼都沒有了.
那時候你就可以和神彼此來相交.
背一下詩篇23篇.
給玉卒聽聽.
約瓦敬服侍他三年.
然後背叛他.

$^{681}$耶和華使到迦勒底軍.
亞蘭軍 摩押軍和阿們人的軍.
來攻擊約瓦敬.
你可想而知.
是不能回頭的.
因為神興起列國來攻擊猶大.
其實猶大是一個很小的地方.
不過是一個歐洲 亞洲 非洲必經的地方.
所以是一個很重要的軍事據點.
現在就要毀滅猶大.
正如耶和華直他僕人眾先知所說的第三節.
「這禍臨到猶大人.
誠然是耶和華所命的.
要將他們從自己面前趕出.
是因嗎哪西所犯的一切罪」.
嗎哪西其實就是以賽亞的兒子.
對不起 是希西加王的兒子.
很慘.
希西加居然有個這樣的兒子.
很難想像.
一個這麼勁權.
除了大衛 所羅門.
就到希西加了.
現在希西加的兒子是嗎哪西.
而且是猶大國最長做王的.
多少年呢.
55年.
最長是他的.
你可想而知.
那些歷史告訴我們.
嗎哪西將他的父親.
很敬重的那個墓者.
就是以賽亞.
擺在一棵樹上鋸死.
擺在那些樹洞上.
一棵樹劃開一個心.
放下去鋸死以賽亞.
可以這樣.
主任牧師生氣成這樣.
將他鋸死.

$^{721}$這是猶太人的歷史告訴我們.
這裡說第四節.
又因他流無辜人的血.
充滿了耶路撒冷.
耶和華決不肯赦免.
約瓦敬其餘的事.
凡他所行的都寫在猶大列王記上.
約瓦敬與他列祖同睡.
他兒子約瓦敬接續他作王.
埃及王不再從他國中出來.
因為巴比倫王將埃及王所管之地.
從埃及小河直到伯拉河.
都奪去了.
留意第八節.
約瓦敬登基的時候年十八歲.
在耶路撒冷作王三個月.
所以很容易的.
三個月十一年.
約瓦敬就是第三位.
最後的尾二的王.
所以他作王是三個月.
他母親名叫尼胡斯他.
是耶路撒冷人.
以利拿丹的女兒.
約瓦敬行耶和華眼中看為惡的事.
效法他父親一切所行.
就是約瓦敬.
那時巴比倫王尼布格尼薩的軍兵.
上到耶路撒冷圍困城.
當他軍兵圍困的時候.
巴比倫王尼布格尼薩就親自來了.
又帶王日雅根和他母親神僕首領太監.
一同出城投降巴比倫王.
巴比倫王便拿住他.
那時是巴比倫王第八年.
巴比倫王將耶和華殿和皇宮裡的寶物都拿去了.
將以色列王所羅門所造.
耶和華殿內的金器都毀壞了.
正如耶和華所說的.
又將耶路撒冷的眾民和眾首領.

$^{761}$就是這裡了.
並所有大能的勇士.
共一萬人.
連一切木匠鐵匠都老了去.
除了國中極貧窮的人以外.
沒有剩下的.
並將約雅根和皇母.
即是他母親太后.
後妃太監.
與國中的大官.
都從耶路撒冷老到巴比倫去了.
又將一切勇士七千人.
和木匠鐵匠一千人.
都是能上陣的勇士.
全老到巴比倫去了.
巴比倫王立約雅根的叔叔馬探雅.
代替他作王.
就是西底家王.
吸馬探雅改名西底家.
即是約西雅的另一個兒子.
或者應該是孫子.
你可以想像到.
這個.
不對,是兒子才對.
你可以說是第二次的被擄.
因為.
但以理那些是第一次的被擄.
現在這個是第二次的被擄.
就是約雅根王的時候被擄.
所以.
這樣計算.
現在去到.
這個被擄歸回的這班人.
在這裡來去去.
很有感觸.
就是可能這些.
很年輕的時候就被擄了去的.
現在過了五十多年之後.
他們長大了.
那就.

$^{801}$真的會有這樣的感覺.
這個就是我們剛才看的第三章那段經文.
就是三章十二節.
就有很多這些曾經見過舊殿的老年人.
即是現在是老年人.
舊的聖殿是很輝煌的.
是所羅門王所起的.
所以他們.
真的很有感觸.
現在能夠親自踏足自己的故土.
很小的時候就被擄了去.
現在長大年老回來.
居然現在是親手.
可以來到一起.
去重建神的殿.
所以這個是.
你可以說那個哭聲.
歡笑聲.
呼喊的聲音.
就是八感交集的聲音.
震動整個耶路撒冷.
記得我猶太人他們真的有一個這樣的信念.
耶路撒冷的聖殿的位置.
即是所羅門的位置.
是最強的5G WIFI.
直上天庭.
所以他們一定要在那個地方建聖殿.
就是現在第三個清真寺的地方.
伊斯蘭的清真寺.
現在在那個地方.
即是聖殿山.
那個地方下面就是所羅門的位置.
所以他們是.
祈禱都要上聖殿山祈禱.
你可以想像得到.
所以他們當時回來.
在這裡立那個殿的根基的時候.
一定是在那個方位.
他們真的有這個信念.
認為可以直接和天上.

$^{841}$其實新約來說是對的.
因為新約希伯來書.
真的是這樣說的.
地上的這一切.
其實是天上的那個真體的一個模型.
天上的是真.
地上的是一個模型.
所以他們猶太人有這樣的信念.
又不可以說是錯的.
當然去到主耶穌告訴我們.
重要的是你的心靈和真理裡面.
去敬拜就可以了.
不是說教會就不需要.
那我不回教會了.
每個星期看著電視.
我那個是高清來的.
一百吋的.
曾牧師的奏文我都看到的.
流汗我都看到的.
不是這樣.
不關這樣的事.
而是要敬拜神.
沒錯.
任何地方任何時間都可以敬拜神.
但是一個公共的聚集.
是神規定的.
是我們要和弟兄姊妹一同敬拜神.
朝見神.
但也有弟兄姊妹之間的相交.
所以嚴格來說.
疫情放緩.
其實就要取消直播.
真的要有這個公共性的學習.
不關乎現在的網速.
視像的穩定.
Zoom的或者YouTube的穩定.
不是這個問題.
而是崇拜本身.
所以現在查經班.
我覺得這樣是好.

$^{881}$大家可以兩個.
你可以回來實體.
你可以在自己家線上.
當然如果遲些我們那個學習.
是有大家的相交或者分享的.
那就另計.
實體就會相對比較容易.
如果不是.
其實兩者是不會有太大抵觸.
但崇拜就不同了.
崇拜就不只是看著這麼簡單.
而是你要整個人的參與.
才是那個需要的.
除非你根本不能夠.
不能夠就不是停止聚會.
是不能夠聚會.
例如你要輪班.
你就不能夠聚會.
例如你病了.
你要在家休息.
你要隔離.
你這個不能夠聚會.
例如你在醫院.
你就不能夠聚會.
你就可以看直播.
但如果不是.
其實應該就是每個禮拜.
我們都要來到朝見神.
好百感交集的哭聲和笑聲.
也有很多人大聲歡呼.
第十三節.
甚至百姓不能夠分辨.
歡呼的聲音和哭號的聲音.
因為眾人大聲呼喊.
聲音聽到遠處.
一個很震撼的畫面.
一個相當大的呼喊.
是他們開始立神的殿.
和唱一三六篇的詩篇.
這樣來贊頌神.

$^{921}$這一次是一個很大的勝利.
有五萬人.
四萬九千多人回來重建.
有秩序.
有自己的崗位.
有領袖索羅巴伯和大祭司約書亞.
來領導他們.
即是說.
關乎宗教的約書亞就會領導.
關乎到那些事工的.
雜務的.
索羅巴伯就會做領導.
這樣來大家一起分工.
首先現在做了的就是.
那個聖殿的根基建好了.
打好了那個根基.
然後在上面建造.
相對很容易.
因為最難就是要打那個樁.
因為那個聖殿是相當巨大的.
需要的就不是那麼簡單.
如果沒有問題.
第四章我們來看看.
聽到遠處是一至三章的結果.
就是他們起義被波斯帝國的古列皇帝.
頒佈的天下的詔書.
各國的人民可以回到你的原居地.
敬拜你們自己的神.
除了不可以叛亂之外.
你們可以有你們自己的敬拜.
即是說.
妥協一種懷柔政策.
不再是高壓強迫.
例如拆了你們的聖殿.
擄走你的精英.
過去同化.
這樣你的國家自己會亡的.
你會慢慢就什麼都沒有了.
文化的侵略.
過去之後.

$^{961}$就是要令到他們吃得好穿得好.
盡快讓他們能夠對巴比倫產生那種愛慕.
自然就會慢慢忘記自己的文化.
這是文化上的滅族.
一些強權的國家就會這樣.
但是波斯真的很厲害.
很聰明.
吸取到巴比倫和馬代的失敗.
很快就滅亡.
所以一來就不用這麼多的行政費用.
因為這麼多國家很大開支.
軍餉.
軍人的福利.
所以就給你們有自治.
給你們有你們自己的敬拜.
我只需要有一些當地的官員.
看著你們不要叛亂.
就已經搞定了.
所以現在猶太人.
剛剛在這個政策底下.
就受惠.
不過聖經一開始告訴我們.
是神激動鼓勵的心.
是神感動一些願意回到上去的人.
只有四萬九千多人.
其實不是多的.
在波斯帝國已經散居的.
真相都有很多人.
譬如我們很快揭一揭.
以斯帖記.
將來我們會看的以斯帖記.
就好像想告訴我們.
做物主不會只是在耶路撒冷.
來感動猶太人.
在外邦.
在波斯帝國的書山城.
就是他們的首都.
都可以感動.
在皇宮裡面的皇后.
原來是一個敬虔的猶太人.

$^{1001}$可以感動她挽救自己的國民.
不過浮現了一個很坎坷的事實.
原來猶太人相當多.
但是不回到上去重見.
那你怎麼解釋呢.
第九章.
以斯帖記第九章.
沒抵改記錄者事.
第二十節.
寫信予阿俠徐老王.
各省遠近的所有的猶太人.
祝福他們每年守阿達月十四十五兩日.
就是這個叫做普爾日.
以這月的兩日為猶大人脫離仇敵.
得平安.
轉憂為喜.
轉悲為樂的吉日.
在這兩日切然歡樂.
彼此攜送禮物.
周濟窮人.
就是這個普爾節.
那就變成猶太人又多一個節日.
又是他們創造的.
因為本來是要被滅族的.
但是現在因為以斯帖.
所以很多人叫做Asta.
美容院就叫Asta.
就是因為這個皇后的原因.
出名.
所以普爾日都是因為這個Asta而來的.
所以猶太人他們又有一個新的節日.
他們要守的.
那些小朋友就要玩一些遊戲.
有一個叫做Quacker.
就是一邊轉.
就好像來打哈曼.
敵人哈曼.
所以他們會送一些糖等等.
來去饋贈的.
二十三節.

$^{1041}$於是猶大人按著末底改所寫的.
與他們的信.
應承照初次所守的守為永禮.
是因猶大人的仇敵阿甲族哈米大他的兒子哈曼.
設謀殺害猶大人.
祭普爾.
就是祭遷.
普爾是希伯來文.
祭遷為要殺盡滅絕他們.
這時報告於王.
王面降旨.
使得哈曼謀害猶大人的惡事.
歸到他自己的頭上.
並吩咐把他和他的眾子都掛在木架上.
照著普爾的名字.
猶大人就稱這兩日為普爾日.
他們因這順常的話.
猶因所看見所遇見的事.
就應承自己與後裔.
並歸附他們的人.
即入猶太教的人.
每年按時必守這兩日.
是兩日的.
永遠不廢.
各省各城家家戶戶.
世世代代紀念遵守這兩日.
使得這普爾日.
在猶大人中不可廢掉.
在他們後裔中也不可忘記.
其實我們也可以自己家族裡面.
設一些獨特的紀念日.
例如你是何時洗禮.
你就告訴你的兒子和孫子.
今天是我洗禮.
我們一起感恩.
或者你送禮物給孫子.
給他們有一個印象.
我們的父輩對神有敬重和愛慕.
其實日曰言表是挺好的.
例如我們中國人.

$^{1081}$每個人都趕著過年.
又要買東西.
要開始買蓮花.
開始要貼那些揮春等等.
慢慢慢慢成為了大家的風俗.
就有那個團聚的意思.
其實都很重要.
但聖經原來都有這些團聚.
本來是要滅俗.
什麼都沒有了.
那個旨意本來已經下了.
但是可以推翻.
如果不是出於神.
怎麼可能是一個哈該.
對不起.
這個沒底改.
就可以來改變呢.
一個皇后就改變呢.
不可能的.
所以一定是有神在其中介入的.
所以他要定兩日.
還要立誓.
後裔一定要記得這兩日.
是他們轉危為安.
轉悲為.
轉憂為喜的日子.
來自送禮物.
所以我們自己都可以.
創造一些心思.
最好是由小朋友開始.
讓他們知道.
奶奶今天洗禮.
奶奶送禮物給你.
你記得了.
記住.
將來你十六歲就要洗禮了.
要敬拜神.
和奶奶一起守聖餐.
其實你這樣自小去教育.
他自小就知道神的愛和存在.

$^{1121}$所以他們有這個普爾日.
第29節.
阿比孩的女兒皇后以斯帖.
和猶大人沒底改.
以全權寫第二封信.
堅囑猶大人守者普爾日.
用和平誠實話.
寫信給阿黑徐老王.
國中127省所有的猶太人.
所以這句說話.
我想讀給大家聽.
其實很多猶太人.
但是很少人回去.
只有四萬九千多人.
所以你可以想而知.
絕對不是很多很多.
不要緊.
但是神的工作.
也不只是限於耶路撒冷.
在以斯帖的時候.
在波斯帝國.
阿黑徐老王的時候.
都有這樣的事.
我們回到以斯拉記的第三章尾.
眾人的呼聲達到遠處.
由他們很遠的地方.
大家知不知道要由巴比倫.
或者波斯帝國.
他們是Travel by foot 行路.
走上耶路撒冷要多久.
其實第七章就有說.
以斯拉用了五個月.
由波斯帝國走上耶路撒冷.
可能有不同的原因.
譬如有土匪.
或者有其他的事情.
但是猶太人他們被老歸回.
都差不多用了七個月的時間.
所以相當長的時間.
我們還有些時間可以繼續讀下去.

$^{1161}$其實這個時候.
開始出現第四章.
就是你必然會面對反對.
因為你浩浩蕩蕩被老歸回回來.
沒錯,是有古列王的文書.
有可能古列死了.
新的王開始興起.
就會帶來如果當地的原居民.
上個星期我們說過.
亞述帝國的時候.
將那些外族人遷徙來到耶路撒冷.
又將耶路撒冷北面的以色列國的人.
擄去遷徙到外地.
將他們放逐.
將他們來充軍.
所以在北面的以色列國.
就不是純種的猶太人.
就是說已經不是和猶大國那些兄弟了.
所以猶大國就有這條刺.
一直都看不起北面的以色列.
所以這裡一直下去.
有些解經家就認為.
到了耶穌時代.
北面的撒瑪利亞.
為什麼撒瑪利亞人.
是被猶太人歧視.
和不會和他們來往.
甚至乎都要繞過不去撒瑪利亞.
來到上去加利利.
原因是因為撒瑪利亞人.
正正就是那些亞述遷徙去的那些外族人.
和北面的以色列人通婚.
所以來到成為了後來的撒瑪利亞人.
所以在耶穌的時代.
撒瑪利亞是一個被歧視的地方.
他們是不會去的.
他們很稀奇.
撒瑪利亞人可以信主.
那個婦人或者那個十個大麻蜂.
未有十個大麻蜂.

$^{1201}$接著耶穌叫他們你們去就會醫好了.
他們一離開就全部好了.
只有一個回來.
路加醫生告訴我們.
那個回來的十個大麻蜂.
九個走了.
只有一個就是撒瑪利亞人.
有麻蜂的時候.
十個人都不理你是什麼民族.
大家都是在患難裡面.
大家都不介意.
我們一起好了之後.
可能那九個猶太人.
突然發現那個和我們一起出入的.
大家都是生滿了麻蜂的爛掉的器官.
突然那個大家都好了.
但那個是撒瑪利亞人.
馬上就歧視他.
這個都反映了我們那種人性.
反映了我們那種學習得不好的時候.
那種兩面.
返完崇拜去到會所坐下喝茶.
職員不經意倒了水滾到你.
馬上罵到他狗血淋頭.
不是想的.
不過突然間就爆火了.
就會開始就悶了.
這些就是有原因.
但十個大麻蜂是沒有原因的.
大家都是很慘很慘等死.
但當神跡醫好了之後.
突然間開始立即.
是歧視那個人.
那個撒瑪利亞人.
你可以想像得到人性的問題.
所以來到第四章就出現這個逼迫.
我們留意這裡很重要.
猶大和汴雅敏的敵人聽說.
被老歸回的人為耶和華以色列的神建造殿宇.
就去見所羅巴伯和以色列的族長.

$^{1241}$對呀這個所羅巴伯.
其實他是猶大的省長.
等下我也有時間.
其實要和哈該書一起讀.
因為等下就要涉及到神的寓言.
默示給哈該.
先知哈該.
所以我們要記得.
在這個以色列記沒有提到先知.
一起回去.
就不代表沒有.
因為等下就會出現哈該和撒迦利亞兩位先知.
第二節他們就去見所羅巴伯.
對他們說.
請容我們與你們一同建造.
因為我們尋求你們的神.
與你們一樣.
自從亞述王以撒哈頓帶我們上這地以來.
沒錯就是這班人.
我們常祭自神.
你怎樣理解呢.
他們這麼有誠意.
我們是原居民住了70年.
你們走了我們在這裡住了.
我們是被亞述王.
即是前幾個朝代.
遷徙到來.
我們已經落地生根在你們猶太的國土.
我們現在開始已經學習敬拜你們的神.
當然那只是其中一個保佑他們的神.
一定是不會放棄他們自己民族的神.
一定是不會單單歸向耶和華.
因為根本就沒有聖殿.
你們何來去敬拜呢.
何來去祭祀呢.
何來你們會懂得去祭祀呢.
你們都沒有祭司沒有利尾人.
怎麼會教導你們那些律法.
那些聖經的真理.
神要怎樣的敬拜呢.

$^{1281}$是不是.
那你可以在家自己做個像.
或者買個觀音.
當這個就是耶和華都可以的.
如果你的真理是錯的.
那你每天拜著觀音就當他是耶和華.
我拜耶和華的.
不可以的.
因為你要根據聖經這樣說的敬拜.
所以那個是要在乎靈.
和在乎真理的敬拜.
是耶穌這樣說的.
所以這裡告訴我們.
這樣的情況.
其實只是一種的.
很想帶來一種的.
大家相愛不好口.
你給我們參與.
我們有份於你們的建造.
有什麼後果呢.
如果所羅巴伯站得不夠硬呢.
他們.
因為他們已經住在這裡.
差不多大半年了.
對當地的這些外族人.
其實都開始療癒指掌了.
還有更大的就是.
已經吸取了所羅門和歷代的王.
他們那個慘痛的教訓.
就是和外族人來往.
慢慢放下那個設防.
來跟隨他們拜當地的偶像.
拜那些摩洛.
拜那些阿斯塔祿.
拜巴力.
阿瑟拉.
慢慢慢慢.
自己的神.
只是最後的防線.
最後的保障.

$^{1321}$但實際上.
拜的是其他的.
那就變成.
會帶來一個.
又是再.
好像慢慢就會回去重蹈覆轍了.
所以所羅巴伯在這個時候.
就站得很硬.
因為如果他們人多.
他們來到.
譬如一個建造的委員會裡面.
如果他們的票數多.
聖殿不應該這樣做的.
沒理由這樣做的.
沒理由致聖所這樣做的.
沒理由那個聖所.
是要這樣做的.
這樣擺設的.
如果他們不同意的時候.
那你就沒得做了.
如果你說要開一個委員會.
嘩,很容易的.
將自己的親信.
全部放進去那個委員會裡面.
接著我們要講民主的.
我們要信任議會的制度.
議會的制度就是.
我們就是大家彼此尊重.
現在十個當地的.
這些原居民.
來投票.
我們認為聖殿不應該這樣做的.
聖殿的門不應該向東面開的.
一定要向北面的.
可能是這樣的.
這樣的時候.
你引進這種的.
好像同心.
實際上不是真正的.
聖經說的同心.

$^{1361}$是會減低了.
那個你們一心一意去建造.
那個神的殿.
是根據神的那個啟示.
是根據摩西的律法.
的安排來建造.
所以在這個時候.
第三節.
所羅巴伯.
耶穌亞.
和其餘以色列的族長對他們說.
我們建造神的殿.
與你們無關.
我們自己.
為耶和華.
以色列的神協力建造.
是照波斯王古列所吩咐的.
所以站得相當硬.
就是.
你們祈禱就行了.
我們足夠人手.
足夠的力量.
足夠的材料.
而且是有足夠的法理基礎.
是波斯王古列.
吩咐我們回來去做的.
所以第四節.
那地的民就在猶大人建造的時候.
使他們的手法軟.
擾亂他們.
從波斯王古列年間.
直到波斯王大利烏登基的時候.
賄賣謀士.
要敗壞他們的謀算.
開始用一些計謀.
用一些陰險的手段.
用一些可能是說一些壞話.
說一些是非.
搬弄一些人.
令到大家猶太人不同心.

$^{1401}$又或者.
來訴諸於君王.
因為一直在過程中.
由古列到大利烏.
已經過了兩三位皇帝.
有可能古列一離開世界.
那個政策未必能夠持續.
就出現了這個.
小小的隔逢真兇.
就會令到這班人.
一心一意要推翻他們.
因為其實他們怕的不是宗教這麼簡單.
有宗教就自然會有團結.
有團結自然會有政治.
就會強大.
猶太人又會復國.
又會有他們自己的勢力.
就會影響當地的原居民.
即是在外地來定居的這些人.
所以他們不只是阻撓他們宗教上的建造.
是擔心他們將來一有宗教.
自然就會帶來力量的對衡.
那種對抗.
所以真的.
宗教是對人很大的維繫.
所以教會如果能夠靈明好.
頂尖會很團結.
很同心.
很合一.
一起去做一個大方向的事情.
就是宗教的核心強才會得.
所以他們現在第二方面就是.
索羅巴伯又意識到.
他們現在就說.
本來是和你講.
中國人講那個叫什麼.
突然間忘記了那個成語.
先禮後兵你都不領情.
變了現在就開始.
本來說一起做.

$^{1441}$我們是又出錢又出力這樣做.
和你們一起協力.
我們都是敬拜你們的神.
很動聽的記住.
所以不要這麼簡單.
越是重要的崗位.
越是要審查那個人的靈明.
是越要嚴格.
初入教的不可作監督.
恐怕他自高自大.
提摩泰前書第五章.
阿保羅教提摩泰.
其實這些是智利名言.
初順很熱心.
拼命就熱心.
還不找他做執事.
你還不找他做堂董.
熱心.
當他慢慢坐穩了位置.
很多世俗觀念就會出來.
那就很大件事了.
又有自以為已經有那個權.
我沒有任何針對.
你千萬不要介意.
講一般的分析.
他有了那個權之後.
就認為我應該用盡那個權.
去守望教會.
實際上可能他有很多世俗的觀念.
世俗的管理學.
大進教會裡面.
就變成了擴行一些時宮的發展.
你看看佐羅巴巴.
他斬釘截鐵.
你們可以繼續居住.
我們沒有意思去阻撓你們的居住權.
但是興建猶太人的聖殿.
是我們這些被擄歸回的.
甚至乎所以你之前我們不是說.
有些人尋找自己的戶籍.

$^{1481}$找不到.
都不准許他做那個官階.
你是利未人.
不行.
我們要等祭司出現之後.
開始事奉之後.
就要來到去有一個手續.
就是用污淩土名.
即是說.
使得那些找不到自己的戶籍的.
都可以得到證明.
他們真的是利未人.
才可以在那個時宮裡面公積.
所以這些嚴格.
原來他們自己都在做.
所以譬如未洗禮.
真的不可以做其他的事.
不關事說我們歧視.
而是他自己的靈明都明白.
自己為什麼要信耶穌.
和在這個教會裡面一路的成長.
等等.
如果不是.
其實就是找石頭打自己的腳.
就帶來很多很多的問題.
所羅巴巴是很有智慧和決斷.
和你們無關.
我們是做我們為神所做的聖功.
所以他們就不會說.
不行.
不要緊.
他們是決心.
是要來到去擾亂他們.
所以這裡就不是說得太多.
尼熙米記就有差不多四至五章.
是令到尼熙米是夭心夭肺的.
所以原來那個攻擊.
一直存在在神的殿的興建.
神的家的建造.
教會的建造.

$^{1521}$那個攻擊.
那個敵人.
是一直的出現.
所以我們記住.
是有教會的存在在地上.
都是會有困難.
是會有那個攻擊的.
而且有時更多是內部的攻擊.
這個是相當可惜.
相當不值得.
所以他們現在賄賣謀士.
要敗壞他們的謀算.
第六節.
在阿特律財登基的時候.
上本控告猶大和耶路撒冷的居民.
果然新皇興起.
那些前朝的政策.
甚至乎那個詔書.
會不會失傳了.
又或者新的太子登基.
覺得要來到去.
給一些可能同輩的大臣.
慫恿影響.
要更改那些國策.
不奇怪的.
真的會有這樣的情況.
譬如羅波安接續了爸爸.
所羅門做王.
同輩的年輕人.
和他一起成長的在皇宮.
羅波安.
現在有一個叫耶羅波安.
是另一個人.
是以色列十國之派.
其中一個公允的領袖.
因為所羅門.
大家知道.
他是其中一個很大的人生的工程.
就是建造聖殿.
建造他自己的行宮.

$^{1561}$和建造很多的積貨城.
儲糧的.
這樣來到去又打仗.
所以他的一生四十年.
是令到以色列人.
南丁是很辛苦的.
不停建造.
所以當他一死.
耶羅波安就直接來找羅波安.
就說你的父親.
令到我們整個以色列.
過去是很辛苦.
我們很想現在暫時停一停.
讓我們大家回到自己的家庭.
自己的地方.
來稍微休養生息.
羅波安就.
不知道怎樣回答.
問下他身邊的人.
其實不是.
這個人其實是看不起你.
借這個機會藉口.
想奚落你.
想下馬威.
羅波安就說是.
我剛剛登基.
你就這樣對我們.
他就拿捏不到時勢的改變.
所以就給他們.
有一個更加重的要求.
更加痛.
要他們更加做苦工.
自然就帶來的分裂.
所以新王興起.
有好有不好.
有新的不同的政績.
但亦都有些可能會過渡的改變.
所以這裡就來到他們上本控告.
猶大和耶路撒冷的居民.
我們去到第四章尾.

$^{1601}$去到大理烏王第二年.
其實不是.
對 波斯王大理烏第二年.
第七節.
就是那些狀師寫的狀詞.
來上書給亞達薛西.
亞達薛西年間.
比斯蘭米特利達他別和他們的同黨.
上本就告波斯王亞達薛西.
本章是用亞蘭文字.
亞蘭方言.
亞蘭文就是希伯來文字.
省長利宏書記申遂.
要控告耶路撒冷人.
也上本就告亞達薛西王.
省長利宏書記申遂.
和同黨的底那人.
阿法薩提加人.
他比提拉人.
阿法薩人.
阿基維人.
這些全部是外族人.
巴比倫人.
書山城的人.
書山家人.
底亥人.
以蘭人.
和尊大的阿斯拿巴.
所遷移安置在薩馬里亞城.
所以那個中心點.
並帶河西一帶地方的人等.
上奏亞達薛西王說.
河西的臣民云云.
王該知道.
從往那裡上到我們.
這類的猶大人.
已經到耶路撒冷重建.
這反叛惡劣的城.
築立根基.
建造城牆.

$^{1641}$如今王該知道.
他們若建造這城.
城牆完畢.
就不再與王進貢.
交貨納稅.
終究王必受虧損.
我們既食御鹽.
不忍見王吃虧.
因此就告於王.
請王考察.
先王的實錄.
必在其上查之.
這城是反叛的城.
與列王和各省有害.
自古以來.
其中常有背逆的事.
因此這城曾被拆毀.
我們謹奏王知.
這城若再建造.
城牆完畢.
河西之地.
王就無份了.
那時王御福.
省長李宏.
書記申遂.
和他們的同黨.
就是住撒瑪利亞.
並河西一帶地方的人說.
願你們平安云云.
你們所上的本.
已經明讀在我面前.
我已命人考察.
得知此城.
古來果然背叛列王.
其中常有反叛背逆的事.
從前耶路撒冷.
也有大君王.
統管河西全地.
人就給他們進貢交貨納稅.
現在你們要出告示.

$^{1681}$命這些人停工.
使到這城不得建造.
等我降旨.
你們當謹慎.
不可遲延.
為何容害加重.
使到王受虧損呢.
亞達斯西王的上諭.
讀在李宏和書記申遂.
並他們的同黨面前.
即是說這些是波斯的官員.
即是說他們賄賣了這些官員.
變成這次上訪.
這次的告發.
真是帶來果效.
亞達斯西王.
好像不去問前朝.
父王或者古列大帝.
為何要這樣做呢.
不去求問.
反而信以為真.
信了這些地方官的密告.
就帶來要立即使到猶大人要停工.
現在是皇帝立命要即刻停工.
你們可想而知.
那個影響是相當沉重.
因為是立即命令.
要所有的工程要停頓.
所以他們就急忙往耶路撒冷去.
見猶大人.
用勢力強迫他們停工.
即是說拿那個王榜出來.
那個上諭.
是亞達斯西王的上諭.
有那個玉璽蓋印.
現在要你們立即停工.
這樣就完全打窒了.
所羅巴伯他們的士氣.
令到他們都不敢再做什麼.
這是第24節.

$^{1721}$於是耶路撒冷神殿的工程就停止了.
直到波斯王大利烏第二年.
第三個皇帝大利烏第二年.
是另一個大利烏.
不是之前的大利烏王.
那個馬代的大利烏.
即是但以理的大利烏.
現在是波斯的大利烏.
這些是他們的國號的名字.
其實他們有他們的年代.
和他們的真名.
稍後我會在群組.
讓大家來看波斯帝國前面的幾位皇帝.
古列亞達斯西.
這個波斯王的大利烏.
跟著有阿赫徐魯.
他們一直在那裡統治.
所以一朝亞達斯西要他們停頓.
一停就停到新王大利烏開始.
你就可以想像得到相當長.
差不多有16年.
整個建造.
其實當時都未建造那個城牆.
只不過是建造那個聖殿.
其實立了根基沒多久.
就立即出現了對抗.
記住有三部曲.
首先就是我們是和你們一起的.
我們敬拜你們的神.
你們離開了故土.
我們一直守住這個地方.
我們敬拜神.
所以現在一起做.
所羅伯伯說不行.
我們是猶太人.
我們回來建造我們的神殿.
所以記得不要讓未信的人.
在教會裡面有過份的權力.
去統領那個教會.
那個教會就會倒塌.

$^{1761}$他未信根本就是沒有神的靈.
沒有那個屬靈的思考的方向.
所以這個是要切忌的.
那他們就怎樣呢.
你不聽我說嗎.
他們敬酒不喝一罰酒.
他們就來恐嚇.
來捏造一些壞話.
攻擊這些領袖.
或者是他們想來到去.
有些部署想來到去立國.
獨立.
推翻古烈.
類似這樣的謠言.
接著就更加訴諸於政治.
直接賄賂一些謀士.
賄賂一些地方官.
直接去向最高的波斯.
阿達舍西王.
來到去寫那個詔書.
寫那個告狀書.
來到去有那麼差.
說到那麼差.
他們如果建好了城牆.
他們就不會再來到去進貢納粹.
就是這樣.
彼拉多被那些猶太人說他.
你如果不殺耶穌.
你就是背叛該殺.
彼拉多就不敢出聲了.
因為他也是巡撫.
所以可想而知.
三部曲.
外邦人他們會這樣來對待猶太人.
首先就是表面好像很好.
實際上是想混雜了那個信仰.
不行就會來到去阻撓那個建造.
再不行.
因為見到他們那麼同心.
不受到影響.

$^{1801}$就直接用政治的介入.
改變了他們那個建造.
最重要就是令到他們立即停工.
誰能夠改變這個現象呢.
猶太人沒有了建造之後的.
那個共同的信念.
上來為什麼呢.
不是上來.
是建一個地方來住.
他們在外邦已經是落地生根了.
所以沒有了這個共同的理想.
那個目的.
就一盤散沙了.
就出現了.
我們揭一揭去很後.
其實不關編排的事.
記得其實同時代.
還有三卷書就是哈該,撒迦利亞和馬拉基書.
這三卷被老歸回的先知書.
哈該書第一章.
其實是馬太福音前三卷.
我們稍微讀一點.
就是將這件事記載了.
大利烏王第二年六月初一.
剛才說一路停到大利烏王的第二年.
現在在六月初一.
接近七月的節期.
哈該的意思就是耶和華的節期.
希伯來文.
哈該這個名字.
耶和華的華籍先知哈該向猶大省長.
猶大省長.
即是說他上到去之後.
在那個地方.
王可能分封他做猶大的省長.
就是撒拉提的兒子所羅巴伯.
和約撒達的兒子大祭司約書亞說.
萬君之耶和華如此說.
這百姓說.
關於百姓的說話.

$^{1841}$記住是已經停工了.
大家都沒有一個方向.
大家都沒有一個焦點.
大家都沒有地方去敬拜神.
自然就沒有人有興趣回教會.
教會都不存在.
還怎樣回教會呢.
回哪間教會呢.
你可以想像得到.
這百姓說建造耶和華殿的時候尚未來到.
那時耶和華的華臨到先知哈該說.
這殿仍然荒涼.
你們自己還住天花板的房屋嗎.
現在萬君之耶和華如此說.
你們要審查自己的行為.
你們殺的種多.
收的卻少.
你們吃卻不得飽.
喝卻不得足.
穿衣服卻不得暖.
不工錢的.
將工錢裝在破漏的囊中.
就是說他們的生活.
一點都不會光彩.
興奮和滿足.
是很多的欠缺.
這是物質上的欠缺.
反映內心的心靈的欠缺.
因為政治沒有了焦點.
上來這麼辛苦.
連根拔起就是想起神的殿.
有宗教的敬拜.
有神同在.
整個民族有尊嚴光彩.
他們不知道孕育彌賽亞的降生.
但是現在沒有了方向.
敵人令到你停工.
在亞達斯西王的年間.
要你停工.
現在一直停到大利烏王的第二年.

$^{1881}$現在在大利烏王第二年的六月初一.
記載得很清楚.
哈該聽到神的預言.
跟他說.
要他們首先來勸勉.
留意第七節.
萬君之耶和華如此說.
你們要審查自己的行為.
你們要上山取木料建造這殿.
我就因此喜樂且得榮耀.
這是耶和華說的.
你們盼望多得所得的卻少.
你們收到家中我就催去.
這是為什麼呢.
因為我的殿方涼.
你們各人卻顧自己的房屋.
這是萬君之耶和華說的.
所以為你們的緣故.
天就不降金路.
即是說那十幾年.
大家都過得很乏味.
大家都覺得好像舉步維艱.
做很多事情都不是那麼順暢.
因為先是心靈不再順暢.
大家都沒辦法.
當然顧自己家庭.
做好自己的本分.
但是也沒有了力量.
一起同心的去建聖殿.
因為你什麼時候才算好呢.
難道等皇帝突然頒佈詔書嗎.
是不可能的.
所以沒有人可以說了事.
即是說連猶大的省長索羅巴伯.
大祭司約書亞.
都不能夠兩個人一鎚定音.
但是神可以興起第三類的人.
第三個職份的.
就是先知.
所以猶太人是聽先知說的.

$^{1921}$先知現在有神的話.
所以耶和華如此說.
是等同於最高權威.
沒有人可以違抗.
這就是告訴我們.
神現在親口這樣說.
說他們的光景.
一而再兩段都說他們的光景.
所以如果近來少了教會.
你自然心裡會很混亂.
很多事情你都很不開心.
或者很多問題.
因為你已經遠離了神.
教會或者當時的猶太人都是這樣.
原來神令到天不降金路.
地也不出土產.
那你多拼命都沒有用.
天不下雨.
地不出土產.
我命乾旱淋到地土.
山江五谷新走和游.
並地上的出產人民生畜.
以及人手一切勞碌得來.
都令到這些全部因為乾旱淋到大地.
其實還是很有趣的.
就是只是猶大那個地方.
那就是說整件事.
神是有目的地要挑動他們的心.
其實最重要的一大段聖經.
是那個「審察」這個字.
說了兩次要審察自己的行為.
那個時候.
其實是一個反思.
為什麼我們這四萬九千多人.
上來耶路撒冷呢.
為什麼要這樣做呢.
我們就是要重建神的殿.
要在這個地方有神的敬拜.
因為神規定了激動我們回來.
是神的喜悅.

$^{1961}$是神的榮耀.
為什麼我們沒有做呢.
因為政治的力量.
難道神的力量不大過政治嗎.
所以第十二節.
那時撒拉鐵的兒子所羅巴伯.
和約撒達的兒子大祭司約書亞.
並盛夏的百姓.
都聽從耶和華他們神的話.
先知夏該奉耶和華他們神.
差來所說的話.
百姓也在耶和華面前.
傳敬畏的心.
就是這個火花.
那個火花就是.
夏該說.
耶和華這樣和我說.
叫我勉勵你們.
還有為什麼你們這十幾年.
在亞達舍西王的年間.
你們會這麼困難.
你們的生活是這麼不如意.
內心是這麼枯竭.
什麼原因呢.
就是因為神的電方良.
這是意象.
其實天天都看到.
我們要上去守節.
去哪裡守節.
還是只有根基.
什麼都沒有.
守什麼節呢.
你就會想像得到.
他們就慢慢慢慢見怪不怪.
因為都沒有這個能力.
哪個人可以一呼百應.
哪個人可以統領群雄.
沒有.
所以神就出手了.
感動夏該.

$^{2001}$所以這些就是神的工作.
神先感動先知.
先知領受了神的話語.
在大家都是很疲乏的心靈底下.
大家其實不是不做事.
他們可能很勤奮的去耕種.
去做其他的修建.
不過是個人自己做個人的事.
一起做神的事情.
所以也是其中一個指標.
如果教會的弟兄姊妹.
對神家的事根本不在意.
不上心.
不一起去做.
大家個個都是獨當一面的.
在世界的工作.
本來是好事.
不過在教會.
個個都闊綽濫利.
那就很大件事了.
就好像回到夏該的年代.
他們就說.
不是時候了.
等亞達四世王先死了.
可能是這樣也說不定.
不過最終帶來的就是.
一路神的殿都是荒涼停工.
但突然之間.
夏該注入的那個生氣.
注入的就是.
耶和華如此吩咐我.
我就這樣去默示給大家知道.
那索羅巴巴和大祭司約書亞.
就有那個異象了.
點燃他們內心的那團火.
就來到好像少林足球一樣.
突然間大師兄或者誰.
就有那個火花了.
立刻有那個心意.
要復興少林功夫.

$^{2041}$現在索羅巴巴和約書亞就是這樣.
來勉勵.
所以都是感動索羅巴巴和大祭司去做.
先知只是代言.
最終是感動那個領袖.
和那些群眾一起做.
所以這個是很重要的一件事.
那個方程式.
神是先來到去感動先知.
先知好像違反那個時代.
違反那個常規.
常規就是還不是時候.
現在那個王阻撓得很厲害.
讓那個王先過去.
但先知好像看到的是再遠一點的那個景況.
就是只要大家同心起來.
就可以改變那個環境.
這些都很重要的.
如何拿捏那個是最好的時機.
那個時代.
又能夠令到整體的百姓.
都能夠大家一起同心.
上下一心.
這個又是需要那個先知來抓緊那個時間.
這個時候好像已經差無可差了.
神要出手了.
神要來到感動哈該了.
其實還有撒迦利亞的.
撒迦利亞十一章聖經.
有很多預言的撒迦利亞.
來到他不斷.
在這段時間.
大利烏王第二年是很重要的日子.
大家記住這個被擄歸回.
大利烏王.
波斯王的大利烏第二年.
是很重要的日子.
因為神終於出手了.
就是說那些百姓.
那個慘況.

$^{2081}$一直這樣來到去.
持續的原地踏步.
十幾年.
神的時間到了.
不可以再這樣了.
神就感動哈該.
由哈該公開的向兩位領袖.
和其他的領袖.
和百姓有這個宣告.
所以這個就是.
第十四節.
耶和華激動猶大省長.
薩拉鐵的兒子.
所羅巴伯.
和約薩達的兒子大祭司約書亞.
並盛下之百姓的心.
他們就來為萬軍之耶和華.
他們神的殿作供.
留意第十五節.
好像有個猶蹉.
這是在大利烏王第二年.
六月二十四日.
六月初一.
先知有神的密使.
感動.
用了二十幾天.
來跟他們做教導.
查經班.
奮興會.
感動他們.
終於在六月二十四日.
所羅巴伯.
約書亞.
他們兩位領袖.
是被擄歸回的領袖.
有這個代表性.
所以要那兩個領袖自己都復興.
其他下面的人都要復興.
要好的靈性.
所以用了二十幾天的時間.

$^{2121}$來勉勵兩位領袖.
二十幾天.
可想而知.
人的心需要慢慢的煲熱.
熱了之後.
他就會有那個力量.
其實應該是回到那個初心.
今天說的初心.
就是他為什麼要回來呢.
為什麼要帶領他的族民回來呢.
就是要重建神的殿.
今晚我們先來到這裡.
我們就說到這裡.
我們一起禱告.
再一次多謝主.
我們能夠將被擄歸回的歷史書.
先知書.
一併的去讀.
我們又能夠結合當時的時代.
那個問題.
和當地人的反對.
這些和我們今天.
或者任何一個時代的教會.
都是完全一樣的.
我們很盼望我們能夠.
真是屬神的子民.
我們剛強奮勇.
我們能夠靠主.
來建造神的國度.
不信的人.
會帶一些世俗觀念.
甚至乎是.
一些的謠言.
一些的攻擊的說話.
又或者是一些的論斷.
批評.
令到那個.
大家的童心會受到瓦解.
更甚的是會用政治.
來迫害猶太人.

$^{2161}$雖然這個聖殿.
一直停頓.
不能夠建造.
但是神你的救恩的工作.
彌賽亞的降臨.
在神的殿裡面.
將來那個.
事公.
是一路會發生的.
所以我們看到神自己.
來感動先知.
來再感動那個領袖.
我們很盼望.
教會都能夠.
吸取到.
這些重要的真理.
我們真的能夠.
常常有神的.
分面教導在教會裡面.
和我們的領袖.
都是敬畏神.
愛主.
愛神的道.
和願意聆聽神的真理.
以致當我們很多的決策.
不是只是一種純粹.
理性的分析.
而是有神的聖靈.
感動我們大家一起做.
和弟兄姊妹.
亦都很願意去跟隨.
一起這樣去做.
求主賜給我們有方向.
有意象.
有清晰的大靈.
求主給我們在過年期間.
我們能夠和家人有美好的相聚.
亦都能夠散發出.
我們基督徒獨有的.
神國度的特質.

$^{2201}$不單純是一些迷信的.
或者是過年只是.
其中的可能是賭博.
可能是百無聊賴等等.
這樣慢慢又蹉跎一年.
盼望我們能夠將神的愛.
透過我們的生活去反映.
我們這樣感恩禱告.
奉耶穌基督的名.
阿們.
\newpage



\section{馬太福音 13:24-30-36-43}
\label{sec:c4O5F1DLbjI}
\textbf{2023-01-15 加恩樓《成為麥子被收割》馬太福音 13\_24-30,36-43 講員:陳麗蟬牧師}
\newline
\newline
連結: \href{https://youtube.com/watch?v=c4O5F1DLbjI}{\texttt{https://youtube.com/watch?v=c4O5F1DLbjI}} ~~~~ 語音日期: 2023-01-17
\newline
\newline
\hyperref[sec:ZwtOgh6Qmm4]{\small{< < < PREV SERMON < < <}}
~
\hyperref[sec:index_chronic]{\small{[返順時目]}}
~
\hyperref[sec:index_scriptual]{\small{[返順卷目]}}
~
\hyperref[sec:ezP3uupXjgk]{\small{> > > NEXT SERMON > > >}}
\newline
\newline
馬太福音 13:24-30-36-43
\newline
\begin{longtable}{cl}
\hline
\hline
章節 & 經文 (和合本修訂版)\\
\hline
13:24 & \begin{tabularx}{0.7\textwidth}{X} 耶穌又設個比喻對他們說：「天國好比人撒好種在田裡， \end{tabularx} \\ \\ \relax
13:25 & \begin{tabularx}{0.7\textwidth}{X} 在人睡覺的時候，他的仇敵來，把雜草撒在麥子裡就走了。 \end{tabularx} \\ \\ \relax
13:26 & \begin{tabularx}{0.7\textwidth}{X} 到長苗吐穗的時候，雜草也顯出來。 \end{tabularx} \\ \\ \relax
13:27 & \begin{tabularx}{0.7\textwidth}{X} 地主的僕人進前來對他說：『主人，你不是撒好種在田裡嗎？哪裡來的雜草呢？』 \end{tabularx} \\ \\ \relax
13:28 & \begin{tabularx}{0.7\textwidth}{X} 主人回答他們：『這是仇敵做的。』僕人對他說：『你要我們去拔掉嗎？』 \end{tabularx} \\ \\ \relax
13:29 & \begin{tabularx}{0.7\textwidth}{X} 主人說：『不必，恐怕拔雜草，也把麥子連根拔出來。 \end{tabularx} \\ \\ \relax
13:30 & \begin{tabularx}{0.7\textwidth}{X} 讓這兩樣一起長，等到收割。當收割的時候，我會對收割的人說，先把雜草拔出來，捆成捆，留著燒，把麥子收在我的倉裡。』」 \end{tabularx} \\ \\ \relax
13:31 & \begin{tabularx}{0.7\textwidth}{X} 他又設個比喻對他們說：「天國好比一粒芥菜種，有人拿去種在田裡。 \end{tabularx} \\ \\ \relax
13:32 & \begin{tabularx}{0.7\textwidth}{X} 它原比所有的種子都小，等到長起來，卻比各樣的菜都大，且成了樹，以致天上的飛鳥來在它的枝上築巢。」 \end{tabularx} \\ \\ \relax
13:33 & \begin{tabularx}{0.7\textwidth}{X} 他又對他們講另一個比喻：「天國好比麵酵，有婦人拿來放進三斗麵裡，直到全團都發起來。」 \end{tabularx} \\ \\ \relax
13:34 & \begin{tabularx}{0.7\textwidth}{X} 這都是耶穌用比喻對眾人說的話，不用比喻，他就不對他們說甚麼。 \end{tabularx} \\ \\ \relax
13:35 & \begin{tabularx}{0.7\textwidth}{X} 這是要應驗先知所說的話：「我要開口說比喻，說出從創世以來所隱藏的事。」 \end{tabularx} \\ \\ \relax
13:36 & \begin{tabularx}{0.7\textwidth}{X} 當時，耶穌離開眾人，進了屋子。他的門徒進前來，說：「請把田間雜草的比喻講給我們聽。」 \end{tabularx} \\ \\ \relax
13:37 & \begin{tabularx}{0.7\textwidth}{X} 他回答：「那撒好種的就是人子， \end{tabularx} \\ \\ \relax
13:38 & \begin{tabularx}{0.7\textwidth}{X} 田地就是世界，好種就是天國之子，雜草就是那惡者之子， \end{tabularx} \\ \\ \relax
13:39 & \begin{tabularx}{0.7\textwidth}{X} 撒雜草的仇敵就是魔鬼，收割的時候就是世代的終結，收割的人就是天使。 \end{tabularx} \\ \\ \relax
13:40 & \begin{tabularx}{0.7\textwidth}{X} 正如把雜草拔出來用火焚燒，世代的終結也要如此。 \end{tabularx} \\ \\ \relax
13:41 & \begin{tabularx}{0.7\textwidth}{X} 人子要差遣他的使者，把一切使人跌倒的和作惡的從他國裡挑出來， \end{tabularx} \\ \\ \relax
13:42 & \begin{tabularx}{0.7\textwidth}{X} 丟在火爐裡，在那裡要哀哭切齒了。 \end{tabularx} \\ \\ \relax
13:43 & \begin{tabularx}{0.7\textwidth}{X} 那時，義人要在他們父的國裡發出光來，像太陽一樣。有耳的，就應當聽！」 \end{tabularx} \\ \\
[1ex]
\hline
\hline
\end{longtable}
$^{1}$請姐妹平安,因為是農曆之前的最後一個星期日,我們在這裡一起用詩歌去讚美上帝,用他的說話我們彼此勉勵..
在馬太第十三章裡面,全部都是有關一些的比喻,這些比喻都是和天國有關的.耶穌去說這個比喻的時候,他都是會用當時他們所熟悉的環境,或者大自然,總之那些人一聽就明白耶穌基督在說什麼..
不過對於我們這一幫在現代的人,可能就有一點點的困難了..
耶穌就藉著這些比喻,就要去幫助一幫信徒,讓他們知道在地上要活好見證.這次這個比喻在十三章,其中一個比喻就是在說墨子和一些雜草..
首先我們就看一下,這個墨子和雜草是什麼東西.墨子我們知道是一些有用的食物,通常在歐洲,在亞洲一些地方,甚至我們中國大陸,我們北面都主要是吃小麥的..
這些麥就不可以像稻米那樣吃的,粒粒的就很難吃的,通常磨成粉,然後製成麵包,製成麵條,甚至通心粉都是用麥做的,啤酒也是用麥做的,是一些有用的食物..
另外一種就是在經文裡面說到的一些雜草,不過在舊的和合本裡面就很清楚指出,這些雜草的名稱就叫做「敗子」..
這些敗子就是一些沒什麼用的植物,甚至是對人的身體有害的,如果你不小心吃了下去就糟糕了,輕的就會暈頭頭,或者你又會嘔吐,會抽筋,如果嚴重的話你就會死的,不過通常很少人死的,為什麼呢?.
因為它是很苦的,所以你不小心吃下去的時候就會立刻苦到吐出來,所以通常都是一些輕的..
這個麥子和這些敗子的樣子是很相似的,但是我想種麥的人就會比較容易去分辨,正如在加州的錦繡也有,或者鄉村的地方,人們一看龍眼樹,荔枝樹就分得出,我們不認識的那些看的都是一樣的..
這個麥子和這些敗子在他們的細芽細苗的時候,或者在他們成長的時候,就好像一群兄弟一樣,樣子是一模一樣的,他們生長的環境都是一樣的,到什麼時候才能夠分辨到呢?.
就是當他們長大,開始有些麥子和敗子的樣子凸出來的時候,就會清楚的能夠看得到一點點,雖然看完之後可能對真實的時候都不太懂得分辨..
這個比喻就說到那些僕人,照常去到麥田那裡可能淋水弄理,他們就發現在麥田裡面多了很多這些雜草,敗子,他們就去問主人,主人,你殺種的時候不是殺好種的,為什麼突然間會有這麼多雜草呢?.
主人就說,不是我做的,我殺的那些是好種來的,這些是受敵做的..
在耶穌的時代,當大家和隔壁的鄰居又說兩句,或者有什麼不對的話,他們通常會怎麼做呢?.
他們就會靜靜的去到隔壁的麥田那裡,可能是一些水輸到麥田那裡,他們就將敗子的種子灑在水裡面,等水一直流到麥田的時候,除了有麥子的成長,敗子也會來去成長..
或者當他們走過隔壁的麥田的時候,他們就靜靜的在袋子裡拿敗子的種子出來,一灑下去..
那時候那些人就是這樣,他們就要去破壞別人的農作物.不過這個報仇的情況,我想不只是在耶穌的時代,現在我們也會有的..
我記得那時候我在一間基督教小學裡面去代課,代一年的那些,那個校長是基督徒,我就很詫異,為什麼他會編一個佛教徒教六年級的聖經?.
我心裡就想,為什麼?後來我就知道原來這位佛教徒老師曾經得罪過校長,所以校長就這樣報復,讓他沒有好日子過..
那這個僕人就看到這個情況之後,他就跟主人說,現在麥田裡有這麼多的敗子,是不是要把它們都拔掉?.
因為真的,那些僕人這麼用心去種這個麥田,但現在竟然有一些雜草在那裡,搶了那些的營養,那當然是很想把它們都拔掉..
這兩三個星期,我在家裡學人種一些盆栽,我又不懂事,我真的不懂事,我以為在碟子裡就多倒一點水給他們,讓他們一直吸收水,有時候真的趕著上班,忘記了,不倒水也可以..
誰知道後來,我前幾天發現有很多小蚊子跑出來,我說為什麼會這樣?於是那些人就教我,總之就是把那些污水清理乾淨,然後洗乾淨碟子就可以了..
我以為真的可以了,洗了一兩天,好一點了,誰知道昨天早上我起來,一走過桌子上的時候,我又看到幾十隻蚊子,很小的,飛來飛去,我覺得很討厭,拍拍拍,拍不到的,他們飛得很快..
我記得有人教我,用一些香茅水噴一下桌面,我看到很多小蚊子就開始在泥裡出來,我又氣又像你,就把泥噴了,然後我看到他們在上面撈來撈去,於是我又把葉子噴了..
我想我很明白那些僕人的心態,就說我這麼用心種這些田地,你們在搶我的營養,很想把我拔走,那些小蚊子在飛來飛去,很滋擾,我很想殺死他們..
耶穌說的這個比喻,他用門徒去明白一件事,在這個世界上一定會有好人,也會有一些惡人..
根據我剛才說的那位基督徒校長,我相信大家都會覺得他是惡人,竟然作為一個基督徒的校長,可以請一個佛教徒去教聖經,是惡人,不過上帝將來都會審判的..
但據我所知,這個佛教徒的老師,他之前曾經在周圍的老師,甚至在外面,到處唱這個校長的一些不好的東西,總之就是把別人的東西揭發出來..
照理來說,這個佛教徒的老師,都是一個惡人,大家都是惡人,所以就要想想,誰是好人,誰是壞人,有沒有一定的標準呢?.
或者,好人是不是永遠都是好人,壞人永遠都是壞人呢?耶穌在解釋這個比喻的時候,他就說這個敗旨是一些敵對上帝的勢力..
又指一些假信徒,又指一些掛名的信徒,說到掛名的信徒,我心想,慘了,我是不是說掛名的信徒呢?.
又想想,這個敵對上帝,有時,雖然我們信了耶穌,甚至做了牧師,但有時我的心態都覺得,嘩,上帝有沒有搞錯呢?.
有時我的心思意念,都會有些反叛上帝的,如果我真的對上帝有些反叛,就問上帝有沒有搞錯的時候,那我就在那一刻,我就成為了對上帝敵對的..
如果我們一旦敵對的時候,上帝就會立即除掉的,即是那些僕人說,是不是現在這麼多敗旨在這裡,我們立即要全都拔掉呢?.
坦白說,我們一旦不聽話,上帝就將我們完全毀滅的時候,這樣在這個世代,就沒有剩下多少真信徒在這裡,是不是?.
那又想想,在這個世代的上高,有多少人能夠真真正正成為了上帝眼裡的墨子呢?.
所以我們就不要去判斷更多的東西,我預備這篇講道的時候,我在網上看到了一篇很有趣的漫畫,.
他說不要去判斷哪些人是墨子,哪些人是敗旨,最重要的是確定自己都是天國的子民,我們要活出天國的特徵來..

$^{41}$因為有時候都會的,我們就會去評論一下,哪個是好的信徒,哪個是不好的信徒..
認識有三位的姐妹,很好朋友的,像唐獅頭一樣,有一次他們就一起去短孫一個月,.
回到家裡的時候,雖然未至於反目成仇,不過大家以後就不抽不彩了,.
沒多久,A姐妹就走來跟我聊天,她就說在短孫一個月很不開心,B姐妹不好,C姐妹不好,.
接待我們那個孫教士都不太好,我聽到,他們那些人是這樣的,孫教士都不太懂得做孫教士,.
沒多久,B姐妹又來跟我聊天,B姐妹又告訴我,A姐妹真的不行,跟她一起合作不行,.
C姐妹又不要提了,那個孫教士都是一般般,就是說A,B,C和孫教士都不是很好,都是一般般..
不過我又對他們整個短孫的生活的圖畫,我又知道多一點了,我又發覺到一件事,.
除了A,B兩位姐妹來跟我聊天之外,我算過了,其實A,B,C和孫教士,他們四個人都是一般般,.
因為跟我們來的那些人都是在說自己的好事,在數其他三個不好的,所以其實全部都不好..
在某一次,這位孫教士回來香港述職,他不認識我的,剛好我去聽他的分享,.
聽到他一直在說在當地的生活,事奉一些處境的時候,我又明白多一點了,孫教的童工,.
他真的按照當地的文化,按照警察會給他的一些規定,一些處理的手法去做的,.
我又明白到為什麼A,B姐妹都覺得孫教士不好呢?他有他的難處在這裡..
整件事情我心裡就想了一下,誰對誰非,誰好誰壞,真的很難判斷,.
就好像這幅漫畫說的,我們不需要去判斷誰好誰壞,我們都相信每個人都有好的地方,也都有不好的地方..
中國人有一句話,一隻手就拍不響的,一定要兩隻手一起拍才有響的,.
所以一個人就不會很激氣,或者有什麼錯事出現,唯有一起合作的時候才會擦出這個火花..
這個就是在這個世界上一定有這個物質,有好的,也有這個敗子,有不好的,.
不過你真的用我們有限的人的眼光,我們分不清的..
我們繼續看一下這個敗子,當僕人很想很想去拔掉這些雜草,.
但主人怎麼回答呢?他說不需要,因為怕你拔這些雜草的時候,會連這個麥子的根都拔出來,會一起死的..
因為他們兩種植物既然一起在這裡成長的時候,這個麥子和這個敗子,那些根一定黏在一起的,.
如果你一會兒把這些電線收起來,你不要把它們都繞在一起收起來,你隨便一棵一棵放在一起,.
你一個星期回來拿出來,你會把它們都繞在一起的,根都是這樣的..
所以主人就說,你千萬不要在這個時間拔掉這些敗子,連那些麥子也拔出來,這個會死的..
按著耶穌基督這個比喻,祂說在這個世界上,好人,惡人,他們都會一起共存的,.
大家都可能認同,在這個世界上,我們都覺得有些人很好的,有些人很惡的,.
不過我自己屬於惡人還是好人呢?可能這邊的人覺得我是好人,這邊的人又覺得我是惡人,.
但無論怎樣也好,上帝說一定是共存的,在這個暫時的世界..
主人就說,這樣拔掉很不容易的,主人說,當收割的時候,.
我會對收割的人說,先把雜草拔出來,菌成菌,留著燒,把麥子收在我的倉內..
你放心,到收割的時候,我就會叫人把這些雜草的敗子拔出來,燒掉它,.
有用的麥子就放進糧倉庫,那些敗子就燒死了..
耶穌基督就解釋這個比喻,收割的時候,就說到世代的終結,.
在裡面,他就會來審判,但在這裡收割的是哪一個人呢?.
就是這個天使,由上帝親自差派人來審判..
所以,對一個絕對公平,公正的上帝,所有的事情,祂都會來審判..
不過,祂未必在現在就來審判,而是到了世界的終結的時候,祂一定來審判..
就在這段經文告訴我們,審判的人不是你,也不是我,.

$^{81}$不過我們很多時候就喜歡做審判的人,但耶穌就說,審判的人不應該是你和我..
我們不能夠將上帝的審判權拿走,我們就自己定了誰是惡人,.
然後就將這個惡人飛出局,不行,不行..
上帝會自己親自到來,派使者來做審判..
上帝容忍著這個敗子,容忍著這個麥子,一起在這個田地裡面去成長,.
究竟為什麼呢?在彼得前書,彼得後書,第三章他們說,.
我們有人以為上帝是在耽言,其實祂不是在耽言,.
說上帝為什麼做事這麼慢?不是的,不是的,而是在寬容,.
祂不願一人沉淪,乃願人人都悔改得救..
所以上帝是很想更多的人去得到這個救恩..
昨天早上很生氣,將香茅噴在桌子上,再噴在泥上,.
再把火噴在葉子上,因為真的到處飛,.
昨天沒有回家,想著想著,糟了,今晚回去,.
今晚回去,那些所謂的植物不知道會變成怎樣..
我想想,如果我只要忍耐一下,我忍耐一下的話,.
可能那些植物不會這麼快死的,蚊子又不會咬人,.
我容忍一下牠飛一下飛一下,忍一下就行了,.
但就是因為我不肯忍,所以我相信我今晚回去看牠們,.
這個蚊子和這個草也可能死了一點..
所以上帝那個遲延到來的審判,.
他為了給更多人機會悔改,可以得到永恆的生命..
蚊子的結局就是被火燒,這個蚊子又怎樣呢?.
在第43節裡面說,.
「義人要在他們父的國裡發出光來,像太陽一樣.」.
這個義人不只是好人在聖經裡面,.
是說到跟隨上帝的人,跟隨耶穌的人,.
跟隨耶穌的人,他們要在父的國裡面發出光,.
像太陽一樣..
這個父的國除了指將來永恆的國度,.
也指到在今日地上高的國,.
因為在主禱文講到,「願上帝的國降臨.」.
怎樣降臨呢?.
就是指我們一班的信徒,.
在地上高,在這個黑暗的世界裡面,.
我們為主來發光,行出上帝的人外的時候,.
就是上帝的國度臨到地上了..
在這個不公義,不公平的世代裡面,.
耶穌給我們有命令,.
「你們在世上的光,成灶在山上,是不能隱藏的.」.
「人點燈,不放在斗底下,而是放在燈台上,.

$^{121}$就照亮一家的人.」.
在古時候,沒有街燈,.
所以人們如果晚上要外出,.
或者工作完,太陽已經收了,.
很黑的時候,他們就要抹黑,.
然後回家,很淒涼的..
但是一些很有心的人,.
他就會在窗台上,.
就放一盞油燈,很少油燈,.
以致到晚上要抹黑回家的人,.
走路的人,就藉著這一點點的光,.
他就可以看得清路上哪裡有石頭,.
去到哪裡要轉彎,.
他就會走得有信心一點,安心一點..
我在今年第三年,.
我在杜豐山收取一些課程,.
一年兩次我們要去退休,.
一月和十二月,.
明天又上山四天去營地..
有幾位同學,他們很喜歡看日出,.
所以每次我都和他們,.
我們應該下星期有三次去看日出..
我們要不就走去大十字架那裡,.
要不就出了杜豐山,.
去信義中再上半山..
因為又是要晚上出,.
五點多,很黑的..
在山那裡又沒有街燈,.
所以又要抹黑..
那些同學很厲害的,.
他們不用拿電筒,.
我一定要拿著手提電筒,.
就照實地走..
走著走著,.
我聽到前面的同學小聲大叫,.
知不知道什麼叫小聲大叫?.
「哎呀,哎呀,好可愛,好可愛,好可愛」.
因為太空礦又不敢大聲..
原來看到有幾隻螢火蟲,.
他們很開心,.

$^{161}$「嘩,有螢火蟲啊」.
在漆黑黑裡面,.
看到兩三點的光,.
很少很少的,這兩三點的光,.
但就是這兩三點的光,.
帶給人那份興奮..
在古時代,.
看到這些光可以平安走路回家,.
帶給人那份力量..
所以,弟兄姊妹,.
在這個世界,.
我們暫時在這個地方,.
是由惡者所統治的,.
是幽暗的,是邪惡的,.
這是約翰一書所告訴我們的..
但不過,他下一句就說,.
因為我們都是屬於上帝的,.
是光明的,.
而且是至宗德性的..
所以,我們千萬不要留在一些黑暗的情況,.
也不要留在一些很多的埋怨,.
很多的不滿,很多的批評裡面..
我們也不要去想,.
誰是好的信徒,.
誰是惡的人,.
我們不需要,.
像剛才那幅漫畫告訴我們,.
我們只要做好我們自己天國指紋的崗位,.
我們也不要覺得,.
我們只是一個很小很小的力量,.
我們自己一個,.
我們經常說我們自己很小,.
小薯仔,.
我不行,我不行的,.
但是,上帝所看重的,.
是我們有沒有努力去為祂做功..
第43節,我再讀一次,.
他說,義人要在富的國裡面,.
發出光來,像太陽一樣..
我們跟隨耶穌基督的人,.

$^{201}$耶穌基督就是要求我們,.
就是在祂暫時由撒旦所管轄的世界裡面,.
我們就算發一點點的光,.
按照祂給我們的恩賜,.
有多少發一點點的光,.
都可以帶給人興奮,.
帶給人力量,.
這樣就足夠了,.
我們一起禱告吧..
真是天乎上帝,.
我們在這個世界裡面,.
我們真的會看到很多不公義,.
很不公平的事情,.
但是,我們在不同的角度裡面,.
或者不同人的體裡面,.
又會有不同的答案,.
我們對人也有很多不同的評價,.
但我相信,沒有一個人,.
所有人都覺得他是惡人,.
或者所有人都覺得他是好人,.
所以我們謝謝你,.
藉著耶穌基督的教導,.
我們看到,.
我們要將審判權交回給上帝,.
幫助我們吧,.
在地上高,.
我們不需要去理會那麼多,.
就將一切這些審判交給上帝,.
我們在地上高,.
幫助我們,.
聖靈你幫助我們,.
給我們力量,.
藉著你給我們不同的恩賜,.
我們就在地上高,.
為主你來發光,.
縱然我們所作的很微小,.
但我們多謝上帝,.
因為你不會介意,.
你只要我們來去忠心,.
我們多謝你,.

$^{241}$再來就恭敬張嘉欣樓,.
每一位弟兄姊妹,.
交到主你的金錢,.
願意每一個弟兄姊妹的生命,.
都為主你來去作精兵,.
時刻的為主你來去作見證,.
為主你來去發光,.
我們多謝你,多謝主的恩典,.
願意我們都是,.
每一個我們都成為了主你所使用的墨紙,.
多謝你瑞廷禱告,.
奉主基督聖名,.
阿們..
\newpage



\section{以賽亞書 40:1-11-27-31}
\label{sec:ezP3uupXjgk}
\textbf{2023-01-15 午堂《如鷹展翅上騰》以賽亞書 40\_1-11,27-31 講員:郭鴻標牧師}
\newline
\newline
連結: \href{https://youtube.com/watch?v=ezP3uupXjgk}{\texttt{https://youtube.com/watch?v=ezP3uupXjgk}} ~~~~ 語音日期: 2023-01-26
\newline
\newline
\hyperref[sec:c4O5F1DLbjI]{\small{< < < PREV SERMON < < <}}
~
\hyperref[sec:index_chronic]{\small{[返順時目]}}
~
\hyperref[sec:index_scriptual]{\small{[返順卷目]}}
~
\hyperref[sec:ApJnSQORNII]{\small{> > > NEXT SERMON > > >}}
\newline
\newline
以賽亞書 40:1-11-27-31
\newline
\begin{longtable}{cl}
\hline
\hline
章節 & 經文 (和合本修訂版)\\
\hline
40:1 & \begin{tabularx}{0.7\textwidth}{X} 你們的神說：「要安慰，安慰我的百姓。 \end{tabularx} \\ \\ \relax
40:2 & \begin{tabularx}{0.7\textwidth}{X} 要對耶路撒冷說安慰的話，向它宣告，它的戰爭已結束，它的罪孽已赦免；它為自己一切的罪，已從耶和華手中加倍受罰。」 \end{tabularx} \\ \\ \relax
40:3 & \begin{tabularx}{0.7\textwidth}{X} 有聲音呼喊著：「要在曠野為耶和華預備道路，在沙漠為我們的神修直大道。 \end{tabularx} \\ \\ \relax
40:4 & \begin{tabularx}{0.7\textwidth}{X} 一切山窪都要填滿，大小山岡都要削平；陡峭的要變為平坦，崎嶇的必成為平原。 \end{tabularx} \\ \\ \relax
40:5 & \begin{tabularx}{0.7\textwidth}{X} 耶和華的榮耀必然顯現，凡有血肉之軀的都一同看見，因為這是耶和華親口說的。」 \end{tabularx} \\ \\ \relax
40:6 & \begin{tabularx}{0.7\textwidth}{X} 有聲音說：「你喊叫吧！」我說：「我喊叫甚麼呢？」凡有血肉之軀的盡都如草，他的一切榮美像野地的花。 \end{tabularx} \\ \\ \relax
40:7 & \begin{tabularx}{0.7\textwidth}{X} 耶和華吹一口氣，草就枯乾，花也凋謝。百姓誠然是草； \end{tabularx} \\ \\ \relax
40:8 & \begin{tabularx}{0.7\textwidth}{X} 草必枯乾，花必凋謝，惟有我們神的話永遠立定。 \end{tabularx} \\ \\ \relax
40:9 & \begin{tabularx}{0.7\textwidth}{X} 報好信息的錫安哪，要登高山；報好信息的耶路撒冷啊，要極力揚聲。揚聲不要懼怕，對猶大的城鎮說：「看哪，你們的神！」 \end{tabularx} \\ \\ \relax
40:10 & \begin{tabularx}{0.7\textwidth}{X} 看哪，主耶和華必以大能臨到，他的膀臂必為他掌權；看哪，他的賞賜在他那裡，他的報應在他面前。 \end{tabularx} \\ \\ \relax
40:11 & \begin{tabularx}{0.7\textwidth}{X} 他要像牧人牧養自己的羊群，用膀臂聚集羔羊，抱在胸懷，慢慢引導那乳養小羊的。 \end{tabularx} \\ \\ \relax
40:12 & \begin{tabularx}{0.7\textwidth}{X} 誰曾用手心量諸水，用手虎口量蒼天，用升斗盛大地的塵土，用秤稱山嶺，用天平稱岡陵呢？ \end{tabularx} \\ \\ \relax
40:13 & \begin{tabularx}{0.7\textwidth}{X} 誰曾測度耶和華的靈，或作他的謀士指教他呢？ \end{tabularx} \\ \\ \relax
40:14 & \begin{tabularx}{0.7\textwidth}{X} 他與誰商議，誰教導他，以公平的路指示他，將知識傳授與他，又將通達的道指教他呢？ \end{tabularx} \\ \\ \relax
40:15 & \begin{tabularx}{0.7\textwidth}{X} 看哪，列國都像水桶裡的一滴，又如天平上的微塵；看哪，他舉起眾海島，好像舉起極微小之物。 \end{tabularx} \\ \\ \relax
40:16 & \begin{tabularx}{0.7\textwidth}{X} 黎巴嫩不夠當柴燒，其中的走獸也不夠作燔祭。 \end{tabularx} \\ \\ \relax
40:17 & \begin{tabularx}{0.7\textwidth}{X} 列國在他面前如同不存在，在他看來微不足道，只是虛空。 \end{tabularx} \\ \\ \relax
40:18 & \begin{tabularx}{0.7\textwidth}{X} 你們究竟將誰比神，用甚麼形像與他相較呢？ \end{tabularx} \\ \\ \relax
40:19 & \begin{tabularx}{0.7\textwidth}{X} 至於偶像，匠人鑄造它，銀匠用金子包裹它，又為它鑄造銀鏈。 \end{tabularx} \\ \\ \relax
40:20 & \begin{tabularx}{0.7\textwidth}{X} 沒有能力捐獻的人，就挑選不易朽壞的木頭，為自己尋找巧匠，豎立不會倒的偶像。 \end{tabularx} \\ \\ \relax
40:21 & \begin{tabularx}{0.7\textwidth}{X} 你們豈不知道嗎？豈未曾聽見嗎？難道沒有人從起頭就告訴你們嗎？自從地的根基立定，你們豈不明白嗎？ \end{tabularx} \\ \\ \relax
40:22 & \begin{tabularx}{0.7\textwidth}{X} 神坐在地的穹窿之上，地上的居民有如蚱蜢。他鋪張穹蒼如幔子，展開諸天如可住的帳棚。 \end{tabularx} \\ \\ \relax
40:23 & \begin{tabularx}{0.7\textwidth}{X} 他使君王歸於虛無，使地上的審判官成為虛空。 \end{tabularx} \\ \\ \relax
40:24 & \begin{tabularx}{0.7\textwidth}{X} 他們剛栽上，剛種好，根也剛扎在地裡，經他一吹，就都枯乾；旋風將他們吹去，像碎秸一樣。 \end{tabularx} \\ \\ \relax
40:25 & \begin{tabularx}{0.7\textwidth}{X} 那聖者說：「你們將誰與我相比，與我相等呢？」 \end{tabularx} \\ \\ \relax
40:26 & \begin{tabularx}{0.7\textwidth}{X} 你們要向上舉目，看是誰創造這萬象，按數目領出它們，一一稱其名，以他的權能和他的大能大力，使它們一個都不缺。 \end{tabularx} \\ \\ \relax
40:27 & \begin{tabularx}{0.7\textwidth}{X} 雅各啊，你為何說，以色列啊，你為何言，「我的道路向耶和華隱藏，我的冤屈神並不查問」？ \end{tabularx} \\ \\ \relax
40:28 & \begin{tabularx}{0.7\textwidth}{X} 你豈不曾知道嗎？你豈未曾聽見嗎？永在的神耶和華，創造地極的主，他不疲乏，也不困倦；他的智慧無法測度。 \end{tabularx} \\ \\ \relax
40:29 & \begin{tabularx}{0.7\textwidth}{X} 疲乏的，他賜能力；軟弱的，他加力量。 \end{tabularx} \\ \\ \relax
40:30 & \begin{tabularx}{0.7\textwidth}{X} 就是年輕人也要疲乏困倦，強壯的也必全然跌倒。 \end{tabularx} \\ \\ \relax
40:31 & \begin{tabularx}{0.7\textwidth}{X} 但那等候耶和華的必重新得力。他們必如鷹展翅上騰；他們奔跑卻不困倦，行走卻不疲乏。 \end{tabularx} \\ \\
[1ex]
\hline
\hline
\end{longtable}
$^{1}$我們一起同心祈禱.
我們親愛的天父上帝.
我們願你的聖靈在我們心裡工作.
藉著今天聖經的說話.
你去教導我們.
我們心裡面都有很多的感受.
很多的經歷.
我們很需要你自己的提醒.
你的安慰.
我們都很需要有鼓勵和盼望.
但我們很需要的是.
從你自己的話語裡面.
得到這些盼望和力量的根據.
讓我們知道我們所信的是實在的.
我們的禱告奉耶穌基督的名.
阿們.
今天很開心來和大家見面.
我不知道今天有師範獻師.
又不知道有聖餐.
所以我預備了講一張.
我想都很長的.
我和周牧師說.
我說今天我會講一場的.
你猜周牧師怎樣回答我呢.
他說你盡情去講.
我就有些.
好,我就盡情去講.
但我不知道大家會不會夠精神呢.
我想大家都很熟悉.
《如英展翅上騰》那節經文.
是不是.
《耳塞書》第40章第31節.
但是我自己最近聽別人講道.
或者自己看聖經的時候.
我就有些新的發現.
我就很想和大家分享.
你要有上文才會有下理.
我就很想大家.
如果你方便的話.
你打開舊約聖經的1058頁.

$^{41}$因為我會一次一次和你們解釋一下.
在舊約聖經的1058頁.
或者你們的手機.
《耳塞書》第40章.
我講道裡面有三個部份.
第一個部份就是說.
以賽亞去做了先鋒.
第二部份就是說.
以賽亞他反問猶太人.
一會兒會講到.
今日讀經的時候是沒有讀到.
因為太長.
整章聖經讀就很辛苦了.
第三部份才是說.
我們等候上帝.
我想我們比較熟悉等候上帝的部份.
特別是《耳塞書》40章31節那段經文.
我們是很喜歡的.
我們首先看一下.
以賽亞作先鋒究竟是什麼.
在《耳塞書》第40章到55章.
這個部份就被稱為以色列的安慰書.
我們要了解一本書裡面.
他大概怎樣分那個結構.
背景是什麼呢.
就是波斯王居魯士.
他將會取代巴比倫的皇帝.
去管治巴勒斯坦地.
就是猶太的地方.
大家都記得.
以色列人曾經被俘虜去巴比倫之後.
可以回去耶路撒冷重建圍牆.
還有重建聖殿.
你都記得有這些歷史.
這段時間就是在說.
這班以色列人在波斯王帝的恩准下.
可以回去耶路撒冷重建城牆和重建聖殿.
當然你讀以色拉記的時候.
是重建整個民族的.
上帝就要先知去傳達他的話.

$^{81}$這件事是很特別的.
上帝要先知說什麼.
他就要說.
但是我要說.
看回那時候的背景.
是有以色列人是不想聽或者不想接受的.
為什麼呢.
就是說你被俘虜了去巴比倫.
現在你有機會回去耶路撒冷.
為什麼你不想回去呢.
你不高興呢.
原因是這樣的.
就是說那些猶太人.
他離開了耶路撒冷的地方.
原來那時候是出現已經有逆權侵佔這件事.
就是你的房子可能被人住了.
而你去了巴比倫那個地方.
人生路不熟.
不過你熬了這麼多年.
你叫定居了.
現在突然間有一個新的呼召回來.
你離開巴比倫.
你回去耶路撒冷.
為神重建聖殿.
你想想當時的以色列人.
不是每個人都這麼高興的.
確實是有人回去的.
但是有些人很矛盾的.
我回去.
你叫我去建聖殿.
但是我自己要建回房子才行.
所以你看回十二小先知書裡面.
哈該書等等.
就說到有些人.
只顧著建你的房子不去建聖殿.
背景就是這樣.
你回去自己都不能住.
我怎麼有空跟你一起建城牆建聖殿.
所以其實當時的以塞夏就面對一個這樣的情況.
所以上帝給他一個很大的任務.

$^{121}$就是鼓勵這班人.
你願意離開一個暫且偷安的地方.
你回去耶路撒冷為神做一些事.
你想想這個其實是不是安慰.
我問你.
安慰書第一樣東西.
原來說這件事.
這班人他不一定覺得是安慰.
他會覺得你在挑戰他.
以及給他一個很大的擔子.
我想我們要看的就是.
以塞夏書第一第四十章第一節是說什麼.
他們的神說我用的是和合本.
你們是用修訂本.
裡面有些東西不同了.
你們要安慰安慰我的百姓.
這裡我也看了希伯來文.
這裡出現安慰這個字.
就和第二節那個字是不一樣的.
第二節說你要對耶路撒冷說安慰的話.
又向他們宣告的話.
他爭戰的日子已經滿了.
他的罪孽赦免了.
他為自己的一切罪.
還說手中加倍受到懲罰.
安慰這個字是另一個希伯來文的.
但無論是怎樣.
當你看一段聖經的時候.
他經常出現的那個字.
你說他重要嗎.
重要嘛.
第一節第二節都有安慰.
雖然他兩個字是不同.
但很明顯第四十章開始是說安慰.
上帝要先宣告你爭戰的日子滿了.
你的罪孽赦免了.
他為自己一切的罪.
還說手中接受了加倍的懲罰.
然後我想大家留心.
第四十章的第二節.

$^{161}$他說要對耶路撒冷說安慰的話.
說話就是這樣說話.
聲音不用太大.
又向他宣告說.
宣告也是這樣.
好像我現在這樣說話.
但我為什麼這樣說呢.
我想大家留心.
到你第三節的時候.
發現有一個字很特別.
是喊叫.
你想一下.
宣告或者報告.
都是很清楚.
不會特意大聲的.
但到了第三節.
我們一起看看.
他說開始有人聲喊著說.
作者很有趣.
聽到有人大聲喊.
喊什麼呢.
在曠野預備耶和華的路.
然後他說.
當你要去預備耶和華的路等等.
這裡你會發現這段經文很熟悉.
在哪裡出現過呢.
哪個人物跟他有關呢.
在新約聖經那裡.
是世若罕.
那就是說你要想一下.
這段經文很有趣.
開頭是說安慰.
然後小聲說話宣告也無所謂.
到了第三節好像轉了焦點.
有人大聲喊.
在曠野大聲喊.
那你就會問.
那第三節和開頭那兩節.
怎麼可以相連呢.
那你要安慰人.

$^{201}$你去曠野大聲喊就行了.
你想一下是不是.
你去曠野大聲喊就能安慰人了嗎.
你會覺得有些問題.
原來你再看下去.
你會發覺這段聖經很特別.
因為上帝叫以賽亞去做的事情.
是你猜也猜不到的.
你不說明的.
我一會兒試著告訴你是什麼事情.
就是說你安慰那些意識的人.
以賽亞你去曠野那裡.
你是不是很不明白呢.
當然你住在錦繡你就會明白.
我不在市中心.
但是你想一下我要去安慰一班人.
最好就是在地鐵站附近.
那裡聚集很多人.
但是為什麼神要以賽亞你去曠野呢.
那裡不一定有很多人住的.
還有他不止這樣.
他要去做什麼事呢.
他要修平上路.
目的是什麼呢.
去幫助那班意識的人.
可以遇見上帝.
這樣就明白多一點了.
上帝要以賽亞做的事情.
我們今天的術語.
抱歉.
就是鋪橋踏路.
曠野裡面有很多沙石.
你要拿走它.
還有你要在沙漠地.
你要做一條路出來.
這樣你就可以明白了.
就是說原來先知.
他去安慰那班意識的人.
不只是我說.
我安慰你啊.

$^{241}$上帝安慰你啊.
你不要這麼煩啊.
不要這麼擔心啊.
不止這樣.
他要拿走那些.
一路阻礙那班意識的人.
走近上帝的雜物.
你想一下.
如果你家裡要弄一條小路.
你也要撿走雜物.
是做那些事情的.
還有你要修平那條路.
不要彎彎曲曲的.
又有沙石這樣.
原來上帝要以賽亞去做什麼呢.
是一個建設的工程.
改變外在的環境.
提供土壤.
去方便上帝的子民.
走近上帝.
安慰就是這個意思.
就是說你要那班意識的人得到安慰.
你就需要有以賽亞這個人物.
做牛做馬.
做這些事情.
到第五節的時候.
以賽亞做了這些事情.
那段經文就說.
耶和華的榮耀會顯現.
凡有血氣的人.
一定一起體健.
這是耶和華他親口說的.
我們要小心一件事.
在以賽亞這件事件裡面.
他做完這些事情.
耶和華的榮耀是顯現的.
但當我們去試試套在我們身上的時候.
就算我們去做了很多這些鋪橋踏路的事情.
不表示上帝就跟著做事.
不一定的.

$^{281}$這個不可以說是等好的.
不過在以賽亞的經歷裡面是這樣的.
但在我們身上可能不是這樣的.
比如有些宣教士.
很多十年在一個地方.
他最大的努力.
最大的貢獻就是讓人知道.
這個世界有一種宗教叫基督教.
可能是沒有什麼人洗禮的.
我們不可以用這樣的角度去看.
那個人的成果.
神顯現他的榮耀.
不是我們去決定.
我們只能夠做到這件事.
我們就是去忠誠去做.
在這裡你記不記得.
師弟約翰曾經.
你記得他的故事.
其實用人的眼光去看.
他都很美得.
他自己的門徒都去跟耶穌.
看一看上帝的羔羊.
拯救世人的主.
那兩個門徒.
可能我就去跟了他.
就是說你越做越縮水.
你明不明白.
是啊.
你知道師弟約翰最後說什麼嗎.
他說他必興旺.
我必衰微.
你明不明白.
原來是這樣的一回事.
原來師弟約翰就是.
將以賽亞那種精神.
實現出來.
上帝就叫以賽亞做這些.
我們俗話叫做.
二打六的功夫.
做牛做馬都只是做節.

$^{321}$好了我們再看到第六節.
有人聲說.
你喊叫吧.
有一個說.
我喊叫什麼呢.
他說凡有血氣的.
盡都如草.
他們的美容都好像野地的花.
這些經文很熟悉吧.
但問題就是又來喊叫了.
但是你想一下處境在哪裡.
在曠野.
大聲叫誰聽.
你會覺得很奇怪.
當然新的修訂本就不是凡有血肉.
不是凡有血氣.
是凡有血肉.
這是絕對正確的翻譯.
我想我跳一跳.
第四章第六節.
他的確應該是易造凡有血肉.
但以賽亞喊叫.
喊叫什麼呢.
他喊叫人是很有限的.
人的光榮都是很有限.
這個是他大聲疾呼的地方.
你要記住.
他大聲疾呼.
即是那些以色列人要被安慰.
是需要以賽亞大聲告訴他.
人是很有限.
很有限.
很有限的.
為什麼是這樣呢.
就不是說你不要生氣.
你不要煩.
你多點仰望神.
好像不是我們想象中的模式.
我們可以說去到第七節.
就要看到了.

$^{361}$虎必枯乾花必吊謝.
耶和華的氣吹在旗上.
百姓成賢是草立.
一會兒你會發覺.
同一節又好像重覆了.
請留心這裡.
是在說上帝吹氣的時候.
那些植物就會生長.
但是我自己看過原文.
其實說上帝的氣吹動.
是更加準確的翻譯.
英文的NASB就說.
The breath of the Lord blows.
即是原來上帝的氣.
即是靈.
他會吹動給人的生命.
上帝是氣.
人只是血肉.
如果沒有上帝的氣.
我們就會怎樣.
我們會死.
你明不明白.
我們不是要黑心過年.
但這個真是事實.
我們第八節了.
他又來了.
草必枯乾花必吊謝.
所以我們神的話.
必永遠立定.
大家要注意有兩樣東西.
第一 上帝的氣.
使人有生命.
第二 上帝的話.
使人的生命立定.
這兩個向度是很重要的.
我們基督徒如果想使靈命.
好像剛才向高處行.
你想加深你對神的認識.
你起碼有兩個方面.
第一 神的靈住在你的生命裡面.

$^{401}$第二 神的話在你的心裡面.
如果只有一樣.
即是兩樣東西都一樣.
是不是功力會差一點.
很自然.
但當你有兩樣.
情況就會很不同.
你就會真的經歷到神.
但最有趣.
去到第九節.
他好像轉了一個題目.
報好信息給石安定.
如果你看現在你們手上.
那個和修本新版本.
報好信息的石安.
現在本來這裡舊的和修本.
是以賽亞報好信息給石安.
但這個新的修訂本.
報好信息的石安.
即是反過來的主角.
這個留待我回去再慢慢研究.
但我講回我的版本.
有趣的地方是.
你要登高山.
起碼我這樣看.
以賽亞要登高山.
已經沒有很多問題.
但我想告訴你.
有問題不是一定差的.
有些人讀聖經.
是讀了很多想知的問題.
你有追尋下去.
你就會有得益.
剛才去哪裡.
去曠野.
去沙漠.
現在上高山.
去上高山幹什麼.
然後就說報好信息給耶路撒冷.
你要極力揚聲.

$^{441}$這又是.
在山上大聲喊.
你明不明白.
在沙漠又大聲喊.
上高山又大聲喊.
還要說.
你揚聲不要懼怕.
你還要對猶大的誠一講.
看看你的神在那裡.
其實講耶路撒冷.
講石安.
講猶大都是那個地方.
我自己看過中文和日本的聖經.
這樣的結構是不清晰的.
看起白萊文.
我就分了四個小句.
第一你要登高山.
第二你報好信息給石安.
你要極力揚聲.
然後第三小句.
報好信息給耶路撒冷.
你揚聲不要懼怕.
第四個小句.
對猶大的誠一講.
看看你的神在那裡.
這個就和剛才講.
人的生命好像花草一樣軟弱.
是很不同的.
而且是要上高山去講.
神在那裡.
我們不用怕.
好 我們下一章.
以賽亞書40章第3到8節.
就講著以賽亞要在沙漠大聲去講.
人真是很渺小.
第49章第40章第9節就說.
上帝要以賽亞在山上講.
上帝在這裡.
大家明不明白我的意思.
你這樣明白.

$^{481}$原來以賽亞在做這些事.
就是在做上帝要安慰他百姓的事.
第10節主要說.
好像大能者那樣臨到.
他的傍臂會為他掌權.
他的賞賜在他那裡.
他的報應在他面前.
然後第10節在這裡講.
上帝是要先知去講.
上帝掌權.
上帝賞賜.
上帝報應.
大家記不記得我剛才講.
他是被擄去了巴比倫.
你明不明白.
他是打書仗的那個國家人.
我們下一章.
第11節不止說上帝在這裡.
不止說上帝有權.
這一節很溫柔.
他的賞牧人牧羊自己的羊群.
用傍臂聚集羊羔抱在懷中.
慢慢引渡那羽羊小羊的.
前面那句是在講什麼.
上帝在這裡.
上帝掌權.
上帝審判.
上帝報應.
第11節說.
上帝在這裡.
是很親切的來安慰我們.
你明白原來.
是以蔡做了這麼多事.
慢慢帶出這個訊息.
上帝掌權.
同時他會親自牧養我們個別的弟兄姊妹.
即你和我都是個別的羊群.
他會聚集羊群.
親切的抱在懷中.
很耐心的引渡我們.

$^{521}$你就明白了.
不只是我們需要聽到安慰的訊息.
而是讓我們真的走近這位.
願意安慰我們的上帝.
我們下一章.
接下來這一段.
我剛才主席沒有讀到.
因為我怕太長.
就是以蔡下她反問猶太人.
這段為何我覺得這麼重要.
第二部分.
她安慰這班猶太人.
她是需要令到他們改變一些思維.
以蔡下書第4章12到26節.
記載以蔡下她反問猶太人.
她反問他們.
令他們開悟.
我們用這個字眼.
領受上帝的安慰.
他們的腦袋不轉.
他們根本領受不到上帝的安慰.
4章12到14節.
重點是.
上帝不需要人指導.
人都不能用人的思想去局限上帝.
好像約伯記所講的一樣.
蔡下書4章12節.
誰曾經用手心去量水.
你用手來量水嗎.
你用手的虎口去量蒼天.
你可不可以這樣量天.
還有你用星斗去盛大地的塵土.
大地這麼多塵土.
你可不可以用星斗.
用壺的殼來盛呢.
你可不可以用瓶去秤山嶺.
你可不可以用天瓶去秤江嶺.
以蔡下去問那班以色列人.
你有沒有辦法比上帝厲害.
他想講這件事.

$^{561}$他反問你為何可以用人的想法去規限上帝呢.
第13節.
誰曾經測度過耶和華的心呢.
或者你做過他的謀士.
指教他 給他意見呢.
很到肉的 你明不明.
然後到以色列去問猶太人.
誰可以去指教上帝.
到第14節.
他和誰去商量.
誰教會他做事呢.
誰將公平的路告訴他呢.
又將知識來訓示他.
將通達的道指教他呢.
即是講一件事.
上帝做事是否需要你去教他.
這樣的意思.
你會發覺這段經文很重要.
我們下一章.
第15至17節.
重點是甚麼呢.
萬民都是虛無的.
我很感恩.
到第17節我看到.
萬民在他面前好像虛無.
被他看為不及虛無.
即是虛無都不夠.
還要是虛無中的虛無.
然後就是乃是虛空.
你會發覺這些東西很高.
首先萬民可以叫All Nations.
即是以前以色列的北角.
被亞述大帝國打敗了.
而南角被巴比倫打敗了.
這些萬民就是在說.
統治以色列的那些外族.
以蔡說這些萬民是虛無的.
然後虛無中的再不及的.
還是虛空.
這段經文我覺得很精彩.

$^{601}$他的意思是.
在這班以色列人的經歷裡面.
他們自己是微塵.
他們會覺得.
亞述,巴比倫很厲害.
但是以蔡說.
萬民在上帝面前是虛無的.
究竟那時的人會怎樣想呢.
古代近東的世界是神話世界.
即是每個民族都有自己的神.
有他的偶像.
你這個民族打贏隔壁的民族.
即是你背後的神打敗了隔壁的神.
原來是這樣想的.
很有趣的.
所以以色列人就會覺得.
我們的神不行.
不夠人打.
而以蔡說不是的.
我們的神是很偉大.
其他民族的神是虛無.
虛無中的虛無還虛空.
你明不明白上帝要去安慰這班以色列人.
是需要令到這班人的思想.
要開通才行.
他不開通是不行的.
所以他說在上帝面前.
這些強大的敵人.
是比虛無更加沒有價值.
所以為何這些經文是虛無.
虛無有些奶事虛空.
其實的意思就是.
很差很差很差很差.
但是他們那時會怎樣想呢.
當然不是.
當然覺得別人厲害.
我們下一章.
講到這裡都未可以.
我們下一章.
未可以這麼快回答到第31節.

$^{641}$慢慢我就會講到給你聽.
以色列人就勸猶太人.
你不要覺得上帝不在才行.
第21到26節.
他說我們的上帝是創造天地.
審判的上帝.
他不只是山神.
不是河神.
不是海神.
很多人的想法是分工的.
你知道我們中國人什麼故事都很厲害.
是有分工的.
我不是研究中國宗教.
但是你發覺.
我好奇問這個問題.
譬如有個人.
他排錯隊.
本來他想問他的兒子.
讀書考試的問題.
但他去了負責管恩怨的偶像那裡.
你猜負責管恩怨的偶像會怎樣呢.
會不會說不好意思.
我們這個部門不處理這些事.
你可能會說.
可不可以幫我internal transfer去隔離.
管學業的那裡.
不好意思 不行的.
我又再問一個問題.
你們這些部門.
你不合作的嗎.
有分工不合作的嗎.
你們有沒有跨部門的聯席會議.
你們的偶像上面有沒有一個大一點的.
你明白我的意思嗎.
其實以前的人的偶像觀念很有趣.
廚房的就不管門口的.
你明白嗎.
那個是我們.
我解釋給你聽.
但很多建東人都是這樣想.

$^{681}$真神不管海.
海的不管山.
所以以斯里蘭卡的領袖.
他的上帝是什麼都管.
跨部門的那個頭.
你明白我的意思嗎.
好了.
第22節說.
神在地球的大圈之上.
當然你用今天的科學思想.
不是圓的.
地上的居民好像蝗蟲.
他鋪張窮倉好像蠻子.
展開豬天.
成為可以住的帳篷.
他用一些形容的方法說.
這個宇宙是上帝做的.
上帝管理大自然現象的.
是有能力駕馭大自然的力量.
很清楚的.
要用他們調整的思想.
我們下一章.
然後上帝是審判的上帝.
強大的力量.
都在上帝面前是虛無的.
第23節說.
他使君王歸於虛無.
使地上的審判官成為虛空.
很重要.
所以我找到一些希伯來文字.
虛空和虛無是兩個很不同的字.
我們下一章.
24節.
他說他們只不過是剛剛種了.
你們有時候種花.
這些這麼偉大的靈人.
他們剛剛下來.
大風一些就吹倒了.
剛才說大風一些.
是給植物有生命.

$^{721}$現在大風一些是會吹散一些東西.
那些是旋風.
很有趣的.
同一樣東西在這段聖經裡面出現.
是有兩個不同的作用.
40章24節是對應40章第7節.
草必枯乾花必凋謝.
因為耶和華的氣吹在旗上.
百姓成賢是草.
第7節是說上帝要吹氣才能生.
這裡說上帝一吹就散了.
你會發覺很有趣.
下一章.
去到等候上帝.
我想你們都等了很久.
想去到這個重點.
其實我也發覺.
為什麼這章聖經要這麼長.
要31節.
簡單一點行不行.
有原因的.
搞了這麼久才可以去到重點.
去到27到31節.
你去等候上帝.
但你在等候上帝之前.
那班以色列人.
是需要有一個正確的態度去等的.
同不同意嗎.
你不是一個正確的態度去等.
你只是亂等.
第27節說.
你為何說我的道路向耶和華引以色列人.
你為何說我的冤屈神並不查文.
開始描入心裡.
以色列人.
為什麼你會說我們的道路很迷糊.
我自己都不知道怎麼走.
為什麼你會說你的冤屈神是看不到.
不理你.
你明白嗎.

$^{761}$他要接受安慰之前.
這些問題是不是要先清理.
你明白了.
到了27節就到了這些.
以塞俠直接就說.
猶太人 以色列人.
你心煩意亂的原因.
是看不見前度.
是不是.
你覺得道路在上帝面前是隱藏的.
同時你覺得上帝不理你那些冤屈.
你那些受苦.
是不是.
這個才是關鍵的地方.
如果以塞俠只是說.
哎呀 別想那麼多了.
當然要看上帝.
上帝要全能.
你覺得這群猶太人或以色列人.
可不可以得到真正的安慰.
或者他可不可以真的走近上帝.
可能未必.
他可能每一次聽完就覺得好一點.
再說 你再說多一點給我聽.
重複開機.
但他根本的問題是他沒有面對.
到了28節.
你豈不知道嗎.
你豈不是聽見嗎.
我們這個永在的神.
耶和華.
他創造天地的主.
他不疲乏.
不困倦.
他的智慧是無法測度.
去到這裡就和以色列人說.
你不要以為上帝睡了覺才行.
你不要以為他太累.
我們上帝是不累的.
他不疲乏的.

$^{801}$他說到這裡是對應我們後面.
我們上帝既然是不疲乏.
不困倦.
所以他可以幫助我們這群.
感到疲乏困倦的人.
如果我們上帝和我一樣疲乏.
他怎樣幫我.
所以開始去到這個位置.
你覺得見到他到位了.
我們下一章.
第29節.
疲乏的.
他賜他能力.
軟弱的.
加他力量.
去到這裡你越來越聽得出安慰的聲音.
我們上帝不是疲乏.
他也不軟弱.
你覺得疲乏.
他會給你力量.
你覺得軟弱.
他會加你的勇氣.
以所謂的形容.
人人都有機會軟弱和疲乏.
譬如30節.
你說那些年輕的.
少年人.
都有疲乏的時候.
他也會疲乏困倦的.
強壯的.
也會全然跌倒.
你這句話聽完之後.
是不是很順.
就是沒有令人歧視.
你明白嗎.
你到這個位置就說.
人人都會覺得軟弱的.
人人都會覺得失敗的.
人人覺得有一種無力感的.
這個是很自然的.

$^{841}$然後.
到第31節.
就是我們經常都很熟悉的.
但是那些等候耶和華的.
他必從心得力.
他們好像鷹一樣.
輾刺上行.
他們奔走卻不困原.
他們行走卻不疲乏.
為什麼我會這麼.
我覺得很想說這麼多東西給大家聽.
因為我們很喜歡.
將這些金句卡給人.
我們自己都收藏起來.
因為這些經文很好.
但是當我看到.
《爾賽亞書》第40章的時候.
我發覺還有很多東西學到.
到最後就要來總結了.
原來就是說.
沒錯.
我們是需要去領受上帝安慰.
但是整段經文.
是有一個很重要的人物.
叫做爾賽亞.
神就要這個.
我用我們今天的術語.
這個傻子爾賽亞.
走去曠野.
走去沙漠.
走去高山.
去講上帝的說話.
然後就去跟以色列人.
去拆解他心裡面一些疑惑.
然後就將上帝.
他不疲乏不困原.
他可以幫你.
這個訊息打進人的心裡面.
然後呼喚以色列人.
你等候上帝.

$^{881}$如果沒有這些人物.
那你怎樣可以真的幫到.
那些以色列的子民.
得到真正的安慰呢.
所以我在這裡就覺得.
這段經文是教我們很重要的道理.
就是說.
我們首先不要覺得自己.
是世界上最慘的那個.
是不是.
我們不要有這種思維.
我們的思維應該對準我們的上帝.
我們的上帝是跟我們同行的.
剛才以色列人說.
上帝在這裡.
上帝掌權.
我們下一章.
有些人他要得到安慰.
是要改了一些他固有的思想.
不要覺得上帝不在這裡.
不要覺得上帝看不見.
有時我們會覺得我們被傷害了.
我們好像上帝不理我們.
我們要認定的是上帝在這裡.
上帝是創造的主.
審判的主.
祂不疲倦.
祂不軟弱的.
祂的智慧超過我們所想的.
還有我們要得到上帝的安慰.
我們真的要有一個很渴慕的心.
有一個正確對上帝的理解.
還有一個很重要的就是.
要有一個很正確的態度.
來等候上帝.
我想我們這幾年都等了很久.
其實很大壓力.
我們現在是盼望.
看何時不用戴口罩.
那種局限有些不舒服.

$^{921}$但是剛才曾牧師說得很好.
通關是我們渴望的.
和上帝打通的關係.
可能是更加根本的.
所以我們除了是.
好像大家心裡願望的那種安慰.
但更加重要的是.
我們有沒有一些的思想.
被現實的世界困住我們.
我們看不到更高的境界呢.
我很希望.
我這麼長氣.
不會令到你們覺得.
整個下午怎麼吃飯.
不要這樣.
還有我希望大家.
真的記住.
神的話 神的氣.
神的話語是真實的.
是生命的糧.
你越去看.
你不會得到這麼表面的東西.
你會得到深層一點的東西.
譬如《耳塞書》第40章.
31節就在這一節已經很寶貴了.
但你從第一節看下去的時候.
你會得到很多東西的.
我就很希望.
我們不只是有人去.
有安慰的訊息.
而是今天我們要記得.
當年的施洗約翰.
他引用以賽亞的說法.
他去修籍主的路.
我們是不是今天都需要多些施洗約翰呢.
對不對.
要多些施洗約翰 多些以賽亞.
你才可以幫到更多的人.
對不對.
現在我們是需要去想一個問題.

$^{961}$神有沒有感動你.
做這個角色呢.
是不是.
我們自己裡面的東西清了沒有呢.
我們有很多東西屈住了.
清了沒有呢.
我們一起同心祈禱.
我們天父上帝.
我們再一次的來到你的面前.
我們懇求你.
讓我們所得到的.
是從聖經裡面而來的那種安慰.
你改變我們的思維.
你讓我們看見.
我們平日看不到的東西.
我們懇求主.
你將這種我們用錢買不到回來的那種福份.
你賜給我們.
讓我們無論是什麼年紀.
我們都會感覺到生命充滿著力量.
我們不疲乏 不困倦.
不是因為我們自己本事.
而是因為你是那位不疲乏 不困倦的神.
讓我們的心跟你有聯繫.
我們等候你 仰望你.
讓我們能夠真正的體驗到.
在一個變幻的世代.
我們最重要的.
不是在地上有什麼保障.
而是我們能夠親近到我們的神.
你是會很重視我們每一隻小羊.
你親自的來將我們抱在懷中.
你去安慰我們 沒讓我們.
你都使用以賽亞.
你都使用施洗約翰.
由你祝福地上的教會.
都能夠好像以賽亞和施洗約翰一樣.
去勇敢的 大聲的宣講.
宣告你是神 你是樹.
我們的祈禱 奉耶穌基督的名 阿們.

\newpage



\section{羅馬書 12:13-21}
\label{sec:ApJnSQORNII}
\textbf{2023-01-22 加恩樓《只要祝福,不要咒詛》羅馬書 12\_13-21 講員:鄧泰康傳道}
\newline
\newline
連結: \href{https://youtube.com/watch?v=ApJnSQORNII}{\texttt{https://youtube.com/watch?v=ApJnSQORNII}} ~~~~ 語音日期: 2023-01-26
\newline
\newline
\hyperref[sec:ezP3uupXjgk]{\small{< < < PREV SERMON < < <}}
~
\hyperref[sec:index_chronic]{\small{[返順時目]}}
~
\hyperref[sec:index_scriptual]{\small{[返順卷目]}}
~
\hyperref[sec:9IcRN7FThks]{\small{> > > NEXT SERMON > > >}}
\newline
\newline
羅馬書 12:13-21
\newline
\begin{longtable}{cl}
\hline
\hline
章節 & 經文 (和合本修訂版)\\
\hline
12:13 & \begin{tabularx}{0.7\textwidth}{X} 聖徒有缺乏，要供給；異鄉客，要殷勤款待。 \end{tabularx} \\ \\ \relax
12:14 & \begin{tabularx}{0.7\textwidth}{X} 要祝福迫害你們的，要祝福，不可詛咒。 \end{tabularx} \\ \\ \relax
12:15 & \begin{tabularx}{0.7\textwidth}{X} 要與喜樂的人同樂；要與哀哭的人同哭。 \end{tabularx} \\ \\ \relax
12:16 & \begin{tabularx}{0.7\textwidth}{X} 要彼此同心，不要心高氣傲，倒要俯就卑微的人。不要自以為聰明。 \end{tabularx} \\ \\ \relax
12:17 & \begin{tabularx}{0.7\textwidth}{X} 不要以惡報惡，眾人以為美的事要留心去做。 \end{tabularx} \\ \\ \relax
12:18 & \begin{tabularx}{0.7\textwidth}{X} 若是可行，總要盡力與眾人和睦。 \end{tabularx} \\ \\ \relax
12:19 & \begin{tabularx}{0.7\textwidth}{X} 各位親愛的，不要自己伸冤，寧可給主的憤怒留地步，因為經上記著：「主說：『伸冤在我，我必報應。』」 \end{tabularx} \\ \\ \relax
12:20 & \begin{tabularx}{0.7\textwidth}{X} 不但如此，「你的仇敵若餓了，就給他吃；若渴了，就給他喝。因為你這樣做，就是把炭火堆在他的頭上。」 \end{tabularx} \\ \\ \relax
12:21 & \begin{tabularx}{0.7\textwidth}{X} 不要被惡所勝，反要以善勝惡。 \end{tabularx} \\ \\
[1ex]
\hline
\hline
\end{longtable}
$^{1}$這段經文其實是很特別的.
因為裡面有很多你不明白的字.
因為在很早期去翻譯.
就算現在來到這個版本也好.
其實也有些大家不太明白的.
譬如在和合本那裡.
要一味的款待.
我們以前經常都說笑.
接待客人你煮一味就行了.
是正經的.
又或者反過來.
眾人以為美的事你要留心去做.
有時這裡有很多不同的解釋.
我想今天嘗試用另一個版本.
我覺得是一個比較現代的版本.
會譯得好一點.
大家會明白多一點.
上帝的子民有需要時.
要樂於幫助他們.
總要熱情好客.
要為逼迫你們的人祝福.
不要就坐他們.
都要祈求上帝祝福他們.
要跟喜樂的人一同喜樂.
跟哭泣的人一同哭泣.
要彼此和睦相處.
不可心高氣傲.
都要樂於跟平凡的人相處.
不要自以為通曉萬事.
絕對不要以惡報惡.
要讓眾人看見你們的行事為人.
是高尚正直的.
也要竭盡全力跟每個人和平共相處.
親愛的朋友.
絕不要自己報仇.
要把他留給上帝的兒奴.
因為聖經上說.
我會報復.
會讓他們得到報應.
相反 你的仇敵餓了.

$^{41}$就給他們吃.
渴了就給他們的水喝.
你這樣做.
就是把羞愧的炭火堆在他們的頭上.
不要讓惡勝過你.
都要以善勝過惡.
這個版本比較容易讓大家了解這段經文的意思.
整個羅馬書裡面的第十二章.
第十一至十一章講了很多論述關於罪的問題.
用很多關於哲學.
用很多不同的角度去講關於罪.
讓當時的人都明白.
無論你是猶太人也好.
你是外邦人也好.
我們都是有罪的.
我們都是在罪的裡面.
我們只可以靠主耶穌能夠成國.
這是保羅在羅馬書裡面.
所以當大家要去研究罪的問題.
你就看羅馬書.
而到了第十二章開始的時候.
你就會發覺其實保羅就開始很實際去教信徒.
我們應該怎樣去過一個信徒的生活.
所有講任何一支詩.
他第一節第二節.
十二章第一節第二節.
他有一個很重要的.
就是要將你全人獻給上帝.
弟兄姊妹.
這個是我們作為基督徒.
我們很重要的起點.
你說你信耶穌.
其實我要再扭正這個字.
不是信這麼簡單.
信就是將你整個人完全奉獻給上帝.
就是說你不再屬於自己.
你是屬於上帝.
這個是起始點.
所以我鼓勵大家.
你不是要花很多的努力.

$^{81}$很大的力氣.
我怎樣扭轉我的生命.
我怎樣去.
我以前有什麼不好的習慣.
我怎樣很努力去戒掉.
我可以告訴大家.
其實很多時候都是徒勞無功.
我以前所認識的.
我以前事奉的教會.
有一位弟兄.
他很坦白和我分享.
他年青又年少.
一路被性這個問題一直困擾著他.
他有他自己一些不好的習慣.
他會看色情的事情.
所以現在很多時候.
他一有這些困惑.
他就會傳訊息給我.
那些感覺又來了.
為我祈禱.
但是在過程裡面.
突然之間我想告訴你.
你真正能夠去改變.
不是靠你自己去忍.
忍是沒有用的.
其實說真的.
你會知道你自己人性裡面是怎樣.
但原來你將你自己完完全全獻給上帝.
你是屬於上帝的了.
當你知道上帝原來在任何的境況裡面.
都看著你.
上帝在任何的境況都與你一起的時候.
你有很多的罪.
有很多的惡.
你不敢去做.
你對上帝有沒有這份足夠的敬畏.
有沒有這份足夠的懼怕.
才能夠讓你可以得著.
所以保羅一開始的時候.
他就告訴你起點是獻.

$^{121}$而接著到三節八節.
他就說了很多關於.
第二個就是請你盡力的.
去將上帝所賜給你的恩賜.
去發揮出來去祝福其他人.
你會發覺他一直在跟大家說.
我不知道大家有沒有留意那個流程.
你首先知道自己屬於上帝.
接著你努力去做好的事.
去將上帝今天所給你的才能裡面.
去發揮出來.
我不知道大家有沒有發覺.
很多時候我們人去困擾的時候.
是因為很多時候我們都將我們的焦點.
不是放在我們現在可以做的事上.
而是有很多我們不開心的事上.
有很多人怎樣對我的事上.
當你聖經能夠消化你明白的時候.
你能夠真的跟著去做.
我想告訴大家.
很多困難你都可以去渡過.
在當我們面對過去幾年疫情當中.
我們有什麼可以做.
我問大家.
我們除了戴口罩.
我們就是可以洗好手.
如果你天天都看著數字.
天天都看著這些情況.
你不會出來的.
但是如果你能夠選擇.
在這個疫情裡面.
原來我還有什麼可以做.
你就會發覺原來我還可以去派口罩.
我還可以去派酒精.
我還可以去幫那些單身.
獨居的老人家.
我可以去買飯給他們.
我可以協助他們.
我可以關心他們.
是不是.

$^{161}$你整個生活就會不一樣.
所以保羅比這個次序是有關係的.
而到了今天跟大家講第九節.
這裡開始的時候.
其實他就講到我們怎樣去處理.
對別人對我.
特別是一些傷害我們的人.
我們對他的態度是應該怎樣.
我可以告訴大家.
其實這個是最不容易的.
因為人受傷害的感覺是很真實的.
是很實在的.
很多時候是可以困擾我們一生.
有很多人放不下的時候.
他困擾我們一生.
你發覺這段經文當中.
他一開始第一句.
我拿出來的是什麼呢.
上帝的子民有需要就要樂於幫助他們.
你總要熱情好客.
要為逼迫你們的人祝福.
不要就坐他們.
倒要祈求上帝祝福他們.
你會發覺.
我們保羅教大家的就是.
我們帶著善意去和其他人相處.
無論那個人是你認識的.
無論那個人是不是得罪你的.
他在這裡的時候.
他沒有說對你好的上帝子民.
你就熱情款待.
不是的.
他這裡是說.
上帝的子民有需要的時候.
你就幫助他.
他這裡沒有分.
他在說的是.
上帝他所去看顧的.
特別是當時的世界.
所有去追隨上帝的人都被迫迫.

$^{201}$他們一定要互相去支援.
所以他特別提到.
就是你要去幫助他們.
你要帶著你這份熱情.
當你很專心去做這些好事的時候.
你就會發覺.
你的正能量自然會出來.
而接著對於那些懷有敵意的人.
上帝用了一個.
就是你去為他求上帝的祝福.
很難的.
他得罪我.
我要叫上帝祝福他嗎?.
但是弟兄姊妹.
這個就是以善勝惡最好的詮釋.
這段經文裡面不是就這樣完結了.
所以剛才我說.
我鼓勵大家多看幾本譯本.
在舊的譯本裡面你看不到.
但在剛才我給大家的譯本裡面.
在這幾句話當中.
全部是用分號.
中文的分號.
就是說它們全部有關.
要為迫迫你們的人祝福.
不要就坐他們.
禱要祈求上帝祝福他們.
分號.
要跟喜樂的人一同喜樂.
要跟哭泣的人一同哭泣.
分號.
要彼此和睦相處.
分號.
不可心裡高傲.
禱要樂於跟平凡的人相處.
分號.
不要自以為通曉萬事.
很有趣的.
我們今天原來做的每一樣事.
我經常都鼓勵弟兄姊妹.

$^{241}$教弟兄姊妹.
我們不是因為聖經告訴你有這個吩咐.
所以我就去做.
你很難的.
你就會受到很多信仰的壓迫.
因為我心裡面明明還在生氣.
你叫我不要生氣嗎.
但這裡其實它是想告訴你.
你不要就坐他們.
祝福他們.
但在裡面.
它突然跳下去.
你要跟喜樂的人一同喜樂.
跟哭泣的人一同哭泣.
什麼意思.
同理心.
你怎樣不生氣這些人.
你有同理心.
很多時候你會發覺.
那個人去傷害你.
原來他本身都可能是一個受過傷害.
或者正在受傷害的人.
家長花園.
很多人家裡都有養狗.
大家有沒有帶狗出去散步.
你有沒有發覺.
吠得最大聲的是什麼狗.
小狗.
我最初沒有留意.
我不懂的.
是養狗的人告訴我的.
小狗.
吠得最大聲的.
為什麼.
大家知不知道.
獸醫說.
寵物有關的人.
專家說.
因為牠小.
牠先去控制你.

$^{281}$你不要走過來.
你會發覺一件事很有趣.
原來當一個人.
為什麼我會走去傷害其他人.
因為他本身的心裡.
可能他都受過傷害.
或者正在受傷害.
你有沒有嘗試了解他背後.
他傷害了我的目的.
你有沒有看出這些人背後的可憐.
其實你會發覺.
原來當你明白的時候.
唯有.
當你這樣做的時候.
你和他的關係,態度.
才會有轉變.
而他才會轉變.
當這個人的心裡.
生命裡會轉變.
這個的冤仇才會化解.
互相攻擊.
永遠只會越來越大.
我很記得.
在早幾年的時候.
當整個香港,整個環境.
連教會都不斷在說.
什麼顏色是什麼顏色的時候.
可能在一家人裡.
很多時候我們坐在一起.
都會有不同顏色.
大家有很多的衝突.
很多的不開心.
而當時我的處理方法是怎樣呢.
我就會跟他們說.
你們的日常生活只有說顏色.
你沒有其他東西要說.
我說其實你們有沒有了解過.
為什麼他在那件事當中.
他背後會這麼生氣.
或者他背後會這麼不開心.

$^{321}$你們有沒有了解過.
不是說對和錯.
不是的.
你所愛的這個人.
他因著一些事情.
他心裡被困擾.
他晚上睡不著.
有些事情他不明白.
我們有沒有從這個角度去看我們身邊.
好像得罪我們的人.
我們有沒有呢.
我們有沒有用這個方法.
聖經所教的方法去做.
所以他最後的結尾是什麼呢.
他就說不要自以為通曉萬事.
你不要以為你知道所有的東西.
跟人相處.
跟這些人相處.
就是謙卑下來.
所以你會發覺保羅很有趣.
保羅其實是一個知道很多東西的人.
但他一直教我們的時候.
他就說謙卑下來.
有很多東西你不是知道一切.
因為他就是在上帝面前.
他有十足的肯定.
我很微小.
儘管我知道多少東西.
我都很微小.
我都不足以成為其他人很大的祝福.
唯有上帝才是祝福.
所以最後的經文就說.
絕對不要以惡報惡.
要讓眾人看見你們的行事為人是高尚正直.
也要盡力跟每個人和平共處.
所以你就會發覺.
讓所有看見其他人的行為.
因為基督徒其中一個很重要的地方.
就是見證.
你要記住見證.

$^{361}$我不是說.
我想大家搞清楚一些觀念.
因為我是基督徒.
所以我就好像要做一個外表給別人看.
不是的弟兄姊妹.
你是屬於神的.
你一舉一動其實是在見證上帝.
我們說一說.
我們一生其中一個就是.
讓所有人去榮耀上帝.
記不記得.
你所做的一切的事.
我們讓更多人去榮耀上帝.
還是讓人去討厭上帝.
不喜歡上帝.
我們所做的任何行為.
說的每一句話.
我們是在做哪一件事.
所以才說.
盡力去跟這些惡人和平共處.
當然.
坦白說在這個世代.
如果你問我.
有時候我們見到一些很不平的事情.
我們在YouTube上.
見到一些事情.
那個人又沒有禮貌.
又沒有什麼.
我們一起批判他們.
然後見到他有些最後的惡果.
我們心裡面很爽快.
我要承認那種感覺.
這個人活該.
但不是的.
弟兄姊妹.
我想告訴大家.
這種感覺.
你慢慢延伸出去的時候.
是會變成心理不平衡.
其實我們就變成一樣.

$^{401}$我們自己站在一個道德高地上.
去看著其他人.
所以保羅在這裡勸我們的時候.
就是將你這個審判的權柄.
交回給上帝.
怎樣能夠令你不要那麼意氣填空.
因為你不是上帝.
你會知道上帝告訴過你.
他應許過告訴你.
他是知道的.
你要相信上帝.
他會行公義的.
當你有這份相信.
你的心才能夠放下.
之後你才能做到下一個行為.
就是怎樣.
你不單止不要去討厭他.
他餓的時候.
你要給他食物.
他口渴的時候.
你要給他水喝.
唯有當你的心能夠放下.
你才能夠做到這個行為.
但原來這個行為的背後是什麼.
我很喜歡這個譯本.
將羞愧的炭火堆在他的頭上.
有很多人可能一聽到這裡.
就覺得我以前試過了.
我都不知羞恥.
你對他好的時候.
他也是這樣.
之後也是這樣.
請弟兄姊妹相信聖經的說話.
人的心裡是會有良知的.
當他一直對你很差很惡的時候.
但你仍然可以對他行善.
對他能夠做一些關心的行動.
羞愧的炭火是會在他頭上.
會令他羞愧.
這個我不是作的.

$^{441}$如果大家有興趣.
所以沒有什麼特別的東西.
都可以參加別人的浸禮見證.
聽聽.
有很多人真的這樣.
有些人以前不知為什麼.
他很討厭基督徒.
見到他都會笑你.
但他的朋友很好.
每次他都會關心他.
他不理他.
當沒聽到.
但他說就是因為這樣.
我都不明白你們這班人為什麼這麼奇怪.
我明明對你這麼差.
但為什麼你還可以關心我.
為什麼你還可以這麼認真地對我.
他說你是不是傻.
但最後當他接觸的時候.
他就因為這份感動.
所以信了耶穌.
所以弟兄姊妹.
我們在神面前不怕靠著上帝.
去做一個在世界體位的傻瓜.
不怕.
因為上帝所說的話.
才是真實又活的.
所以最後一句話就是.
不要讓惡勝過你.
你要用上帝的善去勝過惡.
弟兄姊妹記住.
是用上帝的善去勝過惡.
我今天能夠去原諒你.
是因為有上帝.
不是我大方.
不是我大度.
按人性說一句.
我一樣是很怒氣的人.
我一樣是會見到不公義.
我一樣是很不開心的人.

$^{481}$但是在上帝自己的愛裡.
我們看到自己的軟弱.
所以你會發覺.
其實看這段經文最重要的是看保羅.
因為保羅作為一個.
一直以來受人傷害和逼迫的人.
但你會發覺他的教導.
從來都不是從人的血氣.
或短暫的勝利出發.
從來都不是.
因為我們已經屬於上帝.
而上帝是不會這樣做的.
他不會去即時審判.
惡待做錯事的人.
當你看回聖經歷史就知道.
上帝寧願透過先知.
幾百年不斷提醒.
善勸當時的皇帝和他們的人民.
任何行動都沒有出現過.
只不過是過了幾百年.
沒有辦法.
他最後才將審判淋到當中.
但我想大家留意.
當上帝出手的時候.
你會發覺結果是很嚴重的.
他可以亡國.
他可以毀滅.
上帝是那位說有就有的上帝.
所以弟兄姊妹.
當我們看這段經文的時候.
我很盼望.
我們經常說我們新年就拜年.
去探望親情.
我不是說大家.
我自己也經歷過.
我記得我小時候住在唐樓.
唐樓裡面.
樓上樓下我住在六樓.
對面隔壁的很熟悉.
叫陳師奶李師奶.

$^{521}$同一個年代的.
我們叫師奶.
師奶平時不知道為什麼.
她真的以為自己是我媽媽.
走過的時候罵你.
教你.
然後經常戳背.
我們年輕的時候拍拖.
永遠不是我們說的.
是師奶告訴媽媽.
剛才我在下面碰到你兒子.
看到一個女孩.
很憎恨她.
但新年的時候要怎樣.
要恭喜發財.
見到她的時候要拜年.
那時候我感覺.
我明明很憎恨她.
我想說很多時候我們.
人與人之間總有很多嫌隙.
在我們的親戚裡.
在我們的朋友.
在我們家人當中.
很多時候都是.
在我心裡我不喜歡她.
但問題是沒辦法.
一家人回父母家吃飯.
對著.
我形容的只是.
暫時你掩蓋了她.
你沒有處理過.
但我盼望弟兄姊妹.
其實我們心裡是否真的開心.
我們心裡是否真的放下了.
不一定.
但求上帝用她的話.
去幫助我們.
好像剛才所說.
去嘗試了解她.
我用一個例子.

$^{561}$我相信我大哥也不會介意.
我大哥.
我們以前從小到大.
兄弟姊妹都不喜歡他.
因為我大哥是一個.
很自私的人.
他買很多東西後.
他會在自己的衣櫃.
用鎖鎖住.
你知道我們小時候很頑皮.
我們那時候很喜歡.
你越不給我們就越要拿.
越要去找.
但後來我就很明白.
為什麼他會這樣.
當後來有機會再談的時候.
他有機會說.
回到自己的心底話.
他說他以前是第一個.
他是萬千寵愛在一身.
真的.
我大哥特別疼父母.
我們兄弟姊妹都看到.
他說我看到你們.
弟弟出生.
第三個弟弟出生.
第四個弟弟出生.
妹妹出生.
突然間那種很不安全感.
我就要將我自己.
這些東西納起來.
我要納回來.
這些是我的.
你不要碰 你不要搞.
原來是這樣.
突然間原來每一個人背後.
當你了解了他之後.
你會沒有那麼憤怒他.
你會沒有那麼憎恨他.
但當然.

$^{601}$今天我哥哥也相信耶穌.
他今天已經完全不是一個自私的人.
他很大方.
他很多東西願意和我們分享.
他很好的.
所以唯有在上帝的裡面.
一切的事我們才能夠解決.
求神幫助 報仇我們.
我們就同心祈禱.
差不多我們要去多謝贊美你.
主人.
我們知道在神的面前.
才能夠為我們將一切所有的關係修補.
是你教我們去祝福.
不要走祖.
主人我們知道地上有很多.
讓我們心裡不開心的事.
但同樣的.
你告訴我們.
主人原來這個世界上.
所有人都是受傷的人.
在神的面前的時候.
我們很需要上帝你去醫治.
所以我求你幫助我們.
我們不單止不會走祖.
甚至我們願意將傷害我們的人.
帶到上帝的面前.
去得到醫治.
我們求你親自的保守幫助.
求你親自的保守.
祈禱奉教主耶穌.
得勝名字祈求.
阿們.
\newpage



\section{列王記上}
\label{sec:9IcRN7FThks}
\textbf{2023-01-22 新春崇拜《蒙福秘訣》列王記上 3: 3-9 講員:陳一華牧師}
\newline
\newline
連結: \href{https://youtube.com/watch?v=9IcRN7FThks}{\texttt{https://youtube.com/watch?v=9IcRN7FThks}} ~~~~ 語音日期: 2023-01-26
\newline
\newline
\hyperref[sec:ApJnSQORNII]{\small{< < < PREV SERMON < < <}}
~
\hyperref[sec:index_chronic]{\small{[返順時目]}}
~
\hyperref[sec:index_scriptual]{\small{[返順卷目]}}
~
\hyperref[sec:S5k9W4TGOeU]{\small{> > > NEXT SERMON > > >}}
\newline
\newline
$^{1}$各位大家好 各位大姐妹 早安.
今天是大年初一.
我看到很多人的面容都很寬恩.
雖然戴了口罩.
但是你們的眼神好像有些光彩.
那我們趁這個時候.
大家和隔壁的朋友.
來恭祝他們新年快樂.
大家打個招呼.
新年快樂.
在這裡我要向大家恭祝你們.
新年蒙福.
身體健康 精神奕奕.
就好像套年一樣.
套氣什麼 揚眉.
有一個優良的傳統.
就是我們逢新春年初一.
我們都有敬拜的一個很好的行動.
提醒我們和人拜年之前.
我們先向神來什麼 敬拜.
因為我們要把握最佳的時間.
最好的光陰來親近神.
去敬拜上帝.
完了我們再去拜年.
就綽綽有餘了.
中國人有一句話叫做.
服從天國 大家認同嗎.
原來真福是真的從上天降下來.
證明粵曲對唱粵曲的時候.
都有一句說話叫做.
因澤門庭.
因福臨到你的門口 從上而下.
我就很快想起亞國書這樣說.
亞國書說什麼呢.
他說一切美好的賞賜.
和完備的恩賜都是從上頭來的.
有印象嗎.
都是從重光之父那裡降下來的.
這裡特別用了一個降下來.
就是說服從天國.

$^{41}$所以弟兄姊妹們.
我們有一位是天上的爸爸.
他是真福的源頭.
我們做他的子女.
有沒有福氣.
大有福氣了.
為什麼呢.
因為我們享受的真福.
是真的從重光之父.
重福的源頭賜給我們的.
所以我們在新春一開始的時候.
去敬拜上帝.
是非常非常之有意義的.
是不是.
當我們想起這個福的時候.
今早牧師就和大家分享一個題目.
都離不開那個福了.
是一個什麼呢.
蒙福的途徑.
大家想不想知道.
蒙福的途徑是怎樣可以得到的呢.
又或者上帝.
我們天上的父親.
怎樣將一個真福.
一個重大的福賜給我們呢.
其實真是有跡可尋的.
大家今早聽完主席.
來讀了一段經文.
所羅門王這個名字.
你們都很熟悉的.
在以色列人的歷史當中.
他們感覺到這麼多個皇朝裡面.
最興旺的皇朝.
莫過於是所羅門.
當然他爸爸大衛的皇朝.
一樣都是這麼鼎盛的.
在所羅門王做皇的時候.
國家富庶.
興盛.
國泰民安.

$^{81}$得到的福大不大.
非常之大.
但是他蒙福的原因在哪裡呢.
途徑在哪裡算呢.
今早我們不妨在經文當中.
來看一看.
這些經文剛才都讀了.
第三節告訴我們.
所羅門之所以蒙福.
是因為他遵行上帝的律例.
即是現在來說.
我們聽話.
聽神的話.
遵行神的話語.
這個遵行神的話語.
和愛是有些關係的.
我不知道你們會不會感覺到.
和愛的表現是有密切關係.
即是說如果你愛一個人.
你重不重視他的說話.
重視的.
如果有一天.
媽媽跟你說.
兒子.
現在通關了.
你拿駕陪我回家鄉走走吧.
兒子聽完之後.
如果他愛他媽媽.
他考不考慮媽媽這個要求.
他會的.
為什麼呢.
因為他重視他媽媽的提議.
又或者有一天.
你太太跟你說.
老公現在疫情放緩了.
不如我們去日本旅行吧.
Happy嗎.
Happy.
可能你們都在計劃.
還未定.

$^{121}$或者去完回來都未定.
如果你愛你太太.
你考不考慮她的提議.
你會考慮的.
你開始上網.
看看哪些機票是最便宜的.
為什麼.
因為你重視她.
你愛她.
你就重視她的說話.
所以各位頂尖的節目.
今天我們為什麼會成為蒙福的兒女.
是因為我們對上帝的話語.
是極度的重視.
人的說話都有時候很有道理.
很有文腦.
不過當比起神的話語.
始終差很大.
今天我們有上帝的話.
透過聖經啟迪我們的時候.
提醒我們的時候.
各位頂尖的節目.
我們有沒有像所羅門一樣.
去遵行上帝的話語.
如果你發覺過去一年.
你都很聽上帝的話.
很遵行神的命令.
我相信你過去一年.
是經歷很多神的恩典.
為什麼.
因為遵行神命令.
是蒙福的途徑.
這是第一步.
我們再看第二步.
為什麼.
所羅門蒙福呢.
是因為他樂意去親近神.
樂意親近神.
這個行動.
和記掛的表現.

$^{161}$或者和掛念的表現.
是掛勾的.
你認同嗎.
就是說當你很記掛一個人的時候.
你想不想去接近他.
你很想去接近他.
最近中港兩地通關了.
為什麼那些人排隊排到半夜.
都要做核酸.
因為做完之後.
他就回鄉下.
和家人團聚慶祝新年.
所以這個行動.
告訴我們.
他記不記掛家人.
很記掛.
千辛萬苦都要去做.
做好之後.
他就可以回到自己的家.
親近回.
接近回他的家人.
所以樂意親近神.
和記掛有沒有密切關係.
有密切關係的.
鬼頂子母們.
所羅門為什麼蒙福.
因為他記掛著上帝.
他做了王.
什麼都有齊了.
不過他的心仍然很記掛著上帝.
很想去親近神.
就算當時候聖殿未建好.
上帝沒有一個很明顯的居所.
如果你看回歷史.
原來當時候.
以色列人要去親近神.
要去獻祭.
是要去到一些幽壇的地方.
就是說讀經有沒有出現這個字.
幽壇.

$^{201}$什麼東東來的呢.
幾年前我和大家敢孫的弟兄姊妹.
去以色列旅行的.
幸好當時去得早.
遲些就去不了.
弟兄姊妹們.
我們去到那個地方的時候.
真的見到一些古蹟.
是幽壇的古蹟.
我想告訴大家.
幽壇一般來說.
都是建在半山的地方.
高處的地方.
半山高處的地方.
有一大塊的平原.
以色列人就在那裡建造幽壇.
為什麼呢.
因為他們要在那裡獻祭.
來還他們的怨.
無論是凡祭,平安祭,贖獻祭.
他們有什麼呢.
要做這個行動.
對所羅蒙王來說.
他為了要親近神.
不辭勞苦.
由他住的地方.
去到這個機片的幽壇的地方.
都是長途跋涉的.
但是他照去.
還要帶很多生畜去.
你剛才讀的聖經有沒有看到.
燒的生畜.
獻祭的生畜.
都超過一千的.
整個行動都是.
可以怎麼說呢.
勞師動眾的.
大家認同嗎.
不過.
所羅蒙有沒有覺得很煩.

$^{241}$在家附近自己獻就行了.
有沒有這樣做.
沒有.
不辛苦.
他都願意付出這個行動.
去親近上帝去獻祭.
因為當時候以色列人.
沒有聖殿.
各位電視節目們.
你有沒有很渴求.
很想去親近上帝.
因為你記掛著他.
是吧.
你在家的時候.
記掛著他.
你在上班工作的時候.
你記掛著他.
你和朋友一起有活動.
去旅行的時候.
都記掛著上帝.
就算我玩得很開心.
萬水千山.
風景如畫.
我都記掛著上帝.
各位電視節目.
如果你有這樣的心.
你去到任何的地方.
你都可以來到親近上帝.
認同嗎.
所以.
新年大家都知道.
大家都有些旅行計劃.
請你記得.
旅行不單止自己心情.
放鬆.
很享受之外.
你都要記掛著上帝.
不要忘記親近上帝.
無論你的讀經.
無論你的祈禱.

$^{281}$無論你的默想.
無論你唱的詩歌.
都是表達什麼.
你記掛著上帝.
各位電視節目.
這是一個什麼.
蒙福的途徑.
為什麼很多時候.
基督徒得不到上帝的祝福.
是因為他忘記了上帝.
在他日常生活當中.
他忙得來.
忘記了上帝.
忘記了天上的爸爸.
所以他的生活工作當中.
很多時候.
眼下跳不出煩惱的問題.
好像沒有什麼力量.
人很累.
很耗盡.
今早我們來到新春的崇拜.
大家記得.
好像所羅門一樣.
樂意親近神.
無論付出多大的代價.
在所不辭.
我們再看第三.
第三個蒙福的途徑是什麼.
謙卑在神前.
當上帝和所羅門說.
你可以向我求了.
他怎樣向上帝求.
他就這樣向上帝求.
他說上帝.
你是知道的.
我不過是一個幼少的.
幼少的是和修本的譯本.
如果你看回我們舊陣史和合本的譯本.
那裡說.
我是幼童.

$^{321}$那些頂智慕門學家覺得很奇怪.
為什麼這麼小就做王呢.
實際上如果你看回聖經的記載.
所羅門做王的時候.
都是大人的.
都是成人的.
接著他爸爸做王.
他說自己很幼少幼童.
其實有原因.
因為他感覺到自己在上帝的面前.
是多麼的微小.
再加上他自己真的膚淺.
真的初出茅廬.
沒有什麼人生的歷練.
沒有什麼社會的經驗.
做一個王.
面對這樣的情景的時候.
他看到自己的需要.
所以他謙卑在上帝的面前.
各位頂智慕門.
謙卑在上帝面前也是一個蒙福的途徑.
為什麼呢.
因為聖經告訴我們.
上帝阻擋驕傲的人.
次因級謙卑的人.
有記憶嗎.
這句是段經文.
有沒有印象.
所以當我們都在上帝的面前.
有一顆謙遜謙卑的心.
感覺到自己很幼嫩.
雖然自己很有經驗.
但始終面對這個變幻莫測的世界.
面對種種複雜的人際關係.
要去處理的時候.
我們都謙卑在上帝的面前.
承認什麼.
神啊.
我是很膚淺的.
我很多事情都不懂.

$^{361}$求你幫助我.
我相信我們這樣的禱告.
上帝是喜悅的.
上帝會聆聽我們的禱告.
祂會特別加能賜力給我們.
因為我們謙卑在上帝的面前.
我們再看第四個原因.
第四個秘訣.
為什麼上帝賜福給他.
第八節就告訴我們.
上帝啊.
我今天接受我父親造禍.
你看看.
在我面前的百姓是什麼.
數之不盡.
在我面前的人.
我要管治的人.
多得不可勝數.
上帝.
這是我的真實處境.
各位頂字幕們.
有時候我們都要看看自己的真實處境.
當我們回望自己真實處境的時候.
你會發覺自己.
真是很大很大的挑戰.
或者我們不用困難.
用挑戰這個字眼.
真是很大的.
怎樣可以去管治這麼多的人呢.
怎樣令到國家裡面減少了紛爭.
怎樣令到國家裡面的人.
遇到困難的時候.
我可以跟他.
能夠排解他的問題.
所以上帝啊.
我看到我自己真實的處境.
所以我承認我是力猶不逮.
今天我們都很需要認識自己.
大家都認同嗎.
有很多人他失敗.

$^{401}$不是失敗在別人的手上.
不是敵人強勁過他.
是他自己.
看自己是看得太過的重要了.
或者看自己看得太過大了.
他覺得自己是天下無敵了.
所以有些人真是這樣看自己的.
他覺得自己什麼都可以的時候.
什麼都可以的時候.
他可以目空一切.
所以心很紅.
心紅的時候就盲目了.
最近有一間餅店.
它全盛的時期.
都開了很多間餅店.
不過最近他都要面對破產.
弟子們.
有很多時候我們的心紅.
我們以為自己大有能力的時候.
以為自己什麼都可以的時候.
沒有看回自己真實的處境.
我們真是會碰壁的.
所以我們怎樣可以得到上帝的賜福呢.
我們在上帝的面前都承認自己是力猶不逮.
需要上帝加能賜力給我們.
需要上帝幫助我們.
我們看第五最後一點.
第五所羅門他蒙上帝賜福.
你發覺他的祈禱裡面.
向神的呼求裡面.
他的重心不是為自己去求.
你有沒有留意到.
他的重心是以百姓為本.
上帝我真是很需要你給我什麼.
有智慧.
因為我管治這麼多人.
我要跟他們排難解災.
我要跟他們排解紛爭.
我怎樣可以做到這麼多事呢.
尤其是有時候有些紛爭是家庭式的.

$^{441}$還更加難.
正所謂清官都難審家庭事.
所以審死官.
頂子們.
當所羅門求的時候.
他說上帝我求你給我有智慧.
明辨是非的能力.
這個非常之好.
所以上帝很欣賞他.
因為他發覺原來上帝現在叫他做王.
不是叫他威.
不是要他炫耀自己.
是要他好好地將上帝交託給他的任務.
做得好.
完成得好.
各位頂子們.
你有沒有留意到.
所以有時候上帝賜福給你.
給你生而興隆.
給你什麼都很順利.
是好的.
你記得上帝給你這麼多祝福.
這麼多好處的時候.
其實這個是一個任務.
這個任務就是你怎樣將上帝交託給你的.
即是能夠忠心地去完成.
能夠有效榮神益人地去完成.
所以頂子們.
你現在明白.
所羅門祈禱不是為自己去祈禱.
他是以百姓為本.
他是以上帝的任務為主.
上帝竟然託付這個任務給我.
正所謂有權利都要有義務.
你做得皇帝.
但你有沒有義務.
你要管治好你的國家.
你要讓你的國家不要出事.
你不要讓你的百姓離氣神.
歸向偶像.

$^{481}$這個任務重不重要.
很重要.
所以各位頂子們.
你現在明白.
所羅門為什麼會這樣祈禱.
結果上帝是否欣賞他的禱告.
我知道你們都拿到答案了.
上帝非常欣賞他的禱告.
你向我求的時候.
不是求長壽.
不是求富貴.
不是求打贏敵人的戰爭.
你求的是智慧.
所以上帝說你沒有求的.
我都給你.
送給你.
因為你求得好.
今天我們在新春年初一.
來到上帝面前敬拜的時候.
不得不說各位弟兄姊妹們.
我們都學習所羅門的榜樣.
知道蒙福有它的途徑.
當然所羅門也是人.
他晚年也做錯了事.
但我們引以為鑒.
但他好的地方.
我們應不應該學習.
都要學習的.
所以求主賜福我們.
感孫家眾弟兄姊妹.
套年來到了.
我們要有一番新氣象.
靠主來到.
你們有一番新的挑戰.
新氣象.
叫我們大家都能夠靠主得勝.
願上帝賜福給所有弟兄姊妹.
多謝大家.
\newpage



\section{馬可福音 7:1-23}
\label{sec:S5k9W4TGOeU}
\textbf{2020-10-03 晚堂《用嘴來尊敬神的偽善者》馬可福音 7\_1-23 講員:高銘謙牧師}
\newline
\newline
連結: \href{https://youtube.com/watch?v=S5k9W4TGOeU}{\texttt{https://youtube.com/watch?v=S5k9W4TGOeU}} ~~~~ 語音日期: 2023-01-27
\newline
\newline
\hyperref[sec:9IcRN7FThks]{\small{< < < PREV SERMON < < <}}
~
\hyperref[sec:index_chronic]{\small{[返順時目]}}
~
\hyperref[sec:index_scriptual]{\small{[返順卷目]}}
~
\hyperref[sec:mXSWAPPoVUc]{\small{> > > NEXT SERMON > > >}}
\newline
\newline
馬可福音 7:1-23
\newline
\begin{longtable}{cl}
\hline
\hline
章節 & 經文 (和合本修訂版)\\
\hline
7:1 & \begin{tabularx}{0.7\textwidth}{X} 有法利賽人和幾個從耶路撒冷來的文士聚集到耶穌那裡。 \end{tabularx} \\ \\ \relax
7:2 & \begin{tabularx}{0.7\textwidth}{X} 他們曾看見他的門徒中有人用不潔淨的手，就是沒有洗的手吃飯。 \end{tabularx} \\ \\ \relax
7:3 & \begin{tabularx}{0.7\textwidth}{X} 法利賽人和所有的猶太人都拘守古人的傳統，若不按規矩洗手就不吃飯； \end{tabularx} \\ \\ \relax
7:4 & \begin{tabularx}{0.7\textwidth}{X} 從市場來，若不洗淨也不吃飯；他們還拘守好些別的規矩，如洗杯、罐、銅器、床鋪等。 \end{tabularx} \\ \\ \relax
7:5 & \begin{tabularx}{0.7\textwidth}{X} 法利賽人和文士問他說：「你的門徒為甚麼不照古人的傳統，竟然用不潔淨的手吃飯呢？」 \end{tabularx} \\ \\ \relax
7:6 & \begin{tabularx}{0.7\textwidth}{X} 耶穌對他們說：「以賽亞指著你們假冒為善的人所預言的說得好。如經上所記：『這百姓用嘴唇尊敬我，他們的心卻遠離我。 \end{tabularx} \\ \\ \relax
7:7 & \begin{tabularx}{0.7\textwidth}{X} 他們把人的規條當作教義教導人；他們拜我也是枉然。』 \end{tabularx} \\ \\ \relax
7:8 & \begin{tabularx}{0.7\textwidth}{X} 你們是離棄神的誡命，拘守人的傳統。」 \end{tabularx} \\ \\ \relax
7:9 & \begin{tabularx}{0.7\textwidth}{X} 耶穌又說：「你們誠然是廢棄神的誡命，為要守自己的傳統。 \end{tabularx} \\ \\ \relax
7:10 & \begin{tabularx}{0.7\textwidth}{X} 摩西說：『當孝敬父母』；又說：『咒罵父母的，必須處死。』 \end{tabularx} \\ \\ \relax
7:11 & \begin{tabularx}{0.7\textwidth}{X} 你們倒說：『人若對父母說：我所當供奉你的已經作了各耳板』（各耳板就是奉獻的意思）， \end{tabularx} \\ \\ \relax
7:12 & \begin{tabularx}{0.7\textwidth}{X} 你們就容許他不必再奉養父母。 \end{tabularx} \\ \\ \relax
7:13 & \begin{tabularx}{0.7\textwidth}{X} 這就是你們藉著繼承傳統，廢了神的話。你們還做許多這樣的事。」 \end{tabularx} \\ \\ \relax
7:14 & \begin{tabularx}{0.7\textwidth}{X} 耶穌又叫眾人來，對他們說：「你們都要聽我的話，也要明白。 \end{tabularx} \\ \\ \relax
7:15 & \begin{tabularx}{0.7\textwidth}{X} 從外面進去的不能玷污人，惟有從裡面出來的才玷污人。」 \end{tabularx} \\ \\ \relax
7:16 & \begin{tabularx}{0.7\textwidth}{X} 有耳可聽的，就應當聽！ \end{tabularx} \\ \\ \relax
7:17 & \begin{tabularx}{0.7\textwidth}{X} 耶穌離開眾人，進了屋子，門徒就問他這比喻的意思。 \end{tabularx} \\ \\ \relax
7:18 & \begin{tabularx}{0.7\textwidth}{X} 耶穌對他們說：「你們也是這樣不明白嗎？難道你們不了解，凡從外面進去的不能玷污人嗎？ \end{tabularx} \\ \\ \relax
7:19 & \begin{tabularx}{0.7\textwidth}{X} 因為不是進入他的心，而是進入他的肚子，又排入廁所。」（這是說，各樣的食物都是潔淨的。） \end{tabularx} \\ \\ \relax
7:20 & \begin{tabularx}{0.7\textwidth}{X} 耶穌又說：「從人裡面出來的，那才玷污人； \end{tabularx} \\ \\ \relax
7:21 & \begin{tabularx}{0.7\textwidth}{X} 因為從人心裡發出種種惡念，如淫亂、偷盜、兇殺、 \end{tabularx} \\ \\ \relax
7:22 & \begin{tabularx}{0.7\textwidth}{X} 姦淫、貪婪、邪惡、詭詐、淫蕩、嫉妒、毀謗、驕傲、狂妄。 \end{tabularx} \\ \\ \relax
7:23 & \begin{tabularx}{0.7\textwidth}{X} 這一切的惡都是從裡面出來，且能玷污人。」 \end{tabularx} \\ \\
[1ex]
\hline
\hline
\end{longtable}
$^{1}$各位弟兄姊妹晚安 我們一起祈禱.
慈悲人外的天父 你啟示你的聖言.
讓我們在不同的時空 不同的年代.
我們仍然聆聽 直接聆聽到你的教導.
給我們生命有提醒.
求你驅走我們的黑暗.
讓你的話語再次更新潔淨我們的生命.
幫助孩子宣講你的話語清楚.
我們同心領受 禱告 祈奉耶穌基督的名求.
阿們.
很高興一個月前我來這裡送拜.
和大家一起分享.
一個月後我終於見到有人可以坐在這裡.
和大家一起去敬拜 實在很激動.
因為我過往的三個月.
都是對著鏡頭說話的.
但是對著人說話始終都是安心一點.
看到大家一起在這裡敬拜.
代表建築學院再次問候大家.
我們聖學院的學生都陸陸續續進入宿舍.
他們開始進行實體的教導.
仍然有些學生可能有些困難.
所以我們同步都會進行一個直播.
和大家用一些軟件進行課堂.
實在不容易 很多的應變.
我都很佩服每一位教會的弟兄姊妹.
和教睦同工 搞盡腦汁.
在不同的線上線下進行一個的牧養和教導.
實在不是很容易的一件事情.
這個星期很開心可以見到大家.
而今天的講題.
我就承接上一次我們討論到《馬賀福音》.
若然大家還記得我來過當中的時候.
就是《馬賀福音》的第二章.
為什麼是《馬賀福音》呢.
因為我自己的教會有外堂.
都是在講《馬賀福音》.
所以有外堂講完.
同步我又在這裡和大家思想《馬賀福音》.
今天我們就試一下一起看看第七章.

$^{41}$這二十三節的經文.
今天的講題其實就是用嘴.
來尊敬神的偽善者.
如果你有留意.
其實這個講題都頗吊詭的.
你用嘴來尊敬神.
但卻你的本質是一個偽善者.
所以今天主要處理的.
就是那個嘴和偽善者之間的關係.
期望可以成為我們的教導和反思.
先講一下什麼叫偽善者.
偽善者原來來自一個希臘文.
那個希臘文如果翻譯成英文.
通常都是翻譯成Hippocrate.
Hippocrate這個字.
在希臘文來說.
就是解作演戲者或者冒充者的意思.
也就是說這個人本身不是一個良善的人.
但他的演技超凡.
懂得利用動作,經驗,口述.
來塑造自己一個良好的形象.
在社交的技巧上亦都非常超凡.
容易利用別人認為良善的外表.
去掩飾內心的邪惡.
對自己種種的罪行進行開脫.
他一生都在演戲.
成為一個偽善者.
而背地裡只有一個目的.
就是謀求自己的利益.
教會裡很容易出現偽善者.
因為教會是一個講求良善,公義,愛心的地方.
這樣的質素都期望在我們的弟兄姊妹.
分享互動當中能夠找得到.
但當有弟兄姊妹帶著不良動機.
以美麗的外表來得到認同和讚賞.
卻同時以敬虔的外貌.
來做一些邪惡事情的時候.
教會就會很自然地存在著我們所謂的隱藏的邪惡.
而這一類的邪惡其實有它的感染力.
其他弟兄姊妹亦都不自覺地學習.

$^{81}$造成教會內部的分化.
而其中一樣偽善者用來演戲的地方和技巧.
就是利用那把口和嘴來說話去掩飾自己.
所以嘴和偽善者之間有一個很直接的關係.
而今天的講道其實就《猛學福音》的七章一到二十三節.
就說明了偽善者和嘴當中的關係.
原來有三種偽善者的特質.
而這些特質全部都是和嘴有關的.
所以今天的聽道.
我鼓勵大家不要想隔離有沒有這些特質.
想自己有沒有這個特質.
通常我們喜歡喝隔離.
但想自己這個是重要的.
所以第一個特質和大家分享.
就是這個嘴是會尊重傳統過於上帝的話語.
這是記載在第一到第五節.
剛才我們已經聽過這段經文.
這段一到五節又或者去到最後十到十三節的解釋.
它就說到經文一開始就提到有法利賽人.
和幾位的文士就從耶路撒冷來到耶穌那裡聚集.
並提出他們對耶穌的門徒作出一個指控.
那個指控就是他們沒有跟隨古人的傳統.
在吃飯之前就沒有先洗手.
第二節提到門徒就用了剛才的版本.
就是「不潔淨的手」.
和諧本的版本就叫做「濁手」.
英文的翻譯就叫做「Defined Hand」.
或者「Common Hand」.
這是甚麼呢?.
這個叫做「不潔淨的手」.
是否代表沒有用「初授液」呢?.
又或者沒有用「簡易洗手」呢?.
答案就是原來我們會發現.
在利尾記的十五章第十一節.
有一個這樣的記載.
它就說「患了漏症的人.
沒有用水擦手 無論摸了誰.
誰必不潔淨到晚上.
並且要洗衣服用水洗澡」.
這個在說患漏症的人.

$^{121}$這一節的利尾記本身.
其實應用在患漏症的人身上.
譬如女性的血漏.
男性的不能控制的夢迴.
類似這些的漏症.
就說明了那些患漏症的人的污穢.
其實要留意.
本身就不是一些道德上的污穢.
而是叫做禮制上的污穢.
當他們要抹別人的時候.
他們就必須用水洗乾淨手.
他們本身的污穢就不會傳染到其他人身上.
而患漏症的人的污穢.
本身並不是道德上的問題.
這個污穢只不過停留在禮制層面的污穢.
也就是說.
污染了的人是沒有任何道德的問題.
他們只不過不能夠進入會幕敬拜和獻祭而已.
對他們的日常生活是沒有太大影響.
就算他們真的不小心污穢了別人.
這個被污穢的對象.
也只需要做一些潔淨的禮儀.
譬如洗洗衣服.
這樣就可以了.
成為潔淨.
他們就能夠進入會幕去敬拜.
但是事情的發展就是.
直到耶穌基督的年代.
這種先洗手後摸人的律法要求.
就產生了很多的演變.
一方面洗手的要求就提高了.
第三節就記載說.
仔細地去洗手.
根據一些新約學者的分析.
甚麼是仔細地去洗手呢.
原來有一個文獻記載.
當時人們需要拿水的時候.
就要捧在手心那裡.
捧完一會兒之後.
然後內內外外地洗乾淨.

$^{161}$這樣才叫做仔細地去洗手.
即是說你洗一洗也不算洗過手.
即是說洗手的要求就提高了.
另一方面.
由本身先洗手後摸人的原則.
就改為先洗手後吃飯.
是改變了的.
本來是先洗手後接觸人.
然後現在就是先洗手後吃飯.
但是大家想想.
這樣的改變.
有一種根本上的向導的改變.
就是本來先碰人.
是為了別人著想.
但是現在先吃飯就是為自己著想.
即是說本來的律法是不想污衊別人.
現在的改變.
就變成不想污衊自己的改變.
慢慢地.
這就成為了古人的傳統.
古人的傳統.
基本上不是一個上帝的要求.
它只不過是一個古人傳下來的演變.
而那種演變通常是經過口傳的傳統傳遞下來.
所以為何用那把口去傳.
這是重點.
其實不是上帝的吩咐和命令.
在第十到十三節提出一個例子.
古人建立了一個傳統.
當有人願意奉獻自己的財物.
成為國耳版的話.
就不需要供奉你的父母了.
我解釋一下甚麼是國耳版.
國耳版.
我們原來根據希伯來文的讀音去分析.
國耳版這個字.
在希伯來文裡讀作「過班」.
如果你查考這個字.
你就會發現在利未記的第一章裡能夠找到.
甚麼是過班.

$^{201}$通常在利未記就譯作禮物或是貢物.
例如繁祭是一個過班.
因為繁祭就像一個禮物般獻給神.
又或者素祭也是一個過班.
素祭也像一個貢物般獻給神.
即是告訴你.
你獻了繁祭之後.
那些資源當你貢獻給繁祭.
你就不需要用那個資源去供奉你的父母了.
這就是傳統遺留下來的人.
耶穌就罵他們 截斥他們.
說你尊崇古人的傳統.
就違反了十誡的第五誡.
十誡的第五誡其實就是「當孝敬父母」.
如果大家還記得.
其實古人的傳統的微妙之處.
就是將本來是兩個上帝話語的吩咐.
一個就是獻那個過班.
也就是獻那個禮物.
是不是上帝的心意.
其實也是.
這是利未記第一章的要求.
但那邊廂「當孝敬父母」也是上帝的心意.
你只要兩樣都做其實是沒有問題的.
但古人的傳統竟然將兩個上帝的心意對立起來.
即是說你做了這個可以不用做那個.
又或者你做這個可以不需要做那個.
這種對立正正就是很微妙地.
就將兩個本來是上帝的心意.
形成一個對立的局面.
以致就能夠打架了.
這樣我們就因此去到明白.
那個人的吩咐.
他的地位就好像高過上帝的話語一樣.
的一種吩咐.
就耶穌直斥他們的虛偽.
直斥他們的虛偽.
舉個例子.
我覺得回來崇拜應該穿得好一點的.
如果你家裡有件恤衫穿吧.

$^{241}$又或者不要穿短褲,涼鞋.
可能穿一些好一點的.
有球鞋或者牛仔褲也是好.
又或者來到崇拜穿得好一點.
我都覺得是尊敬神.
不過這個聖經裡面是沒有教導.
穿崇拜的應該怎樣穿的.
大家想清楚一點是沒有的.
但這個就是人的傳統.
我都覺得本來其實是可以的.
但如果一間教會守人的傳統.
守到一個地步高過上帝的話語.
這樣就出現問題.
舉個例子.
若然真的有一個衣衫冷漏的人.
進來教會.
而你這間教會有一個傳統.
不可以穿成這樣進來教堂的話.
你就排他於門外.
豈不是正正違反愛倫社兒童自己.
上帝的吩咐.
我再次重申.
穿得好來到教會是好事.
真是好事.
但當這件事去到一個極端.
那個極端就是你穿得不好就不准回來的話.
其實就去到另一個極端.
但去到那個極端就正正違反了.
上帝對愛倫社的吩咐.
又打個比喻.
準時出席崇拜是一件好事.
我覺得是很應該.
正正就是因為.
那個是你要尊敬的主.
你準時甚至早到.
這是一個非常好的一件事.
但是如果我們對於這件事.
推到一個極致.
就說11點.
譬如11點崇拜.

$^{281}$你11點01分就拉閘.
然後不讓任何人進來崇拜的話.
就去到一個極端.
正正就是違反了不可停止聚會上帝的吩咐.
不過不可停止聚會.
其實他也沒有說清楚.
可不可以遲到來聚會.
不過無論如何.
我們就會能夠看得到.
其實我是贊成準時的.
不過很多時候.
譬如舉個例子.
有些父母.
他有一個很小的小朋友.
臨出門口之前.
才來一個大小二便.
你就很難準時到.
來到的時候就拉閘.
不讓你進來.
其實就是一個這樣的做法.
又舉個例子.
如果有吸煙習慣的.
可不可以入小組.
當然我們不贊成吸煙.
吸煙是對身體不好的.
還有吸煙的煙會影響隔離.
但是仍然地.
我們會發現.
讓他進入小組.
這比他真的懂得吸煙重要.
讓他能夠感受到舊的愛.
沒有人是完美的.
我們每個人都知道.
所以有很多時候.
為何人們這麼喜歡古人的傳統.
而廢棄上帝的說話呢.
我告訴你原因.
原因就是.
因為古人的傳統.
你遵守的成本是低的.

$^{321}$但是他的回報.
對你的明星的回報是高的.
但是如果你要實踐上帝的話語.
成本是高的.
為甚麼呢.
愛一個衣衫藍柳的人.
成本高不高.
高的.
但是可能回報真的很低.
或者你不是要問回報的.
其實是這樣.
為甚麼這麼重視古人傳統的人.
會漸漸成為偽善者呢.
就是因為他想透過這些.
成本比較低的遵守.
而去確立自己是一個.
競權外表的人.
我回教會一定不會遲到.
我回教會一定穿得好看.
我是不吸煙的.
種種美麗的感覺.
好像高人一等一樣.
實際上可能是一個偽善者.
因為他會鄙視一些.
衣著不太光鮮的人.
他會對著那些吸煙的人.
會秒秒醉.
他會對於遲到的人.
直覺就會覺得他不淑齡.
當然我再次強調.
準時回教會是很應該的.
穿得好看是很應該的.
最好不要吸煙.
我們每個人都知道.
不過我們因此去到明白.
古人的傳統和上帝的吩咐.
有可能真的有這樣微妙的差別.
所以第一點跟大家分享.
就是尊重傳統過於上帝的話語.
很有可能是偽善者的表現.

$^{361}$第二點.
我去到下一章.
我去不到下一章.
就是用嘴唇尊重神.
但是心卻遠離神.
這是第六到第九節的經文.
耶穌在回應伯利賽人.
針對門徒不耕耘古人的傳統指控時.
他引用了《以賽亞書》二十九章.
十三節的經文.
我們看看這段經文.
你明白.
他說「主說:因為這百姓親近我.
用嘴唇尊敬我.
心卻遠離我.
他們敬畏我.
不過是領受人的吩咐」.
耶穌引用了《以賽亞書》時.
他引用的版本.
其實就是LXX七十四億本的版本.
所以就和本身希伯來文的版本.
有少許的不同.
主要的就是LXX七十四億本.
強調人的吩咐.
當作神的吩咐.
而本身希伯來聖經MT的版本.
就是心口不一的情況.
就是人的吩咐.
但無論如何.
重點就是那個嘴的問題.
而偽善者就將自己提高成為教師的角色.
所教的內容就是這些人的吩咐.
而不是上帝的說話.
而是由於實踐人的吩咐.
往往只需要較低的成本.
卻換來更高的回報.
所以當他們做得到的時候.
別人就會以為他們是真理的實踐者.
吸引人的羨慕.
慢慢地人就沒有那個心.

$^{401}$去跟從真正上帝的吩咐.
《以賽亞書》十九章十三節.
其實是一句應酬的話.
為甚麼這樣說呢.
特別是「親近」這個字.
本身是禮制的用詞.
通常用來描述人將那個「個班」.
也就是那個「供物」.
去到獻給神的動作.
「親近」這個字有這樣的意思.
這個禮物是甚麼呢.
就不是那隻牛或羊.
而是祂的嘴和祂的唇.
百姓用這些去尊敬神.
但祂的內心卻遠離神.
在這裡我們看到.
內裡的心和外在的口.
形成一個強烈的對比.
就指出外表的敬虔和內在的心.
明顯地南轅北轍.
成為內外不一的人.
由於那個口和那個唇.
所說的敬虔的說話.
未必真是反映你內心的實況.
所以其實是一種應酬性的說話.
大家同不同意.
可能你心裡不太喜歡某個人.
但基於禮貌.
不想讓人覺得自己太沒禮貌的話.
你就應酬的說話對待他.
我們叫這些應酬的說話.
如果你放回以賽亞書的場景去看.
有一句說話是以賽亞書蒙詔的時候所說的.
那句說話就說.
「因為我是嘴唇不潔的人.
又住在嘴唇不潔的民當中」.
解釋給大家聽.
百姓吹嘴唇不潔.
並非是因為他們說粗口.
並非他們是因為說一些不乾淨的話.

$^{441}$反而他們的嘴唇是說一些很文雅.
但是他們的嘴唇尊敬上帝.
心系遠離上帝.
我們叫這些「內外不一」.
非常虛偽.
所以他們的嘴唇成為不潔.
並非因為說粗口.
而是說一些比粗口更加難聽的說話.
就是表面上公公敬敬.
實際上非常陰毒的一種說話.
所以我們會看到.
這群人沉醉在謊言和宗教的術語裡.
卻不認真,正直地去履行上帝的吩咐.
我解釋給大家聽.
我想了很久.
究竟有沒有一些人會是這樣呢?.
我就想到原來喜歡說是非的人.
通常都是屬於虛偽的人.
我不知道大家是否同意.
說是非的邏輯是怎樣的呢?.
說是非的邏輯就是「人地衰你好」.
大家同意嗎?.
這是說是非的邏輯.
說是非的人通常都會說到一個地步.
或者大家知不知道.
原來說是非是比較能夠增進友誼的.
即是和隔壁的人一起說.
然後說某人的壞話.
當然你和那個人覺得自己和他不同.
即是他是最壞的.
我們兩個是最好的.
慢慢你就明白原來說是非是一種甚麼呢?.
一種認清自己和對方是有分別的.
一個這樣的本質.
而這樣的本質我們就會發現.
正正就是看不到原來.
其實自己都可能是一個罪人的情況.
當你面對到上帝的神聖.
就好像以賽亞那樣.
你第一個直覺就不是覺得別人是壞的.

$^{481}$當你面對上帝的神聖.
當你屬靈是很好的時候.
你直覺就會看到.
不是別人的問題是我的問題.
因為以賽亞都看到.
他是活在嘴唇不潔的文當中.
他又是嘴唇不潔的人.
這個這樣的察覺.
並不是任何說是非的人.
懂得察覺的.
說是非的人之所以虛偽.
就是因為他用一把尺去度別人.
而不懂得用同樣的尺去度自己.
這個是其中一樣.
我們能夠看得到.
話說有一位猶太拉比.
他的嘴唇經常都讚美上帝.
敬畏上帝.
但是有一次.
他就見到一個窮人.
那個窮人的樣子.
他看上去是非常醜陋.
就走去跟那個窮人說.
其實你知不知道你不適合敬拜上帝.
因為你的樣子好像反映不了上帝.
那個窮人就反問回這位拉比.
就跟他說.
你說我不適合敬拜上帝.
不如你試試問一下創造我的那個主.
我適不適合去敬拜他.
原來任何的樣子都是上帝所創造.
只不過人是分門別類.
將這些人.
我們看得到.
第八節就指出這樣的人.
是離棄上帝的誡命.
守人的傳統.
第九節進一步說明.
否定上帝的傳統.
用了你的傳統去形容.

$^{521}$由人的傳統成為你的傳統.
將人的吩咐成為自己的吩咐.
這個是偽善者墮落的原因.
第二點就是這一點.
最後.
「嘴中所出的都能夠污衊人」.
在18到23節.
耶穌基督進一步進行教導.
就指出入口的東西不能污衊人.
而祂所指的污衊.
並不是一種飲食衛生的污衊.
而是指一種道德層面的污衊.
也就是說無論你吃魚生.
無論你吃海鮮西餐還是中餐.
那些食物並不能構成你道德上的污衊.
這直指如果我沒有按照古人的傳統洗手.
根本就無關痛癢.
最多你是不衛生.
有些菌的時候就去了茅廁那裡.
污衊而已.
對本身個人生命的道德操守.
它的品格是沒有任何的影響.
反而偽善者所出口的東西.
才能夠污衊人.
而這些污衊.
一系列的清單是道德上的污衊.
帶有強而有力的感染力.
可以傳染每個人都和偽善者一樣一起墮落.
21到21節列出一系列污衊的清單.
當我們細心看.
全部都是道德的問題.
是由一個人心裡面發出來.
根據《以賽亞書》我們剛才所提到的29章13節.
偽善者的心已經遠離神.
就算嘴巴說著一些敬虔的話.
但他們的內心還是這麼邪惡.
但由於他們的心長出一系列道德的失敗和敗壞.
但他們的嘴巴又說得很敬虔.
可想而知其實是非常危險的一種情況.
舉個例子.

$^{561}$在美國有一個情況出現過.
就是有一個男小學的老師.
很疼愛那班學生.
但就出現一位九歲的女小學生.
她很不喜歡這位男老師.
不知為何九歲的時候就學會怎樣說謊.
然後就投訴這位男老師性騷擾她.
這位女同學.
這位女同學也連同幾位不喜歡這位男老師的女同學.
一起就說整班的女同學都說這位男老師.
就是性騷擾這班的女同學.
後來就報警處理.
原來報完警之後就發現.
這位男老師就沒有做過這些事情.
是那班女同學誣衊.
又或者誣告這位男老師.
本來事情就完結了.
應該就不會.
因為體育路電視各方面都沒有做過這些事情.
最後是因為那位女同學的爸爸媽媽.
就接受不了.
或者不覺得自己的女兒真的說謊.
所以她的父母就覺得要繼續告下去.
最後因為壓力的緣故.
那位男老師就無奈地要辭職.
就要離開那間小學.
你就會看到其實人的嘴.
其實很有可能都挺恐怖的.
人的嘴是可以很恐怖的.
人的心就反映在嘴裡.
我們其實求助幫助我們的.
就是離開心中裡面那份邪惡.
就是要離開成為偽善者.
我們的嘴要寬大衛懷.
將疑點利益歸於別人.
舉個例子.
如果你今天進來崇拜.
有人沒有跟你握手.
你就要寬大衛懷.
不要覺得他是不是沒禮貌.

$^{601}$是不是今天看不起我.
是不是覺得我有甚麼不方便握手.
其實很有可能背後最簡單的動機.
就是因為最近肺炎.
可能這樣就算了.
其實可能真的為別人想一個善意的解釋.
又不一定想負面的解釋也不一定.
將疑點利益歸於別人.
這也是令到這個群體更加良善地.
去到對待別人的一個契機.
所以三點跟大家分享.
第一點尊重傳統過於上帝的話語.
第二點用嘴去尊重神.
但心卻遠離神.
第三點嘴中所出的都能護衛人.
這三個特徵成為偽善者的特徵.
最後這是我自己看聖經的時候.
提出的七個結論.
和大家看看.
通常耶穌基督所說的那些假冒偽善的人.
有甚麼特徵呢.
大家可以看看.
第一個特徵.
喜歡功德和回報勝於恩典與不配.
第二.
喜歡世俗名銜勝於神兒女的身份.
第三.
喜歡別人記得自己的善恨多於神記得你.
第四.
喜歡服侍之後計算功績過於榮耀上帝.
第五.
喜歡外表的華美過於內在的敬虔.
第六.
喜歡賣弄知識勝過看到生命的改變.
第七.
喜歡批評別人過於虛心學習受教.
這是我自己看聖經的時候.
那天不知為何心血來潮就寫下這七句.
寫完之後.
原來那些字數是很工整.

$^{641}$都是這麼多個字.
每一句都是這樣.
都成為我的提醒.
都成為我的每個弟兄姊妹的提醒.
求主幫助我們能夠成為一個回歸小孩子的心.
單純地去敬拜上帝.
用善去對待身邊的人.
我們不是要活得完美.
正正我們不完美.
我們更加需要不要充當一個.
叫做帶頭鬼的角色.
去成為偽善的人.
我們要謙卑.
承認自己的不足.
我們一起在這個神聖的路上.
彼此兼併兼同行.
這是很寶貴很真誠.
但一下看不開.
我們就落入了這個圈套.
祈求主幫助我們每一個.
我們一起祈禱.
慈悲仁愛的天父.
多謝你讓我們看到你的教導.
耶穌基督的吩咐.
祈求主你讓我們能夠離開心中的邪惡.
那種的偽善.
以致我們可以真誠地.
愛愛如一地去敬拜榮耀你.
求主幫助.
奉主耶穌基督的名求.
阿們.
\newpage



\section{}
\label{sec:mXSWAPPoVUc}
\textbf{2023-01-26 周四晚查經班《被擄歸回歷史研讀(4)》 講員:曾裔貴牧師}
\newline
\newline
連結: \href{https://youtube.com/watch?v=mXSWAPPoVUc}{\texttt{https://youtube.com/watch?v=mXSWAPPoVUc}} ~~~~ 語音日期: 2023-01-27
\newline
\newline
\hyperref[sec:S5k9W4TGOeU]{\small{< < < PREV SERMON < < <}}
~
\hyperref[sec:index_chronic]{\small{[返順時目]}}
~
\hyperref[sec:index_scriptual]{\small{[返順卷目]}}
~
\hyperref[sec:4sQReCkRl1k]{\small{> > > NEXT SERMON > > >}}
\newline
\newline
$^{1}$我們一起簡單祈禱,多謝主再一次我們過了年之後,又恢復有茶經,.
我們很盼望透過研讀舊約歸回的歷史書和先知書,.
當我們一併來看的時候,希望能夠從其中,.
雖然是已經發生了,已經過去了的歷史,.
猶太人也有他們過去的聖殿和在主耶穌來到的時間,.
和主後70年聖殿被毀,猶太人他們要亡國,.
他們要被分散去到天下,這些都是我們知道的歷史,.
也是教會時代的開始,我們今天已經過了2000年,.
教會仍然屹立,全世界仍然等待主耶穌來塞亞的降臨,.
幫助我們對這一段差不多有400年的歷史,.
接近200年多一些,或者可以是400年的歷史,.
都能夠有所掌握,求主帶領我們,.
我們這樣仰望禱告奉耶穌基督的名,阿們..
好,和大家重溫上一次,就是對上一次我們完之前提到,.
就是當他們興高采烈要來重建這個聖殿的根基,.
希望來到去恢復那個憲制先,至少搭好一個螢幕,.
首先最重要有個祭壇,但是其他的可能就未必有所羅門的規模,.
因為所羅門的規模是非常之大的,所以是前無古人的,.
但是我們上一次說過,有幾方面的攻擊立即就出現了,.
首先就是當地的原居民,其實就是由列王記下第17章那裡,.
當時的亞述將外族人遷徙到北面的以色列,主要是薩馬利亞,.
所以那裡的人就不再是純種的猶太人,.
但是隨著那個歷史,隨著那個王朝,.
到了巴比倫的時候,連猶大國都亡了,.
南面的猶大國都亡了,所以整個以色列都要被擄,.
所以猶太人就被擄到巴比倫,.
亦都是輾轉過了70年,歸回歷史的開始,.
就是古列王頒布天下,所有的民族曾經被巴比倫王擄去的,.
現在可以回到自己的原居地,甚至乎拿國庫的一些資源來在當地,.
因為有很多的行省,好像是20個的行省,.
這個是古列的時候設立的,.
這樣來到這些行政區就可以調撥那些資源,.
當然是那些材料,石,木,來去重建宗教的temple,.
因為不同的民族都有自己的神,.
所以其實不只是猶太人,.
不過在二賽亞書45章,古列就是神的那個僕人,.
某程度是那個被神感動,為神效力的那個僕人,.
來到去重建這個聖殿的,.
所以我們讀這些歷史書和先知書,.
這一段時間是風起雲湧的,.

$^{41}$你越對波斯帝國的早期某程度是他們的法則史,.
他們一直的強大,去到大理烏王是整個波斯帝國最強大的,.
當然後來就是希臘亞歷山大這個軍事天才的興起,.
所以你可以想像得到由主前的五百多年,.
他慢慢將巴比倫,將馬代,來到去吞併,.
成立了自己的波斯帝國,.
那個版圖去到印度河,去到西面,去到地中海,.
很大的版圖,.
整個現在的伊朗的地方,一帶甚至可能去到高加索的地方都不出奇了,.
你可以想像得到是非常大的帝國,.
所以我們看這幾段的歷史,.
我們如果純粹看聖經,.
我們未必感受到在世界那個風起雲湧,.
但是如果我們對古代的波斯帝國的歷史,.
或者說歐洲史,因為亞歷山大大帝,.
這個希臘的軍事天才,.
這樣來到去興起,.
當時的波斯帝國和希臘的城邦,.
特別是那些斯巴達人,.
非常勇悍的,.
那個戰爭,好幾場關鍵的戰爭,.
我們讀中學的時候都有機會讀過的,歐洲史,.
你可以想像到,.
聖經給我們看到有三位比較重要的王,.
其實應該是四位的,.
因為例如古列我們有看過,.
這個班布詔書,.
某程度就是歌舞昇平,.
因為他打拼贏了巴比倫,贏了馬代,.
等等,其實馬代王是他的岳父,.
所以開始懷柔,後來就吞併了,.
跟著他的兒子,.
就是在聖經講的第四章,.
去到波斯王的亞達薛西,.
或者應該叫做阿赫隋魯,.
四章第六節,.
這個阿赫隋魯,.
這些國號好幾代都用的,.
所以嚴格來說,.
看回那個帝系圖,.

$^{81}$第二任的皇帝就應該叫做金拜西二世,.
當然他們的代號就是阿赫隋魯,.
所以你不要混亂了,.
在以斯帖記快要出現的將來的阿赫隋魯,.
那個是另外一個王,.
所以這個金拜西二世,.
其實就是古烈的兒子,.
古烈主前的530年離世,.
跟著就是這一位金拜西二世做王,.
所以在新王登基,.
也是當時,.
上一次我們講,.
猶太的,.
特別主要是耶路撒冷的原居民,.
就開始感覺到,.
這一班上來重建聖殿的人,.
不懷好意,.
甚至乎如果他們真的能夠成功重建了,.
就會對他們有一個威脅,.
所以首先用一些比較懷柔的手段,.
就是我們一起做,.
我們也很想敬拜你們的神,.
我們在這個地方已經住了超過70年,.
已經在跟隨著敬拜你們的神,.
當然又拜耶和華,.
又拜自己的偶像,.
這個就不在話下了,.
所以很決定當時上去的領袖,.
索羅巴伯,.
他怎樣理解這一班人,.
他們就很強烈,.
其實他們很少人,.
我們講過,.
五萬人,.
相對於70年在波斯帝國,.
在巴比倫帝國,.
已經落地生根的猶太人,.
被擄的第二第三代,.
其實是少得可憐的,.
可想而知,.

$^{121}$這五萬人,.
聖經要記載他們每一個家族,.
每一個身份,.
可想而知,.
在神的眼中,.
這班人,.
分外顯得寶貴,.
他們連根拔起,.
回到重建神的殿,.
但是相信這個重建,.
重建包括心靈的復興,.
重建就會帶來,.
一定有你的習慣要被打破,.
你去到別人的地方,.
其實是你自己的國土,.
就要面對,.
當地人對你的不友善,.
所以上一次我們都歸納了,.
有幾方面的,.
慢慢升級那個反對的,.
首先就是一起做,.
如果你認定這一班人,.
最終他們會蠶食,.
宗教的那個純正,.
這個是非常之微妙的事情,.
所以有時覺得,.
再白一點,.
如果你有心想在教會裡面,.
來到去呼風喚雨,.
你找你的同路人,.
即是你能夠操控票數的人,.
進入那個議會,.
哪個權力最大的議會,.
進入那個議會,.
你慢慢來操縱那個議會,.
你就可以操縱整個教會,.
這些不是什麼秘密,.
所以極之極之要小心,.
不要將自己的野心權力,.
帶進教會裡面,.

$^{161}$這個是非常之,.
神是不喜悅的事情,.
所以慢慢慢慢,.
你就會發現,.
如果所羅巴伯,.
容讓這一班外族人,.
明顯沒有家譜,.
不是猶太人的,.
容讓他們進入那個重建裡面,.
他們的人多,.
他們的聲音,.
他們那些外族的敬拜假神的觀念,.
慢慢慢慢就會蠶食淡化,.
猶太人他們要back to the mosaic,.
摩西的律法書,.
那個的Torah,.
那個的摩西五經,.
那個的標準,.
那個的神的,.
在舊約時代的準繩,.
如果不是,.
神就不會賜福,.
這個哈該書也有重申這件事,.
在他的第二次的講論裡面,.
等一下我們也有機會先看哈該書,.
你就可想而知,.
這個是所羅巴伯的堅持,.
獨排眾議,.
可能未必是獨排,.
可能是和一些領袖商討過,.
因為至少他們自己回來的時候,.
因為你既然是要重建聖殿,.
要立根基,.
要來獻祭,.
哪些人可以獻祭呢?.
就是祭司,利未人,.
所以他們就要尋查那些家譜,.
有些人說自己是,.
我也是祭司,.
但是那個家譜遺失了,.

$^{201}$可能記錄了他的家族,.
所以不要緊,.
我們不是不信任你,.
不過要等待有那個祭司興起,.
他能夠用污名土名來決議,.
到時候就看一下那個決斷,.
到時候就可以讓你們去恢復那個事奉了,.
就是說猶太人他們自己,.
很緊張這個這樣的身份,.
所以現在外族的當地人,.
說要和他們合作,.
就是說給材料,.
給那些各方面的資源等等,.
方便的,.
但是猶太人就不肯,.
第四章的第三節,.
和大家回到記憶,.
因為過了兩個禮拜,.
但所羅巴伯耶穌亞和其餘以色列的族長對他們說,.
我們建造神的殿與你們無關,.
我們自己為耶和華以色列的神協力建造,.
是照波斯王古列所吩咐的,.
你可以想像得到很決絕,.
與你們無關的,.
我們自己去做,.
與你們無關的,.
因為我們建造的是神的殿,.
我們是猶太人,.
我們是神感動我們,.
激動我們,.
透過古列,.
神都用當時的歷史,.
神都用當時的列王,.
神都用當時最強大的波斯帝國,.
來使他們因利成便,.
用國家的資源,.
所以我們不要抽離了自己的處境,.
所以真是風起雲湧,.
當時的波斯帝國,.
沒錯是很強大,.

$^{241}$但是去到第二代,.
金拜西,.
接著到大理烏是最強大,.
其實都已經有列國的城邦勇士,.
領袖,.
開始興起挑戰波斯帝國的強權,.
你可以想像得到,.
猶太人的歷史,.
重建的重要歷史,.
都仍然發生在平衡的時空,.
即是神的國度,.
藉著一班聖文,.
神聖的子民,.
回去重建,.
回去聖地,.
重建聖殿,.
重建聖城,.
但是原來在周邊的外圍,.
有外族人,.
有戰爭,.
有列強的吞併,.
你可以想像的,.
戰爭是沒有任何停止過的,.
但是在神的聖經裡面,.
他著重的是這個時代,.
大概如果從以色列記到撒加尼亞書,.
大概是150年左右,.
之後就再沒有任何的密使,.
要到施洗約翰第一個才出現,.
就再有神的密使,.
你就可以想像得到,.
這段時間,.
神是感動古列的時代,.
但是同時間,.
是感動一班人,.
其中是以索羅巴伯,.
大祭司約書亞,.
和兩位先知哈該和撒加尼亞,.
當然後期還有馬拉基先知,.
你就可以想像得到,.

$^{281}$我們讀聖經,.
如果能夠有一個平衡的時空,.
其實我們看《約翰福音》看過,.
要有一個平衡的時空,.
來看我們聖經的世界,.
雖然好像很平鋪直敘,.
但是原來在外面的世界,.
是不斷發現充斥著混亂,.
吞併,.
好像不斷地洗牌,.
這些都是那個時代的特色,.
所以任何一個時代跟隨主,.
都不是容易的,.
要有一個很大的決心,.
但是現在這個這麼大的工程,.
透過五萬人去造成,.
五萬人加幾位領袖,.
去造成這麼大的工程,.
等待真正的聖殿的主,.
就是耶穌基督,.
其實祂自己就是那個,.
來到天上降到地上的那個人,.
任被祂駐在我們中間,.
充充滿滿有恩典,.
有真理,.
我們也見過祂的榮光,.
乃是赴獨生子的榮光,.
所以可想而知,.
神在耶穌未來之前,.
祂都要恢復敬拜的世界,.
那個秩序,.
不容讓猶太人散散地離開自己的故土,.
所以這裡就是歷史的發展,.
就是在甘拜西剛上任的時候,.
他的父王,古列大帝死了,.
就出現了這個波折,.
第四節,四章四節,.
那地的民,就是當地的人,.
就在猶大人建造的時候,.
使他們的手發軟,.

$^{321}$擾亂他們,.
你不讓我們一起去做,.
不容許我們參與你們的敬拜,.
就是這麼決絕的一刀兩斷,.
沒有任何的來往,.
所以他們就不會罷手,.
就擾亂他們,.
第五節,從波斯王古列年間,.
直到波斯王大利烏登基的時候,.
賄賣謀士,大利烏是第三個皇帝,.
是准許重建的那個,.
所以在這裡,.
作者或者編輯者,.
將整段時間,.
由古列開始,.
去到大利烏准許重建,.
中間就是有甘拜西二世,.
原來這麼長的年日,.
530至522,.
有十幾年,.
原來不斷賄賣一些謀士,.
用不同的方式,.
要敗壞他們的謀算,.
就是猶太人的謀算,.
用賄賣洗黑錢,.
務求讓你的工程不能夠成就,.
這個就是你可以說當時的仇敵,.
魯修成路,.
敬酒不喝,喝罰酒,.
說跟你一起做,.
你們不肯,.
現在就要這樣攻擊你們,.
這些都是我們要留意的,.
暫時香港的教會,.
我暫時都沒有聽過,.
不要緊,教會沒有錢,.
叫馬會贊助,.
這麼多錢,沒問題,.
或者在禮拜堂,.
平日不用開教會的時間,.

$^{361}$沒有聚會,整個禮堂空了,.
不如我們找一些生財工具,.
開麻將桌抽水,.
我暫時沒有聽過,.
近乎沒有人會這樣想,.
因為明知道一定是大錯特錯,.
就不會這樣做,.
但是我們都要留意,.
骨子裡面我們要問,.
世俗有沒有進入教會,.
這個後話,.
這個是第七章的以色列記,.
一直延伸到尼希米記第七章,.
都是講這個問題,.
世俗有沒有玷污了聖文,.
聖殿和聖城,.
這三個聖的地方,.
就是人,事物,.
有沒有被外族玷污,.
有,以色列記,.
我們快要看的第八章,.
有,就是intermarriage,.
和外族通婚,.
延伸到第十一章,.
以色列記都是處理這個問題,.
就是世俗不可以進入聖城,.
其實應該進入聖文的婚姻關係裡面,.
所以這裡就告訴我們,.
我們要留心,.
但是有一個很大很大的問題,.
就是到底舊約這種斷絕關係,.
去到主耶穌的新約時代,.
怎樣理解,.
耶穌說,再沒有分猶太人,.
撒瑪利亞人,.
保羅說,無論是猶太人,.
西古提人,化外人,.
圍爐的,自主的,.
或者男或者女,.
在主耶穌裡面都是一,.

$^{401}$那麼這樣的一個轉變,.
轉向,.
這些就是我們一生要學的功課,.
我們如果像舊約的時代,.
黎希米的時代,.
以色列的時代,.
一律的禁止,.
是相對比較容易,.
但是隨著時代和聖經的演變,.
因為舊恩的歷史的進程,.
原來有些東西,.
反過來,現在不是這樣,.
因為去到耶穌時代,.
他們就真的,.
什麼都好像和別人要分開,.
但是耶穌說不是這樣的,.
這些就是,.
當我們進入到那些經文,.
我們再討論,.
現在我們說回那個蘭佐,.
起那個根基,.
都已經有一個阻撓,.
就是說那個聖殿,.
可能只是有那個祭壇,.
但是就沒有其他的那個配套,.
所以出現了這個問題,.
另外就是上個禮拜,.
上一次我們有機會,.
整篇的那個控訴詞,.
是來到去看到,.
作者將整份的,.
控告的那個,.
那個狀,.
記錄在第四章裡面,.
很特別就是,.
以色列有兩份的狀書,.
都是詳細記錄,.
是寫給皇帝的詔書,.
是由那些專業的,.
那些狀官,.

$^{441}$來到去撰寫的,.
但是,.
接收的人,.
是完全那些內容一樣的,.
但是,.
接收的皇,.
居然間會有不同的理解,.
那你可想而知,.
有時都不是說,.
那麼簡單,.
就是說,.
一定我們做了那件事,.
就一定會有那個好的結果,.
第四章第六節,.
在阿克徐老才登基,.
這個就是甘拜西二世,.
就是古列的兒子,.
因為古列死了,.
那他的國號就是,.
阿克徐老,.
就是類似羅馬,.
就是蓋薩,.
埃及的就是法魯,.
法魯是一個名號,.
不同的王都叫法魯王,.
上本控告猶大和耶路撒冷的居民,.
那個本,.
阿達薩西年間,.
比斯蘭米特利達他別和他們的同黨,.
上本就告波斯王阿達薩西,.
本章是用亞蘭文字,.
亞蘭方言,.
省長利宏,.
書記申遂,.
要控告耶路撒冷人,.
也上本就告阿達薩西王,.
省長利宏,.
書記申遂,.
和同黨的抵賴人,.
等等,.

$^{481}$我們之前讀過的,.
去到第十一節,.
上奏阿達薩西王說,.
河西的臣民云云,.
當時他們波斯帝國以西的地方,.
以右法拉底河那一帶,.
來到以西,.
伊朗原來,.
我剛才找資料,.
古烈王的陵墓,.
是伊朗的國境,.
很大的陵墓,.
值得去看一下,.
伊朗原來有很多古代的遺址,.
比如巴比倫的空中花園,.
就近伊拉克的頭,.
所以伊朗如果是政治,.
不是那麼危險,.
其實伊朗是值得去看的,.
有很多聖經時代的遺址,.
歷史,.
那就,.
哪位,.
是不是可以關一下那個,.
什麼先,.
可以,我靜音了,.
第十二節,.
王該知道從王那裡上到我們這裡的猶大人,.
已經到耶路撒冷重建,.
等等,我不詳細讀這裡了,.
因為我們之前讀過了,.
第十七節,.
那時王裕福省長利弘書記身遂,.
和他們的同黨,.
就是住撒馬利亞,.
沒錯了,撒馬利亞一帶的人,.
那裡,因為那些全部是外族的人,.
所以,.
由那個歷史淵源,.
一直延伸到耶穌基督的時代,.

$^{521}$撒馬利亞都是猶太人,.
所看不起,.
鄙視,.
甚至乎是刻意和他們分開的,.
你想一下,.
這樣的一種隔膜,.
是了,.
很難想像那種隔膜,.
在耶穌的面前,.
連自己的門徒,.
看到這些撒馬利亞人,.
都有那個觀念浮上心靈,.
就是,.
哇,.
撒馬利亞人,.
在福音書時時都有這些衝突面,.
所以,.
這些是歷史帶來的包袱,.
盼望我們上一代,.
給我們傳承的包袱,.
不要我們自己去照單傳修,.
我們信了主耶穌,.
應該要有那個包容,.
接納,.
不要再分什麼人,.
政治,.
宗教,.
都不要再分了,.
問題就是,.
他是不是在主耶穌裡面,.
在救恩裡面的人,.
如果沒有呢,.
那他需要救恩,.
每一個人都是得到公平的,.
神的拯救,.
所以我們要盡力帶他進入救恩裡面,.
當然有不同的民族,.
有不同的困難,.
有些是多神的,.
他很難相信耶穌是獨一,.

$^{561}$有些是敬拜他們自己的,.
譬如阿拉真神,.
或者是印度教,.
你叫他改信,.
就很大的困難,.
但是要怎樣幫助這些人去歸信呢,.
這個都是我們大家一起要去研究,.
你可以想像得到,.
一個民族他們要改變,.
需要很多很多,.
很多的一代一代的努力,.
可能踏石仔那樣,.
才能夠踏出來的,.
所以盼望我們都能夠給一個宏觀去看,.
不同的民族需要救恩,.
譬如巴基斯坦,.
現在巴特利亞中學,.
很多巴基斯坦的,.
很難,.
他們聽到靜娟姑娘說話,.
已經塞住耳朵了,.
很困難,.
尼泊爾不困難,.
因為他們多神,.
耶穌說好,沒問題,.
不過到他要,.
哇,不行,.
耶穌才是唯一真神,.
他就要卻步了,.
他就有困難了,.
他就沒辦法暫時接受了,.
因為他們拜的是多神的印度教,.
這些就是不同的特色,.
所以我們很盼望,.
我們能夠不要再對人表面的情況,.
來到他先有了一種心靈的拒絕,.
而是我們要慢慢將他帶進耶穌的救恩裡面,.
當然有些人一生都拒絕耶穌,.
否定耶穌,.
攻擊耶穌,.

$^{601}$攻擊教會,.
那些我們只能夠留給神去管理,.
我們要的是福音和恩典,.
這個是我們應該要有生之年去做的,.
第四章就告訴我們,.
他們就是這樣去控告,.
第十七節的下半節,.
說:願你們平安云云,.
你們所上的本已經明讀在我面前,.
王親自的回覆,.
所以我們有幸看到這個告的預兆,.
而有皇帝親自的回覆,.
當然有他的使觀,.
我以命人考察,.
得知此成,.
古來果然背叛列王,.
其中常有反叛背逆的事,.
從前耶路撒冷也有大君王統管河西全地,.
人就給他們進貢交貨納稅,.
現在你們要出告示,.
命這些人停工,.
使道者成不得建造,.
等我降旨..
他們都估計到建造完聖殿,.
就會建造聖城,.
所以這個阿達舍西,.
其實應該都是甘拜西這個兒世,.
來降旨,.
禁止他們重建,.
等到我降旨,.
你們當謹慎,.
不可以遲延,.
為何容害加重,.
使道王受虧損?.
第23節,.
阿達舍西王的上諭,.
讀在李宏和書記身上,.
並他們的同黨面前,.
他們就急忙往耶路撒冷,.
去見尤大人,.

$^{641}$用勢力強迫他們停工,.
拿個王的禦紙,.
王榜,.
來頒令立即停工,.
比勞工署禁令還要重要,.
好像工業意外,.
立即要停工,.
一切工程,.
延後,.
損失,.
那些不理了,.
因為現在涉及到政治層面,.
即是原來政策改變了,.
第二任的王,.
他可能又想推翻自己父王,.
古烈懷柔的政策,.
即是巴比倫做得不對,.
現在我新王登基,.
我就採取一個懷柔,.
你們可以回到你們自己的原居地,.
建造你們原居地,.
你們所拜的神的廟,.
或者殿,.
但是仍然在那個行省裡面,.
你們要納稅,.
要進貢等等,.
有你們民間宗教上的自治,.
但絕對不是政治上的自治,.
這個其實是很寬容的一個政策,.
即是說一改朝換代,.
連猶大人自己都受虧損,.
受影響,.
透過當地人的所謂女性外合,.
就來攻擊他們要停工,.
所以我們面對疫情,.
面對口罩令,.
面對聚會,.
只有幾多人聚會,.
有時我們都沒有辦法,.
但我們又不要被那件事影響了,.

$^{681}$一方面好像和政府為敵,.
針對教會,.
我聽過這樣的傳聞,.
但怎麼會針對你的教會呢?.
現在不只是教會不准聚會,.
任何的餐廳,.
任何的地方都可以看電影,.
可以去跳舞,.
是沒有這件事的,.
是全香港,.
全世界都一起停,.
怎麼會突然間這樣想,.
這些特別針對教會呢?.
完全不是這樣的,.
所以我們都要記住,.
當時帝王的政策轉變,.
教會,當然是猶太人,.
猶太人他們回歸,.
就面對一個突然間的政策轉變,.
而帶來的影響,.
再加上當然,.
聖經很直接說,.
是有敵人,.
是有不喜歡你回去重建的人,.
來阻撓,.
所以政策是很遠的,.
在帝國那裡,.
他們管不到這麼遠的,.
但是當地人這樣來上書,.
好像上訪,.
那就不同了,.
那就王好像借這個機會去打壓,.
彰顯他的聲望,.
可能讓他爸爸更加收緊,.
不同的王都會有這樣的做法,.
所以第24節就是那個結果,.
於是在耶路撒冷神殿的工程就停止了,.
一直停到波斯王大利烏的第二年,.
就是520年了,.
522年登基大利烏,.

$^{721}$跟著就是和但以理的大利烏又不同,.
那個是馬代王的大利烏,.
這個是波斯王的大利烏,.
兩個是不同的,.
已經相差一段時間了,.
所以這個時候,.
又是有大概12年的時間,.
那個聖殿的工程就完全停頓,.
人不可以建造聖殿,.
不代表他不需要生活,.
所以他們就會繼續有他們自己那個,.
已經沒有了能夠團結的力量,.
就只會分散,.
所以一個群體很需要的是那個凝聚,.
很需要的是有神的道,.
他們有個方向,.
所以道在新約很多時候是說那個way,.
是那個方向,.
很重要的,.
所以這個時間就是他們都不知道可以做什麼的時候了,.
12年,.
那個王金拜西,.
繼續這樣去借著這些,.
可能是借著這些蠶神,.
或者這些地方的諸侯,.
或者這些的陀地,.
當地人,.
這樣的上訪,.
這樣的上書,.
就借這些機會,.
就去打壓這些恆省裡面的那個自治,.
宗教的自治了,.
收緊了那個政策了,.
所以你在那個時代,.
你都要面對那個時代的那個衝擊的,.
你很難說你是絕緣,.
避免,.
沒有的,.
耶穌的時代,.
都是正正式式在羅馬帝國的統治了,.

$^{761}$彼拉多都是分封的王了,.
希律都是加利利分封的王了,.
就是說都是在under這個政治的規模底下,.
但是,.
奇就奇在這一點,.
有信心的人,.
他可以看到神,.
看到神在那個地上的工作,.
有一套電影,.
本來我沒有看過的,.
後來有人介紹我看,.
當然是構想出來的,.
好像類似近乎幻想片那樣,.
劉青雲的神探大戰,.
唯有他看到人看不到的怪獸,.
他破案,.
差不多100\%,.
那件事還沒發生他都看到,.
本來很恐怖的,.
差不多好像恐怖片那樣,.
很懸疑,.
很離奇,.
近乎有精神失常,.
但是他真的看到背後那個怪獸,.
當然這是一個虛構和世俗的電影,.
但是聖經告訴我們,.
每一個時代都有他的政治層面,.
社會的層面,.
不是絕緣,.
突然間有一個地方什麼都沒有的,.
只有一群熟神的人,.
從來都沒有的,.
從來都不會是這樣的,.
所以現在回去重建,.
有帝王的御准,.
是神感動他的,.
但是他的兒子一上台,.
就立即收緊,.
你都是要受那個影響,.
那個影響,.

$^{801}$如果追本族元,.
是所羅巴伯和那個領袖群,.
反對他們參與重建,.
而導致人家老羞成怒,.
去投訴在這個金拜西二世,.
那個王權底下,.
而導致猶太人那五萬人,.
都受到影響,.
不能重建,.
這個時候,.
整個國家,.
未是國家,.
整個猶太人這一班被老歸回的,.
軍心就渙散了,.
12年的時間,.
去到大理烏的年代,.
去到大理烏第五章,.
又再出現那個詔書,.
又有當地人,.
見到猶太人突然之間,.
奮勇地建造,.
不理那麼多,.
突然間建造,.
他們又寫信給當時的王,.
就是大理烏王,.
一樣是這樣說的內容,.
一樣是那個陳述,.
但是大理烏王是另一個角度去看,.
好,停一停,.
找一找有沒有那些典籍,.
去巴比倫省的寶庫裡面,.
找前朝,.
有沒有這些典籍,.
有沒有這些禦寺,.
他們是准許回去重建,.
如果有,.
就不是亂來的,.
這是第五章的以色列記,.
所以這12年,.
聖經就告訴我們,.

$^{841}$在這個,.
可能第五章都要說一點點,.
第五章第一節,.
那時先知哈該,.
那時就是大理烏王第二年,.
那時先知哈該和易多的孫子撒迦利亞,.
奉以色列神的名,.
向猶大和耶路撒冷的猶大人說勸勉的話,.
向猶大的耶路撒冷的居民,.
去講勸勉的話,.
即是預言,.
於是撒拉帖的兒子所羅巴伯,.
和約撒達的兒子耶穌亞,.
都起來動手,.
建造耶路撒冷神的殿,.
有神的先知在那裡幫助他們,.
這個幫助,.
相信不是用手用腳,.
而是用神的道去勸勉,.
令到他們的心靈能夠合一,.
大家看到神原來喜悅我們去恢復這個重建,.
神的殿,.
荒廢了,.
耽誤了,.
這個耽誤,.
表面上就是十二年前,.
這個亞特薩西,.
或者這個甘拜西二世的政策的改變,.
收緊了宗教上的政策,.
但是現在還沒有皇的新的政策出爐,.
為什麼猶大人會繼續自己這樣做呢?.
原來聖經也不需要再轉彎回角了,.
是因為有神感動先知,.
而先知來勸勉這些猶大人,.
要他們重建,.
這是神的心意,.
可能要問的問題就是,.
到底哈該,.
撒迦利亞,.
他們這兩位是不是已經是老人家,.

$^{881}$是七十年前被擄的時候,.
那些小朋友已經長大了,.
還是隨著他們在外地上來,.
而能夠成為先知呢?.
很慚愧地,.
聖經完全沒有交代這兩個先知的背景,.
他們是什麼時候出現,.
什麼時候蒙召,.
什麼時候得到神的默示,.
沒有說到他們的成長背景,.
所以有些人推測,.
他們可能是第一批老到外地,.
而他們被擄歸回的老人家,.
也就是說,.
他們快要入土為安了,.
可能說完這些先知預言之後,.
他們就會離開世界了,.
完全沒有交代他們什麼時候出現,.
什麼時候離開世界,.
只是聚焦在他們的默示上,.
現在我們就去看看,.
上一次我們也有機會看一點點的哈該書,.
有四個預言,.
或者說,.
你可以說是他們的對話,.
是神和哈該的對話,.
哈該書第一章,.
哈該的希伯來文就是節期的意思,.
大理烏王第二年,.
就是主前的520年,.
第二年的六月初一,.
很特別,.
有了exactly的日子,.
王的時間,.
這樣來到很清楚,.
他真的get到一個message,.
是from神的,.
然後他將這個message,.
傳遞給當時的猶大人,.
耶和華的說,.

$^{921}$直先知哈該,.
向猶大省長薩拉鐵的兒子所羅巴伯,.
和約薩達的兒子大祭司約書亞說,.
萬君知耶和華如此說,.
即是一個第三身講,.
這百姓說,.
建造耶和華殿的時候,.
這百姓,.
不是我的百姓,.
其實是神的選民,.
回到來已經是有心的了,.
五萬人,.
來重建,.
但是因為這個亞特薩西,.
逼令他們要停工,.
一停就十二年,.
去到大理烏王第二年登基了,.
等等,.
其實還沒有新的政策出爐,.
但是居然間,.
就是神的時間到了,.
興起,.
我不知道這個興起對不對,.
是不是一直都在群體裡面做先知呢?.
還是突然間出現呢?.
又似乎不是突然,.
是他們所認識,.
所尊敬,.
所聽從的,.
但是為什麼第一天不馬上出來呢?.
要十二年後才出來呢?.
這些又是沒辦法得到答案的問題,.
是那個形勢太強,.
還是現在有新王興起,.
所以神的時間到了呢?.
這樣都是比較自然的,.
因為亞特薩西離開世界,.
金拜西二世離開世界,.
大理烏出現了,.
可能是一個較為寬鬆的政策的王,.

$^{961}$果然是,.
所以有時神的工作都是這樣,.
第二節就是說,.
這百姓說,.
就是神給哈哥先知聽,.
百姓有這樣的想法,.
想法可能是很淺,.
慢慢成為一個信念,.
就很牢固了,.
很強了,.
想法就是一閃就過了,.
但是成為信念呢?.
十二年應該都是一個信念,.
這百姓說,.
建造耶和華殿的時候尚未來到,.
神很緊張那個聖殿,.
一直等待有人去興建,.
神仍然都不會用天使,.
用天君來興建,.
不會的,.
因為地球上是他創造人類,.
所以是用人類去做,.
或者說應該用他們的選民,.
猶太人上來重建,.
第三節,.
那時耶和華的話臨到先知哈該說,.
這殿仍然荒涼,.
你們自己還住天花板的房屋嗎?.
現在萬君之耶和華如此說,.
你們要審察自己的行為,.
你們殺的總多,.
收的卻少,.
你們吃卻不得飽,.
喝卻不得足,.
穿衣服卻不得暖,.
不攻前的,.
將攻前裝在破漏的囊中,.
在描述這十二年裡面,.
他們知道那個形勢,.
有了黃的詔書,.

$^{1001}$逼令停工,.
一停就停到大理烏的第二年,.
所以在這十二年裡面,.
什麼都做不到,.
但是發現,.
好像當地,.
一定是因為天災人禍,.
一定是天災帶來失修,.
帶來事半不是公配,.
帶來很多事都是事與願違,.
很多事都不順利,.
很多民生受到影響,.
他們入不敷出,.
很多很多的困難,.
他們接受這個命運,.
繼續要等待,.
等待那個時間,.
什麼時候重建聖殿,.
可能甚至乎已經沒有這個信念,.
可以重建聖殿,.
只是與患難共存就算了,.
沒有那種雄心壯志,.
相信一定是他們上去重建的時候,.
那種同心,.
那種大家的意象,.
那種心靈的火熱,.
才會上去,.
回到自己的故土,.
但是也發現,.
他們自己居住的地方,.
自己的國境,.
已經是頗破落的了,.
所以很多因素加起來,.
無論是面對的是現實,.
面對的是處境,.
面對的是人心的渙散,.
這是最大的問題,.
最難的就是人心,.
不能夠一起的同心合一,.
所以往往有時候,.

$^{1041}$我們不斷去研究教會,.
哪間教會成功呢?.
其實不需要研究的,.
是大家願意回到神的道裡面,.
大家願意放下自己的身份,.
自己的一些想法,.
願意放在神的真理底下,.
大家有這個單純的心,.
那就是教會的復興,.
只有這個方法,.
所以你不需要去複製別人,.
複製不了的,.
因為如果我們根本不具備這樣的土壤,.
不具備這樣的人心,.
你說什麼都是沒有用的,.
所以是要等待,.
這個等待過程當然就是.
牧者的祈禱,.
探訪,.
背後很多的眼淚,.
很多的積怨,.
來到如何能夠疏解,.
用神的道去安慰不同的弟兄姊妹,.
一定唯有是這樣的,.
所以你要看到,.
在這個現實底下,.
他們不是不想,.
是想的,.
不過就是繼續在他們自己.
面對現實的生活裡面,.
但是他們又知道,.
這一切都令到,.
不可以令到他們的內心,.
有真正的滿足,.
所以其實,.
大家共同都想到,.
我都很想教會好,.
但是就是,.
總是內心很多不滿足,.
所以記得不要用那些活動,.

$^{1081}$去麻醉那個內心的空虛,.
真正的滿足是,.
大家一起擁抱,.
追求,.
來到去傳福音,.
大家一起追求,.
來認識主耶穌,.
在有生之年,.
學習要似耶穌,.
那個心靈裡面的德行,.
讓聖靈來修剪我們,.
所以這裡就是,.
你會看到這個現實困難,.
現在有先知有神的時間,.
先知可能都講了很久了,.
但是現在有神的時間,.
就很不同了,第七節,.
萬君子耶和華如此說,.
你們要審查自己的行為,.
你們要上山取木料,.
建造這殿,.
我就因此喜樂,.
且得榮耀,.
這個就是Glory,.
不只是物質的富庶,.
那個豐富,.
現在是在講心靈上的豐富,.
心靈的豐富就是,.
先是為神的國度,.
沒錯的,.
為神的國度,.
這班是第一批上去重建的人,.
所以陸續,.
一路等待彌賽亞,.
真正的降臨,.
都有不同時代的參與,.
我們是第二批,.
等待第二次降臨的彌賽亞,.
我們不是第一次降臨的彌賽亞,.
是耶穌基督肉身回來的,.

$^{1121}$所以我們今天,.
要帶著這個信念和信心,.
來建造現在,.
應該叫教會時代的聖殿,.
其實就是我們這班人,.
彼得前書二章第七節,.
你們來到主面前,.
乃是像活石,.
被建造成為一個靈宮,.
奉獻耶穌基督所喜悅的祭,.
神所喜悅的祭,.
所以,.
那個殿就是我們這班人,.
我們每一個都是活的石頭,.
砌砌下砌砌下,.
成為一個spiritual house,.
屬靈的家,.
一個地方,.
所以,是說那班有活潑生命靈明的人,.
即是基督徒,.
所以,我們其實都需要預備我們的材料,.
我們有生之年,可能我們是教初信班,.
我們是做詩詩,.
我們是做聖樂,.
忠於我們的托付,.
那你就會慢慢發現,.
我們是會有一個聯絡性的,.
因為我們是有聖靈的,.
我們就去聯絡其他人,.
不知不覺就會成為了聖殿,.
那個spiritual house,.
彼得他真的與大相關,.
他自己是Peter,.
也是那個石仔,.
但是,他說,.
每一個基督徒去到主面前,.
都是一個活石,.
被建造成為靈宮,.
所以,我們求主幫助我們能夠,.
與弟兄姊妹的相交,.

$^{1161}$和人的相處,.
都能夠帶出這種黏性,.
那些英黎,.
將一塊磚一塊磚疊在一起,.
慢慢就是一幅牆,.
再將它連繫在一起,.
就是一個圍牆,.
慢慢就是一個建築,.
所以這些都是神一直在歷世歷代,.
都在找那個磚,.
我們是其中一塊磚,.
神仍然是這樣說,.
第九節,.
你們盼望多得,所得的卻少,.
你們修道家中,.
我就催去,.
這是為什麼呢?.
好像一個提問的方式,.
切問,其實答案是呼之欲出的,.
就是因為我的殿方樑,.
你們各人卻顧自己的房屋,.
這是萬君之爺和華說的,.
的而且確是事實的,.
真是有現實的生活,.
面對重建,.
面對自己的家園,.
可能被人侵佔,.
可能已經破落70年了,.
你想想,你現在去三流國的圍村,.
你去行山,.
你見到那些古老的圍村,.
都已經完全破落,荒廢,.
沒有人住,.
無論現在說的伏耕,.
都要很多的工程,.
你可想而知,70年之後,.
上回來,其實那個熱城被淋息之後,.
一想就是想自己的家庭,.
所以這些全部不是錯的,.
沒有否定的,.

$^{1201}$不過,仍然那個大前提是,.
因為我的殿方樑,.
你們各人卻顧自己的房屋,.
這是萬君之爺和華說的,.
特別強調萬君之爺和華,.
其實要打仗嗎?.
沒得打的,你是一個行省來的,.
你是一個波斯大帝國底下,.
但是萬君之爺和華,.
出現了好幾次,.
在哈該書,再一次提到他們這個民族,.
其實可能第二,第三代,.
根本都沒有讀過舊約,.
很熟悉那些出埃及,.
或者是烈王時代的戰爭,.
但是萬君之爺和華,.
在哈該書特別強調的,.
真正的那個leader,.
那個將軍,.
我們的背後,.
是我們這一位擁有實際權力的神,.
所以,這些說話都是有特別的用意,.
都是哈該奉神的名,.
是神自己親自去改的,.
不是一位很有慈愛的神,.
或者很信實的神,.
是的,完全沒有否定,.
不過要強調萬君之爺和華,.
似乎就是這個時候,.
這五萬多人,.
或者他們都會慢慢生小朋友,.
都是一個很minority,.
很少數的少數,.
即是你做什麼,.
你回來自己那個地方都這麼破落,.
當地的人這麼大聲夾鶴,.
上訪來到去,.
逼令你停工,.
你有什麼power在你背後呢?.
哈該就告訴我們,.

$^{1241}$有這一位萬君之爺和華,.
這個稱呼,.
這個卡片,.
就已經夠了,.
第十節,.
所以為你們的緣故,.
天就不降金爐,地也不出土產,.
我明乾旱臨到地土,.
山崗五谷伸走和遊,.
並地上的出產,.
人民生畜,以及人手,.
一切勞碌得來..
即是將各樣的東西,.
都來到去涵蓋,.
即是說,.
真是一種沒一口好吃的困難,.
我們這幾年真是,.
去到真是疫情帶來的.
民不聊生是很慘的,.
所以稍微有些開放,.
即是大家的內心就很舒暢,.
這個是真是需要的,.
但是我們不要只是看到地上,.
那個平面,.
要看到一個屬靈的興旺,.
要看到人心歸向神的興旺,.
我們要追求的不是,.
嘩,真是好了,通關了,.
我可以去日本了,.
不是錯,.
我們仍然是一個平衡的世界,.
不過我們不要只是看到這件事,.
我認識有人,.
由年初到現在,.
其實都未夠二月,.
已經去了幾次旅行,.
沒有批判的,.
不過在我們的世界,.
在聖經茶經班的弟兄姊妹,.
你的世界,.

$^{1281}$是不是只有地上的活動呢?.
地上個人的你的需要,.
你的享受,.
或者可能是你的工作,.
或者是你的公務等等呢?.
有沒有想到,.
一切原來背後是有神的時間,.
令到你去到地上,.
去到地上,.
背後的時間,.
令到這十二年,.
這麼荒涼,.
這麼慘,.
這麼困苦,.
是因為我的電荒涼,.
這個是,.
真是好像神探大戰,.
這樣的橋段,.
你要看到別人看不到的東西,.
你要看到的就是教會,.
是你最大的,.
那個生命的投放的投資,.
希望神的電,.
希望這個地方有更多人認識神,.
得到救恩,.
得到平安,.
得到家庭的和睦,.
小朋友從小就在教會很開心地長大,.
他認識神,.
他有真理在他的內心,.
這些是我們一生,.
越是信主內的,.
就越是應該著緊,.
我仍然說,.
我真的沒有批判,.
我只是說,.
我仍然說,.
我真的沒有批判,.
每個人為自己所做的向神負責,.
但是如果我們整個的生命,.

$^{1321}$哇,通關了,.
我有錢,.
我去哪裡都可以,.
去日本一個月,.
然後再去其他地方,.
沒有任何的貶義,.
不過,.
如果這本書告訴我們,.
黑蓋書告訴我們,.
我的電荒涼,.
我去哪裡都可以,.
去哪裡都可以,.
而你沒有忘記,.
你忘記了神的電,.
教會這麼需要參與,.
為什麼我們會當看不到,.
或者說看到,.
而我們只是一種犬儒式的謾罵,.
而不是真的下手下腳的,.
捲起衣袖的,.
我就是回來做了,.
我就是回來幫忙,.
有什麼做,.
我就做了,.
低頭默默的去做,.
這個很重要的,.
發現了問題,.
不是很厲害,.
發現了問題,.
而面向那個問題,.
一起去解決那個問題,.
這個才是我們需要.
今年一起去學習的地方,.
第十二節就是,.
這個講論帶來的.
那個結果,.
就是有時間的,.
那個的密使,.
是六月初一,.
距離七月初一一個月,.

$^{1361}$第十二節,.
那時撒拉鐵的兒子.
所羅巴伯和約撒達的兒子.
大祭司約書亞,.
並盛夏的百姓.
都聽從耶和華他們神的話,.
和先知哈該奉耶和華他們神差來所說的話,.
很重要的那些字眼,.
是奉神的名,.
差遣而來,.
所羅巴伯和約撒達的兒子.
大祭司約書亞,.
差遣而來,.
所說的話,.
就是說,這是耶和華如此說,.
不是哈該,.
不是撒迦利亞,.
那個個人的智慧,.
見識,.
或者是老人家,.
我食鹽多過你們食米,.
聽我說吧,.
那些都不是權威,.
權威是.
奉神的名,.
所說的話,.
這兩位領袖就.
接受了,.
為什麼呢?.
百姓也在耶和華面前.
傳敬畏的心靈,.
所以,教會.
很需要大家.
去尊重神的道,.
聆聽神的道,.
聆聽神的指引,.
以致大家.
一起有這個方向,.
去向神.
效忠,.

$^{1401}$這才是教會.
最根本的問題,.
當然,.
有些事情,.
除非聖經很明顯,.
否則,我絕對不敢說,.
耶和華透過我.
這樣說,任何人都要.
聽我這樣說,我絕對沒有.
這樣的權,.
可以分析,可以按過去的.
經驗,告訴你,.
這樣做法會好一點,.
但絕對不是耶和華如此說,.
正正因為這樣,.
不夠堅定,.
不夠這麼均勻,.
很多人眾說紛紜,.
教會如果眾領袖.
都眾說紛紜,.
這個說去A那條路,.
這個說去B那條路,.
教會就不需要.
想了,.
一定是.
不會有方向了,.
有的話,都要經過很多角力,.
這也是問題所在.
的根本,.
所以這些經文其實.
是很重要的,.
大家能不能夠.
對神的.
話語尊重,.
尊重不只是欣賞,.
聽聽很開心,.
而是真是大家身體力行的,.
完全的.
放下自己的自我,.
一起這樣走上,.

$^{1441}$這個過程,是需要.
我們大家都全謙卑的心靈,.
因為再沒有.
哈該,薩加利亞,.
這些耶和華如此說,.
如果有,當然說我都想跟他,.
如果有一位領袖.
在香港說,.
耶和華透過我說就是了,.
那我也想跟他,.
但是,如果有這樣的人.
這樣說,我就寧願你不要跟他,.
因為.
新約不是這樣運作,.
不是這樣靈恩式的,.
一種先知運動.
的運作模式,.
而是仍然都是.
大家謙卑在神面前,.
大家一起去透過神的真道,.
我們等候聖靈,.
等候神的.
指引,.
大家是同心的,.
單純的,謙卑的,.
讓那件事,.
最終都要有一個最後的方案的時候,.
大家都願意.
去服從的,這樣才是.
教會前進的.
根本,.
如果不是,其實就很難,.
現在就是他們,.
那個政治還沒改變,.
第四章那裡說,.
大利烏還沒頒布,.
甚至乎,.
當他們建的時候,.
就立即有當地人.
又再一次上本.

$^{1481}$控告尤大人的建造,.
一樣這樣的.
內容去質疑,.
去攻擊,.
甚至乎,.
很想左右新王大利烏,.
都是好像以前那樣,.
來阻撓.
這些工程,.
但是,神的.
先知,神的道,.
是優先過.
這些反對,.
神的時間來了,.
先知是配合,.
去說這些,.
不是風涼話,不是離地,.
相反就是,.
能夠指出這十二年,.
在尤大,在耶路撒冷,.
那個困難,.
困苦,.
人心的虛怯,.
原因是,不是天災,.
是呀,天災是.
那個表面的現象,.
真正的就是,神不起悅,.
你們回來,.
但是不建造,.
這個就是那個,.
根本的解決的問題,.
將整件事,.
來去撥亂,.
來去只需要他們重建,.
就行了..
你想想,想像得到,.
先知,.
能夠用.
一篇的訊息,.
就可以扭轉整個局勢,.

$^{1521}$你可想而知,.
神在哈該和.
撒迦利亞的心靈裡面,.
做了多少的工程,.
感動,.
而他們能夠講出這些,.
大家都認為,.
好像不是時間,.
不是這個時機,.
的一個訊息,.
但是,.
聖經告訴我們,.
第十三節,.
耶和華的使者哈該,.
奉耶和華差遣對百姓說,.
耶和華說,.
我與你們同在,.
其實,.
不是很深奧的東西,.
但是,.
他們聽下去,.
感覺到憤外的有力量,.
再一次,.
神,.
耳脈來,.
神和你們現在建造的人,.
同在..
所以,第十四節,.
耶和華激動猶大省長.
撒拉鐵的兒子.
所羅巴伯和約撒達的兒子.
大祭司約書亞,.
並盛夏之百姓的心,.
他們就來為萬軍之耶和華.
他們神的殿作供,.
開始.
供呈,.
就是.
六月二十四,.
即是廿三日,.

$^{1561}$即是說,.
這篇訊息講完之後,.
繼續在不同的群體當中,.
傳播了..
神就再激動.
所羅巴伯,.
約書亞,.
他們就做他們的工作了,.
就是聯絡那些領袖,.
分發出去,.
來感動.
其他的百姓,.
大家一起,.
因為神的時間到了,.
神的Calling也到了,.
要我們去做了,.
所以,這是大利烏王第二年.
六月二十四日,.
興宮建造神的殿,.
廿三日就開始,.
這個訊息,.
有時我們大家都明白的,.
上個禮拜講完什麼道,.
我們都因為太多了,.
都會忘記了,.
廿三日,.
更加可能會忘記,.
但是,相信這一篇訊息,.
又不是什麼很強烈的責備,.
只是列舉了,.
那個現況,.
那個現實,.
和猶大人,.
現在面對的,.
好像排除了那些,.
政治上的困難,.
為什麼你不講呢?.
當地人的反對,.
又不講呢?.
為什麼只講我們,.

$^{1601}$日不呼出,.
沒一口好吃,.
心靈困苦,.
為什麼只講這些呢?.
為什麼只講我們猶大人的情況呢?.
零勝的十二年的荒涼,.
大家都沒有心在,.
上去獻祭了,.
因為都沒有那個制度在運作,.
怎麼獻呢?.
上到去那個地方都荒涼的,.
四周都圍住了,.
好像那些警戒線那樣的,.
你怎麼去拿開那些警戒線,.
去獻祭呢?.
接著說你非法獻祭,.
接著又被人抓了,.
你想一下,.
廿三日,.
可以帶來一個復興,.
聖經就很平鋪直敘,.
但是索羅巴伯和約書亞,.
首先自己被激動,.
接著就去不斷的,.
可能開一些研討會,.
不只是研討,.
而是鼓勵百姓,.
一齊的來做,.
就是說兩位領袖,.
一個是祭司的,.
一個是行政的,.
就來到去,.
各自開始去謀劃,.
怎樣能夠再激動.
其他的百姓,.
來到去,.
因為這個message,.
哈該的message,.
來到去,.
大家一齊去建造,.

$^{1641}$所以都幾重要,.
那個訊息,.
基督教,.
真是著重那個訊息的,.
著重那個message,.
所以一定要透過聖經,.
來產生的message,.
所以港台,.
不要講例子,.
講故事,.
講看什麼電影,.
不是講聖經的真理,.
活在我們這個群體的應用,.
所以,.
教會的牧師,.
和神學院的講師,.
是有不同的,.
教會的牧師,.
牧者,.
是真的知道弟兄姊妹的情況,.
而講一些和弟兄姊妹有關連的道,.
因為如果你是一個教師,.
神學院的,.
他是專研聖經,.
他講出來,.
沒錯,.
一定是有足夠的時間的準備,.
和他神學的訓練,.
和原文,.
或者他有很多其他的參考,.
他的時間,.
他投放的精力,.
一定是會比駐堂的牧者好,.
但是,.
兩者是應該相輔相成,.
不可以是彼此的,.
來有抵觸,.
不需要抗拒,.
一些學者的教導,.
但是,.

$^{1681}$更重要的是,.
駐堂的牧者的模樣,.
那個的關心,.
那個的跟進,.
和弟兄姊妹面對的困難,.
的陪伴,.
共同的渡過,.
等等,.
往往這些不是豐收野味,.
而是家常便飯,.
是,.
兩者是不應該有任何的衝突,.
但是,.
一定是會帶來,.
有不同的.
恥笑比掌,.
所以其實是要互補,.
而不是排斥,.
所以這裡告訴我們,.
所羅巴伯他被這一番說話的.
勸勉,.
某程度是一個責備,.
為什麼,為什麼,.
是神問,.
好像第三者,.
看這班選民,.
這班人本來很有心,.
是神激動他們回去,.
重建的,.
但是在那個字裡行間,.
神說,這百姓說,.
就是表達好像有種憤怒,.
但是,.
所羅巴伯,.
其實真的一直都有心,.
但是,好像,.
空有那個理想,.
但是群體在停擺,.
沒有人理睬你,.
大家看到那個形勢,.

$^{1721}$但是所羅巴伯,.
需要的就是那個message,.
需要的,.
所以他不是利用哈該,.
有些人很容易的,.
利用那個牧者,.
講一些說話,.
要做他自己的事,.
這些是野心,.
所羅巴伯不是,.
他同樣都是身受其害,.
他都是在那個等待裡面,.
那團火,.
還沒熄滅,.
當哈該來到燃點的時候,.
就馬上再出來,.
好像很簡單,.
23日,整個群體,.
在耶路撒冷的,.
馬上被勒著,.
勒著回來重建,.
好像都不理,.
新王登基之後會怎樣,.
都不理,.
當地的那些原居民的反對,.
等等,.
大家一起去做,.
好像有些很誇張,.
不理當時的政治體系,.
但是,.
他們現在就是為神去做這個聖殿,.
所以盼望我們,.
在現在還是,.
宗教自由,.
一國兩制,.
我們不要小看自己的福音工作,.
更加不要自己,.
自絕於那個福音,.
我們要全福音,.
不要理其他地方,.

$^{1761}$其他內地,.
18歲以下,.
不准聽福音,.
先不要理那些東西,.
我們現在可以做的就做,.
到時才算,.
現在是別人真的給我們一國兩制,.
真是宗教自由,.
我們就做,.
不要因為內地不准,.
我們就好像不斷批判,.
去罵,去反,沒有意思的,.
現在我們要做的就是,.
我們要和那個福音,.
共存的去傳福音給香港人,.
給來到香港的,.
可能是新香港人,.
可能是在外國,.
來香港居住的,.
譬如有些難民,.
非洲的難民,.
在紅磡一帶,.
都有這些非洲的難民,.
我們要在這個地方,.
既然神要我們在這裡,.
我們就好好地傳福音,.
很多人移民,.
不要緊,神有感動他們去到那些地方,.
去到那些學中心,.
我們在這裡的,.
我們繼續傳福音給香港人,.
我們就不要分哪些不是香港人,.
哪些不是我們自己的,.
你這樣分,.
你就是變成連你的福音的本質,.
你都錯了,.
你根本就,.
你自己的心靈都出了問題了,.
居然你還要緊要神去分了,.
哪些人是需要福音,.

$^{1801}$哪些人就不是香港人,.
我們哪有可能有這樣的資格去分呢,.
任何的人,.
我們有機會都跟他講福音,.
無論他是貧富,.
都應該要跟他講福音,.
讓他能夠成為.
福音改變的那個動力的對象,.
福音好像一頭猛獸,.
當然這樣誇張了,.
好像在一個籠裡面的獅子,.
你困住他,.
那隻獅子沒有用的,.
小朋友可以潑水,.
可以拿支棍子撩他,.
但是你打開籠裡,.
那隻獅子就會出來,.
那個力量就會出現,.
所以其實我們是要宣講福音,.
用我們的生命去見證福音,.
福音的能力就會改變這個社會,.
改變香港的生態,.
所以這裡就是23日一篇信息,.
總共兩章聖經有四篇信息,.
現在就答九,.
先喝口水,.
有些少測試,.
這個就是哈該書的特色,.
只是四篇信息,.
但都帶來很大的改變,.
我們試試還有些少時間,.
我們看第二篇的信息,.
都是有些關聯的,.
由24日去到7月21日,.
7月21日是什麼日子?.
駐彭捷的最後那一天,.
有一段經文和大家參考,.
我們看看歷代志下,.
第七章,.
第七章的第一節,.

$^{1841}$歷代志下,.
所羅門祈禱以筆,.
就有火從天上降下來,.
燒盡凡際和別的際,.
耶和華的榮光充滿了電,.
這個就是第一聖殿的重大開幕典禮,.
因耶和華的榮光充滿了耶和華的電,.
所以祭司不能進電,.
那火降下,.
耶和華的榮光在電上的時候,.
以色列眾人看見,.
就在鋪石地,.
俯伏叩拜,.
稱謝耶和華說,.
耶和華本為善,.
他的慈愛永遠長存,.
這個也是以色列人的共同認訊,.
對於眾物主的慈愛,.
他的善,.
猶太人都是這樣理解的,.
第四節,.
王和眾民在耶和華面前獻祭,.
所羅門王用牛二萬二千,.
羊十二萬,.
諸如此類,.
第七節,.
所羅門因他所做的同壇,.
繁祭,掃祭和滋遊,.
因為這麼多,.
兩萬多,.
只是供牛,.
便將耶和華電前院子當中分別為聖,.
即是說,.
其他地方,.
即是我們的廣場,.
都分別為聖,.
都可以用了,.
在那裡獻繁祭和平安祭,.
生的滋遊,.
第八節留意了,.

$^{1881}$從哈瑪口直到埃及小河,.
即是說,.
他的國土,.
所有的以色列人都聚集成為大會,.
首節七日,.
首什麼節呢?.
似乎就是駐彭節了,.
第八日設立嚴肅會,.
行奉獻壇的禮,.
七日,首節七日,.
接著去到七月廿三,.
王就遣散眾民了,.
果然在最大那天,.
就是七月廿一,.
最大的獻祭的嚴肅會,.
駐彭節最後那天,.
或者叫第八日,.
去到七月廿三,.
就完了,.
即是曲終人散,.
王就遣散眾民,.
他們因見耶和華,.
向大衛和所羅門.
與他們以色列所思的恩惠,.
就都心中喜樂,.
各歸各家去了..
為什麼引這段聖經呢?.
因為哈該,.
他們重建聖殿的時候,.
就是六月廿四開始,.
去到七月廿一,.
就有一個寓言又再出現了,.
所以有時神都幾風趣的,.
第二個聖殿重建了,.
即是被擄歸回重建的第二聖殿,.
又是在七月廿一,.
有神的聖經,.
又是在七月廿一,.
又有神的聖經,.
又是在七月廿一,.

$^{1921}$又有神的聖經,.
又是在七月廿一,.
又有神的聖諭..
我們看第二章,.
第二章這裡,.
七月廿一日,.
耶和華的華臨到先知哈該說,.
你要曉遇猶大省長.
撒拉鐵的兒子.
所羅巴伯,.
和約撒達的兒子大祭司.
約書亞,並盛夏的百姓說,.
你們中間全流的,.
有誰見過這殿.
從前的榮耀呢?.
從前的榮耀,.
就是我們讀的第七章,.
就是那個Glory,.
天火焚燒濟生,.
是很奇特的神蹟,.
即是說,.
那個Glory是讓整個民眾.
是感覺到溫暖,.
感覺到快樂,.
因為見到神月立這個王,.
而這個王有智慧地.
統領著整個國家,.
從哈馬口.
直到埃及小河,.
你開賞而知,.
你是其中的國民,.
你就感覺到那個幸福,.
這個是屬靈的通關,.
不是只是普通的通關,.
而是大家感受到那種歡樂,.
因為見到Glory,.
神的榮耀,.
那個榮光,.
是燃燒那個濟生,.
所以接著就是,.

$^{1961}$你們中間誰見過?.
這句說話,.
這一節,.
很多解經家就認為,.
當然沒有十足的證據,.
就說哈該,.
其實就是第一批小孩被擄的時候,.
現在七十多歲的老人家,.
即是說他自己見過.
舊的聖殿,.
所羅門的聖殿,.
當然這是一個猜測,.
是也不出奇的,.
即是說,.
說完之後,.
七十多歲的可能不止,.
譬如達爾利,.
那些可能已經離開世界了,.
不是,.
不是肯定離開世界,.
因為已經去到大理烏王的第二年,.
你可以想像得到,.
這個的震撼,.
這個叫他們緬懷,.
昔日的光輝,.
所以有些人看完之後,.
就會哭,.
這個榮耀,.
現在你們看著如何呢?.
豈不在眼中看如無有?.
真是沒得比,.
所以神就說,.
你不需要這樣比,.
將來有一個更大的,.
更榮耀的,.
可想而知,.
你可以想像得到,.
主耶穌成為那個聖殿,.
將來甚至乎啟示錄說,.
找不到聖殿了,.

$^{2001}$因為神和羔羊,.
就是神的殿了,.
你可以想像得到,.
將來我們去到永恆的榮耀,.
那種永遠沾養神的榮光,.
那種快樂,.
那種滿足,.
那種不會衰敗的滿足感,.
是永恆地存在,.
所以現在他們看到,.
就一定有比較,.
真是都不知道有什麼好比,.
根本沒得比,.
第四節,.
耶和說,.
「羅伯伯,雖然如此,.
雖然如此,.
肯定不需要比,一定不夠,.
你當剛強,.
約撒達的兒子大祭司約書,.
你也當剛強,.
這地的百姓,.
你們都當剛強作功,.
因為我與你們同在,.
這是萬軍之耶和華說的.」.
這是第二篇的訊息,.
由六月初一講完,.
去到七月廿一,.
駐彭節,.
即暫時,.
其實六月廿四開始開工,.
七月初一就有.
催國節,.
七月初十救贖罪日,.
七月十五就是.
駐彭節,要守八日,.
所以其實整個月,.
好像我們放完聖誕,.
沒多久就倒數,.
新年,.

$^{2041}$接著二十二號就初一,.
其實初一之前,.
由年廿八已經很忙,.
一直到初三,.
然後初四還是假期,.
其實很多事都做不到,.
但其實不要緊,.
節期是要喚起人心,.
去朝見神,.
去敬拜神,.
其實這個才是真正力量的根本,.
所以當他們明白了之後,.
神一而再叫他們.
剛強作功,.
起了名,.
叫所羅巴巴,.
約撒達的兒子,.
約書亞,.
或者耶書亞,.
來到去剛強,.
眾百姓剛強,.
為什麼呢?.
因為我和你們同在,.
是萬軍之耶和華,.
又再一次出現了,.
這樣說的,.
第五節,.
「這是召著你們出埃及,.
你們立約的話,.
那時我的靈住在你們中間,.
你們不要懼怕.」.
好像出埃及的陣容,.
是那個處境,.
出埃及是.
整個國家要逃亡,.
要走去一個新的國土,.
還有追兵,.
埃及在追來,.
神說這樣的規模,.
但現在他們是.

$^{2081}$回到去重建,.
即是新的國土,.
回到去重建,.
即是新的出埃及,.
開道路,.
開江河,.
使他們行乾地,.
現在他們是回歸重建耶路撒冷的聖殿,.
所以,.
那個剛強背後是.
神和他們同在,.
所以如果一個弟兄姊妹,.
他能夠每個禮拜.
很忠心,.
很謙卑,.
在他們面前得到.
心靈的滿足,.
你可以肯定,.
他的工作,他的家庭,.
都會很和睦,.
他的心靈都會很穩健,.
很安,.
安逸不是負面的,.
是他裡面.
很有那個peaceful,.
他就會與人為善,.
因為耶和華本為善,.
他的慈愛.
是永遠常存,.
所以你會感受到那份善,.
你會感受到那份平安,.
所以,.
如果我們香港有.
很多基督徒的領袖,.
在不同的領域,.
都能夠有那份善,.
去改變那個世界的.
殘暴,.
暴戾,.
就很少了這些問題,.

$^{2121}$所以如果香港的罪惡率高,.
某程度上應該是.
基督徒的靈命不好,.
應該是這樣看的,.
如果不是,.
怎樣去改變世界的人呢?.
所以這裡說,.
神與我們同在,.
盼望感受這個社區.
先行,這個社區.
充滿了一份和睦,.
人們來到這個地方,.
不會嘲笑我們,.
會有一種敬畏我們,.
因為.
見到這個地方有神,.
有愛,.
諸如此類,.
這是事實,.
真是應該這樣被人見到的,.
所以,.
要留意,.
神說在出埃及的時候,.
他的靈.
住在他們中間,.
當然透過摩西,.
當然透過律法書,.
當然有些神蹟,.
但原來有神的靈盤石,.
站在哥林多前書,.
普羅也是這樣理解的,.
有時我們很難理解.
他們這樣靈異的解經,.
臨前的第十章,.
弟兄們,.
我不願意你們不曉得,.
我們的祖宗從前.
都在雲下,.
就是那些雲處火處,.
都從海中,就是紅海經過,.

$^{2161}$都在雲裡,海裡,.
受洗歸了摩西,.
雖然是那個領袖,.
摩西帶領他們.
離開埃及,.
過紅海,成為一個.
新的民族,.
就是神的選民,.
其實過紅海.
是怎樣的洗禮呢?.
不過,好像是經過生和死,.
好像紅海的打開,.
他們好像不是.
如沐春風,.
是好像受浸,.
歸入神的名下,.
即是神第二次.
拯救他們這樣,.
保羅就很有那種.
靈異的去看.
過紅海,.
第三節,.
並且都吃了一樣的零食,.
屬靈的食物,.
嗎哪就肯定是,.
可以是神跡地,.
六日有的,.
安息日就沒有的,.
出去撿都不會變臭的,.
這些的零食,.
即是說,一路都是神的.
恩典去牧養.
這個民族,.
二百萬人,.
四十年,.
由他們被神的咒詛,.
他們在曠野,.
都可以繼續的存活,.
如果不是神的恩典,.
能力,.

$^{2201}$怎可能繼續呢?.
所以,.
保羅說,.
也都學一樣的零水,.
所以我們不要亂解經文,.
我們自己製造的神仙水,.
雙氧水,.
不是,.
是零水,.
說的是神.
透過一些盆石,.
出水,.
或者在一個叫馬拉的地方,.
那裡的水是苦的,.
神就叫摩西,.
摩西那棵樹,.
放一塊木下去,.
那些水就會變甜的了,.
你就會見到,.
保羅不斷將這些神蹟,.
看為是零食,.
零水,.
所學的是出於.
隨著他們的零盆石,.
那盆石就是基督,.
追隨著這一班的人,.
這個民族的,.
是基督,.
保羅這樣理解經文的,.
但他們中間多半是神不喜歡的人,.
所以在曠野倒閉,.
這些事都是我們的尷尬,.
叫我們不要貪戀惡事,.
像他們那樣貪戀,.
也不要拜偶像,.
像他們有人拜的,.
如經上所記,.
百姓坐下吃喝,.
起來還撒,.
或者還撒,.

$^{2241}$這是出埃及的三十二章,.
神透過利未之派,.
將他們殺,.
拜金流毒,.
在偶像面前,.
做出不好的事,.
你能想像到,.
保羅這樣理解,.
四十年的出埃及記,.
民數記的曠野,.
是有神的零伴隨,.
甚至乎是有基督的伴隨,.
這樣理解舊約的,.
作者,.
使徒他自己,.
哈該書這裡,.
都是一脈相承,.
都是這樣理解,.
神的零住在他們中間,.
神自己住在他們中間,.
即是說,.
那個民族原來,.
一直都有神在,.
這個就很值得我們去想,.
我們活在香港,.
我們不只是看到,.
那個施政好不好,.
那個社會民生好不好,.
我們要去想,.
那個社會民生好不好,.
我們要看到神在這個社會,.
神在我們的教會,.
神在我們大家的心裡面,.
我們在做的是神國度的事情,.
所以,.
不是這麼簡單,.
只是地面,.
表面的事情,.
我們要有一種,.
人家看不到的那種,.

$^{2281}$屬靈的眼光,.
這個就最難,.
亦都是我們大家要謙卑的去尋求,.
這個就是聖經這樣說,.
第六節,萬君之耶和華如此說,.
過不多時,.
我必再一次震動天地,.
滄海與旱地,.
我必震動萬國,.
萬國的珍寶,.
必都運來,.
我就使到這殿滿了榮耀,.
這是萬君之耶和華說的,.
萬君之耶和華說,.
銀子是我的,.
金子也是我的,.
這殿後來的榮耀,.
必定大過先前的榮耀,.
不是物質去衡量,.
不是只是做一間比所羅門更大的,.
是說那個,.
相信是靈性方面的宏偉,.
那個盛載,.
屬靈的層面,.
相信一定是說主耶穌,.
祂成為那個,.
好像私恩座那樣,.
來到這個地上,.
祂是在這個地上,.
這是萬軍之耶和華說的,.
出現最多的就是萬軍的耶和華,.
他們現在去建造,.
其實一定是不容易的,.
暫時放下家庭的事務,.
大家同心的去建造,.
很快就要面對那個壓力,.
就是來自當地人那個迫擊,.
那個上告,.
是否第二次查封呢?.
這個不知道的,.

$^{2321}$我們後人當然知道,.
但是他們就要面對這件事,.
但是當他們建造這個聖殿,.
神的平安就會臨到那個地方的人,.
所以日後神會賜下平安,.
會賜下那個恩典,.
但是建造的人,.
在那個過程,.
他們是未那麼快感受到這份平安的,.
但是很重要的就是,.
銀紙是神的,.
金紙也是神的,.
所以求主讓我們能夠明白,.
不是要去競爭哪間教會大一點,.
而是神給我們的份是多少就多少,.
知道自己的限制,.
但是要力求的,.
我們是要去傳福音,改變人心,.
我們不是要追求成為全香港最大的教會,.
我們不是需要這樣做,.
每個人有他們自己的神給他的份,.
感宣的份是什麼呢?.
是我們大家一起去尋求,.
一起去等候,.
一起去努力,.
但是神的資源是無限的,.
我們在我們的份裡面去做,.
神就一定會有祝福,.
最大的一定是神會賜平安,.
建造的人和建造好之後,.
會影響後代的平安,.
所以我們真的盼望,.
相信我們在賊景班不斷的教導,.
弟兄姊妹已經感受到神的話語的能力,.
我們保持這份單純,.
教會一定會有復興,.
復興不是一定要用量量化,.
是群體是否忠心,.
群體是否忠心,.
群體是否擁戴耶穌的救恩,.

$^{2361}$這些才是量度的標準,.
不是數字不斷攀升,.
分堂開到幾十間,.
\newpage



\section{哥林多後書 8:9}
\label{sec:4sQReCkRl1k}
\textbf{2023-01-28 晚堂《無比恩典》哥林多後書 8\_9 講員:陳一華牧師}
\newline
\newline
連結: \href{https://youtube.com/watch?v=4sQReCkRl1k}{\texttt{https://youtube.com/watch?v=4sQReCkRl1k}} ~~~~ 語音日期: 2023-01-31
\newline
\newline
\hyperref[sec:mXSWAPPoVUc]{\small{< < < PREV SERMON < < <}}
~
\hyperref[sec:index_chronic]{\small{[返順時目]}}
~
\hyperref[sec:index_scriptual]{\small{[返順卷目]}}
~
\hyperref[sec:65Ibdpx_Vbg]{\small{> > > NEXT SERMON > > >}}
\newline
\newline
哥林多後書 8:9
\newline
\begin{longtable}{cl}
\hline
\hline
章節 & 經文 (和合本修訂版)\\
\hline
8:9 & \begin{tabularx}{0.7\textwidth}{X} 你們知道我們主耶穌基督的恩典：他本是富足，卻為你們成了貧窮，好使你們因他的貧窮而成為富足。 \end{tabularx} \\ \\
[1ex]
\hline
\hline
\end{longtable}
$^{1}$弟兄姊妹晚安.
今晚大家會得到兩個意料之外的收穫.
第一個收穫就是會見到牧師的真人面目.
三年沒有見過了.
今晚晚堂的弟兄姊妹都有這個福氣.
你們可以一睹牧師的真面目.
這是第一個意外的好處.
第二個意外就是如果年初一有新春崇拜的弟兄姊妹.
你有第二個收穫就是七天之內聽到牧師兩次講道.
是不是很難得呢.
要不就不聽 我一聽就連續七天之內聽兩次.
所以今晚牧師跟大家分享一個題目.
牧比恩典.
新年初我們大家都去拜年.
拜年的時候都很多祝福的說話.
要不就說新年蒙福 新年快樂.
又或者祝大家身體健康.
有時候我都聽到有些人的祝福是說什麼呢.
是福杯滿日 你有沒有聽過.
福杯滿日 又或者主恩滿滿.
又或者主恩滿日.
有時候我們都聽到人家跟你來祝賀的時候.
都用這些說話的.
所以今晚我們趁這個新春頭.
我們去思想一下主恩滿日 主恩滿滿.
經常都是怪口唇邊講到主的恩典.
到底是什麼樣的恩典呢.
正如保羅在剛才的經文告訴我們.
臨後第八章第九節.
他說你們知道.
都打出那一節經文出來.
他說你們知道我們主耶穌基督的恩典.
我們會問到底這是個什麼樣的恩典呢.
這個恩典又怎樣得著呢.
得著這個恩典的時候會有什麼好處呢.
所以我們今晚就循著這個思路.
來思考這段經文.
經文告訴我們原來耶穌基督的恩典是怎樣的.
他說本來祂是富足的.
卻因我們成為貧窮.

$^{41}$而好使我們因祂的貧窮.
我們可以成為富足.
你聽到這樣的描述.
有一點點像交叉異形換影的感覺.
首先耶穌基督的富足因我們成為貧窮.
而我們因耶穌基督的貧窮成為富足.
像不像我們小時候看電視劇的片集.
打排球的交叉異形換影.
有沒有這樣的印象呢.
好像你們沒有什麼印象.
你們都是很年輕.
讀起這篇經文的確有這樣的感覺.
這篇經文非常寶貴.
為什麼呢.
因為這篇經文告訴我們.
保羅說你們知道主耶穌基督的恩典.
我想在這裡特別強調.
這個恩典是一個基礎性的恩典.
即是上帝給我們眾多眾多的恩典.
又或者怎樣主因滿滿.
主因滿溢的裡面.
你要記得原來這篇經文是一個基礎性的恩典.
如果沒有了這篇經文的話.
我們發覺我們得到的恩典都不是那麼鞏固或者那麼確實.
但當我們有了這個基礎性的恩典.
就是這篇經文保羅所強調之後.
頂尖的門利擁有這個基礎性的恩典.
在上面再建造的時候.
再得著神的恩典的時候.
這個恩典是會牢牢的.
是穩固的.
所以今晚這節聖經是十分之重要.
因為它是一個基礎性的恩典.
怎樣基礎性的恩典呢.
我想和大家分享.
原來這個恩典有兩個部分.
頭一個部分是耶穌基督降生成人.
這是第一個恩典.
這個恩典第一個部分是耶穌基督降生成人.
當耶穌基督降生成人的時候.

$^{81}$就是說他原本是富足.
卻因你我成為了貧窮.
所以在這個過程裡面.
我們要記得有四項的轉變是很重要的.
第一個轉變是一個什麼樣的轉變呢.
就是位置的轉變.
這個位置的轉變就是說.
耶穌基督他降生成人的時候.
他要從天上補助這個位置.
而降低降下下來.
不斷地下降.
而最後降生最低的位置.
就是他出生在馬槽裡面.
所以各位觀眾.
這個過程是非常不容易.
是非常不簡單的.
正如你讀到菲勒比書.
保羅那裡告訴我們.
耶穌基督本有上帝的形象.
有印象嗎.
不與與神同等為強敵.
一路降一路降.
最後他取了奴僕的形象.
各位聽眾.
我想跟大家分享.
這個位置的轉變.
對耶穌來說.
這個抉擇是要付出很大的勇氣.
或者我簡單舉個例子來說.
如果有一些做生意的朋友.
他做生意做得好的時候.
好像快了一達.
馬上買了一間豪宅.
二千多呎.
開不開心.
快不快樂.
很快樂.
容易不容易適應.
很容易適應.
但是一個疫情.

$^{121}$生意失敗了.
怎麼辦呢.
資不抵債.
在這個時候.
他要連他住的房子都要.
拿出來拍賣.
頂智滿門.
他找到一間房子.
另一間房子只有二百呎.
我想問你們.
由二千呎一下子搬去住.
二百呎的房子.
容易不容易適應.
不容易適應.
很辛苦.
叫苦連天.
我們住房子.
都有這樣的經歷的時候.
何況我們的主耶穌.
他更加艱難的.
就是要將天上的位置.
捨棄.
而取了一個最低的地上位置.
你可以想像過程當中是多艱辛.
但是耶穌甘心這樣做.
因為祂知道.
這是第一個步驟.
第一個轉變.
我們再看第二個轉變.
第二個轉變是一個什麼轉變呢.
是一個身份的轉變.
頂智滿門.
耶穌降生成人的時候.
祂要將祂自己尊貴天上的身份放下.
然後祂來到這個地上.
拿了一個最卑賤的身份.
就是一個竇木老的兒子.
木匠的兒子.
大家可以想像到.
那種身份的轉變.

$^{161}$也都帶來一個很大的挑戰當中.
也都不容易面對的.
大家都認同嗎.
我們的上一任特首.
自從去年六月卸任之後.
七月一日新特首上任.
想問問大家.
你在電視機的畫面當中.
有沒有看到舊的特首的一些行蹤.
有沒有啊.
沒有.
當然她沒有再暴露在媒體當中.
有很多的原因.
但是我相信其中一個原因.
她是要適應.
怎樣由一個七百萬人以上的特首.
一夜之間.
她要變成一個普通的市民.
她要坐公共交通工具.
她要適應.
她要去街市買菜.
她要適應.
頂智慕門.
由一個很尊貴的身份.
忽然之間改變了另一個身份.
而一個普通卑微的身份的時候.
是不容易面對的.
但是我們的主.
因為愛世人的緣故.
他甘願接受這個轉變.
所以這是一個奇異的恩典.
我們再看第三個轉變.
第三個轉變.
是生活的轉變.
頂智慕門.
主耶穌原本在天上的生活.
是一個享受很多的榮耀,尊貴.
有天使,天君來侍候祂的.
一個非常榮耀的生活.
不過因為祂要降生成人.

$^{201}$祂要完成上帝救贖的計劃.
祂就改變祂的生活.
選擇來到地上.
住在一個罪惡的人世間當中.
而且和我們人一起同舟共濟.
各位頂智慕門.
你發覺耶穌在祂付出的恩典當中.
為你為我所做的一個行動.
是一個多大的轉變.
來到世間上.
祂不是享受世界,不是享受世界.
不是坐遊艇,享受咖啡,到處去西施營.
不是的.
祂是什麼?.
祂是關懷有需要的人.
祂是與罪人在一起.
頂智慕門.
耶穌這種美好的榜樣.
就激勵了歷世歷代裡面.
有無數的基督徒都學校主耶穌的榜樣.
來到關愛其他有需要的人.
我很快就想起靈實醫院.
大家聽過靈實醫院沒有?.
你知不知道靈實醫院的創辦人是哪一位?.
靈實醫院的創辦人就是一位西教士.
有少許的護士的經驗.
來到香港創立一個診所.
醫治當時候在調景嶺.
那時候都未叫將軍澳那麼好聽.
叫調景嶺.
說歪一點就是另一個意思.
在那個地方去醫治患了肺病的病人.
由那時候開始.
大家都知道肺病可不可怕?.
很可怕.
尤其是當時候醫學都不昌明.
但是這位孫教士憑著愛心激勵其他的孫教士.
本地的弟兄姊妹.
就創立了靈實醫院的前身.
我就想起為什麼他要那麼辛苦來到香港呢?.

$^{241}$為什麼他不在自己的祖國.
享受他自己的舒適,安舒的生活呢?.
我相信他因為學會了主耶穌的榜樣.
願意在生活上改變他的人生.
改變了他生活的模式.
所以這是生活的轉變.
我們再看耶穌基督降生成人.
第一個部分的恩典.
是他親手造成的.
但是要經過四個轉變.
第四個轉變是什麼呢?.
就是生命的轉變.
一個生命的轉變就是說.
耶穌基督來到這個世界上.
他知道他的生命最後.
是要成為一個獻祭的祭神.
他知道最後他的生命是要犧牲.
要走上十字架的道路.
被掛在十架上.
流出他的寶血.
洗淨全人類的過犯的罪.
因為聖經說.
若不流血,罪就不得赦免.
所以耶穌基督知道他來到這個世界上.
他的生命有一個很大的轉變.
最後他要作出他最大的犧牲的性命.
這種生命的轉變.
是一個很不容易去面對的.
所以耶穌在赫西瑪利園的時候.
大家都記得.
他的挑戰,他的掙扎是非常之大.
但是他最後都信服.
然而不要照我的意思乃是甚麼?.
要照你的意思去成就.
所以我們見到主耶穌.
他生命的擺上.
令我們想起他一出道的時候.
施洗約翰就在眾人的面前.
去介紹耶穌.
你記得嗎?.

$^{281}$他怎麼說?.
他說「看那上神的羔羊,上帝的羔羊,除去世人的罪孽的.」.
大家有印象嗎?.
原來他就是來到這個世間上.
他要成為一個犧牲的羔羊.
去洗淨世人的過犯.
所以他是神的羔羊.
除去我們世人的罪孽.
各位聽眾.
當我們看到這裡的時候.
我們現在明白.
為甚麼保羅說耶穌基督的恩典.
特別強調是耶穌基督的恩典.
是因為他首先要完成第一個部分.
就是他要降生成人.
這個第一個部分.
是原本他是付出的.
他是揀選了.
因你和我他揀選.
他甘願成為貧窮.
所以我們明白.
這一節經文的第一個部分.
原來就經歷了一個這麼重要的轉變過程.
好了.
我們現在來看第二個部分.
下一個部分是甚麼呢?.
下一個部分和你和我有密切關係.
下一個部分就是我們的貧窮.
我們的恩耶穌基督的貧窮.
可以成為甚麼呢?.
付足.
各位頂智者們.
換句話說.
第二個部分.
你和我是要迴轉.
是要悔改.
是要經歷重新的得救.
當你和我願意迴轉悔改.
歸向神的時候.
我們經歷重新的得救.

$^{321}$我們就能夠因耶穌基督的貧窮.
而成為了付足.
所以這裡.
你說是不是多麼奇妙呢?.
凡是每一個人.
不論他的背景.
不論他是一個怎樣好也好.
壞的人也好.
無論他有多少個問題.
他願意迴轉悔改.
歸向上帝的時候.
他就可以經歷這份無比的恩典.
所以因著你和我願意歸向上帝.
我們今天就可以領受四項恩典.
頂智者們.
我們的恩耶穌基督的貧窮.
就可以成為付足.
第一個恩典是甚麼呢?.
是新造的人.
新造的人.
凡是願意接受耶穌基督的救恩的人.
我們在基督耶穌裡面.
都成為甚麼呢?.
都成為新造的人.
所以哥林多後書告訴我們.
他說「舊事已過,都變成新的了」.
你和我因為耶穌基督的恩典.
我們今天都成為一個新造的人.
非常寶貴.
我記得以前有一位.
在戒毒村出了來的弟兄.
在戒毒村我和他都一起生活.
做過他的導師.
他出來了之後.
這位弟兄他分享他的見證.
他原本吸毒又坐牢.
又黑社會裡面打滾.
但當他在戒毒村經歷了.
重生得救之後.
他就分享.

$^{361}$他說他今天的生活.
是一個非常非常奇妙的生活.
是上帝很大很大的恩典.
臨到他身上.
為甚麼呢?.
他說因為我已經成為一個新造的人.
我從來沒有試過這麼被人接納.
現在有很多人願意接納我.
我從來沒有試過被這麼多人尊重過.
但現在人家尊重我.
你知道以前他的生活.
是這麼糜爛的時候.
誰會尊重他?.
但現在很多人尊重他.
他很感恩.
因為他回到教會裡面的時候.
沒有人會歧視他.
他很感恩.
因為他信了耶穌之後.
他的生命得到改變.
人家信任他.
所以給他一份工作.
試一下去做.
他很感恩.
因為他發覺到.
當他也都被人接納的時候.
人家也都願意推薦他.
再去接受一些深造的課程.
當然那些是成人的訓練課程.
因為他已經很大年紀了.
各位弟兄姊妹妹.
這位弟兄的經歷.
令我想起.
新造的人.
是非常非常奇妙.
很大的變化.
很大的祝福.
就臨到我們每一個信主的人身上.
你有沒有想起你自己信了耶穌之後.
是領受了新造的人.

$^{401}$所以你的生活.
你的經歷是與眾不同了.
舊事已過.
都變成新的了.
我們再看第二個恩典.
是一個尊貴的身份.
就是說我們因為成為新造的人.
得到這個救恩.
重新得救之後.
我們原來是在神家裡面的人.
不止這樣.
我們更加是神的寶貴兒女.
這個真是一個非常尊貴的身份.
今天有不少人對於皇親國戚.
或者對於一個國家的元首.
或者對於皇室裡面的人.
都很關注.
他們一舉一動都很關注.
有人寫書.
有人有評論.
等於姐妹們.
一個地上的公主.
一個皇子.
都被人這樣重視的話.
何況我們今天成為上帝的兒女.
我們是一個萬王之王.
萬主之主的寶貴兒女.
我們的身份是更加什麼.
更加高度的尊貴.
難怪彼得前書怎樣說.
他說唯有我們是被揀選的族類.
是有君尊的祭司.
是聖潔的國度.
亦都是屬神的子民.
所以你發覺我們的身份.
是不是無比的榮耀.
第一次妹妹.
不容易得到的.
但是今天我們都得到了.
是因為我們有這位主耶穌.

$^{441}$祂的貧窮就成為了我們的富足.
我們再看第三個的恩典.
是什麼呢.
是一個豐盛的人生.
當我們成為上帝的兒女之後.
上帝不單止給我們一個無比尊貴的身份.
而且聖靈內住在我們內心裡面.
從今之後.
聖靈就是我們的導師.
祂會引導我們人生.
走上每一個階段.
祂會帶領我們.
而且聖靈亦都會將恩賜賜給我們每一個人.
換句話來說.
我們每一個人都有從上帝而來的恩賜.
這樣的緣故.
我們事奉上帝.
就非常之有那份喜樂和憧憬.
你和我都有恩賜的.
各位聽眾.
以前你覺得自己很差.
技不如人.
或者自己很膚淺.
請你不要小看自己.
為什麼.
因為你已經是上帝的兒女.
亦都有聖靈的內注在你心中的時候.
聖靈就會將事奉的恩賜多多賜給我們.
只要我們願意謙卑,學習,追求,成長.
不斷的事奉.
你的恩賜會越來越發揮得更加的淋漓盡致.
難怪有一位姐妹她講見證的時候.
我聽完之後.
她怎樣說呢.
她說自從信了耶穌之後.
我覺得自己每一天的生活都很充實.
都很不夠時間用.
為什麼呢.
她說因為當我信了耶穌之後.
上帝給我有很多服事的機會.

$^{481}$我不單在家中有服事的機會.
我也可以在教會裡面有事奉的機會.
我也可以在醫院裡面做義工.
當我聽到這位姐妹分享的時候.
我為她感不感恩呢.
很感恩.
她的生活我想問大家充不充實呢.
很充實.
因為她說都不夠時間用.
很不夠用.
為什麼呢.
原來她發覺當有了主的生命的時候.
她可以發揮這麼多事奉的機會.
能夠幫到這麼多人.
祝福這麼多人.
那份喜悅.
我相信你和我.
如果有這樣的事奉情況.
你會感受得到的.
如果我們未信耶穌的話.
我想問大家.
你每天的生活會怎樣.
很平平凡凡.
開工,收工,聽電話.
看WhatsApp,睡覺.
你會不會想到原來你的人生.
可以有更多可以祝福人的地方.
就算你想到.
你都可能沒什麼力去做.
有心無力.
但今天很不同.
因為今天你已經是上帝的兒女的時候.
有聖靈的能力.
有聖靈的內注.
有聖靈的推動.
真的好像啪著褲子.
咻一下就去了.
你明白嗎.
所以各位電視節目們.
這是不是一個很大的恩典.

$^{521}$令你的生活活得精彩.
令你的生活活得很充實.
很有意義.
為什麼.
因為我們的生命不再一樣.
最後一點.
我相信人人都明白.
我不需要多費唇舌.
就是當我們成為上帝的兒女.
做神的子女的時候.
第四個恩典是什麼.
有個永生的福祿.
這個更加的寶貴.
在世間.
我們可以有充實的生活.
但是我們也都對將來充滿盼望.
所以當我們去到安息的聚會的時候.
我們的心裡面.
是有一份無比的盼望.
活潑的盼望.
因為我們知道什麼.
我們屬主的人.
地上的生活是一個寄居.
我們有一個永恆的天家.
永恆的將來.
我們在上帝的國度裡面.
我們享受與神的同在.
有這個福祿.
所以各位頂智母們.
你想一想.
我剛才說的那四項的恩典.
你聽了之後.
你覺得這是不是無比的恩典.
真的是無比的恩典.
但是你要記得.
不是平價廉價隨便去得到的.
為什麼.
因為耶穌基督為我們付出了沉重無比的代價.
才可以今天我們得到這一份無比的恩典.
所以我們現在都明白了.

$^{561}$保羅說你們知道這是主耶穌基督的恩典.
什麼恩典.
這是一個基礎性的恩典.
是耶穌基督原本是富足的.
卻因為你和我成為了貧窮.
第二個部分.
我們可以因耶穌基督的貧窮而變成什麼.
變成富足.
所以你說一節經文.
你說奇不奇妙.
很奇妙.
原來裡面蘊含了一個很重要的真理結構在當中.
各位聽眾.
你們知不知道為什麼我今晚要講這個訊息.
知不知道為什麼我今晚要講這個無比恩典.
我想你們記得兩件事.
第一件事.
你要記得.
每一次想起.
我享有無比恩典的時候.
為什麼.
原來上帝要我學習.
凡事者恩.
因為這個恩典我們根本不配得到.
我們是不足的.
我們是貧窮的.
何德何能得到這一份恩典呢.
我們得到這份無比的恩典.
不是我們有什麼了不起.
有本領的地方.
而是耶穌基督完全的犧牲.
擺上祂自己.
所以我們能夠擁有這份恩典.
記得.
任何時候.
隨時隨刻.
我們都要常存感恩.
第二.
我想你們記得.
是要廣傳福音.

$^{601}$因為你今天和我得到這一份恩典.
無比的恩典.
上帝給過我們的.
不是要我們一個人單獨去享受.
廣東話叫「獨食」.
不是叫我們獨食.
第二叫我們與人分享.
告訴一些未認識主的朋友家人.
他們都有資格得到這份無比的恩典.
只要他完成了這節經文的第二個部分.
要做的事情.
他就可以得到這份無比的恩典.
低頭祈禱.
親愛的主多謝讚美你.
我們今晚看了一節經文.
但是裡面有很豐富的教導.
我們真的常存感恩.
因為我們知道.
沒有你的擺上和犧牲.
我們根本得不到這份恩典.
既然我們今天得到這份恩典的時候.
我們的人生不再一樣.
是一個富足的人生.
所以我們要凡事者恩.
我們要廣傳福音.
讓更多的人都得到這份無比的恩典.
祈禱奉耶穌名求.
阿們.
\newpage



\section{馬太福音 18:21-35}
\label{sec:65Ibdpx_Vbg}
\textbf{2023-01-29 加恩樓《接納和饒恕》馬太福音 18\_21-35 講員:謝靄薇牧師}
\newline
\newline
連結: \href{https://youtube.com/watch?v=65Ibdpx-Vbg}{\texttt{https://youtube.com/watch?v=65Ibdpx-Vbg}} ~~~~ 語音日期: 2023-02-01
\newline
\newline
\hyperref[sec:4sQReCkRl1k]{\small{< < < PREV SERMON < < <}}
~
\hyperref[sec:index_chronic]{\small{[返順時目]}}
~
\hyperref[sec:index_scriptual]{\small{[返順卷目]}}
~
\hyperref[sec:fpQjtet_7Iw]{\small{> > > NEXT SERMON > > >}}
\newline
\newline
馬太福音 18:21-35
\newline
\begin{longtable}{cl}
\hline
\hline
章節 & 經文 (和合本修訂版)\\
\hline
18:21 & \begin{tabularx}{0.7\textwidth}{X} 那時，彼得進前來，對耶穌說：「主啊，我弟兄得罪我，我當饒恕他幾次呢？到七次夠嗎？」 \end{tabularx} \\ \\ \relax
18:22 & \begin{tabularx}{0.7\textwidth}{X} 耶穌說：「我告訴你，不是到七次，而是到七十個七次。 \end{tabularx} \\ \\ \relax
18:23 & \begin{tabularx}{0.7\textwidth}{X} 因為天國好像一個王要和他僕人算賬。 \end{tabularx} \\ \\ \relax
18:24 & \begin{tabularx}{0.7\textwidth}{X} 他開始算的時候，有人帶了一個欠一萬他連得的僕人來。 \end{tabularx} \\ \\ \relax
18:25 & \begin{tabularx}{0.7\textwidth}{X} 因為他沒有甚麼償還之物，主人下令把他和他妻子兒女，以及一切所有的都賣了來償還。 \end{tabularx} \\ \\ \relax
18:26 & \begin{tabularx}{0.7\textwidth}{X} 那僕人就俯伏向他叩頭，說：『寬容我吧，我都會還你的。』 \end{tabularx} \\ \\ \relax
18:27 & \begin{tabularx}{0.7\textwidth}{X} 那僕人的主人就動了慈心，把他釋放了，並且免了他的債。 \end{tabularx} \\ \\ \relax
18:28 & \begin{tabularx}{0.7\textwidth}{X} 那僕人出來，遇見一個欠他一百個銀幣的同伴，就揪著他，扼住他的喉嚨，說：『把你所欠的還我！』 \end{tabularx} \\ \\ \relax
18:29 & \begin{tabularx}{0.7\textwidth}{X} 他的同伴就俯伏央求他，說：『寬容我吧，我會還你的。』 \end{tabularx} \\ \\ \relax
18:30 & \begin{tabularx}{0.7\textwidth}{X} 他不肯，卻把他下在監裡，直到他還了所欠的債。 \end{tabularx} \\ \\ \relax
18:31 & \begin{tabularx}{0.7\textwidth}{X} 同伴們看見他所做的事就很悲憤，把這一切的事都告訴了主人。 \end{tabularx} \\ \\ \relax
18:32 & \begin{tabularx}{0.7\textwidth}{X} 於是主人叫了他來，對他說：『你這惡奴才！你央求我，我就把你所欠的都免了； \end{tabularx} \\ \\ \relax
18:33 & \begin{tabularx}{0.7\textwidth}{X} 你不應該憐憫你的同伴，像我憐憫你嗎？』 \end{tabularx} \\ \\ \relax
18:34 & \begin{tabularx}{0.7\textwidth}{X} 主人就大怒，把他交給司刑的，直到他還清了所欠的債。 \end{tabularx} \\ \\ \relax
18:35 & \begin{tabularx}{0.7\textwidth}{X} 你們各人若不從心裡饒恕你的弟兄，我天父也要這樣待你們。」 \end{tabularx} \\ \\
[1ex]
\hline
\hline
\end{longtable}
$^{1}$各位弟兄姊妹平安.
很開心在新年開始.
我們都可以有更多的人聚集.
在這裡祝大家像兔子一樣.
很跳脫.
靈情是步步升.
讓聖靈帶動我們.
每一天我們都能夠靠著聖靈.
我們很有力的為主而活.
在講道之前我們先一起祈禱.
讓我們在新春的時候.
我們可以聚集去敬拜你.
這是你給我們很大的福分.
願主在你的靈眼當中去運轉.
去帶領我們的心.
若然我們有什麼是不要你的.
求主挪開.
以致我們滿身歡喜的.
能夠遇見你.
有願孩子口中的言語.
心中的意念能夠蒙你於立.
我們這樣祈禱交託.
奉主耶穌基督名求.
我們會聽到很多人說.
「報歉」或「我原諒你」.
其實在說「報歉」這句話的時候.
我們的心是怎麼想的呢?.
我們是不是真的報歉呢?.
或者有一個人一而再地犯錯.
你到底要饒恕他多少次呢?.
其實在聖經裡面.
彼得問耶穌.
「主啊,弟兄得罪我了」.
「我饒恕他七次夠不夠呢?」.
但是你會覺得.
他在問這句話的時候.
他應該是覺得自己是很足夠了.
因為拉比對於犯同樣錯的人.
只可以赦免他三次.
七次已經是加倍了.

$^{41}$但是如果你再看聖經文的時候.
你會看到耶穌的印象是不止的.
不是七次.
而是七十的七次.
七是一個完全的數字.
你會看到原來彼得所做的.
似乎不夠.
我們今天來看看.
到底我們饒恕的限度是怎麼樣呢?.
在這章裡面.
仔細看十八章的時候.
你會看到有四次出現了饒恕.
第一次就是二十一節所說的饒恕.
彼得進前來對耶穌說.
「主啊,弟兄得罪我了」.
「我饒恕他七次夠嗎?」.
你會看到這個對彼得來說.
他的饒恕已經是很盡了.
第二個就是在二十七節.
那個人.
二十七節.
請你幫我按一下.
好像不行.
二十七節.
二十七節是說到主人動了持心.
把他釋放.
也是同一個字.
接著三十二節.
是主人叫他來.
對他說你這個惡奴才.
因為他沒有憐憫他的同伴.
所以這裡說到免.
大家想起我們的主禱文.
免了人一債.
免字也是釋放.
也是饒恕的意思.
最後就是我們剛才讀的最後一節.
就是三十五節.
說到你們各人又不從心裡.
饒恕你的弟兄.

$^{81}$我天父也要這樣待他們.
當我們看到耶穌對我們的吩咐.
祂說要饒恕七十個七次.
於是我們就看看.
我們天父的饒恕是怎樣呢?.
當你看二十三節到三十四節.
剛才我們聽到的經文裡.
提醒我們.
原來這裡說到一個和.
和那些縱暴人算帳的時候.
發現有個人欠了王的債.
這個王可能在治理那個地方.
要收一些稅項或者進貢.
人如果欠債不可以償還的時候.
其實是很嚴重的.
你會發現債主強迫他.
和他的家人.
這個人可能就要坐牢.
他的家人可能就要賣他做奴隸.
總之對他是很不利的.
這個人除非他坐牢的時候.
他可以變賣一切.
或者他有親屬去幫他償還債項.
不然他就要一輩子都要坐牢.
你知不知道在這裡說到.
他欠了多少錢呢?.
原來是一萬他年得.
原來一萬他年得就是希臘文的用字.
說到一千萬銀子.
一他年得就等於一千五百塊錢.
一塊錢就是一般勞工的四天工資.
你知道這個數字.
一萬他年得其實在主前60年的時候.
可以說是羅馬大將軍龐貝.
他征服猶太之後.
收取的進貨.
那些進貨就是要這麼多的價值.
你想一輩子怎麼賺這麼多錢呢?.
現在看下去.
這些經文的時候.

$^{121}$你會發現這個同伴欠他多少錢?.
只是十兩銀子.
十兩銀子是一百錢銀子.
就是打一百日的工資.
一日工資就是一錢銀子.
可見這個僕人和這個同伴所欠的錢.
是不是差距很大?.
你會看到這個王竟然去免了這個僕人的債.
是不是很大的憐憫?.
聖經說到他是動了持心.
聖經裡面.
馬太福音裡面出現了五次.
但有四次是說到耶穌基督對人的心腸.
動了持心.
這裡說到這個王動了持心.
他欠了這麼多錢.
他免了他.
不用他還.
但是轉過頭來.
他的同伴欠他多少錢?.
真的太少了.
他的同伴欠他這麼少.
但他竟然和他求情.
他都不肯.
還要把人抓去坐牢.
你說是不是很抵打這個人?.
所以當這個王聽到.
原來這個人這麼差的時候.
難怪他叫他惡奴才.
你會看到原來我們每個人都好像這個僕人一樣.
是夢神很大的恩典和憐憫.
但主耶穌為我們承受很大的刑罰.
為我們死因十字架上.
他沒有犯罪.
但為我們受死.
免了我們的罪.
他用保禦洗了我們的罪.
他今天這樣對我們的時候.
我們怎麼可以不原諒別人的罪呢?.
所以我們看到第二方面就是.

$^{161}$原來饒恕是必須的.
因為神已經寬恕了我們.
我們怎麼去饒恕別人呢?.
原來在聖經十八章裡面.
他描述到神的赦免和慈愛.
鼓勵我們去復和.
如果我們不肯去饒恕別人.
就好像那個僕人一樣.
但神赦免我們.
和饒恕人是拉在一起說的.
他要我們去饒恕一些傷害我們的人.
怎麼做呢?.
很不容易的.
我們有時候不是不想去饒恕.
但心裡總是覺得有些不可以釋懷.
我需要時間去消化.
我怎麼這麼快可以饒恕你呢?.
但這裡是很迫切的話.
即使你得罪了別人.
你也不需要去饒恕.
原來這裡是講天國的教訓.
原來是講到饒恕和天國很有關係.
所以在十八章第三節的時候.
他已經提到.
我們要像小孩子一樣才可以進入天國.
所以在第四節.
他提到如果我們可以謙卑.
就像小孩子一樣.
他在天國裡面是最大的.
小孩子是天真無邪的.
所以在另一個經文裡面.
我們去看對照的時候.
原來看到以弗所知四章31到32節裡說到.
保羅勸我們一切的苦毒,憤怒,老恨,養鬧,誹謗.
和一切的惡毒.
都要從你的心裡面去除掉.
要人慈相待.
全憐憫的心彼此饒恕.
正如神在基督裡饒恕了你們一樣.
原來保羅告誡我們.

$^{201}$千萬不要粗言穢語.
心裡不要有苦毒.
不要用一些憤怒,紛爭,誹謗去惡待別人.
反而要樂於寬恕別人.
好像神寬恕我們一樣.
其實我們想一想.
如果我們用同理心去想一想.
有什麼方法去想一想.
其實我們誰沒有做錯事.
但是你做錯事的時候.
別人會怎樣對待你呢.
可能別人沒有計較就已經寬恕了你.
當我們用同理心去想的時候.
我們可以用這個心情去同情傷害我們的人.
所以當我們想到的時候.
我們又何必那麼根根計較.
我們不可以謙卑一點嗎.
我們謙卑下來去想.
別人怎樣去寬恕我們.
和神怎樣赦免我們.
我們就要感恩了.
曾經有一對姐妹.
這個姐姐跟妹妹是很大的對比.
妹妹是很隨和的.
很會做事的.
而姐姐是很嚴肅的.
很謹慎的.
她很看不慣妹妹那種做事的方法.
而她自己的成績也很好.
所以她看不起她.
覺得自己是最厲害的.
所以她的態度是很驕傲的.
好了.
到她長大的時候.
她就開始很欣賞她妹妹的做法.
她長大後就嫁到馬來西亞.
自己就結婚去了美國.
有一天聖靈光照她.
讓她知道原來她過去是太過份了.
傷害了她妹妹.

$^{241}$所以要她向妹妹道歉.
要道歉是尊嚴的問題.
如何道歉呢?.
我是姐姐.
雖然我是有點壞.
但如何向她道歉呢?.
我過去真的很差.
但聖靈感動她.
要做這件事.
於是她就祈禱了.
求神幫助她.
她真的很不想.
不過她最後還是信服神的意思.
她就祈禱.
趁著她媽媽要做生日.
她們要回家了.
不如在慶祝生日之前.
向她妹妹道歉.
於是她每天都在祈禱.
求神幫助她.
但到媽媽生日那天.
她覺得氣氛不對.
她做不到.
但是神沒有放棄她.
神知道她無法開口.
於是在她臨走之前.
聖靈就催迫她.
讓她知道有件事還沒做.
於是她就很有勇氣.
向她妹妹道歉.
她妹妹就趁機.
說她以前對待她的一些事.
說出來.
當她妹妹說的時候.
數她以前的事的時候.
其實不太好受.
不過她心裡已經預備好了.
因為她經過了祈禱.
心裡已經預備好.
接受了打擊.

$^{281}$但是她心裡很感恩.
雖然她道歉了.
但她總覺得.
她傷害她妹妹太深了.
這個傷痕需要醫治.
她很想向她妹妹祈禱.
但她妹妹拒絕了.
她繼續成為動力.
向神禱告.
為了她妹妹禱告.
她相信神既然感動她.
要向她妹妹道歉.
神一定有辦法幫助她妹妹.
使她們兩人有一天可以復和.
所以其實饒恕是不容易的.
不容易是一件事.
但是需要勇氣.
聖經另一面教我們.
要由心裡去饒恕人.
在聖經裡面說.
你們各人如果不是從心裡去饒恕弟兄.
我天父要這樣待你.
好像那個惡奴才.
他的結局不是很好.
人與人之間都會有磨擦.
都會有爭執.
但聖經教我們要由心裡去饒恕人.
其實饒恕人是基督徒的權力.
不是責任.
是權力.
一種不被妒忌和仇恨捆綁的權力.
是很釋放的.
一個人不肯饒恕人.
這個不被你饒恕的人很痛苦.
其實你何嘗不痛苦呢.
你也很痛苦.
有些人會睡不著.
失眠.
或者他沒有喜樂.
甚至祈禱都不蒙娛樂.

$^{321}$在馬可福音11章25節可以看看.
即是你不去祈禱.
神不會聽你的禱告.
你不會和人和好.
所以原來饒恕人是很重要的.
其實饒恕不是一種損失.
反而是蒙恩的機會.
大衛傷害下屬.
他搶了人家的太太.
一個罪無可恕的惡人.
但為什麼他如此蒙神的賜福.
為什麼他的國家如此強盛.
有道理的.
因為他懂得向主認罪.
他求主赦免.
而且他也饒恕他的仇敵.
這個蘇勒旺.
三番四次迫他.
害他.
但他放棄報仇的機會.
你見到他懂得饒恕他.
不容易的.
所以你發現.
接納不是叫你去包容別人的罪.
或者縱容一個軟弱者.
再而三的犯錯.
反而是要你去引導他.
能夠悔改認罪.
用愛心去饒恕他.
當你接納他的時候.
你就等於在建立他.
就好像耶穌基督接納我們這些罪人一樣.
曾經在1997年4月14日.
有三個人.
聽了你就知道.
在台灣發生的.
在台北市.
陳俊興和林春生,高天明.
他們三個人綁架了.
藝人白冰冰的女兒白曉燕.

$^{361}$她的贖金500萬美金.
但這個綁架失敗了.
於是這個白曉燕的女兒.
就很慘地被殺害了.
然後這班人就繼續犯錯.
犯了很多罪.
警方用很大的警力去追捕他們.
後來他們還大膽.
去到北投的地方.
去挾持了南非駐台灣大使館的武官.
朱茂琪他們全家.
他們想利用外國的事件.
作為人質是最好的.
可以引起全球和政府的關注.
他們希望談判可以換來.
他和他的兒子妻子可以遠走高飛.
在他們的官邸裡.
警方就用槍對他們.
很激烈地對峙.
結果這個武官朱茂琪和他大女兒受傷.
於是就談判.
陳俊興就願意釋放他們兩個.
那家裡剩下誰呢.
有五個人.
有三個是反人.
有兩個就是他太太和小女兒.
小女兒就看著他媽媽.
她說我們會不會死呢.
她媽媽其實也很害怕.
不過她說我們不會死.
我們祈禱神會保守我們.
於是他們兩個就在那裡祈禱.
當他們兩個在那裡祈禱的時候.
他女兒看到桌子上有張紙和筆.
她就拿起來畫.
她就畫了一個十字架.
然後在十字架外畫了一個很大的心.
然後就用筆畫上紅色.
一個很大的紅色的心.
然後在十字架上就寫著「我愛上帝」.

$^{401}$她畫完這個東西.
接著她的目光就望著這個兇神惡煞的反人.
陳俊興.
他這個時候就覺得.
神永遠與我們同在.
神會保守我們一家.
但是他覺得這個人很可憐.
為什麼可憐呢?.
他忽然想起他爸爸跟他說的故事.
就是一個牧羊人.
他有很多隻羊.
但是有一個羊很慘.
在外面風吹雨淋.
這個牧羊人找到他之後.
就抱著他去照顧他.
他覺得這個人就是那個羊.
他很想告訴他.
上帝愛你.
縱然你犯了很多錯.
他很想告訴你.
於是他低頭.
就把這張紙遞給那個人.
他不敢看他.
他很怕他.
因為他樣子很兇.
所以他看不到他的表情.
但是這個人接到他的表情的時候.
這張紙他很感動.
因為他小時候沒有吵.
也沒有哭.
乖乖地在畫畫.
竟然給他一張這樣的紙.
他從來都不會流淚.
但是他這一刻.
他很感動.
透過一個電話.
一個電視台的記者.
跟他通話.
他說一人做事一人當.
他說警察做事都不如一個外國的小女孩.

$^{441}$他送了一幅畫給我.
是他自己畫的.
上面有一個十字架.
他告訴我耶穌愛我.
雖然我不知道耶穌是誰.
但是我知道這個小女孩很勇敢.
我很感動.
他說了這句話.
他不知道是不是逃亡了太久.
或者是他已經知道自己.
已經落在警方的天網地網當中.
或者是他的敏滅的心.
良心被發現.
不知道.
但是這一刻.
他立刻拿起槍.
對著喉嚨.
拿起扳機.
扣住的時候.
在旁邊坐著的.
這個武官的太太.
立刻叫她不要.
她也不理那麼多.
雖然她的普通話說得不是那麼好.
但是她也要叫住她.
叫她不要.
她說Baby.
太太.
她說的很簡單的字.
提醒了她.
我還有家人的.
這個時候她崩潰了.
就跌在地上.
好像一片爛泥一樣.
當然大家也知道.
這兩個人就得救了.
他們一家人就在醫院裡面.
互相擁抱.
因為這個得來不易的得救.
他們向神感恩.

$^{481}$因為台灣和南非斷交.
所以這個武官沒多久就要回國了.
在他回去之前.
他太太和他商量.
我們去監獄探望陳俊興.
因為他太太很想.
在他回國之前.
完成神的託付.
她覺得這個逃犯.
他闖入他們的家.
一定有神奇妙的安排.
所以他們在聖誕節的前夕.
他們兩個人就去探訪他.
他們帶著聖經和聖誕節的糖果.
就去探訪他.
你知不知道這是一個.
被全部台灣人唾棄的人.
覺得他是人渣.
人們很痛恨他.
但是一個被人傷害的人.
竟然來探訪他.
你知不知道他當時的表情是怎樣.
他很驚訝.
一個被他傷害的人竟然來探訪他.
他太太跟他說什麼呢.
他說我們將要回國了.
但是呢.
我很想告訴你.
我們的傷勢都已經好了.
但是我們全家都饒恕了你.
他們的饒恕包含了.
無限的愛和憐憫.
然後他告訴我的小女兒.
特別要告訴你.
他說上帝愛你.
雖然你曾經做了錯事.
但是只要你去尋找耶穌的時候.
向他認罪.
上帝就會原諒你.
大家有沒有想到.

$^{521}$在1998年.
21月的時候.
5號.
在監獄更生團契.
黃明振牧師和義工.
孫悅可能大家都認識的一個藝人.
他們在長期的輔導之下.
這個震驚全球的汪洋大盜.
竟然願意悔改.
認罪.
信耶穌.
然後在監獄裡洗禮.
這是神奇妙的改變.
大家有沒有想到.
這個犯人和這個小女孩.
真的不相干.
但是竟然是因著十字架的愛.
神的大愛.
所以他們兩個可以連在一起.
你看到耶穌基督當年寬恕釘在十字架上的人.
我們每一個在座的.
我們都是蒙主恩,蒙主愛的人.
神寬恕我們.
不是因為我們寬恕人.
而是在我們還是罪人的時候.
祂已經寬恕了我們.
神這麼大的愛.
在我們身上.
如果我們今天都不明白.
神這麼大的愛.
是很可惜的.
神要我們去寬恕人.
如果我們還不願意寬恕的時候.
表示我們.
都不明白神赦免我們.
這個恩典是多麼的大.
盼望我們都願意去回應神的這份愛.
我們願意.
都樂意的.
由心裡去寬恕人.

$^{561}$我們一起祈禱吧.
親愛的主耶穌.
你這份無限的愛.
我們一生都無法去回報.
但是主耶穌求你憐憫我們.
在什麼時候我們跌倒.
我們不想聽你的話的時候.
求你聖靈偏察我們.
幫助我們能夠是真正的.
能夠看到我們.
實在我們都不配的.
你竟然憐憫我們.
去詬罷我們.
如果身邊的人.
我們又怎麼能夠.
這麼驕傲的不饒恕人呢.
求你憐憫我們.
幫助我們.
以致我們真的可以.
按著你的旨意.
能夠去寬恕人.
好像你寬恕我們一樣.
我們這樣祈禱.
教託愛奉主耶穌基督的名教.
阿們.
\newpage



\section{約翰福音 12:20-26}
\label{sec:fpQjtet_7Iw}
\textbf{2023-01-29 午堂《藉門徒見耶穌》約翰福音 12\_20-26 講員:林誠信牧師}
\newline
\newline
連結: \href{https://youtube.com/watch?v=fpQjtet-7Iw}{\texttt{https://youtube.com/watch?v=fpQjtet-7Iw}} ~~~~ 語音日期: 2023-02-02
\newline
\newline
\hyperref[sec:65Ibdpx_Vbg]{\small{< < < PREV SERMON < < <}}
~
\hyperref[sec:index_chronic]{\small{[返順時目]}}
~
\hyperref[sec:index_scriptual]{\small{[返順卷目]}}
~
\hyperref[sec:i9iPPyuXEWw]{\small{> > > NEXT SERMON > > >}}
\newline
\newline
約翰福音 12:20-26
\newline
\begin{longtable}{cl}
\hline
\hline
章節 & 經文 (和合本修訂版)\\
\hline
12:20 & \begin{tabularx}{0.7\textwidth}{X} 那時，上來過節禮拜的人中，有幾個希臘人。 \end{tabularx} \\ \\ \relax
12:21 & \begin{tabularx}{0.7\textwidth}{X} 他們來見加利利的伯賽大人腓力，請求他說：「先生，我們想見耶穌。」 \end{tabularx} \\ \\ \relax
12:22 & \begin{tabularx}{0.7\textwidth}{X} 腓力去告訴安得烈，然後安得烈同腓力去告訴耶穌。 \end{tabularx} \\ \\ \relax
12:23 & \begin{tabularx}{0.7\textwidth}{X} 耶穌回答他們說：「人子得榮耀的時候到了。 \end{tabularx} \\ \\ \relax
12:24 & \begin{tabularx}{0.7\textwidth}{X} 我實實在在地告訴你們，一粒麥子不落在地裡死了，仍舊是一粒；若是死了，就結出許多子粒來。 \end{tabularx} \\ \\ \relax
12:25 & \begin{tabularx}{0.7\textwidth}{X} 愛惜自己性命的，就喪失性命；那恨惡自己在這世上的性命的，要保全性命到永生。 \end{tabularx} \\ \\ \relax
12:26 & \begin{tabularx}{0.7\textwidth}{X} 若有人服事我，就當跟從我；我在哪裡，服事我的人也要在哪裡；若有人服事我，我父必尊重他。」 \end{tabularx} \\ \\
[1ex]
\hline
\hline
\end{longtable}
$^{1}$弟兄姊妹早晨.
我謹代表我們播道會廣福堂眾弟兄姊妹.
在主內向大家問安.
因為今天是農曆大年初八.
我在這裡先恭祝大家.
「吐年蒙恩,福杯滿日」.
我們是屬於上帝的兒女.
我們真是在神的家.
我們在我們這個群體裡面.
就算我們在不同的教會當中.
我們其實在神的眼中.
我們都是一家人.
我們是同屬耶穌基督的身體.
亦都是要透過我們的群體.
雖然我們坐落在不同的社區裡面.
但要讓人看見主.
看見主耶穌在我們的生命裡面.
那個呈現其實不只是個人.
是一個群體.
一個群體在歷世歷代.
在教會歷史裡面.
雖然教會都遭遇過不同的時代.
不同的艱難.
不同的挑戰.
不同的風雨.
但我們卻仍然可以屹立不倒.
原因只有一個.
就是當我們願意以耶穌基督為我們的防國石.
教會立基就在耶穌基督的身上.
我們是要彰顯耶穌基督的真理.
將基督的愛呈現在這個充滿了很多苦難的世界當中.
我為何會有這個題目呢.
是源於我今年教會都準備要在港島南區.
去樹立一個社區中心.
原因是因為我們很想幫助弟兄姊妹.
將我們的信仰實踐耶穌基督的生命.
在我們金鐘區是一個商區.
我們很難叫商區的朋友.
如何與他們建立一個所謂的鄰舍的關係.
我們要進入基層的一個群體當中.

$^{41}$在這裡給大家有兩個分享.
首先在去年9月的時候.
我看過一個12月.
去年12月初的時候.
有一個資料出來是說2022年.
全球的生活成本指數出爐.
全世界各地的平均通脹超過8\%.
很厲害.
是20年來最大的一個升幅.
我想問一下大家.
在他們這個調查裡面.
全球172個大城市當中.
你猜香港的生活成本是位居全球第幾呢.
第幾呀.
第一呀.
那還不至於.
第幾呀.
第五呀.
都很近了.
第四.
第三.
我想跟大家說.
我們排第四.
是與洛杉磯齊名的.
都很貴的.
其實我們的生活.
樓價又高.
平時你去那些所謂茶餐廳.
你發覺過去這一年.
其實所有的東西都貴了.
最貴的就是紐約和新加坡排第一.
這件事也告訴我們.
我們其實是面對一個現實的生活處境.
我們這個城市其實有很多很艱難的貧苦大眾.
生活在水深火熱的當中.
我們今天面對的這個世界.
其實普遍很多大城市.
貧富懸殊是越來越大.
這個是不爭的事實.
在去年12月12日.

$^{81}$我亦都看過一個新聞.
是說慈雲村發生了一宗淪陷的慘劇.
一對年逾70歲的夫婦.
他們懷疑因為一個長期要照顧患病五年以上的.
他們的女兒47歲.
引發到這對夫婦情緒壓力爆煲.
他們晚上趁著這個女兒行動不方便.
在睡覺的時候.
他們就私語獨守.
親手將他們這個47歲的女兒.
用一些鐵鎚打死了她.
聽到這件事.
我覺得心裡頭很難過.
再看一看香港防癌會.
在去年9月到10月.
他們訪問256位癌症患者的照顧者.
發現5成9的受訪者感到抑鬱.
更加有一成人曾經說自己想過.
甚至已經試過自殺.
調查顯示第五波疫情高峰期間.
受訪的照顧者的焦慮指數攀升了3成4.
指數是由疫情前的5.8分升至7.8分.
這個真實的數據和事例.
亦都帶出我們現在所謂國際的大城市.
其實我們有很多人都活在很大的焦慮.
壓力甚至徬徨抑鬱當中.
究竟我們今天作為神的兒女們.
作為教會人們在不只是個別的信徒身上.
而是在整體教會的群體身上.
他們看到什麼呢.
我們經常說我們需要傳福音給別人.
這是絕對沒有錯的.
但今天我們說傳福音給別人的時候.
是要讓周圍的人在我們整個群體裡面看見.
不單止有基督的真理.
更加有基督的愛.
是他們感受得到.
他們是來經歷得到.
實實在在的接觸到人群.
他們生活所需要的.

$^{121}$因為我們所傳的福音.
是關心到整個人的需要.
身心瀉靈.
全人的需要.
因為我們的福音這個使命.
亦都是一個全人的使命.
全人要健康.
家庭要健康.
群體要健康.
社會要健康的話.
社會都有這個.
真的很重要.
我們在這個履行大使命的過程裡面.
讓人看見基督的愛.
是呈現彰顯.
究竟我們怎樣能夠讓人看見.
主在我們身上呢.
我盼望用今天這段經文.
和大家一起思想這件事.
我們做一個簡短的禱告.
天父上帝當我們打開你的聖言.
求主再一次透過你自己真理的聖靈.
在你的兒女的心中.
你去作感動.
作引導.
作剛召和鼓勵.
叫我們真心打開.
我們的眼睛張開.
我們願意回應你自己的說話.
我們實踐我們的信仰.
在我們每一天的生活裡面.
願主保守.
叫講者聽的.
我們都能夠同感一靈.
我們同德造就.
同心祈禱.
奉主明求.
阿們.
我想和大家再一次.
來看這段經文.

$^{161}$其實就是剛才主席已經帶我們看過.
其實很好的一段經文.
在約翰福音十二章.
這裡是講到.
那時上來過節禮拜的人中.
有幾個希臘人.
他們來見加利利的.
伯塞大人肥力.
請求他說.
請你們先生我們想見耶穌.
這件事我們想問一問大家.
那時是講什麼時候呢.
來過節是過什麼節呢.
如果大家看約翰福音的時候.
這是描述耶穌生前.
他最後一次進入耶路撒冷.
是講他釘十字架之前.
最後一次進城.
一個經歷.
這個節期是什麼節期呢.
有沒有人知道.
就是如熱節.
是要紀念.
舊約是說他們紀念神的羔羊.
要塗抹羔羊的血.
在門檻上.
滅命的天使就要越門而過.
是紀念這種請求救贖.
但對於猶太人.
歷世歷代他們等待的.
就是要來的這位彌賽亞.
他們守這個節期.
就是讓他們不要忘記.
上帝救贖的恩典.
而當中我們看到很特別的.
因為在耶路撒冷.
當時因為過這麼重要的節期.
是一年中最重要的這個雨節期.
所以是很多人來到.
而當中我們發覺.

$^{201}$不是猶太人來.
還有幾個希臘人.
這個就很特別.
希臘人屬於愛邦人.
愛邦人他們不能夠進到.
他們猶太人的聖殿裡去敬拜.
他有一個愛邦人的院子.
在外圍的地方.
但他現在很想來見耶穌.
當肥力知道這件事.
他告訴安德烈.
安德烈就和肥力去找耶穌的時候.
就說有人想見你.
你看看耶穌的回答.
是很有趣的.
耶穌說.
人子得榮耀的時候到了.
好像答非所問.
今天如果有人來你家找你.
我想見你.
你的子女或者其他.
或者有工人的話.
他說阿Sir.
有人要找你.
那你說.
我要得榮耀的時候到了.
人家不知道他們說什麼.
但這個是很有意思的回答.
一後文就是耶穌解釋.
這個得榮耀是什麼意思.
這個是指向他為我們釘十字架.
他要為我們眾人的罪死.
而這兩個人物的出現也都很特別.
他們希臘人他們講求智慧.
他們追求這些神蹟.
他們在前文十一章是說.
耶穌是經過伯大尼.
他們才進入耶路撒冷.
伯大尼當時他要探望.
他們有一個他所愛的家庭.

$^{241}$就是馬大馬尼亞.
和他的弟弟拉撒路.
當時拉撒路病重.
有人通知耶穌.
耶穌等多四天才去.
結果拉撒路都死透了.
來到的時候.
馬大馬尼亞他們很傷心.
主耶穌你來了.
不過弟弟其實已經死了.
對於猶太人來說.
死了四天是死透了.
但耶穌說.
復活在我 生命也在我.
信我的人雖然死了.
也必復活.
復活對於我們信主耶穌的人.
我們無懼於死亡.
是因為我們看到.
這個榮耀的指望是實實在在.
耶穌基督從死裡復活.
父神為他的兒子來去引證.
他確實是上帝所差派來的.
為我們眾人去罪犧牲.
贖罪的這位神的羔羊.
聖潔沒有瑕疵的神的兒子.
為我們死而復活.
而且他應許我們.
他一定要再回來.
接我們回到這個榮美的天家.
所以今天無論我們在這個世界.
遇到多大的難處.
我們仍然知道.
將來那個榮耀的盼望.
是一定會應驗.
而我們今天我們需要的.
就是履行我們在地上.
父神託付給我們.
宗教會的使命.
耶穌在這裡說的.

$^{281}$他得榮耀其實還有一個含義.
是因為他為我們死而復活的緣故.
他是打破了.
愛邦人和猶太人之間的隔閡牆.
就是以弗所書第二章所形容的.
原本猶太人和愛邦人是水溝油.
溝不來的.
但是主耶穌一來.
他十字架滅了當中一切的冤仇.
他將猶太人和愛邦人.
都是要用他的寶血去救贖.
凡信靠主耶穌基督的名的人.
無論你是屬於哪個國籍也好.
無論你是屬於什麼民族也好.
你都可以因著信靠耶穌.
你歸為一體.
是十字架叫我們能夠合而為一.
是基督的大愛.
叫我們可以彼此同心.
在這個社區 在這個城市.
在全世界任何一個角落.
我們一起同工.
我們要帶領人進入一個光明榮耀的角度.
這個就是我們歷世歷代.
教會在地上永遠不改變的使命.
神透過耶穌基督這個回答.
他更加要指向一個很重要的一個信息.
他下文跟我們提到.
一粒墨紙這個比喻.
一粒墨紙不落在地裡.
死了仍舊是一粒.
若是死了就結出許多紙粒來.
他跟著說.
我們作門徒的.
就是我們要寫記.
背去十格.
要跟隨主.
愛惜自己生命就喪失生命.
幸護自己在世上的生命.
必要保存性命到永生.

$^{321}$若有人要服事主耶穌.
就要跟隨他.
他在哪裡.
服事他的人都要在那裡.
如果有人是這樣去跟隨主的腳蹤.
去服事他的話.
父神必尊重這個人.
所以今天無論我們身處何方.
我們只要我們清楚上帝的心意.
無論上帝帶領我們在任何一個地域.
任何一個國家.
任何一個城市.
我們都能夠奉主的命.
我們說我們要作主耶穌的門徒的話.
我們必須要順服父神的旨意.
和他的差遣.
做當做的事.
究竟這件事是什麼呢.
我們看看主耶穌為我們犧牲的時候.
耶穌是百分之百是神的兒子.
但同樣他是百分之百是人子.
耶穌面對要上十字架.
為我們眾人寫命這個事實.
其實你不要想很簡單.
他是神的兒子.
甚至有些人曾經在慕道的過程.
跟我說牧師耶穌既然是神的兒子.
他在十字架死.
死得多慘都不要緊.
不會痛的.
他是超人.
我說你這樣想你搞錯了.
耶穌不單止是肉身.
他要受著那些皮肉的鞭打.
和刺傷.
耶穌上十字架是要與父神分離.
因為他背負了我們眾人的罪.
他原本是聖潔沒有任何瑕疵.
他不應該死在那裡.
他是代替你和我死在那裡.

$^{361}$但耶穌他面對這個考驗的時候.
他在厚瑟瑪利園祈禱三次.
去求父神免去這個苦悲.
但最後他都信服.
他說父啊不是照我的意思.
而是照你的意思.
這個就是作門徒一個最關鍵的學習.
我們要真正學好.
我們如何去服侍別人之前.
我們必須要通過這個十字架.
那種破碎和對付.
原來我們的生命沒有被雕琢.
沒有學會如何完全信服父神的旨意.
我們被神雕琢我們很多性情的瑕疵的話.
我們不會去服侍別人.
往往我們去服侍神.
所以我想大家先要看清楚這個圖畫.
原來我們作主的門徒.
跟隨他我們這個效法主.
效法主就是徹底完全的信服.
父神在我們生命的計劃帶領.
最後是成就父的旨意.
是預願上帝的旨意去成就.
讓別人在我們身上看見耶穌.
這樣才可以看見耶穌他的榮耀.
當我們有一天.
我們的安息禮拜.
別人給我們的評價是什麼呢.
如果別人說我們在我們身上看見主耶穌.
我相信你會很感恩很開心.
真的.
這件事.
我記得有一個例子.
是在紐約地鐵很繁忙的時候.
早上上班.
有一個老人家.
一個婆婆.
她的袋子裡有些橙.
但一不小心掉到地上.
那些橙掉到滿地都是.

$^{401}$但那些人趕著上班.
每個人只是走過就算了.
只有一個男士.
他突然來收拾好這些橙.
放回紙袋裡.
給這個婆婆.
婆婆這些是你的橙.
婆婆抬起頭來.
看這個男人.
跟著跟他說.
先生多謝你.
我在你身上.
我看見主耶穌.
看見主耶穌.
因為他看到這位男士.
就好像上帝差來的.
一個及時的天使一樣幫助他.
這個男士再看一看這個婆婆的時候.
原來她是慢的.
難怪她根本沒可能收拾好那些橙.
我們看到有需要的.
就算我不認識的一個陌生人.
就在地鐵站有什麼需要的時候.
我們立時伸出援手.
這個是惻隱之心.
人皆有之.
我們中國古人萬子都會這樣教導我們.
神的兒女們是需要做成這些東西.
人在我們身上.
可以看見主與我們同在同行.
在六年前.
我們教會一個做前台接待的.
一位姐妹同工離職的時候.
她跟我說林牧師多謝你.
讓我在這裡有一段時間去服侍.
跟著她說了一句話.
我當時一方面很感恩.
但回頭一看.
我又覺得很慚愧.
她怎麼說呢.

$^{441}$她說林牧師我在你身上.
我看見主耶穌.
你要問我主席的時候.
如果人家今天看到你.
我在你身上看見主耶穌.
你會怎樣.
你會很慚愧.
我想我們第一個反應可能你會.
都會有點開心.
你會很感恩.
你可能也很好奇.
你究竟在我身上.
你看見哪一部分我是像主耶穌的呢.
看見主耶穌在我身上.
彰顯.
我問這位姐妹的時候.
她說我對人很有恩慈.
她看見我平日處事待人.
是很有忍耐,寬容.
我聽完之後.
我說姐妹.
其實我在教會.
在很多人面前.
可能你會看見我有愛心.
但我也有很多論盡野的.
我的忍耐很多時候不足夠.
我們作牧師.
我們都是人來的.
回到家裡問師母.
我有時候論盡野就出來.
我有時候在預備講章.
或者我要找一些書.
找不到的時候.
可能我有些煩躁.
師母啊.
我需要那本書在哪裡.
有時候我找不到.
她就幫我找到.
我感恩.
我有個言耐座.

$^{481}$但我想告訴大家.
原來作主的門徒.
我們作為一個基督的群體.
我們很真實地呈現.
我們的喜怒哀樂.
我們的擔心.
我們的憂愁.
我們都可以成為這個群體.
一起去禱告我們的經歷神.
我們也可以和我們的街坊.
我們的左鄰右舍.
我們這個社區.
我們讓他們透過不同的活動.
透過不同的聚會.
我們邀請他們來參與.
我知道去年.
曾牧師已經帶領大家.
開展了幸福小組.
其實這個幸福小組.
是一個很有效的小組.
因為幸福小組不是個人報道而已.
是群體與群體.
我們一個小組裡面.
起碼有十個八個人.
有幾個朋友.
加上幾個弟兄姊妹.
當中有組長.
有副組長.
一起用愛心.
特別是我們現在慢慢復常的時候.
可以一起吃飯.
吃完飯.
可以一起分享見證.
一起去查經.
這個過程其實就是建立友誼.
很多時候我們說回過去這麼多年來.
在歷代的橫教會.
我們因為都是源自西方的宣教士.
來到我們中國.
來到香港.

$^{521}$來到東南亞.
來到全道.
我們很著重那種粉輕報道.
我們說報道會.
誰要信耶穌請你舉手決志.
我不是要反對這件事.
這件事可以繼續進行.
但是對於我們中文文化裡面.
我們反而很著重那種一家人吃飯.
與朋友一起吃飯.
在這種交往的過程.
我們就覺得.
我是一個很安全.
一個很友善的環境.
我們一起彼此交流生命裡頭的經歷.
從而在我們身上看到神的恩典.
看到神的作為.
也因而在這個過程裡面.
這班朋友被帶進去教會.
他們在當中就來認識主.
有很多人都是因為在幸福小組.
未完之前.
他們就決志信主.
到完的時候還未信主的話.
他們都繼續可以.
他們就會很願意跟你一起回教會.
所以這個方式是很有效的.
遠比單單只是請他們來報道會.
突然在那裡.
他們第一次來.
然後你說決志信耶穌.
有些人就很不習慣.
他甚至很害怕.
甚至如果有些新朋友.
我曾經聽過.
是不是我不決志你不讓我走.
但我們很多時候.
因為很緊張.
訓練了頂展會做個人報道.
他不肯決志.

$^{561}$你就跟著下一次.
你繼續邀請他.
繼續邀請他不肯來.
你以後就覺得這個可能沒什麼希望.
我就去釣其他魚.
就跟他的朋友都沒有了.
如果我們這樣.
很多時候一些誤導的親友.
他就對基督教有一種誤解.
我曾經進過教會.
不過印象不是很好.
我們要將這種友善.
將這種關懷.
我們要帶進整個朋友圈子裡面.
我們要做到.
甚至一些陌生人.
我們教會有位執事.
在過去這幾年的疫情.
他專門去探訪無家者.
有些流浪漢.
他甚至跟他們成為了朋友.
他每一次探訪完之後.
他都寫一個見證給我看.
我為他一起祈禱.
為他感恩.
他現在有一個日記.
叫做光明隧道日記.
就將他這些經歷.
就跟我們分享.
我跟這位執事一起同行.
今年亦都要裝備頂子妹.
我們要走出我們的舒適區.
我們甚至願意花時間.
在人的身上.
信徒身上門徒身上.
怎樣能夠讓人看見這位主耶穌呢.
其實不容易的.
猶太人要打破他們過往固有的觀念.
要跟外邦的信徒.
要合而為一.

$^{601}$要彼此相愛.
他們起初覺得這些不是神的選民.
直至到耶穌來的時候.
直至保羅在當中竭力來傳講福音的時候.
就說你們不要將神的兒女分門別類.
無論他們屬於什麼民族都好.
他們都是上帝所創造的兒女.
我們都可以彼此同為一體.
我們是天父的兒女.
所以在當中我們知道.
我們要做到這件事.
我們要靠真理的聖靈.
是神的靈幫助我們.
轉化我們的生命.
以致我們能夠讓人在我們身上.
看見主自己活著的事實.
我們都在保羅的生命裡面可以見證.
保羅他坐牢的時候.
他仍然寫信給肥立比的教會.
這個教會從頭一天.
直到他寫給肥立比書.
有12年相隔這麼久的時間.
但保羅為他們感恩.
說你們從頭一天直到如今.
你們都是同心合意興旺福音.
保羅說到無論是生是死.
我總要讓基督在我身上照常顯大.
因我活著就是基督.
死了就有益處.
保羅知道他的人生是為了什麼使命.
為了什麼目標而生活.
而他同樣鼓勵昔日肥立比的教會.
同樣我相信今天我們地上所有的教會.
都需要重新反思.
我們在今天這個同樣彎曲薄謬的世代裡面.
我們要做神無瑕之的兒女.
在這個世代裡我們要將明光照耀.
將生命的道顯明出來.
我們必須將這種生命的道顯明出來的時候.
就是我們需要透過神的靈在我們裡面.

$^{641}$去破碎我們 去感動我們.
以致我們願意降服 我們信服主.
我說主 你願意差遠我去哪一處.
你願意叫我在哪一個崗位.
你願意我接觸哪一個群體.
你願意我在我的朋友圈子.
我的親戚的圈子裡面.
有些可能幾十年沒見了.
或者過去有些家族的一些原因.
甚至你的親人.
有些就是因為一些誤會.
一些小小的誤解.
十年八年沒來往.
弟兄姊妹 讓我們今天.
我們會想神要我們去重新接觸他們.
帶領他們.
這位謝遠文醫生.
不需要我介紹了.
大家記得他是誰嗎.
在03年 他是我們香港的香港女兒.
抗疫的英雄.
他在1986年以百優的成績.
入讀中大醫學院.
他在入學的手冊裡面.
他這樣寫他的志願.
他說我要加入醫生的行列.
為社會和人民服務.
我更加要發揮自己所長.
他的同事是這樣來見證他.
他說謝遠文醫生.
他是擔心其他人多過他自己.
他當時是做了屯門醫院的空肺科的主任.
但是他因為要救治那些SARS的病人.
他願意走到最前線.
進到ICU裡面.
幫那些病人去插喉 去抽痰.
但很不幸 他與另外一位男護士劉永佳.
懷疑因為沾上那些病人的飛沫而染上SARS.
他最後殉職 中年35歲.
我們看見這位姐妹很愛神 也很愛人.

$^{681}$他要知道他的人生是為了什麼來努力.
已經將近20年前的SARS奪去了香港299條寶貴的生命.
謝遠文醫生的離世卻是換來給很多人生命的啟發和激勵.
有很多人立志要做醫生 是因為他的緣故.
有很多很優秀的人才 是因為他願意為病人去寫命.
也願意沒有逃避他的責任和天職.
我剛才說過 他當時已經是空肺科的主任.
他其實可以差遣他的下屬去做.
但他率身先試卒 他沒有逃避.
他很勇敢 是因為他真的愛那些病人.
他是一個好醫生 他也是一個天賦的好女兒.
也是我們香港引以為傲的一位女兒.
這個生命的見證 正是提醒我們今天.
我們都是願意要效法耶穌的忠行的時候.
就是有一天或者某一個moment.
如果聖靈要感動我們去做一些你覺得很難的事.
你甚至覺得這個我不敢去鬥.
甚至可能有危險的 但你能否信服主呢.
在過去這幾年的疫情裡面.
教會我們和弟兄姊妹一起仍然關心這些不同的社區.
去年七月我和夫婦區一班弟兄姊妹.
我們一行五十人去播道會和屯門的山景堂.
我們一起有合作 他們都有五十多位弟兄姊妹.
和我們一起分成很多個小組.
你見我早上在教會講完道.
下午我就脫光我的西裝 我不打領帶.
我穿牛仔褲 穿T恤.
因為我知道我要去服侍這些基層的家庭.
我想讓他們感受到其實牧師都是人.
我們可以很平民化 我要和他們打成一片.
我要讓他們知道我們去探望他們.
不只是帶物資去 我們很想在這個過程.
他們感受到神的兒女們帶著基督的愛去關心他們.
關心不單止他們生活的需要.
同樣關心他們心靈的困苦.
一個人沒有主耶穌的時候.
他們的生命就會終日徬徨 沒有死亡.
那天我被派和師母一起.
和他們山景村裡面的播道會的其他姐妹一起.
我們一起在那裡探訪.

$^{721}$我派去那兩個家庭 我真的覺得很巧合.
我相信一切不是巧合 神一早就有預備.
神竟然差遣我們去探望兩個家庭.
而兩個家庭竟然都是潮州人.
我本身是潮州人 我會說潮州話.
所以我第一個家庭探望完他們之後.
一知道原來他們是潮州人.
我聽到那個婆婆的口音很重潮州話.
我說你是不是潮州人 她說是的.
然後我們用潮州話交談 她很開心.
我還感動 一下去的時候.
那個婆婆這位老人家其實不是年紀很大.
對今天來說 她72歲 穿粉紅色衣服.
她是一位清潔工人 但她三個星期前.
知道我們要來探訪 她就去告定價.
她星期天都要上班.
但那天她特意放假 告假.
就是為了等我們來探望他們.
都是將物資交給她的時候.
我們很感動 她來到這樣和我們一起.
我們來到和她一起分享舊恩.
當我和她談 我問她有沒有聽過耶穌.
她怎樣回答我呢 她說她年輕的時候.
很多年前她已經在潮州她的家鄉汕頭.
進過禮拜堂去聽福音.
她聽過 她知道主耶穌是誰.
我說你有沒有信到主耶穌 她說沒有.
後來因為要來香港結婚生兒育女.
之後這幾十年 她就是將她整個生命.
奉獻給家庭 為了丈夫為了兒女.
為了她一生的苦惱 全都給了他們.
當天我聽到這件事 我立刻很感動.
我說上帝 我很相信今天.
這幾十年的福音的種子.
就是你要拆遷我們這一小隊.
來到就是叫這位老人家.
叫他能夠可以去決志去歸向你.
我再很簡單和她講一遍.
我們主耶穌是誰.
我們的救恩是一回什麼事.

$^{761}$我鼓勵她 要接受基督耶穌.
做她生命的主 做她的救主.
她欣然接受 她說她願意.
完了之後 山景堂的那位姐妹就繼續跟進她.
原來幾十年的福音種子.
不是埋藏在地上就消失了.
這個原來是等候我們每一個神的兒女們.
你今天可能在你的朋友圈子.
在你的親戚朋友裡面.
總有一些人就是這樣.
他今天仍需要你去鼓勵他.
重新帶領他歸入耶穌基督他自己的名下.
我們看回另外一個圖片也很感動.
去年十一月的時候.
我們這個長者團契.
我們去到太古城有一個銀杏館.
這幾張圖片你看到都是一班長者.
就是我們長者團的一班老人家.
去到當中有些是比較年輕的長者.
我想大家留意這四張照片.
都有一個特色.
就是每張照片都只是一個弟兄.
有一個男士.
這個弟兄92歲.
他在疫情前是可以吹瑟瑟風的.
他竟然可以和我們一班年輕的長者.
進到銀杏館那裡.
吃飯 舀湯.
然後做好之後就去分派給那些街坊.
鄰近的長者.
我最感動的是右下角那張照片.
就是這一張照片.
你看到他是92歲的一位長者.
他仍然有健康.
他就因為基督的愛驅使他.
有心有力地.
他仍然去服侍這個社區有需要的人.
是一位92歲的老人家.
將這些飯送來物資.
送給另一位有需要的長者.

$^{801}$我看見的時候.
我就說今天我們服侍社區.
不是青少年人的專利.
是每一個神的兒女們.
我們都能夠作的.
我們只要.
我們看見在我們這個群體裡面.
神的家.
正如彼得前書所說的.
人人要照自己所得的恩賜彼此服侍.
作神各種恩賜的好管家.
彼得在這裡跟我們說的.
這個恩賜有兩大類.
一方面是說話那種恩賜.
講道 教道 輔道.
我們勸勉.
這些是說話有恩賜的.
原來神要教會有很多弟兄姊妹.
都能夠成為他自己的恩賢.
他自己的聖賢的出口.
但是同樣.
另一方面.
我們有很多弟兄姊妹.
要成為手 成為腳.
成為服侍其他別人的.
我們要一起同心去服侍 配搭.
這樣才能夠成事.
讓人在我們身上看見耶穌基督的愛.
曾經有人問一個很成功的企業家.
就說我們接下來這三十年.
究竟我們下一代的小朋友.
最需要的是什麼呢.
他的答案我覺得很有啟發性.
他說其實今天我們那些.
網絡世界的知識是爆炸的.
如果只是講知識.
其實很多小朋友自己上網.
他都學到很多東西.
我們今天的教育.
不只是要他IQ.

$^{841}$不只是那些知識.
也不只是EQ.
EQ是那些所謂的情感智商.
Emotional.
這種的healthiness.
這些都能夠在很多人眼中.
你要成功.
你就必須有這些的因素.
但他提到有一樣東西.
你必須要有的.
但是書本是很難教得到我們的孩子.
這個叫做LQ.
大家知不知道LQ是什麼.
就是Love Quotient.
就是講愛這種的智商.
愛是你讀書讀不到的.
愛也不是你上幾堂愛的功課.
就搞定.
或者你聽幾堂愛的講座.
就搞定.
不是的.
你問很多夫婦.
你只去聽那些恩愛夫婦.
那些教你怎樣彼此相愛.
最重要是什麼.
實踐.
你聽完恩愛夫婦.
你回到家.
哇 真是.
牧師或者哪個專家講得真是好聽.
只是好聽沒有用.
你要實踐.
原來我們今天最重要的.
是要將神的話語.
在我們這個群體呈現之外.
就是神的愛.
去關愛別人.
就要彰顯上帝的心意.
在上個世紀裡面.
有一個很著名的宣教士.

$^{881}$這個是愛爾蘭的教會宣教士.
叫做Amy Kamiko.
假艾梅.
她當時去到印度.
她開設了一個孤兒院.
在那裡一服事.
就是55年.
她整個人生奉獻給那個地方.
有需要的這些孩童.
她也寫了35本書.
是有關宣教士的事奉.
其中一本書是翻譯成中文.
是成為很暢銷的一本小冊子.
不過已經很陳舊.
年青一輩那些弟兄姊妹可能沒見過.
我想問一下他們在當中.
有沒有人看過一本叫做《若》.
《若果的若》.
這一本小冊子.
有啊 有些弟兄姊妹你看過.
這本《若》真的寫得很好.
它其實有很多個小段落.
而每一個小段落.
它寫到最後就是.
若果你沒有這些東西的話.
那我就是絲毫不懂得加勒山的愛.
它說若果我對別人的需要不夠體貼.
或者對他們的感受.
或者他們甚至少少的軟弱.
毫無感覺的話.
若果我沒有留心注意他們所受的少少傷害.
因而失掉機會去撫慰幫助他們的話.
若果我叫家中的甜美運作變得更加困難.
那我就還是絲毫不懂加勒山的愛.
若果我身旁的人在受苦.
我卻毫不覺察.
因為敏銳的靈不在我裡面.
那我還絲毫不懂得加勒山的愛.
弟兄姊妹今天我們有沒有.
因為我們裡頭這一種.

$^{921}$在一個很物質化.
一個很商業化.
一個很功效化的社會裡面.
現代化的社會裡面.
我們連一些最基本的人倫之間.
那種的人愛都沒有了.
我們中國聖人孔子跟我們說.
這一種人愛就是由近至遠.
由親到疏.
這一種就是在我們家庭.
我們家人開始.
我們要學習彼此兄弟有愛.
姊妹有愛.
我們父母和兒女有愛之外.
是延伸到整個社會眾人.
教會就是這一種群體.
我們必須要堅持我們一個純正的福音.
和真理之外.
你就必須要呈現這一種入世卻不屬世.
真正是上帝的愛.
在我們當中散播開去.
求主幫助我們.
今天我們看到基督的愛不是抽象的.
基督的愛就好像魯約翰說的.
我們既然經歷過基督怎樣為我們寫命.
我們如今知道什麼叫愛.
我們都要為弟兄去寫命.
故此如果我們今天有世上的財物.
但見到有一個弟兄缺乏.
我們竟然關閉了惻隱的心.
神的愛怎樣在我們裡面呢.
我們的傷害不可以只是在言語舌頭上.
我們不可以只是有口號.
但沒有實際行動.
我們要在行為和真誠裡面表達出來.
所以今天我們必須愛我們一個見不到的上帝.
但你千萬不要塞住你的心.
你去恨你的弟兄.
然後你說你要愛神.
求主幫助我們.

$^{961}$在這裡我們也看到另外一個見證.
最近在美國的郵報和紐約時報.
都成為一個非常流傳在網絡世界的新聞.
這個退役的空軍.
這位叫做Hody Chydris.
他很愛自己的國家和社區.
他退役之後成為一個農夫.
他每天去耕種.
但後來他看到他的社區越來越艱難.
有很多人因為生活拮据.
就算有醫療保險.
但因為現在的藥物越來越貴.
有些人患了很特別的病.
要買一些特效的藥.
貴得很厲害.
就算有保險也買不起.
一顆藥丸要800美金.
有些人買不起.
他每一個月第一天.
他就將100美金拿去社區的藥房.
他靜靜地跟職員說.
我每個月一號會拿這100美金給你.
但唯一的條件你不可以告訴別人.
你找信封存起這些100美金的鈔票.
如果當中我們的社區有一個居民.
因為沒有錢要買一些特效藥.
你就用我這裡的存款.
他如是者一做.
做了10年.
直至到今年1月.
他安息主懷.
他的女兒知道了這件事之後.
他經過了一些掙扎和祈禱.
最後他都說.
他要將這件事公諸於世.
因為他很想透過這個見證.
去感動和鼓勵更多的人.
都要效法他的爸爸一樣.
要去關心他社區裡面很多貧困的人.
其實不單止在美國.

$^{1001}$在其他的民族.
在其他的國家.
我都涉獵過.
有很多這些非常好有心.
去幫助貧困人的這些人士.
我們少少一個舉動.
其實足以不單止是窩心.
是可以讓別人心痛感受.
我記起我們播道會恩福堂的創辦人.
蘇永志牧師有一次來我們講福堂.
講到的時候他分享過.
他多年前帶領一個.
我們香港一個很出名的企業家.
叫田家炳先生.
順主.
他其實順主之前.
已經是很熱心公益.
做很多製貧和輕學的事.
他順了主之後.
繼續有田家炳的基金.
他就在國內做很多扶貧.
也建了很多學校.
包括在香港.
這些中學他就很想給貧苦的子弟.
有機會可以繼續升學.
以至他們不要被扼殺.
沒有書讀.
如果他們要脫貧的時候.
必須要給他們一個機會.
而這些機會.
你會發覺有些很有心的人士.
他願意去這樣做.
我過去這十年裡面.
我都一直.
當我最初接觸到.
Batoto.
就在非洲的一個Africa.
這個Uganda.
他們一個孤兒院.
他們就訓練那些孤兒.

$^{1041}$鼓勵培育他們.
成為將來國家的人才.
讓他們由一個被遺棄的警方.
栽培他們成為一個.
在社會裡面有用的人.
我自己都助養了兩個孩子.
一直到現在.
看著他們大學甚至畢業.
然後他們又再來跟我說.
他們已經長大了.
他們現在可以自力更生.
你會不會考慮再助養另外兩個呢.
我說沒有問題.
每個月可以助養另外兩個.
最近我亦都.
去年亦都助養了三個殘疾的孩子在中國.
他們有很多這些慈善的團體.
都是將這些一分一毫去幫助那些有需要的.
這些困難的.
一些甚至被人遺棄的小孩子.
在我們今天這個社區裡面.
我相信都有很多有需要的人.
就是等待著.
神的兒女們.
我們不要失著連戌的心.
我們反而好像雅各書裡面的教導.
提醒我們真正的虔誠.
就是我們要看顧在患難中的孤兒和寡婦.
保守我們自己不顛簸世俗.
一個真正的敬虔.
就是不會因著你眼見到一些困乏的需要.
你看到一個人沒有衣服穿.
或者沒有東西吃的時候.
你會說我為你祝福.
願你平平安安去吧.
你講不出.
你都沒有可能這樣去祈禱.
你需要就是為他祈禱之餘.
你要祝福他的話.
你要關心他那個生活的需要.

$^{1081}$今天你會不會成為上帝使用你.
成為一個及時雨.
去幫助你身邊.
幫助在這個社區裡面.
一個有需要的人.
我為你們錦繡堂很感恩.
我知道你們有一個社區中心.
你們做得很好.
我為你們那份傳福音的心志.
為你們感恩.
求主繼續賜福你們.
使用曾牧師和你們的團隊.
帶領眾靈姐妹.
我們一起很好的關愛別人.
領人歸主.
我們一起祈禱.
因為我們看見在你自己的恩典.
你的愛當中.
你賜下你自己的救恩.
也感動你的兒女們.
我們要還這個福音的債.
我們體貼人的整個生命的需要.
求主你顧念我們這個城市.
也顧念在普世很多大城市裡面.
有很多生活在很大貧困的人裡面.
讓各地的教會都能夠去正視這個問題.
而很好的我們做這些社區關懷的需要的過程裡面.
我們將主你自己的救恩.
都帶給我們這些居民.
願主繼續施恩.
賜福錦繡堂.
賜福曾牧師和他們的團隊.
帶領靈姐妹.
成就上帝你自己的微兒.
願主你回來的時候.
都能夠我們每一個被你染中.
成為忠心良善的好管家.
同心立意的祈禱奉主耶穌基督得勝的名給.
阿們.
\newpage



\section{}
\label{sec:i9iPPyuXEWw}
\textbf{2023-02-02 周四晚查經班《被擄歸回歷史研讀(5)》 講員:曾裔貴牧師}
\newline
\newline
連結: \href{https://youtube.com/watch?v=i9iPPyuXEWw}{\texttt{https://youtube.com/watch?v=i9iPPyuXEWw}} ~~~~ 語音日期: 2023-02-04
\newline
\newline
\hyperref[sec:fpQjtet_7Iw]{\small{< < < PREV SERMON < < <}}
~
\hyperref[sec:index_chronic]{\small{[返順時目]}}
~
\hyperref[sec:index_scriptual]{\small{[返順卷目]}}
~
\hyperref[sec:2Us7kVIsCO8]{\small{> > > NEXT SERMON > > >}}
\newline
\newline
$^{1}$多謝天父再一次招聚我們.
能夠在這個平台.
有機會我們有系統地去研讀聖經.
來透過這樣的查考.
我們去明白一些歷史的書卷.
和一段相當長的.
被擄歸回的以色列人的歷史.
都幫助我們能夠在這些.
雖然已經過去了.
但是在那些關鍵的時代.
波斯王朝好幾位重要的君王的時代.
以色列人 猶太人.
他們是怎樣來到去.
在那個被擄的時間.
能夠歸回 重建.
聖殿 聖城 和聖文.
我們都盼望神給我們透過.
這些經卷.
去能夠掌握一些.
屬靈的教導.
今天我們都會多在.
先知書上面來探討.
求主都帶領我們大家的心靈.
我們這樣祈禱.
奉耶穌基督的名 阿們.
上一個禮拜我們開始有機會看了.
哈該書四篇的預言.
其實是默示.
或者是神的啟示給哈該先知.
他和撒迦利亞同時期的先知.
還有沒有其他.
當然就沒有在.
聖經的舊約裡面出現.
有些可能是其他的領袖.
都不出奇的.
其中要問的.
上個禮拜都提過.
到底他們是不是跟隨.
被擄歸回的呢.
他們是怎樣來的呢.

$^{41}$他們在一個關鍵的時候.
因為都拖了.
由古列王回去重建.
去到這個亞達斯西.
強令他們停工.
去到這個大利烏王的時候.
大概都有16年的時間.
停止了建造的工程.
可想而知那個方良.
就會帶來.
大家都不會有一個合一.
有一個一起為神的國度.
聖殿.
這樣來去同心的興建.
自然就會有很多的分化.
有很多的矛盾的出現.
固然十幾年.
那些憲制.
都是你有你做我有我做.
大家都不會有一個統一.
也都不會有一個一起.
守的一個大的節日.
直至到先知的出現.
先知的勸導.
勉勵的說話.
這樣來去激動了.
當時的省長.
同樣都是那個領袖.
索羅巴伯.
即是第七章所說的那個領袖.
一刻不記得他的名字.
就是.
可以說是波斯比的一個領袖省長.
帶來的就是一個民族.
真是一個同心的重建.
這是第六章.
近尾.
我們上個禮拜有機會看了那個.
第二個的詔書.
或者說是一個的上訴書.

$^{81}$是當地的猶太人的敵對.
其中以這個達賴.
這個河西的總督.
和這個士他波斯賴.
這些波斯的名字.
趙直譯.
和那些同黨.
來上書去到大理烏王.
反而這個上訴書.
即是控告書.
其實比第四章那個控告書.
那個告的狀詞.
相對溫和和持平一些.
第四章那篇的上訴書.
是相當針對性.
對猶太人的針對.
所以令到亞達捨西王.
改變了那個策略.
要強行猶太人要立即的停工.
反過來大理烏王收到這個告狀.
相對持平.
他就很客觀.
去尋查那些古籍.
果然找到先王古烈的時候.
頒布天下.
要將這些不同的.
避老的人.
回到他們原居地.
千萬不要得罪他們的神明.
就准許他們拿庫房的錢.
來買當地的材料.
來興功.
很自然的.
你做好那個聖殿.
自然就一定要有那個地方的.
他的禮聚.
那些禮拜的儀式.
很自然就需要那個祭司利未人.
所以第七章也是很合理.
以斯拉似乎都是有一個授權.

$^{121}$真的以一個祭司的身份.
上去耶路撒冷.
可能是做一個交流.
可能是帶著一個使命上去.
來將王的那些義師.
帶去自己的國家.
當然第七章開始就重點介紹.
以斯拉的出現.
他是一個怎樣的人.
他過去的裝備.
在波斯帝國.
這一位的祭司.
他是怎樣來裝備自己.
而帶來第七章.
一直到第十一章.
就是下半段的以斯拉記.
就是聖文的重建.
整個以斯拉就是這樣.
兩個大的段落.
一到第六章就是.
羅伯伯帶來的重建.
聖殿和聖殿的規模做好.
但是第七章就不再講聖殿了.
是講其中一次的節期.
應該是駐彭寺了.
他們來到去守完節期之後.
就有嚴肅會.
接著就有一個復興會.
這個大致上是以斯拉記的脈絡.
上個星期因為有兩個關鍵人物出現了.
在第五章.
就是哈該和易多的孫.
撒迦利亞.
兩位先知.
哈該書就四篇的預言.
應該叫做默示.
vision.
就容易一點.
撒迦利亞就有總共十二章.
撒迦利亞書就比較難解的地方.

$^{161}$因為它有很多很多的擬像.
那些擬像都相當難解的.
但是它有個特色就是.
有八個擬像.
是四對的.
頭尾的首尾呼應.
結構就簡單.
但是要解釋就很難解釋.
但是又給多一點點冷知識.
就是撒迦利亞書.
是啟示錄.
歷史上文術最多的舊約的先知書.
有很多的.
就是說撒迦利亞見到的擬像.
成為將來的所謂天啟文學.
啟示錄.
天啟文學.
他們可以直接套用的.
一些類似金句式的.
帶進他們的教訓.
所以啟示錄.
可能一脈相承.
因為神默示撒迦利亞.
他見到這些擬像.
所以去到啟示錄.
約翰見到的那些擬像.
很多大同小異是和撒迦利亞書有關的.
所以如果我們熟悉撒迦利亞書.
對我們看啟示錄.
相對又會比較容易一點.
今天就很想和大家.
來看第三篇哈該書的默示.
如果可以的.
第四篇也看完.
不單止是免禮重建.
上個星期我們說過.
有16年的時間.
神的殿方樑.
神就說.
這百姓.

$^{201}$已經不是我的百姓.
好像神對這一班的被擄歸回.
這五萬人.
好像有一個怨言.
十幾年都還這樣去蹉跎.
荒廢神的殿.
不需要再問是不是神的旨意.
因為已經被擄歸回.
一定是要重建神的殿.
但是因為政治因素.
當地的土著.
即是佔據了耶路撒冷的外族人.
本來說一起合作.
但索羅巴伯就說.
不可以,我們是猶太人的身份.
才可以為神的殿興功建造.
你們是無份的.
帶來的當然是進一步的行動.
結集政治的力量.
來逼令這個工程要停工.
所以他們就立即去告如莊.
寫信去古列已經過去了.
離世了.
亞達薩西去做王.
或者叫金拜西二世.
來逼令他們要停工.
所以一停就停到大理烏王的第二年.
我們可以看看哈該書這一段.
在馬太福音之前的兩卷書.
小先知書的最後那兩卷.
第三卷,對不起,尾三.
因為還有一卷叫馬拉基書.
哈該書的,這次是第二章.
上一次我們近期的時間就提到.
本來是六月初一.
大理烏王的第二年.
應該是主前的520年,方便去計算.
第二年的六月初一.
哈該就得到那個異象,默示.
來到他有神,好像一個對話一樣.

$^{241}$神跟哈該說話.
第三者就是這百姓.
就是在耶路撒冷的百姓.
他們不肯或者不敢.
為神的殿來奮勇的去建造.
所以後來當那個信息激勵了他們之後.
他們就真的上山取木料.
他們來去預備很多的材料.
不再懼怕.
當然同時間大理烏王發出那個詔書.
就勒令河西的總督.
要提供各方面需要的木材.
所以你會看到.
先是跟神的關係打通了之後.
原來政治那個層面.
又不太需要我們擔心.
不太需要我們去.
死就死吧,殉道就殉道吧.
又不是這樣發展的.
雖然波斯的帝國都不是真正的敬畏神的.
他們只是想有一種懷柔政策.
不得罪各自的神明.
但是不代表他們會有一種對神的新的認識.
在這個時候.
這一篇的信息帶來的就是.
在六月的二十四日.
他們就立即響應.
來興功建造神的殿.
所以這個就是一章的十四節.
哈該書我們在說.
耶和華就激動猶大省長.
撒拉鐵的兒子所羅巴伯.
和約撒達的兒子大祭司約書亞.
並盛夏之百姓的心.
他們就來為萬軍之耶和華.
他們神的殿作功.
這是在大利烏王第二年六月二十四日.
經過了二十三日的醞釀.
可能是教訓.
可能是開一些的發佈會.

$^{281}$帶來就是可能很多的領袖都先受感動.
有了那個默契和心靈的熱誠.
那就再推動那個事功.
所以不只是一篇的信息.
說完之後.
如果領袖都不去有一個動作.
其實很難去到最前線的那些的民工.
那些的平民百姓.
所以現在大家就興功建造.
不過在這個時候.
哈該書就不是停留在那些的實務的操作.
哈該既然是先知.
緊張的就是這一班興功建造的工人.
猶太人.
被擄歸回的這班人.
那個靈性同神關係.
現在就在第二章的第三個.
對不起.
應該是第二個的密使.
就是在七月二十一日.
七月二十一日剛剛完了駐彭節.
他記得很清楚.
猶太人的節期.
七月初一就是贖罪日.
對不起.
是摧國節.
七月初十就是贖罪日.
接著由七月十五一直到七月二十一.
就是駐彭節七天的.
第八天就會有一個大的嚴肅會.
所以去到七月二十一.
在駐彭節差不多最後那天的時間.
耶和華的說.
臨到先知哈該.
開始說釋法了.
對於神的摩西律法.
那個的什麼是潔淨.
什麼是聖潔這個範疇.
他們去為神做工.
才要有一個的認識.

$^{321}$所以哈該不是從那個民間的實務.
應該就是所羅巴伯來處理的.
比如要怎樣分工.
預備什麼材料.
那些圖織是怎樣.
這個就是所羅巴伯應該要做的事.
哈該緊張的就是百姓對神的認識.
緊張這班人.
回到上來.
你還有沒有帶著那些.
在外邦波斯帝國的色彩.
那些價值觀呢.
七月二十一.
耶和華的說.
臨到先知哈該說.
你要曉遇猶大省長.
撒拉鐵的兒子所羅巴伯.
和約撒旦的兒子大祭司約書亞.
記住是由先知發出那個密使.
發出給這兩位領袖.
一個是祭司的約書亞.
即是統管著祭司利未人的那個層次的.
在聖殿裡面服侍的.
負責獻祭的.
負責真是百姓來到去獻祭的時候.
他們要根據摩西的律法.
來去執行獻什麼祭.
有不同的做法的.
怎樣tong 法都不同的.
所羅巴伯就更加是實務性的領袖.
可能周旋在政府.
河西的那些官員.
和當地的土著領袖的長老.
譬如達賴.
那個是他波斯賴.
和其他的一些領袖.
去到黎希米的更多.
更加多這些人.
來衝擊黎希米.
不過好在黎希米本身是一個政府官.

$^{361}$省長.
所以就不懼怕他們多方的挑釁.
和來去試探.
甚至是危言聳聽的攻擊.
這裡告訴我們.
然後就是.
並盛夏的百姓說.
你們中間存留的.
有誰見過這殿.
從前的榮耀呢.
現在你們看著如何呢.
這樣的設問.
即是說就算你見過.
所羅門輝煌的聖殿.
所以有些人估計.
哈該和薩加利亞.
都是一百年多的老人家.
這個當然是估計.
現在神就好像告訴他們.
就算你怎樣做.
都不及得上所羅門當年的輝煌.
這個當然有它實際的意義.
你怎麼可能做一個聖殿.
要輝煌極之奢華.
比波斯帝國的還要奢華.
可能你會令到有關當局.
引起關注.
所以都盡量希望低調.
不要那麼鋪張.
不要那麼豪華.
可能只是基本的擺設.
是符合了摩西律法書的要求就算了.
你可能不需要再那麼貼金.
等等的安排.
所以有可能是這樣的意思.
未必代表哈該是一個很大年紀的.
然後就說現在你們看成怎樣.
其實還沒開工.
剛剛開工就已經說.
一定不會達到第一聖殿的輝煌規模.

$^{401}$那個格局.
但是重點在下面.
豈不看在眼中看如無有.
這個是一個事實.
第四節.
.
非常之大的.
好像一個三重的祝福.
連續說三次.
所羅巴伯你要剛強.
約書亞你要剛強.
你們中間的百姓.
你們都當剛強作供.
原因就是我和你們一起.
就是說神的默許.
神的時間.
神的激動.
不是人.
不是哈該一廂情願.
不是他很想在有生之年.
見到自己的事業.
得到成就.
本來已經很好了.
他的事業就是為神的殿.
已經很好了.
但是不是他個人的野心.
這個有時很微妙.
真的相當之難.
因為人心是很難測的.
是你個人野心.
還是你說為神的國度.
但是這裡就不需要想了.
因為是神三次說.
你們.
首先是個別的領袖.
個別的領袖很重要.
他的內心如果是被神感動.
剛強起來的時候.
他為神做的.
是事半功倍.

$^{441}$因為那個力量很大.
一個人.
他如果能夠得到他的心靈.
是徹徹底底地去降伏.
信服神.
所以今天.
全世界.
我想全個香港.
或者教會.
如果能夠得到一個青年人的心.
那個心是被真理.
被神改變.
感動.
那個就是一個非常大的力量資源.
所以其實是全世界的領袖都想得到.
無論是什麼政權.
都想得到青年人的心.
得到他的心.
就得天下了.
真的說得不是誇張的.
所以我們這一代的年青的信徒.
如果是被神激動復興.
我們根本不需要擔心.
膠棒或者斷層.
根本不需要擔心.
他們會有神在他心靈裡面的激動.
當然現在可能是iPhone.
可能是去旅行.
可能是很多其他的東西.
譬如創業.
都不是錯事.
但是又是他怎麼擺他自己人生的次序.
如果青年的時代就已經對神有一個很強烈的渴慕.
他現在所做的一切.
都是為神準備.
等待神的呼召.
那我們就不擔心下一代教會的處境.
和前途.
我們要得的是人心.
所以現在神就得到所羅巴巴的心了.

$^{481}$透過哈該講的勉勵的說話.
所羅巴巴你要剛強.
若孫雅你要剛強.
百姓你們都要剛強的作功.
即是輔助配合兩位領袖的帶領和牧羊.
然後因為神的同在.
這個是萬軍之耶和華說的.
我都提過了.
哈該書是出了很多次萬軍之耶和華如此說的.
不知道是不是因為特別在這個階段.
要激勵他們的心.
這個時機可一不可再.
大利烏王的第二年.
那個王突然之間又有一個很客觀的想法.
就是去查那些典籍.
查到有先王古列的詔書.
所以就准許他們回去重建.
已經回去了.
准許他們在河西那裡調動那些國庫的資源.
但是會不會突然間又改變政策呢.
那是不知道的.
所以在這個時機.
原來我們從另一個角度.
經常都講就是一個平衡的角度.
天上面的和地上的.
所以不要只是看.
哦現在開綠燈.
因為王比較懷柔的政策.
和比較善待這些被老的人.
這個是一個角度.
但是原來是因為屬靈的領袖.
是因為神的選民.
他們開始認定自己的身份.
就帶來一個復興.
首先就是硬件上他們不再關心自己的房屋.
而是大家一起將精力.
他們的時間.
就放在建造神的殿這個情況上.
這個就是神的話語激動.
所以萬軍之耶和華如此說.

$^{521}$是一切問題的迎刃而解的答案.
如果在教會有萬軍之耶和華如此說.
一切就不用開會了.
也都不需要這樣去.
你又這樣講.
他又這樣講.
大家又不同意.
再嚴重一點.
又學習得不好的.
就拍桌子離場.
根本就不應該是這樣的.
但是這個過程.
隨著舊約.
一直去到新約.
去到聖經的完備.
去到聖靈的降臨.
去到我們經過2000年的教會歷史.
等等去到今天.
萬軍之耶和華如此說.
就要看你怎樣去將它演繹.
現在靈恩派.
如果比較嚴重的靈恩第四波.
他們要恢復先知的即是.
現在有新使徒.
新先知的運動.
所以你可以去他們那裡.
去學新先知的言語和使徒.
其中現在在美國.
多倫多現在就應該沈寂了一點.
多倫多機場那個叫Toronto Blessing.
那個第四波的運動相對比較沈寂.
但是在美國還有一個教會.
都是這樣的.
他們現在如果我沒有記錯.
就是24小時不停的祈禱.
要不停有人輪值的禱告.
本來不是錯事.
不過有時候一些的mindset.
未必是聖經所說的那個祈禱觀.
但是最大的問題就是.

$^{561}$他們要恢復新使徒運動.
即是說我有這個authority這個權威.
我講的就是耶和華如此說.
其實point就是想拿回那個獨裁的權柄.
也都違反了聖經一向以來.
那個真是舊約的神諭.
即是說oracle神諭的默示.
聖經也都不是容許.
哈該,撒迦利亞,以賽亞,以西結.
他們就講完了神諭此說之後.
他們突然之間就會有了一個.
好像高人一等的權威.
從來都不是的.
相反就是很多的舊約先知.
是死於自己的百姓的手中的.
所以沒有.
因為這樣來為自己的身份增值了.
所以神的話和那個先知本身.
沒有因為這樣講了這些神的預言.
他就可以得到一個凌駕於百姓.
或者是那個聚會的人的身份.
這個要搞清楚.
所以這個是其中一個盲點.
你說自己是新使徒.
你說你有先知的言語.
你說的話就是耶和華如此說.
這樣講已經是將自己擺在一個很危險的.
好像將自己擺在神台一樣.
非常危險的.
誰夠膽說在新約還有這樣的耶和華如此說呢.
或者新的使徒.
好像保羅這樣的authority呢.
因為是有他的特殊性的時代的背景.
所以不可以這樣隨便說要復興使徒.
還有那個先知言語.
所以講到耶和華如此說.
就要來給大家有一點點的概念.
其實都是一個複雜的課題.
接著我們就看看下面.
對不起第五節.

$^{601}$(主持人翻譯).
雖然規模不及第一聖殿.
但是這裡講到的珍寶.
這裡是眾數來的.
眾多的珍寶.
是會被神激動運送來.
而且這個殿後來的榮耀.
還多過之前的.
相信這裡講的不是那個物質的規模.
是在講那個屬靈的層面的.
將來的這個聖殿.
要承載的那個永恆的救恩.
所以那個聖殿的硬件.
是必須在耶穌來之前重建的.
建好之後.
猶太人又在那裡來到日籍的聚會.
去到耶穌出現的時候.
大家都記得了.
那個殿就變成一個.
好像人民公社一樣.
大家都可以在那裡出入.
沒有了那種神聖感.
在那裡買賣.
在那裡兌換銀錢.
因為在外地上來守節的人.
其實就很不方便.
等於我們七八十年代.
見到有很多香港人.
就要托一些電視上去.
因為大家那個生活水平相差太遠.
香港就富起來.
內地還是一窮二白的時候.
我們要擔一些油上去.
擔一些米.
然後就托個電視上去.
後來就不用那麼麻煩了.
代用券就OK了.
這個代用券制度.
就在聖殿就實施了.
就是你上到去.

$^{641}$才現成的去買那些牛羊.
一來就是你不會.
那條路途遙遠.
那些牛羊就會受傷.
等等.
還有兌換銀錢原來很方便.
因為你上到去聖殿.
聖殿那些祭司.
守殿官就不理你.
你是哪裡來的.
麻煩你要兌換希伯來的貨幣.
所以就衍生了這些攤檔.
就幫人家兌換了.
等於你要上深圳.
你要換人民幣了.
這個就變成在耶穌的時代.
其實已經不是耶穌的時代才有.
是大家見怪不怪.
大家也都覺得很方便.
所以聖殿就變成了一個公社.
大家都可以出入.
沒有人覺得那個聖殿.
需要有一種的莊嚴.
耶穌就賦予那個聖殿.
其實某程度是很討人厭的.
本來是方便利民的政策.
耶穌就看為是令到他的殿.
不能夠成為萬民禱告的殿.
這樣你就可以想像得到這個情況.
所以這個時候.
神就預言這個殿是充滿了榮耀的.
這個是神這樣說的.
其實在一段聖經裡面.
已經出現了四至五次耶和華如此說.
是萬君之耶和華如此說.
這個殿後來的榮耀.
所以榮耀和平安.
是好像孿生兄弟一樣.
先是經驗神的榮耀.
人就會得到心靈的平安.

$^{681}$我相信弟兄姊妹.
如果你真是帶著一個很虔誠的心.
來敬拜神.
你又透過詩歌.
透過訊息.
給你有一個可能是一個感動.
你走的時候.
你的內心就不是有一份榮光.
而是有一份平安在你的內心.
你總是覺得很舒服.
好像很多的困難.
暫時放下了.
雖然你去吃飯.
但你的腳步都很輕盈.
你都好像有一份內心的平安.
湧入出來.
最好就是希望能夠會所的職員都看到.
嘩!你為什麼這麼平安?.
這樣就證明你有榮光了.
因為人家都感受到.
你有那個內心.
去完教會之後.
那份平安.
這裡聖經是這麼說的.
第九節.
這殿後來的榮耀.
必大過先前的榮耀.
在這地方我必設平安.
這是萬君之爺和華說的.
所以每一個奉神的名聚集的地方.
你不要管它是一個二樓教會也好.
是一個大禮堂也好.
有神的同在.
有神來到這裡的時候.
他就會設平安.
所以我們要記得.
每一次我們是與神聚會.
不是與人聚會.
當然我們不要這麼極端.
曾牧師你請你走到一邊.

$^{721}$我們是與神聚會.
不應該這樣的.
我們是應該帶著一個對神的敬虔的朝見.
朝聖.
所以朝聖之旅是對的.
但是我們要記得.
我們都是要跟人相交.
因為神都希望我們彼此之間的相愛.
這個就是第二篇的訊息.
7月21日.
重點是安慰他們.
你們雖然做一個規模很小.
但是我應許我與你們同在.
我應許這個殿後來的榮耀是很大的.
列國的珍寶都會運來.
當然我們見不到這麼多.
但是相信是說後來不同的民族的人.
可以在神的殿裡面信住.
我的老師龔永謙牧師說.
猶太人對於聖殿有一個很強烈的想法.
就是在所羅門的殿那裡.
是直通天上的天庭.
就是天家.
我們新疆叫天家.
真正的聖殿.
因為地上的只是一個模型.
真正的是天上.
所以為什麼一定要在那個方位.
在所羅門建造聖殿的地方.
原址不可以第二個地址.
因為那裡是直接連通天上.
用我們今天的言語我們就明白了.
就是那個WIFI特別強.
直接可以去到.
所以為什麼主耶穌說.
這個殿是萬民禱告的殿.
不是說那些東西是錯的.
但是如果你將那個重點擺開了.
那就有一個很大的問題了.
這個就是第二個勸勉的這個牧師.

$^{761}$就是7月21大利烏王的第二年.
7月21.
就是哈該領袖的.
就是安慰的訊息.
就是說當百姓有心重建.
神的安慰就藉著先知來.
給他們看到.
第一個的預言那個牧師.
就是責備帶著一份的.
是很無奈的.
這百姓說還不是時候.
其實是神和哈該訴苦.
但是第二個的安慰.
因為他們起來建造.
6月24起來建造.
7月21就立即有第二個安慰的牧師.
或者叫做Vision.
那就去到第三個.
其實跳開第二個之後.
第三個是在薩加利亞書.
就不是哈該書.
就是說神原來好像多工作業的.
又給哈該又給薩加利亞.
我們很快站一站去旁邊那一頁.
因為旁邊那一頁就是薩加利亞書.
第一章大利烏王的第二年8月.
大家都是用回波斯帝國的年號.
就是說今天是農曆十一.
類似吧.
我忘記是不是十一還是十二了.
大家都大概用得到.
不是說你用新曆我就用舊曆.
大利烏王的第二年8月.
耶和華的說臨到亦多的孫子.
比利加的兒子.
先知薩加利亞說.
耶和華曾向你們列祖大大發怒.
所以你要對以色列人說.
萬君之耶和華如此說.
你們要轉向我.

$^{801}$我就轉向你們.
這是萬君之耶和華說的.
不要效法你們列祖.
從前的先知呼叫他們說.
萬君之耶和華如此說.
你們要回頭離開你們的惡道惡行.
他們卻不聽也不順從我.
這是耶和華說的.
你們的列祖在那裡呢.
那些先知能永遠存活嗎.
只是我的言語和律例.
就是所吩咐我僕人眾先知的.
豈不臨到你們列祖嗎.
就是說人可以過去.
神的預言.
神的警告.
當然舊約的先知.
在他未被擄之前.
大部分特別是臨近被擄的時候.
譬如耶利米時代.
不是以賽亞的時候.
譬如阿摩斯的時候.
大部分是警告責備.
再不回轉.
神就會出手.
就會來應罰.
所以這裡再說.
似乎又再透過先知.
給他們一個比較重的說話.
就是說要一鼓作氣.
7月21日的安慰說話.
給他們一個信心的保證.
就是我的同在.
我的榮耀.
我的平安.
和我的名在這個聖殿裡.
在你們這群人.
神的選民的心靈裡.
所以你們就勉勵去做吧.
但是到了8月.

$^{841}$又有另外一個人.
另外一個先知.
另外一個預言.
這次就是重提舊事.
重提列祖過去的失腳.
失敗的地方.
就是說前車之切.
我們就要留心了.
不要再重蹈覆轍了.
就是說那個馬車.
已經走了一條車輪的痕跡.
你就不要再跟著它走了.
這樣的意思.
所以一到第六節.
是神透過撒迦利亞.
去重提這個民族.
過去的失敗的歷史.
作為一個尷尬.
他們就真的要照著.
現在的先知的勉勵的說話去做.
所以這裡就告訴我們.
剛才我讀到第五節.
那些先知是不是能夠永遠存活.
沒錯.
先知那個人本身離開了.
但是神的話語就永存了.
尤其是擺在聖經裡面.
「只是我的言語和律例.
就是所吩咐我僕人眾先知的.
豈不臨到你們列祖嗎?」.
豈不臨到就是說.
你們眼睜睜見到耶利米的警告.
果然被擄去到外邦.
果然第一聖殿.
所羅門的聖殿.
是被拆毀.
移為平地.
什麼都沒有了.
亞述王將那些外族人.
遷居遷徙來到這個撒瑪尼亞.

$^{881}$接著巴比倫王又將那些人.
遷到猶太耶路撒冷.
就是說有70年.
整個以色列地.
神選民的地方.
已經是無論聖殿有沒有了.
聖文有沒有了.
完全應驗了先知的預言.
先知都過去了.
那些預言都還沒有發生.
但現在你們這班回來.
被擄歸回的人.
你們是自己親自經歷了.
有幾代人要在外地飄零.
這樣來等候復國.
所以說了這麼多.
神的話是會善良的.
應許的.
就會追討那個人.
直到千代.
那個祝福.
但是如果那些的咒詛.
那些的審判.
一樣是會臨到.
致富及子直到三四代.
這個是《出埃及記》第20章的.
十誡的經文.
所以這裡就告訴我們.
勉勵的同時.
要他們追上.
哈哥書都提過的.
要追上以前的日子.
要追上我們的列祖.
所以我們真的要常常懷念.
我們的前輩.
他們在創辦錦繡堂的時候.
他們的心靈.
有些還在的.
譬如田先生.
或者還有一些比較元老的.

$^{921}$還在的.
譬如坤勝這些.
都是早期的弟兄姊妹.
他們的心我們要效法.
我們要看到他們.
怎樣單純為神的國度.
為神在這個社區.
將來要有教會.
這樣來擺上.
所以我們要記得.
當然在四十年裡面.
有個什麼問題.
重大的問題.
是神不喜悅的.
我們就不要當看不到.
聽不到.
我們都應該作為一個尷尬.
我們要避免的.
因為這段聖經這六節.
因為下面由一章第七節.
一直到整本書的結尾.
全部是異象.
但是頭六節.
就是一個序言.
就是撒迦利亞書還沒說異象.
他還沒見異象之前.
就再一次.
來跟以色列人.
重提舊事.
過去列祖.
不肯聽先知的預言.
尤其是警告的預言.
就落得這樣被老的下場.
所以現在.
先知又再一次有神諭.
有神的啟示.
那個預令.
來到去.
其實都是勉勵.
叫他們現在要認真聽從神的律法.

$^{961}$這個就是第六節的.
最重要的部分.
指示我的言語和律例.
就是所吩咐我僕人.
眾先知的.
豈不臨到你們列祖嗎.
一定是.
他們回頭說.
萬君之耶和華.
定義按我們的行動作為.
向我們怎樣行.
他已照樣行了.
這裡好像給他們一個立約.
他們願意.
因為在這個時候.
有兩位先知.
有兩位領袖.
在百姓當中來勸勉.
好像現在那些的延習會.
二十大的延習會.
讓到各級的官員.
都能夠明白那個精髓.
所以這裡再一次提醒他們.
要效法先知.
聽從先知奉神的名.
所說的話.
他們就有運行了.
他們就可以.
來到去照著神的祝福.
他們就可以得到榮耀.
榮耀當然帶來的就是.
那個人的平安.
也都得到神的眷顧.
所以沒錯.
周邊的列國.
其實蠢蠢欲動的.
波斯帝國是強盛了一段.
第四任都已經開始.
慢慢衰落了.
有一個很大的勢力.

$^{1001}$興起中.
就是亞歷山大.
這個大帝.
差不多是世界.
其中一位.
簡直是.
神級的領袖.
軍事天才.
不過三十多歲就死了.
又是英年早逝的一個君主.
一個很出色的奇才.
希臘大帝.
又是將馬其頓.
很多的城邦.
又是這樣吞併.
籠絡.
成為一個大的帝國.
所以你看列強.
雖然政治上.
都是風起雲湧.
但是百姓的心.
不要看外面.
看自己內部.
先知的眷面.
就有一個很大的凝聚的力量.
所以這個都很重要.
至少福音派還是相信.
一個崇拜.
最重要的部分.
是講道.
雖然現在用很多禮儀學.
每一個環節.
連那個花擺都有禮儀.
有它的重要性.
我又不反對.
但是.
永遠都是用.
大家聽得明白的.
勸勉的說話.
用聖經的講解.

$^{1041}$來打入.
一群回眾的心靈.
至少大家都仍然認定.
我坐定聽講道.
是.
想明白神.
今天跟我說什麼.
至少還是一個信念的情況下.
崇拜的講道.
仍然是這麼重要的原因之一.
所以.
你可以說.
聖經撒迦利亞.
不是打沈他們的信心.
已經安慰了.
已經說我們同在.
我們現在努力去做.
為什麼還要重提.
舊事.
還要數我們的列祖.
然後透過撒迦利亞.
來勉勵他們.
這一群的百姓.
真的不要.
重蹈列祖的覆轍.
不要來到.
自己去.
理解.
那個局勢.
用自己的手法.
手段.
世界的價值觀.
而不聽先知的說話.
當然你想想.
耶利米.
他得到的意義.
那個預言就是說.
沒用的.
就算你少收少補.
神都不收貨.

$^{1081}$所以他得到的預言.
就是叫那些百姓.
投降.
巴比倫.
投降.
但是君王西底迦.
和西底迦.
身邊的先知.
就不斷在製造.
一些預言出來說一定打贏.
耶利米賣國賊.
那你就可以想像.
到底誰才是.
真的先知.
是很複雜的.
在當下的時代.
能夠應驗到.
所以也可以這樣說.
第五節.
神的說話.
你們有沒有聽從.
你們列祖在哪裡.
那些先知是否能夠永遠存活.
不過神的.
言語和律例.
說了的.
原來是要很後期.
才真正的兌現.
後來才懷念.
耶利米是真的先知.
原來被擄投降.
才是神的旨意.
是神的計劃.
因為是不能再悔改了.
必須要受那個刑罰.
那個就是出路了.
一定要體驗一下.
去到異邦.
生活七十年.
整個耶路撒冷.

$^{1121}$七十年的安息年.
地要得到安息.
你可以想像到.
這個刑罰是多麼的大.
多麼的重.
所以神的說話.
那個震撼就是這裡.
所以其實要祈禱的.
讓教會的.
當然全香港.
或者甚至乎現在.
有些移民去英國.
譬如林姑娘.
其他的牧者.
也為每一個牧者.
每一次的講道祈禱.
真是得到神的心意.
去講.
是講神的心意.
是和一班相熟的群體.
弟兄姊妹.
去苦口婆心的.
去講神的說話.
神的安慰.
神的告誡等等.
希望大家仍然維持在.
這一種.
傳遞的水平.
這樣的一個.
大家的身份.
牧者的身份.
信徒的身份.
大家去領受.
領受神的教導.
今天是不會有人.
告訴你耶和華.
透過我如此說.
必須這樣做.
仍然是將聖經.
透過我們的生活.

$^{1161}$我們的體驗.
我們的經歷.
我們的經文的背景.
背景是古代的.
但如何將古代的經文的背景.
套用在現實的生活裡面.
你和我的實際的應用.
一篇道.
不同的人都可以有.
他自己的領受.
例如我做領受的.
例如我做其他的事情的工作的.
那篇道可能是一個原則性的.
我如何能夠在這個原則裡.
我在我的生活的層面.
去實踐.
去運用.
每個人就有不同的感動.
這就是牧者的責任.
將古時的道.
已經成為了聖經.
來去詮釋出來.
有些是很應那個時代的需要.
要說的.
有些是按部就班.
來去教導.
例如我們現在選取.
被擄歸回這段歷史.
我們就一起進入那個時代.
在其中.
去找回他們的.
種種的脈絡.
希望我們能夠從其中.
吸取到.
所以你會見到神擺設的.
有感動的.
那個領袖.
行政的領袖.
這個叫切巴薩.
我記得了.

$^{1201}$也有大祭司.
因為回去重建的是.
聖殿.
所以所羅巴百是不可以獻祭的.
連仙姿哈該都不可以獻祭的.
仍然都只能夠.
是大祭司代表著.
整個民族去到.
聖殿獻祭.
利未人是負責那些.
樹木預備那些水.
那些柴.
預備那些陳設餅.
或者百姓上來獻祭的時候.
他負責做檢驗等等.
然後就交給祭司.
去做那個宰殺.
和彈血獻祭.
仙姿的角色就是.
神興起的.
來對律法書.
尤其是摩西的律法書.
有認識的.
和他會將神的.
那個啟示.
傳遞給會眾.
這個就是.
我們現在去到這樣說.
已經有第三個的.
默示了.
回到大利烏王的第二年.
第九月又有另外的.
默示.
一共有兩次的默示.
同時間在.
九月廿四.
大利烏王的第二年.
在耶路撒冷原來.
燈火鼎盛.
香火鼎盛.

$^{1241}$這麼多一下子.
由六月初一.
接著六月廿四他們就開工.
當然不用撤銷.
開工.
接著就七月廿一.
到第二年的八月.
現在去到九月廿四.
又有什麼默示呢.
這次是律法的釋法.
是這個群體.
需要明白.
為神做事.
需要有那個聖潔.
大利烏王第二年九月廿四日.
耶和華的話.
臨到先知.
哈該說.
我們回到第二章.
哈該書.
第二章第十節.
今之耶和華如此說.
你要向濟斯問律法.
叫哈該問律法.
濟斯即是.
這個.
若曰答的兒子.
問什麼律法呢.
即是說.
假如若有人用衣襟.
兜聖肉.
這衣襟挨著餅.
湯酒油.
或別的食物.
算不算聖呢.
是否變算為聖呢.
濟斯說不算為聖.
我們一會兒再解釋.
哈該又說.
又再問.

$^{1281}$如果有人摸了死屍.
染了的污穢.
然後挨著這些物的那一樣.
即是上文所說.
例如.
挨著那些餅湯酒油.
或其他食物.
濟斯說.
必算污穢.
留意.
不算為聖.
和必算污穢.
首先是說了兩個.
假設的問題.
真正的重點.
是十四節.
於是哈該說.
耶和華說.
這民這國在我面前.
也是如此.
甚麼叫也是如此.
哪一方面也是如此.
算為聖潔還是污穢呢.
他們手下的.
各樣工作都是如此.
他們在壇上.
所獻的.
也是如此.
出現了三次.
也是如此.
即是要問那個問題.
第一個假設.
就是那件.
制服.
因為是.
對不起.
我搞錯了.
一個普通人.
他拿他的衣襟.
即是那些長袍.

$^{1321}$就繞起.
將那些.
記住是聖肉.
和其他.
似乎是餅湯.
和酒油.
都是可能和聖殿獻祭員有關的東西.
就問.
那些.
食物.
算不算為聖潔呢.
不算.
因為已經是用一個普通人.
的衣襟來.
拿了.
然後.
第二個就是說.
如果有人摸了死屍染了的污穢.
然後再.
挨著那些聖物.
那些聖物.
還算不算為潔淨呢.
算不算污穢應該直接問.
必定算為污穢.
因為那個死屍本身.
碰一切東西都變污穢.
那個凡人.
不是祭司.
什麼的.
用他的衣襟.
來拿著那些聖物.
那些聖物.
都不算為聖潔.
這樣的意思就是說.
似乎這裡也是.
如此應該這樣理解.
怎樣理解呢.
就是.
神就自己在第十五節.
去解釋了.

$^{1361}$現在你們要追想.
在這一天之前.
可能就是九月廿四.
耶和華的殿沒有一塊石頭.
累在.
即是疊著在另一塊石頭上的光景.
在那一切日子.
有人來到谷堆.
賞得二十斗.
只得了十斗.
賞得五十斗.
但是只有二十斗.
在你們手下的各樣工作上.
我以旱風.
霉難,冰雹.
攻擊你們.
這些是天災來的.
那些風是旱風.
由藍地吹上來.
因為大家猜到了.
阿拉伯沙漠,埃及.
吹上來那個叫藍地的風.
就是非常乾旱的.
即是死海那一帶.
霉難.
即是說.
可能有些瘟疫.
令到那些修成都不行.
冰雹即是極端的天氣.
那個時候.
你們仍然都.
不歸向我.
這是耶和華說的.
你們要追想,第二次提到追想.
十五和第十八節.
要追想.
此日以前.
就是從這九月廿四日起.
追想到.
立耶和華電根基的日子.

$^{1401}$就是他們上去的時候.
這個.
古烈王的時候.
班趙書上去的時候.
他們立了電的根基.
這十幾年的時間.
由這一日.
即是說這個時候.
9月24日.
追想回去立根基的時候.
追想到.
立耶和華電根基的時候.
倉裡面有菊種.
葡萄樹.
無花果樹,石榴樹.
橄欖樹.
都沒有結果子.
但是從今天起.
我必賜福於你們.
說了這麼大堆東西.
又解釋那個律法.
即是說.
普通的人.
就算他.
拿了那些制慾.
聖殿獻祭的制慾.
都不會算為聖結.
相反.
污穢的人.
碰任何東西都會污穢.
即是意味著.
如果百姓的心靈.
不是對準神.
那做外表的事.
都不會帶來真正的祝福.
即是說.
如果用我們今天的.
社會.
已經不再是整個群體.
很個人主義.

$^{1441}$特別是後現代的世界.
我就是.
那個話事.
因為我的感覺是最重要的.
現在據聞.
美國就是有.
中性的廁所.
甚至極端到.
我什麼時候覺得.
自己今天是.
另外一個性別.
我就可以去廁所.
即是.
好像要緊張.
你那個人的感受.
大過.
天生的結構.
那個結構都可以改變.
例如變性.
你就可以想像得到.
和這個時代.
古代的時候.
他們始終還是一個群體.
始終還是要依附著.
整個民族.
沒有個人主義可言.
但是.
他們都要面對.
整個群體.
其實神只是說十幾年.
大概十六年左右.
最多頂多去到二十年.
由這個.
興宮建造.
神的殿的根基.
一直到大利烏王的第二年.
九月廿四.
真是十六至二十年.
差不多.
但是原來在他們.

$^{1481}$一直都為自己去繼續.
經營自己的家庭.
那個工作.
其實是勤奮的.
你要周旋.
當地人.
你又要重新去適應.
那個氣候.
重新要殺種.
重新要來改良地圖.
你拿回你自己的地.
是不是真的分到你自己.
祖宗給你的地呢.
這些細節我們又不是太知道.
沒錯.
大家好像去.
伏奸那樣.
不過.
出來的後果.
神就說.
原因是.
剛才說的不算為聖潔.
算為污穢.
原來他們就算.
回到聖城.
建造聖殿.
如有挫折.
停工等待.
原來神說.
你們還沒排好.
次序.
最重要的次序.
當他們現在.
六月初一.
排好次序之後.
六月廿四.
大家一齊.
同心興工建造.
接著去到九月廿四.
三個月.

$^{1521}$神就再一次提他們.
提他們.
其實真正的聖潔.
我們知道.
舊約的時候他們不太知道.
他們每次都看律法書.
原來神說.
真正的神選民.
是要心靈對神.
要有一個專一.
和對神.
真的從心底裡.
對神要有這個.
效忠.
對自己選民的身份.
意味著.
不是這麼簡單一兩條律法.
而是涉及你們這個民族.
有沒有再拜偶像.
有沒有再和.
當地人通婚.
第七章.
開始處理這個問題.
第七章有一個.
很厲害的先知.
應該是.
牧者.
這個以斯拉.
其實他是祭司.
由波斯上來.
接著就發現.
住過彭捷之後就發現不對勁.
他就來到.
自己一個人.
安靜的來到去坐.
一直坐到獻晚祭的時候.
這個第七章的情況.
我遲些再說.
你可以想像到.
9月24日這個第三篇的訊息.

$^{1561}$是.
告訴百姓知道.
很多事情.
是有跡可尋.
可以見到那個問題.
即是離開神.
其他事情都會很混亂.
沒有了.
那個最高的心意.
去尋求神.
一定你只會.
很多事情忙.
沒有事都可以忙.
更多的就是.
熟悉地上的事情.
是會高過神的.
那個時間表.
偶然間你不舒服.
回不了教會.
當然沒有問題.
但如果你的時間表.
全年的時間表.
回教會只不過是.
節期你才回.
我認識一個同學.
他的教會.
去到聖誕節.
12月25日.
是突然間多幾百人.
原來他們逢聖誕節才回教會.
很特別.
我的同學.
都在教會長大的.
在他的母會長大的.
都很熟的.
問為什麼.
很多合理的原因.
我的工作很忙.
星期日我已經很累.
小朋友又要上學.

$^{1601}$現在星期六又要補習.
基本上星期日才有.
少許時間家庭時間.
有些說要進修.
很忙碌.
要相親.
將時間表排到.
有多低就多低.
但你不要理.
聖誕節一定要回.
要和教會有一個連繫.
見一下老友.
還有一樣東西一定會找教會.
怎樣都不會不找.
就是當有喪葬.
一定找教會.
就是這個教會的形態.
我不說哪間教會了.
這個就涉及時間表.
我們所謂.
priority的次序.
人生的次序.
我們那個週期.
是自己定的.
即是說我們一定要.
讓弟兄姊妹看到.
什麼是算為聖潔.
什麼是算為污穢.
所以神就很有趣.
就是這個的寓言.
先是透過哈該.
來問祭司.
那個律法.
先是設局.
祭司你答我.
這樣幫我釋法.
這樣算不算為聖潔.
祭司說一句就即刻.
快而準.
當然不算.

$^{1641}$你那些聖物.
沾染了日常的東西.
就不會再是聖潔.
即是說這個民族本身.
如果他們做.
神的工作.
那個心不是真的為神.
是一種不聖潔的後果.
污穢就是.
這班人.
他們如果帶著這個世俗.
根本心不是這樣來為神.
所做的事都會帶來污穢.
或者他們所做的事.
是會沾染.
其他人的污穢.
變成污穢.
所以這個釋法.
其實就是神終於解釋了.
他們要追上.
所以這個百姓.
這十六年.
也是如此.
出現了三次.
你們所做的一切也是如此.
所以出現了三次追上.
三次的也是如此.
三次的你們要追上.
所以其實我們自己.
每年都真的要.
自己反思一下.
一年過去.
真是.
越來越.
離開見主面.
越來越近.
身邊的人.
可能跟你很近很熟.
都真是突然離開.
其實還覺得是做夢.

$^{1681}$不覺得是現實.
原來他真的不在了.
但是.
有時候自己反思的空間.
不夠大.
但這裡就告訴我們.
神不是要戳穿他們.
而是告訴他們一個事實.
你們這麼賣力的.
栽種.
賣力的.
做這個工程.
但是.
正正就像上文所說.
你們是用.
依侵來接觸.
那些聖物.
你們用你的污穢.
因為你們接觸了死屍.
污穢來碰到其他東西.
就是說.
有一種群體彼此的傳染.
那個傳染.
不是敬虔的感染.
而是.
世俗的互相影響.
這樣帶來的就是.
你們知道那個結果.
你們要追想.
你們能夠有的什麼光景.
神的殿方娘的時候.
你們.
那一切的日子.
你們能夠得到的.
都是好像.
竹筋劍爪那樣.
你們都不能夠.
有真正的滿足.
原因是什麼呢?.
因為你們都沒有歸向我.

$^{1721}$你們要追想這個日子之前.
所以在這個時候.
似乎.
復興.
這個安慰.
是入心了.
真的去到百姓的心靈.
他們真的看到.
追想.
有一個反思.
來跟隨.
神的僕人的教訓.
所以在這個時候.
很重要的.
第十九節的下半節.
從今天開始.
我必賜福.
與你們.
最後第四個.
的默示就是.
向所羅巴百.
同樣都是那天.
第二十節.
這月二十四日.
主話二次臨到哈該.
你要告訴猶大省長.
所羅巴百.
我必震動天地.
我必傾福列國的寶座.
除滅列邦的勢力.
並傾福戰車.
和坐在旗上的.
馬必跌倒.
騎馬的必敗落.
各人被弟兄的刀所殺.
萬軍之耶和華說.
我僕人.
留意一下這個身份.
第一次出現.
我僕人撒拉鐵的兒子.

$^{1761}$所羅巴百.
到那日.
哪日是哪日呢?.
建造好聖殿.
或是遙指將來.
主耶穌的那日.
我必以你為人.
一個很重要的.
通常是一個領袖.
蓋印.
古代就更加明顯.
國王的御骰.
一鎚定音.
或是戒指.
就是那個.
生殺大權.
所以法老王就給戒指.
約瑟.
就等於他把權力.
移交給外邦人.
希伯來人.
掌管在.
豐年的時候.
那個國家的命脈.
即是說他講什麼.
他發佈什麼.
王榜諭旨.
所有埃及人都要聽.
在這個時候.
所羅巴百將來會成為.
神的印.
因我揀選了你.
僕人揀選.
這是萬軍之夜.
和話說的.
四篇的信息.
就這樣完了哈該書.
當然哈該不是說四篇.
就完他人生.
而是應該一直都會在民間.

$^{1801}$陪伴著他們建造.
那個聖殿和一直.
看著他們那個.
靈明的工程.
不過第四篇是特別針對.
所羅巴百.
我們就要對這個人多一點點瞭解.
就會好一點.
還有時間和大家揭一揭.
我來到去揭了.
我們先看歷代至上.
和這個所羅巴百的背景.
有些關係.
歷代至就在.
列王記夏後面.
第三章的歷代至.
頭九章的.
歷代至如果你沒有什麼.
信心沒有什麼.
毅力你都讀不到.
因為全部是家譜.
是悶了一大輪.
悶到你都沒辦法想像.
基本上是很難讀.
第十章才開始.
又開始記載那些.
烈王的事跡.
第三章.
就不可以看太長了.
就和猶大支派.
有關的這個.
所羅巴百.
三章.
或者由所羅門的後裔開始看.
第十節.
所羅門的兒子是羅波安.
羅波安的兒子是亞比亞.
或者叫亞比央.
亞比亞的兒子就是阿薩.
阿薩的兒子就是約沙發.

$^{1841}$然後是若蘭.
若蘭就是去到阿哈謝.
若蘭就是.
娶了耶洗別.
阿哈謝的兒子.
是約阿斯.
阿哈謝的.
太太就是阿他利亞.
就是耶洗別的女兒.
就帶來.
阿哈謝都幾乎.
被殺.
約阿斯本來是被殺的.
他的舅父.
耶和耶大.
把約阿斯.
收藏在聖殿.
八年.
阿他利亞.
這個女王.
猶大唯一一個女王就是.
阿他利亞耶洗別的女兒.
就是猶大的女王.
就像武則天.
到處都去找約阿斯找不到.
原來在聖殿裡面埋藏.
即是說.
他都不屑入聖殿.
第十二節.
約阿斯的兒子就是阿瑪謝.
阿瑪謝的兒子.
就是阿撒利亞或者叫烏西亞.
阿撒利亞的兒子.
就是約坦.
約坦就跟著阿哈斯.
希西加 嗎哪西.
希西加的兒子就是嗎哪西.
全個猶大.
十九位軍王.
最壞的就是嗎哪西.

$^{1881}$誰知最長做王的就是他.
五十五年.
你可想而知嗎哪西.
聖經在舊約曾經說過一句說話.
因為嗎哪西的罪孽.
猶太人是.
不能再回轉.
一定要面對受罰.
這個是嗎哪西.
不過嗎哪西的兒子就是.
阿們.
阿們之後就是約西亞.
最後一位最好的中興的王.
約西亞.
約西亞有一個復興.
約西亞的長子就是約哈蘭.
或者叫約哈斯.
次子就是約瓦敬.
三子就是西底加.
因為約瓦敬的兒子被擄了.
所以叔叔.
又做王了.
第四個兒子就是沙龍.
約瓦敬的兒子就是耶哥尼亞.
耶哥尼亞或者叫哥尼亞.
和西底加.
或者耶哥尼亞也可以叫約哈斯.
十七節.
留意這裡.
耶哥尼亞被擄.
就是約瓦敬的兒子被擄.
所以他的叔叔.
西底加就做王了.
耶哥尼亞被擄.
他的兒子是撒拉帖.
撒拉帖.
你剛才知道.
就是所羅巴巴的爸爸.
瑪基蘭.
比大瓦.

$^{1921}$是納薩.
尼大比瓦.
比大瓦的兒子是索羅巴巴.
是梅.
索羅巴巴的兒子是米舒蘭.
哈拿尼亞.
這個是他們的族弟兄.
等等的家譜.
所以哥尼亞原來也有後裔.
撒拉帖.
因為在麻太福音.
直接說撒拉帖的兒子.
就是索羅巴巴.
我們又要再揭多一點經文.
來看關於.
尤其是提到.
神特別點名提索羅巴巴.
說是他的僕人.
說是他的人.
還有.
當那日.
神的國度就會興起.
所以我們要揭少少.
關於這個僕人.
這個概念.
我們可以看看.
以西結書兩段.
以西結書的三十四章.
開始有大衛是神僕人的概念.
可以看短少少.
以西結書的三十四章.
由十七節這段開始提.
我的羊群呀.
輪到你們.
主耶和華如此說.
我必在羊和羊中間.
公綿羊與公山羊中間.
施行判斷.
你們這些肥壯的羊.
在美好的草場吃草.

$^{1961}$還以為小事嗎.
剩下的草.
你們竟用題踐踏.
你們喝清水.
剩下的水.
你們竟用題攪混了.
至於我的羊.
只得.
吃你們所踐踏的.
喝你們所攪混的.
即是說那些肥壯的.
可能是說那些領袖.
不顧惜群羊.
只懂得去踐踏.
肥幾.
第二十九節.
對不起二十節.
所以主耶和華如此說.
我必在肥羊瘦羊中間.
施行判斷.
因為你們用脅.
用堅擁擠一切瘦弱的.
又用割底足.
以致使你們四散.
所以我必請求你們.
四散.
所以我必請求我的群羊.
不再作掠物.
我也必在羊和羊中間.
施行判斷.
我必立一牧人.
照管牠們.
牧養牠們.
就是我的僕人大衛.
他必牧養牠們.
作牠們的牧人.
我耶和華必作牠們的神.
我的僕人大衛.
必在牠們中間作王.
這是被擄的年代.

$^{2001}$被擄的年代.
大衛已經死了很久.
幾百年了.
大衛是主前的一千一百至一千年.
多一點的時間.
一千年左右.
因為始終是出土的文獻.
不是那麼清楚.
所以.
你想想現在這個是被擄的時候.
是主前的六百多年.
接近五百多.
五百九十多.
所以.
如果以西結是最後那批.
或者可能是第一批.
那就是六零六.
最後那批就是五百八十多.
那你想像得到.
這個.
四百年前的人.
大衛.
在神的眼中心目中.
都是他的僕人.
照管國家的牧者.
又是神設立的.
一個中心的僕人.
這樣來重提.
過去的這些.
屬靈人.
去到三十七章又再有.
似乎是說.
大衛這一條線.
這一條線.
這個脈絡.
這個中派.
或者這個支派.
猶大支派.
一直都是被神揀選.
所以他們一路.

$^{2041}$承繼著這個職份.
成為那個牧者.
特別重要就是成為神的僕人.
三十七章.
這次在廿四節這一段.
我的僕人大衛.
必作他們的王.
眾民必歸一個牧人.
他們必順從我的典章.
謹守進行我的律例.
他們必住在.
我次級我僕人.
那國的地上.
就是你們列祖所住之地.
即是歸回了.
他們和他們的子孫.
並子孫的子孫.
都永遠住在那裡.
我的僕人大衛.
必作他們的王直到永遠.
並且我要與他們.
立平安的約.
作為永約.
我也要將他們安置在本地.
使到他們的人數增多.
又在他們中間.
設立我的聖所.
直到永遠.
我的居所必在他們中間.
我要作他們的神.
他們要作我的子民.
我的聖所在.
以色列人中間直到永遠.
愛邦人就必知道.
我是叫到以色列成為.
聖的耶和華.
這是很雙重的應許.
和安慰.
不過是將來的預言.
神的選民.

$^{2081}$有牧者.
有憤悔.
有神同在帶來的平安.
因為那個永約.
不再遷移.
不再修改.
是一個平安的約.
這就是新約.
耶穌與我們立的.
耶穌是真正的僕人.
真正的我是好牧人.
好牧人為羊.
捨命.
所以耶穌是完全對舊約的概念.
對於神預言的.
好牧人.
他完全好像對號入座.
主耶穌很清楚.
這個群體.
這個選民.
這個神的百姓.
需要的.
牧者就是他.
他就是由天而降的.
大牧人.
他就是那一位.
背負世人罪孽的.
神的羔羊.
又是憲制羔羊.
又是牧羊的.
牧者,又是祭司.
很多的身份.
都在主耶穌的身上.
又是君王.
這個僕人.
這個概念.
在被擄之前.
已經出現.
大衛就是那個代表.
是一個很完美的.

$^{2121}$一個典範.
大衛就是這個.
的後裔,永遠接續.
所以現在去到.
這個被擄歸回.
不是因為.
不是因為所羅巴巴.
因為剛剛的評分.
比特首還高,57分.
100分,所以現在.
給一個新的名詞.
就是我的僕人.
其實就不是.
其實就是一直以來.
大衛這個體系.
是孕育將來的彌賽亞.
所以猶大之派.
神都興起這個所羅巴巴.
在這個這麼關鍵的.
被擄歸回的時期.
他就是神的僕人.
所以.
這是一個很尊貴的身份.
也是.
受苦僕人.
我們揭一揭以賽亞書.
分外明顯是主耶穌.
41章.
以賽亞書.
40章之後.
就是一個很新的.
段落.
當然有.
近代的解經家就認為.
40章到66章.
是叫第二.
第二以賽亞社.
是原著的以賽亞社.
因為那個結構是完全不同.
他們認為是另外一個人.

$^{2161}$寫的,不過還是要用.
以賽亞.
近代已經.
差不多沒有什麼懸念.
解釋,就好像.
推翻了以往.
一個作者的.
那個.
理論,因為.
始終我們現在都找不到.
原著,手稿.
有的是抄本.
有的是.
馬索拉抄本.
都是抄的.
死海古卷有整卷的.
以賽亞書,都是抄本.
第二世紀.
的作品.
最完整就是以賽亞書.
就是在死海發現的.
那個古卷.
當然價值連城.
因為這麼重要的文獻.
可以推翻上去.
而以賽亞書的抄本.
和我們的馬索拉版本.
是接近沒有任何的錯誤.
你想想第二世紀.
出土的死海古卷.
和我們現在.
中文翻譯的.
舊約希伯來文.
是用馬索拉抄本.
的希伯來文.
直譯.
馬索拉抄本是.
十至十三世紀.
的抄本.
猶太人的拉比的抄本.

$^{2201}$推索十多個世紀.
去到死海的.
以賽亞書.
完全一樣的.
很輕微的.
抄寫上的.
瑕疵.
你可以想像到.
猶太的拉比抄聖經.
那個嚴謹.
所以我們接近完美.
又不是用印刷術.
是人手抄的.
都可以保存到這樣的質素.
就是神的話語的.
嚴重的重要.
你可以想像到.
在這個權威裡面.
死海古卷的歷史價值是這麼重要.
當然有很多已經被那些.
牧羊人撿去燒了.
所以就找不回了.
已經沒有了.
當時一個困難.
一個社團在耶穌時代.
某程度上是一個異端.
他們有很嚴謹的.
自己的守則.
在死海附近居住.
其實是不屑耶路撒冷.
已經很腐敗.
他們就要很原教旨.
住在困難的一帶.
來做一些苦修.
他們的工作.
就是抄寫聖經.
就不會和耶路撒冷的人來往.
我們看看.
以柴亞書的四十一章.
(四十一章).

$^{2241}$他們各人幫助淪陷.
各人對帝輕說壯膽吧.
不讀這麼長了.
第八節.
唯有你以色列我的僕人.
瓦割我所揀選的.
我朋友亞伯拉罕的侯裔.
這裡強調的僕人是以色列這個民族.
他們要歷經很多的痛苦.
接著我們看第四十二章.
四十二章的第一節.
一而再講這個僕人.
看啊我的僕人.
表面一直在講的是以色列這個民族.
是神的選民.
神的僕人.
看啊我的僕人.
我所扶持所揀選.
心裡面所喜悅的.
我以將我的靈賜給他.
他必將功利傳給外邦.
他不軒揚不揚升.
也不使在街上聽見他的聲音.
壓傷的蘆葦.
馬太芬九章就引用過.
他不截斷.
將殘的燈火.
他不催滅.
他憑真實將功利傳開.
他不灰心.
也不喪膽.
直到他在地上設立功利.
海島都等候他的憤愧.
似乎一直在講這個僕人.
他會有這些的工作.
以及帶來的公平公義.
甚至乎.
他能夠對一些軟弱的人.
有一個很大的.
那個牧養的情.

$^{2281}$那份恩典.
使得壓傷的蘆葦.
都不會截斷.
這已經成為很好聽的詩歌了.
似乎是在講以色列這個民族.
但是.
當然我們就一直知道.
要明白了.
猶太人.
他們因為是拒絕新約.
是不接受新約.
他們一直都定性這個神的僕人.
就是他們自己這個民族.
是以色列.
他們不能夠接受.
你們新約基督教講的僕人.
理解為耶穌基督.
因為.
是新約自己的猶太使徒.
就是彼得.
保羅他們.
他們再看回這些舊約的時候.
是他們自己認定了.
原來舊約所預言的這個忠心的僕人.
或者叫受苦的僕人.
不是在講以色列這個民族.
是耶穌.
所以.
這個就是.
釋經的問題了.
猶太人的釋經.
就是對舊約的釋經.
和新約.
和基督教的釋經.
是不同的.
所以大家要明白這一點了.
44章再有了.
其實你一直看著就是.
我的僕人.
就是這一大段聖經的主題.

$^{2321}$蜜諾來的.
第44章.
我的僕人雅各.
我所揀選的以色列.
現在你當聽.
造作你.
又從你出胎造就你.
並要幫助你的耶和華.
如此說.
我的僕人雅各.
我所揀選的耶穌倫.
不要害怕.
因為我要將水.
澆灌口渴的人.
將河澆灌乾旱之地.
我要將我的靈.
澆灌你的後裔.
將我的福.
澆灌你的子孫.
他們要發生在草中.
像溪水旁的柳樹.
如此等等.
詳細再不讀下去了.
一路說如何將民族改變.
本來要揭53章就不看了.
不如和大家揭一篇詩篇.
我預備了的78篇.
詩篇78篇.
這篇是很長的詩篇.
如果你想看以色列的歷史.
這篇詩是很好的.
因為作者很有興趣.
他把古代的歷史.
和他們現在的處境.
拉上了關係.
很長的.
我只看一段.
因為是七十多節的.
由65節開始.
65節看完這一段就完了.

$^{2361}$那時主將世人睡醒.
將勇士飲酒呼喊.
他就打退了他的敵人.
叫他們永蒙羞辱.
並且他棄掉弱勢的帳篷.
不揀選以法連之派.
卻揀選猶大之派.
他所喜愛的石安山.
記住.
猶大之派石安山.
代表著耶路撒冷.
蓋造神的聖所.
好仗高峰.
又仗他建立永存之地.
又揀選他的僕人大魏.
從羊圈中將他召來.
叫他不再跟從那些大賴的母羊.
為要牧養自己的百姓雅各.
和自己的產業以息列.
於是他按心中的純正牧養他們.
用手中的巧妙引導他們.
這裡再一次講大魏.
猶大之派.
和神喜愛的石安山.
這樣做一個聯繫.
即是說.
神的救恩計劃.
與猶大大魏石安山.
有關聯.
所以那個僕人.
一直講到呼之欲出.
就是受苦的僕人.
其實就是神.
在末世為人預備的彌賽亞.
或者為猶太人預備的彌賽亞.
當然主耶穌出來.
他就說自己就是這位的僕人.
這個就是索羅巴伯.
所代表著的將來的神的僕人.
就是耶穌基督.

$^{2401}$所以這個預言.
對索羅巴伯來說是一個高台.
因為他現在作為全民的.
輕工建造聖殿的主領人.
他有很重的責任.
當然也有很重的擔心.
因為要周旋在那些政治.
當時的波斯的君王.
和那些地頭蟲.
那些地方的勢力.
要擺平.
要建造聖殿是不容易的.
不過神就說給索羅巴伯知道.
我的僕人.
不需要懼怕.
我以你為人.
必定與你同在.
將來到那日.
萬國的珍寶和那些祝福都會來到.
相信是說的是耶穌基督將來的實現.
所以有他的預言應驗的相關性.
有局部的短暫的.
很簡單的應驗.
不是已經完全應驗的.
這個就給我們有少少印象.
對於索羅巴伯處於重建聖殿的時代.
一個重要的標誌性人物.
我們一起禱告.
多謝主給我們今晚.
能夠梳理好幾個很重要的.
密時的預言.
有哈該在六月初一.
七月二十一.
九月二十四.
也有撒迦利亞.
在八月.
大利烏王的第二年的八月.
一共有五個的密時的預言.
來安慰.
和要他們有一個回想.

$^{2441}$那個日子.
使他們能夠上下一心.
真的將很多的生命的次序排列.
以神為優先來服事.
他們過了那一代人之後.
聖殿也都是重建完畢.
很盼望我們在今天.
當我們在教會來服事.
一起來為神的福音.
我們去傳遞的時候.
都幫助我們能夠在.
我們的家庭工作.
我們的個人興趣.
等等方面.
我們都能夠有.
上主給我們的.
心靈的感動和考慮.
我們能夠在.
不同的範疇.
真的在價值觀上面.
因為一切問題的根源.
都應該在每個人的價值觀.
怎樣去看事物.
怎樣能夠看到.
一般人看不到的.
屬靈的眼光.
求助幫助我們能夠.
真的有.
以至我們的人生的投資.
不是在一些真的會扭壞的事物.
不是那些事情不重要.
而是那些事情是會過去.
唯有遵行神的旨意.
才能夠永遠常存.
求主幫助和祝福.
我們這樣感恩禱告.
奉耶穌基督的名.
阿們.
\newpage



\section{馬可福音}
\label{sec:2Us7kVIsCO8}
\textbf{2023-02-05 加恩樓《是賊窩?還是神的殿!!》馬可福音 11: 12-25 講員:鄧泰康傳道}
\newline
\newline
連結: \href{https://youtube.com/watch?v=2Us7kVIsCO8}{\texttt{https://youtube.com/watch?v=2Us7kVIsCO8}} ~~~~ 語音日期: 2023-02-07
\newline
\newline
\hyperref[sec:i9iPPyuXEWw]{\small{< < < PREV SERMON < < <}}
~
\hyperref[sec:index_chronic]{\small{[返順時目]}}
~
\hyperref[sec:index_scriptual]{\small{[返順卷目]}}
~
\hyperref[sec:k6XACT4Tfao]{\small{> > > NEXT SERMON > > >}}
\newline
\newline
$^{1}$今天跟大家講這段經文的時候,我還沒講之前,我要先做一些聲明.聲明什麼呢?.
因為我發覺我們做基督徒做得久了,有時會玻璃心..
什麼叫玻璃心呢?就是我一講的時候就像我在說你一樣,.
然後裡面有很多反應,就是「不是啊,我不是這樣的,你誤會了我.」.
為什麼我這麼肯定呢?因為我聽過不少這樣的回應..
我想強調,鄧全道是絕對不會用講台變炮台的,OK?.
並不是我見到任何現象,我藉著講道去責備任何人,我從來都不會做這件事..
我只是按著聖經裡面,聖經怎麼教我們,我們就怎麼說..
聖經的信息裡面怎麼說,如果你心裡面覺得責心,.
甚至你會覺得那些經文好像在說自己,請你留意,不是鄧全道見到任何東西,.
然後我在說你,而是聖靈在提醒你..
好不好?這個很重要,不然我們聽再多的道都沒用的,因為你要不斷反彈..
今天跟大家講這些,為什麼我要這麼閃聲明這件事呢?.
因為今天這篇道是重的,為什麼呢?你會發覺題目是「究竟是賊魔還是神的電?」.
是很重要的,因為這段經文其實是直接指向於教會和基督徒生命..
一間教會最讓人擔心的是什麼呢?是虛有其表..
基督徒最讓人擔心的是虛有其表的生命..
所以你就會發覺這段經文其實是很工整,很容易看的..
為什麼呢?因為他用了一個三文字的結構..
你會發覺先是耶穌見到這棵無花果樹,然後發生了一些事情,.
就在它之後,夾在中間的聖殿,之後再講無花果..
這是一個很典型聖經的格式,是在講一個其實重點是在講什麼呢?.
就是中間的聖殿,而透過無花果去講聖殿..
所以核心一定是聖殿所發生的事情..
而這段經文很直接指向於信仰上的虛假..
而在哥林多前書3:16,他說:.
「豈不知你們是神的殿嗎?神的靈就住在你們裡嗎?」.
所以這段經文,我今天也會同樣地去反省,.
今天除了教會之外,我們自己的生命是如何..
因為你要記住,你信了耶穌之後,聖靈住在你的裡面,.
你是上帝的殿..
什麼叫做虛假?從這段經文裡面,不說其他地方,什麼叫做虛假?.
就是「人餓了,找不到吃的.」.
耶穌當時很清楚,我餓了,那我就要吃東西,.
很簡單,大家今天餓了,那去哪?.
去學校?你不是的,你去找餐廳,對嗎?.
你去餐廳當然是想找到你想吃的東西,.
是帶著一個很重要的期望的..
但是耶穌現在發生什麼事呢?這裡一定有一點背景,.
如果你不明白這段背景和講聖殿的話,你解不通這段話..

$^{41}$從來我不知道大家有沒有發覺,這段經文裡面的耶穌突然間變了..
第一樣被人很誤會的就是,他餓了,看到一棵無花果樹,.
沒有果子,還有聖經的作者已經跟你解釋了,.
那個時候不是收無花果的時候,合理嗎?.
其實很合理的,不是收的時候在樹上找不到果子,.
這樣很合理的,但是耶穌就了那棵樹,就了那棵樹死了,.
是不是很不正常?.
這個絕對不是耶穌,然後到聖經,我想問,有沒有見過耶穌這麼暴躁?.
你記住耶穌是直接走去推翻那些桌子,.
你從來沒有見過這樣的耶穌,.
所以你就會發覺這裡整段經文是帶著耶穌跟你說一些很重要的東西,.
而促使這件事發生.為什麼呢?.
原來無花果樹的葉子是代表著有果子..
什麼意思呢?在一個不是收果子的時候,.
所有的無花果樹上都還沒有葉子,但是突然間有一棵上面有葉子,.
然後耶穌就要去找了,.
就是說,如果不是收無花果,就是根本沒有果子給人吃,.
就拜託你不要生出葉子,你明白我的意思嗎?.
就是說這棵樹不知道什麼原因,發生了一個很特殊的情況,.
竟然在一個不是收果子的時候,不應該有葉子的時候,.
他出了葉子,然後吸引了一些餓的人上去找食物,.
而我去到我竟然找不到,明白嗎?.
我想告訴大家,其實這個故事是有一點為了講而講的,.
是的,因為他要講聖殿,.
當你一套進去聖殿之後,整件事就通了,.
每一個尋找上帝的人,做什麼要去聖殿?.
就是要找到聖殿,聖經裡面講得很清楚,.
因為我生命裡面的饑渴,我要去到上面找到上帝,.
而現在這些人,他們是帶著期望來到聖殿當中,.
我要找到上帝,但是現在發生的事情就是,.
整個表面上,在這個聖殿裡面的外表,.
在所有人的衣著,做的任何一件事上,.
都好像告訴你,你能夠找到上帝,.
但是那些人去到裡面的時候,卻是找不到..
其實裡面,大家要明白背景,.
我想大家有些私心信徒,你聽過很多次了,.
兌換銀錢和賣甲者這件事,其實是合法合情合理的,.
為什麼?第一件事,就是當時因為他們在用希臘的貨幣,.
羅馬的貨幣,上面有一個皇帝的頭像,.
他們需要去到當中,他們認為這些是有偶像的東西,.

$^{81}$我們不可以奉獻給神,然後那裡就有一個專門的地方,.
就給人兌換銀錢,兌換他們的貨幣,.
之後才可以將這些奉獻給上帝,合理嗎?.
第二件事,為什麼要賣甲子?.
因為那些人是千里迢迢來到這個地方,.
你想像一下,你帶著生畜,你記住不是今天,.
不是有車,不是有飛機,你帶著牠去山長水遠的時候,.
美的都變死了,然後我就專門在那裡有個地方,.
給你買一些最美的甲子,最簡單,最直接可以獻祭的動物,.
合理嗎?對信仰是非常尊重的做法,.
但是有些學者研究之後告訴大家,.
有一個現象很有趣,原來猶太人相信聖殿是最接近上帝的地方,.
他們很相信在聖殿裡面,我做的任何事情,.
上帝聽的禱告是最聽得清楚的,.
所以我不知道大家有沒有發覺,.
以色列人教他們一定有一個叫朝聖,.
他們一定要去墨家,一生人裡面一定要有一次去墨家,.
他們就要去到最近他們那個神的地方,其實是同一個觀念,.
所以他們當時的人無論住在哪裡都好,.
他們都要回到耶路撒冷去做這個獻祭,.
好了,裡面出現問題了,聖殿這麼大,哪個位置最接近上帝?.
大家能回答我嗎?.
大家知道聖殿的格式,有外邦人院,有聖所,有致聖所,.
致聖所不是每個人都可以進去的,你給錢都沒用,都不能進去,.
只可以大小師一年進去一次,.
但是聖殿哪個位置最接近上帝?.
當然是致聖所門口,.
哪個人可以走到那個位置?.
由誰決定?.
由祭司決定,.
肥沙,你有什麼情況下,我作為祭司,.
我邀請你來到這個最近的前面,有什麼方法?.
給錢吧,明白嗎?.
其實裡面有些不法的事情正在發生,.
你想近上帝?給錢吧,.
你講得直白一點,你賄賂我,.
你就能夠去最近上帝的位置,.
又或者兌換銀錢的時候,你去到怎樣都要換的,.
那怎麼辦?.
對了,排價,有高低之分,.

$^{121}$貴一點,合謀定價,.
然後等你的錢換得不合理,.
為什麼耶穌這麼生氣?.
因為你用一些表面神聖的方法去獲取一些不合理的收入,.
甚至乎,講得直白一點,.
你用神聖的方法,用上帝的名義去達到人的目的..
弟兄姊妹,今天在教會當中,.
每一個人來到教會,我很肯定,.
特別是最初,我一定是想找到上帝,.
首先你要問,這個餓的人,他來到當中,能否找到上帝?.
上帝是誰?上帝是什麼?.
他來到教會裡面的時候,他看到的是什麼?.
是人與人之間的那種從上帝裡面出來的愛?.
是我們裡面上帝所給大家的那種盼望,那種唇亡愛?.
還是只是看到人與人之間裡面很多的爭鬧,.
心裡面很多的挑戰挑動,.
人與人之間很多的不和,.
甚至乎,好像這裡所講的一樣,.
去偷取,為什麼叫做拆窩?.
因為你偷取了上帝的榮耀..
這裡,我們會不會用上帝之名去達到我們自己的目的?.
作為傳道人,其實這個位置是非常危險的..
有機會的時候,我時時都會在上面跟大家講道..
很多時候,我也很感恩,.
有些人說,鄧炫濤,我很喜歡聽你講道,我也很感恩,.
但問題是,如果你只喜歡聽鄧炫濤講道,就大件事了..
因為我只是將你帶到我面前,.
我只是讓大家成為我的門徒,.
我不是讓大家成為上帝的門徒,.
我不是讓大家成為找到上帝的流通管子..
弟兄姊妹,我們今天在教會裡面所做的任何事情,.
為的目的是什麼?.
坦白講,為什麼我們這些弟兄姊妹,.
這些同工,我們會這麼拼命?.
原因就是,我們在這裡不斷思索一個很重要的問題,.
工具,途徑不重要,.
重要的是那個工具是否能夠帶領人來到上帝的面前..
舉個例子,沙發店,大家還記得嗎?.
沙發店出現的,大家知不知道?.
後面的吧檯店出現的,大家知不知道?.

$^{161}$又講回一點歷史,一會兒富哥也會講一點..
世界盃,我想比較少教會敢用禮堂來看世界盃..
白天還在崇拜,晚上大家還在吃著薯片,看著世界盃,一起"Go!".
我們去做啊,但我想告訴大家,.
就是因為這樣的氣氛,吸引了一群未信的街坊弟兄姊妹來到裡面..
曾經有二百五十人以上..
而最奇妙的地方,我想告訴大家,.
有些未信耶穌的人,他真的需要拿著一杯咖啡才能夠參加聚會..
所以你會發覺後面的吧檯店出現了..
有一個時間,疫情之前,我很記得,.
有一些未信主的朋友,街坊,.
他們就是拿著一杯咖啡,坐在後面的吧檯店,全部坐在最後,.
去聽我們講道,和我們一起唱詩歌..
弟兄姊妹,這些沙發,這些吧檯,這些咖啡,將人帶到神的面前,.
是不是很值得感恩?.
所以,在這裡的時候,上帝要的是什麼?.
我們要很清楚這個地方最重要的作用是什麼..
但是,如果我們為了玩而玩,.
就是說,我們這裡製造了很多人氣,.
大家來到的時候,大家都很開心,.
不過,沒有了上帝..
Sorry,不行..
這裡必須要找到上帝..
所以你會發覺,上帝在這裡的時候,他要什麼?.
就是在後面的時候,當他提醒門徒,他講了一句話,.
他要的不是你們這些兌換錢的心意,.
或者買了什麼..
偷取上帝榮耀的意思就是,人以為我兌換了錢,.
買了這些甲子,就等同我愛上帝..
我在教會做了一些事,發生了一些事,.
就等於我愛上帝.不是的..
在二十五節的時候,他說,.
你們站著禱告的時候,若想起有人得罪你們,.
就要饒恕他,.
好叫你們在天上的府也,饒恕你們..
他要的是你們行出上帝的愛,.
上帝有多愛你,你就要怎樣愛人,.
而不是只是回來教會參加聚會..
這個才是這個殿裡面真正的需要..
所以你會發覺這段經文,講到後面的時候,.

$^{201}$有一個桃花果樹的結果,你就會發覺很特別..
你看到彼得對耶穌說,.
耶穌啊,你看到嗎?你就在那棵樹上枯萎了..
其實這段話很難解釋,因為沒有任何情緒..
但耶穌的回應就告訴你,.
其實彼得真正去讚譽的是什麼呢?.
是你的能力,耶穌你看你多厲害..
你就在那棵樹上枯萎了,.
他很想得到這種能力,.
說風就風,說雨就雨,.
其實彼得出現了幾次,.
平靜風浪也是這樣,很多次,他很羨慕這種能力..
但耶穌很有趣,他的回應是什麼呢?.
為什麼可以這麼肯定呢?.
他就告訴他們,其實你們所看的,你們很看沒能力,.
小兒科而已..
所以今天大家的時間,一看這段經文,.
不是叫我們基督徒移山就移山,.
大家一起,是不是完了崇拜之後就出去,.
這間屋變成那間屋,不是的,他不是說這件事..
他用了一個很誇張的比喻告訴你,.
如果你對上帝有一個真正的信服的話,.
你對有真正信心的話,.
這些很簡單,.
什麼叫信服?.
第一樣東西,是對上帝有100\%完全的信心..
你若心裡不疑惑,.
只信他所說的必成就,.
就必給他成就..
所以我告訴你們,凡你們禱告祈求的,.
無論是什麼,只要信,就必得著..
我們今天是因信稱義,.
我們是信這位上帝,弟兄姊妹,.
但是你對上帝的信心有多少?.
我們今天是不是真的完完全全相信上帝,.
無論在任何情況之下,你都掌權,.
任何情況之下你都看著,你都掌管..
如果你真的對上帝有這份信心,.
我可以夠膽告訴大家,.
無論發生任何事,你都會很平安,.

$^{241}$你都會很平靜,.
你會處變不驚,你會處境不變,.
因為你有上帝,你完全相信上帝..
耶穌在說的是這種信心,.
你相信上帝會移山助山就走,.
上帝會用祂的方法,而不是你的方法..
我們今天能夠得到這個救恩,.
都是靠這份信心,弟兄姊妹..
在座的有沒有人見過耶穌?.
沒有,包括我,我都沒有見過..
為什麼今天我們可以這麼肯定?.
因為我相信,.
我經常和未信的朋友說到最後,.
我和你的唯一最大的分歧是什麼?.
我相信聖經和你不相信聖經..
我相信聖經裡面所說的一切,.
你不相信聖經所說的一切,.
就是這麼簡單,.
就正因為我相信,.
所以我們的生命可以不一樣..
所以為什麼昨晚我們團契的時候,.
我們說我們的生命需要調教,.
是因為我們的生命,.
就正因為有上帝的緣故,.
所有的事情都會不同..
你今天做事的優先次序,.
你今天做人的目標,.
你今天所擁抱的東西,.
你今天所相信的東西,.
完全會完全不同..
這個才是你真正生命裡面的能力,.
和你可以感動人的地方..
往往我們最能夠幫助人的時候,.
是我們面對困難的時候,.
有沒有發覺?.
人們會很驚訝的,.
為什麼我們基督徒不同?.
為什麼我們基督徒會有這樣的信心勇氣?.
你就能夠去燃點身邊的人,.
因為他們在這個世界上找不到..

$^{281}$儘管他很有才能,.
一個殖命他就會垮掉..
儘管他很有錢,.
一個風暴他就可以什麼都沒有了..
但是我們卻不是看這些,.
我們看的是上帝,.
所以一切都可以不一樣..
第二樣東西什麼叫信服?.
是聽命..
所以剛才我提過,.
上帝要大家的,.
是神怎樣愛你,.
你就怎樣愛人..
神怎樣為你犧牲,.
你就怎樣為其他人去犧牲..
又一樣,當你的價值觀改變,.
你自然就會改變..
我們在奸瘤裡面經常都有一句話,.
就是苦瓜口頭禪,.
按時分量..
奸瘤是一個這樣的地方,.
來到的人有沒有所需要的東西,.
他能夠有得吃?.
我們經常都不斷思考這個問題..
所以很簡單的,.
邏輯很簡單的,.
為什麼會有咖啡機?.
因為要上課,.
上課之後會怎樣?.
會睏,.
睏需要咖啡,.
明白我的邏輯嗎?.
就是這麼簡單..
大家在這裡的時間只叫你聊天,.
你不會聊,.
為什麼我們要有茶點?.
因為吃著東西來聊天..
每一個在時間裡面都有需要的東西,.
我們就在適當的時間,.
我們就需要讓他們能夠找到..

$^{321}$今天遇見你的人,.
能不能夠得到飽足呢?.
突然之間不會..
來到奸瘤的人,.
能不能夠得到飽足呢?.
究竟教會和我們的生命裡面,.
會不會只是一些宗教上的兌換銀錢,.
和賣夾子的操作,.
只是讓人覺得,.
我已經做了什麼就很安心,.
甚至乎好像祭司和法利賽人那樣,.
去為著自己的利益而去謀損..
第三,我們對上帝的信心如何?.
你是不是真的一點都不疑惑?.
我們有沒有去活出屬於上帝的生命?.
所以弟兄姊妹,我們很盼望,好嗎?.
初心,.
奸瘤最初在這裡,為什麼會開始崇拜?.
是因為有一群弟兄姊妹,.
見到這個地方,.
生命是需要崇拜..
所以很多時候,今天我們一起和敬拜隊,.
我想敬拜隊也很認同一樣東西,.
我們很多時候,我們祈禱也說,.
我們今天不是表演,.
我們今天是上帝透過我們的生命去祝福其他人..
我們的生命是如何影響了我們今天敬拜的整個表現?.
我們今天生命如果有能力,有火,.
你自然帶出去敬拜,就有生命力..
這個才是上帝使用我們,.
所以求神幫助,帶領我們,好嗎?.
讓奸瘤成為真正上帝的聖殿,.
不要墮入去成為賊窩這個危險,.
求神保守我們,我們一起動心祈禱..
因為我們知道上帝你賜教會,.
讓我們一起當中有份,.
你讓我們成為教會,.
你讓我們周派望向這個地方,.
每一個人都能夠找到你,.
能夠帶領更多人去相信你,.

$^{361}$也因為我們的生命,我們去看見這份亮光,.
看見這份力量,.
以致能夠用我們的生命去感染其他人..
你不斷地拖帶我們走過很多不同的路,.
就像剛才所唱的捉人一樣,.
主啊,有很多事我們都不明白,.
但原來你告訴我們,當我們最困難的時候,.
你抱起我們,你帶起我們,.
我們一起向前走..
主啊,我們既然經歷了這份這麼大的恩典,.
我們求主讓我們在這裡同樣可以.
跟我們身邊的人分享這份恩典,.
讓更多人得到這份祝福..
我們求主你親自去保守,.
求主你繼續幫助,.
求主你使用我們,.
其中奉求主耶穌,得勝名至其穆..
阿們..
\newpage



\section{馬可福音}
\label{sec:k6XACT4Tfao}
\textbf{2023-02-05 午堂《有果效的門徒生活》馬可福音 11: 12-26 講員:曾裔貴牧師}
\newline
\newline
連結: \href{https://youtube.com/watch?v=k6XACT4Tfao}{\texttt{https://youtube.com/watch?v=k6XACT4Tfao}} ~~~~ 語音日期: 2023-02-07
\newline
\newline
\hyperref[sec:2Us7kVIsCO8]{\small{< < < PREV SERMON < < <}}
~
\hyperref[sec:index_chronic]{\small{[返順時目]}}
~
\hyperref[sec:index_scriptual]{\small{[返順卷目]}}
~
\hyperref[sec:lrudeppALpc]{\small{> > > NEXT SERMON > > >}}
\newline
\newline
$^{1}$弟兄姊妹主裡平安.
和大家一起祈禱.
多謝主再一次恭謹的等候在您的面前.
我們眾人包括講的.
我們每一位弟兄姊妹的心靈.
都來到去靜心聆聽您的聖言.
我們能夠有信心.
真的好像主耶穌所講的.
要有兒孫的信心.
這樣來到去等候.
我們的禱告奉耶穌基督的名.
阿們.
這段聖經發生的場景時間.
是主耶穌在最後的壽蘭州.
祂騎驢入城.
這是第一到第十節的經文告訴我們.
作者很清楚聚焦了在聖殿裡面發生的.
當然是兩天左右的事件.
然後中間夾雜的.
就是祂每晚都離開聖殿耶路撒冷.
去到伯大利.
伯大利是離開耶路撒冷六公里路.
所以不是很久.
祂就在伯大利.
似乎就是在拉薩路.
馬大.
馬尼亞的家裡居住.
每天早上就入城.
來到聖殿教訓百姓.
如是者.
壽蘭州七天.
最後在逾越節前.
就被捉拿.
然後當然是審問.
然後就被掛十字架.
作者馬可.
將幾段本來分開的經文.
聚在一起.
就帶出有多震撼的訊息.
值得我們今天去想一想.

$^{41}$我們的教會生活.
是不是能夠有果效.
結出果子.
好像主耶穌所說.
主耶穌在廿四節提到.
我們求什麼.
一定可以得著.
廿三節更加誇張.
可以移山的信心.
如果你有.
你就可以移山了.
這些經文的背景.
我們要梳理一下.
變成我們知道.
似乎不可以按字面解釋.
我們再一次看回第十一節.
就有更準確的背景.
因為作者在這裡特記載.
其他三卷福音書都沒有提及這件事.
就是當主耶穌騎驢仔入城的時候.
是以君王的身份.
因為當時有很多猶太人.
尤其是一些很想推翻羅馬政權的.
一些粉瑞黨的人.
或者一些祭司.
他們一直都很耿耿於懷.
自己的民族要在羅馬帝國底下.
成為了神文.
所以有分封的王希律.
和有巡撫彼拉多.
在這裡治理.
只不過是一個行政區.
猶太人自己.
所以是亡國的.
所以有一個這麼重要的.
稱自己為彌賽亞的耶穌出現.
他們就會將他們擁戴的王.
投放在耶穌的身上.
壓住在耶穌的身上.
所以他們認為.

$^{81}$逾越節應該會有大事發生.
所以在這個禮拜裡面.
是相當轟動當時的節日.
在這個時候.
主耶穌在第十一節告訴我們.
靜靜地進入聖殿.
周圍看了各樣的物件.
天色已晚.
就和十二個門徒出城.
往伯大利去了.
尤其是作者強調.
他自己和十二個使徒.
去到聖殿.
不是聚會的期間.
因為已經是黃昏了.
他看了周圍.
似乎摸了一下一些物件.
然後.
當然他沒有進入聖殿裡面.
他不是祭司他不可以進入.
他就離開.
和門徒出城去伯大利.
也似乎像約翰福所說.
他在拉薩盧的家裡作寄居.
在這個時候.
聖經就將第四節.
去到第十四節.
告訴我們有一個很特殊的神蹟.
神蹟在你的內心.
在你的印象.
一定是做出來幫人.
無論是醫治什麼疾病.
或者要帶出耶穌的權柄.
叫聽的人都能夠信服.
但是這個神蹟.
唯一在聖經裡面.
福音書裡面唯一一個神蹟.
你覺得匪夷所思的.
因為耶穌走過.
進城的路.

$^{121}$遠遠看到一棵.
當然是野生的.
無花果樹.
按他們的習慣.
如果有果實.
又可以拿來作為一點點的充飢.
無可厚非的.
因為是野生的.
當然是能吃的野生.
不是有人種的果園.
遠遠看到有葉.
似乎很茂盛.
然後走到樹下.
發現找不到果子.
很正常.
可能有一點點季節性的問題.
也可能在這棵樹.
它有一點點毛病.
有葉.
但是還沒出得切果實.
也很正常.
不正常的地方就出現在這裡.
主耶穌在這裡說的十四節.
耶穌對樹說.
哇.
各位在家中感受應該也有果樹.
你不會因為吃不到果子.
就說這麼負氣的話.
而且這個負氣的話.
彼得對耶穌說.
你詛咒的無花果樹.
真的枯乾了.
彼得是定性這句話.
十四節是詛咒.
從今以後.
永沒有人再吃你的果子.
他的門徒都聽到.
不是心裡面說.
作者很想凸顯這個神蹟.
往後門徒的回應.

$^{161}$很想讀者能夠清晰的知道.
耶穌不是心裡面說.
自己自言自語.
小聲地.
吟吟沉沉.
有沒有搞錯.
沒有果子的.
是指著他這麼說.
與樹何干呢.
強樹所難呢.
有不信的人就在這裡做文章.
哦.
Hungry.
然後就Angry.
然後就說一些Ugly的話.
有人真的不信.
這樣的耶穌說自己是彌賽亞.
咒詛那棵樹.
那棵樹沒有生命不會回應你的.
但是咒死了它.
下面告訴我們.
二十一節.
彼得想起耶穌的話來.
就對他說.
拉比.
請你看看.
這棵樹你所咒詛的無花果樹.
真是這一棵.
已經枯乾了.
各位弟兄姊妹.
為什麼一個這樣的神蹟.
於人毫無幫助.
於今天可能你第一次聽這個經文.
你都會覺得.
哇 耶穌.
你欠我們一個解釋.
一個合理的解釋.
到底你是不是真的Angry.
所以你來到這裡說出這麼重.
這麼務實的話.

$^{201}$一位彌賽亞.
這麼偉大的智者.
為什麼你會說出這樣的話.
你咒那棵樹沒有果子.
沒有果子關你什麼事.
我是有毛病.
你應該幫助我.
而且我是野生的.
沒有人種的.
很多這樣的想法.
但是作者將這些經文拼湊在一起.
似乎有屬靈的含義.
因為主耶穌.
接著在十五節到十九節.
他就好像都是相類的.
對聖殿已經失去了正常敬拜神的功能.
來提出譴責的時候.
這次就不是一棵樹沒有回應.
這次是有人.
而且是有相當影響力的宗教領袖.
對耶穌的回應.
十五到第十九節.
他們來到耶路撒冷.
說完這番話就上路.
接著就去到聖殿.
接著主耶穌就進入聖殿.
趕出殿裡作買賣的人.
似乎主耶穌在之前那一晚.
他已經看清楚在聖殿裡面.
有什麼擺設,運作等等.
有可能也都諮詢過不同的人.
當然更加主耶穌每年都會上耶路撒冷.
因為約翰福音的結構.
每年主耶穌都會上耶路撒冷守禦月節.
這次是最後一年.
第三年由他出來傳道.
第三年他上耶路撒冷.
這次他作出一個很出乎意料人想像的舉措.
就是趕殿裡作買賣的人走.
推倒兌換銀錢的人的桌子.

$^{241}$和賣甲子之人的椅子.
只是將器具推倒.
他有一個權柄去做一件事.
要讓人看到他這樣做.
是要帶出一個很重要的訊息.
這個訊息就是聖殿已經失去它的功能.
要恢復真正的功能是什麼呢?.
這就是關於祈禱.
所以當彼得來叫耶穌留意無花果樹枯死之後.
耶穌在22-26節所說的這段聖經最重要的中心.
就令我們更加吃驚.
我們先說回這裡.
主耶穌這樣做.
他要引述兩段聖經.
分別是《倚塞亞書》的56章.
和《耶利米書》的第7章.
兩段聖經的上下文.
都是譴責當時聖殿的失衡.
第一段聖經.
耶穌引述的.
我的殿真正的功能作用.
是萬國禱告的殿.
你們使得這個殿成為賊窩.
兩段都是和經文的背景是有相同的.
主耶穌好像一錘定音.
現在你們在做的事.
我們說說背景.
可能新來的弟兄姊妹未必太熟當時的背景.
很多外鄉的人.
住在不同的地方.
小亞,細亞甚至遠到歐洲的地方.
他們要回到耶路撒冷守節.
一年有三個節期.
一定要回到耶路撒冷.
逾越節,五旬節和主鵬節.
一定要上去守節.
現在在逾越節的時候.
他們上去守節.
要獻祭.
當然要帶牛羊.

$^{281}$在外邦住.
回到耶路撒冷.
要奉獻.
要兌換希伯來的舍夏勒當地的貨幣.
才可以奉獻到神的錢箱.
這樣的時候.
很多當時在聖殿裡的人.
他們因利成便.
有些便民的政策.
就是設置一些攤檔.
可以買賣牛羊.
現成的去到耶路撒冷.
再去購買.
就可以立即拖到祭司獻祭.
也可以有兌換的錢的地方.
這樣是不是對上來朝聖的人的一個方便呢.
為什麼耶穌說.
這個成為了賊窩呢.
我們就要了解.
原來當時在耶穌的年代.
這些所謂便民的政策.
已經成為用高價去買賣.
去投得那些攤檔.
而那些賣的人.
背後有一個很大的勢力.
操縱著那些攤檔的買賣.
導致出現了這樣的後果.
很多人上來獻祭.
敬拜神,朝見神.
你想想從這麼遠的地方上來.
發現的是聖殿裡面.
當然是外面廣場的地方.
就出現了這樣的後果.
等於我們很快的口罩令除開之後.
我們在廣場就會恢復一個朝聖的茶室.
你想想.
如果有人向我買.
我要在這裡投.
有個攤檔.
可以賣東西.

$^{321}$你給我.
然後你給我黑錢.
我給你.
類似這樣的情況.
在耶穌的年代.
就是這樣發生的.
耶穌看到整個的敬拜.
不能夠帶來百姓一心一意.
單純的去到神的面前.
而又有很多這些事情.
來影響了人的心的敬拜.
所以主耶穌在這裡講的時候.
好像做一個推倒從來的舉措.
留意第十八節.
有回應.
不單止沒有去反思耶穌這個撥亂反正的一個責備.
一個行動.
到底祂有沒有這樣的權柄.
來讓聖殿恢復單純的敬拜.
十八節.
如果相對剛才我們講的那棵樹不懂得回應.
就枯乾了.
而現在在這裡有回應.
就著這兩個推倒的行動.
似乎應該不是笑笑口地去做.
不是嬉笑怒罵地去玩.
很認真很嚴肅的氣氛.
耶穌是推倒責備這些人.
是趕出那些買賣的人.
耶穌展現的是一份嚴肅的聖潔的權威.
招來的回應就是.
不是那些持份者對耶穌講閒言閒語.
留意十八節.
祭司長和民事聽見者說.
就想發旨要除滅耶穌.
卻又怕祂.
因為眾人都欺其他的教訓.
似乎將兩個神跡.
當然趕出聖殿的人不是神跡.
似乎樹的神跡.

$^{361}$是反映著這個聖殿.
不能夠結出神期待的果子.
而來到去有引發了.
這班真正背後勢力的人.
對耶穌的敵視.
敵視到一個地步.
是要除滅耶穌.
你想想.
耶穌用那個神跡.
表面上不是對人有任何的幫助.
就一棵樹.
表面上好像耶穌由憤怒變成咒詛.
實際上.
原來不是.
固然只有十二個使徒.
一直跟著耶穌見到.
群眾是見不到這個神跡的.
耶穌好像展現給他們看.
主耶穌擁有的權柄.
和祂對這個不能夠結果.
表面有翠綠的葉的這棵樹.
我們不去研究到底是什麼原因.
是水土.
是祂的營養.
是季節.
我們不去研究.
聖經不是想我們這樣去研究.
聖經想我們.
作者想我們.
將幾段經文拼湊在一起.
然後.
如何可以恢復聖殿的功用勁大呢?.
相信就是第23節.
主耶穌所說.
我們需要有一份移開大山的信心.
來改變很多很多事情.
教會的功能.
是應該移開很多的大山.
因為.
耶穌基督知道有人這樣來對祂.

$^{401}$祂就不理了.
因為祂知道.
祂是要在耶路撒冷殉道的.
祂要死在耶路撒冷.
所以這一班的祭司長,文士.
是促成耶穌為神的真理.
將人帶到神面前.
唯一的道路打開的那班人.
主耶穌甘心樂意為這個很重要的功能.
來付上祂的生命代價.
亦都是我們每一個基督徒.
我們要具備成熟的生命.
而甚麼山呢?.
甚麼才是阻礙我們向神的山呢?.
這就是我們要一起去尋求.
再下去.
第20節.
就是剛才我們已經點出了.
所以作者真的很巧妙地.
將那個奏坐樹.
和往後神跡的奇妙.
拉開了來看.
中間夾雜了聖殿的譴責.
這樣就有一個很特別的意思在此出現了.
剛才已經讀了.
「早辰,他們從那裡經過,看見無花果樹,連根都枯乾.」.
本來正常就不應該見怪不怪的.
不過真的很怪.
彼得就煞有介事地想起.
不是的.
不是普通的營養不良.
不是普通的不夠水分,季節性.
是耶穌就過這棵樹.
不知道有沒有點下這棵.
不知道當時有多少棵.
在沿途的路上.
這一棵彼得記得.
就告訴耶穌.
拉比.
你看看.

$^{441}$你就過那棵樹.
真的枯乾了.
這是道德問題.
那棵樹.
某程度上是沒有罪的.
你就它.
它死了.
它受你的咒.
這是一個道德問題要解釋的.
是不是你真的生氣到一個地步.
沒有東西吃.
你就要咒它.
不可以的.
你去排隊.
一直都排不到位.
你就罵餐廳走了.
撕掉那張門票.
這樣也不好.
大家都知道不好.
我們不會咒到人要死.
要倒閉.
但是耶穌咒這棵樹.
咒到它死.
還要是很快.
第二天就發生.
我們已經感覺到.
好像有一個.
耶穌要為他自己的身份.
解釋的原因.
但是耶穌.
如何回答彼得呢.
弟兄姊妹.
這裡才是真正整段經文的重點.
由22到26節.
主耶穌根本沒有解釋.
為什麼.
反而還要凸顯.
你們門徒更加要有.
宜山的信心.
弟兄姊妹你來到教會.

$^{481}$這是不是一個.
只是節日的興趣.
因為你很想來親近神.
敬拜神.
而你有一個宜山的信心.
剛才早堂我都說了.
我們過年有一班同學.
我們聚餐.
神學院的同學.
不同的教會.
很多都已經成為堂主任.
我不說哪一間教會.
免得大家去追問一下.
有一個同學.
其實很年輕.
他其實是在教會長大的.
他回到教會.
現在成為堂主任.
他就提出一個很有趣的現象.
每年都這樣發生.
他說他們的教會.
每年的12月25日的聖誕節.
就會突然間多了超過100人.
平時是不來教會的.
平時是52個星期來1個星期.
這麼有趣的現象.
有些是他認識的.
一起長大的.
很奇怪的.
去家房問一下.
什麼原因.
你們為什麼只有聖誕節才來呢.
平日你們又不來教會呢.
不一而足的答案.
傳道你不知道.
我們平日1至5上班都很忙.
有些是做基層的.
但都不是全部是基層的.
你知道小朋友星期六要補習.
星期日只有上午最多時間.

$^{521}$睡晚一點陪家人喝茶.
所以就不行了.
有時候我又要進修.
不同的答案不同的原因.
問為什麼你們25日要回來.
因為那個日子特別不同.
我們一定要回來.
否則你們不認識我們之後.
將來我們的安息聚會.
找誰呢.
果然真的有安息聚會.
就立即找我的同學.
趕去醫院祈禱.
然後就辦安息禮拜.
原來有這樣的習性.
是這間教會.
請姐妹.
主耶穌在這裡的教導.
想群體具備兒生的信心.
相信不是一次.
你就能夠一年來一次.
你可以具備這一份信心.
彼此的認識.
在認識裡面.
一起有一個方向.
能夠知道什麼阻礙人.
去到神面前的山.
留意經文.
耶穌第22節就解釋.
其實沒有解釋給彼得聽.
為什麼我要遷就這棵樹.
背後有什麼神學.
有什麼目的.
更加不需要告訴他.
昨天我真的說重了話.
根本不是.
所以我們就不要幫耶穌解脫.
耶穌根本沒有解釋給彼得聽.
留意22節.
耶穌回答說.

$^{561}$你們當信服神.
我實在告訴你們.
無論何人對這座山說.
特有所指的這座山.
耶路撒冷當時就是他們處身的聖殿.
就是石安山.
到底耶穌是不是有所指呢.
石安山所承載的這個聖殿.
因為當時有些背景.
剛才我已經說了一點.
粉瑞黨的猶太人.
他們是要推翻羅馬政權.
密謀要來革命.
所以耶穌就是一個最好的政治領袖.
他們期望耶穌來做這件事.
所以他們都醞釀要將羅馬人驅逐.
回到地中海.
去歐洲的地方.
這樣的遺山.
原來當時是有這樣的盛行的想法.
耶穌是不是就著當時已經可能密謀的問題.
來提出這麼誇大的一個說法呢.
這座山.
相信這裡不可以按字面去解.
這座山就是石安山.
整個山移開.
我們都看到門徒.
這十二位的門徒.
因為是給得問.
耶穌就答.
這十二位的門徒或許是內圍幾位的使徒.
聽到耶穌這麼噁心的回答.
耶穌說的羅開.
這個吾疑惑.
因為下面就是說.
你們如果心裡不疑惑.
只信他所說的成.
必定成就.
就必定為他成就.
接著二十四節是再雙重確認.

$^{601}$所以我告訴你們.
凡你們禱告祈求的.
無論是什麼.
只要相信是得著了的.
就必定會得著.
二十二,二十三,二十四節.
三節聖經去確定.
祈禱那個的重要.
當然這裡問題就是.
如何去解釋這座山呢.
如果不能按字面去解釋.
也知道神不是一個好像沒有智慧的人那樣.
做一個大石自己都推不動的.
自己的神能.
自己的力量.
是自打嘴巴的.
自圓其說的.
相信神不會是這樣的一位的神.
更何況我們人是有意識的.
是有思想的.
我們不會突然之間.
吩咐對面的山.
新位的雞公嶺挪開.
是沒有任何意義的.
這個不是關乎你有沒有信心的問題.
是根本沒有意義的問題.
神是不會聽我們祈禱的.
但是如何去理解這個字面的這座山呢.
如果當時的聖殿坐落在石安山.
那個敬拜耶穌說不合格.
不合格不是說人不可以去到聖殿.
是聖殿真正的功能.
是人可以來朝見上主.
來得到心靈的平安.
與上主的相和帶來的救贖.
那份恩典.
以及自己認定做物主是掌管萬有的神.
是猶太人所敬拜的獨一真神.
這樣就告訴我們.
似乎這個這座山.

$^{641}$每個人都有不同.
有些人是他人生的價值觀.
那個排序是很混亂.
如果剛才那個例子.
那個現象.
這超過一百人.
一年來一次教會就是敬拜神.
那就是說他人生的優次序是排錯了.
他只是以為那個節慶就是一個罪咎.
不是他想著來敬拜神.
每個星期朝聖之旅來到神的面前.
這個是其中可能是一座山.
更多的似乎是我們內心裡面的自我中心.
對神對人的自我.
有時無限的放大.
透過不同的形式來放大了自己.
如果我們能夠意識到.
我們又為這樣的弟兄姊妹禱告.
慢慢他在人生很多不同的際遇.
發現原來自己是很多的自我自大.
而他能夠慢慢的降低.
豈不是這一座山慢慢的.
因為有很多的恩典.
陪伴和勸導.
而可以來到他移開嗎.
當然也有些是有疾病的.
我們能夠幫助他面對.
昨天的下午.
我們一直在教會祈禱會.
為一個弟兄.
他患了漸凍人症.
我們來到他祈禱探望他.
他說這個症令他感到絕望.
因為是有科學的證據.
只有18個月命.
是不能醫治的.
而且那個肌肉一直萎縮.
那些機能一直壞死.
是看著自己一直衰退.
但是那個頭腦是完全正常的一個人.

$^{681}$你可以想像得到這樣的一種禁閉.
他說他真的有一個朋友.
和他一樣的症狀.
真的18個月就離開世界.
是很恐怖的一件事.
你要面對自己一直看著自己萎縮離世.
但是他的分享反而令到我們兩位同工很感動.
他說不要緊.
我已經試過這麼多的辦法.
既然已經證實了是這樣的疾病.
我就和他共存.
我把握我現在的時間.
來幫助人.
18個月已經過去了.
醫生不斷和他會診.
其實每一次的治療診斷.
他說要睡在磁力共振的機裡面.
是完全不准動三個小時.
只不過是動一動手.
醫生就說不行了.
你剛才那三個小時不算了.
要從頭再來.
這樣只能睡著不准動.
你試想一下那種痛苦.
但是他最後告訴我們.
我選擇靠恩典與他共存.
對我自己來說是很大的震撼.
我們打算來到這裡為他祈禱.
給予安慰.
但是他自己能夠願意面對.
我邀請了他們夫婦.
接下來15號的祈禱會分享.
講述他們的內心的心路歷程.
人有不同的山.
每個人都有你自己的那座山.
主耶穌說你們如果有信心.
你信服神.
你就可以來到去移開那座山.
在教會裡面.
我們看到有很多的和好.

$^{721}$弟兄姊妹彼此之間的和好.
弟兄姊妹彼此之間.
單純的來同心服事.
豈不是這個群體經歷上主的恩典.
牧羊帶來的見證嗎.
豈不是我們同心一起面對很多的困難.
我們來改變很多社會的問題.
最近內地一套很熱門的電影《狂飆》.
其實就很真實地反映內地的結構性的貪腐.
那些保護傘.
拍得很真實.
我們基督教在香港.
很多的連結已經在我們心靈裡面的DNA.
我們能不能夠將我們這個DNA.
在不同的機構.
不同的地方.
來分享為什麼我們會特別的連結.
原因是因為我們相信.
做物主是審判的主.
我們要這樣做.
沒有神的人.
他不知道有一個主宰.
他真的不會有明白.
但是我們是有神.
敬畏神的人.
我們就知道.
這是我們要移開的山.
不同的人.
神賦予你有不同的恩賜.
才幹和你的位置.
如果你是有位置可以改變遊戲規則的.
你就去靠主.
祈求你有這個移山的信心.
改變那些結構性的問題.
求主幫助我們能夠在這段聖經裡面.
找到很重要的中心.
不過還沒停在這裡.
25節26節.
相信一定不是個人的那種.
好高騖遠的.

$^{761}$一種不切實際的神跡性的信心.
而是一個群體的信心.
因為耶穌突然之間.
將那盆冷水潑回門徒那裡.
如果彼得聽到.
耶穌我要這個信心.
我因為成為你的門徒.
我可以移山.
突然之間耶穌說.
你連向神閉上眼祈禱.
你都不夠資格.
你們站著禱告的時候.
如果想起有人得罪你們.
就要饒恕他.
好叫原來有一個因果關係.
你閉上眼祈禱.
聖靈馬上告訴你.
有一個人在你的心裡面.
你容不下.
你要立即先不要祈禱.
你要在當下.
先和神搞定這件事.
神啊.
我饒恕他.
閉上眼.
說完饒恕.
張開眼就要找他有機會復和.
這樣才有一個行動.
如果你閉上眼.
只說我饒恕.
神知道你其實沒有饒恕.
你都沒有行動.
你就不能夠得到神的饒恕.
因為好叫你們在天上的父.
也饒恕你們的過犯.
耶穌在另一個角度再說一次.
你們如果不饒恕人.
你們在天上的父.
也不饒恕你們的過犯.
其實二十五節就已經是這個意思.

$^{801}$不過再說一次.
強調再強調.
很多時候我們沒有兒生的信心.
歸根結底.
不是天上的Wi-Fi不夠強.
是我們沒有開Wi-Fi.
我們塞住自己的心靈.
明知而繼續的不寬恕.
不諒解.
不和好.
我們一起祈禱.
多謝主賜給我們能夠有這個信心.
真的能夠兒生.
不是一些不切實際的.
好像玩的神蹟.
而是我們能夠看到人生.
很多人需要我們去幫助.
有些人可能是基層的.
或者徘徊在一些犯罪的.
一些邊緣的人.
例如一些囚犯.
吸毒者.
槍記.
其實沒有一個人天生下來.
就喜歡做作奸犯科的事.
很需要教會.
很需要我們有神愛的人.
去幫他們納入一個正途.
我們很懇求主.
讓我們能夠在你的面前謙卑.
我們自己明白.
我們有禱告的權柄.
這個權柄不是用在我們個人的.
要怎樣去有好的生活.
富有.
或者只是自己個人的一個福樂.
而是我們看到.
這個是群體的牧羊.
是一個彼此在群體裡面的依治.
尤其是心靈和好的依治.

$^{841}$願主幫助我們.
使到教會能夠在這個時代.
有很多地方需要愛.
需要復和.
求主施恩幫助我們.
我們這樣禱告奉耶穌基督的名.
阿們.
\newpage



\section{}
\label{sec:lrudeppALpc}
\textbf{2023-02-09 周四晚查經班《被擄歸回歷史研讀(6)》 講員:曾裔貴牧師}
\newline
\newline
連結: \href{https://youtube.com/watch?v=lrudeppALpc}{\texttt{https://youtube.com/watch?v=lrudeppALpc}} ~~~~ 語音日期: 2023-02-10
\newline
\newline
\hyperref[sec:k6XACT4Tfao]{\small{< < < PREV SERMON < < <}}
~
\hyperref[sec:index_chronic]{\small{[返順時目]}}
~
\hyperref[sec:index_scriptual]{\small{[返順卷目]}}
~
\hyperref[sec:P577AmrjvF0]{\small{> > > NEXT SERMON > > >}}
\newline
\newline
$^{1}$我們一起祈禱.
多謝主讓我們有時間去研經.
幫助我們進入以色列第二階段的歸回.
幫助我們從其中的經文明白作者的寫作意圖.
也讓我們從幾次的被擄歸回能夠掌握猶太人的心路歷程.
願你親自帶領我們.
我們等候禱告奉耶穌基督的名.
阿們.
我們來看看以色列的第四章.
我想給大家一些背景.
今天我們主要講第七章.
暫時不講撒迦利亞書.
因為上星期我們講了一到第六節撒迦利亞書的背景.
因為是緊接哈該四個的密使.
大家記得那個日期.
六月初一到六月二十四開工.
七月二十一到九月.
總共有四個密使.
同時撒迦利亞先知在同一個時期.
也有一個密使臣臨到他身上.
他來到全港.
由第七節開始.
一直到撒迦利亞書的結尾.
也有八個異象.
一對一對的.
所以我們遲些有機會再補充異象的解釋.
反正讀歸回的歷史書先知書.
一併去講一下也好.
今天是第二個階段的被擄歸回.
這個已經過了一段長時間.
也可以這樣說.
在巴比倫第一次上去的人.
他們回去了.
但是我們也講過在波斯帝國.
有很多猶太人沒有歸回重建耶路撒冷.
那些沒有歸回的同樣在波斯帝國.
有他們的生活.
有他們的繁衍後代.
有他們的工作.
有些甚至在皇室有官職.

$^{41}$例如以色列和尼希米也是.
所以你可以想像同步在耶路撒冷有這些歸回的人.
大概五萬人.
在波斯帝國其實人數更加多.
他們在那裡.
接著第二次的被擄歸回.
就是由一個人.
就是這個作者以斯拉.
由他引發的一個歸回.
這個歸回他回到去不是重建聖殿.
也不是重建聖城.
而是重建聖文.
所以這裡一直涵蓋到.
由第七章開始.
一直到應該是十三章的以色列的完結.
其實是一次的奮興會.
一個奮興會的記錄.
就在第七章開始來到這本書的結尾.
所以也是容易分段的.
我們要給一點背景.
第一次當他們來到去.
537年主前.
他們就是第一批.
由索羅巴伯.
大祭司約書亞.
來帶領他們因為古列王的頒佈詔書.
全國的人無論你是什麼國籍.
你可以回到你的原居地.
為著你們的神明.
你們的信仰來興建聖殿.
不同的民族.
總之敬拜神.
祭祀神的殿.
猶太人就歸回了.
那個時候大概主前的537或536.
一般是537.
接著就過了一段時間.
去到520年.
重新再重建.
原因就是第四章.

$^{81}$我先給大家知道一下年份.
這裡有些關鍵.
第四章第六節.
或者第四節.
那地的民就在猶大人.
即當地的原居民.
即非猶太人的土著.
就在猶大人建造的時候.
使他們的手法軟.
擾亂他們.
從波斯王古烈年間.
直到波斯王大利烏登基的時候.
大利烏是第三任.
第二任就是亞達舍西.
他講的時間.
那個range是由.
古烈可以回到建造.
去到大利烏.
是停頓了.
大利烏王的第二年.
主前大概是520年.
他們就正正式式.
因為哈該先知.
撒迦利亞先知.
來奉神的名.
來講寓言.
默示.
激動他們的心.
所以他們就重新建造.
即是停了16年.
在這個時候.
他們賄賣謀士.
其實是亞達舍西的時候.
要敗壞他們的謀算.
留意第六節.
在阿哈徐魯財登基的時候.
上本控告猶大和耶路撒冷的居民.
這裡就比較難去理解.
因為第七節.
亞達舍西年間.

$^{121}$比斯蘭米特利達他別和他們的同黨.
上本就告波斯王亞達舍西.
這個就是原因.
一直停到大利烏王的時間.
一直阿哈徐魯.
去到亞達舍西.
已經經過了兩個朝代.
其實那個情況都沒有改善.
所以亞達舍西就一直頒那個禹旨.
這個是第23節.
因為第七章有另外一個亞達舍西.
所以就不要混亂.
23節.
亞達舍西王的常遇.
獨在利宏和書記身上.
並他們的同黨面前.
他們就急忙往耶路撒冷去.
見猶大人.
用勢力強迫他們停工.
於是在耶路撒冷.
神殿的工程就停止了.
直到波斯王大利烏第二年.
就是那個背景要交代.
是因為第一份的告的禹狀.
告到亞達舍西的身邊.
亞達舍西王就終於.
因為這個這樣的告狀.
就來調查.
傾向對當地的土著的支持.
而逼令猶大人要停工.
好像不理先王這個古列.
頒佈的詔書.
來建造聖殿.
所以這個亞達舍西.
其實就不是給猶太人.
所以就一直停.
停到大利烏的第二年.
所以這個就是當時大概520年.
520年的時候.
大利烏又來到了.

$^{161}$這個第五章那裡這樣說.
那個背景我們說過了.
就是神激動哈該先知.
和薩加尼亞先知.
來勉勵當時的領袖.
所羅巴伯和約書亞.
和其他的人.
要去勉勵.
為神的殿一起同心的去建造.
所以在這個時候就是大利烏王的第二年.
他們就正正式式.
重新的來到興功建造.
所以這個就是大概520年.
用了四年的時間.
去到大利烏王的第六年.
聖殿就建好了.
這個就是去到第六章.
第六章的十三節.
很特別的.
大利烏又是收到告狀.
又是當地的原居民.
告尤大人.
說他們現在又再來建造了.
那要了解一下什麼情況.
會不會有一些事情.
他們想煽動.
或者想要革命等等之類.
大利烏就要去尋查那些典籍.
不過大利烏都是一樣.
面對的那個告的預狀.
但是他反而去尋查那些典籍.
所以有時神的工作.
又很難去想像.
因為之前我們說過.
有先知在民間.
在尤大人中間.
說了一些勸勉激勵的話.
所以他們就在這樣的感動下.
大家在神的旨意下.
大家一起去做.

$^{201}$接著又是招致當地人.
看到尤大人重新.
好像死灰復燃那樣去建造的時候.
他們又告預狀了.
不過這次同樣都是那些控訴.
但是這次第六章就說.
於是大利烏幹子要尋查典籍.
庫內就是在巴比倫藏寶物之處.
在馬袋省阿瑪他城的宮內.
尋得一卷.
其中記著說.
古列王元年.
他幹子論到耶路撒冷神的殿.
要建造這殿.
為獻祭之處.
建立殿的根基.
殿要高多少爪等等.
接著就是尼泊爾加尼沙朗所努力的器皿.
都要歸還給耶路撒冷的聖殿的用處.
所以現在他們頒佈那個詔書.
同樣都是那個告狀.
但是這一封信.
反而就更加去肯定.
得到一個王的允許.
那個的預准.
那個的審批.
所以就現在索羅巴伯和約書亞.
他們就不只是靈性上去完奮興會.
去完退休.
有一個的熱誠.
而是在神那裡有先知的默示.
他們可以來到去做.
而且要做.
不可以再躲懶了.
接著連在政治上都配合.
記住.
是先有神的默示.
感動.
和先知的安慰.
接著有領袖去get到這個訊息.

$^{241}$他們就有那個責任.
統領著群眾一起去建造.
接著才有那個政治上的開綠燈.
所以不是政治開綠燈.
然後就反過來.
神都要來到去允許.
兩個過程.
那個先後次序.
有一點不相同.
所以這裡就告訴我們.
這個就是當時.
都同樣是索羅巴伯.
都同樣是大祭司約書亞.
但是兩個的情況.
是天淵之別.
這個時候.
就是去到六章的十三節這一段.
這個就是一個背景.
給大家知道第一次的建造.
因為今晚我們要講的是第二次的建造.
那個重建的建造.
第十三節六章.
於是何西總督達賴和士他波斯賴.
並他們的同黨.
因大利烏王所發的命令.
就急速進行.
由大長老恩先知哈該和易多的孫子.
撒迦利亞所說勸勉的話.
就建造者奠凡事亨通.
政治的體系.
就是人間見得到的那一層.
開綠燈.
開綠燈的背後原來是因為.
神的時間.
他們初期.
他們因為那個阻礙.
停了十幾年了.
現在呢.
有先知的勸勉.
當然是奉神的命.

$^{281}$來勸勉.
所以他們就勉勵.
兩位領袖.
而他們就真的這樣去做.
所以呢.
你可以說這個水稻渠城.
是先有神的激動.
然後有民間都配合.
連外邦的波斯帝國.
這個大利烏王.
一樣的配合.
就可以想像得到.
我們就不要.
將我們全副的心思.
放在政治的那個環境.
希望那個君主比較開明一點.
可能是的.
希望不要那麼殘暴.
但是.
不到你們去這樣想的.
一定只能夠我們在神的旨意裡面.
去做我們那個應份的部分.
地上的.
無論是嚴謹.
或者是比較柔和.
都不到我們去操控.
我們應該要來交託給神.
因為這個是聖功.
不是我們人間的一些的功能那麼簡單.
所以他們就.
在這個勉勵底下.
他們就建造這殿.
凡事亨通.
他們遵著以色列神的命令.
和波斯王古列大利烏.
阿達舍西的旨意.
建造完畢.
四年時間.
第十五節.
大利烏王第六年.

$^{321}$阿達月初三.
這殿修成了.
四年裡面.
包括尋查典籍.
包括他們重新.
那個荒廢的.
爛尾工程.
重新來到可能招標.
可能來到預備材料.
可能來到計劃.
劃職.
接著就開工了.
四年時間前後.
由那個預旨開始.
頒布的大利烏王的第二年.
到現在就是第六年.
阿達月的初三.
這殿就修成了.
這個是出於神.
所以.
第十六節.
以色列的祭司和利尾人.
並其如被老歸回的人.
都歡歡喜喜地.
行奉獻神殿的禮.
接著說他們的數目.
接著去到就是.
既然正月初三.
那個聖殿進宮了.
所以呢.
即是完成了.
所以呢他們就第一次.
這次是一個更加完整的逾越節了.
十九節.
正月十四日.
被老歸回的人守逾越節.
原來祭司和利尾人.
一同自結.
沒有一個人不結成.
利尾人為被老歸回的眾人.

$^{361}$和他們的弟兄眾祭司.
並為自己.
來到去宰逾越節的羊羔.
從老道之地.
歸回的以色列人.
和一切除掉所染.
外邦人污穢.
歸附他們的.
可能是行割禮.
要尋求耶和華以色列神的人.
都可以吃這個逾越節的羊羔.
因為是規定了的.
你如果走遠路.
你染了一些污穢.
譬如那些屍體.
那些動物屍體.
或者你是去處理一些屍體.
或者一些污糟的東西.
那你就要自結的.
這個就是他們的規例.
這個是逾越節的規例.
大家都自結了.
就來吃逾越節的羊羔.
第22節.
歡歡喜喜手除教節七日.
因為逾越節是正月十四.
正月十五.
一直到正月二十一.
就是除教節.
第一天是逾越節.
接著就是除教節.
因為他們.
因為耶和華使他們歡喜.
又使亞述王的心.
轉向他們.
堅固他們的手.
作以色列神殿的工程.
這個就是第一次的重建.
順利的完成.
就是在主前大概是516年.

$^{401}$這個是大利烏王的第六年.
留意第七章.
因為突然間出現的亞達舌西.
這是以後.
波斯王亞達舌西年間.
有個以色拉.
他是西來亞的兒子.
開始說他的家譜.
他的出身.
這樣來到他.
這麼長的介紹他自己是哪個家族.
是哪一個的支派.
有什麼原因呢.
原因就是.
他是要合乎那個資格.
來重建聖文的心靈.
因為第七章的十.
等一下.
我看看那個經文.
或者我等一下先讀到.
暫時找不到.
我等一下讀到.
我就告訴大家.
有一個背景的.
就是為什麼以色拉要申報他的家譜.
他的重要的利益.
第七章的第一節.
這是以後.
一個轉接.
一個段落.
其實由第六章和第七章之間.
其實已經相距了一段長時間.
可能接近有60年.
因為這個亞達舌西.
就不是之前那個亞達舌西.
已經是後兩代的君主.
都是波斯帝國.
所以在這個時候.
意味著大家要有一個平衡的想法.
就是以色拉.

$^{441}$肯定是在波斯帝國出生的.
可能第二代第三代的猶太人.
在波斯那裡人育.
所以他的學養.
他對律法的認識.
在帝國裡面.
他的學問等等.
都是幾出類拔萃的一個人.
有個以色拉.
他是西來亞的兒子.
西來亞就是亞撒尼亞的兒子.
亞撒尼亞就是希臘加的兒子.
希臘加就是沙隆的兒子.
沙隆就是薩都的兒子.
薩都就是大衛王年間的祭司.
薩都是阿希特的兒子.
阿希特是亞瑪利亞的兒子.
亞瑪利亞是亞撒尼亞的兒子.
亞撒尼亞是米拉約的兒子.
米拉約是西拉希亞的兒子.
西拉希亞是烏西的兒子.
烏西是布基的兒子.
布基是亞比書的兒子.
亞比書是菲尼哈的兒子.
菲尼哈是以利亞撒的兒子.
菲尼哈就是在民數記24章.
去殺死其中一個族長.
就是西面支派的森尼族長.
這個就是菲尼哈.
菲尼哈就是摩西的孫.
對不起是亞倫的孫.
菲尼哈是以利亞撒的兒子.
以利亞撒是大祭司亞倫的兒子.
是不是很清楚.
他講自己的背景.
即是說他是根正苗紅的大祭司亞倫的後人.
即是說他是利美洲派.
而且不是普通的利美洲派.
還要是亞倫的子孫.
即是說他一出生就已經是大祭司.

$^{481}$當然他要訓練他要學.
但是他一出生他的身份就是大祭司.
因為他是生在利美洲派.
他是亞倫的後人.
所以耶穌不可以在地上做祭司.
但是耶穌基督是根據另外一個體系.
就是默基西德為永遠的祭司.
所以因為希伯來書都講的.
如果耶穌在地上.
他不是亞倫之派.
他不可以做祭司的.
他不可以進去聖殿來代表獻祭的.
但是他獻的是屬靈的祭.
是他的生命.
所以希伯來書第七章就說.
默基西德是無父無母無始無終.
是與神的兒子相似的.
說的是永遠的祭司的職任.
所以這裡是地上的.
他們回去重建都很講求歷史.
支派和律法.
所以他們要尋查那些典籍.
所以這個不用尋查了.
他本身就有個家譜在那裡.
他自己都知道自己是亞倫的後人.
這就很重要了.
因為當他回到那裡.
記住那個時候是聖殿剛剛建好.
沒多久.
這些消息慢慢慢慢滲到波斯了.
波斯帝國可能又經過改朝換代之後.
他們又准許.
有新一批不同地方的人.
可以回去自己的原居地.
繼續的履行那些憲制的條例.
因為我等一下會讀給大家聽.
有一段聖經是波斯王他們比較懼怕的.
所以他們准許這些祭司.
就可以回去做這些聖職的工作.
第六節.

$^{521}$者以斯拉從巴比倫上來.
當然等一下就會說.
用多久的時間.
由巴比倫上去耶路撒冷呢.
是足足用了四個月時間.
那個時候真的用腳.
用駱駝這樣走的.
還有會有土匪等等.
還有晚上會有野獸.
所以是不容易的.
當然記住.
還要是帶著很多的金銀財物.
所以他們真的要步步為營這樣上的.
這裡告訴我們.
者以斯拉從巴比倫上來.
他是敏捷的文士.
通達耶和華以色列神所賜摩西的律法書.
是通達的.
敏捷的文士.
就是說他對摩西律法書.
對舊約是相當的純熟.
那個學習.
那個掌握.
那個運用.
都是很敏捷的.
就是說那個思辨的能力很強的.
能夠舉一反三.
能夠實踐應用的.
接著就是通達神的律法書.
王允准他一切所求.
是因耶和華他神的手幫助他.
在這裡開始出現這個詞彙.
神的恩手幫助他.
出現了多少次呢?你猜猜.
在這兩張聖經出現了六次.
神的手幫助他.
但是其實.
是他求王允准.
他來到翻上去重建.
這裡是讓我們看到.

$^{561}$好像沒有一個底觸.
他看到的是神的手.
在他的計劃.
在他的安排.
但是他又是求王允許他翻上去.
就是說他在皇宮裡面.
是可以立在王面前的一個祭司.
這個是猶太人身份的祭司.
但是他在波斯.
他是不是做一個在王面前的.
可能是歷史官.
可能是其他的文官.
這個我們就不太知道.
但是很特別的地方就是.
由七章第六節開始.
總共有八次他提到.
在兩張聖經裡面.
他提到他一切所做的.
都有神的手幫助他.
當然不是真的看到一隻手.
而是看到那個通達.
那個順暢.
他的計劃.
他的思考.
和他的策劃.
都好像有神的開綠燈的那個恩手.
所以神的手.
就是代表著神的那個Blessing.
就是說他清楚神要他回去.
他的位份.
他的裝備.
他的學養.
足以可以來到去.
一方面就是讓耶路撒冷的猶太人.
不會對他的身份有懷疑.
接納他的事奉.
因為他經過訓練.
是一個敏捷的文士.
既是文士.
同時間又是祭司.

$^{601}$又是一個在皇宮裡面的領袖.
又是一個可以.
來拿到皇的一些資源.
來重建耶路撒冷自己的故土.
自己的民族.
所以這個人好像是.
剛剛在這個時代出現的.
很重要的一個人.
所以得人心者得天下.
這個人的心是歸向神.
因為他自小.
他在巴比倫.
在波斯帝國長大.
但是他形容他自己.
他是一個敏捷的文士.
是通達神的律法.
是有這一種神給他的保守引導.
神的手引導他.
然後他也不抗拒.
王允順他一切所求.
他求什麼呢.
待會後面又會有交代了.
那你就可以看到一個這樣的人.
結合你可以說政治.
他的地上的智慧.
他的裝備.
可能大學畢業碩士.
可能他是讀不知道什麼工程.
或者是文學.
但是他的major是神的律法.
我們試試找找調查.
我們的弟兄姊妹.
很多在外國回來.
讀完一次聖經.
有些碩士.
在一些很出名的大公司工作.
前途無可限量.
有些是博士學位.
我很想問一下弟兄姊妹.
這些年輕的弟兄姊妹.

$^{641}$有沒有自己裝備.
讀通神的華語.
我懷疑連福音書都沒讀完.
不要說讀通.
這樣就很大件事了.
一個時代需要的是.
兩方面都有裝備.
不只是自己躲起來讀聖經.
什麼都不懂.
又不懂思辨.
又不懂得各樣的智慧.
這樣就變得有他的限制.
但是更多的.
基督教現在在香港是極中產.
是富裕階層.
甚至是統治階層.
政府的官都不少.
但是有多少是有很深的學養.
很高的智慧.
但是很屬靈很敬虔.
很對神的華語有一個愛慕.
認識.
而能夠通達神.
透過摩西所頒佈的律法書.
以色列是兩者兼備.
整個時代就需要這個人出現.
這個人就帶來教會.
民族的復興.
所以今天仍然在等待這樣的人出現.
我這個時代.
都是公認.
譬如北宣剛剛換屆.
新的堂主任出現.
其他的譬如全宣也換了一段時間.
張雲生牧師.
不同的大堂會.
譬如沙田的宣道會.
又是堂主任換得很好.
等等.
一定要接上一代的棚.

$^{681}$問題是能否具備.
牧人,領袖,學者.
兼具的智慧.
通達的生命.
這樣的學養.
這個就真的很需要.
大家為不同的教會.
為一個時代的轉變.
去禱告.
去祈求有新的領袖.
具備有學養.
具備有學識.
具備有通達的智慧.
敏捷的思考思辨的能力.
和更加對神的話語.
有一個通達.
第七節.
亞達斯西王第七年.
以色列人祭司利未人.
歌唱的守門的尼提靈.
有上耶路撒冷的.
王第七年的五月.
以色拉就到了耶路撒冷.
用了原來是五個月.
第九節.
正月初一.
他從巴比倫起程.
帶領的就是各個恩賜.
各個不同的部門.
全部是在聖殿裡面.
事奉的那些職位.
崗位的人.
用了五個月時間.
由巴比倫上來.
當然是以那個做一個單位.
其實都是波斯帝國.
巴比倫就已經滅亡了.
可能就是米索波大米.
那一帶地方.
來到上去.

$^{721}$就是兩河流域.
右法拉底河.
和這個叫底格里斯河.
記住是上底下右.
就是兩條河.
上面一條就是底格里斯.
下面一條就是右法拉底河.
就是當時文明的古國的地方.
所以這裡說.
由五月初一.
留意第九節.
正月初一.
他從巴比倫起程.
因他神施恩的手幫助他.
五月初一.
就到了耶路撒冷.
以色列定治.
考究.
遵行耶和華的律法.
有三個形容詞.
定治.
一生的治向.
他來定了這個治向.
然後考究.
考究就是說.
去study.
去做一個嚴謹的study.
不只是背一隻金句.
神愛世人.
就是你在聖經裡面的認識.
那就很慘了.
基督教已經沒落到一種文盲.
對聖經是非常不認識.
非常不純熟.
不是說我們上一代.
很多老人家.
不識字.
信耶穌.
你想他怎樣呢.
他都不識字.

$^{761}$他能夠聽得到聖詩.
能夠跟著去背一些簡單的金句.
學學寫字.
我們已經覺得很偉大了.
但是第二代第三代.
已經是普遍有大學程度了.
但是對聖經從來都沒有真正的定治.
考究.
更加說不上進行.
所以這個以色列是為這個時代而預備的.
一位神的僕人.
大祭司的身份.
有民事的裝備.
有社會的經驗的閱歷.
因為在皇宮裡面.
他是求王.
王允准他一切所用的.
你想想幾個的身份都具備在他的內心.
當然更重要的.
就是他對神的單純敬畏.
跟隨神.
因為這個工程很大.
你要上去離開你的安書區.
上去耶路撒冷.
一定是你聽到耶路撒冷.
有一些的消息.
而令到你感到不安.
感到要為這個地方祈禱.
可能有沒有小時候上過去呢.
好像回鄉下那樣.
那聖經沒有交代.
但是為什麼在波斯帝國.
這些沒有上去.
第一批上去的猶太人.
為什麼可以在波斯帝國的境內.
居然都可以培育到這麼敬虔的猶太人呢.
這個是我們要問的問題.
如果我們有很多人移居到外國.
是香港人.
下一代了.

$^{801}$這一代過去了.
下一代再下一代.
我們能不能夠把握.
在加拿大.
在英國.
在美國.
或者在荷蘭.
能夠有敬虔的人.
還是錦宣家.
要收窄範圍.
是錦宣家的第三代.
在外國的了.
能不能夠有以色列的敬虔呢.
我們擔心他可能連信耶穌都不信.
不要說這麼誇張.
他能夠定志.
考究.
遵行.
又通達.
又敏捷.
這樣的眾多的形容詞.
是在說他那個腦袋.
以色列這個人的心靈.
他的腦袋和他的心靈.
是緊密地結合.
弟兄姊妹是很不容易.
讀到這些經文.
令到我們擔心下一代.
所以我們擔心他能不能找到工作.
擔心他其他的升學就業.
結婚.
生孩子.
會不會將那個次序再擺高一點.
我們擔心他愛不愛主.
敬虔.
那個就更加重要.
因為有了敬虔.
有了愛主.
其他他自己會在神的恩手的帶領下.
他自己會去追求.

$^{841}$我們就不要一味為了地上的.
他的生活的存在.
來過分的擔心.
第十一節.
祭司以色列是通達耶和華.
解明和賜以色列之律例的文士.
剛才已經說了.
用了這麼多篇幅說那個家譜.
說他的背景.
說他在皇室的工作.
然後現在說的就是文士.
是說他真的來到教聖經.
文士是教聖經.
因為離開了耶路撒冷.
沒有了聖殿.
他們要堅持在外邦.
保持自己的民族的獨特性.
有會堂的開始.
在外地他們有幾個家庭.
四至五個家庭.
一般是十五個男丁.
他們就會有自己的會堂.
最重要就是有男丁家長.
來讀聖經.
要保持對神摩西的律法的跟隨.
當然要守節旗.
當然他們會繼續教導自己的兒女.
讀聖經對神的話語的敬重和遵行.
所以就會有這些文士的制度的衍生.
所以這裡再介紹.
他有另外一個身份.
就是文士這個身份.
當然這個身份帶到第八章.
那個復興會.
那個奮興會.
由第八章開始.
一直到後面的最後一章.
是一次的水門的事件.
Watergate.
十三節下面.

$^{881}$亞達舍西王賜給他御旨.
上面寫著說.
諸王之王亞達舍西.
達於祭司以斯拉通達天上神律法大德的文士.
又是大祭司又是文士.
云云.
住在我國中的以色列人.
祭司利未人.
凡是甘心上耶路撒冷去的.
我降旨准他們與你同去.
王與七個謀士既然差你去.
照你手中神的律法書.
查問猶大和耶路撒冷的警方.
猶大金銀.
就是王和謀士.
甘心獻給住在耶路撒冷.
以色列神的.
並帶你在巴比倫全省所得的金銀.
眾籌.
全省裡面的猶太人.
一齊樂意為神的電獻奉獻.
百姓祭司樂意獻給耶路撒冷的.
他們神殿的禮物.
所以.
你當用者金銀急速買.
供牛供綿羊羔.
和銅獻的掃祭.
電祭之物.
獻在耶路撒冷你們神的壇上.
剩下的金銀.
你和你的弟兄看著怎樣好.
就怎樣用.
總要遵著你們神的旨意.
所交給你神殿使用的器皿.
你要交在耶路撒冷神面前.
你神殿裡若再有需用的經費.
你可以在王的府府.
在河西裡面來支取.
其實應該是河東.
我亞特斯西王.

$^{921}$又降旨與河西的一切父官.
說他們在東面.
所以以色列就是西面.
通達天上神律法的文士祭司以斯拉.
無論向你們要什麼.
那個王這樣來稱呼自己的一個神子.
這個以斯拉.
是通達天上神律法的文士祭司.
通達天上神律法的文士祭司.
你想想連王都知道.
以斯拉是通達天上神的律法.
不是普通法.
不是地上的大陸法.
不是其他普通的不成文的法.
是天上神的律法.
傳道人應該是這樣.
我們的裝備.
我們的一生的追求.
就是要這樣來對神的律法認識.
來教導弟兄姊妹一起明白.
這個是我們天職的責任.
現在王這樣來稱呼.
這樣來頒佈.
是叫當地的河西的官.
那些管庫房的官.
要來聽他所說.
無論他向你們求什麼.
你們要速速地被辦.
就是銀子值到一百他年得.
是那些重量.
墨紙一百開爾.
那些度量衡的單位.
就是多少斤多少磅.
酒因為液體.
一百的巴特.
譬如可能是多少公升.
油一百巴特.
鹽就不計其數.
鹽是值錢的.
你在內陸.

$^{961}$也要給他凡天上之神所吩咐的.
當為天上神的殿詳細辦理.
就是這裡了.
就是以斯拉他為什麼這麼緊要.
說他自己的家譜.
就是這裡了.
因為原來當時有一個想法就是.
亞達舍西王擔心自己的國家.
拜一神多神庇佑.
如果沒有人去敬拜.
和去打理那個聖殿.
如果有什麼事情降禍於我們.
那就不值了.
所以關乎到.
如果有一個合格的人.
能夠上去重建.
那個人就是.
要什麼可以有什麼.
因為是關乎到我們的帝國.
怕著招惹了神的憤怒.
所以這個就是第十九節.
我想跟大家讀的這節經文.
為何使憤怒臨到王和王的種子的國呢.
這個二十三節.
所以凡是天上神所吩咐的.
當為天上神的殿詳細辦理.
記住是眾數的天上神.
亞達舍西王不是特別厚待猶太人.
不過他們可能有一個想法.
不知道是不是占卜還是怎麼樣.
就是說突然間大發神心.
就是說天上神頒佈的.
就要根據他們的那個律例.
來去辦理.
現在對了.
有個以色拉來求我.
說要回去耶路撒冷.
他呢.
又被驗證過通達天上神的律法.
天上神是亞達舍西常常說的.

$^{1001}$現在你們猶太人就是拜天上的神.
好了.
他又是通達.
又是祭司.
就是說他的身份.
他的出身.
他的裝備都可以符合這個剛剛的國策.
所以呢.
以色拉就順理成章.
可以帶著百姓.
更重要就是帶著王的那個御旨.
和所需用的物件材料.
不過記住.
一而再.
以色拉說的是天上神的恩手幫助他.
很不簡單.
有那種運籌為額的世界仔的智慧.
但是整個身心靈.
是認定是神的手來去掌管著.
是很難的.
真的很難的.
因為有智慧.
有地上的學養知識.
就容易.
口就說是神的恩典.
但實際上.
說笑一句多謝你.
你都可以不開心的了.
因為你很容易就會將榮耀不知不覺歸給自己的智慧.
自己的人脈.
自己的其他東西.
就很少會歸給神.
所以.
剛才我們都說了.
不單止沒有抵消到以色拉他的地上的智慧.
他的學養.
他的身份.
他的出身.
和他的裝備.
完全不單止沒有抵觸.

$^{1041}$甚至神會大大使用他.
不過更重要的是.
他自己.
這個領袖.
他要看到是神的手.
是一路帶領著他.
這個是七章第六第九節.
和一會兒二十八節都有的.
接著他就曉喻了.
第二十四節.
我就曉喻你們.
至於祭司利未人歌唱守門離提靈.
離提靈就是那些做那些雜役.
就是雜工.
搬搬抬抬的.
沒有固定的崗位的.
守門就很清楚了.
守住耶路撒冷那些不同的聖殿的門.
並在神殿當差的人.
不可以叫他們進貢交貨納稅.
不用給稅的.
神職人員免稅.
在當時的國策.
德國是由政府出糧給傳統人的.
德國的教會.
我不知道是不是免稅.
他做的研究教會都是公務員.
由政府出糧的.
如果香港這樣都不錯.
我們可以專心去事奉.
接著有一間公屋.
簡約公屋.
每一個傳統人都有的.
這樣都不錯.
第二十五節.
主子要照著你神賜你的智慧.
王勸他.
將所有明白神律法的人.
納他做事司.
做審判官.

$^{1081}$治理河西的百姓.
使他們教訓一切不明白神律法的人.
怎麼會有這樣的君王.
是不信主的.
外邦的君王.
突然之間.
好像去完培靈會.
就叫以色拉招募你們的猶太人.
總之你認為他們有智慧.
又是明白神律法的.
你就招募他們.
在上去之後.
就做事司.
做審判官.
原來懂得神的律法.
人就會全面一些.
不會偏心.
不會有一種傾側.
是會持平的審判.
這個是審判官最需要具備的.
就是公平公正.
所以對神的律法.
有一個智慧和掌握.
都是會帶來做事的公平.
所以這裡.
告訴我們.
要教導一切不明白神律法的人.
第二十六節.
凡是不遵行你神律法.
和皇命令的人.
就當速速定他的罪.
或者自死.
或者充軍.
或者抄家.
或者囚禁.
誇張了一些.
運用了政治的手段.
似乎以色拉不需要出到.
這個所謂的上方寶劍.
他只是好好地去教導神律法書.

$^{1121}$就會有復興.
果然第九章就是一個復興.
我們可能今晚不夠時間說.
那裡的了.
不過我們再繼續下去.
第二十七節.
以色拉說.
耶和華我們列祖的神.
是應當稱重的.
因他使得王起這個心意.
修昔耶路撒冷.
耶和華神的殿.
又在王和謀士.
並大能的君長面前.
施恩於我.
留意一下.
第三次出現.
對不起.
第四次.
這裡是第三次.
因耶和華我神的手幫助我.
我就得以堅強.
從以色列中焦聚首領.
與我一同上來.
焦聚的不是他的班底.
是要有具備神的律法的智慧.
那個通達.
那個全面.
還有那個心.
都是很想上去.
去重建神殿的工程.
所以以色列就在這裡.
得到王的允准.
來做.
由正月初一.
亞達茜西的時候.
去到五月初一.
上去.
第八章就說他們歸回的記載.
第九章就是回到上去之後.

$^{1161}$就開始來面對.
那個現況.
帶著這個理想.
帶著這個差遣.
帶著王這個突發奇想的心願.
就是整個耶路撒冷都能夠歸向神.
明白神的律法.
透過以色列和他的陪伴.
他的同伴.
來教導耶路撒冷的猶太人.
和一切歸向神的人.
包括外邦.
進猶太教的波斯人.
或者是當地的一些原住.
原居地的人.
第八章就比較詳細.
說回之前的事情.
他們怎樣來到上去.
這五個月.
所以都很工整的.
先介紹了以色列這個人.
他的身份.
他的為人.
他的裝備.
和當時有一個很特殊的政治氣候.
是王自己頒佈詔書.
邀請以色列上去.
作為代表.
所以以色列真的能夠遊走於皇室.
和又是這麼敬虔.
對神律法這麼純熟認識.
我們都有很多的政務官是基督徒.
忙到一個地步.
有沒有時間讀聖經呢.
禮拜有沒有機會去崇拜呢.
可能都未必有的.
所以都要為這些高官祈禱的.
他們要很多東西去決策.
需要有通達神智慧的.
是為整個香港的福祉的.

$^{1201}$所以不怕為這些官員祈禱的.
真是他們很需要具備.
一個良好的氛圍.
那個政治的氣候又是好.
友善.
制定的政策.
都是能夠促進.
神國度的福音的興旺.
所以我們不怕為官員祈禱的.
所以這些是我們要為不同的.
位份的人來禱告的.
使得他們能夠有一個好的機遇.
機會有神的恩手.
在他們的事奉裡面.
工作的裡面去幫助他.
第八章就記載他們回到上去.
當亞達斯西王年間.
和我從巴比倫上來的人.
他們的族長和他們的家譜.
就記在下面了.
因為他們是要回去建造聖殿.
上到去發現原來聖殿不需要再建造了.
要建造的是聖殿聚會的人.
所以如果他們沒有這個.
祭司的身份.
民事的裝備.
上到去就有困難了.
所以上去不是看熱鬧.
不是回去旅行.
是回去要來去服侍.
要來去有那個復興.
所以如果不是自己有一個根底.
打好對神的律法的認識.
和對神的愛慕敬畏.
其實就很難做到感動人心了.
他們的族長家譜記在下面.
屬菲尼哈的有甲順.
屬以他媽的子孫有但爾尼.
屬大衛的子孫有哈特.
屬巴祿的後裔.

$^{1241}$這些全部是大族.
是大的支派.
就是釋迦尼的兒子.
有撒迦利亞.
同著他按家譜計算.
南丁一百五十.
屬巴哈摩亞子孫.
有西拉希亞的兒子.
以尼約乃.
同著他有南丁二百.
屬釋迦尼的子孫.
有這個瓦哈式的兒子.
有南丁三百.
屬亞丁的子孫.
有約拿丹的兒子.
以別.
同著他有南丁五十.
屬以蘭的子孫.
有阿他利亞的兒子.
耶西亞.
同著他有南丁七十.
屬史法提瓦的子孫.
有米加勒的兒子.
西巴第亞.
同著他有南丁八十.
就是每一個的領袖.
不同支派的.
名門的望族.
一些大的姓.
他們帶著多少人上去.
帶著多少人上去.
有些是一百五十.
有些是一百.
有些是二百.
有些是八十.
差不多是很多的.
整個家族上去的.
這是第二批上去的.
就是以色列做的領袖.
帶上去的.

$^{1281}$第九節.
屬約阿的子孫.
有耶歇的兒子.
俄巴底亞.
同著他有南丁二百十八.
屬士羅密的子孫.
有約西肥的兒子.
同著他有南丁一百六十.
屬比拜的子孫.
有比拜的兒子.
撒迦利亞.
同著他有南丁二十八.
屬厄甲的子孫.
有哈加坦的兒子.
約哈蘭.
同著他有南丁一百一十.
屬阿多尼干的子孫.
就是末尾的.
他們的名字是.
以利法列 耶利 斯瑪雅.
同著他們有南丁六十.
屬比格瓦伊的子孫.
有烏大和薩布.
同著他們有南丁七十.
你稍為計算一下.
就算到多少人了.
我召聚這些人.
再流入阿哈瓦河邊.
我們在那裡住了三日.
講得很細緻.
那個旅行日記.
由初一.
他們怎樣上去.
首先就是查了家譜.
一來就是計算.
他們的家譜.
一來就是計算他們有多少人口.
同著就是來自哪個支派.
哪個家族.
有多少人.

$^{1321}$因為這個很關鍵.
上到去之後人家就會幫你查.
你自己不查.
上到去人家幫你查.
你如果不是祭司.
你連點什麼器命都不准.
你根本就不可以進去聖殿.
所以其實很緊要.
所以我們信主.
初信班.
我們是要問心事.
你是否清清楚楚得救.
你才可以去浸禮班.
要有牧者的推薦.
就是說對你有認識.
然後才可以上浸禮班.
完了就可以去浸禮.
我們都要查心事.
就是你信什麼.
如果你根本都不知道自己信什麼.
那你為什麼還要浸禮.
不是那些水特別神聖.
反而就是你根本就不得救.
你回不到上天.
那更大件事.
所以不關那些水事.
是講你的內心.
有沒有這個.
屬靈的那個位份.
那個身份.
那個救恩已經在你裡面發生.
所以這裡就是他們要查他們的家譜.
所以不是不要緊.
很重要.
然後就是第三天.
對不起.
第二節都有東西講.
十五節這裡.
在那裡住了三天.
我查看百姓和祭司.

$^{1361}$見沒有你們在那裡.
就召首領.
以利爾謝亞烈.
是馬雅.
就是那些人裡面.
其中比較重要的領袖.
是馬雅 以利拿丹.
雅納 以利拿丹.
拿丹 撒迦利亞 米舒蘭.
又召交集.
約雅納和以利拿丹.
打發他們.
到加西肥亞的地方.
去見那裡的首領.
亦多 又告訴他們.
當向亦多.
和他的弟兄來提靈.
說什麼話 叫他們為我們.
神的殿.
帶使用的人來.
就是找一些.
有利美茲派背景的.
又是.
裝備好的.
這樣一齊上去.
神的那個殿.
由巴比倫.
上去耶路撒冷.
來重見.
第十八節.
第四節出現了.
蒙我們神斯恩的手幫助我們.
他們在以色列的曾孫.
利美的孫子.
茉莉的後裔當中.
帶一個通達的人來.
還有士利比.
和他的眾子與弟兄.
共十八人.
又有哈沙比亞.

$^{1401}$還有米拉尼.
這些是利美派的.
子孫耶西亞.
並他的眾子與弟兄.
共二十人.
從前大衛和眾首領.
派來提靈.
服侍利未人.
現在這來提靈中.
也帶了二百二十人來.
都是按命指定的.
來提靈.
那些本來是外邦人.
是欺騙人.
因為他們在約書亞的年間.
曾經欺騙約書亞.
約書亞就.
跟他們納了約.
不會將他們滅絕.
他們願意.
成為以色列人的.
奴隸.
來服侍.
所以那些欺騙人.
後來就成為了來提靈.
就是.
聖殿裡面或者會幕裡面.
的那些.
習翼,習公.
所以現在這些來提靈.
就是後代.
都繼續一起回去重建.
所以很難得.
他們是做奴隸的.
做奴隸都可以回去重建.
那就是後面.
第二十一節.
那時我在阿哈瓦河邊.
宣告禁食.
為要在.

$^{1441}$我們神面前.
刻苦己心.
求祂使我們和婦人孩子.
並一切所有的.
都得到平坦的道路.
我求.
王撥步兵馬兵.
幫助我們.
抵擋路上的仇敵.
本以為羞恥.
因我曾對王說.
我們神.
施恩的手.
必幫助一切尋求.
他的但他的能力.
和憤怒.
必攻擊一切離棄.
他的.
他反而有心有感.
而發.
今次事關重大.
他都不夠信心.
所以他都求王.
給一些步兵.
和沿途有護送.
有些.
保鏢.
上去.
所以他就覺得有點羞恥.
因為他曾經對王說.
我的神.
他的施恩的手.
必幫助我們一切尋求.
他的.
但現在你又說完.
神幫助你.
你又要求軍隊.
即是你兩手準備.
他就覺得羞恥.
有這個羞恥感都好.

$^{1481}$即是說.
是的,我都信的,不過.
都要有兩手準備好.
這就是第五次說.
神施恩的手幫助他.
所以這一章聖經.
第七章.
我們讀到這裡已經有五次.
廿三節裡面有五次說.
神施恩的手幫助.
由一個.
好通達智慧的.
民事.
祭司說的.
由一個在王的面前.
幾位居要職的一個.
可能是文官.
來說的.
可以想像到.
他不是沒有世界的智慧.
但他更加有對神的敬畏.
和愛神的心.
所以他才看到.
是神的手.
管理著他的想法.
那個部署.
那個行程.
那個計劃.
五個月你要帶隊.
上去耶路撒冷.
都不是小事.
現在這麼方便.
所以就是第廿三節.
所以我們禁食.
祈求我們的神.
他就應運了我們.
所以在整個過程.
祈禱是很重要.
所以盼望我們越來越發現.
無論我們多現代.

$^{1521}$多科學.
多有學問智慧.
祈禱從來就不是一種.
好像訴諸於.
我們什麼都不懂.
很無知.
是那些迷信的婆婆.
才會祈禱.
其實不是.
越是有知識.
越是發現.
你很多東西掌控不到.
更加要依靠神.
不是要抹殺我們的腦袋.
我們的智慧.
而是有時.
特別是神的工作.
不是腦袋不是智慧.
是可以改變的.
尤其是人心.
教會的合一.
每個人都有自己的想法.
因為每個人都認為自己是智慧.
更加能夠有一個同心.
必須是要.
大家都回到神的面前.
唯一的就是祈禱.
不是投票.
是祈禱.
廿四節.
我分派祭司十二個人.
就是史利比.
和他們的弟兄十個人.
張王和牟氏.
軍長並在那裡的以色列眾人.
為我們神殿所獻的.
金銀和器皿.
都過了稱.
要先稱一稱.
因為那些是度量行.

$^{1561}$金銀.
不是紙幣.
紙幣你都要數一數.
一萬元.
要點數要記賬.
其實是記賬.
沿途不知道發生什麼.
和運送過程.
會不會有些人.
來到去落格.
你又要稱一稱.
這些都是智慧.
上到上面之後耶路撒冷再重新過稱.
有沒有缺失了.
有沒有一些遺失了.
等等.
過稱了.
第二十六節.
我稱了交在他們手中的銀子.
有六百五十他連得.
很厲害.
一他連得是九十斤.
很厲害.
一車一車的珠寶.
用雷仔這樣托上去.
當然現在順風.
當然方便.
但那時候沒有.
也很不安全.
所以要這麼久時間.
銀器重一百他連得.
金是一百他連得.
金碗二十個.
上等的光銅器銘.
兩個寶貴如金.
我對他們說.
你們歸耶和華為聖.
器銘也為聖.
金銀是甘心獻給.
耶和華你們列祖之神.

$^{1601}$是弟兄姊妹的奉獻.
是王的王父.
裡面那個父父裡面拿出來.
都是甘心樂意獻給神.
所以.
數點獻金的我們需要有.
這種心靈的.
聖潔.
因為是甘心獻的.
是歸神的.
所以要很小心.
所以第29節你們當警醒看守.
所以我們數點奉獻.
至少有三個.
起碼三個.
最好是四個甚至乎要五個.
也要有童工看守.
這樣來數點.
因為需要.
很警醒.
直到你們在耶路撒冷.
或神殿的庫內.
在祭司長當地.
和利美的族長.
以色列的族長.
面前過程.
過了程.
即是說我們帶了這麼多上來.
大家都是很清楚的.
點收得很清楚.
上到去之後交收.
我們全部.
是為神殿奉獻的.
這個很重要.
真是.
一個智慧的安排.
這是以色列對這件事.
這麼嚴肅的看待.
不會這麼.
隨隨便便去處理這些事情.

$^{1641}$偶然間.
有些教會.
有時一兩個人.
數完就算了.
有時有點危險.
將這兩個人全部在試探裡.
如果兩個人合起來.
大家分了就很麻煩.
難免會有人.
有些說病.
所以一定要有智慧.
要很嚴謹看這件事.
很嚴肅很警醒.
30節.
於是祭司利未人按照份量.
接受金銀和器皿.
要帶到耶路撒冷.
我們神的殿.
然後正月12日.
才啟程.
我們從阿哈瓦河邊.
起行.
要往耶路撒冷去.
守如節.
即是在路上守.
我們神的手保佑我們.
這是第六次說.
第31節.
我們到了耶路撒冷.
在那裡住了三天.
即是五個月後.
第四天.
在我們神的殿裡.
將金銀和器皿.
都來到去過稱.
交在祭司當地的.
烏尼亞的兒子.
史理末的手中.
同著他有菲尼哈的兒子.
以利亞薩.

$^{1681}$還有利未人耶穌亞的兒子.
約撒特和賓內的兒子.
羅亞底.
當時都點了數目.
按照份量.
寫在冊上.
很清清楚楚.
登記了誰人奉獻的.
諸如此類.
很浩大的工程.
因為數量都相當多.
這是祭司.
這個以斯拉.
他所做的.
即是他又是領袖又是文士.
又是祭司.
等等的官職.
在一身.
原來第九章.
他更加是傳道人.
是牧師.
開始他自己.
由他禱告.
變成一個嚴肅會.
這個嚴肅會就帶來.
這個民族的復興.
從老道之地歸回的人.
向以色列神獻梵祭.
就是為.
以色列眾人獻.
公牛十二隻.
公綿羊九十六隻.
綿羊羔七十七隻.
又獻公山羊十二隻.
作屬罪祭.
這都是向耶和華奉獻的.
梵祭就是整隻獻了.
梵祭是不可以吃那些祭肉的.
梵祭只能夠留.
那個祭身的皮.

$^{1721}$譬如羊皮牛皮.
那些是歸給祭司的.
其他整隻要燒在祭壇.
屬罪祭和平安祭.
那些祭肉就可以吃.
第三十六節.
他們將王的御子.
交給王所派總督.
和河西的省長.
他們就幫助百姓.
由供給神電裡.
所需用.
首先是去到神的電.
向祭司交代.
交代好.
數點.
過了稱.
記錄在冊子.
做完這些神聖的工作之後.
以色列才去找.
當地的地方官.
河西的.
省長.
將王的御子.
告訴他們.
叫他們撥一些.
貢銀.
幫助神電的工程.
的需用.
所以先是去到神那裡.
好像先來到去.
敬拜神.
然後才做一些地上的手續.
不是一去到.
就巴結當地的.
地方官.
然後才去處理.
聖電的事.
次序.
不是說對還是錯的問題.

$^{1761}$而是他真的對神.
有一個首先的敬畏.
對當地的祭司的尊重.
當地的祭司.
那個.
受人.
尊重的.
水平.
高不高呢.
如果不高,你首先就找祭司.
不是找地方官.
就證明以色列的人.
對神的愛神.
對神的僕人.
當地的耶路撒冷的.
服侍神的.
他有一份首先的尊重.
這個都很重要.
所以盼望弟兄姊妹.
有時我們.
有時都會有弟兄姊妹.
回來香港探親.
他可能一選.
就是禮拜一定要回教會.
先敬拜神.
現在是我的母會.
探望我的弟兄姊妹.
這個是優先.
很不同.
如果不是,你做完.
臨走才去探弟兄姊妹.
或者不方便時間.
都不回教會.
只約弟兄姊妹.
當然我不知道,因為你沒回來教會我怎會知道.
除非一些我認識.
次序不好.
應該就是我們首先去到神的面前.
一切的順利.
或者離開幾十年.

$^{1801}$現在好像.
回來你的母堂.
你的母會.
你的靈明的.
那個母家.
敬拜這位.
獨一的真神.
這個是很重要的一個.
不可以叫姿態.
是一個心靈的舉措.
其實很重要.
一個人迎祖歸宗.
尤其是屬靈的家.
其實很重要.
往後其他的事.
自然就會順暢.
因為神的恩手.
會一路帶領.
所以耶穌就示範了一個.
敬虔的人這樣想事.
很有智慧地要過程.
很公正.
處理.
其實很多.
很重要.
這裡價值連城.
但是他既擔心.
那些土匪.
或者有山洪.
或者有其他的意外.
因為要五個月時間.
但是他更加緊張的.
就是這些是為神.
奠獻的東西.
分毫都不可以有差錯.
要很謹慎.
這些事情.
這個也可以說是我們.
後世處理獻金的.
教會某程度的.

$^{1841}$藍本.
聖經有時候也有教導.
好像故事的形式.
教導我們怎樣處理.
神的獻金.
神聖的物件.
要怎樣處理.
這個也是我們有很大的啟發.
以色拉雖然是古代的人.
住前大概.
四百五十多年.
可能四百八十多.
的一個人.
距離耶穌出世降生.
還有四百多年.
這位的.
以色拉.
都是這麼通達,這麼有智慧.
這麼有那種屬靈的眼光.
這麼有那種.
屬靈的慰憤.
這樣去帶領著.
第二次.
歸回的歷史.
還有三個字.
第九章.
就是那個憤興會.
一直延伸到.
第十章.
兩章聖經.
就是這個.
憤興會.
所以他就看到.
這是作完了.
就是一個.
定下來了.
但是他在耶路撒冷.
停留多久呢.
兩個星期.
一個月.

$^{1881}$這個就沒有交代.
這兩章聖經就是一個.
憤興會.
還有.
眾首領來見我.
說:以色列民和.
祭司利未人.
沒有離絕迦南人.
黑人,比利西人,耶布斯人.
阿滿人,摩訶人.
埃及人,阿摩利人.
就是迦南七族.
仍效法.
這些國的民.
恆可爭的事.
這裡就出現了一個.
不是硬件的問題.
住在耶路撒冷的人.
慢慢慢慢.
習以為常.
就和那些本地的人.
迦南人.
信仰都看齊了.
不知道.
你有沒有.
衝動要買六合彩呢.
好像是過千萬.
有沒有衝動.
試一下買一張十元電腦票.
中了.
我不要,全給教會.
不知道有沒有人這樣想.
慢慢不知.
眼光都跟了.
一般香港人的想法.
這裡告訴我們.
是那些首領.
眾首領來找.
以色列.
告訴他們.

$^{1921}$那個情況.
因為他是一個觀光客.
他是鳳凰的.
遇子.
帶領.
那些聖殿的物資.
這些獻金.
帶上耶路撒冷.
安排好之後.
相信都有那個節期.
手元節.
接著.
眾領袖.
約他.
來講.
現在.
國民.
那個內心的情況.
第二節.
為自己和兒子.
娶了這些外邦女子.
為妻.
以至聖潔的種類.
和這些國的民運.
集.
而且首領和官長.
在這件事情上為聚灰.
首領和官長.
都在這件事情上.
帶頭去做.
第三節.
就是以色列的反應.
我一聽見.
其實以色列住在外邦.
住在波斯帝國.
他自己.
和他的.
親朋好友.
是第二第三代.
不是耶路撒冷的原居民.

$^{1961}$他聽到這件事.
怎麼辦呢.
他就撕裂衣服和外袍.
拔了頭髮和鬍鬚.
驚懼.
幽悶而坐.
很驚慌懼怕.
很憂心.
很悶悶不樂.
這樣來到他自己一個人坐.
凡為以色列神言語.
戰驚的.
都因這被老歸回的人.
所犯的罪.
罪集到我這裡來.
我就驚懼.
幽悶而坐.
直到獻晚祭的時候.
晚祭就是黃昏六點鐘.
那個獻給神的.
凡祭.
第五節.
獻晚祭的時候我起來.
心中受苦.
穿著撕裂了的衣袍.
雙膝跪下.
向耶和華我的神舉手.
跪下舉手.
說我的神.
這個就是他的祈禱.
第九章餘下的全部是他的祈禱.
一路去數.
我們得罪神.
由我們的列祖.
所羅門開始.
我們好像都沒有.
真正去有痛定思痛的.
那個反思.
和這些外邦人.
當地的迦南人.

$^{2001}$來去通婚.
他們有了後代都不一定.
因為都不是一時三刻.
的問題.
大家都知道.
錦宣家或者宣道會.
必須是兩個都是基督徒.
才可以結婚.
我們不會幫他們行婚禮.
我聽一些前輩說.
偶然間都有要求.
我的另一半.
未信.
未是基督徒,但他有心的.
信的,可不可以在教會行婚禮.
好像要我們破例.
這個區聯會定的,不是我們定的.
好像要我們破例.
算了,不要執著.
他都說自己是信的.
你不用他洗禮.
偶然間都會有這些聲音.
是不是我們太執著.
我們現在見到.
以色列.
他聽到,其實已經是第二第三代.
已經是和當地人.
要來去.
與少數民族融和.
合一.
如果不是你可能都和.
國家的政策.
有抵觸.
可能波斯帝國都想你不要.
這麼極端.
例如新加坡.
不可以你很極端針對.
回教,伊斯蘭教.
去詆毀,在報章.
去說別人不好.

$^{2041}$說那些神是假的,不可以的.
新加坡是要做一個.
比較懷柔的.
多元的宗教.
大家互相的尊重.
你不可以這樣.
因為有犯法.
雖然香港未是這麼嚴重.
但是都是.
文風.
不傾向你很偏激.
在報章針對.
什麼教,這樣都不好.
現在社會.
某程度的進步.
進化都不容許.
一個宗教獨大.
說自己是唯一的,都不准.
你可以教內.
但在教外.
你不可以這樣做.
所以這會不會是當時的政治背景.
令到他們.
慢慢慢慢.
有這種融和.
但是在以斯拉.
這個民事的內心.
他就看這件事看得很嚴重.
第六節.
我的神啊,我抱愧蒙羞.
不敢向我神.
讓面,因為.
我們的罪滅頂.
我們的罪惡滔天.
從我們列祖直到.
今天,所說的是所羅門.
我們的罪惡甚重.
因我們的罪孽.
我們和君王.
祭司都交在外邦.

$^{2081}$列王的手中,殺害.
擄掠搶奪,就是他們.
被擄的歷史.
臉上蒙羞,正如.
今天的光景,現在.
耶和華我們的神.
暫且施恩於我們.
讓我們留聲逃脫的人.
使我們安穩.
好像釘子一樣釘在他的聖所.
即是他們被擄歸回.
已經是第二批.
第二個時代.
我們的神.
很光照我們的眼目.
使我們在受.
合制之中稍微的復興.
我們是奴僕.
然而在受合制之中.
我們的神.
仍沒有丟棄我們.
在波斯王眼前.
向我們施恩,叫我們復興.
能重建我們神的殿.
修其毀壞之處.
使我們在猶大和耶路撒冷有長活.
我們的神.
既是如此.
我們還有什麼可以說呢.
因為我們已經離棄了你的命令.
是在說那個.
娶那些當地迦南的婦女為妻.
似乎又生了子女了.
第十一節.
就是你僕人眾先知所吩咐的.
說.
你們要去德遺孽的地.
是污穢之地.
因列國之民的污穢.
和可憎的事.

$^{2121}$叫全地從這邊.
直到那邊都滿了污穢.
所以不可.
將你們的女兒.
嫁給他們的兒子.
似乎好像很極端.
用我們今天的眼光.
看回舊時.
主前的四百八十多年的事.
好像很極端.
極端到就是.
任何猶太人都不可以.
取愛幫女子.
不是啊 路德.
這個摩訶的女子.
就嫁了阿波亞斯.
就有這些先例.
為什麼以色列會這麼奇怪.
你是第二代.
你是竹星仔.
你是波斯大.
為什麼你的思想這麼迂腐.
這麼落後 這麼古板.
這個有些奇怪.
特別的.
對當地的.
耶路撒冷的居民.
難道我去哪裡.
找猶太人的女子呢.
如果不夠.
那怎麼辦呢.
那我的兒子.
都不用娶老婆了嗎.
當地的迦南人.
有這麼好的女子.
是不是真的你所說的.
這地滿了污穢呢.
所以不可以將你們的女兒.
嫁給他們的兒子.
也不可以將你們的兒子.

$^{2161}$去娶他們的女兒.
永遠不可以求他們的平安.
和他們的利益.
那你又求皇.
給你有金銀.
互為兵.
好像雙重標準.
以色列.
好像有一個.
衝突的矛盾.
這樣你們就可以.
強盛吃這地的美物.
並遺留這地.
給你們子孫永遠為業.
一個這麼強烈的.
隔離主義.
存在在以色列的心底.
這是一個難題.
第十章.
領袖已決了一個決定.
下星期再說.
大家回去看看.
那個已決是什麼.
很嚴峻.
其實一點都不簡單.
不外乎.
要不就分.
要不就合.
問題是他肯不肯分.
要他們簽約.
要逼他們要求.
要求他們簽約.
要逼他們要離棄.
自己結了婚的妻子.
趕他們走.
嫁了的就要.
回頭.
娶了老婆的家男女子.
趕他走.
和他離婚.

$^{2201}$這樣要求.
和要立約.
這是以色列和領袖.
最終的決定.
所以這兩章聖經就有很大的爭議.
到底這樣做.
是不是最智慧呢.
奇就奇在.
就停了在立約之後.
又沒有說神.
是什麼心意.
所以這就成為了後代.
這段經文的爭論.
十三節.
神啊.
我們因自己的惡行.
和大罪遭遇這一切的事.
並且你.
刑罰我們輕於我們罪.
所當得.
又給我們留下這些人.
我們豈可再違背你的命令.
似乎要結論.
我們怎可以.
再違背呢.
意味著有行動.
與者行可憎之事.
的文傑稱.
還一路說人家嘉楠人.
沒得救了.
是可憎的了.
污穢的了.
我們怎可以再和他們.
來到一起.
結為恩親呢.
若者讓行.
你豈不向我們發怒.
將我們滅絕.
以至沒有一個剩下逃脫的人.
耶和華以色列的神啊.

$^{2241}$因你是公義.
我們這些剩下的人.
才得以逃脫.
正如今天的光景.
看啊.
在你面前有罪惡.
因此沒有人.
在你面前站立得住.
接著第十章.
以色列禱告認罪哭泣.
俯伏在神的殿.
的時候.
我們今天就停在這裡.
我們一起祈禱.
親愛的主.
我們由第七章看到第八章.
是很詳細講述.
以色列的歸回.
和他的內心.
和他過去在波斯帝國的裝備.
他的.
身心靈.
他的頭腦.
他的智慧.
他的學識的裝備.
更加是他對神的認識的裝備.
對神的律法.
不是一種.
主觀的一種感覺.
而是他很通達神的律法.
能夠查考.
能夠來到他很全面的掌握.
能夠隨時的去教導.
而且有一個很敏捷的.
思想頭腦.
更加他在外邦的君王的面前.
是得到認同.
連君王都知道.
他是通達天上神律法的.
文士.

$^{2281}$我們很盼望.
在這個時代.
求你興起.
這樣的牧者.
僕人領袖.
來到他真的對神的道.
是有心.
是有能力.
這樣去勸勉.
教導.
能夠對時代有一個準確的把脈.
不懼怕於時代的恐嚇.
而是.
真的對神.
有一個全面的掌握.
掌控.
是有這一份信心.
這份信心能夠.
灌注在那個群體.
讓弟兄姊妹.
無論是什麼背景.
什麼裝備.
都不缺欠對神的道.
那個追求.
我們很需要這個時代的年青人.
是翻主一學.
愛慕神的律法.
和多祈禱.
敬虔.
懇求主為錦宣家.
和其他的無論宣道會.
其他的宗派.
預備下一代的領導的.
牧者領袖.
真的有屬靈的眼光.
和智慧.
是真的敬虔愛主.
愛弟兄姊妹.
求主來幫助.
我們下個星期當看那個結果.

$^{2321}$第十章的時候.
讓我們能夠.
看到再去分析.
當時這樣的決定.
是不是最有智慧.
還是本來就是神期望.
以色列帶來的那個復興.
或者是不是這樣.
就是叫做復興.
求主讓我們下個星期.
都一齊的去學習.
我們這樣禱告奉耶穌基督的名.
阿們.
你.
\newpage



\section{馬可福音 10:13-16}
\label{sec:P577AmrjvF0}
\textbf{2023-02-12 加恩樓《讓小孩子「到我這裡來」》馬可福音 10\_13-16 講員:吳景霖傳道}
\newline
\newline
連結: \href{https://youtube.com/watch?v=P577AmrjvF0}{\texttt{https://youtube.com/watch?v=P577AmrjvF0}} ~~~~ 語音日期: 2023-02-14
\newline
\newline
\hyperref[sec:lrudeppALpc]{\small{< < < PREV SERMON < < <}}
~
\hyperref[sec:index_chronic]{\small{[返順時目]}}
~
\hyperref[sec:index_scriptual]{\small{[返順卷目]}}
~
\hyperref[sec:raKQe4zsFjc]{\small{> > > NEXT SERMON > > >}}
\newline
\newline
馬可福音 10:13-16
\newline
\begin{longtable}{cl}
\hline
\hline
章節 & 經文 (和合本修訂版)\\
\hline
10:13 & \begin{tabularx}{0.7\textwidth}{X} 有人帶著小孩子來見耶穌，要他摸他們，門徒就責備那些人。 \end{tabularx} \\ \\ \relax
10:14 & \begin{tabularx}{0.7\textwidth}{X} 耶穌看見就很生氣，對門徒說：「讓小孩到我這裡來，不要阻止他們，因為在神國的正是這樣的人。 \end{tabularx} \\ \\ \relax
10:15 & \begin{tabularx}{0.7\textwidth}{X} 我實在告訴你們，凡要接受神國的，若不像小孩子，絕不能進去。」 \end{tabularx} \\ \\ \relax
10:16 & \begin{tabularx}{0.7\textwidth}{X} 於是他抱著小孩子，給他們按手，為他們祝福。 \end{tabularx} \\ \\
[1ex]
\hline
\hline
\end{longtable}
$^{1}$喂喂喂,各位大英姐妹,早晨,主人平安.
很久沒有在大家當中了.
其實上一次,你有沒有發覺題目是鎖了屬面.
上一次我講的題目都是講養小孩子到我這裡來.
不過我的括號是括在小孩子那裡.
段經文也不同了位置.
其實上次也有跟大家講過.
三卷福音書其實也有講到這個片段.
而三卷福音書也是黏著財主的講論.
所以大家會在三卷福音書也看到.
如果心水清一招,我上次講是有些東西還沒講的.
為什麼要跟財主放在一起呢?.
我會重溫一下上次的講道.
所以天崇一等英姐妹,如果還沒聽過不要緊的.
我有前回的重溫,就像韓劇一樣,播放前面一點.
那在講的,怎樣像小孩子才可以進天國呢?.
傳統我們有一種理解,就是要學習小孩子的樣式.
那你學小孩子什麼呢?.
學他天真,學他單純.
但是天真和單純有沒有得學?.
沒有得學的,你只可以扮天真和扮單純.
很多人都會扮天真,扮不知道.
但是實際上,當你知道了一些東西之後,你是不能再扮的.
你已經是邪惡了,就是這麼邪惡了.
我不像大家這麼年輕就有機會回教會.
我19歲才相信耶穌.
所以之前就過著很邪惡的生活.
我很多壞事我都知道,都明白,都懂得.
所以當我要回頭,要很天真很單純.
我自問我是完全做不到的.
但是,這樣我們是不是可以繼續邪惡呢?.
也不是的,人是可以選擇善良的.
你可以選擇做一些善良的事情,你可以像小朋友一樣善良.
你可以像行為貞直,你可以做一些善行.
但是不是說你要反過來,像小孩子一樣單純.
你做不到的,你只有扮,扮不明白.
邪惡的事情,我聽不明白,不知道為什麼.
但其實現實就是你可以明白,不過你不去做.
是這樣的一種情況.
所以說,其實我們學不到小孩子,我們無法變回頭.

$^{41}$既然是這樣,我們怎樣了解.
耶穌要我們像小孩子一樣才可以進天國呢?.
在當時的背景,小孩子其實是沒有價值的.
不像現在,可能你只是九代單傳,只有一個,然後就七星半月.
記不記得什麼是七星半月?.
爸爸媽媽,爺爺奶奶,公婆婆還有一個工人.
圍著一個來轉,那些叫七星半月.
很寶貴的,但當時來說,小朋友的夭折率幾乎是50\%.
就是說你生得越多,下面能夠繁衍下去的機會率就越高.
所以是量產的.
於是小孩子其實是被輕看的.
他當時是沒有權利的,也是父母親的財產.
小孩子是大人的附屬品.
同時我們看回經文.
你會發覺那些人帶小孩子來到耶穌的前面.
是不受歡迎的.
那些門徒是怎樣的?.
是責備他們,說不要妨礙他們.
大人正在做正經事,你過來幹嘛?.
摸什麼摸,妨礙耶穌講道.
小孩子不要過來,於是驅逐了他們.
耶穌講道這麼正經,這麼重要,這麼神聖的時刻.
當然是不要被小孩子妨礙了,他們這樣的想法.
小孩子是不受歡迎的.
我們看回,不是說這一段.
不是說要學長小孩子,天真無邪.
我們要看回小孩子的特性.
他不受歡迎,他是被拒絕的.
他被趕離耶穌身邊,離開.
小孩子象徵著一個被拒絕的群體.
他是軟弱無力.
如果我們用這個看法的話.
進天國的人是怎樣的?.
就是軟弱無力,就是被拒絕的.
被這個世界所拒絕的一群,偏偏他們就在天國上.
如果我們放在一起,跟年輕財主一起比較.
那個特性又會再明顯一點.
一個年輕的財主.
他作為一個社會上的上流人士.
就是社會的上層.

$^{81}$他做了很多善事.
可能他身光頸靚,或者他做的事情是人盡皆知.
那些大善人,如果有些考古你會發現.
遺跡裡面,那個人捐錢是有個牌子的.
不是現在才這樣,什麼總理.
古代就已經是這樣的.
他捐錢建這條路,有個石碑在那裡.
某某人他捐錢.
可能他是一個人盡皆知的善長仁翁.
他在上流裡面有錢有地位.
又年輕,如果帥的話就無敵了.
他信仰上的規矩可能是毫無瑕疵的.
因為他夠膽在耶穌面前說.
我所有的東西都守好了.
什麼規矩我都做足了.
其實都相當厲害.
如果現在問大家聖經裡面的事情.
你有沒有守足.
十誡你是不是都守好了.
有些人可能都不敢回答.
在這樣的情況下,他以為靠著自己的能力.
我就可以上天國.
他以為耶穌都客客氣氣的.
你什麼都做了,你就可以上天堂了.
誰知道耶穌跟他說什麼.
你要放下一切,變賣所有的離根重我.
於是他就悠悠愁愁地走了.
真的搞不定,買了所有東西.
剛剛才買了Tesla,你叫我放了它.
不行的.
有很多東西黏著他.
令他不能夠放下根重耶穌.
財主和小孩子一放在一起.
你就會發覺一幅天國的圖畫.
世上一些軟弱的群體.
被拒絕的群體.
例如基層的人,很貧困的人.
或者是一些SEN的小朋友.
或者是一些智力有問題的朋友.
一些外傭或者一些露宿者.

$^{121}$這些他們竟然在天國上面.
是有一席位的.
是進得去的.
但是一些可能在社會上的上流人士.
你會發覺他悠悠愁愁地離開耶穌.
因為他的產業很多.
這裡做了一幅天國的圖畫.
但要記住.
經常有人說吳傳徒很仇富.
經常都說這些,我不是仇富.
我先戴上頭盔.
因為從來我不認為有錢是有問題的.
聖經有說過.
財富是上帝對人的祝福.
上帝祝福你可以透過錢財.
所以有錢完全沒有問題.
但是這段經文其實是反對.
如果你以為你的財富,你的能力,你的社會地位.
或者你日常做的好行為.
是可以幫你進入天國的.
是一個進入天國的條件.
那就錯了.
那就錯了,聽得明白了.
不是說你擁有一些東西有問題.
而是說你以為你擁有的東西.
是幫你上天堂的.
那就有問題.
我們作為基督徒絕對不應該武斷地.
以一個人的身份,地位或者他的情況.
去判斷他的屬靈光景.
以為那個人更加接近上帝.
你不要想著一個人斯斯文文穿著西裝進來教會.
他一定是來追求信仰.
他可以是來追女孩.
或者是來認識人.
當這裡是一個信仰的俱樂部.
即是大家都是身光頸靚.
斯斯文文回教會.
應該都是好人去認識一下.
可能建立一下人際關係.

$^{161}$正所謂識人好過識字.
可能是這樣的心態.
相反是有些人他看起來很粗獷的.
整個一看就好像街邊流氓的那種.
但他進來教會.
他真的靜靜地坐著聽一下.
到底耶穌是什麼.
信仰是什麼.
其實很多人.
你看外表你是完全不知道的.
我記得我在錦繡的時候.
我曾經幫一個錦繡的工友.
他上初信班.
一個工友他真的很多東西都不懂.
看字都很慢.
讀得很慢.
聖經他都不是很看得懂.
但是他的信心很單純.
他的信仰很單純.
他回來教會就是想追求上帝.
所以你是不知道.
我們做基督徒不能夠.
按那個人的外表.
或者一些他的行為.
去判斷他的屬靈光景.
經文裡面.
那些門徒其實就是.
為耶穌下了一個閘.
他為耶穌做了篩選.
他們覺得去到耶穌面前的人.
是要經過他們的篩選的.
有些有價值的人才去到耶穌的前面.
那些帶著小孩來的.
聽又聽不明白又吵吵鬧鬧.
走去耶穌前面.
阻礙著耶穌講道.
叫耶穌摸摸他們.
這些是干擾.
於是斥責他們離開.
不要阻礙著.

$^{201}$跟觀那麼久不知道觀性是什麼.
原來耶穌很喜歡小朋友.
也很喜歡他們.
很想他們去到耶穌面前.
我記得我.
剛剛回教會的時候.
我第一位的牧者.
我們叫他陳生.
那時候.
男傳道先生女傳道姑娘.
叫他陳生.
他講道的時候.
他兒子很小我記得大概三歲.
兩三歲.
他講道的時候突然聲音大大.
衝進禮堂.
衝到講台那裡.
抱著陳生的腳.
陳生怎麼做.
難道是甩開的.
他兒子走過來.
那時候講台不是透明的.
是木台.
跟著跟他說.
然後就繼續講道.
兒子就在下面抱著.
弟兄姊妹就照聽.
這個印象很深刻.
我忘記他講道.
講什麼.
我相信大家很多時候都忘記別人講道講什麼.
餵養.
吃了就不記得.
這個OK的.
我忘記他講道講什麼.
但是我記得這個片段.
我覺得好像小朋友.
走向耶穌一樣.
他沒有被拒絕.
沒有被接納.

$^{241}$也沒有說正經事.
兒子你不要搞著.
人家跟我拖他出去去嬰兒房.
我們現在某程度是.
那個小朋友在那裡.
人家跟我拖他出去護幼那裡.
青年人跟我拖他出去.
不要在成崇那裡.
反過來有些上了大專.
跟我拖他出去成崇.
不要留在天崇.
不應該是這樣的.
大家明白是一種接納.
當然我要補充.
我不是說教會是一個沒有秩序的群體.
教會要有秩序的.
正常的.
如果有些東西干擾了秩序.
譬如我兒子衝上台.
我都一定會跟他說.
會提醒他這樣是不禮貌的.
下次不要這樣做.
或者以為總之打開門.
什麼人進來我們都會給他.
有一些可能是帶著不良動機來到教會.
有一些可能是專門.
認識人去推銷.
去完一間就一間.
有些是來教會借錢的.
教會老鼠.
專門來這間借完就去第二間.
也有一些可能是.
不良動機.
來洗底也好.
是需要拒絕的.
是需要留心的.
但實際怎麼操作.
要有智慧.
這些是特別的個案.
聽到這裡的時候.

$^{281}$我意識到一個問題.
到底說來說去.
什麼叫道我者裡來.
怎樣才算道耶穌前面.
是不是參與崇拜就叫道耶穌前面.
是不是.
你參與崇拜是不是道耶穌前面.
有的.
當然.
如果你來崇拜不是道耶穌前面.
我也不知道你到了什麼前面.
是的.
我們有很多不同崇拜.
青少年有青少年的崇拜.
成人有成人的.
兒童有主日學.
我們是不是已經做到要求.
什麼群體我們都帶到.
耶穌的面前.
是去到耶穌面前.
但我們也要.
記得一件事.
耶穌是無所不在的.
耶穌不是被局限在一個建築物裡面.
這裡耶穌就在.
走出樓梯.
耶穌就不在.
也不是局限在一個場景.
星期日上午.
11點半到12點半.
耶穌就在.
然後他就要去另一間教會輪班.
不是這樣的.
耶穌是無處不在.
也不被場景所限制.
所以教會以外的人.
他有沒有機會去到耶穌面前.
有機會的.
是不是.
他是需要去到耶穌面前.

$^{321}$所以我們怎樣令到一些教會以外.
沒有回來的群體去到耶穌面前.
很直觀的.
第一樣就是傳福音.
剛剛我們說到.
小孩子.
是進入天國的條件.
所以就要去傳福音.
也是對的.
但我再問深一點.
什麼叫傳福音.
怎樣傳福音給人.
叫他口裡承認.
相信就可以.
有沒有上朕禮班.
口裡承認心裡相信.
這個應該說得出.
不說的就找AKER補課.
這樣的情況.
是不是他舉手決志.
我相信耶穌.
就叫他可以上天國.
就叫他已經去到耶穌基督的面前.
剛才說的.
經文的小孩子去到耶穌面前.
耶穌有沒有叫他舉手決志.
沒有.
但他去到耶穌面前.
耶穌只是按手摸摸.
為他們祝福.
如果很多弟兄姊妹.
跟了我一段時間.
有聽過我講道.
不多不少都知道我想的方向.
我從來都不認為.
舉手決志.
信耶穌就算完事.
其實就.
感覺就好像.
某年某月某日某分.

$^{361}$某秒你就已經拿到.
天國入場券.
這就像法庭式.
總之那一刻.
就決定了.
我上天堂還是下地獄.
我有決志.
某年某月某日某分某秒.
我決志我上到天堂.
否則就要下地獄.
我從來都不是認為.
一種這樣的信仰.
就是叫做信耶穌.
相反我認為信仰是一個.
時間的.
很考驗你在這個時間.
有沒有繼續停留.
在這個信仰裡面.
有很多人舉手決志.
信耶穌決完之後就走了.
有些決志無限次.
每次報道會都去坐坐.
每次決志都去.
然後第二次又去.
他都不確定自己.
基督徒的身份.
平時又不回教會.
很多人是這樣的.
我們如果真的.
要讓一個人.
有上天國的把握.
其實你需要讓他停留在.
這個信仰裡面.
他保持著一個屬靈的進步.
光景.
他能夠和上帝保持著一種關係.
這樣才叫做.
真的成為門徒.
你記住大使命.
叫人成為門徒.

$^{401}$不是叫人去舉手決志.
是叫人.
要將耶穌所.
傳的教訓.
他們遵守.
這是成為門徒的.
心態.
讓人成為門徒的心態.
今日的經文不是處理.
一個這麼大的問題.
但我覺得基於這個理念.
基於這個理念.
我們要明白我們帶人去到.
耶穌面前,讓小孩子去到.
耶穌面前不是一個點.
不是說好,我現在要帶.
這個群體信耶穌.
你信不信耶穌.
現在舉手決志,決志,好,下一個.
不是一種流水作業的生產.
我們人生每一個階段.
其實都是一個.
尋求上帝的時刻.
你每一個階段.
都是去接觸上帝的時刻.
有機會去接觸上帝.
我用一個肉麻一點的例子.
大家說完覺得冷的話.
就穿上衣服.
我和我老婆剛剛去拍拖的時候.
她有說過.
你不要想著追到手就不用對我好.
結了婚那些.
不要想著結了婚就不用對我好.
我沒有想過.
我覺得拍拖是一輩子的事.
我老婆我要追一輩子.
我剛剛早兩天才買了一個公仔給我老婆.
做情人節禮物.
所以結了婚對老婆好.

$^{441}$笑就結了婚很久那些.
對老婆好一點.
保持住.
不要想著結了婚就不用對我好.
保持住.
不是說婚姻.
不是說那一刻.
簽紙OK搞定.
你不好一點說.
你做得不好.
隨時會離婚.
不是那一刻搞定舉手決戰就結束.
你做得不好.
你會離開信仰.
同樣的例子.
你要保持一種關係和上帝.
你那個淑齡光景.
無論你去到什麼年紀.
要對上帝有一種追求.
你不要想著我都信了二十幾三十年.
你還叫我.
你回來港島那些.
回來福去就是那幾段.
你傳遞人說的那些我懂得說.
我還用什麼回來.
或者靈修祈禱讀經那些.
做到悶了就不用做了.
不是的你要保持住和上帝有一個關係.
是不是.
好了.
回到剛才的問題.
怎樣去到耶穌面前.
叫做.
是不是翻崇拜.
有祈禱有靈修有讀經那些.
叫做去到上帝的面前.
我們看回經文那裡.
耶穌怎樣對他的小孩子.
他說於是他抱著小孩子.
給他們按手為他們祝福.

$^{481}$這句書文.
專門留給我讀的.
他抱著小孩子.
為他們按手為他們祝福.
這是一個很單純.
一種愛的行動.
一種關係性的行動.
一種關係性的接觸.
大家都.
熟我都知道.
我兒子都長大了都十歲了.
我自問我們的關係真的建立得很好.
他.
常常都走過來黏著我.
有時候黏身到.
有一個地步都覺得有點辛苦.
但是.
他怎樣都會走過來黏著我.
黏著我有時候是一些實際需要.
譬如肚餓.
過來找東西吃.
譬如有些東西想要.
譬如想看電視.
走過來叫我想看電視.
這些實際需要.
有時候可能是.
一些.
情緒上的需要.
譬如他很開心他又會喜歡黏著我.
他很不開心.
都會喜歡走過來坐在我大腿上.
跟我說死了明天默書.
有些心靈上的需要.
安慰一下又會走過來黏著我.
有時候純粹是想黏著我.
譬如我在房間寫講章的時候.
他就會拿著書.
坐在我房間裡面.
看書.
我有我做他有他做.

$^{521}$總之是同一個空間.
你趕他出聽他就不高興了.
他就喜歡黏著我.
有時候.
其實關係就是這樣.
我會回應他.
譬如他需要安慰就摸摸他的頭.
他需要我陪他.
我不會趕他走.
他在我身邊我容許他在我身邊.
已經是一種愛的回應.
其實去到耶穌面前.
我覺得是同樣的一件事.
就是一種很簡單的.
父子之間的交流.
上帝和你之間.
有一種交流.
你不用每次都是我有事祈禱我才去祈禱.
找神.
你沒有事祈禱就去祈禱.
你不是有事找你老婆你才跟她說話.
你平時無聊都會說話.
或者在他身邊.
停留一下已經是一種.
交流.
我太太有時候她說.
大家都很忙的時候.
她說最期待就是晚上十點多.
兒子睡了我們在看韓劇.
看韓劇都沒話說.
看著電視.
大家都是同在.
享受那種同在.
上帝享受與你的同在耶穌享受與你的同在.
如果他不.
享受與你的同在.
他為什麼要道成肉身下來.
以前很爽.
上帝高高在上.
人們敬拜他要去聖殿.

$^{561}$燒一隻牛給他.
如果不小心聽到上帝的聲音會死.
見到上帝的面是會死.
這種高高在上.
非常有威嚴的情況.
他不要.
他下來跟人一起住.
一起吃一起睡.
一起哭一起笑甚至被人釘死.
如果他不願意.
與人同在他為什麼要道成肉身下來.
耶穌基督就是.
願意與人同在.
他願意下來當中.
去接觸我們.
他很愛小孩子.
他愛每一個人甚至愛那些.
不被世界所接納.
他愛那些.
每一個社會上的基層.
有困難的人.
他有特殊需要的小朋友.
特殊需要的人.
他都愛這些人.
我們這樣看的時候.
我們要做的.
就是讓這些人去接近耶穌.
我們不要像門徒那樣.
下閘.
你們這些阻礙著的.
你們這些干擾崇拜的.
你們這些聽不懂堂道的.
聽不懂信息的.
你坐在這裡做什麼.
我們有時候是這樣去想.
但原來.
其實他在當中已經可以享受崇拜.
他可以親近.
透過崇拜.
透過這個場景.

$^{601}$透過這些弟兄姊妹.
他可以親近上帝.
可以接觸上帝.
你不要幫他判斷了.
有時候我看到很有趣的場景.
有些小朋友在我們當中.
我們真的會驅逐他.
聽不明白你去那裡參加兒崇.
但是為什麼你可以判斷他聽不明白.
或者他聽不明白又怎樣.
他聽不明白.
他不可以感受那種氣氛嗎.
他不可以在這種屬靈的氣氛裡面感受上帝嗎.
是不是啊.
有時候氣氛的事情很難說.
但是小朋友在你旁邊.
坐在這裡向你學習.
你這樣敬拜他.
他將來就會這樣敬拜.
你這樣守聖餐.
他就這樣守聖餐.
你的信心是怎樣他就怎樣.
但是現在某程度上.
我們的年輕人或小朋友.
是看不到我們是怎樣的.
看不到我們是怎樣.
他們自己的群體是這樣.
他就只看到自己的水平.
只看到自己的水平.
看不到他的父母或長輩.
早信耶穌的人.
他是怎樣信耶穌的.
是不可以的.
所以我覺得要讓人去到.
耶穌面前.
我們就要有一種去接納.
這些群體的心態.
我們要接納每一個.
他真是有心來到耶穌面前的人.
回來回覆他.

$^{641}$又是回覆他.
怎樣才能去到耶穌面前.
怎樣應用呢.
我分為作為信徒.
我們怎樣應用.
或作為一些未信者怎樣應用.
作為信徒.
我們要去到提醒自己.
其實我們經常都需要耶穌.
我們要經常去到耶穌面前.
我們經常都處於一種.
其實我們都是被耶穌所接納的群體.
我們某程度都是.
小孩子.
我們每一次去到崇拜.
或者每一次去到靈修.
查經主日學這些.
我們都是去到耶穌面前.
我們承認自己那種無知不足.
我們要謙卑.
去到耶穌面前.
去到學習去去支取力量.
或者是我們最基本的.
就是和他親近.
不用很功利.
靠近耶穌面前.
讓耶穌摸摸頭.
這件事已經很重要.
但有些情況就是.
我們會拒絕很多東西.
例如.
譬如讀神學.
讀神學也是一種.
去到耶穌面前學習的東西.
但很多時候.
我們給自己很多藉口.
不用學了,都信了這麼久.
讀神學是傳道人才要讀的.
我們不用讀的.
普通信徒不用的.

$^{681}$有時看聖經就足夠了.
不是的.
甚至有些教會阻礙人去讀神學.
知識叫人自高自大.
你班弟兄姊妹.
讀完之後比我厲害,那怎麼辦.
有些是這樣的.
其實是阻礙人去耶穌面前.
甚至有些會覺得.
你出去讀書.
問過我沒有.
我去判斷你的屬靈光景.
你不適合,你不要去讀.
阻礙人去耶穌面前.
或者是.
有些SEN的小朋友.
上一堂課到我有講.
或者有些特殊需要的小朋友.
其實很多時候我們會視他為.
整個程序的干擾.
或者是事實.
他如果在同一班小朋友裡面.
他需要特別多的支援.
才能夠相處.
甚至特別多的支援最終都相處不到.
教會放這麼多資源.
去幫一個小朋友.
我不如去幫一些容易處理的.
很容易.
很功利地去想.
資源分配.
但其實他們是需要.
人幫助,他們都想去到.
耶穌面前.
我們作為信徒要幫他開一條路.
讓他可以.
去到耶穌面前,他可以加入.
我們這個群體.
有些基層的朋友.
對他們來說,跨區去到.

$^{721}$活動是很昂貴的一件事.
所以大部分的基層.
只困在自己的社區裡面.
他們的社區有教會.
或者教會接納他們,他們就有得回.
如果教會不接納他們.
他們就沒有得回.
香港應該沒有什麼地方沒有教會.
但問題是接不接納.
所以在之前.
上個月.
中華基督教播道會.
講福堂的林誠信牧師.
來過錦繡港島,我不知道他有沒有.
過來這邊講,但他說得很好.
他講了一件事,就是他們原來.
他們坐落在金鐘.
在港福堂.
金鐘哪裡有窮人.
金鐘是商區.
他找不到一些.
需要幫助的群體.
他們就聯絡了.
粉嶺和屯門的教會.
去到他們那裡.
去做一些基層的事工.
好像我剛才說的,耶穌不是困在建築物.
不是出去了.
耶穌不在.
他不肯來堂食,我們就去外賣.
我們就去到那些人.
當中,去將他們.
帶到耶穌面前,讓他們有.
機會接觸耶穌.
這些是很重要的.
不要以為我們限制了.
一定要我們的對象.
來到我們當中,其實我們出去.
都一樣.
還有我們做一些基層.

$^{761}$工作的時候,我們又很容易很功利地問.
做這件事能不能傳到福音.
是不是.
傳不傳到福音,經常被人問.
有些弟兄姊妹跟過我.
我出去油麻地派飯.
也會這樣問.
派了這麼久.
傳不傳到福音,有沒有人信耶穌.
有人信又怎樣.
沒有人信又怎樣.
是不是他某年某月.
某時某分某秒,他決了自信耶穌.
他就可以拿到天堂入場券.
上天堂.
我剛才已經推翻了這一點.
所以如果.
很多人我們都是帶著這個想法.
所以我們才會這樣問.
他派了這麼久,他有沒有決志.
你很關心他信不信到天堂.
是不是.
但不是,其實是一種細水長流的關係式.
所以.
當我們去幫他的時候.
讓他知道原來基督徒會幫人.
其實已經很足夠.
他已經是親近了耶穌一步.
他已經對我們的信仰.
多了一分理解.
派飯盒的時候還拿著邊.
你信不信耶穌.
聽過耶穌是誰了嗎,你可不可以聽我說一次完整的福音.
跟他爭飯盒.
不用這樣.
你跟他好好交代.
他知道基督徒原來.
是幫人的,已經很好了.
你已經讓他接近了耶穌一步.
是不是.

$^{801}$譬如我們剛才說土耳其地震.
我們會捐款.
很多弟兄姊妹都很熱心.
就開始想要寫張卡.
寫張卡告訴他.
耶穌愛他,然後寫完整的福音.
我不知道他寫什麼文字.
但土耳其人應該看不懂中文.
總之.
一定要對方知道.
福音.
這樣我才做的事有意義.
不是的,你單純去幫助別人.
去讓人知道.
原來基督徒會關心人.
已經是幫了他們.
已經是讓他們認識耶穌多一點.
在一個漫長的生命裡.
他靠近了耶穌一步.
其實足夠.
不過我不是跟你說.
叫你不用去傳福音.
不是叫人決志.
我說的是不要太功利.
我們去到那個群體裡面.
服侍他.
我們去到那個群體裡面.
如果他有些人.
他經過了這麼多基督徒服侍之後.
他能夠接觸耶穌.
他認識耶穌更多.
由一個潛在的基督徒.
變成一個真正的基督徒.
你出那一分力就已經足夠.
你不要反過來跟他下閘.
你覺得那個人搞了這麼久.
都不相信耶穌.
算了,飯盒留給下一個.
吃了我一年飯.
都不相信耶穌,算了.

$^{841}$下一個.
很容易很功利地想.
不要跟門徒一樣.
像那些門徒下了閘.
你那班小孩都聽不懂耶穌說什麼.
給大人一點機會.
其實邏輯一樣.
有沒有發現?.
我們基督徒每一個行動.
只要能夠榮耀神.
只要讓人在你身上.
看見上帝.
只要讓一個未信的人.
看到你對基督徒.
或者對耶穌基督.
他增加了一些好感.
其實我覺得已經是讓他進入了信仰.
多一步.
慢慢他人生一步一步走.
終於會去到上帝面前.
祝福大家.
希望大家都懂得.
透過你的生活見證.
去養這些小孩子.
去到耶穌的面前.
我們同心圖告.
秋惠天父我們謝謝你.
謝謝你去接納我們這些小孩子.
讓我們成為你的家的一份子.
成為你的手手腳腳.
可以去幫你.
全部去成為你地上的代言.
求你幫助我們.
讓我們排除一些成見.
或者排除一些.
阻攔人們去信耶穌的東西.
相反我們學習.
怎樣開一條路.
讓你所愛的小孩子.
可以很順暢地進到你的面前.

$^{881}$讓你給了一些.
你願意的祝福.
和你的愛去給他們.
求你幫助我們祈禱交托.
奉耶穌基督的名求.
\newpage



\section{詩篇 139:1-24}
\label{sec:raKQe4zsFjc}
\textbf{2023-02-12 早堂《心知.神愛》詩篇 139\_1-24 講員:楊明麗姑娘}
\newline
\newline
連結: \href{https://youtube.com/watch?v=raKQe4zsFjc}{\texttt{https://youtube.com/watch?v=raKQe4zsFjc}} ~~~~ 語音日期: 2023-02-14
\newline
\newline
\hyperref[sec:P577AmrjvF0]{\small{< < < PREV SERMON < < <}}
~
\hyperref[sec:index_chronic]{\small{[返順時目]}}
~
\hyperref[sec:index_scriptual]{\small{[返順卷目]}}
~
\hyperref[sec:la_oxDk2vYY]{\small{> > > NEXT SERMON > > >}}
\newline
\newline
詩篇 139:1-24
\newline
\begin{longtable}{cl}
\hline
\hline
章節 & 經文 (和合本修訂版)\\
\hline
139:1 & \begin{tabularx}{0.7\textwidth}{X} 耶和華啊，你已經鑒察我，認識我。 \end{tabularx} \\ \\ \relax
139:2 & \begin{tabularx}{0.7\textwidth}{X} 我坐下，我起來，你都曉得；你從遠處知道我的意念。 \end{tabularx} \\ \\ \relax
139:3 & \begin{tabularx}{0.7\textwidth}{X} 我行路，我躺臥，你都細察；你也深知我一切所行的。 \end{tabularx} \\ \\ \relax
139:4 & \begin{tabularx}{0.7\textwidth}{X} 耶和華啊，我舌頭上的話，你沒有一句不知道的。 \end{tabularx} \\ \\ \relax
139:5 & \begin{tabularx}{0.7\textwidth}{X} 你前後環繞我，按手在我身上。 \end{tabularx} \\ \\ \relax
139:6 & \begin{tabularx}{0.7\textwidth}{X} 這樣的知識奇妙，是我不能測的；至高，是我不能及的。 \end{tabularx} \\ \\ \relax
139:7 & \begin{tabularx}{0.7\textwidth}{X} 我往哪裡去，躲避你的靈？我往哪裡逃，躲避你的面？ \end{tabularx} \\ \\ \relax
139:8 & \begin{tabularx}{0.7\textwidth}{X} 我若升到天上，你在那裡；我若躺在陰間，你也在那裡。 \end{tabularx} \\ \\ \relax
139:9 & \begin{tabularx}{0.7\textwidth}{X} 我若展開清晨的翅膀，飛到海極居住， \end{tabularx} \\ \\ \relax
139:10 & \begin{tabularx}{0.7\textwidth}{X} 就是在那裡，你的手必引導我，你的右手也必扶持我。 \end{tabularx} \\ \\ \relax
139:11 & \begin{tabularx}{0.7\textwidth}{X} 我若說「黑暗必定壓碎我，我周圍的亮光必成為黑夜」， \end{tabularx} \\ \\ \relax
139:12 & \begin{tabularx}{0.7\textwidth}{X} 黑暗對你不再是黑暗，黑夜卻如白晝發亮。黑暗和光明，在你看來都是一樣。 \end{tabularx} \\ \\ \relax
139:13 & \begin{tabularx}{0.7\textwidth}{X} 我的肺腑是你所造的，我在母腹中，你已編織我。 \end{tabularx} \\ \\ \relax
139:14 & \begin{tabularx}{0.7\textwidth}{X} 我要稱謝你，因我受造奇妙可畏，你的作為奇妙，這是我心深知道的。 \end{tabularx} \\ \\ \relax
139:15 & \begin{tabularx}{0.7\textwidth}{X} 我在暗中受造，在地的深處被塑造；那時，我的形體並不向你隱藏。 \end{tabularx} \\ \\ \relax
139:16 & \begin{tabularx}{0.7\textwidth}{X} 我未成形的體質，你的眼早已看見了；你所定的日子，我尚未度一日，都在你的冊子寫上了。 \end{tabularx} \\ \\ \relax
139:17 & \begin{tabularx}{0.7\textwidth}{X} 神啊，你的意念向我何等寶貴！其數何等眾多！ \end{tabularx} \\ \\ \relax
139:18 & \begin{tabularx}{0.7\textwidth}{X} 我若數點，比海沙更多；我睡醒的時候，仍和你同在。 \end{tabularx} \\ \\ \relax
139:19 & \begin{tabularx}{0.7\textwidth}{X} 神啊，惟願你殺戮惡人；你們好流人血的，離開我去吧！ \end{tabularx} \\ \\ \relax
139:20 & \begin{tabularx}{0.7\textwidth}{X} 他們說惡言頂撞你，你的仇敵妄稱你的名。 \end{tabularx} \\ \\ \relax
139:21 & \begin{tabularx}{0.7\textwidth}{X} 耶和華啊，恨惡你的，我豈不恨惡他們嗎？攻擊你的，我豈不憎惡他們嗎？ \end{tabularx} \\ \\ \relax
139:22 & \begin{tabularx}{0.7\textwidth}{X} 我恨惡他們到極點，以他們為我的仇敵。 \end{tabularx} \\ \\ \relax
139:23 & \begin{tabularx}{0.7\textwidth}{X} 神啊，求你鑒察我，知道我的心思，試煉我，知道我的意念； \end{tabularx} \\ \\ \relax
139:24 & \begin{tabularx}{0.7\textwidth}{X} 看在我裡面有甚麼惡行沒有，引導我走永生的道路。 \end{tabularx} \\ \\
[1ex]
\hline
\hline
\end{longtable}
$^{1}$各位姐妹 早安.
我們一起來聽神的話語的時候.
我們也一起禱告.
七天父上帶我們來到你的面前.
這是我們敬拜你的時候.
還有我們深信你在我們身邊.
還有我們更加求你將你的話語.
來教導我們.
將你的心意告訴我們.
最後讓我們每一個.
無論站在哪裡 坐在哪裡 說的 聽的.
都在你面前蒙恩典.
都能夠得著你的教導 你的指引.
我們仰望你.
奉主耶穌基督的聖名 阿們.
今天跟大家看一段經文.
這段經文是詩篇139篇.
相信大家是很熟悉的詩篇.
也很熟悉這首詩歌.
我們剛才也有唱.
其實這首詩歌寫出人的真相.
我們一起讀的時候.
我們看看這首詩歌.
這首詩篇是誰寫的呢.
這首詩篇是戴衛寫的.
寫作的時間應該是他比較年輕的時候.
可能是未做王的時候.
或者是初做王的時候.
也就是說他可能是在被掃羅追殺的年代.
這首詩篇分開四段.
總共有24節分開四段.
每一段有6節.
6節裡面 第5,6節就好像一個小節.
再到下一段的第5,6節.
就是第11,12節好像一個小節.
第17,18節又好像一個小節.
然後到第23,24節.
就回到第一節的內容.
第23節和第一節是很相似的.
很像一個總結.

$^{41}$一個開頭與結尾牽引得很好的一首詩.
當時應該是交給令長來編曲.
這首詩歌他跟我們說什麼呢.
一開始的時候.
他就說.
人家說.
你監察著我.
你認識我.
我坐下來起身.
你都看著我.
我走到哪裡 我躺下來.
你都在這裡 你都知道的.
我要說的話.
我還沒說一句 你都明白的.
可能是這個原因.
男士少說話.
男士的祈禱可能是天父你都知道的.
女士好像很長.
因為神都知道.
這首詩是男士作的.
是大衛寫的.
神啊 你就監察我.
監察的意思是.
不單只是看著.
是很留心的看著我.
試煉一下我.
盡量逐一說.
好像監察著我一樣.
你是這樣.
看得很實地.
看著我的神.
看到我一個地步.
我走到哪裡.
我躺下來起身 坐下來.
你什麼都知道 你什麼都看到.
我還要說一句話嗎.
我還沒說 你都知道我想說什麼了.
這個神.
你還像什麼呢.
你很保護我的.

$^{81}$你圍繞著我去保護我.
你不單只是像一個盾牌.
來幫我擋箭.
你好像一個帳幕.
包圍著我的身上.
這種知識很奇妙.
我是不能夠很明白.
弟兄姊妹 你的感覺.
大衛想說什麼呢.
你們很多都是為人父母的.
當子女有這種感受的時候.
我們是神的兒女.
當子女有這種感受的時候.
他在表達什麼呢.
媽呀.
早餐我自己會搞定.
天氣冷 我知道的.
我記得我在神學院的時候.
有一個同學是很好的.
不知道大家有沒有知道.
神學院有三層 三格床.
我讀海外神學院的.
我們神學院很出名的一件事.
地方夠小.
地方夠小 大家很親密的.
一間房最小的房間.
可以住六個人.
就是兩張三格床.
大家坐在下格床的時候.
膝蓋對膝蓋.
長腿的那些.
麻煩把腳縮到上面.
雙層.
我睡第三格 最頂的那一格.
中間下面那一格.
是不能夠上到那麼高層的.
那時候我覺得還挺健康的.
可以上高層一點.
中間那一格有一個姐妹.
是我師妹.

$^{121}$但是師妹她總是覺得.
我比她年輕.
她要關心照顧一下我.
我家人又不在香港.
於是我每天早上起來.
我都還沒有下到地.
我第三格還沒有下到地.
她就說.
明麗 多穿一件衣服.
我就很想告訴她.
我知道了.
我知道了 我每天早上都知道了.
夏天她也是這樣說的.
啊 你子女會不會有時候有這樣的感覺.
媽 我知道了.
總之 爸爸.
你記住.
iPhone出的時候我存在就行了.
其他時候.
我去黑暗的地方.
躺一下.
你不用知道 不用見到.
但是我好像在這個時候是想告訴天父.
是你 我知道你看緊我.
你看緊我到一個地步.
你不只是讓一個盾牌.
來幫我擋一下箭.
你簡直就好像一個帳幕.
這樣包圍著我.
你對我這麼好.
但是其實我不是很明白的.
我去到哪裡.
你都在樹上.
我躲不住的.
我躲在哪裡.
我想飛上天嗎.
你都在樹上.
我找一個很低層的地方躲下去嗎.
你不也是在那裡.
等我飛到天底這麼遠的時候.

$^{161}$現在有FaceTime.
去到美國 英國你都看到兒子.
啊.
他說其實.
我有時候想躲到地底.
但是在那裡.
你都好像在幫我鋪排著每一樣東西.
在那個情況下.
為什麼大衛有神看著他.
會有這種反應呢.
其實我們很多時候祈禱就是說.
神啊 你聽啊 你不要不聽啊.
你聽我們禱告啊.
我心裡面的事很重要的.
我為他祈禱的那個人.
他很需要你醫治的.
我們經常都很想說.
神在我們身邊 你聽著啊.
你不要太忙.
忘記了我的祈禱啊.
但是大衛他說.
你給我一點空間 天父.
有時候我們可能是這樣的.
天父我知道了.
恢復實體崇拜.
不過天父你給我一點空間.
都要讓教會給我一點空間.
等我由疫情.
去到後期要完結的時候.
我可以鬆動一下.
這幾年真的是.
我都不知道自己困在哪裡了.
疫情完結都給我一點空間.
天父 有時候這個環境很困難的.
我心裡面都有我的感受的.
我都要躲起來一下.
各位姐妹.
當弟兄他不太回答你.
你問他什麼 他就說.
有沒有事啊 沒有.

$^{201}$沒有表情的時候.
你給他一點空間.
無論他想去到海國.
或者想找一個深山.
你給他幾個小時 一晚.
最多一個禮拜 他會回來的.
人原來有些時候.
需要紓張一下心靈裡面的感受.
大衛他都有這些時刻.
當他被訴羅追殺的時候.
他去了哪裡.
本來這個是訴羅殺死千千的.
大衛殺死萬萬被人擁戴的大將.
他應該是元帥.
應該是下一任的皇帝.
被大衛追殺的時候.
他要將他的父母送去安全的地方.
求人去收留他的父母.
心想我大衛是誰.
為什麼淪落成這樣.
當訴羅追殺他的時候.
他要躲在山洞裡面.
山洞要來做什麼的.
山洞應該是當時的人.
去到山上.
打仗也好 做其他事也好.
要躲起來 埋伏.
或者要來大解.
甚至乎有些山洞是要來做墳墓.
我大衛是誰.
為什麼淪落到今天.
我們覺得他很帥.
當訴羅進去大解的時候.
大衛在同一個山洞裡面.
他剪了他的衣角下來.
多帥.
多英俊.
他的感覺是我手震的.
我怎麼知道是他死還是我死.
他還有大軍在.

$^{241}$陪我躲在山洞裡面有多少人.
哪裡能夠及他的大軍.
是 他也會有這些時刻.
男士在外面工作的時候.
其實現在也是.
女士 專業人士.
在外面工作的時候.
每一天都在面對著不同的壓力.
在疫情期間.
在家工作有在家工作的壓力.
出去見人 見客.
或者要去辦公室工作的時候.
也是當面對著人工作的壓力.
面對不同的工作有不同的壓力.
甚至於疫情當中.
教小朋友 特別幼稚園的時候.
你怎麼跟一個三歲的小朋友.
我跟你在這裡zoom.
哇 老師們 他們真是很厲害.
用很多心血去預備.
他知道他不單是教那個三歲的小朋友.
他教他二十多歲的媽媽.
或者是他四十多五十歲的媽媽.
因為人家要繼續幫小朋友.
啊 你知不知道今天上了什麼課啊.
面對很多不同的壓力.
有些人他做一年的工作.
主要是要簽一張單.
我們就會覺得.
你另外的三百六十四天就放假了.
是那天簽的那張單.
他整年的努力就是為了那張單.
那張單簽不到的時候是怎樣.
他整年的壓力.
特別是到那個月 那個星期的感受是怎樣.
我們應該知道.
有時候你就是這樣 我就是這樣.
我們就是在面對不同的事情.
不同的壓力.
然後身邊的人好像不一定很明白.

$^{281}$但是很明白我們的神.
我們都可能對他同樣的反應.
對家人就會跟他說 你不明白的了.
我沒事 我打一下機而已.
我打一下機而已 打機 打機 打機.
其實是在抒發內心.
不是因為那個遊戲有什麼吸引.
只是那個是一個.
可以跟環境隔開一段時間.
大衛他都經歷過.
大衛他的經歷告訴我們.
他說他感覺到連黑暗都要來壓碎他.
這個是很大的壓力.
他說周圍你覺得是光.
我覺得是很暗.
一點盼望也沒有.
但是天父 原來多黑也好.
你都覺得很有盼望.
多暗也好 你都覺得生機處處.
對你來說 黑暗算什麼.
就是光明.
我是你所創造的.
但是他說 我是你所創造的.
我還沒出生之前 你都已經知道我了.
我的心肺.
我的心肝脾肺腎 你都知道了.
意思是.
我心想什麼 怎麼逃得過你的眼睛.
你一直監察著我.
我心裡面那種恐懼.
被追殺時候那種恐懼.
要送父母去 請人保護他們的時候.
那種無奈和慚愧.
面對父母那種 真是沒面見他們的感覺.
只有你知道.
當我逃難的時候 要裝瘋扮傻.
玩口水 查牆壁的時候.
被人恥笑的時候.
天父.
天父只有你知道.

$^{321}$只有你知道.
在我們最黑暗 最暗中的時候.
我心裡面的形態 真實的形態.
對你都不是隱藏.
都不是隱藏.
你是我的主 我的神.
我未成年的身體 你一早看過.
我未想定的東西.
就算我的感受未出生 我都未知 未明白.
未了解到自己為甚麼是這樣的時候.
你都已經了解我.
是呀 你是我的神.
我的事你全部知道.
不單止這樣 你有一個很奇妙的行動.
就是在你的日記簿裡面.
我每一件事你都寫在裡面.
神呀 你的確沒有忘記我.
你的確沒有忘記我.
我去到哪裡 你都記得我.
我裡面是怎樣 你都知道.
不過沒有因為我心裡面那種恐慌.
那種絕望.
你真的將我當作絕望擺在一邊.
你寧願在你的日記簿裡面.
還有很多事情為我預備著.
每一件事情你都是為我安排到最好的.
甚至乎在我覺得最絕望的時候.
你要為我做的事情還是像海邊的沙一樣.
我心靈清醒過來的時候.
我要找回你 天父.
我要回到你的身邊.
我回到你的身邊做甚麼呢.
我是很難受.
但我要抓緊一條路.
這條路就是.
你覺得好的東西 我都要覺得好.
你覺得不好的東西 我都覺得不好.
我要學習你的價值觀.
去接受面對你處理的方法.
是的 很多事情是去到絕望.

$^{361}$陰暗當中當你懷著一絲絲一點點的盼望.
覺得還有一點點光的時候.
可能你告訴別人 別人還會恥笑你.
沒有用的.
半年前.
阿Mo 李啟賢.
他受傷在醫院的時候.
他父親一直堅持的醫治.
是父愛我們很明白的.
牧師的信心我們都很明白.
但我們有多少真的會相信有一天.
阿Mo可以站起來.
可能我們只是覺得.
天父給李牧師一個時間.
給師母一段時間.
讓他們去接受這件事.
誰知道.
李啟賢真的可以站起來.
走他的第一步.
走他新的生命的第一步.
不是說他剛剛順了.
是他的肉身的生命好像真的重生.
這個是他的新的生命.
不只是走一步 不只是站起來.
他是走了一個圈.
是的.
生命有沒有盼望.
半年前相信很多人都告訴他.
堅持要有個限度.
但半年後.
大家都會告訴他.
繼續努力.
不只是走一個圈.
是的.
我們有面對事情的時候.
我們的信心是.
我們明白神所說的.
我們都會相信無論是什麼情況.
對李啟賢都是最好的.
無論是什麼情況.

$^{401}$阿Mo都是去到天父的身邊.
不要讓他要受那麼多的苦.
但是他不單止可以跟父母有溝通.
他可以站起來.
可以跨出第一步.
可以繼續走下去.
他以後每一天都要繼續學習走下去.
大衛他有面對他的山洞的時候.
面對他要裝瘋賣傻求平安的時候.
但是這段時間裡面.
神也都施恩典給他.
讓他在不同的時候.
他都有妻子.
他娶了幾個妻子.
神給他有兒女.
大衛原來他的生命.
在神的手中一直都是那麼寶貴的.
神安立了他.
鼓勵了他做下一任的皇帝的時候.
這件事一直都沒有改變.
無論他是躲在山洞當中.
剪掃來的衣角的時候.
或者他裝瘋賣傻的時候.
這個英雄一直都存在在神的記事簿當中.
神知道他會走過這段艱苦的道路.
但是幫助他.
幫助他讓他平安經過這個階段.
這個是扶持著大衛的神.
這個也是扶持我們的神.
我們經過過很多困難的情況 是不是.
昨天我去了一個喪禮.
這位長者經歷過第二次世界大戰.
他成長的當中跟人溝通.
產生很多的困難.
但是他安然渡過他整個的生命.
他的過程當中他有機會聽福音.
他有機會祈禱.
他有機會去親近神.
他走過他整個人生的時候.
晚年的時候他有工人去照顧.

$^{441}$他生養的幾位子女 各有所成.
經歷過的 從小時候開始有很多不開心的事.
很多的艱難困苦.
戰爭的時期更加.
但是生命完結的時候.
可以是燦爛的.
回想的時候可以看到很多美麗的地方.
我跟他的一個女兒聊天的時候.
她就說有很多地方以前不明白父親.
原來父親是愛自己.
很疼愛自己.
很想自己多一些陪他.
是的 我們有很多時候.
生命裡面會遇到不同的事情.
人家不一定明白.
父母不一定明白 子女都不一定明白.
但是神仍然為我們有美好的安排.
即使到了他離世了.
但是子女也都已經明白了.
在他晚年的時候多去陪他 多去關心他.
為他預備美好的住處.
幫他請好好的工人 每天去幫助他.
我們的生命有時候會很不同.
要面對很多不同的情況.
但是在不同的情況的時候 神與我們同在.
三年的疫情大家怎樣過?.
網上出了很多形容詞.
「陽腔陽過」都出來了.
不知道大家有什麼名稱.
我家裡人就叫我做「九陰真經」.
因為我一直都沒有機會.
很近距離接觸病毒的.
我會去喪禮 會去醫院探訪.
我家裡人有感染過 我一直都沒有事.
是不是很開心地過呢?.
大家都一定在這三年裡有過不同的心情.
在2020年的3月.
我們就已經有兩位親戚70多歲的.
應該是在Alpha的那一段時間.
在疫情底下離世.

$^{481}$那種傷痛無奈就是.
喪禮是要在網上進行.
你見不到最後一面的.
是的.
我們這幾年有很多不同的經歷.
只不過疫情是其中三年.
已經有很多事情我們不知道怎樣說出來.
最好不要再想它.
不過前面還有很美麗的時光神帶著我們去.
神仍然看顧著我們.
有人問我.
你們教會是不是開放回去崇拜.
當然啦 政府都宣佈了.
是不是少了很多人呢?.
弟兄姊妹 我們不是少了很多人.
在教會信了主的一直都信了主.
缺了子的 缺了子.
受了洗的 是我們的會友.
這個是不變的事實.
我們現在只是什麼呢?.
我們更加有空間.
可以邀請新朋友來到我們當中.
我們更加要有一股的熱誠.
將我們在主的恩典底下怎樣經過的生命.
跟人家分享.
讓他們明白.
當我很想飛到海極的時候.
飛上高天的時候.
神在哪裡呢?.
當我萬事都是黑暗.
光明都好像黑暗的時候.
神為我有什麼預備呢?.
我們都將我們的生命.
我們的經歷.
被神使用.
好像是詩人大衛最後他說.
是的 你經常都看到我神.
天父 你就看緊一點我.
你給了我很多試驗.
監測還有一個試驗的意思.

$^{521}$你繼續試驗我.
還有一樣東西.
我可能裡面都有些不好的.
你既然看得我這麼穿透.
在我還沒做惡事 還沒想惡事的時候.
你就引導我.
引導我走永生的道路.
我們都要學習.
全這種的盼望去生活.
全這種的盼望去禱告.
不是因為不好的環境已經過了.
而是不好的環境來不及都會來的.
我們不會每一餐都是鮑魚雞.
我們有些時候都要有些鹹魚青菜.
都很好吃的 是不是?.
都會很健康的.
為了神 在我們身上的計劃.
好與不好的 我們學習去感恩.
約瑟他為神帶他去埃及.
他坐了很久的牢.
他被人誣衊.
他對他的主母不尊重.
非禮他主母的事上.
用這種方法帶他去埃及.
讓埃及人 讓他自己的父母.
他自己的家庭.
他的兄弟姐妹都能夠生存.
能夠在大旱災當中生存.
他十個兄弟出賣他.
他竟然跟他們說.
我會養你們和你們的子孫.
那種的.
他明白這個是神的帶領 神的預備.
在神的日記簿裡面的一件事的時候.
他願意感恩著.
你去面對這些神在他身上的計劃.
這位呢 大家有沒有看過一套片叫《密室》呢.
有一本書的.
它是在說納粹時期的一個家庭的遭遇.
這個家庭叫做Ten Boom.

$^{561}$隨便翻譯了 叫做Ten Boom.
我不太懂得怎樣翻譯.
當中有兩個姐妹一起被抓去集中營的.
一個叫做Betsy的 碧茜.
她是姐姐來的.
她死了在集中營裡面.
寫回一些回憶錄出來的那個是Cory.
那個是妹妹.
她寫出來的時候.
她覺得這個姐姐簡直就是集中營裡面的牧者.
是她的典範.
是看著她 照顧她的.
雖然她死了.
但是她死了之後幾天.
Cory就被釋放出來.
Cory被釋放出來據說是一些文件上的錯誤.
集中營有文件上的錯誤 她被釋放出來.
總之神就看著她.
但是她在裡面的時候.
她就說這個姐姐Betsy做了很多事.
當時她不單止是不理解.
簡直是覺得不可以跟她一起這樣去祈禱的.
因為這個Betsy她什麼事都要去感恩.
我想我一定不會做這件事的.
如果你在一個困難的環境裡面.
但是她真的在集中營裡面的時候.
她向集中營裡面其他的同營的營友.
那些受迫害的人傳福音.
他們每天讀聖經.
那段時間他們將《新約聖經》讀完兩次.
因為他們開始讀第三次 還沒有完成的時候.
Betsy就去了天父的書院.
她為他們住的地方去感恩.
那Corrie就說.
你有沒有搞錯啊.
為了住在集中營去感恩.
她說《貼瘦萊利加傳書》五章是這樣說的.
凡事謝恩.
你都要跟我一起謝恩.
那Corrie基於對姐姐的尊重.

$^{601}$反正只有兩個人.
這麼難得還在同一個營舍.
那你說讀經就當然讀經了.
你說祈禱就當然祈禱了.
好啊 你說謝恩.
都謝恩了.
謝恩他們住在集中營裡面.
有失馬.
咬到他們很不舒服.
我們可能會為這件事祈禱的.
如果是我們 神啊 拿走這些東西.
我知道是你所創造的.
給他們另一條路去走.
給他們另一個地方去住.
不用跟我們住在同一個房間 同一張床.
那些被鋪啊 木啊的地方.
很多這些生物的滋擾.
你不能夠要求別人幫你清理.
沒有人肯走近他們的營舍.
到後來 差不多他們出來之後.
他們才知道.
就是因為他們的房間有失馬.
神保住了他們.
其他在同一個營舍裡面.
其他房間的女士.
都受到納粹黨的人員很多的侵犯.
到底甚麼是恩典.
不好意思 我們真的不是很看得清楚.
不過神知道.
神知道怎樣對我們是恩典來的.
是幫助我們.
這個是我們一位姐妹的手.
我得她同意跟大家分享.
在我們程序表的當中.
第七頁裡面.
有幾位我們代討的對象.
他叫做廖澤安.
在第二排 打動第二排的第三位.
是他的手.
這個是比較現在好了一點的情況.

$^{641}$之前我收到的照片 我不敢給大家看.
是見了骨的.
腳還是見骨的.
走每一步路是痛苦的.
他自己坐車 星期三回到屯溪.
手痛不痛 痛.
痛多久 不斷痛.
曾律師跟他聊天 關心他 或者祈禱.
好像針刺一樣 整隻手.
整雙腳都是痛.
所以你說我甚麼動作 甚麼位置不痛.
不斷痛.
祈禱我痛到睡不著的時候.
我祈禱求神幫我.
神竟然讓我痛著都睡得著.
這個是一個藥物的副作用.
這個藥物 曾律師問他你要吃多久.
他說有生之年.
除非吃到它無效.
有效我都會繼續吃.
我與祂共存.
是穿得沒那麼痛.
我現在可以出街.
因為這個藥令到我.
間中會不知道自己在甚麼地方.
我來做甚麼 我下一步要做甚麼.
會令到我坐車 下錯車.
下了車不懂回家.
但是當這個時候我很害怕.
我就找個地方可以立刻坐下.
因為站一步我都痛.
可以靠著坐一坐.
我就祈禱.
我一祈禱我的心靈的感覺就開始安靜下來.
然後慢慢我就想到我下一步要怎麼做.
最後我一定回到家.
弟兄姊妹.
你和我的生命.
都可能有些時候是很艱難.
我們有幾位.

$^{681}$我們代祖的對象 他們的生命經過.
現在都是很艱難.
是每一步.
24小時.
他們身上的手腳.
都是在痛苦當中.
照顧他們的人都很辛苦.
對比之下我們所面對的.
可能是.
可能真的小事而已.
小事而已.
過得到的.
而且神知道 神看著的.
神明白的.
不是要我們白白去受什麼.
可能.
是我們自己和別人都能夠得益的.
我們為了一些事情我們都學習去感恩.
不是因為事情順利.
不是因為事事如意.
如果事事如意的時候.
我們年初的時候不會恭賀別人這些的.
那是因為我們的神.
祂一直都看著我們.
祂用障礙來圍繞我們 保護我們.
祂不會令我們受不必要的痛苦.
祂知道我們所承受的.
祂也有足夠的能力保護我們.
而且祂有足夠的心力陪我們去走過.
每一步.
更加.
祂看重我們每一個.
祂重視我們每一個.
祂稱我們每一個是祂的兒女.
祂不會放下了.
在任何情況下祂不會放下了.
不去理會我們.
即使當時.
阿Mo大家都覺得祂只是一絲氣息.
當然祂前路還有很多事情我們需要為祂祈禱.

$^{721}$例如說 祂說.
接下來下一步要拆掉尿管.
不是這麼容易的 不是說拆就拆的.
過程和適應都需要很多的待同.
大衛到了最後.
在詩篇139.
23,24節的時候.
他講回第一,第二節.
我們都會回看我們的主題.
我們心知.
我們的神.
祂是深深地愛著我們.
祂愛著我們.
讓我們因著我們的知道.
我們學習去明白.
學習去接受.
還有.
學習去感恩.
正如Betsy.
她在集中營當中.
餵著那些室拿去感恩.
這些我們不明白為什麼有.
但是總是為我們好的事.
我們一起低頭祈禱.
因為你是我們的神.
因為你重視我們.
因為你愛我們.
因為你.
你願意這樣來照顧我們每一個.
你不覺得疲勞.
你不覺得承擔得太多.
你不覺得太多.
而至於要少創造一個我.
你愛我們每一個.
你是我們的神.
你是我們的主.
你是保護著我們.
監察著我們.
環繞著我們.
好像帳幕一樣保護著我們.

$^{761}$我們的神.
我們的父.
願你的恩典上在我們心靈裡面.
願我們走的每一步.
願我們全著感恩.
全著信心去走.
願我們在你的帶領當中.
我們都能夠明白.
每一步.
可能就是我們的見證.
幫助到見到的人.
願你的恩典上在我們心靈.
我們又上在你的心中.
祈禱奉主耶穌基督的聖名.
阿們.
\newpage



\section{詩篇 139:1-24}
\label{sec:la_oxDk2vYY}
\textbf{2023-02-12 午堂《不分離不拋棄的愛》詩篇 139\_1-24 講員:陳麗蟬牧師}
\newline
\newline
連結: \href{https://youtube.com/watch?v=la_oxDk2vYY}{\texttt{https://youtube.com/watch?v=la\_oxDk2vYY}} ~~~~ 語音日期: 2023-02-17
\newline
\newline
\hyperref[sec:raKQe4zsFjc]{\small{< < < PREV SERMON < < <}}
~
\hyperref[sec:index_chronic]{\small{[返順時目]}}
~
\hyperref[sec:index_scriptual]{\small{[返順卷目]}}
~
\hyperref[sec:oRiBGTYbBfo]{\small{> > > NEXT SERMON > > >}}
\newline
\newline
詩篇 139:1-24
\newline
\begin{longtable}{cl}
\hline
\hline
章節 & 經文 (和合本修訂版)\\
\hline
139:1 & \begin{tabularx}{0.7\textwidth}{X} 耶和華啊，你已經鑒察我，認識我。 \end{tabularx} \\ \\ \relax
139:2 & \begin{tabularx}{0.7\textwidth}{X} 我坐下，我起來，你都曉得；你從遠處知道我的意念。 \end{tabularx} \\ \\ \relax
139:3 & \begin{tabularx}{0.7\textwidth}{X} 我行路，我躺臥，你都細察；你也深知我一切所行的。 \end{tabularx} \\ \\ \relax
139:4 & \begin{tabularx}{0.7\textwidth}{X} 耶和華啊，我舌頭上的話，你沒有一句不知道的。 \end{tabularx} \\ \\ \relax
139:5 & \begin{tabularx}{0.7\textwidth}{X} 你前後環繞我，按手在我身上。 \end{tabularx} \\ \\ \relax
139:6 & \begin{tabularx}{0.7\textwidth}{X} 這樣的知識奇妙，是我不能測的；至高，是我不能及的。 \end{tabularx} \\ \\ \relax
139:7 & \begin{tabularx}{0.7\textwidth}{X} 我往哪裡去，躲避你的靈？我往哪裡逃，躲避你的面？ \end{tabularx} \\ \\ \relax
139:8 & \begin{tabularx}{0.7\textwidth}{X} 我若升到天上，你在那裡；我若躺在陰間，你也在那裡。 \end{tabularx} \\ \\ \relax
139:9 & \begin{tabularx}{0.7\textwidth}{X} 我若展開清晨的翅膀，飛到海極居住， \end{tabularx} \\ \\ \relax
139:10 & \begin{tabularx}{0.7\textwidth}{X} 就是在那裡，你的手必引導我，你的右手也必扶持我。 \end{tabularx} \\ \\ \relax
139:11 & \begin{tabularx}{0.7\textwidth}{X} 我若說「黑暗必定壓碎我，我周圍的亮光必成為黑夜」， \end{tabularx} \\ \\ \relax
139:12 & \begin{tabularx}{0.7\textwidth}{X} 黑暗對你不再是黑暗，黑夜卻如白晝發亮。黑暗和光明，在你看來都是一樣。 \end{tabularx} \\ \\ \relax
139:13 & \begin{tabularx}{0.7\textwidth}{X} 我的肺腑是你所造的，我在母腹中，你已編織我。 \end{tabularx} \\ \\ \relax
139:14 & \begin{tabularx}{0.7\textwidth}{X} 我要稱謝你，因我受造奇妙可畏，你的作為奇妙，這是我心深知道的。 \end{tabularx} \\ \\ \relax
139:15 & \begin{tabularx}{0.7\textwidth}{X} 我在暗中受造，在地的深處被塑造；那時，我的形體並不向你隱藏。 \end{tabularx} \\ \\ \relax
139:16 & \begin{tabularx}{0.7\textwidth}{X} 我未成形的體質，你的眼早已看見了；你所定的日子，我尚未度一日，都在你的冊子寫上了。 \end{tabularx} \\ \\ \relax
139:17 & \begin{tabularx}{0.7\textwidth}{X} 神啊，你的意念向我何等寶貴！其數何等眾多！ \end{tabularx} \\ \\ \relax
139:18 & \begin{tabularx}{0.7\textwidth}{X} 我若數點，比海沙更多；我睡醒的時候，仍和你同在。 \end{tabularx} \\ \\ \relax
139:19 & \begin{tabularx}{0.7\textwidth}{X} 神啊，惟願你殺戮惡人；你們好流人血的，離開我去吧！ \end{tabularx} \\ \\ \relax
139:20 & \begin{tabularx}{0.7\textwidth}{X} 他們說惡言頂撞你，你的仇敵妄稱你的名。 \end{tabularx} \\ \\ \relax
139:21 & \begin{tabularx}{0.7\textwidth}{X} 耶和華啊，恨惡你的，我豈不恨惡他們嗎？攻擊你的，我豈不憎惡他們嗎？ \end{tabularx} \\ \\ \relax
139:22 & \begin{tabularx}{0.7\textwidth}{X} 我恨惡他們到極點，以他們為我的仇敵。 \end{tabularx} \\ \\ \relax
139:23 & \begin{tabularx}{0.7\textwidth}{X} 神啊，求你鑒察我，知道我的心思，試煉我，知道我的意念； \end{tabularx} \\ \\ \relax
139:24 & \begin{tabularx}{0.7\textwidth}{X} 看在我裡面有甚麼惡行沒有，引導我走永生的道路。 \end{tabularx} \\ \\
[1ex]
\hline
\hline
\end{longtable}
$^{1}$各位姐妹早安.
前幾天看新聞的時候.
我相信大家都感到很難過,很震驚.
就是土耳其和敘利亞一帶發生大地震.
新聞網上都流傳了一些很感人的照片.
其中一幅就講到一個爸爸.
一直牽著一個已經死去的女兒的手.
可能在地震的時候這個女兒走不及了.
所以在床上死了.
爸爸去到現場的時候一直牽著她的手.
她不肯離開,她不肯離開.
另外一幅就講到有一個爸爸被救出來.
他在後面又揹著一個小朋友.
救出來的時候無論後面的小朋友.
或者是他自己都已經死了.
不過在他前面還抱著一個BB.
但是我們就看到了.
在這個危急的裡面.
我相信這個爸爸他自己走.
不用理會小朋友,BB的時候.
他可能走得及的.
但是他很想去保護著兩個小朋友.
他為了努力保護自己的家人.
他都沒有理會自己的生命.
在這個詩篇139篇.
我們就看到這個詩人大衛.
他講到上帝是一個極之仁愛的上帝.
是一個很仁慈的上帝.
他為什麼會這樣講呢?.
在詩篇裡面整個24節裡面.
他都講到他很明白每一個人的警惕.
他特別用那些字眼.
都講到人心裡面的黑暗和污穢.
雖然他知道了 他都知道了.
但是他仍然想幫助我們,保護我們.
以致我們可以從一些的幽暗.
一些污穢的情況裡面跳出來.
剛才我們主席幫我們讀了1至6節的經文.
他講到上帝對我們完全的了解.
他是一個無所不知的上帝.

$^{41}$但是他知道這麼多東西.
是不是想來攻擊我們.
或者想來挑剔我們呢?.
不是的 他其實是想將我們救來這個黑暗.
在這6節裡面 大衛就用了6個的詞語.
講明了上帝對我們的認識.
第一個字就說 上帝已經監察我們.
這個監察這個字就是講到那些珠寶商.
或者到今天讓我們去金舖.
你去賣鑽石那些.
就去到將這些珠寶給他看.
他就用一些儀器或者用一些很強的光.
去照這顆珍珠或者鑽石.
以致這顆珠寶裡面如果有任何的一些瑕疵的時候.
他都可以知道的.
他就知道這顆珍珠或者鑽石究竟是不是真的.
或者那個沒有瑕疵的程度去到怎樣.
其實就講到上帝很了解我們.
知道我們裡面的生命.
我們的心靈裡面究竟有多少的幽暗.
多少的瑕疵.
人家不知道 上帝都知道.
最後他就用了認識去認識我們.
我們坐下 我們起來 他都知道.
這個在原文裡面同樣是一個詞語.
就是講到他很清楚知道裡面哪一樣東西是真的.
哪一樣東西是假的.
究竟我們是真好人呢.
還是其實是一個大惡人.
還是說一個好人是一個很差很差的人.
上帝知曉了我們裡面那個人的那個境況.
知道這個字都是了 他很了解很明白的.
都可以看得很清楚 他很微細的東西都看得到.
這個字是講到在農業的社會裡面.
可能如果我們大家在國內裡面去成長的.
我們就可能會清楚多一些.
那些農夫在篩穀物的時候.
因為那些皮和穀物就混在一起的.
只要他放在大篩機那裡.
將這些穀物一放飛上天的時候.

$^{81}$那些輕的那些坑坯那些皮就會被風吹走了.
而這些穀物掉下來的那些穀物.
這些就叫做「篩擦」.
那個意思就是說.
在上帝裡面沒有任何東西可以隱藏的.
那些不是真正的穀粒 那些皮那些東西.
沒有什麼大作用的那些東西.
全部都會被吹走的.
他就講到上帝 同樣講到上帝很知道我們裡面.
在外面我們可能很漂亮.
穿上一件西裝外套在那裡講道.
好像一個很好的小姐,很好的先生.
但是裡面我們多麼的奸狡,多麼的污穢.
你們不知道 不過上帝知道.
這個「心知」這個字眼.
我們無論怎麼樣都好 他都「篩擦」.
他「心知」我們一起.
他就好像一個老朋友那樣.
他知道我們的想法.
有時甚至我們講好朋友好到一個地步.
我們眉毛動一動.
他都知道我們的心在想什麼.
這個就是講到上帝.
源源處處知道我們裡面內心的世界.
有人就形容了.
他說這個人性很複雜的.
但是就好像一個大洋蔥那樣.
在外表上 人看到我今天就這樣.
正正的站在那裡 都覺得我多好.
但是當我們剝開洋蔥的第一層皮的時候.
可能就有一些很懶惰,很驕傲.
第二層就是一些很自大,很不知道怎麼樣的.
他說通常人家不太認識我.
就只是看到我外面.
那些好朋友熟悉一點的.
就知道我說這個人很厲害的.
又自高自大的.
因為再深入到我裡面.
第四,第五,第六層那些.
就說我是一個可能真的很陰險.

$^{121}$很喜歡自吹自擂.
只是講自己好東西.
就講別人壞東西的那些人.
但是他到深入裡面那些.
唯有上帝才知道.
上帝很清楚知道.
無論我們坐下來,睡在那裡.
又知道我走路,我這個起來.
祂都知道.
講到上帝是一個無所不知的上帝.
祂樣樣東西都知道.
有時我們人是這樣的.
我們很喜歡被人知道多一點的 是不是?.
很想在群體裡面被人認識多一點的.
所以我們就會想做多一點的門面功夫.
希望出名多一點.
但是有些人覺得不用的.
只是去追求得到上帝的接納,上帝的認識.
好了 想問一下.
認不認識這個人?.
他叫做何思德.
有沒有人認識他呢?.
沒有啊 昨晚晚堂講都沒有人認識他.
有沒有人認識戴德生?.
認識了 有人點頭給我了 是不是?.
很多人都認識戴德生.
我們知道戴德生先生.
他是內地會的創辦人.
但當60多歲.
當戴德生中了風之後.
他就需要找接班人了.
當時他就找了何思德.
何思德這個弟兄.
他是一個很有錢.
在英國的時候是很有錢的.
另外五個同學.
都是劍橋大學裡面畢業的.
另外還有一位青年軍官.
在歷史上稱他為劍橋七傑.
他們就一起來到中國宣教.

$^{161}$何思德自己一直事奉的一句很出名的說話.
他就說.
Live to be forgotten.
in order that Christ may be remembered.
他的意思就是說.
他的事奉.
他不願意自己很出名.
他不需要自己很出名.
他只願他的事奉.
能夠將耶穌基督的名字來高舉.
原來何思德這個人.
他在內地會做了33年.
那段時間是中國最動盪的時間.
沒有任何的人敢來中國做宣教.
但在何思德做內地會的總幹事的時候.
他竟然招募到200個新的宣教士.
200個年輕的宣教士.
來到中國裡面去事奉.
現在比起我.
多麼威風.
其它的警察都找不到人來宣教.
但是他竟然找到200個人來宣教.
比起我.
我當然想全世界都知道我多有成就.
如果不是看一本雜誌傳書.
我們辦公室那裡都有的.
連續七期介紹他.
我都不認識他.
何思德他一直的想法就是.
他要隱藏在耶穌基督裡面.
他不需要讓人知道.
他只是要耶穌基督的名字被高舉.
他只需要耶穌基督知道他的忠心.
我相信耶穌基督知道的.
並且耶穌基督現在已經在不斷的讚賞他.
人性很複雜的.
我們大家外表.
你不太知道我內心的世界.
我也不知道大家內心的世界.
有一位德國的神學家.

$^{201}$他就這樣形容.
他說人的心靈其實是一部恐怖片.
聽起來都很恐怖.
裡面有色情,有暴力,有仇恨.
是一部三級片.
大家想一下.
是不是去到這麼差呢?.
我覺得每個人的程度不同.
每個人的程度不同.
他就好像強烈了一些.
當耶穌基督的光照進每一個人的心靈的時候.
就知道自己的心靈有多少的瑕疵.
有多少的黑暗.
不過我們感恩.
無論我們的本相是怎樣的時候.
上帝仍然的愛我們.
所以在第五,第六節.
第一至四節就說.
上帝知道我們內心裡面所有的瑕疵.
所有的污穢.
但到第五,第六節就給我們安慰,鼓勵.
他說上帝在我們的前後環繞著我們.
安守在我們的身上.
這個環繞就說一些的型類.
好像出去打仗的時候.
在附近弄了一些型類.
這些型類就是用來保護的.
還說安守在我們的身上.
安守也是一個保護,庇護的象徵.
就是說雖然我們真正的裡面的人.
是一個很污穢的人.
但是我們的上帝不會這樣.
每一次就來挑剔我們.
來審判我們,懲罰我們.
祂為的是要保護我們.
救我們出這個黑暗的內心的生命.
所以接著就說.
這樣的知識是奇妙的.
是我不能測的.
至高是我不能及的.

$^{241}$所以大衛就說.
一般人來說知道有些不好,有些差的話.
當然是死囚,慢囚.
將你們挑出來.
希望打倒你.
但是這個上帝.
祂竟然接納我們.
並且想辦法來將我們保護.
救我們出這個黑暗.
所以人沒有辦法想得到.
這個上帝的愛是多麼的大.
就好像耶穌基督那樣.
在十字架上為我們寫信.
接著我們看第二方面.
上帝是無所不在的.
但是祂仍然保護我們.
是第七至十二節.
我們請嘉文.
第七節.
我往哪裡去躲避你的靈.
我往哪裡逃躲避你的面.
我若升到天上你在哪裡.
我若躺在陰間你也在哪裡.
我若展開清晨的翅膀飛到海極居住.
就是在哪裡你的手必引導我.
你的右手也必扶持我.
我若說黑暗必定壓碎我.
我周圍的亮光必成為黑夜.
黑暗對你不再是黑暗.
黑夜卻如白晝發光.
黑暗和光明在你看來都是一樣.
在這段裡面.
詩人就很想說一件事.
就是有時我們很想逃離開.
這個上帝的眼光.
我們不喜歡被祂經常管著.
所以看到詩人就用了一個很強烈的對比.
我若飛到天上你也在哪裡.
我下塌到陰間你也在哪裡.
天上當時人們就覺得是最高最高點.

$^{281}$去到這個陰間是死人居住的地方.
是最低最低的地方.
所以無論我們想躲起來去到最高的地方.
或者去到最低的地方.
上帝都在那裡.
然後我若與一個清晨的翅膀飛到海的居住.
就算在那裡 上帝都在那裡.
這個清晨的翅膀就說到.
就好像太陽有一對翼一樣.
連太陽都想逃避.
就飛到很高很遠的地方.
去到這個天涯海角.
希望上帝找不到.
不過呢 上帝又在那裡 又看著我.
有一個無神論者就說.
「God is nowhere」.
上帝都不在這裡 根本上帝不存在.
但是那個信徒就將這個「nowhere」這個字改了.
「God is now here」.
上帝在這個時間就在這裡了.
上帝是隨時隨刻都在的.
我記得好像在《五堂》裡面講過這個故事.
一個爸爸和一個兒子回完教會了.
就開車回家.
在路途當中就見到有一個西瓜田.
那些西瓜又大又靚.
這麼多個拿去一個沒什麼所謂.
於是他就停在車旁邊.
就叫兒子「你幫我看著有沒有人來」.
「我現在下去拿個西瓜」.
一有人來就大聲喊過.
他一下車之後 兒子就大叫.
「有人來了」.
那個爸爸就嚇得西瓜都拿不到.
快點就上車.
看看 兒子 兒子 沒有人了 你近視啊.
兒子說「天父在這裡啊」.
「舉頭三尺有神明」.
我們中國人都會說的.
上帝隨處都在.

$^{321}$我們講回一下.
當他講到我們這個人心.
有時候我們真的很想逃避上帝.
余大森博士有這樣的形容.
當我們陷在一個很低潮很低潮的境況.
我們不想見人.
我們也不想祈禱.
我們不想見神.
甚至容讓自己去自暴自棄.
陷在一個罪惡的深淵.
不知道大家試過沒有.
去到一個這樣的境況裡面.
很低沉 很黑暗.
但是上帝仍然在.
只要我們去悔改 回轉的時候.
上帝的光就會將這個幽暗.
這個污穢的心來照亮的.
我剛剛出來做傳道的時候.
我都是做中學生的工作.
做青少年的工作.
每年暑假我就會帶他們.
去到這個狼啟解讀中心那裡去露營的.
在當中就認識了一位弟兄.
他也是村裡面做Head.
處理很多很多的事務.
每次去紅磡一些大型的報道會的時候.
又見到他在那裡主持一些長毛的事情.
他就覺得他都很醒目的.
很厲害的一位弟兄.
他熟悉了之後就跟他聊天了.
原來他都大或小的解讀了十幾二十次.
他最後一次成功就是用福音解讀.
他也很明白那些過來人.
解讀那種心癮的辛苦.
他也很明白有些人解讀了出來又吸回那個辛苦.
所以他就決定留在狼啟村那裡去服務.
也都用他自己的一些經歷去鼓勵一些過來人.
也都用上帝給他的恩賜.
來去管理村裡面大小的事情.
在大型的報道會裡面是需要行政人才的.

$^{361}$他又出來幫忙.
從這位弟兄的身上我就看到了.
黑暗 過往那個幽暗的生命.
現在它不再黑暗了.
黑暗卻如白晝發光發亮.
過往這個幽暗的生命到它之後.
悔改了 繼續在狼啟村.
在大型的方聚會裡面去事奉的時候.
但是這個就成為了它那個生命的見證.
黑夜對它來說就好像發亮的光一樣.
所以這個到處都在的上帝.
它飛到天國底 天腳底 它都在的.
它不是來去偵查著我們.
來去抓著我們.
而是它要以它大能的手.
以它大能的手來引導我們 扶持我們.
它說黑暗對你不再是黑暗.
黑夜卻如白晝發亮.
黑暗和光明在你看來都是一樣.
還沒有得罪的 先回到剛才那裡 謝謝.
講到這個黑夜.
上帝的目的都是它到處都在.
不是想著你什麼時候犯錯就罵我們 罰我們.
而是它想去引導我們走出這個黑暗.
這個扶持這個字 講到它大能的右手扶持著我們.
是在說拖實我們.
這個拖實是很大力的 緊緊的抓著我們.
我時時都有一個情景的.
小時候我們住一間房 72個房客.
但裡面我們的關係很融洽的.
有一家人他們後來就申請到荃灣.
就搬到荃灣的公屋住.
當他大女兒結婚的時候.
他就特意走出來灣仔請我們吃飯.
我記得那晚吃完飯10點多就走了.
我爸爸就一直拖著我的手.
其實我就很想甩掉他的手.
走上去珍惜和其他小朋友一起聊天.
玩玩 跳跳舞的時間.
但是我就發覺我爸爸一直很大力的抓著我的手.

$^{401}$他不讓我走 他怕我走就亂飆出馬路.
我明白他其實是想保護著我們.
上帝都是 在我們這個人生的路上.
有時候我們低到那個境況.
我們是不想見他的.
飛到這個天涯海角不想見他的.
但是上帝一直都抓著我們的手.
他不希望我們去到一個.
全然灰心,全然迷失的境況.
他仍然在等著我們.
到第13到18節.
就說到原來我們的生命是上帝創造的.
我們請家文為我們讀出13到18節.
我的肺腑是你所造的.
我在母腹中你已編織我.
我要稱謝你 因我受造奇妙可畏.
你的作為奇妙 這是我深深知道的.
我在暗中受造 在地的深處被塑造.
那時我的形體並不向你隱藏.
我未成形的體質 你的眼早已看見了.
你所定的日子 我尚未到一日.
都在你的冊子寫上了.
神啊 你的意念向我何其寶貴.
其數何等眾多.
我若數點比海沙更多.
我睡醒的時候仍和你同在.
大衛在這一段裡面進一步的說.
他明白了 明白上帝知道我們.
他不是想罰我們 而是想保護我們.
他時時無處不在.
他為了叫我們不要迷失 不要跌倒.
在這裡就明白我們的生命.
原來上帝早有這個安排.
在無福的當中 他已經編織了我們.
我們未出生的時候.
我們一生的家庭 我們的事業 我們的婚姻.
甚至我們的歲數.
上帝都有特別的心意和計劃.
這個編織這個字.
一看就明白了 這個編織.

$^{441}$就是我們織毛衣 或者我們織毛衣.
當我們買了毛衣回來 織毛衣.
我們就會想 我這次要織一件背心.
還是織一件高頭粒的毛衣.
或者這次我想織一些什麼樣的圖案.
我們一定會在腦裡面大約有這個計劃.
無論你多厲害的織毛衣的人都好.
你一定起碼知道 我織在前面開房子.
還是高頭粒 是背心還是長袖呢?.
按照著我們心裡面的計劃.
一些藍圖呢.
可能我們就開始了 織了一幅 又第二幅.
接著就把它們黏在一起.
按照我們心裡面早已經想到.
大約那個構思的時候.
我們就會可以織到一件可以穿的毛衣.
有沒有人試過織了毛衣出來穿不了的?.
有啊 有人幫我織了一件毛衣.
我還是穿不了 因為頭套套不住.
是不是?.
那你去做的時候呢.
我們會織到一件可以穿的毛衣.
我們可以穿得漂亮的.
並且穿得很清新.
讓你看出來都是漂亮的.
其實在這裡施人大衛就說了.
只要我們按照上帝的那個計劃.
去編織那個藍圖.
我們來去生活的時候呢.
我們是可以有一個光明的生活.
是一個榮神益人的生命.
是一個燦爛的生命.
耶穌基督看到我們每一個人.
都是這個寶貝來的.
我很想講一節的金句.
是哥林多前書第二章第九節.
它是上帝為愛祂的人所預備的.
是眼睛未曾看見.
耳朵未曾聽見.
人心也都未曾想過的.

$^{481}$就是上帝自己早就已經知道了所有的東西.
所以有時候我們今天坐落在這樣的光景.
我們很不明白上帝一些的事情.
不過當日子一天一天這樣過的時候呢.
我們就明白到上帝的那個安排.
我記得在我讀神學第四年的時候.
我們在神學這樣的.
第二年班和第三年班.
就有很多學會要做什麼獎什麼獎的.
到第四年畢業了可以什麼都不用做了.
所以到四年班的時候就想著輕輕鬆鬆.
還有有時候我連暑假也收了一些科目.
就讀完了.
所以我第四年班的時候不用收那麼多的科目.
我就想著勤力一點練琴.
還有我們教會有兩個合唱團.
兩個都參加了.
就開開心心的想著畢業了.
但是怎知道開學了一個星期.
我爸爸中風了.
還說他很危險.
想著他不知道行不行了.
所以那段時間我每天都進出長洲.
所以什麼事情我都不敢報名.
總之要畢業了就算了.
有一天在天坎靈修.
這一節的聖經就出現了.
是上帝為愛祂的人所預備.
是眼睛未曾見過.
耳朵未曾聽見.
人心都未想到的.
當時那一口氣就不太順了.
我就跟上帝說.
哇哇 你有什麼好東西給我啊.
我現在前面就沒有好東西了.
原本可以很開心的過最後一年的神學院生活.
到現在就這樣了.
但是人心都未曾想到的.
我就跟上帝說.
我當時都挺大膽的.

$^{521}$我就放長雙眼.
看看你的應許是短的.
上帝很奇妙的.
在兩年之後.
祂就給我再一次有機會.
進神學院的部分時間讀碩士.
當中還給我主修做教會音樂.
但是其實按我音樂上的知識.
我是不合格的.
祂又要求五級琴.
又不知道八級樂理的.
我沒有五級琴的.
我也沒有八級樂理的.
所以不知道為什麼學院就收了我.
它都說了.
前無古人 後無來者.
以後都不會這樣了.
當我讀完第一個學期.
我說很辛苦.
教會音樂只有三個同學.
一個是中大出來的.
教中學的.
每年都去國際校際音樂節比賽的.
這個又是在美國讀完音樂回來的.
我經常跟他們一起就很自卑的了.
所以三個月之後告訴我的老師.
我不讀了 我要轉系.
老師不給我.
他說 上帝帶你進來.
你就要繼續.
上帝自然會帶你出去的了.
當時我心有一憂.
有一句回應.
是不是踢我出校呢?.
上帝又給我不倒業.
所以我每一次想起這個經歷的時候.
就想起我在四年級.
我那時候我爸爸中風.
但是我很恨上帝.
但是他的應許.

$^{561}$雖然我當時很不順氣.
但是當我在碩士畢業之後.
我就想到什麼叫做.
眼睛未曾看見.
耳朵未曾聽見.
人心也都未曾想到的意思.
主的應許是真實的.
他在我們人生裡面所有東西都編排了.
每一樣都是對我們有益的.
最後一部分.
就是大衛.
他想到上帝的全知.
看到祂的全能.
但是他都很想詬病我們.
這個人心污穢的人.
他怎麼回應呢?.
在第19到24節.
請加文.
神啊 惟願你殺六惡人.
你們好流人血的.
離開我去吧.
他們說惡言頂撞你.
你的仇敵罔稱你的名.
耶和華 恨惡你的.
我豈不恨惡他們嗎?.
攻擊你的.
我豈不憎惡他們嗎?.
我恨惡他們到極點.
以他們為我的仇敵.
神啊 求你監察我.
知道我的心思.
試煉我.
知道我的意念.
看在我裡面有什麼惡行沒有.
引導我走永生的道路.
謝謝.
大衛很清楚知道.
上帝的全知.
他什麼都知道.
是想救我們出這個黑暗.

$^{601}$上帝的周圍都在.
不讓我們亂飛.
祂的目的是保護我們.
大衛明白了之後.
於是他在第19到24節.
就對上帝一個回應.
他要完全脫離過往的一些幽暗.
他說你們好流人血的.
離開我吧.
就是說他不再和惡人同流合污.
他不再頂撞上帝.
過往走過的歪路.
他都要完完全全的棄絕.
大衛知道他自己的軟弱.
心靈的軟弱.
也都知道自己會很容易.
就跌落這個試探裡面.
所以他就謙卑的和上帝祈禱.
上帝啊 求你監察我.
知道我的心思意念.
思念我 知道我的意念.
看在我裡面有什麼惡行沒有.
引導我走永生的道路.
一位年輕人在街裡面.
就拿著手提電話大大聲的說.
喂 A公司啊.
請問你們是不是請文員啊.
哦 請到了.
那會不會再請人或者想換人啊.
哦 老闆說他做得很好.
不用換人 不用請人.
哦 好啊 那麼謝謝 拜拜.
那隔壁一個老人家.
聽到他找不到工作.
就去鼓勵一下他.
年輕人 很可惜啊.
你沒有這份工作.
但不用灰心的.
你繼續慢慢找吧.
上帝會為你預備的.

$^{641}$這位年輕人就和伯伯說.
我就是A公司那個文員.
我想知道 我想從旁邊知道.
老闆是不是滿意我的工作.
所以我這樣假裝打電話.
你覺得有點傻傻的.
我相信不知道真假.
但是在這個小故事裡面.
讓我們知道一樣東西.
我們要定期檢查一下.
我們自己裡面的心思,意念.
究竟是否合乎上帝的心意.
耶穌說 人心裡面的惡念.
有兇殺 有姦淫 有淫亂 有偷渡.
有說謊 作假見證 誹謗.
這些裡面的惡念.
都在污穢我們的生命.
使我們的生命不乾淨.
使上帝不喜歡.
所以我們要怎樣呢?.
求上帝以監察來提醒我們.
怎樣去監察呢?.
我們今天有更多的禱告.
更多的讀經.
我們需要找一些時間去安靜.
去審查我們自己的生命.
以致我們求上帝.
指我們怎樣走一條永生的道路.
我們一起禱告吧.
天父上帝 我們謝謝你.
因為當我們還作罪人的時候.
你已經定下了計劃來拯救我們.
謝謝 多謝天父上帝.
你沒有嫌棄我們.
我們的真實的我.
我們的本相就是污穢的.
是這個真的誤土處理的喜悅.
多謝你一路的接納.
一路的拯救.
多謝你.

$^{681}$祝福一祝福我們在這裡每一位弟兄姊妹.
我們時刻在你面前醒覺.
以致我們走一條.
我們整個人生走一條合理心意的道路.
求主去扶持.
求主指引.
求主幫助.
多謝你垂聽禱告.
是奉靠主耶穌基督得勝靈治而求.
阿 明.
\newpage



\section{詩篇 104:31-35}
\label{sec:oRiBGTYbBfo}
\textbf{2023-02-18 晚堂《新年願望》詩篇 104\_31-35 講員:張美薇牧師}
\newline
\newline
連結: \href{https://youtube.com/watch?v=oRiBGTYbBfo}{\texttt{https://youtube.com/watch?v=oRiBGTYbBfo}} ~~~~ 語音日期: 2023-02-21
\newline
\newline
\hyperref[sec:la_oxDk2vYY]{\small{< < < PREV SERMON < < <}}
~
\hyperref[sec:index_chronic]{\small{[返順時目]}}
~
\hyperref[sec:index_scriptual]{\small{[返順卷目]}}
~
\hyperref[sec:sY5rw9BFGjc]{\small{> > > NEXT SERMON > > >}}
\newline
\newline
詩篇 104:31-35
\newline
\begin{longtable}{cl}
\hline
\hline
章節 & 經文 (和合本修訂版)\\
\hline
104:31 & \begin{tabularx}{0.7\textwidth}{X} 願耶和華的榮耀存到永遠！願耶和華喜愛自己所造的！ \end{tabularx} \\ \\ \relax
104:32 & \begin{tabularx}{0.7\textwidth}{X} 他看地，地便震動；他摸山，山就冒煙。 \end{tabularx} \\ \\ \relax
104:33 & \begin{tabularx}{0.7\textwidth}{X} 我一生要向耶和華唱詩！我還活的時候，要向我的神歌頌！ \end{tabularx} \\ \\ \relax
104:34 & \begin{tabularx}{0.7\textwidth}{X} 願他悅納我的默念！我要因耶和華歡喜！ \end{tabularx} \\ \\ \relax
104:35 & \begin{tabularx}{0.7\textwidth}{X} 願罪人從世上消滅！願惡人歸於無有！我的心哪，你要稱頌耶和華！哈利路亞！ \end{tabularx} \\ \\
[1ex]
\hline
\hline
\end{longtable}
$^{1}$大家好, 各位同學平安.
今天是正月二十八.
我看過日曆了.
是正月的月尾.
在此祝大家新年快樂.
當我在此祝大家新年快樂的時候.
我就想起一張四甲漫畫.
我有可能是給過大家看過的.
不過我都不太記得是不是.
這是一個花生的漫畫.
我們可以看到新年的時候.
主角查理布朗半茶里.
在新年的時候碰到露西.
他就跟她說.
露西, 祝你新年快樂.
我想大家新年或者是農曆年.
見到面都會彼此祝賀.
通常來說.
如果是這樣的話.
露西應該會怎樣回答呢.
應該祝她新年快樂.
但她竟然不是.
露西就說.
你說了就是了嗎.
是不是你告訴我.
說了新年快樂.
我就會在新的一年裡.
真的快樂呢.
你可不可以單單說了.
事情就是這樣成了呢.
這是不是一個保證呢.
是不是是不是.
一直追問下去.
查理布朗唯一可以做的就是.
奧天啊.
我想我們其實很多時候.
我們彼此祝賀新年快樂.
新年進步.
有人說比較.
叫做什麼.

$^{41}$使喚一點.
恭喜發財.
或者說身體健康等等.
究竟我們有沒有事情.
可以做的呢.
我們開始的時候.
我們首先一起祈禱.
天父上帝.
求你使用孩子今天所說的.
叫孩子所說的.
都是你要孩子說的.
幫助會眾.
我們有一個能聽到的耳朵.
有一個能明白你的頭腦.
更加有一個能.
願意將你的話語.
藏在心裡的心.
更加幫助我們.
能夠回應你的手和腳.
好叫我們可以一起.
走出你所吩咐我們的.
聽我們的祈禱.
我們這樣禱告.
和耶穌基督的命求.
所以當我們今天看詩篇.
第104篇的時候.
但我都會用很多時間.
將詩篇的說話讀出來.
讓大家可以被神的話語充滿.
我們在104篇的時候.
看見詩人大衛.
在詩裡面.
刻上宇宙萬物的詞句.
他向神獻上感恩和讚美.
有些聖經學者.
將這首詩定為讚美詩.
讚美創造的讚曲.
也有學者將這首詩.
定為感恩詩.
也有慶賀君王登基的詩.

$^{81}$也有新年時候所唱的詩.
今天我們就讓神的說話充滿我們.
所以今天我們很多都會去看這個說話.
很多時候學者都喜歡將聖經某一段.
分成ABCD.
然後對稱下半部的C'B'A'.
其實這五節.
剛才幫我們讀出的五節.
都可以這樣對稱.
當我們這樣對稱的時候.
我們就會看到.
詩人有三個祝願.
第一個祝願是.
「願耶和華的榮耀傳到永遠」.
所以我會叫做.
「願造物的神榮耀永存」.
那個「願」是「願耶和華的神」.
第二個祝願是.
「願祂閱立我的默念」.
所以我就會叫做.
「願蒙恩的我歌頌默念」.
第三個祝願是.
「願罪人從世上消滅」.
「願世界之上罪惡消除」.
詩人的祝願並不是單單說了就算了.
我們當然知道.
願景當然要由神去成就.
但是人也可以在當中促成或造就.
我們在祝願之後的一年之內.
我們會不會有些事去做.
去跟進.
以致我們的祝願.
可以讓我們祝願的人更加能夠達到.
所以今天就想和大家分享.
這三個祝願.
讓大家可以在裡面看到.
其實我們看到.
詩人大衛在詩篇裡面.
默想宇宙的次序.
他向神獻上感恩和讚美.

$^{121}$所以我們願意.
大家都一樣.
在新年裡面立下心願.
我們看第一個願.
是「願造物之神榮耀永存」.
這兩節.
「願耶和華的榮耀傳到永遠」.
「願耶和華喜愛自己所做的」.
「祂看地 地就震動 祂摸山」.
「山就無意」.
詩人在神所創造的當中.
他看到神的榮耀.
所以他第一個就會看到.
神是願意.
祂喜悅他所創造的.
我們看到他說.
「願耶和華喜愛他自己創造」.
我們如果是慢慢地回家細讀.
詩篇第104篇的時候.
就會發覺好像大衛進入了創造博物館.
其實最近香港有很多博物館.
有故宮博物館.
有M+博物館.
大家都可以進去看.
其實大衛去看世界.
他好像是進入了創造博物館.
「安置次序 由天而地 由萬物而人」.
「逐一觀賞 逐一讚歎」.
所以當詩人看到的時候.
他第一個願造物的神榮耀永存.
裡面的小願.
「願神喜悅他自己所創造」.
我們會看到.
當我們去看的時候.
神的創造.
第一天我們會看到.
祂是做了光.
在詩篇第104篇裡都有說.
在第1的「下」到第2的「上」裡說.
「你以尊榮威嚴為衣 披上亮光 如披外袍」.

$^{161}$作者好像是和創世記的第一章所描述的.
互相呼應.
首先用「以人的筆法」.
描述神好像是一位穿上華麗莊嚴的龍袍君王.
神首先創造了光.
好像披上了一件外袍一樣.
我們又看第二天的時候.
神是做空氣和天.
在詩篇第104篇第2節「下」到第4節.
祂這樣說.
「鋪張窮瘡 如鋪萬紙 在水中立流角的棟樑.
用雲彩為車簾 藉著風的翅膀而行.
以風為使者 以火焰為伏翼」.
我們看到第三天.
第三天神又造了海.
地和其中的植物.
在詩篇第104篇第5到第9節.
詩人這樣說.
「將地立於根基上 使地永不動搖.
你用深水遮蓋地面 猶如衣裳.
珠水高過山嶺 你的戚札一發 水便崩逃.
你的雷聲一發 水就崩流.
珠山上升 珠谷下沉 歸你為它所立定之地.
你定了界限 使水不能超越.
不再轉回 淹沒大地」.
神不單在第三天造了水 海和地.
祂更加使地上的植物生長.
在第14節「上」裡面.
祂使草生長 給生畜吃 使菜蔬生長 供給人用」.
我們看到神在第4天造了光體 管奏夜.
所以在第19至20節「上」裡面.
詩人說「你安置月亮 為定季節.
太陽自知沉落 你做黑暗為夜」.
我們再看第5天 神造天上和海中的生物.
在第19至20節裡面.
「天上的飛鳥立在水旁住宿」.
在之間提叫.
然後在第25節又說.
「那裡有海 又大又廣 其中有無數的動物 大小活物都有」.
所以你會發覺在詩篇裡.

$^{201}$可以不斷將神的創造描繪.
到第6天了.
第6天 神造地上的生物和人.
於是詩人又說.
在第20至23節.
「你做黑暗為夜 林中的白獸都爬出來.
少壯獅子叫 向神尋求食物.
太陽一出 獸就躲避 躺在外面凍裡.
人出去做工 勞碌直到晚上」.
神創造的從光至地上的生物.
每做完神都說甚麼呢.
是好的.
除了做人之外 創世記並沒有這一句.
但最後在創世記第一章第31節.
「神最後都看一切所做的 都甚好」.
這句說話是甚麼意思呢.
即是神做的 神看為好的.
意思是神應當喜悅.
神應當喜悅他自己所做.
但為何詩人要在此祝願耶和華.
喜悅自己所做呢.
難道神不滿意自己的作品嗎.
我們先押下這個問題.
我們遲些會再說.
我們再看另外一個.
「詩人關乎神榮耀」的祝願.
我們留下這個.
詩人為何要祝願神喜悅自己所做呢.
我們再看下一個.
這裡說到.
「願耶和華的榮耀傳到榮元」.
這個細節我會說.
「願神的榮耀可以傳到榮元」.
因為有些人說.
神在做好六天之後.
神就安息了.
我們都知道.
聖經創世記說.
神做了六天之後就安息了.
但在創世記第二章第二節說.

$^{241}$到第七天.
神做的工已經完畢.
就在第七天.
「卸了他一切的工」安息了.
我們當然知道.
事情並不是一般.
有些人這樣說.
神做了世界之後.
就拋下了世界.
讓世界上自行運作.
自己顧自己.
神不是這樣的.
神是創造主.
神更加是維持世界的和.
祂投身在創造的世界中.
看顧,管理,做的世界.
所以我們看到.
《詩人大衛》第104篇裡.
其實不單稱頌神創造的偉大.
還深一層的體見.
隱藏在神創造大能的深處.
還包括了神維護萬有的奇妙設計和愛心.
我們可以看第10到第15節.
「耶和華使全員擁在山谷.
留在山間.
使野地的走獸有水害.
野狼得解其害.
天上的飛鳥在水旁住宿.
在枝幹間啼叫」.
祂從流角當中囂冠山嶺.
因祂作為的功效地就豐足.
祂使草生長.
比生畜吃.
使菜蔬生長.
供給人用.
使人從地裡得食物.
又得酒能悅人心.
得油能潤人面.
得糧能養人心」.
我們看到.

$^{281}$神不是在創造萬物之功完成後就安息.
他看到神仍然在無限的智慧和大愛.
繼續關顧萬有.
維護萬有.
為萬物無論大小都作了最好的安排.
最美的照顧.
所以神創造了天地萬物.
運行不息.
我們在這裡更加可以看到.
神對人很特別.
他說不單給食物和棲息地.
神更使人有酒能悅人心.
油能潤人面.
還有各樣的糧食.
能夠養人心.
很難得.
我們看到神供應和管理他所做的一切.
人和動物也向他回應.
從他的手裡得食物.
我們看第24節.
第27至30節.
詩人說.
「耶和華,你所做的何其多,都是你用智慧造成的.
全地遍滿了你所做之物」.
然後說.
「這些都仰望你按時給他們食物.
你給他們,他們就食起來.
你張手,他們就飽得未食.
你轉念,他們就驚惶.
你收回他們的氣,他們就死亡.
歸於塵土.
你差遣你的靈,他們就受造.
你使地面更換為新」.
所以在詩人大衛看來.
宇宙的王不是高不可攀的.
他乃是一個慈父.
這個宇宙不是空洞無物的.
它是一個家.
是天賦的世界.
大衛默想萬物.

$^{321}$使他的心和神的心融合在一起.
神給予自由意志給人.
讓人有自由去做自己心所喜悅的事.
其實這樣就留下了一個空隙.
讓人可以不跟隨神的心意而活.
當我們去看第三個竹苑.
「願造物之神要永存」的時候.
我們就會看到.
這裡第32節.
詩人這樣說.
「他看地,地就震動,他摸山,山就無意」.
這裡就好像在形容神在這裡憤怒.
神在這裡發怒.
其實我們看到為何詩人要竹苑耶和華.
喜悅自己所做的呢.
這句話是什麼意思呢.
因為神在整個詩篇的104篇裡面.
有幾處其實都是在形容神的.
憤怒.
神在當中是會有憤怒的.
我們去看一看.
第6到第9節.
「你用心水遮蓋地面,猶如衣裳,珠水高過山嶺,.
你的戚澤一發,水就崩逃,.
你的雷聲一發,水就崩流,.
珠山上升,珠谷下沉,歸你為他所立定之地,.
你定了界限,使水不能超越,.
不再轉回,掩沒地面.」.
大地面對神的憤怒.
他這樣說.
但又會看到水都很聽話.
他說神一斥責的時候.
神一定了界限的時候.
水就不能超越.
不能轉回.
掩沒地面.
不能過去.
其實除了水之外.
生物也會面對神的憤怒.
其實這一段我剛才已經讀過了.

$^{361}$可以繼續再讀.
第24到30節.
他說「耶和華,你所做的何其多,.
都是你用智慧造成的,.
全地遍滿了你所做之物,.
那裡有海又大又廣,.
其中有無數的動物,大小活物都有,.
那裡有船行走,.
有你所做的力威亞探游游在其中.」.
即是說這是水怪.
「這些都仰望你按時給他們食物,.
你給他們,他們便拾起來,.
你將手他們就飽得未食,.
你轉臉,他們就驚惶,.
你收回他們的氣,.
他們就死亡,歸於塵土,.
你差遣你的靈,.
他們就受造,.
你使地面更換為生.」.
所以地上的生物.
其實都是活在神的掌管.
如果神向他們發怒.
神收回的時候.
他們都會很聽的.
除了什麼呢.
我們會發覺.
因為神給了人自由意志.
人對神的回應.
會影響神對這個世界的喜悅.
神的怒氣能不能夠止息.
那就要看人是怎樣回應神.
看眾人,世上的人,以前.
我們首先要看一下自己.
所以斯人第二個祝願就是祝願自己.
然後第三個祝願才是祝願世界.
第二個祝願是關乎自己的祝願.
所以我們看到第二個祝願是.
「願蒙恩的我歌頌默念」.
斯人在第33至34節說.
「我一生要向耶和華唱詩.

$^{401}$我還活的時候要向我的神歌頌.
願他閱納我的默念.
我要因耶和華歡喜.
願蒙恩的我歌頌默念」.
所以當我們去看.
他首先說「願我要因耶和華歡喜」.
意思是「願我因耶和華歡喜」.
要回應上主的同在.
我們就要好像斯人所行一樣.
神要喜悅他所做的物.
而我們也要因神所做的物而歡喜.
「因耶和華歡喜」.
其實我們以很多不同的東西去追求.
所以我們新年的祝願.
祝人「恭喜發財」.
即是我們追求金錢.
我們又會恭祝人們步步高升.
我們希望有權貴,有地位.
又或者我們希望他們學業進步.
希望知識又或者身體健康等等.
但斯人在此說.
「願蒙恩的我的歌頌默念」.
是神能夠默念.
所以他說.
「首先要願我因耶和華而歡喜」.
當我們希望神喜悅他所做的物.
和實際神真的喜悅他所做的物的時候.
我們人也要因耶和華而歡喜.
我們因耶和華而歡喜.
然後我們就會看到.
為何我們要因耶和華而歡喜呢?.
剛才也讀過了這一段.
「使草生長 給生畜吃 使菜蔬生長 供給人用」.
「使人能從地裡得食物 得酒能悅人心」.
「得油能潤人臉 得糧能養人心」.
即是說神的供給人的比其他生物都多.
除了食物之外 還有酒.
當然人要懂得如何釀酒.
能夠叫人喜樂.
又有油 人可以提煉油出來去潤膚.

$^{441}$讓人可以滋潤.
神的供應真是無微不至.
所以我們要因耶和華而喜樂.
我們要欣賞神為他子民的作為.
欣賞神的工作.
誇讚他的智慧.
我們除了欣賞要因耶和華而喜樂之外.
我們更要一生向耶和華歌唱.
我還活的時候要向我的神歌頌.
我們要讚美神.
我們要欣賞自己去歌頌神.
當我還活的時候我要歌頌神.
其實這一句是參考詩篇第146篇第二節.
它是很澎湃的敬畏和恩喜.
我一生要讚美耶和華.
我還活著的時候要歌頌我的上帝.
所以詩人在這裡提醒我們.
我們要喜悅耶和華所做的.
我們也要歌頌神.
我們要讚美神.
然後第三個詩人要做的.
就是「願神喜悅我的物念」.
詩人在這裡禱告.
當我們去看見神所作的事情.
神所做的事情之後.
我們只可以時時去默念 去默想.
我們不單止默想神的話.
當我們看見天的時候.
當我們看見地的時候.
當我們看見太陽出來很燦爛的時候.
當我們看見雷電交加的時候.
我們都要去默念神的作為.
我們都要去讚歎神的作為.
而詩人亦都願神是悅納我的物念.
所以這裡讓我們看到.
原來詩人都期盼神喜悅他的百姓.
是欣賞他的作為.
所以我們看到.
詩人期盼做萬物的神的榮耀永遠傳下去.
詩人期盼「蒙恩我的物念」.

$^{481}$能夠讓耶和華喜悅.
然後詩人亦都會做到第三.
「願世界之上罪惡消除」.
當我們去看第三個物念的時候.
第三個願景.
「願罪人從世上消滅 越惡人歸於無尤」.
我的心你要稱頌耶和華.
Hallelujah.
當我們去看這一段的時候.
我們可能就會去問.
世界之上罪惡消除.
甚麼是令到神會有震怒呢?.
剛才最初的時候和大家看過這裡.
因為在第31,32節.
詩人說「他看地 地變震動 摸山 山就無煙」.
這是B.
然後相對的B'就是說.
「願罪人從世上消滅 越惡人歸於無尤」.
那甚麼是罪人 甚麼是惡人呢?.
當我們去看的時候.
罪人 惡人是甚麼呢?.
我們就會問.
究竟神會不喜悅甚麼呢?.
罪人和惡人是甚麼呢?.
世上其他的人.
罪人 惡人.
特別是那些不認識神的人.
那些不能夠一起以感恩來回應神的人.
就好像第二個祝願.
詩人會期盼認識神的人.
就會默念神的偉大.
當他想到其他不懂得向神感恩.
他們就是罪人.
那些不將上主看為他們一切物質豐富來源的人.
是與恩典無份的人.
他們無權享有神的恩典.
因此在神的地上也沒有住處.
但我們卻知道.
神是愛世人的.
神愛世人更加願意世上所有的人.

$^{521}$都可以敬畏耶和華.
福氣遍及萬國萬族萬民.
我們看這裡.
詩人並不是叫人去殺掉罪人.
殺掉惡人.
當耶穌也做過一個比喻.
有杯子 墨紙的比喻.
當祂很初的時候將杯子除去.
祂讓他們繼續留下.
讓他們有機會.
我們知道杯子本來不會變成墨紙.
但在神的恩典下.
杯子有可能變成墨紙.
我們如何可以消除罪惡呢?.
不知道大家有沒有記得.
有一個成語故事.
「周處除三害」.
話說周處是一個義兄的人.
他年少的時候很凶強.
他很豪俠 很豪爽.
但鄉里覺得他不太行.
他經常都在做各種事.
所以就將周處.
另外他們當中有一條蛟龍.
山上有個白鴨虎.
就叫做三害.
就是周處 蛟龍 白鴨虎.
周處自己當然不知道自己是其中一害.
他說既然這裡有一些害.
影響著這裡.
於是他就去山上殺了那隻白鴨虎.
然後他就去水裡和那條蛟龍搏鬥.
但他打得很激烈.
三天三夜都看不到他.
那些鄉民就會說.
看不到他?.
那就是兩個可能性.
一是周處死了.
一是蛟龍死了.
無論如何都是除多一害.

$^{561}$於是他們就在慶祝 很開心.
然後周處就三天三夜.
殺了那條蛟龍之後.
他出水的時候才發覺.
為什麼那些人會慶祝呢.
那些人原來當他都是其中一個惡人.
於是他就知道了.
於是他就改過.
變成他以後就不再哈哈巴巴了.
結果周處就除三害.
當罪人被除滅.
當罪人悔過成為義人的時候.
罪人就可以在世上消滅.
這個就是詩人第三個祝願.
願罪人從世上消滅.
願惡人歸於無尤.
所以他是願罪人在世上消滅.
信主悔改就可以得生.
但是罪人如何可以悔改呢.
所以當我們去問第三個問題的話.
我可以如何做呢.
我如何可以去讓罪人從世上消滅.
願惡人歸於無尤呢.
弟兄姊妹.
其實耶穌基督給我們這個大使命.
叫我們去做見證.
叫我們去傳揚福音.
讓更加多人去認識父上帝.
一人信主.
世上就少了一個罪人.
罪惡就消除一竭.
所以當我們今天去祝願.
去願造物的神榮耀永存.
更加願蒙恩的我的默念得到紀念.
更加願世界之上罪惡能消除的時候.
我是可以做到一些事的.
當我們要去將這個福音傳開的時候.
讓罪人他們可以悔改歸向神的時候.
罪惡就可以從世上消滅.
這個世界就真的成為一個完美的天賦的世界.

$^{601}$這個世界就一直是讓耶和華喜悅的世界.
我最近講到經常用這個來結束.
於我何干.
或者我應該改做於你何干.
當你聽了大概講了三十五分鐘之後.
你會怎樣去讓你的祝願能夠達到呢.
不是單單說恭喜你新年快樂.
就真的好像Lucy那樣問.
你是不是講了就是了呢.
講了就自然會的.
你能不能夠保證.
原來有些事我們都可以有些在當中有所幫助.
願大家都能夠好好的在今年裡面.
就好像大家講的傳揚福音.
好好的去讓地上高的人能夠少一個罪人.
世界就更加趨向更加的完美.
我們低頭而圖.
天父上帝求你使用感宣家的弟兄姊妹.
幫助大家.
真的可以在我們彼此祝賀的時候.
我們都想到原來我們自己都可以在祝賀上有份.
我們就願主親自去幫助大靈.
聽我們的祈禱.
奉耶穌基督的名求.
阿們.
\newpage



\section{詩篇 11:1-7}
\label{sec:sY5rw9BFGjc}
\textbf{2023-02-19 午堂《義路》詩篇 11\_1-7 講員:吳景霖傳道}
\newline
\newline
連結: \href{https://youtube.com/watch?v=sY5rw9BFGjc}{\texttt{https://youtube.com/watch?v=sY5rw9BFGjc}} ~~~~ 語音日期: 2023-02-21
\newline
\newline
\hyperref[sec:oRiBGTYbBfo]{\small{< < < PREV SERMON < < <}}
~
\hyperref[sec:index_chronic]{\small{[返順時目]}}
~
\hyperref[sec:index_scriptual]{\small{[返順卷目]}}
~
\hyperref[sec:jaVgwb2h0Gg]{\small{> > > NEXT SERMON > > >}}
\newline
\newline
詩篇 11:1-7
\newline
\begin{longtable}{cl}
\hline
\hline
章節 & 經文 (和合本修訂版)\\
\hline
11:1 & \begin{tabularx}{0.7\textwidth}{X} 我投靠耶和華；你們怎麼對我說：「你當像鳥逃到你們的山去； \end{tabularx} \\ \\ \relax
11:2 & \begin{tabularx}{0.7\textwidth}{X} 看哪，惡人彎弓，把箭搭在弦上，要在暗中射那心裡正直的人。 \end{tabularx} \\ \\ \relax
11:3 & \begin{tabularx}{0.7\textwidth}{X} 根基若毀壞，義人還能做甚麼呢？」 \end{tabularx} \\ \\ \relax
11:4 & \begin{tabularx}{0.7\textwidth}{X} 耶和華在他的聖殿裡，耶和華在天上的寶座上；他的眼睛察看，他的眼目察驗世人。 \end{tabularx} \\ \\ \relax
11:5 & \begin{tabularx}{0.7\textwidth}{X} 耶和華考驗義人；惟有惡人和喜愛暴力的人，他心裡恨惡。 \end{tabularx} \\ \\ \relax
11:6 & \begin{tabularx}{0.7\textwidth}{X} 他要向惡人密佈羅網，烈火、硫磺、熱風作他們杯中的份。 \end{tabularx} \\ \\ \relax
11:7 & \begin{tabularx}{0.7\textwidth}{X} 因為耶和華是公義的，他喜愛義行，正直人必得見他的面。 \end{tabularx} \\ \\
[1ex]
\hline
\hline
\end{longtable}
$^{1}$各位「頂」姐妹主裡面平安.
很久沒有試過在台上脫口罩.
有些可能是剛剛回來.
這幾年回來的那些可能連我的樣子都不知道是怎樣.
今早我突然間想起都先剃鬍子.
如果不是由頭救臉.
沒有剃鬍子見大家不是很好意思.
我們今天講的是詩篇.
詩篇是一首詩.
跟大家分享一首詩.
這首詩叫《折腐蔭》.
君知妾有夫,贈妾雙明珠,.
敢君前眠矣,繫在紅羅月..
妾家高樓連院起,良人執擊明光里..
知君用心如日月,視夫誓兒童生死..
還君明珠相淚垂,恨不相逢未嫁時..
可能最後那句就比較熟悉.
這首詩是講什麼呢?.
其實按字面解都解到.
有一個人竟然知道這位太太有老公.
但是都送雙明珠,可能是兩顆很貴的珍珠給她.
那個女士知道他有意思.
也都是心領的意思.
但是我家高樓連院起,良人執擊明光里.
即是她老公負責拿武器.
明光里其實就是皇宮裡面的地方.
即是皇宮裡面負責守衛.
都相當之高位.
知君用心如日月.
都讚揚那位男士這麼有心.
但是她已經決定和老公同生共死.
於是就含淚還兩顆明珠給他.
恨不相逢未嫁時.
即是說如果早點遇到就好了.
應該是客氣說話.
這首詩如果按字面解了之後就完成了.
就沒有事了.
有時候我們看聖經都是.
詩篇其實很害怕的.
穆斯比講題就幾節詩句.

$^{41}$然後叫我說半個小時都不知道怎麼辦.
但是這首詩背後其實都有段故事.
作者叫張直.
是唐朝時代的一些文官.
當時唐朝是一個群同割據的時代.
有很多人都是霸山頭.
都是招兵買馬.
希望有多一些的文官武官去到他身邊追隨他.
其實就是有意造反.
當時有一個叫李思道的人.
他就去招攬張直.
而張直就作了這首詩之後送給他.
大家明白背後的意思是什麼嗎.
就是告訴他我是不會背叛朝廷的.
我沒有想過成為你手下的人.
去到你的旗下.
你會發覺原來你知道背後的故事的時候.
這首詩是不同解法的.
是完全兩件事來的.
所以我們看聖經的時候.
我最怕有些弟兄姊妹完全跟字面.
拿個字面出來就跟我爭辯.
就說聖經是這樣寫是這樣解.
但是你不理背後的背景.
或者背後發生什麼事.
你只是按字面解.
很多時候的東西你解得通.
那首詩解得通.
但是是不是這樣的意思呢.
聖經都說了不讓女人講道.
教會的姐妹回來要蒙頭.
你們有多少人沒有蒙頭.
我們都有女牧師.
我們不是罪大滔天.
有些東西不是可以忽略背景去解.
我想大家明白.
用了這個做一個開場.
其實就是想大家去看看.
當我們去看詩篇十一篇的時候.
到底背景是什麼呢.

$^{81}$詩篇十一篇的背景其實就傳統.
我們相信是大衛王給他的兒子亞沙隆追殺的時候所寫的.
大衛是誰呢.
大衛是以色列的第二任君王.
是很出名的一個君王.
他在以色列史上.
聖經裡面都是非常重要.
他的地位很高.
甚至耶穌基督都被稱為大衛的子孫.
是很重要的人物.
是很高地位的.
撒母耳基和歷代志下.
都有很多記載是關於大衛王的.
他是一個很有領導才能的人.
也是一個可以說是一個象徵性的人物.
很重要.
但是他的兒子亞沙隆卻是想推翻他.
為什麼呢.
其實事件的起因.
就是亞沙隆有一個妹妹叫做他瑪.
但是她被他同父異母的兄弟暗戀強暴了.
這件事你可以說是一個家庭糾紛.
但其實也是一個皇帝皇室裡面的事情.
當大衛知道了之後.
你猜他有什麼反應呢.
大衛知道了之後非常憤怒.
但是他沒有做事.
於是這件事就成為了往後的事情的導火線.
亞沙隆懷恨在心.
兩年之後他設計殺死了暗戀.
殺了他的兄弟.
自然不可以當沒事發生.
於是就被迫要逃亡.
要逃走.
但是大衛其實也很愛惜亞沙隆.
最後是被亞沙隆再回來.
這件事就平息了.
但是這樣就沒事了嗎.
不是的.
亞沙隆似乎氣都還沒有順.

$^{121}$他繼續招兵買馬.
他的勢力越來越大.
去到最後連他爸爸大衛都怕了他要逃亡.
因為很可能接下來就會政變.
其實這一切的發生.
我看過電視劇《攻心計》.
看下去更誇張的劇情.
所以大家要追一追.
看回《聖經》.
其實整件事.
你會發覺基於大衛作為一個皇帝也好.
作為一個父親也好.
他沒有好好處理到他的家事.
或者他國家裡面發生的事.
而導致一些這樣的情況.
而亞沙隆.
你不要以為他為了妹妹報仇.
他就是好人.
其實整件事裡面.
你也看到他沒有什麼道德可言.
他不敬畏上帝.
也不敬畏他的爸爸.
他為了權利.
他什麼都做得出.
所以某程度上他可以不折不扣是一個惡人.
在大衛被亞沙隆追殺的時候.
他就寫了這仇詩.
第一節說到.
我投靠耶和華.
你們怎麼對我說.
你當將鳥屠到山去.
大衛在危險之中.
他依然是有信心.
他是信靠耶和華的.
他是投靠耶和華的.
他的說法.
我現在不是逃難.
我是投靠耶和華.
但是有一個人跟他說.
有可能真的有一個這樣的人.

$^{161}$也有可能這是文學的表達手法.
自言自語.
大家知道這首詩.
一首詩就有一點創作在裡面.
他說.
你當將鳥屠到你們的山去.
用廣東話怎麼說.
雞咁腳咁走.
當將鳥屠到山去.
你就是要雞咁腳咁走.
為什麼要雞咁腳咁走.
看吶.
惡人頑弓.
把箭踏在弦上.
要暗中射那.
心裡正直的人.
弓拉得的時候.
把箭放上去.
接著就要射了.
你不要看我這樣.
其實我以前教過一下射箭.
太久沒有做運動.
你拉了弓.
你準備射.
你不能hold很久.
hold得很久手腕就一直開始震.
所以拉弓.
接著下一步就射了.
就是描述那種.
如箭在弦.
一觸即發.
轉個頭就會發生.
這樣的情況.
當然要雞咁腳咁走.
根基若毀壞.
義人還能做什麼呢.
大概意思就是說服大衛要走.
用廣東話就是.
留得青山在.
哪怕無柴燒.

$^{201}$你留住青山.
將來自然就還有機會翻身.
根基若毀壞.
就是沒有了.
還有什麼可以做呢.
這裡提到正直人.
義人正直人.
都是指向大衛.
惡人就指向亞沙隆.
其實聖經裡面關於義人.
這個概念.
義人這個概念是很重要的.
那些是用來描述.
那個人是遵守上帝的誡命.
他在品德行為方面.
都是正直的公平的.
在上帝眼中是什麼呢.
是合理的.
是合義的.
是合乎道德的.
三個合不是三合一.
合理合義合乎道德的.
三個合.
這一種人是上帝祝福的義人.
他們過著幸福的生活.
他們對上帝是忠誠.
是正直.
是其他人的榜樣.
他們的生活裡面.
充滿了公義和真理.
一看到他們就知道.
受到上帝的祝福.
是好事來的.
是值得尊重的.
其實聖經裡面.
都很多關於義人的記載.
例如《箴言》.
但義人的路好像黎明的光.
越照越明直到正午.
惡人的道路幽暗.

$^{241}$自己不知因何跌倒.
將公義人和惡人去比較.
義人行為純正.
他後代的子孫有福了.
肯切求善的.
就求得恩寵.
但求惡的.
惡必臨到他.
依靠財富的.
自己必跌倒.
義人必興旺如綠葉.
這裡又是說義人一定會興旺的.
好的.
不義之財毫無益處.
唯有公義能救人脫離死亡.
這裡說到不義之財.
即是不義的財富是沒有用的.
但是財富是否真的完全不好呢.
再看詩篇.
Hallelujah 敬畏耶和華.
甚喜愛他命令的.
這人有福了.
他的後裔在世必強盛.
正直人的後代必蒙福.
他的家中有金銀財寶.
他的義行傳到永遠.
不是說金銀財寶不好.
金銀財寶都是上帝的祝福.
問題是到底那些金銀財寶是不義還是公義.
新約也有說.
羅馬書一句.
我們一定熟悉的.
義人必因信得生.
義人是因著他的信而得到上帝的祝福.
他永生.
這些經文其實是表達義人是怎樣的人.
彷彿是這個世界上最高的標準.
他們是上帝眼中合.
三個合.
合而合乎.

$^{281}$合理和合乎上帝的道德標準.
得到上帝讚許和欣賞的人.
但是怎樣才能夠合乎上帝的標準呢.
第四節說到.
耶和華在他的殿中.
耶和華在天上的寶座.
所有他的眼睛察看.
他的眼目察驗世人.
上帝是會看著世人.
即是上帝監視著你.
上帝看著你是甚麼呢.
他是為了去看看你到底是否守他的公義.
到底你所做的事是否合乎他的心意.
我們要說到公義.
義人的標準.
一定要回到上帝那裡.
一定不是人的標準.
人的標準是會不停地去改變.
剛才說到唐朝的詩人.
唐朝另一樣令人想起的就是.
他們的審美眼光.
當時圓潤的才叫做漂亮.
大份幾乎看不到頸的那些.
那就叫做漂亮.
瘦猛猛的就不受歡迎.
但是來到現代不是的.
瘦猛猛的腰很幼.
臉很小.
臉尖尖.
那些才叫做漂亮.
美的標準是會不同的.
又或者是人世間的一些標準都會變的.
以前大陸那裡要生很多小朋友才叫做好的.
所以我有個姑媽就生了七個小朋友.
我有七個表哥表姐.
但是後來就不是了.
一孩政策.
要生一個.
所以我那七個表哥表姐.
全部下面都只有一個.

$^{321}$然後去到現在.
原來早幾年有個三孩政策.
又要生三個了.
你會發覺標準是不停地在變的.
你不知道哪些標準才是對的.
有沒有一個永恆的標準.
又或者是大家都明白.
近這幾年.
有些事情在以前做是沒有問題的.
但是到了今天去到做.
你會發覺原來已經是不可以的了.
是不許可的了.
所以說我們要尋求公義.
我們要尋求正義.
我們尋求這個標準是甚麼呢.
我們只能夠回到上帝那裡.
因為人的標準.
人的所謂的公義.
人的看法是會不停地改變的.
但是上帝的標準又是怎樣呢.
上帝是否定了一些規矩出來.
我們跟著去守到完全足夠就算成功呢.
如果是這樣的話.
我們看回舊約的律法.
你跟著做就行了.
你守安息日甚麼都不可以做.
或者是你不要吃不潔的食物.
記得不要吃豬肉.
不要有血的.
那就糟糕了.
我最喜歡吃叉燒的.
牛扒那些通常都有血的.
如果你跟足.
你守回舊約的規矩.
是否就可以讓你是公義呢.
你守回那些規矩.
是否就可以呢.
其實上帝的標準是很特別的.
義和信是分不開的.
一定是關乎信心的.

$^{361}$《雅各書》有說.
亞伯拉罕信了神.
這就算他為義.
這個是《雅各書》說的.
另外《希伯來書》也有說.
我還要說甚麼呢.
若要細說基甸巴拉參宣耶伏他大位.
殺母耳和眾先知的事.
時間就不夠了.
他們藉著信.
制服了敵國.
行了公義.
得了應許.
堵住獅子的口.
滅了烈火的威力.
在鋒利的刀劍下逃生.
從軟弱變為剛強.
征戰中顯出勇猛.
打退外邦的全軍.
這些歌詞也有得唱.
因著信就是在這首歌.
做基礎作出來的.
你會發覺有很多在聖經裡面的人物.
被稱為信心的榜樣.
鬼人.
稱他們為公義.
但是你看看名字.
是不是真的很公義.
這班人平時做的事.
是不是真的很公義呢.
亞伯拉罕.
信心之父.
他也是榜樣.
但是他也試過兩次不認自己老婆.
搞到他老婆被人搶了.
我也奇怪.
其實他老婆應該七八十歲.
還可以被人搶走去做老婆.
應該真的很漂亮才會這樣.
七八十歲也很漂亮.

$^{401}$但是情況的重點是.
亞伯拉罕他沒有誠實地承認那是自己老婆.
為了自己的安全而去到隱瞞了一些事實.
你說他是不是很公義.
又或者是忌顛.
忌顛他晚年其實不是很好.
他晚年有七十個子女.
為什麼會有七十個.
因為他有很多很多的老婆.
他有很多很多的老婆.
他的兒子耶律巴力.
你看回士師記的記載.
是一個很差的人.
耶律巴力.
他是一個很差的人.
為什麼叫耶律巴力.
他的意思就是.
我爸爸是王.
你會不會和你的兒子改一個這樣的名字.
你爸爸是王.
暗示我你爸爸很厲害.
不會的.
你不會這樣去改一個名字.
其實某程度上他就是想做皇帝.
他有很多的妃嬪.
很多的妾侍.
很多的子女.
他其實一點都不是一個好人.
或者是參孫.
參孫就出了名.
出了名是壞孩子.
他在聖經你見到他做很多壞事.
唯一一件好事就是他臨死的時候.
和那些非利士人同歸於盡.
還有大衛.
大衛是不是二人.
剛才說過大衛是歷史上.
以色列史上一個這麼重要的人.
他最出名的是什麼呢.
就是找塊石頭扔死了哥利亞.

$^{441}$最出名的就是這件事.
第二出名的就是裝人洗澡.
他裝人洗澡.
然後去搞人家老婆.
到最後他怕通槍事發.
就推了他老公去死.
最終就名正言順搶了人家老婆.
在現代來說.
這個簡直是不能接受.
這樣的大衛.
你還是二人.
還是耶穌基督的祖先.
怎麼搞.
所以你明白.
你在聖經裡面聽到保羅說.
沒有二人連一個也沒有.
如果按回聖經.
剛才我說的標準.
要是合理的.
要是公義的.
要是合乎上帝的心意的.
這班人沒有一個可以.
真的沒有二人一個都沒有.
但是有沒有自相矛盾.
既然是這樣.
為什麼又可以稱他們為二人呢.
因為在上帝眼中.
義和信基本上是分不開的.
一個二人.
他同時間要有信心.
他要有信仰.
雖然大衛他這麼壞.
他做過很多壞事.
他有很多地方是不足的.
是不好的.
但是他也有很重要的一件事.
就是他會悔改.
他對上帝有一種不離不棄的信心.
他依然相信當自己悔改.
當自己知罪.

$^{481}$當自己認罪的時候.
回到過去的時候.
上帝是依然在.
他依然會對他不離不棄.
因此上帝才會繼續祝福他.
稱他為二.
令他在聖經裡面被記載.
令他成為耶穌基督的一個祖先.
否則其實他就像掃羅一樣.
上帝可以高納一個人.
也可以廢棄一個人.
大衛就是因為他的信.
他能夠回到過去.
我想跟大家說的是.
原來一個二人並不是你可能認為的.
一輩子沒有犯過錯.
虧行舉報.
一點壞事都不做的.
螞蟻都不踩死一隻.
紅燈都不會過馬路的.
是不是這樣就叫做一個二人.
有時候我們以為你沒有做壞事.
你沒有做錯事就叫做二人.
但其實這個世界是很複雜的.
你可以是沒有犯法但你犯罪.
也可以是你犯了法但你沒有犯罪.
這個世界是這樣的.
其實不是說我們沒有做過壞事.
虧行舉報就叫做二人.
相反是我們看看我們的信仰.
我們到底有沒有做錯事之後.
願意回到上帝那裡回轉.
向祂認罪悔改.
你向祂認罪悔改.
才是你能夠重新回到這條二路上的一個方法.
當然不要走到另一個極端.
覺得既然是這樣.
我做什麼壞事都做盡都無所謂.
回家認罪就行了.
平時在公司那裡.

$^{521}$什麼都睜一隻眼閉一隻眼.
說謊也好老闆叫我做什麼都好.
不要緊的回到家認罪.
這樣行不行.
這樣是不行的.
這也不是一種真是有悔意的態度.
但我想說的是.
我們不可以覺得我一生人沒有犯錯.
沒有犯罪虧行舉報才叫做公義.
有些人他很有罪疚感.
我犯了罪我犯了錯.
因此會覺得上帝不是很理我的.
或者我都不知道我自己有沒有得救的確據.
我死了之後能不能上天堂.
我真的不知道.
有時候是這樣.
但我想告訴大家.
一個人人誰無過.
你犯了錯犯了罪不要緊.
最重要是你願不願意悔改認罪.
回到上帝那裡.
願不願意悔改認罪.
回到上帝那裡.
因此我們需要的是有些東西真的要學習.
你不會無端端知罪的.
你不會無端端知道自己這樣不對.
你一定要從聖經裡面找到一些標準.
因此有四點幫助我們去學習.
怎樣成為一個二人.
或者學習聖經的榜樣.
第一就是看聖經.
很多時候我們弟兄姊妹.
真的很依賴崇拜的那一堂道.
一個星期裡面回來的一個小時.
聽了半個小時道.
那就叫做OK的了.
然後剩下的六天23個半小時.
就不知道做了什麼了.
不看聖經了.
什麼都不理了.

$^{561}$回來聽了一堂道.
就叫做學了.
其實不足夠.
你真的應該要自己多看聖經.
其實有很多東西可以幫你.
你安裝手機APP.
它自己會響的.
提醒你要靈修.
每天看一點點.
或者是有些網站.
它已經有注釋了.
你看到經文旁邊有注釋.
講了背景.
已經幫了你很多了.
不會說只是看字面.
然後好像求籤一樣.
有一段經文.
自己解一下.
或者是找曾牧師當他解籤佬.
這段怎麼解呢.
解完之後原來是這樣.
不是的.
你自己回去學習.
去識經.
去知道上帝的心意是怎樣.
他的標準是怎樣.
第二就是要悔改.
要知罪.
當你去審視自己的日常生活.
你平時做慣的事也好.
或者你做的事也好.
你想一下自己其實有沒有錯.
有沒有罪.
有沒有一些事是做了其實是不對的.
有些事可能你做了幾十年.
你覺得沒有問題的.
但是原來聖經的標準告訴你是有問題的.
是不是.
可能是的.
你不知道的.

$^{601}$或者是坦白說.
可能有時候做生意.
有些手法.
你做了幾十年都是這樣.
但是原來其實聖經的標準不是這樣.
那你怎樣呢.
你繼續還是你知罪悔改呢.
這些是我們需要去經常審視的.
我最擔心的就是有些弟兄姊妹.
他們是完全覺得自己沒有犯法就行了.
我合法的就沒有罪.
法和罪是兩個概念.
有很多事是上帝眼中不好.
是犯了罪你不知道的.
可能是沒有愛心.
沒有同情心.
沒有憐憫人的心.
隨便的說謊或者是不肯去原諒人.
不肯給機會人.
或者是對你的員工.
或者對你的外傭非常之刻薄.
我跟足勞工法出足這麼多工資.
沒問題.
洗到盡都可以.
其實很多事是我們整個社會這樣運作.
社會這樣運作不代表它對.
你有沒有去審視一下.
我在上帝眼中我做這些事.
正不正確.
有沒有犯罪.
這件事很重要.
第三就是維持聖潔.
你要保持自己是遠離邪惡的.
是要遠離邪惡的.
千萬不要高估自己的定力.
或者高估自己的意志力.
很多時候就是衰在那些地方.
那些入賭場賭幾百元的那些.
輸光就收手的那些.
通常是輸身家的.

$^{641}$不要高估自己.
有些朋友貪一些職業.
明知那些職業裡面的人是三教九流.
或者是很多很壞的人在裡面.
但是可能貪他人工高福利好.
或者貪他不用怎麼讀書就可以做.
進去後.
他開頭跟我說.
不會的 我不會跟那些人一起.
轉過頭幾個月都沒有再見到他回教會.
不要高估自己的定力.
邪惡是要逃避的.
為什麼教我們不站罪人的道路.
不坐涉萬人的座位.
你是要避開那些人的.
你不要黏在一起.
千萬不要高估自己.
你要遠離邪惡.
還有第四點.
我覺得最容易被人忽略的.
就是你要為他人而去服務.
你要去做一些事.
你要做一些事是顯示到.
你那種與眾不同.
你是有愛心的.
你是基督徒的.
你要做一些事.
其實聖經裡面的教導.
通常都是跟人有關的.
我們經常誤以為.
我們靈修是什麼.
就是找個房間躲起來.
然後看聖經.
然後祈禱.
就是一種修行.
然後你的靈命就會日日進步.
好像是一個很私人的事情.
但是聖經裡面的教導不是這樣的.
你要做的.
最大的誡命是什麼.

$^{681}$愛神 愛人如己.
你要有個人.
你才可以愛人如己.
聖經不是教你將自己收起來.
而是你要對愛.
很多事情就是上帝看你.
你怎樣去接待人.
你怎樣去對人.
這些事情你做在一個最小的弟兄身上.
就是做在耶穌身上.
是不是.
聖經不是教你封閉自己.
在一個暗室裡面修煉.
你就會升仙.
而是教你去到社群當中.
你要愛人.
你要做事.
你要去到彰顯基督的愛.
換句話說.
其實作為一個義人.
不是說看你沒有做些什麼.
我沒有犯罪.
我沒有犯法.
我沒有做壞事.
我沒有大奸大惡.
就叫義人.
相反是看你有做過些什麼.
你有沒有做一些合神心意的事.
你有沒有做回一個基督徒應該要做的事.
上帝給你的資源.
你有沒有用來幫助別人.
有沒有關顧弱勢社群.
有沒有關顧這個社會.
有沒有用你的行動去彰顯上帝的國.
這些才叫義人.
到最後我們又要想了.
做義人為什麼呢.
做義人有沒有好處.
沒有好處.
為什麼叫我做義人呢.

$^{721}$我們看第五到七節.
耀護說考驗義人.
惟有惡人和喜愛暴力的人.
他心裡恨惡.
他要向惡人密報羅莽.
烈火流黃.
熱風作他們杯中的糞.
因為耶和華是公義的.
他喜愛義行.
貞直人必得見他的面.
其實某程度.
我自己覺得.
你作為一個義人不是說有什麼好處.
最起碼就是說你作為一個惡人.
你是一定有壞處.
上帝是會向惡人追討他的罪.
上帝是會向惡人去追討他們生前所犯的.
或者可能不是生前.
可能現世就已經會.
上帝向他追討他的罪.
或者至少死後.
會有烈火流黃熱風這些.
其實就是有一種末世.
或者將來在地獄裡面的一種懲罰的暗示.
所以要相信上帝一定會向惡人追討他的罪.
而我們作為基督徒.
我們今天能夠得到這個救恩.
我們能夠將來不需要去受那些永恆的苦.
而是能夠在天國和上帝有一個永恆的相聚.
永恆的福樂.
其實這個已經是最大的福份.
你沒有理由還說既然是這樣.
我在上帝眼中成為一個義人的時候.
你是不是應該祝福我.
有病就醫病.
沒有病就身體健康.
或者一帆風順.
所有投資都賺的.
走在街上總之一盞燈都是綠燈.
你不會期望上帝是這樣去祝福你的.

$^{761}$你不會覺得上帝會違反了整個世界運作的規律.
就來祝福你.
因為你相信了耶穌.
其實有時候不少情況我們是這樣期望的.
所以當我們作為一個義人.
我想和大家說的就是.
你本身你信耶穌.
上帝拯救你.
你去做合祂心意的事.
成為一個義人是理所當然的.
是應該的.
不過上帝有沒有祝福.
有的.
剛才我和大家說了很多經文.
上帝都是會祝福你的.
你成為一個義人上帝是會賜福你的.
不過那個不是你的目標.
你不是因為我要上帝祝福.
所以我成為義人.
而是你已經信了耶穌.
你願意成為一個義人的時候.
上帝會祝福你.
邏輯不要反過來.
不要以結果為目標.
你不是為了得到祝福.
然後我去做義人.
還有義人其實不是保證你一定免於苦難.
聖經甚至很清楚地告訴我們.
義人是會被考驗的.
義人是會面對苦難的.
例如《雅各書》說到.
我的弟兄們,你們遭受各種試煉時,都要認為是大喜樂,.
因為知道你們的信心經過考驗就生忍耐,.
但要讓忍耐發揮完全的功用,使你們又完全又完整,一無所缺..
神同樣會考驗義人.
甚至都不會手輕.
我想大家知道.
約伯就是一個好例子.
約伯很慘的.
約伯就是所有的財產都被上帝摧毀了.

$^{801}$在試驗中.
所有的子女都沒有了.
連自己都全身生瘡了.
什麼好東西都被上帝拿走了.
剩下老婆.
老婆是不是一個苦難呢?我不知道.
說笑而已.
什麼都沒有了.
還有幾個朋友.
約伯不知為什麼.
他那三個朋友提供了很多答案給他.
但大家都不知道為什麼.
去到最後上帝都親自出來跟他說話.
上帝都沒有告訴他為什麼.
他只告訴他.
我就是創天造地的神.
你們人是不明白很多東西.
然後約伯就在六灰中懊悔.
約伯心悅誠服.
我想大家知道其實有很多東西.
我們在這個世界上會面對不同的考驗.
或者很多東西我不明白.
很多東西會覺得我信了耶穌這麼多年.
為什麼我會這樣面對.
很多東西我都已經是.
甚至我都會.
我出來做了傳道.
我面對的困難.
為什麼上帝不幫我.
幫我拿走呢?.
是不是?很真實的.
有一次我兒子發高燒.
晚上我在醫院裡面陪著他.
在病床旁邊的時候.
我真的這樣問的.
我都已經整個人放了出來幫你做事.
我只是想我兒子退燒.
為什麼這樣都不幫我.
人就是會問這個問題.
是不是?.

$^{841}$很真實的信仰的掙扎.
但經過這些掙扎之後.
你會發覺你的信心是會堅強了.
你的信仰是會紮實了.
這些東西是會讓我們變得更加完全.
更加完整.
所以上帝是會考驗二人的.
而我們要面對這個考驗的時候.
其實你要相信上帝最終是在祝福你.
是在幫你.
他上帝最愛的是你.
今天的講題我用的二路.
路這個概念.
我想大家明白.
我們作為一個信仰裡面.
我們想成為一個二人.
不是說你一出生就是二人.
或者你一出生就是惡人.
或者你的本質就是二人.
你的本質就是惡人.
不是你的本質問題.
而是你有沒有做一個選擇的問題.
你有沒有走在那條二路上.
用一個路的概念就是.
一條路是有終點.
終點就在上帝那裡.
而你走這條路的時候.
你是可能有時候走到這邊.
走到那邊.
有時候跌倒.
有時候失腳.
有時候做錯事.
那怎麼辦呢.
你是不是走出去之後.
哎呀 這樣就.
哎呀 我走了出去.
我這輩子就玩完了.
不是的.
你走回到正路上.
好像大衛那樣.

$^{881}$他是做錯事.
甚至在人眼中是不能原諒的.
但是他願意回去認罪悔改.
承擔他的責任.
承擔他的錯.
上帝依然愛他.
依然能夠繼續走這條二路.
所以弟兄姊妹.
我鼓勵大家.
可能過去有很多事情.
你覺得自己走錯了.
或者有些事情我做得不好.
不要緊.
你今天開始做好它.
勤力去看聖經.
認真去實踐你的信仰.
逃避那些邪惡的事.
行上帝眼中公義的事.
回去那條二路裡面.
繼續活出你的信仰.
成為世人的祝福.
我們一同去禱告.
七夕天父.
我們求你幫助我們.
其實我們每一個人都很軟弱.
我們每一個人本來都是不配.
但是因著上帝你的恩典,憐憫.
我們能夠成為你的子民.
我們求你幫助我們每一個人.
我們真的知罪.
我們知道自己犯罪.
我們知道自己在日常生活也好.
在教會裡面也好.
有些事情我們沒有做到.
我們有些事情沒有去服侍人.
我們沒有去愛人.
我們沒有為上帝你出一分力.
我們求你原諒我們.
讓我們知罪.
並且我們悔改.

$^{921}$好讓我們將來在上帝的眼中.
被你看為義的.
我們祈禱交代.
奉耶穌基督的名求.
阿們.
\newpage



\section{加拉太書 3:21-29}
\label{sec:jaVgwb2h0Gg}
\textbf{2023-02-19 加恩樓《律法之上》加拉太書 3\_21-29 講員:陳韋安牧師}
\newline
\newline
連結: \href{https://youtube.com/watch?v=jaVgwb2h0Gg}{\texttt{https://youtube.com/watch?v=jaVgwb2h0Gg}} ~~~~ 語音日期: 2023-02-22
\newline
\newline
\hyperref[sec:sY5rw9BFGjc]{\small{< < < PREV SERMON < < <}}
~
\hyperref[sec:index_chronic]{\small{[返順時目]}}
~
\hyperref[sec:index_scriptual]{\small{[返順卷目]}}
~
\hyperref[sec:v6FulXmsPK8]{\small{> > > NEXT SERMON > > >}}
\newline
\newline
加拉太書 3:21-29
\newline
\begin{longtable}{cl}
\hline
\hline
章節 & 經文 (和合本修訂版)\\
\hline
3:21 & \begin{tabularx}{0.7\textwidth}{X} 這樣，律法是與神的應許對立嗎？絕對不是！如果律法的頒佈能使人得生命，義就誠然出於律法了。 \end{tabularx} \\ \\ \relax
3:22 & \begin{tabularx}{0.7\textwidth}{X} 但聖經把萬物都圈在罪裡，為要使因信耶穌基督而來的應許歸給信的人。 \end{tabularx} \\ \\ \relax
3:23 & \begin{tabularx}{0.7\textwidth}{X} 但這「信」還未來以前，我們被看守在律法之下，像被圈住，直到那將來的「信」顯明出來。 \end{tabularx} \\ \\ \relax
3:24 & \begin{tabularx}{0.7\textwidth}{X} 這樣，律法是我們的啟蒙教師，直到基督來了，好使我們因信稱義。 \end{tabularx} \\ \\ \relax
3:25 & \begin{tabularx}{0.7\textwidth}{X} 但這「信」既然來到，我們從此就不在啟蒙教師的手下了。 \end{tabularx} \\ \\ \relax
3:26 & \begin{tabularx}{0.7\textwidth}{X} 其實，你們藉著信，在基督耶穌裡都成為神的兒女。 \end{tabularx} \\ \\ \relax
3:27 & \begin{tabularx}{0.7\textwidth}{X} 你們凡受洗歸入基督的都披戴基督了： \end{tabularx} \\ \\ \relax
3:28 & \begin{tabularx}{0.7\textwidth}{X} 不再分猶太人或希臘人，不再分為奴的自主的，不再分男的女的，因為你們在基督耶穌裡都成為一了。 \end{tabularx} \\ \\ \relax
3:29 & \begin{tabularx}{0.7\textwidth}{X} 既然你們屬於基督，你們就是亞伯拉罕的子孫，是照著應許承受產業的了。 \end{tabularx} \\ \\
[1ex]
\hline
\hline
\end{longtable}
$^{1}$各位弟妹平安,今天很高興跟大家一起崇拜.
今天我們會看加拉太書的經文.
加拉太書是一段保羅在早期的書信.
講的都是一些我們信仰裡很根本的東西.
大家都知道的,信耶穌是怎樣.
不要靠行為,我們要憑信心,欣賞菁義等等.
所以加拉太書的第三章其實都是一些很教義的事情.
作為這麼多年基督徒,大家都很清楚我們信仰的基本.
我們的信心的重要,我們不是靠律法行為等等的事情.
不過我們跟大家一起再看這段經文.
去看看加拉太書裡保羅怎樣探討到律法的事情.
一些很根本的事情,我們怎樣看待行為,律法的事情.
剛才Jason也讀了一段經文,有關律法,欣賞菁義的事情.
所以今天講經文之前,我再講一點有關欣賞菁義的信仰.
欣賞菁義是保羅在加拉太書裡的教導.
也是後來在羅馬書第三章裡的教導.
不過更重要的是,欣賞菁義更加是500年前.
馬丁路德對保羅再一次的重新演繹.
為什麼這麼說呢?.
雖然保羅在新約裡,在加拉太書和羅馬書裡.
早就提過欣賞菁義的道理.
但500年前,我們基督教,新教開始重新重視欣賞菁義的教義.
對馬丁路德來說,信心不單單是信仰的起點.
更加是信仰的全部.
所以跟天主教不同,我們將欣賞菁義看為一個很重要的救恩.
再次回到中教改革的根源,500年前新教的起點.
馬丁路德是一個天主教會的修士.
1517年改教,丟掉95條之前.
其實他早就在當時的Wittenberg大學裡教書.
教的是詩篇,羅馬書,加拉太書.
基本上他是一個聖經科老師.
當時他開始醞釀出欣賞菁義的觀念.
但還未能夠很好的形容出來.
未能說出那個字眼,就開始慢慢醞釀出來.
馬丁路德是一個很高要求的人.
一個很認真,很著緊的人.
一個經常要求自己拿A的人.
甚至乎是一個非常火爆的人.
所以今天如果找他來做錄影的話.
應該是萬梓良最適合的.

$^{41}$一個很要求的人.
所以路德對於自己的問題開始有些醒覺.
無論他怎樣做都仍然發現.
無法達到上帝的標準.
所以他發現自己原來人很本身的軟弱.
這種對於人最核心的軟弱的看見.
我們稱之為卑微神學.
他還未有欣賞菁義的道理.
去說出人就是無能為力.
人就是怎樣也不能夠達到上帝的要求.
這個卑微不是我們中國人所說的謙卑.
不是一種禮貌,一種態度.
而是人徹底地對於自己的罪.
發現自己什麼都做不到.
只能夠求上帝的幫助.
這個人對自己的罪的發現.
和一個完全無能為力的態度.
是我們基督教新教的根本.
我們最重要的強調.
人是無能為力的.
唯有上帝幫助我們.
人只能夠去到十年前.
空空的去單求上帝恩典和幫助.
人是完全什麼都做不到.
所以在這個什麼都做不到的前提下.
他就慢慢發現羅馬書第三章所說的信.
那個重要.
因為什麼都做不到.
所以我只能夠信.
信心其實就是我們什麼都做不到.
你只能夠接受,相信上帝為你做了的事.
你只能夠完全地接受他.
所以對老德來說.
信心不是一個條件.
不是因為我信,我很厲害,我信了耶穌.
所以我能夠得救.
信心不是一個交換恩典,救恩的條件.
不是因為你做了什麼,所以換了什麼回來.
而是很強烈地說我什麼都做不到.
我只能夠相信這個已經發生了的事情.

$^{81}$我接受他.
所以我們說唯獨信心.
其實是指向著一個更加重要的四個字.
就是唯獨基督.
因為基督耶穌做的.
我們只能夠去相信,接受他.
信心不是一種我們做了什麼.
來得救的方法.
如果這樣的話.
其實耶穌再厲害也沒有用.
如果我不相信,耶穌怎麼救都救不了我.
所以信心的意思不是我做了什麼.
而是反過來我什麼都做不到.
所以信心不是另一種的行為.
信心是兩手空空地接受上帝的恩典.
所以信心是帶著一種信靠.
就是trust的意味.
不是信某個故事.
信一個方橋或是四律.
那些東西就算了.
而是完全信靠上帝的恩典.
他的幫助來依靠他.
所以英信清二是更好的說法.
中文叫直信清二.
我們是instrumentally直注相信來去清二.
不是因為你信.
因為的原因永遠都是基督耶穌自己.
所以信心不是主因.
純粹是一個過程.
我們來接受那個過程.
我想這正是基督教和天主教最大的差別.
我們不覺得信心是一種行為.
信心不是一種人的功德.
信心就是沒有方法.
唯有依靠上帝提供的方法.
所以有些人說上帝出了雞我出了豉油.
其實我們連豉油都出不到.
完全都是上帝給我們的.
所以這個就是.
我想我們簡單來說說.

$^{121}$這個有關保羅或者是路德所說的.
有關信的清二的看法.
簡單說了幾分鐘神話之後.
我就開始去看這個《迦太書》第三章.
保羅在當中提到有關律法的問題.
我再說一次.
有一節.
(文字).
全一個能叫人得心的律法 而的意就呈現本乎律法了.
但聖經百眾人都圈在罪裡 使所應許的福 因信耶穌基督歸向那信的人.
但這因信得救的理 還未來之先 我們被看守在律法之下.
直圈到那將來的真道顯明出來 這樣律法是我們的糞蒙的師父.
在我們的基督那裡 使我們因信稱義 但這因信稱義的理既然來到 我們從此不再是師父的手下.
保祿說 雖然因信稱義是我們在得救的事情裡 無能為力.
任何的律法 各禮行為 都在得救的事情裡是沒有功效的.
他就問 難道律法就沒有意義嗎?難道律法就跟福音是相違背嗎?.
保祿就很清晰地回答 不是.
他說 這樣律法就與神的應許對立嗎?絕對不是.
我們就可以去想想 到底這個律法行為對我們來說有什麼意義.
如果我們說 清義不是靠律法的話 到底律法是什麼呢?.
保祿就用一個比喻來解釋 寫詞是這樣說的.
這樣律法是我們的糞蒙的師父 引我們到基督那裡 使我們因信稱義.
鑊本就用了翻譯 叫做律法是糞蒙的師父.
什麼叫糞蒙的師父呢?.
如果看回新修訂的鑊本 它就解作 律法是啟蒙教師.
其實兩個版本都不算很清楚 說得明白律法的意思.
什麼叫啟蒙的師父 什麼叫糞蒙的老師呢?.
原文就是Pedagogos的字 如果讀教學的就知道了.
Pedagogos就是教育.
但意思不是純粹高級的老師 不是教授.
在古代Pedagogos是陪小朋友讀書的奴隸.
負責在旁邊拿書寫字的 如果中文來說就是在旁邊磨墨的書僮.
幫少爺拿著豬卷的書僮.
所以更多的日本女性中文本 發展得更加積極.
叫做童年導師.
所以保羅說 律法的意義 律法是重要的.
不過是小朋友的童年導師 是那些書僮 是那些嬰兒的東西.
小朋友的東西.
什麼意思呢? 這些東西是很初級的階段.
是嬰兒學行 101的階段.

$^{161}$你問這些嬰兒的東西重要嗎? 當然重要.
你學某樣東西總是有101的階段.
某個初階階段 某些比較低級的階段.
就像小朋友學游泳一樣 你的子女學游泳.
總是要帶著一個壓制水泡.
這個當然需要 浸水不用嗎?.
小朋友學踩單車也是一樣.
開始的時候 加個輔助輪在這裡 先踩一下.
小時候學走 總是有輔助輪 讓你扶一下 推一下.
所以水泡是重要的 輔助輪是重要的 學走是重要的.
但如果一個人 男人老狗都帶著一個水泡.
Hello Kitty水泡這樣游泳的話.
是一件很醜怪的事情.
一輩子踩單車都用輔助輪.
這樣就不太好看了.
你永遠都學走 這樣就不是那回事了.
所以是這樣的一個理解.
難不成這樣? 當然不是一個醜怪問題.
甚至你去到某個階段 某個水平.
相反這些水泡 這些學走 這些輔助.
反而會阻礙你的學習.
你一直都帶著水泡游泳 學不會游泳.
經常都要抱著輔助輪 你總是學不到走單車.
所以這些律法是重要的 是好東西.
但都是一些低BBB的東西.
你長期抱著的話 這樣就有礙你的成長.
所以等於律法是一個基督徒.
符合上帝標準的一個階段.
而基督的義是另一個新的階段.
所以普通話說 律法是有用的.
但律法並不能夠長遠地幫助到一個基督徒.
去接近上帝的標準.
但你總是有一刻要去放開它.
它還是要領我們去基督那裡.
這是一個很重要的道理.
你的信仰不能夠停留在一個律法BB的階段.
這就成為了我們今天反思的問題.
你的信仰是否仍然停留在一個Hello Kitty水泡.
輔助輪子的QQ學行車的階段.
這些是好的 這些是重要的.

$^{201}$但你是不是仍然停留在一個階段.
作為一個牧師 作為一個在網上寫東西的人.
這幾年我收到很多PM.
即是inbox我的問題.
很多問題我可以總結為一個類型的問題.
就是可不可以的問題.
我收到很多可不可以的問題.
我讀出來給大家聽 因為很長.
基督徒可不可以崇拜遲到.
基督徒可不可以星期日不回教會.
基督徒可不可以不回主約學.
基督徒是否真的要受一封獻.
基督徒是否可以炒股票.
基督徒可不可以和微信主人結婚.
基督徒可不可以直發不祈禱.
基督徒可不可以不靈修.
基督徒可不可以練瑜伽.
參加萬聖節.
基督徒可不可以整容.
基督徒可不可以看鬼片.
基督徒可不可以做心理測驗.
基督徒可不可以看相.
基督徒可不可以賣樂彩.
基督徒可不可以轉教會.
基督徒可不可以看中醫.
基督徒可不可以說謊.
基督徒可不可以說白色的謊.
基督徒可不可以網上崇拜.
基督徒可不可以吃午餐.
參加網上崇拜.
基督徒可不可以吃午餐.
戴著面膜 做facial.
參加網上崇拜.
很多這些可不可以的問題.
當然這些問題都可以有答案.
我都可以回覆他.
講講聖經 講講原則.
不過在這些問題上.
我反而想問.
為什麼會有這麼多問題.

$^{241}$為什麼會問這麼多可不可以的問題.
你問基督徒可不可以吃午餐.
參加網上崇拜.
我可以的話.
那你下次可以問.
基督徒可不可以戴著指甲.
參加網上崇拜.
明白我的意思嗎.
可不可以的問題.
我會問可不可以沒有這麼多問題.
我以前都寫過.
很久以前.
當時是iPhone 6的年代.
當時香港炒iPhone炒得很厲害.
有人問我.
基督徒可不可以炒iPhone.
我就預備了一篇文章.
打算寫一篇洋洋的幾千字.
講基督徒和iPhone的關係.
但當我打完文章後.
就寫著基督徒應否炒iPhone.
我覺得這個話題很噁心.
為什麼會問到這樣.
主耶穌看到會有什麼感覺呢.
基督徒應該炒iPhone.
有人說一買一賣有何不可.
很多生意人都是基督徒.
神都像我們投資的眼光去賺錢.
用管家的比喻.
可以寫一篇很長的文章.
關乎基督徒的經濟倫理.
炒賣對社會的貢獻.
一些倫理問題等等.
為什麼我們會問這個問題.
反而我們問這個問題.
因為基督徒的倫理.
其實不是純粹律法式的可不可以.
不要問可不可以.
這樣做是否可以.
這樣做是否違法聖經.

$^{281}$這樣做是否可救.
這是一些很消極.
很被動很懶惰的做法.
長期保持一個可以過關的要求.
尋求一條不會被上帝擊殺的線路.
這樣更可救.
只要我這樣做都可以.
神都不會擊殺我.
這樣我就做了.
不斷尋求這樣的底線.
可不可以我就做了.
不可以我就試試問第二個問題.
這不是聖經所說的.
有關基督徒的倫理.
這是一個律法式的思維.
這是一個蒙訓師父的方法.
這是Hello Kitty的水泡.
是一些輔助學行車的基督徒思維.
我們如何理解我們的行為.
我們如何理解基督徒的倫理.
其實是很登山寶粉式.
不是消極地尋求最低的要求.
而是竭力去最高的水平.
因為耶穌說你們要完全.
將你們的天賦完全一樣.
不是問既然不可以殺人.
那我可不可以做些什麼.
而是問我不但不會.
更加會愛仇敵.
是這樣的一個思維方式.
所以保祿教徒我們要離開這個律法的嬰兒.
離開這個憤蒙的師父.
去到基督耶穌的義裡面.
才是基督徒應該要去持守.
去追求的倫理的榜樣.
基督徒會問離開了這些嬰兒.
去到基督耶穌那裡做.
那是什麼階段呢?.
那會是什麼事情呢?.
當然我們不是再去.

$^{321}$有人會嫌我們是恩上仇義.
所以怎樣都得救.
不關行為事.
反而我們就不會守行為.
行為也不重要.
但這不是的.
他們以為既然恩上仇義.
那我就什麼都不用管.
什麼都可以做.
這絕對不是這個意思.
但也不是那種律法式地問.
怎樣做怎樣做怎樣做問題.
但如果又不是那種恩上仇義.
什麼都可以做.
什麼都做得.
又不是那種問律法式的.
怎樣做的話.
我們怎樣去理解呢?.
這正正是我們今天再想一步的問題.
所以說如果不是律法嬰兒.
那些可不可以.
又不是什麼都可以.
怎樣才是更高層次的基督的義呢?.
我們記得在單一部分裡.
其實是這樣說過的.
耶穌這樣說.
第五章裡面說.
莫想我來要廢掉律法和先知.
我來不是要廢掉.
乃是要成全.
我實在告訴你們.
就算天地都廢去.
律法的一點也不能廢去.
也要成全.
然後他就說.
我告訴你們.
你們的義若不勝於民事和法理裁人的義.
斷不能進天國.
耶穌直接要求門徒.
你們的義如果不勝於民事和法理裁人的律法的話.

$^{361}$斷不能進天國.
這是一個非常嚴謹.
非常認真的命令.
我們問究竟門徒的義.
基督的義.
怎樣才能勝過律法嬰兒的義呢?.
其實原文裡面的勝過.
其實不是解作勝過.
勝過是解作perazo.
不是贏輸的問題.
勝過是解作多過,大過,厲害過.
都可以.
超過.
但不是贏輸不是勝過.
可以解作豐盛.
即是abundant.
豐富的意思.
over的意思.
所以我試著說明.
即是你明白的.
perazo的意思是在五餅二魚的故事裡面.
當耶穌行完五餅二魚之後.
食物多到爆.
東西滿到滿血.
十二個男子都裝不完.
然後才用perazo這個字.
去形容食物多到滿血的狀況.
這個就是perazo.
耶穌說你們的義.
要perazo.
過門徒的義.
耶穌要求門徒的義.
是一種不斷向上的超越.
不是問去到哪裡就夠.
可不可以,給不給.
而是問要去到不斷追求更加好.
更加合乎上帝的心意.
更加好的義.
所以你覺得登山寶訓的教導很難.
即是愛受的.

$^{401}$不要罵弟兄.
被人打完右邊再打左邊.
很難.
如果是這樣的話.
其實你就誤會了耶穌的意思.
如果你以為登山寶訓所說的教導.
是一種更加嚴苛,更加高的.
可不可以律法的時候.
其實你就誤會了耶穌的意思.
耶穌不是要做壞事.
不是說他們叫你不可以兼任.
我連你想都不讓你想.
別人叫你愛人如己.
我連你連受都愛敗.
他們叫你不可殺人.
我連你連生氣都不讓生氣.
如果耶穌的義只不過是更加嚴苛.
更加高要求的可不可以的話.
耶穌只不過是更加高要求的法利賽人.
耶穌不是這樣.
耶穌是想我們去超越.
不斷地超越.
不是找一條得救的底線就過關.
所以耶穌的義和法利賽人的義.
分別不是一個在Level 5.
一個在Level 100.
重點不是Level 100.
重點是你不要以為Level 5就能過關.
只要我輕鬆到Level 5的話.
那我就算是OK了.
就算是基督徒也OK了.
重點不是Level 幾.
而是不要以為Level 幾.
就是那個要求.
這正正就是文獻法場的問題.
他們不是不好.
他們都是做好事.
但他們以為那些好事做到某個位置.
就算是OK了.
那我就下班了.

$^{441}$主啊我一個星期禁食兩次.
凡我所得都捐十分之一.
這個思維.
如果你以為耶穌的義.
只不過是更加高的律法的話.
你就當是一種這樣的.
輕鬆了要求的話.
其實都是沒有分別的.
耶穌說凡向大興動勞的.
都要受審判.
那我就不生氣了.
我就暗算他.
耶穌說要你多走兩公里路.
那我陪你走兩公里路.
那我就馬上就走了.
耶穌說要你受禮禱告.
那我奉主名求安門.
那你就糟糕了.
不是這樣的.
耶穌叫我們不斷的去Paraiso.
不要覺得某個要求做了就算是收工.
不斷的超越.
不斷的去超越他.
可能大家會覺得.
是不是聽過這個道理.
免費的東西是最貴的.
耶穌基督的恩訊清義是免費的.
不過卻是耶穌負上沉重的代價.
都要求你去嘗試.
因為這個免費的東西.
是要你負上自己最大最大的所有.
沒有人逼你的.
沒有人監你的.
你不回來崇拜.
不會被上帝擊殺的.
你不靈修.
你不祈禱.
你發覺.
都是不會影響你的得救.
不過耶穌基督的義.

$^{481}$是要求你自動自覺.
安靜的.
甘心情願的去親近他.
回來教會去敬拜他.
這些不是儲積分的東西.
不是儲了多少次就有些什麼.
而是你甘心情願的.
一心一意的去做.
所以保羅在後來.
在哥林多後說.
原來基督的愛去催逼我們.
因我們如此判定.
一個人替我們眾人死.
眾人就都死了.
所以我就去為這個替我們而死的基督去活.
這是一個完全基督耶穌的催逼.
因為我們很知道.
耶穌很愛我們.
我們就節日各國的星期日.
九點鐘起床.
預備好自己.
回來崇拜裡面.
沒有人逼你的.
沒有人監你的.
沒有人不准你的.
不是問可不可以.
而是你自動自覺.
這樣去做一個基督徒.
不是因為律法的原因.
這個不是一個律法的交換.
而是純粹因為你知道.
耶穌基督在十字架上為你死.
並且為你的生命得付上代價.
你就願意做基督徒.
你就去聽一些完全沒有什麼特別要求.
意思是沒有什麼特別要監你.
沒有什麼懲罰.
但你都願意去做的一些病靈.
所以我們稱之為天父上帝的兒女.
不是老僕.

$^{521}$老僕或者高人姐姐都一樣.
放假就放假.
這些事情做多就要補中.
但是天父上帝的兒女.
絕對不是這些思維.
將來你兒女找你喝茶.
可以分到多少錢.
你就糟糕了.
原來是為了分你的身家.
完全不是天父的兒女的思維.
經文到最後.
保羅就用了一個很重要的說話來做總結.
保羅說.
所以你們因信基督.
耶穌都是神的兒子.
你們就受洗歸入基督.
都是披戴基督.
第七節經文.
這個所以非常可圈可點.
保羅說既然你們不是在律法之下.
不是在低B的律法主義之下.
保羅就怯面之下說.
所以你們就要取而代之地去做.
你們要披戴基督.
我們都知道.
做基督徒是很難的.
不是在律法之內.
也不是在律法以下.
要長期保持在一個高要求.
沒有什麼特別要監理的要求之上.
是一件極高難度的事情.
你做了有bonus.
做了不做你罰你.
那你就要做.
但這麼高要求的事情.
完全沒有人監理.
是很不容易的.
在一個律法以上.
登山寶訓是很難的.
耶穌基督的義是很難的.

$^{561}$但你們要不斷地去超越.
不斷地去做好.
不斷地去做更好的基督徒.
保羅說面對著基督的義.
恩順的義.
你們要披戴基督.
不是不做.
但你們要披戴基督地去做.
這個字我們覺得是一個很淑女的字眼.
披戴基督.
不過其實這個字是一個很搞笑的字.
什麼是基督?.
英文叫做put on Christ.
中文叫做coating on.
即是穿衣服的用字.
一個穿衣服的字眼.
大家想想.
所謂披戴基督的秘施.
其實是一個cosplay.
你想想明明你不是基督.
但你要一套甲痴衣.
去扮基督的cosplay.
一個custom.
其實是很特別的.
所以基督徒我們不參加萬聖節.
因為我們每天都參加萬聖節.
每天都在做cosplay.
都套上基督耶穌的衣服.
去做基督徒.
這個就是原文的意思.
你們要做一個cosplay.
明明你不是.
但你都要做到好像是.
說一個秘密給大家聽.
其實我中五的時候做過這些事.
有時候在商場裡扮過老虎.
穿著很熱的衣服.
暑氣功穿著這些衣服.
很臭的.
都不是我穿的.

$^{601}$很多人都穿過.
套上去.
最怕的不是中五.
中五真的會被小朋友襲擊.
真的會有人在後面踢我.
扯我的尾巴.
拉我的衣服.
我就想了一個方法去報復.
握手的時候.
輕輕捏小朋友的手.
就哭了.
小朋友就問我發生什麼事.
我就在那裡.
就這樣.
當時我還未被信主.
不要怪我.
後來老闆就換了一件衣服給我.
一件翠兒的衣服.
就沒有人來攻擊我了.
因為那件衣服是很恐怖的.
頭是人身形的翠兒.
是不可愛的.
整個人體翠兒.
這就是皮特基國的意思.
明明你不是耶穌基督.
但你仍然用盡力.
穿上基督的外衣.
去做好一個基督徒.
你沒有能力在律法上.
但你都用盡你的全力.
去嘗試去做.
我想這是我們今天的態度.
基督徒不是聖人.
基督徒也不是超然的人.
但我們卻願意.
每天都選擇.
穿上基督的外衣.
來嘗試去追求.
更加好的律法.
面對一些很難受的事.

$^{641}$面對一些很辛苦.
甚至付出的事.
我們都仍然付出.
明知道這個禮拜.
上崇拜不上.
沒有人會說你.
但我仍然付出一些努力.
明知道你不去做.
不去關心這個人.
其實都可以的.
沒有人會指責你.
但仍然甘心情願.
去做多一點.
我自己信主的時候.
三兄妹信主.
我也是錦繡信主.
所以我們三兄妹.
都在錦繡堂上教會.
我們三兄妹信主.
父母信主.
我跟妹妹說.
我們三個信主.
一定要做好一點.
要做好一點.
讓父母知道.
我們真的改變了.
我們中學生.
我妹妹是中三.
暴風少女的階段.
她穿著校服.
不換衣服.
吃零食.
看電視.
不換衣服.
不收拾東西.
我媽媽經常罵我們.
回家要幫忙洗衣服.
洗碗.
做些什麼.
那時候我跟妹妹說.

$^{681}$我們一定要做好一點.
要做好家務.
回來換好衣服.
做好家務.
很乖的.
這個其實不是我們.
這個是我們硬做出來的.
但這個是為了主的緣故.
頂硬上做出來的.
你說這些是不是很假呢.
當然不是假.
都是我們認真真心想去做的事.
這種投入一些努力.
明明自己沒有能力.
但仍然用盡力.
去裝扮基督的外衣.
讓人看到祂的存在.
這種的頂硬上.
正正就是披戴基督.
一個很重要的教導.
所以等下我們一起想想.
我們作為基督徒.
我們不在律法以下.
也不在律法以內.
耶穌叫我們不斷地超越.
不斷地盡力穿上基督的外衣.
明明裡面是很差的.
我們叫做什麼.
沒有材料裝四條.
但我們真的這樣盡力穿上外衣.
上帝有恩典的.
讓我們能夠用我們這個外衣.
用我們這個行為去見證祂.
開一隻祈禱.
十課上的救恩.
你叫我們每個人因為信的緣故.
接受就可以能夠被拯救.
但你卻無條件地在我們時候.
更加呼召我們去付上一切.
為你做到最好.

$^{721}$在我們的生命裡.
在我們的行為裡.
做一些超過我們所能做到的事情.
超水準的表現.
我們求你恩典.
讓我們願意穿上你的外衣.
成為你所喜悅的門徒.
求你幫助我們.
奉主名求 阿們.
\newpage



\section{}
\label{sec:v6FulXmsPK8}
\textbf{2023-02-23 周四晚查經班《被擄歸回歷史研讀(7)》 講員:曾裔貴牧師}
\newline
\newline
連結: \href{https://youtube.com/watch?v=v6FulXmsPK8}{\texttt{https://youtube.com/watch?v=v6FulXmsPK8}} ~~~~ 語音日期: 2023-02-25
\newline
\newline
\hyperref[sec:jaVgwb2h0Gg]{\small{< < < PREV SERMON < < <}}
~
\hyperref[sec:index_chronic]{\small{[返順時目]}}
~
\hyperref[sec:index_scriptual]{\small{[返順卷目]}}
~
\hyperref[sec:jYaZ2w15SEQ]{\small{> > > NEXT SERMON > > >}}
\newline
\newline
$^{1}$再一次多謝主.
我們能夠有機會研讀聖經.
我們很盼望有系統地去讀歸回的歷史書.
還有兩卷先知書的時候.
都讓我們可以探索.
在這個差不多有接近50年的.
一個被擄歸回的歷史.
我們可以掌握裡面的一些很重要的神學.
因為既然擺在聖經裡面.
不單止是給當時代的猶太人.
也都是在這個重建.
他們回去.
不只是硬件的重建.
而是聖文.
猶太的選民.
以色列人.
神選民的身份的重建.
這個似乎都是以色列記.
和尼希米記.
比較重要的信息.
求主都幫助我們.
雖然我們已經過了2400多年.
我們能夠透過這些文本的真理.
是沒有變的.
又讓我們能夠.
一方面回到那個時代.
他們那個關注.
但是當來到我們自己新約.
和在我們現在香港的現實.
我們能夠怎樣應用這些經文.
求主在我們中間帶領我們.
我們這樣禱告.
奉耶穌基督的名.
阿們.
我們上一次是去到第九章.
是以色列的祈禱.
一個個人的凝聚祈禱.
引發一個小規模的復興.
接著第十章.
這個極具爭議.

$^{41}$就是他們做出一些.
這些領袖.
在當時耶路撒冷的領袖.
聽以色列的教誨.
有什麼行動.
那個行動就是.
要將所有娶了外邦女子的.
全部要休妻.
因為有一個很大的關注.
所以這個行動本身.
就引發兩派的爭議.
到底他們這樣做是很衝動.
還是神所默許呢.
因為他們引用的是申命記的經文.
所以這裡就值得我們去思想一下.
到了尼希米記.
同樣出現了這個.
到了第十三章.
甚至連他們的梵安昔日.
尼希米也干涉.
也不准他們買賣.
就似乎都相當嚴峻.
他們似乎是想這兩位的領袖.
一個是祭司.
又是先知.
但是尼希米.
兩者都不是.
領袖.
在波斯已經做到很高級的一個官員.
亞特薩西王的一個酒政.
這樣有心.
很想回耶路撒冷.
做一個重建的工作.
建好城牆了.
有自己的框.
就是說我們這個城牆內.
就是我們神的選民.
就是聖民.
但是發現.
只是在這個框裡面的選民.

$^{81}$都有問題.
所以又進入尼希米記第十章開始.
又引發一個.
你可以說是一個認罪帶來的復興.
所以似乎兩卷書.
重點都不是那個硬件的重建.
的確一到第六章的以色列記.
是很清楚的.
是那個聖殿.
輾轉.
本來回去重建.
但是遇了挫折阻礙.
就停工.
一停就停了十五年.
後來就真的有先知的興起.
哈該和撒迦利亞先知的興起.
就勉勵他們.
他們就重新的去建造.
去到第六章其實已經做完的了.
但是第七章.
以色列就出場了.
以色列就是說他自己的背景.
尤其是要數他自己的家譜.
他是如假包換的.
祭司亞倫的後人.
就是說.
他是很想照做一些.
和他的那個輩分.
和是在摩西的律法書.
規定了的.
聖殿只能夠是神職人員.
就是祭司利未人.
才可以來到.
在裡面來服侍的.
所以這個就似乎是一個很大的關注.
他要找.
但是找不到.
不是那麼多人願意回去.
普羅大眾已經不多了.
我再查一些資料.

$^{121}$在波斯帝國.
以色列的年代.
差不多有五十萬.
粗略當然.
但是回去的四萬都不夠.
那麼你就可想而知.
十分之一的人.
願意連根拔起.
回去重建自己的地方.
就是不可以經常這樣旅行.
沒有飛機.
不會有船.
全部是靠騎駱駝這樣走.
是不可能你周車勞頓.
這樣周圍走的.
當然利未人.
他回去都要半個月時間.
那個令計.
因為他是有官職在身.
其他的普通平民百姓.
是不會這樣上去做完.
就回到自己的地方.
所以一回歸.
就是徹徹底底地回去.
我們看過第一章.
其實只有四萬多人.
來到第九章這個祈禱.
就引發了一個小小的波動.
就是第十章開始.
可能我們看多一點.
就是去到第九章的第五節.
當時是一個住棚節.
第五節獻晚祭的時候.
早晚的獻祭.
就是他們常規的.
獻的就是繁祭.
心中愁苦.
穿著撕裂的衣袍.
雙膝跪下.
向耶和華我的神舉手.

$^{161}$就是這樣站起來舉手祈禱.
說:我的神啊!.
我抱愧蒙羞.
不敢向我神仰面.
因為我們的罪孽滅頂.
我們的罪惡滔天.
從我們列祖.
直到今天.
我們的罪惡甚重.
因我們的罪孽.
我們和君王祭司.
都交在外邦列王的手中.
殺害,擄掠,搶奪.
捻上蒙羞.
正如今天的光景.
就是說將過去的被擄.
歸咎於得罪了神.
所以神就藉著外邦的巴比倫.
來懲罰,潔淨我們這個民族.
雖然我們被完全夷為平地.
但是神愛的不是那些.
對不起.
不是愛那個地方.
而是愛這一班人.
這班人需要有一個潔淨的行動.
就是70年.
在外邦徹底被煉淨.
所以刑罰滿了.
他們就會回去重建.
現在以斯拉.
來將過去的背景.
列祖所干犯的.
和面對神的審判的時候.
他份外看到.
猶太人本來是神的選民.
捻上蒙羞.
這是以斯拉在神面前的一個祈禱.
第八節.
現在耶和華我們的神.
暫且施恩於我們.

$^{201}$給我們留些逃脫的人.
使我們安穩如釘子.
釘在他的聖所.
我們的神好光照我們的眼目.
看到我們受核制之中.
稍微復興.
我們是奴僕.
然而在受核制之中.
我們的神仍沒有丟棄我們.
在認罪裡面都有盼望.
其實是在支取所羅門的那個祈禱.
歷代至下的第七章.
應該是十四節.
哪個人如果在聖殿向你禱告.
求你按照你的名去寬恕.
這是當時第一個聖殿.
所羅門獻殿的時候.
神和他應許的.
所以現在以斯拉真的在聖殿.
真的回到上來.
他就看到一些問題.
其實有些解經書說.
其實那些人數絕對不是少.
不是三兩個人娶了外國人.
即當地的迦南人.
或者其他族裔的那些外邦女子.
就不是猶太人.
而且是祭司都帶頭娶外邦女子.
不知道是不是實際上.
回到上去根本不夠女性.
不夠姐妹.
分配不夠.
所以沒辦法.
都要找外邦的女子.
這些是很現實的.
你上到去之後.
而且現在已經過了超過20年.
以斯拉是第二批上來的.
所羅巴巴那一代差不多都過去了.
所以現在他們又在這樣的情況下.

$^{241}$其實絕對是時勢.
實際的人情.
情有可原.
為什麼會看得這麼重這個問題呢.
這個是戴萊第十章.
就有這樣的具體的行動.
當然就是有領袖支持.
以斯拉這樣做.
他才有這樣的能力.
決心這樣去做.
現在仍然在這個祈禱裡面.
他在支取一個盼望.
雖然我們是奴隸.
在受壓制之中.
我們的神都沒有丟棄我們.
在波斯王眼前.
向我們施恩.
就是古列王的時候.
准他們回去重建.
叫我們復興.
能重建我們神的殿.
修其毀壞之處.
使我們猶大耶路撒冷有長門.
這個好像洩漏了一點點.
會不會其實尼希米.
都是在這個.
祈禱的那個延長呢.
因為我們如果將來看尼希米第九章.
尼希米又在那個水門的地方.
來祈禱認罪.
所以可能在不同的角度.
做著同樣的事.
當然因為那個困難.
有些認為是亞特斯西的一世.
有些可能以為是另外一個朝代.
但是因為聖經不講得太詳細.
所以這個重疊性.
都不能夠有一個一致的定案.
但是多數都認為.
他們是有重疊一起去復興的.

$^{281}$這個復興不只是以色列.
而是有尼希米的.
這個似乎是一個共識.
很重要的.
這裡告訴我們原來.
不是建好了聖殿就滿足了.
原來以色列回來.
本來一心是祭司.
他要帶領祭司的家族.
回去事奉.
回去公職.
回去聖殿.
帶領自己的國民.
去歸向神.
有回那個.
我們中國人叫香火鼎盛.
有回那種熱鬧.
好像剛剛初一.
又准可以公開聚集.
在黃大仙廟等等.
都很多人.
類似這樣.
但是回到上去.
所見到的情況.
就嚇一大跳.
就是發現.
很多的猶太人.
娶了外邦的女子.
第十九節.
我們的神啊.
既是如此.
我們還有什麼話可說呢.
因為我們已經離棄你的命令.
就是你直理僕人.
眾先知所吩咐的.
說你們要去得遺孽之地.
是污穢之地.
這些是摩西五經的經文.
相當重的.
這樣引述出來.

$^{321}$就是說和這些的.
外國人通婚.
就一定會.
被他影響.
就不會影響他.
來歸向真神.
是很負面的.
某程度是.
沒有信心.
為什麼一定不可以帶外邦女子.
和她的家族.
去歸向耶和華呢.
有路德就是這樣的案例.
但是.
似乎那個影響.
是對方.
會影響那個.
男丁.
來去玷污了聖潔的身份.
你們要去得遺孽之地.
是污穢之地.
因列國之民的污穢.
和可憎的事.
叫全地從這邊.
直到那邊滿了污穢.
是說那個.
對假神的敬拜.
假神敬拜.
帶來的生活的習俗.
和價值觀的取向.
所以敬拜假神.
就是將我們自己.
內心的慾望.
來無限的放大.
不過用一個神的名字.
來去敬拜.
所以拜財利.
譬如觀音開庫.
是要滿足人的.
內心的無盡的貪婪.

$^{361}$對錢財的.
那個.
發狂的執著.
因為錢財.
是最大的安穩.
保障.
所以不同的.
宗教都是將.
人性想得到的.
心靈裡面的.
內心裡面的.
那個自私.
那個貪慾.
放大出來.
所以這裡說的污穢.
不是那些髒髒的.
很吵架的.
烏煙瘴氣.
而是說的是.
心靈上.
有對造物主的專一.
單純的敬拜.
和沒有以祂是神.
這個.
這樣的.
價值觀來去.
敬拜.
而是自己創造一些.
自己想出來的.
神來去拜.
其實是滿足自己那個.
自我的膨脹.
所以基督教都要小心.
我們不是一個俱樂部.
不是一個我們在這裡.
打發我們的時間.
唱一些很好聽的聖詩.
使到我們的心靈能夠有一種滿足.
甚至乎有些人喜歡唱歌.
喜歡唱歌.

$^{401}$唱得很好聽.
其實他入師班只不過是想滿足.
他的歌喉.
要讓會眾聽到他的歌聲.
本來.
完全不是錯的.
基本的運作操作.
都是對的.
但不是唱給.
會眾聽.
來滿足你的虛榮心.
而是.
我們作為一個領導者.
是領導會眾.
一起去敬拜神.
唱詩讚頌的是神.
所以其實是一線的.
一線之隔.
端扶我們的心態.
我們是表演我們的慾望.
彈風琴很厲害.
彈鋼琴很厲害.
絕對沒有挖苦任何一個人.
說得很厲害.
說得很好聽.
很多人喜歡聽.
我每次都很想別人喜歡聽.
這樣就大件事了.
因為原來整個焦點.
是那個表演者.
大家要明白.
那種的污穢就在這裡.
所以如果我們事奉完.
而我們覺得.
有一種虛榮心和飄飄然.
那就真的要很小心.
當然別人讚你.
你又說不要讚我.
又不需要.
是一份的欣賞.

$^{441}$彼此的祝福.
我們的責任是應該這樣做.
但最終榮耀是要歸給神.
真的不要客套.
真的不要.
自表謙卑.
變成假謙卑.
是真的沒有人.
可以褫奪.
神的榮耀.
但要弄清楚.
每一次來敬拜.
是榮耀神.
這很重要.
所以現在這裡.
可以說以色列都是講到人性的焦點.
他看這個背後.
那個精神.
心靈那個.
污穢和性結的問題.
所以他在這裡.
其實他是.
找對那個焦點.
我們這一班人這麼辛苦.
上去.
遭遇那些困難.
已經排除了那些困難.
最終聖殿已經做好了.
最終.
現在我們本來開始.
有我們的敬拜.
70年沒有了.
重建又用了十幾年.
現在有回這個敬拜的生活.
是我們猶太人獨有的敬拜的模式.
所有的外族.
人都不會的.
現在我們跟回.
摩西的律法書的教導.
有聖殿設計好了.

$^{481}$擺好了器皿.
有祭司.
有證書.
真的.
那個家譜.
他們真的是祭司.
利未人.
不是找其他那些凡人.
已經做好了.
但是問題在這裡出現了.
是整個群體.
原來對於聖結.
是已經很模糊了.
是一個很離開.
摩西的標準.
相當遙遠.
剛才都說了.
是很多實際的現實.
和情有可原.
是令到你沒辦法.
覺得這件事.
是大錯特錯的.
是令到你沒辦法.
覺得這件事.
是大錯特錯的.
因為回到上去.
四萬的猶太人.
怎樣去解決.
已經是第二代.
甚至有些人是第三代.
他們的娶妻生子.
我是猶太的男丁.
找一個當地的外族人.
為什麼一定會帶來身份的污穢呢.
為什麼會這麼迂腐呢.
因為在聖殿裡面.
經過了很多的討論.
在聖殿裡面的領袖.
以色列還沒上來之前.
一定會引發一個討論.

$^{521}$其實對不對的.
神喜不喜悅.
但最終都是一個群眾的輿論.
和那個現實下.
沒有人再提這個問題.
沒有人再覺得這個問題是問題.
但以色列現在一回到上來.
自己看不過眼.
自己跪下.
自己去認罪.
自己去流淚.
自己去祈禱.
這裡說驚懼幽悶而坐.
拔自己的頭髮.
扯自己的鬍鬚.
這是九章第三節的經文.
再下去第十二節.
所以不可將你們的女兒.
嫁他們的兒子.
也不可為你們的兒子.
娶他們的女兒.
永不可求他們的平安.
和他們的利益.
這樣你們就可以.
強盛吃這地的美物.
並遺留這地.
給你們的子孫.
永遠為業.
這是摩西五經的經文.
申命記的經文.
第十三節.
現在以色列在祈禱.
神啊,我們因自己的惡行和大罪.
遭遇了這一切的事.
並且你刑罰我們.
輕於我們罪所當得的.
又給我們留下這些人.
我們豈可再違背你的命令.
與這行可憎之時的民結親呢.
整件事是在說.

$^{561}$草老婆.
嫁那些外邦男子.
是外族人來的.
不是敬拜耶和華的.
可能已經是迦南的原居民.
不是真的很遠的.
可能就是大家同住的附近的鄰居.
那怎麼辦呢.
第十四節.
我們豈可再違反這個命令.
與這些可憎的民結親呢.
如果這樣行的時候.
你,即是神.
豈不向我們發怒.
將我們滅絕.
以至沒有一個剩下逃脫的人.
耶和華以色列的神.
因你是公子.
而被人家迫害.
你不如去祈禱.
以色列的神.
因你是公義的.
我們這剩下的人才得逃脫.
正如今天的光景.
看啊,我們在你面前有罪惡.
因此無人在你面前.
站立得住.
其實是他自己一個人在祈禱.
不過呢.
第十章.
以色列禱告認罪哭泣.
俯伏在神殿前的時候.
有以色列中的男女孩童.
聚集到以色列那裡.
成了一個大會.
眾民無不痛哭.
即是說.
以色列的這個舉措.
這位領袖.
自己單方面.

$^{601}$一個人在聖殿.
已經完成了獻祭.
其實當時已經是黃昏.
獻晚祭的時候.
下午三點鐘的時間.
獻晚祭的時候.
他們就見到以色列在那裡.
不知道有什麼事情.
在那裡呼天搶地.
在那裡認罪.
哭泣,祈禱.
眾民就成了一個大會.
大家靠近.
聚集.
去到以色列的面前.
見到他.
聽到他.
這樣為自己的民族.
來痛哭.
當然他自己上來的時候.
很多人就見怪不怪.
因為已經是習以為常了.
已經是無可奈何.
都是要這樣做.
但是以色列上來的時候.
憤外感覺到那種.
呼吸的困難.
因為很多人.
已經是這樣.
去娶了外邦的女子.
又將自己的女兒.
猶太的女兒.
嫁給這些外邦的男子.
所以第二節.
蜀以蘭的子孫.
耶歇的兒子.
釋迦尼.
對以色列說.
這些都是祭司的家族.
我們在此地.

$^{641}$娶了外邦女子為妻.
干犯了我們的神.
然而以色列人.
還有子望.
現在當與我們的神立約.
休這一切的妻.
離絕他們所生的.
不單止要拋妻.
還要棄子.
離絕他們所生的.
照著我主.
和那因神命令.
賤經之人.
所議定的.
按照著律法而行.
這個時候出現了一個.
很大的一個.
猶太人的一個轉變.
其實這裡有一個很劇烈的轉變.
我們過去一兩堂有機會說過.
神仍然用先知.
來去.
講語言.
給一些意象.
來去.
讓以色列人.
得到那個啟發.
先知講到明.
是神吩咐我這樣講.
這樣講.
是語言來的.
不是先知個人的見解.
所以是耶和華如此說的.
但是在以色列這裡.
開始成為了一個.
重要的分水嶺.
他們不再有.
很強的先知的出現.
和那個先知的.
語言.

$^{681}$更加神跡都少了.
不多的.
相反取而代之.
就是有.
這些的文士.
的興起.
以色列就是其中的表表者.
文士的興起.
來去將.
舊約的.
律法書.
來去將它.
應用在現在.
我們這個民族.
這個群體.
那個身份和.
那個實際的處境.
來去詮釋.
摩西已經是.
二千多年前的那些律法.
怎樣能夠.
用在我們這個群體身上.
所以開始這裡.
不再是一種語言式的.
啟示.
而是透過.
文士對聖經.
的那個.
了解.
那個研讀.
和那個分析.
而帶出的那個應用.
就在這裡開始.
往後我們就見到.
無論是尼希米,無論是以色列.
他都是這樣來套用.
所以有些解經家.
都認為這是一個.
新的時代的開始.
標誌著以色列.

$^{721}$再沒有先知了,再沒有.
那些很轟天動地的先知.
講預言的先知了.
所以去到.
耶穌基督出現的之前.
就有這個施洗約翰.
就猶太人.
當時的猶太人.
就感覺是非常之奇特.
因為他居然間.
奉神的名.
來去開始講預言.
講一些啟示.
施洗約翰.
所以就帶來一個新的時代.
的開始,也都因為.
耶穌出現,又有很大量.
的神蹟,那就是說.
又有一個變成一個.
很特別的時代,又是.
和以色列的時代不同.
所以在.
這個耶穌基督.
降生之前的.
推翻上去大概400.
至,當然沒有500年那麼多.
這400年.
是再沒有先知預言.
的時代,而是有.
文時的.
對聖經,當然是舊約.
聖經的那個傳息的時代.
所以這很重要.
他們現在要.
按律法而行.
哇,摩西是2000多年前.
你們當時.
出埃及,進了迦南地.
你們當時.
是一個聖戰的一個群體.

$^{761}$你要將他們趕走.
當時他們在拜.
那些偶像,此一時彼一時.
歷史是不是.
可以這樣來到去照板煮碗.
這樣應用呢?.
這是一個問題.
現在已經是過了那麼多年.
有兒子,有孫子.
你現在.
以色列就說要他們全部拋妻棄子.
是不是.
很.
過了火位的一個.
舉動呢?.
很不容易的.
所以你讀這些經文都相當.
難去理解.
如果偶然間給一些.
不信的人看到這些經文.
就覺得基督教.
極之.
原教旨的塔利班.
跟足那個經文的字面.
去應用.
就是說.
難怪那些.
極端的回教徒要斬首.
偷東西要斬你的首.
真是斬的.
如果你做了一些不對的事.
道德上就要將你.
譬如同性戀者.
是要將你從二樓推下來.
跌死的.
這樣的一個刑罰.
真是按他們舊約古代的字面.
這樣應用理解.
那就真是沒有轉圜的了.
你可以想像得到.

$^{801}$可能在美國.
一半的美國人都不再.
可以在教會裡面.
來生活.
還有如果他真是.
這樣跟足去做的時間.
你可以想像得到那個問題的嚴重性.
這就告訴我們.
要按律法而行.
第四節.
你起來就是這個.
史嘉尼.
這位領袖.
就叫以色拉.
你當辦的事.
我們必定會幫助你.
第四節.
你起來就是你當辦的事.
我們必幫助你.
你當奮勉而行.
以色拉便起來.
使得濟斯長和利未人.
並以色列眾人起誓.
必照這話去行.
他們就起了誓.
以色拉從神殿前起來.
進入以利亞實的兒子.
約哈蘭的屋裡.
不吃飯也不喝水.
因為被老歸回之人.
所犯的罪心裡面悲傷.
他們通告猶大.
和耶路撒冷.
被老歸回的人.
叫他們在耶路撒冷聚集.
凡是不遵守令和長老所議定.
三日之內.
不來的.
就必抄他的家.
使他離開.

$^{841}$被老歸回之人的會.
離開我們這個群體.
被老歸回的人.
其實已經回來了.
超過二十年.
現在.
雖然未有成長.
但是.
這個群體.
已經成形了.
被老歸回.
上來的猶太人.
現在好像集中.
要處理我們自己.
被老歸回的猶太人.
這個身份的問題.
到底我們認不認.
我們自己是屬神的子民.
我們認不認.
我們有這個.
聖潔的本份.
我們需要有這個.
在神眼中的.
聖潔的要求.
不過那個要求是摩西的時代.
是不是現在.
仍然適合應用呢.
和在現在這麼少人.
回到被老歸回的.
這個光景.
是不是要這麼認真的.
落實執行呢.
在波斯帝國.
其實他們已經取了.
都大有人在.
說的是五十萬人.
在波斯帝國的國境.
現在這麼少的人.
又如何呢.
譬如這個.

$^{881}$Esther 即是以斯帖.
她是嫁給阿克徐老王.
是皇后.
又如何呢.
即是說.
其實是很普遍的.
在波斯帝國.
已經不應該是問題.
但是現在.
以斯拉就很介意.
我們這個.
小的群體.
由四萬人衍生.
多少呢 我就沒有計算了.
也沒有給我們有數據.
有多少人其實在耶路撒冷呢.
但是.
你堅持要這樣做的時候.
是一個完全沒有.
彎轉的一個.
極之嚴重的.
一個舉措 差不多帶來.
暴動了 三天之內.
誰不來就抄家.
我們看一個後續發展.
要離開.
即是說要將你們驅逐.
不可以再成為.
我們的召會裡面的.
一窟一分子 即是說.
你不是我們聖潔的身份.
第九節.
於是猶大.
變雅敏眾人 這兩個支派.
是最明顯跟隨.
祭司上來.
重建的 歸回的.
因為他們是.
大衛的子孫.
大衛的皇朝.

$^{921}$不是北面的十個支派.
即是說.
原來都是有背後的原因.
為甚麼只是猶大.
變雅敏支派會歸回呢.
和有祭司會歸回呢.
因為他們本身一直都是.
這個.
一脈相承的敬虔的系統.
所以真的有關係.
即是你的太公如果.
很敬虔 你都會有機會.
是不知不覺.
你的靈命等等.
都有幫助 你的父母.
如果經常祈禱.
而不是經常打牌.
其實差很遠的 你的.
那個敬虔 所以猶大.
變雅敏支派的人.
大部分跟上來重建的.
即是歸回的.
於是猶大.
變雅敏眾人三日之內都聚集.
在耶路撒冷.
那日正是九月二十日.
過了所有節期.
三日之內都聚集.
在 啊 對不起.
眾人都坐在神殿前的寬闊處.
因這事.
又因下大雨.
就都賤敬 祭司.
站起來 強調祭司.
這個身份 因為現在處理這件事.
是關乎到.
律法 關乎到.
祭司.
和他的選民.
神的選民.

$^{961}$這個是神學上.
宗教上.
不是民生上.
不是民間的生活.
的問題 是.
靈性方面的問題.
所以 祭司 以色列站起來.
對他們說 你們有罪了.
因你們取了.
外邦的女子為妻.
增添以色列人的罪惡.
現在.
當向耶和華你們列祖的神認罪.
遵行他的旨意.
離絕.
這些國的民和外邦的女子.
其實.
你想想.
其實是很.
很惡 很惡 很惡.
你想想 想像得到.
肯來.
聽你講.
三天之內 大家都願意.
來聚集.
當時是有這種.
的服從性.
因為大家都見到這個領袖.
那個的.
可能是威望.
可能是他那個的敬虔的.
權柄.
那種的影響力.
但大家都願意來的時候.
你是很.
清清楚楚.
一點都不含糊.
你們是有罪 現在要做什麼呢.
就是要離絕.
你們的.

$^{1001}$其他的國民和這些.
外邦女子 即是說可能是你的太太.
可能是你的媳婦.
那你就可以想像得到.
那個問題的嚴肅性.
會眾都大聲回答說.
我們必照著.
你的說行.
只是百姓眾多.
是大雨的時令.
即是他們有秋雨.
秋天的雨.
我們不能站在外頭.
這也不是一兩天.
辦完的事 原來.
很多人.
基於一個現實.
基於一個當時的處境.
基於自己只有四萬多人.
上來重建.
一路生都不夠那個人.
而.
實際上真的要找了.
外邦的女子.
外邦的女子不一定是壞人.
不一定是拜偶像.
但是現在.
以色列一回到上來就發現這個問題.
果然真是.
原來很多人都是這樣.
所以不是一兩天.
就辦完這件事.
因我們在這事上犯了大罪.
不如為全會眾.
派首領辦理.
凡我們成一中.
娶外邦女子為妻的.
就當按所定的日期.
即是一個.
好像類似叫做.

$^{1041}$寬限期.
過了這個限期.
如果都不做的就有首尾根.
可能是這個意思.
就按照著.
那個日期.
同著本城的長老和事司.
而來 不單只是長老.
還有那些審判官事司.
即是說有他們.
的宗教警察也說不定.
來到即是說.
可能有機會用武力.
根據摩西的律法.
來施以一些可能是刑罰.
可能是他們自己制定了.
的律法也不一定.
一些法律 宗教的法律.
即是去到耶穌時代.
那些的守殿官.
祭司長.
他們都有他們的法律.
誰做了什麼事.
就要打多少下.
如果是不知道什麼罪.
39下.
就可以免一下 如果猶太人.
這裡告訴我們.
直到辦完.
這件事之後 神的烈怒.
就轉來我們了.
是這一班的領袖.
特別是由.
這個史嘉尼和當時.
猶大變亞敏茲派.
的眾人來叫.
以色列認真地.
頒佈這件事 第十五節.
惟有阿薩克的兒子.
約拿丹 特瓦的兒子.

$^{1081}$阿哈謝.
阻擋這事.
並有米舒蘭.
和利未人沙比泰.
幫助他們.
真是有阻力.
而且利未人都帶頭.
阻礙這件事.
那個難度.
就是在這裡了.
首先就是有沒有法律的效力.
二千多年前.
有西方的律法.
那個時候出埃及.
進迦南.
要驅逐.
所有的迦南七族.
是神吩咐.
但現在不單止.
不是我們要驅趕.
是我們現在回到這裡.
人家才是主場.
我們現在是寄居.
當然是以迦南人的想法.
實際上猶太人.
是回到自己的地方.
就是能不能夠一廂情願.
將所有的.
這些外邦的人.
趕走.
所謂趕走就是不再來往.
不可以再有任何的.
很緊密的一個.
接觸.
當然是指婚姻上的那個.
要將他們趕走.
你想想這件事.
那個情況.
再加上有阻力.
甚至乎是由.

$^{1121}$迦南人那裡提出.
一個阻擋.
第十六節.
被擄歸回的人如此而行.
祭司以色拉和奢族長.
按照著宗族.
都指明建派.
在十月初一.
上文講的九月二十.
已經過了十天.
一同在座.
查辦這事.
到正月初一又過了兩個月.
十一月十二月.
現在去到正月.
正月初一才查清.
取外邦女子的人數.
開始原來是調查.
用了差不多三個月時間.
由九月二十.
那個的.
祈禱認罪.
痛哭.
引發的.
這個士嘉尼.
這個領袖叫.
以色拉免力而行.
因為這個是.
你作為祭司的責任.
你要教我們律法書.
什麼是神的心意.
你要教導我們.
以色拉就來宣讀.
律法書.
然後也有一些領袖.
不如這樣.
既然大家都願意.
我們就三天之內聚集.
承諾.
起誓.

$^{1161}$我們會跟你的指示去做.
所以就有這個.
引發了他們就開始.
成立了一個調查委員會.
去調查.
每一個家庭.
你是什麼支派.
你是什麼人.
現在你的家庭背景.
你的太太.
你的丈夫.
是什麼的支派.
原來不是猶太人.
原來是外邦人.
就出現了這個問題.
記錄在案.
怎樣籌辦呢.
十八節.
在祭司中查出取外邦女子為妻的.
特別要強調的就是.
祭司.
特別強調的就是.
聖殿裡面事奉的.
如果他們都不能夠.
來到去聖結.
就令到會眾更加.
無所適從.
就是耶穌亞的子孫.
約撒達的兒子.
和他弟兄.
馬西亞.
以利爾謝.
雅納.
基大利.
相當多人.
他們便應許.
必休他們的妻.
他們因有罪.
就獻群中的一隻.
公綿羊贖罪.

$^{1201}$他們自己立即.
來到去承諾.
簽了協議書.
說我們承諾了.
必定會休我們的妻子.
休書.
以前摩西說可以休書.
休妻子.
就是怕.
那些男丁.
太過硬心.
已經不再喜歡這個太太.
但是如果規定說.
不准休妻.
那就寧願虐待到她死.
因為.
你不准我休妻,那我就虐待她.
不准她吃飯.
要她去柴房.
要她去淡水.
要她很辛苦.
甚至在眾人面前.
奚落她.
一定會有這樣的侮辱.
和欺凌的事件.
所以摩西就說.
因為.
這是耶穌說的.
因為你們的心硬.
所以准許你們休妻.
其實神的心意是不想的.
所以.
現在這個休妻又不同.
現在是另外一個處境.
不是太太有問題.
不是太太.
現在只是拜偶像.
自己星期六.
不去安息日.
不去跟丈夫,其實他們是不能跟的.

$^{1241}$因為是在家裡.
來自己去拜其他的神.
可能不是這樣的.
所以現在這裡.
這麼介意.
聖潔的民族.
所以這個就是.
以色列的核心部分.
原來這兩卷書.
想探討的.
不是硬件的重建.
是.
這一群被擄了這麼多年的人.
現在回到聖地.
似乎.
領袖們想.
達至的就是.
要回到.
神起初.
這個民族的要求.
就是要有.
聖潔的體現.
體現就是.
這個就是入心入肺.
個個要將自己.
外邦的女子趕走.
那你家家戶戶.
都真的很慘.
現在不是集體的.
結婚,是集體離婚.
還要是猶太人.
可想而知.
到底這是復興還是.
暴動呢.
不過很有趣的就是.
整卷書停留在家譜.
就沒有再說下面的後續.
就記載了.
這麼多人.
娶了外邦女子.

$^{1281}$而且是.
祭司的家譜.
真的挺有趣.
好像讀者自己評斷.
整件事.
到底是一個很魯莽的.
舉動.
還是摩西的期望.
在現在.
被擄回來的處境.
要這樣做.
還是有這個理想.
我們這個民族.
是屬神的.
我們現在在聖殿裡獻祭.
我們這個民族.
是歸神為聖的.
但是是不是應該每個人.
個別自己去反思.
自己按照著.
自己的力量.
去做.
心甘情願的.
去做.
這個.
行動.
而不是一個集體的被強迫的.
要簽協議書的.
在群眾.
壓力底下的.
這就是兩者的對比.
我們再讀那個家譜.
第20節.
陰默的子孫當中.
有.
哈林的子孫中.
有.
(念句).
烏西亞巴斯護爾的子孫當中.
有.

$^{1321}$(念句).
利未人當中.
剛才是祭司.
現在是利未人當中.
有.
(念句).
歌唱的人當中.
有.
大衛連耕.
他沒有機會建聖殿.
但是他分派利未人.
有些是唱歌恩賜的.
例如有.
這個.
一陣子想不起那些名字.
耶道頓.
那些人.
安排他做師班.
歌唱的有.
(念句).
守門的人都是在聖殿裡面.
有沙龍.
提倫.
烏尼.
以色列人當中巴祿的子孫.
有拉米.
耶西亞.
馬基雅.
米雅文.
以利亞.
馬基雅.
比拿雅.
以南的子孫當中.
有馬他尼.
薩加利亞.
耶歇.
馬哈尼.
耶利米.
阿西薩.
比拜的子孫當中.

$^{1361}$有約哈蘭.
哈蘭尼亞.
薩拜亞納.
巴尼的子孫當中.
有米舒蘭.
馬祿.
阿大雅.
阿述.
耶利米.
巴哈摩亞的子孫.
有比撒列.
登內.
嗎哪西.
哈林的子孫.
有以利爾謝.
伊士雅.
馬基雅.
是馬雅.
西緬.
馬祿.
是馬尼亞.
哈順的子孫當中.
有馬特乃.
馬達他.
比底雅.
基祿.
米利末.
以利亞實.
馬他尼.
馬特乃.
巴尼.
是美.
是米雅.
阿大雅.
沙瀨.
是米雅.
是馬尼亞.
沙隆.
阿拉伯的子孫.
有耶利.

$^{1401}$馬他提雅.
薩拔.
西比拿.
阿大.
約爾.
比拿雅.
這些人.
全部.
還有這些.
大的家族.
是被老.
歸回.
以派佔大多數.
現在.
因為以斯拉的舉措.
簽協議書.
承諾.
這些人都娶了.
外邦女子為妻.
其中都已經生了.
兒女.
就停留在這裡.
所以以斯拉.
尼希米很多時候.
一併去讀.
雖然尼希米.
突然之間又記載.
尼希米的生平.
他的.
上去重建.
城牆的背景.
一到第二章.
第三章就上去.
然後籌劃.
第六章就建造好.
第七章又開始引發一個.
所謂水門事件.
Watergate.
那裡.
又是這樣見到.

$^{1441}$外邦女子的問題.
所以.
之前很早說過.
希伯來文的.
以斯拉尼希米.
是一卷書.
是歸類.
因為.
原文的.
希伯來文的舊約聖經.
他們沒有新約.
舊約聖經有.
三分法.
摩西格雷法書.
斯哥哲偉書.
或者應該叫.
前先知書.
後先知書.
和其他的書卷.
其中有.
被擄歸回的.
以斯拉尼希米的一卷書.
所以.
雖然這樣斷開了.
有些人.
將它.
你可以說.
將它從.
很大的次序.
變成第七到第十章.
是.
以斯拉的一個回憶錄.
尼希米.
這一大卷書.
本身就是用回憶錄的手法.
去寫成的.
所以.
如果這樣理解.
我們將它再由第七章的尼希米記.
那裡再這樣看.

$^{1481}$就很銜接到.
整件事.
這個復興.
關於外邦女子的問題.
現在.
不是特別針對外邦女子.
是針對.
我們這個民族.
因為現在.
心慶得人.
以斯拉她一開始.
就說她自己的身份.
背景和她的學習.
她的學養.
她的訓練.
她甚至準備好.
是要有聖殿的事奉.
她是.
絕對有資格.
有學歷.
來帶領自己的國民.
來敬拜神.
但是.
未說敬拜之前.
就發現了我們這個民族.
距離神的.
那個聖潔標準.
無論精神上.
或者那個叫做靈性上.
實際的操作上.
就是生活的層面.
都和聖潔.
那個距離相差很遠.
所以.
以斯拉由她自己.
一個自發的.
認罪祈禱.
而引發了這個.
可以說是一個小的復興.
來到她.

$^{1521}$大家要尋找一個問題的答案.
就是神的律法.
怎樣說.
我們應該怎樣.
行為我們的生活行為.
要符合.
神的律法的教導.
所以就.
不需要有先知的預言.
或者說神都沒有再興起.
反而.
大家就透過.
以斯拉的解經.
那個對律法的詮釋.
大家就開始有一個感動.
這件事.
不是個人.
這麼簡單.
是整個群體.
面對一個.
極之大的尷尬的事件.
差不多引發.
好像暴亂一樣.
因為會引起.
波斯帝國.
因為這個都是屬下的一個河東的.
一個管轄地.
有他們的河東的省長.
有他們的巡撫.
這個.
猶太這個民族.
現在說要來做一個.
休妻的行動.
是因為以斯拉.
慫恿我們這樣做.
等於好像在美國的邪教.
現在集體.
我們要去哪個地方.
避世.
因為就快有大災難.

$^{1561}$FBI很緊張.
因為集體離開等於集體自殺.
又真的發現.
有些邪教.
真的一起服毒自殺死.
所以有些救不及.
有十多人死了.
真的是新聞.
所以一有邪教的教主.
這樣來慫恿.
世界末日要來了.
我們這個支派.
我們這個教派.
就會被提了.
就會得到神的拯救.
FBI一直在追蹤這些邪教.
就很緊張.
就來到.
甚至破門進去.
後來有些已經集體服毒死了.
現在.
在波斯帝國的眼中.
在那些.
持份者.
外邦的家庭.
那些親家.
那些長老.
他們怎樣看.
猶太人.
以斯拉這個激進的.
舉動呢.
所以這個問題就停留在這裡.
奇妙就是.
為什麼不給我們答案呢.
承諾了.
簽了協議書.
但是沒有行動.
去到.
尼欽米記十三章.
也是一樣.

$^{1601}$沒有行動.
最後就是尼欽米.
神求你聽我祈禱.
施恩給我.
就完結了.
好像留了一條很大的尾巴.
不少人.
我們剛才讀一讀那些人名.
所牽連的家族.
其實很大.
他們承諾了.
會處理自己家庭的.
這些娶了.
外邦女子.
甚至乎生了子女的問題.
怎能夠.
在這麼嚴肅的道德倫理.
來到去.
說離開就離開呢.
這樣拋妻棄子呢.
這樣來到.
去休妻呢.
答案.
不知道.
當然.
我個人認為.
既然.
以色列.
說到自己的身份.
和抱負.
和現在.
回到重建.
神的帶領.
和他的決心.
他的敬虔.
他真是一個.
看到.
面前的.
一個社會現象.
說的是猶太人這個社會.

$^{1641}$的現象.
大家是無可奈何.
大家都是看到這樣的情況.
但因為太多了.
已經很難再個別處理了.
所以根本就變成.
不能夠再處理了.
就好像70年代.
葛柏事件.
很多警務人員曾經涉及貪污.
就定了.
有一個寬限期.
之後就真的會追究了.
很多人.
以前做過的就算了.
都是要這樣.
抓了一些.
特別明顯的.
有些殺雞儆猴.
但是.
現在就是在這個情況下.
所以以色列這個舉措.
因為是要.
回到摩西的律法.
那個對於聖潔.
的標準.
那個的尋求.
還有.
由這個階段開始.
尤其是.
以色列.
尼希米兩位領袖.
來自不同的界別.
他是.
靈性上因為是祭司.
尼希米肯定就是.
一個很能幹做了高官的.
一個領袖.
很有心要.
重建自己的.

$^{1681}$國家耶路撒冷.
的那些城牆.
希望不要被欺凌.
最終又是發現.
又是有外邦.
結盟結親的問題.
所以似乎.
這兩個領袖.
在這個時期.
引發出一個運動.
但是這個運動.
不是停留在這裡.
而是一直蔓延.
去到四百年之後.
耶穌出來之前.
就有一個真真正正的.
法利賽運動.
就是大家經過了.
這麼多的年日.
一路一路去尋索.
終於大家都有一種.
好像痛定思痛.
就是我們這個民族.
真的不應該和一般的.
迦南人外邦人.
這樣的生活質素.
和那個靈性.
的那個眼光.
應該是有我們自己獨特的.
屬靈的眼光.
應該有我們神選民.
的身份認同.
的.
所以就導致.
有這個應該是.
主前的一百六十幾年.
一百六十二年.
那個法利賽運動.
就真真正正.
要back to the bible.

$^{1721}$尤其是舊約摩西.
的律法書的標準.
所以.
兩位的領袖.
可能是.
這個舉動是極之激進.
但是.
是要將整個民族.
扭轉回去.
那個敬虔的.
一個準則.
類似如果能夠很近似的.
就是1517年的馬丁路德.
因為整個當時的基督教.
已經是極之的.
低落.
那個靈性已經是很.
已經是很失敗.
所以馬丁路德.
就將大家的.
焦點帶回去聖經.
所以他強調的是唯獨聖經.
還要是唯獨信心.
因為靠信心不是靠你的.
行為不是靠你的功德.
去換取救恩.
是唯獨耶穌基督.
唯獨基督.
唯有主耶穌是我們.
袋屬的唯一的途徑.
接著就是.
唯獨那個恩典.
是恩典來的不是換取的.
是白白的.
這個是馬丁路德.
那個時代的其實應該是.
五個唯獨還有就是人生.
是要唯獨榮耀神.
所以是由馬丁路德.
發現了因為他自己.

$^{1761}$是一個天主教的.
一個修士.
他自己一直都是在教廷裡.
去生活.
但他自己在年輕的時候.
他完全沒有得救的.
信心.
很怕神.
很怕自己不能夠.
得到神的永恆的永生.
一直都靠自己.
但是他以為.
去再神聖一點的.
應該是.
羅馬的.
主教.
大教堂.
但他去朝聖的時候.
發現原來.
都是改變不了他內心.
那個的.
無論是他內心的痛苦.
和他內心的掙扎.
和他內心面對自己的.
罪.
他自己不能夠解決.
就是他自己的自傳.
這樣說.
他一直想嘗試.
在那個.
威登堡教堂.
自己跪上去.
靠這些功德.
來去一直練正.
自己的心思意念.
誰知道.
不跪又折鶴.
一路跪的時候一路背聖經.
接著突然間.
羅馬書一張十七節.

$^{1801}$就閃入他的腦裡面.
因為他是熟讀聖經的修士.
羅馬書.
原來自己一路讀.
這些聖經.
居然一路都不知道原來.
保羅已經說了.
為義人是恩信.
得生.
恩信就可以得生.
他突然間閃入去.
是聖靈光照他.
所以他自己回家.
再看羅馬書.
終於發現自己得到的亮光.
和整個當時的天主教.
大公教會.
所主張的.
所高舉的.
例如賣贖罪券.
例如大家爭相的.
透過不同.
離世的人.
留下的聖物.
去買那些聖物.
來保佑你.
很多這些變成.
信仰很簡單的恩信.
就變成一個.
嫁床疊屋的.
依靠行為.
很多層層的代理.
這樣就帶來.
一般的平民.
是耕田的農民.
什麼都不懂.
你說什麼我們就做.
就帶來.
根本就不是真正的.
一般平民百姓的信心.

$^{1841}$的釋放.
而是更多更多的代理.
代理人就變成.
甚至有些貴族.
買了贖罪券.
就可以肆無忌憚的犯罪.
因為有贖罪券.
就導致本來那件事.
是很好的.
帶來的就是.
有錢的人就可以有這些贖罪券.
沒錢的人根本不明白.
什麼叫信耶穌.
根本得不到良心的.
清潔和.
信心帶來的自由.
釋放.
這個就是馬丁路德.
他的後繼者就有嘉爾文.
還有比如慈雲尼.
這些神學家.
就更加對聖經的.
詮釋.
回到對聖經的解釋.
所以那個時候.
就湧現了.
大量的聖經翻譯.
因為天主教.
只能夠說.
你必須要讀的就是.
拉丁文,不是拉丁文.
不可以讀.
因為自己解釋聖經.
所以每一次的崇拜.
他們叫彌撒.
或者叫微撒.
就是由神父讀.
拉丁文,不是每個人都懂.
拉丁文的,那些根本不懂.
就當去.

$^{1881}$看看,去聖餐.
都不知道你聽完說什麼.
就走了,根本聽不明白.
也都不會這麼有系統.
這樣來解經.
也都在早一些的時間.
如果有人試圖.
要翻譯聖經,有的.
天道,天條.
就是將聖經翻譯.
就會.
有這些宗教警察.
就會去捉拿.
就會找到這些聖經.
就要來沒收,因為怕你.
自己解經.
就會變成異端.
如是者,就在那個時代.
就因為馬丁路德的改革.
其實有政治因素.
因為很多當時的.
城邦的皇.
不滿.
教皇的權力.
無限大.
教皇說,你不聽我的話.
你將來上不到天堂.
那些皇帝都怕.
但那些皇帝又越來越.
有自己的分封的地.
有自己的權力.
很自然,開始不再.
這麼聽話,但有什麼辦法呢.
現在有個馬丁路德在.
原來在神學上.
和天主教.
可以來到去.
有一個分庭抗禮.
所以就出現了.
支持這個馬丁路德.

$^{1921}$對抗.
大公教會,馬丁路德.
很自然.
無可厚非地都會.
借助政治的勢力.
去保護自己的安危.
所以馬丁路德一生都.
差不多近乎逃亡.
因為都追殺他.
後來晚年就不需要.
安定樂利.
甚至最大的舉措就是.
一個神父.
為了要履行聖經.
他自己在晚年的時候就娶老婆.
很難接受,很難想像.
我們不是這樣的人.
但是馬丁路德.
為什麼呢.
因為你引發的宗教改革.
很多天主教徒.
很多修女.
離開了教會.
年老,沒有人要,怎麼辦呢.
馬丁路德要自己.
埋尾.
他自己主動地.
娶其中一個修女.
就是說其實可以結婚.
是要照顧這些.
離開了大公教會.
真是無依無靠的人.
你要幫助他.
這就引發了一個神學的解釋.
讓其他人都知道.
其實你帶來的復興.
是要執手尾的.
你不可以說.
我又收視,沒錯.
其實他可以不結婚的.

$^{1961}$但是他真的自己去履行.
這個責任.
就帶來當時的宗教改革.
真是勢成水火.
所以今天我們的基督教.
和天主教.
雖然我們是分離的弟兄.
天主教稱呼我們.
我們其實是.
有些翻譯叫抗羅宗.
抗以羅馬的.
宗教的一個宗教.
就是我們這些基督徒.
基督教.
更多的就是我們是.
歸回教.
歸正教.
更正教.
更正原本聖經的.
真理的.
唯獨信心這個課題.
所以.
如果這樣理解的時候.
以斯拉其實是.
好像開了第一槍.
引發這個.
復興的一個.
焦點.
那個理想是極之高的.
就是我們這個民族.
回到聖地.
敬拜的是.
這位耶和華神.
但是我們的.
實際的生活應該是什麼標準呢.
他們就要找標準.
找什麼標準呢.
必然就是摩西的律法書.
那個標準.
所以整件事.

$^{2001}$是非常之神聖.
但是做起來.
就似乎.
很黑社.
非常的嚴峻.
好像逼人們.
要簽約.
簽協議書.
必定要在.
第二年的初一.
要離絕.
自己的已經娶了的.
外邦女子.
那真是家家戶戶都沒有.
女士.
我覺得.
就是聖經停在這裡.
就是說.
其實這一次的復興.
不是很多人.
跟進去做.
當然沒有答案.
不是我有其他的經文支持.
而是正正就是.
那個理想.
不可以一下子跳到那個階段.
一步就登天.
去到最理想的那個境況.
似乎是.
仍然要回去.
如果既然我們是尋求.
聖經律法那個標準.
那個解釋.
和大家按照著.
這個理想.
去達到的時候.
似乎又是一個新的時代.
開始神容許.
每個人.
按照著自己的良心.

$^{2041}$去達到這個.
理想.
這個民族.
屬神的.
不要在頭腦上認知是屬神.
以色列期望的.
就是你要做給人看.
這樣去做.
似乎.
還是舊約的那個背景.
就是.
他們當然要.
最高的準繩.
一定是摩西五經.
所以去到新約.
耶穌出現的時候.
已經不是這樣的狀態.
沒有任何一個地方.
沒有任何一個人.
你做了的事就等同於什麼.
而是新約有.
新約的倫理準則.
又是很不同的.
新約就更加說.
一定是你個人.
信心的程度.
你要合乎中度.
這樣去看你自己.
那個擺相.
所以呢.
這裡你可以說是一個.
挺突出的.
一個引發.
復興的.
那一點的火花.
以色列真的在.
她的責任和.
她的.
理想上.
就是想做這件事.

$^{2081}$然後真的被她發現到.
原來不需要去周圍.
明查暗訪.
因為原來比比皆是.
個個都是.
你根本就是一個.
你沒可能再改變的.
但是她都要停止.
她都要畫一條線.
那條線就是神的選民.
是不是應該有神的.
那個要求的標準呢.
她畫了這條線之後.
就出現了這個.
大家按照著.
你自己的良心.
你自己去做了.
簽了協議書了.
做與不做呢.
就你自己去問你自己的良心.
似乎那個答案.
是停留在這裡.
就沒有一個.
清清楚楚有多少.
決心離開.
他們的外邦女子的太太呢.
所以呢.
這個是有趣的地方.
我要呈現給同學知道.
我們要來去學習.
新約一定會對我們.
有這個要求.
特別是新約的書信.
他們處於的那個.
情況就是希臘化.
羅馬帝國底下.
的一個帝國.
的一個當時的.
生活.
當時他們那種的.

$^{2121}$思潮.
思想.
他們的民生.
所以當時的情況.
每一卷書都是針對他們.
那個現實的處境.
其中又是這個最大的課題.
你成為了基督徒.
你在神的眼中.
應該怎樣.
你就要想你自己了.
那就有很多實際的教導.
譬如從前偷竊的.
這個是二分所說第五章.
可能第四章.
就不可以再偷了.
好像若有所指.
有些人以前真的偷竊.
坐牢.
現在不可以再偷了.
怎樣呢.
都要努力.
親手的作正經事.
即是你要工作了.
你要賺錢了.
就可以有餘分給有需要的人.
很實際的教導.
信耶穌都要教別人不要偷東西.
就是有人偷東西.
當時.
譬如提多書.
就說那些.
格里底的人.
格里底即是克里的那個島.
希臘其中一個很大的島.
那些格里底人.
就說.
很饞嘴的.
很胃蝕的.
很懶惰的.

$^{2161}$保羅說.
這些傳聞是真的.
然後提多說.
你要這樣警戒他們.
信了主耶穌就不要這樣.
這樣教.
很實際.
即是說原來那些.
可能剛剛不是很開發.
那些島民.
還是很.
橫蠻無理.
還是很胃蝕的.
海島的人.
你去到傳福給他.
他信了主耶穌.
但他還有很多的生活未改變.
你要慢慢教他.
可能要一代兩代.
才慢慢改變這些的習慣.
譬如我爸爸的年代.
我們出街.
隨便就吐口水在街.
吐痰.
是經常性個個都是這樣.
見到長輩全部都是吐痰.
你見怪不怪.
你自己又習慣了吐.
信了耶穌之後當然整個社會.
那個生活改善.
大家又有了教育.
有了群眾壓力.
下一代就說.
爺爺你不要吐痰好噁心.
慢慢他都收斂.
在有生之年一樣吐.
信了耶穌都照吐.
因為習慣了.
你就知道原來有些東西是要一代一代.
慢慢慢慢這樣改變.

$^{2201}$所以呢.
這裡這麼困難的.
一個簽約.
協議書.
實際去進行.
都是要個別.
有些很強的心智.
有些是良心.
於心不忍.
太太原來對他很好.
你怎能夠說我要跟大隊.
我要趕走你.
現在不是文革.
文革就很恐怖.
所以似乎留一個空白.
讓我們好像.
不知道當時的情況.
又好像讓我們有個空間.
很多東西是.
我們的內心.
其實我們自己的選擇.
我們是否想愛神多些.
我們是否想行神的旨意多些.
是每個人個人的選擇.
所以往往就是.
每一次聽到.
有些人越聽就越.
深入.
越來越他有感動.
有些人越聽.
他可能就會覺得.
嘩又這麼多要求.
就會有很多的.
反彈.
而他慢慢忘記最重要的.
那個大前提.
意思是想帶他們看.
神到底期望.
我們被擄回來的這班人.
我們是.

$^{2241}$什麼的一個.
屬靈的.
那個位份.
其實這個point導致.
有這些舉動.
所以如果他們明白了.
慢慢朝向.
可能這一代不行.
他們下一代就真的.
這樣做了.
又發現原來.
這樣走又好像work.
又不是帶來.
真的家散人亡.
又真的可以不會令到.
那些外邦女子.
帶著很大的.
眼淚被趕走.
可能很多.
這些人情世故.
而以色列其實都明白.
他自己都完全知道.
帶來的代價是很大的.
所以呢.
這裡聖經好的地方.
就是留下一個空白給我們.
相當大的空白.
到底有多少人做.
不告訴我們.
讓我們自己去一個揣摩.
所以去到新約.
就更加告訴我們.
我們一信主耶穌.
就是聖潔的種類.
已經是.
我們是要.
這是彼得前書的二章.
第七節.
唯有你們是被揀選的.
種類是聖潔的.

$^{2281}$祭司是屬神的.
子民要叫我們.
能夠.
歌頌這一位救我們.
出黑暗進入.
奇妙光明者的美德.
我們是歌頌神的.
那個救恩的美德.
要我們時時.
引發我們的歌頌.
如果我們裡面都不是.
很honor(榮耀)這個身份.
那何來歌頌呢.
譬如你都不喜歡.
你自己這個國民身份的.
你還怎麼會說歌頌的說話呢.
你都不喜歡自己這個民族的.
那很難的.
所以其實傳道人.
是要提升弟兄姊妹.
你不只是回來教會.
而是到底我們這一群人.
我們稱為錦宣家的這一群人.
我們是否明白.
我們是屬神的子民呢.
我們在我們的契約.
在我們的生活.
我們有沒有這個屬神.
子民的追求呢.
我們是不是.
被揀選的族類呢.
被揀選.
就是說在眾多的人當中.
揀了你出來.
就是說有背後的.
那個目的的.
但是能不能夠.
反映這個身份呢.
所以其實都是身份的認同.
以色列就是要帶進.

$^{2321}$身份的認同.
實際上的操作.
似乎聖經保留.
就是說其實容許.
每個家族.
每個祭司利未人.
他們自己去.
做出他們自己良心的決定.
今天就將.
這兩章.
來做一個比較詳細的解釋.
也引發.
新約.
所以如果任何一個群體.
如果能夠.
大家當然是.
對準真理.
又願意服從真理.
又願意對自己.
這個身份認同.
不是邪教那些.
那些洗腦式的.
蠻蟲的那些不算.
如果我們越來越有.
這個這麼強烈的身份認同.
那才是一股.
很大的香港的勢力.
可以帶來基督教的復興.
如果不是.
那就.
要等待每個人自己的選擇.
那個力量.
就不強.
所以這個一定是那個問題.
為什麼教會這麼能夠.
有一個這麼大規模的增長.
原因就是.
那個領導層的.
屬靈的長輩.
一講.

$^{2361}$弟兄姊妹就很認真地.
去效法跟隨.
就會見到很多的神蹟.
很多在她們的生命裡面的改變.
很多生活上.
她們見到的祝福.
就越發一個正面的效應.
這些都是我們一路期待.
希望弟兄姊妹有.
這些的學習和成長.
所以如果神的話語.
不能夠很深化.
在弟兄姊妹的內心.
你都沒有導在你的內心.
何來你有力量去追求.
那個的身份認同呢.
你都沒有導.
你有的就是.
世界的價值觀.
很難你會認同.
神的這個標準.
譬如第十章.
大家讀完之後.
其實效不效法.
以色列你是否想被人炒魷魚.
或被人用石頭打死你.
因為帶來暴動.
涉及到人們.
拋妻棄子.
但是以色列是想.
大進回到.
神這個選民的身份.
雖然被擄.
雖然亡國.
但是神期望我們這個身份.
今天就.
我們講到這裡.
我們就祈禱.
再一次我們將.
剛才的經文都仰望著.

$^{2401}$我們很盼望這一段的講解.
和聖經的說話.
聖經的教導.
以色列的初心.
和要恢復這個.
聖文聖潔的子民.
這個.
神期望的身份.
這個的理想.
在當時.
我們不知道具體的落實.
但是我們肯定知道.
我們成為基督徒之後.
是應該在我們的身份上.
是要貼近神的心意.
因為我們是有永恆盼望.
和我們等待主耶穌.
再降臨的一班人.
盼望我們能夠在.
我們的思想行為.
我們的生活.
都能夠有主的.
豐厚恩典.
你的感動你的帶領.
使我們能夠.
在我們的生活.
去影響我們身邊的人.
影響其他未信的朋友.
讓他們看到.
我們不是故意作怪.
而是我們真是一班.
不是因為信邪教.
而是我們因為明白了.
神的真理.
我們真是一班被神拯救出來.
的一班神的子民.
以致我們真的要.
將救恩.
針對基督徒的要求.
我們敬拜的是獨一.

$^{2441}$永存.
聖潔和恨惡罪惡的神.
叫我們就能夠在.
主你的面前.
來有這個聖潔的.
渴求和追求.
求主幫助感宣家在新一年.
我們要成長.
我們在疫情之後.
都真是要來到去成為明亮的.
燈台傳福音.
我們盼望的不是有.
很多弟兄姊妹.
只是安於自己的生活.
現狀.
本來這些不是錯.
有回自己的生活的調劑.
工作的能力.
這些本來就一定是.
應該的.
但我們這個群體.
是應該反照救恩的真實.
讓世人.
讓街坊看到.
我們真是有.
一個永恆的真理.
在我們的心裡.
在我們的價值觀裡面.
求主幫助我們.
我們將今天晚上.
我們的教導仰望.
在主的手中.
很盼望聽到的人.
都能夠得到那個啟發.
和對自己的內心有反思.
在主面前.
來到去有這個等候.
禱告.
我們這樣仰望奉耶穌基督的名.
阿們.

$^{2481}$下一個禮拜.
我們就開始彌熙禮記.
如果大家有時間.
可以先看看.
我再保留.
撒迦利亞書還有八個儀仗.
我們還沒看.
其實已經橫跨了那個時代.
我們遲些都解釋.
讓大家有一個小小的概念.
\newpage



\section{使徒行傳 8:14-24}
\label{sec:jYaZ2w15SEQ}
\textbf{2023-02-25 晚堂《屬靈沒有捷徑》使徒行傳 8\_14-24 講員:陳一華牧師}
\newline
\newline
連結: \href{https://youtube.com/watch?v=jYaZ2w15SEQ}{\texttt{https://youtube.com/watch?v=jYaZ2w15SEQ}} ~~~~ 語音日期: 2023-02-28
\newline
\newline
\hyperref[sec:v6FulXmsPK8]{\small{< < < PREV SERMON < < <}}
~
\hyperref[sec:index_chronic]{\small{[返順時目]}}
~
\hyperref[sec:index_scriptual]{\small{[返順卷目]}}
~
\hyperref[sec:Lu3g7hRuGdE]{\small{> > > NEXT SERMON > > >}}
\newline
\newline
使徒行傳 8:14-24
\newline
\begin{longtable}{cl}
\hline
\hline
章節 & 經文 (和合本修訂版)\\
\hline
8:14 & \begin{tabularx}{0.7\textwidth}{X} 在耶路撒冷的使徒聽見撒瑪利亞人領受了神的道，就打發彼得和約翰到他們那裡去。 \end{tabularx} \\ \\ \relax
8:15 & \begin{tabularx}{0.7\textwidth}{X} 兩個人下去，就為他們禱告，要讓他們領受聖靈， \end{tabularx} \\ \\ \relax
8:16 & \begin{tabularx}{0.7\textwidth}{X} 因為聖靈還沒有降在他們任何一個人身上，他們只奉主耶穌的名受了洗。 \end{tabularx} \\ \\ \relax
8:17 & \begin{tabularx}{0.7\textwidth}{X} 於是使徒按手在他們頭上，他們就領受了聖靈。 \end{tabularx} \\ \\ \relax
8:18 & \begin{tabularx}{0.7\textwidth}{X} 西門看見使徒一按手，就有聖靈賜下，就拿錢給使徒， \end{tabularx} \\ \\ \relax
8:19 & \begin{tabularx}{0.7\textwidth}{X} 說：「請把這權柄也給我，使我手按著誰，誰就可以領受聖靈。」 \end{tabularx} \\ \\ \relax
8:20 & \begin{tabularx}{0.7\textwidth}{X} 彼得對他說：「你的銀子和你一同滅亡吧！因為你想神的恩賜是可以用錢買的。 \end{tabularx} \\ \\ \relax
8:21 & \begin{tabularx}{0.7\textwidth}{X} 你在這道上無份無關；因為你在神面前心懷不正。 \end{tabularx} \\ \\ \relax
8:22 & \begin{tabularx}{0.7\textwidth}{X} 你要為你這樣的惡而悔改，祈求主，或者你心裡的意念可得赦免。 \end{tabularx} \\ \\ \relax
8:23 & \begin{tabularx}{0.7\textwidth}{X} 我看出你正在苦膽之中，被不義捆綁著。」 \end{tabularx} \\ \\ \relax
8:24 & \begin{tabularx}{0.7\textwidth}{X} 西門回答說：「請你們為我求主，使你們所說的，沒有一樣臨到我身上。」 \end{tabularx} \\ \\
[1ex]
\hline
\hline
\end{longtable}
$^{1}$弟兄姊妹晚安.
我們今晚可以敬拜天父.
真是一個很大的恩典.
因為在香港的生活裡面.
大家都非常忙碌.
我們都在繁囂的社會裡面.
有時聽到很多不同的聲音,消息.
很多事情令到我們的心.
有時都會有很多的繁擾在當中.
我們能夠抽到這段時間.
來去親近天父,去敬拜天父.
這一刻我的形容.
就好像我們進入了一個心靈的綠洲.
大家認同嗎?.
叫我們心靈裡面可以藉著敬拜天父.
讓我們可以重新得力.
所以各位弟兄姊妹們.
我們能夠敬拜上帝.
感不感恩?.
很感恩.
今晚牧師和大家看的一段聖經.
半生熟都未定.
熟的原因是因為他是士徒行傳.
大家都讀過士徒行傳.
三普的地方就是.
可能這段經文你曾經立過.
但是立完之後就忘記了.
好像水過鴨背那樣.
為什麼呢?.
因為你不覺得有什麼很特別的地方.
但是如果我們讀這段聖經.
去看下去的時候.
你會發覺是有一點點不同的.
這段經文.
就是陳弟兄幫我們已經讀了.
其實這裡說的一個人物叫西門.
你記得.
此西門不同彼西門.
這個並不是耶穌門徒的西門.
是另一個西門.

$^{41}$這個人住在撒瑪利亞城.
所以我們今晚看這段聖經.
你就會發現到.
其實出現這個人物.
可以說是一個小插曲.
因為這個背景就是.
如果大家印象還是很清晰的話.
就知道當史提反被逼迫受害殉道之後.
記不記得耶路撒冷這個地方.
就起了一個很大的迫害.
於是門徒返了教會.
或者親近神的人就開始受迫迫了.
他們就四散了.
其中有一位執事.
大家聽過肥力的名字嗎?.
肥力被聖靈帶領.
就下去一個地方叫撒瑪利亞城.
這個城根本是猶太人不喜歡去的.
因為這個城裡的人.
在猶太人的觀念當中.
就覺得這座城裡的人是污煙瘴氣.
或者裡面可以說是一個罪惡之城.
是很多問題的.
猶太人覺得他們不應該進入這座城.
就算由南到北.
由北向南也好.
他們寧願繞路都不進入這座城.
但這次很奇妙.
有一個突破.
就是竟然肥力進入了這座他們討厭的城市.
就是撒瑪利亞這個城市.
當聖靈引導肥力去的時候.
肥力就很信服,很聽話.
就按著聖靈的帶領.
就踏入了撒瑪利亞城.
進入城之後.
肥力就做了三件事情.
第一 他宣講基督.
告訴城裡的人聽.
早一段時間在耶路撒冷發生的事情.

$^{81}$你們聽到有人被釘十字架.
其實這位被釘的那位人物耶穌.
就是神所差派來的救世主.
就是我們的彌賽.
所以肥力宣講基督.
將福音傳開去.
第二件事情肥力做的是什麼呢.
就是他醫治有疾病的人.
第三 他做的事情.
就是港族污鬼.
就是對一些被邪靈,被污鬼.
去綑綁的人.
肥力就奉耶穌的名港族污鬼.
他做了三件事之後.
結果整座城就起了一個很顛覆性的表現.
原來因為肥力做這三件事的緣故.
令到整個城裡的人就開始很有興趣和心.
願意去接受耶穌做救主.
所以當城裡的人相信主.
接受耶穌做救主之後.
聖經就記載在城裡的反應是什麼呢.
就是大有歡喜.
用嘉克的說話來說.
就叫做皆大歡喜.
就是說城裡的人因為聽到耶穌的訊息.
耶穌的福音.
大家都非常之的喜悅.
這是一個背景.
現在一個小插曲出現是什麼呢.
就是有一個人物叫西門.
這個西門的人物就見到彼得.
他們來到撒瑪利亞城.
因為他被教會差派進入城裡之後.
就為城裡的人祈禱.
去求聖靈降臨在信主的人身上.
很可能當時的背景就是.
那些人信是信耶穌.
不過是頭腦或者知識上去認識耶穌.
在認知的方面.
但是他們都未經歷聖靈的感動.

$^{121}$未經歷聖靈的充滿.
未經歷聖靈將恩賜賜給他們.
所以當彼得去到的時候.
就為他們祈禱.
為他們按手.
讓他們都接受聖靈.
那些人一接受到聖靈之後.
就覺得很喜悅,很有魄力,很有衝勁,很歡欣.
來過他們的生活.
很有能力的來見證上帝.
這些聖靈的功在.
當然對西門來說.
它看在眼內覺得很與別不同.
為什麼有這麼大的分別呢?.
你未祈禱,未按手.
那些人未有那些經歷.
但是你一祈禱,按手.
聖靈就降臨.
聖靈就充滿他們.
以致他們有這麼大的轉變.
所以西門看到目定口愛.
這麼厲害的?.
於是他就心懷不軌.
他想想大家都知道.
西門如果剛才你讀聖經.
你不知道他做哪一行?.
他是老騙門的.
因為他是行邪術的.
所以他就心懷不軌.
他想想不如這樣吧.
我拿一筆錢給使徒.
然後叫使徒過幾招給我.
讓我為人祈禱,為人按手.
以致我也有能力讓人得到聖靈.
大家明白這個背景吧.
你知道西門作為上帝的僕人.
使徒一聽到這番話.
真的火都燒起來了.
我說怎麼搞的?.
你以為上帝的恩賜.

$^{161}$或者聖靈的恩賜.
是用錢來買?.
你有沒有搞錯?.
所以他又說.
怎麼說呢?.
他又說了一句非常嚴厲的話.
他說你的銀子跟你一同滅亡.
這個很厲害.
這個警告很厲害.
特別發生在使徒行傳.
說出這些厲害的話.
特別要什麼?.
特別要警醒.
你記不記得沒多久前.
一對夫婦欺控聖靈.
有沒有印象?.
欺控聖靈之後.
然後彼得又斥責.
然後這對夫婦就被擊倒.
離開世界.
所以當彼得這樣說的時候.
是一個很嚴重,很嚴厲的警告.
西門就開始害怕了.
知道自己說錯話了.
於是就希望求彼得為他們祈禱.
讓他們可以免於這個刑罰.
我這樣和大家有個小小略略.
將經文的過程和大家分享.
其實我想和大家說.
就是當你聽到一個這樣的小插曲.
一個這樣的人物出現的時候.
不如我們試一下看看.
首先第一件事.
我們去了解一下西門的背景是怎樣的.
聖經告訴我們.
西門這個行邪術的西門.
第一件事是他很厲害.
他很喜歡在人面前很厲害.
因為他要人尊重他.
所以他是妄自尊大.

$^{201}$這是第一個描述他的特徵.
第二個描述他的就是迷惑人.
他用一些邪術.
旁門左道.
一些邪門的旁門左道的東西.
來叫人害怕他.
聽他的話.
控制人.
這個就是他第二個特別的地方.
他被人聖經描述他.
第三個描述他.
他也是一個機會主義者.
因為這個人他發覺哪裡有利益的時候.
他就靠近那裡.
肥膩這麼厲害.
可以行神蹟.
所以他經常靠近肥膩那裡.
為什麼呢.
希望跟他得到一些功力.
因為他自己本身是行邪術的.
所以彼得來的時候.
他又靠近彼得那裡.
換句話說.
他是一個很看準機會的人.
哪裡有利益有好處.
他就靠近那裡.
這是第三個機會主義.
第四.
聖經直接說這個人是一個心懷不正.
心術不正的人.
這個很清楚我們見到這位西門的情況.
接著下去.
當你明白整個過程之後.
你自然會問一條問題.
問一條什麼問題呢.
就是這條問題.
為什麼這個行邪術的西門.
聽到肥膩的講道又信了.
還洗禮了.
但他仍然會表現出心術不正呢.

$^{241}$大家會問這個問題吧.
為什麼他還會在彼得的面前.
說這樣的說話.
表現出這樣的情況呢.
這是你和我都很想問的問題.
就好像今天.
當大家有時候在社交場合.
或者在公司.
或者在你的寫字樓.
或者在你所接觸的左鄰右里.
有時候你見到一些人.
他說他是基督徒.
但似乎他做出來的行為.
有些不多似.
有時候都問一個問題.
就是說這個人的信仰和行為不相稱.
為什麼呢.
甚至有些人說.
我都差點被他叛賭.
為什麼呢.
因為可能見到一些.
好像說是信耶穌.
但他的表現好像不像.
各位電視節目的朋友們.
我們問這個問題的時候.
就是說我們都懂得去提醒自己.
去警惕自己.
為什麼會出現這樣的問題呢.
我想答案大家都很想知道.
原來最基本的問題是什麼呢.
就是信耶穌的過程裡面出了問題.
信耶穌的過程裡面出了問題.
有些人說他信耶穌.
但他什麼原因要信耶穌呢.
為什麼要信耶穌呢.
可能第一個問題出在什麼呢.
可能他有不同的目的.
不同的目的.
正如西門他想像肥利或者彼得那樣.
他們這麼能夠行神蹟騎士的.

$^{281}$我都要很辛苦地行邪術.
才令到那些人來聽我說.
如果我黏在一起.
我又說信.
我又說接受耶穌.
又洗禮.
那我也有些得益的.
所以我們可能想一下.
會不會他的動機.
是因為這樣的緣故而去信耶穌呢.
這是第一個分析.
第二個分析.
會不會他發覺到.
信耶穌是什麼呢.
是因為看到整個城裡的人.
大部分都願意決志.
個個都很熱心去信耶穌.
會不會這個時候我也要去信呢.
如果不信我好像撿書一樣.
因為人人都信.
所以好像臭葉撈一樣.
這麼多人都願意信.
我也信了.
這個可能是第二個問題也未定.
第三個問題.
可能是一個為人而去信.
有時候有些人.
他信耶穌的過程是什麼呢.
是因人而去信.
可能他要給面子某個人.
他會信.
譬如舉個例子來說.
可能我的老闆.
他是基督徒.
所以我在那條團隊.
我都覺得信耶穌多好.
都沒有撿書.
你明白我的意思嗎.
可能是因為人而去信.
又或者可能你的老友記.

$^{321}$四十年同學.
五十年朋友.
原來他都是信耶穌的.
那我不如我也去信吧.
可能這是一個什麼.
一個因人而去信.
各位我們就要留意一下.
這個西門.
他是不是因為某些人.
他很崇拜某些人.
以致他這樣去信耶穌.
值得我們檢討一下.
第四個原因.
可能還有其他的原因.
我都沒有寫出來.
還有可能其他的原因而去信耶穌.
以致他在人面前說.
我都是基督徒.
我都去信耶穌.
不過原來在他信主的過程當中.
第一步就已經出了問題.
可能是因為這樣的緣故.
所以我們現在明白.
為什麼西門他在彼得的面前.
會說這樣的說話.
而有種這樣心術不正的表現.
這個是幫我們找到答案.
所以各位電影節目們.
最後我們都會問.
我怎樣去信耶穌.
是一個合上帝的心意.
或者我怎樣信耶穌.
令我能夠真的重生得救呢.
我們要回到聖經最基本的基礎上.
聖經給我們的基礎是很重要的.
就是說一個人信耶穌.
他要經過什麼.
他要經過神的話語.
還有聖靈的感動光照之下.
他在上帝的面前認罪悔改.

$^{361}$決志信主.
在這樣的過程裡.
才是一個真正的重生得救.
說到這裡.
大家或者都很快去想起.
約翰福音.
記得有一次.
耶穌對尼哥底姆和他的對話.
你們有沒有留意.
耶穌知道尼哥底姆.
他是猶太人的官.
他是有學識的.
他是有他的專業的.
他是與眾不同的.
他懂得很多東西的.
所以當耶穌和尼哥底姆對話的時候.
他沒有跟他行圈幾句.
沒有跟他兜了兩個圈.
然後有了開場白就說.
不是.
耶穌跟尼哥底姆說話的時候.
劈頭就說了.
他說什麼.
他說人若不重生.
就不能進神的國.
大家記得嗎.
尼哥底姆怎樣回應.
他的理解就是說.
這樣啊.
他說我都已經老了.
我還要回去媽媽的肚子裡.
去生出來.
那叫重生嗎.
然後耶穌怎樣回應.
耶穌回應他.
耶穌說.
從肉體生的就是肉體.
從靈生的就是靈.
人若不是從水和聖靈生的.
就不能進神的國.

$^{401}$是不是說得很清晰.
換句話來說.
一個人他能夠重生得救.
他首先會聽到神的話.
聖經上帝的道.
在他聽到的過程當中.
得到聖靈的光照.
感動.
以致他看到.
神的話是一面鏡子.
反照他的人生.
他發覺自己真是很有問題.
所以他就願意謙卑.
接受耶穌基督成為他的救主.
所以各位弟兄姊妹們.
這個過程.
就是一個真正去決志.
信耶穌的人.
他能夠經歷重生得救.
是必須具備的條件.
所以各位弟兄姊妹們.
當我們明白之後.
我想和大家分享.
今天你和我.
信了耶穌 決了志之後.
應該怎樣.
今晚的講道.
那個戲玉就在這裡.
就是說.
我們決了志 信了耶穌.
成為上帝的兒女之後.
記得 我們千萬不要單打獨鬥.
就是說我們不要以為.
我信了耶穌就行了.
很安全了.
我拿到天堂的簽證.
護照都拿了.
我就很淡定了.
千萬不要這樣想.
當你決了志 信了耶穌之後.

$^{441}$你是離開了教會的群體.
離開了神的家.
生活的時候.
沒有了教會生活的時候.
其實我們是很危險的.
所以弟兄姊妹們.
就好像一個BB.
他出生之後.
好像很健康.
很胖.
很可愛.
又多撕幾下.
但是如果再久一點.
再久一點.
他沒有奶喝的.
沒有水給他喝的.
沒有食物供應給他的.
我想很多人都知道.
這個小朋友 這個BB.
是很難再生存下去的.
為什麼.
因為他沒有食物.
沒有養分 沒有養料.
去餵養他 去成長.
所以同樣.
我們每一個決了志.
重新得救的基督徒.
我們都需要有教會的生活.
在神的家裡面.
我們一齊地去成長.
一齊地互動.
一齊地提醒 勉勵.
這個是很重要的.
所以我們不要單打獨鬥.
否則我們很容易.
我們的信仰的根基.
就會被這個世界 潮流.
各樣的風氣侵蝕了我們.
所以我記得以前.
在警察團契的時候.

$^{481}$當我做團務的時候.
我經常都提醒弟兄姊妹.
記得你不要單打獨鬥.
不要以為自己可以.
因為你要明白.
你要知道我們身處這個世界.
這個世界是一個魔鬼撒旦.
掌控的世界.
記不記得《伊忽所書》怎麼說.
《伊忽所書》說.
我們並不是與屬血氣的爭戰.
記不記得.
乃是什麼.
乃是與屬靈氣的惡魔來爭戰.
它是管轄這個幽暗世界的.
所以我們要加倍地.
在靈命上面.
來接受各方面.
給我們的培訓和栽培.
我現在就想把它分為三個階段.
我希望弟兄姊妹們.
你都能夠進入這三個階段.
第一就是初階.
初階是什麼呢.
初階就是當我們重新得救之後.
我們需要在教會當中.
繼續學習神的話語.
所以我們上主教學.
上茶經班.
我們上祈禱會.
上小組.
甚至上團契.
我們種種的參與活動.
這個行動就好像在吸收屬靈的養分.
我們不斷地有這個培訓.
有這個裝備.
這個是很重要的.
所以弟兄姊妹記得.
我們不是說.
信了耶穌就完成了我的責任.

$^{521}$也不是說我洗了禮.
完了我的洗禮班.
我就畢業了.
不是.
記得.
那個是一個開始.
所以有人說.
畢業其實是什麼.
其實是開始.
另一個階段的開始.
我們在教會當中.
特別我們在錦宣家.
上帝給我們很多很多屬靈的栽培.
我們有不同的主教課程.
我們有不同的一些茶經聚會.
很寶貴.
我很希望弟兄姊妹們.
很努力地學習.
努力地吸收屬靈給我們的養分.
不是單單我們一次的崇拜這麼簡單.
這是一個很重要的.
是一個初階.
好了.
中階是什麼呢.
中階就是說.
當你和我在教會裡面.
得到很多的栽培.
很多的訓練裝備之後.
記得.
你和我是要行動出來.
就是要去實踐.
怎樣去實踐呢.
原來天父真的很疼我們.
當我們接受栽培之後.
其實聖靈老早已經內住在我們的心裡面.
聖靈就是那個恩賜分配的源頭.
聖靈按照我們個人的需要.
將恩賜給你和給我.
所以你和我就可以透過上帝給我們的恩賜.
去服侍上帝.

$^{561}$去服侍人.
去到不同的地方.
可能在教會.
可能在你的家.
可能在你的社區.
你能夠將神給你的恩賜財綱.
去好好地操練出來.
各位聽眾.
如果我們信了耶穌之後.
沒有事奉.
沒有將恩賜操練發揮出來的時候.
其實我們的靈命成長都會有些限制.
這個限制是什麼呢.
就是在中途就窒礙了.
唯有你真的經過實踐.
經過去事奉.
經過來和人.
和弟兄姊妹一起的合作配搭.
這樣來學習.
將恩賜互相配搭建立耶穌基督的身體的時候.
你就會發覺在靈命上面.
你會有很明顯的進步.
因為你將你所學的東西.
實踐了出來.
發揮你的行動.
我記得以前在醫院裡面.
我們都有一個課程.
是給那些遠目義工去上的.
上完課程.
上完八堂之後.
我還記得好像上八堂那麼多.
上完之後.
還要有一個月時間.
進病房去實習.
兩樣加起來.
合格了.
才會給一份義工證書給他.
我記得有一位義工.
有一次在病房.
實習完下來.

$^{601}$心情不是很開心.
看得出的.
一個人開心不開心的樣子看得出的.
為什麼他去完病房下來.
那麼多重擔呢.
為什麼好像遇到一些很不開心的事呢.
我就跟著問這個弟兄.
你剛才是不是去了病房.
他說是.
你是不試探了病人.
他說是.
我說為什麼你探完病.
好像很多擔子.
很多重擔那麼不開心呢.
接著他就告訴我.
他說牧師我剛才去探一個病人.
我都猜不到.
那個病人問我一條問題.
問到我張開嘴巴就有個洞.
不懂回答.
那時候我覺得很醜怪.
希望有個洞挖到我們就鑽進去就好了.
原來他發覺.
他想著去傳呼音給病人.
誰不知病人反問他一條問題.
他完全不懂回答.
張開嘴巴就有個洞.
那條是什麼問題呢.
我現在不告訴你.
有機會我才告訴你那條什麼問題.
這個弟兄告訴我之後.
我知道他有些挫折.
所以要鼓勵一下他.
叫他不要灰心.
繼續努力.
我也鼓勵他.
不如這樣.
你試一下回教會繼續上主一學.
或者回查經班.
希望你能夠掌握多些聖經的教導.

$^{641}$以致到別人再問你問題的時候.
你懂得怎樣去回答別人.
這個弟兄真是好聽話.
果然他由那次開始.
再回教會的主一學.
回教會學習聖經的課程.
給我的一個經歷就是.
原來這個弟兄.
何時知道他自己學的東西不夠呢.
何時他知道自己的靈明有不足呢.
不是坐在這裡嘆世界的時候.
而是他真的去實踐的時候.
當他真的去事奉去實踐的時候.
他就遇到原來我只有半桶水.
或者原來我真的有很多東西都不懂.
原來我都是很不足.
當一個人發覺自己不足夠.
那就好了.
他就會開始去想.
我怎樣可以去補足我的不足呢.
各位聽眾們.
所以上帝給我們第二個階段的成長.
是一定要透過你的事奉.
如果你沒有去事奉上帝的時候.
你所學到的東西.
亦都是停於一個知識的層面.
你的腦袋的層面.
因為你沒有將它拿出來.
經過考驗,經過實踐.
所以我們的靈明成長.
第二個階段是要實踐.
是要將我們的恩賜洗出來.
和弟兄姊妹互相配搭.
一起去事奉.
我們就發覺自己.
有哪些更加需要去學習的地方.
最後.
第三個階段.
進階.
進階就是.

$^{681}$原來當我們經過一輪事奉的時候.
我們發覺自己還有一些不足的時候.
我們就會繼續去追求.
我們的靈明的培訓.
所以進階就是告訴我們.
原來每一個教徒.
我們的屬靈生命的本質.
在彼得後書已經告訴我們.
有一個很好的樣板.
大家知不知道.
彼得後書一章五到十一節告訴我們.
原來基督徒有一個美好的追求生命的本質.
是有八個階段的.
有沒有印象?.
我說一點點你就會記得.
聖經怎樣說?.
彼得後書說什麼?.
有了信心就要加上什麼?.
德行.
有了德行又要加上知識.
知識又要加上節制.
有沒有印象?.
一路操練下去.
是有八個階段這麼厲害.
我叫那段經文叫天國八部.
有沒有聽過天國八部?.
天龍八部你聽過.
天國八部你可能沒聽過.
我今晚介紹你認識.
學一下天國八部.
當你能夠掌握天國八部的時候.
你的靈命的進階是更加豐盛.
更加有光彩的人生.
來見證去榮耀上帝.
所以各位電視節目們.
今天晚上我們看到這段經文的時候.
我們再一次提醒自己.
原來屬靈有沒有捷徑?.
沒有捷徑.
屬靈的生命一定要什麼?.

$^{721}$要不斷地學習.
不斷付出時間.
付出我們的行動.
付出我們的努力去裝備.
然後在碰過釘之後.
我們的靈命就會開始有更加好的成長.
在座當中可能有很多弟兄姊妹.
都很喜歡聽一些音樂會.
或者聽一些管弦樂團的演奏.
喜不喜歡?.
喜歡的舉舉手.
有多少人喜歡聽管弦樂團的音樂會?.
大家都經常說的.
台上十分鐘台下多少?.
台下十年功.
其實沒有捷徑的.
你看人家表演得這麼精彩.
拍爛手掌.
拍按叫絕.
但是你知不知道人家在私底下的操練.
是用了多少精神,多少時間呢?.
各位電視節目沒有捷徑的.
又或者我說些動態的.
不要說這麼文靜.
動態一點.
有沒有弟兄姊妹在這裡喜歡跑步的?.
舉舉手.
喜歡跑步的.
特別跑馬拉松的.
有沒有?.
馬上縮手下來.
弟兄姊妹們.
跑馬拉松的人.
用不著操練嗎?.
不是拿個號碼掛在這裡就去跑.
不行的.
拿個號碼掛在這裡就去跑.
很危險的.
跑到中間可能會休克都不一定.
要不要練習?.

$^{761}$要練習的.
所以記得.
任何的靜態,動態.
所有的活動.
都是沒有.
你要拿到好的成績.
都是沒有捷徑的.
同樣我們靈明的生命的成長.
屬靈的生命的那種精彩.
那種豐盛.
都是沒有捷徑的.
求上帝幫助我們.
我們抵抗我們祈禱.
親愛的主.
多謝你讓我們今天從一個經文.
寄了一個反面的教材.
而又提醒我們.
真是靈明的成長.
是需要付出代價.
是需要不斷地努力學習.
我們知道是沒有捷徑的.
求上帝讓我們眾弟兄姊妹.
不斷地努力.
在面前日子當中.
我們繼續來好好地.
把握不同的機會.
裝備自己.
以致我們能夠被主好好地使用.
成為多人的祝福.
奉主名求.
阿們.
\newpage



\section{帖撒羅尼迦前書 5:16-18}
\label{sec:Lu3g7hRuGdE}
\textbf{2023-02-26 加恩樓《時刻經歷主的同在》帖撒羅尼迦前書 5\_16-18 講員:李適清博士}
\newline
\newline
連結: \href{https://youtube.com/watch?v=Lu3g7hRuGdE}{\texttt{https://youtube.com/watch?v=Lu3g7hRuGdE}} ~~~~ 語音日期: 2023-02-28
\newline
\newline
\hyperref[sec:jYaZ2w15SEQ]{\small{< < < PREV SERMON < < <}}
~
\hyperref[sec:index_chronic]{\small{[返順時目]}}
~
\hyperref[sec:index_scriptual]{\small{[返順卷目]}}
~
\hyperref[sec:d9OO9rhlq8k]{\small{> > > NEXT SERMON > > >}}
\newline
\newline
帖撒羅尼迦前書 5:16-18
\newline
\begin{longtable}{cl}
\hline
\hline
章節 & 經文 (和合本修訂版)\\
\hline
5:16 & \begin{tabularx}{0.7\textwidth}{X} 要常常喜樂， \end{tabularx} \\ \\ \relax
5:17 & \begin{tabularx}{0.7\textwidth}{X} 不住地禱告， \end{tabularx} \\ \\ \relax
5:18 & \begin{tabularx}{0.7\textwidth}{X} 凡事謝恩，因為這是神在基督耶穌裡向你們所定的旨意。 \end{tabularx} \\ \\
[1ex]
\hline
\hline
\end{longtable}
$^{1}$太英姐妹早安.
很高興來到這裡跟大家一起敬拜.
先說說我有家族遺傳的手震.
所以拿著咪高峰你看到我震.
如果你在醫護界你不用想我什麼病.
這是我爸爸遺傳給我的.
我會形容是保羅所說的一根刺.
小時候我拼命祈禱.
神都沒有拿走.
我擁抱他.
他就是我的一部分.
不過上帝很特別.
我小時候覺得我沒可能是拿咪高峰的人.
我沒可能是教書.
我小時候很害羞.
我很怕在很多人面前說話.
教書還要天天跟很多人說話.
我覺得我一定不行.
但是我的年輕人告訴我.
上帝很屈機.
輾轉地.
我今天做了一個信學院老師.
還要經常拿咪高峰.
我手震不要影響大家聽上帝的話就行了.
我想這兩三年無論你是什麼環境崗位都好.
你都經歷了很多變化.
最近跟弟兄姊妹聊起.
她們都說有不少同事可能離職.
有些去外地.
有些轉工.
有一位姐妹她做人力資源.
她跟我分享請人請到怕.
因為請一個就走一個.
有一次老闆問她.
你請這個會不會走.
她心想我怎麼知道.
不過另一邊廂又有很多弟兄姊妹因為疫情.
沒有工作.
又轉職.
甚至是失業.

$^{41}$在這些時候.
我不知道你會不會想上帝的大令在哪裡.
我們其實都很想回到以前.
很想復常.
我們開始都說我們復常了.
但我經常想復常其實不是回到以前.
而是我們迎接新的一頁.
我經常覺得上帝好像在搞動這個世界.
我們都知道如果不是上帝允許.
這一切都不會發生.
但他就是在搞動這個世界.
以致人可以從我們安樂和安恕的狀態走出來.
上帝就是這樣.
在我們的人生平台.
在我們的工作,家庭.
不同的處境裡面.
我們就是會遇到不同的衝擊.
當我們聽到我們的信仰受到衝擊的時候.
這些時候才是挑戰我們是否真的相信.
我們相信得多多.
是否真的.
這幾年我們都知道世界突然間可以完全不同了.
今天的經文我選了三節.
很簡單的經文.
我覺得可能你都會背.
聖經說要上上喜樂.
不住禱告凡事謝恩.
因為這是神在基督耶穌女向你們所定的旨意.
經文裡面有三個很明顯的重點.
喜樂,禱告,謝恩.
其實這三樣在我們的生活裡面.
是環環緊扣是分不開的.
首先我們看第十六節.
上上喜樂.
早期的信徒.
因為他們的信仰和四周的民族其實都很不同.
因為這樣他們經常都被人騷擾,排斥,奚落.
他們周圍的人都會覺得.
這班人不可能快樂.
不過他們又好像與別不同.

$^{81}$因為他們內心裡面好像有一股力量.
令到他們面對困難,面對苦難的時候.
仍然可以有喜樂.
就好像在耶路撒冷的使徒.
當他們受到刑罰.
甚至當時的人捉了他們在會堂裡面.
他們被鞭打完.
《士堂傳》五章四十一節形容.
「他們離開公會,心裡歡喜,因被算是配為這名受辱」.
這班人很特別.
喜樂就是聖靈所結的果子之一.
就是我們作為天國的兒女其中一種特質.
第十七節是不住的禱告.
意思是無論何時何地.
我們都在聖靈裡面.
我們都有與神的互動的禱告.
這裡也令我們想起耶穌的教導.
在路加福音十八章一節裡說.
「要常常禱告,不可灰心」.
不住地禱告.
不是甚麼事都不理會.
每天就在這裡祈禱.
而是在我們所有日常的活動.
無論你做著甚麼.
都可以有這個禱告的心,禱告的精神.
又或者說無論做著甚麼.
你都能夠經歷到神與你一起.
有這種經驗和感覺.
究竟在我們每一天的生活工作如此忙碌.
怎樣敏銳有上帝同在呢?.
怎樣有這種喜樂去經過呢?.
十八世紀在法國.
有一位基督徒弟兄.
他叫做勞倫斯(Lawrence).
他可以說是這方面的專家.
他小時候家裡很窮.
所以他很小的時候就加入了軍隊.
在他十八歲的冬天.
年輕的勞倫斯.
他看見一棵光滑的樹.

$^{121}$沒有果子也沒有樹葉.
他就想.
春天來臨的時候.
這棵樹會變成怎樣呢?.
就在那一刻.
他心裡好像突然經歷到一種改變.
他突然間很清楚地看見.
原來上帝就是這樣看顧大地的.
就是這樣上帝是有充滿能力的.
春天的景象就會完全不同.
這種他對於上帝臨在的一種敏銳感.
成為他餘生裡面不停的一種經驗.
通常這些經歷.
我們叫做信主得救的經驗.
其實我們每個人都不同的.
每個人都會有的.
但是當這位弟兄經過這件事之後.
不久後.
他就因為當時有一個嚴重的重傷.
令到他有坐骨神經痛.
被迫離開了他軍隊的生活.
你也知道當時的醫學沒有今天那麼發達.
所以他的坐骨神經痛.
令到他走路都好像跛下跛下.
他餘下的一生.
都是活在這種痛苦的底下.
過了幾年.
勞倫斯就以一個普通信徒的身份.
因為他很窮又不知道怎樣生活.
他就進了一間修道院.
他餘下的一生就在這間修道院裡面.
做甚麼呢?.
因為他沒有學歷.
他做不到修士.
他就在廚房裡面煮東西,清潔.
另外補鞋,造鞋.
就在他這個工作的過程裡面.
他就建立了經常敏銳.
上帝同在的一種習慣.
雖然他在這個修士的群體裡面身份卑微.

$^{161}$但是他晚年的時候.
很多人從不同的地方裡面.
去找他接受靈修指導.
聽他分享怎樣不住的禱告.
時刻經歷上帝的同在.
勞倫斯說很可惜.
其實我們的信心是很少的.
我們做很多宗教活動.
而且不斷找一些新的活動.
但是沒有一天是我們要以信心來生活.
大家知道他的環境就是這樣.
修士每天早午晚.
很多時段去祈禱.
但是在他心裡面.
他覺得為何這些時段才叫做祈禱呢?.
他都很馴服.
他都會去參加.
但是他參加之餘.
他就會覺得其實每一天.
他實踐上帝的心意.
就會被他滿足.
所以無論是所謂的世俗.
或者所謂的屬靈的事情上.
他都很堅持將自己交給上帝.
無論他很舒適.
或者是經歷痛苦的時候.
他都是一樣的.
你記得他經常坐骨神經痛.
所以他是很不容易的.
他說在屬靈枯乾的日子.
就是我們感受不到上帝.
覺得他很遠的日子.
甚至我們不想祈禱.
那些時候是我們更加要用心用力去盡忠和毀身的時候.
因為這些情況會讓我們成長.
勞倫斯說他是被愛所管理.
做什麼都是對著上帝的愛.
不會想自己.
他為了這樣.
他覺得他可以為主去做很多不同的事情.

$^{201}$包括在地上撿起一件垃圾.
都是為主去做的.
因為沒有其他事情.
不是為了得到神的恩典才這樣做.
他說當他這樣想的時候.
是令他很快樂的.
意思是說.
他要將平時我們以為為自己做的事情.
改變成為為上帝去做的事情.
每一天他在廚房工作的時候.
他都會先讓自己的心思意念充滿著對上主的毀身.
然後他就想他的工作要怎麼做.
什麼時候做什麼.
然後在工作的前後過程當中.
他都在禱告的裡面.
我想起我自己.
或者對我們來說.
工作的時候很自然.
就是想想該事情要怎麼解決.
想想要怎麼做.
最好是快點做完才可以下班.
遇到困難的時候.
我們通常都會很理性地想想有什麼方法.
如果搞不定怎麼辦呢.
搞不定心裡有些埋怨.
為什麼會這樣.
然後就祈禱.
祈禱叫上帝給我另一份工作.
我們的祈禱很少會沒有事沒有幹.
我們去祈禱讓上帝介入.
有時候在困難當中.
我們不太覺得上帝可以解決前面的問題.
的確上帝也不是要來解決問題的.
但是上帝可以給我們力量去面對這個問題.
所以我們也可以嘗試一下.
勞倫斯的這種不住的禱告.
他工作之前的禱告是這樣的.
主啊,你與我同在.
我現在必須要順服你的命令.
去做這些外在的事務.

$^{241}$求你賜給我有恩典.
去繼續經歷你的同在.
為此求你幫助我.
閱納我做的工作.
佔有我一切的情緒.
然後在工作過程當中.
他就會繼續求神去賜下恩典.
將他所能做的來呈獻.
工作完了.
我覺得這件事也很特別.
因為他工作完之後.
他也會想一想.
去檢視一下自己.
可能今天做得很好.
有很多經歷很好的事.
他就會向主感恩.
但是有時候.
就不是這樣.
覺得自己做不到.
做錯了.
他就會求主寬恕.
寬恕完之後.
他會放下.
然後他說.
繼續經歷主的同在.
所以勞倫斯說.
當他每一次跌倒起來.
其實他都在不斷更新.
自己的信心和愛心.
這時候.
很難不去想神.
他就是這樣與神同行.
在過程當中.
找到很大的祝福.
所以他很樂意和人分享他的看法.
尤其是他的生命實例.
更加有影響力.
他的好朋友怎樣形容他呢.
這個人的生命中.
有一種美麗和平靜的特質.

$^{281}$他很自然就會感染他身邊的人.
就算他在廚房工作的時候.
很忙亂的時候.
他都表現到很穩陣.
有一種平靜自然的感覺.
弟兄姊妹.
你平時上班,上學.
買菜,煮飯.
與家人相處的時候.
是一個怎樣的人呢.
上帝在你的生命中.
令你與其他人有什麼不同呢.
我就想起.
有一次我去一間教會分享.
講完道之後.
有一位中年女士就來跟我聊天.
她說:真是好啊.
教會都講一下我們怎樣面對工作的困難.
她跟我說:我是一間公司裡.
做茶房和清潔的.
即是那些倒茶的姐姐.
我就想.
咦 現代版的勞倫斯.
大概都是在做我們都市人覺得很平凡的工作.
我經常說這一類的工作有一個特點.
就是你做了.
沒有人會察覺.
沒有人會覺得你做得很好.
但你不做.
就有很多人罵你.
所以他們面對的都是.
一些這樣的狀況.
我們就知道.
在過去這兩三年.
疫情底下.
其實清潔工人實在是值得我們敬重的.
只不過她說.
她的公司很大.
有幾層樓.
她的上司就很好.

$^{321}$就訂了一個更表.
你什麼時候去哪層樓工作.
你跟著更表就可以了.
因為有一次.
我去到一個地方.
正在工作的時候.
有一個經理來找我.
叫我立刻去另一層.
幫我做這些工作.
她就說.
現在我還沒有得著.
不如我之後幫你做吧.
她就說.
誰知道這個經理立刻板起臉.
還整個月都不理她.
其實不少這類的事.
真的有苦自己知道.
不過她繼續說.
雖然她不理我.
不過我是一個基督徒.
我堅持每天早上.
都會繼續跟她說早晨.
因為我覺得人不是為了什麼.
就是為了要做神喜歡我做的事.
姐妹很清楚她自己的身份.
就是在她日常工作裡.
其實她做了一些.
沒有什麼人會覺得.
沒有什麼我們認為偉大的事.
但是我覺得很偉大.
因為她能夠認定.
她是上帝的女兒.
她就是生活在上帝的心意當中.
常常喜樂.
不住禱告.
還有什麼.
凡事謝恩.
不是口頭禪.
而是當你不住禱告.
你發現上帝經常都跟你一起.

$^{361}$又經常都這麼喜樂.
很自然心裡會感恩.
在聖經裡.
感恩和祈禱.
常常放在一起.
意思是當我們向上帝祈求的時候.
我們經常都會先向祂感恩.
認定這一切都是上帝所賜的.
是我們內心很自然地.
對上帝恩典的一種反應.
這三個小句還未完.
因為還有一個小節的句子.
「因為這是神在基督耶穌裡.
向你們所定的旨意」.
上帝在我們身上所定的旨意.
「這是」其實是什麼呢.
原來這三節的經文.
在經文裡.
是順著上一章所說的.
是貼在第四章.
在第三節裡.
一早就直接說了.
什麼是上帝的旨意.
那裡說「神的旨意就是要你們成為聖傑」.
然後第四章的段落.
就處理很多不同的生活範疇的應用.
到了第四至五章.
在當時的教導裡說完很多不同的生活範疇.
到了第五章23節是這樣說的.
「願賜平安的神親自使你們傳言成聖」.
所以我們知道第四至五章.
是說生活所有的範疇.
都包括在成為聖傑的裡面.
而「常常喜樂,不住禱告,凡事謝恩」.
就在中間第五章16至18節這裡.
這裡是聖靈在我們身上的工作.
就是要我們在生活的經歷裡成長.
目標是要我們成為聖傑.
屬於上帝.
弟兄姊妹.

$^{401}$我不知道你們是否想每一天都經歷主的同在.
我們都很需要很自然很簡單地.
就在跟隨主耶穌的路上.
去經歷主耶穌的大靈.
是很自然的.
不是特意去找什麼.
因為是耶穌已經應許了.
祂是與我們同在.
記得當我在神學院讀道學碩士.
想在教會做傳道人的時候.
我都跟其他神學生一樣.
很想找找上帝你做的這個我.
你在我身上有什麼心意呢?.
當時我只是想我完全奉獻給神.
神你用我吧.
所以祂擺我去哪裡呢?.
是否去宣教呢?是否去牧養呢?.
我都不是很清楚.
神你就作主吧.
當時的老師告訴我們.
其實我們每個人的生命裡都有一個主題.
就好像亞伯蘭的名字叫做亞伯拉罕.
後來神告訴我們他的名字叫做以色列.
我們在上帝裡面的名字.
是幫我們發現祂做我們原本的獨特的意思和使命.
有時候這一類所謂的負擔.
好像一團火一樣.
我們不知道怎麼去做.
甚至覺得做不到.
但還是很想去做.
因為上帝給了我們這個呼喚.
叫我們去做一件這樣的事.
所以當我進了神學院.
我就一直在尋問.
我曾經用了一個月的時間.
上帝,我在你裡面的所謂的名字是什麼呢?.
當時大概都祈禱了一個月.
有一天,在神學院的一個早會上.
有一位專門研究信仰與藝術的老師.
他分享了舊約聖經裡出埃及記的一位人物.

$^{441}$叫做比薩列.
比薩列就是建會幕的人.
他被神揀選建會幕.
神告訴他有很多人會和他一起.
很多的工匠都有精湛的手藝.
來和他一起建會幕.
經文在《出埃及記》36章1節裡是這樣說的.
那裡說:比薩列和阿荷尼亞伯.
並一切心裡有智慧的.
就是蒙耶和華賜智慧聰明.
叫他知道做聖所各樣使用志工的.
都要照耶和華所吩咐的做工.
這個時候很奇怪.
我心裡就湧出一個感動.
就像是說你的名字就是比薩列.
我心裡想:比薩列?根本不可能的.
我自小家族遺傳的手震.
所以我拿著咪高峰震震.
沒辦法做任何的手工藝.
從小開始畫畫.
尤其是那些毛筆字都搞不定.
所以我從來都不是一個懂得做藝術的人.
很快我就下了一個定論.
一定是我想錯了 想多了.
我的名字不可能是比薩列.
沒多久 吃完中午飯.
下午又上課了.
今天上課的是《舊約聖經》.
《道論》 另一位老師說的.
神學院精彩的地方就是有不同專長的老師.
很多不同的角度去看經文.
看一件事情 令事情都立體起來.
那位老師就告訴我們.
比薩列不是藝術家.
因為他手裡所做的事.
大至繪幕整體的設計.
小至每一件很微小的物件.
每一個細節都是按照上帝的吩咐去做的.
《聖經》裡很清楚地形容了起繪幕的要求.
而比薩列就是一個執行的人.

$^{481}$他要如何呢?.
他要分毫不差地完成上帝的指示.
所以比薩列是一個按照上帝的心意去做的人.
而不是一個藝術家.
這個時候我心裡有種茅塞頓開的感覺.
因為上帝如果要我做一個藝術家.
我做不到.
但是如果要我按照你的心意去做每一件事.
要分毫不差.
這個我可以努力去試.
我終於明白了比薩列這個名字的意思.
不過原來這個答案也是一連串問題的開始.
一個建造繪幕的人.
那是怎樣呢?.
仍然很抽象.
後來我畢業了.
有一年時間我回到自己的舞會做傳道.
有位認識我的姐妹就跟我說.
很開心你回來做傳道.
我聽過你的故事.
你就是一個建立教會建造聖殿的人.
當時我就想.
建立弟兄姊妹的生命和大家同行.
也是我很想做很開心的.
不過建造繪幕不是建造聖殿.
好像還有點距離.
直到十幾年前.
我開始在神學院教書.
我教的是職場神學.
在工作中我們怎樣去活出信仰.
我終於明白什麼叫做建造繪幕.
原來我是要做一個建造生命繪幕的人.
在曠野漂流的日子.
以色列人是沒有固定的居所.
上帝的繪幕就在他們當中.
繪幕上有雲柱,火柱.
帶著他們何時停,何時行.
要去到哪裡.
以色列人建造繪幕是要潔淨繪幕.
然後在那裡獻祭.

$^{521}$全部都是為了耶和華要臨在他們的群體當中.
在《出埃及記》說到.
耶和華說:我要在哪裡與你們相會和你們說話.
我要在哪裡與以色列人相會.
繪幕就要因我的榮耀成為聖.
所以以色列人和耶和華相遇的地方.
臨在的繪幕是要成聖的.
繪幕不是一個地方的建築物不動的聖殿.
繪幕就好像保羅所說的帳棚一樣.
一個為敬拜上帝而建的帳棚.
可以建造,拆卸,四處遊走.
在《新約聖經》中.
保羅也用帳幕來比喻我們的身體.
說:我們原知道我們在地上的帳棚若拆毀了.
就必得神所造的天上永存的房屋.
保羅用帳棚來形容我們在地上的身體.
這個舊我是要過去的.
以至我們穿上新的房屋.
另一位仕徒彼得.
在他晚年快要離世的時候.
他也這樣做過比喻.
在《彼得後書》第一章說.
我以為應當趁我還在這帳幕的時候提醒你們.
因為我知道我脫離這帳幕的時候快到了.
帳幕不只是一個地方.
它可以四處移動.
就像舊約聖經中的繪幕.
四處移動.
但上帝的榮光,榮處,火處就臨在其中.
繪幕遊走的地方是曠野.
一個乾涸無水的地方.
在那裡上帝會帶領前行.
上帝臨在哪裡呢?.
就是按照祂的心意來建造的繪幕當中.
原來是這樣.
原來弟兄姊妹的生命就像繪幕一樣.
聖經很詳細地記載繪幕裡要擺放甚麼東西.
建造的方法,尺寸等.
就好像聖經給我們上帝的心意,祂的話語.
我們的生命就是要按照上帝的說話來建造.

$^{561}$然後上帝要親自臨在我們的生命當中.
聖靈就是這樣與我們同在.
帶領我們去到社會不同的角落.
各個行業的裡面.
那些地方是人心靈的曠野.
人被聚聖扭曲.
是乾涸無水.
我們遇到很多與我們價值觀不同的人.
但是上帝說我要在那裡與以色列人相會.
繪幕就要因我的榮耀成為聖.
要讓祂的榮耀去潔淨繪幕.
就好像我們的生命一樣.
被耶穌的寶血潔淨.
要成為聖潔.
弟兄姊妹.
上帝就是要讓你的生命被寶血潔淨.
讓你分別為聖.
成為聖靈光照和臨在的居所.
游走在這個世界好像很乾涸無水的人群當中.
成為上帝彰顯榮耀的出口.
你願不願意這樣做呢?.
你願不願意讓你的生命屬於上帝.
好像神聖的繪幕一樣.
在黑暗的世代.
讓上帝的榮光臨在.
讓祂的運處火處帶領你前行.
保羅的工作就是織帳棚.
其中一件事.
他織帳棚織出了一個道理.
或者你的工作也給你一些不同的啟示.
曾經有一位從事保險的弟兄跟我說.
他的使命就是要幫人反思和規劃人生.
很多做保險的目的都是想賣到帳單.
不喜歡做售後服務.
因為很麻煩,沒錢賺.
但我認識很多基督徒的保險從業員.
他們跟不同的家庭同行.
他們不單著重賣到適合的保險.
而且當這個家庭有患難,有危機的時候.
如何讓他們的保障能夠幫助到一個家庭.

$^{601}$與他們同行.
有一位姐妹跟我說.
她的使命就是要擁抱身邊的人.
讓上帝的愛可以彰顯出來.
她後來已經成為一位牧者.
現在做了堂主任.
有一位金融界的弟兄跟我分享.
他說信仰是生命攸關的最重要的事情.
所以追求屬靈生命的成長.
就是你最有價值的投資.
我想你也有你的行業.
有你的使命和信仰的演繹.
信徒的生命就是這樣.
成為聖靈臨在的回眸.
讓我們都跟著上帝的魂處,火處的帶領.
祂放我們在不同的崗位上.
無論是在你的家庭,在你的堂會,在你的工作.
或者在你的社區裡.
願大家都繼續時刻經歷聖靈的同在.
活出這個常常喜樂,不住禱告,凡事謝恩的生命.
我們同心禱告.
因為你就是這樣讓你愛子耶穌基督的寶血.
首先去潔淨我們.
以致我們因著你而成為聖傑.
我們知道我們自己是做不到的.
但是你卻以為我們成就.
願主你都差遣我們.
讓我們在社會裡面.
不同的崗位.
在我們家庭,在我們的社區,在我們的堂會裡面.
讓我們都時刻經歷你的同在.
讓我們以喜樂,禱告,感恩的心.
常常都呈獻.
以致我們的生命能夠成為你話語的出口.
讓我們成為你喜悅的人.
我們多謝你.
我們祈禱奉耶穌基督得勝名求.
阿們.
\newpage



\section{詩篇 73:1-28}
\label{sec:d9OO9rhlq8k}
\textbf{2023-02-26 午堂《當信念被動搖》詩篇 73\_1-28 講員:鄧泰康傳道}
\newline
\newline
連結: \href{https://youtube.com/watch?v=d9OO9rhlq8k}{\texttt{https://youtube.com/watch?v=d9OO9rhlq8k}} ~~~~ 語音日期: 2023-03-01
\newline
\newline
\hyperref[sec:Lu3g7hRuGdE]{\small{< < < PREV SERMON < < <}}
~
\hyperref[sec:index_chronic]{\small{[返順時目]}}
~
\hyperref[sec:index_scriptual]{\small{[返順卷目]}}
~
\hyperref[sec:E8bPEPFjp7c]{\small{> > > NEXT SERMON > > >}}
\newline
\newline
詩篇 73:1-28
\newline
\begin{longtable}{cl}
\hline
\hline
章節 & 經文 (和合本修訂版)\\
\hline
73:1 & \begin{tabularx}{0.7\textwidth}{X} 神實在恩待以色列那些清心的人！ \end{tabularx} \\ \\ \relax
73:2 & \begin{tabularx}{0.7\textwidth}{X} 至於我，我的腳幾乎失閃，我的步伐險些走偏； \end{tabularx} \\ \\ \relax
73:3 & \begin{tabularx}{0.7\textwidth}{X} 因為我嫉妒狂傲的人，我看見惡人享平安。 \end{tabularx} \\ \\ \relax
73:4 & \begin{tabularx}{0.7\textwidth}{X} 他們的力氣強壯，他們死的時候也沒有疼痛。 \end{tabularx} \\ \\ \relax
73:5 & \begin{tabularx}{0.7\textwidth}{X} 他們不像別人受苦，也不像別人遭災。 \end{tabularx} \\ \\ \relax
73:6 & \begin{tabularx}{0.7\textwidth}{X} 所以，驕傲如鏈子戴在他們項上，殘暴像衣裳覆蓋在他們身上。 \end{tabularx} \\ \\ \relax
73:7 & \begin{tabularx}{0.7\textwidth}{X} 他們的眼睛因體胖而凸出，他們的內心放任不羈。 \end{tabularx} \\ \\ \relax
73:8 & \begin{tabularx}{0.7\textwidth}{X} 他們譏笑人，憑惡意說欺壓人的話。他們說話自高； \end{tabularx} \\ \\ \relax
73:9 & \begin{tabularx}{0.7\textwidth}{X} 他們的口褻瀆上天，他們的舌毀謗全地。 \end{tabularx} \\ \\ \relax
73:10 & \begin{tabularx}{0.7\textwidth}{X} 所以他的百姓歸到這裡，享受滿杯的水。 \end{tabularx} \\ \\ \relax
73:11 & \begin{tabularx}{0.7\textwidth}{X} 他們說：「神怎能曉得？至高者哪會知道呢？」 \end{tabularx} \\ \\ \relax
73:12 & \begin{tabularx}{0.7\textwidth}{X} 看哪，這就是惡人，他們常享安逸，財寶增多。 \end{tabularx} \\ \\ \relax
73:13 & \begin{tabularx}{0.7\textwidth}{X} 我實在徒然潔淨了我的心，徒然洗手表明我的無辜， \end{tabularx} \\ \\ \relax
73:14 & \begin{tabularx}{0.7\textwidth}{X} 因為我終日遭災難，每日早晨受懲治。 \end{tabularx} \\ \\ \relax
73:15 & \begin{tabularx}{0.7\textwidth}{X} 我若說「我要這樣講」，就是愧對這世代的眾兒女了。 \end{tabularx} \\ \\ \relax
73:16 & \begin{tabularx}{0.7\textwidth}{X} 我思索要明白這事，眼看實係為難， \end{tabularx} \\ \\ \relax
73:17 & \begin{tabularx}{0.7\textwidth}{X} 直到我進了神的聖所，思想他們的結局。 \end{tabularx} \\ \\ \relax
73:18 & \begin{tabularx}{0.7\textwidth}{X} 你實在把他們安放在滑地，使他們跌倒滅亡； \end{tabularx} \\ \\ \relax
73:19 & \begin{tabularx}{0.7\textwidth}{X} 他們轉眼之間成了何等荒涼！他們被驚恐滅盡了。 \end{tabularx} \\ \\ \relax
73:20 & \begin{tabularx}{0.7\textwidth}{X} 人睡醒了，怎樣看夢，主啊，你醒了也必照樣輕看他們的影像。 \end{tabularx} \\ \\ \relax
73:21 & \begin{tabularx}{0.7\textwidth}{X} 因此，我心裡苦惱，肺腑被刺。 \end{tabularx} \\ \\ \relax
73:22 & \begin{tabularx}{0.7\textwidth}{X} 我這樣愚昧無知，在你面前如同畜牲。 \end{tabularx} \\ \\ \relax
73:23 & \begin{tabularx}{0.7\textwidth}{X} 然而，我常與你同在；你攙扶我的右手。 \end{tabularx} \\ \\ \relax
73:24 & \begin{tabularx}{0.7\textwidth}{X} 你要以你的訓言引導我，以後你必接我到榮耀裡。 \end{tabularx} \\ \\ \relax
73:25 & \begin{tabularx}{0.7\textwidth}{X} 除你以外，在天上我有誰呢？除你以外，在地上我也沒有所愛慕的。 \end{tabularx} \\ \\ \relax
73:26 & \begin{tabularx}{0.7\textwidth}{X} 我的肉體和我的心腸衰殘；但神是我心裡的力量，又是我的福分，直到永遠。 \end{tabularx} \\ \\ \relax
73:27 & \begin{tabularx}{0.7\textwidth}{X} 看哪，遠離你的，必要死亡；凡離棄你行淫的，你都滅絕了。 \end{tabularx} \\ \\ \relax
73:28 & \begin{tabularx}{0.7\textwidth}{X} 但我親近神是於我有益；我以主耶和華為我的避難所，好叫我述說你一切的作為。 \end{tabularx} \\ \\
[1ex]
\hline
\hline
\end{longtable}
$^{1}$大家好 各位節目的早晨.
你的聲音夠不夠我的精神.
我們都準備好要在神的面前.
我們一起去敬拜.
我們心裡面每一次來聽神的說話的時候.
我們都要準備好 準備好些什麼呢.
準備好的就是你的心願意軟化下來.
承載上帝的說話給人.
今天和大家分享的這段詩篇.
其實有一個稱呼叫做憤悔詩.
憤悔詩是什麼呢.
就是一個悔改者向上帝的懺悔.
去懇求寬恕的禱告.
這些詩篇可以幫助人認識到自己的罪過.
去感到悔恨和懊悔.
並且尋求上帝的面容和寬恕.
這些詩篇亦都是提醒人.
只要通過悔改和認罪.
才能得到上帝的恩典和憐憫.
他們鼓勵人放下驕傲和自我.
謙卑地向上帝懺悔.
和尋求祂的指引力量.
以繼續前行.
總括來說.
聖經裡面的憤悔詩.
是對人悔改和再一次回到上帝的道路提醒和鼓勵.
是應該對自己的罪過和挑戰有效的工具.
這些是一些學者告訴我們.
憤悔詩有一個很重要的作用.
所以當今天和大家分享的時候.
你會發覺這段經文其實不深.
很淺白.
特別我請今天.
剛才主席讀得不容易.
這麼長.
但是用了一個叫做新普及版.
比較容易的版本.
去讓我們能夠.
其實很淺白.
但是裡面其實整個架構.

$^{41}$我形容的又很豐富.
為什麼呢.
你就會發覺其實.
這個作者.
他對這個世界.
他是有一種負擔或者責任感.
他不是作為一個局外人.
他是在這個世界當中.
他覺得他有他的責任.
以至於他看到這個世界的時候.
他會覺得在他很多.
他自己眼中看來很多的不公.
很多的惡人.
可以這麼說.
無論今天你站在哪個角度.
你站在哪個位置.
在你的眼中都一定會有一些惡人.
你都會覺得不開心.
你會為了一些事情.
你會忿忿不平.
所以今天這篇詩篇.
是非常好的.
去幫我們處理這些問題.
但是剛才我回到.
請大家今天聽這篇道.
記得.
這篇道最後的時候.
是想告訴我們.
無論今天你怎樣看都好.
請你回到上帝的面前.
原來我們有很多的.
是需要和上帝悔改的.
講個題外話.
大家喜歡看《約伯記》嗎.
《約伯記》的時候.
很簡單就是.
約伯遇到一件苦難.
他身邊有三個摯友.
你叫做摯友.
他們很關心約伯.

$^{81}$但是他們用了三個非常之神學.
非常之有深度的解釋.
去解釋約伯這次這件事.
為什麼會發生這件事.
而最後甚至出現了一個聰明人.
去解釋這件事.
但是最後的結果是什麼.
是上帝說.
你們全部都錯.
你們只是從你的角度.
去看我的作為.
我們今天去解讀.
這個世界上所發生的事.
很多時候請你記得.
我們只不過是一個人.
我們不是上帝.
我們看的東西非常片面.
所以當今天這個作者帶我們去看的時候.
我很希望帶回我們內心重溫的一次.
所以說很淺白.
裡面的事情大家都經歷過.
這段經文開車的時候.
其實這個作者是帶著一個.
很堅定的信念去成為基礎.
他一開始第一句是.
上帝實在恩待以色列人.
恩待內心純潔的人.
這個是他很清楚的信念.
但是他現在下面所面對的.
正好是挑戰他這個信念.
就是說你內心.
你信耶穌的日子.
你有你自己一套神學.
甚至這個信念是絕對沒錯.
是不是.
我想大家看完這句話.
都不會有任何異議的.
上帝是實在恩待以色列.
恩待我們這個民族.
也都恩待一些內心純潔的人.

$^{121}$但是他的經文很快.
第二節就是然而.
明明我是這樣想的.
但是當我去看見的時候.
就有一些現象.
我想大家留意.
這段經文開始那幾句是.
我險些失足.
我腳下一滑幾乎跌倒.
腳下一滑就是一不小心.
但是當我一不小心.
我就會在這些現象裡面.
我就會跌倒.
這些是什麼現象.
就是剛才那段經文裡面.
講得很清楚.
當你去見到有一些惡人.
你會見到明明已經是.
很完全表面化的了.
看到了.
這個人他所做的一切的時候.
是邪惡的.
他所做的事的時候.
是不可一世的.
他對人是指出惡言的.
一個很差.
非常之差.
明明是非常之不好的人.
為什麼他活得比好人還要好.
為什麼他沒有報應.
為什麼今天我看見這個人的時候.
我沒辦法相信.
剛才我第一段所講.
上帝你去愛你所愛的人.
甚至他講到身邊有人問他.
上帝知道什麼.
究竟這個至高者發生什麼事.
你看不看見這個惡人.
當我們在這個世界上.
我去看見的時候.

$^{161}$但他用了一個很有趣的字眼.
因我妒忌狂傲的人.
他所有的一切出現.
是什麼.
是因為我妒忌.
妒忌的意思是.
我看不過眼.
我心懷不平.
中國人有一句很有智慧的說話.
殺人放火金腰帶.
修橋補路無屍骸.
但是我想告訴大家.
這個正好就是我們的問題.
人總是很喜歡做裁判者.
最近這幾天的時候.
整個新聞的篇幅.
一定是在講大埔的兇殺案.
如果大家有機會去到酒樓.
去到餐廳.
你一定多多少少會聽到.
隔離檯在談這個話題.
甚至你和其他人都在一起.
在談這個話題.
我不知道大家聽到什麼.
有些人對於殺他的疑兇.
是滿譜不平.
我也有看新聞.
按照暫時新聞所出現的時候.
這個女死者對於你們家的照顧都算是這樣.
你們還想怎樣.
之後我聽到很多時間.
貪心.
他們有很多的問題.
很多很多.
很多批評.
活該你死刑.
最好判死刑.
很多這些.
但同樣的我也聽到一些評論.
是對這個女死者的批評.

$^{201}$你好像是活在她旁邊.
你知道她的心裡想什麼.
你知道她一世人過什麼生活.
之後當然這個女孩子.
她不方是好人.
這麼年輕的時候.
嫁兩次.
有時候嫁了個有錢老公.
就這樣那樣的時間.
弟兄姊妹當我們聽到這些的時候.
你可以說.
我不知道.
人真是很喜歡的.
茶餘飯後.
如果不說這些還有什麼說.
所以這些是我們可以去維持小圈子的材料.
但大家就說到自己就好像專家一樣.
說到我們是裁判官一樣.
說到今天我們是很有論點一樣.
好像知道我已經掌握大局.
甚至你說到自己就好像上帝一樣.
當我們今天去看見很多這些不公的時候.
其實我們是不是真的應該要做一個裁判者呢.
我經常都會覺得.
我很怕我們的人生將所有東西中擺兩邊.
要麼這樣.
要麼就這樣.
要麼我們就很熱心.
茶餘飯後一起去說.
一起去談.
一起去想.
扮專家.
要麼就說到我們對這個世界不聞不問.
什麼都不關我事.
不是的 弟兄姊妹.
我們是活在這個世界.
上帝讓我們活在這裡的時候.
總是有他的目的.
上帝都說過我們在地上成為他的使者.
當今天我們面對這些事情的時候.

$^{241}$第一 我們很見到那個現象.
我們很碰到那個事情.
這個是事實.
我們會不會不開心呢.
我們會不開心.
但是我們應該怎樣去處理這個不開心.
或者我們應該在這個不好的現象裡面.
我們應該做些什麼.
這個才是我們需要去說.
我經常都說.
吃花生 做花生油 做鍵盤盒.
很容易的.
因為你不需要付任何代價.
你講了一些話的時候.
你不需要負任何責任.
但是你又知不知道.
當我們一不小心的時候.
我們會傷害了很多無辜的人.
在這個過程裡面.
你會發覺作者一樣.
他看見這些不公義.
他看見這些的時候.
他出現一個情況.
他開始懷疑自己的善.
保持內心單純難道全然圖演.
我持守自身清白難道毫無緣有.
他又提出了一個問題.
他說當我看見世間上這麼多惡事的時候.
為什麼我還要保持我自己在上帝面前.
這個單純 這個善良.
接著之後其實他想做什麼.
他想找出行善的原因和價值.
他想找個答案出來.
我們人都是一樣.
當我們遇見事情的時候.
我們總是想找個答案出來.
所以剛才說約伯那三個朋友.
就去給他一個答案.
其中一個就是什麼.
其實很無辜的.

$^{281}$因為我們是作者我們是讀者.
我們很知道其實約伯是一生裡面都非常之愛主.
他們都教那些子女教得非常之好.
對於他們來說的時候.
很無辜的無端端就倒下了.
我的子女全部死了.
但是為什麼死剩下老婆.
老婆就不斷罵他.
你快點離開上帝.
你不要再親近上帝.
就不斷的去blame他.
再加上這個朋友突然還跟他說什麼.
一定是你犯罪.
你不要驕傲.
你有些事情一定做錯了.
上帝是公義的.
上帝如果不是的話他不會這樣的.
我們人總是很喜歡在裡面找答案.
我們很想找個答案出來.
這個作者他很想找個答案出來.
因為他到了一個地步.
已經違背了他之前所自己堅信的神學.
就是第一句.
所以他在這裡說.
我若真是這樣對人說話.
就是背叛了你的子民.
因此我竭力想要明白.
為什麼惡人得勢.
但這是多麼困難.
因為他作為一個傳道者.
他說當我面對這些事情的時候.
我都不知道怎樣去面對人.
怎樣跟他說.
上帝啊原來你是慈愛.
你是公義.
我不知道怎樣做.
大弟兄姊妹.
聖經裡面其實你知不知道.
告訴你.
如果你總是在事情上要找答案的話.

$^{321}$最終你只會帶來痛苦.
傳道書其實他直接就跟你說.
這個世界上根本就充滿著矛盾.
這個世界上.
這個沒有上帝的世界裡面.
根本就是虛空.
是荒謬.
所以傳道書給你的出路.
不是答案.
他沒有跟你解答每一件事.
他給你的出路就是.
你是被造者.
請你好好的按照上帝所給你的份.
好好地去過你的生活.
好好地去成為其他人的祝福.
好好地去做上帝要你做的事.
到最後上帝必審問.
所以在這裡你會發覺.
原來行善是甚麼一回事.
原來行善完完全全是神對贖他的人的命令.
就是今天我告訴你.
你在地上我是做你的.
你要行那個道路.
去祝福其他更多的人.
行善的目的不是選擇性的.
不是環境性的.
不是因為你想不想.
或者不是因為環境的緣故.
你決定做不做.
甚至乎我回到這段經文所說.
不是因為你行善帶來你任何好處.
所以你去行善.
弟姐妹明白嗎.
所以無論這個世界是怎樣也好.
弟姐妹你要堅持繼續行善在上帝的善當中.
因為這個是上帝的要求和命令.
你是屬於他.
我可以告訴弟姐妹.
根本這個世界在我們人生當中.
很多事情你都沒有辦法找到答案.

$^{361}$最近我一直和弟姐妹.
我沒有辦法和大家分享.
我和我的團體分享.
最近在鄧泉道大家都為我祈禱.
好好的時間.
跟著的時間一直都很正常.
我又很喜歡運動.
你看我身形都不像坐著不動的那些.
我都會小心飲食.
小心健康.
但最近的時間我的腳突然間麻木了.
我的右腳突然間對冷熱沒有感覺.
我的左腳開始影響我走路.
肌肉突然間很緊.
我都莫名其妙.
我都很想上帝給我答案.
上帝你玩什麼.
我很努力事奉你.
但上帝不會在這裡給你答案.
他不會給你答案.
在我心裡我只知道一件事.
他有沒有拿走你的命.
他沒有.
那應該怎樣.
繼續每天你要做的事你就做.
那我不是望天打卦.
我即時看醫生.
我都會找醫生治療.
我都會尋求一些積極的幫助.
但你不要停留在這個困難當中.
你不要停留在疾病裡面.
你還有很多事要做.
所以上星期我去了狼茄街道村.
大家有沒有去過.
大家有沒有去過那條山路.
我都很擔心.
我的腳都不知道能不能走路.
當我祈禱的時候.
日日有弟兄.
我陪你去.

$^{401}$陪我一起.
兩個人同行就真的很好.
後來還借了一支爬山杖給我.
我終於領略到爬山杖的用處.
原來是這麼重要.
很平平安安去到之後.
我和他們那班的戒毒朋友說完.
之後很平平安安回到.
我還跟弟兄開玩笑.
其實昨天我還在淘寶.
看看網上有沒有一些很輕巧的電動輪椅.
萬一行不到的時候.
還可以動得到.
弟兄姊妹.
究竟我們今天作為一個基督徒.
我們有我們的信仰.
有上帝.
我們是不是要被我們身邊的事情.
打倒我們.
我們今天應該要專注關注什麼.
不是找答案.
我們是要繼續在神面前.
履行上帝所給我們的即時.
所以《論經文》裡面的轉捩點.
是我很認同.
因為我也是同樣.
走回同一條路.
他就說上帝啊.
我進入了你的聖所.
終於明白惡人的結局.
他的轉捩點就是他來到神的殿裡面.
神的聖所當中.
來到神的聖所當中做什麼.
弟兄姊妹.
敬拜.
屈膝.
當我來到上帝的面前.
我集中焦點看回你的時候.
我要敬拜你的時候.
我屈膝的時候.

$^{441}$我就知道我是誰呢.
我只不過是一個極之微小的人.
今天我就是因為見到你的緣故.
我相信你在聖經裡面.
你所給我的應許.
所以我見到這些事情的時候.
我就會回想起.
上帝你所跟我們說過.
聖經裡面寶貴的話.
你說你是公義的.
所以他即時你發覺.
他整個轉變了.
他說我即時就看見惡人的結局.
沒錯.
你已經將他們放在滑漏的道路上.
讓他們順著迫一路向下滅亡.
傾一間他們就滅亡了.
徹底被恐懼一掃而光.
主啊 當你起來的時候.
你就必嘲笑他們愚蠢的想法.
就像人在清早笑話自己的夢一樣.
即時就轉變了.
他就會見到.
是喲 這些人原來在你面前的時候.
這些人以為自己不可一世嗎.
原來不是 是一場夢來的.
虛幻的夢來的.
今天那些惡人.
他們在行惡的時候.
請他小心.
你不要以為自己是不可一世.
因為上帝什麼都知道.
你一路走 你以為自己是順利嗎.
你明不明白剛才的圖像.
一個順滑的波動.
很暢順啊 明不明白.
但他一路滾去哪裡.
滅亡.
這是鄧全道的深深經歷.
我每次分享我以前的見證的時候.

$^{481}$我都說在我讀神學之前.
一個很重要的見證的轉捩點是什麼.
當我選擇離開上帝的時候.
上帝給了我五年的時間.
是什麼 是順風順水.
我一生人裡面所有最重要.
在人世間最有價值的事.
我都發生了.
在那五年裡面升職 結婚 買房 買車 生孩子.
全部在那些時間都已經發生了.
人生勝利組.
但不是啊.
我覺得誇張一點.
原來上帝給你那麼順利.
是推你去死的.
這是我自己很深心裡面很明白的.
我很明白當時我自己的驕傲.
為什麼當時我分析的問題.
因為我自己當時是什麼.
嘩 我這麼聰明.
我回到公司的時候.
我這麼聰明.
我見到這麼多這麼漂亮的女同事的時候.
我老婆算什麼.
心的變異是這樣來的.
你會發覺為什麼後來我整個工作裡面一仆一祿.
因為我自以為自己很靈.
沒理由的.
那個傻瓜都可以做經理.
我不做得.
上帝就是這樣讓我們.
推到人生的懸崖當中.
我是哭著向上帝回轉回來.
有人問我為什麼做傳道人.
如果沒有上帝的話.
今天我已經不在了.
所以弟兄姊妹.
原來當你一回到上帝面前的時候.
你相信上帝公義的時候.
一切都不會再成為你心裡面的困惑.

$^{521}$你就能夠有出路.
你就能夠繼續有力量的向前走.
而這個轉化不單止.
你不要以為停留在這裡.
看到外人心涼.
不是的.
第二個才更加重要.
此時我才知道.
我曾內心苦澀.
我曾肺腑拘裂.
我曾如此愚昧無知.
那時我在你眼裡就像毫無理智的畜生.
詩人為什麼要罵自己是畜生.
就是會覺得自己為什麼這麼愚蠢.
為什麼我會這麼不相信上帝.
為什麼我會這麼無知.
不相信上帝你會在當中.
為什麼我不相信你的公義.
為什麼我要在之前受這些內心的苦.
我真是蠢.
但接著他說什麼.
但我仍然屬於你.
你緊握著我的右手.
你以你的訓言引導我.
要將我領到榮耀裡面去.
他說在這裡發現到一樣東西.
原來自己對上帝這麼無知.
這麼愚昧.
但是上帝竟然去體諒自己.
和對自己的不離不棄.
當我見到其他人是惡人.
原來當我一不小心.
我在神面前同樣可以變成惡人.
我不知道弟兄姊妹.
你們對於人能夠決志信耶穌.
你有什麼感受.
在聖經裡面說.
當有人願意相信上帝得著恩典的時候.
人不去讚美天上的天使都會讚美.
但我卻在過去的日子裡.

$^{561}$我覺得這種審判很可怕.
曾經有人對於.
當然我經常說.
我不知道過去幾十年裡.
是不是整個基督教裡.
很習慣了一種批判的思維.
去批判大型報道會的不足.
去批判有些人信耶穌之後的不是.
弟兄姊妹.
說有些人的觀察是對的.
很簡單 大型報道會.
去到之後.
然後決志說舉了手.
然後都不見了人.
都不知道去了哪裡.
所以他就開始去批判.
大型報道會的價值的用處.
弟兄姊妹.
我要問一句.
你是上帝嗎?.
大家有沒有聽過很多.
我太太也是從業其中一個.
當年去聽國培理報道會.
是不是.
的時間.
然後就舉手決志.
決志之後有沒有回教會?.
沒有.
但之後過了很多年之後.
他相信耶穌.
他上台講見證的時候.
他才講.
其實他很相信上帝是將種子放在心裡.
聖經教我們.
災眾有時.
修葛有時.
不是我們去決定.
是上帝只是想你用盡一切方法.
去將福音給對方.
而至於對方的內心是怎樣.

$^{601}$你不會知道的.
究竟那粒種子何時萌芽.
你是不會知道的.
如果在我們信仰裡面.
我們有那麼多估算的話.
弟兄姊妹.
我們傳不到福音.
你亦都不會去傳福音.
因為你會覺得某些事情成效太低了.
但弟兄姊妹不是.
聖經的教導是.
讓你盡心盡性盡義盡力去愛神和愛人.
所以你會發覺.
在這時候.
這個作者到最後.
我稱之為他的內心改換一生.
除你以外在天上我還有誰呢.
除你以外在地上我別無眷戀.
我的身體會衰老.
我的精神會衰弱.
但上帝仍是我心裡的力量.
祂永遠屬於我.
多漂亮多美麗的宣告.
從一個困苦.
從一個很大的失落.
不開心的裡面的時候.
因為在上帝裡面的轉化.
然後內心裡面充滿著這種對上帝的盼望.
對人生的盼望.
對人生的肯定.
所以到最後他的宣告是.
那離棄他的必定滅亡.
因你消滅那些離棄你的人.
至於我親近上帝多麼美好.
我以至高的上主為庇護所.
我要向眾人述說你奇妙的作為.
去到最後的宣告就是.
我只知道到最後.
我要和所有的人.
去繼續述說上帝的大浪.

$^{641}$弟兄姊妹當我們看完這段詩篇的時候.
我們就會發覺.
整個環境有沒有變.
沒有變過的.
但是問題有沒有答案.
沒有答案的.
但是你會發覺.
完完全全是一個心靈的轉化.
弟兄姊妹.
基督信仰.
基督信仰最強大的力量.
一直都是我們的恩信 情宜.
是因為你信.
所以你相信上帝的應許.
你有今天這個得救的盼望.
因為你相信上帝的應許.
所以無論你走到什麼位置的時候.
你仍然不會灰心.
你仍然會奮力向前.
直至到你見主的面.
其實最重要的就是.
今天上帝要改喚我們的心.
上帝從來沒有要你去.
在這個世界上做任何的裁判者.
但是我也說.
如果有一天上帝要你成為裁判者.
請你就按照上帝給你的信仰的勇氣.
好好去做.
如果你做裁判官.
請你不要因為任何的緣故.
按照你所見到的事實去做好的裁判.
如果你有機會做陪審團.
請你按照你所聽到的.
所了解的內容.
好好去做你的判斷.
但是我們今天有更多的時間.
沒錯,我們每一個人都有不同的看法.
我們每一個人都有不同的思想.
但是請你在裡面的時候.
好好回到神的面前.

$^{681}$我們一起去求民.
就好像我剛才說的.
有人很火熱的大埔兇殺案.
當你看到這件事.
其實會不會讓我們內心有更多的反省.
我們在過程裡面去想想.
無論對人.
對我們身邊所擁有的東西.
或者其他東西的時候.
會更多令到我們人生會多一些的更新.
讓我們更加要去把握時間.
讓我們更加知道原來有些事.
當我們處理得不好的時候.
我們心裡面的惡魔會一路成長.
所以弟兄姊妹.
求上帝幫助我們,好嗎?.
我們同心企圖.
秋天會求你幫助我們.
最後當我們在神的面前.
我們的信念受到挑戰的時候.
我求主你自己親自去保守.
讓我們來到你的面前.
讓我們能夠來到你的當中.
我們謙卑下來.
我們再一次去肯定上帝你自己的主權.
肯定上帝你自己的英雄.
以至於我們能夠內心得著釋放.
繼續有力量去向前行.
繼續走上帝你要我們走的路.
去走祝福人的路.
去走傳揚福音的路.
最後我求主你親自幫助.
求主你引領,帶領我們.
也求主你讓我們成為更多人的光明之子.
和平之子.
以至於能夠成為更多人的安慰.
將一切所有的我們交在你面前.
其中奉靠主耶穌,得勝名自祈禱.
阿們.
\newpage



\section{}
\label{sec:E8bPEPFjp7c}
\textbf{2023-03-02 周四晚查經班《被擄歸回歷史研讀(8)》 講員:曾裔貴牧師}
\newline
\newline
連結: \href{https://youtube.com/watch?v=E8bPEPFjp7c}{\texttt{https://youtube.com/watch?v=E8bPEPFjp7c}} ~~~~ 語音日期: 2023-03-03
\newline
\newline
\hyperref[sec:d9OO9rhlq8k]{\small{< < < PREV SERMON < < <}}
~
\hyperref[sec:index_chronic]{\small{[返順時目]}}
~
\hyperref[sec:index_scriptual]{\small{[返順卷目]}}
~
\hyperref[sec:M8De3gUONVA]{\small{> > > NEXT SERMON > > >}}
\newline
\newline
$^{1}$.
多謝主慈恩賜給我們今天這本信.
歡迎我們來到聖經廣告.
顯然很舒服,我們可以來到一場聖經.
我們很盼望,無論是現場的聽眾.
或之後成員錄音的同學.
有約11月份,甚至會有教會.
可以透過演奏去分享,去得到祂的祈禱.
求主的聖在聚會當中.
在廣告中,去到祂在聖經的心裡.
我們感同身受,彈劾在你的面前.
獻上好報,奉耶穌基督的名義.
我們上星期談到的課題.
都是相當重的一個問題.
我們看看這裡.
特別是第十章.
以色列所做的激進手段.
就是要強迫歸回的.
其實已經是第二代的形態.
可能有第三代.
歸回的猶太人.
他們基於很多的限制因素.
其中的限制就是不夠老婆.
不夠猶太回歸的.
或者出世,生的.
很現實地,嘉臘地.
原居民很多的女子.
很自然地會由婚嫁.
娶外邦女子.
也會嫁給嘉臘的男子.
是每個家庭.
雖然聖經沒有這些資料讓我們知道.
以色列和當時的猶太領袖.
所擔心的.
就好像以前的所羅門.
娶了那些嘉臘女子之後.
就會引誘所羅門.
在晚年的時候.
去拜他們的假神.
就偏離了神的真理.

$^{41}$再單純地敬拜同一真神.
但是聖經沒有什麼案例.
讓我們看到在以色列的時代.
因為現在是法賢.
這個法賢也不是一個早期.
而是一個已經有子女的.
已經是很普遍的.
而且數量也不少.
單是有記錄那些年輕的.
全部是領袖.
以及利美茲派.
作為重要的一個指標.
也就是說很現實地.
是真的需要.
因為不夠猶太的女子.
所以以色列回到這裡.
看到的現象.
就是這個情況.
他就採取一個很激進的手段.
首先就是普查.
到底有哪一個家族.
哪一個事奉的支派.
娶了外方女子.
就要求他們在十月初一.
要來到一同的在座.
去查辦.
然後去到正月初一.
就查清楚.
娶了外方女子的人數.
這是十章的十六,十七節.
首先就是以色列他自己一個人.
痛哭,祈禱,認罪.
後來就感動了一些祭司.
一些領袖就說.
以色列只是自己感傷口是沒有用的.
你應該聚集民眾.
願意的一起來去履行.
重點就是他們要回到一個很重要的身份.
就是我們要歸回成為.
神在我們這個民族.

$^{81}$一個特別選擇的民族.
那個聖潔的身份的背景.
所以這個就是以色列在第七章.
一直到第十章.
整本書其實這裡才是真正這本書的.
寫作的重點.
以色列記的第七到第十章.
是第二個部分.
很清楚.
首先就是說有一個人.
這個人就是拉比.
是祭司.
其中特別是他.
他是祭司.
但是他更加專門的就是.
那個的民事.
教導聖經的民事.
所以他的回歸就帶來一個.
你可以說是一個小型的復興.
至少有人肯說出自己的情況.
有人願意承諾.
去離棄自己的太太.
自己的子女.
因為都聽從以色列所說.
我們所做的.
是與舊約.
波西部經.
申命記所吩咐應許的.
是有一個距離.
所以他們就承諾.
不過我們上個星期就點出過.
就是整張聖經這麼激進.
似乎聖經就沒有在這裡給我們一個結果.
到底這樣的安排.
是不是真的全民都要去效法.
會跟隨這個命令.
一個行政的指令.
強迫所有娶了外方女子的.
都要去休妻.
趕走自己的子女.

$^{121}$你想想後續的問題就很大件事.
迦南女子你不娶.
猶太的女子可以娶的都娶了.
那你還哪裡有女子來嫁.
這些有婚的猶太男丁呢.
是不是在外方.
波斯帝國運送一些猶太女子回來呢.
我們又看不到這裡有一個後續的情況.
我們看到引發這個問題.
是以色列.
當然是她一廂情願.
是她一個心很大的感覺.
就是希望我們的同胞.
既然我們居住在有聖殿的聖城.
我們當然沒有城牆.
我們是神的選民.
我們是有這個心志.
由很遠很遠的波斯帝國.
回歸回來.
已經是第二代第三代了.
開始第三代了.
我們應該怎樣去面對神的聖經教導呢.
所以以色列引發這個.
你可以說在她的角度是一個復興.
但是在民眾百姓的角度.
這是一個革命.
是要將你自己的情況.
如實地告訴以色列.
要記錄在案.
然後什麼時候.
她有一個定義.
承諾了一定會凌奪自己的妻子.
到最後的44節.
這些人都娶了外方女子為妻.
計算的數字.
這裡只是領袖.
都超過40個.
這樣目測的數字.
然後他們都已經有生了兒女.
其實是很大的事情.

$^{161}$是帶來.
差不多是民族之間.
你想想很多的外方女子.
哭著回家.
回外家.
外家的人怎樣看待猶太人呢.
怎樣看待以色列呢.
怎樣去講到民族的共融的含義呢.
為什麼會這麼介意這個身份呢.
因為生下來.
她是家立女子.
就一定不可以跟她有任何接觸.
跟她有任何交往.
因為他們是愛宗的.
是敬拜耶和華的.
不是每個家庭.
都是這樣自己藏起來.
去拜九像的.
可能他們入國之後.
他們是敬拜神.
跟隨自己的丈夫.
敬拜神.
怎樣去計算這些情況呢.
但是聖經就完全.
一點交代都沒有給我們.
只是說了.
記錄了.
有這樣的人.
這些靈就.
是利其之派.
過去他們都是有參與.
試過的.
現在因為以色列的出現.
就要將他們全部.
給了去.
如果不做.
就可能不讓他們試過.
但是聖經就沒有交代那個結果.
更加重要的是.
這裡的分水嶺.

$^{201}$很多神學家都說.
這裡會不會有過分的異端.
過分的.
一種好像.
宗教的發明.
好像中世紀的.
西班牙.
那個叫宗教財經所.
就是基督教對基督教.
或者基督教.
如果你自己收藏聖經.
在家裡自己讀.
或者你跟教會的做法.
去做任何事.
你就是一個異端.
就要裁判你.
如果發現是事實.
當時他們有那個權.
要燒死一些靈體.
譬如說牧師.
丁道主.
他們都是這樣信的.
所以這裡.
是不是一個異端.
異端裁判所的.
一個.
這樣的一個舉止.
往後.
很多很多的這些問題.
如果你不是出於情願.
心甘樂意的.
會不會被強迫.
去做一些違反了你個人意願的事情.
個人的意願.
與舊日的聖經.
到底那個比例應該要多少呢?.
聖經全部沒有交過.
所以有些神教家就認為.
當然如果又是認為.
是聖經盲舉了.

$^{241}$以色列這樣做是對的.
是沒有錯的.
但如果有一些神教家就說.
其實會不會是.
聖經的沉默呢?.
就是以色列過往激進的一個手段.
這個手段.
好像一個行政的手段.
去逼明了.
真是家殺人亡.
那些避信.
那些外邦的女子.
就是這樣去頒外加.
趕走的.
甚至有他們那個.
宗教上的指令.
登記了.
逐一逐一登記了.
可能真是連帶自暴.
這個家庭.
有外邦女子的.
好像變成了.
在那個政治上.
那個行政上.
來去干預了.
人生的情況.
是不是也是一個指慰呢?.
去到新約就完全沒有這些的教訓.
也沒有這些手段.
但當然代表新約.
沒有這些的要求.
不過新約就處在一個.
真是嚴重大偶像的社會.
即是多神的主義.
所以反而就是.
在那個寧靜的層面.
要求信徒.
你處身在一個這樣的社會.
你應該.
你的那個道德.

$^{281}$要比信的人再高一點.
這樣你是一個性別的種類.
你是有君尊的祭司.
這些身份的.
所以這些就.
因為你沒有了一個具體的行政指令措施.
就變成了.
很容易就是每個人按自己的選擇去選擇.
這裡就給大家.
有一點點這個背景.
因為我們可以很快去看看.
尼希米基的經文.
特別我們不看開始.
而是看最後那一段.
就會比較長一點.
尤其是十三章那裡.
兩個領袖.
都是有這樣的體度.
就是很激進地.
要求那些.
合乎摩西古經的那個教導的.
就要做出一些很果斷的決議.
還有那個行政的措施是很重的.
就是說行政的手段.
凌駕了個人的選擇.
這個就是.
譬如我們用十三章.
其實很小事.
不過.
尼希米這位的省長.
他是波斯國的省長.
就是幫外邦人打工.
波斯國的亞達薩西.
就是他的老闆.
他在王面前是一個酒精.
十三章.
我們很快看看這裡.
因為和剛才的以斯拉的十章有關.
而且.
當然有兩個理論.

$^{321}$不過比較大部分的就是.
以斯拉和尼希米是同時的.
曾經有十至二十年.
或者十多年.
因為不太清楚.
就是同時的.
所以就不是另外一個時期.
是同時的.
大家一起去做過.
第九章那裡就有了.
現在我們看第三章.
當日.
人而摩西的律法書.
給百姓聽.
御見書上寫著說.
阿滿人和摩納人.
這個是第二十章的新歷史.
永不可入神的路.
這是很舊時.
很過氣的律法律.
但是如果時代轉變.
我們還要實行特首承規.
在經文的詞典裡.
似乎尼希米就是這樣理解的.
第二節.
因為他們會有拿食物和水來迎接以色列人.
而且還請了巴拿.
就是那個假先知.
其實他不是假的.
是真先知.
不過他的信息.
是要米飯先知.
誰給米飯我就幫他說話.
不是講神要他說什麼.
這些就是賊閒底先知.
因為他們會有拿食物和水來迎接以色列人.
而且僱了巴拿就在他們.
但我們的神.
使那就在.
其實就是巴拿的就在.

$^{361}$見羅足公.
以色列民聽見這律法.
就與一切閒雜的人絕交.
然後.
先是同派.
這段時間在哪裡傳?.
因為有一件事.
就是在安息日的星期六.
當然這個月介乎是星期六的黃昏六月.
太陽落山的標準.
到第二天的黃昏.
就是為了安息人.
所以就不是星期六的上午傳.
而是由星期五的黃昏.
到星期六的黃昏.
就叫做安息日.
在這個時候.
他們就在安息日買賣.
所以就激到李米就突然跳了.
先是無派第四節.
先是無派管理我們神殿中庫房.
祭司以利馬臣和多比瓦.
多比瓦就是以前的三張房子的領袖.
一個迦南的領袖.
與多比瓦結親.
為他們預備一間大型的屋.
就是從前收存素祭與香器人和照病人.
供給李米榮歌唱首門的福新酒.
或者神明之恩.
並歸祭司舉祭的屋.
即是那個地方是成份庫人的.
是那些百姓獻祭.
獻祭完之後有不同種類的祭.
就會收窄了.
可能是一個難度.
因為要雪條肉.
或者有其他的工具.
讓那些獻祭的祭身.
就不會被他破壞.
是分派給祭司和李米榮.

$^{401}$那些叫家族.
即是說這些是祭司和李米榮的下屬.
第六節.
「那時我不在耶路撒冷.
因為巴比雷和亞達薩西三十二年」.
這邊有些奇怪的.
突然間說巴比雷和.
卻是波斯王.
所以有些抄寫上的說法.
「我回到王那裡.
過了多日.
我向王報答.
我來到耶路撒冷.
就知道以利亞臣為多比亞.
占用了祭司的一個成份庫.
給了他的出家多比亞的住.
就帶來」.
李兄是怎樣呢?.
「臣以裡面預備屋子的那件惡事.
我就拾起來.
把多比亞的一切家具.
從屋裡拋出去.
是一個神聖的手段」.
很神奇.
當然我們還未看李米.
第一章我們稍後會看.
李米是波斯的省長.
所以是直接王委派他.
去到耶路撒冷.
在那裡拜.
所以他是和東的省長.
所以他真的有權有實權.
現在他用一個很大的神聖的手段.
把多比亞佔據的這間屋子.
裡面的家具全部被扔出去.
分庫人見證這間屋子.
當然不是一定不見證.
而是這間屋子是屬神的地方.
這個地方本來是分段為聖.
那些食物是百姓獻祭.

$^{441}$是給祭司庇護人.
作為他們的薪金.
是他們的享用和家庭的需求.
但是現在多比亞佔據了那個地方.
知道會不會拿那些食物.
拿東西拿出去賣等等.
那些是獻祭的.
不是普通的祭用品.
在街上在市場上賣的.
所以分段為聖所.
只有祭司才可以吃.
所以在這個時候.
他們又要吩咐下人.
見證這間屋子.
把神殿的器皿和掃祭.
禹香又搬進去.
或者說其實很大件事.
祭司以利亞神.
把神的.
總之我講笑而已.
沒有人知道.
那張禮堂後面這兩間房.
好像是人走的.
租了給人做tong 房.
接著收租.
類似這樣.
現在以利亞神就是這樣.
把神殿裡面的地方.
根本是祭司才可以.
可能那間房是他的.
他自己出入的.
他自己有鎖匙.
他原來給了多比亞去住.
沒有人知道.
你就知道.
其實當時有這個情況.
第十節.
我見利未人所當得的房.
沒有人供給他戶.
甚至供給利未人與歌唱的.

$^{481}$開國分僚成員的田地.
本來田地是分給利未人的.
可以耕種.
有些婆婆他們可以住.
例如Cotton.
是神批認.
以色列人要每一個支派.
都要分僚成一些地方.
有些田地是給利未人.
是不可以買回來的.
給他們去住.
因為他們是無份無業的.
所以他們要留這些.
就好像宿舍.
但是利未人公積無糧水.
所以他們就要分配自己的田地.
就是說他們變成了自由傳播.
不過是地盤的.
要成為自由傳播.
因為祭司.
聖殿沒有糧水給祭司.
原因是有些人在補償.
第十一節.
「我只甲官揚說.
為何你棄神的天」.
因為祭司你們人都是需要承擔的.
你沒有糧水給我.
我就自然要耕種.
我自己要去找.
讓我的兒女和孫子.
都需要有兒女的承擔.
所以就沒有人再.
在尼希米回到波斯帝國之後.
那個盛宴就出現了這個混亂.
首先就是.
以利亞神居然間自己建了那個礦場.
給了多比亞.
是他的親家.
就將你們祭司應該得到的.
每個月的供糧.

$^{521}$就停止發放.
之後令到祭司你們人.
都不再專心地可以在聖殿裡示眾.
你想一下這個情況.
就會出現什麼樣的難處.
所以就不是真正的自由戰鬥.
平常的不可以再傳統了.
這個是當時的情況.
第十一節.
「我只甲官揚.
貴學你棄神的天.
我們招聚你們人.
使他們照舊功績.
由大眾人就打.
谷身走和.
由其神之恩送人護臺」.
就是神的洞穴.
再次送回到那個護臺.
「我派祭司.
使你米馬.
弘時薩諾和利未人.
給大馬作副官.
管理護臺.
副官是哈蘭.
哈蘭是薩克的人.
薩克是瑪他尼的人.
這些人都是忠信的」.
很重要的.
需要有那個智慧.
因為現在已經有人發生.
在聖殿裡面的爭拗.
所以要令到那件事.
能夠恢復那個秩序.
首先就是要讓那個副官.
有那個秩序.
和這個奉獻的制度.
要有那個秩序.
就是鼓勵猶太人.
你們應該盡你們的本分.
去奉獻給神的殿.

$^{561}$之後讓神殿裡面的事故的人.
有那個所謂無後顧之憂的.
那個願望.
願望的需要.
然後他們就會在聖殿裡面.
去建制 普告.
教導律法書.
審判等等.
就會令到那個秩序.
會有那個影響.
這個時候就要找幾個人.
去發揮以前信心的.
自發性的.
現在可能就需要有這個行程.
讓那件事有那個規矩.
他們的職責就是.
將所供給的職責.
給他們的文件.
即是他們的文件.
我的神啊.
求你恩者使記念我.
不要遺忘我為神的願望.
而做忠於文件所行的事.
你不是看到一個特色的這本書.
就是你做這件事.
你不是也好像有一個個人的團隊.
在這裡教授.
接著又來另一件事.
這件事大家可以看一看.
對手不知.
那些日子我在猶太.
見有人在安息人家走.
當然是踩走駕.
搬運著和坐在路上.
有把酒 葡萄 無花果和各樣的丹.
在安息人家監禁著我撒冷.
我就在他們派食物的那一天.
警戒他們.
似乎未必是猶太人.
是外方人士.

$^{601}$又有推羅人駐在我撒冷.
他們把魚和各樣的貨物運進來.
在安息人家大吸猶太人.
我就立刻猶太人跪下.
你們怎麼行者惡事.
煩了安息人呢?.
從前你們的列祖豈不是也約.
以致我們的神使一切的災禍.
臨到我們和者成安.
現在你們說煩安息人.
使憤怒越發臨到而逝.
一件可能很小的事.
今天我們很少會覺得有人.
我們今天已經沒有分所謂安息人.
因為有教會是安息人家.
現在已經比較適應了.
即是三獄中學就是安息人家.
星期日才要吃飯.
星期日就不吃飯.
他們認為安息人才是最快事.
即是跟人家人家推不了.
即是現在好像有些人離險離貧.
好像有安過.
但其實到世後.
主耶穌說是哪個日子.
在哪個地方.
那個精髓是主耶穌才是安息人的主.
我們要敬拜這位主的神.
所以星期日其實是約定要吃飯.
等於是星期日一定是大錯特錯.
也不是.
所以現在有些教會因為需要.
星期天不吃飯.
有些教會星期三的早上.
都有他們的吃飯.
因為有些父女.
星期三可能遲到.
有家庭的人.
他們可以來吃飯.
那個訊息是跟星期日一樣.

$^{641}$所以教會有那個因時制宜的安排.
所以去到新約那個智慧.
所得到的那個見諒.
其實任何一天都可以吃飯.
不過現在從宗教改革開始.
星期日已經是大家的習慣.
其實過去通常.
應該說去到主耶穌的時代.
離世之後.
七日的第一日星期日.
他們是聚集的.
所以已經是差不多二千年的文化習慣.
所以一般人就算信的都知道.
星期日是去禮拜敬拜神.
所以就是用這個做一個大家普遍的想法.
所以任何教會的那個設定.
一定是以禮拜做那個上午.
做那個聚餐的基本的單位.
有些教會.
例如歐洲的皇教會.
他租不到 借不到.
有時是禮拜的下午才吃飯.
有時是因時制宜.
當然大部分人都是星期日不放假.
星期日的上午.
如果先去崇拜.
是方便他們其他的時間.
家庭心悅.
這些全部不是錯的.
所以要視乎那個地方.
那個教會的需要.
所以新日就已經沒有這個問題.
但是回到紐約.
我們看這些經文的時候.
在我們新日的地方.
很迂腐.
很馬過.
有什麼理由要推廣那些人.
不過你要記得.
他們是猶太人的社會.

$^{681}$大家要記住.
《兩表書要談》這個文章.
兩位的領袖.
以色列 尼希米.
都很想帶領自己的朋友.
要成為神聖節的氣候.
是和外方人的想法.
那個做法.
是有不同的.
所以我們就不難理解.
他們要回到.
起初摩西人的禮拜六.
是安息日.
一切的功都不可以做.
要回歸那個安息.
安息本身是極之重要的.
所以我們又要回頭.
我們是否真的.
每個禮拜.
至少一天.
甚至可能半天.
我們心裡是預備好.
是真的有所準備.
來去朝見神 敬拜神.
是很重要的.
不要按你自己的時間.
按你自己的喜好.
這樣去想自己的.
一個禮拜的作息.
工作 娛樂 或者運動.
應該是以七天的第一天.
作為我們想事情的單位.
然後最重要的我們做了.
其餘的時間.
代表我們要敬拜神.
不過我們在不同的處境.
神安排你人生的崗位.
來到那裡.
在那裡為神去工作.
為神去事奉.

$^{721}$為神去作劫.
所以當我們回到這裡的時候.
我們記得大前提.
他們不是為了要考密卓賤.
而是那個時代的需要.
更宏觀去看.
那個時代需要這樣表達.
自己是屬神的猶太人.
因為現在我們是回到.
猶太耶路撒冷.
這個聖城.
所以他們就有這個.
那個時代的需要.
要成為一個時代的工人.
就是要帶領自己的民族.
猶太的民族.
要讓外邦人見到.
我們是屬神的.
如果我們自己都安身立法.
記得因為耶穌還沒有來到這個世界.
那個時候是雨前的四百多年.
還有四百年耶穌才來.
所以這四百年.
其實是需要保護這個民族.
是要認真地看自己是神的選民.
這個大前提.
所以是有他那個時代的.
一個重要的領袖的角色.
這個這樣的里程碑.
所以因為這樣的緣故.
一開始這樣做是很極端.
但是當他們明白了初心.
就是要表白給外邦人知道.
我們是神聖節的選民.
這樣就情有可原了.
不過當時的手段是很激進很專門.
所以要把那些破墳.
不要把這些外邦人趕走.
准他們入城.
第十七節.

$^{761}$「我就斥責猶大的貴族.
你們敢不敢行邪惡事呢?」.
然後等等.
第十九節.
「在安息人的前一日」.
就是星期四.
耶路撒冷城有黑警的時候.
我就吩咐人將門關鎖.
不過安息人不准開放.
即是一個手段就是.
將城關了.
遊貪人可以出去.
在自己的城裡面.
外邦人做生意的.
可以進來買賣.
我們今天就不會有這些問題.
我們根本已經.
有時去到另一個極端就是.
我們都不把星期日自己吩咐來吃.
你自己吩咐來吃.
你自己變成了國家.
在進行爭取法定的國家.
你都真是需要來到安靜.
與敗事.
與姐姐們相交.
這個是標籤很重要的.
生命的循環.
所以你差不多不用想了.
一個人他如果慢慢離開教會.
再過個星期返回教會.
再緊張返回教會.
他的力量一定會穩縮.
甚至他的信仰.
就是信一定會.
你不需要再和我見面.
一定他會離開了那個中心.
他自己就去中心化.
他就沒有了那個生命的焦點.
他就會慢慢自己的人生.
代表他不開心.

$^{801}$他可以有很多的積蓄.
很多的時間.
可以去哪裡旅行.
可以去吃什麼.
可以是很開心的.
家庭都可以很快樂的.
不過他的生命裡面.
那個熟悉焦點在哪裡.
如果更加快讓他發現.
我這樣做是為了什麼呢.
我的人生為什麼會這樣做呢.
很多時候就是家人不和.
自己有重病.
事業不順利.
他們才會想起神.
所以不是每個階段.
會想起神.
是人生九頭無路.
才會想起神.
往往就是這樣的.
某程度上不可以代表神.
某程度上是神.
其實你要回到那個焦點.
人生的生命焦點.
很多小時候回教會的人.
去到有兒女.
事業.
人生很多事情都順利.
他們才開始慢慢.
發掘自己童年的那個共識.
有時發現.
年年過了幾十年.
就浪費了很多很多的時間.
很多的青春.
所以這裡告訴我們.
李希米這樣做.
是很激進.
很極端的.
極端到新日.
我們覺得這些是一些熱端.

$^{841}$但是我們記得.
這些經文是由他的時代來寫的.
是要讓.
這個被老歸老的民族.
要有回那個.
自己的身份的形態.
對自己的身份的形態.
要有一個很不同的意識.
所以.
以斯拉和黎童文.
他們不是一些很淵源的.
好像很白譜的一些老人家.
說什麼都真的要.
要那個街這樣.
都不是這樣.
其實他們是一個很強烈的.
一個改革家.
改革家.
他們是一些貧窮家.
他們是要將神的.
摩西的律法書.
來去灌注在.
黑社的內心.
當然了.
你只是這樣治理.
你就會有很多的衝突.
一定是要那個民族自己本身.
他自己變化了自己的身份.
形容自己的身份.
是甘心忘義地去做.
那就是一個真正的.
黃興民.
所以.
今天我們幾多教會.
是不是我們連教.
鬧我們都沒有了.
每次都是很開心的.
回來唱詩歌很平安.
那這些完全不是錯的.
不過是不是我們已經沒有了.

$^{881}$屬於的.
那個主人.
在我的教導.
以致領子們.
他都不知道.
什麼是神的心意.
和他自己的屬靈的身份.
到底他有多少的奇形.
我們先到下面.
第20節.
「於是商人和販賣.
各樣貨物的人.
兩次住宿在耶路撒冷城外.
我就警戒他們.
你們為何在城外住宿呢?.
若在者我便下手難辦.
從此以後他們在安息日不再來.
我吩咐你們人見證自己.
在守城.
在安息日不吃.
我的臣啊求你.
因這事記.
照你的大慈愛.
回信我兩次」.
都是有一個主題.
首先就是.
讓利未人.
能夠有那個.
公益.
那個富國.
有那個秩序.
不要成為多比亞的湯館.
另外就是安息日.
發現了.
人家來買賣.
你可以叫他不要買.
但是這個就好像.
變成了.
引誘了.
引誘了猶太人.

$^{921}$深思思.
很多東西買.
安息日沒有了那個專心.
安息日.
都專心去敬拜神.
去朝天神.
安息日就只顧著買賣.
只顧著那些外邦王.
又賣魚又賣.
可能賣肉.
還有其他的各樣的國民.
商人和白賣各樣國民.
他們很聰明.
在城外住.
讓安息日.
安門就進來.
好像.
不關我的事.
我們在城外住.
我們起碼說城外住都不行.
引誘猶太人.
深思思.
每天都顧著去買東西.
其實是買賣.
你可以想像得到.
黎希敏就當這些是小朋友.
大支科板糖在桌面都不行.
看到都想吃.
沒有了那種自制的自律.
好像是真的沒有了.
就令到黎希敏要做這些.
這麼厲害的手段.
說連那些外邦王.
在外面做生意.
在外面租地方住.
來等安成.
即是說猶太人購買力都很強.
可以這樣來買賣.
其實買賣都是正常.
現在純粹是安息日.

$^{961}$這樣都不行.
但在黎希敏的角度.
正正就是安息日.
表明猶太人.
這個分別行乘的新文化.
平日.
六日你可以買賣.
你可以做你的生意.
是沒有問題的.
但安息日就應該大家一起.
要仍然真真地.
安靜下來.
敬拜神.
安息日不是甚麼都不做.
這是去到靈學耶穌的時代的結論.
所以他就繞往過程.
但是.
安息日最後最終的定義.
最終結就是.
要安靜地.
去到神的面前.
去思想神父.
做物主和受道之民的關係.
最好就是不同的人.
有不同的背景.
不同的那一輪的際遇.
和我們中國人說的那一輪的運數.
那一輪的運程.
當然這些是民術.
每個人都不同.
你.
六旺的時候.
你又去教會.
有些人三衰的時候又去教會.
現在就有這個互動.
如何能夠彼此的安慰.
彼此的鼓齒.
諸如此類.
有些人很順暢.
有些人很倒楣.

$^{1001}$有些人一直這樣衰運.
你就要明白.
人生就是這樣.
你就要過個星期六.
安息日就要去到聖殿.
敬拜神的時候.
你就會去看自己的人生.
因為有更大.
創造的主義.
某程度上是表現你人生.
表現你命運的時候.
它會來指點你.
你人生的終極.
你的方向.
你的追求的目標.
你就有那個方向.
所以不是每一個星期都很清晰.
而是不斷不斷的潛行.
潛行.
是會讓我們有一個新的平安.
如果我們內心真的有個指南針.
是不斷告訴你.
其實我們自己人生.
都是向天加的路.
人生.
其實不信人都相信這件事.
都知道這件事.
但是他只不過是不信.
人離開世界.
就繼續存在.
但我們要求如果是信有這個存在.
那我們就應該要.
真的要不斷的.
去索取我們自己.
要有一個定位.
我們自己的方向.
否則就擺來擺去很大問題.
我們到下面.
石章聖經.
是想帶出.

$^{1041}$黎希米.
大大小小的行政手段.
好像一個.
獨裁的領袖.
如果用黎希米記.
有些坊間.
就是用黎希米記.
教你做領袖.
坊間很多人都很喜歡黎希米記.
因為他一除靜音.
做很多事.
有時他會有很多角度去想.
有時不用這麼多.
一除靜音就要這樣做.
是一個很powerful的領袖.
很多人喜歡他.
所以黎希米記.
在坊間很hit的.
管理學的過程.
不過.
你又考上他.
黎希米都沒有想過自己是一個.
管理學的領袖.
他不是一個EMBO的.
一個工商部隊.
的高層生.
這本書.
是想帶出.
黎希米的領袖.
其實他外方高官.
科斯第二代是糾正.
但是.
他是以右派的身份.
很想建立.
自己的民族.
對自己神殉民一身.
的認同.
所以這本書.
其實是講到一個有信心的領袖.
怎樣去和國民的.

$^{1081}$坑民.
用他的智慧.
用他的熱情.
去感染他那個時代.
的右派的同胞.
既然我們已經.
身處科斯帝國的.
內部.
我們已經回到.
我們回到城城.
回到這個有成年人化的地方.
我們.
是不是應該.
有這個成年.
國民身份的認同呢.
所以就會出現了.
這樣的案例.
所以兩位的領袖.
其實都是要帶盡.
認識他們的同胞.
當然那個人不是很好.
那個人要跟進.
是帶來很大的衝突.
但是聖經.
只是說他有這樣做.
但是那個後果是怎樣.
聖經就很交代.
我們可以知道.
人是變動的.
領袖是會有抗衡的.
有些外邦人是會有.
可能是會.
大兇殺他.
可能是會用.
其他的手段去抗衡.
所以.
如是者經過了.
四百年了 到了耶穌的時代.
猶太人對於.
自己的身份.

$^{1121}$又過分了.
見到耶穌是很形式主義的猶太人.
所以為什麼那些法利賽人.
成為了當時.
人人都不怕.
人人都不敢得罪的宗教領袖.
就是因為.
正正就是這個法利賽主義.
就是將.
喜米 以斯拉.
他們這些.
激進的元素.
真的照單全收了.
所以什麼都不理會.
我們認為安息人要這樣做的.
所以他們訂一些很清楚的條件.
你的鞋差不多爛了.
安息人就不准穿了.
因為你爛了.
你的辦公室就是他的了.
你現在穿的像玩鞋.
你回家就犯了安息人的罪.
因為你拿鞋的動作.
就是工作了.
所以你就可以.
在安息人身上穿一些差不多爛的鞋了.
然後.
你又不可以穿超過幾千蚤.
因為聖殿.
這個好像是由舒雅記.
那個約棍離開.
那個民族.
過了大門的時候.
他們應該要穿一千蚤.
現在都要距離.
所以他們認為二千蚤就是一個合理的距離.
你走路要走遠路.
你就是犯安息人的罪.
所以說.
如果我們今天跟足了.

$^{1161}$安息人就不要出門了.
就不可以坐飛機.
因為你已經離開了超過那個地方.
甚至你可能已經住在廣州了.
你不可以住得太遠.
因為你不住太遠.
你就犯安息人的罪.
所以到了耶穌的時代.
已經是那個.
精神背後的那些.
法輪功學者已經超過了.
聖經所說的那個.
初心的原則.
變成了很重的擔.
所以耶穌罵那些.
派系上的人.
你們將那些冷淡的擔子.
壓制在那些.
不怕死的人身上.
但是自己自由都不夠.
就是說.
解釋得很難.
很絕緻.
安息人要這樣這樣這樣.
耶穌就不斷挑戰.
甚至否定他們.
這樣的指令.
所以你會看到那個時代.
一直來演進.
到了耶穌的時代.
差不多是一個.
律法主義到極致的時代.
誰如果.
違反了猶太的法規.
好像是613條.
那個法規.
你們就是犯了神的禁忌.
這樣去讓.
百姓其實活在一個.
律法主義的枷鎖底下.

$^{1201}$現在起碼.
表面就是極之嚴峻的.
獨裁領袖的.
說了是.
但是如果我們這樣看.
又是有一個誤會.
其實他是要帶領這個民族.
有關神的選民的身份的形容.
但是當然.
是不是可以勸說.
是不是可以有一個比較溫柔的手段.
甚至乎一個過渡.
這就不是這本書裡面所談到的.
到下面這個比較長的.
23書.
馬錫日子.
「我有一件猶大人取了亞實達那官.
汪洋的女子為妻.
他們的兒女說話.
一貫是亞實達的說話」.
即是那個方面.
已經不懂得說猶大的話.
「所說的是照著各族的謊言.
我就斥格貪婪.
就坐貪婪.
打掉貪婪幾個人.
拔下貪婪的頭髮.
叫貪婪以著神之勢.
必不將自己的女兒.
嫁給外方人的兒子.
也不為自己和兒子.
娶貪婪的女兒.
我又說.
以色列王.
所羅門不思在野弱的事上犯罪.
在多國中並沒有一王著他.
且無他臣所愛.
實立他作以色列全國的王.
然而連他也被外方女子人有貧戾.
如此.

$^{1241}$我豈聽你們含冤大惡.
取愛當女子肝膀我無敵神.
表面看就是.
千多年前.
沒有千多年.
因為那時候主持四百多年.
所羅門的時代.
是主持了大概九百多年的時間.
你想想.
五百年前.
我們的太公所羅門.
他這樣做.
在晚年的時候.
是一個政策.
即是說.
我娶了周邊的國家的民族.
迦南的民族.
娶了他們.
即是說.
是一個政治根因.
不需要打仗.
我已經得到他們的地方和尚.
大家相安無事.
大家和親.
好像匈奴.
復漢保帝.
結親.
變成了嫁進去.
你不需要打仗.
大家就可以有個民主.
可能不可能這樣想.
但是.
殊不知.
他那些的秘訣.
就引誘他去拜神.
拜不同的國家的神.
直到所羅門.
慢慢的.
他.
民間.

$^{1281}$都可以拜他.
因為你自己拜.
別人就不出拜.
保佑香港.
所以就帶來.
整個國家.
傾向於再拜神.
所以第一個就已經.
爸爸也是大位.
這麼敬佩.
但是所羅門就立即.
還有一個很聰明.
很有智慧的領袖.
非常出色.
神都愛他.
他是佛事家所生.
但是.
居然在晚年的時候.
所羅門自己.
他跟隨他的秘訣.
就拜那些當地的假神.
例如科羅.
阿斯塔諾.
亞瑟拉.
等等這些假神.
這些教會出名的.
還有巴力.
這些神.
就使得整個國家.
傾向於再敬拜神.
所以.
你就找回這個案例.
五八年前.
當時的形勢.
國家.
人數這麼少.
但是.
他也是看這件事.
看得很近.
就是希望有派的民族.

$^{1321}$要有承認的國民身份.
這是一個很傳統的家常.
但是.
似乎聖經是說.
這是他們兩位領袖.
在那個時代.
有他的需要.
這樣去承認這個民族.
以這個身份.
去承認.
那個人的體格.
那個人的經歷.
所以.
那個手段.
是非常之嚴重.
非常之根深柢固.
所以我們今天.
很難以這些舊的經文出來.
參照.
我們可以拿他的精神.
背後的精髓是很重要的.
是跨越任何時代的.
但是.
那個手段.
就已經不可能了.
你說.
我們是宗教財團所.
誰如果妄大錦繡堂的那個奴隸.
就要怎樣怎樣.
豬狗做鴨子.
不可能的.
當然是所有的人.
他自己去做一個選擇.
但是我們要將那個天若神奇.
放的.
新若的.
在耶穌還沒回來之前.
我們那個屬靈的身份.
我們怎樣來形容這個身份呢.
怎樣讓那些未信的人.

$^{1361}$能夠知道.
分別到我們內心.
是有神.
是有信仰.
是有福音.
是有眼望的.
這個就是那個分別.
這個就是我們需要.
來去.
保持守護的.
接著就是28節.
大祭司爾利亞神的孫子.
耶和尼大祖一個兒子.
就是和尼.
參巴拉的女婿.
我就從我這裡.
把他趕出去.
我的神啊.
求你紀念他們的罪.
因為他們.
玷污了.
玷污了祭司的職任.
違背了你.
與祭司你們人所立的約.
這約.
我見證他們.
使得他們來傳一切愛光華.
派遣祭司和你們人的班次.
使得他們各盡同職.
我又派百姓.
按定期獻茶.
和初熟的土茶.
打弄他們自己的功行.
他們自己的身心.
我的神啊.
求你紀念我.
施恩衣服.
這個就是尼天麥.
他好像有一些自傳.
有一些.

$^{1401}$好像自傳式的傳記.
來講他自己所做的.
所以整個天典是他自己.
作為領袖.
的核心.
帶動著.
這個民族.
在他現在的現狀.
一開始.
這個民族其實已經沒有了.
那個廣義的身份.
與愛光人已經沒有任何的.
分別.
所以他們的.
宗教是受到.
愛光人的教授.
和利用.
很多時候都割低到很低的標準.
安聖人就跟那些.
愛光人買賣.
沒有任何的界線.
就是說其實沒有什麼宗教可言.
方便自己就可以了.
那就.
變成沒有了.
宗教的一個.
太民族身份.
有的已經沒有.
儀式的.
一種很僵化的儀式.
的一個.
宗教.
帶不到一個很.
良麗的一個.
去令到.
自己.
這個民族可以.
一路這樣去.
生火相傳.
讓自己的同胞.

$^{1441}$下一代感染到這份極權.
相反.
外光人都.
嘲笑猶太這個民族.
就是說.
實際上可能是.
拜金錢.
拜.
娛樂.
拜其他的東西.
就會帶來.
在這個時代.
就會有一個很大的.
一個保佑.
我想做一個小小的簡介.
就是對.
以色列第七到第十章.
這個.
復興.
由以色列.
去感受到.
回到上來.
感受不到猶太人.
有一個的屬靈.
當然當時他們不叫屬靈.
就很實實在在.
猶太人.
就是應該這樣做.
因為當時還是在律法的時代.
他們還是受著.
摩西律法的規範.
當然在波斯帝國那些.
可能已經沒有規範了.
但是.
現在你上了.
這個是一個中心點.
因為將來救世主.
在這裡降生.
所以神某程度上.
是有一個存在性.

$^{1481}$讓地老歸原.
不是人歸原.
是要人心歸原.
什麼是人心呢.
就是猶太選民.
的一個心.
要歸原成聖.
在那個聖殿裡面.
是無懼以色列的建制.
而是建制的人.
要知道自己.
是在摩西律法之下.
守神界的一個.
神臘那樣的.
以色列的民族.
這個就是神的.
自己存在的那個神.
所以如果這樣的時候.
去到新的 去到耶穌的.
回歸之前 沒錯.
我們不再需要紐約那樣.
去死守一些的界定.
但是那個新民.
已經是由紐約的猶太人的新民.
轉化成為神選民.
神子民的新民.
所以譬如.
《歌羅西書》和《三章》.
就有這些經很多的.
保羅在每一個.
教道裡面.
都強調在這個新民裡.
就是三章.
《歌羅西書》.
三章就是神的意志.
《歌羅西書》.
所以你們既然是神的選民.
完全不是猶太人.
哥羅西就是奴隸家族.
七教會.

$^{1521}$即是古前的那一代.
你們是神的選民.
即是說.
外邦人現在因為耶穌.
我們都成為神的選民.
Crucifed people.
是Crucified.
是神選的.
成為聖傑國外的人.
很明顯就是.
哥羅西十誡裡面.
十九章.
三誡神將這個民族.
好像鷹那樣.
背負在艾林.
進入迦南.
成為一個新的民族.
是神和紐約的民族.
現在保羅.
忽然間對外邦人.
這班哥羅西教會的人.
都是這樣理解.
你們每一樣都是神選民.
是聖傑的.
是蒙神所愛的選民.
這又如何呢?.
當然要穿著.
猶太的衣服.
穿得很圓.
要留鬍子.
可以看鬍子.
要戴隱形眼鏡.
遮住自己.
要戴光面具.
要在天上.
不得地要堪稱.
在這些儀式裡.
不是在通常的儀式裡.
這樣就能夠救主的位置.
救主已經在這個時候.

$^{1561}$就是這樣.
在這裡感受救主的光.
希望救主大於我們.
問題就不是這個.
新人.
很清楚的教會.
救主要穿上.
穿上的意思.
穿上的意思是.
穿著衣服.
猶太的外表的衣服.
更加不是我們基督教.
有制服,有衣服.
那些是在福利堂很需要.
但是真正在人生裡面.
在微信的眼前.
我們要穿上這個制服.
就是要穿著.
有恩慈的.
有慷怜的.
有青春世紀的.
很有誠意的.
或者是.
有誠意的.
在這些的.
青春世紀.
恩怪氣.
這些人.
心裡面.
是很有那些.
很有那些.
很感恩的心.
所以必然.
要穿上這個衣服.
這件事的一個.
就一定是正確的.
即是你做錯了.
你都運氣好.
但這些人有一個.
可以讓他們成為一名警察.

$^{1601}$有恩慈.
有謙虛.
有溫柔.
有隱惡的心.
接著具體的有案例.
倘若這人.
與那人有誤會.
怎樣發揮你內心.
選民的特質.
就是你本來就去.
恨他.
成為一個犯法的人.
有些人的命.
你和我都有.
各不在此.
但有了現職之後.
就要做犯法.
所以要成為一個盡快洗到容不下.
容不下的人.
同樣的事情.
會成為大家的衝突.
這些是很難學的.
但本來就說.
選民就要這樣學.
所以是勸說的.
是給一個方向給的.
是禁止你這樣做的.
誰和誰就立刻和好.
立刻放棄.
是這樣的.
是你自己去選擇.
起了一個案例.
倘若.
這人與那人.
大家知道的.
你見過誰與誰.
是不和的.
他們後來和好了.
但如果成了一個案例.
可以說曾毓思.

$^{1641}$這個會友,這個領袖.
是不和的.
後來大家都知道.
他向對方認錯.
因為犯下性侵.
大家自己很不熟悉.
保證了大家的關係.
而且那個關係更加和和.
更加的良心.
這些就是我們受選民的.
身份.
就好像紀密上.
不過紀密.
他的時代就需要用.
一個特別的一重一的.
一個不同的方式.
一個手段.
但是保羅.
保羅是使徒.
他都不會迫使.
人們你叫他立即.
都不可能.
但是他會講了.
基本上.
將我們的內心.
處理了這些神.
就延闢.
就變成.
會加增.
就變成增濃.
之後就會變廣.
那個榮譽.
就會增進.
繼續這樣.
延伸下去.
就會變成.
變成一個保存的.
一個廣泛之增.
可能這樣就是意義的意思.
總要給錢包容.

$^{1681}$給錢要恕.
主怎樣要恕你們.
你們又要怎樣要恕.
在這一切之外.
要存著愛心.
愛心就是.
聯絡傳德.
所有的德行是由愛心那裡.
來滋養的.
沒有愛心你都做不到.
沒有愛心你怎樣能夠對自由.
可以有一個平台呢.
你怎樣可以不認同.
他可以有改變呢.
沒有愛心你看到他的樣子.
你都走開了.
你怎樣能夠建立.
大家會更加好的關係呢.
所以就要叫.
基督的平安.
在你們的心裡面坐著.
基督的平安.
基督如果被人.
侮辱.
嘲笑,被毆打.
他都保持他的平安.
他知道他們這個世界的使命.
就是要這樣受.
他知道神.
是不會看錯的.
人可以看錯他.
說他稱呼自己是奴才.
但是要很多年之後.
包括保羅自己也是.
很多年之後.
才慢慢認定.
耶穌是他的主.
是他的救贖.
所以基督的平安.
就是說就算我被人侮辱.

$^{1721}$我也仍然要保持.
我和神之間.
這一份的.
往生的相通.
那份平安.
以致我能夠所做的.
是神所要求.
因為.
你們是為此無恥.
歸為一體.
所以這裡是說.
上面的有言的.
盡快保持.
有誤會的.
盡快去保持.
盡快保持大家的誤因.
大家一起.
所謂何事神之一舉.
就是與人間第一次的一件事.
一同建立神的國度.
你可能有不同的意見.
不同的意見.
但是最終.
如果還是要盡快有一個決定的時候.
那你就做那個決定.
那個決定如果不是你那個決定的意見.
那麼你就跟隨.
那個大量的意見.
這很重要的.
總之大家都是行前一步.
大家都是為神.
這個很重要.
剛陽放若的旨在大眾心.
凡事都是要一念一念的.
全部都是你自己辛苦了.
你能夠生活.
你能夠有這個機會.
其實都是人.
你能夠可以起床上班.
可以有四肢的健全.

$^{1761}$可以行道.
但是每個人生來的那一天.
有機會帶大家去看看.
那些真是險重的.
傷殘有知的人.
那些人很慘的.
很可憐的.
我們要用那種很大的愛心.
去關愛這些人.
但是我們自己是正常的.
不要假設我們.
不應該是這樣的.
如果我們是基督徒.
你要去到你的內心越來越簡單.
那個心是越來越簡單.
就是恩典.
沒有其他的解釋.
是恩典味道你能夠.
嘩!吃這種味道.
嘩!牛肉是這樣的味道.
雞肉是這樣的味道.
蝦肉是這樣的味道.
這就是恩典.
你英文的那些.
你去到那裡.
吃了任何的.
食物.
都會咬著膠.
咬著雞.
咬著膠.
想想想.
但是你要知道.
一切是恩典.
我們有很多的.
千萬的那條地.
有感應的.
我們能夠感應到.
我們是基督徒.
所以這樣的口味.
嘩!其實是薑有原味的.

$^{1801}$你想想多麼的.
我們的享受.
沒有了這個之道.
你就知道多麼的難.
所以那些真的有難.
全程都是一樣的.
不一樣的.
很痛苦的.
所以我剛才說了.
你要知道.
要穿上這個新的身份.
新的衣服.
就是前愛.
前恩.
恩情.
天卑.
溫柔的心.
是你的心的能力.
接著就是剛陽割藥的指導.
基督的道理.
空空自然的存在.
心如融獅將具.
情如蚵 彼此交融.
彼此交融.
無論作神光或說話.
或行事都要同主耶穌的名字.
極著他.
好像將這個.
薑土人生.
是一個歌唱的人生.
歌唱.
內心充滿了.
詩歌 充滿了神的真語.
你對人生.
越來越透徹了認識.
你知道人的短暫人生.
是為將來的一刻.
所預備的.
你就要珍惜.
自由的時間.

$^{1841}$自由的金錢.
自由的精神.
去盡力去做事.
不要再浪費.
雖然不是說.
我現在需要休息.
需要調劑.
是走舊意的路.
是對的.
正常性只是娛樂.
你都沒有想過事故.
其實是一個很難容忍的事.
所有這一切.
是斷斷續續的.
是平靈的.
是要你有意識地.
知道你自己已經有了.
這個新的選民身份.
同阿彌奈.
所做的.
是極端的.
局才領袖.
這樣的一個人.
你就知道原來.
那個時代真是有意義的.
耶穌未來的自然事後.
又是地老歸回.
因為神已經群了這個民族.
在這個民族的再重新開始的時候.
神就興起了.
以斯拉和尼歇米.
這兩個人.
很強烈地.
對自己的身份認同.
然後.
就在他們的日內.
在他們的人生.
可以事故的那個年代.
就在他們的平台.
帶神感動他們.

$^{1881}$回到耶路撒冷.
那個應該就是.
神選民要表現.
他們那個身份的地方.
來表達.
對多西六法.
那個的尊重和尊嚴.
這個又帶來了.
那個.
再有新意識的.
顯示.
新的時代.
百姓.
怎樣來聽取.
這些領袖的教導.
就是宣講.
多西的六法書.
因為當時神還沒有來.
一些新的都沒有來.
耶穌都還沒有來.
根本都沒有子徒.
在以前四百多年.
要記得.
那時候只有多西的六法書.
是他們的寵成.
所以他們要有.
這個身份認同.
就需要來跟隨.
這些領袖的教導.
即是說突顯了一樣東西.
開始就是有.
宣講的同時.
宣講什麼呢?.
宣講一些尊嚴.
一些啟示.
將先有的六法.
有系統的去傳承.
去做那個教導.
做那個啟示.
所以有一個.

$^{1921}$將來的人學制度的前世.
那個初年.
所以猶太人一直就是有教導的.
宗教.
所以猶太人是很有文化.
小時候就必須要學子.
成都民族.
很緊張時度.
所以一定會從小就要讀六法書.
所以他們第一本.
在小學打到大學.
他們的心很悶.
很濕.
但是猶太人從小就不這樣.
你都不知道多開心.
所以猶太人是有很多機會.
去接觸神的靈體.
當然慢慢慢慢就變成了.
六法主義是不同的.
不過在尼西蘭的時候.
在以斯拉的時候.
他們是需要.
這個改容的復興.
這個改革是跟隨著.
所羅巴巴第一批的猶太人上去.
是很艱難地.
建立了成年.
建立了根基之後.
就開始那個勁發.
慢慢當然是純粹建制.
有他們自己的民間的生育.
就是復元.
開始有他們自己的地方.
有他們自己的大圈古墓.
接著每個星期的安息日.
就上聖殿.
預算都過了大約八十年.
第二批就是以斯拉上來.
就發現了.
他們猶太人已經是.

$^{1961}$沒錯.
還是猶太人.
每個星期都去聖殿.
但是有些事情.
已經習以為常地在做.
其中有幾件事就是違反安息日.
安息日都不在聖殿.
安息日就是要求猶太人.
去做其他事情.
這個是一個情況.
另外就是最嚴重的.
就是以斯拉的起源.
似乎是他剛剛發言.
剛剛開始.
而是已經有全面的.
即是說你要真正.
需要很大的力量.
但是他們的口都認為.
做不成.
做呢.
再做下一代.
再下一代.
就會變成這樣.
又會再招致神的神會.
所以由以斯拉和自己的弟弟.
做完之後.
某程度上.
猶太人真的不再是.
完全不再是安的領事.
即是說他們都拉了兩位領袖.
那個行政權力.
即是說因為言辭的警告.
如果再這樣做.
你們小心.
好像所羅門的結局.
神都可以再一次.
將我們起源.
將我們拋到外國.
所以猶太人都害怕.
所以你可以想像得到.

$^{2001}$雖然兩位領袖.
用這麼重的行政手段.
和以斯拉.
是叫人犯約.
而希米呢.
因為他是官員.
他真是.
哥斯拉是和工的長者.
所以他真是用一個政治.
應該說是白人的手段.
就趕那些人走.
特別是哥比亞.
佔據了那個富國.
做他的tong 房.
就趕走.
又將那些利益.
招聚回來.
你們放心.
做事是會有收入的.
我要叫百姓.
將那些憲制的十分之一.
給你們做下流的.
你們是可以有足夠的感性.
你們要專心的事務.
不要再做假經濟.
做當勞經濟.
學了我工業本身的當勞經濟.
陶清潔在後欄洗碗.
如果有就告訴我.
因為反映了他們不夠錢.
或者他們被勞.
要處理這些問題.
到底是什麼原因出現這個情況呢.
當然不是開玩笑.
香港的全部人都能夠.
能夠解釋這個問題.
除非你是一個官職.
其他事當然是自己的事.
我都說德國的.
你又搞又認不出來.

$^{2041}$因為一定要登記.
認不出來工業.
香港希望將來的這件事.
如果是好的人不好.
這些要研究.
好的地方就是專心.
就好像李健偉所說.
但不好的就是.
政府要研究這個地方.
又不可以說一些違反憲法的東西.
所以我們都要看看時代是什麼.
所以如果能夠我們在.
在那些精密之後.
再用新約的這個眼光.
去聯防.
和新約怎樣來達到.
這些帶來的律法主義呢.
新約其實是想怎樣呢.
其實怎樣去信徒.
更加封了自己這個.
新的選民身份呢.
而且這個選民.
再走兩派的國旗.
是任何.
只要在耶穌的救恩裡面.
就是神選民.
那就精彩很多.
也難很多了.
因為不同的民族.
自己民族的特性.
或者那個習俗.
怎樣將那些習俗.
轉化成合適於真理的.
新的信徒的真理的.
看法,教導呢.
這些就是需要有.
不同的民族去詮釋.
這個真理.
若信徒在他自己.
那個時代的處境.

$^{2081}$怎樣去.
自己去揣摩.
自己去舉一塊衣服.
甚至最好就是.
用其他企業主義者去討論.
譬如我還可以去抽煙嗎.
抽煙不代表健康.
但是我沒事.
是不是我可以去煙呢.
我自己在家裡吃的.
我不影響別人的.
我回來教會就不吃了.
是不是行不行呢.
很多這些邏輯的問題.
其實都是要.
就可以帶來一個討論.
不過其實都是這樣說的.
新約似乎已經沒有了.
一個企業主義者.
來之後.
可以這樣去討論.
很多時候都是一個.
亂說一個的.
一個告訴你.
這個身份.
才是十幾月的.
我們要視乎這個身份.
我們要不斷地.
掃這個身份的.
比例和純明.
以致大家嚮往.
這個身份.
那麼你就自然會.
自己達到.
這個身份.
那個的.
就是.
譬如你很想.
躋身一些很有文化的街坊.
你沒有理由.

$^{2121}$在他們面前站在褲子上.
去試.
你的身份都是.
黑背的.
甚至不要衣著.
你都不應該是.
有太過的.
這些我們都明白.
所以基督教其實我們.
不是講文化.
我們是要告訴人知道.
如何我們現在在神.
耶穌的救恩裡面.
已經擁有這個.
一般人都不可能達到的.
這種純明.
就是神選擇.
這個身份.
還有聖靈已經在我們裡面.
讓我們的生命.
隨著那個.
神的生命.
學習.
我們是會轉變.
我們學習.
生命的轉化過程.
所以這個就是.
基督教為什麼和.
其他人的宗教.
對比的地方.
就是我們不是學了道理.
我們是有了新的.
神的生命.
和聖靈.
在我們裡面更新我們.
使我們能夠.
可以對神的.
期望的一個道理.
這麼美的基督徒.
真正可以散發出.

$^{2161}$那種氣質.
就是純明主的氣質.
你想想一個人.
就有神的氣.
就有愛.
在我們的內心去做.
愛是.
言諾全德.
所有的真行是由愛.
來出發.
所以我們為什麼會追求.
成信.
做一個天兵.
做一個去派對.
原因就是我們是愛人.
我們是.
希望得罪人.
我們是希望.
能夠成為別人的.
奴婢.
我們是希望成為一個使人.
的奴婢.
這不只是聖經.
是一個國神.
天國的國神.
所以.
弟子們.
這其實是非常寶貴的教育.
所以教會.
其實是要多些.
我們這個新的神.
在這個地上.
信的人.
當然一定有衝突.
但是應該更多是發揮.
我們那個純明的.
那個意念.
純明的那個.
香氣.
要人試著.

$^{2201}$我們屬神的.
怎樣去追求成信.
今天.
看看我們有沒有機會.
去開兩分鐘.
不如我們.
下一次我們才正式.
進入幾個記錄.
今天晚上是總結第十章.
的以色列記錄.
以色列的行政手段.
很不錯.
不過他只是登記.
作策.
逼那些人立即做.
但是你希望他.
因為他是省長.
他是有和給他的.
媒體預知.
甚至有他的軍隊.
運送那些文資.
所以他本身.
就是一個這樣的.
地方官.
省長.
但是他更加.
在自己是猶太的領袖.
是神.
他有這個慰問.
感動.
去做這個見證.
要使自己的同胞.
嚮往.
當然是立即.
立即.
要體現自己是神選民.
猶太人.
才是神選民.
因為這樣.
外邦人就是外邦人.

$^{2241}$神選民這個神.
所以.
兩段地老歸回的.
字孫.
都是想帶出一個信心的功課.
就是怎樣.
可以盡快成為.
神選民的身份.
這個就是以色列.
希望的主要宗旨.
當然加上的.
就是建造聖地.
建造城牆.
這個是一個次要.
重要的是那個選民的重建.
人心的回歸.
重建.
那個自己是猶太人.
神立約的選民身份.
最後我想和大家講一講.
第十九章.
《純凡人記》.
重要的.
相當重要.
大家要明白要熟悉.
雖然我們不是猶太人.
但是那個精神.
是很可能和你們立約的.
以色列人出埃及以後.
滿了三個月.
那些人就來到.
西乃去抗美.
你現在收到的信息.
他們從鳥離飛地.
來到西乃去抗美.
就在那裡的山下安營.
波西到神那裡.
然後說從山上呼喚他.
說你要這樣.
報數大國家.

$^{2281}$表面以色列人說.
埃及人所行的事實在.
你們都看見了.
而且看見.
好像鷹一樣.
將你們背在翅膀上.
大龍歸鷹.
背在翅膀上.
就是鷹仔.
小雞.
鷹揹著自己的小雞.
叫他學飛.
你們有沒有聽過這件事.
這件事是真的.
那隻老鷹.
在一個程度.
從頭插到尾.
到那些鉛牌.
那些鉛牌.
飛的時候.
突然間打開了口.
應該是咬.
咬起自己的小雞.
然後就在高的鉛牌.
扔它下去.
說那隻小鷹.
因為遠在遠山.
異域之風.
有毒有疫.
但未有毒未有疫.
所以有了.
但不識用.
就急速掉.
急速掉的時候.
那隻老鷹就回來了.
當然知道.
看著你.
差不多掉到不能再掉.
它就會急速飛下去.
用整個翅膀.

$^{2321}$蓋住那隻老鷹.
再回到上面.
再一次掉下去.
那隻老鷹就.
終於識用了.
有一次有一天.
那隻花圃又是這樣扔下去.
這次.
不停地.
開始自己識用.
自己識飛.
然後就變成了.
叫做「二打花圃」.
原來毒物界都是這樣.
這個磁界.
與生俱來.
當時的學者.
都是看到這個現象.
才會明白.
所以神就說.
「我好似一塊肉」.
「將你埋住」.
所以你們經歷到的.
實在,有眼識.
有遺處,有可處.
這些是神職.
因為你們還要長大.
因為你們還要.
心智成熟.
你們還是小孩子.
對宋人的事.
仍然是一知六道.
所以不要緊.
但始終你不可能.
每次都要開一隻無邊的花.
埋住你.
讓你自己都要學飛.
永遠你都需要神職.
所以以色列人.
在這個階段就好像一棵木.

$^{2361}$他們就好似.
時空的這個階段.
一沒有東西吃.
就回不去了.
我們要回埃及.
我們要怎樣怎樣到達.
但我們都不怕了.
有些負氣喘不過來.
但神就叫他們.
有一個神.
神將他們背負著.
出埃及.
神期望他們成為一班.
很有意識.
很有意識的.
以色列神的學者.
所以就開始.
有一個律法的單據.
要揣摩那些律法.
要怎樣去實行.
第一就是要團結.
「今你如入神」.
聽從我的話.
遵守我的約.
就要在大圓剛中.
割鎖我的指紋.
因為傳來都是不正.
你傳不來.
即是有一個傳世代的宇宙.
都是我們要團結.
我們要拼班.
我們要成團.
即是同性.
這個性.
差不多無差都可以變成同性.
所以可能.
人人都會變成同性.
你就知道.
我們容大師祖的頭大人稱.
現在發大人稱.

$^{2401}$都是順序項目走大.
何況幾個月前.
還未計算那個順序.
就增加了.
增加了我們知道.
神父是願意和我們同行.
願意和我們同行.
但不知道你們要歸我作祭.
使我國民主光光.
這些說是要對順序.
所以容大師祖.
很重要.
神和奴隸是同行的.
他們何來安息人.
他們都是做奴隸.
做那個.
給他們埃及人.
做那個神祭.
埃及人怎樣給他們安息.
但神和奴隸.
安息人那一件事.
他們慢慢就.
變成好像很殷勤.
或者好像無所謂.
他們沒有去想.
你們是奴隸.
你們是主公.
你們是一個奴隸.
是我揹你們出門.
所以這一段是一個.
保持著大人心聲的.
一個男人.
一個奴隸.
是特別在.
亞當中特別有.
天主之位.
特別有得天王權.
也是神的祖先.
因為亞伯拉罕的信心.
神就.

$^{2441}$定義亞伯拉罕.
要祈求.
王臣的祝福.
所以現在時間到了.
神就差點摩西.
帶他們回亞當.
希望進入迦南.
成為一個很能夠.
遇神的一個神明.
一個神明.
就是這樣.
這就是起始點.
所以尼希米.
是回到這個原地.
回到這個情況.
所以保持著.
是神聽完這件事.
以斯拉,尼希米.
兩位的主人.
兩位都是當時的.
烏瓜伊尼的一個領袖.
在.
宗教上的.
以斯拉.
在政治或行政上.
帶來的.
尼希米.
希望能夠將.
短文的新聞.
帶回那個高度.
所以我們就能夠理解.
這樣的一個.
激進的神父.
有他的福音和一些信息.
今天就不可以了.
今天.
永遠都是你.
港人的一個尊嚴.
原者.
覺得失望.

$^{2481}$聽著繼續可行可說.
自己要遇到神.
離開神的那個情況.
是非常不妙.
你一定要過你自己.
那個過程的生命.
你就與他遇到神.
那個關係就很疏遠.
我們今晚就完結.
多謝主給我們.
今天文上可以再一次.
來到他們面前.
歸聊歷史書.
我們很盼望雖然時代轉變.
但是選民的身份.
其實.
有救若.
有主耶穌降生前.
已經神賦予猶太人.
但當主耶穌降生.
你是萬民的救主.
是全地的每一個人的救主.
所以我們今天.
成為基督徒.
我們有救恩.
主耶穌的寶血.
將我們見證.
使我們成為神的選民.
是聖見是猶太人.
幫助我們能夠在一生.
我們去追求對這個身份的.
認識和認同.
使我們能夠在有生之年.
都能夠反映.
救恩帶給我們的生命改變.
為我們的行為.
是救恩帶給我們.
心靈改變的力量.
聖靈在我們內心.
那個光照.

$^{2521}$那個人道.
那個力量.
使我們能夠.
不是標旗立意.
而是我們能夠.
可以散發一個崇明的香氣.
使更多人認識.
這個救恩.
我們這樣的感恩普告.
共耶穌基督的名.
有一件事要先說一說.
讓大家知道.
下星期要完一次.
因為要準備.
因為我十日深夜.
就要去英國.
有一個工作.
要去荷蘭.
和聖徒森道.
看宣教的工場.
去談宣教事.
代表教會.
所以九日我就要去.
第一次.
回來後.
我們就正正式式.
開始了這本書.
之後我們會再看.
撒迦利亞書.
我們所看的八個語言.
還有最後的.
以知絕記.
是另一個課程.
大家就要到.
三月底.
我才能夠再和大家.
分享正義.
就到這裡了.
拜拜.
\newpage



\section{歌羅西書 1:21-29}
\label{sec:M8De3gUONVA}
\textbf{2023-03-05 午堂《全力以赴的福音執事》歌羅西書 1\_21-29 講員:曾裔貴牧師}
\newline
\newline
連結: \href{https://youtube.com/watch?v=M8De3gUONVA}{\texttt{https://youtube.com/watch?v=M8De3gUONVA}} ~~~~ 語音日期: 2023-03-07
\newline
\newline
\hyperref[sec:E8bPEPFjp7c]{\small{< < < PREV SERMON < < <}}
~
\hyperref[sec:index_chronic]{\small{[返順時目]}}
~
\hyperref[sec:index_scriptual]{\small{[返順卷目]}}
~
\hyperref[sec:BXGw_xefZws]{\small{> > > NEXT SERMON > > >}}
\newline
\newline
歌羅西書 1:21-29
\newline
\begin{longtable}{cl}
\hline
\hline
章節 & 經文 (和合本修訂版)\\
\hline
1:21 & \begin{tabularx}{0.7\textwidth}{X} 從前你們與神隔絕，心思上與他為敵，行為邪惡； \end{tabularx} \\ \\ \relax
1:22 & \begin{tabularx}{0.7\textwidth}{X} 但如今，他藉著他兒子肉身的死，已經使你們與他自己和好了，把你們獻在他的面前，成為聖潔，沒有瑕疵，無可指責。 \end{tabularx} \\ \\ \relax
1:23 & \begin{tabularx}{0.7\textwidth}{X} 只要你們持守信仰，根基穩固，堅定不移，不致動搖，離開了你們從前所聽見的福音的盼望；這福音也是傳給天下一切被造之物的，我—保羅作了這福音的僕役。 \end{tabularx} \\ \\ \relax
1:24 & \begin{tabularx}{0.7\textwidth}{X} 現在我為你們受苦，倒很快樂；並且為基督的身體，就是為教會，我要在自己的肉身上補滿基督未盡的苦難。 \end{tabularx} \\ \\ \relax
1:25 & \begin{tabularx}{0.7\textwidth}{X} 我照神為你們所賜我的職分作了教會的僕役，要把神的道傳得完滿； \end{tabularx} \\ \\ \relax
1:26 & \begin{tabularx}{0.7\textwidth}{X} 這道就是歷世歷代所隱藏的奧祕，但如今向他的聖徒顯明了。 \end{tabularx} \\ \\ \relax
1:27 & \begin{tabularx}{0.7\textwidth}{X} 神要讓他們知道，這奧祕在外邦人中有何等豐盛的榮耀；就是基督在你們心裡成了得榮耀的盼望。 \end{tabularx} \\ \\ \relax
1:28 & \begin{tabularx}{0.7\textwidth}{X} 我們傳揚他，是用諸般的智慧，勸戒各人，教導各人，要把各人在基督裡完完全全地獻上。 \end{tabularx} \\ \\ \relax
1:29 & \begin{tabularx}{0.7\textwidth}{X} 我也為此勞苦，照著他在我裡面運用的大能盡心竭力。 \end{tabularx} \\ \\
[1ex]
\hline
\hline
\end{longtable}
$^{1}$我們一起同心祈禱.
再次多謝主.
我們打開聖經.
盼望這段經文.
成為一個很重要的提醒.
我們在疫情復常之後.
教會是有能有力.
這樣來見證福音的真理.
將人帶進主耶穌救恩的裡面.
我們這樣的禱告.
奉耶穌基督的名.
全力以赴.
是今天講題的一個重點.
什麼是福音.
我們才知道.
如何做一個福音的執事.
在保羅.
如果收集範圍.
24到29節.
是他個人.
對自己的存在.
服侍的一個自述.
講他自己如何去領受這個職份.
和他如何去盡心.
盡他的能力去講明白這個奧秘.
就是耶穌基督.
這是24到29節的自述.
在這個段落的上下文.
就有一個很重要的鋪排.
由一章保羅講述.
尤其是第15節開始.
當然我們可以從21節開始看.
這一段聖經.
告訴我們.
首先他要告訴.
哥羅西教會的弟兄姊妹.
你們所領受的福音.
帶領你們去到.
什麼屬靈的階段.
在神的裡面.

$^{41}$你們是怎樣的一個光景.
很重要.
不是任何人可以做到.
只是神自己的救贖工作.
留意21到23節.
你們從前是與神隔絕.
心思上與他為敵.
所以這個隔絕.
是在日常的行為.
生活價值觀.
我們思考的方式.
我們做人的處事.
是與神為敵.
行為上的趨近邪惡.
22節.
但如今.
他直著他兒子肉身的死.
即是耶穌真的來過這個世界.
替每一個民族的人死了.
已經使你們與他.
他就是神自己和好.
這是第一個很重要.
主耶穌的死帶來的救贖功效.
就是與神不再隔絕.
不再生疏.
與神和好了.
神看我們是一份友好的關係.
在耶穌的救恩寶血裡.
留意22節.
這裡第一次講奉獻.
是神獻上我們每一個人.
去到他自己面前.
把你們獻在他的面前.
保羅有這樣對人生的透視.
原來每一個信了主耶穌.
由你信那一刻.
你就是被呈獻到神的面前.
是怎樣的狀態.
在神的眼中是怎樣的.
成為聖潔沒有瑕疵.

$^{81}$無可指責.
這是第一個的獻.
留意23,24節.
23節.
只要你們持守信仰.
根基穩固.
堅定不移.
不致動搖.
離開了你們從前.
所聽見福音的盼望.
意思是告訴我們.
這個福音是傳給我們.
用不同的形容.
要穩固.
要不要動搖.
要不移.
即是說有一個危機在隱藏著.
保羅在第二章.
由第七節開始.
他就會講.
開始有異端的毒素.
進入哥羅西教會.
所以他要來到.
首先告訴哥羅西教會的弟兄姊妹.
其實你們已經完備了.
你們在神的眼中.
是已經被呈現的.
一個聖潔的制.
當然人在有生之年要做的.
就是對這個信仰.
對你所獲得的救恩.
主耶穌.
你要有充分的傳遞的認識.
和你在你的生活的實踐裡.
去體驗救恩帶給我們的生命的潔淨.
這個是相輔相成的.
不是我們做很多的事.
去加增我們在救恩裡的成果.
成果是耶穌為我們成就的.
這一段就是一個很重要的開始.

$^{121}$所以保羅就以這個人生的責任.
來成為他自己事奉的一個治術.
就是24到29節.
接著第二章開始就是入正題.
他來到去勉勵警告哥羅西教會.
不要偏離救恩的元素.
不要添加一些副產品.
甚至乎有害的物質.
第二章第六節開始.
你們既然接受了主基督耶穌.
就當遵他而行.
在他裡面生根建造.
信心堅固.
就是要一直有一個.
由初信到成長的成熟的生命.
是以信心和耶穌結連的一個穩固的基礎.
在這裡長大.
以至於你在地上讓人看到的生命.
是基於內心的信仰的穩固.
正如你們所領的教訓.
留意第八節開始.
留意這裡.
正正式式的提出一個警告.
你們要謹慎.
恐怕有人用他的理學和虛空的妄言.
即是一些傳說.
即是一些假的道理.
甚至是似是而非的道理.
例如近兩年比較普遍的電話騙案.
你猜一下我是誰.
等等諸如此類.
有學問的沒有學問的都容易受騙.
有一個朋友.
他只是根據WhatsApp的短訊.
按一下按鈕.
立即就發現信用卡已經被人盜取.
太多這些令我們在地上不斷要小心防範.
但是保羅這位使徒緊張的是.
心靈上對耶穌救恩真理的防範.
要怎樣保證我們每一個信徒對神的真理認識.

$^{161}$和在我們實踐過程.
是不會偏離了主給我們的教訓和成長呢.
他說這些人就是不照著基督.
乃照人間的遺傳和世上的小學.
就將你們老去了.
這個老去不是捉你去柬埔寨.
是使你偏離了你人生應該成長的屬靈的階梯.
捉去柬埔寨都容易分辨到.
而是我們自己可能嘗試.
去將本來的神聖的生命成長.
那些步驟.
那些必須我們自己用神學去解脫去詮釋.
導致我們成長的生命.
未能成為一個呈現的祭.
因為去到一章.
剛才我們沒有看的.
就是24到29節.
就是保羅這個自述.
第28節這樣說.
我們傳揚他.
是用諸般的智慧.
勸誡各人.
教導各人.
要把各人在基督裡完完全全的獻上.
第二次提獻情.
奉獻.
這是祭司獻祭的術語.
保羅都看自己是一個呈獻者.
將每一個信徒.
這裡提到三次各人.
不是侵罪的意思.
是強調在保羅的內心.
宣教達到萬民.
每一個民族都在他的心裡.
他不會只是偏向猶太人.
過去他的確是這樣.
但現在在他的內心.
救恩讓他有一個新的.
對人觀的認識.
普天之下在神眼中都是罪人.

$^{201}$是沒有辦法靠自己的.
但是在耶穌裡面.
成為聖傑.
來呈獻給神.
已經是一個很清楚.
保羅看到.
驗證過.
他自己就是這樣的一個過來人.
他自己都是可以.
靠著耶穌.
而得著這份聖潔.
得著這份與神真正的和好.
所以他現在一生的努力.
就是要將人.
這裡是說要很完備.
完完全全要成長成熟.
這樣來獻上.
剛才早堂因為很多老人家.
可能有一點點共鳴感.
果然真的有個姊妹說.
你小時候有沒有試過.
我剛才說的故事.
就是我自己一個成長.
就是我家開農場.
有些人都聽過.
我在相當多的一段青年期.
18歲前後.
相當長的時間要幫忙家裡.
來養豬.
原來那個過程.
現在回想幾十年前.
40年前都有了.
再回想的時候.
有些東西都可以共鳴.
跟牧羊都有關.
首先就是.
我們去到冬天之前.
必須要幫每一隻豬.
打疫苗.
就算沒有病毒出現.

$^{241}$如果你不打.
就很難切.
不是那年最新的病毒資料.
因為那些疫苗.
要在藥房買.
會有新的研發.
你要打不斷要更新變種等等.
如果不打.
可以很嚴重.
去到冬天.
你不知道什麼原因.
後來發現原來那些候鳥.
飛下來.
不知不覺帶著病毒.
去到農場.
又不知怎樣.
就進入豬欄裡面.
那就很嚴重.
我試過親眼看著.
很健康的.
很重的大豬.
差不多200斤重.
一感染了病毒.
通常是口蹄腳蹄.
不出三日.
它會奄奄一息.
由一隻本來200斤重的.
突然間變成只有幾十斤.
是三日內.
接著它的情況嚴重.
就會死了.
是大規模地死.
是很恐怖.
所以我們很緊張.
冬天我們要做好防禦措施.
例如打疫苗.
如果有陌生人來農場.
來看.
我們就要做消毒.
要噴很濃的消毒水.

$^{281}$開到很濃.
如果車經過.
就要有水灘經過.
要做好這些預防.
但有時都很難避免.
這個過程.
我自己嘗試再回想.
我試過好幾次.
有些豬已經是奄奄一息.
但你想想.
你要將牠分開.
去到另外的地方.
就不是正常地餵養.
每一次牠們需要食物的時候.
因為牠已經不想.
和沒有能力食.
我看著我哥哥.
進去豬欄裡面.
抱著牠.
拿了一些暖和的食物.
給牠們吃.
幫牠塗牠的口.
好像當牠是嬰兒一樣.
200斤重.
但已經很瘦.
原來這些過程.
我現在回想.
和保羅講的.
這個自述.
其實是很重要.
保羅看到了什麼.
是人生最重要的成長.
剛才提到兩次.
28節是獻上給神.
22節是神.
將每一個信耶穌.
在耶穌救恩裡面的弟兄姊妹.
他就帶到自己的面前.
即是說已經不看.
這是他的敵對.

$^{321}$但28節.
保羅參與奉獻的工程.
看到每一個民族.
每一個人.
最需要的是福音.
最需要的是耶穌.
所以在這裡.
他來到講這些.
這麼重責備的話.
你們要謹慎.
恐怕有人會迷惑你們.
是不照著基督的.
照著的是人間的遺傳.
世上的小學.
就會將你們擄去.
擄去即是說.
剛才已經講過.
你不能夠有健康的成長.
所以我們要檢視一下.
有沒有屬靈的病毒.
進入了我們的心靈.
以至你根本沒有興趣去追求成長.
以至你根本沒有想過.
基督教期望的是你整個生命成長.
成熟.
對神的真理.
保羅講的是一個奧秘.
就是基督這個奧秘.
歷世歷代隱藏的奧秘.
就是基督.
過去這兩個星期.
城中的熱話.
一定是那個無頭公案.
大家很想去追查.
聽查看查.
是甚麼事.
越揭就越多.
好像我們很有興趣對奧秘.
但基督的奧秘.
我們有沒有興趣去研讀.

$^{361}$來去認識.
和我們來去真的有一個決心要成長.
上兩個月我有機會.
參加我們新界西宣道會的.
童工.
堂主任的一個祈禱會.
我有一個神學院的學兄.
他在屯門服事.
他說他今年.
第20年事奉.
他有新的神給他的領袖.
他放下事奉的工作.
他進神學院進修一些課程.
準備他走下一條人生的路.
過下一個如果還有20年.
他就下半場繼續在神的國度裡面服事.
他換一個工場.
我想下要很大的信心.
要很大的毅力.
這樣來將自己事奉了20年.
毅然的改變這個跑道.
他很清楚神在他生命裡面的帶領和呼召.
也讓我自己有一個反思.
有一個思想.
我今年事奉35年了.
由1988年7月.
突然間聽到.
其實當時去下令會.
只是信主半年.
去到建道神學院參加下令營.
第一晚就聽到有聲音說.
「銳貴 你為我做傳道」.
晚上就將這個message.
因為分組.
聽完就分組.
導師聽到之後就很緊張.
馬上潑我冷水.
不要那麼緊張 不要那麼激動.
你先等一下.
怎麼會這麼離奇呢.

$^{401}$一路按著按著.
自己一路都不應該那麼倉促.
自己一路繼續工作等等.
去到1992年.
就正正式式放下工作.
全時間的裝備和學習.
如是者走到今天2023年.
但是相對於保羅.
他對於這個積分的敏銳.
那份的擺上.
我們可以進入24到29節.
看看這位使徒.
其實他不認識哥羅西教會.
是因為在坐牢的時候.
在羅馬的監獄裡面.
認識一個叫做爾巴弗.
這位信徒.
第12節.
有你們那類的人.
四章的12節.
有你們那類的人.
作基督耶穌僕人的爾巴弗.
問你們安.
他在禱告之間.
常為你們竭力的祈求.
願你們在神一切的旨意上.
得以完全信心充足.
能夠站立得穩.
他為你們和老弟家.
並希拉波立的弟兄多多的勞苦.
這是我可以給他作見證的.
保羅特別提到爾巴弗.
是哥羅西教會的弟兄姊妹.
現在都因為一些信仰的原因.
在羅馬坐牢.
認識了保羅.
保羅也知道.
他所描述的哥羅西教會.
現在的情況.
保羅就憂心了.

$^{441}$保羅也知道.
裡面蘊含著.
一些添加了的元素.
是會將人.
如果接受這些的.
例如敬拜天使.
例如受一些猶太的節旗.
例如你覺得自己.
不是一個生下來的猶太人.
你不能夠受割禮.
很多這樣的教導說法.
就會混淆了一個外邦信徒.
對主耶穌救恩的原備.
這個添加就帶來危險.
我們就有責任.
作為每一個時代的領袖.
牧者.
我們就有責任.
傳遞一個原備的救恩.
和不斷灌輸給弟兄姊妹知道.
你自己的身心靈.
要裝備對耶穌的認識.
是真知確信的認識.
而這個知識是帶來你的信心穩固.
然後在你的實踐裡面.
經過不斷的鍛鍊.
而來到成長.
最終是成熟的獻情給神.
例如用剛才那個例子.
我自己人生的成長.
我們養豬最重要的就是催谷牠.
谷到牠能夠足夠的重.
就是要出欄.
要來去賣.
最重要的了.
就是這樣獻情的了.
所以差不多到豬商.
差不多要來買的時候.
大概兩個星期.
那時候是私人屠宰的.

$^{481}$所以通常是半夜才有出豬的.
我們就在那兩個星期.
知道差不多要出欄.
就用最好的飼料去催谷那些豬.
不斷地鼓.
鼓到牠們很突然的脹.
但是那些瘦肉是不漂亮的.
因為是鼓出來的.
勉強地把牠們鼓大.
所以這些就是那些技巧.
大家都知道的了.
所以就變成了.
但是在保羅的事奉牧養裡面.
期望不是這樣的.
期望每一個弟兄姊妹知道.
由你信主在耶穌的救恩裡面.
你已經在神的眼中是完備的.
但是這個完備就要在我們有生之年.
我們要具備一個很強烈的心智.
為主耶穌的福音來學習成長.
最終是我們在不同的崗位.
領受神給我們的恩賜來服侍.
所以留意第24節.
25節和28節.
對不起29節.
現在我保羅為你們受苦.
賭博歡樂.
並且為基督的身體.
就是為教會.
要在我的肉身上.
補滿基督患難的缺陷.
你可以說這個就是他事奉的.
很重的內心價值觀.
快樂地為主的教會而受苦.
不會有任何怨言.
我們有沒有這樣的弟兄姊妹.
有這樣的價值觀呢.
教會有沒有足夠的渠道.
教導你能夠有這樣的價值觀.
為主耶穌.

$^{521}$你是歡喜快樂的為別人擺上呢.
要在我的肉身上.
接著25節.
我召神為你們賜給我的積分.
作了教會的執事.
這個執事是做什麼呢.
責任就是將神的道理.
全避地去傳開.
全避地.
這個就是保羅看自己的積分.
不是他的才幹.
不是他的過去的背景.
不是他的成就.
不是他的其他東西.
不是他的學歷.
那些可以被神用.
但最根本他看到的.
這個積分是神賜給我的.
我是為教會.
我要做的就是在宗教會裡面.
將神的道理傳得全避.
這個道理是什麼呢.
就是這個隱藏在歷世歷代的奧秘.
現在就顯明了.
神願意任何一個人知道.
這個奧秘在外邦人當中.
有何等豐盛的榮耀.
這個奧秘就是基督.
在每一個人的心靈裡面.
成了有榮耀的盼望.
是榮耀的盼望.
弟兄姊妹.
我們有沒有一個很強烈的心.
希望改變很多人的命運呢.
希望讓很多的人.
有榮耀的盼望.
而這個榮耀.
是在神永恆裡面的盼望.
不是今生今世物質的盼望.
是在神裡面.

$^{561}$香港的教會.
到今天都沒有大規模的受苦逼迫.
這個時代會不會就是呢.
我們是否生在這個時代的信徒呢.
我們要怎樣傳遞基督呢.
我昨晚都講過一個例子.
我想都可以再講一次.
我們神學院的同學.
我們有群組.
他提到他的宗派.
現在正在服侍的宗派.
因為是美國的差會.
美國現在正在開會.
短期內可能會通過.
就是可以安立同性戀的傳道人為牧師.
如果美國通過了.
香港的宗派的教會就要follow了.
如果不follow會怎樣呢.
就要問.
很簡單.
香港的物業歸回聯會.
你就沒有了.
你就要離開那個宗派.
你就要自己去尋找那個服侍.
這個應該怎樣去面對呢.
弟兄姊妹.
這個是不是一個逼迫呢.
是不是要為信仰而付代價的受苦呢.
在保羅這個自述裡面.
他認為是.
因為已經將基督的真理扭曲了.
是不是我們有責任.
要捍衛聖經的創造論的重要的人觀呢.
是不是這個就是一個時代.
我們要學習來受苦呢.
沒錯.
不是每一個信徒.
不是每一個教會.
都有直接的肉身的逼迫苦難的.
但是在保羅現在坐牢的時候.

$^{601}$他歡樂地願意為教會.
基督的身體而來受苦.
他要講的是傳避的救恩.
主耶穌是這個隱藏的奧秘.
他要來講明白.
務求的就是每一個信徒.
都能夠成長成熟地.
最終是被呈獻.
好像我們養豬.
很開心就是看到他很肥肥白白.
拜拜了.
他要上路了.
當然好像很好笑.
其實都很不開心.
有些有感情的.
你就會有這樣的情況.
我們作為牧者.
我們要守住這個無添加的基督真理.
和信徒的生命成長.
我們有沒有自己來改自己的版本.
符合你自己的信仰成長呢.
早前我都提過.
作為牧者我們完全不會對一個指桑罵槐.
是點出一些問題.
希望弟兄姊妹自己去反思.
一點都不重的.
只是一個勸勉.
在我們的教會.
有些只是回團契.
飲飲食食.
不回崇拜.
不回教會.
不回聚會.
只是有團契.
是他自己很內圍的團契.
他就去了.
不回教會.
不回崇拜.
崇拜是最基本的.
聖餐是最基本的.

$^{641}$完全不去.
當然我沒有資料.
他不去我們去其他例外.
但是可能不是.
這就不是一個基督徒成長的生命.
即是說在他內心裡面.
有自己一個自己度身訂做的信仰成長的階梯.
這個嚴重.
即是說他達至不了成為一個完備被呈獻的祭.
牧者有責任.
勸說.
幫助.
了解.
和禱告.
不一而足.
亦都很坦誠.
受浸完的弟兄姊妹.
很多不再回主一學.
甚至有些已經不再回崇拜.
不在我們中間來聚會.
受浸完是做什麼呢.
受浸就是你真正認識自己與耶穌的關係.
的一個公開的認順見證.
整體的再成長的.
怎樣來承載你的生命.
你要聽明很多的真理.
你都不開放自己給自己一個空間去接受.
這是一個限制.
會不會有你自己訂制的基督信仰的成長.
而不是聖經講的呢.
這是很需要我們大家一起彼此的勉勵.
剛才提過24,25節.
保羅一直將自己放在與基督信仰的傳遞裡面.
他要盡心來勞苦.
留意第29節.
再一次他說.
我都為這件事來勞苦.
這件是什麼事呢.
就是完備地將基督在人的心裡傳揚.
照著他在我裡面的大能來盡心竭力.

$^{681}$原來保羅的盡心.
一方面是有盼望的.
因為是神的力.
聖靈在他裡面.
使他能夠盡他人生的力.
去將人帶進那個完備的裡面.
不過留意第28節.
三次講在這個自述裡面.
他自己 他的身份.
他領受的恩賜.
和他用的那個力.
不過留意第28節.
他加一個重述.
而這個重述是很特別的修辭.
我們傳揚他.
是用諸般的智慧.
然後提到三次各人.
勸誡 教導.
要將各人在基督裡面.
完完全全的奉獻到神的面前.
你可以想像到這樣的修辭結構.
有保羅的很重的心思在裡面.
我們不僅僅是坐牢的那幾個人.
是整個保羅的時代.
認識主耶穌的救恩的完備的每一個信徒.
都應該參與這個改變人類.
唯一的方法的這個命運.
就是每個人都參與.
在這個成長的服侍裡面.
然後就開始去講.
每個人按照著不同的時代.
諸般的智慧.
來勸誡 教導.
務求的是要將各人.
這個各人 相信在保羅寫這封信的時候.
他的內心已經沒有任何的民族的偏差 偏見.
和對某些民族有特別的歧視.
很清楚地每一個民族都應該.
都可以在救恩裡面達至那個完備.
你想想這個工程有多大呢.

$^{721}$我有一個同學 神學院的同學 同班的.
在西藏隱密地做宣教士.
連講都不准講的.
因為被知道是真的有生命危險的.
他的代禱信都是要很婉轉地寄回來.
這樣做的.
你可以想像到不同的民族.
你怎樣能夠帶他來得到這個完備的生命.
成為一個祭可以呈獻給神.
是需要我們每一個弟兄姊妹.
首先我們自己要有這個意識.
我們可以傳.
和已經給了我們去傳的這個完備的基督.
就是人心無基督 就只會有混亂.
世界無基督 只會帶來種種的殘暴.
有時我們祈禱會為烏克蘭 俄羅斯祈禱.
是東正教的地方.
大家是信神的.
為甚麼都要因為種族 因為各種的政治.
而最終帶來的很低層次的人與人之間的殺戮呢.
為甚麼會有這麼多諷刺的事情呢.
相信就是人心沒有真正的傳備的擁有基督.
而導致到這些情況出現.
個別的個人的.
你就可以想像到為錢可以做出一些簡直滅絕人性.
都不應該再看那些內容太深入的新聞.
聽下標題就已經想揭開它 不要看了.
很恐怖的事情 是沒有了基督.
不是沒有錢 是沒有了基督.
整個社會沒有了基督 就會有很嚴重的混亂.
我們如果能夠看到 這是一個保羅告訴我們.
他看到的大前提.
所以這個三次提「個人」是很寶貴的.
保羅再沒有任何對其他人的偏見.
教會人與人之間.
你能不能夠這樣學習呢.
可以不同意見 但不會有偏見.
不會有脈絡這樣的狀態.
有仇恨 耶穌說你要放下你的憲制.
首先與對方和好.

$^{761}$你想想基督徒就是應該這樣的層次成長成熟的.
作為牧者我們都見到很多.
見到你前面我馬上就暈路走.
你這個暈路走本來沒有問題的.
但你的心如果是因為見到那個人而暈路走.
你的態度動機就錯了.
你在神的眼中其實你已經是一個虛偽的人.
不知道在座的有沒有因為某些人.
你不喜歡回冷堂 不喜歡回祖堂 不喜歡回五堂.
你要特意避開.
如果是的 牧者勸你不要避開.
直接面對那件事.
直接面對那個人.
尋求彼此的復和.
內心能夠有那份真正的安然.
才是我們人生完備的呈獻的祭.
苦心自問我內心到今天.
沒有一個人是我內心我放不下而我恨的.
苦心自問是沒有的.
很感恩.
真的要這樣學習.
求主讓我們大家一起看到.
保羅這個自述.
28節是向每一個基督徒的挑戰.
你的人生有沒有將這件事.
擺為我們最重要的責任.
一起在復償之後為主努力.
我們每個人都有不同的崗位.
一定是神安排的.
一定是在神裡面.
祂會給你有足夠的大能.
讓你可以盡心竭力.
求主讓我們每一個都成為一個.
全力以赴的福音執事.
我們一起祈禱.
多謝主.
但願你的說話.
保羅的自述.
成為我們大家的激勵提醒和動力.
一起去追求.

$^{801}$連一個稱自己是罪魁的保羅.
過去以猶太人拉比來優越.
完全看不起耶穌和跟隨耶穌的人.
但是他現在能夠在他內心.
充滿了各人.
充滿了對每一個人的那份關愛.
我們獻上禱告求主幫助.
奉主耶穌基督的名.
阿們.
\newpage



\section{詩篇 46:1-11}
\label{sec:BXGw_xefZws}
\textbf{2023-03-05 加恩樓《上帝是我們患難中隨時的幫助》詩篇 46\_1-11 講員:鄧泰康傳道}
\newline
\newline
連結: \href{https://youtube.com/watch?v=BXGw_xefZws}{\texttt{https://youtube.com/watch?v=BXGw\_xefZws}} ~~~~ 語音日期: 2023-03-08
\newline
\newline
\hyperref[sec:M8De3gUONVA]{\small{< < < PREV SERMON < < <}}
~
\hyperref[sec:index_chronic]{\small{[返順時目]}}
~
\hyperref[sec:index_scriptual]{\small{[返順卷目]}}
~
\hyperref[sec:GHx9ZnTIMjk]{\small{> > > NEXT SERMON > > >}}
\newline
\newline
詩篇 46:1-11
\newline
\begin{longtable}{cl}
\hline
\hline
章節 & 經文 (和合本修訂版)\\
\hline
46:1 & \begin{tabularx}{0.7\textwidth}{X} 神是我們的避難所，是我們的力量，是我們在患難中隨時的幫助。 \end{tabularx} \\ \\ \relax
46:2 & \begin{tabularx}{0.7\textwidth}{X} 所以，地雖改變，山雖搖動到海心， \end{tabularx} \\ \\ \relax
46:3 & \begin{tabularx}{0.7\textwidth}{X} 其中的水雖澎湃翻騰，山雖因海漲而戰抖，我們也不害怕。（細拉） \end{tabularx} \\ \\ \relax
46:4 & \begin{tabularx}{0.7\textwidth}{X} 有一道河，這河的分汊使神的城歡喜，這城就是至高者居住的聖所。 \end{tabularx} \\ \\ \relax
46:5 & \begin{tabularx}{0.7\textwidth}{X} 神在其中，城必不動搖；到天一亮，神必幫助這城。 \end{tabularx} \\ \\ \relax
46:6 & \begin{tabularx}{0.7\textwidth}{X} 萬邦喧嚷，國度動搖；神出聲，地就熔化。 \end{tabularx} \\ \\ \relax
46:7 & \begin{tabularx}{0.7\textwidth}{X} 萬軍之耶和華與我們同在，雅各的神是我們的避難所！（細拉） \end{tabularx} \\ \\ \relax
46:8 & \begin{tabularx}{0.7\textwidth}{X} 你們來看耶和華的作為，看他使地怎樣荒涼。 \end{tabularx} \\ \\ \relax
46:9 & \begin{tabularx}{0.7\textwidth}{X} 他止息戰爭，直到地極；他折弓、斷槍，把戰車焚燒在火中。 \end{tabularx} \\ \\ \relax
46:10 & \begin{tabularx}{0.7\textwidth}{X} 你們要休息，要知道我是神！我必在列國中受尊崇，在全地也受尊崇。 \end{tabularx} \\ \\ \relax
46:11 & \begin{tabularx}{0.7\textwidth}{X} 萬軍之耶和華與我們同在，雅各的神是我們的避難所！ \end{tabularx} \\ \\
[1ex]
\hline
\hline
\end{longtable}
$^{1}$好,大家好,早安.
今天我們終於平等了.
我不用戴口罩,有些人也不用戴口罩.
今天跟大家分享這段經文的時候.
其實是很有種感覺.
在這段時間我們也一直.
我看到很多不只是我們,有很多教會也一樣.
在過了這個疫情.
我們叫做大家都開始,整個香港叫做取消口罩令之後.
之後大家就有一種感覺.
過去我們也要處理我們自己的一些感受.
我們一些壓力,或者其他的東西.
而放在監牢裡面的時候.
其實在過去日子裡面.
我們也不斷地經歷各樣不同的困難.
特別是我們在這段時間也一直在清打,不斷地禱告.
監獄接下來究竟是怎樣,我們也不知道.
接下來的時候談租約.
究竟最後的時候新鴻基是否繼續租給我們呢.
他會不會有一個很大幅的加租.
加到我們應付不到,說不定呢.
我們不知道.
我形容的就是在我們身邊.
總是,特別是香港人.
我想也沒有誰敢在我面前舉手.
我一點壓力也沒有的,有沒有.
應該不會的.
香港是一個很急速的地方.
香港是一個整個環境.
我明明想自命清高.
但總是身邊的人都會給你壓力.
我用一個很簡單的例子.
譬如昨天我們派餐.
我不知道大家有沒有留意我們派餐.
特別是在大堂那邊.
我們有一隻手拿著餅盤.
另一隻手拿著鉗子去夾那塊餅.
是不是已經很衛生.
之前疫情的時候我們也做到足.
兩隻手都戴上手套.

$^{41}$接著就拿著鉗子,戴上口罩.
我想是差一套防護衣.
到了昨天的時間.
口罩也不需要戴了.
不過不要緊,我們這麼近.
同一間店都要講租金.
我們戴上口罩吧.
我覺得不需要再戴手套.
就只戴這隻吧.
特別是拿杯子的手套也要.
昨天還可以.
今早我聽到同工.
因為我們也邀請了一些執事幫忙派.
昨天我們決定.
只戴一隻手就行.
因為那隻手是用來拿餅和杯子.
執事說不要,戴上吧.
戴上吧,會給弟兄姊妹一個安全.
我想強調這件事不是對錯的問題.
你明白我的意思嗎.
但問題是當我們見到.
更多人在說戴吧,戴吧,戴吧.
那怎麼辦.
那就戴吧.
你會發覺原來很多環境裡面.
不是你自己.
只是單獨你自己認為是怎樣.
而是我們身邊有很多不同的壓力.
有很多不同的人的不同聲音.
會令我們有這種壓力.
所以在今次的裡面.
我就會發覺.
當我們再去打開聖經去看的時候.
特別是在這些環境當中.
其實我們很鼓勵弟兄姊妹.
去看詩篇.
詩篇不是讓大家就這樣看.
不是平鋪直敘.
詩篇在當日的聖典當中.
有一群詩班.

$^{81}$他們是會唱出來的.
如果大家看過,記得一定要留意.
為何詩篇.
其實今天我們已經不唱了.
但他仍然用了很多的標記.
去告訴你是怎樣的.
其實我舉個例子.
就好像今天這篇詩篇一樣.
他會強調在某些位置.
特別是我想大家留意.
他裡面很澎湃的地方.
他著名是用女高音.
為何會這樣.
其實他很想.
你去參與這篇詩篇.
他會用一些方法.
令你一起進入.
在這篇詩篇的裡面.
而大家會看到.
有些經文,其實有兩幅圖畫.
一幅圖畫是甚麼.
是很混亂.
令你很害怕.
大家有沒有見過.
鄧紫棋唱標高音.
我想問大家,你記住.
他唱歌不是為了標高音.
而標高音,大家要搞清楚.
不是,不是顯示實力.
而是你會發覺.
原來去到某些位置的時候.
我需要有一個.
澎湃的唱法.
而是將你的情緒.
帶到最高點.
你有沒有留意.
將你帶到最高點.
所以你會發覺這段經文.
裡面的時候.
這個作者很刻意用一個女高音.

$^{121}$要將那些地方.
帶你進入裡面,是怎樣的呢.
他說.
所以地雖然改變.
這裡說的是大地震動.
即是地在震動.
接著的時間.
他說山須搖動到海深.
意思是.
這些山是崩塌到一個地步.
山泥傾瀉.
湧入海中.
你記住那一句話.
震,湧入海中.
其中令波浪翻騰.
他用了一個字,是咆哮.
即是轟.
即是震.
震完之後,衝來之後.
轟一聲的時間,之後怎樣.
就令到海上.
所以那些山因為海浪的拍打.
而去震動.
原來是一連串.
這麼厲害的圖畫.
大家有沒有試過在海邊看浪.
其實我不知道大家有沒有見過.
一些很誇張.
當然我不贊成大家.
爆號風球走去海邊.
千萬不要做這種事.
我鼓勵大家.
但如果你有機會去到一些.
國外,一些地方專門.
可能那裡是太平洋.
和大西洋,我不知道.
是打下來的,我們香港很少會見到.
但你會發覺那種.
那種是厲害的.
如果大家暫時都沒有想出國的話.

$^{161}$去一個地方.
去學咀.
那裡專門有一個地方.
有個地方是浪湧進來.
有些浪聲很厲害的.
去聽一下.
其實這個作者他就是想告訴你.
當外面的這個環境.
他借助大自然.
這個場景.
告訴你很害怕.
很震撼,令到人會感覺到.
很混亂.
但他用了一句說話.
很快就去做這個總結.
但是我們.
是不害怕.
很平淡.
但他就是說儘管在外面的環境.
是多麼混亂也好.
但是我們在上帝的裡面.
我們是不害怕.
不害怕什麼?.
是安靜.
對比就是這樣.
外面是平平白白.
但是作為我們知道有上帝的人.
我們在上帝的裡面.
我們能夠安靜.
弟兄姊妹我想告訴大家.
基督教裡面最大.
對於你今天最實用的價值.
就是這件事.
從來在聖經裡面.
都不會強調和高舉神蹟這件事.
你會發覺在聖經裡面.
所有的神蹟都是為了.
上帝的一些特別的目的而發生.
舉個例子.
我們經常說笑.

$^{201}$一去我們有機會.
喝宴會怎樣?.
整個神蹟隨便走.
但大家要記得.
隨便走這個神蹟背後.
有場景.
是想達到目的.
不是為了做而做.
你不會無緣無故.
想演戲給老公看.
老公我就將這裡.
隨便走給你看.
看我們的神蹟有多厲害.
你不會的.
但是弟兄姊妹.
當整個世界.
平靈白蘭的時候.
你卻是表現出來.
你讓人感覺到.
作為基督徒一個屬神的人.
我竟然可以這麼平靜.
我竟然可以這麼安靜.
這個就是我們信仰裡面.
最大的價值.
為什麼我這樣說呢?.
按照香港的統計.
香港是越來越多人去看情緒課.
很多抑鬱症.
很多鬱躁症.
很多原因就是.
因為我們不斷面對環境的時候.
大家努不努力.
大家很努力.
我想強調.
其實我所認識的大部分的人.
都是勇往直前.
你面對困難的時候.
他們都是會盡力.
但最近的時候.
有很多人.

$^{241}$特別是讀心理的輔導員.
都會跟你說.
不要再跟對方說.
加油.
不要再跟對方說加油.
因為他已經很努力.
很加油.
你還想我怎樣呢?.
他們需要的不是.
更加多的比例.
反而更加需要的是.
平靜.
這裡很清楚.
為什麼我們不害怕?.
因為有上帝在當中.
第四節就告訴你.
有一道河.
這道河就帶來.
將歡樂帶到上帝的城裡.
這個城是至高者的住所.
你就會發覺.
河是神帶來祝福的一個象徵.
在啟示錄22章第一節.
在當這個世界末後的時候.
有一道生命水的河.
光亮如水晶.
從神和陽高的寶座流出來.
原來在上帝的城裡.
屬於的地方有什麼呢?.
是有從上帝而來的祝福.
進入上帝的城裡.
我們今天來到這裡.
一起敬拜.
弟兄姊妹不是為了唱歌.
我們不是為了一齊唱一首.
好聽的詩歌.
甚至乎不是為了.
聽一個傳道人.
在這裡和大家說一番說話.
我們今天來到這裡的時候.

$^{281}$是去領受上帝而來的祝福.
詩歌的背後.
有上帝嗎?.
為什麼我們的詩班.
要花這麼多時間去練習.
我們很想將最好的.
呈現給大家.
然後在詩歌當中.
能夠找到上帝.
我們今天如果.
唱詩歌唱到「甩甩咳咳」.
我們唱詩歌.
唱到走音.
大家整個過程裡.
只會專心去想.
為什麼Michael會唱錯.
為什麼Michael會走音.
我們就不會在裡面找到上帝.
我們在裡面努力的目的.
是為了如此.
我很記得當日.
我去學港道的時候.
我們當然要在聖經裡面學裝備.
我的老師.
他當時是.
牧養藝人之家.
他們是一起有份.
去牧養很多參與演員的人.
他跟我們說.
他教我們去港道學.
他告訴我一個很重要的事.
你要記住每一次.
當你站上台的時候.
其實是一個你為上帝.
而去的演出.
為什麼他要用這個詞.
有些人未必接受這個詞.
但是原因是什麼.
我想問問弟兄姊妹.
大家喜不喜歡.

$^{321}$看子華臣.
我不知道.
有人喜歡他,有人不喜歡他.
但你不要以為.
特別在以前.
他在台上說的棟篤笑.
他不是就這樣說的.
是精心設計的.
是花了很多心機.
怎樣讓你在他的說話裡面.
除了笑之外.
想到一些東西.
他提醒我們.
我們今天除了要有神的話語之外.
你所去做的.
跟大家說的時候.
你不要讓台下的人睡著.
他說睡著了.
無論你的東西多精彩.
他都聽不到.
所以弟兄姊妹.
你會發覺.
原來我們是一起來到神的面前.
去領受從神而來.
他在這裡.
他的恩典給我們.
我們知道為什麼我們不懼怕.
因為裡面有上帝的祝福.
他說.
神在其中成就必不動搖.
這是一個很重要的信念.
只要有上帝.
什麼都不用怕.
然後他說.
道天一亮.
神必幫助者成.
道天一亮.
弟兄姊妹.
上帝的時間.
當今天你在一個難處困境的時候.

$^{361}$我們都很想.
我們都想快點過去.
快點出來.
上帝為何祈禱了這麼久.
你都不回應.
我們都會這樣給上帝.
但你要記住.
上帝的時間.
是上帝認為天光的時間.
你放心.
是一定會幫助這個成的.
所以他最後才說.
儘管外邦圈養一片混亂.
列國搖動.
他用的字很有趣.
儘管外面的土地瓦解也好.
他用了一個很有趣的.
只需要神一出聲.
地便融化.
到今天我仍然會看弟兄姊妹.
上帝留給大家.
就是這本聖經.
就是神的話.
神的話.
是帶著能力的.
所以他的說話.
說有就有.
命立就立.
然後他說.
融化這個字很有趣.
他翻譯出來又是平靜下來.
由很混亂的裡面.
然後平靜下來.
原來很有趣.
整段經文.
不斷跟你說.
在神裡面我們最重要的是.
安靜下來.
你用神的說話.
安靜下來.

$^{401}$而到了這段經文的第三段.
他就說.
你是來看看耶華的作為.
看看他屎地方涼.
很有趣.
原來這個翻譯是什麼呢.
看看他怎樣摧毀這個世界.
一聽到我想大家馬上想起.
就是全部不要.
原來這個摧毀世界不是這個意思.
再看下面.
他會結束全地所有的征戰.
接著.
將所有戰爭裡面的器材.
都毀滅.
這個世界不需要摧毀.
因為是屬於神的.
但這個世界需要摧毀的是一切的戰爭.
我經常說.
我們今天看到有些地方在打仗.
我們心裡面當然會非常不安.
弟弟兄姊妹.
今天在我們的生活當中.
都是在打仗.
沒有煙燒的戰爭.
有沒有聽過.
人與人之間的關係.
很緊張.
讀書被整個教育制度扯到.
你想他不溫書都不行.
他沒有辦法達標.
我們被整個生活環境裡面.
扯到我們沒有辦法不繼續下去.
你想停下來嗎.
所以有時候.
以前我聽過有些人說.
你比老闆早離開.
就是你有問題.
你要比老闆遲下班.
代表你勤力.

$^{441}$有很多很多.
其實你會發覺原來不斷在我們整個生命裡面.
我們都在打仗.
弟弟兄姊妹.
但是聖經裡面上帝在做什麼.
我要停止這一切.
我要將這一切所有在戰爭當中用的器材.
都將它全部銷毀.
我要將它回復在神的面前.
回復到我的面前.
所以第十節.
就是大家好中的那段經文.
你們要休息.
要知道我是神.
在天中被尊崇.
在遍地上也被尊崇.
在患難外面混亂的時候.
我們要休息這個字.
安靜.
我們要安靜下來.
認清祂才是上帝.
要敬拜這個上帝.
因為祂要被尊崇.
你不靜下來.
你不會看到這個上帝的偉大.
我們今天都很習慣去解決問題.
我們很習慣一發生一些事的時候.
我們就用很多的方法.
我經常都說我們真的很厲害.
我們真的很厲害.
我不知道大家知不知道.
Aka是一個其實我很急的人.
其實我走路很快的.
我走路很快的.
就是打衝鋒的.
我在上一間教會.
有第二次機會跟他熟悉.
他跟我說.
你講得可不可以慢一點.
他說你一卵嘴我們聽不及.

$^{481}$所以他送了一些什麼.
他送了一些貼紙給我.
是蝸牛.
他說麻煩你貼在螢幕上.
提醒我要慢一點.
我以前不覺得.
我以前會覺得.
剛剛爭取時間.
盡快做完.
但最近我跟大家分享.
我腳的問題.
你不能不慢.
我現在想快都快不來.
我想跑都跑不到.
在這段時間很有趣.
就是我繼續要我的事奉.
我都是要繼續我的事奉.
但上帝就要你重新再去思考.
究竟你的事奉是什麼.
當你身體的所有機能都慢了.
你的事奉是不是可以繼續呢.
在一個這麼混亂的環境.
當上帝透過這段經文.
跟我們做總結的時候.
我想提醒姐妹們.
就是一個安靜的功課.
我們今天其實透過靈修.
透過靜觀.
我們去尋求自己的內心.
其實你是需要去找回.
你自己是什麼.
你要找回你自己究竟在想什麼.
我預備的時候.
最近很流行AI.
我就問了.
AI一個問題.
我就問從心理的角度.
一般人是用什麼方法去面對壓力呢.
他給了幾點.
他說一般人面對壓力.

$^{521}$可以用以下的方法.
第一.
深呼吸和放鬆練習.
第二.
健身的運動.
有益身心的運動.
第三.
社交的互動.
跟你的朋友和家人交流.
提高你的情緒.
第四 調整你的思想方式.
改變一些負面的思維和語言.
第五.
創造放鬆的環境.
創造一個舒適的環境.
例如.
聽一下音樂.
泡一下腳.
看一下戲.
當然他很會戴毛.
他就說.
以上的方法一般都可以有效地幫助人.
面對壓力.
但如果你的情況沒有改善的話.
記得找專業人士輔導.
我看完之後.
我笑了出來.
連機械人.
他搜集.
外面.
不是信仰的環境.
他都提醒.
大家要什麼.
安靜.
靜下來.
慢下來.
所以弟兄姊妹.
原來靜下來.
我們都教一個很重要的.
傳統和遺產.

$^{561}$其實今天有很多.
外面不是我們信仰.
所以大家要小心.
不要一見到零修兩個字.
你就走過去.
外面有很多不同的信仰都是用零修.
用了我們這個字.
零性修煉.
他們用了這些字眼.
你會發覺原來連其他人.
都開始意識到這個重要.
但弟兄姊妹.
在我們的生活裡.
我們有沒有去.
靜下來.
讓自己好好地.
在我們安靜當中.
尋找上帝.
我看到今天很多人.
其實真的很怕靜下來.
當我們跟他們做一些零修練習的時候.
不要吧,會睡著的.
坐著不做事的時候.
總是不知道怎麼回事.
因為我們習慣了做事.
但弟兄姊妹.
原來這種安靜是需要習慣.
你要安靜.
要知道.
祂是神.
因為你要發覺.
這段經文其實有一句說話.
是不斷出現.
不斷出現.
原來是這個主題.
祂說.
神是我們的避難所.
出現了三次.
在第一節.
第七節,第十一節.

$^{601}$即是祂的開首.
小結和祂的總結.
祂都不斷地告訴你.
萬軍的耶和華.
是跟我們一起的.
祂是我們的避難所.
而避難所直接.
譯為寶女.
當我們整個世界.
大家有沒有一點點戰爭意識.
整個世界現在被轟炸機炸著這個地方.
你唯一可以做的是什麼.
大家知不知道.
防空洞.
記不記得.
為什麼會有防空洞.
就是當整個地方.
地戰式轟炸的時候.
你唯一可以做的.
是你躲進去防空洞.
你唯一可以做的就是那裡.
現在上帝告訴你.
我就是你的防空洞.
當你感覺到.
現在你外面.
像打仗一樣.
像大浪一樣的時候.
請你第一時間躲進上帝的防空洞.
請你找回上帝.
原來.
信仰這個信念.
是很重要的.
昨天我們.
跟一個團體.
我們去講的時候.
在看一本書.
那本書是關於禱告.
和神學的整合.
有一位叫曾金發牧師.
去到第20課的時候.

$^{641}$最後他給我們的提醒是什麼.
原來我們一個屬於神的群體.
最重要的是一個信念.
究竟禱告.
是不是.
你今天都同意.
是基督徒最重要.
我們今天的事.
如果你是有這個信念的時候.
你就會花盡一切方法.
改變一切.
你今天生活的所有東西.
而去達到這個目的.
就是說.
所有你今天認為很重要的東西.
都要讓路.
所以弟兄姊妹.
其實你會發覺.
這段詩篇就是想告訴大家.
我們今天面對於.
上帝的裡面的時候.
外面所發生的很多事情.
我們不是做一個脫離世界.
只是躲在山上.
什麼都不理的信徒.
我們都有我們的責任.
但請你首要回到上帝的房空洞當中.
找回上帝.
去確定上帝.
他在你生命的裡面.
在你生活當中的重要性.
你之後.
才去繼續走下一步.
所以求上帝.
幫助我們.
不要在一個越是忙碌的世代當中.
我們越是失落了.
這個這麼重要.
我們的位置.
所以我真的很鼓勵.

$^{681}$有些事情不要浪費.
大家有多久沒有拿一本書來看.
其實看書.
大家知不知道是一個很好的.
安靜的時間.
大家知不知道.
在我們教會的書架.
基督書樓的書架上.
有一本好書.
大家知不知道.
這本書就叫做.
跟自己.
身心靈對話的七堂課.
我們都忙到.
忙了休息.
好好休息.
大家知不知道作者是誰.
我經常都說.
我經常都掛在嘴邊.
我們的鑽石就在我們的後花園.
這本書的作者是.
使徒永富先生.
董愛莉姊妹.
還有他們兩位.
很熟悉的朋友.
杜淑貞小姐.
還有曾子怡姊妹.
我不是特意幫她買廣告.
我只是想.
讓她知道.
這本書的作者.
是我特意幫她買廣告.
是因為我認識了富哥.
大家都知道.
他什麼都與她無關.
這個世界上.
很多事情他都會幫忙.
一定是忙到要脫身.
但我想告訴大家.
正好是弟兄告訴大家.

$^{721}$在他的領袖當中.
他就是太忙了.
他會和你分享很多.
他自己心底的故事.
很想讓大家知道.
我們需要休息下來.
弟兄姊妹重拾.
不要讓你的書架.
這本書在風塵.
找些時間拿著這本書.
去看一看.
沖一杯好咖啡.
好茶放在自己前面.
好好地讓自己休息.
和自己對話.
好嗎?.
所以求上帝幫助我們.
願意我們在神裡面.
重拾我們本來.
應該有的半想.
求上帝讓我們.
心靈和神有更多的對話.
讓我們在越大的風浪當中.
我們越能夠安穩.
很盼望這個.
就是佳恩家.
我們一起同心祈禱.
上帝我們求你.
去保守和帶領我們.
最後.
當我們在身邊波濤洶湧.
我們看到很多事情.
我們都沒有辦法.
去知道結果.
我們知道很多事情.
都令我們心裡擔憂的時候.
你就告訴我們.
回到你們面前.
你要為我們心裡.
去止息這一切的戰爭.

$^{761}$你要去幫助我們.
重新歸到上帝.
來的當中.
去能夠在你裡面.
重新得力.
因為我們知道.
當你在裡面的時候.
就正如你透過詩篇去提醒我們.
我們在這個地方.
有你就足夠.
我們的心裡面.
有你就足夠.
主啊,我們求你幫助我們.
讓我們重新學習.
安靜的功課.
求你讓我們.
更加看見自己.
看見我們身邊人的需要.
將一切我們交託給你.
多謝你今天.
所給我們分享的時間.
祈禱奉教主耶穌,得成名字祈求.
阿們.
\newpage



\section{創世記 18:18-19}
\label{sec:GHx9ZnTIMjk}
\textbf{2023-03-11 晚堂《上帝子民要秉公行義》創世記18\_18-19 講員:黃頌言姑娘}
\newline
\newline
連結: \href{https://youtube.com/watch?v=GHx9ZnTIMjk}{\texttt{https://youtube.com/watch?v=GHx9ZnTIMjk}} ~~~~ 語音日期: 2023-03-14
\newline
\newline
\hyperref[sec:BXGw_xefZws]{\small{< < < PREV SERMON < < <}}
~
\hyperref[sec:index_chronic]{\small{[返順時目]}}
~
\hyperref[sec:index_scriptual]{\small{[返順卷目]}}
~
\hyperref[sec:0LbO5gq5ujg]{\small{> > > NEXT SERMON > > >}}
\newline
\newline
創世記 18:18-19
\newline
\begin{longtable}{cl}
\hline
\hline
章節 & 經文 (和合本修訂版)\\
\hline
18:18 & \begin{tabularx}{0.7\textwidth}{X} 亞伯拉罕必要成為強大的國；地上的萬國都必因他得福。 \end{tabularx} \\ \\ \relax
18:19 & \begin{tabularx}{0.7\textwidth}{X} 我揀選他，為要叫他命令他的子孫和後代家屬遵行耶和華的道，秉公行義，使耶和華所應許亞伯拉罕的話都實現了。」 \end{tabularx} \\ \\
[1ex]
\hline
\hline
\end{longtable}
$^{1}$各位弟兄姊妹晚安.
我們一起祈禱吧.
親愛的天父 多謝你.
我們很渴望你的話語.
願你的話語來更新我們的心.
多謝你 奉耶穌名求 阿們.
今晚晚堂很開心.
我看到很多熟悉的面孔.
小朋友的面孔在當中.
竟然晚堂也有小朋友可以在當中.
這是一個恩典.
想跟大家分享一件事.
在我的大學時代.
大約十年前.
我跟同學討論過一個問題.
因為他說.
究竟我們為什麼要結婚.
為什麼要生小朋友.
他就說.
在這個時代裡.
生小朋友真的辛苦他們.
因為又難找工作.
又要供車又要供樓.
但是我從來沒有這個疑問.
我從來沒有想過這件事.
也沒有這個疑問.
因為我這樣回答他.
因為神要我們在地上生養眾多.
養育敬虔的後裔.
所以我沒有想過不生小朋友.
他就回應.
啊 是啊.
原來這次的對話.
改變了他的想法.
因為原來他一直不想生小朋友.
變成原來跟我聊天之後.
他有一個誘因可以生小朋友.
雖然他到現在還沒有小朋友.
因為想要也未必這麼容易會有.
不過上一代似乎沒有這麼多困惑.

$^{41}$因為我媽媽也有七個兄弟姐妹.
爸爸也有八個兄弟姐妹.
當時上一代經濟也未必很好.
但是他們也真的生養眾多.
我自己也有三個兄妹.
我們從小到大都在教會裡成長.
婆婆那一代已經信主了.
所以很感恩.
這位神不單只是我婆婆公公的神.
也不單只是我爸爸媽媽的神.
更加是我的神.
然而我們應該怎樣去看待我們的孩子呢.
以致我們可以一代接一代.
都行在神的心意裡.
現在教會也在說一個傳承的問題.
究竟要怎樣才可以得到神的祝福.
可以傳承下去呢.
今天我們一起去看這段經文.
在《創世記》18至19節.
這段是神直接對亞伯拉罕的說話和祝福.
這句話很短.
所以我們一定要看上下文.
就會更加明白背景.
故事是發生在羅亞方舟之後.
洪水毀滅大地.
因為人犯了罪.
洪水毀滅大地之後.
人建立了巴別塔.
然後神就將他們分散各地.
亞伯蘭 即是亞伯拉罕改名之前的亞伯蘭.
其實是羅亞的兒子閃的後代.
閃到亞伯蘭是10代.
即是相隔10代.
即是292年.
大概30年就一代.
60歲就要做公公婆婆.
我想現在也差不多.
到了亞伯蘭這一代.
他的妻子薩萊是不生育.
沒有小朋友.

$^{81}$亞伯蘭的爸爸塔拉.
其實也是70歲才生他三兄弟.
其實也很像生小朋友 70歲.
因此他爸爸有很長的生命.
205歲才死.
亞伯蘭就帶著他的侄子.
即是哥哥的兒子.
他的兒子叫做羅德.
他就帶著他的侄子一直相依為命.
後來其實神很多次祝福亞伯蘭.
一會兒我們就再看看神怎樣很多次祝福他.
來到這一次更加成為訪客.
去到訪他那裡.
亞伯蘭似乎都知道他是誰.
因為耶和華曾經多次向他說話.
所以他其實非常熟悉耶和華.
他看到這三位訪客的時候.
就知道兩位是神的使者.
其中一位是耶和華.
因此亞伯蘭很盡心思去服侍這三位訪客.
有宰牛 有預備食物.
跟獻祭的時候一樣非常用心.
這三位訪客有兩個目的.
第一個目的是預言明年撒拉就會生二撒.
亞伯蘭要有小朋友.
所以是宣告生命的生.
第二個目的是宣告死.
要審判所多瑪和.
所以在二十節我們看到耶和華說.
所多瑪和 阿摩拉的罪惡甚重.
聖文於我 我現在要下去察看他們所行的.
果然盡像那達到我耳中的聲音一樣嗎.
若是不然我也必知道.
在宣告生和死兩件事之間.
其實就是今天我們要看的經文.
這個已經是神第五次對亞伯拉罕的祝福.
是《創世記》十八章十八至十九節.
十七節其實講到耶和華說.
我所要做的事豈可瞞著亞伯拉罕呢.
我們一起讀第十八十九節.

$^{121}$亞伯拉罕要成為強大的國.
地上的萬國都必因他得福.
我眷顧他.
為要叫他吩咐他的眾子和他的眷屬.
遵守我的道 秉行公義.
使我所應許亞伯拉罕的話都成就了.
原來耶和華認識二人.
在第十九節說「我眷顧他」.
橙色的字「我眷顧他」.
在另一個女真中的譯本.
也有「我和他之交」.
也有其他的翻譯說.
「我選他 我認識他」.
用詞原來跟詩篇一章六節都有點像.
一章六節說.
因為耶和華知道二人的道路.
惡人的道路卻必滅亡.
原來二人的道路.
就是耶和華的道路.
所以耶和華認識他.
第十七節.
耶和華說.
我所要做的事豈可瞞著亞伯拉罕呢.
「瞞著」這個字也可以翻譯為「隱藏」.
原來耶和華不僅沒有向亞伯拉罕.
隱藏神在他身上的計劃.
甚至不停地提醒他.
亞伯拉罕認識耶和華.
多次向耶和華提出祝福和預言.
亞伯拉罕後裔的未來.
告訴他未來會成為大國.
在《創世記》十二章開始已經去祝福他.
他說「我要賜福給你.
叫你的名為大.
祝福你的 我會祝福他.
咒詛你的 我會咒詛他」.
紅色那個字.
「地上的萬族都要因你得福」.
我們很快再看.
他第二次的祝福.

$^{161}$在《創世記》十三章.
神去祝福他的後裔.
如同地上的塵沙般多.
第三次耶和華祝福他.
是在《創世記》十五節.
他說「其實耶和華」.
他不太明白耶和華的祝福是甚麼.
因為他說「我都沒有後代」.
「怎樣可以有這個祝福呢?」.
「應該是生在我家裡的工人 僕人」.
「我就當他是我的兒子」.
「我的東西就給他承繼」.
耶和華說「不是這樣的」.
「是出自你的先 是你的後代」.
「就叫他看看天空」.
「你的後裔會像星星般多」.
亞伯蘭就相信了耶和華.
耶和華就以此為他的兒子.
然後下格就生了以撒瑪利.
這個小朋友真的出自亞伯拉罕.
第四次耶和華再來祝福他.
是在《創世記》十七章.
耶和華要幫亞伯蘭改名做亞伯拉罕.
要跟他立約.
要幫他改名做多國之父.
亦預言撒拉會生以撒.
他說亞伯拉罕要在神面前作完全的人.
耶和華就要與他立約.
耶和華認識二人.
而且跟他們之交.
他們是彼此認識的.
亞伯拉罕知道神是信實的 是公義的.
並且相信他所應許的.
當亞伯拉罕未見到自己有子孫.
他已經相信.
甚至相信.
所以他後來願意將他的獨生子.
以撒獻上.
當神要他去獻以撒的時候.
他也是獻上.

$^{201}$因為他的神竟然可以將不生育的人.
在不生育的人身上創造生命.
即是在他的太太身上創造生命.
他也相信神一定會為他預備.
他所需要的一切.
神會親自成就這個應許.
直至他也認識神是公義的神.
所以他可以向神禱告.
亦為所多瑪去禱告.
耶和華認識二人.
認識每一個相信他的人.
我們也很喜歡如此認識我這一首歌.
因為我們也相信神很認識我們.
然而我們是否也認識神呢.
我們是否也知道聖經.
神的應許有多少呢.
我們有沒有盡力去讀聖經.
我們有沒有盡力去讀我們的聖經.
去盡力去認識這位神呢.
神究竟要賜多少的應許給我們呢.
原來我們抓著這些應許去祈禱的話.
我們的生命會很不一樣.
你有沒有研究過.
究竟聖經有多少個應許呢.
我上網搜尋過.
大概有八千多個應許.
聖經裡面也有另一個版本.
叫做應許版.
我曾經很喜歡看這個應許版.
因為太貴了.
所以我當時沒有買.
不過它會降低了.
聖經裡面所有神的應許.
讓我們可以更加知道神的應許.
有一本書叫做.
神的120個應許.
感恩有學者將神的應許分為13大類.
讓我們可以認識這位神更加多.
祂很認識我們每一個.
祂也很樂意賞賜我們.

$^{241}$然而祂賞賜的背後.
是需要我們持有信心.
和進行祂的教訓.
第一點要和大家分享的是.
義人為惡神去代求.
這三位訪客.
耶和華來找亞伯拉罕.
向亞伯拉罕透露了他的計劃.
其實是有雙重的目的.
首先耶和華要不生育的亞伯拉罕的家庭.
成為大國神.
去選擇他.
去認識他.
亞伯拉罕將要成為萬國的祝福.
所以當一個聖所多瑪.
他即將失去這個祝福的機會.
耶和華就要向亞伯拉罕交代.
第22節說.
義人轉身離開.
就是那三位訪客.
其中兩個使者.
他就向所多瑪去.
但亞伯拉罕仍舊站在耶和華面前.
亞伯拉罕近前來說.
無論善惡你都要殲滅嗎?.
假若城裡有五十個義人.
你還要殲滅那地方嗎?.
不為城裡這五十個義人.
饒恕其中的人嗎?.
將義人與惡人同殺.
將義人與惡人一同看待.
這斷不是你所行的.
參觀全地的主豈不行公義嗎?.
這裡提到義人與惡人.
原來憑著信心連繫於耶和華.
信服他的標準去行的.
這些就是義人.
亞伯拉罕就是因為他有信心.
所以被稱為義.
相反與立約無分.

$^{281}$無意信服耶和華的.
所以被稱為惡人.
萬國如何因為亞伯拉罕得福呢?.
亞伯拉罕率先為索多瑪代求.
他提出如果將義人與惡人一同毀滅.
是不公義的.
他大膽向耶和華求.
求他因著這五十個義人.
去饒恕其中的人.
後來他一求再求.
向神求到十個義人.
在十八章三十二至三十三節.
亞伯拉罕說.
求主不要動怒.
我再說一次.
假若在那裡見有十個呢?.
他說.
為這十個的緣故.
我也不毀滅那城.
耶和華與亞伯拉罕說完話就走了.
亞伯拉罕又回到自己的地方去了.
亞伯拉罕很紀念自己的侄兒羅德.
因為羅德在索多瑪城當中.
也相信那個城.
最少也有十個義人吧.
因此他求完就安心地回去了.
最後神的使者真的救了羅德一家.
然後再毀滅那個城.
因為真的十個義人都沒有.
但是因為亞伯拉罕的代求.
救了羅德一家.
他也以真正謙卑和敬畏神的態度去懇求.
原來一個義人的存在.
神也會去顧念.
當義人有信心.
傳遞公義的審判者.
不會將義人和惡人一同毀滅.
而且義人有能力保存社會.
不受毀滅.
所以不要小看自己的存在.

$^{321}$當我們一個社會變壞了.
義人的存在代求.
在神的眼中.
神一點也不稀罕.
神不是那些濫磨殘暴的上善罰惡支配者.
相反.
神積極尋求使我們全體可以脫離死亡的一條路.
神也不是獨裁者.
他會聽取義人的禱告.
讓人可以參與在當中.
甚至願意為義人的緣故去放過惡人.
因此義人更應該發揮他的作用.
在生活的團體當中去推行公義的美德.
如果你是社會當中僅餘的十個義人.
你會不會堅持自己的信念呢?.
許神也可以用那十個義人的生命去改變整個城市.
最近我們也看到.
有很多令人心痛的新聞.
其實這些惡事距離我們很近.
惡事就是發生在我們的身邊.
我們其實可以做些什麼呢?.
原來我們可以向神去代求.
求主的公義臨到.
求神去改變人心.
讓人知罪悔改.
讓執法者有智慧.
我們不要輕看自己的祈禱.
我們可以成為社會的守望者.
神會親自調動萬有去拯救困苦的人.
我也鼓勵大家去參與星期三的祈禱會.
因為每一個星期三晚上八點鐘.
我們都會為病患的肢體祈禱.
八點半開始我們就會為教會,社會,國家大使一起去祈禱.
或者其實神可以用十個義人的生命去改變這個城市.
只要我們是遵照神的旨意去生活.
並且照祂的旨意去選擇對的事.
第二點要和大家分享的是.
遵照神的旨意去生活.
並且去照祂的旨意去作出對的選擇.
耶和華向亞伯拉罕透露了他的計劃.

$^{361}$有雙重的目的.
首先是要萬國因為亞伯拉罕得福.
讓義人為惡城去代求.
第二原來是要亞伯拉罕教導他的子孫行公義.
照神的旨意去生活.
照他的意思去選擇對的事.
原來是成為了繼承應許的條件.
他要得到祝福.
原來是有條件就是行公義.
今天我們都因著基督.
因著信.
我們都成為了神的子民.
遵照神的旨意去生活.
其實都不容易的.
我們不應該只是一個星期日的信徒.
或者星期六晚的信徒.
我們所信的應該是我們平日的生活的細節裡面去活出來.
不單止要持守真理.
更加在生活裡面做一個對的選擇.
更加要教導兒女和我們的眷屬.
提醒我們的配偶要去行在合神心意裡面.
要做對的事.
不知道大家有沒有聽過蝴蝶效應呢.
前幾天有一個家長團契裡面.
有家長分享.
他說他受到別人的幫助.
他身邊有一個很好的基督徒去幫他.
他覺得很感動.
因此他都這樣去幫其他人.
當你用愛去對人的時候.
別人原來都會在你身上學習.
都會這樣去對其他人.
慢慢社會就可以改變了.
因此不要輕看自己每日堅持對別人好的舉動.
也不要小看小孩子的教育.
特別是遵從主的教導.
因為很多習慣都是由小的地方去建立.
記得我小時候第一次學的土文就是主土文.
每一晚我媽媽都會在床上.
跟我哥哥妹妹一起去禱告.

$^{401}$背主土文.
到現在我都很記得.
甚至當我還未懂得背主土文的時候.
我記得所有教會的小朋友都懂得背.
於是我就要求媽媽.
可不可以給我一些主土文卡.
一邊走一邊背主土文.
有很多字我又不懂.
但我都要逼自己背了它.
甚至我要求媽媽去解釋給我聽.
其實我不明白為什麼.
我都不明白那些字怎麼解釋.
不過媽媽真的解釋給我聽.
我都發現原來我對神的話很渴慕.
小朋友從小都會對神的話很渴慕.
小學之後我就建立了晚上祈禱的習慣.
因為媽媽以前都會晚上跟我們一起祈禱.
當媽媽要回自己房間睡.
我就去了妹妹的床代替媽媽.
我就捉著妹妹的手一起去祈禱.
原來父母為小朋友的祈禱都很重要.
你有沒有為自己的孩子祈禱.
或許你的孩子已經成人了.
或許他們已經有自己的家庭了.
但原來父母的提醒,代禱都很重要.
你的孩子孫兒你為他們的祈禱.
原來神都會聽到.
有朋友問我.
你為甚麼對神那麼有信心.
你的父母是不是經常為你祈禱.
這個問題我覺得很有趣.
但我回答他.
是的.
那為甚麼我會知道我媽媽為我祈禱呢.
因為小時候我靜靜的走進我媽媽的房間.
我看到書架上有三本簿.
分別寫了哥哥,我和妹妹的名字.
我很好奇就拿了我名字的那本書來看.
我一打開第一頁就流淚了.
因為其實我媽媽爸爸很少表達愛.

$^{441}$但我在這本簿裡面.
我很感受到媽媽的愛.
原來我媽媽向天父祈禱.
她寫信給天父告訴他我的情況.
我怎樣令到媽媽傷心.
或許給我感恩我的進步.
原來我媽媽全部都記下了.
還有為我的祈禱都記下了.
她這個習慣感染了我.
遺傳給我.
後來我從小學開始就這樣寫信給天父.
我到現在已經寫了很多本.
我到現在都在寫.
到今天媽媽依然會在她的家庭群組.
發送金句和長輩金句圖給我們.
不過我都會繼續懷著感恩和敬畏的心.
去接受這些提醒.
不要輕看你的身教.
首先你要成為一個遵守神的道.
秉公行義的人.
無論你的兒女有多少歲.
他們都會從你身上學習.
在社會當中.
秉公行義是要非常大的勇氣.
特別在我們的職場當中.
我之前聽到最多的引誘就是.
做帳戶的人.
他說老闆要我做假數的時候.
究竟要不要做呢.
會不會選擇說不呢.
在職場上.
或許我們也要踩低別人才可以升職加薪.
你又會選擇怎樣做呢.
我們要在社會當中去分別出來.
我們所做的事要別人看得出.
我們是跟從神的.
別人可不可以從你身上.
看見你是神的子民呢.
我們成為上帝的子民.
神要使用我們讓萬國得福.

$^{481}$神要怎樣使用你在你的位置上.
去彰顯祂的國呢.
神設立婚姻.
為了賜福給我們的後代子孫.
要我們遵照神的心意去養育他們.
以至於他們能夠成為經前的後代.
使全地都充滿了好像神的兒女.
所以生兒養女.
養這一班子孫可以成為神的子民.
一直世世代代都像神一樣.
最後我以一個故事去作結.
這個故事的背景.
是一個小學的班房.
大概是一個社會的縮影.
這個故事的書名叫做.
「不是我的錯」.
我們可以看見這個班房裡面.
有小朋友 最小的那個.
在哭泣.
有一班小朋友站在後面.
眼睛旁觀.
而且逐個逐個上前來.
告訴我們發生了甚麼事.
這個小朋友就說.
「那件事是下課後才發生的」.
「不關我的事」.
另一個小朋友就說.
「我都沒有看到那件事經過」.
「所以我不知道他為甚麼哭」.
「雖然我看到 也知道是甚麼事」.
「但是不是我的錯」.
「我很害怕 又幫不了忙」.
「所以只是站在一旁 就這樣看了」.
「有很多人都欺負他」.
「我一個人沒有辦法阻止」.
「那也不可以怪我」.
「有很多人打他」.
「其實所有人都打了他」.
「雖然我都打了」.
「不過我打一下而已」.

$^{521}$「不是我先打他」.
「是別人先打」.
「所以不是我的錯」.
「難道是我錯?」.
「我覺得他真的有點古怪」.
「這件事一點也不怪」.
「他會被欺負 也是怪他自己」.
「他一個人孤伶伶地站著哭」.
「男生這麼喜歡哭」.
「真的要收價」.
「雖然應該告訴老師」.
「但是我不敢」.
「更何況這件事也不關我的事」.
「他靜悄悄地在一旁哭」.
「甚麼都沒有說」.
「大家好像甚麼事都沒有發生過」.
「他甚麼都沒有說」.
「所以我們只是站在一旁看」.
「他自己應該大聲求救的」.
「雖然我也打了他」.
「但是我覺得沒甚麼」.
「因為所有人都有打他」.
「所以不可以怪我」.
「跟我有沒有關係呢?」.
「如果你也在這個班房」.
「你又會怎麼做?」.
「我們有沒有其他的選擇呢?」.
有些事我們會選擇沉默.
或者覺得事不關己.
但是神的子民在地上.
是要遵行主的教導.
特別我們跟別人是有點不同的.
而且需要有勇氣去表達和面對.
就像亞伯拉罕他也很有勇氣向神表達.
向神求公義和悔改的機會.
求神也讓我們有智慧.
可以做適合的選擇.
能夠在地上與神同工.
我們一起禱告吧.
親愛的天父.

$^{561}$你認識我們每一個人.
你讓我們在地上生養眾多.
成為蒙福的一群.
我們在地上是有使命的.
求你幫助我們有智慧.
將我們的孩子帶到你的面前.
求你讓我們有能力.
教導我們的子女子孫.
去走他們當走的路.
就是道路都不偏離.
也讓我們有愛別人的能力.
讓我們為我們的社會代求.
為萬國代求.
為我們的子女代求.
願我們世世代代都有愛主敬虔的後代.
奉主耶穌基督明求.
\newpage



\section{使徒行傳 9:26-30-11:22-26}
\label{sec:0LbO5gq5ujg}
\textbf{2023-03-12 午堂《生命師傅的典範》使徒行傳9\_26-30,11\_22-26 講員:陳一華牧師}
\newline
\newline
連結: \href{https://youtube.com/watch?v=0LbO5gq5ujg}{\texttt{https://youtube.com/watch?v=0LbO5gq5ujg}} ~~~~ 語音日期: 2023-03-14
\newline
\newline
\hyperref[sec:GHx9ZnTIMjk]{\small{< < < PREV SERMON < < <}}
~
\hyperref[sec:index_chronic]{\small{[返順時目]}}
~
\hyperref[sec:index_scriptual]{\small{[返順卷目]}}
~
\hyperref[sec:1tknD6q3Urw]{\small{> > > NEXT SERMON > > >}}
\newline
\newline
$^{1}$弟兄姊妹早晨.
我們可以回到教會親近神.
敬拜天父.
我們心裡面很感恩.
特別今早上聽到靈唱的弟兄姊妹帶我們唱一首詩歌.
我相信你們都很投入的.
那首就是《除你以外》.
這首詩歌告訴我們.
我們人是很軟弱的.
我們的肉體,我們的心腸.
有時會漸漸的衰敗.
但是上帝是我們的力量.
是我們持久的力量.
所以當我們唱這首詩歌的時候.
我們很得安慰.
也很得到鼓勵在當中.
是的,我們教會最近都多了一些肢體.
身體不舒服.
有牧者,有弟兄姊妹.
我們都很明白到.
我們人是有限制的.
有我們脆弱的地方.
但是感恩.
因為上帝是我們的幫助.
是我們的依靠.
是我們的力量.
所以今早崇拜的裡面.
很感受到這一份感受在當中.
跟大家分享一下.
今天跟大家分享一個題目較為特別一點.
就是生命師傅的典範.
上個世紀.
我相信有不少人都聽過一句話.
就是你學師不如做學師仔.
有沒有印象.
原來我們為了要有一技之長.
要能夠可以賺到錢.
所以上個世紀很流行.
你學甚麼都好.
就是學師仔.

$^{41}$跟著師父一起去學師,做事,生活.
來到這個世紀.
來到這個千禧的世紀的時候.
有另一個名詞出現了.
就叫做使徒仔.
Mentor 和 Mentee.
大家都有印象.
所以我們都很流行.
用這句話去形容使徒的關係.
無論你是在教育界,專業界.
或者在各行各業裡面.
你總是找到一些經驗之深.
或者很有豐富的學識的人.
做你的師父.
一個 Mentor 在當中.
如果你說講到我們的宗教.
我們的信仰各方面.
有一個名詞更加是很貼切去用的.
就是生命師傅.
今早我跟大家探討一下.
你和我有沒有想過.
我們去扮演生命師傅的角色呢?.
剛才主席也讀了一段頗長的經文.
你發覺經文裡面有兩個人物.
一個叫做巴拿巴.
一個叫做素羅.
後來就叫做保羅.
如果你提起這兩個人.
很快就會想起.
不用問阿貴.
馬上彈出來的就是素羅保羅.
因為保羅的知名度高.
多人認識他.
他的成就相當大.
大家都知道.
保羅不單止是一個報道家.
他建立很多的教會.
他也是一個神學家.
他的神學思想很豐富.
很淵博.

$^{81}$他寫的書信裡面揚入了.
他的神學思想的思維很豐富.
所以保羅他是神學家.
是報道家.
不單止.
原來保羅也是一個生命的師傅.
大家有沒有留意到?.
有沒有注意到?.
我提起人物.
可能你會有少少印象.
提一些比較三寶的.
譬如有沒有人聽過.
居拉和伯基拉有這個名字?.
史陀行傳記了這對夫婦的名字.
保羅也是生命的師傅.
有沒有聽過爾巴忽提?.
這個就更加陌生一點.
也是他的生命的師傅.
大家最熟的莫過於提摩太.
因為有提摩太前後書.
所以我們每個人都很認識提摩太的名字.
提摩太的生命師傅也是保羅.
大家可能做夢也沒有想到.
原來保羅這個素羅.
他的生命師傅是誰?.
他的師傅就是巴拿巴.
名不經傳.
不是很多人認識他.
但是竟然他可以做了保羅的生命師傅.
真是認真不簡單.
我想起有些power point.
現在可以給大家看看.
這個經文大家讀過了.
大家知道嗎?.
為何巴拿巴做了保羅的生命師傅呢?.
其實他留下了一個很重要的榜樣給我們學習.
因為巴拿巴的身上.
我們看到做生命師傅有三大特質在裡面.
這三大特質是甚麼呢?.
第一就是我們發覺巴拿巴是甚麼?.

$^{121}$他樂意去接待人.
他很樂意去接待素羅.
大家還記得當時候的素羅.
他信了耶穌之後.
好像你和我一樣很想回教會.
很想找一個基督徒的群體.
於是他毅然去到耶路撒冷的教會.
誰不知一進到這間教會.
他就吃了不少檸檬.
一進到這間教會.
沒有人理會他.
為甚麼呢?.
因為素羅的名字在耶路撒冷是非常出名的.
不過他在教會的圈子.
素羅的名字是甚麼呢?.
這個「朵」是非常不好的.
因為素羅曾經用他的計謀.
想盡辦法用大洗禮師的手令去捉拿基督徒.
大家有印象嗎?.
所以大家一聽到素羅這個名字都很害怕.
而且很抗拒.
不過現在聽到他信了耶穌.
是否真的呢?.
你明白我的意思嗎?.
大家可能會想.
他信耶穌是否在耍花樣呢?.
是否做臥底呢?.
你明白我的意思嗎?.
是否要在我們當中搜集更多證據呢?.
是否又要有第二次,第三次去捉拿我們呢?.
其實有這些思想的人是很正常的.
為甚麼呢?.
因為他們對素羅的背景太熟悉了.
太知道這個人是抗拒基督的.
與基督作敵的.
所以當時基督徒聽到素羅說信了耶穌.
還要來參加教會的團契,小組,崇拜.
每個人的心有甚麼呢?.
很害怕.
也很抗拒.

$^{161}$是否真的呢?.
這個時候對素羅來說.
其實他是遇到一些危機的.
有甚麼危機呢?.
大家看看便知道了.
因為他做夢也沒有想過.
原來當他信了耶穌之後.
竟然在教會裡沒有人跟進他.
沒有人去栽培他.
大家有沒有留意到?.
因為每個人都避開了他.
每個人都不敢靠近他.
每個人都很怕他另有目的.
但是一個初信耶穌的人.
好像你和我.
回到教會當然很想有人帶我們去主日學.
帶我們去初信班.
跟我們有個人的聊天,栽培.
我們都很期盼這些的.
但是現在剛信耶穌的素羅.
雖然靈明又幼嫩.
但是幾年間沒有人跟進栽培他.
你可以想像他的信仰的根基.
很快便失落了.
很容易放棄自己的信仰也不一定.
這是第一個危機.
第二個困難是被人抗拒.
跟進栽培已經不止了.
你知道已經不利了.
靈明沒有進步.
更加難受的是.
人們用奇異的眼光去看他.
人們看到他的時候.
都不知道他是真信還是假信.
對他有保留.
所以在這樣的情況下.
我想問大家.
這個素羅容不容易投入教會.
是很困難的.
為什麼呢.

$^{201}$因為這麼多人都抗拒他.
用很奇異的眼光去看他.
他都覺得自己是不是應該要回教會呢.
人們這樣看我.
這是造成了不利的環境.
第三個危機.
就是孤立無援.
他回到教會的時候.
沒有人理睬他.
沒有人跟他聊天.
沒有人接納他.
所以他回到教會.
好像隱形了.
在這裡沒有人看到你.
不在這裡更加看不到你.
隱形了.
沒有人當他存在.
你說慘不慘呢.
如果你回到教會.
回到隱形的.
當你沒有存在.
那就真的很笨.
不過幸好.
現在我們回到敢選家.
很多人都理睬你.
很多人都請你喝咖啡.
門口很多人跟你聊天.
很多人關心你.
回到教會應該是這樣.
但是他回到教會.
遭遇一連串的挫敗.
這樣的對待.
所以這個對素羅來說.
真的是一個非常不利的處境.
正當這樣的環境下.
上帝預備了一個很忠心的兒女.
就叫做巴拿巴.
大家留意到.
巴拿巴竟然是獨派眾議.
他不理其他人的看法.

$^{241}$但是對巴拿巴來說.
素羅是一個很寶貴的.
屬靈的果子.
他能夠信耶穌.
真的是一個很大的恩典.
一個天大的罪人.
現在信了耶穌.
你說珍不珍貴.
很珍貴.
就好像現在有些人成功.
福音解讀.
成為新造的人.
回到教會.
或者大家都聽過保叔的名字.
聽過保叔的名字嗎.
他是一個死囚.
現在不但信了耶穌.
還可以放出來.
繼續去事奉上帝.
見證上帝.
我們感恩宣家也有聽過保叔見證.
所以一個這麼大改變的人.
我們是不是應該很要擁抱他.
各位.
所以在這樣的情況下.
獨派眾議.
你們不理他.
但是我就覺得.
這是上帝給我們教會一個很大的祝福.
恩典.
他就去接待素羅.
各位.
當時候初期教會.
你接待一個人的時候.
你要記得.
不是跟他聊幾句話.
跟他有些活動.
就唯依接待.
不是的.
原來當巴拿巴接待他的時候.

$^{281}$他是甚麼.
跟素羅一起同食飯.
一起同住.
甚至招呼他住.
而且一起同行.
更加一同事奉.
大家有沒有留意到.
原來那個接待的意義.
是涵蓋了整個人的生活.
到甚麼.
都被接待.
所以這是給我們一個很好的啟發.
各位聽眾.
當我們教會有團契是很好的.
有小組是很好的.
但是你要記得.
我們要去栽培一個人的時候.
不是單單一次的活動.
一次的周會.
完了就完了.
如果你是這樣的情況下.
你發覺你的團契小組裡面的人.
他感受很少.
是甚麼.
是在屬靈上對他的偽陽.
大家有沒有留意到.
當然我們現在香港的環境.
也不方便跟人同食同住.
特別有疫情的情況.
所以對我們來說.
另一個意思就是說.
當我們要去跟進關懷.
栽培一個人的時候.
我們有沒有想過.
原來我們需要留意對他平日的關心.
他平日的需要.
我們有沒有留意到.
所以如果我們平時都有聯絡對方.
給他電話.
跟他聊天.

$^{321}$WhatsApp他.
跟他有些聯繫.
甚至有些活動.
我們一起出來吃飯.
一起出來聊天.
有交流.
有分享.
我記得在沙田祈禱會裡面.
我們除了每個星期有祈禱之外.
原來我們發覺.
我們平時的交流.
分享都是非常重要的.
所以我記得在疫情之前.
我們很常有的就是去家訪.
去家訪團友.
又或者我們會搞一些生日的聚會.
又或者我們發覺有人不舒服的時候.
我們會去醫院探病.
或者有其他活動的時候.
我們都一起來參與.
特別留意到弟兄姊妹.
平日在生活上他有什麼需要呢.
各位弟兄姊妹們.
原來一個生命的師父.
給我們學到一個榜樣.
對人的關懷.
對人的培育.
是需要一個全面性的.
整全性的來關心.
所以我們見到巴拿巴.
給我們一個很好的榜樣.
他很樂意去接待素羅.
特別當他初信主的時候.
他用很多的時間.
很多的精力.
來關心,培訓這位素羅.
這個讓我們見到.
作為一個生命師父.
我們是要付代價的.
大家有沒有留意到.

$^{361}$我們是要對人的關心.
不是一會兒就完結了.
而是我們要恆常地.
要持久地來關心.
好.
我們再看第二個特質是什麼呢.
作為一個生命的師父.
除了樂意接待之外.
其實他還需要以身作則.
他自己要展露一個好的榜樣.
給他的跟隨者.
所以巴拿巴自己本人.
就做了這件事.
為什麼.
因為有一次教會.
就差派巴拿巴.
要他去到另一個地方的教會.
來做一些特別工作.
所以巴拿巴就信服教會的差派.
那次的行動不是一批人.
好像短孫那樣.
十個八個.
十幾二十個過去.
不是的.
原來那次的行動是一個人.
One man band.
走去另一間教會.
代表母會來關心那間教會的弟兄姊妹.
所以當巴拿巴要去接受這個任命的時候.
他是很信服地被差派.
就算一個人要擔任這個任務.
他都很盡心盡力去做.
當他去到另一間教會的時候.
他就去關心.
去鼓勵其他的弟兄姊妹.
當他去到的時候.
也跟他們一起有彼此的勉勵.
分享,交流.
讓教會的弟兄姊妹.
感受到原來母會是很關心他們的.

$^{401}$原來母堂是很支持他們的.
所以還特別差派代表來到我們這裡.
跟我們有交流,分享,鼓勵.
各位弟兄姊妹們.
這個就是你看到巴拿巴.
當他被差派的時候.
他的態度是一種很積極的態度.
他很樂意去事奉.
他沒有問為什麼你只找我一個人去.
你有沒有留意到.
他沒有去質疑這次的行動.
是不是太倉促.
所以沒有人去.
就派了我去.
你知道我們平時也會這樣想.
沒有人做.
叫你做.
你很快就想.
當然沒有人做.
你叫我一個人去做.
你會發覺巴拿巴不是這樣的.
他是很積極,很樂意.
很正面去看這件事情.
果然他去到那間教會的時候.
能夠鼓勵那間教會的弟兄姊妹.
更加愛主,更加積極.
弟兄姊妹們.
其實有沒有想過.
當你信了耶穌之後.
你原來也領受上帝給你很多的恩典.
是不是.
上帝給你的祝福多不多.
很多.
再加上聖靈內在我們的心裡.
祂也賜給我們很多不同的恩賜.
叫我們可以去事奉上帝.
所以有沒有想過.
原來當我信了耶穌之後.
為什麼.
原來上帝不只是單單讓我決志信耶穌.

$^{441}$這麼簡單.
而且祂比我的生命更加豐盛.
比我的生命更加精彩.
我現在可以透過上帝給我的恩賜.
去做更多服侍的工作.
各位弟兄姊妹們.
你有沒有留意到.
原來當我們信耶穌之後.
我們努力積極的事奉.
其實也留下一個好的榜樣給人看到.
大家有沒有留意到.
當你很投入事奉的時候.
別人發覺你的生命當中.
是展現出那種屬靈的生命力.
我記得有一個姐妹.
她曾經這樣分享.
當我信了耶穌之後.
我覺得我的生活很充實.
為什麼.
因為我發覺我自己信耶穌之後.
很忙.
除了我要自己工作之外.
原來很感恩我可以在教會參加師班.
可以參加關故部的事奉.
可以去探訪.
所以她覺得自己信了耶穌之後.
她的時間是很珍貴.
很不夠用.
為什麼.
因為她發覺上帝給她很多事奉的機會.
所以我聽到這位姐妹分享的時候.
我為她很感恩.
因為她發覺自己信耶穌之後.
生活得很充實.
這個我們發現到.
原來你和我都可以好像巴拿巴一樣.
我們在事奉方面.
來流露出生命的見證.
不但我們可以讓人得到幫助.
甚至因為你和我的事奉.

$^{481}$更加可以榮耀上帝.
這個就是我們看到巴拿巴第二個做生命師傅.
他的特徵,特質的地方.
我們再看巴拿巴第三個特質是什麼.
原來他是主動去找一些人.
主動去找什麼人呢.
這個就告訴我們.
巴拿巴是主動去找素羅出來.
其實巴拿巴這個人都很有敏銳的觸覺.
他很看得到形勢.
他知道當時的形勢對素羅不妙.
為什麼呢.
因為素羅信了耶穌之後.
他很熱心,很積極.
很肯去講見證.
很肯來跟人分享.
結果那些猶太人很妒忌他.
你有沒有留意到.
甚至那些猶太教的人.
以前跟他的同行兄弟.
很不喜歡他.
覺得他叛變了猶太教.
現在歸順基督教.
所以這班舊底的兄弟們.
很生素羅的氣.
想要剷除他.
想要殺了他.
結果教會的弟兄姊妹就勸素羅.
不如這樣吧.
你現在避一避風頭.
因為現在風頭火勢.
你不要在我們教會.
你最好回鄉下.
一聽到回鄉下就淪盡了.
回鄉下避一避.
大家知道素羅的鄉下在哪裡嗎.
就是大數.
就是大數這個地方.
你回鄉下避一避.
遲些環境好些,穩定些.

$^{521}$你才再回來.
大家明不明白這個意思.
結果弟兄姊妹很熱心安排下.
將素羅送回大數自己家鄉的城市.
大家知不知道.
當素羅被人安排送回鄉下的時候.
會發生什麼事呢.
可能會發生以下這五件事.
第一件事.
就是素羅感覺到自己沒什麼用.
被人扔回鄉下.
感覺到很蹉跎歲月.
在那裡避一避.
真的不知道避多久.
大家有沒有留意到.
從此之後.
可能素羅消失在教會的圈子都未定.
為什麼呢.
因為他回到鄉下的地方.
沒有其他可以發揮的餘地.
所以你看到這個.
可能是第一個情況會出現的.
是打擊了素羅的心智.
你不需要在耶路撒冷的地方.
這是第一個情況.
第二個情況.
就是很容易被人忘記了.
大家認同嗎.
特別沒有出現一段時間.
在教會裡面容易被人忘記嗎.
很容易.
可能教會都覺得多一個又不多得幾多.
少一個素羅又不少得幾多.
就任由他消失在我們當中.
再加上他的底又不是很好.
所以這個對素羅來說是第二個問題.
第三個問題是什麼呢.
第三個問題就是素羅不出現.
反而對我們教會更加好.
你有沒有想過這個問題.

$^{561}$因為素羅這個人是野火人物.
他是風眼人物.
因為這麼多猶太教的人都要追殺他.
所以如果他不在我們的教會就更加好了.
少了很多問題.
少了很多煩惱.
多一事不如少一事.
所以你會發覺到素羅是更加不會被人.
接納他,歡迎他.
這個是第三個問題.
第四個問題是什麼呢.
第四個問題就是素羅也不再被人重視.
為什麼呢.
因為他已經在鄉下的地方.
不知是耕田還是種花.
總之在鄉下的地方沒有人再去賞識他.
沒有人再去找他.
所以日子久了之後.
再加上他剛剛出現在教會不久.
更加隨著日子的消逝.
沒有人再提起素羅這個人.
這個可能是第四個情況會出現.
他不再被重視.
最後第五個情況是什麼呢.
最大的打擊就是這一項.
就是第五項.
因為素羅的記錄不好.
底子很差.
既然他以前在教會有這麼多壞的聲譽.
不如藉這次機會讓他永不翻身.
大家明白我的意思嗎.
因為這是一個問題人物.
他去到哪裡都很惹火.
去到哪裡都有很多問題爭議.
所以這個人最好不要用他.
所謂一次不忠就八次不用.
讓他就這樣消失.
各位弟兄姊妹們.
我想提出這幾項事.
有可能真的發生在素羅的身上.

$^{601}$如果真的發生.
從此之後.
你和我都沒有機會在新約看到這麼多卷.
十三卷的書信.
是由保羅親筆寫成的.
你說對教會來說.
是不是一個很大的遺憾.
是不是 弟兄姊妹們.
有沒有想過.
也沒有想過將來初期的教會.
一間又一間被建立.
是因為素羅那種積極認真報道傳福音的心智.
這個人已經消聲匿跡.
這個人已經消失在我們的教會當中.
各位弟兄姊妹們.
你有沒有想過對教會的損害大不大.
是難以估計的.
正是最關鍵的時候.
當素羅好像被人遺忘的時候.
各位弟兄姊妹們請你記得.
我們永遠都在工作.
上帝祂掌管一切.
祂會在適當的時候.
做適當的工作.
差派適合的人選.
來完成祂的計劃.
誰去擔任這個角色呢.
巴拿巴.
上帝對他說.
你再去找素羅.
雖然他回到鄉下.
沒有人理他.
你再去找他.
巴拿巴有這種敏銳的觸覺.
所以他幾艱辛苦.
都要去大掃的地方.
再找回素羅出來.
再跟他江湖重出.
各位弟兄姊妹們.
這幅相就讓我們看到.

$^{641}$再找回素羅出來.
弟兄姊妹們.
有巴拿巴這位獨具慧眼.
被聖靈充滿大有信心的人物.
你知不知道這幾句說話.
形容巴拿巴是很不簡單的.
你有沒有留意到.
這三個的形容.
有一個發生在我們身上.
都相當不錯.
你明白我的意思嗎.
上帝就用這位巴拿巴.
不是很多人認識他.
但是他遠而被神使用.
結果他就成為了.
素羅,保羅的生命的師傅.
你說厲不厲害.
各位弟兄姊妹們.
將來對教會的貢獻是否很大.
是你做夢都沒有做夢得到.
原來一個這麼簡單的小人物.
既然上帝用他的時候.
後面出現一個什麼.
屬靈的巨人.
就是素羅.
弟兄姊妹們.
我想問你們當中.
大家有沒有聽過李思敬博士這個名字.
聽過他舉手.
他說全部都舉手.
你知不知道李思敬博士的背景.
李思敬博士的爸爸是哪一位.
就是李飛吾牧師.
這個是上個世紀的人物.
我想李飛吾牧師的名字.
你們很少聽到.
很少人聽過李飛吾牧師.
宣道會的人應該聽過他.
因為他是很被神使用的.
在宣道會拍著.

$^{681}$藤勻輝牧師.
好像走馬步一樣.
早期的教會.
怎樣來開很多的分堂.
一位就是藤牧師.
另一位就是李飛吾牧師.
他是希伯倫堂的堂主任.
為什麼我今天要提他們的名字呢.
弟兄姊妹們.
你知不知道李飛吾牧師.
實整他的背景.
是非常淒涼的.
因為李飛吾牧師.
他是一個孤兒.
當他在廣西的地方做孤兒的時候.
上帝就預備了一位.
沒有什麼人認識的宣教士.
他的名字就叫何爾斯宣教士.
有沒有聽過.
每個人都沒有反應.
沒有反應就是沒有聽過.
何爾斯宣教士.
這位西教士.
不知道為什麼.
上帝感動他.
就把這個孤兒撿回來.
在家裡養大這個孤兒.
何爾斯宣教士.
又把很多其他的孤兒撿回來.
因為當時中國是國亂的時候.
很多孤兒.
何爾斯宣教士就把孤兒撿回家之後.
很好的.
不單止是來養育他們.
更加是訓練他們.
培訓他們.
結果就出現了一位李飛吾牧師.
上帝使用的一位.
很重要的人物.
李飛吾牧師.

$^{721}$由李飛吾牧師再出的後人.
就是李思敬博士.
各位聽眾.
有沒有想過.
上帝的工作就是這麼奇妙.
你說神啊.
我是很微小的.
我在感受堂都沒有什麼人認識我的.
我上教會日子不是很久.
我恩賜不是很多.
我口才不是很厲害.
聽眾.
不需要擔心這一切.
我今早想告訴大家.
我想鼓勵你們.
到學校巴拿巴的榜樣.
做別人的生命的師父.
可不可以.
OK嗎.
你在教會裡.
找一個人.
找兩個人.
你發覺他是初信主的.
你發覺好像沒什麼人關心.
你發覺好像他來到教會好像很陌生.
你就不妨學一下巴拿巴.
主動一點.
找一下他.
跟他聊聊天.
關心他.
跟他祈禱.
跟他查經.
跟他有些靈命上的慰養.
頂智滿滿.
你不要輕看你所做的工作.
分分鐘.
你培訓出來的人.
可能是一個保羅的人物.
可能是李思敬博士的人物.
是你完全想像不到的.

$^{761}$上帝是很奇妙.
只要你願意被神使用.
你說神啊.
我在這裡求你容我.
讓我從今天開始留意一下.
在教會當中有沒有哪些弟兄姊妹.
或者在我的職場裡.
在我的公司裡.
有沒有同事是初信耶穌.
沒有人栽培跟進的.
我去關心一下他.
各位頂智滿滿.
如果你願意這麼做的時候.
今天早上這篇道.
就能夠進入你的內心.
成為上帝要使用你的一個印證.
讓你能夠祝福其他的人.
成為別人生命的師父.
我們低頭祈禱.
主啊.
因為你拯救我們之後.
成為你的兒女.
是蒙福的.
是蒙恩的.
但是你想我們的恩典.
得到的福氣.
是更加傳承到其他有需要的人身上.
今天早上多謝天父.
透過巴拿巴.
留下美好的典範.
榜樣.
讓我們看到.
原來上帝你的奇妙的作為.
令我們讚歎.
令我們驚訝.
就讓巴拿巴的榜樣.
今天激勵我們.
都立志成為別人生命的師父.
求主使用我們.
祈禱奉耶穌名求.

$^{801}$阿們.
\newpage



\section{創世記 18:16-21}
\label{sec:1tknD6q3Urw}
\textbf{2023-03-12 加恩樓《守主道‧行公義》創世記18\_16-21 講員:陳麗蟬牧師}
\newline
\newline
連結: \href{https://youtube.com/watch?v=1tknD6q3Urw}{\texttt{https://youtube.com/watch?v=1tknD6q3Urw}} ~~~~ 語音日期: 2023-03-15
\newline
\newline
\hyperref[sec:0LbO5gq5ujg]{\small{< < < PREV SERMON < < <}}
~
\hyperref[sec:index_chronic]{\small{[返順時目]}}
~
\hyperref[sec:index_scriptual]{\small{[返順卷目]}}
~
\hyperref[sec:BtpgDcAKeBQ]{\small{> > > NEXT SERMON > > >}}
\newline
\newline
創世記 18:16-21
\newline
\begin{longtable}{cl}
\hline
\hline
章節 & 經文 (和合本修訂版)\\
\hline
18:16 & \begin{tabularx}{0.7\textwidth}{X} 三人從那裡起程，面向所多瑪觀望，亞伯拉罕與他們同行，要送他們一程。 \end{tabularx} \\ \\ \relax
18:17 & \begin{tabularx}{0.7\textwidth}{X} 耶和華說：「我所要做的事豈可瞞著亞伯拉罕呢？ \end{tabularx} \\ \\ \relax
18:18 & \begin{tabularx}{0.7\textwidth}{X} 亞伯拉罕必要成為強大的國；地上的萬國都必因他得福。 \end{tabularx} \\ \\ \relax
18:19 & \begin{tabularx}{0.7\textwidth}{X} 我揀選他，為要叫他命令他的子孫和後代家屬遵行耶和華的道，秉公行義，使耶和華所應許亞伯拉罕的話都實現了。」 \end{tabularx} \\ \\ \relax
18:20 & \begin{tabularx}{0.7\textwidth}{X} 耶和華說：「所多瑪和蛾摩拉罪惡極其嚴重，控告他們的聲音很大。 \end{tabularx} \\ \\ \relax
18:21 & \begin{tabularx}{0.7\textwidth}{X} 我要下去察看他們所做的，是否真的像那達到我這裡的聲音一樣；如果不是，我也要知道。」 \end{tabularx} \\ \\
[1ex]
\hline
\hline
\end{longtable}
$^{1}$弟兄姊妹平安.
近日我聽到一個故事.
其實很久之前.
講到有一個牧師和一個司機.
他們死了.
去到天堂門口等候.
看看能否進入天堂.
結果最後怎樣呢.
講得故事當然相反.
結果居然是牧師不能進入天堂.
司機是可以進入的.
天使就解釋為何.
他說因為牧師講道都幾悶.
所以人們一聽他講道.
個個樣子好像在祈禱.
不過是在睡覺.
所以牧師一點熟能力都沒有.
相反司機原來是開小巴.
還要是開亡命小巴.
所以那些人一上了小巴.
就會立即去祈禱.
為司機祈禱.
求司機開車開得好一點.
並且為自己的生命有平安去祈禱.
所以開亡命小巴的司機熟能力非常大.
可以上天堂.
聽到的時候好像沒有什麼公義.
但感恩只是一個故事.
如果天堂還是沒有道理.
那就慘了.
但聖經告訴我們.
很清楚告訴我們.
我們所信奉的耶和華.
我們的上帝.
他是絕對的公平公義公正.
他是用公義去掌管世界.
審判每一件事.
剛才所讀出來的經文.
是講述上帝和亞伯拉罕.
他對亞伯拉罕的一個吩咐.

$^{41}$無論他自己或他的子子孫孫.
都要遵守他的律例典章.
都要行公義.
我們看看整件事的情況.
在十八章一開始的時候.
亞伯拉罕在第十八章一第一節這樣說.
剛才沒有讀出來.
但他看到的情況.
當時天氣很熱.
亞伯拉罕就坐在帳篷門口.
在那裡休息.
人來人往.
他就看左看右.
突然間看到對面有三個人.
這三個人的臉很陌生.
他知道一定是客人.
他按照過往聖經的教導.
一定要熱情地招待客人.
所以他就立即跑去.
跑去到這三個人面前跪下.
請求他.
請你接受我對你的款待.
到我那裡坐一會.
他就帶著帳篷進入地主之寓.
用水幫他洗腳.
叫僕人用最好的麵粉.
做一些餅給他吃.
自己又快點跑去牛欄.
找一隻又漂亮又嫩的牛.
叫僕人熬湯給這三個客人吃.
他真的出錢出力.
當他吃完之後.
希望他們能夠有氣流有力地繼續上路.
當時這個遊牧民族.
他們有一個習慣.
當客人吃完東西.
休息完繼續走的時候.
主人就要去送客.
現在通常我們飲宴完之後.
主人的送客就站在門口握手.

$^{81}$但在遊牧民族的時候.
他們的送客就要陪客人走一段路程.
走這段路程走多點路.
還是走一會走一條街就算了.
就表示到主人和客人的關係.
主人對客人的尊敬.
當時學者就計算過.
就和當時亞伯拉罕所處的地方.
在希伯倫的萬利這個地方.
然後就記載去到死海南邊的一個高山.
在這個高山就會看到.
這個索多瑪和俄摩拉這兩處的地方.
估計整個旅程就有50公里左右.
一般人大約走多久呢.
走12個小時左右.
所以我們看到亞伯拉罕.
他對於客人的款待是多麼熱誠.
可以說是出錢出氣出力.
不是發脾氣.
是出錢出力去招呼客人.
亞伯拉罕按照上帝給他們的誡命.
要招呼這些客女.
因為整個旅程他們辛苦了.
亞伯拉罕可以這樣說.
他把自己擁有的東西.
把最好的東西給這些客人.
我們想到亞伯拉罕也成為我們的榜樣.
當我們每一次去接待教會.
接待一些家人,親戚和周圍朋友的時候.
我們有沒有像亞伯拉罕一樣.
把最好的東西擺上.
我知道佳欣樓你們平時也有小食.
有咖啡喝的.
今天沒有看到的.
他們都是用最好的咖啡.
我看到錦繡那邊.
他們上個星期開始恢復朝聖茶室.
我在他們的群組裡面.
看他們一直商量的時候.
他們真的自己出錢出來.

$^{121}$買了新的咖啡機.
買最好的咖啡回來和你們分享.
這個就是聖經裡面的教導.
我們要把最好的和別人分享.
同樣亞伯拉罕也是按照聖經的教導.
去信服上帝的教導.
去接待客人.
若不是三不識七.
我為什麼要這麼好對待他呢?.
我們看到亞伯拉罕.
為什麼能夠和上帝的關係這麼好.
就是因為那份信服.
在喜伯來書裡面.
他記載了第十三章.
你們務要常存弟兄相愛的心.
不可忘記用愛心接待客人.
因為曾有接待客人的.
不知不覺就接待了天使.
喜伯來書的作者就在描述這個場景.
當時亞伯拉罕不知道他在接待的.
其中一個可能是耶和華.
另外兩個就是這個使者.
不過他就在這麼熱誠的接待下.
原來他就接待了上帝.
接待了他的使者.
帶來了很大的祝福.
所以就再一次去祝福亞伯拉罕.
向亞伯拉罕必要成為強大的國.
地上的萬國必因他得福.
其實這次不是第一次說的.
當第一次上帝和亞伯拉罕相遇的時候.
他呼召亞伯拉罕.
從他自己的家鄉吳爾.
要起行.
當亞伯拉罕信服他呼召的時候.
他會怎樣呢.
他說我必使你成為大國.
我必賜福給你.
使你的名為大.
你要使別人得福.

$^{161}$他再一次說到.
他要成為大國.
他的生命要成為別人的祝福.
在這裡他再一次說.
在19節.
他就說.
我揀選了他.
這就指亞伯拉罕.
為要命令他的子孫和後代家屬.
遵行耶和華的道.
秉行公義.
使耶和華所應許亞伯拉罕的話都實現了.
第一句就說.
上帝揀選了亞伯拉罕.
揀選這個字.
在舊的和合本就譯作眷顧.
也有些譯本就說到.
上帝認識亞伯拉罕.
這個眷顧這個字.
就說到很愛惜他.
在這裡就說到上帝和亞伯拉罕的關係.
是很親密的.
他很愛惜他.
愛惜到一個地步.
就是很想亞伯拉罕能夠明白他的心意.
人與人之間的關係都是這樣.
當我們是好朋友的時候.
或者我們是夫婦之間的關係的時候.
我們很想對方.
明白自己的心意.
我有時候看到姐妹們吵架.
是因為什麼呢.
大家是老友.
你都不告訴我秘密.
不告訴秘密就是不告訴他心意.
所以就生氣了.
或者是好朋友.
或者夫婦的時候.
我們會怎樣呢.
眼眉擦一擦.

$^{201}$對方就知道我的心意.
就是一個這麼好的關係.
所以上帝和亞伯拉罕很愛惜他.
他當他是好朋友.
那個關係就是很想亞伯拉罕明白他的心意.
甚至想到不單單是亞伯拉罕明白.
還要他的子子孫孫都要明白.
要一代傳一代教他們.
要秉行公義.
遵守耶和華的律例.
點章.
在上兩個星期.
香港發生一宗很兇殘的兇殺案.
看到都很害怕.
這宗兇殺案竟然是由爸爸帶領.
然後策動全家人.
老婆兩個兒子.
後期才發現原來他的情婦也有份殺人.
社會學者指出.
這是一個家庭價值觀徹底崩壞的現象.
通常良好的價值觀從哪裡開始呢?.
就是由家庭父母開始.
父母除了言教之外.
還會以身作則影響下一代.
小時候父母怎樣教你呢?.
我父母教我們教我們教我們.
小時候父母怎樣教你呢?.
我父母教我們教我們.
花錢時要看好錢.
千萬不要花在未來錢.
所以很吝嗇.
我快要落底了.
不要再花錢了.
有些父母教我們.
吃飯時千萬不要飛漲過河.
很沒禮貌.
這些就是家庭教育.
上帝揀選了亞伯拉罕.
他要一代傳一代.
將他的吩咐傳到後代.

$^{241}$他盼望藉著亞伯拉罕的家庭.
這個家族.
將他的公義.
盼望藉著行公義時.
將愛和和平一直傳下去.
以至這個家族是使人蒙福的家族.
在《申命記》中.
摩西已經秉承這個說法.
上帝透過摩西定立了一件事.
每七年的時候.
他們會在駐彭節裡.
叫所有的傳會眾.
無論是男女老的.
或者很小的.
他們要集合在一起.
再一次將律法書宣讀出來.
然後又要來講解.
命令所有的會眾.
都要學習如何敬畏上帝.
如何進行律法書所講的一切說話.
在這裡我們也要反省.
我們要多讀聖經.
聖經教我們.
我們每一天都要將上帝的說話.
看了進入我們的心靈.
然後我們要怎麼走出來.
每一天我們看了之後.
到集合第七天的時候.
我們就在這裡.
我們一起有這個敬拜.
一起去學習上帝的話語.
雖然有時候我們姐妹之間聊天.
我們會說我們沒有記性.
上帝看完之後.
不過一個小時.
我已經不太記得.
我今天零修看過哪一段經文.
但不要緊.
每天儲一些.
一定潛移默化.

$^{281}$在我們心靈裡面.
是有功效的.
接著到第20到21節.
他就說.
索多瑪和俄摩拉的罪惡非常嚴重.
控告他們的聲音也非常大.
我要下去察看他們所做的.
看看情況是否差到.
我聽到居民的哀聲.
是否真的這麼差.
如果不是的話.
我就不會施行這個審判.
上帝是全知全能.
他不可能不知道.
不過他這樣去和亞伯拉罕對話.
他要亞伯拉罕知道.
他是一個很公義的上帝.
他一定會查問到底.
然後才作出審判.
不守上帝的道的人.
終必得到審判.
我夢召之前是教書的.
其中一年我去代課.
在一間小學.
教特殊班.
這個特殊班是將五年班和六年班.
在英文和數學很差成績的.
出來我負責教的.
通常很差成績的.
在行為上都一般般.
不過都有一些例外.
其中一個五年班的男同學.
他是插班進來的.
為什麼會插班進來呢.
因為在另外一區.
全部學校都知道這個人非常麻煩.
所以沒有人肯收.
他說你這一區有些特別班.
就將他放在我代課的學校.
我對他一年.

$^{321}$真的辛苦了.
他上課不斷倒彈.
其實他還很聰明.
他知道說髒話就會被罰.
所以他就說一些諧音的髒話.
不過又令到你聽到.
又說一些三級的笑話.
他又不會正式說.
又會令到你明白.
那班同學五六年班都會明白.
還有他又聰明到.
他去同學的家.
竟然很聰明到.
可以哄到同學拿了他爸爸的傳設出來.
又看看他的成績表.
他爸爸怎麼簽名呢.
他竟然可以冒認他爸爸簽名.
去了銀行拿了三次錢.
到第三次的時候就被發現了.
所以被抓了.
然後他又很勤力上學.
上下午班.
他上午一早就上學了.
就過隔離學校上學.
就看哪些人上體育課.
他又進去一起踢球.
結果就踢斷了學生的腳.
手腳忘記了.
晚上又帶三四個同學去外面.
打爛了人家車的玻璃窗.
到最後到學期尾的時候.
有三位同學就被抓了.
上法庭.
主任聽完審回來.
進到教員室.
他就大拍桌子.
這個世界是沒有公理的.
這樣都可以.
原來最後審判的結果.
其中兩個同學.

$^{361}$被判入兒童院.
唯獨這個五年級的男同學.
他停課兩星期之後.
可以繼續上學.
為什麼會這樣?.
因為他爸爸有錢.
請了律師求情.
所以他兩個星期之後.
就繼續上課.
真的沒有公理.
到暑假前一天.
最後一堂.
上課後.
他爸爸打電話給我.
他爸爸就罵他.
這次不是罵我.
是罵他兒子.
他兒子是個壞人.
早上偷了我二千元.
然後不知道去了哪裡.
二十多年.
二千元已經很大了.
偷了之後.
還不知道跑到哪裡去了.
我就發抖.
他這樣溺愛他的兒子.
不教導.
有空就打電話回來罵老師.
但到現在這樣.
他自己自食其果.
受到審判.
所以我不知道現在兒子怎樣.
這樣家庭教育.
這樣下去的時候.
兒子不會有出路.
我們需要.
正如聖經所說.
我們需要以真理.
好好教導下一代.
以致他們走上一條.

$^{401}$正直幸福光明的路.
走一條正直幸福光明的路.
對於自己的生命.
是有益處的.
也對整個社會.
整個世界.
是有益處的.
使人得福.
聖經裡說到.
神絕對不會以有罪.
成為無罪.
我想在座當中.
無論我或大家.
都一定有罪的.
是不是很絕望呢.
在詩篇第89篇第14節這樣說.
公義和公平.
是你寶座的根基.
上帝管治世界.
在審判萬事的時候.
都會以公平公義來處理.
但是感恩的是.
有慈愛和信實在祂的面前.
就像剛才我們詩歌裡說到.
有上帝的愛.
祂的寶血.
來潔淨我們的罪.
在約翰福音第一章也說到.
耶穌基督到,成了肉身.
耶穌基督來到這個世界裡.
住在我們中間.
有恩典.
就是慈愛和信實.
也有真理.
住在我們中間.
有恩典.
有真理.
就是有公義,公平.
也有上帝的慈愛和信實.
這個公義是什麼呢?.

$^{441}$我們常說公義.
我們每個人都愛惜自己.
一定為自己爭取最好的.
不過這個公義.
是尊重別人.
看其他人就像看自己.
所以這個公義.
就是看其他人的事情.
看其他人的看法,觀點.
去想自己的一象.
我絕對不會想著一定要贏.
我自己是最適合的.
最大的.
我來跟別人爭奪.
我來跟別人打.
我們要想的就是.
每一件事情.
我們冷靜下來想的時候.
我們的目標.
都是要建立我們自己.
建立我們的家庭.
建立整個社會.
能夠有一個和諧,和平,幸福的人生.
一個和諧,和平的社會.
所以上帝一定要來剿滅.
這個索多瑪和俄摩拉.
因為他一定要跟亞伯拉罕說.
他要顯明他的公義.
因為這兩個地方是罪惡滔天的.
我們在創世記第十九章的時候.
我們就看到.
在人的關係上.
他們是很亂來的.
總之我喜歡那兩個男人.
我就要出來.
就要去強行去拿人.
在以西結書裡面.
以西結才說到.
這兩個城市有什麼罪惡呢.
他們說他們心驕氣傲.

$^{481}$一群人很驕傲.
不可一世.
每個人都說自己最厲害.
每個人都說自己最對.
然後就說到.
他們說糧食飽足.
享安逸.
卻不願意去向貧窮人.
施以援手.
不願意去幫助其他人.
就好像我剛才在聖餐裡面說到.
上帝要我們分享.
幫助有需要的人.
他們那群人是怎樣呢.
恨可憎的事情.
令人討厭.
那些不是耶和華所閱立的事情.
所以上帝要來審判.
當耶和華說到.
要你去審判的時候.
亞伯拉罕仍然有那份.
對這兩個城市那份愛.
人那份愛.
他就去跟上帝談判.
按照上帝的屬性.
去跟他談判.
上帝你是公義的.
你是不是真的要將一些義人.
和惡人一起毀滅呢.
如果我在這兩個城市裡面.
我找到有50個義人的時候.
你是不是也會將那些義人.
和惡人一起全部剿滅呢.
上帝說你找到50個的時候.
我就不滅這兩個城了.
他就一直減下去.
50 45 40 30 20.
10個.
他說如果去到10個.
你找到10個的話.

$^{521}$因為我是公義的主.
我也有慈愛和信實.
我也不會將全城毀滅.
亞伯拉罕說到10個之後.
他就沒有再求了.
因為他知道他找不到.
我們看回亞伯拉罕向上帝的祈求.
他一直按照上帝的屬性.
按照他的應許.
他不厭其煩地求他.
憐憫這兩個城市不要滅絕他們.
在聖經裡面.
他對我們的教導.
就是要行切的禱告.
什麼叫做行切的禱告.
恆久熱心的祈禱.
人很容易就會灰心.
相信很多人都會這樣.
天天為我們家人.
信耶穌來祈禱.
祈禱一年兩年.
我們就會覺得好像什麼氣息都沒有.
很灰心不想求下去.
我自己也是.
我每天早上和晚上.
我都會提名我家人.
他們未信的人.
為他們來信主祈禱.
其中一年兩年三年.
真的會灰心.
不過又可以在這裡見證.
我爸爸媽媽在他們離開世界的時候.
他們都信了耶穌.
我爸爸我親自幫他們洗禮.
在他們灰心的時候.
上帝在他們出其不意的時候.
你會發覺他們有改變.
甚至我弟父.
我弟父在三月.
上年八月初的時候.

$^{561}$他就離開了世界.
我都說他還沒有癌症的時候.
我已經開始.
就是一個禱告的名單.
很灰心的.
去到醫院探望他的時候.
他很冷淡的.
他祈禱又好像尷尬的.
不想又不想推我.
但他癌症差不多四年.
每次去探望他.
我都為他祈禱.
他就說很喜歡有人為他祈禱.
不只是我.
其他人也很喜歡.
最初想他客氣.
慢慢多去為他祈禱.
很喜歡.
感覺心靈很平安很舒服.
我覺得他客氣.
直到有一天.
他晚上跟弟弟說.
他想拿十字架出來.
那個十字架.
是他入院沒多久.
一個信耶穌的朋友.
送給他的.
我們是一起的.
拿完我們走了.
我就知道他把十字架放在抽屜裡.
不再拿出來了.
但那天他拿出來.
跟著他跟弟弟說.
他想信耶穌.
所以我想邀請大家.
有時候去為一些事情.
為我們家人的信主.
我們千萬不要說.
祈禱這麼久.
都是差不多的.

$^{601}$沒什麼看的.
但我們知道.
上帝邀請我們.
恆久熱切地不斷禱告.
到時候.
上帝必定會來成就.
我們也看到.
上帝最後都滅了.
這個索多瑪和阿摩拉.
因為在當中.
他找不到十個異人.
他真的要按真理去審判他們.
但他先救出了羅德他們一家人.
羅德他們一家人都不是好人.
絕對沒有完全的好人.
不過他們都是信奉耶和華的人.
在聖經裡說到異人就是信奉耶和華的人.
所以上帝就用他的慈愛.
用他的應許.
去救了他們出來.
上帝是滿有恩典.
也都滿有真理.
公義公平是上帝寶座的根基.
但在他前面.
有慈愛和信實在當中.
所以弟兄姊妹.
弟兄姊妹.
我們要去努力行出上帝的公義.
但有時候我們都會軟弱.
但我們知道.
上帝的慈愛來包底.
但上帝的心意.
是要我們每一個人的生命.
和我們的後代.
都能夠去秉行公義.
以至我們的生命.
和我們周圍人的生命.
可以因著我們去得福.
求主賜恩.
我們一起禱告.

$^{641}$天父上帝.
多謝主你給我們的話語.
求你給我們力量.
以至我們自己的生命.
行出主你的公義.
以至我們的生命.
也都能夠好好地教導下一代.
讓周圍的人都能夠過一個和諧.
幸福的生活,生命.
求主賜恩.
多謝你垂聽禱告.
奉靠主耶穌基督得勝的名字而求.
阿 明.
\newpage



\section{歌羅西書 3:1-17}
\label{sec:BtpgDcAKeBQ}
\textbf{2023-03-19 午堂《天地之間》歌羅西書3\_1-17 講員:吳景霖傳道}
\newline
\newline
連結: \href{https://youtube.com/watch?v=BtpgDcAKeBQ}{\texttt{https://youtube.com/watch?v=BtpgDcAKeBQ}} ~~~~ 語音日期: 2023-03-21
\newline
\newline
\hyperref[sec:1tknD6q3Urw]{\small{< < < PREV SERMON < < <}}
~
\hyperref[sec:index_chronic]{\small{[返順時目]}}
~
\hyperref[sec:index_scriptual]{\small{[返順卷目]}}
~
\hyperref[sec:W3jetg3vhiA]{\small{> > > NEXT SERMON > > >}}
\newline
\newline
歌羅西書 3:1-17
\newline
\begin{longtable}{cl}
\hline
\hline
章節 & 經文 (和合本修訂版)\\
\hline
3:1 & \begin{tabularx}{0.7\textwidth}{X} 所以，既然你們已經與基督一同復活，就當求上面的事；那裡有基督，坐在神的右邊。 \end{tabularx} \\ \\ \relax
3:2 & \begin{tabularx}{0.7\textwidth}{X} 你們要思考上面的事，不要思考地上的事。 \end{tabularx} \\ \\ \relax
3:3 & \begin{tabularx}{0.7\textwidth}{X} 因為你們已經死了，你們的生命與基督一同藏在神裡面。 \end{tabularx} \\ \\ \relax
3:4 & \begin{tabularx}{0.7\textwidth}{X} 基督是你們的生命，他顯現的時候，你們也要與他一同在榮耀裡顯現。 \end{tabularx} \\ \\ \relax
3:5 & \begin{tabularx}{0.7\textwidth}{X} 所以，要治死你們在地上的肢體；就如淫亂、污穢、邪情、惡慾和貪婪—貪婪就是拜偶像。 \end{tabularx} \\ \\ \relax
3:6 & \begin{tabularx}{0.7\textwidth}{X} 因這些事，神的憤怒必臨到那些悖逆的人。 \end{tabularx} \\ \\ \relax
3:7 & \begin{tabularx}{0.7\textwidth}{X} 當你們在這些事中活著的時候，你們的行為也曾是這樣的。 \end{tabularx} \\ \\ \relax
3:8 & \begin{tabularx}{0.7\textwidth}{X} 但現在你們要棄絕這一切的事，就是惱恨、憤怒、惡毒、毀謗和口中污穢的言語。 \end{tabularx} \\ \\ \relax
3:9 & \begin{tabularx}{0.7\textwidth}{X} 不要彼此說謊，因為你們已經脫去舊人和舊人的行為， \end{tabularx} \\ \\ \relax
3:10 & \begin{tabularx}{0.7\textwidth}{X} 穿上了新人，這新人照著造他的主的形像在知識上不斷地更新。 \end{tabularx} \\ \\ \relax
3:11 & \begin{tabularx}{0.7\textwidth}{X} 在這事上並不分希臘人和猶太人，受割禮的和未受割禮的，未開化的人、西古提人、為奴的、自主的；惟獨基督是一切，又在一切之內。 \end{tabularx} \\ \\ \relax
3:12 & \begin{tabularx}{0.7\textwidth}{X} 所以，你們既是神的選民，聖潔、蒙愛的人，要穿上憐憫、恩慈、謙虛、溫柔和忍耐。 \end{tabularx} \\ \\ \relax
3:13 & \begin{tabularx}{0.7\textwidth}{X} 倘若這人與那人有嫌隙，總要彼此容忍，彼此饒恕；主怎樣饒恕了你們，你們也要怎樣饒恕人。 \end{tabularx} \\ \\ \relax
3:14 & \begin{tabularx}{0.7\textwidth}{X} 除此以外，還要穿上愛心，因為愛是貫通全德的。 \end{tabularx} \\ \\ \relax
3:15 & \begin{tabularx}{0.7\textwidth}{X} 你們要讓基督所賜的和平在你們心裡作主，也為此蒙召，歸為一體。你們還要存感謝的心。 \end{tabularx} \\ \\ \relax
3:16 & \begin{tabularx}{0.7\textwidth}{X} 當用各樣的智慧，把基督的道豐豐富富的存在心裡，用詩篇、讚美詩、靈歌，彼此教導，互相勸戒，以感恩的心歌頌神。 \end{tabularx} \\ \\ \relax
3:17 & \begin{tabularx}{0.7\textwidth}{X} 你們無論做甚麼，或說話或行事，都要奉主耶穌的名，藉著他感謝父神。 \end{tabularx} \\ \\
[1ex]
\hline
\hline
\end{longtable}
$^{1}$各位早晨.
麻煩主席讀了這麼長的經文.
因為我想信息比較完整一點.
我不會之後再和大家讀一次.
於是辛苦他再讀一次.
今天講歌羅西書.
其實之前都有講.
不同的傳達人都有講過.
今天講到第三章.
可能都要講一講歌羅西書的背景.
哥羅西教會.
其實很可能是保羅在爾忽所的時候.
那三年的時間.
由他的同工爾巴弗所建立的.
就是說哥羅西教會並不是保羅親手建立的.
但是由他的同工建立.
並且保羅其實也很關心哥羅西教會.
所以他是清楚哥羅西教會的情況.
哥羅西教會大部分是由外邦人組成的.
並且他們本來信了主之後.
情況是好好的.
信得都挺好的.
保羅是讚賞他們的.
但是後來就出問題了.
因為有一些異端走到他們中間.
大家知道什麼是異端.
異端就是其實他都是在用基督的信仰.
說的是很相似的.
大家都是看聖經.
但是他有一些核心的問題.
核心的教義是出了偏差的.
就為之異端.
表面上都是基督宗教.
但其實實際上已經出了問題.
這個異端叫做什麼呢.
叫做洛斯底主義.
你大概知道名字就可以了.
不用學術地記住.
洛斯底主義.
他很強調一件事就是.

$^{41}$他是靈肉二分的.
靈性和肉體是分開兩件事.
是沒有相關的.
他很強調這件事.
而強調這件事是靈是好的.
肉體是徹底的不好的.
完全是分開了這兩件事.
出了什麼問題呢.
他的主張就是說.
總之最重要的就是你的靈.
你的肉體那些東西是不好的.
是和靈相對的.
於是乎有一種洛斯底主義就是.
肉體是要克服它的.
克制它的.
千萬不要對它那麼好.
知道嗎.
就是說你所有東西都是惡的.
身外物都是惡的.
除了你的靈性之外.
你的身外物是惡的.
你的錢是萬惡的.
知道嗎.
然後你坐的椅子也是惡的.
所有東西都是惡的.
最重要就是靈.
他說的是要有一種苦待己生.
於是乎吃東西是不可以吃好東西的.
吃一些最差的.
吃白飯喝清水.
因為我要去做身體檢查.
我這兩天都是這樣.
我就覺得很痛苦.
剛才在外面的咖啡檔.
那些咖啡那麼香.
但是我不能喝.
就是要苦待自己.
然後上班不要坐車.
知道嗎.
由錦瑟走出去.

$^{81}$就是你苦待自己來到黑性的肉體.
就是操練自己的靈魂.
就是有一種這樣的想法.
但是洛斯底主義.
其實也是會有另一個極端的.
怎樣才算是另一個極端.
因為靈魂和肉體是不相干的.
兩樣東西不相干的話.
我的肉體怎樣做都好.
是不會影響我的靈魂的.
於是乎出現一種怎樣的情況呢.
我的肉體怎樣放縱都好.
對我的靈是沒有關係的.
其實我們不多不少有這樣的想法.
我有了一個救恩.
我的靈魂可以上天堂.
那我身體做甚麼都好.
不緊要的.
犯罪得罪神.
晚上就騎個肚子.
不影響我的靈魂的救恩.
其實是有些問題的.
所以這種極端的洛斯底主義.
它是會放縱自己的情慾的.
因為反正和靈魂都是沒有相干的.
兩樣東西沒有相干.
那就做甚麼都可以.
說得好像有點.
不是很清楚我說甚麼.
但在中國傳統佛教有個故事.
叫做《濟公活佛》.
我小時候看亞視的時候有套電視.
就是《濟公活佛》.
它有句名句就是.
酒肉穿腸過佛祖留心中.
它是一個和尚來的.
它應該不能吃肉不能喝酒的.
但它照吃照喝.
它就說酒肉穿腸.
就出來了.

$^{121}$佛祖留心中.
這個說法其實有多少洛斯底主義.
就是靈性和肉體是可以分開的.
在他們眼中.
大概記得我這些東西就記得了.
待會說到的時候.
可能有些東西會稍微提起洛斯底主義.
我們要說回今天的經文.
上面的事和地上的事到底是甚麼來的.
保羅提醒我們要把心思意念放在上面的事.
他沒有清楚說到上面的事是甚麼來的.
但是他有暗示到.
哪裡有基督坐在神的右邊.
你可以從這個去推斷.
就是上面的事一定是在說關於基督的.
關於耶穌的.
關於可能是將來的新天新地的國度.
因為耶穌基督現在已經坐在神的右邊.
右邊是一個很尊貴的位置.
右邊是一個很重要的.
你在神的右邊.
就是說你已經和神一樣有權柄.
就是基督已經掌權.
也可能是說將來新天新地的事.
而地上的事相對是甚麼.
就是第五節或者第八節說的那些.
淫亂,污衊,邪情,惡慾和貪婪.
惱恨,憤怒,惡毒,毀謗.
或者口中的污衊言語.
這裡保羅強調我們要想的.
心思意念要放在上面的事.
上面的事優先.
為甚麼.
原因是因為你們已經死了.
你知道不是實際死了.
實際死了就看不到那封信.
也不需要思念甚麼上面,下面.
因為你已經去了上面.
但他在用甚麼呢.
他在用死作為一個比喻.

$^{161}$一個比喻.
你要當自己已經死了.
這樣用這種態度去面對我們的信仰.
事實上當我們信耶穌的時候.
可能大家都有記得.
或者有不同的說法.
就是說我們的舊我已經死了.
今日的我們是新我.
我們是新造的人.
舊事已過 一切都變成新的了.
我們的信仰是相信我們信了之後.
我們的生命是會得到一個更新的.
這個是保羅一向的主張.
他不同的書卷裡面都有講.
不過我們華人就很忌諱講死的.
不會講死的.
就算講死都會加句.
大家理事.
鬼佬都是的.
其實迷信在哪裡都有.
但很忌諱去講.
事實上死亡的確是負面的.
死亡的確是令我們畏懼的.
令我們不是很想提的.
但是其實我們也都明白.
在我們的信仰裡面.
死亡某程度就是講我們基督徒.
由地上的人生轉移到新天新地的一個過程.
我們地上的生活已經完了.
我們只是去到新天新地那裡.
開始另一個新生活.
所以某程度我可以說.
我們移民了去天國.
是移民了去天國.
雖然今天信徒裡肉身活著.
但是我們的靈性其實是經歷了一次死亡.
並且一次更生.
我們是在等待著我們移民的那一刻.
等待著我們移民去新天新地的那一刻.
用等死好像難聽一點.

$^{201}$不應該是等死的.
我們在等著去到新天新地的那一刻.
是用一種預備移民的心態.
去面對我們的人生.
當你真的要移民的時候.
不知道大家有沒有經歷過.
有些人可能就算沒有移民過.
可能都會搬過屋.
當你要去到移民的時候.
最頭痛的就是到底我要帶些什麼走.
我要帶些什麼走.
看著整個屋都是東西.
我都搬過幾次屋的時候.
真的非常頭痛.
到底要帶什麼走呢.
去一個新地方.
如果那個地方是你完全陌生的時候.
你是真的不知道你帶些什麼是有用.
帶些什麼是沒有用的.
到底帶些什麼是浪費的呢.
帶些什麼是一定要拿走的呢.
如果你明知你要離開了.
離開這間屋.
譬如你可能在錦繡搬出去.
你會不會說家裡的電視都舊舊的.
先換個80吋.
然後再弄一個三角琴放在這裡.
或者你不會覺得這間屋舊舊的.
裝修一下它.
你裝修它來租給別人.
可能你會.
但如果你說你不要的.
你賣了的了.
你不會在這一刻.
你還在那間屋裡弄個三角琴在這裡.
因為你知道有些東西是帶不走的.
你可能會預備一些東西是將來有用.
譬如你預備一個新的行李箱.
或者買些東西是會搬得走的.
你接下來新屋或者移民.

$^{241}$你會用到的.
你就會帶.
保羅提醒我們.
基督徒你已經死了的原因是.
其實他要我們用一種新的眼光.
去審視你現在的生活.
你今天你做的一切.
到底你對於將來.
你上天家有沒有用.
你現在做的事情.
到底你將來上到天堂的時候.
有沒有價值.
今時今日我們覺得很多事情是很重要的.
譬如有些肢體是做生意的.
可能你的生意就是最重要的.
生意就是養妻活兒的.
或者養家的.
很重要的.
或者是我們香港人的核心價值是什麼.
買樓.
很多人經常都說核心價值買樓.
最重要就是上車.
供兩代不要緊的.
你爸爸媽媽賣了一層樓給你做首期.
再幫你供也不要緊的.
彷彿人生的價值就是要買樓上車.
我也不太知道為什麼.
但是有些人就是這樣.
一層樓不就是最重要的嗎.
或者是你的事業.
事業是最重要的.
有些人為了自己的事業.
是可以整個家都不理的.
白天就是七點出門.
晚上就十一點回來.
一個星期做了七天.
為了自己的事業.
有一些成功感.
又或者是子女的學業.
聽不少這樣的情況.

$^{281}$專門搬去名校區那裡.
去到那裡就是為了子女入名校.
我昨天聽到高明軒博士說過才奇怪.
有間名校只收一月的小朋友.
於是你怎樣令到小朋友一月出生.
就要計算時間.
我告訴你我試過.
我不是為了一月出生入名校.
因為我和我老婆都是八月生日的.
所以我想夾到兒子一次的慶祝.
結果失敗了.
七月.
所以他早了一個月.
不能一起慶祝.
其實就是這樣.
有些人覺得很重要的東西.
你就會將你整個心機放進去.
你將你所有的時間精神都用在那裡.
但是現在保羅說的就是.
當你去到人生盡頭的時候.
這些東西是不是真的重要.
這些東西是不是真的幫到你在天國那裡.
你現在你的生意.
你的事業.
你的學業.
或者是你的房子.
你的靚車.
這些東西到底在天國對你有沒有用.
如果耶穌基督今天回來.
回來之後和你聊聊天.
你一個星期的時間是怎樣分配的.
你覺得你說出來的答案.
耶穌滿不滿意.
耶穌會不會願意聽到你.
我一天七點鐘出門.
晚上十一點回來.
基本上就是上班.
你覺得耶穌滿不滿意這個答案.
當然我們其實要好好去衡量.
我們自己怎樣去為神工作.

$^{321}$另一方面我們更加要小心的就是.
我們要逃避一些東西.
我們要逃避這個世界的邪惡.
我們要逃避一些上帝不喜悅的東西.
例如剛才說過的那些.
淫亂 污穢 邪情 惡慾.
貪婪 惱恨 憤怒 惡毒.
毀謗和口中污穢的言語.
我們要逃避這些東西.
一方面追求合上帝心意的東西.
一方面是逃避這些東西.
而且這些東西不是保羅隨便寫出來的.
不是隨便寫出來的.
他寫出來是有些目的的.
這個世界有很多很邪惡的東西.
為甚麼專門挑這些出來說呢.
其實你會發現這些東西.
某程度是一種追求.
一種放縱和一種享受.
剛才我開始說的.
諾斯底主義的異端.
他們侵入了哥羅西教會.
而其中一個特徵就是.
他們認為靈魂和肉體分開是沒有關係的.
所以你的肉體你如何去放縱它.
是沒有所謂的.
在這樣的情況下.
那些人就大件事了.
他會放縱自己做一些邪惡的行為.
保羅就是提醒他.
不是的 不行的.
不是這樣的.
上帝會追究你的.
你不要以為你做完這些事情是沒有責任的.
上帝會追究你的.
這一點很重要.
這個很值得我們深思.
特別是我剛才提到.
我們有一種想法就是.
我們好像一決志信耶穌.

$^{361}$就拿到天堂入場券.
有沒有聽過這個說法.
信耶穌先拿了天堂入場券.
拿了入場券之後怎樣.
接著回不回教會那些事情.
接著再說吧.
我現在叫你決志而已.
不是叫你回教會.
我偶爾也會拿出來講這個例子.
不應該是這樣的.
我們信仰不應該是這樣的.
你覺得好像拿了入場券之後.
我就等死等上天堂那樣.
這個有多少諾斯底主義.
就是覺得靈魂和肉體是兩回事.
我靈魂得救就行了.
接著我的肉體就繼續做平日的我.
不是這樣的.
信仰是一個整體.
你的身體可以說是影響你的靈魂.
你的靈魂也都影響你的身體.
也是一個整體.
接著就說一下穿上新人.
什麼叫穿上新人呢.
其實呢.
本來用穿上新人來比喻.
是說我們基督徒.
我們已經是和基督徒同死的.
我們已經是舊我已經死了.
所以我們要有一個新的我.
他用穿上的某程度.
就好像我們穿衣服一樣.
把衣服勒上去.
一件衣服其實是在反映著你的人.
或者是在表達一些很重要的訊息.
譬如呢.
我經常都有一些靈體姐妹見到我.
就問我為什麼這麼斯文.
我也要檢討一下.
Thomas昨晚又問我為什麼這麼斯文.

$^{401}$我說我回來崇拜.
我要回來當值.
我當然要斯文的.
是不是.
勒上一件西裝就斯文.
可能我要反省.
我平時是不是太膽小.
平時太過污穢.
所以大家覺得我穿西裝就很斯文.
其實是一個語言來的.
一個語言.
在告訴別人.
我在這個場合有些特別的東西.
你看到師班的各位靈體姐妹.
你穿上那件師班袍就有一個訊息在.
你是一個在事奉的人.
在服侍的人.
甚至你平時是紅色的.
但是安息禮反轉的是紫色的.
你是有一些訊息在給人.
穿上新人.
其實都是在給一些訊息.
給外間的人.
當時保羅的社會.
不同民族的人自然都有他們不同民族的衣著.
而問題是猶太人和外邦人.
其實是大家很不和的.
大家都看對方不順眼的.
看對方不順眼.
然後我一眼就知道你是什麼人.
我看你的衣著.
就知道你是什麼人的時候.
就很容易出現一種分歧.
標記了那個人.
哦 猶太人來的.
哦 外邦人來的.
這樣的情況.
除此以外.
很多人都講到有些不同的人種.
有很多不同的人種.

$^{441}$特別是講到有些是自主的和為奴的.
奴就是奴隸.
自主就是自由.
就是他屬於不是別人的奴隸.
甚至他可能是擁有奴隸的人.
這個是在講社會階級的問題.
奴隸就是沒有人權的.
他只是別人的財產.
他甚至可能是衣不閉體都不出奇的.
他沒有自己的財產.
主人如果死了的話.
奴隸是要賠償的.
當時的情況.
所以他的社會地位是非常低的.
在這裡講的就是.
原來不論人種.
不論你的身份.
不論你在社會的地位.
你有錢沒錢好.
其實你都應該要穿上新人.
你都要穿一個新人.
去代替你舊有的身份地位.
其實就好像學校裡面的校服.
你平日是身光頸靚都好.
還是T恤短褲都好.
你回到學校你都要穿校服.
你都是那套校服.
不會因為你家裡有錢一點.
你那套校服就傷金的.
而老師對著這麼多的同學.
全部一清一色穿校服的.
就要對他們一視同仁.
而學生穿著校服.
亦都代表了學校.
所以他的行為舉止.
就要被人看到.
他有一個學生的典範.
我兒子經常跟我投訴.
我現在不可以在街上吃東西.
因為老師說穿著校服不准在街上吃東西.

$^{481}$那我唯有陪他一起遵守規矩.
因為他代表著學校.
他穿著校服.
同樣地.
你穿上新人.
你代表著耶穌基督.
你穿上新人.
你就是一種外顯的.
肉眼可見的形式.
去到讓人見到.
你是一個基督徒.
讓人見到你是一個基督徒.
有風險的.
我告訴你.
有風險的.
早前我看到一段新聞.
有一對年輕的情侶.
他們在公開的場合.
即是海旁的公園那裡.
在那裡跳舞.
然後就被捕.
被捕了.
為什麼呢.
他們被控兩條控罪.
第一條就是鼓勵墮落和公開賣淫.
第二條就是從事以破壞國安為目的的聚會.
你聽到的時候.
難道他們穿著很少保險.
難道他們不穿衣服跳舞.
不是的.
他們跟我一樣.
除了頭和兩隻手指外.
什麼都看不到.
長袖衫長袖褲.
為什麼呢.
因為他們在伊朗.
他們是伊朗人.
伊朗的女人是要蒙頭.
他們除了臉和手指外.
所有的地方都要穿著保險.

$^{521}$那對情侶.
他們就是在公開場合.
女生脫下頭巾跳舞.
因此他們被捕.
坐牢十年半.
為什麼他們要在公開場合跳舞呢.
為什麼明知道這樣的情況.
身處這樣的國家.
你還要這樣去做呢.
因為這樣是不公義的.
這樣是不合理的.
在座有哪位姐妹.
你願意你的人生就是用布勾著頭.
只是露面出來.
略為文明的人都會覺得.
這是非常野蠻的一件事.
他們用了他們的自由.
用了他們.
犧牲了他們的青春.
去挑戰這個野蠻的制度.
他們犧牲了自己.
基督徒穿上新人.
某程度上是一樣的.
危險的.
因為你是用一些可見的形式.
去讓人見到.
我們是和這個世界不同的.
當然我不是叫你奇裝異服.
不是叫你.
我彰顯我是基督徒.
我穿著一件斯班圖.
搭西鐵回家.
你去到村巴站.
已經每個人都看著你.
我不是說這一種.
你的人生是會讓人見到.
你是與別不同的.
你是一個基督徒.
你不出聲.
人家都覺得你是基督徒.

$^{561}$應該是這樣的一種模式.
不是說我要坐車的時候.
我也有個朋友是這樣.
拿著十字架這樣祈禱.
然後旁邊那個拍他的手.
叫他放下.
不是這樣古怪的行徑.
有時候我們是要做一些事.
讓人知道.
為什麼你是不同的.
當這個世界說.
人不為己.
天誅地滅的時候.
我們基督徒是說.
我們犧牲自己.
用我們的愛去愛其他人.
我們愛其他人好像自己一樣.
這一種就是基督徒.
奇怪嗎.
奇怪的.
當這個世界的人.
可能覺得我們很傻.
很傻的.
會拿我們好處的.
借錢專門找基督徒借的.
有些所謂教會老鼠.
就是專門去到不同教會.
進去認識人.
認識完人之後.
借錢借錢就走人了.
貪你們什麼.
貪你們善良.
在這個世界上.
你的生活.
你去到做見證.
你去到對一些不公義的東西去發聲.
去到反抗.
你就是有風險的.
但是我們的信仰告訴我們.
你要穿上新人.

$^{601}$聖經的教導.
你要用你的身體.
身體力行地走出來.
我們的信仰從來都不容易.
只不過有時候.
坦白說.
我們將聖經的信息弱化了.
弱化到一個什麼地步.
就是讓大家舉手決志.
然後等上天堂.
記得每個星期回來教會那裡坐一坐.
在家裡就讀經靈修.
我不是說不重要.
這些不是待會被人投訴的.
重要的.
但是我說的不單止是這樣.
你的人生難道就是.
拿了天堂入場券之後等死.
我想你都不會這樣想.
我們一定是有些事情要做.
我們穿上新人.
是不是就等於是一種離地的生活呢.
我們是不是就是一種與世隔絕的生活呢.
第十二到十七節要讀一讀.
所以你們既是神的選民.
聖潔蒙愛的人.
要穿上憐憫恩慈.
謙虛溫柔和忍耐.
倘若這人與那人有嫌隙.
總要彼此容忍彼此寬恕.
主怎樣寬恕了你們.
你們也要怎樣饒恕人.
除此以外還要穿上愛心.
因為愛是貫通全德的.
你們要讓基督所示的和平在你們心裡作主.
也為此蒙召歸為一體.
當用各樣的智慧.
把基督的道豐豐富富的存在心裡.
用詩篇讚美詩靈歌彼此教導.
互相勸誡.

$^{641}$以感恩的心歌頌神.
你們無論做什麼.
或說話或行事.
都要奉主耶穌基督的名.
要有些勒卡也行.
這裡說的是一種生活態度.
我們信耶穌的人.
當你穿上新人.
你就是用一種新的態度去生活.
圍繞著的主要是你怎樣與人相處.
你見到憐憫恩慈謙虛溫柔忍耐.
你一定是與人相處的.
你不會是對神的.
你不會說對神要憐憫他.
是上帝憐憫你.
不是我們憐憫上帝.
這些一定是對人的.
你都不會自己去對.
需要對自己去溫柔忍耐.
通常就是你旁邊那個人忍耐你.
我老婆忍耐我.
我不會說我忍耐我自己.
這些一定是對人.
你做基督徒不是說.
我們可以將信仰變成了一種個人化或內化的事.
以為信耶穌就是回到自己的房間裡.
關上門 蓋上被子.
然後躲在被窩裡祈禱 讀聖經.
我們做基督徒當然這些是需要的.
祈禱 讀經 靈修那些我經常強調.
但是是不是做基督徒就是唯一這件事.
躲在房裡修煉就是.
不是的 我們要對外的.
我們要對人的.
我們要在這個世界裡生活的.
我經常舉一個例子就是.
當你要愛人如己的時候.
你一定要有個人在.
是不是.
愛人如己的時候.

$^{681}$你不可以說我也是人.
你愛人如己的話就一定要有一個對象.
你才能夠實踐基督的信仰.
你在一個荒島裡只有你自己一個人.
你是實踐不了這個信仰的.
我講回 又記得諾斯底主義.
諾斯底主義講到靈欲二分.
他認為世上所有東西都是邪惡的.
靈性才是最高峰的.
你的錢是邪惡的.
你穿的衣服是邪惡的.
你坐的椅子也是邪惡的.
既然這個世界這麼邪惡那怎麼辦.
不要理它.
不要理它 躲起來.
自己自成一角.
我們基督徒聚在一起就行了.
我們大家一起做靈性的操練就行了.
我們等上天堂就行了.
某程度上其實又是有些基督徒是這樣.
我好像港獨有講過.
我認識一個肢體.
他說從來都不看報紙的.
看什麼書 看屬靈書和聖經.
聽歌只聽詩歌.
這個世界好像所有東西都和他無關.
其實又是有多少洛斯底主義的味道.
這個世界是邪惡的.
我不是否定這個世界是邪惡的.
很多的惡行.
但是上帝依然將我們放在這個世界裡面.
我們在這個世界裡面有一個責任.
就是我們穿上身人進入這個世界.
讓人見到我們是不同的.
我們除了讀經 靈修 祈禱 傳福音之外.
我們還要有一些行為.
是讓人見到我們是這個世界的一份子.
我們的耶穌基督是怎樣.
所以我們會做很多的工作.
譬如我之前都帶過人去扶貧.

$^{721}$去探訪一些基層的家庭.
或者是如果你看到這個社會上有些不公義的東西.
其實我們有責任去出聲的.
不是說我們關上教會門在裡面聚會就OK.
我們如果又用移民做例子.
如果你將來知道你要上天國.
你要去到一個新天新地的時候.
今天你怎樣去應付你的生活呢.
剛才說過我不會做些什麼.
不會買一個新的80吋電視.
我不會買一個新的鋼琴.
但是相對地我什麼都不理.
你就快要搬離開錦繡的房子了.
發現花園有條蛇.
你理不理它.
當然要理.
你不可以不理.
你看到有老鼠.
當然要理.
不可以不理.
沖沖涼涼地.
突然之間浴室的天花板整塊石屎掉下來.
你會不會不理.
不可以不理的.
甚至你會不會說我遲些移民去英國.
然後走出去的時候.
看到有個婆婆掉進湖裡.
你救不救她呢.
當然救了.
你不會說我都移民英國了.
Sorry.
你看看有沒有誰還在那裡住.
救一救你.
不會的.
應該做的事情就要做.
應該做的.
你現在生活裡面要面對的.
你要做的.
你的工作也好.
你在讀書也好.

$^{761}$你作為別人的子女.
或者作為別人的父母.
你應該要盡的責任.
在這一刻你應該要做.
我們穿上新人是一種實實在在.
讓人看到的見證.
你是透過你的生活.
透過你的生命.
去影響這個世界.
讓人看到你.
就看到基督.
這個過程是有風險的.
就好像我剛才說伊朗那對男女.
他們真的犧牲了自己.
去讓人看到這個國家的問題.
他們讓全世界都看到.
一個這麼荒謬的新聞.
是有問題的這件事.
同樣地.
我們有時候可能做一些很傻的事.
但是讓這個世界看到.
這個世界有問題.
耶穌基督才是正確的.
才是這個世界的出路.
我們用行動去彰顯自己的價值觀.
今天的講題是天與地之間.
天地之間.
其實我們這個位置是很尷尬的.
是嗎?很尷尬的.
剛才我教主學的時候.
有個弟兄都問我.
他說一方面我們又想祈求耶穌基督快點來.
但是耶穌基督回來的時候.
有很多人就沒有機會得救.
或者是世界末日的時候.
災難就臨到.
其實真的有一種這樣的張力在裡面.
我們是很尷尬的.
一方面我們在等天國的來臨.
但是另一方面.

$^{801}$我們要在這麼邪惡的世界裡面生活.
但是正正就是耶穌要我們在這麼尷尬的位置.
這個天與地的合縫中間.
我們要成為祂的代表.
在裡面去做一些事.
我們不可以完全不理地上的事.
但是我們也不可以將自己完全投放在地上.
到底比例是多少呢?.
我就不懂回答你.
因為這個是你和上帝之間去交代的事.
你和上帝之間.
你認為你花多少時間去為祂作見證.
是你對得起上帝的.
那你就去做吧.
不可以不理地上的事.
也不可以不理天上的事.
最後想提醒大家一件事.
知不知道最後我們唱的那首詩歌叫什麼詩?.
使命詩.
通常叫使命詩.
那為什麼是使命詩呢?.
為什麼是使命詩呢?有沒有想過?.
有想過的就舉手吧.
看看你夠不夠膽.
你是帶著使命離開這個禮堂的.
在這裡聚會聽到得著偉仰.
你是帶著使命離開這個禮堂.
成為這個世界的炎和光.
出去的時候要讓人見到你.
就見到耶穌基督.
所以出去的時候留心自己的見證.
譬如去到會所不要和別人搶位.
留意一下.
去到大吉樂不要和別人鬆爭.
我先排隊.
在你的日常生活裡你要做好你的見證.
用愛去愛這個世界.
用愛去彰顯耶穌基督.
我們同心圖告.
讓我們有機會去服侍你.

$^{841}$其實能夠為你作見證是很光榮的一件事.
能夠為你去受苦也是很光榮的事.
雖然我們可能離受苦還是很遠.
但是我們都期望我們可以在日常的生活.
很普通的生活裡都可以為你做一些見證.
可以讓我們身邊的朋友,同事,親人.
見到我們的時候就認出我們是基督徒.
從而彰顯到你那份榮耀.
求你去保守.
帶領我們祈禱交託奉耶穌基督的名祈求.
\newpage



\section{歌羅西書 3:7-14}
\label{sec:W3jetg3vhiA}
\textbf{2023-03-19 早堂《信主之後》歌羅西書3\_7-14 講員:楊明麗姑娘}
\newline
\newline
連結: \href{https://youtube.com/watch?v=W3jetg3vhiA}{\texttt{https://youtube.com/watch?v=W3jetg3vhiA}} ~~~~ 語音日期: 2023-03-21
\newline
\newline
\hyperref[sec:BtpgDcAKeBQ]{\small{< < < PREV SERMON < < <}}
~
\hyperref[sec:index_chronic]{\small{[返順時目]}}
~
\hyperref[sec:index_scriptual]{\small{[返順卷目]}}
~
\hyperref[sec:wPsCGHG1pLA]{\small{> > > NEXT SERMON > > >}}
\newline
\newline
歌羅西書 3:7-14
\newline
\begin{longtable}{cl}
\hline
\hline
章節 & 經文 (和合本修訂版)\\
\hline
3:7 & \begin{tabularx}{0.7\textwidth}{X} 當你們在這些事中活著的時候，你們的行為也曾是這樣的。 \end{tabularx} \\ \\ \relax
3:8 & \begin{tabularx}{0.7\textwidth}{X} 但現在你們要棄絕這一切的事，就是惱恨、憤怒、惡毒、毀謗和口中污穢的言語。 \end{tabularx} \\ \\ \relax
3:9 & \begin{tabularx}{0.7\textwidth}{X} 不要彼此說謊，因為你們已經脫去舊人和舊人的行為， \end{tabularx} \\ \\ \relax
3:10 & \begin{tabularx}{0.7\textwidth}{X} 穿上了新人，這新人照著造他的主的形像在知識上不斷地更新。 \end{tabularx} \\ \\ \relax
3:11 & \begin{tabularx}{0.7\textwidth}{X} 在這事上並不分希臘人和猶太人，受割禮的和未受割禮的，未開化的人、西古提人、為奴的、自主的；惟獨基督是一切，又在一切之內。 \end{tabularx} \\ \\ \relax
3:12 & \begin{tabularx}{0.7\textwidth}{X} 所以，你們既是神的選民，聖潔、蒙愛的人，要穿上憐憫、恩慈、謙虛、溫柔和忍耐。 \end{tabularx} \\ \\ \relax
3:13 & \begin{tabularx}{0.7\textwidth}{X} 倘若這人與那人有嫌隙，總要彼此容忍，彼此饒恕；主怎樣饒恕了你們，你們也要怎樣饒恕人。 \end{tabularx} \\ \\ \relax
3:14 & \begin{tabularx}{0.7\textwidth}{X} 除此以外，還要穿上愛心，因為愛是貫通全德的。 \end{tabularx} \\ \\
[1ex]
\hline
\hline
\end{longtable}
$^{1}$大家早晨.
大家很久都沒有這樣的機會.
大家真的揭開了罪免層.
大家可以很真心地彼此相望.
互相祝福的時候.
很感恩的 今天大家可以一起見到面.
願主也賜福給我們.
願主賜福給我們.
讓我們無論聽的說的都一起蒙福.
我們一起禱告 好嗎.
親愛的天父上帝 我們來到你的面前.
我們感恩我們可以一起見面.
我們又求你紀念我們.
主人我們很需要你的話語 你的慰養.
求主幫助我們.
讓我們無論是說的聽的.
都一起在你的面前得到你的看顧.
得到你的話語 滋養我們 建立我們.
求主幫助.
我們祈望你 逢主耶穌基督的聖名 阿們.
弟兄姊妹今天跟大家講一個題目.
叫做信主之後.
可能你說30年前了.
可能你信了主很久了.
不過我們又重溫一下.
可能有一些弟兄姊妹.
你是這兩三年 或者去年.
洗了很久.
我們又看看信主之後我們要怎樣.
因為我們講的經文是.
歌羅西書第三章七到十四節.
但我又想跟大家看前面一點.
前面一點是第一節.
第三章的第一節說.
所以既然你們已經與基督一起復活.
就要求上面的事.
那裡有基督坐在神的右邊.
你們要思考上面的事.
不要思考地上的事.
一開頭第三章就講一個所以.

$^{41}$第三章有三個所以.
一個在第一節.
一個在第五節.
一個在第十二節.
一路講所以.
那因為甚麼呢.
有因為才有所以.
所以我們就不要想地上的事.
要想天上的事 為甚麼呢.
原來因為我們已經有一個新的生命.
這是寫給信主的人.
信主之後.
我們與基督已經有一個復活的生命.
神這麼早講這些.
我都未甚麼甚麼.
不過由我們信主那一刻.
我們就有一個新的生命.
舊的生命.
我們未信主之前的情況.
就好像一個完結.
畫了一個句號.
現在我們有一個新的開始.
這個新的開始就好像一個復活了新的生命.
這個新的生命.
是在基督裡面.
與基督一起復活.
不是.
因為基督說我們完成救恩.
所以信了主之後.
我們得著他的生命.
就好像與他一起復活.
剛剛受洗不久.
你很清楚.
上完浸禮課程.
由一課講到浸禮的時候.
是與基督.
好像同死同復活.
所以浸禮就是完全浸在水裡面.
然後起來.
代表我們好像死而復活一樣.

$^{81}$我們有一個新的生命.
所以.
所以要想上面的事.
不要想地上的事.
地上的事是甚麼呢.
原來講地上的事.
我們不是很想想.
污穢淫亂邪情.
惡慾貪婪.
拜偶像.
大家沒有做.
好像這些事離我們很遙遠.
這樣就好了.
因為這些事.
令到神很憤怒.
神會很憤怒.
做這樣的事情的人.
我們沒有做這些事.
但以前.
未信之前可能有些事.
我們都做了.
所以現在我們要想一想.
多想一想上面的事.
上面的事.
是講我們已經是一個復活的人.
我們的心思意念.
應該是復活了.
有更新.
與以前不同.
復活了的生命與以前不同.
但我們又會覺得.
我還沒有去到那個地步.
那個地步可能去到天父的那裡的時候.
第三第四節.
其實告訴我們一件事.
告訴我們一件事.
我們一個新的生命.
是與基督一起.
隱藏在天父的那裡.
基督都隱藏.

$^{121}$基督一直存在.
我們都奉祂的名禱告.
沒錯.
其實是存在的.
不過大家還未見到祂.
我們這個復活的生命.
都是存在的.
不是去到我們離開這個世界的時候才出現的.
我們復活的生命.
永恆的生命.
是由我們信著的那一刻.
已經得到的.
不過人家又不是很覺.
其實是藏在天父那裡.
藏在那裡的意思是有一天會顯現出來.
現在暫時放在那裡.
收起來.
遲些就會拿出來.
為什麼要收起來呢.
你看看你的手.
你兩卡的鑽石戒指就在你的手上.
或者在你的頸上.
但是你把家裡的東西.
你不會帶著它出來養的.
你一定要收起來.
最好的東西.
你會收起來的.
到需要的時候.
適當的時候.
才拿出來.
我們永恆的生命都是很重要的.
這個很重要的永恆的生命.
暫時我們就收在天父那裡.
不過總有些東西人是看到的.
當你有很多張地契的時候.
你都會有一些東西在手上閃一閃.
在耳朵上閃一閃.
我們大家看到的.
在你的衣服上閃一閃.
我們看得到.

$^{161}$很欣賞.
我們的生命雖然說永恆的生命.
我們未完全將它顯露出來.
但是這個生命的一些本質.
是逐漸在我們的生命當中.
慢慢每一天表達出來.
就好像保羅.
一直在鼓勵我們要做一件事.
就是信主之後.
我們要心意更新而變化.
我們的生命是逐漸逐漸地有些改變.
這個逐漸逐漸有些改變的生命.
一直表現出來的時候.
一些我們覺得不好的東西.
我們就會不要.
我昨天整理今天要穿什麼衣服的時候.
拿了件衣服出來.
然後發覺.
為什麼把衣服遮起來.
拿出來會有這麼多顆粒呢.
一件衣服再漂亮也好.
有顆粒就不會穿出街.
特別不會穿出來崇拜.
更加不會穿出來講道.
因為覺得最漂亮的衣服.
最新的衣服.
總是要穿回教會.
為神而穿的.
這個是我自己的想法和做法.
最新的最好的.
一定是為主而去用的.
一件有顆粒的衣服.
無論我喜不喜歡它.
還是要讓它休息一下.
在其他場合才穿.
我們舊有的生命.
有一些東西是有瑕疵的.
那些是我們的性格行為.
有一些瑕疵的地方.
我們就要放下它.

$^{201}$不會穿著它出來.
拿著這些舊有的有瑕疵的行為去見人.
然後告訴別人.
我是基督徒.
這些行為表現出來.
我們會覺得不是很恰當.
不是很配合一個基督徒.
所以這裡經文也告訴我們.
當我們還在這些事當中活著的時候.
當我們還是一個信主之前的境況的時候.
我們的行為都會有一些令人不舒服的地方.
但現在我們應該改變過來.
因為我們已經去到信主之後.
在這個情況下.
我們就放下了這些有瑕疵的行為.
什麼是有瑕疵的行為呢.
惱恨 憤怒 惡毒 誹謗 污穢的言語.
我不說的.
後面一句.
不要彼此說謊.
我想問一下大家.
有沒有發過脾氣.
我不是說今天.
今天你可能真的沒有發過脾氣.
這個月已經過了18天了.
今天是第19天.
這個月有沒有發過脾氣.
如果沒有的話你還是一個很好的人.
今年呢.
今年有沒有發過脾氣.
我們會的.
發脾氣的時候會說什麼.
說一些你平時不想說的東西.
發完脾氣之後.
你也要找一些解釋去解釋.
為什麼我會說這些東西.
因為接受不到自己說的那些東西.
是的.
有些東西是我們不想說的.
但在生氣的時候.

$^{241}$在我們的情緒上來的時候.
不想說的東西就說出來.
為了解釋不想說的東西.
就再說幾句更加不想說的話.
是的.
我們已經有一個新的生命.
我們知道這些話不好.
我們不想說.
我們已經懂得不想說.
但聖經就提醒我們.
要除掉這些言語.
這些行為.
或者這些習慣.
要放下它.
我們嘗試做一個沒有這些.
我們不想說的言語.
有一些更加美麗的言語.
的基督徒.
更加有一些話.
可能大家又不想說.
有時候又不知道怎樣去避免說.
就是說謊.
原來說一些不真實的話.
都不是那麼好.
這樣可能比發脾氣更加難.
但是我們又要去學習.
學習不要做這些事.
好像放下一件舊的衣服一樣.
我們要除掉這些行為.
求主去幫助我們.
當我們在那個時候.
我們會做這些事.
現在就要除掉這些行為.
因為我們已經過了那個.
用舊陣時的自己的時候.
讓我們能夠靠著主的緣故.
我們有一個新的行為.
新的生命表達出來.
這個生命.
開始發覺了.

$^{281}$這個生命我們擁有了.
但是要表達出來.
不發脾氣.
有些事不說謊.
原來不是很容易.
每樣事都是隱讓.
關心.
很有愛心.
不是不想去做.
是有些困難.
好像我們這個生命.
真的隱藏在天府.
未能完全表達出來.
求主幫助我們.
幫助我們能夠從神的話語.
和從神的.
和神的關係更加密切的時候.
這些舊有的習慣和語言.
我們可以放下.
就好像放下了另一種語言.
讓我們學習一種新的語言.
讓我們的語言能夠多了愛心.
多了關心.
多了彼此的見諒.
這種語言就像.
穿上一件新的衣服一樣.
這個新人的生命要表現出來.
我們穿上一件新的衣服.
穿上一件新的衣服.
能夠表達出.
我們已經有一個新人的生命.
這個生命是真的擁有的.
因為在基督裡面.
我們得著了這個救恩.
神的形象再次表現出來.
當然我們要去學習聖經.
所以我們在知識上要不斷地更新.
那些人要做這些事情呢.
哪些人要做這些事情呢.
他說是受國禮和未受國禮的人.

$^{321}$即使是猶太人.
信了主都應該是這樣的.
有沒有文化的人都要是這樣的.
為奴的自主的都要是這樣的.
哇 這是一個很大的挑戰.
不知道在座的大家有沒有同事.
或者是上司下屬都一起是基督徒的呢.
大家是怎樣相處的呢.
他講到為奴的自主的.
在當時的哥羅西信徒當中.
他們就是這樣.
有些是奴隸.
有些是這些奴隸的主人.
他們都回到同一個教會.
受同樣的教導和牧養.
接著就說你們要學習彼此尊重.
彼此關心愛護.
不但這樣.
不要罵人.
主人不要罵人.
奴僕聽到當然很開心.
你不要罵我.
你是基督徒.
不過千萬不要這樣說我.
你是奴隸.
到時他不罵你.
不過會賣掉你的.
但保羅就很忠正地跟他們說.
是真的.
無論你是做主人的.
或者是做奴隸的.
或者是奴隸當中你是管理另一些奴隸的.
對不起.
信了主之後.
大家都是神的兒女.
大家都是基督徒.
大家都有永恆的生命.
有新的形象.
在你的下屬.
你的奴隸面前的時候.

$^{361}$你不要說一些傷害他的說話.
從另一方面看.
你作為奴僕的時候.
你口不說出來.
心也不要這樣想.
不要不在主人面前的時候.
在背後就說他什麼不好.
去咒詛他.
跟你一樣.
都是弟兄.
都是姐妹.
都是神的兒女.
這一說.
好像保羅就在說.
誰誰誰.
每個人都要做的.
到真的要做的時候.
在這群哥羅西的信徒當中.
這是一個很大的挑戰.
生活在同一間屋當中.
一個是吩咐人.
一個是被吩咐的.
一個應該被尊重的.
一個是奴隸.
大家要受同樣的教導.
不要罵人.
不要發脾氣.
很不容易.
很不容易的.
我們自己雖然不是.
在一個這麼大型的家庭當中生活.
但是在家中.
有兩個人都可能會發脾氣.
不知道有一個人只有自己的時候.
會不會發脾氣.
有兩個人的時候.
都有機會發脾氣.
只要有兩條鎖匙就會響.
有三個有四個.
兒子怎樣這樣.

$^{401}$是你的責任.
是甚麼.
會有不同的說話.
我們會表達出來.
你最著緊那個.
這麼著緊他.
他應該很完美.
跟著對他會有很多的要求.
還有言語上的要求.
是會有要求.
會有期望.
最期望他有些東西不變.
不是完全不變.
有些東西你想他變.
你叫他都不變.
嘩又是這樣.
上次都說了你.
今次又是這樣.
上年都說了你.
十年前都說了你.
勞氣.
有些東西你又想他不變.
三十年前他追求你的時候.
等你兩個小時.
遲到都沒有問題.
很早而已.
我沒事沒事.
沒事做.
你站得住.
你走走就站得住.
現在約了他出門.
喂 等了你十五分鐘.
我還未走得.
為甚麼換了一個人.
是喔.
我們有些東西又想他改.
有些東西我們又想他不改的時候.
我們其實應該怎樣呢.
我們其實原來應該有一種.
恩慈的心.

$^{441}$憐憫的心.
謙虛的態度.
彼此忍耐.
再忍兩個小時.
不一定.
不過忍耐他不肯等兩個小時.
然後忍耐他會遲十五分鐘.
有些東西我們進到一個階段的時候.
進入一個新的生命.
原來大家有些東西要不同了.
那樣東西就是.
我們不再用地面上的眼光.
去彼此接觸.
是用天上的眼光.
天上的眼光那種包容和愛.
神體我們身邊的人.
你很重視那個人.
是一個蒙愛.
需要被愛的人.
在全宇宙幾十億人當中.
神揀選了在你身邊的人.
他坐在你身邊.
他和你一起吃飯.
或者他煮飯給你吃.
大家有不同的性情.
你喜歡的口味不是他喜歡的.
我家裡很特別.
我又喜歡吃少許辣少許酸的東西.
我先生兩樣都不行.
兩樣都不行.
吃了會出事.
我弟弟喜歡吃很辣的東西.
很酸的東西.
我弟弟也是兩樣都不行.
所以我們四個出去吃飯.
好像叫兩類的東西.
一種是我和我弟弟吃得很開心.
終於有盤子陪我吃了.
另外兩種就是我們吃清淡一點的.
那些不關我們事.

$^{481}$但是就是不合味道的人才一起.
不同恩賜的人也一起.
除了施班元.
因為要彼此扶持.
彼此幫助.
彼此關心.
因為他需要幫助的就是你得的那些.
你需要幫助的就是他得的那些.
大家就是要這樣學習去相處.
除了飲食上的味道之外.
可能有很多不同的地方.
我初訓的時候有一對牧師夫婦.
師母沒有車牌.
牧師開車很好的很厲害.
牧師沒有方向感.
師母很有方向感.
是不是很風趣.
神就是這麼風趣.
永遠就是牧師開車.
師母坐在旁邊.
這個配搭.
令我們覺得.
神就是這麼特別.
我們的生命裡面就是這樣.
在你的家庭裡是這樣的.
在你的兒女當中是這樣的.
甚至你請過的工人也是這樣.
你有要求的事情他做不到.
不過他做得好的事情.
你不一定知道.
或者發掘了出來.
你會生活得很開心.
神給我們.
有這個恩典.
我們不要用地面的眼光去看.
我覺得你應該怎樣.
但是用神的恩典去看.
祂怎樣給我們相遇.
但我為什麼經常遇到一些.
不好脾氣的人在我身邊.

$^{521}$如果我們用上面的角度去看.
跟我們用下面的角度去看是不同的.
我們用下面的角度去看.
可能會覺得很不開心.
你時時都遇到不開心的事情.
不開心的人在你身邊.
無論是同事也好.
或者是鄰居也好.
都會令你覺得不開心.
但是從神的角度去看.
祂可能在鍛鍊我們.
讓我們學習心意今生而變化.
讓我們每天明白多一點.
什麼叫做忍耐.
讓我們明白多一點.
什麼叫做憐憫.
什麼叫做愛心.
其實大家都很可愛的時候.
愛心是不用練習的.
因為都沒有地方讓你發脾氣.
但是當大家有一些不同的時候.
愛心就要去操練.
我們要操練愛心.
不是因為身邊每一個人都很有愛心.
很可愛.
是因為我們生活在一個邪惡的世代.
太多事情是我們估計不到的.
我相信這兩個月發生的一段.
我們都不想提的新聞.
大家深刻印象.
不說都知道那件事.
甚至乎我想有一段很長的時間.
大家都不想喝青紅蘿蔔湯.
因為我們都覺得很邪惡.
世界就是這麼邪惡.
邪惡到不可信地邪惡.
但是神就說我們從上面的角度看下去.
是有人邪惡.
有人做了很多神不高興的事情.
因為他們還沒有認識神.

$^{561}$但是當我們認識了主之後.
我們要學習除去這些邪惡的行為.
去把美好穿上新的行為.
新的生命上身去表達出來.
這些邪惡的行為提醒我們.
我們有一個更美麗的家鄉在天上面.
我們由那個角度去看我們的生命.
看我們的社會.
我們學習這樣去相處.
這是神給我們的機會.
信主之後.
要怎樣呢?.
活出愛.
誰要做?.
旁邊那個.
信主之後.
你要活出愛心.
你已經洗禮了.
你信主久了.
或者你剛剛洗完禮.
你要學.
其實難道還沒有信主的人不用嗎?.
不過做不到.
他們都應該要的.
好的行為有言語.
有愛心的行為.
有愛心的言語.
信主和未信主的人都應該有的.
神做人的時候.
祂不想人被仇視.
大家有怨恨.
不過人因為犯罪的緣故.
一路傾向.
做一個隨心所欲.
我想說什麼就說什麼.
於是就做了很多不合神的心意.
和彼此傷害的事.
我們不是特別要做有愛心的人.
是神想每一個人都是一個有愛心的人.
只不過不是每個人都做得到.

$^{601}$我們有這個機會.
我們要去學去做.
因為我們有一個新的生命.
誰能夠做得到?.
有新生命的人才能夠做得到.
有新生命的人.
因為有一個新生命的動力.
這個力量是由神而來.
我們去親近神.
祈求神幫助我們.
天父說這個人我很難面對他.
求主幫我.
賜力量給我.
天父說我不滿意我自己的說話.
求你幫助我.
讓我的說話能夠更加有愛心.
更加得體.
讓人的生命能夠得到幫助.
得到鼓勵.
讓人在我的言語當中.
能夠聽到神的存在.
能夠見證到你.
誰能夠活出神的愛呢?.
一個敬畏神的人.
怎樣去活出來?.
怎樣去活出神的愛來呢?.
我們剛才讀了《愛的詩篇》.
我們唱了詩歌.
但我們怎樣去生活出這種神的愛呢?.
這段經文說.
當你們還生活在那個時候.
但現在你們已經不是那個時候.
我們不再在信主之前.
我們已經去了一個信主之後的時候.
既然去了一個信主之後的時候.
我們有那個力量.
那未信的人又說.
他要生活得有愛心.
也是的.
難道你認識的每一個未信主的人.

$^{641}$或者有其他宗教信仰的人.
他都是毫無愛心的嗎?.
不是的.
有些人.
有些人你覺得你跟他傳福音很有困難.
你為什麼?.
因為他很好人.
他都不覺得自己有罪.
他好人都不覺得自己有罪.
是的.
真的有些人都是很有愛心.
但是我們要做得有愛心.
跟別人本身很好的性情有什麼分別呢?.
最低限度有一個分別.
我們要做得很有愛心.
要為神而去做.
無論我是不是一個有愛心的人.
為神的緣故.
我應該有愛心.
這是我穿上新人的意思.
這個新人是屬於神的.
這個新人是神給我的.
因為我是屬於神的.
是神給我的.
所以我應該要有愛心.
我還沒做到.
我要去學.
這個是神的生命在我裡面.
我是屬於神的.
所以我是一個基督徒的時候.
我要這樣去做.
一個有愛心的人.
是.
這個就是神給我們的恩典.
給我們的教導.
什麼人要有愛心呢?.
約翰福音告訴我們.
主耶穌基督說.
我賜給你們一條新的命令.
就是要我們彼此相愛.

$^{681}$愛也要命令嗎?.
給我們這條命令有沒有愛呢?.
給我們這麼艱難的命令.
是不是很有愛呢?.
是.
他說我們死在十字架上.
剛才詩歌告訴我們.
我們沾染他的腳.
他的手.
他的胸的時候.
看到他的釘痕.
看到刀刺進他的痕的時候.
我們很看得到.
主耶穌基督很愛我們.
他很有愛心地.
為我們不單離開天上.
來到地上.
地上也很好享受.
什麼啊.
在天上他又不用走來走去.
那時候沒有車.
哇.
他要去不同的地方.
他要去傳福音.
告訴別人.
要人應該怎樣生活.
他去醫病.
去趕鬼.
去建立門徒.
做很多事.
然後還要為我們.
被釘十字架.
為我們放下他.
尊貴的身份.
他是神的身份.
不是他死了之後.
他會復活的.
不只是這件事.
這麼尊貴的一個神.
要去到一個地步.

$^{721}$被釘十字架.
這麼恥辱.
是.
主耶穌基督就是這樣.
他放下了自己是神.
被釘十字架.
他也可以恥辱到這樣的地步.
去愛我們.
我們為他表達一些愛心.
在對我們.
不是這麼有愛心的人面前.
我們不要看那種.
被沒有愛心對待的恥辱.
去欣賞主耶穌基督.
為我們負上.
承擔了恥辱的.
大愛的緣故.
我們做一個有愛心的人.
做一個有愛心的人.
在言語.
在行為上.
好像穿了一件新衣服.
舊的行為.
不相襯.
讓我們有新的行為.
表達出來.
我們這樣做.
是讓基督賜給我們的.
愛的身份.
榮耀的身份.
更加表現出來.
這個身份.
我們需要一樣東西.
一樣東西.
讓我們能夠更加表達出來.
就是那種平安的心.
和平的心.
和平和平安.
其實在聖經裡是同一個字.
意思就是我們沒有一種爭吵.

$^{761}$怒恨.
是平和地去看一件事.
這個平安的心.
這種和平的心.
會令到我們怎樣呢.
這種平安的心.
令到我們遇到困境的時候.
內心安寧.
無論是疾病的困境.
或者現在是經濟的困境.
面對不同的困境的時候.
我們都欣著有這種平安.
我們可以安然去面對.
另一方面.
就是因為這種平安的心.
我們看的角度不同了.
我們看事物的角度不同了的時候.
於是乎我們.
從上面看下來的角度.
我們會明白神.
祂面對這些問題.
這些恥辱言語的時候.
祂想我們怎樣.
祂想我們仍然是以平和的心去面對.
仍然是去忍耐和包容.
這樣表達的人.
正如主耶穌基督在十格上的時候.
釘十字架的人和在下面的人.
正在恥笑祂.
兩個釘十字架的人.
有一個都是在恥笑祂.
都一起釘十字架了.
還要去恥笑耶穌.
但另一個不是.
另一個得到救恩.
去拯救耶穌 拯救祂.
我們要有一種.
我們是得夢上帝的恩典.
我們從上面看下來的時候.
我們自己都是一個蒙恩的人.

$^{801}$身邊有些人可能還沒有蒙恩.
有些人蒙恩之後.
還未懂得時時穿著蒙恩的新衣服.
所以在言語行為上.
還未能表現出自己有新生命.
我們有時候也是.
求主幫助我們.
既然是這樣.
我們就一起學習.
去過一個從上面看下來的人生.
有愛心的人生.
為了什麼呢.
不是為了我們自己會過得舒服一點.
少一點罪疚感.
而是神喜歡.
我們是為神去做.
包容 忍耐 饒恕.
不是因為我們有沒有能力.
而是因為神喜歡.
你穿衣服出來.
也希望別人讚賞你.
我們穿上這件新人的衣服.
天父讚賞 欣賞.
和加力量.
讓我們可以長久地穿著.
第十七節說.
我們無論做什麼事.
或說話 或行事.
都要奉主耶穌的名.
求主讓我們有這個恩典.
常常都是這樣做.
其實第三章一到十七節.
就告訴我們.
因著舊時.
不久之前.
我們是一個罪人的時候.
在中段我們信了主之後.
現在我們應該有信主之後的生命.
有信主之後的生命的時候.
第一件事我們是為神而做.

$^{841}$每一件事我們都為神而做.
最後一件事我們都為神而做.
為神而做.
我們才能與神有一個美好的關係.
我們靠著自己做不到.
一個和善的人.
唯有神給力量我們.
我們藉著禱告.
困難的地方告訴天父.
求主幫助我們.
我們藉著讀聖經.
明白多些神其實真的想我們怎樣呢.
我們有不明白的.
可以求問.
又可以上一些課程.
無論是星期四晚的一些查經課程.
或者是一些神學院.
現在都很多的網上的課程.
我們如果不知道上什麼課程.
問一問身邊的傳道人牧師.
他一定有一些覺得適合你.
接著之後.
過了十七節之後.
接著的經文就是說.
我們怎樣與神建立關係.
和彼此怎樣互相鼓勵.
然後二十二節之後就說.
在家裡你怎樣做.
回去慢慢看.
一定有些是令你覺得.
驚喜呀 我有件新的衣服要穿著.
原來在家裡我應該是這樣的.
最近我學了一樣東西.
就是家庭當中大家都有使命.
做父母有父母的使命.
不過使命不是長憂九十九.
是建立兒女而已.
帶領他歸向神.
做夫妻都有使命.
有使命不是讓你只是去嘆對方的愛.

$^{881}$有使命你去愛對方.
建立對方.
讓對方的生命在神的面前更加完整.
怎樣可以這樣做呢.
是你對他的尊重.
他一生面對很多批評.
不用多於我們去批評對方.
我們需要的是尊重他 欣賞他.
這世人稱讚他的人很少.
他的父母未必懂得稱讚他.
正如你的父母也未必懂得稱讚你.
身為配偶的時候我們學一下這件事.
這是我今年或最近一年後.
最近在想的時候.
我要更加加強這件事.
學習去尊重 欣賞 讚賞我丈夫.
在這個世界生活.
很少人會讚賞你的丈夫 你的太太.
這個責任留給你.
你令到他走每一步出來.
都更加像一個基督徒.
因為基督的生命在他身中.
而他很看到.
由你的說話中得到幫助 得到鼓勵.
大衛遇到一個人叫拿伯.
大衛保護了拿伯和他所有的財物.
拿伯侮辱了大衛.
引來大衛想去殺死拿伯.
但是拿伯的妻子很妥善地處理了這件事.
令到大衛沒有殺死拿伯.
不過後來大衛就娶了拿伯的妻子.
阿比蓋.
如果大衛能夠從上面的角度看下去.
他就不會殺死拿伯.
阿比蓋就是提醒了大衛.
他從神的角度看下去這件事的時候.
大衛覺得自己真是愚蠢.
大衛也說了拿伯是愚蠢.
如果我還要去回應一個愚蠢的人.
去殺一個愚蠢的人.

$^{921}$令到我的手尖鬱.
所以他就沒有去殺拿伯.
阿比蓋其實他救了拿伯.
雖然最後他離開了拿伯.
不過當時拿伯已經死了.
因為他被自己的行為可能引起的結果嚇到.
可能那種叫心臟病發.
有時間大家看看聖經.
哎呀 是呀.
有些事情原來換個角度去看.
自己都會覺得自己做的事情是愚蠢的.
於是就不會去做.
另外有一個比喻.
是主耶穌基督說的比喻.
浪子回頭.
浪子回頭當中.
父親給了很多東西給小孩.
大兒子一直都在他身邊.
大兒子覺得我這麼孝順.
對你這麼好.
你什麼都給小孩.
現在他一回來.
他已經拿了他應該有的一副身家.
他一回來你又再好好照顧他.
又不見你幫我慶祝一下.
只是兩兄弟.
一生人只是兩兄弟.
是不是要這樣計較呢.
好像耳掃和啞角都去計較.
是呀 要計較.
怎樣都要計較一天.
那句成語騙我們的.
人哪是計較一天的長短.
一天你就不計較.
一年你都不計較.
你計較多少.
你計較一世.
你為何要承認自己這麼辛苦一世呢.
放下它.
從神的角度看.

$^{961}$我們好好地建立自己的生命.
我們的生命是天上神去欣賞.
去賞賜我們的時候.
給我們地上所能爭到的.
多多 多好 多榮幸.
我們一起努力.
親愛的天父上帝.
我們來到你的庚前.
天父說我們很多事未必做得很好.
我們很多事需要你的恩典.
你的幫助.
我們不是不知道基督徒應該怎樣.
而是天父說.
要做到我們缺乏力量.
我們都很想由一個天上的角度.
由上面看下來.
這樣去處置.
但有時我們缺乏這種智慧.
求主幫助我們.
幫助我們.
讓我們能夠因著你的緣故.
因著你的教導.
我們能夠慢慢去明白多一些.
慢慢可以建立到愛心.
求主你紀念我們.
幫助我們.
讓我們有了信心可以加上特幸.
有了特幸可以加上知識.
有了知識可以加上節制.
有了節制可以加上忍耐.
有了忍耐可以加上虔敬.
然後我們可以有愛弟兄的心.
可以有愛眾人的心.
最後我們需要你逐步建立我們.
不是因為我們可以.
而是因為神你可以.
求你幫助我們.
紀念我們.
讓我們每一天得到提醒.
好好穿好這件新衣服.

$^{1001}$求主幫助.
我們仰望你奉主耶穌基督的聖名.
阿們.
\newpage



\section{馬太福音 25:1-13}
\label{sec:wPsCGHG1pLA}
\textbf{2023-03-19 加恩樓《心理預備,宏觀資訊的裝備,行動準備》馬太福音25\_1-13 講員:郭鴻標牧師}
\newline
\newline
連結: \href{https://youtube.com/watch?v=wPsCGHG1pLA}{\texttt{https://youtube.com/watch?v=wPsCGHG1pLA}} ~~~~ 語音日期: 2023-03-22
\newline
\newline
\hyperref[sec:W3jetg3vhiA]{\small{< < < PREV SERMON < < <}}
~
\hyperref[sec:index_chronic]{\small{[返順時目]}}
~
\hyperref[sec:index_scriptual]{\small{[返順卷目]}}
~
\hyperref[sec:rNmdtGHU_n0]{\small{> > > NEXT SERMON > > >}}
\newline
\newline
馬太福音 25:1-13
\newline
\begin{longtable}{cl}
\hline
\hline
章節 & 經文 (和合本修訂版)\\
\hline
25:1 & \begin{tabularx}{0.7\textwidth}{X} 「那時，天國好比十個童女拿著燈出去迎接新郎。 \end{tabularx} \\ \\ \relax
25:2 & \begin{tabularx}{0.7\textwidth}{X} 其中有五個是愚拙的，五個是聰明的。 \end{tabularx} \\ \\ \relax
25:3 & \begin{tabularx}{0.7\textwidth}{X} 愚拙的拿著燈，卻沒有帶油； \end{tabularx} \\ \\ \relax
25:4 & \begin{tabularx}{0.7\textwidth}{X} 聰明的拿著燈，又盛了油在器皿裡。 \end{tabularx} \\ \\ \relax
25:5 & \begin{tabularx}{0.7\textwidth}{X} 新郎遲延的時候，她們都打盹，睡著了。 \end{tabularx} \\ \\ \relax
25:6 & \begin{tabularx}{0.7\textwidth}{X} 半夜有人喊：『看，新郎來了，你們出來迎接他。』 \end{tabularx} \\ \\ \relax
25:7 & \begin{tabularx}{0.7\textwidth}{X} 那些童女就都起來挑亮她們的燈。 \end{tabularx} \\ \\ \relax
25:8 & \begin{tabularx}{0.7\textwidth}{X} 愚拙的對聰明的說：『請分點油給我們，因為我們的燈要滅了。』 \end{tabularx} \\ \\ \relax
25:9 & \begin{tabularx}{0.7\textwidth}{X} 聰明的回答：『恐怕不夠你我用的；你們還是自己到賣油的那裡去買吧。』 \end{tabularx} \\ \\ \relax
25:10 & \begin{tabularx}{0.7\textwidth}{X} 她們去買的時候，新郎到了。那預備好了的，與他進去共赴婚宴，門就關了。 \end{tabularx} \\ \\ \relax
25:11 & \begin{tabularx}{0.7\textwidth}{X} 其餘的童女隨後也來了，說：『主啊，主啊，給我們開門！』 \end{tabularx} \\ \\ \relax
25:12 & \begin{tabularx}{0.7\textwidth}{X} 他卻回答：『我實在告訴你們，我不認識你們。』 \end{tabularx} \\ \\ \relax
25:13 & \begin{tabularx}{0.7\textwidth}{X} 所以，你們要警醒，因為那日子，那時辰，你們不知道。」 \end{tabularx} \\ \\
[1ex]
\hline
\hline
\end{longtable}
$^{1}$我們一起同心祈禱.
我們親愛的天父上帝.
我們今天來到你的面前去敬拜你.
我們的心是要常常想起你的恩典.
也想起我們要靠著你才可以明白這個世界.
也明白我們怎樣在當中去自處.
怎樣去回應.
我們更加求主幫助我們在這個變化的環境裡面.
明白你的心意.
能夠去關心有很多找不到自己生存目標的人.
或者很多人他們感覺到生活都很大擔子.
很大壓力.
或者他們失去了喜樂.
失去了盼望.
求你的說話來成為我們的力量.
你的靈在我們當中運行.
幫助我們看到神給我們看到的異象.
我們這樣的祈禱是奉耶穌基督的名.
阿們.
讓我來到嘉欣樓和大家分享聖經的說話.
是很開心的事情.
有一次港道.
講完就提起.
在2015年有一套美國的電影.
名字叫The Big Shot.
在台灣翻譯成大賣兇.
在香港翻譯成孤主一擇.
這套電影得到第88屆奧斯卡金像獎最佳電影.
最佳導演等五個獎項提名.
就贏得最佳改編劇本獎.
我不是說跟隨今屆奧斯卡.
是說過去2015年的事情.
這個故事是改編自真人真事.
在2007至2008年環球金融海嘯期間.
有四位用專業手法去投資金融的財進.
去進行孤主一擇獲得利潤的事情.
當時在美國還未暴露次按危機的時候.
有一個基金經理叫Michael.
他通過大量數據來分析.
就想當時雖然好像很興旺的地產市場.

$^{41}$其實是隱伏了很多的危機.
那些人過度借貸.
所以隨時都會出現爆破的泡沫的情況.
由於他不顧其他金主的反對.
他就孤主一擇.
向多間投資銀行作出了一個信貸違約的對賭.
結果那些的行行家就說這個傻子.
有一個行行家叫Jarrett.
他知道這個消息之後.
那些做生意的人很聰明.
他就說我都不要怕吃虧了.
但他為了要減低風險.
在機緣巧合之下.
他就哄到另一個保險經理叫Mark.
來跟他合作.
另外有兩個年青人就做房地產的投資.
他在當時處處碰壁的情況下.
很意外的發現了一份內容.
關於房地產保沒的那些棄置了的計劃書.
這兩個人就說.
看來都應該大有可為.
於是就找隔壁的舊鄰居.
這個退休的銀行家Ben.
來出手幫助.
因為他們兩個年青人不是很夠資本.
希望以小博大.
這套電影是描述什麼呢.
美國普羅大眾是那種投資的心態.
具體的比喻就是說.
大家都是貪心的.
還有很慾望的.
上至銀行家.
下至證券評級員.
甚至那些分析員審批的人.
全部都是向錢看的.
可為大家合力的炮製了這一次的金融危機.
影片是反映了當年那種借貸氾濫的情況.
還有政界經濟合力的打造一個叫做.
你賺緊沒虧的美好假象.
讓那些民眾失去了危機感.

$^{81}$結果就損手爛腳.
銀行家和港大投資者.
他們深信不疑的一件事.
就是樓市是不會倒的.
但是他們實際上的假設是什麼呢.
就是誰不會去付那些按揭.
這是一個假設.
他們沒有對實際的情況進行客觀的分析.
Mark和他的團隊.
聽到樓市快要倒的這樣的預言.
起初都半信半疑的.
他只是想.
怎麼會沒有人知道這件事.
沒有人發現的呢.
他對這個說法作出質疑.
不過他們有一樣東西.
就是這個團隊去找找資料.
去訪問那些人.
向那些借貸人.
或者其他的證券評級員去問問他們.
你們怎麼看的.
經過小心求證.
他們都相信了樓市是有泡沫的.
就決定買入了那些信貸違約的調期.
來去進行從中的獲利.
當眾人以為一切如常的時候.
這幾個人就做了一些與眾不同的事情.
他們就在逆流而上的環境下.
就叫做早著先機.
這套電影是有很多條主線的.
不過我今天只是想用這個例子去說什麼呢.
是一個危機感.
和獨排眾議的看法.
回到經文.
我們可以請音響的弟兄.
把經文放出來.
就是馬太福音第25章第1到第13節.
說的十個童女的比喻.
我準備這次的講道的時候.
題目比較長一點.

$^{121}$叫做心理準備.
知識和技能的裝備.
和行動的準備.
其實是說三樣東西的準備.
講道自然會分三個部分.
第一就是心理上的準備.
第二就是宏觀資訊的裝備.
第三就是行動的準備.
首先我們看看心理的準備.
就是在馬太福音第25章第3節.
就說到有十個童女裡面有五個.
被稱為愚拙的.
他們是怎樣呢.
他們是拿著燈卻沒有帶油.
是什麼意思呢.
他們是知道要去接人的.
不過他們有帶燈.
那時候是黃昏的時候.
不過他沒有帶多一點油.
你明不明白.
他有帶燈.
燈裡面應該有油的.
在那裡著的.
不過他沒有帶多一個器皿去裝油.
經文就說這些童女出去接新郎.
其實去接人的時候.
要預備些什麼呢.
你要試一下.
不要在我們今天的時代去想事情.
耶穌的時代和現在的時代是不一樣的.
我們習慣交通工具都會準時的.
如果你坐西鐵都會照計準時的.
但是耶穌時代的交通工具可以不準時的.
可以在沙漠這樣走的時候怎麼辦呢.
就算你知道對方什麼時候出發日期.
他什麼時候到達和什麼時候到達.
大家都是只有一個估計.
不會這麼準的.
可能你還要去.
他說今天到達.

$^{161}$但是你要等的.
等多久都不知道.
經文形容其中有五個童女是愚蠢的.
有五個是聰明的.
愚蠢的五個是怎樣呢.
拿著燈但又不預備多一點油.
但是經文就沒有解釋為什麼會說他蠢.
你慢慢去揣摩的了.
然後第25章第4節.
聰明的是拿著燈.
又勝了油在器皿裡.
什麼意思呢.
他拿著燈.
燈裡面有油.
又有燈芯.
他穿著了.
但是他有另外一個瓶子.
或者一個瓶子.
載著油.
就是當現在他手上的燈用完油之後.
他可以添加下去.
就是這樣.
就是做多一件事帶多一件事.
那你帶多一件事.
那就被人叫聰明了.
是不是.
可以想像是什麼呢.
那就試一下做個例子.
我們每天出門都看看天文台.
今天氣溫是怎樣.
當今天涼一點.
那我只可以推論一下.
老人家當然是這樣.
帶回衣服.
晚上冷.
你當是這樣的意思.
你帶的你就聰明.
你不帶就不聰明.
就是這個意思.
在這個經文裡面.

$^{201}$就是在說有兩批不同的同類.
我們可以說愚蠢的同類.
他根本是沒有想過.
我帶的燈是會.
點完的時候新郎還沒到.
他沒有想過這件事.
沒有這個意識.
沒有這個心理準備.
沒有想過去找個瓶子.
裝另外一些油.
傍身的.
沒有想過這件事.
經文是提醒我們.
我們應該要有危機的意識.
我們要有心理準備.
如果事情有變的時候.
你有什麼計劃呢.
在一個一切如常的時代.
我們覺得很多事情是沒有問題的.
事情照樣運作.
正常情況下.
羅亞時代的人都是這樣想的.
所以看到羅亞在建方舟.
你會覺得他在做什麼.
羅亞你是不是太空閒.
你病了嗎.
你為什麼要建方舟.
這幾年的變化.
令我們開始加強了一些危機意識.
我們都會有心理準備.
有些事情可能會變差.
Worst Scenario這個詞彙.
我們經常都聽.
如果再壞的情況出現的時候.
我們有沒有Plan B.
我們經常都這樣說.
我問這個問題.
意思就是說.
你會不會有兩手準備.
在外在環境無法控制的情況下.

$^{241}$我想我們大家擺出防守陣容.
都很正常.
叫做保本 止蝕.
這個是首要任務.
肯定是這樣做的.
我們香港雖然不算好.
但是總比很多地方好.
大家同不同意嗎.
這件事我們真的很感恩.
值得我們去想的問題就是.
我們有沒有心理準備.
我們的假設.
是可能不能夠對應真實的情況.
我們其實要問自己.
你正在想的事情是不是真的對應實況.
或者我們的想法是否建立在錯誤的基礎上.
這個更加嚴重.
可能你從頭到尾都想錯了.
我自己很喜歡去想事情.
所以我知道很多人不想去想事情.
因為平時都很煩.
還要想事情就不太好.
不過居安思危是很重要的.
同樣的我們在黑暗裡面.
預備黎明的出現都是很重要.
一定是這樣.
特別就是說.
現在我們都不說疫情.
疫情這兩個字都已經淡化了.
我可以說形容現在我們的心境.
現在是誰鬥快彈起身.
你做生意的話.
鬥快彈起身你就知道怎樣去回應.
早著先機.
不是現在去想彈起身.
是幾個月前你已經部署下.
你今天怎樣彈起身.
剛剛這樣說.
同一個道理.
我們作為基督徒.

$^{281}$其實我們有沒有去調整我們的心態.
按照著情況來做好準備.
或者我們在神的角度裡面.
我們去服事.
我們有沒有這個思維呢.
當然有些人很著重實用.
他只是問敲的問題.
你講方法你先告訴我ABC.
我自己覺得方向和方法同樣都是重要的.
我們要知道我們去哪裡.
同樣我們要知道我們用什麼方法去達到我們想要去到的目標.
有人形容在疫情之後.
教會是出現大洗牌.
所以我不是駐堂的牧師.
我在神學院裡面.
但是我除了喜歡看書之外.
有時也有空.
那些牧者我問一下他們兩句.
問一下教會的情況.
所以也有一點路邊的消息.
我就不能夠什麼都講.
不過我舉一些令我覺得很重要.
我不能夠忽視的例子.
有一些叫超大型的教會.
他們叫Mega Church.
人數下跌.
非正式的統計.
其中有一間真的不講名.
是由五千人左右跌到二千人.
有一間大概由一萬人跌到四千人.
不講名是真的.
不開名.
去年我去其中一間教會的成人團體裡面.
他的訓練班去分享.
我發現他們除了有牧者之外.
是有一些弟兄姊妹是做導師.
有一位姊妹她是擔任超過一個團體的導師.
她已經分享弟兄姊妹之間的張力.
對信仰產生很多疑惑.
對教會有失望.

$^{321}$很多屬靈的衝擊.
她盡心盡力去鼓勵弟兄姊妹.
我就想如果他們不是有這班導師.
這班忠心的弟兄姊妹守護著.
他們有可能更多的人是很冷淡.
基於危機的意識.
我就向上帝禱告.
我就善用時間投入事奉.
上帝給我的感動是做什麼呢.
就是用上帝的道去牧養弟兄姊妹.
所以我就努力講道.
寫靈修文章.
用神的話語去牧養弟兄姊妹.
我就把握機會寫了一本書.
叫做基督信仰連結與整合.
是一本基於信仰的書.
我相信你們有部份弟兄姊妹.
在這本書未出版的時候.
都用過那些材料.
錦繡堂好像有百多人參加.
其他的書都是陸續出版.
我就反問我自己.
你做這麼多事有什麼用呢.
你可以改變到目前教會的困局嗎.
我覺得一定不可以.
這是很坦白的說話.
那為什麼我還要去做呢.
我只可以說.
這是我們對上帝召命的一個理解.
這是上帝的呼召.
如果這是上帝對你的呼召.
我相信你都會不計代價去投入.
即是有一樣事你想做.
你都會奮身去做.
每個基督徒都有召命的.
每個基督徒都應該有忠於使命的態度.
即是說這是你人生一定要做的事.
我們要履行我們的天職.
信徒皆為祭司.
大家都很熟悉.

$^{361}$這是一個很重要的信念.
無論你是牧師.
傳道人.
或者你不是.
你都應該去事奉上帝.
這是很肯定的.
在香港教會現在只是轉型的階段的時候.
最老套的說一句就是教會需要你.
這個去到哪裡都對的.
第二我想說的是宏觀資訊的裝備.
在馬太方第25章6到9節.
就是在說新郎遲到.
他遲到的時候那班人就睡了.
夜半就有人哭了.
新郎來了.
你們出來去迎接他.
這個真的令我們覺得很奇怪.
為什麼呢.
其實在這個經文記載的時候.
那班童女在等新郎.
老實說等得久都會睡覺.
很正常.
不過你要想.
那五個有帶備後備遊的童女.
睡了都沒所謂.
你明不明白.
她睡醒就說來了.
我又有後備遊.
但那五個沒有帶備後備遊的童女.
一聽聞就大件事.
等了很久.
睡覺.
那些遊燈已經沒什麼油了.
她就癱瘓了.
那怎麼辦.
然後他們就要跟隔壁那個說.
你可不可以分些油給我.
我的燈沒油了.
如果你就是看這兩節的經文.
或者看第九節.

$^{401}$聰明的就說不行.
恐怕我不夠.
我給了你.
我和你都不夠用.
不如你自己去買.
如果就是看這兩節經文.
你會覺得用不用這麼吝嗇.
有的就兩份用.
你明不明.
我們可能就這樣看這兩節經文.
就會有個感受.
但你再看整個背景.
你就覺得不是.
第一你又不帶油.
第二你在等的時候在睡覺.
如果你知道不夠油.
你就拿走.
你在睡覺的時候去買.
到現在新郎來了.
你就馬上說.
不是啊我不行.
你來救我.
你明不明.
那個故事你看那個脈絡的時候.
你發覺.
難怪經文說這五個童女是愚蠢.
他不止愚蠢.
他還有個假設.
就是老奉人家會幫他.
就好像奉旨人家會幫你.
我不是要說黑薄說話.
而是說你再想一下.
經文是在說.
這五個童女的心態是怎樣.
好了你說.
其實都不關我的事.
最壞是新郎.
不準時來.
那是不是啊.
他遲到是不是不對.

$^{441}$不過你可以解釋.
那時候交通工具是這樣的.
他不知道騎牢也好.
騎什麼也好.
走著走著就休息一下.
沒氣就要喝水.
我們不可以將問題推到別人身上.
如果你知道去等人的時候.
在中東那個環境.
你不帶備後備的東西.
那你自己是不是真的.
好像準備得不好呢.
如果其他人幫不到你.
分享不到的時候.
那你就早一點去買油.
你不是你去睡覺.
疲倦是自然的.
老實說個個都疲倦.
不過如果你想到新郎還沒到.
你遲早會用完的時候.
為什麼你不早點去行動呢.
你不這樣做是表示什麼.
我們中文說叫無知.
很正常的說.
在過去的事奉經驗裡面.
我自己就引伸了一個想法.
就是說這段經文提醒我們.
要有心理準備迎接神的來臨.
同樣是引伸我們事奉的人.
我們要預備我們自己.
要預備我們自己來給神所用.
我們有沒有去了解一下.
宏觀一點去看一下.
香港的教會的情況呢.
所謂宏觀一點的資訊包括.
全世界教會不要只說香港.
有些有代表性的教會.
或者我為什麼要這樣做呢.
是要確立我自己的信仰的看法.
和我自己的教會觀.

$^{481}$每個人都不一樣.
但是你知不知道你自己傾向什麼.
為什麼傾向那樣東西.
是你要自己去想的.
除了對香港教會的認識.
還有對海外教會的認識.
當然對有代表性的教會的認識.
你要留心不同時期的發展.
和經過了什麼.
其實都要花很多時間的.
很多時間.
但是你想一下.
如果我對教會是上心的時候.
你很自然會追下去的.
你想它好你自然要想辦法.
去明白到究竟是一回什麼事.
很正常的.
如果是牧者.
一定要有他自己認同的教會理念.
如果你是信徒領袖.
你都要確定你自己.
所認同的教會理念.
你不只是扮一個事工一個活動.
現在是一個轉變的時代.
沒有一套方法永遠有效的.
你要知道這樣一件事.
不同的牧者正在嘗試新的方法.
有些有效果.
有些是沒有效果.
所以其實呢.
試就齊齊試的.
成功就一起成功.
不成功就要想辦法.
這個是一個事實.
現在信徒很多有訴求的.
是會選擇自己喜歡的模式.
我們以前不懂得這樣想.
我們只有一種方式.
但在疫情期間.
很多人可以上網去崇拜.

$^{521}$突然間發現原來隔離那間.
他的領唱都幾好.
好像那間弱一點.
領唱那個走音.
你明不明.
以前這些不會暴露出來的東西.
現在是完全都見到.
或者你會.
大家我不懂彈結他.
你一網上崇拜.
那間那些會甩的.
預備得不夠.
你就會開始有你自己的想法.
我記得有些弟兄姊妹就說.
不是啊我們要精益求精.
如果你的講完.
早點給了講章來.
我還可以自磨出街.
是不是很高.
講一兩年前.
這些很佩服.
因為只是這些.
有些人是按下按下按下.
看你幾秒.
看了幾秒就決定了他的取向.
這間不用看.
這間可以不用看.
這間值得看.
當然我不是鼓勵大家要追求那個遊戲.
但你會預算到很多弟兄姊妹.
他會有取向的變化.
他以後會有改變.
這個才是問題.
現在教會發展的地步.
就會形成某些文化.
某些文化就會慢慢形式化.
組織化定型.
是很普遍的.
求變是很自然的.
不過怎樣變呢.

$^{561}$最終的目標是什麼呢.
是未必有明確的答案.
不過呢.
又有很多人.
就好像這段經文所講.
他仍然在沉睡中.
在沉睡中.
一直停留在自己的安書區裡面.
他從來都不會關注其他教會.
或者其他普世教會的經驗.
那有些人又怎樣呢.
他又不能夠等下去.
他就很快嘗試新的東西.
當然如果可以不斷修正的時候.
就有很大的可能.
殺出一條新的血路.
這是自然的.
不過在社會事件和疫情之後.
很多教會的人數下跌.
奉獻數字都下跌.
牧者信徒的領袖都移民.
是很普遍的事.
我從其他牧者的口中知道.
在英國,美國,加拿大,澳洲.
是有很多華人教會是發展的.
我都參考不同的同道.
他的信息.
來建立我自己的看法.
以下我沒有寫下.
不過我都要和你說.
其實很多人都找我寫推薦信.
很多傳統人找我寫推薦信.
我手下都寫過不知多少封.
即是推薦去外國.
讀書,工作等等.
不評論了.
我自己的看法是.
我們面對本土化和國際化的宣教理念.
本土化是什麼呢.
就是在現在很多資深的牧者退休潮.

$^{601}$成熟的牧者移民潮.
信徒領袖開始流散的情況下.
我們怎樣重整教會.
固本培元.
這是肯定的.
國際化就是宣教在香港人口流散的時候.
去到世界各地.
我們怎樣可以和這班人同行.
和幫他們建立教會.
做還是不做.
在一個急變的時代.
我們需要有新的思維.
要看清楚世界的新面貌.
去問上帝的心意是什麼.
要有新思維.
同時要有深度的內容.
現在的危機就是說.
所謂新思維.
就是我打倒你所有舊的東西就是新的.
另一方面可能潮流流行什麼我們做什麼.
隨著時代變化.
信徒面對信仰和生活之間的張力.
和怎樣落實信仰.
怎樣活出信仰.
是一個很重要的課題.
從事教導和牧養的牧者和信徒.
不單是需要在基本的聖經.
教會歷史.
教義倫理知識.
同時他需要對時代處境的分析.
還有需要悉切的神學的詮釋.
技能的裝備是需要與時並進.
知識的裝備也都需要與時並進.
我們需要一群是願意學習新科技.
同時願意好好學習聖經歷史神學的信徒.
願意嘗試將所學的應用在急變的社會裡面.
很多人想嘗試新的東西.
但那個餡是有很多問題的.
結果其實都是成功了一半.
因此我們焦點就是仍然做訓練神學人才.

$^{641}$傳承基督信仰.
如果要更新傳統.
首先你要認識傳統.
我鼓勵大家你要按著上帝的感動.
去做你應該做的事.
去榮耀上帝.
第三要講的就是行動的準備.
第25章第8到第13節.
就是重複那些愚蠢對聰明的說.
你分些油給我吧.
因為我們的燈快要滅了.
聰明的就回答說.
恐怕不夠你和我用.
不如你去賣油的地方去買吧.
但後面就說.
當他們去買的時候.
剛好新郎就到了.
那些預備好的那些.
那五個聰明的同類.
就由新郎邀請他進去.
共赴婚宴門就關了.
這個完全是象徵的說法.
因為古靈精怪這麼晚還開飯.
你明白嗎?很奇怪.
不過這個只是比喻.
其餘的同類就跟著來了.
哎呀他找到油了.
跟著就說主啊主啊.
開門給我吧.
那麼這個主就說了.
我實在的告訴你們.
我不認識你們.
哇是不是很重啊.
這個回應很重啊.
所以後面那句才是最重要的.
你不要管它重不重.
願景民說.
所以你們要警醒.
因為那日子那時辰.
你們不知道.

$^{681}$所以整個故事裡面.
你要追它的情節.
你可以有很多的問題.
但是最重要就是說.
如果我們有醒覺有危機意識.
我們知道去等新郎的時候.
除了有油燈還有另外的油.
栽在旁邊.
那麼我等久一點就沒事.
這個故事就像是在說危機意識.
你是不是真的用心去處理那件事.
不是說主人這麼嚴厲.
沒有油而已.
遲了來你關門說我不認識你.
如果我們這樣想.
就會走錯了方向.
誤解了這段的比喻.
耶穌說這段比喻.
告訴我們.
我們要學習的是居安思危.
有危機的意識.
如果你是一個信徒.
你沒有危機意識.
傷害不是很大.
你是一個信徒領袖.
沒有危機意識.
你對教會的傷害可以很大.
如果你是一個牧者.
你沒有危機意識.
你對教會的傷害情況是更大.
這段經文告訴我們.
沒有預備額外油的童女.
買油回來的時候.
新郎就來了關門.
可能會覺得.
新郎你太不負責任了.
你遲到還要禁我.
不過.
有人問.
為什麼這麼重罰這五個童女.

$^{721}$我也不懂回答.
但這段經文是告訴我們.
時間不等你.
也不等我.
機會一失就沒有.
如果是一件小事還可以.
但這段是說.
你去做神的事情.
你沒有這個意識的時候.
就很難做到耶穌的門徒.
我用這個說法.
這個.
比如有時候我們要明白.
是要突破我們人的思維.
我們可能覺得.
不用這麼認真.
不用這麼嚴重.
但我相信你再想深一層.
當我們面對這麼大的挑戰的時候.
上帝是要求我們.
為祂做事的人.
是要有更高的水平.
所以我們要想的.
是否要像那幾個聰明的童女.
好好預備自己.
在每件事上.
我們都盡我們的力量.
所以我們應該行動起來.
不過我想跟你說.
行動起來也有一些限制.
其實我都是行動起來.
我的行動起來.
是因為在疫情的時候.
其實是安息年.
所以我沒有怎麼去講道.
我也有網上幫忙講道.
其實去年我比較多時間留在長洲去崇拜.
但是我計錯了.
今年2023年.
我就擁漢得多.

$^{761}$擁漢到一個地步.
我懷疑除了復活節崇拜不用講道之外.
我差不多每個星期都要去講道.
這個就是擁漢過渡的結果.
我不是說我聰明.
我只是說因為以前.
覺得有一種停擺的感覺.
我就鞠行乎己.
我就去到盡 衝到行.
但這個不是好的做法.
因為現在三月不是去十二月.
我已經覺得很累.
整個星期沒有休息.
很難受.
所以我奉勸大家.
我們真的要想.
你有一個感覺.
我要把握我的時間.
我要用盡我的時間去為神做事.
但是你要比我聰明一點.
就是你要編排時間好一點.
如果不是你不會走得更遠.
我的身體都不會好.
我自己都會出訊號.
所以我們行動是要準備.
但不是一下就完.
我和我太太說.
我打算為主消盡.
接著她說.
那你很快就消完 玩完.
你明不明.
我們用這樣的態度來看.
我們是要有心理準備.
我們是為主拜上.
我們是為主死就死吧.
但很多長者都會告訴你.
死你都要找一個適合有價值的方法.
你這樣是什麼意思.
你一時高興是沒有意思的.
燒一陣煙花是沒有用的.

$^{801}$應該細水長流.
耐心一點.
我們要考慮一件事.
行動我們要有一個心態.
是全然為主.
但是你要計算一下你是人.
你要適當地去做.
所以我都有一個決定.
明年我就不會像今年那麼勇悍.
起碼我都要給自己休息.
不是這樣下去.
我很可能已經是躺下了.
見到上帝.
說到這些之後.
我很想和大家分享.
我們初訊的時候很熱心.
其實我初訊的時候都很熱心.
是一個很激情的宗教份子.
我不知道你們是不是很激情的宗教份子.
我算是一個很激情的宗教份子.
所以我很明白初訊的時候是很有熱誠.
教會林林審審各樣東西.
你都會很熱心去做.
但去到一個階段.
就會有一種.
好像為什麼人們都不做.
我舉一個例子.
我本來都沒有寫下.
當年我的教會是在一間.
在駟字區樓下的越南室.
我們就租了.
決定我們那時候不是很多人.
可能20多個人.
青年人.
租了一間天台.
荒廢的學校去搞報道會.
大家都很開心.
大家分工作.
我就領了一份工作.
是去清理場地.

$^{841}$清理場地.
你明白嗎.
但你要想.
在一間荒廢的天台學校.
你們知不知道駟字區天台學校是什麼樣的.
中間空的.
兩邊有課室.
一頭一尾的.
是我領了這份差事.
叫統籌.
那班弟兄就說.
你統籌我們來幫你.
大家都很明白.
我就呼召那些人.
我一個人怎樣拆那些鐵枝.
爛窗.
怎樣整.
當然是約人.
發覺.
為什麼我們分工的時候.
大家都很豪氣.
為什麼我們約人的時候.
好像經常都約不齊.
結果.
最後.
我去找外援.
找了我一個同學來幫忙.
我還記得他的牛仔褲都弄壞了.
那時候他還沒有信主.
他現在正在回某間教會.
我都很擔心.
他會不會因為我.
去幫我搬台搬椅搬垃圾.
他從此就覺得教會是很笨的.
誰知道神就使用他.
我這樣說.
我都會經歷這些事.
如果你有類似的經歷.
我都很明白.
我就會問.

$^{881}$講的時候個個都很豪氣.
做的時候不見了影子.
我都會想.
是不是這樣做.
所以在那個階段的時候.
我都很多不滿.
我認為這樣都可以.
現在這樣當然成熟了很多.
知道你傻仔.
你當然不要接那個位.
你去接到時候做招待的.
你接那些做什麼.
這個是我現在的智慧.
但就算你接.
你都要多一個人來做一齊.
這些不說了.
但是我想說.
其實我們每個人多多少少都有一些經歷.
你在事奉的過程中.
你也是經歷的.
在你的信仰中也是經歷的.
不過我想說.
行動起來.
你有沒有處理這些經歷.
我剛才說給你聽很輕鬆.
其實我已經忘記了.
其實我已經毫無憤怒了.
你明不明.
但是如果我們在我們的記憶中.
有一些東西令你覺得卡住.
其實都不是一件好事.
是會阻礙你去發揮.
阻礙你去投入.
坦白說.
其實我是很知道人性的軟弱.
和在教會裡面的那種活動的表現生態.
我又舉個例子給你聽.
我以前在教會是宣教士辦的.
是西人宣教士.
師母很喜歡關心人.

$^{921}$因為人數不多.
二十多三十個人.
他很留心每次崇拜誰人沒有來.
他會在那個星期後誰人沒有來.
他會打電話去問候一下人.
我發覺我一次是忘記了跟師母告別.
我在那裡.
不過我沒有跟她拜拜.
接著她的電話就到了.
我說其實你這麼關心我是好.
不過我不是很想.
下一次回去教會我就聰明了.
我臨走跟師母.
故意在她面前兜兜兜.
師母再見了.
這些年輕人的時候是這樣.
我不想你再經常這樣來關心我.
有些人是不喜歡的.
我是很怕的.
所以我又知道弟兄姊妹的表現.
他們的動作是什麼含義.
在你面前仰仰之後.
謝牧師我今早見了你.
你不用來問候我.
那回事.
我是不想這麼多.
其實心知肚明.
不用你這麼張揚.
拜拜都不用這麼明顯.
你的動機很明顯是想人知道你到了打了卡.
我這些一眼看得出.
另外一些.
當我去.
剛才這樣說經歷了那些事情之後.
我又想隱藏一下自己.
我都是沒有什麼力量.
我很明白.
有些弟兄姊妹.
上教會盡量隱性賣命.
其實她很厲害.

$^{961}$她以前是很能幹的.
不過她說我都是受傷的小鳥.
我都不想再出江湖了.
我去到發現兩下.
原來自己人來的.
我當年也是這樣.
所以我不會煩你.
我不會去騷擾你.
我很想說.
其實我遇過不少這些弟兄姊妹.
當然我不會刻意去騷擾她.
但是我想說.
行動的準備.
我是一個足球員.
我受傷之後.
我都想出球場.
你明不明.
我是夢想去球場.
我現在都想的.
我受了傷現在出不了人.
等於一個運動員你受了傷.
你都想再去比賽.
你不會甘心在牆邊.
同不同意嗎.
所以我很希望我們的弟兄姊妹.
行動起來不是喊口號.
而是說如果你們有各類各類的.
你的經歷.
你的感受.
任何大道理你都聽過.
培靈會你都去過很多次.
但是為什麼你仍然選擇.
做一個抽離的弟兄姊妹呢.
這樣大家就明白了.
我不知道可以說什麼給你聽.
不過我只想說.
這個比喻告訴我們.
機會不會經常在等我們.
或者這樣說.
上帝做事是不需要我們.

$^{1001}$我也不是什麼特別的人.
只不過上帝做事.
有機會預了我一份.
這樣說吧.
其實如果上帝做事預了你一份.
你的想法是怎樣呢.
你的想法會是.
上帝你是不是想多了.
其實我不是很想做的.
還是說.
上帝做事預了我一份.
其實我很光榮.
究竟你的選擇是哪一種呢.
你明白我的意思嗎.
好了.
我知道我不想說太長了.
但總的來說.
現在是做事的時期.
我今天其實是不經意地.
答應了另一次港島.
我去得很勇猛的原因就是這樣.
我下午四點半.
就會用Zoom去幫德國一間華人教會港島.
初時我跟肥沙他們說.
可不可以借用這裡.
全都很看得起就來.
其實那個教會是五點鐘崇拜.
四點半我要去試音.
很嚴格的.
提早半個小時來搞事.
我想五點鐘說到六點多都還沒下班.
這樣會阻礙人吃飯的.
還有我要回長洲.
這樣搞得我有事.
這裡的傳道人又收不到工.
所以我就轉了另一個方法.
我想說為什麼我去年太勇猛的原因.
我覺得應該多做點事.
但我自己不懂得收掣.
但我想說.

$^{1041}$你有沒有想過原來疫情期間.
你可以用Zoom去很多地方港島.
我真的有做過.
不是假的.
我去年都試過幫德國好幾個城市港島.
我幫英國的Stock on Trench.
史督城教會港島.
我幫愛丁堡的教會港島.
我幫法國Toulouse的教會港島.
用普通話的.
我從來沒有想過有這個機會.
你明不明.
以前你要搭飛機過去.
現在你不用.
你抽時間.
我試過最晚那次一點半.
他要教主日學.
那一點多鐘.
沒所謂在宿舍講完睡覺.
現在的機會就在我們面前.
你做還是不做.
當然我們不可以要求每間教會都擁抱所有的東西.
你沒可能的.
但我會覺得.
我覺得可以做就盡做.
這是我去年的想法.
所以就自己不知不覺接了很多單件回來.
但我想說.
我會將這個體驗和大家分享.
機會是放在你們面前.
每個弟兄姊妹都可以為神做事.
我希望你們有這個心態.
從這個故事這個比喻裡面.
你學聰明的同類.
你好好預備你自己.
心態上預備自己.
認知上預備自己.
行動上預備自己.
好好的在這個疫情之後.
好好地生活.

$^{1081}$好好地把時間.
發揮如光如熱.
我想用這四個字都很對.
我們一起同心祈禱.
我們天父上帝我們向你禱告.
因為你自己就是差遣我們的神.
是你要做這個時候.
去更新去復興的工作.
我們看到很多的人.
他的人生有很多的困難.
他們找不到他的目標.
他們雖然有好的發展.
但是內心很空虛很空洞.
因為他們沒有神.
如果神差遣我們去做這些工作的時候.
如果我們的內心有這些掙扎的時候.
有很重的擔子的時候.
就求你預備我們的心.
從同類的比喻裡面去得到教訓.
我們真的要在心態上調整.
在知識上要調整.
在行動上要調整.
讓我們真的能夠不枉此生.
我們好好地在你的面前.
去領受你的召命.
去做你要我們做的工作.
我們願意你的靈在我們當中.
心裡面去工作.
在我們心底裡面有不少的那些困鎖.
或者是那些不開心的事.
但是我們的主願意你.
首先來處理我們內心的困擾.
讓我們心裡面充滿你自己的愛.
充滿著你自己的力量.
讓我們真的把握我們每一天.
因為我們真的不知道.
我們活在這個地上的日子.
是有多久.
我們也知道沒有人可以有把握.
一個疫情可以將一個健康的人奪去生命.

$^{1121}$我們求主讓我們有這種心智.
好好地珍惜在世的日子.
來做神理的工作.
我們的祈禱是奉耶穌基督的命.
阿們.
\newpage



\section{以弗所書 3:14-21}
\label{sec:rNmdtGHU_n0}
\textbf{2023-03-26 加恩樓《按神的心意成就一切》以弗所書3\_14-21 講員:鄧泰康傳道}
\newline
\newline
連結: \href{https://youtube.com/watch?v=rNmdtGHU-n0}{\texttt{https://youtube.com/watch?v=rNmdtGHU-n0}} ~~~~ 語音日期: 2023-03-28
\newline
\newline
\hyperref[sec:wPsCGHG1pLA]{\small{< < < PREV SERMON < < <}}
~
\hyperref[sec:index_chronic]{\small{[返順時目]}}
~
\hyperref[sec:index_scriptual]{\small{[返順卷目]}}
~
\hyperref[sec:pzC3e8fjAp8]{\small{> > > NEXT SERMON > > >}}
\newline
\newline
以弗所書 3:14-21
\newline
\begin{longtable}{cl}
\hline
\hline
章節 & 經文 (和合本修訂版)\\
\hline
3:14 & \begin{tabularx}{0.7\textwidth}{X} 因此，我在父面前屈膝— \end{tabularx} \\ \\ \relax
3:15 & \begin{tabularx}{0.7\textwidth}{X} 天上地上的各家都是從他得名的－ \end{tabularx} \\ \\ \relax
3:16 & \begin{tabularx}{0.7\textwidth}{X} 為要他按著他豐盛的榮耀，藉著他的靈，使你們內心的力量剛強起來； \end{tabularx} \\ \\ \relax
3:17 & \begin{tabularx}{0.7\textwidth}{X} 又要他使基督因著你們的信住在你們心裡，使你們既在愛中生根立基， \end{tabularx} \\ \\ \relax
3:18 & \begin{tabularx}{0.7\textwidth}{X} 能夠和眾聖徒一同明白基督的愛是何等的長、闊、高、深， \end{tabularx} \\ \\ \relax
3:19 & \begin{tabularx}{0.7\textwidth}{X} 並知道這愛是超過人的知識所能測度的，為要使你們充滿神一切的豐盛。 \end{tabularx} \\ \\ \relax
3:20 & \begin{tabularx}{0.7\textwidth}{X} 神能照著運行在我們心裡的大能充充足足地成就一切，超過我們所求所想的。 \end{tabularx} \\ \\ \relax
3:21 & \begin{tabularx}{0.7\textwidth}{X} 願他在教會中，並在基督耶穌裡，得著榮耀，直到世世代代，永永遠遠。阿們！ \end{tabularx} \\ \\
[1ex]
\hline
\hline
\end{longtable}
$^{1}$頂智梅早安.
今天有機會和大家一起分享新的說話.
今天和大家分享的一段經文其實又是不複雜的.
很簡單的.
其實大家從字面上都很容易會catch到裡面的意思.
但重要的是背景.
監獄書信 大家知道吧.
有四卷的監獄書信是保羅在獄中.
他仍然記掛著這些頂智梅.
然後他就寫信去鼓勵他們.
其實我不知道大家會不會好奇這件事.
我們都信了耶穌 甚至我都信了三十多年.
我們信耶穌一段日子之後.
我不知道大家有沒有一種好奇.
為什麼在聖經裡面的人物這麼厲害.
保羅為什麼這麼厲害.
他自己都曾經說過.
他不知道被人打了多少次.
他不知道受過什麼屈辱.
但他竟然一直以來.
他仍然和這些弟兄姊妹表現到一樣東西.
就是非常有信心.
非常有喜樂.
非常有一種屬靈的能力.
我不知道大家羨慕不羨慕.
我經常說到.
有些片段我會不斷重複.
大家不要覺得我慢慢都五十多歲.
快六十歲的話.
我就會是一個老人家.
會不斷地和大家說一件事.
記住保羅不是神仙.
保羅不是頭上有光環.
他和你和我都是一個普通人.
那就是說保羅能夠得到的東西.
是我們都應該可以得到.
這個邏輯對不對.
但是為什麼很多時候.
今天在我們身上.
我們總覺得我們只有羨慕的份.

$^{41}$但好像不是在我們身上發生的.
弟兄姊妹.
看聖經就是看這些.
看聖經不只是一個習慣.
看聖經是讓我們能夠從聖經當中.
從上帝留下的這些人身上.
我們去學究竟我們欠了什麼.
特別今天這段聖經.
其實奇妙的地方是什麼呢.
這段聖經是保羅受感動的一個禱告.
這段經文就是到了第三章的時候.
保羅一直去鼓勵以法所信徒.
第三章第一至十三節的時候.
保羅跟他們說.
其實為什麼會發生現在在監獄裡的事.
因為我在上帝那裡領受了一個奧庇.
而這個奧庇是什麼呢.
就是愛邦人能夠得著救恩.
這件事在猶太人來說是不可能發生的.
因為猶太人一直都覺得.
上帝只是跟我們立約.
所以為什麼當保羅傳道的時候.
最初他被人攻擊得這麼厲害.
意思就是你不要在這裡亂說話.
上帝的救恩只有我們有.
其他人不會有.
但保羅很清楚說.
我是在上帝裡面直接的.
如果大家從史上傳得很清楚.
是上帝直接的跟他說.
要讓這個救恩能夠去到愛邦人身上.
而這些是有舊約的聖經才知道.
他們曾經預言過的.
所以你就會發覺.
保羅既然很肯定.
這是上帝的託付之後.
他就一直堅持下去.
就算坐牢都在所不惜.
弟兄姊妹.
那就是說要決是什麼呢.

$^{81}$你只需要知道.
那件事是上帝的心意.
就一定會成就.
是不是.
所以聖經裡面保羅不斷鼓勵大家.
你要尋求神的心意.
我現在很相信弟兄姊妹.
很多時候都是一個問號.
所以經常都當我們是廟祝.
鄧傳道究竟我這個是不是神的心意.
我經常都會回應我怎麼知道.
因為我起碼在我的信念裡面.
我只可以說我從來沒有經歷過.
是上帝真的開一把聲音跟我說話.
我是沒有經歷過.
這樣的情況下.
我不會知道上帝跟你的傾談.
跟你的回應是怎樣.
我不會完全清楚的.
我可以用什麼.
作為一個傳道人.
我用回聖經.
所以為什麼我最重要的是.
弟兄姊妹學聖經.
原因就是.
究竟你那個是不是神的心意.
唯有聖經可以給你一個答案.
所以今天有其他人說.
我看到神的一個超乎於聖經以外的異象.
或其他東西的時候.
弟兄姊妹小心啊.
我不知道大家有沒有興趣.
我還沒有這樣看.
最近在Netflix有一套劇集.
在韓國有一些邪教.
我還沒有勇氣去看.
但我看裡面的時候.
我會發覺其中一個很重點的.
就是他們所說的是超乎聖經的.
我現在是一個新的基督.

$^{121}$我就是那個彌賽亞.
你就會發覺你小心啊.
所以學聖經.
從聖經裡面出的.
是非常重要的.
當然我不會排他.
因為我經常說我不是上帝.
上帝可以有他自己各樣的方式.
去跟我們相處.
但在我個人的經歷裡面.
我極其量的就是心裡面的.
一份上帝的感動.
或者心裡面的一份微聲.
但這份微聲我都仍然要回到聖經那裡.
聖經告訴我.
上帝告訴我是不是真的.
因為當你一不小心的時候.
可以是魔鬼的聲音.
所以你就會發覺這件事很重要.
保羅他就是很肯定知道.
這個是神的心意.
他很確定的是一定會成就.
所以他在那裡的時候.
他就跟他們說.
他很想成為一個榜樣.
去告訴以弗所的信徒.
神怎樣使用他.
鼓勵他.
以至於他能夠鼓勵這些.
以弗所的信徒.
不要因為看見現在保羅的處境.
而害怕,懼怕.
而要繼續堅持下去.
所以你會發覺這段經文.
在上面的時候.
保羅用了他自己的經歷去提醒他們.
所以只要是上帝的心意.
就一定成就.
你會發覺第十四節開始的時候.
其實保羅只是做了一件事.

$^{161}$聖靈感動他裡面做了一件事.
就是為這群以弗所的人去代求.
去即時跪下去禱告.
那就是說.
如果這件事是神的心意.
弟兄姊妹你不需要擔心.
我經常告訴弟兄姊妹.
我們今天上帝叫大家做什麼.
是傳道者,是不是?.
你做傳道者.
弟兄姊妹或是你的朋友.
他信不信.
與你有沒有關係.
與你沒有關係.
上帝叫你做傳道者.
他的信不信是在上帝自己的旨意.
和上帝的感動當中.
再加上他自己的回應.
但今天很多時候我們是怎樣?.
我們就是看到沒有果效.
我們講完之後.
然後我們就不傳.
氣餒,灰心.
原來做了這麼多都沒有用.
都沒有人信耶穌.
不是啊,弟兄姊妹.
保羅當日他自己在獄中的時候.
寫了書信出來.
他都不會知道書信.
最後會成為這麼強大的力量.
是不是?.
他純粹就是很清楚.
我是去建立這些教會的.
我作為牧者.
我不要因為今天坐牢的緣故.
而阻礙了我今天的事奉工作.
所以我就寫信去鼓勵他們.
就是這麼簡單.
但上帝就使用.
所以回到我們監牢的處境中.

$^{201}$同樣道理.
我們之前一直和大家分享.
請大家禱告.
究竟嘉恩樓未來如何?.
大家要努力.
我們要一起努力.
但無論將來如何.
如果我們確信嘉恩樓是神的心意的話.
神一定會按照他的心意成就.
所以為什麼我上次說.
當我們看到上帝在動工.
我們就覺得我真是Hallelujah.
有些不是我們做到的事.
有些是上帝自己的工作.
但究竟最後如何?不知道.
我們繼續仰望上帝.
但我們要做的事還是要做.
所以你會發覺這段經文的基礎是很重要的.
第十四節開始的時候.
因此我就在父母面前屈膝.
他即時的感動就跪下來.
為他們祈禱.
第一句是什麼?.
天上地上各家都是從他的名而得的.
即是說他是天地萬物的創造主.
今日你這麼確信的時候.
因為你很相信這是創造天地那位的主.
他遠超乎你的能力.
遠超乎你的想像.
他要成就的事是從來沒有不成就的.
這是一個起點.
你今日對於在神裡面.
這一份這麼重要的確信.
你已經確定了嗎?.
即是說無論你在任何處境.
為何保羅這麼厲害?.
因為我在監獄裡面.
我仍然確信這個上帝.
他是會成就他的事.
所以我滿有盼望.

$^{241}$我仍然可以唱歌.
我仍然可以喜樂.
力量是這樣出來的.
不是看到我面前的監獄.
不是看到我面前的宅打.
不是看到其他人的迫害.
我看到的唯有就是上帝的大能.
弟兄姊妹,這就是基礎.
第二件事,你會發覺.
所以屬於他的人.
生命就需要改變和更生.
今日似乎更多的情況就是.
不是上帝出了問題.
很多時候人遇到一些問題的時候.
就不找上帝.
上帝,你看不到嗎?.
有沒有搞錯?.
你去了哪裡?.
是上帝出了問題.
但不是的,弟兄姊妹.
這裡保羅很清楚地說.
上帝從來都不會有問題.
是人出了問題.
這裡其實是提醒我們.
我們有沒有去改變我們的生命.
去盛載或配合上帝的計劃.
所以我們作為信徒更加要問.
究竟我們今天的生命有沒有阻礙了上帝.
有阻礙,請你把那些東西拿走.
用一個最簡單的.
大家有沒有靈修?.
又問一些老套的問題.
為什麼不靈修?.
為什麼我沒有靈修?.
忙.
早上起來就累了.
聖經很難看.
是不是?.
弟兄姊妹.
忙,那怎麼辦?.

$^{281}$令自己不要太忙.
不是的,我真的很忙.
我知道,我也很忙.
但如果你問我.
那我有沒有時間讓自己放鬆一下.
去看Leviticus?.
當然有.
那你有沒有將你看Leviticus的時間去靈修?.
我今天很累.
為什麼累?.
為工作?為勞累的事情?.
我們有沒有傳言將這些交給上帝.
由上帝去掌管?.
我們今天很累.
工作的,公司的繁忙很累.
然後晚上.
既然這麼累,那怎麼辦?.
跟朋友去聊天,喝東西,吃東西.
然後甚至很晚才睡.
但我想強調.
我很明白.
有些人特別是做父母.
晚上不是你不想睡.
是你的子女不肯睡.
有時是很多忙碌,很多事情.
但是弟兄姊妹.
如果當你確信在神面前.
看祂的話語是重要的話.
這個必定會是優先的.
早一段時間.
我們跟一個團契去查經.
查了一本書.
叫做《禱告的實學》.
這位真金法牧師提到一樣東西.
其實不斷提到同一樣東西.
究竟你認為祈禱是不是你生命中最重要?.
如果你是最重要的話.
你就要改變你自己的生活方式.
你就要改變你的心態.
以至於能夠配合這個禱告的發生.

$^{321}$弟兄姊妹,有沒有改變?.
所以你會發覺.
在這裡的時候.
他就看了我們要怎樣改變.
第一件事.
靠著上帝的榮耀.
靠著聖靈.
心裡面剛強起來.
剛才我提過.
這是創造天地的主.
這是掌管一切的主.
難道祂沒有辦法改變我們嗎?.
首先你要確信.
在理不能.
在神凡事都能.
你不會比耶穌所說的厲害.
我對著這座山.
我只需要有信心.
祂就可以移開.
你比這座山硬吧.
弟兄姊妹.
如果你肯讓上帝去改變.
上帝一定可以改變你.
一定可以.
第二件事.
他一直強調.
靠著聖靈.
弟兄姊妹.
今天我真的很懷疑.
我自己也很擔心.
究竟我們今天.
我們作為基督徒.
我們有沒有被聖靈充滿.
而你會發覺.
在《使徒行傳》裡面.
我們也說得很清楚.
這是聖靈行傳.
弟兄姊妹.
我想列舉一點.
大家把印象拼出來.

$^{361}$第二章的時候.
因為他們被聖靈充滿.
所以他們說起別國的話.
來做什麼呢.
翻譯.
因為當時彼得在說.
是很多世界各地.
來自不同的猶太人.
但已經不是說那種語言.
來到那裡的地方一起過節.
接著他們被聖靈充滿.
他們就說起別國的話.
來幫他們去做翻譯.
即時的.
所以聖靈就像舌頭一樣.
去淋到那個人身上.
第四章.
彼得因為被聖靈充滿.
就對眾人講道.
所以以至於很多人都被感動.
接著聖經說.
因為那些人.
聽的人都被聖靈感動.
他們就放膽去傳道.
第六章.
被聖靈充滿的人.
他才能夠可以做執事.
當他們再選執事的時候.
其中一個很重要的條件.
必須被聖靈充滿.
是不是.
第七章.
屎堤反他被迫害.
他說去到最後的時候.
他說因為聖靈感動的緣故.
所以他不懼怕.
他看到天開了.
就是說人在困難.
患難當中.
因為聖靈的緣故.

$^{401}$而他能夠看到天國.
他能夠超越困難.
看到未來.
第十一章.
巴拿巴被聖靈充滿.
所以他是一個很有能力的人.
因為他的緣故.
很多人信了耶穌.
第十二章.
保羅要被聖靈充滿之後.
才開始事奉.
弟兄姊妹.
其實我們究竟今天.
我們有沒有成為.
被聖靈充滿的人.
我想強調的.
我們被聖靈充滿.
不是我們有什麼特別的能力.
不是的 弟兄姊妹.
而是我們整個生命會改變.
大家有沒有發覺.
裡面有很多生命被改變.
他們很專注於.
我今天所做的每一件事.
是不是在做神的工作.
我是否在改變其他人.
我是否在幫助其他人.
我今天是否一個.
充滿著聖靈能力的人.
如果我是充滿著聖靈能力的人.
我應該可以讓人看到的是什麼.
像保羅一樣.
在任何患難當中的時候.
我仍然會有白目.
當我今天在任何處境裡的時候.
我仍然相信上帝的同在.
你的臨在.
而是更少看見自己的困難和處境.
被聖靈充滿的人.
應該有這份屬靈的生命.

$^{441}$有一種屬靈的力量.
弟兄姊妹.
今天我們回來教會.
我們回來崇拜.
我們相信上帝的時候.
我們有沒有.
我們嘗試一起去想.
平時我們的生活是怎樣.
平時我是一個怎樣的人.
我是一個仍然為自己去活的人.
為自己的想法去活的人.
我們是一個.
別人看見我們的時候是得了祝福的人.
還是別人感覺我們的時候是得了一種討厭和咒詛的人.
別人看見我們面容的時候是得了平安.
得了喜樂.
還是看見我們面容的時候是充滿著憤怒.
憂愁.
主啊.
我要求神幫助帶領我們.
所以為何保羅在那一刻的時候.
他即時的感動.
要為他們禱告.
所以你會發覺.
如果我們被聖靈充滿的話.
你就會看見.
他們就會因為信基督的緣故.
他們就會經歷到.
基督住在他們的裡面.
而讓基督這份無法測度的愛.
能夠完全充滿他們.
我想大家留意最後的圖畫是甚麼.
是一幅很漂亮的圖畫.
即是說無論任何處境.
你隨時隨地都會感受到.
主啊,原來你這份愛的大能.
是這麼厲害的.
是你整個人都會充滿著基督的這份愛.
是你時時都會感受到這份甜蜜.
是你時時都會感受到很有力量.

$^{481}$原來保羅想說.
為何他能夠這樣呢?.
因為他在監獄裡面.
他就是這樣的狀態.
他在裡面的時候.
他只看見上帝.
他只看見上帝還有事要他做.
所以弟兄姊妹.
原來是我們自己需要去改變.
因為到最後一句他這樣說.
神能照著運行在我們心裡面的大力.
充充足足地成就一切.
超過我們所想所求.
原來當你有上帝這份愛的力量的時候.
你會去看到你自己看不到的東西.
你會超越你自己.
你會以為我以前就是這樣.
我們人最喜歡.
我就是這樣.
你不要逼我.
但不是的.
我想補充一點.
我們在講台前經常提醒大家.
我們要按恩賜事奉.
但弟兄姊妹我不知道你有沒有經歷過.
通常在教會裡面.
很多時候要你去幫忙事奉.
都未必是你的恩賜.
但有些人反過來說.
這樣教會就不對了.
你都沒有按我的恩賜事奉我.
但不是的弟兄姊妹.
超越.
你有沒有想過有些情況下.
不是你本身有.
但因為上帝的緣故.
你可以超越自己.
我想在這個時間和大家分享一首詩歌.
我相信大家很熟悉.
如果大家懂得唱就一起唱.

$^{521}$這首詩歌曾經很感動我.
是令我從低谷重新再起的一首詩歌.
謝謝.
可以再大聲一點.
謝謝.
Face me up, you won't betide me..
You make me love, oh, I can't stand the mountains..
You make me love, you are the stormy sea..
I am strong, when I am on your shoulder..
You make me love, oh, if I could be..
You make me love, oh, I can't stand the mountains..
You make me love, you are the stormy sea..
I am strong, when I am on your shoulder..
You make me love, oh, if I could be..
You make me love, you are the stormy sea..
(音樂結束).
好,謝謝.
其實我們每一個人都一定會有低潮的時候.
這首詩歌是每一次當我低潮的時候都會響起的.
是什麼情況出現呢?.
就是當我前妻她癌症離開的時候.
我記得那段時間我問了上帝很多問題.
我問了上帝很多,喂,上帝,什麼意思?.
我不明白.
我都問,究竟是不是有很多事情是我做得不好.
我問了很多問題.
但我記得我哥哥將這首歌寄給我.
然後我發現這首歌裡面提醒我一件事.
有很多問題我未必能解得通.
但我仍然能夠帶我再一次起來.
我仍然相信上帝在當中,祂是慈愛.
所以我會發覺這首歌說,當我遇見你的時候.
我就能夠超越自己.
弟兄姊妹,你相不相信上帝真的可以超越你?.
我們今天其實每一個人都很需要在這個時代裡.
我們能夠有這份能量.
我們有這份作為基督徒的能力和心意.
因為我們才能夠讓更多人找到這個盼望.
為什麼我很喜歡這個短片呢?.
我不知道大家有沒有留意.

$^{561}$我剛才在街上街頭街道.
你會發現,我剛才很留意.
當他唱的時候,很多人可能只是路過.
你不知道他是什麼心情.
但當你願意走出一步的時候.
你去向別人見證,你去告訴其他人的時候.
原來我們生命能夠遇見的時候.
你不知道會感動了誰.
弟兄姊妹,我們今天的人生都一樣.
我們不會知道你在下一刻會遇到誰.
但弟兄姊妹,聖經的教導就是.
請你因為上帝的緣故,持續將善行出來.
你將糧食撒在河上.
你不知道什麼時候這些事情會開花結果.
你不知道的,弟兄姊妹.
但請你堅持去做.
上帝會成就一切大事.
所以請大家在今天這段經文當中相信.
我很盼望,弟兄姊妹.
上帝的愛,我不知道你如何體會.
上帝對你的付出和犧牲,我不知道你如何體會.
但在這裡告訴我們.
上帝的愛是你無法測度的長,闊,高,深.
是超越你自己所能夠想到的.
所以弟兄姊妹,我很盼望我們作為基督徒.
我們能夠活出一個合乎基督徒的生命.
我們能夠求上帝保佑我們.
當我們不行的時候,改變我們.
求神興起我們.
讓我們能夠因為遇見祂的緣故.
我們可以站在高山.
因為相信祂的緣故.
祂會帶我們經歷這些狂風巨浪.
我們仍然可以跟其他人說.
你放心,原來我們的上帝在這裡.
因為祂是愛.
所以求神保守帶領我們.
很願意今天我們一起有這個提醒.
我們同心一意祈禱.
主愛求主,你讓到我們的生命當中.

$^{601}$我們能夠更加體會上帝的愛更多.
求你讓我們小看自己的限制.
讓我們能夠在神的面前的時候.
我們能夠求你的聖靈.
親自充滿我們每一個.
主讓我們能夠有更多的力量.
我們活出屬於你的生命.
我們願意將整個間流.
裡面也包括了我們每位弟兄姊妹.
全然交在神的當中.
我們求你親自使用我們.
求你成就這一切.
我們也很盼望.
我們繼續看著上帝你的聖靈.
怎樣去引領我們一起走這些路.
我們獻上感恩,多謝你.
其樂福,靠主耶穌,得勝名至其及.
阿們.
\newpage



\section{提摩太前書 6:3-10}
\label{sec:pzC3e8fjAp8}
\textbf{2023-03-26 早堂《敬虔與知足》提摩太前書6\_3-10 講員:陳靜娟姑娘}
\newline
\newline
連結: \href{https://youtube.com/watch?v=pzC3e8fjAp8}{\texttt{https://youtube.com/watch?v=pzC3e8fjAp8}} ~~~~ 語音日期: 2023-03-28
\newline
\newline
\hyperref[sec:rNmdtGHU_n0]{\small{< < < PREV SERMON < < <}}
~
\hyperref[sec:index_chronic]{\small{[返順時目]}}
~
\hyperref[sec:index_scriptual]{\small{[返順卷目]}}
~
\hyperref[sec:yk5Oh8RBHLY]{\small{> > > NEXT SERMON > > >}}
\newline
\newline
提摩太前書 6:3-10
\newline
\begin{longtable}{cl}
\hline
\hline
章節 & 經文 (和合本修訂版)\\
\hline
6:3 & \begin{tabularx}{0.7\textwidth}{X} 若有人傳別的教義，不符合我們主耶穌基督純正的話語與合乎敬虔的教導， \end{tabularx} \\ \\ \relax
6:4 & \begin{tabularx}{0.7\textwidth}{X} 他是自高自大，一無所知，專好爭辯，擅於舌戰，因而生出嫉妒、紛爭、毀謗、惡意猜疑， \end{tabularx} \\ \\ \relax
6:5 & \begin{tabularx}{0.7\textwidth}{X} 和心術不正與喪失真理的人不停地爭吵，以敬虔為得利的門路。 \end{tabularx} \\ \\ \relax
6:6 & \begin{tabularx}{0.7\textwidth}{X} 其實，敬虔加上知足就是大利。 \end{tabularx} \\ \\ \relax
6:7 & \begin{tabularx}{0.7\textwidth}{X} 因為我們沒有帶甚麼到世上來，也不能帶甚麼去； \end{tabularx} \\ \\ \relax
6:8 & \begin{tabularx}{0.7\textwidth}{X} 只要有衣有食，我們就該知足。 \end{tabularx} \\ \\ \relax
6:9 & \begin{tabularx}{0.7\textwidth}{X} 但那些想要發財的人就陷在誘惑、羅網和許多無知有害的慾望中，使人沉淪，以致敗壞和滅亡。 \end{tabularx} \\ \\ \relax
6:10 & \begin{tabularx}{0.7\textwidth}{X} 貪財是萬惡之根。有人因貪戀錢財而背離信仰，用許多愁苦把自己刺透了。 \end{tabularx} \\ \\
[1ex]
\hline
\hline
\end{longtable}
$^{1}$早安 弟兄姊妹早安 聽得清楚嗎?.
好.
這個大家應該很熟悉的這個標誌.
馬上就會聯想起金錢.
我也是在Google搜尋去找的時候.
發現原來2022年.
全世界有18億的人口.
他們每天的收入只有不超過30元港幣.
18億人口.
其實現在全世界有多少人口?.
是80億 80億人口.
想不想像得到有18億人口.
每天的生活是這麼低微.
接著我又上網去搜尋一下全世界最有錢的人.
根據2023年頭5位的.
他們擁有的 不是每天.
是說他們自己的身家.
有超過1000億.
頭5名都超過1000億美金.
很難想像這個貧富懸殊.
但是這個貧富懸殊的情況的確不只是今時今日.
是在發生在2000多年前.
在保羅的那個年代.
當時他是帶著.
他視他為一個屬靈的兒子提摩泰.
去到以弗所的時候.
傳道傳道到那裡.
建立了教會.
接著就留下了提摩泰在那裡繼續默會.
他自己就繼續進行宣教的行程.
但是他都不是很安樂.
所以到了某個時候他就要寫一封信.
回去給他的提摩泰.
原因就是當時的以弗所教會那個地方.
以弗所是一個都是貧富懸殊.
因為是一個商業.
各樣的貿易都是很發達的一個地方.
所以不難發現會有這個.
出現貧富懸殊的情況.
不過今天的講題是提到知足.

$^{41}$如果你問知足的人.
是不是一定是有錢人.
或者一定是貧窮的人呢?.
又不一定.
有錢人可以很知足.
貧窮人也不一定是不知足.
反過來如果用知足對比就是貪財.
那是不是有錢人就一定很貪財呢?.
又或者一定不貪財呢?.
貧窮人是不是一定很貪財呢?.
又或者不貪財呢?.
所以最後的結論其實就不關貧富的問題.
知足或者貪財.
不等於富有或者貧窮.
所以保羅他寫這封信給提摩泰.
他都預了周邊的以弗所教會信徒都會看得到的.
而這封信就去到第六章的時候.
即是今天看的六章三至十節.
已經是這封信的最後一個段落了.
在這個時候他特別提出了兩種人.
和提出了兩種的心態.
大家要提醒.
保羅為什麼特別要提及知足.
去到這封信的尾端的時候呢?.
原來知足和敬虔又有很密切的關係.
「敬虔」這個字英文是Godliness.
在提摩泰前書裡面出現的次數是達到九次這麼多.
其實應該說在新約裡面的所有書信裡面.
出現敬虔最多的次數就是在這本書裡.
而「敬虔」其實是什麼意思呢?.
「敬虔」是代表一個尊崇或者敬畏.
尤其是指的是對上帝的敬畏.
既然「敬虔」是這麼重要.
但是竟然在當時的已忽所教會.
就出現了一批人.
他們的想法是什麼呢?.
他們在六章的五節記載.
他們以敬虔為得利的門路.
這句話其實是什麼叫以敬虔為得利的門路呢?.
意思就是說利用敬虔來發財.

$^{81}$英文就譯作得利是一個financial gain.
但是保羅就想告訴你 你搞錯了.
你以為你想發財是透過利用敬虔可以令你發財嗎?.
搞錯了.
接著他在第六節就馬上指出.
他說敬虔要加上知足.
才是真正的得到最大的利益.
真正的獲利.
英文就譯作大利是great gain.
真正的利益就是來自敬虔加知足.
所以保羅就想告訴你.
其實第五節裡面所指的那些.
想利用敬虔的那些是一幫什麼人呢?.
然後第六節裡面說.
你要做一個敬虔加上有知足的人.
其實又是指什麼人可以有敬虔有知足呢?.
所以在六章的三到十一節.
今天看的經文裡面.
你就會分別看到這兩種人.
一種被稱為背離信仰的人.
一種叫做屬神的人.
另外兩種的心態.
一個就是貪財的心態.
和知足的心態.
好 讓我們先祈禱.
上帝啊 我們真的來到你面前.
求你賜下你的恩賢.
幫助提醒我們.
主啊 雖然這個好像很熟的.
老生常談的一個主題.
但主啊 求你仍然對我們講說話.
保守我們的心.
在你面前清潔.
聽到你要我們聽到的.
幫助我們走你要我們走出來的道.
奉耶穌基督的聖命而求.
阿們.
這段經文其實是教會的安排.
按我自己的選擇.
我未必會考慮選擇這段經文.

$^{121}$因為一看的時候.
我起初覺得.
「知啦」.
聽過很多次了.
我相信在在的弟兄姊妹.
都未必是第一次聽這段經文.
未必是很新鮮.
而且其實它講得很白的.
阿保羅很直接講出來.
但是當我一路看的時候.
就發現原來真的指出兩種人.
兩種的心.
第一種人就是一個背離信仰的人.
在六章的三至五節.
不如邀請大家讀出來好嗎?.
第三節到第五節.
一二三.
若有人傳別的教義.
不符合我們主耶穌基督純正的話.
與合乎敬虔的道.
他是自高自大.
一無所知.
專好爭辯.
善於竊賤.
因此紛爭,偽旁,惡意猜疑.
和心術不正.
與喪失真理的人.
不停爭吵.
以敬虔為得利的門路.
好 謝謝.
你有沒有留意呢?.
他說指出有些人是「傳別的教義」.
「傳別的教義」.
其實「傳別的教義」有另一個和合本.
他就譯作「傳異教」.
又或者有些聖經在這段經文有一個主題.
就說他們是一個「假教師」.
False teachers.
這班假教師.
就是說他們是教師.

$^{161}$你想像得到嗎?.
這班人是回教會的.
在教會裡擔當教導的角色.
但偏偏他們所說的教導內容.
很有問題.
出現什麼問題呢?.
至少有三個問題.
第一就是.
不符合我們主耶穌基督純正的話語.
意思就是說.
不是照主耶穌的教導來教導.
也就是說.
同生的信經裡記載的那些說話.
使徒的教導.
是背道而馳的.
然後.
第二個問題就是不合乎經見.
意思就是說.
他教的東西.
不會幫助到信徒去敬畏上帝.
更加認識更加愛上帝.
第三個問題.
他的教導會帶來紛爭.
帶來教會的分裂.
信徒之間關係破壞.
出現一些不信任.
互相猜疑.
甚至爭吵.
即是吵架.
這是極度嚴重的問題.
所以保羅就是要提醒提摩太.
如果有人傳的教導.
會帶來這三個後果.
就是不是說耶穌的教導.
會使人不去敬畏上帝.
或者是引來教會信徒之間的分裂.
紛爭 嫉妒 毀謗.
這些教導要小心.
絕對這些人.
是假的教師.

$^{201}$不是真的.
所以第五節再進一步.
他說那些心術不正.
喪失真理的人.
不停地爭吵.
以為這樣就顯得他們很敬虔.
也以為這樣.
使他們獲得一些很大的利益.
但是保羅就指出.
搞錯了.
這是錯的.
其實保羅不單在第六章提到這些人.
其實他在第四章時.
已經提到了.
聖靈明說在末後的時期.
一定有人是離棄信仰.
去聽從那些有惑人的邪靈.
和鬼魔的教訓.
所以今天我們真的要很小心.
無論你選擇聽什麼教導.
真的要用開頭保羅那三個指標.
去幫助你去釐清.
到底我今天聽到的,學到的.
是否符合聖經,耶穌基督的教訓呢?.
是否可以帶來我們合一呢?.
是否令我們更加愛上帝,敬畏上帝呢?.
接著到第十節.
保羅也有再提一次.
他說這些人是貪戀錢財.
而背離信仰.
原來錢財和信仰.
可以成為一個.
當你愛錢財的時候.
你就會離開你的信仰.
所以這些人.
是背離信仰的人要小心.
如果你說要怎樣才不會背離信仰呢?.
自己呀.
當然你要熟悉聖經.
無論在香港或者我記得.

$^{241}$我那幾年在尼泊爾生活的時候.
都發現教會很快就會分裂.
它不是因為成長而去分堂.
而是因為有些人聽了之後.
不想再聽你們的教會.
自立門戶.
這樣就導致教會分裂.
所以如果要避免這樣.
唯一一個最重要的途徑.
就是自己看好聖經.
自己去勤力去明白聖經.
而不是靠聽其他人的教導.
自己一邊聽的時候.
也要看聖經,讀聖經.
明白裡面真正的意思.
因為聖經裡面也有一句話提到.
它說「聖靈的寶劍就是上帝的道」.
要有聖靈的寶劍.
就是上帝的道.
才可以對付的惡魔.
所以要跟隨真道.
你才可以離開這些假的教導.
好 說完背離信仰的人.
保羅又指出.
剛剛相反.
有一種人叫「屬神的人」.
第十一節.
又請大家讀出來.
一二三.
「但你這屬神的人啊.
要逃避這些事.
要追求公義,敬崎,信心,愛心,忍耐,溫柔」.
原來這個「你」.
就是保羅想特別要提.
保羅要特別提.
特別提「提摩太」.
既然教會裡面已經出現了一些假教師.
但是你提摩太.
你是屬神的人.
所以你就要逃避這些事.

$^{281}$而且你要追求有以下的這些品質,品格.
同樣其實這句話.
不單止保羅想提醒提摩太.
也都提醒當時的信徒.
同樣今天提醒我們.
如果我們要成為一個屬神的人.
其實我們就要追求以下這些品質,品格.
包括要公義,敬崎,有信心,有愛心,忍耐,和溫柔.
所以再一次提到敬崎.
在四章七節.
保羅提到要棄絕世俗的言語和老虎荒謬的傳說.
要在敬虔上操練自己.
一開始提到有一班人背離信仰.
但是剛剛相反.
我們不是背離信仰.
而是我們要背離或棄絕這些世俗的言語.
這些無稽的傳說.
有時很多人負責教導.
但如果只是有好的口才.
或者說得很動聽的故事.
但是離開了基督本身的道理.
這個就要小心了.
所以這裡提到.
如果你要成為一個敬虔的人.
必定是要棄絕假的教導.
而去跟隨真理.
要成為一個屬神的人.
其實是一個跟神有很密切的關係.
那種密切的關係.
當你想到你是屬於誰的.
我是屬於上帝.
即是那個關係就像葡萄樹和枝子一樣.
我是枝子.
我是屬於這棵樹的.
我是屬於葡萄樹的.
主在我裡面.
我在祂裡面.
其實這是一個很深的神學觀念.
但同一時間也是一個很奧妙的關係.
神很大很偉大.

$^{321}$但祂大到願意住在你的心裡面.
而我們其實是罪人.
我們很渺小.
但神願意赦免我們的罪.
讓我們住在祂裡面.
有時你再想深一層.
其實真的很不配.
神的聖潔高崇.
但容許我們住在祂裡面.
祂的偉大.
願意住在我們這麼微小的人裡面.
這種枝子和葡萄樹的密切關係.
其實我不知道你每次想到的時候.
會不會覺得很不配.
或者很榮幸.
的確是很榮幸.
我們應該很感恩.
因為我們是神很寵愛.
很寶貴.
是用寶血救贖回來.
對不起,今早不正確唱歌.
非常感動.
特別在想到主耶穌為我們寫明.
想到主耶穌為我們寫明.
其實不是一些.
不只是唱歌.
不只是一些好像很熟悉的道理.
而是當你用心去感受.
主對你的愛.
祂讓我們屬於祂.
這件事要好好珍惜.
亦都很寶貴.
我們都要提防的.
就是要有兩個心態.
一個是貪財.
一個是知足.
貪財其實是指甚麼呢?.
請大家讀出來.
第九節和十字.
請你們大聲讀出來.

$^{361}$1,2,3.
(念念不忘).
這裡提到離開真道.
其實貪財的心會令人離開真道.
背離信仰.
其實背離信仰離開真道.
其實是離開了上帝.
原本你可以屬於上帝.
但你偏偏因為財利.
令你離開了上帝.
因為你不服那個信仰.
你不服從上帝的話.
聽不入耳.
你就離開了上帝.
貪財的意思.
其實就是貪圖錢財.
意思就是貪婪或貪心.
你都會問.
怎樣才算是貪財呢?.
如果我是做生意的.
我沒理由不想賺錢.
我應該賺多少才算是.
剛剛OK.
不是貪.
而是合理地賺取錢財呢?.
我很喜歡投資.
或者我很懂得投資.
我怎樣利用投資.
令自己得到最多的利益.
這樣又算不算貪財呢?.
其實不是在說財富的多與少.
而是在說你的財富在哪裡.
聖經也有提到.
當你的財富在你的心裡.
和在天上是兩樣不同的事情.
或者當你的財富.
只可以停留在地上.
這些財富其實是短暫.
會被蟲咬.
會被蟲蟲偷去.

$^{401}$會消失.
但是當這些財富積存在天上.
是永存,永恆的.
其實如果你真的是一個懂得投資的人.
你更加應該認真考慮.
我應該怎樣將我的財富,財寶.
去投資在一個可以存到永恆的帳戶裡.
將心思全部放在錢財上的時候.
你很容易所有人生的決定.
就是就著那些錢財來做決定,做抉擇.
這樣問題就出來了.
第十節那裡很直接指出.
貪財是萬惡之根.
我記得很多年前我母會有一個姐妹.
她新年要封利是.
她突然間看到這段經文.
她就說封利是好像.
封多少錢才合理,合適呢?.
最後她就說.
她開玩笑地說.
不如就弄一張書簽裡面寫著貪財是萬惡之根.
就放進這些利是封裡面.
等人收了之後馬上很醒覺.
當然這是她開玩笑.
我也不相信她真的會這樣做.
不過真的是貪財是萬惡之根.
意思就是說有很多惡事是從貪財開始的.
貪財下面是根部.
然後萌生上去就出現很多罪惡.
其實今時今日很多詐騙案件.
或者利用別人的信用卡去盜取錢財.
又或者簡單點去街市,商店買東西被人騙.
其實有很多這些為了賺錢的方式.
其實都是源於因為貪財.
所以就去犯罪.
或者最過分的.
你會發現最近有一個兇案.
是為了幾千萬的樓宇而去設局.
去謀財害命.
這個案件連我在非華語生的群體裡面去服侍.

$^{441}$有一個學生跟我說.
「哎呀 什麼事啊 竟然可以拿骨頭來煲湯啊」.
「真的很恐怖啊」.
我說「是啊」.
你想像不到原來貪財到一個地步.
可以做出這樣的事.
或者你說 是不是只是成人會貪財呢?.
其實在學校裡面.
盜竊的事件都會來來回回出現.
有同學說明明袋了十塊錢.
放在抽屜裡 怎麼轉眼間就消失了呢?.
其他同學就說「你當然是忘記了」.
我說「不是啊 真的不是忘記了」.
原來貪財可以是不分年齡 性別 形式.
總之各司各法 你如我詐 呃呃哄哄.
其實目的只是為了自己得到最大的利益.
就不惜奪去別人原本所擁有.
其實貪心的意思就是說.
你把一些不屬於你的東西歸於自己.
在十誡裡面很清楚地說.
不可貪心 不可貪戀別人的房子.
不可貪戀別人的妻子.
其實貪就是貪了.
那貪的結果是什麼呢?.
第十節這裡說得很清楚.
他說貪的結果就是使人沉淪 敗壞 滅亡.
用很多的仇苦令自己刺透了.
其實貪是會越貪越深.
去到泥足深陷不能自拔的時候.
帶來就是這個永遠的滅亡.
因為離開了上帝.
相反 另外一種心就是一個知足的心.
第七第八節請大家讀出來.
一二三.
(念經).
有衣有食 就該知足.
知足的意思 我看回一些解釋.
它的意思就是說.
對已經得到的事情 事物感到滿足.
這樣就是知足了.

$^{481}$英文也提到.
(念經).
就是你現在擁有什麼.
你就是為了那些你已經感覺到快樂.
這種感覺就是一種滿足 知足的快樂.
我也提到.
我是在非華裔學生的學校工作裡面去做的.
所以我很多時候.
其實差不多星期二到五.
除非真的出去探訪或者做其它事情.
不然的話 我都會在學校的小息.
或者午飯的時間.
或者放學之後就去接觸學生.
有一天放學之後.
我走過操場.
就看到學生在玩什麼呢.
他們三個是尼泊爾的中一的學生.
尼泊爾人.
他們就將一盒橡筋 很大盒橡筋.
捲起來就成了一個球體.
我就想 我小時候好像沒有這樣的玩法.
他們就用這個球體.
就可以拋高 或者去扔人.
總之你扔到就扔不到人 或者扔到人.
我也不明白他們的規則是什麼.
總之他們都會哈哈大笑.
笑得很開心.
開心到有些老師走過.
就會站在那裡看著他們三個兒子扔來扔去.
當然第一個想法就是.
這個好像小學一樣.
好像小學生的玩意.
但另一方面你會看到.
其實這個玩具 我問他多少錢.
他說9塊錢買一盒橡筋.
很便宜的 自製的玩具.
不是很多 什麼技巧.
就拋來拋去 笑聲不斷.
也不是在玩什麼競賽 要分輸贏.
沒有很單純 就在那裡拋來拋去.

$^{521}$這樣就很開心了.
那個開心真的很簡單.
那個滿足也是來自一些很簡單的.
不是要一些什麼遊戲機.
或者一些很密集的享受.
知足就是要明白自己所擁有的.
本來就不是來自自己的.
也不能夠永遠屬於自己.
這裡提到我們沒有帶什麼來到世界上.
也不能夠帶什麼走.
不知道你們有沒有聽過一個故事.
關於亞歷山大帝.
他可以打很多仗.
很厲害的一個皇帝.
可以擁有很多的土地.
但是據聞當他要離世之前的時候.
他就吩咐他的下屬要幫他做一個棺材.
而這個棺材的左右兩旁都要開一個洞.
原因就是將他的雙手.
可以從這個棺材裡伸出來.
他的意思是想叫世人知道.
就算他生前可以得到多少的土地.
財富,名利 什麼都好.
但是當他一離世的時候.
一樣都帶不走.
兩手空空的就離開.
在舊約聖經有一個人物叫做約伯.
大家都很熟悉.
他原本很富有.
但是因為某些原因之下.
可以在很短的時間之下.
他失去了所有的財富.
但是在這個時間.
他不僅不是埋怨.
反而他講出了一句名言.
就是說上次是耶和華.
收取的也是耶和華.
是應當稱重的.
可以這麼豁達,這麼豁然.
看得開.

$^{561}$其實就是因為他是一個知足的人.
他知道他所擁有的.
不是他生出來就應該有的.
也都知道不是永遠屬於他.
所以知足的人看的.
不是現時眼前短暫.
看得見好像有益的東西.
知足的人是看得到在永恆裡面.
永遠長存有益的東西.
有人為了心頭好.
會不停地去賺錢.
努力賺錢 追求 追求 追求.
在這種動力之下.
原來只會快的話.
越來越不夠 越來越不滿足.
有時候都會想想去旅行.
去多少次旅行才叫做足夠呢?.
要買車 買多少輛車才叫做足夠呢?.
要換更大的房屋.
要去到多大才叫做足夠呢?.
原來越推動這種追求心頭好.
或者物慾 或者享受的推動力之下.
發現越來越不夠 越來越不滿足.
慢慢地就變成了事奉馬門.
其實聖經裡面都很清楚直接告訴我們.
當你追求馬門的時候.
你就不會追求上帝.
當你的心放在財富上.
你的心就用不下上帝.
但相反 當你的心被上帝充滿的時候.
這些財富物慾就放不下你的心裡.
其實我們的心很小的.
你是二選一的.
你放在哪裡就放在哪裡.
聖經裡面提到.
其實真正的滿足就是知足.
只要我們以上帝為樂成為我們人生的目標.
你就會發覺上帝在你裡面是越來越大.
你也會在上帝裡面發現自己其實真的很渺小.
葡萄樹和枝子這種緊密相連的關係.

$^{601}$就是一種黏著葡萄樹就已經是滿足快樂.
最後 知足的結果是甚麼呢.
「經乾加上知足就是大利」.
原來敬畏上帝加上滿心的快樂知足.
才是最大的得益 最大的獲利.
讓我們一起祈禱.
親愛的主 我們真是願意以你為我們的手.
為我們的主 為我們生命的王.
求你在我們心中居首位.
讓我們事事以你為最大的滿足和快樂.
因為神啊 你是那位口齒白密的神.
在你裡面我們一無缺乏.
求你幫助我們事事數算我們所擁有的.
就為了這樣我們可以凡事謝恩 時時的喜樂.
主啊 如果我們真的有缺乏.
有現在真的陷入財政的困難裡面.
主啊 其實你知道 求你幫助.
你幫助弟兄姊妹 看到那個出路.
不是靠我們怎樣去努力 去計劃 去想 去做些什麼.
求你叫我們只要專心去仰望 投靠你自己.
你就可以幫助我們面對困難 渡過難關.
同樣如果我們現在活在一個很風俗的境況裡面.
求你就更加叫我們看到我們身邊有需要的人.
幫助我們不單只顧自己 也顧別人.
求天父繼續去帶領 幫助我們聽道 行道.
祈禱 交託 奉靠主耶穌基督的聖名 阿們.
最後有一節經文送給大家的.
就是你們要全心不可貪愛錢財.
要以自己所有的為足.
因為主曾說 我總不撇下你 也不丟棄你.
\newpage



\section{}
\label{sec:yk5Oh8RBHLY}
\textbf{2023-03-30 周四晚查經班《被擄歸回歷史研讀(9)》 講員:曾裔貴牧師}
\newline
\newline
連結: \href{https://youtube.com/watch?v=yk5Oh8RBHLY}{\texttt{https://youtube.com/watch?v=yk5Oh8RBHLY}} ~~~~ 語音日期: 2023-04-01
\newline
\newline
\hyperref[sec:pzC3e8fjAp8]{\small{< < < PREV SERMON < < <}}
~
\hyperref[sec:index_chronic]{\small{[返順時目]}}
~
\hyperref[sec:index_scriptual]{\small{[返順卷目]}}
~
\hyperref[sec:x06IlO4FHos]{\small{> > > NEXT SERMON > > >}}
\newline
\newline
$^{1}$我們一起禱告.
多謝天父給我們.
雖然停了幾個禮拜.
但是我們都能夠繼續.
有系統的研讀聖經.
尤其是一些舊約的書卷.
我們平時都.
可能在港台都比較少.
能夠去整卷這樣看的.
都幫助我們能夠明白.
特別是一個被老歸回的時代.
我們能夠掌握.
聖經給我們的資料.
我們看到有.
不同的時代.
他們對神的.
效忠.
的擺上.
和帶來的是.
神的家.
當時的以色列這個民族.
重新能夠有一個.
憲制的制度.
有那個.
國度的等候.
其實就是為了耶穌基督降生.
來去預備的.
可能當時代的人.
都沒有想得那麼遠.
但是當我們新約.
看回幾千年前的歷史.
我們就看到.
背後是有神去管理的.
今晚我們看.
尼希米記.
盼望聖經給我們.
的這個人物.
如果是這麼重要的時候.
他一定是有一些.
心靈的價值觀.

$^{41}$值得我們去效法.
去跟隨的.
求主都帶領我們明白.
我們這樣祈禱.
奉耶穌基督的名.
阿們.
有幾個禮拜.
因為公幹.
去了歐洲和開會.
和探望宣教士.
所以就不在大家中間.
大家不知道.
有沒有在以色列記.
去重溫.
和在哈該書.
那個被擄歸回的.
那個先知書.
來到去有少少的重溫.
今天我們要開一個新的書卷.
大部分的解經家.
都認為尼希米.
這位聖經人物.
是和以色列.
是同時期的.
尼希米記.
相對比較容易看.
因為都是工整的.
一到第六章.
比較像是他的自傳.
他為什麼.
要回去重建.
這個聖城.
和第七章.
之後.
去到第十三章.
結束.
他怎樣來到去.
透過.
宣讀律法書.
來到去可以叫重建聖文.

$^{81}$即是神的子民.
這個是.
尼希米記給我們看到的.
那個分段.
當然夾雜著的.
就有很個人性的.
他的自轉式.
的祈禱.
和面對一些的難處.
他怎樣來到去周旋.
在那些.
耶路撒冷那些外邦人.
對他的.
這個宮城的敵對.
阻撓.
他怎樣能夠來到去化解.
這些裡面都是.
涉及了很多他.
倚靠神又見到有神保守.
當然更重要.
就是有.
百姓認同他的.
事奉.
所以來到去.
跟隨他.
配合他的策略.
有很多去到第六.
第七章.
都有一些值得我們去想一想.
他面對一些敵人對他的攻擊.
可能是用一些批評.
或者是一些獵造.
一些惡的.
說話.
他都來到去.
很有效地.
很有智慧地.
去解除.
這些作為領袖都需要學習.
當然這本書.

$^{121}$主要不是教我們怎樣去領導.
不過就很多人.
可能神學.
會用這本書.
來教領導學.
怎樣做一個領袖.
其實這本書不是講.
教人做領袖的.
這本書是一個.
信徒領袖.
他.
用我們現在新約能夠明白的說話.
一位領袖.
其實他本來就不應該.
背這麼重的重機.
來到去.
耶路撒冷做那個工程.
但是.
似乎在耶路撒冷的.
宗教領袖.
或者其他的領袖.
在當地在聖城的.
都已經失去了那個功能.
所以神就在外邦.
在皇帝身邊的.
一個紅人.
一個.
叫做幾高的官職.
一個酒政.
在皇帝面前.
來到去.
試那些酒菜.
很信任的一個神子.
才可以來到去.
服侍皇的飲食.
所以酒政.
都不是普通的一個位置.
所以.
都不是普通的一個位置.
差不多好像大內總管.

$^{161}$但是居然神就興起.
一個在外邦.
書山城.
似乎是波斯帝國.
的那些行宮.
來到去.
服侍皇.
這個皇帝亞達薩西王.
他知道了耶路撒冷的事情.
他就有一個.
的決心.
祈禱了四個月.
一路等待時機.
然後就上耶路撒冷.
他都不會.
因為他自己這個官職的身份.
他就放棄了.
而是他仍然.
用皇帝給他的.
權和.
那個俸祿.
和他的.
拿到皇帝給他的.
那些材料.
來建造耶路撒冷的城牆.
這個就是尼希米基.
一到第六章.
那個經文大概的意思.
當然裡面就涉及.
他是怎樣.
有智慧地去.
一個都幾陌生人.
上到去耶路撒冷.
到底人怎樣看他呢.
是看他是一個高官.
還是看他是一個弟兄.
是和自己這個民族.
有份的一個.
弟兄來信任他.
信託他.

$^{201}$跟隨他.
去建造那個城牆.
周圍連絡.
將那些倒塌了的城牆.
全部連絡.
就是剛剛一個surroundings.
就是一個四圍的圈.
做好了.
將那些門安好.
totally是52日.
就做好了.
這是第六章給我們的經文.
告訴我們.
但是那52日其實是.
不容易的.
又要拿刀又要拿那些材料.
工具.
因為敵人隨時就.
上來去威嚇.
上去阻攔.
但你起碼就.
勉勵他自己的同胞.
你可以想像得到.
一個第三代.
在外邦出世的.
不是耶路撒冷的人.
但在外邦.
已經是做到.
比較高官.
但都能夠取得.
耶路撒冷的一些.
鄉親父老的信任.
你可以想像得到.
尼欽米無論在外邦.
人的名聲.
品格.
信任度.
他上去之後.
又能夠很有條理地.
講解他的.

$^{241}$初心.
他的熱心.
去鼓勵.
耶路撒冷的弟兄姊妹.
一起和他.
興功建造.
這個不容易的.
因為你是一個空降的.
高官.
我們讀讀經文.
尼欽米記的第一章.
看第七章的.
以色列記.
幫大家有少少的記憶.
第七章就是開始.
記載.
耶路撒冷.
上去.
極相似.
就是描述.
又是講.
他自己是何許人.
和為什麼會上去.
之後會帶來.
耶路撒冷的一個復興.
就是看到.
那個情況.
不是如他所期望.
因為他是一個很嚴格.
受訓練的一個宗教的領袖.
一個文士.
特別強調他是文士.
這個身份的.
以色列第七章的.
過了一段時間.
聖殿.
已經建好了.
已經完成了.
但是過了一大段時間.
之後似乎在聖殿裡面.

$^{281}$有那個硬件.
但是那個軟件.
就是說憲制的規模.
憲制的.
那個.
心靈不是.
以色列所看到的情況.
所以勾起.
他發現原來.
很多的祭司都還沒有.
準備好.
去準備,所以他就要.
尋查那些家譜.
所以你會看到第七章第八章.
都是想帶出.
哪些真是具備.
這些宗教上.
神規定了的.
那些祭司.
利未人,他們就要.
回到自己的崗位.
那就會帶來.
那個憲制.
就會慢慢有那個熱誠.
有那個.
真心,百姓就會.
慢慢來投入.
去奉獻.
事奉神的,所以.
都是仍然回到.
事奉人員的內心.
是不是真的對準了神.
的恩典和.
那個積分,這個就是.
告訴我們,以色列.
是在波斯王亞達舍西.
的年間,雖然就沒有.
強調是哪一年.
他回去,接著很快.
去第七節,就是亞達舍西.

$^{321}$王的第七年.
以色列人.
祭司利未人,歌唱.
的守門的尼提靈.
有上耶路撒冷的.
王第七年的五月.
以色列到了耶路撒冷.
正月初一.
他從巴比倫起程.
走了五個月,原來.
第七年的時間.
就是亞達舍西王的時間.
我們如果回去看.
這個.
尼希米記.
就很清楚.
告訴我們,原來.
以色列是,你可以叫.
第二批.
過了很長時間.
他們上去,以色列.
就帶來一個復興,因為.
他是一個宗教領袖.
他是一個文士.
當然他自己本身都是祭司.
但他更加強調.
他是一個很敏捷的.
文士,即是說.
很熟讀摩西的律法書.
很能夠融會貫通.
很能夠將一些.
重要的真理.
講清楚,講明白.
給百姓知道.
知道哪些是重要.
哪些是必須要去進行.
來到尼希米記.
這裡就強調是二十年.
即是說,大概.
十三年之後.

$^{361}$又是完成了.
那個復興.
但是,慢慢慢慢.
好像那個復興.
又不能夠再持續.
又見到.
在第二章尼希米記.
開始告訴我們.
尼希米上到.
耶路撒冷的時候.
比他想像的困難.
更加嚴重.
因為.
不只是表面的.
破損,表面的.
日久失修.
真正的問題.
是外邦人.
對猶太人.
的阻礙.
以致猶太人.
都慢慢.
不能夠.
很同心地.
來到去,大家.
一起去將.
本來屬於我們自己的環境.
即是說.
四圍的城牆.
來聯絡.
即是說,外邦人.
他們可能都會在.
不合適的時間.
進入耶路撒冷.
來做買賣.
可能和猶太人.
有很多的來往.
商業的貿易.
又或者可能談婚論嫁.
的情況,就會帶來.

$^{401}$猶太人慢慢.
又會落在.
外邦人的.
思想水平.
或者信仰.
的混亂裡面.
就會.
帶來一個比較根本性.
對可能是.
下一代,或者再.
稍為延後的.
猶太子女.
會有影響.
所以,這都是.
我們要自己思考.
當然,拆開來說.
這次我在.
英國開會.
在開會的地方附近.
就已經是很.
因為英國曼城,人們喝酒的文化.
當然他們喝啤酒.
的文化.
很強烈.
酒店附近.
曼城的市中心.
已經是一個.
整個酒吧區裡面.
是一個同志村.
一個Gate Village.
就是.
根本是那些同志的.
組織的總部.
大家已經見怪不怪.
整個環境.
用了彩虹的顏色.
來佈置.
甚至乎有很多是遊客.
去拍照.
可想而知.

$^{441}$那種不能夠.
對其他人.
的這個.
做法.
有任何的可以歧視.
已經塑造.
這樣的文化.
如果你的子女在外國長大.
在英國.
長大,他可能.
反過來跟你說.
爸爸媽媽.
不能歧視我的朋友.
或者.
不同的取向的人.
他們的思想.
這些就是.
一定要在一個這樣的環境.
你就會受到.
影響.
我們看經文裡面.
一到第六章.
是尼熙米.
幾重的個人的自轉式.
他有那個心意.
要為耶路撒冷.
來做一些事.
做一些事是很空泛的.
他要祈禱.
他要等候.
因為他的正職是.
做王身邊的一個.
管理飲食的一個酒政.
就是說.
要令到王飲食.
至少不會吃壞了他.
可能很多宮廷的.
國宴.
他要去準備.
所以可能要試酒菜.

$^{481}$因為.
皇帝通常被其他敵國的人.
隨時都想買起他.
暗殺.
所以那個酒政其實.
幾重要的古代.
我相信現代一樣的.
普京的飲食.
其他皇帝的飲食.
習近平的飲食.
都應該是有專人去看的.
我想不會普通的.
隨便買個飯盒.
一定不可能的.
因為我們知道.
這些都是重要的.
在王身邊的崗位.
如果不信任.
其實很難做到這些位.
記住他是一個外邦的猶太人.
第三代的猶太人.
哈加利亞.
聖經沒有說到.
這個哈加利亞.
到底是什麼背景.
只提到尼希米的爸爸.
哈加利亞.
當然希伯來文.
是解釋大概的含義.
例如耶和華.
是怎樣的.
諸如此類.
尼希米就是耶和華的安慰.
哈加利亞的兒子.
尼希米的言語如下.
言語就像一個回憶錄.
就像寫一個自傳的日記.
亞特什西王二十年.
基斯留月.
我就在網上.

$^{521}$很容易找到的.
就是一個月份.
猶太人自己有一個月份.
我們都有.
正月,我們有潤月.
當然潤月,虧月就是.
月圓月缺.
現在這裡告訴我們.
基斯留月.
相對於陽曆.
就是我們現在的新曆.
就是十一,十二月.
因為他們是用月頭的.
二十九,三十日.
是用農曆的.
猶太人.
所以就有這個.
可能在十一月.
或者十二月.
所以這個基斯留月.
就特別強調.
在這個年間.
十一,十二月.
就是哪個日子.
會有一些的節日.
就是普爾節.
或者是.
對不起.
可能是哈路卡.
就是那個縱光節.
會有在這個月份.
會出現的.
當然.
當時就沒有縱光節.
因為現在是主前的四百多年.
四百四十五年.
或者四百四十四年.
相差一兩個月.
很難計算.
但是縱光節.

$^{561}$或者叫修殿節.
就是主前的一百六十五年.
所以根本沒有出現.
所以就沒有縱光節.
我在書山城的宮中.
那時.
有我一個弟兄哈拿尼.
同著幾個人.
從猶大來.
我問他們那些被老歸回.
盛夏逃脫的猶大人.
和耶路撒冷的光景.
他們對我說.
那些被老歸回.
盛夏的人.
在猶大省遭大能.
受凌辱.
並且耶路撒冷的.
城牆拆毀.
城門被火焚燒.
到底有什麼的.
暴動.
有什麼的問題.
導致到為什麼會日久失修.
為什麼會有一個.
大能凌辱.
是不是外邦人.
借一些機會.
來去羞辱.
攻擊猶太人.
聖經雖然沒有記載.
哪個具體的事件.
但是.
在他們.
一問的情況下.
他們就立即.
講述了.
最近現在.
耶路撒冷出現的問題.
和破損.

$^{601}$的情況.
就是說在尼熙米.
的內心裡面.
他一直記掛著.
耶路撒冷.
這個問題要出現.
到底在外邦的.
猶太人.
他們沒有上去耶路撒冷.
他們是不是從來都沒有上去.
這是他們的故鄉.
未必是他們出生.
有沒有上去耶路撒冷.
還有是不是所有在外邦.
已經落地生根.
去到第二第三代的.
猶太人.
他們會有這個很強的.
耶路撒冷聖殿的意識.
如果不是.
那就很特別了.
尼熙米是一個很特殊的人.
因為他.
生的地方.
現在工作的地方.
不應該讓他有.
這麼重這麼強的.
對耶路撒冷的.
依戀的心情.
如果你說普遍.
外邦人都是有一個很強烈.
他很想.
回復興.
復國.
重建.
第二次世界大戰之後.
猶太人已經分散到.
列國.
加上他們有.
大屠殺.

$^{641}$所以回到耶路撒冷.
的意識.
是很強烈的.
其實是由第一次世界大戰.
之後.
猶太復國主義.
已經有這個很強烈的.
意願.
透過這樣.
來和一些當時.
世界的領袖.
似乎是英國.
去謀求.
希望讓他們.
能夠回到巴勒斯坦地.
重建他們的國家.
一直拖.
都不能夠成就.
後來到二戰.
當然後來以色列可能很多科學家.
亦都是.
很重要的人.
研發原子彈等等.
所以.
促使.
當時英國美國的領袖.
准許.
以色列人回到.
耶路撒冷.
重建他們的國家.
1948年.
是否5月14日.
如果你們找到就幫幫我.
我記得好像是.
5月14日.
你就會感受到.
那個時代.
每個都猶太人.
他們自己家裡會有一個盒子.
大家就是家家戶戶.

$^{681}$都存錢.
回去復國.
你想想那個人心的.
準備.
是全世界.
無論你是埃塞俄比亞.
蘇聯.
不同國籍的猶太人.
因為已經散居在.
各個地方.
大家都有很強烈.
要回去復國.
這樣的意欲.
所以你想想不是很久.
七八十年前的事.
你可以想像得到.
這個就是猶太人.
但在尼希米的時候.
其實不是很強.
所以就變了尼希米.
就是一個很特殊.
被神.
用的一個人.
神可以在外邦.
在一個本來.
應該是服侍王.
要王時時都在宮廷的飲宴.
飲得開心.
食得滿足.
的一個官.
你想想.
你可能真是每天都要品酒.
哪些的.
紅酒.
哪些的白酒.
哪些的新的菜式.
你要.
見的是最好的東西.
那個擺設.
是要怎樣的.

$^{721}$美麗.
王喜歡什麼的口味.
是鹹一點.
還是他喜歡吃甜的.
多一點.
吃甜的就不會覺得飽.
一般人都是.
就是吃完飯之後.
要不要糖水.
要啊 又說很飽.
糖水就不會飽.
因為一般人都會這樣.
但是亞達斯西王是不是呢.
你作為酒精.
你又要揣摩.
你的主子的胃口.
他近來會不會.
有些腸胃不舒服.
他又要吃清淡的東西.
這個是酒精的工作.
你的心應該完全擺在.
那個皇宮裡面.
的每一個細微的細節.
但是.
這裡很快就告訴我們.
他很緊張.
就是.
他自己的.
這個民族.
最重要的地方就是耶路撒冷.
最重要的當然就是.
維繫於那個聖殿.
當然聖殿.
只是那個硬件.
最重要的就是神的選民.
能不能夠在那個.
聖殿裡面去敬拜神.
所以這些是.
我們大家都應該見微知著.
有一段時間.

$^{761}$就是.
去年的雨季.
很嚴重的大雨.
第二天早上.
我一早就馬上回來.
打完風之後馬上回來.
小玲很快就回來了.
一回來就看看風琴.
有沒有入水 果然全都浸了.
就馬上一起.
拿毛巾不停地吸.
因為入了那些線位.
就很麻煩.
所以有時很緊張的事情.
都馬上就會想起.
教會會怎樣.
很正常的.
你就可想而知.
李希米他.
那個內心.
其實他不是一個很空閒的人.
然後第四節.
我聽見者說.
坐下哭泣悲哀幾日.
在天上的神面前.
禁食祈禱說.
耶和華天上的神.
大而可畏的神.
你向愛你.
守你誡命的人.
守施慈愛.
願你睜眼看.
側耳聽.
你僕人晝夜在你面前.
為你眾僕人.
以色列民的祈禱.
承認我們.
以色列人向你.
所犯的罪.
我與我父家都有罪了.

$^{801}$突然之間.
不單止是以色列人的罪.
我與我的父家.
都好像有份於.
耶路撒冷的破損.
他沒有上耶路撒冷.
他可能是因為礙於他的成長.
他的工作.
等等,他可能只是一年.
三次上耶路撒冷首節.
不是.
全身而退的.
移民回到.
一個他自己不是出世的地方.
但是.
沒有阻礙到他.
對現在聽到.
神的電.
神的聖神.
那個破損的消息.
而令到他有一個悲哀.
再加上他將他自己.
他的父家.
都好像擺在.
那個有份於.
令到這件事情發生的.
可能是不團結.
可能是猶太人.
離開了神的律法.
所以就變成.
這個有罪.
就不是一種官式的.
一個隨便這樣說.
第七節.
我們向你所行的.
甚是邪惡.
沒有遵守你直著僕人.
摩西所吩咐的.
誡命,律例,典章.
又是強調.

$^{841}$回到十誡.
回到摩西的律法書.
回到神的聖經.
摩西五經的規範.
那個教導.
就是說.
就算你幾代人.
離開耶路撒冷.
仍然看自己.
好像是基督徒.
是對於神的聖經.
是有對我的生命的規範.
是神.
給我們的誡命.
我們要遵守的.
那你就可想而知.
那就很難得.
已經不是在.
耶路撒冷.
每個星期.
有安息日,有獻祭.
有律法師.
或者是當時的利未人.
來去教導.
但是呢.
黎希米就是.
有神的道在他的心靈裡.
所以他很緊張.
就是我們沒有遵守到.
你直著摩西.
所吩咐.
所頒佈的誡命.
是直著摩西.
就是說那些道是神的.
不是已經過期的一些條文.
而是真的有效.
在人心靈裡面的.
第八節.
求你紀念所吩咐你僕人摩西的話.
說你們若犯罪.

$^{881}$我就把你們.
分散在萬民中.
但你們若歸向我.
謹守,遵行.
我的誡命.
你們被趕散的人.
雖在天涯.
我也必從那裡.
將他們焦罪回來.
大道我所選擇.
立為我名的居所.
這都是你的僕人.
你的百姓.
就是你用大力和大能的手.
所救贖的.
主啊.
求你則以聽你僕人的祈禱.
和喜愛敬畏你名.
眾僕人的祈禱.
使你僕人.
言今亨通.
在王面前蒙恩.
為什麼不是在神面前蒙恩呢?.
為什麼要在.
王的面前蒙恩呢?.
你可想而知.
他原來都是為了下一章.
他有一個伏筆.
他現在祈禱.
不是在神面前蒙恩.
而是他.
得到神的.
祈禱神的垂聽.
是在亞達舍西王的.
面前蒙恩.
神的恩典.
就是說他祈禱.
不單是感動於天.
感動於神.
而是他的祈禱是求神.

$^{921}$去感動亞達舍西王.
那就難很多很多.
你怎麼可以.
你的祈禱可以感動到.
一個不信神的.
外邦波斯帝國的.
一個君王呢?.
如果君王是一個很敵對神.
很反對.
耶和華的神.
的一個君主.
其實你怎麼能夠去.
感動他呢?但是.
看到這件事不是.
他個人那麼簡單.
這件事是要有一個.
清晰的印證.
一個開路.
就是王都是被神.
感動而來到去.
蒙恩.
當這個王被感動蒙恩的時候.
就可以印證.
尼熙密的禱告.
他的為著民族.
為著他自己.
為著他的家族.
來向神認罪.
認罪之後.
他不是停在那裡.
他也相信神有.
挽回的恩典.
就是只要我們.
向你來去.
懇求.
你會有一個挽回的恩典.
使我們重新.
被神.
安置在耶路撒冷.
為什麼呢?.

$^{961}$因為我們是你的.
選民.
這是很重要的一個.
摩西五經給我們的.
信息.
申命記28章.
這章聖經相當長.
所以尼熙密記.
是很熟悉摩西五經.
他將精髓.
都已經擺了出來.
第28章.
這章聖經很長.
進行誡命的就.
王福.
不好意思.
因為回來就有幾天病了.
可能長期.
天炸.
又不夠睡.
回來香港可能.
一熱一點.
可能細菌會多過歐洲.
所以就會有喉嚨不舒服.
第28章大家留意.
你如果留意.
聽從耶和華.
謹守遵行.
他的一切誡命.
就是我今天所吩咐你的.
他必使到你超乎天下萬民.
之上.
你若聽從耶和華你神的話.
這以下的福必追隨你.
臨到你身上.
你在城裡必蒙福.
在田間也必蒙福.
你身所生.
地所產.
生畜所下.

$^{1001}$以及牛獨羊群.
也必蒙福.
你的糠子和磚面盤.
也必蒙福.
就是說你整個廚房.
打開就不是什麼都沒有.
不是發霉.
你看到什麼罐都裝滿了.
就是很豐足.
好像五穀豐收.
就是出有蒙福.
入也蒙福.
這些真是中國人最喜歡的.
仇敵.
起來攻擊你.
耶和華必使到他們在你面前.
被你殺敗.
他們從一條路來攻擊你.
必從七條路逃跑.
在你倉房裡.
並你手所辦的一切事上.
耶和華所命的福.
必定臨到你.
耶和華你神也要在.
你所給的地上賜福於你.
你如果謹守.
耶和華你神的誡命.
遵行他的道.
他必照著向你所起的誓.
立你作為自己的聖文.
天下萬民見你.
歸在耶和華的名下.
就要懼怕你.
你在耶和華向你列祖起誓.
應許賜你的地上.
他必使到你身所生的.
身足所下的.
地所產的.
都綽綽有餘.
多到滿到寫.

$^{1041}$這樣的意思.
耶和華必為你開天上的虎腹.
按時降雨.
在你的地上.
要下雨的時候就有雨.
要有好太陽的時候就有好太陽.
這樣就很重要.
你想想.
如果收成的時候.
不斷傾盆大雨.
東西都爛了.
菜都浸到爛了.
那就不是.
季節合時的雨.
那就不及時的雨.
但是這裡說.
按照著時令.
方便的時候就下雨.
不方便的時候就不下雨.
這樣就可以想像到.
多麼的蒙福.
在你手上所辦的一切事.
都賜福於你.
你必借給許多的國民.
多到一個地步.
可以借錢給外國人.
他不至向他們借貸.
你想想.
你家裡的子女.
不用打卡,打爆卡.
每個月只給免費.
甚至.
他根本不需要打卡.
還有很多現金在他的戶口.
大概是這個意思.
你可以想像到.
多麼的有福,多麼的蒙福.
現在.
月光族.
我的兒子說他的同事.

$^{1081}$出糧.
吃好一點.
阿明仔說不要,我帶飯.
以為他很慘.
其實不是,是節儉.
然後去到月底.
見到阿明仔又吃飯.
慘了,是不是不夠錢,沒錢了.
不是.
其實根本就不是那件事.
你可以想像到.
這裡說的是.
進行.
神旨意的一個蒙福.
但是後面.
這裡只是用了.
十四節,但是那個奏座.
是由十五節.
一直到.
我再站多一點.
一直到.
六十八節.
是說奏座.
如果不聽.
神的話語,當然.
最重要就是我們剛才看的.
那個.
尼希米基一個重要的.
Point就是,什麼時候回轉.
神就會收回.
他的憤怒.
會使到神的選民.
重新.
可以回到他們的地方.
這個.
就是神的.
祝福和奏座.
所以第28章.
其實是相當重要.
相當我們值得去思想.

$^{1121}$和去看的.
也是整個申命記的精髓.
所在,是因為.
臨離開世界的.
一位老人家.
望著一班二十歲.
的青年人.
全部等著.
入迦南,在摩亞.
境內.
摩西一百二十歲.
快要死了,對著.
這一班年青人.
因為和他一起.
出埃及的都死了.
所以告訴他們.
你們快要入迦南.
我就快要見主了.
摩西升天了.
就來到.
講述神的.
祝福和神的奏座.
就是說將神的屬性.
一定是定義私恩.
一定是留意著.
奏座的一位.
神,你們.
是要記掛這件事.
看得為重要.
這個是摩西.
他的這個雨中心祥.
的一個.
臨終的一個祝福.
一個長者的叮嚀.
因為他真的見識過神的憤怒.
不過當時的實在對象是.
那個對象是.
埃及王法老和埃及人.
和那些軍隊.
那個憤怒沒有.

$^{1161}$遷向於神的選民.
摩西也都.
經驗過.
神什麼時候.
降那些憤怒給.
以色列人.
他都經歷過,譬如可拉的事件.
就是挑戰摩西.
阿倫的權威的時候.
可拉是利美的.
一個支派.
就這樣突然間被神.
降和.
地.
突然間裂開.
他們整個可拉族.
掉下去.
地合起來.
是一個很超自然的神蹟.
我們回來看尼希米記.
尼希米他.
搞緊了一件很重要的事.
將28章濃縮了.
作為一句的禱告.
就是第八節.
你如果犯罪.
我就把你分散在萬民中.
但你們若歸向我.
即是再次回轉歸向神.
謹守.
遵行我的誡命.
你們被趕散的人.
雖在天涯,我也必從那裡.
將他們招聚回來.
大到我所選擇.
立為我名的居所.
當然就是耶路撒冷.
即是說他們可以回去重建.
神的聖城.
回去.

$^{1201}$有了那個城.
自然就有那個民族.
有他們那個根據摩西的.
律法書來去獻祭.
的那個操作.
這個宗教的操作.
其實這樣生生不息.
其實是為著.
真正的救贖的救主.
彌賽亞降臨.
鋪路的.
所以他們為何在.
主前的四百多年.
尼希米要這麼緊張.
回去重建.
那個聖殿.
聖城,聖文.
當然是分三批人.
那個原因.
就是他們.
都是有一個歷史很重要的責任.
就是彌賽亞.
何時降臨呢.
他們知道的.
他們回去就是要預備那個氛圍.
那個氣氛.
那個宗教的虔誠.
等待彌賽亞最終降臨.
這個就是.
尼希米那一代人.
他們因為在.
耶穌降生之前的人.
他們就是要等待彌賽亞的降臨.
所以他們每一個時代.
都會問.
有童女降生呢.
這個是應許.
二次亞書第七章.
那個應許.
等等.

$^{1241}$何時有那個被.
那個無罪的.
那個僕人.
公義的僕人.
來出現呢.
他們是這樣一路等候.
所以.
你可以這樣說.
這個就是尼希米他們那一代人.
他們肩負著.
一個他們要傳承的.
那個責任.
就是一定要讓耶路撒冷.
有一個氛圍.
有一個等待.
就是彌賽亞.
要在我們這個民族.
來出現.
彌賽亞那個拯救者.
是要在耶路撒冷來出現.
說得近一點.
就是在伯利恆降生的.
這個是尼加書.
告訴我們的.
尼加書的第五章.
尼加書是小先知書的其中一卷.
以摩訥拿尼.
很多人就對.
尼加書的六章.
九節就會比較熟.
但是第五章.
成群的民.
現在你要聚集成隊.
因為仇敵圍攻我們.
要用仗擊打.
以色列審判者的臉.
伯利恆以發他.
你在猶大諸城中為少.
將來.
必有一位從你那裡出來.

$^{1281}$在以色列中.
為我作掌權.
他的根源.
那個根源.
他的origin.
是從耿古.
從太初就有的.
這兩節經文.
很重要.
因為.
這個大希律.
就是希律王.
在耶穌降生的時候.
那個希律王.
當然他的兒子就是.
施洗約翰那個.
這個希律王就見到.
有幾個博士.
來找他.
就問他.
希律王就.
打發這幾個博士.
去伯利恆.
為什麼呢.
因為希律王.
召見那些.
那些.
當時的.
當時的那些.
教律法書的人就說.
基督.
生在哪裡呢.
那些文字就說.
生在猶大的伯利恆.
馬太福音第二章就引這一段.
是很清楚.
這樣來講的.
所以.
這個希律王就打發那些.
那幾個博士.

$^{1321}$在東方來的幾個博士.
所以他們是很對的.
這幾個博士.
從東方來到耶路撒冷.
就去找希律.
就問希律王.
生下來做猶太人的王.
在哪裡啊.
我們在東方.
遙遠的東方.
看見他的星.
特別來拜他.
希律王聽見之後.
這幾個外訪的客人.
突然間進宮.
見我.
問我.
那個生下來的.
猶太人的王.
是在哪裡出生.
希律王聽見就.
心裡不安.
會不會帶來那個混亂呢.
其實希律王.
自己是分封的.
當時是羅馬皇帝.
該殺的.
希律王是一個分封的王.
當時的巡撫.
是彼拉多.
這裡就是說.
耶路撒冷合成的人.
都不安.
第四節馬太封第二章.
他就召齊了祭司長.
和民間的.
文士,即教聖經.
教律法書的,就問他們.
基督當生在何處.
他們就回答.

$^{1361}$在猶太的伯利恆.
因為有先知記著說.
就是剛才我們說的.
尼加書第五章.
第二節.
猶太地的伯利恆.
你在猶太諸城中.
並不是最少的.
因為將來有一位君王.
要從你那裡出來.
牧養我以色列民.
有少少不同的出入.
因為大家用的版本.
是不同的.
新約用的.
是一個叫做七十時.
希律文的翻譯.
舊約,七十時.
是一個叫做Septuagint.
是出名的.
在耶穌時代是一個出名的抄本.
但是舊約的.
是叫做馬索拉版本.
即是我們今天.
現在我們手上這個翻譯的舊約.
我們和合本的舊約.
是取材著希伯來文的.
馬索拉抄本.
來做翻譯.
所以大家用的版本.
是不同的.
所以是會有出入的字眼.
但是沒有影響到.
整體的文字.
的表達和結構.
你可以想像到.
尼希米.
都不例外.
尼希米都是.
等待彌賽亞降臨的一個人.

$^{1401}$那個時代.
每一個敬虔的猶太人.
他們都是會被教育.
要等候彌賽亞的降臨.
而且.
如果深水清.
熟讀律法書的人.
就會知道神已經定了.
基督.
彌賽亞.
其實是同義詞.
基督即是救主.
就是要.
降生在伯利恆.
所以耶穌.
的媽媽瑪利亞.
好像這麼巧地.
來到去.
懷了新人.
那麼巧.
這個該撒.
當時的大帝.
即是羅馬的第一位皇帝.
因為大.
大獲全勝.
剷平了當時.
的敵對的安東尼.
那麼就.
要來到去.
頒令天下.
重新做一次.
全國的人口普查.
所以.
每一個民族的人.
回到自己的出生地.
所以.
這個約瑟.
就帶著太太.
瑪利亞.
回到伯利恆.

$^{1441}$其實他們本來.
是拿撒勒的人.
所以就回到伯利恆.
在伯利恆那裡.
來做一個人口普查.
所以這麼巧地.
周圍的人都回到伯利恆.
就變成當時.
好像旅遊旺季一樣.
你想有.
酒店,有客館,旅館.
都沒有了.
所以也臨盆在即.
就在馬槽.
可能是乾的.
當然不會是.
都不會那麼.
慘到.
那麼髒.
可能就是因為太緊急.
沒有地方.
就可能馬上安置一個.
馬停放的地方.
就讓瑪利亞來生.
主耶穌.
這個就好像各樣巧合.
但原來聖經舊約.
已經預言了彌賽亞.
是要在伯利恆那裡.
降生.
所以.
尼希米.
他雖然緊張.
現在耶路撒冷的光景.
但是心底的裡面.
他們都是.
要等待.
他們的救主.
他們的彌賽亞.
什麼時候會降生.

$^{1481}$這個都是.
一個虔誠的猶太人.
必定會有的.
那個內心的一個態度.
就算是今天的猶太人.
當然比較虔誠.
他們.
仍然都是這樣.
等候彌賽亞降臨.
不能夠接受.
耶穌是他們的彌賽亞.
所以你會看到.
如果你去聖地.
你去到哭場.
你會看到那些很虔誠.
傳統的極之傳統的那些猶太人.
穿上黑色西裝.
戴了一頂黑色帽子.
頭上有一個盒子.
綁住的.
一條帶.
然後手上纏著那些帶.
是一些經文.
在哭場裡.
這樣來叩頭.
一直在等候.
彌賽亞降臨.
懇求彌賽亞降臨.
但是奈何.
諷刺地.
我們基督教我們基督徒.
很清楚的就是.
彌賽亞主耶穌第一次已經降臨了.
我們等待的.
是第二次的降臨.
但是.
現在的猶太人.
等候第一次降臨的彌賽亞.
因為他們拒絕.
相信和接受.

$^{1521}$耶穌是他們的彌賽亞.
所以這是一個.
最大的問題.
所以他們是不接受.
新約的聖經.
所以我們.
為什麼印新約的希伯來文.
新約是希臘文.
用希伯來文印新約.
其實是很難的.
猶太人是不接受.
所以如果平心靜氣.
用希伯來文去讀新約.
那些基督徒.
即是猶太人的基督徒.
其實是很大的感動.
因為如果他真的用希伯來文.
自己的語言去讀新約.
他就真的會被.
新約的對主耶穌.
特別是福音書.
那個人生.
那個.
他的位格.
會來到感動.
和改變.
所以我們真的要為猶太人.
禱告.
因為最難就是他們拿起本.
新約的希伯來文的聖經.
去讀.
他們自己平心靜氣.
去看一下新約的.
作者怎樣來.
去講耶穌.
真是我們的彌賽亞.
真是我們的救主.
我們不需要再.
等候第二個.
頭神去.

$^{1561}$差遣彌賽亞降臨.
他們就是不接受.
彌賽亞是耶穌.
所以我們回來.
這個尼希米記的第一章.
的時候.
對不起,有點乾.
我們回來這裡的時間.
我們就看到.
尼希米很簡單.
他的表達,因為是一個自轉式.
他個人的一個.
聽到一個消息.
很緊張弟兄姊妹.
由猶太下來.
就是說這個下來.
如果以色列用了五個月.
當時的運輸.
還有.
走這些遠路.
這些遠路經過沙漠.
經過一些障礙.
更加不會有那麼多.
驛站,那麼多旅館.
等等.
都會有盜賊.
那個危險.
雖然是一個.
王道.
上去.
但是都有一定的.
旅途的風險.
所以.
那些消息絕對不是即時的.
短訊.
現在根本沒有任何的.
界限.
你在任何一個世界的地方.
只要那個時差你不會影響.
根本上就是即刻通訊.

$^{1601}$所以有時候我深夜的時間.
都是剛剛香港的天光.
有時候會處理一些.
譬如程序表.
都有一些東西要回答.
都是即時的.
沒有任何的阻隔.
但是在黎希米的世界.
主前的四百多年.
距離今天.
二千四百多年.
可想而知.
那個通訊.
是滯後得很緊要.
所以.
那個期待.
又變相是更加的.
熱切.
因為你要幾個月才收到一個消息.
那個消息就是靠那個人.
去表達.
又沒有相片給你拍到.
又不會即刻給你看到.
現在等等那個現場.
所以.
更加有短片.
根本沒有這些.
那個期待.
那個焦急.
是可想而知.
可以來理解.
這個就是黎希米.
他一路.
他自己強調.
我是作王的酒精.
而且第二章就.
立即告訴我們.
有一個宮廷的宴會.
的一個場景.
來到去.

$^{1641}$神.
應允了.
黎希米的.
那個祈禱.
沒有理由那個王.
要三個月至四個月.
才偶然間留意到.
黎希米有不開心.
即是由那個王.
觀察到你的一個內神.
那個的.
面色是有.
很多的掛慮.
有很多的擔憂.
應該不是.
突然之間的.
可能一路積聚著.
黎希米都是正常的工作.
不會是茶飯不思.
不會.
即是說經常因病告假.
經常來到去.
他的內心.
好像有一個很大的.
一個痛苦.
因為在王的面前.
不應該.
也不可能有.
這樣的一個情緒.
這麼明顯的表達.
所以.
你可想而知.
他在自己的內室.
那個的祈禱.
因為這些的.
逐弟兄.
到底這個關係.
是這些人有使命下來.
還是不同的.
經濟的貿易.

$^{1681}$往來.
來到去有這些的交往.
還是這些人.
下來可能是做一些.
宗教上面的一些的教導.
帶一些的可能是聚會.
這個都沒有清楚.
告訴我們.
似乎黎希米他想表達的.
不是由那方面.
而引發他.
而是他自己的內心.
一直對耶路撒冷.
所知道的.
所了解的.
而激發他自己的擔心.
而這個擔心.
使得他.
聽到這個.
消息.
就是這幾個人.
下來所說的.
本來想問他.
照計就不應該.
說得這麼清楚這麼白.
就是由大省.
遭大能受凌辱.
耶路撒冷的城牆.
被拆毀.
城門被人用火焚燒.
你可想而知可能是一次的事件.
導致到外邦人.
有一些的武力.
來對猶太人.
有一個的.
借一些機會.
借一些的藉口.
來攻擊.
所以這些不是可能經歷了幾十年.
可能是突然間的一次性的事件.

$^{1721}$也都剛好在這個情況下.
雖然我們不知道.
確實是什麼事件.
但這個交往.
到底是官方的.
還是宗教的.
這些在耶路撒冷下來的人.
本身是帶著一些工作.
可能就是下來書山城.
來教導.
在書山城的猶太人.
來守節旗.
或者是有一些聖經教導.
這個都不知道.
作者都不是太感興趣.
說那個.
背後的操作.
重點是將那個消息.
放在黎希敏的心裡.
黎希敏聽到.
這個消息.
而她的內心.
有這種刺痛.
消息帶來的.
一個人的內心.
來去.
看為是自己的責任.
是自己都要感同身受的.
是自己都需要.
來去.
停下來.
當然是用她自己私人時間.
下班的時間.
一個酒精.
有沒有休息時間.
因為皇室的東西.
皇的酒席.
其實什麼時候要吃.
就吃了.
雖然有時那些宮廷的宴會.

$^{1761}$都會安排什麼時間.
有什麼的.
外賓來到訪.
你都會編排好的.
有那些內務的人.
會幫你編排.
但是有沒有自己的時間呢.
這裡告訴我們.
黎希敏在她自己.
回到自己的家裡的時間.
當然她需要休息.
她需要有自己的.
和其他的.
從大陸來的弟兄.
來交往.
即是說.
有她自己的放假時間.
她就將這個消息.
放在心裡.
來不斷的醞釀.
在她的內心.
來看到.
這個和她自己.
這個民族.
她是神的選民.
這個身份是息息相關.
是有一個很密切的.
關係.
導致她.
要回到自己的家裡.
來禱告.
然後日常.
上班.
上朝.
她又繼續在皇的面前.
來侍當.
第二章是記載.
一個宮廷宴會.
而且是開心的.
因為為什麼呢.

$^{1801}$到第六節的時候.
作者特別告訴我們.
那個時候皇后都坐在皇的旁邊.
特別強調.
那個宮廷宴會.
皇后都在.
即是.
要不就很開心的.
宴請那些外賓.
要不就皇.
這個不是重要的.
即是很重要的會議.
即是.
那些太太團.
只是拍照.
其餘時間她就要撤退.
因為是談一些世界.
國家大事.
但現在似乎皇后.
是坐在皇的旁邊.
即是說.
可能是一個喜慶.
是一些外賓來.
進見的.
那個宴會.
可能比較隨意一點.
隨意就更加不會那麼嚴肅.
反而就凸顯了.
皇就.
發現.
尼熙米.
比平日已經是三個月.
祈禱了至四個月.
由十一月.
就是基斯留月.
去到第二章的.
這個彌散月.
就是大概.
一至二月.
總括前前後後.

$^{1841}$就差不多有四個月.
那你就可以想像到.
皇.
這次很認真發現了.
尼熙米是心事重重.
那這個心事重重.
嚴格就不是.
每天都苦口苦面.
而是他一路的禱告.
一路的在神面前.
去等候.
一路等待那個蒙恩.
原來是在.
皇的面前.
尼熙米能夠蒙恩.
這樣的禱告.
好像是很迂迴.
他沒有直接開口.
跟皇說.
似乎是.
似乎是八章打後.
他回到.
原來他做著做著.
他會回到.
書山城.
他會回到.
亞達斯西王那裡.
一段時間.
有些功績他要履行.
後來他才再回來.
再回來又發現了一些問題.
又再來處理.
就是尼熙米記.
給我們的.
一個這樣的教導.
我們在這一章裡面.
當我們還沒讀第二章之前.
尼熙米記就有很多個人的.
自轉式的記載.
他的言語.

$^{1881}$既然是他的.
自轉式的言語.
每一章都講他自己的禱告的內容.
有時是.
閉上眼企圖.
有時他就很認真.
來禁食.
哭泣.
祈禱.
等等.
有時他就.
圍到哪些人去祈禱.
有時就求神.
去禁止哪些人.
對他有一些的攻擊等等.
就是求神去.
禁止哪些人的計謀.
所以尼熙米記.
是充斥著這個.
差不多每一章都有的.
那個自轉式的禱告.
所以這本書.
很多人就將它.
尼熙米的性格.
又來做分析.
不是錯的.
不過整本書是想講.
一個敬虔的猶太人.
面對彌賽亞降臨.
他有什麼可以擺上他自己呢.
他.
後來他回去.
是將他自己的俸祿.
完全投放給猶太人.
因為當時.
他在猶太人的.
教育環境.
是不可以給猶太人.
因為當時他們要建城牆.
他們又受到外邦人.

$^{1921}$的阻攔.
所以那些百姓.
特別是做工的.
那些工人.
沒錢.
尼熙米就將.
他自己的俸祿.
給了這些工人.
尤其是利未人.
所以他就是.
真是一個.
盡心盡力.
去為神.
耶路撒冷.
神的國度.
記住神的國度.
神的耶路撒冷.
在舊約只有一個目的.
就是準備.
聖民.
即是神的選民.
去預備彌賽亞降臨.
去預備救主降臨.
所以尼熙米.
他都知道.
當然他們不知道.
主前的.
差不多第六年.
主耶穌降生.
不是零年.
就在大希律的時候.
所以.
四百多年後的事.
尼熙米他只知道.
彌賽亞會降臨.
在我們的地方.
是否在巴里行呢?我不知道他知不知道.
有可能知道.
但後期的民事.
一定知道.

$^{1961}$所以很戲劇地.
他們知道.
但似乎那些的民事.
那些的祭司長.
就沒有興趣去巴里行.
去找那個聖嬰.
反而那幾個博士.
真是去到就見到.
有個小孩子包著布.
在馬槽.
向瑪利亞.
和這個聖嬰.
耶穌.
敬拜和獻上.
禮物黃金乳香.
活躍.
作為禮物.
接著他們就被指示.
回到他們的原居地.
沒有回去告訴希律王.
所以.
希律王就很生氣.
就要殺光.
兩歲以內的那些的嬰兒.
即是說.
那些人都要用了一些時間.
才去找到耶穌.
所以對希米爾來說.
他不只是看到.
那些城牆破落.
自己的.
阿公的故居.
要收集.
要去做一些土木工程.
好像移居在英國.
其他地方的.
荷蘭.
很多的原居民.
新界的,譬如元朗.
我去到認識的.

$^{2001}$有一個在餐館.
荷蘭餐館做老闆的.
原來就是牛潭尾.
的原居民.
很小的,十幾歲就已經.
移民了去荷蘭.
後來就是盧森堡.
在盧森堡,現在開了兩間餐廳.
都很好生意.
他的餐廳就是開放.
做教會.
這樣的一個弟兄.
有機會去見到他.
又請我們吃飯.
你可以想像到一個原居民.
如果想回到.
自己的牛潭尾的祠堂.
已經是很破落.
其實在那裡長大.
你也有一個心.
很想回去收集.
很想回去出錢重建.
這個祠堂.
因為是和你們這個村.
你們這個姓.
的一個人.
一個支派.
一個村.
的原居民.
你也有這個心.
何況尼熙米就是.
所以又是回去問一個問題.
在書山城的.
宗教.
教育.
他們的主一學.
是怎樣的.
發展呢?.
為甚麼可以孕育到一個尼熙米.
這麼敬虔.

$^{2041}$這麼對神的事情.
來擺在心裡面呢?.
我們有沒有.
一些年青人.
已經去了外國留學.
仍然對.
神的家.
對教會.
他有很重的心呢?.
在外國.
這麼多文化的衝擊.
這麼多他要自己面對的.
不同的國籍.
不同的朋友.
很多的分途.
慢慢就會.
收集了.
思想文化.
能夠保留中國人文化.
已經難了.
不僅是中國人文化.
香港人的文化已經難了.
再想想.
要有基督徒的.
價值觀的文化就更加難了.
但是.
這個也是最重要的關鍵.
如果.
有一批.
一代的人.
他能夠保留.
對神的敬虔.
對神的敬畏.
對神的追求.
這個相信就是.
就算在任何的地方.
神都可以.
感動他.
猜險他.
好像給他一個GPS.

$^{2081}$他可以返回上去.
有神給他的.
calling.
他可以來做.
神要他去復興的工作.
李希敏就是一個.
很典型的.
在聖經裡面.
主持四百多年的一個領袖.
一個.
他的成長.
是在王的面前.
得到王的信任.
得到王的.
近身的.
一個重要的官員.
但是他可以.
是拿大家.
用王的.
給的.
即是國度.
波斯國度給的資源.
可以在河西.
那裡來到去.
任意.
去拿王的詔書.
來到去申請.
一些建造的材料.
河西的.
總督是要.
給他的.
那你就可想而知.
是可以借助.
外邦的國家的資源.
來到去.
用在神的殿裡面.
的那個建造.
還有整個耶路撒冷.
是要有一個.
規模.

$^{2121}$有自己的城牆.
有自己的宗教的制度.
有自己那個.
至少是恆常.
不是成為.
例行公事.
他的心是很想.
一到星期日就去教會.
你想想在外國.
更加難.
教會已經少了.
你去到星期日你很緊張要去教會.
你就可想而知.
那個.
虔誠如果能夠延續.
就是最重要.
的那個資產.
最重要讓.
下一代無論他在哪個地方.
他自己都能夠.
對神有這個.
的渴慕.
這個就是很困難.
所以尼希米和.
以斯拉他們要.
令到那個可以.
敬拜的那個氛圍.
能夠維持.
不是一種僵化.
而是一種的.
復興而是讓.
更多的猶太人能夠.
排除那個.
的困難.
沒錯是困難的.
困難從來都在這裡.
所以尼希米上到去才.
發現那個困難比他.
想像的比他.
流淚來到.

$^{2161}$哭泣的時間想像.
的更加困難.
因為他真的要計算.
有什麼什麼的困難.
第二章.
後半段開始到.
第三章就真的列舉.
每一個困難出來了.
但是沒錯.
領袖的出現.
就是要解決難題.
一個領袖.
應該就是他要.
很清楚他的出現.
就是解決那些難題.
一個領袖.
如果是有承擔的.
是應該面向那個難題.
就不要走去.
避開那個難題.
至少你想辦法.
去嘗試.
去解決.
嘗試去提出.
那個警告.
現在有個洞缺題.
荷蘭那個提吧.
如果你現在再不塞住.
這個洞這麼小.
越衝就越大.
所以領袖.
就是要在教會裡面提出.
一些.
如果你繼續下去可能不合真理.
我們一起要停在這裡.
大家要小心一點去反思.
譬如今年11月.
佳音樓續約.
大家要同心一起面對.
不可以大家在這裡.

$^{2201}$你這樣想.
我又這樣想.
這樣就不行.
大家一定要一起同心.
領袖就有這個責任.
去帶動大家的討論.
在神裡面的討論和思想和合一.
這裡起碼都不例外.
他現在聽到之後.
就發現了.
有一個很大的難處.
人.
特別男性.
我不知道女性.
一聽.
有些人就馬上火滾.
滾得很厲害.
當場就很厲害.
那個熱心.
但是第二天就沒有了.
馬上淡了.
原來他講完就當做完.
但是.
有些人就很沉實.
他會.
認真看這件事.
他會好像.
不溫不火那樣.
但其實他放在心裡面.
慢慢慢慢去沉澱.
沉澱完之後.
應該怎樣做.
他會來計算.
所以真正的領袖就不是當下.
爆一爆.
就好像全世界你真是厲害.
這麼熱心.
他發現原來.
困難還更加大.
他當下的衝動.

$^{2241}$馬上就蓋過了.
這些都是領袖.
但不能帶動改變.
帶動改變.
第一件事就是.
首先就是面向困難.
首先就是要.
沉澱自己的心情.
我自己.
都已經二十多年了.
我以前都是很衝動.
我很深刻印象初訊的時候.
開始學習.
傳道.
真是有時不知不覺.
在港台會很重的罵人.
不知不覺.
有時真的會有些事情.
你會馬上有一種.
很強烈的衝動.
很想馬上解決問題.
但發現原來不是這麼簡單.
馬上就解決到.
不是你想像的這麼容易.
經過了很多年.
很多年.
人就真的沉靜了.
有些事情一來.
事物多了.
你會給空間.
自己也要.
慢慢消化.
做一個很嚴格的分析.
來去研判.
然後才去梳理.
譬如哪些.
可以短期就解決.
哪些是中長期要解決.
和要動員哪些人.
一起去幫忙.

$^{2281}$這些過程其實是幾十年的.
醞釀的一個改變.
所以.
很多的長輩說.
一個傳道人.
他真的要去到五十多歲後.
才會被神開始用.
因為頭幾十年.
是他自己裝備自己的性格.
你才體會.
人的心靈的軟弱.
對別人你會多一種的體恤.
多了一種的諒解.
多了一種的同行.
和你會感同身受.
這個.
我想傳道牧者.
就是做一個模樣.
所以黎軒蜜在這裡的時候.
她的確.
立即坐下哭泣.
悲哀了幾天.
一個男性能夠悲哀幾天.
真是一個.
很寶貴的一個操練.
女性就馬上哭.
就很容易了.
有些哭完就沒事.
也都很好的.
但是男性.
很難去馬上哭.
我不知道猶太人.
是不是特別情感豐富.
中國人就真的很難.
難到一個地步.
就是我們某程度.
是將那個情感的那個區.
關閉了.
要剔除那個區.
通常有時很嚴重的.

$^{2321}$事情才會激發那個.
哭泣的意欲.
但很多時候都不容易的了.
但是.
不代表.
不認真和不賞心.
情感的.
表達有時只是一個.
的宣洩.
宣洩過後那個承擔.
宣洩過後仍然.
看這件事是嚴肅的.
要跟進那個才是領袖的.
承擔.
哭泣是暫時代表著.
那個情感豐富.
如果男性.
能夠哭泣完幾天.
都再沉澱那件事.
認真去梳理.
這些情操.
這些性格.
和這樣的一個.
領袖的風範就很重要了.
就很重要了.
所以黎希敏就真的具備.
這個她自己感同身受.
之後.
她真的以大局.
男性如果胸襟夠大.
眼光夠闊.
就是以大局.
這個大局就是神國度的.
彌賽亞她這樣來看.
自己的時候.
我是一個領袖.
我是一個酒精.
我是一個在皇身邊的.
飲飲食食的一個官.
我有沒有份於耶路撒冷的修造.

$^{2361}$這個重建.
有沒有份於.
那個外敵.
能夠做一點點事.
可以幫忙改變.
所以每一個領袖都應該有.
這樣的心智.
是為大局.
不是為他自己個人的聲譽.
或者是野心.
如果野心.
他是一個酒精.
他不需要他已經有了.
他已經在皇的身邊.
所以真的.
到處都可以爭.
其實更加不應該爭.
到處都不應該爭.
我們做人應該是.
要很謙厚.
按神的時間.
尤其基督徒應該要有這樣的信念.
但是教會就絕對.
不可以爭.
誰爭神一定會出手.
所以要小心.
不要做神的功.
想將功勞歸給自己.
千萬不要.
所以.
什麼時候有人不贊成你.
神告訴你.
功勞不是你的.
現在人家不贊成你.
那你就更加應該.
講雅門.
因為功勞是神的.
就算做得不好.
都是神自己會負責.
一個領袖就應該要有.

$^{2401}$這個胸襟和學習.
如果不是.
其實就很不應該.
悲哀完幾天之後.
他.
加起來是禁食祈禱.
所以悲哀加祈禱.
才是一個.
出路.
如果只是悲哀.
而沒有真正的禱告.
在神面前.
來向神講.
尋求.
神的出路.
所以祈禱.
很快猶太人.
就會馬上想起.
申命記的聖經.
摩西的律法書.
馬上想起.
這些的應許.
馬上想起.
什麼時候我們回到神的寶座.
向神.
表達我們的需要.
我們的困難.
和我們的不足.
我們的罪惡.
什麼時候我們回到.
神就是有恩典.
有出路.
甚至他應該會忘記.
其實不是由.
他開始.
由所羅巴百開始.
由當時的古烈王.
的時候.
一百二十幾年前.
已經有了那個被擄歸回.

$^{2441}$所以.
現在只是在那個.
根基上面.
再繼續建造上去.
不是由零開始.
所以尼希米.
更加有這個信心.
神一定會.
為他這個禱告.
這個發現.
來給他有出路.
給他有.
去學習.
所以盼望弟兄姊妹我們都成熟了.
年紀的成熟.
大家都應該超過五十歲.
聽我講聖經.
如果有年輕人十幾歲.
喜歡聽.
很感恩.
我這次去荷蘭.
發現有些人有聽查經班.
很感恩.
盼望又祝福到不同的人.
因為放上網有好處.
其實只是講聖經.
有系統大家一起去讀聖經.
可能有些趣味性.
大家一起有些背景.
講一講.
幫助大家去讀.
還要真的.
我們真的要讓.
大家去聽.
尤其是.
不只是金句式.
好像尼希米.
他一祈禱.
摩西五經這麼重要的應許.
就立刻閃入他的內心.

$^{2481}$他在祈禱裡.
不只是就坐.
惡人就坐外敵.
就坐完又怎樣.
他居然.
很大膽.
他會不會是神用的.
一個工具一個器皿.
他祈禱就是.
神啊如果我們.
回到你的面前.
我們去回.
那個無助的.
本相去到你的面前.
你是真的有應許.
會使我們能夠回到.
那個應許之地.
和那個.
最高的你給我們的.
那個Blessing.
尼希米真的這樣的心.
等待三個月四個月.
由這個.
基斯留月一路等待到.
彌撒月一路這樣.
祈禱一路這樣祈禱一路等候.
就是說.
我這裡就是第一章.
我這樣祈禱.
這幾個月我都會這樣祈禱.
就是這樣的意思第一章.
但是原來他是.
日間的工作是做酒精.
日間.
他是要令到黃.
不要吃錯東西.
開心和如果是國宴.
就令到那個國家.
不會被人看不起.
這些是酒精的工作.

$^{2521}$但是尼希米就能夠.
平衡他自己的.
工作和.
他的教會生活.
是很平衡.
現在香港人.
忙到一個地步.
一個星期他如果將.
這個星期日的上午.
分別為聖由他自己.
起身吃早餐.
洗完澡和子女.
一起回教會一家人.
很認真的來到去.
放下他在天崇.
兒校跟住他們.
回來坐下.
很認真的敬拜神.
很認真的參與我們.
都覺得很寶貴很難能可貴.
因為一個星期香港人真的很忙的.
但是有沒有這個意識.
你想想.
七天裡面的半天.
對神要有一個很強烈的意識.
但是尼希米.
他不單止是這樣.
不只是平常的崇拜.
而是現在.
有一件這麼特殊的事.
他很想將自己放在.
參與裡面.
祈禱了三個多月.
等候一直等候.
第一章他繼續祈禱.
第二章.
他離散月.
恆常照常.
這次是皇后.
在席宴會.

$^{2561}$可能是一個國宴.
皇后都坐在這裡.
這樣來到.
一個宴會.
但是皇后居然.
發現尼希米.
他的心情.
可能不是面容.
而是他的內心.
真是難能可貴.
你沒有病.
為什麼你會面帶愁容呢.
一定不是.
第二樣東西.
不是吃錯東西.
不是有其他東西.
國家沒事.
最近沒事.
沒有煩惱的事.
沒有馬英九訪華.
說什麼.
沒有這些.
完全是.
國家很正常.
裡面有困難.
有什麼的愁悶.
有什麼的愁煩的事.
現在皇后.
很關心這個臣子.
他的內心.
所以.
祈禱了三個多月.
這就是尼希米.
所以如果一個人.
對臣的心.
持續地.
來關心.
聽到主一學.
聽到儀校.
聽到天崇.

$^{2601}$聽到非印語.
聽到現在的奸流.
聽到孫德.
聽到洪恩.
聽到老牧師.
他繼續祈禱.
都不用問阿貴.
神感動這個人.
可能就要做全道.
通常都是這樣.
因為你不斷地關心.
聽完之後.
但是祈了三個多月.
一定是神.
有感動在他內心.
給他一些東西.
他有什麼可以參與.
我自己個人都是這樣.
初信的時候很熱心.
一路開始.
那時二十多歲.
真的.
回教會祈禱會.
那時在屯門醫院做.
星期日.
我的上司很憎基督徒.
特意編我.
不能回教會.
很激氣.
後來就.
聽了一堂道.
傳到人教我們.
我們不能回聚會.
不是代表我們不能回聚會.
是不能回聚會.
是通了.
禮拜那時在屯門醫院.
沒有什麼人.
剛剛開新院.
星期日當值.

$^{2641}$突然在哪裡教我.
我想到.
我不能閉上眼.
當值的時候.
我用支筆在本堡.
一路寫一路祈禱.
為聚會為弟兄姊妹祈禱.
這樣回了半天.
很開心.
一放工馬上趕輕鐵.
衝回元朗.
那時剛剛散會.
人們開始走.
抓住人一起吃飯.
大家融洽.
補假.
星期三下午放假.
下午.
可以回去.
在教會.
那時我開始幫忙錄音.
整理錄音帶.
差不多吃飯.
跟傳到人一起吃飯.
回到祈禱會.
如是者.
越來越對教會的事.
很關心.
後果就是.
神感動了傳到.
一定是這樣的後果.
所以祈禱三個月.
要小心.
神感動你.
要為神的工作去事奉.
這是榮耀的為神去做.
所以人生.
在在有些人已經退休.
在線上.
已經退休.

$^{2681}$事奉神是最崇高.
永恆.
為神去做.
當然不同崗位.
可以默默去做一些.
現在退休更多時間.
默默觀察.
哪個童工已經燃燒.
為祈禱.
甚至乎買瓶水給他.
鼓勵他.
這些已經是很大的事奉.
不需要有一個很出名.
當然出名不是罪.
不是因為有不同崗位.
有些人要顯耀.
有個名字.
有些人不需要有個名字.
都要這樣做.
因為是神安排的.
一定要有尼希米這個心智.
真的祈禱幾個月.
為那件事.
一路一路去沉澱.
我們要祈禱了.
多謝主帶領我們今晚.
開始的尼希米記.
第一章.
我們將一些背景都講了.
希望神自己繼續去施恩.
感動.
亦都盼望我們大家.
這樣來到去.
很簡單的去研讀的經文.
可以成為很多人的祝福.
不單止是感宣家.
而是有機會透過網上.
聽到這個系列的.
都能夠蒙主的祝福.
求你藉著聖靈.

$^{2721}$去感動不同的弟兄姊妹.
在世界各地.
都能夠為神的家.
為福音.
為主的降臨.
我們去齊心努力.
亦都為猶太人的信主.
我們繼續去禱告.
盼望他們.
真的有更多的猶太人.
可以接納.
新約的希伯來民.
來去研讀靜心的.
打開他們的心靈.
這樣去研讀.
我們這樣禱告奉耶穌基督的名.
阿們.
這個往往是一個非常的.
私心一點的不誠實的.
安靜一點.
然後你看她們.
那麼紳士的.
就是這個.
到底對啊.
那我覺得.
在聽完這個.
我覺得.
我可以坐在這.
5-10分鐘.
\newpage



\section{馬太福音 21:1-11}
\label{sec:x06IlO4FHos}
\textbf{2023-03-26 午堂《看你的王》馬太福音21\_1-11 講員:李思敬牧師}
\newline
\newline
連結: \href{https://youtube.com/watch?v=x06IlO4FHos}{\texttt{https://youtube.com/watch?v=x06IlO4FHos}} ~~~~ 語音日期: 2023-04-03
\newline
\newline
\hyperref[sec:yk5Oh8RBHLY]{\small{< < < PREV SERMON < < <}}
~
\hyperref[sec:index_chronic]{\small{[返順時目]}}
~
\hyperref[sec:index_scriptual]{\small{[返順卷目]}}
~
\hyperref[sec:ythaF42qv_4]{\small{> > > NEXT SERMON > > >}}
\newline
\newline
馬太福音 21:1-11
\newline
\begin{longtable}{cl}
\hline
\hline
章節 & 經文 (和合本修訂版)\\
\hline
21:1 & \begin{tabularx}{0.7\textwidth}{X} 耶穌和門徒快到耶路撒冷，進了橄欖山的伯法其時，打發兩個門徒， \end{tabularx} \\ \\ \relax
21:2 & \begin{tabularx}{0.7\textwidth}{X} 對他們說：「你們往對面村子裡去，會立刻看見一匹驢拴在那裡，還有驢駒同在一處，解開牠們，牽到我這裡來。 \end{tabularx} \\ \\ \relax
21:3 & \begin{tabularx}{0.7\textwidth}{X} 若有人對你們說甚麼，你們就說：『主要用牠們。』那人會立刻讓你們牽來。」 \end{tabularx} \\ \\ \relax
21:4 & \begin{tabularx}{0.7\textwidth}{X} 這事發生是要應驗先知所說的話： \end{tabularx} \\ \\ \relax
21:5 & \begin{tabularx}{0.7\textwidth}{X} 「要對錫安的兒女說：看哪，你的王來到你這裡，謙和地騎著驢，騎著小驢—驢的駒子。」 \end{tabularx} \\ \\ \relax
21:6 & \begin{tabularx}{0.7\textwidth}{X} 門徒就照耶穌所吩咐的去做， \end{tabularx} \\ \\ \relax
21:7 & \begin{tabularx}{0.7\textwidth}{X} 牽了驢和驢駒來，把他們的衣服搭在上面，耶穌就騎上。 \end{tabularx} \\ \\ \relax
21:8 & \begin{tabularx}{0.7\textwidth}{X} 許許多多的人把自己的衣服鋪在路上，還有人砍下樹枝來鋪在路上。 \end{tabularx} \\ \\ \relax
21:9 & \begin{tabularx}{0.7\textwidth}{X} 前呼後擁的人群喊著說：「和散那歸於大衛之子！奉主名來的是應當稱頌的！至高無上的，和散那！」 \end{tabularx} \\ \\ \relax
21:10 & \begin{tabularx}{0.7\textwidth}{X} 耶穌進了耶路撒冷，全城都驚動了，說：「這是誰？」 \end{tabularx} \\ \\ \relax
21:11 & \begin{tabularx}{0.7\textwidth}{X} 眾人說：「這是從加利利的拿撒勒來的先知耶穌。」 \end{tabularx} \\ \\
[1ex]
\hline
\hline
\end{longtable}
$^{1}$弟兄姊妹主內平安.
很開心可以不用戴口罩講道.
也見到你們的燦爛笑容.
今天我們所讀的經文.
相信你也熟悉的.
耶穌騎驢進耶路撒冷城.
跟我們有什麼關係呀.
我們又不是禮儀教會.
所以復活節快到了.
早一個星期就是中樹主日.
有沒有聽過呀.
但我們不會那天要排隊.
每個人手上拿一支中樹枝.
一路搖著一路出去走一個圈.
有沒有呀 沒有的.
如果真是聖公會.
要做這件事的話.
還要很嚴謹的.
那些中樹枝搖完之後.
不是扔進垃圾桶的.
是要儲起明年的大齋期第一天.
聖灰日.
將去年的中樹枝燒成灰燼.
然後在崇拜.
星期三晚上有一個陶灰禮.
我不知道大家有沒有參加過呢.
宣導會沒有的.
我在神學院參加過一次.
一位信義會.
對不起 崇真會的老師帶領我們.
很莊嚴的一個崇拜.
我們排隊來到台前.
那位牧師在我的額頭上.
就用爐灰畫一個記號.
最震撼的是一路畫的時候.
他一路跟我說.
司徑 你本是塵土.
仍要歸回塵土.
我說 啊 以前我聽到這句話.
都是說那個而已.

$^{41}$我們很慶幸我們不是塵土.
我們未歸回塵土.
我們走出去.
不是啊.
這個事實第一次在我的心中.
留下很震撼的一個提醒.
所以牧師接著說.
所以你要背起十字架.
專一跟從主.
我們不是這樣的禮儀教會.
我們叫做free church.
所以我們沒有這個傳統.
於是乎.
我們讀這段聖經的時候.
很多次我都聽過.
我們的焦點.
就會集中在.
主耶穌所騎的那一匹路的身上.
講到的題目.
我們會引用其中一句金句.
要用他.
效法這一匹謙卑的路區.
讓主來騎在我們身上.
用我們.
他得到榮耀.
很感動的訊息.
完了之後.
我們心裡都說.
主啊 我就是那條路.
但是這一段聖經.
那個焦點.
那個主題.
在那條路上.
不是.
耶穌騎路進入耶路撒冷.
其實坦白說.
耶穌走進耶路撒冷都可以.
他又沒有腳痛 膝蓋不舒服.
所以那條路.
其實可有可無.

$^{81}$為何要安排一條路.
還有路區.
小路.
在旁邊.
我們很少會去.
尋根究底.
打爛沙盤.
問道督.
但其實值得這樣做.
因為.
經文.
福音書記載.
耶穌入耶路撒冷城.
都特別引用舊約聖經.
的典故.
先知的說話.
來到我們中文有一句成語.
畫龍.
甚麼呀.
點睛.
那個最精彩的地方.
我們一定要明白掌握的.
原來就是福音書引用.
舊約先知撒迦利亞書.
撒迦利亞書在哪裡.
很容易找到.
雖然它是十二先知.
我們叫做小先知書.
放在一邊.
其實它在馬拉基書之前.
新約聖經.
馬太福音.
剛才我們讀了.
早兩卷.
第幾章呀.
是第九章.
第幾節呀.
撒迦利亞書九章九節.
福音書引用這句說話的時候.
就說錫安的民.

$^{121}$你們要大大喜樂.
為甚麼呀.
耶路撒冷的民.
要歡呼.
看呀.
你的王.
來到你這裡.
它是公義的.
施行拯救.
謙謙和和.
騎著驢.
就是驢驅子.
這句說話.
引起了一些西方的學者.
拗了頭.
耶穌騎著兩隻驢.
一齊進入耶路撒冷.
一腳踏兩船.
不是.
你這樣就不明白.
其實因為我們離開了.
古代的社會.
一匹的牲口.
作為坐騎.
一個可能性.
預備多一匹.
以備不時之需.
這麼隆重的日子.
這麼重要的場合.
有沒有contingency plan.
你開會.
你都會懂得問這個問題.
一下下傾盆大雨.
急症室就漏水了.
不是.
我們有backup.
馬上載他們去天水圍.
去哪裡.
這個很重要.
所以驢驅在旁邊.

$^{161}$是甚麼意思呢.
一定不是耶穌擺大的.
他的雙腿.
要騎著兩隻驢.
沒有這回事.
另外一個可能性.
告訴我們.
耶穌騎著的那匹驢.
是媽媽.
是不是.
有沒有見過崇拜師班獻師.
牽著小朋友出來.
他哭了.
媽媽如果自己走出去.
他不肯.
乖呀.
牽著出來站好.
有些教會就不給.
我見過給.
帶著小朋友出來.
他又很精靈.
站在那裡.
媽媽唱詩.
眼尾看著他.
這隻驢驅跟著母驢.
一齊出來.
強調這隻驢的性別.
不是後備方案.
都有這個可能.
為甚麼有這個可能呢.
因為我們要明白.
這裡所用的又是一個典故.
福音書引用先知書.
先知書引用哪一段聖經.
我們叫引上引.
不容易.
我們不是很熟悉舊約聖經.
撒迦利亞書是引用.
列王記上第一章.
你有沒有聖經.

$^{201}$問你有沒有iphone.
順便查一下whatsapp.
不要專心聽到.
列王記上記載的是大衛晚年.
其實應該說他已經離樓之際.
在病床上起不到身.
為甚麼證明在一個這麼危急的關頭.
因為他的四王子亞多利亞.
已經籠絡了軍隊.
約翰元帥的全力擁戴.
就催國宣佈繼位登基.
大局已定了.
未.
列王記上第一章就記載.
拔士巴和那丹先知就來見大衛王.
提醒他老人家.
其實你答應過是所羅門繼承你的王位.
大衛自己親口說的.
不過不是記載在撒母耳紀或列王記.
記載在歷代至上最後的一章.
不要看了.
已經遺失了.
大衛很清楚在百姓面前宣佈.
他說耶和華在十二支派當中.
揀選了猶大支派.
在猶大支派當中揀選了我的父家耶西.
在耶西我父的家中揀選了我.
大衛有多少個兄弟你有沒有印象.
他有七個哥哥.
他是第八.
石到lone 了.
他不是潮州人.
他是伯利恆人.
所以父親忘記了這個兒子.
面色紅潤.
會彈琴.
看羊.
和他的哥哥高大衛猛.
英明神武.
不是我們都要有些中文文化.

$^{241}$氣宇軒昂.
打仗走出來.
所以撒母耳去到耶西的家.
耶和華上帝揀選的一定是這個了.
這麼高.
上帝說我不是看外貌.
我看內心.
七個哥哥走過.
撒母耳迫著要問.
你還有沒有藏起一個不見得人的.
我告訴你.
叫齊你的兒子在那裡.
耶西很尷尬.
真的漏了.
這樣都可以.
我們很難想像.
大衛就是一個被父親遺忘了的.
面色紅潤.
即是baby face.
會彈琴.
即是男孩子不是學空手道.
是學芭蕾舞.
都不知道是直還是彎的.
大衛說上帝揀選我.
然後在我的兒子當中.
揀選所羅門.
你知道大衛有多少個兒子.
不知道.
知道就去神學院教書.
不考慮了.
十九個.
你乘下份數.
上帝的揀選.
大衛清楚.
所以大衛不是老懵懂.
雖然他身體起不了床.
但他頭腦清醒英明.
他就吩咐他的僕人.
拉我平時所騎的那一匹雷出來.
希伯來文還講到明.

$^{281}$是母雷.
不是公雷.
然後讓所羅門騎在我這一匹雷的背上.
拉他出去.
百姓見到就明白了.
你認不認識這個暗號.
認不認識.
能夠騎在大衛王的座騎上的.
那個才是真命天子.
你再有軍隊.
再有什麼巴閉摧國.
不算數.
果然.
歷史人心裡又問問題了.
舊約難提.
哪認得的.
你認不認識這輛車.
在停車場看到後面的角落.
你已經知道是自己的車.
你認不認識你太太.
在超級市場只見到前面額頭半分.
你認得到的.
是嗎.
不要講這麼深奧.
車牌一字你認不認得到.
是誰的.
走過去踢兩腳.
你作死.
認得到的.
我們不是生活在那個社會上.
但是這個記號.
百姓明白.
所以撒迦利亞書不是大衛王朝.
是亡國被擄之後.
回歸耶路撒冷重建聖殿.
已經事過程遷了好幾百年.
但是仍然用歷史的事實.
來跟耶路撒冷的百姓說.
你的王來了.
誰知道哪有王.

$^{321}$我們都亡了國.
他騎著驢進來.
為什麼不是驢.
第一典故.
不寫一定要找那隻動物出來.
什麼叫典故.
今天我們出國留學.
臨行那一晚.
媽媽好好心機.
幫我們撿衣服進行行李箱.
爸爸在旁邊看到.
就微微一笑.
哎呀 真是慈母守中線.
你明白她說什麼嗎.
不是啊 媽媽都沒有拿著針線.
不是 不用要這樣去.
那個不是寓言.
孟浩然在唐朝.
寫了一首詩 寓言你媽媽.
但是她寫了這五個字.
大家就明白了.
垂燕寸草心.
一起說 步得三春暉.
所以慈母守中線的典故.
重點不是那條線.
不是那個動作.
明白嗎.
所以騎著那匹驢還是驢.
其實不是最重要.
但是一個君王騎著這樣的狀況.
走進來的歷史上.
就是所羅門那一次.
百姓就認出.
他是以和華所揀選的.
大衛王的心意.
你有軍隊站在一邊.
這個歷史的根據.
我們必須要心知肚明.
這個是猶太人的暗號.
亡國被擄.

$^{361}$不再有自己的國家.
你造反.
你找個王出來.
不行的.
耶穌寫在十字架上.
上面的罪狀.
寫什麼你記得嗎.
猶太人的王.
所以撒迦利亞書.
用了所羅門這件事情.
百姓完全明白.
耶穌安排騎驢,母驢.
進入耶路撒冷.
百姓明白.
這個聖經的雙重典故.
撒迦利亞書.
列王記上.
百姓心知肚明.
前面那句應該叫做.
雞食放光蟲.
有沒有聽過.
不知道為什麼.
心裡面只要吃一條蟲.
放光的.
肚子就光了.
心知肚明.
這樣都行.
這些是中文的歇後語.
耶穌這樣的安排.
留意福音書記載.
合成都震動起來.
唯一沒有震動的是哪一班人.
羅馬軍隊.
懂不懂得這樣做聖經.
羅馬人不知道發生什麼事.
有探子來報.
報有一個猶太人.
騎著一隻驢走進來.
都不是騎馬.
都沒有軍隊.

$^{401}$不用理他.
那些猶太人是這樣的.
所以外邦人沒有讀聖經.
完全不明白發生什麼事.
神的百姓熟悉舊約聖經的典故.
先知書.
大衛王.
現在真的發生在我們眼前了.
合成都震動.
衝出去.
禍殺.
奉主名來的.
是應當稱重的.
我們不要做外邦人.
讀完聖經都只看著驢.
主要用他謙謙卑卑.
不是.
主角是主耶穌.
耶穌入耶路撒冷城的時候.
他是透過這個行動.
告訴不是全世界.
羅馬人蒙在鼓裡.
但是他自己的百姓猶太人.
完全知道發生什麼事.
行不行.
行嗎.
還有一句.
還有一個可能.
牧師.
騎著驢謙謙和和.
所以耶穌很謙卑.
人家問他是不是彌賽夏.
不是不是不是.
不是你說的.
所以耶穌騎驢的意思.
不是說他自己.
是很謙卑.
不承認.
推搪.
不如你做主席.

$^{441}$多做一堂.
不行不行不行.
這樣才謙和.
好啊.
這麼驕傲.
你又讀錯聖經了.
請問.
聖經當中.
最謙和的人是誰.
你知不知道.
耶穌.
一定沒有錯.
去面試考中一.
一定是這樣答.
錯了.
你記不記得民數記.
十二章.
是什麼.
我從來都不讀.
摩西的哥哥姐姐.
毀謗弟弟.
第二個人攻擊摩西無所謂.
摩西的姐姐.
是她做B的時候.
在尼羅河救過弟弟.
你記得吧.
摩西的哥哥亞倫.
是上帝跟摩西說.
你絕口不說.
未出世我就知道了.
所以你未出世之前三年.
你哥哥我預備了.
他做你的先知.
在法老面前.
他講說話.
你動在那裡就行了.
摩西身邊最親的兄長姐姐.
現在來攻擊他.
你讀回去.
今天下午.

$^{481}$給你有機會讀一下.
列王記上第一章.
又再讀五經民數記十二章.
很短的經文.
摩西沒有為自己辯護.
為什麼事.
那些我們不說了.
時間關係.
摩西一句話都不出.
在這個時候.
作者應該不是摩西自己.
最保守的聖經學者都認為.
最謙和的人不會寫在聖經裡.
說我自己是最謙和的.
是吧.
那就不謙和.
所以這句話是上帝為他的僕人做見證.
因為上帝出聲.
責備哥哥和姐姐.
摩西從頭到尾一句話都沒有出.
聖經稱讚他是世上最謙和的人.
這個字很重要.
不要翻譯錯.
所以八福溫柔的人有福了.
因為他們必承受地土.
說話不要太大聲.
不然買不到地.
不是.
謙和的人有福了.
同一個字.
什麼叫做謙和.
不出聲.
不是.
不出聲只是表面.
心裡的那個不是謙和.
是吧.
嘴角都已經歪了.
冒汗了.
心裡在說髒話.
就坐著人.

$^{521}$不是.
這不是叫做謙和.
什麼叫做謙和.
摩西在這件事情上.
示範了什麼叫做謙和.
他雖然是領袖.
受到攻擊.
百詞莫辯.
不是很容易解釋.
但他選擇信上帝.
不靠自己.
聽得明白嗎.
很難的.
要大件事開會.
做領袖一定懂得倚靠上帝.
裝模作樣都要倚靠上帝.
吃飯.
無端端被人砰一句.
你裡面謙不謙和.
馬上就出了.
是不是這樣.
所以摩西在這個情況下.
他完全安靜.
代表他從頭到尾.
他是一個怎樣的人.
他是一個讓上帝用他.
任憑主用.
不是那隻驢.
是摩西.
是主耶穌.
聽得明.
看得清楚.
所以這裡用的.
其實不止兩重.
三重.
回到摩西.
回到民數記.
謙謙和和的君王.
不是靠自己.
非勢力.

$^{561}$非才能.
是靠耶和華.
方能成事.
這個是進耶路撒冷城.
百姓明白.
見到騎著驢.
知道這個是上帝.
所預備的彌賽亞.
聖經再加多一個典故.
除了撒迦利亞書.
騎著驢驢驢驢子.
還有謙謙和和的摩西.
他不是自己話事.
自己掌權.
自己控制.
計劃一切.
他是懂得將自己.
擺在神的手中.
上帝要怎樣用摩西.
上帝要怎樣用主耶穌.
一脈相承.
在這個行動的過程當中.
不是說耶穌就頭凹凸.
眼濕濕.
驢又咪咪下.
總之就很淒涼.
這樣就謙和了.
你完全解錯聖經了.
耶穌安排這個行動.
就是要向耶路撒冷的百姓.
宣佈漢雅.
你的王來到了.
但這個王.
不是百姓期待的彌賽亞.
你知道的.
百姓期待的是一個.
勇武攬炒的粉瑞黨領袖.
和羅馬軍隊死過.
不推翻羅馬不義政權.
起碼將駐軍逐出耶路撒冷聖城.

$^{601}$好了.
上帝差遣.
真命天子來到耶路撒冷.
起義革命.
百姓群情洶湧.
但很快就你知道.
在彼拉多的衙門前面.
又是這班人喊什麼.
釋放巴拉巴.
將耶穌釘十字架.
彼拉多都不信自己的耳朵.
真的嗎.
釋放巴拉巴.
我們又是不明白.
不過這個就不關外邦人不外邦人的事.
我們又不讀聖經.
巴拉巴是一個強盜.
瑞斯案.
這麼兇殘.
怎麼會釋放他.
強盜.
我們的印象就是強盜.
你讀清楚聖經.
四本福音書.
馬太,馬可,路加,約翰.
沒錯.
說巴拉巴是一個強盜.
約翰福音十八章四十節.
但是馬太福音呢.
馬太福音說.
巴拉巴是一個很出名的囚犯.
二十七章十六節.
被羅馬關了的囚犯.
但很出名.
哪些囚犯出名的.
政治犯.
懂嗎.
普通囚犯.
抓了一個在街邊偷東西的.
不出名的.

$^{641}$沒人認識的.
但政治犯.
很多報道的.
以前沒有報紙的.
沒有WhatsApp的.
沒有TikTok的.
都有口傳的消息.
怎麼證明被拉多.
被拉了巴拉巴.
是因為政治的原因.
馬太福音十五章第七節.
告訴我們.
殺人作亂.
聽到這四個字.
殺人放火就強盜.
殺人作亂就暴動.
路加最後加多了一個關鍵.
二十三章十九節.
巴拉巴做這些事.
是在城裡面做的.
你可以將整幅圖畫湊在一起嗎.
所以當百姓發覺.
耶穌原來這麼軟弱的.
不是他們心目中來救他們脫離羅馬的軍隊的遊律的.
不能夠解放耶路撒冷的.
不能夠為猶太國盡一分力的.
他們就放棄耶穌.
換誰.
換一個他們心目中佩服的政治領袖巴拉巴.
耶穌來到.
不是要做巴拉巴做的事.
但他是王.
被拉多再三問他.
你是否猶太人的王.
耶穌怎麼回答.
我們幾乎閉上眼就.
耶穌說不是.
耶穌不回答他.
我給你讀一讀.
記載得最詳細的約翰福音.

$^{681}$十八章.
三十三節.
被拉多進了衙門.
叫耶穌來升堂.
如果你看以前中國的電視劇集.
叫完就審問耶穌.
你是猶太人的王嗎.
你明不明白你的罪狀.
法庭第一件事問.
告你什麼.
耶穌回答說.
這話是你自己說的.
還是別人論我對你說的呢.
耶穌沒有沉默.
耶穌沒有說不是.
耶穌說是誰說的.
是誰告我的.
被拉多說.
我都不是你們猶太人.
你自己國家的人和祭司將你交給我.
你究竟做了什麼.
你快點說.
耶穌回答說.
我的國不屬這世界.
我的國若屬這世界.
我的臣僕必要爭戰.
使我不至於被交給猶太人.
此時我的國不屬這世界.
被拉多對他說.
這樣你是王嗎.
耶穌回答說.
你說我是王.
我為此而生.
也為此來到世間.
特為給真理作見證.
凡屬真理的人就聽我的話.
耶穌否認嗎.
耶穌只是否認.
他的國不屬這世界.
所以他說我的臣僕.

$^{721}$不會和你爭戰.
那是另一個層次.
那是低層次.
這個世界上的政治鬥爭.
低層次.
我的國不屬這世界.
我為此而生.
我是為這件事來到世間.
為什麼.
被拉多剛剛問他.
你是王嗎.
耶穌不單沒有否認.
耶穌沒有沉默.
耶穌承認.
不過他說.
我不是你們心目中的王.
是這麼多.
我們就要聽清楚.
我們所跟隨的耶穌.
原來不是我們心目中的那種王.
你跟不跟他.
那不跟了.
福音沒有了公義.
那就不是福音了.
你說的.
我曾經說過.
耶穌來到這個世界上.
他的確是王.
但是他說得很清楚.
我不是你們這樣想的.
你們這樣想錯了.
耶穌那條路.
他怎麼做王.
我們怎麼跟他.
我們肯不肯跟他.
時間差不多了.
我和你讀西門彼得.
預估整件事情的時候.
彼得這位使徒.
怎樣提醒我們.

$^{761}$彼得前書第二章二十一節.
彼得前書二章二十一節.
我讀給大家聽.
你們蒙召原是為此.
蒙召做什麼的.
不是去跟巴拉巴.
是跟耶穌的.
清楚了.
因基督也為你們受過苦.
給你們留下榜樣.
叫你們跟隨他的腳蹤行.
耶穌基督留下的腳印.
我記得了.
在沙灘上一行腳印.
他抱著我.
那個是海報.
那個不是聖經.
你找到有一次聖經這樣說.
我什麼都給你.
你找不到.
不過我們相信.
我們很喜歡.
全躲避.
就算真是一行腳印.
是耶穌先行.
我們一步一腳印.
踩著耶穌走過的腳印走下去.
這樣合聖經好多.
他留下腳蹤.
給我們跟從.
不要走到旁邊.
不要自己走了.
對.
但耶穌的腳印.
耶穌的腳蹤.
耶穌留下的榜樣.
是一個受苦的榜樣.
我知道了.
十字架.
為誰受苦.

$^{801}$為我.
你是誰.
他愛我.
我又愛他.
不是.
聖經不是這樣說.
是我們幻作罪人的時候.
耶穌為我們死.
神的愛是這樣顯明.
做罪人的時候.
我們是怎樣對耶穌的.
希伯來書很清楚.
告訴我們.
十二章第三節.
忍受罪人頂撞的.
我們頂撞他.
你都不是彌賽亞.
你都不剷起不義政權.
你都不幫我們.
你都不公義.
你的福音都是林北北.
只是教人上天堂.
我們頂撞主耶穌.
你跟不跟他.
你跟不跟.
他的路我找個看清楚才好.
他是受苦的.
為誰受苦.
為那些頂撞他的罪人受苦.
所有宗教都不是這樣說.
所有宗教都是上帝愛你.
保佑你.
幫助你.
拯救你.
安慰你.
扶持你.
聖經怎樣說.
上帝愛你.
呼召你.
叫你來受苦.

$^{841}$你跟不跟.
信耶穌這樣嗎.
不信.
很多人就是這樣.
計算代價清楚之後.
悠悠愁愁就走了.
不單止這樣.
我們求神拜佛.
就希望上帝幫我們打小癮.
50元一個.
有沒有見過.
我見過.
不過不拍得上.
很便宜.
拍片他.
打殘他.
下釘不下已經算給他面子.
上帝做什麼.
上帝幫我去懲罰惡人.
審判罪人.
是我心目中的惡人罪人.
我多次講這篇道的時候.
有人問.
那神的仇敵呢.
上帝說他死了.
差遣獨生兒子來.
都不是拍扁他.
都不是打殘他.
都不是剷起他.
都是虐失.
虐失講得好.
他和十個哥哥講什麼.
創世記第五十章.
從前你們的意思是要害我.
但神的意思是好的.
為你們保存生命.
成就今天的光景.
我們跟隨主講得出這句話嗎.
我又是講完這篇道之後.
有人問.

$^{881}$神的意思怎樣好.
我覺得不好.
為害我的人而死.
為害我的人受苦.
不值得.
上帝幫我拍扁他.
你拍扁他我信你.
那個不是上帝.
那個是鬼仔.
你養他就是去拍扁人.
今晚拍這個.
明晚拍那個.
不是.
耶穌騎驢進入耶路撒冷.
群眾知道他是誰.
他拒絕跟從.
寧願跟巴拉巴.
因為大家政見相同.
耶穌說我為此而生.
我的國不屬這個世界.
我為那些頂撞上帝的罪人而死.
你們來效法我的榜樣.
跟從我的腳蹤.
你願意嗎.
我願意嗎.
我沒有問你是甚麼顏色.
我不管你氣誰 喜歡誰.
不過主耶穌的國.
是這樣的國度.
主耶穌留下的腳蹤.
清清楚楚一點都沒有含糊.
我們今天見到他騎驢進入.
看守你的王.
我們轉頭走了.
上帝你跟我.
主耶穌你幫我.
搞定他.
不是呀 錯了.
這樣不是跟從耶穌.
求主幫助我們.

$^{921}$請同心低頭祈禱.
如果不是你的愛.
我們今天還是死在罪惡過犯當中.
沒有盼望的死人.
多謝你.
你沒有計較我們怎樣對你.
就算今天也好.
就算我們說我們信了.
耶穌的基督徒也好.
你都沒有計較.
我們原來真的有這麼多的反叛.
我們要你跟我們.
多過我們心悅誠服的去跟你.
求主赦免.
求主憐憫.
求主幫助.
用你的說話成為我們腳前的燈.
路上的光.
讓我們認清我們的主是誰.
信得過 好像摩西.
我們看不到結果.
我們不知道上帝怎樣成就你的救恩.
但我們信得過.
既然你是獨一的真神.
除他以外就別無拯救.
你有永生之道.
我們就不需要去跟第二個.
多謝你.
多謝你聽我們不配的祈禱.
奉主耶穌基督成名求.
阿們.
\newpage



\section{約翰一書 5:20-21}
\label{sec:ythaF42qv_4}
\textbf{2023-04-02 午堂《信仰求真》約翰一書5\_20-21 講員:曾裔貴牧師}
\newline
\newline
連結: \href{https://youtube.com/watch?v=ythaF42qv_4}{\texttt{https://youtube.com/watch?v=ythaF42qv\_4}} ~~~~ 語音日期: 2023-04-04
\newline
\newline
\hyperref[sec:x06IlO4FHos]{\small{< < < PREV SERMON < < <}}
~
\hyperref[sec:index_chronic]{\small{[返順時目]}}
~
\hyperref[sec:index_scriptual]{\small{[返順卷目]}}
~
\hyperref[sec:_BOzzljBDkA]{\small{> > > NEXT SERMON > > >}}
\newline
\newline
約翰一書 5:20-21
\newline
\begin{longtable}{cl}
\hline
\hline
章節 & 經文 (和合本修訂版)\\
\hline
5:20 & \begin{tabularx}{0.7\textwidth}{X} 我們知道，神的兒子已經來到，並且將悟性賜給我們，使我們認識那位真實者，我們也在那位真實者裡面，就是在他兒子耶穌基督裡面。這是真神，也是永生。 \end{tabularx} \\ \\ \relax
5:21 & \begin{tabularx}{0.7\textwidth}{X} 孩子們哪，你們要遠避偶像。 \end{tabularx} \\ \\
[1ex]
\hline
\hline
\end{longtable}
$^{1}$一齊祈禱.
多謝主.
我們再一次聚集在你的面前.
我們很盼望.
我們的信仰能夠活在我們的主.
你的榮耀.
你的同在.
和你的愛裡面.
常常保守自己.
而且我們的信仰的神學是穩固的.
是真的照著聖經的啟示.
我們這樣信的.
我們一齊在你面前.
我們很盼望.
我們能夠有一種不知不覺的屬靈的影響力.
在改變這個世界.
因為當我們待會去詮釋經文的時候.
我們真的會發現.
這個世界是臥在惡者的手下.
那個惡者是會侵害每一個人的.
幫助我們能夠明白.
這樣的祈禱.
奉耶穌基督的名.
阿們.
在這個月我們會集中有四次.
來講關於主耶穌.
一些很重要很重要的真理.
今早我們就會先去探討.
第20節這裡.
「這是真神,也是永生」.
是在說誰呢.
是在說來這個世界做過人的耶穌.
這是很斬釘截鐵的告訴我們.
這是我們信仰的核心.
我們就在這一段經文.
當然說笑了.
孩子們啊.
你們要自首.
遠避偶像.
為什麼突然間一直在說真神.

$^{41}$又在說我們要逃避拜偶像呢.
到底作者想帶出什麼樣的分別呢.
如果我們再看回上一點的經文.
19,20節.
應該是18,19節.
突然間加插了魔鬼在地上的工作.
和對人的侵害.
我們就能夠了解.
作者是在描繪.
一個世人.
如果他沒有神給的智慧.
聖靈.
是不會看到這個世界以外的另一個世界.
就是屬靈的世界.
就是有永生的弟兄姊妹才會掌握.
明白和來到去.
在這個世界裡面尋求的.
就是我們所說的信仰.
如果我們從13節開始看這一段.
相信就更加清楚.
可以將這兩個世界或兩個國度.
來描繪出來.
13節就是整本書.
約翰一書最重要寫作的目的.
作者告訴我們.
這一位年老的約翰.
在這個時候他應該是最後一位.
還沒有離開世界的.
跟隨耶穌的使徒.
11個已經分別離開世界了.
最早第一個殉道的.
不是殉道.
是自殺.
是猶大.
所以等一下.
那裡16節說的.
有些人會致於死.
是不是在說猶大呢.
當時寫這封信的時候.
他說有些人是致於死的罪.

$^{81}$你不要為他祈禱.
為他祈禱都沒有用.
他致於死.
在神面前好像永不超生一樣.
相當恐怖的.
所以這一本書現在告訴我們.
13節先描繪了一個很重要的目的.
就是信主耶穌的人.
掌握到什麼呢.
是真的可以掌握的.
我將這些寫給你們信奉神兒子之名的人.
要叫你們知道自己是有永生的.
有永生.
所以這裡說的是.
離開這個物質世界之後的生命.
聖經給我們清楚的就是.
那個叫永恆生命.
你要知道現在這一刻.
你再生的年日.
因為你信奉神兒子的名.
你就是有永生的了.
你要掌握著這件事.
那怎樣能夠可以來到牢牢的掌握呢.
如果大家稍微熟悉一點教會的歷史.
清教徒和宗教改革之後.
那個時候叫新教徒.
他們唯一認為自己是蒙神保守的.
其中他們可以拿捏的就是.
必須在有生之年.
殷勤的做工.
殷勤的致富.
然後殷勤的奉獻.
將自己常常擺在神的那個作為.
那個福音的事工裡面.
他們就肯定自己是蒙神的賜福和保守.
這個是新教徒的信念.
這個是清教徒一個很重要的.
他們的邏輯思維.
和這一段聖經似乎都相當吻合.
我們可以來到從其中找到一些.

$^{121}$我們怎樣可以拿捏.
我們是有永生的弟兄姊妹呢.
所以今天講的.
作者想講的是對著信了主耶穌.
你已經有了那個悟性.
知道自己有永生.
認識了神的.
當然在座如果聚會.
你還沒有信耶穌呢.
也是為你而設了.
因為這個地方叫做禮拜堂.
是有永生的.
世界的物質的世界.
你已經足夠掌握了.
你的生老病死.
你的起居飲食.
每個人都重視.
每個人都緊張.
但是來到禮拜堂.
不是和你探討世界的試驗程.
在這個世界.
很真實的世界以外.
是有一個屬靈的世界.
是聖經告訴我們.
甚至比世界現在的這個物質的世界.
更加的真實.
更加我們應該要掌握拿捏.
當我們在這個目的之下.
作者就很想告訴我們.
永生就是一份很重要的神.
給我們無條件的禮物.
在這個無條件的禮物裡面.
誰人能夠牢牢的珍而重之.
和保存.
不會被人偷.
這個就是重點.
所以十四十五節.
甚至延伸到十六節.
講述有會被偷走信仰永生的人.
我們要怎樣幫助.

$^{161}$包括我們自己.
怎樣保守自己在神的裡面.
十四節有一個很大的應許.
就是除非你不祈禱.
你祈禱.
神一定會聽.
但只有一個條件.
沒有第二個條件.
這個條件就是你要在神的旨意底下去祈禱.
看看作者怎樣講這個應許.
十四節.
我們有神生命永生的兒女.
我們就會祈禱.
所以基督徒如果不祈禱.
不懂祈禱.
要學習祈禱.
因為這是一個向我們每一個的應許.
所以這個世界不會有些特殊的人.
唯有他祈禱才有效.
我是不相信這一套.
所以靈因.
我們某程度是不認同的原因.
就是他們有些特殊的行列的人.
稱自己是使徒.
稱自己是先知.
他們的祈禱就特別靈驗.
可以為你的病祈禱.
一定會好返.
不好就不是他問題.
是你沒有信心.
輸打贏要講完.
我們傳統福音派是不相信這件事的.
這裡告訴我們.
每一個神兒女都可以禱告.
原則只有一個.
有沒有照神旨意祈禱.
第十四節.
我們若照他的旨意求什麼.
即是無論求什麼.
那個寬闊度是無論求什麼.

$^{201}$你想想.
我們每天有很多東西需要.
如果我知道明天有什麼會發生.
已經不得了.
不得了.
我們有很多人想知道明天會發生什麼.
再簡單一點.
如果我真的有一種透視.
開什麼號碼.
今天下午如果跑馬.
我不知道有沒有.
因為已經很久沒看了.
如果我知道哪場哪場大義門是否真的贏.
你想想如果能夠有這個.
但現在告訴我們.
無論求什麼.
你留意聖經怎樣的應許.
他就是神.
他就聽我們.
就算我們向他所傳坦然無懼的心.
還要是坦然無懼.
即是你不是奸.
你的祈禱他會聽.
甚至是一種坦然無懼的聽你的祈禱.
而他會來成就.
甚至還要再說多一次.
既然知道他聽我們一切所求的.
即是無論求什麼.
就知道我們所求於他的.
無不得著.
弟兄姊妹.
有神而屢生命.
信奉神而子之名的人.
即是信耶穌的人.
只要開聲祈禱.
無論求什麼.
必定得著.
記住其中一個條件.
這個條件就是.
十四節已經說了.

$^{241}$我們若照他的旨意求什麼.
是他的旨意.
是有一個框.
在這個框裡面.
無論求什麼.
這個框就是.
你可以這樣說.
就是基督教最寶貴.
在有生之年.
我們一生一起去追求的.
認識這個框裡面的內容.
你就是那個可以求什麼.
他就會聽我們.
而且既然知道他聽我們.
一切在這個框裡面求的人.
所求的.
我們就知道無不得著.
作者給我們一個這麼美麗的圖畫.
原來這個永恆世界.
這個屬靈的世界.
是可以這麼精彩.
是可以這麼實在.
是可以我們能夠拿捏到.
因為你是坦然無懼的心.
去到神的寶座的面前.
不是一種抽象.
有些人不懂祈禱.
有些人很怕祈禱.
他好像覺得自己的祈禱.
穿過不了天花板.
當然是開玩笑.
但是其實這裡告訴我們.
是可以更加越來越坦然無懼.
感覺到一種神人的那種相交.
是可以在神的寶座去求什麼.
而且你的求是在他的旨意的框裡面.
我們很想第16節應該問.
什麼是你的框.
作者約翰.
你可不可以給一些指示我們.

$^{281}$什麼是神旨意裡面的框裡面.
是不是地上的物質.
我們求到某一個階段.
不可以再求.
很有趣的地方.
16,17和18節.
開始探求同樣有神永生生命的兒女.
有哪些人是已經離開了神的恩典.
那個護衛保守.
而我們就需要去祈禱.
好像作者沒有告訴我們.
12345就是那個框裡面的.
求姻緣是不行的.
求學業OK.
這些是我們去拜黃大仙的時候.
會幫我們分類.
現在是簽什麼簽.
你想求什麼.
我求自身而已.
或者求一些事業.
作者沒有讓我們這樣分類.
他再一次帶我們去看.
在教會裡面.
我們現在是追求這個信仰之旅.
我們對永恆的世界.
是有一個認知.
我們就要來到.
在我們的生活裡面.
開始要探求一件事.
就是其他的神兒女.
他的靈明的光景是如何呢.
16至17節和18節開始說一件事.
就是在這個實在的世界裡面.
作者在19節20節就會再說.
特別是19節.
我先點出.
我們知道我們是屬神.
但是全世界都臥在.
這個惡者的手下.
我們知道我們有信仰的人.

$^{321}$是屬神的.
不過有另外一個存在的敵對勢力.
敵對神.
就是全世界.
當然不是說花草樹木.
是說住在這個世界的人類.
臥在這個惡者的手下.
我們回頭看第16,17,18節.
我們就發現這件事的重要和寶貴.
第16節.
人如果看見弟兄犯了不至於死的罪.
就當為他祈求.
如果照上文所說的.
無不得著的祈求.
照神旨意去祈禱.
他就會聽我們.
既然他聽我們.
就知道我們所求於他的.
是無不得著.
這樣的思路回應下去.
作者的意思就相當的清楚.
作者現在描繪的一個框.
就是神兒女的屬靈生命.
就是神旨意底下.
我們需要互相扶持.
互相代求.
人如果看見弟兄犯了不至於死的罪.
這裡又要分兩樣罪.
至於死.
不至於死.
我們可能很好奇.
什麼是至於死.
他說不至於死的罪.
就當為他祈求.
神必將生命賜給他.
相信大家都不會誤解.
是肉身.
是要找醫生來搶救的.
那個肉身生命的生命.
相信是說永恆在神眼中的生命.

$^{361}$神會將生命賜給他.
然後突然又說了另外一樣東西.
有些是至於死的罪.
我不說當為者罪祈求.
就是說至於死.
怎樣去分辨呢.
你講一些例子.
偷東西超過九百元美金就不是至於死了.
因為可能根本都不會抓你.
我講在美國.
不是香港.
怎樣叫做至於死.
怎樣是不至於死呢.
為何作者可以這樣界分.
有分兩種嗎.
不至於死就可以祈禱.
至於死.
我不說當為者至於死的罪祈禱.
作者講得好像很玄妙.
誰去拿捏.
誰去判斷至於死不至於死呢.
不至於死就可以祈禱.
接著17字再補充.
凡是不義的事都是罪.
也有不至於死的罪.
即是說繼續將那件事清清楚楚地點出.
原來祈禱和陪伴那個已經和神的永生.
慢慢有距離的人.
我們要去幫助他.
因為18節開始就講那個惡者的出現.
即是說在這個世界有一個敵對的勢力.
是與神的永生為敵.
18節我們知道凡是從神生的人.
是必定不會犯罪的.
從神而生的必定會保守自己.
這個保守是在魔鬼底下的保守.
接著說那惡者也就無法害他.
你想想信耶穌的人.
原來相信這個世界不只是物質.
這個世界原來有敵對神的魔鬼.

$^{401}$作者一而再講.
那個惡者是會害他.
這個害相信不是一些很離奇怪誕.
經常見鬼的害.
相信是和上文和下文有關連的.
就是那個智慧失去了功能.
失去了一種悟性.
失去了對神的愛的同在.
而慢慢地.
這世界是我的家.
人一世物一世.
我做什麼.
只要我自己能夠得到真正的滿足.
其他東西都次要了.
講永生太遙遠了.
剛才我也舉過一個例子.
我自己剛剛出來傳道.
在我的教會在我的團契.
已經二十多年前了.
所以現在變成怎樣我不知道.
我真的不知道.
有一個弟兄.
教會就來跟他說.
因為他的信仰非常之模糊.
開始被一些哲學影響得很厲害.
不斷問是不是真的有神呢.
其實他自小就回教會了.
是不是真的有信仰呢.
是我的團契.
那時候我們很少人.
那時候我二十多歲.
就和他對話.
其實是熟的.
但是就聊天.
叫他不要這樣.
那時候也不是很懂得處理.
有時候說話都很重.
覺得有沒有理由不信神呢.
後來他就清楚和教會說.
你們不要再找我了.

$^{441}$我不信了.
不要再干涉我的生活了.
自此他真的信守諾言.
不再回教會了.
現在大家都五十多歲.
不知道現在變成怎樣了.
這樣一位弟兄.
其實起初信得很好的.
慢慢就真的離開了.
相信在我們一直成長.
如果我們回到青年團契.
來到今天.
石果僅存的.
一定是少數.
但是我們自己的天崇.
我們的兒校.
能夠長大去到青年團.
如果留學回來之後.
還繼續回教會.
相信一定是少數.
有時候也不用說得這麼白了.
我們自己坐在這裡的弟兄姊妹的兒女.
現在他們在床上.
或者是昨晚喝完通宵酒.
看完七欖之後就去玩.
接著就沒什麼興趣回五堂了.
相信大有人在.
絕對不是譴責.
而是聖經告訴我們.
有兩個很重要的陣營.
同時間存在著.
教會的功能.
如果我們不能遵守.
下一代是會失落的.
我們會不會有很多.
我們已經固有的觀念.
強行灌輸給下一代.
導致他們對信仰起了反感.
導致他們每次有些嘗試.
都被人譴責.

$^{481}$導致慢慢自己對信仰.
不再感興趣.
甚至變成敵對信仰.
相信也會有.
我們大家有時候也會感到一種痛.
我的兒女不回教會.
不再像我們小時候.
在天崇的時候.
那種活潑那種熱誠.
所以這裡給我們一個方向.
作者很想告訴我們.
原來我們要意識到.
是可以改變.
要由我們開始.
所以十八節.
我們知道凡是從神生的.
是必不犯罪.
從神生的是保守自己.
惡者也就無法害他.
然後再詳細說這個意思.
十九節.
我們知道我們是屬神的.
全世界都臥在冷惡者手下.
然後再一次說我們知道.
屬神的人知道什麼呢.
就是一個的吾性.
就是二十節.
我們知道神的兒子已經來到.
每個人都知道.
四月七號就是壽苦節.
十二月二十四就是平安夜.
就是基督降生.
每個人都知道.
連猶太人都相信知道.
有誰在這個地球上.
會不知道耶穌說自己是神呢.
何況甚至無神的社會.
阿拉伯的世界一樣相信.
知道耶穌是他的身份.
但是這裡說的不是知識的知道這麼簡單.

$^{521}$第二十節說的是.
且將智慧賜給我們.
特別強調不是知識性的知道.
不是你小時候讀聖經.
在基督教學校有機會讀Bible.
這裡說的是一個吾性.
這裡說的吾性.
這裡說的智慧.
相信是說聖靈進入一個信主的人的心靈裡.
就會讓他認識到.
這個認識就是說他從此不再只是屬於這個世界.
他有永恆世界.
就是另外一個世界的那個.
儘管用正見.
他有另外一個世界的那個意識.
這個就是聖經說的智慧.
剛才喝咖啡的時候.
不經意聽到兩位弟兄在分享.
你有聖靈.
很重要.
就開啟你內心對神的認識.
不經意偷聽到.
就是這裡說.
聖經告訴我們就是二十節這裡.
原來一個有永生的人.
為什麼他會這麼貼近神呢.
原來就是神賜了.
透過耶穌來這個世界.
他的救贖原來賜了一樣東西給信耶穌的世人.
大家都是世人.
但是有永生的就多了一樣東西.
這個叫做智慧.
屬靈的智慧.
也是新的和合本說的悟性.
這個悟是領悟了.
不是你個人的一種參透.
不是你個人的一種讀書很厲害.
哇 還厲害的 查GPT.
什麼都研讀了.
是神賜給我們的這個悟性.

$^{561}$帶來的就是使我們認識這一位的真實者.
我們是在這一位真實者的裡面.
就是在他兒子耶穌基督的裡面.
這個悟性就是.
作者用了兩句修昔句.
告訴我們這個悟性.
不是我們突然之間可以變成一種超人.
這個悟性就是對耶穌的一個真確的認識.
因為他是一位真實者.
我們是住在這一位真實者裡面.
就是他的兒子耶穌基督裡面.
這一位的耶穌就是真神.
就是永生.
這是20節講的悟性.
原來是神賜的 弟兄姊妹.
所以這個是真正的無價的禮物.
所以回過頭來.
我們要保存這個禮物.
也要回過頭來.
18,19節告訴我們.
有一個惡者會偷我們這個禮物.
會偷走的.
各種不同的思潮.
各種不同的慾念.
會偷走我們有永生的這個悟性.
以致我們失去了最終對神的認識.
和自己可以享用的.
在神面前的身份.
和與神相交帶來的那個福樂.
會偷走我們的.
就是說如果我們的信仰不是真實.
我們就是在拜偶像.
所以很特別的就是21節突然間出現拜偶像.
孩子們啊.
你們要遠避偶像.
舊的和諧本就是說.
小子們啊 你們要自我謹守.
遠避偶像.
突然間講偶像.
雖然第20節講的是耶穌是真神.

$^{601}$是永生.
就是說.
作者的內心.
寫到最後這句.
他緊張的就是.
我們信奉神而子之名的人.
有沒有看自己這個永生是這麼重要的禮物.
你真而重之.
你參與在那個祈禱的服事裡面.
這個祈禱的服事.
涵蓋的是人際關係.
你慢慢修飾自己.
不再是所謂翹口參手.
不再是得罪人多 稱呼人少.
你的內心.
剛才的詩歌說.
對不起 等一下的回應詩.
要學校主耶穌的謙柔.
學習主耶穌那份對別人的仁愛.
我們應該這樣才發揮到.
為人祈禱的功效.
如果別人看到你都已經轉身離開.
不可能你能為這樣的人祈禱.
不至於死的罪.
神可以來聽你.
來改變.
相信這是在教會裡面很重要的人際關係.
相處 團契的生活.
令到我們真是在神的裡面.
我們真是豐富.
不是一個概念.
是需要我們大家一起著力的去參與.
所以嚴格來說.
善上達不到這個效果.
何來你能夠見到有弟兄犯不至於死的罪.
你為他祈求呢.
剛才我都說過了.
昨晚的港道都說過了.
我們星期五晚執事會.
我們在研究一個問題.

$^{641}$發現了在晚堂直播.
我們在實體的崇拜有超過八十多人.
回來崇拜實體.
但是同時間發生.
而且是很平均地每個月都是這樣.
就是有一半人.
四十多人.
在線上看著我們的晚堂直播.
我們就討論了.
什麼原因呢.
是偶然間患病回不了教會呢.
還是他還是怕去外面呢.
還是有原因要看著小朋友.
等等的原因.
還是怎麼樣呢.
我們沒有數據.
我們很想去了解弟兄姊妹.
為什麼仍然要透過一個螢光幕去看崇拜.
而不是參與一個實體的崇拜.
如果根據經文告訴我們.
如果這件事是這麼嚴肅.
這麼重要.
和這麼寶貴的.
參與在神的國度裡面.
將弟兄姊妹這個永生的禮物.
能夠常常保存.
不要被惡者侵害.
不要慢慢被這個惡者搶走.
就需要人與人之間的互相關愛.
所以團契的生活.
彼此的關愛.
就是很重要.
其實不知道大家有沒有這個感覺呢.
我們這個月開這個茶水檔.
有那個人氣.
有那個溫度.
大家寒暄一兩句.
已經感覺到肢體之間的相交.
那份甜蜜.
真的不是那杯咖啡.

$^{681}$不是那杯茶.
是大家願意彼此的相交.
所以我們要準備自己與人相交.
成為一個真是良善的人.
來將神教導我們.
我們自己保存我們這個永生.
記住作者描繪兩個世界.
屬神的神會保護.
我們是屬神的.
但是全世界臥在那個惡者的手下.
就是說這個世界是有一個敵對神的.
一個這樣的勢力存在.
不同的方式會偷走我們的下一代.
會偷走我們自己對永生的熱誠.
所以今早第三堂.
很重要的就是我們的信仰.
是否尋求真實的神的敬拜呢.
還是我們自己只是在拜我們自己的內心的慾望.
那個的神那個的偶像呢.
有些人拜的可能是成功.
可能是權力.
可能是他自己可以無拘無束的.
喜歡做什麼都可以.
那這個是不是聖經說的對真神的敬拜呢.
求主幫助我們能夠彼此有一個反思.
彼此有一個在神的裡面.
我們現在我們的袋裡面.
這個這麼重要的禮物.
是否已經慢慢地失去了呢.
這個就是永生耶穌賜給我們的.
我們一起祈禱.
多謝主帶領我們今早在你面前的一個思想.
盼望我們的信仰是尋求真實.
我們的主你就是真神.
更加是永生.
求主幫助.
這樣的祈禱奉耶穌基督的名.
阿們.
\newpage



\section{提摩太前書 6:6-10}
\label{sec:_BOzzljBDkA}
\textbf{2023-03-25 晚堂《知足人生》提摩太前書6\_6-10 講員:盧思源姑娘}
\newline
\newline
連結: \href{https://youtube.com/watch?v=-BOzzljBDkA}{\texttt{https://youtube.com/watch?v=-BOzzljBDkA}} ~~~~ 語音日期: 2023-04-05
\newline
\newline
\hyperref[sec:ythaF42qv_4]{\small{< < < PREV SERMON < < <}}
~
\hyperref[sec:index_chronic]{\small{[返順時目]}}
~
\hyperref[sec:index_scriptual]{\small{[返順卷目]}}
~
\hyperref[sec:lYwkS21XVXg]{\small{> > > NEXT SERMON > > >}}
\newline
\newline
提摩太前書 6:6-10
\newline
\begin{longtable}{cl}
\hline
\hline
章節 & 經文 (和合本修訂版)\\
\hline
6:6 & \begin{tabularx}{0.7\textwidth}{X} 其實，敬虔加上知足就是大利。 \end{tabularx} \\ \\ \relax
6:7 & \begin{tabularx}{0.7\textwidth}{X} 因為我們沒有帶甚麼到世上來，也不能帶甚麼去； \end{tabularx} \\ \\ \relax
6:8 & \begin{tabularx}{0.7\textwidth}{X} 只要有衣有食，我們就該知足。 \end{tabularx} \\ \\ \relax
6:9 & \begin{tabularx}{0.7\textwidth}{X} 但那些想要發財的人就陷在誘惑、羅網和許多無知有害的慾望中，使人沉淪，以致敗壞和滅亡。 \end{tabularx} \\ \\ \relax
6:10 & \begin{tabularx}{0.7\textwidth}{X} 貪財是萬惡之根。有人因貪戀錢財而背離信仰，用許多愁苦把自己刺透了。 \end{tabularx} \\ \\
[1ex]
\hline
\hline
\end{longtable}
$^{1}$弟兄姊妹平安.
願意開始之前我們都有個祈禱.
參與我們做一次恭敬的去到你的面前.
一同去敬拜 一同去聽你的話語.
一膝奸孩子要分享你的信息的時候.
求真理的聖靈都與我們同在.
願意我們無論說的聽的都進入在你的真理當中.
避免你自己的話語得著提醒 得著詛咒 得著怨念.
我們恭敬的將以下時間交上.
奉主耶穌基督明治祈求 阿們.
弟兄姊妹平安.
開始之前想跟大家說一個故事.
大家都知道魔鬼的工作是什麼.
魔鬼的工作就是很喜歡搞破壞.
從前有一隻老魔鬼.
他看到人間的生活太過幸福.
他就想我要去擾亂一下.
我不想人們過得那麼開心.
於是他就派一隻小魔鬼出去擾亂一個農夫.
因為他看到這個農夫每天很努力地工作.
雖然他所得的很少.
但他仍然過得很開心 很知足.
這隻小魔鬼就在想.
怎樣才可以將這個農夫變得壞一點.
他就想到了.
他就將這個農夫的田地變得很硬.
讓這個農夫可以知難而退.
但這個農夫雖然看到地很硬.
但他仍然很努力地去挖地.
挖了大半天 做得很辛苦.
但他只是休息一下.
他又繼續挖地.
他一點抱怨都沒有.
這隻小魔鬼看到這個計劃失敗了.
他就垂頭喪氣地走了.
這個老魔鬼又派第二隻小魔鬼去.
這次這個小魔鬼就在想.
既然這個農夫被他這麼辛苦工作也沒有用.
好 不如這次我拿走他擁有的東西.
這隻小魔鬼就想.

$^{41}$我先將他的午餐拿走.
就將他的麵包和水都拿走.
他就想 這個農夫你做得這麼辛苦.
又累又餓.
如果連麵包和水都不見的話.
你一定會覺得很不開心.
一定會發脾氣.
但這位農夫當他又肚餓又口渴的時候.
想著休息一下的時候.
看到麵包和水不見了.
他就在想.
不知道哪個人這麼可憐.
比我更加需要這些麵包和水.
如果這些麵包和水可以讓他得到溫飽的話.
那就好了.
第二隻小魔鬼又覺得失敗了.
這個老魔鬼就覺得很奇怪.
難道不可以有辦法令這個農夫變得更壞嗎.
於是在這個時候.
剛好第三隻小魔鬼出現.
出來就跟這個老魔鬼說.
我有辦法.
我一定可以將他變得更壞.
大家想不想到他怎麼做呢.
於是這個小魔鬼.
原來他跟農夫做朋友.
農夫很開心跟他做朋友.
因為這個小魔鬼有預知的能力.
他就告訴農夫.
明年這個時間會有乾旱.
就教農夫把稻種種在濕的地方.
不要種在乾旱的地方.
農夫就照做.
結果當第二年所有人都沒有收成的時候.
這個農夫就滿滿的收成.
因此他就富裕起來.
這個小魔鬼每年都跟農夫說.
應該種些什麼不要種些什麼.
三年過去了.
這個農夫就變得很有錢.

$^{81}$他又教農夫把種到的米和麥拿去賣.
拿去釀酒.
賣一賣這些酒就可以賺錢.
更加賺得多.
慢慢地這個農夫就開始不再耕種.
他開始靠這些販賣的方式去生活.
而且賺到很多很多的錢.
有一天老魔鬼就來了.
就見到小魔鬼.
小魔鬼就說.
現在我給你看看我的成果.
這個農夫已經開始變了.
有一天農夫在家裡就舉辦了一個宴會.
請了很多很多有錢的人來參加.
喝最好的酒.
吃最好的東西.
而且有很多工人去侍候他.
但是他們就很浪費.
大吃大喝.
甚至喝到不省人事.
而這個農夫原來已經開始變得很貪圖這些的燕樂.
也很貪吃很貪吃.
小魔鬼就跟老魔鬼說.
等一下.
你看看.
他還有些不同的.
在這個時候.
有一個工人就拿了一些葡萄酒出來.
不小心.
掉了.
這個農夫怎麼做呢.
竟然就在那裡破口大罵.
你做事有沒有這麼不小心.
你知不知道這些東西很貴的.
接著那個工人就說.
主人.
我到現在都沒有吃過飯.
餓到已經沒有力氣了.
但是這個農夫卻在那裡說.
你的東西都還沒有做完.

$^{121}$你怎麼可以吃飯呢.
這個農夫原來開始變得很兇很自私.
老魔鬼看到的時候就很開心地跟這個小魔鬼說.
你這麼厲害.
你怎麼可以將它變成這樣.
這個小魔鬼就說.
很簡單.
我只不過讓它擁有比它需要的東西更多.
就可以慢慢引發它內心的貪婪.
慢慢慢慢我就可以誘惑它.
弟兄姊妹.
當然我們只聽到這個只是一個故事.
但是有時候它在我們人的心裡面.
很多時候不其然都會想要更多.
舉一個例子.
可能有沒有試過給一些糖果或者一些零食給小朋友呢.
如果你給他一顆糖果的時候.
他看到你手上原來你還有很多糖果.
又或者你正在用很多的時候.
可能有一部分小朋友都可能想要多一顆.
人大了當然我們現在就不是想要糖果.
可能我們上班我們去賺錢的時候.
假如我們賺到十塊錢.
可能我們都會想.
怎樣可以賺到二十塊錢呢.
甚至想賺到更加多.
這些很多時候都是很正常.
也都是很多時候我們大家都會想到的東西.
因為今天我們活在這個社會的時候.
活在這個環境.
我們就是需要金錢.
我們衣食住行.
每樣都要錢.
沒有錢又怎樣去過生活呢.
所以聖經一向都不是說我們基督徒.
我們不是說我們不食人間煙火.
不是說我們不需要錢去過生活.
而是問題我們怎樣去看待金錢.
正如今天我們剛才所讀到的經文.
保羅所說的不是金錢是萬惡之根.

$^{161}$他不是這樣說.
他是在說貪財是萬惡之根.
我們看到財富本身是沒有錯的.
但是假如人對他有錯誤的渴求的時候.
這個才是錯.
而且如果我們看回經文.
裡面說貪財不是萬惡唯一的根源.
因為如果在原文我們看根.
他是沒有一個貫詞.
沒有一個der.
就是英文簡單一點.
他是各種萬惡裡面其中的一個根源.
如果我們看保羅.
為甚麼他會說這段話呢.
當我們翻開聖經.
我們可以看的時候.
看前一點.
原來保羅在第六章的段落裡面.
他提到當時有一些假教師.
原來這班人所教導的.
不是符合主耶穌基督的話語.
也不是說一些合乎敬虔的教導.
甚至他們挫靶敬虔.
當作一些盈利的手段.
估計當時的人.
那班自稱為敬虔的人.
他們有達到靈情高峰的秘訣.
所以他們就向一些想從他們手中.
獲得這些秘訣的人.
拿一些好處.
所以在這段經文裡面.
保羅提醒提摩泰.
要小心啊.
不要學這些人.
保羅提醒提摩泰.
敬虔本身已經可以帶來益處.
但是只要在知足的配合下.
才會發生.
接著如果你看聖經的時候.
很有趣.

$^{201}$保羅不是即時說我們要知足.
因為有耶穌基督就是我們的滿足.
保羅在這裡不是這樣說.
雖然我們知道保羅的生命.
他的滿足的而且確是因為耶穌基督.
因為保羅在《菲勒比書》四章都說.
「我並不是因缺乏而說者說.
因為我已經學會在無論什麼境況都可以知足」.
所以我們看到保羅的生命.
他是因著耶穌基督而滿足.
但是他在這裡並不是這樣說.
在第七節他所表達的.
是一句很多時候我們中國人經常都聽到的格言.
就好像說生帶不到什麼.
死也帶不走什麼.
保羅在這裡說.
「因為我們沒有帶什麼到世上來.
也不能帶什麼去」.
而在第八節更加說了.
滿足是要有衣有食.
只要有衣有食.
我們就應該知足.
其實在這裡看到.
保羅是很體現到當時耶穌基督的教導.
如果在《馬太福音》六章.
我相信大家都很熟悉的.
六章二十五節那裡.
「昔日耶穌教導門徒.
不要為了衣食而憂慮」.
當時耶穌這樣跟門徒說.
他說「不要為了生命憂慮吃什麼喝什麼.
不要為了身體憂慮穿什麼.
我們的生命不勝於飲食嗎.
我們的身體不勝於衣裳嗎.
我們看看天上的飛鳥.
也不種也不收.
也不積蓄在倉庫.
但是我們的天父尚且養活它.
我們不比飛鳥更加貴重嗎.
還有野地的百合花.

$^{241}$也不枋善也不勞苦.
但是就是所羅門極榮華的時候.
也不及這朵花.
所以當時耶穌教導門徒.
不要憂慮吃什麼喝什麼.
這些都是愛邦人所求的.
你們需要的一切.
我們的天父都知道.
我們要先求神的國和神的義.
這些都交給我們.
所以我們看到.
當時耶穌也是在說衣食.
祂知道衣食就是我們人的需要.
而這些需要.
神應許祂必然會去供應我們.
所以保羅在這裡提醒.
只要我們有衣有食.
我們就應該知足.
那什麼為知足呢.
知足的定義究竟是什麼呢.
我上網看關於知足的定義.
知足其實就是指對我們自己的生活.
我們的財富.
我們的地位.
我們擁有的東西.
我們感到滿足感.
我們有感激之情.
當中強調我們珍惜已經有的東西.
我們不貪心.
我們不嫉妒.
我們滿足現在有的境遇和條件.
我們不追求過多物質的財富.
而知足的人.
就是懂得去欣賞生活的美好.
懂得去感恩身邊的人和事.
也不會被外界的壓力.
或者輕易動搖內心的平靜和滿足感.
那問一下.
知足和滿足有沒有分別呢.
其實知足和滿足.

$^{281}$其實都是在說感到滿意滿足.
但是當中的側重點是有些不同的.
剛才說知足是強調我們對已經有的東西.
我們感到珍惜和感激.
我們不會追求過多的物質或者財富.
而滿足就是側重強調.
當你達到某個目標.
或者獲得某種願望而感到滿足.
所以知足和滿足雖然很相似.
但是知足是強調內心裡面的平靜和滿足感.
而滿足是強調外在的成就和目標達成.
剛才一開始的時候提到.
我們人很容易就會追求更多.
我們很容易得一.
我們很容易想得二.
那怎樣才學會真正的知足呢.
還是我們很多時候都是在追求一種外在的滿足呢.
我看了一本書.
我覺得也挺好的.
它裡面說的也挺好的.
可以跟大家分享一下.
它說很多時候我們很想要很多東西的時候.
提議我們想買很多東西.
或者想要很多東西的時候.
它提議我們可以找一張紙去寫下來.
一欄就寫著「需要」.
一欄就寫著「想要」.
我們想一下.
究竟我們的東西是我們想要的.
還是我們真正需要的呢.
到後來我們可能發覺.
可能我們想要的東西.
是多於我們實際需要的東西.
而聖經也有很好的提醒.
為甚麼我們要知足.
保留在這裡說.
因為就算今天我們每個人想要很多東西也好.
其實到最後是帶不走的.
聖經裡面說.
因為我們沒有帶甚麼來世.

$^{321}$也不可以帶走甚麼.
今天我們在出生的時候.
我們是赤身而來.
到我們每一個人.
我們面對人生終結那一刻.
無論我們擁有多少東西也好.
我們都真的帶不走的.
分享一件事給大家聽一下.
最近可能有人知道我剛剛結婚.
我就要搬家.
搬家當然要收拾家居.
在收拾家居的過程裡面.
我發覺自己真的有很多東西.
大家要知道.
我家不是像錦繡那樣的一間大屋.
我只是一個小單位.
一個兩房的小單位.
有兩個衣櫃.
但是在收拾的過程裡面.
有很多不同的東西.
特別是我發覺自己有很多衣服.
很多都是在我讀神學之前.
就是工作的時候買的.
當在收拾衣服的過程中.
我發覺有些衣服還沒有穿過.
連牌子也沒有剪過.
就有很多在這裡.
但是如果大家知道或者認識我.
其實現在我經常穿來穿去.
都是那幾套衣服.
因為接下來搬家沒有那麼多地方放.
最後我是扔掉了超過十大袋的衣服.
去救世軍的回收箱.
是十大袋.
在過程裡面我發現.
原來我擁有很多東西.
但是其實可能實際需要的只是很少.
鼓勵大家有時間都可以整理一下.
自己的衣櫃 自己家裡的東西.
又或者好像剛才那本書提議.

$^{361}$我們可能想一下有些什麼是我們需要.
或者有些什麼是我們想要的.
可能發現原來我們擁有的東西.
比我們實際需要的多了很多.
回到經文裡面.
保羅在第九節提到.
但那些想要發財的人.
就陷在誘惑,羅莽和許多無知有害的慾望中.
使人沉淪以至敗壞和滅亡.
貪財是萬惡之根.
有人因貪戀錢財而背離信仰.
用許多仇苦把自己刺透了.
保羅在這裡說.
剛才也提及過不是財富有錯.
而是那些想要發財.
或者貪財.
或者貪戀錢財.
這個才是問題.
財富本身是沒有錯的.
而是人對它有錯誤的渴求的時候.
這個才是錯.
我們作為基督徒.
我們怎樣去看待錢財呢.
我們要知道.
擁有錢財.
即是有錢不是罪.
有時很多基督徒.
對於主耶穌關於錢財的教訓.
有時有一種錯誤的領會.
以為是否主耶穌排斥有錢的人.
因為在路加福音時說.
有錢財的人.
進神的國是何等的難呢.
但實際上.
其實耶穌沒有歧視過有錢的人.
也擁有錢財不是有錯,不是有罪.
甚至在耶穌的門徒裡面.
如果我們看馬太福音27章.
耶穌的門徒當中.
也有財主.

$^{401}$所以問題不是錢財本身.
而是如果我們可以善用我們的錢財.
它可以產生良好的作用.
相反.
如果我們對錢財的態度不正確.
是會帶來很可怕的後果.
我想最近大家都看報章.
看到一宗很轟動全城的謀殺案件.
我們也很難想像.
為何一個人.
可以對另一個人下毒手.
甚至肢解.
消息報道.
是因著一些錢銀的瓜葛引起.
在裡面我們看到.
原來錢可以讓人失去理性.
可以讓人做出一些完全違反人性的事.
所以我們看到聖經說得很對.
它說那些想要發財的人.
是陷在誘惑,羅莽,慾望裡面.
甚至是使人沉淪,敗壞,滅亡.
所以保羅為何要教導提摩太.
要知足.
因為原來當我們為著我們所擁有的去感恩的時候.
我們不會貪戀,追求更多的東西.
不過在這裡可能有些人會這樣說.
知足的人就是很甘於現狀.
是不是很停滯不前.
那就是沒有進步.
是不是這樣.
事實.
事實是不是這樣呢.
我們知道因為有些人可能不滿足現狀.
沒錯.
因為他不滿足現狀.
所以他會不斷地努力.
不斷地可能去改變.
只不過如果一個人.
他是不知滿足.
其實想深一層.

$^{441}$他的慾望是沒有盡頭的.
他是有無窮無盡的慾望.
甚至乎有無窮無盡的追求.
或者這個追求可以帶來物質的豐裕.
但是帶來不了精神上的滿足.
如果一個不知足的人.
就算他得到很多東西也好.
若然他不可以因著他所擁有的而滿足.
而不可以因著他所擁有的而獲得快樂的話.
那又有什麼用呢.
所以我們看到在聖經裡面.
很刻意地教我們要小心一點對錢財.
因為錢財是有影響力的.
在馬太福音.
耶穌也說了第六章.
你的財寶在哪裡.
你的心在哪裡.
耶穌也說過一個人.
不可以事奉兩個主.
不可以護這個愛那個.
就是輕這個重那個.
我們不可以又事奉神又事奉馬門.
大家都知道馬門就是金錢.
而且當耶穌說殺種的比喻.
大家都很熟悉殺種的比喻的時候.
大家記得殺種的比喻.
耶穌說的時候殺了四個地方.
第一個就是在路旁.
當他殺在路旁的時候.
就是被鳥獸去啄了.
第二個當他殺在石頭的時候.
因為殺在石頭裡面沒有根.
所以太陽一曬那些種子就已經枯萎了.
第三個就是殺在荊棘裡面.
不過因為荊棘阻礙種子生長的時候.
所以它也生長不好.
而耶穌說這個比喻的時候.
他怎樣形容荊棘呢.
他說荊棘就好像世上的憂慮.
也是財富的迷惑和各種的私慾.

$^{481}$所以我們看到財富是會迷惑人的.
它很容易將神的道去制住.
令到神的道生長不了.
而且我們要知道魔鬼也很會利用錢財去霸佔人的心.
正如一開始的時候我說的那個故事一樣.
龍夫的故事一樣.
魔鬼很會利用金錢.
叫人在不知不覺當中受它支配奴役.
而且是不自覺.
所以保羅在這裡說貪財是萬惡之根.
有人貪戀錢財背離信仰.
用許多仇苦把自己刺透.
但是很可能我們自己也會想.
我們也不算是貪戀錢財.
起碼我們在座的也不會為錢甚麼都去做.
我想我們大部份人也會覺得.
我們有錢救就行了.
我們也不是為了錢甚麼都去做的人.
雖然我們也未至於去到很貪戀錢財.
但是會不會有時不其然.
將我們的安全感放在金錢上呢.
我一個朋友也是一個基督徒.
他有回教會也有奉獻有事奉等等.
但是因著一些原因.
家裡需要一些現金去周轉.
他的積蓄當然因著這些原因減少.
當然他沒有告訴我剩下多少.
不過他有少許分享.
他說當眼見自己的積蓄越來越少的時候.
他內心那份安全感也慢慢少了.
另外我不知道在座有沒有投資的朋友.
我有些投資的朋友他跟我分享.
他們說他的心情很多時候是隨著市的起跌而起跌.
市升他的心情就會好一點.
市跌他的心情就會跌一點.
有時可能我們的心情很容易被金錢牽引著.
所以如果我們今天看聖經.
今天沒有讀到的.
如果看回提摩太前書六章十七節.
保羅也昔日提當時的基督徒.

$^{521}$不要倚靠靠不住的錢財.
要倚靠像百物給我們享受的神.
如果在新譯本那句.
他說不要盼望寄託在無常的財富上.
卻要寄託在賜給我們百物享用的神上.
聖經當中教導我們金錢是靠不住的.
相反我們是要靠賜給我們萬物的神.
弟兄姊妹.
今天我們怎樣去看金錢呢.
沒錯在這個世界上我們是需要金錢去生活.
但會不會我們有時不自覺的.
就是將我們的安全感投放在金錢上.
而忘記了神才是我們最大的倚靠呢.
求主幫助我們.
幫助我們學習定精在祂的身上.
我們不被財富所迷惑.
最後盼望今天的經文.
今天的訊息都成為我們彼此的勉勵.
求主幫助我們去學習我們過一個知足的人生.
為著我們所擁有的我們去感恩.
也求主給我們智慧去善用神給我們的一切.
願意我們一同去同心圖告.
藉著昔日保羅寫信給提摩太.
你再一次去教導我們.
我們要有知足的心.
為著你所給我們一切我們獻上感恩.
因為我們知道今天我們在這個世界上.
我們只不過是寄居.
我們將來有永恆的天價.
求你幫助我們.
讓我們有這份屬靈的眼睛.
我們懂得去看我們今生所擁有的東西.
我們也懂得去珍惜今天你給我們的所有東西.
願意我們真正定精在你的身上.
知道在這個世界上金錢真的有很多誘惑.
但求主給我們智慧.
懂得去運用也善用你給我們的金錢.
以致在我們這個人生裡面.
我們過一個榮耀神利的人生.
過一個滿足神利自己心意的人生.

$^{561}$過一個行在你心意的人生.
求你幫助.
天父多謝你聽我們的禱告.
奉主耶穌基督名字祈求.
阿們.
\newpage



\section{使徒行傳 9:26-30-11:22-26}
\label{sec:lYwkS21XVXg}
\textbf{2023-04-02 加恩樓《生命師傅的典範》使徒行傳9\_26-30,11\_22-26 講員:陳一華牧師}
\newline
\newline
連結: \href{https://youtube.com/watch?v=lYwkS21XVXg}{\texttt{https://youtube.com/watch?v=lYwkS21XVXg}} ~~~~ 語音日期: 2023-04-05
\newline
\newline
\hyperref[sec:_BOzzljBDkA]{\small{< < < PREV SERMON < < <}}
~
\hyperref[sec:index_chronic]{\small{[返順時目]}}
~
\hyperref[sec:index_scriptual]{\small{[返順卷目]}}
~
\hyperref[sec:Am9T43vQMKQ]{\small{> > > NEXT SERMON > > >}}
\newline
\newline
$^{1}$弟兄姊妹,早安!.
今早很高興和大家一起去敬拜.
我在唱詩歌和你們一起敬拜的時候.
很有家的味道.
家的味道.
為什麼呢?.
由台上打鼓到台下敬拜.
都有很多小朋友在當中.
很多小朋友都和大人一起去敬拜.
所以這種感覺是非常濃厚的.
今天也很感恩和你們一起.
在四月第一個星期一起去敬拜.
我們先開始有一個祈禱.
天父,願你僕人口中的說話.
心裡的意念在你的面前蒙月.
也願今天我們聽到,講到的當中.
能夠同感一靈.
以致我們同德造就.
我們同心仰望,祈禱.
奉耶穌的名求.
阿們.
和大家今天思想的題目.
就是生命師傅的典法.
在上個世紀.
如果你想有一門手藝.
能夠希望有一技之長.
磅下身的話.
我相信你會跟隨一個師傅.
跟著和他一起工作,一起住,一起吃.
很多事情你都跟著他.
一腳踢地做.
這個就叫做學師仔.
大家聽過這個名字.
來到現在這個千禧年代.
這個世紀.
如果你說我很想找一個資深一點的.
有經驗的人去跟隨他.
學習一下.
希望他能夠過幾招給你.
現在這個名稱叫做使徒仔.

$^{41}$就是mentor和mentee.
我想在當中也有人做過mentor.
也有人做過mentee.
沒錯.
如果你說將這個概念.
應用在宗教和信仰那方面.
我相信有一個更加貼切的名字.
是很適合我們去形容的.
就是生命師傅.
我們讀聖經的時候.
我們發覺有兩個人.
大家有沒有留意到?.
這兩個人.
一個叫巴拿巴.
一個叫素羅.
合在一起的時候.
如果問大家你熟誰呢?.
就不用問阿貴了.
大家都很熟素羅.
因為素羅就是後來的保羅.
大家知道嗎?.
為什麼我們會這麼熟這個素羅呢?.
是因為素羅的名字的知名度是高的.
大家也知道他在教會當中.
做了很多事情.
包括了他去很多地方去傳福音.
建立很多教會.
所以有人說.
素羅是一個報道家.
是當之無愧的.
因為他到處去傳福音.
建立教會.
同樣地.
如果有人說.
素羅是一個神學家.
我相信大家都不會否認的.
因為他在聖經新約當中.
有十三封書信.
裡面揚入著很豐富.
很淵博的一個神學思想.

$^{81}$所以素羅保羅是一個神學家.
大家都公認的.
不過可能少人知道.
原來保羅也是一個生命師傅.
大家知道嗎?.
或者我試試提一兩個名字.
看看你們有沒有印象.
先提一些比較三普的.
大家有沒有人聽過一個名字叫做阿居拉和伯基拉.
有沒有聽過?.
沒有什麼反應.
即是沒有聽過.
有沒有人聽過阿尼西姆這個名字呢?.
好像沒有什麼反應.
不如說一個你們比較熟悉的.
提摩泰聽過沒有?.
每個人都點頭了.
可能你都叫Timothy都未定.
提摩泰.
好,大家很熟悉.
原來保羅都做過他們的生命師傅.
接著你會問.
哪一個人是保羅,是素羅的生命師傅呢?.
答案剛才已經說了出來.
就是巴拿巴.
就是素羅的生命師傅.
巴拿巴怎樣可以成為保羅的生命師傅呢?.
是因為他擁有三項特質.
在他的生命當中流露三項條件.
這個特質是很重要的.
不如說到現在我才知道.
我有一個PowerPoint在這裡.
給大家看看.
巴拿巴成為素羅的生命師傅.
他有三項特質.
第一個特質是什麼呢?.
就是巴拿巴樂意去接待.
當然你知道是接待素羅.
情況是這樣的.
大家剛才讀經文都知道.

$^{121}$素羅信了耶穌之後.
在大馬色的路上遇見主.
他就迴轉了.
他就悔改了.
信耶穌.
當他信了耶穌之後.
他就很熱心.
他就很想繼續追求靈命成長.
於是他就回到一間教會.
叫做耶路撒冷的教會.
但誰不知當他回到那間教會之後.
就好像潑冷水.
為什麼呢?.
因為當他進入耶路撒冷的教會的時候.
人們就害怕了他.
不敢接近他.
覺得素羅信耶穌是真是假.
你明白我的意思嗎?.
素羅你真的講清楚.
你是真信還是假信.
不要耍花樣.
你明白我的意思嗎?.
因為素羅的名字.
在當時耶路撒冷來說.
是臭名遠播的.
大家有沒有印象.
因為他經常捉拿基督徒去坐牢.
難怪在教會的人一聽到這個名字.
嘩!現在還見到一個真人.
還怕多幾次.
在這樣的情況之後.
人們都彈開.
一見到他就閃開.
個個都不想來接觸他.
所以當時素羅去到這間教會的時候.
就真的好像一盤冷水潑下來.
人們有這樣的反應都很正常.
為什麼呢?.
因為都不知道素羅是否真的信耶穌.
如果他來到我們教會.

$^{161}$搜集資料.
收集名單.
交給大祭司.
再抓我們的話.
那就不得了了.
所以他們有這種反應.
也是很正常的.
但是對素羅來說.
的確遇到三個困境.
因為人們不接受他的時候.
他遇到哪三個困境呢?.
第一個困境就是.
素羅的靈命是很幼嫩.
一個剛信耶穌的人.
剛信主的素羅來到教會.
其實他也很想靈命上有栽培.
有跟進.
但問題是人們覺得他好像傳染病.
個個一聽到他的名字就彈開了.
不敢靠近他.
沒有人夠膽去栽培他.
沒有人夠膽去培訓他.
其實對於素羅來說.
真是一個可以說是先天不足.
因為他回到教會.
剛信耶穌的時候.
他應該要去參加教會的初信班.
或者是回教會的主日學.
但現在教會裡沒有人來跟進栽培他.
所以對於靈命幼嫩的素羅來說.
他的信仰根基是很容易動搖的.
這是其一.
其二就是素羅回到教會的時候.
人們看他的時候是奇異的眼光.
又或者人們看他是一種懷疑的眼光.
到底你是真信還是假信呢?.
各位聽眾.
當你回到教會.
人們用這麼怪的眼光看你的時候.
你覺得很不自然.

$^{201}$你覺得渾身不自在.
為什麼人們會這樣看我呢?.
因為當時候的人的確對素羅是很抗拒的.
所以在這樣的情況下.
素羅就落在一個單人匹馬.
即是孤立無援的光景當中.
因為人們都不敢來接近他.
所以當時候我相信.
素羅回到教會.
受到很大的打擊.
各位靈體姐妹們.
這種情況如果發生在我們身上.
真的會很笨拙.
你有沒有想過.
如果你回到教會.
沒有人理你.
沒有人跟你聊天.
你坐在哪裡都沒有人關心.
你是不是在這裡.
都好像隱形了.
你明白我的意思嗎?.
如果你說我回到教會是這樣的情況.
你就一定不會再回到第二第三次.
為什麼呢?.
因為當時人們好像很冷漠.
我這樣形容.
你現在是不是開始明白多一點素羅的光景.
他的處境就是這樣.
很孤立,很孤單,模糊.
在這樣的情況下.
上帝竟然安排了一個人出現.
這個人是誰呢?.
就是巴拿巴.
巴拿巴的出現.
對素羅來說非常重要.
為什麼呢?.
因為巴拿巴當時是獨排眾議.
他去接待素羅.
因為他感覺到.
這個素羅與別不同.

$^{241}$在他的眼中可能覺得.
這個素羅會是我們教會將來一個莫大的祝福.
這個素羅可能就是我們教會一個很寶貴的屬靈的果子.
所以當時巴拿巴看素羅的時候.
他的眼光截然不同.
所以他就獨排眾議去接待素羅.
等於節目門說起接待.
你都要有一個概念.
原來當時在新約的時候.
你要去接待一個人.
即是說你要去栽培一個人的時候.
是怎樣的呢?.
原來你要跟他一起吃.
一起住.
一起生活.
一起事奉.
大家明白這個意思嗎?.
即是說你去接待一個人.
你去培訓一個人的時候.
你差不多要放上你整副的心機.
用你很多精神時間.
這樣為止是甚麼呢?.
你只要去接待一個人.
可能你覺得牧師這樣搞法.
我們香港地不可以.
為什麼呢?.
我們人多地細.
家裡又小.
怎麼可以接待一個人來住呢?.
你明白我的意思嗎?.
雖然我們有些客觀的困難.
不過我覺得巴拿巴的做法.
給我們有些啟迪.
就是當你去培訓一個人.
你想去栽培一個人的時候.
你要記得.
你不只是跟他吃個包.
喝杯咖啡.
這樣就完了.
你也不是.

$^{281}$你跟他單單開一次團契.
這樣就完了.
一個月一次團契.
或者兩個星期一次團契.
或者你跟他開一個小組.
開一次小組.
這樣就完了.
不是.
頂智慕們.
如果我們要去培訓.
栽培一個人的時候.
請你記得.
千萬不要忽略平日的關心.
你平日對他的關懷.
是很重要的.
你平日對他的那種關心.
你怎樣跟他聊天.
跟他聯繫.
你發一些金句給他.
你跟他一起祈禱.
你對他的關懷.
是非常重要的.
我記得在錦宣家.
我自己就參加星期三晚的祈禱會.
叫做沙田區祈禱會.
顧名思義.
這個祈禱會裡面.
大部份的人都是沙田友.
我們在這個祈禱會當中.
是有些不同的經歷.
就是我覺得這個祈禱會.
對我來說.
是一個很強的凝聚力.
因為我們不是單單祈禱這麼簡單.
我發覺這個祈禱會.
原來在平日.
他是會很留心關心.
祈禱會裡面的弟兄姊妹.
他會用電話聯絡弟兄姊妹.
如果遇到有弟兄姊妹不舒服.

$^{321}$病了.
或者家人病了的時候.
這個祈禱會.
會有兩三個人去醫院探病.
如果他知道自己的團友.
有些人家裡有些問題的時候.
同樣這個祈禱會.
也有兩三個人去家訪弟兄姊妹.
換句話來說.
當你在這個祈禱會當中.
無論你是什麼時候.
平日你的處境是有什麼需要.
甚至知道你快要做生日的時候.
他們會開個生日會和你慶祝.
是否開心?.
很開心.
弟兄姊妹們.
換句話來說.
一個對人靈命成長的栽培.
其實很重要的是什麼呢?.
是你平日對他的關愛有多少.
你有沒有留意到他在靈性上的需要呢?.
你有沒有和他一起讀聖經.
和他一起祈禱.
一起分享呢?.
各位弟兄姊妹們.
當我們這樣做的時候.
其實就好像巴拿巴接待素羅.
對他有一個很全面的關懷,關心.
以致更加可以鼓勵到他靈命的成長.
我們說完了.
作為生命師父的第一個特質.
是我們願意接待,栽培.
生命師父第二個特質是什麼呢?.
第二個特質.
我不是很習慣用PowerPoint.
說完之後才按給你們看.
第二個特質就是巴拿巴是很積極的事奉.
因為不知道是不是當時教會發生一些事情.
所以母會耶路撒冷教會.

$^{361}$就差派巴拿巴去到安提亞的地方.
去處理一些問題.
那你知不知道這次的任務是多少人去?.
原來這次的任務是一個人去.
那你說孤不孤單?.
淒不淒涼?.
弟兄姊妹們你們當中有沒有人去過短孫?.
去過短孫舉手給我看看.
有沒有人去過短孫?.
有沒有說參加過短孫是你一個人去的?.
沒有的嘛.
如果你一個人去.
通常都是幾倫盡的.
你明不明白我的意思?.
其實巴拿巴當時被教會這樣委任差派.
他心裡可以去想.
這次教會是否籌備得很倉促?.
你明不明白我的意思?.
只有我一個人去的?.
這個可能是第一個問題會問的.
可能去問第二個問題.
這次為甚麼教會只派我一個人去呢?.
是不是這次的差事沒有人喜歡去呢?.
通常沒有人去只是派你一個人去.
你有沒有留意到?.
沒有人喜歡做的事通常都是派你一個人去.
你有沒有留意到?.
如果派你一個人去.
有時你心想不妨是好得多了.
怎麼搞的?.
每個人都不派?.
只是派我一個?.
其實很多人心裡都有這些問號.
是很自然的.
不過對巴拿巴來說.
他作為一個生命的師父.
他就在他的事奉上流露了一個美好的見證.
他是甚麼?.
他是順服差派.
他沒有質詢教會.

$^{401}$為甚麼你這樣去做.
沒有懷疑教會.
為甚麼你這樣派我出去.
他就是甚麼?.
他順服差派.
單獨一個人去到安提亞的地方.
當他一個人去到安提亞的地方的時候.
各位弟兄姊妹.
他就很用心的.
跟安提亞教會的弟兄姊妹.
大家有談有交流有分享有了解.
明白知道教會有甚麼需要.
遇到甚麼問題.
結果他跟安提亞教會的人.
是非常融洽.
而且都能夠在主的裡面.
大家互相的勉勵.
互相的鼓勵.
就算發生甚麼問題都好.
都讓這些問題可以找出解決的方法.
各位弟兄姊妹們.
你有沒有留意到原來.
當巴拿巴這樣去服侍.
以一種順服的心態.
積極的去事奉的時候.
其實就帶給教會一個非常美好的結果.
就是令到安提亞教會的弟兄姊妹.
能夠跟母會的關係建立得更加緊密.
你有沒有留意到原來.
作為一個生命的師父.
我們的事奉必須要是一個.
很投入很積極的事奉.
在我們錦宣家.
有一對年長的夫婦.
不過我都不說他們的名字.
讓你知道是哪對夫婦.
他們有一個特別的習慣.
就是通常你知道.
年紀大的弟兄姊妹.
回來教會.

$^{441}$通常都會一進教會.
找回他們最熟悉的兩張椅子.
坐下來.
然後有五分鐘時間讓他們癱頭.
癱頭完之後就開始敬拜.
來進入崇拜.
不過這對夫婦不是.
這對夫婦很特別.
每一次崇拜之前.
她都會站在大門口的地方.
她不是站在很正中.
為什麼呢?.
因為她知道牧師傳道人在中間招呼人.
但這對夫婦都是從旁協助牧師傳道人.
去歡迎一些新朋友.
你明白我的意思嗎?.
去跟一些新人聊天.
去關心,去問候.
令到新朋友或新的弟兄姊妹來到的時候.
感覺到有一種很溫暖的感受.
有一種賓至如歸的感受.
所以後來有個弟兄.
他就講見證了.
他說我來到這間教會的時候.
我很感恩.
我可以享受到這對夫婦對我的接待.
所以我就決定留在這間教會來參加崇拜.
他還多說一句.
他說我亦都決定.
由今年開始我會參加崇拜部的事奉.
各位弟兄姊妹們.
當他這樣去分享的時候.
我就發現到.
原來當你投入事奉.
積極事奉的時候.
你未必會發覺.
但別人會發覺.
你那種表現.
令他留下了很深刻的印象.
而且你的事奉也留下了美好的榜樣.

$^{481}$不知不覺間.
你就能夠影響到其他的弟兄姊妹.
所以各位弟兄姊妹們.
作為一個生命的師父.
我們的事奉必須是以身作則.
聖靈內在我們心裡.
將不同的恩賜賜給我們.
各位弟兄姊妹們.
你有沒有留意到.
原來當上帝給我們不同的恩賜的時候.
是要叫我們拿出來用的.
當你肯用你的恩賜去事奉的時候.
我們屬靈的生命才會成長.
如果你藏起來不用.
你的靈命不會有成長.
第二.
當你能夠用上帝給你的恩賜去事奉的時候.
你的積極的事奉.
其實就留下了一個很美好的榜樣.
給身邊的人來學習和跟隨.
所以各位弟兄姊妹們.
巴拿巴能夠成為一個生命的師父.
因為他自己以身作則.
所以他一種積極的事奉.
就帶給身邊的人.
有一個非常美好的典範.
我們再去看第三.
巴拿巴擁有的特質是什麼呢?.
巴拿巴第三他的生命師父的特質是什麼呢?.
就是他主動去尋找.
找誰呢?.
去找塑羅.
可能你會問.
為什麼他會無緣無故去找塑羅呢?.
各位弟兄姊妹們.
一定不會是無緣無故的.
為什麼呢?.
多謝弟兄幫我彈一些泡泡泡出來.
為什麼他要找塑羅呢?.
是因為當時候.

$^{521}$你有印象嗎?.
塑羅信了耶穌之後.
他很積極.
他去傳福音.
他甚至和猶太教的兄弟去辯駁.
在這樣的情況下.
就危險了.
因為這一批猶太教的兄弟.
非常不喜歡塑羅.
覺得塑羅是叛徒.
到底你和我坐一艘船.
大家都是什麼?.
逼伯基督徒.
逼伯基督教.
但現在你跳船.
你不單止這樣.
信了基督教.
你還更加大發熱心去傳道.
去傳福音.
所以這一班猶太教的兄弟.
看在眼內是怎樣呢?.
火都來了.
你有沒有留意到?.
於是他們就要想辦法.
要殺這個塑羅.
在這樣的光景下.
塑羅真的不知道怎樣去做.
當時竟然有一班兄弟姐妹.
想到一個方法.
告訴塑羅.
既然現在處境這麼危險.
不如這樣.
我們提議你.
安排你回鄉下.
大家剛才讀聖經.
都讀到.
塑羅的鄉下在哪裡?.
大塑的地方.
各位頂智妹妹.
你知道嗎?.

$^{561}$回鄉下.
這三個字.
可圈可點.
正面地去想.
回鄉下.
即是說你很有成就.
這就叫做衣錦榮歸.
但另一個意思.
回鄉下是甚麼?.
是一個貶義.
即是說.
不要說執甚麼.
總之.
回鄉下的意思是.
你一生人的成就.
到此為止.
你明白我的意思嗎?.
即是說.
你從此之後.
就是回鄉下耕田.
意思就是.
你一生人玩完.
大家明白這個意思嗎?.
所以.
很不幸.
這個回鄉下三個字.
對塑羅來說.
是一個後者.
是一個貶義.
你知不知道.
當塑羅真的被弟兄姊妹送回鄉下之後.
其實對塑羅來說.
他面對五個危機.
在當中.
我們看看.
弟兄姊妹們.
第一個危機是甚麼呢?.
第一個危機就是.
初同歲月.
其實.

$^{601}$塑羅這樣回鄉下避一避.
不知道避多久.
等一等.
又不知道甚麼.
等多久.
大家都明白.
因為你經歷過甚麼叫等.
尤其是.
當你等一些事情.
是很渺茫的時候.
當你等一些事情.
都沒有把握的時候.
但你又要等這個消息.
因為你寄了一百多封信出去.
找工作.
都不知道何時有封信.
你明白我的意思嗎?.
所以人們說.
等到頸都長了.
為甚麼呢?.
因為等是很辛苦的.
現在塑羅避一避.
不知道避多久.
這樣的光景.
不知道可以捱多久.
對塑羅來說.
是蹉跎他的歲月.
一個怎樣心智頑強的人.
就最怕在這樣的光景裡.
在百無聊事.
都不知道要等多久.
所以很擔心.
從此會不會燒沉下去.
這是第一個危機.
大家認同嗎?.
第二個危機是甚麼呢?.
第二個危機就是被人遺忘.
換句話來說.
在我們的教會裡.
這麼多弟兄姊妹.

$^{641}$好說不好聽.
多你一個.
不多得多少.
你明白我的意思嗎?.
少你一個.
又不少得多少.
現在你回到鄉下.
更加甚麼呢?.
被人忘記了.
被人遺忘.
就成為了第二個危機.
我們再看.
第三個危機是甚麼呢?.
第三個危機就是不受歡迎.
因為塑羅這個人.
臭名遠播.
所以大家都知道.
他是一個野火的人物.
他是一個瘋眼的人物.
是一個爭議性的人物.
所以一個這樣的人.
在我們的教會.
只會麻煩多多.
你看猶太教的人都追殺他了.
現在他回到鄉下.
好了.
一天都光了.
從此教會太平了.
大家明白我的意思嗎?.
所以你不要回來.
最好.
大家就知道.
這就是塑羅的處境.
我們再看第四個危機是甚麼呢?.
第四個危機.
就是不被重視.
因為他剛剛才開始回教會.
在耶路撒冷教會.
有幾長日子.
不是很長.

$^{681}$多不多人認識他?.
多,不過認識他的.
通常都是不好的.
以前的惡行.
所以塑羅對我們來說.
對教會來說是負資產.
是嗎?.
不需要那麼計較.
不需要那麼認真.
他在不在都無所謂.
反正他現在不在更好.
少了一些負資產.
大家明白我的意思嗎?.
這是第四個危機.
我們再看第五個危機是甚麼呢?.
是最淪盡的.
最淪盡的就是這個危機.
就是從此之後.
因為塑羅這樣回到鄉下.
被人遺忘不被重視之後.
我想問大家.
他翻身的機會大嗎?.
從此之後.
不會再有機會.
回到耶路撒冷教會.
所以是一個永不翻身.
因為基於他過去的臭名遠播.
他的惡行.
哪裡有機會再讓他翻身呢?.
大家明白我的意思嗎?.
好了,請你想一想.
假使塑羅真的消失在教會的時候.
你知不知道會帶來甚麼後果呢?.
請你想一想.
我想告訴大家.
次低限度有兩個後果.
是最明顯的.
第一,如果塑羅真的消失在教會之後.
我想告訴大家.
在早期新約時代的教會.

$^{721}$所有新的禮拜堂被建立.
一定不會這麼迅速的發展.
大家認同嗎?.
因為塑羅.
保羅曾經在《使徒行傳》有三次的旅行報道.
有沒有印象?.
他傳了福音,帶了很多人信主.
建立了很多教會.
所以如果保羅消失之後.
我相信新的禮拜堂被建立.
一定會有很大的阻滯.
這是第一個結果.
第二個結果是甚麼呢?.
第二個結果就是.
我相信因為塑羅消失在教會.
你和我手上這一本的新約傳書.
就沒有了十三封的書信.
你說是不是很大的損失呢?.
各位電視節目的觀眾.
如果塑羅真的從此消失在教會的時候.
其實帶給教會一個很大的遺憾.
大家認同嗎?.
不過上帝很奇妙.
各位電視節目的觀眾.
想都沒有想過.
當人到了盡頭的時候.
當人落在絕望的處境的時候.
想著塑羅從此消失人.
在教會的時候.
竟然上帝的作為很奇妙.
上帝的作為叫我們拍案叫絕.
為甚麼呢?.
原來在這樣的惡劣環境下.
上帝原來再次感動巴拿巴.
你記得嗎?.
剛才經文說.
原來巴拿巴又聽到上帝給他的感召.
那一刻巴拿巴心裡有很敏銳的觸覺.
他知道上帝要找他做些事.
因此聖經記載.

$^{761}$巴拿巴原是個好人.
被聖靈充滿大有信心.
有印象嗎?.
他冒著一個很大的危險.
冒著一個被母會可能會排擠的危險.
他也不理會那麼多.
他聽從聖靈的指引.
去大塑的地方.
將塑羅再一次找出來.
你有沒有留意到?.
當他再一次將塑羅找出來的時候.
頂智滿們.
這次巴拿巴聰明了.
他沒有帶塑羅回耶路撒冷教會.
你知不知道他帶了他去哪裡?.
他帶了塑羅去安提亞的教會.
聖經記載.
在那個地方.
巴拿巴和塑羅兩個人同心似縫.
一起彼此互相勉勵.
在安提亞的教會.
和弟兄姊妹互相鼓勵.
開拓很多的培訓.
很多的栽培.
令到安提亞的教會.
足足一年的時間.
就馬上.
教會就脫胎換骨.
聖經記載.
有一句話.
在最後那句.
大家剛才讀的時候有沒有留意到?.
最後那句說什麼?.
最後那句說.
門徒稱為基督徒.
是從哪裡開始?.
是從安提亞開始.
各位頂智滿們.
我想全世界所有的教會.
都很羨慕這個安提亞教會.

$^{801}$因為沒有一間教會.
竟然在未信主的人眼中.
是首先稱他們為.
名符其實跟隨基督的人.
這些是基督徒來的.
被未遂的人.
這樣去稱讚他們.
這個稱譽.
這個讚譽.
從哪裡開始?.
是從安提亞開始.
各位頂智滿們.
你從這裡看到.
原來當巴拿巴和保羅.
在安提亞教會.
同心配搭事奉.
來關懷教會.
來培訓弟兄姊妹.
鼓勵弟兄姊妹的時候.
教會就非常之興旺.
結果教會就能夠榮耀上帝.
各位弟兄姊妹.
在這裡我想問問大家.
有沒有人聽過.
李思敬博士的名字?.
聽過你舉舉手.
差不多全部人都聽過.
我想問大家多一個問題.
有沒有人聽過.
李思敬的爸爸.
叫做李飛吾牧師的名字?.
聽過你舉舉手.
好像一隻手都沒有.
頂智滿滿.
因為都是上個世紀的人.
我為什麼今天要在大家當中.
提一提這個家庭呢?.
我想告訴大家.
李飛吾牧師.
即是李思敬的爸爸.

$^{841}$李思敬的爸爸.
李飛吾牧師的身世.
是非常之可憐.
因為他是一個孤兒.
流落街頭.
沒有人理會他.
為什麼呢?.
因為當時候的中國.
兵荒馬亂.
但是很奇妙.
上帝就興起了一個宣教士.
他的名字就叫做何義思宣教士.
一個西教士.
他眼見當時候的中國.
這麼淒涼這麼慘.
這麼多小孩子.
沒有父母.
流落街頭的時候.
他就一個一個把他們撿起來.
回家來去養大他們.
撿著撿著的時候.
竟然就被他撿到一個.
李飛吾牧師回到家.
去撫養他.
去培訓他.
結果這個李飛吾牧師.
竟然在上個世紀.
大家都是宣道會的弟兄姊妹.
我想你們一定要知道這個歷史.
在上一個世紀.
宣道會的堂會.
為什麼發展得這麼迅速.
這麼厲害呢?.
比其他宗派的堂會都發展得這麼快呢?.
當時候有兩個人是非常之關鍵.
一個就是滕近輝牧師.
大家聽過他的名字.
另一位就是李飛吾牧師.
這兩位牧師走馬步.
在上個世紀的時候.

$^{881}$他們很同心.
大家意象一致.
大家就很想教會不斷地開新堂.
不斷地直堂出去.
不斷地開分堂.
所以在上一個世紀.
很多宣道會的教會就建立了出來.
弟兄姊妹們.
可能何義師宣教士.
都沒有想過他撿了這麼多孤兒回來.
既然有一個.
不止一個.
還有另一位牧師.
也被神大大使用.
既然撿到一個李飛吾牧師.
而他在上個世紀被上帝重用的時候.
我相信何教士自己都不知道.
正如巴拿巴當年都更加不知道.
他手上培訓的這個掃羅,這個保羅.
突然間是一個屬靈的巨人.
他完全不知道.
最後弟兄姊妹們.
今天的講道.
到結束的時候.
我想和你們說的是什麼呢?.
就是每一位家恩樓的弟兄姊妹.
你們都不要小看自己.
不要小看自己.
上帝用得著你.
上帝可以用你.
只要你願意謙卑忠心地被上帝使用的時候.
上帝一定用得著你.
你只要留意一下.
在我們的教會當中.
有哪些是初信的弟兄姊妹.
有哪些是沒有人關心的弟兄姊妹.
而你願意主動地去關懷他.
去關心.
去栽培.
去培訓他的時候.

$^{921}$弟兄姊妹們.
又或者你在職場上.
你留意一下.
有沒有哪些同事是初信的.
他好像沒有什麼人關心他的.
你嘗試一下去關心他.
和他聊聊天.
和他讀讀聖經.
一起祈禱.
一起分享.
各位弟兄姊妹們.
我想告訴大家.
可能你所栽培的一個人.
將來是一個屬靈的巨人也不一定.
為什麼呢?.
因為上帝的作為.
是非常令人讚嘆的.
是出乎人意料之外的.
只要你願意被上帝使用.
上帝一定可以讓你成為別人生命的師父.
最近全世界的人.
包括香港人.
都有一個熱潮.
這個熱潮就是.
大家都很喜歡用這個Chart GPT.
這個聊天的機械人.
在當中.
我想你們也在用.
你們也在用這個Chart GPT也不一定.
各位弟兄姊妹們.
我想告訴你們.
這個Chart GPT真的很厲害.
是高手.
沒得頂.
發展的驚人效果.
連馬斯克都要出聲.
停一停.
休息一休息.
為什麼呢?.
因為很怕第四第五代的GPT再出來的時候.

$^{961}$不知道會發生什麼改革.
這個Chart GPT真的高手.
很可怕.
很厲害.
有人問.
從今之後人工智能會不會代替我們人類呢?.
各位弟兄姊妹們.
今早我想將確實的答案告訴大家.
Chart GPT怎樣厲害都好.
但是生命師父永遠都不能夠被替代.
Amen.
我們同心祈禱.
親愛的天父.
求你使用我們眾異女們.
今天我們學校巴拿巴的榜樣.
在教會當中.
在職場裡面.
求上帝用得著我們.
成為別人的生命師父.
因為我們知道.
當我們願意謙卑忠心.
積極的來事奉你的時候.
你就要再一次建立我們所做的工作.
讓我們所栽培的.
所培訓的弟兄姊妹.
將來更加被上帝大大的使用.
願主祝福我們的生命.
祈禱奉耶穌名求.
Amen.
\newpage



\section{}
\label{sec:Am9T43vQMKQ}
\textbf{2023-04-06 周四晚查經班《被擄歸回歷史研讀(10)》 講員:曾裔貴牧師}
\newline
\newline
連結: \href{https://youtube.com/watch?v=Am9T43vQMKQ}{\texttt{https://youtube.com/watch?v=Am9T43vQMKQ}} ~~~~ 語音日期: 2023-04-07
\newline
\newline
\hyperref[sec:lYwkS21XVXg]{\small{< < < PREV SERMON < < <}}
~
\hyperref[sec:index_chronic]{\small{[返順時目]}}
~
\hyperref[sec:index_scriptual]{\small{[返順卷目]}}
~
\hyperref[sec:Pndq2oLIHnI]{\small{> > > NEXT SERMON > > >}}
\newline
\newline
$^{1}$弟兄姊妹,我們可以開始今晚的茶經班了.
因為現在用Mixer錄音功能,所以要按一下鍵才能回來.
大家留意就行了.
好像音質會比用外置錄音機好.
我們一起祈禱.
再一次多謝主.
我們真的有機會對歸回的歷史有研究.
很盼望神都讓我們透過以色列記和尼希米記.
這兩卷是很相近時代的作品.
還有兩位領袖他們帶來的.
在當時猶太人的復興.
都讓我們從其中得到一些重要的原則.
使我們能夠在自己的時代.
怎樣作主的僕人.
因應時代的需要.
作出向神的禱告和具體的實踐進行.
今天晚上我們仰望在神的手中.
我們盼望已經成為了錄音的製成品.
都可以去祝福不同地方的弟兄姊妹.
透過神的話語的宣講和言讀.
我們大家都一起得到那個建立.
我們這樣的祈禱.
奉耶穌基督的名.
阿們.
我要靜一靜音.
我們說過尼希米記第一章.
大概是主前的444或445.
與以斯拉的時代相差不遠.
可以說是第二批第二代的領袖.
或者是當時在波斯帝國.
亞達斯西王的年間.
神所興起的僕人.
以斯拉我們知道他是一個文士.
一個很敏捷的文士.
由通達神籍著摩西的律法書.
所頒佈的.
即是說他是一個祭司.
經常有機會去研讀摩西五經.
和都是一個合資格.
大家都公認的一個教法師.

$^{41}$教導律法的.
但是他是生在波斯.
他來到上去.
發現了一些內部的問題.
所以不只是硬件的重建.
而是神的選民的重建.
這是來了一層.
我們上一次說到以斯拉.
他某程度是強迫自己的國民.
要來遵守符合摩西的律法.
外邦的女子要休他.
甚至乎要很清楚.
有一個時間.
大家要來表態.
但是聖經就好像在這個.
比較辣的一個手段裡面.
似乎沒有什麼之後的跟進.
當然有些人因為已經登記了.
哪些是祭司的行列.
本來是可以事奉的.
他們都登記了.
但是到底具體有沒有執行.
聖經就沒有交代了.
一到就突然停在第十章就完結了.
好像留給讀者.
來到去有一個反思.
就完結了.
突然間進入了彌軒禮記.
在原文以往的希伯來文的原文.
兩卷書是一起的.
以色列的原著就是彌軒禮記.
彌軒禮記的頭六章.
因為第七章是另外一個故事.
頭六章都是彌軒禮個人的回憶錄.
所以你會看到這六章聖經.
有很多我啊我啊我祈禱啊.
我神施恩的手幫助我.
很多的.
他好像一個個人的自述.
好像自言自語.

$^{81}$那個過程.
重要的是他為什麼會起義.
要上耶路撒冷.
這個過程.
他是什麼身份.
他有什麼資源.
第二章就說那個資源.
他怎樣來籌劃.
他作為一個那個時代的領袖.
他怎樣做.
這個都是值得我們今天可以去想一想的.
等一下可能我會用另外一個例子.
給大家有一個對比.
大家都是領袖.
大家都是生在一個混亂的時代.
但是到底上主.
就是神.
在兩位的領袖身上.
有什麼帶領呢.
有什麼特殊的神蹟呢.
這些都是一些指標.
我們可以來到去.
從其中去想一想.
某程度上真的要對號入座.
我們要問就是.
我們現在生在什麼時代.
我們在這個時代.
神在做什麼工作.
我們就要問了.
我們就能夠拿捏到.
我們自己的那個位份.
和怎樣去運用神給我們的資源.
第二章我們回到尼希米記.
都是要讀經文的.
因為經文自己都承載了一個故事.
既然是一個個人的自述.
很多的字裡行間.
一定是他有很深的領袖.
這樣來說述.
很多的浮言.

$^{121}$特別這位領袖.
因為屬世.
他是一個酒精.
如果古代的酒精.
真的是皇帝信任.
不是普通的一個AO的政務官這麼簡單.
可能不一定是投考.
而是皇帝信任.
是自己最靠近飲食的一個官.
可能不只是遞酒這麼簡單.
而是要試那些飲食.
會不會有人下毒.
因為皇帝都會很容易被人行食.
在古代.
所以變成我們都要留意.
他這個身份.
他這麼強調.
自己這個屬世的.
在地方的政府.
當時的波斯帝國.
帝國裡面的皇身邊的一個酒精.
一個得到皇帝信任的大臣.
臣可以在這個時代.
在這樣的位置.
興起一個這樣的領袖.
而且整本書.
其實就是圍繞著他個人的回憶錄.
來講述.
所以又可以讓我們.
從他身上.
又可以看到一些原則.
他是怎樣來依靠神等等.
我們上一個星期講過.
他有一個很強的心.
就是他雖然在波斯.
尤其是在書山城.
即是說就算不是守都.
都是一個重要的行宮.
即是說他其實都是日理萬機.
很忙碌的.

$^{161}$而且他的工作.
經常都是對人要歡笑.
就要背人仇.
你不開心.
你不可以在皇帝面前不開心.
因為這些是宮廷的重要的公開場合.
你不可以自己經常苦口苦面.
會切換你.
甚至有生命危險.
如果你真的得罪了皇帝.
這些我們遲些在以斯帖記.
就會看到的.
這個沒底改.
就是得罪了一個哈曼.
哈曼見到他.
又不跪拜.
又不給他面子.
又不可能卑躬屈膝.
所以就會引發一些.
就算是同埋一起同朝為臣.
都會受到.
被人這樣誣衊.
被人甚至想剷除他.
和剷除他的猶大族.
這些是很冷酷和背後的大陰謀.
即是說叛君就如叛婦.
你在皇帝身邊自然有人會眼紅.
自然會有人想借機會來攻擊.
所以尼熙米一定是在這樣的位置.
他自己也有一個很大的運籌遺額.
他也是一個自己步步為營.
和也是一個很謹慎行事.
更加也知道官場的險惡.
這些等等的東西.
在他的內心其實已經有了.
但是既然現在是聖經.
他就表露無遺他的內心.
對神的聖殿.
對耶路撒冷.
神的聖城.

$^{201}$那個繫鸞.
那個重視.
這些可以反映出.
一個已經是第二第三代.
你不是第一代上來的人.
你不是看著這個聖殿建造的一個人.
你是已經看到這個聖殿.
有他的一點點的輝煌.
已經建好了.
重建好了.
而且是第一批回去回歸的人重建好了.
但是慢慢不知不覺又破落了.
那個破落是當地的外族人.
來踐踏.
來遊輪和嘲笑.
所以對神的憲制.
還有敬拜.
沒有了那個最好的時代.
還有一定的困難.
去到安息日去敬拜神.
等於你在一些回教的國家.
你說是宗教自由.
但是偶然間會有炸彈襲擊.
有些極端的回教徒.
他會放炸彈來你的教堂.
你都真的會有提心吊膽.
你是不是都堅持要回教會信住呢.
有些人可能很堅持.
但是也有可能更多的人.
都會擔心,懼怕,不敢.
或者不會太公開.
這樣來聚會.
這個一點都不奇怪.
所以我們就要了解一下.
黎希米他自己的內心.
所以感恩就是.
透過這本書頭六章.
我們有很多他的自述.
我們可以看到一個敬虔的領袖.
他怎樣來運用他的智慧.

$^{241}$他的膽量.
他的情感.
他的意志.
這樣來做好一件.
他那個時代需要做的事情.
所以某程度就是他的盡心盡性.
盡力的愛主.
他的神.
在他的人生.
他能夠成就一番的.
為神的事業.
第二章我們看.
亞特薩西王的二十年離散月.
距離上一次的基斯留月.
三至四個月.
他一直的祈禱.
一直的部署.
可能是心靈的部署.
可能是和那幾個弟兄.
來到去有默契.
有一起祈禱.
有一起的去籌劃.
但是沒有行動.
因為他在官場.
他得到的是天上的神.
和地上的老闆.
就是亞特薩西王.
兩者都要來到去有那個開綠燈.
他才能夠去做的.
這個工程也不是小的.
他現在是在很遠的.
波斯的舒山城.
如果以斯拉記告訴我們.
要走路都要五個月.
不知道尼希米的時代.
有沒有盜賊的危險.
或者有其他的.
州居的運輸的問題.
所以很多事情真的不是這麼簡單.
要籌劃的.

$^{281}$心靈的感動和具體的落實.
其實是有一個很大的距離.
當然最重要的是.
首先要有心靈的感動.
有了這個感動之後.
就開始認真的.
看神怎樣來開路.
現在最重要的路要開的.
就是皇帝的思路.
那個心靈.
和是不是要讓皇帝.
用什麼樣的方法去開口.
去爭取.
去拿大價.
無薪價.
或者是怎樣呢.
而且涉及的有很多的行政省.
現在在河東.
以約旦河.
在河東.
波斯帝國.
其實很遠.
來到河西.
河西的地方的行政官.
那些省長會不會接受呢.
因為始終帝國很大.
不是香港.
不是說你出九龍.
深水Po 區.
這麼小這麼小的地方.
是很遠的.
你想想那個地方的行政官.
會不會採納呢.
這些全部都要想的.
在經文裡面.
今天我們都會見到.
李軒蔚是怎樣去梳理.
這些具體的問題.
李三月.
在王面前擺酒.

$^{321}$因為她說她是王的酒精.
我拿起酒來奉給王.
我素來在王面前沒有愁容.
王對我說.
你既沒有病.
為什麼面帶愁容呢.
素來沒有愁容.
現在王今次.
發現了李軒蔚.
她面帶愁容.
可以帶來一個危機.
無端端你不開心.
是不是你.
如果王不是愛錫她.
就以為她是在那裡.
有一些工作上的態度.
但如果王愛錫她.
就可能揣摩.
是不是她有病有不舒服.
或者有心裡的疑難.
需要有人幫助她.
所以你看到.
在這本書裡.
亞達塞西王是一個很通達.
對李軒蔚是非常愛戴信任的王.
能夠在一個外邦.
拜假神.
雖然不是直接由波斯去征服猶大.
但始終整個猶大的國.
即以色列國.
都是亡了.
在這個行政省裡統治.
即是波斯王.
都是統治猶大.
現在會不會有些蠶神.
在背後指著背脊.
說猶大如果現在重建.
會不會有機會.
他們會來要獨立.
要推翻波斯的當地行政省長.

$^{361}$這些可能涉及第三第四章.
當地的外族人.
會借這些機會.
用這些理由.
在王面前打小報告.
所以李軒蔚.
其實他做這件事.
吃力不討好的.
他自己可能會招致.
政治的報復.
帶來可能水洗都不清.
因為始終離開王的勢力很遠.
所以這些.
李軒蔚應該都有思想.
有自己的掙扎.
然後他覺得.
有一樣東西更加重要.
就是神的殿.
那個榮耀.
比他自己面對的困難.
更加重要.
所以他就決心.
真的排除萬難.
來決心在王的面前講述.
所以現在有機會來了.
第二節.
王對我說.
你既沒有病.
為什麼面帶愁容呢.
這不是別的.
必是你心中愁煩.
於是我甚懼怕.
很正常.
王居然主動問你.
如果當時是一個公開的場合.
似乎是.
因為第六節說.
王後都在一起.
一個公開的場合.
不是說皇帝現在下班了.

$^{401}$在後宮吃飯.
只是自己人.
不是.
似乎是一個公開的場合.
這不是別的.
必是你心中愁煩.
於是我甚懼怕.
我對王說.
願王萬歲.
我列祖墳墓所在那城方涼.
城門被火焚燒.
我豈能面無愁容呢.
講那個原因.
就是他的鄉下耶路撒冷.
當然是聖殿的範圍.
他沒有提到聖殿.
他提到的是他的故鄉.
即是說他有機會都回過去.
可能是數目.
可能是祭祖.
可能是首節.
因為猶太人一年三次.
都要上耶路撒冷首節.
他那個時代.
是否那麼賣力.
殷勤地去進行.
就沒有任何的資料.
但是一個敬虔的猶太人.
就是在逾越節.
在五旬節.
在鑄坪節.
就要上耶路撒冷首節.
所以這裡告訴我們.
他就說出那個原因.
他怎能夠沒有愁容呢.
因為城門都被焚燒了.
那個地方荒涼了.
第四節.
是不是他出人意外的.
奇妙的工作呢.

$^{441}$他祈禱有沒有求神.
去感動去改變.
亞達斯西王的內心呢.
有沒有祈禱.
希望亞達斯西王.
主動去詢問他呢.
還是他已經寫了架子.
他要回去一段長時間.
去視察和重建.
那些破損的地方呢.
但是要多久的時間呢.
只是走上去都要五個月了.
你可以想像得到.
那個困難不是那麼簡單.
更加不是他一個人的人力.
就可以去造成的.
所以他有沒有想到.
王是最關鍵呢.
如果他支持他.
又給材料他.
又頒佈詔書准許他.
那就是說.
好像是由國家的層面.
來做神的工作.
那就很特別了.
你可以想像得到.
這些情況就是說.
真的很明顯的一個印證.
他可以這樣上去重建了.
所以第四節就出現了.
王問我說.
你要求什麼.
於是我默禱天上的神.
似乎不是在這一刻.
他才突然之間祈禱.
剛才說的就是.
基斯留月到彌撒月.
有三至四個月的時間.
如果這三至四個月.
他都是闊綽難離.

$^{481}$很感動啊.
聽完那幾個弟兄哈拿尼.
來到波斯的書山城.
說現在耶路撒冷的光景是怎樣.
馬上很感動.
感動完之後又哭又祈禱.
然後就忘記了.
因為每天工作那麼忙.
排到滿了.
然後慢慢又給其他的公務.
個人的一些生活.
去霸佔了思想.
慢慢又淡忘了.
一淡忘了.
那個情緒就慢慢會跌了.
相信你和我都是.
就算你看完很好笑.
很好看的電影.
慢慢慢慢都會淡忘的.
慢慢慢慢他都會那個波動.
不會有人波動三個月都不停的.
他要看醫生的.
沒可能的.
因為真的隨著.
你很高漲很高漲很高漲.
但是可能有些人就很快.
幾個小時就已經忘記了.
有些人就還有一點印象.
還會回味.
會有一些的思想.
但是那個思想始終.
如果你不持續的禱告.
那件事慢慢慢慢.
你好像覺得都沒什麼出路.
我都是一個人丁單薄.
還要是在外地那麼遠.
說一下什麼怎麼可能呢.
那你那個祈禱.
慢慢慢慢你就會淡化的.
人沒可能你是.

$^{521}$經常在這種高漲的情緒波動裡面的.
除非那個祈禱是.
由那個激動.
到慢慢是認真的.
已經進入那個意志裡面.
已經進入了那個決心裡面.
已經不是一種感情.
一下子的感動.
而是你已經開始.
神啊.
我是不是有份於這件事呢.
我過去的成長.
我的訓練.
我的籌劃的那個能力.
有可能是不是神你要我.
突然間知道這個這樣的情況.
這個現象.
這個光景.
還有更重要的就是.
為什麼在耶路撒冷這麼近的.
都沒有人有這個感動.
來去決心的去做.
甚至乎來去找.
當時的河西的地方官.
來去申請重建.
為什麼沒有呢.
那會不會真的.
透過這個現象.
尼熙米所見到的聽到的.
由這個現象慢慢一直祈禱.
祈禱就成為尼熙米的那個意象呢.
雖然我沒有用powerpoint.
去解釋這個現象.
一個phenomenon.
一個你觀察到的.
已經發生的事實.
但這個事實不是一下子.
三兩天的.
是持續了很多年的.
但是在耶路撒冷.

$^{561}$大家見怪不怪了.
大家都沒有能力.
可以去改變.
去左右那些當地的外邦人.
的嘲笑踐踏.
也都沒有那個能力.
訴諸於波斯王會不會開綠燈.
讓我們重建.
在耶路撒冷已經沒有人.
有那個心力.
有那個心智.
甚至有那個行動.
來去重建了.
那會不會現在這個現象.
尼熙米聽到幾個.
從猶大下來的弟兄.
其中有哈拿尼.
這樣來去描述.
耶路撒冷的破落.
他就很有感動.
哭了一晚.
然後悲哀.
幾天了.
悲哀了幾天.
幾天之後.
人始終是不可能.
不斷的每天都悲哀的.
你沒理由每天都.
好像想死那樣.
不行的.
你要做事.
你那個感情.
你多澎湃.
你都會自己冷靜下來的.
但是如果幾個月.
持續的這樣祈禱.
就不再是感情上.
那種感性.
而是已經轉化成為一個意志.
成為一個心智的祈禱等候.

$^{601}$就是說開始進入那個思想裡面.
我要怎樣籌劃.
我要怎樣準備我自己.
繼續在那個收工之後的祈禱.
就是他可能就開始.
和那幾個弟兄.
和他自己的一些同伴.
一些土伴.
就來認真去祈禱.
祈禱完之後.
用他僅有的時間.
他就來籌劃那個建造.
如果真的可以的.
他就會來籌劃那個建造.
所以在這裡.
他的脈絡天上的神.
不是在這個時候才祈禱.
如果在這個時候才祈禱.
他的感動.
只不過是一個.
剛剛突然之間.
而是他已經是三至四個月的時間.
他一直去為自己.
那個神的帶領去禱告.
他自己都準備.
可能一些計劃書.
來在他自己的工餘的時間.
來籌劃.
如果要重建聖殿.
哪一個部分.
當時也沒有什麼圖則.
更加不會有一些攝影.
拍攝影片.
高清4K.
沒有這些東西.
是靠哈拿利他們的描述.
所以對於黎希米來說.
他本身不是一個建築家.
更加不是一個耶路撒冷的領袖.
可能跟裡面的人.

$^{641}$是素未謀面.
他是一個普通的猶太人.
是在波斯帝國裡面的一個官員.
但是這裡告訴我們.
他這一下緊急的脈搏.
其實就是他過去三個多月.
持續的等候.
持續的尋求神的旨意.
持續的由一個他觀察到的現象.
變成一個vision.
一個的異象.
異象就是揮之不去.
在他內心要為神做一些事.
揮之不去在他的內心.
要來到他在他的有生之年.
要為神的聖功.
他要擺上他自己.
所以很多的傳道者.
他做傳道.
一定是要經過這個過程的.
他觀察到一個現象.
由這個現象.
慢慢在他內心一路去醞釀.
去等候去裝備去禱告.
去尋問問一些長輩.
各方面甚至乎周邊的人.
都觀察到他有這個.
要為神來工作的那個心智.
這樣這樣等待.
很明顯見到神一路開通達的路.
那個綠燈.
讓那個等候的人.
是由那個現象轉變成一個異象.
這個是相當重要的.
今日我們很多的現象.
剛剛昨晚和同學.
神學院同學吃飯聊天.
有一位老師.
他說這麼多人去英國去移民.
香港的教會.

$^{681}$有沒有人現在去學俄文.
俄羅斯如果停戰.
烏克蘭停戰.
就需要大量的傳道者.
會講俄文的.
就要進去幫手.
有沒有一些香港的教會.
有這個意識.
現在就學定俄文呢.
內地很多人會俄文的.
因為那個關係是很密切.
有專門俄文俄語的學校大學.
但是內地的人不是個個信主.
更加不是傳道人.
香港有多少人有這個意識.
主動我們來學俄文.
去學習俄語.
將來我們進去傳道呢.
那個老師這樣講.
都勾起一些反思.
我們是絕對沒有批判的.
每個人.
我等一下會講.
每一個傳道人.
每一個的領袖.
有自己的大力猜險.
不過為什麼.
譬如我去荷蘭.
我去盧森堡.
去看宣教士.
這麼大的地方.
這麼大的國家.
宣教士差不多一個都沒有.
文姑娘.
一個人走幾個地方.
隨隨便便.
要開一個多小時車.
為什麼香港的傳道人.
我又是沒有批判的.
為什麼全部的感動.

$^{721}$都要去英國呢.
他們就笑了.
因為荷蘭文.
沒有人會講.
沒有人會聽.
完全是另外一個語言.
完全你學都有困難.
是很難的語言.
盧森堡是講法文的.
那就沒有了.
講德文的.
原來如果是這樣.
這麼低層次的原因.
那就.
難怪這麼大感動.
就是去英國.
因為我們香港是英文的世界.
至少不含淡都是英文.
聽得懂.
寫得到.
只是講得不流利.
這樣的.
現在我們回來.
不是要批判什麼.
而是要告訴我們.
由現象變成異象.
是要醞釀的.
是要準備的.
是要等候的.
好像那個機會.
是留給那些有準備的人.
這是世界的價值觀.
都這樣認定的.
你準備好.
機會來.
你就可以來得到.
我們現在不是說什麼買樓.
現在樓價低.
不是這些機會.
給那些有錢.

$^{761}$有現金的人.
Cash is king.
不是說這些.
李希米是準備了三個多月.
語言當然對他來說沒問題.
他是猶太人.
他懂希伯來話.
亞蘭話.
但問題是.
他要等待神.
感動亞達薩西.
但這個感動不是大了.
他是準備自己.
由那個現象.
變成一個異象.
所以弟兄姊妹.
我們有沒有裝備自己.
去讀好聖經.
我們來到祈禱上.
我們自己去殷勤學習.
我們去傳福音.
你不介意跟一個陌生人講福音.
神是否用你這方面的恩賜等等.
我們中產的教會.
我們是否慢慢關上門.
我們只是享受.
我們的崇拜.
頌唱.
這類真好.
就已經滿足呢.
還是我們的聚集.
都是為了我們被差遣.
去幫助有需要的人呢.
弟兄準備了三個多月.
所以現在這個脈絡.
是他三個多月的.
正式的蒙英韻.
所以他現在突然間.
閉上眼的脈絡.
不是因為他懼怕.

$^{801}$不敢說話.
而是他現在.
突然之間.
他感動神給他的機會.
然後你看他怎樣去解釋.
第五節.
我對王說.
僕人若在王眼前蒙恩.
王若喜歡.
求王差遣我往猶大.
到我列祖墳墓所在的那城.
即耶路撒冷.
我好重新建造.
很大的膽量和胃口.
要求王差遣他.
即是出於王的召命.
這件事提升成為一個國策.
由國家動用資源.
即是公務回到耶路撒冷.
不是他個人拿大架上去.
回耶路撒冷探親.
不是的.
是提升到這樣的層次.
才是神.
可能過去三個多月.
祂的祈禱.
這樣來求神.
這樣是最好的.
王看到.
當年.
其實是很多年前.
古列王的時候.
國策就是准許.
不同被吞併的國家.
不可以自治.
但可以有自己宗教的禮儀.
令到大家都服從帝國.
但有自己國家的崇拜.
敬拜神.
只要你們不作反就行了.

$^{841}$這是一個很大的國策轉變.
和亞述帝國.
要將每一個國家.
每一個神都禁止.
這就帶來很大的反對.
你想想.
現在我們內地.
雖然是無神思想.
但也准許基督教信仰自由.
也准許佛教不要過份.
可以裝香.
可以來拜神.
也准許.
可想而知.
他也知道.
如果完全禁止.
全中國十四億人.
沒有一個人可以信有神.
這樣就不行了.
這樣就會帶來.
衝擊更加大.
所以波斯帝國就有這個得證.
知道禁止他們拜神.
就帶來很嚴重的反對.
所以他們的國策.
是容許你們有限度.
拜自己的神.
有你們的禮儀.
有你們的祭祀等等.
所以現在.
黎氏不覺得.
世俗的東西.
用在神的裡面.
好像不太適合聖經.
好像很熟悉.
不熟悉靈.
但黎氏的內心裡面.
他不覺得是這樣想的.
他覺得外邦.
我用一個例子大家會明白.

$^{881}$信者一聽就會明白.
教會要為新的服償.
需要很多的經費.
來運作.
來傳福音.
不如這樣.
更加簡單.
馬會每年都會找一些機構.
他們認為值得支持的.
你有好的計劃書.
清楚的.
馬會就會撥款.
為什麼好像這麼少.
香港的福音機構.
福音派的教會.
會申請這些馬會的撥款.
是慈善的.
是政府規定馬會要做的.
為什麼不多申請呢.
大家一聽就會明白.
這些錢是靠主要的收入來源.
是那些賭徒.
可能傾家蕩產.
輸光了馬.
現在還搞足球.
這些錢的來源本身.
不是這麼道德的.
雖然香港說賽馬是合法化.
但是賭博的問題.
其實都是嚴重的.
所以教會.
大家幾乎肯定知道.
不敢這樣做.
想都不敢這樣想.
拿這些慈善.
每年都穩定的.
來贊助.
贊助什麼.
這裡就不要叫黃振輝社區中心.
政府給一個地.

$^{921}$我們可以蓋房子的.
我們就透過馬會贊助.
蓋新一座大樓.
當然是亂說的.
由馬會的哪一個主席命名.
大家都知道絕對不可以.
這些就知道那些是熟世的.
那些真的不是神心意的金錢.
但是很奇妙的就是.
尼希米基大量運用.
波斯帝國的資源去重建.
甚至還要求皇.
來頒下一個法令.
詔書.
這次上來是皇帝准許的.
所以地方官是一定要配合的.
這樣的好像中央的一個專員.
現在上去.
不過蓋的不是波斯帝國裡面的建築.
是猶大的耶路撒冷那個聖殿.
和周邊的圍牆.
重新的來去建好.
和安好那些門.
因為有一個很大的.
有不同的地方的出入的門.
所以那些門已經破損了.
就是說外邦人隨便都可以進來.
就是說武演雞籠.
安息日.
本來只是猶太人可以來安靜.
敬拜神.
外邦人又走進來買賣.
做生意.
說你安息日我們要做生意.
就變成了令到自己的猶太人.
一來很尷尬.
二來就破了摩西的那個律法.
對那個時代的情況.
就帶來某程度是很褻瀆.
所以李希米和他們那一代的人.

$^{961}$就看到這件事的問題所在.
但是現在當他祈禱的時候.
他來求王去差遣他.
似乎這個不是一時之間的靈感.
一時之間的急才.
而是他經過深思熟慮之後.
他如果真的有神的帶領.
他持續的看到的現象.
這個現象出現了很久了.
但是李希米不是普通的猶太弟兄.
這個現象在他的內心.
變成了一個異象.
這個時候他就知道神要用他.
相信一個弟兄姊妹.
如果你能夠認定神在哪一個地方用你.
在這方面你已經清楚自己人生的路向.
你為神所用是非常榮幸.
因為做神的聖公.
不是那麼多人稀罕.
我儘管用這麼重的話.
稀罕來到教會事奉.
那就錯失了能夠為神事奉的一份榮幸.
一份恩典.
然後裝備自己來被神用.
那是另一個層次.
感覺自己有一個差遣.
由一個現象變成一個異象.
而這個異象他自己繼續的去裝備.
繼續的去磨練自己.
準備自己.
希望神的時間到.
他可以做了.
三個多月祈禱.
現在的默禱不是臨時的抱佛腳.
不是一個突然之間的感動.
突然之間的急才.
想到可以叫王給我一些材料.
不是突然之間經過幾個月的祈禱.
沉思商量和籌劃.
而來作出的這個祈禱和決定.

$^{1001}$所以這個默禱是讓他的心神.
當下很平靜地說他的計劃給王知道.
所以這本書給我們看到一個很沉實的領袖.
似乎看不到有什麼神跡.
在整本書.
看到一個這樣的領袖.
是很恭敬地.
稍後他要回去述職.
回去皇帝那裡述職.
過了一段時間.
回來之後發現了一個問題.
這個我遲些再說.
如果你有興趣.
你回去自己看一下.
你就會知道.
他要動用他自己個人的.
他自己的俸祿.
來解決當下的危機.
你就可以知道.
他是出心出力的一個領袖.
在這裡的時候.
最重要的是有王的准許.
所以現在經過這麼多個月的祈禱.
現在王主動地來邀請他.
讓他看到.
讓他知道.
神已經感動亞達舍西王.
這樣就很不簡單.
所以神的作為.
有時真是.
你可以說出神入化.
不知道哪個時間.
和不知道怎麼做的.
尼希米都不可以限制神的工作.
現在神可以感動一個.
外邦帝國的一個王.
亞達舍西.
來准許自己的愛的神.
大神尼希米.
可以回去重建.

$^{1041}$我們看一下後果.
第六節.
那時皇后坐在王的旁邊.
特別強調皇后在這裡.
似乎是一個宮廷的飲宴.
不是談一些很重大的公務.
或者是一些很危急.
很機密的一些公務.
好像現在馬克龍.
法國總統訪問北京.
開始的時候給記者拍下相.
雙方的夫人都會一起拍下相.
然後談內部機密的事.
那些夫人就不可以在這裡.
就要離開.
記者一定要離開.
就談一些可能大家有什麼協議.
大家有什麼讓步.
你要幫忙勸俄羅斯.
你又要幫忙勸烏克蘭.
大家這樣一人走一步.
有這些秘密的協議是重要的.
一切的人都要離開.
在這個時候.
夫人都坐在這裡.
都很開心的時間.
所以皇后就很感受到.
尼熙米的仇容.
而不是掃興.
而是反而是怎樣幫助自己的臣子.
解決他的仇容.
解決他內心的仇苦.
即是皇后現在很感受到.
尼熙米這樣的策略.
甚至發現他能夠這樣做.
我相信她未必會想到.
我可以事奉神.
我可以事奉你們猶太人的神.
我相信她不是這樣想.
完全是因為尼熙米.

$^{1081}$完全是因為看到尼熙米過去的忠誠.
過去辦事的殷勤.
辦事的準確.
沒有甩鍋.
沒有吞噗.
不會隨便做一些事.
一直都是這麼勤奮.
做的事都是這麼準確.
所以這次一定是有些事不尋常的重大事.
所以當問完清楚之後.
皇都來準他.
第六節的下半節.
皇就問我.
不斷用他的第一身來說.
一個自轉式的自述.
尼熙米好像寫日記一樣.
皇問我說.
你去要多少日子.
何時回來.
我就定了日期.
所以不是突然之間.
是籌劃了幾個月.
於是皇歡喜.
差遣我去.
強調是差遣的.
差遣即是由皇.
由波斯帝國.
由國家.
現在好像委派尼熙米.
作為一個專員.
上去視察.
現在河西的光景.
這個恆省裡面的聖殿.
出了什麼問題.
是我的一個相信.
信任的一個官.
回家鄉.
回故鄉.
去視察.
這樣的一個情況.

$^{1121}$第七節.
我又對皇說.
皇如果喜歡.
求皇賜我詔書.
所以不只是口頭的差遣.
是要有詔書.
因為去到另外一個恆省.
這麼遠的河西.
那些地方官信不信.
有沒有那個掣肘.
不拜尼熙米的那個上來.
都可能會的.
但現在有詔書.
就有白紙黑字.
就沒得抵賴了.
河西的總督就要配合.
國家的那個政策.
即是尼熙米玩到一個地步.
就是去到最高層次的幫助.
完全沒有介意.
這些是熟悉的材料.
熟悉的金錢.
完全不會說這些不屬靈.
他沒有這樣分不屬靈.
但是.
他的內心深處.
是一直是為事奉神.
不是要來為自己.
去建功立業.
不是要在皇的面前.
來討好皇.
動機完全是純全的.
是他內心一個很單純的動機.
為神.
為神的殿現在的荒涼.
他悲哀幾日.
他祈禱天上的神.
他認罪.
為自己為同胞認罪.
和他去支取摩西律法書.

$^{1161}$我們上個禮拜說過的28張申命記.
只要有在外邦.
被神刑罰的子民.
謙卑認罪.
回到神的應許.
神就會赦免.
和會有那個祝福.
那個Blessing又會回來.
所以現在黎希敏.
不只是祈禱完就算了.
而是他要得到神這個祝福.
就是用他的人生.
和其他弟兄姊妹.
即猶太人.
在回去之後重建.
即是說他內心的這個計劃是很大的.
當然我們也說過.
可能猶太人他們想的就是.
為著彌賽亞將來的降臨.
他們去準備.
因為如果他們一直去計算.
那個時間由但以理書第九章.
由耶路撒冷重建.
有42.7.
然後再有1.7.
即是說.
他們可能這樣.
如果按字面去計算.
那個情況.
很可能不是那麼久.
就會有彌賽亞降臨.
當然剛才我們也說了.
主前的445年.
黎希敏的時代.
即是說也要400多年.
他們未必知道.
要那麼遠的時間.
但是他們也清楚.
彌賽亞是會降臨的.
所以這個重建耶路撒冷.

$^{1201}$是有這個重大的.
屬靈的原因.
不過神現在就在外邦.
在君王面前.
得到信任的一個酒精.
來在他的內心.
有一個那麼強的感動.
揮之不去的感動.
這個一定是異象.
揮得去的.
就是一時的衝動.
每個人都有的.
你看完某些難民.
好像上星期.
詠儀姐妹分享她在香港.
赤鷗洪磡區.
那些非洲人的需要.
我們大家都很有感動.
但是這個感動如果要持續.
要不斷在你的內心.
成為一個很重大的異象.
就需要持續關注.
和有具體的行動等等.
否則那個感動.
隨著時間生活.
就慢慢淡忘淡化.
就不會再有了.
但是現在尼希米.
她開始已經不再是.
內心的感動的波動.
而是很具體.
我怎樣去籌劃.
現在她甚至訴諸.
王啊你頒詔書.
這是第七節的下半節.
求王賜我詔書.
通知大河西的省長.
准我經過.
直到猶大.
又賜詔書.

$^{1241}$通知管理黃源林的亞薩.
就是那些木材的.
使得他給我木料.
作熟電營樓之門的橫樑和城牆.
與我自己房屋使用的.
因為他要不是短暫的時間.
你去一個工地.
首先見到的就是那些工程師.
那些的職師.
先有一些貨櫃.
有冷氣.
又有臨時洗手間.
因為那些工程師.
那些畫跡的.
他們要先進駐.
要來畫圖跡.
接著那些工人.
就搭了一個簡單的棚.
摘鐵或者機件.
打樁等等.
需要有他自己的房屋.
即是說尼欽米.
不是上去短暫的.
三天的一個短孫.
是一段時間的進駐.
來建造聖殿.
所以他自己都知道.
是一個大工程.
居然間這麼大的工程.
皇帝准許.
現在甚至乎他要求的是詔書.
所以你會看到.
一而再.
不是一下子的.
在這個基斯留越的感動.
而是這個感動.
已經在尼欽米的內心.
轉化成為意象.
成為他要為神.
以他的人生.

$^{1281}$為聖殿.
為將來的彌賽亞準備.
所以他能夠有.
他這個名字.
留下在聖經.
讓我們看到.
他為彌賽亞的降臨.
那個準備.
準備聖殿的重建.
聖城的重建.
那個神的祝福.
使他成為.
名留在聖經裡面的.
一個聖經人物.
是我們大家都需要學習.
都需要來去.
欣賞和學悟.
我們自己有生之年.
都有我們自己的意象.
我們要有我們自己的意象.
否則你的人生.
可能會很庸碌.
庸碌不是說你要做大事.
做小事.
而是那個庸碌是你沒有意象.
你根本就從來沒有感動.
其實是很嚴重的.
因為就算你做一些很簡單的事.
你都需要有那個感動.
有那個意象.
為神去做.
你會準備自己.
譬如不可以隨便亂說話.
不可以說是說非.
裝備自己成為一個可以被神用.
而能夠幫助到教會.
幫助到人的人.
這個就你的時間.
你的各方面都不是再好像.
你自己喜歡怎樣做就怎樣做.

$^{1321}$的一個人.
其實你這個過程.
本身就已經是一種.
屬靈之旅.
教會.
其實不怕人多人少.
教會最怕的是沒有意象的群眾.
沒有了意象的群眾.
喜歡就回教會.
不喜歡就一頭半個月都不回.
走過教會都完全沒有興趣.
因為已經沒有了意象.
這些是嚴重的.
如果甚至感受就更加是.
社區教會.
你沒理由為了不回這個教會.
你搬走的.
那就更加凸顯那個沒有意象.
走過.
就好像.
當這個是一個普通的建築物.
不是你的屬靈的家.
不是令到你的靈明.
令到你與神的相交.
越來越甘甜密切.
那就嚴重到更加凸顯了.
因為你已經沒有了意象.
沒有了內心的感動.
如果一個教會有很多弟兄姊妹.
但卻是沒有意象.
沒有熱心.
沒有感動.
那個教會就有很大問題.
很嚴重的一個危機.
不單止是年齡老化的危機.
而是內心沒有了意象.
所以我們很盼望.
真的弟兄姊妹能夠在.
我們的有生之年.
每個人都有.

$^{1361}$每個人的意象都不同.
大小.
本來在羅馬書的十二章.
講完了.
他用神的慈悲去勸我們.
將我們的身體.
當作活祭獻給神之後.
第三節就立刻講.
憑著他所得到的恩典.
勸我們看自己.
要看得合乎中度.
不要過激.
不要過冷.
我們要合乎中度.
神給我們的人生是什麼.
我們就接受.
神給你的人生.
很可能是一個很小的部分.
去中心.
你就很中心地.
在那個部分去中心.
那你就合乎中度了.
如果不是的話.
我要做使徒.
我要做才知道.
我要做什麼.
因為下面就立刻講.
其實一個身體.
有不同的器官的例子.
告訴我們.
每一個器官.
都有它的用處.
去建立整個人.
不只是手.
不只是腳.
是每個不同的身體.
都有它自己各個器官的功用.
很簡單的例子.
但是實際套用在教會.
我們就要留心了.

$^{1401}$神給我們有不同的帶領.
有不同的背景.
有不同的學識.
有不同的家庭.
但是每個人都應該.
看自己合乎中度.
和看自己是一個已經.
我奉獻了給神的活祭.
神怎樣用我.
是神的主權.
重要的是我要維持.
我內心裡面的感動.
我見到的現象.
要轉化成為意象.
這個是重要的一個學習了.
所以這裡告訴我們.
黎熙敏就是這樣.
你再次我詔書.
通知管理黃元林的阿薩.
他給我有木料.
我們可以建造聖殿的營樓.
和那些橫樑城牆.
和我自己用的房屋.
黃就允准我.
為什麼呢.
因為黎熙敏很信任得到.
因為黎熙敏很聰明的開會.
他的會議記錄寫得很好.
因為黎熙敏得到黃的歡心.
因為黎熙敏很好.
因為我神施恩的手幫助我.
完全將榮耀歸給神.
但是他做了三個多月功夫.
籌劃 禱告 準備.
一路等待.
居然現在不是他想到的契機.
就是黃主動找他 問他.
現在整件事完全水到渠成.
因我神施恩的手幫助我.
到底是51\%是神施恩的手.

$^{1441}$49\%是他的籌劃.
是他的口才.
說服黃各方面.
是100\%神施恩的手幫助他.
但是又是100\%他的祈禱.
他的等候.
他將他見到的那個方良.
轉化成為他自己的意象.
繼續在他內心來禱告.
等候神的時間.
100\%是他的準備.
是他的籌劃.
不是200\%.
都是100\%.
那你怎樣加呢.
100加100就是200.
即是完全是神管理著萬事.
我們見不到的那個屬靈世界.
神怎樣管理著.
這些不關我們事.
因為做物主有他的主權.
他不讓我們見到那個.
他怎樣在黃的心裡面來搞動.
這些不關我們事.
但是關我們事的.
就是黎希米.
他要在這個感動裡面.
在這個內心的觸動裡面.
他要轉化沉澱.
禱告和籌劃.
和有一個決心.
未來我會很艱難的.
我做得這個功績.
為神重建聖殿.
會有人眼紅.
會有人嫉妒.
會有人攻擊.
會招致批評.
會招致不必要的論斷.
甚至乎很難聽的說話.

$^{1481}$都會在背後講的.
陸續出現的.
所以你不要擔心.
一定會有的.
你做特首.
都大把人來插你.
任何一個政策.
都是有不同的角度去看的.
一定有對有不對的.
問題就是.
你只能夠吸取那個教訓.
不斷去完善那些東西.
沒有人一推出一個政策.
是百分之一百對的.
所以.
過去那五年我已經.
無數無數這樣.
有時是有些人.
很善意告訴我.
有人這樣說你.
是嗎?我完全不是這樣的.
有時很笑笑就算了.
如果你每次都上心.
你十樓跳十次都跳不完.
那何苦呢?.
我是為神做的.
那你就說.
如果我這樣表達.
這樣報告.
這樣可能會好一點.
原來又吸取到教訓.
我又是一個願意改變的人.
還有根本不會.
上心在那個心裡不開心.
因為.
你發現太多了.
你開心,你都不開心.
你真的經常都不開心.
你不用做正經事.
你每天坐在那裡.

$^{1521}$哭就行了.
那些就不是異象.
是你自己在那裡感懷身世.
其實很難.
你繼續事奉.
你的內心都已經有很多苦水.
一個領袖就應該.
明白了很多東西.
所以現在.
這麼順利.
要什麼有什麼.
皇帝又給詔書.
先不要這麼快開心.
大把困難陸續來.
甚至要陷害黎希米.
那個外邦人.
洗還手.
叫一個祭司.
叫黎希米去聖殿開會.
黎希米多聰明.
他的內心.
我如果進了聖殿.
我就大件事了.
因為那個地方我們是領袖.
我們是不可以進去.
他不會權傾一時.
去到自己突然間.
自大到以為自己.
就是那個絕對的領袖.
連祭司耶路撒冷的祭司.
都要聽他講.
原來是被人陷害.
所以原來他不斷.
保持自己那種謙卑.
我們一路讀下去你就知道.
很多很多的陷阱.
所以.
各位做公職.
其實不要做.
做傳道不要做.

$^{1561}$唐主任現在在香港教會.
就是個個都不做.
真的不是我說的.
是我們在英國.
都是這樣說的.
如果是為自己做.
有什麼好做.
回來聚會都很開心.
聽一下道 唱一下聖詩.
跟頂尖的聖母聊天 跟著祈禱.
都很開心 為什麼要做一些.
好像吃力不討好的說話.
的工作.
又有這麼多人會閒言閒語.
但是你要記得.
你要為神 如果是神安排的.
那你就要.
承擔那個責任.
不過你要一定要.
心理準備.
所以讀尼希米記.
是一個可以真的不斷.
潑一個事奉者的冷水.
的一個裝備.
潑完.
你都繼續.
有團火都沒有熄.
怎麼可能不是神.
揀選你 使用你.
人不是傻的.
潑完你都.
繼續有團火.
那還是你為自己 沒可能的.
所以.
求主幫助我們能夠常常.
都有我們自己被人.
兜頭林的冷水.
我們不是咒詛 是真的.
潑一次.
你都不熄 潑你.

$^{1601}$繼續有這個負擔.
你就很清楚.
教會有一個很.
重要 很寶貴的.
資產 就是這個地方.
有很有心的弟兄姊妹.
是很忠誠地.
為神的國度.
來裝備自己.
那教會就有.
希望了 更希望.
就是下一代的年青人.
有這些的.
負擔.
潑什麼冷水 都不會.
淋濕他.
他都是繼續這樣去事奉.
第九節.
王派了軍長.
和馬兵護送我.
我到了.
河西的省長那裡.
將王的詔書.
交給他們 雖然沒有講到.
多長的時間.
似乎重點不再是時間.
重點是怎麼能夠.
有效地去完成.
這個使命.
現在他一而再說.
這個詔書是.
王頒發給他.
在最高的.
政府的.
層級.
得到王的.
那個預準來做.
這件事.
真的不是一件小的.
簡單的工程.

$^{1641}$而是一個相當大的工程.
居然能夠得到外邦.
軍王的允許.
原因似乎就是.
尼希米.
他這個身份.
能夠周旋在王.
的身邊 以致.
得到信任 而能夠.
可以來動用.
國家的資源 去為.
神的殿去重建.
不過留意 第十節.
和倫人參巴拉.
繼續會在這十幾章.
的聖經出現.
這些外邦人.
當地的土著.
和倫人參巴拉.
並圍爐的阿滿人.
多比亞 聽見有人.
來為以色列人.
求好處 就甚擾路.
就是說.
作者都不等我們去猜.
立即告訴我們.
他一上到耶路撒冷.
那些外邦人.
當地的可能是.
土豪領袖.
意識到這個人上來.
不是那麼簡單.
因為他本身已經帶著官職.
帶著王的詔書.
還有馬兵護送.
就是說.
高規格的一個重建.
可想而知.
你可想得到.
當地人.

$^{1681}$看這件事.
不只是熱心猶太人.
希望來重建那麼簡單.
現在動用到王的.
國家的資源.
所以他們.
開始只是憤怒.
但是他們一定.
在明處.
攻擊不到.
尼欽米 就一定會用.
暗處 所以.
明和暗.
作為一個領袖都要小心.
第十一節.
我到了耶路撒冷.
在那裡住了三日.
我夜間起來.
我一個人也同起來.
但神使我心裡.
為耶路撒冷.
作什麼事 我並沒有告訴人.
所以原來.
我們現在看這一節.
我們就知道原來神.
在尼欽米的心裡.
幾個月一直.
給他這個負擔.
這個負擔越來越重.
重到他.
深深感受到自己.
有那個責任.
他要為神的.
殿來做一些事.
某程度上就是神.
將他在這個位置.
可以被王.
信任 可以.
動用到國家的資源.
所以神的事情.

$^{1721}$是我們好像.
很難去揣摩的.
我們千萬不要.
在我們自己的水平.
去評斷其他人.
否則就真的.
會變成論斷.
因為神怎樣用他.
是我們大家都不知道的.
所以我們要有一個大一點的胸襟.
不要以為.
你用國家 這些外邦人.
這些是不敬拜.
真神的波斯人.
你都用他的東西.
神可以用得著.
那些資源.
資源本身是中性的.
問題就是你的動機是怎樣.
所以這裡告訴我們.
神一直在他內心.
所作的事.
他沒有告訴人.
當然就是這幾天.
這三天在耶路撒冷.
剛剛初到貴境.
不是一個普通的.
回去好像馬英九那樣.
回到鄉下湖南.
祭祖.
而是這次真的是一個.
高格調 是有詔書.
重建的一個.
工程團隊.
李希米就是那個領袖.
但是更重要.
就是在他的自傳.
的自述裡面.
他是清楚神要他做的.
兩者完全沒有衝突.

$^{1761}$和很和諧.
這個真是我們大開眼界.
一個人可以.
被神這樣用.
是他內心的感動.
但是用他的影響力.
影響著皇帝.
動用國家的資源.
去重建的是神的殿.
你就可以想像得到.
這本書是多麼有趣.
現在這個就是.
表面在明處.
人見得到的.
安排.
但是.
在他建造的.
這52天裡面.
由他開始建造.
到完成52天.
但是.
真的帶著眼淚.
帶著驚恐.
和謹慎.
因為敵人無處不在.
敵人可以是這一班.
這個的.
參巴拉 多比亞.
基善.
或者是其他的人.
但是有些是.
猶太人的祭司.
已經是.
變節了.
已經那個靈性.
是破壞了.
是和這些的外邦人.
裡應外合的.
來到去想.
來到去.

$^{1801}$陷害.
黎希米的.
不是見到有人上來重建.
就很開心 很興奮.
因為神用這個.
弟兄在外邦.
有影響力.
來幫助我們的聖殿.
他甚至乎是.
要陷害黎希米的.
所以要小心.
事奉是要小心.
不要得意忘形.
更加不要.
以為自己.
真是不可一世的.
手頭上有詔書.
是皇給我的.
詔書.
我有這個上方寶劍.
要小心.
不是那麼簡單.
做神的工作.
魔鬼是會阻撓.
魔鬼是會攻擊.
好像上星期說的三唐道.
那個惡者.
是會害.
那些人.
所以我們要怎樣.
去將他帶進.
救恩的真理裡面.
三日的.
籌劃.
第.
十二節.
我夜間起來.
有幾個人和我一起.
我並沒有告訴人.
除了我騎的牲口以外.

$^{1841}$也沒有別的牲口在我那裡.
當夜.
我出了谷門.
往野狗井去.
到了糞廠門.
所謂糞廠門.
就是倒垃圾的地方.
猶太人認為那些地方.
是污穢的.
那個叫恩亂子谷.
在新約耶穌的時代.
就是一些.
因為他們是憲制的.
那些倒灰那些垃圾.
等等就是一個.
恩亂子谷.
這個叫糞廠門.
就是這樣.
去察看.
耶路撒冷的城牆.
似乎是夜間.
沒有電筒.
沒有照明.
不像香港那樣到處都亮燈.
地方大.
外國也是.
香港那麼特殊.
城市太小.
基本上你去什麼街都有街燈.
不是到處都是這樣的.
所以其實是完全沒有街燈.
你可以想像得到.
他晚上.
可能用一些很微弱的蠟燭.
他們所說的油燈.
來去視察.
就是說.
他為什麼要這樣做呢.
因為晚上.
可能真的要暗中.

$^{1881}$不要那麼張揚.
其實他有步兵.
有很多官員.
工程的人.
來陪著他.
但是他好像很低調地.
一個人先來做.
有兩方面的好處.
自己先做了功課.
胸有成竹.
哪些地方.
真的破損得嚴重.
他自己可能畫了一個簡單的圖.
沒有iPad.
不會有手寫筆.
不會立刻畫圖.
可能很簡單的.
有一個概念.
方便他第二天.
去召見那些猶太領袖的時候.
他不是一個.
好像.
擺一下官威上來.
什麼都不知道.
這樣的一個.
好像.
大巫師一樣的一個.
的人上來.
人家耍一下你.
根本就沒什麼可以做.
而是他很認真的自己.
做了功課去籌劃.
夜間有另一個情況.
就是正正是多比亞那幫人.
參巴拉這一幫人.
基善這一幫人.
他們的憤怒和攻擊.
就低調地.
自己.
做好自己的內部的安排.

$^{1921}$因為重要的.
不只是用外力.
重要的是要動用.
自己在猶太耶路撒冷的弟兄姊妹.
來去重建.
這樣才會得到.
那個真正的.
當地的猶太人的.
信服.
所以也都是這樣說.
一個傳道人剛剛.
新入職.
什麼都看不順眼.
什麼都要改.
這個時候.
你只會帶來很大的誤會.
可能你看到的是一個.
很真實的現象.
但是難道周圍的人.
看不到?可能是的.
新的傳道人就會是這樣的.
我當年也是這樣.
好像周圍的人看不到.
是你最明察秋毫.
最具雪亮的眼睛.
看到.
你就很想什麼都改.
但問題就不是這麼簡單.
你要讓人看到.
這個現象其實是可以改變的.
這個改變是可以.
大家一起去動員的.
這樣就真的事半功倍了.
還有重要的.
不是改變那個現象.
是改變那些人的.
內心.
他們自己重新火熱.
去依靠神.
這個才是真正的改變.

$^{1961}$那個的.
治本的方法.
改變那個現象.
不是你將奉獻箱放在哪個位置.
而是改變大家對神學的認識.
和他都認同.
你這個改變背後的神學的原因.
這個就真是可能要.
首先得到大家彼此的.
慢慢的互信.
彼此建立的信任的默契.
這些是需要很多的時間.
需要很多的.
大家的工作上的交流.
才可以慢慢建立到的一個信任.
最難就是這個了.
所以你可以看到.
尼熙米他.
一路這樣描述自己.
他看到真的看到城牆拆毀.
城門真的被火焚燒.
第十四節.
我又往前.
到了全門和王祠.
但所騎的牲口就沒有地方過去了.
於是夜間沿著.
可能是吸輪溪.
一路而上.
就察看城牆.
又轉身去到谷門.
就回來了.
我往哪裡去.
我作什麼事.
官長都不知道.
即是說不會好像來到打鑼打鼓.
現在有一個中央巡撫.
要來到去.
有詔書.
有王的步兵.
有工程的專業團隊.

$^{2001}$來和我一起上來建造.
這樣.
只會令到猶太的弟兄.
可能就有一個誤解.
不單止不參與.
可能還會帶來.
一個反對尼熙米的建造.
所以他現在.
先自己做好功課.
不是一個.
根本是.
即是.
在長老面前.
在那裡大搖大擺.
擺官威.
有王的詔書.
人家不接受你這一套.
而是他自己.
很低調地.
很勤力地.
自己在夜間.
去做了這些功課.
每一處地方去視察.
每一處地方.
去分析.
那個嚴重性.
那個損毀程度.
是怎樣做完這些功課之後.
第十七節.
以後我對他們說.
我們所遭的難.
耶路撒冷怎樣荒涼.
城門被火焚燒.
你們都看見了.
這個看見不是今時今日.
這個看見就是.
尼熙米未上來之前.
其實都已經看見了.
沒有能力改變.
那個現象.

$^{2041}$不能夠令到.
猶太的弟兄姊妹.
成為異象.
那個現象已經是見怪.
第一次見.
很不開心.
很激氣.
很感動.
但是每次上來.
都不見有什麼改變.
那個現象慢慢就.
慢慢慢慢就放棄了.
是改變不了的.
因為當地的人太過.
強橫了.
我們都沒有辦法改變了.
希望有新的領袖興起.
可能是這樣想.
但是.
自己就沒有能力.
組織自己內部的.
耶路撒冷的.
一個復興小組.
所以教會.
如果有弟兄姊妹自發.
為神的家.
兩三個人祈禱.
不用教會求.
自己自發的去禱告.
甚至一個人自己自發祈禱.
這樣去將那個現象.
看到的現象.
祈禱之後.
不是衝動的.
去跟別人分享.
去跟別人一起去做.
譬如我都給一點點credit.
Chris 余健峰.
很想快要復上了.
祈禱 一直等候.

$^{2081}$領袖肥沙.
我們要搞好茶水檔.
牧師你不用理會.
我們會做的.
我們就會找人.
大家回來可以相交.
大家可以喝杯咖啡.
不是想喝咖啡.
是想大家有那個溫度.
大家有那個相交.
既然有這麼好的廣場.
除了下雨.
我們都盡量用.
希望弟兄姊妹回來.
大家就一起坐下聊聊天.
互相扶持.
這些就是現象.
轉化成一個異象.
其中一個例子.
你就可以想像到教會.
如果有弟兄姊妹這樣祈禱.
神的需要.
神家的需要.
不需要等牧者去呼籲.
去禱告.
去要求你.
你都已經準備好了.
這些就是異象.
變怪不怪的現象.
更嚴重的就是.
看到這些現象.
就罵傳道人.
罵行政同工.
什麼什麼.
廁所沒有廁紙.
沒有廁紙你就下去告訴他.
我上去幫你安.
都是可以一個服侍.
現在你不是回來.
這個不是你回來.

$^{2121}$就可以好像坐下.
你就是那個消費者.
不是這樣的.
這個是你的家.
屬靈的家.
大家一起參與.
可以我們再請多十個清潔.
我們再請多十個人.
來做這些事.
不是這樣的.
那個way.
而是我們應該是.
大家都在神的家裡.
一起去參與不同的服侍.
所以你看到.
現在他就動員.
這一班的弟兄.
動之以情.
動之以那個.
的聖靈感動和真理.
不是動之.
以他的法律.
我有詔書.
我有足夠的材料.
我們人強馬壯.
我們有步兵.
我們有工程師.
所以你們要聽我說.
現在我做了.
我命令你們這樣.
領袖.
他能夠自己.
很低調地.
很謙卑地.
苦咒那個事工.
他去視察.
不是等那個報告.
不是找那些波斯的工程師.
要很高調地.
白天去視察.

$^{2161}$現在貼了.
國家現在要.
愛戴猶太人.
我們現在要為猶太人.
去訓練他.
要再漂亮一點.
不是這樣.
你們都看見了.
內霸.
我們重建耶路撒冷的城牆.
免得再受凌辱.
我告訴他們.
我神私恩的手.
怎樣幫助我.
說的不是.
君王的詔書.
君王的國家的材料.
說的是天國.
神國.
是猶太人聽得明白的.
祈禱的蒙英韻的.
神私恩的手.
外邦人說.
你的神有沒有手.
但是阿彌起瑪說.
神私恩的手.
伸到出來.
我們看得到.
跟這些領袖說.
猶太人一定聽得明白.
我們教會的.
不同宗教.
我們都會說.
我們有什麼屬靈交通.
不是交通的交通.
是彼此的.
在聖靈裡面的互相的.
分享建立造就.
這叫交通.
所以你會看見.

$^{2201}$這個字交通.
對不信耶穌的人.
都不知道你說什麼.
什麼交通.
怎樣交通.
但是基督徒就會明白.
這個交通.
大家彼此的心靈是通的.
因為大家都有神.
大家都有聖靈.
大家都有感動.
所以現在就說.
我們一起建造吧.
所以你看到.
彌起瑪能夠對.
王對外邦.
對國家.
對屬世的.
他能夠拿捏.
但是他的心底.
不是突然之間.
輕飄飄.
覺得自己.
位高權重.
而來到去.
看不起這些.
都看死你們.
幾十年在這裡.
都是這樣荒廢.
都是這樣沒用.
再差一點就是.
你們廢成這樣.
等我來吧.
現在我拿材料上來.
是很謙卑的.
因為他知道是做神的工作.
是要復興的.
是這一班的領袖.
是這一班的猶太人.
當地的沉靜了的猶太人.

$^{2241}$是要復興他們的心靈.
所以真正的.
事半功倍的復興.
是說一班弟兄姊妹的復興.
不是一堂道的陪靈那麼簡單.
是一班人.
被神的話語.
被神的愛去感動.
所以.
將這一件大家都見了.
幾十年的事告訴他們.
背後就有一份.
很大的感動.
現在神差遣我上來.
和你們一起去建造.
這個就是神施恩的手.
已經為我們.
開路了.
十八節就是了.
神施恩的手.
怎樣幫助我.
並王對我所說的話.
他們就說.
我們起來建造吧.
於是他們奮勇作者善功.
開始開工了.
但是十九節.
但和倫人參巴拉.
並圍爐的阿滿人多比亞.
和阿拉伯人基善.
聽見就癡笑我們.
藐視我們說.
你們作什麼呢.
要背叛王莫.
我回答他們說.
天上的神必使到我們亨通.
我們作他僕人的.
要起來建造.
你們卻在耶路撒冷.
無份無權無紀念.

$^{2281}$這樣來到去分得很清楚.
你們一路.
不懷好意.
而且是要排除.
在那個建造以外.
因為你們那個動機.
是極之不遵循.
是要攻擊我們這個建造.
而尼希米已經定性了.
這個是神施恩的手.
開路的運作.
就是說.
這個是很清楚.
是神的聖功.
現在尼希米.
只不過是一個神.
感動他去執行的人.
所以他不是.
現在突然之間高高在上的一個.
要統領著整個耶路撒冷的.
所有的人.
都要服他.
聽他 聽他的命令.
完全不是這樣.
但對於外邦人.
對於這些不是有心建造.
反而在這裡.
去阻攔的.
尼希米就很重的說話.
來對他們一點都不給面子.
你們是無份的.
你們是無權.
阻止我們的.
你更加不被紀念的.
一些外邦人.
當然這個是時代.
我們新約.
聖經教導我們.
我們已經.
不可以再這樣分.

$^{2321}$福音以外.
他就算不信.
我們都不可以用一些很重的.
咒詛的說話.
其實真的只要祝福.
不可咒詛.
我們可以講述他們的問題.
譬如他們怎樣反對基督教.
我們可以講述他們的理由.
但我們就不可以因為他們反對.
而咒詛那些不信的人.
我們是為那件事.
為那個真理.
去澄清.
不是為了針對人.
而去講一些.
咒詛 惡毒的說話.
基督徒新約.
其實就不是這樣分的.
在耶穌裡面.
每一個人都應該是我們所愛.
當然我們有限的.
人生和能力.
所以我們只能夠.
安排我們最需要.
去關心的人.
我們去表達那份愛.
今天就停在這裡.
本來想預備新約有一個人物.
另一個領袖.
但我都不夠時間.
下個禮拜我們再補充.
我們一起禱告.
多謝主.
我們第二章看到.
李希米他作為一個地上的領袖.
他能夠得到.
王的信任 委託.
更加重要就是.
這個被委託之前.

$^{2361}$他的動機不是為自己.
要去謀求更高的位置.
而是他看到.
他的內心是屬神的.
他要為神去發熱心.
他看到的那個現象.
他看到的那個破落.
激發他內心的禱告.
和他自己的裝備.
以致他能夠.
在幾個月之後.
可以被神感動.
神的手帶領他.
讓他排除很多的困難.
很多根本已經是.
只是一兩句說了.
整個行裝.
整個行旅.
那個困難都沒有講述出來.
因為重點可能就是.
第三章之後.
重建的那個真實的困難.
所以我們看到.
一位的領袖.
是怎樣來到.
有智慧,有見識,有膽量.
更加是有.
靈性上的.
與神同工同行的裝備.
求你在每一個時代.
都有你自己的工作.
我們看到的現象.
是有很多的弟兄姊妹.
有很多的甚至傳道人.
都移民了去.
加拿大,英國.
我們很盼望神.
感動他們每一位.
在當地有很大的復興.
也求主讓我們在香港.

$^{2401}$我們不是慘淡的.
因為神的工作.
根本就跨越任何的領域.
讓我們能夠.
忠心的在地上.
香港我們繼續.
讓更多未信主的人.
得救.
還有弟兄姊妹.
不管什麼年紀.
留在香港都是有神的帶領.
我們繼續忠心的去牧養.
彼此的去相交.
一起的去同心的事奉.
以致神的國度.
因為我們都.
等待,準備.
當內地各方面都.
配合的時候.
裡面真的有很多的人未信主.
更加有很多的少數民族.
是需要福音的.
盼望神自己.
感動,施恩.
剛才我們提過.
俄羅斯和烏克蘭.
相信始終都要停火.
但是我們盼望.
現在有這麼多的難民.
或者俄羅斯有這麼大的困局.
都需要有福音進入.
盼望神興起.
有更多.
新一代的弟兄姊妹.
同工能夠.
裝備自己.
準備自己.
這麼大,這麼困難的和場.
求主來施恩幫助.
我們這樣禱告.

$^{2441}$奉耶穌基督的名.
阿們.
\newpage



\section{}
\label{sec:Pndq2oLIHnI}
\textbf{2023-04-13 周四晚查經班《被擄歸回歷史研讀(11)》 講員:曾裔貴牧師}
\newline
\newline
連結: \href{https://youtube.com/watch?v=Pndq2oLIHnI}{\texttt{https://youtube.com/watch?v=Pndq2oLIHnI}} ~~~~ 語音日期: 2023-04-15
\newline
\newline
\hyperref[sec:Am9T43vQMKQ]{\small{< < < PREV SERMON < < <}}
~
\hyperref[sec:index_chronic]{\small{[返順時目]}}
~
\hyperref[sec:index_scriptual]{\small{[返順卷目]}}
~
\hyperref[sec:7t8JFmETTzs]{\small{> > > NEXT SERMON > > >}}
\newline
\newline
$^{1}$很感恩,又到星期四晚的茶經班.
現在是七點半,我們一起開始祈禱.
天父,我們再一次聚集.
無論是在禮堂這裡,或是在家中.
不同弟兄姊妹,又或者是有些透過.
網上的錄音,來聆聽聖經的話語.
我們都很盼望,我們所說的關於聖經的事情.
大家都能夠明白,也都和我們的屬靈生命拉上關係.
讓我們能夠經歷我們的主,你的教導,你的牧養.
今天晚上我們會進入,尼希米的建造工程的過程.
和遇到的困難,阻力.
都幫助我們能夠從其中明白.
在不同的時代,神的工作都有不同的.
也都幫助我們,不是當作歷史去看.
而是我們的主怎樣在我們這個時代.
成為我們的幫助和拯救.
也都讓我們能夠有一個信心.
可以為主的真理去作見證.
我們仰望等候禱告,奉耶穌基督的名,阿們.
其實大家在有關祈禱會都知道.
以巴的衝突,阿拉伯人和以色列人.
一直都從來沒有停止過的爭戰.
無論是為了領土的完整.
又或者是兩個民族之間,信仰帶來的種種的衝擊.
這些都在往後的日子都不能改變的.
整個大前提就是.
以色列人得到神應許的.
我們知道的新月形的以色列地圖.
即是猶太人現在回歸1948年.
他們回去得回自己的國家.
被整個阿拉伯世界圍繞著中東的局勢.
也是猶太教的聖殿所在.
所羅門的聖殿.
你可以想像得到.
那裡的地方現在也有一個清真寺.
所以其實是一個很大的諷刺.
也是一個歷史一直拖著解決不了的問題.
表面上就是不斷的衝突.
但實際上我們要看到.
在某些敬虔的猶太人的內心.

$^{41}$等候第三聖殿的重建.
即是說所羅門的聖殿沒有了.
耶穌基督的那個是希律王.
大希律王的那個是一個第二聖殿.
他們在等待一個新的將來的.
就是尼塞爾瓦會降臨的一個新的聖殿.
據聞很多的材料都已經預備好了.
但問題是那個局勢仍然在短暫.
我們眼見的都改變不了.
當我們現在回歸歷史的時候.
說到李希米他這個時代.
是一個神感動這些領袖要去重建的.
我們讀過這個《以色列記》一到第七章.
就是以色列的技術.
又是一個好像自傳式的.
技術第一次.
真是神的刑罰猶太人70年的時間期滿了.
神感動在當時的波斯王古列大帝.
用一個很懷柔的政策.
使得不同國家的被擄的人.
都可以回到他們的原居地.
所以不只是猶太人.
是其他國家的亡國的人.
都可以回去重建.
尤其是給他們有宗教上的自由.
他們的自治.
所以猶太人就回去重建聖殿.
既然是聖經的歷史.
是整個舊約一直前進.
是說和舊恩.
神的救贖計劃.
人類有關連.
所以歷史只會記這一筆.
不會記其他國家其他民族.
所以大家不要誤會了.
以為只是猶太人可以回去重建.
但是因為聖經是要記載猶太人.
他們的復國.
神的審判完結了.
他們就回去重建.

$^{81}$輾轉輾轉在那個時間裡面.
他們也是面對很大的阻力.
他們怎樣去克服.
接著就去到第八章開始.
其實應該是第七章.
以色列是第二批回去重建的領袖.
他特別強調.
他自己是一個很敏捷的文士.
也是通達神摩西律法書的教師.
所以他回去發現那個情況.
不只是硬件出現了破損.
而是猶太人出現了一個.
好像心靈的真空的狀況.
所以他就很想鼓勵自己的同胞.
要回去神摩西的律法書的教導.
所以就有一個復興.
在以色列記第八章.
一直延伸到第十章.
有一個記載.
就是一次的復興.
在整個耶路撒冷的蔓延.
他們要有一個決心.
甚至乎好像要簽一個協議書.
就是遵照摩西的律法.
在他們那個時代.
在他們那個有限對真理的詮釋底下.
他們要來將自己娶過的.
外邦的當地女子.
要來跟他們離開離婚.
生了小朋友都要趕走.
就成為了以色列記在後半段.
一次頗噁心的復興.
但是作者以色列.
就沒有交代到最後的結果是怎樣.
只是告訴我們.
有一件準備要這樣做的事.
也有一些相當具影響力的領袖.
就簽了名承諾要這樣做.
然後就停在那裡.
接著我們進入尼希米記.

$^{121}$我們用了一至兩堂的時間.
講了一點點背景.
尼希米的身份和他的熱心.
今天我們來思想一下.
當他上到耶路撒冷的時候.
他是一位怎樣的領袖呢.
在一到第六章.
又是尼希米個人的回憶錄.
第七章又是一個完整的.
一個尼希米帶動的復興.
所以那個分段也是相當清楚的.
在這個個人回憶裡面.
夾雜的就是.
你一提起作為領袖的崗位.
那個肩頭就很重了.
那些咒罵的聲音就來自四面八方.
作為領袖必須要有心理準備.
一定會面對這個重建.
要恢復摩西時代的面貌的盛典.
規模,憲制,敬虔.
你就要面對當地的外邦人.
人家已經盤踞在這裡好幾十年了.
盤根錯亂的.
這樣來做一個革新.
近乎革命.
就帶來很大的阻力.
尼希米很特別的就是.
我們說過.
他是在亞達薩西王身邊的一個酒精.
他是一個遺居要職的地方官.
你可以說他真是波斯帝國的一個.
如果我們用比較俗的.
在一些很敬虔的.
不知就裡的猶太人的內心.
說你是波斯帝國的走狗.
你現在回來重建.
到底你是什麼居心呢?.
你是否想借助王的勢力.
來令到我們的猶太教世俗化呢?.
可能會有這些初步的誤會的衝擊.

$^{161}$如果這些由內部而引伸的問題.
很難梳理的.
第五章很快就出現了.
有那些工程.
但是沒有工資.
那些工人的太太.
怨聲載道.
又要我們工作.
又不給我們工資.
可能如果再加上.
因為都不是短時間的.
再加上有其他的.
可能是一些糧食上的困難.
就會更加怨聲載道.
本來那個工程不關那些事的.
都賴了你這個工程.
賴了你會不會你們在其中.
有些人又落格.
有些官員鍾飽私囊.
貪污.
很多這些的謠言.
很多這些的誤傳.
就會繼續蔓延.
最直接的一定是.
由外邦阻撓這個工程的.
那些既得利益者.
就是我們上個星期提過的.
參巴拉.
多比亞.
阿拉伯人的基善.
這些領袖.
就是阻礙工程最直接的.
因為他們是受害者.
當建好之後.
他們就不能夠隨隨便便進入耶路撒冷.
現在就是冒險雞籠.
很多的城牆都破壞了.
就是說.
外邦人自出自入.
那就是說給我們知道.

$^{201}$在舊約耶穌還沒有出現的世代.
猶太人仍然認為.
自己和外邦人.
是不可以有太緊密的來往的.
猶太人自己是神的選民.
所以那個來往.
都是很有限度.
是要在禮儀的律法底下.
是有很嚴格的規限的.
所以如果現在.
一向都沒有那些城牆的時候.
大家都不能夠做得很足這些禮儀.
但是當做好之後.
就會出現了這個對立.
很快.
譬如我們去聯想到.
在路加福音曾經有一個這樣的神蹟.
就是有十個帶麻風的人.
來找耶穌.
耶穌就叫他們回去.
他們一去就馬上康復了.
但是那個故事.
不是告訴我們.
神的恩典.
耶穌的能力.
可以醫治麻風.
那個故事的重點不是在十個麻風.
那個故事的重點.
是十個帶麻風的人.
有九個是猶太人.
一個是撒瑪利亞人.
當他們康復之後.
那耶穌就慨歎.
那九個猶太人.
你可以很清楚的聯想到.
大家都落難了.
大家都無可救藥.
就是沒得醫.
大家都患了帶麻風.
根本不會再介意你是什麼人.

$^{241}$我是猶太人又怎樣.
我都生帶麻風.
你是撒瑪利亞人.
你又生帶麻風.
我們是九個猶太人.
和這個撒瑪利亞人.
這樣來到去.
根本大家已經變成難兄難弟的.
大家的病友.
根本就不會再在意.
大家是什麼國籍.
什麼的種族.
根本就沒有種族歧視.
一康復之後.
九個猶太人突然間彈開.
馬上就發現.
這個撒瑪利亞人和我們是不可以來往的.
馬上就走這麼遠.
馬上那個歧視就立即出現.
其實你康復.
很明顯是耶穌醫好你們十個.
就剩下這個撒瑪利亞人.
他不會介意自己是撒瑪利亞人.
因為我是一個患病的.
患這個重病的大麻風的一個病人.
我不是一個撒瑪利亞人.
再加上如果他更加意識到.
真的意識到了.
因為開始有人歧視他了.
有病的時候沒有人歧視他.
一起生活.
但是當康復之後.
那九個馬上歧視他.
說他是撒瑪利亞人.
大家沒有來往的.
就剩下這個撒瑪利亞人回來.
所以你看到在耶穌的年代.
那個歧視.
那個看自己猶太人身份優越.
都是很嚴重的.

$^{281}$所以在尼希米他們這個時候.
他們已經沒有這種能力.
這樣來和外邦人有任何的分別.
因為分別不了.
你自己的地方亂七八糟.
那些城牆都破損.
城門都隨便打開.
人家喜歡進就進.
安息日照進來買賣.
所以去到第十二章.
十三章就有關於安息日的問題了.
所以你就可以想像到.
這個關於分別為性.
在尼希米的時代.
一直延伸到耶穌的時代.
都是嚴重的課題.
亦都是摩西的律法書.
比較重的對猶太人的規範.
他們是屬神的子民.
要分別出來的.
但他們很多時候.
就是將人在宗教上的.
他們做了就以為分別了.
其實真的去到耶穌的時候.
真正的分別不是這個物質上的禮儀.
不是你去完市集買東西.
你回家要做好清潔.
那個清潔不是洗手.
不是仔細地去洗乾淨手.
免得有病菌.
而是要行一些禮儀的潔淨.
即是我證明外邦人他們的污穢.
是不會沾染我這個猶太人的.
這樣的一種自欺欺人.
但已經成為了一種禮儀了.
所以禮儀就大過天了.
所以其實重點就是.
那個分別.
要回來我們討論.
所以一開始進入重建.

$^{321}$外邦人是沒有份來重建.
所以這是上個禮拜第二章說過的.
上個禮拜本來我預備了另外一段聖經.
在新約和大家去看一看.
不同的時代都有不同的面對的困難.
特別是為神的工作來事奉的時候.
有不同性質不同程度的困難.
但是怎樣去面對呢.
到底我們應該用什麼態度去看.
神的那個恩典和那個拯救呢.
我和大家去看看第十二章的《仕途行傳》.
整章聖經值得我們和黎熙敏做一個比對.
時代是完全不同的.
正正是時代不同.
我們就可以透過舊約新約.
我們去看到同一個心智去事奉神的領袖.
面對的遭遇都不同.
我們就可以有少少結論.
我們來到去.
在我們自己面對的時代.
我們自己要拿捏得準確.
尤其是對神那個全面的認識.
所謂神學.
我們要來到去準確.
千萬不要偏向其中一方面.
你偏向了其中一方面.
我們的神學都不夠平衡.
十二章.
那時就是一個新的開始.
因為第十一章結尾.
就是告訴我們安提柯有大復興.
有很多人在安提柯信主.
甚至他們已經成為一個很有力量的教會.
是外邦教會.
在敘利亞的安提柯.
不是比西底的安提柯.
比西底的安提柯就是土耳其的內陸地方.
這是敘利亞的地方.
在以色列再上一點點的.
北面的敘利亞的安提柯的地方.

$^{361}$我們詳細不看第十一章結尾.
第十二章.
那時希律王下手.
苦害教會中幾個人.
這個希律王不是要殺耶穌的大希律.
是他的兒子希律安提柯.
那時希律王下手苦害教會中幾個人.
用刀殺了若翰的哥哥雅各.
用刀殺了雅各.
是使徒來的.
是跟隨耶穌十二使徒的其中一個.
原來耶穌離開世界升天之後.
他是第一個殉道的使徒.
是有記載在聖經裡面的一個殉道者.
是第一個.
他見猶太人喜歡者事.
你就可以想像到.
作者路加醫生告訴我們.
當時耶穌的復活.
沒有帶來整個耶路撒冷.
整個以色列.
整個當時的巴勒斯坦轟動.
你說一個人復活又不死.
而且是永活.
甚至這樣來到一路升天.
耶穌升天記是不會帶來一個震撼的.
全球的大新聞.
猶太人喜歡這個雅各殉道.
第二節的中間.
接著就抓了彼得.
那時正是除教的日子.
除教即是正月十五.
接著就是.
希律拿了彼得.
即是將他抓了.
收在監獄.
交付四班兵丁看守.
每班四個人.
四十一十六.
意思要在逾越節後.

$^{401}$即是正月二十一.
因為除教節是要七天的.
把他提出來.
當著百姓來扮他.
扮他即是動他.
即是審問他.
定他的罪.
或者根本不需要再定罪了.
根本就有罪了.
或者不需要審訊的定罪了.
於是彼得被囚在監獄.
教會卻為他切切的禱告神.
開始有教會.
所謂教會.
真的不是在政府那裡登記.
當時根本不准登記.
根本就不承認.
是一班人.
一班在猶太領袖的眼中.
在希律王的眼中.
是一班烏合之眾.
說是跟隨耶穌的一班門徒.
那個耶穌.
那個首領已經被我們殺了.
被彼拉多定罪了.
釘十字架了.
居然他們現在在民間去謠傳.
說那個死囚耶穌.
他復活了.
其實是找不到屍體的意思.
不過沒有人再去考究.
一定要找屍體出來.
好像我們解剖外星人那樣.
今天看真一點.
解剖了五天都解剖不了.
那個好笑的笑話.
其實找到耶穌的屍體出來.
那個謠言就打破了.
耶穌是沒有復活的.
但是沒有人夠膽.

$^{441}$和有興趣去周圍.
找耶穌的屍體出來.
更加不會強迫那些使徒.
說你們一定要承認.
你們要好像拍一條片.
你現在認罪.
坦白我們以前是做錯的.
其實很簡單.
但是沒有這樣做.
反而他們就要禁止彼得他們.
不要再講復活.
但是彼得就完全不聽他們這樣講.
所以到了第五節.
於是彼得被囚在監獄.
教會卻為他設設的禱告神.
第六節.
希律將要提他出來的前一夜.
那就是說囚禁了七天.
彼得被兩條鐵鏈鎖著.
睡在兩個冰釘當中.
就是在中間.
看守的人也在門外看守.
一班是四個人.
四班就是24小時嚴密看管.
就是說在監獄裡面的ICU.
因為ICU就是24小時的護士.
全程看著你的.
你想想有個姑娘.
有個男護士.
24小時看著你.
接著輪班交更.
都繼續看著你.
現在彼得就不是在ICU.
在監獄裡面.
有四個人裡面和他一起睡.
兩個在外面看著.
輪班四班.
你想想沒得走了.
第七節.
忽然之間.

$^{481}$有主的一個使者站在旁邊.
屋裡有光照耀.
天使拍彼得的肋旁.
拍醒了他.
但那兩個冰釘就睡著.
說快快起來.
那鐵鏈就從他手上脫落下來.
天使對他說.
束上帶子穿上鞋.
就是說當時可能在監獄裡面.
可能都相當熱.
不會有通風.
你不用想了.
一定是很惡劣的環境.
他就那樣作.
天使遊說.
披上外衣跟著我來.
彼得就出來跟著他.
不知道天使所作的是真的.
只當見了異象.
他曾經在第十章見過異象.
他在哥塞尼亞.
在一個做牛皮的地方.
又叫西門.
在那裡人們招呼他.
他又見到異象.
見到天上有一塊大布.
包著不同的飛禽走獸.
一直拋下來.
他在異象裡面聽到聲音.
神跟他說.
你來吃這些食物吧.
彼得就說我是猶太人.
我不可以吃這些不潔之物.
可能有些魚蝦蟹那些.
是不潔的.
不可以吃的.
叫了幾次.
後來那些異象.
那些布包著的食物.

$^{521}$收回上天.
他一直想.
為什麼我會見到異象呢.
是不是我太餓了.
所以有幻覺呢.
又不是.
跟著有人來找.
一直從很遠的地方.
找到西門的家.
可能是工廠.
前舖後居.
有沒有彼得的西門.
我們從哪裡來的.
羅馬.
在哥塞尼亞來的.
我們很想找他.
我們的主人.
這個哥尼劉.
就見到異象.
輾轉.
彼得是見過異象的.
現在他自己身歷其境.
在監裡面的內監.
有兩個小孩.
睡在他左右的.
大家都可能綁著自己.
好像下了手鐐那樣.
那些腳鏈鐵鏈.
是自己解開.
他就被天使拍醒他.
叫他穿好衣服.
跟他一起出去.
一路出去的時候.
那些門是自己開的.
不是那些間諜片.
或者很誇張的.
湯姆·高魯士的.
Mission Impossible.
不是那些.
是真是天使.

$^{561}$令到那些監牢.
首先就是那些肉質.
能夠非常含碎.
大教訓.
好像吃了迷藥那樣.
另外那些監門是可以自己開的.
彼得被鐵鏈鎖住.
又是自己會解開的.
為什麼會有這麼奇妙.
是想告訴我們.
一而再.
有神蹟令到彼得.
是不可以死的.
是神蹟令到彼得.
不可以死的.
因為他還有來日方長的.
那個傳道的生涯.
他要去傳道.
所以神就保住彼得的命.
他不可以死.
第十一節.
對不起.
還有些經文.
第九節.
彼得就出來跟著他.
不知道天使所作的是真的.
只當見了異象.
過了第一層第二層監牢.
就來到臨街的鐵門.
那門自己開了.
他們出來.
走過一條街.
天使便離開他去了.
天使自己走了.
天使完成了他的工作.
神要他去釋放彼得.
帶他離開監牢.
天使完了.
出去了監外.
天使自己走了.

$^{601}$即是說.
突然間消失.
飛上天.
好像我們看的那些鬼怪電影.
聖經的作者沒有興趣去追究.
天使是大新聞.
為什麼不再追蹤下去呢.
他走了就是走了.
重點是彼得才是主角.
所以大家不要看錯了經文的含義.
第十一節.
彼得醒悟過來.
整個人醒了.
突然有了意識.
因為之前他都迷迷糊糊的跟著走.
現在有了意識之後.
我現在真知道主.
差遣他的使者.
救我脫離希律的手.
和猶太百姓一切所盼望.
當然這個相片.
可能就是他當時的意念.
本來他自己跟自己說話.
說了一句出來.
所以這個我們都要留意.
不好意思我要關掉音.
有點吵.
繼續.
他就去到稱呼馬可的約翰.
他母親瑪利亞的家裡.
家庭祈禱會.
當時因為沒有教會.
更加你不會註冊到一個教會.
我們要拿免稅.
我們要合法團體.
當時根本就是.
你信這位耶穌.
你就要面對隨時的逼迫.
根本就不會跟你說.
好像今天我們這個社會.

$^{641}$你不註冊才是犯法.
這樣的基本的商業登記.
就不是的.
他去到這個地方的時候.
在那裡有好些人聚集禱告.
彼得敲外門.
有一個侍女名叫羅大.
出來探聽.
看看誰在晚上這麼晚來.
聽得見彼得的聲音.
就歡喜得忘記了開門.
就跑回去告訴眾人.
彼得站在門外.
他們就說你是瘋的.
侍女就說不是的.
真是他.
他們就說一定是他的天使.
為什麼會有這麼奇妙的說法呢.
為什麼那個人會說是他的天使呢.
所以有些天主教就說.
我們每一個信徒.
有一個守護天使.
大概是這個經文來的.
有一個守護天使.
但是是不是真的呢.
我們不知道的.
真的不知道.
是不是真的有一個守護天使.
會這樣保護我們呢.
都不知道.
但是偶然之間.
聽某些前輩說.
有時候有些神蹟.
好像真的.
這個我們很難去說.
但是.
這個是他們的一個認訊.
仍然都是.
大家都不敢信得這麼真.
實在太誇張了.

$^{681}$在祈禱會裡面.
就馬上神應運祈禱.
這個時候.
彼得不斷敲門.
接著第十六節.
彼得不住的敲門.
他們就開門.
看見他.
就甚為驚奇了.
彼得就擺手.
就說不要他們作聲.
都不要這麼誇張.
不要這麼張揚.
現在我是被人這樣去解救.
那些的兵丁.
可能就說我越獄.
有人來偷偷的去監牢裡面.
搶犯.
釋放了我們.
好像交換手帶.
放錯了犯.
就令到其他人可能都會擔心.
也不想牽連.
就告訴他們.
留意一下這裡.
十七節.
他們開了門.
看見他.
就甚驚奇.
彼得擺手.
不要他們作聲.
就告訴他們主.
怎樣令他出監.
又說.
你們把這事告訴雅各和眾弟兄.
於是出去.
告訴雅各.
就不是裝著香.
那個雅各.
因為雅各已經死了.

$^{721}$這個雅各是教會的領袖.
是耶穌的肉身的弟弟.
所以不要混淆了.
十二章頭那個雅各已經死了.
即是約翰的哥哥已經斬了頭.
升了天.
靈魂已經在神那裡.
不過有沒有辦安息聚會就不知道了.
但是這個雅各是當時的教會領袖.
就叫他們告訴雅各和眾弟兄.
即是當時的領袖.
然後就自己走了.
去到別處.
第十八節.
到了天亮.
丁丁就擾亂得很.
因為害怕.
四班都走失.
難辭其罪.
希律王就會追究.
就不知道彼得去了哪裡.
第十九節.
希律找他找不著.
就審問看守的人.
吩咐把他們.
十六個都殺了.
拉去殺了.
後來希律離開猶太.
下蓋撒尼亞去.
住在那裡.
聖經就是這樣記載了一件這樣的事.
我想說的就是.
在當時是一個.
很明顯的.
是一個耶穌升天之後.
不單止沒有令到耶路撒冷的猶太人.
大規模的信主.
謙卑.
感恩.
因為神的彌賽亞.

$^{761}$神的兒子.
已經降臨.
在我們的這個國家.
在我們這個時代.
所以現在是一個復興的時代了.
他們完全不能夠接受.
耶穌是死而說自己復活.
所以彌賽亞是未降臨.
所以是逼迫是合理.
是一個很應該要加入那個逼迫.
這些信徒的行列.
所以聖經的記載.
在外邦的安提阿.
在外邦的這個.
撒瑪尼亞.
蓋撒尼亞.
是有大復興.
但是在耶路撒冷.
是有逼迫.
大家記住這些士道行傳的地區性的情況.
但是總體而言.
耶穌的離開.
都是大逼迫的開始.
所以如果你生在那個時代.
你是跟隨耶穌.
你自己要預計.
你洗禮.
未必有人送花.
送殯就可能會有了.
你就可能要面對死亡了.
雅各都已經殉道了.
那你就可以想像到.
不同的時代.
我們要有心理準備.
居然在新約在耶穌的時代.
有一件這麼大的神蹟.
真是神蹟.
我們不要當這個是天方夜譚的虛構故事.
彼得神不要他死.
你怎樣捆綁他.

$^{801}$十六個還是羅馬的兵丁.
訓練有素.
大家去過機場吧.
復活節應該大家都去過機場了.
最近.
那些機場特警.
個個都很大隻.
裝備全部是隨時都可以開火的.
重裝備.
你想一下.
這些才叫訓練有素.
就不是那些衣服不整齊.
又年紀大.
不是那些.
不是那些城管.
而是真的.
看守著彼得的十六個兵丁.
希律王手下的兵丁.
但是天使何時來.
何時走.
天使做了什麼.
一個兵丁都不知道.
以為他們集體失職.
甚至收了黑錢放彼得.
所以希律王都不知道怎麼辦.
就算了.
殺了十六個就算了.
這樣的一個處境.
我們想說.
當我們比對不同的時代.
我們稍後就會見到.
你要有心理準備.
尼熙米很多時候是訴諸於.
他有祈禱.
但是更多是他的囚惑.
他的智慧.
他的謹慎.
因為他步步為營.
他知道裡外不是人.
你是空降下來.

$^{841}$重建聖城.
當地的土著原居民.
那些阿拉伯人.
一定會抵擋你的.
就算你是河西的官員.
河東的對不起.
因為他們過回這一邊.
你是阿達舍西王的一個大臣.
但是當你來到這些山旮旯的地方的時候.
人家未必會靠近你.
而且是切身利益有關的.
即是你是中央的官員.
來到地方.
地方的官員.
是會有阻撓的.
這是我們幫大家做了一點分析之後.
你進入經文.
你就會發現這些情況.
再加上猶太人.
其實還是很半信半疑的.
是不是真的接納這個尼希米.
作為我們一個屬靈的領袖.
去帶領我們去重建.
以致我們能夠回到摩西律法的要求呢.
這個是兩面不是人的.
作為一個尼希米在他那個時代的領袖.
最好的就是.
好像阿彼得那樣.
有天使出現.
天使來到他身邊不斷保護著他.
讓尼希米能夠有光環.
外邦人不敢動他.
猶太人敬重他.
因為這個是神的使者.
現在來拯救我們.
但是我們在整卷的尼希米.
見不到一個的神蹟.
一個都沒有.
但是在新約.
我們又發現.

$^{881}$在《士徒行傳》初期教會.
很多很多的神蹟.
我們在第幾章.
阿彼得的手帕.
不是什麼牌子.
他的手帕可以醫病.
人們拿他的手帕.
他的影都可以醫病.
可以神蹟成這樣.
你就記住.
你就不要拿阿彼得的遺物.
我要搶他的手帕.
那我就沒事了.
沒用的.
不是那些物件.
是他所信的主.
在那個時代.
神要彼得.
來行這麼大的神蹟.
其實是印證那個福音.
你就要記得.
初期教會是有很多這些神蹟的.
所以今天.
我們絕對不是要批評.
靈因派如果過度強調神蹟.
沒有神蹟的教會.
就不是教會了.
沒有神蹟的教會.
就是軟弱的教會.
可能是對的.
但是我們要很仔細去看.
是不是沒有神蹟的地方.
就等於沒有神的工作呢.
現在我們強調的是神的工作.
神自己透過一個時代.
或者一些人去工作.
這個是最重要的.
那個大的神學.
所以現在說的是那個時代.
是一個什麼時代.

$^{921}$是不是一個有大量神蹟的時代.
和神是不是只能夠.
透過神蹟去工作.
如果是.
當然沒問題.
我們要自認.
如果香港有二千間教會.
其實不止二千.
如果只有十間是有神蹟.
那對不起了.
其餘的一千九百九十.
都是不可以存在了.
因為沒有神蹟.
行不到神蹟.
就要麻煩你自己除牌了.
那很合理.
問題就是.
為什麼你堅持.
你們行不到神蹟.
才叫做是教會呢.
才叫做是有工作.
有神工作的教會呢.
問題就是.
我們不可以隨便這樣去說.
誇耀自己.
要看整個時代.
要看神的工作的作為.
這樣才是公允和客觀.
如果不是.
我們自說自話都沒有意思.
更嚴重就是.
你假借了以為神.
是真的這樣做.
可能原來那一千九百九十間.
才是真正的教會.
而這十間.
不是神要的.
認可的教會.
這樣就更嚴重了.
你個人吹噓沒有問題.

$^{961}$但原來你帶領著弟兄姊妹.
不是在神的正道上.
那個就更嚴重了.
所以我們這樣.
兩個領袖.
兩個時代.
做一個比較.
本來上一個禮拜.
要和大家先做一個比較.
我們進入第三章開始.
到第六章.
聖城建好之後.
你就會看到.
有些事情.
我們要來做一個歸納.
和一個小小的分析.
第三章.
那時.
你閒謎第三章.
那時大祭司以利亞實.
和他的弟兄眾祭司.
起來建立陽門.
那個圖跡.
大家上網都找到.
他有不同的門.
其實是一個.
不規則的多邊形.
他一路圍繞著.
整個耶路撒冷.
那個舊城.
這樣來建造.
當然就是圍繞著聖殿去建造.
重點是聖殿.
聖殿就要在.
整個圍牆裡面.
最重要的地方.
不可以有另外的地方.
必須是在.
所羅門建立的地方.
重建.

$^{1001}$所以這裡告訴我們.
他們一路由陽門開始.
分別為.
安納門線.
即是有神的祈禱.
祝福.
由竹城牆去到.
哈米亞樓.
直到哈南葉樓.
因為有不同的樓房.
即是有門.
亦都有看守的.
在城門上.
來做守衛.
所以有樓房.
其次就是.
耶利哥人建造.
其次是欽利的兒子.
薩克建造.
開始記載那些公頭.
哈西那的子孫.
建立漁門.
嫁橫樑.
安門線和那個山鎖.
即是開始.
安好門框.
將門擺好.
還加上那些鎖.
可能是木鎖.
其次是哈哥斯的孫子.
烏利亞的兒子.
米利末收造.
其次是米氏薩別的孫子.
比利加的兒子.
米舒蘭收造.
其次是巴拿的兒子.
薩督收造.
其次是提哥亞人.
收造.
但是他們的貴就.

$^{1041}$不用兼擔他們主的工作.
這樣讀是不是很悶啊.
剛才看第十二章.
嘩又神蹟.
比特自己都很詫異.
嘩為什麼這麼戲劇性的.
我是不是在做夢啊.
大家情願看哪些的記載.
當然是看神蹟的記載.
但是這些的記載.
很悶啊.
但是告訴我們.
尼熙米用了.
多少的時間.
和這班的領袖開會.
說服他們.
神怎樣.
來聽我的祈禱.
我在波斯的.
這個書山城的籌劃.
準備.
那些圖則.
我已經計算好.
要用多少材料.
大概我已經做了.
前期的工程.
那個預算.
要用多少的木料.
用多少的石.
要用多少的錢等等.
我已經做了.
那個早期的.
那個預算.
我們之前都說想.
打通那個.
106去那幢牆.
出來就是側門.
再從.
翻新101至103.
變成整個地方.

$^{1081}$如果活化之後.
可以吸納很多這個區的青少年.
我們本來的心.
是想這樣.
做了前期的工程.
找了弟兄免費幫我們做了報價.
那些主力牆要打.
要多少錢.
要用什麼材料.
連那個圖都可以畫出來.
讓我們看到一個.
電子版的PDF版的.
這樣的工程.
不需要用這樣的.
一個時間來規劃.
可想而知.
李希敏.
不是兩手空空的.
上來的.
在這裡頭那五節經文.
能夠說服.
那麼多的領袖.
大家做不同的地方.
在不同的門.
你做你的工程.
我做我的工程.
大家彼此的聯絡.
要用多少的.
心力.
要用多少的.
口才去說服.
猶太人的領袖.
支持這個.
的事工.
大家一起去做.
就不是我們搬了材料上來.
我們找科斯人.
幫手去做.
而是猶太人.
當地自己的猶太人.

$^{1121}$自己一起去做.
你想想.
這些經文何來有神跡可言呢?.
何來你見到.
這樣讀這些五節經文.
全部是李希敏.
個人的魅力.
她的籌劃.
她的用心.
她的熱心.
她要放大架.
去為神的工作.
去擺上她的心血.
她的心力.
她要畫圖則.
來計算.
這個地方要用多少材料.
那些工人.
就要搬運那些磚要多少.
是很多很多的工程.
第六節.
巴西亞的兒子.
耶和耶大與彼所代的.
兒子米舒蘭修造.
古門,駕橫樑.
安門線和山鎖.
不斷都是說這些.
重複的字眼.
就是欺騙人.
米拉提.
欺騙人是奴隸.
在約書亞年代.
他們是騙了約書亞.
說自己是在很遠的地方來.
後來約書亞就不殺他們.
因為他們歸行.
就以為.
我們多厲害.
都不用打就歸行了.
殊不知原來在近的.

$^{1161}$神規定一定要殺光.
但因為他們.
下了一個誓.
我們不會殺你們欺騙人.
那怎麼辦呢.
求問神.
雖然現在不殺你.
是因為我們太大意.
所以欺騙人你們.
以後你們的子子孫孫.
都要做奴隸.
要在聖殿裡面.
來做斬柴.
淡水下賤的工作.
欺騙人就說.
沒問題.
不用死.
本來要死的.
被神滅絕的.
但現在我們可以歸入神的名下.
所以欺騙人來到這一代.
被老歸回.
回來.
仍然都是做奴隸.
他們都不介意.
世世代代感恩不盡.
這樣來服侍.
所以.
第一代的弟兄姊妹.
是非常非常.
愛著神的恩典.
但很多代之後.
你的子孫.
還是不是像欺騙人一樣.
都不介意自己的身份.
繼續在猶太的地方.
在神的殿裡面.
做這些下賤的工夫.
沒什麼人做的工作.
這個就是欺騙人的故事.

$^{1201}$接著有米倫人.
雅頓和欺騙人.
並屬河西總督.
所管的.
其次是銀杖.
哈海雅的兒子.
烏蝕修造.
一直都在說修造.
即是說有陶側.
有工頭.
有監工.
有材料的運輸.
有一些不同的專業.
去參與.
安門不是那麼簡單.
要安平水.
要整好那些位置.
都不是那麼容易.
不是說.
你整就能整.
現在我們有一個叫恩典維修組.
久不久他們的弟兄.
就會來.
阿兵就會報告在群組裡.
教了現在哪裡門脫了.
他們就安排時間來整.
我們現在的風扇.
又整好了.
你就可想而知.
很多東西是要規劃的.
這裡告訴我們.
尼熙米用他的公務的時間.
來為神的殿.
去做這個領袖的工作.
去規劃.
去安排不同的人.
這個過程.
就是他要說服.
他要感動他們.
材料我有.

$^{1241}$亦都有亞達薩西王給我的預旨.
是開綠燈的.
但是真正要做.
其實是要猶太人.
大家一齊做.
所以如果我們看到這些工程.
教會的弟兄姊妹.
如果我們每個禮拜.
當然你剛剛回來.
又未相信.
你每個禮拜回來坐在這裡聽道.
已經是你的工作.
但如果你相信了很久.
你每個禮拜回來又是站在這裡聽道.
那你就有辜負了.
你自己所夢的恩典.
其實是應該你要想一想.
其實不是想一想.
問哪個人.
我甚麼都懂 你告訴我.
不是這樣 是你坐下來.
去求問神.
過去你給我這麼多恩典.
我有這麼多恩賜.
或者我有這麼多資源.
我在哪一方面.
可以在神你的家裡.
去服事呢.
記得那個次序.
是影響你的心靈.
你的心靈.
是要服在神的面前.
先有一個奉獻給神的心靈.
我有甚麼可以.
貢獻給這個地方呢.
那你就開始要問傳道牧者.
啊 牧者.
我有甚麼可以服事呢.
那牧者就會安排.
這樣的次序是很重要的.

$^{1281}$就不是你一個.
高高在上 甚麼看不上眼.
但又不幫忙參與.
甚麼你都批評.
這樣的弟兄姊妹.
其實不要介意說得這麼重.
真是你如果.
越信主內 而你繼續高高在上.
小心.
神阻擋驕傲的人.
賜恩給.
謙卑的人.
很重要的.
因為你有這個眼光.
證明你已經是驕傲.
你有一個謙卑的心.
哇 樣樣東西我都可以幫忙.
如果你需要的話.
這樣的態度就已經完全不同了.
記得這個不是俱樂部.
這個是神的家.
這裡是有神.
所以你看到.
這裡就是尼希米.
他一直可以聯絡這麼多的領袖.
這麼多的領袖都願意.
和他一起配合.
你看到就是 去到第九節.
其次是管理耶路撒冷.
一半.
護爾的兒子利法雅修造.
其次就是.
哈路末的兒子耶大雅.
對著自己的房屋.
修造.
在自己附近的地方.
有些人不是那麼近的.
都不介意.
其次就是哈沙尼的兒子.
哈達修造.

$^{1321}$第十一節.
哈林的兒子馬基雅.
和巴哈摩亞的兒子.
哈述.
修造一段.
並修造牢留.
其次就是管理耶路撒冷那一半.
剛才說的另一半.
現在就是其餘的一半.
哈羅克的兒子.
沙龍和他的女兒們.
修造.
為什麼要強調女兒們呢.
一般就不會拋頭露面.
去做這些戶外工程.
也就是說.
又有很特別的地方.
神又記載在這裡.
讓我們看到.
他的女兒可能都是工程師.
有幫忙做.
第十三節.
又或者可能不夠藍丁.
可能也不奇怪.
第十三節.
哈嫩和薩羅亞的居民.
修造這個谷門.
納門,安門線和山鎖.
又建造城牆.
一千爪.
一直到這個糞廠門.
就是那個像是倒灰的地方.
就是處理一些垃圾的地方.
所謂垃圾.
是他們獻祭.
然後倒的灰.
清理的灰.
管理伯哈基林利甲的兒子.
瑪基亞.
修造糞廠門,納門.

$^{1361}$安門線和山鎖.
主要這次.
這個前期工程.
就是要安好門.
讓那些城牆.
雖然還沒有建好.
但是先有門.
有山.
可以晚上鎖的.
外邦人進來.
是會被人發現.
發覺到的.
這樣的安排.
要圍到整個聖殿.
要有一個規範.
他們很強調.
猶太人是神的選民.
其他.
如果你不歸入猶太教.
你是不可以來聖殿的.
當然這個和我們今天說的一樣.
這個不是一個游泳池.
你不可以租這個地方.
我想在這裡讓小朋友玩水.
很熱.
這是浸池.
誰可以進去呢?.
一定要信主.
經過牧者.
執事去研判.
透過上課.
他又講到自己的信仰.
在一個特定的時候.
所謂特定的時候.
就是一個隆重.
眾弟兄姊妹聚集的時間.
不是私下.
我們隨隨便便.
沒有人知道.
就找個人幫他洗了.

$^{1401}$不是這樣的.
是要很公開的見證.
這樣洗禮.
所以你可以想像.
整件事是在說.
那個屬神.
所以猶太人.
他們很緊張分別為性.
新約的分別為性.
就是你有沒有聖靈.
你有沒有信仰.
即是耶穌的救恩.
是不是在你裡面.
所以真正的分別為性.
就在這裡.
不是你穿什麼制服.
或者什麼特別的.
不同的儀式.
就不是那些.
十五節.
管理米斯巴葛河西的儀子.
沙崙修造全門.
納門蓋門頂.
安門扇和山鎖.
又修造靠近.
黃猿西羅亞池的長門.
直到那從大衛城下來的台階.
西羅亞池就是大衛的舊城.
我們去過聖地.
你就知道.
真是大衛舊城.
有一個希西加的水道.
一路走出來.
走出來那個一見到光的池.
因為現在已經乾了.
那個就是西羅亞池.
駐彭哲.
祭司就在那個西羅亞池打水.
就抬上去.
聖殿.

$^{1441}$抬上去聖殿.
所以是向上的.
大斜路來的.
上回聖殿.
在聖殿那裡就倒水.
這個就是.
駐彭哲最大那一天.
第八天.
他們就會這樣做.
所以在駐彭哲.
耶穌就說我才是那個活水.
這是約翰福音第七章.
我們那時也說過.
你可以想像得到.
就是在不同的地方.
我們現在這樣讀.
好像已經有一個平面圖.
已經有一個工地.
擺在我們的面前.
我們見到的就是.
既然現在是工程.
自然就是有工程師.
自然就是有策劃的那個有心人.
就是尼熙米.
回到去之後.
他來到去用很多的時間.
和領袖去開會.
去遊說.
我一直都說.
大前提.
是神施恩的手.
幫助他的行程.
幫助他的心願.
達成這個心願.
就是我希望.
在我有生之年.
見到神的殿.
周圍的城牆.
都不要破落.
城門不要來到去被人踐踏.

$^{1481}$又被火焚燒.
希望聖殿.
有一個有規有矩的敬拜.
就是這麼多.
這個就是尼熙米的.
禱告和神聽.
祈禱的成就.
你有沒有見到.
有神蹟在其中.
連一隻字都沒有.
連我們.
一直這樣讀的時候.
都完全感受不到.
有神蹟.
完全見到一個人.
在這裡走來走去.
真正的奔跑.
為神.
大發熱心的去做.
他是一個酒精.
往後再讀下去.
再多的困難.
再出現.
他甚至要自己自掏腰包.
自己拿自己的私房錢.
去給那些猶太人.
出糧.
你可想而知.
在尼熙米記.
你見不到任何一句.
是有神蹟.
但他都是為神.
忠心一代的領袖.
出現了我們剛才問的一個問題.
到底.
是不是有神蹟.
才是神工作的教會呢.
沒有神蹟.
是不是教會.
神在不在其中工作呢.

$^{1521}$如果我們用這樣的邏輯.
這樣記載在聖經.
聖經就更加.
真實的啟示.
我們就可以有一個.
很清楚的.
其實現在都已經可以做到一個歸納.
肯定是錯.
不可能.
神只是容許神蹟.
才是神的工作.
如果.
容許神蹟才是神的工作.
那就相對簡單.
我們的信仰.
也相對極之落後.
因為我們今天眼見到的.
整個宇宙.
整個世界.
那個結構.
那個工整.
那個奇妙.
那個自然的運作.
真是自然運作.
樣樣東西.
你知也好,不知也好.
它都是生生不息地發生的.
不同的.
動物界.
有他們自己的秩序.
我們人類的世界.
也有我們的秩序.
科學越是進步.
越發現人體的.
那個結構的奧妙.
天文學.
宇宙那個望遠鏡.
越來越準確.
越來越高清的時候.
越望得遠.

$^{1561}$用那些電子來射線.
得出的結論.
是說幾十億年後的事.
那個這麼大的發現.
只是告訴我們.
讓我們更加了解.
人類是極之渺小的渺小.
但是渺小得來.
人類自己都是有結構的.
是有秩序.
是可以累積過去的.
那個知識.
而令到下一代可以再進步的.
你就可以想像得到.
神的那個偉大.
所以單是不信的人.
都可以否定了.
只有神跡才能經歷神.
是一個迷信的信仰.
甚至乎是一個.
很錯的信仰.
因為神.
他可以讓神跡出現.
正正就是他更加喜歡.
沒有神跡.
而那些信他的人.
透過信心.
透過倚靠.
而能夠做到一些事情.
那個就應該叫神跡.
那個沒有神表面的介入.
而能夠做到一些事情.
這個應該更加理解為.
是神工作的一個痕跡.
就是神跡.
所以不是超自然.
每次都要.
今天糟了約了六點吃飯.
遲到,馬上祈禱.
神讓我衝過去那個湖.

$^{1601}$走過那個湖.
每個居民走在湖上都看到.
嘩,為什麼這個人走在湖上.
為什麼不會沉.
他們是什麼人.
不是.
神不要這樣.
衝過去,為什麼不說一聲對不起.
我只顧著看一些文件.
忘記了.
遲了來.
說一聲對不起就行了.
為什麼要走過湖,走神跡.
衝過來.
為什麼每次都要這樣衝.
為什麼下次不可以.
六點多就要準備過去.
很多這些事情.
我們千萬不要有一種迷信的觀念.
沒有神跡.
沒有醫治.
就不是神的工作.
千萬不要這樣做一個這樣的結論.
相反就是.
我們看到這個時代.
其實很需要有.
神的恩典.
很需要有.
這些的領袖.
很需要有奇妙的神跡.
來介入.
讓他們排除很多困難.
外邦人.
只是有記錄的.
多比亞,阿拉伯人.
基善,參巴拉.
這個和倫人.
已經是這樣處處去.
阻撓.
這個尼希米.

$^{1641}$接著還很陰毒.
買通了一個猶太的祭司.
叫尼希米.
大家可能還沒有讀.
叫尼希米.
去聖殿裡面開會.
我有些秘密.
想告訴大家.
尼希米說.
我不是祭司.
我不可以進入聖殿.
你不要在這裡.
看出是騙案.
你可以想而知.
尼希米是不斷用他的智慧.
周旋在.
這些敵對勢力底下.
看不到有神跡.
所以尼希米.
是沒有一處地方.
一句是說到神跡.
但是我們看到.
一個領袖.
有沒有神跡.
他都是這麼殷勤的.
按照神給他的恩典.
那個感動.
那個熱心.
和他的能力.
和他的.
在過去訓練他的.
而且那些訓練是波斯帝國給他的.
那個訓練.
讓他能夠有智慧.
有分析力.
有親和力.
可以和.
不同的人.
講他的計劃.
可以說服對方.

$^{1681}$參與這個建造.
這個過程本身.
已經是很高的藝術.
是很高的難度.
要人.
認同你所說的.
一起去做.
首先可能是觀望.
觀望已經是很好的了.
他開始有了解.
很想看多一點.
觀望完之後.
如果發現不對勁.
就不參與.
甚至阻撓你的工程.
這樣經常都出現.
你有多好的政策.
如果沒有經過很好的梳理.
都會被人有很多.
示意的批評.
甚至.
完全.
可能不合理的.
那種單打.
一定會令到.
整個政策的初心.
體無完膚.
但是黎希米.
其實可能他真是身經百戰.
他很清楚.
人性的軟弱.
所以他自己都來到.
在這個時候.
他都有一個很倚靠神的心.
在這個時候.
開工順利的時候.
突然下雨.
突然之間有阻礙.
第四章就出現了.
所以你不用擔心.

$^{1721}$既然這麼.
仔細的說他的回憶錄.
在這六章裡面.
由第四到第六章.
那些困難就陸續湧現.
所以現在你看到這麼簡單.
誰誰做哪裡.
誰誰又做哪裡.
這個是說他有很強大的.
號召力.
但是這個號召力.
一遇到難處就會.
信心動搖而瓦解.
我們再讀一點.
第十六節.
其次是管理白戚一半.
這個壓卜的兒子.
黎希米收造.
直到大衛墳地的對面.
又到挖成的池子.
並擁視的房屋.
其次是利未人巴尼的兒子.
利弘收造.
其次就是管理基伊拉一半.
哈沙比亞.
為他所管的本景收造.
其次就是利未人的弟兄中.
管理基伊拉那一半.
開始又祭司利未人.
希拿達的兒子巴瓦伊收造.
其次是管理米斯巴耶穌亞的兒子.
以謝收造一段.
對著母父的前輩.
以謝收造一段.
對著母父的前輩.
對著母父的上輩.
城牆轉彎之處.
其次是撒拜的兒子巴祿.
揭力收造一段.
從城牆轉彎.

$^{1761}$直到大祭司.
以利亞實的城門.
大祭司開始.
由他記載第三章一節.
他帶頭.
因為他是當地的領袖.
你一定要說服到他.
認同你這個建造.
你一定要說服到他.
所以現在回頭說他.
說以利亞實.
再做這一段.
第二節.
其次是駐平原的祭司收造.
其次是變雅敏哈述.
其次是駐平原的祭司收造.
對著自己的房屋收造.
其次是亞蘭尼的孫.
馬西亞的兒子.
阿薩利亞.
靠近自己的房屋收造.
其次是希臘的兒子.
賓內收造一段.
從阿薩利亞的房屋.
直到城牆轉彎.
開始有一個輪廓.
因為是不規則的圓形.
即是說開始有轉彎的地方.
需要一點技巧.
如果你看過人砌磚牆.
打直沒有問題.
但有些轉彎位.
有些巧妙的.
那個位置的力學.
可能要弄一個柱.
一個主力的牆.
的磚位.
我小時候都見過人弄.
自己都幫手弄過.
自己砌的磚永遠都不會直線.

$^{1801}$你見人家很容易.
好像很直.
你砌上去就知道.
很考功夫的.
原來放每一塊磚壓平.
不要說外國那些現在用機器弄.
是整行機器放下去.
那些不算.
說回本地那些三行.
弄那些外牆.
你自己學人家砌.
你就知道原來不是那麼簡單.
小時候見到那塊磚明明好.
他要剖開一半.
原來那塊磚.
不可以一塊塊地放.
是會倒的.
那個力一推就倒.
是要不規則每一行.
是不一樣的.
一來線條好.
原來是力學的問題.
你不懂就告訴你沒用.
因為你看人弄你就會知道.
但我自己弄過.
那時十多歲.
家裡要砌磚.
砌不到.
那裡那些磚.
歪歪斜斜的.
又很浪費那些英泥.
人家弄得很工整.
很漂亮.
但你也知道.
背後的技巧是經年累月.
累積的經驗.
每一樣東西.
都是有師父的.
我想告訴大家.
這些一路轉彎的地方.

$^{1841}$開始有這些城角.
即是城的角落.
第25節.
烏塞的兒子巴拉修造.
對著城牆的轉彎.
和皇上宮突出來的城樓.
靠近護城樓.
靠近護衛院的那一段.
其次是巴六的兒子.
比大雅修造.
尼提靈住在俄斐勒.
直到朝東水門的對面.
和突出來的城樓.
很長.
其次就是提哥雅人.
又修造一段.
對著突出來的大樓.
直到俄斐勒的這個牆.
從馬門往上.
眾祭司各對自己的房屋修造.
其次就是.
欽默的兒子薩都.
對著自己的房屋修造.
其次就是守東門.
是加尼的兒子.
是瑪雅修造.
其次就是.
是利米雅的兒子.
哈拿尼雅和薩拉的第六兒子.
哈嫩又修造一段.
其次就是.
比利加的兒子.
米舒蘭對著自己的房屋修造.
其次銀杖.
瑪基雅修造.
到尼提靈和商人的房屋.
尼提靈就是那些看門的電翼.
是一個織銜來的.
等於我們的.
詩詩.

$^{1881}$或者詩琴.
是一個織銜來的.
尼提靈不是一個人來的.
對著哈米弗甲門.
直到城的閣樓.
銀杖和商人.
在城的閣樓和陽門.
中間修造.
大概將一個不規則的圓形.
安排了不同的.
工人不同的.
工頭有些是有.
名字的門有些是.
自己房屋的.
這樣來到去好像佈置.
好那個修造.
好像我們壽苦節.
我們在開同工會的時候.
哪個哪個.
就是負責.
做什麼工作了行政同工.
又做什麼工作了安排好了.
大家有一個小小的.
預言大概的流程.
整份流程表我們對一次.
終於又發現了.
一個小瑕疵.
就是師班站在那之後.
Elsie她是負責那個音響.
開咪關咪.
因為如果你開著了有時她自己會關掉.
有時會有干擾.
所以到誰拿咪就要說.
終於發現了.
有一個瑕疵我們沒有想到的.
就是師班站在那裡.
好像人牆.
後面領唱.
遮住了.
當她開咪我們就見得很清楚.

$^{1921}$她不停的轉那個按鈕.
為什麼沒有聲音呢因為.
Elsie在後面根本被所有.
師班員遮住了.
我們那個清單就漏了這裡.
領唱會說話.
但是.
Elsie就需要開咪.
她用的iPad控制那個.
控制台的panel.
在後面看不到.
所以延遲了兩秒鐘.
才那個領唱.
可以說話.
很多東西是要有一個預演的.
原來都漏了那個部分.
可想而知.
尼熙米在第三章.
全部是.
自己人建造.
全部是.
那些領袖.
大家.
一呼百應.
這樣去做.
大家可以自己.
做哪個部分.
每人做一個部分.
記住第三章是一個.
現在是好天建造的時候.
所說的好天.
不是真的好天.
是還沒有開始有難處的時候.
所以真正見到一個.
弟兄姊妹.
那個心靈.
仍然面對挫折.
面對那些批評.
都仍然繼續堅持.
那個是真正的敬虔.

$^{1961}$為神而做.
所以我們需要.
進入下一次我們說的第四章.
我們就會發現.
那些批評就陸續而來.
第四章第一節已經說了.
參巴拉.
聽見我們修造城牆.
就發怒.
大大的怒恨.
癡笑尤大人.
對他弟兄和薩馬.
對不起我先喝口水.
都是有些乾.
所以不要小看第三章.
好像是很順利.
很多人參與.
沒有什麼難處.
外邦人.
很接受.
你們去建造.
因為有皇的詔書.
因為有.
中央的大臣.
親自來督工.
好像更加有.
皇下了諭旨.
河西的總督.
要給那些材料.
因為你沒理由在這麼遙遠.
運這麼多材料上來.
一定是去到河西.
那裡去徵用.
當地的材料.
是皇批准的.
但是.
第四章就馬上給我們看到.
所以尼欽米.
他先是說了.
他自己的心裡的計劃.

$^{2001}$就是.
不同的地方.
他已經視察過了.
他已經巡查了.
規劃了.
畫了圖織.
雖然當時沒有iPad.
不可以馬上給電子PDF.
大家看.
是沒有的.
可能是透過.
很粗糙的圖織.
來給每一個的.
公頭,每一個支派.
來做你自己的部分.
你可想而知.
這個是一個本身.
令到內部團結.
已經是很大的神職.
是神的工作.
所以.
因為沒有神職這些字眼.
就不是神的工作.
教會到今天.
都仍然.
見證耶穌是神的兒子.
沒有.
消失到這個真理.
是神保守.
不是人的.
一代一代的堅持.
是神在.
不同時代的領袖.
弟兄姊妹的心裡.
來去謹守.
有些時代是大逼迫的.
譬如由耶穌升天.
去到君士坦丁皇帝.
信主.
這三百年.

$^{2041}$是大逼迫的.
很多人殉道的在羅馬.
很多人的骨頭.
就在羅馬的.
亂葬崗的了.
輾轉後來.
教會.
能夠上水面了.
可以見天日了.
因為皇帝信主頒佈.
基督教是國教.
就開始不會再有逼迫了.
這些是.
當時的三百一十八.
還是三百二十八年.
我先溫習歷史.
已經很久沒有再溫習了.
君士坦丁皇帝信主.
這樣來令整個基督教.
開始蓬勃地去傳遞.
一蓬勃.
自然就慢慢慢慢又會.
衍生一些腐化.
你可以想像得到.
在這個時代.
總是會有困難.
會將基督教的教義扭曲.
將真理.
可能會來去竄改.
所以.
由三百幾年去到大概五百幾年.
是教父很盛行的時代.
也是那些信經.
很重要的時代.
很多的信經.
都是在這大概二百年裡面.
出現的.
使徒信經是每個星期都會背的.
還有其他的信經.
大家應該見都沒有見過.

$^{2081}$尼西亞信經.
加克敦信經.
阿他拿修信經.
還有其他的.
甚至乎那些.
要理問答.
那些其實就有專書.
去記載這些歷史.
我們今天只是選取.
使徒信經.
使徒信經是最基礎的.
真理的辨別.
在這個時代之後.
有很多異端.
所以有新的信經.
更加仔細的描述.
救恩.
耶穌和神.
三位一體的神是怎樣的.
更多的信經.
是在使徒信經.
之後再出現的.
那些時代為什麼會有信經.
正正就是異端的出現.
是自己基督教內部的異端.
不是外面攻擊我們.
是基督教內部的異端.
還要是那些領袖.
那些叫做教父.
錯誤了的神學.
所以需要有另外的領袖.
要澄清.
要開一些大公的會議.
來確立.
哪個領袖是異端.
要將他趕出教會.
然後就要有一個.
教會的領袖.
趕出教會.
然後就要有那個信經頒佈出來.

$^{2121}$來大家.
教會要遵守.
這樣的歷史發展.
所以現在我們來到.
尼希米這個時代.
是一個重建的時代.
神不給他們神跡.
讓一個領袖.
可以獨挑大樑.
首先就是.
得到皇帝.
那個.
差不多絕對的信任.
給他有放大架.
給他有一個.
詔書.
即是皇帝的.
聖旨.
然後准許他.
去到河西.
他由河東.
去到以色列.
可以在當地.
拿物資.
來去建造.
城牆.
往後就是尼希米.
他自己要發揮他的小宇宙.
他自己.
要怎樣能夠上到耶路撒冷.
之後和那些.
當地的領袖.
其中最大的就是.
這個大祭司.
一定要取得.
他的信任.
支持才可以.
去說服其他的祭司.
其他的地方的領袖.
或者一些做生意的人.

$^{2161}$或者其他不同的職位的.
譬如有利緹靈.
有那些手電的電役.
有其他的支派的人.
你可以想像得到.
第三章.
其實是有很多的.
純洩.
用我們再簡單一點.
就是喝了很多次茶.
摸著茶杯.
酒杯.
談了很多東西.
談了那個大綱.
大家都明白了.
大家都知道.
大概有個計劃.
可以多少天就能建好.
這些工程.
我想做這個比較.
彼得他生在一個.
有大神蹟的世代.
因為神.
要透過使徒的見證.
印證.
他的愛子耶穌的福音.
我看看希伯來書.
第一章.
希伯來書.
很清楚地說.
這個話.
可能第二章會更準確.
希伯來書的第二章.
第二章.
所以我們要.
越發鄭重.
所聽見的道理.
當然是主耶穌是誰.
恐怕我們會.
隨流失去.

$^{2201}$隨流.
水流.
水流即是會有水漲水退.
我試過在.
中一暑假.
那時中學.
中文大學.
現在.
中大的火車站.
對出大家知不知道.
是一個海灘.
很多馬尿水.
爬艇.
我中一的年代.
現在新一代應該完全沒有見過.
現在就是一個大巴士站.
接著再出一些.
中大的酒店.
辦公樓.
等等.
當時沙田進去.
馬尿水站.
未叫中文大學站.
馬尿水站.
就是一個.
爬山板的地方.
有個同學在那裡.
開山板的店.
暑假我們一班同學.
去玩.
那時.
有些板一直出到海中心.
大家如果有印象.
爬艇仔那些.
好像有條橋.
一直走出去.
由地下一直走出去.
坐那隻山板.
有些繩.
綁住那些木.

$^{2241}$很深刻.
我們爬了個艇出去.
有條波褲.
污糟了.
出去游水.
掛在那些木上.
不知道.
回來後原來水漲.
那條褲子沖走了.
因為中一都算窮.
那條波褲是ADDI.
都幾貴.
沖走了.
我們就發現.
原來有水漲這件事.
有潮汐的漲退.
作者告訴我們.
信仰慢慢的.
隨著潮流.
水流.
會失去.
什麼潮流呢.
即是說.
如果這個世界.
信耶穌是很落後很老土.
可能慢慢你都發現.
你不再熱心.
不再回教會.
恐怕.
隨流失去.
不是突然失去.
如果是洪水.
人人都害怕.
人人都準備.
溪流下大雨.
山洪暴發.
人人都知道.
人人都害怕.
人人都會防備.
人人都不會走去.

$^{2281}$我們知道.
這個叫恐怕.
所以我們就要留心.
我們會不會有一個.
又用另一個例子.
有個水桶.
有少少漏水.
我就是.
我的房.
我有一樽.
洗頭水.
不知為什麼.
整個地下.
都是濕了.
原來它漏.
但你找不到有洞.
但你放在這裡.
過一會兒.
就有一噸洗頭水.
原來是漏.
找不到有什麼洞.
這裡告訴我們.
那個隨流可能就是.
你裝了一桶水.
聽完很開心.
很感動.
但下星期你回來.
原來漏了.
你都不知道.
為什麼會漏了.
因為你的水桶有個洞.
所以這裡告訴我們.
要越發恐怕.
要鄭重.
聽見的道.
要鄭重.
第二節就講道.
藉著天使所傳的話.
既是確定的.
即是在講摩西.

$^{2321}$頒布十誡.
凡是肝犯背役的.
都已經受了該受的報應.
就是那些以色列人在抗疫倒閉.
第三節.
我們如果忽略.
這麼大的救恩.
即是用舊約.
摩西和福音.
耶穌.
做一個強烈的比較.
你想想.
十戒人都害怕.
在西乃山.
但希伯來書的作者就說.
這個更大的救恩.
如果我們忽略了.
怎能逃罪呢.
接著開始講了.
這救恩起先是主.
即是耶穌基督親自講的.
後來聽見的人.
就是使徒.
給我們證實.
我們就是第三代的信徒.
即使徒已經死了.
現在來到我們新一代.
我們承繼福音的.
教父的年代.
承繼那些Big Bang.
這個時候也是Big Bang時期.
希伯來書尤其是強調.
猶太人.
對基督徒的Big Bang.
對信了主耶穌的猶太人的Big Bang.
希伯來書更加強調這個Big Bang.
所以.
告訴我們後來.
聽見的人.
給我們證實.

$^{2361}$第四節.
神又按自己的旨意.
留意這一節是很重要的.
神按祂的旨意.
按祂的時間.
按祂的計劃.
用神的旗幟.
和百般的異能.
你想不到的.
那些異能.
手帕也是異能.
彼得的影也是異能.
你可想而知.
神可以用.
不同的.
爐仔也會說話.
巴蘭的爐仔.
你可以想而知.
那個神蹟有多大呢.
我們有個姊妹說.
她家裡的狗會祈禱.
一吃飯就自己合手.
我不說是誰.
是否真的就很難說.
是一個條件反射.
每次都捉住它.
還是怎樣呢.
但巴蘭的爐真的說話.
是真的有意思的說話.
罵巴蘭.
那就是神蹟.
所以是百般的異能.
還有留意.
聖靈的恩賜.
和他們作見證.
主權在神那裡.
神按祂自己的旨意.
用的方式.
就是神蹟騎士.
那就是說.

$^{2401}$我們將這節聖經再清楚的演繹.
就是神用.
神用.
我的演繹就是.
神用各種的方法.
來到去.
特別是聖靈的恩賜.
和那些為耶穌.
作見證的人來到去同在.
那就是說.
當我們有.
信心或者有膽量.
向一些我們熟悉的人講福音.
其實那個已經是一個恩賜.
已經是神在我們.
裡面的內心的感動.
來到去做了.
當然神如果要這個時代.
例如可能一個新時代開始.
需要大量的神蹟.
那你和我就有機會.
和我們參與到.
那個神蹟的運作.
但記住.
不是讓你和我.
是高人一等.
出類拔萃.
重點是.
和他們作見證.
見證耶穌是.
真神.
所以有沒有神蹟.
那個是其次.
作見證.
才是那個真正的主軸.
今天我們就停在這裡.
天父.
我們很盼望.
我們看第三章.
和剛才的小道行傳十二章.

$^{2441}$我們能夠透過.
不同的時代.
神你自己.
是怎樣來到去用不同的方法.
在那個時代.
作供.
我們要問的.
這個時代是要怎樣去傳福音.
是不是透過.
高科技的GPT.
或是其他的.
仍然是用人.
用我們很卑微的人.
向我們的親戚朋友.
來作規勸.
友誼的報道.
或者我們要用很多的心血.
用很多的心思.
來慢慢的.
說服我們的朋友.
使他們對我們沒有抗拒.
對我們有信任.
以致當我們講福音的時候.
他們是很自然地.
會用心的聆聽.
和信納.
接受.
以致他們經歷神的醫治.
拯救.
如果這個時代是進入.
需要壁吧.
需要神蹟的時代.
讓我們有共同的.
認識.
我們的信心.
我們的內心謙卑能夠配合.
而不是我們.
興奮的.
只著重神蹟.
而忘記了神蹟背後.

$^{2481}$是為主耶穌作見證.
這件重要的事情.
我們仰望主.
禱告奉耶穌基督的名.
阿們.
consistency.
好啦 我再說我小意思.
辦法也是差不多.
助手可能有些訓練.
是需要第二個.
阿姨 Song.
那我就是用上一個.
風馬牛馬.
\newpage



\section{路加福音 10:38-42}
\label{sec:7t8JFmETTzs}
\textbf{2023-04-16 加恩樓《世界不能奪去的福份》路加福音10\_38-42 講員:蔣文忠傳道}
\newline
\newline
連結: \href{https://youtube.com/watch?v=7t8JFmETTzs}{\texttt{https://youtube.com/watch?v=7t8JFmETTzs}} ~~~~ 語音日期: 2023-04-18
\newline
\newline
\hyperref[sec:Pndq2oLIHnI]{\small{< < < PREV SERMON < < <}}
~
\hyperref[sec:index_chronic]{\small{[返順時目]}}
~
\hyperref[sec:index_scriptual]{\small{[返順卷目]}}
~
\hyperref[sec:cuPMYAewAe8]{\small{> > > NEXT SERMON > > >}}
\newline
\newline
路加福音 10:38-42
\newline
\begin{longtable}{cl}
\hline
\hline
章節 & 經文 (和合本修訂版)\\
\hline
10:38 & \begin{tabularx}{0.7\textwidth}{X} 他們繼續前行，耶穌進了一個村莊。有一個女人，名叫馬大，接他到自己家裡。 \end{tabularx} \\ \\ \relax
10:39 & \begin{tabularx}{0.7\textwidth}{X} 她有一個妹妹，名叫馬利亞，在主的腳前坐著聽他的道。 \end{tabularx} \\ \\ \relax
10:40 & \begin{tabularx}{0.7\textwidth}{X} 馬大伺候的事多，心裡忙亂，進前來，說：「主啊，我的妹妹留下我一個人伺候，你不在意嗎？請吩咐她來幫助我。」 \end{tabularx} \\ \\ \relax
10:41 & \begin{tabularx}{0.7\textwidth}{X} 主回答說：「馬大，馬大，你為許多的事操心煩惱， \end{tabularx} \\ \\ \relax
10:42 & \begin{tabularx}{0.7\textwidth}{X} 但是不可少的只有一件。馬利亞已經選擇了那上好的福分，是沒有人能從她奪去的。」 \end{tabularx} \\ \\
[1ex]
\hline
\hline
\end{longtable}
$^{1}$好,弟兄姊妹早安.
很感恩在大家中間好像都一段時間沒有跟大家有機會分享訊息.
或者各種安排等等.
今天能夠跟大家一起分享實在是上帝的恩典.
也都求主有足夠的力氣去跟大家分享今天的這篇訊息.
這個經文相信大家都很熟的.
我相信沒有什麼弟兄姊妹沒有讀過這篇經文.
可能你初信,你查經,你栽培都會讀到馬大和瑪利亞的這段經文.
不知道你去學習這段經文的道理的時候.
會不會都是我們其實要學習瑪利亞的洋式.
是不是?就是馬大很忙.
頭上有些資深的會友都知道.
這個都是一般我們理解的訊息.
我們就要放下馬大那種忙碌的洋式.
專心好像瑪利亞那樣在主面前聽.
可能大家都聽過這個道理.
但究竟這段經文其實想說什麼呢?.
這段經文其實是從一個什麼角度去理解.
希望今天從一個重新的角度跟大家分享這段經文.
你可以回想一會兒.
如果耶穌對著瑪大和瑪利亞.
叫瑪大你學學妹妹.
我們做家長有時候人性軟弱就會這樣.
學學弟弟乖一點.
學學妹妹,姐姐聽話一點.
其實這是一種分化.
如果說媽媽偏心或者爸爸這樣看我.
其實這種人性的教導我們應該戒掉.
每個孩子都有自己應該走的路.
我們不應該將人和人之間.
小朋友之間孩童成長的路上.
有一種榜樣式.
然後叫另一個去學習的模式.
這個應該不是聖經.
或者耶穌應該不會這樣去理解瑪大和瑪利亞之間的關係.
如果你這樣放下這個道理的時候.
我們可以重新讀一讀這段經文.
我想跟大家從一個服侍的團隊的角度來讀這段經文.
因為服侍的團隊裡面有每個人獨特的角色.
而也有一個團隊之間的那群人.

$^{41}$或者起碼是兩個人以上.
亞大,瑪利亞有兩個人.
他們有一個團隊合作默契之間.
其實應該怎樣彼此去配搭.
彼此去體會對方的存在呢?.
這個我希望在完結的時候.
跟大家去到一個這樣的進度去明白這段經文.
我想從一個……我們先想一想我們服侍的態度.
開始進入.
譬如我們每個人都會有一定的服侍.
可能是一些崗位上的服侍.
或者是身邊的鄰舍,家人等等各種.
有些有需要的人在身邊.
我把它們泛稱為別人的需要.
這是很正常的.
因為你有一種惻隱之心.
你看到別人的需要.
譬如街邊有個婆婆她過不了馬路.
你幫忙去扶她一把推過去.
這是很正常的反應.
我們應該勇於樂於去做這些事情.
這是一個別人的需要去回應.
一個需要能夠幫助到她.
但是大家要留意我們的服侍.
也不可能完全是看別人的需要去做.
因為世界上需要很多.
如果你遇到什麼需要你都去回應.
你都去做的時候.
你會發現你很快會burn out(消失).
你會發現自己力不從心.
你會發現自己做不完這麼多事情.
你記得耶穌自己也是.
當去到一個地方去醫病講鬼,行神蹟等等的時候.
然後那些人想留下耶穌.
耶穌你不如留下來繼續幫我們吧.
我們有很多需要.
整條村子有病的人都跑來找耶穌.
但是記得耶穌說.
我是以恆我猜我來者的事為念.
是負猜來這個世界.

$^{81}$他人生的一個推動力或者指引.
是來自天父給他的帶領.
而並非完全看人的需要.
所以就離開了村落.
再去到另一個地方傳讀.
又或者我們有自己的理想.
我們有很多dream.
這個也是很正常的.
我們覺得教會有些缺失,缺乏.
你看到某些需要.
然後就可能跟村民說.
不如讓我去嘗試一件事.
去做一件事.
希望能夠做到一些別人看不到的東西.
這個就是所謂的自己的理想.
個人的理想.
但是你要明白.
弟兄姊妹.
個人的理想不一定跟上帝的國相等.
有時候我們有很多未完成的夢.
可能你自己心裡有很多未滿足到.
想成就的事.
可能有時候我們分辨不到的時候.
就以為那個就是要服事的事.
這個是要重新我們不住地去判斷一下.
究竟我們想推動.
我們想引導他去做的事情.
那些服事.
是不是真的上帝給你自己.
或者是給教會的心意呢?.
所以對第三點那個.
我們跟上主的關係是非常重要.
我們需要常常都回到神的關係裡面.
大家都很熟悉.
我們愛因為神先愛我們.
我們的服事是出於一份領受主的愛.
然後以愛去回應.
我們是跟人分享上帝給你的愛.
所以我們無論是別人的需要.
或者自己的理想.

$^{121}$其實我們走一段路的時候.
我們都要重新檢視.
所以都要退休.
需要退醒.
問自己其實我們所做的事.
是不是走遠了呢?.
是不是越走越遠離開上帝更遠.
還是繼續維持在出於上帝的恩典.
出於上帝的愛.
而去有一個動力去作你的服事呢?.
這個是我們個人上服事.
我們可能有一些這樣的想法.
譬如我們是一個團隊.
或者一個教會的群體.
我們又會有什麼的考慮呢?.
當我列了這三點的時候.
不是全部的.
其實我只是點題.
引起大家的刺激思考.
有時候我們都會有一些目標導向.
我們每個教會都會有年提.
很多教會都會這樣做.
可能全團就會開會.
定下一個教會的方向.
然後就會和頂尖會分享.
希望整個教會未來一年都跟著這個方向走.
但是頂尖會的留心就是.
其實我們生命的服事.
也都不是純粹用數字.
用一些量化.
用一些統計.
用一些指標可以量得到的.
那些是一些參考.
但是卻不一定.
或者不能夠絕對.
就是我們用的數字.
用的統計數字.
或者用的走勢.
這些表達.
就等於我們做了我們的服事.

$^{161}$那個目標是很小心的.
生命的目標不是能夠量化得到.
或者我們很想會眾齊心.
或者我們很想眾志聖城一起做一些事.
表達一個合一的精神.
我們希望一起.
全個教會一起推動一些運動.
一些參與一些大型的聚會.
我們要留心就是.
其實有時候在一個教會群體當中.
都有不同的人.
可能其實有些頂尖姐妹.
未必會在一些很邊緣的位置.
或者在一些隱而未見的位置.
在那裡默默地服事.
所以我們都要小心.
其實在上帝國度當中.
每個人都很多元.
有時候我們想推一些合一的事情.
千萬不要忽略了.
其實不要成為那件事是唯一要做的事.
而忽略了很多頂尖姐妹.
其實她們都在做上帝國的事.
而她們未必能夠.
或者時間上也好.
或者各種原因也好.
做教會所推動.
希望能夠一起去做那件事.
這個都要大家.
有一種開放的心懷.
來去理解每一位頂尖姐妹所做的事.
或者最後一點我想說的是.
我們想立綱見形.
什麼意思呢?.
就是說我們很想很短時間就有結果.
很短時間就有成果.
譬如我們推一些提供一些資源.
人力.
希望過了暑假.
接著我們就可以有清晨的增長.

$^{201}$我們有些希望.
希望透過一些program.
一些很密集的活動.
就能夠帶來一些新的景象.
大家我們.
剛才我和鄧全濤聊天都說.
我們說到靈命的成長.
其實有時候我們不能急功近利.
或者不是一些可以即時.
可以見到有果效的事.
是一些心間細作.
長遠的工作.
或者這個更加.
在教會的成長內更加重要.
所以我們有時候也不能太快.
用我們有時候在社會裡面.
在工作裡面.
那種很想有一個果效.
很想有即時的結果.
那些指標量度.
來看教會的工作.
我講了這麼久.
和我們今天的經文有什麼關係呢.
今天的經文其實.
我想從一個團隊的一幅事.
和大家分享.
就是說你看馬大和馬尼亞的時候.
他有他自己.
好像我剛才第一張slide說的那件事.
他有自己個人的領袖.
有自己可以服侍的事情.
但亦都是一個配搭.
都是一個在服侍著這件事.
是一起去完成這件事.
我的背景就當是38到39節.
螢幕上那個其實是一個.
新漢語譯本來的.
和剛才主席讀的那個.
有少少不同的譯法.
不如我們一起讀.

$^{241}$38到39節在螢幕上的那個譯本.
預備 1 2 3.
他們繼續前行.
進了一個村莊.
有一個名叫馬大的女人接待他.
馬大有個妹妹.
叫馬尼亞坐在主的腳前聽他說話.
大家聽得出那個翻譯上的分別嗎.
聽得出來嗎.
馬尼亞在做什麼呢.
聽他說話.
如果記得的話怎麼翻譯呢.
大家都不太熟悉聖經吧.
聽他的道好像是.
聽道.
以為耶穌坐下來就講道了.
沒有其他東西要講了.
不會講生活習慣.
講其他的風花雪月.
講其他聊聊天.
在這個背景.
我就說是一個服侍的初心.
其實某程度上是馬大.
去主導的一件事.
其實整段經文38到42節.
基本上馬尼亞沒有出過聲.
如果你留心看一看經文的結構.
最後後面那半部.
其實都是馬大和耶穌的對話.
某程度上馬大是主角.
馬尼亞某程度上是配角.
基本上.
好像是一個背景或是一個在旁邊.
在參與.
但又不是很特別.
要回應他的情況.
留心看經文.
耶穌來到這個村莊.
大家都知道.
如果你看《讓福音》第十一章.

$^{281}$就是耶穌和馬大馬尼亞和拉薩路的家庭是很親密的.
當拉薩路死了.
其實那段經文就說了耶穌也哭了.
大家都很熟悉.
其實是很親密.
彼此的關係是很緊密的.
是很愛惜三姊弟.
所以大概你想像.
耶穌全都經過到馬大的村莊.
其實馬大一知道.
第一時間就接待耶穌.
你想像馬大和馬尼亞.
其實都有不同的性格.
可能你有兩姐妹.
可能你和你姐姐或妹妹.
也未必是同一個性格.
姐姐大概是一個.
好像是一個管家式的性格.
她讓你打點.
讓你安排.
妹妹你做這件事吧.
早就已經計劃好了.
是一個這樣很行動型的人.
馬尼亞說她是一個很善解人意.
或者是一個很善於聆聽人.
分享,做聆聽者.
在這裡陪伴別人.
有一種舒服的感覺.
所以她的角色就是.
在客廳裡陪伴耶穌聊天.
如果你這樣想不好聽.
將它用港道來看的時候.
其實是一個接待客人.
我們中國人也一樣.
有客人來到你.
你會做什麼呢?.
煮一頓好的東西吃.
你一定會煮一些他喜歡吃的東西.
可能你早就去了超級市場.
去了超級市場買東西.

$^{321}$又要買海鮮,又要買牛扒.
要準備很多不同的配菜,調味料.
各樣的飲品.
馬大大概就是這樣運籌為壓.
準備好一切.
預備接待耶穌.
然後就安排妹妹.
有一個角色.
我進了廚房工作的時候.
你不如在客廳裡陪伴主耶穌聊天.
讓她也覺得不會被冷落.
一種這樣的感覺.
想讓她煮有種賓至如歸的接待.
休息完之後重新得力.
讓耶穌再上路去服侍.
這個就是服侍的初心.
應該是很有默契的.
或者這麼說的時候.
其實馬大和瑪利亞是在做不同的角色.
和不同的服侍.
沒有說誰高一點,誰低一點.
誰重要一點,誰不重要一點.
你想像如果兩姐妹都衝進廚房.
留在耶穌的廳裡.
孤零零的時候.
也不是很受的.
其實這樣很沒人性.
怎麼會冷落客人呢?.
全部出來廳裡也不行.
其實是一種分工.
彼此各按其職.
在這個服侍裡去配搭.
或者他們兩個最初是有一個很清楚的目標.
都是很想服侍他們的主.
很想他們的主能夠.
在他們這次接待的服侍裡.
能夠重新得力.
有一個很飽足的一餐.
也可以跟他們分享.
有一種很滿足的團契.

$^{361}$這個就是他們最初很想能夠做到的事.
但是姐妹.
我想我們的服侍.
很少會從頭到尾都沒什麼難度.
有沒有人試過?.
由最初的計劃到完成這件事.
就如你的計劃,如你的想法就完成了.
其實我們很多時候都會遇到很多挑戰.
有些挑戰其實是我們控制不了的.
譬如你說本來有些探訪隊.
有幾個姐妹.
很齊心,很熱誠.
不如我們逢星期幾去做探訪吧.
但後來你知道.
每個星期都要去.
可能要做三個月至半年.
但中途可能又會有不同的情況出現.
突然間有些人說不好意思.
我自己近來可能要照顧家人了.
不好意思,我突然身體有點不舒服.
或者有些病突然出現了.
要一段時間休養.
然後整個事奉可能就不能夠貫徹始終地完成.
但這些是一些客觀,外在.
其實都與你無關.
但我們都要明白.
有時我們的規劃事奉都不是那麼容易.
逐步逐步上帝就大力.
所以我們每一步都要求主恩典.
但有更多時候,像馬大這個情況.
其實是一種自己的問題.
自己內心出現了一些變質.
大家可以留心一下.
到底馬大服侍的內心和初心.
是如何變化的呢?.
你見到我都特意將一個很重要的,很小的字.
去highlight出來.
其實這幾個經文的短短幾節.
有幾個but.
有幾個但是.

$^{401}$有幾個相反的情況.
就好像暗示給他們知道.
故事不是這樣就發生了.
有些事情正在發生.
有些事情正在轉變.
馬大卻為許多服侍客人的事憂心.
心裡煩躁.
他上前對耶穌說.
主啊,我妹妹撇下我一個服侍.
你不管嗎?請吩咐她來幫助我.
馬大衛了.
剛才說了他要做很多事.
那些各種服侍的材料.
煮菜.
各樣的服侍客人的.
要去餵奴他的食物等等.
令到他.
經文就說心裡煩躁.
這個心裡煩躁是很有意思的.
其實有兩個逆法.
第一個就是.
worry.
憂慮.
就是產生一種憂慮出來.
為了那些東西憂慮.
你可以想想有什麼會憂慮呢.
我都想到可能性就是.
覺得煮得不好吃.
怕主人不喜歡.
是嗎?.
就是一個完美主義者.
糟了,下錯了料.
或者時間不對.
然後憂慮可能牛扒太韌了.
煮得過頭了.
海鮮又蒸過頭了.
這些.
為了那件事想很完美.
想做得很好.
但卻是失手了.

$^{441}$憂慮那件事質素不好.
接著服侍出來的時候.
覺得主不開心.
不喜歡那件事.
這些大概可能是他憂慮的東西.
另一個字更精彩.
distracted.
distracted的意思就是.
分心.
分心就是說他本身是專心.
才有分心.
不知道你有沒有試過.
我經常都是這樣.
我專注力很弱.
拿著手機.
你試過沒有.
你本來想搜尋一件事的時候.
我希望大家不要像我一樣分心.
你滾來滾去的時候.
嘩,挺好看.
滾來滾去的時候.
大家笑,原來和我一樣.
不會專心.
不如記住.
你做這件事.
我經常都是有時.
嘩,不然浪費時間在這裡.
搜尋.
就帶來了你原初的初心.
做了另一些東西.
這個我想大家生活上.
每天都面對.
這個分心的挑戰.
但在馬大的情況就是.
他本來專心想服侍主.
這個是他最終的.
本來的目標.
但現在把那些東西分心了.
服侍了那些東西上.
想做好那些東西.

$^{481}$帶好那首歌.
本來大家一起敬拜上帝.
那首歌走音了,不好意思.
或者怎樣.
我希望.
這是另一個問題.
我們有沒有盡力去做.
但不是被那件事帶來.
我們本身想服侍主.
那個心懷意念.
所以當他內心.
出現了一種變化之後.
你不要看低.
你看看這件事都挺嚴重的.
接著,下面第二句.
他上前來抱怨.
你不用太氣的.
馬大的語氣應該很差的.
出來說話.
主,其實.
如果你看回英文.
就倒過來.
你不理我嗎?.
這是第一句.
你不理我嗎?.
妹妹丟下了我一個人在服侍.
這是整句.
有一點含沙射影.
他對著耶穌說.
又說了妹妹.
一箭雙鵰.
說了兩個.
請吩咐他來幫助我.
是不是?.
這個不是.
他原初的想法.
馬大因為服侍的事太多.
而感到煩躁不安.
剛才解釋了.
有一種憂慮.

$^{521}$有一種分心.
接著.
這種情緒令到他.
內心出現了一種.
對主和妹妹的不滿.
接著就是一種.
叫做受害者的心態.
Victimize.
可能大家聽過.
你們全都丟下我.
孤立了我.
離棄了我.
你們就很開心.
在外面.
然後就留下我.
孤單地服侍.
大家有沒有試過這種感覺?.
自我Victimize.
自我成為一個受害者.
好像別人.
每個人都針對著你.
其實不是這樣的.
有世界不是這樣發生.
但你主觀的情願.
馬大就是這樣.
沒有猜到.
最初大家都說定了.
你在廳裡.
我在廚房裡.
其實沒有產生任何變化.
是你自己的內心.
出現了這份變化.
然後最重要是第三點.
接著有一種價值觀的轉變.
一種價值觀的轉變.
本來剛才我說.
彼此分工.
沒有分高低的.
是吧?.
大家都幹奇蹟.

$^{561}$服侍著一起.
在配搭著.
現在就說.
我做的事是最重要.
其他人的工作就沒有那麼重要.
還不快點叫其他人來幫我做事?.
其他人都好像要繞著你的地球來轉.
放下手上的工作.
當我讀到這裡.
經常想起一個.
我自己服侍的見證.
不是見證.
是一個經歷.
我自己也有.
很多時間在校園裡.
在大學裡服侍.
其中一個經驗是怎樣呢?.
其實發生在一年多前.
也不遠.
那時候在疫症高峰的時候.
一群大學團體同學.
很想舉辦報道會.
他們也很勇敢.
其實也有很多挑戰和限制.
大學裡已經不准借一個.
借到場地.
但坐200人的地方.
只能坐100人.
或者坐很少人.
不讓他們坐滿場地.
所以就要找教會.
很感恩.
教會已經肯借一個.
坐150人的地方給他們.
很感恩.
我先說結果.
結果是很感恩的.
單從結果來看.
因為他們在星期二辦完一個報道會.
用歌劇形式.

$^{601}$他們很厲害.
不只是唱歌.
而是有劇集.
有不同的.
又唱又跳又說.
然後又會轉背景.
一個一個多小時的.
有唱和有劇情的表演.
訊息在那裡演繹出來.
透過同學自己.
而且找人也是他們自己找.
在IG自己分享出去.
鼓勵團友.
你一個帶兩個.
你可以帶你的中同.
大家知道什麼是中同.
中學同學.
也有帶同班同學也可以.
我有機會和他們其中一些人聊天.
證實了他們是其中一些同學帶來的.
像那個男生這樣說.
我從小就在教會學校.
上教會學校.
也聽很多福音的.
到中學的時候就開始沒有上教會了.
是誰呢?.
他說有一個月節目.
邀請我來.
我就來.
聽完之後也令他很觸動.
也令他重新想想信仰的問題.
這個也是很好.
啟發他再思考信仰.
已經是很成功.
而且不單止這樣.
他還要百五天很快滿了.
因為他要實名報.
然後再開一場星期六.
也滿了.
所以基本上整個Weekend的報道會.

$^{641}$有二三百個大學生參與.
而且有一半是真的.
未信或者已經離開信仰的年輕人.
是一個很難得的事.
他每年都會做.
所以大家也為他事公祈禱.
發生什麼事呢?.
那個事情就是.
可能你知道疫情有很多限制.
後來我記得.
那些負責音樂的弟兄姊妹.
在團體上呼籲.
有人幫我們忙.
因為我們人手不足.
因為需要壯丁.
來幫我們搬搬陀陀.
很多沙發和椅子都要搬.
我感受到當時他這樣呼籲的時候.
已經有很多人回應.
其實很好.
大家已經即時應對回應他的需要.
後來就發生了一些變質.
是怎樣變質呢?.
我發現聽到兩邊的投訴.
團契的職員.
有一點點對搞音樂的弟兄姊妹投訴.
搞音樂的弟兄姊妹.
也向團契的職員投訴.
是怎樣呢?.
團契的職員就覺得.
他們只是故意唱歌.
故意練習.
不太熟齡.
怎麼辦呢?.
他們不只是表演.
接著我聽完他這樣說的時候.
我的回應是什麼呢?.
我說你自己做功課.
講限期的時候.
你也沒有空祈禱.

$^{681}$你講1159的時候.
你也在衝功課.
你也不會祈禱.
你幫他祈禱吧.
不要等了.
你為他的熟齡生命祈禱吧.
我這樣提醒他.
你幫他承擔他的需要.
第二就是那邊又說.
我們又不夠時間.
又怕到時會失去動作.
又怕走音.
又怕記不到台詞.
接著我跟他們負責人說.
其實不要緊.
我們只要將最好的獻上.
上帝在整個聚會.
最重要的是那些新的朋友.
他們透過你們的服侍.
能夠被觸動.
能夠重新對信仰產生興趣.
我們將一切我們做的事.
仰望上帝.
做一個中間人.
兩邊調和一下.
其實最後慢慢消化了這件事.
我感受到沒有怎樣將它發酵成為一個很大的傷口.
傷害或分裂.
是很感恩的.
你想像一下.
如果馬大和馬尼亞的小小的割裂.
源自於馬大.
我不知道大家有沒有經歷過一些教會分裂的情況.
如果耶穌不去處理的話.
就很大件事了.
那個破口會繼續發下去.
好像癌症一樣一直發下去.
不會就這樣解決的.
所以到最後一段.
你會看到耶穌的回應是多麼有意思.

$^{721}$寫一節.
主卻對他說.
馬大馬大你為許多事憂心煩惱.
但只有一件事是必須的.
馬尼亞已經選擇了最好的一份.
是每個人從他那裡奪去的.
留心我有個卻字在這裡.
不過卻字其實可惜.
因為新漢日本.
和本都沒有翻譯.
其實我自己看原文翻譯出來.
那個字其實是有連接詞在這裡的.
不過有時候中文聖經.
可能我懶一點.
就不會把很仔細的每個字都翻譯出來.
其實有一個相反的意思.
什麼意思呢.
就是說馬大和耶穌投訴.
好像我們打投訴熱線一樣.
我們經常打.
服務不開心就打投訴熱線.
你想想你電話那個人.
很少會罵你.
不會說是是明白明白.
我會跟進這樣.
多數都是順你的意思.
但耶穌不是這些接線生.
耶穌卻回應他.
就是耶穌要他明白一些東西.
沒有順著馬大的請求.
來去馬上call.
不好意思我馬上叫妹妹來幫你.
沒有這樣做.
這個就是耶穌即時介入.
這個問題裡面.
避免兩姐妹之間的紛爭.
或者那個破口.
越來越嚴重.
你看耶穌怎麼說.
馬大馬大.

$^{761}$你為許多事憂心煩惱.
我經常覺得來到這裡.
耶穌就像一個輔導員.
大家有沒有讀過輔導.
輔導員其中一個最重要的就是.
將對方的情況告訴你.
其實你是這樣.
馬大說我這樣很慘.
其實你是這樣.
將真實的情況告訴你.
馬大將一個很語重心長.
好像要他清醒.
要他回來.
不要在自己的世界裡轉.
回到真實的世界當中.
面對你自己內心出現了變質的情況.
你裡面為了許多事憂心煩惱.
要明白自己在這種狀態裡面.
這是耶穌即時和馬大分享.
要他明白.
我不是按照你的意思去要求去做.
有時候我們也會.
我們內心有很多雜亂.
如果我們做了馬大很大的事.
求主的時候.
主不聽.
其實不是想不理我們.
而是我們在想的不是主的心意.
都不理.
接著再有另一個彈字.
這個彈字是清楚一點.
因為很清楚.
他為許多事憂心.
但只有一件事.
是許多事和一件事的對比.
這個就很清楚.
你做了許多事.
忙了許多事.
但其實只有一件事.
你捉緊一件事就可以了.

$^{801}$就夠了.
接著來到我.
我開頭說的.
是不是馬地亞那件事呢.
大家留心一下如何演繹.
馬地亞已經選擇了.
最好的一份.
是沒有人能從他那裡奪去的.
我要先解釋一下這一句經文.
首先不好意思.
其實你看到.
新漢語又是最好的一份.
和本就應該是上好的福份.
所以詩歌唱最好的福份.
但其實我題目也說福份.
其實是有少少誤導大家.
因為其實經文是沒有福字的.
是個Portion.
那個字是The Good Portion.
你可以譯作那好的份.
Portion是沒有祝福的意思.
是那一份.
是好的一份.
好的一份即是不是所謂最好.
是Best沒有這個意思.
也不是Better.
不是好一點.
是好.
是對於馬地亞來說是好的一份.
什麼叫一份呢.
這個字在舊約最多應用的時候.
就是在一些婚姻遺產的時候.
譬如一個族長年老了過世.
然後他的子子孫孫排隊來婚家產.
大家可能聽過舊約裡面.
長子分雙份.
然後其他兒子也有一份.
女兒也有.
女兒有半份.
你不要搶了女兒的份量.

$^{841}$因為你是長子.
你有權勢.
就奪去其他的子孫孫女性.
那些本身應該得到的那一份.
所以不能被奪去.
意思就是這樣.
那個遺產得到.
你是一個得到遺產的繼承者.
你是那個族長本身預設了給你的時候.
沒有其他人.
這個世界都不能夠被奪走.
所以那句其實是不能被奪去.
是一個解釋那樣的本質.
無論任何事情.
那是他的.
他要去預計那件事.
你用多強勢.
你用多怎樣的身份也好.
我不可以被奪走.
所以耶穌是守護著瑪利亞那一份.
不知道大家明白嗎?.
他不能被奪走.
瑪利亞已經選擇了.
他已經選擇了.
所以當放回今天這段經文裡面.
這個舊約的婚遺產的意思來到.
新約瑪利亞的經歷就是.
他選擇了.
不是分給他.
是他選擇了.
而且都是屬於他的.
還有是其他人拿不走的.
這一份.
這個就是他新定的那一份.
我們經常說的那種照命的意思.
屬於他的.
在主角度當中.
屬於他的位置.
需要承擔.
要做的事.

$^{881}$沒有人能夠拿走.
所以當你這樣去演繹的時候.
那馬大利呢?.
是不是?.
那耶穌可能就再加一句.
馬大利自己那一份是什麼呢?.
是.
馬大利應該.
是不是應該把握回.
同樣.
你應該選擇回.
好像瑪利亞那樣選擇回.
你是本來屬於你.
不能被奪去那一份呢?.
這個大概猜是.
耶穌來到這裡.
告訴你.
瑪利亞做了他要做的事.
這是一個例子.
不好意思.
他不能夠轉移給你.
你呢?.
那你應該找回你在上帝國當中.
給你的那一件事.
經文沒有在這裡.
經文在這裡就完結了.
有時候讀聖經故事.
都挺有趣.
就是.
我叫做懸空了它.
什麼意思呢?.
沒有一個大團圓結局去解釋.
譬如你讀約拿書.
最後約拿怒爆爆走了.
是不是大兒子的故事.
他有沒有回去抱著小兒子呢?.
你不知道的.
就是他留給讀者.
你代入去的時候.
你代入大兒子的身份.

$^{921}$你代入去約拿的時候.
你代入去馬大的時候.
你問回自己.
你是怎樣呢?.
那我也給那些.
我估計的答案給大家知道.
我估計.
有兩個可能性.
馬大聽.
馬大不聽.
是不是?.
我們都是一樣.
聽話不聽話的.
你聽耶穌.
順服聽回耶穌那句話.
不聽話繼續.
很野蠻都可以.
我當他聽.
我當他聽的話.
我估計他會回到廚房裡面.
做回一個稱職.
和重新定焦點的一個廚師.
這個是他的本性.
因為他是一個專長.
因為他是一個聖像.
他做回自己.
但是重新讓自己生命的焦點對著.
這個我相信.
是馬大應該要選擇的道路.
照命是什麼呢?.
我很喜歡一句名句.
我太太經常引用這句話.
就是布克立的一句對照命的意思.
上帝呼召你到一個地方.
就是你那極大的喜樂.
和世界那深切的渴求.
兩者相遇的地方.
大家可以去看一看那句句子.
有兩個層面.
一個是你裡面很想有的東西.

$^{961}$這個要很小心.
剛才說的可能是個人理想.
你要分辨得到.
你的想法和渴望.
是不是從神而來的呢?.
是有一個重疊的意味.
你會得到滿足感.
你會覺得開心.
做了這件事.
其實是在服侍的.
但同時後半句很重要.
後半句不是叫你拯救地球.
不是叫你將整個溫室效應的問題解決.
其實是解決不了的.
你不可以出來做環保戰士.
就將溫度降低.
其實是不行的.
但意思是.
這個照命不是一個內在的問題.
是一個在世界當中.
外顯的事情.
這個很有意思.
就是說我們每一個信徒.
實踐上帝給我的照命.
是在世界當中有位置的.
可能傳遞人就在教會裡的模樣.
但大家在自己的光位裡.
有自己的位置.
你那裡是在實踐最深切的渴望.
和最大的需要.
那是怎樣呢?.
就是說如果你真的做好.
上帝給你的那一份的時候.
其實就是對世界最大的祝福.
不要看低自己.
我不是做一個二打六.
我做一個文員.
我又不是CEO.
又不是什麼.
其實不一定的.

$^{1001}$你知道一個公司這麼說.
上上下下都有自己的角色.
每一個人.
如果你做好那一份.
其實對於整個團隊.
對於文化.
對於效益.
是很大影響力的.
我真的越來越相信這件事.
所以每一個人都不要輕看.
自己的角色的位置.
做好.
這個就是你最大的滿足.
也是對世界最大的祝福.
和最大的帶來上帝的喜樂在裡面.
所以.
照明我們人生最大的渴望.
對象是我們的主.
在我們國當中有份參與.
是我們生命裡面最大的喜樂和滿足.
好像馬大和瑪麗一樣.
我們和弟兄姊妹一起去服侍的時候.
彼此配搭.
當我們的焦點都是上主的時候.
人和人之間的比較和疾度.
就得以消解.
這個也是我相信的一個信念.
馬大出生的時候.
出現了一種比較.
出現了一種疾度.
出現了一種變質.
問題就像剛才的情況.
不是在瑪麗亞身上.
而是在馬大自己身上.
如果調教自己的焦點.
在主的時候.
整件事就很不同.
他在看著妹妹的時候.
他懂得欣賞妹妹做的事情.
他在做一些我做不到的事情.

$^{1041}$他在做一些幫我和我配搭的事情.
他和我配搭的事情.
他和我一起.
一家等於三.
整個服侍裡面.
產生了一種很美麗的效應.
這是當我們每一個一起合作的弟兄姊妹.
都對準自己的焦點的時候.
都是把著上帝的照明的時候.
一切我相信很多的紛爭.
很多人和人之間的問題.
其實你都會成熟地去疏解它.
因為你不會被那件事困擾著你.
你都不會被那些人和人的溝通上面.
蓋過了你體會對方.
在上帝國度當中.
他把握著那份參與的角色.
所以一切的問題.
其實可以消解得到.
我最後想用一個見證.
去結束我今天的講道.
一間教會的目的就是認識他.
其實是一間很小的教會.
就是他.
最初的時候.
我不記得怎樣走出來.
總之有大概二十位弟兄姊妹.
出來.
成立一間很獨立.
很小的教會.
可能之前有些問題之類的.
於是他們就租了一個地方.
大約八百呎在深水Po 區.
就開始聚會.
開始教會的運作.
然後目前帶領他們的弟兄姊妹.
就問了.
我們租了一個地方.
租了一個地方用了.
然後可以崇拜.

$^{1081}$一群弟兄姊妹一起.
去找一些事奉的機會去服侍.
但我們人手這麼有限.
我們不如一起祈禱.
一起去尋找一個獨特的群體.
我們就用兩年時間.
一起去服侍這群人.
不要把資源打散.
也不要把大家互相.
在這裡消磨大家的.
那種服侍.
一起同心去服侍這個群體.
有些人就在前線做.
有些人就在支持等等.
於是他們就決定.
這兩年時間.
就去做青少年工作.
然後就去印第安.
因為他們租的地方.
剛好.
就是隔壁那間.
非基督教的中學.
有間中學.
想做青少年.
都要去拍門.
看看有沒有機會.
做一些計劃.
可能是全年的計劃.
希望做一些功課輔導.
或者周末的時候.
帶一些學生去遊歷一下.
或者一些交遊.
遊戲的時間.
原來.
你不找也不知道.
一找的時候.
校長是基督徒.
於是當牧志仔去找校長的時候.
校長就很快答應了他.
讓他試行這兩年計劃.

$^{1121}$就開始他們教會的第一個.
一群二十人一起.
同心去做的工作.
他們有一個.
很有趣的.
互相的共識.
可能聽了你都會嚇一跳.
他們說他們這群人的共識就是.
在這件事裡面.
一起都要守這個原則.
是甚麼來的呢.
他們說我們無論如何.
都不可以主動.
跟這群青少年說.
叫他們回教會.
有沒有這麼笨.
你既然叫他們.
當然是為了教會的增長.
很多人都是這樣想的.
為甚麼牧志仔會這樣想呢.
他覺得是想避免大家很功利地.
純粹是為了拗青少年回教會.
而不是一種全心全意的服事.
所以就大家守一個原則.
不准說.
你全心全意去做.
去服事他們.
用你的愛心去關懷他們.
於是大概一年之後.
很感恩.
其中一個青少年主動問他們.
你們是不是有教會背景的.
他不是我們.
不是自己去拗他們.
於是就從那一刻開始.
就跟他們分享.
我們在隔壁有一個地方.
其實有間甚麼教會.
其實是甚麼時候有崇拜的.
這群大概中二的.

$^{1161}$十多個學生.
成為了教會第一個青少年團體.
十年之後.
有些職稱.
成為了教會第一批青年的執事.
到今天教會由二十多人.
到六七十人.
不是很多人.
但你看到他們能夠在主裡面.
很好的選擇.
和把握他們的一份.
一起的團隊.
一起的服事.
將祝福,將世界的需要.
能夠可以將.
透過他們去滿足得到.
去實踐得到.
傳達人說過一句說話.
現在是牧師.
他說在這個時候漢始甚麼都做不到的時候.
教賽總的比喻.
鼓勵著我們.
上帝的大能和要使一切.
為他所做的工作成就.
在這隻漢始短缺.
人手又沒有,錢又沒有,資源又少.
的情況之下.
懷著堅定的信心.
憑著單純,渴望.
去關愛論社的初心.
與主同行,上帝必然.
祝福,這是他自己很深切.
體會到主對他的.
祝福.
等之會怕我們都不要輕看.
我們內心.
那一份.
我們學習瑪利亞.
好好地去把握.
可能有很多壓力.

$^{1201}$可能好像馬大會,怎麼搞的.
你不來幫我.
但是他很清醒.
很清楚自己的角色是甚麼.
但馬大雖然迷失了.
他都可以透過主的呼喚.
重新.
回轉.
回到他的正途.
這張畫.
我都介紹書給大家.
這張畫來自一本書叫《獨著》.
一位作者將自己.
《物想經文》裡面.
想像得到的一個情境.
他畫下來.
這個情境是他想像自己.
像一個少女向著夕陽奔跑.
他覺得自己是一個.
擁抱著自己服侍的初心.
的時候.
很有.
目標.
很有衝勁.
很清楚自己走的方向.
然後他在《物想經文》裡.
我都截了出來.
我們今天的狀態.
又去到哪裡呢.
會不會離開了這個情境.
很遠呢?會不會在我們初相的時候.
有這個情境?或者是曾經試過.
很熱心或者是很知道自己.
如何為主服侍.
但到今天好像已經.
沒有動力.
或者失去了一份盼望呢?.
作者問我們會不會.
祈禱?會不會想想我們.
有甚麼行動重新開始.

$^{1241}$可以讓我們重拾.
這份初心呢?我想有些時間.
大家安靜.
先為自己祈禱.
然後帶大家一個禱告結束.
低頭禱告.
七夕多謝你讓我們今天藉著.
馬大和瑪麗亞服侍.
你的這個故事.
讓我們重新去思想.
去生命的.
一份.
我們.
在這個國度當中.
你為我們預備了.
但它是要我們.
好好地去選擇.
好好地去把握.
若然我們給很多的事情.
分心,憂慮.
願你幫助.
願你好像挽回.
馬大一樣去挽回我們.
願我們再在你的國當中.
尋回一個適合的位置.
也看到.
你在你的國當中.
如此的豐富.
有弟兄姊妹一起同行.
一起服侍.
為天國,為上帝.
你自己有需要地上的人.
我們.
能夠.
可以帶著一份.
很有喜樂,很有力量.
去參與,去回應.
就好像我剛才讀的那個.
是一個內心.
一個很大的滿足.

$^{1281}$也祝福到這個世界.
主祈求你幫助眾弟兄姊妹.
若然離開了.
我們這份初心的時候.
祈求你讓我們今天禱告立志.
能夠透過我們的.
選擇,透過我們.
向你的那份堅持.
可以尋回.
馬大應該尋回的那一份.
也是我們自己生命應該尋回的那一份.
多謝你.
求你繼續帶領我們.
隨聽我們祈禱奉主名求.
我這個很簡單的介紹.
我剛才說讀著的那本書.
就是我這裡.
右上角那本.
是一本真的不錯的.
適合成人去拿著來靈修.
和物相的書.
我們FS Press也出了.
有聆聽卡.
其實是一個詩篇.
可以作為一個卡遊戲.
給弟兄姊妹在小組運用.
透過圖像去分享生命的故事.
原來我有些詞彙.
讓你去.
認識自己的情感.
同工也有心設計.
每個詞彙背後.
都有一些經文詩篇.
希望我們透過詩篇的經文.
去給自己的服侍.
另外我們也有青少年的工作.
不知道大家有沒有留意.
有一個叫不靈修的IG.
如果沒有的話可以追蹤一下.
其實我們每天都有一篇貼文.

$^{1321}$給青少年透過一個.
靈性修煉.
習慣地每天都去.
讀上帝的話和發覺上帝的話.
和他們生活上很有關係.
所以歡迎大家去追蹤.
這個不靈修的IG.
也可以介紹給你認識的青少年.
都希望建立一個.
每天都親近上主.
和他覺得是一種貼地的.
生命的問題.
是相關的.
上帝的話原來是很有意思.
和對他們很有對應的幫助.
祝福大家.
我們所做的工作.
好,我將時間交給Ada.
\newpage



\section{5:1-14}
\label{sec:cuPMYAewAe8}
\textbf{2023-04-16 午堂《敬拜那又能又配的羔羊》啟示錄5\_1-14 講員:陳麗蟬牧師}
\newline
\newline
連結: \href{https://youtube.com/watch?v=cuPMYAewAe8}{\texttt{https://youtube.com/watch?v=cuPMYAewAe8}} ~~~~ 語音日期: 2023-04-18
\newline
\newline
\hyperref[sec:7t8JFmETTzs]{\small{< < < PREV SERMON < < <}}
~
\hyperref[sec:index_chronic]{\small{[返順時目]}}
~
\hyperref[sec:index_scriptual]{\small{[返順卷目]}}
~
\hyperref[sec:Yd_an_Ish1s]{\small{> > > NEXT SERMON > > >}}
\newline
\newline
$^{1}$各位姐妹平安.
在二十多年前.
我參加了一個旅行.
叫做「阿西亞訪韓盛會」七天.
去到韓國有七天的時間.
不過其中的四天.
我在韓國的祈禱山禁食了四天.
到星期天的時候.
我們就去到韓國當時最大的教會.
我想現在也是.
也曾經成為世界最大的教會.
這個羽爾島的純福音教會.
在當時的教會.
它應該的崇拜人數.
差不多接近八十萬人.
很大的一個地方.
這個在網上拿出來的.
因為二十多年前都沒有手機.
沒有這些東西.
當中的外形就好像一雙手.
一雙手一同舉手.
向我們的主耶穌基督去敬拜.
去歌頌.
他們每堂的崇拜有萬多二萬人.
所以早上七點鐘就有崇拜.
不過星期五晚的通宵祈禱會.
我們都去了一晚.
沒有全部去.
去了三個小時.
很厲害很多人.
他們在崇拜的時候.
他們的領唱.
在台上的領唱.
有二十四個人.
穿上整齊的西裝.
打了領帶.
剛剛我們師班都唱得非常震撼.
但是他們有二百多三百人.
是一個非常非常震撼的場景.
在當天參與崇拜.

$^{41}$我就心想將來在天上高.
我們的敬拜都是這樣.
因為在祈禱山那四天.
我們早上六點鐘起床.
七點鐘就去到大會.
開始有一個多小時的唱詩.
然後有一個多小時的聽道.
然後又有一個多小時的個別祈禱.
或者分組祈禱.
連續四天.
最初聽都覺得.
不知道能不能熬下去.
因為很悶.
但是.
不過發了四天之後.
我們覺得不夠厚.
因為你覺得越唱詩.
越聽道.
越一起祈禱的裡面.
你會覺得整個人.
真的很興奮.
到主日的時候.
趙容基牧師.
當時的堂主任.
就是趙容基牧師.
他講道就是講詩篇23篇.
今天我忘記他講什麼了.
不過在當天就記得.
很紮心.
在我還沒有去旅行團的時候.
在網上都介紹過.
這個牧師.
他每一年都會去到.
一些很窮困.
落後的地方.
去傳福音.
去為人祈禱.
他去到每一處.
都很多人跟著他.
整個好像明星款式.

$^{81}$他一祈禱.
那些人就得到醫治.
拍攝那些機器.
就拍攝他.
他一為人祈禱的時候.
原本那一面.
很多很多的.
不知道濕疹還是什麼.
立即祈禱就光滑了.
低的那些又立即走得了.
很欣賞這位牧師.
我也覺得.
他說那些靈力.
那些氣場很厲害.
但是.
在2014年的時候.
新聞有這個報導.
調查到這個牧師.
在2002年的時候.
他利用他的職權.
他是一間這麼大的教會裡面的主任牧師.
他就去要教會.
用一個高於當時的市價.
去買入他大仔公司的一些股份.
令到教會.
蝕了1200萬的美金.
一個很龐大的數字.
結果這位牧師.
被罰坐牢三年.
就緩刑五年.
也都要罰款500萬美金.
當我看到這段新聞.
那段時間都有點影響.
不過很快就禁了它.
當人在敬拜的焦點.
就和為我們贖罪的主耶穌.
這是我們敬拜的焦點.
但當你將這個焦點轉移了.
轉移了一個自己的權位.
轉移到金錢的賺取.

$^{121}$為了自己家庭的金錢賺取的時候.
上帝一定來審判.
在啟示錄第四第五章.
是天上一個很完整的敬拜場面.
在當中就告訴我們敬拜的對象.
在第四章.
在第一部分.
就是敬拜那位坐寶座上的.
是創天造地的上帝.
他是萬有的主宰.
到第五章的時候.
就是這個羔羊.
約翰就看到在這個寶座上.
他的右手就拿著一卷書卷.
就說到這卷書卷正面反面都寫了字.
很嚴謹的就有七個印.
我們現在叫做Label紙.
你將紙捲起來了.
那怎樣可以黏著呢.
就用七個印.
就是說非常秘密的書卷.
當時第五章就描述到約翰.
他看到這麼重要的書卷.
他又不知道他怎麼知道.
都是關於人類將來的命運.
將來的歷史.
他覺得這麼重要的書卷.
沒有人能夠打開.
將來的歷史是怎樣呢.
他就緊張到哭起來了.
他非常緊張國度的事情.
教會的事情.
我們都需要學校的榜樣.
我們都需要教會.
緊張我們教會的事情.
因為這個就是一個屬靈的家.
當時大力的天使就在那裡說.
沒有人能夠打開這個書卷.
約翰就哭了.
他哭了就說.

$^{161}$因為這個羅馬皇帝不能夠打開.
這些祭司長 文士 法律蔡人.
甚至現在的總統 總理 特首.
都不能夠打開.
約翰很緊張.
因為當時看到教會的情況.
基督徒被逼迫到不得了.
有些人就被扔到鬥獸場.
一定輸給獅子的.
又或者有些人被扔到油鍋裡.
燒滾油鍋來炸死他.
也有些人被釘在十字架上.
所以當時整個警方.
那些人心 特別是基督徒.
很害怕 很害怕.
約翰看到這個警方.
就說 糟了 這些逼迫.
去到什麼時候.
之前就有尼祿王.
現在又有第二個皇帝.
所以他又知道這個書卷.
裡面寫著的.
都是人類將來的歷史.
他很盼望能夠有人來打開.
但是當天使說.
沒有人能夠打開.
他覺得 糟了.
上帝這個救贖計劃.
那怎麼辦.
人類是不是絕望呢.
是不是慘到永遠呢.
感恩.
在當中有一位長老.
他帶給我們一個很大的安慰.
他說 誰能夠揭開這個書卷呢.
他說 不要哭 看啊.
猶大之派中的獅子.
大衛的根 他已經得勝了.
能夠剪開這個書卷.
能夠揭開這個七印.

$^{201}$當長老宣佈了之後.
這個四活物和二十四位長老.
就俯伏敬拜.
當耶穌基督一拿起這個書卷.
即是說 當他拿了的時候.
即是說 他將會可以打開這個書卷.
原來當羔羊介入了人類的歷史.
介入了我們生命的時候.
可以將一個悲傷的境況.
轉化為歡樂.
所以四活物和二十四位長老.
就俯伏敬拜讚美.
這個羔羊為甚麼能呢.
因為他是猶大之派中的獅子.
這個名稱是源自於創世記第四十九章.
當時這個雅各.
他在臨終前為了他的十二個兒子去祝福.
祝福到猶大的時候.
他說 猶大啊 你的兄弟必讚美你.
你的手必勒住你仇敵的頸銜.
你父親的兒子要向你下拜.
就說了這個猶大之派裡面.
是一個領導的角色.
接著他繼續說這個祝福.
猶大是隻小獅子.
我兒啊 你要捕獲獵物就上去.
他進服他討餓.
如公獅猶如母獅誰敢惹他呢.
是描述到那些獅子.
原本就在山上居住.
但當他肚餓了 要吃東西.
他就下到平原那裡.
他就要看準一些獵物.
他就會.
有一點點像這隻.
他就會整個人向前一點.
看著哪裡有些動物可以殺的時候.
就撲前他.
接著他說他進服他討餓.
吃飽了.

$^{241}$他就睡在那裡休息.
不過無論他準備去獵這個食物.
又或者去討餓休息.
又無論是公獅還是母獅.
不過當你看見有隻獅子在你前面的時候.
你敢不敢走過去.
撩一下他 撼一下他呢.
沒有人夠膽.
他就說到猶大派裡面.
將會出一位君王.
他滿有威嚴和權柄.
這位羔羊就像獅子一樣.
更有威勢.
這就是耶穌基督.
另外這位羔羊又是誰呢.
他說他是大衛的根.
這也是出自《二賽書》第十一章的描述.
這個根就是根源.
這裡是說到耶西的殘幹必長出亂芝.
耶西就是大衛的爸爸.
他這個平民百姓沒有什麼特別.
但說這個平民百姓沒有什麼特別.
好像殘幹一樣沒有什麼用.
但竟然可以長出亂芝.
他有生命的.
並且他將會能夠結出果子.
有果實的.
他非常有用的.
然後他繼續說.
耶和華的靈必住在他身上.
就是智慧和聰明的靈.
謬略和能力的靈.
知識和敬畏耶和華的靈.
這個大衛的謬略.
他的後代是很有智慧.
很有聰明.
很有謀略的.
並且是一個公義和平的統治者.
在彌賽亞.
在大衛的後裔.

$^{281}$這個君王.
在他的統治底下.
是充滿了和諧.
安詳.
力量和喜樂的.
如果讓我們看到.
在我們的生命裡面.
只要有耶穌基督在我們裡面的時候.
我們理應就有耶穌基督那份和諧.
喜樂和力量的.
還有約翰他自己看到了.
那隻羔羊是站著的.
不過好像被殺的.
所以讓我們知道.
這個就是直指耶穌基督.
他有七個角.
七隻眼睛.
就是上帝的七靈.
代表了羔羊.
他是充滿權柄.
充滿能力.
充滿智慧的.
他有七隻眼.
七這個字在聖經裡面代表了完全.
所以說.
他這個洞察力.
觀察力.
是很清楚.
很清楚.
沒有任何一些不好的東西.
或者不乾淨的東西.
可以逃避過去那隻眼睛.
這隻羊.
是非常威嚴.
非常有能力.
耶穌基督能.
能夠打開這個書卷.
展開這個人類命運的前途.
因為他是.
充滿權柄.

$^{321}$尊貴.
智慧.
能力.
他也是掌管著人類的歷史.
掌管著你和我人生的每一步的道路.
剛剛出版的那個中信.
有一篇文章.
是說到有一位威爾斯裡面的女醫生.
她說到她小時候讀書已經是非常厲害了.
所以當她讀醫的時候.
對她來說不是那麼困難.
有一次.
她要為一個紅斑狼瘡的女孩子去做手術.
她非常有信心.
她一定可以的.
所以她就對這個女孩子說.
不用怕.
手術之後見我們.
但很可惜.
在那個手術的裡面.
這個女孩子就在手術的途中離世了.
她之後反省.
她說她只是一個醫生.
她對這個醫學上升很多很多很多.
不過她不是一個掌管生死的上帝.
所以她之前那麼大膽跟那個女孩子這樣說.
她覺得她驕傲了.
她做了一段時間之後.
她就去到大學那裡去做教書.
做副教授.
很快之後她就升了做正教授.
原本是一個合約的員工.
逐年簽約的.
現在有些教書都是逐年簽約的.
弟兄就告訴我.
很煩的.
做了一年.
做了半天.
不知道下一年幫不幫你簽約.
很沒有穩定.

$^{361}$這位女醫生同樣都是.
她就去申請轉為長工.
怎知道人事部不知道為什麼.
將她的申請信放在正教授的文件櫃裡.
她就想一下.
剛剛升做正教授一年.
現在立即跟其他正教授爭.
她想著自己一定不行.
死定了.
她怕不單止不行.
還怕可能當那個約滿了一年之後.
她不幫你續約.
那就沒得做了.
但她說.
我又很喜歡教書.
她一直切切的祈禱.
這次她真的切切的祈禱.
但上帝一直都沒有回答她.
她在裡面就反省.
反省她自己在過去的日子裡面.
她為了她自己那份聰明伶俐去驕傲.
過往她在醫病.
在教學.
在很多病理的研究裡面.
她時時只是去榮耀自己.
就是想表揚自己是多麼的聰明.
但她從來沒有想過.
要藉著她這個醫學.
藉著她這個工作的成就.
去榮耀上帝.
當她一直安靜安靜的時候.
她終於明白.
上帝要在這段時間.
要將她生命裡面的垃圾拿走.
將她生命裡面.
經常都要自己是最聰明的.
自己是聰明伶俐的.
要給很多人去讚賞的.
她說.
她要將她自己.

$^{401}$要榮耀自己.
這些驕傲的光環.
拿走它.
賜給她生命的氣息.
賜給她有智慧.
賜給她讀書可以這麼聰明.
以致她可以在這個崗位裡面去工作.
她上通之後.
她告訴上帝.
她以後都不需要這些外來的光環.
不需要靠周圍的人去讚她.
去肯定她.
各位電影姐妹.
我們的上帝.
在充滿能力.
充滿尊貴權能.
掌管一切.
掌管歷史.
掌管人生.
所以將我們那個的生命.
我們每一樣事情.
交給祂.
並且記得.
要將榮耀.
頌讚.
都歸還給祂.
因為祂是一位萬事都能的上帝.
第二方面.
我們看.
好了.
我們看完這兩幅漫畫.
這位牧師.
蔡頌陽牧師.
他畫了很多幅.
很多幅的.
淑寧漫畫.
給我們很多鼓勵.
在左手邊那一幅.
他就說.
人生有很多的事情.

$^{441}$當經歷的時候.
很慘很慘.
就好像約翰那樣.
打開這個書卷.
很慘很慘.
但是.
他過後.
過後的時候.
他看到.
原來上帝的手.
在當中來去引領著.
幫助著.
一會兒我們就會說了.
好了.
另外一幅.
在我右手邊那一幅.
他就有些事情.
我們覺得我們被人傷害了.
我們身體或者心靈都受傷.
但是.
我們的生命.
如果是這樣下去的時候.
就會很頹廢.
但是.
靠著這位大能的耶穌基督.
我們可以變作剛強的人.
第二方面.
我們看這位耶穌基督.
是怎樣去配受領導者的.
當這個長老宣佈了.
這個高揚又拿了這個書卷.
這個四活物.
和24位長老.
就俯伏敬拜了.
他就說了.
向你配拿書卷.
配揭開七印.
為甚麼呢.
因為你曾被殺.
用自己的血.

$^{481}$將各支派.
各語言.
各民族.
各幫國中賣了人來.
使他們歸於上帝.
這個是第一個師班.
四活物和24位長老.
常樂師班.
接著就是這個千千萬萬的天使.
就是靈界裡面的天使.
被殺的高揚.
配得權能.
豐富智慧.
力量.
尊貴.
榮耀和頌讚.
這是第二個師班.
可能是常青師班.
接著就是這個.
在天上地上地底下.
滄海裡和天地間.
一切所有被造之物.
就是這個全世界.
所有被造之物.
都在這裡去頌讚上帝.
這個在天上地上地底下.
滄海裡天地一切被造之物.
即是包括了你和我.
就來去唱詩贊頌了.
我們就一起讀出這裡.
好 一二三.
願頌讚尊貴榮耀權勢.
都歸給坐在寶座上的那位和高揚.
直到永永遠遠.
阿們.
這個曾經被殺的高揚.
主耶穌基督.
他是配受極大的讚美的.
當時這個羅馬皇帝叫做竇米先.
其實那時候做基督徒很慘.

$^{521}$在幾年之前就是這個尼祿王.
瘋癲的尼祿王.
接著這個竇米先就更加的瘋癲.
他是一個窮凶極惡的皇帝.
他執政的時候.
他就弄了自己的那些像.
就強迫那些人民.
特別是基督徒.
他說我羅馬王.
這個竇米先王.
就是這個神.
所以每逢人們走過.
看到他的像的時候.
就要來跪那個像.
接著就要大聲的說.
竇米先王是佩德敬拜的.
哇 很誇口.
很誇口.
其實有些老人家.
他經歷過大戰的.
他們都經歷過.
我爸爸那時候都經常跟我們說.
那時候走過那些地方.
看到那些日本人 那些兵.
你記住要跟他們鞠躬.
如果你不鞠躬.
還有九十度鞠躬.
如果你不鞠躬的時候.
就回來回來.
接著就掌摑你.
摑到你牙血又流 鼻血都流.
當時.
約翰不肯跪拜他.
唯有主耶穌基督.
這個被殺的羔羊.
才配受這個頌讚.
他不肯.
所以竇米先就判了他.
就將他放進這個.
燒到火熱辣的油鍋.

$^{561}$要將他燒死.
但不知道為什麼.
燒他不死.
有不同的說法.
有其中兩種.
有人說天突然間下大雨.
你知道那些油鍋.
燒到熱熱的.
大雨下來怎麼辦.
熱油又不滾又不飛出來.
沒有人夠膽走過去.
第二個說法.
那個油鍋燒得.
熱過頭.
竟然裂了.
油都滲出來 怎麼燒.
不要理那麼多.
總之約翰.
就能夠活下去.
竇米先沒有辦法.
炸你不死.
就判你去這個拔魔島.
去充軍.
將他流放到拔魔島.
這個竇米先.
藉著這個權柄.
要人去讚他.
拜他.
但沒有什麼意思.
因為不是甘心情願.
當你拿著主槍.
逼我讚我靚.
讚我靚.
讚你一百句都沒用.
因為不是真的.
但是這個羔羊.
他是配得尊崇的.
他為我們的世人.
被殺.
死在十字架上.

$^{601}$他用自己的血.
從各族各元議各民族.
各邦國當中.
將我們所有的世人.
從這個罪惡.
黑暗撒旦的權勢.
拯救過來.
洗到飯.
信服他的人.
都可以歸入上帝的國度.
這個竇米先.
就想著將約翰.
流放到拔魔海島.
就想著.
去到這個孤島裡面.
你不能再.
出來周圍的地方.
全都.
你一定孤獨死在這個島上.
天上高.
但是上帝的工作.
沒有人可以攔阻的.
只要他心意的時候.
有他的心意的時候.
沒有東西可以攔阻的.
在這個海島上.
這個九十幾歲的約翰.
他竟然就看到.
在天上高.
敬拜那個澎拜的場景.
並且將來.
這個審判.
這個四刃揭開.
然後就.
七號角的.
七刃揭開.
七個號角的吹響.
七碗的憤怒的淋度.
全部被約翰看到.
所以.

$^{641}$剛才就說.
當時很慘很慘.
拉了他這個海島上高.
但是.
原來上帝給這個.
這麼大的榮耀.
部分就被約翰看到.
天上的敬拜.
人類將來覓死的.
那個命運.
剛才我們在聖餐的裡面.
唱那首的詩歌.
十字架 十字架.
曾牧師都告訴我們.
這個是一個盲眼的.
詩人.
他出生了.
六個月之後.
他的眼睛有些東西.
開過醫生就找不到.
有一個巡迴的.
醫生.
他說他會醫.
就找些不知什麼東西.
敷熱了他之後.
就打上他的眼睛.
敷了幾次之後.
他就慢慢看不到東西.
更加慘的.
就是他爸爸在童年.
離開了世界.
所以他的家庭裡面.
就剩下了他的媽媽.
他媽媽當時是21歲.
和他.
這個六個月的.
BB女.
人都覺得他很慘.
看了六個月的東西.
其實都不知道是什麼.

$^{681}$BB 其他什麼都.
沒看過了 都不知道什麼叫.
看東西 是吧 都不知道什麼叫.
看.
但他那首詩歌.
他寫了.
有福的卻具.
耶穌屬我 我甘得先上.
主榮耀喜樂.
他有.
他有一個很有福的人生.
他有福的人生.
為什麼 他因為上帝.
使他成為一個盲的人.
嘩 這樣說話.
都有的 他說.
因為他從來都沒看過東西.
所以根本不知道看東西.
是些什麼東西來的 所以他從來.
都不覺得自己有什麼損失.
他沒什麼損失過.
不過他說他只是.
知道.
在做夢當中.
他曾經看過一些很美的.
景象 一些很美的.
眼睛 很動人的.
面孔 他也說.
他知道將來.
終有一天.
他會開眼看到東西.
他應該當他.
第一眼看到東西的時候.
就會看到的是.
主耶穌基督.
自己.
所以他說.
這是他自己生命中的福氣.
有福的確具.
耶穌屬.

$^{721}$他另外一首詩歌.
就是我們剛才聖餐中說的.
十字架.
永遠是我們的榮耀.
在這個.
屬靈的定律中.
耶穌基督.
他先經歷.
苦難.
然後才進入.
榮耀裡面.
在《肥勒給書》第二章.
他說耶穌基督.
全心信服.
信服上帝的旨意.
以至於死.
且死在十字架上.
一個受苦.
一個極殘酷.
極羞恥的.
受苦.
然後是怎樣.
神將他升為至高.
使一切在天上的.
地上的.
地底下的.
無不屈膝.
無不口稱耶穌基督.
是主.
使榮耀歸於.
父上帝.
這是耶穌基督.
他受苦之後.
他就得著至高的榮耀.
和頌讚.
他的名字是超乎.
萬名之上的名.
所以我們的祈禱是奉主耶穌基督的聖名.
我們處理靈界所有的事情.
都是奉.

$^{761}$耶穌基督的聖名.
因為他就是.
勝過了死亡.
勝過了殺彈全世.
那位被殺的高弱.
他先付出.
美國的開國公臣.
富蘭克林.
在美國開國的時候.
他需要很多人.
去參與.
他有很多.
鼓勵人的說話.
他說.
大多數的男人.
在25歲之後.
就死於.
頸部以上的事情.
什麼叫死於頸部以上的事情呢?.
他說.
現在的人不再有夢想.
不再努力了.
所以國家不行了.
一定要有更多的人.
有更多的理想.
有更多的努力.
他說吃飯是為了生存.
但生存.
不是為了吃飯.
所以.
肯做事的人.
才是幸福的.
最可憐的人.
是懶惰的人.
這是一句名句.
他說.
No pain, no gain.
他沒有.
這個心路.
哪裡會有.

$^{801}$這個成就呢?.
沒有這個心路.
怎麼會有這個成就呢?.
近日就非常忙碌.
因為我姐姐有癌症.
又做了手術.
要幫忙照顧.
教會又很忙碌.
所以沒有時間.
去行山,運動.
有位姐妹說.
你每天做15分鐘.
拉筋就行了.
15分鐘.
於是我就立志了.
幾個星期前.
每天起床.
用腳踩單車.
哇.
不行.
一個星期做了三天已經很感恩.
每天.
就累了.
人們說我樣子很睏.
但不單止.
身體是這樣.
你不去做運動.
就不可以.
經歷到.
在運動辛苦之後.
那份釋放.
同樣.
在這個屬靈的生命.
都是這樣.
如果當我們要屬靈生命有長進.
和主耶穌基督.
有經歷的時候.
我們一定要放上這個時間.
放上一些.
讀經的時間.

$^{841}$一些安靜的時間.
一些默想的時間.
當我們在生命很混亂.
或者處理一些事情.
我們總是不肯.
拿出一些時間來祈禱.
不肯拿出一些時間.
來安靜.
我們就去找一些東西.
找一些人來講.
但當我們發覺到.
只要.
只要靜下來.
靜下來很辛苦.
當我們在學的時候.
很慘.
靜下來的時候.
你會想死了好像什麼都沒做.
死了好像又不知道怎樣.
但當我們.
不經歷.
不經歷覺得很辛苦.
一些屬靈的.
操練上的學習的時候.
我們是經歷不到.
在主裡面.
那份力量.
那份平靜.
那份喜樂.
最後我們來看看.
這位高揚是怎樣的呢.
他是去.
從這個高天.
來到這個地上.
為了要使萬民.
歸入.
父上帝的國度.
在第六節.
對他的形容.
就是從天上.

$^{881}$奉差遣來到.
在這裡再加一句.
主耶穌基督就是.
來到這個幽暗的世代.
將人從這個.
黑暗的地步.
帶進去光明的.
道路.
那他是怎樣的呢?第八節.
他是曾被殺的.
用他的血來救贖人.
救贖人.
將世人.
從撒旦的權勢.
歸入上帝的國度裡面.
使萬人.
歸入上帝的國度.
是上帝的心意.
也是.
教會存在在地上的.
目的.
也是你和我.
生命當中.
上帝給我們的.
大使命.
因為被殺的高揚.
他要將這個.
這個.
愛傳到.
給所有的人.
不單止香港的人.
韓國的人,英國人,美國人.
非洲的人使到普天下.
都歸於上帝.
今天耶穌基督.
已經升了天.
他坐了父神的寶座.
右邊.
你和我都是耶穌基督的門徒.
我們就要延續.

$^{921}$這個使命.
奉差遣.
到普天之下.
不過呢.
我又化到當年.
幾萬萬大的時候.
不知你們是不是.
我告訴你我就是.
我老想.
差不多退休一年多.
之後又怎樣.
休息一下.
停一下.
但上一個禮拜蔡院長給我一個.
很大很大的提醒.
他說要上帝的.
召命繼續的.
臨到.
所以奉差去到.
這個普天下裡面.
仍要繼續在我們的.
思想裡面.
梵克林這樣說.
千萬不要死在這個警暴.
以上的頭顱.
要繼續為上帝有使命.
我們未必走在.
這個前線.
但我們可以在後方.
裡面繼續去禱告.
用金錢.
用關心前線的同工.
一些信結.
我們繼續同心的.
完成上帝給我們的使命.
是要直到.
見主滅.
最後.
都是蔡頌揚牧師的.
一幅漫畫.

$^{961}$我們給你紙去勉勵.
作為信徒.
不單要熱心.
教會的活動.
我們教會的活動.
因為我們屬靈的家.
一起去關心.
但我們同樣都要尋求.
上帝在我們人生裡面的使命.
不同的階段.
有不同的使命.
我們就向上帝問一問.
我們這個.
時間上帝要我們做些什麼.
我們同心的禱告.
天父上帝.
我們多謝你.
讓我們能夠先上到主恩的.
之命.
主願意繼續去祝福我們在座當中.
每一位弟兄姊妹.
我們願意.
繼續去承擔主的使命.
使萬民.
都能夠得到.
福音 求主使用.
幫助我們.
千萬不要懶惰.
讓我們經歷.
為你擺上的仇.
我們就經歷到.
屬靈的福氣.
求主恩待.
多謝你瑞聽禱告.
奉主基督的聖名 阿們.
\newpage



\section{5:1-14}
\label{sec:Yd_an_Ish1s}
\textbf{2023-04-16 早堂《異象中的羔羊是世人惟一的希望》啟示錄5\_1-14 講員:謝靄薇牧師}
\newline
\newline
連結: \href{https://youtube.com/watch?v=Yd-an_Ish1s}{\texttt{https://youtube.com/watch?v=Yd-an\_Ish1s}} ~~~~ 語音日期: 2023-04-19
\newline
\newline
\hyperref[sec:cuPMYAewAe8]{\small{< < < PREV SERMON < < <}}
~
\hyperref[sec:index_chronic]{\small{[返順時目]}}
~
\hyperref[sec:index_scriptual]{\small{[返順卷目]}}
~
\hyperref[sec:SLUSAYk3lm4]{\small{> > > NEXT SERMON > > >}}
\newline
\newline
$^{1}$各位弟兄姊妹平安.
在我們看神的話語之前.
我們一起去仰望主.
昔在今在永在的父神.
榮耀的主耶穌基督.
我們真是不配.
因著你.
我們自有自由.
自有盼望.
在這個時候.
都願意聖靈在你眼中運行.
讓我們心中.
是有些神力不起穴的.
求聖靈來剔除.
讓孩子心中的意念.
口中的言語.
夢裡穴立.
我們這樣祈禱交託.
奉主耶穌基督的名叩.
阿們.
今天我們看這一章的時候.
其實第四章.
都已經是說到很奇妙的.
是約翰提到天上的異象.
他看到天開.
他被感動.
於是他看到寶座.
因為天上的寶座他看到了.
然後就跟他說.
你上來吧.
來這裡.
我要將你發生的事.
去指示你.
然後他還看到.
這個寶座上.
坐著的神.
是手裡拿著的書卷.
這個書卷.
被稱為.
封賢或者密封的書卷.

$^{41}$這個書卷.
是什麼意思呢.
其實.
這個書卷.
是誰可以揭開呢.
你想想這個書卷.
本來不是現在的書這麼簡單.
是用羊皮卷去製造的.
它要捲起來.
然後用繩子或者線綁起它.
然後燒得很熱的鈉.
就放上去了.
然後有一個指環的.
好像戒指一樣的印記就吸上去.
這樣就完成了過程.
但是有七印的時候.
它表示它要捲七次.
如是者都要這樣做.
然後才可以封起來.
將來要有人去啟封.
才可以看裡面的內容.
我們看到這個書卷裡外都有字.
曾經在以色列的二章九到十節.
都同樣好像啟示錄五章一節所說的.
它說有一隻手向我伸來.
看那手中有一書卷.
它在我面前展開書卷.
它內外都寫著字.
上面所寫有哀號 嘆息 悲痛的話.
同樣都有字.
到底封嚴和密封是什麼意思呢.
就像我剛才告訴你的.
是封住了的.
封住了之後.
究竟這本書是說什麼的呢.
有些人就說可能是語言書.
可能是歷史書.
或者是一些過去的.
或者現在正在發生的事情.
經過學者的研究.

$^{81}$就認為是一些基業的文契.
內外都有字的.
這些文契根據猶太人的規矩.
就是文契上的條款或者價值寫好了之後.
就會封上印.
而文契上面就會寫上一些見證人的.
如此者內外都有字.
耶利米斯三十二章曾經記載了.
利布格尼撒王要去奪取耶路撒冷的時候.
上帝指示耶利米.
你要告訴百姓.
他們將要被擄到巴比倫.
然後七十年之後.
你們可以回到這個地方.
他有一個親屬.
耶利米斯有一個親屬.
在這個親屬有一塊地.
他想想.
既然被擄了.
這塊地就沒有用了.
那我要賣掉它了.
那些哈拿蜜呢.
就想辦法.
自己去想.
不如賣給耶利米斯吧.
可能熟的人都是這樣想的.
他賣給了耶利米斯.
然後又指示他.
你儘管賣吧.
這塊地將來是很有價值的.
同時會被人重用的.
七十年之後.
你看到不久之後哈拿蜜.
就真的.
又要稅于耶利米斯了.
然後他就在這個契上簽字.
然後封好它.
封了鑒之後.
因為他自己將要被擄.
所以他就交給這個書記巴魯.

$^{121}$將來伊斯利人被擄歸回的時候.
七十年之後.
耶利米斯的後裔.
一個代表的繼承人.
他就告訴他.
在公堂的公證人面前.
說我是耶利米斯的後裔.
是代表人來的.
可以繼承的.
於是他就可以有權去啟封.
得到這塊地.
所以約翰所看見的神手中的書卷.
可能就是這個契約的書.
但是四章明明是很清楚的.
去指示告訴他.
將來發生甚麼事.
但是來到一章的時候.
第五章.
第五章是繼續第四章的.
第四章是講到他的寶座怎樣.
然後第五章就見到有書卷出現了.
但是你看到就是.
在經文裡面.
講到一個天使大聲地叫.
誰是可以展開書卷的呢.
天上地上地下.
沒有人是可以展開的.
在這個時候.
約翰忽然哭了出來.
很悲傷.
他很難過.
到底誰可以配這書卷呢.
但是有一個長老就去勸他.
這個長老說.
不用哭 你不要哭.
看 猶大支派中的獅子大衛的根.
他已得勝 能展開那書卷.
揭開那七個印.
不要哭.
其實你想起就是.

$^{161}$那因城的寡婦.
她的兒子死了.
耶穌都曾經叫她不要哭.
但是現在.
可以展開書卷的是誰呢.
是耶穌基督.
所以.
耶穌基督曾經在路亞夫音.
六章二十一節.
說到現在哭泣的人有福了.
因為你不將要歡笑.
愛的應許是會實現的.
所以長老中.
有一個說不要哭.
猶大支派.
說到耶穌基督.
就是大衛的根.
不只是表明他是一個得勝的君王.
他還是一個永遠在寶座上做王的.
耶穌基督.
所以.
肯定主耶穌.
既然是大衛的根.
就表明大衛是由耶穌基督而生.
即是未有大衛之前.
已經有了耶穌基督.
而在《羅馬書》一章三到四節都說到.
「論到他兒子我主耶穌基督.
按肉體說.
是從大衛後裔生的.
按神聖的靈說.
神因從死人中復活.
用大能顯明他是神的兒子」.
所以主耶穌基督.
可以坐在大衛的位置上.
做王直到永遠.
耶穌基督是猶大支派的獅子.
其實約翰當時看到的情景.
就是看到的羔羊.
獅子原本是很威嚴的.

$^{201}$很嚴厲的.
很有榮耀的.
很有能力的.
但是羔羊.
是很謙卑的.
很溫柔的 很軟弱的.
獅子表明了耶穌基督的神聖.
而羔羊表明了耶穌基督的人性.
主耶穌就具備這兩個.
神人異性.
才可以完成這個救贖的功勞.
他降世成為人.
成為羔羊.
他由羔羊這個溫和的特性.
為我們而降生.
為我們而受死.
為我們而復活.
升天.
將來他再來的時候.
是以獅子的特性成為王.
在萬國裡面去掌權.
現在主對這些忠心的聖童.
很溫柔 很謙遜的.
因為他是曾經被殺的羔羊.
去除去世人的罪孽.
所以.
這裡就說到.
我在《漢見》第六節說到.
漢見寶座和四個活物.
以及長老之中有羔羊站著.
將是被殺的有七個角.
七隻眼睛.
就似神的七靈.
奉差遣往普天下去的.
這個羔羊是誰呢.
羔羊在《約翰福音》出現了兩次.
在《約翰福音》第一章的29節和36節.
看神的羔羊.
大家要記憶.
而在《啟示錄》出現了27次.

$^{241}$《聖經》裡每逢說到羔羊的時候.
都是說到祂待死.
祂贖罪的.
所以.
在《耳塞夏斯文》53章7節說到.
祂被欺壓受苦卻不開口.
祂將羔羊被牽去宰殺.
有像羊在指模的人手下無聲.
祂也是這樣不開口.
大家看到就是.
這個羔羊.
祂是被殺的.
憑著耶穌基督的血.
是如同無下無始的.
無顛覆羔羊的血.
既然你們是無口的面.
要把舊口除淨.
好使你們成為新團.
因為我們御節的羔羊.
基督已經被殺獻為祭神了.
所以羔羊.
約翰死的羔羊的時候.
原文裡面是用了「小羔羊」.
去形容羔羊是多麼的溫柔 謙遜 信服 可愛.
更加將整個愛和旨意彰顯了出來.
因為它就是曾經被殺的羔羊.
它只可以剪開書卷.
揭開七個印.
羔羊就是指耶穌基督.
因為祂為我們付上極重的代價.
祂曾經被殺.
用自己的血買熟了我們.
羔羊的形狀是怎樣的呢.
它可以揭開七印.
羔羊的形狀.
首先我們看到的就是有七個角.
角是這樣的.
有七個.
這個角其實是代表能力和權柄的象徵.
它詩寫.

$^{281}$州制貧窮.
他的義行傳到永遠.
他的角必被高舉.
大有榮耀.
在聖經看到角的時候.
都可以看到是很有能力的.
很有權柄的.
七個眼.
一個眼睛.
有七隻.
如果身上有七隻眼.
不知道你會怎樣想像呢.
七字在聖經是表示完全的數字.
表示神是無所不知的.
祂全部都知道的.
七靈.
神的七靈指的是聖靈.
指的是完全的聖靈.
主耶穌有完全的智慧和知識的靈.
與祂同在.
所以在《二乘二》第11章1-2節.
說到.
從耶西的殘幹必長出聯知.
在根所抽的枝子必結果實.
本來已經是不行的了.
果實.
耶華的靈必住在祂身上.
就是智慧和聰明的靈.
謬略和能力的靈.
知識和敬畏耶華的靈.
所以聖靈很厲害.
主耶穌說.
「我若不去保衛施,就不會到你們這裡來」.
「我若去,就猜他到你們這裡來」.
所以聖靈代表耶穌的完全的智慧和知識.
可以參透萬事.
所以祂配可以開的書卷.
可以揭開七印.
所以耶穌基督由寶座上的神.
手中拿出書卷的時候.

$^{321}$就代表由神那裡得到權柄.
國度榮耀.
去奉命行事.
所以高揚就是指著耶穌基督.
因為祂負上極重的代價.
說我們被殺.
用自己的寶血去買贖我們.
高揚的救贖又是怎樣呢.
高揚的救贖就是指著耶穌基督用他的血.
在聖經裡面.
是說到.
祂是由各支派.
各語言.
各民族.
各邦國.
真的指著甚麼呢.
就是萬世.
萬族.
萬民.
因為主願萬人得救.
不願一人沉淪.
所以祂是要救贖我們.
祂說我們死也是要救贖普天下的人.
所以.
你看到.
主耶穌是多麼愛世人.
有一個寫暗室之後的蔡素娟女士.
她曾經在1890年出生.
她爸爸是做官的.
家裡有18個孩子.
她的小名叫做太多的意思.
她的家住在南京王府的大街上.
她家就佈置得像怡和園一樣.
你可以想像她家裡有這麼多人.
那個價格是很富貴的.
很有派頭的.
她是這樣成長的.
她看盡榮華背後的虛假.
她很早就已經體會到生活的空虛.
面對未來的恐懼.

$^{361}$她小小年紀就要出家為例.
她很想離家.
慢慢長大後就想離家.
離開家裡去基督教學校讀書.
她有三次交了學費.
都被父母攔住不讓她去.
她就想辦法.
這次我不遠行.
近一點就在南京.
去南京遊說她媽媽.
去南京教會學校讀書怎樣.
她媽媽就讓她去.
這麼近.
16歲的時候.
她進入孫教士所辦的學校.
這個孫教士叫李萬女士.
影響她一生.
她在明德女子書院讀書.
她一心想學英文.
學音樂.
她對基督教很反感.
她沒什麼興趣.
因為孫教士內在的光芒.
他那種沉靜的力量.
是很吸引她的.
在早會裡.
她聽到一個美國傳道人.
戈登先生.
他說耶穌是世上的光.
她馬上就體會到.
神就是接受耶穌基督.
帶來的救恩.
於是她如意接受耶穌做救主.
由那天開始.
她的生命改變了.
雖然她在經裡.
看到滿清被推翻.
又軍閥的時代.
抗日內戰.
種種顛簸的生活.

$^{401}$震盪的時代.
她都經歷過.
但是她和孫教士李萬女士.
是有深厚的感情.
做了她的屬靈母親.
她們一起配搭.
去傳福音.
她就在她家裡.
帶了五十五個人.
她是大家庭.
帶了五十五個人信主.
她英文中文都好.
所以她和孫教士.
是最佳拍檔.
經常帶她去美國做見證.
她的事故越來越多.
在1931年的時候.
那天是一個冬天的早晨.
忽然間她的眼睛很赤痛.
痛得不得了.
找人醫治一年.
醫生都束手無策.
沒辦法.
她不見得光.
16年之後.
終於發現她的病症是.
眼睛怕光.
因為她是受到虐疾.
引起這個症狀.
沒辦法失控.
所以.
沒辦法平衡的時候.
發現太晚了.
沒辦法醫治16年.
所以她的生活改觀了.
在一個動亂的年代.
她的身體受到限制.
但是神有奇妙的帶領.
保守她.
她開始寫作.

$^{441}$1949年的時候.
李邁女士帶她去賓州.
寫作成為她的另一個事奉.
她的著作裡面就是.
《暗室之後》.
她裡面說到.
床蝶不是她的監獄.
而是她受訓的學校.
聲明是她的導師.
訪客是她的功課.
她到了1984年才離世.
她足足有半個世紀.
在這個暗室裡面生活.
生命卻影響了不少人.
弟兄姊妹.
耶穌基督救贖了你.
你的生命有沒有改變呢.
高揚的祝福.
祝福我們甚麼呢.
在第六節說到要奉差遣往普天下.
主耶穌為我們受死.
他復活.
臨升天的時候.
祂就祝福門徒要往普天下傳播音.
給萬民聽.
大使命就是你們要去.
使萬民做祂的門徒.
大家很清楚這個大使命的經文.
今天.
不要說那麼遠.
宣教很偉大.
又要學語言.
那麼麻煩.
在本地也有宣教.
本地那麼多外籍人士.
我們教會也有印尼姐姐.
有菲律賓.
很多功夫可以做.
只要你的英文好.
你可以放膽去做.

$^{481}$同時你也可以安排.
你的工人.
可以去聽福音.
你也是在幫助他.
甚至年輕一點.
伯地利中學.
陳靜娟姑娘.
她就需要人手.
這些尼泊爾的學生.
他們要學中文.
不懂廣東話.
你教教他吧.
只要懂一點英文.
就可以跟他溝通.
這些學生很單純.
他們想學中文.
你教教他.
他可以通過.
成績就可以好.
他又滿足.
你又開心.
家長又歡喜.
你跟他建立關係.
有時在加州有團契.
你懂一點英文.
就可以陪他一起玩.
跟他一起成長.
了解他.
建立了關係就可以傳福音.
其實宣教不要想得那麼遠.
在本地就可以做.
高揚.
他配蘭.
這個書卷.
你看到.
耶穌基督.
因為他曾經為我們受死.
所以他.
可以拿這個書卷.
可以揭這個七印.

$^{521}$施活默.
和廿世偉長老.
他們俯伏在高揚的面前.
獻上祈禱.
和唱新歌.
祈禱.
你看到在經文裡.
他講到.
城滿鄉的甘露.
看到禽.
是指甚麼呢.
崇拜裡面的樂器.
在敬拜裡面.
頌讚的時候都會加上不同的配樂.
其實是表示甚麼呢.
表示崇拜的人.
向神敬拜.
盡心盡性盡義盡力去愛我們的神.
所以他要將最好的獻給神.
所以這個都是我們敬拜的時候.
應該有的態度.
是預備充足.
然後祭司的特性.
他的特權就是.
要獻上新的香.
香就是.
象徵向神的祈禱.
這個香.
都是眾聖徒他們的禱告.
去達到神的面前.
同時是.
蒙神.
月立.
很歡喜的.
很香的香.
是達到神的面前.
神去聽你的祈禱.
唱新歌.
你會看到在啟示錄裡面.
經常有新天新地.

$^{561}$新耶路撒冷.
唱新歌.
當耶穌基督揭開了七印.
他們就唱新歌.
他還唱新歌.
因為耶穌基督.
這個羔羊基督.
好轉了這個形勢.
所以.
不得不你要去.
歌頌他 去頌讚他.
此情此景.
你怎可以還可以.
借用其他詩歌去頌讚他.
一定是.
由你心底發出一些.
最誠心去頌讚他的.
詩歌 很應景的詩歌.
所以叫新歌.
看到.
頌天使.
他們去歌頌羔羊.
由千千萬萬的天使去頌讚耶穌基督.
頌讚這個曾經被殺的羔羊.
被殺的羔羊.
配得權能豐富智慧.
力量尊貴榮耀.
頌讚.
七重的頌讚.
去頌讚我們的神.
所以整個崇拜裡面.
去參與的.
不是只是天上的.
都包括全宇宙的.
這些受造之物.
包括了你和我.
大家都是.
願意去頌讚.
尊貴榮耀權勢.
歸給坐在寶座上的神.

$^{601}$和耶穌基督.
直到永遠.
那崇拜的中心就是神.
無論你這個世界有多麼的分岸.
你的際遇有多麼的不好.
溫慧耀博士.
他寫了一本書.
《上帝與人間的苦難》.
他描述了人體.
為絕路.
是上帝奇妙.
作為的起點.
他曾經申請全家移民.
去加拿大.
他各方面的條件都夠.
於是很快就批准了.
但是他的小女兒有嚴重的.
智障的病.
所以政府就沒有批准.
當時加拿大政府解決的方法.
就是全家可以.
永久居留.
他女兒.
就批准了一張.
移民部發出的部長特許證.
需要每年續期.
他覺得合理.
因為.
你不知道你女兒.
去到那裡.
在醫療的花費.
健康照顧.
會不會動用很大的開支.
就這樣批准了.
要看察五年.
他的許可證.
是文件八十六.
五年就很快就過去了.
他發現他們.
都沒有用太多.

$^{641}$政府的資源.
全部都是自己去承擔一切的開支.
然後他們就.
等著通知.
但是收到的是一個.
紅色的打印的公函.
內容竟然是叫他女兒.
三個月就要離開了.
永遠都不可以再回來.
他們很沮喪.
都不明白.
什麼原因.
後來才知道過中的原因.
是這個政府.
加拿大政府搞錯了.
一個很荒謬的錯失.
他女兒文件是八十六.
五年之後會轉為永久居留.
在最後一次續期的時候.
不知道是哪個糊塗的官員.
就將文件八十六.
抄錯了.
文件八十六就剛剛相反.
就是要離開的.
他們在完全不知情之下.
做夢都沒有想過.
會碰到一個這麼糊塗.
荒謬的事發生.
他女兒是在英國出生的.
她在香港.
還沒有住滿七年.
那怎麼辦呢.
加拿大又不收她.
這個時候.
很困難.
她女兒已經十幾歲了.
她的病情很複雜.
需要有特殊的宿舍.
給她照顧.
但是再申請.

$^{681}$要等九年.
真的走投無路.
又跟教會一起去祈禱.
後來有人建議她去英國.
既然她又出生在英國.
很快就幫她找到倫敦.
有一個特殊照顧的宿舍.
她又符合資格.
她認為神為她開一條.
不可能的路.
繼續祈禱.
很多人建議她去嘗試.
這一年她剛剛好.
又答應瑞典.
在北歐.
有一個聯合退休會.
她說完.
剛剛好我去倫敦.
然後兩個星期就去退休會.
剛剛好.
臨行之前.
這個特殊照顧的宿舍.
女主管很友善的.
她說不用擔心.
經過三天之後.
女兒一定很開心.
然後一年探望一次就好.
她很期待.
於是很艱辛的坐了九個小時的飛機.
你知道女兒是特殊的.
不是正常的.
所以兩夫婦在飛機上.
要抓住她.
她經常叫.
這個時間很難過.
終於到達目的地.
但是.
這個主管.
電話中叫她不要擔心的主管.
接見她的時候.

$^{721}$她說我們不會接受你的女兒.
我有沒有聽錯.
為什麼會有這樣的反應.
真的很難置信.
趕他們走.
兩三個人推她出去.
這麼奇怪.
她說.
我沒有說過.
這麼奇怪.
叫我不要擔心.
我這麼辛苦來到這裡.
你趕我走.
她說你不走我們就要報警.
這麼嚴重.
她苦苦哀求兩天都不成功.
後來她靜悄悄的問清楚.
什麼原因.
這個女主管跟她說.
剛剛有福利署的官員.
他們是巡視.
剛剛你的檔案放在桌子上.
他打開來看.
原來你既然不是本土的人.
他們拒絕.
不准他們留下.
因為他們會用英國很多資源.
幾經波折.
國際保護兒童的人.
又來調查.
很奇妙的事發生.
這個機構.
我們不需要計較兒童的國籍.
就可以收留她.
於是找到一個庇護宿舍.
給這個女兒.
一天就能夠收留她.
一個宿舍給這個女兒.
一住就住到26歲.
她離開世界.

$^{761}$而她女兒就是最後一個.
被收留的.
外籍人士.
最後一個.
因為當她收留了之後.
政府又改了.
他說不再收留.
不是加拿大籍的人.
他們想起.
當年.
很絕望.
很多質問.
很多不明白.
但回頭看.
神的作為.
神的計劃.
祂的部署.
奇妙到不能形容.
問題就是.
我們今天.
相不相信神.
我們相不相信神.
是不會搞錯的.
我們身體.
會有疲乏.
但神的手.
仍然繼續工作.
儘管公司的成績.
是多麼的標青.
但有一天.
我們終於在神的寶座面前.
去面見神的審判.
這本書卷.
代表人類歷史的終局.
包括審判.
和救贖.
神最後的旨意.
就是一定會執行.
既然高揚耶穌基督.
是我們唯一的盼望.

$^{801}$你是否真的相信祂.
你是否願意聽祂的話.
要往普天下.
傳福音給萬民聽.
使人能夠同得福音的好處.
既然主耶穌是唯一.
可以開啟書卷的.
祂救贖了我們.
我們是否要尊主為大.
我們是否要將榮耀頌讚.
歸給坐在寶座上的神.
和耶穌基督呢.
我們一起祈禱.
掌管萬有的主.
你是配得我們頌讚的.
今天我們生活在一個.
不可以.
明天是怎樣.
這個疫情.
又不知道.
有怎樣的變化.
但總是.
主.
你透過聖靈.
帶給我們.
內心有從來以來的平安.
還有.
我們有時.
得罪你.
可能說的語言.
是失控.
但主你仍然在.
等待著我們回轉.
還有今天.
若然我們的心遠離你.
求聖靈.
在當中督促我們.
要回轉到神的面前.
因為將來我們要交賬.
我們要面對神的審判的時候.

$^{841}$我們有甚麼可以交給神呢.
主又求你.
幫助我們要快點作供.
幫助我們真的要.
快點去傳呼音.
因為你快要回來.
主又求你幫助我們.
加我們能力.
我們何時退後.
求聖靈你鞭策我們.
要快點去作供.
以致我們能夠有所.
呈獻給主.
求主在當中帶領我們.
幫助我們.
以致我們所做的.
能夠合乎你的旨意.
我們這樣祈禱交待.
奉主耶穌基督的名教.
\newpage



\section{}
\label{sec:SLUSAYk3lm4}
\textbf{2023-04-20 周四晚查經班《被擄歸回歷史研讀(12)》 講員:曾裔貴牧師}
\newline
\newline
連結: \href{https://youtube.com/watch?v=SLUSAYk3lm4}{\texttt{https://youtube.com/watch?v=SLUSAYk3lm4}} ~~~~ 語音日期: 2023-04-21
\newline
\newline
\hyperref[sec:Yd_an_Ish1s]{\small{< < < PREV SERMON < < <}}
~
\hyperref[sec:index_chronic]{\small{[返順時目]}}
~
\hyperref[sec:index_scriptual]{\small{[返順卷目]}}
~
\hyperref[sec:qxsLxZod7kk]{\small{> > > NEXT SERMON > > >}}
\newline
\newline
$^{1}$弟兄姊妹,我們可以開始今晚的茶經班.
我們一起祈禱.
再一次多謝主.
我們能夠有系統地去研讀.
歸回的歷史和先知書.
我們來到尼希米記第四章.
很盼望這些第三批回去重建.
成長和面對的困難.
都讓我們透過歷史.
有一個了解.
更重要的讓我們摸索.
在這個時代上主的工作.
是要怎樣開展.
讓我們能夠在我們所處的環境.
一方面我們要找到.
什麼是我們的工程.
什麼是我們要去建造.
同心的建造.
還有什麼是我們的仇敵.
在我們的附近.
或者是一些.
可能和屬靈有對抗的價值觀.
都讓我們能夠去明白.
一起來去.
好像彼此的守護.
能夠最終這個工程能夠完成.
我們人生都有我們自己的工程.
都是不知何時我們要息怒歸主.
我們要面對死亡的一刻.
盼望神自己都來幫助.
我們今天晚上等候在你的面前.
我們這樣禱告奉耶穌基督的名.
阿們.
偶然間都會有新朋友加入.
來查經.
我們是有足夠的空間.
可以將一些過去的背景.
每次都深一點.
讓新朋友他們會覺得.
不會太過陌生.

$^{41}$同時也要有些背景了解.
經文再讀下去的時候.
又會多一點點立體感和印象.
剛才在祈禱的時候提過.
尼希米其實是在波斯帝國.
亞達薩西王年間的一個酒政.
在當時古代的帝國裡.
酒政是一個很重要的崗位.
是在帝王身邊.
來試酒.
其實是一些重要的宮廷宴會.
和日常的飲食.
都是要靠酒政來預備的.
我們稍後將來看的《耳絲帖記》.
原來宮廷是會有暗殺事件的.
本來這個阿黑隋老王.
有人不喜歡他.
真的想暗殺他.
剛好被那個叫茉底葵的猶太人.
發現了報告果真有其事.
就抓了那兩個太監.
所以原來在古代暗殺事件都有的.
這個大家可以留意.
但是我們一開始就說尼希米的時候.
就會發現第一到第六章.
是尼希米他的一個回憶錄.
他講述他自己的夢召.
他自己身處的當時的世代.
那個處境.
猶太人已經可以在自己的故土耶路撒冷.
來到去安居了.
但是不能夠有一個.
回到摩西時代的身份.
那個屬靈的遺訓.
就是因為耶路撒冷.
還不能夠有足夠的自治的能力.
一來就是硬件的破損.
整個尼希米都是想告訴我們.
那個城牆很早已經說是城牆被火焚燒.
城門拆毀.

$^{81}$外邦人他們就可以隨隨便便.
來到去耶路撒冷.
一來就是令到猶太人.
沒有了分別為性的那種觀念.
那種的執行.
很多時候就會有這種跟外邦人.
當地人.
有撒馬利亞人.
有阿拉伯人.
這樣來到去做生意.
來到去通婚.
令到其實他們那個國民身份.
叫做Sense of Belonging.
歸屬感.
對自己身份的歸屬不強烈.
慢慢慢慢很多東西就不能夠堅持.
不可以來到去執行.
有時祭司都變了好像形勢.
大家都不是有那個心向著神的聖殿.
或者是去到節期.
要守節.
又有很多的可能五花八門的節慶.
或者是他們那些的什麼遊藝會.
或者是他們一些的.
跟外邦人一起的.
可能是做生意.
那就變了.
大家慢慢就淡忘了自己聖徒.
就是神選民的身份.
所以導致到這個情況.
似乎持續了一段長時間.
慢慢外邦人根本不會尊重你猶太人.
安息日就來跟你做生意.
可能讓你看到便宜靚正.
你就會慢慢疏於對聖殿的敬拜.
等於譬如星期六星期日.
那些香港的節目.
每個節目都豐富.
那就很容易搶了我們信徒的心靈.
來去敬拜.

$^{121}$暫時好的就是.
早上十一點前.
都沒有什麼大型活動.
個個都還沒醒.
還沒轉身.
所以回教會就很正常.
都不會令到大家很尷尬.
但是偶然間都有一些.
譬如馬拉松等等.
不是經常性每個月.
一年一次都尚且.
是可以接受.
某程度都有一定的運動量.
那就無可厚非.
但如果好像當時.
尼欺米的時代.
已經沒有了所謂安息日的界線.
很快我們又會見到.
去到近尾十章就會出現.
你就會發現.
那個城牆就變成一個.
必需要的界線.
即是要有一個borderline.
是和外邦人要隔絕.
這個就好像哲理行間.
我們第二堂課講過.
就是他們要保持這個身份.
因為神的救贖的計劃.
很快彌賽亞就會降臨.
就會出現.
必有童女懷孕生子.
人要稱他的名為爾瑪萊尼.
和在伯利恆爾法他.
會有一位君王會出現.
所以誰是童女懷孕生子呢.
其實猶太人他們.
尤其是大姐仔未結婚的.
他們都期望自己是童女.
懷孕生子.
我是不是呢.

$^{161}$所以去到瑪利亞.
他真的問.
我就是那個蒙福的女子.
天使說是呀.
你就是蒙大恩的女子.
所以其實對猶太人來說.
特別是那些女孩來說.
未出嫁的那些.
他們就期待自己會不會.
在自己的肚腹.
生這一位彌賽亞耶穌救主.
所以其實某程度.
對猶太人來說.
這個救贖其實是在將來的.
這個救贖.
當然他們就說彌賽亞.
一次就是救了一切.
整個國家都可以改變了.
包括我們可以執政了.
在地上真的作王了.
好像我們可以抬頭了.
但是他們就是一直都猜不到.
彌賽亞是兩次.
一次是來受苦.
登十字架贖罪.
第二次才是榮耀的再來.
也是我們基督徒.
等候的另外一次.
彌賽亞的再來.
那個復興.
所以可以很肯定就是.
猶太人在他們被擄歸回.
為什麼那麼緊張.
那個聖殿.
神都感動一些有心人.
由老遠的波斯帝國.
回到去重建聖殿.
那個原因.
不是只是想有敬拜.
燒香而是這個民族.

$^{201}$是要孕育將來.
當然他們不知道是四百年後.
的彌賽亞降臨.
所以呢.
這個也是我們去看.
這個整個的.
回歸的歷史的時候.
因為我們是後代的人.
所以我們就能夠.
從後代回望.
的時候.
我們就知道原來神的計劃.
是一路一路的實現.
所以在這個時代裡面.
這是一個.
神的憤怒.
刑罰.
的時代已經過去了.
就是說.
如果用我們現在比較能夠.
容易明白的.
疫情過去了.
現在是一個所有的東西都.
好像爆發式的復常的.
開始了.
大家就要動起來了.
對的,每個人都想.
真的出去走走.
三年屈門住很辛苦.
身心靈都累.
那是其中一樣.
對整個人都會有一個.
恢復,但是要想想.
屬靈的.
如果三年那個沉寂.
現在我們重新開始的時候.
我們都要想想.
我這個神選民.
基督徒的身份.
的歸屬.

$^{241}$那個sense of belonging.
怎樣去活躍出來.
這些是要想的.
牧者是要推動弟兄姊妹.
要真是令自己去想.
我們大家都很強烈.
我們是香港人.
我們是新一代的香港人.
或者我們是什麼的.
但是我們要有這個歸屬.
譬如你是一間名校的.
假設你是黃人的.
或者你是牛津的.
我是哪間學校的.
當然我也不敢說我自己是哪間.
因為牛津道最差的那間.
就是我那間.
無犯.
我上個禮拜.
前個禮拜報告的時候.
有姐妹跟著認我,我是隔壁那間.
類似這樣.
那你就可以知道.
身份的歸屬.
是要透過.
有心人建造聖殿.
有以色列的.
復興.
認罪的復興.
有現在我們看的黎希米.
他在很大的困難.
底下他都要來到.
重建整個的城牆.
就是說四圍.
是有界限.
界限人就是這樣.
慢慢去到耶穌的年代.
那個界限就變成.
很強烈的.
那種民族主義.

$^{281}$很強烈的排外.
所以撒瑪利亞人.
在耶穌的年代.
是被歧視的.
是不可以有來往的.
不可以有大家的交際的.
原因就是.
他們又變成矯枉過正.
現在就是需要有一個界線的時候.
這個界線.
只不過是讓自己那個.
聖徒.
神選民的身份.
有回這一種的歸屬.
當你有了.
神的工作就會在這個民族.
慢慢慢慢地開展.
好奇妙地.
這四百多年.
其實是神慢慢地工作.
他們是慢慢地等待.
一直都是受著.
幾個強的帝國.
埃及.
接著慢慢興起的.
就是波斯之後.
就是這個.
希臘帝國.
就興起了.
那你就可想而知.
其實猶太人都是被這幾個帝國的歷史.
左右著的.
他們就是好像扮演著.
一個很友善的.
這樣的角色.
但是你不要理神的工作.
就是透過這個小小的地方.
這一班的選民.
來顯示.
彌賽亞將來.

$^{321}$降臨的獨特之處.
所以.
我們有這個背景的時候.
當我們fit回到經文的世界.
你就知道.
雖然尼希米他沒有那麼偉大.
想到四百多年之後.
就會有他的後代.
耶穌基督出現.
他不是我們.
因為他是發生在耶穌之前的人類.
我們是之後.
所以我們就能夠看到.
這二千多年的歷史.
但是當我們剛才一直這樣說的時候.
第三章.
這一位的酒政.
有相當.
大的領導才能.
那個行政的.
主導權.
那個運籌維握.
和能夠連絡很多.
他那些其實不認識的.
他的猶太人的同胞.
第三章其實是相當悶的.
哪個家族.
哪個支派.
哪一班人.
在哪個方位.
來做工.
記載了三十多節.
就是這樣的一個描述.
好像先告訴我.
現在這個工地的藍圖.
給大家.
來看到完全是.
很悶的.
很平鋪直敘的.
哪個家族.

$^{361}$做哪些事.
他們是在靠近自己的家.
那個地方.
來輕工建造.
有些就是被抽調.
去到不同的城門.
來做的.
就是記載這些事.
就是告訴我們.
當時李希米真有其人.
真是帶著那個詔書.
達舍西王的那個.
御旨.
由河東去到河西.
約旦河西邊.
這樣來去.
重建這個城牆.
所以他.
預備好他的計劃.
他的藍圖.
因為他之前那幾晚.
就視察好.
有了心理準備.
好像開會之前.
做定一些功課.
然後就焦聚.
猶太耶路撒冷的長老.
祭司長.
來講他的計劃.
講他的那個.
初心.
講他在神面前.
得到的異象.
是由他自己禱告.
求回來的.
他自己真是.
好像沒有事找事來做.
本來在那個官職裡面.
真是風馬牛都不相及.
沒理由他要自己.

$^{401}$這麼辛苦.
長途跋涉上去做一些工程.
應該在皇宮裡面.
繼續都可以這樣做.
你安息日.
自己慢慢.
遊說到皇帝.
安息日一定要休息.
要敬拜我的神.
我都為皇帝祈禱.
其實沒問題.
安息日自己休假.
來向著耶路撒冷祈禱.
已經是很好.
很敬虔.
你可以想像得到.
但是.
他有更大的.
一個心空.
心靈.
是裝載著.
神的電.
裝載著神家的芳良.
這些就沒得.
強求的.
每個人 保羅都教我們.
講完.
大家熟悉的羅馬書十二章.
所以弟兄們.
我以神的慈悲.
勸你們.
因為我們是被主耶穌拯救.
將你們的身體獻上.
給神.
好像一個活祭.
獻上給神.
是聖潔的.
是神所喜悅的.
你們這樣事奉.
將自己獻上.

$^{441}$就是一個事奉.
事奉這個字是stand.
等候.
等候神的差遣.
你們如此的等候.
你們如此的事奉.
是理所當然的.
因為理所當然.
是因為我們要回到.
神透過他的愛子.
無條件的犧牲.
將我們每一個人.
拯救 所以我們的獻上.
不是很偉大.
向這位造物主獻上.
獻上你的生命 他願意收納.
是聖潔 是神喜悅的.
是理所當然.
古羅說 他本來就是罪魁.
他說他自己.
所以理所當然.
就是說我們每一次來到教會.
我敬拜神.
我做唱詩歌敬拜神.
是理所當然的 是我的福份.
你這樣想.
你來崇拜都開心些.
你看到.
不同的環節.
你的內心都充滿了.
感動 因為.
理所當然.
不是必然.
更加不是奉旨.
沒有你是無所謂的.
沒有你 沒了我.
那個敬拜是不會影響的.
如果這樣想.
那我就緊張了.
因為我一定要在事奉裡.

$^{481}$要合神的心意.
否則沒有我 是可以的.
是可以的.
神可以.
不要你 而興起另外的人.
這是其中一個問題.
羅馬書.
接著說.
我憑著所賜給我的恩典.
保羅說.
第三節.
對你們各人說.
不要看自己過於所當看.
各人要按.
信心的程度.
看得 看自己.
這個獻上.
是要合乎忠道.
合乎忠道的獻上.
就是說不要過急.
不要過慢.
要按神給你的恩典.
你的背景.
來看自己.
當然更大的就是.
要在救恩下看自己.
你要合乎忠道.
當你越蒙恩.
你就覺得自己的獻上越不足夠.
這個就隨著我們對主的愛的認識.
就會有一個與日俱增.
或者與時俱進的那種心情.
所以這個合乎忠道.
永遠都停留在一種張力.
永遠都停留在一種.
不是我獻上了.
就這樣了.
有一條線畫了.
由這裡開始就不關我的事.
其實不是的.

$^{521}$應該是不斷你的人生.
都有新的經歷.
你自己與神之間.
與舊恩的愛.
裡面的合乎忠道.
不是與隔離位的合乎忠道.
兩個人都不同.
他有他的殉教.
殉道.
你有你的坐在這裡.
去到白頭佳老.
每個人都不同.
但仍然都是講合乎忠道.
所以現在神給了尼熙米.
某程度上是一個念頭.
放在他裡面.
這個念頭是從何時開始.
他沒有交代給我們知道.
他只是告訴我們.
他已經一開場.
他已經是在亞達舍西王.
面前的酒精.
他如何爬到這個高位.
他如何在波斯帝國成長.
受教育.
而他能夠保持自己猶太人的血統.
那個歸屬感.
sense of belonging.
他如何能夠.
既是一個外邦.
君王的.
面前的酒精.
但同時間保持自己的敬虔.
自己的熱心.
這個拿捏.
這個平衡就很難了.
他如何能夠拿捏得這麼平衡.
而導致他現在.
一見到有幾個猶太人.
在耶路撒冷下來.

$^{561}$他就追問他們.
耶路撒冷的光景.
我祖宗所在的地方現在怎樣.
然後第一章第二節就說.
他們告訴他們.
哈拿尼就是那個弟兄.
就告訴他們耶路撒冷的城牆.
破壞.
城門被火焚燒.
猶太人受盡凌辱.
那個凌辱不是真的被虐待.
行刑.
可能就是看你們猶太人不起.
我來做生意.
你又跟我做.
你又不守安息日.
很多時候就會慢慢.
他們就看不起你耶路撒冷.
又說重建.
又說神是你們自己唯一的神.
人家根本不會對你們.
有這樣的尊重.
但是.
聖經沒有交代.
給我們知道.
尼兄米也沒有告訴我們.
他在什麼情況下.
哪一次的陪靈會.
哪一次的讀經.
他有這份很大的感動.
導致他將這一點的火.
內心的一個這樣的火花.
燃點成為一個事工.
而這個事工.
經過很多難處.
他仍然能夠繼續.
可以繼續到一個地步.
就是真的看著.
那個工作完成.
完成之後.

$^{601}$再有另外的課題.
第七章.
第六章告訴我們.
總共用了52天.
在很多嘲笑聲下.
完成了這個工程.
這個就是第六章.
最後那段經文.
告訴我們.
成祥修完了.
共修了52天.
我們一切的仇敵.
四圍的外邦人.
聽見便懼怕.
就愁眉不展.
因為看到這個工作完成.
是出乎我們的神.
這個就是在.
智慧行間.
近乎每一章.
尼兄米一面對困難.
馬上就仰天祈禱.
馬上就來到.
這個回憶錄.
差不多是他的自傳日記.
記載了那些事情.
到第四章.
我們上兩堂課都說過.
那些困難.
一面對神的工作.
四方八面湧過來.
一般其實.
你可以歸納.
危言聳聽.
嘲笑你.
想你和他同流合污.
給你一個西瓜皮踩.
我們上次也提過.
有些猶太人.
受了賄賂.

$^{641}$叫尼兄米.
找人.
密告給他知道.
有件暗中重要的事.
約你入聖殿談.
尼兄米一聽就知道.
是陰謀來的.
如果進去他就慘了.
就會被人用另外的說話.
來說他.
擅自無大無小.
他就去聖殿裡面.
他不是祭司.
他不是像以色列那樣.
又是文士又是祭司的身份.
他只不過是一個.
普通的猶太人.
不過是一個酒精.
不過是亞達斯西王.
差派他回去.
重建他祖宗的.
聖城.
他很.
看自己的位份.
我是一個猶太的弟兄.
你叫我入聖殿.
一定是陰謀.
我是不可以進去的.
你就可以想像到.
他思考的文節.
對自己的.
位份的步步為營.
是很值得.
我們去小心和留意.
我們去到第四章.
希望我們一路看完.
這十幾章聖經.
能夠再回頭再問.
你那點的火花.
是怎樣拼出來的.

$^{681}$你是由嬰兒餵奶的時候.
和你背誦聖經.
還是你上主一學.
阮媚教你讀主一學的時候.
唱兒歌的時候.
你來認識.
有那種對神的愛.
對神的家.
那個熱心.
但你長大之後.
世界很現實.
我現在在皇帝的面前.
伴君如伴虎.
我還要放長假.
放安息年假.
上去重建聖城.
其實很難的.
開始說那個難處.
第四章第一節開始.
參巴拉聽見我們修造城牆.
就發怒了.
大大怒行.
癡笑猶太人.
其實不是開玩笑的.
小時候上完主一學.
長大之後可能很出類拔萃.
不過就不再在教會裡出類拔萃.
很多人都不上教會.
很多人你問過他.
我小時候都上過兒童主一學.
那些詩歌我懂得唱.
不過我們千萬不要灰心.
因為我們作為主學老師.
當年為這些小朋友做過祈禱.
將來聖靈工作.
在他們的內心.
那粒種子就會發芽.
什麼時候發就沒有人知道.
所以千萬不要灰心.
做兒童工作最大的回報.

$^{721}$就是你自己.
你種的那粒種子.
將來在他內心.
就會發芽.
這件事就很奇妙了.
第二節.
對他弟兄和薩瑪利亞的軍兵說.
這些軟弱的猶太人.
作什麼呢.
薩瑪利亞的軍兵.
阿拉伯人.
參巴拉.
這些軟弱的猶太人.
作什麼呢.
在他們的眼中.
你們猶太人.
不過是軟弱的.
我們叫你做什麼你就做.
你們都沒有自己的身份.
沒有自己的國籍認同.
沒有自己那種的持守.
正正因為沒有了.
城牆的界限.
所以無牆.
是一個問題.
現在他們不是無牆.
是牆爛了.
所以現在要盡快連絡.
所以第三章千萬不要覺得悶.
是不同的領袖.
自己.
自己的部份.
一起去做.
如果大家都守住.
團契部又守住.
栽培部又守住.
猜拳又守住.
聖學又守住.
那麼教會就連絡得很整齊.
就不是其中一個很大.

$^{761}$其他就很弱.
那就不好.
你可以想像得到.
但這是你自己.
一廂情願的藍圖.
第四章.
作者就集中講.
記載了準備開工.
但是你一開工.
人家就要你停工.
你留意下第二節.
下半節.
要保護自己要憲制.
完成之後他們就會有.
那個聖殿摩西.
要求的那種規模.
那種分別為.
同一類型的憲制.
就是說外邦人.
你只能夠在外邦人院.
那裡去等候.
如果你想憲制就行.
你只能夠遠遠的.
你和我們是沒有份的.
我們猶太人的男丁.
才可以進去裡面去憲制.
女院就是.
姐妹.
這裡就告訴我們.
他們的嘲笑是.
是不是要一天就成功.
要從土堆裡面.
拿出火燒的石頭.
再立牆.
土堆就是說那些根本就爛了.
那些廢堆.
撿些磚再重新燒著.
就是說根本連材料都沒有.
阿滿人.
多比亞站在旁邊.

$^{801}$有阿滿人.
有薩馬利亞人.
有阿拉伯人.
有這幾個的.
外邦的領袖.
多比亞.
基善等等.
然後再嘲笑.
你一言我一語的嘲笑.
他們所修造的.
城牆石牆.
就是狐狸上去.
留意中文這裡和合本.
看不到.
英文是很清楚.
一隻跳上去都爛了.
這樣的意思.
就是狐狸上去也必.
癡倒.
這個是黎希米的自傳.
我們的神啊.
求你垂聽.
因為我們被藐視.
求你使他們的偉旁.
歸於他們的頭上.
使他們在老道之地.
作為掠物.
不要遮掩他們的罪孽.
不要使他們的罪惡.
從你面前逃沒.
因為他們在修造的人.
眼前惹動你的怒氣.
這樣來到.
當下向神祈禱.
有沒有舉手.
向天開聲祈禱.
有沒有這種姿態.
這種對抗式的祈禱.
如果在今天.
突然我們站在彌敦道.

$^{841}$大聲喊祈禱.
可能我們要立刻被人詢問.
是不是有思覺失調.
是不是什麼.
但當時黎希米就不是了.
因為很明顯這個就是聖城的範圍.
現在我們就在做.
神要我們做的工作.
更加是有.
亞達舍西王的詔書.
我們是合法的建造.
就是說批出了工程的合約.
我們現在是真的合法地去做的.
所以黎希米現在面對這些.
可能真的當面.
就是.
兜口兜面.
單單打打的說話.
他的回覆.
就是不激動.
不過向神.
大聲呼籲禱告.
一方面作為領袖.
又是讓.
一起建造的.
同胞.
弟兄姊妹.
看到你的勇氣.
看到你的持守.
看到你的肩負擔.
大家都不會退縮.
因為如果連黎希米自己都退縮.
自己都去.
怕這班人嘲笑.
那件事就很難做得成功.
所以.
他就好像化做一個公開.
向神去舉手禱告.
似乎挺生效.
第六節.

$^{881}$這樣.
我們修造城牆.
城牆就都連絡.
高至一半.
因為百姓專心作工.
就是說大家暫時.
在這樣的嘲笑下.
因為始終有.
正式的皇帝的.
詔書.
這些外邦人.
都不敢太明目張膽.
的挑戰.
或者用武力.
只是現在用一些.
很陰毒的.
嘲笑.
令那個士氣.
低沉.
令他們就好像覺得.
其實是.
我們搞這麼多事.
是不是大龍鳳.
沒有意思的.
最終都不可以成功.
又或者可能以前曾經都有人.
帶領過去建造.
但最終.
工程實在太大了.
阻力也太大了.
所以就做不成.
不過今次不同.
今次是真的有欽差大臣.
由河東.
阿達薩西王那裡而來.
所以今次百姓.
仍然是帶著.
一種可能是半信半疑.
不過已經現在.
初步的那個.

$^{921}$輪廓就出現了.
這是第六節.
所說的.
城牆都連絡好.
高至一半.
似乎可能用一個.
當時人的高度的一半.
可能就是兩三呎.
已經做好了.
有一個小小的初型.
但是第四章說的.
第五章說的.
就是那個困難.
怎樣來到成為.
真正的阻力.
第五章似乎是內部阻力出現了.
就是百姓.
沒有糧出怎麼辦.
因為個個都要.
做事.
可能耕種.
又會出現困難.
第七節開始.
參巴拉多比亞.
阿拉伯人 阿滿人.
又再多一些 阿什特人.
聽見修造耶路撒冷城牆.
著手進行.
堵塞破裂的地方.
就甚法路.
上文都已經在說.
嘲笑.
法路.
是參巴拉聽見.
現在參巴拉已經找到.
自己的外族.
聯繫的.
這些.
聯盟的人.
越說就越多.

$^{961}$不同的.
種族.
阿拉伯人 阿滿人.
或者阿滿人.
阿什特人.
很多都是.
猶太人的.
一些世仇.
過去約書亞的年代.
曾經臣服於.
以色列的.
現在成為了敵人.
繼續在周邊.
來騷擾.
不是安居.
大家的和睦.
共處.
很容易就會受他影響.
會沾染他們.
那些拜賈神的.
風氣習俗.
但現在因為見到.
如果被猶太人成功.
建造好城牆.
他們就在這個界限裡.
裡面.
來去有那個敬拜.
他們就會上下一心.
這個就是.
外邦人最大的困擾.
和.
最大的那個懼怕.
因為當.
這個民族有回自己的.
敬虔的時候.
很多事就可以做到.
亦都可以和外邦人.
正式的更加沒有任何的.
聯繫.
所以他們就.

$^{1001}$要想辦法.
怎樣來去拆毀.
但聖經用了兩次.
是甚為法路.
即是說他們.
不只是普通批評那麼簡單.
通常那些吃花生.
看熱鬧那些.
看完第二天就沒事了.
所以做那些公關.
你一定要守住那一天.
人們就慢慢不再留意.
報紙都不再賣.
是那幾天.
你守得住那幾天.
那個問題.
就會慢慢迎刃而解.
因為香港人很善忘.
第二第三天.
根本就沒有人再理會.
但是現在.
他們這個情況.
不只是嘲笑那麼簡單.
而是.
真是去到一種憤怒.
即是說.
對他們的民族.
是會帶來一個.
如果那個聖神.
是能夠有那個光彩.
敬拜.
是對外邦人一個威脅.
所以他們不只是嘲笑完.
就走了去就算了.
大家和睦共處.
而是他們會想辦法.
怎樣來去.
對那個士工的攻擊.
所以他們這個.
淫法路.

$^{1041}$就不只是說說.
所以第八節.
大家同謀要來攻擊耶路撒冷.
使到城內擾亂.
confuse.
用英文.
同謀.
跟著就是一個.
conspiracy.
陰謀.
英文聖經是這樣說.
這個擾亂是令到.
城內有confusion.
大家是.
因為一些.
他們用一些有組織的陰謀.
來去擾亂.
猶太人.
我們就逐樣去拆解.
怎樣來去影響.
然而我們禱告.
我們的神.
又因他們的緣故就派人.
看守就夜防備.
這個很現實地.
黎軒蔚都要保障.
還沒有建好的城牆.
那個.
安穩.
不會被人.
夜晚.
來去拆毀.
或者無緣無故的用一些推土機.
來去推倒它.
猶太人說.
灰土常多.
江台的人力氣.
已經衰敗.
所以我們不能建造城牆.
第十一節.

$^{1081}$我們的敵人且說.
不知不見我們進入他們中間.
殺他們.
使工作止住.
那靠近敵人居住的猶大人.
十次.
從各處來見我們說.
你們必要回到我們那裡.
所以你看到.
有幾方面.
黎軒蔚正在面對.
就是猶太人就說.
有些已經在做工作的人.
已經是一種.
很緊繃的.
狀態.
白天又要做.
晚上又要保持警戒.
所以導致.
開始已經有一種.
他們的內心.
很忐忑和很惶恐.
再加上.
有些地方.
有些工地.
有些住宅根本不夠人.
所以十次來.
叫派人調派人.
去.
某程度上.
是一個工地.
的地方.
好像臨時的.
貨櫃.
給工程師畫跡.
來看圖等等.
交代給.
管工.
派工作給.
不同前線的紮鐵.

$^{1121}$或者是做砌磚.
三行.
做電.
這個情況.
尼熙米坐鎮在一個中央的地方.
現在不同.
因為耶路撒冷都頗大.
不同工地的人.
來找尼熙米.
描述情況.
因為那些.
警告.
那些似乎.
開始的陰謀.
和有些可能是暗中的行動.
仍然是暗中的.
因為不敢.
公然的.
挑戰尼熙米.
其實他是一個.
省長.
來對尼熙米.
有一個.
明顯的攻擊.
所以尼熙米怎樣去拆解.
這些.
無論是必需的訴求.
又或者是.
百姓都現在是在一種.
很疲倦的.
心態,心靈開始有疲倦.
怎樣去幫助他們.
我們留意第十三節.
所以我使得百姓.
各按宗族拿刀.
拿槍拿弓.
站在城牆後邊低窪的空處.
我察看了.
就起來對貴就官長.
和其餘的人說.

$^{1161}$不要怕他們.
當紀念主事大而可畏的.
你們要為.
弟兄兒女.
妻子家產.
去爭戰.
很寶貴.
勉勵.
這些領袖.
叫這些領袖.
看到的.
不是表面的現實.
不夠人.
還有大家都疲倦.
仇敵是有陰謀.
要擾亂.
使得混亂.
使得我們來到這個工程.
自己內部.
都已經沒有了軍心.
但起碼叫他們.
沒錯現實是真的.
但不要只看現實.
當紀念.
主.
我們的神.
是大而可畏的神.
而且現在我們.
這一場的建造.
是為我們的兒女.
為我們的家庭.
為我們的家產.
更加為我們.
屬靈的身份.
神選民猶太人的身份.
我們是分別為聖.
我們是繼承.
我們的祖宗亞伯拉罕.
給我們的神的土地.
這些是.

$^{1201}$猶太人看得很重要.
我們就不是很強.
因為我們都是一些移民的.
後代.
1949年我爸爸也是這樣走難.
來香港.
那些故居被人沒收.
那些都沒有什麼感情.
我都不知道自己現在.
在哪個地方.
上去都不知道在哪.
因為從來沒有上去.
所以沒有這些感情.
猶太人就不是.
他們四百年復國.
一落飛機就親吻自己的土地.
知道是神的.
那個應許.
他們終於可以復國.
所以這個.
都是一個.
有智慧的領袖.
不要看人.
不要看現狀.
通常我們傳媒都說不要看人.
正正就是不開心.
才會看人.
但是有時.
應該可能.
用另一個角度.
你這個.
感覺和遭遇.
被人這樣傷害.
無論說話.
具體的一些事件.
但是都要想念神.
大而可畏.
真正可以.
使我們來面對.
那個.

$^{1241}$懼怕的是這位造物主.
所以我們真的要怕神.
敬畏就是怕.
這樣的意思.
對神真的要有一種.
敬畏的心靈.
很久之前.
可能我都說過.
陳友遞兄.
來港見證.
一起吃飯在會所.
如履崇拜.
好像他和太太一齊來.
在席間.
我坐他旁邊.
聊天.
他很尊重牧者.
因為他太太真的令到他.
信心 信主等等.
他說他們一班朋友.
溫那那些.
很熟很熟很熟.
一說到神.
個個都不敢放肆.
他們玩得很厲害.
很熟.
什麼都說.
但是一說到神.
他們一班朋友.
馬上很嚴肅.
阿倫.
阿B.
真的很認真.
不少說.
神不少說.
我們想像得到.
有些不信神.
有些信佛.
但陳友遞.
表達到一種心態.

$^{1281}$讓他們看到.
所以連他的朋友都知道.
不要說笑.
一說到神.
大家都不敢開玩笑.
這都是一個很重要的信念.
所以越小.
他有那個種子.
就越緊要.
對神的認識.
敬畏.
不是呼之則來.
揮之則去.
這種敬畏.
包含了一種愛.
他永遠是造物主.
我們永遠是受造之物.
我們最好擁有越來越多的知識.
對這個世界.
宇宙的結構.
越了解越好.
我們就知道創造這一切的造物主.
就更加偉大.
更加超越.
有時詩歌盛載的歌詞的內涵.
正正告訴我們.
已經不能用言語.
用文字.
盛載造物主的偉大.
有時只能用很偉大.
盛載盛載盛載.
盛載何其大.
要這樣來到.
只能用這樣.
嘩這樣表達.
其實我們來崇拜.
健康的.
嘩這樣敬拜神.
其實是健康的.
死氣沉沉的.

$^{1321}$快點睡覺.
快點說完我喝茶.
餐餐堂堂都是這樣.
這樣才不健康.
偶然間.
令到你有一個震撼.
神現在在聚會當中.
我感覺到.
這種震撼.
是我們福音派有時需要的.
所以絕對不是靈恩派.
有時過多.
又變成一種造作.
很假.
但真正發自內心的敬拜.
其實偶然間我們需要有.
譬如唱到某些詩歌.
如果有個弟兄.
平時很嚴肅.
很理性.
不停流眼淚.
你不要以為他有什麼問題.
他是敬拜.
唱詩歌勾起他內心的敬拜.
由腦筋的理性.
轉化成感性.
這些就是敬拜.
這些就是聖經告訴我們.
我們偶然間應該要這樣.
不是造作.
不是裝模作樣.
不是要擠眼淚出來.
而是發自內心的.
對神的敬拜.
你不是告訴我們.
不要怕他們.
那些是現實.
我們當然要拿刀拿槍拿弓.
守護著我們的城牆.
因為還未完成.

$^{1361}$我都留意了.
等等.
不過更加寶貴就是.
我們要為主而戰.
我們要為.
國民身份.
神選民身份.
我沒有意思批評.
黃絲藍絲.
但一定追求的.
如果只是地上的國度.
其實不足夠.
我們是神的國度.
我們作為牧者就更加.
義不容辭.
告訴大家我們有神的國度.
我們追求的是永恆.
真正聖潔.
真正可以讓人.
進入平安.
盼望的國度.
所以我沒有意思批評地上的國度.
我要爭取什麼.
因為我不懂.
我們是超越.
應該讓黃的.
藍的請大家坐定.
敬拜真神.
追求神國度.
這些永遠都對的.
必須的.
如果不是的話.
教會有人叫過我.
不如讓這個教會.
有一堂崇拜是專給年青人.
如果是黃的就在那裡聚會.
我不可以這樣.
教會應該任何人.
任何政見都要.
坐在一起敬拜.

$^{1401}$因為我們是弟兄姊妹.
我們是有神的生命的弟兄姊妹.
絕對不是.
同一個政見坐在一起.
不是.
所以我絕對不贊成我們的教會.
人家說我們是藍營.
我們不是.
我們是神國度的教會.
這個是很重要.
大家要知道.
我們不是為政治去效力.
我們是神的國度去效力.
所以這些錄音.
被人聽到就更加好.
被人知道我們追求的是.
神的國度的福音.
去改變.
但不等於地上有不同的人.
他有不同的神的帶領.
那個就不是.
我可以去干預.
我的責任就是要讓.
我們的同工我們的弟兄姊妹.
一起去追求福音.
改變人的生命.
這個就是我的國度.
所以當時黎兄姊妹他們.
在這麼大的困難下.
沒錯.
仍然是復興的開始.
因為最大的大前提.
神的刑罰已經結束.
七十年的奏坐.
七十年的刑罰已經過去.
現在是要重建的時候.
所以黎兄姊妹.
就在這個復興的浪潮.
第三代.
來到去.

$^{1441}$重建.
最後的城牆.
然後輾轉.
神選民身份的.
歸屬就慢慢形成.
就去到耶穌的年代.
真是.
他們會對自己的身份.
相當強烈.
當然又是另外一個極端.
就是法利賽主義的興起.
第十五節.
仇敵聽見我們.
知道他們的.
心意.
見神也破壞他們的計謀.
也破壞他們的計謀.
甚至他們也感覺到.
有神與他們同在.
我們都回到城牆那裡.
各作各的功.
從那日起.
我的僕人一半作功.
似乎是他們自己的.
一些親信.
他們帶上來的一些僕人.
即是說他們真是一個.
顯赫的酒精.
是可以動用.
國家的庫房.
去到河西才去購買.
去到河西的省長那裡.
去調度那些資源.
木材,石頭.
去建造.
這裡告訴我們.
我的僕人一半作功.
一半拿槍,拿盾牌,拿弓.
穿那些盔甲.
官長都站在.

$^{1481}$猶大眾人的後邊.
修造城牆的.
鋼台材料的.
都一手作功.
一手拿兵器.
修造的人都腰間配.
都修造.
吹角的人在我旁邊.
吹角即是說.
似乎是祭司.
才可以做吹角.
亦都告訴我們.
他們都會有一種.
彼此的呼應.
那時不會有WhatsApp.
不會有對講機.
所以用吹號.
來做一個通知.
令到大家.
都可以.
步伐能夠一致.
似乎這裡都是一個.
很緊張的狀態下.
所以那個張力.
不是說現在人家給你.
整個地方.
你喜歡怎樣做就怎樣做.
你現在完全圍封了.
你多少年都來做.
都沒有問題.
只有張力.
所以事奉都有很大張力.
如果你說.
每樣東西都要有一個.
很平衡的平衡點.
有時事奉就未必能夠.
提供到.
即是.
大家如果堅持.
某個程度.

$^{1521}$我要commit我自己.
參與星期三晚的.
教會祈禱會.
某程度在你的工作上.
一定有張力.
如果你還是上有高堂.
下有兒女.
在讀書.
其實比較難.
到你退休.
你的子女都長大了.
你的時間比較彈性自由.
相對又沒那麼難.
所以其實每個事奉.
只要你commit.
一定有張力.
但這裡你會見到.
他們在這個建造裡面.
充滿了事奉.
參與上的張力.
是很不容易的.
隨時都好像要面對爭戰.
要拿著刀拿著盾牌.
即是說可能有些.
是負責守衛的.
有些是本身.
對於木工.
石工能夠有.
那個恩賜的.
可能他們就是做建造.
即是說有旁邊的人.
來成為他們的護衛等等.
但也仍然告訴我們.
在這一章裡面.
仇敵.
是正正式式地出現在.
那些侍工裡面.
所以侍工計劃.
其實應該要計劃.
那個挫敗.

$^{1561}$那個損耗的時間.
心靈的.
那個大家的低沉.
怎樣來刺激他們.
這些.
應該是侍工計劃裡面.
好像從來都沒有什麼著墨.
好像.
一寫完就好像快完成了.
就好像.
做到了.
但其實往往很多東西都是這樣.
來到去.
很漂亮很亮麗的.
侍工計劃.
但很快可能都不用一個月.
大家都不了了之.
開完會就當完成了.
原因相信.
就是我們很少從第四章.
那裡找到那個靈感.
就是那個張力.
是由.
參與事奉.
用的時間.
用的奉獻.
無論金錢.
時間物質.
再加上敵人.
這裡有明顯的就是外邦人.
有實實在在的城牆.
但參與做神.
的工作的新約.
我們去傳福音.
希望大家同心.
到底哪些是.
我們的朋友.
我們的弟兄姊妹.
哪些是嘲諷.
在這裡看熱鬧.

$^{1601}$自己一點都不參與.
不做還要說人做的人.
我們很難表面分到.
也不應該這樣分.
但.
實實在在的真的面對很多.
的閒言閒語.
又怎樣能夠讓他.
來到去將那些閒人.
的閒語.
改變他成為.
參與者事奉者.
這個又是很高.
很高的那個技術.
要來到去我們大家.
來到去一齊的同心.
相信.
這張聖經李起米一而再.
他不是計劃.
怎樣動用軍隊.
動用他自己的.
走政這個身份.
有的.
因為後面第七章.
他就說他自己是走政.
那個身份他就回來.
諸如此類.
後面他就會說.
但我們看第四章見到的.
一味就是祈禱.
望著天和.
同胞說.
我們要紀念我們的上主.
是大而可畏的.
不要只是看.
現在的困難.
不要只是看那些敵人嘲笑.
他一直都存在.
就是說他們算什麼.
這樣就真的會.

$^{1641}$令到和他一起做工的人.
有個方向.
和有那個.
大家的同心.
這個是李起米很成功的地方.
他將這些回憶.
擺在他的六章.
一到第六章的回憶錄.
裡面成為我們後世.
都是可以.
從其中學習.
李起米記不是教我們做領導學.
李起米記是說一個有心人.
他內心的一團火.
成為了一個事工.
就是神的.
選民.
真是有了界限之後.
就慢慢.
build up 自己.
sense of belonging.
選民的歸屬感.
所以我們.
錦宣家.
要建立我們自己的歸屬.
真是.
很容易.
引自己是錦宣家為榮.
這個就很成功.
當然.
不算什麼的.
因為比起全個香港的.
宗教會.
全個世界的宗教會.
古往今來的聖徒.
我們算得什麼呢.
但至少這是一個起點.
至少這個起點就是.
現在這個教會有純正的福音.
有救人靈魂.

$^{1681}$為榮.
我又能夠參與.
我以此為榮.
成為一個起點.
再去看其他教會.
其他的弟兄姊妹.
甚至古往今來的.
神的工作.
你就會更加與有榮焉.
我們是做物主的.
在耶穌裡面拯救的.
祂的兒女.
人生就是這麼多.
很快就過去.
我們就進入永恆.
我們就成為在天上.
看著你們修成的行列.
所以現在.
我們未去見主之前.
我們就應該.
要以自己是神選民為榮.
要有這個歸屬感.
所以黎兄姊妹.
就是迫切要建造城牆.
就是要建立這個歸屬感.
第二十節.
對不起,第二十節.
你們聽見各星在那裡.
聚集到我們那裡去.
我們的神.
必為我們爭戰.
於是我們作弓.
一半拿兵器.
從天亮直到星宿出現的時候.
那時.
我又對百姓說.
各人和他的僕人.
當在耶路撒冷住宿.
好在夜間保守我們.
白晝作弓.

$^{1721}$不要分散在.
其他地方的.
小小的城市.
主要集中建造耶路撒冷.
這個聖城的城牆.
所以要招聚多些.
的工人在耶路撒冷.
留宿.
這個也可以說是很現實地.
看到那個實際的需要.
晚上可以多些.
的人可以來到.
看守著那個未完成.
的工作.
這個告訴我們.
黎希米是一個很務實的.
一個領袖.
這樣.
我和弟兄僕人並跟從我的護兵.
刀不脫衣服.
出去打水.
也帶兵器.
即是說長期在這個階段.
都是一個.
很備戰的狀態.
這些都是.
我們必須要.
學習.
弟兄姊妹千萬不要.
每個禮拜日.
那一個半小時.
是基督徒的身份.
要有一個警告.
由我敬拜完.
離開會堂.
我應該帶著基督徒的身份.
和別人交往.
我的談吐.
都是表達出一份的.
聖徒的體統.

$^{1761}$我的價值觀.
都是要滲透基督的真理.
這個就是要不斷.
慢慢擴散在.
生活的其他環節.
這個是重要的.
一個備戰.
否則我們每個禮拜回來.
見到一群.
六天都是世界仔.
六天半.
坐在那裡一個半小時.
是很莫名其妙的.
好像對神很陌生.
其實很難.
對話.
現在他的心都不得而撤.
但如果我們已經在那個氛圍裡面.
大家都是一種.
備戰狀態.
那個備戰不是神經質.
那個備戰就是你的內心.
真是很強烈.
慢慢越來越有一種歸屬.
我現在已經.
去到這年紀.
不是什麼流行歌星.
MIRROR那些粉絲.
根本什麼都不是.
但如果你.
很簡單的.
你見到現在的青少年.
他是哪一個粉絲.
他的談吐衣著.
根本就完全追之若無.
我都講過我是張國榮的.
粉絲.
那時真的學他.
穿白西裝.
其實很搞笑.

$^{1801}$你又不在別人的環境生活.
你又要扮到.
好像那麼時下.
因為.
他是偶像.
你又可以知道.
李希敏.
她的成功.
就是.
雖然未有那個硬件.
但已經在滲透給.
猶太人知道.
其實你們是屬於神的.
建造城牆.
連絡他.
有我們自己的敬拜.
有合乎神心意的敬拜.
這個是.
真的是重中之重.
這樣讓猶太人.
有這種身份的認同.
那就不只是.
做工程那麼簡單.
打散工那麼簡單.
而是.
其實牧者就是很想.
無論你去哪裡.
這是你的帶領.
但你需要有對自己的基督徒.
選民身份的認同的歸屬.
這就是最重要的.
就是說.
如果有一天.
那個時代神安排.
我們這個時代是要坐牢的.
不要緊.
那個監就是你敬拜身份認同的地方.
沒有問題.
或者你要去到南極.
或者你要去到南美洲.

$^{1841}$諸如此類.
當然多數去英國或者去加拿大.
不要緊.
你一定要有身份的認同.
和歸屬.
你既然已經去到別人的國家.
你不要只是看自己是香港人.
只是一班香港人.
追在一起.
這樣就不好的.
你應該將福音帶遠一點.
帶更多的人認識我們的主和福音.
或者認識我們.
美好的.
德行的價值觀.
如果我們說香港人值得驕傲的.
就將真正的特質.
譬如廉潔.
守法.
就要推銷出來.
不只是你這班人.
自己要走在一起.
而是你要讓內心的美.
散發出去.
所以牧者其實要.
這樣去推動.
所以其實要求問的.
現在是什麼時代.
現在的時代.
是不是尼希米神的.
那個就坐刑罰.
過去.
他們要被感動.
其實每個人都應該要感動的.
猶太人要回去重建的.
但是.
已經不是.
耶利米那個年代.
你怎樣流眼淚.
怎樣祈禱神都不會聽.

$^{1881}$因為那個是神決心.
刑罰就坐.
猶太選民的時代.
所以耶利米蘇.
的情況就是.
叫那些人不要再祈禱了.
不要再抵抗了.
要快點被老歸回.
去巴比倫.
這樣教訓會眾的.
每個星期叫你投降.
接著就被人.
說他是賣國賊.
扔了在膚井.
這是耶利米蘇的經文.
我忘記了是.
第幾章了.
我回去告訴大家.
你可想而知.
有時代的背景.
尼希米就是.
生在一個復興的世代.
他就是被神選中的.
一個領袖.
來帶領.
復興的浪潮.
這就是重建的意義.
這就是.
要將猶太人.
已經是.
分散得很厲害.
沒有了很強烈的.
身份認同的歸屬.
尼希米就可以在他那個時代.
透過他的影響力.
屬靈的影響力.
帶領他們仰望神.
開始建立了.
這個歸屬感.
第七章就開播了.

$^{1921}$真的看到一些具體的.
內部問題.
逐樣逐樣跟進.
所以就要那個歸屬感.
要慢慢這樣來到.
才可以開始由那個建造成長.
開始.
今晚就要早一點.
因為我等下要開會.
九點鐘.
多謝天父給我們.
今天雖然只有一個小時.
我們可以來到第四章.
簡單的來講解.
很盼望這個.
屬靈的歸屬感.
我們是神選民.
是主耶穌用寶血.
買贖回來的.
基督徒的身份.
能夠很強烈的認同.
這個認同.
不是對其他未信的人.
一種排擠.
和侮辱.
反而我們善意的.
表達我們那份.
心裡面專注為聖的.
那種.
上上準備.
來分享福音.
的這份熱誠.
不是對人有一種.
感覺很陌生.
或者是壓迫感.
的熱誠.
而是發自內心的對別人的.
關愛和尊重.
因為我們知道上主.
都給人有一個自由意志.

$^{1961}$去揀選福音.
所以我們都需要求主.
給我們能夠有這一份.
對自己身份認同.
的強烈的歸屬.
但是對別人.
我們有一種空間.
給人有一個感到.
舒服的慢慢的.
去思考去迎順.
願主幫助.
多謝你給我們一起.
有機會去研讀.
這些舊約古代的經文.
經卷.
讓我們能夠透過.
他們那個時代的見證.
我們都能夠有.
自己的思想.
在我們今天香港.
現在我們要.
尋問我們現在.
在疫後大家一起.
要躍動的時候.
幫助我們在.
參與.
可能不同的活動.
或者是一些旅遊.
但是我們更深的.
那個生存.
都不只是為了旅遊.
而是我們內心的生命.
都是為了可以.
見證主的福音.
旅遊就成為了一個.
我們可以心靈更新的.
一個活動是極有意義.
求主幫助.
我們仰望你.
禱告奉耶穌基督的名.

$^{2001}$阿們.
(無人回應).
\newpage



\section{哥林多前書 15:12-34}
\label{sec:qxsLxZod7kk}
\textbf{2023-04-22 晚堂《基督 ─ 我們的初熟果子》哥林多前書 15\_12-34 講員:黃頌言姑娘}
\newline
\newline
連結: \href{https://youtube.com/watch?v=qxsLxZod7kk}{\texttt{https://youtube.com/watch?v=qxsLxZod7kk}} ~~~~ 語音日期: 2023-04-25
\newline
\newline
\hyperref[sec:SLUSAYk3lm4]{\small{< < < PREV SERMON < < <}}
~
\hyperref[sec:index_chronic]{\small{[返順時目]}}
~
\hyperref[sec:index_scriptual]{\small{[返順卷目]}}
~
\hyperref[sec:rnbJeDm_hDc]{\small{> > > NEXT SERMON > > >}}
\newline
\newline
哥林多前書 15:12-34
\newline
\begin{longtable}{cl}
\hline
\hline
章節 & 經文 (和合本修訂版)\\
\hline
15:12 & \begin{tabularx}{0.7\textwidth}{X} 既然我們傳基督是從死人中復活了，怎麼在你們中間有人說沒有死人復活的事呢？ \end{tabularx} \\ \\ \relax
15:13 & \begin{tabularx}{0.7\textwidth}{X} 若沒有死人復活的事，基督就沒有復活了。 \end{tabularx} \\ \\ \relax
15:14 & \begin{tabularx}{0.7\textwidth}{X} 基督若沒有復活，我們所傳的就是枉然，你們所信的也是枉然。 \end{tabularx} \\ \\ \relax
15:15 & \begin{tabularx}{0.7\textwidth}{X} 這樣，我們甚至被當作是為神妄作見證的，因為我們見證神是使基督復活了。如果死人真的沒有復活，神就沒有使基督復活了。 \end{tabularx} \\ \\ \relax
15:16 & \begin{tabularx}{0.7\textwidth}{X} 因為死人若不復活，基督也就沒有復活了。 \end{tabularx} \\ \\ \relax
15:17 & \begin{tabularx}{0.7\textwidth}{X} 基督若沒有復活，你們的信就是徒然，你們仍活在罪裡。 \end{tabularx} \\ \\ \relax
15:18 & \begin{tabularx}{0.7\textwidth}{X} 就是在基督裡睡了的人也滅亡了。 \end{tabularx} \\ \\ \relax
15:19 & \begin{tabularx}{0.7\textwidth}{X} 我們若靠基督只在今生有指望，就比所有的人更可憐了。 \end{tabularx} \\ \\ \relax
15:20 & \begin{tabularx}{0.7\textwidth}{X} 其實，基督已經從死人中復活，成為睡了之人初熟的果子。 \end{tabularx} \\ \\ \relax
15:21 & \begin{tabularx}{0.7\textwidth}{X} 既然死是因一人而來，死人復活也因一人而來。 \end{tabularx} \\ \\ \relax
15:22 & \begin{tabularx}{0.7\textwidth}{X} 在亞當裡眾人都死了；同樣，在基督裡眾人也都要復活。 \end{tabularx} \\ \\ \relax
15:23 & \begin{tabularx}{0.7\textwidth}{X} 但各人是按著自己的次序復活：初熟的果子是基督；然後在他來的時候，是那些屬於基督的。 \end{tabularx} \\ \\ \relax
15:24 & \begin{tabularx}{0.7\textwidth}{X} 再後，終結到了，那時基督既將一切執政的、掌權的、有權能的都毀滅了，就把國交給父神。 \end{tabularx} \\ \\ \relax
15:25 & \begin{tabularx}{0.7\textwidth}{X} 因為基督必須掌權，等神把一切仇敵都放在他的腳下。 \end{tabularx} \\ \\ \relax
15:26 & \begin{tabularx}{0.7\textwidth}{X} 他要毀滅的最後仇敵就是死亡。 \end{tabularx} \\ \\ \relax
15:27 & \begin{tabularx}{0.7\textwidth}{X} 因為經上說：「神使萬物都服在他的腳下。」既然說萬物都服了他，那使萬物屈服的，很明顯地是不在其內了。 \end{tabularx} \\ \\ \relax
15:28 & \begin{tabularx}{0.7\textwidth}{X} 既然萬物服了他，那時，子也要自己順服那叫萬物服他的，好使神在萬物之中，在萬物之上。 \end{tabularx} \\ \\ \relax
15:29 & \begin{tabularx}{0.7\textwidth}{X} 不然，那些為死人受洗的，能做甚麼呢？如果死人不會復活，為甚麼替他們受洗呢？ \end{tabularx} \\ \\ \relax
15:30 & \begin{tabularx}{0.7\textwidth}{X} 我們為甚麼要時刻冒險呢？ \end{tabularx} \\ \\ \relax
15:31 & \begin{tabularx}{0.7\textwidth}{X} 弟兄們，我在我們的主基督耶穌裡，指著你們—我所誇的極力地說，我天天冒死。 \end{tabularx} \\ \\ \relax
15:32 & \begin{tabularx}{0.7\textwidth}{X} 從人的觀點看來，我當日在以弗所同野獸搏鬥，對我有甚麼益處呢？如果死人沒有復活，「讓我們吃吃喝喝吧！因為明天要死了。」 \end{tabularx} \\ \\ \relax
15:33 & \begin{tabularx}{0.7\textwidth}{X} 不要被欺騙了；「濫交朋友敗壞品德。」 \end{tabularx} \\ \\ \relax
15:34 & \begin{tabularx}{0.7\textwidth}{X} 你們要醒悟為善，不再犯罪；因為有人不認識神。我說這話是要使你們羞愧。 \end{tabularx} \\ \\
[1ex]
\hline
\hline
\end{longtable}
$^{1}$各位弟兄姊妹平安.
我們一起祈禱.
親愛的主.
多謝你.
我們很渴望你的話語.
願你也跟我們說話.
讓我們有一個謙卑的心.
多謝你.
奉耶穌名求.
阿們.
我有一個問題想問大家.
我想問.
你相信有死人復活嗎?.
如果我告訴你.
真的有一個死了的人復活.
可能你會想.
「蛤?思辨?很難以置信」.
不如我統計一下.
如果覺得有死人復活.
請舉手來看看.
也有不少.
如果覺得沒有.
不會的.
請舉手來看看.
也有.
那麼.
你相信基督復活嗎?.
剛剛大家聽了主席讀的經文.
所以大家舉手也非常小心.
因為本月的月題.
是「再施基督的生命」.
究竟基督的生命跟我們有什麼關係呢?.
究竟.
為什麼復活這麼重要呢?.
為什麼要慶祝復活節呢?.
你有沒有想過如果.
基督沒有復活的話.
世界會變成怎樣呢?.
又會有什麼不同呢?.
除了復活節可能只是慶祝復活兔和復活蛋之外.

$^{41}$如果基督沒有復活的話.
我們所傳的福音又會有什麼不同呢?.
基督的故事會不會只停留在釘十字架呢?.
如果是這樣的話.
其實這是一個悲慘的故事.
或者.
以利亞可能比基督更厲害.
因為她沒有經歷過死亡.
她就直接被提起.
在我小時候大概中學.
我真的有想過這個問題.
如果沒有基督復活的話.
我們真的很可憐.
因為我們在信什麼呢?.
如果基督沒有復活的話.
就是沒有永生的盼望.
生命的終結就是死亡.
死之後就是完結.
死就是滅亡.
其他的宗教也都會說輪迴.
他們會說做好事的話.
就會轉世.
下一世就會好一點.
在道德上.
如果這樣的話.
都可以警惕一下人.
但是如果死了就這樣完結.
我們整個人生觀就會很不一樣.
可能我們都會覺得.
其實都不需要那麼節制.
不如放任.
玩盡一點吧.
生命完就沒有了.
也都不需要為自己做過的事情去負責任.
今天我們一起來看看哥林多前書.
保羅寫信給哥林多教會.
因為當時教會充斥著不同的學說.
保羅要寫信給當時的教會.
保羅特別在書信最後.
其實他用了一章.

$^{81}$是第十五章.
去講論基督復活的事.
不如我也講一講.
講一點哥林多教會的背景.
原來當時哥林多這個地方.
是被羅馬帝國統治的.
哥林多是位於希臘半島.
是重要的商貿城市.
在公元前44年.
就是耶穌來之前44年.
羅馬重建哥林多做殖民之用.
所以人口很多.
很急升.
哥林多從城市的設計.
律法.
管治.
慢慢變成了羅馬帝國的都會.
所以稱為羅馬哥林多.
社會發展得很快.
在第一世紀.
哥林多是一個提供社會上行流動的場所.
不同的階級.
有很多不同階級.
他們往上爬.
去增進.
去互相踐踏.
哥林多不單是多元的社會.
更加是多元宗教.
有很多像保羅這樣的人.
來到哥林多.
他們都會去聆聽不同的人分享.
所以在部分的哥林多人眼中.
其實保羅是其中一個來尋找機會的人.
所以我們可以想像得到.
哥林多人的信仰.
其實參雜了很多不同的學說.
以及不同的社會文化.
都參雜在裡面.
保羅在公元的50至51年.
去到哥林多.

$^{121}$基督復活升天之後的18年左右.
哥林多教會被社會影響.
因此很多人都會踩低人.
去抬高自己.
有很多等等的倫理問題.
分黨分派.
而且他們中間也會有.
不信死人復活這樣的事.
原來這樣的想法.
是非常影響他們的倫理價值的.
今天我們一起來看哥林多前書第15章.
12至34節中間.
我都截錄了一部分.
讓我們一起去反思基督復活的意義.
我們先看12至19節.
12節這樣說.
既然我們傳基督是從死人中復活了.
怎麼在你們中間有人說沒有死人復活的事呢.
若沒有死人復活的事.
基督就沒有復活了.
基督若沒有復活.
我們所傳的就是謊言.
你們所信的也是謊言.
基督若沒有復活.
你們的信就是徒然.
你們仍活在罪裡.
就在基督裡睡了的人也滅亡了.
我們靠基督旨在今生有子望.
就比所有的人更可憐了.
剛剛我也問大家.
會不會相信死人復活呢.
原來不信死人復活.
是等同於否定基督復活.
第17節說.
你們信的就是徒然.
這句意思是指虛空無用.
對人毫無價值.
舊約聖經用這個詞語來形容.
假神和偶像.
雖然基督的死和復活.

$^{161}$在歷史事件來說可以區分為兩件事件.
死和復活可以當作兩件事.
但作為拯救的行動.
基督的死和復活兩者是不可以分割的.
假如基督只是死了而沒有復活.
祂的救恩就沒有完成.
福音也沒有力量去拯救我們.
信的人和不信的人.
都有相同的命運.
就是滅亡了.
基督被釘十字架的功效.
是在神是基督耶穌律復活的前提下.
才有意義的.
因為基督的死代表祂是完全的人.
祂會死亡的.
同時祂也是完全的神.
因為祂是從死裡復活的.
因此我們在《使徒信經》裡面.
用非常大的篇幅.
我們剛才也讀過了.
去形容基督的人性.
我們也看看.
剛才我們讀過的《使徒信經》.
我們一起讀紅色的部分.
也看到紅色.
我們一起讀紅色的部分.
1 2 3.
我主耶穌基督.
上帝獨生的子.
因聖靈感恩.
由童貞女瑪利亞所生.
在本雕彼拉多手下受難.
被釘於十字架.
受死埋葬.
降在陰間.
第三天從死人中復活.
升天.
坐在全能父上帝的右邊.
將來必從那裡降臨.
審判活人死人.

$^{201}$我信身體復活.
我信永生.
原來祂和人.
祂是人來的.
祂死了.
降在陰間.
最重要的是祂三天從死人中復活.
因此我們每個星期都回來.
聚會崇拜.
我們其實是慶祝基督的復活.
這是我們對基督的信心的宣告.
原來我們回來去敬拜祂.
我們是一種宣告.
是我們一起慶祝祂的復活.
因著基督.
我們能夠和神和好.
所以.
末後會有審判.
我們都相信身體復活.
我們相信永生.
這是我們所相信的.
我們回看剛才第17節.
它說.
否則我們仍活在罪裡.
沒有力量得勝罪.
就是在基督裡.
碎了的人也會滅亡.
死後就永遠和神分開.
如果我們的盼望只是今生.
其實就等於沒有盼望.
假如將來沒有復活的話.
信基督的人就是全人類當中最可憐的.
我們去想想.
如果我們用一生去背十假.
我們一生都要為二壽逼白的話.
為甚麼我們要為福音去受死呢.
為甚麼我們要這麼堅持呢.
為著福音的緣故.
我們這麼堅持去背十假呢.
為甚麼我們不只是盡情享受今生.

$^{241}$如果我們不是有未來的盼望.
相信基督未來有公義的審判的時候.
基督徒的堅持就會是最愚蠢的.
保羅在這裡用了非常大的力度.
去表明自己所傳的復活的基督.
是千真萬確的.
我們去看看第三節.
它說當日所傳給哥林多教會的.
是照聖經所說的.
基督為我們的罪死了.
而且埋葬了.
又照聖經所說的第三天復活了.
基督復活是救藥.
基督復活其實是救藥.
以賽亞書對彌賽亞的預言的成真.
我們看看第六節.
它提出基督復活之後.
向五百多位弟兄去顯現.
其中一大半到現在都在這裡.
在他那時候.
就是基督死了之後的十八年.
那時候那些人還在這裡.
被顯現給五百個人看.
有部分他還在這裡.
但是有部分已經死了.
第七節.
之後基督顯現給亞國看.
再顯現給眾人看.
眾使徒去看.
第八節最後也顯給我看.
保羅很強調自己親身見過基督的顯現.
他說是真的.
保羅指出基督事實上已經復活了.
否定死人復活.
和否定基督復活是沒有分別的.
因此否定死人復活是不合理的.
保羅論證的核心是基督復活意味著.
信徒將來都要復活.
因為復活的基督是首先獻給神的.
我們看看二十節.

$^{281}$我們一起讀.
1 2 3.
其實基督已經從死人中復活.
成為歲了之人初熟的果子.
這裡的初熟果子.
原文翻譯的意思是獻祭的開始.
意思是投生的濟生.
或者我們可以理解為供品的第一批收成.
即是成為獻給神的各種禮物.
有上好的意思.
這個初獻的部分很有代表性和保證性的意義.
表示其餘部分的都是聖潔屬於神的.
因此用農作物收成的意象.
將基督的復活和信徒的復活串連起來.
復活的基督是首先獻給神的.
信徒就好像其餘的收成一樣必會復活.
基督是初熟的果子.
就是首先復活上好的祭物的意思.
第二十一節說.
既然死是因一人而來 死人復活也因一人而來.
在亞當裡眾人都死了.
同樣在基督裡眾人也要復活.
第二十三節.
但各人是按著自己的次序復活.
初熟的果子是基督.
然後在他來的時候是那些屬於基督的.
亞當成為人類的代表和元首.
經歷到人人都經歷的事.
就是因犯罪而死亡.
同樣耶穌是所有屬於他的人的代表和元首.
他經歷了復活.
也是屬於他的人會經歷的事.
全人類都是在亞當裡.
但不是都在基督裡.
保羅不是說所有人都會在基督裡得到生命.
而是說在基督裡所有人都會得著生命.
重點是在基督裡才可以得著生命.
第二十三節.
表明基督來臨的時候.
只有那些屬於他的人才會復活.

$^{321}$而且所有的國度權柄都是歸在基督的權下面.
第三十二節.
那裡說.
如果死人沒有復活.
讓我們吃吃喝喝吧.
因為明天要死了.
不要被欺騙了.
濫交朋友.
敗壞品德.
你們要省悟為先.
不再犯罪.
因為有人不認識神.
我說者說是要使你們羞愧.
那些否定死人復活的信徒.
他們的心態是將來都沒有盼望的.
因此要盡情於今生的歡愉.
做了很多敗壞道德的事.
原來若果神學錯了的時候.
是非常之嚴重的.
整個生命的方向.
活著的方向都是錯誤的.
保羅這裡提醒哥林多的信徒.
不要受騙.
他們會與一些不信的人為伍.
他們會跟從他們.
以及一些社會的風俗.
以及會拜偶像.
保羅這樣說.
主要是針對當時社會地位高的.
年輕的男子.
因為他們會參加私人的派對.
逢場作慶.
吃喝嫖妓.
當時沒有教會建築物.
因此他們在富裕的家庭裡.
聚會愛賢.
原本這樣看來其實也不錯.
我們一起去愛賢.
但是原來因此.
教會當中成為了.

$^{361}$因主受助的關係.
教會背後是不同的因主.
角力較勁.
去獲取榮譽.
他們公認窮困.
為的是確立自己的地位和榮譽.
愛賢和公認窮困原本是好的.
所有表面上都是好的東西.
但是他們的背後動機.
就出現了很大的問題.
父母以盛餐愛賢.
去說明自己的高位財富.
甚至在外邦的祭祀裡面去吃肉.
為了得到社會的肯定.
當時也會有亂倫.
等外邦的價值.
這些是教會不可以接納的東西.
但是他們不單止沒有將他們趕出教會.
因為這些人是有財富地位的.
因主.
就是他們有錢的人.
在他們的世界叫做因主.
這些因主為教會提供需要的地方.
在那裡聚會.
對因主他們會選擇.
不去指責他們的罪行.
這樣就顯示了哥林多教會.
是一個以社會連結.
和實際利益作為考慮為前提的.
他們的信仰同情和倫理行為.
不是首要的考慮因素.
他們認為贊助人受保護.
免於批評.
不至於逐出教會.
這個是最好最合宜的.
保羅形容他們是.
有些人是不認識神.
原文直譯就是.
有對上帝的無知.
這些人.

$^{401}$或者會自誇自己有知識.
但是.
從他們否定.
復活這件事.
就知道他們其實對神一無所知.
不知道大家有沒有看一套紀錄片.
叫做《以神之名》.
我看了幾集.
那時候是由70至80年代開始.
原來有很多不同的邪教出現.
他們都是拿著聖經.
和我們一樣的聖經.
但是曲解了聖經.
到現在仍然有人繼續相信這些邪教.
他們讓人陷入罪裡面.
犯奸淫.
欺騙回眾的金錢.
控制不了他們就會去傷人.
甚至有一些邪教是教唆人去殺人.
他們的主教會說.
他們自己就是基督.
會用很多方法想去控制人心.
我看到非常生氣.
很生氣.
要用一些時間才可以平息.
因為有些人原來他不認識神.
不敬畏神.
還會用神的名去放肆.
基督第二次回來的時候是會審判的.
這個是真實的.
因此再來警惕我們.
我們要認真自己所相信的東西.
以致我們可以分辨.
我們要知道自己所信的是誰.
所信的是什麼.
以致末後其實會有很多的敵基督出現.
我們都要懂得去分辨.
如果有人問.
你所信的東西和那些邪教有什麼分別.
你會怎樣回應呢.

$^{441}$我們甚至是看同一本聖經.
因此我們更加要認真地讀經和上主一學.
我們要找緊神所賜的恩典.
因為神不單叫基督從死裡復活.
得勝了死亡.
而且他賜下聖靈住在我們心裡面.
聖靈是滿有能力的.
在我們心裡面的聖靈.
是叫基督復活的聖靈.
這個復活的大能能夠超越時間的限制.
去到2023年後的今天.
也都可以繼續.
這個聖靈.
使基督復活的聖靈在我們心裡面.
在我們的生命裡面.
讓我們有力量不再犯罪.
以致我們有活著的力量去愛.
這個復活的大能.
是我們現在能夠活出的生命.
昨晚我去了崑聖弟兄的安息禮.
不知道大家有沒有出席呢.
謝牧師說道.
弟兄的生命不是終結.
而是在永恆的神的國裡面剛剛的開始.
復活的盼望就是相信基督在.
復活的盼望就是相信在基督裡面.
已經睡了的人都可以得著永生.
昨晚有一首詩歌.
叫做《應許之詩》.
很觸動我.
有一句其中一句的歌詞是這樣說的.
它說.
從疑惑到真心相信.
從神處尋到永生.
奮凝裡再生出光芒.
到最終歡呼得勝.
這一份得勝死亡的力量和盼望.
讓忠心愛主的人的安息禮.
是充滿著盼望和愛的.
另一首歌是《因他活著》.

$^{481}$歌詞是這樣的.
當我走完人生路程.
面對死亡,征戰,痛苦.
救主為我戰勝了死亡.
在榮耀中看見救主.
他是活著.
神戰勝死亡的力量.
讓我們有力量去走完人生的路程.
就算人生是艱難的.
因基督的活著.
讓我們有力量去面對明天.
就算走完人生的路程.
我們都是懷著平安的.
不知道大家有沒有參與過.
非基督教的一些喪禮呢.
有一次我去了一位.
未信主的親友的喪禮.
當中有一個叫做破地獄的儀式.
他們會不停地哭.
不停地叫和打鑼打鼓.
哇,我很難受.
我除了為了這位我已經過世.
和我所愛的親人難受之外.
我更加為了在生的親人很難過.
原來沒有復活,盼望的喪禮.
面對著死亡只有恐懼.
和難過.
求主憐憫.
讓我們可以更加努力地傳福音.
讓我們所愛的親人.
都可以得到福音的好處.
最後.
哥林多教會有人不信死人復活.
保羅其實想指出.
哥林多人拒絕死人復活.
其實是否定福音.
甚至是否定神的大能.
保羅強調.
活在當下體會復活的能力.
復活是今生盼望的盼望.

$^{521}$在綜活的救贖裡面.
今生的身體和生活.
與將來的.
將臨的復活不是沒有關係的.
今生的身體和生活.
和將臨的復活不是沒有關係的.
兩者是有關係的.
復活的信念不單止注目在天堂.
而忽略了當下的生活.
而是以一生的生命去見證.
和承認復活的耶穌是主.
讓我們一起祈禱.
親愛的天父.
你是那位又真又活的神.
讓基督成為.
那個先獻上的初熟果子.
讓我們有復活的盼望.
有得救的確具.
讓我們的生命可以不致滅亡.
反得永生.
無論現在還是將來.
都可以.
以永遠和.
永遠與天父在一起.
求你保守我們的信心.
以致我們活在世上.
可以不被社會文化同化.
懂得有智慧的分辨是非.
你是復活的大能.
你的復活大能.
會不停地更新我們.
讓我們可以心意更新.
信仰不勞舊.
能夠充滿活力.
我們都宣告我們的生命是屬於你的.
願你賜給我們有勇氣.
以一生去見證你.
奉主耶穌基督名求.
阿們.
\newpage



\section{希伯來書 12:1-3}
\label{sec:rnbJeDm_hDc}
\textbf{2023-04-23 加恩樓《奔跑吧!》希伯來書 12\_1-3 講員:周弘毅牧師}
\newline
\newline
連結: \href{https://youtube.com/watch?v=rnbJeDm_hDc}{\texttt{https://youtube.com/watch?v=rnbJeDm\_hDc}} ~~~~ 語音日期: 2023-04-25
\newline
\newline
\hyperref[sec:qxsLxZod7kk]{\small{< < < PREV SERMON < < <}}
~
\hyperref[sec:index_chronic]{\small{[返順時目]}}
~
\hyperref[sec:index_scriptual]{\small{[返順卷目]}}
~
\hyperref[sec:l5hkKLKHhA0]{\small{> > > NEXT SERMON > > >}}
\newline
\newline
希伯來書 12:1-3
\newline
\begin{longtable}{cl}
\hline
\hline
章節 & 經文 (和合本修訂版)\\
\hline
12:1 & \begin{tabularx}{0.7\textwidth}{X} 所以，既然我們有這許多見證人如同雲彩圍繞著我們，就該卸下各樣重擔和緊緊纏累的罪，以堅忍的心奔那擺在我們前頭的路程， \end{tabularx} \\ \\ \relax
12:2 & \begin{tabularx}{0.7\textwidth}{X} 仰望我們信心的創始成終者耶穌，他因那擺在前面的喜樂，輕看羞辱，忍受了十字架的苦難，如今已坐在神寶座的右邊。 \end{tabularx} \\ \\ \relax
12:3 & \begin{tabularx}{0.7\textwidth}{X} 你們要仔細想想這位忍受了罪人如此頂撞的耶穌，你們就不致心灰意懶了。 \end{tabularx} \\ \\
[1ex]
\hline
\hline
\end{longtable}
$^{1}$我相信在希伯來書裡面,也有一些經文是大家很喜歡的,.
也最容易想起的可能是第十一章,裡面有很多人物,.
他們所生命經歷的,他們被上帝所讚賞的,都是我們嚮往的..
這個信心英雄榜的結尾在四十節是這樣說的,.
因為神給我們預備了更美的事,叫他們若不與我們同德,就不能圓全..
你見到這裡有第三身的他們,即是第二句那裡,.
又有第一句的我們,作者很期盼透過這些不同的弟兄姊妹,.
幫助我們回顧過去的歷史,然後就將焦點去到十二章開始,.
就不斷我們我們我們..
所以整個希伯來書來到這裡的轉折位,.
都很期盼在我們的生命當中,都能夠繼續很努力的去跑前面的路..
我在大家的教會網頁裡面知道,.
今年你們的年題叫做生命更新,傳揚福音..
相信背後的想法都是期盼我們在生命的歷程當中,繼續邁進向前走..
我相信無論是個人還是教會都是這樣,.
希望我們可以在我們短暫的人生歷程當中,好好地將它活出來,.
將上帝的福音的托付來承傳下去..
今天跟大家分享的經文很短,只有三節,主席的姐妹也幫我們讀出了..
但在開始分享今天這一段的訊息之前,.
我想先跟大家來看一看一個很短的短片,一分鐘的..
這個短片其實就是去年在歐洲田徑錦標賽,地點在慕尼黑..
這一場賽事是4乘100的女子接力賽..
我想大家去留意一下,究竟在這一場賽事裡面發生了什麼事..
請專心去看一看..
當中發生了一個小意外,看看大家是否夠眼力看到..
如果不知道發生什麼事,再來一遍慢鏡..
為什麼抱著頭?.
仍然不知道發生什麼事,看著紫色背心,黃色跑鞋..
還有兩個選手..
看到嗎?還是嘔吐了..
他們掉了一根棒..
我相信如果你玩過田徑,或者你看過田徑,你也知道,.
接力賽掉了一根棒,基本上你哪裡都不用去,對不對?.
你撿一根棒再跑,就差了一條街..
所以接不接到棒,其實是很重要的..
那你說,播放這個片段,跟我們今天的訊息究竟有什麼關係呢?.
我們看回剛才的第一節,看看他們怎麼說..
我們既有者許多的見證人如同雲彩圍住我們,就要…….
其實在這一節的經文的開頭,有一個連接詞,叫做「由此看來」..
畫本就認不出來,不過如果你看新漢語譯本就能看到..

$^{41}$「由此看來」是什麼意思?.
原來今天我們所經歷的,跟前人是有些關係的..
我們是承接著那個由上一手給我們的棒..
我們的先言先信走了一節,接著輪到我們了..
在第十一章這個信心英雄榜裡面,很多人物,對不對?.
這裡說好像雲彩一樣,那麼美..
他們全部都是夢神紀念的..
他們的共通點是因為信心,上帝讚賞他們..
這也是神應許得以實現的一個名證..
所以作者在第一節就已經提醒我們了..
原來我們跟我們身邊的弟兄姊妹,甚至在歷史當中走過這個信仰歷程的,.
在這個信仰路上,我們有一個時空結連的關係..
作者正如他一開始的時候說,他會一直說「我們我們我們」,.
繼續我們下去的,就要告訴我們原來我們這個屬靈群體,.
我們是有承接這個信仰的那一棒..
請我們不要將這些棒掉到地上,好嗎?.
當我們今天仍然在世,帶到我們手上的那一棒是未完的..
我們總是有一些未完成的事,是嗎?.
正在上班的弟兄姊妹,你手上一定有一些工作未完成,.
明天上班的時候要繼續..
正在讀書的,可能你的功課還未完成..
又或者現在來到四月底了,我不知道大家有沒有一個習慣有讀經計劃,.
不知道你會不會落後了幾條街,甚至你說我減肥大計..
總是這個是畢生的事業,總是未完的..
生活上我們有很多未完成的事情,信仰其實也一樣..
究竟我們要做好什麼準備,以致我們能夠在往後的日子,.
調教好我們的心態,以致我們能夠完成上帝給我們的託付呢?.
我們再看多一次第一節,.
經文這樣說,我們既有者許多的見證人,如同雲彩圍著我們,.
就當放下各樣的重擔,脫去容易纏累我們的罪,.
全心忍耐,奔那擺在我們前頭的路程..
作者用了「奔」這個字,即是說有速度的,.
英文是run,不是jog,即是說不是萬無目的,.
隨隨便便四處逛逛,累了就休息一下,.
躺平一下,不是的..
而且在最後那句,就說奔那擺在我們前頭的路程,.
除了有速度,還有方向,前面,對嗎?.
而路程在呂鎮忠譯本和詩高譯本,.
將它譯作賽程,即是說這是比賽,.
可能是馬拉松之旅,鼓勵我們要完成上帝給我們的長途賽,.

$^{81}$究竟我們能否完成神給我們的託付和目標呢?.
在信仰的路上跑了這麼多年,我不知道在座的弟兄姊妹,.
你信了耶穌有多久?.
跑著跑著,不知道我們會不會掉進一個不太想跑的狀態,.
可能我們沒有傳福音很久,也沒有和身邊的未信的朋友分享.
上帝在我們生命裡的奇妙作為,.
我們沒有打算去服侍留給年輕人,.
甚至因為疫情已有一段時間,.
所以我們隱藏了很久,實體聚會也不太願意回來,.
究竟我們是否正在跑呢?.
又或者我們正在跑,我們是否全心全意地跑呢?.
到覺得有點累了,就坐在牆邊,坐下來,收工了,.
是否這樣呢?.
現在踏入了疫後,福常,很多事情,.
都是時候回來參與不同的教會生活團契當中,.
知道大家在接下來的星期又有聚餐,.
剛才在感受堂那邊聽到報告說,.
弟兄姊妹踴躍參與,是好事來的,.
踴躍我們真的能夠重拾我們和身邊弟兄姊妹的關係,.
有時我這樣想,如果我們的生命不更新的話,.
就有點像我們平時吃的東西,現在都挺好的,.
包裝上會有一個最佳食用日期,.
Best Before Date,不論何時食用,.
如果我們真的過了這個食用日期,.
我相信自然我們不會享受到最美味的地方,.
可能會有一點失色,甚至你說再等久一點,.
令到它變壞了,就吃到肚子痛,.
回想我們的信仰歷程,我不知道大家覺得哪一個階段,.
你覺得最精彩,是否剛剛信耶穌的時候,.
得著福音,很熱切追求成長,.
其實今天我們仍然有可能有力量繼續跑下去,.
為主跑得更好,不過總是心不知去哪裡,.
期盼今天透過希臘書短短這三節聖經,.
幫助我們重新激勵我們的心智,.
讓我們的生命可以被上帝的話更新,.
要更新生命其實有三方面的準備,.
聖經告訴我們,第一方面就是要放下各樣的重擔,.
重擔是甚麼呢?可能是你放不下的東西,.
你過往的歷史,受過的傷害,面對的痛苦,.
心裡面的重擔,你現在掛心的東西,.

$^{121}$可能是你的工作,可能是你和身邊的人的關係,.
甚至是金錢,經濟上遇上壓力等等,.
總之那些拖累你,令你跑不快的東西,.
我相信無論你是喜歡做運動,跑步,.
甚至你不喜歡做運動,平時走路的時候,.
相信大家都會遇過,可能有粒沙走進你的鞋裡,.
當沙走進你的鞋裡,你一定會覺得不舒服,.
甚至你會覺得很痛,那你會怎樣辦呢?.
停下來,脫下你的鞋,先倒掉那粒沙,.
這樣就可以舒舒服服地繼續走下去,或者跑下去,.
以致我不會走到半路,我放棄了,.
在我的信仰歷程當中,我連上帝都放下了..
第二,這裡提醒我們要脫去容易纏累我們的罪,.
香港現在有很多跑步的運動,.
有時你會見到不同的人除了很正規,正正經經地跑,.
有些人都會視這些跑步的活動,.
好像嘉年華那樣,穿著闊袍大袖,扮鬼扮馬,.
如果你真的想正正經經地跑,.
我估計你一生的裝束一定不是這樣的,.
你會輕裝上陣,將那些阻礙你的東西脫下,.
這裡告訴我們有什麼會阻礙我們呢?.
是容易纏累我們的罪,.
中文譯本是說與我們糾纏不休,.
這裡的意思是告訴我們,是他糾纏我們,.
即是說他不會輕易放過我們,.
好像我們平時早上外出換衣服,.
我相信你不會在換衣服的時候,.
只脫下一件衣服,然後第二件就不脫,.
然後穿上你要外出的衣服,做不到的,.
你一定要好好地脫下你的衣服,.
兩件衣服都脫下了,.
然後能夠穿起你需要穿上,上班外出的衣服,.
所以保羅在《異邦所書》第四章第七節提醒我們,.
不要讓魔鬼留地步,.
原來在我們的信仰歷程裡,.
罪會與我們糾纏不休,.
究竟有什麼東西是在攔阻我們,.
而你還沒有放下呢?.
會不會那些東西真的收得很緊,.
在這個階段連自己也被騙了?.

$^{161}$第三是講心態,.
作者提醒我們要全心忍耐,.
有一個人在路加福音9章57節,.
跟耶穌這樣說,.
他說你無論往哪裡去,我要跟從你,.
你去哪裡我都會跟從你,.
但58節耶穌怎樣回應他?.
耶穌說,狐狸有洞,天空的飛鳥有窩,.
只是人子沒有枕頭的地方,.
意思是如果要跟隨耶穌是有代價的,.
作者在這裡好像在用馬拉松作為比喻,.
對於馬拉松跑手來說,.
其實要開始相信比繼續下去容易,.
不過有一個危機,.
就是開始的時候你控制不了自己的步速,.
跑得太快,.
當你越跑越快,.
去到後半節的時候就沒有氣,也沒有力量,.
所以完成不了,.
另外一個需要跨越的危機就是極點,.
我沒有跑馬拉松,.
不過有人跟我說,.
去到差不多30公里,.
你會覺得很累,累到不想跑,.
你必須要跨越到這個極點,.
你才能夠繼續下去,.
這就是忍耐和堅持..
報導家施暴真說過一個很有趣,.
但也很真實的比喻,.
他說在洪水的時候,.
那隻蝸牛一定要全心忍耐才能上方舟,.
By perseverance, the snail reached the ark..
如果跑一跑就算了,.
結果就淹死了..
要跑好我們的人生馬拉松,.
除了要留意身邊像雲彩般的見證人,.
看看別人,.
也要留意自己怎麼跑,.
不過最重要的不是這兩樣,.
而是看耶穌怎麼跑,.

$^{201}$因為經文在第二節說,.
仰望為我們信心創始成中的耶穌,.
他因那擺在前面的喜樂,.
就輕看羞辱,.
忍受了十字架的苦難,.
便坐在神寶座的右邊..
如果你信主的日子夠長,.
我都很肯定你一定會想起.
那些你昔日的戰友,.
今天不在你身邊,.
不在教會裡,.
甚至離開了信仰,.
一定有這樣的面孔在你腦海裡浮現..
為什麼呢?.
原因是他們沒有仰望耶穌,.
仰望這個動詞,.
如果我們看其他聖經中文的譯本,.
包括這裡,.
現代中文譯本,.
呂鎮忠譯本,.
詩高譯本,.
他們通通都將這個詩譯為「注時」,.
即不是像畫本那樣,.
「我高個頭看看吧」,.
不是,是注時,.
專注不單只是作為運動員必需要的要求,.
在我們信仰信心的路途上都是一樣的..
我不知道在座的弟兄姊妹,.
你喜不喜歡打籃球,.
或者你看籃球,.
美國職業籃球聯賽,NBA,.
你看我這樣的身形,.
我都不是打籃球的人,.
我不懂打,.
不過我喜歡看..
如果你打籃球,.
或者你喜歡看籃球,.
又讓你猜猜,.
究竟在美國職業籃球聯賽,NBA裡面,.
射罰球射得最準的籃球運動員,.

$^{241}$你猜猜是誰?.
這些名字可能你在新聞都聽過,.
是否史提芬居里,.
Stephen Curry,.
這個人還在打,.
勒邦詹士,.
勒邦詹姆斯,.
米高佐敦,.
Michael Jordan,.
甚至是離開了世界的Gobie Bryant,.
這個高比拜仁,.
是否他們呢?.
答案是,.
以上皆非,.
全部都不是,.
他們全部都不是射罰球最準的,.
這些打球很厲害的,.
射籃好像這樣扔下去就中了,.
都不是,.
那誰是?.
我保證你猜不到,.
因為我預備這篇訊息之前,.
連我也不知道是誰,.
你又說,.
在資料上找到,.
是這個人,.
叫湯姆安柏瑞,.
Tom Embry,.
這個好像是伯伯,.
因為在健力士世界紀錄大全,.
最近被破了,.
不過這個人之前保持了世界紀錄,.
他的紀錄是怎樣呢?.
2017年他去世了,.
他除了喜歡打籃球之外,.
他是一個足科的醫生,.
即是扁平足那種,.
幫這些病人的,.
他在1993年,.
他破了世界紀錄的時候,.

$^{281}$他71歲,.
你說這個伯伯都可以了,.
71歲,.
他的紀錄是怎樣呢?.
用了12個小時連續射入2750球,.
這樣都可以嗎?.
你說他怎樣可以射得這麼準呢?.
他的指定動作就是這樣,.
站在罰球線上,.
拿著球拍拍拍,.
三下,然後就起手,.
然後再射一球,再射兩球,再射三球,.
2750球,.
當然有人問他,.
你都射得這麼準嗎?.
這樣都可以嗎?.
有什麼秘訣呢?.
他的答案就是專注,.
當然你要射球射得好,.
技術是必須的,對不對?.
你要好好操練,改善射球的技巧,.
你有技巧之後,.
接下來就是專注,.
專注都是專注,.
我深信在過程中,.
什麼叫專注呢?.
就是將其他東西全部放下,.
這就叫專注,.
這一節的經文教導我們,.
我們需要專注,.
我們需要注目在耶穌的身上,.
你當然會問,.
具體是怎樣呢?.
作者有說,.
他告訴我們,.
這位主耶穌是為我們信心創始成中的那位,.
這裡你會看到,.
原本有「我們」這個字,.
但原本聖經是沒有的,.
翻譯的時候才加上去,.

$^{321}$所以如果我們看回呂鎮忠譯本和施高譯本,.
更加準確的譯法就是注視著信仰上的開創者和完成者,.
耶穌,弄清楚,.
這裡不是說我們的信心有多大,.
有多小,不夠,.
不是說我們的信心,.
是說耶穌的信心,.
創始,當然大家都可以猜到,.
開始的人,冠軍領袖,先驅者,發起人,諸如此類,.
成中就更加容易解釋,.
就是完成那件事,.
這裡說的是,.
耶穌是甘心降世為人,.
他走上十字架,.
就滿足了這個救贖的要求,.
他坐在神的右邊,.
所以這裡是說耶穌的信心,.
你說耶穌當然有信心,.
為什麼他那麼有信心?.
經文有說,第二節,下半節,.
他因擺在前面的喜樂,.
就興旱羞辱,忍受了十字架的苦難,.
便坐在神寶座的右邊..
十字架本來是用來懲罰奴隸和罪犯,.
所以每一個走上十字架的人,.
都是受盡了藐視和羞辱..
這裡說他興旱羞辱,.
忍受了十字架的苦難,.
這麼容易?.
很難的!.
但是為什麼耶穌做到?.
他還能夠將這樣的經驗反轉,.
興旱羞辱,.
最終得以坐在神寶座的右邊,.
前面有一個很重要的秘訣,.
因為他看到擺在前面的喜樂,.
因為他知道他最終能夠坐在神寶座的右邊..
這是他的目標,注目在耶穌身上,.
就是向前望,期盼見到祂,.
和祂一起分享這份永恆的盼望和喜樂..

$^{361}$所以在現實生活裡,.
很多東西令我們沮喪,.
能夠有機會認識耶穌,.
我很期盼這個永恆的盼望和喜樂,.
我們真的能夠抓得住,.
抓得穩..
在《讓福音》差不多完結的時候,.
21章,主耶穌重複問彼得三次:.
「你愛我嗎?」.
這段經文大家熟悉得不得了..
每一次主耶穌去吩咐他餵養,.
或者牧養我的羊之前,.
都會問,問完又問,問完又問:.
「你愛我嗎?」.
耶穌要讓彼得很清楚知道,.
他要去牧養,他要去服侍之前,.
一定要搞清楚,確認他愛耶穌..
我們有能力這樣做,.
甚至我們在信仰上有能力繼續跑下去,.
這也是很重要的原因,.
是源於我們知道耶穌愛我們,.
我們也認定我們愛耶穌,.
不過現實很多時候我們不得不承認,.
我們搞錯了焦點..
當我們參與服侍,甚至教會服常,.
我們去打個電話,.
或者關心一下我們身邊的弟兄姊妹,.
我們覺得我們為神做了一些事情,.
對不起,神不需要我們為祂做一些事情..
其實在這個過程裡,.
我們要問的就是,.
究竟我們是否覺得耶穌愛我們?.
如果我們都不是很覺得耶穌愛我們,.
我深信我們更加談不上我們愛耶穌,.
做到累,.
面對生活裡面很多的難處和挑戰,.
我們就會覺得耶穌離場了,.
不知道去了哪裡..
所以當我們今天面對這個處境,.
我們覺得困難,.

$^{401}$我們覺得沮喪,疲倦,甚至耗盡,.
如果是在這個處境裡,.
我相信我們需要先做一件事,.
騰出空間和上帝建立這份關係,.
這份愛的關係,.
以致我們能夠享受裡面的喜樂..
所有事情都是關係先行,.
但是如果我們忘記了這件事,.
相信我們就很容易掉進沮喪的陷阱裡,.
當我們能夠重建和耶穌的關係,.
我們就不會失去這個心智,.
亦都能夠給予我們力量,.
繼續向前跑,.
不要放棄這段人生的才晴..
除了注目耶穌之外,.
其實接下來的第三節,.
還有少許補充提醒,.
幫助我們走完人生馬拉松,.
經文這樣說,.
那忍受罪人這樣頂撞的,.
你們要思想免得疲倦灰心..
走路走得久,走得遠,.
累是正常的,.
有起有跌,.
屬靈的道路亦都是一樣..
所以作者提醒我們要思想,.
那位忍受罪人這樣頂撞的,.
這當然是在說耶穌,.
當我們面對困難的時候,.
我們會容易跌入處境,.
轉入黑洞裡面..
很慘,覺得身邊的人不明白自己,.
為何這麼艱難,.
難過身邊的人這麼多..
這段經文提醒我們要思想,.
要當然是一個命令,.
思想就是想清楚,.
考量一下,.
我們就做一個判斷或決定..
但這裡講清楚告訴我們,.

$^{441}$如果我們要仔細思量比較,.
我不是在看其他人,.
這裡要思想的是在說,.
那忍受罪人這樣頂撞的,.
即是要思想耶穌,.
祂面對四面八方而來的凌辱,.
挑戰,苦待,被唾棄,.
鞭打,傷害等等..
我們就知道,.
祂過渡,因為我愛祂,.
祂也愛我們,.
所以祂願意甘心樂意走上十字架,.
忍受了十字架的苦難,.
坐在神的右邊,.
今天祂已經從死裡面復活,.
靠著祂,我們也能過渡,.
我們不需要疲倦灰心..
今天和大家分享這段經文很短,.
只有三節,.
作者提醒我們幾個重點,.
第一,他告訴我們,.
我們真的要好好承接信仰這一棒,.
將信仰承傳下去,.
如果在座的你是父母的話,.
我也是兩個兒子的爸爸,.
彼此勉勵,.
我們真的不要讓信仰的棒掉到地上,.
第二,我們真的要跑好我們這一場人生的馬拉松,.
有幾方面我們需要準備好,.
第一,我們要放下那些.
拖著我們後腿,.
使我們跑不快的重擔,.
和罪惑清界線,.
並且要有忍耐的心,.
第三,跑的時候,.
是專注在和耶穌的關係上,.
期盼見到祂,.
可以和祂一起分享這份喜樂,.
還有第四,.
真的不要看其他人,.

$^{481}$人比人比死人,.
這裡告訴我們,.
如果我們真的要看,.
看看耶穌,.
祂所承受的一切,.
祂已經得逞,.
以致我們能夠靠著祂,.
有力量奔跑前面人生的路程,.
最後我想用一個,.
同樣是關於馬拉松,.
很美麗的故事,.
作為總結,.
在1986年11月6日,.
有一個跑手,.
名叫爆拍威蘭,Bob Whelan,.
他以19,413名,.
都很厚,.
去跑完紐約的馬拉松賽事,.
你問跑19,413,.
即是跑第幾,.
即是跑到最後,.
跑到最後,.
不過當他跑完,.
去到經過終點線的時候,.
仍然受到旁邊的人,.
夾道歡呼,拍手掌,.
歡迎他,.
為什麼呢?.
因為他是第一個,.
用他的手去代替他的腳,.
去跑完,.
跑完他人生很重要的關口,.
這個馬拉松賽事,.
到他完成了,.
當然報紙媒體會來訪問他,.
嘩,這麼辛苦,.
你為什麼仍然要跑紐約馬拉松?.
他就回答他,.
第一,其實我也想試一下自己的體能,.
究竟能不能撐下去,.

$^{521}$第二,其實他就是要為一個健康組織,.
他的服務去做宣傳,.
第三,最重要,.
因為他要見證上帝,.
他這樣說,.
其實他的體會就是,.
成功不是在於你的起點,.
今天我們經常聽到,.
贏在起跑線,.
不是,不是,我覺得反過來,.
是在於我們的終點,.
他說今天我完成了,.
他說開始的階段,.
他覺得最難,.
剛才不是說了,.
在很多馬拉松選手來說,.
開始的時候,.
充滿力量,應該比較容易跑,.
他說很難,.
為什麼開始的時候這麼難呢?.
有原因的,.
因為賽會知道他,.
是唯一一個用他的手代替腳來跑,.
所以就讓他早兩個小時開跑,.
人家開始之前,.
多給他兩個小時,.
當他一路跑,一路跑,.
你也想像得到,.
其他跑步的人,.
健全的運動員,.
他就看著一個一個超越自己,.
越跑越後,.
他就是這樣,.
不過他說跑著跑著,.
他覺得自己好像回到家裡一樣,.
他平均的步速,.
一個小時就是1.6公里,.
我不是跑步的人,.
如果你有跑開步,.
1.6公里,.

$^{561}$你大概知道有多快,.
有多快,.
他這樣的優勢,.
他怎麼跑呢?.
是這樣跑的,.
他坐在馬鞍上,.
然後緊握著拳頭,.
拳頭用軟墊包著,.
然後用雙臂的力量撐著自己跑,.
所以他開玩笑地說,.
他說我穿的是一號跑鞋,.
很小的,.
用拳頭撐著去終點,.
就是靠他雙臂的力量,.
他沒有放棄,.
結果他跑多久才能跑完呢?.
他用了四天兩個小時,.
48分17秒,.
來完成這場賽事,.
四天多,.
這樣撐著去終點,.
求主幫助我們,.
我們身邊有很多美麗的見證人,.
在聖經裡面給我們很多的鼓勵,.
容讓這些激勵,.
真的幫助我們繼續跑下去,.
很期盼我們每一個弟兄姊妹,.
都能夠像有一本書叫做Finishing Well,.
我們真的好好地去完成我們人生的歷程,.
讓上帝得著當得的榮耀,.
我們低頭有個祈禱,.
天父上帝我們知道在我們生命的歷程當中,.
我們少不免遇到難處和挑戰,.
在這個生命的歷程當中,.
我們深深知道不是天色常藍,.
花香常蔓,.
但我們也知道,.
你在我們身邊,.
與我們同行,.
與我們同跑,.

$^{601}$以致我們不需要放軟我們的手腳,.
我們仍然可以充滿力量地,.
繼續奔跑我們人生的歷程,.
幫助我們,.
有很多時候我們都需要被提醒,.
被身邊的弟兄姊妹和你的說話來激勵,.
以致我們每一個都能夠成為微小的器皿,.
去完成你在我們生命當中,.
各個人不同的託付,.
多謝你聽我們的禱告,.
是否愛我們主耶穌基督的命求,.
阿們..
\newpage



\section{路加福音 16:1-13}
\label{sec:l5hkKLKHhA0}
\textbf{2023-04-23 午堂《不義與忠心》路加福音 16\_1-13 講員:楊慶球牧師}
\newline
\newline
連結: \href{https://youtube.com/watch?v=l5hkKLKHhA0}{\texttt{https://youtube.com/watch?v=l5hkKLKHhA0}} ~~~~ 語音日期: 2023-04-25
\newline
\newline
\hyperref[sec:rnbJeDm_hDc]{\small{< < < PREV SERMON < < <}}
~
\hyperref[sec:index_chronic]{\small{[返順時目]}}
~
\hyperref[sec:index_scriptual]{\small{[返順卷目]}}
~
\hyperref[sec:J_KPU0uZysw]{\small{> > > NEXT SERMON > > >}}
\newline
\newline
路加福音 16:1-13
\newline
\begin{longtable}{cl}
\hline
\hline
章節 & 經文 (和合本修訂版)\\
\hline
16:1 & \begin{tabularx}{0.7\textwidth}{X} 耶穌又對門徒說：「某財主有一個管家，有人向主人告管家浪費他的財物。 \end{tabularx} \\ \\ \relax
16:2 & \begin{tabularx}{0.7\textwidth}{X} 主人叫他來，對他說：『我聽到了，你做的是甚麼事？把你所經管的交代清楚，你不能再作我的管家了。』 \end{tabularx} \\ \\ \relax
16:3 & \begin{tabularx}{0.7\textwidth}{X} 那管家心裡說：『主人辭我，不用我再作管家，我將來做甚麼呢？鋤地嘛，沒有力氣；討飯嘛，怕羞。 \end{tabularx} \\ \\ \relax
16:4 & \begin{tabularx}{0.7\textwidth}{X} 我知道怎麼做，好叫人們在我不作管家之後，接我到他們家裡去。』 \end{tabularx} \\ \\ \relax
16:5 & \begin{tabularx}{0.7\textwidth}{X} 於是他把欠他主人債的，一個一個地叫了來，問頭一個說：『你欠我主人多少？』 \end{tabularx} \\ \\ \relax
16:6 & \begin{tabularx}{0.7\textwidth}{X} 他說：『一百簍油。』管家對他說：『拿你的賬，快坐下，寫五十。』 \end{tabularx} \\ \\ \relax
16:7 & \begin{tabularx}{0.7\textwidth}{X} 他問另一個說：『你欠多少？』他說：『一百石麥子。』管家對他說：『拿你的賬，寫八十。』 \end{tabularx} \\ \\ \relax
16:8 & \begin{tabularx}{0.7\textwidth}{X} 主人就誇獎這不義的管家做事精明，因為今世之子應付自己的世代比光明之子更加精明。 \end{tabularx} \\ \\ \relax
16:9 & \begin{tabularx}{0.7\textwidth}{X} 我又告訴你們，要藉著那不義的錢財結交朋友，到了錢財無用的時候，他們可以接你們到永遠的住處去。 \end{tabularx} \\ \\ \relax
16:10 & \begin{tabularx}{0.7\textwidth}{X} 人在最小的事上忠心，在大事上也忠心；在最小的事上不義，在大事上也不義。 \end{tabularx} \\ \\ \relax
16:11 & \begin{tabularx}{0.7\textwidth}{X} 若是你們在不義的錢財上不忠心，誰還把那真實的錢財託付你們呢？ \end{tabularx} \\ \\ \relax
16:12 & \begin{tabularx}{0.7\textwidth}{X} 如果你們在別人的東西上不忠心，誰還把你們自己的東西給你們呢？ \end{tabularx} \\ \\ \relax
16:13 & \begin{tabularx}{0.7\textwidth}{X} 一個僕人不能服侍兩個主；他不是恨這個愛那個，就是重這個輕那個。你們不能又服侍神，又服侍瑪門。」 \end{tabularx} \\ \\
[1ex]
\hline
\hline
\end{longtable}
$^{1}$我第一次來見到很漂亮的呂貴堂.
為你們很感恩.
我已經不是加拿大環球學院的教授.
我退休了.
我是加拿大宣導會的退休牧師.
回來香港退休.
現在在我舞會 藝人會大會堂.
做顧問牧師.
不過我一直都是他顧問牧師.
退休後很開心見到大家.
剛才讀的那段經文.
我不知道你們怎麼看.
字比較小 都是剛才那段經文.
我不知道你們讀的時候會不會覺得有點奇怪.
這個人這麼差勁.
在外面做得不好.
然後找主人的客人來.
私自來一個減50一個減20.
等他將來退休後有好地方去.
其實是不合理的.
如果在今天.
你一定會找廉署抓他.
不坐十年八年 你的氣都不順.
很奇怪.
但耶穌當時說.
人們有沒有質疑耶穌.
有這樣的事.
更奇怪的是主人讚你聰明.
你做老闆.
底下的人幫你倒米.
你會說聰明 你倒得好.
你讀了是不是不舒服呢.
但耶穌說完.
那班門徒又收到.
很明顯.
耶穌當時說.
那班人是明白的.
而且他們接受.
並且和我們今天的了解有差距.
究竟發生什麼事.

$^{41}$一個這麼不合理的管家.
在外面敗壞了主人的名聲.
或者財務管理不好.
或者得罪了所有人.
諸如此類.
總之他有問題.
然後主人才去解僱他.
我們看看這個管家怎麼做.
這本經文.
這個管家.
他幫主人做大生意.
不是雞毛蒜皮的事.
100劉油.
大概4000公斤.
減50.
大概2000公升.
在當時.
即是今天.
大概6萬加幣.
大概40萬港幣.
加幣跌了一點.
當你38萬.
100石墨紙.
4萬公升.
然後去減那個數目.
也是差不多40萬.
而這個這麼大筆數減了之後.
是當時一個工人.
三年的薪酬.
即是說.
這樣的做法.
是會為他得來很大的益處.
但我們想不明白.
為什麼主人會稱讚他.
明明主人損失.
我們就明白.
當時的管家和今天的管家.
是有不同的地方.
因為當時的管家.
他和主人的關係.

$^{81}$不單止是僱傭關係.
他是一個代理關係.
他代表主人出去接生意.
而他接生意是有風險的.
而他所賺到的這個管家.
他是有很高的利潤.
這個利潤是歸管家的.
另外利息.
是歸那個主人.
現在他將這個人.
他所賺到的.
他減去.
學者相信.
他減這麼大筆錢.
就應該將他的佣金和利息.
一次退給那些人.
所以才會這麼大筆.
退給那些借債做生意的人.
換句話說.
那個主人損失了利息.
管家有沒有損失呢.
有 他損失了佣金.
但是主人損失利益.
損失了利息.
為什麼主人都會稱讚他呢.
原來主人雖然損失利息.
但是他是有賺的.
為什麼呢.
按照當時的猶太人.
他們很遵從聖經的教導.
去做生意是不收利息.
你還記不記得.
在他們做生意的時候.
那個利息是不可以收的.
但是你不收利息.
你怎麼做生意呢.
所以潛規則就是.
在當時所有的人都收利息.
只有那些很高尚的人.
很遵從上帝的說話的人.

$^{121}$才不收利息.
因此這個管家.
他將利息歸了給借債的人.
這個主人不是完全損失的.
他是獲取一個很高的社會地位.
一個很高的社會聲譽.
所以這個管家稱讚他.
其實也不是完全沒有道理.
因為這個主人稱讚他.
完全不是沒有道理.
因為這個管家是有損失的.
他犧牲了自己的利潤.
而這個主人雖然損失了利息.
但是他得到社會給他的稱讚.
得到很高的社會地位.
恢復回去.
因為之前這個管家.
可能令到他在社會上失去了聲譽.
現在恢復了.
相對來說.
這個主人他也有賺.
所以這個主人就這樣稱讚他.
就說他用了一個例子.
這個例子是很奇怪的.
就是今世之子做事比光明之子厲害.
這個說話真是可圈可點的.
耶穌是用這句說話.
來對教導門徒.
就是說這個管家.
他上方設法為了自己退休.
將來有地方去.
他用了這麼極端.
這麼前衛的方法.
就要保障他將來可以有人接濟他.
有人幫助他.
所以他不惜犧牲主人的利息.
犧牲自己的佣金.
以致他退休無憂.
所以耶穌就借這個主人說.
今世之子做事比光明之子更加聰明.

$^{161}$什麼叫今世之子.
什麼叫光明之子.
如果按照聖經一貫的做法.
今世之子.
就是世上的人.
很簡單的說.
不是基督徒 非基督徒.
光明之子.
就是耶穌所稱讚的追隨的人.
在這裡說就是基督徒.
基督徒 非基督徒.
耶穌說那些不信耶穌的人.
做事比你們信耶穌的人厲害.
你明白意思嗎.
他說那些不信耶穌的人.
他可以為了自己退休.
為了自己的生活.
所以上方設法來預備未來的日子.
我們信耶穌的人.
我們屬靈的人.
是不是應該要更加好好的.
去考慮我們將來屬靈的日子.
我們未來的在天上與父神一起.
我們的處世.
是不是要比外面的人更加積極呢.
所以這句話是有價值判斷的.
事實上在這個世界上.
在外面工作的人是很辛苦的.
外面工作的人.
朝九晚五不單止.
還很多憂慮 用很多方法.
事奉上帝的人.
我們也有要像他們那樣積極.
在工作上.
因此事奉主人.
我們在教會事奉也好.
我們事奉神.
我們更加不要馬虎.
要像外面的人.
外面的人為世上的生活.

$^{201}$為暫時的生活都這麼拼命.
我們是不是要比他們更加拼命呢.
所以其實這句話背後是有這樣的勸導.
不過這段經文.
其實我鼓勵大家.
看著這段經文.
有聖經打開.
我今天是和大家講聖經.
這段經文一路講下去.
我們就開始慢慢發現.
耶穌用的字是不容易明白的.
這個比喻.
我們要像外面的人.
外面的人他們做事是為了什麼呢.
就出現這個字.
這個字就是一個不義的錢財.
記不記得剛才那段經文.
大事上的忠心.
小事的忠心.
如果在大事上不義.
在小事上不義.
這個不義很明顯有不好的意思.
但耶穌用這個不義的時候.
他的意思是轉化了.
是什麼意思呢.
我們事奉上帝.
屬靈的事情.
我們努力為神去事奉.
是會產生永恆的效果.
帶人認識主.
得到永生.
在世上所做的事情.
雖然賺很多錢.
但那些錢是不能讓你得到永生.
所以就用不義.
但這個不義不是不好的意思.
為什麼呢.
因為如果這個不義是不好的話.
耶穌就不會叫你在不義的錢財努力.
第一段經文是這麼說的.

$^{241}$我們再讀一次第九節.
我又告訴你們.
要藉著那不義的錢財結交朋友.
到了錢財無用的時候.
你會接你們到永生.
這字太小.
所以我眼睛看不清楚.
耶穌怎麼說.
他要藉著不義的錢財結交朋友.
耶穌不會說那些不義錢財不要.
那些販毒回來的殺人放火的.
扔掉它.
不是啊.
不是這個意思.
我們要透過那些不義的錢財.
那些不義就是沒有產生永恆的果效.
是這個意思.
要將不產生永恆果效的錢.
轉化成做屬靈的事情.
以至產生永恆的果效.
就是這個意思.
所以耶穌就叫我們要在不義的錢財上中心.
有沒有看到那一節經文.
我們看前面.
要在不義的錢財中心.
就在路加科那裡.
我不知道你們有沒有看到.
要在不義的錢財中心.
然後要藉著不義的錢財.
去結交朋友.
這篇經文很有趣.
很多人就將世界看成不義.
就將世界的東西.
是不好的.
扔掉它.
只是屬靈的好.
我在Calgary的時候.
那時候我在Calgary的宣導會神學院.
在那裡事奉的時候.
在Calgary有一個牧師.

$^{281}$就邀請他去家裡吃飯.
你知道加拿大的人.
那些房子很大.
地上又有放電影.
有樂隊房.
上面有很多大房.
房子很便宜.
幾百萬港幣.
看的時候.
突然經過一間房.
那位牧師.
師母就說.
每間房都要看.
他說這間房你不能看.
其實我沒有什麼心機.
其實我不是很有興趣.
要看別人的房子.
不過他說了不要看.
我就很想看.
為什麼房子不能看.
吃完飯之後經過.
我就撞頭.
看一看.
為什麼不能看.
一看之後就嚇了一跳.
房子很骯髒.
我猜一定是女生房.
衣服穿到沒有掛.
東西放到處都是.
吃完杯麵又在床邊.
又在桌面.
師母見到我.
牧師說不好意思.
這間房是我女兒住的.
我不敢批評人.
他自己說.
他說很骯髒.
我說不是.
因為在教會裡.
他的女兒很decent.

$^{321}$穿得很漂亮.
大唱詩.
不是你.
大唱詩.
很好.
我說為什麼.
他房子這麼骯髒.
後來吃完飯.
我們喝茶.
師母就說.
牧師說女兒很奇怪.
她很愛煮.
教會事情.
她做到百分百.
但是他說.
俗世的事.
她不放心.
她覺得收拾房間不重要.
她說媽媽你不值得.
你收拾吧.
所以她不收拾房間.
她氣死我.
所以收拾到那裡.
我都不幫她收拾了.
由得她.
狗鬥就狗鬥.
她說可不可以教她.
收拾房間也很重要.
她的觀念就是這樣.
事奉上帝是最好的.
收拾房間.
這些熟悉的東西.
就不重要.
分得很清楚.
是不是這樣看.
上帝是不是這樣看.
是不是兩個世界.
一個是屬於上帝的.
一個是屬於這個世界的.
後來我都沒什麼機會和她的女兒聊天.

$^{361}$我相信她媽媽和爸爸.
足夠可以教導她.
不過真的有些人.
就將世上的事情成為不義.
有時候.
保羅都教導我.
看世界如糞土.
於是糞土就全部不要.
保羅的意思不是不要.
我們不要將它看成.
比永生更重要.
所以不義就不是.
你賺的錢是臭的.
不好的.
不是這個意思.
所以在不義的事情.
要轉化.
產生永恆的價值.
接著這段經文.
又是很難解釋.
他說.
這段經文怎麼說.
他直接不義的錢財.
結交朋友.
就是說在這個世界上.
我們賺到錢.
我們和外面的人來往.
不要將錢看成臭的.
不要東東的不是販毒.
正當當賺回來的.
這些錢.
到了錢財無用的時候.
就是將來.
我們見天父.
世上的錢財無用.
他們可以接你們到永遠的住處去.
你們是誰呢.
前面就說.
結交朋友.
很明顯.

$^{401}$他們就是那些朋友.
你結交了那些朋友.
那些朋友可以接你去.
永生的住處.
那些人接你去.
不是你接他去嗎.
耶穌有沒有說錯了.
你留意經文是這樣.
沒有改經文.
當耶穌說這句話.
祂說你要好好在外面賺錢的時候.
你要透過在外面的金錢.
透過在外面的勞力.
來幫助那些人.
或者是結識那些人.
不要和他們斷絕關係.
然後透過你的勞力.
和他們一起工作也好.
賺到錢也好.
來使他們能夠認識主.
所以到錢財無用的時候.
他們會接你們去永生的住處.
那個你們.
他們接你們.
他們就是那些人.
是什麼意思呢.
原來耶穌是用當時很流行的諺語.
而這個諺語在馬太福音也說過.
就說.
第九節.
不易戰就結交朋友.
這個諺語.
你們就指門徒.
就是指耶穌所說的.
對著門徒說.
他們就指外面的人.
如果你做在那些人身上.
就做在主的身上.
所以他們就代表上帝.
接你們永生的住處.

$^{441}$原來是很大價值.
當我們大人信主.
當我們幫助人.
就做在上帝的身上.
是上帝很喜歡的.
所以經文就說.
在這個世上.
我們很忠心做好.
它會產生一個永恆的效果.
永遠的住處.
產生永恆的效果.
我在溫州教書的時候.
過去我每年都會在中國大陸教書.
在溫州教書的時候.
有一次有一個弟兄.
做見證.
他在那裡見證.
他說他在一個外資公司工作.
那間外資公司每年都賺錢.
一直發展得很好.
他只是中層.
他不是很高級.
很中層.
公司很大.
人們做得很開心.
不過員工就有一個.
不是很好的習慣.
很喜歡拿公司的東西回家.
因為公司很多東西.
兒子跟媽媽說.
我沒有筆.
媽媽說不用了.
公司有很多.
我拿幾支筆給你.
沒有什麼紙用.
公司有很多.
不是髮夾.
文具那些.
很習慣的.
公司又沒什麼所謂.

$^{481}$公司賺錢.
那些人很努力地工作.
所以是習以為常.
這個弟兄就說.
他有一天信了耶穌.
他信了耶穌之後.
他覺得這樣拿公司的東西.
是不對的.
那些東西不是很值錢.
很好笑.
不是很值錢.
他就決定.
他就跟耶穌說.
他說我不再拿這些東西.
雖然不是值很多錢.
他說我不再拿.
他不再拿.
對他有沒有幫助.
沒有.
人們就繼續拿.
他就繼續工作.
他不會舉報人.
他沒有那麼壞.
那些人也習以為常.
過了差不多大半年.
有個評估.
就是上面經理.
見下那些中層的領導人.
說說工作.
都說他工作很好.
表現都不錯.
然後最後上頭經理就跟他說.
他說我覺得你好像跟其他人有點不同.
他說有什麼不同.
他說我發覺.
我留意了很久.
你從來不拿公司的東西走.
他當時就害怕.
他說你知道的.
他說知道.

$^{521}$不過那些東西很閒.
我們也不會計較.
他說不過.
他說你從來都不拿.
他說一張紙.
一眼針你都不拿.
他說你是有點不同.
他說你怎麼會有這樣的不同.
這個時候.
他說他的見證.
他說沒有.
他說我也知道這些東西不值錢.
他說不過我覺得.
我信了主之後.
我相信耶穌之後.
我不應該拿這些東西.
那經理就跟他說.
他什麼叫信耶穌.
他沒有聽過.
他就有機會跟他傳福音.
他沒有想過.
他完全沒有想過.
他經理信了主.
後來經理對他很好.
因為他工作做得好.
很少的錢財.
很少的東西.
是可以轉化的.
如果我們很努力去工作.
不要將這個世界移分.
透過在世界上.
道德的益處和好處.
與人的相處.
原來可以轉化成為永恆的價值.
因此經文就特別要提這個字.
不義是沒有道德判斷.
不義不是說世上的錢不道德不好.
我們不要有這樣的觀念.
也不要看不起那些人.
說他們在外面拼命.

$^{561}$是應該的.
不過他要知道.
他要拼命.
不是為世界.
是為上帝.
這個錢財就可以被轉化.
所以11節.
若是你們就是不義的錢財上不忠心.
所以耶穌叫我們.
要在不義的錢財上忠心.
有沒有留意到.
換句話說.
你在外面的不義中要忠心做好你的不義.
那不是一個價值判斷.
誰還把真實的錢財託付給你們呢.
真實的錢財就屬靈了.
如果在世界上我們都不做好本分.
做好一個人.
我怎樣能做好基督徒呢.
怎樣能在天國上為神做大事呢.
就是這個意思.
所以我們要在生活上.
無論是家庭或是世界上.
我們的本分原來是有神聖的意義的.
是可以轉化的.
然後神就將真實的.
就是屬靈的.
產生永恆效果的.
交託給我們.
這個就是耶穌基督當時跟門徒所說.
做好我們的本分.
因此這個經文.
最後就是一個總結.
你記不記得最後的總結是什麼.
就是一個人不能夠事奉兩個主.
我看看有沒有打出這段經文.
就是這個屬靈.
不單只是屬靈的表面.
我們再看下去.
我不說這個了.

$^{601}$我講最後一段經文.
剛才那段經文就是.
如果你記得.
一個人不能事奉兩個主.
不時互謝過 恨那過.
不能事奉上帝 又事奉馬門.
很多人就將這句放在另一段.
其實不是的.
因為這句和前面是一段的.
所以這是結尾.
耶穌基督跟門徒說.
用這個做結尾.
一個人不能事奉兩個主.
不時恨謝過 愛那過.
受謝過 慶那過.
你們不能又事奉神 又事奉馬門.
是總結上面的.
是什麼意思.
因為上面太多人.
將這個世界二分.
將這個世界變成兩個主人.
一個就是上帝.
一個就是金錢馬門.
禮拜一到禮拜六.
我們為金錢就金就奴隸.
世界是我們的主人.
禮拜日就回來敬拜上帝.
耶穌說你這樣將世界二分.
就有兩個主.
我們只得一個主.
因為當你一二分了兩個主.
你的生命就會分裂.
你會重罪慶那過.
如果我們認定這個世界是一個主.
連世上出糧給你的主人都是上帝.
他不是主人.
統管萬有的才是主人.
這樣我們生命就不會分裂.
我們也不會隨便輕看這個世界.
輕看你的工作 輕看你的本分.

$^{641}$一個主.
我在屋泰華宣導會的末會的時候.
教會有一個很好的長老.
很熱心愛主.
他在加拿大政府國防部工作.
他的見證很感動我.
當年他在Alberta大學.
一個很出名的大學.
他讀太空飛彈.
博士學位.
他成績非常高.
他畢業的時候申請國防部的工作.
而國防部的工作只有他一個人.
有資歷去申請.
於是他就申請了工作.
很快就見他.
然後他就填寫了所有東西.
填到最後就有一行.
就是你做什麼都會問.
你有沒有犯過刑事案.
你有沒有坐過牢.
因為國防部政府一定要審查的人.
沒有 沒有 交上去.
聖靈就責備他.
你為什麼要說謊.
你犯過刑事案 你坐過牢.
當他心中不平安.
他第一認識.
他曾經在上帝面前.
祈禱立誓.
一個主.
加拿大都不是他主.
只是上帝的主.
他和那個官.
你給我申請表.
我有些東西填錯了.
他查了.
我犯過刑事案.
我坐過牢 遞上去.
接收到 你搞錯了.

$^{681}$你填這件事 你想不做.
他說 我填了.
結果呢 失去工作.
過了三個月.
沒有消息.
他在溫哥華找了一份工作.
做電話的 電子的.
做了半年.
突然收到政府的電話.
做了半年.
他看回政府的網頁.
即是招聘.
職位仍然空缺.
他覺得 還沒有人做.
於是他再申請.
這個時候 政府就馬上請他.
申請他.
但因為他填了那兩份.
我坐過牢 我的刑事案.
政府就說 你可不可以解釋一下.
你為什麼犯過刑事案.
為什麼你坐過牢.
於是呢 他就解釋.
他說 他當年是難民.
由福建 走到印尼.
在印尼那裡.
讀中學 十八歲那年.
他想出國去加拿大讀書.
但因為他是難民 他沒有護照.
他沒有護照 去不了加拿大.
他申請加拿大的大學收他.
去不了.
他就要申請什麼.
因為中國沒有護照.
他要用各種方法.
但都沒有辦法.
人們就說 他在印尼很容易.
找一個假的護照.
給錢 官員看錢就不會看真假.
教了一門路 他十八歲那年.

$^{721}$於是就用朋友教的方法.
買一本假的護照 去出境.
誰知一出境在機場就被抓了.
就是他用假的護照.
於是一判 坐十九年牢.
當時他絕望.
他說我十八歲 坐十九年牢.
他說我的人生都沒有了.
他在獄中祈禱.
他是基督徒.
他覺得他不應該這樣做.
他將生命交託給上帝.
他說從今以後 我只認一個主.
就是耶穌基督 世界不是他的主.
在獄中他坐了三個月.
後來不知什麼原因.
可能覺得判得太重等等.
放了他出來.
放了他出來 他仍然沒有護照.
於是他就申請了一個CIA.
申請了台灣 台灣給他一個CIA.
他就可以去香港 很短時間.
只有三個月 馬上離開去香港.
怎知道中央阻滯.
去到香港就過期了.
滯留在香港.
他機票已經買了去加拿大.
他在這裡祈禱.
他說主啊 真是絕路.
怎知飛機航空公司說.
不要緊 他說加拿大很寬鬆.
特別對於難民各方面.
他們照相機 去到加拿大下飛機之後.
你才去解釋.
他就憑了過期的CIA.
台灣給的CIA 去到溫哥華.
加拿大政府真的很好.
知道他的情況.
就發了一篇文 收了他.
給了他入境.

$^{761}$現在給了他讀書.
因為他考了大學.
就開始了.
這個就成為他認定在這個世界上.
能夠掌控生命 是上帝.
所以沒有兩個主 只有一個主.
在他以後的工作上.
他都認定這件事.
只有一個主.
只有天上的一個.
後來他還有很多經歷.
就是他那份工作.
後來他再申請的時候.
原來有一個人去爭的.
一個英國很出名的劍橋博士.
跟他讀一樣的東西.
那個人已經上了工.
本來他沒有份的.
但怎知那個人去滿地可的時候.
遇到大雪.
他太太就罵他.
他那些鬼地方不是人住的.
要回英國.
他那份工作的人就走了.
所以那個位置就空了出來.
後來他見到他之後.
劍橋那個英國人就說.
我不做 所以你有得做.
他就跟我說.
不是 是上帝讓我做.
不是你讓我做.
很清楚.
只有一個主.
整段經文雖然好像很複雜.
但耶穌教導我們要很盡本份.
因為盡本份.
我們都是面對上帝.
不是面對世界.
我們今天事奉主.
是很有福氣.

$^{801}$神給我們很多很多的報酬.
我們一同祈禱.
因為你是我們的主.
你救贖我們 你是我們生命的主.
在世上雖然有很多的人自認為主.
無論是政府.
無論是外面的老闆.
無論是各樣的才幹的人.
他們都不是.
只有你是我們主.
你是我們生命的主.
並且是我們一切的主.
求你祝福我們.
當我們盡本份在世界上.
這個本份會轉化成永恆的價值.
讓人得到你永恆的生命.
讓我們能夠去愛護我們身邊的人.
愛護主給我們世上的一切.
因這些都是你給我們的 是美好的.
我們祈禱奉主名求.
\newpage



\section{創世記 22:1-14}
\label{sec:J_KPU0uZysw}
\textbf{2023-04-29 晚堂《我在這裡》創世記 22\_1-14 講員:褚永華牧師}
\newline
\newline
連結: \href{https://youtube.com/watch?v=J_KPU0uZysw}{\texttt{https://youtube.com/watch?v=J\_KPU0uZysw}} ~~~~ 語音日期: 2023-05-02
\newline
\newline
\hyperref[sec:l5hkKLKHhA0]{\small{< < < PREV SERMON < < <}}
~
\hyperref[sec:index_chronic]{\small{[返順時目]}}
~
\hyperref[sec:index_scriptual]{\small{[返順卷目]}}
~
\hyperref[sec:bNffFxPN8ik]{\small{> > > NEXT SERMON > > >}}
\newline
\newline
創世記 22:1-14
\newline
\begin{longtable}{cl}
\hline
\hline
章節 & 經文 (和合本修訂版)\\
\hline
22:1 & \begin{tabularx}{0.7\textwidth}{X} 這些事以後，神考驗亞伯拉罕，對他說：「亞伯拉罕！」他說：「我在這裡。」 \end{tabularx} \\ \\ \relax
22:2 & \begin{tabularx}{0.7\textwidth}{X} 神說：「你要帶你的兒子，就是你所愛的獨子以撒，往摩利亞地去，在我指示你的一座山上，把他獻為燔祭。」 \end{tabularx} \\ \\ \relax
22:3 & \begin{tabularx}{0.7\textwidth}{X} 亞伯拉罕清早起來，預備了驢，帶著跟他一起的兩個僕人和他兒子以撒，劈好了燔祭的柴，就起身往神指示他的地方去了。 \end{tabularx} \\ \\ \relax
22:4 & \begin{tabularx}{0.7\textwidth}{X} 到了第三日，亞伯拉罕舉目遙望那地方。 \end{tabularx} \\ \\ \relax
22:5 & \begin{tabularx}{0.7\textwidth}{X} 亞伯拉罕對他的僕人說：「你們和驢留在這裡，我和孩子要去那裡敬拜，然後回到你們這裡來。」 \end{tabularx} \\ \\ \relax
22:6 & \begin{tabularx}{0.7\textwidth}{X} 亞伯拉罕把燔祭的柴放在他兒子以撒身上，自己手裡拿著火與刀；於是二人同行。 \end{tabularx} \\ \\ \relax
22:7 & \begin{tabularx}{0.7\textwidth}{X} 以撒對他父親亞伯拉罕說：「我父啊！」亞伯拉罕說：「我兒，我在這裡。」以撒說：「看哪，火與柴都有了，但燔祭的羔羊在哪裡呢？」 \end{tabularx} \\ \\ \relax
22:8 & \begin{tabularx}{0.7\textwidth}{X} 亞伯拉罕說：「我兒，神必自己預備燔祭的羔羊。」於是二人同行。 \end{tabularx} \\ \\ \relax
22:9 & \begin{tabularx}{0.7\textwidth}{X} 他們到了神指示他的地方，亞伯拉罕在那裡築壇，把柴擺好，綁了他兒子以撒，放在壇的柴上。 \end{tabularx} \\ \\ \relax
22:10 & \begin{tabularx}{0.7\textwidth}{X} 亞伯拉罕就伸手拿刀，要殺他的兒子。 \end{tabularx} \\ \\ \relax
22:11 & \begin{tabularx}{0.7\textwidth}{X} 耶和華的使者從天上呼喚他說：「亞伯拉罕！亞伯拉罕！」他說：「我在這裡。」 \end{tabularx} \\ \\ \relax
22:12 & \begin{tabularx}{0.7\textwidth}{X} 天使說：「不可在這孩子身上下手！一點也不可傷害他！現在我知道你是敬畏神的人了，因為你沒有把你的兒子，就是你的獨子，留下不給我。」 \end{tabularx} \\ \\ \relax
22:13 & \begin{tabularx}{0.7\textwidth}{X} 亞伯拉罕舉目觀看，看哪，一隻公綿羊兩角纏在灌木叢中。亞伯拉罕就去牽了那隻公綿羊，獻為燔祭，代替他的兒子。 \end{tabularx} \\ \\ \relax
22:14 & \begin{tabularx}{0.7\textwidth}{X} 亞伯拉罕給那地方起名叫「耶和華以勒」 。直到今日人還說：「在耶和華的山上必有預備。」 \end{tabularx} \\ \\
[1ex]
\hline
\hline
\end{longtable}
$^{1}$各位弟兄姊妹晚安.
我們今天晚上所讀的經文.
是大家都聽到的創世記第22章.
經文裡面將我在這裡一共出現了三次.
我希望大家在中讀的過程中間.
都能夠有這樣的領悟.
我們在家裡領修讀聖經的時候.
見到經文出現一些重複的字眼.
很容易幫助我們記憶.
那些重複的字眼.
不單止是讓我們知道分段是怎樣.
讓我們知道整個事件的進程.
事件中間的來龍去脈.
讓我們要留心的是甚麼東西.
我在這裡這個字.
希伯來文裡面叫做Heneli.
Heneli的意思是甚麼呢.
他比我們所謂點名.
點到你的名字.
Jason 到.
他比我在這裡的到.
是要多一點的意思.
就好像你在家裡叫兒子出去買東西.
不要嫌棄了.
他幫我買東西.
你兒子不會說到的.
兒子通常都是說.
等一下.
實際上他等待的意思是.
我聽到了.
但我不會馬上做的.
當在經文裡面出現.
我在這裡的時候.
在現在這個篇章裡面.
當耶和華上帝考驗亞伯拉罕.
對他說.
亞伯拉罕在這裡的時候.
他的意思不單止是.
亞伯拉罕你在哪裡.
亞伯拉罕說我在這裡.

$^{41}$但當經文再演繹下去的時候.
我們會發現到.
那個我在這裡的意義.
是有說.
我準備好了.
I am ready.
我準備好了.
又或者說我在這裡的意思是.
我準備好了.
有什麼吩咐呢.
每一個字眼的後面.
當我們自己用心去讀.
靜心去揣摩.
上帝的話語.
會向那些安靜揣摩的人.
向他顯示更多的意義.
急急腳滾水六腳那些.
你永遠都覺得經文很淋沉.
很沉氣.
你摸不到他的意思.
所以我們經常都說.
我們今天做信徒的人.
要安靜自己的心.
安靜在上帝的話語面前.
看他到底要跟我們說什麼東西.
在那個時候你會發現到.
原來上帝的話語是甘甜.
秘密甘甜.
這些話很多時候.
只是在我們的主學裡面.
或者我們上團契.
或者我們查經的時候.
讀到這些金句詞的經文.
我們背了它.
但是在我們自己的生命.
在哪一個時刻.
在哪一些時段.
在我們的生命.
去到什麼事情.
發生在我們的生命裡面.

$^{81}$與上帝自己的話語有交接的時候.
讓我們真正的體會到.
上帝的話語是秘密甘甜.
你再想深一層.
是不是上帝的話語是甜.
好像我們喝咖啡.
加糖的甜呢.
你再想深一層.
原來那個甜.
不是那些甜.
是你真是有一種.
很被照顧的.
幸福的感覺.
原來我在上帝自己的手中.
沒有東西逃得過你的眼睛.
我的困難.
沒有一樣是你不關心的.
哎呀真是甜啊.
你的話語.
在這一刻.
形容在我的生命裡面.
好像咖啡加了糖一樣.
不是加了白糖.
就是加了咖啡黃糖.
很香.
很甜.
很信心.
上帝啊.
Thank you.
你真是知道我.
所以我們每逢讀到一些經文.
有些字眼在出現的時候.
它就要我們去揣摩它.
到底那個是什麼意思.
有時候經文裡面.
給我們的意思傳回到.
是在我們喜樂的情況之中.
那個兒子考那個女兒考DSE.
哇 幾多封五星星啊.
一有經文到.

$^{121}$一到的話.
Jason.
我兒子有五星星.
I am here.
但是有時候.
在你的生意不順利.
上帝的話語.
臨到你的時候.
賞賜的是耶和華.
收取的又是耶和華.
耶和華的命是應當稱重的.
你能不能夠說得出.
I am here.
當你感情生活不順利的時候.
又失戀了.
真言裡面有句話.
房屋財產是祖宗所遺留.
惟有嚴厲的妻.
是耶和華所賜的.
耶和華上帝.
I am here.
你能不能夠在那些生命.
最困難的關節口上.
能夠說得出.
我在這裡呢.
我猜在我們這一代的人裡.
從來都猜不到.
要去到現在這樣的階段.
到我們這樣的年紀.
要想自己生命的去留.
我以為我爸爸那一代已經定了.
為何來到我們這一代的時候.
我們還要做這麼痛苦的抉擇呢.
我們在這些抉擇裡面.
我們能不能夠跟上帝說.
上帝啊.
I am here.
當這段經文.
我們這樣思考的時候.
我們就說.

$^{161}$上帝啊.
你跟亞伯拉罕.
說這些話.
上帝你不覺得你很殘忍嗎.
亞伯拉罕75歲離開加拿大這麼不容易.
按照專家的統計.
年代的推算.
他這一刻.
耶和華上帝吩咐他.
亞伯拉罕.
將你的兒子.
將你獨生的兒子.
將你所愛的兒子.
大摩利亞地那裡獻給我.
他已經110歲了.
中間經歷了35年的時間.
他在等待.
等待.
等待上帝的約.
亞伯拉罕.
我要使你成為大國.
我要將一個子孫後代賜給你.
你的子孫.
要使萬人得福.
亞伯拉罕75歲離開加拿大這麼容易.
一路等.
等到他等不及了.
就下甲生了20馬利.
上帝說20馬利不是我要的.
是我的位置.
再等.
然後有一天.
上帝去到帳篷裡.
與亞伯拉罕相遇.
跟他說.
明年這天.
你的太太.
哈.
撒拉.
將要生一個兒子.

$^{201}$他終於等到了.
直到115歲.
以撒也已經十幾歲了.
上帝竟然吩咐.
他說.
將你自己的兒子.
將你獨生的兒子.
將你所愛的兒子.
帶上去摩尼亞地.
我想亞伯拉罕在那一刻.
跟上帝說.
上帝啊.
你的英文不是這樣的.
你的約不是這樣的.
為什麼你的約給了我.
到這一刻.
我75歲離開加勒底的吾爾.
你告訴我.
我要將他成為大國.
我的子孫就像天上的星.
海邊的沙一樣多.
為什麼你會這麼殘忍.
在我115歲的時候.
110歲的時候.
你叫我做這麼殘忍的決定.
上帝難道你忘記了嗎.
你忘記了你和我納的約嗎.
我想我們每一個人.
在我們生命的某些時段裡.
都會和上帝問這樣的話.
我納至跟從你.
我相信你.
我倚靠你.
你的生命到了這一刻.
是這麼難忍受的嗎.
上帝啊.
我想不通.
上帝啊.
為什麼你可以這麼殘忍.
說過的話不算數嗎.

$^{241}$你答應我的話.
為什麼還會這樣改變.
叫我將我的兒子.
我所愛的兒子.
拿去摩利亞地獻為飯祭給你.
上帝啊.
在這一刻.
你會發現到.
有一個很奇特的現象.
如果你留心讀聖經.
留著它對著它久.
不要滾水六腳一下就過.
留心聽它說.
留心看著它.
你就會發覺到.
阿北海沒有說話.
你說他說了話.
你沒有記下.
如果每樣都記下.
聖經就長得不得了.
在這裡我們相信.
上帝有很多時候記載一些事.
上帝有很多時候選擇.
不記載一些事.
但在這裡我們看到.
當耶和華上帝跟他說.
耶和華上帝要考驗.
他就跟他說.
將你的兒子.
將你的獨生兒子.
將你所愛的獨生兒子.
熱殺到摩利亞地獻給我.
當上帝跟他說.
阿北拉罕就說.
我在這裡.
生命去到這一刻.
或者我們自己的生命.
我們的生活.
生命裡面的困難.
生命裡面的喜樂.

$^{281}$去到那一刻的時候.
那個我在這裡.
是很難說得出的.
又或者說.
在我們自己的生命之中.
這個沉默了的阿北拉罕.
我們怎樣理解他呢.
有一位哲學家.
丹麥的哲學家.
奇克果.
在他寫的一篇.
哲學的短的論文裡面.
恐懼與戰略裡面.
就很長篇講這個故事.
他用他一貫的.
西方人存在主義的一個立場.
來講.
你想想一個做爸爸的人.
一個跟上帝納了約的人.
當上帝說.
將你的兒子拉去摩利亞地獻給我的時候.
他走了三天的路程.
沉默的三天.
他說作為一個爸爸.
怎會不說話.
不想事呢.
他的內心裡澎湃到不得了.
他內心裡充滿了很多的.
想不通的問題.
千百萬條問題.
但是每條都沒有答案.
他很想問上帝.
直呼蒼天.
天啊.
為何你會這樣對我呢.
為何會這樣呢.
我想我們每一個人.
都可以代亞伯拉罕.
問這些話.
上帝啊.

$^{321}$你真的不講道理的.
上帝啊.
為何你不可以通情達理一點呢.
上帝啊.
你可不可以對我好一點呢.
上帝啊.
為何你會這麼殘忍的對我.
對這個耳塞呢.
十幾歲的青少年.
他做錯了什麼.
你要將他獻為煩祭呢.
叫爸爸自己來獻為煩祭呢.
要將自己的兒子殺在祭壇上.
上帝啊.
你太殘忍了.
其中一個在他的論文裡.
用很長的篇幅來解釋.
來講.
亞伯拉罕為何會這樣做呢.
一聲不吭就帶著兒子離開了.
我們不想推論太多.
他很怕.
撒拉的.
撒拉叫他做什麼他都會做什麼.
這次我們試了一下.
他生的兒子可以成為我的兒子.
亞伯拉罕很聽話.
就拿了夏革做妾侍.
生了二十馬.
我們不知道他為何會這樣.
但是當上帝呼喚.
亞伯拉罕.
他就說我在這裡.
我預備好了.
你有什麼要說呢.
請說.
如果你用這樣的心理去看.
你會發覺到.
在舊約經文裡有很多人.
存著這樣的心來應對上帝.

$^{361}$撒母耳啊.
在聖念裡事奉學習.
事奉上帝.
然後到夜間聽到聲音.
三次之後.
伊利都跟我說.
我沒有叫你.
伊利跟我說.
如果再有聽到聲音.
肯定是夜話上帝跟你說話.
你就說.
主啊.
請說.
僕人.
敬聽.
I am here.
以賽亞呢.
誰因為我沒去.
我沒可以浪費誰呢.
以賽亞就說.
我在這裡.
I am ready.
在這一刻.
亞伯拉罕的心情是極之複雜.
要將自己的兒子.
死在自己的刀下.
見為凡制.
為了應驗這位上帝.
立約的上帝.
答應我的上帝.
被見為凡制.
我真的想不通啊上帝.
亞伯拉罕會說.
我想不通.
但是中西裡.
都有很多人.
用不同的方法.
用不同的文字來表達.
到底這裡的經文.
我在這裡.

$^{401}$手起刀落.
要將以賽見為雞蛋上帝.
那個心路的歷程.
屬靈的教訓.
到底是怎樣的呢.
有很多中外的作家.
借著他們的作品.
借著他們的文字.
來對我們講解.
其中一個就說.
這是一個信心的調柔.
我想不通.
極之矛盾.
paradox.
互相矛盾到天翻地覆.
我想不通.
但是我憑著信心跳過去.
死就死了.
希伯來信的作者就說.
原來亞伯拉罕在那一刻.
完全明白.
就算我將以賽.
牽在雞蛋上.
儘管他死了.
上帝都可以叫他.
從死裡復活.
應驗上帝自己的.
對我所立的約.
Kurt Giger就說.
亞伯拉罕有這種.
Lip of Faith的勇氣.
Lip of Faith.
信心的調柔.
明知不行.
明知矛盾.
但是上帝說.
我信你.
死就死了.
就是這樣一步.
我如何去理解這些心情呢.

$^{441}$有時候我可以領略到一點點.
就是什麼呢.
有時候我做事都做得.
就是.
就是心思很.
入了去.
不要想其他事.
不要做其他事.
朋友發WhatsApp來給我.
我都聽不到聲音.
又寫了東西之後.
這樣寫不是很好.
等一等.
等一等就忘記了.
有時候就說.
不要再想了.
就這樣篤了他就算了.
死就死了.
他喜歡也好.
不喜歡也好.
就這樣篤了他.
稍微體會得到.
死就死了.
不要再想了.
想多多多.
夜長夢多.
都不知道怎樣住.
在這個WhatsApp的世界.
WeChat的世界裡面.
上限是一分鐘兩分鐘.
你不回答他.
就被人罵了.
所以養成我們每一個人.
都是有一種很急的性格.
就快快速速回答.
快快速速回答.
回答完後.
後悔的用那些字眼.
回答完後後悔.
我不是這樣寫.

$^{481}$我不是這樣的意思.
回答完後才算.
當其他人用這樣的心境說.
在你的困難.
在你極端的困難.
極端矛盾.
想不到的過程之中.
你的信心仍在.
死就死了.
就這樣了.
但是後面最終結的意思是.
上帝啊.
我比你還要信.
就算我有錯.
但是你一定沒有錯.
我比你還要信.
你的帶領是最好的.
我這一刻想不通.
但是我知道.
我是要這樣做的.
你的話語鼓勵我這樣做.
我禱告的時候.
你四年感動.
我要我這樣做.
我身邊的知心熟練的朋友為我祈禱.
他鼓勵我要這樣想.
要這樣做.
上帝啊.
我雖然想不通.
但是我仍然這樣做.
就算撞了板.
不要緊.
因為我知道.
長遠地.
你知道就夠了.
當故事一路進行的時候.
就發覺到.
他起程去了.
走了三天的路程.
內心洶湧澎湃.

$^{521}$但是口裡一聲都不出.
我們讀經的人讀到.
亞伯拉罕你真是厲害.
很信服.
這種很痛苦的.
很radical的信服.
不顧一切地.
帶著自己的兒子.
往終極的目標去.
知道去到那裡.
他的兒子的性命會不保.
但是他說.
上帝啊.
我比你還要信.
我相信你.
所以我會這樣做.
三天靜靜的旅程.
內心洶湧澎湃.
終於兒子就問他.
阿伯.
憲制所有的用具都預備好了.
憲制的陽光在哪裡呢.
很多時候.
在我們生命裡無助的一刻.
我們都會問這樣的問題.
上帝啊.
一切都預備好了.
但是出路在哪裡呢.
我好像找不到出路.
難道這裡就是一條.
這樣的死直的道路.
我要將二十大將的一個死亡.
獻在死亡的祭壇上.
放在凡濟獻過給各位.
不講道理的上帝.
似乎是違背了他的約的上帝.
我想不通.
上帝啊.
求你在我的靜默之中.
在我的沉默之中.

$^{561}$在我的黑暗之中.
你給我一點點亮光.
你點醒我一點點.
讓我能夠看到一點點的亮光.
讓我心裡知道.
我走在正路上.
你的話語不是答過嗎.
在後面有聲音說.
你或向左或向右.
後面的聲音都說.
這是正路.
你要行在其間嗎.
上帝啊.
求你給我一點點的亮光.
讓我看到.
原來我踏在正確的路上.
上帝啊.
但是上帝在你最需要的時候.
上帝總是沉默不語.
這是不是我們很多人的經驗呢.
我們很羨慕一些人.
一禱告的時候.
內心裡面澄明一片.
光明一片.
上帝跟他說話.
有些人一看到聖經.
一看到一些地方就戳下去.
就是了.
你千萬不要這樣嘗試.
你發覺到.
上帝跟他說話.
生命裡面有一種很沉靜.
就是了.
但是絕大部分的信道的經驗.
都是在沉默之中.
上帝在沉默之中.
在我摸索的裡面.
上帝好像沉默不語.
當耳塞這樣問的時候.
亞伯拉罕說.

$^{601}$當耳塞一叫.
阿爸.
亞伯拉罕說.
我在這裡.
我不知道他用什麼心境來應兒子這句話.
當耳塞問.
阿爸.
所有的事情都預備齊全了.
憲制的陽光在哪裡呢.
那句我在這裡.
上帝必有預備.
是很痛苦的答案.
那裡是兩重的衝擊.
當上帝跟他說.
亞伯拉罕.
把你的兒子帶去摩利亞地.
獻埋繁殖給我.
亞伯拉罕說.
上帝你的應兒到底怎麼樣.
你給我立的約.
到底全部都翻轉了.
你要我怎麼做呢.
走了三天的路程.
然後兒子又問.
憲制的用具全部有.
憲制的陽光在哪裡呢.
那個我在這裡的時候.
他回答不了.
兒子.
我要怎麼回答你才好呢.
我跟你說長篇大論的神學道理嗎.
我跟你說長篇大論上帝沉默嗎.
我跟你說長篇大論上帝沒有說話嗎.
我要怎麼回答你才好呢.
但是那個答案亞伯拉罕已經知道了.
三天的旅程裡面.
他還沒有起程之前.
他都知道.
上帝早就告訴他.
要將你的兒子.

$^{641}$將你所愛的獨生的兒子.
以撒牽在祭壇上.
但是亞伯拉罕在那一刻.
他只能夠跟以撒說.
上帝必有預備.
那句話我在這裡所表達的就是.
我相信上帝.
我肯定相信上帝.
在我的無知裡面.
在我的無奈裡面.
在我的不肯定裡面.
在我的矛盾裡面.
在我對上帝懷疑的時刻裡面.
上帝啊.
我仍然選擇相信你.
我仍然選擇信服你.
雖然我不明白.
世界的聲音告訴他.
你這麼相信是盲的嗎.
是盲信啊.
但是其中一個就會在我們身邊說.
No.
A leap of faith.
明知前面有一個很大的心願.
想不通的心願.
矛盾之極的心願.
不可不忘.
最近看到.
幾本書.
現在都要跟基督教書局說.
做一些廣告.
我們現在的人很不喜歡買書.
我鼓勵大家買書.
光顧一下.
基督教的出版.
因為印了書.
沒有人買.
很痛苦的.
做一些決策.
要印哪本不印哪本.

$^{681}$都是很大的決策.
印一本書要用很多錢.
買不到的位置.
全倉放在倉裡又要很多錢.
但是很多基督徒不喜歡讀書.
這是一個很大的矛盾.
實際上在書本裡.
你可以將別人的經驗.
別人的屬靈的經驗.
言經的經驗.
透過文字向你傳遞.
讓你可以跟作者深認心地.
了解上帝在他的生命裡所做的事.
所以讀多一點書.
幫助我們體會到.
他人經驗裡的上帝經驗.
提升我們自己屬靈的觸覺.
是好的.
鼓勵大家有時間去走走.
基督教書局.
七月的時候又有基督教書展.
可以去買書.
很多人在基督教書展最後一天.
去撿便宜的書.
拿個大行李箱去撿.
鼓勵你做這些傻事.
可能你家裡沒有那麼多地方.
來等那些書.
等在書架上看一看都很開心.
買些書回來等.
讓你看.
其中有一本書.
很觸動我的.
那本書是絕版的書.
不過很幸運.
有個猶太的網站叫Safari.
Safari.
一個網站.
是一個猶太人他們雜誌建立的網站.
將一些猶太的哲學家也好.

$^{721}$拉比也好.
寫的書.
將它掃描之後.
放上網上.
然後給我們免費閱讀.
在網上可以讀.
其中這本書.
就是由一個經歷過大屠殺.
Holocaust 大屠殺.
幸存的.
生存下來的一位作家.
叫做Elias Berkowitz.
B-e-r-k-o-v-i-t-s.
Berkowitz.
Berkowitz寫的書叫.
With God in Hell.
當我看到這本書名的時候.
已經覺得.
咦,不太適合我們基督教.
我們基督教相信.
耶穌同在就是什麼.
你沒唱過這首歌嗎?.
小玲小芬應該唱過這首歌.
耶穌同在就是天堂.
為什麼With God in Hell.
所以一看到這本書的書名的時候.
咦,怎麼這麼邪門.
邪邪的.
我們相信.
上帝耶穌與我們同在就是天堂.
但這個人所寫的.
乃是將他生命裡.
在集中營裡遇到的.
那個痛苦.
比喻是他自己的地獄.
他說在我的地獄裡.
上帝仍然在那個地方.
與我相遇.
那整本書裡.
就在說那些困難的地方.

$^{761}$在書裡會說到很多的例子.
因為寫大屠殺的作品.
有很多不同的觀點與角度.
一會兒我會提到另外一本.
我為什麼要提這些人的作品呢.
就是這些人可以幫助我們提升我們.
當我們遇到生命裡極端的困難.
極端的不解.
好像無法解的困境裡.
那些人是怎麼想的呢.
怎麼幫他脫離這個苦難呢.
我們透過他們的文字.
他們怎麼解釋他們的經驗.
來提升幫助我們.
讓我們遇到我們生命裡.
這樣的困難的時候.
我們可以多一條出路.
多一道想法.
伯克利斯就說.
在集中營裡.
每天都是與死亡搏鬥.
每天都是與死亡鬥爭.
我們想不通.
在一個充滿死亡氣息的集中營裡.
竟然有人會這樣做.
我們知道猶太人祈禱的時候.
一定會有些器具幫助他們.
有些猶太人中文就帶一個經文盒.
帶一個頭像一個泵一樣.
就把東西拔出來.
叫做Tifoline.
希伯來文叫做Tifoline.
是一個經文盒.
因為舊約教導我們.
要把耶和華上帝的話語.
繫在額上為經文.
所以他們祈禱的時候.
一定要守這樣的規矩.
有一條帶子.
帶子中間有個盒子.

$^{801}$盒子裡就放了一個書.
放了一個經文卷.
用手寫的經文.
放在盒子裡.
然後把帶子綁在自己的頭上.
不想自己.
不願意自己.
不希望自己.
強迫自己.
一定要記得.
上帝的話在這裡.
還有他們要繫在手上為經文.
所以有條帶子綁在手上.
一邊綁一邊打一圈.
要打七個圈.
然後綁在自己的中指上.
上面又有一個經文盒.
他說我想不通.
他說在那個集中營裡.
因為那些警察.
那些禁衛軍.
會來不及來守他們的集中營.
守那些非法物品.
譬如說有些人.
把吃東西的刀.
這個沒有刀了.
那些匙羹.
磨成尖銳的刀子.
可以打架用來刺人等等.
他們經常都會來.
來他們的營房裡去守.
有時候守到他們不喜歡.
他們就收.
所以這些祈禱的用具.
就被人收走了.
但是那些人就在想.
原來負責收.
收到去燒的那個.
是一個波蘭籍的猶太人.
所以想盡辦法.

$^{841}$就跟那個波蘭籍說.
喂 通融一下.
給我那些東西.
我祈禱要用的.
我帶著你的那個.
是在手上的經文.
戴在手上的經文.
那些器具.
我一戴上的時候.
很多人圍著我一起祈禱.
是我們猶太人祈禱.
你也知道.
是猶太人祈禱的時候.
必要的器材.
請你幫幫忙.
不要燒那些.
那些人就說.
行啊 用什麼來換呢.
找東西來換.
在書裡面記載.
他們用一天的糧食.
一天的糧食.
來換那兩件東西.
握上為經文的盒子.
手上為經文的經文盒.
你說 嘖.
就我一天就算了.
我們今天我一天就算減肥.
他們我一天就少了.
得救的一天的機會.
你就說缺少那一天的糧食.
熬不過那一天.
第二天就解放了.
那怎麼辦呢.
每一個人都想盡辦法.
保護著自己那一天的食物.
一碗菜湯.
一塊麵包.
維持著自己的生命.
第二天早上又要出去.

$^{881}$起路也好.
捉鐵路也好.
做一些很辛苦的工作.
然後晚上回來.
又是一片麵包.
一碗很稀的湯.
就靠這樣維持著自己很脆弱的生命.
讓自己可以熬到有一天.
盟軍可以進入他們.
讓他們知道原來有一班這麼悲慘的人.
在悲慘的人性之下.
讓人受這麼多的痛苦.
我一定要保持我的生命.
直到那一天.
但是人們就說.
你來換吧.
找什麼來換呢.
找你的糧口來換吧.
將你的糧口拿了之後.
給了我之後.
我就可以多一天的機會.
生存.
讓我可以活著.
走出你的集中營.
他說我不明白.
為什麼還有一個這麼昏天黑地的人.
他還這麼堅持.
要帶著那些經文盒.
手上的經文帶來祈禱.
而那些人圍著他們來祈禱.
將就一點吧.
上帝在這裡嗎.
聽你祈禱嗎.
算了吧.
他就說.
他所見到的人.
很執著地.
就是這樣帶著經文盒.
逼著祈禱.
用自己寶貴的糧口.

$^{921}$來換取祈禱的用具.
令我讓自己的生命.
少一天獲救的機會.
不放棄祈禱的權利.
我希望我們有這樣的祈禱心.
但是另外一個作者.
或者我再說另外一個.
Berkowitz.
我們知道在那些困境裡.
他看到猶太人生了嬰兒之後.
他們要為嬰兒行割禮.
然後要做一些禮儀.
生活這麼艱苦.
生出來等於送入火爐焚化爐燒.
你還生來做什麼呢.
一生出來就有這樣的苦難.
你不如不要生就好了.
但是他發現有一群猶太人.
看著嬰兒.
看著歡欣高敬念頌土文.
讓這個死寂的世界.
可以有一線的因因向榮.
Berkowitz說我真的想不通.
為什麼會有這樣的人.
在信仰裡可以這樣極致地.
我信上帝呢.
好幾年前我看過.
另外一個猶太作家的信.
一個書叫做《耶》.
Berkowitz那本《Bit God and Hell》.
就沒有翻譯.
只有英文沒有中文翻譯.
但這本就有中文翻譯.
我鼓勵大家如果有機會.
可以去書局買來看.
又或者去公共圖書館.
可以借的.
由一個作者叫做維素.
Eli Vizzo.
他1987年.

$^{961}$諾貝爾和平獎的得主.
很有幸我在外國讀書的時候.
聽過他一次演講.
是很會演講的人.
因為他十幾歲的時候.
進入集中營裡.
他媽媽在集中營裡死了.
姐姐在集中營裡死了.
爸爸後來都在集中營裡死了.
一家幾個人.
七七八八都死在集中營裡.
但他將他集中營裡所見到的東西.
寫了出來.
他是一個很虔誠的人.
他說生命中我都不明白.
為什麼會遭受到這樣的苦難.
他說有一次他見到一個警衛.
將一個十元仔.
很年輕的十元仔.
掛在橫樑上.
將他吊死了.
所有猶太人就在下面看.
迫著看.
看一個生命的終結.
已經受到皮包骨.
但就是吊在橫樑上.
看著他氣絕身亡.
Eli Vizzo就說.
為什麼這麼殘忍呢.
人性為什麼可以去到這麼黑暗呢.
為什麼人可以這樣這樣黑暗呢.
在他自己的書裡面.
他問一個這樣的問題.
上帝在哪裡呢.
然後Eli Vizzo自己回答.
上帝和那個小孩.
已經一起死在橫樑上.
當我看這些書的時候.
當然我那時候還沒有看過.
Berkowitz的書.

$^{1001}$但當我預備的時候.
想起這些書.
想起裡面的文字.
我就說一個在集中營裡面.
差不多全家都死絕的人.
他對上帝說.
上帝已經死了.
一個就說他觀察到.
有很多的人在死亡的寂靜裡面.
仍然要相信上帝.
絕對依靠相信上帝.
丹麥的哲學家.
Kiakor就說.
積蓄信心跳過去.
Berkowitz他的經驗就說.
我站在這裡.
我就會站在這裡.
看看上帝在我的苦難.
在我的地獄裡面.
上帝要怎樣對我呢.
I stand here.
所以在他Vicodin Hell的書裡面.
他就這樣寫.
他說在我的矛盾裡面.
在亞伯拉罕的矛盾之中.
最終使他仍然堅持下去的就是.
我信服你.
我仍然相信你.
到最後亞伯拉罕上到摩利亞地.
那裡不是有祭壇的了.
是要自己還要建一個祭壇.
你猜在建一個祭壇的過程中間.
以撒就手旁觀嗎.
看著他110歲的爸爸.
在自己搬石頭建一個祭壇嗎.
我想以撒也會跟著他爸爸搬石頭.
砌自己的祭壇.
當爸爸要將祭壇建好.
柴火放在上面.
用繩子將他綁在.

$^{1041}$我不明白亞伯拉罕.
怎樣在110歲將一個童子.
搬在那個火壇上面.
可能綁好了之後.
就跟兒子說.
請你自己上去吧.
很痛苦.
那些圖畫.
這些在墓地.
很心痛.
上帝啊.
我不知道你為什麼會這樣做.
兒子就問.
軒仔的陽光在哪裡.
很心痛.
內心傷痛到不得了.
只說上帝必有預備.
到祭壇砌好了.
將以撒放在祭壇上.
手去刀下的時候.
上帝就說亞伯拉罕.
亞伯拉罕就說我在這裡.
到那一刻.
亞伯拉罕的心思.
亞伯拉罕說我在這裡的時候.
意義在哪裡呢.
我相信我們可以這樣理解他.
當上帝初次跟他說.
亞伯拉罕將你的兒子.
將你所愛的兒子.
獻給我的時候.
亞伯拉罕說我在這裡.
I am here.
I am ready.
當兒子問他.
軒仔的陽光在哪裡的時候.
也說必有預備.
I know.
I am here 兒子.
But I know.

$^{1081}$上帝必有預備.
但最後要手起刀落.
最後真的要執行.
上帝對他的吩咐命令的時候.
手起刀落的時候.
也說上帝對他說.
亞伯拉罕我知道你真的信服我.
在那一刻亞伯拉罕說.
我在這裡的意思是甚麼呢.
I am still here.
我仍在這裡.
在我的苦難裡.
在我的痛苦之中.
在我如地獄一般的思念之中.
I am still here.
我仍在這裡.
我希望我們的苦難.
或者生命裡的困局.
不會像亞伯拉罕那樣沉重.
但也夠沉重.
我知道大家心裡一定有很多沉重的思維.
生命的各方面.
從教會來的.
從家庭來的.
從兒女來的.
從父母來的.
各種的困難.
不知道如何處理.
但我希望你在各種的困難裡.
你能夠對上帝說.
I am here.
我預備好了.
在我的苦難.
在我的困難之中.
在走過的困難之中.
我仍然要堅持下去.
I am still here.
我知道這個說話是很難說的.
做一些很困苦的抉擇.
很痛苦的抉擇.

$^{1121}$是非常之難的.
但願上帝憐憫我們.
實際上我還有另一個故事.
不過今天時間不夠了.
我下次再告訴你.
不過下次也要幾個月之後.
不過如果你有機會的話.
我鼓勵你去看一本書.
這本書挺好看的.
只有九章.
加上一個結尾.
是一個日本作家.
叫做遠藤周作.
這個名字很熟悉吧.
遠藤周作.
遠藤是一個很好的小說家.
他寫一本書.
但實際上那本書.
也有歷史在裡面.
叫做《沉默》.
有沒有聽過這本書.
有沒有看過那部電影.
講一班耶穌會的宣教士.
在葡萄牙去到日本宣教.
去到日本長崎.
在加沙瓜那裡宣教.
有一個神父.
做了二十年的宣教.
做了一班信眾出來.
但後來政策改變了.
很多人要受很多的苦難.
到最後這個神父.
是棄教的.
叫做棄教.
就是叛教.
放棄信仰.
他的徒弟跟著來找他.
不相信他的師父.
會做這樣的事.
我下次再講給大家聽.

$^{1161}$你有沒有發覺到.
一個天主教的.
耶穌會的宣教士.
面臨到生命裡面.
生死的抉擇.
他做了什麼抉擇呢.
那些人將信眾一個個排在那裡.
將他們倒轉插在糞坑上面.
然後在頭上畫了刀.
讓血慢慢流.
讓他們慢慢死.
除非你放棄你的信仰.
叛教.
叛教的方式.
就是將耶穌的聖像.
踩在聖像上面.
踩它.
表示你的放棄.
我願意在我們的日子之中.
生命經歷信仰的挑戰.
能夠跟上帝.
有親密的回應說.
I am here.
我們一起祈禱.
用你自己的話語來幫助我們.
亞伯拉罕的生命.
成為我們的照鏡.
在極端不解.
不知如何是好.
生命遇到極大的困難與挑戰的時候.
他選擇.
跟你的回應.
我在這裡.
最後仍然堅持地說.
我仍在這裡.
願意我們的生命.
願意我們的信仰.
同樣成長.
同樣成熟.
在我們的教會.

$^{1201}$在我們與人相交.
整個過程中間.
經歷上帝的同在.
經歷我們的成長.
多謝主隨聽我們的禱告.
奉耶穌基督的聖名.
阿們.
\newpage



\section{加拉太書 6:12-16}
\label{sec:bNffFxPN8ik}
\textbf{2023-04-30 加恩樓《我們所誇的是什麼?》加拉太書 6\_12-16 講員:鄧泰康傳道}
\newline
\newline
連結: \href{https://youtube.com/watch?v=bNffFxPN8ik}{\texttt{https://youtube.com/watch?v=bNffFxPN8ik}} ~~~~ 語音日期: 2023-05-02
\newline
\newline
\hyperref[sec:J_KPU0uZysw]{\small{< < < PREV SERMON < < <}}
~
\hyperref[sec:index_chronic]{\small{[返順時目]}}
~
\hyperref[sec:index_scriptual]{\small{[返順卷目]}}
~
\hyperref[sec:sHtcf0yp8MI]{\small{> > > NEXT SERMON > > >}}
\newline
\newline
加拉太書 6:12-16
\newline
\begin{longtable}{cl}
\hline
\hline
章節 & 經文 (和合本修訂版)\\
\hline
6:12 & \begin{tabularx}{0.7\textwidth}{X} 那些想要炫耀外表的人才勉強你們受割禮，無非是怕自己為基督的十字架受迫害。 \end{tabularx} \\ \\ \relax
6:13 & \begin{tabularx}{0.7\textwidth}{X} 他們那些受割禮的，連自己也不守律法；他們要你們受割禮，不過是要拿你們的肉體誇口。 \end{tabularx} \\ \\ \relax
6:14 & \begin{tabularx}{0.7\textwidth}{X} 但我絕不以別的誇口，只誇我們主耶穌基督的十字架；因這十字架，就我而論，世界已經釘在十字架上；就世界而論，我已經釘在十字架上。 \end{tabularx} \\ \\ \relax
6:15 & \begin{tabularx}{0.7\textwidth}{X} 受割禮或不受割禮都無關緊要，要緊的就是作新造的人。 \end{tabularx} \\ \\ \relax
6:16 & \begin{tabularx}{0.7\textwidth}{X} 凡照這準則行的人，願平安憐憫加給他們，和神的以色列民。 \end{tabularx} \\ \\
[1ex]
\hline
\hline
\end{longtable}
$^{1}$好,弟兄姊妹早安.
今天跟大家分享的經文其實一點也不難解釋.
只要你對聖經有一點歷史的認識.
你就已經很明白究竟在說什麼.
其中在引發的時候,特別是在四個保羅比教會的書信.
不斷地重複重複,就是因為當時有一個很重要的議題.
就是究竟外邦人是否需要受割禮.
這個我們叫做割禮之爭.
其實在保羅他們傳道初期的初期教會是一個非常大的議題.
也是猶太人和外邦信徒之間他們一個很大的高牆.
甚至乎曾經會引發一些很嚴重的事件.
所以你就會發覺當大家看回《士徒行傳》.
有一段時間是很急於要立即召阿保羅回去耶路撒冷去解釋這件事.
你就會發覺其實為什麼猶太的信徒他們這麼重視呢?.
因為割禮是他們整個信仰或者是猶太文化的一個很重要的象徵和儀式.
代表了上帝和人之間立約.
如果大家明白他們整個的文化都是和上帝立約.
因為這個割禮是從亞伯拉罕開始的.
就是他們的先祖開始的.
所以在這個情況下,這一定是他們最核心的一件事.
他們就要去做.
所以你會發覺在舊約的聖經裡面.
當亞伯拉罕和上帝立約的時候.
上帝是說得很清楚.
外邦人如果要成為以色列人的一份子的話.
就必須要遵從猶太人的傳統和習俗.
所以就包括受割禮.
你唯有接受割禮才能夠是他們的一份子.
所以很自然當他們伸約的時候.
他們覺得既然這個是我們所說的舊約所延伸的彌賽亞是我們的救主.
所以你們要得到這個救主的救恩.
自然就要這樣,回到核心裡面.
所以這個思想是自然延伸到基督的救恩當中.
這個對於猶太人的信徒來說是非常合情合理的.
但是保羅在羅馬書其實是解釋得很清楚.
當基督的救恩來到之後.
他將一切這些東西突破.
你記住不是刪除而是突破.
而是這個我們跟神的立約裡面.
再不是透過割禮而是透過信.

$^{41}$所以你就會發覺在羅馬書三章21到26節.
有一段很重要的經文.
但如今神的義在律法以外已經顯明出來.
有律法和先知為正者就是神的義.
因信耶穌基督加給一切相信的人並沒有分別.
因為世人都犯了罪,虧缺了神的榮耀.
如今卻蒙神的恩典.
因基督耶穌的救贖就白白地稱義.
神設立耶穌作挽回制是憑著耶穌的血.
藉著人的信要顯明神的義.
因為祂用忍耐的心寬容人先前所犯的罪.
好在今時顯明祂的義.
使人知道祂自己為義.
也稱信耶穌的人為義.
這是一個很重要的突破.
我們再不是憑著什麼去得著這個救恩.
就單單因為義.
所以你就會發覺保羅在二弗所書補充了一句很重要.
你們得救是本乎因,也因著信.
這並不是出於自己,乃是神所賜的.
也不是出於行為,免得有人自誇.
這是大家很熟悉的一句金句.
意思就是說,所以我們經常跟大家說.
洗禮重要嗎?洗禮很重要.
但你不會因為洗禮的緣故而得救.
你要搞清楚這個觀念.
所以這個問題其實你可以說.
去到新約是一個非常重要的信仰核心.
所以同樣道理,為什麼保羅盤了命.
他也要守護著這個核心.
而今天這段經文其實是想帶出.
當時有些人表面上好像在維護.
現在猶太人傳統下來的事.
我要繼續維護下去,所以會產生很多矛盾.
但保羅在這裡是一點即破.
他想告訴你,這群人根本不是維護.
而是維護這個傳統.
他們這群人其實是為了自己.
可能另一個版本,去讀剛才那段經文12-13節會更加清楚.
新普及版裡這樣說.

$^{81}$那些企圖強迫你們受割禮的人.
無非是想在別人面前炫耀.
他們不去教導只有基督十字架才能拯救人的道理.
是因為不願受到迫害.
即便是那些提倡割禮的人.
自己也沒有遵從所有的律法.
他們想要你們受割禮.
只是要藉此誇耀並宣稱你們是他們的門徒.
夠清楚嗎?.
所以我想大家弄清楚.
即是說現在發生的一件事.
就是有些人他借助著一個他們猶太傳統的律法或是要求.
但其實他最終的目的不是要達到這個要求.
他是想為自己.
在這裡的時候說得很清楚.
這裡有兩點是人為了自己.
第一他怕受迫.
怕受什麼迫?.
就是這些猶太人會迫他們.
他們本身是猶太人.
他們是這個群體.
其實對於他們來說.
這些所謂的核心信仰他們不在乎.
但他們在乎的是什麼?.
那些人怎樣對我?.
我很怕因為我不做這件事的時候.
你不只是迫他,你迫我.
火燒連環船.
所以在這種情況下我就要迫你去受割禮.
因為你肯受割禮我就沒事了.
第二件事是誇耀個人的成就.
因為這個某程度上對於他們來說.
如果今天大家有去做一些工作是需要額度.
你明白嗎?.
一個月裡面有30個額度.
其實那個客人到底是否申請了信用卡.
關你什麼事?.
但老闆要我去完成這30張就很重要.
我還能夠完成30以上的額度.
在這種情況下我就可以得到榮譽.

$^{121}$對於他們來說.
這也是他們很在乎的.
就是說他們兩個人.
兩方面就算不同的原因.
他們是很在乎人的讚譽.
他們很在乎人的眼光.
當然我經常說我們作為基督徒.
就算聖經也會教我們.
我們很多時候跟人的相處跟他們的不同.
為了讓他們能夠得著救恩.
為了讓他們能夠得著福音的好處.
我們作出一些退讓.
我們盡力與人和諧.
其實是應該的.
所以有些弟兄姊妹.
在這裡我好像戴了圖盔.
我立場不是官方立場.
有些人曾經問我一些很有趣的問題.
如果我媽去旅行.
她要一起去參觀寺廟.
我作為基督徒.
我應該拉她走還是一起去參觀呢?.
我個人的立場是陪她去參觀.
你進入一個寺廟的裡面.
因為你不是在參與裡面的儀式.
你是去陪著寶寶貴的媽媽.
她有這個需要.
你陪著她一起去參觀.
我沒有人叫你進去裝香.
沒有人叫你進去做這件事.
你明白我的意思嗎?.
這是我個人的立場.
這些事情我覺得是可以退讓的.
但是如果反過來.
升一個級別.
媽媽說既然進入寺廟.
進入寺廟就拜神.
進入家裡就叫人.
所以進來就要裝香.
這樣就不行了.

$^{161}$我可以陪你去參觀.
但我不會因為其他人眼光的緣故.
這真的發生過.
我記得有一次我陪我媽媽回鄉下掃墓.
因為我們家裡是很被尊重的.
回到鄉下的時候.
所以他們當然會做很多祭祀儀式.
我們就站在旁邊.
我媽媽就說.
平時她很小聲地說.
在香港不裝香就算了.
現在這裡有這麼多親戚.
你做一下樣子也裝香吧.
我想說是很大掙扎的.
因為我是背負著整個家族的榮譽.
你明白我的意思嗎?.
但反過來.
可以嗎?.
我跟媽媽說.
不可以.
我就說.
不如你跟他們說.
我信教的.
我媽媽就很無奈地跟親戚說.
他們信教的.
他們就問什麼是信教.
他們不明白.
他們說是拜基督.
拜十字架.
她就這樣說.
不理會他們怎麼解釋.
之後他們就說.
好.
我們就沒有參與.
我想說在一些無關重要的事情裡.
是可以的.
我們做退讓.
我們盡力和和諧.
也是應該的.
但是在真理的面前.

$^{201}$不可以.
在真理的面前不可以.
現在他們因為群眾的壓力的緣故.
最近幾年有一個字經常出現.
我覺得很有趣.
就叫做政治正確.
他們因為群眾壓力的緣故.
去選擇做一些政治正確的事情.
就是說他們完完全全在乎的.
只是人的眼光.
你怎樣去看我.
你怎樣去評價我這個人.
但我經常說.
去到這個位置我就會問.
那你的上帝呢.
你有沒有在乎過你的上帝.
他怎樣看.
所以保羅曾經在羅馬書的教導.
他說因為隨從肉體的人.
就體貼肉體的事.
隨從聖靈的人就體貼聖靈的事.
體貼肉體就是死.
體貼聖靈乃是生命平安.
原來體貼肉體的.
就是與神為仇.
因為不服神的律法.
也是不能服.
而且屬肉體的人.
不能得神的歡喜.
原來這裡分得很清楚.
就是究竟今天.
你整個信仰生命裡面.
你是體貼人還是體貼神.
這是我們信了一個很重要的分水嶺.
很多時候我們其實在這裡是很混淆的.
因為我們正在生活.
在一個很人性很制度化的社會裡.
我們很多人對我們的要求.
很多這個世界上的價值觀.
是不斷地對我們有很多提出不同的希望和盼望.

$^{241}$但原來聖經裡面在這個位置.
是很清晰的.
非黑即白的.
究竟你今天所做的每一樣東西.
其實你只能當你自己問自己.
你才會知道究竟你是體貼人還是神.
剛才我也說了.
剛剛所發生的這幅圖畫.
其實表面上他們是很光明正大的.
他們在守護一個猶太人的傳統.
甚至我剛才也說了.
原來有舊約聖經的支持.
反而這樣的時候.
我可以說他們還尚算是守護.
他們認為自己在裡面守護上帝的東西.
但不是的.
最可怕的地方.
你就是找了最神聖的東西.
去滿足自己心裡面的需要.
這是最可怕的.
所以你會發覺.
耶穌曾經和彼得發生了一件事.
這件事是極之危險的.
就是當耶穌很清楚地跟他們說.
說接下來我將會死.
但之後我會三日復活.
彼得即時地拉著耶穌說.
耶穌你不要.
這可以是出於彼得對耶穌的愛.
我不想你死.
我不想你受苦.
但耶穌的表現是什麼呢.
耶穌就轉過來對彼得說.
撒旦退我後面去吧.
你是伴我腳的.
因為你不體貼神的意思.
只體貼人的意思.
你會發覺這是一段很奇妙的經文.
彼得明明是用愛去抓住你.
但原來當人只是體貼人的意思.

$^{281}$我只是想見到你沒事的時候.
原來是會破壞了上帝的工作.
因為彼得從來沒有從神的角度去看這件事.
耶穌已經說得很清楚.
我來到世界上就是為了完成這個使命.
所以弟兄姊妹.
原來我們要很小心.
當你一不小心的時候.
你將體的意思放在人身上的時候.
上帝整個的計劃就可以破壞了.
其實當我們今天一路去走很多不同的改變的時候.
這是我們不斷反省不斷思想的事.
你會發覺我們有很多隔身.
很多教會都不會發生的事.
我們會做很多的事.
你會發覺很生活化.
我經常都說你什麼時候會見過禮堂女妹的幾張沙發.
你就會發覺其實我們在做很多的事.
不過我們不斷要問一個問題.
究竟我們在做的事是在繼續去完成上帝的計劃.
還是我們只是滿足人的需要.
我們越來越變成的只是一群人來到很開心.
很開心,很好玩的一個環境.
還是我們是一步一步透過這些軟性的工具.
帶人來到神的面前讓他認識上帝.
這個說真的一線之差.
我剛剛舉個例子.
這個星期四大家有多少人去了吃飯.
很多人去了.
其實我想告訴大家.
那晚我不斷提醒弟兄姊妹.
我們去不只是為了吃飯.
去到不只是見到大家嘩一聲.
最後大家還有驚喜見到大小姐.
但問題是大家會發覺一件事.
我們看到整個場景裡面.
我不斷提醒弟兄姊妹.
你是去見證上帝的奇蹟.
因為這件事在最初發生的時候.
是沒有人有信心.

$^{321}$是可以發展成這樣.
你會發覺整件事裡面很多元素.
我們是很感恩的.
譬如有一個很好的機會.
我看到Barry邀請了一群神魂朋友.
整個台很開心.
有些家人從來都沒有來過教會的.
來參加一家大小一起吃飯.
我們在裡面的時候聽了很多分享.
一些奇妙的事.
但如果我們只是聚焦.
嘩 厲害 好玩 精彩.
那你跟歡樂滿東華有什麼分別.
但不是這樣.
我們在裡面的時候.
我們真的在裡面見到很多人.
我聽了很多弟兄姊妹的分享.
他們覺得整個過程.
有些地方他們很觸動.
他們很感動.
所以弟兄姊妹.
原來當我們的心一定要放正.
究竟今天我們為了什麼.
而保羅在後面的回應.
正好講明了一個信徒應該有的生命核心.
但我斷不以別的誇口.
只誇我們的主耶穌基督的十字架.
因為這個十字架就我而論.
世界已經釘在十字架上.
就世界而論我已經釘在十字架上.
在主耶穌面前.
世界的價值觀和世人的目光已經死了.
不重要.
你今天信了耶穌.
其實你最應該去問的是.
上帝你認為怎樣.
不是其他人.
這裡他講得很清楚.
這個世界價值觀.
人的目光對你來說.

$^{361}$應該已經死了.
應該不再重要.
但對於我來說.
這個世界也不重要.
保羅在這裡的時候.
他講得很清楚.
現在活著的不再是我.
乃是基督在我裡面活著.
但弟兄姊妹.
我又想下地點.
看聖經不要只是去到這麼高.
而是我想下地點.
你問我.
我需不需要別人的肯定.
其實每一個人我們都需要別人的肯定.
是不是.
所以弟兄姊妹.
我不是叫大家走上山去做退隱.
這些是其他宗教信仰的觀念.
與這個世界好像絕緣.
我們不需要.
我不是在說這件事.
我是想說健康一點.
當我們做每一件事的時候.
回到聖經當中.
你不要為了追求其他人的讚譽.
你真的只為了追求上帝的喜悅去做.
但是.
其他的人.
我們其他身邊的人.
請你去學會對人多一點的肯定.
當大家一起去做的時候.
不要把尺的標準放到這麼高.
我想強調.
特別是香港.
在香港這個投訴之都.
已經不再需要投訴.
已經足夠了.
你們在外面的世界已經足夠了.
很多人對你們的評論和投訴.

$^{401}$回到神的家裡.
請大家一起多點去欣賞上帝的作為.
就等於剛才我說的.
在晚會裡面有沒有掉落的地方.
當然有.
怎麼會沒有.
七八人的.
我弟兄姊妹智慧很有趣.
走來問我.
七八人忙不忙.
然後.
會不會混亂.
你說會不會混亂.
當然有.
怎麼會沒有問題.
是不是.
但是問題就是.
整件事為什麼我總是要針對那些問題呢.
是不是.
我可以告訴大家.
其中一幕的時候.
是我開了.
原因就是我負責節目.
然後那裡的場地設計.
又不知道為什麼這麼神化.
他們的控制室.
正常應該是我在控制室.
看得到外面.
外面看不到我的.
是不是.
那裡的設計的控制室.
是裡面看不到外面的.
裡面是一點聲音都聽不到的.
但是只有一個門口.
我負責那裡.
Yoyo坐在裡面.
我要跟裡面溝通.
我不知道大家有沒有印象.
有一個心心牌.
擺在門口那邊.

$^{441}$我是完全進不去.
我走過去.
我跟弟兄說.
真的阻礙著工作.
那個弟兄呢.
我還以為.
他說是嗎.
不是.
他請我走.
說真的.
在那一刻我著了.
我想強調.
我真的著了.
但是我想告訴你.
為什麼我後來這麼快.
可以冷靜下來.
原因後來我是說.
我整件事為了什麼.
整個晚會為了什麼.
為了一口氣.
為了一點點瑕疵.
不是.
我們是為了歡樂.
然後很感恩的.
上帝的靈就在心裡面提醒我.
那個弟兄還做什麼.
他就是很努力幫人拍照.
他很努力的幫每一個.
來到的嘉賓留下他的.
他的善影.
是不是.
OK.
那你最重要是做到這件事.
就改.
所以那晚我和Yoyo.
整晚我只是問一個問題.
你的電話的通話.
有沒有上限.
整晚和她保持通話.
所以頂智妹.

$^{481}$我想大家一起去學的.
一樣東西就是什麼.
我們每一個人做的一樣東西.
不是為了其他人喜不喜歡.
而是上帝你認為怎樣.
你一定要做回這件事.
因為這個世界裡面.
這個世界怎麼看你.
這個世界的價值觀對你不再重要.
但是我們需要去活出.
對其他人生命的祝福.
對人更多的肯定.
對人更多的祝福.
讓每一個人都在裡面.
都能夠快樂.
原本這段經文去到這段可以完的.
但是補落在後面.
他再加了第十五節.
什麼才重要.
接著他說第十五節.
受割禮不受割禮都無關緊要.
要緊的就是作新造的人.
對於我們外邦的信徒.
你是覺得很正常的.
是啊.
因為我們沒有受割禮.
你明白我的意思嗎.
受不受割禮沒有問題.
但是頂智妹.
我想告訴你.
受不受割禮無關緊要.
你想清楚.
這是對於猶太人非常大的冒犯.
這個受割禮.
不是跟神立玉的記號嗎.
但是其實保羅作為.
加瑪烈的門下.
他不清楚.
他清楚.
但是他想提醒大家.

$^{521}$割禮是為了什麼.
為了割而割.
割禮是為了提醒人.
去記得你跟上帝的關係.
你是屬於上帝的.
所以才需要一個這樣的儀式.
這樣的過程提醒你.
讓你一生裡面有記號.
是這樣的.
但是當時的人就是.
我有割禮.
但我沒有生命的改變.
跟耶穌對敵的所有人.
不是全部都受割禮了嗎.
但他們卻不是在走上帝的生命.
所以保羅在這裡再強調一件事.
不要去做一些.
只有我們信仰表面的事情.
我們今天做的每一件事.
都是有目的的.
你回來敬拜.
是為了提醒你.
記得這個世界有一個創造的上帝.
你是被造的.
有一個上帝在這裡.
你今天奉獻是提醒你.
今天萬物都是從主而來.
你今天去事奉.
是想提醒你.
你今天最應該做的是什麼.
為什麼我們做事奉.
我經常說義工.
我很欣賞弟兄姊妹.
你不單是義工.
你還貼錢做.
我看到我們很多人.
但裡面的就是告訴我們.
我們今天存在生活的意義.
這個才是那個的意思.
是我們能夠祝福更多的人.

$^{561}$我們的生命.
這樣才是屬於上帝的生命.
我們今天做的每一件事.
都是為了提醒我們.
上帝他在這裡.
所以割不割禮的原因.
為什麼不重要.
因為你為了割而割.
你沒有意思.
說真一句.
我經常說我鼓勵弟兄姊妹去浸禮.
但反過來浸完.
特別我每一次.
見弟兄姊妹受浸的時候.
我都提醒他們.
浸禮不是一個結束.
有些人說浸完就搞定.
不是,是一個開始.
是一個新的開始.
是你在神面前.
你納若告訴全世界.
你屬於上帝.
所以你的生命.
需要有一個屬於上帝的表現.
你需要做更多.
在神裡面的工作.
這是對你的提醒.
所以保羅斯說.
要緊的就是做新做的人.
弟兄姊妹.
最後我想做一個很簡單的總結.
保羅在這裡經常強調一件事.
我們今天信了耶穌.
究竟我們的生命和未信之前.
有沒有改變.
我們只是在做很多宗教的活動.
還是我們生命裡面真的有改變.
你恕我這樣說.
我經常說到這裡.
我會恥辱埋怒.

$^{601}$鄧全道做傳道友做了二十年.
我見其他東西不多.
我見信徒的生命最多.
但我想告訴大家.
我覺得最可惜的那件事.
我早幾天.
我和我們的團體.
有些弟兄姊妹信了耶穌.
十幾二十年.
但接著問他.
福音書裡面某一個故事.
你知道嗎?.
他們是領頭的.
他們不知道.
為什麼不知道?.
因為沒有看聖經.
有些弟兄姊妹.
每個星期回來的時候.
很快樂.
很開心.
自己的家裡亂七八糟.
兩夫婦的關係很緊張.
和子女的關係很緊張.
弟兄姊妹,這是我們很心痛的.
我們很心痛的.
因為我們真的知道.
上帝創造的大能可以改變一切.
但如果我們只是停留了.
我們回來教會.
我們只是停留了回來聚會.
但我們從來都沒有去處理過.
我們自己生命的問題.
沒有去處理過我們自己家庭的問題.
我們和朋友之間的問題.
弟兄姊妹,這是非常可惜的.
當然我明白.
鄧全都是做輔導的.
我很明白每一個人.
特別我是男人.
我到這個年紀.

$^{641}$我都知道.
我很怕將我自己裡面的東西.
我去跟別人說.
我很怕打開自己被人指罵.
但弟兄姊妹.
求上帝讓你能夠有一個.
在上帝當中的平安和勇氣.
去處理我們生命這些很核心的問題.
以至於讓我們生命才能夠真正去改變.
說一句心底話.
我作為傳道人.
我經常提醒自己.
我很清楚.
如果我生命裡面.
有些問題未被醫治.
我是會影響大家.
我講的訊息.
我講的道一定偏差.
我一定會有很多憤怒.
我一定會有很多埋怨.
我一定會有很多不開心.
我會在我的講道.
在我的教導裡面出現.
所以弟兄姊妹.
我們的生命改變才是最重要.
所以今天講題就是.
你誇的是什麼.
你的生命能不能像保羅一樣.
誇我們的上帝.
誇耶穌基督的十字架.
能夠跟別人很從心地說.
基督的十字架的大能.
真的很偉大.
真的很厲害.
如果祂能夠改變我們的生命.
用你的生命去見證上帝的大能.
我們就同心努力.
我們就同心祈禱.
參加這次我們求你幫助我們.
你是創造的神.

$^{681}$你是醫治的神.
主讓我求主.
你自己親自的保守幫助帶領我們.
我們很需要你.
我們求你像保羅所說的一樣.
我們不願意只是做一個.
我們只有去參加禮儀的基督徒.
我們很願意我們生命被改變.
讓我們能夠去讓更多的人.
我們身邊的朋友.
見到我們的時候.
他們被吸引.
他們願意回到你.
主在你的面前.
我們知道唯有這樣.
教會才會慢慢慢慢.
被上帝你自己去使用.
我們求你轉化我們每一位弟兄姊妹.
願意我們在你面前的時候.
成為新造的人.
將一切榮耀頌讚.
我們交在神的面前.
祈禱奉教主耶穌.
得勝平次祈求.
阿們.
\newpage



\section{加拉太書 3:21-29}
\label{sec:sHtcf0yp8MI}
\textbf{2023-04-30 午堂《律法之上》加拉太書 3\_21-29 講員:陳韋安牧師}
\newline
\newline
連結: \href{https://youtube.com/watch?v=sHtcf0yp8MI}{\texttt{https://youtube.com/watch?v=sHtcf0yp8MI}} ~~~~ 語音日期: 2023-05-02
\newline
\newline
\hyperref[sec:bNffFxPN8ik]{\small{< < < PREV SERMON < < <}}
~
\hyperref[sec:index_chronic]{\small{[返順時目]}}
~
\hyperref[sec:index_scriptual]{\small{[返順卷目]}}
~
\hyperref[sec:FRK_krDh8hE]{\small{> > > NEXT SERMON > > >}}
\newline
\newline
加拉太書 3:21-29
\newline
\begin{longtable}{cl}
\hline
\hline
章節 & 經文 (和合本修訂版)\\
\hline
3:21 & \begin{tabularx}{0.7\textwidth}{X} 這樣，律法是與神的應許對立嗎？絕對不是！如果律法的頒佈能使人得生命，義就誠然出於律法了。 \end{tabularx} \\ \\ \relax
3:22 & \begin{tabularx}{0.7\textwidth}{X} 但聖經把萬物都圈在罪裡，為要使因信耶穌基督而來的應許歸給信的人。 \end{tabularx} \\ \\ \relax
3:23 & \begin{tabularx}{0.7\textwidth}{X} 但這「信」還未來以前，我們被看守在律法之下，像被圈住，直到那將來的「信」顯明出來。 \end{tabularx} \\ \\ \relax
3:24 & \begin{tabularx}{0.7\textwidth}{X} 這樣，律法是我們的啟蒙教師，直到基督來了，好使我們因信稱義。 \end{tabularx} \\ \\ \relax
3:25 & \begin{tabularx}{0.7\textwidth}{X} 但這「信」既然來到，我們從此就不在啟蒙教師的手下了。 \end{tabularx} \\ \\ \relax
3:26 & \begin{tabularx}{0.7\textwidth}{X} 其實，你們藉著信，在基督耶穌裡都成為神的兒女。 \end{tabularx} \\ \\ \relax
3:27 & \begin{tabularx}{0.7\textwidth}{X} 你們凡受洗歸入基督的都披戴基督了： \end{tabularx} \\ \\ \relax
3:28 & \begin{tabularx}{0.7\textwidth}{X} 不再分猶太人或希臘人，不再分為奴的自主的，不再分男的女的，因為你們在基督耶穌裡都成為一了。 \end{tabularx} \\ \\ \relax
3:29 & \begin{tabularx}{0.7\textwidth}{X} 既然你們屬於基督，你們就是亞伯拉罕的子孫，是照著應許承受產業的了。 \end{tabularx} \\ \\
[1ex]
\hline
\hline
\end{longtable}
$^{1}$鄧小妹 早安.
我們今天一起去看保羅在《迦太叔》第三章裡面.
有關律法的教導.
講這個律法課題之前.
我們都要花一點時間.
講保羅在第三章裡面提到恩順清二這個題目.
事實上恩順清二是保羅很重要的教導.
早期在《迦太叔》後來就在《羅馬書》第三章裡面.
很清楚的講述.
更加重要的是五百年前.
我們基督教新教.
路德在五百年前重新強調.
保羅在聖經裡面所講的恩順清二的教導.
大家很熟悉的.
不過都仍然講得很清楚.
保羅在《迦太叔》和《羅馬書》裡面恩順清二的教導.
馬丁路德在五百年前裡面更加強調.
他認為信心不單單是信仰的起點.
更加是我們整個信仰的一切.
你知道路德本來是一個天主教的修士.
在1517年之前.
當他改革之前.
其實他在大學裡面教聖經.
教詩篇,教羅馬書,教迦太書.
在這裡他開始醞釀出恩順清二一個很重要的道理.
他認為路德很深深地體會人的罪.
路德是一個很認真的人.
都是一個很出色的人.
對自己很高要求的人.
都是一個很火爆的人.
如果今天拍一套無線劇的話.
以前做的萬子亮都做得對.
很火爆,很高要求,很認真的人.
所以路德是一個很高要求的人.
都發覺自己完全符合不了上帝的要求.
所以他發覺人是一個很徹底的無能為力.
我稱之為悲媚神學.
這種悲媚不是純粹謙卑.
不是在禮貌上客氣謙卑.
而是人徹底地發現自己的心心不配和做不到.

$^{41}$所以在一個完全不能做什麼的情況下.
路德突然發現原來保羅所說的信是一個非常重要的道理.
信心是我們唯一能夠回應上帝恩典的方法.
所以信心對路德來說不是一個條件.
不是我們用信心來換取上帝的救恩.
不是因為你做了信心這個行為.
所以來換上天堂的入場券.
這樣說的話其實信心只是一個行為.
你相信能夠換上上帝給你的禮物.
這個不是路德的看法.
路德認為信心是一種無能為力.
徹底地做不到.
你只能夠相信接受.
所以路德所說的信心是一種完全悲媚無能為力的回應.
所以我們說上帝出雞出豉油.
其實完全不是.
我們連豉油都出不到.
全部都是基督耶穌的恩典和救恩.
所以如果我們將信心看為一種我們能夠做到的事.
這個就不是保羅所說的教導.
信心是一種只能夠我們無能為力之下.
只能夠單單相信接受上帝的恩典.
所以說的時候我們發現.
這個大家都熟悉的.
這個就是保羅在恩訊稱義所說的道理.
所以我們救恩不在乎行為.
不在乎我們能夠做到什麼.
包括相信.
都不是我們能夠認可做到的事.
一切都是完全出於上帝的恩典.
說了一點點有關恩訊稱義的實話之後.
我們就開始要去明白保羅在第三章裡面提到的律法問題.
我們再看一次 經文裡面第三章有一節.
保羅說 這樣律法是與上帝應許反對嗎?.
斷乎不是 若曾傳一個能叫人得心的律法.
亦就承認本乎於律法了.
但曾經百眾人都圈在罪裡.
使所應許的福恩信基督耶穌歸給那信的人.
但這恩信得救的理還未來之先.
我們就被看守在律法之下.

$^{81}$直圈到那將來的真道顯明出來.
這樣律法是我們蒙訓的師傅.
或者你所說的啟蒙的教師.
引我們到基督那裡 使我們恩信稱義.
但這恩信得救的理既然來到.
我們從此就不再是師傅的手下了.
這樣說 雖然恩信稱義是我們得救裡面所說的.
我們是完全沒有事情可以做到的.
任何的律法 割禮行為都在得救的事情裡面沒有功效.
但這樣說難道律法就沒有意義嗎.
難道律法就是一件違背福音的事情嗎.
保羅就很清楚地回答 他說不是.
保羅說既然律法是同神 應該對那媽絕對不是.
我就去想一下這個律法對於我們的生命有什麼意義.
如果稱義不是靠律法的時候.
那律法是什麼來的呢.
這個就保羅用了很生動的比喻來說明這個道理.
他就說 律法是我們的蒙訓的師傅.
我本的翻譯就叫做蒙訓的師傅.
新的版本就叫做啟蒙的教師.
想讀到經文裡面都說.
什麼叫做啟蒙教師 什麼叫做蒙訓的師傅呢.
什麼說律法就是我們的啟蒙教師呢.
原文就叫做Pedagogos.
Pedagogos就是教育的關係.
但大家不要誤會 原文的意思不是解作教授.
或者是高級的教師.
所以在聖經裡面都未能夠好好說出原本的意思.
在古代羅馬社裡面Pedagogos不是解作高級教授.
而是相反是陪小朋友讀書的奴隸.
即是負責拿書幫手去服侍孩童讀書的奴隸.
所以中文來說叫書童.
在侍衛旁邊就有望脈.
幫手拿著竹架的書童.
所以保羅說律法就是這樣.
律法是有意義的.
不過律法只是小朋友童年的導師.
是書童 是低級初階的小朋友.
即是保羅說律法就是低級初階的小朋友.
你問那些小朋友有沒有重要性.

$^{121}$中文重要是任何東西都要學.
MNLM開始的.
小朋友學游泳 開頭總是要帶個水泡.
在泳池游來游去.
學踩單車 以前都有兩個輔助輪在旁邊.
幫著你去踩單車.
BB學走路 都要學走路.
扶著扶著扶著.
水泡是重要的.
輔助輪都是重要的.
學行車都是重要的.
因為它能夠幫助我們去開頭去學習那件事.
變成一個去最初階的階段.
如果一個男男老狗在泳池裡面還背著水泡游泳.
那就很醜怪.
一個人踩單車一生都是用輔助輪的話.
那就比較醜怪的事情.
不僅僅是醜怪問題.
其實他在某個階段裡.
相反這些小鴨水泡 學行車 輔助輪.
不僅不能夠幫助我們去學習.
更加是怎樣呢.
更加去到某個階段就會去阻礙我們去學習.
總是有個位置要扔開水泡.
拆掉兩個輔助輪去嘗試去到更加新的階段.
去正式學習那件事情.
這個正正就是律法的位置.
普羅強調 律法的意義是一個階段.
它是重要的 有意義的.
但是基督耶穌的意義是一個真真正正的新的階段.
律法是有用的.
不過普羅告訴你.
律法只是我們上帝關係裡面的那些小事.
一些初階的小事 低B的程度的事情.
你不能夠一輩子都靠著這些律法的階段去做基督徒.
他們的目標都是一個 就是引我們去基督那裡.
律法就是為了去引導我們.
做一個踏腳石去引導我們去基督二那裡.
我們信仰不能夠長期停留在這些BB學校.
這些啟蒙教師 小朋友學童的階段裡面.

$^{161}$我們就去反思這個問題.
我們的信仰究竟是不是還在鴨仔水泡.
是不是仍然都是一些輔助輪胎 學行車的階段呢.
作為一個牧師 作為一個網上寫文章的人.
間中都會收到一些弟兄姊妹的奇難雜症.
很多網上的弟兄姊妹PM我問很多不同的問題.
我發覺這些問題可以歸納出一些問題的方式.
就是一些可不可以的問題.
我們基督徒有很多可不可以的問題.
我想我讀出來.
基督徒可不可以返崇拜遲到.
基督徒可不可以星期日不返教會.
基督徒可不可以不返主約學.
基督徒是否真的要受一奉獻.
基督徒可不可以炒股票.
基督徒可不可以和未信主人結婚.
基督徒可不可以吃飯不祈禱.
基督徒可不可以不靈修.
基督徒可不可以練瑜伽.
基督徒可不可以參加萬聖節.
基督徒可不可以整容.
基督徒可不可以看鬼片.
基督徒可不可以做生理測驗.
基督徒可不可以看相.
基督徒可不可以買樂彩.
基督徒可不可以轉教會.
這個很奇怪.
基督徒可不可以看中醫.
基督徒可不可以說謊.
基督徒可不可以說白色的謊.
基督徒可不可以只返網上崇拜.
基督徒可不可以吃著午餐.
參加網上崇拜.
基督徒可不可以吃著午餐.
戴著面膜做facial.
參加網上崇拜.
以上的問題全部都可以有答案.
可以給答案.
告訴你他們不可以有什麼經文.
不過我想問.

$^{201}$為什麼這麼多問題.
為什麼這麼多可不可以問題.
你問基督徒可不可以吃著午餐.
參加網上崇拜.
如果可以的話.
那你下次可以剪著腳踏腳甲.
返回網上崇拜嗎.
可不可以問題的背後.
可不可以不要這麼多問題呢.
曾經我寫過一篇文章.
那時候還是iphone7的年代.
寫了一篇文章叫.
基督徒應否炒iphone.
這樣一篇文章出來.
那時候炒iphone的年代是很熱門.
寫了一篇文章叫.
基督徒應否炒iphone.
預備了很多經文.
基督教倫理.
管家的比喻.
很多很多不同的問題.
寫了千多字出來.
解釋了這些出來.
發現這個title.
基督徒可不可以炒iphone.
這個title.
可以寫醜怪.
基督徒應否炒iphone.
當然沒有說不可以.
但這個不是我們基督徒.
應該要去想的思維.
明白我的意思嗎.
基督徒的倫理.
不是一些可不可以的倫理.
不斷地扯著牧師問.
這樣做得不得.
這樣做有沒有違法聖經.
這樣做得不得救.
不斷地很被動.
消極地不斷地去問.

$^{241}$上帝的底線.
只要上帝不擊殺你的話.
我就要做了.
只要我得救的話.
我就繼續做下去.
不斷地去尋求一條.
不會被擊殺的底線就做了.
這個不是我們基督徒.
倫理的方式思維.
這個不是基督徒的倫理.
這是一個律法的思維.
這是啟蒙老師的律法思維.
這就是那些鴨仔水泡.
低B輔助粒那些東西.
我們的倫理問題不是問可不可以.
可以就做.
做盡它.
不是不斷地尋求倫理的底線.
而不被擊殺的位置.
而是怎樣呢.
應該是要去竭力去追求最高的義.
曾祖母訓裡面教導我們.
你們要完全將你們的天賦完全一樣.
不是問既然我不可以殺人.
那我可以撕掉它嗎.
不是問我可不可以.
隨便地走一條路.
保羅教導我們說.
我們要離開這些律法主義的BB東西.
離開這些啟蒙的教師.
要去基督耶穌的義那裡.
然後我們問離開這些BB東西.
我們怎樣叫基督的義呢.
我們是否叫恩順清義.
不是靠行為.
不要靠行為的時候.
我們就沒有行為.
所以我們不是要靠行為.
但又不是完全沒有行為.
都不是問可不可以這樣做.

$^{281}$這樣就算了.
而是不斷地去追求基督耶穌的義.
登山寶訓是一個很高要求的義.
這就是我們今天嘗試思考的地方.
耶穌有教的.
耶穌有教我們應該怎樣做.
不是法理賽人那些律法的時候.
我們應該怎樣呢.
耶穌怎樣說呢.
登山寶訓裡面.
耶穌說:莫想我來要廢掉律法和先知.
我來不是要廢掉乃是要成全.
我就是要告訴你們.
就算天地都廢去.
律法的一點不滑也不能廢.
他也要成全.
然後他說什麼呢.
我告訴你們.
你們的義若不勝於民事法理賽人的義.
斷不能進天國.
一個非常非常嚴謹認真的命令.
我們問什麼叫做門徒的義.
要勝過法理賽人的義.
什麼叫勝過他們的義呢.
原文的意思不是解做勝過.
那個字是periso不是勝過.
不是贏輸的問題.
periso是什麼呢.
是可以解做多過大過好過勁過超過都可以.
總之不是贏輸的問題.
耶穌不是要你贏他們.
重點不是勝負.
而是可以是一種質數上的改變.
或者是量數上的改變.
periso這個字是解做豐富.
豐盛的意思.
abundant over的意思.
一種超越.
我上次說給你們聽是明白的.
periso是一種五餅二魚裡面用過的字.

$^{321}$當耶穌在五餅二魚裡面行神蹟的時候.
當他的食物多到爆.
多到滿舍的時候.
十二個男子都裝不完的時候.
所以就用了periso這個字去形容.
食物多到爆的情況.
完全的over超過滿舍.
多到爆滿到不得了的情況.
就是periso.
耶穌要求我們的義.
正正就是periso.
不是去到多少可不可以.
那就叫做quota夠了.
而是不斷的超越你能夠做到的.
不斷的往上尋求更加高的極限.
所以你覺得登山寶訓的教導很難遵守.
愛仇敵.
不可罵弟兄.
被人打完右邊再打左邊.
很難.
其實你可能誤解了耶穌的意思.
如果你以為登山寶訓只是一個更加高難度的律法的時候.
你可能誤解了耶穌的意思.
耶穌不是純粹想做難的事.
他們說你不可姦淫.
我叫你連想都不讓想.
他們叫你愛人如己責.
我連仇敵都愛賣.
他們叫你不可殺人.
我連你生氣都不讓生氣.
如果耶穌的義只是更加高更加嚴苛的律法主義者的時候.
耶穌只是更加嚴苛的法律裁判.
耶穌是想我們去periso.
耶穌的義不是要你從level5升到level100裡面.
不可以level5要level100不是這樣.
而是叫你不要停留在level5.
可不可以level5.
可以就做level5.
而是不斷地尋求你可以做到更加多更加好的極限.
不要以為尾線做到某個位置就可以收工.

$^{361}$不可以這樣做就可以.
不要以為去到level幾.
牧師說level幾就可以了.
那就做.
而是不斷地超越我們能夠做到的極限.
periso.
這個正正就是法律裁判的問題.
他們不是不好他們都是義的.
他們都是教人學義的人.
但他們覺得那個義是有quota的.
做完就做了.
我一個禮拜禁食兩次.
我得的都是十分之一捐上療.
如果你以為耶穌的義是一個更加高的律法的話.
我們其實和法律裁判沒有分別.
不斷地問可不可以禁.
這樣行不行.
耶穌說凡向底洞動腦的都難免受審判.
ok 那我不生他的氣.
那我暗串他.
耶穌說有人強迫你走一里路走出兩里路.
那我就走完之後我就走了.
耶穌說叫你為你的仇敵禱告.
奉傳命求阿們.
然後你就對付他.
不是這樣的.
耶穌是叫我們不斷地去periso.
不斷地超越.
不是問可不可以禁.
那就停留在這裡.
而是不斷地超越自己能夠做到的極限.
所以是的.
大家都聽過.
免費的東西是最貴的.
耶穌的恩典是免費的.
恩典是不需要靠行為換回來的.
但你卻是同時間因為這個恩典的緣故.
不斷地超越自己能夠做到的極限.
這個不是律法問題.
這個不是做了就夠的問題.

$^{401}$而是你甘心情願地無限地去回應基督耶穌的愛.
在哥林多後書這樣說.
原來基督的愛催迫我們.
人忌替眾人死.
眾人就都死了.
好要叫那些活著的人不再為自己而活.
去為那個替他死而復活的主而活.
做一個好的基督徒.
不是問可不可以給不給.
而是沒有人提你.
沒有人點你.
沒有人去監住你.
你都自動自覺地獻上你的最好.
今天大家回崇拜回來這裡.
沒有人監你的.
沒有人逼你的.
你不回牧師不會打給你問你為什麼不回.
但你卻是甘願地自動自覺甘心情願地獻上你自己的一切.
這個不是律法的問題.
這些不是交換的問題.
這些是我們自己甘願地獻上最好最好的回應.
這個正正就是上帝兒女和天父上帝的關係.
將來你們的子女都是這樣.
將來他們來找你.
你老了之後.
他們來找你喝茶.
問你喝茶多少錢.
喝茶之後分多少身家的話.
那你就麻煩了.
他們因為你很愛他.
所以他願意甘心情願地做到最好的回應你.
這個不是律法.
這個是我們因為回應基督的愛.
所以是很難的.
因為不是換積分.
不是儲coupon.
是你甘心情願自動自覺地不斷地超越自己能夠做到的極限.
所以經文裡面很好.
最後經文用保羅這樣的一句話來總結.
保羅說:所以你們因信基督耶穌都是神的兒子.

$^{441}$你們受洗歸入基督都是披戴基督了.
這個所以這個字很可圈可點.
保羅說:既然你們不在律法以下.
都不是在低B的律法主義那裡.
你們要保持在律法以上這麼高要求的時候.
超越律法.
不斷地勝過他們的律法的時候.
保羅就怯面等於門.
所以你們要這樣做.
你們要怎樣做.
你們要披戴基督.
是很難的.
長期保持自己在律法以上.
不在律法以下.
也不在律法當中.
保持自己在律法以上的回應是很高難度的.
登山寶訓的回應是很難的.
不斷地超越自己極限.
所以保羅說:面對基督的兒子.
面對這個律法.
你們要披戴基督.
這個是不是大家都聽過很多次.
一個很熟悉的字眼.
披戴基督.
這個字其實是一個有趣的字.
什麼叫披戴基督呢.
英文就是Put on Christ.
就是穿衣服的比喻.
保羅叫我們穿起基督的外衣.
所以其實披戴基督好像一個cosplay的概念.
你明不明白.
你不是耶穌.
但你要穿著一件耶穌的cosplay.
來扮演成像耶穌一樣.
所以基督徒我們不會參加萬聖節.
因為我們每天都是萬聖節.
每天都穿著一件耶穌的cosplay來做人.
這個就是披戴基督的意思.
我們要穿著一件耶穌的cosplay.
你明不明白.

$^{481}$你不是耶穌.
但你要穿著耶穌的外衣.
給大家聽一點秘密.
我中五的時候做暑期工.
我真的試過做這些工作.
在商場裡面扮演一隻老虎.
穿著很厚的公仔衣服.
那時候20元一個小時.
穿著一件很厚的老虎衣服.
很臭很熱.
因為不是我穿的.
穿著一件很臭很熱的衣服.
後面有些朋友會攻擊我.
後面的人會突然扯我的尾巴.
打我.
又捉我的手.
後來我就想了一個方法去報復.
有些朋友來的時候.
我就輕輕捏一捏他的手.
他們立即哭了.
什麼事.
我扮得很可愛.
還在那裡.
當時我還未信主.
悔改.
不應該這樣做.
後來老闆給了我一個翠兒的cosplay.
之後就沒有人搞我了.
因為翠兒是一個非常噁心的人體翠兒.
頭是這麼小.
但翠兒是這麼高.
完全不可愛的翠兒.
保羅告訴我.
我們要披大基督.
就是這個意思.
你明明不是基督耶穌.
不過你要用盡你的力量.
穿起這個披基督的外衣.
去做到基督耶穌一樣.
讓人覺得你真的是耶穌基督一樣.

$^{521}$但你知道你沒有能力.
你沒有能力去實踐律法.
但你用盡你的力量.
來裝扮.
去披大基督的外衣.
去做到最好.
所謂沒有能力裝扮四條.
我們明明不是基督耶穌.
但你仍然用盡我們的力量.
將自己推到最高.
去做到基督耶穌的義.
我記得我自己在錦宣信署.
我是第一個.
後來我帶了兩個妹妹去信署.
我那時候是中四中三.
那時候我們就跟妹妹說.
我們三個既然是基督徒.
就要做好見證給爸爸媽媽看.
那時候我兩個妹妹都是中三.
她是有點不懂事.
回家之後就不換校服.
吃零食 看電視 不做家務.
我媽媽經常罵我們.
為什麼不做家務 為什麼那麼不懂事.
就跟妹妹說我們三個一定要.
做基督徒之後.
我們一定要做好見證.
我們一定要做好家務.
完全保持家裡很乾淨.
做得很好.
令爸爸媽媽覺得我們信基督徒之後.
就真的完全不同了.
後來我們真的這樣.
三兄妹一起合力.
在家裡做好見證.
清理好家裡 做得很乖.
你問這些是不是我們骨子裡的本性.
當然不是.
我們當然不是這樣.
但她願意為基督耶穌的緣故.

$^{561}$來將自己突破極限.
我們去披大基督.
你說這些是不是很假呢.
這些不是假.
這些是我們願意在有限的裡面.
來做到最好的.
來見證上帝.
我們每個人都是一樣.
面對著律法.
面對著很多上帝給我們的規矩.
我們沒有能力去做.
但願意披大基督.
將我們最好的獻上.
嘗試去做.
在你的工作裡.
在你和家人的關係裡.
你按著你的本性就會炒.
但你仍然披大基督.
願意去超越自己的界線.
來不斷地做到最好.
不是問可不可以.
可不可以 給不給.
而是不斷地去突破我們的極限.
我們祈禱.
在我們一班錦村家姐妹.
我們一起來去呼求你的幫助.
我們知道我們是一個罪人.
我們很多時候都是有限.
我們都有自己的腦我和本性.
但求主讓我們在有限的裡面.
不斷地去追求你的義.
不斷地尋求更好的行為.
不是因為能夠做到什麼救恩的換取.
而是我們單單地願意回應你.
幫助我們等姐妹的生命裡.
我們可能有各自的不同的困難.
很多不同的功課來我們學習.
讓我們能夠超過我們所能做到的.
能夠突破我們自己的界限.
來做更加好.

$^{601}$一個合乎更加你心意的基督徒.
你的兒女.
幫助我們.
當我們去披戴基督的時候.
當我們願意去學像你一樣.
去讓人看到你的樣式的時候.
給我們力量去做到最好.
但我們知道我們不是一個超人.
我們是一個軟弱的人.
但我們願意不斷地去突破自己.
一生地跟從你.
奉主名求.
阿們.
\newpage



\section{加拉太書 6:11-16}
\label{sec:FRK_krDh8hE}
\textbf{2023-04-30 早堂《怎麼誇口?》加拉太書 6\_11-16 講員:楊明麗姑娘}
\newline
\newline
連結: \href{https://youtube.com/watch?v=FRK_krDh8hE}{\texttt{https://youtube.com/watch?v=FRK\_krDh8hE}} ~~~~ 語音日期: 2023-05-02
\newline
\newline
\hyperref[sec:sHtcf0yp8MI]{\small{< < < PREV SERMON < < <}}
~
\hyperref[sec:index_chronic]{\small{[返順時目]}}
~
\hyperref[sec:index_scriptual]{\small{[返順卷目]}}
~
\hyperref[sec:AGfClw0jrYY]{\small{> > > NEXT SERMON > > >}}
\newline
\newline
加拉太書 6:11-16
\newline
\begin{longtable}{cl}
\hline
\hline
章節 & 經文 (和合本修訂版)\\
\hline
6:11 & \begin{tabularx}{0.7\textwidth}{X} 你們看我親手寫給你們的字是何等的大！ \end{tabularx} \\ \\ \relax
6:12 & \begin{tabularx}{0.7\textwidth}{X} 那些想要炫耀外表的人才勉強你們受割禮，無非是怕自己為基督的十字架受迫害。 \end{tabularx} \\ \\ \relax
6:13 & \begin{tabularx}{0.7\textwidth}{X} 他們那些受割禮的，連自己也不守律法；他們要你們受割禮，不過是要拿你們的肉體誇口。 \end{tabularx} \\ \\ \relax
6:14 & \begin{tabularx}{0.7\textwidth}{X} 但我絕不以別的誇口，只誇我們主耶穌基督的十字架；因這十字架，就我而論，世界已經釘在十字架上；就世界而論，我已經釘在十字架上。 \end{tabularx} \\ \\ \relax
6:15 & \begin{tabularx}{0.7\textwidth}{X} 受割禮或不受割禮都無關緊要，要緊的就是作新造的人。 \end{tabularx} \\ \\ \relax
6:16 & \begin{tabularx}{0.7\textwidth}{X} 凡照這準則行的人，願平安憐憫加給他們，和神的以色列民。 \end{tabularx} \\ \\
[1ex]
\hline
\hline
\end{longtable}
$^{1}$各位電影姐妹早安.
你好嗎.
好 感恩.
那為什麼我這樣問呢.
我先自首一下.
我沒有.
我不是患新冠.
不過咳得很厲害.
應該是氣管敏感.
那再說幾句我就會咳了.
那如果一會兒我咳了.
大家不用擔心了.
都是這個緣故的.
所以我還是戴著口罩比較好.
免得大家覺得有點害怕.
第二樣我覺得很需要說一句.
其實是我們對繼光弟兄不主.
我們應該早一點跟繼光弟兄說清楚.
我們的程序是怎樣.
免得一上來的時候.
因為繼光弟兄主要應該是監流聚會.
(繼次弟兄).
對不起 繼次弟兄.
所以就有一點點.
跟我們這裡的程序有一點點不同.
大家不用介意.
我們大家都是宣道會錦繡堂的電影姐妹.
我們以後會學習做好一點.
我們一起祈禱.
親愛的天父上帝.
我們來到你的面前.
最後我們很想你賜給我們恩典.
賜給我們智慧.
讓我們能夠明白你的話語.
還有賜給我們的恩典和看顧.
讓我們能夠明白你的話語之後.
我們能夠學習去做.
求主幫助.
我們祈禱仰望你.
奉主耶穌基督的聖名.

$^{41}$阿們.
今天我們說的題目.
怎樣去誇口呢.
剛才讀了一段經文.
保羅說到一些事情.
有些事情就是.
他覺得一些法利賽人.
做了一些不適當的誇口.
我想問一下.
我們會不會誇口.
有沒有誇口.
你會為什麼事情去誇口呢.
其實有些事情可能大家都會覺得是很好的.
我讀了整本聖經《新舊約》.
或者我讀了《新舊約》聖經多少次了.
我很記得每一首詩歌.
第幾首呢.
《生命聖詩》的第四首就是什麼了.
或者我們有可能誇口其他的事情.
會不會誇口.
我們的學歷呢.
還是誇口.
我雖然沒有什麼學歷.
但是我做得很好的.
我可以照顧得一家大小很好.
供給子女去不同的地方.
讀他們想讀的.
我給他們很大的自由.
讓他們選擇他們想讀的科目.
還有一些事情我們有機會去誇口.
這些事情不一定是房子.
不一定是錢.
但是總是一些你覺得很重要的.
以前會說.
哎呀 我的兒子真的很蠢的.
這個可能是四五六十年前的事情了.
就一定是說自己的兒女很蠢的.
其實某程度上都是想誇讚一下.
現在不用了.
現在真的會讚賞一下自己的兒女.

$^{81}$有時候也是好的.
我們之所以誇獎我們身邊的人.
無論是誇獎父母.
或者是誇獎兒女.
都是因為我們重視.
重視的我們會去誇獎.
保羅這次就是說.
法理賽人.
他們有些什麼誇獎了.
是不合義的呢.
保羅在要跟他們說的時候.
他特別提了一件事.
他說這段是我自己寫給你們的.
你看我寫的字多大.
為什麼要特別寫大了.
還告訴他們.
我寫這麼大的字.
因為我自己親手寫的.
其實他不是想說.
你看我的筆跡多漂亮.
我寫的東西是很不同的.
我特別愛你們.
所以加拿大人.
我特別自己寫幾個字在書.
而他是想告訴他們.
我親手寫給你們.
我接下來要說的話是很重要的.
所以我告訴你們.
是我親手寫的.
所以字這麼大.
因為他的眼睛不是很好.
他的眼睛有點傷.
所以他就要寫得大字一點.
看得清楚一點.
他特別告訴加拿大人.
我到現在雖然好像要收筆了.
說到最後了.
但我最後有些很重要的東西.
一定要告訴你們.
告訴你們什麼呢.

$^{121}$有些人是炫耀外表.
炫耀外表就是叫你們一定要受割禮.
無非是怕自己為了基督十字架的緣故而受迫害.
原來是受割禮.
是猶太人一出生第八天就要做的事情.
但這班加拿大人他們不是猶太人.
他們是當時的外邦人.
或者好像我們一樣.
我想我們在座應該主要都不是猶太人.
應該沒有哪一位是.
我們都是屬於外邦人.
他就說到其實外邦人是不需要受割禮的.
信主受洗就足夠了.
但是就給一班法利賽人.
就要他們去受割禮.
目的就是要炫耀.
炫耀什麼呢.
他說這班人其實他們自己都不受割禮.
他們自己都不是.
他們自己受割禮.
但是這班法利賽人他們受割禮.
他們都不是整個律法他們都守的.
但他要你們去受割禮.
是要在你們的身體上去誇口.
保羅說到他們是要去誇口.
要利用這些不是猶太人.
信了主之後去受割禮這件事來誇口.
誇口什麼呢.
他在表示我們叫這班不是猶太人的基督徒.
他們都要受割禮.
表明他們都很願意好像猶太人一起生活.
但這不是一個信主之後需要的條件.
在《使徒行傳》第十五章.
在耶路撒冷的大會當中.
特別講到不是猶太人.
外邦人他信了主之後什麼要做呢.
當中都講到他們不需要守律法.
受洗就可以了.
不需要行割禮.
不需要戒除什麼食物是潔淨才能吃.

$^{161}$什麼食物是不潔淨不能吃.
碰到什麼東西是不潔淨要怎樣.
不需要守這些的律法.
不需要為自己做過錯事.
有過犯而要再去獻祭.
在《使徒行傳》第十五章已經講了很多.
去到這個時候.
加拉太的時候已經是《使徒行傳》第十五章之後.
一段是超過十年之後的事.
在這個時候.
當保羅再跟他們講的時候.
他們法利賽人要加拉太人為什麼要去受割禮呢.
他要告訴人們我們傳福音.
不單叫他們做基督徒還要叫他們做猶太人.
這樣來誇口.
還有一樣就是.
我們這次去.
加拉太人當中有357人受了割禮.
他在數那個數字.
357是我加進去的.
我的意思是.
法利賽人他要炫耀一個數目.
我們洗到多少人受了割禮.
這個都是一個炫耀.
一個誇口.
保羅告訴他們我們不是按著這些來誇口.
最重要的是什麼呢.
最重要的是我們信了主受了洗之後.
要做一個新做的人.
新的生命要在我們裡面去成長.
我們要活出一個新的樣式.
這個才是重點來的.
誇口.
割禮並不是一個我們需要去誇口的重點.
法利賽人有什麼事呢.
他們不是很守全律法呢.
法利賽人都有不守律法的時候的.
他們有什麼特別不守律法的時候呢.
其實人總有些地方做錯的.
所以他們一直都有繼續憲制.

$^{201}$去認罪要憲制.
總有些地方他們會錯的.
特別是還有一些人.
他們會在星期六的時候.
做了一些不應該做的事.
按著他們的律法.
星期六很多事不能做的.
他們都還有人做的.
另一樣就是他們利用了割禮板.
來減少對父母的供養.
但是十屆裡面.
有一屆是沒有.
第五屆告訴我們什麼呢.
我們要孝敬父母.
我們需要供養父母 照顧父母.
他們利用了一件事.
那件事就是奉獻給神的錢.
就是奉獻給神的.
不可以奉獻給神.
然後你又在那裡拿一些出來.
於是他們設立了一個叫做割禮板.
這份錢我就放在這個割禮板.
意思是我要拿來奉獻的錢當中.
雖然我還沒有拿去聖殿奉獻.
還沒有拿去會堂奉獻.
但是我這一份的錢是要來奉獻的.
所以爸爸媽媽.
你們不可以在這裡拿錢去花.
但是又沒有給別的.
他要來供養父母的錢給父母.
他將要來供養給父母的錢.
他就放在裡面 叫做他要去奉獻的錢.
其實這份錢他還沒有拿去奉獻.
不過他告訴父母.
我將你的錢拿去奉獻了.
於是你就想一想.
我預告了.
這不是一個很好的榜樣.
但是以法利賽人他們很熟悉聖經.
又很教導人要守律法來說.

$^{241}$這是非常不好的榜樣.
反而他們還要去叫加拉太人要受國禮.
而以這個來誇口的時候.
保羅就告訴加拉太人.
其實他們都不守全律法的.
他們都有很多其他的問題.
就為這個的緣故的時候.
他們都不能夠要求你怎樣去守律法.
那麼保羅是不是因為他自己不守律法.
按照律法他沒有東西可以誇口.
所以他這樣說呢.
他說其實不是.
他說如果要誇口.
我都不是沒有東西可以誇口的.
不過我只是誇主耶穌基督的十字架.
可能他真的按照律法看.
他沒有東西可以誇口.
所以只有誇主耶穌基督的十字架.
其實在《使徒行傳》的第22章.
保羅他都說了.
他說什麼呢.
他說自己是猶太人.
意思是說.
他自己是會說希伯來話的希伯來人.
他是最純正血統.
又能夠說自己民族語言的猶太人.
是很有得誇口的.
當時都不是那麼多猶太人.
懂得說希伯來話的了.
因為呢.
不是當時流行的語言來的.
弟兄姊妹 你們懂不懂說自己鄉下的語言.
我想不多的.
特別是居住在香港的外省人.
如果在香港出生.
在香港成長.
在香港受教育.
或者在外國受英語教育的時候.
對自己鄉下的語言其實不熟的.
就算在廣東省.

$^{281}$在不同的縣市.
可能你的語言有些不同的.
例如說潮州話就有些分別.
真是能夠說流利潮州話的香港人.
出生的香港人其實很少的.
當時代猶太人.
在當時的羅馬政府的統治底下.
出生的人.
他們學的語言.
都會是當時流行的亞蘭語.
或者是文字上再高深教育的.
學會了當時的希伯來文.
對於自己本土的語言.
已經是沒有什麼人認識的.
保羅其實是會說希伯來話的希伯來人.
他自小不單有受割禮.
有被在聖經上的昏陶.
而且他還是一個很專.
很嚴謹的去守這個律法的一個法理賽派別當中.
是來受教育的一個.
他當時應該都是法理賽人來的.
他是很受到嚴謹的教育的.
那個中派的派別應該是.
讓我找一找 對不起.
是加瑪列的門下.
加瑪列是當時的法理賽的名校來的.
是數一數二的.
而且他們守律法.
為什麼是名校呢.
因為他們守律法是非常非常嚴謹的.
保羅他怎樣嚴謹守律法呢.
他嚴謹守律法到他按著那個律法.
去向羅馬的政府申請這個.
是要鎖基督徒的.
因為他覺得基督徒就是.
他當時覺得基督徒就是一個叛變了.
不是信耶和華的神的人.
所以他要這樣去做.
其實他後來就是在那次的過程當中認識主耶穌.
就信了主耶穌.

$^{321}$領受了十格的救恩.
所以他再去看的時候.
其實他沒有什麼可以誇口的.
拿律法來誇口.
誇什麼呢.
是我們對神的認識少了一點.
應該要更加明白基督的救恩.
其實他是一個.
他是猶太人.
在基里加的大數道成長的時候.
在加瑪列當中他是一個很嚴謹的宗派的受教.
另外還有一樣東西.
他是不是.
他沒有什麼可以誇口了.
除此之外.
所以他就誇十字架.
如果我們看新約我們現在知道.
他後來是加上.
他應該寫了十三卷新約的聖經.
但是當年代不知道.
不過他做了一些事情.
他去了以弗所傳福音.
去了哥林多傳福音.
他用了很長時間在以弗所和哥林多那裡.
去培養他們.
幫助他們認識神.
為他們做栽培的工作.
建立他們成為門徒.
還在不同的地方.
他栽培了教會的領袖.
好像當時的亞居拉,伯基拉.
提摩泰,提多.
這些我們比較熟悉的人.
西拉也和他一段很長的時間.
一起去宣教.
路加也是他建立的.
他帶著他信主認識神.
帶著他出去傳福音.
其實他要誇口的東西多的是.
是不是.

$^{361}$他要誇口他有多少人聽過他講道呢.
真的多的是.
他要誇口他講過多少個鐘頭道呢.
真的計算不到.
但是他就說.
其實我沒有什麼可以誇口.
我要誇什麼呢.
我真的要誇口的時候.
我是要誇主的十字架.
為什麼誇主的十字架呢.
他說我誇主的十字架.
因為主的十字架在我面前的時候.
世界就好像被釘在十字架上面.
而對於世界來說.
我就好像被釘在十字架上面.
因為十字架在他和世界的中間.
他的意思是表達.
世界上面的隱憂.
我們每個人都會面對過的.
最好早上晚一點起床.
最好工作不用那麼辛苦.
但是可以加薪.
最好假期長一點.
現在最適合去旅行.
到處都開關了.
但是.
保羅他想的不是這些.
他說對著十字架的時候.
我去看回這個世界.
世界對我的隱憂.
是當它是死了的.
我和這個世界的關係已經死了.
因為它被釘在十字架上.
對我不成隱憂.
反而我的生命是要在基督的恩典當中去成長.
在世界來說.
它看著我的時候我是死了的.
它應該知道它對我的誘惑是沒有作用的.
沒有什麼東西是吸引到我的.
也不能阻止我的工作.

$^{401}$這是保羅他自己的看法.
他覺得他就要這樣去誇口.
世界對他的誘惑.
是對他沒有作用的.
因為他和世界已經沒有了關係.
和一個沒有關係的地方.
它的誘惑對我們有什麼用呢.
對於這個世界來說.
再沒有什麼東西可以引誘它.
最好這個世界不要再來引誘它.
最好這個世界不要再來引誘我們.
我們能夠朝著這個方向學習是很好的.
雖然是有些挑戰.
不過是很值得去努力的.
讓主耶穌基督的十字架.
在我們和世界的中間.
讓世界再不能夠引誘我們.
而我們只能看到主基督的十字架.
這個就是他所要誇口的.
他說我已經被釘在十字架了.
在世界看起來.
於是乎是沒有什麼東西.
它能夠再來影響到我.
受割禮和不受割禮.
最重要是什麼呢.
他說是不重要的.
最重要的是做新做的人.
凡照著這個準則的人.
願平安憐憫加給他們.
和神的以色列民.
原來保羅他講這麼多.
他不是要將以色列人猶太人.
和非以色列人的外邦人間開.
他說願平安和憐憫.
都臨到兩者的.
大家都得到神的恩典.
不過大家都要清楚一件事.
不要用割禮來逼迫對方.
不要用律法去控制著對方.
反而要是.

$^{441}$大家知道什麼最重要.
無論是猶太人或不是猶太人.
最重要原來是信主之後.
要做一個新做的人.
什麼是一個新做的人呢.
原來是加拉太書一開頭的時候.
它就說.
信主之後其實我們就是一個自由人.
自由的意思是.
不再受律法去控制.
信主受洗就夠了.
我們就可以得到主的拯救.
我們的生命要在這個自由的.
信仰的土地當中去成長.
而是律法不能夠控制我們.
反而我們要學習一樣東西.
不要讓世界去控制我們的時候.
世界被釘在十字架上的時候.
它不要令到我們失去自由.
沒錯 原來.
令我們失去自由的未必是律法.
反而是很多其他的東西.
它提過有些是迂肘.
有一些不合適的生活的習慣.
對我們自己身體有害.
也都可能傷害我們身邊的人.
有時是言語上的不恰當.
有時我們不說出來.
我們心裡面的憤怒 嫉妒.
這些都會的.
反而它教導我們.
怎樣可以得到自由呢.
我們要得到自由.
我們就不要在這一方面.
被世界去吸引我們.
或者被我們自己去順從了.
讓我們自己的自由去順從了.
這個社會的制度或者是價值.
或者是社會的人是這樣說.
可以這樣去處事的.

$^{481}$而是要以我們有的自由.
去馴服聖靈.
讓聖靈來帶領 感動著我們.
大家這段經文一定很記得的.
聖靈的果子 人愛喜樂和平.
忍耐恩慈良善信實溫柔節制.
在哪裡呢.
在《迦拉太書》五章二十二二十三節.
沒錯.
《迦拉太書》一開始就說到.
我們信主之後得到自由.
得到自由之後.
就不要放縱私慾.
不要使我們的思想或者生活習慣.
因為自由而沒有受到控制.
反而我們要馴服聖靈.
對我們的帶領和提醒.
讓我們凡事上都能夠結出美好的果子.
做一個新做的人.
它不只是在《迦拉太書》這樣說的.
它在《爾弗所書》都提醒信徒.
我們信主之後.
我們要有新的生命.
好像把舊人脫下了一件衣服.
穿回一件新的衣服.
我們有一個新的角色.
我們是神的兒女.
做了神的兒女之後.
弟兄姊妹要怎樣相處.
家裡當中父母我們要怎樣對待.
子女怎樣對待.
夫妻怎樣相處.
和傭人怎樣相處.
很多的教導都很確實地.
來告訴《爾弗所》.
告訴哥羅西.
其實它沒有特別只是叫迦拉太的人是這樣.
而是提醒我們.
信主之後我們的生命要有自由地成長.
不是以律法去誇口.

$^{521}$而是我們以十格的恩典去誇口.
做一個新做的人.
我們會誇口 是不是.
一定會誇口的.
剛才問過大家你會誇口些什麼呢.
可能真的有時候要想一下我們會誇口什麼.
我想我沒有什麼可以誇口的.
我沒有什麼是特別很厲害的.
應該是今年開始有一樣東西我做得不好的.
就是我很想我先生語言上學習多一些.
我經常去糾正他的語音.
因為我們是不同省份成長的.
我們會有不同的語言.
我特別喜歡去糾正他的普通話.
其實我自己都沒有去考牌.
我都不是說得很好.
不過總有些字我們會知道要怎樣讀.
好像有些特別的字.
瓊瑤 是讀瓊瑤的 瓊瑤.
其實這樣都是不好的.
雖然我沒有特別告訴他這樣多好.
其實都是間接地在他面前去誇口.
他是不是一定會很開心呢.
雖然他都很開心跟我聊天.
但是這些事我都要小心去做.
請問有沒有一些事你覺得自己會誇口的呢.
但總有些事你是厲害的.
總有些事是你的強項的.
甚至不少的.
有些人犯犯佳精的.
讓我們將這些事放在我們的自由當中.
去事奉我們的神.
讓祂成為我們向十格低頭.
向世界去誇口十格的一個媒介 好不好.
我們一起禱告.
上天 父上 神 我們來到你的面前.
我們要多謝你.
多謝你看顧我們.
主人 我們本來就是罪人.
我們要誇的話.

$^{561}$可能應該誇我們做了多少錯事.
但是主你有恩典 上苑同在看顧我們.
讓我們能夠有十格成為我們的誇口.
我們是有永生的人.
我們可以將眼界放得遠一點.
主要是我們放向永生.
放向天國的恩典.
放向你對我們的愛.
主要是求你幫助我們.
求你次次威給我們.
讓我們時時看到上帝你的預備.
你的教導.
讓我們時時有你的提醒.
領受到你的感動.
讓我們能夠學習.
在生活當中我們多一份謙和.
但是時時高舉十格.
求主幫助我們.
主人 我們真的有時會有我們的喜好.
而至於在運用自由上面.
我們未懂得以十格為我們的中心.
求主憐憫.
常常提醒我們.
不要厭棄我們.
我們仰望你.
奉教主耶穌基督的聖名.
阿們.
\newpage



\section{}
\label{sec:AGfClw0jrYY}
\textbf{2023-05-04 周四晚查經班《被擄歸回歷史研讀(13)》 講員:曾裔貴牧師}
\newline
\newline
連結: \href{https://youtube.com/watch?v=AGfClw0jrYY}{\texttt{https://youtube.com/watch?v=AGfClw0jrYY}} ~~~~ 語音日期: 2023-05-06
\newline
\newline
\hyperref[sec:FRK_krDh8hE]{\small{< < < PREV SERMON < < <}}
~
\hyperref[sec:index_chronic]{\small{[返順時目]}}
~
\hyperref[sec:index_scriptual]{\small{[返順卷目]}}
~
\hyperref[sec:eLIEUCC43fA]{\small{> > > NEXT SERMON > > >}}
\newline
\newline
$^{1}$頂子妹晚安,我們開始今晚的茶經班,一起禱告先.
多謝主,我們暫時停頓了一晚.
今晚又可以恢復茶經班.
尼希米記第五章.
盼望神都給我們透過尼希米這六章裡的自述.
他的個人的一些回憶,讓我們看到要建立一個很大的工程.
其中需要克服的,尤其是作為領袖要面對要克服的.
好像裡面都有一些大的原則.
必然會發生可能不同的方式.
但是必然的是要面對,是要克服,是要去承擔的.
求主都讓我們能夠從其中找到一些很重要的線索,方向和出路.
使我們在神的家裡面去服事的時候.
一方面都是小心翼翼,真的要來到去.
在很多重大的事情都是要謹慎,求主帶領我們.
今天晚上我們等候,亦都知道近來多了,無論是弟兄姊妹還是香港.
多了確診的宗數,盼望主都保守我們大家的身體健康.
雖然不是那麼嚴重,但是盼望不影響大家的生活工作.
我們一起在你面前等候禱告,奉耶穌基督的名,阿們.
回過頭來補充一點,開始這個重建城牆的工程.
一方面固然是那個硬件在耶路撒冷的破落,受敵.
一直都隨隨便便,任意的去欺凌.
那個欺凌可以透過大家很投契,星期六安息日.
大家一起去市集買東西,看一些新穎的東西.
安息日大家都不再著重,這個不等於是逼迫.
但是在屬靈的角度,這個就是一個靈性低落的原因.
因為猶太人安息日就不打算回教會去會堂.
這個就變成我們都要舉一反三,就是我們是不是一定要長假期.
我們就一定要去長的旅行去放電.
一定要逼去日本,我們是不是可以有其他的想法.
這些都是我們要去認真想的,每個弟兄姊妹都要明白的.
你起碼看到自己有這個托付,有個人擺上自己的承擔.
沒有必然它是要做的,因為它只不過是一個領袖.
一個你可以說世俗的酒精,但是它就很有自己內心的感動,負擔.
將自己放在那個崗位,你就知道那個崗位衍生的壓力就自然是倍增的了.
一方面它要從老遠的河東來到去河西耶路撒冷.
要對當地的省長來尋求協助.
固然是有亞達西西王的御靈,但是山高皇帝遠.
再加上那些本地的外邦人,即是已經盤踞在耶路撒冷的外邦人.
就一而再出現,首先就是第四章,是冷嘲熱諷.
希望動搖百姓的童心,挑剔,嘲笑.

$^{41}$令到尼熙米在動員的時候都有一定的難處.
就是要梳理百姓的心靈,會不會被潑冷水之後.
就會低落那個情緒,那個事奉.
這個是領袖必須要有這個考量.
第四章叫做勉強可以避開了那個外邦人.
尤其是以幾位領袖,參巴拉,多比亞,阿拉伯人,基善.
幾個都是外邦人,外族人,但是和猶太人相當密切的.
遲些我們就會知道有其中一個,他娶了祭司的女兒.
這些都是我們要留意的.
到第五章,今天我們講的第五章.
就出現了一個內部,因為由來已久的.
就算尼熙米沒有上來,都已經有的問題.
不過就透過現在這麼繁重的建造工程,一下子就凸顯出來.
也可以這樣說,一下子就因為這個新的領袖,一下子就爆出來了.
就是內部原來很多的猶太人,過去因為種種原因.
因為農業社會不是每年都是豐收的.
有失收,有荒年的時候,他們差不多窮到要賣了自己的田地.
甚至連自己都賣了給那些外邦人,或者給官府.
所以很多的人就會借這些機會.
有錢的,有土地的,有資產的,只會越來越富有.
譬如你在Yohamitown有一個單位,你在香港哪個地方又有一個單位.
隨著香港的樓價通貨膨脹,富者就越富.
貧者就越貧,如果你不小心十年前賣了一層樓.
是97的時候賣了,最低價的時候.
你以後想買一層樓,隨著樓價不斷狂升,你的工資根本追不上的時候.
你就會發現這種壓力就會帶來很嚴重的問題.
剛才提到,猶太人賣了田產,而且還要給很高的利息.
所以不斷的還利息都已經解決不了他們日常的溫飽.
所以藉著現在這個建造工程一下子就爆出來.
他們可能是遷怒於利希米,但更加希望利希米秉公的去處理.
那怎樣去拆解,如果發現不只是外邦人,而是猶太的領袖,猶太的貴袖,即是有錢人.
那怎麼辦呢?這是很棘手的倫理問題.
譬如我們的教會有很富有的弟兄姊妹,很多單位的,租給我們自己的弟兄姊妹.
很貴租,不交租就沒情講,逼你走,逼遷,加很貴租.
這些真的變成了倫理問題,這次租客又向教會的目的投訴,來尋求協助.
諸如此類,這些很現實的問題,如果發生在內部就相當困難了.
第五章,百姓和他們的妻,大大的呼號,埋怨他們的弟兄猶太人.
有的說,例如我們和兒女人口眾多,要去得糧食度命,度命即是說找朝不得晚.
差不多連基本的糊口都不夠,不要說要吃好的和牛.
有的說,我們點了田地,即是說另外一些人說,即是說另外一個case.

$^{81}$我們點了田地,即是說點檔,葡萄園房屋要得糧食充飢.
到底這些投訴,到底是單方面還是事實的全部呢?.
領袖首先聽了投訴,不要聽一面之詞,都要來秉公的.
事實是怎樣呢?這個就是要領袖,首先不要這麼快動怒.
有沒有搞錯?這樣搞到百姓,原來去到這樣的田地.
哪有心,哪有力去做公職呢?為神的聖殿範圍來重建城牆呢?.
所以不要這麼快動怒,先聽完和調查故事的雙方面.
是不是真的有這樣的無良領袖,貴袖和有錢人.
這樣來對自己的同胞可能加重他們的壓力.
例如加很重的稅,你想借錢?沒問題,我們照一般的利息再加一點.
即是說更加是在別人的困難上面,百上加斤.
這些都是我們必須要去反思的.
第四節,有的說,再有另外的case,有的說.
我們已經紙著田地葡萄園,借了錢給黃立稅,原來甚至乎要交稅.
你有沒有錢是你的事,你有沒有收成也是你的事.
黃的田地就要交稅了,即是說你失業,但你有層樓.
麻煩你也要給差餉,地租,地稅,諸如此類,可能有其他的稅項.
這裡就告訴我們,我們的身體與我們弟兄的身體一樣.
我們的兒女與他們的兒女一般,現在我們將要使到兒女作人的瀑皮.
我們的女兒已經有圍皮的,我們並無力拯救.
因為我們的田地葡萄園已經歸了別人.
不單止賣田賣地,現在要賣了子女.
他們的訴求就是說,大家都是猶太人同胞.
為什麼因為他有資產,他就越來越有錢.
我們因為沒有資產,而天災或者失收,導致我們越來越窮.
貧者就越貧,富者就越富.
大家都是神的兒女,猶太人,可能同支派就更加嚴重.
為什麼會變成這樣呢?有沒有天理呢?.
猶太人就說,大家都是神的兒女,為什麼我們會變成這樣呢?.
已經有我們的女兒,做了別人的妹仔.
我們都沒有力去拯救,就是說繼續要賣.
那怎麼辦呢?這樣的社會結構性的不公.
可能不是一個個人的賭輸錢,賣了房子.
可能不是,是一個結構性的一種塵壓式的欺凌.
商業上的一個不健全的體制,那怎麼辦呢?.
當然我們要記住,聖經不是教我們理財.
不是教我們在投資上要怎樣有眼光等等.
不是這樣的,但是反而聖經就給我們.
透過這些很現實的社會的實際倫理處境.
看看神的兒女怎樣用愛超越那個社會的標準.

$^{121}$來讓那些貧富懸殊,本來必定會出現的.
可以減到傷害最低,低到一個地步.
可以甚至給到那個本來是書家,他已經賣了田地.
要賣掉自己的兒女,怎樣能夠有一個.
在人可以製造的制度底下.
可以令到他有一個改變的所謂受人耳語.
給他有一個契機,可以扭轉這個可能永遠都改變不了的局面呢?.
第六節,我聽見他們呼號說這些話便甚發怒.
有黑勒,尼欺米,這個民怨不是說笑.
因為幾個向導,涉及的人,涉及的是田產.
涉及的是到底幕後是不是單純是失收這麼簡單呢?.
是不是有計劃的,有一些背後的,好像集團式的.
來做到這些窮人永遠都窮,繼續更加貧窮呢?.
那就是說,失收只不過是表面的原因.
如果這個表面原因,人為是操控不了的.
但是一些人為性的操控的,其實如果那個社會.
特別這個是猶太人的社會,有摩西律法的社會呢.
能不能夠回到摩西律法裡面,尋求到一些智慧.
可以解決這個結構性的貧富懸殊呢?.
所以尼欺米這個時候,他要面對一個很大的炸彈.
他要處理民間的民怨,怎樣能夠釋放呢?.
現在不是外邦人嘲笑,第四章就是集中外邦人.
對這個事公的批評,論斷,嘲諷,攻擊,甚至恐嚇.
要外邦人,要猶太人停工,不要再做.
因為做了,就規範了你不是猶太人.
你就不可以進來我們耶路撒冷了.
你是外邦人,你就是外邦人了.
即是沒有染金,雞籠呢,就是以前外邦人希望的局面.
但是現在不行了,所以他們一定會反抗的.
所以就見到尼欺米反而能夠鼓動百姓齊心.
去擁抱自己這個身份,就是我們是神的子民.
我們才配在耶路撒冷來居住,敬拜的.
外邦人他們可以透過我們去敬拜神.
但是外邦人不可以隨便進來聖殿的.
這個是很清晰的,百姓也能夠明白.
也有領袖現在鼓勵他們,他們都不介意學習這個身份.
但是涉及到現在這個是一個根本性.
尼欺米還沒有在的時候就已經出現了.
就不是這一個多月,因為他們建了五十多天那個城牆.
不是這一個多月的時間才有這些問題.

$^{161}$是過去年年失修,或者耶路撒冷可能有一些官.
特別欺壓猶太人來討好其他的外邦人.
導致會有這些官商的勾結.
導致那個田地不知不覺就被變賣.
歸入了那些官或者貴族的手.
所以尼欺米不到他不處理.
因為現在那個工程還沒有完.
你就要面對現在幫你重建的猶太人.
他們那個民怨,他們在這個時候引爆這個炸彈.
到底是陰謀還是積壓了太久.
必須要透過這個新的領袖尼欺米來處理呢?.
尼欺米聽到這幾個情況.
一來他身處在河東非常遙遠的波斯的舒山城.
他一直都在黃的俸祿底下生活.
還有作官.
對他來說耶路撒冷是他的家鄉.
不是他現在的原居地.
現在居住的地方.
所以他現在要面對這個問題.
首先就是要了解背後有沒有結構性的制度上的一個剝削.
一個是很容易令到窮人陷入一個危機.
類似大家認識的信用卡.
你稍為一次忘記了交你用了的錢.
跟那個月結單來交錢.
你那個所謂的minimum pay.
你以為由你付的那一下開始.
就馬上追加你的利息.
不斷地double上.
所以是一個陷阱.
我很多年前試過一次.
無意中.
挺有趣的.
用信用卡當提款卡.
按進去櫃員機.
然後又試試玩玩.
又可以原來套.
所謂信用卡套現.
試一下按了.
一按完之後才發現不行.
原來是馬上要追加不知道20\%多的利息.

$^{201}$還有你的本金都已經要給一個charge.
馬上拿回銀行.
那裡已經收了我忘記了多少錢.
才發現原來.
對不起.
曾先生.
這個提款功能是不可以取消的.
你可以不要張卡.
但是你不可以.
這個信用卡提款的功能.
透支套現功能.
是不能夠取消的.
這是一個陷阱.
你按錯一個按鈕.
你就馬上當是借錢.
而且利息是相當高的.
這些就是一些.
你可以說是表面方便你去買東西.
可以來到去.
不需要付那麼多現金.
但實際上是帶來很多.
可能是青年人比較多.
來到去會落入消費的陷阱裡面.
一個消費的誘惑裡面.
胡亂的.
自己根本沒有錢的.
就強行來去洗了一些.
和自己未必有直接急需要用的物件.
這些是消費的誘惑.
接著就是消費的陷阱.
結構性的.
這些是我們真的很難避免的.
所以現在這裡.
人們就很生氣.
所以原來領袖.
都要有時有這些的生氣.
發怒.
發怒之後怎樣處理呢.
就是那個技巧.
和那個屬靈的情操.

$^{241}$首先就是你要承擔.
因為現在你必須要處理.
因為永遠都是那個軟件.
重要過那個硬件.
守了那個城牆又如何呢.
百姓在裡面原來個個都賣田賣地.
完全自己是為人打工.
自己的兒女原來是賣了做瀑皮.
還要那個主人.
是自己的猶太的同族的弟兄.
這樣就很諷刺.
原來大家都不去守摩西的律法書.
原來摩西的律法書一早就有教導.
有講的.
那就變成了名存實亡.
建好城牆又如何呢.
圍住猶太人又如何呢.
猶太人.
但不是真的猶太人.
就是心靈裡面.
大家都沒有真正看自己是神的子民的.
這個同胞的身份.
和大家都很敬重神的律法.
這個就是軟實力的問題.
是軟件的問題.
是百姓的心靈的服不服的問題.
可能起初表面的原因就是.
誰叫你失修啊.
誰叫你的田地沒有生產啊.
我的田地很多生產.
我爸爸投資眼光獨到.
很懂得找哪些地方.
那些收成很高.
慢慢一代一代.
接著一個結構就是.
你賣東西給我們.
田地我們收購的.
不過價值你的田地都沒有生產.
很低價的.
那就很大件事了.

$^{281}$那就變成了.
真的結構性的.
對自己同胞的剝削.
首先就是那些田地不能永賣.
去到稀年.
要無條件給業主.
即是賣給你的弟兄姊妹.
神有這樣的結構.
不會永遠都是一個支派.
去贏的.
你可以想像到.
尼熙米的發怒.
如果這個發怒只是.
一種脾氣.
而沒有真理.
沒有其他的應對.
這個發怒只會帶來一種破壞性.
例如鼓動這些貧窮的人去革命.
鼓動這些人去鬥貴族地主.
用武力去改變那個結構.
在世界的歷史不斷出現這些問題.
你看尼熙米怎樣能夠.
又不是刻意去做一個社會改革家.
他的責任是.
現在出現了一個這麼明顯的.
內部同胞的問題.
在神的聖經的律法書摩西五經裡面.
因為當時他們只是有摩西五經.
能不能夠在摩西五經裡面.
找到一些原則.
來去很簡單地.
只要大家將心比己.
這樣來去行事.
其實可以很簡單.
第七節.
我心裡籌劃.
斥責貴族和官長.
相信這個是經過研究.
經過分析.
不是聽一面之詞.

$^{321}$所以真的要學習.
不要聽一面之詞.
有些童工我真的很欣賞.
他聽完某些東西.
他自己.
有時和我分享.
他說我不會只聽他一面的說話.
會聽對方的.
我覺得是很有智慧.
我們有時作為領袖.
都要做雙方面的研判.
如果我們一面都傾向某一個單位.
有時就很容易先入為主.
甚至我們會被他誤導.
我們那個決定.
有時就很災難性.
對方可能不是這樣.
有些人是加鹽加醋.
想借我們出手.
這樣就很麻煩.
所以真的需要有一個客觀的研判.
可能未必是那些貴族全錯.
可能真的不是那些窮人所想像的.
欺壓你們.
所以一定要做一個很仔細的分析.
不過作者都長話短說.
斥責這些貴族和官長.
果然貧富懸殊.
富者越富.
一定是貴族官長.
用一些手段.
表面上好像對窮人很好.
但實際上長遠.
其實那些窮人就更窮.
因為他們變賣了那些資產.
第六節的下半節.
對不起第七節.
你們各人向弟兄取離.
哦,這裡是問題的癥結.
收離息.

$^{361}$相信是令到百常家庚.
出現了一個很大的問題.
律法書有教導.
借錢給弟兄.
窮乏的弟兄.
不可以收取離息.
和他那個.
唯一的衣服.
或者他真是僅有的.
日用的.
無論日用品或者他的衣服.
你不可以進他家裡.
來取.
他給你什麼你做抵押.
你就拿什麼.
有這樣的厚度.
這個是律法書摩西五經教導.
於是我招聚大會.
攻擊他們.
批鬥大會.
相信那個動機.
不是想批鬥.
不是想站在窮人這一邊.
去做批鬥大會.
領袖.
那個動員力.
領袖那個影響力.
領袖那個煽動性.
要很小心拿捏.
因為你一成為領袖.
其實你已經擁有權力.
你一擁有權力.
帶來的就是.
到底你是想你自己的野心.
還是你真是想.
那個用回聖殿.
耶路撒冷.
聖城重建.
這個最大的.
神的選民.

$^{401}$要有自己的.
被老歸回的重建.
重建最終是要重建聖文.
等候彌賽亞的降臨.
他們都大概計算.
可能真是四百多年後.
耶穌基督會降臨.
當然他們不知道叫耶穌基督.
他們知道叫彌賽亞.
所以這個是猶太人的心願.
雖然是幾百年後才發生.
但是總是有他們的心願.
神要他們被老歸回.
有自己的聖殿.
有自己的敬拜.
有自己神聖選民的身份.
為著的不是只是.
給他們一個地方住.
而是為著將來彌賽亞的降臨.
神的救恩的成就.
所以尼希米現在在這裡的批鬥.
一定不是出於他自己.
一種他現在對了.
哪些貴就曾經不是很滿意.
現在我借這些窮人來批鬥他.
一線之隔.
都是一個星土.
但是一線之隔.
還有你要用民意.
要用得很小心.
不要有一種內心很大的野心.
這些只能夠神監察的.
因為包裝其實是太美麗了.
我不是想告訴大家.
領袖都是黑暗.
不是這個意思.
而是一去到這些批鬥.
總是會群情洶湧的.
總是會人會有一種偏見.
還有會有先入為主.

$^{441}$這些一定要避免.
盡量不要帶來這種.
因為這些永遠都是重大議題.
藉著窮人的處境.
去爭取一些東西.
永遠都沒有人會反對.
唯一反對的就是.
貴就官長是不是一定是錯呢.
所以這裡告訴我們.
我對他們說.
我們盡力贖回我們弟兄.
就算賣與外邦的猶大人.
你們還要賣弟兄.
使得我們贖回來嗎.
這裡有一個point.
黎希米也很重要.
就是賣了給外邦人的.
同胞猶太人.
不同的支派.
應該盡力去幫助.
這個支派的窮乏的弟兄.
等到他被賣了的外邦的.
可能是瀑被子女.
都要贖回來.
有錢去贖回來.
使大家不會落在一個.
更加長遠.
惡性的結構的循環裡面.
這是民族的一個智慧.
一個童心.
讓這個民族.
大家都有這一份的尊嚴.
就是說出得起錢的.
就幫那些窮乏的.
好像嗎哪的分配那樣.
所以如果猶太人他們能夠.
這樣去看這些事情.
是真的會帶來.
大家有一個很強的凝聚.
在那個十二支派分地圖.

$^{481}$因為你是哪個支派.
分到什麼地圖呢.
你就不可以永遠賣的.
你只能夠賣給你同族的弟兄.
你不可以賣給另外的支派.
更加不可以賣給外邦人.
沒錯 雖然那個商業世界.
一路一路演變.
猶太人都不可以再獨善其身.
和他們根本不能夠.
保護自己的地圖.
要被擄.
甚至馬代.
波斯.
這個巴比倫.
都來到去.
將他們完全的攻陷了.
將他們瓦解了.
完全被擄到去外邦.
但是當他們.
當他們回來的時候.
這個時候就要回去.
是否能夠可以慢慢.
大家去.
在摩西的律法書裡面.
找到那個.
很重要的原則.
就是那些地不可以永遠賣.
那些地如果賣了給外邦人.
其他的支派.
要幫忙.
讓那個支派的弟兄.
可以有機會贖回那些地.
贖回那些兒女.
不要再永遠成為僕僻.
但現在不單止這樣.
連這些猶太人的貴就.
都來到去.
慫恿他們繼續去賣.
那就是說更加助紂為虐了.

$^{521}$更加就是說.
自己就越來越富有.
自己就越來越多田地.
田產.
那些支派基於.
不知道什麼天災.
或者什麼原因.
或者什麼小小的投資的失利.
種那批的東西.
沒有人買.
或者是特別有失收.
就導致經營不善而帶來的.
可能破產.
就賣了這些田地.
就對不起了.
是你投資失利不關我的事.
我沒有義務幫你了.
說這句話就已經犯了.
教導.
雖然不是一個命令.
但是如果大家都願意.
這樣去遵行的時候.
就很美麗了.
所以這裡李希米就這樣.
振振有詞的批鬥他們.
甚至是說一些很重的話.
是不是你們要加入.
去販賣弟兄的行列呢.
他們就靜默不語.
無話可答.
這個是當時.
可能群情.
但是似乎更加是.
他們都感受到.
自己其實一路內心.
都是做得不夠好的.
其實可以做好一點的.
為什麼不做呢.
就是.
具體一點.

$^{561}$譬如你真的有五六個單位.
你看到有.
譬如全套童工.
租得很貴的.
他都買不到樓了.
很象徵式的給童工住.
等於幫你看樓.
是的沒錯.
你收不到市值的租金.
是的的確是的.
但是你有沒有從另一個角度想.
我能夠有一個.
你想一下.
我都是沒有任何的數據.
沒有任何的人名.
是一個誇張的.
你有五套單位.
你給一套奉獻給神.
而且是有傳道人.
為你的家庭去祝福.
他真的感受到神的恩典.
用這麼低廉的.
收入.
我的租金.
就可以在這裡居住.
我可以專心事奉.
又不會逼遷.
這個業主是這麼好的.
主裡面弟兄姊妹.
甚至你不要告訴別人.
我是誰.
這樣的服侍.
你想一下是多麼的蒙福.
你有五套房.
五個單位.
你將一套奉獻給神.
你不是送給他.
是租給他.
但是租金是很廉宜.
是和他計算過的.

$^{601}$他真的能夠負擔得起.
他的家庭.
相對來說就輕省了.
教會又可以因為這樣.
又令到同工.
可以很專心服侍.
你想一下這樣的循環.
是帶來彼此的祝福.
黎希米就立即.
教導他們.
待會你會看到.
他們就真的不再收利息.
第九節.
我又說.
你們所行的不善.
你們行事不當敬畏.
我們的神.
不然難免我們的仇敵.
外邦人都毀破我們.
我和我的弟兄與僕人.
也將銀錢糧食.
借給百姓.
我們大家都當免去利息.
黎希米.
看到百姓的困難.
他自己.
和他那些上來的.
和他一起.
同心建造的.
未必是猶太人.
可能是他的屬下.
追隨黎希米的.
他們將自己的.
自己的俸祿.
自己的有的.
因為這個工程.
他來借貸.
甚至後來.
來給工人的工錢.
這樣的身體力行.

$^{641}$的付出.
導致那些官長.
更加無話可說.
這個敬畏神.
原來是有一個具體的行動.
在事實裡面.
黎希米強調.
可以借.
但不可以收利息.
甚至乎是.
分期的攤還.
借貸.
讓窮乏的弟兄.
慢慢慢慢.
有那個力量.
不會被那些.
只是利息壓扁了.
永遠.
那個本金還不了.
現在中美洲的國家.
南美洲的國家.
起碼不斷還利息.
導致國內.
大家就有權力.
有錢的更加貪污.
更加盡快拿了錢.
就走出國.
那個國家就永遠.
在一個萬劫不復裡面.
窮人就不斷在動盪裡面.
又不能夠讀到書.
又繼續社會動盪.
繼續製造很多很多的.
暴徒很多的.
人君子.
那個國家.
出現大量的妓女.
導致這些情況.
根本就不要再說.
那些生產.

$^{681}$發展.
教育.
醫療.
根本就沒有辦法了.
只會是罪惡.
有錢有權的人.
繼續壟斷整個國家的.
命脈.
他們可能.
最早期的時候.
那些不是太有智慧的領袖.
高度的借貸.
來解決國家的問題.
但是.
去到一個回不了頭的.
債務的危機.
阿根廷如是.
這麼多的.
天然的.
礦產或者是這麼多的物資.
為什麼.
國家都不可以有一個.
這樣的改變呢.
第十節.
對不起.
第十一節.
如今我勸你們.
將他們的田地葡萄園.
橄欖園房屋.
並向他們所取的.
銀錢糧食新酒.
和油.
百分之一的利息.
交給他們.
眾人說.
我們必歸還.
不再向他們索要.
必照你的話行.
我就照了祭司來.
叫眾人起誓.

$^{721}$必照著所應許的而行.
我也斗著空前的衣襟.
說凡是不成就.
者應許的.
願神照樣斗他.
離開家產.
和他勞碌得來的.
我斗空了.
這是一個很特殊的.
批鬥大會.
對貴就.
很無情地.
揭露了他們過去.
你可以說是結構性.
但這個結構性.
可能合法的.
因為外邦人都是這樣做的.
借錢就是要給利息的.
誰可憐你啊.
有免費午餐的嗎.
這些說話.
只是換了不同的方式.
這樣呈現.
但猶太人的社群.
摩西的律法書一早就說明.
你可以借給.
你窮乏的弟兄.
因為人.
意味著人.
總會有不同的.
三衰六旺.
很難去計算.
每一次你都是很正確的.
諸如此類.
但是.
有錢的.
幸運的.
有眼光的.
你不應該只是因為你的眼光.
不斷地自負.

$^{761}$就沾沾自喜.
就繼續.
將你擁有的.
成為了一個很巨大的資產.
當有人有.
極度困難的時候.
你將你的資產拿出來.
對不起.
你跟一般人一樣.
收取利息.
沒得還就要沒收你的田產.
這樣就帶來一個.
惡性循環.
現在尼希米面對的.
這個結局.
所以尼希米就.
感恩他有這樣的.
勇氣.
所以不只是發牢.
他有這個很強的勇氣.
勇氣涉及的是.
承擔.
面對的困難.
知道的困難.
但是承擔的困難.
就是要去解決.
最好的辦法就是.
有錢的人將你的資產.
你自己本身.
已經擁有很多的資產.
現在你擁有的是.
對方的弟兄.
其他支派的資產.
能不能夠慷慨地.
可能是有一些.
協議.
將這些資產.
給回他.
讓他能夠成為一個有資產的.
重新開始的.

$^{801}$一個這樣的規範.
他們答應.
這個是.
最奇妙地解決.
貧富懸殊的.
根本的一個問題.
出於雙方面大家有一個.
同心的愛.
很難.
但是真的可以解決.
政府.
國家甚至乎.
諾貝爾獎都解決不了的問題.
是有錢的人.
能夠這樣這種受人.
耳語的一種.
方式我將你.
賣給我的資產.
現在我給回你.
讓你去經營讓你去.
生產讓你慢慢.
致富將來你才還.
那些可能.
協議了的那個多少錢.
這個是很大的.
智慧.
當然永遠不是人想不到.
這些智慧.
是人根本不會實踐.
世界是不會這樣實踐.
我為什麼會將你.
給了我.
是我的那層樓已經入了我的名.
我現在給回這層樓.
你讓你慢慢在這裡居住.
你的子女讀書.
將來出人頭地.
你就再到時還給我.
怎麼去計.
沒得計.

$^{841}$你可以想像到.
尼熙米在這個時候.
他行使的是摩西的律法書.
那個精義.
那個精髓.
就是地不可以永賣.
是屬耶和華的.
就算你現在.
不去.
給回他.
去到希年的時候.
7749年的.
第50年.
就是希年.
希就是開心.
Happiness的意思.
希年的時候.
所有的.
地.
是要歸還給.
本身的業主.
就是說.
不可以永賣.
當然未到希年之前.
都要計算.
就是說.
你去查那個數.
距離希年多久.
按借貸的.
可能是那層樓.
那間屋的價值.
用了多少年.
去計算.
還有多少年是希年呢.
就去計算.
似乎雙方面都算合理.
不會偏向某一方面.
一百萬的.
那間屋.
還有很久才到希年.

$^{881}$你要什麼時候還給他.
可能又不可以一元.
如果是一個合法.
合理的賠償.
就可能要按價值.
這些似乎審判官都會做的評斷.
是律法書有教的.
你可想而知.
讀民數記.
利未記.
其實不沉悶的.
很多案例.
有些原則.
是很人道的.
神愛這些選民.
他們知道.
神知道將來一定會有.
這些訴訟的問題.
貧富懸殊的問題.
社會的商業.
貿易.
或者糾紛的問題.
不是每次都訴諸於.
公堂.
去對駁公堂.
可以找一個屬靈的領袖.
來仲裁.
大家在愛裡面.
怎樣來處理.
有錢那個人.
不單止.
不會突然間窮了.
不是好像.
打倒地主.
要打倒這些有錢人.
有錢就是邪惡的.
這些控訴很容易說了.
但是聖經從來不是這樣.
相反就是.
你在沒有影響.

$^{921}$你自己本身的價值.
之餘.
你輸出的是你很大的愛心.
和你給到一個動力.
窮乏的.
可以有機會重新的.
開始.
會激發他的動力.
他會繼續的去供書.
給子女讀書.
勤力讀書.
將來記得報答我們這個恩人.
那個愛的一個循環.
是一直在發生.
可想而知.
聖經那個奇妙.
就是在這裡發生.
希望有他那個很強的.
影響力去推動.
這個變革.
去推動這個.
推倒這個社會的結構.
不是說外邦人.
他現在不是.
將王的一個命令.
新的法則.
要打倒官商.
打倒有錢.
不是.
他現在是要.
搞共產不是這個意思.
他現在是要幫助.
這些投訴的窮人.
沒錯.
他們真是.
不知道什麼原因.
賣了他們的田地.
賣了他們的房屋.
賣了他們的兒女.
是鐵一般的事實.

$^{961}$那個藥那個賣身契.
都是不是被人.
恐嚇著來簽的.
是真是他們自願簽的.
但因為那個社會.
是這樣的方式.
但原來.
起碼可以不是這樣的方式.
你讓他從頭開始.
第十二節.
我就照了祭司.
叫眾人起誓.
必照著所應許的而行.
這裡我讀了.
眾都說阿們又讚美耶和華.
百姓就照著.
所應許的去行.
第十四節.
自從我奉派作猶大地的.
省長.
原來不單止是.
重建聖城的.
一個的遇子.
原來給了他.
猶大的省長.
河西省長.
李希米.
你可想而知.
他有實權.
但他沒有提.
他自己的身份.
相反在這個時候.
當他自己.
作為一個領袖.
作為一個.
某程度都是一個貴袖.
因為都是很高的.
政府都有高官.
和普通的.
前線的低級的民員.

$^{1001}$等等都有很不同的.
連薪金的.
技巧都不同.
李希米明顯是省長.
但他從來都沒有提.
去到這一刻第五章.
他才提他自己是猶大的省長.
這裡告訴我們.
從亞達士西王的.
二十年直到.
三十二年.
其實他一開波的時候.
第一章就已經說.
亞達士西王二十年.
幾時流月.
我在書山城的宮中.
我與我弟兄哈拿尼.
你可以想而知.
原來一開始.
他由二十年去到三十二年.
共十二年之久.
我與我弟兄.
都沒有吃省長的俸祿.
這樣就.
給我們一個疑問.
到底.
他這個的祈禱感動.
是因為他本身.
是省長.
不過很少上去.
還是由他.
被委派去.
重建聖城.
省長.
他的祈禱.
他的懇求.
而亞達士西王.
被他的真情感動.
就差派他.
無薪假.

$^{1041}$上去重建.
你祖宗的聖城.
我就傾向後者.
美麗很多.
然後.
因為都擔心.
河西的地方官.
未必會.
接納.
同時間.
委派他的.
愛臣.
這個酒政.
尼熹米.
作猶大省的省長.
當然在河西.
仍然有他們的.
河西的首長.
即是外邦的首長.
總督.
就是.
波斯帝國的總督.
就不是猶大.
猶大省.
即是說他們猶大.
猶太人.
可以有少少自治的.
宗教上.
當然不是政治上.
宗教上他們可以有一個.
自己的自治.
尼熹米都被委派.
作這個省長.
原來不只是52日做好成長.
原來這裡又再揭露了.
有12年的時間.
尼熹米.
做好了成長之後.
第六章就說做好成長.
尼熹米原來有這12年.

$^{1081}$穿梭兩地.
他來到去.
主政於猶大.
那個聖文的.
靈性的重建工程.
這個就.
真是.
在這個回憶錄裡面.
這裡就真是揭示.
原來尼熹米.
他上到去看到的.
不只是硬件的成長.
他看到的是.
神所愛.
亞伯拉罕的自己的同胞.
猶太人.
他們那個靈性.
所謂靈性就是有沒有照摩西的律法書.
來到去生活.
可能買賣商業的活動.
這裡告訴我們.
一涉及到金錢.
他說我和我的弟兄.
都沒有再吃.
省長的俸祿.
什麼意思呢.
不出糧還是拿股票.
認股權還是怎樣.
原來不是.
在我以前的省長.
加重了百姓的擔.
每日索要.
糧食和酒.
並銀紙四十赤勒.
就是他們的僕人.
也合濟百姓.
但我因敬畏神.
不者讓行.
即是說皇帝原來准許.
你是駐猶大的省長.

$^{1121}$你就可以拿那些俸祿.
就是由百姓提供給任命的省長.
在他前面原來也有省長.
真的沒錯,真的有省長.
所以尼熙米是一個.
隨著他回去重建.
而被委派做省長.
我自己這個假設就會成立.
就不是他一直都是省長.
所以他才會這樣祈禱.
我就認為不是.
問題就在這裡.
他上到去之後.
觀察到民情.
感受到那種壓力.
當然有些東西是結構性的.
例如省長.
根本不理那麼多.
我要發工資是天經地義的.
你們不能給是你們的事.
我就會用不同的方法.
來追收百姓的糧食.
酒,銀紙.
這裡說44下.
甚至乎.
核製那些百姓.
獨人.
省長的獨人.
都去核製那些百姓.
更加令到.
那個關係破裂.
更加令到.
這班人只是養尊處優.
懂得來開會.
開會完之後就去.
公費的去吃什麼.
開什麼酒,茅台.
但是對百姓的死活.
就完全不理.
你起碼上來不是做省長.

$^{1161}$他上來是重建.
姓王委派他這個官職.
方便他辦事.
所以我們知道原來這個不出糧.
是這樣的原因.
固然是這樣.
還有就是.
他將自己本身有的錢.
拿來給百姓.
第十六節.
(念句).
(念曲).
都聚集在那裡作供.
除了從四圍外邦中來的猶大人以外.
有猶大平民和官長.
150人在我席上吃飯.
每日預備一隻公牛,六隻肥羊.
又預備些飛禽.
每十日一次多預備各樣的酒.
雖然如此.
我並不要省長的俸祿.
因為百姓服役甚重.
兩次提到他不要省長.
應份的.
是皇帝給的.
合法的.
不是貪污,不是枉法.
不是來強加徵收的.
但兩次提.
在這樣的時候.
要共度時艱.
我不要省長的俸祿.
但有150人因為工程.
在他家裡的大宅.
即省長的官邸.
來一起吃大鑊飯.
預備的物資.
即是意味著.
一定不是百姓.
怎樣沒有你都要給我.

$^{1201}$一定是他自己來預備.
他自己來買出來等等.
例如拆開說明.
明晚我們一班.
上星期四晚晚宴時.
逢的弟兄姊妹.
明晚會有一個簡單的分享會.
在會所和同宮等等.
誰付錢.
我問不如教會出.
幾位領袖說不要.
不要教會出.
反正都是籌款晚宴.
我們幾個分了算了.
這些是一些.
你要看到背後.
本來教會出沒問題.
都是慰勞弟兄姊妹這麼辛苦.
但意義就不好.
教會出就變了.
又用多了一兩千元.
可能不止.
如果幾位三十幾人的話.
幾位領袖都不用說.
二話不說.
我們幾個分了算.
不需要.
你就可以想像到.
李兄妹就是這樣.
一百五十人.
是真的大家在做工程.
我專心工程.
我沒時間去置買田地.
就是說不會陷人在那個.
痛苦危險的時候.
你就再落井下石.
嘩現在這麼便宜.
我就買多一點.
反正我有錢.
我有比其他的百姓多的那個.

$^{1241}$譬如現金.
我有多的物業資產.
我去買就很容易了.
現在這麼便宜.
那些田地.
那個房屋.
嘩這麼便宜.
我就現在買多一點了.
你就可以想像.
李兄妹她表達的.
她不是想告訴我們.
她多麼淑齡.
而是她多麼緊張神的事工.
她看到的是更加長遠.
如果我建成那些城牆又如何呢.
多漂亮.
裝飾完.
很好的那個結構又如何呢.
裡面更嚴重的就是.
百姓叫苦連天.
不是現在沒得喝紅酒.
我連基本謀生都沒有.
我怎麼還說什麼.
恢復聖文的那個.
淑齡啊.
我連我的子女都賣了.
做了一幅被子.
我還說什麼.
星期六.
安息日去敬拜神.
和富貴的弟兄一起坐在那裡.
敬拜神.
我是欠他錢.
他正在追我的利息.
我怎麼可能來到.
我能夠很平安地.
去敬拜這一位.
說自己是很公義.
很信實的神呢.
很富貴的弟兄就說.

$^{1281}$你的供應很豐富.
我一無所缺.
我很開心.
所以我每次來聖殿.
我都很開心.
我敬拜你.
你想想那個情況.
出現的問題是多麼嚴重呢.
其實新約主人是有教的.
如果你明明白白地.
得罪了人.
不肯來尋求饒恕.
那個和你在一起.
事奉的弟兄姊妹.
一見到你.
心裡就火都來了.
他都不能夠事奉了.
因為他已經發脾氣了.
在心裡.
你就繼續.
心裡吹著口罩.
大搖大擺.
還要在講台上.
振振有詞.
可能做領唱.
可能做其他東西.
你想想那個落差.
主耶穌說.
你去獻禮物給神的時候.
如果想起有弟兄向你懷怨.
你要扔下那些禮物.
首先去和那個弟兄和好.
然後才來獻禮物.
獻禮物是對的.
是好的.
是神預納的.
但是在那個同場裡.
不和好是帶來.
神都不能夠賜下那個祝福.
雙方面都沒有了祝福.

$^{1321}$懷怨的那個繼續懷怨.
不肯來寬恕.
不肯來和好的.
繼續停留在那種自我感覺良好.
那個教會就變成.
或者那個聖殿就變成.
一個充滿了苦毒的地方.
所以這裡就已經告訴我們.
尼熙米他看到這些問題.
內部.
所以第五章完全沒有了.
外邦人的那些問題.
他不斷凸顯的就是.
聖文怎樣能夠讓他們.
走回上那個誠信的路上.
尼熙米首先就是和他的下屬.
他的僕人.
他們首先就是真的以身作則.
一心的為建造神的成長.
知道那麼重的那個搖曳.
那些伏曳的艱難.
他就不要那些俸祿.
就是說免去百姓再上繳的那些需要.
留給他們自己家庭的需要.
然後他們就可以專心的去替.
自己的聖殿.
神的聖殿.
來一起同心的去建造.
所以第六章很快就見到那個成效.
第五章就告訴我們.
內部這個問題怎樣能夠梳理.
首先就是尼熙米他觀察到.
意識到又真的有具名的.
三個案例的投訴.
具名了.
有的說有的說.
有的說.
雖然他沒有說那個人名.
他們當時一定知道.
然後不是特意去為反而反.

$^{1361}$而是焦聚大會.
就是說這件事已經不是個別的事件.
是一個可能社會結構性的問題.
怎樣可以解決這個結構性的情況呢.
就真的焦聚了所有的同胞一起.
貴就明顯就是知道自己過去這樣做.
不過是好像跟隨那個潮流.
跟隨那個世俗.
外邦人都是這樣.
為什麼我不行.
我真是真金白銀借給我的弟兄.
為什麼你還給外邦人.
你就可以還.
還給我你就可以吵.
可能他們在靈性低的時候.
就會這樣振振有詞.
但當尼熙米點出那個關鍵.
要有一個讓人有翻身的機會的恩典.
就不是外邦人有了.
外邦人就不會跟你講人情.
但是在神的聖殿的選民.
就應該要講那個.
怎樣能夠可以彼此的均平呢.
這個就是很高度的智慧.
其實是很簡單的.
將心比己的一份愛.
尼熙米就能夠動之以情.
他又是以身作則.
他又是不會逃避這個問題.
也都不是要碰外邦人.
他們的那個商業貿易.
他不是要現在在這裡.
你刻意來推翻一些認為.
凡是借錢都是罪.
都是不對的.
聖經都有教.
舊約摩西都有教借的.
但是你借錢給你窮乏的弟兄.
今晚大家回去找找.
應該是在埃及記和申命記都有的.

$^{1401}$我忘記了申命記第幾章了.
有教導怎樣來處理的.
但是外邦人是可以處理.
但是猶太人是不可以處理.
背後的原因就是不希望.
那個貧者越貧.
繼續結構性的.
去到一個完全是.
跌落一個陷阱裡面.
是出不來的一個深淵裡面.
但是有的人就繼續不斷.
透過他的資產無限增值.
那就帶來.
真的會出現這個問題了.
所以我現在一而再看到.
中世紀的資產家.
基督徒.
那個時候是叫新教徒.
因為剛剛宗教改革.
16世紀.
社會主義學家馬克思.
就觀察了歐洲.
很多的新教徒.
尤其是新教徒.
基督徒.
因為當時很多的是清教徒.
他們有一個很強烈的盼望.
就是我們是不是神揀選的子民.
如果是神揀選.
必然可以上天堂.
和一生蒙福.
就導致了一個這樣的.
由一個這樣的信念.
一個的信仰.
導致了他們的生活行為就改變了.
怎樣改變呢.
殷勤工作.
這個是神.
差不多是命令每一個人.
在地上.

$^{1441}$你有開始工作能力.
你就要殷勤的工作.
當然.
什麼工種.
是每個人都不同.
拼命工作.
自然就會致富了.
致富之後.
新教徒就有一個很強烈的信念.
錢是神的.
我是管家.
所以我要努力去幫人.
和奉獻給教會.
就帶來他們在這個過程.
就致富.
致富他們又不會隨便亂花錢.
所以那個社會學家就寫了一本書.
新教的倫理.
是強調新教.
他是不信耶穌的.
新教的倫理.
就是剛才說的這個價值觀.
就是誰才發現自己是蒙神揀選.
就是我要在人生.
努力工作.
努力為神.
我就會是蒙神揀選.
這個過程原來不知不覺間.
就由小康變成為資本家.
變成為企業家.
這個就是新教倫理.
和資本主義的精神.
有一本書.
是讀社會學必須讀的.
經濟學都必須讀的.
句誅.
就是這個point.
就是看到這班基督徒.
集體的勤力工作.
集體的致富之後.

$^{1481}$不是大魚大肉地花.
而是怎樣累積他們的財富.
可以繼續去投資.
繼續去有錢.
是怎樣去為神事奉.
所以新教徒.
是有這個被外邦人都看到的.
一種這樣的特質.
你可以想像得到.
十六世紀的宗教改革.
那個大的覺醒.
導致十八世紀的大復興.
約翰·維斯利.
美國的就是.
我突然間忘記名字了.
英國就是維斯利·約翰.
和他的哥哥.
瑪利·維斯利.
然後就引發.
美國歐洲的大復興之後.
十九世紀開始.
就是普世的宣教運動.
因為這麼多的歐洲國家.
已經是基督教了.
但是非洲,亞洲,印度.
全部都還沒有信主.
大量的猜派年青的宣教士.
由大復興,大覺醒.
去到普世宣教.
是一個一層一層地發生的漣漪.
一直是源於宗教改革.
Back to the Bible.
唯獨信心,唯獨聖經,唯獨基督.
而帶來的大的扭轉.
當然馬丁路德只是其中一個.
點著這場火的奮興的人.
但是整個歐洲當時因為有印刷術.
大量印聖經.
也開始有翻譯.
普通平民百姓都可以自己讀聖經.

$^{1521}$就出現這種雨後春筍對聖經的研讀.
研讀帶來的重大啟發.
所以尼欽米在古代二千四百多年前.
已經是這樣的對神的話語.
摩西五經的精神.
他自己身體力行去用.
他上來拿省長俸祿.
如果大家都是很開心的.
大家都是很順暢的.
國泰民安的.
是沒有問題.
但是首先就是已經外邦人冷嘲熱諷.
批評,攻擊.
甚至動用一些輿論諸如此類.
但是尼欽米知道這些是難避免的.
外邦人一定是會反對的.
但是到第四章的時候.
對不起.
到第五章的時候.
就出現了這些結構性財富不均勻.
財富不均勻你要找個原因.
但是你要解決的那個辦法.
好像很翻天覆地.
從來不會這樣的.
他叫那些貴就.
你將那些給他們.
重新協議過一些協議書.
不要收取利息.
你們是將你們多餘的給他們.
不是將你們本身有的.
去被人搶來.
好像鬥那些地主那樣.
不是這樣.
不是要這樣來到.
不是均平.
不是這樣均平.
是你能不能夠在你的好友餘裡面.
你去祝福到其他的人.
你可以想像得到.
如果我們能夠每個月.

$^{1561}$我們自己在我們銀行的財富裡面.
我們自己有一個個人的心願.
我拿一些錢出來.
真的要幫某些人.
這個是神很大的祝福.
因為你可以繼續一直去儲.
只不過可能多了幾個零.
但是有些真的有很嚴重需要的.
可能和時間競賽中的.
你能不能夠幫助呢.
這些就是教會為什麼能夠可以動員.
和鼓勵弟兄姊妹去參與的.
那個的聖公.
和那個的工作裡面的背後的意義.
我們就能夠可以將那個視野.
幫弟兄姊妹可以看得更加遠更加闊.
這個是很重要的.
我們要將我們內心.
錢是會讓我們越來越緊閉.
越來越看得不夠長遠的.
錢是會令我們更加的.
所謂守財奴.
會令我們更加對人.
沒有了那個憐憫的心.
這個是我們需要求主幫助.
但是要有智慧地.
有策略地去行駛的.
尼希米就是讓我們看到.
他這個真是以身作則的一個領袖.
自己這樣去做.
你還有什麼可以這樣說呢.
貴就官長.
祭司.
他是招聚祭司來處理.
叫祭司來向眾人起誓.
必定照著所應許.
就是他所建議的.
來去執行.
接著他就要.
好像根據摩西的律法書.

$^{1601}$就抖著胸前的衣襟.
就是說.
誰不照這樣做的.
就好像這件衣襟一樣.
被神趕逐離開這個聖地.
就是和這個民族無關.
將來有的Blessing是沒有你的份.
這個是非常嚴重的.
我們要緊張神的Blessing.
不然就很嚴肅.
所以猶太人他們知道.
這件事不是開玩笑的.
但是尼希米大前提.
不是為批鬥而批鬥.
不只是他站在窮人那邊.
窮人永遠是對的.
這個本身是有問題的.
而是告訴我們.
要解決一個現在根本內部的問題.
他不是要來到去.
透過鬥地主.
你有錢你是錯你是罪.
你不應該這樣.
那又是過分了的一種簡單主義了.
但是怎樣能夠.
大家在一個困難.
在一個富裕裡面.
能夠可以透過一個動力.
可以幫助到那個的循環.
不是一個惡性的循環.
而是一個愛的善意的舉措.
而帶來的是一個愛的回應.
可能要下一代他們的子女讀完大學.
才有機會去報答這個恩人.
當年的幫助.
下一代才報恩.
那個循環就不是一個惡性的.
苦毒的循環.
是一個愛的循環.
你可以想像得到.

$^{1641}$當年投放的資源.
換取的是那個家族的祝福.
一個人對你的恩典的報答.
我提過了.
我是入建道的時候.
05年.
是有19封個人推薦信.
因為沒有教會推薦.
沒有牧者.
執事會的推薦信.
很特殊.
當年建道都看我是很特殊.
但知道是有特殊的原因.
所以就破格收了我.
接著又是破格.
為我開了一個戶口.
在建道.
我是一毛錢都碰不到的.
每年的學費.
就有愛心弟兄姊妹.
擠到那個pool裡.
不夠.
他們繼續侵.
四年的神學.
我是不用給學費.
你可以想像得到.
這一份的愛心.
我跟明仔說.
這些是當年.
我能夠再讀神學的恩情.
這些弟兄姊妹.
有些是認識我的.
有些是不認識的.
知道聽到我的故事.
就這樣放進去.
神學院有很多這些故事.
所以這些是.
你怎樣去幫助.
不同有需要的人.
讓他能夠可以被建立.

$^{1681}$這些就是將來的恩情回報.
甚至可能要去到天家.
你就會得到很多.
不知道是誰走過來.
多謝你.
將來是怎樣.
真的不太知道.
但是至少在神裡面.
那個恩典的戶口.
一定有你那一份.
所以求主讓我們能夠.
在貧富懸殊.
看我們能夠去幫助人.
看得闊一些.
在永恆去看.
不只是短暫的去看.
那就會更加值得我們去看.
錢原來這麼有用.
原來可以用得這麼開心.
我們就要求神.
讓我們能夠在這裡去學習.
譬如我的弟兄.
知道我要去神學.
給我一間屋住.
住到畢業.
一毛錢租金都不需要交.
他自己搬去天水圍住.
這樣的一個恩典.
四年.
你可以想像得到.
他每個月可以收多少錢租金.
你可以計一計那裡.
多少十萬.
如果四年.
這些都是很大的恩情.
到今天還是保持一個很愛的關係.
一段時間他不用再做執事.
他就走出來.
我們做聚會.
但因為很忙.

$^{1721}$沒有機會和他聊天.
你可以想像得到.
大家在神裡面.
有這種彼此的服侍.
是一件很美麗的事.
差不多了.
我們第六章我不說了.
大家回去再看.
因為再次涉及到一些陰謀.
甚至乎是一些很陰險的攻擊.
黎希米.
都是黎希米要排除.
所以作為領袖.
心胸要闊.
肩頭要夠硬.
要有承擔.
自己都是以身作則.
被人看到.
我是為那個事公.
完全沒有所謂針對哪個人.
人家都會慢慢.
就算那些貴就.
官長都會服.
黎希米自己.
你說要將我們清算.
但你自己首先清算你自己.
大家都會明白.
是想整個猶太的選民.
重新上回那個屬靈的起步點.
來前行.
所以慢慢慢慢.
當他們見好之後.
他們就會有他們那個.
真正的嚴肅會.
真正的.
大家回到神的摩西五經裡面.
去思想.
我們一起祈禱.
多謝主.
今天晚上我們能夠透過第五章.

$^{1761}$去思想一下一個領袖.
他的那個內心.
他的那個情操.
我們見到他發怒.
他都是日於言表.
在他的自傳的回憶錄.
來講述.
真是讓我們感受到.
他不認同.
猶太人要跟世俗的價值觀.
將這種買賣合理平常化.
令到猶太的窮乏的弟兄.
要永遠成為奴隸.
甚至乎要永遠在那個自賣產業上.
受到那個結構性的剝削.
甚至乎可能那個租金的利息.
或者那些要償還的借貸的利息.
根本是他們都已經解決不了.
所以我們看到.
好像這一章是極之.
如果我們錯誤的解讀的時候.
以為是一些鬥爭.
以為是一些清算.
一些批鬥.
甚至乎是動員一些窮人.
來刻意去做一些.
可能是反社會的事件.
因為當我們看到尼熙米.
終於他講自己.
十二年都是做省長.
他看到那個光景.
他好像自願.
甚至乎是不要任何的俸祿.
來完成這個工程.
真的自己不單止是.
由這個舒山城.
上到去耶路撒冷.
路途的遙遠.
他這樣來到去.
身體力行的.

$^{1801}$在百姓的當中.
來到去工作.
調度和鼓勵.
甚至乎受盡的.
都是外邦人的嘲笑.
攻擊和對他的動機的.
一種懷疑等等.
尼熙米就是這樣.
看到自己是一個敬畏神的人.
他不敢的來得罪主.
他來到去做的事情.
內心的動機出發點.
都是要幫助百姓.
避免在一種互相的埋怨底下.
整個的事工.
其實根本就已經瓦解.
所以我們看到這個是很寶貴的.
一個領袖的生命的榜樣.
成為我們大家都去學習的功課.
幫助我們能夠回來教會.
或者我們在不同的地方.
我們都能夠看到.
我們能夠服侍都已經是恩典的.
以致我們真的可以在神的面前.
來到去做一個良善.
忠心有見識的管家.
幫助我們能夠看到.
我們管家的這個身份.
以致我們.
好像主耶穌所說的.
所求於管家是要他忠心的.
幫助我們能夠在我們的職份上.
忠心盡責.
多謝您聽我們禱告.
奉耶穌基督的名.
阿們.
\newpage



\section{約翰一書 3:11-18}
\label{sec:eLIEUCC43fA}
\textbf{2023-05-07 加恩樓崇拜《共融》約翰一書 3\_11-18 講員:吳景霖傳道}
\newline
\newline
連結: \href{https://youtube.com/watch?v=eLIEUCC43fA}{\texttt{https://youtube.com/watch?v=eLIEUCC43fA}} ~~~~ 語音日期: 2023-05-09
\newline
\newline
\hyperref[sec:AGfClw0jrYY]{\small{< < < PREV SERMON < < <}}
~
\hyperref[sec:index_chronic]{\small{[返順時目]}}
~
\hyperref[sec:index_scriptual]{\small{[返順卷目]}}
~
\hyperref[sec:_3DxJVmmiXg]{\small{> > > NEXT SERMON > > >}}
\newline
\newline
約翰一書 3:11-18
\newline
\begin{longtable}{cl}
\hline
\hline
章節 & 經文 (和合本修訂版)\\
\hline
3:11 & \begin{tabularx}{0.7\textwidth}{X} 我們要彼此相愛。這就是你們從起初所聽到的信息。 \end{tabularx} \\ \\ \relax
3:12 & \begin{tabularx}{0.7\textwidth}{X} 不要像該隱；他是屬那邪惡者，殺了自己的弟弟。為甚麼殺了他呢？因為自己的行為是邪惡的，而弟弟的行為是正直的。 \end{tabularx} \\ \\ \relax
3:13 & \begin{tabularx}{0.7\textwidth}{X} 弟兄們，世人若恨你們，不要驚訝。 \end{tabularx} \\ \\ \relax
3:14 & \begin{tabularx}{0.7\textwidth}{X} 我們知道，我們已經出死入生了，因為我們愛弟兄。沒有愛心的，仍住在死中。 \end{tabularx} \\ \\ \relax
3:15 & \begin{tabularx}{0.7\textwidth}{X} 凡恨自己弟兄的，就是殺人的；你們知道，凡殺人的，沒有永生住在他裡面。 \end{tabularx} \\ \\ \relax
3:16 & \begin{tabularx}{0.7\textwidth}{X} 基督為我們捨命，我們從此就知道何為愛；我們也當為弟兄捨命。 \end{tabularx} \\ \\ \relax
3:17 & \begin{tabularx}{0.7\textwidth}{X} 凡有世上財物的，看見弟兄缺乏，卻關閉了惻隱的心，神的愛怎能住在他裡面呢？ \end{tabularx} \\ \\ \relax
3:18 & \begin{tabularx}{0.7\textwidth}{X} 孩子們哪，我們相愛，不要只在言語或舌頭上，總要以行為和真誠表現出來。 \end{tabularx} \\ \\
[1ex]
\hline
\hline
\end{longtable}
$^{1}$大型姊妹早晨.
狀態不太好 因為上星期剛中完第二次COVID.
所以戴口罩保護大家.
今早肚子餓了兩次.
所以今天狀態不太好.
但都不影響 內容早就預備了.
今次講題其實就配合教會的方向.
就會講關於團契.
我十九歲的時候信耶穌.
那時候是在一個細胞小組教會.
沒有團契的 細胞小組教會分小組.
在裡面我開始服事.
當初回教會的時候 我是一個比較討人厭的人.
為什麼討人厭呢 因為我很喜歡問問題.
我經常問問題 問到有些比我高年紀的不懂回答.
當然是惱羞成怒.
到了這個年紀我依然很喜歡問問題.
我會經常問別人的問題.
剛才我說我那時候是在細胞小組.
然後一開始去了另一間教會就回團契.
我都會問團契和小組有什麼分別.
又是十幾二十人圍在一起.
又是卓詩,查經,祈禱,分享.
又是那些東西 偶爾出去吃飯.
然後又去勾挺 有什麼分別呢.
有些有傳道人 有些沒有傳道人.
有些是由信徒領袖帶領.
到底團契是什麼呢.
到了這一刻我問你們團契是什麼呢.
為免我繼續令人討人厭.
我就不要大家回答.
不然等下被我問到口啞不知怎麼辦.
團契是不是一個組織呢.
是不是一個組織呢.
我們教會有很多團.
溫柔良善 信實和平.
忍耐節制 恩慈 古靈精怪.
獅子坑都有一個.
很多不同名字的團契.
是不是一個組織呢.

$^{41}$團契與團契之間是否有關係呢.
是不是你的團和我的團不合.
或者不同的團契不一起去玩呢.
其實我們看聖經裡面說的.
團契這個字.
Colonia 是一個有正面含義的詞語.
意思是相交 分享 有份 一同.
意思就是與某個人分享某樣東西.
那種為之團契.
是一種分享.
如果我們這樣看的時候.
團契不是一個組織.
我們看一些經文.
多倫多前書一章九節這樣說.
「神是信實的 他呼召你們好與他而子」.
「我們的主耶穌基督共享團契」.
這個團契的意思.
我斷估也不是說耶穌基督.
來與我們一起去開團契.
他所說的共享團契的意思是.
神呼召我們與耶穌基督一起去分享.
你與耶穌基督可以一起去分享一些東西.
分受一些東西 共同去擁有一些東西.
然後在肥勒比書三章十節說.
「使我認識基督 知道祂復活的大能」.
「並且知道和祂一同受苦效法祂的死」.
好 那考考你了.
這裡哪個字是Colonia.
剛才解釋過 一同.
如果用那個字眼.
就是「和祂團契受苦」.
你原來是可以與祂一起去受苦.
團契的意思是一起 一同去受苦.
所以團契的意思是.
相交分享 一同去分受一些東西.
有另一個翻譯說.
將團契翻譯為「共融」.
共融 可能更加表達到團契的意思.
團契就是一種共融的關係.
融入了 不是說我融入了你.

$^{81}$大家融入 大家融洽.
大家有一種混在一起的關係感覺.
這樣就更加表達到團契.
所以說團契並不是一個組織.
但是我們教會這麼大.
總要有組織的.
總要有些東西去分.
譬如某個年紀的一班.
或者某種特性的一班.
我們有婦女團契的.
其實應該可以有男士團契.
不過通常都是失敗的男士團契.
因為男人不喜歡說話.
但是如果是男士團契.
可以去喝喝吃吃 又可以組織起來.
真的很有趣.
但是不同的年紀.
甚至有些商人團契.
可能是做生意的一班.
比較合得來.
有些弟兄姊妹說生意經的時候.
我是不能回答的.
所以不同的團契有它的功用.
但是說真的.
你是不是說落戶了某一個團契.
你就跟旁邊那個沒有關係呢.
這個團契跟那個團契根本就不認識.
是不是這樣呢?.
就不是的.
你只要是基督徒.
不論你是什麼年紀.
什麼人種.
你是在哪裡信的.
你都應該是跟你的弟兄姊妹.
是有一種共融的關係.
是一個團契來的.
OK.
在二次大戰期間.
德國有一個很出名的神學家.
他叫潘霍華.

$^{121}$他寫了一本書叫做《團契生活》.
很薄的.
我們讀神學的都要看的.
你們都可以看的.
不是什麼高深莫測的東西.
很薄.
但是《團契生活》這本書.
就是說了一些很重要的東西.
什麼是團契?.
團契就是透過耶穌基督.
並在耶穌基督裡面那種甜美和諧的生活.
所謂的團契不多不少.
正是這樣.
無論是短暫的.
一次性的.
經年累月的.
還是每天都有的團契.
都是這樣.
我們唯獨透過耶穌基督.
也是在耶穌基督裡面.
我們才彼此相熟.
團契最重要的本質.
就是共融.
共融的最重要的本質就是.
透過耶穌基督在耶穌基督裡面生活.
所以耶穌和生活這兩點在團契裡面是分不開的.
你有團契你就要在團契裡面有生活.
你的團契也是融入了你的生活.
不是說星期一至六.
我就在外面生活.
星期日回來.
我就掛著另一個面孔.
或者是認識你十幾年了.
我也不知道你做哪一行.
或者我不知道你日常生活是怎樣的.
這些就是生活和信仰分割了.
但其實是要生活.
有些人幻想來到教會的時候.
進來教會應該是怎樣的呢?.
每個人都像天使一樣的面孔.

$^{161}$你隱約看到頭頂有個光圈.
每個人都是好人.
進到教會裡面沒有鬥爭的.
沒有不開心的.
每個人都像小天使一樣.
人畜無害的.
一群羊咩咩.
但事實上教會是不是這樣的呢?.
我回到教會二十多年.
我的很大的感覺就是基督徒都是罪人.
其實很多人都是三尖八角.
我自問我也是三尖八角的.
我經常都會刺到不少人.
我相信有很多人對我都很生氣的.
但原來教會就是這樣的.
教會都是人聚集的地方.
有人的地方就有問題.
剛才說我信主初期就在細胞小組成長.
細胞小組的特色就是很快就可以開始事奉.
其實浸信會也是很快就可以洗禮.
洗禮可以很快就可以事奉.
我開始做組長帶茶經.
後來轉了教會去了另一間教會.
我又很快就做了團長.
到現在成為了牧者.
我有很久的帶領性崗位的事奉經驗.
但我會發覺帶領一個團體是真的很麻煩.
在團體裡面例如安排周會.
團體的本質是什麼?.
耶穌嘛.
那你就一定要茶經的了.
要帶茶經.
好啊,你茶了幾次.
然後開始有弟兄姊妹投訴.
你每次都茶經很悶.
可不可以有些活動.
好好好.
然後我就搞一些outing.
出去玩玩.
玩了幾次有些就說.

$^{201}$每次都玩不如坐下來吃飯.
好好好,坐下來吃飯.
吃了幾次飯然後又有人投訴.
你每次都是玩和吃飯.
可不可以茶經.
其實很難服侍的.
真的有些弟兄姊妹.
坦白說.
不停地有不同的要求.
還有有些弟兄姊妹.
每個人的性格都不同.
有些好動有些好靜.
有些喜歡玩有些喜歡聊天.
有些喜歡坐下來吃飯.
最慘的就是有些很追求真理的.
每次茶經都會說一些問題.
有些就在釣魚.
那怎麼辦呢?.
還有.
有些就在做很穩定的工作.
收入不錯.
一至五.
有些就上輪班.
收入沒有那麼高.
你安排的活動要走到兩邊是很困難的.
是很困難的.
你去那邊.
這邊又說嫌貴.
你去那邊.
那邊又說有些這麼低級的東西.
不要叫我去.
我又說搞短孫.
坐飛機去的那些.
有些嫌貴.
有些說沒有假期.
好了.
我不在公眾假期搞了.
我在閒日那裡搞了.
然後那些就說你在閒日搞.
你怎麼參與呢?.

$^{241}$總之是很多很困難很麻煩.
我想起摩西帶著以色列人出埃及就是這樣.
最近說沒有肉吃.
那裡說沒有水喝.
接著這班就說你霸佔了先知位來做.
還要造反.
所以其實是很困難.
頭也戴上了.
你會發覺團契.
Colonial本身就是一種生活.
你的生活是融入了教會的生活裡面.
並且是會發生很多的衝突.
很多的麻煩.
很多的東西.
如果一個團契裡面是完全沒有衝突的.
完全是大家沒有意見的.
其實團契又很死.
可能大家都是不將他的生活.
或者不將他所有東西帶回來.
只不過是人道的靈魂不在.
其實我們在這個群體裡面.
我們這麼困難.
要融入是這麼困難.
或者有這麼多東西的時候.
我們是不是應該抽離一點呢?.
正所謂你不放這麼多心機進去.
就不會有這麼大的傷害.
也不是的.
其實團契我們說的是一種共融的關係的話.
你一定要將你的生活帶進去.
並且你不要怕會有衝突.
你不要怕有衝突.
有衝突是正常的.
大家有事情要談.
有事情要討論.
這是正常的.
在潘博華的團契生活裡面.
另外有一套這樣說.
他說一個團契.
除非遭到各種的不快和醜惡.

$^{281}$因而陷入極大的失望.
就不知道在上帝面前應該怎麼走.
但經歷了這些事之後.
就開始憑著信心找上帝給的應許.
所以這些失望越早臨到個人和團契當中.
對兩者就越有利.
剛才也是潘博華說的.
甜美和諧的生活.
轉過頭來說.
一定要遭受到各種的不快和醜惡.
為什麼呢?.
因為現實就是你不遭受過這些不快和醜惡的話.
其實你是去不了甜美的生活的.
是不是?.
可能兩夫妻一樣.
你不吵架的話.
你怎麼會進入甜美的生活呢?.
你總要先吵架.
如果結婚之前還是很沮喪的.
大家相敬如賓.
很和諧的樣子.
你結婚之後吵架了你就知道味道.
有時候有些事情是.
你經歷過衝突.
你經歷過矛盾.
你才知道原來那個人是這樣的.
或者你知道怎樣去到.
和他保持一個恰當的距離.
或者是你了解他的真性情.
團契之中就是需要這樣.
我們要面對那些不愉快的醜惡.
我們要面對失望.
一開始對那個人就很有期望.
發覺原來他都不是我想像中的那樣.
是好事來的.
打破了你的幻想之後.
其實你更加明白.
原來這個人是怎樣.
他在你面前是更加真實.
更加是一個活生生的人.

$^{321}$不是一個你想像中的圖畫.
我剛才都說了.
我都三尖八角的.
我都三尖八角.
我不怕被人知道.
有時候我是有一些軟弱.
或者是有一些掘雷錘的地方.
或者是有一些我做不到的地方.
有很多時候.
特別是牧者.
人家總會對牧者有一個完美的期望.
你是牧者來的.
為什麼會發脾氣.
為什麼會生氣.
為什麼會軟弱.
為什麼會犯罪.
為什麼你會累.
為什麼你會生病.
差不多是這樣.
為什麼你回不來.
但是你明白他是人.
然後打破了這些時候.
我覺得我和弟兄姊妹之間的關係是比較真實的.
我不怕告訴他我軟弱.
我不怕告訴他我失望.
我不怕告訴他我都不夠信心.
然後大家一起去祈禱.
大家一起互相去支持.
所以這些東西原來是越早看到的.
就越好.
越早看到就越好.
當我們遇到不愉快的事.
當我們遇到不好的事.
有些人可能會選擇逃避.
有些人會選擇忍一步.
就是忍一忍.
大事化小.
小事化無.
忍一時風平浪靜.
退一步海闊天空.

$^{361}$有些人會選擇離開.
有些人選擇不再回這個團計.
這麼多團計我慢慢選擇.
是不是這樣的呢.
不是的.
我覺得是要面對.
我們要面對團計發生的事.
今天的經文說到.
我們要彼此相愛.
先找回經文.
我們要彼此相愛.
這就是你們從起初所聽到的訊息.
不要像該人.
他是屬那惡者.
殺了自己弟弟.
為什麼殺了他呢.
因為自己的行為是邪惡的.
而弟弟的行為是正直的.
弟兄們.
世人若恨你們不要驚訝.
因為我們知道我們已經出死入生了.
因為我們愛弟兄.
沒有愛心的仍住在死中.
凡恨自己弟兄的.
就是殺人的.
你們知道凡殺人的.
沒有永生在他裡面.
基督為我們寫明.
我們從此就知道何為愛.
我們也當為弟兄寫明.
其實約翰一書的作者.
他用一種很二元對立的方式去寫作.
什麼叫二元對立呢.
其實我們了解這個世界.
你總是將一些東西分成兩邊的二元.
譬如男人和女人.
高和矮.
肥和瘦.
有很多東西是分成二元的.
相對的.

$^{401}$他用一種非黑即白的方式去描述.
他在描述弟兄姊妹之間的關係.
他也在用舊約的一個人物.
象徵性的人物就是該人.
他用該人來象徵一種邪惡.
其實一說就是了.
就像我們說的踩浪.
踩浪象徵什麼呢?.
兇殘.
狐狸呢?.
狡猾.
我們有一種對象徵性是很肯定的.
譬如我稱讚一個姐妹.
你很像狐狸.
我想她是很不高興的.
你會知道她是有一種負面的象徵.
一說比如說秦始皇或者紂王.
你就知道是一個暴君.
如此類推.
該人就是象徵著屬於魔鬼的邪惡的惡者.
他用一個很極端的方法去比較.
愛弟兄.
對比著恨你的弟兄.
殺你的弟兄.
徵籍.
對比著邪惡.
就是蓋因和阿伯.
你是屬於耶穌還是魔鬼呢?.
你是為弟兄捨命還是你殺了他拿他的命呢?.
這是一種很極端的表達方式.
他告訴我們作為基督徒你有沒有選擇呢?.
沒有選擇的.
你一定要愛弟兄.
如果你不是愛弟兄.
你就殺他那個了.
是不是很極端呢?.
除非你不做基督徒.
否則你沒有選擇.
不過你可能會覺得不是吧.
我要麼就要愛他.

$^{441}$愛到為他捨命.
要麼就等於殺了他.
你要明白這是一種寫作手法.
看聖經從來都不是看一頁.
你要看很多不同的經文.
它有對應不同的場景.
例如有些人坦白說十惡不赦.
很討厭的.
真的你看到他前面討厭他後面的.
或者他真的做了很差的壞事.
我是不是要愛到他為他捨命呢?.
如何處理呢?.
耶穌有教的.
例如那個人犯了罪.
你要私底下跟他談.
談不攏的你找兩三個人跟他談.
再談不攏的就放上教會.
再談不攏的就可能要視他為外邦人.
聖經有不同的地方去解釋不同的脈絡.
所以今天說的.
你要記住是用一種言對立的極端方式去描述.
大家回去多看一些聖經.
如果我在這裡最後有什麼說話.
作為牧者要跟大家說.
我都是想跟大家說.
你看多一些聖經.
繼續看下去.
第十七節說.
凡有世上財物的.
看見弟兄缺乏卻關閉了惻隱的心.
神的愛怎能住在他裡面.
孩子們啊!我們相愛不要只在言語或是頭上.
總要以行為和真誠表現出來.
這裡就說了一個實際的例子去說.
當你在弟兄姊妹之間的那種關係.
你如何去到真的表達出你是愛他的呢?.
就是我們如果知道對方有一個實際的需要的時候.
你不可以去到當沒事當看不到.
你不可以當沒事當看不到.
雖然其實在教會裡面我們很少討論金錢.

$^{481}$我們會覺得錢是邪惡的.
但現實上我覺得錢是很忠誠的.
看你怎麼用而已.
你要知道雖然事實上很多罪惡的東西.
很多邪惡的東西是和錢有關.
但你做很多善事都和錢有關.
你沒有錢你怎麼去做好事呢?.
就算早前我曾經和不同的弟兄姊妹去油麻地派飯.
你派飯派物資你都要錢的.
你都要…就是那些東西不會憑空變出來.
你做好事你都要給錢.
最直接的方法都是給錢.
就算我在街上看見一個乞丐.
最直接的幫助可能都是給錢.
但問題是這個邪惡的世代對我們來說實在太多影響.
用乞丐的例子來說.
現在你給錢他的時候你都要想一想.
那個是不是大陸那些有些集團經營的.
集團式經營的帶他下來的.
有很多你一看就知道他不是香港人.
然後你給他會不會是便宜了他背後的利益集團呢?.
然後會不會其實害了他呢?.
他被人操控下來香港行黑這樣.
現在搞到你連做一件好事你都要去想的時候.
的而且確呢.
說到錢我們做基督徒原來真的要很小心.
我們做每一件事都要很小心.
大家在教會裡面盡量都不要有金錢上的瓜葛.
如果有需要的話透過教會去處理.
這個要提醒.
但是呢.
小心處理不代表我們什麼都不做.
或者關閉了我們的惻隱之心.
有句諺語就是說.
不要把嬰兒巾洗醋水一起倒掉.
就是說你幫寶寶洗澡的時候.
倒水的時候不要連寶寶也倒掉.
是不是?.
就是有些事情我們要小心.
但是不代表我們要關閉自己的惻隱之心.

$^{521}$經文裡面說到如果你有財物.
而你的弟兄有缺乏.
你是不可以視而不見.
不可以當沒有事情發生.
而且這裡說的是世上的財物.
世上說明了.
世上就是說你能不能帶走.
錢是你帶不走的.
這是很坦白的事情.
錢帶不走.
但是那個是你的弟兄.
就是你的弟兄姊妹.
將來會在天國見到的.
見到你的時候.
喂你啊?.
曾慧漢啊?.
你不是上次見到我沒飯吃都不借錢的那個.
你在天國見到他的時候多尷尬.
你要知道.
世上的金錢.
其實它帶不走的時候.
你是更加應該將你的焦點放在能夠將來在天國裡面有永恆價值的東西上.
當你的弟兄姊妹有需要的時候.
請你記住你是跟他一個共融的關係.
是共融的.
你跟他相交分享是有份的.
他的窮.
他的缺乏.
可能就是正正是需要等著你的幫助.
上帝放他在你眼前.
不是放在第二個眼前.
其實可能就是需要你.
另外就是.
其實幫助不一定是每一次都拿錢出來的.
你是有很多東西是可以去到實際的去幫助和支持一個人.
有時候我也是比較多事幹的人.
就是我老婆形容我是那種見到別人有問題就扯開衣服就衝出去的那種超人.
但是我又沒有這樣的能力.
但有時候實在看到有很多東西你是不能夠當沒事發生.
譬如我知道有一個弟兄失業了很久.

$^{561}$他找不到工作.
我無法只幫到他.
我就找一些做生意的弟兄姊妹.
我真的打去問他.
你請不請人.
我有一個弟兄真的幾個月沒工作搞不定.
這樣就是幫他領取.
如果說弟兄你找不到工作.
我為你祈禱.
祈完禱然後收線.
你是不是真的幫了他.
不是的.
其實我們真的需要有行為.
所以聖經說我們相愛不是言語和舌頭上.
是要以真誠的行為去表現出來.
有做和沒有做是差很遠的.
行為是反映你的內心.
反映你的信仰.
你有沒有真正的信仰其實就是看你怎麼做.
而且第十六節說到我們為弟兄寫命.
我想如果讓你選擇錢重要還是命重要.
如果有個賊拿著刀巷著你.
要錢要命.
那你當然要命.
但這裡說的是你要為你的弟兄寫命.
連命都可以擺出來的時候.
你世上擁有的財物金錢其實是不足掛齒.
你更加應該要運用你的弟兄姊妹的缺乏.
而去行動.
而去付出相對的努力.
當然不是說一下子就要做到的.
事實上也沒有人說一信耶穌.
他就能夠做到完全的和弟兄姊妹做到共融.
或者是能夠為弟兄姊妹去到寫命.
但其實是需要以這個為目標.
慢慢一步一步地讓自己去成長去成熟.
我很多時候都會面對很多困難和掙扎.
有很多事情去幫助一個人.
其實自己要捨棄一些東西.
或者是要犧牲一些東西.

$^{601}$到底自己願不願意呢.
自己要當自己要捨棄的時候.
很多時候都是靠祈禱.
靠上帝幫助.
然後我能夠去做到.
所以鼓勵大家弟兄姊妹.
你多些積極去參與事奉.
和參與一些幫助人的活動.
或者是你看多些別人的需要.
關心一下他.
見他板著臉的時候問一問他.
喂,你今天沒事吧?.
有什麼不開心的?.
可能就在這些過程裡面.
你就發現,你就看見他的需要.
然後你就可以幫助他.
我還記得我請大家.
去到開始崇拜.
臨崇拜之前有祈禱.
有分組祈禱.
這個是我過去的學習裡面.
一件很重要的事情.
你透過跟對方的分享,祈禱.
你知道他的需要.
你才可以有進一步的幫助.
祈禱只是第一步.
如果沒有了這些相交認識.
你可能一輩子都不知道他有什麼問題.
你可能一輩子都不知道他有什麼困難.
鼓勵大家.
透過建立Colonial團體的共融關係.
你是能夠繼續在神的國度裡面有份.
去建立神的國度.
我們同心去禱告.
七位天父我們多謝你.
多謝你呼召我們成為你的門徒.
去跟從你.
又多謝你賜下弟兄姊妹.
你讓我們可以透過.
跟弟兄姊妹之間的相交,共融.

$^{641}$我們真的見證到你在我們這個群體裡面.
你真的在我們這群人裡面去作供.
讓我們成為一家人.
成為真正的弟兄姊妹.
求你幫助我們.
讓我們不要關閉我們的猜念之心.
我們多些去看.
去看一看身邊的弟兄姊妹有什麼需要.
然後真真正正切實地去幫助.
好讓我們真的能夠實踐到你的教導.
我們也祈禱交託奉耶穌基督的名頭.
阿們.
\newpage



\section{}
\label{sec:_3DxJVmmiXg}
\textbf{2023-05-11 周四晚查經班《被擄歸回歷史研讀(14)》 講員:曾裔貴牧師}
\newline
\newline
連結: \href{https://youtube.com/watch?v=-3DxJVmmiXg}{\texttt{https://youtube.com/watch?v=-3DxJVmmiXg}} ~~~~ 語音日期: 2023-05-12
\newline
\newline
\hyperref[sec:eLIEUCC43fA]{\small{< < < PREV SERMON < < <}}
~
\hyperref[sec:index_chronic]{\small{[返順時目]}}
~
\hyperref[sec:index_scriptual]{\small{[返順卷目]}}
~
\hyperref[sec:A6GWJTLfDEo]{\small{> > > NEXT SERMON > > >}}
\newline
\newline
$^{1}$好,弟兄姊妹晚安.
我們又可以來到去研讀歸回的歷史書.
尼欺米記.
我們今天會看第六章.
我們一起祈禱.
天父多謝你.
多謝你讓我們可以每個星期四晚.
大家能夠用心去研讀聖經.
能夠從經文裡面去思想.
作為一位領袖.
在當時他要帶領重建聖城.
很盼望神都讓我們透過領袖的擺相.
我們從其中學習.
我們今天晚上等候在你的面前.
我們這樣禱告奉耶穌基督的名.
阿們.
是不是有些雜音.
讓我看看.
還是要靜音.
我們進入第六章.
又是重溫一點.
第三章你見到.
尼欺米他講不同的四圍的城牆的方位.
安置哪些猶太人.
哪些領袖.
怎樣做法.
第三章很集中地講.
他能夠動員有多少人可以協助他.
因為沒有人根本就做不到.
大家都認同那個事公.
一齊來到去.
重中之重的.
現在是要趕緊讓城牆和門能夠連絡好.
去到第七章就可以有事做了.
如果不是安息日.
就可以外邦人進來做普通的買賣.
但一到安息日.
這個民族.
這個神的選民.
就要守安息日.

$^{41}$就很嚴肅地暫時不可以買賣.
外邦人就進不來了.
所以第七章你會見到.
當城牆做好之後.
百姓就要來把關.
有些地方是不准外邦人進.
這樣就很清楚了.
你可想而知.
如果沒有了那個基本的基建.
很難做到要讓神的選民.
慢慢回到那個標準.
就是摩西的律法書的標準.
這個標準都不只是為了做而做.
而是這個民族.
是要孕育彌賽亞降臨.
所以你可以看到.
雖然尼希米未必很清楚.
何時耶穌基督.
這位彌賽亞會降生.
但是我們現在回望的時候.
神在歷史裡面.
就揀選了一位這麼重要的領袖.
他是一個外邦的酒政.
但是他有感動要回到.
一回到自然就是.
當你肩負起一個這麼大的工程.
各方面的指責.
挑剔.
就會四面八方湧來.
今天第六章我們就會看到.
外邦人他們就會來到去.
表面上很善意.
其實是一些陰謀.
我們說了第四章.
就是正正式式.
多比亞,參巴拉,基善,阿拉伯人.
來到去要攻擊這個工程.
很自然的.
他們是從外圍嘲笑.
潑冷水.

$^{81}$讓到神的選民.
自己來到去.
好像有一種自愧不如.
覺得工程很大.
就不能夠.
那個士氣.
能夠堅持下去.
這個時候就很需要領袖.
來到去有一個很堅強的意志.
去抵擋這些外面的攻擊.
就是對猶太人嘲笑侮辱的攻擊.
需要來到去擋了它.
到第五章.
全章就不說外邦人.
上個星期我們說了.
是說內部的問題.
其實一直都本身存在的.
就是神的選民.
原來都很窮困.
面對這樣的情況.
恥笑比掌底下.
有錢的貴袖,領袖.
就越來越有錢.
他們又借錢給猶太人.
又買他們的田地.
就是說慢慢慢慢.
那個一來一回.
那些普通的百姓.
就很大的生活困難.
再加上甚至連那些兒女.
都被賣為瀑被.
就更加怨聲載道.
所以這個不是重建才有的事情.
而是尼欽米去到之後.
一路建造.
有機會和當地的原居民去聊天.
去同工共事的時候.
就會慢慢浮現.
這些結構性的問題.
所以尼欽米都很生氣.

$^{121}$他就要來指出這些問題.
糾正這些問題.
就是要confirm那些領袖.
就是說你們不可以.
來借錢又要收取利息.
違反了摩西的律法書.
你應該與百姓共度時間.
來到困難的時候.
怎樣能夠幫助他.
例如免息借貸給他.
當他將來慢慢的生活改善.
或者可能要下一代.
改善之後他們就可以.
能夠有自己的謀生.
和不會越來越貧者越貧.
所以這個其實是聖經有這樣教導的.
讓社會不會完全貧富懸殊.
不是用一些極端的手段.
而是出於能夠感受到.
自己是神祝福的一個領袖.
可能是有錢人.
他能夠帶著一份憐憫的心.
一切都是上主給你的.
這個我們上個星期說過.
很重要的新教的倫理.
基督徒尤其是加爾文主義的.
簡選論的弟兄姊妹.
你怎麼知道你自己是蒙神簡選.
將來一定可以上天堂.
就是你在各方面都蒙神的賜福.
怎樣賜福?.
就是你的生活.
你慢慢勤奮.
你富裕.
你積蓄.
你不會亂花錢.
這樣你就能夠一直這樣自負.
你就繼續這樣幫助有需要的.
有些人他們可能是命不好.
或者是生活有些其他的天災人禍.

$^{161}$各方面一定會有的.
但是既然你是這麼蒙福的時候.
你就應該要來將那個恩典和別人分享.
這個是新教很重要的倫理.
新教舊教也有教導.
譬如保羅突然間閃過來.
在《迦太史》2章.
保羅猶太人的使徒彼得·約翰.
認識到神託付他.
我和大家揭曉.
怕大家不熟悉.
《迦太史》2章.
留意這段經文很特別.
《迦太史》的第二章.
過了十四年.
我和巴拿巴又上耶路撒冷去.
並帶著提多銅去.
我是奉啟示上去.
我在外邦人中所傳的福音.
對弟兄們陳說.
卻是背地裡對那有名望之人說的.
可能這個不是鬼鬼祟祟的背地裡.
而是很低調地去向耶路撒冷的領袖.
當然是以彼得·約翰作為最高的單位.
因為他們還是留守在耶路撒冷的使徒.
跟他們說在外邦神的福音.
有什麼可以改變到不同的外邦人.
建立教會等等.
唯恐我現在或從前徒然奔跑.
但與我同去的提多.
雖是希利尼人希臘裔的.
也沒有勉強他受割禮.
因為有偷著引進來的假弟兄.
私下窺探我們在基督耶穌裡的自由.
要叫我們作奴僕.
你就可以感受到.
保羅在這裡不是為了攻擊人.
針對哪個人.
而是他現在看到一件很嚴肅的事.
就是加拉太教會.

$^{201}$可能在基本的恩訊清義上.
加了一些額外的東西.
是不需要加的.
而加多了.
也可能有人影響他們.
例如你信耶穌是不夠的.
你應該要受割禮.
好像我們猶太人那樣.
就好像給加拉太人一個混淆.
以為恩訊清義是有些欠缺.
好像有些二等公民的感覺.
所以保羅就要捍衛了.
他就能夠舉到一些例子.
他當年帶領提多上耶路撒冷的時候.
提多是外邦人.
他沒有強迫他受割禮.
因為外邦人信耶穌.
受不受割禮是沒有問題的.
但反過來提摩泰.
他是猶太人.
因為他的媽媽是猶太人.
他的外婆也是猶太人.
他的爸爸是希臘人.
所以當保羅想帶提摩泰去傳道.
尤其是向猶太人傳道的時候.
猶太人就很關心.
這個提摩泰也是猶太裔的.
為什麼不受割禮呢?.
所以保羅就要提摩泰受割禮.
所以你看看他.
這個是考慮到人情世故.
不是一個跟恩訊清義.
來到去相違背的那種.
說就是這樣.
實際做不是這樣.
保羅認為不可能的.
不能夠接受的.
恩訊清義就是恩訊清義.
每個人.
不論你是什麼民族.

$^{241}$都是恩訊清義.
但是對於保羅要帶提摩泰去傳道.
他就覺得.
你應該要有一種這樣的禮儀.
受割禮.
不是因為你未信耶穌.
不是因為你受割禮才得救.
是因為你要在猶太的叔伯兄弟面前.
來證明.
你在他們眼中.
你都是循規蹈矩的.
但是和救恩無關的.
這是很大的學問.
值得我們真的要去傳道.
我們都要去.
透過保羅這些教誨.
我們都要有些揣摩.
有些東西原來都要講禮節.
但是不關救恩的事.
就沒什麼所謂.
但是關乎救恩.
就沒得講了.
沒情講的了.
所以很重要的那種拿捏.
這裡其實又拆開很遠了.
真理和傳統.
是兩個很大的.
要小心的.
真理可以是絕對.
但是傳統呢.
不一定是絕對的.
所以有時慢慢我們做慣了某些的東西.
那個慢慢變成一個教會的傳統.
不同的教會有不同的傳統.
那個傳統不是真理.
所以我們要搞清楚.
我不想拆得太遠.
尤其是講到關於第五章的尼希米記.
講到對窮人的問題.
是涉及到整個內部的心靈.

$^{281}$那個童心.
如果解決不了.
其實尼希米.
整好這些很漂亮的城牆.
塗上油.
很漂亮.
裝飾完都沒用.
因為百姓的心根本就不服.
根本就不會童心.
根本上做不到事.
所以童心是最重要.
怎樣能夠令到那個群體童心呢.
這個也可以這樣說.
第五章.
尼希米.
我們上個星期看過了.
他真是以身作則.
將自己的俸祿.
全部拿出來.
或者說.
不像以前上一任的省長那樣.
他都是被封為省長.
這個河西的省長.
當然是接管猶太人耶路撒冷.
因為他是本身猶太人.
他就不像上一任的省長那樣.
加重百姓的稅收的勞苦.
所以他自己就不要省長的俸祿.
你就已經釋放了一個很重要的訊息.
他是和百姓一起面對困難的.
所以他不是高高在上.
這樣做一個領袖的.
所以親民是相當重要的.
第六節.
至於那些有名望的.
不論他是何等人.
都與我無關.
似乎若有所指.
說的是那些人自命清高.
自命是很淑齡等等.

$^{321}$保羅說如果不是涉及真理救恩的.
我都不需要和他們有什麼來往.
神不以外貌取人.
所以大家要記得.
那些有名望的並沒有家憎我什麼.
反倒看見了主.
託我全福音給那未受國禮的人.
那些有名望的.
看到保羅這麼多年的默默耕耘.
他一心一意.
他不是特意要和猶太人為敵.
他自己是一個拉比.
他自己是一個很高的法利賽的教法師.
受於名門加瑪列的門下.
但是他被猶太人終於認同他.
他向外邦人傳福音.
所以原來需要一段時間.
讓人去認識你的內心.
那個初心.
你才能夠讓那些反對的人無話可說.
所以這裡告訴我們.
看到他是神的託父.
他向外邦人傳福音.
第七節的下半節.
正如託彼得傳福音給那受國禮的人.
所以這裡你會看到.
原來不同的僕人千萬不要比較.
彼得是向猶太人受國禮的傳福音.
他有他自己特定的神的託父的工場.
那個牧羊的群體.
但是保羅.
其實他絕對有資格在猶太人當中做的.
但是神要他遠遠的去到整個歐洲.
去向外邦人傳福音.
所以現在終於被猶太人.
特別是一些有影響力的猶太人.
來去所印證和看出是神託父他.
所以如果我們作為傳道者.
能夠被弟兄姊妹看出是神的託父.
這個是很重要的.

$^{361}$相信都是很多年很多年的默默的牧羊.
來去封住弟兄姊妹的心靈.
不是只有一句話.
是真的要默默地去做.
所以現在人們真的認同保羅.
是向外邦人傳福音的神的揀選.
差派和託父.
這是很高的身份.
第八節.
那感動彼得叫他為受割禮之人作使徒的也感動我.
叫我為外邦人作使徒.
又知道所賜給我的恩典.
那稱為教會儲石的雅各.
姬法若翰.
雅各不是若翰的弟弟.
這個雅各其實是耶穌的肉身的弟弟.
因為那個雅各已經死了.
所以這些告訴我們.
教會的儲石.
在耶路撒冷裡面.
這些使徒為主耶穌作生命的見證.
所以所謂儲石.
不要以為坐在一個王的寶座上.
誰都不敢動你.
其實當時是逼迫時期.
你說跟隨耶穌.
你就快要被人斬頭.
尼祿將彼得.
將約翰終老而死.
彼得很快就要釘十字架.
所以所謂教會儲石.
是保羅自己對他們的一個尊重.
在外邦人的眼中.
誰理你是儲石.
誰理你是教會領袖.
你就更加要首當其衝.
去面對那個逼迫.
面對那個羞辱.
面對你這麼喜歡跟隨耶穌.
好呀,耶穌釘十字架.

$^{401}$那你們一起吧.
彼得就是這樣死了.
這裡告訴我們.
留意這裡.
我讀了這麼多.
我想說這裡.
就向我和巴拿巴.
用右手行相交之禮.
叫我們往外邦人那裡去.
他們往受國禮的人那裡去.
所以這裡都是很值得我們去細味.
使徒們和保羅行右手相交之禮.
到底是怎樣呢?.
是不是握手呢?.
還是怎樣打工作泣呢?.
我們又不是太知道.
但是這個相交之禮.
是保羅的眼中體系很嚴肅的.
就是說猶太的領袖.
這些教會的儲石.
認同保羅和巴拿巴.
有神的託付.
向外邦人傳道.
所以某程度上.
是好像他們去開綠燈.
就是說由他們的口.
再來到去耶路撒冷.
替保羅去背書去宣傳.
就是說這兩位使徒.
是有神的印證.
他們向外邦人傳道的.
所以猶太的弟兄姊妹.
千萬不要攔阻他們.
他們在做的事情.
未必和我們所做的.
有一點完全一樣.
可能傳道的方式不一樣.
因為人家是外邦人.
完全沒有猶太歷史背景.
完全沒有舊約的禮儀背景.

$^{441}$對於什麼叫憲制.
對於什麼叫聖殿.
他們是完全不懂的.
外邦人根本就不會有那些飲食條例.
魚蝦蟹是照吃的.
猶太人就不行了.
所以他們一定要給到一個空間.
保羅,巴拿巴.
所以這個相交之禮.
就是說一個很正式的.
認同保羅和巴拿巴所做的.
就是印證著神有主權.
差遣不同的僕人.
做不同的聖功.
所以我們就要有一個很闊的眼光.
去看不同的牧者.
在不同的工場所做的.
當然如果是違反真理.
就當然不行了.
因為做法不同.
要包容.
胸襟要夠闊.
你未必認同.
甚至你未必要這樣做.
譬如用一些不是太過敏感的例子.
譬如林以樂牧師在紅館.
開洞篤的講座,報道會.
你可以不認同他.
欺少勞猛.
但他背後承載著一些未信的朋友.
聽得懂他的語言.
能夠聽得明,聽得入心.
又覺得你們的耶穌多好.
又何妨呢.
我以前就真的不接受.
但是如果慢慢對這位神的僕人.
人家不是亂來的.
人家做的時候.
雖然表面上就好像.
有很多的.

$^{481}$和一般的報道會不同.
但是他有神的託付.
你不要這樣學就行了.
你學也不應該.
因為神沒有這樣託付你.
你經常在講台講笑話.
這樣就沒有意思了.
神不是叫你講笑話.
叫你講道.
所以每個人都有不同的託付.
我們要有一種這樣的相交之禮.
這樣的一個胸襟.
去對不同的牧者.
記住他的背後是有神.
他的內心有神的帶領.
我們就不要指指點點.
否則很容易就變成一種批判.
和一種論斷.
是不好的.
這個大家都要學習的.
接著就留意下面第十節.
最重要這一節.
只是願意我們紀念窮人.
這也是我本來熱心去走的.
很特別突然間.
講一傳福音.
彼此來接納.
彼此來欣賞.
彼此來祝福.
大家去做神要託付你去做的.
這次使徒他們很焊傷.
阿保羅要紀念窮人.
阿保羅因為他在這裡寫他的自傳.
寫他自己對於真理的捍衛.
他就close了兩個案例.
來support他在第三章.
應該第二章的中間已經講了.
第三章就最詳細講了.
因順清義就可以得到聖靈.
無需要做任何的手續.

$^{521}$包括行國禮.
好行為.
不需要.
是憑信心.
接受耶穌.
你就可以因順得著所應許的聖靈.
第三章就已經這樣講了.
他就用一些案例.
提多.
我都沒有勉強的行國禮.
即是說.
信耶穌的外邦人是不需要行國禮.
意思就是這樣.
甚至他做到一個地步.
連彼得.
雅各.
約翰都來去印證.
神真的託付他和巴拿巴.
向外邦人傳道.
重點就是說.
很多人以為保羅在外邦教猶太人.
可以不需要遵守摩西的律法.
這個《小圖行傳》21章.
那裡這樣去講的.
所以拆開.
遲些我們11月去這個土耳其.
我們就真的去到保羅的足跡了.
都是我算相當熟的《小圖行傳》的經文.
到時可以有很多的講解.
給大家來認識.
回到這裡就是.
突然之間出現了這一句這樣的說話.
《小圖》很想保羅.
就算去到外邦.
都要顧念那些窮人.
保羅就自己在這裡.
加了一個他自己的注釋.
其實都是我一直所關心的.
熱心去走的.
就是說一個敬虔的猶太人.

$^{561}$他熟讀摩西的律法書.
他就會知道.
我上個星期引述過.
《申命記》的.
我應該是看回是不是15章.
就應該會知道.
在律法書裡面.
神的那個心意.
就是說如果一個.
跟隨主耶穌的人.
我們都要讀舊約.
新約.
我們都要來到去慢慢發現.
我們如果真的越來越成熟.
你的心底裡.
是應該有對貧窮的理解.
對貧窮人的那個惻隱之心.
對貧窮這件事.
我們要.
在我們人生.
盡我們的力量去幫助.
15章沒錯在這裡.
第七節.
我不能夠成章地說.
在耶和華你神所賜你的地上.
無論那一座城.
你弟兄中如果有一個窮人.
你不可以忍著心.
這個叫.
可能是那個舊字還是什麼.
我要回去查一查這個字.
這個字很深.
就是收起你的手.
不幫補你窮乏的弟兄.
總要向他鬆開手.
照他所缺乏的.
借給他.
借.
補他的不足.
你要謹慎.

$^{601}$不可以心裡起惡念.
一開始就見到.
原來心裡出了一個問題.
就是.
很多事情要想.
第七年的豁免年快到了.
你變惡眼看你窮乏的弟兄.
什麼都不給他.
以至他因你.
求告耶和華.
罪便歸於你了.
你總要給他.
給他的時候.
心裡不可受煩.
因為耶和華你的神.
必在你這一切所行的事情上.
並你手裡所辦的事情上.
賜福於你.
就是暗中的父.
在暗中會賞賜你.
第十一節.
所以如果我們能夠對舊約.
摩西的律法書.
能夠有一些掌握.
認識.
原來那地上的窮人.
是永遠都不會斷絕的.
所以我吩咐你說.
總要向你地上困苦窮乏的弟兄.
鬆開手.
下面就說對待奴隸.
我們就不再詳細說了.
說了兩次鬆開手.
和.
似乎告訴我們.
我們是大地的管家.
我們是神託管我們去管理.
我們的一生.
我們的衣食住行.
我們自己和我們的財富.

$^{641}$那就是說.
如果我們能夠見得到.
神的感動在我們身上.
我們去幫一個惡念.
心靈.
和一個感動一線之隔.
所以一個成熟的基督徒.
他應該是那個心靈是柔軟.
柔軟不是愚蠢的賴衣褲.
柔軟是心靈的柔軟.
但是那個頭腦是要很複雜.
因為世情是複雜.
所以我們心靈要柔軟.
我們要有一種隨時.
我們是神的管家去幫助.
因為神自然就會再將更多的.
來給我們去做那個行善.
所以舊約新約.
一而再這樣提我們.
要對這些貧窮要上心.
所以回到這裡的時候.
其實就是上個禮拜我們還沒說夠多的.
就是第五章的經文.
這裡就已經告訴我們.
黎希米就看到.
我們先說那個結構.
一到第六章是一個自轉式的自述.
第二章正正式式.
由波斯的書山城.
上到耶路撒冷.
開始就視察.
有幾個隨從去視察.
視察完之後.
第三章就焦聚所有的.
首先就是跟那些領袖說.
接著由領袖去焦聚群眾.
大家在自己的房屋附近.
這樣來修理建造.
可能有些人就搬運材料等等.
大家在自己城門附近.

$^{681}$來修造.
既減省旅行費.
但也可以對於自己的地方有感情.
有熟悉的地形.
當然有些公投.
負責在其中調度人手.
第三章就是這樣.
第四章就是參巴拉的介入.
代表著當時原居民的外邦人.
已經跟耶路撒冷的人同住.
耶路撒冷就是沒有染雞籠.
他們可以隨便進來.
當然真正要去守節.
去聖殿.
他們未必會來.
未必有興趣.
但是因為耶路撒冷的人.
已經沒有那種力量.
去抗衡世俗化的影響.
所以聖殿的敬拜.
變得可有可無.
有時變成繁文縟節.
變成一個阻礙.
變成去不去問題不大.
變成沒有那種敏銳.
童宮有時也會說.
那些小朋友上兒校.
大考了.
叫兒子千萬不要上學.
要溫習.
由信主的父母.
叫兒子不要上教會.
香港不只是我們自己一間.
每間教會都是這樣普遍.
不爭取一兩次.
不要緊.
有時價值觀的問題.
一兩次反而不是問題.
但是價值觀灌輸.
有重要的東西.

$^{721}$上教會就變成不方便.
過多小時.
溫習就拿不到A.
有時有錯誤的價值觀.
慢慢就潛移默化.
自小的培養.
就帶來.
現在信耶穌還要死守戒條.
真的不行.
怎麼辦.
變成反過來.
牧者不對.
你迫小朋友上教會.
溫不了書怎麼辦.
真是大考.
很多這些就變成.
次序的問題.
優次的問題.
有時建立了一種.
個人對神的敬虔.
敬畏神的心.
這些是老生常談的事.
你想想.
你能夠在他的心靈裡.
建立了一個種子.
可能對他來說.
小時候.
什麼這麼律法主義.
什麼這麼形式主義.
真的很不舒服.
你都要請我上教會.
真是這樣.
當然如果真的有大病.
當然不要.
那是另一個問題.
只不過是一個所謂不方便.
那不方便去到什麼程度.
會不會用另一種價值觀.
我們不如先上教會.
反正溫書溫到這麼悶.

$^{761}$照常上教會.
完了之後.
簡單吃點東西.
父母陪你.
我們不要逛街.
繼續溫書等等.
其實都可以的.
我認識一位傳道人.
真是很好.
他不是守律法.
他自己.
禮拜日.
他什麼都不做.
就是上教會.
事奉.
其實很忙的.
他應該是在神學院裡面.
或者說他.
對不起再早一點.
他讀大學的時候.
他覺得.
星期日是應該要給神的.
我讀了六日書.
他在港大.
後來就是做傳道.
後來就是神學院的院長.
很深刻的.
是他自己親口說的.
他讀大學的時候.
他認為.
禮拜日為什麼不給神呢.
不把時間給神.
其實在教會聚會.
和教主一學等等.
他覺得.
六日已經很忙了.
為什麼禮拜日都要這麼忙呢.
他就將自己.
不是要死守那些戒條.
而是星期日是要敬拜神的.

$^{801}$他真的將那天分別出來.
又不是讀書差了.
一樣很好.
很厲害.
所以有時.
看看你的價值觀.
怎樣能夠在小朋友的內心.
種下這個內心的種子.
到將來有一天.
連他自己都知道.
我敬拜神是優先的.
無價之寶.
你也不得不將來他是這樣.
他不需要你扯他回去.
帶他回去.
他自己回教會.
其實就是那個種子.
那你就功成身退.
他將來自己會敬畏神.
你想想.
是真的要潛移默化的.
是真的要慢慢.
每個禮拜這樣來到.
由他不想.
很睏很累.
你都要強行扯醒他.
很不人道.
我六天都已經上學這麼辛苦了.
禮拜你都不讓我睡晚一點.
是沒錯.
六天都辛苦.
不過回教會.
我們敬拜神很開心的.
看你怎樣去演繹那件事.
慢慢由他小時候.
一直長大.
他能夠如果建立了.
這一種的態度.
沒有人可以再奪去的那個態度.
那你就是真的.

$^{841}$做了你做管家的角色.
如果不是.
教會原來.
有事就可以不用回教會.
端乎看你那件事是重大還是不重大.
那他就會分.
我今天很不開心.
跟男朋友吵架.
我不回教會了.
有個大考.
我又不回教會了.
那他到大學的時候.
他去到外國的時候.
其實他自己能夠去拿捏那些時間的時候.
那他就自己會分.
什麼是大事小事.
回教會原來是小事.
那其實就不好.
你可以想像得到.
這個就是.
為什麼我們要留意.
其實我剛才說遠了.
就是猶太人.
在沒有離希米上來的時候.
對於敬拜已經是恩秦了.
所以去到第七章.
成祥喜好之後.
由第七章去到第十三章.
是開始說回敬拜.
開始說回一些嚴肅會.
開始說回七月初一的贖罪日.
對不起.
催國節.
跟著一路延伸去到贖罪日.
跟著再延伸去到鑄盤節.
其實是七月.
其實已六月就是六月.
第六章的那個已六月.
就是六月.
但去到第七章.

$^{881}$就是七月.
整個月裡面.
其實有很多的聚會.
很多的節期.
由第七章一路延伸去到十三章.
即是說還有一半我們都未讀.
就是在七月裡面發生的.
那些敬拜的生活開始有東西看了.
離希米開始跟他們一起.
當然加上以色列.
一起來去研讀摩西的律法書.
大家就被神的道再一次感動了.
那就開始有自己民族的尊嚴了.
就是守摩西的律法書.
所以你可想而知.
它的結構都是幾工整的.
第三章全民皆兵.
第四章外面的攻擊.
接踵而來.
是由一些嘲笑侮辱.
即是一些難聽的說話.
來讓百姓.
甚至離希米知難而退.
不敢強硬攻擊.
因為真的有王的詔書.
有波斯書山城的文件.
真的准他們重建的.
可以來拿些材料的.
在河西那裡.
但問題就是說.
當地的土著.
即是原居民.
是會諸多的阻撓的.
因為他們看到.
如果建好城牆.
就是說我們就會被排除在聖殿以外.
我們真的變成真正的外邦人.
就不可以跟猶太人.
有一個繼續很密切的來往.
所以離希米要建造城牆.

$^{921}$一定要將世俗和敬虔.
聖殿分開.
不好意思.
我先聽聽誰開著.
這樣來做一個界線.
所以由第七章到十三章.
就開始說猶太人.
神的選民.
他們要有那個敬虔的界線.
不是城牆的四邊.
是心靈裡面.
讓神的道在他們的心靈裡面.
第五章就是處理了.
內部其實結構性的問題.
就是有嚴重的貧富懸殊.
有百姓已經是窮的了.
還要我們做這些苦差.
那就更加苦了.
就向離希米來反映和投訴.
甚至乎他們的太太.
是大大的呼號.
埋怨他們的弟兄猶太人.
所以這個就是離希米都要處理.
他怎樣處理就上個星期我們說了一點.
這個是第五章的十四節.
離希米的五章十四節.
自從我奉派作猶大地的省長.
就是從亞達斯西王二十年.
直到三十二年共十二年之久.
我與我弟兄都沒有吃省長的俸祿.
在我以前的省長.
加重了百姓的擔.
每日索要的非常之多.
所以呢.
離希米他一而再強調.
但是因為我敬畏神.
不者讓行.
並且我恆心修造城牆.
並沒有置買田地.
我的僕人也都聚集在那裡作供.

$^{961}$除了從四圍外邦中來的猶大人以外.
有猶大平民和官長.
一百五十人.
即是那些可能是公頭.
在我席上吃飯.
每日預備一隻公牛六隻肥羊.
又預備一些飛禽.
即是一些可能是.
班溝 鋤甲 諸如此類.
每十日一次.
多預備各樣的酒.
雖然如此.
我並不要省長的俸祿.
兩次說了.
因為百姓服的役甚重.
即是說.
真的要有那個意識的共度時間.
上來呢.
不是一種居高臨下.
即是高高在上.
一個領袖.
應該是捲起衣袖.
和百姓面對一起的困難.
應該是感受到那種困難.
所以牧者傳道人.
應該是要.
心靈時時都要有一種敏銳.
時時都要有一種痛.
你感覺到弟兄姊妹的痛.
你才能夠真正去做到牧羊.
雖然你未必是有切身處地.
譬如有些.
很多年前.
一位和我差不多大的姊妹.
很多年前.
我的母會.
小產.
她.
看著她整個人.
低得很厲害.

$^{1001}$我是男性.
那時未結婚.
怎會明白這些事情呢.
她的感受是很痛.
即是你可想而知.
如果你是牧者.
你要安慰她.
你真的要感受到那種痛.
譬如.
如果我是男性.
比較容易明白的.
一個弟兄突然被解僱裁員.
突然失去工作.
那個打擊.
對一個弟兄來說.
即是男士.
特別是香港人.
都是男主外.
是一件很大的事情.
你試過站在地上.
那個地好像浮沙一樣.
我相信就有這種感覺.
我不是說過.
我被教會解僱.
我由深水Po 搭68X回到元朗.
我真的感受到.
我下面好像浮沙一樣.
我坐在車上.
我不是有幻覺.
當我打電話給太太的時候.
她一直哭.
我跟著也哭.
你無法想像那種.
那種突然間的.
一下子.
你去一個深淵.
所以牧者要有這種.
感同身受的痛.
對別人.
對牧羊的群羊.

$^{1041}$要有這種.
你真的要陪伴.
面對.
等等等等.
所以有時我們牧者.
為什麼有時會給一些空間.
他突然間要站在急症室.
或者在病房外.
站幾個小時是閒閒的事.
所以童工我們是彈性時間.
是需要的.
難道說對不起.
五點鐘你在ICU.
我明天來探你.
這樣是不行的.
所以其實牧者真的有這種.
要感同身受.
要陪伴.
頂尖會面對不同的人生的事情.
和遭遇.
如果不是.
有時我們沒有那種傷痛的感覺.
第五章就是.
我的神求你紀念.
我為這百姓所行的一切事.
施恩與我.
所以記得是在神的裡面.
蒙神紀念.
雖然我們覺得.
什麼這麼演戲.
說出來.
但因為這是一個回憶錄.
和神都感動他.
是一個很率直.
很真誠的一個領袖.
他可以來和他的會眾分享.
大家都知道的.
紀開誠佈公的.
不是他私下做這樣的事.
第六章.

$^{1081}$又有另外一種.
結構性的攻擊.
是用一些很陰險的手段.
而且還有一些是猶太的領袖.
你就看到.
原來牧者.
是一個重大的.
時公的一個領袖.
就要有這樣的心理準備.
新約我都引述了很多.
等一下和大家可以揭一揭.
第六章的第一節.
都是一個引言.
我們說了九個字.
參巴拉多比亞.
阿拉伯人基善.
和我們其餘的仇敵.
聽見我們已經收完了城牆.
這是第十五節.
以六月.
猶太年曆的六月.
六月到七月.
是一個準備期.
所以六月很多時候.
是一個悔改的準備.
因為七月初一是摧國.
七月初十是贖罪日.
七月十五是很開心的鑄棚節.
所以是一連串的節期.
所以六月收好了.
都是有很重要的意義.
所以你會看到.
第七章立即說.
七月的節期開始.
我們先看這裡.
第二節.
參巴拉和基善.
打發人來見我.
說請你來.
我們在阿羅平原.

$^{1121}$一個村莊裡相會.
他們卻想害我.
於是我差遣人.
去見他們說.
我現在辦理大功.
不能下去.
焉能停工下去見你們呢.
他們這樣四次打發人來見我.
我都如此回答他們.
參巴拉第五次.
打發僕人來見我.
手裡拿著未封的信.
信上寫著說.
外邦人中有風聲.
加斯姆也說.
你和猶大人謀反.
很嚴重的指控.
收造城牆.
這件事是對的.
不過將這件事扭曲.
你要作他們的王.
假借王的名.
原來上來是想自立為王.
描繪得似塵塵.
真是這樣結集這麼多猶太人.
建好城牆之後.
是不是圖謀不軌呢.
第七節.
你又派先知在耶路撒冷.
指著李宣說.
軍權神授.
即是說有先知來講預言.
李希米是我們一個領袖.
他現在很值得.
更加有神的引導.
我們受感動.
他會是我們的王.
由先知的口講出來.
即是說在這裡.
找一些的妹妹.

$^{1161}$透過先知的口.
可想而知.
參巴拉和基善.
他們是處心積慮.
要用一些陰謀.
去攻擊底位.
和毀謗李希米的名聲.
希望用一些很強硬的陰謀.
去拉他落馬.
第七節的下半節.
說在猶大有王.
現在這話必傳於王之.
所以請你來.
與我們彼此商議.
即是說我們現在幫你.
怎樣可以禁止這些謠言.
這些由先知的口講出來的.
其實可能不關你的事.
你現在和我們去商議.
因為我們都是民間的領袖.
我們都可能幫你.
在亞特薩西王的面前.
或者他的那個.
鑫差大臣來到的時候.
我們幫你說好話.
第八節.
我就差遣人去見他說.
你所說的這事一概沒有.
是你心裡捏造的.
他們都要使到我們懼怕.
二司說他們的手必軟弱.
以至工作不能成就.
神啊求你堅固我的手.
我到了米希大別的孫子.
蒂萊瓦的兒子.
是瑪雅的家女.
那時他閉門不出.
似乎是領袖來的.
是猶太的領袖.
他說我們不如在神的殿裡會面.

$^{1201}$將殿門關鎖.
因為他們要來殺你.
就是夜裡來殺你.
要將尼希米窩藏在聖殿裡.
不是敬拜的時間.
是晚上將他藏在聖殿裡.
聖殿外邦人不可以進.
猶太人固然都不可以進.
聖殿裡只有祭司可以進.
所以都是一個很大的叛逆.
要尼希米愉悅了自己的身份.
沒錯他是亞達斯西王委派的猶太省長.
記住這是政治上的安排.
是由王安排你.
但是祭司是神安排的崗位.
你是一個領袖.
走政猶太省長.
但你不是祭司.
你如果進聖殿就糟糕了.
你就是神都不怕.
人更加不怕.
你就不知不覺間.
跌入了這個圈套裡.
由一個領袖.
似乎都是祭司的職級領袖.
這樣來勾結了這些外部勢力.
參巴拉 多比亞 阿拉伯人 基善.
這樣來同謀.
要來陷害尼希米.
所以都要很小心.
將電門關鎖之後.
因為他們要來殺你.
第十一節.
我說.
像我這樣的人豈要逃跑呢.
像我這樣的人.
其實很重要的.
沒有他.
做不到整個聖殿的城牆的修復.
就不會再有第七章之後的聖文的重建.

$^{1241}$是心靈的重建.
像我這樣的人豈能進入殿裡面保存生命呢.
我不進去.
第十二節.
很敏銳的思想.
我看明神沒有差遣他.
是他自己說話攻擊我.
是多比亞和參巴拉賄賣了他.
賄賂.
不知道有什麼好處.
那些延後的利益.
賄賣他的緣故.
是要叫我懼怕.
依從他凡事.
他們很全然惡言.
毀謗我.
先是善意.
當你下了答.
就是惡意.
你想想是一個陰謀.
我們先讀完.
再引一引新約.
主耶穌都被人這樣攻擊和毀謗.
第十四節.
我的神多比亞參巴拉.
呂先知羅亞底.
講先知的預言.
說李希米作王.
神感動我.
告訴大家李希米做王.
呂先知羅亞底.
和其餘的先知.
都要叫我懼怕.
求你紀念他們所行的這些事.
每一次講完一些案例.
李希米就祈禱.
在神面前禱告.
求主建立他.
一方面給他很強的意志.
更重要是要有一個很清晰的腦袋.

$^{1281}$不要被驚恐.
被一些似是而非的善意.
捉了你下去.
我自己有一個很深刻的反省.
一定不能說是人.
因為在這個教會發生.
很多年了.
因為我在這裡十多年了.
我想都有六七年前.
有一個姐妹.
我完全不能說名字.
很深刻的.
跟我說了一些事.
我就告訴她某些人某些事.
突然之間她一句說話.
令我有警惕.
突然說Bingo.
之後我跟張牧師說.
這個姐妹.
張牧師馬上提點我.
這個人你要小心.
你要慢慢觀察她.
我自己當下跟她對話完.
我絕對沒有意思要挖苦誰.
她突然間.
我說了一大堆事之後.
她說Bingo.
我就有一個警戒.
突然內心有一個警戒.
為什麼她這樣說話呢.
我就回去跟張牧師聊天.
她叫我留心這個姐妹.
叫我留心.
這樣我就知道.
我日後當然再跟她有不同的事工.
也知道多一點她的為人.
就是這樣.
不是想自相貌壞.
而是我想說.
我也要很小心.

$^{1321}$看不同的事情.
不同的人.
你起碼真真正正面對.
在她這麼大的事工.
所以你看到那個結構.
也挺工整的.
她真的很有思路.
第三章她知道要怎麼做了.
就找了那些人.
那些人也無所謂.
我們也真的想幫忙.
但是那些困難來.
就考驗那個人.
到底是不是真正忠心的.
在神那裡.
就開始那些困難出現了.
第四章就集中說.
從一個外面的嘲笑.
侮辱.
言語的攻擊.
對負責的人.
辦事的人.
諸多的攻擊.
希望知難而退.
這個地盤耶路撒冷.
是我們阿拉伯人已經盤踞了.
你們有聖殿已經.
掩著嘴笑了.
不要再建好城牆了.
就是說他們也意識到.
如果建好整個耶路撒冷的舊城.
那就是說他們外邦人.
就有一個地方.
好像城寨那樣.
三不管.
外面的人進不來.
沒錯的確是這樣的.
但是黎希敏不是想跟他們分別.
而是想自己的民族.
神的選民.

$^{1361}$有那個聖潔的國度.
是屬神的範疇.
所以當然了.
自然就是.
你如果要這樣的時候.
外邦人就會.
很多事情不方便了.
第七章就是講家譜.
第八章就開始說.
安息日就不准那些人來買賣.
就是說.
開始要有一個界線.
所以外邦人就不希望有界線.
所以第四章就是用一個.
這樣的言語的攻擊.
第五章就真的出現了內部的問題了.
內部的那個不同心.
是因為可能真是天災.
可能真是恥小比長的貧富懸殊.
出現了那個.
不對稱的分佈的財富.
就令到耶路撒冷的一些人.
現在參與在工作上面的.
就直接向黎希米反映那個現實.
投訴那個情況.
就是說那個奴役很重.
黎希米就用他的那個影響力.
當然根據聖經的律法書.
來去直斥那些領袖.
你借給那些猶大人.
是一份的善意.
是一份的愛心.
讓他能夠慢慢慢慢可以自力更生.
所以這個聖經都是很有智慧的.
可以借,但是就不要收利息.
讓他慢慢自力更生.
他就可以來去改善他自己的生活.
就不只是一種施捨.
施捨只會製造一種貪婪.
一種不勞而獲.

$^{1401}$那就不健康的.
所以第五章就真的告訴我們.
要看一個領袖他自己怎樣能夠面對這些困難.
他都真的迎向那個困難.
和他自己都讓弟兄姊妹看到.
他自己都是這樣來去.
一起面對那個難處.
他不要那個俸祿.
未必代表他很富貴.
但是真的代表他真的看到那個需要.
所以領袖是應該時時先天下之憂而憂是對的.
後天下之樂而樂也是對的.
第六章就是告訴我們.
有好幾個範疇的攻擊.
寫信原來已經有內線的人.
是耶路撒冷的領袖.
是一些甚至是祭司.
又起了名字,有呂先知羅亞底.
這些原來已經和外邦人千絲萬縷的關係.
你上到去你就知道原來不是那麼簡單.
為什麼不是欠缺領袖去重建.
不是欠缺那些材料.
也都不是欠缺百姓的決心.
是有很多的事情.
是令到那些重建的重大計劃付之於流水.
是因為太多更多背後的利益的瓜葛.
就是外邦人根本不想你做成這件事.
所以他們就會找一些可能對真理一知半解的.
來籠絡了他.
告訴我們原來有賄賂的.
有一些是陰謀的.
要這個善晚雅.
這個似乎是祭司.
要黎希米進入聖殿裡面去窩藏.
躲起來.
其實可以引聖經的.
因為在無論是列王記上.
或者是歷代志下.
有一個君王約阿斯.
小時候他都是藏在聖殿裡面的.

$^{1441}$當時因為是一個很危險的時候.
當時是耶誓別做王.
對不起.
耶誓別的女兒亞他利亞做女王.
猶大的女王.
要誅殺猶大的太子等等.
那些王子.
所以約阿斯的舅父耶和耶大.
就將這個外甥約阿斯.
收藏在聖殿.
令到亞他利亞找不到.
所以收藏了六年還是七年.
你可以想而知.
其實有案例的.
但是黎希米知道.
現在不是這樣窩藏的時候.
現在是他給人一個.
很陰險的一個局.
水勢都不清.
你是否恃著自己是亞達舍西王.
委派你的省長.
你不是猶大的祭司的身份.
你就進到聖殿裡面.
你就可以想而知.
你進去之後.
人家就會說你壞話.
人家就會唱你.
所以我們逐段都可以去看看這裡.
當我們看這裡之前.
和大家看一些新約.
福音書大家可能.
有時這樣聽完未必會那麼賞心.
但現在我們特定的主題.
就是一些陰險的說話.
我們就知道原來福音書一樣有.
主耶穌那麼屬靈.
一樣會被人這樣去唱.
被人這樣誣蔑.
我們先看馬太福音第九章.
九章的.

$^{1481}$由第九節開始.
耶穌從那裡往前走.
見到一個人名叫馬太.
是一位學者.
坐在稅關上.
即是他正在收稅.
即是他是稅吏.
對他說你跟從我來.
他就起來跟從了耶穌.
接著第十節.
耶穌在屋裡坐直的時候.
當然就是這個稅吏馬太.
宴請他的同行.
他的朋友他的家人.
來招呼耶穌.
即是公開宣告.
我現在要跟隨耶穌了.
因為耶穌呼召我.
他就不再做稅吏這個職業了.
有很多背景大家都聽過.
稅吏當時是替羅馬人打工.
收了基本的羅馬人要求的稅之後.
他自己就會另外再收猶太人其他的稅.
即是說自己就要來到中保斯郎.
亦都是羅馬人默許的.
所以作為一個稅吏.
其實就是等於一個奸細.
欺壓自己的百姓同胞.
所以當時稅吏和罪人是統稱.
即是說猶太人都知道在耶穌的時代.
做的稅吏不是稅局的人這麼簡單.
做稅吏其實是畫爛臉做惡事的了.
然後告訴我們有很多稅吏和罪人.
與耶穌和他的門徒一同坐直.
留意一下11節.
法利賽人到處找耶穌的錯處把柄.
法利賽人看見.
留意一下.
就對耶穌的門徒這樣說.
不是跟耶穌說.

$^{1521}$是跟耶穌的門徒說.
即是說希望說這些說話給門徒.
你們自己去思想一下.
你的老師是這樣的.
和那些稅吏罪人吃飯的.
你還跟他?.
類似這樣.
其實是用一些好像很婉轉的說話.
其實是背後攻擊和毀謗耶穌的.
不是直接跟耶穌說.
可以直接跟耶穌說.
這些罪人他們做了多少壞事.
你現在為什麼要跟他們同流合污呢?.
可以直接問耶穌.
不問問他的門徒.
給一些懷疑給一些輕視侮辱.
給一些門徒心裡面對耶穌的觀感.
重新有另一個看法.
就對耶穌的門徒說.
你們的老師為什麼和稅吏罪人一同吃飯呢?.
很特別就是第十二節.
耶穌聽見.
就說.
其實不是跟耶穌說.
你可想而知一種這樣的指桑罵槐.
單打在門徒的面前.
耶穌就說.
康健的人用不著醫生.
有病的人才用得著.
經常說我喜愛連戌.
不喜愛制治.
這句話的意思.
你們且去揣摩.
我來本不是召義人.
乃是召罪人.
可圈可點的主耶穌.
不是要去反抗.
回敬這些法利賽人.
你們只顧著道德的道貌.
有沒有想過其他人的需要.

$^{1561}$耶穌不是要和他們挑戰.
耶穌是想給門徒知道.
有病才需要醫生.
我不是要飲飲食食.
我是要接近這班人.
讓他們有機會感受到真正的聖潔.
讓他們感受到.
可以被神接納的悔改.
他們就會改變他們的生命.
所以法利賽人其實很陰險.
在這裡跟他的門徒說.
耶穌的門徒說.
耶穌有機會可以澄清.
將真理說出來.
他背後的動機.
他的內心原來是這樣的.
藉著法利賽人.
他更加有機會告訴門徒.
告訴世人.
其實牧者去接觸.
其實就是能夠將真理帶出來.
神是愛世人.
每一個人都是需要救恩.
有病才需要醫生.
沒有病的人就不需要.
我來這個世界不是去拯救義人.
根本就沒有.
是要拯救罪人悔改.
所以藉著這樣的機會.
主耶穌就將真理.
更加牢固地.
落印在門徒的心靈裡.
但不是每次你都聽到.
聽不到又怎麼辦呢.
就在門徒的內心.
會有一個很壞的印象.
人就是這樣.
很容易會有這種先入為主.
當有一個人不斷在你面前.
說誰誰的壞話.

$^{1601}$聽下一次你不相信.
兩次三次.
可能慢慢聽多了.
你就會覺得那個人是這樣的.
所以我們一定要有一個.
尋求那個證據.
我們自己都要對真理.
有一個全面的掌握認識.
不要一面之詞.
更加不要聽信一些的殘言.
否則門徒都不健康.
聽完聽完就不健康了.
這是其中一段.
你成為了耶穌的門徒.
你就會成為被人統戰的對象.
記得了.
你不是耶穌的門徒.
就沒有人理睬你.
現在你是耶穌緊密的跟隨者.
那些人就要挑戰耶穌.
毀謗耶穌.
他就會跟你說.
所以大家都要很小心.
不同的人.
我們應該要很客觀地去聆聽.
不是要塞住他.
但是聆聽完之後.
我們要消化.
我們要好像有主耶穌的陳述.
這個分析.
為什麼呢.
其實你要幫助耶穌說.
如果門徒成熟.
去到一個真正能夠獨當一面.
他們可以將真理說出來.
耶穌這樣做.
你知道為什麼嗎.
我告訴你.
法利賽人的意思.
其實就是要接近這些人.

$^{1641}$不是為了吃.
而是為了將真理.
神的愛.
帶進他們的心靈裡面.
難道這些人是不需要救恩嗎.
難道這些人就是被神拒絕嗎.
沒錯.
我們是被法利賽人拒絕.
但是沒有被神拒絕.
反過來.
法利賽人覺得自己是義人.
不需要救恩的時候.
其實就被神拒絕了.
所以我們要記得.
永遠都要看自己是一個.
虧欠神的一個罪人.
是一個已經蒙了主耶穌拯救的罪人.
我們就要記得.
時時有這樣的虧欠和感恩的心.
我們就健康地成長.
如果不是.
我們就會慢慢變成法利賽人.
這是第一段.
另外十七章.
瑪太.
這個比較後了.
就是耶穌進入耶路撒冷.
對不起,還沒有進入.
去到加伯隆.
有收丁稅的人來見彼得.
就說.
你們的先生是不是不納丁稅的.
彼得就說納.
然後他就回到屋子裡.
耶穌就首先問彼得.
西門.
你的意思如何?.
世上的君王向誰徵收關稅?.
丁稅.
是向自己的兒子呢?.

$^{1681}$是向外人呢?.
當然就不會向自己的兒子.
太子,王子來徵收關稅.
彼得就說向外人.
所以很巧妙地.
幫彼得之前的衝口而出.
為什麼耶穌會知道呢?.
耶穌知道彼得.
這麼快地衝口而出.
去答收稅的人.
說一定會納的.
耶穌就說.
我是神的兒子.
既然如此.
兒子就可以免稅.
但恐怕觸犯了他們.
即地上的法例.
你且往海邊去釣魚.
把先釣上來的魚.
拿出來.
拿起來.
開了他的口.
就必定會有一塊錢.
可以拿去給他們.
作你我的稅銀.
好像慰藥彼得這個無失.
無大無小.
這麼快就衝口而出.
幫耶穌答了.
就好像收不到科.
表面上其實要納稅.
是人之常情的.
彼得其實不是.
好像一個很無失的說話.
其實每個男丁.
在羅馬要收稅.
是沒有問題的.
一定要做的.
耶穌都說了.
怕別人誤會.

$^{1721}$觸犯了法律.
就叫彼得去打一條魚.
用一個神蹟.
去幫彼得解這個尷尬.
再印證自己是神的兒子.
是不需要納稅的.
因為全宇宙世界都是主耶穌的.
都是神的.
這樣的一種情況.
就讓我們看到.
耶穌基督對自己身份的認識.
也透過這樣.
很想彼得都能夠認識.
耶穌基督的身份.
當然這裡就是彼得.
被收稅的人質問的時候.
就來問關於耶穌的事.
當然彼得不會了解得那麼深.
所以他就說得很快.
一定要納的.
因為是必須要做的事.
這裡又讓我們看到.
稅官到底他是想.
在這件事情上.
對主耶穌有一個影響.
問題.
還是對彼得有一個影響呢.
表面看得不是那麼清楚.
但是好像都有一重阻隔.
為什麼不直接問耶穌呢.
為什麼要透過彼得去問呢.
好像又是一個值得我們再去.
有機會再思想的.
比較深入一點的一個案例.
不過我們回頭看馬太的十二章.
回頭看.
有一段經文可能大家有機會聽過.
但是當我們這樣去看.
很容易會有人.
在我們耳邊說很多話.

$^{1761}$我們要很小心.
因為今天我們看的黎謙敏的六章.
正正就是那些陰謀論.
很多很多的說話.
我們要很小心的去拿捏.
那耶穌又面對一個.
這次是耶穌被人去投訴.
被耶穌知道.
十二章.
那時.
馬太福音的十二章.
那時.
耶穌在安息日從墨地經過.
他的門徒餓了.
就撕一點點墨水來吃.
就是在田的墨水.
番里賽人看見.
這次就不是跟門徒說了.
就對耶穌說.
看啊.
你的門徒作安息日不可作的事了.
向耶穌投訴他的門徒.
在安息日.
安息日就是什麼都不可以做.
在安息日.
耶穌帶著門徒在墨水經過.
墨田經過.
那些門徒真的餓了.
就撕一點點墨水.
來放進口裡.
很基本的充饑.
就不是一路像修葛那樣.
去打那些禍菌.
而是撕一點點墨水.
這樣去充饑.
有個律法書有教的.
是可以的.
但番里賽人.
就向耶穌投訴.
第三節.

$^{1801}$耶穌對他們說.
經常記著大衛和跟從他的人.
飢餓之時.
所作的事.
飢餓.
我就說兩個舊約案例給你們知道.
給你們這班番里賽人知道.
你們是不是熟讀舊約.
應該是的.
我就告訴你.
飢餓都有這些特殊的案例.
是.
即是.
不是常規.
但是.
是有過這些個案.
是神都接納容許的.
哪兩個案例呢.
第三節那裡說.
經常記著大衛和跟從他的人.
就是逃避.
素羅王追殺的時候.
撒慕爾記上的.
應該是十七章的經文.
所作的事.
你們.
對不起.
應該是二十一章.
撒慕爾記.
你們沒有念過嗎.
他怎麼進了神的殿.
當時是回末.
吃了陳設餅.
因為當時有一個祭司.
很害怕.
大衛就跟他說.
為什麼你會出現.
祭司問他.
因為王的事很緊急.
所以我們就.

$^{1841}$手頭上沒有那些.
走遠路的.
那些的食物.
你可不可以給一些食物我.
這樣.
祭司說沒有.
只有剛剛更換的陳設餅.
是祭司吃的.
外人是不准吃的.
這是祭司的下欄.
供奉完神之後.
就是祭司吃的.
他的兒女都可以吃的.
這些告訴我們.
吃了陳設餅.
大衛和他跟隨的.
他的兵丁.
吃了這些陳設餅.
這餅不是他和.
跟從他的人.
可以吃得.
唯獨祭司才可以吃.
主耶穌引這個案例.
當然他絕對不是.
來讚賞大衛說謊.
欺控祭司.
但是說一個飢餓的案例.
是神都某程度.
接納大衛.
在當時.
他有這樣的舉措.
接著還有另外一個案例.
第五節.
再者.
律法上所記.
當安息日.
祭司在殿裡.
煩了安息日.
還是沒有罪.
你們沒有念過嗎.

$^{1881}$其實是什麼意思呢.
最忙的就是安息日的祭司.
個個來上聖殿獻祭.
首節.
最忙的就是祭司你們的人.
幫忙來驗證那些祭生.
有沒有殘疾.
教你們怎樣獻法.
那個百姓首先要按手.
在那個祭司.
認同這個祭生.
稍後就會代表著你.
在神的祭壇獻上.
如果你是獻贖罪祭.
你就要根據你獻的這個祭生.
如果你是獻公牛.
或者是羊.
綿羊.
或者是班溝.
或者是初甲.
甚至乎你可以獻一些麵粉.
即是那些細麵.
諸如此類.
那祭司就會告訴你.
那個手續要怎樣怎樣怎樣做.
那你們的人就會幫你去宰殺.
宰殺完之後.
祭司就會幫你去彈血.
贖罪祭所以就不可以吃.
那些祭生的肉.
就會分給祭司吃.
又是祭司的下欄.
贖罪祭.
但如果是獻平安祭.
就完全不同了.
做法完全不同了.
就是將那些內臟的精華獻上.
然後整個祭生.
這是利美的第三章.
平安祭.

$^{1921}$整個祭生就可以祭司吃.
祭司的家人可以吃.
利美獻祭的人可以吃.
再加上那些上來獻祭的.
例如這個伊利加那.
他們是吃祭肉的.
就是全家都可以吃的.
因為是平安祭.
酬恩祭.
所以祭司.
獻祭者.
和奉獻的人.
都可以一起吃.
這個就是神人共享祭.
神可以吃.
神就是吃了那些肌肉.
那些肝.
那些精華.
精華就獻上了.
其實那些油.
那些油脂.
但整個祭生.
是給回獻祭者的.
和那些事奉的祭司.
都一起享用的.
所以是神人共享祭.
就是平安祭.
所以是有不同的.
所以這裡主耶穌說.
最忙碌的那個.
凡安息日的就是祭司.
但是沒有罪的.
因為是神安排的.
所以你想像得到.
主耶穌說了兩個案例.
不斷地發生的.
祭司安息日工作.
甚至他自己法理上也是.
一些特殊的案例.
不應該發生的.

$^{1961}$但是都有某程度的神的閱立.
因為大衛是一個很危急的逃亡.
不是讚賞那件事.
成為一個常規的案例.
就是說大衛如果第二次.
肆無忌憚的.
做了君王之後.
說他又去聖殿.
隨便拿一些陳設餅來吃.
那就大件事了.
那個就不是危急.
他就變成了.
他真的無大無小的.
不將神放在他的心靈裡面.
所以求主都幫助我們.
來到神的殿.
聖殿或者是教會.
教堂.
我們都真的要有那個敬虔.
相信我們暫時.
香港人都不會的.
突然衝上去港台.
趕傳達人下來.
我講話.
應該就沒有人敢做的.
沒有人敢這樣做的.
暫時整個社會.
都很尊重宗教.
暫時沒有這樣的情況.
不過要很小心.
真的要很小心.
這些事情絕對不能說笑.
神的殿絕對不是開玩笑的.
這裡告訴我們.
兩個案例.
讓法利賽人.
向主耶穌.
抵毀他的門徒.
其實背後都是想抵毀耶穌.
你怎樣教你的門徒.

$^{2001}$安息日這麼重要.
居然間來去壓那些墨水.
其實是用手指這樣捏一捏.
很少的.
但他們就將這件事無限的放大.
其實背後真正的問題.
不是他們真的這麼緊張.
那些墨水.
或者是那個安息日.
是緊張要怎樣來推倒耶穌.
和他的門徒.
那個德行等等.
其實背後真的很多的.
內心的很多的.
可能是玻璃心.
可能是內心裡面.
總之是對主耶穌所做的.
都是一概看不順眼.
所以一路增加.
這一種的這樣的底位.
接著路加福音就很強烈了.
路加福音就21章.
對不起20章.
路加的20章.
可能我們看完這裡.
就差不多要完了.
第20章.
很重的說話.
去到這裡.
差不多是主耶穌.
在聖殿裡面.
最後的受難周.
第20章的第19節.
路加福音.
文士和祭司長.
看出這個比喻.
就是那個所謂的空樂園戶的比喻.
是指著他們說的.
他們就是那些空樂園戶.
又霸佔著.

$^{2041}$又不能夠結出果實.
去交給那個田主.
所以他們就是.
不單只是霸佔地圖.
是阻礙人去得到神的恩典.
接著就告訴我們.
當時就想要下手.
捉拿耶穌.
不過就懼怕百姓.
第20節.
於是窺探耶穌.
打發奸細.
裝作好人.
大家可能沒有留意到這些經文.
有時福音書太長.
我們未必那麼留意.
我們如果集中那個主題.
原來作為主耶穌的門徒.
你是會被統戰的對象.
說你老師的壞話.
在老師面前說門徒的壞話.
現在已經去到最後受難周.
他們的內心已經去到一種.
整個仇恨充斥他們的內心.
已經不理那些真理.
不拉奸的都拿走.
要找一些很巧妙的.
又要隱藏在問題的背後.
讓人不知道他內心的動機.
是要衝擊耶穌的.
但是作者回憶這些事情的時候.
就告訴我們.
原來他們因為事情已經揭發了.
作者是回憶那些事情.
就將他們內心的態度.
內心的動機.
都描繪出來.
當時是沒有人知道的.
就很杜猛黯然地說.
打發奸細.

$^{2081}$裝作好人.
要在耶穌的說話上得到把柄.
好將他交在巡撫的政權之下.
奸細就問耶穌.
夫子啊.
我們曉得你所說的所傳的.
都是正直的.
也都不取人的外貌.
乃是誠誠實實.
傳神的道.
整句話.
差不多是對主耶穌的.
那個欣賞.
那個羨慕.
一個正直的老師.
在面前.
但是作者告訴我們.
那些是煙幕.
人家這樣誇讚你.
你就要很小心.
一來他的為人是怎樣.
你要知道.
原來真正的重點就是.
第22節.
是很詭詐的一個問題.
我們納稅給該撒.
可不可以.
真是.
好巧妙地想到一個.
這麼惡毒的一個說話.
其實你問不問.
你都要納稅的.
給該撒.
你就說了.
那個稅收是要給的.
但是神的兒子其實不需要.
就是說如果要答是.
猶太人其實就很不喜歡.
我們是亡國.
我們其實根本不應該納稅給該撒.

$^{2121}$我們是神的子民.
我們應該.
來到去納稅.
奉獻給神.
但是.
為什麼會斷線.
哦.
因為沒有了網.
網斷了.
原來這樣.
突然間沒有了那個.
因為是106所以.
不要緊了他們都斷了現在.
過得了.
他們上了.
不好意思.
剛才突然間斷了線.
教會那個網.
最後這一段.
我們看完就算了.
就完結了.
我們就想到.
就是.
要很小心.
每個人.
對你.
或者在這件事情上.
對主耶穌的提問.
納稅給該撒.
是一個.
現在真的很風頭火勢.
就快逾越節.
猶太人都好像等著.
有一個很強烈的.
很高漲的復國的情緒.
就是說.
可能已經醞釀著.
耶穌可能就是.
我們的釋放者.
釋放我們離開羅馬的政權.

$^{2161}$所以大家.
你叫我不納稅.
我們真的不納了.
所以很容易煽動的.
我們的說話要很小心.
就是.
不要讓群眾.
成為了我們的炮灰.
主耶穌很巧妙地.
解釋了.
你可以想像到.
要你馬上就要回答.
不讓你想.
可想而知.
耶穌對這些事情.
一路祂自己的研判.
都是非常具備真理.
和很有sense.
去知識這些問題.
祂知道現在.
猶太人的復國情緒.
推翻羅馬的高漲.
就是說.
耶穌可能就是一個革命家.
可能被人推上去成為革命家.
要推翻羅馬.
但是.
如果不是的時候.
又怎樣能夠不犯法呢.
所以你可以想像到.
主耶穌怎樣來.
拆解這個問題.
24節.
對不起23節.
耶穌看出.
他們的詭詐.
對他們說.
你拿一個銀錢.
來給我看.
這個賬和這個號.

$^{2201}$可能就是由這些.
所謂的奸細.
拿給耶穌.
不知道耶穌原來.
是有一些再深度一點的教導.
好啊.
就給你一個銀錢.
是羅馬的貨幣.
這賬和這號是誰.
當然就是.
蓋薩.
如果大家知道在耶穌的時代.
蓋薩.
這個阿古士督因為打勝仗.
就鑄幣.
鑄一些幣紀念.
他的勝仗.
羅馬的承平.
羅馬的和平.
這樣來鑄一些貨幣.
就好像回歸.
紫荊花.
類似這樣.
一定是會這樣.
對他自己的政權.
歌功頌德.
所以他們不如有詐.
當然耶穌不是有詐.
他們拿錢給耶穌.
這個是羅馬的貨幣.
耶穌說.
這讓蓋薩的物.
當歸給蓋薩.
神的物.
你們是屬神.
當歸給神.
所以不只是地上.
有天上.
我們活在一個.
羅馬統治的世界.

$^{2241}$但是我的差遣.
是將百姓.
帶進神的國度.
所以你們是神的選民.
你們是歸神屬神.
表面好像.
講了好像沒有講.
但實際上是很有智慧地.
將那些奸細.
一些鬼詐.
直接.
衝擊耶穌.
不再掩飾.
但其實是掩飾的.
怕百姓不知道.
好像突然間很好學.
每天就來找一些把柄.
希望耶穌講錯話.
就吃住它.
你說什麼都錯.
你說只是給神.
即是說不用納稅.
即是現在你煽動我們.
去反羅馬.
很難解決.
但處耶穌.
在這個時候.
仍然不是.
掉在地上那堆沙.
而是很有智慧地.
將兩件事.
不單沒有抵觸.
而且更加帶進神的國度.
來到這個地上.
即是說.
意味著.
他不是一個地上一次的革命家.
第一次來.
是要拯救百姓的罪惡.
第二次來.

$^{2281}$是要拯救百姓的罪惡.
與神和好.
沒有叫你.
不去納稅.
但如果你只有納稅.
你沒有了神的國度.
永生.
即是彌哥底母.
你就見不到.
永生.
人如果不是藉著我.
就不能夠見神的國.
你就可想而知.
主耶穌又將一個平衡的世界.
告訴我們.
有你見到的.
現在你受苦的羅馬的壓制.
但是.
心靈的罪惡.
你見不到.
但實在是需要神的救恩.
藉著我的救贖.
你就可以解除.
這個壓制.
成為.
不再是神的奴隸.
罪人.
而是成為神的兒女.
所以可以想像到.
主耶穌是多麼有力量.
智慧.
來幫助.
當時的處境.
我們說的就是.
那個是一個.
陰險的詭詐.
詭計.
所以我們都很需要求神.
讓我們有這些學習.
真的.

$^{2321}$不要有一個.
單純的頭腦.
複雜的心靈.
這樣就大件事了.
信耶穌就是要讓我們有心靈.
簡單.
簡單到.
我們的心時時都有感動.
但腦袋是要複雜.
我們去研判不同的事情.
但是記住.
心靈是需要有一份.
神的上上的感動.
我們一起禱告.
多謝主讓我們能夠在第四.
第五第六章的.
尼希米記.
尼希米個人的自述.
他在面對.
這個建造工程.
那些困難是.
內外都這樣來到.
接踵而來.
其實每一個問題.
都會有的.
多謝主讓我們能夠透過聖經.
將那些大的方向.
原則都講述出來.
讓我們看到尼希米.
他絕對不會被勝利.
充分了頭腦.
他看到有很多人支持他.
去建造城牆.
但他更加看到.
外邦人他們的.
拒絕.
攻擊是不會改變的.
因為正正他們就是.
神的選民以外的人.
到第五章的時候.

$^{2361}$我們又見到.
內部的矛盾如果不同心.
解決.
其實根本談不上什麼叫做合一.
和復興的.
多謝主讓我們能夠看到尼希米.
他面對這些.
的難處.
他直接了當的對貴族.
領袖來到去.
作出挑戰使得他們能夠.
回到神的律法書.
願意來到去.
跟隨不再來到去.
向一些窮乏的弟兄.
收取暴利.
來到去面成一種欺壓.
多謝主讓我們.
能夠有第六章.
我們見到有層出不窮的陰謀.
來到去衝擊尼希米.
雖然可能是.
只是.
隨便舉例幾個.
案例.
但都已經足夠讓我們看到.
我們要很小心.
真是有人.
是會有.
計劃地這樣去攻擊.
神的工作.
我們懇求主都讓我們能夠.
很有恩典地.
來到去看.
對人我們有一份的寬大.
愛心,但對一些的事情.
我們的研判.
也都要有智慧和恩典.
求主幫助我們.
多謝你聽我們禱告.

$^{2401}$奉耶穌基督的名,阿們.
\newpage



\section{申命記 6:1-9}
\label{sec:A6GWJTLfDEo}
\textbf{2023-05-14 加恩樓《一家之主》申命記 6\_1-9 講員:謝靄薇牧師}
\newline
\newline
連結: \href{https://youtube.com/watch?v=A6GWJTLfDEo}{\texttt{https://youtube.com/watch?v=A6GWJTLfDEo}} ~~~~ 語音日期: 2023-05-16
\newline
\newline
\hyperref[sec:_3DxJVmmiXg]{\small{< < < PREV SERMON < < <}}
~
\hyperref[sec:index_chronic]{\small{[返順時目]}}
~
\hyperref[sec:index_scriptual]{\small{[返順卷目]}}
~
\hyperref[sec:noei9nAUNuE]{\small{> > > NEXT SERMON > > >}}
\newline
\newline
申命記 6:1-9
\newline
\begin{longtable}{cl}
\hline
\hline
章節 & 經文 (和合本修訂版)\\
\hline
6:1 & \begin{tabularx}{0.7\textwidth}{X} 「這是耶和華－你們的神所吩咐要教導你們的誡命、律例、典章，叫你們在所要過去得為業的地上遵行， \end{tabularx} \\ \\ \relax
6:2 & \begin{tabularx}{0.7\textwidth}{X} 好叫你和你的子孫在一生的日子都敬畏耶和華－你的神，謹守他的一切律例、誡命，就是我所吩咐你的，使你的日子得以長久。 \end{tabularx} \\ \\ \relax
6:3 & \begin{tabularx}{0.7\textwidth}{X} 以色列啊，你要聽，要謹守遵行，使你可以在那流奶與蜜之地得福，人數極其增多，正如耶和華－你列祖的神所應許你的。 \end{tabularx} \\ \\ \relax
6:4 & \begin{tabularx}{0.7\textwidth}{X} 「以色列啊，你要聽！耶和華－我們的神是獨一的主。 \end{tabularx} \\ \\ \relax
6:5 & \begin{tabularx}{0.7\textwidth}{X} 你要盡心、盡性、盡力愛耶和華－你的神。 \end{tabularx} \\ \\ \relax
6:6 & \begin{tabularx}{0.7\textwidth}{X} 我今日吩咐你的這些話都要記在心上， \end{tabularx} \\ \\ \relax
6:7 & \begin{tabularx}{0.7\textwidth}{X} 也要殷勤教導你的兒女。無論你坐在家裡，走在路上，躺下，起來，都要吟誦。 \end{tabularx} \\ \\ \relax
6:8 & \begin{tabularx}{0.7\textwidth}{X} 要繫在手上作記號，戴在額上作經匣； \end{tabularx} \\ \\ \relax
6:9 & \begin{tabularx}{0.7\textwidth}{X} 又要寫在你房屋的門框上和你的城門上。 \end{tabularx} \\ \\
[1ex]
\hline
\hline
\end{longtable}
$^{1}$各位聽眾平安.
一家之主.
問問你們.
在家裡誰有話事權呢?.
很多時候.
我發現的情況就是.
做丈夫的那個.
會問老婆意見.
有時候會聽到.
女士們說.
大事就交給老公.
那些小事就他自己解決.
不知道你們是不是這個情形呢?.
其實.
我們家裡應該有一個金字塔.
應該是.
以神為首.
一家之主應該是我們的神.
讓我們一起在看神的話語之前.
我們一起祈禱.
沒有恩典沒有慈愛的主.
我們真的歡喜快樂.
因為今天是你的日子.
你讓我們.
因著你的愛招聚我們回來.
盼望申命記這段話.
你同樣的呼召我們.
去教導我們.
以致我們願意.
順服你的旨意.
去走出來求你保守我們的心.
我們這樣祈禱教託.
奉主耶穌基督命求.
阿們.
我們剛才說到.
這個西乃之約.
是最早的時候.
和以色列人去立約的.
現在摩西已經是年紀大了.
當他們年紀大的時候.

$^{41}$新一代將要進入.
這個很美麗的.
迦南地.
但是這群人.
他們並不是那麼了解.
當時的西乃之約.
他們又不明白.
作為神的僕人.
對神應該有什麼態度呢?.
要忠信.
謙遜 愛心.
所以摩西就再次.
和他們重新宣佈.
這個律法的意思.
所以你在.
《申命記》十二章裡.
就會看到他們有定時的敬拜.
他們去遵行神的律法.
去更新他們和神.
所納的約.
因為將來他們.
進入這麼美麗的地方.
當他們願意遵守神的話語.
就會享盡.
神給他們一切的豐盛.
首先.
我們看到在經文裡.
他就說到要聽.
要聽就是一個命令的語氣.
我們為何要聽呢?.
原來.
我們要聽這位神.
祂是我們獨一的真神.
在第四節.
在《申命記》第六章的第四節.
就說到.
以色列你要聽.
耶和華的神是獨一的主.
聖經很清楚地.
擺明給我們知道.

$^{81}$當你要認識這個獨一的主.
就是說.
除了我們這位神.
耶和華神之外.
其他的都不是真神.
就是這個意思.
基督教最基本的信仰就是.
違背了的.
就是異端了.
就是假先知.
就要從民眾杜絕.
所以.
陸續在《申命記》第十三章.
一到三節.
就告訴我們.
你中間若有先知或是.
做夢的人起來.
向你顯神跡其事.
他對你說的神跡其事.
應驗了.
說我們去隨從別神.
事奉他們吧.
那是你不認識的.
你不可聽那先知或是.
那做夢之人的話.
因為這是耶和華.
你們的神考驗你們.
要知道你們是否盡心盡性.
愛耶和華你們的神.
真的先知的憑據.
不是看他有沒有能力.
去行這些神跡.
最重要的是.
他所說的話.
他的信息.
是否合乎神的啟示.
和合乎基本的信仰.
基督教的信仰也有一個不變的準則.
就是我們相信.
耶穌基督的超越性.

$^{121}$和祂的神性.
神跡其事.
和傳教的工作.
這不是我們真實信仰的基礎.
真假的信仰的區別.
在於你是不是真的相信.
神是藉著耶穌基督.
去完成救贖.
承認耶穌基督.
是三位一體的神.
古時候認為夢是神的啟迪.
因為神經常透過夢.
去向人啟迪.
所以後來這些夢.
就成為了假先知.
他們自欺欺人的手段.
他們騙人.
神跡是為了指示將來的預兆.
其事.
即是奇妙的奇.
其事即是要證明先知.
或者做夢的.
他確實有能力行其事.
但試驗真先知的一個方法.
就是他所說的預言.
要完全的應驗.
假先知.
即是說他也會應驗的.
不過有時候要敲.
即是一些花巧的語言.
他說他自己可以預測將來的.
但他的目的是.
令人離開神.
當你看到這些人.
帶你離開神的道的時候.
你要小心了.
要防備他們.
要反對他們.
因為神讓這些人.
去顯出神跡的奇事.

$^{161}$是為了試驗以色列人.
是否真的專靠耶和華神.
考驗一個人的靈性.
就是在於.
他是不是真的.
遵從神的話語.
他有沒有去事奉耶和華神.
有沒有去愛這位神.
有沒有去倚靠他.
我們如果是真的愛神.
我們就會敬畏我們的神.
同時我們也會遵守.
他所說的話語.
就在聖經裡面.
那我們要聽從甚麼呢?.
聽從甚麼呢?.
就是在第六章的開頭.
就會看到.
就是誡命,律例,典章.
然後在第二節.
就被人看到.
你的子孫.
在一生的日子.
都要敬畏耶和華神.
要謹守他.
又是說到律例,誡命.
就是神吩咐的.
還說你的日子可以長久.
也就是說你的壽命會很長.
敬畏神.
遵守他的律例,誡命.
不單止你長壽.
這個國家.
也有福祉.
然後在第三節.
第三節那裡.
就會看到.
又是叫你要謹守進行.
要聽.
將來在神應許的地方.

$^{201}$禮物之地.
你可以得福.
同時人數會增多.
這是神應許我們的.
謹守.
神的話語.
神應許他的子孫.
你的下一代.
再下一代.
都可以在這個福地裡面.
去享受.
同時人數會增多.
因為那個地方是神所眷顧的.
最重要就是.
聽命.
誠於憲制.
所以你會看到.
神很願意.
這班屬於他的子民.
要遵從他的話語.
要去聽.
但我們聽了很多.
叫你做.
其實不容易.
叫你每個星期回來.
事奉.
有時會覺得.
好像是一個重擔.
為什麼呢?.
每個星期都回來.
好像那段時間.
捱得很辛苦.
不知道要多久才能看完.
不知道何時才能做完.
可能那天時間長了.
那時才可以收拾.
有時看完收拾.
人家會覺得好像趕客.
還沒關門.
收拾好東西.

$^{241}$關燈.
不是歡迎人的態度.
最後才能走.
好像很痛苦.
為什麼呢?.
因為我們不是用愛去事奉.
用愛去事奉又不同.
聖經教我們.
要盡心盡性盡力.
愛耶和華是你的神.
因為愛在這裡.
愛是生命的重要主題.
神對子民的愛.
我們不懂得愛.
聖經說.
是神先愛我們.
我們才懂得.
愛人和關心人.
所以.
接著就會看到.
原來神愛你.
所以你就很樂意.
去進行.
因為神很關心你.
關心你的大事.
小事.
他都留意的.
這麼愛你的神.
當你親近這位神.
這麼愛你的時候.
你都願意去做.
因為他很愛惜你.
你很樂意去進行.
是因為有愛在裡面.
因為神先愛你.
所以在《申命記》第七章七節.
就看到.
說是專愛你們.
揀選你們.
當時對這些以色列民說.

$^{281}$原來他們.
不是有什麼厲害的地方.
他們人又不多.
人又多過他們.
但神偏偏要揀選他們.
做他們的子民.
神偏偏要愛他們.
專愛他們.
你會看到.
在萬民裡面.
原來他最愛他們的子民.
雖然人數少.
但他揀選他們.
他的後裔.
一代一代.
他都揀選他們.
因此.
神還吩咐他們的子民.
你得到這個享受.
但你還要.
關心孤兒寡婦.
或寄居.
要憐愛他們.
所以種的麥田.
麥碎都不要撿走.
留一些讓孤兒寡婦.
可以撿回去.
有一個憐愛的心.
神常常提醒他們.
你們曾經是.
圍爐之家.
出來的.
你們沒有什麼比別人厲害.
你們曾經是奴隸.
所以你們對人.
都不要偏心.
今天可能很風光.
但你們的過去.
曾經都有.
很痛苦的日子.

$^{321}$你回想那個痛苦.
你就有一個憐憫的心去對待人.
所以盡心盡性.
去愛我們主.
在《馬虎音》裡面.
耶穌基督教導我們.
當我們.
很樂意去做的時候.
聖經又跟著.
告訴我們.
要去教導下一代.
要殷勤去教導.
你的兒女.
要教.
要付代價.
將來在座你們.
有兒女的.
要知道那種痛苦.
由他們出生.
開始沒有自己.
可能養大他們也不容易.
還要教他們.
教他們還要生生性性.
不單是他去打你嘴巴.
更加激氣.
教導是不容易的.
同時在一個家庭裡面.
自己人很難教的.
不聽你的.
中國人早傳秘方.
有時候也沒有傳下去.
以前我事奉的教會有一個團友.
他爸爸是做蝦膏生意.
賺很多錢的.
在大陸設廠.
後來他爸爸去世之後.
他一直沒有.
告訴別人.
他只是獨自經營.
當他離世的時候.

$^{361}$就翹起來了.
怎麼辦呢?.
出現很多問題.
到底要怎樣追溯呢?.
怎樣維持呢?.
到底是否要結束呢?.
無從著手.
因為他也沒有交代.
還有這個教育.
不單是要殷勤.
還要隨時教育.
無論你是.
坐在家裡.
走在路上.
原來是一個家庭的時間.
在路上.
或者躺下來.
都要談論.
隨時教育.
大家有沒有想到.
原來你的小朋友就在你家.
隨時都會看到.
無論睡覺還是起床.
你都要.
看著他.
小朋友分分鐘不看.
就會怎樣呢?.
可能他撞到哪裡.
就很麻煩.
坐在家裡.
走在路上 躺下來.
原來都要淫訟.
對於整個話語.
你要去讀.
讀著讀著 小朋友又聽到.
是隨時的教育.
但今天我們.
多數是談論什麼呢?.
是八卦新聞嗎?.
或者馬經呢?.

$^{401}$或者股市上升?.
或者下跌呢?.
我們是不是經常談這些呢?.
但是要教他認識神.
不如交給主要學位的老師.
你想想.
他一個星期來一次.
他能記住多少呢?.
下星期可能又忘記了.
又交給你.
大家做老師的.
應該有這個經驗.
因為他一個星期.
就算他上五天學.
也有五天.
但你一個星期來一次.
怎麼記得呢?.
可能今天跟他說的.
他都忘記了上星期說了什麼.
問你上星期說了什麼.
可能你都忘記了.
可能連經文都忘記了.
他又學了很多.
芭蕾舞,琴,畫.
可能有些人還要學珠算.
身邊有很多東西.
電腦啊.
他怎麼記得那麼多呢?.
但父母會問他.
你今天學了什麼呢?.
他學琴也要練琴.
你不會讓他坐在家.
因為你交了錢.
他交了學費.
那你就會緊張.
但回主一學.
隨便吧.
有沒有人會溫習金句呢?.
我說你很熟悉.
你會記住今天學什麼金句.

$^{441}$可能是.
自己弟兄姊妹.
或者你是傳道人.
可能會緊張.
但如果是普通弟兄姊妹.
他又回來已經算少了.
你會不會跟他溫習金句呢?.
我想應該是.
你應該這樣做.
因為你沒有那麼多時間.
請問你.
你到底有多少時間陪你的小朋友呢?.
陪你的兒女呢?.
就算他長大了.
你有多少時間跟他聊天呢?.
你用多少時間?.
這裡說隨時教育.
或者你見他的時候.
他已經在睡覺.
可能我們自己都沒有祈禱.
也沒有讀聖經.
我也不知道有什麼話題.
要講到神的事.
要淫頌.
你經常去讀.
我自己都沒有做的時候.
很難去要求他.
很多時候你會聽到.
有些人打電話.
找他父母.
可能他告訴他父母不在.
新教.
比現教重要.
你讓他看到的就是負面的榜樣.
他每天都在學習你的脾氣.
學習你的談吐.
他全部都在方言淚.
我記得我兒生.
他很小的時候.
他已經會背.

$^{481}$哥林多前書.
十三章.
那個愛的詩章.
三歲.
做父母其實.
沒什麼的.
以前的錄音帶.
很有趣的.
以前的故事講得很生鬼.
小朋友覺得很有趣.
就按來按去聽.
一天聽了幾百次.
他已經熟悉了那個故事.
他會告訴你.
他覺得很好玩.
背金句一樣.
大人只要留意他.
幫他的時候.
原來他讀得到.
這麼小的時候他都讀得到.
老人家.
你說我沒有記性.
你到老的時候.
你就記得神的話語.
很有意思的.
羅雅一家在.
豐州上面.
他們一家八口.
得救.
羅雅帶領一家人去敬拜神.
有原因的.
他曾經在地上.
有120天.
他都有唐人傳福音的.
不過人們不太理會他.
不相信他講的話.
他就死在大洪水的浩劫裡面.
羅雅當他得救之後.
離開了豐州.
他就去竹壇.

$^{521}$他拿一些潔淨的牲畜.
和潔淨的飛鳥.
去竹壇獻祭.
獻祭繁祭.
於是.
羅雅就這樣做.
然後亞伯蘭.
所謂的亞伯拉罕.
他就跟著做.
還有他.
以撒 雅各.
都會跟著這樣做.
因為前人是這樣做.
後代的都會跟著這樣做.
因為他們學習.
如何去敬拜神.
你會看到.
父母的一舉一動.
都是成為他學習的模範.
是會影響他們一生的.
雅各.
他從小就.
有著長者的名份.
當你看聖經的時候.
他很懂得長者的名份.
長子的名份.
雖然他用.
鬼詐的手段.
去拿了長子的名份.
神早就知道.
他會這樣做.
神早就預言.
他在武福裡面.
已經預言大的要伺世.
是不是很奇妙呢?.
神已經預備了這樣.
如果你看回.
《伊節經文》.
下集說到.
要謹慎.

$^{561}$免得有人失去了神的恩典.
免得有毒更新出來.
擾亂你們.
因而使許多人.
沾染互為.
免得有人淫亂.
華不敬虔.
把自己長子的名份賣了.
後來你們知道.
他想要承受.
父親的祝福.
竟被拒絕.
雖然流著淚苦求.
卻得不著門路.
使他父親回心轉意.
看到嗎?.
到他很想他爸爸.
祝福的時候.
怎麼哀求他.
但他爸爸已經將最好的.
祝福了.
來接受祝福的.
所以他講了最好的.
以為是給了他.
原來不是.
一個人失去神的恩典的時候.
會生毒根.
樹木的根.
如果是苦的.
果子就是苦的.
人心如果有.
毒根的時候.
就會生什麼?.
仇恨.
驕傲.
憎恨.
這樣的果子.
如果是整個團體.
這樣就恐怖了.
如果教會是這樣.

$^{601}$是不是會發生很多問題.
很恐怖.
今天很多年輕人.
無法無天.
為什麼這樣呢?.
因為沒有人教他們.
沒有人愛他們.
沒有人帶領他們認識主.
去敬畏神.
這個耳掃.
他失去.
長子的名分之後.
他就不知道悔改.
他要生出什麼?.
要殺他的弟弟.
所以我們如果遭遇逼迫的時候.
我們如果不願意.
走正道.
犧牲正道.
然後苟且行事的時候.
是很恐怖的.
所以你會看到.
在一個家庭裡面.
父母一生.
是足以影響子女.
所以他自己都要願意.
信服神.
願意聽神的話.
這個小朋友他什麼都不懂.
他需要你去帶領他.
你不要當他全都懂.
他不知道的.
很多事情他不懂分辨.
他需要你去帶領他.
要有耐性去教導他.
讓他在人生裡面.
每時每刻都有一個.
生命的教導.
這是隨時的教育.
你隨時看到.

$^{641}$不太妥你們要教他.
讓他能夠知道耶穌基督是一家之主.
讓他懂得去敬畏.
最後這個經文就說到.
我們去實踐.
要繫在手上做記號.
戴在額上做經盒.
又要寫在你房屋的門框上.
和你的城門上.
為什麼這麼重要呢?.
原來以色列人.
當他們被擄到巴比倫.
然後回歸自己的土地的時候.
這些以色列人.
他們就把金句.
外表化.
只是做外表的.
抄了一些經文.
然後放在一個小盒子裡.
盒子比較漂亮.
可能是用皮的.
然後就戴在額頭上.
或者在鼻子上有一個盒子.
代表他.
我經常都記住經文.
但是這個只是留於表面.
只是裝飾品.
但是對他內在的生命.
是沒有什麼影響的.
在《基督教的福音》.
耶穌都說過.
叫我們不是徒有外面的形式.
而沒有實質的行為是不行的.
因為信心沒有行為是死的.
所以你看到.
人們戴十字架.
他不一定是相信的.
不相信也能戴的.
很漂亮的.
但如果他不明白十字架是什麼意思.

$^{681}$對他的靈性.
對他的生命是沒有什麼意義的.
當我們看到他們的額上和手上有這個記號.
其實是什麼意思呢?.
當時神寫了這個金句是什麼意思呢?.
其實就是告訴你.
你要記得神的話語.
神的話語是很重要的.
你同時要有心裝作.
要記在心裡面.
不只是做個樣子.
讓人看到的.
要去實踐這個誡命.
最近聽說.
有些幼稚園的家長說.
他們的小朋友.
聽到國歌播放的時候.
他們竟然會叫家人.
「靜一點」.
「不要這麼吵」.
他們為什麼會這樣表現呢?.
很小我的.
大人叫你靜一點.
小朋友叫家人靜一點.
為什麼呢?.
因為他們學到的態度就是.
唱國歌.
嚴肅一點.
莊嚴一點.
這個時候.
其實要有一點點的恭敬.
尊重.
原來他們學到這一點.
雖然他們這麼小.
同樣的.
愛頓神的誡命.
神的律法.
是不是成為我們的眼睛.
成為我們的手呢?.
去指引你.

$^{721}$去看到耶穌基督.
要我們做什麼呢?.
我們看到什麼呢?.
我們有沒有一個熟齡眼光呢?.
當我們生活中.
將神的話.
深深的錄印.
在我們心中的時候.
我們的家就等於貼滿了很多金句.
是很美麗的.
掛在門框上.
還有在城門.
城門是做什麼的呢?.
城門就是來來往往.
一個市集.
一個中心.
很多人走的.
好像你看到一些廣告.
你也要多看幾眼.
提醒你.
又有什麼節目.
所以我們在外面.
也有宣傳.
人們經過.
來往的人就會看到.
你有什麼消息.
買賣.
社交的活動.
都會在那裡進行.
所以在城門.
是很重要的.
曾經在撒瑪利亞.
他們也發現.
有一個石碑寫著十誡.
他們看到.
原來當我們有些裝飾的時候.
我們不是只是外表.
我們要實質的.
真的要實行.
要實踐出來的.

$^{761}$我們要教弟弟.
我們要教他.
明白了.
最難做的就是啟迪.
如何啟發他去做.
是不容易的.
不過交給家長.
每一個.
你想想.
有什麼可以啟發你家裡的人.
可以去實踐.
我們一起努力.
我們一起祈禱.
每一天.
每一天都會有慈愛的主.
每一次我們來到你面前的時候.
是真的無限的恩惠.
你臨到我們.
我們真的是不配的罪人.
但你竟然給我們有兒女的身份.
求你幫助我們.
真的好像《申命記》裡說的.
不單止我們得到你的福分.
我們要告訴下一代.
讓我們世世代代.
都能夠高舉主你的名.
求你幫助我們.
給我們聰明智慧.
以致我們可以實踐你的話語.
可以榮耀你的名.
我們這樣祈禱.
教託.
奉主耶穌基督的名下.
阿們.
\newpage



\section{出埃及記 2:1-10}
\label{sec:noei9nAUNuE}
\textbf{2023-05-14 早堂《我們這一家》出埃及記 2\_1-10 講員:黃頌言姑娘}
\newline
\newline
連結: \href{https://youtube.com/watch?v=noei9nAUNuE}{\texttt{https://youtube.com/watch?v=noei9nAUNuE}} ~~~~ 語音日期: 2023-05-16
\newline
\newline
\hyperref[sec:A6GWJTLfDEo]{\small{< < < PREV SERMON < < <}}
~
\hyperref[sec:index_chronic]{\small{[返順時目]}}
~
\hyperref[sec:index_scriptual]{\small{[返順卷目]}}
~
\hyperref[sec:PwPuQ_V9PRU]{\small{> > > NEXT SERMON > > >}}
\newline
\newline
出埃及記 2:1-10
\newline
\begin{longtable}{cl}
\hline
\hline
章節 & 經文 (和合本修訂版)\\
\hline
2:1 & \begin{tabularx}{0.7\textwidth}{X} 有一個利未家的人娶了一個利未女子為妻。 \end{tabularx} \\ \\ \relax
2:2 & \begin{tabularx}{0.7\textwidth}{X} 那女人懷孕，生了一個兒子，見他俊美，就把他藏了三個月， \end{tabularx} \\ \\ \relax
2:3 & \begin{tabularx}{0.7\textwidth}{X} 後來不能再藏，就取了一個蒲草箱，抹上柏油和樹脂，將孩子放在裡面，把箱子擱在尼羅河邊的蘆葦中。 \end{tabularx} \\ \\ \relax
2:4 & \begin{tabularx}{0.7\textwidth}{X} 孩子的姊姊遠遠站著，要知道他究竟會怎樣。 \end{tabularx} \\ \\ \relax
2:5 & \begin{tabularx}{0.7\textwidth}{X} 法老的女兒來到尼羅河邊洗澡，她的女僕們在河邊行走。她看見在蘆葦中的箱子，就派一個使女把它拿來。 \end{tabularx} \\ \\ \relax
2:6 & \begin{tabularx}{0.7\textwidth}{X} 她打開箱子，看見那孩子。看哪，男孩在哭，她就可憐他，說：「這是希伯來人的一個孩子。」 \end{tabularx} \\ \\ \relax
2:7 & \begin{tabularx}{0.7\textwidth}{X} 孩子的姊姊對法老的女兒說：「我去叫一個希伯來婦人來作奶媽，替你乳養這孩子，好嗎？」 \end{tabularx} \\ \\ \relax
2:8 & \begin{tabularx}{0.7\textwidth}{X} 法老的女兒對她說：「去吧！」那女孩就去叫了孩子的母親來。 \end{tabularx} \\ \\ \relax
2:9 & \begin{tabularx}{0.7\textwidth}{X} 法老的女兒對她說：「你把這孩子抱去，替我乳養這孩子，我必給你工錢。」那婦人就把孩子接過來，乳養他。 \end{tabularx} \\ \\ \relax
2:10 & \begin{tabularx}{0.7\textwidth}{X} 孩子長大了，婦人把他帶到法老的女兒那裡，就作了她的兒子。她給孩子起名叫摩西，說：「因我把他從水裡拉出來。」 \end{tabularx} \\ \\
[1ex]
\hline
\hline
\end{longtable}
$^{1}$各位弟兄姊妹平安.
我們一起有個祈禱開始吧.
親愛的主多謝你.
讓我們可以一家人.
一起去在你的面前敬拜讚美.
我們學武理的話願你開啟我們的心.
讓我們有聆聽的耳朵.
多謝你 奉迎蘇明求.
阿們.
好 想問大家一個問題.
你怕不怕蛇?.
怕不怕?.
怕的話舉手來看看.
哦.
我看到很大部分都不怕.
很勇敢.
不過如果那個場景是這樣.
如果有一條蛇.
在你的小朋友旁邊張開口.
你怕不怕呢?.
怕.
我想問一下.
我想訪問一下.
我想問一下如果有一條蛇.
在你的小朋友旁邊張開口.
你會怎樣做呢?.
我請吳全道幫我選一個家長.
讓我問一下.
我會尖叫大聲嚇走牠.
你會尖叫大聲嚇走那條蛇.
你都是想嚇走那條蛇.
趕走牠 用你的方法就是尖叫.
我想問一下小朋友.
你又覺得你的爸爸媽媽會怎樣做呢?.
如果有一條蛇在你旁邊張開口.
你爸爸媽媽看到.
你爸爸媽媽會怎樣做呢?.
有沒有小朋友勇敢可以舉手告訴我?.
好.
永遠不嚇死牠.

$^{41}$牠會怎樣?.
嚇死牠.
ok 其實小朋友都覺得.
爸爸媽媽一定會保護牠.
我想告訴大家.
接下來我會說的是真人真事.
這件事是發生在我一個朋友身上.
有一天他很不開心.
所以我這個朋友和他媽媽一起.
聊天 平排坐.
突然有一條蛇在我朋友旁邊出現.
然後他媽媽做了什麼呢?.
一隻手拿著那條蛇.
然後你猜他媽媽發生什麼事?.
被那條蛇咬了一口.
然後他媽媽不怕.
一下子就把那條蛇扔到海裡.
這是一個新聞.
變成了一個新聞報導.
沒錯 其實媽媽真的很偉大.
當小朋友遇到危險的時候.
媽媽會竭盡所能去保護自己的孩子.
今天我們會一起去看聖經的故事.
我們會看一位媽媽.
她的小朋友都將要遇到危險.
不知道她會怎樣做呢?.
大家還記得摩西這個故事嗎?.
當時以色列人住在埃及這個地方.
埃及的法老王.
以色列人越來越多.
就開始很害怕.
很害怕他們會反抗.
於是就想了一個辦法去對付他們.
他要以色列人做苦工.
甚至下了一個很可怕的命令.
要埃及的士兵.
將所有以色列剛出生的嬰兒.
扔到尼羅河裡.
殺死他們.
剛出生的男嬰.

$^{81}$女嬰就可以保存.
當時有一個希伯來人的家庭.
生了一個很可愛的男嬰.
叫做摩西.
摩西的爸爸,媽媽,哥哥,姐姐.
都很喜歡他.
摩西的媽媽叫做若姬別.
她每天都很小心照顧摩西.
但是媽媽很擔心.
她擔心摩西的哭聲會被埃及人發現.
所以她每天.
全家人都會一起祈禱.
求神保護摩西.
全能的父上帝.
求你保守我們這個小小的小朋友.
真是的主啊.
願你與他同在.
保守他的腦.
神啊.
希望你保護摩西.
讓埃及士兵發現.
神啊.
希望你賜我們一家平平安安.
但是摩西一天一天長大.
轉眼間已經三個月了.
摩西長得很帥很可愛.
但是有時候他會放聲大哭.
哇哇.
媽媽若姬別聽到之後.
很害怕.
因為她很擔心.
很擔心被埃及人聽到.
所以她就用布蓋住摩西.
又會將摩西放入廚房或其他地方.
很想收起來.
最後摩西的父母想到一個方法.
他們就用蒲草做了一個箱.
上面塗上防水的石漆.
將摩西放入蒲草的箱子裡.
再將這個箱子放在尼羅河旁邊的蘆葦草.

$^{121}$摩西的姐姐米利奄在旁邊觀察.
她看著摩西會怎樣.
結果埃及公主在河裡洗澡的時候.
她發現了可愛的摩西.
公主還很喜歡她.
當公主預備收養摩西的時候.
機智聰明的米利奄就站出來了.
勇敢地大聲建議一個希伯來人做奶媽.
公主不僅答應.
她還說BB長大後可以將牠送回皇宮.
其實這個奶媽就是摩西的親生母親約基別.
神在奇妙的安排下.
摩西平平安安地回到自己的家裡.
約基別在家裡教導摩西有關以色列人的歷史.
神怎樣帶領以色列人的祖先.
亞伯拉罕,以撒,雅各.
天父對以色列人的應許等等.
爸爸媽媽還會不斷地為摩西祈禱.
摩西長大後進入皇宮成為王子.
後來神太他帶領以色列人逃出埃及.
成為偉大的民族領袖.
我們來看看大兒子.
大兒子他是阿倫.
阿倫是摩西的哥哥.
神揀選的第一位祭司.
為百姓去獻贖罪祭禱告.
他的口才很好.
是一個雄辯滔滔的演說家.
神安排了阿倫協助摩西和法老據理力爭.
阿倫和摩西互補不足.
一起帶領以色列人出埃及.
而姐姐米利安不單機智勇敢.
她也是一位女先知一位音樂家.
當摩西帶領以色列人過紅海之後.
米利安率領以色列婦女擊鼓跳舞.
歌頌耶和華領導才能表露無遺.
他們約基別的夫婦敬威神.
將神的孩子都帶到神的面前.
所以他們三位孩子都為神所用.
約基別這個名字的意思是.

$^{161}$耶和華的榮耀.
雖然聖經對媽媽約基別的描述只有很少.
但是神也藉著媽媽的信心.
讓神的榮光照耀在她的子女身上.
他們各有恩賜成為神的團隊.
為神的子民效力.
各位弟兄姊妹.
今天我的講題是「我們這一家」.
在神的家裡面.
我們在神的家裡面.
有神成為我們的避難所.
我們在神的家裡面就可以得到平安.
就像約基別的一家人.
他們面對很大的困苦和困難.
但是因為有神成為他們的依靠和幫助.
他們都很勇敢.
特別是姐姐米莉暗.
我們剛才看到她很聰明.
年紀很小 她只有十一,十二歲.
也可以拯救弟弟.
今天我們遇到困難的時候.
我們都可以到神面前尋求幫助.
我們剛才演出的一家人.
海祥,Wubi.
還有九歲的呂一心.
六歲的兒子定一.
還有海澄依依.
好.
我們也請我們一起拍照.
我們也請海澄可以先下去.
我現在趁現在要訪問一下.
我們的海祥一家人.
我看到你們的子女非常穩定出席主日學.
爸爸你是天崇的導師.
媽媽又是初小的導師.
相信你們都會非常清楚.
將小朋友帶到神面前其實很重要.
我也想問一下.
除了星期天回來教會之外.
平日你們的信仰生活是怎樣的.

$^{201}$我先說一下我們.
我太太逢星期三八點半會準時開看YouTube.
不是看YouTube 是看祈禱會.
感受堂線上祈禱會.
有時候我也會懶惰.
但我看到她這樣也令我沒法懶惰.
所以我會和她一起祈禱.
平日有祈禱會的舉手來看看.
也有的 很多小朋友有舉手的.
還有呢.
除了我們自己的靈修生活之外.
我們也會開車上學.
在那段時間最好當然是播耳道自見.
一路聽.
小朋友聽著聽著也知道.
另外也在紅綠燈位.
不一定是紅綠燈位.
有很多紅綠燈位.
一定會叫小朋友鼓勵他們一起祈禱.
通常包括五樣東西.
感恩 讚美 認罪.
小朋友告訴我們原來我們最常說的就是祈求.
剛才我們也聽他們說的時候.
哇 厲害 在紅綠燈位就立刻可以認罪.
認罪也很快.
也說兩句而已.
也不容易的.
因為我們日常生活也很忙碌.
要抓緊時間跟子女相處.
在裡面跟他們一起祈禱.
一起成長也是很重要的.
我也想問一下一心.
知道你平時也有跟小朋友一起祈禱.
我想知道你平時怎樣跟小朋友一起祈禱.
你們是什麼時候祈禱的.
睡覺前.
你們會一起祈禱.
你們會祈禱什麼.
感恩的事和不開心的事.
也會跟小朋友一起.

$^{241}$你們是睡覺時一起祈禱嗎.
不是.
你們會怎樣祈禱.
跪在那裡.
跪在那裡 你跪在哪裡.
跪在床頭.
跪在床頭一起祈禱.
謝謝你們的分享.
可以先回去坐.
各位弟兄姊妹.
剛才我們聽了很多大家分享的時候.
我們也要去明白一件事.
就是我們不要小看自己的角色.
特別是為人父母的角色.
其實神正在使用你.
在祂的神的國裡.
是培育著一位被神重用的領袖.
今天我們在約基別一對的父母身上.
我們看到他們對神很有信心.
也帶給孩子有平安和勇氣.
特別是在培育小孩子去敬畏神.
愛神的心.
讓小孩子在每一件小事上.
都可以為主去忠心.
今天其實.
神也要再一次呼召我們每一個家庭.
不知道你們是否已經認識神.
今天或許你們是被邀請來到中間.
來到神的教會當中.
我想問有沒有第一次來錦繡堂.
我們被邀請來的第一次來錦繡堂.
我們舉手來看看.
看來也有.
謝謝.
是呀 我們是歡迎大家的.
來到錦繡堂當中.
其實我們每一個人都很有限制.
我們的愛心也很有限.
能力也很有限.
教導孩子其實很不容易.

$^{281}$但神允許我們.
養育的過程我們不孤單.
因為有神和我們一起去培育我們的孩子.
如果你願意.
這是一個很慎重的邀請.
如果你願意.
將你的家庭放在神面前.
被神使用.
被神帶領.
願意認識神.
讓神成為你的一家之主.
我現在也請你.
可以在你的座位裡面勇敢地舉手.
如果你願意將你的家庭.
一同帶到神面前.
求神成為你隨時幫助的.
請你也是舉起你的手.
Amen.
好 我們一起去祈禱.
我們同心合意一起去祈禱.
親愛的主.
你愛我們每一個.
及至我們的家庭當中的每一個.
願我們的家庭在你的手裡.
願你親自成為我們隨時的幫助.
願你祝福剛剛我們很勇敢地舉手的每一位.
或者是沒有勇氣在心裡有舉手的每一位.
主我們都求你在每一個家庭當中.
成為我們隨時的幫助.
成為我們每一個家庭的主.
祝福每一位家長都能夠有敬畏神的心.
及至能夠在你的愛裡面.
去培育滿有信心的下一代.
能夠在世界當中為你作見證.
將愛帶給身邊的人.
願你復興你的家 復興你的國.
多謝你 奉主耶穌基督名求.
Amen.
\newpage



\section{腓立比書 4:4}
\label{sec:PwPuQ_V9PRU}
\textbf{2023-05-13 晚堂《喜樂再思》腓立比書 4\_4 講員:陳一華牧師}
\newline
\newline
連結: \href{https://youtube.com/watch?v=PwPuQ-V9PRU}{\texttt{https://youtube.com/watch?v=PwPuQ-V9PRU}} ~~~~ 語音日期: 2023-05-17
\newline
\newline
\hyperref[sec:noei9nAUNuE]{\small{< < < PREV SERMON < < <}}
~
\hyperref[sec:index_chronic]{\small{[返順時目]}}
~
\hyperref[sec:index_scriptual]{\small{[返順卷目]}}
~
\hyperref[sec:_W_gd9p_W_Y]{\small{> > > NEXT SERMON > > >}}
\newline
\newline
腓立比書 4:4
\newline
\begin{longtable}{cl}
\hline
\hline
章節 & 經文 (和合本修訂版)\\
\hline
4:4 & \begin{tabularx}{0.7\textwidth}{X} 你們要靠主常常喜樂。我再說，你們要喜樂。 \end{tabularx} \\ \\
[1ex]
\hline
\hline
\end{longtable}
$^{1}$弟兄姊妹.
首先要祝大家母親節快樂.
我們跟旁邊的左右說一聲.
母親節快樂.
特別是旁邊的母親.
就更加要恭祝她.
OK.
今晚這麼高興.
母親節的晚上.
在分享訊息之前.
也說一點見證.
就是今天其實也很刺激.
因為我能夠平安順利.
回到教會講道.
也是神很大的保守.
因為大家知道.
我今天下午五點鐘.
就在大埔浸信會.
主講一場的報道會.
六點半就完畢.
大家知道.
只有半個小時.
能夠平安回到教會.
特別多謝志偉弟兄在當中.
雖然他開車不是狂野奇速.
但是他開得很鎮定很穩定.
很順利回到教會.
特別沒有塞車.
這是最感恩的.
明天他們還會有四場報道會.
所以大家也繼續為我祈禱.
希望整個母親節的福音主日.
能夠一氣呵成.
很順利的完成.
今天他們五間黃昏的報道會.
也有很多新朋友出現.
雖然下雨.
我們希望在天氣不穩定的情況下.
也希望能夠多些新朋友.
可以出席母親節的福音主日.

$^{41}$今晚牧師跟大家分享的訊息.
其實只有一節經文.
當你看到這節聖經的時候.
你心想牧師今晚講一節聖經.
有些人說不夠厚.
講這麼少.
有些人很開心.
一節就真的好了.
不同人有不同的感覺.
我們一起讀一次.
這節經文你們也很熟悉.
預備 1 2 3.
你們要靠主常常喜樂.
我再說你們要喜樂.
大家知不知道.
2022年有一個調查報告.
是全世界最幸福快樂的國家.
大家知不知道是哪個國家呢.
有沒有人猜到是哪個國家呢.
全世界這麼多個國家.
哪個國家是最幸福快樂呢.
有沒有人可以猜到是哪個國家.
沒錯了.
有人答到答案了.
這個國家真的很冷門.
就是芬蘭.
原來芬蘭是全世界最幸福最快樂的國家.
接著你可能會問牧師.
香港排第幾呢.
大家知不知道香港排第幾.
報告出來是香港排第84.
比前年跌了4位.
你心想低不低一點.
可能24 34都可以.
原來不是.
原來香港是排第84.
你可能會覺得為何會這麼低呢.
都可以解釋到.
因為去年香港疫情很嚴峻.
又不能出街又要quarantine.

$^{81}$又要做個人檢測.
所以香港人很多的重擔不開心.
都可以解釋.
不過今年開始香港復上了.
又通關了.
又可以四圍走.
到底香港人的幸福快樂指數.
開心的指數會不會提升呢.
我希望我們都能夠見到香港有一個進步.
希望更加開心快樂.
一個這樣的城市.
其實你想想.
要一個人開心都這麼困難.
如果叫一個人喜樂就更加困難.
為甚麼.
因為快樂和喜樂.
你猜有沒有差別呢.
有沒有差別呢.
我敢問你.
你知道一定有差別在這裡.
為甚麼呢.
因為原來快樂是很受你外在的環境影響你的.
舉個例子來說.
你很喜歡踢球.
你踢球的時候.
一腳就散入龍門.
開心嗎.
是不是開心.
一腳就散入龍門.
但你要知道.
不是每場球賽你都能進球.
不是每場球賽你都能踢進龍門.
如果不射進龍門又如何.
那就變成好像思朗一樣.
有時候就黑面.
很不開心.
這是甚麼.
環境影響了我們快不快樂.
但喜樂就不是.
原來喜樂是從內在的心態影響的.

$^{121}$即是說你內在心態如果能夠有一個喜樂的心.
你就超越了環境的因素.
所以今天我們發覺到.
保羅在菲勒比書教導我們.
說到你們要靠主常常喜樂.
其實真的非常之寶貴.
大家知不知道.
原來在菲勒比書這本書裡.
四章雖然不是很多章節.
但竟然裡面記載.
喜樂,歡喜或者歡歡喜喜.
這個字眼有二十次這麼多出現.
而當中跟福音掛勾的.
因為有全福音而歡喜快樂而喜樂的.
都有九次這麼多.
所以你可以想像到.
這本書亦都可以叫做一本喜樂的書信.
是名符其實的.
因為菲勒比書整本書裡面.
充滿了保羅流露出來的喜樂,歡喜,感恩的心.
所以大家讀菲勒比書的時候.
你就為保羅感恩.
為什麼呢?因為保羅能夠有這份喜樂的心.
實證真的很不容易.
為什麼呢?因為當時候的保羅.
他在面對外在和內在的問題,壓力.
都影響著他.
外在是什麼問題呢?.
外在的問題就是.
因為保羅寫這本書信的時候.
他知道他在羅馬坐牢.
所以這封信又叫做監獄書信.
你知道保羅在監獄裡面.
完成這本菲勒比書.
是真的不容易的.
為什麼呢?因為你可以想像.
坐監的環境一定是很惡劣.
衛生程度一定會很差.
還有那些光線.
我相信是很微弱.

$^{161}$保羅什麼時候寫菲勒比書呢?.
我相信一定是趁著還有光線的時候.
趁著有自然光的時候,太陽光的時候.
他就把握寫這封信.
大家都知道其實保羅很喜歡讀書.
因為他很喜歡讀書卷.
無奈監房裡面的光線很微弱.
要靠自然光.
我不知道當時有沒有油燈在當中.
無論如何大家都知道.
保羅的眼睛都不多不少.
因為坐監受光線的影響.
導致保羅患了眼疾.
各位弟兄姊妹們.
一個坐監的保羅.
不單止他面對這些惡劣的情況.
而且因為坐監的緣故為基督而被囚.
大家更加發覺他是一個囚犯.
失去了尊嚴.
也失去了自由.
你可以想像到一個這樣的處境的保羅.
又怎能夠心裡面.
不要說喜樂.
能夠快樂都已經很困難.
大家認同嗎.
還更加描述自己不斷有喜樂,歡喜,喜樂,再喜樂.
所以你發覺真是令人拍案叫絕.
這是一個外在的因素.
剛才跟大家分享.
但如果說到內在的因素.
保羅更加難以喜樂出來.
因為你看到《肥臘比書》.
保羅寫這本書的時候.
他發覺原來教會都有不少問題.
特別是肥臘比教會.
雖然肥臘比教會是一間很興旺,很愛主,很有愛心.
很願意傳科的教會.
人看來是相當不錯的教會.
很棒.
這間教會很好.

$^{201}$不過雖然這間教會是非常好.
但也有一些裡面的問題.
大家是否察覺到這間教會裡面有甚麼問題呢.
保羅在第一章已經告訴我們.
原來教會裡面有些是靈性倒退的弟兄姊妹.
這些靈性倒退的弟兄姊妹.
他們雖然說去傳福音.
不過他們的傳福音.
保羅說有些是出於嫉妒.
有些是出於紛爭.
有些是出於結黨.
有些是出於結派.
你有沒有看到這些經文.
即是說有些弟兄姊妹表面上是去傳福音.
但他們帶著的心態是跟聖經的教導違背了.
不過保羅在他的分享當中.
他說不論人們帶著甚麼心態傳福音.
無論如何福音都已經被傳開了.
有沒有印象.
這就是保羅在菲勒比書描述當時.
有些靈性倒退的弟兄姊妹.
他們雖然說去傳福音.
但他們的目的動機是錯誤了.
這對一個牧師來說.
其實心裡面是相當難過的.
因為見到自己教會裡面的人.
原來有些靈性倒退了.
忘記了愛主傳福音.
是因為要遵行上帝的大使命.
來做這個工作.
各位不單止是這樣.
我想告訴大家.
當保羅寫這封菲勒比書的時候.
他又遇到第二種情況.
第二種問題.
就是教會裡面竟然有些假師父的出現.
第三章經文告訴我們.
這些假師父是甚麼呢.
這些假師父在當中告訴其他弟兄姊妹.
你們信耶穌不用信得那麼執著.

$^{241}$你們以為靠信心得到拯救.
不一定是這樣.
你們想想看.
你們看看自己的樣子.
你們都是甚麼.
你們都是猶太人.
你們要記得最重要的是甚麼.
我們的祖宗教我們要守律法.
要行國禮.
我們要爭取有好的行為.
讓人來見到.
亦要有好的行為讓上帝喜悅我們.
換句話來說.
不一定是因順稱義.
是要因我們的行為來稱義.
各位當教會有些人.
有些假師父去傳這些訊息的時候.
坐牢的保羅心裡擔不擔憂.
是很擔心的.
亦都很難過.
但他都知道甚麼.
邊長莫及.
但是不要忘記.
我所看的就是看上帝的作為.
看上帝怎樣去動功.
他的信心仍然建立在上帝的大能身上.
上帝必然會工作.
所以他不單單看眼前的困難.
難怪保羅的心都能夠超越這些困難.
還有第三.
保羅見到當時的教會出現第三個問題是甚麼呢.
就是假弟兄的出現.
甚麼是假弟兄的出現呢.
他都是弟兄.
不過姓假.
他表面說信耶穌.
不過保羅斥責他們.
保羅在第四章怎樣說呢.
說這些人的行事為人.
是基督十家的仇敵.

$^{281}$厲害吧.
原來他們的行事為人是甚麼呢.
基督十家的仇敵.
他們的神就是自己的肚腹.
他們原來在教會裡.
假借教會的名義.
利用教會.
去建立自己的利益.
利用教會去爭取自己的好處.
所以這些人表面說是基督徒.
其實是一個假弟兄.
如果你讀到這些經文的時候.
各位弟兄姊妹.
你可以想像到.
教會都有這些問題出現的時候.
你的心是否非常焦慮.
你的心也會非常難過.
在這樣的光景裡.
何有喜樂呢.
弟兄姊妹們.
我們真的不能不佩服保羅.
保羅作為神的僕人.
當他在監房裡.
無能為力的時候.
但仍然見到有問題出現的時候.
他沒有放棄.
沒有灰心.
因為他知道上帝是全能的.
上帝是掌管一切的.
上帝必然會親自去動功.
所以保羅在這樣的困難當中.
仍然擁有一個喜樂的心.
不單止他自己有喜樂的心.
你會發覺保羅更加告訴我們.
每一個做神子女的.
原來我們都可以在困境,逆境當中.
擁有一個喜樂.
這麼大的恩典.
不要失去這個恩典.
各位弟兄姊妹們.

$^{321}$不要因為我們見到眼前的困難.
人就灰心了下來.
不要因為我們見到一些不能解決的逆境.
我們就放棄了自己.
弟兄姊妹們.
我們需要學習保羅.
他提醒我們.
不單止保羅以身作則.
他亦勸勉我們.
他亦給我們有一個喜樂的秘訣.
這個喜樂的秘訣是甚麼呢?.
就是今天我們讀的經文.
首先他告訴我們.
這個秘訣是甚麼呢?.
你們要.
這個要是一個命令式.
換句話來說.
神的心意,神的旨意.
是要我們每一個屬祂的兒女.
都歡喜快樂的.
都擁有喜樂的心.
因為這個是神的旨意.
既然是神的旨意的時候.
弟兄姊妹們.
我們就需要信服神的旨意.
這個是上帝的命令.
甚至我們需要下定決心.
追求能夠信服神的心意.
以致我們在任何的境況當中.
我們都擁有一個喜樂的心.
這個就是保羅給我們第一個秘訣.
記得我們要信服.
要努力去爭取得到這個喜樂的心.
接著第二個秘訣就是.
保羅說要靠.
靠主.
這個秘訣就是甚麼?.
靠著主.
就能夠擁有這個喜樂了.
所以這裡更加提醒我們.

$^{361}$不是看人.
不是看環境.
看人,看環境.
大家都知道.
很多時候都令我們失望.
因為人五時花,六時變.
大家都很明白.
因為我們人自己.
也是有限公司.
我們很脆弱.
我們想做好一點的事.
有時候做不到出來.
大家都有這樣的經歷.
不想做的不好的事.
偏偏做了出來.
難怪保羅告訴我們.
納至行善是由得我們.
只是行出來不由得我們.
所以這是人的矛盾.
那種掙扎,那種痛苦.
保羅告訴我們.
不是靠這些.
靠這些我們肯定會很失望而回.
記得要甚麼?.
要靠主.
這個靠主是很重要的.
所以我們看到.
保羅提醒我們.
靠主為甚麼能夠擁有這個喜樂呢?.
首先因為我們的主.
祂是喜樂的源頭.
所有一切喜樂的源頭.
都是由祂而賜給我們的.
只要你和我.
我們的生命.
接駁在主的生命裡.
接駁得穩陣.
就能夠擁有喜樂.
就好像你有一個插蘇.
你想插在電掣裡.

$^{401}$你插得好.
你就能夠源源不絕.
有電力去供應.
所以當我們說靠主的時候.
因為主是喜樂的源頭.
我們要接駁在主的身上.
就好像楊福音告訴我們.
枝子需要常在葡萄樹上.
就能夠結出果子.
還有原來主也擁有豐盛的喜樂.
所以當我們和主的關係緊密密切的時候.
我們就能夠經歷從主而得那份喜樂.
各位靈子母們.
靠主.
如何來到靠法呢?.
記得原來靠主不是抽象的.
是可以很具體.
這個具體就是告訴我們.
我們追求在主的裡面.
然謙.
當我們能夠每日的生活都在主的裡面的時候.
在主裡的時候.
我們就能夠將我們的生活.
我們的生命接駁在主的身上.
弟兄姊妹們.
如何知道我們是在主的裡面呢?.
其實有兩個檢測是很明顯的.
第一個檢測就是.
有沒有學母親近神.
如果你說我的生活裡面.
從來我都不會讀聖經.
不會好好祈禱.
也不會唱詩歌.
也不會去默想.
更加很少上教會的時候.
弟兄姊妹們.
明顯我們就不是能夠學母親近神.
但如果相反.
你說我雖然很多事情要做.
很繁重.

$^{441}$但我怎樣都要每一天爭取一些時間.
我要去好好去靈修.
親近主.
我很珍惜每一次能夠有機會回教會的時候.
我都盡量去爭取.
哪怕天氣怎樣不穩定也好.
我都盡量爭取去親近神.
弟兄姊妹們.
甚至你說牧師.
我就算去外國旅行的時候.
我都記掛著上帝.
這樣就很棒了.
大家明白這個意思嗎?.
原來一個人很歡喜快樂的時候.
很容易會忘記上帝.
但你說牧師.
我就算去到任何地方.
或者我有機會去旅行的時候.
我都記掛著上帝.
當你研究有這種心境的時候.
弟兄姊妹們.
這個第一個檢測.
我們就可以過關了.
因為我們知道什麼?.
我們真的學母親近神.
我們真的知道自己是在主裡.
第二個檢測是什麼呢?.
第二個檢測.
怎知道自己是在主裡.
就是我們有沒有積極的事奉神.
弟兄姊妹們.
原來事奉上帝是基督徒的權利.
大家有沒有想過?.
不是人人都可以事奉上帝.
特別是未認識主的人.
他更加不懂什麼叫事奉.
也不會事奉上帝.
是唯有當你和我重新得救之後.
我們又發覺到.
原來我們所擁有的一切.

$^{481}$都是上帝賜給我們.
無論是我們的時間.
我們的金錢.
我們的恩賜.
我們的財綱.
都是上帝賞賜給我們的時候.
而又上帝讓我在生活裡面.
遇到有不同的服侍的機會.
那我就什麼?.
努力了.
積極的來什麼?.
去事奉上帝.
不論在教會裡.
不論你在家裡.
或者你在職場裡.
或者你在社交圈子裡.
因為上帝給你有機會去服侍的緣故.
你就爭取在不同的範疇裡.
你都要來事奉上帝.
去見證上帝.
去榮耀上帝.
等於是每一門.
當你有了這種.
很積極事奉上帝的心的時候.
我先告訴大家.
你就在主裡面.
所以這兩個檢測.
對我們每一個人來說.
是一個很明顯的測試.
所以我們在我們的生活當中.
有時候我們見到.
為什麼是信了耶穌一段時間了.
但是好像人很多時候都沒有什麼喜樂.
整個人都很.
怎麼說呢?.
愁眉苦臉.
或者整個人都很多憂心掛慮.
臉上都很憂鬱.
鬱鬱不歡.
等於我們有時候.

$^{521}$如果發覺自己失落了.
這些喜樂的時候.
可能我們都要檢視一下.
我們有沒有在主裡面.
跟主有一個緊密的關係.
我們是不是好像.
資治經常都在葡萄樹的身上.
當我們這樣去檢測的時候.
很快我們就知道.
我們是不是在主裡面.
如果我們真的在主裡面的時候.
我們可以很放心.
我們隨時隨刻.
隨時隨地都能夠領受.
從主而來的喜樂.
就好像保羅那樣.
他遇到無論外在的因素.
內在的因素.
他都能夠突破.
都能夠擁有喜樂的心.
因為他已經跟主的關係.
是密不可分.
所以弟兄姊妹們.
這個是保羅給我們一個.
非常好的提醒在當中.
還有這節聖經告訴我們.
你們要靠主.
然後是什麼.
常常喜樂.
這一招最難做.
常常喜樂.
偶爾都可以做到.
或者有很多事情都很順風順水的時候.
就容易做到.
但不是這樣的.
這個常常喜樂的意思就是說.
當你找不到有什麼喜樂的原因,理由.
或者當你經歷不到.
有什麼喜樂的感覺的時候.
但你仍然在你的意志上面.

$^{561}$你是追求喜樂的心.
這個行動就能夠令我們.
常常的喜樂.
換句話來說.
我們專心堅持去依靠上帝.
這一種心態.
是比我們的感覺來得更加重要.
所以我們就能夠什麼.
常常的喜樂.
這一種喜樂是一種意志上面.
決心上面.
就算客觀的環境都得不到.
感覺不到的時候.
但我仍然深信.
我在主裡面.
我一樣可以什麼.
擁有喜樂.
一樣讓主喜樂的力量.
充滿我整個人.
叫我無論在任何地方.
我都能夠可以活出一個喜樂的心.
貴鼎子們重要嗎.
很重要的.
如果我們有一個喜樂的心的時候.
其實你不單止自己個人得益.
我們還更加可以去見證上帝.
大家相不相信.
我以前認識一個義工.
這位遠目義工年紀都很大.
我認識他的時候都七十幾歲.
現在可能年紀更加大.
這位義工去到病房探病的時候.
是很受歡迎的.
人見人愛的.
下面那句不用說了.
為什麼呢.
原來這位義工.
他有一個花名.
他的花名是什麼呢.
就是病人契媽.

$^{601}$因為他每一次去到病房的時候.
當然年紀大.
他的衣著都很特別.
他經常都穿著長衣.
然後梳著髻.
進到病房探病.
因為這位義工的嘴唇.
他的面容常常都好像蓮子蓉的口面.
有一種喜悅的心.
當他這樣去探病的時候.
我想問大家.
病人見到喜不喜歡他.
當然喜歡他.
因為平時在病房都已經悶到不得了.
見到病人都沒有笑容.
但這位義工進來.
竟然是很有笑容.
很有喜樂的心.
所以當他去探病的時候.
其他不是他探開的都入手請他過來.
為什麼呢.
因為原來當我們有一個喜樂的心.
有一個很好的笑容的時候.
那些人喜歡靠近你嗎.
很喜歡的.
反過來你試試黑口黑面.
有沒有人靠近你.
今天你上班一進到辦公室.
老闆黑口黑面.
你馬上就自己醒目了.
找地方躲走.
做自己的事.
知道老闆心情不好.
但當你能夠有一個喜樂的心.
有一個歡愉的心.
有一個笑容.
我相信你就能夠吸引人.
能夠讓人更加樂意的去接近你.
你就能夠見證上帝.
你不需要說很多話.

$^{641}$但你的行為你的行動.
就已經讓人見到你有上帝的同在.
弟子們.
你現在明白為什麼病人契媽.
帶病人信耶穌這麼成功.
就是因為什麼.
因為他有一個喜樂的心.
人見人愛.
所以當他在病房裡面穿梭的時候.
他就能夠為住留下美好的見證.
今天你和我.
想起保羅這裡提醒的時候.
我們再一次去反省一下.
會不會我們因為太忙碌.
或者會不會我們因為太過的多焦慮.
擔心以致我們近日.
近期人好像沒什麼喜樂.
沒什麼歡愉的心.
如果是的時候.
求主幫助我們.
我們是可以擁有的.
我們不需要失去的.
只要我們記得.
今天保羅給我們一個很好的秘訣.
最後我們一起讀這節經文一次.
四章四節.
預備起.
你們要靠主常常喜樂.
我再說你們要喜樂.
我們同心祈禱.
親愛的主多謝你.
因為你已經將一個這麼大的福分.
喜樂的恩典.
放在我們的心中.
我們學校保羅的榜樣.
無論順境逆境.
無論我們身處任何的困難.
我們都相信能夠靠主得成.
靠主喜樂.
所以求天父賜給我們常常有一個喜樂的心.

$^{681}$不單止我們個人得益.
我們更加能夠榮神益人.
祈禱奉耶穌名求.
阿們.
\newpage



\section{使徒行傳 8:1-8}
\label{sec:_W_gd9p_W_Y}
\textbf{2023-05-14 午堂《福音帶來的生命改變》使徒行傳 8\_1-8 講員:曾裔貴牧師}
\newline
\newline
連結: \href{https://youtube.com/watch?v=-W_gd9p_W_Y}{\texttt{https://youtube.com/watch?v=-W\_gd9p\_W\_Y}} ~~~~ 語音日期: 2023-05-17
\newline
\newline
\hyperref[sec:PwPuQ_V9PRU]{\small{< < < PREV SERMON < < <}}
~
\hyperref[sec:index_chronic]{\small{[返順時目]}}
~
\hyperref[sec:index_scriptual]{\small{[返順卷目]}}
~
\hyperref[sec:MtaDL5MWzm0]{\small{> > > NEXT SERMON > > >}}
\newline
\newline
使徒行傳 8:1-8
\newline
\begin{longtable}{cl}
\hline
\hline
章節 & 經文 (和合本修訂版)\\
\hline
8:1 & \begin{tabularx}{0.7\textwidth}{X} 掃羅也贊同處死他。從那一天開始，耶路撒冷的教會遭受到大迫害，除了使徒以外，眾門徒都分散在猶太和撒瑪利亞各處。 \end{tabularx} \\ \\ \relax
8:2 & \begin{tabularx}{0.7\textwidth}{X} 有些虔誠的人把司提反埋葬了，為他大大哀哭。 \end{tabularx} \\ \\ \relax
8:3 & \begin{tabularx}{0.7\textwidth}{X} 掃羅卻殘害教會，挨家挨戶地進去，拉著男女關在監裡。 \end{tabularx} \\ \\ \relax
8:4 & \begin{tabularx}{0.7\textwidth}{X} 那些分散的人往各地去傳福音的道。 \end{tabularx} \\ \\ \relax
8:5 & \begin{tabularx}{0.7\textwidth}{X} 腓利下撒瑪利亞城去，向當地人宣講基督。 \end{tabularx} \\ \\ \relax
8:6 & \begin{tabularx}{0.7\textwidth}{X} 眾人都聚精會神，同心合意地聽腓利所說的話，一邊聽他的話，一邊看他所行的神蹟。 \end{tabularx} \\ \\ \relax
8:7 & \begin{tabularx}{0.7\textwidth}{X} 因為有許多人被污靈附著，那些污靈大聲呼叫，從他們身上出來；還有許多癱瘓的、瘸腿的都得了醫治。 \end{tabularx} \\ \\ \relax
8:8 & \begin{tabularx}{0.7\textwidth}{X} 那城裡，有極大的喜樂。 \end{tabularx} \\ \\
[1ex]
\hline
\hline
\end{longtable}
$^{1}$各位弟兄姊妹 早安.
母親節快樂.
跟大家一起有一個祈禱.
親愛的主.
我們等候仰望在你的面前.
的確是一個特別的日子.
紀念我們很多的母親,姊妹.
他們在家庭裡的付出.
一生的苦苦.
孕育我們每一個的弟兄姊妹.
能夠健康快樂.
尤其是在神的家裡面成長.
大家擁有共同的信仰和盼望.
今早我們更加感恩.
我們能夠有份於神國度的一個差遣.
有穎怡宣教士和她的家人.
今日都受差遣.
來香港有福音工作.
天父我們很盼望剛才的經文.
我們能夠真正體驗到.
聖靈,福音.
藉著我們面對困難.
仍然我們勉勵的宣講.
那個道帶來的生命改變.
是意想不到的.
是我們需要去體驗的.
盼望這段聖經成為我們彼此的提醒.
我們帶著的是福音.
去不同的地方.
在不同的處境.
我們務要的去傳道.
我們這樣仰望等候禱告.
奉耶穌基督的名.
阿們.
由於待會有一個相當嚴肅和興奮的差遣禮.
所以今早只能夠很短的將這段聖經.
來做一個好似一個引子.
我們今天很開心邀請到孔慶明醫生.
他是代表著非洲內地會的香港分會.
他會來簡單介紹和講一些他的經歷.

$^{41}$讓我們知道這個分會在香港的服事.
當然更加重要的是.
榮兒,攝師.
她待會都要在大家面前.
我給了她五分鐘述詞.
講一下她為何對非洲人這麼有負擔.
待會我們就會看這裡.
這段聖經其實很簡單.
當然不簡單的地方.
這一班猶太人.
令到薩瑪利亞人全個城都大有歡喜.
甚麼令到這班薩瑪利亞人歡喜呢.
是宣講基督.
猶太人去到薩瑪利亞城.
即是說他們面對逼迫.
面對耶路撒冷大規模的逼迫.
他們要離開.
他們不會選擇性去猶太的地方.
不去薩瑪利亞城.
他們不同的地方都去.
肥利是其中一位的執事.
第八章的第三節已經告訴我們.
斯提反已經殉道.
已經死了.
被石頭打死了.
但是在恐懼.
在徬徨.
但是沒有影響到這一班弟兄姊妹.
沒有試圖和他們一起.
但是他們分散.
有秩序地去到不同的地方.
能夠腳行的地方.
盡量行去到其他地方.
暫時沒有耶路撒冷的有組織的逼迫.
但是他們帶著的.
這個是我們大家真的要去問.
帶著什麼離開一個地方呢.
有沒有組織地帶著福音.
無論面對什麼困難.
去到那個地方.

$^{81}$見到那個群體.
去看這班人是自己的骨肉.
當然是神兒女的骨肉.
但其實不是的.
猶太人看馬尼亞人.
是異類.
是雜種來的.
福音書是很清楚告訴我們.
包括這些弟兄姊妹.
可能都有這個觀念.
已經是很先入為主.
但是迫迫沒有令到他們.
選擇性地向他們喜歡的群體去做.
反而他們去到那個地方.
不會覺得馬尼亞人是不需要福音.
他們將福音帶給了馬尼亞城.
很言簡意黑的作者路加.
說只是在那個城裡面宣講基督.
那個城就大有喜樂.
中間當然是夾雜著很多神蹟奇事.
有醫病有趕鬼等等.
但是重點仍然是宣講基督帶來的城市的改變.
這是很重要的一個元素.
這段聖經不可以隨便地引用.
因為第四章開始已經是大逼迫的開始.
使徒都要被坐牢.
監禁甚至要死.
但是沒有影響到他們上上下下.
三千人信主的五千人信主的.
都來到大家共同擁有一個信息.
就是福音可以改變生命.
福音是每一個民族都需要的.
他們就秉承著主耶穌升天之前的差遣和命令.
使萬民可以成為主耶穌的門徒.
所以逼迫只不過是處境的少少改變.
是這班弟兄姊妹不知道有多長時間.
已經來到去裝備好了.
在教會裡面大家一起抹餅.
大家一起敬拜.
學習真理.

$^{121}$相信不會是十幾二十年.
可能半年都沒有.
那個逼迫就要他們立即離開.
但是有執事帶著去到不同的地方.
就是這樣來到去傳揚上主的福音.
還要宣講這位改變人生命的基督.
這樣就帶來一個很戲劇性的改變.
整個城的人都來到去悔改歸主.
甚至去到一個地步就是.
連當地一個很重要的行邪術的領袖都改變.
你就可以想像福音的能力是多麼的大.
今日真是適逢其回.
我們弟兄姊妹.
我們有我們的肢體穎儀.
她自己過去無論神給她的成長背景.
她的抱負 她的負擔.
她好像情有獨鍾.
對一些我們未必經常有機會接觸的黑人在香港.
她能夠有這個愛心.
很想將基督傳給他們.
稍後八章.
其實聖靈就差遣菲利.
去接近一個埃提亞伯的太監.
正正就是非洲裔的人.
你就會想像到整段聖經是在說聖靈.
怎樣在教會裡面.
有一班人被主的愛感動.
被主的道去感動和裝備.
他們就好像隨時可以整裝待發去傳道.
逼迫只是令他們更加的.
可以去到周邊的地方.
加快那個速度.
向不同的族群去傳福音.
頂智我們要擁抱這件事.
環境可以改變我們的.
但是如果我們的內心.
大家都擁有基督這一份很深的愛.
和很穩固的那個真道.
環境只會加快我們.
將福音傳給更大更遠的地方和群體.

$^{161}$這是其中一個.
很想來到去讓大家知道的.
特別疫情過去了.
我們應該要來到去.
一起去將這個福音傳給有需要的人.
最近大家看香港的新聞就知道.
有陸陸續續不同的人.
各種的原因.
選擇了結自己的生命.
丁九橋.
大家有聽到那個消息.
那個男士因為自己患重病的兒子.
終於離世.
車子去到中途.
叫同事停下給他.
他說他想吐.
一下車就已經越走越出.
那個司機同事捉住他都捉不到.
他就一身跳出去.
就這樣了結了那個生命.
很多的痛.
很多的故事.
很多的不同的事情.
需要我們基督徒.
我們要看到生命的寶貴.
尤其是永恆生命靈魂的寶貴.
所以這段聖經是要挑戰我們.
將基督給到有需要的人.
他們就會經歷基督.
而得到那個心靈的喜樂.
我今早的講道就停在這裡.
我們有一個簡單祈禱.
親愛的主.
很盼望你的恩典.
在我們生命裡面.
不是一種我們慢慢習以為常.
而是恩典帶給我們一個挑戰.
我們要分享.
我們要傳揚基督.
你的救贖.

$^{201}$你的大愛.
你賜給人的平安.
使更多人因為認識你.
而得到這份心靈平安.
永恆盼望的福樂.
我們這樣禱告奉耶穌基督的名.
阿們.
\newpage



\section{}
\label{sec:MtaDL5MWzm0}
\textbf{2023-05-18 周四晚查經班《被擄歸回歷史研讀(15)》 講員:曾裔貴牧師}
\newline
\newline
連結: \href{https://youtube.com/watch?v=MtaDL5MWzm0}{\texttt{https://youtube.com/watch?v=MtaDL5MWzm0}} ~~~~ 語音日期: 2023-05-19
\newline
\newline
\hyperref[sec:_W_gd9p_W_Y]{\small{< < < PREV SERMON < < <}}
~
\hyperref[sec:index_chronic]{\small{[返順時目]}}
~
\hyperref[sec:index_scriptual]{\small{[返順卷目]}}
~
\hyperref[sec:We85QzJyMNQ]{\small{> > > NEXT SERMON > > >}}
\newline
\newline
$^{1}$我們一起同心祈禱,多謝天父再一次招聚我們能夠有茶經班的時間.
很盼望我們這個講授,雖然不是一個互動式的討論.
但是盼望能夠將神的話語大家一起用心的認真去研讀.
能夠很仔細的去思想當時經文的背景.
作者寫作的企劃和要帶出的目的.
我們盼望透過這樣的研讀.
可以能夠重塑整個書卷要帶給我們歸回的歷史裡面.
一些重要的訊息.
未必是我們後世的人會很賞心或者很迫切.
但是當我們看到神的救恩的進程.
我們就不得不覺得奇妙.
神是有各方面的安排.
特別是我們想到領袖的出現.
這個絕對不可能是偶然的.
有心的人為神的國度這樣擺上.
是經年累月的禱告尋求.
和將自己的生活放下.
在神的時間成為一個時代的見證人.
盼望神幫助我們.
我們仰望禱告奉耶穌基督的名.
阿們.
因為上一個星期我們講到第六章.
不得不我們看到那個結構.
我們一直都分析了第二章尼希米記.
他拿了亞特赦西王的譽旨.
那個詔書之後.
也有那些材料.
還有去到大河以西.
來去支取那些材料.
還有那些工人幫手來去督工建造.
第三章就是朝聚那些領袖.
在不同的地方.
因為耶路撒冷的四圍相當大.
所以整個構成要將它重建一個四面的圍牆.
有自己的門山.
是很大的工程.
尼希米就是先有這個譽旨.
當然更重要的是他在神那裡領袖的差遣.
接著就沒有否認到地上的政權對他的支持.
因為都仍然活在亞特赦西王之下.

$^{41}$他是一個政府官公務員.
波斯帝國的一個高官.
一個酒政.
他運用這些權的.
運用這些智慧的.
運用這些人脈的.
這個是尼希米記.
很多時候有他個人的祈禱的呼喚.
但更多是很務實的.
他的智慧.
他的策略.
他的謀略.
還有第四章我們又說到.
正正式式出現了好幾位地方的領袖.
土豪惡霸.
外邦人.
最重要是外邦人.
不是猶太人.
我們也發現.
在第五第六章.
有些猶太人因為過去很多的原因.
跟這些外邦人已經打龍通.
所以在這個藍左裡面.
第六章我們就發現.
有猶太人主動跟這些外邦.
參巴拉.
阿拉伯人的多線.
基線.
等等的人.
來阻攔尼希米.
他們知道有了城牆.
他們就真真正正是外邦人.
他們就真的跟神的聖城無分無關.
所以這個劃線的規範.
是尼希米現在要做的.
做好之後.
第七章就開始.
再一次他要找一樣東西.
就是第一次回去重建聖殿的家譜.
所以其實跟以色列記的第二章.

$^{81}$是一個.
你將併排去看.
其實差不多一樣.
經過了一百多年.
拿回那個家譜出來.
再去.
其實重點仍然是想說.
在耶路撒冷有多少合資格.
可以事奉的祭司.
利未人.
和一些在聖殿裡面工作的人.
和也發現.
數點完之後.
很現實地.
當然我說遠了.
第七章來的.
我想給大家看到那個分段.
第七章就是告訴我們.
他發現原來在耶路撒冷.
根本是人丁單薄.
所以他就要想辦法.
就去做一些遷徙.
好像長江要建水壩.
就將一些人遷走了.
因為淹過了.
那個八帝城.
很出名的.
淹過了就將.
重慶某些地方的人.
好像是上游.
淹過了就要將他遷徙了.
一樣.
現在他發現.
建好城牆之後.
本來完了.
第六章就是他個人回憶錄的完結.
但是原來利未人真正的戲肉.
就不是重建城牆.
是重建這個群體.
重建聖文.

$^{121}$所以聖結的指紋.
是由第七章.
一直到十三章.
所以開始就會記載很多的嚴肅會.
首先就是他要遷徙那些.
在周邊地方的猶太人.
回到耶路撒冷.
這個就是大概的分段.
第六章就夾雜在這裡.
其實很多東西需要大家一起去思考.
要去用一個很清潔的心靈.
去思考.
今天我們的場景就是教會.
當時他們就是在聖城耶路撒冷.
現在他們要做的一件事就是.
很清楚的.
神要.
神祂的指紋猶太人.
要回到來.
不只是重建聖殿.
不只是建好那個硬件.
而是要有一個真正的敬拜的生活.
所以利未人就首先.
第一件事.
都是做好那個硬件.
就是要規範那個聖城周圍.
周邊.
是要有自己的城牆.
有自己的門樓.
接著.
第一時間就立刻想到.
安息日.
不可以隨便開門.
這個是很實在的.
你去第七章第八章就立刻浮現.
為什麼這麼緊張呢.
為什麼這麼緊張安息日呢.
就是要希望帶領神的指紋.
回到那個敬拜的焦點和中心.
就是真真正正.

$^{161}$心靈誠實的敬拜.
所以你會看到.
原來真正利未想告訴我們.
不是他的領導力很強.
而是告訴我們.
他有一個他自己很清楚的使命.
不只是重建的一個builder.
而是他配合.
這個以斯拉文士.
來去重建.
怎樣重建聖文的心呢.
很簡單的.
就是宣讀律法書.
在一些的嚴肅會.
因為第七章開始.
就立刻發現.
七月初一.
記載這七月份三個重要的節期.
七月初一.
催國節.
七月初十.
贖罪日.
七月十五號開始.
一直延伸到二十一號.
二十二號就是特別的.
第八天的嚴肅會.
就是駐彭節.
所以有這三個大的節期.
很多事情可以做.
就是在這些嚴肅會裡面.
用神的道去感動.
挑戰神的選民的心靈.
就真的帶來復興了.
就真的帶來他們有立約了.
所以第七章.
去到第十三章.
其實是幾個的嚴肅會.
但是第六章.
我們上個星期點出了一點.
也開始引了新約.

$^{201}$主耶穌都面對那些誹謗.
挑剔.
惡意的重傷.
在搬弄是非的挑撥離間.
你就可以想像到舊約新約.
都是這樣出現的.
所以作為領袖.
必須要有一個心理準備.
當你肩負起一個很大的使命.
責任的時候.
四方八面的那些攔阻.
尤其是透過言語.
對事公的批評.
即雞蛋裡面挑骨頭的.
那種惡意的重傷.
我們讀第六章就會發現.
都幾恐怖和幾令人痛心.
尤其是猶太人和多比亞.
本身是有親戚關係的.
所以你會看到.
我們現在開始看第六章.
再一次看這裡.
參巴拉多比亞.
阿拉伯人基善和我們其餘的仇敵.
聽見我已經收完了城牆.
其中沒有破壞之處的時候.
當然那時還未安門線.
參巴拉和基善就打發人來見我.
說請你來.
我們在阿羅平原的一個村莊相會.
他們卻想害我.
即是說引尼歇米不設防地.
去到一個偏僻的村莊.
可能就是洗還手暗殺他.
即是說殺害波斯帝國的官員.
可能希望他不設防的.
沒有戒心的.
來跌落那個陷阱裡面.
這個可能是陷害.
或者是有其他的問題.

$^{241}$於是我差遣人去見他們說.
我現在辦理大功.
不能下去.
應能停工下去見你們.
他們這樣四次打發人來見我.
我都如此回答他們.
很堅定地否決了他們這個善意的.
說尋求會議.
秘密的協商.
甚至乎趁他下去的時候.
陷害他甚至乎害他.
保羅試過的.
在哥林多.
《小徒行傳》十八章.
不是十八章.
十八章他還是在哥林多.
應該是二十二章的經文.
在《小徒行傳》.
他被遞解到哥林多的時候.
有一些猶太人.
密謀趁他去哪裡的時候.
就要來下手害他.
但是被一個八夫長知道.
所以就來回報給千夫長知道.
所以有千夫長和他的兵兵.
就壓戒保羅.
上去千夫長的衙門.
免於這一次暗殺的行動.
一個傳道人都要被暗殺.
你可以想像得到.
多不容易.
尼熙米是位居亞達舍西王的.
一個酒政.
不過現在山高皇帝遠.
你去到河西.
你是一個猶大省的省長.
有這個官職.
我們講過了第五章.
他不要那些省長的俸祿.
不想加增當地人.

$^{281}$太重的那個謠譯.
要招呼這些地方官的飲食.
可能就是他們法例規定.
由皇帝轄下的那些行省.
那些省去供應.
這些情況.
所以尼熙米對內.
他安撫百姓的矛盾.
因為各樣的貧富懸殊.
第五章就是那個問題.
他解決了.
現在第六章.
已經去到白熱化.
好像垂死的掙扎.
對尼熙米的一個攻擊.
你可以想像得到.
尼熙米運用他的智慧.
運用他的常識.
運用他的謀略.
來避免了這次暗殺的行動.
第五節.
參巴拉第五次打發僕人來見我.
手裡面拿著微風的信.
信上寫著說.
外邦人中有風聲.
這些捏造的偽名的.
冒名的.
甚至乎是不署名的信.
不知道大家有沒有收過.
我在任內都有幾封收過.
信上寫著說.
外邦人中有風聲.
加斯慕基善也說.
也說即是認同這封信裡面的內容.
你和猶大人謀反.
修造城牆.
你要作他們的王.
你又派先知在耶路撒冷.
指著你宣講.
有君權神授.

$^{321}$好像有天意.
有一些神蹟.
諸如此類.
甚至乎有些先知在說話.
尼熙米你是我們的真命天子.
你是我們猶大省的救星.
類似這樣.
這樣令尼熙米.
你不可以不害怕.
你不可以不停工.
現在涉及的你背叛亞達斯西王.
等於保羅也是這樣.
被人這樣質問他.
你教導外邦人.
不守摩西的律法書.
不去做一個敬虔的.
應該說教導外邦的猶太信徒.
不需要去守摩西的律法.
這是一個很大的誤解.
但其實是一個很嚴重的指控.
對保羅來說.
所以你會看到.
無論舊約新約.
你都見到這是第五次.
這次還有一封無名的信.
沒有封口的信.
來到他呼之欲出.
我們四次找你去談.
希望幫助你解決現在傳得很熱門的.
這樣的一個指控.
你是不是真的這樣情況.
是我們一起幫你去梳理.
這樣的一個.
特別是謀反.
你可以想像得到是多麼嚴重的事.
說他們現在收造成長.
即是說自己在那裡做山寨王.
在耶路撒冷不是猶大省長.
猶大省長是由亞達薩西王所委派的.
因為是亡國的.

$^{361}$委任你的.
但現在你不僅不甘願只是做省長.
你希望成為以色列猶太人的一個領袖.
一個真正有權有勢力的一個領袖.
這樣可以是很大件事.
要告訴我們.
第七次你又派先知在耶路撒冷指著你宣講.
說在猶大有王.
沒錯了.
真的這樣.
好像來到去假借先知的口.
來到去君權神授.
即是說得到民心.
百姓的心完全歸向你.
對外邦波斯帝國的君王來說.
是一個威脅.
我派你去做省長.
建造省長.
你現在去聯絡你的同胞.
要來推翻波斯帝國.
是很嚴重的指控.
所以第六章.
你看到李敏的回憶錄.
真的很齊齊的.
第三章是他的同胞.
他的領袖.
聽他說.
你這個意象.
你這個願景.
我們都很認同.
很接受你這個想法.
我們一齊跟你做.
第四章就開始出現那些單單打打.
那些阻撓.
但是都不敢明目張膽.
但是現在去到最後的階段.
他們好像最後的垂死一擊.
所以怎樣能夠瓦解這樣的謠言.
是不是燒了封信就行呢.
到底李希敏在百姓的心靈裡.

$^{401}$是怎樣呢.
還有一些領袖.
是不是真的認識李希敏.
對他的為人.
有一個認同和信任.
而來聽到很多謠言.
可以幫他去說.
你們這樣說.
一方面是你們捏造.
另一方面就是.
我認識李希敏的為人.
應該直接.
如果你說不是.
有背後的信.
有信.
是真的有人這樣說.
最好就是.
跟你拿一封信去問當事人.
那就最好了.
所以周邊的弟兄姊妹.
特別是信徒領袖.
是必須要站在教會這個層面.
不是要包庇任何一個人.
而是要為整個團隊的守望.
如果不是就很嚴重.
我都聽聞是.
我都覺得是.
上到來很氣焰.
又要我做這裡做那裡.
做哪個城牆.
我本來都不夠什麼.
很多這些說話.
突然間從信裡浮現出來.
你可以想而知.
所要面對的.
不只是李希敏一個人.
而是整個團隊.
整個耶路撒冷.
一起跟他同心.
才能夠做得起那件事.

$^{441}$第八節.
我就差遣人去見他.
就是這個參巴拉.
說.
你所說的這事一概沒有.
是你心裡裂做的.
他們都要使得我們懼怕.
意思說.
他們的手必軟弱.
記住他們使得我們懼怕.
不只是李希敏.
是整個團隊.
是整個建造工程的.
上上下下的不同的團隊.
我們有的是唐董執事團長.
主學老師.
或者有師事長.
諸如此類的.
有師班長.
有很多不同的領袖.
有些資深的弟兄姊妹.
其實我們需要定期的.
有分享.
有溝通.
有一起在愛裡面的.
彼此的坦誠.
不要有很多的留言.
不是留下的言語.
而是那些是是非非.
就要澄清.
甚至乎.
是要.
新約告訴我們.
甚至乎要督責那些.
說是非裂做壞話的人.
你可以想像在這裡.
在這裡.
李希敏.
將整個團隊擺出來.
你想恐嚇我.

$^{481}$我們是齊心去建造的.
我們沒有想過.
我們整好城牆之後.
他們要推舉我做王.
這個指控.
其實某個程度.
應該你都要預計到的.
因為真是山高皇帝遠.
阿達什西王.
真是很偏遠的地方.
我們信任你這個.
亡了國的猶大人.
李希敏.
給你材料.
上去讓你和你的民族.
可以心安理得的敬拜神.
所以准許你們建造.
其實原本就已經有的城牆.
不是額外加多給你.
不過給你們能夠.
真正按摩西的律法書.
君王不知道你什麼律法書.
總之按你們認為.
什麼是最正確的敬拜.
就准許你們有敬拜.
其實其他民族一樣.
所以李希敏停在這裡.
他沒有一個歪念.
沒有一個野心.
他假借王的命令.
實際上他自己想立王.
如果有這個歪念.
很容易有的.
因為你上去做領袖.
人家就是上了床.
想拿被cum 的.
那些猶太人就想.
我們一直都沒有王.
沒有王管.
現在有你在這裡就更加好了.

$^{521}$我們推舉你做王.
有什麼所謂呢.
所以是一個很微妙的.
領袖和權力.
永遠是一線之隔.
亂斬權力.
往往是領袖一個很嚴重的致命傷.
一個致命的一個.
因為你有權.
自然就會有這個野心.
你沒有那種野心.
野心可以是一種推動力.
野心可以是一種很強的.
自發性的一種力量.
但是野心可以是.
把心一橫的一種自己.
要做那個真正的獨裁領袖.
所以很多時候是一線之隔.
李希敏就是微妙地.
在這個處境裡面.
他自己拿捏得很好.
這個是在字裡行間浮現出來.
李希敏他自己這個人的品格.
那種的高尚.
連他的跟隨者.
就是那些跟班.
你等一下留意一下.
上個星期我們讀過一次.
是忠誠的.
不偏不倚的.
不是為了權力.
不是為了拆鞋的跟班.
而是敬畏神過於一切.
等一下我讀給大家聽.
我們看看第九節.
他們使我們懼怕.
其實沒有.
意思就是他們的手必定會軟弱.
這個是多比亞的想法.
以至工作不能夠成就.

$^{561}$神求你堅固我的手.
他的祈禱.
加上的.
所以你看到李希敏的記.
他不是徹夜禱告.
只是一個呼喊.
但是更多的就是.
他祈禱完之後.
他的思想是很精密.
是很謹慎.
是很小心.
是步步為營.
是知道作為領袖.
這個這麼大的工程.
不是那個工作難.
是敵對的勢力很大.
他很知道.
他不是現在有一個上方寶劍.
的確是的.
大肆西王命令他.
河西的省長.
當地的總督.
是要給材料的.
是要開綠燈的.
但是他不是想用這個姿態.
而是他希望自己同胞.
用皇給的資源材料.
來建造.
所以你看到聖經.
沒有否定到人為性的參與.
到底是神的時間.
神的旨意.
神職.
還是人的百分之一百的參與.
其實應該是百分之二百.
神的工作.
人的參與.
就是百分之二百.
但其實也是百分之一百.
就是沒有否定到要祈禱.

$^{601}$要依靠神.
但是利謙偉記就更多.
就是神的感動.
這個領袖.
他有一個很健全的心態.
他的心靈.
是很清潔.
很純潔的動機.
心靈就管那個動機.
心靈是管我們內心的計謀.
但是頭腦呢.
就是管你的人生閱歷.
你的智慧.
你的知識.
所以很多時候就是經一事.
你應該要掌一智.
經一事你都不掌一智.
那你就是很愚蠢的了.
我們應該是要吸取教訓.
很深刻的就是朱牧師經常提醒我.
要找一些板來撞.
不要亂撞板.
有些輕的沒所謂.
但是有些嚴重的就不行.
這是很深刻的提醒.
他有時候每次回來都會約我去喝咖啡.
當然我不喝.
我喝茶.
但是很重的.
很重要的.
所以李希敏呢.
他是很健全的思想.
就是腦袋和心靈.
很知道自己欺負其分.
他只是暫時被賦予猶大的省長這個官階.
實際上他知道.
其實他是神的僕人.
重建聖文的心靈.
恢復一個真正的敬拜.
所以這樣就已經是最好了.

$^{641}$最足夠了.
他如果想多了.
哇.
左右逢源.
王又信任我.
百姓又擁戴我.
我在這個建造好之後.
有時候那種沾沾自喜.
飄飄然的成功.
哇.
你看看.
百多年都沒有人做到.
我就做到了.
二十年都不用.
我就搞好了整個耶路撒冷.
小心.
這個就是權力會令人的思想腐化.
我們回頭看這裡的時候.
第十節.
我到了米希大別的孫子.
蒂雷瓦的兒子.
是瑪雅嘉里.
那時呢.
他閉門不出.
他說.
我們不如在神的殿裡會面.
第二次的危機.
就是祭司居然邀請尼希米.
去到在聖殿裡面.
密謀一些重要的會議.
好像現在李輝這個特使.
去烏克蘭開和談會.
希望拉攏.
俄烏能夠坐下來談和平方案.
現在呢.
這個瑪雅.
突然間在他家裡不出來.
就說不如我們不要在家裡.
去聖殿裡面談.
將電門關鎖.

$^{681}$因為他們要來殺你.
尼希米.
就是夜裡他們會來殺你.
第十一節.
像我這樣的人.
豈要逃跑呢.
像我這樣的人.
豈能夠進入殿裡保存生命呢.
我不進去.
進聖殿會死的.
有沒有記得.
鎖羅門的一個將軍.
突然間忘記了.
大衛的那個.
藥丫.
走進去聖殿.
抓住那個彈角.
但是鎖羅門說.
都要可以拉出來.
將他殺了.
聖殿不是給我們這樣用的.
現在尼希米.
他一方面他根本不怕.
因為這個只是一個恐嚇.
而且又是這樣.
納粹吸該殺.
可以不可以.
兩難的問題.
現在尼希米就要進去聖殿.
你不是祭司.
更加不是憲制的日子.
你走去聖殿.
還要過夜.
是死罪.
對他來說.
他覺得他不能夠逾越摩西的律法書.
對於祭司的那個教導.
就是說.
他很尊重.
祭司是神的聖殿裡的.

$^{721}$唯一的合法的工人.
他是一個信徒的領袖.
他是很厲害.
他是肩負起一個很重大的使命.
他來到去召集群眾.
重建神的聖牆.
重點不是他是一個建築師.
重點是他知道.
這班人有了敬拜才是最根本.
是最重要.
所以最終.
其實跟著他就沒有什麼角色.
他的角色就是交給以色列.
來召開嚴肅會.
你可以想像得到.
他不會逾越進入聖殿.
這是一個很大的.
很陰險的一個托詞.
由一個祭司的口中說出來.
和我一起去聖殿.
我們保護你.
他們不敢在神的殿裡.
來得罪神.
諸如此類這些包裝的說話.
黎希敏又是洞悉了他們的問題.
然後直接的直斥其非.
他這個人怎麼可以進到聖殿裡面.
來保命.
我不進去.
第十二節留意一下.
我看明神沒有差遣他.
是他自己說者說攻擊我.
是多比亞和參巴拉.
賄賂了他.
收買了他.
來這樣做.
賄賣他的緣故.
是要叫我懼怕.
依從他犯罪.
他們很全然惡言.

$^{761}$毀旁我.
進了聖殿你就死了.
你是一個領袖.
現在你沒有大沒小的去做這件事.
記不記得有一個君王呢.
這個我記得了.
烏西亞王.
他最成功最輝煌的時候.
突然間心高氣傲.
打仗很厲害.
他很喜歡農耕的農地農事.
種花.
搞很多果園.
突然告訴我們.
他心高氣傲的時候.
進去聖殿燒香.
進聖殿敬拜.
是敬拜的日子.
沒有問題.
所謂燒香.
是應該說.
拜託祭司幫我.
來行憲制的條例.
然後代表著我去敬拜神.
他去擔當祭司的職位.
來燒香.
那些祭司很害怕.
叫他不要.
他不聽.
真的繼續在燒.
然後突然間.
兩隻手生了大麻風.
那些祭司就很害怕.
馬上拉他出聖殿.
當然不可能.
過個時代.
過個禮拜.
馬上會有人生大麻風.
這樣就容易了.
突然間見到阮微.

$^{801}$生大麻風.
就是神擊打.
不會的了.
一般現在都不會的了.
很少會有這樣的情況.
一定是.
我們要有一個.
教會要明白.
有專職的神的僕人.
是經過了很嚴格的訓練.
甚至是蒙召.
和肩負起牧羊的聖職.
大家一定要在崗位身份上.
是越來越模糊的.
因為現在越來越多信徒領袖.
去讀神學.
神學院也開辦了很多晚間的神學課程.
現在是晚間課程.
很穩定.
很多收入.
去補貼日間的.
全時間蒙召的神學生.
是越來越少的.
有些小的神學院.
根本連修生都不夠.
你就會看到.
有時那個界線是很模糊的.
但黎希米沒有這個模糊的自己.
他對這個界線很清楚.
我這個人.
怎可以進到聖殿呢.
因為他馬上想到.
你們一會兒就毀謗我.
一會兒就在我背後.
就說黎希米你真的很嚴重.
你居然間進到聖殿裡面.
這樣來放肆的一個這樣的人.
你可以想像得到.
原來有時真是.
聖經說我們不要兩邊說話.

$^{841}$快點來聖殿我保護你.
然後轉過頭來說.
黎希米你進了聖殿.
你原來真是將你的內心.
那種的心高氣傲.
完全表揚出來.
聖殿在你的眼中.
原來你是可以很隨便的.
人就是這樣兩邊說話.
所以聖經是很禁止這樣的情況.
一來就是因為他.
收了賄賂.
律法書有說的.
收了賄賂之後.
就會令到一個人.
那個眼睛就不夠明亮.
因為你已經有.
已經有了那個打架波.
如丸打架波.
你一收了錢.
是要放軟手腳的.
今場一定要輸的.
輸得好像很盡力的.
這樣整件事就變成一個混亂.
所以這裡的.
好意善意的叫他進聖殿.
所以黎希米又是要有.
很清醒的頭腦.
原來你身邊的人.
還要是在你身邊的一些領袖.
你沒有理由想著領袖會害你.
正正他就是害你.
轉頭他就說你.
你這個領袖黎希米.
居然這麼大膽.
進到聖殿.
在那裡來去.
真是沒有大沒小的.
當我們這些祭司.
好像沒有身份和沒有那個夢照.

$^{881}$你可以想像得到.
黎希米這樣記載這些.
其實只不過是例子.
列舉幾件.
其實每天都應該收到很多的投訴.
匿名信.
單打的說話.
很多這樣的明刀暗箭.
這樣來對他攻擊.
我們還做不做這個領袖呢.
所以我都說笑了.
其實不是說笑.
大部分香港的教會.
最難找的都不是剛出來的傳道.
是沒有人希望和喜歡做堂主任.
是真的.
真的沒有人喜歡做.
有些是去到第二個梯隊.
要升他做堂主任.
講明我不做.
一定不做堂主任.
我認識的都是這樣.
你開想而知.
要面對的就是.
真的好像千呼所指的.
這樣又說不對.
那樣又說做得不好.
很多的無謂的指控.
是令到本來鞠躬盡瘁去承擔的.
都灰心了.
都算了.
我一來都不是.
你有突然間十倍的待遇.
不是的時候.
我為什麼要招致無憂子的.
那種甚至是背後的那些謾罵呢.
這是一個惡性的循環.
這裡告訴我們.
他們就是全愕然去毀謗黎希米.
十四節.

$^{921}$我的神.
我很喜歡黎希米其中一個.
當然我們講的整套歸回歷史.
和歸回的先知書.
但是黎希米其中一個特色.
大家都讀到的.
每一個風波之後.
我的神啊我的神啊.
他不只是用他的思維去應對.
應對其實很累的.
你的心靈一次又是這樣.
兩次又是這樣.
每次都說我不對.
每次都是亂說的.
心靈真的很累的.
但是他的出路就是.
不只是看心理輔導.
要吃安眠藥.
而是他真的回到神那裡.
我的神啊我的神啊.
神是他的神.
是神你要我做.
不是我想做.
是神你要我安排我這個位置.
本來我在波斯帝國是做酒精.
我現在不是做其他事.
我做酒精.
我每天都應該是一個開心的場合.
你想想.
試一下用兩個公種.
大家就明白了.
你喜歡郵差多.
還是喜歡一些其他的公務員多.
有些是譬如說稅局的.
或者做警察的.
曾經社會事件的時候.
做警察就.
其實我沒有做錯任何事.
都被人針對的.
但是郵差呢.

$^{961}$永遠不會有人針對他的.
我們不算現在有電郵.
有WhatsApp的年代.
算以前.
沒有WhatsApp.
沒有電話的年代.
你看到郵差其實很開心的.
我小時候住在我屋區.
那個郵差是每一間.
每一戶都認識的.
我們走到上山腳的.
由九龍塘的沙福道的地鐵站.
走到上去城市大學.
斜路的達芝路.
你想想多遠啊.
但是那時候是木屋區.
是九曲十三彎.
那個郵差做到退休.
是認識我們家家戶戶的.
所以我們每個人見到他都很開心的.
給他水喝.
給他水果吃.
跟我們很熟的.
所以你想想.
你做那個工作.
你每天一上班就很開心.
是重的.
辛苦的.
背著那些郵包.
那些郵包很長的.
有些狗等等.
但是一定是每天很有使命.
很有開心.
就是我將那些信.
誰誰.
我根本都認識的.
我派給那個家家戶戶.
很開心的工作.
但是你可以想像得到.
尼欽米面對的.

$^{1001}$就是這些的毀謗.
但是他有的就是神.
我的神.
多比亞 參巴拉.
呂先知 羅亞底.
這個說.
似乎是猶太人的呂先知.
和其餘的先知.
要叫我懼怕.
求你紀念他們所行的這些事情.
每一次提出一個風波.
提出一個指控.
馬上有一個最後的註腳就是.
他的出路就是我的神.
我的神.
所以真的很多時候.
有時候都百詞莫辯的時候.
回到神那裡.
神是我們的出路.
神是知道的.
他知道.
你還介意什麼呢.
這是一個很重要的領袖的價值觀.
領袖需要去學習的一個.
你可以說是一個.
一個心靈裡面的一個氣節.
等等.
第十五節.
以六月.
就是六月.
二十五日.
六月就要大家有些mindset.
你要用猶太人的想法.
去到猶太人的六月近尾.
所以以六月.
六月是猶太人的.
好像阿拉伯人的齋戒月.
因為準備等候七月初一.
就是他們的催國節.
所以七月是很重要的.

$^{1041}$一個猶太的節期的循環.
所以六月收好城牆.
準備禁食.
準備等候神的恩典和祝福.
雖然不是像所羅門那麼戲劇性.
第一次的聖殿建好之後.
記不記得有什麼特別事情發生.
所羅門獻祭完之後.
天火焚聖殿.
燒那些祭物.
還有聽所羅門的祈禱.
第一次當然是.
一定有一個里程碑.
不是每次都是這樣.
所以這裡說以六月的二十五日.
城牆收完了.
共收了五十二日.
我們一切的仇敵.
四圍的外邦人聽見了.
就懼怕了.
愁眉不展.
因為見這個工作完成.
是出乎我們的神.
沒有將榮耀歸於自己.
是做了很多事.
不是沒有做事的.
不是坐著翹著手.
神叫你做事.
是感動你們這一班的領袖.
里程碑和當地的領袖.
來做很多的抵抗.
很多的陰險.
你要逐樣逐樣的拆解.
所以領袖就是每天都好像.
做一個排憂解難的工作.
Problem solving.
和危機處理.
這樣來將那個問題.
最不受影響地解決了.
希望更加是法理程解決了.

$^{1081}$所以這裡告訴我們.
因為見這個工作完成.
是出乎我們的神.
所以不是一些恭維的話.
是真的彌起彌彼.
深深的感受到.
這是神的作為.
在他的身上的工作.
第十七節.
在那些日子.
猶大的季就屢次寄信於多比亞.
多比亞也來信於他們.
在猶大有許多人與多比亞結盟.
這些稍後就會有首尾跟的.
就是涉及到那些和外邦人通婚那裡.
也是以色列寄信.
曾經都出現過.
就是以色列要他們立誓.
簽那個好像悔過書那樣.
其實是休書.
決定休他們的妻子.
和叫他們的子女生了.
即是半堂回家.
都要離開.
以色列我們之前說過了.
就是第十章那裡.
現在前面後面都會有這個.
因為現在就先點出.
原來問題的癥結.
原來過去因為沒有人理.
猶太人和外邦人.
那個過從甚密.
就有許多的結盟.
許多的瓜葛.
可能商業的瓜葛都有.
安息日就很混亂.
你沒有了冒險雞籠.
外邦人就守什麼節.
不要那麼執著.
安息日就安息日.

$^{1121}$睡晚一點有什麼所謂.
少一次的敬拜有什麼問題呢.
安息日就是最重要的買賣.
許多這些似是而非的東西.
慢慢猶太人就沒有了界線.
根本都沒有城牆.
都沒有了山門.
沒有了大家一起守嚴肅會.
那個氣氛.
那你何來.
怎樣能夠將這些外邦人趕走呢.
所以慢慢慢慢.
在過去的日子.
你想想就是說.
冰封三尺.
非一日之寒.
現在不是突然間上來才結盟.
是一直以來.
猶太人因為沒有了自己的神學.
就沒有了界線.
譬如我們昨晚的祈禱會.
我都留意到.
因為我會找那些新聞.
台灣的同性婚姻已經過了.
但是現在呢.
以前就是要由其中一方的子女.
他可以收養.
但是現在呢.
兩個他們所謂的配偶.
就可以收他方.
就是和自己沒有任何血緣的.
就是正式開綠燈的收養.
這樣就禍害很大了.
我儘管這樣說.
那個影響性是很大了.
很深遠了.
你兩個一定會灌輸.
你們認為的婚姻價值觀給下一代.
就是說慢慢慢慢的潛移默化.
下一代已經沒有這些界線.

$^{1161}$現在說的那個界線的問題.
其實真的很不容易了.
你可以想像得到亞洲人.
應該說中國兩地三岸.
台灣就走得很前了.
對於這個所謂的要持平.
不要有歧視.
就走得很前了.
對下一代的那個價值觀的影響.
是很深遠的.
這些是我們擔心的東西.
這個也不是不人道的問題.
為什麼我不可以收養.
因為首先要將倫理打沉了.
打低了倫理.
沒有倫理.
不需要倫理.
我們兩個是好人.
說什麼倫理呢.
說什麼聖經的倫理呢.
打斷了這個東西之後.
什麼都說人權了.
為什麼我們不可以做呢.
你們可以做.
為什麼我不可以做呢.
我們都很想用我們的愛去收養.
為什麼不可以呢.
現在又可以了.
又可以開通了.
你可以想像得到.
將那個策略一代一代地滲進去.
不要有價值觀.
不要有界線.
不要有難阻.
人是每個人都是一樣的.
不要聽聖經的那個界線.
意思就是這樣的詮釋.
所以這裡告訴我們.
多比亞的結盟.
因為他是阿拉的兒子.

$^{1201}$是加尼的女婿.
並且他的兒子約哈蘭.
這些是猶太人的領袖.
這些是.
祭司應該是這麼說.
娶了比利加兒子米舒蘭的女兒為妻.
他們常在我面前說多比亞的善行.
剛剛說的是那些惡行.
但是這些領袖.
因為原來跟他們結了盟的.
有了一些可能是親戚關係的.
其實不是另眼相看.
是已經有偏頗.
是有一個自己的前設.
人最怕就是有了自己的前設.
有了自己的很強的那個執著.
和背後的利益.
利益行先.
大於一切.
就什麼都當看不到.
我很喜歡一句說話.
當然是一些古書.
他說好而知其惡.
善而知其.
惡而知其善者.
天下鮮而古文.
你喜歡一個人.
但是你知道他的惡.
你很喜歡他.
你和他是好老友.
但是你能夠指出他的問題.
不會因為是老友包庇他.
惡而知其善.
你不喜歡那個人.
但是你從他身上找到優點.
稱讚他.
不會有色眼鏡.
問題是人心有色眼鏡.
古人都告訴我們.
天下鮮而知其少之又少.

$^{1241}$太罕有.
這樣的人太罕有.
就是說對事不對人.
是我們中國人說的.
很難對事不對人.
往往是對人不對事.
更好的就是.
對人對事.
更好的就是對事不對人.
所以好而知其惡.
惡而知其善.
天下鮮而.
古人都提醒我們.
我們要真的像主耶穌那樣.
罵門徒罵得很多.
真的愛他們.
點出他們的問題.
給得撒旦退我後面去.
給得拉住耶穌.
你不要上十字架.
你可以想像到.
主耶穌對不同的人.
他真的出了一種很深的愛.
要提攜他.
要點出他的盲點.
這個往往有時我們就是.
面子的文化.
中國人的面子文化.
報喜不報憂.
講好話.
背後講壞話.
太多這些我們的價值觀.
有時很扭曲的人性.
所以其實很需要我們幫助.
我們能夠.
不是對人講人話.
對鬼講鬼話.
不是這個意思.
當然要有智慧.
與外人交往.

$^{1281}$保羅說.
我們要隨時說造就人的好話.
叫聽見的人得到益處.
造就人的好話.
其實可以是責備.
出於一種有智慧的責備.
令到對方.
不只是覺得喜降.
而是他感覺.
你這樣講.
我真的沒有這樣想過.
你這樣提我這個點.
是很對的.
為什麼這麼遲都沒有人提我呢.
你真的幫我很多.
你一點我就完全明白了.
要這樣說好話.
要隨時說造就人的好話.
叫聽見的人得到益處.
得益處不是說.
你真是厲害.
你真是好.
你沒有什麼問題.
這樣其實是害死他.
明明很多問題你都不告訴他.
但如果我們能夠隨時.
即是說你很有智慧地.
在一些不適合的場合.
你就不應該講.
適合的場合.
你就要找他坐在一旁.
平心靜氣地告訴他.
讓他真的學到.
他的內心.
這是很高的智慧哲學.
是我們一生都要學習.
但我們希望我們自己.
就是一個這樣的知音人.
別人這樣跟我們說.
我們先不要有一個捉迷藏.

$^{1321}$我們應該要去聆聽.
人越大越有地位.
就越難聆聽.
越難去聽.
和苦心自問的那種揣摩.
當然說的人是需要有智慧.
不是說你值得.
人都聽不明白.
聽不入.
應該真是很善意.
所以聖經真是很好.
隨時隨著那件事.
隨著那個時間.
去講一些可以造就他的好的說話.
叫到那個人得到那個益處.
益處就是他可以改變.
他可以有生命的更新.
要這樣.
所以我們應該要小心我們的說話.
如果你經常都說笑.
每個人聽你說話.
都不會覺得你是認真的.
你就不夠份量.
含金量不夠.
但如果你是一個很認真的人.
又一個很謹慎.
平時你表達自己.
是很愛對方.
愛其他弟兄姊妹.
你講話.
講一些很嚴肅的說話.
別人就會覺得不是小事.
和你是認真的.
其實我們需要在一個認真的裡面.
這個就是.
李希敏給我們見到的.
這個外邦人.
垂死的最後一擊.
希望將李希敏的信心打沉.
用一些謠言.

$^{1361}$用一些無名的偽造的信件.
托持.
說他要做王.
又說有先知來報夢.
說李希敏就是我們的真命天子.
又叫祭司.
籠絡了一些祭司.
叫李希敏進去聖殿.
躲藏一段時間.
避開現在的暗殺等等.
再加上這些祭司.
他們有很多的惡言.
來毀謗李希敏.
但見到這個工作完成.
但未真正完全的.
連絡好.
城門都.
還有一些要修理.
第七章.
我先喝口水.
大家聽完之後.
有沒有什麼.
特別是要舉一反三.
在教會裡面.
我們和不同的兄弟姐妹的交往.
我們是要以教會.
作為一個很重要的.
神的居所.
所以有不同的意見.
其實更加好.
是真是意見.
大家拿出來傾.
去溝通.
有什麼不清楚.
直接找當事人.
是善意的.
不是要挑戰的.
大家一齊溝通好.
最近聽到你這樣.
有些人這樣說你.

$^{1401}$給機會別人申辯.
這是一個很好的溝通.
不要將一些東西.
好像炒在一起.
一定是這樣.
因為人很容易.
先人為主的觀念.
越積越厚.
其實很容易.
你的眼光又會不同了.
這個就再牽涉到.
涉及到你的內心.
你就抒發不到你的內心.
聽著聽著的東西.
就會對那個人.
不單止有偏見之後.
就開始有一個惡的印象.
那個印象一路累積.
就變成一個憤怒.
所以教會很多時候.
就是人際關係的問題.
令到那些問題瓦解.
都是由很多很多的誤會.
一路這樣累積.
還有新約告訴我們.
末世那個不解怨.
不解怨就是說.
沒有想過和你解除那個誤會.
你想想新約這樣告訴我們.
教我們提摩太.
你要小心.
這是一個不解怨的世代.
所以你更加要將神的道講好.
你想想.
我沒有想過.
我要和任何人解除那些誤會.
其實是不是很可惜呢.
本來這個地方.
神的家.
任何一個教會.

$^{1441}$應該是最好.
不可能有任何瓜葛.
任何糾紛的地方.
世界沒辦法.
因為有利益衝突.
我們經常這樣認為.
但往往教會令到很多人.
帶著很多傷痕來聚會.
甚至有些人選擇離開.
原因就是離開最容易解脫.
但其實沒有解除你內心的鬱結.
你去到另一個教會.
起碼都好.
但都解除不了.
因為你根本釋放不了自己.
你可以想像得到.
這些結只會越扭就越實.
令到我們有時失眠.
令到我們不開心.
令到我們受傷害.
心靈的傷害.
其實很難去醫治.
很難見得到.
對人歡笑.
背人愁.
怎樣醫治呢.
內心燙血的都沒人知.
你受傷了.
包著傷口.
戒了還會好.
大家遷就著.
但心靈的傷口就無法知道.
和無法醫治.
有些人為何無端端不再回祈禱會.
不再崇拜.
原因就是心靈一定受了很大的鬱結.
帶來的鬱悶.
所以第六章是一個很嚴肅的.
神的聖聖.
用今日就是神的家.

$^{1481}$我們怎樣能夠捉好圍牆.
大家彼此守望.
大家有一些謠言.
就要溝通.
就要大家能夠在神的愛裡.
坦誠.
甚至知道之後.
都尋求對方的歡恕.
或者彼此的接納.
包容.
主怎樣去接納我們.
我們都要彼此去接納.
這樣的一個學習.
教會的生活才是真正的這種味.
第七章 城牆修完.
我安了門扇.
守門的歌唱的和利尾人都已經派定了.
開始發現.
數一數人口就發現不夠人.
你就可以很快想像到.
你修好四圍的城牆.
你就會發現原來這麼少猶太人.
尤其是事奉的人.
就立即要去提升猶太居住在耶路撒冷的比例.
第二節 我就派我的弟兄哈拿尼.
和仍留的宰官哈拿尼雅.
管理耶路撒冷.
留意一下這裡.
因為哈拿尼雅是忠信又敬畏神.
過於重仁.
不是跟班.
不是拍馬屁的認性蟲.
更加不是一個阿諛奉承的.
拍尼希米巴結他.
奉承他的下屬.
他是一個敬畏神的猶太人.
你可以想像到尼希米有他這兩個手下.
把耶路撒冷的管理權交託給他們兩個.
重點是他們忠信.
非暴.

$^{1521}$忠心.
信實.
忠心就是他知道自己做什麼.
他就要不是吞噩的.
假傳聖旨.
做得不夠徹底.
忠信就是他領受了尼希米的吩咐.
他就忠心的去辦事.
不會有異議.
不會有自己的偏心.
不會有自己的野心這樣去做.
然後就是敬畏神.
最重要的.
過於重人.
可能他也會有一些人對他.
連尼希米也被人批評.
被人攻擊.
這些手下.
哈拿尼和哈拿尼亞.
都一樣受到攻擊的時候.
敬畏神過於重人.
我怕神.
不是不怕人.
而是我的心.
你又要討好.
他又討好.
我哪有時間討好你們每一個人呢.
每個人都說不同的意見.
不如我簡單一點.
我是敬畏神.
我是忠心辦事.
所謂忠心和敬畏是一個平衡.
不是只有天道就沒有人道.
他是一個很忠信的領袖.
對尼希米的吩咐.
是很忠實地去跟進的.
和執行的.
但更重要的大前提.
不是擁戴巴結尼希米.
尼希米叫我做.

$^{1561}$我就什麼都做.
別人叫我做.
我什麼都不做.
不是這樣.
而是對神的敬畏.
這是很重要的.
有時候大家都是.
雖然網上聽這些錄音.
都會有很多外人聽到.
但我覺得不要緊的.
張牧師臨走之前.
大概2016,2017.
他知道自己很清楚.
自己要離開錦宣.
唐統輝也開始嘗試.
試用我.
2017年的3月.
張牧師就安息年.
正式由2015年的3月.
做代唐主任.
有半年張牧師不在.
他要讀博士論文和宣教.
他經常跟其他領袖說.
要支持我.
支持張牧師.
好像支持我一樣.
是為教會為神.
不是為誰人.
這是很重要的.
胸襟和領袖是要跟隨的.
是領袖.
不是因為我喜歡張牧師.
就跟隨他.
是他很坦誠地跟傳道同工.
跟領袖說.
我覺得是很大的胸襟.
去說這些話.
你不要因為他就不支持.
你支持我的原因.
過去一直都是對神的愛和敬畏.

$^{1601}$所以將來要轉.
你也要支持.
這是他的事工.
是神給他一個新的開展的時候.
你也要這樣支持.
我覺得領袖的胸襟.
是我自己很大的見識.
很敬佩張牧師.
他這些最後一兩年.
不是最後一兩年.
他對其他領袖的勸勉.
我覺得是很重要的.
不應該自己留尾巴之餘.
更加你們要看準的是.
對事工對神的家的參與服事.
你不是我的粉絲.
我不需要粉絲.
我需要的是大家一起去繼續服事.
這是一個很重要的角度教會觀.
這些是智慧行間.
你會看到內心價值觀的情操.
不是說你不要支持他.
你繼續支持我.
不應該是這樣.
這樣想也是錯的.
一定要很坦誠地.
我們是為神的工作.
繼續熱心地去事奉.
誰人的轉變.
都不應該有任何的影響.
沒錯.
感情上19年的牧者和主任.
一定感情是比剛加入對公的牧者深.
那是沒有問題的.
感情上是完全沒有問題.
但是在教會裡面的事奉.
就不應該是我們勒住肩膀.
我們不要做了.
就不應該這樣看.
所以這些其實是很重要的價值去學習.

$^{1641}$我們每個人.
其實職場是一樣的.
我們真的要對事不對人.
仍然是這麼樂意.
這麼全心地為神去擺上精力去服事.
這裡告訴我們.
哈拿尼亞和哈拿尼是忠信的人.
所以不要只看敵對勢力.
應該是培植自己認同你的價值觀.
你又可以委以重任.
讓他有一個很強的授權.
記住不是你的班底.
不是你的操縱者.
是你要授權給忠信的人.
提摩太後書二章第一節.
你要交託忠心能夠傳道的人.
能夠教導別人的人.
保羅教提摩太.
你要交託給忠心能夠教導別人的人.
即是那些的屬靈領袖.
你要給他們訓練.
給他們有信心.
給他們有授權.
他們慢慢來建立成長的領袖成熟.
這些都是很需要一代一代的傳承訓練.
所以我們見到有兩個忠信的人.
尼希米有時會回去述職.
遲些我們會見到.
他會回去波斯帝國亞達赤西王的王宮述職.
過了一段時間才回來.
所以有這兩位領袖在耶路撒冷把守住.
所以那個把守是慢慢恢復摩西五經.
摩西的律法書的標準.
敬拜的規範.
規模.
合神心意的敬拜.
回到這個問題.
尼希米想告訴我們.
聖文的重建.
先是有成長硬件的.

$^{1681}$接著就是要排除外邦人的影響.
接著就是要在人心裡恢復合神心意的敬拜.
這樣才能孕育四百多年後的彌賽亞降生.
所以是一個一個的時代.
但需要有這個基礎.
需要有尼希米的這個意象.
去看到這個問題.
接著可以講第三節.
我吩咐他們說.
等到太陽上升才可開耶路撒冷的城門.
人上看守的時候就要關門上山.
也當派耶路撒冷的居民.
各按班次看守自己的房屋對面之處.
一方面是守衛.
另一方面就是開始有一個.
真的要分別為勝的界線.
外邦人不可以隨便的自出自入.
甚至他們還有一種敵對的.
更加不容許他們這樣去示意的.
進來破壞或者影響.
這是平日.
到了安斯拉就更加需要謹慎.
第四節.
城是廣大,但其中的民卻稀少.
房屋都無人建造.
我的神感動我心.
這些是很重要的.
慢慢沉澱下來.
有些事是會想到的.
所以又是拆開說.
我們這個茶經班都是多成年人.
大家來到敬拜神.
不要習以為常.
真的每一次都要看自己是新鮮的敬拜者.
早點回來安靜.
當然可以和別人聊天,沒問題的.
但是心靈的安靜等候神.
預備好自己的心靈.
會有感動.
會有突然間聖靈的啟示給你的.

$^{1721}$這裡我都說了.
猶太人說.
聖殿的WIFI是最強的.
5G 6G直上天庭.
猶太人的信念.
我們新約也是.
一個真正敬拜的地方.
就是心靈誠實.
為什麼我們要強調禮拜堂要安靜呢.
敬拜的時候要安靜.
就是沉澱我們的心靈.
讓我們的心靈倒空.
沒有雜念.
預備好自己.
神會和你說話.
可能透過指鈴領唱.
和你說話.
可能直接和你說話.
可能直接真的.
甚至會讓你看見異象.
我們不要抗拒這些.
但我們不是追求靈恩那些.
而是我們的心靈準備.
是重要的.
很多時候.
福音教會慢慢都覺得自己.
不會有經歷神蹟的機會.
都覺得自己是例行公事.
慢慢真的發現自己是例行公事.
又沒有出路了.
不回教會又會被人說.
慢慢真的.
集議為常.
坐那個位.
幾點鐘到.
10點58分就來到.
有些人等我準備進電的時候.
就快腳衝進來.
基本上每個星期都幾乎是這樣.
這樣才有趣.

$^{1761}$從一個說笑的角度.
但其實大家有沒有想過.
聽到的.
茶經班弟兄姊妹.
可能有幾百個like的點擊率.
如果有四百多人聽到.
你有沒有想過.
不是因為曾牧師說.
早15分鐘大家要連唱.
而是你自己真的有一個意識.
人最怕就是你有意識地做一件事.
你有意識地.
我今天真的好想等候神.
和我講話.
我準備我自己.
今天我特意找一個.
我平時不習慣坐的位.
我坐到第三四行.
是沒有人騷擾到我的.
我真的準備自己.
尼希米就是說.
我的神突然感動我的心.
招聚那些貴就官長和百姓.
要照家譜計算.
我找著第一次上來.
即是以色列記第一章的家譜.
是古典的典籍.
他計算那些家譜.
然後就開始數點.
我們覺得很悶的事.
尼希米就看那個家譜.
他就是說.
比對一下現在我們建造好聖城的城牆.
現在有現在公職的家族.
哪個家族哪個家族.
現在有哪些人原來是渾水摸魚.
不是祭司的家族.
又在做呢.
那就要可能他們不適合.
他們原來不合規格.

$^{1801}$那就要將他們剔除.
你看看第61節.
很快讀一讀這裡.
從特米拉.
特哈薩.
基祿.
阿頓.
音墨上來的.
他不能夠指明他們的宗族譜系.
是不是以色列人.
他們就是地來瓦的子孫.
多比亞的子孫.
尼哥大的子孫.
有642名.
接著第63節.
祭司當中.
哈巴瓦的子孫.
哈哥斯的子孫.
巴西來的子孫.
因為他們的先祖.
娶了基列人巴西來的女兒為妻.
所以起名就叫巴西來.
這三家的人.
在族譜之中.
尋查自己的譜系.
都找不到.
因此.
算為不潔.
不可以再供祭司的職任.
當然說笑.
突然間發現.
誰傳了人.
原來他的神學畢業.
證書是假的.
他一直都是用假證書.
當然說笑.
沒有這些事.
那怎麼辦呢.
大概這樣.
尼熹米.

$^{1841}$他說.
我的神.
感動我的心.
其實你一眼就看到了.
當你建好之後.
非猶太人.
不是買賣的日子.
拜託不可以進來.
也就是說耶路撒冷.
人丁很稀少.
說明稀少.
你不要那麼執著.
隨便就算了.
都不夠人事奉.
尼熹米說不行.
一樣說一樣.
不夠人是另一個問題.
你不合資格事奉就是問題.
大家分明到這個問題嗎.
一眼就看到.
是人丁稀少.
第七章已經說了.
城是廣大.
其中的民稀少.
房屋沒有人建造.
接著.
感動他的心.
人口普查.
查第一次上來的族譜.
查完之後.
比對現在.
參與事奉的.
就發現有些人.
不可以展示他.
那個戶籍.
那個支派.
那個身份.
但是問題就是.
他們參與那個憲制的事奉.
尼熹米就說.

$^{1881}$不准公濟師的職任.
因為是算為不潔.
當然.
之後就會行一些手續.
什麼手續.
留意.
第65節.
省長就對他們說.
省長就是尼熹米.
不可吃自盛的物.
什麼叫自盛的物.
平安祭.
贖罪祭.
如果不是全年的贖罪祭.
濟師都可以吃那些祭肉.
平民百姓就不行.
平安祭.
就是神人共享祭.
就是你憲制完之後.
只是將那些精華.
那些腎.
肝.
少少的脂油.
燒在壇上.
整隻.
是可以.
斬件.
濟師可以吃.
那些憲制的平民百姓都可以吃.
所以以尼加拿上來憲制.
他兩個太太都可以吃祭肉.
分給哈拿.
就是雙份.
記不記得.
我們說過哈拿書.
說過《十武二記》上.
哈拿分雙份的祭物.
但是都不開心.
都經常哭.
因為.

$^{1921}$他的對頭比利娜.
經常激動他.
激他.
就是那個.
姨奶激哈拿.
激到他不開心.
留意一下這裡.
尼希米再尋求摩西的律法書規定的條例.
雖然我們不知道什麼叫污淩土明.
但是.
去到尼希米的時代.
他要解決這些疑難.
當然不是驗DNA.
驗血.
你的血統.
而是說.
要等到有可以用污淩土明.
去決疑.
解決那些疑問的濟師興起.
這個是舊約摩西規定的.
將來你有疑難.
那個濟師就用污淩土明.
好像不是抽籤.
是不知道是什麼.
可能放在那個.
濟師有一個盒子.
在前面.
裝十月一寶石.
好像是將他放在裡面.
好像是一個.
不知道是不是紮飾仔.
我們不知道.
因為已經失傳了.
但是.
到尼希米的時候.
他都仍然訴諸於濟師.
就是說.
尊重濟師有這個職份.
濟師來解決.
濟師來安排.

$^{1961}$如果濟師查過之後.
認為這個家族.
真是對的.
雖然那個家譜已經遺失了.
但是都可以讓他們恢復.
那個在聖殿裡面的功績.
就是說.
對事奉是嚴格的要求.
寧願沒有人.
正在做的.
巴西來的子孫正在做.
濟師的職任.
暫時不要做.
不是對你有偏見.
很容易就是.
尼希米會被人說.
正在做你不讓我做.
有什麼事.
尼希米就說.
不是.
是.
人都不夠了還不讓我做.
都不是.
而是.
摩西的律法書規定的.
怎樣的事奉的資格.
要濟師.
要利未的子孫.
現在你沒有那個家譜去submit出來.
證明到你是.
暫時你停止這個功績.
直到有濟師可以解決這個疑難.
他們一定有他的辦法.
就是說.
一個客觀的印證.
原來沒問題了.
原來你們一直都在做.
原來你們的家譜.
真的可以追溯的.
這樣.

$^{2001}$尼希米不是.
在此來到此地無銀三百兩.
更加不是.
好像來到去獨攬的大權.
他說了事.
不是.
而是.
我們回到摩西的律法書.
神的指引.
所以重建聖文.
是要在方方面面.
他們回到神的話語那裡.
在這些.
其實挺矛盾的.
不夠人.
你不要那麼挑剔.
你不要那麼要求.
未洗禮.
他做堂董有什麼問題呢.
他很厲害的.
他是EMBA來的.
他是管理學.
很出類拔萃的.
為什麼不可以做.
他是不喜歡洗禮.
他是不想加入我們教會.
有什麼問題呢.
他很熱心的.
現在沒有人了.
很多的理由.
我們是不是要將那條線.
一直有那麼低.
弄那麼低呢.
甚至他不信的.
我們都不要理他了.
他每個星期來教會.
信不信由他.
他信不信到天堂.
將來是他的事.
不關我們的事.

$^{2041}$是不是我們要將那條線.
一直降低到越來越低呢.
還是原來前人告訴我們.
洗禮某程度是一個.
很認真的信念來的.
為什麼他不做呢.
有什麼難言之隱呢.
有難言之隱就要堅持.
做一些領袖崗位呢.
其實很矛盾的.
你既然承擔這個領袖崗位.
你應該以身作則.
你就應該更加.
你不叫我做我也做.
這樣才對的.
為什麼我們想事情.
有時候會反過來呢.
越來越遷就那些人.
降低那些標準呢.
應該就是.
我們更加要讓他看到.
黎希米就是這樣.
直到有能夠決議的祭司.
可以運用污淩土明.
來解決這個疑難.
就可以公職了.
暫時因為真的看不到你的家譜.
所以就要停止你的公職了.
所以第七章.
你可以想像得到.
就是黎希米.
很特別的.
你看到建好了.
外邦人不可以再隨便.
進入之後呢.
城很大.
人很少.
但是他突然間有感動.
少還少.
我要現在看回.

$^{2081}$在聖殿裡面.
公職的祭司.
利未人.
到底有多少.
有多大的欠缺.
還有下一步就是.
首先第一步就是.
哪些有疑難的暫時停下來.
然後下一步呢.
就是遷徙了.
就是鼓勵更多.
在周邊地方住的.
遷入耶路撒冷.
更多在耶路撒冷裡面.
去參與那個.
重建的工作.
你就可以想像.
一直去.
好像做旺了.
整個耶路撒冷.
有足夠的人.
一齊.
更多呢.
就是在第八章.
正正就是七月初一的節期開始.
所以你看到.
以六月的六月二十五日.
接著有大概七天的時間.
就有一個這樣的準備.
接著.
大家都準備好了.
就開始嚴肅會了.
嚴肅會就帶來那個靈性的復興了.
所以你會看到.
原來有些東西都是一個進程的.
先是將外邦.
外邦人.
其實他們是.
以色列地的原居民了.
因為.

$^{2121}$所謂原居民就是猶太人被擄了.
沒有人了.
只剩下那些老弱殘兵.
現在被擄歸回了.
都大概百多年了.
一百二十五年左右.
現在來到黎天敏的時代了.
其實外邦人已經真的成為.
當地的原住民了.
但是.
他們不是猶太人.
也沒有想過進入猶太教.
對不起了.
有界線.
界線就是那個圍牆.
你是不是我們的族群了.
你要離開了.
今天我們用什麼.
作為我們的界線呢.
就是信不信.
這是一個基本的界線.
天堂和地獄是界線.
所以我們真的要看到.
我們的兒女.
我們的兒女的另一半配偶.
我們的孫子.
我們的下一代.
下兩代.
那個界線就是.
信主在耶穌的救恩裡面.
和不在耶穌的救恩裡面.
那個天淵之別.
所以你看到黎天敏就.
當然他們是因為舊約.
是一個約的子民.
所以那個界線就相對簡單.
就是你是一天猶太教.
你可以進入猶太教.
你敬畏神.
但是如果你都不做.

$^{2161}$你就是我們的所謂的.
非猶太人的外邦人.
外邦人這個名字.
不是在黎起米的時候那麼普遍.
慢慢慢慢就普遍了.
再過一段時間.
猶太人外邦人這個界線就很強烈了.
所以現在開始分猶太人外邦人.
當然後期去到耶穌的年代.
又變成過度的分開了.
看撒瑪利亞人又是外邦人.
希臘人其他非猶太人又是外邦人.
就變成了一種歧視排擠.
就不是黎起米這個時候的需要.
那個問題的精結.
所以現在就是要將聖文分別.
先用城牆去分別.
第八章開始.
就是開始在神的話語上的分別為聖.
對不起我漏了一件事.
下個星期四晚.
因為有一個聖樂出盡會.
他們有一個晚宴.
我就被邀請去.
所以下個星期四晚要停一次的查經班.
我們一起禱告.
多謝主的恩典.
我們能夠一起認真思考.
第四第五第六章.
我們看到在第六章的時候.
真的有很多事情.
我們需要一起去把關.
心靈裡面我們要對準神的救恩.
公義和信實.
我們要敬畏神.
在神的家裡面我們做一個建設者.
一個和平之子.
來讓更多人在神的愛裡面去經歷.
和得到神自己的照管.
也讓我們能夠在一些我們發現.

$^{2201}$可能真的只有我們發現的破口.
我們去堵塞.
可能是用愛.
可能是用眼淚.
可能是用禱告去堵塞.
或者可能是用言語去堵塞.
這些的破口.
不容讓仇敵.
在教會裡面搬弄是非.
挑撥.
甚至乎惡意的去打壓.
一些的牧者.
或者侍工等等.
我們很盼望.
我們能夠有更多忠信.
敬畏神.
過於眾人的一些領袖.
可以被興起.
在主的感動底下.
常常都有主的感動.
在聚會當中.
以致我們能夠.
真的可以成為很多人的指路的燈.
一些腳前的光等等.
我們這樣禱告.
奉耶穌基督的名.
阿們.
\newpage



\section{腓立比書 3:7-16}
\label{sec:We85QzJyMNQ}
\textbf{2023-05-20 晚堂《丟棄與得著》腓立比書 3\_7-16 講員:盧思源姑娘}
\newline
\newline
連結: \href{https://youtube.com/watch?v=We85QzJyMNQ}{\texttt{https://youtube.com/watch?v=We85QzJyMNQ}} ~~~~ 語音日期: 2023-05-23
\newline
\newline
\hyperref[sec:MtaDL5MWzm0]{\small{< < < PREV SERMON < < <}}
~
\hyperref[sec:index_chronic]{\small{[返順時目]}}
~
\hyperref[sec:index_scriptual]{\small{[返順卷目]}}
~
\hyperref[sec:9tgCzMr4KBk]{\small{> > > NEXT SERMON > > >}}
\newline
\newline
腓立比書 3:7-16
\newline
\begin{longtable}{cl}
\hline
\hline
章節 & 經文 (和合本修訂版)\\
\hline
3:7 & \begin{tabularx}{0.7\textwidth}{X} 只是我先前以為對我是有益的，我現在因基督的緣故而當作是有損的。 \end{tabularx} \\ \\ \relax
3:8 & \begin{tabularx}{0.7\textwidth}{X} 不但如此，我已把萬事當作是有損的，因我以認識我主基督耶穌為至寶。我為他已經丟棄萬事，看作糞土，為要贏得基督， \end{tabularx} \\ \\ \relax
3:9 & \begin{tabularx}{0.7\textwidth}{X} 並且得以在他裡面，不是有自己因律法而得的義，而是有信基督的義，就是基於信，從神而來的義， \end{tabularx} \\ \\ \relax
3:10 & \begin{tabularx}{0.7\textwidth}{X} 使我認識基督，知道他復活的大能，並且知道和他一同受苦，效法他的死， \end{tabularx} \\ \\ \relax
3:11 & \begin{tabularx}{0.7\textwidth}{X} 或許我也得以從死人中復活。 \end{tabularx} \\ \\ \relax
3:12 & \begin{tabularx}{0.7\textwidth}{X} 這不是說我已經得著了，已經完全了；而是竭力追求，或許可以得著基督耶穌所要我得著的。 \end{tabularx} \\ \\ \relax
3:13 & \begin{tabularx}{0.7\textwidth}{X} 弟兄們，我不是以為自己已經得著了；我只有一件事，就是忘記背後，努力面前的， \end{tabularx} \\ \\ \relax
3:14 & \begin{tabularx}{0.7\textwidth}{X} 向著標竿直跑，要得神在基督耶穌裡從上面召我來得的獎賞。 \end{tabularx} \\ \\ \relax
3:15 & \begin{tabularx}{0.7\textwidth}{X} 所以，我們中間凡是成熟的人，總要存這樣的心；若在甚麼事上存別樣的心，神也會把這些事指示你們。 \end{tabularx} \\ \\ \relax
3:16 & \begin{tabularx}{0.7\textwidth}{X} 然而，我們達到甚麼地步，就當照這個地步行。 \end{tabularx} \\ \\
[1ex]
\hline
\hline
\end{longtable}
$^{1}$大兄姐妹平安.
願意我們開始之前.
我們都有一個禱告.
出去之後我們再次去到你的面前.
讓我們可以回到教會當中.
一同去敬拜.
一同去聽你的話語.
一息奸孩子要分享你自己話語的時候.
求聖靈與我們眾人同在.
打開我們的心.
願意我們一同的.
無論講的 聽的.
我們都進入你自己的真理當中.
被你自己的真理得著幫助.
得著勸勉 得著造就.
我們恭敬的將以下時間交上.
奉主耶穌基督明智祈求.
阿們.
大兄姐妹.
斷捨離三個字.
我相信大家都有聽過.
原來這個詞語是由一個日本人.
叫做山下英子去提出的.
他的意思就是要斷絕一些不需要的東西.
捨去一些多餘的東西.
也都脫離對一些物品的執著.
但當中不是只是說.
你丟掉了一些東西就算了.
而是透過斷捨離的方式.
由整理一些物件.
捨棄一些沒用的東西.
從而令到人的內心的心靈.
可以專注在一些重要的事上.
所以當中就好像.
你要去捨棄一些事情.
但另一個層面.
你是得著的一些東西.
而剛才我們讀的肥勒比書第八節.
保羅說.
我為他已經凋棄萬事.

$^{41}$看作糞土.
為要得著基督.
很多時候我們都會很容易很直接.
去解這節的經文.
就在想.
是不是叫我們作為基督徒.
我們就什麼都不要了.
就算有很多東西都好.
都要捨棄.
都要當作糞土.
最好就不要那麼多東西.
有的都丟掉了.
這樣就可以得著基督.
這樣就好.
是不是這樣解釋呢.
但保羅其實在這裡不是這個意思.
如果我們翻開肥勒比書.
大家可以打開聖經.
剛才沒有讀到.
肥勒比書第三章開頭的時候.
其實保羅去勸勉當時肥勒比的信徒.
要小心原來當時有一些假師傅.
在三章二節那裡這樣說.
應當防備犬類.
防備作惡的.
防備妄自行覺的.
保羅稱當時有一些猶太主義的信徒為犬類.
就是狗.
為什麼他們會是犬類呢.
原來當時因為那些信徒.
他就靠自己行國禮.
去證明他自己得救.
以一種外在宗教的行為.
來證明自己比其他人強.
比其他人健健.
他們以為這些所做的.
用這些行為來證明他們自己.
但是保羅告訴他們.
真正蒙神越立的是什麼呢.
所以在第三節那裡說.

$^{81}$因為真受國禮的.
乃是我們以神的靈敬拜.
在耶穌基督裡面誇口.
不靠肉體的.
原來真正受國禮的.
是在內心裡面接受.
並且以神的靈去敬拜.
而不是按照人所定的規條.
接著一段很熟悉.
保羅在第四節之後就說.
其實保羅他也可以靠肉體.
如果他其他人想靠肉體.
保羅說他自己更加可以.
因為相對於猶太教徒.
保羅是更有資格去說這句話.
他誇口的.
大家看聖經的話.
第五節第六節有七件事去誇口.
頭四件事是關於他自己的出身.
他說第一.
他是作為亞伯拉罕的子孫.
按照第八天的標準已經受國禮.
第二他是一個以色列人.
他不是歸化入籍的.
第三他是屬於變顏問太.
當時是一個名門望族.
也是有家譜的.
第四他說他是希伯來人.
山的希伯來人.
代表保羅的血統是很純正的.
之後的三件事就記載保羅自己的努力成就.
就律法來說他是法理賽人.
大家知道法理賽人是非常專注和熟悉律法的.
而且是當時頂尖的教派.
而且保羅當時是一個最出名的門下.
就是在加瑪列的門下.
是一個高材生.
所以當時他比同年紀的人更加長進.
更加有成就.
之後他提到.

$^{121}$他就熱心迫不及待教會.
因為信仰正統的猶太人.
當時認為耶穌派是異端.
所以保羅大家知道.
他曾經是很竭力去消滅.
迫不及待教會.
最後保羅就說在律法裡面說.
他是無可指責的.
就是說保羅說.
沒有人可以在律法的義行上.
可以找到保羅任何的錯處.
所以保羅說如果要誇口.
他更加可以誇口.
去到今天我們剛才所讀的經文第七節.
那裡有一個轉接.
他說只是.
但是他先前以為對他有益的.
我現在因基督的緣故而當作有損.
原來從前保羅他以為.
剛才說的那七樣他所誇口.
很優越的東西.
他現在認為是一個虧損.
原因是因為基督.
使他在觀念和看法上有所改變.
但是如果我們看清楚一點.
在這裡說有益和有損的東西.
他不是在說兩樣東西.
是在說同一樣東西.
剛才那四節六節.
那些保羅曾經很引以為榮.
或者很努力追求.
甚至乎他很自誇的東西.
然而到了現在.
對同樣的東西.
這些東西他反而當作有損.
所以在這裡他不是將兩樣東西對比.
不是有價值的東西和沒有價值的東西對比.
而是同樣的東西.
有了完全不同的看法.
完全是因為他個人的價值觀去轉變.

$^{161}$大家可能都曾經有過這樣的經歷.
大家有沒有試過曾經很喜歡一件東西.
可能是一件衣服.
一個袋子.
一件物件.
或者小朋友可能很喜歡一個公仔.
甚至乎大一點可能很喜歡一個人.
曾幾何時.
可能你會覺得這件東西.
這件衣服.
這個公仔.
這個袋子.
這件物件.
這個人.
是很重要.
沒有了是不行的.
不過可能時間久了.
因著一些原因.
同一件衣服.
同一個袋子.
同一件物件.
同一個公仔.
甚至乎是同一個人.
你對他的感覺不同了.
但不是因為這件東西變了.
而是我們的思想改變了.
所以令我們對同一件事有改變的時候.
原來一個人的價值觀是很影響我們一個人的看法.
而我們每一個人的價值觀.
很可能我們都是受著我們家庭.
我們的成長.
我們的朋輩.
甚至乎我們隨著這個世界的價值觀去影響.
不過今天我們作為基督徒.
我們的價值觀不是根據剛才所說.
單單說我們的家庭.
我們成長的世界價值觀去影響.
而是當我們成為基督徒的時候.
我們的價值觀是有變的.
我們是應該按著聖經所說的價值觀去看每一件事.

$^{201}$例如我們要重新去看我們的人生.
重新看對金錢的看法.
重新看我們的成功失敗.
重新看我們的婚姻家庭朋友等等.
最近有機會和一些不同的基督徒去分享去聊天.
當中都提到一些在信仰上的掙扎.
原來在分享裡面.
大家都發現有時候我們信仰的掙扎.
未必是表面做到還是做不到的問題.
舉個很簡單的例子.
可能只是說讀經靈修.
十一奉獻等等.
未必是說我做到還是做不到.
原來還原在基本上.
原來是根深蒂固的價值觀問題.
有時候可能是因為我們看不到事情的價值本身.
以致當我們遇到一些事情的時候.
或者做一些決定的時候.
有時候可能和聖經所教和我們所做的有一段距離.
回到聖經.
保羅說.
「只是我先前以為對我有益的.
我現在因基督的緣故而當作有損」.
我們在這裡看到.
保羅因為遇見耶穌基督的時候.
他整個人的價值觀有一個很大的改變.
而且在第八節裡這樣說.
「不但如此,我已把萬事當作是有損的.
因我已認識我主基督耶穌為至寶.
我為他已經丟棄萬事,看作糞土,為要贏得基督」.
保羅不單止將以前覺得有益的東西.
現在當作有損.
甚至他說連萬事都當作有損.
一開始我們提到.
有人很容易誤解這句.
基督徒是否要將所有東西萬事都丟棄.
我們才可以得著基督呢.
事實上不是這樣去解釋.
在聖經的原文裡.
丟棄和有損兩個字其實是同一隻字.

$^{241}$不過它們是動詞和名詞的分別.
丟棄是動詞,有損是一個名詞.
如果直接按字面上去解釋.
其實這句話應該這樣說.
保羅說「我為主已經將所有的事當作無益處,看作垃圾.
避免因為受任何事的影響而失去基督」.
這是保羅那天的觀念.
原來保羅任何會影響他和主的關係的東西.
他都可以丟棄.
即是說如果金錢是影響他和主的關係的話.
保羅說他可以放棄金錢.
如果感情是影響他和主的關係的話.
他可以放棄感情.
如果名譽地位等等影響他和主的關係的話.
保羅說他都可以放棄.
保羅不是說為了主的緣故.
他就應該要放棄金錢,感情,名譽,地位等等所有的東西.
大家聽不聽到分別在哪裡.
保羅他不是一個苦行僧.
不是一個苦修的主義者.
他都是有人性的需要.
我們知道在聖經裡.
他曾經因為身體上一根刺帶給他痛苦.
他三次求神可以磨開它.
保羅在羅馬的時候.
可能因為天氣冷.
他寫信給提摩太.
叫他拿件衣服給他.
我們看到保羅他不是主張要受苦才懂得愛主.
有很多其他的宗教.
可能主張人要抑壓你的慾望.
比如佛教.
你要四大皆空.
你要六根清淨.
你要無慾無求.
這樣才是好的.
但是在我們基督教的信仰裡.
我們不是不可以有這些東西.
在聖經保羅這裡說.
如果有什麼東西會影響我們的信仰.

$^{281}$阻礙我們和耶穌基督建立關係.
叫我們減少對耶穌基督的追求的話.
這些事我們就應該要放下.
今天我們可以想一想.
有沒有什麼人事物.
是阻礙我們和神建立的關係.
我自己在想的時候.
很普遍也很普通.
我想手機有時候也是其中一樣.
當然我們知道現在手機很好.
給我們很多資訊.
資料.
很多東西我們都可以即時找到.
但是有時候手機也挺消耗我們的時間.
不知道大家有沒有覺得.
或者自己有沒有試過.
滾動手機.
或者年輕人打機的時候.
時間很快就過了.
我自己也是這樣.
半小時.
一個小時.
好像時間慢慢就已經流走了.
但相反.
反過來.
叫你打開聖經.
我們去讀經靈修.
可能叫你坐定定看的時候.
可能15分鐘.
你已經覺得.
15分鐘而已.
甚至可能10分鐘就已經睡著了.
又或者.
當你真的想讀經靈修的時候.
手機突然間有個訊息.
可能我們都已經被打差了.
打斷了.
今天.
聖經不是說我們要放棄我們所有的東西.
而是如果當有些東西.

$^{321}$阻礙我們和神的關係的時候.
叫我們減少了和神的追求的時候.
我們就要學習放下.
學習丟棄.
在聖經裡面.
保羅為何能夠丟棄這些東西呢.
因為保羅說.
我以認識我主耶穌基督為至寶.
這句話說出了為何保羅有這麼大的轉變.
因為他說他以認識耶穌基督為至寶.
所以保羅將世界和他所認識的耶穌一比之下.
他發覺世界的價值就像糞土.
就像垃圾.
保羅因為要得著耶穌基督的時候.
他甘願地去丟棄世界很多的事情.
或者我們可能會問.
保羅得救的時候.
不就已經得著耶穌基督了嗎.
不就已經在神的裡面了嗎.
為何在這裡還要得著基督呢.
沒錯.
當保羅信主.
或者在我們信主的每一刻.
我們已經得著基督.
但多多少少.
很多時候都是著重在基督上的恩典.
我們信主的時候.
我們得著神的赦免.
得著他恩順稱義.
我們得著子女的名分.
但在這裡.
保羅所渴求的是得著耶穌基督自己.
單單得著基督的恩典.
不可以滿足保羅的心.
保羅渴望的是主耶穌基督自己.
就像宣導會的創辦人.
孫迅博士.
他有一首詩歌.
大家有沒有聽過.
叫做《主自己》.

$^{361}$當中有一段歌詞是這樣的.
前所要是福祉.
前就是以前.
今就是現在.
今所要是主.
前所要是感覺.
今要主話語.
前所要是恩賜.
今要賜恩主.
前我尋求醫治.
今要主自己.
這份心志.
也是我們今天每一位基督徒.
我們應該去追求的.
我們不是單單停留在.
我們信主得救的階段.
而是在我們信仰的路上.
我們更深的去認識.
我們的主耶穌.
我們要更深的去和祂建立關係.
要去得著耶穌基督自己.
回到經文第十節我們看看.
保羅繼續說.
使我認識基督.
認識基督就是知道祂復活的大能.
並且知道和祂一同受苦.
效法祂的死.
在這裡保羅所說的認識.
並不是單單說頭腦上的認識.
而是在生命上體驗到.
基督復活的大能.
對保羅來說.
認識基督復活的大能的經驗.
不是說行神蹟.
說方言.
或者講預言.
見異象等等.
而是效法耶穌基督.
獨人的形象.
悲慘受苦甚至受死.

$^{401}$因為當他受苦的時候.
仍然可以坦然面對的話.
如果不是因為有耶穌基督復活的大能.
在他的身上.
其實是做不到的.
接著在十二節.
保羅說.
他這裡不是說自己已經得著.
雖然我們知道.
保羅寫這封信的時候.
他在監獄.
他為主坐牢.
而且甚至剛才看到.
保羅已經可以為主.
凋棄萬事.
但是保羅仍然說.
他不是說自己已經得著.
已經完全.
反而他是什麼.
他說而是揭裂的追求.
當中意味著.
保羅不只是單單告訴我們信徒.
信仰是一個歷程這麼簡單.
而是保羅他自己.
仍然在這個歷程當中.
就是因為尚未完全得著的時候.
保羅說他只有一件事.
而這件事對於他來說.
是優先的.
是必須的.
是重要的.
保羅引用當時的田徑競賽.
他就說.
描述他這個追求.
就是要忘記背後.
努力.
面前.
向著標竿直跑.
就好像一個選手.
在一個競技賽中.

$^{441}$他不可以瞻前顧後.
他一定要專注在賽程中.
專注在他的標竿中.
這是一個終點的標誌.
而一個選手為何可以這麼專注呢.
就是因為他跑.
是為了希望可以得到勝利.
可以贏得獎牌.
在以前就是一個.
一個用業績成的冠冕.
在當中保羅說.
他的追求是有目標的.
他是希望他可以得到基督.
而耶穌基督也有獎賞給他.
就是要得神在基督耶穌.
上面照他來得到的獎賞.
而這個獎賞就是今世完成.
上帝給他的託付.
而在永生中有一個永恆不朽的冠冕.
而且我們再看.
在第十五節.
保羅不只是說自己去奔跑.
不是他自己說.
他在勉勵我們一群信徒.
和他一起奔跑.
所以我們中間凡是成熟的人.
總要存這樣的心.
這裡的成熟不是在說年齡上的成熟.
不是在說年紀上的成熟.
而是在我們屬靈生命上成熟的基督徒.
保羅相信一個靈命成熟的基督徒.
一定會像他這樣.
察覺到信仰是一場永無止境的旅程.
而同樣是一場永無止境的旅程.
是會有這個心.
就是全這樣的心.
就是一個不斷追求.
不斷前進的心.
就算真的尚未有這種心智的信徒.
保羅也相信.

$^{481}$在什麼事上存別的心.
神也會把這些事指示你們.
當中聖靈會親自教導啟示.
他一群信徒.
開始追求得著基督的旅程.
弟兄姊妹.
今天我們在信仰的路上.
我們是什麼情況呢.
或許今天我們每個人信主的年日都不同.
各有長短.
但今天我們在我們的生命裡.
我們追求著什麼.
是我們的工作.
是我們的事業.
是我們的家庭.
是感情.
正如剛才我們也說過.
聖經不是說這些事有問題.
而是若然我們所追求的事.
影響了我們與神的關係.
我們如何選擇呢.
在經濟學裡有一個詞.
叫做機會成本.
可能大家也有聽過.
叫做opportunity cost.
這個詞在經濟行為裡是一個很關鍵的觀念.
簡單來說就是將人力物力.
用在某一件事上.
例如你去看電影.
你去看電影就會放棄其他機會.
那些機會可能是你去看書.
逛街 睡覺.
那些機會就是你去看電影的機會成本.
到底看書還是睡覺是看電影的機會成本呢.
答案是哪一樣較大價值呢.
這個就是機會成本.
例如這個時間看書比睡覺有價值.
如果你選擇了去看電影.
那你的機會成本.
看電影的機會成本就是看書.

$^{521}$所以機會成本簡單來說.
就是其他最佳選擇的價值.
就是the best alternative.
我們很多時候正所謂.
大家都會聽有得就有失.
事實上失就是得的機會成本.
大家聽得明白嗎.
所以機會成本提升了.
我們不只是記得得到的東西.
而忘記了失去的東西.
舉個例子.
好像有些大學生他們走堂.
不過走堂沒有被老師點名.
心裡想著賺到了.
沒有take attendance.
但可能他的機會成本都頗大.
就是因為那堂課的內容沒有聽到.
導致之後接不下去.
那堂課的內容就越來越聽不明白.
可能最後不合格.
或者是得不償失.
回到我們的信仰.
今天我們的選擇是什麼呢.
有時候會不會在我們的選擇裡.
我們得到我們在世很想要的東西.
但無意間耶穌基督變成機會成本呢.
如果是這樣的話.
這個機會成本的價值都頗大.
因為損失了我們與神建立關係.
求主讓我們智慧懂得去選擇去分辨.
正如今天經文保羅所教導.
盼望我們在信仰的路裡.
我們以得著耶穌基督.
我們追求與神的關係為首要的事.
以致我們在我們的生活.
在我們各項的事情上.
我們知道我們應該要放下什麼.
丟棄什麼.
願主的話成為我們的幫助.
也成為我們的勉勵.

$^{561}$我們一同去禱告.
抽離我們再一次到你的面前.
昔日你藉著保羅所教導.
你教導我們.
今天我們的信仰路.
我們不是單單停留在一個地方.
而是我們真的要有追求.
我們要與耶穌基督建立那份關係.
求你幫助我們.
幫助我們眾弟兄姊妹.
我們有一份得著基督.
追求耶穌基督你這份心.
因為我們知道.
耶穌基督你是最寶貴最重要.
求你幫我們有這份看見.
有這份眼光.
以致當我們在生活行事的時候.
我們懂得怎樣去選擇.
我們選擇有價值.
在你眼中有價值的事.
就是與你建立那份關係.
求你幫助我們.
願意我們各人在你的裡面.
在你的愛裡面繼續去學習.
去成長.
以致我們活出真正基督的樣式.
求主幫助.
以上的禱告奉主耶穌基督明智祈求.
阿們.
\newpage



\section{腓立比書 3:7-16}
\label{sec:9tgCzMr4KBk}
\textbf{2023-05-21 午堂《你得著珍寶了嗎?》腓立比書 3\_7-16 講員:鄧泰康傳道}
\newline
\newline
連結: \href{https://youtube.com/watch?v=9tgCzMr4KBk}{\texttt{https://youtube.com/watch?v=9tgCzMr4KBk}} ~~~~ 語音日期: 2023-05-23
\newline
\newline
\hyperref[sec:We85QzJyMNQ]{\small{< < < PREV SERMON < < <}}
~
\hyperref[sec:index_chronic]{\small{[返順時目]}}
~
\hyperref[sec:index_scriptual]{\small{[返順卷目]}}
~
\hyperref[sec:vWqPqvo2gRk]{\small{> > > NEXT SERMON > > >}}
\newline
\newline
腓立比書 3:7-16
\newline
\begin{longtable}{cl}
\hline
\hline
章節 & 經文 (和合本修訂版)\\
\hline
3:7 & \begin{tabularx}{0.7\textwidth}{X} 只是我先前以為對我是有益的，我現在因基督的緣故而當作是有損的。 \end{tabularx} \\ \\ \relax
3:8 & \begin{tabularx}{0.7\textwidth}{X} 不但如此，我已把萬事當作是有損的，因我以認識我主基督耶穌為至寶。我為他已經丟棄萬事，看作糞土，為要贏得基督， \end{tabularx} \\ \\ \relax
3:9 & \begin{tabularx}{0.7\textwidth}{X} 並且得以在他裡面，不是有自己因律法而得的義，而是有信基督的義，就是基於信，從神而來的義， \end{tabularx} \\ \\ \relax
3:10 & \begin{tabularx}{0.7\textwidth}{X} 使我認識基督，知道他復活的大能，並且知道和他一同受苦，效法他的死， \end{tabularx} \\ \\ \relax
3:11 & \begin{tabularx}{0.7\textwidth}{X} 或許我也得以從死人中復活。 \end{tabularx} \\ \\ \relax
3:12 & \begin{tabularx}{0.7\textwidth}{X} 這不是說我已經得著了，已經完全了；而是竭力追求，或許可以得著基督耶穌所要我得著的。 \end{tabularx} \\ \\ \relax
3:13 & \begin{tabularx}{0.7\textwidth}{X} 弟兄們，我不是以為自己已經得著了；我只有一件事，就是忘記背後，努力面前的， \end{tabularx} \\ \\ \relax
3:14 & \begin{tabularx}{0.7\textwidth}{X} 向著標竿直跑，要得神在基督耶穌裡從上面召我來得的獎賞。 \end{tabularx} \\ \\ \relax
3:15 & \begin{tabularx}{0.7\textwidth}{X} 所以，我們中間凡是成熟的人，總要存這樣的心；若在甚麼事上存別樣的心，神也會把這些事指示你們。 \end{tabularx} \\ \\ \relax
3:16 & \begin{tabularx}{0.7\textwidth}{X} 然而，我們達到甚麼地步，就當照這個地步行。 \end{tabularx} \\ \\
[1ex]
\hline
\hline
\end{longtable}
$^{1}$好, 兄弟姐妹午安.
今日跟大家一起分享.
我都有一個習慣.
我每一次都想花一點時間.
去跟大家讀一個另外的版本.
目的是第一.
這些版本會令大家更加明白這段英文講什麼.
第二是我想不斷去引誘大家.
請大家用三至四餐快餐店的平價午餐價錢.
多買一本聖經.
讓你讀聖經時不要經常覺得看不明白.
而是我想講上帝給了很多恩典讓我們可以看得明白.
花一點時間專心聽我讀.
有一個版本叫新普及版.
是剛才那段英文.
從前我認為這些事肯定有價值.
但現在我因基督所成就的.
就把他們都看為毫無價值了.
是的, 相比之下.
認識我主耶穌基督才是無價之寶.
其他的事都毫無價值.
我為了祂的緣故已經拋棄了一切.
把這一切都看作是垃圾.
好讓我能得著基督.
與祂合而為一.
我不再依靠自己通過遵守律法而得來的義.
而是通過相信基督而成為義人.
因為上帝使我們跟祂恢復正確關係的方法.
取決於信心.
我希望認識基督.
並且體驗那使祂從死人中復活的莫大能力.
我希望能跟祂一同受苦.
一同赴死.
好讓自己無論怎樣都能經歷從死人中復活.
這並不是說我已經做到.
已經達至完全了.
我只是竭力要得著完全.
基督耶穌先得著了我.
就是為了讓我能得著這完全.
親我的弟兄姊妹.

$^{41}$我並不是已經達到了這個目標.
我只是專注於一件事.
忘記背後.
期待那擺在前面的.
竭力奔赴塞城的終點.
要得到那屬天的獎賞.
就是上帝藉著基督.
耶穌召喚我們去得的.
願所有靈性成熟的人.
都在這事上看法一致.
要是你們對這某一點不同意.
我相信上帝會讓你們明白.
然而我們取得了怎樣的進步.
就務必照樣進步下去.
我不知道大家聽完之後會不會更加明白多一點這段經文.
很多時候舊的版本會令我們有很多誤會.
當每一次講這段經文的時候.
都很容易會讓人很快就會.
降起之降.
我知道神是重要的.
但是地上的東西都重要.
你可不可以不賺錢.
你可不可以曲高和寡.
其實不是.
弟兄姊妹.
很多時候我們都誤會了聖經.
誤會了很多保羅的意思.
其實我會稱之為這段經文.
是叫做保羅的豪情壯語.
為什麼會這樣講呢.
剛才當他講的時候.
我已經將以為以前有價值的事.
現在看來沒有價值.
如果大家看回上面的經文就會記得.
其實當時整個的腓立比書.
仍然在討論關於究竟我需不需要行割禮.
而得到這份救恩的問題.
而保羅在上面的經文就說.
如果你要講他合乎這個資格的話.
無論在身份上.

$^{81}$無論在當時我對於信仰的認真.
我都比這些人更加有資格.
但是這些在他以前.
看為最高價值的東西.
到今天來說他會覺得是沒有價值.
為什麼呢.
今天的主題叫做.
你得著真寶了嗎.
保羅並不是無端端放下的定志妹.
在這裡的時候他講得很清楚.
是因為我已經尋找到這個無價之寶.
所以我願意去放下.
今天我們作為基督徒.
我們得著了這個救恩.
但究竟這個救恩對於大家來說.
是不是一份無價之寶.
是可以去到一個程度.
讓你去將從前你認為很重要.
不斷要去追求的東西去放下.
其實是一個比較.
很多時候就是.
當然啦 我今天上耶穌會很認真的.
我都參與很多的事奉.
我都做了很多的事.
但定志妹.
其實這個問題只可以靠著你自己去問你的內心.
我其中有一個事奉.
現在我已經可以退到很後.
我很感恩.
一位非常好的童工Debra.
盧姑娘來了.
之後她做一些清善的工作.
我都會給一些意見.
做清善工作裡面.
我想不斷提及.
今天是一個成人的崇拜.
我很想告訴大家.
做兒童工作.
做清少年工作.
其實有一半的.

$^{121}$顯出來的能力.
其實是從大家作為家長來的.
是你平時在家裡面.
你怎樣去讓他看見這個信仰.
在以前我一直去做.
特別當時我來到.
剛剛感受我去做很前線的清善工作的時候.
永遠去面對一個掙扎.
就是有些年青人.
我們在教會裡面.
跟他們一起好好過生活.
我們跟他們一起好好聊天.
慢慢地對上帝有些信心.
但往往當他回到家裡.
第二個星期他就回來跟我們說.
Aka.
為甚麼回到家裡.
我父母都是回教會.
為甚麼他們不是像你所說的那樣.
其實他們是受大家的影響.
而他們有一個.
究竟是怎樣的水準生活.
是由你們去構成的.
有一次我記得我們一班同工.
去到觀塘.
參觀有一間宣道會觀塘塘.
他們的清善工作做得很好.
他們的同工都很直言.
他說是呀.
我們這裡的弟兄姊妹.
基本上是他們在家裡面已經做了一半的清善工作.
其實本身這些青年都是帶著一個.
都幾成熟的生命來到當中.
所以他們就能夠很早在他們初中.
在他們中四中五的時候.
都已經能夠可以開始做門訓.
而你會發覺一件事.
其中一個很彰顯的地方.
就是今天都仍然出現.
如果今天你作為父母.

$^{161}$到你的孩子.
特別是剛過去考DSE.
在人生裡面一個非常之重要的考試.
在這個考試裡面.
你會怎樣跟他說呢.
星期日的時候.
你是會跟他說.
無論怎樣都好.
先回教會吧.
敬拜上帝是很重要的.
還是你會跟他說.
先讀書吧.
回小廁不用死的.
弟兄姊妹.
其實這個是反映在我們生命裡面.
這個主耶穌基督.
他是不是真的我們的至寶.
而去到一個地步.
我們願意為了他的緣故.
而放下這些世界裡面.
你認為很有價值的事.
你記住我要再強調.
我不是叫你不賺錢.
我不是叫你不過一個優質的生活.
我不是叫你過一個一定要很貧窮.
我不是在說這件事.
我是在說.
當這些事相對於主耶穌的救恩的時候.
是什麼.
所以你會發覺.
保羅在這裡的時候.
他就說因為我在裡面的時候.
我願意放下.
我很想做什麼.
因為我很想達到一個跟基督完全合一的生命.
其實我們卡住了.
什麼叫跟基督完全合一.
因為原來後面他解釋了.
因為我很想去體驗什麼叫做.
耶穌基督復活的大能.

$^{201}$在剛才的版本裡面他講了兩次.
我很想去體會耶穌基督復活的大能.
當講到這裡.
你就應該要慢慢揣摩明白主耶穌基督.
你要明白保羅他自己當時的心情.
是 沒錯.
這個世界上有很多很有價值的事.
但是有更多不是這些物質.
這些學歷.
不是這些榮譽可以換到回來.
當你病的時候.
沒錯.
你有能力 你有經濟的條件.
你可以比其他人有更優質更快的醫療.
他可以讓你盡快去看醫生.
但是弟兄姊妹.
你的心呢.
你的心裡面.
當你在疾病當中.
你是很恐懼的.
很憂心的.
很憂慮的.
很沒有盼望的.
這些物質這些東西可不可以幫到你.
當你在人生的低潮低谷的時候.
究竟昔日裡面你所追求的這些東西.
能不能夠幫到你.
而保羅在講耶穌基督復活的大能就是.
當你在這些低谷.
甚至乎去到要面對死的時候.
竟然透過主耶穌基督.
你竟然可以復甦過來.
其實這個就是當大家很有興趣.
甚至有時候我們會邀請一些見證人.
來當中和我們分享的見證.
當我們的婚姻出了一個很大的卡住的時候.
當我們去面對我們都不知道怎樣算的時候.
竟然我們見到主耶穌基督.
我們的婚姻就得到突破.
竟然我們可以去重修於好.

$^{241}$當我們去面對疾病那種威脅的時候.
我們是有多麼的辛苦.
但是當我遇到主耶穌基督的時候.
祂竟然可以在我疾病當中.
讓我看見出路.
當我面對我的另一半離開的時候.
我心裡面不知道怎樣去面對.
但是主耶穌的同在.
祂的永生盼望竟然可以告訴我.
我不需要擔心.
因為祂告訴我有一天我們可以再見.
所以弟兄姊妹.
保羅就是非常渴慕.
他能夠得著這樣的大能.
而他很清楚昔日的裡面.
他所去追求的是換不到回來.
所以他很願意和耶穌一同受苦.
一同受死.
為著的我要得著.
因為他很清楚如果我不是和耶穌一起去走的時候.
我不會經歷這些事情.
昨天我們團契的時候.
我們才說了其中一個點.
就是當上帝想和我們同工的時候.
他不是命令你.
小芬 我現在有件事叫你去做.
你做吧.
然後就推你出去.
推你出去自己做.
好了 到你搞不定.
不要緊 你回來吧 求我吧.
求我的時候我幫你.
我們經常都以為是這樣.
為上帝工作.
但其實不是.
上帝更重要的 更重視的那個是你.
他會在你的生命裡面出現.
然後他會告訴你.
有些事情我很想和你同工.
但卻是反過來告訴你.

$^{281}$我想怎樣做.
而帶來你生命裡面的一些矛盾.
原來當我要和上帝同工的時候.
我要放下一些東西.
那些可能是我現在追求的東西.
可能是我的工作.
可能那是我的性格.
那個可能是我待人接物的態度.
當你作出改變調教.
上帝就說 好啊 你和我同工.
而當你和他同工的時候.
你就會越來越體驗到主耶穌基督.
原來是你的真實.
原來我們的信仰是這樣出來.
所以保羅他為什麼願意和耶穌.
一同受苦 一同受死.
他知道到最後.
我就是因為這個緣故.
我能夠經歷到耶穌基督復活的大能.
所以才會出現之後的保羅.
沒有任何東西可以讓我和基督的愛隔絕.
是不是 記不記得.
是因為他們經歷了很多.
他越來越在裡面經歷到.
原來主耶穌所說的話是真實的.
原來你可以令我的生命復活的.
弟兄姊妹.
這個才是你今天真正要去追求信仰的力量.
所以上帝他不單止將物質的東西賜給你.
而更加重要的是你心靈的裡面.
你可以因為有他的緣故.
有不一樣的你.
所以保羅後來所說.
為什麼他要這樣做.
因為主耶穌要先得著我.
主耶穌用他的救恩去得著大家.
他的目的就是要讓你去得著這種的完全.
這是保羅清清楚楚和大家說的.
他拯救你不是只是給你一張很廉價的入場券上天堂.
而是他很想你去得著.

$^{321}$從你去信他那一天開始.
你就慢慢慢慢去得著這種屬靈的能力.
你可以好像保羅一樣.
這麼豪情壯語去和其他人說.
我為著耶穌的緣故.
我可以去看從前的東西是什麼.
垃圾.
所以你就會發覺.
才有我們那段精彩的經文.
忘記背後努力.
面前向著標竿直跑.
因為這樣的緣故.
他人生有一個很清晰的目標.
他知道自己要向著什麼去做.
他知道他做的每一樣東西為了什麼.
他知道做每一樣東西的時候.
到最後上帝究竟是喜悅還是不喜悅.
這個就是保羅在這段經文當中很想送給大家.
究竟你今天這個信仰.
是不是真的你的真寶.
是不是真的值得你放下這些東西.
所以我要形容他最後那句是很搞笑.
為什麼我說他的豪情壯語呢.
第一他說.
我希望靈命成熟的人都有一致的目標.
首先挑戰你.
究竟今天你的靈命成不成熟.
如果你成熟的話就應該好像保羅一樣.
但是如果你今天的整個心裡面不是這樣的時候.
你記不記得在聖經當中曾經說過.
有些人在最初沒錯.
剛剛信耶穌的人是要吃靈奶的.
但是你不可以吃一世的靈奶.
你要吃肝糧的.
你要長大的.
所以不要讓自己停留了一部分.
所以保羅對大家的挑戰就是請你成為一個靈命成熟的人.
而在這段經文當中給大家的試劑就是.
究竟你有沒有我這一份對上帝這麼大的看重.
接著他下面說.

$^{361}$我覺得很搞笑.
你不同意不要緊.
上帝會讓你明白.
他的意思就是.
他在這裡我不會再花心機和你去爭論這個問題.
你明白就是明白.
你不明白就是不明白.
但是他很相信上帝會讓你明白.
但是接著他說什麼.
但是如果你在這裡取得到一些進步的話.
你就繼續努力去做.
就是說當你今天在靈命裡面.
你有些成長.
你感受到信仰和你近一點.
上帝和你近一點的時候.
你就不要疑惑了.
你不要再想來想去了.
就繼續向前了.
我多麼相信今天當我們信耶穌的時候.
我們總是會有一些時刻.
是你會感受到上帝和你很近.
但是保羅的提醒就是.
當你有的時候.
你抓緊它.
你不要放手.
接著繼續向前走.
你不要走回頭.
繼續向前走.
所以弟兄姊妹.
我很盼望今天這段經文.
我只是很想和大家分享.
我們有沒有保羅對上帝的看重.
主耶穌基督的救恩.
上帝自己.
祂究竟是不是大家的真報靈.
求神幫助我們.
讓我們這個啞氣.
真的有寶貝在我們裡面.
我們要衷心祈禱.
七天佛我們求主你幫助帶領我們.

$^{401}$主恩我求主你帶領.
有很多時候我們自己在這個社會上.
我們真的會有迷失.
我們看到有很多很多.
我們人生的當中.
正因為我們不知道要追求什麼.
所以我們就放了很多心機.
在我們的成就,地位,讀書.
我們的生活,金錢.
我們的生活物質各方面.
但是主啊.
我知道這些一切都是你賜給我們的.
但是我們更需要的.
我們要在我們人生當中.
我們能夠有這個目標.
主啊我求主你讓我們.
去真的明白.
好像保羅所說的一樣.
你是我們人生的至寶.
當這一切去比較於上帝你的愛.
你的救恩的時候.
其實是一文不值.
因為我們有著你的緣故.
我們才能夠得到一個真真正正的完全滿足.
但是如果弟兄姊妹他們的心裡面.
有這種疑惑,有這種掙扎.
好像在他們信仰的歷程裡面都沒有經歷過.
主啊我求你幫助他們.
求你用你的方法.
讓他們去經歷.
讓他們能夠得著你.
將一切所有我們要交託給你.
其樂逢靠主耶穌.
得勝名至其琴.
阿們.
\newpage



\section{腓立比書 3:13-15}
\label{sec:vWqPqvo2gRk}
\textbf{2023-05-21 早堂《與基督同跑》腓立比書 3\_13-15 講員:楊明麗姑娘}
\newline
\newline
連結: \href{https://youtube.com/watch?v=vWqPqvo2gRk}{\texttt{https://youtube.com/watch?v=vWqPqvo2gRk}} ~~~~ 語音日期: 2023-05-23
\newline
\newline
\hyperref[sec:9tgCzMr4KBk]{\small{< < < PREV SERMON < < <}}
~
\hyperref[sec:index_chronic]{\small{[返順時目]}}
~
\hyperref[sec:index_scriptual]{\small{[返順卷目]}}
~
\hyperref[sec:wCpZ46fk4yY]{\small{> > > NEXT SERMON > > >}}
\newline
\newline
腓立比書 3:13-15
\newline
\begin{longtable}{cl}
\hline
\hline
章節 & 經文 (和合本修訂版)\\
\hline
3:13 & \begin{tabularx}{0.7\textwidth}{X} 弟兄們，我不是以為自己已經得著了；我只有一件事，就是忘記背後，努力面前的， \end{tabularx} \\ \\ \relax
3:14 & \begin{tabularx}{0.7\textwidth}{X} 向著標竿直跑，要得神在基督耶穌裡從上面召我來得的獎賞。 \end{tabularx} \\ \\ \relax
3:15 & \begin{tabularx}{0.7\textwidth}{X} 所以，我們中間凡是成熟的人，總要存這樣的心；若在甚麼事上存別樣的心，神也會把這些事指示你們。 \end{tabularx} \\ \\
[1ex]
\hline
\hline
\end{longtable}
$^{1}$各位早晨.
張麗芳姐妹.
我沒有提 張瑞芳姐妹.
真是老花了.
今天我們要.
原來我們不是要與基督同行.
今天我們要學一樣東西.
與基督同跑.
我們剛才唱了詩歌.
練定了一些戲.
我們一起祈禱.
我們一起禱告.
參加挑撥上帝.
我們需要你的恩典.
我們自己也會跑 也會走.
但我們很想你教我們.
如何與你同跑.
求你帶領我們.
讓我們聽的 說的.
都一起在你的面前.
蒙你的恩典.
蒙你的教導.
我們仰望你.
奉靠主耶穌.
基督的聖名.
阿們.
你喜歡跑嗎?.
我不知道你喜歡不喜歡跑呢.
可能每個人都有一些不同的運動.
可能你喜歡跑.
可能你喜歡高球.
可能你喜歡練習呼吸.
唱歌 示範.
練習手.
做樂師.
或者.
你喜歡做手部運動.
我相信有些人.
久不久在家中會做手指部運動.
做什麼運動的目的是什麼呢?.

$^{41}$有部分人是為了健康.
但可能很多人是為了開心.
還有.
原來有些人喜歡做眼部運動.
閱讀.
讀書很重要.
現在我們不要只做電話的手部和眼部運動.
我們也要看看神怎麼說.
保羅也是一個很有運動的.
很有全福音運動的人.
他對於開心都很重視.
他寫了很多關於喜樂的事.
特別他寫了一本書.
腓立比書.
可以稱之為喜樂的書信.
因為裡面講了很多喜樂.
用了很多字.
或者一個字的不同形態來表達喜樂.
其中他就說到.
Kara.
這個字他用了五次.
表示是一個名詞.
一起喜樂的動詞又用了兩次.
用另一個動詞.
Kara又用了九次.
總共就是在腓立比書.
他用了16次.
還有在其他書信他也說.
喜樂.
他應該在運動場.
寫這本書.
還是他在做手部運動.
摸摸那隻是紅中.
還是一索的時候所寫的.
原來不是.
他是在監獄寫這封信.
他不是去探監.
覺得自己自由所以很喜樂.
他是寫監獄書信.
在監獄當中.

$^{81}$他寫給菲律賓人的時候.
他竟然說喜樂.
我都不知道多喜樂.
其中一個喜樂的原因.
就是有人因為我坐監.
他努力去傳福音.
我喜樂.
另一個原因就是.
有人因為我坐監.
他堅定地相信.
我喜樂.
還有一個喜樂的原因.
你猜也猜不到.
有一班人去傳福音.
不過他詆毀我.
我也喜樂.
第一有人在你背後說了什麼.
無論好的不好的.
我們都學學喜樂.
他在肥樂比書中說到喜樂的時候.
他說一件事.
他完全忘記以前的野怒.
他不是中風.
他不是失憶.
他不是跌倒.
他不是被人打傷.
他只是將他以前覺得很珍貴的東西.
他忘記了.
他不是忘記了.
所有聖經的教導.
他只是將他覺得很重要的.
以往.
以往覺得很重要的.
很貴重的東西.
他全部.
我將它忘記了.
我只要最真正珍貴的.
現在最真正珍貴的.
就是我跑步的目標方向.
十格.

$^{121}$於是他就向著標竿直跑.
他這樣做.
和耶穌基督的教導.
都很符合.
很吻合.
耶穌基督說.
天國是什麼呢.
天國就好像寶貝一樣.
藏在地裡面.
誰發現了就儘快儲起錢.
買了塊地回來.
掘起它.
這樣才能拿到寶貝.
天國又好像什麼呢.
好像商人.
他很懂得看珠寶.
不會被人騙到.
他看到真正漂亮的珍珠的時候.
他就會去賣掉所有的.
要買這一顆珍珠.
放棄其他的權利和財富.
為了要得到這個寶.
寶來又說.
其實我這樣說.
我要得到寶.
其實不是說我已經得到了.
不是說我已經得到了.
我是很努力地追求得到.
或者可以得到.
真是耶穌基督.
所要我找得到的.
他就說.
弟兄啊.
弟兄我不是以為自己已經得到.
這件事.
其實是.
其實是怎樣的.
是不是楊姑娘你抄錯了.
原來不是.
他說我不是以為自己已經得到.

$^{161}$他是一節又一節.
連續兩節都要重複這件事.
我們有時覺得.
我信了主洗了禮.
我還不是得到.
保羅說.
我傳福音.
我為福音而坐牢.
我還不是自己得到.
不要這麼快以為自己得到.
我都沒有以為自己已經得到.
我只有一件事.
就是忘記背後.
努力面前.
向著我去飆竿直跑.
我們通常是背聖經.
背到這裡.
忘記了後面.
要得神在基督耶穌裡面.
從上面召我來得的獎賞.
所以我們中間.
凡是成熟的人.
總要傳這樣的心.
如果在什麼事上.
傳別的心.
神都會在這些事上指示你們.
我們都很怕神不指示我們.
好 就算走錯方向.
祂都會指示我們.
不過原來有一樣東西.
我們要忘記背後.
努力面前.
向著飆竿直跑的時候.
其實.
這個背後.
是什麼呢.
不用看後面那個.
你不是要忘記後面那個位置.
要忘記背後.
努力面前.

$^{201}$本來說要忘記背後.
就是將以前覺得很有價值的東西.
他以前覺得是無價寶的東西.
他都放下了.
放下到覺得他從來沒有發生過.
沒有擁有過.
於是他就要去得到更加寶貴的.
這個是他的目標方向.
他要去得到最寶貴的.
放下以前的.
他以前一定很窮.
不是.
其實當時的信徒.
給猶太的猶太信徒.
猶太教徒.
是一路鼓勵他們要做一樣東西.
守律法 行國禮.
這樣才可以.
但是保羅就告訴他.
這些我都做了.
我又不是說這樣才可以.
他說我不是因為我沒有行國禮.
不守律法.
所以我這樣說.
要大家來遷就我.
其實如果要用律法來誇口.
我也很能誇的.
我是第八天受國禮的.
意思就是說.
第八天受國禮是很重要的.
為什麼.
不是所有的猶太人 以色列人.
都應該是第八天受國禮的嗎.
可能還有些人.
因為當時的羅馬政府的壓迫.
還有些人會偷偷算算.
因為要讀名校的時候.
要做運動.
他們的運動服裝和我們很不同.
他們的運動服裝是一絲不掛.

$^{241}$有一些當時的羅馬政府.
對猶太人的身體特徵.
很快認出.
也會給他們一些民族上的欺凌.
他說我是第八天受國禮的.
也有另外一些人.
他不是第八天受國禮的.
是什麼呢.
外邦的信徒.
信猶太教.
不是說基督教的信徒.
外邦想進猶太教的信徒.
他是要受國禮的.
他要受國禮.
他就一定不是第八天.
因為他成年人的時候.
他成年之後.
他在社會上的經歷.
他認識了猶太教.
他的抉擇.
他要進猶太教的時候.
他要做的一件事.
就是受國禮.
所以他們不是第八天受國禮的.
本來說我是最按照傳統.
律法的指示.
正統的猶太人.
第八天受國禮的.
我是以色列族民.
我不是龜化的.
他也有另一件事.
他不是薩滿尼亞人.
他不是十三歲受國禮的薩滿尼亞人.
他是正統的猶太人.
是以色列人.
他是變雅敏支派.
變雅敏是皇族的支派.
是蘇羅王的後代.
變雅敏這個人.
是以色列雅各第十二個兒子.

$^{281}$在聖地出生的.
他是貴族的後裔.
他是希伯來人所生的希伯來人.
希伯來人這個名稱是用眾數的.
後一個希伯來人的名稱是單數的.
是指保羅茲基.
意思是說.
他是父族母族都是希伯來人.
世代都是希伯來人.
是家族說希伯來話的人.
不只是有希伯來血統的以色列人.
言語上血脈上他都是希伯來人.
按法說.
他不只是希伯來人.
以色列人.
第八十壽國禮.
他還是一個敬虔的宗派.
法利賽人.
他還是法利賽人當中.
他在迦拉塔書有寫.
他就是以加別的馬加列的學生.
他是預備要去承受.
這個宗派這個組織的負責人.
但是他以前的熱心.
是熱心到猶太人都覺得他很熱心.
基督徒都怕了他.
他就是要申請一道羅馬政府的命令.
可以令到他可以抓所有的基督徒去坐牢.
然後才處理他們.
你要和我按律法去看.
是不是真的有人比我更加守律法呢.
按律法看我是無可指責的.
按律法看他是無可指責.
意思是說.
每天要祈禱多少次.
我祈禱的.
每個禮拜要禁食多少次.
我禁食的.
我不是要祈禱的時候.
站在十字路口告訴人.

$^{321}$我現在時間到了我要祈禱.
在街上我也要祈禱.
他不是要禁食.
然後走出街告訴人.
哎呀今天好像沒什麼氣沒什麼力氣.
讓大家都知道他要禁食.
他很認真地.
很認真地去守律法的人.
所以就算你用律法去看我.
都是無可指責的.
不過他說.
以前的事忘記它.
以前的事.
我怎樣在律法上面看.
我應該是很貴族的.
我應該是很敬虔的.
我應該是很認真的.
我現在都將它放下了.
因為信主耶穌不是要我去守律法.
信主耶穌是要我依靠主耶穌.
而去蒙恩得救.
所以現在我只有一件事.
就是忘記背後努力面前.
向著標竿直跑.
要得神在基督耶穌裡面.
從上面召我來得的獎賞.
這個好像很不容易.
怎樣可以得到這件事呢.
如果他忘記前面.
按著律法.
他不用坐牢了.
他不需要自己再受第二次的國禮.
因為他第八天已經受了國禮.
守律法對他沒有什麼影響.
他一向都守開了.
他只是說別人不用守律法.
不用守國禮.
新的信徒.
外邦的信徒.
不需要守國禮.

$^{361}$不需要守律法.
飲食上不需要.
行為上都不需要.
這樣他就被抓去坐牢.
其實抓他去坐牢.
那個人是想保護他.
如果不抓他去坐牢.
他應該被人殺了.
打死了.
用石頭打死.
為什麼這個人坐在監獄.
他會喜樂呢.
還要忘記背後努力面前.
如果一個人.
他說忘記背後努力面前.
坐在監獄裡面.
你當然想著他.
要挖一條地道走出來.
但保羅不是.
他勸導大家要做一件事.
就是要有一個目標方向.
向著那個目標方向走.
那個目標方向是什麼呢.
以前守律法.
他們都是要得赦免.
有永生的盼望.
現在不守律法.
目標是什麼.
得恩典.
得永生的恩典.
得永生的恩典的時候.
他就可以.
可以拿著這個恩典.
他可以得到主耶穌賜賜他的永生.
而別人也可以拿著這個恩典.
他向別人傳福音.
不用別人守律法.
這群人也可以做一件事.
靠著主耶穌.
全心相信主耶穌.

$^{401}$向著標竿直跑的時候.
得著神在基督耶穌裡面.
照著我們來得的獎賞.
我們可以不需要守律法.
去得蒙這個獎賞.
其實守律法不單止是說.
男士受國禮.
女士也有很多影響.
例如說.
他們是守星期六是安息日.
星期五晚之前.
預備了所有星期六日要的東西.
就算你要吃東西的時候.
你不要說.
我一會兒去錄一個公仔麵就行了.
是不可以的.
這樣就是工作.
你應該要準備好東西.
擺好在某些地方.
太陽曬下來就會溫暖.
還要指定什麼地方可以擺.
什麼地方不可以擺.
什麼植物旁邊不能擺.
什麼植物旁邊可以擺.
什麼石頭可以擺.
很煩.
對我們來說.
不用守律法的確是很喜樂.
但是保羅來說.
什麼叫做不用守律法.
什麼叫做忘記背後.
其實他做的事情.
不只是他做過的.
應該在舊約也有人忘記背後.
努力面前的.
我們今天坐在一起崇拜的人.
都比較年輕.
沒有哪位應該是很熟悉.
粵語殘片.
是否有很多.

$^{441}$某時某段.
前一代的胡楓.
不知道跟哪一位一起.
之後胡楓.
到他後來被仇家殺了之後.
誰知道兒子可以在仇家成長.
到第二代胡楓出來的時候.
18年之後的時候.
媽媽來找他.
告訴他他的身世.
然後他要離開.
那個被他富富貴貴.
開開心心.
開著跑車生活的家庭.
去跟他工廠當中.
幫人洗衣服的媽媽生活的時候.
那種忘記背後.
舊約有一位.
記不記得是誰.
他過著豪華的生活.
他有他的車.
有他的隨從.
有權柄.
他有一個高位.
他是公寓給他的.
他是摩西.
是誰呢.
摩西是以色列人.
摩西是利未人的後裔.
他父母都是利未人.
他是利未人所生的利未人.
但是呢.
他是被法老的女兒收養了.
在河邊救了他回來.
然後他就過著皇族的生活.
得到皇族的教導.
熟悉當代多國的語言和智慧.
但是一旦他知道自己是以色列人的時候.
他最後就是要忘記背後.
不做法老的孫.

$^{481}$不是將來要坐上法老的位.
他去了哪裡呢.
他要逃亡.
去到曠野.
後來幫他外撫牧羊.
幫他外撫牧羊.
就是說那些羊是他外撫的.
不一定是他.
他所得的是什麼呢.
他所得的是好像奴隸一樣逃亡的身份.
還要是比奴隸更加不可以出來見人.
啊 忘記背後.
忘記背後原來不是我們做錯了事.
那道遠協忘記了他就算了.
我們都忘記別人做錯的事.
得罪我們的事.
原來忘記背後是說忘記你輝煌的身份.
以往的價值觀.
以往令你覺得很珍惜的.
甚至別人對你很欣賞的.
無論摩西或者保羅他們都有這樣的背景.
別人都很欣賞 很羨慕他們.
但是當日的時候.
摩西在山上.
他在荊棘蟲當中.
真正認識到耶和華是他的神.
保羅拿著一張拘禁的命令.
他可以拘禁所有基督徒的時候.
在那個時候.
他見到向他顯現的耶穌基督.
他在意象當中見到耶穌基督.
他向祂顯現的時候.
他知道祂是神.
這個神就說.
保羅啊保羅.
你為什麼逼迫我.
讓我重新看自己的生命.
自己的信仰.
我這麼努力這麼認真去守律法的時候.
神說我逼迫他.

$^{521}$神說我要走的一條路.
和他要指示我走的一條路.
是相反方向.
不是他要我走這條路.
不是他召我來要得的獎賞.
他對我另有一條路安排.
弟兄姊妹你們都有你們的身份.
你們都有你們的背景.
有你們珍惜的地方.
甚至在你們信仰上的成長.
在你們教會裡的成長.
都有一些你珍惜的.
我們不是沒有感情的人.
我們一定有我們珍惜的事物.
但是如果這些事物.
好像保羅一樣.
好像摩西一樣.
這些這麼好的quality.
這麼好的背景.
這麼好的經歷.
這麼好的身份.
阻礙我們信耶穌.
阻礙我們跟隨耶穌基督的教導去做人.
阻礙我們去得上帝召我們要得的獎賞.
去做上帝要我們做的.
我們能不能夠學習保羅一樣.
我們只有一件事就是忘記背後.
我相信絕大部分人.
如果你的情況是很好的時候.
也沒有人說你這樣做是錯的.
你現在的信仰情況是錯的.
沒有人這樣說的.
保羅如果他守以前那套.
每個人都很歡迎他.
沒有人說他錯的.
甚至可能他帶來信主.
信耶穌基督領受救恩的人.
都未必說他錯的.
這樣就很好.
兩邊都搞定.

$^{561}$但是他說這件事.
不單止阻礙他信主.
跟隨主得獎賞.
還阻礙其他人信主.
所以這件事.
他覺得是很寶貴的事.
他就放下了它.
因為要得更貴重的東西.
就像主耶穌基督所說.
變賣一切所有的.
要得著最美最珍貴的珍珠.
是.
要得著最美最珍貴的珍珠.
我們要向著標竿直跑.
這個向著標竿直跑.
是一件什麼事呢.
他在《耳窟所書》當中也說.
我們不要想自己做了什麼.
我有多好有多厲害.
或者任之乎多厲害.
但只有一件事是最重要的.
就是我有基督的恩典.
有基督的恩典.
不是靠行為.
但是我們都有好的行為.
我們都不是.
未信主之前.
每個人都做壞人.
他說.
就算我們有做了什麼好事.
其實我們是基督的工作.
我們原始他的工作.
是基督耶穌裡面造成的.
為了叫我們行善.
就是神所預備叫我們行的.
我們是神在基督耶穌裡面造成的工作.
我們做得好的事.
是神預備給我們做的.
你剛好走到那裡.
有個人掉了個錢包.

$^{601}$你可以告訴他.
這位你掉了個錢包.
我幫你撿回來.
行善的機會.
不是被騙的機會.
都是神預備的.
我們能夠行善.
是因為基督的工作.
神的恩典.
我們可以怎樣呢.
他說我們就是要努力.
向著這個目標前進.
這個目標.
保羅說.
是我們能夠放下所有東西.
只是望著十字架的道路前進.
好像大家都在做.
我們應該全部都起步了.
在跑道上.
但跑道上.
時時都會有些分叉位.
可以跑錯路.
你不會跑錯隔離道路.
凡是願意跟隨主耶穌基督的人.
你都會發現.
生命裡面.
時時都不喜歡你做下決定.
到底向左還是向右.
為何正路總是這麼窄.
好像走過去都會被扯到.
但路不是正路.
那條又闊闊大大.
舒舒服服.
我是可以這樣走的.
很多時候.
我們都有這些抉擇的時候.
甚至有一件T恤都告訴我們.
到底耶穌基督在這樣的情況下.
他會想什麼呢.
他會怎樣決定呢.

$^{641}$我們要學習耶穌基督的抉擇.
學習耶穌基督的抉擇.
記不記得約莫三十年前.
或者二十年前.
如果你有兒女在這個年紀的話.
或者當時你是青少年的話.
當時的父母.
是很不想兒女去做一個事.
那樣東西叫做任天堂.
現在不是.
現在的父母是很怕兒女下地獄.
連電子遊戲機都由天堂去到地獄.
我們都要轉變.
我們的價值觀都要轉變.
我們的生命成為別人的一個榜樣.
看下去都不知道到底是要上天堂還是下地獄呢.
我們的生命被基督去帶領的時候.
我們怎樣去跑前面的路呢.
目標很清晰.
我們是要得恩典.
我們沒有人是進來不是想得永生.
我們得永生.
要得上次.
得基督為我們所預備的.
不過每一天我們可能做了一些.
除了目標方向跟著基督之外.
還有很多其他事我們攬上身都想去做.
跟著停在中間甚至乎退後兩步.
好像老人家所說.
我走一步就退兩步.
怎樣都去不到目標.
我們的生命不要變成這樣的情況.
不要讓我們所喜好的事.
你所喜好的運動.
我們就與基督一同去跑.
不要讓你所喜好的運動.
或者習慣.
成為一些與神的隔遠了一點.
不要讓我們一些隨著時代.
讓時代的價值觀影響了我們的時候.

$^{681}$覺得這樣我都算是依靠基督吧.
我每天都有祈禱的.
起碼一次.
每個星期我都有崇拜的.
起碼在家中Zoom一下吧.
我聽到的.
不過我把時間調校了兩倍.
不是我的錯.
說完說得慢.
現在這個時勢我們要快.
我不是沒有傳福音的.
我時時在麻將桌傳福音的.
是不是呢.
我不是沒有帶人回教會的.
教會有什麼表演.
節目.
旅行.
我都找人來的.
兄弟姐妹.
我們再做多一點.
我們朝著這個目標方向.
我們再做多一點.
我們找我們生命裡面.
要被神更加得著的地方.
有什麼是我們不是靠基督的.
我記得讀神學的時候.
神學院的同學和老師.
都會說一些當時的情況.
有些人是靜看醫生不吃藥.
有些人是靜祈禱不看醫生.
有些人是三樣都做的.
三樣都做的意思.
他不只是借飯.
都為吃藥這件事去禱告.
為疾病也去禱告.
為找醫生看來去禱告.
原來有人是不看醫生.
不吃藥.
只祈禱讓疾病好回來的.
可能有機會的.

$^{721}$但不一定每次都是這樣.
我們吃飯.
祈禱.
我們都可以借飯.
和紀念一些身邊.
你想紀念的人.
我們跟隨主的腳步走的時候.
身邊有哪一位.
你很想他今年信主嗎.
你什麼時候開始邀請他.
不是.
邀請他之前你做下鬆禱.
先祈禱.
你什麼時候開始為他祈禱.
他什麼時候進入你祈禱的那本寶.
是不是排在最後.
你會一個星期為他祈禱多少次.
對不起.
應該是一天.
你打算為他祈禱多少次.
然後帶他來祈禱會.
你打算帶他來祈禱會之前.
為他做些什麼.
讓他思想一下生命.
雖然你不是直接和他傳福音.
都要讓他思想一下生命.
然後帶他回來報道會.
你預備了一個土壤給他.
讓他做些信主.
特別是家人.
家人不容易向他傳福音.
因為你很多時候要接觸他.
你很多時候要和他談其他的事情.
他看到你整個的生命.
要邀請是很困難.
但是他看到你的生命.
看到你的轉變.
你要邀請他又可以很有力量.
為他去禱告.
你知道講完是哪位.

$^{761}$為他去禱告.
為那個講完祈禱.
不是那個講完出不出名.
偉不偉大.
而是我們是否靠著主去開一個報道會.
我們是否靠著主.
以基督的大能大力去引領人回教會.
這個講傳福音不是因為他常常坐在監中.
而是他常常倚靠神.
當去到哪個時候.
神要他.
你不是去這個地方.
你去另一個城市的時候.
他是願意轉他的旋.
轉方向.
跟隨主的引導.
求主幫助我們.
跟隨他的引導.
讓我們能夠在他的引導下.
好好地去跑他為我們預備的路.
讓我們朝著十字架.
朝著他的方向去努力去跑.
有些人.
他身體很健康.
他都跑這條路.
不一定是很喜樂.
有些人.
他身體不是很健康.
他都跑這條路.
好像很夠氣很喜樂.
我認識有幾位傳道人.
其實他們身體健康是比我差.
我不能跑步.
不夠氣.
他們是帶著一個癌症患者的身體.
經常進醫院.
然後可以跟隔離位的病人傳福音.
我心想你已經進了醫院.
氣都沒有多一口.
不留一口氣給自己休息.

$^{801}$還可以去傳福音.
原來人是這樣的.
我們沒有的時候.
你真的禱告以靠神的時候.
好像是有的.
又真的有機會.
真的隔離床的人會聽他們傳福音.
會被他帶到信主.
有什麼地方.
什麼事情.
是阻攔著我們.
事奉神.
不是地方.
不是健康.
可能是我們怎樣看待我們的事奉.
我記得有姐妹告訴我.
她會跟一位黃先生學跑步.
她上了幾堂課就告訴我.
黃先生叫我忘記以前的跑步方法.
重新學跑步.
我不知道應該怎樣跑.
我沒有學過沒有跟過.
但原來跑步都是這樣.
我們很多事情都是這樣的.
要真的學回基本.
事奉我們的基本是什麼.
事奉的基本我們不是表現自己.
事奉的基本我們是忘記自己.
好像保羅所說.
忘記我們自己是誰.
讓神可以使用我們.
事奉的基本.
我們是要做到神要我們去做到的.
耶穌基督想.
你在這個事奉裡面做到什麼.
我們就朝著那個方向去.
這個是我們每一件事應該有的抉擇.
有些事情.
有些決定你以往已經做了.
你改變不了以往的決定.

$^{841}$但就是今天的路途上.
你開始在你現在這個狀況當中.
全部交給神的時候.
你怎樣去決定前面的道路.
包括你的婚姻情況.
包括你教子女的方法.
我都說了.
以前的父母不讓兒童上去任天堂.
但現在的父母不讓兒童去下地獄.
怎樣都好.
在你的情況當中.
耶穌基督會想你的婚姻怎樣.
我們今天下午有一個團契都會說這件事.
我們.
在我們和子女的關係當中.
祂會想你們家和子女的關係是怎樣.
我們一起學.
學習.
跑步.
我們是向著主為我們預備的方向.
不要調轉方向.
不要跑到旁邊.
讓我們向著前面跑.
放心.
你需要水的時候.
需要一顆橙的時候.
主耶穌基督在前面遞給你.
有沒有困難呢.
普羅他說.
所以我們中間凡是有成熟的人.
在信仰上已經成熟的人.
我們要存著這樣的心.
他不是說我們要做到.
存著這樣的心和做到還有一段距離.
但我們要存著這樣的心.
如果在第二次世上我們有第二樣的心的時候.
我們上上禱告將它交給神的時候.
讓神去指示我們.
然後.
普羅說了一句很有溫情的話.

$^{881}$他說但是呢.
我又不是逼你們.
然而我們去到什麼地步.
就照著那個地步走.
我們不需要回到以往.
我們不需要要求我今天開始一百分.
但我今天開始學習存著這樣的心.
我要去做決策的時候.
要去做計劃的時候.
我就向著新的方向去學習.
有哪裡走錯了就按那個時候的地步.
走回一個正確的方向.
進入大道.
朝著標竿直跑.
我們一起禱告.
親愛的天父上帝我們謝謝你.
天父我們謝謝你.
謝謝你.
你看顧我們.
你沒有計較我們的以往.
你沒有計較我們的身份.
我們的背景.
你沒有計較我們現在的心思.
但是你就是願意照我們.
叫我們能夠得著永生.
你就願意照我們.
讓我們在你的計劃當中.
可以有份.
讓我們可以在你的安排當中.
我們可以得著美好.
讓我們因著依靠你.
我們可以得到喜樂.
主耶穌基督他一生跑得太好了.
無論去到哪裡.
他正在治病.
人們讚賞他.
想他留下來繼續治病.
他都會朝著正確的方向.
可能去第二座城.
可能開始傳福音.

$^{921}$主耶穌我們不懂.
我們迷惘.
我們時時會被人影響.
又被自己影響.
求你建立我們.
幫助我們.
讓我們能夠在你的引導下.
我們遲時朝著正確的方向.
向著標竿直跑.
得著你在基督耶穌裡面.
照我們來得的獎賞.
求主憐憫.
我們仰望你奉靠主耶穌基督的聖名.
\newpage



\section{路加福音 1:5-25-57-66}
\label{sec:wCpZ46fk4yY}
\textbf{2023-05-21 加恩樓《不相稱的禱告》路加福音 1\_5-25,57-66 講員:陳韋安牧師}
\newline
\newline
連結: \href{https://youtube.com/watch?v=wCpZ46fk4yY}{\texttt{https://youtube.com/watch?v=wCpZ46fk4yY}} ~~~~ 語音日期: 2023-05-24
\newline
\newline
\hyperref[sec:vWqPqvo2gRk]{\small{< < < PREV SERMON < < <}}
~
\hyperref[sec:index_chronic]{\small{[返順時目]}}
~
\hyperref[sec:index_scriptual]{\small{[返順卷目]}}
~
\hyperref[sec:aPVZvLNc78I]{\small{> > > NEXT SERMON > > >}}
\newline
\newline
路加福音 1:5-25-57-66
\newline
\begin{longtable}{cl}
\hline
\hline
章節 & 經文 (和合本修訂版)\\
\hline
1:5 & \begin{tabularx}{0.7\textwidth}{X} 在希律作猶太王的時候，亞比雅班裡有一個祭司，名叫撒迦利亞；他妻子是亞倫的後代，名叫伊利莎白。 \end{tabularx} \\ \\ \relax
1:6 & \begin{tabularx}{0.7\textwidth}{X} 他們兩人在神面前都是義人，遵行主的一切誡命和條例，沒有可指責的。 \end{tabularx} \\ \\ \relax
1:7 & \begin{tabularx}{0.7\textwidth}{X} 只是他們沒有孩子，因為伊利莎白不生育，兩個人又年紀老邁了。 \end{tabularx} \\ \\ \relax
1:8 & \begin{tabularx}{0.7\textwidth}{X} 撒迦利亞按班次在神面前執行祭司的職務， \end{tabularx} \\ \\ \relax
1:9 & \begin{tabularx}{0.7\textwidth}{X} 照祭司的規矩抽籤，進到主的殿裡燒香。 \end{tabularx} \\ \\ \relax
1:10 & \begin{tabularx}{0.7\textwidth}{X} 燒香的時候，眾百姓在外面禱告。 \end{tabularx} \\ \\ \relax
1:11 & \begin{tabularx}{0.7\textwidth}{X} 有主的一個使者站在香壇的右邊，向他顯現。 \end{tabularx} \\ \\ \relax
1:12 & \begin{tabularx}{0.7\textwidth}{X} 撒迦利亞看見，就驚慌害怕。 \end{tabularx} \\ \\ \relax
1:13 & \begin{tabularx}{0.7\textwidth}{X} 天使對他說：「撒迦利亞，不要害怕，因為你的祈禱已經被聽見了。你的妻子伊利莎白要給你生一個兒子，你要給他起名叫約翰。 \end{tabularx} \\ \\ \relax
1:14 & \begin{tabularx}{0.7\textwidth}{X} 你必歡喜快樂；有許多人因他出世也必喜樂。 \end{tabularx} \\ \\ \relax
1:15 & \begin{tabularx}{0.7\textwidth}{X} 他在主面前將要為大，淡酒烈酒都不喝，從母腹裡就被聖靈充滿。 \end{tabularx} \\ \\ \relax
1:16 & \begin{tabularx}{0.7\textwidth}{X} 他要使許多以色列人回轉，歸於主—他們的神。 \end{tabularx} \\ \\ \relax
1:17 & \begin{tabularx}{0.7\textwidth}{X} 他將有以利亞的精神和能力，走在主的前面，叫父親的心轉向兒女，叫悖逆的人轉向義人的智慧，又為主預備迎接他的百姓。」 \end{tabularx} \\ \\ \relax
1:18 & \begin{tabularx}{0.7\textwidth}{X} 撒迦利亞對天使說：「我怎麼能知道這事呢？我已經老了，我的妻子也年紀老邁了。」 \end{tabularx} \\ \\ \relax
1:19 & \begin{tabularx}{0.7\textwidth}{X} 天使回答他說：「我是站在神面前的加百列，奉差遣來對你說話，把這好信息報給你。 \end{tabularx} \\ \\ \relax
1:20 & \begin{tabularx}{0.7\textwidth}{X} 到了時候，這些話必然應驗；只因你不信我的話，你會成為啞巴，不能說話，直到這些事實現的日子。」 \end{tabularx} \\ \\ \relax
1:21 & \begin{tabularx}{0.7\textwidth}{X} 百姓等候撒迦利亞，詫異他在聖所裡遲延那麼久。 \end{tabularx} \\ \\ \relax
1:22 & \begin{tabularx}{0.7\textwidth}{X} 到他出來，卻不能和他們說話，他們就知道他在聖所裡見了異象；他直向他們打手勢，因為他成了啞巴。 \end{tabularx} \\ \\ \relax
1:23 & \begin{tabularx}{0.7\textwidth}{X} 他供職的日子一滿，就回家去了。 \end{tabularx} \\ \\ \relax
1:24 & \begin{tabularx}{0.7\textwidth}{X} 這些日子以後，他的妻子伊利莎白就懷孕，隱藏了五個月； \end{tabularx} \\ \\ \relax
1:25 & \begin{tabularx}{0.7\textwidth}{X} 她說：「主在眷顧我的日子，這樣看顧我，要除掉我在人前的羞恥。」 \end{tabularx} \\ \\ \relax
1:26 & \begin{tabularx}{0.7\textwidth}{X} 到了第六個月，天使加百列奉神的差遣往加利利的一座城去，這城名叫拿撒勒， \end{tabularx} \\ \\ \relax
1:27 & \begin{tabularx}{0.7\textwidth}{X} 到一個童女那裡，她已經許配大衛家的一個人，名叫約瑟；童女的名字叫馬利亞。 \end{tabularx} \\ \\ \relax
1:28 & \begin{tabularx}{0.7\textwidth}{X} 天使進去，對她說：「蒙大恩的女子，你好，主和你同在！」 \end{tabularx} \\ \\ \relax
1:29 & \begin{tabularx}{0.7\textwidth}{X} 馬利亞因這話就很驚慌，又反覆思考這樣問候是甚麼意思。 \end{tabularx} \\ \\ \relax
1:30 & \begin{tabularx}{0.7\textwidth}{X} 天使對她說：「馬利亞，不要怕，你在神面前已經蒙恩了。 \end{tabularx} \\ \\ \relax
1:31 & \begin{tabularx}{0.7\textwidth}{X} 你要懷孕生子，要給他起名叫耶穌。 \end{tabularx} \\ \\ \relax
1:32 & \begin{tabularx}{0.7\textwidth}{X} 他將要為大，稱為至高者的兒子；主神要把他祖先大衛的王位給他。 \end{tabularx} \\ \\ \relax
1:33 & \begin{tabularx}{0.7\textwidth}{X} 他要作雅各家的王，直到永遠；他的國沒有窮盡。」 \end{tabularx} \\ \\ \relax
1:34 & \begin{tabularx}{0.7\textwidth}{X} 馬利亞對天使說：「我沒有出嫁，怎麼會有這事呢？」 \end{tabularx} \\ \\ \relax
1:35 & \begin{tabularx}{0.7\textwidth}{X} 天使回答她說：「聖靈要臨到你身上；至高者的能力要庇蔭你，因此，那要出生的聖者要稱為神的兒子。 \end{tabularx} \\ \\ \relax
1:36 & \begin{tabularx}{0.7\textwidth}{X} 況且，你的親戚伊利莎白，就是那素來稱為不生育的，在年老的時候也懷了男胎，現在懷孕六個月了。 \end{tabularx} \\ \\ \relax
1:37 & \begin{tabularx}{0.7\textwidth}{X} 因為，出於神的話，沒有一句不帶能力的。」 \end{tabularx} \\ \\ \relax
1:38 & \begin{tabularx}{0.7\textwidth}{X} 馬利亞說：「我是主的使女，願意照你的話實現在我身上。」於是天使離開她去了。 \end{tabularx} \\ \\ \relax
1:39 & \begin{tabularx}{0.7\textwidth}{X} 在那些日子，馬利亞起身，急忙前往山區，來到猶大的一座城， \end{tabularx} \\ \\ \relax
1:40 & \begin{tabularx}{0.7\textwidth}{X} 進了撒迦利亞的家，向伊利莎白問安。 \end{tabularx} \\ \\ \relax
1:41 & \begin{tabularx}{0.7\textwidth}{X} 伊利莎白一聽到馬利亞問安，所懷的胎就在腹裡跳動。伊利莎白被聖靈充滿， \end{tabularx} \\ \\ \relax
1:42 & \begin{tabularx}{0.7\textwidth}{X} 高聲喊著說：「你在婦女中是有福的！你所懷的胎也是有福的！ \end{tabularx} \\ \\ \relax
1:43 & \begin{tabularx}{0.7\textwidth}{X} 我主的母親到我這裡來，為何這事臨到我呢？ \end{tabularx} \\ \\ \relax
1:44 & \begin{tabularx}{0.7\textwidth}{X} 因為你問安的聲音一入我耳，我腹裡的胎就歡喜跳動。 \end{tabularx} \\ \\ \relax
1:45 & \begin{tabularx}{0.7\textwidth}{X} 這相信的女子是有福的！因為主對她所說的話都要應驗。」 \end{tabularx} \\ \\ \relax
1:46 & \begin{tabularx}{0.7\textwidth}{X} 馬利亞說：「我心尊主為大； \end{tabularx} \\ \\ \relax
1:47 & \begin{tabularx}{0.7\textwidth}{X} 我靈以神我的救主為樂； \end{tabularx} \\ \\ \relax
1:48 & \begin{tabularx}{0.7\textwidth}{X} 因為他顧念他使女的卑微；從今以後，萬代要稱我有福。 \end{tabularx} \\ \\ \relax
1:49 & \begin{tabularx}{0.7\textwidth}{X} 因為那有權能的為我做了大事；他的名是聖的。 \end{tabularx} \\ \\ \relax
1:50 & \begin{tabularx}{0.7\textwidth}{X} 他憐憫敬畏他的人，直到世世代代。 \end{tabularx} \\ \\ \relax
1:51 & \begin{tabularx}{0.7\textwidth}{X} 他用膀臂施展大能；他趕散心裡妄想的狂傲人。 \end{tabularx} \\ \\ \relax
1:52 & \begin{tabularx}{0.7\textwidth}{X} 他叫有權柄的失位，叫卑賤的升高。 \end{tabularx} \\ \\ \relax
1:53 & \begin{tabularx}{0.7\textwidth}{X} 他叫飢餓的飽餐美食，叫富足的空手回去。 \end{tabularx} \\ \\ \relax
1:54 & \begin{tabularx}{0.7\textwidth}{X} 他扶助了他的僕人以色列，不忘記施憐憫， \end{tabularx} \\ \\ \relax
1:55 & \begin{tabularx}{0.7\textwidth}{X} 正如他對我們的列祖說過，『憐憫亞伯拉罕和他的後裔，直到永遠。』」 \end{tabularx} \\ \\ \relax
1:56 & \begin{tabularx}{0.7\textwidth}{X} 馬利亞和伊利莎白同住，約有三個月，然後回家去了。 \end{tabularx} \\ \\ \relax
1:57 & \begin{tabularx}{0.7\textwidth}{X} 伊利莎白的產期到了，生了一個兒子。 \end{tabularx} \\ \\ \relax
1:58 & \begin{tabularx}{0.7\textwidth}{X} 鄰里親屬聽見主向她大施憐憫，就和她一同歡樂。 \end{tabularx} \\ \\ \relax
1:59 & \begin{tabularx}{0.7\textwidth}{X} 到了第八日，他們來給孩子行割禮，並要照他父親的名字叫他撒迦利亞。 \end{tabularx} \\ \\ \relax
1:60 & \begin{tabularx}{0.7\textwidth}{X} 他母親回應說：「不！要叫他約翰。」 \end{tabularx} \\ \\ \relax
1:61 & \begin{tabularx}{0.7\textwidth}{X} 他們對她說：「你親族中沒有叫這名字的。」 \end{tabularx} \\ \\ \relax
1:62 & \begin{tabularx}{0.7\textwidth}{X} 他們就向他父親打手勢，問他這孩子要叫甚麼名字。 \end{tabularx} \\ \\ \relax
1:63 & \begin{tabularx}{0.7\textwidth}{X} 他要了一塊寫字的板，寫上：「他的名字是約翰。」他們就都驚訝。 \end{tabularx} \\ \\ \relax
1:64 & \begin{tabularx}{0.7\textwidth}{X} 撒迦利亞的口立刻開了，舌頭也鬆了，就開始說話稱頌神。 \end{tabularx} \\ \\ \relax
1:65 & \begin{tabularx}{0.7\textwidth}{X} 周圍居住的人都懼怕；這一切的事就傳遍了猶太山區。 \end{tabularx} \\ \\ \relax
1:66 & \begin{tabularx}{0.7\textwidth}{X} 凡聽見的人都把這事放在心裡，他們說：「這個孩子將來會怎麼樣呢？」因為有主的手與他同在。 \end{tabularx} \\ \\
[1ex]
\hline
\hline
\end{longtable}
$^{1}$.
今天是一個很好的敬拜.
我們今天是來求主幫助我們.
祂將祂的說話跟我們說.
我們一起來禱告.
天父我們懇求你的聖靈在我們當中.
在你的說話,在你的課堂書裡.
你昔日的教導.
即是撒卡利亞的故事.
我們一起來聆聽.
求主幫助孩子.
雖然孩子不配.
但是你這樣引導我們.
讓我們能夠明白.
讓我們能夠在土地生命裡回應你.
求主你這樣的幫助.
讓我們今天一起來恭敬,聆聽.
我們期待你對我們每個人的生命裡的說話.
求主你這樣去臨界.
在當中親自告訴我們.
我們當行的路.
求主你這樣的幫助.
奉主名求,阿們.
我們會說一篇有關禱告的課題.
剛才敬拜隊有一首歌很好的歌詞.
說到我們在疑惑的時候.
我們仍然去禱告.
禱告是我們最基本的渣拿.
當我們什麼都沒有的時候.
當我們的生命裡.
去到多麼不濟的時候.
我們多麼差的時候.
我們仍然可以禱告.
就算我們是疑惑.
就算我們有很多困難的時候.
仍然可以禱告.
有人問我們.
罪人的禱告有沒有效呢?.
如果我們犯罪的時候.
也都禱告沒有用的話.

$^{41}$我們就沒有辦法了.
所以禱告就算我們是罪人也好.
我們仍然可以向上帝去禱告.
是我們最基本.
最後我們最終的.
去上帝親近的方式.
所以我們一起來思考一個有關禱告的課題.
今天用的經文是《路加福音》第一章.
這段經文可能大家都比較熟悉.
是一段聖誕節的故事.
我們平日在聖誕節裡都會看到.
瑪利亞和耶穌出生的故事.
但我們比較少留意另一段的故事.
第一章裡有兩個很平衡的故事出現.
不單單是一個聖誕節的故事.
同時有另一對的夫婦.
懷孕的故事也記載在第一章裡.
其實兩段經文是一個很大的對比和平衡.
無論是瑪利亞生孩子也好.
還是以利沙生孩子也好.
兩段經文是一個很大的對比.
第五到二十五節是天使加伯列.
對撒迦利亞顯現的記載.
二十六到三十八節是加伯列向瑪利亞顯現的記載.
四十六到五十六節是瑪利亞的禱告.
是專注在這段經文.
六十七到七十九節是撒迦利亞的禱告.
其實是兩段很平衡的故事.
結合了兩個媽媽 兩個生孩子的故事.
我們就嘗試去看另一段的故事.
不是約瑟瑪利亞生孩子.
而是撒迦利亞和以利沙伯生的約翰的故事.
今天我們嘗試用禱告這個關鍵字來理解這段故事.
其實你會發現路加福音是一個非常強調禱告的小說.
小說裡也有很多禱告.
但路加福音是有很多有關禱告的主題.
舉個例子.
路加福音第十一章是一整章都是禱告的篇章.
特別是故意編排著主禱文.
門外求餅的比喻.

$^{81}$祈求祖國的教導.
都放在這章第十一章裡.
路加第十六章更是講到耶穌揀選門徒.
當耶穌揀選門徒的故事裡.
唯獨路加福音是加了耶穌通宵祈禱的片段.
只有路加福音有.
路加福音特別強調耶穌是特意通宵禱告來揀選門徒.
路加第十八章也有很多獨家有關禱告的比喻.
法拉塞人和瑞利禱告的比喻.
今天我們嘗試用路加禱告的主題來看看第一章的部分.
經文中有很多部分是故意編排著禱告的關係.
說的不是純粹是專注於禱文.
而是整個故事都在禱告的篇幅裡被描寫.
看看第五章的第十節.
經文說在希律作猶太王的時候.
阿俾亞班裡有一個祭司名叫撒加尼亞.
他妻子是亞倫的後代名叫伊利沙伯.
他們兩人在神面前都是義人.
進行著一切戒命和條例.
沒有可指責的.
只是他們沒有孩子.
因為伊利沙伯不生育.
兩個人又年紀老了.
第八節撒加尼亞就按班次在神面前執行祭司的職務.
按照祭司的規矩抽籤.
進到主的殿裡燒香.
燒香的時候眾百姓在外面禱告.
有主的一個使者站在香壇的右邊.
向他顯現.
首先我們知道撒加尼亞是一個什麼樣的人.
根據路加記載他是一個祭司.
祭司是做什麼的.
祭司的職責基本上就是在聖殿裡燒香.
祭司是唯一一個能夠進入聖殿聖所的人.
他是代表著眾百姓的禱告.
《瑪俐伯論文》裡的禱告.
Tafina這個字就是如煙往上升的意思.
是一樣的字.
所以燒香不是一個宗教禮儀.
燒香是代表著他們的禱告.

$^{121}$是藉著祭司的手往上升到上帝的面前.
所以根據猶太律法.
祭司的角色就是一個中保角色.
他是代表著眾百姓.
來將他們的禱告.
藉著燒香呈獻給上帝.
簡單來說.
祭司是一個中介者.
是一個專業的幫人.
將他們的禱告獻給上帝的一個人.
所以發覺祭司是一個這樣的人.
是一個專業將人的禱告獻給上帝的職責的人.
因此傑文特別強調.
天使是向撒迦利亞顯現的時候.
正正就是這個禱告的時刻.
當眾百姓在禱告的時候.
就如剛才所說.
燒香的時候眾百姓在外面禱告.
就在這個時候.
有主的使者站在香壇的右邊向祂顯現.
專業地專門為眾百姓的禱告呈獻給上帝的人.
在眾百姓禱告的時候.
天使顯現並告訴他們.
你的禱告中了.
很多年前你曾經的禱告已經應允了.
所以用一個我們現代教會的語言.
你想想星期三晚教會祈禱會.
當每個人在祈禱的時候.
牧師自己的禱告就應允了.
所以一個非常之禱告的時刻.
一個專門幫人代禱的人.
一個祈禱會的時候.
眾百姓祈禱的時候.
他自己的禱告就蒙英文.
所以是一個非常之跟禱告有關連的時刻.
上帝應允了撒迦利亞自己個人的.
很多年前的一個這樣的祈禱.
他就猜顯他的天使加伯列.
告訴你的禱告蒙英文.
不知道弟兄姊妹你的感覺是怎樣.

$^{161}$如果你是撒迦利亞.
你的工作就是將人的禱告帶到上帝面前.
你自己的禱告就在這個禱告的時刻.
知道應允了.
你有什麼感覺.
是不是一個非常開心很興奮的時刻.
所以發覺一個這樣的時刻.
一個非常之刻意去強調.
一個極度跟禱告有關係的人.
和這個禱告的時刻.
時地人都是非常吻合他們的禱告蒙英文.
於是天使就一五一十.
就將後來約翰BB出生的事告訴他.
他說你的妻子伊瑟伯就給你新一個兒子.
你要將她的名起做約翰.
你必歡喜快樂.
有許多人因她出生必歡喜喜樂.
她在主面前將要為大.
淡酒濃喝就有多不可.
從無福裡就被聖靈充滿等等.
是一個非常之好的消息.
你將要生一個兒子.
這是一個眾人祝福的禱告的英文.
不過令人震驚的是.
撒迦利亞的回應.
經文一文記載.
路加的第十八節這樣說.
撒迦利亞對天使說.
我憑著什麼可知道這事呢.
我已經老了.
我的妻子也年紀老了.
撒迦利亞似乎不但忘記了.
她曾經這樣祈禱.
明明是上帝告訴她.
她的禱告亡英雲.
但她似乎不太記得自己這樣祈禱.
更加不太相信這個禱告會發生.
甚至說我有這樣祈禱嗎.
我好像完全忘記了.
這樣的事情發生在我身上.

$^{201}$對於這個天使忽然的驚喜.
撒迦利亞的反應是.
我有這樣祈禱嗎.
她完全不記得.
不相信.
不覺得會發生這樣的事情.
所以撒迦利亞似乎忘記了.
N年前曾經有這樣的禱告.
不過神記得.
但她只是忘記了.
不過其實是正常的.
大家都是這樣的.
我們一輩子說了那麼多話.
通常都不太記得自己說了什麼.
通常聽者有心.
說者無意.
無論是鼓勵人一輩子說話.
還是一句傷心的話.
我們說過就不太記得這些話.
聽過就比較容易記得一輩子.
大家也是一樣.
大家的禱告都是一樣.
大家記不記得你今早祈禱了什麼禱.
你今早的早餐是怎樣祈禱.
或者你昨晚睡前的祈禱.
如果有祈禱的話.
是祈禱什麼.
記不記得.
你昨晚.
你和我聊天.
分享一下你昨晚祈禱什麼禱.
你記得你昨晚祈禱什麼禱嗎.
可能都不記得.
大概方向都記得.
但什麼都忘記了.
一輩子祈禱那麼多.
哪會那麼多記得呢.
所以似乎又很合理.
十年前你祈禱都不記得祈禱什麼禱.
所以我明白.

$^{241}$為什麼我們發現.
似乎都是這樣.
撒迦利亞似乎都忘記了.
他自己一輩子祈禱那麼多禱.
曾經其中一個這樣祈禱的祈禱過.
而上帝又聽了.
又應允他.
但撒迦利亞自己作為一個.
專業的祈禱代禱者.
他自己似乎不太記得.
曾經有這樣說過.
所以似乎是這樣.
所以天使就這樣說.
到了時候.
這些話必然應允.
你這個禱告仍然是會蒙英文的.
上帝不會落空的.
不過因為你不相信那個禱告.
你就將會成為啞巴.
不能說話.
直到這些事情實現的日子.
我們問對於撒迦利亞的回應.
似乎伽巴烈的回應.
是出乎我們意料之外.
天使說似乎要去懲罰他.
要他成為一個啞巴.
不能說話直到實現的日子.
天使是要去懲罰撒迦利亞.
成為一個啞巴.
日子不是太長.
但也不是太短.
十月懷胎這麼久.
他就要啞巴不能說話.
說了多久.
十個月不說話.
其實也不是很長.
但也不是很短的日子.
你想想作為丈夫.
你太太懷孕.
你不能說話.

$^{281}$也挺麻煩的.
到時接電話也接不到.
生孩子也不能說話.
也挺麻煩的.
所以他的太太十個月懷胎.
他就要成為啞巴十個月的時間.
也是一個不少的懲罰.
我們問為何天使要這麼嚴厲去懲罰撒迦利亞呢?.
為何會這樣呢?.
其實我們在聖經裡也見到不少生子不順的故事.
例如舊約裡有些什麼故事呢?.
很多這樣的故事.
生不了孩子,聖帝就要孕育他們.
其中一個就是哈拿.
哈拿也是生不了孩子.
他祈禱耶和華,他許願說.
萬貫之和說:你若顧念婢女的苦情,顧念不忘婢女.
賜我一個兒子,我必使他終生歸於耶和華.
不用剃刀剃他的頭.
這個就是後來的薩姆爾的故事.
亞伯拉罕和撒拉也是這樣.
如果你熟悉舊約的話.
亞伯拉罕和撒拉也是這樣.
撒拉的月經已經斷絕了.
撒拉就心裡暗笑說.
我記爾雖敗,我主也老邁,豈能有這喜事呢?.
對亞伯拉罕說,撒拉為什麼暗笑?.
想想,撒拉是自己偷笑了.
說:我記爾年老果能生育嗎?.
耶和華豈有難成的事呢?.
到了日期,明年之時,我必回到這裡.
撒拉必生一個兒子.
撒拉就害怕,我暗笑被上帝聽到.
不承認,說:我沒有笑.
耶和華說:不然,你實在笑了.
很嚴重的,你偷笑上帝.
你笑我,撒拉沒有笑,你笑我.
撒拉也沒有被耶和華懲罰.
雖然撒拉笑了,但也沒有被上帝懲罰.
撒拉不相信.

$^{321}$為什麼撒拉這次來到撒迦利亞.
作為一個丈夫.
不記得自己曾經祈禱.
就被天使懲罰.
為什麼天使會這麼玻璃心呢?.
以前撒拉也沒有懲罰他.
為什麼這次會這麼玻璃心懲罰他呢?.
我認為其中一個可能的答案是.
這可能是對專業幫人祈禱的人的一個很大的諷刺.
一個專業地幫人禱告的祭司.
將人的禱告獻給上帝的一個祭司.
他自己的禱告也不太相信,甚至不記得.
可能這是一個原因.
在牧師傳娘的位置是比較嚴重.
自己禱告不相信可能會比較嚴重.
所以一個經常專業幫人禱告的人.
他自己連禱告也不相信.
這可能是一個很大的諷刺.
所以天使要去警戒他.
要他好好地去重視自己的禱告.
實際上十個月不能夠祈禱是一個很有意義的事情.
大家有沒有試過?當然大家不要試.
十個月不祈禱.
但十個月不祈禱是一個很大的體驗.
你說我說不出話來就要密討.
以前不是很流行密討.
在舊時代的人是怎樣祈禱的?.
是張開手,開眼,大聲地禱告.
很少密討.
所以十個月不能祈禱.
似乎是一個很大的懲罰.
也是一個很大的禱告反省.
你既然這麼不珍惜你的禱告.
我就要你十個月不能去禱告.
讓你嘗試一下不能開口禱告的滋味.
所以這是一個很特別的十個月的屬靈體驗.
當然大家不要試.
不要試十個月期不祈禱.
最後經文的結束.
五八七節就關掉了.

$^{361}$伊利沙伯就產其道了.
生了一個兒子.
倫舍親屬聽見主向他大師憐憫.
就和他一同歡樂.
到了第八天.
他們來給他孩子行國禮.
並要他父親名字叫撒加尼亞.
他母親就回應說不要叫他約翰.
他們對他說.
你身上中末有名字的.
他們就向他打手勢.
問他自己的孩子叫什麼名字.
於是撒加尼亞就叫人家寫板.
因為說不出話來.
寫到他的名字是約翰.
他們都驚訝.
突然撒加尼亞就開口.
立刻開了.
嘟嘟嘟.
舌頭就鬆了.
就開始說話稱頌神.
撒加尼亞十個月不能說話.
不能祈禱.
當他儲足十個月不能祈禱.
一開口就說什麼.
就是一個Hallelujah的.
頌讚的禱告.
所以這個撒加尼亞十個月不能祈禱.
一開口就是一句讚美上帝的禱告.
是一個很有意思很特別的故事.
這個就是我們今天去看經文.
是一個很有趣的.
是一個有關禱告.
一個專業的禱告者.
自己忘記禱告.
十個月不能祈禱.
一開口就是禱告這樣的故事.
這個禱告故事和我們有什麼關係呢.
今天我們就和大家思想一下.
這個禱告和我們之間的關係.

$^{401}$我認為撒加尼亞最大的問題是什麼呢.
其實就是他太過熟悉禱告.
他是一個每天都在禱告的人.
每天都把禱告獻給上帝.
一個祭司的身份.
有時候我們太過熟悉.
反而令我們越來越不覺得是一回事.
以致他的生命和自己的禱告.
是有一點點脫鉤.
今天我不知道大家怎麼祈禱.
今天大家可能基本上都經常祈禱.
特別是在WhatsApp裡.
在群組裡大家有什麼事.
大家怎麼祈禱.
大家睡覺的時候就按按按按.
把手伸出來.
按幾個都可以.
不知道你們有沒有祈禱.
你們按完之後有沒有祈禱.
還是按了就當祈禱了.
這個科技很大.
我都很難去反思.
我們按他的手是否真的有祈禱.
我經常都這樣.
在教會裡看到頂尖會說什麼.
我說了可能是為你祈禱.
我就很怕自己說完就忘記祈禱.
所以我經常都會一轉頭先補好.
主要是問他.
起碼先祈禱.
不要忘記.
不要說完之後算了.
今天我們很多時候.
禱告成為了一件.
怎麼說呢.
一件很熟悉的.
是一件很低成本的事情.
為你祈禱是一件非常好的事情.
但是一件你不需要.
投入任何努力的事情.

$^{441}$為你祈禱.
不知道有沒有祈禱.
不過都做了.
事實上為你禱告是一件很好的事情.
但也是一件不是最低成本的事情.
你想想.
教會當中有個頂尖戰輩跟你說.
他患了癌症.
那你會怎麼做呢.
你就為他祈禱.
為你祈禱.
但你想想.
如果是你爸爸媽媽.
或是你子女患癌症.
那你會怎麼做呢.
你就不只是為他祈禱.
你會付出很多的心力.
很多的時間.
去跟他同行.
籌算很多的東西.
所以為你祈禱是很好的.
但本來是令到我們的生命和禱告脫勾了.
我們的生命跟不上我們禱告中的自己.
特別是信主的人.
可能是信主三十年以上的那些.
你做主席.
做團契.
做很多的時候.
你經常都會祈禱.
我們禱告到可以天花亂墜.
在黑門山滴下來的滋油.
很多這些很美好的禱文.
說出來.
說得很美.
但我們有沒有很認真去看待.
我們每一個的禱告.
甚至我們每一句的禱告.
這個不知道是否這麼大的經驗.
有時候你禱告.
有時候你不知道自己在說什麼.

$^{481}$很多這些很屬靈的魔鬼.
說出來可能是為了買時間.
不知道有沒有感覺.
說出來之後我不知道下一句說什麼.
就說出來.
一邊想一邊說出來.
有些人笑.
你明白我說什麼嗎.
有時候禱告是很多東西.
但其實那些說話都未必是這麼多.
出於你自己的內心.
只不過是一些禱告的說話.
所以我們很好的去反省.
撒迦利亞就是這樣的一個崗位.
一個專業的祭司.
一個專門把弟兄姊妹的禱告獻給神的人.
他自己的禱告都不記得.
甚至不相信.
問題不是不相信.
而是他太過熟悉禱告.
禱告成為了一個他經常做.
但不是很重視的事情.
今天我們去反省我們的禱告.
是我們怎樣祈禱.
我經常說我們什麼人就祈什麼禱.
你是一個怎樣的人.
就一五一十地在神面前.
將你自己呈現出來.
禱告不能騙人.
因為神知道你是怎樣.
都不用怎麼雕花.
將你最真誠的一面獻上給神.
第二個值得我們反省的地方.
好好珍惜我們每一次禱告的機會.
我們看看撒迦利亞的故事.
當他十個月不能祈禱的時候.
原來是很饑渴的.
十個月不能祈禱.
一祈就湧流爆出來讚美上帝.
我們很多時候發覺原來可以禱告.

$^{521}$我們不要以為是必然.
不知道大家是不是這樣.
我們小時候團契.
我們團契禱告是怎樣祈禱的.
要猜贏了就要祈禱.
團契裡的禱告.
猜贏了就要祈禱.
好像很不想這樣.
圍著吃飯的時候.
大家猜就好像祭籤出來.
輸了就要祈禱.
我們很少這樣.
好 我來吧.
好像不太想這樣.
我自己是傳道人.
我自己讀神學.
或拍拖的時候.
我還沒結婚的時候.
去我岳父家吃飯.
家裡就是這樣.
那時候我還沒去岳父家.
現在我去岳父家.
吃飯的時候.
他叫我John仔.
John仔 祈禱.
每次吃飯的時候.
他都叫我祈禱.
十年二十年都是這樣.
每次都是我祈禱.
由拍拖到結婚.
現在都是我祈禱.
吃飯的時候.
又是我.
可不可以叫別人.
你明白我的感覺嗎.
但我發覺.
我自己覺得有一點認罪.
不行的.
做傳道人沒理由.
要有那種.

$^{561}$好像很不想這樣的感覺.
有人叫你祈禱.
應該怎樣做.
應該是能碰到他.
我就祈禱.
真人.
每次我們被邀請去祈禱的時候.
無論是團契.
或者任何機會.
去領土吃飯都簡單.
應該是一種榮幸.
一種很快樂的事情.
一個很值得我們珍惜的事情.
因為可以開口.
代表著你的弟兄姊妹.
去獻上一個謝飯禱告也好.
這是一個很值得我們.
來去珍惜和光榮的事情.
不要覺得這是一件推搪的事情.
我覺得很醜怪.
突然間猜輸了那一期.
不知道是什麼事情.
所以大家去好好地.
珍惜每一次禱告的機會.
特別是牧者邀請你去祈禱的時候.
應該覺得是一件很好的機會.
能夠開口.
讚美,感恩,代禱.
好好地去反省我們今天禱告的生命.
我想大家都懂得禱告.
從第一天回教會開始.
你懂得怎樣去祈禱.
但我們今天想祝這麼一年.
我們祈禱太多.
以致我們的禱告.
沒有好好地去和我們的生命相稱.
我自己聽過.
曾經聽過最好的禱告是什麼呢?.
就是我媽媽的祈禱.
那時候我剛回錦繡堂.

$^{601}$上了珠沒多久.
我就開始帶我媽回教會.
後來我媽也上了珠.
就缺了字.
那時候錦宣在天水圍有一個祈禱區.
天水圍區.
那時候星期三晚帶我剛上珠的媽媽回祈禱會.
那時候我媽媽上了珠大概一個月兩個月左右.
在分區裡.
她在家裡分區祈禱.
她不懂得祈禱.
慢慢地去說.
她就為我祈禱.
我記得當時我媽媽為我祈禱.
她就將她平時說的說出來.
就是這樣祈禱.
不要太晚睡.
不要捱壞自己.
早點回來記得喝湯.
她將這些說話說出來成為她的禱告.
就是煮.
叫她不要太晚睡.
叫她不要捱壞自己.
記得喝湯.
但這是我聽過最好的禱告.
因為這個禱告很簡單.
但每一句都是她內心100\%的說話.
不多不少.
每一句都是阿們阿們的禱告.
真的向上帝祈求.
讓自己的兒子不要太晚睡.
我們好好地回到我們禱告的基本.
我們今天想煮這麼久.
回到我們禱告的基本.
禱告就是將你內心的那句說話.
說給上帝聽.
有時我們說得太快.
一句說話太多.
blah blah blah.
這樣說出來.

$^{641}$我們都不知道自己在說什麼.
都忘記自己在說什麼.
我們試試.
不妨我們真的回到我們基本.
我們一起閉上眼.
我們一起向上主禱告.
將你內心的渴求.
將你內心的需要.
將你這一刻的無論是感恩.
或者是認罪.
或者是代求.
我們去向上帝去尋明.
向我們天父尋明.
不需要很多.
將你此刻的心靈.
變成一個禱告.
你自己就是這個禱告.
我們一起去祈禱.
向上帝祈禱.
\newpage



\section{羅馬書 15:22-29}
\label{sec:aPVZvLNc78I}
\textbf{2023-05-27 晚堂《五餅二魚的團契》羅馬書 15\_22-29 講員:陳麗蟬牧師}
\newline
\newline
連結: \href{https://youtube.com/watch?v=aPVZvLNc78I}{\texttt{https://youtube.com/watch?v=aPVZvLNc78I}} ~~~~ 語音日期: 2023-05-30
\newline
\newline
\hyperref[sec:wCpZ46fk4yY]{\small{< < < PREV SERMON < < <}}
~
\hyperref[sec:index_chronic]{\small{[返順時目]}}
~
\hyperref[sec:index_scriptual]{\small{[返順卷目]}}
~
\hyperref[sec:hSecMeCWH4s]{\small{> > > NEXT SERMON > > >}}
\newline
\newline
羅馬書 15:22-29
\newline
\begin{longtable}{cl}
\hline
\hline
章節 & 經文 (和合本修訂版)\\
\hline
15:22 & \begin{tabularx}{0.7\textwidth}{X} 因此我多次被攔阻，不能到你們那裡去。 \end{tabularx} \\ \\ \relax
15:23 & \begin{tabularx}{0.7\textwidth}{X} 但如今，在這一帶再沒有可傳的地方，而且這許多年來，我迫切想去你們那裡， \end{tabularx} \\ \\ \relax
15:24 & \begin{tabularx}{0.7\textwidth}{X} 盼望到西班牙去的時候經過，得見你們，先與你們彼此交往，心裡稍得滿足，然後蒙你們為我送行。 \end{tabularx} \\ \\ \relax
15:25 & \begin{tabularx}{0.7\textwidth}{X} 但如今我要到耶路撒冷去，供應聖徒的需要。 \end{tabularx} \\ \\ \relax
15:26 & \begin{tabularx}{0.7\textwidth}{X} 因為馬其頓和亞該亞人樂意湊出一些捐款給耶路撒冷聖徒中的窮人。 \end{tabularx} \\ \\ \relax
15:27 & \begin{tabularx}{0.7\textwidth}{X} 這固然是他們樂意的，其實也算是所欠的債；因為外邦人既然分享了他們靈性上的好處，就當把肉體上的需用供給他們。 \end{tabularx} \\ \\ \relax
15:28 & \begin{tabularx}{0.7\textwidth}{X} 等我辦完了這事，把這筆捐款交付給他們，我就要路過你們那裡，到西班牙去。 \end{tabularx} \\ \\ \relax
15:29 & \begin{tabularx}{0.7\textwidth}{X} 我也知道去你們那裡的時候，我將帶著基督豐盛的恩典去。 \end{tabularx} \\ \\
[1ex]
\hline
\hline
\end{longtable}
$^{1}$各位姐妹平安.
今天很高興可以跟大家一起用聖經的說話.
來彼此勉勵一下.
今天是5月27號 明天是5月28號.
明天是什麼日子?.
不知道吧.
明天是耶穌基督復活之後的50天.
所以在聖經裡面來說是一個很重大的日子.
明天是聖靈降臨的日子.
所以剛才在詩歌和體驗經文裡面.
特別是講到關於聖靈的工作.
在二千多年前.
就是耶穌基督復活了50天之後.
這個聖靈就像風,像火一樣降臨到地上.
祂就要賜給信徒有這個力量.
祂要來更新信徒的生命.
在利未記第26章是這麼說的.
「我要安示降雨給你們.
使地長出農作物.
田地的樹就會結出果實」.
那些乾旱的田地.
一定需要有雨水來灌溉,來滋潤.
那些果樹才能夠結出甜美的果實.
同樣地 人的心靈,人的心田.
需要聖靈的灌溉.
然後才能夠生出生命一個很美的果子.
《以弗所書》就這麼說.
「從前你們是暗昧的」.
講到未信之前.
「但如今在主裡面是光明的.
行事為人就要像光明的子女.
光明所結的果子.
就是一切的良善,公義和誠實」.
要經歷這個聖靈的風.
好像火.
不是 要經歷的聖靈好像火,好像風.
這樣來更新我們,滋潤我們.
除了要聖靈來充滿我們之外.
也要藉著弟兄姊妹的相交來經歷.
大約五年之前.

$^{41}$我跟印尼姐姐回到她的鄉下.
認識一下她的家人.
和認識一下她的教會.
因為印尼是這樣的.
他們要整條村落都是信耶穌.
大家可以互相保護.
不然周圍的回教徒就會來攻擊.
所以整條村落都是信耶穌的.
有一天我就去到她耕田的地方.
去到的時候我聽到有幾個人嘩嘩聲.
然後又看到前面有個人追追追.
原來他們很開心.
有個人從香港而來.
又去到他們耕田的地方.
那個很高的弟兄.
他就爬了上一棵椰子樹裡面.
摘了幾個未成熟的椰子.
就是椰青.
就扔了在地上.
特意扔在灌水道那裡.
然後這個很矮的弟兄就去追追追.
追回那幾個椰青來.
說第一次喝新鮮從樹上摘來的椰青.
很清甜很好喝的.
其實他們家庭很窮困的.
剛才這兩位弟兄身體都有缺憾的.
高的那個是單眼的.
他的眼睛看不到東西的.
矮的那個就有癌症.
在兩年多之前他都回天父家了.
他們很開心.
就覺得有人肯去到他們這麼落後的地方.
還去到他們那個田地.
覺得田地骯髒骯髒的.
還有沒有椅子的.
他們那裡沒有椅子的.
他們覺得很開心.
為了歡迎我去.
他們就去分享.
他們在田地裡最好的東西.

$^{81}$就是樹上高的那個椰青.
所以他們就爬了上去摘下來.
他們跟我去分享.
分享在環境裡最好的東西.
在五月份裡面.
所有教會童工.
本堂教會童工講到的主題.
都是「團契」.
所以你們都會聽到.
之前的童工都講了一個希臘文的字.
就是這個「Kononia」.
是一個「團契」的意思.
不過這個字的主要意思就是.
和其他人去分享.
一起合作去分享.
原來這個就是一個「團契」.
我們看看羅馬書的背景.
保羅在哥林多的時候.
他就寫了這個羅馬書.
他在羅馬書第一章.
或者到最後那一章.
有很多問候的說話.
他問候很多很多很多的人.
並且好像跟他們很熟悉.
但保羅其實沒有去過羅馬教會的.
這些人可能是他的親屬.
可能是他的朋友.
又可能是他的一些童工.
不過這些人全部都是在保羅.
他去到馬其頓或者亞細亞那一帶地方的時候.
在傳道裡面認識的.
在裡面很寶貝的淑齡分享.
後來這幫人就搬到了羅馬.
保羅很想跟他們能夠再見面.
在這封書信裡面他都再一次講.
很想去到羅馬教會.
跟他們在裡面有彼此的相交.
這個團契.
在我們教會裡面有四十多個團契的.
如果沒有記錯的話.

$^{121}$團契的基礎是什麼呢?.
就是我們跟我們的主耶穌基督一體的工作.
在哥林多前書第一章第九節.
他這樣講.
上帝是信實的.
祂呼召你們.
就是上帝呼召我們.
好與他的兒子.
就是耶穌基督.
共享團契.
上帝呼召了我們.
成為了他的子民.
是做什麼呢?.
跟主耶穌基督一起有這個團契.
有一個緊密的關係.
互相的分享.
在《約翰福音》14章20節.
耶穌自己這樣講.
到那一天你們就會知道.
我在父裡面.
你們在我裡面.
我也在你們裡面.
我很怕讀這段經文.
很多的裡面都不知道誰在誰裡面.
第一個.
我在父裡面.
這是耶穌講的說話.
耶穌說.
我在父的裡面.
即是說.
祂與父上有一個很緊密的關係.
主耶穌基督在地上所作的一切.
都是按照著天父的心意去做的.
第二個.
第一個就是記住.
耶穌基督與父上的關係.
第二個.
你們在我裡面.
就指那些門徒.
即是今天的你和我.

$^{161}$我們要在耶穌基督的裡面.
又即是說.
我們要與主耶穌基督有一個很密切的關係.
一個愛的關係.
我們在地上的行事.
我們的價值觀.
我們的準則.
都要以耶穌基督的心意為主.
所以唯有這樣.
我們在主耶穌基督的裡面.
唯有我們這樣做的時候.
就可以出到第三個了.
我也在你們裡面.
主耶穌就會在我們裡面.
我們在主耶穌裡面.
跟著祂的旨意來走.
主耶穌基督就會在我們裡面.
意思就是說.
耶穌在我們裡面.
就是說.
在我們遇到什麼事情都好.
祂會給力量我們.
祂會保護著我們.
去與我們緊密的.
緊緊的相連.
所以這個就是一個愛的團契來的.
與耶穌基督一起.
我們基督教是一個群體的信仰來的.
所以不可以做一個孤單人.
自己一個人的.
用一個在愛心生活裡面來實行出來的.
不同我們去黃大仙拜神.
黃大仙拜神可以不認識其他的人的.
特別在過新年的時候.
那些人爭著去黃大仙廟上頭炷香的嘛.
去爭這個上頭炷香.
你可以不用理其他人的.
為了自己快快的.
裝完香之後就可以早點回家睡覺.
你可以做什麼都可以.

$^{201}$鬆綁 踩人家的腳.
以致自己快快碎碎的搞定它.
就可以回家了.
你可以不用理其他人.
也都不需要給面子給其他人.
但是在耶穌基督的團契裡面.
在教會裡面.
在教會的團契裡面.
或者教會和教會之間.
或者整個屬於上帝的教會裡面.
是一個實踐耶穌基督愛心的地方.
我們就看看這個怎麼樣實踐.
剛才讀到這個《羅馬書》15章的時候.
就講到這個保羅已經完成了.
大概完成了.
在阿哥亞和馬其頓的宣教工作.
在這幅圖裡面.
在這個紅色的橢圓形.
就是耶路撒冷.
去到中間的那個大的藍色的框框.
那裡就是這個馬其頓和阿哥亞.
就是保羅一次,兩次.
三次去到那裡來宣教.
一直都很掛念著這個羅馬.
這個羅馬就是那個綠色的橢圓形.
在中間的地方.
保羅的計劃.
因為他就在馬其頓那裡.
因為哥林多城都在那裡了.
馬其頓和阿哥亞那裡.
他就在那裡.
想著都存了一些地方了.
就完成了這裡的部分的工作了.
但是他又收到他們一部分的捐獻.
給耶路撒冷教會的.
因為當時耶路撒冷教會比較貧困.
因為一方面就是.
在裡面的一些猶太教的信徒.
就不值得他們信耶穌.
又發展得這麼快.

$^{241}$給他們很多的逼迫.
而另一方面.
也在耶路撒冷當時都有這個饑荒.
所以裡面都有很多窮困的信徒.
這個馬其頓和阿哥亞的教會.
他們籌了一些錢.
就托保羅拿去給耶路撒冷的教會.
保羅就想著.
拿著這筆錢就回去耶路撒冷.
接著他又會再起行的了.
他的目的地是哪裡呢?.
就去到另一方面.
那裡就叫做西班牙.
其實就是西班牙.
當時他們認為.
這個就是地極的地方.
最遠的了.
那裡福音還沒有到.
所以保羅就很想去西班牙那裡宣教.
他就想著從耶路撒冷開始起行.
他就會路經羅馬.
在羅馬的時候.
又會再和那些信徒.
一起在靈裡面彼此鼓勵.
才再去西班牙.
也盼望在羅馬停頓的時候.
也都看看會不會籌一些金錢.
一些實際的需要.
以至成為了去西班牙裡面.
一些的需用的經費.
在羅馬書15章27節.
剛才我們讀了的了.
他就這樣說的.
他說這些人 指的外邦人.
固然是樂意的.
就說到剛才馬其頓和阿蓋亞.
他們很樂意的.
其實也算是所欠的債.
因為外邦人既然分享了他們靈性上的好處.
就當把肉體上的需用供給他們.

$^{281}$說到這幫外邦人.
即是說馬其頓和阿蓋亞那幫人.
他們在靈性上得到這個好處.
什麼叫靈性上得到好處呢?.
就是得到這個福音.
所以現在耶路撒冷教會窮困.
你們都應該去供給他們.
這個「供給」這個字.
在希臘文原來又是團契這個字.
Condonia 又是那個字 供給.
《羅馬書》另外一部份.
在12章13節.
「聖徒缺乏要幫補 客要一味地款待」.
不是請人吃飯得一味送.
不是.
就要專心努力地去接待他.
款待這個字.
又是同一個字根.
是團契有緊密的關係.
保羅就說.
這個馬其頓和阿蓋亞教會.
很樂意去幫助耶路撒冷教會的人.
那馬其頓和阿蓋亞教會.
是不是很有錢呢?.
在《哥林多後書》他又這樣介紹.
兄們 我要把上帝賜給馬其頓宗教會的恩惠告訴你們.
告訴羅馬教會的人知道.
馬其頓教會的人怎麼做.
他們在患難中 受大考驗的時候.
仍然沒有喜樂.
原來馬其頓教會和阿蓋亞教會.
一樣受到很多的逼迫.
但不過在這個逼迫考驗當中.
他們仍然感到滿足.
在另一個裡面有喜樂 有平安.
他說 在極度貧困中.
還額外顯出他們樂觀的慷慨.
雖然他們很困難.
但仍然很樂意地和耶路撒冷教會的信徒.
分享他們的金錢.

$^{321}$他繼續說 我可以證明他們是按著能力.
每一個信徒都是按著自己的能力.
而且超過了能力來捐助.
他看自己有多少錢.
他可能捐了30元 裡面只捐了3元.
不過他們竟然有大膽的走出多一步.
超過了能力來捐助.
可能捐了5元.
他還主動地求保祿.
求當時他們教會.
准許他們在供給耶路撒冷教會貧困的人的善事上.
有份.
原來有份這個字.
又是團契.
那個字根又是團契這個字.
馬其頓教會 阿哥亞教會.
他們自己在一個很窮困的境況.
自己都不足夠.
但仍然樂意這樣去奉獻.
不單止按著自己的力量.
還要超過了他的力量來幫助其他的人.
我們看到 剛才我們只是讀了1,2,3.
三段的經文.
我們看到這個團契字根有關的字.
有這個「供給」,「幫補」,「有份」.
全部都是一些行動的字.
是幫人的字.
在一個很窮困的山區裡面.
有一個A信徒叫許赫.
他有一天在牧師那裡拿了一大包米.
抬著那包米走回家.
從牧師的家走到他家.
要走三個小時.
是一些山路來的 走三個小時.
回到家裡跟家人說.
我們終於有米吃了.
然後他就拿了這袋米裡面.
少許煲了一些稀粥.
一家十口就有稀粥吃.
看完這個A信徒的家.

$^{361}$我們又看看牧師的家.
在牧師的家裡面.
原來家裡只剩下兩匙米.
只剩下兩匙米.
我們就叫師母煮了它.
吃了它 煮了它.
我們吃飯吧.
當借飯祈禱的時候.
這個牧師說.
他說多謝天父賜給我們飲食.
叫A信徒的一家和我們一家都有足夠.
這個A信徒拿了牧師那包米.
他們十幾個人.
那包米可能吃了他們十幾天.
或者兩個星期.
他們很足夠.
但我們看回牧師的家.
他們只剩下兩匙米.
他們那天煲了飯吃.
所以明天不知道有什麼吃.
但這個就是牧師的信心.
他覺得我今天吃.
有得吃已經是很足夠的了.
明天會怎樣呢?.
他對上帝有信心.
其實就好像我們經常說.
當我們見到一杯水.
只有半杯水的時候.
有時候我們會覺得很擔憂.
糟了 只剩下半杯水.
喝完就沒有了.
但是呢.
但是有些人就看到.
所以我們就要怎樣看每一件事情.
當我們那時候在神樂院讀書的時候.
糟了 還有一個星期就要交功課.
又要考試了.
但有些同學就會說.
哇 還有一個星期可以做功課.
可以去讀書.

$^{401}$才到期限交功課.
才可以去考試.
還有時間 還有時間.
所以我們就是要學習.
當我們有的時候.
我們不是看著那半杯水.
沒有那些東西來優秀.
而是我們為了有的.
我們怎樣去分享.
為了有的時候.
我們要去獻上感恩.
我們繼續看羅馬書.
「這固然是他們樂意的.
其實也算是所欠的債」.
那這個馬其頓和阿哥亞欠了什麼債呢?.
欠了耶路撒冷教會什麼債呢?.
就是這個福音的債.
在耶路撒冷教會.
他們受到逼迫.
就分散去不同的地方.
去到馬其頓 去到阿哥亞.
也有些像羅馬他們那樣.
組成宣教隊.
去到那裡的地方去宣教.
所以這個馬其頓教會.
他欠了耶路撒冷.
不是 他不是欠了他們.
對不起 應該從頭再來說.
他們就是這個福音的債.
就是當他從耶路撒冷教會那邊的人.
拿了這個福音.
認識了這個福音.
相信了耶穌之後.
他就有責任將這個福音.
繼續傳開去.
但第二方面.
這個馬其頓教會.
他要做的就是.
分享在物質上的善事.
就是他在基督教群體裡面.

$^{441}$他就要學習.
看看周圍的人.
他知道他們有需要的時候.
他就有責任.
去跟其他人.
分享在物質上的擁有.
無論是福音的債.
或者分享物質的事情上.
都是我們一班信徒.
所應該要做的.
供給 款待 有份.
都是一些很實際的行動.
信徒要來去做的.
就是表示我們要做一個僕人.
要做基督的僕人.
去服侍上帝.
服侍我們周圍的鄰舍.
團契就是一個愛的地方.
是一個實踐愛的地方.
所以無論我們有錢也好.
沒錢也好.
出多也好 出少也好.
在這個團契裡面.
我們都要學習彼此的分享.
分擔.
而不是坐著.
等別人來服侍.
我們不能夠坐著.
我只是回團契.
我只是聽別人分享見證.
我不會說的.
我沒有話好說.
我回來團契.
我不會做任何的事情.
我也不會去幫忙任何的事情.
團契就是一個愛的.
實踐的地方.
近日有位姿體.
她的丈夫癌症.
最近醫生告訴她.

$^{481}$已經擴散了.
也告訴她.
已經沒有適合的藥物.
可以幫助她丈夫.
大家都很慌張.
在那段時間.
她的印尼姐姐確診了.
她覺得簡直好像.
天塌下來.
我們中國人說什麼.
屋漏兼逢夜雨.
很悽慘.
要遷就一些小提琴聲.
很慘.
有位姊妹.
知道這個情況的時候.
義不容辭.
主動地說.
由我每天幫你買菜.
買一些東西.
去送到她家裡.
很感動.
這位姊妹.
就在實踐一個團契.
一個愛的分享.
很多弟兄姊妹都聽過.
這個「五餅二魚」.
在聖經四卷福音書.
都有記載耶穌基督的神蹟.
當日耶穌基督.
聽到他的表哥施洗若翰.
被殺.
被斬頭.
很難過.
本來他想去到山上高.
安靜.
但那些人不知道耶穌在哪裡.
五千多人.
只是男子.
只是男人.

$^{521}$跟著來.
所以耶穌基督的心情雖然很沉重.
他仍然要趕到黃昏.
去到黃昏的時候.
門徒就告訴耶穌.
「天黑了」.
「不如你叫那些人群回家吃飯休息吧」.
耶穌就說了一句什麼話.
「不用叫」.
「你幫他們找吃的吧」.
我想那些門徒一聽到.
嚇得飛起來.
五千多人.
去哪裡找這麼多東西吃呢.
去哪裡有這麼多錢呢.
但就在這個時間.
一個小孩.
一個很單純.
很慷慨的小孩出現了.
他就將身上所帶著的五個餅.
和兩條魚拿出來問耶穌.
「在這裡嗎?夠不夠啊?」.
耶穌說「夠了」.
「你已經擺上了所有了」.
結果那個神蹟.
餵飽了不單止那五千人.
還有很多婦女,小孩.
還有剩下十二個男子的食物.
大家都很熟悉這個神蹟.
我想你和我都絕對相信.
在這五千多人裡面.
絕對不會只有這個小孩.
有食物在身的.
不過就只有這個小孩.
肯將他所擁有的食物.
拿出來給主耶穌基督.
也就是因為他肯拿出.
他自己的五個餅,兩條魚.
奉獻出來的時候.
耶穌基督就藉著.

$^{561}$他所奉獻的出來.
行了這個很偉大的神蹟.
所以弟兄姊妹.
我們每一個人的手裡.
都有五個餅,兩條魚.
有些人有十個餅,十條魚.
有些人只有兩個餅,一條魚.
但不過無論我們的能力.
我們的時間,我們的智慧.
我們的恩賜才幹有多少都好.
我們都要盡力的去服侍上帝.
服侍人.
上帝會透過我們.
小小的.
我們覺得自己微不足道.
覺得自己說出來都醜死人的恩賜.
上帝可以藉著我們這個小小的奉獻.
來行出這個偉大無比.
祝福很多很多人的神蹟.
這個團契這個字.
Colonia.
就是一個愛的分享.
我們先和耶穌基督有一個緊密的關係.
然後我們就懂得去愛其他的人.
在見證上給予彼此的分享,鼓勵.
在密切上給予彼此的分享,分擔.
希伯來信說到.
我們要彼此相顧,激發愛心.
勉勵行事.
不可停止聚會.
好在那些停止慣了的人.
都要彼此勸勉.
希伯來書的作者.
他就鼓勵當時的信徒.
不要因為逼迫.
來輕易放棄教會裡面的聚會.
放棄和弟兄姊妹的相交.
同樣在今天.
我們彼此勉勵弟兄姊妹.
我們千萬不要因著任何的事情.

$^{601}$因著任何的人際關係.
以致我們從主的群體抽離了出來.
我們一定要在基督裡面.
活出這個團契的精義.
彼此分享生命.
在密切上我們都彼此分擔.
彼此的支持.
這是上帝所喜悅的.
我們就唱一首詩歌.
分享主愛.
他說我們要分享主的愛.
我們將從上帝所領受到的愛.
我們無論經歷.
我們的物質和周圍的人分享.
(曲:張碧晨).
讓我共分享主恩愛.
讓主愛充滿生命.
願主愛心救眾教會.
使我會心願服你.
求主更深深顯耳熟.
使我會真摯相望.
讓我共分享主恩愛.
誠實相愛.
(再來一次).
讓我共分享主恩愛.
讓主愛充滿生命.
願主愛心救眾教會.
使我會心願服你.
求主更深深顯耳熟.
使我會真摯相望.
讓我共分享主恩愛.
誠實相愛.
\newpage



\section{使徒行傳 2:1-21}
\label{sec:hSecMeCWH4s}
\textbf{2023-05-28 加恩樓《聖靈充滿的教會》使徒行傳 2\_1-21 講員:曾浩賢院長}
\newline
\newline
連結: \href{https://youtube.com/watch?v=hSecMeCWH4s}{\texttt{https://youtube.com/watch?v=hSecMeCWH4s}} ~~~~ 語音日期: 2023-05-30
\newline
\newline
\hyperref[sec:aPVZvLNc78I]{\small{< < < PREV SERMON < < <}}
~
\hyperref[sec:index_chronic]{\small{[返順時目]}}
~
\hyperref[sec:index_scriptual]{\small{[返順卷目]}}
~
\hyperref[sec:H85tdsHmqKI]{\small{> > > NEXT SERMON > > >}}
\newline
\newline
使徒行傳 2:1-21
\newline
\begin{longtable}{cl}
\hline
\hline
章節 & 經文 (和合本修訂版)\\
\hline
2:1 & \begin{tabularx}{0.7\textwidth}{X} 五旬節那日到了，他們全都聚集在一起。 \end{tabularx} \\ \\ \relax
2:2 & \begin{tabularx}{0.7\textwidth}{X} 忽然，有響聲從天上下來，好像一陣大風吹過，充滿了他們所坐的整座屋子； \end{tabularx} \\ \\ \relax
2:3 & \begin{tabularx}{0.7\textwidth}{X} 又有舌頭如火焰向他們顯現，分開落在他們每個人身上。 \end{tabularx} \\ \\ \relax
2:4 & \begin{tabularx}{0.7\textwidth}{X} 他們都被聖靈充滿，就按著聖靈所賜的口才說起別國的話來。 \end{tabularx} \\ \\ \relax
2:5 & \begin{tabularx}{0.7\textwidth}{X} 那時，有從天下各國來的虔誠的猶太人，住在耶路撒冷。 \end{tabularx} \\ \\ \relax
2:6 & \begin{tabularx}{0.7\textwidth}{X} 這聲音一響，許多人都來聚集，各人因為聽見門徒用他們各自的鄉談說話，就甚納悶， \end{tabularx} \\ \\ \relax
2:7 & \begin{tabularx}{0.7\textwidth}{X} 都詫異驚奇說：「看哪，這些說話的不都是加利利人嗎？ \end{tabularx} \\ \\ \relax
2:8 & \begin{tabularx}{0.7\textwidth}{X} 我們每個人怎麼聽見他們說我們生來所用的鄉談呢？ \end{tabularx} \\ \\ \relax
2:9 & \begin{tabularx}{0.7\textwidth}{X} 我們帕提亞人、瑪代人、以攔人，和住在美索不達米亞、猶太、加帕多家、本都、亞細亞、 \end{tabularx} \\ \\ \relax
2:10 & \begin{tabularx}{0.7\textwidth}{X} 弗呂家、旁非利亞、埃及的人，並靠近古利奈的利比亞一帶地方的人，僑居的羅馬人， \end{tabularx} \\ \\ \relax
2:11 & \begin{tabularx}{0.7\textwidth}{X} 包括猶太人和皈依猶太教的人， 克里特人和阿拉伯人，都聽見他們用我們的鄉談講論神的大作為。」 \end{tabularx} \\ \\ \relax
2:12 & \begin{tabularx}{0.7\textwidth}{X} 眾人就都驚奇困惑，彼此說：「這是甚麼意思呢？」 \end{tabularx} \\ \\ \relax
2:13 & \begin{tabularx}{0.7\textwidth}{X} 還有人譏誚，說：「他們是灌滿了新酒吧！」 \end{tabularx} \\ \\ \relax
2:14 & \begin{tabularx}{0.7\textwidth}{X} 彼得和十一個使徒站起來，他就高聲向眾人說：「猶太人和所有住在耶路撒冷的人哪，這件事你們要知道，要側耳聽我的話。 \end{tabularx} \\ \\ \relax
2:15 & \begin{tabularx}{0.7\textwidth}{X} 這些人並不像你們所想的喝醉了，因為現在才早晨九點鐘。 \end{tabularx} \\ \\ \relax
2:16 & \begin{tabularx}{0.7\textwidth}{X} 這正是藉著先知約珥所說的： \end{tabularx} \\ \\ \relax
2:17 & \begin{tabularx}{0.7\textwidth}{X} 『神說：在末後的日子，我要將我的靈澆灌凡血肉之軀的。你們的兒女要說預言；你們的少年要見異象；你們的老人要做異夢。 \end{tabularx} \\ \\ \relax
2:18 & \begin{tabularx}{0.7\textwidth}{X} 在那些日子，我要把我的靈澆灌，甚至給我的僕人和婢女，他們要說預言。 \end{tabularx} \\ \\ \relax
2:19 & \begin{tabularx}{0.7\textwidth}{X} 在天上，我要顯出奇事，在地下，我要顯出神蹟，有血，有火，有煙霧。 \end{tabularx} \\ \\ \relax
2:20 & \begin{tabularx}{0.7\textwidth}{X} 太陽要變為黑暗，月亮要變為血，這都在主大而光榮的日子未到以前。 \end{tabularx} \\ \\ \relax
2:21 & \begin{tabularx}{0.7\textwidth}{X} 那時，凡求告主名的都必得救。』 \end{tabularx} \\ \\
[1ex]
\hline
\hline
\end{longtable}
$^{1}$各位弟兄姊妹早安.
午安,原來過了12點.
多謝上帝,亦多謝大家的邀請.
讓我今日有機會第一次來到大家當中一起崇拜.
亦都做港獨服侍.
剛才稍稍,你們叫富哥是嗎?.
Ricky就帶我走了一圈.
這個教會很特別,很多thing out of the box的設置.
我都是第一次坐在沙發崇拜.
非常欣賞,因為有時我們真的要用珠胖的智慧去傳揚福音.
在港獨儀仙都容許我代表信義忠誠學院一眾師生同工向各位問安.
願因為平安,從我們的父上帝並主耶穌基督.
閨譽參與崇拜的每一位.
剛才敬拜隊都表達今天是五秦節.
或者教會稱為聖靈降臨節.
就是復活主日之後的第五十日.
紀念聖靈的降臨.
所以今日都特別選了一個題目叫聖靈充滿的教會.
亦都選了一段聖靈降臨的經文.
就是剛剛敬拜隊的主令讀了的經文.
就是《使徒行傳》二章一到二十一節.
這段經文其實就是主耶穌在受難和升天的時候.
向門徒所應許的聖靈.
在五秦節那天降臨在門徒身上.
而這段經文之後.
使徒行傳亦都一步一步記載.
這些被聖靈充滿的門徒.
後來被稱為教會的這一班人.
如何將福音從耶路撒冷開始.
一直傳到不同的地方.
當然耶穌所應許的.
被聖靈充滿的這個應許.
並不是單單給第一代的信徒.
或者初期教會的信徒.
主是應許給我們歷世歷代跟隨基督的人.
即是包括你包括我.
不過一個被聖靈充滿的基督徒.
或者一個被聖靈充滿的教會.
是一個怎樣的經驗呢?.
曾經有一位青年人問過我.

$^{41}$因為他有朋友回靈恩派的教會.
都有跟他朋友回他的教會經驗一下.
他就發覺那些聚會有些人會講方言.
有些人會講聖靈激導.
會講神蹟醫治等等.
這些很奇妙的超自然的聖靈工作經驗.
不過在香港其實大部份的教會.
包括我自己所服侍的教會.
都是比較少提及這些很特別的聖靈充滿的經驗.
如果在這些教會的信徒在那裡聚會.
在那裡事奉的信徒.
沒有這些特別的聖靈充滿的經驗.
那麼這個聖靈充滿的體驗究竟是怎樣的一回事呢?.
是不是我們沒有聖靈的同在呢?.
我相信大家都明白不是的.
當我們去信耶穌,當我們去洗禮的時候.
上帝就應許了有聖靈成為我們的引經帶領著我們.
如果是這樣的話,我們都是一班有聖靈的基督徒.
我們都是一個聖靈充滿的教會.
那麼究竟應該的體驗是怎樣的呢?.
或者我們可以透過今天這個《使徒行傳2章》1到21節.
我們去描述一下這份經驗.
那麼繼續講下去,我們先祈求聖靈的幫助好嗎?.
我們一起動心祈禱.
天父上帝多謝你讓我們有機會聚集在這裡一起去敬拜.
一起去聽你的話語.
父啊求你,差遣聖靈臨到我們當中.
使無論孩子口中的言語.
弟兄姊妹的心懷意念都蒙你的閱立和引導.
奉耶穌基督的聖名.
一個聖靈充滿的教會會是一個怎樣的體驗呢?.
今天就著《使徒行傳》和大家分享三樣.
第一,被聖靈充滿的教會.
他們是會得著從上頭而來的能力.
今天的經文第一節開始.
他說五秦節都一到.
他們門徒聚集在一處.
忽然間有響聲從天上下來.
好像一陣大風吹過.
充滿他們所坐的整座屋子.

$^{81}$又有蝕頭如火焰向他們顯現.
分開落在他們每個人身上.
他們都被聖靈充滿.
按著聖靈所賜的口才說起別國的話來.
經文中被聖靈充滿的門徒.
獲賜了說起別國的話來的能力.
他們突然間懂得用不同的香壇.
向帕提亞人,馬代人,以蘭人等等.
十五個地區的人講論上帝的大作為.
而這種聖靈而來的能力.
是一種可以令人感到詫異驚奇的能力.
就好像經文描述.
那些從各地各方來的人.
為著門徒有這個能力而感到詫異驚奇困惑.
當我每一次讀到這段經文的時候.
我心裡都很想上帝賜給我.
這個說起別國的話來的屬靈的恩賜.
如果大家認識信義中神學院.
我們這個神學院和其他神學院有少少不一樣.
我們的老師,同學其實都不只是本地香港人.
我們有來自東南亞國家的.
有來自歐洲的.
有來自美國的.
也有來自內地.
老師群體是這樣.
同學群體都是來自不同的地方.
所以我們神學院的生活和教學都是兩文三語.
我在那裡事奉自然都需要兩文三語.
對於我覺得自己有些語文不是很厲害.
其實有些障礙.
嚴格來說.
有些語文障礙的人來說.
要用兩文三語在那裡生活和事奉.
其實有時都頗挑戰的.
所以更加很想上帝賜我這個說起別國的話來的恩賜.
當然我既然說得很想.
即是還沒有.
還沒有這個恩賜.
我仍然都不是說能夠兩文三語.
這樣非常流利和學院的同工同學一起生活.

$^{121}$不過話說回來.
相信大家都明白.
聖靈所帶給人的能力.
帶給信徒的能力.
不是每個人都說起別國的話來.
聖靈是會給不同的信徒.
有不同的恩賜和能力.
哥林多前書十二章.
他告訴我們.
聖靈是按著他的心意.
賜給不同的基督徒.
有不同的恩賜和能力.
經文說.
恩賜有很多種.
不過是同一位聖靈所賜.
事奉有很多種.
卻是事奉同一位主.
工作有很多種.
卻是同一位上帝在萬人中運行萬事.
聖靈彰顯在個人身上.
是要叫人得益處.
有人得到智慧的言語.
有人得到知識的言語.
有人領受信心.
有人領受醫病的恩賜.
有人作先知.
有人能夠辨別諸靈等等.
全部都是同一位聖靈所持.
同一位聖靈運行.
隨他的旨意分辨各人.
即是說原來一個聖靈充滿的教會.
不是每個人都能夠說起別國的話來.
相反.
聖靈會按著他的心意.
賜給每一位信徒.
有不同的恩賜和能力.
而不同的恩賜能力.
其實有一個目標.
目標是甚麼?.
是要使人得益處.

$^{161}$接著如果我們再看下去.
哥林多前書十二章的時候.
保羅在最後的時候就說.
你們要追求更大的恩賜.
我現在將最妙的道告訴你.
他講完告訴哥林多教會的信徒.
原來聖靈會賜下不同的恩賜.
接著他說.
不過你們要追求更大的恩賜.
我現在將最妙的道告訴你.
接著就下到第十三章.
是甚麼?.
哥林多前書十三章熟不熟?.
他說甚麼?.
你能夠做甚麼都好.
你說話好像天使一樣都好.
如果沒有了甚麼?.
愛都成為命的錄響的筆.
接著他就講了一個我們很熟悉的愛篇.
愛是恆久忍耐又有恩賜.
原來聖靈不單止賜給每一個信徒.
是有不同的恩賜和能力.
不過更重要.
更大的恩賜.
我們也要祈求上帝給我們的.
是一個愛的恩賜.
是一個能夠可以用愛來串連.
不同的能力的一個恩賜.
為要叫使人得益處.
如果沒有了愛.
就算我們曉飛.
能夠行神蹟醫治.
能夠講方言.
多麼奇妙都好.
其實都是命的錄響的筆.
甚至乎是會破壞上帝的工作.
所以因此這麼說.
一個聖靈充滿的教會.
就是我們要明白.
聖靈是會給我們每個人.

$^{201}$每個基督徒.
是有不同的恩賜.
有些人有教導的恩賜.
有些人有陪伴的恩賜.
有些人有關心人的恩賜.
有些是懂得弄音響的.
唱歌的.
玩音樂的.
等等很多很多不同的恩賜.
不同的恩賜最重要的是.
我們能夠有愛.
上帝的愛將它串連起來.
然後是叫人得益處.
愛很重要.
或者我分享一個故事給大家聽.
剛才說的聖學院是兩文三語的.
有一次密集課程.
什麼叫密集課程呢?.
密集課程就是連續一個星期五天.
早上九時到五時.
上午五日.
這些叫做密集課程.
那班是有普通話的同學.
即是有內地的同學.
亦都有海外的同學.
即是說那班是需要用普通話和英文上課.
那位老師是說普通話的.
那時候我被邀請去做翻譯.
由普通話翻譯做英文.
由廣東話翻譯做英文已經覺得很困難.
還要是普通話翻譯去做英文.
但是人手不夠也要硬著頭皮去.
其實五天也不太知道自己說了什麼.
盡力而為去做.
盡力去說.
五天就過了.
過了之後.
我真的感覺到聖靈的工作.
工作不是我能夠說起別國的話來.
我沒有這個恩賜.

$^{241}$因為我自己也知道自己的事.
自己也翻譯得一舊舊.
但是恩賜在哪裡呢?.
聖靈給我的很強的感覺是什麼呢?.
是那班海外的同學.
走過來跟我說辛苦你了.
你翻譯得很好諸如此類.
其實也知道都是一個愛心的話.
基本上是一種愛心的恩賜.
原來就是這樣一班人互相鼓勵.
縱然有時未必做得很好也好.
不要緊.
我們一起去學習.
一起去經驗.
一起去面對.
那種愛的能力.
其實也是聖靈在我們當中的工作.
親愛的兄弟姐妹.
一個聖靈充滿的教會.
就是一個愛的教會.
這個愛是讓我們願意.
拿出聖靈給我們個別的恩賜.
然後會聚.
然後一起去事奉.
一起去生活.
要的是讓到神的名得榮耀.
叫人得益處.
第二.
一個被聖靈充滿的教會.
是能夠想起和明白神的話.
《約翰福音》14章25到26節.
他記載主耶穌和門徒這樣說.
「我還與你們在一起的時候.
已對你們說了這些事.
但保衛斯」即是聖靈.
「就是父因我的名所要差離的聖靈.
祂要將一切的事教導你們.
並且要使你們想起.
我對你們所說的一切話」.
這是耶穌給門徒的應許.

$^{281}$被聖靈充滿的教會.
有聖靈同在的門徒.
他是能夠想起耶穌基督的說話.
想起神的說話.
今天所讀的《使徒行傳》2章.
讓我們看到被聖靈充滿的門徒.
他們不再因為耶穌的死亡復活升天.
感到困惑.
而是各自用不同的香談.
去跟別人講論上帝的大作為.
當群眾覺得很奇怪.
甚至有些人開始去譏笑他們的時候.
經文告訴我們.
彼得就跟十一個使徒站起.
然後彼得就開始了.
相信是第一次對著這麼多人的講道.
這一次臨時臨急的講道.
就不跟我有講章.
突然間彼得沒有講章.
臨時臨急的講道.
聖靈就幫助他想起.
《約珥書》2章28到32節的經文.
用這段經文去解釋當時發生的事.
就是上帝其實已經在先知書上表明.
在末後的日子.
會將聖靈囂貫在他的僕人和侍女身上.
會有這些奇妙的事情發生.
當然想起明白神的話語.
除了這一次的講道.
在往後的《使徒行傳》的記載裡面.
我們都多次讀到門徒,使徒.
他們引用當時他們的喜伯來聖經.
又或者主耶穌基督的教導去傳福音.
去處理他們面對的挑戰和問題.
就好像《使徒行傳》記載.
第一位殉道者斯提反.
在他被石頭打死之前.
他和當時的大祭司和群眾.
就從亞伯拉罕一直講到耶穌基督的事情.
各位弟兄姊妹.

$^{321}$對於基督徒來說.
聖經相當重要.
我們亦都很重視去研讀和學習聖經當中的內容.
不過我都會去問.
究竟各位弟兄姊妹.
為什麼我們這麼重視聖經?.
為什麼我們這麼用心去研讀,去學習聖經?.
什麼原因呢?.
其實在世界上.
不只是基督徒去研讀聖經.
很多非基督徒,即是不信的人.
他們都很努力去研讀聖經.
所以在很多大學裡面.
都有他們的宗教系.
宗教系裡面的學生很多都不是基督徒.
但是很多人很努力,很厲害地去研究.
基督教的信仰和經典.
去研究聖經.
究竟兩者有什麼分別呢?.
基督徒研究聖經和非基督徒研究聖經.
兩者有什麼分別呢?.
或者我這樣說.
非基督徒研究聖經或者研究基督宗教.
他當它是一件東西去研究.
聖經是一個歷史上的文獻寫了出來.
他們就去研究聖經裡面的內容.
是一個死了的,已經完成了的東西.
就去研究它.
得出一些學術的結論.
基督徒研究聖經或者去研讀聖經.
其實是不一樣的.
我希望大家都是用一個這樣的心態去看聖經.
基督徒看聖經是希望能夠聽到上帝的話語.
明白上帝的心意.
而我們想聽到,想明白上帝的心意.
上帝不是已故的,不是死了的.
不是歷史的一個上帝.
他是一個到今天都仍然和我們同在.
一個活的上帝.
我們相信一個活的上帝.

$^{361}$我們也相信他的話語都是活的.
什麼叫活?.
活的意思就不是聖經在歷史上寫好了.
就是這樣的了.
而是說上帝都透過這些文字.
和我們今天的人去說話.
上帝今天有什麼跟我們說?.
永活的上帝用一個活的道.
去向我們傳講他的心意.
所以上帝設立教會.
讓教會有崇拜.
崇拜有講道.
講道不是只是教導聖經之息.
講道是上帝透過講道.
去將他的心意向弟兄姊妹.
向他的子民去宣講.
在神學院裡面我是負責講道學的.
經過三年的疫情.
很多教會.
不只是香港.
全世界教會都已經很熟悉直播.
網上直播.
我們這些不懂對鏡頭說話的人.
現在經過三年的訓練.
我都懂得跟鏡頭打招呼.
看一下鏡頭.
這裡有個鏡頭那裡有個鏡頭.
要看一下.
已經很習慣了.
講道學的同學其實他們.
預備講道是很大壓力的.
覺得講道很難.
另一邊廂.
偶爾都聽到一些匯眾.
不知道大家有沒有.
可能這裡沒有.
這裡應該OK的.
可能其他教會有些匯眾.
有時候會說.
講道很悶.

$^{401}$沒什麼質素.
我就跟同學說.
你們有這麼大壓力.
他們又覺得悶.
現在直播又這麼方便.
很簡單.
香港就有1300個.
約莫1300間華人的堂會.
我當每一個星期日.
有1300個講章.
1300個講道.
其實都很浪費人的靈魂力.
你又覺得很困難.
你又覺得很悶.
又有直播.
不如我們找一些明星講完.
一個星期一個就行了.
全港1300個堂會直播.
那就很簡單.
又不用預備這麼多.
又不會悶.
有明星講完.
又省功夫.
那就好了.
行不行?.
好不好?.
我想一次幾次都可以.
短期都可以.
但是長遠不OK.
為什麼不OK?.
沒錯明星講完.
可能他真的講得很好聽.
很吸引.
很有內容.
但是明星講完.
他不認識1300間堂會裡的弟兄姊妹.
所以他的講道只能夠.
在一個比較大眾化的層面去講.
但是上帝就是差遣牧者.
去到每一個堂會.

$^{441}$去牧養他的弟兄姊妹.
所以講道的.
在堂會裡的牧者講道.
其實是認識弟兄姊妹的需要.
上帝就是透過牧者.
在主日的崇拜裡.
去將上帝的道.
告訴弟兄姊妹.
是按照弟兄姊妹的需要.
按照弟兄姊妹的情況.
去跟他們宣講.
這些明星講完是做不到.
因為他不認識那些羊.
但是堂會的牧者.
應該是要認識那些羊.
上帝也透過堂會的牧者.
去牧養這群羊群.
上帝聽弟兄姊妹的禱告.
然後就會透過講道.
或者在教會裡的一些活動.
將上帝的道帶給你們.
一個聖靈充滿的教會.
其實就是能夠有上帝.
能夠讓我們明白.
能夠記起上帝的話.
是一個永活的神.
跟我們講一些活的道.
是能夠對應大家的生活的道.
第三.
被聖靈充滿的教會.
是有能力面對別人的.
飢餓意志逼迫的.
一個聖靈充滿的教會.
是有能力去面對別人的.
飢餓意志逼迫.
今日經文第十二節.
眾人他們覺得很驚奇.
聽到他們講到別國的話來.
彼此就說這是什麼意思呢.
接著也有人就譏笑他們說.

$^{481}$他們是灌滿尿生酒.
這一班門徒受到譏笑.
不過彼得和這十一個使徒.
沒有因為這些譏笑.
覺得不開心不安.
或者有退縮.
反而他們站起來.
然後用約珥書去講道.
彼得和十一個使徒.
沒有因為這些譏笑.
而停下來而繼續傳道.
使徒行傳繼續告訴我們.
這一班使徒他們面對的.
不只是饑草.
他們面對的是監禁.
他們面對的是鞭打.
他們面對的是威嚇.
甚至是一個殉道.
但是他告訴我們.
這一班使徒仍然是很勇敢地.
去面對這些苦難.
他們仍然勇往直前.
不怕困難,不怕苦難.
面對著當權者的逼迫.
他們繼續將福音講絕.
彼得前書四章十二節.
他這樣說.
他說:親愛的.
有火一般的考驗臨到你門.
不要奇怪.
似乎是遭遇非常的事.
我重複.
親愛的.
有火一般的考驗臨到你門.
不要奇怪.
似乎是遭遇非常的事.
這一節經文.
作者跟信徒說.
你們現在面對著火一般的考驗.
不過你不要覺得奇怪.

$^{521}$似乎是一個.
似乎以為自己遇到非常的事.
意思是什麼?.
他只是說.
其實我們作為信徒.
遇到火一般的考驗.
不是非常.
是非非常.
即是正常.
現在我們疫情過後.
每樣都講復常.
復常是什麼意思呢?.
在基督徒來說.
我覺得大家值得深思.
復常是否就是.
大家不用戴口罩.
大家可以回來聚會.
大家可以去外面吃飯.
不用隔間板之類.
或者經濟可以自由一點活動.
旅遊業界又可以復甦.
不過對於很多底英姐妹來說.
復常最重要的是什麼呢?.
是我可以去旅行.
復常是否說這個常.
是大家期望的常.
但是另一邊廂我們看到香港.
或者看到世界.
其實面對著的困難依然存在.
香港的氣氛.
仍然都是灰灰濛濛.
世界有很多戰爭搏奕.
我們都面對很多困難.
移民潮等等都讓教會.
整個香港都有很多的.
低氣壓存在.
我們面對著這些情況.
究竟是非常還是正常呢?.
其實聖經常常都告訴我們.
我們作為基督徒.

$^{561}$我們面對的是一個聚世人生.
聚世人生有兩樣.
一是試探.
叫我們跌倒.
一是叫我們受苦難逼白.
這些才是正常.
原來世界面對這些是正常.
基督徒面對這些是更加正常.
問題是我們面對這些正常.
火一般的考驗的時候.
我們是怎樣面對的?.
我們的生命是怎樣面對?.
今天的經文告訴我們.
聖靈充滿的這班門徒.
他們是不怕飢潮.
不怕困難.
今天聖靈充滿的教會.
我們會不會都好像門徒一樣.
不怕飢潮.
不怕困難.
就算是灰灰濛濛.
就算是很多的試探和困難.
我們仍然都有能力去面對.
在一二世紀.
教會發展得很快.
不只是因為那班弟兄姊妹基督徒.
很努力去傳福音.
這個固然是.
不過發展得很快.
有些社會學者告訴我們.
其中一個很重要的原因.
是那班基督徒活出了一個.
和其他人不一樣的生命.
大家記得第一二世紀的時候.
基督徒是少數.
是小眾.
是被打壓的小眾.
但是這班基督徒.
他們的影響力為何這麼大?.
因為他們活出了一個.

$^{601}$很不一樣的生命.
例如是什麼?.
例如當時有瘟疫.
很嚴重的瘟疫.
很多人很怕染病.
很怕死亡.
甚至家人染了瘟疫.
第一時間就將他們趕出家門.
死了的人.
在街上死了的人.
沒有人敢去碰.
沒有人敢去埋葬.
在那個時間.
教會的一班基督徒.
願意去照顧這些病人.
願意去埋葬這些屍體.
活出了一個不一樣的生命.
為什麼他們能夠這樣做?.
就是因為他們.
聖靈告訴他們.
他們的生命不是這一生.
他們的生命是有永生的.
因為我們已經是屬基督的.
我們已經是一班重生的人.
我們已經是有永生的人.
這一生的生命其實不是一個完結.
我們有天家的生命.
就是這個既得.
這個信念.
讓這班人活出了不一樣的生命.
親愛的弟兄姊妹.
在我們生活裡面.
面對一些低氣壓.
困難,試探等等.
有沒有開放你的心.
讓聖靈提醒你.
關於上帝的說話.
我經常在想.
如果一個地方.
充滿埋怨.

$^{641}$或者很多怨恨.
如果在那個社會裡面.
基督徒的埋怨和怨恨.
都和其他人一樣多.
甚至乎多一點.
我們怎樣能夠彰顯上帝那種光明.
在那個社會裡面.
如果大家都覺得灰暗無望.
我們作為基督徒.
我們和別人一樣灰暗無望.
甚至乎黑色.
黑黑黑,比他們還差.
我們又怎樣彰顯上帝的榮耀.
親愛的弟兄姊妹.
我們需要一個被聖靈充滿的教會.
就是需要被聖靈開放我們的心.
被聖靈提醒我們.
關於上帝,關於聖經的心意.
以致我們能夠面對饑草.
甚至乎是逼白.
在這個社會裡面作炎,作光.
好嗎?.
我們一齊祈禱.
天父上帝,我們多謝你.
因為你常常用你的話語陪伴我們.
你的話語其實就是有聖靈在當中作供.
主上帝,我懇求你給我們.
每一個弟兄姊妹.
在今天這個聖靈降臨節.
再一次被聖靈充滿.
給我們有愛心.
給我們能夠嚮往你的話語.
記得你的話語,明白你的話語.
給我們面對饑草困難.
面對社會那種低氣壓,灰暗.
很多人想要去逃避.
我們能夠在這個社會裡面作炎,作光.
讓你的名得榮耀.
人得幫助.
我們這樣禱告.

$^{681}$奉耶穌基督的聖名.
阿們.
\newpage



\section{傳道書 11:1-6}
\label{sec:H85tdsHmqKI}
\textbf{2023-05-28 午堂《不知道中的信仰》傳道書 11\_1-6 講員:蔣文忠傳道}
\newline
\newline
連結: \href{https://youtube.com/watch?v=H85tdsHmqKI}{\texttt{https://youtube.com/watch?v=H85tdsHmqKI}} ~~~~ 語音日期: 2023-05-30
\newline
\newline
\hyperref[sec:hSecMeCWH4s]{\small{< < < PREV SERMON < < <}}
~
\hyperref[sec:index_chronic]{\small{[返順時目]}}
~
\hyperref[sec:index_scriptual]{\small{[返順卷目]}}
~
\hyperref[sec:C1vy6YO3h0g]{\small{> > > NEXT SERMON > > >}}
\newline
\newline
傳道書 11:1-6
\newline
\begin{longtable}{cl}
\hline
\hline
章節 & 經文 (和合本修訂版)\\
\hline
11:1 & \begin{tabularx}{0.7\textwidth}{X} 當將你的糧食撒在水面上，因為日子久了，你必能得著它。 \end{tabularx} \\ \\ \relax
11:2 & \begin{tabularx}{0.7\textwidth}{X} 將你所擁有的分給七人，或八人，因為你不知道會有甚麼災禍臨到地上。 \end{tabularx} \\ \\ \relax
11:3 & \begin{tabularx}{0.7\textwidth}{X} 雲若滿了雨，就必傾倒在地上。樹向南倒，或向北倒，樹倒在何處，就留在何處。 \end{tabularx} \\ \\ \relax
11:4 & \begin{tabularx}{0.7\textwidth}{X} 看風的，必不撒種；望雲的，必不收割。 \end{tabularx} \\ \\ \relax
11:5 & \begin{tabularx}{0.7\textwidth}{X} 你不知道氣息如何進入孕婦的骨頭裡；照樣，造萬物之神的作為，你也無從得知。 \end{tabularx} \\ \\ \relax
11:6 & \begin{tabularx}{0.7\textwidth}{X} 早晨要撒種，晚上也不要歇手，因為你不知道哪一樣發旺；前者或後者，或兩者都一樣好。 \end{tabularx} \\ \\
[1ex]
\hline
\hline
\end{longtable}
$^{1}$弟子們早安.
我今天在大家中間分享上帝的訊息.
特別能夠聽到阿拉伯宣教的分享.
.
去帶給我們一個鼓勵.
或者大家可能都認識的話.
可能知道在一個福音的顏土.
去能夠傳遞福音.
或者是用不同的方法.
讓未得之民得到福音.
其實是一個很不容易的事.
一個很多挑戰,很多難度的事.
如何配合我今天跟大家分享的訊息呢?.
其實就是對於我們身處的一個現實.
其實倒過來想.
我們活在香港.
其實可能你要去傳福音都不是容易.
我們都有另一種的挑戰.
可能是一個物質豐富.
或者是一個好像很多的知識.
很多的事情都在自己的控制裡面的時候.
你要向一個朋友分享福音都是很不容易.
所以今天的傳道宣教訊息.
可能許許我們就從一個.
更加去了解我們生活的現實.
進入到今天上帝的話當中.
求主幫助讓我們能夠在今天的道當中.
可以有主給我們的啟發.
有主給我們的觸動.
我們先一起有個禱告.
然後開始聽今天的訊息.
能夠有宣教的分享.
也讓我們能夠繼續去.
讓我們的眼光不會停在自己身上.
我們看到這個世界.
這個地方這個社會.
其實有很多的需要.
也都有很多的艱難很多的挑戰.
按從人的角度.
其實是很多問題都解決不了.

$^{41}$或者從我們的力量的角度.
我們覺得人很有限.
但求你幫助我們.
讓我們特別今天從傳道書的訊息.
去讓我們得到一個鼓勵.
或者得到一個啟發.
能夠從傳道書的一些智慧文學.
去幫助我們去明白主你的工作.
是更加讓我們更加對於生命.
有更多的體會.
也對我們向別人去分享福音的時候.
能夠有更闊的角度去明白.
怎樣去從一個人生.
一個生命的一些限制的角度.
去跟人分享福音.
願你親自的去與我同在.
保守我們 帶領我們以下的時間.
同心祈禱奉救主耶穌.
得勝命自求 阿們.
我估計如果不是我講到的話.
不知道大家有沒有機會看過這段經文.
可能我們揭傳道書都不多機會.
我記得我自己入神學院.
讀一科傳道書的時候.
老師跟我說那時候我還不夠30歲.
傳道書是要40歲以上的人才明白的.
隨著年華歲月增長.
我發覺真的.
偶爾我都會拿傳道書來讀.
特別在這幾年的香港環境.
也更加有時會讀多了傳道書經文.
我偶然就讀到這段經文.
令我有一個很深的體會.
想跟大家分享.
不知道在座的你還不夠40歲的時候.
希望你都能夠讀得明白傳道書今天的訊息.
不知道大家是否認同.
我們都會在一個比較逆境的情況下.
我們今天疫情.
一種環境.

$^{81}$我們都說一種叫抗疫.
一種對於整個香港或整個世界的大洗牌.
各種的生活節奏,習慣,經濟.
各樣的環境關係.
都給我們面臨一個很大的挑戰.
所以產生一個對我們來說是一個艱難挑戰的一種逆境.
在一種逆境當中.
其實我們在信仰當中.
更加需要主一同行.
有主一同行裡面.
更加需要有弟兄姊妹.
身邊有人支持去同行.
因為大環境的衝擊.
令到我們生命裡面可能都很乏力.
對前景感到不光明.
我不知道大家的生活節奏如何.
或者睡得好不好.
有時我們從自己的生活裡面.
就可以檢視一下自己.
有時我自己都睡得不好.
各種原因.
有時都很容易會累.
不知道大家覺得會不會這幾年的生活裡面.
除了年紀增長之外.
是喔 容易累了.
或者休息得好像休息都不夠.
不知道有沒有這種狀態.
所以大家要好好鍛鍊好自己.
希望在神裡面我們可以重新得力.
重新在神裡面捉緊一番.
對生命的一種妥善的節奏.
讓我們休息有時 工作有時.
讓我們能夠沉默有時.
回應上帝的使命有時.
我們需要有主的力量.
特別是群體的支持.
讓我們分享內心的感受和掙扎.
這個都是很重要的.
如果我們有什麼感受.
或者有什麼情緒.

$^{121}$你埋在自己裡面.
其實大家都明白.
不會是一個解決的良方.
我們需要一個適切的渠道去抒發自己.
甚至最好每天都去抒發.
你就不會積累了一些.
壓抑了的一些感覺或者感受等等.
你最好有一個渠道.
有身邊一些信任的朋友.
家人 教會的弟兄姊妹.
讓你可以定期地將自己生命裡面.
體會到一些可能是負面的東西.
可以有機會分享.
這個是非常重要的.
所以我們需要同路人.
我自己都是特別做年輕人工作.
特別覺得年輕人都需要同路人.
雖然可能他們覺得不需要.
或者他們覺得我自己能夠應付.
其實不是的.
他們定時定候都需要有導師.
有人陪伴他們 開解他們.
面對他們生活的問題.
我們需要有信任的關係.
去真誠面對不同的問題.
這份真誠是很重要的.
直接去坦誠地將我們生命的掙扎.
我們不需要強裝自己是很有能力.
很有把握可以處理得到.
有時候我們需要儘早向人求助.
去讓人知道我們生命有軟弱的時候.
需要有弟兄姊妹代禱 守望.
甚至需要一些專業的援助.
去支持我們的生命.
以致我們能夠身心靈.
才能夠重新回復看見.
我們就從傳道書去智慧文學裡面.
我們去體會一下主的訊息.
智慧文學是一個很特別的一種文體.
讓我們更加從一個生活和社會世界的現實.

$^{161}$去參透人生的處境和狀況.
它不是不講信心.
不過它會提示我們.
我們的信心都不可以麻木.
我們的信心都需要衡量我們的處境之下.
來去調節好.
以致我們投放在主的信心是更加準確.
我相信 特別是Mr. Quat 弟兄.
我們在阿拉伯的世界裡面.
可能在做福音工作裡面.
其實有很多挫折.
譬如我們在香港.
至少我們還有一個宗教自由.
譬如我們教會搞街頭報道.
出元朗 起碼你都會接觸到一些人.
可能你回來跟牧師說.
我起碼跟了三個人做福音.
他信不信就另一個回事.
但起碼我們不會有困難.
你走過去接觸他.
他聽不聽就對方問題.
我們起碼不會有一個障礙去接觸到人.
但你在一個敵對基督教環境裡面.
就完全不是一回事.
可能我們很艱難.
又要衡量很多的環境.
是否安全.
校園裡面就不可以做.
要在一個私人的暗角裡面.
找到一個方向對象.
暗暗才能跟他說.
可能找到一個已經很難得 很感恩.
那個現實的困難.
是不能跟我們相比.
所以某程度上在傳道書裡面的經文.
就是一種讓我們去看透現實.
但又卻不被現實打倒.
來繼續為主而活.
繼續為主發光發熱.
相信這個都是.

$^{201}$特別對於宣教工作是更加需要的時候.
我們除了念讀大使命的經文.
對於智慧文學其實都需要很有認識.
以致我們能夠不被現實打倒.
不會被困難所令自己有挫敗和洩氣.
能夠可以持續地為主去工作去而活.
這本經文我的講題是.
「不知道中的信仰」.
很有趣.
我這本《采信經文》裡面有三至四個「不知道」.
我都在經文上用紅色讓大家提示得到.
我不再重讀一次.
剛才多謝主席幫我讀了一篇經文.
在第一節裡面有一個「不知道」.
將來有什麼的災禍臨到.
在第五節.
「不知道」風和懷孕的婦人的胎兒怎樣成長.
行萬事的神的作為我們就更加不知道.
是一個比較.
那些事你都不知道.
何況行萬事的上帝你更加不知道.
第六節很有意思.
亦是我最喜歡的一節.
不知道怎樣都好.
你就要繼續好好把握機會去殺種.
這就是我們到最後的時候會去到一個很鼓勵性的訊息.
這三個的「不知道」.
是不知道災禍何時來到.
不知道創造主隱密的作為.
它是真實的.
但其實不在我們能夠眼見.
就立刻捉到上帝的工作.
隱密的作為.
第三個就是不知道效果怎樣.
但卻是對於我們來說.
它是要盡力去回應.
做好自己的本份.
做好自己把握機會的工作.
既然因為我們的效果不在我們手上控制.
這就給大家一個總結的訊息.

$^{241}$知道自己的限制的時候.
其實我們的回應是怎樣呢.
我們不是要去否認.
有時候我們不認自己有限制.
譬如我剛才說了.
可能體力下降了.
我說不是的.
怎樣怎樣的時候.
我們希望自己十幾二十小時都可以不做工作.
其實是不行的.
那就是一個人生的限制.
我們要面對現實就是這樣.
我們要擁抱我們生命裡面的一些限制.
然後就能夠找到自己的範圍之內.
就在那裡盡力為主而做.
上帝就會在其中為我們成就大事.
這個就是大概在傳道書裡面.
我們可以從這個角度去理解.
我們面對現實之下.
也都擁抱現實.
也都能夠不被現實所困.
我們先從第一段去分享.
第一段裡面講到.
其實在傳道書裡面.
或者很多真言裡面.
都是很多的.
好像一些謎語.
或者很多的.
可能從一些當時的猶太社會當中.
一些日常的生活當中.
就拿一些.
隨意拿一些身邊發生的事.
拿出來講.
所以可能你聽下去的時候.
有些爭執.
你不知道他為甚麼要講這些東西.
但是你多看幾次.
你讀多幾次.
你體會那一節裡面.
他其實是想有甚麼寓意.

$^{281}$想帶出來的時候.
你就會得到當中的味道.
我們不會將糧食灑在水面.
這是第一件事.
你不知道他在講甚麼.
他說為甚麼要這樣做呢.
因為日久必得著.
然後就將你有的東西分給七份.
分給七個人.
分給八個人.
我們今天就講分散投資.
就是這件事.
如果你有錢.
你去買一些股票.
叫做甚麼.
我不懂得說.
叫做甚麼.
風險大一點.
風險小一點.
你都會分頭去做.
一邊你蝕了.
一邊你會賺回來.
其實就是講類似這個概念.
或者這個行動.
分開分攤去給不同的人.
有些人可能會風生水起.
有些人可能會倒楣.
但不要緊.
有部份人你給了他一個禮.
他將來會回饋給你.
其實是講這件事.
你不知道將來有甚麼禍患.
會淋到你身上.
所以可能你說.
嘩 這麼現實.
你這麼講一些風險評估.
不是我們的生命全憑信心嗎.
第一節我想跟大家分享一點.
就是說我們的主.
如果只是叫我們一味憑信心的時候.

$^{321}$但卻不去了解我們真實.
面對一個甚麼處境的時候.
其實可能我們信心.
剛才講了都是麻木的.
我們要了解現實那種虛實.
傳道書來說那種虛空.
虛空的意思就是說.
是一些我們不能夠把握得到的事情.
我們想盡辦法.
用盡人的有的知識.
其實你都把握不到.
你不能夠為將來控制得很好.
這個就是那種生命.
傳道者講到的那種虛空.
他能夠看透那種虛空.
當我們能夠看透虛空的時候.
我們第一時間.
或者生命裡面有甚麼反應呢.
就是一種不會過份樂觀的反應.
不會單單以為憑信心.
就一切就可以解決到.
當我們感覺好像很小心.
就是說其實我們的信心同時間.
是了解我們真實現實的環境.
我們分配我們擁有的東西.
相反有時候我們憑信心的時候.
就是上帝工作的.
我們就會反而忽略了人的工作.
容易跌落一種守株待兔的心態.
我們反而會傾向了自己.
應該要盡力的那一部分.
其實我們不是只想等上帝工作.
我們要盡我們那一份.
做好自己的那一部分.
原因是因為事情發展不在我們掌握裡面.
隨時都要迎接不同的結果.
不同的結果是隨時發生.
甚至是一些事與願違的結果都會發生.
所以傳道士第一節就講.
其實我們不知道災禍何時會來到.

$^{361}$所以我們要有心理準備.
我們要為自己的生命做好準備.
當這些事情發生的時候.
我們能夠處變不驚.
能夠靠著上帝.
可以應對到這些災難或者苦難.
來到自己的生命當中.
這就是第一點他所表達的意思.
所以那種分散.
去擺上你自己的生命的資源.
其實是講一種我們自己要為自己做好準備.
既然我們不被災禍.
不能夠知道何時來到.
或者這樣說.
其實沒有一個定律.
沒有一個定律或者沒有一些規則.
我們能夠所謂預測到災禍何時來到.
其實是怎樣呢.
是不是要悲觀.
我們控制不了.
等那些禍患能到.
還是一種過份樂觀.
不如我們今早有酒今早醉就算了.
也不是.
也不是一種過份悲觀.
也不是一種過份樂觀的心情.
其實傳道書很現實又很務實.
既然苦難來到何時你都不知道.
那就不如豁出去吧.
自由地將自己的生命.
擺給上帝去使用吧.
不要被那些經驗困住.
不要被聽到別人的生命的經歷.
嚇怕自己.
別人是這樣我也會是這樣.
每個人都獨特的.
你自己有上帝帶領你走的路.
你就豁出去.
痛痛快快為上帝活出你的生命.
第二段是最難解的.

$^{401}$第三第四節.
不要被現實嚇怕而行動.
它說什麼呢.
我將經文都分開了一點.
它說雲若滿了雨傾倒在地上.
樹若向南倒向北倒.
倒在那裡它就在那裡了.
這兩節說什麼呢.
它說的一種就是.
所有的結果必有因.
當場大雨的時候.
因為雲再滿了雨水.
所以就淋下來.
你看到那棵樹掉在那裡.
因為有風吹向著它.
它就掉在那裡了.
這是一個每次你看到的現實結果.
其實一定有一個原因.
這個就是第三節上半節.
傳道者從觀察現實去看到的東西.
但後面就很有意思.
看風的就必不殺種.
望雲的就必不收穫.
什麼意思呢.
其實是在說天文台的第一句.
就是說它一直留意著風向.
颱風正在殺過來的時候.
你就什麼都不會做了.
你要在這裡把家裡都貼了膠布.
躲在家裡什麼都不做了.
接著就是看著那些雲一塊塊衝過來的時候.
你就不會出外玩了.
你一定保護小朋友在家裡.
因為預測的時候.
你被預測的現實嚇怕了.
你就不會敢於行動.
那想說什麼呢.
他想說什麼呢.
就是告訴他.
沒錯.

$^{441}$當你遇到一件事情.
就是淋一場大雨的時候.
它當然有個原因.
雨雲不會突然間走出來.
但是不一定每次有雨雲的時候都會下雨.
其實想說一件事.
那就看我們了.
那你的信心有多大呢.
究竟這雲一塊塊來到的時候.
你會不會出去行山呢.
有些人可能會說.
明明昨天天文台已經報告了下雨.
今天我就取消我的行程.
有些人不會的.
先看到明天早上.
出到門口之後.
才決定雲可能已經走過了.
沒有下雨.
那他就可以繼續放到陰天.
可以有一個很清涼的一天去行山.
是在說這件事.
每一件事.
沒錯.
這樣發生.
但是不會說每一次有雲的時候都會下雨.
不是每一次有風來的時候都會打風.
那個打風可能會離開了香港.
走去別的地方.
所以你不要被那些定律綁死了你自己.
不要被你過去的經驗綁死了你自己.
不知道大家明白我的意思嗎.
就是說.
我們可以不受一個現實嚇怕.
而什麼都不做.
所以傳道書有時候是超現實.
傳道書有時候是更加鼓勵我們.
不要被現實去困擋住我們.
不要被我們過去的經驗困擋住我們.
每一天都是新的.
你猜不到主在今天.

$^{481}$是會給你一個特別的天氣.
給你一個特別的際遇.
是可以神跡奇事地經過了一個很奇妙的日子.
可能那個條件那個因素和昨天是一樣.
但是今天你的經歷是可以不同.
因果不是必然.
有時候我們會將那個因果.
太過看到好像一種宿命的說法.
是每次可能都會有不同的結果.
每次的條件是一樣.
但可能是帶出了不同的結果.
所以傳道者用大自然的定律.
告訴他人生都是這樣.
你可以改變的.
每次做到一個很疲憊的時候.
結果可能就是一種枯竭的生命.
burn out.
但如果你及時將自己的生命改變的時候.
你不一定會跌入一種身體脆弱的狀態.
我們每一次都有機會為主去改變.
和可以經歷到主恩典更大.
不要被現實所嚇怕.
所以既然我們不知道災禍來到.
其實反而傳道書提醒我們.
你豁出去.
不要被各種好像風雨欲來的事嚇到.
你來到的時候.
你還有可能性為自己計劃自己的人生.
所以其實都很現實.
而且是叫我們把握到最後的機會.
去到他不知道主的作為這一點.
傳道者也是在用一些當時.
或者他自己體會的一些真實的情況.
風從哪裡來你不知道.
即是說其實一刮過來.
你被一場風刮到頂毛都吹掉了.
但究竟那場風在山上怎樣醞釀.
然後最後那個氣候.
那個龍捲風或者那些熱帶風暴怎樣形成的呢.
其實你是不能夠很仔細地掌握到的.

$^{521}$那個源頭是怎樣來的.
你是不知道的.
接著再說骨頭在懷孕的婦人當中.
究竟她是怎樣成長的呢.
有時你找醫生看著.
照超聲波的時候.
你見到她有心跳.
有所有的輪廓出現了.
你已經看到結果了.
但究竟在過程裡面.
其實是怎樣發生的呢.
你不能夠那麼仔細地看著那個小朋友的生命.
逐個部分逐個部分地形成出來.
你是不知道的.
是一個隱密的工作.
可能你每次檢查的時候.
你都見到一個分別.
但中間的過程.
你是不知道很奇妙地.
她是怎樣成長下來的.
我們尚且不知道這些生命.
或者是大自然的東西的變化.
何況行萬事的神.
我們更加不知道.
這個行萬事的神.
意思是.
我將它譯成一個.
Who makes the whole.
將事情圓滿成就的神.
你更加不知道.
即是說.
其實每一件事情.
都是上帝讓它圓滿發生成就.
有時候我們只是看到結果.
而是不知道上帝在當中.
背後是怎樣工作的.
弟子們.
有更大的神在掌管.
有更大的智慧.
在讓這個世界去運行.

$^{561}$如果我們只是按著.
我們人自己的智慧.
人自己的經驗去掌握的時候.
其實我們會失去很多.
但矛盾就是.
其實聖經說.
其實有更大的智慧.
更大的計劃.
其實我們不能夠完全掌握.
那怎麼辦呢.
原來面對上帝的時候.
我們不是要掌握住祂.
在我們的手中.
有時候我們很想一切事情都運籌為額.
一切事情都是好好自己管理得好.
連上帝的計劃都是.
但傳道書告訴我們.
我們要承認我們不知道.
甚至它用一個字再進一點.
是用無知的.
我們無從知曉.
上帝在隱而隱密的作為.
那人有什麼反應呢.
人可以怎樣呢.
面對一個你完全不能夠掌握的上帝.
第一.
你相信祂是次因的上帝.
你掌握不到祂.
你不要緊.
祂每次為你所預備的.
是出乎你意料之外.
更加美好.
更加似更大的祝福.
首先你相信.
我們上帝是次因的上帝.
第二.
既然祂不能夠在我們掌握之中.
我們就保持一種開放性.
保持對祂常存敬畏的心.
知道我做完.

$^{601}$剛才說很現實.
我做完那一部分.
其實你知道還有一個更大的上帝.
在這裡瞭望著整個地球發生的事.
瞭望著生命當中下一步會怎樣經歷.
祂知道.
但我們還不知道.
你就放手.
敬畏祂.
來面向上帝為你預備的道路.
這就是傳道書帶給我們一種對上帝的關係.
你記得鄒賢經常說.
敬畏而又說是知 或者開端.
其實道理就是這樣.
那種fear the lord.
那種真正對上帝的敬畏.
是在於祂是超乎現實中我們所把握的一切.
祂往往會在生命當中.
在不同的時候介入.
賜給我們出乎意料的驚喜.
一些祝福.
一些恩典.
所以不知道隱密上帝的作為.
是對於我們在世上過活.
是一個很重要的體會.
我們永遠記得有一個更大的上帝在掌管.
我們自己不是最大.
就算我們人生當中已經過了很多的經歷.
可能對於你的工作.
對於你很多的事情.
你很有把握也好.
但有時你知道健康的問題.
很多的突如其來發生的事情.
都不在我們掌握之中.
我們這份常存敬畏的心是很重要.
無論我們在人生中有什麼閱歷.
什麼年資也好.
其實盼望我們都能夠更深體會到.
這個我們不知道,不掌握到上帝這個隱密的作為.
就算我們在人生中可能已經以為最圓滿了.

$^{641}$其實可能上帝隱密作為.
更加繼續賜給我們.
所以我們繼續保持對上帝的開放.
去迎接祂賜給我們的作為.
到了最後第三段.
不知道效果如何.
當我們不被現實嚇倒.
然後也要敬畏我們不能夠掌握祂工作的上帝.
那麼我們在我們的日常中.
我們的生活中如何過活呢.
早晨要殺你的種.
晚上也不要揭你的手.
因為你不知道哪一樣發旺.
或是早殺的.
或是晚殺的.
或是兩樣都好.
祂說的是什麼呢.
你知道農夫有不同的節奏.
或者對於節氣的掌握.
早上出去田中殺種.
或者晚上去也不停.
可能有不同的鼓有不同的想法.
哪些會容易發芽成長得更加茂盛等等.
祂說你不要停手.
早上你出去.
晚上你出去.
為什麼呢.
因為你有些東西你都不知道.
因為你不知道哪一樣會帶來結果.
你不知道早上出去殺了之後.
會被蝗蟲吃掉.
會有一個濕潤的雨露.
讓你的種子可以發芽成長得更好.
你不知道.
所以祂說.
既然你大自然的世界掌握不到的時候.
你就盡量做好自己.
盡量做足你自己手上能夠做到的事.
你早上也要去.
晚上也要去.

$^{681}$因為或是早殺的有結果.
或是晚上殺的有結果.
或是兩樣都好也不一定.
你雙倍的結果有豐收的成果.
這個就是第六節祂所表達的意思.
弟兄姊妹原來當我們不知道.
效果如何的時候.
我們有什麼反應呢.
有時候我們會.
不如有結果才算吧.
有時候我們都會採取一種消極的態度.
因為我掌握不到.
做來也不知道有沒有結果.
那我不如不要做.
不知道大家有沒有聽過.
不同的人有不同的態度.
未必是一種很積極的回應.
大學聲講其實傳道者裡面.
他鼓勵我們.
其實是當我們真實面對.
承認現實的時候.
反而是一種真實.
合乎信仰的一種積極.
也不是過於看重自己.
也不是將自己放大.
而是明白到.
你要回應上帝.
上帝不知道的作為.
是隨時在我們各樣的機會當中.
是其中一樣.
顯明他的作為也不一定.
幾種可能性都會發生.
都不在人的手中.
這個就是農夫的工作的真實情況.
可能我們今天.
未必會在農夫的工作.
那麼不能夠.
即是要面對大自然.
所以很多事情.
他真的要望天打卦.

$^{721}$要在這裡看著天來行事.
但可能我們今天工作.
我們很多都是商業社會裡面.
可能很多都在辦公室當中工作.
或者我們未必明白.
這種現實的一種變數.
但其實我們有時只不過是被一些.
表面的工作環境.
其實是遮蓋了一些現實的變數.
現實的變數是隨時都會出現.
在我們的生命當中.
所以最後這個殺種的一個比喻.
其實是在提醒我們.
要盡所能.
不要有多少種可能性的時候.
你都用盡你的方法去把口.
你不知道哪一種有結果.
最近我兒子開始找工作.
找教師工作.
如果你是老師.
你都知道這幾年的移民潮都很嚴峻.
教師很多人都走了.
他很挑剔.
我要找南校的.
我要找基督教學校的.
我說兒子.
有多少間.
你又要把南校的資料.
你又要把基督教的資料.
100間裡面.
有個廣告的時候.
可能只剩下5間.
結果是一種很大的.
你不試過.
甚至你試了一間的時候.
你見過一份工作.
你拿了經驗.
其實你才到機會來到的時候.
你才能辦得到.
然後我說我不喜歡約見.

$^{761}$我不喜歡面試.
你不喜歡面試.
你就可能要克服心理的障礙.
其實你可能面試了第三次.
第四次.
最後到第五次的時候.
你那份工作.
上帝預備的那份工作.
才在你手上得到.
你不會從天而降.
一擊擊中.
就是這樣.
所以其實就算回到我們.
這個都市社會當中的時候.
其實這個.
無論早殺夜殺這種心態.
其實都是可以應用在我們生命之中.
弟兄姊妹.
我們能不能夠.
好像傳道書所說.
我們只管去盡力.
雖然結果不在手中.
你相信上帝一定會為你去預備.
就好像《先秘的百二十六篇》那樣說.
大家都熟悉的.
第六節.
那殺種.
那帶著腫瘤類出去的.
必要歡歡樂樂地帶禍菌回來.
很辛酸的.
帶著流著淚出去.
很艱難.
現實有很多限制.
甚至現實很殘酷.
去見工.
分分鐘打了十封信.
都沒有回音.
不會說你是哪裡畢業.
就一定會見你.
不一定.

$^{801}$你不要以為自己的履歷很好的時候.
那份工作一定會等著你.
不是這樣的.
但是你相信.
你繼續去嘗試的時候.
終有一天是歡歡樂樂地.
帶著禍菌回家.
這是上帝給我們的一個.
就算傳道書也好.
或者詩篇也好.
都是給了我們這樣的一個應許.
最後我想用一首詩去結束.
一個.
這個不說了.
一位耶穌會的修士.
他有一首的祈禱.
詩叫.
Trust in the slow work of God.
即是.
有趣的不是快工作.
即是慢的工作.
即是怎樣Trust in慢的工作呢.
很快和大家去看一次.
這個就是英文的版本.
他說有時候我們是很沒有耐性的.
對上帝.
甚至我們很想skip了一些intermediate stages.
我們很想即刻到結果.
上帝有嘛.
快點給我結果.
我不想經歷中間的interview.
被人問到.
臉都糊了.
人家又用了眼光看著.
我不想經歷這些過程.
我想即刻有.
上帝有給我.
就即刻有那個結果.
我們很想.
不想去經過一些不穩定的階段.

$^{841}$這個就是中文翻譯.
我們走到未知的時候總會缺乏耐性.
我們不想面對那些不穩定的階段.
因為這個疫情可能費事很多.
很長的時間.
我們不想經過那個process.
不想經過那個過程.
但傳統書告訴你.
人生的過程是真實的.
是必須經過.
反而你踏進那個過程.
你最終一定會見到上帝.
到中間的時候.
祂就說.
其實你也是一樣的.
你的生命是要假以時日才成熟的.
你的思想.
是不是想想你.
你今天可能你能夠做到很好的判斷.
你不是二十歲就立刻有的.
可能你經過二十年,三十年.
你到今天才思想成熟.
你才懂得處變不驚去做一些決定.
那個過程,那個成長是必須的.
也都不可以急功近利.
不可以催逼你生命.
不可以催化它成長.
這個就是你今天.
你是一個明天的我.
你不可以催逼到今天.
就即時出現了一個明天才有的你.
時間是明天才出來醞釀的你.
你不可以以為今天就不可以催成.
你不可以揣描助長.
催逼生命很快地成長.
所以只有上帝在你裡面逐漸形成.
一個新的靈在裡面帶給你成長.
我們相信上主的大令吧.
就算處於一個焦慮的感受.
在一些不完整的狀態當中.

$^{881}$你都接受上帝繼續他的恩手.
拖帶著你去經過.
求主幫助我們.
讓我們在我們剛才說的逆境.
不知道大家對於這個狀態有多深的體會.
盼望傳道書的訊息.
是當中的智慧.
當中對人生的現實那種的汗頭.
對上帝的信心重新有一種.
有深度的一種的體會.
能夠幫助我們今天能夠繼續.
憑信心也能夠活得豐盛地.
去走人生的路.
我們一起低頭禱告.
七姐妹謝謝你給我們今天.
藉著傳道書的訊息.
裡面有很多的智慧.
很多的格言.
很多的平時在身邊出現的一些事物.
大自然的東西.
原來都能夠參透到一些.
對人生的領受.
對你的作為的體會.
對你的態度.
一份敬畏.
一份等候.
一份讓你在生命當中.
去拖帶我們走.
讓我們去接受.
我們一些需要經過的.
必須經過的過程.
求你當中親自的報仇.
我們可能在今天.
面對的現實都是充滿挑戰.
我們盼望不被現實所打倒.
所o下怕.
也都不需要掉落一個消極.
什麼都放棄的狀態.
也都不需要過份樂觀.
好像見不到現實那種的殘酷.

$^{921}$卻是能夠很真實地.
腳踏實地在當中.
把握我們人生的機會.
積極地去回應.
相信你.
忍而未現的作為.
必然讓我們有一天能夠看得見.
必然能夠像詩篇所說.
我們必然會歡呼收割.
你所要我們所可以看見的成果.
願你保守.
特別祝福.
在阿拉伯的科研工作.
求你當中.
在科研的硬土裡面.
宣教士,宣教機構都用盡很多的方法.
務求讓未得之民能夠聽聞福音.
求主你讓我們面對一些現實.
一些困難裡面.
都不被打倒.
能夠繼續在裡面.
去滋養一份信心.
看透你.
繼續慢慢地.
繼續但沒有停止的作為.
讓我們每一個都能夠.
去堅持下去.
榮耀你的名.
同心祈禱奉主名.
阿們.
\newpage



\section{}
\label{sec:C1vy6YO3h0g}
\textbf{2023-06-01 周四晚查經班《被擄歸回歷史研讀(16)》 講員:曾裔貴牧師}
\newline
\newline
連結: \href{https://youtube.com/watch?v=C1vy6YO3h0g}{\texttt{https://youtube.com/watch?v=C1vy6YO3h0g}} ~~~~ 語音日期: 2023-06-03
\newline
\newline
\hyperref[sec:H85tdsHmqKI]{\small{< < < PREV SERMON < < <}}
~
\hyperref[sec:index_chronic]{\small{[返順時目]}}
~
\hyperref[sec:index_scriptual]{\small{[返順卷目]}}
~
\hyperref[sec:WLN153QoHOE]{\small{> > > NEXT SERMON > > >}}
\newline
\newline
$^{1}$同大家一齊祈禱先.
再一次我們感恩.
因為又去到星期四晚.
我們可以來到去研讀經卷.
《尼欺米記》其實是歸回歷史書的其中一卷.
我們很盼望我們能夠這樣仔細.
在這些大的歷史的那個時代.
來到去研讀的時候.
都幫助我們能夠從其中.
吸取到一些屬靈的教訓.
讓我們能夠看到在過去.
特別是被擄之後.
這一段相當長的歸回的歷史.
神你怎樣來到去透過聖靈.
去感動一些的領袖.
無論是祭司的.
或者是酒正的.
又或者是其他一些的利未人等等.
能夠上下一心.
來到為神的聖殿.
為神的選民.
來效力.
這樣去讓那個時代.
能夠有一個復興.
我們今天看的這幾張經文.
很肯定是一個很大的復興.
是帶來一個相當影響力的一個時候.
我們不知道一路可以延續多長.
但是我們看到.
神的工作就是這樣.
不斷在神的選民當中來工作.
求主帶領我們.
我們今天晚上能夠有幸去看這個復興的時候.
都給我們自己有一些反思的空間.
我們這樣仰望禱告.
奉耶穌基督的名.
阿們.
無論我們看到那個編排.
第三章正正式式.
黎希米的計劃.

$^{41}$由他在亞達舍西王.
他做酒政的時候.
他已經有這個規劃.
或者他未必想到將來他上到耶路撒冷.
是會遇到這麼大的困難.
那個困難其中是外力.
就是當地的原居民.
參巴拉.
阿拉伯人基善.
或者是其他的人.
再有另外就是.
在耶路撒冷本身的一些領袖.
有些是祭司.
這樣來阻撓他的工作.
因為要變是一件很大的事.
作為一個領袖.
其實就是你要看到那個求變的原因.
不是為變而變.
是實在有需要要變.
但是其中最大的困難就是.
人的心未準備的時候.
他就會無緣無故或者是有意地.
去阻撓那個改變.
因為習慣了是很難改變的.
無論那個習慣是一個好的還是不好的習慣.
所以很自然的就是面對第四章的攻擊.
詆毀.
張力.
第五章就是內部的矛盾.
就是連吃都沒得吃了.
你還說什麼重建呢.
原來沒得吃的背後.
是他們習以為常.
已經見怪不怪的就是.
那些省長抽很重的稅收.
百姓就很辛苦.
有時候如果失收的時候就更加百上加根.
再加上就是很多經年累月.
就是他們慢慢由本來有土地的慢慢都變賣.
而來到是一個這樣的結構性裡面的貧窮.

$^{81}$就只是帶來那些有資產的就更加富有.
那黎希米就見識到這個情況.
他自己就以身作則.
將他省長的俸祿都來到去不要.
然後將他有的去招呼那些幫手重建的那些工人.
那些的工頭.
來做一個以身作則.
第六章就更加明顯了.
這次就是針對黎希米他的人格來攻擊.
這個很自然的了.
如果有人用一些圖像在一個part point上面.
星期日打曾牧師.
勵善牧師.
鄭牧師.
是一個怎樣的樣子呢.
可能是三個頭六隻腳.
就是一個怪人來的了.
每個人都會將一些自己不喜歡的東西.
就投射在一些你不喜歡的領袖身上.
那黎希米就在第六章就明顯.
那些人就設個陷害他.
邀請他進去聖殿避難.
因為那些風聲已經傳到黃那裡.
現在我是幫你的.
你去到聖殿避難.
接著就說他自己一個平民百姓走進聖殿.
再加上就是很多的誣衊.
很多的匿名的告密信.
這樣去針對.
變成是無從辯解.
也就是想挫你的銳氣.
和你的信心.
希望你能夠停工.
我們看到的就是在四五六章.
其實由第二章都已經開始.
其實應該是第一章開始.
每一章都散發在每一個的事工.
每一個的困難裡面.
都是我的神求你幫助我.
我的神求你施恩的手看顧我.

$^{121}$都是這樣不住在尼希米的祈禱裡面.
當然這個是一個寫作的手法.
那就是說意味著他每天起床都會祈禱.
今天我的行事曆.
我規劃到晚上之後.
他就去反思今天我做過什麼.
我遇到什麼事情.
我有什麼值得感恩的事情.
有什麼我可以做得再好一點的事情.
每天這樣去反思.
其實都很重要的.
如果大家有用行事曆.
網上的或者是一些軟件.
其實都很好用的.
幫你去做規劃.
你每天就不用放在腦袋裡面.
你記住所有的東西.
但是你有這樣的機會去反思的時候.
將來你能夠看回.
差不多等於你每天的日記.
一些札記.
這樣很重要的.
尼希米記頭六章.
根本上就是她的個人回憶錄.
第七章就是一個很長的家譜.
再一次來數點百姓.
再一次來盤點一下.
在耶路撒冷到底有多少人.
原來有很多是散居在其他地方以色列的.
這樣的準備.
其實就是為第八章這個水門事件.
水門的不是醜聞.
不是尼克遜的醜聞.
就是這個水門處一個寬闊處.
一個集會的地方.
也是叫做Watergate.
水門.
就是一個復興.
這裡有三張聖經.
八,九,十三張聖經.

$^{161}$那個條例都很簡單很清晰的.
就是當領袖被邀請.
因為當時他們有的材料就是摩西五經.
就是律法書.
根本不會有新約.
根本都沒有其他的經卷等等.
特別是馬拉基書可能都沒有.
都還沒有出現.
可能撒迦利亞書都有.
因為都是那個時候重建的.
你可以想像得到.
他們就是拿著摩西五經來研讀.
去中讀.
有可能是讀利未記也說不定.
來讀的時候.
百姓就受感動.
不知道他讀多長.
讀哪一卷.
就帶來他們有一個集體的凝聚.
就是為什麼過去這麼有恩典.
特別是從亞伯拉罕和摩西.
兩位的領袖.
作為主要的悔改對象.
就是說應該要回到那個時候.
神賜福我們的民族.
我們希伯來人這個民族.
為什麼今天我們要被擄.
現在回來之後.
又受周邊鄰居的欺凌.
最主要有一個很大的問題.
在這三章聖經.
他們開始發現自己失落了.
那個sense of belonging.
歸屬感.
對自己身份的歸屬.
對自己身份的認同.
個人都已經少了.
更何況是一個群體.
所以這三章的復興.
就是回到這個身份的認同.

$^{201}$有這個身份認同.
很多事情就容易辦了.
所以第八章開始.
我們又再重新比較認真去看這幾章.
可能今天集中是第九章.
因為第八章我們都講過一些.
記住是七月是一個重要的猶太人的節期.
三個重要節期.
七月初一是催國節.
七月初十是贖罪日.
是一天的.
全年最重要的一天.
所以不是新年.
不是逾越節.
最重要是七月初十的贖罪日.
接著就是七月十五.
一直到七月二十一的.
就是住棚節.
等下第十章會有講住棚節.
所以七月是他們最重要的日子.
所以六月完成整場.
給他們有這個機會.
準備他們的心靈.
開始在催國節的時候.
他們就不只是一個儀式.
而是有一個嚴肅會.
大家都很想聽一下.
神既然招聚了我們這班人.
有了我們的四面牆.
那想我們這個民族.
現在我們這個大的群體.
在聖地.
神賜給我們的土地.
想我們又有什麼作為呢.
到了七月.
以色列人住在自己的城裡.
那時他們如同一人聚集在水門前的寬闊處.
就是這個Watergate.
請民是以色列將耶和華.
藉摩西傳給以色列人的律法書帶來.

$^{241}$似乎摩西已經帶來了.
記住帶那些卷軸很麻煩.
因為那卷軸是整卷.
真的是卷軸.
是很神聖.
因為是手抄的.
而且是相當重的.
所以不是那麼容易隨便帶出來的.
這裡告訴我們.
到了七月初一.
果然是這個催國節.
祭司以色列將律法書帶到聽了能夠明白的男女會眾面前.
聽了能夠明白.
就是說懂事的.
不是那些手抄的BB.
就是說他們真的傳會眾都認真能夠聽得明白.
當然是用希伯來話.
這個可能有一點技術性.
因為被擄歸回的人.
可能未必懂得希伯來話的方言.
可能他們已經講慣了一般的亞蘭話.
就是說可能你廣東話很厲害.
但是可能書寫就用那些書面語.
未必每個人都懂得.
所以這個可能也需要翻譯.
所以有很多利未人.
稍後就會來做那些組長.
來做一個傳譯.
或者是將那些講完的經卷.
來到去再可能小組的研習.
來告訴會友.
這裡讀下去你們會看到.
第三節.
在水門前的寬闊處.
從清早到響午到十二點.
在眾男女一切聽了能夠明白的人面前.
讀者律法書.
眾民則以而聽.
是單向的.
只是宣讀.

$^{281}$還未解釋.
當然不是不解釋.
所以之後的研習班就會解釋.
第四節.
上文就在講這個祭司以斯拉.
現在又再講文士以斯拉.
就是說兩個身份都兼備.
我們在以斯拉記第七章已經介紹過這個以斯拉.
就是說.
他是有機會見識過.
他自己第一次上去的時候.
現在去到尼希米.
可能過了二十多年了.
尼希米又上來重建城牆.
這次又帶動一個復興.
所以尼希米他做完那些硬件之後.
其實重點就是能夠配合.
神感動以斯拉祭司.
已經是具有身份作為教導.
作為屬靈領袖.
所以一個尼希米就是.
這個重建的總工程師.
是亞特薩西王.
差遣他.
當然尼希米有他的感動.
這樣上去.
但原來在耶路撒冷.
已經具備了.
有一定的屬靈的制度.
就是祭司以斯拉.
其實一直都在耶路撒冷裡面.
做著零零星星的.
可能每個主日的講道.
他們是星期六 安息日.
這樣去做訓練.
教導 工作等等.
但現在有一個很好的契機.
就是大家工作完成了.
大家很想享受一下.
工作完成之後的成果.

$^{321}$就是開嚴肅會.
是開一個.
研習神話語的機會的聚會.
第四節.
文是以斯拉站在為者是.
特備的木台上.
即是我們的講台.
接著就是瑪他提瓦 士瑪 阿奈瓦.
烏尼瓦 希臘加 和馬西瓦.
站在他的右邊.
比大瓦 米沙利 瑪基瓦 哈遜.
哈拔大拿 薩加利亞 和米舒蘭.
站在他的左邊.
特別強調這些的領袖.
相信因為這次的復興.
是帶動由祭司層級的階層.
這樣從上而下.
即是在他身邊.
現在你能夠數到有12個.
就是左右兩邊.
即是一併的列隊排成一列.
中間的就是講台上的以斯拉.
來宣讀神的話語 律法書.
表示這個復興.
首先是神的工作.
先是感動祭司利未人.
他們這個階層.
來到去很想既然我們重建好城牆.
重點仍然是城牆裡面.
敬拜神的這班選民.
神的選民 以色列人.
所以其實每一個來教會的.
都是神想感動你的工作的.
內心的重要的內容.
就不是那個硬件.
所以那個硬件美輪美奐.
但是如果人心不歸向神.
群體沒有自己那個歸屬感.
那個sense of belonging.
一切都徒然.

$^{361}$硬件神隨時可以不要.
所以有些教會如果真的沒有硬件.
沒有就沒有.
我們這班人才是最重要.
我們這班人才是神的群體.
這些有時很抽象.
你看現在你見到.
他們準備好了.
即是說在一切.
原來第八章才是真正.
尼希米基的那個中心思想.
一到第六章其實是一個回憶錄.
一個準備.
第七章就是數點文數.
有什麼人 他們是什麼身份.
是什麼支派等等.
這樣數完之後.
再一次來做一個盤點.
跟著就是一個嚴肅會的開始.
所以第十章之後.
就是他們大家願意.
來簽約行摩西的律法書的心情.
第五節.
以色列站在眾民以上.
在眾民眼前展開這書.
他一展開.
眾民就都站起來.
即是意味著可能大家坐下.
好像五餅二魚那樣.
大家坐下很安靜的.
一排一排的.
來等待.
都不理那些天氣等等.
因為遲些有些變化.
現在一見到神的律法書展開.
大家都站起來.
這些可以是一個.
慢慢習以為常的一個儀式.
一定會的.
人是很想將那樣東西.

$^{401}$習慣化.
開始可能是很好的.
但慢慢慢慢習以為常.
就變成一個形式.
例如就是說.
今天還未祈禱.
事奉人員還未祈禱.
很緊張很緊張.
馬上祈禱.
其實大家都知道.
只不過是一個儀式.
不是今天.
真的不是那麼重要.
這個是慢慢一路一路遞減的.
無論你起初多麼有心.
要這樣設立.
慢慢慢慢都習以為常.
所以要不斷深化.
我們背《主禱文》.
或者我們背《小圖信經》.
不僅僅是一個儀式.
你會見到是一個儀式.
因為有弟兄姊妹.
你如果不提她.
她又插袋.
又這樣.
完全口不對心.
你就會見到.
其實人就是這樣.
慢慢慢慢會習以為常.
你怎樣能夠讓他.
在一生他能夠看到.
這些那麼重要的真理.
是承載每一個基督徒的生命的.
那個信仰.
那個不是信條.
是信仰.
是真是我們整個人.
整個生命的皈依.
就是因為這些那麼重要.

$^{441}$不是一些格言那麼簡單.
是濃縮了的.
這些真理.
透過《聖經》的啟示.
我們濃縮了.
我們知道.
為什麼我們知道我們會有永生.
因為《聖經》說了.
真是這位主耶穌.
會再回來審判活人死人.
我們相信我們能夠有永生.
我們每個禮拜這樣背.
其實是要讓我們不斷去.
確認我們自己的信仰.
所以你可以想像到.
如果我們慢慢變成一個形式.
我們必然會這樣做.
我們只能夠盡量.
我們自己有生之年.
是要令到那樣東西不要急速的惡化.
和要維持那個屬靈的生命力.
和更重要的是.
我們要傳遞給下一代.
下一代就更加遞減了.
再第三代第四代.
更加遞減的.
除非有些大的事情.
讓他們有一個大的復興.
一路一路就會變成例行公事.
歐洲人慢慢現在他們.
因為世俗的個人主義.
已經強烈到沒有東西.
可以干預他們做任何的事.
包括信仰都沒有那個力量.
所以去到教會.
只能夠甚至乎都不再去.
可能如果天主教徒出世.
結婚死.
才和教會有關係.
其餘時間教會和他們.

$^{481}$基本上就沒有關係.
這是必然會一路一路出現的問題.
所以現在我們透過這個復興.
又讓我們可以有些輪廓.
去看看人心.
怎樣可以被神的靈去感動.
而來到去有那個被挽回.
仍然就是他們回到.
他們自己這個身份的認同.
沒錯.
他們是以色列人.
有身份證.
真是希伯來人.
真是說亞蘭話的.
真是他們有時偶然間.
他們在家可能找到律法書的.
但是仍然那個問題就是說.
他們對自己這個.
很榮耀的神選民的身份.
就不是一個很強烈的認同.
如果不是.
他們就不會安息日.
就根本不介意.
外邦人來買賣.
他們就很喜歡這樣去買賣.
出現這些情況.
我們讀下去.
我們就一定看到的.
第六節就是.
以色列的稱頌和祈禱.
以色列稱頌耶和華.
至大的神.
眾民都舉手應聲說.
阿們 阿們.
一個和應.
就低頭面伏於地.
敬拜耶和華.
接著耶穌亞巴尼.
史利比 雅文 阿谷 沙比泰.
荷蒂雅 馬西雅 基利他.

$^{521}$阿薩利亞 約撒伯 哈蘭 比萊雅.
和利未人.
使得百姓明白律法.
剛才我說了.
需要翻譯.
就是說意味著.
不是就這樣讀完敬拜.
是以色列可能用了一個多兩個小時.
將律法書有系統地.
慢慢地 大聲地.
因為是一個聚集的地方.
沒有咪高峰的.
大聲地宣讀.
宣讀完之後.
就分小組的研習.
利未人將剛才讀的經卷.
有意識地解說.
可能要做翻譯.
有些聽不明白.
以斯拉這位祭司.
這位文士.
剛才用希伯來話讀的.
他們聽不明白.
利未人就幫他們翻譯.
做一個他們可以聽得明白的.
重點仍然是聽得明白.
明白所聽的律法書.
才會帶來一個影響.
聽得明白是重要的.
因為是一個溝通.
所以港台是要講明的.
又要給一個悉切的形容.
諸如此類.
這裡我再讀落馬路下面.
你可以想像到.
一而再提這麼多利未人的名稱.
這一次是很有意識地.
不是Ed Hart臨時的即興.
是有計劃地準備好的.
即是說.

$^{561}$黎希米和以斯拉聚集的大會.
有很多的籌備.
有很多的邀請.
有很多的禱告.
預計可能會分組.
來去研習.
才會有一個傳聞.
帶來的復興.
所以你會留意到.
接著就是百姓.
這是第七節的下半節.
百姓都站在自己的地方.
他們清清楚楚地念.
神的律法書.
講明白意思.
使得百姓明白所念的.
清清楚楚.
講明白.
所以百姓的接收.
不是聽不明白.
更加不知道他講什麼.
更加是又小聲.
又不知道他講的語言.
是和我們有相近.
那就變成完全聽不明白就沒意思了.
現在是要強調.
大家都能夠聽清楚明白.
明白心靈能夠開始有一些.
神的道在心裡的醞釀.
不只是讀了那句經文.
而是我的牧者.
大家讀完就算了.
不要去思索.
到底牧者給我們什麼意義呢.
在什麼處境牧者是我的牧者呢.
這個牧羊.
我是不是只看自己是羊就夠了.
一定有很多處境的一個.
研習.
分析.

$^{601}$形容.
使得大家都能夠明白.
第九節.
省長尼希米和作祭司的文士以色拉.
並教訓百姓的利未人.
所以是很清楚.
是有意識.
有準備的一個召會.
呼召的一個大會.
有這個水門的嚴肅會.
水門的一個研習的茶經.
讀經日營.
這裡差不多是全日的.
接著就第九節的下半節.
他們全體對眾民說.
今日是耶和華你們神的聖日.
是摧國節.
不要悲哀哭泣.
因為這是因為眾民聽見.
律法書上的話就都哭了.
聽見甚麼律法書而哭呢.
有甚麼令他們感懷身世地哭呢.
是一種開心的哭.
還是發現自己的身份和聖經有一個距離.
而來到有一個內心的震撼.
這個讀下去就知道了.
反而這裡就告訴我們.
似乎現在去到吃中午飯的時間.
開始他們感覺到現場的氣氛.
是有些不尋常的哭泣.
這個哭泣就要有一個出路.
這個出路要怎樣呢.
這個領袖就開始在呼喚.
第十節.
又對他們說.
你們去吃肥美的喝甘甜的.
有不能夠預備的就分給他.
因為今日是我們主的聖日.
你們不要憂愁.
因靠耶和華而得的喜樂.

$^{641}$是你們的力量.
這是金句.
利未基的金句.
因靠耶和華而得的喜樂是力量.
在這個時候他們的哭泣.
他們的哭泣不是因為沙塵滾滾.
有沙入眼.
是因為聽.
講了兩次了.
領袖利未人.
是清清楚楚的念神的律法書.
不講其他東西.
不是一個開玩笑.
很開心.
我們大家歌功頌德.
誰誰做誰的成長.
我們在這裡數算.
大家一起感恩.
不是一個感恩會.
大家願意聆聽神的律法書.
祭司利未人在這裡宣讀.
宣讀完之後.
有意識地講解清楚明白.
記住是百姓的心靈能夠明白.
明白之後.
他們就開始有一個內心的觸動.
其實《小徒行傳》第二章就是了.
彼得講完十個故事之後.
其實很平鋪直敘.
講到你們所釘死的耶穌.
神已經立祂為主為基督.
眾人三千的南丁.
在各處上來的首節.
首五場節的這些以色列的南丁.
就突然間內心窄心.
就問使徒.
不只是問彼得.
是問十二使徒.
現在我們應該怎麼做呢.
彼得就說你們各人要悔改.

$^{681}$奉耶穌基督的名受洗.
就叫你們的罪得了赦免.
和得到神在過去.
一路應許要賜給我們的.
這一位的聖靈.
住在我們的心裡面.
所以是很清楚的.
聽明了彼得的講道.
就窄心.
其實彼得很平鋪直敘.
至少他有勇氣.
這個曾經三次說謊.
否認耶穌的彼得.
現在就站在眾人面前.
不是那些人不認識他.
還有很多耶路撒冷人.
是還沒死的.
不是過了很多年.
去了很遠的地方.
譬如你在香港傳道.
現在你去了很遠.
去到新西蘭.
沒有人認識你的.
那你說什麼都沒有人知道.
你過去的事情.
但現在彼得只不過是.
耶穌四十天的顯現.
再加十天.
接著就升天.
現在是五十天之後的彼得.
曾經否認耶穌的彼得.
現在站在眾人面前.
來講述.
所以他有一個很大的.
悔改的信心.
被神重新建立的.
屬靈的位份.
屬靈的影響力.
來將十誡的故事.
他自己至少不會很悍然地.

$^{721}$再講述.
他現在真的將主耶穌.
這位榮耀的主.
介紹給上來首五巡節的.
以色列的男丁去聽.
現在以色列也是.
他們很清楚地讀完之後.
百姓全體都受到觸動.
內心有一種悲哀.
就是為什麼我們的民族.
現在我們去到一個.
這麼低的層次的.
那個感受的懷念.
就是那個歸屬.
自己身份的歸屬.
為什麼會這麼低呢.
如果我們真的做一個.
香港的調查.
基督徒你敢不敢承認.
自己是基督徒.
為主耶穌.
覺得是一份榮耀.
等等諸如此類.
可能那個.
sense of belonging.
是很弱的.
你能不能夠對自己.
是哪間教會.
你引以為榮呢.
可能又再低一點.
你能不能夠承認.
我是一個很commit的基督徒.
我又不單止有奉獻.
我有事奉.
如果主耶穌要我.
差遣我.
我都願意去哪個地方服侍的.
我想又再低一點.
這個就是.
那個身份認同的問題.

$^{761}$所以他現在.
全個會眾.
在有了自己的城牆之後.
開始他們的內心.
似乎很想.
和神有一個.
再近一點的關係.
所以這個嚴肅會.
是一個很好的契機.
所以.
大家都有了自己的群體.
你都不去做這個群體的.
一些再嚴肅一點的要求.
還是很lose.
你有沒有讀聖經靈修.
不要緊的.
最重要是你回來幫忙.
你有沒有回教會.
都不要緊的.
總之輪到你做私事.
你要回來.
你要領唱.
你一定要回來.
你另外三個禮拜不回教會.
不崇拜不要緊的.
我們是不是這麼低的.
sense of belonging.
那個歸屬感.
如果是的話.
似乎我不太想.
好像我借題發揮.
自相罵槐.
但是我自己.
每個星期.
我三堂.
如果不用教主一學.
早堂我就不在了.
我要教主一學.
教朕禮班.
真的作為堂主任.

$^{801}$我有這個肩負的責任.
是要勸導.
我沒有權去篤責.
現在這個世代.
即便是篤責.
其實都是.
用愛心說誠實話.
我的目測.
不用事奉的.
有時我真的見不到祂.
對不起.
事奉的那個禮拜祂在.
這些是要.
嚴肅的討論的事.
當然.
我不是神.
祂原來其他那幾個禮拜.
去另外的教會.
或者有另外的原因.
是我完全不知道.
至少要告訴我.
讓我不要有誤解.
或者誤會.
否則.
這正正就是.
可以看到我們教會.
那個淑寧的探熱針.
那個溫度.
這個就是.
彼此之間.
大家的身份歸屬.
記住是對神的兒女.
身份歸屬是最高的.
接著就是.
你有神的兒女身份.
當然你有地上的一個教會.
是一個很實在.
你能夠摸到那張椅子.
你坐下去.
那個位置.

$^{841}$是每個禮拜.
你的出現.
我們不是已經去到靈界.
我們不是.
雲遊在一個三層天.
我們每個禮拜實實在在.
是要從家裡走來.
這個教堂.
所以.
沒有假的.
必須是.
我們要對這個身份認同.
就一定是由於我們是神兒女.
這個最崇高的身份.
被十字架的愛拯救.
然後我們就有這個地上的身份的歸屬.
我們是屬於這個教堂.
教會的.
我覺得每個禮拜.
就要承諾在這個地方.
來到去.
有那個參與.
有那個事奉.
有那個敬拜.
等等.
這些已經是.
放諸於四海之內.
都已經是很重要的.
一個真理.
如果連這些都沒有.
就要問.
我們的歸屬性.
所以這裡開始告訴我們.
他們就是要因靠耶和華而得的喜樂.
就是他們的力量.
靠神而來的喜樂.
所以現在是一個月立的時候.
神現在不是要來追究他們的罪的時候.
現在是要將這個民族.
再一次帶回去一個.

$^{881}$好像摩西時代.
那種榮耀的身份的時代.
好像是剛剛出埃及的時候.
那個身份.
神很榮耀地宣告.
這個民族是屬於我的.
我是將他們帶來歸我自己的.
十九章的出埃及記.
那裡.
頒佈十誡之前.
有一段這樣的說話.
好像那些麻鷹.
將他們背負著來歸我的.
所以現在.
當他們開始有心靈的觸動的時候.
以色列就提醒他們.
現在這個時候.
不是一個哭泣的時候.
是一個要喜樂的時候.
因為神會有那個恩典.
來為我們這個民族準備.
所以在這個時候.
就是放飯了.
他們就要彼此來到很開心地.
享用那些食物.
就是說下午.
再有另外的聚會.
所以這裡.
那個神給我們的喜樂.
才是真正的心靈的力量.
這個真的要求主給我們.
每個人都要經歷.
靠神而來的喜樂.
就是說是神所閱立的.
是神所賜予的福氣.
我們是很體重的.
我們來到去領受了.
這個是我們的出路.
是我們的力量.
因為神已經將我們收納.

$^{921}$我們不是被撇棄了的一班人.
所以我們要求主.
讓我們能夠有這種.
很健康的眼淚.
這個健康的眼淚.
就是因為被神的愛.
神的真理.
去感動了我們的心.
那我們就怎麼辦呢.
祭司說.
以色列說.
不要哭泣.
因為這個時候是神私恩的時候.
你回到神的話語裡面.
你開始對自己這個身份.
有一個很強烈的認同.
這個就是神的時間.
工作的時間.
後面我們就繼續會見到.
他們自己主動去守節旗.
那你可想而知.
是多麼的寶貴.
不用求 不用哄.
不用威逼利誘.
沒有任何獎品.
大家都很想去守節.
接著就說.
今天是聖日.
不要作聲.
也不要憂愁.
眾民都去吃喝.
也分給人.
大大快樂.
因為他們明白.
所教訓他們的話.
第三次提到.
他們.
不是祭司.
祭司固然明白.
自己要準備.

$^{961}$是那些百姓.
明白了自己的身份.
自己在神的話語裡面.
原來神這樣看我們.
我們不需要這樣.
和外邦人這樣.
來到一般的見識.
我們應該有神.
給我們這個民族.
我們有我們那份榮耀.
所以如果珍惜自己.
很寶貴自己.
基督徒的遺訓.
基督徒的身份.
你能夠在這個遺訓裡面.
能夠常常與主有相交.
來看到自己.
那種內心的平安.
喜樂.
日於言表的.
這個就是基督徒.
我們應該要有的.
對自己身份的歸屬.
第二天.
次日.
眾民的族長.
祭司和利未人.
都聚集到民事.
以斯拉那裡.
要留心聽律法上的話.
他們見律法上寫著.
耶和華籍摩西.
吩咐以色列人.
要在七月節住棚.
即是快了.
並要在郭城和耶路撒冷.
宣傳報告說.
你們當上山.
將橄欖樹.
野橄欖樹.

$^{1001}$番石榴樹.
棕樹.
和各樣茂密樹的枝子取來.
照著所寫的答行.
於是百姓出去.
你可以想而知.
這是一個濃縮了的.
不是每天的日記.
不是一個札記.
是一個濃縮了.
即是說他們.
七月初一.
中午放飯之後.
他們很開心地去用飯.
之後就不再記這麼仔細了.
接著再聚集的時候.
就是第二天了.
接著又再有那個嚴肅會了.
接著這次是領袖.
利斯 利未人.
都聚集到以色列的面前.
來留心地去聽.
接著他們自己也發現了.
其實好像很尷尬.
這些都不用說了.
因為一定是那個節期都不用說.
一定是要守的.
即是說這次是極之有意識地去守.
大家是很踴躍地去守.
不是例行公事.
今天年三十晚了.
明天年初一了.
照例都要跳下輝春了.
如果是一些可能是從外國回來的人.
他越來越意識到.
原來自己中國人守這個年三十晚.
去做那些事情.
和媽媽一起做湯圓 做餃子.
原來是這麼開心的事.
他真的對自己是中國人這個身份.

$^{1041}$那個強烈的認同.
和很渴慕有壓歲前.
初一就講功績發財.
他有這個很強烈的意識.
很想快點到年初一.
你可想而知.
這個就是猶太人.
他們這一刻是他們民間.
很想守住棚節.
不是一個 只是一個的習俗.
而是好像是從來都沒有過的.
有我們自己四面圍牆的住棚節.
是屬於只是我們這個民族.
我們神選民的住棚節.
其實是失落了千多年.
摩西在利未記頒佈.
神的七個節旗.
你們進到迦南之後就要遵守了.
然後聖經就很有趣的.
他不會逐個逐年逐年.
告訴我們他們怎麼守的.
那意味著進到迦南之後.
可能開始還會有做的.
但做著做著一兩代之後.
就真的變成例行公事了.
為做而做了.
不做都不行了.
諸如此類了.
那就變成慢慢慢慢.
不感覺到住棚帶來什麼興奮.
有什麼屬靈的含義了.
因為住棚就是要紀念他們在埃及寄居.
你們是奴隸.
就是說時時要告訴自己.
你是奴隸.
現在你在以色列安居樂業了.
記住你爺爺你的太公是奴隸.
是神將我們拯救.
所以我們每年的7月15日到7月21日.
我們要守住棚節.

$^{1081}$不是因為我們很想去露營.
去看星星.
去大自然那裡.
讓小朋友去看大自然.
是因為你是奴隸.
在利未記講得很清楚.
拆開說明.
今年的培靈會.
我很期待.
因為連明仔都很期待.
周曉輝牧師.
北宣的堂主任.
今年會講利未記.
上一次來很精彩.
我們主日聚會講.
他會講利未記.
所以大家期待8月1日到10日.
的培靈研經會.
很少人會講利未記.
講得好.
你可以想像到.
他們主動.
渴慕守住棚節.
告訴祭司二師奶.
可想而知.
有時我們作為牧者.
很等待.
等姐妹幾時自己好自發.
好想聽多些聖經.
因為你可以想像到那個設計.
主日講台上.
講超過40分鐘.
大家已經受不了.
現代人已經不能集中精神.
還有其他環節.
變得很緊湊.
有時我最急的是.
講完你講快些.
因為還有很多事要後面等等.
要報告 要牧土.

$^{1121}$有時最急的是我.
其實會眾應該.
如果大家不介意.
聽久些.
當然有各樣原因.
譬如坐得太久.
或者完了又趕著有其他事.
這是明白的.
所以主日只能夠.
一定在最有限的時間.
去編排每一個聚會.
有不同的儀節.
但是重點仍然在這裡.
告訴我們.
原來有系統地.
有時間.
有心靈的準備.
去有一些的嚴肅會.
其實是很大的.
對我們的生命成長.
有一個很強烈的需要.
現代人其實.
應該最需要就是休息.
睡眠.
或者是七月暑假.
和家人去旅行.
這些是強烈的需要.
相信錦書家的弟兄姊妹.
很多都已經計劃好.
要去旅行.
但是強烈很想在.
比較長一點的時間.
自己讀一卷書.
或者如果有一些很大的聚會.
我能夠在這些聚會裡面.
認真的去研讀.
其實這個是.
每個人內心基督徒的一個呼聲.
不過這個呼聲.
可能已經被我們太忙碌的.

$^{1161}$日常工作已經壓低了聲音.
所以有時候就是.
聖靈需要很大的工作.
才能夠喚起我們.
內心的呼喚.
內心對神話語的渴慕.
那個呼喚.
所以接著他們就主動跟.
以色列說.
告訴我們第十六節.
於是百姓出去.
取了柴之來.
各人在自己的房頂上.
或院內或神殿的院內.
或水門的寬闊處.
或以法連門的寬闊處.
搭棚.
可以在附近搭棚.
所以你猜七月.
有時候真的鼓勵你.
飛去以色列.
你看看人家猶太人他們駐棚節.
當然他們就是一個螢幕.
這樣疊出來.
有沒有再追溯.
那個.
剛才說的那個原因.
相信不是太多猶太人.
已經這樣想.
已經不會再覺得自己是.
奴隸.
駐棚那種痛苦日子.
要追本族元.
想到自己這個民族的起源.
就是神將我們.
拯救離開這個.
叫為奴之家.
奴隸之家.
很少你將自己一些太公的相.
拍張相出來.

$^{1201}$那個是奴隸來的.
來紀念我們很少的.
應該是哪個博士哪個名人.
我們就很想跟他有關係.
但是.
你的爺爺是太公.
你的太公是奴隸.
我們盡量都不希望被人知道.
但是神就叫這個民族.
如果你真的想健康.
你要記住.
守住棚節.
是歡樂的.
歡樂的背後是因為.
你們這個民族曾經是被壓制.
本來根本是一無所有的.
是埃及人的.
那些妹仔.
做勞工.
在柴房這樣來做.
做那些.
那些草的磚.
但是現在.
神將我們抬舉.
將這個這麼好的.
奴隸與物之地賜給我們.
所以在這一刻.
就是他們自發的守住棚節.
不用大力的推廣.
已經他們自己去做.
所以你會見到.
第十七節.
從老道之地歸回的.
全會眾.
就踏行住在棚裡.
很誇張.
從亂的兒子約書亞的時候.
直到這一天.
以色列人沒有這樣行.
於是眾人大大喜樂.

$^{1241}$從頭一天直到末一天.
7月15日到7月21日.
七天.
第22日第8日.
就是嚴肅會.
以色列每日念神的律法書.
眾人守節七日.
第8日照例有嚴肅會.
每日都念神的律法書.
來勸勉.
教導他們.
他們不只是一種.
歡樂的踏行.
而是知道節期.
背後的含義.
所以今天.
我們連耶和華的七個節期.
我們都數不到.
數手指都數不到.
我們基督徒.
其實每個月.
因為我們不一定要計那個次數.
因為聖餐會是每個星期.
我們就每個月.
每個月我們都在聖餐.
來做一次.
深化我們自己身份的認同.
Sense of belonging.
這個很重要.
我們要再一次.
因為十價.
因為這個餅這個杯.
曾經代表著主耶穌的犧牲.
再一次和我們生命的關係.
每個月.
讓你去沉澱.
這個生命.
因為這樣而來的快樂.
因為這樣而來的.
嚴肅.

$^{1281}$是因為我們明白了.
律法書.
對自己這個信仰.
完全生命的認同.
和你的皈依.
所以我們今天.
當然.
香港的基督徒是很少.
真的佔全個香港的.
七百八十萬人.
那個人口.
那個比例.
可能二十萬都還沒有到.
可能已經計算了所有天主教.
基督教甚至東正教.
甚至有些人計算了.
摩明教諸如此類.
政府的官方.
但是真正commit的有多少呢.
真正願意對自己的身份.
有強烈的意識.
有強烈的.
願意來遵守神的話語.
這樣的.
基督徒.
這樣的家庭就更加少.
現在這裡就出現一個復興的開始.
就是.
靜心.
認真去聆聽.
透過.
有預備的一個聚集.
一個水門的.
復興 奮興會.
來去聆聽.
上主藉著祭司.
利未人的講解.
去對舊約的研讀.
摩西五經.
真是最重要的猶太人的.

$^{1321}$生命的藍本.
重要的生命的教材.
就是律法.
就是律法背後.
是神期望這個民族.
那個behavior.
你的品格應該是怎樣.
所以十誡是擺得這麼重.
就是近乎將神的屬性.
來和這個民族.
拉上一個直接關係.
所以和其他民族.
是完全不相同的.
第九章.
這一月的二十四日.
去到住棚節的.
過了住棚節的.
宴宿會之後.
有可能去到第二個的安息日.
因為.
可以有一個特別的安息日.
現在可能是另一次的安息日.
以色列人聚集.
禁食.
吃螞蟻.
頭蒙灰塵.
以色列人就如一切外邦人離絕.
暫著承認.
自己的罪惡.
和列祖的罪孽.
那日的四分之一.
站在自己的地方.
念耶和華他們神的律法書.
由四分之一認罪.
敬拜.
就是說有.
聖經的教導.
研習的時間.
是將分成.
全日的分四節.

$^{1361}$可能是這樣.
接著有另外一節.
讀完聖經.
大家去默想.
去反思去悔改.
去認罪.
敬拜耶和華他們的神.
第四節.
又再一次說那些利未人.
耶穌亞巴尼革滅.
是巴尼布尼.
是利比巴尼基拿尼.
站在利未人的台上.
大聲哀求.
耶和華他們的神.
利未人耶穌亞革滅.
巴尼哈沙尼.
是利比荷第亞.
是巴尼比他希雅說.
你們要站起來.
稱頌耶和華你們的神.
永世無盡.
耶和華.
你榮耀之名是應當稱頌的.
招呼一切.
稱頌和讚美.
利未人帶領著大家.
敬拜.
大家唱詩歌.
詩篇來讚頌神.
現在分了.
有讀經的時間.
聽道的時間.
四分之一.
大家來.
集體性認罪.
然後他們.
有一齊敬拜.
可能用詩篇來敬拜神.
利未人站在.

$^{1401}$桌子上.
好像是詩班的.
領唱帶動敬拜.
雖然沒有說到.
但是大衛的年間.
已經設立了.
二十四班的.
詩班長的詩班.
即是說他們是輪流的.
有不同的詩班.
來做.
帶領.
會眾來唱詩篇.
或者代表著.
百姓.
來頌揚神.
第六節開始.
開始有認罪的禱告.
收錄在.
第九章.
到底是誰.
帶動的祈禱呢.
是以色拉.
還是尼欽米.
我們看看.
唯獨你是耶和華.
你做了天和天上的天.
並天上的萬象.
這是創世記的經文.
這是對.
創世記的深信不疑的.
一個宣告.
認同.
做物主一開始.
你就是做物主.
所以權柄.
一定是從你那裡而來.
因為唯有你去創造萬有.
唯有你是.
自有永有的存在.

$^{1441}$你是無始無終.
沒有開始.
沒有人創造你.
你是本身自我存在.
自我去啟示.
啟示給我們知道.
所以我們這個民族.
能夠現在這樣宣認.
是很偉大的.
所以我們的下一代.
盡快要對創世記有認識.
對做物主.
他掌管萬有.
從小就要在他的心底.
要有這個價值觀.
他長大一定面對很多衝擊.
很多的事情.
但是如果他知道.
做物主最大最大的.
那一位就是他的天父.
那個心靈就很平安.
很能夠穩妥.
將來他面對各樣的洪水.
很多的衝擊.
他都能夠有一個很穩妥的心靈.
不是窮得只剩下錢.
他內心真的有這個撈.
拋撈在榮耀.
在榮耀.
拋撈在永恆裡面.
和神有這個.
最重要的關係.
就是認他是做物主.
所以地和地上的萬物.
海和海中所有.
這一切都是你所保存的.
天君也都敬拜你.
你是耶和華神.
曾簡選.
開始和人有關.

$^{1481}$開始和我們的老祖宗.
亞伯蘭有關.
所以一掃珠.
我們希伯來人.
就是由大河.
即是右法拉底河.
還是底格里斯河.
兩個大河.
那裡發源地.
輾轉那裡被簡選.
來到這個.
以色列地的就是亞伯蘭.
所以.
亞伯拉罕後來就改名.
你.
唯獨你耶和華神.
曾經簡選亞伯蘭.
令他離開了這個迦勒底的吾爾.
給他改名叫.
亞伯拉罕.
是多國之父的意思.
你見他在你面前.
心裡誠實.
就與他立約.
應許.
把迦南人,黑人,阿摩尼人,比利西人,.
耶布斯人,甲加薩人.
即是這個叫迦南.
七族.
即是這些比較大的.
迦南的土著.
比較大的國家.
的國民.
所以很容易背的.
就是迦南七族.
將他們的地.
賜給他的後裔.
就是亞伯蘭的後裔.
而且應驗了你的說話.
因為你是公義的.

$^{1521}$你曾經看見.
我們的列祖在埃及.
將歷史突然間由亞伯蘭.
立即就跳到去.
這個民族的起源.
亞伯蘭的立約.
神與他立約.
接著神就再與.
摩西立約.
這個民族.
就是我的子民.
從以色列人立約.
是什麼約的記號呢.
就是安息日.
與亞伯蘭立約.
是什麼記號呢.
上個創世記的.
就可能知道了.
就是葛禮.
以色列人是安息日.
作為那個記號.
這裡告訴我們.
神啊你曾經在紅海.
誰聽他們在紅海邊的哀求.
就施行神跡奇事.
在法老.
和他一切的神僕.
一個認罪祈禱.
先是講一講.
想當年.
當年我們是怎樣開始.
這個民族的起源.
就是由一個叫亞伯蘭.
神你在那個.
吾爾的地方揀選他.
他就有信心跟隨.
一路輾轉.
去到迦南地.
落了地寄居.
後來你幫他改名叫亞伯拉罕.

$^{1561}$接著由這個民族.
由這個人的後裔.
就產生十二個先祖.
其實這些是斯特恆傳.
斯特反的.
正道.
所以現在.
這個.
祈禱的人他還沒講是誰.
就開始由.
亞伯蘭一跳就跳到去.
摩西因為就是這個民族.
的起源.
第十節.
的下半節.
他國中的眾民身上.
你也得了名聲.
因為殺了埃及人.
正如今天一樣.
因為你知道他們向我們列祖.
行事狂傲.
你又在我們列祖.
面前把海分開.
使他們在海中.
行走乾地.
這個就是紅海的分開.
將追趕他們的人.
拋在深海.
記住這些的經文.
正正就是.
印證著出埃及記.
的紅海十五章.
十二十三章的紅海事件.
是真有其事的.
因為是過了很久.
很多年之後.
他們這樣去描述.
所以有些人就去解.
出埃及記新派的學者.
去解的時候.

$^{1601}$那個紅海其實原文叫盧偉海.
盧偉海就不是深海.
其實是淺水.
所以那些埃及人.
所謂浸死.
其實就不是這個意思.
諸如此類.
所以就不是分開紅海.
但是.
人家後世的.
以色列人自己都去推翻了.
這些的講述.
都告訴我們.
是一個神蹟.
令到埃及人浸死.
以色列人是.
行乾地而過.
如石頭.
拋在大水中.
並且白晝用雲柱.
引導他們.
黑夜用火柱照亮他們.
當行的路.
在西乃山.
從天上與他們說話.
賜給他們正直的典章.
即十誡.
真實的律法.
美好的條例和誡命.
這樣理解律法書.
所以.
這一節聖經是不是一個誘因.
我們大家都應該用心.
去讀一下舊約的摩西五經.
讀一下.
這些律法書.
我們能不能夠.
用這樣的形容詞.
去形容律法呢.
我們今天新約的基督徒.

$^{1641}$其實我們絕對不是.
沒有學問的基督徒.
但一對讀舊約.
聖經就灑手擰頭了.
真的沒有興趣.
完全揭的意欲都沒有.
可能真的要調查一下.
有多少弟兄姊妹來崇拜的.
你的手機.
是有聖經的apps呢.
如果沒有的.
其實這都是一個恥辱.
已經不帶硬本聖經了.
叫你帶硬本聖經.
你說又重又醜又老土.
你出街.
搭車.
拿著聖經有沒有搞錯.
被人笑你是什麼人.
好像已經去到一個.
你帶本聖經好像很不好意思.
現在進化了.
手機的apps可以用了.
有沒有呢.
沒有.
有沒有時事體呢.
又是一個值得要問的問題.
如果聖經已經去到一個.
變成大家都是揭都不懂揭.
按都不懂按.
的一種文盲.
聖經的文盲.
如果苦十皆是.
其實是一個時代的悲哀.
和需要復興.
這裡告訴我們.
神的道.
在猶太人.
特別現在是復興的浪潮.
的這個月.

$^{1681}$裡面.
他們現在被神的話語.
感動.
他們能夠這樣理解.
正直真實美好.
是神的律法.
是神頒佈給我們.
由使他們.
知道你的安息聖日.
開始講到立約的.
一個指標.
就是猶太人.
你們怎麼知道神.
約納你們.
就是他們要當守安息日.
就是一個立約的.
規條.
所以在這個時候他們就告訴我們.
這裡就說.
並你們的安息聖日.
並直理僕人摩西傳給他們.
誡命條例律法.
從天上賜下糧食.
充他們的基.
就是嗎哪的事件.
四十年.
由他們背叛開始算起.
他們開始反應.
就是.
天天都吃這些.
天天吃這些很苦悶的東西.
就很苦澀.
有沒有新的東西.
類似他們開始.
有嗎哪的供應.
在這個時候.
第十五節.
賜下的就是.
糧食充基.
從盤石洗到水流出.

$^{1721}$解他們的渴.
吩咐他們進去.
得你喜誓應許.
賜給他們的地.
流奶與麥之地.
有激盤出水.
讓他們可以有飽足.
記住那個叫瑪拉的地方.
那些水是苦的.
神就叫摩西.
你在樹上.
放一支樹枝下去.
那些水就會變甜的.
由瑪拉.
就變成.
甜.
就是.
拿鵝米.
這個原文.
這裡告訴我們.
第十六節我們再看.
但我們開始祈禱認罪.
但我們的列祖.
行事狂傲.
硬著頸項不聽從你的誡命.
不肯順從.
也不紀念你在他們.
中間所行的歧視.
竟硬著頸項.
居心背逆.
自立首領.
要回他們遺牢之地.
這個就是.
十三章.
十六章.
就是可拉的事件.
其實就應該.
十四章開始.
都已經有一個背叛.
出埃及記的十四章.

$^{1761}$已經有一個背叛.
要回到埃及.
就是他們現在說的這個事件.
未入迦南之前.
已經想背叛.
已經覺得在曠野很辛苦.
還有迦南裡面的人.
是他們不可以.
打得贏的.
就開始有一個.
表面是厭戰.
實際上是一個不信.
跟著告訴我們.
提到這個.
硬著頸項.
都很全神.
不聽.
就一種.
好像.
很堅持自己.
那個意思.
不能夠去.
透過神的恩典和工作.
來看到那個出路.
但是你是樂意.
饒恕人,有恩典有憐憫.
不輕易發怒.
有風聲慈愛的神.
這個豈不是摩西.
在出埃及記的三十二章.
跟神說的.
那個頌讚.
這些又是.
很重要的金句.
在律法書裡面經常出現.
往後的先知書.
又是經常出現.
每一次他們回去認罪.
禱告都會引述.
這個神的屬性.

$^{1801}$所以.
你奉旨一定要赦免你.
你是誰.
一定要赦免你.
所以我們要學習認罪祈禱.
都要學希伯來人.
他們真的再一次.
回去不是我的罪.
你必須要赦免我.
好像那些二世祖.
我爸爸是李剛那些.
我是做錯了,怎樣呢.
我爸爸是李剛.
不是這樣,我是做錯.
但是神,你是有恩典.
有慈愛,你是不輕易發怒.
你有豐盛的憐憫.
所以我現在願意.
回轉回到你面前.
求你赦免我.
主客的身份.
對調.
不是你為主.
神你可以不饒恕我.
但是你的屬性,我明白.
我現在是.
其實有這樣的心態,你的心靈已經是謙卑.
反過來.
完全是.
不介意.
我是犯錯,那又怎樣.
下次不敢.
不是這樣的心態.
是神,你是有恩典.
有憐憫.
不輕易發怒的神.
我知道.
我正正就是知道你是這麼公義的神.
現在我就是很有意識地.
我知道自己得罪了你.

$^{1841}$我們這個民族得罪了你.
所以其實大家要抄下這節金句.
真是他們以色列人.
一有重大的.
尋求.
就要這樣來祈禱.
不是特意作文.
而是他們有感而發的.
內心的謙卑態度.
所以.
其實這個就是所羅門教我們的.
歷代志下的.
七章十四節.
不如很快和大家揭一揭.
怕大家不算是熟.
那就會好一點.
歷代志下的七章.
讀少少.
其實就是.
所羅門發夢.
當晚.
第七章.
歷代志下.
所羅門祈禱二筆.
就有火從天上降下來.
燒盡凡祭.
和別的祭.
即是這一次.
是歷史性的里程碑.
因為第一個君王.
起好聖殿.
第一個君王.
帶領全民族獻祭.
所以那次的祭是相當嚴肅.
有天火焚祭.
那就證明.
神.
閱立這個民族.
閱立這一次的獻祭.
很重要的.

$^{1881}$以後當然不會每次.
年年都有天火焚祭.
除非再有另外的.
特殊性的事件.
所以.
因為耶和華的榮光充滿了聖殿.
因為耶和華的榮光充滿了聖殿之後.
祭司就不能夠進殿.
那火降下.
耶和華的榮光.
又有光又有火.
在殿上的時候.
以色列人眾人看見.
就在鋪石地.
俯伏叩拜.
稱謝耶和華說.
這些主題曲.
耶和華本為善.
他的慈愛永遠長存.
如何唱呢.
一定是相當優美.
真是發自內心對造物主.
他的.
善.
他的屬性的善.
不是普通的.
很有人情.
那麼簡單.
這種善是包含了公義.
有真理的.
這樣的屬性.
那種kindness.
你看見一個很嚴肅的人.
他憐憫和慈愛.
只有這麼嚴肅.
你就怕他.
但造物主的嚴肅.
公義是一種善.
而且他的慈愛永遠長存.
所以我們這個民族.

$^{1921}$希伯來人不至於滅亡的原因.
就是因為他是善.
他不是一個惡神.
迦南就有很多惡神.
摩訶的神.
就很惡的.
你要奉獻.
那些屍骸的像.
中間是一個焚化爐.
就是把那些嬰兒放下去燒.
也有出土一些.
叫亂葬崗.
是那些嬰兒骸骨.
集體.
可想而知.
你拜迦南的神.
你是買他怕的.
你要把那些嬰兒放下去獻祭.
我想是很殘忍.
但是造物主耶和華是善.
所以我們要教導.
我們小朋友千萬不要搞錯.
你不聽話神就不喜歡你.
你不要嚇他.
我們要聽話.
神愛我們.
但是神是真理.
我們不要說謊.
我們要學習神的那種屬性.
就不要當警察抓你.
將那樣東西.
妖魔化了就不好.
所以猶太人他們.
是很明白這件事.
不是因為他們是.
什麼優質的民族.
是因為神.
他的本性是善.
他的慈愛是永遠.
常存.

$^{1961}$輾轉他們獻祭之後.
第十一節.
所羅門造成了.
耶和華殿和皇宮.
再耶和華殿和皇宮.
凡他心中所作.
都順順利利的作成.
第十二節.
這裡開始重要了.
沒有人了.
只有他自己在皇宮.
耶奸因為他剛剛就職典禮.
成為以色列的國君.
耶奸耶和華.
向所羅門顯然對他說.
我已經聽了你的禱告.
也選擇了這個地方.
作為祭祀我的殿羽.
我如果使得天.
閉塞不下雨.
或使得蝗蟲吃.
這地的出產.
或使得瘟疫流行在我民中.
這稱為我名下的子民.
即是說你們要有意識地.
有少少對號入座地.
就是如果你看到這些事情.
你就要有意識地.
那個心靈的意識.
就是要反思一下.
神是不是有些事情.
想我們去處理呢.
這裡告訴我們.
這稱為我名下的子民.
這就是七章十四節.
關於禱告的呼喚.
如果他們自卑禱告.
就不求我的面.
所以是兩條腿的.
首先自己.

$^{2001}$有一個很強烈的意識.
我要謙卑自己.
準備好我的心靈.
所以古代的人就.
沐浴更衣.
代表著對現在這件事.
很嚴肅地看待.
所以我們真的有時.
如果你能夠得的.
星期三來祈禱會.
早點來安靜.
這個聖殿裡面.
自己準備好心靈.
去尋求神.
這裡告訴我們.
尋求我的面轉來他們的惡行.
我必從天上垂聽.
赦免他們的罪.
醫治他們的地.
我必睜眼看則以聽.
特別強調那個形容詞.
其實神是沒有眼睛的.
祂是不需要有聽覺的.
但是祂是全地的主.
特別留心.
這樣的人.
作這樣的禱告.
即是說一切都在乎.
我們的心靈有沒有準備好.
有沒有準備好.
有意識地來敬拜神.
我真的很盼望.
錦宣家的弟兄姊妹.
你都已經很自由來崇拜了.
要有意識地來崇拜.
這個意識是.
你真的很想去朝見.
去拜神.
來到這個地方不是疑神疑鬼.
是你的心靈準備好.

$^{2041}$我見這位主.
我有一個準備.
這樣的意識.
來朝見神.
是很重要的.
沒有意識.
即是心懷雜念.
可有可無.
給面子派對.
這樣是不行的.
這樣是不會令到我們的靈命.
有一個明顯的改變.
重點仍然是.
真是那個.
心靈.
第十七節.
你如果在我面前.
效法你父大位所行.
等等.
我不讀這麼長了.
第十九節.
倘若你們轉去雕器我.
使到你們的律例戒命.
等等去事奉別神.
我就會使到.
以色列人從我賜給我的地上.
拔出根來.
並且我為我明.
等等就一定會.
去拆毀.
路過的人.
路過.
都會去.
好像都在暗笑.
為什麼.
耶和華.
為何使到這個地和這個殿.
變成這樣.
那些人就說.
離棄了神.

$^{2081}$這裡就很特別.
沒有記載所羅門醒了之後的祈禱.
列王記有記載.
因為我們不夠時間.
就不看這麼長了.
我們回來看這個離棄米記.
那個情況.
所以我們要有意識地對於神的慈愛.
是要.
有這樣的賞心.
所以我們剛才.
說遠了.
其實就是這個禱告的第二個部分.
開始要來對.
這個豐盛慈愛.
有一個很強烈的意識.
再讀一次第十三節.
離棄米記的第九章.
應該是第十七節.
不肯順從.
也不紀念你在他們中間.
所行的騎士.
竟硬著心.
諸如此類.
居心背逆自立首領.
要回他們遺留之地.
但李耶和華.
是以饒恕人.
有恩典有憐憫.
不輕而發怒.
有豐盛慈愛的神.
並不凋棄他們.
所以他們雖然注了一隻金牛毒.
彼此說.
這是領你出埃及的神.
你想想.
好像連一個神密教材都在.
叫阿倫.
因為摩西上山四十天.
四十天都還沒有回來.

$^{2121}$一個人上了山.
一個人也不回來.
又沒有通知.
怎麼辦呢.
打電話下來也好.
怎麼辦呢.
百姓就鼓噪.
迫阿倫.
將他們有預備好的金.
去求神問卜.
給我們一些指示.
終於出現了金牛毒.
他們就拜金牛毒.
不過他們就.
很婉轉的.
這就是.
帶我們出埃及的神.
就是說他們已經將.
種物主沒有形狀的.
不可以有形狀的.
來變成.
埃及的神.
降格到.
變成牛毒.
神就是金牛毒.
很大的罪.
十誡才剛聽完.
四十天前才剛頒佈了.
十誡除我以外.
你不可以有別的神.
不可以將.
朱天的.
你見到的地上的滄海裡面.
你見到的.
來變成神.
這樣犯了兩條十誡.
這樣來到去說.
現在我們要趕著怎麼辦呢.
就求問了.
回到埃及好不好.

$^{2161}$好像來到去.
他們製造.
一個神出來.
亞倫被迫於無奈.
要接受這樣的安排.
當時很特別.
沒有WhatsApp.
沒有飛甲傳輸.
摩西突然間聽到.
因為神就跟摩西說.
你回去吧.
山下的子民.
已經犯罪.
離棄我了.
他們不是離棄你,是離棄我.
做他們的神.
所以摩西原來這麼先進.
都完全不需要.
有任何的科技.
就已經知道.
山下發生什麼事.
你可想而知.
多特別的救藥.
所以我們今天.
我們能夠隨時WhatsApp電話.
打去加拿大.
打去盧森堡,打去不同地方.
但是已經是可以.
摩西都不需要.
有任何亞倫的一個.
信息,他立即就知道.
山下已經是.
來到他背叛了.
硬著頸汗的居心背役.
這就告訴我們.
大大的.
野動神的怒氣.
所以《出埃及記》的32章.
摩西下山之後.
就吩咐了,所以那段聖經.

$^{2201}$大家如果能夠留意到.
利未人從此.
就可以翻身了.
因為利未人站在摩西那邊.
摩西吩咐.
凡屬耶和華的.
就站在我這邊.
所以利美之派就站在.
摩西那邊,接著摩西吩咐.
未收太陽之前.
你們.
為神去殺盡.
這些以色列人.
利美之派就真的去.
行殺戮了,所以.
神就透過這樣.
就說利美之派.
本來利美是受咒助的.
那個咒助是什麼呢?.
《創世記》49章.
利美的咒助.
因為她和她的哥哥.
在樓邊去殺了.
事件全家.
本來人家.
行完國禮.
還是很痛,傷口還沒有埋口.
你就去殺了人家的男丁.
就是要.
竊行,那麼陰險.
那就.
雅各臨終.
十二個兒子.
都有不同的報應.
有些是祝福,有些是咒助.
利美是咒助.
那個咒助就是.
我的民不要和他們一起.
利美要分散在以色列當中.
就是說他們.

$^{2241}$沒有自己固定的地圖.
很慘.
變成一個沒有自己的.
家鄉固定的地方.
一個咒助.
後來神就.
使咒助變成祝福.
所以這個就是.
利美之派.
站在神這邊.
好像一個翻身的機會.
當然還有後來.
就是菲尼哈.
一個很忠誠的.
殺了.
心利之派.
和那些魔壓的女子.
那就更加挽回.
神的怒氣.
所以這兩件事令到利美之派.
沒錯,那個分散就.
不能改變了,但是.
給一個很崇高的身份.
他們成為了祭司.
利未人,他們可以在.
不同的支派.
有他們自己的屠城.
他們沒有自己的永久.
地方,但是可以有住的地方.
而且那個住的地方.
還要是有一個.
很尊貴的一個.
的職份,就是.
教導律法書.
和在聖殿裡面.
去公職,事奉.
但是分散在.
十二個支派的地方.
來作為那個.
好像主學老師,或者是傳道人.

$^{2281}$這樣的身份.
這就是那個.
怒氣的典故.
第十九節.
李環是大師憐憫.
在曠野不丟棄他們.
白晝雲柱不離開他們.
仍引導他們.
行路,黑夜.
火柱也不離開他們.
仍照亮他們當行的路.
你也賜下你良善的靈.
教訓他們.
是在說出埃及的時候.
未嘗不似嗎哪.
洗他們湖口.
並賜水.
解他們的渴.
主耶穌就說.
我們日用的飲食.
不要掛心.
你看那些飛鳥.
又不種又不收.
又不會有庫房.
又不會有戶口.
又不會有電子支付戶口.
拿消費券.
那些飛鳥.
耶穌這樣說.
你們比飛鳥起碼.
還要貴一種.
你看野地的百合花.
又不會紡織.
又不會什麼.
但是在所羅門的時候.
最榮華最風光的時候.
他們的裝飾.
都不夠野地的百合花那麼漂亮.
那麼燦爛.
開得那麼鮮艷.

$^{2321}$雖然開一天就可能已經枯萎.
但是神都會尚且.
來到去有這樣的裝飾.
何況你們呢.
所以你們不要掛心.
吃什麼喝什麼.
這是外邦人所求的.
你們要先求神的國.
和神的義.
這些東西.
很肯定.
就是上下文說的.
基督徒的日用飲食.
必需品.
所以大家要問一問.
我們剛剛.
上個禮拜教完主一學.
我就問同學.
你認為的必需品.
是不是別人的奢侈品呢.
這是一個.
大家要自己請問的問題.
我的必需品就是.
我每天要怎樣怎樣.
這個級數我才OK.
才是我必需.
但是原來是別人.
一年收入.
認為的奢侈品.
這個就.
大家要思想.
所以盡量.
基督徒要簡樸.
盡量我們能夠有的.
我們分給一些沒有的人.
真的能夠有神.
那個善良的心靈.
所以這裡告訴我們.
日用飲食.
未嘗不似嗎哪.

$^{2361}$洗他們湖口.
二百萬的以色列人.
在這個曠野.
並賜下水.
解他們的渴.
這些是主耶穌說的.
日用的必需的飲食.
在曠野.
四十年.
你養育他們.
他們就一無所缺.
衣服沒有穿破.
也沒有腫.
你可以想而知.
真是神蹟.
你想想在曠野.
那種炎熱.
四十多度.
他們可以這樣生活.
近乎沒有瘟疫.
沒有死人.
一定是在背疫的時候死.
不計他們在大自然.
因為暴露.
在這麼炎熱的曠野.
這裡告訴我們.
第22節.
並且你將烈國之地.
照分.
照份.
因為每一個支派有自己的份.
是由摩西.
抽籤來給他們.
臨終之前.
他們就得了.
西弘之地.
希實本之地.
和巴山王岳之地.
這是劉弼和加德兩個半支派所得到的話.
是在約旦河東的地方.

$^{2401}$你也使他們的子孫.
多於天上的星.
帶他們到你.
所應許.
他們列祖進入得為業之地.
這樣他們進去.
得了那地.
你在他們面前.
制服那地的居民.
就是迦南人.
將迦南人和其君王.
並那地的居民.
都交在他們手裡.
讓他們任意而待.
他們得了堅固的.
誠一肥美的地圖.
充滿各樣美物的房屋.
作成的水井.
葡萄園.
橄欖園.
這些是經過悉心的研發.
耕種.
開墾.
你才有深水井.
你才有灌溉.
很多的系統慢慢越來越先進.
那個豐收.
那些提煉的果實.
越來越更肥美.
鮮甜.
諸如此類.
都是神給的智慧.
神給的地圖.
神給他們能夠有生存的.
民族的尊嚴.
無論是.
頭一段時間.
約書亞帶他們南征北伐.
這樣打贏了.
最終他們得到這個地圖.

$^{2441}$最終都是神的恩典.
讓這個民族.
能夠吃得飽.
身體肥胖.
因你的大恩.
心中快樂.
再一次.
回憶那個美好的記憶.
然後就是延移.
就是26節.
我們下一個禮拜.
就和第十章一併去看.
因為不夠時間看完.
不知道為什麼突然斷線.
沒有了那個網絡.
不要緊.
我們差不多完結了.
他們讀不完一段.
原來我還沒轉回那個網絡.
在辦公室那邊的宿舍.
所以不怪得.
所以我有些大意.
我一直都看著.
不知道他們是不是很早就沒有網了.
都不理了.
到時我們去聽錄音帶.
多謝主.
讓我們能夠在今晚.
有第八,第九章,半章.
一起研讀.
我們將經文這樣讀.
簡單的講解.
都讓我們看到.
我們一個群體.
很需要有集體的意識.
很需要對我們的存在.
為什麼我們是教會.
我們不是一個普通的團體.
我們不是一個政治的團體.
我們不是一個興趣班的團體.

$^{2481}$我們是一班.
對神.
對永生.
對天堂.
認真相信的群體.
以致我們在地上.
我們真的要為福音而活.
我們很盼望.
我們這個群體.
因為第四十二年.
我們的長輩.
前輩建立了.
很重要的根基.
現在去到第三,第四代的時候.
第二代都仍然.
在這裡的時候.
我們真的很盼望.
有天主來感動我們.
讓我們能夠在整個疫情.
三年半之後.
我們大家都有一個.
很重要的心靈的一個躍動.
而且這個躍動不是一個口號.
是心靈裡面.
大家很強烈的意識.
我們要重拾.
我們是神兒女.
這個身份.
這個身份的歸屬.
很重要的一個sense.
在我們內心裡面.
求主幫助.
我們仰望禱告.
奉主耶穌基督的名.
阿們.
\newpage



\section{馬太福音 28:11-20}
\label{sec:WLN153QoHOE}
\textbf{2023-06-04 加恩樓《後真相年代下的信仰反思》馬太福音 28\_11-20 講員:吳景霖傳道}
\newline
\newline
連結: \href{https://youtube.com/watch?v=WLN153QoHOE}{\texttt{https://youtube.com/watch?v=WLN153QoHOE}} ~~~~ 語音日期: 2023-06-06
\newline
\newline
\hyperref[sec:C1vy6YO3h0g]{\small{< < < PREV SERMON < < <}}
~
\hyperref[sec:index_chronic]{\small{[返順時目]}}
~
\hyperref[sec:index_scriptual]{\small{[返順卷目]}}
~
\hyperref[sec:jJiE2EQ4f48]{\small{> > > NEXT SERMON > > >}}
\newline
\newline
馬太福音 28:11-20
\newline
\begin{longtable}{cl}
\hline
\hline
章節 & 經文 (和合本修訂版)\\
\hline
28:11 & \begin{tabularx}{0.7\textwidth}{X} 她們去的時候，看守的兵有幾個進城去，把所發生的事都報告祭司長。 \end{tabularx} \\ \\ \relax
28:12 & \begin{tabularx}{0.7\textwidth}{X} 祭司長和長老聚集商議，就拿許多銀錢給士兵， \end{tabularx} \\ \\ \relax
28:13 & \begin{tabularx}{0.7\textwidth}{X} 說：「你們要這樣說：『夜間我們睡覺的時候，他的門徒來把他偷去了。』 \end{tabularx} \\ \\ \relax
28:14 & \begin{tabularx}{0.7\textwidth}{X} 若是這話被總督聽見，有我們勸他，保你們無事。」 \end{tabularx} \\ \\ \relax
28:15 & \begin{tabularx}{0.7\textwidth}{X} 士兵收了銀錢，就照所囑咐他們的去做。這話就在猶太人中間流傳，直到今日。 \end{tabularx} \\ \\ \relax
28:16 & \begin{tabularx}{0.7\textwidth}{X} 十一個門徒往加利利去，到了耶穌指定他們去的山上。 \end{tabularx} \\ \\ \relax
28:17 & \begin{tabularx}{0.7\textwidth}{X} 他們見了耶穌就拜他，然而還有人疑惑。 \end{tabularx} \\ \\ \relax
28:18 & \begin{tabularx}{0.7\textwidth}{X} 耶穌進前來，對他們說：「天上地下所有的權柄都賜給我了。 \end{tabularx} \\ \\ \relax
28:19 & \begin{tabularx}{0.7\textwidth}{X} 所以，你們要去，使萬民作我的門徒，奉父、子、聖靈的名給他們施洗， \end{tabularx} \\ \\ \relax
28:20 & \begin{tabularx}{0.7\textwidth}{X} 凡我所吩咐你們的，都教導他們遵守。看哪，我天天與你們同在，直到世代的終結。」 \end{tabularx} \\ \\
[1ex]
\hline
\hline
\end{longtable}
$^{1}$大英姐妹午安.
我還有兩次在大家當中講道.
所以辛苦大家忍耐一下.
今天跟大家講道是一個非常困難的任務.
希望上帝閱立我講的道.
何謂後真相呢.
我們首先要講什麼叫做真相.
真相可以說是事實的實際情況.
我是一個很喜歡看偵探的人.
由神探伽利略看到金田一和明偵探柯南.
我都看的.
明偵探柯南有一句說真相只有一個.
但真相是不是只有一個呢.
我舉個例子.
見Michael笑得這麼大聲拿他來做例子.
舉個例子有一天你見到Michael打了Penny一巴.
打老婆.
打老婆是他不對.
事件就是這樣.
但這個是不是所謂的真相呢.
原來我推前一點時間線去看.
原來不是.
原來之前Penny已經又插眼又抓又拿刀.
跟著Michael打了他一巴讓他冷靜一點.
跟著Penny不對.
舉例子大家不要以為是真的.
原來不是這麼複雜.
原來Michael爛賭拿了全副身家去賭.
跟著Penny就瘋了.
於是又出刀.
跟著Michael就打他.
跟著Michael不對.
你明不明白將時間線一路推的時候.
發覺為什麼事情看來看去都不是找到真相.
如果由其他人加入去看這件事可能會更加複雜.
譬如我們認識Michael.
於是我說不是的他平日不是這樣的人.
會不會有些事情誤會了.
但如果一個女士去看的時候.
當然是男人不對.

$^{41}$我也聽過別人這樣說.
總之就是男人不對.
都不理那件事是男人不對.
不同人看不同角度去看.
似乎真相不是只有一個.
我甚至可以說100個人看同一件事可以有100個真相.
事實就只有一個.
某年某月某日Michael打了Penny一巴.
就是這樣.
跟著怎樣去解釋或者怎樣去看呢.
可能隨著不同的人不同的角度會出現不同的事情.
甚至我只需要將一些事情.
那個事實去到挪用或者是不同的角度去到詮釋.
就會產生不同的結果.
這個就是所謂的後真相.
真相本來就是浮動.
後真相就是透過一些事實的操弄.
就可以讓事情完全不同.
舉個例子我什麼都不說.
只是說Michael打了Penny一巴.
其他的全部不說.
於是他就變成了壞人.
是不是.
但我反過來我又可以說.
我只是強調Penny出刀.
在這樣的情況下.
他打了一巴.
那些輕描淡寫過了就算了.
那你就變成了Penny是壞人.
事情發生的時候.
你發覺原來我們透過一些事實的操弄.
我只是說事實.
但是就可以讓真相完全不同.
後真相年代就是沒有客觀的真相.
只有被篩選下來的部分.
對於操弄者他認為他想要的東西.
他就用那些虛假的東西.
或者一些真相去到挪用.
去到支持他自己.
你上網經常都會看到一張照片.

$^{81}$然後加一些說話.
然後發出來.
大家WhatsApp應該收不少.
是不是.
這樣就好像去到支持那句說話.
好像特別有說服力.
最好就是黑白照片.
一個你不認識的外國人.
然後寫一些說話.
就會好像那句說話有些說服力.
你要知道其實這些東西.
挪用了或者修改了.
不一定是歷史的真相.
譬如有句說話.
宗教是人民的鴉片.
有沒有聽過.
其實有句說話.
宗教是人的心靈的鴉片.
但是其實整篇東西你回去Google一下.
前文後理不只是這句.
你以為是在毒害人心靈.
其實是在說人對鴉片那種痛苦的麻醉.
去到安慰人心.
你回去Google一下.
你要知道透過這樣的操作.
是很容易去處理的.
就算我現在在錄影.
你中間截一句出來我就死定了.
我說是我砌Michael生豬肉.
說他打老婆.
你只要cut了我前面.
我舉個例子這幾個字.
我就死定了.
是不是.
整件事你會發覺原來所謂的真相.
其實在後真相年代.
是會不斷地被操弄.
在很多年前發生了一件慘劇.
死了萬多人.
有四萬多人受傷.

$^{121}$萬多人死亡.
失去他們寶貴的生命.
這場絕對是一場人禍.
當時我只是一個小學生.
十二歲.
我在電視新聞裡看到發生什麼事.
我大概知道什麼事.
學校老師都有說.
當時電視裡面有很多名人出來.
又跳舞又唱歌.
又譴責又登報紙.
過了很多年.
這件事.
有些人可能記憶已經模糊了.
隔了這麼久.
有一些是未出世的.
我都很肯定他們當時是未出世.
但是有很多相片和影片流傳了.
在今天.
但是亦有人開始去指責.
這些相片和影片是假的.
是經過修改的.
是外國人去誣衊我們.
去扭曲一些事實.
我發覺.
好像這些現實.
為什麼開始和我的記憶出現距離.
有些人說的話.
和我的記憶好像對不到.
到底是什麼一回事.
是不是什麼曼德拉效應.
到底是什麼.
我試過有一個比我小的朋友.
他只是比我小一點.
我都相信他當時還是一個小學生.
但他跟我說.
當年那些不是這樣的.
我覺得很奇怪.
記憶是可以被人操縱去修改的.
所以我們回到聖經.

$^{161}$我們看回.
這些事其實日光之下無新事.
在聖經裡面依然出現這樣的事.
第十一節開始說.
看守的兵丁.
他們去城裡面將所發生的事告訴祭司長.
第十一節說什麼.
原來他們見到耶穌的墓穴是打開了.
有天使出現.
說耶穌已經復活了.
馬上回去報告祭司長.
嚇得鼻子都沒肉了.
回去.
祭司長怎麼辦.
是不是馬上悔改.
我擔心死了神的兒子.
我要悔改了.
不是的.
他拿錢去收買兵丁.
給錢他們.
你要說是晚上睡覺的時候.
門徒偷走了耶穌的屍體.
是不是.
直接用錢去收買.
如果總督聽到.
即是你老闆聽到.
我勸他保你們沒事.
這個完全是一個非常明顯的修改事實.
用謊言去修改事實.
而且成功了.
他說士兵收了錢.
就照他吩咐去做.
這話在猶太人中間流傳直到今天.
馬太科罕成書的時候.
應該都是主後六十年.
即是說耶穌都死了三十多年.
這件事一直都這樣流傳.
都還在說門徒偷走了耶穌的屍體.
原來世界就是這樣.
我發覺聖經永遠都適合看.

$^{201}$為什麼呢.
因為幾千年前已經是這樣操弄事實.
他們嘗試去修改事實.
讓人的記憶和現實出現出入.
接下來發生什麼事呢.
門徒就要面對一個處境.
當他們的教主.
即是他們的宗教領袖.
被釘死之後.
他們要面對的可能是被清算.
他們要被清算.
於是那些門徒就躲起來.
你看到聖經繼續說.
他們是躲起來在一個地方聚集.
然後耶穌復活的時候.
是不用開門都能進入屋子.
他們很驚訝.
然後到士徒行傳.
聖靈未降臨之前.
他們都是躲起來一起祈禱.
他們是擔心自己被清算.
而事實上是會的.
因為你看到耶穌在整個馬太科罕裡面.
都不斷提出.
我是會被釘死的十字架.
而我的門徒是會一起被迫.
這些看聖經你就會發覺.
作耶穌基督的門徒其實很大代價.
於是他留下了一個吩咐.
天上地下所有的權柄都次級我了.
你會看到其實.
接著那個大使命是很熟悉的.
我們讀完它吧.
第十九節那裡.
所以你們要開始萬民作我的門徒.
奉父子聖靈的名給他們施洗.
凡我所吩咐你們的都教訓他們遵守.
看我天天與你們同在直到世代的終結.
這個大使命.
我相信我們各位的基督徒很熟悉.

$^{241}$但通常你從第幾節開始背?.
第十九節.
所以你們要去.
一背就.
前面那句天上地下所有的權柄.
沒有的.
再之前我剛才和你們看那裡.
關於士兵被人收買那裡有沒有看過?.
很多人根本就不看的.
我想說的是.
為什麼他一開始.
耶穌要說天上地下所有的權柄都給了我.
因為耶穌基督他要肯定一件事.
就是我現在是有權的.
我透過釘死在十字架上.
我已經建立了天國.
耶穌基督已經登基成為天國的王.
他是有實權的.
你要知道其實.
耶穌他由出世.
就是馬太福音我們解經的原則.
我們看回同一本書.
馬太福音從第二章開始已經在說希律王.
因為怕這個猶太人的王出世.
影響到他的地位.
所以怎樣?.
追殺他.
我們唱平安夜的時候很平安.
你覺得勝誕節的時候平安夜好像很寧靜.
但實際上耶穌出世之後.
凡兩歲以下的猶太小朋友被殺.
這是一場血淋淋的政治大屠殺.
明白嗎?.
你要知道耶穌出世其實已經被人追殺.
然後追殺到最後.
他釘十字架.
耶穌為什麼要釘十字架?.
有些基本要知道耶穌為什麼要釘十字架.
為我們的罪.
贖我們的罪.

$^{281}$但是釘十字架的人.
是不是耶穌過來.
你要預備贖我們的罪.
請自己上去.
釘他的人是不是想著要耶穌贖他的罪?.
不是的.
耶穌基督被人釘的時候.
那些人他們有自己的意志.
不是說上帝在後面用無形的手控制他們.
然後就去釘耶穌十字架.
你看回聖經.
他們全部是因為自己的利益.
他們的地位.
他們作為宗教領袖.
他們被耶穌基督影響了.
影響了他們的地位.
影響了他們的利益.
所以他們要釘死這個眼中釘.
要殺了他.
所以門徒接下來.
耶穌死了之後他們要面對的就是.
要擔心被人做最後的清算.
所以耶穌就要告訴他們.
你不用怕.
天上地下所有的權柄都賜給了我.
如果你作為一個基督徒.
你成為了天國的子民的時候.
你成為了天國的一份子的時候.
你是不需要擔心地上的政權.
對你的迫害.
或者是對你的一些操弄.
因為上帝耶穌基督是有權柄的.
我們作為一個基督徒.
我們有一個高於一切的身份.
一個天國子民的身份.
你作為任何一個人.
你作為你的.
譬如香港人也好.
你作為你兒子的父親.
兒子的母親.

$^{321}$或者你父親的兒子的女兒.
或者作為公司的職員.
作為一個老師.
作為一個廚師.
你在其他的身份之前.
你第一個身份一定是基督徒.
你最首要最重要的身份就是基督徒.
有一個國家傳說中.
說他的護照很厲害.
在全世界任何一個地方都可以將你帶回去.
我不知道他的護照是否那麼厲害.
但是我卻可以說.
當你作為天國子民的時候.
你在全世界任何一個地方.
上帝都可以照著你.
這個是跟大家說之餘.
其實這句話也是跟我自己去說.
無論我們在世界任何一個角落.
世界上任何一個地方.
你首先是天國的子民.
並且耶穌基督一定會在那裡照著你.
所以耶穌說天上地下.
所有權柄都次級可了.
在傳福音去其他人耳中傳出去之前.
他要確立這個權柄.
你是有資格.
你是有身份.
有權去到傳福音的.
就算是大祭司希臘彼拉多.
或者是凱撒大帝.
都不可以阻止人去傳福音.
說到這裡的時候.
我覺得有一個更重要的問題.
就是什麼是福音.
福音是什麼?.
你們接受了福音.
坐在這裡信了耶穌.
福音是什麼?.
如果聽我講得多.
我從來都非常反感的一件事.

$^{361}$就是認為福音是舉手決志信耶穌.
這樣就算了.
當時傳福音的實際操作就是叫你舉手決志信耶穌.
然後搞定之後就下一個了.
是不是這樣的呢?.
如果是這樣的話.
你決志信了耶穌.
那福音對你來說還有沒有用?.
就是入場券.
我進了迪士尼.
那張票就放在這裡.
我還需要福音來做什麼?.
是不是這樣的?.
不是的.
你想想.
如果你是只是拿了張入場券上天堂.
如果上帝也是這樣想的話.
那耶穌還需要在這裡做這麼多事情嗎?.
三年來不斷地去報道.
行神蹟.
去講話.
然後等門徒寫下來.
我們有這麼厚的福音書.
為什麼要留下這麼多東西給我們?.
有些人說.
耶穌做這麼多事情是要告訴我們他是神的兒子.
那是不需要的.
他出生然後一下子行了神蹟.
然後從嬰兒變成大人.
然後走十幾二十個神蹟.
然後就完成了.
上十字架.
這樣就可以了.
你不要將信仰或者耶穌基督弱化.
簡化成為.
信耶穌上天堂得永生.
搞定.
我們在信仰裡面其實有很多東西.
我經常都在想福音是什麼.
我覺得我信了二十多年.

$^{401}$我到今天依然要去想.
福音是什麼.
對我來說有什麼意義.
在這一刻對我來說.
最大的意義就是與神復和.
與上帝復和.
我們人類得罪了上帝.
離開了伊甸園.
我們那個破損的關係.
你透過耶穌基督.
你可以重新修補.
與上帝復和是我們這個信仰.
一樣很重要的東西.
而本身這件事就是充滿意義的.
與上帝復和不是一個進入天堂的手段.
我記住了.
與上帝復和不是進入天堂的手段.
舉個例子.
譬如你結婚.
你結婚不是得到一個人的手段.
你跟他簽婚書.
你還不死簽了.
你是我的人.
不是一個手段.
而是你願意跟他建立一段關係.
因此你就去跟他結婚.
然後過程也好.
結婚之前的過程也好.
結婚之後跟他相處也好.
這一切都是有意義的.
都是存在意義的.
而我所說的.
你與上帝復和不應該是一個進入天堂的手段.
當然進入天堂也很重要.
但是你本身能夠成為基督徒.
今天你可以在這裡.
過著與神復和的人生.
上帝已經原諒了你.
你成為一個天國子民.
這一切本身就有意義的.

$^{441}$我覺得最基本.
我們都應該這樣去理解福音.
我們在世.
是要去進行耶穌.
他留下的教導.
我們學習一套天國的生活方式.
你作為一個天國子民.
你如何活在這個世界.
你如何在這個世界裡.
讓人見到耶穌基督.
你如何在這個邪惡的世代裡繼續生活.
我們要在這個.
加入一個群體.
我們加入一個天國的時候.
我們要去學習如何進入一個群體.
其實你去到任何一個新的地方.
你都要學一些新的規矩.
你都要學一些新的東西.
你知道那裡是如何運作的.
你知道那裡是如何生活的.
那裡是怎樣的.
舉個例子.
很多人在灌輸我將來去到加拿大.
我如何去生活.
譬如說加拿大的人很慢.
一天只做一件事.
你千萬不要想著一天做多過一件事.
我說不是的.
我平時一天出去九龍.
我做四五件事.
然後才坐車回來.
他說加拿大不行的.
一天你預做一件事.
不然你會爆血管的.
很多這些心理準備給我.
例如叫我學懂如何鏟雪.
諸如此類.
或者你的衣服不用帶過去.
那邊不適合穿.
很多這些東西.

$^{481}$你加入天國都一樣.
在天國裡面如何去生活.
你這個群體是如何運作的呢.
在上帝的眼中是如何的呢.
耶穌基督在那三年裡面.
透過祂自己的言訓.
祂在教給我們.
你如何去活出一個天國群體的模式.
雖然我們是天國群體.
但是我們都依然活在這個世界裡面.
大家明天都要上班.
我就放假的.
其實全世界人星期一放假都不太享受.
因為星期天很忙.
但是我們依然要活在這個世界.
你離開了教會.
你要面對這個世界.
在這個邪惡的世代.
在這個複雜混亂的世代裡面.
你要按天國的價值觀去生活.
去過出一個基督徒的人生.
其實不簡單.
你是會被逼白的.
而且你要面對的是.
這個世界會不斷嘗試將你拉回過去.
回到今天的題目.
就是後真相年代.
在後真相年代裡面.
其實有很多東西已經沒有真相了.
人是會透過不同的方法.
無論真假的資訊.
即是真實的資訊都可以去洗腦.
去做選擇性的修改一些記憶.
去改寫你的價值觀.
你是不能避免的.
你作為現代人.
除非你躲在一個山洞裡.
否則你會不斷被一些東西去騷擾.
譬如我已經算是很小心去看一些資訊.
但是有時候看Facebook.

$^{521}$不經意都是一張照片.
一段說話這樣看進去.
然後我發現.
入了我的腦洗了腦.
而我還沒有fact check.
不知道那句說話是不是真的.
這些東西是不斷地在修改我們的記憶.
但是我要提醒大家.
你不可以隨波逐流.
有些東西對就對.
錯就錯.
是不容許去改寫.
譬如希特拉屠殺猶太人.
這個一定是不可以改寫的.
一定是一個邪惡的事情.
但是那時候其實仍然有人支持.
因為用回一個說法.
猶太人釘死耶穌.
你作為基督徒.
你當然不能夠容許這麼邪惡的人民存在.
真的有這樣的說法.
有些東西對就對.
錯就錯.
你不可以去修改.
例如俄羅斯入侵烏克蘭.
我真的覺得你不應該說是烏克蘭咎由自取.
但是有些媒體都是這樣說.
事實上.
不斷地在網上有很多的言論的修改.
嘗試去修改人的想法.
我知道有很多弟兄姊妹都移了民.
有些在線上看我們的聚會.
我想提醒就是.
去到世界上任何一個角落.
都一樣發生著這些事情.
有權有勢的人是會去修改歷史.
歷史是由勝利者去撰寫的.
很容易的.
出一些電影或者教材.
或者改教科書.

$^{561}$禁網.
建一個網絡圍牆.
其實去到世界上任何一個角落都一樣.
不要說有權有勢.
就算只不過是有錢都可以做到.
Facebook無端端都會鎖人的帳號.
Twitter.
Elon Musk不喜歡你.
你就可以封了你的帳號.
這個世界是這樣運作的.
其實去到任何一個角落.
你去到任何一個角落.
都有人嘗試用這些方法.
在後真相年代.
去改變你的思想.
因此我鼓勵大家.
首先就是.
你看什麼.
首先你要查證.
查一查.
看清楚是不是一件這樣的事.
不要看到一張相片.
看到一堆字你又相信了.
之前呢.
早幾天鑽石山不是有一單慘劇.
真的是.
慘劇之後兩三個小時.
就有媒體出新聞.
說是情殺案.
但是跟著再查下去.
根本就是一個慘劇.
是一個隨機的殺人.
那兩個女孩應該跟他沒有關係.
但是一個大媒體.
他可以隨便出一段這樣的新聞.
到今天我都沒有看到他有道歉.
你會發覺香港.
其實今天那個.
坦白說那個媒體的層次是非常低.
我們要習慣要查證.

$^{601}$不要立即就相信.
所謂不要跟車太貼.
另外就是.
看不同角度的東西.
不同角度都看.
你不要只看你自己的角度.
坦白說我們有時候是會的.
只看自己角度的東西.
甚至看不同語言的東西.
看英文的報紙.
不用很厲害.
Google翻譯很快.
我也是一個網頁.
翻譯成中文就可以看了.
你試試從不同的角度去看這件事.
你會發覺你個人是客觀很多.
另外就是我們要回到耶穌的言訓那裡.
第20節說.
凡我所吩咐你們的都教訓他們遵守.
耶穌基督在世這麼久.
他這麼多的言訓.
這麼多的以身作則的東西.
你要教導人去遵守.
你自己也要遵守.
你要走一套天國的價值觀.
你要按著耶穌基督教你的方法去生活.
當你的生活態度.
你做人.
你學似耶穌基督的話.
你很自然會拒絕一些謊言.
因為跟你裡面的價值觀已經出現矛盾.
當你被耶穌基督塑造你裡面的時候.
那些假的消息或者謊言.
或者錯誤的價值觀來到.
你是懂得去拒絕的.
所以我們是需要不斷地學習.
你不要說傳福音就是教他.
你也要教自己.
你要多看聖經.
其實聽我說到.

$^{641}$我發覺說到最後我也是那句.
我叫你回去看聖經.
多看聖經.
我說來說去三幅被.
都是請你回去多看聖經.
去學習聖經的真理.
讓你自己過一個合乎上帝的生活.
有上帝的價值觀.
你錯極也有個譜.
錯極也有個譜.
我可以很放心.
所以耶穌基督他應許了.
他說你遵守他的教導.
他就天天與你們同在直到世代的終結.
耶穌基督會看著你.
他會陪著你.
還有這裡暗示了一件事.
這個世代是怎樣.
會終結的.
世代會終結的時候.
你要回去天國跟耶穌基督交代.
這個世界末日的時候.
你要跟他交代.
然後去到審判的台前.
耶穌基督問你.
你這輩子做了什麼.
你過了什麼日子.
你做過什麼.
你怎樣跟他交代.
我希望大家透過活出天國的價值觀.
去對抗這個世界給我們的虛假資訊.
對抗那些邪惡的勢力.
去洗我們腦.
接下來我邀請大家.
有一段時間.
一分鐘.
去到安靜.
自己去禱告.
去到跟神說.
你剛才聽到的東西.

$^{681}$裡面有什麼是最觸動你的.
一分鐘之後.
我們就請主席去到唱會影師.
\newpage



\section{路加福音 5:27-32}
\label{sec:jJiE2EQ4f48}
\textbf{2023-06-04 午堂《悔改的筵席》路加福音 5\_27-32 講員:楊明麗姑娘}
\newline
\newline
連結: \href{https://youtube.com/watch?v=jJiE2EQ4f48}{\texttt{https://youtube.com/watch?v=jJiE2EQ4f48}} ~~~~ 語音日期: 2023-06-06
\newline
\newline
\hyperref[sec:WLN153QoHOE]{\small{< < < PREV SERMON < < <}}
~
\hyperref[sec:index_chronic]{\small{[返順時目]}}
~
\hyperref[sec:index_scriptual]{\small{[返順卷目]}}
~
\hyperref[sec:KzM_XrVrUvg]{\small{> > > NEXT SERMON > > >}}
\newline
\newline
路加福音 5:27-32
\newline
\begin{longtable}{cl}
\hline
\hline
章節 & 經文 (和合本修訂版)\\
\hline
5:27 & \begin{tabularx}{0.7\textwidth}{X} 這些事以後，耶穌出去，看見一個稅吏，名叫利未，在稅關坐著，就對他說：「來跟從我！」 \end{tabularx} \\ \\ \relax
5:28 & \begin{tabularx}{0.7\textwidth}{X} 他就撇下所有的，起來跟從耶穌。 \end{tabularx} \\ \\ \relax
5:29 & \begin{tabularx}{0.7\textwidth}{X} 利未在自己家裡為耶穌大擺宴席，有一大群稅吏和別的人與他們一同坐席。 \end{tabularx} \\ \\ \relax
5:30 & \begin{tabularx}{0.7\textwidth}{X} 法利賽人和文士就向耶穌的門徒發怨言說：「你們為甚麼跟稅吏和罪人一同吃喝呢？」 \end{tabularx} \\ \\ \relax
5:31 & \begin{tabularx}{0.7\textwidth}{X} 耶穌回答他們：「健康的人用不著醫生；有病的人才用得著。 \end{tabularx} \\ \\ \relax
5:32 & \begin{tabularx}{0.7\textwidth}{X} 我不是來召義人悔改，而是召罪人悔改。」 \end{tabularx} \\ \\
[1ex]
\hline
\hline
\end{longtable}
$^{1}$提交文件:環境事務及庫務局局長的發言 - 由梁美芬議員發言.
自上次會議舉行至今發出的資料文件.
請各位一起禱告.
親愛的天父上帝我們來到你的面前.
一起學習你的話語.
求你用你的話語來教導我們每一位.
讓我們聽的 說的 都能夠有你的恩典.
一起來學習 以及一起來領受.
主也求你賜福.
讓一群為我們辛苦來事奉.
主也得到你額外的私恩.
我們仰望你奉主耶穌基督的聖名.
阿們.
這段聖經我給大家一個簡單的分段.
今天跟大家分享.
想看一看可以看一看這個分段.
我想跟大家分享另外一個是什麼呢.
另外一個就是.
這個是猶太人的航班的飛機餐.
好像下一次搭飛機都要要求這個餐.
好像很豐富.
也有很多不同的色彩.
很吸引.
這個應該不是當日利美.
利美另一個名字叫馬太.
書寫馬太 福音的各位.
應該不是吃這些的了 他們當日.
不過這個就是現在猶太人他們吃的餐.
所以在飛機當中如果是猶太人就可以要求這個餐.
我們說的這件事.
耶穌基督就去到瑞關那裡的地方.
呼召了利美來跟隨他.
這個場景裡邊就是瑞關.
瑞關不是在灣仔.
他的瑞關應該是在罵市當中.
不是一個特定的辦公室.
大家去交費.
而是一個大家走過會看到的地方.
方便收稅.
收稅是沒有一個叫做薪酬的制度.

$^{41}$為什麼無心地要去收稅呢.
由羅馬發出的一個收稅的制度.
當時的稅利他們只是做了這個位置的時候.
就要在那個地方代表羅馬政府去收稅.
然後把錢交給羅馬政府.
為了方便羅馬政府就很聰明.
不會用羅馬人去不同的地方去收稅.
他讓當地人.
他同民族的人.
去向同民族的人收稅.
所以這個人也是猶太人.
但是做了收稅的人.
因為他做收稅的人.
所以很多猶太人就不高興他了.
覺得他是一個罪人.
你向我們自己民族的人.
拿錢去給一個不是我們民族的羅馬人.
羅馬的政府.
所以就視他是罪人.
這是第一個原因.
覺得他是一個罪人.
第二個原因.
羅馬政府每一個地方他需要收稅.
他怎麼收呢.
這一次在猶大的城裡.
猶大的城就要收多少錢.
加利利的城就要收多少錢.
他是定了額的.
他定了額要收多少錢.
稅你要收多少錢就稅你自己去定.
總之你每年要交足稅.
定期要交足稅去羅馬政府那裡.
其他你收到的就是你的工資.
所以是沒有一個薪酬機制.
但是有錢收.
而且被視為收到的錢是挺豐富的.
因為沒有限定.
你收到就行了.
但是如果你收不夠.
你就要補足.

$^{81}$如果你還想坐那個位.
否則的話就不用坐那個位.
因為羅馬政府要收錢.
不是因為你收不足.
政府可以豁免你官員交的稅是不足.
所以他就要代表羅馬政府.
在猶太人的地方去收稅.
他要收稅的時候.
每個人就會覺得政府都沒有收那麼多.
你為什麼收我那麼多.
所以更加視他們為罪人.
也不喜歡他們.
這是第二個原因.
為什麼把稅當作罪人.
現場當中耶穌路過.
呼召利美的時候.
除了路人之外.
一定有一些耶穌在.
有利美這個人在.
也有一些和他一起的官員.
有一些耶穌基督之前呼召了的門徒.
還有文士和法利賽人.
文士和法利賽人都經常在鬧事.
其實一些重要人物都會鬧事.
他們就一起在這個地方.
耶穌基督應該都在這裡看了一會.
看著利美怎樣收.
過了一輪.
耶穌就去和利美.
他去到利美面前不是去交稅.
他就跟她說.
來吧 你跟著我來吧.
應該收稅的人就說.
來吧 你給稅吧.
他竟然說不是.
這個利美做得不是很足夠.
他反而說.
哦 那我跟著你去吧.
他就把所有東西都放下.
其實如果這幅畫.

$^{121}$有一個現代人去到當地的時候.
有時光隧道去到.
有現在的智能電話.
應該打幾張卡.
然後就放在Instagram.
為什麼呢.
稅務員不收稅.
工作時間擅離崗位.
沒有收妥收回來的稅項.
他明天就不用上班了.
但他可能都沒有想上班.
他就跟著耶穌去.
撇下所有的.
他真的就這樣放下了.
走了.
起身就跟著耶穌.
這是什麼原因呢.
耶穌怎樣令到這個人.
怎樣令到這個人.
正在工作 正在開工.
叫他一聲.
他就走.
跟著耶穌走.
耶穌的衣著並沒有帶著皇冠.
沒有拿著金像.
可以戳一下戳一下.
是什麼吸引他.
比他桌上的金錢更加吸引.
比他的官職更加吸引.
但利未就是撇下所有的.
跟隨了耶穌.
還做多了一件事.
他自己拿錢出來請耶穌吃飯.
他是大排筵席.
還請了一班人來.
這班人不是民事.
不是法理賽人.
我想他請民事法理賽人不肯來.
但他請了一班會來的.
這班人都是跟他一起做稅吏.

$^{161}$還有一些所謂別的人.
都是被視為罪人的人.
你陪耶穌一起吃飯.
跟耶穌一起吃飯.
但很多人在.
有一大群稅吏都會來.
我想利未這個人是不錯的.
他很多同事.
可能是他的上司或者下屬.
都會來一起吃飯.
都接受他的邀請.
他來了.
吃飯了.
法理賽人和民事法院.
我想他們沒有去到利未家.
我將他視為第二天.
或者遲些發生的事.
他們知道這件事之後.
耶穌去了利未家.
一個稅吏的人家裡吃飯.
不止一個稅吏.
還有一大群稅吏的人.
跟他一起吃飯.
跟一些罪人吃飯.
你怎麼可以這樣.
他跟耶穌的門徒說.
他問耶穌.
為什麼我們要這樣做.
他跟耶穌的門徒說.
為什麼你們會跟一班稅吏.
跟一班罪人吃飯.
耶穌的回答.
是另一個層次的.
我不是去吃飯.
我是去工作.
健康的人不需要醫生.
我去的是一班身心靈有不舒服的人.
我來不是為了義人.
我是為了罪人.
我來就是招罪罪人.

$^{201}$對了.
利未這個稅吏.
坐在關口也不是一天.
這麼多法利賽人民事經過.
這麼多去交稅的人經過.
有誰叫他去教會呢.
耶穌去就說了一句.
來吧.
跟我來吧.
他連錢都放在那裡.
就走了.
其實他真的很預備好了.
看誰來叫他.
不知道.
但耶穌叫他他就跟了.
我不知道民事叫他.
他會不會跟了呢.
但耶穌來叫他.
他就跟了耶穌去.
他應該都聽過.
或者見過耶穌做一些事.
知道耶穌是誰.
他為人怎樣.
耶穌叫他就起身.
跟著去.
不是這樣.
還馬上做了一件事.
那天就宴請耶穌.
跟著就不只是跟耶穌聊天.
他大排筵席.
叫了很多他的同事.
很多他的朋友一起來.
他只是我撿到寶了.
我把這件寶給你們看看.
你們也要來撿到寶.
他是我得到好處.
我認識了一個視我為朋友.
叫我跟隨他的拉比.
你們也來認識拉比.
他也像是拉比你呼召了我.

$^{241}$叫了我跟隨你.
我還有一群朋友.
他們也需要你.
他像是要做一個中間的媒介.
讓耶穌都能夠向他其他的朋友.
同事都可以做一些事.
不要說他們當晚做了什麼.
但耶穌基督就是這樣回答.
他說.
對不起.
回答他就是說.
我是去為了一些罪人.
要令他們得悔改而去.
我不是為了吃飯而去.
我是為了心靈有病的人而去.
我不是為了健康的人.
我們去暢飲.
很開心地吃東西.
我去是有原因去.
猶太人他們有一樣東西.
他們不跟某些人一起吃東西.
外邦人他們不跟他一起吃東西.
有罪的人他們也不跟他一起吃東西.
為什麼呢.
這幅圖畫有一些污水和清水不可以共流.
污水不可以流入清水的池裡.
這是廣告.
在他們的觀念當中.
其實猶太人他們早就有這種觀念.
清水會因一滴污水.
而整個清水變得污糟.
清水不能因為一滴清水而變得清純.
所以他們要避開污水.
一起吃飯一起聊天.
猶太人吃飯他們的東西會一邊吃一邊吃.
大家現在經過疫情之後.
如果喝茶需要一碟辣椒醬.
可能變成一人一碟.
或者自己舀開碗裡.
他們很重視.

$^{281}$我們一起用這碟醬汁.
要蘸這個醬汁來吃東西.
但是有罪的人他蘸過的東西都污糟.
因為他是罪人.
他蘸過的醬汁都污糟.
令我變成罪人.
於是他們跟罪人要分開.
好像一個家庭.
現在如果有新官的時候就會分開.
他們就要分開.
他們分得有多開呢.
他們分得很開.
就好像我們很嚴重時期的新官病症的時候.
不單是隔著一塊膠板.
可能隔幾張桌子才有客人.
大家隔得很遠.
他們是這種心態去面對罪人.
更加是以這種心態去面對外邦人.
如果是外邦人他們吃豬的.
他們可能吃的時候又說話的時候.
他們噴了出來有豬汁.
有豬汁噴了出來的時候.
就污染了我們.
不是那塊豬肉.
所以他們吃飯不會跟外邦人同一胃.
同一胃都令他們覺得我污穢了.
他們有這樣的情況.
所以法理賽人和文士對耶穌.
什麼不是飲食.
你經常都想飲飲食食.
你都不禁食.
耶穌基督就告訴他.
什麼時候才禁食.
禁食是有時刻的.
原來法理賽人他們逢禮拜一和禮拜四.
禮拜一和禮拜四他們就禁食.
其他日子他們照常吃東西.
不過禁食的時候他們喜歡裝作告訴別人.
我今天禁食.
耶穌基督都教門徒不要這樣做.

$^{321}$但是他們還有一條條例.
如果你禁食那段日子.
有親友結婚.
你是不用禁食的.
你是可以參加婚宴.
跟新郎他們一起飲宴.
將禁食的日子改期.
禁食是可以改期.
所以耶穌基督就說.
我在這裡他們就不會禁食.
就好像新郎在這裡他們就不用禁食.
到新郎被帶走的時候.
他們就要禁食.
其實耶穌基督是預言有一天他會被帶走.
耶穌基督末世的時候就說.
耶穌基督就是新郎.
教會我們每一位就好像是耶穌基督的新娘.
他預備來接領我們去到天國.
耶穌基督就說去到那個時候.
新郎會被帶走的那段日子他們就要禁食.
這句話他對著這班文士和法理賽人說.
他可能知道一點又不知道一點.
就是有新郎就不用禁食.
不知道一點的就是.
你當自己是誰啊耶穌.
但是耶穌基督都對著他們說.
因為他來不是為了健康的人.
是為了有罪的人.
在這班人當中.
法理賽人和文士.
耶穌基督並沒有省住沒有向他們傳講.
耶穌基督都跟他們說.
明天新郎會被帶走的.
當那個時候如果他再想起這句話的時候.
在耶穌基督死而復活之後.
他們應該悔改的.
耶穌基督沒有放棄法理賽人和文士.
反而其實這個宴會當中.
耶穌基督向了一班文士和法理賽人.
視為罪人的人.

$^{361}$跟他們親近了.
一起藉著飲食的場面.
可以跟他們聊天分享.
可以幫他們做一些問題的處理.
困境的處理.
關心他們的需要.
但是面對著這班人.
這一班是文士和法理賽人的時候.
他都是藉著被他們質疑質問的場景.
都向他們傳講了福音.
如果他們領受到.
他們是得到的.
這個筵席.
利美為他們擺的這個筵席.
利美可以更加清晰福音.
他都開始了.
他可以帶領人去親近主這個使命.
做耶穌基督的門徒.
另一方面這個筵席.
也釋放了法理賽人和文士的質疑.
讓文士和法理賽人都有機會聽聞福音.
讓他們有機會知道.
耶穌基督是誰.
來做什麼.
這一班人.
這一班被視為罪人的瑞麗.
他們需要什麼.
他們不是需要被隔開.
他們需要身心靈的藥物.
他們需要被關心.
他們需要福音.
既然文士和法理賽人.
他們是敬虔的族類.
是很熟悉聖經的人.
應該他們承接著這一棒.
繼續來向這一班罪人.
這一班瑞麗來傳講福音.
關心他們的需要.
就算不和他們一起吃飯.
對他們的態度和和他們相處.

$^{401}$都應該有改進.
不過似乎他們沒有這一種心.
所以耶穌基督再和他們講一個比喻.
耶穌基督和他們講一個比喻.
就是如果衣服破了.
你不會去買一件新的布回來.
剪一塊去補舊的衣服.
這樣的話新的布就會縮水的時候.
將撕爛的情況變得更加不可收拾.
亦都釀酒的時候.
不會將新酒放在舊的皮囊裡面.
因為舊的皮袋已經膨脹過了.
不會再膨脹.
它已經變得硬又脆.
新酒下去.
新酒會膨脹.
它會發酵.
在這個情況下.
舊的皮袋就會爆.
新酒就會浪費.
他這樣講.
其實這段話不是要告訴利美他們.
反而他的聽眾.
他第一個聽眾就是文士和法利賽人.
對著這一群來質疑耶穌基督的人.
耶穌基督竟然是存在一個教導他們的心.
猶太人教導的方式.
時時都是用比喻的.
他們生活當中就是用比喻去教導.
對著他們的子女他們也是.
耶穌基督都是用比喻跟他們說.
其實他們很聽得懂比喻.
他們習慣了聽比喻.
特別是文士.
他們更加熟悉.
但是似乎他們真的沒有去明白.
不願意去接受.
耶穌基督是在告訴他們.
新的和舊的不可以參摘在一起.
要有一種新的觀念.

$^{441}$耶穌基督還沒有叫他們要怎樣去新.
但是就告訴他們一件事.
每一天的生命應該有新的改進.
你不能夠每一天都喝舊酒.
你每一年都有新釀的.
舊酒會喝完.
你一定要用一個新的皮囊.
去承載今年的新釀.
新釀有一天都會變舊釀.
你一路要有新的皮袋.
要一路要有新釀.
否則有一天你就會喝盡家裡所有的無酒的了.
人對於釀酒.
猶太人他們很熟.
他們應該知道很明白.
但是總有一件事.
有一件事就是不想更新.
有些時候我們真的會怕更新.
有一個新的制度.
我們都會打個疑問.
和舊的不就很好.
很多東西我們會懷念.
吃東西都會想古早味.
加了些什麼.
為什麼現在味道不同了.
我喜歡以前那個.
但其實不是的.
很多東西你要新的.
買車的時候買一輛新車.
過年的時候穿一件新衣服.
其實不用過年.
換季,特價會買一件新衣服.
未到2024年.
我們就會留意2024年流行什麼時裝.
我們每一樣東西其實都有新的期望.
不過做事或者對人.
我們時時抱著舊的那個.
耶穌基督說你美同化你才華.
你們的生命要成長.
你這麼熟悉聖經.

$^{481}$要讓聖經的話語在心靈當中.
讓自己的生命可以更新.
讓自己的生命更新的時候.
更加認識神.
其實耶穌基督和利美這班.
瑞利吃飯這件事.
既然民事和法理賽人這麼留意.
又知道耶穌基督說.
我去不是為了吃飯.
我是做事.
民事和法理賽人.
都可以想一些做事的方法.
為了關心這班需要被關心的人.
他可以面對瑞利這班罪人的時候.
從另一個角度去看.
好像耶穌基督那樣去看他們.
他們是有需要的人.
他們是身心靈有需要的人.
他們是在社會當中被歧視的人.
用這個角度去看.
去接納他們的時候.
可以幫到這班人.
不但幫到他們.
還能夠令到.
他民事教導的圈子可以闊了.
他可以教導這班.
需要被教導的瑞利和罪人.
難道他真的坐在那裡.
讓罪人和瑞利自己悔改.
而不把神的話語教導給他們嗎?.
耶穌基督去賦予這個筵席.
利未為他預備的筵席.
是一個真的悔改的筵席.
是讓利未可以在那個時候.
告訴他的同事.
他改變了 他悔改了.
而且他要轉行了 他跟耶穌.
他可以讓這一班朋友.
知道為什麼他有轉變.
認識耶穌基督是誰.

$^{521}$耶穌基督也藉著這個機會.
因為瑞利這件事.
引起民事和法理錯人的關注的時候.
再教導民事和法理錯人.
如果他們願意.
都可以悔改.
藉著這個筵席引起的原因.
令到他們可以悔改.
不過可惜.
罪人聽了的.
有些應該是悔改了.
那晚他們可以很歡宴.
去聽到主耶穌基督的教導.
但是民事和法理錯人似乎聽不明白.
他們似乎每次都要聽不明白.
他們好像守著原生古時的道理.
以至於耶穌基督要說他們.
你們好像要把新的布.
剪一塊出來補到舊的那裡.
令到新的又被你剪爛了.
舊的又用不了.
他們持守著一些以往.
他們覺得很重要的東西.
但是不要求進步.
其實舊的不一定是不好.
但是我們需要進步.
在信仰上面.
民事和法理錯人他們是一個教導的身份.
他們也很熟習聖經.
更加需要進步.
進步到能夠看到神的那種愛.
看到神怎樣來去.
藉著牛羊代替人去獻祭的時候.
而赦免了人的時候.
去明白到神的愛.
看到耶穌基督為了罪人.
而和他們歡宴的時候.
去看到神的那種愛.
都去將那種愛在他們當中.
來去伸展延展.

$^{561}$讓每一個人都能夠得到.
若果民事都能夠彼此相愛.
法理錯人能夠彼此相愛.
還有去愛他們身邊不是民事.
不是法理錯人的人.
可能整個社會都會有轉變.
這個筵席是為了愛的緣故.
帶來悔改.
我想大家都會的.
待會崇拜完之後.
大家都會有一個歡宴的時候.
這個歡宴的時候.
都是大家可以分享.
可以關心代祖的時候.
可以彼此支持鼓勵的時候.
這個歡宴的時候.
是可以約一些.
以前有一起回來崇拜.
現在少見了他們的人.
崇拜完一起吃飯.
你今天沒有來嗎.
不要緊 吃完飯你下次再來.
下個星期一起崇拜.
是可以我們關心人的時候.
是可以我們去關心一些.
他未能夠妥善去處理.
星期日膳食的人.
是可以我們去關心一些.
家人可能有疾病.
他們不能夠得到好好照顧的.
我們可以吃飯的時候.
預備一份帶到家給他們.
可以關心一些不能夠利到灼上.
和我們一起飲食的.
可以在你放假的時候.
和一些未認識神的人.
吃飯喝茶.
介紹一下他.
我們七月八月都有報道會.
我們有父親節.

$^{601}$家庭崇拜的慶典.
我們有不同的聚會.
我們接下來十號星期六.
我們有一個很奇妙的音樂會.
你可以在飯桌上.
將你的愛言展下去.
讓到未得到赦免的人.
得到悔改的機會.
求主幫助我們.
你也可以因著你的才幹.
去為在神預備給你的地方.
再將那份愛擴展開去.
又藉著你忙碌的時間.
再去計劃一下怎樣使用.
好好去看聖經.
讓聖經成為你的提醒.
每一天讓你的生命.
可以再有更新.
也讓著你的話語.
你得到的話語.
你可以幫到一些.
都需要神話語的人.
有沒有一些福音對象.
你可以每一個星期.
送一節金句給他.
試一下.
一年也只是五十多次.
不是很多.
我們一起祈禱吧.
親愛的天父上帝.
我們來到你的面前.
我們很想一件事.
我們知道主耶穌基督來.
是為了我們罪人.
為了我們身心靈.
都有需要的人.
主耶穌多謝你.
讓我們能夠被赦免.
多謝你.
讓我們能夠逐漸變得健康.

$^{641}$主耶穌我們有時候.
會有需要求主你.
醫治我們.
我們又求你看顧.
主耶穌我們很需要你的看顧.
讓我們能夠.
帶著你給我們的恩典.
去關心和我們身邊.
不是最健康的人.
不是完全得到赦免.
還在最終生活的人.
讓他們能夠來到你的面前.
得著你的愛.
得著你的恩典.
得著你的照顧.
求主憐憫.
主我們不是最美好的器皿.
但是求你使用我們.
求你用我們.
讓我們有機會服侍你.
和你同工.
去關心有需要的人.
服侍有需要的人.
祈禱仰望你.
奉主耶穌基督的聖名.
阿們.
阿母.
\newpage



\section{路加福音 5:27-39}
\label{sec:KzM_XrVrUvg}
\textbf{2023-06-04 早堂《悔改的筵席》路加福音 5\_27-39 講員:曾裔貴牧師}
\newline
\newline
連結: \href{https://youtube.com/watch?v=KzM_XrVrUvg}{\texttt{https://youtube.com/watch?v=KzM\_XrVrUvg}} ~~~~ 語音日期: 2023-06-07
\newline
\newline
\hyperref[sec:jJiE2EQ4f48]{\small{< < < PREV SERMON < < <}}
~
\hyperref[sec:index_chronic]{\small{[返順時目]}}
~
\hyperref[sec:index_scriptual]{\small{[返順卷目]}}
~
\hyperref[sec:qmt5oGwLkF4]{\small{> > > NEXT SERMON > > >}}
\newline
\newline
路加福音 5:27-39
\newline
\begin{longtable}{cl}
\hline
\hline
章節 & 經文 (和合本修訂版)\\
\hline
5:27 & \begin{tabularx}{0.7\textwidth}{X} 這些事以後，耶穌出去，看見一個稅吏，名叫利未，在稅關坐著，就對他說：「來跟從我！」 \end{tabularx} \\ \\ \relax
5:28 & \begin{tabularx}{0.7\textwidth}{X} 他就撇下所有的，起來跟從耶穌。 \end{tabularx} \\ \\ \relax
5:29 & \begin{tabularx}{0.7\textwidth}{X} 利未在自己家裡為耶穌大擺宴席，有一大群稅吏和別的人與他們一同坐席。 \end{tabularx} \\ \\ \relax
5:30 & \begin{tabularx}{0.7\textwidth}{X} 法利賽人和文士就向耶穌的門徒發怨言說：「你們為甚麼跟稅吏和罪人一同吃喝呢？」 \end{tabularx} \\ \\ \relax
5:31 & \begin{tabularx}{0.7\textwidth}{X} 耶穌回答他們：「健康的人用不著醫生；有病的人才用得著。 \end{tabularx} \\ \\ \relax
5:32 & \begin{tabularx}{0.7\textwidth}{X} 我不是來召義人悔改，而是召罪人悔改。」 \end{tabularx} \\ \\ \relax
5:33 & \begin{tabularx}{0.7\textwidth}{X} 他們對耶穌說：「約翰的門徒常常禁食祈禱，法利賽人的門徒也是這樣，惟獨跟你在一起的又吃又喝。」 \end{tabularx} \\ \\ \relax
5:34 & \begin{tabularx}{0.7\textwidth}{X} 耶穌對他們說：「新郎和賓客在一起的時候，你們怎麼能叫賓客禁食呢？ \end{tabularx} \\ \\ \relax
5:35 & \begin{tabularx}{0.7\textwidth}{X} 但日子將到，新郎要被帶走，那些日子他們就要禁食了。」 \end{tabularx} \\ \\ \relax
5:36 & \begin{tabularx}{0.7\textwidth}{X} 耶穌又講一個比喻，對他們說：「沒有人把新衣服撕下一塊來補在舊衣服上，若是這樣，會把新的撕裂了，並且所撕下來的那塊新的和舊的也不相稱。 \end{tabularx} \\ \\ \relax
5:37 & \begin{tabularx}{0.7\textwidth}{X} 也沒有人把新酒裝在舊皮袋裡；若是這樣，新酒會脹破皮袋，酒就漏出來，皮袋也糟蹋了。 \end{tabularx} \\ \\ \relax
5:38 & \begin{tabularx}{0.7\textwidth}{X} 相反地，新酒必須裝在新皮袋裡。 \end{tabularx} \\ \\ \relax
5:39 & \begin{tabularx}{0.7\textwidth}{X} 沒有人喝了陳酒又想喝新的；他總說陳的好。」 \end{tabularx} \\ \\
[1ex]
\hline
\hline
\end{longtable}
$^{1}$請姐妹主女們平安.
和大家一起有祈禱吧.
多謝主.
我們一起打開聖經.
很盼望這個延籍的故事.
成為我們一個反思的空間.
特別在這個時代.
我們需要學習主耶穌.
對別人的呼喚,呼召和一個接納,拯救.
我們仰望在你的面前.
禱告奉耶穌基督的名.
阿們.
想大家留意第十七節.
和今天我們要講的悔改延籍.
有相當大的關係.
有一天耶穌教訓人.
有法利賽人和教法師在旁邊坐著.
他們是從加利利各鄉村.
猶太並耶路撒冷內.
即是全個以色列的法利賽人.
文士.
教導律法的教法師.
留意一下這裡路加醫生.
特別顯著提到突然間有一樣東西.
其實是不是當時的人感受到呢.
很特別的就是.
這麼多的法利賽人坐在他旁邊.
但是當他要講以下的神蹟之前.
就講了一句話.
是路加醫生特別強調的.
就是主的能力與耶穌同在.
使祂能夠醫治病人.
這個就是段落的開場白.
就講到之後有一個癱瘓的人.
有四個好朋友抬到耶穌面前.
耶穌第一句話是發人心醒.
出人意表的.
所以立即引起法利賽人那種敏感的神經.
那句話是什麼呢.
小子你的罪赦免了.

$^{41}$一個癱瘓的人完全行不通.
即是說他不良於行的殘缺.
耶穌看到的不是他的可憐.
不是那四個朋友的熱心.
是對這個病者一個呼喚.
好像一個應許的命令.
現在來到發生.
就是小子你的罪赦了.
這個就是第二十節.
耶穌見他們的信心.
那四個朋友的信心.
就對癱子說你的罪赦了.
真是挑動了法利賽人的神經.
文士和法利賽人就議論.
這個講淺妄話的人是誰.
當然知道是耶穌.
除了神以外.
誰能夠赦罪呢.
就引進討論.
這個其實不是今天集中要講的經文.
是27至35節.
之後不夠時間再講新郎.
和新舊報的比喻.
這裡告訴我們.
有新的一個段落.
就是這個癱子好之後.
很明顯一個神蹟在他生命裡面發生.
他到底是感受到那個能力是雙重的.
是他的罪赦了.
還是他的身體完全復原呢.
相信大家都不需要想太久.
一定是知道是後者.
他的身體有一種新的能力.
進入他裡面.
馬上就好.
馬上就拿著玉紙自己離開.
和一路走歸榮耀給神.
然後眾人見到這個場景.
又是歸榮耀給神.
路加醫生特別強調.

$^{81}$主的能力進入這個癱子的生命裡面.
固然就是醫治的能力.
更重要就是這個癱子都感受到.
主耶穌對他的生命的接納.
如果一個這樣的人.
尚且能夠起來行走.
尚且能夠歸榮耀給神.
路加很想告訴我們.
還有另外一個人.
這個人那種起來.
更加巨大的能力進入他的生命裡面.
這個是當時在社會上.
當時耶穌的年代.
福音書年代的一個階層.
叫做瑞利階層.
是聲名狼藉的.
是普遍的猶太人都不喜歡的.
大家如果知道少少背景就會明白.
因為他們是替羅馬人打工的.
收自己的猶太同胞不同的稅項.
稅收本來就定了.
是不能更改的.
羅馬人是要多少就是多少.
但是因為是用高價買來的官職.
所以羅馬人也都是容許.
他們在例定的數目以外.
自己喜歡收多少就收多少.
所以瑞利往往就是內心良心.
都會有時會不安的.
因為有時會收多.
那個酌情權是由他自己來行使的.
所以很多時候都會用一個強權.
來壓制自己的同胞.
所以久而久之.
這個行業就是被猶太人所唾棄的.
不知道在座有沒有澳門的居民.
有澳門籍的.
你來不及回過澳門的.
你想想如果你家裡很多人是做賭場的生意.
或者在賭場裡面工作.

$^{121}$可能都未必是這麼光彩的.
雖然賺到很多錢.
如果你是放那些疊碼的.
可能更加會令到其他的居民.
對你會有一些意見.
但是耶穌的時代就是這樣.
我們看看這段剛才主席已經讀了的經文.
有一個人有一個很大的能力.
進入他的生命裡面.
他能夠立即起來跟隨耶穌.
但是招致的是很大的反對.
第27節.
這是以後耶穌出去看見一個瑞尼.
名叫利美.
利美是猶太希伯來的名字.
利美就是聯合的意思.
與神聯合.
他有另一個名字叫馬太.
不知道有沒有叫Matthew的弟兄.
就是馬太.
他就是另一個名字.
他在這裡告訴我們他是一個猶太人.
很肯定他是一個猶太人.
不過他現在在上班的時間.
坐在稅關收稅.
即是低級官員.
因為路加分第19章.
有一個在耶利哥的稅務長.
很高級的.
所以這裡告訴我們利美.
他坐在稅關上如常的工作.
收取路過的人的稅.
雖然沒有凸顯是甚麼類型的稅.
總之他現在在工作裡面.
如是者繼續生活.
但是當耶穌今次經過的時候.
刻意走到他的稅關那裡.
可能是一個奸廷.
就跟利美說.
你跟從我來.

$^{161}$第28節.
其實我們要思想一下.
聽了一句說話.
有一個陌生人跟你說了一句話.
你就撇下所有.
起來跟從耶穌.
有一個很大的力量.
令到利美好像不假思索.
要將他現在做的事.
不是暫時吃個下午茶.
休息一下3.3.
不是.
是撇下起來跟隨.
是一個連續的動作.
即是說他現在是跟他的過去了斷.
要過一個新的生活.
是跟隨耶穌的新生活.
因為第5章是說耶穌.
開始出來傳道.
去呼召門徒跟隨.
第一節已經說了.
彼得約翰跟隨耶穌.
又是撇下所有起來跟隨耶穌.
但那些是打魚的基層.
撇下的是魚網.
撇下的是魚船.
是正當職業.
但是利美.
沒錯.
在羅馬人眼中是正當職業.
在他自己的良心中也是正當職業.
不過在他自己過去做過的一些事情.
絕對不能夠被自己的同胞所包容接納.
所以他自己的良心.
總是會有一種有愧在他自己的內心.
但在這一刻.
有一個時間.
耶穌來到他的面前.
正在工作的時候.
他有很大的一份能力.

$^{201}$進入這個不是身體有殘缺.
是他的心靈.
沒有力量面對自己的生活.
他的生命.
和跟同胞的關係.
但是主的能力進入他裡面.
主耶穌這句話.
是給他一個很肯定的改變的開始.
就是他能夠跟隨耶穌.
進入一個新的生命.
接著我們看到第29節.
利美在自己的家裡面.
為主耶穌大擺筵席.
有錢的人.
他現在決心開一個頌讚感恩會.
邀請的是耶穌.
更重要邀請的是一大群的瑞利.
就是他的WhatsApp群組.
你可以肯定知道.
他的瑞利群組.
他就發出邀請.
哪天哪天的晚上.
我邀請了主耶穌一同來坐席.
邀請你們一起來見證.
我會在現場講我的見證.
我怎樣脫離瑞利的行列.
我現在不是派散水餅.
我是邀請你們一同.
感受一下耶穌的這份能力.
在我的生命裡面的改變.
一個坐席的歡樂.
這是盧家寬第一次提到筵席.
所以我們在未來都有好幾次.
應該有兩個月的時間.
我們會集中講.
在寬恩書裡面的筵席.
今天是第一講的悔改的筵席.
盧家福音作者盧家.
很有興趣對瑞利這個行業的人的留意.
有四次提到瑞利.

$^{241}$都是很正面的.
這一次第五章的利未悔改.
第十五章.
聖經告訴我們.
眾瑞利和聚人.
來去哀求耶穌.
要聽他講道.
然後快來錯人就不高興了.
這個人和那些瑞利聚人為伍.
吃喝.
十五章的第一第二節.
所以耶穌就講了三個比喻.
失淺失陽.
和小兒子的失落.
離開爸爸的父家.
是在講那個尋找.
所以在主耶穌的眼裡.
瑞利是被尋找回來的小兒子.
去到第十八章.
有一個很重要的關於祈禱的教導.
法利賽人和瑞利的祈禱.
十八章第九節.
有兩個人去到聖殿祈禱.
法利賽人就一路祈禱的時候.
就是說我整個禮拜都沒有違背律法.
我奉獻十分之一.
我怎樣怎樣的.
我禁食.
耶穌說另外一個人祈禱.
就是瑞利.
他祈禱的時候.
遠遠的低頭.
不敢望聖殿.
然後就隨空去求神.
哀聲的懇求.
神赦免我這個罪人的不義.
這樣的禱告.
似乎主耶穌很正面看瑞利這個階層.
當然剛才說了十九章.
有這個薩該瑞利獎的信主.

$^{281}$就是他很想見一下耶穌.
但周圍的人都不喜歡他.
阻礙他.
他爬上一棵桑樹上.
想看看耶穌經過.
但耶穌見到他就抬頭跟他說.
薩該你下來.
我今晚就住在你家.
薩該很開心迎接耶穌.
然後那些法利賽人又在吵.
他立即怎樣說呢.
我將我的一半的身家.
來分給窮人.
我曾經訛詐過誰的.
我要還他四倍.
主耶穌說今天救恩來到這個家.
人子來為要尋找.
拯救失喪的人.
四次不是少了.
盧家福音的作者很有興趣告訴我們.
哪些人越快承認自己的不足.
承認自己的軟弱.
越快得到救恩帶來的歡樂.
你見到利微的歡樂.
薩該的歡樂.
你見到一個哀求的瑞利.
可以在神的面前得到接納.
十五章又是一個歡樂的言籍.
又是說眾瑞利和聚人.
來哀求耶穌.
要聽耶穌講道.
所以這裡告訴我們.
主耶穌有一個很大的能力.
就是醫治新心靈的能力.
所以如果耶穌只是醫治疾病.
沒錯他是一個神醫.
但更多更多的是心靈的醫治.
更多的是主耶穌對這個已經被拒絕的階層的人.
來接納歡愉和諒解.
使利微的計算很厲害.

$^{321}$誰要他多少歲就由他說了算.
這個婆婆這麼可憐.
算了收沒那麼重.
這個多收一點.
什麼都是他說了算.
總之他每個月交給羅馬人的稅收.
其他就任由他.
還有就是他要維持這個官階.
他就要繼續賄賂.
所以很多很多的問題都是由此而衍生的.
久而久之自己是有錢了.
不過就沒有了朋友.
沒有了能夠被自己的同胞接納的工作.
但是這個耶穌特意來找我這個人.
他難道不知道這個行業的人.
是被猶太的拉比所拒絕的嗎.
法利賽人時時都說.
耶穌你知不知道的 你懂不懂得看的.
看一看就知道這個人是稅吏來的.
稅吏和別人 稅吏和罪人 稅吏和娼妓.
是這樣來的去同義詞的.
好像要提著耶穌.
這些人其實不應該這麼快把他納入我們的聖殿.
好像要做一些品格檢查.
耶穌好像不是這樣理解.
耶穌感受到的就是他的能力進入他裡面.
這個人就有力量改變了.
這個稅吏利微.
他就有一種愛的力量進入了他的內心.
他就會改變.
我們星期六晚.
現在我們在完了崇拜之後.
逢一三週.
我們的大團小組.
越來越多人.
而且是回來一兩年都不夠的新朋友.
來我們錦宣家.
慢慢他們感受到教會對他們的接納.
一個一個地越來越開始想分享.
願意分享他們的故事.

$^{361}$感受到他們在這個群體裡面被接納.
被接納之後他們慢慢就自己很願意去分享.
他們的信主 他們的故事.
弟兄姊妹教會是一個時時要準備筵席的一個地方.
一個是能夠打開門.
是讓人得到接納.
被一份能力來改變的一個地方.
相信我們要來學習主耶穌一樣.
對這些已經好像被人label了.
沒得救了.
稅吏為錢.
他們可以無惡不作.
在法利座人的內心.
這班人注定要下火壺.
但是耶穌那種改變的能力.
可以大到一個地步.
就是這個人日後成為一個很重要的猶太的一個重要的寫作的一個人.
在主耶穌的十二使徒裡面.
這個人是有份.
從來在福音書沒有出現過一次他講話的.
但是他默默聽著耶穌的教導.
跟隨耶穌的十二使徒的其中一個.
這個就是曾經有長時間坐在稅關的利微.
主的能力進入他的內心.
使他能夠起來撇下跟隨耶穌.
親愛的弟兄姊妹.
教會的牆.
不只是拆了四面牆.
教會的牆.
是要讓人看到.
我們真的會接納不同的人.
當然這個過程可能是有風險的.
因為有人還沒有認識大愛之前.
他會有一些可能有的企圖.
例如我很想我的兒子入名教.
基督教.
我就回教會.
曾牧師會幫我寫推薦信.
或者我很想回教會認識朋友.
認識女孩子.

$^{401}$我又回教會了.
又或者諸如此類.
有些人再想到就是.
我希望我的家人將來離開世界.
可以葬在基督教的墳場.
這些全部某程度是基督教帶來的副產品福利.
但是人有時只是看到表面的福利.
看不到的就是其實神的能力.
進入你的生命改變你.
使你能夠與神相和.
和在一個群體裡面.
大家一起去幫忙改變不同人.
他們的生命.
就是稱為悔改的筵席.
不過我們看到第30節.
法利賽人和文士就向耶穌的門徒發怨言.
你們為何要與稅吏並罪人一同吃喝呢.
猶太人特別拉比的傳統就是.
悔改應該與披麻蒙灰禁食.
拉上一個等號的關係.
為何要這麼開心快樂的筵席.
還要與這些階層稅吏.
一群稅吏坐直呢.
耶穌不會被影響嗎.
還是耶穌可以影響到他們嗎.
法利賽人就不相信這件事.
認為會受玷污身份.
身份是不值錢的.
改變這群人才是真正的重點.
一個拉比.
他會覺得我去這些場合.
就會玷污我的身份.
所以這些傳統.
或者這些制服.
會令到一個人有了階級.
有了一個錯誤的迷思.
以為要與這些人割席.
但主耶穌讓我們看到.
有能力進入這些人的內心.
就有改變.

$^{441}$有改變的可能.
有改變的機會.
使利微.
第31節耶穌告訴我們.
好像要幫利微做一個開脫解釋.
因為我是醫生.
不單止身體的疾病得到醫治.
而且是心靈的疾病.
而且這句話與大相關的對法利賽人.
是一個當頭棒喝的指控.
你們健康.
你們以為宗教的外衣.
使到你健康.
其實你離開神越來越遠.
有病的人才會感受到醫生的愛.
所以瑞利在盧卡醫生的眼裡.
寫作的裡面.
一而再強調.
這個行列原來是很近神的.
很得到神的醫治.
因為他們知道自己是被社會所氣絕的一個行列的人.
所以他們是有病的人.
耶穌說我就是醫生.
耶穌說我就是新郎.
新郎在筵席裡享受快樂.
所以弟兄姊妹你們是不是已經很久沒有笑了.
來到教會的弟兄姊妹都不笑.
你怎麼期望世界的人會笑呢.
來到教會在主耶穌的大筵席裡.
我們都不是從心裡有一份平安喜樂.
而是帶著憂愁來教會.
所謂何事呢.
為什麼呢.
是不是很多恩怨情仇.
是不是因為有很多你的標籤.
你的包裝.
你的宗教的俄衣.
令你不能夠喜樂.
在這個筵席裡.
那你要問一個問題.

$^{481}$到底你是有病還是健康.
耶穌說你們這些健康的人就不需要醫生了.
不是要說一些這麼負氣的話.
而是給我們一個反思.
每次耶穌和這些利微.
瑞利一起的時候.
快利卓人就跳出來質問.
無論是耶穌.
或者是耶穌的門徒.
正正因為他們慢慢被自己的傳統.
被自己的宗教的學識.
慢慢影響了他們的心靈.
是模糊了.
不發現其實每一次.
我們大家都是帶著不同的困難來教會.
主耶穌就是這樣來醫治我們.
我們就能夠互相分享神豐厚的恩典.
所以記得我們能夠事奉是恩典.
我們能夠繼續在教會裡.
仍然都是恩典.
用恩典的眼的視力.
那個眼鏡來看其他人.
我們就能夠充滿了包容.
充滿了接納.
充滿了彼此的相愛.
不需要帶著什麼包袱來看其他人.
這個是利微的筵席.
這是我們每一位敢宣的家園.
我們要來反思的.
否則我們永遠就是一個言直以外的旁觀者.
一些論斷的人.
我們應該就是言直裡的座上客.
我是不同程度的稅吏.
但是主耶穌的能力進入了我裡面.
我就能夠起來.
我就能夠有我的撇下.
有我的跟隨.
我們一起祈禱.
多謝主.
願你在我們大家的心靈裡.

$^{521}$讓教會真的有一份愛的能力.
進入我們大家的心靈裡.
以致我們能夠在教會裡.
我們這個群體.
彼此的寬恕.
彼此的接納.
彼此的來欣賞.
一同去見證.
神豐厚的言直的那份大恩.
\newpage



\section{}
\label{sec:qmt5oGwLkF4}
\textbf{2023-06-08 周四晚查經班《被擄歸回歷史研讀(17)》 講員:曾裔貴牧師}
\newline
\newline
連結: \href{https://youtube.com/watch?v=qmt5oGwLkF4}{\texttt{https://youtube.com/watch?v=qmt5oGwLkF4}} ~~~~ 語音日期: 2023-06-09
\newline
\newline
\hyperref[sec:KzM_XrVrUvg]{\small{< < < PREV SERMON < < <}}
~
\hyperref[sec:index_chronic]{\small{[返順時目]}}
~
\hyperref[sec:index_scriptual]{\small{[返順卷目]}}
~
\hyperref[sec:Fl2IkOLpxSI]{\small{> > > NEXT SERMON > > >}}
\newline
\newline
$^{1}$我們同心祈禱.
多謝主讓我們可以在平安的晚上來到你的面前.
雖然天氣異常炎熱.
很難適應.
而且濕度非常高.
很多地方都沒有香港潮濕滿熱.
但我們想到.
你是創造萬有的.
既然你容許香港這麼富庶密集的地方.
也是非常先進的地方.
但卻這麼炎熱的時候.
都盼望不會讓一些收入比較低的基層家庭.
他們在一種很大的困苦裡.
求主看顧他們.
也保守教會.
我們能夠在這個時代成為一個登台見證.
來祝福周圍不同的人.
用我們可以想到的方法.
來將神的福音有效力地告訴微信的朋友.
使得神的恩典能夠傳開去.
不是讓我們自己去享用.
甚至是享受.
而是我們有責任將福音傳開.
求主帶領我們.
我們透過經文.
特別是提到納若這個課題.
都是我們需要思想的.
盼望神在我們中間帶領.
特別是一些可能在晚上來不了.
但會在網上聽這些記載錄音的時候.
都求聖靈去祝福.
因為可以慢慢聽的時候.
可以聽得明白詳細的.
思想和細著神的話語.
我們今晚等候在你的面前.
我們這樣祈禱.
奉耶穌基督的名.
阿們.
我們一直第八章.
第九第十章.

$^{41}$是一個粉兄會.
似乎是連續七天.
因為是主棚節.
由七月十五.
一直到七月廿二.
因為有一個特別的日子.
就是主棚節第七天之後的第八天.
是一個更加大的安息日.
在約翰福音.
主耶穌都有守.
約翰福音第七章.
在主棚節最後那天.
也是最大那天.
主耶穌說祂自己就是那個活水.
要去到祂那裡.
就好像流出活水江河那樣.
是指到信祂的人受聖靈.
這裡就是舊約.
尼希米的時代.
和以色列帶來的一個粉兄會.
因為我們說過.
就是城牆終於都來到修建好了.
但是重點就是.
都不是做好了那個硬件.
是神的選民.
或者說這些的摩西之約裡面的約民.
就是以色列人.
他們需要重新.
這一幫的選民.
要對神的真理.
律法.
是要有那個的意識.
所以他們就很好.
召開一個粉兄會.
就是.
真的將神的律法書.
來去中讀.
當然他們強調.
是要講解的.
甚至乎可能要翻譯的.

$^{81}$因為可能有些言語上是不通的.
有些只會說亞蘭話.
他們未必會聽希伯來話.
所以就有一些的利未人.
幫手去做一些的傳遞.
將那個真理講解明白.
透徹的可能用一些例子去解釋.
用一些的案例.
來到去在不同的處境.
將摩西的律法書來解明.
所以讀經的.
以色列.
可能是將那個大方向的原則.
講了.
但是他們會有小組的研習.
有利未人.
分了很多的小組.
這樣來到去有那個查經班.
這樣來到大家一起去研讀的.
所以.
神的話語.
在很多人的心靈裡面.
就產生一個很大的力量.
所以.
去到第九章.
他們有一個的特別的嚴肅會.
就是.
這個的住棚節.
由七月十五.
去到七月廿一.
接著七月廿二.
就有一個特別的那個的安息日.
這個第八日.
就是最大那一日.
但是其實已經完了.
但是第九章.
我們上一個禮拜.
只是去到第二十五節的.
就是說.
去到這個月的廿四日.

$^{121}$有可能就是.
他們又是剛好.
過多兩日就是一個安息日.
所以又是一個可以聚集的時間.
所以.
他們就聚集.
來到禁食.
大家.
好像意猶未盡.
大家都很想.
這個自發性的聚集.
就.
非比尋常.
大家有一個很強烈的意識.
要聚集成為一個大的聚會.
或者他們叫一個聖會.
很神聖的一個召會.
那就.
身穿麻衣.
聚集禁食.
頭蒙灰塵.
不是好像印度那些.
殺那些彩粉.
而是大家.
在一個很心靈裡面.
被神的道.
神的話語.
在過去的七日.
來到去感動.
所以.
不是一個儀式性的一個守節.
是一個自發性的.
一個大家.
很想去親近神的.
一個禁食.
的一個披麻蒙灰.
披麻蒙灰.
他們就不是說.
笨藏那些.
我們中國人就.

$^{161}$很不吉利.
說你無端端.
披麻.
就.
非同小可.
一看就會覺得.
很不老理.
但是猶太人就不是.
他們是一些.
很哀哭的事.
重大的事.
又要急切要祈禱的時候.
他們會有這樣的裝束.
但是.
可以裝扮的.
耶穌說.
你要在祈禱的時候.
有沒有印象.
馬太福音第六章.
主耶穌說.
祈禱要在暗中.
向暗中的父祈禱.
如果你禁食.
是沒有問題的.
不過.
就不要在人面前.
扮得很可憐.
扮得很辛苦.
扮得幾天沒吃東西.
很慘.
而是.
你要梳頭洗臉.
來到去.
不要讓人在明處.
看到你禁食.
就是說.
做了.
都不需要告訴別人.
這樣的一份.
因為是.

$^{201}$心靈與神的一個相交.
但是這裡就不是說.
去到已經失去意義的禁食.
在這個時候.
以色列人.
大家意猶未盡.
被神的道感動.
經過了八天.
現在去到24號.
7月24號.
本來是沒有節期的.
但是自發性的聚集.
很難得的.
如果弟兄姊妹.
很感受到.
教會.
神的道.
那個的感動.
很自發的.
例如很早回來.
自己安靜.
自己來到去.
等候神.
自己在那個聚會裡面.
默默的為聚會祈禱.
為講者祈禱.
為附會的來到去祈禱.
如果很多弟兄姊妹.
有這種這麼大的自我的感動.
其實是復興的一個開始.
因為神的道.
很容易在一些很遙遠的心靈.
來播種.
如果不是.
你想想你連差不多上台.
都還在講笑話.
大家完全沒有任何的準備.
這樣的一種談笑風生.
不是不對.
而是個個都是.

$^{241}$在這樣的狀態.
其實很難一下子.
人會有一個這樣的.
突然之間的轉向.
所以這告訴我們.
第二節.
以色列人就與一切外邦人離絕.
接著就有一個很長的祈禱.
這個祈禱.
就是訴諸神.
在亞伯拉罕.
神的那個簡選.
神的那個呼召.
叫祂離開這個加拿大的梧椅.
去到迦南地.
輾轉就將迦南.
賜給了先祖亞伯拉罕.
然後輾轉輾轉.
因為長話短說.
是一個公土會.
所以就去到摩西.
就是這個民族.
終於脫離那個圍爐之家的埃及.
就去到西奈曠野.
接著神就與摩西頒佈這個律法.
和正正式式將這個民族.
立為自己的選民.
神的選民.
所以在當時出埃及記的20章.
神與以色列人的立約的記號.
就不是葛禮.
葛禮是與亞伯拉罕納的記號.
17章的創世記.
與以色列人所立約的記號.
就是安息日.
所以大家記得.
安息日是很重要的.
猶太人他們知道這一天是神與他們納的.
但是慢慢慢慢就是矯枉過正.
慢慢慢慢就是為了守安息日.

$^{281}$而不是為了安息去敬拜神.
兩者是不同的.
和一個奴隸的民族.
7天24個小時都要工作的.
當然他可以睡覺.
因為不睡覺哪有力去工作呢.
埃及人都不是這樣奴役到不能睡覺.
不能吃的.
因為要幫他做磚.
做那些草的磚.
但是呢.
是沒有假期的.
沒有法定假期這一回事的.
但是神與以色列人立約.
居然就用一個法定假期.
作為那個立約的記號.
你想想神是多麼的有恩典呢.
你試想想.
如果你是打工的.
你是連13天都不能放假.
人都會瘋的.
真的不行的.
神安排人真的需要有一個調整的.
是有一個週期的.
就是說你去到差不多星期五.
去到差不多星期六.
你就很想快點到星期日的了.
你就知道有一個調整.
以色列人就是這樣.
有一個安息日做記號.
輾轉就開始講述.
他們怎樣去違背了摩西的那個律法.
一路來到去.
在神的面前.
特別是提到那個金牛毒的事件.
那個是一個很大的分水嶺.
就是說呢.
講到明不可以有形象.
他們就刻意去做一個形象出來.
表面就說不關我們事.

$^{321}$我們將那些金.
其實亞倫這樣講.
這個是出埃及的三十二章.
我們將那些金啊.
金耳環啊.
那些東西呢.
拿了收集之後.
就放進去住出來.
它變了一個形象就是.
金牛毒的形象.
就是牛毒的形象.
就是牛仔.
牛毒就是牛仔.
其實你沒有想到.
怎會出現這樣呢.
根本就是人的內心.
你們在埃及.
見到這些偶像.
所以有這些形象.
在你們的腦海裡面.
到你們要住出來的時候.
就變成一個金牛毒的形象.
這樣來將神的無限.
來降格.
成為一個.
你可以呼喚.
你可以去猜險.
個個就跪拜那個金牛毒.
就說這個就是帶我們出埃及的神.
牛毒就是神.
這樣來將神.
好像形象化了.
其實是將.
造物主無限的.
這一位的真神.
降格成為一個偶像.
其實是極其的.
一個大的轉捩點.
這個民族的背叛.
所以.

$^{361}$這裡告訴我們.
一而再提到.
民族的背叛.
現在所以.
那個轉回呢.
就是由26至今晚我們看的.
一路第十章呢.
就一個很強烈的.
憤興會.
禁食的一個祈禱會.
帶來的一個.
大家願意.
同神更新那個摩西之約.
因為沒有新約.
沒有耶穌來這個世界.
仍然都是等待耶穌來.
彌賽亞降臨之前.
神就感動尼希米.
發動一個.
表面就是重建硬件的工程.
原來.
做著做著發現.
那個硬件都不是那麼重要.
是重要的.
如果沒有那個四圍的城牆.
外邦人就隨便的踐踏的了.
是需要有一個圍牆的.
但是真正的問題.
原來是.
以色列.
在耶路撒冷.
在以色列定居.
回歸的以色列人.
需要有一個很強烈的意識.
我們是同神立約的子民.
我們現在.
有什麼要做呢.
原來就是要更新那個約.
需要和神重新更新摩西的約.
怎樣更新呢.

$^{401}$十章就有說了.
現在首先就是二十六節.
一到九章尾.
都是不斷的凝聚.
就是說我們已經干犯了神很多的事情.
二十六節開始.
然而他們就是以前的先祖.
他們不順從.
竟背叛你.
將你的律法雕在背後.
殺害那勸他們歸向你的眾先知.
大大的惹動你的怒氣.
所以你將他們交在敵人的手中.
就是當地的迦南人.
和不同的迦南的那些七族.
磨爛他們.
苦難某一個程度.
是神的試煉.
是神藉著一些困難.
給我們有一個反思的空間.
如果太順暢.
你的人生.
如果你每天都是星期日.
你很難有一個警醒.
真的很難的.
每天你都沒有什麼需要擔心.
樣樣東西你都順暢.
這樣的人很難愛主.
當然我不是神.
我不知道每個人的心.
我從我的目測.
從我這麼多年的目會.
如果弟兄姊妹.
每天都是星期日.
沒有任何的憂慮.
沒有任何的困難.
家宅又平安.
什麼都順暢.
這樣很難愛主.
大家可以去觀察一下.

$^{441}$所以苦難是會逼使我們去親近神.
去求問神.
去倚靠神.
去有弟兄姊妹的陪伴.
和代禱.
苦難會令到我們謙卑.
你以為你自己什麼都可以嗎?.
很少的事情都令到你很麻煩.
你就知道原來很多事情.
我們是需要恩典.
恩典那樣東西是你經歷過.
經驗過.
你就發現真的很甘甜.
所以這裡告訴我們.
是磨爛他們.
好像磨爛了那些磨爛.
是又磨你.
又給你艱難.
這樣令到他們.
希望他們有一個迴轉.
他們遭難的時候.
就哀求你了.
這個似乎是事先的循環.
你就從天上垂聽.
照你的大憐憫.
賜給他們拯救者.
就是事先了.
他們脫離敵人的手.
但是他們得到平安之後.
一個循環的loop.
就是一沒事了.
一太平了.
那些敵人離開了.
特別是那個拯救者.
事先又離開世界了.
他們又顧太復萌了.
又是這樣重蹈覆轍.
再走回舊路.
等他們得平安之後.
就在你面前行惡了.

$^{481}$所以你丟棄他們.
在仇敵的手中.
使到仇敵核制他們.
然而他們轉回.
哀求你.
每一次的轉回.
其實都是令到我們吸取到人生的教訓.
你仍從天上垂聽.
屢次照你的憐憫.
去拯救他們.
又警戒他們.
要使到他們歸服你的律法.
使他們卻行事狂傲.
不聽從你的誡命.
干犯你的典章.
你有沒有留意到.
整個禱告集中的焦點.
都是against the law.
違背神的那個律法.
那個的典章.
那個的真理.
不聽從.
干犯.
人如果遵行.
就必定會因此活著.
扭轉肩頭.
硬著頸看.
不肯聽從.
但你多年寬容他們.
又用你的靈.
直眾先知.
勸戒他們.
他們仍不聽從.
所以你將他們交在.
列國之民的手中.
然而你大發憐憫.
不傳言滅絕他們.
即是說留下一些所謂的愚民.
剩餘之民.
很大部分的人.

$^{521}$就會在.
這些一而再的審判裡.
在那個循環裡.
而且是一個.
螺旋形的循環.
每放月下的那個循環.
來到被淘汰.
只是一些真正敬虔.
無論怎樣艱難.
都仍然跟隨神的.
以色列人.
就能夠留下成為愚民.
所以你會看到.
這個.
好像神.
不斷的來懲罰.
但其實每一次懲罰之後.
都是有.
那個大憐憫.
都是當以色列人.
一呼求一哀告的時候.
就立即.
又有那個恩典.
所以這個叫做一個挽回的恩典.
這個就是在.
史詩時代和進入.
迦南的王國的早期.
就是這樣的.
但到後期就算.
他們再哀求.
其實是由嗎哪西王.
之後你讀列王記.
上.
列王記下開始.
去到這個嗎哪西王.
即是喜西家的兒子.
你就知道.
無論他們怎樣悔改.
神的懲罰.
是不會減少的.

$^{561}$就是那個被擄.
就是加比倫的被擄.
所以這個還是一個.
早期的一個循環.
但其實這個循環.
都已經告訴我們.
是一個每放月下.
好像走那些螺旋樓梯.
加州就比較明顯.
四層.
一路這樣螺旋地下.
如果你走得快.
你會暈的.
一路這樣螺旋地下.
現在這裡就是.
其實每次他們去到.
某些關鍵位去呼求.
神又聽.
再呼求又聽.
但其實是每放月下.
的呼求.
每一次的得罪.
又更加的嚴重.
這個就是以色列人的民族.
的歷史.
這個祈禱的人.
似乎可能都是以色拉.
雖然他沒有說到是哪一個.
帶領會眾去祈禱.
但第一節那裡是說.
以色列人站起來.
要稱頌神.
接著他們就開始呼求.
似乎可能都是以色拉.
做那個帶領.
來引導百姓.
回想那個歷史.
回想那個每放月下的.
那個循環.
那個沒事.

$^{601}$平安,敵人離開了.
繼續又離開神.
又拜當地的假神.
接著神的怒氣.
又發作,又興起那些敵人.
來攻擊以色列人.
令到他們面對這些戰爭.
這些攻擊.
他們又可能.
那些農作物又失收.
如是者,去到走投無路.
又一起呼求神.
這樣的每放月下的.
螺旋.
就是在這一段的祈禱裡面.
他們緝錄了.
的熱士師記的那個.
那個時代.
這個時候.
告訴我們知道,很強調.
神仍然有那個.
大發憐憫的心.
不存而滅絕.
不丟棄他們.
因為神是有恩典.
是有憐憫的神.
不是因為百姓.
特別可愛,不是因為他們的.
罪很少,而是.
神有恩典,有憐憫.
很盼望弟兄姊妹我們能夠.
對這個神的屬性.
真的要有一個.
很強烈的經歷和意識.
我們的主.
是真的不會永遠.
懷怒的,只要我們.
幾時,好像我們.
上一個星期提過.
這個歷代治下的七章十四節.

$^{641}$這稱為我名下的.
子民,若是自卑.
禱告,哀求.
尋求我的面.
我必從天上.
垂聽.
所以我們何時謙卑自己.
來哀求神呢.
但這個大段落.
是講一個集體.
有意識地.
被神的道感動.
很想.
親近神.
很想這個民族有沒有什麼.
盼望,現在我們.
有了城牆,有了.
領袖,這個領袖.
現在帶我們去看復興.
這件事,其實回到.
律法書,我們有沒有.
什麼可能,我們有.
那個經驗,神的.
在我們這個民族.
這班人的這個時代.
帶來一個復興.
這個復興,是有什麼.
具體的內容呢.
第九,第十章.
一定要一併去讀的,所以.
現在就是他們.
回顧歷史的.
螺旋的每況愈下.
之後,留意第三十二節.
我們的神,這個.
就是現在開始,就是他們.
現在祈禱會的時候.
那個現在式.
你是至大至能.
至能就是.

$^{681}$有絕對的能力.
至可畏.
守施慈愛的神.
所以這些,如果大家不懂祈禱.
學一下黎希米.
學一下以色列祈禱.
至大至能.
至可畏.
去到極致.
極致,就是最厲害的了.
做物主是至大的.
最有.
最有權能的.
至可畏的.
守施慈愛的.
我們的君王首領.
祭司先知列祖.
和你的眾民.
從亞述列王的時候.
直到今天,所遭遇的.
苦難,現在.
求你不要以為少.
在這一切.
臨到我們的事情上面.
你卻是公義的.
就是說神不會因為.
他的刑罰,他的指紋.
就不義,打重了.
有時我們都會.
生氣得太多.
罵自己的兒子.
通常都不用.
怎麼罵的,如果要罵.
有時說大了.
惡了,很容易.
變成自己不義了.
因為你發了脾氣.
本來你對的.
開始的時候,但是神就.
不需要,從來不需要.

$^{721}$向任何一個人.
說他是不義的.
他所有做的事,都是有.
他的主權,是有他的.
公義,這些是我們.
很難想像的,我們看到.
太多不義的事.
鵝窩戰爭,整個水壩.
到底是誰炸的呢.
驗一下那個炸彈.
驗不驗到呢,那些碎片.
炸了之後.
現在整個村子都淹沒了.
聖靈塗炭.
這些是不義來的.
到底是誰做的呢.
大家又你推,我又推.
好像不了了之.
好像都沒什麼所謂.
下次就可能是鵝窩.
很多這些我們根本.
都看到的不義.
但是,祈禱的人.
真的很厲害,很敬虔.
很屬靈.
他能夠告訴.
當時的群眾.
我們的神,我們現在信的.
跟隨的,敬拜的.
這位神,是至大至能.
至可畏.
又是守施慈愛的神.
就是說,如果你和他.
有那個約在身.
你就有福了,他又是至大.
但他又是守約.
堅守他和.
你所納的約.
那你不用擔心,說明.
我一定會兌現我的約的.

$^{761}$這樣的神.
我們最怕那些人不講信用.
講好了你都不背.
但現在神說,他既是.
至大至能,更加是.
守施慈愛.
是以慈愛去守.
他的約.
那你可以想像得到.
我們基督徒為什麼.
這麼蒙福,我們基督徒每一次.
在聖餐.
為什麼這麼嚴肅.
就是每一次我們看到這個餅.
其實就要你去.
不是叫你幻想,是叫你回想.
昔日主耶穌.
在這個.
樓房最後晚餐.
他自己親自.
設立的一個新約.
就是和我們所有.
相信他的人.
所納的這個永約.
就是說,你的命仔.
一定得到保守.
因為主耶穌.
神的愛子.
是為你這個.
有限的人釘十字架.
所以.
所有的神學家很誇張.
全世界都沒有罪人.
只有你這個這麼小的人.
只有你這個人是罪人.
耶穌都會為你死.
當然這個是不可能的.
是誇張了,但是某程度講到.
那個約的精神.
這個神就是守施慈愛.

$^{801}$所以.
現在在祈禱裡邊,為什麼有盼望呢.
這個民族為什麼有盼望呢.
首先就是.
很久很久的時候.
神感動一個酒精尼希米.
上去.
看到那個國家.
這麼破落.
這麼寒酸.
希望他們能夠有一些尊嚴.
回去重建那些城牆.
做著做著.
神的計劃.
再在這些敬虔的領袖的心靈裡邊.
再做感動.
不要只是做城牆.
做下人心.
那個的內心.
更加重要.
可以碰到人的心靈.
人的心自己願意打開.
向神.
這樣就非同小可了.
這個民族就變成.
一個真是有力量.
可以抗衡周邊的.
外邦人的.
一個民族.
因為大家的心.
都被神的道.
和神的靈.
和領袖的呼喚.
來感動.
所以這個領袖.
這個祈禱的人.
他現在不是咬文嚼字.
是將做物主的屬性.
化在祈禱裡邊.
一個祈禱的人.

$^{841}$帶領公土.
一個盼望.
需要給弟兄姊妹有方向.
告訴我們我們信的是.
怎樣的神.
很重要.
要給我們這個民族有出路.
給我們有一個方向.
現在是應該要復興的時候.
做好了城牆.
大家要開始了.
你可能仍然要耕種.
所以第十一章.
要抽籤.
每個支派就要抽些人上來.
老撒冷住.
要千骸.
所以我們下一次.
再講十一章.
我們看這裡.
這裡說的是一個有氣質.
有屬靈氣質.
有屬靈的動見.
有屬靈的承擔的領袖.
尼希米.
以斯拉和在他身邊的濟斯.
這些重要的領袖.
已經被神的靈.
感動.
所以他們現在.
帶領回眾.
呼喚回眾.
大家一起仰望這位神.
他是怎樣的神.
所以兄姊妹我們要問一問.
現在你信的.
在錦繡堂.
你信的是甚麼神.
可能大家都不知道.
甚麼神.

$^{881}$就信耶穌.
可能只知道這個名字.
這裡就已經有教育.
已經告訴我們.
眾物主是這樣的性情.
或者是神性的屬性.
這個叫神性的屬性.
神性的屬性.
我們每個人.
都有不同的性向.
譬如我是比較.
我不知道是屬於慢還是快.
我是屬於感性.
還是屬於理性.
可能都不是.
我又覺得自己很慢.
但又有人覺得我很快很急.
我又真的好像有時很急.
但其實我很慢.
我是不是一個很感性的人.
又不覺得是.
但其實我很感性.
我看現在的文藝片都會哭.
我又是很理性.
有時是將整個感性的系統.
刻意關掉.
因為如果.
每次都感性真的不行.
真的會.
做不到事.
你想想這麼複雜.
但是眾物主.
他有他自己的attribute.
所以我們讀這些舊約聖經.
祈禱的領袖.
很多時候.
都帶回眾去看.
眾物主的屬性.
因為不是他.
我們這個民族根本是沒有用的.

$^{921}$沒有意思的.
我們沒有能力去抗衡周邊的人.
他們的影響.
只是你想想.
一個安息日買賣.
他們都已經抵受不住這些誘惑.
安息日不買住.
星期日.
日落之後.
你就去買賣.
都不行.
都要被那些外邦人任意擺佈.
擺佈完之後.
嘲笑你們猶太人.
都沒有自己的.
立場.
所以很重要祈禱.
要帶領回眾.
看到眾物主榮耀的屬性.
帶給這個民族.
有盼望有出路.
所以公土會.
很重要是要帶出路.
帶出盼望.
給弟兄姊妹一個方向.
領袖應該就是vision.
應該是要這樣的.
留意下面.
第35節.
丹門在本國女.
尖利大恩的時候.
在你所賜給他們.
這廣大肥美的地上.
迦南美地.
不事奉你也不.
轉離他們的惡行.
我們現今作了奴隸.
奴僕.
至於你所賜給我們列祖.
享受其上的土產.

$^{961}$並美物之地.
看啊我們在這地上作了奴僕.
這地許多的出產.
歸了列王.
就是你因我們的罪.
所派.
合制我們.
他們任意合制.
我們的身體和生畜.
我們租了大能.
一路去回顧.
在很多的恩典裡面.
我們白白的租撻.
很多的美物.
地圖的.
或者是很豐富的.
但是我們.
享受這些恩典.
敬畏你不聽你的道.
不來去遵行你的旨意.
這個就是我們.
很容易.
我們很緊張團契生活.
很重要.
彼此相交.
但是不要成為了.
這個是我們返教會的主菜.
我們返教會.
無論如何是敬拜神.
和神的關係.
的一個復和.
和對眾物主的認識.
所以團契的相交.
是要促進我們一班人.
一起在靈命上.
成長.
如果沒有了神的部分.
只有我們每個禮拜喝茶.
很熟悉的了.
無所不談的了.

$^{1001}$不過每次都是無所不談.
只是除了神例外.
什麼都談.
那就不是真正的團契了.
團契應該是要將我們.
帶進去和神的.
相交裡面.
所以這裡告訴我們.
他們現在就是面對著.
享受著神的.
祝福的地圖.
那個豐富.
的國泰民安.
但是誰給國泰民安呢.
無所謂了.
到真的有困難的時候.
求這位神.
又不聽.
不如我們求一下那些.
迦南的亞瑟拉.
還有.
這個.
巴力.
巴力就是.
掌管.
風調雨順的神.
所以他第一時間就.
拜巴力.
你就可以想像到.
這個情況就是這樣.
要兒子不要母親.
不要神.
要神的祝福.
有時很可惜.
有神自然就會有神的祝福.
記不記得耶穌說過.
你們要先求.
他的國和.
他的兒.
這些東西就是.

$^{1041}$上文說的日用的飲食.
不需要掛類的.
外邦人才會這樣掛類.
這些東西都要.
加級你們.
你們要先求他的國.
和他的兒.
這是六章三十三節的馬太福音.
所以和這裡完全一樣.
我們要緊張神的事情.
我們緊張永恆的.
我們整個心靈.
緊張永恆的事.
自然.
身體物質的.
不是不要緊.
但真的沒有那麼要緊.
還有神給我們什麼.
我們帶著感恩的心靈.
這樣去享用.
去享受.
這樣就已經是我們基督徒.
真是.
有衣有食就當知足的.
那個註解.
本來就是這樣勉勵.
這個.
耳筆鎖的教會.
就是要這樣去做.
這裡告訴我們.
第三十八節.
因這一切的事.
我們現在.
立確實的約.
寫在冊上.
我們的首領利未人.
祭司.
都簽名.
主動.
以色列祈禱員.

$^{1081}$既然我們的歷史.
是.
一而再的.
享受恩典.
忘記次恩典的神.
現在我們要有一個改變.
我們不是忽視恩典.
但是我們更加重要.
就是和這位次恩典的神.
和我們主動.
立約的神.
來立約.
所以現在他們弄了一本.
冊子.
大家要簽名作實.
都挺.
demanding.
挺扯.
要你簽名作實.
承諾.
我要和摩西的律法.
來重新.
有那個.
關係.
去遵守.
最簡單.
星期五的黃昏.
我就要開始和我的家人.
和我的親戚朋友.
開始要守安息日.
開始要來.
什麼工都不做.
要來準備.
去敬拜神.
安息日就去到聖殿.
來聽律法書.
當時應該都有會堂.
不好意思.
再後期才有.
就去到會堂.

$^{1121}$因為聖殿都被毀了.
就去到會堂.
讀律法書.
這是每個安息日必須要做的工作.
有十五個.
南丁猶太人.
分散在當時的地方.
就會有會堂.
設立.
這個是後話.
但現在他們在.
復興的浪潮.
七天的讀經運動.
七天的培靈會.
和研習班.
查經班.
在這個七月的二十四日.
第九章第一節開始.
在這個長的祈禱裡面.
然後.
有一個行動.
決定大家.
誰願意的就.
簽名,作實.
我們有什麼做呢.
當然就是和摩西的律法.
有一個回調.
第十章開始.
幾個系統地.
先是從那個.
屬靈的領袖.
然後就是祭司利未人.
和在聖殿裡面.
公職的例如尼提靈.
即是那些電役.
看守電門.
那些職位.
其實很低級的.
是當時.
那些機騙人做的.

$^{1161}$第十章簽名的.
就是哈加利亞的兒子.
這個我們知道.
第一章就說了黎希米就是省長黎希米.
就是那個酒政.
哈加利亞的兒子.
省長黎希米.
第一個簽名.
首先.
既然是由你和以色列.
引發的一個.
獨經運動.
培靈會.
現在這個祈禱之後.
就有一個很大的感動.
大家誰願意去.
祈禱.
如果大家對號入座.
一起去.
這樣來看的時候.
用我們自己的場景.
如果我們的同工.
我們的領袖.
堂董.
執事,一些部的部長.
一些的師班長.
團契的團長.
或者其他的領袖.
如果我們能夠.
這樣來簽名.
作實.
大家在神的靈.
的感動.
不是人為的壓力.
底下,所以要叫人事奉.
不用求人.
求是一次過得.
持續的感動.
是那個人.
的心被神的愛.

$^{1201}$神的靈.
和神的真道感動.
持續自己去做.
沒有人做.
不要緊,我做.
我去承擔.
我默默的去做.
做了也不需要告訴別人.
也不需要告訴別人我很辛苦.
而是要這樣來引發.
一個復興的.
一個漣漪的效應.
所以這裡告訴我們.
真誠簽約受感動的.
你看到這個.
家譜.
簽名的人.
很多是領袖.
是由祭司帶動的復興.
弟兄姊妹.
我們不妨.
為我們的領袖祈禱.
現在做什麼呢.
看東張西望.
還是很忙碌.
在工作.
還是和朋友.
在喝酒.
在南貴坊.
諸如此類.
沒有批判的.
而是出於神的靈的感動.
大家願意為神的家.
為這個時代的復興.
我們去尋求.
去守望.
每個人都說.
裡面的.
收緊了.
那我就更加要祈禱.

$^{1241}$更加我們要站在.
一個時代.
去為真道見證.
更加我們第一時間.
為我們的後來的弟兄姊妹.
還是很幼嫩的.
去守望.
去培育他們成長.
他們將來就可以來到.
獨當一面.
當我們過去的時候.
甚至還沒過去而要坐牢的時候.
他們就可以承擔.
全福音的工作.
為什麼一想就立刻走.
我又是沒有批判的.
因為我自己不走.
我這裡不走.
沒有批判的.
每個人有神的帶領.
我不是神.
我沒有權說這個是這樣那個.
我只能夠干預.
我自己的心靈.
但這裡你會看到.
是一件很美麗的事.
很多的領袖.
所以如果你是當時的以色列人.
在這個城牆重建好之後.
在這個.
嚴肅會的祈禱裡面.
你看到這些簽名的人.
你的心是很受感動.
你看到他們.
真的很認真.
為自己的民族尊嚴.
擺上.
自己的生命.
基督教往往.
是願意犧牲.

$^{1281}$為耶穌犧牲的人.
才能成就一些事功.
基督教不是跟你說.
你應該這樣做.
沒有應該.
是感動.
帶來復興.
是感動造成一個事功.
感動就會成就一個事功.
所以你會看到.
很多人都知道應該要祈禱.
但有多少人.
真的認真.
一齊回來祈禱呢.
甚至乎內心裡面.
很緊張.
希望多些人祈禱.
的那種力量呢.
這個就是告訴我們.
今天我們將.
神最重要.
看為最重要的事情.
那個次序放得.
有這麼低就這麼低.
我們自己的事有時候放到無限這麼大.
有這麼高就這麼高.
偶然間沒問題.
但是一而再.
周而復始都是這樣排序.
那個次序.
一定是有些事反映到的.
簽名的就是.
哈加利亞的兒子省長.
哈利希米和西底加.
第二批就是.
祭司了.
西來亞阿薩尼亞.
耶利米巴斯胡爾.
阿瑪利亞.
馬基亞哈特.

$^{1321}$斯巴黎馬六.
哈林米利末.
俄巴底亞.
但以理近頓巴六.
米舒蘭阿比亞.
米瓦文.
馬西亞碧蓋.
瑞士瑪雅.
祭司都這麼多.
如果我們能夠將那些名字.
對號入座.
曾牧師.
麗時牧師謝牧師陳牧師.
然後一直這樣數下去.
魯牧師一直數下去.
鄧全道.
吳全道楊姑娘.
還有其他的.
曾姑娘.
陳靜娟姑娘.
林牧師.
然後到唐董.
使徒永富弟兄.
劉振強弟兄.
一直這樣讀下去.
你想想多麼美麗的一件事.
是多麼美麗的一件事.
第九節.
又有利未人.
就是.
阿薩尼的兒子耶穌亞.
希臘的子孫.
賓來革滅.
還有他們的弟兄.
是巴尼,荷第亞,基利他,比萊瓦,哈蘭,米加,尼合,哈沙比亞,薩克,史利比,是巴尼,荷第亞,巴尼,比利路.
又有文的首領.
就是巴祿.
這些就沒有官職了.
不是在聖殿裡事奉的了.
祭司利未人就肯定是在聖殿裡供職事奉的.

$^{1361}$亦都是曾經在過程裡.
陪伴尼希米.
來建造城牆.
鼓勵百姓.
和之前的七人.
在那個陪靈會裡面.
來教導百姓.
律法書的這些.
祭司利未人.
現在去到民間的領袖.
一些可能在教會很久的.
告訴我們知道.
就是巴祿,巴哈,摩亞.
以蘭,薩土.
巴尼,布尼.
雅甲,比拜,亞多尼亞.
比格瓦伊.
阿丁,阿特,希西加.
亞素,荷第亞.
哈遜,比賽.
哈拉,阿拿特.
尼拜,莫比亞.
米舒蘭,希西.
米氏,薩別.
薩督,亞道亞.
比拉提,哈蘭.
阿萊瓦,荷西亞.
哈拿尼亞.
哈述,哈羅克.
比利哈,素伯.
尼宏,哈沙拿.
馬西亞,阿興亞.
哈蘭,亞蘭.
馬祿,哈林.
巴拿.
接著還有其他普通的.
神的選民.
其餘的文祭司利未人.
守門的歌唱的尼提靈.
和一切尼訣.

$^{1401}$律法書的一些.
和一切尼訣.
倫邦居民.
歸復神律法的.
記住這個形容詞.
他們這個簽名的內容.
是要和.
外邦人,但是是我們的.
鄰舍來的,是有分別.
是要尼訣.
這個尼訣就不是說.
你經常板著臉.
見到人你都要生氣.
像村人一樣,是要和他們有分別.
是要讓人知道.
他們是屬神的子民.
他們是不屬神的.
這樣的一個分別.
所以他們是要歸復神的律法.
並他們的妻子.
兒女,凡是有.
知識能夠明白的.
就是說對神的律法書.
他們是聽明白了.
是願意.
就是說有代價要付.
聽不明白什麼都不知道.
你這樣說,我也不知道你說什麼.
那就沒有什麼代價.
你聽明白,有要求.
而你又願意簽名.
就是說你很想在神.
面前有一個更新.
你就願意這樣去做一個.
行動.
這個行動就是.
有知識能夠明白的.
都隨從他們的貴就弟兄.
這裡很重.
法就起誓.

$^{1441}$必定遵行神.
直他僕人摩西所傳的律法.
謹守遵行.
耶和華我們主.
的一切誡命.
典章律例.
重重複複都是想說.
要回到摩西的.
那個律法書.
那個精神.
那個原先的.
那個要求.
就是說加了很多人為的.
因素,令到摩西的.
律法書沒有了那種殺傷力.
沒有了那種的影響力.
現在他們清楚了.
我要回到神.
透過摩西.
和我們這個民族所立的.
約那個.
最.
原先的.
那個要求的.
水平.
來有一個.
簽名,法就起誓.
因為沒有新約.
的出現.
他們最高的就是回到.
神的摩西律法書.
十誡和十誡.
衍生的不同的道德戒律.
當然有些禮儀的.
例如安息日他們要去守.
有些節期,有些憲制.
但是更重要就是.
同神,同人.
的關係,所以主耶穌.
最後在新約.

$^{1481}$逼耶穌,誡命.
最大是什麼.
耶穌說你們有沒有讀過.
他們立刻說了.
生命的六章四節.
你們要聽.
耶和華,我們的神.
是獨一的主,所以要.
盡心盡性盡意盡力.
愛耶和華.
你的神,其次呢.
你還沒記的十九章.
要愛人如己,不緊張是十九.
還是二十九,其次就是.
要愛人如己.
耶穌說,這個就是律法的總綱.
所以.
保羅也是這樣說的.
愛是.
包括.
聯絡,傳德.
所有的.
德行,都是在愛的.
總綱裡面.
有愛你就不會.
頂枝門你有沒有試過.
被人轉過面.
不看你,其實.
如果你是成年人.
很難受的.
真的很難受的,在一個文明.
社會,不是一些食人族.
人們真的不看你.
走開,故意不看你.
很難受的,你知道嗎.
那是一種侮辱.
但是我們要記得.
在愛裡面,我們不會有.
這些眼神,愛裡面.
我們不應該這樣藐人.

$^{1521}$愛裡面我們不應該.
這樣去.
有這些眼神的.
敵視,因為愛.
是聯絡傳德.
就是說,你慢慢慢慢.
其實你不會再這樣.
鬥雞眼地射人的.
你不會再很兇地.
去藐人的.
你自己的面容都會改變的.
這個就是愛.
保羅說,愛是.
聯絡傳德的.
羅馬書應該是十五章.
你就可以想像得到.
耶穌的教訓.
他說將整個律法書.
其實摩西律法書的精神.
就是愛造物主神.
盡心.
盡,這個字真的很重要.
我們整個人生.
其實不斷是盡.
每個人都.
懂得什麼叫盡.
你為你的小朋友盡.
沒有人會.
兒子,我只是10\%.
90\%是我自己的.
我們不會這樣.
我們一有小朋友,我們就知道不是這樣.
但是.
聖經就告訴我們.
沒錯,我們要盡力去照管.
因為是神給我們的產業.
但是我們是管家.
祂是一個個體.
祂是屬於神的.
你要教導祂,祂自己要敬畏神.

$^{1561}$祂自己要100\%.
盡祂的一切.
愛神,但是我們盡了.
我們做父母的本份.
其實最重要就是.
我們自己盡心愛神.
和愛其他人.
所以,這個就是.
新約律法書的精神.
就是要回到律法裡.
所以,未有新約.
他們就要回到.
舊約的摩西律法書.
就不是再另外立一個.
由李希米,由以色列.
自己再立一些的戒條.
其實不需要.
回到神.
透過摩西和這個民族.
所立的約就已經足夠了.
所以,你看看.
舊約原來這麼偉大.
一直承載著.
主耶穌的出現.
才是新的約.
因為舊約是不能再遵守了.
他們不能恆心地遵守.
和大家看看.
希伯來書第八章.
就說了舊約和新約.
希伯來書的八章.
其實是引耶利米書的32章.
希伯來書.
比較後一點.
說到約.
就一定要看看這段經文.
作者.
都是很熟悉舊約的人.
第八章的第十節.
主尤說.

$^{1601}$那些日子之後.
我與以色列家所立的約.
乃是這樣.
我要將我的律法.
放在他們裡面.
當然不是那個.
USB手指.
插進你那裡製造一個晶片.
是在我們的心靈裡面.
有聖靈.
有一個新的生命.
新的心靈.
我們在我們的內心.
我們就能夠遵行神的道.
在愛.
在神的真愛.
在耶穌的新的生命裡面.
重生的生命裡面.
我們就能夠遵行神的道.
這裡說.
我要將我的律法寫在他們裡面.
寫在他們心上.
我要作他們的神.
他們要作我的子民.
他們不用各人教導.
自己的鄉鄰和自己的弟兄.
就說你該認識主.
其實.
當然這裡不是叫我們不要上主一學.
不用上主一學了.
不是這個意思.
而是告訴我們.
有聖靈在我們內心.
我們就會認識神.
所以信了主耶穌.
你要祈禱.
向天父祈禱.
初時是不熟悉的.
是很陌生的.
但慢慢地你知道.

$^{1641}$這位真是你的天父.
所以你跟他說.
我的天父是聽的.
所以越是不懂得祈禱.
越是要感受一下.
你現在有個爸爸.
是做物主.
這是一個很寶貴.
很美麗的生命成長階段.
由一個嬰兒.
他開始第一個呼吸.
一路長大.
他越來越懂事.
越來越知道這是他的家.
他的父母,兄弟姊妹.
開始有朋友.
開始對這個世界有認識.
知道自己的存在.
慢慢地成長.
基督徒一樣.
不斷地呼吸.
祈禱.
向神祈禱.
不是長短的問題.
是你的內心和神的關係的問題.
跟著告訴我們.
第十一節的下半節.
因為他們從最小到最大都必認識我.
我要寬恕他們的不義.
就是新約的精神.
就是十字架的寶血.
不再紀念他們的罪騙.
完全忘記.
其實忘記是錯的.
不是忘記,是塗抹.
忘記即是說我還記得.
其實不是.
是真的寶血已經完全塗抹.
真的好像沒有發生過.
第十三節.

$^{1681}$既說新約.
就以前約為舊了.
但那漸舊漸衰的.
就是舊約.
就必快歸無有了.
既說新約.
耶穌和祂的子民.
信祂的人所立的.
是新約.
所以我們今天不是守摩西的律法.
信耶穌.
我們是信耶穌的袋淑.
寶血而得救.
所以我們是和神立了一個新約.
但是在這個新約裡.
我們就是神的選民.
神的子民.
屬天的子民.
但是我們仍然要愛神.
同時要保護神.
我們仍然要愛神.
和愛人.
和我們要來遵行神的真理.
這個就是那個約.
我們回到剛才說的第十章這一段.
尼兄米記的第十章.
這裡告訴我們.
去到三十節.
並不將我們的女兒.
嫁給這地的居民.
也不為我們的女兒.
嫁給這地的居民.
也不為我們的女兒.
嫁給這地的居民.
也不為我們的兒子.
娶他們的女兒.
也不為我們的兒子.
娶他們的女兒.
也不為我們的兒子.
娶他們的女兒.

$^{1721}$也不為我們的兒子.
娶他們的女兒.
也不為我們的兒子.
娶他們的女兒.
也不為我們的兒子.
娶他們的女兒.
也不為我們的兒子.
娶他們的女兒.
這一刻他們要有一個具體的.
這個民族的一個.
就是要告訴外邦人.
要不你們歸向我們的神.
如果不是我們因為神的緣故.
我們寧願有一個.
因為你少了很多選擇.
只有自己猶太人.
那個限制.
開始不太覺.
但慢慢就會浮現.
所以這一刻的這麼決絕.
是有它的特殊的原因和意義.
就是很清楚.
現在我們在這樣的一個情況.
我們要有一個界線.
其實就是要立一個.
神選民的界線.
外面是圍牆.
但真正的界線是在說那個敬虔.
不是在說那個歧視.
所以猶太人去到耶穌的年代.
就真是去到一個歧視.
就是以自己.
極端的優越.
排擠撒瑪利亞人.
排擠那些瑞利昌記.
罪人.
排擠那些外邦人.
但自己就不是真正的.
有這個界線.
所以真正的界線.

$^{1761}$就不是對其他人的歧視.
真正的界線是自己有一個.
分別為性.
一個問題.
所以這個永遠都是一個.
誠性的.
你可以說是一個張力.
和一個交戰.
還有其他.
有些輕微的.
第31節.
這地的居民若在安聖日.
或者什麼聖日.
帶了貨物.
這地的居民就是外邦人.
迦南人.
即是參巴拉多比亞那些.
我們必定不買.
即是說.
真是要有一個.
分別就是.
安聖日是安息.
我們要敬拜神.
我們要有聖會.
我們要專心敬拜神.
我們真是實行的.
心靈的安息.
是從身體力行的.
一個改變.
所以真是奉勸.
我不知道有多少人做生意.
有多少工作的弟兄姊妹.
如果你自己可以去選擇的.
真是盡量.
選擇.
回教會.
有你自己半日至一日.
的一個休息.
安息.
有意識地帶領.

$^{1801}$如果你是家長.
帶領你的兒女.
現在要敬拜神了.
什麼都要放下.
特別是他小時候.
你不任他的時候.
他自己要任自己.
可以的時候.
所以我們都曾經說過.
我們那時候在那些小教會.
我太太捉住阿明仔.
很小的時候.
手抱他.
那時候很得意.
坐在第一行.
聽到不准動.
要敬拜神.
是大人崇拜.
不斷動.
媽媽夾住他.
有個小玩具放在旁邊.
給他玩.
慢慢再大一點.
都不行.
後來有寺廟房就讓他進寺廟房.
如是者是這樣.
一方面其中一個原因.
寺廟房太嘈.
每個禮拜回去就交換情報.
打牙教.
好像鳳子望著前面.
因為通常有鏡.
現在那些玻璃.
通常都是這樣.
不會密封.
望著講完講到.
但就在那裡打牙教.
在教會都是這樣.
所以不是批判.
很習慣.

$^{1841}$所以有些人真的受不了.
出來坐.
感受都有.
原來不是他不想進去.
而是太嘈.
出來聚會.
其實我這麼忙.
我真的想回禮拜堂.
聽道.
但一見到姊妹.
很多事想講.
一個小時就講完.
就是一個.
如是者的循環.
所以.
這裡告訴我們.
他們要有意識地守安息日.
如果有人.
不知道那些東西是不是特別好賣.
還是特別重要.
要到安息日才特意引誘這些意識的人來買.
帶了貨物和糧食.
來賣給我們.
我們必不買.
要有這種決心.
我們真的.
如果你是家長.
你有權的時候.
你就有權要用盡.
我們要盡心盡性敬拜神.
你六天都上學這麼辛苦.
第七天就要一起成家回教會.
我們敬拜神.
給了這個印象.
這半天.
什麼都不可以做.
敬拜神.
你患重病有特別原因.
當然是回不了教會.
回不了教會和不回教會是兩回事.

$^{1881}$所以.
如果從小就有這個意識.
很開心.
我們回教會一起敬拜神.
將來他自己一個人的時候.
他自己長大.
他自己都要面對自己的信仰的時候.
你當年的.
缺絕.
就成為他很大的益處.
如果你當年.
你自己都隨隨便便.
隨隨便便.
不夠睡覺.
就大件事了.
他的崇拜生活.
就因為被你影響.
他就可有可無.
所以求神幫助我們.
小朋友小.
不知道是不是十二歲以下.
做家長應該缺絕一點.
一定要回教會.
一起回教會.
和你一起去飲茶.
買東西給你吃.
當然不是賄賂.
而是令到這件事很美麗.
很開心.
所以必不買.
所以我們今天來到香港.
我們沒有什麼逼迫的時候.
我們自己去到敬拜神.
我們都要大家.
有一個很強的精神.
梳頭洗臉.
穿最漂亮的衣服.
敬拜神.
雖然我都已經沒有那麼嚴格.
但我都.

$^{1921}$有時都受不了.
現在夏天.
又短褲又拖鞋.
敬拜神.
你去其他場合.
你穿得很隆重.
禮拜日你又特意穿成這樣.
都不好.
其實是一個形象.
基督教是不是.
被人覺得你是街市.
外邦人.
未信的人怎樣看我們回教會.
你帶著愁眉苦臉.
回教會.
這家人真是很特別.
禮拜個個都很莊重.
很開心.
他們去哪裡.
回教會.
回教會都這麼開心這麼莊重.
是一個見證.
我們不是律法主義.
而是當我們被神感動.
我們應該有.
這樣的態度.
有時我都擔心.
你想想我星期六.
已經忙又睡得不好.
我已經三四點起床.
五點多準備就已經差不多.
吃完早餐要回教會.
我時見到弟兄姊妹十一點.
還污眉黑睡的樣子.
我已經做了很多事.
我五點多做到十點.
早堂就教完聖經.
準備五堂.
大家回來都還未睡醒.
有沒有搞錯.

$^{1961}$禮拜十一點.
你都還未睡醒.
回教會.
我又不是批判.
而是將一些現象.
給大家一起在這段經文裡面反思.
現在他們很認真地.
我們要回到摩西的律法書.
其實有六天都可以買賣.
耶穌就說.
六天都可以醫治.
為什麼.
安息日.
是那些法利座的人.
來罵耶穌.
六天都可以醫治.
為什麼安息日才醫治呢.
其實耶穌不是想打破安息日.
是那些人已經將那件事扭曲了.
變成一種很沉重的枷鎖.
但回到這段經文.
是復興的時候.
是復興的.
這個嚴肅會.
在這個祈禱會裡面.
之後大家簽名.
來納藥.
記住是與神納藥.
不是和尼希米和以色列納藥.
所以尼希米是與神納藥.
以色列是與神納藥.
所以現在他們納藥.
他們都是有這一份.
這樣的尊嚴.
其實是尊嚴來的.
是看你怎樣想的.
我們是不買你們外邦人的禮物.
不是我們不和你做生意.
而是明天先.
現在是安息日.

$^{2001}$我們要敬拜神.
所以每逢第七年.
必定不耕種.
要順靠神的供應.
讓地土安息.
讓地土自己生的植物.
那些果實.
是給那些人呢.
是給那些寄居窮苦和外邦人.
這樣的安息年.
是給那些寄居窮苦和外邦人.
這樣的安息年.
接著就是.
凡是欠我們債的.
我們不追討.
因為你是有餘.
安息年就是在神的恩典裡.
你將神給你的恩典.
少少.
給那些有極端困難的人.
我們不會追債.
不要說收很貴的利息.
而是去到安息年.
我們就免了這些人的債.
神就會加倍的.
來祝福放債的人.
所以你看看這個.
你看到那個.
對神的一個.
感動就會引發一個人的.
善心對人都有感動.
對人都有一個.
對人都有一個憐憫.
所以是相輔相成的.
盧家福音曾經有一個.
這樣的比喻.
有一個欠債的僕人.
欠債主一千萬兩.
寬恕了他.
他很開心.

$^{2041}$出去有另外一個.
的僕人欠他十兩.
他捏著他的頸.
你今天就要還給我.
剛剛才興高采烈.
一千萬兩不用還.
有些僕人就告訴那個主人.
剛才那個僕人.
這樣對他的債主.
債主.
那個主人就很生氣.
命令要將這個僕人.
來到去收監.
要他還清那些債項.
有了神的.
後恩而我們對人.
來到斤斤計較.
和沒有了那個憐憫.
這些是很大的罪.
所以我們應該是越來越.
有一份慈愛.
善良的心.
所以你看到.
保羅在講完.
真理之後.
就教那些.
譬如哥羅西教會.
倘若者人與那人有嫌隙.
總要彼此包容彼此饒恕.
主怎樣饒恕你.
你都要這樣.
饒恕他這樣的教導.
是因為主.
而我們就有這個能力.
有這樣的胸襟.
去做那些事情.
所以一切都是因為有救恩.
一切都是因為有一位主.
看著我們.
所以如果我見到.

$^{2081}$回來茶經班的弟兄姊妹.
如果你用一些很愚蠢的眼神.
看人.
你記住你聽過那些堂道.
我自己我真的自問.
我真的沒有了.
我對弟兄姊妹.
我是沒有任何的恆意.
很感恩.
不是我這個人是聖人.
是因為救恩.
越來越經歷到神的愛.
根本不需要.
人會誤解你.
誤會你.
甚至會有一些眼神.
我都講過.
剛才都講完.
看著你特意轉開臉.
你是很難受的.
但不要緊.
一點都不介意.
你總是會明白的.
也都沒有辦法解釋.
但神是有恩典的.
我們就要求主讓我們能夠.
有這種必不買的決心.
因為記住是同神納若.
我們在神的面前.
我們同神納若.
我們有這種具體的.
身體力行的行動.
所以其實每天.
起身你能不能夠.
有些記事簿.
或者現在用電子版的.
Apps.
今天我有一件事可以做.
我決心.
同神納若我做.

$^{2121}$譬如我每天一定要和一個人.
講一些恩典的說話.
等等.
每天我都能夠自己感恩.
什麼事情我都感恩.
你的生命就很不同了.
你會有一個.
很強烈的愛神愛人的心.
星期日.
我一定和一個不認識的人.
打聲招呼.
回來教會.
你就可以想像到.
我是誰.
還是用阮眉.
我是阮眉你是誰.
我在這個教會很多年了.
很多恩典的.
打開了那個話盒子.
你又認識一個新朋友.
代表教會又歡迎了一個新朋友.
他又可以慢慢融入.
這個教會的生活.
如果你能夠每個星期都決心.
神納這個若.
一年過去.
你寫下這些名字.
我回家我為你祈禱.
你就成為一個.
代禱者.
神給你的內心的感動.
就從一個感恩開始.
每次.
一個星期記一個.
我這個星期就為A祈禱.
你想想你的人生.
多麼美麗.
其實你已經是參與一個.
教會就需要.
這些表面.

$^{2161}$就好像.
在現在的社會每個人都很想說.
個人空間.
你這麼熱情做什麼.
你為什麼要這樣來問我.
其實不是.
教會生活就是大家彼此.
扶持彼此認識.
的而且確絕對不應該.
進入別人的私人空間.
來侵犯別人的私隱.
絕對不應該的.
我們是大家一起在神的國度裡面.
去追求.
所以我們需要有這些這麼認真.
願意聽了又明白.
願意有具體.
去實行.
與神立約.
回去摩西的約裡面.
剛才他們說的幾樣東西.
有戒善我們是屬神的.
我們離絕.
外邦的女子.
不准娶這些太太.
我們不會嫁.
給這些外邦的男子.
要有戒善.
你們歸入猶太教.
成為敬畏神的人.
你就是我們的子民.
你就可以和我們通婚.
但有戒善.
接著就是說.
其實不守安息日.
在他們來說好像很大件事.
因為又是已經沒有了.
那個戒善.
曾幾何時.
就是無所謂了.

$^{2201}$安息日有和沒有.
都沒有分別.
人們都照樣進來.
都沒有人守住門口.
又不讓他們進來.
那些東西又真是很吸引.
慢慢慢慢大家就.
不是太緊張敬拜的生活.
就出現了這種情況.
就好像此消彼長.
如果我們的敬拜生活.
是開心活潑.
大家就不會太過介意.
彼此的相交.
所以其實是需要.
有心人.
有意識地.
自己以身作則行先.
所以沒有人做.
不要緊,我做.
沒問題的.
我們自己去行這一步.
行出這一步.
總會有人看到.
總會有神的恩典.
來伴隨著.
所以很盼望.
這個就是.
他們具體的行動.
對其他人.
有一份憐憫.
地圖.
可以得到安息.
對工人我們又會多了一份.
體恤.
不只是跟勞工法例.
我們都能夠盡量有一種體恤.
跟著.
我們如果有那個恩典.
能夠輸送的.

$^{2241}$我們能夠輸送給其他人.
這個都是一個很緊要的.
一個心態.
第32節.
我們又為自己定例.
每年各人捐.
銀一射下.
三分之一.
為我們神殿的使用.
所以你會看到.
奉獻.
不是逼,不是求.
是感動.
是神的靈.
的感動.
被神的愛感動.
被神的福音去感動.
所以沒有.
不要緊,神給我們那麼多,我們就先做.
做到有神的感動.
弟兄姊妹出現,就繼續做.
這個是真正的.
教會的理財.
至少我自己這樣認為.
我們不應該去勸.
不應該去借.
我們應該感動弟兄姊妹.
這個是為神殿的.
功用,所以你會看到.
他們每人現在就立了.
一個定例.
為陳設餅,常獻.
的素祭,梵祭.
安息日,月素.
節期,所獻的聖物.
並以色列人的贖罪祭.
以及我們神殿.
裡面的一切費用.
祭司利未人他們的.
生活開支.

$^{2281}$都是透過這個奉獻.
如果沒有人再獻贖罪祭.
他們就不會有.
祭生.
他們不會有那些物件.
來做陳設餅.
他們就不會有這些聖餅.
可以吃,更加不會有.
其他的下欄.
就是說.
如果他們都忽視獻祭.
那聖殿就很荒廢.
好像現在歐洲的教堂.
被人一定要.
進去就要收門訣.
你信不信耶穌都要收.
進去不是敬拜神.
是參觀,是用來維修.
收緝.
接著都不夠了.
收緝都沒有人看了.
就賣給那些阿拉伯人的回教.
或者其他的地方.
做咖啡店.
教堂現在就變成這樣.
變成一個景點.
變成可以去.
或者不去都沒有問題.
正正就是因為.
大家已經沒有了那個感動.
是弟兄姊妹的.
神賜給弟兄姊妹的.
一個屬靈的家.
這個屬靈的家是需要大家一起.
在感動底下去.
維護.
奉獻的意義就是.
在這裡是很清楚.
不是逼,不是勸.
也都不需要經常提.

$^{2321}$弟兄姊妹按感動.
按你的能力去做.
神要你做的工作.
這個就是.
他們有意識地.
看到.
使我們靈命持續的.
奮興.
其實重要就是.
神的殿要興旺.
祭司利未人是要無後顧之憂.
的服侍.
是要他們在神的殿裡.
是很殷勤的.
來參與每個禮拜.
的獻祭.
和其他的工作.
平日他們做教導.
和那些訴訟.
他們來做一個審判官.
諸如此類排難.
解那些糾紛.
就是利未人他們做的.
所以.
你可以想像得到,第34節.
我們的祭司利未人和百姓.
都祭簽看每年.
哪一個族.
按定期將獻祭的柴.
奉到我們神的殿裡.
照著律法上所寫.
燒在耶和華.
我們神的壇上.
各支派都受感動.
大家很甘心樂意.
來將聖殿裡.
要用的東西.
就是銅祭壇.
要用的乾柴來預備.
可能購買.

$^{2361}$很優質的.
燒起來沒有那麼多煙的.
那些柴.
來成為聖殿裡.
每個禮拜,其實是每天.
因為一起床.
他們就要在銅祭壇裡.
來管理那些火.
銅祭壇的火.
是不可以熄滅的.
首先整理那些灰.
然後就釀那些柴燒旺.
可能還有其他的東西.
他們要做的.
可能要煲水,要潔淨.
諸如此類.
第35節.
又定每年將我們地上初熟的土產.
和各樣樹上.
初熟的果.
都奉到耶和華的殿裡.
又照律法書上所寫.
將我們投胎的兒子.
守生的牛羊.
都奉到我們神的殿.
交給我們神殿裡.
公職的祭司.
每樣東西都做回.
摩西律法書所要求的.
覺得這份是感恩.
所以.
奉獻沒有了感恩.
就很大件事了.
一方面.
有時看那些電影.
那些黑手黨的大哥.
都會奉獻給天主教.
奉獻很多錢.
好像來者不拒.
但基督教就不是.

$^{2401}$一定是從感恩開始.
你看到他們現在.
是被神的靈,神的道.
感動,在感恩裡.
首先要給自己.
身份,神選民.
的身份,要有一個.
界線,然後有一個.
對人的憐憫.
然後就開始想起神的殿.
之前的荒涼.
大家都沒有盡心.
去為神的殿去想.
現在想到很多.
具體的行動.
就是安息日.
我們要怎樣.
他們的衰庸又是怎樣.
現在就有一個感恩制.
就是我們大家都願意.
將我們所獲得的.
地圖所種的.
我們自己所擁有的.
來奉獻給神的.
殿的衰庸.
然後連我們自己.
投生的兒子.
投胎的.
那些牛羊生畜.
生,對不起.
出埃及記的十二章.
凡是投生的.
都是神的.
因為神在埃及.
擊殺長子.
以色列的長子.
是因為有寶血的.
來倖免於難.
所以神就說.
投生的兒子和.

$^{2441}$那些生畜其實是屬於我的.
所以是要歸神為聖.
後來就是他們.
當然不是神要他們奉獻.
燒了那些.
藍丁.
而是他們要用那個贖銀.
來代贖.
即是意味著要有一個很強烈的.
給下一代.
特別是投生的兒子.
做奴隸.
你在埃及.
後來你生了那個兒子.
那個兒子本來是奴隸.
但是他現在是屬神.
的兒女.
你要灌輸這個首生.
你是神的兒女.
你要將他歸神為.
聖,要這樣的.
所以是神.
叫你代他去.
養你的子女.
不是你的,是你托管.
他,所以.
如果我們有這個意識.
我們就會有一種既緊張.
但同時間亦都放手.
讓神在他的.
生命裡面去管理.
但要有這個意識.
他們現在將投生的.
奉獻給神.
你可以想像得到.
這個是他們一個很強烈的.
復興的那個.
感覺,很盼望.
今年的培靈會.
帶動全香港.

$^{2481}$有一個復興.
我都提過.
中堂.
不記得名字.
北宣的堂主任.
他會說利未忌.
所以是.
值得聽的,盼望.
帶來全香港的基督徒.
有一個復興,或者全世界.
的基督徒,因為現在都在線.
可以用Zoom來看.
或者YouTube,盼望.
真的帶動一個復興.
你看看,一個復興裡面.
他們自己很.
甘心樂意為神.
做這一切.
甘心樂意.
為神去奉獻.
自己所有的.
甘心樂意.
自己多走一步.
不是去望周圍隔離.
而是我們有這個.
的份,能夠服侍神.
是我們的一個.
福氣.
這樣理解,這樣灌輸.
不是洗腦,是事實.
這是你能夠.
事奉神,是你的福氣.
不是.
每個人都可以有這樣的福氣.
如果我們能夠事奉.
所以回來祈禱會.
真是一個很大的.
福氣,是一個.
事奉,是一個神.
很體重的.

$^{2521}$其中一個聚會.
因為一班人.
很認真地奉耶穌的名祈禱.
在禮拜堂.
WIFI特別強勁.
我們真是向神.
為一些很認真的事.
為現在.
我們認真祈禱.
為最近的罪案.
很多的.
家庭問題,很多的.
一些精神的問題,健康的問題.
我們認真的祈禱.
神會在我們中間.
感動和垂聽.
這個就是.
第十章.
和第九章.
這個禁食.
披麻.
蒙灰,認罪.
和回到神起初.
對這個民族的期望.
他們現在很想.
回到那個狀態.
他們現在很主動.
自己.
來填上.
這些奉獻.
即是說.
好像我們現在有一個月的.
恩賜調查表.
其實是鼓勵弟兄姊妹一起.
事奉.
我們感動弟兄姊妹.
一起按神給你的.
恩典,恩賜.
來服侍.
所以我們求主讓我們能夠.

$^{2561}$在人的心裡面做.
感動的工作.
所以他們現在很願意.
將這一切都奉獻到.
神的殿裡面.
還有很多,第37節.
並將初熟的物質.
所磨的面.
和舉祭.
各樣樹上初熟的果.
新酒和油.
奉給耶和華,奉給祭司.
收在我們神殿的庫房裡面.
將我們地上.
所產的十分之一.
奉給利未人.
因利未人在我們一切的.
土產裡面.
當取十分之一.
這是利未記的十八章.
的教導.
民數記的十八章.
那裡有很清楚的教導.
以色列人一切的.
土產的十分之一.
是歸利未人的.
利未人的十分之一.
是歸給祭司的.
很清楚的教導.
現在他們就主動.
要回到摩西的律法書.
去教導.
奉到神殿的屋子裡面.
收在庫房中.
第三十九節.
以色列人和利未人要將五穀新酒.
和油為舉祭.
奉到修存聖所.
器命的屋子裡面.
就是沒錯了.

$^{2601}$就是公祭的祭司.
守門的歌唱.
所住的屋子.
這樣我們就不離棄.
我們神的殿.
即是說.
很強烈的意識.
他們是以神的殿.
作為他們認真的.
考慮.
和關心.
不會說.
我都很忙了.
我有我自己的生活了.
你們自己搞定它.
慢慢就去到一個.
惡性的屬靈的循環.
帶來.
他們在被擄的期間.
或者接近亡國的時候.
就出現這種.
逆轉的問題.
但現在.
將這個問題.
由一個逆轉的惡性循環.
變成一個良性的循環.
就是我們關心神的殿.
神的殿得到復興.
神的殿得到那個供應.
神的殿就興旺.
興旺就帶來.
我們的靈性又會有感動.
我們的修成又會有倍增.
你就會看到.
一個良性的循環.
是因為禁食.
是因為披麻蒙灰.
是因為有那個嚴肅會.
是因為之前有七天的.
奮卿會.

$^{2641}$是因為他們有摧國節.
是因為他們守住彭節.
是因為他們大家回去聽.
神籍著.
尼希米和以斯拉.
宣讀律法書.
的那個階段.
所以這裡是一個.
真正尼希米紀的.
中心的最重要的部分.
原來重建.
聖殿的城牆.
聖城的城牆.
那個只是一個開始.
需要有一個具體的界限.
但原來真正的.
界線是聖和.
族的那個界線.
他們要回去神.
那邊的神期望.
的那個分別為聖.
盼望這個也是我們錦宣家.
在這個心系吾家的.
齊躍動裡面不是一種口號.
而是大家看到.
人生真的很多的.
無常.
我們更加要把握機會.
一起去傳福音.
緊張神的家.
神的殿.
就會令到香港可能.
有得救了.
所謂有得救就是更加多人信耶穌.
得到救恩.
不管什麼環境.
什麼時代.
不管怎樣收緊.
緊張的是我們將福音講明白.
甚或連一些.

$^{2681}$有關的官員都信主.
又為什麼不行呢.
所以盼望我們都能夠.
有這樣的追求.
這樣的學習.
我們已經完了.
我們去到祈禱.
天父我們感恩讚美你.
因為我們用了一段時間.
讀了一段時間.
我們自己都很嚮往.
這個復興不是一種人為的製造.
是真是神的靈.
感動.
也都真是有尼欽米.
這些屬靈的領袖.
當年不辭勞苦排除萬難.
來到去重建.
成長.
他要付出的代價是很大.
很多時候.
面對的那個反對.
攻擊是很大.
但當真是看到現在.
那個修成期.
就是神的靈.
感動百姓的時候.
那份滿足感.
是蓋過那些困難.
和那些.
敵對的困苦.
多謝主讓我們看到.
第十章.
這個復興.
這個認真的大家簽名立約.
來到去回到神的.
摩西律法書裡面.
這樣的一個狀態.
願你感動我們.
帶領我們感宣家.

$^{2721}$一起在不同的崗位.
同心的去服事.
奉耶穌基督的命.
阿們.
小龍哥.
你是怎麼樣.
他撐下去了.
那我們就別看他了.
我們還是回家 好不好.
好.
\newpage



\section{路加福音 7:36-50}
\label{sec:Fl2IkOLpxSI}
\textbf{2023-06-10 晚堂《罪人的筵席》路加福音 7\_36-50 講員:吳景霖傳道}
\newline
\newline
連結: \href{https://youtube.com/watch?v=Fl2IkOLpxSI}{\texttt{https://youtube.com/watch?v=Fl2IkOLpxSI}} ~~~~ 語音日期: 2023-06-13
\newline
\newline
\hyperref[sec:qmt5oGwLkF4]{\small{< < < PREV SERMON < < <}}
~
\hyperref[sec:index_chronic]{\small{[返順時目]}}
~
\hyperref[sec:index_scriptual]{\small{[返順卷目]}}
~
\hyperref[sec:A_K04hyrSP8]{\small{> > > NEXT SERMON > > >}}
\newline
\newline
路加福音 7:36-50
\newline
\begin{longtable}{cl}
\hline
\hline
章節 & 經文 (和合本修訂版)\\
\hline
7:36 & \begin{tabularx}{0.7\textwidth}{X} 有一個法利賽人請耶穌和他吃飯，耶穌就到那法利賽人家裡去坐席。 \end{tabularx} \\ \\ \relax
7:37 & \begin{tabularx}{0.7\textwidth}{X} 那城裡有一個女人，是個罪人，知道耶穌在法利賽人家裡坐席，就拿著盛滿香膏的玉瓶， \end{tabularx} \\ \\ \relax
7:38 & \begin{tabularx}{0.7\textwidth}{X} 站在耶穌背後，挨著他的腳哭，眼淚滴濕了耶穌的腳，就用自己的頭髮擦乾，又用嘴連連親他的腳，把香膏抹上。 \end{tabularx} \\ \\ \relax
7:39 & \begin{tabularx}{0.7\textwidth}{X} 請耶穌的法利賽人看見這事，心裡說：「這人若是先知，一定知道摸他的是誰，是個怎樣的女人；她是個罪人哪！」 \end{tabularx} \\ \\ \relax
7:40 & \begin{tabularx}{0.7\textwidth}{X} 耶穌回應他說：「西門，我有話要對你說。」西門說：「老師，請說。」 \end{tabularx} \\ \\ \relax
7:41 & \begin{tabularx}{0.7\textwidth}{X} 耶穌說：「有兩個人欠了某一個債主的錢，一個欠五百個銀幣，一個欠五十個銀幣。 \end{tabularx} \\ \\ \relax
7:42 & \begin{tabularx}{0.7\textwidth}{X} 因為他們無力償還，債主就開恩赦免了他們兩個人的債。那麼，這兩個人哪一個更愛他呢？」 \end{tabularx} \\ \\ \relax
7:43 & \begin{tabularx}{0.7\textwidth}{X} 西門回答：「我想是那多得赦免的人。」耶穌對他說：「你的判斷不錯。」 \end{tabularx} \\ \\ \relax
7:44 & \begin{tabularx}{0.7\textwidth}{X} 於是他轉過來向著那女人，對西門說：「你看見這女人嗎？我進了你的家，你沒有給我水洗腳，但這女人用眼淚滴濕了我的腳，又用頭髮擦乾。 \end{tabularx} \\ \\ \relax
7:45 & \begin{tabularx}{0.7\textwidth}{X} 你沒有親我，但這女人從我進來就不住地親我的腳。 \end{tabularx} \\ \\ \relax
7:46 & \begin{tabularx}{0.7\textwidth}{X} 你沒有用油抹我的頭，但這女人用香膏抹我的腳。 \end{tabularx} \\ \\ \relax
7:47 & \begin{tabularx}{0.7\textwidth}{X} 所以我告訴你，她許多的罪都赦免了，因為她愛的多；而那少得赦免的，愛的就少。」 \end{tabularx} \\ \\ \relax
7:48 & \begin{tabularx}{0.7\textwidth}{X} 於是耶穌對那女人說：「你的罪都赦免了。」 \end{tabularx} \\ \\ \relax
7:49 & \begin{tabularx}{0.7\textwidth}{X} 同席的人心裡說：「這是甚麼人，竟赦免人的罪呢？」 \end{tabularx} \\ \\ \relax
7:50 & \begin{tabularx}{0.7\textwidth}{X} 耶穌對那女人說：「你的信救了你，平安地回去吧！」 \end{tabularx} \\ \\
[1ex]
\hline
\hline
\end{longtable}
$^{1}$好 弟兄姊妹 主佬們平安.
最後一次在這裡晚堂跟大家講道.
天下無不散之筵席.
希望將來有機會再跟大家分享.
今天這個故事是一個筵席的場面.
是一個很特別的故事.
為什麼這麼說呢.
因為在《馬太福音》和《馬可福音》.
其實都記載了耶穌去一個叫做長大麻風的西門.
家裡吃飯的.
而在《路加福音》 即是今天這一段.
就說到去一個法利賽人家裡.
法利賽人家裡.
不過我們在對話裡面就知道.
原來那個法利賽人又是叫西門.
好了 《約翰福音》就沒有記載.
耶穌去誰的家裡吃飯.
不過在其他福音書裡面.
即是四卷福音書.
都是同樣地有一個女人.
拿著一瓶香膏闖入了筵席裡面.
馬太和馬可就說.
她用這個香膏去勾耶穌的頭.
然後馬可和馬太.
馬可和馬太的高頭是路加和約翰.
就勾耶穌的腳.
三卷福音書 對觀福音.
即是馬太 馬可和約翰.
都說這個女人勾完耶穌之後.
那些門徒就有反應了.
就說這麼貴的一瓶香膏.
你浪費了它.
可以拿來周祭窮人.
那麼耶穌當然就不OK了.
他就去支持那個女人.
就說這個女人所做的.
是預先為我安葬而做的.
所以大家都看到.
在這麼多卷福音書裡面.
雖然那個記載是有出入的.

$^{41}$但我們估計它其實是在說同一件事.
是同一件事來的.
今天這個經文.
我們說路加福音的記述.
我想問一下這個故事.
這個敘事裡面.
誰是主角.
你不要回答我耶穌.
耶穌當然是主角.
但是主要角色除了耶穌之外.
一般來說我們將焦點放在女人身上.
我們會說放在女人身上.
她在耶穌身上做了這麼多事情.
然後耶穌又稱讚她.
當然主角是她.
但是我想問一下.
假設你是一個作者.
你寫一本書.
你去記載一些事情.
一個最主要的主角.
你會不會不給他名字.
不給他名字.
為什麼呢.
為什麼西門就有名字.
但是這個女人就沒有名字.
是不是門徒不知道呢.
不知道可以問的.
就算真的不知道.
就給他一個名字也可以.
如果你是作者的話.
我印象中我看這麼多故事.
看這麼多電影.
看這麼多電視劇.
我真的四十多年.
我只看過一套戲.
女主角是沒有名字的.
然後其他全部.
基本上一定有名字的.
因為很難去處理的.
如果一個角色完全沒有名字.

$^{81}$聖經裡面有沒有故事是沒有名字的.
有的.
例如哪個呢.
年輕的財主.
記不記得.
跑去耶穌前面跪下的那個.
他就沒有名字.
拉薩路和財主.
拉薩路在財主家門口黑吃的故事.
財主是沒有名字的.
格拉森那個被鬼附的人.
我們很怕說這些唐道的原因.
就是每次說就是格拉森被鬼附的人.
被鬼附的人都很繞口.
為什麼這些角色沒有名字.
他沒有名字的原因.
很可能就是他的角色.
已經是最具代表性的.
舉一個例子.
例如我去學校的家長日.
老師要見家長.
他叫我做什麼呢.
就叫天林爸爸.
你發覺你生了個兒子之後.
你沒有了自己的身份.
老師他不care你叫什麼名字的.
他不會叫我做吳全道.
他就是天林爸爸.
而那一刻他都只不過是期望.
我用天林爸爸的身份去回應他.
他不會希望我跟他說唐道.
或者說福音.
他就是想跟我說我兒子的成績.
那個角色你的身份就是最重要的.
所以說.
財主和拉薩路的故事.
財主叫做Simon 還是Peter.
到底叫做Patrick 還是Andrew.
是對故事沒有影響.
對故事沒有影響.

$^{121}$財主兩個字已經是他最重要的身份.
是不是.
如果我說你要找一個適合的身份去形容你自己.
例如有人剛認識.
要介紹一下.
認識一下.
你會用什麼去介紹自己呢.
你給自己一個身份.
你會叫自己做.
我是一個爸爸.
或者一個媽媽.
或者我是別人的兒子.
我是一位老師.
我是一個健身教練.
我是一個音樂人.
例如形容我做雷偉基.
感受張學友.
這樣你就大概知道他的特性是什麼.
是不是.
你會說他是音樂人.
他唱歌好聽.
或者有些人是.
他是一個巴士迷.
他是一個攝影發燒友.
他是一個貓奴.
他最喜歡養貓.
這些東西就可以足以用一個很短的片語.
去形容那個人的特性.
好 聖經說.
那個人是財主.
那是什麼.
對他來說有錢就是他最重要的一個特質.
如果別人介紹你.
例如用雷偉基來說.
雷偉基啊 財主啊.
這樣就慘了.
他可以說窮得只剩下錢.
沒有其他特質.
就是有錢.
這個人這樣.

$^{161}$你要明白聖經裡面.
這些這樣的象徵性的形容.
所以你日後看回財主的比喻的時候.
不是你有錢難上天堂.
而是你生命裡面只有錢的話.
你就難上天堂了.
這些東西最後希望長氣一點.
教大家怎樣去看聖經.
所以我們去判斷.
用這些方法.
你會發覺這個女人.
雖然耶穌對她是很正面的說法.
但是在這個敘事裡面.
她不是最重要的.
她用女人這兩個字已經足以形容她了.
我不是小看姐妹.
不然等一下被人告我歧視.
這個故事反過來.
男人發生的話.
可能就男人這樣已經解決了.
所以說這個女人.
她引起了一場風波.
她在言直上面引發了一些事件.
但是她並不是最重要.
我們不是說一看這個聖經.
然後就學習一下那個女人.
相反在這個段落.
就是這個故事裡面.
最不值得學習的那個.
有名有姓的就是西門.
我們發覺她才是最重要的角色.
我們在這個敘事裡面.
我們要留意的就是西門.
其實其他福音書都是.
你看看重點都不是那個女人.
重點是那些門徒.
那些門徒不是關心那些窮人.
他關心那個香膏可以賣多少錢.
約翰福音更加不給面子.
講明就是猶大去問的這句話.

$^{201}$因為他是個賊.
你懂得看這些說法.
我們看回經文.
第36節開始說.
法利賽人請耶穌吃飯.
耶穌就到法利賽人家裡坐直.
那城裡有一個女人是個罪人.
知道耶穌在法利賽人家裡坐直.
就拿著聖滿香膏的玉瓶.
站在耶穌背後.
挨著他的腳哭.
眼淚滴濕了耶穌的腳.
就用自己的頭髮擦乾.
又用嘴巴連連親他的腳.
把香膏抹上.
當時的筵席是容許一些乞丐.
去到裡面不請自來地乞討一些食物.
不是像我們現在這樣.
擺酒門口有幾個人叫你簽名.
又對你坐在哪裡.
是容許一些乞丐去到裡面.
去討食物.
那個女人這樣不請自來.
其實就是說她自己是困在乞丐裡面.
她是一種相當謙卑.
她視自己為乞丐.
去到筵席裡面.
而她是一個出名的罪人.
當時一個女性還要出名.
多數的情況可能就是她犯了姦淫.
有可能她是一個妓女.
不過在原文裡面.
這個字其實是一個過去式.
就是說有可能她是過去犯過姦淫.
或者她是過去的一個妓女.
但已經是停止了.
這個女人去到耶穌面前.
做了一連串的行為.
那些行為代表著她那種內心的懊悔.
她對自己罪的後悔.

$^{241}$39節就這樣說.
「請耶穌的法利錯人看見者是心裡說.
『這人若是先知,一定知道我他的是誰』.
是個怎樣的女人?她是個罪人?」.
這段是心裡想.
即是內心獨白.
又奇怪.
我心裡面在想.
為什麼作者會知道呢.
所謂無所不知的作者.
作者是知道所有事情的.
因為他寫了一個故事.
但他寫下去這個內心獨白.
其實真的不知道為什麼他知道.
可能是後來他們問回耶穌.
因為耶穌一定知道.
但是不要被這些東西阻礙我們理解經文.
聖經這樣寫就是了.
你當作已經知道.
這個內心獨白其實反映兩件事.
第一就是.
他對耶穌的身份.
先知的身份質疑.
他的邏輯是怎樣呢.
如果耶穌是先知.
他沒有理由不知道這個女人是罪人.
如果他知道這個女人是罪人.
那就沒有理由被他摸.
傳統上他們都覺得罪人是不潔的.
是不應該接觸的.
姑勿論那個女人的動作是多表示自己的懊悔.
你是罪人就罪人吧.
法利賽人就是這樣.
你不應該和罪人有接觸的.
要保持距離的.
所以就用這個邏輯.
耶穌被那個女人接觸.
他就不是先知了.
他的邏輯就是這樣.
另外第二點就是反映他對那個女人的傾向態度.

$^{281}$剛才也說過.
他過去是色尼亞可能是以前犯罪的.
但是如果當場被人抓到他犯姦淫.
可能被扔死了.
找石頭.
他可能是一個過去犯罪的人.
過去是一個妓女.
但是重量了.
不過在西門的眼中.
罪人就是罪人.
犯法就是犯法.
總之就不值得原諒.
總之這個就是罪人.
不應該出現在這裡.
他對這個女人所做的懊悔的態度和行為.
是完全視而不見的.
罪人就是罪人.
西門這樣想而已.
他沒有說出口.
他沒有說出口其實也算給面子.
我不攤出來說.
不過我就發覺耶穌很壞.
經常都是這樣.
你每次喜歡隱藏隱瞞的事情.
耶穌就撩出來說.
他就沒有放過西門.
他就跟她說西門我有話對你說.
一聽到這些.
我們做老公的一聽就害怕了.
一跟你說.
雷偉基有話跟你說.
就出事了.
如果老婆這樣跟你說就大件事了.
總之他這樣說不放話.
他就是要西門.
就是那個法理罪人.
接下來我有話要教你.
要說你.
41節.
耶穌說有兩個人欠了某一個債主的錢.

$^{321}$一個欠五百個銀幣.
一個欠五十個銀幣.
因為他們無力償還.
債主就開恩赦免了他們兩個人的債.
那麼這兩個人哪一個更愛他呢.
西門回答.
我想是那多得赦免的人.
耶穌對他說你判斷得不錯.
當時一個人一天的工資大概就是一個銀幣.
那麼想想五十個銀幣就差不多兩個月工資.
一個多兩個月工資.
五百個就已經欠了人家一年零四個月的工資.
就是一年零四個月不吃不用才還得到.
當然相差很遠.
但問題是其實兩者都是無力償還的.
他都是沒有錢還的.
於是債主就開恩赦免了他們.
那麼耶穌問哪一個更加愛債主呢.
要留意一件事就是希臘文其實是沒有一個字去感激這個字.
他們就是用愛來代替了.
愛就包含了感激的意思在裡面.
所以他這裡說哪一個愛多些呢.
其實就是哪一個感激債主多些.
西門就回答多得赦免的人.
44到50節.
於是他轉過來向著那女人.
就是看著那女人但是跟西門說.
你看見這女人嗎.
我進了你的家你沒有給我水洗腳.
但這女人用眼淚濕了我的腳.
又用頭髮擦乾.
你沒有親我.
但這女人從我進來就不住你親我的腳.
你沒有用油膏抹我的頭.
但這女人用香膏抹我的腳.
所以我告訴你.
他許多的罪都赦免了.
因為他愛的多.
而那少得赦免的愛的就少.
於是耶穌對那女人說.

$^{361}$你的罪都赦免了.
同直的人心裡說這是甚麼.
竟赦免人的罪呢.
耶穌對那女人說.
你的信救了你.
平安地回去吧.
耶穌將一連串三個動作.
那個女人的三個動作.
去到和西門的接待態度去比較.
第一個通常當時的人.
他們是會預備一缸水在門口給人洗腳的.
因為不像我們這樣.
有磚頭有石屎的路.
走那些泥路是沙塵滾滾.
所以門口一定有水缸.
水便走的神蹟裡面.
為甚麼有水缸呢.
水缸的水可能是不令我們洗腳的.
沒有預備這樣算不算犯法呢.
也不是的.
沒有甚麼大錯的.
只是沒有準備.
就好像你家裡請人來.
你沒有預備拖鞋.
也不至於是很過分.
但是招待不周.
但是親嘴是當時一種基本的問候方式.
就好像我們見到一些貴賓.
你跟他握手甚至擁抱.
這些差不多.
耶穌受邀請去別人家裡吃飯.
按道理就是賓客.
無論如何都是客人.
如果你請客人來你家.
人家一進門口.
你就坐在沙發裡.
繼續用電腦按電視.
就這樣看著他.
這樣很沒禮貌.
是不是應該起身應一應.

$^{401}$接待.
其實西門的這個動作.
是相當之沒禮貌.
對耶穌不是很客氣.
用油膏的客人.
對客人表示尊重.
一般是用橄欖油.
橄欖油我們現在覺得很珍貴.
用橄欖油好像貴一點.
但是當時是很普通.
很普遍的一件事.
反而西門沒有做之餘.
女人還用很貴的香膏去膏耶穌.
三個比較.
你看得出西門其實不是很尊重耶穌.
你不要想請他吃飯就很尊重.
你看聖經第六章開始.
其實文士和法利賽人.
已經是非常之不滿耶穌.
他們在窺探耶穌.
希望找到耶穌的錯處.
去對付他.
所以這個飯局其實未必是懷著善意.
可能普通的招呼.
甚至是看看飯局裡面有什麼問題都未定.
我們特別留意47節.
剛才我說希臘文是沒有感激這個字.
用愛這個字.
所以你看回這句可能意思是.
你從他感激的態度.
你可以明白他有很多的罪得去赦免.
你從他感激的態度.
你可以明白他有很多的罪得去赦免.
反過來西門的態度是這樣的時候.
你就可以理解.
你都沒有什麼感激.
甚至你不懷好意.
你得到的赦免本來就是少.
大家不要搞錯了.
很多時候去到這些位置就下結論了.

$^{441}$哦 這樣我感激耶穌多一點.
我就得到多一點赦免.
邏輯不是這樣的.
不是說你感激多一點.
你的赦免就多.
這個是說你表達出感激的時候.
反映著你裡面有很多的東西得到赦免.
不是說你去感激耶穌多一點.
這些是不能裝的.
也不能故意去感激.
這是一種自然的流露.
感激是一種自然的流露.
當你是一個知罪的人.
你知道我有犯罪.
我是一個罪人.
而耶穌赦免了我的時候.
你很自然地流露出一個感激.
事實上很多人是不知罪.
是不知自己有罪.
是不覺得自己需要被赦免.
有很多人我們去到傳福音.
或者跟他決志的時候.
卡住的都是這個位置.
我沒有罪.
為什麼我要認罪.
我相信耶穌可以永生.
我沒有犯罪.
我規行舉報.
我一輩子都沒有犯過法.
螞蟻我都不踩死一隻.
我沒有罪.
你要說的那些妒忌.
小器 貪心.
那些有的那麼小事.
怎麼算罪呢.
很多人是懷著一種這樣的心態.
他們不覺得自己有罪.
也不覺得自己要被赦免.
剛剛我剛才在這裡說的時候.
在說這個故事.

$^{481}$最重要的角色是誰.
西門.
西門這個有名有姓的人.
不是那個女人.
是說西門.
這個故事.
耶穌是在教訓西門.
對於我們來說.
其實我們怎麼應用.
到底這個故事是對什麼人說的.
是不是對一些未信耶穌的人說的.
叫你要懊悔要痛悔.
不是的.
你看回你想一想.
他是對一個在信仰裡面的人說的.
西門是一個法利賽人.
法利賽人是有知識的.
在當時來說.
是一個文盲社會.
那時候是文盲社會.
但是法利賽人是要讀很多書的.
你看寫這麼多保羅書信.
叫我們寫一封都搞不定.
這麼多字.
還要寫這麼多書信.
還要有很多流失了.
其實那只是一部分被記錄下來.
你想想在當時的社會裡面.
法利賽人是非常高知識的.
放回現代.
他可能有兩三個博士學位.
一個人的時候.
他又要在信仰裡面.
又知書識禮.
又要有差不多的社會地位.
原來西門是一個這樣的人.
而耶穌在教訓一個這樣的人.
對於我們來說.
原來我們越信得久.
我們越有知識.

$^{521}$我們越有地位.
我們就要想一想.
段經文其實是不是在跟你說.
段經文是不是在跟你說.
你回了這麼多年教會.
這件事是耶穌對西門發起的一個挑戰.
也是作者要告訴我們.
原來一個人在信仰裡面停留得久.
或者他知識高.
或者他有這麼少的身份地位.
他可能很容易落入一種自以為是的情況.
怎樣才算自以為是呢.
其實同樣地在路加科罕我們會見到.
有一個經文18章9-14節.
耶穌向那些自以為是而藐視別人的人講了這比喻.
有兩個人去聖殿禱告.
一個是法利賽人 一個是遂利.
法利賽人獨自站著自言自語地禱告說.
也不像這個遂利.
我每週禁食兩次 凡我所得的都獻上十分之一.
你說自言自語而已.
這麼清楚 還要擺明是要求一下旁邊的遂利.
我才不像他.
或者祈禱的時候.
我才不像旁邊的港島睡覺.
擺明是要求旁邊的人.
但是那個遂利又怎樣.
那遂利遠遠地站著 連舉目望天也不敢.
只隨著胸說 神啊 開恩可憐我這個罪人.
我告訴你們這人回家去.
被那人倒算為異了.
因為凡至高的必降為卑.
凡今卑微的必升為高.
其實同一本書裡面.
很一致的邏輯.
法利賽人他們以為自己很高尚.
事實上他們在社會上有很高的地位.
有很多的學識.
在宗教信仰裡面.
他們是非常認真和非常投入.

$^{561}$做得很好.
但是其實那個遂利.
或者在西門飯局裡面的那個女人.
他們是更加貼近信仰.
他們知道自己是罪人.
面對上帝他們只是夠膽.
很卑微 頭都不敢抬起.
就在那裡扔心口祈禱.
那個女人她根本就不夠膽去接觸耶穌.
就是靠著腳.
你要知道腳是人最低的地方.
象徵著一個最卑微的.
她靠著腳 親她的腳.
其實就是代表自己那種卑微.
她用頭髮.
一個人最珍貴的.
女士我相信.
你的頭髮應該是最珍貴的.
不像我們男人隨便找個人剪就行了.
女士可能是最重視自己的頭髮.
用來抹別人的腳.
她那種謙卑.
她那種對自己的.
認為自己是不配的那種對罪的懊悔.
是上帝 耶穌最欣賞的.
相反西門她表達出一種.
她覺得自己高人一等.
她以為自己可以用我的標準去審判那個女人.
你這個女人都不懂避忌.
這個飯局是你來的嗎.
耶穌是你碰的嗎.
然後她還審判了耶穌.
審判了耶穌.
你都不知道是不是先知.
明知的.
你都不知道那個女人什麼樣子.
是不是.
然後可能覺得耶穌沒什麼.
接待耶穌請他吃飯不代表悔改.
有本書我覺得名字很棒.

$^{601}$叫《穿起西裝不會變高層,翻了教會不會變好人》.
是不是很棒.
我還沒有機會看.
但是書名很棒.
我看過目錄都很棒.
它說的是一個事實.
人翻了教會不一定會變好.
在教會裡面不一定會進來頭頂有個光圈.
其實好像西門一樣.
他對著耶穌請耶穌面對面和他吃飯.
其實他不知道悔改.
他內心沒有一個感覺.
他覺得我需要赦免嗎.
不需要.
他沒有一個懊悔在裡面.
他會用自己的標準去審判人.
女人是罪人.
根本不應該出現在這裡.
耶穌你也根本不應該接觸這個女人.
但是他沒有想過.
他其實現在是五十步笑百步.
欠五十個銀錢的就欠五百個銀錢.
五十這個詞語很好用.
我們作為一個信徒.
我們翻了教會這麼久.
其實我們應該要很反省.
我們有沒有自以為是.
我們有沒有像法利賽人那樣自以為是.
用我們自己的眼光去審判其他的人.
審判一些犯了罪的人.
或者是他給我們卑微的人.
今天的社會不是這樣.
社會和聖經的教導完全背道而馳.
對一些人他做錯了事.
喊打喊殺.
幾乎是用一句往死裡打.
一定要想辦法把他打死.
一定要想辦法追殺他.
去到天腳底.
要他永不翻身.

$^{641}$為什麼人犯了錯.
我們不可以去原諒.
我們不可以去寬恕.
甚至有些教會他們願意去接待罪人.
去接待一些做錯事的人.
我們不欣賞.
相反我們會覺得為什麼會容納他們.
教會什麼都不是.
就是讓罪人去到罪集.
我們每個人都是罪人.
我們欠50銀.
你很意思去取笑別人欠500銀.
為什麼不可以讓他們找個地方去親近耶穌.
為什麼不可以讓他們有個地方透一透氣.
今天的講題是罪人的言直.
指責別人犯罪很容易.
很容易就說那個人不對那個人不對.
特別是現在網絡世界.
很恐怖.
你上網看他們那些不需要出樣的言論.
是可以很恐怖.
你約他對質他又不夠膽跟你面對面說.
但是網上很多人說的東西是很恐怖.
指責別人犯罪犯錯是非常容易.
但是相反要承認自己錯.
承認自己有罪是非常困難.
是非常困難.
如果一間教會我們擺出一種高高在上的姿態.
我們只是一種.
每個人都好像聖人一樣.
每個人都沒有犯過錯.
每個人都有身份有地位.
每個人都好像很光鮮.
這樣的姿態.
結果就令到一些罪人不夠膽再回來教會.
每個人都比我好.
我實在自尋認愧.
於是離開教會.
又或者是有罪的人.
他不夠膽在教會裡面誠實.

$^{681}$不夠膽告訴大家.
其實我也有些軟弱.
我有些做錯.
我有些犯了罪.
不夠膽去展示.
因為每個人都那麼好.
於是這個教會成為一種50個銀幣的教會.
每個人都被人感覺犯了一點點錯.
每個人都是好人.
每個人都是聖人.
這樣的教會其實出現一種情況.
我們可能經常崇拜聚會.
我們去接待耶穌.
但事實上我們只是為了自己去享受一個現職.
為自己去做一個現職.
但其實耶穌在裡面.
他根本就不自在.
他也見不到他最想見到的一些懊悔痛悔的人來跟前.
每一個崇拜.
我們要想清楚我們到底在接待誰.
我們當然最重要的是敬拜上帝.
但教會的門打開的時候.
我們是要接納每一個犯罪的人.
每一個做錯事的人.
他只要願意進來教會去找耶穌.
我們應該打開門歡迎他.
我們要盡我們的能力讓他感到自在和舒服.
讓他可以在耶穌面前去表達他自己的痛悔心情.
我希望大家去思想一下.
我們要理解耶穌的心腸.
這個罪人的現職.
如果我們只顧著自己.
最終犯罪得罪上帝的可能是我們自己.
我想大家現在有一分鐘的時間.
我們安靜.
我們閉上眼向上帝禱告.
如果我們自己本身有犯罪.
我們向上帝認罪.
包括如果我們有不接待人的罪.
我們有自異的罪.

$^{721}$或者我們有將人推離教會的罪.
小看一些犯罪的人的罪.
我們都一一求上帝赦免.
我們有一分鐘的時間去安靜.
\newpage



\section{馬太福音 9:35-38}
\label{sec:A_K04hyrSP8}
\textbf{2023-06-11 加恩樓《建立無牆的教會》馬太福音 9\_35-38 講員:鄧泰康傳道}
\newline
\newline
連結: \href{https://youtube.com/watch?v=A-K04hyrSP8}{\texttt{https://youtube.com/watch?v=A-K04hyrSP8}} ~~~~ 語音日期: 2023-06-13
\newline
\newline
\hyperref[sec:Fl2IkOLpxSI]{\small{< < < PREV SERMON < < <}}
~
\hyperref[sec:index_chronic]{\small{[返順時目]}}
~
\hyperref[sec:index_scriptual]{\small{[返順卷目]}}
~
\hyperref[sec:lzufA9tLqUc]{\small{> > > NEXT SERMON > > >}}
\newline
\newline
馬太福音 9:35-38
\newline
\begin{longtable}{cl}
\hline
\hline
章節 & 經文 (和合本修訂版)\\
\hline
9:35 & \begin{tabularx}{0.7\textwidth}{X} 耶穌走遍各城各鄉，在他們的會堂裡教導人，宣講天國的福音，又醫治各樣的病症。 \end{tabularx} \\ \\ \relax
9:36 & \begin{tabularx}{0.7\textwidth}{X} 他看見一大群人，就憐憫他們；因為他們困苦無助，如同羊沒有牧人一樣。 \end{tabularx} \\ \\ \relax
9:37 & \begin{tabularx}{0.7\textwidth}{X} 於是他對門徒說：「要收的莊稼多，做工的人少。 \end{tabularx} \\ \\ \relax
9:38 & \begin{tabularx}{0.7\textwidth}{X} 所以，你們要求莊稼的主差遣做工的人出去收他的莊稼。」 \end{tabularx} \\ \\
[1ex]
\hline
\hline
\end{longtable}
$^{1}$(陳志全) 各位節目早安.
今天跟大家分享.
其實這是最近一份很大的感動.
而出來的講章.
是一篇很簡單的講章.
在聖經裡也很直白的講章.
在近幾年的時候.
其實在教會當中.
特別是一些我們去思想.
教會突破的時候.
其中有一個經常出現的字.
叫做無牆教會.
所以你會發覺.
當我們講到教會要無牆.
其實是一個很具體化的.
告訴大家為什麼.
平時給人的印象就是.
堂會一個教會就是四縫牆.
一個崇拜地方.
人家要進來.
人家要來到我這個地方.
但我不知道大家知不知道.
其實有一個觀念.
其實這是一個跨文化.
是人家不是你的文化.
然後人家要走進來你的文化.
你明不明白這個分別.
最近有很多人都會講移民.
你移民到去其他地方.
你需不需要適應.
適應容易嗎.
所以我想強調的就是.
在過去我們經常有一個觀念.
就是怎樣.
我想大家留意.
不是錯的.
我想說這是一個歷史的進程.
過去我們教會.
或者是普遍香港的教會.
我們都會講一件事.

$^{41}$我們開報道會.
然後大家邀請朋友來到當中.
去參加報道會.
然後我們跟他們講一些訊息.
但其實在這個過程裡.
他們要跨過這道牆.
是一點都不容易的.
第一.
要跟你有多友好.
他才肯來.
有些人都說.
認識了很多年.
很多次回教會他都不想回.
通常有兩個情況.
第一是給面子派對.
我真的給你面子.
認識了你這麼多年.
你叫了我很多次.
我不好意思所以我來.
第二就是他真的碰巧在人生裡遇到一些問題.
他嘗試尋求幫助或答案.
通常是這樣.
所以剛才我說你要適應很不容易.
特別是我成長的時候.
有很多字眼在教會講.
是很普遍的.
很正常的.
教會甚麼叫交通.
交通.
甚麼意思.
在教會裡面.
你講交通.
交通當然不明白.
在教會裡面分享.
在世界裡甚麼叫交通.
就叫搭車.
你為甚麼叫我搭車.
甚麼叫靈修.
你明白嗎.
其實我們想說.

$^{81}$當我們今天在講無牆的教會的時候.
首先我們要拆除我們基督徒.
我們自己心裡的那一幅牆線.
其實我們是不知不覺之間.
我們會建立了很多我們信仰上的牆.
我們會覺得來到這個地方的時候.
就會有一些東西是有規矩的.
你記住.
今天試聽我們講的.
我不是回來教會變成沒有規矩.
我不是在講拆牆就是沒有規矩.
我不是在講這件事.
而是我們想說的很多時候.
我們慢慢慢慢的時候.
我們慢慢變成了一個甚麼地方.
是一個宗教場所.
是一群信耶穌的人.
每個星期來到的時候聚集的人.
所以很簡單的觀念就是.
我不是信耶穌的.
我就不來.
其實這句說話.
我在嘉欣樓.
這個地方.
在近幾年我聽得最多.
有很多.
現在我很感恩.
因為我早上透過一些講座.
透過機會其他緣故.
認識了一些街坊.
街坊上來之後.
很奇怪的.
他們的反應基本上都是一致的.
我在這裡住了很多年.
我從來都不知道.
上面有一個這樣的地方.
突然間會發現寶藏.
原來上面可以坐的.
這麼多東西可以學的.
原來這裡是這樣的.

$^{121}$好像發現一個新大陸一樣.
我問他:這麼多年你也經過.
為甚麼你不想來.
教會來的.
你們是教堂.
我又不是教徒.
那我就不來.
你會發覺原來這些就是一些場.
我們不只是宗教場所.
更加我們不是私人會所.
我們是神的家.
我們是神的殿.
這裡是可以找到上帝的.
我們這裡不只是心靈安慰的地方.
我們這裡是找到寶貴的舊恩.
所以我們要搞清楚.
很多時候我們回來之後.
慢慢地我們就會開始.
當然我經常說傳福音當然有一定的困難.
但問題是.
很多時候我們心裡就建立了這幅牆.
慢慢地就變成我們越來越佛系.
來的話就來的.
不來就不來的.
大弟兄姊妹.
我們沒有辦法逼人來的.
大弟兄姊妹.
你們有沒有禱告.
有沒有為了這些朋友禱告.
有沒有多走一步.
正因為這裡是神的家.
神的殿.
我們所給予的是寶貴的舊恩.
所以我們不惜一切代價.
都要完成上帝交給我們的使命.
將這個舊恩交到微信者的手上.
而去拆去一切讓人無法得到舊恩的阻礙.
其實我要為了嘉欣樓.
我是已經很感恩的.
我們已經不斷地向這個方向向前走.

$^{161}$大家知不知道嘉欣樓最厲害的地方是什麼.
就是拆門拆牆.
大家平時在打乒乓球.
不是乒乓球.
是打氣墊球.
夾公仔機的位置.
通往辦公室中間.
那裡有一個洞.
現在叫一個洞.
大家知不知道以前有一個門在那裡.
不知道吧.
但後來的時間就會發覺.
那裡阻礙了人進來.
你明不明白.
因為一封了之後.
前面有一個接待處.
人們就覺得辦公室是不會進來的.
然後時間就拆了它.
打通它.
你就會發覺我們做了很多的事.
是不斷地將所有阻礙人進來的地方.
我們都將它拆掉.
所以你會發覺這個工程.
其實一直到今天都沒有停過.
最近我們很努力地在做的咖啡事工.
大英姐妹.
咖啡事工真的不是做生意.
你弄清楚.
咖啡事工不是為了將咖啡給別人.
我一直和未來的咖啡大使不斷告訴他們.
就是沖一杯咖啡.
其實學了之後不是很難.
我沖一杯咖啡.
由我拿粉到出現咖啡的時候.
特別那本叫什麼.
我要給你的時候.
前後兩分鐘就給了你.
所以我鼓勵大英姐妹.
學一下拉花.
知道為什麼嗎.

$^{201}$除了感覺咖啡好喝.
另外就是拖延時間.
你明白我的意思嗎.
你慢慢幫她弄.
你用心去弄.
你拉一下花線.
然後慢慢給她.
拉花的時候可以聊天.
可以聊天.
所以我以前第一次接觸咖啡.
我是先接觸手沖咖啡.
因為手沖咖啡是不會用機器的.
我們是用手慢慢磨的.
然後就有很多道理說.
因為機器熱度要這樣那樣.
然後慢慢磨就要時間.
磨完之後倒出來.
水又要煲又要時間.
慢慢沖又要時間.
然後再倒出來.
都是時間.
為了咖啡可以跟對方打開話盒子.
也因為咖啡是一個國際語言.
所以無論你信不信.
你都會突破第一件事.
所以大家會發覺.
到遲一點.
我們那個house cafe正式開業的時候.
我們會長期將玻璃門打開.
玻璃門是會打開的.
我們會想盡方法.
現在人們說來這裡是難的地方.
我也不知道誰進來.
其實這裡很尷尬那邊也很尷尬.
遲一點就會想方法.
怎樣讓人們知道原來這裡有咖啡喝.
我跟富哥一起聊的時候.
其中有一個異象.
什麼叫異象?.
大家有沒有發覺.

$^{241}$這裡才是一個主要集中的地方.
以前那裡是走廊.
但現在其實很希望.
咖啡閣很奇妙.
當它設置之後.
成為一個聚腳的地方.
我們的期望就是.
那些朋友.
由一杯咖啡.
一直走走走走.
走進神的殿堂中.
頂展門我想說.
我們心裡需要拆去一些很傳統的過往的牆.
我們才能夠將福音交到人的手上.
今天跟大家看的這段經文.
其實我簡單的原因就是.
因為只有幾句.
是耶穌的很重要的榜樣.
在上面他還受著其他人的批評.
但這段聖經好像突然間沒有理會.
然後他形容耶穌是.
走遍各城各鄉.
頂展門第一件事是走出去.
耶穌不是在耶路撒冷.
設立了一間教會或會堂.
然後叫人來來來.
不是,耶穌是走遍各城各鄉.
是走出去.
福音是主動去尋找.
如果大家讀過一點神學的基礎.
你會知道今天我們得著救恩.
是因為上帝主動啟示自己.
不是你去尋找祂.
是上帝主動啟示自己.
讓你去得著這份救恩.
主耶穌在這裡是主動去尋找人回來.
頂展門.
所以我們今天不是等.
不是在這裡等.
我們心裡沒有目標.

$^{281}$而當你尋找到這個人回來之後.
其實有些很清晰的目標要做.
第一件事,教訓.
用上帝的話語去提醒.
我們今天的教會有沒有上帝的話.
這是我們的基礎.
你記住我們不是罵人.
我們不是說人.
我們不是批評人.
我們是提醒.
因為上帝有些東西.
在教會當中要我們生活出來.
最直接的,對於微信者最受惠的是什麼.
是行出愛,活出愛.
他來到的時候.
感受到的是愛.
還是被拒絕.
其實在教會裡面很需要的是什麼.
是一份欣賞的文化,弟兄姊妹.
我們每一個都是罪人.
我們會不會做錯事.
我們會.
但是我想頂展門學一件事.
不要放大錯誤.
而是我們學會在裡面去欣賞每一個人的特色.
小朋友是在欣賞和鼓勵裡面成長的.
甚少在責罵和責罰當中成長.
記不記得大家在小時候.
小朋友的時候.
由他開始坐在那裡學走路.
我想問大家,你是怎麼做的.
沒錯,是鼓勵.
你不會拿著藤條,很起來走路.
你還不走.
然後一跌下來,你想死嗎.
我打你,這麼差.
是不是這樣.
不是.
是他開始起來走路的時候.
哇,厲害.

$^{321}$快點,過來吧.
你是鼓勵他.
起來吧,再起來吧.
弟兄姊妹,今天其實我們每一個人.
在最原先的小朋友都沒有長大過.
我們都需要鼓勵.
我們都需要欣賞.
教會是很需要和外面很不同.
我們需要這種欣賞文化.
欣賞每個弟兄姊妹在裡面的付出.
欣賞弟兄姊妹在裡面的我們每一個人的盡力.
做得最後在你的標準裡面做得好不好是其次.
第二件事,我們要宣講福音.
馬福音,第一章第一節.
「這是上帝兒子耶穌基督的福音,這福音的開始」.
你今天的宣講福音不是講三福四律.
你是將耶穌整個的生平去介紹給人聽.
做什麼?.
耶穌怎樣去愛人.
耶穌怎樣去改變人.
耶穌怎樣去改變我.
整個福音是整個福音書.
這個不只是靠嘴巴說.
弟兄姊妹,你自己對於主耶穌基督的救恩.
你有多深的體會.
耶穌基督的救恩是不是真的令你流淚動容.
耶穌基督的拯救是不是真的你覺得是非常偉大.
第三件事,醫治釋放人.
來到教會的人,我們都是受傷的人.
我們被罪惡所受傷.
我們在世界上被有罪惡的人去令到我受傷.
來到的人,無論在肉體上,無論在心靈上.
都需要被釋放.
釋放人是需要學習的,弟兄姊妹.
是要有同理心的.
釋放人是需要站第一級的.
去感受對方的需要.
所以弟兄姊妹,第一件事就是我們有沒有走出去.
有沒有做了我們應該做的事.
第二件事是漢見.

$^{361}$耶穌漢見那些人困苦流離,之後憐憫.
你看不見,你不能憐憫.
弟兄姊妹,今天在你的生活裡.
在你生活的圈子群體裡,在你的鄰居當中.
你看見些什麼.
我只看見鄰居經常把垃圾放在門口.
我只看見鄰居的狗很吵.
弟兄姊妹,你看見些什麼.
為什麼鄰居的家裡經常都是老人和狗.
為什麼鄰居的家裡經常都是家族屋幣.
我們今天在我們的生活裡,我們看見什麼.
而跟著這件事,有沒有令你產生憐憫.
弟兄姊妹,福音工作是需要這份憐憫的心.
你都覺得他沒有需要,你何來想將福音給他呢.
第三件事,是懇求.
你會發覺耶穌最後一句.
莊家多,作公的人少,請你不要走得太遠.
請求莊家的主打發工人.
你今天不是單打獨鬥的,弟兄姊妹.
你今天是和上帝同工的.
你是和上帝同行的.
從來耶穌告訴你什麼.
我所擔的,是輕省的.
我的擔子是輕省的.
而耶穌也告訴你為什麼是輕省的.
因為耶穌擔了大部分.
你今天去傳福音給人,你今天跟人分享福音.
其實上帝在,耶穌在,不過你知不知道.
你有沒有邀請主耶穌和你一起來.
所以我們經常說,說到快要老掉牙的一句話.
就是請你去關心別人之前,好好禱告.
好好為對方禱告.
求上帝讓你去看見那個位置在哪裡.
弟兄姊妹,原來當我們今天說要做無牆教會的時候.
不是一句口號,是生命的轉化.
我整個生命裡面,我有沒有被改變.
我有沒有從神的眼光重新去看我的生活,看我的生命.
我能不能夠像這裡所說的,主動去看見一些人.
我跟著時間,我看見他們的需要困苦流離.
所以我去憐憫他們,到最後我為他們祈禱.

$^{401}$我求上帝打發我們能夠去到當中成為他們的幫助.
所以這是大家生命裡面,我們拆去平時生活的長線.
沒錯,香港人是很忙的.
我認同香港人是很忙的.
但我也聽過一個大學的老師,也是同工.
他說大學生一定很忙,究竟忙什麼.
所以很多時候,大家有沒有發覺.
每逢大學在他剛剛招生的時候,都會有一個節目.
全部學會都擺出來,之後怎樣?.
過來參加,為什麼?.
就是想他用時間忙在我們認為有價值的東西上.
今天為什麼團契越來越有人不想回?.
因為回到來的時候,我找不到我認為很有價值的東西.
我回到來的時候,如果我真的能夠去找到同路人.
我找到聽得懂自己說話的人.
我找到明白的人,我找到鼓勵支持的人.
我很相信你是會去的.
所以求神幫助我們,我們生命裡面都需要一些轉化.
所以我說今天這篇不複雜,聽姐妹.
建立無牆教會,需要弟兄姊妹走出的第一步.
我想介紹給大家一本書.
大家記不記得之前有一位叫陳淑娟牧師.
記得嗎?有一次她來跟我們分享.
她送了一本禮物給我們.
她那本書寫著叫做教會無牆的震撼轉化.
弟兄姊妹,這本書是非常容易看的.
裡面的字不多.
但我想今天用一點時間讀裡面的內容給大家聽.
陳淑娟牧師的教會是在昔日裡面.
這間教會差不多有三十多年的歷史.
當她的教會還沒有轉化的時候.
教會的聚會人數大概是三十多人.
到今天的時候,當時是2019年,是三百多人.
她這樣說,我們的教會就在這個地區.
奇怪的是,為何教會有著可以容納250人的禮堂.
過去三十多年的崇拜人數卻停滯不前.
為何教會吸納不了附近那麼多新搬來的中產人士.
為何身邊貧窮的人多的是.
他們卻沒有從我們這裡認識福音.
儘管教會的禮堂能夠容納多人.

$^{441}$但直到2014年.
恆常出席教會崇拜就只有寥寥五六十人.
主日的禮堂座位被空置不在話下.
星期一至六,全教會更是被限制著.
成為人跡罕至的地方.
可以說,當時的宣道會會友主要是主日教友.
他們通常在崇拜那一天才回來.
崇拜過後,教會各自各出外午膳.
平日,堂弟的教會只有一位幹事和牧者.
先有信徒和其他人到訪.
究竟是信徒無心事奉,不關心教會.
還是教會沒有好好栽培他們.
給予他們事奉的機會.
以往教會每年都會舉辦一次大型聚會.
在報道活動舉辦的幾個月.
我們過往不停,或忙於開會準備.
或為了同心祈禱等.
在那幾個月,教會彷彿充滿朝氣.
更好像有著共同奮鬥的目標.
但結果是怎樣呢?.
這些活動猶如年初二的煙花會宴一樣.
瞬間即逝.
當一切靜止下來後.
我們的心靈好像更凝空虛.
這是因為信主的人數不多.
就算有街坊回到教會.
很快他們又流失了.
或許那時我們跟進工作不足.
又或是那時我們仍然是中產教會的模樣.
新來者自覺格格不入.
但更重要的是.
當時我們確實不太懂得接納背景特殊的陌生者.
可惜更悲哀的是.
當我們一班弟兄姊妹,包括我自己在內.
彷彿只活在自己的世界中.
自我感覺良好便是.
或許用自我興奮.
來形容當時的我們.
倒視貼切.
心信上主到了,他不容許我們這樣下去.

$^{481}$在一次主人的崇拜中.
他親自跟我說話.
為何我的殿如此荒涼.
我望著台下可以容納250人的禮堂.
但當中只有不到30人的回眾.
這句話令我很責心.
原來你現在擁有的東西.
下一秒神便可以收回.
如果做教會事工.
只會拿著計算機計算周詳材做.
其實是甚麼都不必做了.
其實我想介紹這本書給弟兄姊妹的原因.
我不是要大家copy and paste.
因為我們的地點,我們的領袖不同.
但弟兄姊妹,我想大家一起在裡面.
我們一起重新思想上帝對他的提醒.
究竟今天上帝讓我們在這裡的時候.
上帝想我們做甚麼.
剛才我看這段的時候.
我發現其實很多一般教會的影子.
我們都會繼續努力地舉行報道會.
定時做一些事.
但果效卻很少.
但為何陳牧師他們的教會會不一樣呢.
我鼓勵大家,如果你有多餘的錢.
這本書不貴,88元.
你可以在基督教書局買一本.
如果你說我也想看一看.
教會有一本.
但我歡迎大家借回去看.
大家要給我,因為有陳牧師的簽名.
所以弟兄姊妹,我很盼望.
其實我們做神的工作,我們需要一份勇氣.
需要一份完全依靠上帝的勇氣.
是,沒錯.
戒恩流到今天,我們都不知道未來會怎樣.
不知道.
但弟兄姊妹,我們做到最後一分一秒在這裡,好嗎.
如果上帝開路給我們.
我們為了上帝祂奇妙的.

$^{521}$我們繼續在這裡走下去.
如果上帝不讓我們在這裡做工作.
但我們的工作是不應該停下來的.
我們的工作是要繼續的.
所以我先給大家看漢見,好嗎.
知不知道在錦繡外面.
遲些有二千多戶的公屋.
簡約公屋.
知不知道在我們附近敦裕學校裡面.
有很多很多來自國內和南亞裔的朋友.
知不知道在我們這個區裡面.
我們其實有很多人的長者.
他們在當中.
求神幫助我們,好嗎.
希望我們能夠一起去努力.
去為上帝去作供.
我相信上帝不想我們成為一個只是閒置的教會.
我相信上帝想讓我們成為無牆的教會.
\newpage



\section{路加福音 7:36-50}
\label{sec:lzufA9tLqUc}
\textbf{2023-06-11 午堂《愛和謙卑的事奉》路加福音 7\_36-50 講員:陳麗蟬牧師}
\newline
\newline
連結: \href{https://youtube.com/watch?v=lzufA9tLqUc}{\texttt{https://youtube.com/watch?v=lzufA9tLqUc}} ~~~~ 語音日期: 2023-06-13
\newline
\newline
\hyperref[sec:A_K04hyrSP8]{\small{< < < PREV SERMON < < <}}
~
\hyperref[sec:index_chronic]{\small{[返順時目]}}
~
\hyperref[sec:index_scriptual]{\small{[返順卷目]}}
~
\hyperref[sec:eaFKBNr6R3g]{\small{> > > NEXT SERMON > > >}}
\newline
\newline
路加福音 7:36-50
\newline
\begin{longtable}{cl}
\hline
\hline
章節 & 經文 (和合本修訂版)\\
\hline
7:36 & \begin{tabularx}{0.7\textwidth}{X} 有一個法利賽人請耶穌和他吃飯，耶穌就到那法利賽人家裡去坐席。 \end{tabularx} \\ \\ \relax
7:37 & \begin{tabularx}{0.7\textwidth}{X} 那城裡有一個女人，是個罪人，知道耶穌在法利賽人家裡坐席，就拿著盛滿香膏的玉瓶， \end{tabularx} \\ \\ \relax
7:38 & \begin{tabularx}{0.7\textwidth}{X} 站在耶穌背後，挨著他的腳哭，眼淚滴濕了耶穌的腳，就用自己的頭髮擦乾，又用嘴連連親他的腳，把香膏抹上。 \end{tabularx} \\ \\ \relax
7:39 & \begin{tabularx}{0.7\textwidth}{X} 請耶穌的法利賽人看見這事，心裡說：「這人若是先知，一定知道摸他的是誰，是個怎樣的女人；她是個罪人哪！」 \end{tabularx} \\ \\ \relax
7:40 & \begin{tabularx}{0.7\textwidth}{X} 耶穌回應他說：「西門，我有話要對你說。」西門說：「老師，請說。」 \end{tabularx} \\ \\ \relax
7:41 & \begin{tabularx}{0.7\textwidth}{X} 耶穌說：「有兩個人欠了某一個債主的錢，一個欠五百個銀幣，一個欠五十個銀幣。 \end{tabularx} \\ \\ \relax
7:42 & \begin{tabularx}{0.7\textwidth}{X} 因為他們無力償還，債主就開恩赦免了他們兩個人的債。那麼，這兩個人哪一個更愛他呢？」 \end{tabularx} \\ \\ \relax
7:43 & \begin{tabularx}{0.7\textwidth}{X} 西門回答：「我想是那多得赦免的人。」耶穌對他說：「你的判斷不錯。」 \end{tabularx} \\ \\ \relax
7:44 & \begin{tabularx}{0.7\textwidth}{X} 於是他轉過來向著那女人，對西門說：「你看見這女人嗎？我進了你的家，你沒有給我水洗腳，但這女人用眼淚滴濕了我的腳，又用頭髮擦乾。 \end{tabularx} \\ \\ \relax
7:45 & \begin{tabularx}{0.7\textwidth}{X} 你沒有親我，但這女人從我進來就不住地親我的腳。 \end{tabularx} \\ \\ \relax
7:46 & \begin{tabularx}{0.7\textwidth}{X} 你沒有用油抹我的頭，但這女人用香膏抹我的腳。 \end{tabularx} \\ \\ \relax
7:47 & \begin{tabularx}{0.7\textwidth}{X} 所以我告訴你，她許多的罪都赦免了，因為她愛的多；而那少得赦免的，愛的就少。」 \end{tabularx} \\ \\ \relax
7:48 & \begin{tabularx}{0.7\textwidth}{X} 於是耶穌對那女人說：「你的罪都赦免了。」 \end{tabularx} \\ \\ \relax
7:49 & \begin{tabularx}{0.7\textwidth}{X} 同席的人心裡說：「這是甚麼人，竟赦免人的罪呢？」 \end{tabularx} \\ \\ \relax
7:50 & \begin{tabularx}{0.7\textwidth}{X} 耶穌對那女人說：「你的信救了你，平安地回去吧！」 \end{tabularx} \\ \\
[1ex]
\hline
\hline
\end{longtable}
$^{1}$弟兄姊妹平安.
在這裡去敬拜上帝.
原來不是必然的.
剛才子堅為我們讀出那段經文.
大家都可能聽過很多次了.
就說到法利賽人西門.
請耶穌回到他家吃飯.
通常來說你請人吃飯.
你當然和那個人的關係很好.
如果關係不好你一定不會請他吃飯.
還要請他去你家吃飯.
關係就會更加緊密.
如果不是的話.
普通去大家樂.
去哪裡吃飯就算了.
但當我們想清楚.
法利賽人和耶穌基督的關係.
一般般.
我們知道.
法利賽人不太抵得耶穌的門徒.
因為法利賽人覺得自己很守律法.
自己很規矩也很聖潔.
不過他們覺得耶穌教出來的門徒.
很不認真.
經常不守律法.
譬如他們吃飯也不洗手.
有時在安息日走過稻田.
肚子餓就去摘和水來吃.
他們覺得自己很不守規矩.
就算耶穌基督自己.
也不喜歡.
因為耶穌說明安息日.
很多事情不能做.
耶穌竟然公然在安息日.
去醫好一些病人.
所以法利賽人一直覺得.
耶穌很一般般.
那麼耶穌又怎樣看法利賽人呢?.
耶穌在聖經裡面一直說.
說他們是假冒偽善的人.

$^{41}$這些人經常自以為是.
自己以為自己做得比別人好.
永遠看人都是不好的.
所以基本上兩者的關係.
一般般.
但究竟在這個一般般的關係裡面.
為何法利賽人西門.
會請耶穌吃飯呢?.
因為當時請人去他們家裡吃飯.
是一件很大的事.
他們要先前有很多籌備.
就好像我們今天.
家人結婚.
很多籌備.
是很大的事.
他們在吃飯裡面.
看他們的布置.
他們在這個側臥.
他們叫側臥不是坐在那裡.
側臥是一些長的椅子上.
前面就放了一些食物.
那些人就會靠在那裡.
用左手抬著頭.
右手就在外面拿食物.
在我們現代的人.
習慣了坐在那裡吃飯.
很難想像.
這樣靠在那裡.
拿食物是很舒服的.
我都不明白.
不過以前的人就是這樣.
還有當時的習俗.
當他們請人吃飯的時候.
他們當然會請一些很出名的人來吃飯.
當時就請了耶穌.
耶穌在那時候.
都是一個叫做很出名的拉比.
很出名的老師.
所以通常他們就會大開中門.
他們通常坐在中間的那個.

$^{81}$就是主人西門.
對著門口.
看著邊邊的人進來.
而這個客人就坐在旁邊.
剛才那個事跡.
我們就看到.
有個女人靠近了耶穌那裡.
他們就在這樣的狀態裡面來吃飯.
這個故事一開始.
就說到耶穌在西門家裡吃飯.
大家客客氣氣地吃東西.
在第37節裡面這樣說.
在那裡有一個女人是個罪人.
很普通的一句句子.
不過原來這句子前面有兩個字.
在原本的文字裡面是一個字.
就說Honor.
什麼時候呢?.
當一個字出現.
譬如有人跟你說.
喂喂喂 看一看.
你們沒有人看的吧.
喂喂喂 看一看.
就是告訴你.
那裡有些很特別的東西.
給大家看.
這句經文就是這樣.
Honor.
那裡有一個女人是個罪人.
就說本來在客客氣氣地吃飯的時候.
就有一個有罪的女人.
就進了西門的家.
她就拿著一瓶香膏.
要送給耶穌.
這個Honor這個字.
其實是在表達西門.
一看到這個有罪的女人進來的時候.
就腎上腺上升.
非常的緊張.
非常的不開心.

$^{121}$那她為什麼這麼不開心呢?.
因為她覺得.
她自己是法律裁人來的.
她是嚴守律法來的.
很聖潔的.
但這個不正經的女人.
她竟然進了我這個家.
還有你看看.
她靠近了耶穌那裡.
她在想.
耶穌究竟是不是先知.
枉你這麼出名.
你知不知道靠近你那個女人裡面.
是個很不正經的女人.
全個城都知道的.
在《利未記》裡面教導過的.
只要有一些不乾淨的東西.
和你接觸過.
或者你摸得到的時候.
你就會成為不聖潔的了.
你究竟知不知道.
這個不正經的女人.
她正在接近你.
耶穌基督是那位知曉人心靈的神.
祂知道西門在想什麼.
又或者祂和西門那個.
很緊張激到臉都紅了的裡面.
祂知道西門在想什麼.
祂就講了一個故事給西門聽.
有兩個債仔.
一個欠債主500個銀幣.
一個欠了50個銀幣.
當時在羅馬政府裡面.
他們每一個人的工資.
就是一個銀幣.
所以那個大債仔.
就欠了人家兩年多的債.
小債仔都要做一個多月才可以還債.
不過當時這個社會裡面都是很窮困的.
和沒有任何的福利.

$^{161}$不是像現在這樣又有這個綜援.
或者又有這個傷殘金.
什麼都沒有的.
人民就是餐餬餐食.
所以餐餬餐食的時候.
你根本就做多少都好.
你沒有力氣去償還給債主的.
但不知為何有一天.
這個債主就大發熱心.
他說看你們兩個都沒辦法還得了.
不用你還了.
所以那兩個債仔聽到的時候.
全身都鬆了.
接著耶穌就問西門.
他說你覺得兩個債仔.
哪一個感受到債主的寬容和仁愛呢?.
哪一個呢?.
不用問你們都知道了.
西門都知道.
當然差很多那個.
但其實無論你差多差少都好.
當有人他做了一些事不是很好.
或者欠了人債.
但如果不再被追究不再被責備的時候.
心裡面是會好感動好感恩的.
上星期五我就去俱樂部吃飯.
完了師班之後去俱樂部吃飯.
我就看到隔幾張桌子.
就有一位老太太好老友記的.
我就走過去跟她聊幾句.
聊完之後就回到台上繼續吃飯.
那天晚上我確診了.
接著到這個星期日早上.
我知道這個老太太又發燒又咳嗽.
我心知不妙.
星期一家人告訴我老太太又確診了.
心裡面好不舒服.
死啦 當然是我惹她.
於是就打電話問候一下人家家人.
說一聲不好意思.

$^{201}$家人家人給我的回應好感恩.
令我很釋懷.
那家人怎麼說呢.
他說現在這麼多人確診.
差不多全港的人都要每人中一次.
不是今天中就是明天中.
所以沒什麼大不了.
老太太今天中了明天就不用中了.
我很感恩 家人沒有怪我.
最壞是你和老太太坐在一起.
還要兩個人坐一張椅子.
我覺得不關我的事.
老太太說牧師不可以站著跟她聊天.
她借半張椅子給我坐.
我覺得心裡面好感動好感動.
所以我就WhatsApp給她家人.
疫情穩定之後 個個都康復之後.
我一定請他們全家吃飯.
心靈覺得釋放了.
沒有被人怪責的裡面.
你就會有一種感覺.
你很想去做一些事回饋他們.
耶穌為什麼要說這個故事呢.
祂的目的是哪裡呢.
祂要告訴人知道.
祂就是這個債主.
西門和這個女人都是債主.
耶穌基督是債主.
你和我都是債主.
在《羅馬書》裡面說了.
世人都犯了罪.
虧欠了上帝的榮耀.
所以在這個世界裡面.
我們所有人都是虧欠了上帝.
所以無論犯罪多或少.
我們沒有人能夠償還給上帝.
每一次聖餐的時候.
我們都要記得.
我們的生命是主耶穌基督用祂的寶血.
用祂的生命.

$^{241}$將我們從罪惡的權勢裡面.
重價地買贖回來.
而你和我呢.
我們靠著耶穌基督的寶血.
我們的罪才能夠得到赦免.
這段經文就要告訴我們.
我們要怎樣去回應.
耶穌基督救贖的恩典.
我們繼續看看這個女人.
為甚麼她一見到耶穌的時候.
眼淚就不斷地流下.
很激動.
我們看看路加福音前幾節的經文.
說了一些答案給我們聽.
在第29至30節.
「眾百姓和瑞麗已受過約翰的洗.
聽見者說就以上帝為義」.
當時施洗約翰.
他在不同的地方講悔改之禮.
「天國近了,你們要悔改」.
所以經文說.
有些百姓,有些瑞麗就代表了罪人.
他們都聽到.
聽到天國近了,你們要悔改.
他們遭到悔改.
並且接受了約翰的洗禮.
這個女人可能就是其中的一個.
但接著這段經文說.
「但法理賽人和律法師沒有受過約翰的洗」.
他們硬心不肯接受約翰的洗.
他們覺得這不是上帝的心意.
竟廢棄了上帝為他們所定的旨意.
他們覺得他們是聖潔的.
是守守所有的律法.
這個女人在施洗約翰的裡面.
她聽到,她知罪,悔罪.
並且洗禮了.
所以當她知道.
耶穌就是那位救世人的彌賽亞的時候.
又知道他在西門的家請吃飯的時候.

$^{281}$他就帶著一瓶香膏來送給耶穌.
但一到西門的家.
她見到耶穌的時候.
她很感動.
眼淚就很快地流出來.
就好像有時我們來到教會.
很自然地進到禮堂.
有時我們心裡面就會覺得.
有一份激情不斷流淚.
或者聖餐.
或者有時聽到一些詩歌的時候.
很自然地不知為何.
是聖靈的感動.
不斷流淚,控制不了的.
這個感恩,感動,感激的淚水.
就在這個女人的裡面不斷地流.
她怎知道不小心滴濕了耶穌的腳.
她覺得真的很不好意思.
很不好意思.
於是情急之下.
她就鬆開了自己的頭髮.
就去抹著耶穌基督的腳.
就親一下耶穌的腳.
又再將她帶來送給耶穌的香膏.
來抹耶穌的腳.
其實在當時的社會來說.
這些行為不是很合大體的.
一般娘家婦女.
在大庭廣眾裡面.
她絕對不會鬆開頭髮的.
但這個女人在情急之下.
就鬆開頭髮.
去抹耶穌基督腳裡面的淚水.
她這樣做.
其實就是表示她對耶穌的愛.
她很想服侍耶穌.
很想快點來抹乾淨耶穌基督的惡.
她知道有人批評她.
但她不介意別人的眼光.
她也不介意別人的說話.

$^{321}$不介意別人的批評.
在事奉當中.
有時我們都會聽到別人不同的一些說話.
一些批評.
一些令人很深入沉下去的說話.
一些比較.
以致我們沒什麼力氣再去事奉.
但在這樣的時間裡面.
我們就需要去反省一下.
反省一下.
我這個事奉的動機.
是不是對準耶穌基督.
或者在事奉的過程裡面.
是不是真的有些事情我自己做得不是很好.
需要去改善一下呢?.
所以經過反省之後.
我們就要去有一個立志.
無論怎樣.
我們的事奉動機是對準神的.
是要這個的.
用一個的愛.
來回應耶穌基督救贖大愛的動機的時候.
我們就要立志.
忠心事奉.
堅持到底.
直至見主面.
我之前在一間基層教會裡事奉了八年多了.
當中弟兄姊妹的學識都不是很高的.
有中學畢業的那些已經是很厲害的了.
當時有一位姊妹.
她就是在執事會裡做總務的.
做了兩年之後.
見到教會財政的位置沒有人做.
那次的執事會.
她就想轉一下這個崗位.
做這個財政.
因為她對這個數字都很敏感的.
特別是那些婆仔數她算得很清楚的.
她上任之後.
她就請教一下之前的財務.

$^{361}$哇 真是慘啊.
她被人罵得體無完膚.
我甚至形容一句廣東話.
很貼切的.
被人罵得狗血淋頭.
那人怎麼罵她呢.
你不要教錯了.
你小學還沒畢業.
讀過那兩科的會計.
有沒有做證書.
學什麼人去做財政.
你有什麼資格去做.
教會的數目讓你這樣做下去.
一定搞到錯了.
亂了.
有膽量去做財政嗎.
告訴你.
基層的人說話就是這樣放的.
罵人的那位姐妹.
她在夜校完成了她的中學.
她又讀了兩個會計的課程.
有證書的.
她就覺得.
哇 我都做了這麼多年財務.
現在沒空就不做.
你現在接手上來.
你只是小學還沒畢業.
又沒有證書的.
你做什麼.
你有什麼資格去做.
新上任的這位姐妹.
她真的是小學還沒畢業.
她因為家裡很窮.
所以很早就出來做車醫.
但到18歲的時候.
她想想車醫收入不穩定.
她就想想去醫院裡面做清潔車.
收入會多一些的.
但上帝給她那段時間又認識了一個牙醫姐妹.
那位牙醫姐妹又挺勇敢.

$^{401}$不如就請她做牙醫助護.
那位姐妹一直做到她退休.
她認識很多牙科的藥物.
或者器材的英文.
不過其他的她就真的認真一般般.
不過她很深的經歷到.
經歷到上帝在她的工作.
在她的生命裡面不斷的提升.
所以在教會裡面很熱心的去事奉上帝.
我看著她的.
最初做組長.
接著就在執事會裡面做總務.
接著就做財務.
最後她就做了執事會副主席.
在當天她做財政的時候.
被人真的淋了一頭冷水.
但她經過這個反思.
她反思到她想做這個位置.
不是想炫耀自己.
而是她真的很想回應上帝的愛.
她覺得在數字上高.
她都可以掌握得到.
她也確知那次是上帝呼召她去做這個崗位.
她堅持到底.
做完一屆又一屆.
在我離開教會的時候.
整條數都是整整齊齊.
出於愛的事奉.
愛上帝的事奉.
不會因著任何人的冷水去退縮.
所以頂子們.
在我們事奉的裡面.
一定堅持.
以愛上帝.
因為真的給了一個很大很大.
救贖的恩典的愛給過我們.
第二方面我們看看.
你的事奉態度.
我們看看西門.
怎樣對耶穌.

$^{441}$在44節.
他就說.
於是他轉過頭.
向著那個女人.
耶穌對著那個女人.
但他不是對那個女人說話.
是對西門說話.
很特別.
他說西門.
你看這個女人.
我進了你家.
你沒有用水幫我洗腳.
但這個女人用眼淚滴濕了我的腳.
又用頭髮去擦出.
以前也是.
進到別人家裡.
一定有人幫他洗腳.
耶穌只是說.
西門沒有親自為他洗腳.
也沒有親自出來接待他.
一定有奴隸幫耶穌洗腳.
因為那時候他們穿的草鞋.
又沙漠地.
骯髒的.
所以每進人家裡.
有錢的人家一定有奴隸.
幫他們去洗腳.
但西門沒有.
不是親自幫他洗.
也沒有親自接待他.
就說來了.
洗腳.
幫他洗腳.
有一次.
有位知體請我和他家人一起喝茶.
去到的時候.
這位知體就說.
跟我介紹.
這些是我家人.
十個八個家人.

$^{481}$這位陳牧師和我們一起喝茶.
坐下了.
陳牧師你自己點餐.
你喜歡吃什麼就自己點.
他和家人個個都在那裡.
嘟嘟嘟嘟嘟.
我沒理由喜歡.
我想吃雞扎.
我沒理由自己點.
但那次吃飯真的很淡.
很淡.
走過來吃飯.
不是和人家聊天.
但耶穌可能去到就是這樣.
那時候.
阿A會幫你洗腳.
但你們應該禮貌.
你走去.
耶穌歡迎你.
來和我們一起吃飯.
沒有.
還有一個.
然後.
他說.
西門你不要來和我親吻.
但這個女人.
進來的時候就不住的.
親吻我的腳.
我們就覺得很不習慣.
但原來這個.
在當時的禮儀.
在現在中東裡面都有.
是吧.
一個很普通的禮儀.
見到的時候就擁抱一下.
又左親吻.
又親吻.
甚至有機會離去親嘴.
一個很普通的禮儀.
但你說到親吻腳的時候.

$^{521}$是一個特別的禮儀.
是一個跪拜的禮.
是表示極度的謙卑.
跪拜是統治者.
通常都是去做.
向一些國王,貴族.
才去做這個儀式.
西門沒有做.
什麼都沒有做.
但相反.
那個女人.
她有做這件事.
但在表達.
這個女人.
非常地愛耶穌.
多謝祂.
因為祂就是那位.
拯救世人的彌賽亞.
還有46節就講到.
西門就說.
我來到你家.
你沒有用油來抹我的頭.
但這個女人.
就帶著一瓶香膏.
來抹我的腳.
這個有錢的人家.
在當時中東地帶.
當那些客人.
來到家裡身體汗水.
當然是全身臭.
有些家人.
就會拿一些香膏出來.
好像香水一樣.
灑在客人的頭上.
讓他精神一點.
爽一點.
在疫情裡面.
我們都經歷到了.
當去到別人家裡的時候.
通常會怎樣.

$^{561}$站在那裡.
噴一下.
消毒一下.
但這個禮儀.
灑一些香水.
在客人的頭上.
可以做或不做.
不過通常有錢人家都會做.
不過耶穌就說.
這次西門沒有幫他做.
不過倒過來.
是女人做了.
你抹她的腳.
耶穌在這裡.
就在說西門.
那些最基本的禮儀.
你都沒有做到.
你究竟請她回來吃飯.
為了什麼呢.
最基本的禮儀.
你都沒有幫耶穌做.
你請她回來吃飯.
是不是想炫耀一下你自己.
你這次可以請到.
這個很出名的.
當時很出名的.
拉比老師.
來你家裡吃飯呢.
你是不是想讓人知道.
你家裡有多麼有錢.
你請了這麼出名的人來的時候.
可以去炫耀一下自己呢.
弟兄姊妹.
當我們事奉的時候.
我們都要想一想.
我們究竟為了什麼呢.
還是想表現一下.
我自己有多麼厲害.
我自己的恩典有多麼多呢.
因為真正的服事.

$^{601}$是要帶著一個謙卑的心.
才是蒙上帝的月亮.
我們看一看.
當大衛迎接約桂.
進入大衛城的時候.
他很興奮.
他和周圍的人一起去唱歌跳舞.
在《撒母耳記》裡面說到.
大衛當時是極力地跳舞.
很興奮地跳舞.
但結果帶來.
是怎樣呢.
他的太太米甲燒他.
米甲燒他.
嘩 你真的有沒有搞錯.
作為一個皇帝.
穿著這些跳舞服.
這些跳舞服.
是奴隸或者是低下階層的人去穿的.
你作為以色列第二任的皇帝.
你自己穿著這樣的服裝.
還有什麼行為.
大跳大唱.
完全不合乎一個君王的體統.
但其實呢.
大衛是很感受到上帝的恩典.
他自己本來是一個小鴻鴻.
但現在被提升成為皇帝.
一切都是上帝的恩典.
他雖然是皇帝.
他不介意自己的身份.
放下了這個身段.
放下皇帝的身段.
他就和周圍跳舞的人.
盡情的歡呼.
盡情的唱歌.
耶穌進了西門教.
西門沒有幫他.
用油膏來刮他的頭.
我是有錢人來的.

$^{641}$沒有令耶穌讓他清爽一下.
但是呢.
女人就用貴重的香膏.
亂撥他的腳.
在事奉當中.
我有沒有將最珍貴的東西.
給過耶穌呢.
我最珍貴的東西.
我生命裡面最珍貴的東西.
給了甚麼呢.
給了我的家庭.
給了我的事業.
給了我去旅行去吃東西.
還是給了耶穌呢.
有時呢.
當我們去邀請人頂子女去事奉的時候.
他們會這樣回答我們.
教會那麼多人.
為甚麼要找我來事奉呢.
我通常都會回答.
你不用理會其他人.
最重要的.
就是你願不願意.
去回應耶穌基督給你救贖的愛.
你自己願不願意.
你不需要理會其他人.
還有這個西門.
一見到那個女人.
你認為有罪的女人.
進了她的家.
就很激動.
嘩 這個有罪的女人.
進了我這個那麼聖潔的家.
哎呀 慘了慘了.
會不聖潔的了.
她在批評人.
她在批評人有罪.
但是她沒有看見.
她自己眼裡面那個動量.
她在服侍耶穌的事情裡面.

$^{681}$她是多麼冷淡.
多麼無禮.
她在批評人.
耶穌基督是唯一一個.
在這個世界上高的大寨主.
只有祂才有權.
去做判斷.
做這個定案.
來去說出.
誰真是不好.
第一 教會裡面.
我們彼此的服侍.
我們真是很清楚知道我們的動機.
並且全一個謙卑.
捨己的態度.
我們來去服侍.
有人就說.
這個生命就好像一條河流一樣.
如果我們的生命.
只是不斷有水進去.
但沒有水出來的時候.
這個生命就像一條死水.
如果我們的生命.
我們在教會或一個群體裡面.
我只是不斷要人去服侍我.
但我不肯去服侍其他人的時候.
這個生命就會慢慢的萎縮.
甚至會停頓.
所以 弟兄姊妹.
我們的生命.
要為主所使用.
以致我們能夠成為一個流通本子.
榮臣益人.
我們來唱一首詩歌.
紅色的字.
「耶穌你拋棄萬有.
施恩愛將我拯救」.
這個就是我們事奉的動機.
回應耶穌基督的大愛.
我們要一生的.

$^{721}$遵從耶穌基督的教導.
事奉祂.
我們一起唱吧.
「耶穌你拋棄萬有」.
我要上耶路撒冷的.
遵從耶穌.
我要上耶路撒冷的.
遵從施福.
耶穌你拋棄萬有.
施恩愛將我拯救.
我要一生事奉你」.
我們一起有安靜的時間.
要在十字架上高.
為我們的罪孽來捨身.
因著耶穌基督的拯救.
我們的生命可以從.
罪惡的權勢裡面.
被釋放出來.
我們要努力的來事奉祂.
將祂的這份大愛.
傳揚出去.
以致有更多的人.
得著釋放的生命.
得著被捨的生命.
我們一起再唱.
我要上耶路撒冷的.
遵從耶穌.
我要上耶路撒冷的.
遵從施福.
耶穌基督起萬有.
施恩愛將我拯救.
我要一生事奉你.
我們同心禱告.
賜福給我們的生命.
我們每一個人都願意.
將從你而來的恩賜.
恩典回饋你的大愛.
求你使用.
使用我們的生命.
成為流通的管子.

$^{761}$榮耀主你的聖命.
禱告,感恩,祈求.
奉主基督的聖命.
阿們.
OK.
\newpage



\section{提摩太後書 3:16-17}
\label{sec:eaFKBNr6R3g}
\textbf{2023-06-18 加恩樓《有故事的人》提摩太後書 3\_16-17 講員:吳景霖傳道}
\newline
\newline
連結: \href{https://youtube.com/watch?v=eaFKBNr6R3g}{\texttt{https://youtube.com/watch?v=eaFKBNr6R3g}} ~~~~ 語音日期: 2023-06-20
\newline
\newline
\hyperref[sec:lzufA9tLqUc]{\small{< < < PREV SERMON < < <}}
~
\hyperref[sec:index_chronic]{\small{[返順時目]}}
~
\hyperref[sec:index_scriptual]{\small{[返順卷目]}}
~
\hyperref[sec:cNC4ab19lyE]{\small{> > > NEXT SERMON > > >}}
\newline
\newline
提摩太後書 3:16-17
\newline
\begin{longtable}{cl}
\hline
\hline
章節 & 經文 (和合本修訂版)\\
\hline
3:16 & \begin{tabularx}{0.7\textwidth}{X} 聖經都是神所默示的，於教訓、督責、使人歸正、教導人學義都是有益的， \end{tabularx} \\ \\ \relax
3:17 & \begin{tabularx}{0.7\textwidth}{X} 叫屬神的人得以完全，預備行各樣的善事。 \end{tabularx} \\ \\
[1ex]
\hline
\hline
\end{longtable}
$^{1}$各位弟兄姊妹早晨,主對面平安.
最後一次和大家講道,很不捨得大家的,大家不用不捨得我,大家又是很好的牧者.
不過我離開一個崗位,我不是不做牧者,我永遠都是大家的牧者.
如果你們有需要找我聊聊天,隨時都歡迎的.
今天和大家講道,有些人都會問,最後一次講什麼訊息?我又不懂得講什麼驚天動地的東西.
我只是按照聖經講,或者講一些我認為有價值的主題.
今天的講題是講有故事的人,首先什麼是故事?故事是什麼來的呢?.
有時候我們會想,故事就是老作那些,童話故事,或者是一些創作故事.
其實除了童話故事之外,電影都是一些故事來的.
有時候有些遊戲都有背景,有故事.
甚至其實在我們生命中真實發生過的事情,都是一些故事來的.
只要是一些過去發生過的事情,都是一些故事.
所以今天我會和大家講很多不同的故事.
第一個就是關於我的故事.
我小時候,其實我是跟著我祖母去成長.
有些人覺得看起來好像很健康成長,其實不是的.
我是一個單親家庭成長,我爸爸媽媽在我小時候已經分開了.
現在有些人都問我,你移民不用理家裡嗎?.
我爸爸和媽媽各自都有家庭,所以我是不用理他們,他們各自有家庭.
我有很多弟弟妹妹,但是同父同母,只有一個弟弟.
我跟著我祖母長大,所以小時候沒有什麼人去理.
他開始照顧我的時候,已經是六十歲.
那時候小時候,他就帶著我去哪裡都跟著他.
去街市也是跟著他,去晨運也是跟著他,一直黏著他.
有一次,說說我的背景,我們住在一個唐樓.
唐樓就是十層高,沒有樓梯的那些,叫唐樓.
我家就在五樓,所以每次回家都要走五層樓梯,很累的.
我媽媽有一次帶我去街市買菜.
去完街市買菜,就抽著菜走五層樓梯回家.
回去的時候,她看到我手上多了一樣東西,多了一個算盤.
知不知道算盤是什麼?知道的吧?我怕年輕人不知道.
但是其實我們的算盤都不是用來算的,當滑板而已,小時候.
就是多了一個算盤,為什麼會有一個算盤呢?.
想想,剛才經過一間雜貨舖,可能就是我順手在那裡拿了.
拿了就回家了,一般人,小孩拿了東西回家.
又沒有人發現,沒有人追出來,那就是沒事了,是不是?.
那你會怎麼處理?有些可能當沒事發生.
但是我奶奶是一個很誠實的人,她就拿了一個算盤下樓給人家.
然後又要再走五層樓梯上來,很辛苦的.
這些我當然不記得了,我那時候很小.

$^{41}$但是這些故事是我奶奶她的複述給我聽的.
她讓我知道過去曾經發生過這樣的事.
那有什麼用?她告訴我,其實是讓我去到知道.
她是一個怎樣的人,她是怎樣去生活的.
這個成為了我人生的一個路標,我學會了一些東西.
然後在我生命裡面是有一些影響的.
有一次我和我女朋友逛書室,和女朋友逛街.
你的集中力當然是在女朋友身上,是不是?.
於是選到一整手東西,一直顧著在那裡聊天的時候.
聊著聊著就走了出去,走了幾條街.
然後才發覺,咦,一整手都是東西,這是什麼?.
剛才拿了沒有給錢,那沒有給錢又怎樣?.
又沒有人追出來,沒有人發現,那就是沒事了.
而且不是很多,不是很貴的東西,那些書籤,那些小東西.
但是我的第一個反應就是,回頭去給錢.
去到跟別人說,對不起,我沒有給錢,忘記給錢走過去.
然後別人又跟我說,對不起,我忘記收你錢.
大家很幽默地去處理了這件事.
其實這些我會發覺是一個人生裡面的一些路標.
當然不會說我去到我的時代,突然間我又拿了一個算盤.
你不會說有些東西是完全重複了.
但是當有類似的事情發生的時候.
這些故事就在你生命裡面產生了一些果效.
就讓我跟隨了我媽媽的腳蹤.
我成為了一個誠實的人,我會去到給錢.
所以說其實故事是可以承載很多東西.
承載的訊息是多於你去找口說.
我教小朋友的時候,你跟他說你要做一個誠實的人.
你要做一個乖的小朋友.
小朋友哪知道你說什麼,什麼叫乖的小朋友,什麼叫誠實.
你要透過一個故事去到承載那個訊息.
你跟他說誠實,他不知道是什麼.
你跟他說狼來了,他就知道了.
你跟他說狼來了,他就知道原來這樣就是不誠實.
會有什麼後果.
所以故事其實是一個平鋪直敘去到說一些東西.
用故事承載你要的訊息.
你去到讓那個故事入心.
所以很鼓勵大家你要多點去分享故事.
分享故事不是說只是大人要跟小朋友分享.

$^{81}$其實彼此之間我們是可以去分享我們生命裡面的故事.
事實上耶穌都經常用故事.
他教的門徒都已經是大人了.
一樣可以用故事去承載一個訊息去告訴他.
所以大家多點去到跟你身邊的人分享故事.
不過我也要提一提父母.
分享你的故事不是藉故去教訓他們.
不是藉故說我好像你這麼大的時候.
我已經出去做事幫補家計了.
你躺在沙發那裡.
不是乘機用你的生命的經歷去教訓他.
而是去分享一些曾經發生過在你身上的事情.
好讓他能夠明白一些道理.
我又繼續講故事.
這次就講關於我兒子.
我變成大人了.
我生了個兒子.
那時候他K3.
我就帶著他上學.
我經常都遲到.
帶小朋友上學要很早的.
幼稚園.
趕著上學拖著他走的時候.
看到前面街市那裡有個婆婆.
她在垃圾堆前面彎腰在做這個動作.
在做什麼呢?你猜她在做什麼?.
她在垃圾堆裡撿一些沒那麼爛的水果來吃.
是一些黃皮.
摘一粒,捽了一粒肉,然後扔掉.
那堆真的很噁心.
有些豬肥膏,爛的菜,垃圾在裡面撿.
在這個時候我兒子問婆婆在做什麼?.
你是家長你怎麼回答呢?.
婆婆在做什麼?.
我告訴他,婆婆可能很窮.
所以在垃圾堆裡撿一些東西吃.
好了,講完了.
是不是事情結束了?.
你怎麼可以….
我心裡覺得還沒結束,你不可以就這樣算了.

$^{121}$於是我就跟他說.
我們不如買一些水果請婆婆吃.
就跟他去了街市買水果.
又跟他想,婆婆可能沒牙齒.
我們要買一些軟一點的,買香蕉.
我當然不會買那些香引提子.
那麼貴.
我認為我可以負擔得到.
而且婆婆是可以吃的.
蘋果那些我又怕她咬不到.
所以就選了香蕉,然後帶回兒子去送給婆婆.
這件事其實對我兒子來說是印象很深刻的.
我偶爾問他,他是記得的.
問他記不記得那次送一些東西給婆婆,他是記得的.
當然有時候他不記得,我就要重述一次那個故事.
讓他知道一些道理.
對他來說有什麼影響呢?.
有很多人都認識我兒子.
我都可以很自豪地說我兒子真的很乖.
我覺得是真的.
其實你不是跟他說你要乖一點,他就會乖.
你是要透過你人生的一些經歷.
或者你親自帶領他去做一些事情,去讓他入心.
他真的很乖,乖到一個地步.
有一次我跟他去廁所.
然後去完廁所,看到有位伯伯拿著袋子進來廁所.
他就跟伯伯說,伯伯我幫你拿著.
伯伯就站在旁邊等伯伯去廁所.
我跟他說,伯伯你不用站這麼近.
但是他很乖,乖到一個地步.
我都有點驚訝.
這些事情不是你跟他說你要這樣做就這樣做.
一定是透過一些故事,透過一些你的經歷.
去讓他明白一些做人的道理.
其實我們成年人一樣,我們要透過不同的故事.
我們去學習,我們去看.
聖經裡面記載很多很古老的故事.
那些很舊的故事,有些人會問幾千年前那些以色列人.
他們發生的事情關我什麼事.
那些是古代人.

$^{161}$他們那個時代的事情這麼舊,關我什麼事呢?.
有什麼教訓可以吸收呢?.
不是的,其實是很重要的.
有些人很輕視舊約.
在聖經的正典還沒定型的時候.
有人鼓吹過我們要新約就行,我們要耶穌就行.
那些不合時候的,不要他.
但其實不是.
那些前人的經驗對於我們來說是很重要.
那些前人的經驗對於我們來說是教導我們.
我們如何面對試探,我們如何面對.
好像以色列人一樣,面對困難,面對試探的時候.
我們如何跟隨上帝的腳蹤.
我們如何去做,我們不用再失腳一次.
我們不用再跌一次,我們透過別人的經驗.
我們可以吸收到一些經驗去幫助我們成長.
那我又再講一個故事.
我帶兒子去超級市場的時候.
看到那些日清意麵,知不知道那是日清意麵呢?.
有隻雞印在上面的,不用煮的,煮滾水就可以吃的.
那些日清意麵都挺好吃的.
我就跟他說,我告訴他.
我試過吃了一個星期的日清意麵.
一天三餐都是吃日清意麵.
但其實不是我想的,不是因為好吃.
話說那時候我家人是很喜歡回大陸的.
很喜歡回鄉下,不是日本.
回鄉下去大陸住,一住一個星期.
那時候很悶,已經沒有手機,沒有WIFI.
去到那裡,外面全部都是草,田,牛就有幾隻.
最多就是新年回去玩炸牛屎.
不知道有沒有試過,找炮仗炸牛屎.
但是是很悶的.
小時候你不能出聲.
長大了,長大就不同說法了.
死都不肯去.
死都不肯去那怎麼辦呢?.
我家人就給了我一個星期的飯錢.
就叫我自己搞定.
於是就回了鄉下.

$^{201}$有錢身癢.
我一下子就買了我想要的模型.
和一個星期份量的日清意麵.
那怎麼辦呢?.
第一天我就把模型都組裝好.
然後開始吃日清意麵.
吃了三天我開始後悔了.
開始受不了了.
早知道就多買一排雞蛋.
就是這樣.
但是沒辦法,家人不在.
那你被迫要繼續吃,就吃了一個星期的日清意麵.
這個我告訴我兒子.
這件事其實對我有些影響.
我現在很孤寒.
我很孤寒的一個人.
因為很怕,不敢月頭花光錢.
很怕要吃日清意麵.
但是這個幫助他什麼呢?.
幫他認識到爸爸曾經犯過的錯.
我現在會給零用錢給兒子.
我都會教他很小心地用.
你自己好好地分配.
不要一次過用完.
你要按著你的需要.
你想想如果你花光了錢會怎樣.
這些是透過一些故事.
爸爸曾經失敗過的經驗.
去讓他有所成長.
所以我覺得有時候和別人分享自己做錯的事.
是可以的.
沒有人是24小時都英明神武的.
你一定是做錯過的.
你一定有做錯過的事.
當然很多人是不喜歡將自己的黑歷史說出來的.
但是你講出來其實是讓人去學習.
不要犯同一個錯.
不需要他在同一個坑裡卡住.
你讓他懂得避開.
我覺得這件事是很重要的.

$^{241}$大家有沒有想過聖經裡面其實很多人物.
他們是做錯過很多事的.
很壞的.
譬如說亞伯拉罕.
他經常說自己是壞老公.
說謊說不認自己的老婆是他的妹妹.
結果老婆被人搶了.
還要兩次.
我那天說的時候有個姐妹很大反應.
兩次啊.
證明了姐妹很生氣.
他的兒子也是這樣.
說過同一個謊話.
還有摩西.
摩西是殺人犯.
逃犯.
大胃王.
搶人老婆殺人老公.
這麼壞的事記在聖經裡面.
這麼醜的事為什麼要記在.
為什麼不刪除.
刪除了聖經就很好看.
但是問題就是.
如果我們沒有看過這些故事.
我們不能夠從他們身上去吸取經驗.
最基本的.
起碼我們知道不要這樣.
另外就是我們明白.
原來上帝的恩典真的很大.
大到一個地步.
是這樣的人他依然去接納.
當他願意悔改的時候.
上帝依然接納這樣的罪人.
正因為這些故事我們才知道.
上帝有多偉大.
他的恩典有多大.
甚至上帝犧牲自己的獨生子.
去拯救我們.
回到我們今天的經文.
其實我今天的經文.

$^{281}$好像說了這麼久.
沒有什麼關係.
只不過我想告訴大家.
其實聖經裡面所說的東西.
都是為了我們好.
上帝為了我們好的緣故.
他記得.
他寫下.
他願意我們去學習.
所有事情都是為了我們好處而去完成.
因此我們應該要去到.
很好的學習聖經.
去看聖經裡面的故事.
但是其實看故事.
也有一些要學習.
幾個層次.
我覺得看故事也好.
看聖經故事也好.
甚至看電影看影評.
我覺得都是適合的.
你看故事其實是要.
有三個層次.
最表面的層次是什麼呢.
就是看表層所要表達的.
舉個例子.
我們說一個故事說狼來了.
狼來了的意思是說什麼.
要誠實.
有一個表層的故事就完成了.
然後兒童主入學.
都是這樣的進度.
說薩戈的故事.
瑞利薩戈.
記得瑞利薩戈是誰嗎.
很矮的爬上樹上.
然後他見到耶穌.
接待他的時候.
他就說我要.
把財產捐一半給窮人.
然後我騙了誰.

$^{321}$我就給他四倍.
就是說什麼.
故事最表面那層.
做人要誠實.
耶穌喜愛誠實的人.
要做好事要捐錢.
但是.
這個就是表層.
另外我們看故事的時候.
其實如果大家記得.
或者可以重溫我在奸流.
講過兩次是關於.
耶穌.
責備門徒不讓小孩子進前.
和少年財主的故事.
這兩個故事.
如果在說耶穌.
責備門徒不讓小孩子進前.
這個故事是在說什麼.
表面那層.
耶穌愛小孩子.
是不是 耶穌愛小孩子.
就結束了.
兒童主義學那裡就可以散會了.
少年財主的故事.
記不記得是什麼.
少年財主的故事.
少年財主他跑去.
耶穌面前跪下問他.
我怎樣可以進入天國.
耶穌就叫他.
變賣所有財產跟從我.
然後他悠悠愁愁地走了.
耶穌就說.
有錢人進天國比.
下陀個針眼更難.
是不是.
表層的故事就說.
原來有錢人進天國很難.
再講深一點就是.

$^{361}$你要捐錢.
做善事不要貪財.
不過.
我們再看第二點就是.
這個故事應該在更加大的脈絡底下去看.
記得聽回那兩堂課.
其實.
小朋友和少年的官.
在其他福音書裡面.
都是連著去講.
小朋友是被拒絕.
小朋友是無權無勢.
是被阻擋的.
但是耶穌說進天國的.
就是好像這些小孩子.
相反少年財主有權有勢有錢.
衝去耶穌面前跪下.
這麼大動作都沒有人阻擋.
但是耶穌叫他.
你要捐財產.
他就悠悠愁愁走了.
進天國的不是這些人.
兩件事是比較.
小孩子和財主.
有錢和無錢.
有權和無權.
被人阻擋和沒有人阻擋.
這些全部是一個對比.
你用一個更大的脈絡.
去看聖經故事.
你的體會很不同.
其實好像一個.
連續劇.
你中間抽一集出來看.
你都看到一些東西.
你都看到一些故事.
但是你可能錯過了整個脈絡.
你就會有些東西.
你是不明白.
或者是有些重要的訊息.

$^{401}$是沒有了.
到最後.
講的是從故事.
去了解講故事的主人.
耶穌.
為什麼他要講這樣的故事.
上帝為什麼要.
將他的故事講給你聽.
為什麼要講不同的.
耶穌說不同的比喻.
那個比喻背後是講什麼.
如果又說回小孩子和財主.
我們如果用.
我們自己人的眼光去看的時候.
我們人類中心主義.
我說這些就是.
我們繼續關心的就是.
我可不可以上天堂.
我學小孩子還是學財主.
我應該怎樣才可以上天堂.
這個是我們從人的角度去看.
很擔心的.
很多人都說我怎樣才可以上天堂.
但如果你想想耶穌的角度.
用耶穌的眼光.
他為什麼要講這個故事.
他在責備什麼人.
他是在說.
這些人才是小孩子的人.
才是進天國的人.
為什麼你阻礙他.
為什麼你會阻礙他.
我們是不是應該要學習的是.
我要有一個心腸.
我要開一條路給小孩子.
我要開一條路給他.
可以進到耶穌面前.
我不是阻礙他.
阻礙他.
這些人不適合來我們教會.

$^{441}$我們不適合你.
你找找其他教會.
或者是有些人根本沒有機會.
去到上崇拜.
星期六日都要上班.
不好意思.
你可不可以請假.
你可以開一個晚上的晚課給他.
你可以開一些其他日子的崇拜給他.
你作為基督徒.
作為耶穌的門徒.
我們是要想辦法.
讓人去到耶穌面前.
這個就是.
我們要透過聖經的故事.
去了解耶穌的心腸.
然後跟著去做.
焦點不再屬於我們自己.
不是再看著我們自己.
而是講故事的主人.
他講這個故事給我聽.
到底有什麼意思.
耶穌這麼多的教導.
到底有什麼意思.
我看到大家開始釣魚了.
所以最重要記住.
三個層次.
看故事表面.
最簡單的.
表面的那個.
大家不要停留在這裡.
看深一點.
在一個更大脈絡下看故事.
你起碼整本書去看.
整本聖經去看.
你起碼看那一卷裡面.
前前後後有什麼.
跟著了解.
這個故事的主人.
耶穌他講出來.

$^{481}$是為了什麼.
他想教導你什麼.
不要再將眼光集中在自己身上.
鼓勵大家用完三個層次.
去看聖經.
對大家將來來說.
我相信一定會有幫助的.
有得益的.
最後一次講到.
在聖經裡面不斷進心.
多謝大家.
我們有一分鐘的時間.
我想大家安靜默討.
去思考剛才講到的內容.
然後交給警法隊.
\newpage



\section{提摩太後書 3:16-17}
\label{sec:cNC4ab19lyE}
\textbf{2023-06-18 早堂《說故事的爸媽》提摩太後書 3\_16-17 講員:吳景霖傳道}
\newline
\newline
連結: \href{https://youtube.com/watch?v=cNC4ab19lyE}{\texttt{https://youtube.com/watch?v=cNC4ab19lyE}} ~~~~ 語音日期: 2023-06-20
\newline
\newline
\hyperref[sec:eaFKBNr6R3g]{\small{< < < PREV SERMON < < <}}
~
\hyperref[sec:index_chronic]{\small{[返順時目]}}
~
\hyperref[sec:index_scriptual]{\small{[返順卷目]}}
~
\hyperref[sec:I3SP1geCVE8]{\small{> > > NEXT SERMON > > >}}
\newline
\newline
提摩太後書 3:16-17
\newline
\begin{longtable}{cl}
\hline
\hline
章節 & 經文 (和合本修訂版)\\
\hline
3:16 & \begin{tabularx}{0.7\textwidth}{X} 聖經都是神所默示的，於教訓、督責、使人歸正、教導人學義都是有益的， \end{tabularx} \\ \\ \relax
3:17 & \begin{tabularx}{0.7\textwidth}{X} 叫屬神的人得以完全，預備行各樣的善事。 \end{tabularx} \\ \\
[1ex]
\hline
\hline
\end{longtable}
$^{1}$各位弟兄姊妹 各位家長 各位小朋友 早安.
主女們平安.
我覺得很開心 我看到家庭崇拜.
其實我很久之前已經覺得.
我們一家人是應該有時間一起崇拜的.
不只是小朋友 有些人誤會.
以為這是小朋友崇拜 不是.
你什麼年齡的 總之你一家人.
都應該有時間聚在一起 一同敬拜上帝.
其實是豐富了我們.
聖經有很多地方都記載著.
有些人是唱歌跳舞去敬拜上帝的.
但是我們平時的崇拜是很少跳舞的.
但是剛才我自己都好像不知道怎麼做.
我不像黃姑娘那麼靈活.
但是你有些肢體動作.
跳舞去表達對上帝的敬拜.
其實是豐富了我們整個講道.
或者整個崇拜.
所以我希望接下來我們會有家庭崇拜.
大家一家人一起去敬拜上帝的時間.
小朋友有些聲音不要緊 不用理他.
正常的 小朋友活潑的.
今天的講題是說故事的爸爸媽媽.
免得那麼複雜 說故事的爸媽.
我想跟大家說一個故事.
小朋友聽一下關於吳全道的一個小時候的故事.
我小時候是住在一些唐樓.
小朋友可能不知道什麼是唐樓.
但是一些比較舊的樓.
那些樓是沒有電梯的.
大家可能覺得感受到沒有電梯.
但是我住在唐樓的五樓.
所以我是要走五層樓梯回家的.
很高 很辛苦.
那時候我是跟著我奶奶去住的.
奶奶就是爸爸的媽媽.
我跟著我奶奶住的.
我奶奶那時候帶著我去哪裡都會帶著的.
去公園也帶著 去晨運也帶著.

$^{41}$去街市買菜也會帶著.
我小時候就跟著她四處去.
有一次我奶奶又帶我去街市買菜.
拖著我去買菜 買完菜回家.
很辛苦抽著菜上了五層樓之後.
回到家發現.
咦 為什麼我手上多了一個算盤.
可能小朋友不知道什麼是算盤.
你回去問爸爸媽媽什麼是算盤.
總之多了一個東西叫算盤.
看到以為為什麼多了一個算盤在手上.
家裡沒有這樣的東西.
她想了想 哎呀 剛才是經過街市的時候.
應該有個檔口 門口就有個算盤.
可能是我經過的時候順手一拿.
然後就拿回家了.
哎呀 偷了東西.
好了 嘩 去街市買完菜回家上了五層樓.
才發現我手上多了一個算盤.
但是既然都回家了.
那就是有沒有人發現.
有沒有人發現.
沒有人發現了 是不是.
沒有人發現的時候.
那你覺得我奶奶她怎麼做.
當沒事發生就行了 是不是這樣.
不是的 我奶奶她就故意拿回那個算盤.
拿回五層樓去街市給別人.
然後又要扔回五層樓上樓.
她那時候都不年輕了 六十多歲了.
所以這件事給我一個印象.
就是我奶奶是一個很誠實的人.
她是一個很誠實的人.
當然有其他的事情.
我都知道她很誠實.
當然我是忘記了.
這件事是誰告訴我的.
當然不是算盤的店舖老闆.
是我奶奶她告訴我的.
她告訴我當時有一件這樣的事.

$^{81}$我拿了別人的算盤.
然後她又要走下去.
這些生活的一些片段.
其實都算是一個故事.
我們說的故事是一些過去發生了的事.
不是說只是童話故事叫故事.
聖經故事叫故事.
其實在日常生活裡面.
或者一些過去的片段.
都是屬於你的故事.
這些故事當你從述去告訴你的下一代聽.
告訴你身邊的人聽的時候.
其實有什麼意義呢.
剛才我奶奶她將她的故事告訴我.
她成為了我心裡面的一些記憶.
和一些指引 一些路標.
我後來我自己都算是一個誠實的人.
有一次我又和別人逛街.
實不相瞞是和我之前的女朋友.
去逛街 去到一間書局那裡.
一邊逛一邊聊天 一邊選東西.
當然因為注意力都在女朋友身上.
然後一邊聊一邊走就出去了.
走了很遠 走了幾條街.
接著看一看 咦 手上為什麼有些東西.
哎呀 剛才在別人的書局拿了一些東西.
在手上沒有給錢 走了.
那怎麼辦.
在這個時候我奶奶的路標就發生了作用.
她告訴我做人要誠實.
我就走回頭 我去到跟別人說.
哎呀 對不起 我剛才忘記給錢.
別人書局的職員也好.
哎呀 對不起 我忘記收你錢.
其實誠實是有它的好處.
原來你將一些過去的故事.
一些片段去到告訴你的小朋友的時候.
其實是成為了他人生的一個路標.
讓他以後去做人的時候能夠去學習.
我們透過故事 我們透過童話也好.

$^{121}$生活片段也好 聖經的故事也好.
其實是能夠盛載很多的資訊在裡面的.
盛載的資訊是多於你能夠就這樣告訴他.
我想每一個爸爸媽媽.
你都試過跟你的小朋友說.
做人要誠實 要乖一點.
通常你叫他乖一點 他都不知道什麼叫乖一點.
你很難去到將一個很抽象的概念.
去到就這樣告訴小朋友.
做人要誠實 什麼叫誠實呢.
但是當它變成了一個故事.
好像我剛才說的我媽媽的故事.
就讓我知道這一種就叫誠實了.
這種行為就叫誠實了.
於是就讓我記住了.
並且我能夠去到跟隨他的腳蹤去走.
我又說一個故事了.
小朋友又多一些故事讓你們可以留心.
這次又是我做主角.
不過已經長大了.
就輪到我牽著我的小朋友了.
那時候我的小朋友是幼稚園高班.
幼稚園高班 有點大.
那時候我就帶著他上學.
不過就快遲到了.
於是就牽著他急急腳走.
走到一個街市.
我看到一個老婆婆.
那個老婆婆彎下身在一個垃圾堆上.
一直在做一個這樣的動作.
知道是什麼嗎.
她在撿垃圾堆裡面一些爛的水果來吃.
在那些垃圾裡面撿一些沒那麼爛的來吃.
我看到那個場景.
我當然覺得心酸.
而那一刻我的兒子就問我.
婆婆在做什麼.
我那一刻就告訴他.
我說婆婆可能很窮.
所以她要撿垃圾去吃.

$^{161}$說完之後是不是就算了.
說完解釋了.
他回答了她的問題.
有人問問題你回答了問題.
好像事情結束了.
但是我心裡面告訴我.
事情還沒結束.
不應該是這樣完結的.
我就再跟我的兒子說.
我說婆婆這麼可憐.
不如我們去買些水果送給她吃.
因為很急急腳快遲到.
就去買了一梳香蕉.
選的時候也有技巧.
我當然不是選富士蘋果那些.
告訴她婆婆可能沒有牙齒.
我們要選一些軟軟的.
有什麼水果是軟軟的.
不用很容易吃到的.
香蕉沒錯了.
真是聰明.
就是買了一梳香蕉.
然後就拿過去送給婆婆.
我說婆婆我請你吃香蕉.
然後就又急急腳地上學.
然後就走了.
這個片段其實發生了之後.
我兒子開頭的記憶也很深刻的.
不過後來就要我不斷地告訴他.
但是這些事情.
是讓他成為了一個記憶.
也成為了一個很重要的品格培育.
到了今天我都自豪地說.
其實我兒子很有愛心.
每逢在街上看到黑衣的時候.
他都會問我拿幾塊錢去給人家.
然後他很樂於助人.
有一次我也覺得很有趣.
我們去廁所.
大家男生一起去男廁.

$^{201}$去完之後有一個伯伯.
他抽著袋子送進去廁所.
我兒子跟他說伯伯我幫你拿著.
你去吧.
然後他就拿著袋子站在伯伯旁邊.
就在那裡等.
我跟他說你不用站那麼近.
你站遠一點.
但是很有愛心.
我兒子是很有愛心.
這些我相信不是我跟他說.
你做人要有愛心.
不是這樣就可以培養到的.
而是要透過一些故事.
透過一些片段.
透過我們作為家長以身作則.
去告訴他你應該是怎樣做人.
當他長大了.
他記得這些片段的時候.
這些片段就會成為他生命的一些路標.
每一個小朋友其實都是跟隨著家長的腳踵去成長.
你怎樣做人.
其實很影響他將來怎樣做人.
我有些做老師的朋友他們都說.
你看到小朋友大概知道他的爸爸媽媽是怎樣的.
所以的時候我鼓勵各位家長.
你都是要以身作則.
去教導你的小朋友.
你試試將你生命的一些故事告訴他.
你小時候發生過什麼事.
你告訴他.
好讓他能夠跟隨你的腳踵去成長.
我又說一個故事.
小朋友聽著.
知不知道你有沒有去過超級市場買東西.
有沒有見過一隻叫日清伊麵.
是啊 有隻雞仔在上面.
不用煮的.
用滾水滾就能吃了.
好不好吃.

$^{241}$小朋友喜歡吃麵嗎.
不喜歡吃麵.
有些喜歡有些不喜歡.
我覺得日清伊麵很好吃.
但是我告訴你.
我試過吃了一個星期日清伊麵.
三餐.
日清伊麵.
但不是因為它好味.
為什麼呢.
那時候我已經是一個中學生.
我家人他們很喜歡回鄉下.
不是日本.
他們喜歡回大陸.
我很討厭.
因為那裡沒有東西.
那時候當然已經沒有電話.
甚至全部外面看出去就是田和草.
很悶.
於是乎我不肯回.
我去到中學.
我終於成功地反抗.
我說我不回.
於是乎我家人沒有辦法.
就給下錢.
叫我自己吃飯.
然後他們回了大陸.
我拿了錢.
我第一天就去買了模型.
還有足夠我吃一個星期的日清伊麵.
於是乎.
那些模型第一天就給我砌好了.
但是當我吃日清伊麵.
吃到第三天的時候.
慘了.
辛苦了.
吃到第三天的時候.
我還想我還有四天要吃.
怎麼辦.
但是都沒有辦法.

$^{281}$沒有錢.
唯有就要繼續吃日清伊麵.
小朋友這樣好不好.
不好.
有什麼教訓.
我不應該先花未來錢.
我不應該這麼沒計劃.
我也不應該貪玩.
把錢用來買模型.
應該用來買飯吃的東西.
買模型.
結果就要自作自受.
其實各位家長.
你不要怕把自己一些過去做錯的事情.
或者是一些失威的事情.
或者是失敗的事情.
告訴你的小朋友.
其實可以說的.
我經過超級市場.
看到日清伊麵.
我就會跟兒子說.
爸爸試過吃了一個星期.
他是可以從你的失敗裡面吸收經驗.
然後避免了他在同一個位置跌倒.
所以我又是再鼓勵大家.
你要多點跟你的子女分享.
去講你的故事.
去講你的經歷.
好讓他們能夠成長.
而今天的經文說到.
聖經都是神所默示的.
於教訓督責使人歸正.
教導人學而都是有益的.
叫屬神的人得以完全預備行各樣的善事.
如果我們說我們人生裡面發生的事情.
叫做真人真事.
聖經裡面的就叫做真神真事.
因為是上帝怎樣創造這個世界.
他怎樣去到介入這個世界.
去拯救犯罪的人.

$^{321}$甚至把耶穌基督賜給我們.
叫我們能夠相信耶穌.
回到祂那裡.
整個故事是在講上帝如何去救贖人類.
那個故事裡面不斷有很多的教訓.
讓我們能夠去到學習去到成長.
所以故事不是只是講給小朋友聽的.
其實耶穌在聖經裡面.
你看到耶穌教的門徒都是講故事的.
都是用很多的東西.
所以我剛才講的故事.
我相信大家可能中間會有些東西去領受到.
去學習到.
千萬不要小看故事.
讓我們的子女在故事裡面成長.
你帶他們回來兒童聖經學校.
讓我們的導師去多點跟他們講故事.
你不要小看這件事.
其實培養了他們很多品格.
弟兄姊妹.
你也是多點去看聖經的故事.
看聖經的教導.
作為家長的.
你看多點聖經.
你的品格建立得好的時候.
你都一定會幫到你的子女.
有一個良好的成長.
我們現在一同去禱告.
我們祈禱.
親愛的天父.
讓我們能夠透過這些故事去學習.
去成長.
好讓我們能夠跟隨你的腳蹤.
變得更加像你.
更加認識你.
求你在當中去拖帶.
保守我們每一位的家長.
每一位的小朋友.
每一位的弟兄姊妹.
透過聖經我們可以學習更多.

$^{361}$認識你更多.
多謝你.
祈禱交託.
奉耶穌基督的名.
祈求.
阿們.
\newpage



\section{傳道書 11:1-6}
\label{sec:I3SP1geCVE8}
\textbf{2023-06-17 晚堂《不知道中的信仰》傳道書 11\_1-6 講員:蔣文忠傳道}
\newline
\newline
連結: \href{https://youtube.com/watch?v=I3SP1geCVE8}{\texttt{https://youtube.com/watch?v=I3SP1geCVE8}} ~~~~ 語音日期: 2023-06-21
\newline
\newline
\hyperref[sec:cNC4ab19lyE]{\small{< < < PREV SERMON < < <}}
~
\hyperref[sec:index_chronic]{\small{[返順時目]}}
~
\hyperref[sec:index_scriptual]{\small{[返順卷目]}}
~
\hyperref[sec:9x4ci5x2GZ0]{\small{> > > NEXT SERMON > > >}}
\newline
\newline
傳道書 11:1-6
\newline
\begin{longtable}{cl}
\hline
\hline
章節 & 經文 (和合本修訂版)\\
\hline
11:1 & \begin{tabularx}{0.7\textwidth}{X} 當將你的糧食撒在水面上，因為日子久了，你必能得著它。 \end{tabularx} \\ \\ \relax
11:2 & \begin{tabularx}{0.7\textwidth}{X} 將你所擁有的分給七人，或八人，因為你不知道會有甚麼災禍臨到地上。 \end{tabularx} \\ \\ \relax
11:3 & \begin{tabularx}{0.7\textwidth}{X} 雲若滿了雨，就必傾倒在地上。樹向南倒，或向北倒，樹倒在何處，就留在何處。 \end{tabularx} \\ \\ \relax
11:4 & \begin{tabularx}{0.7\textwidth}{X} 看風的，必不撒種；望雲的，必不收割。 \end{tabularx} \\ \\ \relax
11:5 & \begin{tabularx}{0.7\textwidth}{X} 你不知道氣息如何進入孕婦的骨頭裡；照樣，造萬物之神的作為，你也無從得知。 \end{tabularx} \\ \\ \relax
11:6 & \begin{tabularx}{0.7\textwidth}{X} 早晨要撒種，晚上也不要歇手，因為你不知道哪一樣發旺；前者或後者，或兩者都一樣好。 \end{tabularx} \\ \\
[1ex]
\hline
\hline
\end{longtable}
$^{1}$我們兄姊妹晚安.
很感恩在晚堂中間跟大家分享上帝的訊息.
不知道大家今天的.
就是那個.
就是都是一個很.
陰霾密佈.
就是全日下雨的一天.
因為我剛才聽主席讀經文的時候.
就好像我們剛才說到.
就是滿了雨.
可能你都不知道今晚花開都崇拜.
卻是很感恩在我們.
就是.
之前的時間就停了雨.
也是很清涼的一個天氣.
讓我們都能夠很安然的.
回到教會當中一起敬拜著.
這個都是一個很特別.
都是一個呢.
好像配合到經文今天我們所說的那個.
就是有時我們很多的事情都不知道.
不知道的就是說我們不能夠.
就是.
預測得到或者掌握得到.
或者我們能夠好好地把握.
以致其實我們都要.
看看天如何轉變.
如果看一下我們.
有沒有一種準備好的心.
來去隨時應變.
原來天氣是雖然全日的下雨.
其實都可以瞬間會轉好.
以致我們都可以有一些即時的計劃.
即時的行動.
這個都是.
在傳道書的經文裡面.
都是一個讓我們基督徒的生命有一個.
提醒.
讓我們能夠隨時有一個好的準備.
我不知道大家有沒有機會多讀.

$^{41}$傳道書的經文.
或者其實這段經文都是今天我們一起.
聽到的時候才能讀得到.
盼望我們都藉著傳道書的信息.
讓我們在今天晚上裡面都能夠一起.
在神面前有一個.
特別好像我今天所說的那種.
很多事情我們不掌握.
不知道裡面我們怎樣好好地.
去活出我們的信仰.
我們先有一個禱告然後去聆聽上帝的讀.
(聆聽上帝的讀法).
漆黑天護 多謝你讓我們今天晚上.
都能夠有一個.
即是風暴風雨過後的.
一個平靜的天氣.
讓我們頂枝目能夠聚集.
讓我們能夠在你的殿當中.
去親近你 去敬拜你.
亦都是聆聽你的道.
願你保守我們每一個的心靈.
讓我們能夠完全地向你躺開.
讓聖靈當中親自地去運行中間.
讓我們有一個謙卑的心.
能夠讓你的道進入我們的內心.
引領我們繼續前行.
同心祈禱奉救主耶穌.
得勝名至二求 阿們.
不知道大家會不會都覺得.
在我們一種生活的環境.
譬如我們經過了三年的COVID.
我們有了疫症.
很多人都會用疫症之後.
其實都是一個逆境.
逆境就是信和逆之間的那個.
不是和我們那麼如意的逆境的意思.
都會令到我們很多的弟兄姊妹.
或者令到我們很多的身邊的朋友.
都是充滿挑戰 在這幾年的生活當中.
很多的轉變 很多的即時應變.

$^{81}$或者很多的都不能夠計劃得很長遠的狀態.
我就說在一個大環境的衝擊之下.
其實我們弟兄姊妹每一個生命.
我們每一個身體的狀態.
不知道大家有沒有覺得有些變化.
或者多了疾病 或者多了疲累.
或者是睡得不是那麼好等等.
我們身體生命的狀態.
可能都會有很多的疲倦 乏力.
或者是憂心 對前景未必看得清楚等等.
在這樣的狀況之下.
這個現實的之下.
無論環境或者我們自身的狀態之下.
其實我們更加需要去體會到.
其實我們需要主的力量 主的幫助.
我們在上帝的話當中.
有沒有令到我們成為一個得力的源頭呢.
更加就是我們需不需要一些身邊的人.
成為我們的支援呢.
特別教會是一個群體.
會不會需要教會的弟兄姊妹.
去支援我們 陪伴我們.
讓我們能夠分享到我們生命的那種感受或者掙扎.
這個都是我們需要承認.
原來我們已經不是能夠.
自己可以很本事地.
去為自己去計劃一切.
去掌握一切.
這個對我們信仰是一個很大的.
雖然是衝擊.
但也是一個提醒.
告訴我們我們不是一個人.
可以足夠應付到.
我們面對的那些挑戰.
面對的那些壓力.
我們需要同路人.
我自己也是做年輕人的工作.
更加體會到青少年.
其實都很需要不同的.
導師 牧者 家長 同伴等等.

$^{121}$去陪伴他們走過一些不同的人生.
一些迷惘.
一些新的階段等等.
需要一個信任的關係和群體.
需要學習真誠地去面對不同的問題.
弟兄姊妹說到這個情況.
其實就是一種現實.
一種生命的現實.
可能我們讀聖經的裡面.
我們很多時候讀新約.
譬如福音書 或者保羅書信的教導.
可能我們很多時候會從一個真理的角度.
去認識信仰.
去學習真理.
去明白信仰說什麼.
但是我們另外也有譬如舊約的傳統當中.
有智慧文學.
譬如浸賢傳道書等等.
他們都是一些從一個創造的本質.
或者我們身處的世界.
那個生命的現況.
來去開始進入.
去在當中尋求一個信仰的意義.
所以我覺得都是很適合我們今天.
弟兄姊妹去讀多些智慧文學的經文.
我自己也是 特別這幾年時間.
其實很多時候都是翻閱傳道書.
更加傳道書很多的內容.
都成為我自己無論是服侍.
或者面對一些生活的壓力等等.
生活的抉擇等等.
一個很重要的提醒.
一些鼓勵.
所以我就讀到11章的這段經文.
所以就跟大家分享今晚的訊息.
在11章這段經文裡面.
其實有6節的經文.
我就會從3個角度跟大家分享.
就是三方面的不知道.
我在螢幕上.

$^{161}$紅色的都給大家看到.
在第一節 第五節和第六節.
其實分別都有幾個位置.
去提出人其實很多東西都是不知道.
包括什麼範疇呢.
在第一節說到一個不知道.
將來有什麼災禍來到地上.
一種不知道災禍什麼時候來到的不知道.
有時候很意外.
有時候是一些疾病.
你都沒有為意.
其實你一發現的時候.
可能就已經是一種很大挑戰的狀態.
第二個就是第五節.
是一個不知道風從哪裡來.
或者一個胎兒在母親的腹中.
怎樣形成 怎樣成長.
這種是我們不知道.
但是有一個比較的第五節.
何況行萬事的神.
我們更加不知道.
一種對於上帝的工作.
那種隱密的工作.
我們不知道.
待會我們都會詳細再分享.
最後一個就是我都很喜歡的.
如果今個會選第六節.
這段經文.
是一個就是我們不知道.
那個效果.
我們做事有時候都是嘗試.
有時候不是每個人都能夠保證.
你做一件事情一定有效果.
一定有好的結果.
但是怎樣算呢.
所以當我們不能夠肯定有沒有效果的時候.
可能你會想到有兩種的反應.
一種就是消極一點.
那我也不知道什麼反應.
有什麼效果.

$^{201}$那我索性不如就懶做.
有時候會這樣說.
就是用一種很躺平的狀態去面對.
可能會這樣說.
但倒過來全圖書不是從這個角度來分享.
他說既然不知道的時候.
反而是叫我們積極去盡力去做.
待會我們去到第六節.
這個總結的時候.
我們就會一起去.
深入一點去想.
或者體會多一點.
究竟全圖書講的積極是什麼意思呢.
所以有三種的不知道.
不知道災禍什麼時候來.
不知道創造主隱密的作為.
所以我講了全圖書的文學當中.
多從一個創造主.
就是我們的主.
有很多不同的對人的生命的影響.
工作.
救贖主.
救贖我們的罪.
但全圖書就很甚少提罪.
不是表示他不理會人的罪.
只不過是從另一個角度.
從人是一個受造創造的質地.
那種限制.
來去講人的生命的本相.
而再次體會到.
創造主和受造之物.
就是我們每一個人.
或者受造界.
之間其實是有一個.
質地上的距離.
所以產生一種敬畏.
創造我們那個的主宰.
第三個就是.
剛才也提了不知道果效的時候.
是可以盡力去回應.

$^{241}$盡力去做.
我的.
先有一個小結.
希望我們去到最後的時候.
都能夠帶到一個訊息.
既然效果不在我們人手中控制.
我們可以知道.
自己的限制.
我們的boundary.
我們的limit.
我們的限度.
那我們要怎樣回應呢.
我們可以在我們的能力範圍之內.
就盡力去做.
因為相信上帝.
同時間在我們工作的中間.
有他自己的工作.
上帝會成就.
他的大能騎士.
這個也是傳道書.
有些人會覺得傳道書好像不是很強調.
上帝的工作.
但其實你自旅行間體會的時候.
你會看到他帶我們.
進入一個更深.
我們比較少的角度去了解上帝的工作.
我先在第一段.
大家留意的就是.
傳道書這些字的文學.
裡面其實會有很多.
好像謎語的表達方式.
他用很多大自然的景象.
人生百態等等.
好像我們有時候.
中國傳道中的猜謎.
那個故事.
那些諺語.
你猜背後.
下一句是怎樣.
背後的意思是什麼.

$^{281}$傳道書也有一種這樣的.
東方文學的.
這種味道.
所以你見到第一節.
裡面第二節也是說類似的東西.
當然他是一個現實.
當時的人會有機會.
這樣行動的意思.
但我們今天當然不會這樣做.
不會將糧放在水面.
意思是說什麼呢.
你要將一些糧食.
順著水流.
流到水下面的時候.
其實是.
可能你不知道.
會落在什麼人的手上.
落在一些你有需要的人手上.
或者會落在報恩的人手上.
你不明白.
你要分你擁有的財富也好.
你擁有的資源也好.
分成七份或者八份.
好像很不確定.
其實這個就好像我們.
今天如果大家.
有投資就分散投資.
一些高風險的你就投資一下.
低風險的你就投資一些.
類似一個這樣的.
好像是.
如蹤殺網.
都去讓都試.
這個表達方式.
因為你不知道將來有什麼.
禍患來到地上.
這個因為留心.
這個是分散.
投資的原因.
我相信傳道者應該.

$^{321}$都很難做到保險從業員.
如果你和你的客戶.
你分成七份八份.
因為其實你隨時禍患.
會來不到你身上.
可能你的客戶都不會幫你買保險.
通常那個保險會怎樣說.
整個市一定會暢旺.
一路這樣.
整個預測的話.
你去到幾年之後賺到多少錢.
可能你買保險都知道.
我不是說要.
講那個保險界的不重要.
但你發覺往往那些預測.
都不會如常發生.
因為經濟會大蕭條.
經濟會.
很多的危機.
大家過去二三十年都會體會過.
不是在一個.
分析裡面.
一個預測裡面就會知道.
就會掌握到.
所以它不是.
從這個角度.
從我們保險的角度來說.
但它是從一個很真實.
的角度來說.
所以傳道書的邏輯.
不是一個經濟學.
傳道書的邏輯是一種從人生裡面.
一個很現實.
當中.
它是提醒我們.
要擁抱這些現實.
要預留有空間.
其實會事如意為.
預留有一些.
跟我們的意願.

$^{361}$會甚至是相反的.
一些結果.
會發生在我們生命當中.
但當我們這樣想的時候.
或者傳道書這樣描述.
一個這麼現實的時候.
反而是叫我們.
不要守諸待到.
叫我們去盡量.
按理所能.
去分散.
去盡量去試不同的事情.
因為我們相信.
上帝在當中.
上帝會使用.
我們這些不同的.
投放資源在不同的事情上.
說不定其中一件.
你會見到上帝.
有回應給你.
這個就是傳道書在.
在它的智慧文學當中.
對於.
人生那種.
不能夠.
掌握的現實當中.
去給予我們一個.
仍然相信上帝同在.
仍然相信上帝是我們的確信對象的時候.
讓我們.
有一種很特別的人生態度.
一種積極.
但又不是純粹相信你自己.
所做的一切的人生態度.
而是相信上帝.
明明在我們當中.
那些可能性當中.
去顯明他的作為.
當我們不知道.
災禍何時來到的時候.

$^{401}$那傳道者.
他其實.
想提醒我們什麼呢.
他想提醒我們.
不要被這個現實.
嚇怕而.
不去行動.
不被這個現實嚇怕.
而不去行動 其實說什麼呢.
他的邏輯.
是說什麼呢.
既然我們都不知道禍患何時來到.
就是說.
其實沒有一些.
宿命律.
沒有一些因果定律.
人生不是一些.
會.
發生一個情況就一定會.
去到一個結果.
我們都不相信一個這樣的.
宿命的人生 命運是控制了我們.
我們不會這樣去想.
我們的人生.
是一件新的事情.
人生是在上帝當中有很多的可能性.
他既然不是一個.
宿命論的時候.
那不如就叫你豁出去.
你不要被那些事情.
嚇怕你.
不要被過去的經驗.
不要被一些過去的認知.
來去捆綁住自己.
這個就是.
我想在第三第四節.
和大家去帶出這個訊息.
他怎樣看得明白第三第四節呢.
第三第四節其實是一個.
對照的.

$^{441}$第三節是講雲.
然後去到最後第四節最後一句.
也是講雲.
第三節雲滿了雨.
那就自然倒在地上.
下雨 今天大家明白.
經常都是這樣經歷.
第四節最後就說.
望雲的.
就永遠不會去收到.
收割.
那我就再解釋 但是那個圖像.
是雷同的 中間那兩句.
都是講風的.
就是樹一向南倒下.
倒下.
或者向北倒下的時候.
樹就會倒在那裡.
就會存在那裡.
一棵大樹躺在那裡.
如果你行山走過的時候.
突然間看到.
整片樹很粗很大的樹.
都倒下了 橫向.
在你的路上 你就知道了.
之前可能早幾天.
打一場大風.
令到那麼大的樹都敗在.
路面上.
就算一個月之後走回路上.
沒有人可以搬得動它.
仍然敗在那裡.
其實是講一個這樣的情況.
第四節就好像.
天文台一樣 如果你是.
看風向的.
你就不會去殺.
你的種子.
所以當你看到講雲.
講風 依大自然的.

$^{481}$狀態的時候 究竟其實它.
借著這些大自然的現象.
想告訴我們一些什麼.
信息 我們知道呢.
中間分開了 先看.
第三節的意思 然後就看.
第四節的意思 第三節這個.
講雲和講風的這兩個現象.
其實是類同的.
都是講一個什麼意思.
就是說你看到結果.
那些雨都下來了.
那一定有一塊很重的.
黑的雨雲在你頭頂.
聚在水裡的時候.
才會淋在你身上 淋在你頭上.
一樣 剛才說樹倒在那裡.
好像剛才說行山的例子.
你想像可能早前.
就有風吹.
吹倒了它.
就是說每個結果.
都必有因.
這個自然界會這樣.
有果就必然有個原因.
去達成那個結果.
但傳說中講什麼呢 它想講反話.
它想告訴我們.
沒錯 每次的結果.
都有一個原因.
但是呢.
不是發生那些原因的事.
都必定.
達成那個結果.
不知道你明白我說什麼.
可以留心想清楚.
我們每一個結果的事情 一定有一個原因.
但發生這個原因的事.
可以有很多次的 可能只有一次.
會達成那個結果.

$^{521}$所以就是說如果我們每一個人.
都相信了這個因果律的時候.
出現了A的時候.
哪天發生B呢.
就會好像第四節那樣.
你看到風 又是刮那個風.
你就不會做事 你就不會殺種.
又是那個雲又來了.
你就不會出去收割.
因為淋到你 可能是.
狗淋頭都不一定.
傳統書書告訴你.
千萬不要被這些因果律.
自然律 宿命律.
去限制了.
我們的人生.
就算看著雲.
你都有機會.
可能是會出去的.
剛才我今早才.
跟弟兄有個WhatsApp.
他說6月12號.
大家有沒有留意 星期一.
可以說好像是.
整個月的.
唯一一天的好天.
這個弟兄呢.
都是我們教會的弟兄姊妹.
他很想帶我們出去玩那些.
KAYAK 大家可能都玩過.
我們很開心 王太太.
就全日去西貢.
風和日麗 跟他們一起.
經過了一個很美好的一天.
讓我自己都很大的鬆弛.
在那種.
好像一個.
跟那個.
海洋上面.
有一個很親近的感覺.

$^{561}$一種經歷.
我們都是在星期一的時候.
都不知道第二天會怎樣.
誰知一早起床.
我們先準備好東西.
一到起床 八點鐘 沒雨.
那就出發了.
如果我覺得不會的.
你留心看天氣報告.
全部都是下雨天.
你不會覺得有機會做到那件事.
所以今早的時候.
我自己有一個很深的感動.
覺得上帝很大的眷顧.
剛好收到弟兄的訊息.
他自己心裡面都體會到一個這樣的感受.
但如果只是.
看風望雨的.
那種預測 那種動靜的時候.
可能我們就什麼行動.
都不會做.
可能我們就不會經歷到上帝很大的作為.
也都不會對人生上面.
有很多新的發現.
新的經歷.
我們要.
離開一種.
固步自封的一種生命觀.
不會以為.
覺得命運是改變不了.
注定沒有彎轉.
而是我們要.
去嘗試想我們的現實.
我們的眼前.
不會被困住.
也都不會因為這樣就不敢去突破.
我們的人生.
這個是傳道書第一段.
這段經文裡面.
給我們的提醒.

$^{601}$原來禍患一些逆境的事.
不是什麼時候來到的時候.
不知道什麼時候來到的時候.
不會成為一個我們的陰霾.
不會成為我們的陰影.
以致我們不敢.
不願意去做其他人生.
可以試的事情.
這個原來是好像聽到的時候.
好像一個純粹.
一個很正面很樂觀的想法.
但其實不是的.
原來有一個在傳道書的.
作者裡面.
是一個很強的信仰的基礎.
因為你知道上帝有保守.
上帝有恩典.
上帝有出其意外.
新的作為.
你相信就算天氣的預測.
他都可以.
不在那個框框.
不在那個.
預計的出現的情況下.
去框得住.
上帝要預備的工作.
第二段.
第五節.
不知道創造主.
隱密的作為.
風從何到來.
骨頭在懷孕的婦人的胎中如何長成.
你尚且不得知道.
這樣行萬事之神的作為.
你更不得知道.
你見到我都講了.
行萬事的上帝.
其實是一個.
Who make the whole.
將事情圓滿.

$^{641}$成就那個上帝.
即是說.
你好像剛才那兩個例子.
那個實例.
你覺得.
是一個什麼例子.
我講第二個.
第一個大家都封了.
不再用了.
即是一個胎兒.
姐妹你是母親.
胎中曾經懷孕的時候.
可能你去醫生那裡照超聲波.
你三個月照一次.
你半年又照一次.
你每天都要照.
每個照的一段時間.
突然間多了手.
突然間心臟在動.
很奇妙.
你只會見到階段之後的改變.
但你不會看著它.
怎樣逐樣逐樣.
很細微的怎樣發生.
每條血管怎樣產生你不會知道.
即是我們在講一樣東西.
我們沒有辦法知道.
那個是超乎人的能力.
你不會24小時看著那樣東西.
你沒有那樣的精神.
你沒有那樣的狀態.
你睡了覺的時候都在成長.
你不會知道.
一夜之間一醒的時候.
開了花.
整個晚上怎樣發生你都不知道.
但雖然這樣講的時候.
這是人的限制.
人沒有辦法全知.
但不表示上帝不知道.

$^{681}$不表示上帝不在當中掌管.
整個嬰孩的生命.
小生命的成長.
是上帝在當中隱密地.
讓他一步一步.
你每次只見到一個成果.
三個月看到的時候.
就是上帝一段階段的成果.
半年的時候又是另一個階段的成果.
你到出生的時候.
在哭的時候.
有一個很奇妙的生命.
離開母體.
出來了一個新的生命.
一個是上帝.
讓這些事情.
歧視完滿發生.
的一位上帝.
這個就是隱密的上帝的意思.
當我們這樣體會的時候.
聖經也用到一個比較.
更加不知道.
更加深奧.
更加複雜.
你怎會知道呢.
你怎會觸摸到.
你估計到上帝一定會這樣做呢.
這個就是.
人和上帝之間的距離.
當我們這樣體會上帝的時候.
人有什麼反應呢.
如果我們真的去.
誠心的.
去了解這種.
人的限制.
和上帝的超越的時候.
我相信我們的反應.
是一種.
嘩 這是一個.
讚嘆 只會有一種讚嘆.

$^{721}$上帝的智慧.
或者是只會更加.
敬畏上帝.
你會更加.
知道上帝是興邁不得.
你會更加躺開自己.
去等候上帝.
有沒有一些更加奇妙的作為.
會發生在你身上呢.
剛才.
剛才和金穆斯在門口.
也和丁子妃說 啊對 明天是父親節.
做傳代牧者也沒有.
沒有什麼慶祝父親節的.
意識 基本上.
我和她也沒有.
什麼時候有慶祝父親節呢.
兒子有東西送上去就感恩接受就算了.
不會說.
晚上吃飯.
很少安排.
我自己父親也九十多歲了.
也感恩住在安老院裡面.
我記得2020年.
的時候.
他四月中風.
剛剛那時候 記得一月開始是COVID.
所以.
到中風的時候.
你也知道2000年那個年度.
有時候我們還以為.
沒事的嘛.
分階段 以為.
看著那幾個案件.
降到某個數量.
以為可以不用戴口罩.
復常 當時也會這樣想.
大家記得2020年是這樣.
遇到他中風.
進醫院的時候也很擔心.

$^{761}$然後還要說.
因為中風.
又不是一個很嚴重的狀態.
要出來.
我們照顧不到他.
所以一定要安排在安老院裡居住.
這是環境的趨勢.
迫著我們一定要這樣安排.
但很感恩.
也有一個很好的地方.
讓他照顧.
讓他慢慢康復.
有時候我想.
如果靠家人照顧也做不到.
這樣的環境.
服侍的質素.
但他一直也沒有相信主.
一直也很反對.
就算我媽媽相信了主.
回到天家 其實她也不願意.
純粹記得當時.
我們安我媽的靈位的時候.
鄭牧師也去探望.
一起去看 和他分享的時候.
他也不會覺得有機會相信主.
簡單來說就是這樣.
去到八月的時候.
2020年的時候.
有一次機會.
當時我們也不可以去老人院.
但唯一可以見到.
你的家屬.
就是帶他去看病.
大家有沒有試過帶他去看病.
安老院也不敢帶他去看病.
你家人會不會去.
你去就不用安排一個.
不知道要穿好很多.
去醫院又擔心很多感染.
那時候我和弟弟談.

$^{801}$不如我們帶他出來.
也把握了一次機會.
鄭牧師有機會.
在安老院樓下探望.
探望他 一個很短的時間.
帶他下來的時候 其實我們要送他去醫院.
把握一段短短的十分鐘.
時間和他分享福音.
很感恩.
在這麼短暫的.
相遇.
帶了我爸相信主.
也即時和他.
談了一個水禮.
正式去歸入教會.
很奇妙 其實只有一個機會.
我會想起這兩年機會.
後來去到不知道什麼時候.
才開始開放給家屬去探望.
進去老人院.
只有這個機會 我自己把握了這個時間.
其實我和陳牧師 鄭牧師說.
不好意思 陳牧師 雖然這麼熟.
也說未必能帶他出來 未必能帶他相信主.
但我只是把握.
能夠掌握的時間.
把握我的機會.
你明白我的心急.
和我的那種自私的心.
我帶他出來.
但他能不能給你.
信主或洗禮.
我不知道 陳牧師也很幽默.
不要緊 我先把工具帶來.
先帶水禮.
到時祈禱 他肯信.
肯洗禮 就做吧.
其實就像傳道書有這樣說.
那位隱密.
作為上帝.

$^{841}$我真的不知道他的時間是何時.
要家人信主.
時間是何時.
我們家人都預備.
到了六月尾.
約一個平日.
帶他出來.
雖然遲了 但也要慶祝父親節.
的一個時間.
想像一下 能夠有這些.
見面的時候 也是很難得.
而且是信了.
一家都信主的一個.
在桌上的一個團契.
一個數說主因的時候.
上帝的作為.
上帝讓事情成就.
其實我不知道 我真的不知道.
全都人都不會知道.
但是只能夠.
用人的.
盡量你是預備好.
自己 迎接到.
上帝的作為 上帝就在適當的.
時候 讓我們.
見到祂的恩典 所以當我.
分享這個見證的時候 也都順應.
到最後一節.
其實都是給第一節節目一個鼓勵.
我們不知道效果如何.
但我們可以怎樣做呢.
我們是翹起雙手.
我們在這裡等.
上帝的事情什麼時候來 等他來.
其實不是 原來.
不是一種被動的等.
去到第六節的時候.
用一個殺種 大家都很熟的.
殺種比喻 新約也有.
去到詩篇126篇.

$^{881}$都有 是不是.
聖經裡面多用的一個imagery.
用殺種這件事.
這裡說早晨要殺.
晚上都要殺 不要停.
因為你不知道.
早上殺還是晚上殺.
有果效.
因為早上殺都會下雨.
晚上殺都會有.
黃蟲吃掉你的東西.
你不知道 盜賊 所以你不知道.
太多因素不是農夫.
可以控制 今天我們.
每個人工作.
有沒有人是務農的 舉手.
當然有 沒有.
我相信都沒有 可能你會去的.
都是週末去休閒一下.
relax的時候 租一塊田.
那種都是都市人.
都是辦公室工作.
在四步牆之中工作.
我們很少面對.
那些天氣來望天打卦.
工作 我們又不是漁夫.
我們又不是農民.
我們不在那些.
環境之中工作 未必會明白.
但我們嘗試放下.
我們那套 去想一下.
那個環境 真是望天打卦.
真是要聽天由命.
真是望著天.
才看有沒有東西吃.
其實那個就是農夫那種的生活.
或者那種的.
生活的狀況.
他怎樣做呢 難道他不做嗎 不會.
他就告訴我們.

$^{921}$或是早殺的.
或是晚殺的 或是兩樣.
都好 那就開心.
有商量出眾開心.
一份都ok.
但兩份都有收成就更加開心.
殺種的.
這個都不是比喻.
這個現實 其實是一個工作.
是提醒我們 儘管盡力去做.
雖然結果不在控制之中.
但其實是.
要去.
盡量去嘗試.
說不定.
好的結果.
就在我們盡量嘗試的那些.
地方當中會出現.
這個就是農夫的工作.
即是.
我兒子大學畢業了.
然後就找教書工作.
那.
可能有些在座的叔叔阿姨.
以前小時候都教過他兒童主義學.
很快就已經.
預備成為英文老師.
他跟我說的時候.
因為他以前是南校的.
他說要找南校.
基督徒 大家都是基督徒.
要找基督教學校.
雖然很多 你知道移民潮之下.
有很多工作的機會.
但你將你.
即是filter了又要南校.
又要是基督教學校.
可能一隻手都數完.
難道你其他都不找嗎.
即是說.

$^{961}$我又不喜歡interview.
即是說我不喜歡跟人問問題.
可能你第一次interview.
是預備你到上帝給你機會的時候.
你準備好自己.
懂得怎樣對待的時候.
你份工來到的時候 你就可以迎接這份工作.
是嗎.
即是可能有時.
跟他在討論這些問題.
他不明白其實你要.
不住這樣去嘗試.
然後你要準備自己.
最後的時候 上帝機會一到.
上帝要收成的時候.
你先要準備好 你先要ready.
去收穫上帝給你的成果.
這個就是.
詩篇126篇第六節.
那個大家都很熟的.
流淚殺種的那個比喻.
殺種的人帶著種子流淚出去.
必定會歡歡喜喜地帶著禍菌回來.
出去帶著淚水.
是一個很辛酸.
很不容易.
很複雜的心情.
去做這些工作.
但你相信.
其實上帝在當中會保守.
上帝在當中會看顧.
即是你相信.
你相信終有一天.
會歡歡喜喜地.
這些收成.
你背在自己的肩膀上.
自己去享受.
這個就是在詩篇126篇.
都給我們一個提醒.
我們那種要去.

$^{1001}$積極去多方的嘗試.
多方的回應.
是一種積極的等候.
上帝才會在當中.
我們會顯明到在我們生命裡.
自己會見到上帝的作為.
我最後用一首詩句.
一首禱告.
去完我今天的講道.
這是一位耶穌的神父.
他有一個禱文或禱告.
我覺得很有意思.
他在生命裡有一個很深的體會.
他說Trust in the slow work of God.
即是上帝較為慢的工作.
即是要信任上帝.
上帝較為慢的工作.
有時我們很想很快.
我們很想上帝按我的節奏.
按我的期望來成就.
原來我們做基督徒最重要的第一步.
是要和上帝慢慢和諧.
一種這樣的關係.
你剛才說了很積極.
不是迫使自己很急速.
你只是按時做應該做的事.
不要輕易放下機會.
剛才我舉了我兒子找工作的例子.
他不可以急速.
你放一百封信.
其實一分鐘都沒有工作.
你不知道.
你不知道哪封信在哪裡.
才有成果.
但不能急的.
不能說我這次發了一百封信.
下星期一定要找工作.
是不可能的.
你不知道上帝的慢動作.
何時實現在你的生命當中.

$^{1041}$所以這首詩歌我想和大家最後分享.
我寫了一篇詩句或討文.
有大約三段.
有英文和中文讓大家看.
第一段大概是說.
你要有信心對上帝的慢動作.
因為我們很喜歡跳過了一些中間的步驟.
跳過了中間的階段.
我們特別是香港人.
這麼快速的城市.
速度的生活.
我們有時沒有耐性走向未知.
但卻不明白.
人生必須經過這些不穩定的階段.
就是第一段的內容.
我們要培養一份耐性.
向著未知前進.
慢慢你會見到上帝在你面前出現.
上帝越顯越明.
我們要進入一個不穩定的狀態.
你都願意去進入.
原來很多東西不在我們人生絕對的把握裡面.
我們今天很想.
下一步問專家意見.
又問AI.
大家會想知道所有的資訊.
當遇到什麼時候你會怎樣應對.
原來是不行的.
因為你有時你的心情怎樣.
你的反應怎樣.
你都不能夠絕對做得到.
應該有反應都不一定.
到第二段.
裡面再講的時候.
他告訴我們.
原來我們每個人生命的成長成熟都是一樣的.
我們不可以太過催逼去成長.
最後都說我們不可以今天很想即刻急速地成為了明天的你.
在這段時間醞釀的時候.
孕育你的生命成長的時候.

$^{1081}$你千萬不可以迫上帝或祈禱.
我今天就要有一個很足夠的成熟度.
去面對人生一切的事情.
大家明白嗎.
人生其實都是隨著你的成長.
你越來越參透萬事.
你越來越按你的經驗累積之下.
你知道事情發生的時候你怎樣面對.
當你未發生過經驗的時候.
可能你真的要諮詢不同的人的意見.
但當你隨著你自己人生經驗累積了的時候.
你就懂得用你的處變不驚去回應.
也能夠提攜下一代教導他們怎樣去面對那些事情.
所以我們不要嘗試百妙全妙助長.
因為今天就即時可以變成明天才能夠孕育出來的你自己.
所以到最後一段的時候.
祂提醒我們只有上主知道我們新的生命的過程是怎樣經歷.
那個forming 人生的成長.
成熟在裡面是怎樣形成.
所以作者就鼓勵我們.
相信我們的主.
讓祂的手引領我們前進吧.
你相信祂的手是拖引著我們走.
接受那些焦慮的感覺.
接受那種好像還沒完成的狀態.
這就是作者提醒我們.
當我們能夠擁抱信靠上帝的slow work.
slow work不是代表上帝很慢很懶的意思.
只不過是說上帝的工作並非按我們自己的期望.
很想很快有成就那種性急的速度去實踐.
上帝有祂自己的plan 有祂的時間.
反而我們要放下我們自己.
更多去用一種開放的心靈.
用一種放下自己很多一切以為掌握的事情.
來去將生命投放給上帝.
接著你會發現上帝會逐步引領我們.
上帝的作為是很奇妙.
你只會最後讚嘆去敬畏我們這位主.
原來真是一位我們生命當中.
這麼獨一無二.

$^{1121}$這麼有自己一個這麼美好的計劃.
在我們人生裡面.
好 我們一起低頭禱告.
讓我們能夠更深的體會到.
你作為我們生命的主 創造我們的主.
我們作為你受造之物.
我們作為你所按著你的形象.
所做男做女的每一個獨特的你所愛惜的生命.
你都在我們生命當中有你的作為.
有你的安排 有你的旨意.
有你自己按你的時間表.
在我們生命當中實踐的各種的恩典.
各種的大靈.
就求主幫助我們眾弟兄姊妹.
我們的生命當中常常向你躺開.
向你保持一份的耐性.
保持一份等候的心.
去營建你一切奇妙的作為.
願你親自的保守大靈.
祝福我們每一個 同心祈禱奉主明求 阿們.
\newpage



\section{路加福音 12:1-24}
\label{sec:9x4ci5x2GZ0}
\textbf{2023-06-18 午堂《基督徒的日常》路加福音 12\_1-24 講員:曾裔貴牧師}
\newline
\newline
連結: \href{https://youtube.com/watch?v=9x4ci5x2GZ0}{\texttt{https://youtube.com/watch?v=9x4ci5x2GZ0}} ~~~~ 語音日期: 2023-06-21
\newline
\newline
\hyperref[sec:I3SP1geCVE8]{\small{< < < PREV SERMON < < <}}
~
\hyperref[sec:index_chronic]{\small{[返順時目]}}
~
\hyperref[sec:index_scriptual]{\small{[返順卷目]}}
~
\hyperref[sec:keP18GcFGPs]{\small{> > > NEXT SERMON > > >}}
\newline
\newline
路加福音 12:1-24
\newline
\begin{longtable}{cl}
\hline
\hline
章節 & 經文 (和合本修訂版)\\
\hline
12:1 & \begin{tabularx}{0.7\textwidth}{X} 這時，有幾萬人聚集，甚至彼此踐踏。耶穌就先對門徒說：「你們要防備法利賽人的酵，就是假冒為善。 \end{tabularx} \\ \\ \relax
12:2 & \begin{tabularx}{0.7\textwidth}{X} 掩蓋的事沒有不顯露出來的，隱藏的事也沒有不被人知道的。 \end{tabularx} \\ \\ \relax
12:3 & \begin{tabularx}{0.7\textwidth}{X} 因此，你們在暗中所說的，將要在明處被人聽見；在密室附耳所說的，將要在屋頂上被人宣揚。」 \end{tabularx} \\ \\ \relax
12:4 & \begin{tabularx}{0.7\textwidth}{X} 「我的朋友，我對你們說，那最多只能殺人身體而不能再做甚麼的，不要怕他們。 \end{tabularx} \\ \\ \relax
12:5 & \begin{tabularx}{0.7\textwidth}{X} 我提醒你們該怕的是誰：該怕那殺了以後又有權柄把人扔在地獄裡的。是的，我告訴你們，正要怕他。 \end{tabularx} \\ \\ \relax
12:6 & \begin{tabularx}{0.7\textwidth}{X} 五隻麻雀不是賣二銅錢嗎？但在神面前，一隻也不被忘記； \end{tabularx} \\ \\ \relax
12:7 & \begin{tabularx}{0.7\textwidth}{X} 就是你們的頭髮也都數過了。不要懼怕，你們比許多的麻雀還貴重！」 \end{tabularx} \\ \\ \relax
12:8 & \begin{tabularx}{0.7\textwidth}{X} 「我又告訴你們，凡在人面前認我的，人子在神的使者面前也必認他； \end{tabularx} \\ \\ \relax
12:9 & \begin{tabularx}{0.7\textwidth}{X} 在人面前不認我的，人子在神的使者面前也必不認他。 \end{tabularx} \\ \\ \relax
12:10 & \begin{tabularx}{0.7\textwidth}{X} 凡說話干犯人子的，還可得赦免；但是褻瀆聖靈的，總不得赦免。 \end{tabularx} \\ \\ \relax
12:11 & \begin{tabularx}{0.7\textwidth}{X} 有人帶你們到會堂、官長和掌權的人面前，不要擔心怎麼答辯，說甚麼話； \end{tabularx} \\ \\ \relax
12:12 & \begin{tabularx}{0.7\textwidth}{X} 因為就在那時候，聖靈要指教你們該說的話。」 \end{tabularx} \\ \\ \relax
12:13 & \begin{tabularx}{0.7\textwidth}{X} 人群中有一個人對耶穌說：「老師！請你吩咐我的兄弟和我分家產。」 \end{tabularx} \\ \\ \relax
12:14 & \begin{tabularx}{0.7\textwidth}{X} 耶穌對他說：「你這個人！誰立我作你們的判官，或給你們分家產的呢？」 \end{tabularx} \\ \\ \relax
12:15 & \begin{tabularx}{0.7\textwidth}{X} 於是他對他們說：「你們要謹慎自守，躲避一切的貪心，因為人的生命不在於家道豐富。」 \end{tabularx} \\ \\ \relax
12:16 & \begin{tabularx}{0.7\textwidth}{X} 然後他用比喻對他們說：「有一個財主，田地出產豐富。 \end{tabularx} \\ \\ \relax
12:17 & \begin{tabularx}{0.7\textwidth}{X} 他自己心裡想：『我的出產沒有地方儲藏，怎麼辦呢？』 \end{tabularx} \\ \\ \relax
12:18 & \begin{tabularx}{0.7\textwidth}{X} 就說：『我要這麼辦：要把我的倉庫拆了，另蓋更大的，在那裡好儲藏我一切的糧食和財物， \end{tabularx} \\ \\ \relax
12:19 & \begin{tabularx}{0.7\textwidth}{X} 然後要對我自己說：你這個人哪，你有許多財物積存，可供多年享用，只管安安逸逸吃喝快樂吧！』 \end{tabularx} \\ \\ \relax
12:20 & \begin{tabularx}{0.7\textwidth}{X} 神卻對他說：『無知的人哪！今夜就要你的性命，你所預備的要歸誰呢？』 \end{tabularx} \\ \\ \relax
12:21 & \begin{tabularx}{0.7\textwidth}{X} 凡為自己積財，在神面前卻不富足的，也是這樣。」 \end{tabularx} \\ \\ \relax
12:22 & \begin{tabularx}{0.7\textwidth}{X} 耶穌又對門徒說：「所以，我告訴你們，不要為生命憂慮吃甚麼，為身體憂慮穿甚麼。 \end{tabularx} \\ \\ \relax
12:23 & \begin{tabularx}{0.7\textwidth}{X} 因為生命勝於飲食，身體勝於衣裳。 \end{tabularx} \\ \\ \relax
12:24 & \begin{tabularx}{0.7\textwidth}{X} 你們想一想烏鴉：牠們既不種也不收，既沒有倉又沒有庫，神尚且養活牠們。你們比飛鳥要貴重得多呢！ \end{tabularx} \\ \\
[1ex]
\hline
\hline
\end{longtable}
$^{1}$請我們一起同心祈禱.
多謝主再一次我們能夠聚集.
我們可以有聖餐 紀念十家的大恩.
我們都能夠有神的話語.
彼此的勸勉.
求你在我們中間.
引導我們大家的心靈.
去明白你的真理的教誨.
一起去學習 一起去持守.
我們這樣禱告.
奉耶穌基督的名.
阿們.
用這個題目.
很想引發大家一起去思想.
我們基督徒.
有什麼是我們經常.
要來到去持守.
在我們的心靈裡面.
在一到廿四節.
如果我們再延伸到三十二節.
有六次提到.
懼怕 不要懼怕.
如果是集中在第四到第七節.
是有五次提到.
要怕 要懼怕.
和第七節的不要懼怕.
為什麼耶穌要這樣來到去說呢.
在第十一章的五十三 五十四節.
是整段聖經主耶穌對法利賽人.
民事那個嚴厲的譴責.
你可以想像到那個氣氛.
和那個時代.
是有一個急劇的轉變.
耶穌不再是人見人愛的夫子.
而是有一群的人.
準備有組織地要將他除滅.
十一章的五十三至五十四節.
耶穌從那裡出來.
民事和法利賽人就極力的催逼他.
引動他多說話.

$^{41}$不是要聽他的道.
是要尋找機會去除滅他.
第五十四節.
聽完之後就私下的規定.
要拿他的華柄.
你可想而知.
作者路加將這段經文的氣氛帶進來的時候.
突然間出現一個場景.
有幾萬人聚集.
幾萬人是什麼概念.
大球場有多少人.
幾萬人聚集要聽耶穌的講道.
耶穌就來到首先提及和他的門徒.
因為第一節這樣說.
十二章的第一節.
耶穌開講的時候.
首先對他的門徒說.
即是很內圍在耶穌身邊的少少一部份人.
來提醒他們警告他們要小心.
自己不要陷落在一個假裝裡面.
屬靈的假裝.
就是聖經說的假冒為善.
就是法利賽人他們現在看不到自己的危機的假裝.
就是假冒為善.
而且門徒要留心的不單止不是假裝.
而是他會發酵的.
他會擴大的.
要留心這件事.
即是在心靈裡面.
要來審查一下自己有沒有一些假裝.
而以為自己已經具備.
這是耶穌對門徒說的.
很快就已經到第三節.
以及去到第七節.
就已經開始來探視.
門徒是需要持續為耶穌的真理.
這個救恩來宣揚.
無論有多大的阻礙.
都要來宣揚.
其中的原因就是第九節.

$^{81}$在人的面前如果不認我.
人子將來在神的面前.
都不會認這樣的人.
就是門徒持續要無論如何.
反對也好.
不反對也好.
歡迎也好.
都要來將真理.
透過不同的方式去說出來.
讓人認識耶穌.
因為第九第十節.
是相當恐怖的.
不過是將來.
神不會認這樣的人.
所以耶穌就勉勵門徒.
在眾多幾萬人的面前.
來這樣告誡門徒.
剛剛十六號.
康文署和中文大學.
出了一個十年的研究報告.
就說有一個隱憂.
他們用了很多的年齡層.
其中收窄了範圍.
中年的成年人.
超過一半以上.
沒有每個星期的運動.
不足夠的運動量.
很多的理由.
其中有幾個的理由.
相當累.
相當忙碌的加班.
天氣不好.
很多的具體的理由.
收窄的一個結論就是.
能不能將運動生活化.
日常化.
例如泊車.
我喜歡泊在很遠很遠的爛地.
然後走回市中心買東西.
買日用品.

$^{121}$為什麼要泊到這麼遠呢.
已經成為了我的習慣.
那個走遠路.
是我要做運動.
我要出汗.
會不會我們能夠在我們的生活裡.
將主耶穌的這麼嚴肅的警告.
將來才發生.
將來有人不認你.
因為你不見證耶穌的真理.
我們有沒有辦法將我們每日的生活.
在我們能夠見到的環境.
人事物.
去將這些真理講出來呢.
叫人認識我們是跟隨耶穌的門徒呢.
不單止是運動的日常化.
生活化.
是信仰的生活.
已經是我們的一部份.
已經不知不覺間.
我們就已經讓人知道我們是屬神的兒女.
這樣的一個生活化.
是我們需要學習.
不過更重要的是第四到第七節.
主耶穌叫我們要小心.
我的朋友啊.
我對你們說.
那些殺身體以後.
即是說第二第三節.
你要去見證真理有代價.
可能你會受逼迫.
如果我們的主在十一章的五十四節.
有人私下規定他.
要拿他的把柄.
即是說有組織地要準備除滅耶穌.
你是他的門徒.
你不要以為很簡單.
是會面對那個代價.
不過我的朋友.
我們要學習.

$^{161}$怕一些真的要怕的東西.
那些最多能夠殺你的身體.
而不能夠在你的靈魂做什麼的人.
你不需要怕他.
怕是我們這個人類社會不斷教導我們.
你需要怕紅綠燈的車衝過來.
我的兒子很小的時候.
手指小到可以入電掣.
第一次被媽媽很嚴厲地責罰他.
以後怕了不敢了.
要教他的.
他不懂的.
他的手指真的穿進去穿到的.
很嚴重的.
去年大家都知道有些嚴重的新聞.
柬埔寨的KK園區.
將人這樣擄走.
禁錮要他做一些非法的事.
大家如果有看過的新聞.
就是在台灣的機場.
有些警務人員舉牌.
不要去騙你的假的.
但人們都不聽.
因為幾倍的人工.
在台灣沒有.
真的有人去了.
不是假的.
是朋友打給我叫我去.
一去到下機.
在泰國就已經被人捉去柬埔寨了.
這些是屢見不鮮的事.
最近警方主動介入一些鵝片的網上.
或者電話的.
特別是一些重的災區.
就是老人家.
來警告他們.
勸勉他們千萬不要聽.
不要相信.
是誰來的.
說你的名字你都不要相信.

$^{201}$這樣來提出警告.
這個世界是告訴我們.
是有很多危險的.
我們要接受這些警告.
但我們很難想像主耶穌.
警告他的門徒.
跟隨他的門徒.
講真理.
你要真正怕的是這位造物主.
因為第五節.
我要指示你們當怕的是哪一個.
就是這位殺了你之後.
又有權利將你丟在地獄裡面的.
我實在告訴你們.
正要怕他.
怕完之後.
又叫我們不要怕.
因為第六節.
立即就免禮門徒.
原來我們這位造物主.
是已經對我們完全的了解.
認識.
和我們的生命的掌握.
都是在他的手中.
五個麻將.
不是賣二分銀子嗎.
但在神的面前.
一個也不忘記.
就是你們的頭髮.
也都被數過了.
不要懼怕.
第一次出現了.
主耶穌叫我們怕神.
他有權的.
是造物的主宰.
跟著叫我們不要怕.
因為你們這班門徒.
是神所寶貴.
所認識.
和所掌握的生命.

$^{241}$這裡給我們一個很重要的.
在這大段聖經裡面的.
第一方面的訊息.
就是我們要建立一個對神的敬畏.
對神的認識.
認識了神之後.
我們就不會輕易去懼怕.
一些敵黨真理的人.
我們應該要去傳揚神的真理.
第二第三節.
剛才我們還沒有讀.
他說掩蓋的事是不會露出來.
隱藏的事都不會被人知道.
即是說是會的.
是一個反面的說法.
是會的.
始終都會被揭露出來.
這裡說的不是個人的私隱.
是在說真理的問題.
第三節.
因此你們暗中所說的.
將要在明處被人聽見.
在內室附耳所說的.
將要在房上被人宣揚.
似乎不是說那些是是非非.
暗中的論斷低評.
是在說我們要持續.
為神真理作見證的媒介.
那個功夫.
是會被人認識的.
所以才有第四節.
叫我們來到去.
不需要懼怕針對真理的人.
所以盼望我們能夠分得清楚.
我們是為神的真理去作見證.
感宣家的弟兄姊妹.
我們需要來到去.
看到這段聖經的嚴肅性.
耶穌說我們要看到.
在神的眼中的我們.

$^{281}$是否一班為神的國度.
來到去效力.
認真不會理會那個環境.
不會遷就那個環境.
像保羅所說的提摩太后書第四章.
無論是得時.
福音是很容易傳的時代.
還是不得時拒絕福音的時代.
仍然去見證福音呢.
這個是我們需要來到去思考的.
這個世界提供太多的便利.
很多東西都是user friendly.
每樣東西都要持續.
從這個角度來到去發展.
信仰就不能夠太過user friendly.
這段聖經已經是很嚴肅的.
要求每一個信徒.
應該要來到為真理.
為耶穌基督去作見證.
而準備付代價.
記住有更大的後台.
做物主.
他了解認識我們.
第二段聖經.
我們由第十三節到二十四節.
本來沒有太大的直接關連.
是有一個在人群裡面的人.
就有一件事.
很想耶穌幫他去做.
就是涉及到兩兄弟的分家產.
似乎有些講論就說.
長子是分三分之二.
那個小孩不是分一半.
是分三分之一.
會不會是因為這樣的緣故.
這個兄弟聽到耶穌這麼正氣.
這麼公義.
這麼有那種威嚴.
就覺得耶穌是一個很正直.
可以很公平的幫他處理分家產的事.

$^{321}$作者很巧妙地將這件事件.
引進去上文這段經文裡面.
然後帶出耶穌另外的教訓.
但是整個段落.
都是帶出一個重要的聖經真理.
這個是作者巧妙的地方.
就告訴我們.
這個人就把握機會在幾萬人當中.
他不是耶穌的門徒.
來去求耶穌請你.
吩咐我的兄弟和我分家產.
耶穌即刻地去責備他.
十四節你這個人.
誰來立我作你們地上的審判官.
幫你們分開家業呢.
第十五節.
大家留意這裡.
主耶穌藉這個機會.
又或者作者非常巧妙地.
將這段本來和上文無關的事.
拉在一起.
留意下主耶穌這裡說的一個生命的透視.
你們要謹慎自守.
明明是那個人私下的事情.
和眾人無關.
主耶穌把握一個機會.
向眾人講述一個很重要的真理.
又是關乎到你自己和神之間的.
永恆的生命問題.
就要躲避一切的貪心.
因為真正的生命豐富.
不是家道豐富.
物質的豐富.
是真正的生命豐富.
其實去到第21節的總結.
就是之後耶穌用了一個比喻.
接著21節的總結就是.
凡是為自己積財.
在神的面前卻是不付足.
都是這樣的下場.

$^{361}$那個下場就是第20節.
就是比喻裡面的一個勤奮的財主.
在他整個的生命裡面.
因為既然今夜他就要離世.
即是說蓋棺可以定論.
比喻裡面的這個財主.
他一生不知是多少歲被收回靈魂.
他的一生沒有為自己的生命豐富.
在神面前的準備.
來積蓄.
一生為他自己去張羅的是.
物質世界裡面的真正的富足.
物質世界的富足.
我們很容易可以去數算.
整個社會.
整個電視媒介.
不斷告訴你.
如果你有足夠的富足.
你就可以這樣吸數的享受.
出海或者各樣的享受.
各樣的美食.
這個世界已經有足夠的教導我們.
甚麼叫做今世的富足.
耶穌在這裡把握機會.
剛才說的就是第14節.
要謹慎免去一切的貪心.
因為人的生命.
不是在乎物質.
家道的豐富.
所以再次回去.
和這位主.
那個生命的關係.
我們有沒有這個生命的財富.
我們去不斷的積蓄.
積蓄神.
看我們這個人的生命.
因為第19節.
這個人居然會和自己的靈魂說話.
他說靈魂啊.
你有許多的財物積存了.

$^{401}$可以作多年的費用.
只管安安日日.
吃喝快樂吧.
對話的和自己.
主耶穌的比喻裡面.
祂不是坐在這裡.
祂真的不斷用祂的腦袋.
去運籌遺額去發展.
即是說祂真的勞動.
但是居然有一個聲音.
從來沒有聽到.
但最後聽到的時候.
已經沒有機會了.
就是第20節.
神在這個時候卻對他說.
無知的人啊.
今夜就要你的靈魂.
你所預備的要歸給誰呢.
人類有不同的民族.
有不同的階層.
人類唯一公平的就是.
我們沒有帶什麼到世上來.
保羅告訴我們.
提摩太前書6章.
應該是第10節.
沒有帶什麼到世上來.
也不能帶什麼去.
人類唯一最公平的.
就是這個point.
當然了.
一色憾就決定了.
你是富裕家庭.
窮人家庭.
你是黑色人種.
你是另外的人種.
如果你在印度.
如果你是賤民.
那你就一世是賤民.
這是在人類的社會裡面.
但是聖經告訴我們.

$^{441}$物質世界以外.
當我們來到這個世界.
和離開這個世界.
有一個富足.
是耶穌警告我們要盡快發展的.
就是心靈的富足.
就是要聆聽到神和你的對話.
所以很快24節之後.
主耶穌又再一次勸勉門徒.
你們很貴重.
你們很重要.
你們去服事.
神是不會忘記.
由24節一直這樣去勸勉.
一直到第31節.
你們只要求他的國.
這些東西都要加給你們了.
很清楚主耶穌在這24節裡面.
告訴我們.
神的國度是否在我們的心裡.
我們的生命是否因為神的國度.
我們忠心的效力.
而我們的生命豐富呢.
上星期見到神學院的一位.
已經離開了.
去到加拿大服事的老師.
魏健志老師.
是我的同班同學.
其中一組的關心小組.
4月我兩位同學飛去加拿大.
他是宣教士.
探望魏健志老師Tim Sir.
探望完之後.
上星期突然之間.
同學的群組很快就收到訊息.
突然離世.
一個月多的時間就離開世界.
很安詳地離開世界.
他的太太.
昨天寫了一個最後的對話.

$^{481}$給所有關心Tim Sir的人.
是Tim Sir親自講的.
他說我不知道為什麼.
因為他差不多離世了.
生命的長短不重要.
重要的是否作好了主的功.
他不知道為什麼.
這個癌症.
突然之間淋到他身上.
但是說我不需要知道.
他吩咐他的家人.
和所有認識他的弟兄姊妹.
不要質疑上帝.
他說我苦心自問.
我的人生已經盡了我的力.
來為主作功.
星期三的晚上.
我們同工會祈禱會完了.
吳全道在車上跟我說.
你知不知道見到神學院的.
魏健志老師突然離開世界.
我說知道 今天收到短訊.
他說早一段時間.
他這位魏健志老師.
原來是十幾年前.
他婚前輔導的老師.
十幾年沒見了.
知道弟兄要去加拿大.
他還沒有病之前.
不知道有病之前.
主動透過不同的朋友.
聯絡了Jason全道.
看看有什麼需要.
可以幫這位十幾年前.
幫他做婚前輔導的輔導員.
就是我的老師Tim Sir.
這樣的一個人品.
這樣的一個為神國度.
跟自己有什麼關係.
居然間上心.

$^{521}$很想聽聽有什麼可以幫到這位全道者.
快要去加拿大了.
去溫哥華了.
這樣一件事.
這麼立體讓我知道.
上午收到的訊息.
晚上同工告訴我.
問我知不知道.
而他自己親身經歷.
這位牧者.
這樣關心他們的夫婦.
弟兄姊妹.
我們有沒有這樣的.
為神的國度效力.
我們有沒有這樣.
在我們的生活裡面.
對做物主有一個敬畏.
以致我們願意在神的國度裡面.
我們努力忠心的去傳揚.
神的真理.
不計較的代價.
這樣來傳揚.
我們一起祈禱.
多謝主帶領我們.
今早能夠一起聚集.
敬拜事奉我們的做物主.
我們的救贖主耶穌基督.
幫助我們能夠.
真是在我們的生命裡面.
看看什麼是我們真正.
永恆重要的.
沒錯地上的.
是重要.
是我們每天都足手可及的事情.
我們需要去擔心.
需要很多的將來.
需要很多的計劃.
但是如果我們不曉得.
為永恆計劃.
這才是最可怕.

$^{561}$最可憐的.
求主幫助.
這樣禱告奉耶穌基督的名.
阿們.
\newpage



\section{路加福音 12:35-48}
\label{sec:keP18GcFGPs}
\textbf{2023-06-24 晚堂《儆醒僕人筵席》路加福音 12\_35-48 講員:黃頌言姑娘}
\newline
\newline
連結: \href{https://youtube.com/watch?v=keP18GcFGPs}{\texttt{https://youtube.com/watch?v=keP18GcFGPs}} ~~~~ 語音日期: 2023-06-27
\newline
\newline
\hyperref[sec:9x4ci5x2GZ0]{\small{< < < PREV SERMON < < <}}
~
\hyperref[sec:index_chronic]{\small{[返順時目]}}
~
\hyperref[sec:index_scriptual]{\small{[返順卷目]}}
~
\hyperref[sec:i86dI_9QW94]{\small{> > > NEXT SERMON > > >}}
\newline
\newline
路加福音 12:35-48
\newline
\begin{longtable}{cl}
\hline
\hline
章節 & 經文 (和合本修訂版)\\
\hline
12:35 & \begin{tabularx}{0.7\textwidth}{X} 「你們要束緊腰帶，燈也要點著， \end{tabularx} \\ \\ \relax
12:36 & \begin{tabularx}{0.7\textwidth}{X} 好像僕人等候自己的主人從婚宴上回來。他來叩門，就立刻給他開門。 \end{tabularx} \\ \\ \relax
12:37 & \begin{tabularx}{0.7\textwidth}{X} 主人來了，看見僕人警醒，那些僕人就有福了。我實在告訴你們，主人會叫他們坐席，自己束上腰帶，前來伺候他們。 \end{tabularx} \\ \\ \relax
12:38 & \begin{tabularx}{0.7\textwidth}{X} 他或是半夜來，或是天亮之前來，看見僕人這樣，那些僕人就有福了。 \end{tabularx} \\ \\ \relax
12:39 & \begin{tabularx}{0.7\textwidth}{X} 你們要知道，一家的主人若知道賊甚麼時候來，就不容賊挖穿房屋。 \end{tabularx} \\ \\ \relax
12:40 & \begin{tabularx}{0.7\textwidth}{X} 你們也要預備，因為在你們想不到的時候，人子就來了。」 \end{tabularx} \\ \\ \relax
12:41 & \begin{tabularx}{0.7\textwidth}{X} 彼得說：「主啊，這比喻是對我們說的呢？還是也對眾人呢？」 \end{tabularx} \\ \\ \relax
12:42 & \begin{tabularx}{0.7\textwidth}{X} 主說：「那麼，誰是那忠心又精明的管家，主人要派他管理自己的家僕，按時定量分糧給他們的呢？ \end{tabularx} \\ \\ \relax
12:43 & \begin{tabularx}{0.7\textwidth}{X} 主人來到，看見僕人這樣做，那僕人就有福了。 \end{tabularx} \\ \\ \relax
12:44 & \begin{tabularx}{0.7\textwidth}{X} 我實在告訴你們，主人要派他管理所有的財產。 \end{tabularx} \\ \\ \relax
12:45 & \begin{tabularx}{0.7\textwidth}{X} 如果那僕人心裡說『我的主人會來得遲』，就動手打僮僕和使女，並且吃喝醉酒， \end{tabularx} \\ \\ \relax
12:46 & \begin{tabularx}{0.7\textwidth}{X} 在想不到的日子，不知道的時候，那僕人的主人要來，重重地懲罰他，定他和不忠心的人同罪。 \end{tabularx} \\ \\ \relax
12:47 & \begin{tabularx}{0.7\textwidth}{X} 僕人知道主人的意思，卻沒預備，又未順他的意思做，那僕人要多受責打； \end{tabularx} \\ \\ \relax
12:48 & \begin{tabularx}{0.7\textwidth}{X} 至於那不知道而做了當受責打的事的，要少受責打。多給誰，就向誰多取；多託誰，就向誰多要。」 \end{tabularx} \\ \\
[1ex]
\hline
\hline
\end{longtable}
$^{1}$各位弟兄姊妹晚安.
我們一起祈禱吧.
親愛的天父 多謝你.
我們很渴望你的話語.
願你讓我們有聆聽的耳朵.
讓我們有好的心田.
可以聆聽你的話.
栽種你的種子.
多謝你 奉耶穌名求 阿們.
我想問大家.
你們有沒有試過有等門的經驗呢.
等人回來家裡.
有沒有試過.
應該有吧.
小時候.
我小時候.
經常等爸爸媽媽回來.
很想念他們.
經常看著門口.
覺得時間過得很慢.
尤其是不懂看時間表的時候.
覺得時間過得特別慢.
我小時候很黏人.
我放學回家.
馬上打電話給媽媽.
告訴她我很想念她.
因為我媽媽要上班.
就算已經懂得看時間表的時候.
依然覺得時間過得很慢.
特別是父母去喝酒的時候.
把我們丟在家裡.
我就會覺得.
爸爸媽媽為什麼這麼晚回來呢.
每當爸爸媽媽一回來家的時候.
看到我們乖乖地吃飯.
做好功課.
他們就會很開心.
我也感恩每當爸爸媽媽回來.
都會帶一些東西回來.
可能是宴會吃剩的東西.

$^{41}$他們會帶回來給我們.
爸爸媽媽很節儉.
有一次他們抽中了.
在宴會裡面抽中了.
一個樂園的門券.
我超開心.
因為我第一次去樂園.
就是這一張免費的門票.
相信做父母的話.
也有很多等門的經驗.
或者更加多.
我大專的時候.
就跟同學出去玩.
爸爸就很生氣.
因為我太晚回家了.
媽媽也很晚睡.
因為要等我們回來.
以前我沒有交代.
所以不回家吃飯.
也沒有說.
因此令到煮飯的媽媽很生氣.
當然後來我就學習.
要有交代.
不要令父母擔心.
所以無論父母等子女.
或子女等父母.
相信都是因為有對方在心裡.
也聽過一些家長和小朋友分享.
明明他們已經睡了.
但是一聽到門的聲音.
開門的聲音.
他馬上就彈起來.
因為休息的時候.
也沒有熟睡.
因為心中有些東西掛念.
掛念著對方.
原來聖經路加福音.
耶穌也曾經說過.
一些比喻是關於等候和警醒.
今天我們一起從路加福音.

$^{81}$十二章三十五至四十八節.
裡面有三個比喻.
看看耶穌對等候他回來的人.
有什麼要求.
也從而反思一下.
我們現在的生命的情況.
我們不如先看看.
為什麼耶穌忽然之間.
會無端端說一個.
關於警醒的比喻.
原來在路加福音.
十二章三十五至四十八節.
裡面有三個比喻.
原來在路加福音.
第九和第十八節之間.
第二次和第三次.
預言自己死亡的中間.
當時耶穌.
第二次預言自己死亡之後.
他.
猜險了七十個人.
去傳福音.
後來耶穌面對著.
一些質疑.
有人覺得他靠著鬼王.
說鬼.
耶穌就回應.
因為神的國臨到.
所以他就有能力.
去趕鬼.
而當時大家依然.
很想在耶穌身上看到神蹟.
而耶穌卻是.
勸他們要悔改.
因此我們看到.
第十二章的時候.
主題是警戒和免厲.
耶穌要門徒.
去傳福音.
同時要他們有警醒的心.

$^{121}$要從世界裡面分別出來.
不要像那些無知的財主一樣.
後來十二章有一個比喻.
是無知財主的比喻.
他在說.
財主花了很多時間.
在地上積聚財寶.
但是下一天.
他就死了.
提醒我們要積聚財寶.
在天上.
也不要憂慮去吃什麼.
穿什麼.
因為未信主的人.
是會這麼憂慮的生活.
我們只要.
去相信.
也要將自己.
所擁有的變成周濟窮人.
我們就可以.
慢慢積聚財寶在天上.
意思是.
我們要照顧身邊有需要的人.
然後.
耶穌就說了.
這個警醒的比喻.
提醒門徒要警醒.
後來就知道原來耶穌的價值.
和世界不同.
是天國的價值.
因此耶穌的國.
來到的時候.
不是太平.
而是有糾紛爭戰.
因為耶穌.
因為救贖的日子.
正式的來臨.
第十三章提醒.
大家要悔改.
十四至十七章有很多的教導.

$^{161}$然後耶穌.
就在第十八章.
第三次預言自己會死.
因此我們看到.
耶穌在這裡要提醒.
門徒要警醒.
提醒門徒.
在耶穌死了之後.
耶穌其實是隨時會回來的.
在這個期間.
他的門徒要隨時做好準備.
我們不要活得.
像沉睡了一樣.
不要只顧自己的享樂.
不要顧著積聚財寶.
而忘記了自己的使命.
所以第一點要和大家分享的是警醒的僕人.
警醒的僕人.
第35節說.
你們腰裡要束上帶.
燈也要點著.
這裡好像出埃及的.
逾越節晚餐的時候.
預備隨時就要走了.
所以腰那裡要束著帶.
燈也要點著.
即是預備好.
是一個工作模式.
不要懶惰.
第36節說.
自己好像僕人.
等候主人從.
婚宴的筵席上回來.
他來到.
溝門就立刻.
給他開門.
耶穌在婚宴回來.
有學者認為是指向.
耶穌在天上已經掌權的時候.
耶穌再來的時候.

$^{201}$應該有的態度.
是警醒.
不是叫我們不要睡覺.
不要誤會.
不是叫我們平日都不要睡覺.
而是這是一個比喻.
是說.
我們現在在地上.
等候耶穌再回來的日子.
我們的心態是要預備好.
自己有警醒的心態.
當耶穌再來的時候.
第37節.
主人來了.
看見僕人警醒.
就有福了.
耶穌吩咐門徒要警醒.
以致他們可以在.
末日審判的時候.
都站立得住.
真正的福是神在.
末世當中賜予的.
救贖.
所以門徒警醒等候主再來的時候.
必要在末日.
得蒙拯救.
然後繼續說.
我實在告訴你們.
主人必在他們.
必叫他們坐直.
自己束上帶.
進前侍候他們.
這裡耶穌.
用筵席去描繪.
神國的景象.
那些警醒等候主.
耶穌回來領受福分的.
就是那些可以.
在天上一同坐直.
竟然.

$^{241}$耶穌自己.
親自倒轉了.
主人的位置.
去侍候.
那些僕人.
這裡不單只強調僕人要警醒.
更加描述了主人的態度.
主人.
取了奴僕的角色.
自己束上帶.
進前侍候他們.
這個唐突的扭轉.
指明了神國.
裡面領導.
領導眾人的僕.
是僕人.
領導者是一個僕人.
主耶穌自己就是這樣.
他自己親自服侍.
門徒.
這裡37節.
我寫了粗體.
因為後面43節.
有一個重複的句式.
所以我們這裡特別留意一下.
第38節.
說:或是二更天來.
或是三更天來.
看見僕人者讓.
那僕人就有福了.
這裡指的是大概.
晚上九時到早上三時.
就是晚上的時間.
大家都知道當時的人.
是日出而作.
日入而息.
所以大概五六時就天黑了.
因此晚上九時.
其實就是大家已經休息的時間.
這個比喻強調.

$^{281}$沒有人能夠知道.
人子再來的時間.
因此我們要保持警醒.
不要分碎.
究竟分碎是什麼意思呢?.
好像是.
活得萬無目的.
根據常聞.
他在說每天在地上.
勞碌為著積聚財寶的人.
每一天生活就是為著.
賺錢而且是更多更多的錢.
多到自己.
都花不完.
當然我們是要有正常的工作.
但不是沒有.
紫警的渴望.
最近有一個報導.
說日本和南韓.
都有一個小子文化.
很少的小.
兒子的子.
小子文化.
意思是大家都不願意結婚.
更加不願意生小朋友.
不希望下一代.
都活在一個競爭的.
社會當中會受到傷害.
現在都活得毫無盼望.
當然我們.
看到其實韓國.
非常蓬勃的演藝圈.
他們的男團女團.
都非常耀眼.
韓劇也很出色.
很多劇集都在說貧富懸殊.
例如.
《上流寄生族》.
這套戲.
是一個社會悲劇的故事.

$^{321}$就是在說貧富懸殊.
韓劇都有很多.
以富二代來作為.
一個題材.
所以我們看到.
奢華美好的背後.
其實有很多受傷的生命.
因此日本.
也有躺平文化.
是一個沒有盼望.
社會的生存之道.
讓我們都不要分歧.
不要被社會帶走.
我們是不需要.
和別人競爭.
也不需要活得.
比其他人好.
也不需要踩著別人.
得到聖亡.
我們要警醒.
反而我們要關心身邊的人.
去分享自己所有的.
神會賜給我們所需要的.
讓我們都一起積聚財寶.
在天上面.
第二點.
要和大家分享的是.
警醒的家主.
第39節說.
家主若知道賊.
什麼時候來.
不容賊或偷房屋.
這是你們所知道的.
40節.
你們也要預備.
因為你們想不到的時候.
人子就來了.
耶穌這裡用了.
主人要.
警醒防盜.

$^{361}$作為比喻.
免禮門徒.
警醒等候他再來.
門徒扮演著家主.
耶穌再來的時間.
就像強盜來的時候.
沒有人能夠知道和預料到.
耶穌突然.
就會回來.
40節說你們也要預備.
提防.
小偷要常作準備.
最近我在.
德國回來.
頭兩天我在羅馬.
意大利的地方.
聽說意大利的小偷是最多的.
於是整隊人.
我們大概有40人左右.
大家都非常小心.
我的同房.
他買了一個很貼身的包包.
還要放在衣服裡面.
我就沒有他那麼小心.
真的很小心.
所以我只是把袋子拉起來.
再穿上外套.
我覺得已經很小心.
我感恩我們這14天.
都沒有被人偷東西.
一來我們互相提醒.
而且我們集體行動.
令到警察.
都不會有機會.
有機可乘.
最後兩天.
導遊提醒我們.
他說旅行團到最後幾天.
因為大家太放鬆.
所以最容易會不見東西.

$^{401}$當時我們.
都特別警惕.
因為導遊再提醒我們.
當我們知道.
那個地方.
最多小偷的時候.
就更加警醒.
也不會被小偷.
那麼容易偷東西.
我們的生命也要那麼警醒.
我們的生命.
到最後.
都是最容易.
前功盡廢.
因為最容易放鬆.
聖經裡面也有幾個人物是這樣的.
例如所羅門王.
他是一個非常優秀的王.
然而他到了晚年.
他娶的外邦妻子.
讓他迷失信仰.
敬拜別神.
讓他陷入罪裡面.
希西加王.
他是神眼中.
看為好的王.
他也依靠神.
修復聖殿.
做城牆.
改善水源.
減稅.
他真是一個很好的王.
後來他竟然患了一個重病.
在重病的時候.
他就去求醫治.
神一直看他做的事情.
都是好的.
在神眼中看為好的.
因此神憐憫他.
再給他有十五年的命.

$^{441}$就在他生命到最後的十五年間.
他竟然毫無警戒之心.
展示他的寶藏.
軍備.
給巴比倫的使者去觀賞.
即是炫耀財富.
讓他的行為.
間接導致.
猶大國的毀滅.
因為他死了之後.
他十二歲的兒子.
嗎哪西.
即是他十五年間.
生了一個兒子.
因為他死的時候.
這個兒子十二歲.
這個兒子是一個邪惡的君王.
敬拜偶像.
屠殺先知.
讓猶大沉淪在混亂和罪惡之中.
以至這樣的結果.
其實因為希西加王.
生命到最後的十五年.
他沒有好好警醒.
也沒有好好管教他的孩子.
各位弟兄姊妹.
我們要警醒.
不要像希西加王那樣.
縱然他曾經是一個很好的王.
他也知道自己.
其實還有十五年命.
但是他不警醒.
我們未必有他的豐功偉績.
也不知道自己生命.
何時完結.
我們更加要警醒自己的生命.
因為我們真的不知道.
自己的生命何時完結.
也不知道耶穌何時回來.
退休的時候.

$^{481}$更加是我們最容易放鬆的時候.
所以我們更加要警醒.
謹慎自己的生命.
不要分碎.
第三點要和大家分享.
是忠心的僕人.
這個時候彼得問了一個問題.
41節彼得說.
主啊.
這個比喻是為我們說的.
還是為眾人呢.
彼得問.
僕人是指門徒.
還是廣泛的信徒呢.
彼得這樣問.
大概是他以為自己有一個特殊門徒的地位.
連耶穌後來也要束上帶.
前來侍候他.
耶穌下文的比喻.
就是要指著門徒.
要作忠心的管家.
否則他們就會失去原先的地位.
耶穌下文的比喻.
就是要指著門徒.
耶穌這樣回答.
第42節.
主說.
誰是那忠心有見識的管家.
主人派他管理家.
家裡的人.
按時分糧給他們呢.
一說到.
忠心有見識的管家.
就會知道.
用了舊約忠心管家的榜樣.
就是約瑟.
約瑟在主人的面前蒙恩.
事後主人派他去管理家.
主人派他去管理家務.
將一切的都交在他手裡.

$^{521}$自從主人派約瑟.
管理家務.
和他一切的所有的時候.
耶和華就因為約瑟的緣故.
賜予他埃及的家人.
賜福給他埃及的家人.
原來耶穌.
原來神.
是會因為管家的忠心.
而賜福給他的家庭.
而神的子民.
都有管家這個身份.
第43節.
主人來到.
看見僕人這樣行.
那僕人就有福了.
44節.
我實在告訴你們.
主人要派他管理.
一切所有的.
這兩句.
我剛才說和37節.
非常相似.
是要強調僕人的.
警醒和忠心.
要得到神的賞賜.
警醒.
不是空等的.
而是忠心主人所委託的.
事情.
主人不在的時候.
仍然會繼續去完成主人的工作.
派他管理一切.
所有的這一句說話.
不在於榮華富貴.
是在於參與.
在神的計劃裡面.
意思是耶穌要給.
他忠心的門徒.
一個的應許.

$^{561}$在神的國彰顯的時候.
有份在裡面.
第45節.
那僕人.
若心裡說.
我的主人必來得遲.
就動手打他的僕人和子女.
並且吃喝醉酒.
46節.
在他尚不到的日子.
不知道的是神.
那僕人的主人.
要來重重地.
處置他.
定他和不忠的人同罪.
原來怎樣去處理食物.
就反映著他們心裡面的.
對自己的地位評估.
一個按照.
按時分糧.
給人的管家.
明白自己只是一個僕人.
所以他們忠心等候主人回來.
但是一個吃喝醉酒的僕人.
將自己當作是主人.
所以他也不會期待.
他的主人回來.
大家不知道.
會不會想耶穌早點回來呢.
主耶穌再回來.
就是世界末日.
你想不想耶穌要回來.
我聽過很多次.
大家都是不想的.
大家不想耶穌回來.
因為有很多事.
很想做 很多事還沒做.
但是我們想想.
我們想做的事.
是為了神的國.

$^{601}$我們要檢視.
我們自己的心.
一個忠心的管家.
會按時分糧.
一個不忠心的僕人.
只會吃喝醉酒.
按時分糧.
就是要忠心於神委託給你的使命.
關心有需要的人.
將神所賜給你的.
你也分享出去.
盡力發揮自己的恩賜.
去祝福有需要的人.
不忠心的僕人.
回來審判他.
將心思意念放任在.
這個燕洛當中.
而且還要打他的僕人.
使女.
他不善待身邊的人.
第47節.
「僕人知道主人的意思.
卻不預備.
又不信他的意思行.
那僕人必多受責打」.
48節.
「惟有那不知道的.
做了多受責打的事.
必少受責打.
因為多給誰.
就向誰多取.
多託誰.
就向誰多要」.
知道自己的.
照明.
神給你很多東西.
或神放你在某一個地方.
其實是要你.
祝福那裡的人.
當你明白.

$^{641}$神要你做某件事的時候.
卻不好好警醒.
好好警醒.
把握機會的話.
你就要小心了.
因為神回來的時候.
祂會要你去交賬.
我認識一位朋友.
他告訴我.
原來他十年前已經.
他知道神要他.
去讀全職的神學.
但十年前.
其實這個心.
已經在他心裡很久了.
但他一直在逃避.
因為他一直在教會裡工作.
他也在教會裡成長.
他有很多掙扎.
神給他看到.
教會很多的需要.
同時他也看到很多人的黑暗面.
是一個很不容易的事.
但這個呼召.
在他心裡已經十年了.
相信每個人.
都有自己的時間.
但若果我們知道.
自己的呼召.
我們就要好好把握機會.
最近也有另一個朋友.
跟我分享.
他很清楚明白.
神要他做的事.
他很明白神的呼召.
是要關心某一個群體的人.
但是因為工作上的調動.
讓他不能繼續.
他的呼召.
因此他禱告神.

$^{681}$非常平安.
他決定要離開現在的工作崗位.
繼續等候神.
讓神帶領他.
讓他用不同的方式.
完成他的使命.
原來當我們清楚明白.
自己的召命的時候.
我們也不需要害怕自己會去到哪裡.
其實也不用說得那麼深奧.
我們只要在小事上.
忠心.
在我們已經有的位置上.
做好每一件事就行了.
這樣已經是一個忠心的管家.
若果神要你成為.
一個公司的老闆.
你就好好善待自己的員工.
讓他們也感受到基督的愛.
也要好好運用.
自己的管理能力.
在教會裡謙卑.
以自己的恩賜去服侍.
若果神要你在某一間公司裡工作.
你就忠心做好一個員工.
盡你的能力.
在你的位置上忠心.
若果神要你在家庭裡做父母.
你就好好做好一個榜樣.
將子女帶到神的面前.
若果神有恩典.
讓你的父母.
依然在生的時候.
我們就要好好孝順父母.
各位弟兄姊妹.
聖經裡說.
末世的時候.
民要攻打民.
國要攻打國.
多處必有饑荒地震.

$^{721}$且有好些假先知會起來.
謎或多人.
只是不法的事情會增多.
許多的人的愛心.
才會漸漸冷淡.
唯有忍耐到底的必然得救.
耶穌很快就會回來.
我們要警醒忠心.
不要以為耶穌尚未回來.
是擔言.
其實不是擔言.
是寬容我們.
祂不願有一人的沉淪.
因此我們要警醒.
在耶穌未回來之前.
我們要一直做忠心的管家.
一開始我都提到.
父母等待子女.
等待父母回來.
父母等待子女回來.
都是因為有對方.
在我們心裡.
但我們也要.
同樣地將神放在.
我們的心裡.
我們不要沉睡.
當主耶穌回來的時候.
當日.
未來主耶穌回來的時候.
我們都可以一同共享.
警醒僕人的筵席.
讓我們一起祈禱.
親愛的天父.
為我們釘在十字架上.
換取我們有新的生命.
如今我們活著.
都有盼望.
你也非常信任我們.
將傳福音的重任.
交託給我們.

$^{761}$求你幫助我們.
成為警醒忠心的僕人.
在這個世界上.
不要沉睡 活得沒有盼望.
祂是要常常記住我們的使命.
要在小事上忠心.
做一個好的見證.
也顧及身邊的人.
將我們所有的都分享出去.
不要只顧自己的享受.
珍惜我們在地上的每一天.
奉主耶穌基督的名.
求.
阿們.
\newpage



\section{列王記上 18:1-6}
\label{sec:i86dI_9QW94}
\textbf{2023-06-25 加恩樓《掌管萬有的神》列王記上 18\_1-6 講員:陳一華牧師}
\newline
\newline
連結: \href{https://youtube.com/watch?v=i86dI_9QW94}{\texttt{https://youtube.com/watch?v=i86dI\_9QW94}} ~~~~ 語音日期: 2023-06-27
\newline
\newline
\hyperref[sec:keP18GcFGPs]{\small{< < < PREV SERMON < < <}}
~
\hyperref[sec:index_chronic]{\small{[返順時目]}}
~
\hyperref[sec:index_scriptual]{\small{[返順卷目]}}
~
\hyperref[sec:tu4SyHPIS_U]{\small{> > > NEXT SERMON > > >}}
\newline
\newline
列王記上 18:1-6
\newline
\begin{longtable}{cl}
\hline
\hline
章節 & 經文 (和合本修訂版)\\
\hline
18:1 & \begin{tabularx}{0.7\textwidth}{X} 過了許多日子，到了第三年，耶和華的話臨到以利亞，說：「你去，讓亞哈看見你，我要降雨在地面上。」 \end{tabularx} \\ \\ \relax
18:2 & \begin{tabularx}{0.7\textwidth}{X} 以利亞就去，要讓亞哈見到他。那時，撒瑪利亞的饑荒非常嚴重。 \end{tabularx} \\ \\ \relax
18:3 & \begin{tabularx}{0.7\textwidth}{X} 亞哈召來他的管家俄巴底。俄巴底非常敬畏耶和華。 \end{tabularx} \\ \\ \relax
18:4 & \begin{tabularx}{0.7\textwidth}{X} 耶洗別殺耶和華先知的時候，俄巴底把一百個先知藏了，每五十人藏在一個洞裡，拿餅和水供養他們。 \end{tabularx} \\ \\ \relax
18:5 & \begin{tabularx}{0.7\textwidth}{X} 亞哈對俄巴底說：「我們要走遍這地，到一切水泉旁和一切溪邊，或者能找到青草，可以救活馬和騾子，免得喪失一些牲畜。」 \end{tabularx} \\ \\ \relax
18:6 & \begin{tabularx}{0.7\textwidth}{X} 於是二人分地巡查，亞哈獨自走一路，俄巴底獨自走另一路。 \end{tabularx} \\ \\
[1ex]
\hline
\hline
\end{longtable}
$^{1}$弟兄姊妹早安.
我今天一進來嘉欣樓.
就聞到很濃郁的香味.
這個香味除了是咖啡的香味之外.
因為我上次來都沒有.
都是最近才有的.
除了咖啡的香味之外.
更加重要的是.
我聞到弟兄姊妹基督徒的香氣.
因為基督徒的香氣很重要.
我們的說話,態度,行動都很重要.
所以我見到弟兄姊妹.
見到牧師出現的時候.
很熱切的招呼.
我立即好像被愛心融化了.
我也很感受到弟兄姊妹.
你們的聊天,分享,交通.
真的有那種基督的香氣.
今天來到你們當中.
給我一個很深刻的經歷.
今天牧師和大家去看的經文.
屬於冷門的.
弟兄讀出來的時候.
大家對那段經文是否都沒有印象.
好像冷門的.
是真的,比較少接觸到.
但這段經文的背景是什麼呢?.
就是當時以色列國出現一個王.
這個王是相當令神失望的.
為什麼呢?.
因為他追隨他先前的王一樣.
帶領以色列人神的子民離棄了真神.
去敬拜偶像.
而且在國家裡面建造很多祭壇.
很多偶像給以色列人去敬拜.
換句話來說.
陷以色列人在深淵裡.
在這樣的光景裡.
大家就知道.
上帝作為以色列人天上的爸爸.

$^{41}$以色列人是地上的子女.
是神的子民.
一定不能夠有目無睹.
一定不能夠當作沒事發生.
所以在一個這樣最低谷的時候.
或者在以色列國歷史最黑暗的時期.
上帝就出手了.
上帝怎麼出手呢?.
上帝出手就發出一個大自然的災難.
去提醒以色列人這是什麼災難.
剛才你讀了幾節經文.
你的印象這是一個什麼災難?.
這是一個旱災.
大家有沒有留意到?.
即是天不下雨.
這種情況今時今日.
我們見到有些非洲國家沒有雨水.
淒不淒涼?.
真的很淒涼.
上帝就用一個這樣的大自然災難.
去提醒祂的子民.
各位電影節目.
在一個這樣的背景下.
今天這段聖經雖然短短的.
只有六節.
但我想告訴大家.
起碼我們見到上帝在整件事當中.
祂是掌管萬有的.
祂掌管起碼有四件事你可以見到.
第一 上帝掌管什麼呢?.
上帝是掌管大自然的.
你要記得.
上帝不單是創造大自然.
大家都很明白.
而且祂也掌管大自然.
換句話說.
祂要天不降雨.
那天是什麼?.
沒有雨水.
如果祂要天降雨.

$^{81}$那天一下雨.
那地就得到滋潤.
這個現象是表現什麼呢?.
是要讓儀式的人知道.
哪一位才是真創造者.
哪一位才是真神.
有一個背景是這樣的.
就是當時以色列人.
拜的其中一個偶像.
一個當時在迦南地很出名的偶像.
他的名字叫巴力.
大家聽過沒有?.
有沒有聽過這個名字?.
巴力是什麼呢?.
原來在以色列族裡面.
或者在迦南族裡面.
巴力這個偶像這個神明.
是當時告訴大家.
是要掌管什麼呢?.
掌管大自然.
大家明白這個意思吧?.
原來他們心目中的巴力是很威的.
很有實力的.
掌管大自然的.
要風得風.
要雨得雨.
是一個這樣的背景.
現在上帝就跟以色列人說.
我要正面來給大家看到.
哪一位才是真神.
哪一位才是掌管大自然的上帝.
所以就跟他有一個正面的比較.
大家看不看到?.
所以在以色列人遇到旱災的那兩年.
怎樣呼求巴力?.
怎樣懇求巴力?.
怎樣大聲疾呼巴力?.
有沒有回應?.
沒有回應.
天不下雨.

$^{121}$這個就是第一.
要讓以色列人知道.
原來上帝是那位真正掌管大自然的上帝.
大家看到嗎?.
這是第一.
第二.
上帝掌管什麼呢?.
上帝是掌管時間.
讀這段經文.
有沒有發覺有一個時間彈了出來?.
有沒有印象?.
如果剛才留意的話.
第幾年?.
第三年.
有沒有留意到?.
第三年.
上帝就出手了.
有時候我都會想.
為什麼上帝的時間這麼有趣.
他又不是第一年.
又不是第二年.
偏偏是第三年.
又不是第四年.
大家有沒有留意到?.
就是說上帝要讓以色列人知道.
他是掌管時間的上帝.
各位電視節目.
第三年是有些意思的.
什麼意思呢?.
就是說以色列人第一年沒有雨水.
能不能撐下去?.
能撐下去.
因為還有老本.
第二年能不能撐下去?.
都怕撐不下去了.
起碼沒有今早的麵包那麼好吃.
沒有大麥沒有小麥.
都難撐下去了.
第三年也是滴水不下的時候.
慘了.

$^{161}$那塊地三年都沒有東西出.
電視節目.
以色列人就開始進入一個嚴重的饑荒階段.
我相信都會死不少人.
正如大家有時看新聞.
看到一些國家沒有雨水.
大饑荒.
死很多人.
不過有些國家禁止.
沒有把數字告訴大家.
大家都明白我說什麼了.
原來上帝是很奇妙的.
他是掌管時間的.
他說第三年有.
就是什麼?.
第三年開始有了.
各位電視節目.
你和我.
有時候都很.
怎麼說呢?.
我們都有自己的主觀決定.
你有沒有留意到?.
我的計劃.
我的時間.
我的安排.
應該是什麼?.
最fit.
最好.
最沒得頂了.
誰不知人算不如天算.
我遇到一些弟兄姊妹.
很想他父母信耶穌.
祈禱.
傳道.
帶他們回教會.
一年又一年又一年.
都不肯信耶穌.
有沒有這樣的情況?.
有.
但很奇妙.

$^{201}$上帝的時間一到.
他就是老人家的心軟了下來.
願意接受耶穌基督.
有沒有聽過這樣的見證?.
弟兄姊妹們.
是不是很奇妙?.
上帝的時間.
難道我們控制的?.
我們認為最好的時間.
上帝認為未必是.
我們認為最差的時間.
但上帝會告訴你.
這個是最好的時間.
去年香港疫情很嚴重.
我一個很深刻經歷上帝的時間.
準到無法忍受.
因為我去年年尾的時候.
答應了水警總部.
有一個報道大會.
即是說整個水警.
他們籌備了一年.
希望有這個報道大會.
我就被邀請.
去到報道大會之前的星期.
就發生事了.
當我去到星期四的時候.
心想星期日就是.
全港水警的報道大會.
但是星期四發生什麼事呢?.
大家認不認識梁傳道?.
聽過梁傳道沒有?.
沒聽過嗎?.
以前錦繡的童工梁月華傳道.
不好意思我把他的名字也說出來.
我不是說他.
我等一下要告訴你.
我要稱讚他.
認不認識梁傳道?.
他星期四原來已經開始確診.
但他星期四還仍然上班.

$^{241}$還跟我面對面.
面對面相處一個半小時.
到星期五他給我電話.
牧師不好意思我確診了.
你說對我大不大打擊?.
我星期日就說到.
他星期五就確診了.
所以我星期五就趕快去中心做檢測.
做完檢測星期六通知我.
原來我沒事.
即是核酸檢測告訴我沒事.
是不是Hallelujah?.
星期六沒事.
星期日可不可以去?.
可以嘛.
精彩的不是這件事.
精彩的是什麼?.
星期六晚上十點鐘.
我的工人跟我說.
Sir我中招了.
哎呀我確診了.
我差點說有沒有搞錯?.
這個時候才告訴我他中招了.
我就變成一個密切接觸者.
密切接觸者那時候可不可以出街?.
不可以出街.
星期六晚上告訴我他確診了.
我星期日就去港島.
你說我是不是雙重打擊?.
真的好奇妙.
上帝的時間是妙到你拍案叫絕.
我就打去衛生署問.
我的工人確診了.
其實我都掙扎打不打好.
你明不明白我的意思?.
如果我不打.
我星期日就捐出去.
沒有人知道.
然後港島回家.
我可以這樣做.

$^{281}$但我心裡不太平安.
結果我打去衛生署.
1510告訴職員.
你知道他最後問我一句什麼話嗎?.
他說先生.
你最近有沒有做過核酸檢測?.
我說有啊有啊.
剛剛昨天星期五.
我才拿到什麼?.
星期六拿到報告.
我星期五做.
星期六拿報告.
說我什麼?.
沒事.
這位職員說.
因為你星期六有報告沒事.
你明天可以出街.
星期日你可以出街.
你說奇不奇妙?.
各位節目.
我們張榮耀歸給天父.
所以你說我是不是要請梁傳道喝茶?.
我原先都有點生氣.
你流鼻水.
你不要上班.
星期四還上班.
但心裡想想.
上帝真的很奇妙.
如果不是他中招.
我就出不出街?.
我出不出街.
上帝的時間奇不奇妙?.
很奇妙.
人算不如天算.
底子沒活.
你記得.
你的時間怎麼好都好.
都不及上帝的時間更好.
因為上帝說.
我的道路高過你的道路.

$^{321}$我的意念高過你的意念.
所以底子沒活.
我們要學習信服上帝的時間.
第三.
上帝掌管什麼?.
上帝掌管人事.
你有沒有留意到一個這麼惡的王.
一個十惡不赦的王.
阿俠.
不敬畏上帝.
作惡多端.
你讀過他的歷史.
竟然上帝在這個王的身邊.
安插了一個人.
他的名字叫做俄巴底.
大家有沒有聽到?.
俄巴底.
上帝找他做什麼?.
上帝安插這個人.
竟然做了阿俠王的大內總管.
現在叫大內總管.
你說奇不奇妙?.
底子沒活.
竟然加一個這麼不敬畏上帝.
這麼壞的王.
但上帝的作為.
真的令我們非常驚訝.
他安排一個敬畏神的俄巴底.
做了阿俠的左右手.
你說上帝的掌管是不是很奇妙?.
底子沒活.
上帝是什麼?.
掌管人事.
所以各位底子沒活.
你的職位無論升無論降.
你的職位無論在哪個地方的姿勢.
請你記得.
如果你是上帝的兒女.
上帝一定有最好的安排.
特別有時候上帝很奇妙.

$^{361}$他要安排幾個基督徒.
他的兒女都在同一個辦公室出現.
你有沒有聽過這樣的見證?.
奇不奇妙?.
很奇妙.
我也說一個見證給你聽.
大家聽過天水圍醫院.
是不是?.
四年前新開的天水圍醫院.
沒有人想到這間醫院一開始的時候.
不是拜四個國.
這間醫院一開始的時候.
這間政府的醫院一開始的時候.
竟然是用一個祈禱感恩.
做開始這間天水圍的醫院.
奇不奇妙?.
很奇妙.
為什麼?.
我想告訴大家.
因為上帝是掌管人事的.
想不到當時候天水圍的醫院.
院長,副院長.
即是副行政總監.
護士的頭.
行政的頭.
四個都是基督徒.
Hallelujah.
加上院務五個.
五劍合璧.
天下無敵.
各位弟兄姊妹妹.
沒有想到有一間醫院.
最高層的幾位.
同一時間都是基督徒.
所以天水圍醫院的開始.
是由祈禱開始的.
各位弟兄姊妹妹.
即使你去到天水圍醫院探病.
你想去飯堂的地方.
你看到旁邊有一個很大的牆壁.

$^{401}$竟然掛了一節金句在那裡.
奇不奇妙?.
各位弟兄姊妹妹.
因為上帝是掌管人事的.
各位弟兄姊妹妹.
所以今天無論你在什麼公司.
在什麼崗位.
無論什麼職位.
高低.
如果你是神的兒女.
你信服上帝.
上帝安排,差派你去到那裡.
一定有祂的什麼?.
有祂的心意.
有很多人都不喜歡調去崗位.
叫做什麼位呢?.
叫做豬頭骨位.
為什麼?.
因為很難頂的.
但是若果我們願意信服上帝的安排帶領.
上帝一定有祂的心意在當中.
最後牧師和大家說的就是.
在這段經文當中.
有第四個掌管.
你知道是什麼嗎?.
上帝是掌管命運.
各位弟兄姊妹妹.
經常聽人說.
我的命真的好什麼?.
不是很棒,是好什麼?.
很棒,相反是什麼?.
很苦啊.
我真的命苦啊.
有沒有聽過這樣說過?.
有些不信的.
有些迷信的.
就拿自己的生辰,八字.
去計算這個計那個.
各位弟兄姊妹妹.
有很多人覺得自己是生不逢時.

$^{441}$是不是?.
有很多人覺得自己的命運是什麼?.
遭遇比其他人都慘.
有沒有這樣的人?.
有沒有?.
有,我想告訴大家.
若果你是神的兒女.
你的命運不會出於偶然.
因為上帝已經在你人生的際遇上.
他掌管了什麼?.
掌管了一切.
你想想.
屬神的先知.
好生不生.
就生在亞俠的年代.
你說是不是生不逢時?.
還有一個耶駒別.
你不要忘記.
還有一個皇后比亞克還厲害.
不知道多少,二十倍.
是作惡多端的.
離棄上帝的.
如果你做神的先知.
做神的僕人.
生在那個年代.
真的是什麼?.
生不逢時.
可以說是什麼?.
人生的路上.
是九死一生.
差不多死定了.
沒有機會回頭.
奇怪.
既然上帝的僕人.
上帝的先知.
在這麼惡劣的情況下.
出現了一個人.
叫做俄巴底.
剛才不是說過.
上帝所揀選的俄巴底.

$^{481}$很有聰明智慧.
將一百個先知.
上帝的先知.
五十個人一組.
藏在山洞裡.
你有沒有留意到?.
五十個一組.
然後每一天.
供應食物和水給這些先知.
保留他們的生命.
不至於死.
也逃避了亞克和耶駒別的追殺.
各位聽眾.
如果當你看到這一節經文的時候.
你覺得.
命運是由誰操縱?.
命運不是我們操縱.
你不要覺得自己命苦.
我怎樣奮鬥.
怎樣努力.
我都沒有出息.
有些人是這樣的.
你有沒有遇到這樣的人?.
我無謂去任何的.
這麼勤力工作.
為什麼?.
因為我的命是這麼苦.
沒有希望.
底子沒有門.
如果你這樣想.
或者你有些朋友.
家人這樣想.
你告訴他們.
你想錯了.
為什麼?.
因為我們的命運.
不是操縱在自己手上.
我們的命運.
我們的際遇.
操縱在什麼?.

$^{521}$上帝的手中.
上帝掌管一切.
所以傳道書說得很好的.
有時候我們人生的路程上.
以為這樣一定可以.
以為一定自己成功.
不是.
原來是很在乎.
當時候上帝有沒有預備機會給你.
如果你預備到有這個機會.
你就發覺上帝將你的人生際遇.
可以改變過來.
所以當我們見到這段聖經的時候.
各位弟兄姊妹.
今日牧師跟大家分享.
互相提醒.
記得.
你有一位這麼有能力.
這麼偉大的天上的爸爸.
做我們的父親的時候.
人生路上你所遇到的逆境困難.
有何懼之有呢?.
大家認不認同?.
為什麼?.
因為你是屬於上帝.
最重要的是你屬於上帝的時候.
你無論遭遇任何的境況.
你的際遇.
你走的路.
上帝都會看顧著.
上帝都會帶領著.
上帝都會將你的人生.
完成了一個更美好的計劃.
所以弟兄姊妹們.
讓我們大家努力.
活得精彩.
活得更加的憧憬.
更加為神作供.
作見證.
榮耀上帝我們低頭祈禱.

$^{561}$親愛的主.
一段很短的聖經.
亦都很陌生的經文.
卻是穩陷著上帝你很豐盛的教導.
你給我們有很多屬靈的啟發.
屬靈的帶領.
願意我們的人生活得精彩.
我們的人生更加活得有挑戰.
能夠榮神益人.
被你所使用.
能夠祝福更多的人.
祈禱奉耶穌名求.
阿們.
\newpage



\section{路加福音 19:12-23}
\label{sec:tu4SyHPIS_U}
\textbf{2023-06-25 早堂《未來的風險》路加福音 19\_12-23 講員:陳韋安牧師}
\newline
\newline
連結: \href{https://youtube.com/watch?v=tu4SyHPIS_U}{\texttt{https://youtube.com/watch?v=tu4SyHPIS\_U}} ~~~~ 語音日期: 2023-06-27
\newline
\newline
\hyperref[sec:i86dI_9QW94]{\small{< < < PREV SERMON < < <}}
~
\hyperref[sec:index_chronic]{\small{[返順時目]}}
~
\hyperref[sec:index_scriptual]{\small{[返順卷目]}}
~
\hyperref[sec:WZaLHtVToyI]{\small{> > > NEXT SERMON > > >}}
\newline
\newline
路加福音 19:12-23
\newline
\begin{longtable}{cl}
\hline
\hline
章節 & 經文 (和合本修訂版)\\
\hline
19:12 & \begin{tabularx}{0.7\textwidth}{X} 說：「有一個貴族往遠方去，為要取得王位，然後回來。 \end{tabularx} \\ \\ \relax
19:13 & \begin{tabularx}{0.7\textwidth}{X} 他叫了自己的十個僕人來，交給他們十錠銀子，說：『你們去做生意，直到我回來。』 \end{tabularx} \\ \\ \relax
19:14 & \begin{tabularx}{0.7\textwidth}{X} 他本國的百姓卻恨他，打發使者隨後去，說：『我們不願意這個人作我們的王。』 \end{tabularx} \\ \\ \relax
19:15 & \begin{tabularx}{0.7\textwidth}{X} 他得了王位回來，就吩咐叫那領了銀子的僕人來，要知道他們做生意賺了多少。 \end{tabularx} \\ \\ \relax
19:16 & \begin{tabularx}{0.7\textwidth}{X} 頭一個上來，說：『主啊，你的一錠銀子已經賺了十錠。』 \end{tabularx} \\ \\ \relax
19:17 & \begin{tabularx}{0.7\textwidth}{X} 主人對他說：『好，我善良的僕人，你既在最小的事上忠心，你有權柄管十座城。』 \end{tabularx} \\ \\ \relax
19:18 & \begin{tabularx}{0.7\textwidth}{X} 第二個來，說：『主啊，你的一錠銀子已經賺了五錠。』 \end{tabularx} \\ \\ \relax
19:19 & \begin{tabularx}{0.7\textwidth}{X} 主人也對這個說：『你管五座城。』 \end{tabularx} \\ \\ \relax
19:20 & \begin{tabularx}{0.7\textwidth}{X} 又有一個來說：『主啊！看哪，你的一錠銀子在這裡，我把它包在手巾裡存著。 \end{tabularx} \\ \\ \relax
19:21 & \begin{tabularx}{0.7\textwidth}{X} 我向來怕你，因為你是嚴厲的人：沒有放的，也要去拿；沒有種的，也要去收。』 \end{tabularx} \\ \\ \relax
19:22 & \begin{tabularx}{0.7\textwidth}{X} 主人對他說：『你這惡僕，我要憑你的話定你的罪。你既知道我是嚴厲的人，沒有放的也去拿，沒有種的也去收， \end{tabularx} \\ \\ \relax
19:23 & \begin{tabularx}{0.7\textwidth}{X} 為甚麼不把我的銀子存在銀行，等我來的時候，連本帶利都取回來呢？』 \end{tabularx} \\ \\
[1ex]
\hline
\hline
\end{longtable}
$^{1}$定智妹早晨.
今日我們會一起思考我們的未來.
藉著這篇訊息我們一起去想想我們的未來.
不單單是我們自己.
我們香港的未來.
我們這個世界的未來.
特別這幾年裡面.
大家都知道很多朋友親戚都移民.
移民言下之意.
其實可能對香港的未來.
似乎就不是太過明朗.
所以就離開.
就算我們今天不說香港社會也好.
今天我們面對自己的生命的時候.
想一下自己幾歲.
未來這十年之後怎樣.
想起都會有一點點.
都不容易想的這個未來.
所以今天我們一起來思考我們前面的未來.
不知道當中有沒有弟兄姊妹.
二十多年前回錦村.
我記得我小時候.
剛剛是2000年的時候回錦村的時候.
記得我們那時候是1999年除夕倒數.
有沒有去過這裡.
那時候我們還在幼稚園那邊.
一起一班弟兄姊妹.
倒數迎接2000年的時間.
不知道大家有沒有人在場.
記得那時候其實.
不知道大家都是這樣.
那時候2000年倒數.
是心裡面有一種興奮.
又少少的擔心.
我們那時候還記得2000年的時候.
有一樣東西叫做千年蟲.
那時候很奇怪.
很怕千年蟲令到我們世界沒有日子.
銀行系統電腦會壞掉.
那時候我們懷著一種心情.

$^{41}$迎接2000年.
特別是我們基督徒.
不知道你那時候2000年.
就算你不在錦村也好.
那時候你基督徒.
我不知道你是否在暗地裡有一份.
想過的.
2000年的時候會不會主在來呢.
你那時候倒數的時候.
5 4 3 2 1.
然後主就回來了.
有沒有想過這樣.
那時候都不敢跟別人說.
我這樣想是否會這樣.
但其實都有這樣的想法.
那時候我在錦村幼稚園.
當我們5 4 3 2 1.
Happy New Year的時候.
我都會看一看窗戶.
看看外面有沒有什麼發生.
所以我們對未來.
似乎是有一定的盼望.
但也有很多的未知之數.
很多的擔心.
所以究竟我們的未來是怎樣呢.
我們怎樣去看我們的未來.
無論是我們自己的未來.
香港的未來.
甚至是我們錦村家的未來.
所以我們就用這篇.
耶穌基督昔日在路加福音的比喻.
來一起思考.
一起去討論.
主你求你幫助我們.
讓我們在當中去思想.
你昔日對門徒的比喻.
讓我們能夠掌握裡面的意思.
能夠應用在我們當下生活的處境.
無論是我們自己的生命.
無論是香港.

$^{81}$無論是我們教會.
求主你讓我們能夠在當中學習.
你教導我們的功課.
求主幫助孩子.
孩子不配.
但願意你自己來宣講你的說話.
奉主明求 阿們.
今天經文是路加福音十九章十二到十三節經文.
經文如果看和本的標題的話.
就叫做十定銀子的比喻.
在這段經文有一段平衡的經文.
記載在馬太福音二十五章.
一個才幹的比喻.
你會發覺兩段經文其實都很相似.
差不多的故事橋段.
差不多的故事結構.
差不多的結局.
兩個很相似的比喻.
兩段經文都要說主人離開他的地方.
分別將他們的錢留給他的僕人.
叫他去投資.
回來之後第一個僕人賺了很多錢.
第二個僕人也賺了不少錢.
最後一個僕人收了主的錢.
最後被主人責備.
又惡又懶的僕人.
耶穌稱讚那些做生意投資得利的僕人.
叫做既良善又忠心的僕人.
兩個比喻都是差不多的套路.
不過如果你去留心兩段經文的話.
你會發現兩段經文有不少不同的地方.
我今天嘗試去關心兩段經文不同的地方.
今天特別去選擇路加的版本.
是因為要說出一些很特別的東西.
基本上我們不詳細看.
但我們看到.
如果你看回今天所說的經文的時候.
路加福音的故事和馬太的版本.
兩個人的僕人用錢的方法.
和他們所賺錢的回報.

$^{121}$是一個很特別的不同.
大家可能都知道投資.
我們說的reward ratio.
即是回報率是多少.
你會發現路加福音的回報率.
即是投資的回報率.
是對比起馬太是瘋狂地高.
大家還記得嗎.
馬福音那個是五千銀紙.
二千銀紙.
一千銀紙的比喻.
一倍的比喻.
即是一倍的回報率.
五千就賺五千.
二千的就賺二千.
一千的就賺一千.
不過相比之下路加福音的回報率是瘋狂地高.
是多少.
是一定銀紙就賺十定銀紙.
大家不知道有沒有這樣的概念.
這是一個一倍十的概念.
即是基本上你用到一萬元.
可以賺到十萬元.
這樣的概念.
一個非常非常瘋狂的回報率.
就算是路加福音所說的第二步人也好.
一定銀紙賺到五定.
其實都是五倍的回報率.
一個非常非常不尋常的回報率.
大家銀行修飾其實都是最多幾厘.
但你說的是一個五倍十倍的概念.
所以要問究竟怎樣才是.
為何路加福音所說的回報率這麼大.
一倍十一倍五的概念.
不知道大家有沒有這樣的概念.
我們明白十倍利潤是什麼意思.
但你知道十倍利潤即是言下之意是什麼意思.
你要承擔一倍十的風險.
大家都很明白.
回報越多風險就會越大.

$^{161}$你所承擔的風險就越高.
利潤越高風險就越大.
這是一個世間的定律.
無論你說的是股票期權大小.
甚至乎任何的都是一個這樣的定律.
你要去承受更多的獎勵的時候.
就需要去承受相對更多的風險.
這幾年大家的NPF都知道.
都已經虧蝕了不少.
大家知道你越進取.
虧蝕就越進取.
所以這都是一個世間定律.
路加福音似乎是說一個賺十倍賺五倍的股票.
但這也是在承受十倍和五倍風險的.
十倍風險是什麼一回事呢.
大家都不會賭錢.
但你也知道如果你看回一些足球的話.
一倍十是什麼概念.
如果大家有看波弟兄就知道.
一倍十大概是什麼呢.
大概就是我之前也看過一些網站.
一倍十是什麼呢.
就是哈佛柏林對拜博尼克.
你也明白如果看波的話就知道.
哈佛柏林贏到拜仁大概是一倍十的概念.
或者大家不熟悉的可能看世界盃的時候.
英格蘭對伊朗那場.
伊朗能夠贏英格蘭大概就是一倍十的概念.
可能你不看波弟兄也知道我意思.
大概是一個又不能說贏的.
英格蘭也有時候會爆爛的.
但你會發覺大概就是這樣.
也挺難贏的.
贏了就賠十倍給你.
不是不可能.
但是有些困難.
這個就是十倍風險的意思.
原來聖經強調耶穌所稱讚的.
忠心良善的僕人.
這班能夠一錠銀紙賺十錠銀紙的僕人.

$^{201}$正正都是一班承受著一倍十.
十倍風險的人.
真人沒有一個上帝所交託給我們的責任.
是不需要承擔風險的.
每一件美事的事情.
你都一定要承受一定程度的風險.
憑著信心去做.
任何一個向上主負責的人.
無論是基督徒,教會領袖.
都是要為一些美善的事情.
去承擔一定程度的風險.
尤其是今天的社會.
每一件的事情,每一件好事.
其實都是有一定的風險的.
教會要做一件應該要做的事情的時候.
都是要承擔一定的風險.
以前我們在國內的短訊也都明白.
你要帶聖經上去.
或者你要做一些聚會.
都可能有一定的風險要承擔.
相反耶穌譴責那些自以為忠心的僕人.
他們以為自己能夠守得住家業.
可以很安全.
Play safe.
這樣去做一個忠心的人.
耶穌是譴責這班人.
當我們去留心這個比喻的時候.
聖經怎樣去描述這班被斥責的僕人.
聖經說當主人去讚頭那兩個僕人的時候.
第三個僕人就馬上去做一件事.
他說主啊你看.
你看那一錠銀子.
我一直保存在手筋裡面.
這是一個很有圖畫的畫面.
他可能馬上藏在一個洞裡.
然後拿出來.
然後用手筋包得很緊.
拿出來.
主啊這個就是你的錢.
他很安全地放在隔萬裡.

$^{241}$用手筋包住.
用一條手筋包住.
既安全又穩陣.
感覺非常良好.
不怕弄濕不怕弄髒.
不怕弄得不好不見.
你問這個僕人細不細心.
非常細心.
你問這個僕人是不是很穩陣.
非常穩陣.
這個人做事可不可靠.
都是很可靠.
但這個不是天國福音的態度.
甚至乎聖經是稱呼這個僕人做惡僕.
一個wicked servant.
我們問為什麼一個穩陣又安全的人.
為什麼會被主人稱之為惡呢.
為什麼用手筋包住錢.
那麼好保存固息的人.
為什麼會稱之為惡僕呢.
這個人是多麼的caring.
多麼的謹慎.
多麼的愛惜主人的錢財.
多麼的想虧了主人的錢.
究竟這個人惡在哪裡呢.
答案非常簡單.
這個僕人沒有回應主人的要求.
這個正正是耶穌基督對我們的要求.
說真一句.
在比喻裡面.
這個主人的要求是挺高的.
甚至乎是有點harsh.
說我原是怕你.
因為你是一個嚴厲的人.
即是這個僕人都清楚的說.
這個主人是一個嚴厲的人.
所謂嚴厲.
其實是一個挺harsh的人.
事實上連主人自己都承認.
他是一個要求很高的人.

$^{281}$主人的要求有多高呢.
主人要求的僕人沒有放下的.
還是要去拿.
沒有種下的還是要去收.
主人要求的僕人.
要求他們追求一個這樣的風險.
來承擔一個高的匯報.
要求他們面對一個這麼大的風險.
主人說你明知我harsh你都不去做.
親子弟弟好姐妹.
我們要這樣去說.
對於我們基督徒來說.
承受風險.
憑著信心.
迎向未來.
承擔一定程度的風險.
不是一個option來的.
也不是一個純粹的選擇.
這是耶穌基督對我們每個人的一個命令.
你為了福音為了上帝.
為了你自己面對將來的時候.
你要承擔一定程度的風險.
先不說早期教會的先賢先成.
他們怎樣來承擔風險去傳開福音.
也不說我們200年前帶一艘船來中國報道的風險.
作為基督徒你或多或少都要學習.
一定程度的信心.
承擔一定程度的風險.
去迎面我們的未來.
一個或多或少的不確定性.
我們坦白說我們不知道我們前面的前景.
無論是我們自己的生命.
今日的香港.
甚至今日錦繡堂.
我們面對的前面.
但我們知道的就是我們始終都要承擔一定程度的風險.
來走下去.
今日這個社會裡面.
很多時候我們都習慣將風險轉嫁給別人.
做了自己的部分.

$^{321}$那就叫做很安全.
我就不承擔風險.
那份過於謹守.
過於怕什麼.
然後place.
這樣的文化.
但這個不是我們作為教會.
應該要有的態度.
教會作為上帝的一個僕人的時候.
我們應該一定程度要去被要求承擔這個命令.
承擔上帝耶穌基督給我們的風險.
很多人都以為面對風險.
第一個反應就是怎樣去避免它.
他們以為面對風險要做的事情.
就是要避免風險.
我們都覺得我們似乎是面對著很多這樣的事.
我們總之是不要惹禍上身.
總之就做好一些公關.
做好一些一切.
這樣就叫做好.
很真的.
耶穌基督這個銀支給我們.
告訴我們風險是必然存在的.
問題關鍵不是你如何去避免風險.
而是你為著誰去承擔風險.
事實上面對著未來.
風險是必然存在.
沒有一個沒有未來的風險.
做事有風險.
不做事也有風險.
進取有風險.
其實play safe也有風險.
不過我們看不到那些沒有做事的風險.
風險是你選擇承擔什麼風險.
站在什麼角度去理解風險.
如果我們詳細去探討風險這回事的話.
風險可以大概分為三類.
看很多恩人這樣說的時候.
一個就叫做已知風險.
風險就是出現在一些既定的規律裡面.

$^{361}$作為一個已經知道多少風險你會承擔到的.
譬如大小風險就是多少.
50 50這樣的情況.
這個是已經知道的風險.
第二就是一些未知的風險.
這個風險的規律你是不知道的.
不能夠準確地告訴你.
很多人只是關注於一些已知的風險和未知的風險.
他們以為面對著已知和未知的風險.
最後的方法就是play safe.
做多少風險的時候.
多做事就多風險.
少做事就少風險.
不做事就沒有風險.
不過教會往往是規定了第三類的風險.
就是一個叫做隱藏的風險.
隱藏風險是隱藏的.
是長遠的.
是慢性的.
是未來的.
隱藏風險是在今天我們純粹是平庸無為的背後.
當教會以為我們要保持穩定.
先不要亂來.
不要錯的時候.
隱藏的風險其實是靜靜地不斷地去累積.
當教會沒有做一件要去做的事情.
當教會沒有做當下要做的最好的事情.
不去追求最好的時候.
選擇不斷地去play safe的時候.
不斷地去累積自己應該看不到的風險.
即是在惡毒上都是這樣.
用自己的手帕包著自己主人的錢的時候.
這個似乎是最play safe的方法.
這個最安全的方法.
當下來說.
但是當他一直用手帕包著的時候.
他就不斷地去累積著一個慢慢長遠的隱藏的死亡.
大家不要誤會.
大家不要覺得這樣說純粹是叫大家去搏.
純粹是大膽地去不要理風險.

$^{401}$或者是不斷地去冒險.
不要理這樣.
我們仍然是要去謹慎風險的問題.
仍然要好好去管理風險.
今天都很流行這個詞.
Risk management.
我們要認清我們面對的是一個什麼風險.
為了誰而去冒這個風險.
事實上今天有很多人都說.
有一樣風險絕對是要避免的.
就是所謂的滅頂風險.
滅頂之災的風險.
即是你一鋪身家過大海.
賣大伯加諾一賣就贏.
就是輸的那種.
炒股票一炒就新身家那種.
這些風險是一定要避免的.
因為這些你是沒有機會再翻身的.
一鋪清袋再沒有翻身的機會.
這個我們一定要去避免風險.
將來你不能夠再去有任何機會翻身的時候.
這個我們一定要避免.
不過對我們基督徒來說.
世界上唯一不會出現滅頂之災風險的是什麼.
我們任何的投資.
任何的投資的事項.
都可以令你一鋪就輸身家.
唯一一樣不會面對滅頂之災風險的是什麼.
正正就是上帝的國度.
上帝的國度不會給你一下.
給我們教會一下就輸清光.
任何的事情都可以面對一鋪清袋的風險.
任何大的股票.
騰訊任何的東西都可以突然間會沒有了.
但是上帝的國度不會因為某一個人.
或者某一班人所做的一個錯誤決定.
令到他會有任何的滅頂之災.
所以天國的事情我們是要承擔風險的.
都值得我們來承擔這個風險.
我們面對著這樣的思維的時候.

$^{441}$我就知道我們怎樣來做教會.
我怎樣來思考我們教會.
怎樣來承擔我們的將來.
怎樣來將自己的將來參與在上帝的國度裡面.
天國的風險是值得我們為主去承擔這個風險.
大家不要忘記.
老國福音這個風險的比喻是一個關於什麼的比喻.
這個耶穌基督接受皇位.
上帝國度得勝的比喻.
另外一個福音和這個馬福音很不同.
在老國福音的比喻裡面.
就是僕人用錢的比喻裡面.
有一條很不自然的支線說出來.
好像不是很關事的.
明明就是投資.
但是背景是什麼.
是說主人去接受皇位.
這個似乎不是很有關係的背景.
主人要接受皇位.
離開他的家.
這個正是告訴我們.
在我們承擔風險的時候.
其實那幅大圖畫.
那幅長期的圖畫是在畫什麼.
上帝的國度正要降臨.
主人正是去接受皇位.
這個是一個正要勝利的圖畫.
耶穌說因為他快要到耶路撒冷.
就以神的國顯現.
就說有一個出身顯貴的人.
往遠方去接受皇位.
然後才會來.
你發現和馬太很不同.
路加坊裡面的僕人的錢.
無論是賺還是虧都好.
和主人得皇位.
似乎好像沒有什麼關係.
無論這幫僕人承受多大的風險.
無論這幫僕人是賺還是虧.
主的國度都必然會降臨.

$^{481}$總之就是贏.
所以聖經說這個主人很harsh.
因為這個主人就是要贏.
還要是他一定會贏.
所以這個主人才吩咐僕人.
一賠十.
英格蘭對伊朗.
你買伊朗.
哈佛柏林對拜莫尼克.
你買柏林.
當然不是叫大家賭博.
是賭注.
為我的緣故.
不妨去承擔一定程度的風險.
因為我就是會贏的.
耶穌基督說我的國度必定會來臨.
在這個國度的將來下.
你的一個決定.
是正確也好.
錯誤也好.
最後是虧了也好.
賺了也好.
不會搞垮我的國度.
你要做的事是一件事.
就是去忠心良善的.
好好運用你現在有的東西.
為主去承擔一定程度的風險.
做一些美善的事情.
你應該做的事情.
憑著少少的信心.
和一定程度的風險去嘗試.
這樣去做下去.
忠心不是謹慎.
忠心是承擔一定程度的風險.
承擔上帝耶穌基督.
必勝的這個國度.
去承擔這個風險.
請問你今天怎樣去和你的未來共處.
面對著未來.
自己當下眼見的未來.

$^{521}$無論你的事業.
無論是教會.
無論是香港的未來.
我們要學習去勇敢.
忠心去承擔風險.
重點是你為誰承擔風險.
為著什麼東西承擔什麼風險.
每個基督徒每間教會.
都要向主耶穌基督交賬.
這個賬正正就是一個承擔風險的賬.
為主獻上最好.
最好肯定是有風險的.
那些敷衍的東西完全是沒有風險的.
你都明白我的意思.
任何敷衍的東西都是不需要風險的.
你要搞一個活動.
搞一個大的項目.
哪有說必贏的.
哪有說一定會成功的.
這個就是風險.
為主去做最好的事.
永遠都是獻上風險.
作為基督徒的我們.
對於未來應該有一種天然的興奮.
主在來這三個字.
主未來.
主將要來.
未來.
耶穌基督正正在將來裡面.
來看著我們.
你怎樣來懷著風險去走.
走到祂面前呢.
我們一起來思考一下.
怎樣來叫自己的生命.
在教會的事奉裡面.
都可以去學習一點點的風險.
信心地踏出一步.
我們祈禱.
主你幫助我們.
讓我們知道.

$^{561}$忠心不是純粹的穩陣.
忠心是我們要去承擔一些不好受的風險.
一些令我們心裡面覺得不太舒服的東西.
但這些東西可能是一些你要求我們做的事情.
無論是今天我們面對自己的生命.
無論是教會.
無論是怎樣去為我們自己定位.
我們求主你幫我們能夠學習去承擔一些風險.
這個不關我們性格的問題.
而是你對我們每一個人的命令.
求主你幫我們能夠分辨.
我們應該要做些什麼.
去承擔什麼風險.
讓祂做最好的獻上給你.
求主你幫助我們.
奉主的名求 阿們.
\newpage



\section{羅馬書 6:1-11}
\label{sec:WZaLHtVToyI}
\textbf{2023-06-25 午堂《克勝罪惡》羅馬書 6\_1-11 講員:褚永華牧師}
\newline
\newline
連結: \href{https://youtube.com/watch?v=WZaLHtVToyI}{\texttt{https://youtube.com/watch?v=WZaLHtVToyI}} ~~~~ 語音日期: 2023-06-28
\newline
\newline
\hyperref[sec:tu4SyHPIS_U]{\small{< < < PREV SERMON < < <}}
~
\hyperref[sec:index_chronic]{\small{[返順時目]}}
~
\hyperref[sec:index_scriptual]{\small{[返順卷目]}}
~
\hyperref[sec:G4jQV76in2M]{\small{> > > NEXT SERMON > > >}}
\newline
\newline
羅馬書 6:1-11
\newline
\begin{longtable}{cl}
\hline
\hline
章節 & 經文 (和合本修訂版)\\
\hline
6:1 & \begin{tabularx}{0.7\textwidth}{X} 這樣，我們要怎麼說呢？我們可以仍在罪中使恩典增多嗎？ \end{tabularx} \\ \\ \relax
6:2 & \begin{tabularx}{0.7\textwidth}{X} 絕對不可！我們向罪死了的人，豈可仍在罪中活著呢？ \end{tabularx} \\ \\ \relax
6:3 & \begin{tabularx}{0.7\textwidth}{X} 難道你們不知道，我們這受洗歸入基督耶穌的人，就是受洗歸入他的死嗎？ \end{tabularx} \\ \\ \relax
6:4 & \begin{tabularx}{0.7\textwidth}{X} 所以，我們藉著洗禮歸入死，和他一同埋葬，是要我們行事為人都有新生的樣子，像基督藉著父的榮耀從死人中復活一樣。 \end{tabularx} \\ \\ \relax
6:5 & \begin{tabularx}{0.7\textwidth}{X} 我們若與他合一，經歷與他一樣的死，也將經歷與他一樣的復活。 \end{tabularx} \\ \\ \relax
6:6 & \begin{tabularx}{0.7\textwidth}{X} 我們知道，我們的舊人和他同釘十字架，使罪身滅絕，叫我們不再作罪的奴隸， \end{tabularx} \\ \\ \relax
6:7 & \begin{tabularx}{0.7\textwidth}{X} 因為已死的人是脫離了罪。 \end{tabularx} \\ \\ \relax
6:8 & \begin{tabularx}{0.7\textwidth}{X} 我們若與基督同死，我們信也必與他同活， \end{tabularx} \\ \\ \relax
6:9 & \begin{tabularx}{0.7\textwidth}{X} 因為知道基督既從死人中復活，就不再死，死也不再作他的主了。 \end{tabularx} \\ \\ \relax
6:10 & \begin{tabularx}{0.7\textwidth}{X} 他死了，是對罪死，只這一次；他活，是對神活著。 \end{tabularx} \\ \\ \relax
6:11 & \begin{tabularx}{0.7\textwidth}{X} 這樣，你們也要看自己對罪是死的，在基督耶穌裡對神卻是活的。 \end{tabularx} \\ \\
[1ex]
\hline
\hline
\end{longtable}
$^{1}$各位聽眾平安.
見到大家能夠一齊回來敬拜上帝.
我們說敬拜上帝是我們基督徒生命中一件非常重要的事.
千百年以來我們說信了耶穌基督的人在這個世界裡面.
除了我們要敬拜我們的上帝之外.
我們還要以上帝為樂.
我們每天都能夠歡喜高興地過著一個依照信仰指導我們的生活.
但是在我們的生命中很多時候我們都發覺到力不從心.
而我們很多時候將力不從心這個感覺當成為基督徒的日常生活.
就是這樣了.
面對著罪惡我們勝過不過它.
面對著一個依靠上帝.
我們上帝生命而行的生活.
似乎我們無能為力.
我們就好像經常保羅所說的.
我真是苦呀.
我們在兩極之間到底我們要怎樣做呢.
我不想過著一個這樣被罪惡克制我的生活.
次次所思所想所做的都是被罪惡來影響的.
我不想這樣.
但我沒有能力去抗拒它.
《四經》裡面教導我們很多正面的道理.
我們要彼此相愛.
要愛我們的仇敵.
在世生活怎樣將榮耀歸給上帝.
讓我們好行為讓不信的人看到.
就能歸榮耀給上帝.
我們似乎沒有能力過這樣的生活.
但我們每個星期都回來教會聚會.
我們每個星期都聽這樣的道理.
所以在我成長的過程裡.
我有些矛盾.
我相信這個矛盾是我們很多基督徒的矛盾.
而矛盾慣了之後.
我們就當成這種矛盾的生活.
就是基督徒的正常生活.
到底我們在這個世界裡怎樣生活.
怎樣過著一個能勝過罪惡的生活.
到底勝過罪惡是不是很正常呢.
或者是一個不切實際的想法呢.

$^{41}$如果你說不切實際的話.
為什麼《四經》裡面有這麼多的篇幅.
要告訴我們不要愛世界.
我們在這個世界裡不愛世界.
怎樣過活呢.
為什麼《四經》說得這麼絕對.
不要愛世界.
人要愛世界.
愛天父的心就不再貪慮.
我們就在兩極之間.
每天的生活和世界脫離不了關係.
《四經》這裡說.
你不可以愛世界.
愛世界.
愛父的心不在裡面.
我們就將這種兩極生活.
變成慣性地.
就是這樣了.
人怎可以脫離這個世界呢.
《四經》是這麼說的.
我們很多時候就將《四經》的道理.
越說越低.
低到每個人相信就行了.
低到每個人以為信耶穌.
很容易就行了.
原來當我們讀《聖經》讀久了之後.
我們就發覺到.
原來信耶穌是很艱難的一件事.
而我們在我們日常的生活裡面.
將這種不艱難.
變成生活裡面的正常.
克服罪惡.
克服不了.
就是這樣了.
克服不了.
就是這樣了.
在成長的過程裡.
有些話.
令我.
剛才說的很矛盾.

$^{81}$在我們上一代或上兩代.
有個報道家叫葛培尼.
可能很多人都不認識他.
已經過世很多年了.
他是被世界上很多的信徒.
公認為一個很可信的牧師.
這個對我們牧師來說.
真的很忌諱.
好像只有他可信.
其他牧師都不可信.
他在世界裡面被公認.
是一個講老實話的人.
所以當美國總統.
很多時候就職的時候.
都找他坐在那裡.
不管他到底真的要找他做些什麼.
找他出來做他的招牌.
不管他那麼多.
但在他的講章裡.
他曾經說過一件事.
他說在我們每一個基督徒的生活中.
在我們的生活裡面.
我們每一個人都有棵樹.
有棵小植物.
這棵小植物.
我們每個星期天都會回來.
讓它澆水.
當我看了這段之後.
我就說這棵小植物是什麼呢.
要麼就是獅人杖.
獅人杖可以長過一個星期不澆水.
要麼就像蘭花一樣.
一個星期不澆水都可以.
原來他說的意思是說.
我們每一個基督徒.
心裡面都有一棵小植物.
我們每個星期都要回來澆水.
不回來的話又不安心.
又怕它死了.
所以回來的時候.

$^{121}$又是這樣.
不澆水.
澆少一點水.
它照樣活的.
所以每個星期都回來澆水.
回來的時候又很安心.
當我在年輕的時候看到這句話.
我說你為什麼會這樣說呢.
我們不是全心全意愛上帝.
我們全心全意跟隨他嗎.
為什麼一個星期才回來澆水呢.
但是到慢慢長大的時候.
就發覺到.
又有另一個世界了.
我用我最近看到一段話.
來跟大家說.
有一個牧師.
設計了一條問題.
什麼問題呢.
他說每個星期他講完道之後.
下個星期的時候.
他就回來問會友.
他說如果問一個問題.
你記不記得我上個星期講道.
講什麼呢.
如果他說.
我記得.
哪一章哪一字.
一個星期.
星期一到星期六裡面.
我來不及想起它.
如果一個會友這樣回答你的話.
你就給那邊講章十分.
如果他回來只是說.
我只記得你說過什麼.
上個星期沒有想到.
給自己五分.
如果你問一個人.
他說我什麼都不記得了.
只記得你說過的笑話.

$^{161}$那就倒扣十分.
他沒有說得到他得到多少分.
不過引致我問另一個問題.
我不知道在我們會友中間.
記不記得牧師上個星期講哪一道理.
如果你記得的話.
那你真的給牧師十分.
但是很多人聽完道理之後.
最多記他兩個小時.
喝茶的時候已經忘記了.
到星期一的時候.
你發覺到你是另一個人.
你想的東西.
用的詞彙.
跟人的交往.
不是跟教會裡的那些人.
所以你說的東西又不同了.
用的詞彙又不同了.
你發覺在教會裡.
都是弟兄都是姐妹.
你到外面那些人.
壞人 這麼懶.
懶蟲.
那些東西我們在教會不會講.
你怎麼會對你的組長說.
你怎麼這麼懶.
我們發覺到.
為什麼在我們的生活裡會這麼兩極化.
今天這篇的道理.
我想了很久.
到底應不應該.
跟錦奏堂的弟兄姊妹講.
但是在我們的生命裡.
我們來不及提醒自己.
到底我們屬靈的生命.
去到什麼層次.
我們不能夠.
當自己.
一種慣常的.
不是這麼追求.

$^{201}$嚮往屬靈生命成長的.
那個美麗的.
與上帝同行.
與主耶穌同行的那種生活.
我們只是埋首在我們日常的.
無力的.
無奈的生活裡.
日過一日.
我相信大家都很熟悉.
《勞家灰心》裡所講的.
那個狼子的故事.
問爸爸把自己的家業都拿走.
然後離開.
爸爸就走了.
浪費他所有的東西.
然後回家.
爸爸看到他遠遠回來的時候.
就跑上去抱著他.
然後回來開派對.
給他戴上戒指.
給他穿上悉心的衣服.
給他和朋友一起歡喜快樂.
有個牧師就這樣說.
Suppose.
這個狼子經歷過離開家裡.
再最終翻滾的生活.
然後回到家裡.
因為他醒不過來.
回到家裡.
爸爸就不計前言地.
擁抱著他.
歡迎他回家.
然後回到家裡的時候.
那個大哥仍然和他沒有兩句.
在這樣的生活裡.
當他習慣了之後.
他自己會不會再想.
我可不可以再來一局呢.
我再來一局.
問爸爸說.

$^{241}$你剩下的東西我再給我.
我就走了.
拜拜.
爸爸到底會不會.
等我用完所有的錢.
到我潦倒到無可再無以復加的時候.
他又回到家裡的時候.
會不會站在遠處.
看到我就擁抱著我.
會不會呢.
你猜他會不會蠢蠢欲試.
我看到這次的時候我就說.
我希望不會吧.
但是在我們每一個人的生命裡.
最潛意識的地方.
我們每一個人都在過這樣的生活.
就是這種以為很無力.
以為很無能.
以為生活本是如此.
離開上帝給我們的豐盛的生活.
豐盛的生命.
離開千百萬里那麼遠.
但是我們覺得我們現狀很好.
就像浪子回家之後.
再來一次家.
再離開家一次.
再回來的時候.
上帝仍然會擁抱著我們.
我們可以犯罪.
我們可以在罪中翻滾.
但是上帝始終會接納我.
上帝始終會擁抱著我.
我們就過這樣的生活.
在我們今天所讀的聖經裡.
在羅馬書第五章結尾那一句.
他說罪惡在哪裡增多.
罪惡在哪裡增多.
恩典在哪裡增多.
但是保羅就說.
在我們中國一定有些人這樣想.

$^{281}$既然罪惡增多的話.
恩典就越多.
恩典多於罪惡.
那我就可以犯多點罪.
讓恩典顯多一點.
你說我不會這樣想的.
但是實際上在我們的生活中.
我們就是這樣生活的.
我們不覺得我們自己這樣生活.
但是實際上.
這是我們生活原則的指導原則.
罪就犯.
不適應的生活照樣過.
在罪惡裡翻滾照樣過.
總之懂得認罪就行了.
在哪裡罪惡多.
哪裡恩典多.
我們在這樣的兩極之中.
扯拉扯扯拉拉扯扯.
一直這樣過下去.
但是保羅在第一章那一句說.
這樣.
這樣是什麼呢.
就是上文那一句說的.
將罪惡增多.
恩典又多.
越多罪惡.
越多恩典.
可不可以這樣呢.
保羅就說.
這樣可不可以呢.
我們可以認在罪中.
使恩典增多嗎.
他說絕對不可.
我只希望大家記得這一句話.
我們可不可以在罪惡之中.
使恩典顯多呢.
絕對不可.
如果今天講完之後去到12點.
你還記得這一句話就好了.

$^{321}$當然你記得其他就更好.
在整個星期裡你能夠記得這一句話.
我可不可以活在罪惡之中.
使恩典顯多呢.
No 絕對不可.
因為這樣的緣故.
我們就在生活裡追求上帝給我們的聖潔.
難道你不知道我們受洗歸入基督耶穌的人.
就受洗歸入他的死嗎.
一路打後就在說.
他用洗禮這個道理.
洗禮這個道理和洗禮的形式.
那個形象.
來講明生命裡的改變.
這裡我要解釋一個問題.
這個問題或者這樣的想法.
是令我們基督徒很迷糊的.
我們很著重一個人的洗禮.
是的我們應該信耶穌之後.
我們應該洗禮.
但是洗禮有另外一種意義.
不是洗禮那樣東西.
好像我們進入神馬故事那樣.
我們進入河裡.
浸一浸起來.
整個人變得年輕了.
洗禮不是這樣的作用.
洗禮是說.
在我內心裡有一個極大的轉變.
我要用外面見得到的動作.
見得到的形式.
表明在內心裡發生.
你見不到的東西.
內心裡見不到的是什麼.
原來當我信耶穌的時候.
生命有一個很大的轉變.
我的身份改變了.
從前我是屬於這個世界.
從前我是屬於這個魔鬼.
但是現在這一刻開始.

$^{361}$直接洗禮公開表明.
我和他一刀兩斷了.
我現在站在上帝這邊.
我和你沒有關係了.
魔鬼不能再控制我了.
我不屬於這個世界.
我又不屬於魔鬼.
現在我全部屬於上帝.
屬於主耶穌.
洗禮是用外表的形式.
證明我內心所發生的改變.
不是那個形式讓你有改變.
是那個形式告訴所有人.
我內心有改變.
但是很多時候.
我們只著重外面的那個形式.
沒有留心到.
原來當我接受洗禮的時候.
我的生命 我的身份.
隨時有個改變.
我們仍然活在那個洗禮的形式裡.
沒有真正拿捏到那個洗禮.
代表了我心裡所發生的強烈的改變.
當我信耶穌的時候.
表明我撇開這個世界.
表明我不再受他的限制.
我願意一生跟隨主耶穌基督.
我願意我的生命.
彰顯上帝的榮耀.
我願意跟隨主耶穌基督.
讓世人看到我的行為.
歸榮耀給我的上帝.
我和那個舊的生活完全沒有關係.
我願意過著一個新的生活.
我願意我們的弟兄姊妹.
在我們的生活裡.
能夠彰顯這種因洗禮.
被人看到外面的那個形式.
彰顯裡面的實在.
保羅在羅馬書第六章第三節.

$^{401}$打後再解釋.
他說好像主耶穌基督那樣.
他是用耶穌基督的那個洗禮的形式.
耶穌下到水裡.
然後水裡代表了他死了.
埋葬了.
從水裡上來.
他又活著了.
表明我們又活著.
我們以為耶穌基督的洗禮.
只是一個外表的形式.
讓我們看到他是這樣的.
所以我們都會這樣.
但我們只是著重他的形.
而沒有著重他的實.
那個形是這樣.
但是那個最重要的就是說.
你看看耶穌下到水裡.
我們就浸在水裡.
耶穌死了.
我們就在水裡死了.
但當耶穌從水裡上來的時候.
我們又從.
耶穌從死裡復活的時候.
我們從水裡上來.
他一舉一動有新生的樣式.
我們一舉一動也有新生的樣式.
不是那個洗禮.
讓我們一舉一動有新生的樣式.
那個新生的樣式.
在我們相信主耶穌基督的時候.
我們的生命裡面有很大的轉變.
而那個洗禮告訴所有的人.
有些東西在我們裡面發生了.
那個就是我的生命.
從水裡上來的那一刻.
一舉一動有新生的樣式.
各位弟兄姊妹.
願意我們能夠在生命之中追求.
追求和上帝自己的話語更加貼近.

$^{441}$剛才我們所讀的啟應經文.
說明我們要按照上帝自己的話語來生活.
我們怎樣在生活之中熟讀上帝自己的話語.
讓他的話語成為我們生命裡面的指引呢.
我們很多時候覺得勝經的話語.
和我們生命裡面的那個現實.
似乎相差很遠.
聖經裡面說.
上帝的話語是我們腳前的燈.
是我們路上的光.
我們很多時候不覺的.
我有我行 我有我做.
我想這樣就這樣.
我們剛才不是有讀到你的話語.
在我口裡好像風放下的蜜 這麼甘甜.
在我上堂有回甘.
我們喝咖啡 喝功夫茶 有回甘.
上帝的恩典可以在我們口裡.
在我們心裡有回甘的作用.
你問一下 你什麼時候領受了上帝的恩典呢.
而什麼時候領受了上帝的恩典之後.
有回甘的感覺呢.
怎樣才算回甘呢.
我們很多時候捕捉那個感覺.
我用一個現實的一句話.
我想我們成長的過程中間.
我們很多時候都看日本劇.
日本劇裡面有句很有名的東西.
有一種幸福感.
我記得很多年前 三十年前.
我那年開刀做手術.
在家裡看了很多日本劇.
其中一套就是那個叫做長態的壽司.
幾十年前.
那個人說我一定要決定.
要找全日本最有名的壽司專家.
誰知道那個人就說.
我想吃壽司 你做不到.
你一定做不到.
他沒有說到 他自己誤會到.

$^{481}$原來那就是媽媽的味道.
原來他們生活很窮困.
住在鄉下地方.
在那個壽司裡面沒有什麼特別好的東西.
就是用野菜.
所以他誤會到這個道理.
找人去那裡割野菜.
找人回來店裡.
把野菜變成壽司的餡.
把他給了.
那個小孩說.
這就是媽媽的味道.
吃完之後.
全身充滿幸福感.
我們對上帝的恩典.
有沒有這種幸福感.
在上帝的恩典裡面.
我被他照顧得無微不至.
這樣他都知道.
這個真的需要在我們的生命裡面.
慢慢去體會 慢慢去捕捉.
你想一下.
最近看到一段這樣的說話.
上一代的爸爸和這一代的爸爸.
有很多不同的地方.
這一代的爸爸會洗碗.
會幫忙看十元.
上一代的爸爸不幫忙洗碗.
吃完飯之後 丟下筷子丟下碗.
就去那裡看電視.
今宵一邊看電視.
等太太幫忙洗碗.
如果有一天.
先生突然間吃完飯之後.
去洗碗.
本來是太太洗的.
可能是姐姐幫你洗.
洗完碗之後弄得妥妥貼貼.
太太打算進來洗碗一看.
不洗我洗.

$^{521}$今天又是這樣.
希望明天都是這樣.
過了一個星期之後都是這樣的話.
幸福感.
真的很幸福.
少了一樣東西可以做.
上帝的恩典在你的生命之中.
聖經的語句就說.
好像風防下敵的物禁金團.
拿我們現實的語言來說.
上帝的恩典在我們的生命之中.
讓我們充滿了被上帝愛護.
被上帝看顧.
被上帝護佑保護的那種幸福感.
在她眼中.
在她的保護之中.
我絕對安全.
她知道我所需要的一切.
而我華視我的牧者.
我必不自缺乏.
那種幸福感.
你要追求一下.
你要體會一下.
如果在我們的生命之中.
我們可以追求到.
原來一個被上帝所照顧.
被上帝所愛的生活.
是這樣的幸福.
你有沒有藉著這樣的幸福.
可以跟最後說.
我不想要這樣的生活.
我要這樣的生活.
你是不是追求這樣的生活.
追求這樣的生活.
保羅說的.
我們可以靠上帝給我們的力量.
聖靈給我們的力量.
跟最後說No.
然後靠聖靈給我們的力量.
跟上帝說Yes.

$^{561}$是在乎你.
在乎你一念之間.
很多年前我看一本書.
魏諾德寫的.
Dallas Wheeler寫的.
《屬靈操練的禮讚》.
《屬靈操練》裡面說了一件事.
他就說.
我首先介紹一下.
Dallas Wheeler是一個怎樣的人物.
他是一個浸信會的評述道牧師.
他自己的正職.
在南加州大學裡面做哲學教授.
副教授.
在南加州大學裡面做哲學教授.
但他是一個浸信會的評述道牧師.
他做牧師.
他寫了很多書.
幫助了很多人.
認識屬靈生命裡面.
很多很多很基本的東西.
有一次在書裡面他這樣說.
他說當我們自己的生命被上帝自己所使用.
我們的生命可以怎樣榮耀上帝呢?.
我們的生命可不可以去到一個層次裡面.
被上帝自己所使用.
讓我們緊緊跟隨主耶穌.
Dallas Wheeler說可以的.
可不可以在我們的生命之中.
繼續不斷地追求.
讓上帝被彰顯.
他說可以的.
以一個哲學教授來說.
令我們很奇特的.
很多人覺得.
哲學教授怎會這樣說.
講哲學就講理性的東西.
為何講這麼神秘的東西.
要我們追求屬靈生命裡面的成長呢?.
他是少有的.

$^{601}$其中一個這樣的人.
會帶領在洛杉磯裡面.
有很多教會的牧師.
一起過著一個很簡樸的屬靈生活.
他住的地方很偏僻.
但有一個地方.
除了他日常教書之外.
他又潛修聖經.
讀聖經.
默想.
很多牧師和他一起過著一個簡樸的生活.
我們可不可以靠賴.
聖靈的能力.
讓我們可以撇離罪惡.
跟隨主耶穌呢?.
他說可以的.
他說了一段這樣的說話.
絕大部分的罪惡.
我們犯罪.
是一個理性思考的結果.
當我看到這句說話的時候.
我不是很明白的.
我們犯罪的時候.
不是哎呀.
我又犯了.
這樣才叫犯罪.
這些是我們不知不覺犯了的罪.
但這位哲學教授就說.
不是的.
你想想.
我們每一個人犯罪的時候.
那個罪惡的意念.
在我們的腦海之中生出來的.
或者從心裡面的情慾而生出來的.
進到哪裡?.
進到這裡.
然後你自己的理性跟你說.
我要這樣做.
絕大部分的罪惡.
是人理性思考的行動.

$^{641}$你回去想一想是不是.
絕對絕對少有.
哎呀.
我犯了.
還不知道自己犯了什麼罪.
絕對絕對少有.
因此他就說.
我們可不可以在生活之中.
在我們生命成長的過程裡面.
求聖靈幫助我們.
讓我們的思維可以指導我們的生活.
當罪惡的意念.
在我們心裡面生出來的時候.
那個情慾在我們心裡面生出來的時候.
我們能夠靠上帝的恩典.
理性的他就說.
No.
No呀.
我不要.
你可以的.
不過在我們慣常的生活之中.
我們慢慢慢慢.
沒有這樣操練.
沒有這樣來生活.
以為當罪惡一來的時候.
就迎刃而上.
是不是.
聖經裡面有沒有這樣的例子.
有.
藥室.
那些藥室哥哥.
年輕貌美的藥室哥哥.
被賣到波提佛的家.
波太.
波提佛的波太.
每天在他旁邊見到他.
喂.
上床了.
年輕人.
整間房子裡沒有人.

$^{681}$只有你和我.
你不出聲 我不出聲.
就瞞天過海.
怎不知那個藥室哥哥.
怎麼說呢.
我豈能作者大惡.
得罪誰呢.
得罪上帝.
No.
我們經常不是這樣想的.
嘩.
這麼著數的東西.
送到上門.
邪惡了都沒有人.
做完都沒有人知道.
藥室哥哥就說.
No.
理性的決定.
誰知道.
上帝知道.
放了一些東西進去.
你以為只有你知道.
沒有人看到.
上帝看到.
當你過著這樣的生活的時候.
你會發覺到.
原來過著一個基督徒的生活.
按照上帝的話語而生活.
他像一個腳尖的燈.
路上的光.
他的話語好像.
風浪下的一樣.
金電.
原來需要你付出一些代價.
在你的生活之中.
經常不斷去想他.
第二件事.
我要和大家思想的.
就是這裡有句話.
在這裡第五節.

$^{721}$我們若與他合一.
經歷與他一樣的死.
我們也將經歷與他一樣的復活.
我們知道我們的救人和他.
銅釘十字架.
使罪身滅絕.
叫我們不再作罪奴隸.
這句滅絕的兩個字.
讓我們產生很大的矛盾.
不是說他翻譯錯誤.
是當我們中國人一看到滅絕的時候.
以為我們相信耶穌.
我們就可以和罪惡絕緣.
他現在告訴我們一個道理.
在加拿大書裡面也有這樣的道理.
世界對我而言.
世界釘在十字架上.
對我毫無作用.
他已經死了.
對世界而言.
我已經掛在十字架上.
我又死了.
所以世界和我沒有什麼關係.
他看我死了.
我看他死了.
當我們讀到這句.
罪身滅絕的時候.
我們產生一個概念.
我們相信耶穌之後.
我們理應和罪惡絕緣.
滅絕的兩個字.
不是滅絕師太.
但是在這句字.
在新約裡面出現了27次.
有24次在保羅書信出現.
而這個字的翻譯.
我鼓勵大家回去早上查經.
有時候來不及查.
看看不同的譯本.
你會得到不同的概念.

$^{761}$不是一些經文翻譯得不對.
不是.
原來這個字眼裡面有些意思.
NIV的翻譯.
有些聖經的翻譯.
是使罪身paralyzed.
癱瘓.
使罪身癱瘓.
不是死了.
他什麼時候會活動.
引致很多人說.
我們這個人.
相信了耶穌之後.
仍然有這個身體.
只要他存在一天.
我們就和罪惡脫離不了.
因為罪惡經常用這個身體.
來引誘我們.
讓我們落在罪惡之中.
這是在屬靈生命裡.
掙扎的兩極.
我們都經歷到的.
原來發覺到.
我們不是完全死了.
我們又不是完全癱瘓.
我們有時候會活動.
罪惡會在我們的生命最脆弱的地方.
讓我們可以重新攪動起來.
讓我們犯罪.
我們又不是像保羅所說的.
對世界而言.
我已經死在十字架上.
滅絕.
對我而言.
世界被掛在十字架上.
又死掉又滅絕.
在我們每一個人的生命之中.
在兩極之中掙扎.
但Dallas Wheeler就說.
是的,你在兩極之中掙扎.

$^{801}$但只要你靠賴上帝給你的恩典.
當這件事蠢蠢欲動的時候.
你會對他說.
NO.
當這件事讓你對罪惡產生不了作用的時候.
你說上帝啊,我要,I want it, YES.
當你越來越這樣做的時候.
你就慢慢慢慢發覺到.
你做人的生命不斷不斷地成長.
然後對罪惡你可以輕易地說.
NO.
對上帝你可以很輕易地說.
YES.
生命在上帝手中.
我們的成長同樣都在上帝自己手中.
上帝會用他自己的恩典幫助我們.
讓我們可以過著一個上帝喜悅的生活.
在這裡他就說.
因為知道第九節.
因為知道基督既從死人中復活.
就不再死.
死也不再作他的主了.
他死了是對罪死.
只這一次.
他活著是對上帝活著這樣.
你們也要看自己對罪是死的.
在基督耶穌裡對上帝卻是活的.
靠賴著恩典.
我們可以對罪說.
NO.
因為我已經死了.
靠賴恩典.
我可以對上帝說.
YES.
因為在你裡面我是活的.
這樣的生活是需要我們慢慢去操練的.
這樣的生活是需要我們記得上帝自己的話.
從星期一到六.
到下個星期一.
你都要記得的.

$^{841}$願主幫助我們.
讓我們在生命之中繼續的成長.
我們會遇到一些的挑戰.
我們會遇到很多的挑戰.
不同的挑戰.
但是生命裡面的挑戰.
會讓我們無所適從.
不知如何是好的.
有時候我們發覺到我們的難處是大的.
大到我們自己是克服不了的.
但是不要讓難處擾亂了你的視線.
靠賴恩典.
聖靈的幫助.
對罪惡說NO.
對上帝說YES.
永遠都是可行的.
願主賜福給大家.
\newpage



\section{}
\label{sec:G4jQV76in2M}
\textbf{2023-06-29 周四晚查經班《被擄歸回歷史研讀(18)》 講員:曾裔貴牧師}
\newline
\newline
連結: \href{https://youtube.com/watch?v=G4jQV76in2M}{\texttt{https://youtube.com/watch?v=G4jQV76in2M}} ~~~~ 語音日期: 2023-06-30
\newline
\newline
\hyperref[sec:WZaLHtVToyI]{\small{< < < PREV SERMON < < <}}
~
\hyperref[sec:index_chronic]{\small{[返順時目]}}
~
\hyperref[sec:index_scriptual]{\small{[返順卷目]}}
~
\hyperref[sec:W6aFxBk_o0s]{\small{> > > NEXT SERMON > > >}}
\newline
\newline
$^{1}$好,我們一起祈禱先.
感恩天父給我們.
雖然停了兩個星期.
我們可以又恢復.
讀聖經的時間.
盼望我們今晚能夠.
去到尼希米記.
最後兩章.
聖經的時候.
都能夠掌握裡面的經文.
特別是讓我們能夠看到.
尼希米.
他處在自己.
工作的崗位.
還有在這個.
這麼大的事公裡面.
他能夠.
有一個智慧.
還有能夠.
去動員不同的.
猶太的同胞.
特別是在耶路撒冷裡面的.
祭司長.
一些當地的領袖.
一同的同心協力.
更加引發了一個.
很大的規模.
的培靈的.
民俗會.
使很多的猶太人.
心靈都被甦醒.
來投入事奉.
我們都盼望神.
在我們的教會.
在香港的教會.
都來做感動.
做引導的工作.
使我們更加在一個新的時代.
無論我們怎樣看.
那個環境.

$^{41}$我們都能夠.
很明白我們.
在一個需要復興的.
一個社會.
我們能夠將神的福音.
是很有效.
很有力的去傳揚.
使更多人得到.
神的福音的好處.
今天晚上我們都等候在你的面前.
有些舊約的背景.
讓我們能夠容易的明白.
和來到去.
實際在我們的生活.
新約的時代.
我們去學習應用.
上望禱告.
奉耶穌基督的名.
阿們.
又開始.
簡單的一個半小時的讀經的時間.
要和大家.
重溫一些的背景.
第八章.
就是開始.
完成了那個.
城牆之後.
那個嚴肅會的開始.
延展到.
第十章.
經過了鑄棚節.
七月有好幾個大的.
節期,有吹角節.
有七月初十的贖罪日.
有七月十五.
連續有七天的.
鑄棚節.
都是一個.
重要的七月份.
城牆做好.

$^{81}$而進入一個.
重要的節期的開始.
然後.
有以色拉和.
尼欽米一同的同工.
去召集.
在這個.
又是那個水門.
來到去.
有一個讀經.
有很多的利未人做一個.
研習.
真的引發了一個嚴肅會.
引發了他們很多人.
一齊的來到去.
禱告.
還有禱告完之後他們有行動.
那個行動.
就是他們來到去.
要簽約.
要答應.
遵守摩西的律法.
這個是.
我們在這幾章裡面看到的.
今天我們看第十一章.
我們先讀了經文.
第十一章開始.
百姓的首領.
住在耶路撒冷.
其餘的百姓.
制籤每十個人.
中使到一個人.
住在聖城耶路撒冷.
這個抽籤.
不是生死籤.
是大家都願意.
讓神來到去.
感動引導.
好像十一奉獻.
十個.

$^{121}$可能是.
以一個的平均數.
因為你不是那麼齊整.
每一個支派.
每一個家族都不是那麼齊整.
所以他們就可能.
都不計那個家族.
大家都願意.
派人出來.
湊夠十個.
十個這樣.
就抽一.
很特別就是.
在這章聖經提到很多住.
住在哪裡.
為什麼那麼奇怪.
住在哪裡都要.
關心的嗎.
快九月.
不知道那些人會不會.
在名校區.
九龍塘九龍城.
北角.
這些傳統名校區.
或者赤柱.
就植入.
當時是怎樣呢.
第二節.
凡是甘心樂意住在耶路撒冷的.
百姓都為他們祝福.
終於我們發現了.
那個問題所在.
原來是傳聖經.
在這裡開始.
就要安排.
人們住在耶路撒冷.
第三節.
以色列人濟斯利未人.
尼提靈和所羅門.
卜人的後裔都住在.

$^{161}$猶大的城邑.
即不是耶路撒冷.
這裡有五個類別.
不同的階層的.
濟斯利未人.
尼提靈是那些電役.
在聖殿的地方.
做一些孤雷.
擔台的工作.
有所羅門的卜人的後裔.
他們就住在猶大的城邑.
各在自己的地業中.
本省的首領.
住在耶路撒冷的.
記載下面.
就被記錄.
住在耶路撒冷.
其中有些猶大人.
和汴雅敏人.
猶大人中有法勒斯的子孫.
烏西亞的兒子.
阿他亞.
烏西亞是撒迦利亞的兒子.
撒迦利亞是.
阿瑪利亞的兒子.
阿瑪利亞是士法提亞的兒子.
士法提亞是.
馬勒列的兒子.
又有巴祿的兒子.
馬西亞.
巴祿是谷河西的兒子.
谷河西是哈塞亞的兒子.
哈塞亞是.
阿大亞的兒子.
阿大亞是約瓦納的兒子.
約瓦納是撒迦利亞的兒子.
撒迦利亞是.
士羅尼的兒子.
住在耶路撒冷.
法勒斯的子孫.

$^{201}$共有468名.
都是勇士.
特別要強調住在耶路撒冷.
在這張聖經.
聖城.
做好了.
整場的聖城.
聖殿其實一直都有.
120多年前.
已經整好了.
但一直沒有成長.
有的都是破爛.
那些門山.
都是沒有染雞籠.
所以尼欽米回到上去.
就發現了.
外邦人隨便.
就可以來聖城.
來到去挑撥離間.
在安息.
就特意來買賣.
令到猶太人的信仰.
都是.
途具形式.
可有可無.
變了大家都.
不是那麼關心.
安息日要敬錢的.
敬拜神.
獻祭諸如此類.
亦都是太過.
頻繁的交往.
令到外邦人看不起你的.
聖殿敬拜.
自從尼欽米和那些領袖.
做好了整場之後.
第四到第七章.
就是說那個.
批評論斷.
那個難阻.

$^{241}$都做好了.
第七章就是.
數點文數 家譜.
哪些人是.
被揀選出來.
可以事奉.
其實是分別為勝.
第八到第十章.
就是在一個有圍牆的.
新的形勢.
底下的.
節期的.
可以說是第一個駐盤節.
是有圍牆.
耶路撒冷回歸.
回來亡國之後.
回歸回來.
第一次有自己的城牆.
第一次趕走那些外邦人.
第一次沒有分別為勝.
就不可以來到首節.
第一次.
在整個耶路撒冷.
是自己的猶太同胞.
來到歡樂 禁食.
來到去.
聽摩西的律法書.
來到去友賢習班.
來到去被神的靈感動.
大家都願意手鹽縮會.
然後第十章的結論就是.
大家願意簽約.
大家願意.
謹守安息日.
大家去到安息日就不做買賣.
以前就忍不住.
手 現在就.
都因為靈性感動.
就真的看安息日為一個很重要的日子.
去敬拜神.

$^{281}$然後他們又納祭.
不讓他們的兒女.
娶外邦的當地迦南女子.
或者嫁給迦南的外邦人.
就做了這個簽約.
接著就是第十一章.
李希米就去數點.
現成的.
有哪些敬虔的猶太人.
是選擇住在耶路撒冷.
我們一會兒.
就講解一下.
這個結論.
為什麼住在耶路撒冷.
我們一會兒讀完之後.
我們一路去看一下.
住在耶路撒冷有什麼好處.
有什麼重要.
為什麼這張聖經會講.
這麼多關於.
這麼實無的住在哪裡都關你事.
你就可以想像得到.
他們徹底的.
去重建聖文.
重建城牆.
是為著.
重建神的選民.
那個靈性.
與神的立約.
回到最.
出埃及的時候.
那個神和.
以色列這個民族立約的時候.
那個的光景.
那個水平.
現在李希米.
不是甘於早了.
城牆就收工了.
我攻城新退.
我收了那些尾款就走了.

$^{321}$不是的.
他講的是反而是神的選民.
要怎樣能夠恢復.
那個敬虔的.
分別為聖的那個水平.
那就.
和居住在耶路撒冷有密切的關係.
我們記住.
剛才他特別提到.
有五個類別的人.
這些全部是在聖殿事奉的人.
現在就記載了.
有其中一個家族的支派.
就是.
這個法拉斯.
大支派的後裔.
就居住在耶路撒冷的.
都是勇士.
就468名.
就是說.
不只是打得這麼簡單.
在民間裡面有一定的影響力.
都是國家裡面.
一些重要的.
重要的家族.
去到第七節.
辯雅敏人.
中有米舒蘭的兒子.
薩路.
約葉的兒子.
對不起是約葉.
約葉是彼大雅的兒子.
彼大雅是哥拉雅的兒子.
哥拉雅是.
馬西雅的兒子.
馬西雅是.
以鐵的兒子.
以鐵是耶西雅的兒子.
其次就是有嘉拜.
撒萊的子孫.

$^{361}$共有928名.
世基里的兒子.
就是約爾.
的長官.
哈西路雅的兒子.
猶大是耶路撒冷的.
副官這裡說的就是.
有官階有官職的.
是辯雅敏支派的人.
現在有.
928名.
就是說又是.
另外一個支派.
住在耶路撒冷的.
而且在耶路撒冷裡面.
有官階的就有這些人.
只是舉例.
比較大的家族.
那些男丁的名字.
就不是928個.
名字都落齊了.
就是說.
這班的領袖.
譬如我們約書亞團.
團契的裡面.
肥沙就是其中的.
領袖.
肥沙之下有十幾個弟兄姊妹.
類似這樣.
第十節.
祭司中有雅根.
雅納的兒子.
講祭司.
還有管理神殿的西萊雅.
西萊雅是希臘家的兒子.
希臘家是米蘇蘭的兒子.
米蘇蘭就是薩都的兒子.
薩都就是米拉約的兒子.
米拉約就是阿希特的兒子.
這些大衛時代的.

$^{401}$祭司來的.
還有他們的弟兄.
在殿裡面公職的.
共有822名.
由今次就記載.
祭司的家族.
有822名.
就在耶路撒冷居住.
又有耶羅罕的兒子.
亞大雅.
耶羅罕是比拉利的兒子.
比拉利是暗洗的兒子.
暗洗是撒迦利亞的兒子.
撒迦利亞是巴斯胡爾的兒子.
巴斯胡爾是瑪基雅的兒子.
還有他的弟兄.
作族長的.
就有耶路撒冷的兒子.
耶路撒冷的兒子.
弟兄作族長的.
就有242名.
又有阿撒列的兒子.
阿瑪瑞.
阿撒列是阿哈塞的兒子.
阿哈塞是米實利末的兒子.
米實利末是.
陰默的兒子.
還有他們的弟兄.
大能的勇士.
共有128名.
哈基多林的兒子.
薩巴第葉.
是他們的長官.
這裡是說.
祭司的行列.
有哪些比較大的家族.
居住在耶路撒冷的.
接著到利未人.
利未人中有哈述的兒子.
是瑪雅.

$^{441}$哈述是厄利金的兒子.
厄利金是哈沙比亞的兒子.
哈沙比亞是布尼的兒子.
又有利未人的族長.
沙比泰和約薩特.
管理神殿的外事.
外面的神殿.
分開內的就是聖所.
外面的就是會幕.
聖所裡面.
聖所還有一層的就是致聖所.
致聖所只是擺了一個藥櫃.
有一個簾.
是記住聖所和致聖所的.
聖所外面和會幕外面.
又有一個門簾.
祭司才可以進入會幕裡面的聖所.
致聖所是那年的大祭司.
一年一次的.
七月初十才可以進入致聖所.
分得很清楚.
現在原來有些利未人.
是做外面的工作.
外面也有很多工作.
例如那些煩祭壇.
每天都要燒柴.
每天要更換柴的時候.
添柴的時候.
會有很多灰.
大家有沒有燒過柴?.
應該有.
你要拿走那些灰.
不然積著.
那些木蓋住.
燒得不夠旺盛.
所以要拿走那些灰.
拿走那些灰之後.
其實不單止.
不要忘記.
殺身獻祭.

$^{481}$是很血淋淋的.
屠宰完之後.
很多那些血.
可能變成豬紅.
不是豬.
是那些牛的.
或者是羊的血.
要清理的.
要拿去一個地方.
固定一個地方.
來倒垃圾.
變成很多的東西.
這些下欄.
是由利未人做的.
管理外電.
神的電的外務工作.
在戶外的.
就算醫院也是.
管事部.
因為那時候我在醫院管理局.
管事部.
他手下有很多Armor那些.
做不同工作的.
有些是殯儀館.
對不起,是殮房.
有些是在.
不同的辦事處.
做的,不同的病房.
做的.
有些可能是.
大的管理處.
隨時派人去的.
我又認識有些.
被人燉冬菇.
他的長期工作.
就是戶外洗那些外牆.
做到傻.
日曬雨淋.
是沒有人做的.
但是.

$^{521}$哪些人不喜歡他.
就跟那個管事.
講通了.
就派他去.
原來很黑暗的.
內幕.
就是做那些最辛苦的東西.
有些輕鬆的.
我那時候部門就輕鬆了.
我那個Armor.
說每天回來第一時間.
就幫老闆洗杯.
沖咖啡.
沖完之後就沒事了.
搞定了.
九點半就拿文件.
去不同的部門.
十一點就回來了.
拿文件去.
十一點就回來了.
那你知道他到處去哪了.
打牙膠.
原來不同的部門.
都有不同的崗位.
有些很油的油位.
看你懂不懂得做.
但這些告訴我們.
你每人都分了不同的崗位.
可能負責哪些哪些.
特別是去到大衛的時候.
因為其實已經有聖殿.
不需要朝行晚策.
其實就不需要那麼多人力.
不需要那麼多利未人幫忙.
大衛很有智慧.
他提升利未人的級數.
記不記得是什麼.
利未人變了24個班的師班.
當然有其他的樹木.
提升利未人.

$^{561}$誰可以唱詩.
就有其他的職位.
利未人誰可以唱詩.
就看你是什麼聲部.
就有師班長.
例如耶道頓.
有.
一時想不起.
有三個.
師班長.
訓練利未人.
幫他們安排不同的班次.
來去頌唱.
大衛臨終之前.
來去設立的師班.
就是利未人做的.
這班都是要蒙召.
都不是隨隨便便.
口才歌聲漂亮就行.
原來是利未人做的.
所以你可以想像到.
其實是的.
因為你每個星期去領唱.
如果你是誇耀你自己喉嚨的口才.
那個聲線漂亮.
就是漂亮歌王.
那就要小心了.
因為你唱給人聽.
不是唱給神聽.
那就長遠不健康.
慢慢會飄飄然.
會有一種不知名的驕傲.
不是說錯的.
但聖經告訴我們.
不是這樣.
這裡是說.
在耶路撒冷裡面.
都有不同的階層.
全部.
這五個階層.

$^{601}$都是在聖殿裡面來去公職.
即是說讓聖殿的事奉.
能夠成為.
常規的.
運作.
能夠.
不再是一步十行.
突然間有感動.
我們就去獻祭.
沒有感動就一年都不回教會.
一年都不去禮拜堂.
一年都不做奉獻給神.
你可以想像到.
過去.
黎希敏還沒有回來之前.
還沒有成長之前.
一直在說.
他們的敬拜.
是可有可無.
他們的奉獻.
更加是一個這樣的例行公事.
甚至去到一個地步.
都沒有興趣去做.
因為.
越是沒有規模.
越是沒有一個準備.
充足的準備.
來去安排.
自然就是大家都沒有心.
大家都沒有.
mood來去提升.
對神的敬拜.
所以你看黎希敏.
她真是很.
偉大的一個工作.
就是她想到.
一件很重要的事.
就是你住在哪裡.
我們先讀讀.
我們說的是那個外事.

$^{641}$第十六節.
利未人族長沙比泰的約撒拔.
管理神殿的外事.
第十七節.
祈禱的時候.
為清者領首的是米迦的兒子.
瑪他尼.
米迦是撒底的兒子.
撒底是亞撒的兒子.
又有瑪他尼.
弟兄中的伯布加為父.
還有沙姆亞的兒子.
亞大.
沙姆亞是加拉的兒子.
加拉是耶道頓.
剛才我漏了.
耶道頓就是其中一個在大衛的時代.
的兒子.
在聖城的利未人.
共有284名.
利未人都分了幾個崗位.
原來.
連帶領祈禱.
都有利未人的領袖.
現在這裡又.
數了284個.
記住.
這次再一次.
第一次提這個名.
就是住在聖城.
其實一開始第一節已經說了.
住在耶路撒冷的人.
現在來到這裡.
就來記載.
在聖城居住的.
這些事奉者.
這些領袖.
真的很重要.
整張聖經.
是想交代.

$^{681}$五個工作的階級.
來住在聖城的人.
他們是甚麼支派.
他們有多少人.
他們是分開在不同的崗位.
是甚麼崗位.
第十九節.
果然去到守門.
守門阿谷和達門.
並守門的弟兄.
有172名.
其餘的以色列人.
祭司利未人.
都住在猶大的一切城邑.
各在自己的地業.
即是.
分地的時候.
分了給不同支派.
他們就可以.
來到去.
那個支派住在哪一個部份.
他們就去.
拿回自己的原居地.
去居住.
但是在這一章.
比較重要的部份就是.
尼希米.
要求.
人.
住在耶路撒冷.
首先他就是用了18節.
來記載.
這些階層的人.
住在耶路撒冷的.
有多少人.
這樣來編配.
第22節.
在耶路撒冷.
利未人的長官.
管理神殿事務的.

$^{721}$是歌唱者.
亞撒的子孫.
巴尼的兒子.
烏西.
巴尼就是哈沙比亞的兒子.
哈沙比亞.
是瑪他尼的兒子.
瑪他尼是米加的兒子.
王.
為歌唱的出命令.
每日供給他們.
必有一定之糧.
就是說他們要.
專職去專心.
去well prepare.
去做他自己應該做的工作.
例如他要做師班.
他要專職.
專心的去準備.
在重要的節期.
特別是.
每個安息日.
猶太人.
安息日應該要上去.
敬拜神.
他們就去.
利未人就安排了.
你們不是做擔台的.
你們也不是做整理柴火的.
你們是要做.
也不是做殺生的.
你們是要做唱頌的.
他們就要準備了.
所以又分了不同的.
牧區.
不同的部來去做事.
所以.
這都是黎希米.
去配合當時的領袖.
去使.

$^{761}$不同的領袖.
不同的恩賜.
不同的輩分.
例如他們出生.
就是利未人.
你就要長大之後.
去到幾多歲.
利未人.
要離開家裡.
去聖殿學習.
30歲.
就去公職.
25歲就開始.
去神學院.
5年學徒制.
在聖殿.
30歲就開始公職.
去到50歲就要離開.
可能你可以做顧問.
但是正常的工作.
不關你的事.
這是利未記的記載.
原來.
是有那個制度的.
神在未入迦南之前.
已經安排好.
祭司要做什麼.
利未人要做什麼.
所以其實大家要讀舊約.
就知道原來.
神的回望.
神的聖殿.
都是有規有矩.
有制度的.
不是一時感動的.
感動是應該.
但不是一時的.
是要有一個規模的.
神是喜歡秩序的.
神不是叫人混亂的.

$^{801}$是叫人安靜.
和有次序的.
這是哥林多前書.
十四章.
最後那節經文.
神原不是叫人混亂.
乃是叫人.
突然間忘記了那個意思.
我很快就讀給大家聽.
是.
神是喜歡秩序的.
神不是叫人混亂的.
是.
沒錯.
講方言的問題.
是.
沒錯.
三十三節十四章.
因為神不是叫人混亂.
乃是叫人安靜.
上文說的是先知的靈.
是馴服先知的.
就是說.
如果你是先知.
你裡面有感動.
是你的性靈感動你.
你去講.
但保羅說不好意思.
先知的靈.
都是要馴服先知.
即是說先知.
是有規有矩.
是有秩序.
是有sense.
他講完之後.
就不會霸住咪不下來.
不繼續講.
靈因的致命問題就在這裡.
感動不是大哂.
誰都要聽我講.

$^{841}$我兩個小時不停我繼續講.
保羅說不是.
感動不是大哂.
感動都是好的.
但是感動都需要有sense.
都需要有神的.
一路的帶領.
更好的就是.
性靈會在我們的心靈裡面.
感動.
我們在安靜底下被感動.
感動完之後.
我們仍然不是叫人混亂.
是叫人安靜.
仍然是很humble的.
講完神要我講的寓言.
我坐下讓其他.
可能在在有另外的人.
有另外的感動.
他就可以來到去起身講.
所以是一個很和諧的聚會.
不是一個混亂的.
一個好像完全失控的狀態的一個聚會.
所以是一個盲點.
如果一個聚會失控是盲點.
當然每個人都在睡覺.
也是一個盲點.
這個安靜就不是這個意思.
那個是.
大家是很認真地.
心靈裡面是敬拜神.
聆聽神.
有什麼教誨.
到我們要去到唱詩.
頌揚神的時候.
大家是從心裡面.
去頌揚神.
盡力去頌揚神.
每一個環節.
我們都是全情全心.

$^{881}$去投入參與.
去敬拜.
這樣叫聖靈同在的.
敬拜的聚會.
我講完了.
回來這個時候.
我們要看回.
尼希米基給我們很重要的啟示.
沒錯.
已經不是我們的時代.
我們似乎都好像.
很難去.
干預.
大家的生活.
作息.
工作或者是居住.
但是.
要我們從其中找到一些.
很重要的大方向原則.
我們就不只是.
咬文嚼字.
而是那個大方向.
那個時代.
為什麼尼希米這麼緊張.
他們的居住地方.
他們也知道.
有其他的支派.
是應該.
是應份住在.
神安排你們哪個支派.
住哪個地方.
你們是要住在那個地方.
尼希米是很清楚的.
你就可以想像到.
你也可以這樣說.
尼希米.
都是很有分寸地.
凸顯了一個.
待會我們要講的重要的訊息.
就是為什麼住在耶路撒冷.

$^{921}$是這麼重要.
我們回來看這裡.
11章的.
24節.
猶大兒子謝拉的子孫.
米氏撒別的兒子.
米他希雅.
輔助王辦理猶大民的事.
這個王是說尼希米.
是亞達西西王.
將他分封.
是河西的總督.
或者是猶大的分封王.
當然他不是做王.
是輔助.
尼希米去管理.
就是說.
原來就算你不是住在耶路撒冷.
你都是被納入.
在現在.
這個新的時代裡.
被管理的.
其中的一個分子.
你是猶大.
各處散居的.
而且那個散居.
是神安排了.
哪個支派住在哪個地方.
好,第25節.
至於村莊和屬村莊的田地.
有猶大人住在.
基隆.
有猶大人住在.
基列.
基列阿巴和屬基列阿巴的鄉村.
底本和屬底本的鄉村.
葉夾舌.
和葉夾舌的村莊.
耶穌亞.
摩拉大.

$^{961}$伯伯列.
哈薩書亞.
別士巴.
全個以色列最南面的地方.
再下一點點就是.
埃及.
和屬別士巴的鄉村.
洗格拉.
米哥拿和屬米哥拿的鄉村.
欽林門.
索拉.
耶末.
撒羅亞.
阿杜蘭.
就是大衛逃難的時候.
猶大山洞.
阿杜蘭洞.
就是一個比較出名的地方.
和屬拉吉的田地.
亞西加和屬亞西加的鄉村.
他們所住的地方.
是從別士巴直到恩聯谷.
變雅敏人從迦巴起.
住在密末.
密末是一個資種營.
就是軍糧地方.
一個叫阿雅.
伯特利和屬伯特利的鄉村.
是出現很多先知的.
第32節.
阿拿特羅伯.
亞蘭雅.
夏索拉瑪.
基他音哈迭.
洗編尼伯拉.
羅德阿羅.
匠人之谷.
利未人中有幾班.
曾經住在猶大地.
歸於變雅敏的.

$^{1001}$很巨細無遺.
由耶路撒冷.
一直說到猶大的.
城邑.
各處.
是神分了地的.
不同支派.
因為重點不是想那麼詳細.
所以只是一個.
宏觀地說了.
大部分人.
住在猶大各處.
但是呢.
有些再細的鄉村.
都不是大城市.
列舉一大堆的名字.
我們都不太熟悉.
來到說.
有他們散居的地方.
在那些鄉村.
在那些細小的地方.
偏遠的.
距離耶路撒冷相當遠的地方.
當然相當遠.
都不是說要走半年.
只是很細的.
只是步行.
幾天就走完.
由開始.
耶路撒冷.
很緊張住在耶路撒冷的人.
然後.
要求他們.
每十個抽籤.
抽一個出來.
那個人就要住在耶路撒冷.
然後呢.
其他九個.
按十二個支派.
這樣算.

$^{1041}$都有很多人.
因為有五萬多人曾經上來.
他再繁衍的後代.
就沒有再強調有多少人.
就分散在.
猶大的地方.
開始來到我們回頭去看.
一個很重要的訊息.
如果第八到第十章.
是有意識地.
記載.
幾次的嚴肅會.
第七章就很重要了.
因為第七章.
建好城牆之後.
尼希米就發現.
這個已經有門有山的.
耶路撒冷聖城.
欠缺一件事.
就是欠缺居民.
第七章的第一節.
這是以後.
波斯王.
阿達薩西年間.
有個以斯拉.
對不起我記錯了是以斯拉記.
尼希米記的第七章.
因為太相近了.
第七章.
沒錯了.
城牆收完了.
尼希米.
安了門扇.
守門的.
歌唱的.
和利尾人都已經派定了.
我就派我的弟兄.
哈拿尼和.
仍留的宰官哈拿尼亞.
管理耶路撒冷.

$^{1081}$因為哈拿尼亞.
是忠信的.
又敬畏神.
過於眾人.
我吩咐他們說.
等到太陽上升.
才可以開.
耶路撒冷的城門.
人上看守的時候.
就要關門.
上山.
也當派耶路撒冷的居民.
各按班次.
看守自己的房屋.
對面之處.
要警醒.
要按班次.
有些當夜的.
就要起來守住.
重要的是.
因為耶路撒冷有很多門.
東南西北.
不同方位的門.
所以.
提防外邦人.
會來騷擾我們.
甚至.
要來引起騷動.
或者.
一些.
其他的.
不好的勾當.
不過留意第四節.
城是廣大.
其中的文卻稀少.
房屋.
都沒有建造.
我的神.
感動我心.
朝聚貴就官長.

$^{1121}$和百姓.
要照家譜計算.
我就找著第一次上來之人的家譜.
然後就開始記載.
實際上.
和第一次上來的.
這個家譜.
相差多遠呢.
有哪些支派.
特別好生養呢.
有哪些支派.
人丁減少了呢.
或者已經基於.
那個困難.
都移民走了.
都不再在耶路撒冷了.
甚至乎都不再住在以色列了.
經過了百多年.
真是.
人事變遷.
所以你起碼都很有智慧地.
拿回上一次的家譜.
來比對一下.
當然他也會有一個智慧.
就是做一次.
人口的普查.
看一下實際上.
現在有多少人.
更重要就是.
要找回合資格.
無論是.
基於家譜.
基於.
祭司.
的家族.
可以來去公職.
有些如果他.
不知道自己是不是利未人.
是不是利美支派呢.
就要由祭司興起之後.

$^{1161}$就用污名土名.
來去解決.
那個疑難.
如果他OK的到時就給他做.
但如果現在還不行的.
就暫時不要住.
不是針對你.
是要解決了.
希望我們這個.
新的時代的聖殿敬拜.
是要有專職的.
祭司利未人來去做.
這樣就能夠符合.
摩西的律法書和.
神的那個吩咐.
某程度上就是說.
都是在這個過程.
有一個的恩選.
有一個的提升.
就是讓那些.
無關重要的人.
就不要再霸佔住.
祭司事奉的崗位.
如果以前.
因為根本大家都.
沒有所謂.
你是不是祭司都做了.
根本都沒有人做.
你喜歡做就給你做了.
就好像沒有問題.
可以運作.
但是利未人就認定.
這樣是不適合聖經的.
不適合摩西的律法書.
必須要有.
那個專職.
來去做.
所以你可以看到.
第七章近尾.
就是第63節.

$^{1201}$第七章近尾63節.
祭司當中哈巴雅的子孫.
哈哥斯的子孫.
巴西來的子孫.
因為他們的先祖.
娶了基列人巴西來的女兒為妻.
所以起名就叫巴西來.
這三家的人.
在族譜當中.
尋查自己的譜系.
但是尋不到.
因此算為不潔.
這個不潔就不是不乾淨.
而是在那個身份上.
你不是祭司.
不是利未人.
所以就有一個規矩.
除非神有新的突破安排.
否則.
我們要恢復.
摩西的律法書的教訓.
恢復神的心意的時候.
你們就.
無論你們以往做的.
現在都算為不潔.
所以就是說.
不准供祭司的職任.
然後.
第65節.
省長就跟他們說.
不止暫時不可以做祭司.
不可以吃至聖的物.
什麼叫至聖的物.
就是在聖殿.
獻祭完.
祭司吃剩的.
這些下欄.
或者是祭司.
如果你是祭司.
你就是獻祭的祭司可以吃.

$^{1241}$祭司的家人也可以吃.
平安祭的祭肉.
或者有贖罪祭.
都可以吃.
這些就是至聖的物.
或者你獻給神的.
一些掃祭.
又是可以來到.
分給祭司和利未人.
所以.
不是不讓你做.
是直至有用污靈和.
討命,缺夷的.
祭司興起你才可以做.
所以.
比對一下之後.
第66節.
會眾共有42,360名.
另外.
他們的瀑被有.
7,337.
所以總人數.
就是這裡加起來的數.
接著有歌唱的男女.
有245名.
這裡就告訴我們.
現在.
在現存的.
耶路撒冷就有這樣的人數.
數點了.
當然不是一兩個小時就做好.
就是你起碼要.
頒佈了.
一個普查.
就分散了.
不同的族長.
祭司利未人.
來幫忙按照各個支派.
這樣做一個普查.
得到的結論都要寫下.

$^{1281}$來作為.
參照.
所以第七章原來.
是有這個背景的.
不過完了這裡登記之後.
原來突然之間.
引進的是.
第八到第十章的.
嚴肅會.
利希米不是馬上做一些.
好像吃力不討好的工作.
如果.
大家都.
不感受到.
聖神的那個.
福分.
神聖和.
聖殿的獻祭.
需要常規化的重要.
大家.
很勤力地做好一個盛場.
但是大家對.
聖殿的敬拜.
大家對自己神的.
選民的身份.
都不是那麼打緊.
都不是那麼賞心.
都不是覺得是一個榮耀.
是不是移民.
因為我沒有移民.
是不是移民.
如果通過了.
是要拍照.
身份證.
你是屬於美國.
英國的公民.
如果那個國家.
是強大.
是好.
天氣各方面都好.

$^{1321}$你就很開心你能夠成為那個國家的國民.
大家都夢寐以求.
但是現在.
利希米就要提升.
現在.
我們建好盛場.
不要以為那麼簡單.
建好盛場之後.
真正和大家有關的.
現在好像要發一個新的身份證.
大家是神國度的子民.
大家一齊.
來到去對神的.
聖殿.
要賞心.
要敬拜.
要恢復那個秩序.
要恢復那個常規的敬拜.
全民要參與.
這個時候.
好重要.
就是你剛才數了.
能夠.
真的祭司利未人.
真的來提靈那些看門的.
真的那些歌唱的.
其實很少人.
就是說意味著.
第八到第十章.
是透過嚴肅會.
透過憤興會去感動全民.
那個心靈復興.
好像如果真的能夠.
發一張新的身份證.
就告訴你.
你是耶路撒冷的居民.
你是神的子民.
你現在.
是有福的.
因為你在.

$^{1361}$這個聖城.
在這個聖殿裡面.
你才有份來敬拜.
外邦人.
那些想進來都不行.
他們如果不加入.
猶太教他們是不可以進來.
這樣來.
做一個分別為勝.
這樣來.
準備國民的心.
現在.
很多人要拿特區護照去旅行.
排到滿了.
排到聖誕節.
都沒有你的份.
如果你不拿.
疫情又不拿.
復常你又趕快去拿.
回鄉證特區護照都拿不到.
去不了旅行.
但現在.
你給的是.
新的時代.
一個很重要的.
神子民的福份.
那個身份.
你可以想像得到.
如果要恢復.
聖殿的敬拜.
不只是祭司你們的人.
是神.
即是意味著神.
喜悅那些.
甘心樂意.
將自己.
擺在耶路撒冷的人.
神願意.
那些.
不是專職的.

$^{1401}$事奉者.
參與事奉.
如果你能夠事奉.
你要看為是一份honour.
不是責任.
那麼簡單.
不是你喜歡做就做.
不做就不做的可有可無.
是一個很大的honour.
是神給事奉的人.
的一個honour.
首先就是.
如果你說自己是祭司.
對不起.
家譜尋查不到你.
你就不可以公職.
省長說.
你連致勝的物.
都不可以吃.
上個星期吃了.
上個星期的事.
現在人口普查.
你現在不准吃.
警告了.
你有什麼頭暈身癢.
就不要賴我.
這樣說了.
真的很嚴肅.
有一個分水嶺.
尼熙米建好城牆.
之後就數點文數.
他的智慧.
就在這裡.
他不是不知道.
數點完之後.
他自己說了一句話.
這個城.
居民稀少.
房屋都未建造.
即是說.

$^{1441}$以前就沒有掩雞籠.
自出自入.
現在這個聖城.
已經分別為聖.
但是.
居民還沒有這個.
意識.
但是一個領袖.
一個原來一心要重建.
聖文的一個領袖.
他就看到.
百姓看不到的事.
就是.
需要居民.
怎樣需要呢.
是.
好像希特拉那樣.
抓那些人去哪個地方.
強抓猶太人.
離開柏林.
還是怎樣呢.
曾經是.
長江要截流.
在三峽大壩.
有些人.
在三峽大壩.
有些地方將會.
水淹了.
白帝城淹了.
要將那幾十萬居民.
搬去別的地方.
這些是有需要的.
強迫搬遷.
你就會明白.
一定有人不喜歡.
也有人會配合.
但沒辦法.
這是工程的原因.
但是現在.
在什麼形態下.

$^{1481}$尼希米其實空有成竹.
他上來也是一段長時間.
那些人住在那裡.
其實他也開始有一個.
的福祉.
但要強制那些人住在耶路撒冷.
而你不提升.
耶路撒冷的.
神聖的祝福.
而你不提升國民.
擁有.
你是神子民的身份.
這個新時代的身份證.
這個榮耀.
其實是很難的.
那怎樣能夠讓國民.
感受到自己是神選民.
是一份優越呢.
是和外邦人.
是不同的.
怎樣能夠讓他們.
甘心樂意住在耶路撒冷呢.
很重要的就是.
不是立刻開始.
強制那些人.
上去居住在耶路撒冷.
你可以想像得到.
你已經是第三代.
第四代住在.
猶大的不同的村莊.
要你突然離開.
其實很難.
但十一章已經說了.
真的有一個.
這樣的要求.
大家逢十秋一.
就要住在耶路撒冷.
那就.
如果你的支派真是.
多人的就多一些.

$^{1521}$支派小的就少一些.
似乎都是一些藍丁.
那你就可以想像得到.
第八第九第十章.
就是先做好.
百姓那個.
再一次回到神的話語.
去重拾.
到底我們.
建造城牆.
是所謂何事呢.
我們這個民族.
被老歸回回來.
到底神是不是真的眷顧.
我們這個以色列.
這個希伯來的民族呢.
還是我們任外邦人.
繼續嘲笑我們呢.
要恢復.
有敬拜的幸福感.
就需要.
首先就是.
原居民.
不是趕你們走.
而是恢復你們.
對自己的身份的認同.
重拾這個.
我們是神百姓.
我們不是和外邦人.
一般的神.
的見識.
我們是屬於真神的.
子民來的.
所以.
這個七月的嚴肅會.
就分開來得.
珍貴重要.
就是鬆開百姓的.
心靈的.
那個.

$^{1561}$又不是枷鎖.
鬆開他們的內心.
讓他們大家.
感受到自己是神的選民.
那個.
的身份.
因為有了這個身份認同.
很多東西就方便說.
很多東西再教導.
或者有些要求.
逢十進一就有要求.
我就不願意了.
你叫我不可以住在那裡.
要住在這個地方.
又不是和很多人熟的.
不行.
個個都繼續這樣.
我都不要.
關我什麼事.
我又不是什麼人.
但是當那個靈命.
受神的話語.
那個鼓勵.
感動.
和心靈的復興之後.
很多東西就.
真的相對簡單.
所以原來第七章.
尼希米.
看到這個現況.
他和.
以色列都不立刻.
去做一個.
搬遷.
做了人口普查就應該搬遷.
聖殿.
要恢復那個敬拜.
就要抓些人來住.
等到耶路撒冷多人.
多人自然就會熱鬧.

$^{1601}$原來不是.
從一個行政的指令.
原來先從.
那個屬靈的.
方向去做那個工作.
原來先是.
從禱告開始.
這些好像很老套.
老生常談.
個個都是這樣說的.
但是真是身體力行.
真是大家.
來到這個水門事件.
水門前的寬闊處.
在那裡由早上到中午.
來到大家.
守嚴肅會.
大家來到去.
有感動讀摩西的律法書.
有阿弼神的話語感動.
有利未人的教導.
有一起查經.
講解.
做那個詮釋.
解釋摩西的律法書.
怎樣應用.
接著百姓的心就熱了.
去到駐彭捷都意猶未盡.
去到7月24日.
這是第九章.
7月24日繼續守節.
這次就再更加禁食.
再更加就是祈禱認罪.
就是簽約.
大家願意簽約.
來遵守摩西的律法書.
大家來到去.
這是第38節.
第九章的第38節.
因這一切的事.

$^{1641}$我們立了一個確實的約.
寫在冊上.
我們的首領利未人.
和祭司.
都簽名.
這樣的一個準備.
原來我們看利希米記.
就很有立體感.
很有那種.
淑寧的眼光和智慧.
利希米原來不是硬梆梆的一個酒精.
原來他很淑寧.
原來他排除了肉麻.
那些的責難.
那些的誤會.
來去起成場.
起好之後.
是要將外邦人.
隔離開神的選民.
讓神的選民.
開始自己分別為聖.
不是離地.
而是你應該有的身份.
你就要具備這個身份的氣質.
那個的操守.
既然是安息日.
那你為什麼還要跟外邦人買賣呢?.
為什麼你不專心去敬拜神呢?.
就是這麼多.
那你也不行.
不是不行.
不是不買就會死.
是那個心已經枯乾了.
對淑寧的事已經不感興趣.
所以什麼都有興趣.
因為對淑寧沒有興趣.
因為淑寧沒有興趣.
但如果靈性真的有神的感動.
慢慢就不太介意其他東西.
今晚大結局.

$^{1681}$一定要先看完.
其實可以不需要.
有錄影.
有YouTube更加好.
喜歡多看幾次都可以.
但這個最重要.
心靈受神聖靈的感動.
這是最難的.
大家都知道教會要復興.
但香港的教會有多少間.
有多少人.
有多少年能夠持續的復興呢?.
有多少人能夠持續的復興呢?.
真的很少.
我們以為人數過千就要復興.
也是的.
沒理由只有十幾個人就說要復興.
數字有時是反映一些事實.
但數字不一定是.
而是那些人被神的愛.
被神的真理.
被神的福音感動.
是甘心樂意去做事.
這個甘心樂意是很重要.
真的不是開玩笑.
教會從來不是沒有人就請人.
教會是一班被神的愛感動的人.
甘願回來.
不是叫洗淨圓華.
是放下自己的身段.
一到星期六.
你是大機構的老闆.
沒問題的.
你放下你的身段.
你是弟兄姊妹.
來服侍.
按神給的恩賜.
大家彼此信服的服侍.
就不是你繼續回來指指點點.
就不是叫服侍.

$^{1721}$是你繼續要來佔據一個平台.
去做你自己內心的私心.
就不是真正的事奉.
真正的事奉就是在愛裡.
彼此信服的服侍.
所以這個是最困難的.
所以我真的看見.
尼熙米記真的令到我很感動.
一個領袖.
不是運用他的強權.
去壓迫百姓.
為了現在我們要抗衡.
那些外邦人.
我們一定要這樣做.
原來都是需要大家.
那個靈明要配合.
那個靈明還沒有去到那個水平.
講的那些東西.
變成很高大的.
高言大字.
聽的聽不下.
沒興趣聽.
更加不屑.
你的安排.
怎樣就怎樣.
這樣就很麻煩.
教會是一個彼此信任.
互相相愛.
的形態.
之下.
才做到聖公.
教會是這樣的.
要大家交心出來.
在一個互相信任底下.
去服侍.
不問那個代價.
不介意那個.
得與失.
不介意那個辛苦.
那個不辛苦.

$^{1761}$不介意那個出名.
那個不出名.
不介意那個位置.
我一定要爭.
大家就是要這樣.
那個聖公才能夠做到.
所以你看看第八到第十章.
講完沒有人.
什麼都不理.
就開嚴肅會.
一路開到7月24.
接著做的.
就很簡單.
我們剛才第十一章所讀的.
就真的很簡單.
就是百姓的首領.
他們甘心樂意住在耶路撒冷.
他們.
有他們的官職.
有他們的專責的事奉.
他們是甘心樂意去住.
其餘的.
百姓就.
祭籤了.
祭籤是因為甘心樂意.
馴服尼希米.
和以色列.
他們的聖靈的感動.
的那個帶領.
尼希米很早就已經意識到.
其實百姓自己都意識到.
耶路撒冷都沒有人.
聖殿需要.
恆常的服侍.
聖殿的周圍.
需要有不同的守門.
又需要有不同的職位.
幫手去處理.
聖殿的大大小小的工作.
這些就是.

$^{1801}$很合理的事.
但是.
尼希米就是現在.
在百姓.
願意服從神之下.
來服從尼希米.
現在尼希米就說.
公平一點.
凡是十個就一早.
抽籤抽一個出來.
就是大家都很願意馴服.
尼希米的.
這樣的安排.
就是一些.
不是叫做自由工作者.
是按照著.
恩慈.
按照著興趣.
來到他在神家裡面服侍的.
弟兄姊妹.
是一份最大的.
那個需要.
就是甘心樂意.
新約出現好幾次.
和大家稍微揭一揭.
哥林多前書.
第十六章.
如果我們沒有.
這樣的心態去事奉.
遲早你都會.
枯乾和.
會出現.
各種的問題.
例如你不肯放下你的崗位.
這就是亂斬你的權力.
第十六章的.
第十五節.
或者第十四節.
對不起.
第十四節開始.

$^{1841}$你們務要警醒.
在真道上站立得穩.
要作大丈夫要剛強.
凡是你們所作的.
都要憑愛心而作.
弟兄們.
你們曉得斯提反一家.
是亞該亞初結的果子.
並且他們專以服侍聖徒為念.
我勸你們.
要馴服這樣的人.
並一切.
同工同勞苦的人.
斯提反.
和福圖拿都.
和亞該古.
來了我這裡.
我就很喜歡.
很歡喜.
因為你們待我有不及之處.
他們已經補上了.
斯提反一家.
第一章好像有提到.
第一章提到.
先看看哥林多前書.
沒錯.
一章的.
第十節.
讀多一點讓大家對上下文有掌握.
弟兄們.
我藉著我們主耶穌基督的名.
勸你們都要說一樣的話.
因為有分門結黨.
你們中間.
都不可以分黨.
只要一心一意.
彼此相合.
因為格萊士嘉里的人.
曾經對我提起.
弟兄們來.

$^{1881}$就說你們中間有紛爭.
我的意思是什麼呢.
就是你們各人就說.
什麼叫紛爭.
在保羅的眼中.
我是屬保羅的.
我是屬阿波羅的.
說笑.
我是屬曾牧師的.
我是屬勵士牧師的.
我更厲害.
我是屬謝牧師的.
那就大件事了.
保羅就說我是屬阿波羅.
我是屬基法我是屬基督.
第十三節就很重了.
基督是分開的嗎.
保羅為你們釘十字架嗎.
你們是奉保羅的名受洗嗎.
除了基利斯堡和.
該由我沒有給你們.
一個人施洗.
即是有其他同工.
免得有人說.
你們是奉我的名受洗.
我也給斯提反加斯.
一個洗.
即是說只有三個在哥林多的.
弟兄是保羅.
他那一年半幫他們.
施洗.
即是說有些還沒上完洗禮班.
保羅就要走了.
基於那個壁壩.
該留.
保羅就要離開哥林多.
回去耶路撒冷.
這是第三次的旅行報道.
那就.
留下的人.

$^{1921}$就有阿波羅在.
阿波羅就幫他們施洗.
也有當時的教會領袖.
所以斯提反.
都有幫他們施洗.
即是告訴我們.
現在去到第十六章.
阿波羅林潔未.
勸勉.
哥林多教會的弟兄姊妹.
你們要.
懂得看人.
留意一些.
表面不起眼的人.
他們在神眼中.
是神所稱許.
是神所悅納.
是神所喜歡的人.
這些人做過什麼呢.
再讀一次第十五節.
你們曉得斯提反一家.
是亞該亞省.
初結的果子.
即是說.
很早期保羅去傳道.
已經.
這個斯提反就信主了.
不是那個殉道的斯提反.
同名同姓.
就已經跟隨.
保羅在哥林多.
這裡告訴我們.
並且.
他們.
即是他們的家人.
是專以服侍.
聖徒為念.
放在心裡.
時時記掛叫做為念.
專門.

$^{1961}$專心這樣去做.
即是說.
那雙眼是留意.
我有什麼可以再做得好一點.
去服侍聖徒呢.
服侍聖徒即是說.
為了神去服侍人.
有些人受迫迫.
有些人可能有欠缺.
諸如此類.
在教會.
斯提反一家.
他們默默的.
在教會裡有一個念頭.
什麼念頭呢.
服侍為念.
服侍神的兒女為念.
專心留意.
有哪些人需要我們服侍.
為念.
這樣為念.
一句而已.
已經在保羅的口裡.
去誇讚.
你可以想像得到.
譬如我們站一站.
馬太福音第六章.
耶穌的鐵律.
告訴我們.
哪些人可以得到.
神的賞賜呢.
第六章的第一節.
你們要小心.
不可以將善事.
走在人的面前.
故意要得到.
人看到.
如果是這樣的時候.
就不能得到.
最重要的賞賜.

$^{2001}$就是你們天父的賞賜.
所以主耶穌教.
你施捨的時候.
就不可以在你前面.
吹號.
好像那些假冒為善的人.
在會堂.
在街上.
鎖閒的.
為什麼呢.
故意.
和默默地做.
一善之格.
是端乎你的.
心底的動機.
耶穌一針見血.
服事.
完全不是錯.
施捨是很對的.
但原來施捨.
這些宗教行為.
是可以帶來沽名釣譽.
可以帶來假冒為善.
可以帶來不單止.
得不到天父的賞賜.
甚至乎會有.
嚴重的後果.
故意要得.
人的榮耀.
因為就是說緊張人.
多於緊張神.
怎麼看這件事.
主耶穌說.
我實在告訴你們.
他們已經得到他們的賞賜.
你施捨的時候.
就不要叫左手.
知道右手所做的.
左手不要.
讓右手知道.

$^{2041}$所以你不需要.
遮住.
不需要.
去奉獻.
遮住.
不需要.
這是開玩笑的.
告訴我們.
我們做了也不需要.
放在心底.
整晚睡不著.
很開心.
誇耀我.
今天誇我.
不要讓人誇耀.
而是要.
尋求天父的喜悅.
所以主耶穌第四節說.
要叫你施捨的事.
行在暗中.
理乎在暗中察看.
心靈的察看.
所謂暗中不是躲起來.
而是你做了.
一般人都看不到.
是普通人留意的.
但天父知道的.
因為我們的心靈動機.
是天父察看的.
必然的報答你.
這裡連續出現了三次.
第四,第六,第十八節.
都出現了三次.
天父暗中察看.
必然報答你.
講到三個宗教行為.
施捨,禱告和禁食.
所以我們要.
禱告和禁食.
所以.

$^{2081}$我們回看.
《喜米記》.
這個甘心樂意做.
是需要培養.
需要一個.
靈性的復興.
才能做到甘心樂意.
很開心.
當然不是百成千人.
填那個恩賜調查表.
但填的人.
是很有心.
而且.
我不是太認識的.
是新來的弟兄姊妹.
很開心.
牧者我們很開心.
我真的認真祈禱.
很想在這個教會.
一起服侍.
我們要尋求的.
就是這樣的甘心樂意.
我們不介意.
有些事沒有人做.
因為如果一個不稱職的人.
所謂不稱職.
是他的心靈.
未真正被聖靈感動.
未真正的.
甘心樂意做.
其實做下去.
都不會帶來一個.
影響力.
一個很有心的人.
做的聖功.
是會帶來一個很正面的.
影響力.
弟兄姊妹我們要培養這一種.
的甘心樂意.
我們要像.

$^{2121}$斯提法一家.
保羅在這.
事奉保羅.
跟隨保羅.
因為看到這個神的僕人.
是講純正的福音.
他們去落力.
一起去傳.
但開始有人嘲笑.
攻擊保羅.
保羅就離開.
去到提多尤士都的家.
那裡來.
事奉.
那是會堂.
管會堂.
離開會堂.
去到其他地方.
很多外邦的哥林多人信主.
後來保羅就在那裡.
住一年半.
自然才離開.
等等.
但這個斯提法.
一家人.
繼續默默的.
不是打響鑼鼓的.
是來到甘心樂意的事奉.
去做.
保羅寫這封哥林多前書.
有這麼多分門結黨的人.
不斷罵.
但他來到.
在臨尾結尾.
叫哥林多教會的弟兄姊妹.
你要懂得看.
哥林多教會裡面.
有很低調.
很敬虔的弟兄姊妹.
他們這樣來事奉.

$^{2161}$是專以.
服侍.
聖徒為念.
一句說話已經將.
那些深度.
散發出來.
你可想像得到.
在神的眼中的斯提法.
一家.
他們是多麼的.
蒙神的祝福.
他們是專以聖徒為念.
弟兄姊妹.
這是我們真的需要.
在一起彼此.
互勉和學習.
所以你難怪.
見到尼希米基提到.
第十一章是說.
真的住在耶路撒冷的人.
是蒙福的.
所以第二節.
凡是甘心樂意.
住在耶路撒冷的.
百姓都為他們祝福.
就不是說.
你做了我不做了.
不是這樣.
是真的為這些在聖城裡面.
事奉.
如果你是有聖殿的.
專職事奉.
即是說你不是那麼輕易.
隨便喜歡放幾天假.
就放幾天假.
某程度我們作為牧者.
有時.
同工拿假.
都盡量禮拜六日都不拿.
這個某程度上.

$^{2201}$好像是大家的約定.
俗成這樣.
除非真的拿大假要離開香港.
我都鼓勵的.
有時最好是拿兩個星期.
拿完整年一次過放.
就可以真正休息.
但有時同工都是一日一日放.
又不是開會的日子.
通常星期四.
或者星期五.
沒有聚會的時間.
為甚麼呢.
因為始終你的心是聯繫著.
聖殿的事奉.
有專職的工作.
都有你的限制.
所以.
事奉神的人.
其實是有一個犧牲的.
但這個犧牲.
在神的眼中會有賞賜.
會有一個祝福和回報.
所以.
如果沒有第八到第十章.
首先就有那個幸福感.
首先就是.
自己是一個神的選民.
你現在.
在新時代裡.
參與這個.
靈性的復興.
你配合.
尼希米 以斯拉.
那個教導.
你有份參與.
只是你搬一塊石頭去建那個城牆.
你都已經有份.
你已經有福了.
現在完整的工程.

$^{2241}$大家.
各自歸家.
但是尼希米.
釋放一個很重要的訊息.
就是.
聖殿不應該.
做好像以前.
我未回來之前的光景.
是應該有.
準備.
是應該有專職的人.
既然是專職.
弟兄姊妹 百姓.
就要按照住.
他們的收入.
十份書一.
這是第十章的教導.
十章的.
第三十五節.
又定了每年.
將我們地上初熟的土產.
各樣樹上.
初熟的果.
都奉到耶和華的殿裡.
又照律法上所寫.
將我們投胎的兒子.
守生的牛羊.
都奉到我們神的殿.
交給我們.
神殿裡.
公職的祭司.
並將初熟的物質.
草磨的棉和.
舉祭各樣樹上.
初熟的果子.
伸走和游.
奉給祭司.
收在我們神殿的庫房裡.
將我們地裡.
所產的十分之一.

$^{2281}$奉給利未人.
因為利未人.
在我們一切的誠邑的土產.
當取十分之一.
利未人.
取自己的十分之一的時候.
亞倫的子孫.
就當由一個祭司.
與利未人同在.
利未人也當從十分之一中.
取十分之一.
奉到神殿的屋裡.
收在這個庫房裡.
以色列人.
和利未人.
要將五屋伸走.
和游.
全部是神賜福百姓的.
豐收.
工作的成果.
收入的來源.
萬物.
是從上主所賜的.
你要感受到.
這個幸福感.
你要感受到.
你要對號入座.
今次的豐收.
不是因為今年的氣候特別好.
你要在氣候上.
上主的賜福.
你有勞動力.
你能夠在那個季節.
不是突然間.
收成的時候大旱.
或者要播種的時候.
傾盆大雨.
下冰雹.
你要將這一切.
完全歸於.

$^{2321}$神賜福.
我們的土地.
這樣來將這些舉掣.
即是舉起.
呈獻給神.
其實是摩西律法吩咐的.
他們就是簽約有這個.
幸福感.
感受到自己.
回到神選民的.
這個水平.
他們就願意.
既然一個專職的.
聖殿.
專職的事奉者.
就需要有我們的配套.
有百姓的奉獻.
有百姓的.
不多了,十分取一.
來奉獻給神的殿.
那裡來衰庸.
當你上到.
聖殿有好的敬拜.
你們家庭和睦.
牽著手去飲茶.
回家.
週日的,週間的工作又順暢.
兒女又不是.
大病大痛.
各方面的讀書.
升學又有智慧.
又有才幹.
各方面.
接著放假又可以去.
簡單的旅遊.
接著又有很滿足的家庭的幸福.
一切源於.
聖殿的敬拜.
得到那個祝福.
然後將這個祝福的.

$^{2361}$幸福感化做.
他們的生活的動力.
然後將生活的成果.
一個循環.
奉獻給神.
去到神的面前.
所以何時教會.
又虧欠又欠缺.
都是大家那個幸福感.
不足夠.
導致大家都.
不願意奉獻給神.
不甘心樂意的.
按照神給你的.
恩典.
祝福來奉獻給神.
再嚴格的說.
就是.
不感受到神.
給我的.
是祂的祝福.
是我的辛苦,是我的努力.
是我的投資智慧.
是我的幸運.
是我的家族的富有.
不是神的祝福.
就沒什麼好說的了.
所以牧者要做的.
是開懇大家的心靈.
讓大家感受到.
我們在福音裡面.
我們是有很大的幸福.
我們不勉強.
也都不需要去.
募捐.
不需要去求你.
你不奉獻.
不緊要.
你自己將來面對.
做你的造物主.

$^{2401}$你不是要面對我.
我沒權去審你的.
不關我的事.
我有責任勸導弟兄姊妹.
我們要奉獻給神.
向神.
有我們終生的奉獻.
所以.
這裡就是讓我們見到.
去到第十一章.
就是百姓.
真的為住在耶路撒冷.
抽籤抽到的.
來去祝福.
他們才是一班.
受祝福的人.
有專職的.
有甘心樂意的.
被抽中原來不是那個生死籤.
原來是常常籤.
真是開心.
我抽到.
我們抽居屋就是.
我真是很開心.
中六合彩.
抽到住在耶路撒冷.
所以.
今天我們教會需要的是.
幸福感.
我們教會需要的.
不是你家裡擁有多少部車.
我知道有弟兄姊妹.
不止一部的代用車.
甚至乎都不是.
那些很基本的代用車.
不同的款式的車都有.
每年交牌費都不少.
我不會仇富的.
每個人按照住.
自己的生活的條件.

$^{2441}$自己的需要.
那些.
那個不是對錯.
內心不感覺到.
神的幸福感.
那個就真是問題的所在.
那個就需要.
我們大家一起彼此的護面.
彼此的鼓勵.
所以百姓為這些人.
甘心樂意來去祝福.
然後就是剛才我們講的.
其他的不同的.
即時.
不同的.
事奉的.
他們就住在猶大的各處.
有配套.
有其他的工作.
就是說他們有他們的耕種.
有他們其他的.
可能是打獵.
有其他的可能是做捕魚的.
有其他的是做.
工匠的.
大家按照著神給的恩賜.
大家的努力.
大家的恩勤.
然後每個主日.
安息日.
大家來朝見神.
分別為聖.
在整個國家.
先是重要.
這個聖神的分別為聖.
是和外邦人.
但是.
那個還是一個太抽象了.
不只是那個物質.
而是百姓的心靈.

$^{2481}$自己要分別為聖.
要感受到神子民的.
新身份的幸福.
那就分別為聖了.
由你自己.
自發不想.
去守安息日.
不想在安息日.
去做買賣.
去和外邦人.
寧願打牙架都不回去.
敬拜神.
你就可想而知.
那個幸福感就要提升了.
盼望這個是教會.
我們的.
一起去成長一起去追求.
有時星期日早上.
見到弟兄姊妹不回教會.
又知道他們不是回晚課.
回了早課而是真的不回教會.
你好像都不太好意思.
喂回教會今天星期.
有時好像都不太好意思說.
因為始終香港人都是比較.
即是.
禮貌.
但其實某程度我們的身份.
是應該勉勵的.
只顧著跑步又不回教會.
這樣不行的.
敬拜神.
所以要提升這個幸福感.
要讓弟兄姊妹.
感受到回教會是.
蒙福.
感受到自己是神的選民.
是蒙福.
然後能夠.
去到在教會裡面.

$^{2521}$在這個場景.
感受到自己是分別為聖.
和開始參與.
神感動你的那個事奉.
所以司提反一家.
阿保羅都能夠.
離開了很久.
他現在應該是.
在二佛所.
來寫哥林多前書.
你可以想像得到.
他還記掛著.
司提反這一家人.
他們那種敬虔.
愛主.
蒙神的祝福.
蒙神的紀念.
這樣的一個家族.
我們一起禱告.
再一次多謝主.
我們可以停了兩次之後.
一起再去看.
第十一章的尼希米記.
我們真的很盼望.
主的聖靈感動我們.
讓我們感受到.
我們是蒙福的.
我們不是自己說自己是蒙福.
而是我們是神的選民.
是主耶穌用.
重價買贖.
我們真是神的兒女.
這個身份.
我們是蒙福的.
幫助我們感受到.
這個幸福感.
幫助我們能夠投身在.
神的國度裡面去事奉.
願主幫助.
透過這一章聖經.

$^{2561}$成為很多人的提醒.
祝福和勸勉.
獻上感恩禱告.
奉耶穌基督的名.
阿們.
\newpage



\section{列王記上 19:1-18}
\label{sec:W6aFxBk_o0s}
\textbf{2023-07-02 加恩樓《當以為走到絕境時》列王記上 19\_1-18 講員:鄧泰康傳道}
\newline
\newline
連結: \href{https://youtube.com/watch?v=W6aFxBk_o0s}{\texttt{https://youtube.com/watch?v=W6aFxBk\_o0s}} ~~~~ 語音日期: 2023-07-05
\newline
\newline
\hyperref[sec:G4jQV76in2M]{\small{< < < PREV SERMON < < <}}
~
\hyperref[sec:index_chronic]{\small{[返順時目]}}
~
\hyperref[sec:index_scriptual]{\small{[返順卷目]}}
~
\hyperref[sec:Vwm387KuCtc]{\small{> > > NEXT SERMON > > >}}
\newline
\newline
列王記上 19:1-18
\newline
\begin{longtable}{cl}
\hline
\hline
章節 & 經文 (和合本修訂版)\\
\hline
19:1 & \begin{tabularx}{0.7\textwidth}{X} 亞哈把以利亞一切所做的和他用刀殺眾先知的事都告訴耶洗別。 \end{tabularx} \\ \\ \relax
19:2 & \begin{tabularx}{0.7\textwidth}{X} 耶洗別就派使者到以利亞那裡，說：「明日約這時候，我若不使你的性命像那些人的性命一樣，願神明重重懲罰我。」 \end{tabularx} \\ \\ \relax
19:3 & \begin{tabularx}{0.7\textwidth}{X} 以利亞害怕，就起來逃命，到了猶大的別是巴，把僕人留在那裡。 \end{tabularx} \\ \\ \relax
19:4 & \begin{tabularx}{0.7\textwidth}{X} 他自己在曠野走了一日的路程，來到一棵羅騰樹下，就坐在那裡求死，說：「耶和華啊，現在夠了！求你取我的性命吧，因為我不比我的祖先好。」 \end{tabularx} \\ \\ \relax
19:5 & \begin{tabularx}{0.7\textwidth}{X} 他躺在羅騰樹下睡著了。看哪，有一個天使拍他，對他說：「起來吃吧！」 \end{tabularx} \\ \\ \relax
19:6 & \begin{tabularx}{0.7\textwidth}{X} 他觀看，看哪，頭旁有燒熱的石頭烤的餅和一壺水，他就吃了喝了，又再躺下。 \end{tabularx} \\ \\ \relax
19:7 & \begin{tabularx}{0.7\textwidth}{X} 耶和華的使者回來，第二次拍他，說：「起來吃吧！因為你要走的路很遠。」 \end{tabularx} \\ \\ \relax
19:8 & \begin{tabularx}{0.7\textwidth}{X} 他就起來吃了喝了，仗著這飲食的力走了四十晝夜，到了神的山，就是何烈山。 \end{tabularx} \\ \\ \relax
19:9 & \begin{tabularx}{0.7\textwidth}{X} 他在那裡進了一個洞，在洞中過夜。看哪，耶和華的話臨到他，說：「以利亞，你在這裡做甚麼？」 \end{tabularx} \\ \\ \relax
19:10 & \begin{tabularx}{0.7\textwidth}{X} 他說：「我為耶和華－萬軍之神大發熱心，因為以色列人背棄了你的約，毀壞了你的壇，用刀殺了你的先知，只剩下我一人，他們還要追殺我。」 \end{tabularx} \\ \\ \relax
19:11 & \begin{tabularx}{0.7\textwidth}{X} 耶和華說：「你出來站在山上，在耶和華面前。」看哪，耶和華從那裡經過。在耶和華面前有烈風大作，山崩石裂，耶和華卻不在風中；風後有地震，耶和華也不在其中； \end{tabularx} \\ \\ \relax
19:12 & \begin{tabularx}{0.7\textwidth}{X} 地震後有火，耶和華也不在火中；火以後，有輕微細小的聲音。 \end{tabularx} \\ \\ \relax
19:13 & \begin{tabularx}{0.7\textwidth}{X} 以利亞聽見，就用外衣蒙臉，出來站在洞口。聽啊，有聲音向他說：「以利亞，你在這裡做甚麼？」 \end{tabularx} \\ \\ \relax
19:14 & \begin{tabularx}{0.7\textwidth}{X} 他說：「我為耶和華－萬軍之神大發熱心，因為以色列人背棄了你的約，毀壞了你的壇，用刀殺了你的先知，只剩下我一人，他們還要追殺我。」 \end{tabularx} \\ \\ \relax
19:15 & \begin{tabularx}{0.7\textwidth}{X} 耶和華對他說：「去吧，從原路回去，往大馬士革的曠野去。到了那裡，你要膏哈薛作亞蘭王， \end{tabularx} \\ \\ \relax
19:16 & \begin{tabularx}{0.7\textwidth}{X} 又膏寧示的孫子耶戶作以色列王，並膏亞伯‧米何拉人沙法的兒子以利沙作先知接續你。 \end{tabularx} \\ \\ \relax
19:17 & \begin{tabularx}{0.7\textwidth}{X} 將來逃過哈薛之刀的，必被耶戶所殺；逃過耶戶之刀的，必被以利沙所殺。 \end{tabularx} \\ \\ \relax
19:18 & \begin{tabularx}{0.7\textwidth}{X} 但我在以色列中留下七千人，是未曾向巴力屈膝，未曾親吻巴力的。」 \end{tabularx} \\ \\
[1ex]
\hline
\hline
\end{longtable}
$^{1}$(孫)好,電影節目早安.
辛苦書文了,唱完這麼久,讀了這麼長的一段經文.
這段經文其實大家如果肯看聖經,你很熟悉的.
其中一個聖經裡面會不斷講的故事,特別是小朋友的聖經故事更加會講.
但是當你看十九章的時候,你必須要和十八章一起看.
為什麼呢?你會發覺十八章到十九章其實是充滿著好戲劇性的.
你可以說十八章的以利亞是一個什麼樣的人呢?.
是一個意氣風發,因為當時其實他聽命於上帝.
他很以自己為先知去作為他的榮譽,所以上帝跟他說什麼他都相信.
之前發生什麼事呢?就是因為阿哈他們做得不好.
上帝要透過一些方法去懲罰或者叫他提醒他們.
有三年的時間沒有下雨,但是上帝就透過一個寡婦去養活了以利亞這段時間.
到了上面的時候,上帝就叫以利亞去找阿哈.
之後他們碰到,也按照上帝所對他們的應許下雨.
好了,走到高峰了,就是他竟然跟阿哈說.
你們不是有四百個去拜巴黎的先知嗎?.
還有四百五十個去拜阿斯拉的先知嗎?.
你叫了他們來,如果大家還記得,這是一場非常驚天動地.
你可以想像一下,八百五十個先知在山上,然後是一個很虛暗的情況.
然後就發生了一場很轟天動地的事.
就是叫他們,既然你們的神那麼厲害,你們就叫他們.
現在我們每人預備一個祭身,你們就叫祝壇,我們一起求他降火下來燒起祭身.
然後他們形容,當時他們跳舞了,大家可以想像一下,非常混亂.
以利亞還很厲害,還跟人說,你們的神是不是睡了覺?大聲一點吧.
但是都沒有反應.
到最後的時候,他竟然將自己的祭身淋了一桶水下去.
還要是濕的,然後就祈禱,上帝就降火燒了祭身,是不是很震撼?.
接著的一時之間,八百五十個先知全部被殺.
你可以想像到,他對於上帝那種信,他都有一個怎樣的人生的高峰.
但是竟然第十九章一開始的時候,是急轉直下.
如果大家剛才有留意的時候,只不過是阿哈告訴他老婆.
然後他老婆就差遣人跟他說了一句,記住是一句.
如果明天這個時間,我不要你像剛才那群人死的話.
我就不升爺,大概就是這個意思.
在這個情況下,他就走了.
我不知道大家裡面有沒有很多問號,看聖經一定要找這些問號出來.
剛才那個以利亞去了哪裡?.
為什麼耶駛別一句話,他竟然跌到一個這樣的谷底?.
你會發覺有趣的地方就是作者沒有跟我們說究竟耶駛別的可怕是怎樣,從來沒有提.
其實這裡可以給你很多的想像空間.

$^{41}$但是裡面只強調一件事,總言之以利亞去怕耶駛別多於怕上帝.
就是這麼簡單.
或許耶駛別可以是一個很可怕的人.
如果大家以前都會記得,我不知道是不是真的,我不是讀歷史的.
但是以前的譬如是武則天,他們當時的皇帝.
他們那時候用那些殘酷的酷刑,那些皇帝用什麼殘酷的酷刑.
令到那些人聞風喪膽,令到那些人一見到他就會害怕.
我不知道是不是這樣,又說他耶駛別同樣都是,我一是就不說,我一說就會發生.
一定會應驗的,所以他就很害怕,我不知道,作者沒有說.
但是以利亞他就是走.
當我們去看聖經的時候,我一直去想這個問題的時候.
我只可以說似乎是以利亞對上帝的認識毀了他.
或者,我不知道大家有時候聽神學這兩個字,大家就會覺得很高深莫測.
神學的意思就是你對於上帝的認識.
你對於神的認識.
所以你就會發覺其實這裡就是以利亞他對上帝的認識而毀了他.
為什麼呢?因為他一直所認識的上帝.
就好像在十八章所認識的上帝.
你是說了下雨就下雨,不下雨就不下雨.
我也很有面子的,我的上帝說什麼就什麼.
所以我就夠膽去面對著,就算是皇帝也好.
或者是先知也好,我也不怕,我的上帝很厲害的.
他的神學是這樣的.
所以你就會發覺當時的神學是怎樣的.
當他走到這個地方,他即時的就求死.
我想大家留意他求死的內容.
耶和華,算了吧,求你取去我的性命,因為我不勝於我的列祖.
我沒有我的祖先那麼好.
就是說他在這件事上,為什麼會這樣?.
我做得不夠好.
就等於今天我邀請弟兄姊妹你們去事奉的時候.
我還不夠好,你不要叫我了.
我們自己也有一套神學,自己也有一套標準.
我們要有多好,上帝才會選擇我,是不是?.
所以很多時候我們都會覺得我們不夠好.
現在我不知道他到底是真的覺得自己不夠好,還是他嘲笑上帝.
其實這裡語句上沒有很清楚.
上帝,算了吧,你取我性命吧,我都不夠我以前的列祖好.
但是當他一轉到下面,你就找到答案了.
似乎他是嘲笑上帝.

$^{81}$為什麼呢?因為當上帝問他的時候,他說什麼?.
我為耶和華萬軍之神大發熱心.
因為以色列人背棄了你的約,毀了你的壇,用刀殺了你的先知.
只剩下我一個.
你見不見到我多麼為你啊,上帝.
你見不見到現在那群人多麼不濟啊.
是我一個這麼勇敢走出來.
你都不保護我,現在他們還要追殺我.
當耶誓別說這句話的時候,雖然聖經沒有交代.
但是他們心裡,這是什麼人啊?.
耶誓別而已,上帝為什麼不降火,像你剛才一樣,給他們製藥燒死他們就行了.
沒有,沒有發生.
耶誓別仍然可以這麼大權.
所以他就會突然間,他的神學對不上.
我以前所認識的那個上帝去了哪裡.
我想大家留意,他講了兩次.
就是說這是一個非常固有的觀念.
弟兄姊妹,今天我們同樣對上帝的相處,我們這個神學同樣一樣固有.
但其實今天這些你很固有的神學,請你回去重新再好好認識.
究竟上帝是不是像你心中的上帝一樣.
今天這個上帝不等於你心中的上帝.
所以弟兄姊妹,努力一點去讀聖經.
我們的上帝是聖經裡面出來的.
是聖經裡面告訴你他是一個怎樣的上帝.
不是你認為他是一個怎樣的上帝.
而那個開始點在哪裡呢?所以要留意這些.
很多人剛剛信耶穌,下星期又有報道會了.
每一個人信耶穌都有自己的理由.
我想我們最喜歡聽的是什麼呢?.
在一件很大的困難裡面,上帝就出現神蹟幫你解決了.
這個才是上帝,我們的上帝是這麼厲害的.
所以有些人就說,我以前很灰心失意.
但是上帝,我認識你的緣故,而你帶領我的時候.
可以去解決這些問題.
我們的上帝是這樣的.
但是弟兄姊妹,你搞清楚.
這個只不過是上帝在那一刻.
和這個分享見證的弟兄姊妹所相處的那件事.
他在之前曾經和你分享,我有癌症.
上帝怎樣醫治我,我的癌症康復.

$^{121}$那是不是每次當他有癌症他都會康復.
是不是每一個弟兄姊妹當你有癌症的時候你都會康復.
你搞清楚,我想說這些就是神學.
這些就是毀了我們的神學.
當他不是這樣發生在你身上的時候.
你就會倒下.
你就好像以利亞在這裡.
無論你以前多麼相信上帝,多麼跟隨上帝.
你就會倒下.
所以你就會發覺,這裡有趣的地方就是.
上帝在這裡和以利亞有一個完完全全重新的相處.
以利亞在這一刻,他真的覺得他走到人生的絕路.
連上帝你都不再看我了,你再不是我認識的上帝了.
所以上帝你不如讓我死了算了.
但我說很戲劇性的地方,他說是死的.
但上帝說起來吧,吃東西吧,那他又怎樣.
他又吃的.
如果他想死沒理由吃的,你明白我的意思嗎.
吃完之後他又躺平.
上帝又說,你起來吧,吃東西吧,你還有很多路要走.
他又吃,吃完之後又怎樣.
上帝叫他去哪裡,他又去.
這段經文很有趣的.
那我就想問,以利亞你想怎樣.
其實今天我們作為基督徒,我們和上帝的相處很多時候都是這樣的.
我們和上帝的相處裡面,我們充滿著很多很多不明白的埋怨.
覺得很苦,但是某程度上有些事情,上帝你要我做,我又做的.
你叫我事奉,我又繼續事奉.
但是弟兄姊妹,真正幫到以利亞的方法,上帝和她相處的方法.
也是今天可以幫到你的方法.
我不知道大家有沒有發覺這段經文裡面有趣的地方.
是以利亞有以利亞的講,上帝有上帝的講.
大家之間好像沒有什麼溝通,我不知道大家有沒有發覺.
其實在這段聖經裡面,我們要搞清楚的有一個觀念.
就是上帝就是上帝.
其實你會發覺無論以利亞怎麼埋怨他.
第一件事,上帝都是繼續對人施行慈愛的.
他不管你現在發生什麼事.
他要對你的慈愛,他會因為自己的緣故不會間斷的.
所以就算到了那一刻,以利亞已經去到自己想放棄.

$^{161}$上帝所叫她去做的就是什麼?.
你現在最需要的是吃東西.
他就預備食物給你.
弟兄姊妹,這個才是上帝.
無論在那一刻你覺得自己多絕景也好.
上帝從來沒有缺席的.
上帝一直都在.
上帝一直都在為你預備.
上帝很知道你的需要.
所以他講了兩次,你起來吧,吃東西吧.
第二次他做了一個註腳.
因為你要走的路還有很遠.
弟兄姊妹,當你有一些時間突然覺得自己很消沉的時候.
你最需要的是不是這句話.
我作為一個男性,講真一句.
有時候我很煩的時候.
我都不是要別人給我很多意見.
可能我只是想跟他聊聊.
通常我看到就是我跟你說的時候.
突然間你高談闊論,然後給我很多建議.
然後我就走了.
但可能最理想的反而是什麼?.
反而可能就是.
不如我們吃一餐好了.
或者甚至乎我經常提一些.
做太太,看到你丈夫回到家的時候.
他的臉色不是太好的時候.
其實不用問太多的.
倒杯水給他.
或者倒杯他平時喜歡喝的東西給他.
他就明白了.
他就知道了.
其實上帝在那一刻.
他繼續對人施行慈愛.
第二件事.
上帝會繼續去問什麼呢?.
他不理以利亞你跟他埋怨什麼.
他就不斷去問以利亞.
以利亞你在這裡做什麼?.
弟兄姊妹你不要小看這句話.

$^{201}$這句話對以利亞那一刻.
其實回正了焦點是很重要的.
以利亞你為什麼要做先知?.
先知是什麼?.
先知是我的代言人.
也就是說.
由你去蒙召做先知開始.
一直很清清楚楚的就是.
你是跟著我的路去走的.
我要你以前走一條很.
很厲害的路.
贏的路.
你就走吧.
但我現在要你走一條.
面對困難的路.
那怎麼辦?.
繼續走吧.
繼續去吧.
所以我作為傳道人也是的.
有時候做著做著也會有灰心的.
但很救我的那句話是什麼?.
問你的初心.
為什麼你要做傳道人?.
傳道人為了要跟弟兄姊妹拗帳.
為了得著一些權力.
為了得著一些好處?.
不是的.
我們為著的是事奉神.
我們為著的是牧養弟兄姊妹.
究竟今天你那個牧養.
你所做的方法是阻礙你的牧養.
還是幫助你的牧養.
回到你的初心裡面.
重新再思想.
在這裡上帝幫助以利亞.
重新回到你的初心裡面.
以利亞你在做什麼?.
你在做什麼?.
所以這個上帝有趣的地方是什麼呢?.
有一段經文.

$^{241}$也是其中一段難解的經文.
當耶和華出來叫他站口.
去到洞口的時候.
你出來站在山上的面前.
然後耶和華說.
那是耶和華從哪裡經過?.
在他面前有烈風大作.
崩山碎石.
耶和華卻不在風中.
風後有地震.
耶和華卻不在其中.
地震後有火.
耶和華也不在火中.
火後有微小的聖淫.
什麼叫上帝?.
這樣就是上帝.
什麼意思呢?.
剛才他去到的地方叫荷列山.
荷列山就是摩西.
他去到那裡的時候.
他上山.
山上就是有這些景象.
烈風大作.
崩山碎石.
大風.
地震.
但當時這些現象.
就是表示什麼呢?.
所有的以色列人.
尤其是他們是先知.
他們一定知道這些現象.
代表上帝在當中.
所以他們看到山上有煙霧.
圍著的時候.
他們就不敢上山.
這是他們的傳統.
這是他們的神學.
但現在告訴你.
你們以前以為上帝在那裡嗎?.
不在那裡.

$^{281}$你搞清楚.
上帝可以選擇在微星裡面.
微星是什麼?.
沒有人注意.
沒有人留意.
上帝的出現可以不是那麼顯赫.
頂智媒.
不是過去你所認識的上帝.
他就是這樣.
不是.
上帝有他的方法.
所以在我們的人生裡面.
你才會發覺很多時候.
這些片段才會出現.
為什麼在我這麼艱難的時候.
上帝為什麼好像不見了.
不是.
上帝在這裡.
他用著另外的方法.
和你相處.
而你是需要的.
在那一刻.
上帝要用微星和你說話.
你才聽得清楚.
但最有趣的地方.
就是耶和華.
微星一出現的時候.
是什麼呢?.
就是下面.
耶和華對他說.
你回去吧.
你回到原路.
我想大家留意.
你要高哈舍作亞蘭王.
亞蘭是外邦的王.
你要高靈士的孫子.
耶和華作以色列的王.
外邦.
以色列.
還有最特別的地方.

$^{321}$你要高邁就是以利沙接續你.
這個就是上帝給以利亞的回應.
你以為.
你很偉大嗎?.
每一個興起.
每一個廢掉.
都是上帝決定的.
弟兄姊妹.
為什麼我們來到神面前.
我們唯一可以選擇的.
就是謙卑.
就是因為這樣.
我們以為今天自己很厲害嗎?.
最近的時間.
你可以這樣說.
如果.
有些弟兄姊妹.
最近就說你多才多藝.
一時回來就咖啡師.
一時回來就港道.
不是.
我媽的形容是什麼呢?.
你全身都沒有把利.
我每樣都有興趣.
我每樣都有興趣學.
但沒有一件事是真的.
所以我經常都跟弟兄姊妹說.
不要鋼琴.
不要緊,我支持住.
你不要叫我彈旋律.
如果你猜我可以.
我只可以給你我的節奏.
你明白我的意思嗎?.
我給不到很厲害的.
所以我搞不定.
但我想說的情況就是.
不是我們厲害.
是上帝.
給你們這些才能.
的時候.

$^{361}$讓我能夠在那一刻.
去做這個服事.
是上帝.
上帝可以不用你.
上帝隨時隨地可以換人.
我們沒有什麼厲害.
所以當我們走到絕境的時候.
弟兄姊妹.
你不用怕.
你不用怕原因不是你厲害.
還是你不厲害.
我這裡所說的就是.
我都沒有以前的人那麼好.
不是,不關以前的祖先事.
不關其他人事.
是上帝.
他用不用你.
曾經有一個童工.
他跟我說.
只需要一句話.
就可以毀了一個人的一生.
有一個童工.
他以前在教會很好的服事.
他信耶穌的時候.
跟著一班很好的肢體.
跟著一路成長.
他一直以那些為大姐姐.
大哥哥,他很信任他們.
在一路事奉,一路做的時候.
突然間不知道什麼原因.
我不知道什麼原因,他沒有說得很具體.
但就是這班大哥哥姐姐裡面.
其中有一個他很信任的人.
跟他說了一句.
上帝是不會用你的.
我不知道什麼事.
但他就帶著傷痕.
一路來到今天.
他讀完神學,現在就做傳道人.
但他心裡面.

$^{401}$永遠都有這條痣.
他一直在想.
每次都問上帝是不是真的會用我.
弟兄姊妹,我想告訴大家.
上帝.
是會用每一個人的.
除非.
你離棄上帝.
除非.
你.
不跟隨上帝.
所以在神的面前的時候,當我們走到某些位置.
我們很多時候.
人生裡面有很多的計算.
就像這裡所說,上帝對於.
以賴尼亞來說,祂只可以去做.
那個令我.
可以那麼有面子.
令我可以那麼義氣風發的.
上帝.
其他事情祂不會做的.
但弟兄姊妹.
我想告訴大家,不是,我們的上帝.
不是這樣.
所以當有些弟兄姊妹.
現在已經是七月了.
我們坦白說,我們現在也要.
面對我們要面對續租的問題.
一些最後的限期.
你可以這麼說.
我想很多時候.
我們都仍然關注於.
譬如無論怎樣.
最簡單的事情,他加多少租.
他加多少租.
我們能不能應付.
是不是這些問題.
沒錯,我們是需要去想的.
不過弟兄姊妹.
請大家一起去學習.

$^{441}$你首要想的不是想這件事.
我想弟兄姊妹.
你回到在監牢裡面.
你今天作為能夠成為監牢的一份子.
上帝在這裡.
祂有沒有給我們意象.
祂有沒有在這裡.
有這個工作.
祂要我們做.
至於是否能夠.
應付弟兄姊妹.
我想告訴你,是沒有人知道的.
我記得以前有一個蜀靈師傅.
現在是牧師.
他太太是一個非常祈禱.
很好的人.
對上帝很有信心.
甚至他的信心比我這個師傅多.
他不斷為.
每個月不斷奉獻.
去猜拳.
有一段日子.
小朋友在長大,當時的工資又不會高.
很掙扎.
究竟每個月裡面.
如果我自己不斷計算.
應該是不可能奉獻給那個機構.
做不到的.
但他的見證.
他太太只提了一句.
放下你的計算機.
不要再計算了.
你越計算就越沒有.
你越相信上帝,上帝就會準備.
他就聽他太太說,放下計算機.
不要緊,憑信心去吧.
每個月照樣奉獻.
但弟兄姊妹,我想告訴大家.
他的見證是告訴我們.
在那段日子裡面.

$^{481}$他所得到的供應.
是比之前更多.
所以弟兄姊妹.
先求他的國.
和他的義.
他就會將你所需要的一切.
都加給你們.
是不是?.
這不是上帝自己的應許.
所以我鼓勵弟兄姊妹.
一起學這個功課.
首先要去問.
究竟上帝你所給我們.
在這裡還有什麼工作.
我們很想為你去做.
我們還沒完成.
我們還要繼續去走.
求神讓我們.
當我們好像走到絕境的時候.
不要用我們過往.
自己的神學.
去猜度這個上帝.
我用最近的.
其實說真的,我說.
咖啡侍工,上次Yummy.
他們在的時候.
他們都記得.
我和他們當時.
我出去做那間餐廳.
我回來的時候.
最初我不是想做餐廳,我想開咖啡店.
因為咖啡店容易接觸到人.
是不是?.
現在咖啡店開了之後.
我和弟兄姊妹多了一些話題.
多了一些.
容易聊得來的地方.
其實是有.
他的目的.
但我想告訴你,我從來沒有想過.

$^{521}$會有一切這樣的東西在這裡.
但我想告訴你,上帝就是這麼幽默.
明明在這個時刻.
教會都不知道前景是怎樣.
你都不知道會有什麼情況.
竟然做咖啡侍工.
你們不是吧?.
傻的嗎?.
弟兄姊妹,我想告訴你.
在這段日子裡.
我也和Eddie說.
其實這段時間大家真的要.
很感恩Eddie.
一直很投入幫忙做這些事情.
我想告訴大家.
就是有一個傻子.
在過程裡.
我們大家一起.
我們都不知道為什麼.
我們自己都會大發熱心.
我們那次在聊天.
我們都在說這件事.
為什麼會突然間.
去得這麼快?.
原因就是爭取時間.
就是不知道11月行不行.
我們就在這段時間爭取.
要做我們的侍工.
但竟然我們用了一個這麼低的成本.
去做這個事.
竟然我們用了一個這麼低的成本.
運作到今天.
大家都會覺得.
嘩!.
行啊!.
弟兄姊妹,這是神職.
是上帝的工作.
所以弟兄姊妹.
我求上帝只是在當中.
聖靈和大家一起同工.

$^{561}$讓大家在這裡.
找到一個你自己很清清楚楚的意象.
上帝究竟.
我要怎樣為你而活.
我要怎樣為你的工作.
在這個時間我還有什麼人.
我是特別關心.
所以我剛才也說了.
我的呼籲很簡單.
幫忙咖啡侍工的大事.
三樣要點,第一樣東西,你先喜歡咖啡.
你不喜歡咖啡,你不要做.
不要為了沖咖啡而沖咖啡.
第二樣東西.
你真的想轉服任.
第三.
願意奉獻你的時間.
是要毀身的.
所以求神幫助我們.
在上帝的當中.
拿著上帝所給我們的意象.
我們一起向前行.
上帝在這裡.
所應許的是什麼呢.
但我在以色列人中.
為自己留下了七千人.
是未曾向巴力屈膝的.
未曾向巴力和巴力親嘴.
和巴力親嘴.
不要以為我們自己有多厲害.
上帝才是最厲害.
好不好.
所以我求上帝在你們的生活中.
都是一樣.
當你面對困難,你面對轉變.
請你記得.
上帝從來對你的愛.
都沒有變過.
求神一起.
幫助我們同心祈禱.

$^{601}$在我們人生總會有低谷的時候.
有時候我們都不明白.
就好像以利亞一樣.
明明我昨天還在山上.
去禁義氣風法.
為什麼轉頭的時候.
上帝就好像消失了一樣.
但上帝告訴我們.
從來都不是這樣.
當以利亞去了解這幅圖畫之後.
一秒事情就發生了.
他明知要換掉它.
但他就按照你的旨意.
去到每個地方.
特別是去高納伊利沙.
去接觸它.
最後這個就是生命的轉變.
拿得起,放得低.
最後我求主.
你親自保守幫助.
很願意我們每一個人生命中.
都有這份由願.
給上帝塑造我們.
我求你時時刻刻提醒我們.
歡迎我們在神的面前.
我們只走你的路.
主我求主你親自保守.
幫助.
祈禱奉教主耶穌得勝,明知其求.
阿們.
\newpage



\section{路加福音 14:1-24}
\label{sec:Vwm387KuCtc}
\textbf{2023-07-02 早堂《筵席的再思》路加福音 14\_1-24 講員:曾裔貴牧師}
\newline
\newline
連結: \href{https://youtube.com/watch?v=Vwm387KuCtc}{\texttt{https://youtube.com/watch?v=Vwm387KuCtc}} ~~~~ 語音日期: 2023-07-05
\newline
\newline
\hyperref[sec:W6aFxBk_o0s]{\small{< < < PREV SERMON < < <}}
~
\hyperref[sec:index_chronic]{\small{[返順時目]}}
~
\hyperref[sec:index_scriptual]{\small{[返順卷目]}}
~
\hyperref[sec:LQfxIAo88QM]{\small{> > > NEXT SERMON > > >}}
\newline
\newline
路加福音 14:1-24
\newline
\begin{longtable}{cl}
\hline
\hline
章節 & 經文 (和合本修訂版)\\
\hline
14:1 & \begin{tabularx}{0.7\textwidth}{X} 安息日，耶穌到一個法利賽人的領袖家裡去吃飯，他們就窺探他。 \end{tabularx} \\ \\ \relax
14:2 & \begin{tabularx}{0.7\textwidth}{X} 這時在他面前有一個患水腫病的人。 \end{tabularx} \\ \\ \relax
14:3 & \begin{tabularx}{0.7\textwidth}{X} 耶穌回答律法師和法利賽人，說：「安息日治病合不合法？」 \end{tabularx} \\ \\ \relax
14:4 & \begin{tabularx}{0.7\textwidth}{X} 他們卻不說話。耶穌扶著那人，治好了他，叫他走了。 \end{tabularx} \\ \\ \relax
14:5 & \begin{tabularx}{0.7\textwidth}{X} 耶穌對他們說：「你們中間誰有兒子或有牛在安息日掉在井裡，不立刻拉他上來呢？」 \end{tabularx} \\ \\ \relax
14:6 & \begin{tabularx}{0.7\textwidth}{X} 他們對這些事不能反駁。 \end{tabularx} \\ \\ \relax
14:7 & \begin{tabularx}{0.7\textwidth}{X} 耶穌見所請的客人選擇首位，就用比喻對他們說： \end{tabularx} \\ \\ \relax
14:8 & \begin{tabularx}{0.7\textwidth}{X} 「你被人請去赴婚宴，不要坐在首位上，恐怕主人請了比你尊貴的客人， \end{tabularx} \\ \\ \relax
14:9 & \begin{tabularx}{0.7\textwidth}{X} 請了你和他的那人前來，對你說：『請讓座給這一位吧。』你就羞羞慚慚地退到末位去了。 \end{tabularx} \\ \\ \relax
14:10 & \begin{tabularx}{0.7\textwidth}{X} 你被請的時候，去坐在末位上，好讓主人來對你說：『朋友，請上座。』那時，你在同席的人面前就有光彩了。 \end{tabularx} \\ \\ \relax
14:11 & \begin{tabularx}{0.7\textwidth}{X} 因為凡自高的，必降為卑；自甘卑微的，必升為高。」 \end{tabularx} \\ \\ \relax
14:12 & \begin{tabularx}{0.7\textwidth}{X} 耶穌又對請他的人說：「你準備午飯或晚餐，不要請你的朋友、弟兄、親屬和富足的鄰舍，免得他們回請你，你就得了報答。 \end{tabularx} \\ \\ \relax
14:13 & \begin{tabularx}{0.7\textwidth}{X} 你擺設宴席，倒要請那貧窮的、殘疾的、瘸腿的、失明的， \end{tabularx} \\ \\ \relax
14:14 & \begin{tabularx}{0.7\textwidth}{X} 你就有福了！因為他們沒有甚麼可報答你。到義人復活的時候，你要得到報答。」 \end{tabularx} \\ \\ \relax
14:15 & \begin{tabularx}{0.7\textwidth}{X} 同席的有一人聽見這些話，就對耶穌說：「在神國裡吃飯的有福了！」 \end{tabularx} \\ \\ \relax
14:16 & \begin{tabularx}{0.7\textwidth}{X} 耶穌對他說：「有人擺設大宴席，請了許多客人。 \end{tabularx} \\ \\ \relax
14:17 & \begin{tabularx}{0.7\textwidth}{X} 到了坐席的時候，他打發僕人去對所請的人說：『請來吧！樣樣都已齊備了。』 \end{tabularx} \\ \\ \relax
14:18 & \begin{tabularx}{0.7\textwidth}{X} 眾人異口同聲地推辭。頭一個對他說：『我買了一塊地，必須去看看。請你准我辭了。』 \end{tabularx} \\ \\ \relax
14:19 & \begin{tabularx}{0.7\textwidth}{X} 另一個說：『我買了五對牛，要去試一試。請你准我辭了。』 \end{tabularx} \\ \\ \relax
14:20 & \begin{tabularx}{0.7\textwidth}{X} 又有一個說：『我才娶了妻子，所以不能去。』 \end{tabularx} \\ \\ \relax
14:21 & \begin{tabularx}{0.7\textwidth}{X} 那僕人回來，把這些事都告訴了主人。這家的主人就發怒，對僕人說：『快出去，到城裡大街小巷，領那貧窮的、殘疾的、失明的、瘸腿的來。』 \end{tabularx} \\ \\ \relax
14:22 & \begin{tabularx}{0.7\textwidth}{X} 僕人說：『主啊，你所吩咐的已經辦了，還有空位。』 \end{tabularx} \\ \\ \relax
14:23 & \begin{tabularx}{0.7\textwidth}{X} 主人對僕人說：『你出去，到大街小巷強拉人進來，坐滿我的屋子。 \end{tabularx} \\ \\ \relax
14:24 & \begin{tabularx}{0.7\textwidth}{X} 我告訴你們，先前所請的人沒有一個可以嘗到我的宴席。』」 \end{tabularx} \\ \\
[1ex]
\hline
\hline
\end{longtable}
$^{1}$我們一起同心祈禱.
多謝主.
我們要打開聖經.
幫助我們能夠明白這十四節聖經的真理.
我們去思想關於賢侄主耶穌.
你對當時參加父會的每一個賓客的教訓.
以及成為聖經之後對後世的每一個信徒的重要教導.
我們一起在你面前等候.
禱告奉耶穌基督的名.
阿們.
你感受到第一節.
作者已經告訴我們.
賢侄充滿了有火藥味.
充滿了冒犯性.
因為第一節已經說.
有一群人他們準備去窺探耶穌.
窺探祂什麼呢.
一定是想找祂有什麼錯處來把持.
然後來控訴.
但是在這個時候.
其實耶穌基督來參加一個被邀請的賢侄.
大家可能要一點背景.
你就會明白.
主耶穌自己也說的.
稍後那裡.
祂說你如果去邀請一些賓客的時候.
在當時的習俗.
一定是請和自己最適合的階層.
不會低於你的階層來去宴請.
又會彼此輪流.
你請我 我請你.
如果能夠找到一些高於你的階層的人.
就更加有光彩.
這個在希羅的文化.
在當時耶穌的時代.
是相當普遍的.
所以稍後我們感受到那種冒犯性.
你就會看到是很特別的.
不過這裡的冒犯性.
是耶穌基督被邀請.

$^{41}$是一個法利賽人.
來邀請耶穌.
在安息日.
到底是星期五的黃昏.
去到星期六的時間.
還是過了安息日的時間.
如果是安息日的時間.
為什麼要走遠路去飲食呢.
這些技術性的問題.
作者沒有交代給我們知道.
而是告訴我們.
在安息日是有一個聚會.
一個宴會.
而且這個宴會很多都是同階層的人.
即是說這些他們.
就是主耶穌在第三節所說的.
就是耶穌和這些在座的賓客.
窺探耶穌的賓客.
他們是律法師 法利賽人.
為什麼要這樣和他們對話呢.
原因就是他們在窺探的時候.
主耶穌刻意.
來將一個表面傳統.
是不可以做的.
就是安息日不可以治病.
耶穌就刻意在這個宴會的時間.
將一個患水腫的人.
來去醫治.
這個醫治.
作者也沒有交代是什麼方式.
相信是一句話.
他就好了.
耶穌就打發他離開.
醫好之後.
主耶穌就和這些律法師和法利賽人說.
說什麼呢.
在安息日可不可以治病.
第四和第六節都提到他們是無言以對.
不代表心靈裡面是心悅誠服對耶穌.
他們是無言.

$^{81}$表面不出聲.
不回答.
本來應該回答的問題.
可不可以.
如果答可以.
就是犯安息日.
耶穌就明顯犯了安息日.
其實在他們的內心.
已經有一個判斷.
耶穌是不對的 是錯的.
所以耶穌在這裡這樣做的時候.
耶穌的解釋.
好像再一次將人以為鐵律的傳統.
重新帶上真理的部分.
來重新解釋什麼叫做安息.
第五節.
你們中間誰有牛奴.
或者是有兒女有病.
在安息日掉進井底.
你不會立時將他拉上來.
那意味著他們的律法師和法理人.
所謂無言以對.
其實是有雙重的標準.
自己安息日.
面對一些很大的危機.
你自己就做.
別人安息日就什麼都不准做.
這樣的雙重標準.
原來在耶穌的眼裡.
自小他在會堂成長.
他現在出來傳道.
他現在撥亂反正.
將一些人以為是鐵律的傳統.
一下子重新將神的標準再賦予.
在主耶穌的眼裡.
真正的安息.
不單純是死守的字面.
傳統本身不是錯的.
傳統一定有它的來源.
有大家都來行之有效的約定俗成.

$^{121}$但是傳統不是等同於真理.
真理才是真正的真理.
如果我們說傳統是一切.
可能我們有時會不知不覺間.
會拒人於千里.
甚至有時會誤解了神的心意.
我昨晚都講了一個例子.
我在元朕有兩年的實習.
應該是2007,2008年的時間.
都是一樣.
跟隨主音牧師.
要在外面.
你都知道大唐道的路口.
我們就要當值.
弟兄姊妹回來.
我們就跟她打招呼握手.
我是一個神學生.
我忘記了是晚堂還是早堂.
每次都近乎每個禮拜.
有個弟兄都是上一年級的.
每次就離遠近油站.
吃完煙仔才回教會.
吃完很快就衝進去.
又不跟你握手.
又不打招呼.
我見了很多次.
心裡面都有些想法.
回教會還吃煙.
慢慢聊著聊著.
牧者很好.
他說 沒錯.
教會是應該告訴人.
有些事是不對的就不應該做.
但不等於每個人都做到.
他慢慢慢慢.
他雖然自己煙癮都未戒到.
但他都願意吃完煙仔回來教會.
即是說他對神有心.
可能慢慢就改變到他.
所以如果這個教會定了.

$^{161}$傳統不准抽煙.
可能令到很多人有卻步.
這個可能是其中一個小小的例子.
但在這裡.
耶穌在這裡很嚴肅的.
來去confront這些快利賽人.
鐵律法師.
不要將傳統成為一個鐵律的枷鎖.
而且是加諸在其他人的生命裡面.
令到人以為.
宗教信仰有很多很多的條例.
我要一路一路.
好像過關斬將那樣去滿足了.
我才能夠符合到上主的心意.
原來不是.
原來在主耶穌的眼裡面.
安息人不是一種死守的戒律.
有需要的時候.
就可以去改變那個所謂的傳統.
所以在這裡你可以感受到那個氣氛.
這個筵席.
那種的冒犯.
本來耶穌感覺冒犯.
因為很多人窺探著他.
但反過來.
耶穌好像反客為主.
在安息日裡面.
來將一些可能的傳統.
活亂反正.
這次被冒犯的是.
這些法利賽人和律法師.
所以你覺得很奇妙地.
我們的主把握他每一次的機會.
原來吃東西不是重點.
宴會可以在一班人的面前.
將神的真理去教導.
才是主耶穌的一個重點的服事.
接著還有兩類人.
耶穌來勸勉和教導.
首先就是第七節.

$^{201}$耶穌見到一個情景.
就是有些賓客.
揀擇守衛.
這是第七節.
耶穌見所請的客人.
選擇守衛.
接著用了一個比喻.
來跟他們說.
重點是回到第十一節.
在這個段落.
對賓客來說.
明明是被請去付席.
可能當時不是安排了你坐什麼位.
可能有些位是特別顯赫.
而導致出現這樣的效應.
耶穌見到這樣的情況下.
表面借題發揮.
其實不是.
其實是勸勉當時付會的人.
他們要來學習生命的功課.
由一個坐席.
爭取守衛.
轉化成為你的生命.
到底在神的眼中.
你是一個怎樣的人呢?.
這個重要.
因為十一節.
這是很難學習的功課.
基督教差不多是一個很主流的教導.
每一個基督徒.
要在上主的眼中.
你是一個甘心樂意.
自甘卑微的人.
的而且確.
在做物主的眼中.
我們是一無所有.
我們是徹徹底底的罪人.
假如他要收我們的生命.
我們連一句都不能說.
因為凡自高的必降為卑.

$^{241}$自甘卑微的必升為高.
就是這個比喻的重點結論.
就是看到有些賓客.
爭奪那個手座.
耶穌就告訴他們.
當然這個就是當時的氣氛.
當時他們的風俗等等.
耶穌不是看不過眼.
耶穌是借這樣的場景.
用一個比喻來講述人生.
如果宴會是一個人生的盛宴.
我們是用什麼態度.
面對我們的人生.
以致我們的人生在神的眼中.
是一個能夠.
他稱許的生命.
這個是相當要緊的事.
有哪些人真的可以.
在生命裡面.
在上主的眼中.
是一個謙虛而能夠不斷獲得的人.
要自己降卑.
是極之不容易的事.
我們不會有這些情況出現.
去到宴會爭奪的位置.
但是我們的內心.
如果我們是一個驕傲.
我們的內心不肯來到.
在上主的眼裡面.
來學習真正自己的本質.
本相的時候.
這是一個很難醫治的心靈.
有一位在美國的弟兄.
他小時候就問爸爸.
基督教,耶穌是什麼.
爸爸跟他說.
宗教,基督教.
你最好盡快遠離.
不要想,不要問.
他如是者長大.

$^{281}$很不容易完成了他的高中的學習.
讀書不厲害.
輾轉進了美國的陸軍.
最低階的士兵.
很辛苦,很艱苦,很難.
很多的掙扎,很多的挑戰.
後來輾轉就進了空降師.
82空降師.
很出名的美國.
很重要的空降師部隊.
信仰對他們來說.
完全是空白的.
不知如何得到一本聖經.
基甸會贈送的聖經.
這本聖經對他們來說很重要.
因為他們是空降師.
很多危險.
他們很多時候要面對.
這本聖經就成為他們可以放在背囊裡.
每次跳降傘.
他們就感覺到一份平安.
那本聖經原來是用來拍照的.
後來輾轉.
他居然自己拿本聖經來看.
一邊看一邊看.
一年時間.
他算是看完這本聖經.
因為他讀書不多.
很艱難去讀這本聖經.
一年後.
他去找軍隊的牧師.
牧師呀 我信了耶穌.
牧師說跟他聊聊天.
聽他說原來還未信.
就跟他解釋基本的真理.
這個弟兄就學習.
他越來越對聖經有興趣.
就回到一間教會.
一間教會是很認真查經的教會.
他很有興趣地學習.

$^{321}$如是者.
剛才還沒說一件事.
就是他在軍隊的時候.
因為很艱難.
所以很多軍隊裡面的人.
會沉溺在那些違禁藥裡面.
希望從其中找到刺激.
來逃避生活的艱苦.
輾轉他信了主之後.
他靠神的恩典.
這些禁藥他都離開了.
所以沒多久之後.
他的那些同袍.
很多被發現服食違禁藥.
而受到懲處.
他因為自己已經戒掉了.
就沒有了這個問題.
後來他離開軍隊了.
他有機會進大學.
為什麼我會知道這件事呢.
因為有一位美國舉世知名的學者.
D.A. Carson.
這個年輕人.
就是他的學生.
在Trinity神學院.
這個學生.
居然讀完大學之後.
感覺自己是蒙召.
就進神學院做傳道.
讀神學.
Carson說他的希臘文是全班最優秀.
一個本來讀不到什麼書的.
對信仰完全是空白.
家人都幫不到他的.
聖經只是為了辟邪.
因為太危險了.
慢慢慢慢.
神的道開到他的心靈.
讓他能夠謙卑的學習真理.
而越來越有一個提升.

$^{361}$好像一個愚昧愚頑的人.
比他的師父更加通達.
很難想像神的愛.
神的道.
將一個人徹底的改變.
弟兄姊妹.
在座的信主.
有沒有這種生命的經歷.
際遇呢.
越是沒有環境的.
你不是信二代.
你是自己一個人信耶穌.
這個情況是會出現的.
我自己都說過很多次.
以前我吃兩包煙.
有很多很多.
我試過最高紀錄是15天離家出走.
不回家.
當然是明仔不知道.
還沒出生.
15天不在家.
去了朋友家.
什麼時候會回來呢.
沒錢了就回來.
如是者令媽媽流著眼淚.
叫我不要走.
這樣的浪子.
突然之間.
神的愛和神的恩典的感動.
自己來到去.
決心的信主悔改.
弟兄姊妹.
主耶穌表面好像confront那些人.
不要爭啊.
其實不是.
我們的生命.
能不能夠在神的眼中.
算為真正的付足.
就是我們的生命.
其實是不是一個謙虛學習真理的生命.

$^{401}$驕傲的人.
就是裝滿了自己以為對的東西.
就不會再聽其他.
不同的角度的客觀的評語.
謙卑的人.
就好像一個海綿一樣.
不斷願意吸收不同的真理.
和不同人對你的善意的勸勉.
令到自己能夠不斷地有成長.
這個是生命功課.
到第三個場景.
同樣都是那個宴會.
突然之間主耶穌就和那些請人的.
法利賽人.
那個筵席的主人家.
這裡就告訴我們第九節.
對不起應該是十二節.
耶穌又對請他的人說.
你擺設午飯或者晚飯.
不要請你的朋友弟兄親屬.
剛才說了.
付足的輪謝.
恐怕因為大家交換.
他們也請你.
你就得了報答.
你擺設筵席.
都要請貧窮的.
殘廢的.
狂配的.
黑眼的.
你就有福了.
好像真的很冒犯.
在座的每個人都感覺自己的身份.
是匹配和法利賽人.
被邀請的一定是同階層的.
甚至高於主人家的.
包括耶穌.
都是被禮遇的來被邀請的.
但耶穌這樣跟法利賽人說.
叫他請人.

$^{441}$一定不是這次.
這次來不及了.
下次要這樣請.
對當時的社會習俗是很冒犯的.
怎麼會請一些像殘廢宴一樣.
全部是社會上.
好像是邊緣的人.
怎麼會這樣做.
第十四節將這個問題點出來.
因為他們沒有什麼可以報答你.
突然之間宴會就還宴會.
為什麼突然之間說到二人復活.
你就會得到報答.
那就是說.
不只是宴會的飲食文化的關係.
而是主耶穌透過這樣的一個宴會.
不是教導一個餐桌禮儀.
而是透過宴會去講生命的功課.
怎樣可以得到在神上主的眼中.
去到祂面前.
你是付足呢.
曾經路加醫生在十二章.
講過一個叫做無知財主.
一生的為自己去積蓄一切.
突然之間有一天.
神跟他說.
我今夜收回你的靈魂.
你所得的歸誰呢.
接著耶穌講了一個結論.
第二十一節.
十二章的二十一節.
如果你會背就很重要.
凡為自己積財.
在神眼中卻不付足.
也是好像這個無知財主的如此.
如此的下場.
物質一生都豐富.
但是心靈在做物主的眼中一無所有.
這裡不是讓我們不要仇恨富貴.
不是這個意思.

$^{481}$不是富貴就錯.
不是這個意思.
是你擁有這麼多的資源.
你有沒有為自己的生命.
在神眼中.
你的時間.
你的金錢.
你的財幹.
你的生命.
去成為別人的祝福呢.
而最終你在神的眼中.
到二人復活的時候.
你就會得到報答.
所以是一個很大的回應.
勞家醫生不是叫我們仇恨富貴.
因為富足表面是忠誠的.
是我們的內心如何看待金錢.
看待我們自己.
才是最重要.
所以很奇妙地.
一個好像很冒犯的宴會.
耶穌基督藉著這樣來教訓這幾類賓客.
和主人家.
其實帶出的是生命的功課.
親愛的弟兄姊妹.
我們有沒有真的在我們有限的人生.
盡力的為主.
來去積財在天上呢.
有沒有積蓄我們的內心.
成為別人的祝福.
而將來.
聖經說到二人復活的時候.
我們就會得到報答.
不單純是神來報答我們的人生.
更加我們去到神的面前.
將會有很多很多的二人.
會來對我們一生所做的.
你想想.
雖然我們都在這個物質的世界.
想不到靈界將來會是怎樣的.

$^{521}$怎樣排隊來到你的面前.
多謝你所做的一切善行.
我們很難用現在物質的概念.
去想將來的永恆.
但是耶穌已經說得很清楚.
原來我們的一生.
其實每分每秒.
都是為著我們的永恆.
來去準備的.
這裡告訴我們.
主耶穌不是要我們去.
好像施捨給這些貧窮.
殘廢 瘸腿 乞訝.
是很想告訴我們.
我們有的是上主給我們的生命付足.
我們就要和別人去分享.
今天不夠時間說的是.
15到24節的大言直.
因為待會五堂.
就會有充足的時間.
就會說那個大言直.
很奇妙.
原來有三份的名單.
第一批的人一個都不來.
第二批的Plan B的名單.
全部就是.
這個是第21節.
14章的21節.
就和那個僕人說.
好 現在你快點出去.
去到城裡大街小巷.
去將那些人.
貧窮的 殘廢的 乞訝的 瘸腿的.
都請來這個言直.
很特別的.
耶穌基督教主人家.
可能下次你要這樣做.
其實應該說每一個人.
生命都要這樣學習.
但是跟著耶穌說的這個大言直.

$^{561}$很重要就是第15節.
或者是點出.
不可以再講解了.
就是同直的人.
這次突然間出現一個同直的人.
聽到耶穌這樣說.
很羨慕.
就和耶穌說.
在神國裡面吃飯就有福了.
在神國裡面.
就是說他們已經是有了那個.
好像是他們有份的.
一定有份的.
耶穌說你真的有份.
不一定.
跟著耶穌說一個.
更加特冒犯的比喻.
就說到被邀請的第一批的人.
全部不來.
還有是堂而皇之的理由.
這個如果你有時間.
等一下午堂可以聽一聽.
今天我們總結了這三個部分.
表面就有三類的人.
法利賽人這個主人家.
有窺探耶穌的律法師.
法利賽人.
有主耶穌和其他所有在場的賓客.
包括成為今天我們來讀聖經的.
每一個弟兄姊妹.
來探索.
筵席.
生命的功課.
我們在神的真理底下.
去再透視什麼是傳統.
那個價值和它的限制.
我們再來去反思.
我們的生命在主的眼中.
到底是驕傲.
還是謙卑的人.

$^{601}$不斷地學習.
更加重要.
在上主眼中.
是不是真正的富足.
一起禱告.
多謝主.
我們很盼望你打開我們大家的心靈.
一同學習你的真理.
讓我們能夠在你的面前.
我們不斷透過認識真理.
學習真理.
成為真正的富足.
奉主耶穌基督的名.
阿們.
\newpage



\section{}
\label{sec:LQfxIAo88QM}
\textbf{2023-07-06 周四晚查經班《被擄歸回歷史研讀(19)》 講員:曾裔貴牧師}
\newline
\newline
連結: \href{https://youtube.com/watch?v=LQfxIAo88QM}{\texttt{https://youtube.com/watch?v=LQfxIAo88QM}} ~~~~ 語音日期: 2023-07-08
\newline
\newline
\hyperref[sec:Vwm387KuCtc]{\small{< < < PREV SERMON < < <}}
~
\hyperref[sec:index_chronic]{\small{[返順時目]}}
~
\hyperref[sec:index_scriptual]{\small{[返順卷目]}}
~
\hyperref[sec:oxDrdn_rwTg]{\small{> > > NEXT SERMON > > >}}
\newline
\newline
$^{1}$我們一起同心祈禱.
多謝天父讓我們今天晚上再一次聚集.
我們去看最後這幾章的尼熙米記.
我們在上一個禮拜都提到.
真是甘心樂意服侍神的人.
願意住在已經建造好城牆的耶路撒冷.
準備來守護這個城和神的聖殿.
也有一些甘心樂意被抽籤的猶太人.
能夠聚居在耶路撒冷裡面.
也在周邊的地方有其他猶太人.
我們想到這一切都是作者尼熙米悉心的安排.
也是聖靈的工作感動不同的百姓.
來到齊心的努力的參與.
我們看到整件事情是經過復興之後.
帶來的願意配合一些行政的指令.
我們真的很盼望透過這些舊約的經文.
讓我們看到一個恢復敬拜的地方.
是充滿了神恩典同在聖靈工作的地方.
讓我們今天渴望的不是只是人數坐滿.
這個已經是很難了.
在今天這個世代因為太多選擇.
有時宗教信主返教會.
可能是一個我們有時間才會考慮的一個節目.
我們很盼望我們很多人能夠.
很快就將一個生命的次序去搞清楚.
特別在我們今天晚上.
我們看到第十二章的經文的時候.
讓我們真的有一個感動.
我們很願意能夠被神分別出來.
成為了別人的祝福.
今天晚上我們等候.
我們要打開《離險米記》去讀的時候.
都讓我們能夠明白掌握.
有很多舊約的背景.
都讓我們不要怕悶.
而是能夠梳理背後作者要寫作的原因.
使我們能夠得到那個啟發.
我們一起在你面前等候禱告.
奉耶穌基督的名,阿們.
各位我們在第七章曾經說過.

$^{41}$城牆建好之後.
其實《離險米記》立即就發現了一件事.
七章的第一到第六節.
就是告訴我們人口稀少.
耶路撒冷很荒涼,沒什麼人.
防禦聖城需要有人去守護.
這個是自然不在話下的.
但是我們看到作者他不急於用他的行政命令.
他是省長來的.
來去強迫,好像捉壯丁那樣.
將那些人遷徙去耶路撒冷.
有那個實際的軍事需要.
打完仗,建好城牆.
那些人回到自己的地方.
很自然就回到周邊的地方.
耶路撒冷較為疏落,很合理.
但是我們看到的就是.
有一個自發性的嚴肅會,復興會.
一個大家一起研讀聖經的一個奮興會.
第八第九章.
接著就是有很多人願意歸入神的名下.
來遵守摩西的律法書.
出於一份真誠,出於一份感動.
出於一份神的愛.
在很多百姓的心靈裡面.
然後我們就看到第十一章.
正式講在這裡的時候.
哪些人本身住在耶路撒冷.
然後尼熙米就和眾民有一個協議.
就是第二節第三節.
這個就是我們上個禮拜主力講的經文.
就是尼熙米做了一件事.
就是逢十抽一抽籤.
一般以色列人第十章就告訴我們.
他們答應我們在我們的收入.
我們的耕種,我們的田產.
我們的十分之一會奉獻去到聖殿.
來讓聖殿的敬拜祭祀.
是可以有那個物資.
有那些的良享給祭司利未人.

$^{81}$我們承諾這樣做.
現在第十一章就不是那個物資.
不是那些的財物.
是那些人口.
逢十抽一.
就離開自己的猶大的地方.
就住在耶路撒冷.
就是上去耶路撒冷.
就負責可能就是要由省長.
由祭司安排.
來到去做什麼工作.
住在哪個位置.
諸如此類.
都聚集了一定的人數.
就在耶路撒冷.
原來就為第十二章.
有一個奉獻城牆之禮.
我們等一下看看.
四十三節那段經文.
就是第十二章的一個.
就是說這個奮興會之後的一個結束.
由第七章.
那個城牆建好.
第八第九章.
就是駐彭捷.
和七月初一的摧國節.
的奮興會.
然後一路去到差不多七月底了.
因為駐彭捷守七天.
到第八天他們又繼續自發的.
守嚴肅會.
接著就第十章他們就簽名.
願意來跟隨摩西的律法.
來生活.
去到第十一章.
全章就是說.
一個好像人口普查的分佈圖.
哪些人住在耶路撒冷呢.
當然就是祭司.
他們一定是強制要住在耶路撒冷的.

$^{121}$他們是差不多是生生不息.
由他們長大進去聖殿事奉.
他們就要來到去.
在聖殿那裡工作的了.
這個是祭司的天職來的.
神安排了的了.
那利未人呢.
都是一個輔助的角色來的.
所以在這幾章.
第十一章和十二章.
是一個結束來的.
就是這個奮興會之後.
他們首先做了一個人口普查.
他們住在哪一個地方.
就不是調查哪個地方是貴族區.
哪個是貧民區.
而是那些人散居的地方.
那份外感覺到.
耶路撒冷需要多些人.
所以十一章就是說.
那個分佈.
其中就要強制.
但這個強制是出於感動的.
是出於復興的奮興會之後的.
大家願意.
不會抽了生死簽之後.
咦不關我的事.
我不去的.
就不是的.
大家很樂意.
所以這一章其實很重要.
甘心樂意事奉的人.
必定蒙神的祝福的.
甘心樂意默默地來到去做.
任勞任怨.
不問那個回報.
最終在神的眼中.
神一定會報答.
一定會賞賜的.
所以有很多的百姓.

$^{161}$就為這些被抽中的人.
想耶路撒冷住的人.
為他們祝福.
對了你們去吧.
我們不用了.
我們可以甩爛了.
就不是這樣的氣氛.
這樣的氣氛就很慘了.
但是是很榮幸被抽中.
是很榮幸地.
這些人是能夠住在耶路撒冷.
能夠在聖殿的範圍來服侍.
可能做一些一般的工作.
不是一些很重大的工作.
因為憲制一定與你無關.
你更加不能進入聖殿.
因為你不是利未人.
利未人都不可以進入.
是祭司才可以進入.
所以這些其他的百姓.
他們是做一些其他的庶務的.
不過十二章.
作者再一次告訴我們.
提升那個敬拜.
接著有一個圍到周邊的地方的城牆.
來行潔淨之禮.
接著又有兩隊.
來一路的吹奏著那些聖樂.
來得浩浩蕩蕩的.
上耶路撒冷.
守這個奉獻城牆之禮.
所以這是一個大復興之後的.
一個大的慶典.
這個大的慶典.
是源於以色列人.
整個的心都是被尼希米和以色列.
他們的教導.
聖經的重讀.
和很多的研習班.
而帶來他們的心靈.

$^{201}$是很柔軟.
很願意服從神的帶領.
所以這樣做事就很容易.
也都很簡單.
因為很多百姓的心都是柔軟.
相信最難的就是.
講道的時候.
每個人坐在下面.
好像刀槍不入一樣.
你說什麼都無動於衷.
差不多快點看手錶.
喂 時間到了.
訂了位.
喂 要去了.
諸如此類.
講道的人都感受到.
好像道會反彈一樣.
這樣就很麻煩了.
因為大家都不被神的愛.
神的說話去感動.
這是一個最困難的.
所以第八到第十章.
是我們一個很寶貴的一個勸勉.
教會需要復興.
神的家 神的國度需要復興.
很多的事情.
沒錯是理所當然應該做的.
但是如果那個人未具備.
那個靈明的復興.
所謂理所當然.
在他來說是一個壓力.
可能是一個他不想去參與.
所以這是我們盼望.
我們不要混亂了那個次序.
不是行政指令就可以做到事.
我是很深明白.
就算你說什麼都沒有用.
如果他根本未準備.
甚至會聽錯你所說的.
是很可惜的.

$^{241}$神的家最大的問題.
就是分化.
就是彼此有嚴規.
就是大家不合一.
那就什麼都做不到.
有時候說多好都沒有用.
大家都不是那個心靈去準備好.
很柔軟地讓神在我們中間.
來復興 去工作.
有時候是事倍功半.
反過來如果大家的心靈都很柔軟.
那就事半功倍.
這些是屢見不爽的.
盼望我們能夠在其中去明白.
我們第十二章.
頭十二節就比較悶.
因為是將那些家譜.
特別是祭司的家譜.
來複述一次.
不過原來有個脈絡.
是由第一批.
所羅巴伯上來.
對不起.
回來重建聖殿的那一批.
其實已經死光了.
一百二十幾年前.
已經死光了.
但是這一批人.
他們的下一代.
現在在尼欺米數點的時候.
他們是哪一班的祭司呢.
那就很有意思.
原來那條線.
一路延續.
是可以一路追溯上去.
原來他的太公的太公.
是跟著所羅巴伯上來.
由原居地.
即是被擄的巴比倫.
接著就是馬代波斯.

$^{281}$波斯帝國的時候.
古列王的時候.
准許第一批猶太人.
被擄歸回.
其實已經過了七十年了.
他們就被擄歸回.
回去重建聖殿.
這一批人.
現在全部已經回天家了.
但是現在尼欺米.
找回他們的家譜出來.
去數一下.
數就不只是數他們.
是數一下他們的後代.
現在在哪個位置事奉.
那就很有味道了.
第一節.
同著撒拉鐵的兒子.
所羅巴伯.
第一批上來的.
和耶穌.
大祭司.
所羅巴伯就是省長.
回來的祭司和利未人.
記載下面.
第一批就記載祭司了.
為什麼要這樣記載呢.
就是說.
再一次安頓好那些人之後.
十二章.
就開始再一次.
強化聖殿的事奉者.
他們的家族.
他們的一代一代的傳承.
那個熱心.
一路一路來到第二代.
如果尼欺米.
已經可能去到第三代了.
他們這一班的祭司.
曾經當年跟隨.

$^{321}$由一個被擄之地.
七十年已經安居樂業了.
到第三根了.
現在神說要回去了.
什麼都不理了.
就回去耶路撒冷了.
你想想可歌可泣.
真的帶著妻兒.
帶著整個家族的家當.
就這樣跟隨上去.
回到去其實都是幾人離下.
因為當時耶路撒冷已經被人佔據了.
就變成其實很困難的.
由一個富庶的波斯帝國.
回到一個殘破不堪的.
沒錯是自己的鄉下.
但是一個殘破不堪.
什麼都沒有的耶路撒冷.
所以這樣讀的時候.
有這些背景.
我們再看看這些名字.
雖然已經離開世界了.
但是很值得我們去想一想.
譬如我突然間想起.
田園坤先生.
有其他的.
譬如朱文四牧師.
還有就是我們更加是.
幫我們寫那個迴章的.
John他的爸爸.
我突然間忘記名字了.
這些前輩.
有些可能已經離開世界了.
他們曾經在當年.
有他們那個高瞻遠矚.
來為錦宣家訂那個迴章.
其實很複雜的.
但是他們當年用他們的智慧.
一起這樣去做.
為我們日後教會的發展.

$^{361}$一路這樣蓬勃的時候.
都有一個基礎.
我們看看這裡.
銅雀撒拉鐵的兒子.
所羅巴伯和耶穌亞.
回來的濟斯利未人記在下面.
我記得了.
就是儒也魯先生.
已經離開世界的一位.
叔寧的前輩.
就是Mary的丈夫的爸爸.
張瑪麗姐妹的丈夫的爸爸.
就是儒也魯先生.
是一個學者.
然後告訴我們.
濟斯是西來亞.
耶利米.
以斯拉.
亞馬利亞.
馬祿.
哈特.
史嘉尼.
利弘.
米利末.
易多.
近頓.
阿比亞.
米瓦文.
馬底亞.
碧加.
是馬亞.
約瓦納.
耶大亞.
撒路.
阿穆.
希臘加.
耶大亞.
這些人在耶穌亞的時候.
作濟斯和他們的弟兄的首領.
耶穌亞就是和所羅巴伯.

$^{401}$一起上來重建聖殿的大濟斯.
是他們的族兄弟.
就是這一班有名有姓的濟斯.
記載了.
接著第八節.
利未人就是耶穌亞.
賓內革滅.
士利比.
猶大.
馬他尼.
這馬他尼和他的弟兄管理清澈的事.
後面就會很詳細講.
關於敬拜那件事.
因為下半章是說.
為耶路撒冷建好的城牆.
濟斯好像為每一道門.
很特別.
因為耶路撒冷有很多門.
有糞廠門.
有魚門.
諸如此類.
很多門.
很多名.
每道門都有一個名.
就是那些城門.
為每一道門去祝福.
去祈禱.
好像去高沒它.
我不知道有沒有.
挺特別.
我還沒有試過.
有人叫我為他的金魚.
為他的狗仔.
貓祈禱.
有沒有人為一部車.
我買了一部Tesla.
幫我祝福.
幫我祈禱.
暫時我沒有.
等下我們看完那些經文.

$^{441}$可能給我有些啟發.
濟斯為每一道門.
去祝福祈禱.
走完整個耶路撒冷一圈.
很有趣.
很有特別.
第十二章叫做奉獻城牆之禮.
不知道你有沒有聽過.
所以這些給我們好像很搞笑.
但是原來其實是由大復興之後.
他們立約願意歸神為神的子民.
就是說我們嘗試在我們的生活.
行摩西的律法書.
接著起碼說不行.
最重要就是耶路撒冷要有足夠的人口.
耶路撒冷要有足夠的敬拜.
耶路撒冷要恢復恆常敬拜的秩序.
如果不是百姓慢慢.
一個星期可能好一點.
兩個星期可能不用那麼久.
一個月可能已經打回原形.
沒有了敬拜之後.
慢慢就完全不會再鄭重神的律法.
所以要不斷互相提醒.
不斷被感動.
不斷要透過不同的禮儀.
去讓你對自己的身份.
神子民的身份.
要有這個歸屬.
要有這個認同.
你才能夠一路一路感受到.
神的愛的力量.
在我們這個群體.
在我們大家的心靈裡面.
所以十一章人口普查之後.
就有這個奉獻成祥之禮.
希望我們今天能夠說得完.
因為十三章最後那章.
好景不常.
當黎希回到波斯帝國.

$^{481}$辦理公務之後.
馬上就打回原形.
你不是回來生氣.
把多比亞的房子都扔掉了.
那個房子原先.
十二章和第十章已經說了.
那個是將所有奉獻收到的.
是給聖殿用的物資.
那個房子給了多比亞.
居住霸佔了.
神的殿的庫房.
他都敢住.
真的不怕死.
就是不怕死.
黎希米走了.
回到做公務.
過了一段很長的時間.
黎希米回來之後.
發現整個耶路撒冷.
好像他未曾復興過.
你就可以想而知.
十二十三章是有一個很大的gap.
這個是一個很大的教訓.
我們原來弟兄姊妹.
人類是需要不斷被提醒的.
沒有人坐在那裡自己會熟悉的.
沒有人坐在那裡.
哎呀這麼悶.
先拿本聖經讀一讀.
這麼悶.
現在根本不用擔心了.
只看YouTube.
你看一輩子看不完.
還是你自己有興趣的那些都看不完.
你不感興趣的那些很多的.
你都不可能看的.
那個資料庫多到.
海量這麼多的.
已經不說其他的東西了.
現在就沒有什麼雜誌了.

$^{521}$那些八卦雜誌都沒有人買了.
因為上網太方便了.
但是悶是不會令人讀聖經的.
悶只會更加不讀聖經的.
但是如果有神的愛感動.
他就很自然會覺得.
神的懷孕那個味道.
那種的生命.
對那個人的心靈的改變.
有回平安.
有回盼望.
有回愛.
人不會這麼悶.
內心裡面有一個力量去寬恕.
剛剛有個弟兄的見證.
很戲劇性.
我從來都沒有聽過.
原來他有一個弟兄.
他一直都放在心裡面.
他走這道門我就走第二道走.
真人真事.
他不是開玩笑.
他跟我說.
當然他已經勝過了這些東西.
他跟我說.
真的.
那真的是神的愛.
神的恩典感動.
他最不喜歡那個人.
其實沒有人知道.
是他自己知道的.
居然走過來跟他握手.
幾年沒見了疫情.
他又不好意思.
當下有個掙扎.
伸不伸手出來呢.
終於伸了.
一伸一握.
內心裡面的釋放.
幾年沒見了.

$^{561}$很多年內心的嚴拘一下子解除.
你可以想像得到.
原來很多事情.
神工作帶來的是一個.
沒錯.
表面上好像你要認低威.
其實不是低威.
其實是彼此在神的裡面.
神願意我們來和好.
神願意我們每一個人.
都能夠在愛裡面彼此相交.
這樣帶來的.
其實是你自己囚禁的監獄裡面的我.
那個你自己.
突然間放監.
你突然間不怕了.
我可以釋放我自己了.
他馬上跟太太分享.
跟著跟我分享.
我也覺得很Amazing.
原來男性也會有.
一般我們以為女性容易妒忌.
生氣等等.
但原來是一樣的.
所以復興帶來的就是生命的改變.
這個利未人就是管理清姐的事.
我們等一下就會詳細講清姐.
第九節.
他們的弟兄百步家和烏尼.
照自己的班次與他們相對.
已經說了有二十四班.
這是大衛的時代已經編好了.
就是說大家各按照自己的月份.
那個禮拜的週期.
來到上去聖殿去服事.
能夠服事是一個恩典.
是一個榮耀.
一定要這樣想的.
任何一個崗位.
不要小看那個崗位.

$^{601}$無論多重要的.
也都是或者多次要的.
在神眼中就沒有所謂重要不重要.
是那個人本身的態度最重要.
所以如果輕佻.
如果不放在心上.
哪個崗位都被你搞到.
沒有那種榮耀.
是因為你令到那個崗位不榮耀.
所以記得我們要留心.
神給我們每一個崗位.
每一個恩賜.
都是一個祝福.
我們能夠去服事是一個祝福.
我們能夠服事是一份的恩典.
還有記得是一個榮耀.
所以十一章被抽中的人是榮耀.
所以求神幫助我們.
讓我們每一個弟兄姊妹.
能夠聽到這個訊息.
都能夠被感動.
被神的愛感動.
能夠被神去揀選事奉.
是一個很大的榮耀.
第十節.
耶穌亞生若拿金.
若拿金生以利亞實.
以利亞實生耶和耶大.
耶和耶大生若拿丹.
若拿丹生亞道亞.
在若拿金的時候.
這些已經回到古代的人名.
祭司作族長的西來亞族.
有米拉雅.
耶利米族.
有哈拿尼雅.
整個家族作祭司.
當然他們是亞倫的後人.
一定是作祭司.
不過他們都分了不同的族.

$^{641}$以斯拉族有米舒蘭.
亞瑪利亞族有約哈蘭.
米利古族有若拿丹.
士巴尼族有約瑟.
哈林族有阿拿.
米拉約族有希臘海.
亦多族有撒迦利亞.
近頓族有米舒蘭.
阿比瓦族有.
細基利未拿民族.
某.
這裡特別一點.
米拿民族某摩亞底族.
有比勒泰.
碧加族有沙姆瓦.
善瑪雅族有約拿丹.
約瓦納族有馬特乃.
耶德雅族有烏西.
撒萊族有加萊.
阿穆族有希伯.
希臘加族有哈沙比亞.
耶德雅族有拿坦葉.
這個就是個別的祭司.
第一批上去的.
和他們的家族.
即是意味著家族.
他們是會有一個傳承的.
接著就說第二批.
就是說利未人.
第22節開始.
至於利未人.
當以利亞什.
耶和耶大.
約哈蘭.
即是剛才那個.
祭司的家族那裡.
亞道雅的時候.
他們的族長記在冊上.
波斯王大利烏在位的時候.
作族長的祭司也記在冊上.

$^{681}$利未人作族長的記在歷史上.
直到以利亞什的兒子.
約哈蘭的時候.
即是說約哈蘭這個.
祭司利未人.
是有相當巨的代表性.
所以以這個人作為一個指標.
第24節.
利未人的族長是哈沙比亞.
史利比革滅的兒子耶穌亞.
與他們的弟兄班次相對.
照著神人大衛的命令.
一班一班的讚美.
稱謝.
很特別的.
大衛被提升到神人.
不是因為他被築聖.
是在神的眼中.
這一位君王.
是無可指責.
除了在這個拔士巴的事件.
他是.
甘心樂意去遵行神的旨意.
現在他.
當年臨死之前.
聖殿都未做好.
但他已經為聖殿將來的規模.
指示他的兒子所羅門.
再加上不只是硬件的預備材料.
是預備規模.
就是關乎敬拜.
關乎到聖殿的.
生生不息的敬拜的服侍.
就將利未人.
因為你進到聖殿之後.
那個聖殿就是永久性的.
除非被擄.
除非有敵人來拆.
不然就是永久性的.
不需要擔台的.

$^{721}$不需要再像在曠野.
來到一個不斷的漂流的行宮.
不斷的走.
所以既然是永久性.
利未人就沒有事做.
不用搬了.
因為規定了.
哪些支派搬的.
例如聖殿寺聖所裡面的.
藥櫃.
金燈台.
香壇和餅桌.
哪些人抬的.
就是哥核子孫抬的.
其他就是有.
米拉利子孫.
另外那些外圍的.
帳幕.
那些木筍.
等等.
就由另外的家族.
利美的家族的人.
去抬的.
所以分好了.
所以其實.
固定的聖殿.
就是大量利未人失業.
但是多有智慧.
不關他的事.
他都沒有權去建一個聖殿.
所以這個就是.
眼光夠不夠.
胸襟夠不夠大.
你一世都沒有份在榮耀的了.
但是你會不會為下一代.
那個聖殿將來真的建好之後.
你的那個智慧.
神給你的感動.
能夠去傳承下去呢.
能夠根據你的班次.

$^{761}$你定好了的.
這樣去實踐出來.
所以他就將利未人.
分做不同的班次.
而且按照著.
他的恩賜.
來設立了這些詩歌的訓練.
的詩歌師班長.
所以讀敬拜隊.
可能都需要讀一下這一段經文.
這裡提到大衛.
神人大衛.
就將他們分成班次.
是一班一班的讚美稱謝.
就是說.
在聖殿其中比較顯耀的工作.
就是引領群眾.
來獻祭的時候.
是敬拜神.
朝見神.
用希伯來的詩歌.
用不同的樂器.
來去頌揚神.
或者他們可能會繼續作曲.
有當時的.
先知.
來寫書.
然後他們可能會.
引導百姓.
在那個靈感裡面寫詩篇.
或者是有其他的.
一些智慧的文學.
諸如此類.
一定是祭師利未人.
他們要研習.
利未人不是敬拜的時候.
有一些是做教導的.
開研習班的.
教導聖經的.
祭師在平常日子.

$^{801}$是要審判的.
做一些像和事佬.
長老的角色.
大家都信任.
大祭師.
大家都將那些紛爭.
可能你對他對.
他也對他對.
由祭師去做評判.
做一個調和.
調解.
其實在舊約都出現過.
譬如路得記出現過.
波亞斯.
本來.
那個人是可以娶路德的.
就一起去到.
那些長老.
城門口的長老的面前.
來處理.
現在路德.
到底可不可以.
嫁波亞斯呢.
波亞斯說不行.
因為有一個更近的親屬.
比我更近的.
他就優先可以娶你.
他不娶你.
但波亞斯還是很厚道.
自己喜歡的女子.
因為那個人如果說.
我要啊.
我多大都要啊.
那就沒有了.
這些愛情故事.
不是作出來的.
你可以想像得到.
波亞斯和那個人.
都去到那些城門的長老.
面前去求.

$^{841}$那個排解的糾紛.
所以這裡就告訴我們.
其中一個部分.
是清謝.
第25節.
瑪他尼.
達布加.
俄巴底亞.
米蘇蘭達門.
阿谷是守門的.
就是在庫房那裡守門.
這都是在約撒達的孫子.
耶穌亞的兒子.
再一次提約瓦金.
我今天不夠時間.
因為下午出去了一個緊急探訪.
在醫院有個尼留病人.
我下星期前.
我就將約瓦金.
因為這個人很重要.
我找了資料之後.
我就在群組給大家看.
一而再提到他.
好像一個指標性的時代.
一個人.
和省長尼希米.
就是尼希米的時代.
並濟斯文氏以色列的時候.
有職任的.
不知不覺間.
一講一講.
數那些濟斯利未人.
數那些家族.
雖然很悶.
但數數就數到省長尼希米.
文氏以色列的時代.
不知不覺間.
在這26節裡.
由第一代的索羅巴伯.
由古列王的時候.

$^{881}$不理那麼多.
我們現在能夠回去自己的地方.
耶路撒冷.
我們就舉家回去.
濟斯我更加要帶頭回去.
現在尼希米記.
在這12章裡.
好像都寄一筆.
神的honor.
給這些濟斯.
由第一代回來.
到今天.
現在我尼希米.
在這裡來做一個記錄的時候.
都為這些第一代.
移居回到耶路撒冷的濟斯利未人.
都給一些credit.
給一些honor.
神都紀念他們.
當年什麼都不理.
毅然的.
因為濟斯的工作.
唯一的工作就是要在聖殿工作.
他要回去.
以身作則.
回去建造聖殿.
可想而知.
神原來.
一直留意著.
這些家族.
到今天.
尼希米120多年之後.
第三批.
重建聖文的.
這個城牆建造.
尼希米.
這樣來到.
回去的時候.
其實已經是人面全非.
但是家譜.

$^{921}$就是一個很好的.
一個懷念.
家譜就不能騙人.
第一批上去.
有記錄的.
現在他們的家族.
已經還在一起.
在這個城牆建造裡面.
和在現在.
那個聖殿的服侍裡面.
這樣就很美麗了.
就好像大德生.
第六代第七代的子孫.
現在還可以在.
中國人的地方事奉.
這樣就是一件很美麗的事.
有些家族.
可能已經失去了.
但是如果他的後人.
我不知道.
明仔會不會結婚生子.
如果我已經不在了.
離開世界回天家.
到第十代.
有人去做那個普查.
追下追下追下.
數一下那些家譜.
數到明仔的爸爸原來是牧師.
十代之後又有一個.
在培靈會講道的.
姓曾的.
這樣就很美麗了.
尼希米就是這樣記載.
這些人.
其實很想講.
就是由第一代.
透過那個家譜的追索.
可以連繫到現在.
尼希米自己的時代.
你羨慕不羨慕.

$^{961}$你的後代.
能夠這樣被神使用呢.
原來太公.
你的太公的太公原來是這麼敬虔.
為十代之後祈禱.
是為將來.
我的後裔都能夠.
繼續很謙卑的服侍神.
不希望他們飛黃騰達.
但希望他們能夠.
在神的愛裡面.
去服侍.
為神的福音去努力.
齊心.
這就是最美麗最好的遺產.
最好的長輩的叮嚀.
我很深刻.
當然我不是直接參與.
有一個.
弟兄.
真的講家族的.
是一位很敬虔的弟兄.
他有很多兄弟.
都是很熱心事奉.
他的兄弟還在事奉.
還沒死.
他很早就發現有癌症.
很年輕.
很有為的一個弟兄.
他叫黃恩華.
很早期已經死了.
我初信的時候.
已經是二十多年前.
已經離開世界了.
他臨死.
我聽前輩說.
臨死,年輕人.
但他的家族是很熱心.
在教會事奉的.
到今天都是這樣.

$^{1001}$勉勵他的兄弟.
他的太太.
他的子孫.
一定要熱心事奉神.
這樣離開世界.
癌症.
這一位弟兄.
我不認識他的,因為我沒見過他.
但他的兄弟我認識.
有黃炳華.
有黃國華.
幾位兄弟.
某程度上在我們的武堂.
是很出名的.
是一些執事或長老.
相當熱心服侍.
這個哥哥.
黃恩華.
就是這樣熱心服侍神.
這些就是.
你不要小看這個省長.
原來他這麼細微.
由第一代的祭司.
一直追溯.
當然我們不是猶太人.
我們沒什麼感情可言.
我們看到的就是一大堆人命.
關我們什麼事呢.
但是原來.
由第一代來到第三代.
那個感情線就很深了.
那個的敬虔.
就呼之欲出了.
就會看到.
一個家庭教育.
用神的話語.
去勉勵自己的後代.
那個的敬虔.
如果這樣能夠延續.
那就很寶貴了.

$^{1041}$真的是.
默默的.
你做多少.
你怎樣栽種你的兒女的心田.
你是不是幫他除草.
是不是讓他的心田.
是一個好土.
來真的依靠神.
到有一天.
你都不可以再教他的時候.
或者他都不需要你再教的時候.
他自己會去敬虔的.
依靠神.
他自己已經傳承了.
上一代長輩給他的生命.
那就是我們.
可以含笑離世的時候.
真的.
如果不是.
我們留的是什麼資產.
給我們的下一代呢.
我們留下的是什麼給他呢.
或者.
我們在他的心裡面.
到底是怎樣的一個評價呢.
我的父母.
真是很敬虔的.
那你一世人就.
真真正正是最.
最值得安難.
因為是由你自己的兒女.
講你自己.
人家講都是比較外圍.
如果由你的兒女去講.
這個是很重要的.
弟兄姊妹.
我們去到下面.
27節.
開始講這個奉獻禮.
所以你看到.

$^{1081}$尼熙米.
這28節聖經.
26節聖經.
記載這些第一批上來的.
一路追溯他們的後代.
然後.
強調.
祭司利未人.
他們的家族.
他們的事奉的延續.
其實很想再凸顯.
就是聖殿.
如果沒有了敬虔的.
祭司利未人.
就算那個聖殿多麼的金碧輝煌.
如果那個內容.
不是獻祭帶領回眾.
敬拜神.
真是去到神面前.
是帶著一個很單純的.
心靈敬拜.
那個很美的聖殿.
最終有一天還是會被神.
毀滅.
所以你會看到.
尼熙米很緊張.
不是他自己說我是一個.
這麼有功勞.
在阿達舍西王的面前.
在猶大百姓的.
長老面前.
在其他的祭司利未人面前.
我是最大功勞.
如果不是我.
怎麼能夠撥亂反正那些外邦人.
不會再搞事.
又能夠令到上下一心.
去重建這個城牆.
是很大工程.
尼熙米他不記他自己.

$^{1121}$那個功勞.
他現在要讓百姓看重.
最重要的.
就是整個聖殿.
最核心就是敬拜.
帶領敬拜的人.
他們那個家族.
他們那個敬虔.
是否能夠一路延續.
所以你看到尼熙米.
很緊張的就是.
在他身邊的祭司利未人.
他們到底有沒有一個.
氣質.
好像他們的父輩們.
來看到自己的事奉.
是很尊貴的.
是神揀選我們家族的.
來服侍神.
27節開始.
先喝口水.
(喝口水).
耶路撒冷城牆.
構成的時候.
眾民就把各處的利未人.
招到耶路撒冷.
要清澈歌唱.
敲拔鼓瑟彈琴.
即是說多元的樂器.
做物主是真的榮耀的.
是他創造.
多元化.
所以有鋼琴.
有風琴 有結他.
有其他的樂器.
所以我們要開放一點.
讓不同的敬拜.
只要是大家.
well prepare.
敬拜的人一心一意是敬拜神.

$^{1161}$不是炫耀自己的技巧.
不是自己的歌喉.
自己的聲線.
那些一定是重要的.
如果你每個走音的.
美麗可言.
所以兩者.
本來就不應該有排斥.
但是.
往往人很容易注視.
人怎樣看你的敬拜.
而你忘記了.
其實應該是緊張神.
怎樣看我們大家的.
心靈的敬拜.
這個是重要的.
歡歡喜喜的.
行告城之禮.
歌唱的人從耶路撒冷的周圍.
和尼陀佛的村莊.
與白吉甲.
又從加巴和.
阿瑪弗的田地聚集.
因為歌唱的人在耶路撒冷.
四圍為自己.
納了村莊.
祭司和利未人就結證自己.
也結證百姓.
和城門並城牆.
很特別.
怎樣結證城門城牆呢.
是不是真的用一些膏油.
來抹一下.
分別為聖.
這個整個城牆.
整個城門.
是要奉獻給神.
現在是第一次.
好像我們中國人那樣開光.
或者新舖開張.

$^{1201}$一般就燒豬.
剪綵.
猶太人就.
這個城牆是神的.
即是說這個聖城是神聖的聖城.
所以.
那些城牆.
是要幫它抹乾淨.
是行結證之禮.
雖然沒有詳細說.
是不是真的用一些膏油來抹.
分別為聖.
這樣來看.
而且.
又強調第十節.
就是事奉的祭司利未人.
首先自己結證.
即是自己明白.
現在我的事奉.
是事奉神的.
我現在自己首先整理.
我的身心靈.
調教我自己的心態.
來服侍.
準備好自己.
然後帶領整個百姓.
一起.
意味著分開不同的城門.
有不同的班次.
來一起.
去結證.
第三十一節.
我就是黎謙慧.
帶猶太的首領上城.
使清謝的人分為兩大隊.
排列而行.
好像我們進殿一樣.
第一隊在城上.
往右邊.
向這個糞廠門.

$^{1241}$其中一個門.
糞廠門走.
在他們的後頭.
有荷沙瓦與猶太首領的一半.
又有亞撒利亞.
以斯拉,米蘇蘭.
猶太,卞瓦敏.
善瓦瓦,耶利米.
還有些吹號的祭司的子孫.
約拿丹的兒子.
撒迦利亞.
約拿丹是善瓦瓦的兒子.
善瓦瓦是瑪他尼的兒子.
瑪他尼是米該亞的兒子.
米該亞是撒克的兒子.
撒克是亞撒的兒子.
又有撒迦利亞的弟兄.
善瓦瓦,亞撒利.
米拉萊,基拉萊.
瑪艾,拿坦葉.
猶太,哈拿利.
都拿著神人的領袖.
善瓦瓦.
都拿著神人大位的樂器.
文氏,以斯拉.
引領他們.
有時記載一些重要的領袖.
就要說他是誰的兒子.
誰的兒子.
要追溯他的家譜.
即是說.
他們幾代都是在.
這個殘破不堪的聖殿.
聖城.
來留守.
來事奉.
不過現在.
是一個大復興之後的.
奉獻之禮.
他們就.

$^{1281}$真的非常感恩.
他們等待到這一天.
他們的家族延續到今天.
現在尼欽米.
和以斯拉就帶領著群眾.
一齊的來敬拜.
奉獻整個聖城.
就以.
真的很實物教材.
抹下那些牆.
我見過阮妮.
在建這個教堂的時候.
一幅幅那些牆.
頂紙梅.
奉獻在那裡畫畫.
類似.
這些相片是很珍貴的.
你開想而知.
這些就是當年.
這個教堂是整個很標誌性.
在整個感受這個社區.
屹立一個十字架.
屹立一個燈台.
是證明神的.
祝福在這個社區.
是很標誌性的.
所以某程度上就像這一段經文.
分別為性.
所以現在祭司就帶領著.
分兩隊.
可能由南北去.
尼欽米就帶一隊.
另外一個就帶另一隊.
這樣來到整個城牆.
其實不是很大.
可能.
幾公里就已經完了.
但是.
這個意義就是.
外邦人你們是沒有份的.

$^{1321}$整個城牆是屬神的.
我們這班人.
現在住在耶路撒冷的.
在周邊的我們都是神的選民.
現在.
我們將整個神聖的.
這個地方.
耶路撒冷.
以抹牆.
抹城門做一個奉獻.
是屬神的.
接著就是.
這些參與事奉的.
原來他們的祖先.
是第一批上來的人.
或者是早期.
被擄之前的後裔.
即是他們的.
長輩.
就是現在他們的後裔.
我們繼續讀下去.
後面就會再有一些經文.
給我們提醒.
神人的.
惡氣現在都預備了.
是複製.
還是被擄的時候帶過去.
聖經就好像沒有這個線索.
不知道是不是照那個.
藍圖去再做.
來複製大衛的惡氣.
第二次提神人大衛.
所以大衛.
他在神眼中是合神心意.
他為神.
真的做過很多的事.
他真的.
仆心仆命地為神.
他不計較任何事.
你說他是王當然不計較.

$^{1361}$什麼都是他的.
他的心意是神預立的.
所以他為神.
想了很多很多日後.
人們說是.
遺臭萬年.
他是流芳百世.
他的建議,他的智慧.
和他創立的東西.
到今天.
尼希米帶領這個.
時代.
都用回他的惡氣.
用回他的班次,利未人的分班.
可想而知.
神都預立.
大衛.
所編排.
利未人的新崗位.
不是.
墨守成規.
不理會,我是抬的.
你怎麼知道要給我抬.
都不用搬,抬什麼.
都沒有東西要抬.
還有那個聖殿都是永久的,怎麼抬.
沒得抬.
轉做第二個事奉崗位.
學習新的東西.
隨著時代不同.
你的身份,你的工作.
都要有些調整.
你是利未人,沒錯,但你可以學習.
不是搬台.
你學習敬拜.
學習在聖殿裡事奉.
學習吹不同的惡氣.
你就可以想而知.
原來那種與時並進.
利未人都要學習.

$^{1401}$上一代的傳道人.
穿長衣.
三十年代還穿.
可能到.
七十年代就少了.
西裝開始洋化.
以前的上一代.
國內的.
真的長衣馬褂.
時代不同.
所以有時不要墨守成規.
真的要與時並進.
很多東西是有轉變的.
利未人都應該要轉變.
摩西五經的時候.
從來沒有說過他們要做這些惡氣的敬拜.
但是大衛的年代.
他們就設立了.
這些敬拜的師班長.
來訓練利未人.
你就可以想而知.
原來真的有不同的時代.
敬拜可以是.
多元化的.
但敬拜的是同一位真神.
由文氏以色列引領他們.
這是三十六節.
第三十七節.
他們經過全門往前.
從大衛城的台階.
隨地勢而上.
大衛城在山上.
由.
希西加的船.
出來.
西羅亞池.
最低.
是很斜的.
就是耶路撒冷上.
你就可以想而知.

$^{1441}$他們一路隨著山勢地勢.
上來.
一路走到.
朝東的水門.
一道一道門.
分別為聖.
來去潔淨.
第三十八節.
第二對稱謝的人.
要與那一對相迎而行.
我和民的一半.
跟隨他們.
在城牆上過了.
牢樓直到寬牆.
又過了以法連門.
古門 魚門.
南葉樓 哈米亞樓.
有些門是有樓房的.
有些就沒有.
直到羊門.
羊門大家記得.
在耶穌的時代也有.
接著就在護衛門站著.
於是這兩對稱謝的人.
連我和官長的一半.
站在神的殿裡.
還有祭司.
以利亞金.
瑪西亞.
米該亞.
以利約乃.
撒迦利亞 哈南尼亞.
又有瑪西亞.
是馬亞.
以利亞薩 烏西.
約哈蘭 瑪基亞.
以蘭 以謝.
奏樂 歌唱的就.
大聲歌唱.
伊斯拉希亞管理他們.

$^{1481}$那日眾人.
獻大祭而歡樂.
意味著.
不只是敬拜.
唱詩歌.
沿著整個城牆.
整個不同的城門.
來去敬拜.
浩浩蕩蕩分兩隊.
雙形.
大家穿梭.
大家過對面.
他又過這邊.
這樣來到.
整個城充滿了.
歡樂的敬拜.
用大衛的樂器.
有師班長引領.
來歌唱.
好像整個耶路撒冷.
在大復興的時候.
這就是高潮.
整個聖城.
充滿了對神的歡樂.
即是說那些歡樂聲.
已經去到上達於天.
當然他們整個耶路撒冷.
在城牆內.
充滿了.
這份敬拜的.
恩恩喜樂和讚美.
在這個時候.
特別提到.
第四十三節.
那日眾人獻大祭而歡樂.
因為神使到他們.
大大歡樂.
連婦女帶孩童也都歡樂.
甚至耶路撒冷中的歡聲.
聽到遠處.

$^{1521}$你可以想而知.
那個歡樂聲.
一路一路隨著那些歌唱.
飄到很遠的地方.
很遠的地方.
也意味著.
有些人上不到聖殿.
因為太小了.
其他的猶大人.
住在不同的周邊城市.
都聽到敬拜歡樂的聲音.
有些在耶路撒冷的婦人.
孩童.
本來他們沒有參加.
但是他們也有份於.
這份歡樂.
某程度上.
又不可以完全拒絕.
婦孺和孩童.
他們也有份.
於敬拜.
這個就是.
大復興帶來的.
第二個訊息.
人口普查.
安置人口.
甘心樂意.
是被祝福的.
這是十一章第二節.
百姓為這些被抽中的人.
祝福.
你們真的有福了.
能夠住在耶路撒冷.
所以甘心樂意事奉的人.
神是會有祝福的.
第十二章.
告訴我們.
祭司利未人.
而且是尼熙米.
一起事奉的.

$^{1561}$祭司利未人.
原來可以追溯他們的長輩.
曾經被老歸回.
也是以前一直在這個行列.
事奉.
生生不息.
代代相傳.
來服侍神.
很美麗的圖畫.
如果我們懂得看.
可以看到很美麗的圖畫.
不要只看很負面的.
教會不要只看負面的.
看看默默去做.
有時他回祖堂.
回家他都跟我說.
太少年青人回祖堂.
有時看看老人家.
多忠心.
你不要理會他回來是否只是睡覺.
他回來敬拜神.
年青人就是要學這些.
學一下別人真是這樣的.
一個心.
不識字的他都回來.
這些是一些年青人的內心.
其實是一個感動.
應該要這樣看這些敬虔.
有時敬虔.
是隨處可以.
看到的.
你看人家這麼困難.
推著輪椅都要回來敬拜.
譬如二妹等等.
你可想而知原來.
沒錯你是一個年青人.
神對你的要求有另外的要求.
可能你有更大的事情.
神要用你.
但是有時你要感覺到.

$^{1601}$這種立體的感動.
如果不是.
你只看一些很負面的東西.
教會在你眼中.
不行不行不行.
你的負能量.
就會帶來整個教會.
受影響.
你看人家這些事奉者.
看這些.
甘心樂意住在耶路撒冷的.
看尼欽米韻.
這些事奉人員.
他們的祖先原來.
都是一直參與事奉.
有些時代很艱難.
被那些外邦的.
迦南人奚落.
嘲笑.
城門破損.
又被火焚燒.
百姓受盡欺凌.
但是在這麼荒涼的.
警方底下.
祭司利未人仍然.
繼續留守.
來到這個聖殿.
繼續為那個時代.
去守望祈禱.
這些是我們要欣賞.
我們要尊重.
你就可以想像到.
教會是真的需要.
有老中青兒童.
一起敬拜.
不可以只有小朋友.
不可以只有青年人敬拜.
如果沒有了.
上一代下一代.
其實不健康.

$^{1641}$因為大家都在想事的時候.
有時他不懂得想得闊一點.
這些是我們設計的時候.
應該要幫助不同的階層.
不同的年齡層.
去看到教會.
是來自不同的弟兄姊妹.
不同背景.
不同的人生遭遇.
帶來.
救恩的.
榮美救恩的.
偉大.
求神幫助我們能夠這樣看.
這裡就告訴我們.
寬聲聽到遠處.
第十四十四節.
開始說.
敬拜.
這個留了伏筆.
因為十三章.
有人夠膽將庫房.
霸佔庫房.
神的庫房.
當日派人.
管理庫房將舉祭.
初熟之物和所取的.
十分之一.
在按國承田地.
照律法所定.
歸給祭司.
和利未人的份都收在裡頭.
尤大人因.
祭司和利未人.
功績就歡樂了.
即是說.
有這些.
服侍主的祭司利未人.
專心的事奉.
帶領我們敬拜.

$^{1681}$帶領我們奉獻.
帶領我們去聆聽.
摩西的律法書.
他們就安心.
就歡樂.
百姓的責任就是要供奉.
這些專職的事奉人員.
他們能夠在.
神的安排下.
他們都能夠.
稱心的.
富之憂的去專心的.
服侍神.
這個很重要.
他們就歡樂了.
四十三節.
對不起四十五節.
祭司和利未人遵守神所吩咐的.
並守潔淨的禮.
歌唱的守門的.
照著大衛和他的兒子.
所羅門的命令.
都如此行.
古時在大衛和亞薩的日子.
有歌唱的令長.
並有讚美.
稱謝神的詩歌.
當所羅巴伯和尼希米的時候.
第一代和第三代的時候.
以色列眾人.
將歌唱守門.
每日所當得的份.
供給他們.
又給利未人當得的份.
利未人又給亞倫的子孫.
當得的份.
出現了好幾次了.
就是說.
簡單來說.
由大衛.

$^{1721}$最古老的日子.
最好的日子.
因為他設計聖殿.
他的兒子建好之後.
聖殿所羅門就行奉獻禮.
接著就有祭司利未人.
來事奉.
在神的殿裡工作.
事奉.
現在尼希米想告訴我們.
他最敢榮耀的.
不是建好城牆.
他最敢榮耀的.
是提升百姓的幸福感.
那個身為神子民的幸福榮譽.
因為已經有城牆.
和外邦人隔開了.
我們才是屬神的.
我們是有那個的.
不是責任.
是我們有這個的福分.
可以去到耶路撒冷.
朝見敬拜和事奉.
現在更重要的是.
聖殿需要恢復它的最輝煌.
最輝煌就是.
不同的崗位.
有專職的祭司利未人.
來提令守門的.
來到去在他們的崗位.
去盡忠職守.
使到聖殿恢復昔日的榮耀.
那個光彩.
那個功能.
這樣百姓就安心了.
因為他們知道.
整個地方有神同在.
他們就能夠.
可以很專心地.
回去敬拜完之後.

$^{1761}$回到他們自己的原居地.
回到他們自己的家庭.
來勤奮地為神工作.
勤奮地努力.
按照著當德的份.
帶上去耶路撒冷.
來讓神的聖殿有那個供應.
神是透過這個系統.
來祝福聖殿.
聖殿有一個正常和活潑的敬拜規模.
百姓就會感受到那份福氣.
感受到那個祝福.
感受到他們是蒙福.
他們就會很安心地.
忠心地在他們的崗位工作.
和為聖殿的事奉人員打算.
和為他們去祝福.
因為神的功能.
是當德的份.
這裡已經再一次提到.
所以這是一個很重要的事功.
曾經也有些人說.
教會不需要有傳道人.
大家都業餘就行了.
簡單很多.
可以肯定是不可行的.
一定不可行的.
因為到了真正事奉的時候.
原來是沒有準備.
大家都不是想著這樣的運作.
這樣就會帶來.
只會是混亂.
是不會有任何可以改變到的事情.
所以舊約已經是一個很重要的參考.
新約的時候.
其實舊約去到某個階段.
由這個殺母儀的時候.
已經開始有先知門徒的制度.
有系統地訓練先知.
讓他們能夠成為神的代言人.

$^{1801}$來將神的心意啟示.
來告訴百姓.
這個制度正正就是傳道人的系統.
所以到了新約.
很清楚神是呼召一些傳道人.
比如十二使徒.
就是耶穌的呼召.
他們就成為了教會的第一批的牧者.
的長老.
的使徒.
當然所有都是殉道.
因為是一個Big Bad時期.
所以一直就是這樣.
後來神再呼召保羅.
保羅在他的事奉.
他就選了很多有為的年青人.
不同的民族.
來跟隨他去事奉.
比如有提多.
提摩泰.
特羅非摩.
更有路加還要跟隨他.
成為一個好像記者一樣.
追隨保羅去傳道.
你會看到.
有這個使徒制的訓練系統.
來到有專職的傳道人.
來到教會做宣教士.
在不同的地方.
去開闢新的工場等等.
你會看到這些初型.
一直在發展.
嚴格來說就是.
有專職的.
在教會裡事奉.
守望.
禱告.
所以傳道人太忙都不好.
應該是多傳道.
多祈禱.

$^{1841}$多探訪.
不要做那麼多行政的事.
真的好好地預備.
跪下預備祈禱.
為頂子們的需要禱告.
我們童工會.
下半段我一定要引導童工.
有什麼頂子們的需要我們禱告.
大家不同的部.
就說一下一些情況.
我們要祈禱.
所以不只是開會.
談那些事.
突發性的重大事.
才談多一點.
一般我們都是為頂子們的需要禱告.
牧者是要有守望者的角色.
星期三晚祈禱會.
是我自己很看重.
真是牧者我們要聚在一起.
為有需要的頂子們祈禱.
為世界守望禱告.
為教會的發展.
聖功去禱告.
很重要的.
頂子們忙碌.
回不來不要緊.
回到來就是恩典.
我們一起參與一起祈禱.
就是這樣.
教會就是要有這些.
登台的專職事奉人員.
一生奉獻給神.
來服侍.
那個專心.
你喝水就飽了.
神會供應你的.
你不要搞錯.
神到今天還是透過頂子們的奉獻.
所以如果馬會說.

$^{1881}$曾牧師我聽聞錦繡堂.
好像也有些困難.
我在馬會奉獻一億.
給你們要不要.
你不用想一定不要.
因為那些不是我們價值觀相同的東西.
不是那些錢有問題.
是機構的宗旨和我們不同.
我們不是要抗拒那些錢.
而是那個性質我們不要.
我們要的是頂子們.
憑愛心神感動的奉獻.
我們暫時也不是要政府的錢.
德國的教會是要的.
是國家教會.
香港沒有的.
我們是頂子們的奉獻.
就是說神感動頂子們給多少我們.
我們就做多少.
神是無限.
但是我們的資源是有限.
所以神給我們有多大的規模.
我們就做多少.
就不要浩高無遠.
這是我個人的事奉信念.
所以這裡很重要.
有專職的事奉的祭司利未人.
讓聖殿有秩序.
有規模.
有活潑.
活潑就是說有生氣.
那些百姓就歡樂了.
然後黎希米是省長.
他要行一些措施.
就是鼓勵感動眾百姓不同支派的.
你按神給你的恩典.
你的修成.
你來奉獻給神.
其實就是奉獻給神的聖殿的使用.
記住這個庫房.

$^{1921}$去到第十三章就被人霸佔了.
這個可能我們下個星期再說.
今天就去到答八.
我就說完這一章聖經了.
有沒有什麼想問的.
或者想分享的.
當然是難的.
因為在網上又不能直接分享.
所以我一直去.
多謝大家給我有機會.
真的能夠有時間去預備這些茶經.
如果不是和大家茶經.
我都可能未必會這麼仔細再看.
看來看去.
就發現原來.
這些事情不是這麼簡單.
原來第七章發現了人口的問題.
其實每個人都看到.
但沒有人願意住在耶路撒冷.
可想而知.
尼希米和以色列.
他們才是被神感動.
有那個智慧.
才能夠和人一起讀經.
奮興.
然後才再看.
大家一起有那個眼光去看.
耶路撒冷的荒涼.
遲早都會有問題.
遲早會出事.
你們建好城牆之後.
每個人就回到自己的居所.
耶路撒冷繼續荒涼.
這樣就很凋零.
就算聖殿建好.
都要有人去打理守衛.
這樣又沒有人管的時候.
人口這麼少的時候怎麼辦呢.
先有第八.
第九.

$^{1961}$第十.
這個復興會.
然後十一章.
舊事重提.
第一節.
就是說.
百姓的首領住在耶路撒冷.
其餘的百姓.
每十個人當中.
需要一個人.
來住在聖城耶路撒冷.
那九個人.
就住在自己的城邑.
凡是甘心樂意住在耶路撒冷的.
百姓都為他們祝福.
你就可想而知.
大家都看到了.
大家都願意.
接受神聖的差遣.
就以抽籤的形式.
抽到原來是好像中了居屋一樣.
開心到死.
就不是說抽中我就慘了.
我要現在整個家庭搬來耶路撒冷.
很荒涼的.
什麼都不方便的.
原來不是這樣.
這個就是需要百姓的心靈.
受到開懇.
需要有憤青的聚會.
使百姓的心靈得到建立.
那就水到渠成了.
大家是很感覺榮耀.
去到聖殿周圍居住.
看守事務或其他工作.
他們都不介意.
要搬走自己的家庭.
都不介意.
這是相當重要的.
十二章再一次告訴我們.

$^{2001}$那個事奉是要有一個延續.
不要一步十漢.
是要一個傳承的系統.
在自己的家庭.
在自己的教會.
一起真心實意的事奉神.
不要計較困難.
有些時代可能是復興時代.
有些時代是很暗淡的.
很艱難的.
但是都默默的緊守住.
不要輕言說我要離開.
我要和這個教會一刀兩斷.
不要輕言這樣.
總會有時人會有彼此的嫌隙.
那些是很小事.
總是在神恩典裡尋求復和.
所以不要看自己這麼大.
好像膨脹一樣.
反而要緊張的是.
我的下一代.
是否和我一樣熱心事奉.
和他的熱心是否他自己.
發現神的愛.
在他內心裡.
他自己明白.
他自己參與.
小朋友當然容易.
父母帶他回家.
所以到了中二,三.
他可以不回家的時候.
就不回家了.
可能藉著一次.
發高燒不回家了.
原來不回教會這麼多時間.
原來父母回了教會之後.
我這麼自由.
他開始有一個試探.
慢慢慢慢一次,兩次,三次.
就變成一個.

$^{2041}$好像一個約定俗成.
其實就不好了.
應該我們要怎樣能夠.
從一種勸勉.
從一種鼓勵.
不要是一種威逼.
更加不要用利誘.
應該就是讓他看到.
回教會是蒙福的.
父母是很渴慕回教會的.
這樣就很不同了.
所以盼望我們都能夠.
透過這五章聖經.
表面是歷史書.
原來是一次的奮興會.
奮興會之後的.
聖城正正式式的奉獻之禮.
一個這樣的禮儀.
原來曾經有這樣的前輩.
為聖城的城牆來高滅過.
所以上一次如果我們有機會讀完.
才去耶路撒冷.
就很不同了.
因為那個西牆.
西面的牆.
就是所羅門的牆.
現在還在的.
曾經有人抹過.
活過.
用一些我不知道是用什麼的膏油.
曾經有人抹過.
是不是黎希敏抹過呢.
這樣就很特別了.
告訴我們原來.
從前的人為神的家.
擺下過很多他們的心血.
他們的生命.
所以我們不要輕言.
我們隨隨便便就說.
我要離開教會.

$^{2081}$我要離開神.
不要這樣.
看看不同的人.
看看不同的艱難時代.
人們是怎樣為神.
去守望著教會.
譬如文革時代.
很多的婦女.
留守在他們自己的鄉村.
教會.
傳道人很多都要坐牢.
那些婦女.
就成為傳道人.
來守望著那條鄉村.
很多這些無名的傳道者的故事.
所以求神幫助我們.
讓我們看到.
十二章真的很美麗.
那個歡聲.
一個潔淨之禮.
祭司利未人浩浩蕩蕩的.
在不同的門.
吹奏那些樂器.
大衛的樂器.
又有些祭司.
來高沒那些城門.
城牆.
來奉獻整個耶路撒冷給神.
百姓就感覺到福氣.
百姓感覺到自己是蒙福.
利未人還沒有上來之前.
他們身為猶太人就很痛苦.
因為被人欺凌.
但是現在他們的抬頭.
不是因為有強大的軍隊.
不是因為現在他們能夠獨立.
脫離波斯王.
不是.
而是他們感受到.
自己是神選民的福氣.

$^{2121}$那就很重要.
很抽象.
仍然都是那個人.
仍然都是那個地方.
仍然都是那個地圖.
不會特別下多些雨.
豐收十倍.
但是當你自己感覺到.
你自己是神選民.
這個榮耀的時候.
你做這些事情上來.
都會充滿了智慧.
充滿了平安.
充滿了愛.
充滿了對別人的同理心.
充滿了自己內心的盼望.
我相信做事.
就會更加信得益彰.
還有.
盡快將自己的事業.
是神的事業.
自己的公司.
是神的公司.
我賺大錢是為神的.
那很不同的了.
你整個心態.
做事的法則.
都不同的.
所以.
教會是為神做的.
不是為我的野心.
我的事業.
我要.
每個星期要坐滿三百人.
沒有三百人我都不收貨.
那是你的野心.
有時那個時代.
可能只剩下一百人.
不要緊的.
你忠心就行了.

$^{2161}$神有他自己的工作.
我們要做的就是.
勉勵弟兄姊妹.
感覺自己是神選民的福氣.
你有.
你自然就會充滿了歡樂.
所以第十二章.
很美麗.
很多的歡聲.
達到遠處.
不只是周圍.
周邊.
是向上的.
神都聽到這些歡樂的聲音.
很盼望.
我們來到禮拜堂.
大家唱詩讚美.
記住.
不只是去到天花板.
是上達於天的.
所以.
記得來到禮拜堂.
你不要想起誰誰.
你就不開心.
你就不敬拜.
千萬不要.
反而就是.
神令你想起.
盡快在祈禱裡.
甚至有具體的行動裡.
去和好.
你就能夠.
真的帶來整個新心靈.
有一個很大的釋放.
我們今晚就講到這裡.
我們一起祈禱.
透過第十二章.
讓我們見到.
在一個復興之後.
黎希米.

$^{2201}$他這麼多年的心血.
帶來的.
不是他想建立他的野心.
他的事業.
這麼大的城牆.
他一個人可以.
命令那些人去做.
就做好.
他不是秦始皇.
他是一個.
憤卿家.
他是一個.
鼓勵百姓.
感受到自己是神選民的.
這個福氣.
這個身份.
有一位.
這麼敬虔的領袖.
在這個被擄歸回的時代.
出現.
是一個.
走正.
我們以為一定是很壞的人.
在外邦的波斯帝國.
好像是.
喪權辱國那樣.
貪慕虛榮.
但我們發現黎希米.
是一個這麼敬虔.
倚靠神的人.
他能夠感動百姓.
能夠一步一步.
建造好城牆之後.
重要的是.
百姓的心靈.
再無圍牆.
百姓的心.
是充滿了對神.
盼望主.
繼續保守我們.

$^{2241}$下個禮拜.
第十三章.
是很諷刺的.
幫助我們能夠明白.
裡面的經文.
我們這樣仰望禱告.
奉耶穌基督的名.
阿們.
\newpage



\section{路加福音 19:1-10}
\label{sec:oxDrdn_rwTg}
\textbf{2023-07-08 晚堂《奇妙的改變》路加福音 19\_1-10 講員:陳麗蟬牧師}
\newline
\newline
連結: \href{https://youtube.com/watch?v=oxDrdn-rwTg}{\texttt{https://youtube.com/watch?v=oxDrdn-rwTg}} ~~~~ 語音日期: 2023-07-11
\newline
\newline
\hyperref[sec:LQfxIAo88QM]{\small{< < < PREV SERMON < < <}}
~
\hyperref[sec:index_chronic]{\small{[返順時目]}}
~
\hyperref[sec:index_scriptual]{\small{[返順卷目]}}
~
\hyperref[sec:87awTrKJFrY]{\small{> > > NEXT SERMON > > >}}
\newline
\newline
路加福音 19:1-10
\newline
\begin{longtable}{cl}
\hline
\hline
章節 & 經文 (和合本修訂版)\\
\hline
19:1 & \begin{tabularx}{0.7\textwidth}{X} 耶穌進了耶利哥，要從那裡經過。 \end{tabularx} \\ \\ \relax
19:2 & \begin{tabularx}{0.7\textwidth}{X} 有一個人名叫撒該，作稅吏長，是個財主。 \end{tabularx} \\ \\ \relax
19:3 & \begin{tabularx}{0.7\textwidth}{X} 他要看看耶穌是怎樣的人，只因人多，他的身材又矮，所以看不見。 \end{tabularx} \\ \\ \relax
19:4 & \begin{tabularx}{0.7\textwidth}{X} 於是他跑到前頭，爬上桑樹，要看耶穌，因為耶穌要從那裡經過。 \end{tabularx} \\ \\ \relax
19:5 & \begin{tabularx}{0.7\textwidth}{X} 耶穌到了那裡，抬頭一看，對他說：「撒該，快下來！今天我必須住在你家裡。」 \end{tabularx} \\ \\ \relax
19:6 & \begin{tabularx}{0.7\textwidth}{X} 他就急忙下來，歡歡喜喜地接待耶穌。 \end{tabularx} \\ \\ \relax
19:7 & \begin{tabularx}{0.7\textwidth}{X} 眾人看見，都私下議論說：「他竟然到罪人家裡去住宿。」 \end{tabularx} \\ \\ \relax
19:8 & \begin{tabularx}{0.7\textwidth}{X} 撒該站著對主說：「主啊，我把所有的一半給窮人；我若勒索了誰，就還他四倍。」 \end{tabularx} \\ \\ \relax
19:9 & \begin{tabularx}{0.7\textwidth}{X} 耶穌對他說：「今天救恩到了這家，因為他也是亞伯拉罕的子孫。 \end{tabularx} \\ \\ \relax
19:10 & \begin{tabularx}{0.7\textwidth}{X} 人子來是要尋找和拯救失喪的人。」 \end{tabularx} \\ \\
[1ex]
\hline
\hline
\end{longtable}
$^{1}$請姐妹平安.
我們看看這幅圖像.
有一隻黑色杯.
但這不是一隻普通的杯.
因為加水之後.
它可以變出一個彩色繽紛的.
滿佈星球的畫面.
但不是加任何水都可以的.
加冷水就不會變了.
唯有加熱水的時候.
它才會變的.
有不同的杯就變出不同的圖案.
當上帝創造人類的時候.
他就按照他的形象去做人.
所以每一個人應該很有朝氣.
很有力量,很有光彩的.
亦都人與人之間相處是和平,和諧的.
但當人犯了罪之後.
生命就變得烏煙瘴氣.
但怎樣能夠回到上帝創造我們的時候的情況呢?.
只有我們回到上帝的面前.
上帝工作的時候.
會恢復我們原本生命的光彩,和平.
但我們要找上帝.
也不是隨便隨地找一個不知道什麼的上帝來填補.
我們一定需要找到那位唯一的,獨一的上帝.
就是耶穌基督.
祂才能夠化解我們裡面的污濁的生命.
我們就用撒該這一段經文來彼此勉勵.
耶穌基督進入耶利哥城.
然後進入耶利哥城.
然後祂就會出去耶路撒冷.
就開始踏上十字架的道路.
祂故意進入耶利哥城.
其實祂有些事情要做的.
祂就要去找撒該.
幫助撒該的.
我們認識一下這個撒該.
撒該這個名字就是聖潔.
撒該這個名字是聖潔.

$^{41}$不過他的生命就一點都不乾淨.
因為他所賺的錢全部都是骯髒的.
因為他是一個稅吏長.
為什麼一個稅吏長說他賺的錢是骯髒的呢?.
因為當時羅馬政府統治猶太地的時候.
就採取一種形式叫做「包稅制」.
就是將他管治的地方分成一區一區的.
這個區裡面哪個可以負責做收稅呢?.
他就不是去因財來找人的.
而是他就將這個位置來做拍賣.
價高者得.
得到出得最高錢的人.
他就可以做稅吏長了.
稅吏長拿了這個位置之後.
他就沒有人工了.
那他怎樣可以賺取到金錢呢?.
這個羅馬政府有規定.
每一年裡面每一個人要收多少稅.
假設收100元.
但這個稅吏長在他自己管轄的地方.
他可以在100元以上去加的.
你喜歡加到每人收300元也行.
900元也行 2000元也行.
但只要你是每年將那個人的100元.
就交了給羅馬政府.
那其他剩下來的稅吏長就可以代代平安的.
薩戈在耶利哥城裡面做這個稅吏長.
經文特別說他是一個財主.
在聖經裡面很多時候說財主都是一些負面的.
說他是不好的.
但現實生活裡面不是每個人有錢都是不好的.
不過在聖經裡面是專門說到某一些人.
他賺的錢是一些不是很好的途徑.
其實這個薩戈剝削了很多猶太同胞的錢.
所以這個薩戈雖然是有權有勢.
他又有這個地位.
而且羅馬政府會給一些兵他們的.
跟著他來保護他們的.
所以說他真的是一個很威風的人物.
但他一點都不快樂.

$^{81}$而且他因為太過討人厭了.
沒有人喜歡跟他做朋友 他沒有朋友的.
特別那些猶太人對他恨之入骨.
因為覺得說他是羅馬人的走狗.
是賣國賊 是罪人.
因為薩戈在耶利哥城那裡完全不被接納.
再多一重慘就是薩戈是一個矮子.
所以去到哪裡他都要把腳挾高去看東西.
看到他整個無論外在還是裡面.
他都是一個很慘的人.
他外表很風光 有錢 有權勢 有地位.
他這個地位又高人一截.
但是在裡面他是孤單的 是寂寞的.
是不感到生命有什麼好的東西的.
有些心理學者就說.
每一個人在裡面都有一個小孩.
他叫Inner Child.
每一個人裡面都有一個小孩.
但當人在這個童年 或者當他一直長大的時候.
如果他經歷一些不如意的事情.
或者當他一直長大的時候.
被人看不起 或者被人排斥 被人討厭的時候.
他整個生命就是一個創傷的生命.
他一定要去看怎麼樣去被醫治.
這些心理學者就說.
人又要學習接受自己每一個人.
在裡面或多或少都會有一個不完全的小孩.
一個曾經受傷的小孩.
但是我們以信徒來說.
我們需要 唯有耶穌基督在我們裡面的時候.
才可以將我們這個受傷的小孩去醫治他.
再一次讓這個受傷的小孩.
可以能夠得到生命.
我們知道這個薩蓋裡面.
他做了這個瑞利獎之後.
一直都被人排斥 被人拒絕.
被人罵他是賣國賊.
再加上一個男子漢在那裡.
他真的很自卑.
他需要找耶穌 需要找耶穌來幫他.

$^{121}$在這個曾經有一段很長的日子.
不知道有多少年了 我都不知道.
我發覺到我自己每一天起床的時候.
我就會有一種莫名其妙的憂傷.
是很難過 很想哭的.
這種情緒有時候會很短.
做完事 吃完早餐就沒有了.
但是有時候又會很長.
長到當我出門上班的時候又不想上班.
又想著 糟了 回去一會兒可能要帶茶經.
或者要開會.
我覺得我很沒有力量 我很不想去做.
這些情緒 我又不知道為什麼.
但是我覺得有時候很辛苦.
這兩年在杜豐山讀一個靈修寺的課程的時候.
每一年我都要見八次我們的導師.
有一次我就和導師說.
我有這樣的經歷 其實困擾了我好幾年.
其實很影響我的情緒和工作.
他就為我祈禱.
當他為我祈禱的時候.
我就想起我曾經在讀五年級的時候.
我一個的情景.
當時在家裡 我哥哥在那裡.
我姐姐在那裡 我的弟弟都在那裡.
我小時候告訴你 我很容易哭的.
所以我哥哥就和我的姐姐的哥哥說.
他說他是精一大罐的.
我只要數一二三 他就會哭了.
我覺得豈有此理 這樣說我.
不過當我哥哥一數一二三.
我接著就哭了一聲.
我覺得他在欺負我.
為什麼無端端說我是哭鬼呢.
我在那個情景裡面.
我看到哥哥在那裡大笑.
因為成功了 他說我是哭鬼.
他說一二三我就哭了.
他在那裡大笑.
我看到我的姐姐和弟弟.

$^{161}$又一起在那裡笑 哭鬼來的.
我看到我自己當時覺得我很醜.
為什麼哥哥說一兩句.
我就哭了 我覺得自己很沒用.
我覺得我很生氣.
為什麼弟弟和姐姐都在那裡一起笑我.
我又覺得我很無助.
為什麼沒有人幫我呢.
沒有人出來幫我罵一下我哥哥呢.
接著我的老師就叫我.
你看看 看耶穌在不在.
我安靜了一會之後.
我又感覺在那個五年級的我那裡.
我看到我伏在飯桌上哭.
看到耶穌站在我旁邊.
祂很高的.
祂就拍我的頭.
祂就好像在跟我說.
傻丫頭來的.
很小的事來的.
如果人告訴我很小的事.
我就可能罵他.
你又覺得很小的事.
我又不知道為什麼.
當耶穌這樣跟我說很小的事的時候.
我又覺得 也不是很大件的事.
接著我又看了.
其實那些姐姐 那些弟弟.
他們沒有人在笑我.
他們只是用一些很不值得的眼神.
看我哥哥 為什麼又哭了.
然後他們個個都在自己做功課.
我突然間去想到.
感覺到 原來其實.
我的哥哥 姐姐 弟弟他們支持我的.
不過免得跟我哥哥吵架.
我又覺得 耶穌提醒我一件事.
這麼小的事情 做什麼 生氣這麼久了.
原來 我不知道怎樣.
是將那天的場景.

$^{201}$我覺得我很委屈.
我覺得我被人欺負這種感覺.
就進入了我的腦 進入了我的心.
所以我不知道怎樣一直影響著.
我在早上起來的時候.
我就覺得不開心.
我就覺得很無力.
我覺得很冤屈.
但不知道怎麼說.
但在那次當我感覺到.
耶穌站在我旁邊.
祂又拍我的頭的時候.
我覺得祂在支持我.
再去說很小的事的時候.
我就好像一個個提醒.
我自己將一件很小的事情.
就化成很大的事情.
影響到我一直的長大.
我覺得有什麼事.
是令到我自己這麼委屈呢.
這麼小的事.
我就在上帝的面前.
求祂幫助我.
從這些自己覺得很委屈.
很冤屈的感受.
跳出來.
很奇妙 經過那次之後.
我每一次醒來.
那種莫名其妙的憂傷.
就沒有了.
所以現在過了一年多.
我記得只是出現過兩三次.
但很快就消失了.
我相信那次那位牧師.
其實他就在幫我做一個心靈的醫治.
無論那個感覺.
其實一切都是感覺.
是否當日我五年級的真實場景.
我也不知道.
不過他藉著他的祈禱.

$^{241}$藉著耶穌在我旁邊拍我的頭的時候.
他將我裡面曾經受傷的小孩醫治了.
弟兄姊妹.
在信仰的歷程裡面.
我們有沒有經歷過.
在耶穌裡面.
我們經歷過釋放那份力量.
帶來我們整個生命的興盛呢?.
又或者我們有沒有經過.
聖經裡面說的.
聖靈在每一個信主的人裡面.
可以活水江河的湧流.
很有力的.
也很平靜的.
我們有沒有經過.
在生命當中.
那份清澈的河流.
在我們裡面的湧流.
帶給我們的力量.
帶給我們的平靜呢?.
又或者.
在我們的信仰裡面.
我們是不是覺得.
好像停滯不前.
好像吃雞肋那樣.
吃那個雞翼尖那樣.
上教會或者不上教會.
都好像沒有什麼分別那樣.
但如果我們在這樣的情況下.
好像信仰對我們來說.
都沒有什麼幫助.
習慣了不上又覺得有點可惜的時候.
弟兄姊妹.
就請你快點去找傳道牧師來聊聊天.
只要我們肯去祈禱.
上帝一定會在我們生命裡面.
有奇妙的帶領的.
又說回這個薩埃.
這個薩埃.
他很希望見到耶穌.

$^{281}$因為他覺得他自己有地位.
又有那麼多錢.
但是他沒有朋友.
但為什麼這個竇木仔的兒子耶穌.
他竟然去到哪裡.
都有大群人這樣跟著他.
所以他就很想看看這個耶穌.
究竟有什麼吸引力.
但我都說過.
他很討厭黑人.
那麼壞.
那麼壞可能那些人又迫不及待不讓他看.
又可能是故意撞他.
但是因為他對耶穌這一份盼望.
所以他就不理人.
不理人對他有什麼歧視的眼光.
不理他只是這個瑞麗長的身份.
當時這個瑞麗長他們有官服的.
很漂亮的.
比其他平民百姓的衣服漂亮很多.
他也想到辦法壓迫不開那些人.
他就爬上這個桑樹.
他也不介意爬上這個桑樹的時候.
他可能會掉下來.
會掉死的.
他什麼都不理.
他就是出於一個對見耶穌這一份的盼望.
我們這個信仰裡面.
我們不是說我們每一個星期六.
我們就回來崇拜.
每一天我們吃飯.
早餐午餐晚餐吃飯的時候.
記得這個祈禱.
做足這些禮行公事就行了.
我們這個信仰更加需要的.
就是有一個盼望.
盼望我們向上帝禱告的裡面.
祂會來去成就.
盼望在將來這個天國的新天新地.
盼望在天國裡面那份的喜樂.

$^{321}$我們盼望我們能夠去得到.
所以他對耶穌基督的盼望.
就能夠成為了我們今天.
就算活在這個北風巨浪.
或者任何的艱難裡面.
我們有力去面對這一份基礎的力量.
當撒該他很渴望見到耶穌.
耶穌就和他一起.
耶穌和他一起的時候.
就恢復回上帝創造他本來應該有那份的朝氣.
應該有那份的光彩.
應該有那份和人之間的和諧.
以致他不再被這個權勢.
再被這個金錢所捆綁.
去得到釋放.
去得到這個自由.
那麼耶穌進入耶利哥城.
他一見就見到這個撒該.
其實我之前沒有見過他.
不過因為這個瑞利獎.
他有這個官服.
這個官服很漂亮的.
他爬上樹上看到閃亮亮的衣服.
加上這個耶利哥城的瑞利獎.
出名是矮仔.
所以一看就知道這個就是瑞利獎撒該.
他就去到這個桑樹前對撒該說.
撒該快點下來.
今天我要去你家裡住.
其實第一次見人.
如果我第一次見你.
我就說 肅閒 今天請我喝茶.
沒理由這麼沒禮貌的.
肅閒 今天可不可以請我喝茶.
很禮貌的.
但是耶穌是直接了當的說.
今天我必須住在你的家裡.
必須這個字.
表明他必須一定要做那一樣東西.
他一定要做某一樣東西.

$^{361}$就是做什麼呢.
就是將罪人拯救.
在路加方第五章裡面.
當有一天耶穌和門徒和瑞利在吃飯.
法利塞人就去和耶穌的門徒.
就在那裡說耶穌是非.
為什麼你們竟然去和瑞利和這些罪人.
一起在那裡吃飯呢.
耶穌聽到的時候.
耶穌就和他們說.
健康的人不需要醫生.
有病的人才需要去看醫生.
我來到這個世界上.
我不是去幫助那些義人.
我來的是要叫那些有罪的人.
去迴轉悔改.
所以剛才我說.
當耶穌基督將要走上耶路撒冷.
將要面對釘十字架的路的時候.
他特意繞一繞就進入了耶利哥城.
他要特意找一找被人排斥.
被人定了他是大罪人的這個殺戒.
難道這個殺戒一直都被人說是罪人.
是賣國賊來的嗎.
所以那些人要接受這個.
他們一直都覺得不好的人的時候.
是很難的.
很多懷疑的.
所以當耶穌說.
他要去殺戒那裡的家裡去住的時候.
這幫群眾就竊竊私語.
在那裡就說.
這個所謂的夫子老師這麼受歡迎.
他究竟知不知道.
他要去住這個稅利長的家裡.
是甚麼人來的.
是壞人來的.
是大罪人來的.
這個夫子又說自己是先知.
又說自己是怎樣的聖潔.

$^{401}$他究竟有甚麼動機.
是否想去分一份.
當人無端端說我要去他家裡.
這個人當然是跟這個人有關係.
起碼你去認同對方.
覺得可以做朋友的.
或者大家可以彼此去聊天.
為甚麼耶穌要跟這個壞人去住.
去那裡跟他聊天呢.
耶穌這樣做.
他必須要這樣做.
表明他要去尋找,拯救.
那些心靈受創傷的人.
所以當他進入耶利哥城.
那麼多人跟著.
他甚麼人都不看.
只看在桑樹上的薩蓋.
他要幫助他.
那群群眾都不知道發生甚麼事的時候.
其實耶穌這樣做.
有時我們都不明白耶穌這樣做.
不過耶穌所做的一切.
他的意念是高過我們所有人的意念.
他一定有他的目的.
耶穌很快就去了薩蓋的家.
經文就說.
薩蓋就改過自新.
不知道住了多久.
聖經沒有說.
不過是一陣子.
不知道有沒有一天.
但是在這個路加方.
在這段第十九章之前.
耶穌已經說了.
在人所不能的事.
在上帝是可能的.
是可以的.
所以上帝早就說了.
耶穌早就說了.
你們以為有甚麼是不可能的.

$^{441}$上帝可以.
所以弟兄姊妹.
有時我們覺得有些事.
不可能發生的.
但上帝可以.
耶穌去了薩蓋的家的時候.
薩蓋就站著.
他站著對著主耶穌說.
主啊 我將會將我所有的錢.
分一半給窮人.
然後我過往曾經勒索過那些人.
我會四倍還給他們.
聖經故意說.
薩蓋站著對耶穌說.
站著對耶穌說.
本來可能只是坐著聊天.
站著對耶穌說.
他有一個下定了決心.
所以有時你坐著吃飯.
你突然站起來.
我決定這樣做了.
薩蓋我決定這樣做了.
他要在眾人面前.
有決心的宣佈.
他要將他過往的錢財的一半.
賺到的金錢一半.
分給窮人.
剩下一半.
過往勒索過那些人.
我四倍按著出埃及記的指示.
勒索過那些人.
我四倍還給他們.
可以想到薩蓋過往.
刮了人家多少錢.
一半給窮人.
然後四倍還給曾經勒索過的人.
這個薩蓋的改變.
他的悔改是真實的.
是需要有行動表示他的悔改.
才能看到他是真心的.

$^{481}$在路加福音第十八章裡面.
講到有一個少年官.
認為很有錢的.
又有地位的.
他學識也很高的.
他也說在十誡上的要求.
他全部做足了.
但當耶穌說.
你這麼有錢.
你將你所有的錢財.
那些財物.
你去分給窮人的時候.
他聽到就.
悠悠愁愁的走了.
他不捨得.
他不捨得自己的錢財.
所以他找不到這個永生.
當你看到.
這個原本很貪錢的薩蓋.
他本來是賺到盡的人.
他竟然到現在.
是分錢給窮人.
所以我想整個耶利哥城的人.
不知道多麼高興.
所以我們政府派那六千元.
我們不知道多麼高興.
他現在還良心發現.
過往勒索過的人.
四倍償還.
耶穌說過.
在第十八章路加方說.
耶穌說過.
財主得救是極度的困難.
所以那個少年官.
他就得不到這個永生.
但是耶穌又說過.
他說在上帝裡面.
沒有東西是不可能的.
只要他說可以的時候.
就一定可以.

$^{521}$這個原本那麼貪錢的薩蓋.
當他遇到耶穌的時候.
整個生命就有很大的改變.
耶穌在最後一節說.
人子來.
是要尋找拯救失喪的人.
這個失喪這個字.
原本文字的意思就是.
擺錯了位置.
我們的眼鏡.
要擺在鼻樑上.
眼睛的前面.
我們才可以看清楚.
但有時我們又知道.
我們喜歡將眼鏡擺在上方.
就當頭盔用.
擺錯了位置的時候.
看東西就不清楚了.
眼鏡也不是一個眼鏡.
耶穌來到這個世界上.
目的是要做什麼呢.
就是要拯救.
拯救這個字原本的意思就是.
將它放回正確的位置.
所以當這個眼鏡.
放回在鼻樑前面.
在眼睛的前面.
我看東西就清晰了.
蒙查查.
人耶穌基督來到這個世界上.
將我們在生命裡面.
有什麼走歪了路.
殺該就走了一個貪財的路.
以致生命不光彩.
沒有朋友.
自卑.
但當他找到耶穌的時候.
他就很慷慨.
很慷慨.
在被耶穌的接納裡面.

$^{561}$他發覺原來人間有溫暖.
人間有安全.
所以找到耶穌的時候.
他就能夠恢復上帝原本創造他人的本性.
他和人之間應該可以和諧地相處.
是可以彼此相愛.
因為他經歷了耶穌基督.
對他的接納的恩典.
他就很樂意將他的錢和其他人分享.
其實他損失非常慘重.
但他甘心.
他樂意.
他喜樂.
耶穌到這個世上來.
他要將一些走錯了人生路的人.
幫助我們行回正軌.
我信了耶穌.
但有時在我們生命裡面的小孩.
有些小孩仍然在搞擾我們.
我們覺得就算我們信了耶穌.
我們也沒什麼特別.
我們仍然的生命都好像沒什麼活力.
沒什麼朝氣.
和人之間的相似仍然處處都不和諧.
可能就是裡面一個小孩.
曾經的不好經歷.
或者一些錯誤的思想在搞擾我們.
所以回到耶穌基督的裡面.
他可以幫我們走回人生的正路.
我們可以活得喜樂.
活得自由.
活得和諧.
活得豐盛.
我們唱一首詩歌.
是生命聖詩的.
不過你在這裡多看一會兒就行了.
我生命.
自從有耶穌住在我們生命裡面的時候.
我們就有那份喜樂不斷湧出來.
謝謝.

$^{601}$我們一起來唱.
我生命已經有奇妙的改變.
至耶穌來住在我心.
我渴望的燈光今朝要我心間.
至耶穌來住在我心.
彼羅朝人和聯誼觀魂至開頭.
至耶穌來住在我心.
我們再唱一次.
我生命已經有奇妙的改變.
至耶穌來住在我心.
我渴望的燈光今朝要我心間.
至耶穌來住在我心.
彼羅朝人和聯誼觀魂至開頭.
至耶穌來住在我心.
我們一起禱告.
天父神弟 我們謝謝你.
因為你差派你的獨生兒子來到世界.
是為了你去拯救我們這一班的人類.
將我們走錯了生命路的人.
能夠走回正軌.
幫助我們.
我們在座當中.
我們每一個.
我們生命或多或少.
都有一些走錯路的小孩.
在當中令到我們.
生命好像很阻滯.
好像沒有光彩,沒有力量.
雖然我們信了耶穌.
但我們仍然覺得這個信仰好像可有可無.
主耶穌基督求你幫助我們.
以至我們每一個人.
我們都勇敢的去面對我們生命裡面.
一些缺口.
一些不是很光彩的東西.
在我們知道只要我們肯回到耶穌基督面前的時候.
耶穌基督一定有這個能力幫助到我們.
求聖靈幫助我們.
以至我們有勇氣去面對.
以至我們能夠真真正正的經歷到.

$^{641}$在耶穌基督裡面.
生命的豐盛.
生命的光彩.
生命的力量.
榮耀我們的天父上帝.
多謝主耶穌上天禱告.
是奉靠主基督的聖名.
阿們.
\newpage



\section{提摩太前書詩篇羅馬書 1:19-20}
\label{sec:87awTrKJFrY}
\textbf{2023-07-09 加恩樓《平衡空間的人生》羅馬書 1\_19-20,提摩太前書 2\_5-6,詩篇 27\_13-14 講員:李盛林牧師}
\newline
\newline
連結: \href{https://youtube.com/watch?v=87awTrKJFrY}{\texttt{https://youtube.com/watch?v=87awTrKJFrY}} ~~~~ 語音日期: 2023-07-11
\newline
\newline
\hyperref[sec:oxDrdn_rwTg]{\small{< < < PREV SERMON < < <}}
~
\hyperref[sec:index_chronic]{\small{[返順時目]}}
~
\hyperref[sec:index_scriptual]{\small{[返順卷目]}}
~
\hyperref[sec:QuJcmG3w_7k]{\small{> > > NEXT SERMON > > >}}
\newline
\newline
$^{1}$各位基督徒的朋友,我跟阿Mo說是有心人,有些很關心他,有些重建教會,有些有其他信仰.
我相信在去年七月,有兩個星期,剛剛一個年頭,在紅館當日有萬多人,萬多雙眼見證和看到那個場景.
有時我去教會分享為了港道,有些朋友特意知道我去教會分享,上來跟我說,牧師,我當晚在紅館看到了.
有些到今天,心理關口和陰影還未過到,特別是我一直有接觸那班舞蹈員,跟牧師一起跳舞.
他們是有些到今天一提到這些事的時候,在他們的工作室,因為我定期跟他們一起坐.
他們在觀塘比較多的工作室,在那裡找個地方坐在一起.
有些一談到這個題目和當場的情況,有些受不了,這些是Tensor.
有時我回想,台上12個人加24個,包括我的小兒子阿Mo,除了阿Mo,即是23個,即是加12個,35個人.
如果其中一個在當晚,不是阿Mo,是其中一個人,他是被螢幕打中的,他會怎樣?他的家人會怎樣?.
我很難想像,因為我知道他們沒有一個信耶穌,有時我回想,如果我也不信耶穌,我怎樣去面對一樣所謂的創傷和災難?.
如果我的小兒子或我家人都不是信耶穌,或是沒有信仰的,他會怎樣去面對這一年?.
他頭腦清醒,說話清晰,邏輯我相信那個意識可能表達得比很多人好,但是很遺憾,四肢不能動彈.
有時我會嘗試去體驗一下他的情況,我睡在床上,我想像我四肢不能動彈,我會怎樣?.
過了15分鐘,我開始忍不住,想動,但是我自己跟自己說,我要去體驗一下阿Mo的情況,再堅持多10分鐘,已經很辛苦了.
30分鐘已經忍不住要動,但是他24小時是這樣,特別是頭三個多月,我在《代禱信》也說過,他唯一的視線就是一個很冰冷的天花板.
就是那個ICU的天花板,當時他仍然要用一個頸箍,動不了,所以他一醒,睜開眼睛,就是那個天花板.
雖然他身旁有一些玻璃窗,他沒有辦法像我們這樣去看外面的天空,看外面的景色,他唯一的就是這樣.
他能夠爭取看外面的,就是幫他轉身的時候,稍微轉一轉,他可以看到,但是那個不是他主動,那個是他被動去看.
昨晚我就發了第50封的《台禱信》,我不知道你有沒有留意一些,我相信當中有一些弟兄姊妹,一直都有留意,特別我熟悉一些以前的弟兄姊妹在這裡.
他們跟我說他們一直都有,50封台禱信,不容易寫的,50個星期不容易過的.
前幾天,阿Mo他在WhatsApp,他可以用Siri去講電話,用Siri去溝通WhatsApp.
因為我們定期在他原先住的單位打掃,剛剛打掃完我就拍了一個影片給他,他看完之後,他在WhatsApp回應我,I miss home.
我很記掛我的家,所以在台禱信我解釋他一個心情,他的記掛,他的懷念,是因為那個地方能夠讓他回想到昔日的經歷.
其實當他出事的時候,很多他的同事跟我分享,他說阿Mo在那個時候是最頂尖的舞蹈員,在他們當中是最頂尖的.
所以當他寫I miss home的時候,我很難想像他回到自己熟悉的家,家裡是一片鏡子,他常常在家裡練舞,一勒頸.
他想到他跟女朋友一起練舞,他想到跟家人,跟朋友過去的一些景象,他更加I miss home.
他想回到家裡,但是只可惜他回不了,一年,在這一年裡,我也很多在信仰上反省,特別我是信了耶穌超過半世紀.
我在教會和教神學上是超過四十年,信仰對我來說是什麼呢?我教很多的學生遍佈全世界.
所以當我發台獨信給他們的時候,無論中國大陸,北美,歐洲,澳洲,紐西蘭,非洲,南美,全部都有我的學生,他們就再轉發給他們的群組.
有些在香港的朋友,他說牧師,我收你最新的台獨信,不是在香港收的,是中國大陸發給我的.
在這個特別艱難的一年,我再重新思想我的信仰是什麼呢?.
很多時候我們和別人分享我們信耶穌,我們得到的回應就是,信仰都是導人向善,信什麼都不要緊.
但是信仰不只是導人向善,如果信仰單單是導人向善的時候,你信什麼都不要緊.
早上我通常在我屋苑游泳,偶爾都聽到泳友在自己聊天.
信仰就像一棵樹,你信什麼宗教都不要緊,都是一棵樹,只是不同的枝,前幾天游泳聽到他們.
都不想走過去搏,通常游泳跳下去,不理會周圍有人走,就這樣游.
信仰如果是同一棵樹,周圍的不同枝節就是不同的信仰的話.
這棵樹的生命,如果只是叫人導人向善就是了.
但是信仰不只是道德的信仰,導人向善就是你要做好人,從道德的角度來看.
我們做人要自己有良心,自己有交代,自己對得起別人.

$^{41}$這是放諸四海而皆準,你信任何一個信仰都好,你去信吧.
但是信仰不只是叫你做好人,信仰不只是頭腦上的,我教了神學四十年.
信仰不只是你頭腦能夠分析到這位神.
信仰更加不是一本經書的聖經,你在那裡去研讀.
又或者是一些佛經,你在那裡去念.
甚至現在不同的信仰都有所謂的查經班,去查他們的經.
甚至是一些現代的音樂,配合他們的信仰,這樣來去.
但這些只不過是表面,信仰在這個世界,在歷史裡是唯一一個信仰是講關係.
就是我們所信的基督教和耶穌,沒有一個信仰跟你講關係的,沒有.
但是這個信仰是跟我們講到我們跟這位神的關係.
其他的信仰,你可以說,特別上一代,現在年紀一代比較少.
上一代就是我家裡有問題,我兒子出生的時候他有什麼不舒服.
又或者有什麼災難,搞來搞去都不行.
有些人就說,你不如契給王大仙,不如契給如來佛,不如契給哪個神哪個神.
當你契了他成為兒子的時候,他一定保護你的.
這個是我們華人在世界各地的一種觀念.
我們華人有一種文化叫疏通文化,什麼叫疏通文化呢?.
就是我用一些物質,用一些錢,在實際環境上高.
可能在一些官場上高,在一些手續上高,裝得不好,疏通疏通.
不單止是現況疏通的文化,對於一位神,對於信仰都是一種疏通文化.
就是我給你一些物質,我給你一些製品,疏通疏通.
你照顧一下我,你照顧一下我的家人,就是疏通.
但是我們這位神主動地和我們建立一種神和人的關係.
你不需要特意走過去,耶穌你契我吧,不需要.
他已經主動地用了一種關係,用了一個限制來到地上高.
用那一份親情去吸引我們,去建立那個關係.
前幾天有一個基督教機構,世界性的主管.
他給了一個訊息,他說牧師你幫我給一個祈禱文給阿Mo吧.
我都發了給他,那個祈禱文中間提到.
他說阿Mo,你要思想到耶穌好像你一樣.
或者你好像耶穌一樣,你頭腦清晰,你活動自如.
但是現在你是受制於你四肢不能動.
就好像耶穌來到地上高,他甘願用了我們這個人的形態出現.
他的神聖能力是受我們這個軀殼和環境所限制著.
阿Mo,你現在就是這樣.
他看完之後,因為我用訊息傳了給他,傳了給他.
他看完之後,他用了一個心回應我.
因為他深深體會到他的限制.
他深深體會到耶穌同樣是因為我們的緣故受限制.
他可以不需要肚餓,但是他肚餓.

$^{81}$他可以不需要有疲倦,但是他成為我們這個樣式的人.
他有疲倦是因為他要按照聖經說.
是他體驗我們.
他知道我們的限制,他知道我們的軟弱.
以致他能夠凡事去明白我們.
所以我們的信仰不是單求一個道德做好人的信仰.
我們的信仰更加不是只是一本經書的信仰.
很多時候我現在提到,在今天很多基督徒是沒有一個屬靈的經驗.
當你一說在實際生活裡面怎樣去經驗神的同在的時候.
很多時候被人來說,你迷信了.
很多時候我們在這個環境是說數據,你給我數據看吧.
你給我統計吧.
很多時候是用一個科學的角度去處理我們的信仰.
其實不同的一些學科是有不同的方法的.
你用科學的方法去讀歷史,你永遠讀不通歷史.
因為科學的方法是可以重複出現.
但是歷史你不能夠重複出現.
不要說歷史那麼長遠的時間,是說你的人生一樣.
你十歲的時候所發生的事情,在今天不能夠同樣發生.
你說不是一樣發生,但是你的心境,你的智商,你的經驗已經不同了.
聖經裡面很清楚地告訴我們.
原來在地上有不同的空間.
只不過我們人覺得眼見這個空間為實.
在我們讀書的時候,我們第一次去認識這個空間.
不是一個平面,而是一個立體的空間,就是三個層次.
三維的空間,就是XYZ.
再下去的時候,我們就看到原來有更多不同的空間.
今日科學家讓我們了解到,在他們研究的底下.
是有14度的空間,14Dimensions.
對我們這些平民百姓來說,真是不知是甚麼來的.
不過,原來不是我們單純眼見到的空間.
簡單來說,我們所能夠了解認識的.
同一個平行線,我們已經有不同的空間.
譬如我們現在在一個堂會,在我們家是另一個空間.
在我們家是另一個空間,可能發生了另一些事.
這只是我們現在平行的空間裡面.
但是平行的空間不單止是一個平線,可能還有不同的層數在裡面.
特別,我跟大家分享一兩個你也有的經驗.
你就會看到,當我們去殯儀館的時候.
現在分散了,不只是在紅磡.

$^{121}$有時你不能夠不出席,可能有親人等等.
有兩個地方,你是一個禁忌.
逼不得已,你也不會去的.
第一個地方就是靈堂背後擺放棺木和仙人的地方.
我們做牧師,我們也要去衝破這件事.
我們做牧師,要負責第一,第二個安息禮拜或喪禮的時候.
都要膽怯的,為何膽怯呢?.
如果按現象來說,我已經不能夠動了,為何你會膽怯呢?.
有時我看到尖仰圍庸的時候,有些人也不夠膽望仙人.
尖仰圍庸,有些人就不出去了.
有些人就算出去了,面也不是望的.
直線,走過就算了.
為何呢?我要問,你怕甚麼呢?.
他也不能動,他已經躺在那裡了,為何你會怕呢?.
第二個地方,你白癮成金才去的就是廁所.
殯儀館的廁所,你去過嗎?.
你去的時候,你一踏入廁所,你的感覺是怎樣呢?.
一進去的時候,你先看看有沒有人.
你先看看那塊鏡有沒有突然出現一些影子.
為何你會這樣呢?你怕甚麼呢?.
在這個情況下,可能內心已經有一種意識.
不單止是你眼見到的這些人.
可能特別在這些地方,不知道有沒有陰雲滿佈的東西出現.
為何會有呢?是因為不單止是一個現象的世界.
還有一個靈界的空間在這裡.
這個靈界的空間,十四道的空間裡的其中一個.
你怎樣面對呢?.
聖經裡的三段經文,其實這三段經文在這一年裡.
給我們很多的思想,特別是我的家人和阿Mo.
聖經裡提到,神的事情,人所能知道的,願顯明在人心裡.
因為神已經向他們顯明了.
自從做天地以來,神的永能和神聖是明明可知的.
天地背後原來有神的永能和神聖.
這是科學的數據無法解釋,是人的經驗無法了解的.
在發生事件的頭一星期,第一個月的時候.
那些醫生跟我們說,甚至跟阿Mo說.
他說,你在這個狀態,你沒有一年都需要最快半載,就是半年.
才能離開ICU.
這個是憑著他們的醫學經驗和見到一個這麼嚴重創傷的人.
其實他非常之嚴重,第四節骨,頸骨是爆了.

$^{161}$他要用這裡的骨挖出來,然後補上去,然後再用支架固定.
所以他現在在頸上有支架和螺絲頂著,鑲著.
經過一年,當然已經復合了,不過支架還在.
前面的治療方案未定,有機會拆支架做治療.
他現在頭可以動,可以拱.
但後面有支架不能拱在後面.
他現在肩膀可以動.
他現在三頭肌,三頭肌,肩膀,後面的腿是可以控制的.
這是那些醫生一開始憑藉著經驗數據是不會覺得他能動的.
他現在的二頭肌,本來是零,因為我們肌肉的指數正常是五.
中間有些醫生在,他現在由零去到1.5至2.
但仍然未夠力拉起,但他的反應是1.5至2.
雖然如此,但他的手一跌下來,他拿不起.
我們做很容易的,一拿就起.
所以他現在坐在輪椅,是靠著他的手肘在armrest平衡自己.
不會掉出來.
但他一抽筋的時候,一掉下去就糟糕了.
他不能平衡自己會掉出來.
所以我們幫他加了安全帶.
醫生說你沒有一年都要半載留在ICU.
但我們求神很迫切祈禱.
很多兄弟姐妹,很多有心人為他祈禱.
他兩個月可以離開ICU.
這個醫生都解釋不了.
一個這麼嚴重的創傷,他可以在他們的經驗下超速離開.
他要一個助口幫助他抽痰.
這個助口在他當時的情況下.
起碼要半年及更多時間才可以把氣喉全拔出來吸痰.
但他由助口做開了大概一個多兩個月已經可以不需要了.
等他復原.
這些又是在數據經驗上解釋不了.
丁字妹是關心阿Mo的人.
在她身上我們見到神的工作.
有時我問如果台上那35位其中一個.
他們被打中.
他們本身和他們的家人是怎樣.
我不敢想像,因為我嘗試和阿Mo嘗試.
如果我們沒有這個信仰,我們怎樣面對這個衝擊呢?.
很無助,甚至到一個地步崩潰.
頭兩三個月我直接在神面前.

$^{201}$我不想影響家人,不想影響周圍的人.
我自己在洗手間裡,我哭了.
我問神為何會這樣?.
神說我寧願這些事發生在我身上.
我不願意見到一個年輕人.
他好端端一秒前一秒後完全判若兩個人.
他的情況改寫了他的生命軌跡.
同樣改寫了我們家人的生命軌跡.
神為何要我們面對這些?.
有些人跟我說牧師你還敢斜視神?.
他沒有保護你的兒子.
我通常跟他說,因為在想體會一個信仰的過程.
我就回應,我說人為的錯誤導致他面對這個災難.
不是神給他的災難.
大家都知道什麼叫人為錯誤.
但是神是在保護他的.
因為當他跳舞的時候.
舞姿是頭部向上,剛剛好.
如果不是這樣,他站直跳的時候.
螢幕就直接打在他的頭上,即場爆炸.
是沒有命的.
但舞姿是這樣的時候,就剛剛好跌在他的頸上.
第四節,他當然馬上躺下.
開始時那些醫生都很掛心.
因為第四節其中一個去控制的器官就是橫隔膜.
他能不能呼吸呢?.
如果不能呼吸,他一輩子都要用呼吸機.
到現在他都是靠著橫隔膜呼吸.
因為他這些肌肉暫時動不了.
所以他咳不起來.
他不能受感染,不能受新冠或感冒的感染.
因為一感染他咳不起來,咳不起來的時候就很容易肺炎.
兩件事醫生害怕的,肺炎和尿道炎.
因為他一直到今天一年都要用尿道排尿.
用尿袋嘗試醫生幫他脫掉尿喉.
上個月是不成功.
怎樣幫他脫掉尿喉呢?.
就是嘗試脫掉尿袋,夾住尿喉.
嘗試讓他的尿不會自然流出.
在排尿的膀胱能夠儲尿的容量多一些.

$^{241}$上個月一,兩個月,一儲尿的時候.
他就全身發抖,冒汗,很不舒服,不行.
我們都迫切祈禱,這兩個星期他沒有這些反應.
他可以儲到五百,通常我們人到四百多就要吸尿.
就要去洗手間,他可以儲到五百.
然後他自己的感覺可以幫他脫掉尿喉.
讓他再流出來.
重複,重複,都是可以的.
所以再多訓練兩,三個星期.
到八月一日,鼻療科醫生會幫他再檢測.
然後可以把尿袋,尿喉脫掉.
讓他自己用另一些方法去排尿.
這些都是神的恩典.
有時我問信仰是否只是道德的,做好人的.
信仰是否只是叫我們有難的時候才來到神的面前.
我在這一年,我都一直去禱告神.
信仰不是只有一個我有難來到你面前.
或者叫我做好人,我和人際的關係如何.
我這個生活等等是如何被你祝福去過.
但是我的生命,我的身體你是如何去照顧的呢?.
在這一年的經歷.
看到阿Mo神的恩典在他身上.
所以聖經裡面提到.
自從做天地以來,神的永恆和神聖是明明可知的.
這個我特意給大家看的.
如果你是讀物理的時侯,你是看過光譜的覆蓋.
我們人眼可以看到周圍的東西.
我會在那個小的片段再拉下來給大家看,放大.
我們是那麼小的片段.
上方的那個疤,周圍的是我們看不到的東西.
包括了我們中間這裡有很多的電波.
我怎知道有電波?你的電話一響一震.
那些是電波的一個接收器.
我們看不到,但是很多在我們穿穿插插.
所以有些人說.
你看到大廈有很多這樣的天線樹起.
有很多輻射,電波影響你的身體的確是.
所以在外國住,通常電纜下方是不會有房子的.
因為磁場是影響你的人.
我們看不到的.

$^{281}$但是在一個現實的世界.
很多時我們都說眼見到我才相信.
眼見到我才相信.
但是聖經說眼見到背後原來還有神的永能和神聖在.
是我們可以知道.
是從那些受造之物就可以了解.
在事發第一個月.
有一班很有心的醫學教授.
特別是做復康的.
和一些物理治療的教授和我們談.
他說牧師其實在我們人腦那裡.
由出世到死亡中間有大部分是沒有用到的.
有機會這些沒有用到的腦細胞可以取代受損的細胞.
我們見到很多,但是他接著再說一句.
但是那個機制他們不懂得做.
怎樣將這些沒有用到的腦細胞轉化成為取代受損的神經呢?.
他說我們不知道,我們不懂.
唯一就是這位創造主.
所以今天治療阿Mo的團隊有兩個教授帶著全個團隊.
前天他才給我一個訊息.
他說牧師我們盡了力.
因為他是基督徒,非常有心的基督徒.
我們盡了力去醫他.
但是有些位置他們都不知道.
真是唯有仰望倚靠這位神.
一個醫學教授去抗醫他.
在治療阿Mo的那班在QE的醫生都有聯絡的,也是基督徒.
因為他是給我們的恩賜,他是給阿Mo的恩賜.
一個富有經驗的主管,基督徒.
他跟我們說這些話.
所以,弟兄姊妹,朋友.
天地背後原來還有神的永恆和神聖.
從這個現象的世界我們就可以了解到這些情況是怎樣的.
第二段的經文是給我們很多的鼓勵.
其實這些經文聖經裡面.
我一直記錄了有26段的經文.
有個影音的機構在加拿大.
我上個月剛回去有10天.
就一直幫他們的機構錄影了這26段.
當然不能全部錄,錄了4段的經文.

$^{321}$是怎樣幫助我們走過這個困難的路呢?.
這3段是其中26段裡面的一些經文.
第二段就是提摩太傳書.
因為只有一位神,在神和人中間只有一位宗寶.
乃是降世為人的耶穌,基督耶穌.
他寫字跡,作萬人的贖價.
到時候這個事變成顯明,什麼叫做宗寶呢?.
我用現今的一個大家明白的詞彙了解.
宗寶就等於我們現在的調解員或調解人.
以前在我們這個世界或香港沒有這樣的行業.
一有什麼問題就對簿公堂,大家找律師處理.
現在你看電視廣告,不要找律師,找調解員去幫你們解決糾紛.
耶穌就是我們和神之間的調解員.
我們有什麼問題?.
我們可以經過耶穌來到神的面前,這位創造主的面前.
很多人跟我們說,信耶穌都是那些軟弱的人.
都是那些面對不如意的事情的人去找的.
當你不如意的時候,你才正在你人生被迫到牆角.
你才正在想.
你有少許空間都不會想這些.
但你沒有空間的時候,被迫到牆角.
你才正在想,究竟我的人生是不是就是這樣呢?.
究竟有什麼能夠幫我去跨越面對這些困難呢?.
所以尋求信仰不是說你是失敗者.
你是無助,你是軟弱.
因為當你來到人生的絕望,被迫到牆角.
你才真真正正誠實地去面對你自己.
你未被迫到牆角的時候,你都不會真真正正去面對你自己.
你還是用很多其他東西為自己解釋.
用很多其他東西來看可不可以去到.
你同樣沒有把握.
但聖經裡面說,我們有一位調解員就是耶穌基督.
所以一開始的時候,我跟大家分享.
沒有一個信仰是跟你講關係.
唯有基督的信仰跟你講關係.
神和人之間的關係.
並且這個關係不是疏通文化底下的關係.
我給你一些東西,你就幫我吧.
他不需要你給任何東西.
他是甘願來到這個世界上.

$^{361}$為我們付出一切.
主動地和你和我去建立這個關係.
不需要像上一代那樣.
你不能夠去治好你自己.
面對很大的困難,不如你契了給神.
當你契了給神,你就是他的子女.
他不會不理你的.
信耶穌不需要這些.
信耶穌更加不需要疏通.
他會主動地願意去承擔你一切.
這是我信了耶穌,甚至我交神這麼多年.
特別是今年,我常常鼓勵阿Mo和他的女朋友.
有很多人問我,他的女朋友信了耶穌嗎?.
他的女朋友是信了耶穌的.
很多時候我們都是牽著手一起去祈禱.
一起去呼求這位耶穌.
其中詩篇27篇.
是我們在這一年裡很艱難走過的路.
一個對我們很大幫助的經文.
這段經文原本是以色列人冰林城下被圍困.
沒有辦法離開城中.
他們被缺水缺糧等待死亡.
但是在他們和神的祈禱裡.
最後有這兩字.
我若不信在活人之地得見耶和華的恩惠.
我就早已喪膽.
要等候耶和華,當喪膽堅固你的心.
我再說要等候耶和華.
那些醫護憑著統計.
憑著經驗,憑著醫學數據.
只是說它是這樣.
但是原來現象活人之地的背後.
是有耶和華的恩惠.
是有神的永恆和神聖在.
當我們能夠超越了活人之地現象的世界.
那一個先證給我們是不會喪膽的.
所以活人之地是講求統計.
是講求數字.
是講求數據.
是講求現象.

$^{401}$是講求經驗.
如果當我們的眼光只是放在這裡的時候.
我們就絕望.
正如我剛才和大家分享的.
如果我沒有了這個信仰點.
如果我沒有了這位神.
我怎樣?.
剛剛開始崇拜的時候.
我收到另一個舞蹈員給我的訊息.
因為我一直和他溝通.
你現在怎樣?.
那個舞蹈員是受傷.
除了我,有兩個舞蹈員受傷.
其中一個去百販廟.
他說現在沒有那麼痛了.
不過我是減少了出秀.
減少了教跳舞.
他未信耶穌.
他同樣面對很大的掙扎和艱難.
見到他們的時候.
我都會和他分享我和阿Mo.
我們家和阿Mo的信仰.
是怎樣鼓勵我們繼續下去.
給我們有能力.
這個信仰不只是在心靈上.
自己那種正能量去想.
而是聖經的說話.
可以用在我們現時的生活裡面.
是超越了經書的字句.
是超越了教你向好道德好的一面.
而是真正有血有肉的.
是你可以經歷下去的.
這是過去那麼多年的信仰.
不是說過去沒有經歷.
而是今天更加有血有肉.
去經歷聖經的真實.
所以沒有了耶和華的恩惠.
看不到原來在活人之地.
有耶和華,有神的永恆和神聖.
全部是喪膽的,很害怕.

$^{441}$所以我曾經有幾個星期.
我說我們的心情是擴散車.
因為阿Mo的治療是經過很多關.
經過甚至死刑的幽禁.
很絕望.
但唯有再一次來到神的面前.
上兩三個月.
有一個治療阿Mo的主管.
是物理治療的主管.
他跟我說.
他說:李牧師.
阿Mo現在的情況.
是我讀書都沒有讀過.
讀書都沒有讀過.
什麼意思呢?.
就是他讀書.
他的治療只可以將他帶到這個位置.
但阿Mo現在已經在這個位置以外.
過了這個位置.
所以他說.
接下來我不知道.
不懂的.
他唯有用他的處理和修正的方案.
再看看是怎樣.
從阿Mo現在的情況再去看.
他說已經超越了他讀書的範圍.
他讀書是活人的經驗和現象.
但現在阿Mo的處境.
是面對著神的永恆和神聖.
是神的恩惠.
所以在過程中.
我們得到的安慰.
我們坐在神面前.
捉住神的話.
我們壯膽下去.
所以到最後.
我先跳過這些.
因為不需要跟大家說這些.
這個白板寫下來.
是他女朋友寫的.

$^{481}$I can do all things in him.
that strengthens me.
就是弗勒比書.
我靠著那加給我力量的凡事都能做.
有一些經文.
他寫在白板上.
寫在紙上.
貼在他面前.
去鼓勵他.
要看回聖經.
捉住神.
這是其中一段較早前的經文.
所以朋友.
現象很真實.
但如果現象只有這個處境.
我們很多時候會有絕望.
信仰不是單單是有難才找到.
在耶穌基督的聖經.
信仰是將我們帶回到.
神做我們的時候的原意.
是因為人的反叛離開了神.
沒有了不能活在神的原意下.
但今天當我們回到神的時候.
我們可以重新再經歷.
這位創造主的能力.
阿Mo正在經歷.
我們作為父母也正在經歷.
很盼望你能夠打開你的心.
很多不信耶穌的.
他沒有信仰的.
他跟我說牧師.
我可不可以認識你的信仰.
為什麼你們在苦難中.
仍然有那種動力和堅持.
我跟他們分享就是因為聖經.
因為耶穌.
很盼望你能夠嘗試一下.
如果你沒有信仰.
你未認識耶穌.
如果你是有信仰.

$^{521}$但你覺得信仰是導人向善.
我請你繼續你自己的信仰.
但是信仰不只是經書的信仰.
信仰不只是導人向善的信仰.
信仰是關乎你的生命.
關乎你的有血有肉現實的人生.
我很盼望你嘗試去認識.
經歷這位耶穌.
請我們一起祈禱.
我邀請你和我一起來祈禱.
特別是你還沒有信仰.
或者你未認識耶穌.
你很想能夠經歷到.
跟你生命貼切.
跟你有血有肉的人生.
連繫起來的信仰.
你和我一起祈禱.
親愛的主.
面對著人生.
有時我不敢正面面對.
有時我也不知道如何處理.
或者解決方案是怎樣.
因為在一個痛苦的經歷裡.
耶穌幫助我們和阿母.
在我們中間的朋友.
同樣讓你知道面對的困難是怎樣.
願意你幫助我們.
願意你幫助我更認識原來.
在現象世界的背後.
有神秘的永恆和神聖.
這個不只是說.
這個真的能經歷到.
我盼望你打開你的心.
勇敢地面對你自己.
如果你有和我一起這樣祈禱.
我邀請你讓我知道.
我會和阿母說.
因為她知道我來和你們分享.
我會和她說你的經歷.
耶穌的信靠.

$^{561}$同樣有人去認同,感受.
和和你一起走.
如果有這樣的人.
我盼望你表示給我知道.
我一會兒就會去醫院.
我都會和他分享.
如果有請你在座位舉手.
讓我知道我看到有.
還有沒有這樣的朋友.
你可以放下.
你願意勇敢地去面對你自己.
信仰,信耶穌不是一個失敗者.
不是你有困難.
而是因為這種困難.
逼你去面對你自己.
最後還有沒有這樣的朋友.
你願意打開你自己去認識.
和接受這位耶穌去引領你的人生.
還有沒有請你.
可以在你座位舉手.
我一會兒會和阿母分享.
還有沒有.
或者一會兒.
你都可以和我們中間的一些朋友.
基督徒來分享.
你自己內心的掙扎.
請我們一起低頭祈禱.
親愛的主.
願意你在我們心靈裡面工作.
幫助我們打開我們的心.
能夠去讓你引領我們的前路.
幫助我們中間的基督徒.
真的能夠體會到.
我們一生在耶穌你的帶領下.
幫助我們每一天.
能夠有那一份真真正正的屬靈經驗.
以致我們的信仰.
不是淪為一個只講求道德.
只講求經書的信仰.
願意神的靈在我們當中.

$^{601}$我們的禱告是奉基督耶穌名字.
阿們.
\newpage



\section{路加福音 19:1-10}
\label{sec:QuJcmG3w_7k}
\textbf{2023-07-09 午堂《愛在生命改變前》路加福音 19\_1-10 講員:楊明麗姑娘}
\newline
\newline
連結: \href{https://youtube.com/watch?v=QuJcmG3w-7k}{\texttt{https://youtube.com/watch?v=QuJcmG3w-7k}} ~~~~ 語音日期: 2023-07-11
\newline
\newline
\hyperref[sec:87awTrKJFrY]{\small{< < < PREV SERMON < < <}}
~
\hyperref[sec:index_chronic]{\small{[返順時目]}}
~
\hyperref[sec:index_scriptual]{\small{[返順卷目]}}
~
\hyperref[sec:IDYmD1EXhB4]{\small{> > > NEXT SERMON > > >}}
\newline
\newline
路加福音 19:1-10
\newline
\begin{longtable}{cl}
\hline
\hline
章節 & 經文 (和合本修訂版)\\
\hline
19:1 & \begin{tabularx}{0.7\textwidth}{X} 耶穌進了耶利哥，要從那裡經過。 \end{tabularx} \\ \\ \relax
19:2 & \begin{tabularx}{0.7\textwidth}{X} 有一個人名叫撒該，作稅吏長，是個財主。 \end{tabularx} \\ \\ \relax
19:3 & \begin{tabularx}{0.7\textwidth}{X} 他要看看耶穌是怎樣的人，只因人多，他的身材又矮，所以看不見。 \end{tabularx} \\ \\ \relax
19:4 & \begin{tabularx}{0.7\textwidth}{X} 於是他跑到前頭，爬上桑樹，要看耶穌，因為耶穌要從那裡經過。 \end{tabularx} \\ \\ \relax
19:5 & \begin{tabularx}{0.7\textwidth}{X} 耶穌到了那裡，抬頭一看，對他說：「撒該，快下來！今天我必須住在你家裡。」 \end{tabularx} \\ \\ \relax
19:6 & \begin{tabularx}{0.7\textwidth}{X} 他就急忙下來，歡歡喜喜地接待耶穌。 \end{tabularx} \\ \\ \relax
19:7 & \begin{tabularx}{0.7\textwidth}{X} 眾人看見，都私下議論說：「他竟然到罪人家裡去住宿。」 \end{tabularx} \\ \\ \relax
19:8 & \begin{tabularx}{0.7\textwidth}{X} 撒該站著對主說：「主啊，我把所有的一半給窮人；我若勒索了誰，就還他四倍。」 \end{tabularx} \\ \\ \relax
19:9 & \begin{tabularx}{0.7\textwidth}{X} 耶穌對他說：「今天救恩到了這家，因為他也是亞伯拉罕的子孫。 \end{tabularx} \\ \\ \relax
19:10 & \begin{tabularx}{0.7\textwidth}{X} 人子來是要尋找和拯救失喪的人。」 \end{tabularx} \\ \\
[1ex]
\hline
\hline
\end{longtable}
$^{1}$我們一起低頭祈禱 好嗎?.
親天父上帝 我們來到您的面前.
主也是您的恩典.
您的恩典 您看顧我們.
賜給我們救恩 賜給我們力量.
讓我們能夠認識您.
讓我們能夠有勇氣決志.
信主 讓我們現在又能夠.
有力量回來崇拜.
最後 我們很願意將您的話語.
放在心靈當中 求主您賜福.
幫助我們 讓我們得到.
您的聖靈的感動.
能夠在心靈裡領受了您的話語的時候.
最後我們能夠.
可以裝載.
可以實踐 求您幫助我們.
我們仰望您 求您紀念.
奉靠主耶穌.
基督的聖名 阿們.
第一節我想問大家一個問題.
你還記不記得呢.
你和現在的丈夫或太太.
是何時.
如何開始戀愛的呢.
是見了就瞪眼.
那麼有吸引力.
還是何時.
如何.
令到你受吸引.
我聽過一些.
在教會這麼多年.
我聽過一些.
他21歲結婚.
戀愛了7年.
很幸福.
彼此認識後.
進入婚姻.
然後過了一段時間.
弟兄就進入神學院.

$^{41}$太太就照顧.
兩個兒子.
支持他進入神學院.
照顧老爺奶奶.
和他自己的媽媽.
很偉大.
其實愛是何時開始.
可以有這麼大的力量.
你對你身邊.
每天在你身邊的那一位.
你的愛何時開始呢.
你又何時開始.
愛神呢.
是你先愛他.
還是他先愛你呢.
主耶穌基督和撒該的情況.
我們又看看.
他們誰先愛誰.
首先耶穌基督要進入耶利哥城.
耶穌基督這次進入耶利哥城.
應該是他最後一次.
然後在第19章28節.
他就進入耶路撒冷.
那星期他也是.
耶路撒冷出出入入.
然後他為我們.
釘十字架.
刺救因給我們.
在他未進入耶路撒冷之前.
他要經過耶利哥城.
其實不用.
耶利哥城不是路過.
需要經過.
但他來這裡做了兩件事.
一件事就是.
使盲眼的.
能夠看得見.
另一件事.
就是撒該信了主.
他到底是為誰而來呢.

$^{81}$他是不是為盲眼的來.
治好他呢.
因為他從出世開始就盲眼了.
讓他看見.
既然撒該爬到上樹.
救了他.
還是怎樣呢.
我們看看.
耶穌基督要進入耶利哥城的時候.
撒該就早知道.
因為這麼重要的消息.
就好像阿Moba.
李聖林牧師來到我們當中.
講道的時候.
他們很早就知道了.
他去做邀請的工作.
你也可以自己騰出早晨的時間.
來崇拜一樣.
耶穌基督要來耶利哥城.
聽聽他說話.
想他幫身邊的人.
做一些醫治的事.
很多人都會傳說.
他什麼時候來.
他大概什麼時候就走.
只有這麼短的時間.
不僅如此.
還要是耶利哥城的人.
圍著整條大道.
站滿了人.
站滿了人 逼到一個地步.
撒該覺得自己沒有位置.
撒該.
他能夠沒有位置.
我們可以知道他真的有很多人.
這就是Lish這樣形容他.
介紹撒該的時候.
說他是稅利長.
是個財主.
介紹他的原因.

$^{121}$因為他要做一件事.
他要看看耶穌是怎樣的人.
因為人多.
他又不是很高大.
所以他就看不到.
他也知道.
Lish對他的介紹.
這幾節在介紹撒該.
撒該是誰呢.
撒該的意思.
他的名字有一個意思.
他的名字的意思是善良.
純潔.
稅利長.
是很有作為的人.
稅利長在當時的制度當中.
他要懂得衡量.
他是耶利哥城的稅利長.
他要懂得衡量耶利哥城.
可以收到多少稅.
然後向羅馬政府.
好像出標書一樣.
有塊地的話.
就有人出標書.
建築商.
他要標對位置.
夠高.
他才能拿到塊地.
標得太高.
那叫什麼呢.
虧本.
撒該就要有這樣的智慧.
我要做耶利哥城的稅利長.
我應該出一個什麼價.
我就要有一個什麼價.
我就十拿九穩.
拿到這個位.
而我又賺到錢.
不虧了給羅馬政府.
其實他很聰明.

$^{161}$稅利長是很聰明的人.
很懂得思考.
很有商業頭腦的人.
還有一個介紹.
他是財主.
其實聖經.
介紹人是財主.
是很少的.
大家還記得.
有浪子回頭的比喻嗎.
耶穌基督在比喻裡.
一字都沒有說.
父親很富裕.
很有錢.
是個財主.
是個大財主.
沒有.
只是說一個父親.
他有兩個兒子.
不需要用錢財來介紹他.
但是.
說到撒該的時候.
他可能很窮.
他窮到只有錢.
所以就用財主來介紹他.
還有一件事.
介紹他.
他不高.
不過這不是罪.
是他的特徵.
還有一件事.
他不高.
所以他以為自己看不到耶穌.
其實他是財主.
耶利哥城裡面.
是一個很有話事權的.
有話事權.
誰交多少稅的人.
應該很有面子.
很有面子.

$^{201}$每個人都可以讓位.
你來前面坐.
我幫你搭椅子.
我幫你洗澡.
但是他的人際關係.
這裡介紹了.
他是看不到的.
沒有人會給他位.
所以用財主來介紹.
他可能很對.
他很窮.
但是耶穌基督怎樣說呢.
耶穌基督介紹.
介紹.
這個浪子的時候.
沒有這樣介紹.
這個父親是財主.
但是介紹.
這個撒該的時候.
路加就介紹了.
這個撒該是很有錢的.
他想要見耶穌.
他就想辦法.
要見耶穌.
我們知道他爬了一棵樹.
但是你要見耶穌.
就好像進來崇拜.
你要知道幾點鐘.
如果你特別想坐哪個位.
你就知道幾點鐘.
我就坐那個位.
今早有些人就在青衣過來.
有些人在屯門.
是不是有個亞美尼亞的花園.
就是近黃金海岸附近過來.
他要坐在中間.
中間這裡前面一點點的位.
他要很早出門口.
你要知道幾點鐘開始.
幾點鐘你會進到這個位.

$^{241}$這個撒該他要見耶穌的時候.
還要爬上一棵樹的話.
他要知道.
哪一棵樹他可以見到耶穌.
耶穌經過哪一條路.
耶穌什麼時候會到.
無謂這麼熱.
天時中東地帶.
上一棵樹三幾個鐘.
他要有足夠的資料.
用什麼方法可以見到耶穌.
他每樣都預計了.
這個是撒該.
於是他就上了一棵樹.
去預備等耶穌來.
誰知道耶穌就很特別.
耶穌好像一早就知道他在哪棵樹.
走到下面就看上去.
喂你快點下來.
我今晚住在你那裡.
撒該.
原來還要介紹多一樣東西.
他很懂得選擇地方.
選擇一棵可以爬的樹.
知道怎樣去爬.
這棵是當時的桑樹.
不是現在的桑樹科的桑樹.
是無花果樹科的桑樹.
意思是無花果桑樹.
有些果子是黏著樹幹樹枝的.
它不是很高.
所以很容易爬上去.
最高大概20呎.
為什麼很容易爬上去.
大家看到它有很多腳位.
讓你可以爬上去.
離地的腳位不是很大.
爬上去的腳位不是很高.
雖然它不是很高.
不是很高的時候.

$^{281}$就是說它的腳不是很長.
它不是比人高容易爬上去.
是這棵樹適合它爬上去.
它就爬上去.
也有很多樹枝樹葉.
遮蓋它.
讓它飛逝.
不用說它一早就爬上去.
但耶穌基督來到.
就好像沒有東西可以遮蓋.
他知道他在哪裡.
走到那裡.
匆匆的走到那裡.
跟他說:薩埃.
我認識你,你叫薩埃.
快點下來.
我今晚一定要住在你家.
我問一個問題.
薩埃當時愛耶穌嗎?.
我想也有一點.
就算不愛.
也想認識一下.
了解一下.
仰慕一下.
仰慕耶穌.
他有聽過一些關於耶穌的東西.
才去的.
他不會說:我都不知道是什麼水.
我就去看一看.
不會的.
他有收到一些.
聽到一些關於耶穌基督的事.
有打動他的心.
你說他很愛耶穌嗎?.
還沒有.
那耶穌呢?.
耶穌愛他嗎?.
耶穌進入耶利哥城.
經過了整個耶利哥城.
這個應該是近距離要出去.

$^{321}$要離開耶利哥城的城門附近.
然後他對薩埃說.
薩埃,我知道你叫薩埃.
薩埃,你下來.
下來.
我今晚就住在你那裡.
有什麼呢?.
我們有很多時間去談.
在你家裡談.
你預備好,我們就在你家裡談.
是不是?.
耶穌愛他嗎?.
耶穌有沒有要求他?.
不過記得,由你幫我洗腳.
沒有.
耶穌沒有要求什麼.
我要吃龍蝦的.
多弄兩隻鮑魚.
沒有.
耶穌沒有要求.
他只是說,在你家裡住.
耶穌沒有地方住嗎?.
多走兩步,他就出去了.
他將最後的時間.
給了薩埃.
給了這個.
被視為很窮.
窮到只有錢的薩埃.
他給了這個.
在整個耶利哥城.
沒有人喜歡他的人.
薩埃.
耶穌愛他嗎?.
我相信很愛,很有他心.
耶穌基督的鼻釘和十字架.
都有為他.
於是薩埃.
做了一件事.
立刻下來.
開心地.

$^{361}$接待耶穌.
他有沒有預料要接待耶穌?.
沒有.
他預料,能夠看到也好.
這麼高度.
我都預料,看不到.
不過,無論如何都要看一眼.
我爬上去.
但現在耶穌說.
我去你家.
我真的去你家過夜.
薩埃的心.
真是好.
整個耶利哥城都沒有人會和我聊天.
喝杯茶.
吃個包,你來我家.
你真是對我很好.
其實我們不知道.
耶穌何時見到薩埃.
何時決定要去他家住.
還有薩埃.
他決定怎樣招呼耶穌?.
他會不會像馬太一樣.
馬太都是瑞利.
他當晚邀請耶穌.
馬太主動邀請耶穌.
去他家.
然後請了很多.
他的同事來吃飯.
讓耶穌和他們講福音.
薩埃會不會都邀請一大群人去他家呢?.
不知道.
不過都不奇怪.
他的地位他要邀請.
他所有下屬都會來.
他是瑞利長.
他還有很多幫他工作.
收瑞的瑞利.
他可能當晚.
他家的奴隸和工人.

$^{401}$婢女.
要很忙.
但他就很開心.
馬上下來,接了耶穌回去.
那些人怎樣說呢?.
那些人是不是說.
好啊,耶穌去了薩埃家.
以後薩埃不會收重稅了.
不是啊.
那些人說,因為耶穌是誰來的?.
連薩埃這個罪人的家.
他都去?.
罪人家裡過夜?.
我們這裡任何一個人.
都比薩埃更加配去招呼他過夜.
我們這裡任何一個人.
都更加值得他愛和拯救.
我們任何一個人.
都更加配和他喝杯茶.
吃個包.
他來我這裡.
我還宰了兩隻羊.
他們是不是很想.
薩埃可以回到主那邊.
可以重新歸向主呢?.
不是.
但是呢.
薩埃就說了一堆話.
他說了一堆話.
但是呢.
薩埃就說了一堆話.
薩埃站在那裡說.
主啊.
你將我一半的財產.
給窮人.
我勒索過誰.
意思是我強迫過誰交重稅.
我還他四倍.
耶穌基督就說.
今天救恩.

$^{441}$來到薩埃家裡.
因為其實他也是.
亞伯拉罕的子孫.
我們留意一下.
他殺了那棵茄樹.
他應該在那棵樹下來了.
不是坐在地上.
他應該是站在樹上走.
到底突然間說他站在樹上.
他不走了.
站在樹上說還是怎麼樣呢?.
我們不知道他是.
在耶利哥城的街上站在樹上.
還是在接了.
耶穌基督去他家裡吃飯的時候.
他們攤在樹上.
很享受地吃飯.
是不是吃著吃著飯的時候.
突然間站起來說呢?.
我們不知道.
不過.
他有一段很重要的話說.
他說他財產的一半.
他分給窮人.
如果他勒索過.
向誰強取了財物的話.
他就還他四倍.
施洗約翰.
在路加福音.
三章十一節都有說.
施洗約翰說過.
什麼是愛是悔改的心呢?.
你有兩件衣服的.
就分一件給人.
有人沒有的.
有食物的也是這樣.
分一半給人.
他真的做了這件事.
他將他財物的一半.
他說分給人.

$^{481}$《拉吉記》.
摩西的記載.
以色列人要這樣做的.
如果他偷了人家一隻牛.
或者一隻羊.
他不是問人買的.
也不是人家送給他的.
他強搶回來的.
宰了或者賣了也好.
牛要以五隻賠一隻.
羊要以四隻賠一隻.
這些說話其實已經告訴我們了.
薩戈很好家教的.
他們也很勤力看聖經.
很好家教的就是.
這些經文他們做兒童的時候.
幾歲大就會背.
背了說真的.
你記不記得小學讀了什麼.
我只記得第一課.
因為第一課.
幼稚園第一課.
因為那課讀了也不懂.
但那課我時時都記得.
那課就是開學了.
以前我們小朋友的時候.
幼稚園.
很簡單的只有幾個字.
現在不是了.
去我們幼稚園看碧布板.
嘩你寫了這麼多字.
小時候學了的東西.
不記得了.
但薩戈還記得.
他知道的.
知道有什麼他應該做.
就在那天.
在耶穌基督的面前.
他就這樣說.
他不需要騙耶穌.

$^{521}$他不知道耶穌基督快要被釘十字架.
我星期五.
剛好兩天前.
跟一個姐妹一起去九龍探病.
探訪全福音.
四點左右.
我們就在九龍.
坐在她的車上.
我就看到.
為什麼這麼熱.
星期五的四點左右.
整條街都很熱.
金黃色.
但整條街就有很多.
大概是金靈人士.
蹲在那裡坐在那裡排隊.
我想什麼事.
四點多不是銀行齊提.
是什麼事呢.
那條隊很長的.
過了半分鐘左右.
就看到一個龍頭.
原來有一個機構.
一會兒不知道幾點.
細節就看不到.
就看著派飯.
一會兒有飯派.
原來一會兒有飯派.
在這麼炎熱的天氣.
還沒到開始派飯.
就有人排隊.
耶利哥城的居民.
知道殺雞.
是派錢.
他財產的一本.
派錢.
飯都不吃 覺都不睡 排隊.
是不是.
誰認自己很有錢呢.
可以不要殺雞派的錢.

$^{561}$殺雞說還款.
中小企.
馬上派代表你.
他收我們這麼高稅.
還不來.
殺雞說這句話.
說完之後.
他就會窮到連錢都沒有.
但他真的認真地說.
他真的很認真地說.
耶穌基督在這裡.
和他同桌的人在這裡.
或者如果在街上.
剛下來看到人們議論.
這個耶穌是什麼人.
去殺雞家裡吃飯.
於是乎他馬上反應的話.
那整條街都在反應.
他馬上反應的話.
那整條街的人都聽到.
無論他在哪裡說.
都會有很多人知道.
殺雞要派錢.
沒有人會收稅.
誰讓他收過重稅.
那些中小企.
還不去排隊.
就算要爬上樹才行.
都會爬上去.
是.
他這一句話.
他真的要做.
耶穌基督也有一個.
馬上和他有一個說話.
他說人子類.
是要尋找和拯救.
失喪的人.
我們先看一看.
這句話和誰說的.
是耶穌基督對殺雞說的.

$^{601}$今天救恩來到這個家.
因為你也是亞伯拉罕的子孫.
耶穌基督和殺雞說.
救恩來到你這裡.
你得救.
不但如此.
耶穌基督也說.
我們來是要尋找和拯救失喪的人.
殺雞是不是.
肯定他是失喪過.
現在他被主耶穌基督找到了.
耶穌基督經過耶利哥城.
除了說讓孩子能夠見到.
是很重要之外.
令殺雞能夠被尋找到.
被拯救到.
生命有改變是更加重要.
耶穌基督路過耶利哥.
路過耶利哥是有目的的.
他要去拯救殺雞.
耶穌基督何時開始愛殺雞.
原來是很早的.
他是有計劃地進入耶利哥城.
是有計劃進去的.
他進去要令殺雞真的能夠得拯救.
而且不止殺雞.
他是一個瑞利獎影響很多人.
耶穌基督愛殺雞.
所以他進入耶利哥城.
殺雞對耶穌基督的愛.
可能一開始很微.
但因為耶穌基督進了城之後.
他會一路增長.
他對村民的愛.
何時愛到令自己破產.
他將一半的財富拿出來派.
然後我騙了誰.
他的稅錢我還四倍.
他有很多機會成為一個很貧窮的人.
他甚麼都可以沒有了.

$^{641}$他的愛何時開始.
就是見到耶穌基督.
耶穌基督對他的接受.
說要去他家吃飯.
要去他家過夜.
要跟他聊天.
就令到他很大的轉變.
這個是耶穌基督對他的愛.
令到他很受感動.
人子尼要尋找.
要拯救失喪的人.
我們呢.
你是怎樣被拯救的呢.
是你尋找神.
還是神尋找你呢.
我都是一個很難搞的人.
其實我自小是基督教學校讀書.
和天主教學校.
所以我就不是很清楚.
是兩樣我要回哪樣教會.
我有一個同學很可愛的.
我在菲律賓讀中學.
他自小就去過菲律賓.
他從香港過去.
又會說廣東話.
他由我中二到我中四.
每個星期一來我家.
要帶我回教會.
可能是少年人.
你要來我家帶我就不回.
其實他可能不來我家帶我就會回.
三年啊.
那種愛有多奇妙.
用了三年.
每個星期一來.
要陪你回教會.
誰帶你回教會.
有沒有令你有這樣的愛.
其實他這位姊妹很好.
三年都來照顧我.

$^{681}$之後大家就上了不同的大學.
弟兄姊妹.
你跟進一個人你跟多久.
好像今早有報道會.
下星期六有報道會.
你會跟進多久.
邀請人回教會.
或者是怎樣幫助一個信了主的.
好像撒該一樣.
明明是阿帕拉罕的後裔.
但是走到很遠.
跟神的關係很疏離.
你會用多少心機時間去愛他.
關心他.
帶領他回來.
回來天父的身邊.
你願意向他付出多少愛.
或者你願意什麼時候開始愛他.
現在開始吧.
如果還沒開始.
愛他多一點.
其實我們經常都會這樣想.
不行的.
等他有夠先.
等他肯聽.
等他態度好一點.
我要約他.
他不會理我的.
要約他.
他有很多宗教背景.
他不會理的.
他有很多原因.
他住得很遠.
不要想太多.
想想有什麼方法.
帶領他信主.
你可以幫他找近他家的教會.
你可以開車去.
接他過來報道會.
等他信了主再幫他找教會.

$^{721}$你可以早.
我想都不用三年.
早一段時間.
跟他聊天.
吃飯.
看新聞也好.
看電影也好.
打球也好.
做一些事.
讓他知道你愛他.
關心他.
其實你關心他.
說到信主.
你大人信主.
你想起帶不到他信主.
其實你心裡也記住他.
神也讓你心裡記掛他.
神也很想你能夠做多一點.
多走一步.
多發一點你心裡的熱能.
融化他.
讓他能夠得到神的愛.
他現在還沒有.
你覺得他都沒有欠缺什麼.
很平安.
很難去信主嗎.
成年人不會將他心裡的擔憂.
牽掛和不安樂.
說給人聽的.
他多不安樂.
多沒有信心.
他都要表現得說.
我支持你.
有什麼事.
你不行的時候我也支持你.
但是.
他需要你.
他可能很富有.
但是又窮得自有富有.
只有在人家面前.

$^{761}$好像很平安.
但是心靈裡面.
很有需要.
晚上三點.
他需要你的關心.
你的電話.
需要你放消息給他知道.
什麼時候他可以回教會聽福音.
知道有一個很愛他的.
願意為他而死的神.
這個神還可以拯救他.
拯救他讓他心靈得到真真實實的平安和快樂.
是我們生命真實的依靠.
你願意嗎.
耶穌愛我.
耶穌愛你.
耶穌願意我們.
都將這份愛.
好像一個流通的管子那樣.
帶給我們身邊的人.
讓他們真的能夠得到神的愛.
我們不要等他生命改變後才愛他.
不要等他有空才愛他.
不要等他心裡面有些病.
身體有些病.
我們先去愛他.
我們現在就去愛他.
關心他多一點.
一起低頭祈禱.
將天父上帝我們來到你的面前.
我們要多謝你.
你讓我們能夠在生命當中.
平安的時候得到你的恩典.
你的愛.
得到之後我們又繼續有平安的日子.
主要是獻上感恩.
祈求你看顧.
賜福給我們.
讓我們能夠恩著你.
主要是我們能夠恩著你.

$^{801}$我們有愛.
有力量去關心和愛我們身邊的人.
讓我們能夠令他們感受到神你都愛他們.
讓我們能夠令他們感受到他們需要你.
需要你的愛.
需要你的同在.
求主紀念幫助我們.
我們禱告仰望你.
奉主耶穌基督的聖名.
阿們.
\newpage



\section{尼希米記 11:1-3}
\label{sec:IDYmD1EXhB4}
\textbf{2023-07-16 加恩樓《別叫我英雄》尼希米記 11\_1-3 講員:曾裔貴牧師}
\newline
\newline
連結: \href{https://youtube.com/watch?v=IDYmD1EXhB4}{\texttt{https://youtube.com/watch?v=IDYmD1EXhB4}} ~~~~ 語音日期: 2023-07-22
\newline
\newline
\hyperref[sec:QuJcmG3w_7k]{\small{< < < PREV SERMON < < <}}
~
\hyperref[sec:index_chronic]{\small{[返順時目]}}
~
\hyperref[sec:index_scriptual]{\small{[返順卷目]}}
~
\hyperref[sec:vtGjR9atL3U]{\small{> > > NEXT SERMON > > >}}
\newline
\newline
尼希米記 11:1-3
\newline
\begin{longtable}{cl}
\hline
\hline
章節 & 經文 (和合本修訂版)\\
\hline
11:1 & \begin{tabularx}{0.7\textwidth}{X} 百姓的領袖住在耶路撒冷。其餘的百姓抽籤，每十人中選一人來住在聖城耶路撒冷，另外九人住在別的城鎮。 \end{tabularx} \\ \\ \relax
11:2 & \begin{tabularx}{0.7\textwidth}{X} 凡甘心樂意住在耶路撒冷的，百姓都為他們祝福。 \end{tabularx} \\ \\ \relax
11:3 & \begin{tabularx}{0.7\textwidth}{X} 以色列人、祭司、利未人、殿役和所羅門僕人的後裔都住在猶大的城鎮，各在自己城內的地業中。本省的領袖住在耶路撒冷的如下： \end{tabularx} \\ \\
[1ex]
\hline
\hline
\end{longtable}
$^{1}$各位弟兄姊妹,早晨,非常感恩今天能夠因為鄧全道出去講道.
所以我要過來講道,因為平時在錦繡那邊比較忙.
所以今天很難得有機會跟大家分享神的話語.
也有一個簡單的祈禱,多謝主.
當我們打開聖經的時候,幫助我們認識我們這個在天上的主眼中.
這個宿命的遺訓身份,讓我們寶貴之餘.
我們去行駛在我們有生之年,為主,為這個群體,基督徒的群體.
來努力效忠,使更多人透過我們的服事,得到基督的愛和生命的平安.
我們仰望感恩禱告,奉耶穌基督的名,阿們.
一段聖經,很簡單的經文.
就是有十分之一的猶太人,因為他們被抽中,就居住在耶路撒冷.
就是這樣,記得題目,不要搞錯.
近期有一套比較好的電影,別叫我賭神.
我們今天講的是英雄,但是我們要講的不是少數的英雄.
是在講一些默默為神效力,效忠的人.
這段經文是一個嚴肅會之後的抽籤.
因為第七章就是所有的城門,城牆都已經完全做好了.
(我會走的,不用了,夠聲就行了,謝謝).
這個就是當時第三批以色列人回去重建聖城的時候的結果.
第七章尼希米建好城牆之後,就發現了一件事.
我很快讀給大家聽,城牆修完,我安了門扇守門的歌唱的利未人都已經派定了.
我就派我的弟兄哈拿尼和營流的宰官哈拿尼雅管理耶路撒冷.
因為哈拿尼雅是忠信的,忠心信實的一個弟兄.
又敬畏神,過於重仁.
我就吩咐他們,等到太陽上升才可以開耶路撒冷的城門.
人尚且看守的時候就要關門上山.
就是不容許外邦的人,不是猶太人的進入耶路撒冷.
也派耶路撒冷的居民,各按他們的班次,看守自己的房屋對面之處.
不過留意第四節,第七章的第四節.
城是廣大的,其中的居民卻稀少.
房屋都沒有人建造.
然後,尼希米第五節這樣說.
我的神就感動我的心去招贅貴咒.
然後就有一個人口普查.
第七章就做一個詳細的人口普查之後.
去到第八章,一直去到第十章.
其實就是一個延展著二十多天的嚴肅會.
七月初一,猶太人的摧國節.
第八章的第一節這樣說.
到了七月,以色列人就住在自己的城裡,就是國處.

$^{41}$那時他們如同一人聚集在這個水門前的寬闊處.
有一個水門的粉兄會,Watergate的粉兄會.
就不是水門事件.
就請民事以色列將耶和華籍摩西傳給以色列人的律法書帶來了.
在七月初一,祭司這個以色列就將律法書帶到聽了能夠明白的男女.
聽了能夠明白,就是懂得聽到又有回應.
又能夠消化,又明白.
小朋友暫時未必可以了.
他們能夠聽到,明白的.
然後就在這裡延展了一天的宣讀律法書.
然後有利未人開一些研習班.
讓他們在這個節期感受到神話語的能力.
在他們的生命裡.
然後第八章的第九節.
省長尼希米和作祭司的民事以色列.
並教訓百姓的利未人對眾民說.
今天是耶和華你們神的聖日.
不可以哭泣.
為什麼呢?.
這是因為眾民聽見了律法書上的話就哭了.
然後又對他們說.
你們去吃肥美的,喝甘甜的.
如果有不能夠預備的就分給他.
為什麼呢?.
因為今天是我們主的聖日.
你們不可以憂愁.
因為靠耶和華而得的喜樂就是你們的力量.
要百姓聽完之後.
感覺內心有一種傷痛,悲哀,哭泣.
一方面這個時候其實應該是一個很開心的日子.
因為城牆用了53天收好了.
所有的城門都連絡好了.
不是再是冒險雞籠.
這些外邦人隨隨便便進來.
來混亂猶太人自我的身份.
慢慢他們的身份就會認同了.
所以尼希米和以色列.
他們見到一件現象.
很正常的.
當你趕走那些外邦人之後.

$^{81}$很自然你就會發現.
一個這麼大的耶路撒冷.
最需要什麼呢?.
最需要的是人居住在耶路撒冷.
居住在耶路撒冷又如何呢?.
就是要恢復聖殿的敬拜.
聖殿的恢復會帶來什麼後果呢?.
如果每個崗位.
祭司,利未人,電役,守門,歌唱.
能夠各就各位.
各盡他們的職份.
會帶來什麼呢?.
就是百姓恢復向神聖殿的敬拜,獻祭.
在這個獻祭過程.
帶來對自己神選民的身份.
有一個強烈的歸屬和認同.
這時候就會帶來一個戲劇性的轉變.
就是當人的心對自己的內心的身份.
有這個認同之後.
他們就有一份很大的力量.
就是能夠願意.
因為是神的選民.
就願意遵行神召摩西的律法.
來守神的誡命.
所以你看到第八,九章這兩個嚴肅會.
到了第八章的第十三節.
開始他們就發現.
首節有一個意猶未盡.
第十三節,八章的十三節.
次日,眾民的族長,祭司和利未人都聚集.
去到民事,以色列那裡.
要留心去聽律法上的說話.
他們就見律法上寫著.
耶和華籍摩西分庫以色列人.
要在七月十五一直到二十一.
就是駐彭節.
第十五節.
並要在郭城和耶路撒冷宣傳報告說.
你們當上山將橄欖樹,野橄欖樹,翻石榴樹,棕樹.
和各樣茂密樹的枝子.

$^{121}$取來趙雀所寫的答彭.
所以百姓就出去了.
意味著很快他們就要守住彭節.
充滿了一份歡樂.
準備去迎接住彭節.
這個時候.
去到第十七節,八章的十七節.
從老道之地歸回的傳會眾就搭彭了.
住在彭裡.
從亂的兒子約書亞的時候.
直到這個時候.
以色列人都沒有試過這樣來進行.
於是眾人大大的喜樂.
從頭一天,七月十五.
直到末一天,七月二十一.
以色列每天念神的律法書.
眾人守節七日.
第八日,即是七月二十二.
就照例有一個嚴肅會.
親愛的弟兄姊妹.
十一章剛才我們讀的經文.
尼熹米做了一個全文的普查.
之後他就現在跟會眾說.
現在抽籤,不是生死籤.
不是要進去那些dirty room(髒房).
來做那些病房.
是抽籤.
十個支派,逢十抽一.
一個猶太人被抽中.
就要住在耶路撒冷.
百姓就會為被抽中的來祝福.
使他們能夠居住在耶路撒冷裡.
來蒙福.
但是剛才我們已經看了背景.
原來建好城牆,城門之後.
尼熹米完全不是不意識到.
其實耶路撒冷根本沒有人.
需要有被神感動的自己的同胞.
甘心樂意進到神的殿裡服侍.
來讓聖殿的敬拜恢復真正的規模,秩序.

$^{161}$當有了這樣的屬靈的服侍之後.
百姓對自己神的選民身份的認同.
就會慢慢地提高和落實.
大家留意第九章的第一節開始.
「這月的二十四日.
以色列人聚集禁食.
身穿麻衣,頭蒙灰塵.
以色列人就與一切外邦人離絕.
站著承認自己的罪惡和列祖的罪孽.
那日的四分之一站在自己地方.
念他們神的律法書.
由四分之一認罪.
敬拜耶和華他們的神」.
原來尼希米和以色列這兩位領袖.
如何能夠動員自己的同胞.
在耶路撒冷被抽中在那裡定居.
在那裡守護神的聖殿.
不是靠一個行政的指令.
不是靠尼希米是省長.
是亞達斯西王委派他回去重建城牆,城門.
他是有一個法定的權力.
他來做,他來行駛.
沒錯,由第四章到第六章.
他面對外邦人.
去攔阻他重建城牆,城門.
他有時需要用君王給他的行政權.
但是更多的.
這位領袖是禱告和聯絡其他的祭司長.
一起同心去抵禦這些外敵.
但對於自己的同胞.
如何能夠提高他們這一份已經住在耶路撒冷.
已經住在被擄歸回的聖城.
如何能夠提升他們的歸屬,身份的認同.
和帶來的幸福感.
宣讀律法書.
用神的話語去安慰,勉勵你們真是神的選民.
首先就是家譜的普查.
哪些人是祭司的.
就可以回到他自己的崗位.
祭司來的人就肯定.

$^{201}$根本不需要想了.
就要住在耶路撒冷.
在聖殿裡面有足夠的安排,物質的供應.
你們專心去服侍.
但是對於其他支派的猶太人.
在這兩個星期多的嚴肅會.
神的話語感動他們的心靈.
讓他們看到自己和神摩西的律法書的標準.
原來相去甚遠.
其中有幾樣東西.
第十章他們就大家都決心簽約.
跟隨尼熙米的幫助.
決心的去執行.
就是和外邦人離絕.
首先就是安息日.
不幫他們買東西.
不准他們安息日進來做生意.
這個是猶太人.
他們再一次看回自己.
神選民的身份與神納了的約.
在這個約底下.
他們願意分別為聖出來.
另外就是和外邦已經結婚的.
不再有婚姻的文約.
甚至乎不將自己的兒女.
無論是娶外邦的女子.
或者嫁給當地的人.
第十章就記載簽約的大家族.
逐個之派記載.
第十章第一節開始.
簽名的就是一路讀一個家族.
一個家族一路讀下去.
第十八節.
其餘的民.
祭司利未人守門歌唱的離提靈.
和一切離絕淪卸居民的.
歸服神律法的.
並他們的妻子兒女.
凡是有知識能夠明白的.
都隨從他們的貴就弟兄.

$^{241}$來發奏喜誓.
必定遵行神.
藉他的僕人摩西所傳的律法.
謹守遵行耶和華我們主的一切誡命典章律例.
並不將我們的女兒.
嫁給這地的居民.
也不為我們的兒子.
娶他們的女兒.
這地的居民如果在安息日.
或者什麼聖日帶了貨物和糧食.
來賣給我們的時候.
我們必不買.
每逢七年必不耕種.
凡是欠我們債的.
必不追討.
我們又自己訂了例.
每年各人捐銀一四下.
勒三分之一.
為我們神殿的使用.
親愛的弟兄姊妹.
沒錯.
舊約很多的背景.
和我們今天未必拉上關係.
新約已經有一個更大更高的法則.
愛人如己.
我們要來到將人帶進神的舊恩裡.
但是在被擄歸回主前的四百多年的時候.
彌賽亞未來這個世界之前.
他們猶太民族.
仍然要等候這位彌賽亞的降臨.
他們要守住在聖城.
他們的身份.
他們這個選民的身份.
所以尼希米就有被神的靈感動.
要重建這個聖城.
為著的不是整理好耶路撒冷.
為著的是聖殿.
有一個恆常健康正常的敬拜生活.
最終受惠的是在聖城裡面的居民.
對自己神選民的身份.

$^{281}$要有一個認同.
他們當決心來遵守摩西的律法書.
即是和他們有關的律法的時候.
他們就會蒙福了.
他們就會一直守住這個身份.
直至耶穌基督彌賽亞的降臨.
四百多年.
所以在這個時候.
是一個很關鍵的歷史的一個轉向.
百姓回到上去已經是第三代了.
但是那個光景就是.
有城牆但是破爛的.
有門但是都是破損的.
有安息日有恆常的敬拜.
但是是儀式求奇的.
因為外邦人隨隨便便.
就可以來到去嘲笑.
遊論你們的敬拜.
所以百姓的敬拜生活.
是可有可無.
求求其奇的.
沒有了那份真正的.
屬神子民的那份尊嚴.
更加說不上一份榮譽.
但是當尼熙米做好了城牆之後.
吩咐了不同的中心的祭司利未人.
一些守門的弟兄.
來到去看守著.
整個耶路撒冷不同的城門之後.
就賦予一個重要的責任.
就是要讓聖城有所的分別.
要讓外邦的人.
當地的土著.
譬如有其他的.
迦南的人.
譬如參巴拉.
多比亞.
阿拉伯人多士這些外邦人.
這些領袖.
不可以再隨隨便便.

$^{321}$來進入我們敬拜的體系裡.
他們是外邦人.
不敬畏神的.
沒有份沒有關.
要被排除在城門以外.
當猶太人慢慢意識到.
尼熙米不只是緊張硬件.
緊張的是百姓的心靈.
是否真的有神選民的氣質.
當律法書被這兩個星期宣讀之後.
他們自己意識到.
不僅僅是被擄回來的第二第三代後人.
而是他們看到自己.
是有一個很尊貴神選民的身份.
如何落實這個身份呢.
如何讓弟兄姊妹有這個認同呢.
尼熙米和以斯拉甚麼都不做.
宣讀律法書.
開嚴肅會.
大家聚在一起.
祈禱.
尋求.
到底神如何看待我們這個民族.
帶來一個戲劇性的.
居然第十章.
自己主動簽約.
發奏喜誓.
跟隨摩西的律法書去行事.
弟兄姊妹.
教會.
可能很少很少的尼熙米.
有這種遠大的眼光.
他自己本身在波斯帝國.
亞達薩西王的王宮裡面的一個酒精.
是一個很高職位的一個外邦的官員.
但是有神的靈感動他.
他看到耶路撒冷的破落.
他的內心很想為神做一些事.
他就祈禱.
起碼四個月.

$^{361}$直至到君王.
看到他有一些的內心的困難.
就是他內心是有很大的.
那個心靈的那種的憂慮.
所以就問他.
為甚麼你有心事呢.
然後他就說.
他內心和他列祖的神的居所.
現在的破落.
他很想為神去做那個重建.
這樣輾轉就是他回去了.
所以一到第六章.
就是他建好了城牆.
城門的結果.
但是到第七章.
一直去到第十一章.
就告訴我們.
他們還沒有去到奉獻那個城門.
城牆之禮的時候.
很重要的第八,第九章.
就是用神的道.
感動自己的同胞.
讓他們被神的愛.
再一次的去撫摸.
讓他們感受到自己.
神的選民的身份.
這樣就帶來一個很戲劇性的.
當百姓自己被神的愛感動.
看到自己原來是神摩西律法書裡面.
有份的人.
更加是亞伯拉罕的子孫之後.
尼希米以斯拉.
所需要做的就很簡單.
現在聖殿需要有人居住.
要守衛,要服侍.
要有不同的崗位.
一起為神的殿努力.
宣告之後.
就是十一章我們看到的一個很簡單的抽籤儀式.
不是抽居屋.

$^{401}$是誰可以進去耶路撒冷居住.
你想想有了剛才那個背景.
有了那個嚴肅會的24天裡面的感動.
然後尼希米現在做的是非常簡單.
就是每十個人當中.
需要一個人來住在聖城耶路撒冷.
其餘的九個人.
就是那些家庭,那些家族.
就住在自己的城邑.
就是這麼簡單.
不過這裡帶出的就是.
有一群人被百姓祝福.
就是被揀中的.
就是說以能夠被揀中.
住在耶路撒冷為榮.
能夠去到神的殿裡面.
參與服侍是一份榮耀.
這份榮耀慢慢會散發.
去到其他住在其他城市的猶太人的內心.
因為十二章就是一個奉獻城門城牆的大禮.
這是十二章的經文告訴我們的教導.
十二章的二十七節.
就是這個奉獻之禮.
你可以想像得到.
原來建造城門,建造城牆.
為著的是百姓的聖民的身份.
要重建整個耶路撒冷的居民.
對自己的榮耀的身份的認同.
這是很重要.
就是說每一個的百姓.
在神的眼中.
都有你自己的位置.
有你的重要.
有你的尊貴.
在神的眼中.
看每一個聖民都是極其的寶貴.
教會不可能每一個都是英雄.
不可能每一個都是高瞻遠矚的領袖.
如果是固然是好.
但也可能未必是好的.

$^{441}$因為每一個都有自己高瞻遠矚.
就帶來很多不同的方向.
永遠一個群體.
可能很少部分是這樣的領袖.
例如尼希米.
例如以色列.
例如一些利未人.
但是更多的.
更需要的.
就是每一個百姓.
對自己的身份的認同.
每個人參與在神的建造裡面.
剛才也說了.
要等待四百多年後的彌賽亞降臨在伯利恆.
這是神的計劃.
你可以想像得到.
所以尼希米就有這份感動.
參與在這個建造工程裡面.
為著的不是那個硬件的聖城.
為著的是那個時代.
找到一群猶太人.
對自己的身份有一個很強烈的認同.
對自己能夠參與.
服侍.
來為榮.
親愛的弟兄姊妹.
你覺得自己是很平凡嗎?.
還是你覺得.
你是可以參與在神國裡面的英雄呢?.
其實在神眼中.
每個人都不平凡.
每個人實際上都是平凡的.
包括我們經常站在港台的.
我們都是非常平凡的人.
但是不要讓你自己的平凡.
來慢慢地去意識.
蹉跎了你在神國裡面的參與和服侍.
每個人都應該有自己的價值.
我嘗試做一個例子.
有一位女作家.

$^{481}$台灣的.
她叫做劉繼榮.
她寫了一篇文章.
她說她家裡有一個中等的學生.
所謂中等是說成績.
她的女兒就是.
每次考試.
不多不少.
全班有50個人.
她是考第23名.
每次無論大考或是奇考.
都是23名.
爸爸媽媽都很厲害.
也很想女兒成材.
所謂成材就是要很厲害.
要很頂尖.
就雞精班.
補習班就抓她去了.
就幫她補習.
不多不少.
補也好.
不補也好.
都是23名.
所以爸爸就很生氣.
更加生氣就是過時過節.
每個親戚朋友都問.
B仔你有什麼志向.
大個我要做什麼.
通常爸爸就很尷尬.
因為女兒原來胸無大志.
我大個就想做一個很普通的幼稚園教師就行了.
第二個志願再高一點.
我想做一個家庭主婦.
這樣就令爸爸很生氣.
又很沒面子.
怎麼辦呢.
再高一點.
更慘.
就是因為這樣高.
女兒就受不了.

$^{521}$疲於奔命地去學習.
輾轉.
不單止不是23名.
33名.
因為又發高燒.
後來父母徹底放棄.
不再高她了.
接受那個現實.
女兒是中等的23名學生.
爸爸就很生氣.
女兒為什麼你不是神童呢.
女兒就很自然地說.
爸爸因為你不是神父.
那我就不是神童.
不過沒多久之後.
班主任要接見兩父母.
爸爸媽媽就死了.
什麼事呢.
我的女兒是不是犯了什麼事呢.
是不是她的成績一落千丈呢.
班主任就跟她父母說.
我從來教書沒試過.
沒見過這樣的學生.
沒錯.
你的女兒仍然是23名.
不是突然間高了低了那些成績.
但是班主任就說.
我給全班一個奇巧.
以外的一個問題.
就是你找全班50個人.
你能夠說出.
為什麼你要推舉她做全班的領袖.
要說原因.
班主任就跟她的父母說.
跟劉繼榮說.
我從來沒見過.
49個學生.
全部來提名你的女兒.
做這個領袖.
原來這個女兒.

$^{561}$是一個很強烈的親和力的人.
每次家庭的聚會.
她會帶來上上下下的人的歡樂.
是沒有人不會因為她的笑話.
她自己不笑.
就令到整個家庭的宴會.
充滿了快樂.
充滿了一份祝福.
那些小朋友在爭取.
她很簡單.
就可以令到兩個堅持不下的男孩.
可以握手分那個食物.
她可以很輕鬆化解這些危機.
回家之後.
兩父母就很有榮譽感.
女兒原來你是英雄.
女兒說我不想做什麼英雄.
不過我看過一本書.
英雄的出現.
周圍有很多人拍手掌.
我很想做那個拍手掌的人.
怎麼可能的EQ.
怎麼可能的這麼強的親和力.
對自己的身份這麼大的認同.
這些真的值得我們去思考.
值得我們去想想.
在新約我找到一節經文.
史陶保羅在哥林多前書.
十六章第十五第十六節.
差不多結尾了.
保羅處理了哥林多很多很多的紛爭.
接著他講了一個人.
你知道斯提反一家.
是阿該亞省初結的果子.
當時信主很早期信主.
他以事奉聖徒為念.
由保羅的口中講出這句話.
這個人是保羅所敬重.
這個人專心以事奉聖徒為念.
回到教會想什麼呢.

$^{601}$就想想我怎麼去服侍人.
想想哪些人有需要.
我就默默地服侍.
保羅說這樣的人無要敬重.
親愛的弟兄姊妹.
今天的訊息很想說.
這十分一被抽中的人.
不是中了巨獎.
不是抽中居屋.
是他們與有榮焉的被選立.
進入耶路撒冷居住.
服侍來帶動聖殿的敬拜.
讓聖殿恢復健康活潑的敬拜.
就會帶來什麼呢.
百姓感受到自己的榮耀.
神的事公就可以一路一路等待.
彌賽亞我們的救主.
在公元第六年降世來到這個世界.
成就人類的救恩.
我們一起禱告.
多謝主.
帶領今早我們的敬拜.
我們的聖餐.
再再告訴我們.
十字架的救恩將我們拯救.
我們是神的兒女.
這是何等大的福氣.
幫助我們每一個弟兄姊妹.
都看到我們在神的眼中.
我們都是寶貴的.
沒錯我們未必是出類拔萃.
但是我們可以貢獻我們的內心.
來服侍我們的主.
在教會裡更加是一個重要的場景.
我們可以為福音去效力.
求主的聖靈.
感動我們.
讓我們看到.
到處去留意.
有什麼我們可以去幫忙.

$^{641}$有什麼我們可以服侍.
我們就盡心盡力.
做在神的眼裡.
這樣的感恩禱告.
奉耶穌基督的名,阿們..
\newpage



\section{使徒行傳 1:8-28:31}
\label{sec:vtGjR9atL3U}
\textbf{2023-07-16 午堂《放膽傳道,無人能禁止》使徒行傳 1\_8,28\_31 講員:郭鴻標牧師}
\newline
\newline
連結: \href{https://youtube.com/watch?v=vtGjR9atL3U}{\texttt{https://youtube.com/watch?v=vtGjR9atL3U}} ~~~~ 語音日期: 2023-07-23
\newline
\newline
\hyperref[sec:IDYmD1EXhB4]{\small{< < < PREV SERMON < < <}}
~
\hyperref[sec:index_chronic]{\small{[返順時目]}}
~
\hyperref[sec:index_scriptual]{\small{[返順卷目]}}
~
\hyperref[sec:PIfMPVh5oUA]{\small{> > > NEXT SERMON > > >}}
\newline
\newline
使徒行傳 1:8-28:31
\newline
\begin{longtable}{cl}
\hline
\hline
章節 & 經文 (和合本修訂版)\\
\hline
1:8 & \begin{tabularx}{0.7\textwidth}{X} 但聖靈降臨在你們身上，你們就必得著能力，並要在耶路撒冷、猶太全地和撒瑪利亞，直到地極，作我的見證。」 \end{tabularx} \\ \\ \relax
1:9 & \begin{tabularx}{0.7\textwidth}{X} 說了這些話，他們正看的時候，他被接上升，有一朵雲彩從他們眼前把他接去。 \end{tabularx} \\ \\ \relax
1:10 & \begin{tabularx}{0.7\textwidth}{X} 他升上去的時候，他們定睛望天，看哪，有兩個人身穿白衣站在他們旁邊， \end{tabularx} \\ \\ \relax
1:11 & \begin{tabularx}{0.7\textwidth}{X} 說：「加利利人哪，你們為甚麼站著望天呢？這離開你們被接升天的耶穌，你們見他怎樣升上天去，他也要怎樣來臨。」 \end{tabularx} \\ \\ \relax
1:12 & \begin{tabularx}{0.7\textwidth}{X} 有一座山，名叫橄欖山，離耶路撒冷不遠，有安息日可行走的路程。那時，門徒從那裡回耶路撒冷去， \end{tabularx} \\ \\ \relax
1:13 & \begin{tabularx}{0.7\textwidth}{X} 他們一進城，就上了所住的樓房；在那裡有彼得、約翰、雅各、安得烈、腓力、多馬、巴多羅買、馬太、亞勒腓的兒子雅各、激進黨的西門，和雅各的兒子猶大。 \end{tabularx} \\ \\ \relax
1:14 & \begin{tabularx}{0.7\textwidth}{X} 這些人和幾個婦人，包括耶穌的母親馬利亞，和耶穌的兄弟，都同心合意地恆切禱告。 \end{tabularx} \\ \\ \relax
1:15 & \begin{tabularx}{0.7\textwidth}{X} 那時，有許多人聚會，約有一百二十名，彼得在弟兄中間站起來，說： \end{tabularx} \\ \\ \relax
1:16 & \begin{tabularx}{0.7\textwidth}{X} 「諸位弟兄，聖經的話必須應驗。聖經中，聖靈曾藉大衛的口預先說到那領人來拿耶穌的猶大； \end{tabularx} \\ \\ \relax
1:17 & \begin{tabularx}{0.7\textwidth}{X} 他本來算是我們中的一個，並且得了這一份使徒的職任。 \end{tabularx} \\ \\ \relax
1:18 & \begin{tabularx}{0.7\textwidth}{X} 這人用他不義的代價買了一塊田，以後身子仆倒，肚腹崩裂，腸子都流出來。 \end{tabularx} \\ \\ \relax
1:19 & \begin{tabularx}{0.7\textwidth}{X} 住在耶路撒冷的人都知道這事，所以按著他們當地的話把那塊田叫亞革大馬，就是「血田」的意思。 \end{tabularx} \\ \\ \relax
1:20 & \begin{tabularx}{0.7\textwidth}{X} 因為《詩篇》上寫著：「願他的住處變為廢墟，無人在內居住。」又說：「願別人得他的職分。」 \end{tabularx} \\ \\ \relax
1:21 & \begin{tabularx}{0.7\textwidth}{X} 所以，主耶穌在我們中間出入的整段時間， \end{tabularx} \\ \\ \relax
1:22 & \begin{tabularx}{0.7\textwidth}{X} 就是從約翰施洗起，直到主離開我們被接上升的日子為止，必須從那常與我們一起的人中，立一位與我們同作耶穌復活的見證。」 \end{tabularx} \\ \\ \relax
1:23 & \begin{tabularx}{0.7\textwidth}{X} 於是他們推舉兩個人，就是那叫巴撒巴，又稱為猶士都的約瑟，和馬提亞。 \end{tabularx} \\ \\ \relax
1:24 & \begin{tabularx}{0.7\textwidth}{X} 眾人禱告說：「主啊，你知道萬人的心，求你從這兩個人中指明你所揀選的是哪一位， \end{tabularx} \\ \\ \relax
1:25 & \begin{tabularx}{0.7\textwidth}{X} 去得這使徒的職任；這職位猶大已經丟棄，往自己的地方去了。」 \end{tabularx} \\ \\ \relax
1:26 & \begin{tabularx}{0.7\textwidth}{X} 於是眾人為他們搖籤，搖出馬提亞來；他就和十一個使徒同列。 \end{tabularx} \\ \\ \relax
2:1 & \begin{tabularx}{0.7\textwidth}{X} 五旬節那日到了，他們全都聚集在一起。 \end{tabularx} \\ \\ \relax
2:2 & \begin{tabularx}{0.7\textwidth}{X} 忽然，有響聲從天上下來，好像一陣大風吹過，充滿了他們所坐的整座屋子； \end{tabularx} \\ \\ \relax
2:3 & \begin{tabularx}{0.7\textwidth}{X} 又有舌頭如火焰向他們顯現，分開落在他們每個人身上。 \end{tabularx} \\ \\ \relax
2:4 & \begin{tabularx}{0.7\textwidth}{X} 他們都被聖靈充滿，就按著聖靈所賜的口才說起別國的話來。 \end{tabularx} \\ \\ \relax
2:5 & \begin{tabularx}{0.7\textwidth}{X} 那時，有從天下各國來的虔誠的猶太人，住在耶路撒冷。 \end{tabularx} \\ \\ \relax
2:6 & \begin{tabularx}{0.7\textwidth}{X} 這聲音一響，許多人都來聚集，各人因為聽見門徒用他們各自的鄉談說話，就甚納悶， \end{tabularx} \\ \\ \relax
2:7 & \begin{tabularx}{0.7\textwidth}{X} 都詫異驚奇說：「看哪，這些說話的不都是加利利人嗎？ \end{tabularx} \\ \\ \relax
2:8 & \begin{tabularx}{0.7\textwidth}{X} 我們每個人怎麼聽見他們說我們生來所用的鄉談呢？ \end{tabularx} \\ \\ \relax
2:9 & \begin{tabularx}{0.7\textwidth}{X} 我們帕提亞人、瑪代人、以攔人，和住在美索不達米亞、猶太、加帕多家、本都、亞細亞、 \end{tabularx} \\ \\ \relax
2:10 & \begin{tabularx}{0.7\textwidth}{X} 弗呂家、旁非利亞、埃及的人，並靠近古利奈的利比亞一帶地方的人，僑居的羅馬人， \end{tabularx} \\ \\ \relax
2:11 & \begin{tabularx}{0.7\textwidth}{X} 包括猶太人和皈依猶太教的人， 克里特人和阿拉伯人，都聽見他們用我們的鄉談講論神的大作為。」 \end{tabularx} \\ \\ \relax
2:12 & \begin{tabularx}{0.7\textwidth}{X} 眾人就都驚奇困惑，彼此說：「這是甚麼意思呢？」 \end{tabularx} \\ \\ \relax
2:13 & \begin{tabularx}{0.7\textwidth}{X} 還有人譏誚，說：「他們是灌滿了新酒吧！」 \end{tabularx} \\ \\ \relax
2:14 & \begin{tabularx}{0.7\textwidth}{X} 彼得和十一個使徒站起來，他就高聲向眾人說：「猶太人和所有住在耶路撒冷的人哪，這件事你們要知道，要側耳聽我的話。 \end{tabularx} \\ \\ \relax
2:15 & \begin{tabularx}{0.7\textwidth}{X} 這些人並不像你們所想的喝醉了，因為現在才早晨九點鐘。 \end{tabularx} \\ \\ \relax
2:16 & \begin{tabularx}{0.7\textwidth}{X} 這正是藉著先知約珥所說的： \end{tabularx} \\ \\ \relax
2:17 & \begin{tabularx}{0.7\textwidth}{X} 『神說：在末後的日子，我要將我的靈澆灌凡血肉之軀的。你們的兒女要說預言；你們的少年要見異象；你們的老人要做異夢。 \end{tabularx} \\ \\ \relax
2:18 & \begin{tabularx}{0.7\textwidth}{X} 在那些日子，我要把我的靈澆灌，甚至給我的僕人和婢女，他們要說預言。 \end{tabularx} \\ \\ \relax
2:19 & \begin{tabularx}{0.7\textwidth}{X} 在天上，我要顯出奇事，在地下，我要顯出神蹟，有血，有火，有煙霧。 \end{tabularx} \\ \\ \relax
2:20 & \begin{tabularx}{0.7\textwidth}{X} 太陽要變為黑暗，月亮要變為血，這都在主大而光榮的日子未到以前。 \end{tabularx} \\ \\ \relax
2:21 & \begin{tabularx}{0.7\textwidth}{X} 那時，凡求告主名的都必得救。』 \end{tabularx} \\ \\ \relax
2:22 & \begin{tabularx}{0.7\textwidth}{X} 「以色列人哪，你們要聽我這些話：拿撒勒人耶穌就是神以異能、奇事、神蹟向你們證明出來的人，這些事是神藉著他在你們中間施行，正如你們自己知道的。 \end{tabularx} \\ \\ \relax
2:23 & \begin{tabularx}{0.7\textwidth}{X} 他既按著神確定的旨意和預知被交與人，你們就藉著不法之人的手把他釘在十字架上，殺了。 \end{tabularx} \\ \\ \relax
2:24 & \begin{tabularx}{0.7\textwidth}{X} 神卻將死的痛苦解除，使他復活了，因為他原不能被死拘禁。 \end{tabularx} \\ \\ \relax
2:25 & \begin{tabularx}{0.7\textwidth}{X} 大衛指著他說：『我看見主常在我眼前，他在我右邊，使我不至於動搖。 \end{tabularx} \\ \\ \relax
2:26 & \begin{tabularx}{0.7\textwidth}{X} 所以我心裡歡喜，我的舌頭快樂，而且我的肉身要安居在指望中。 \end{tabularx} \\ \\ \relax
2:27 & \begin{tabularx}{0.7\textwidth}{X} 因你必不將我的靈魂撇在陰間，也不讓你的聖者見朽壞。 \end{tabularx} \\ \\ \relax
2:28 & \begin{tabularx}{0.7\textwidth}{X} 你已將生命的道路指示我，必使我在你面前充滿快樂。』 \end{tabularx} \\ \\ \relax
2:29 & \begin{tabularx}{0.7\textwidth}{X} 「諸位弟兄，先祖大衛的事，我可以坦然地對你們說：他死了，也埋葬了，而且他的墳墓直到今日還在我們這裡。 \end{tabularx} \\ \\ \relax
2:30 & \begin{tabularx}{0.7\textwidth}{X} 既然大衛是先知，他知道神曾向他起誓，要從他的後裔中立一位坐在他的寶座上。 \end{tabularx} \\ \\ \relax
2:31 & \begin{tabularx}{0.7\textwidth}{X} 他預先看見了，就講論基督的復活，說：『他不被撇在陰間；他的肉身也不見朽壞。』 \end{tabularx} \\ \\ \relax
2:32 & \begin{tabularx}{0.7\textwidth}{X} 這耶穌，神已經使他復活了，我們都是這事的見證人。 \end{tabularx} \\ \\ \relax
2:33 & \begin{tabularx}{0.7\textwidth}{X} 他既被高舉在神的右邊，又從父受了所應許的聖靈，就把你們所看見所聽見的，澆灌下來。 \end{tabularx} \\ \\ \relax
2:34 & \begin{tabularx}{0.7\textwidth}{X} 大衛並沒有升到天上，但他自己說：『主對我主說：你坐在我的右邊， \end{tabularx} \\ \\ \relax
2:35 & \begin{tabularx}{0.7\textwidth}{X} 等我使你的仇敵作你的腳凳。』 \end{tabularx} \\ \\ \relax
2:36 & \begin{tabularx}{0.7\textwidth}{X} 故此，以色列全家當確實知道，你們釘在十字架上的這位耶穌，神已經立他為主，為基督了。」 \end{tabularx} \\ \\ \relax
2:37 & \begin{tabularx}{0.7\textwidth}{X} 眾人聽見這話，覺得扎心，就對彼得和其餘的使徒說：「諸位弟兄，我們該怎樣做呢？」 \end{tabularx} \\ \\ \relax
2:38 & \begin{tabularx}{0.7\textwidth}{X} 彼得對他們說：「你們各人要悔改，奉耶穌基督的名受洗，使你們的罪得赦免，就會領受所賜的聖靈。 \end{tabularx} \\ \\ \relax
2:39 & \begin{tabularx}{0.7\textwidth}{X} 因為這應許是給你們和你們的兒女，並一切在遠方的人，就是給所有主—我們的神所召來的人。」 \end{tabularx} \\ \\ \relax
2:40 & \begin{tabularx}{0.7\textwidth}{X} 彼得還用更多別的話作見證，勸勉他們說：「你們當救自己脫離這彎曲的世代。」 \end{tabularx} \\ \\ \relax
2:41 & \begin{tabularx}{0.7\textwidth}{X} 於是領受他話的人，都受了洗；那一天，門徒約添了三千人。 \end{tabularx} \\ \\ \relax
2:42 & \begin{tabularx}{0.7\textwidth}{X} 他們都專注於使徒的教導和彼此的團契，擘餅和祈禱。 \end{tabularx} \\ \\ \relax
2:43 & \begin{tabularx}{0.7\textwidth}{X} 眾人都心存敬畏；使徒們又行了許多奇事神蹟。 \end{tabularx} \\ \\ \relax
2:44 & \begin{tabularx}{0.7\textwidth}{X} 信的人都聚在一處，凡物公用， \end{tabularx} \\ \\ \relax
2:45 & \begin{tabularx}{0.7\textwidth}{X} 又賣了田產和家業，照每一個人所需要的分給他們。 \end{tabularx} \\ \\ \relax
2:46 & \begin{tabularx}{0.7\textwidth}{X} 他們天天同心合意恆切地在聖殿裡敬拜，且在家中擘餅，存著歡喜坦誠的心用飯， \end{tabularx} \\ \\ \relax
2:47 & \begin{tabularx}{0.7\textwidth}{X} 讚美神，得全體百姓的喜愛。主將得救的人天天加給他們。 \end{tabularx} \\ \\ \relax
3:1 & \begin{tabularx}{0.7\textwidth}{X} 下午三點鐘禱告的時候，彼得和約翰上聖殿去。 \end{tabularx} \\ \\ \relax
3:2 & \begin{tabularx}{0.7\textwidth}{X} 一個從母腹裡就是瘸腿的人正被人抬來，他們天天把他放在聖殿的一個叫美門的門口，求進聖殿的人施捨。 \end{tabularx} \\ \\ \relax
3:3 & \begin{tabularx}{0.7\textwidth}{X} 他看見彼得、約翰將要進聖殿，就求他們施捨。 \end{tabularx} \\ \\ \relax
3:4 & \begin{tabularx}{0.7\textwidth}{X} 彼得和約翰定睛看他，彼得說：「看著我們！」 \end{tabularx} \\ \\ \relax
3:5 & \begin{tabularx}{0.7\textwidth}{X} 那人就注目看他們，指望從他們得著甚麼。 \end{tabularx} \\ \\ \relax
3:6 & \begin{tabularx}{0.7\textwidth}{X} 彼得卻說：「金銀我都沒有，但我把我有的給你：奉拿撒勒人耶穌基督的名起來行走！」 \end{tabularx} \\ \\ \relax
3:7 & \begin{tabularx}{0.7\textwidth}{X} 於是彼得拉著他的右手，扶他起來；他的腳和踝骨立刻健壯了， \end{tabularx} \\ \\ \relax
3:8 & \begin{tabularx}{0.7\textwidth}{X} 就跳起來，站著，又開始行走。他跟他們進了聖殿，邊走邊跳，讚美神。 \end{tabularx} \\ \\ \relax
3:9 & \begin{tabularx}{0.7\textwidth}{X} 百姓都看見他又行走，又讚美神， \end{tabularx} \\ \\ \relax
3:10 & \begin{tabularx}{0.7\textwidth}{X} 認得他是那素常坐在聖殿的美門口求人施捨的，就因他所遇到的事滿心驚訝詫異。 \end{tabularx} \\ \\ \relax
3:11 & \begin{tabularx}{0.7\textwidth}{X} 那人正在稱為所羅門的廊下，拉住彼得和約翰，大家都覺得很驚訝，一齊跑到他們那裡。 \end{tabularx} \\ \\ \relax
3:12 & \begin{tabularx}{0.7\textwidth}{X} 彼得看見，就對百姓說：「以色列人哪，為甚麼因這事而驚訝呢？為甚麼定睛看我們，以為我們憑自己的能力和虔誠使這人行走呢？ \end{tabularx} \\ \\ \relax
3:13 & \begin{tabularx}{0.7\textwidth}{X} 亞伯拉罕的神、以撒的神、雅各的神，就是我們列祖的神，已經榮耀了他的僕人耶穌，這耶穌就是你們交付官府的那位，彼拉多決定要釋放他時，你們卻在彼拉多面前棄絕了他。 \end{tabularx} \\ \\ \relax
3:14 & \begin{tabularx}{0.7\textwidth}{X} 你們棄絕了那聖潔公義者，反而要求釋放一個兇手給你們。 \end{tabularx} \\ \\ \relax
3:15 & \begin{tabularx}{0.7\textwidth}{X} 你們殺了那生命的創始者，神卻叫他從死人中復活；我們都是這事的見證人。 \end{tabularx} \\ \\ \relax
3:16 & \begin{tabularx}{0.7\textwidth}{X} 因信他的名，他的名使你們所看見所認識的這人健壯了；正是他所賜的信心使這人在你們眾人面前完全好了。 \end{tabularx} \\ \\ \relax
3:17 & \begin{tabularx}{0.7\textwidth}{X} 「如今，弟兄們，我知道你們做這事是出於無知，你們的官長也是如此。 \end{tabularx} \\ \\ \relax
3:18 & \begin{tabularx}{0.7\textwidth}{X} 但神藉著眾先知的口預先宣告過基督將要受害的事，就這樣應驗了。 \end{tabularx} \\ \\ \relax
3:19 & \begin{tabularx}{0.7\textwidth}{X} 所以，你們當悔改歸正，使你們的罪得以塗去， \end{tabularx} \\ \\ \relax
3:20 & \begin{tabularx}{0.7\textwidth}{X} 這樣，那安舒的日子就必從主面前來到；主也必差遣所預定給你們的基督耶穌來臨。 \end{tabularx} \\ \\ \relax
3:21 & \begin{tabularx}{0.7\textwidth}{X} 他必須留在天上，直到萬物復興的時候，就是神自古藉著聖先知的口所說的。 \end{tabularx} \\ \\ \relax
3:22 & \begin{tabularx}{0.7\textwidth}{X} 摩西曾說：『主—你們的神要從你們弟兄中給你們興起一位先知像我，凡他向你們所說的一切，你們都要聽從。 \end{tabularx} \\ \\ \relax
3:23 & \begin{tabularx}{0.7\textwidth}{X} 凡不聽從那先知的，必將從民中滅絕。』 \end{tabularx} \\ \\ \relax
3:24 & \begin{tabularx}{0.7\textwidth}{X} 從撒母耳以來和後繼的眾先知，凡說預言的，也都曾宣告這些日子。 \end{tabularx} \\ \\ \relax
3:25 & \begin{tabularx}{0.7\textwidth}{X} 你們是先知的子孫，也是神與你們祖宗所立之約的子孫，就是對亞伯拉罕說：『地上萬族都將因你的後裔得福。』 \end{tabularx} \\ \\ \relax
3:26 & \begin{tabularx}{0.7\textwidth}{X} 神既興起他的僕人，就先差他到你們這裡來，賜福給你們，使各人回轉，離開你們的邪惡。」 \end{tabularx} \\ \\ \relax
4:1 & \begin{tabularx}{0.7\textwidth}{X} 彼得和約翰正向百姓說話的時候，祭司們、守殿官和撒都該人來了， \end{tabularx} \\ \\ \relax
4:2 & \begin{tabularx}{0.7\textwidth}{X} 就很煩惱，因為使徒們教導百姓，傳揚在耶穌的事上證明有死人復活， \end{tabularx} \\ \\ \relax
4:3 & \begin{tabularx}{0.7\textwidth}{X} 於是下手拿住他們；因為天已經晚了，就把他們押在拘留所到第二天。 \end{tabularx} \\ \\ \relax
4:4 & \begin{tabularx}{0.7\textwidth}{X} 但聽道的人有許多信了，男人的數目約有五千。 \end{tabularx} \\ \\ \relax
4:5 & \begin{tabularx}{0.7\textwidth}{X} 第二天，官長、長老和文士在耶路撒冷聚集， \end{tabularx} \\ \\ \relax
4:6 & \begin{tabularx}{0.7\textwidth}{X} 又有亞那大祭司、該亞法、約翰、亞歷山大，和大祭司的親族都在那裡。 \end{tabularx} \\ \\ \relax
4:7 & \begin{tabularx}{0.7\textwidth}{X} 他們叫使徒站在中間，問他們：「你們憑甚麼能力，奉誰的名做這事呢？」 \end{tabularx} \\ \\ \relax
4:8 & \begin{tabularx}{0.7\textwidth}{X} 那時，彼得被聖靈充滿，對他們說：「民間的官長和長老啊， \end{tabularx} \\ \\ \relax
4:9 & \begin{tabularx}{0.7\textwidth}{X} 倘若今日我們被查問是因為在殘障的人身上所行的善事，就是這人怎麼得了痊癒， \end{tabularx} \\ \\ \relax
4:10 & \begin{tabularx}{0.7\textwidth}{X} 那麼，你們大家和以色列全民都當知道，站在你們面前的這人得痊癒，是因你們所釘在十字架、神使他從死人中復活的拿撒勒人耶穌基督的名。 \end{tabularx} \\ \\ \relax
4:11 & \begin{tabularx}{0.7\textwidth}{X} 這位耶穌是：『你們匠人所丟棄的石頭，已成了房角的頭塊石頭。』 \end{tabularx} \\ \\ \relax
4:12 & \begin{tabularx}{0.7\textwidth}{X} 除他以外，別無拯救，因為在天下人間，沒有賜下別的名，我們可以靠著得救。」 \end{tabularx} \\ \\ \relax
4:13 & \begin{tabularx}{0.7\textwidth}{X} 他們見彼得、約翰的膽量，又看出他們原是沒有學問的平民，就很驚訝，認出他們曾是跟耶穌一起的； \end{tabularx} \\ \\ \relax
4:14 & \begin{tabularx}{0.7\textwidth}{X} 又看見那治好了的人和他們一同站著，就無話可駁。 \end{tabularx} \\ \\ \relax
4:15 & \begin{tabularx}{0.7\textwidth}{X} 於是他們吩咐他們兩人從議會退出，就彼此商議， \end{tabularx} \\ \\ \relax
4:16 & \begin{tabularx}{0.7\textwidth}{X} 說：「我們當怎樣辦這兩個人呢？因為他們誠然行了一件明顯的神蹟，凡住在耶路撒冷的人都知道，我們也不能否認。 \end{tabularx} \\ \\ \relax
4:17 & \begin{tabularx}{0.7\textwidth}{X} 但為避免這事越發在民間傳揚，我們必須威嚇他們，叫他們不可再奉這名對任何人講論。」 \end{tabularx} \\ \\ \relax
4:18 & \begin{tabularx}{0.7\textwidth}{X} 於是他們叫了兩人來，禁止他們，再不可奉耶穌的名講論或教導人。 \end{tabularx} \\ \\ \relax
4:19 & \begin{tabularx}{0.7\textwidth}{X} 彼得和約翰回答他們說：「聽從你們，不聽從神，在神面前合理不合理，你們自己判斷吧！ \end{tabularx} \\ \\ \relax
4:20 & \begin{tabularx}{0.7\textwidth}{X} 我們所看見所聽見的，我們不能不說。」 \end{tabularx} \\ \\ \relax
4:21 & \begin{tabularx}{0.7\textwidth}{X} 官長為百姓的緣故，想不出任何法子懲罰他們，只好威嚇一番就把他們釋放了；這是因眾人為了所行的奇事都歸榮耀與神。 \end{tabularx} \\ \\ \relax
4:22 & \begin{tabularx}{0.7\textwidth}{X} 原來經歷這神蹟醫好的人有四十多歲了。 \end{tabularx} \\ \\ \relax
4:23 & \begin{tabularx}{0.7\textwidth}{X} 二人既被釋放，就到自己的人那裡去，把祭司長和長老所說的話都告訴他們。 \end{tabularx} \\ \\ \relax
4:24 & \begin{tabularx}{0.7\textwidth}{X} 他們聽見了，就同心合意地高聲向神說：「主宰啊！你是那創造天、地、海和其中萬物的； \end{tabularx} \\ \\ \relax
4:25 & \begin{tabularx}{0.7\textwidth}{X} 你曾藉著聖靈託你僕人—我們祖宗大衛的口說：『外邦為甚麼擾動？萬民為甚麼謀算虛妄的事？ \end{tabularx} \\ \\ \relax
4:26 & \begin{tabularx}{0.7\textwidth}{X} 地上的君王都站穩，臣宰也聚集一處，要對抗主，對抗主的受膏者。』 \end{tabularx} \\ \\ \relax
4:27 & \begin{tabularx}{0.7\textwidth}{X} 希律和本丟‧彼拉多，同外邦人和以色列民，果然在這城裡聚集，要攻打你所膏的聖僕耶穌， \end{tabularx} \\ \\ \relax
4:28 & \begin{tabularx}{0.7\textwidth}{X} 做了你手和你旨意所預定必成就的事。 \end{tabularx} \\ \\ \relax
4:29 & \begin{tabularx}{0.7\textwidth}{X} 主啊，現在求你鑒察，他們的威嚇，使你僕人放膽講你的道， \end{tabularx} \\ \\ \relax
4:30 & \begin{tabularx}{0.7\textwidth}{X} 伸出你的手來，讓醫治、神蹟、奇事藉著你聖僕耶穌的名行出來。」 \end{tabularx} \\ \\ \relax
4:31 & \begin{tabularx}{0.7\textwidth}{X} 他們禱告完了，聚會的地方震動；他們都被聖靈充滿，放膽傳講神的道。 \end{tabularx} \\ \\ \relax
4:32 & \begin{tabularx}{0.7\textwidth}{X} 許多信徒都一心一意，沒有一人說他的任何東西是自己的，都是大家公用。 \end{tabularx} \\ \\ \relax
4:33 & \begin{tabularx}{0.7\textwidth}{X} 使徒以大能見證主耶穌復活；眾人也都蒙了大恩。 \end{tabularx} \\ \\ \relax
4:34 & \begin{tabularx}{0.7\textwidth}{X} 他們當中沒有一個缺乏的，因為凡有田產房屋的都賣了，把所賣的錢拿來， \end{tabularx} \\ \\ \relax
4:35 & \begin{tabularx}{0.7\textwidth}{X} 放在使徒腳前，照每人所需要的，分給每人。 \end{tabularx} \\ \\ \relax
4:36 & \begin{tabularx}{0.7\textwidth}{X} 有一個利未人，名叫約瑟，使徒稱他為巴拿巴（巴拿巴翻出來就是安慰之子），生在塞浦路斯。 \end{tabularx} \\ \\ \relax
4:37 & \begin{tabularx}{0.7\textwidth}{X} 他有田地，也賣了，把錢拿來，放在使徒腳前。 \end{tabularx} \\ \\ \relax
5:1 & \begin{tabularx}{0.7\textwidth}{X} 有一個人，名叫亞拿尼亞，同他的妻子撒非喇，賣了田產， \end{tabularx} \\ \\ \relax
5:2 & \begin{tabularx}{0.7\textwidth}{X} 把錢私自留下一部分，他的妻子也知道，其餘的部分拿來放在使徒腳前。 \end{tabularx} \\ \\ \relax
5:3 & \begin{tabularx}{0.7\textwidth}{X} 彼得說：「亞拿尼亞！為甚麼撒但充滿了你的心，使你欺騙聖靈，把賣田地的錢私自留下一部分呢？ \end{tabularx} \\ \\ \relax
5:4 & \begin{tabularx}{0.7\textwidth}{X} 田地還沒有賣，不是你自己的嗎？既賣了，錢不是你作主嗎？你怎麼心裡會想這樣做呢？你不是欺騙人，是欺騙神！」 \end{tabularx} \\ \\ \relax
5:5 & \begin{tabularx}{0.7\textwidth}{X} 亞拿尼亞一聽見這些話，就仆倒，斷了氣；所有聽見的人都非常懼怕。 \end{tabularx} \\ \\ \relax
5:6 & \begin{tabularx}{0.7\textwidth}{X} 有些年輕人起來，把他裹好，抬出去埋葬了。 \end{tabularx} \\ \\ \relax
5:7 & \begin{tabularx}{0.7\textwidth}{X} 約過了三小時，他的妻子進來，還不知道所發生的事。 \end{tabularx} \\ \\ \relax
5:8 & \begin{tabularx}{0.7\textwidth}{X} 彼得對她說：「你告訴我，你們賣田地的錢就是這些嗎？」她說：「就是這些。」 \end{tabularx} \\ \\ \relax
5:9 & \begin{tabularx}{0.7\textwidth}{X} 彼得對她說：「你們為甚麼同謀來試探主的靈呢？你看，埋葬你丈夫之人的腳已到門口，他們也要把你抬出去。」 \end{tabularx} \\ \\ \relax
5:10 & \begin{tabularx}{0.7\textwidth}{X} 她立刻仆倒在彼得腳前，斷了氣。那些年輕人進來，見她已經死了，就把她抬出去，埋在她丈夫旁邊。 \end{tabularx} \\ \\ \relax
5:11 & \begin{tabularx}{0.7\textwidth}{X} 全教會和所有聽見這些事的人都非常懼怕。 \end{tabularx} \\ \\ \relax
5:12 & \begin{tabularx}{0.7\textwidth}{X} 主藉使徒的手在民間行了許多神蹟奇事；他們都同心合意地聚集在所羅門的廊下。 \end{tabularx} \\ \\ \relax
5:13 & \begin{tabularx}{0.7\textwidth}{X} 其餘的人沒有一個敢接近他們，百姓卻尊重他們。 \end{tabularx} \\ \\ \relax
5:14 & \begin{tabularx}{0.7\textwidth}{X} 信主的人越發增添，連男帶女都很多， \end{tabularx} \\ \\ \relax
5:15 & \begin{tabularx}{0.7\textwidth}{X} 甚至有人將病人抬到街上，放在床上或褥子上，好讓彼得走過來的時候，或者影子投在一些人身上。 \end{tabularx} \\ \\ \relax
5:16 & \begin{tabularx}{0.7\textwidth}{X} 還有許多人帶著病人和被污靈纏磨的，從耶路撒冷四圍的城鎮來，他們全都得了醫治。 \end{tabularx} \\ \\ \relax
5:17 & \begin{tabularx}{0.7\textwidth}{X} 於是，大祭司採取行動，他和他所有一起的人，就是撒都該派的人，滿心忌恨， \end{tabularx} \\ \\ \relax
5:18 & \begin{tabularx}{0.7\textwidth}{X} 就下手拿住使徒，把他們押在公共拘留所內。 \end{tabularx} \\ \\ \relax
5:19 & \begin{tabularx}{0.7\textwidth}{X} 但在夜間主的使者開了監門，領他們出來，說： \end{tabularx} \\ \\ \relax
5:20 & \begin{tabularx}{0.7\textwidth}{X} 「你們去，站在聖殿裡，把這生命的一切話講給百姓聽。」 \end{tabularx} \\ \\ \relax
5:21 & \begin{tabularx}{0.7\textwidth}{X} 使徒聽了這話，天將亮的時候就進聖殿裡去教導人。大祭司和他一起的人來了，叫齊議會的人和以色列人的眾長老，然後派人到監牢裡去把使徒提出來。 \end{tabularx} \\ \\ \relax
5:22 & \begin{tabularx}{0.7\textwidth}{X} 但差役到了，不見他們在監裡，就回來稟報， \end{tabularx} \\ \\ \relax
5:23 & \begin{tabularx}{0.7\textwidth}{X} 說：「我們看見監牢關得很緊，警衛也站在門外，但打開門來，裡面一個人都不見。」 \end{tabularx} \\ \\ \relax
5:24 & \begin{tabularx}{0.7\textwidth}{X} 守殿官和祭司長聽了這些話，心裡困惑，不知這事將來如何。 \end{tabularx} \\ \\ \relax
5:25 & \begin{tabularx}{0.7\textwidth}{X} 有一個人來稟報說：「你們押在監裡的人，現在站在聖殿裡教導百姓。」 \end{tabularx} \\ \\ \relax
5:26 & \begin{tabularx}{0.7\textwidth}{X} 於是守殿官和差役去帶使徒來，並沒有用暴力，因為怕百姓用石頭打他們。 \end{tabularx} \\ \\ \relax
5:27 & \begin{tabularx}{0.7\textwidth}{X} 他們把使徒帶來了，就叫他們站在議會前。大祭司問他們， \end{tabularx} \\ \\ \relax
5:28 & \begin{tabularx}{0.7\textwidth}{X} 說：「我們不是嚴嚴地禁止你們，不可奉這名教導人嗎？看，你們倒把你們的道理充滿了耶路撒冷，想要叫這人的血歸到我們身上！」 \end{tabularx} \\ \\ \relax
5:29 & \begin{tabularx}{0.7\textwidth}{X} 彼得和眾使徒回答：「我們必須順從神，勝於順從人。 \end{tabularx} \\ \\ \relax
5:30 & \begin{tabularx}{0.7\textwidth}{X} 你們掛在木頭上殺害的耶穌，我們祖宗的神已經使他復活了。 \end{tabularx} \\ \\ \relax
5:31 & \begin{tabularx}{0.7\textwidth}{X} 神把他高舉在自己的右邊，使他作元帥，作救主，使以色列人得以悔改，並且罪得赦免。 \end{tabularx} \\ \\ \relax
5:32 & \begin{tabularx}{0.7\textwidth}{X} 我們是這些事的見證人；神賜給順從的人的聖靈也為這些事作見證。」 \end{tabularx} \\ \\ \relax
5:33 & \begin{tabularx}{0.7\textwidth}{X} 議會的人聽了極其惱怒，想要殺他們。 \end{tabularx} \\ \\ \relax
5:34 & \begin{tabularx}{0.7\textwidth}{X} 但有一個法利賽人，名叫迦瑪列，是眾百姓所敬重的律法教師，他在議會中站起來，吩咐人把使徒暫且帶到外面去， \end{tabularx} \\ \\ \relax
5:35 & \begin{tabularx}{0.7\textwidth}{X} 然後對眾人說：「以色列人哪，對於這些人，你們應當小心怎樣處理。 \end{tabularx} \\ \\ \relax
5:36 & \begin{tabularx}{0.7\textwidth}{X} 從前杜達出現，自命不凡，附從他的人數約有四百；他被殺後，附從他的人全都散了，歸於無有。 \end{tabularx} \\ \\ \relax
5:37 & \begin{tabularx}{0.7\textwidth}{X} 此後，登記戶籍的時候，又有加利利的猶大出現，引誘百姓跟從他，他也滅亡，附從他的人也都四散了。 \end{tabularx} \\ \\ \relax
5:38 & \begin{tabularx}{0.7\textwidth}{X} 現在，我勸你們不要管這些人，任憑他們吧！他們所謀所為若是出於人，必要敗壞； \end{tabularx} \\ \\ \relax
5:39 & \begin{tabularx}{0.7\textwidth}{X} 若是出於神，你們就不能敗壞他們，恐怕你們倒是攻擊神了。」議會的人被他說服了， \end{tabularx} \\ \\ \relax
5:40 & \begin{tabularx}{0.7\textwidth}{X} 就叫使徒來，把他們打了，又吩咐他們不可奉耶穌的名講道，然後把他們釋放了。 \end{tabularx} \\ \\ \relax
5:41 & \begin{tabularx}{0.7\textwidth}{X} 他們歡歡喜喜地離開議會，因他們算配為這名受辱。 \end{tabularx} \\ \\ \relax
5:42 & \begin{tabularx}{0.7\textwidth}{X} 他們就每日在聖殿裡，在家裡，不住地教導人，傳耶穌是基督的福音。 \end{tabularx} \\ \\ \relax
6:1 & \begin{tabularx}{0.7\textwidth}{X} 那些日子，門徒增多，有說希臘話的猶太人向希伯來人發怨言，因為在日常的供給上忽略了他們的寡婦。 \end{tabularx} \\ \\ \relax
6:2 & \begin{tabularx}{0.7\textwidth}{X} 十二使徒叫眾門徒來，說：「我們撇下神的道去管理飯食，是不合宜的。 \end{tabularx} \\ \\ \relax
6:3 & \begin{tabularx}{0.7\textwidth}{X} 所以弟兄們，當從你們中間選出七個有好名聲、滿有聖靈和智慧，我們派他們管理這事。 \end{tabularx} \\ \\ \relax
6:4 & \begin{tabularx}{0.7\textwidth}{X} 至於我們，我們要專注於祈禱和傳道的事奉。」 \end{tabularx} \\ \\ \relax
6:5 & \begin{tabularx}{0.7\textwidth}{X} 這話使全會眾都喜悅，就揀選了司提反—他是一個滿有信心和聖靈的人；他們又揀選了腓利、伯羅哥羅、尼迦挪、提門、巴米拿，並皈依猶太教的安提阿人尼哥拉， \end{tabularx} \\ \\ \relax
6:6 & \begin{tabularx}{0.7\textwidth}{X} 叫他們站在使徒面前，使徒禱告後，就為他們按手。 \end{tabularx} \\ \\ \relax
6:7 & \begin{tabularx}{0.7\textwidth}{X} 神的道興旺起來；在耶路撒冷門徒數目增加得很多，也有許多祭司聽從了這信仰。 \end{tabularx} \\ \\ \relax
6:8 & \begin{tabularx}{0.7\textwidth}{X} 司提反滿有恩惠和能力，在民間行了大奇事和神蹟。 \end{tabularx} \\ \\ \relax
6:9 & \begin{tabularx}{0.7\textwidth}{X} 當時有從稱為「自由人」會堂，並古利奈、亞歷山大會堂來的人，還有些從基利家、亞細亞來的人，起來和司提反辯論。 \end{tabularx} \\ \\ \relax
6:10 & \begin{tabularx}{0.7\textwidth}{X} 司提反是以智慧和聖靈說話，眾人抵擋不住， \end{tabularx} \\ \\ \relax
6:11 & \begin{tabularx}{0.7\textwidth}{X} 就收買人來說：「我們聽見他說褻瀆摩西和神的話。」 \end{tabularx} \\ \\ \relax
6:12 & \begin{tabularx}{0.7\textwidth}{X} 他們又煽動百姓、長老和文士，就突然來捉拿他，把他帶到議會去， \end{tabularx} \\ \\ \relax
6:13 & \begin{tabularx}{0.7\textwidth}{X} 設下假見證，說：「這個人不斷地說話，侮辱神聖的地方和律法。 \end{tabularx} \\ \\ \relax
6:14 & \begin{tabularx}{0.7\textwidth}{X} 我們曾聽見他說，這拿撒勒人耶穌要毀壞這地方，也要改變摩西所交給我們的規矩。」 \end{tabularx} \\ \\ \relax
6:15 & \begin{tabularx}{0.7\textwidth}{X} 在議會裡坐著的人都定睛看他，見他的面貌好像天使的面貌。 \end{tabularx} \\ \\ \relax
7:1 & \begin{tabularx}{0.7\textwidth}{X} 大祭司說：「果真有這些事嗎？」 \end{tabularx} \\ \\ \relax
7:2 & \begin{tabularx}{0.7\textwidth}{X} 司提反說：「諸位父老弟兄請聽！從前我們的祖宗亞伯拉罕在美索不達米亞，還沒有住在哈蘭的時候，榮耀的神向他顯現， \end{tabularx} \\ \\ \relax
7:3 & \begin{tabularx}{0.7\textwidth}{X} 對他說：『你要離開本地和親族，往我所要指示你的地去。』 \end{tabularx} \\ \\ \relax
7:4 & \begin{tabularx}{0.7\textwidth}{X} 他就離開迦勒底人的地方，住在哈蘭。他父親死了以後，神使他從那裡搬到你們現在所住的地方。 \end{tabularx} \\ \\ \relax
7:5 & \begin{tabularx}{0.7\textwidth}{X} 在這裡神並沒有給他產業，連立足的地方都沒有，但應許要將這地賜給他和他的後裔為業，雖然那時他還沒有兒子。 \end{tabularx} \\ \\ \relax
7:6 & \begin{tabularx}{0.7\textwidth}{X} 神這樣說：『他的後裔必寄居外邦，那裡的人要使他們作奴隸，苦待他們四百年。』 \end{tabularx} \\ \\ \relax
7:7 & \begin{tabularx}{0.7\textwidth}{X} 神又說：『但我要懲罰使他們作奴隸的那國。以後他們要出來，在這地方事奉我。』 \end{tabularx} \\ \\ \relax
7:8 & \begin{tabularx}{0.7\textwidth}{X} 神又賜他割禮的約。於是亞伯拉罕生了以撒，在第八日給他行了割禮；後來以撒生雅各，雅各生十二位先祖。 \end{tabularx} \\ \\ \relax
7:9 & \begin{tabularx}{0.7\textwidth}{X} 「先祖嫉妒約瑟，把他賣到埃及去，神卻與他同在， \end{tabularx} \\ \\ \relax
7:10 & \begin{tabularx}{0.7\textwidth}{X} 救他脫離一切苦難，又使他在埃及王法老面前蒙恩，又有智慧。法老派他作埃及國的宰相兼管法老的全家。 \end{tabularx} \\ \\ \relax
7:11 & \begin{tabularx}{0.7\textwidth}{X} 後來全埃及和迦南遭遇饑荒和大災難，我們的祖宗絕了糧。 \end{tabularx} \\ \\ \relax
7:12 & \begin{tabularx}{0.7\textwidth}{X} 雅各聽見在埃及有糧，就打發我們的祖宗初次往那裡去。 \end{tabularx} \\ \\ \relax
7:13 & \begin{tabularx}{0.7\textwidth}{X} 第二次約瑟與兄弟們相認，法老才認識他的家族。 \end{tabularx} \\ \\ \relax
7:14 & \begin{tabularx}{0.7\textwidth}{X} 約瑟就打發人，請父親雅各和全族七十五個人都來。 \end{tabularx} \\ \\ \relax
7:15 & \begin{tabularx}{0.7\textwidth}{X} 於是雅各下了埃及，後來他和我們的祖宗都死在那裡； \end{tabularx} \\ \\ \relax
7:16 & \begin{tabularx}{0.7\textwidth}{X} 他們又被遷到示劍，葬於亞伯拉罕在示劍用銀子從哈抹子孫買來的墳墓裡。 \end{tabularx} \\ \\ \relax
7:17 & \begin{tabularx}{0.7\textwidth}{X} 「當神應許亞伯拉罕的日期將到的時候，以色列人在埃及人丁興旺， \end{tabularx} \\ \\ \relax
7:18 & \begin{tabularx}{0.7\textwidth}{X} 直到另一位不認識約瑟的王興起統治埃及。 \end{tabularx} \\ \\ \relax
7:19 & \begin{tabularx}{0.7\textwidth}{X} 他用詭計待我們的宗族，苦待我們的祖宗，強迫他們丟棄嬰孩，使嬰孩不能存活。 \end{tabularx} \\ \\ \relax
7:20 & \begin{tabularx}{0.7\textwidth}{X} 就在那時，摩西生了下來，神看為俊美，在父親家裡被撫養了三個月。 \end{tabularx} \\ \\ \relax
7:21 & \begin{tabularx}{0.7\textwidth}{X} 他被丟棄的時候，法老的女兒拾了去，當自己的兒子撫養。 \end{tabularx} \\ \\ \relax
7:22 & \begin{tabularx}{0.7\textwidth}{X} 摩西學了埃及人一切的學問，說話辦事都有才能。 \end{tabularx} \\ \\ \relax
7:23 & \begin{tabularx}{0.7\textwidth}{X} 「他到了四十歲，心中起意去看望他的弟兄以色列人。 \end{tabularx} \\ \\ \relax
7:24 & \begin{tabularx}{0.7\textwidth}{X} 他見他們中的一個人受冤屈，就庇護他，為那被壓迫的人報仇，打死了那埃及人。 \end{tabularx} \\ \\ \relax
7:25 & \begin{tabularx}{0.7\textwidth}{X} 他以為他的弟兄們必明白神是藉他的手搭救他們，他們卻不明白。 \end{tabularx} \\ \\ \relax
7:26 & \begin{tabularx}{0.7\textwidth}{X} 第二天，他遇見有人在打架，就想勸他們和好，說：『二位，你們是弟兄，為甚麼彼此欺負呢？』 \end{tabularx} \\ \\ \relax
7:27 & \begin{tabularx}{0.7\textwidth}{X} 那欺負鄰舍的人把他推開，說：『誰立你作我們的領袖和審判官呢？ \end{tabularx} \\ \\ \relax
7:28 & \begin{tabularx}{0.7\textwidth}{X} 難道你要殺我像昨天殺那埃及人一樣嗎？』 \end{tabularx} \\ \\ \relax
7:29 & \begin{tabularx}{0.7\textwidth}{X} 摩西聽見這話就逃走了，寄居於米甸地，在那裡生了兩個兒子。 \end{tabularx} \\ \\ \relax
7:30 & \begin{tabularx}{0.7\textwidth}{X} 「過了四十年，在西奈山的曠野，有一位天使在荊棘的火焰中向摩西顯現。 \end{tabularx} \\ \\ \relax
7:31 & \begin{tabularx}{0.7\textwidth}{X} 摩西見了那異象，覺得很驚訝，正往前觀看的時候，有主的聲音說： \end{tabularx} \\ \\ \relax
7:32 & \begin{tabularx}{0.7\textwidth}{X} 『我是你列祖的神，就是亞伯拉罕、以撒、雅各的神。』摩西戰戰兢兢，不敢觀看。 \end{tabularx} \\ \\ \relax
7:33 & \begin{tabularx}{0.7\textwidth}{X} 主對他說：『把你腳上的鞋脫下來，因為你所站的地方是聖地。 \end{tabularx} \\ \\ \relax
7:34 & \begin{tabularx}{0.7\textwidth}{X} 我的百姓在埃及所受的困苦，我確實看見了；他們悲嘆的聲音，我也聽見了。我下來要救他們。現在，你來，我要差你往埃及去。』 \end{tabularx} \\ \\ \relax
7:35 & \begin{tabularx}{0.7\textwidth}{X} 「這摩西就是有人曾棄絕他說『誰立你作我們的領袖和審判官』的，神卻藉那在荊棘中顯現的天使的手差派他作領袖，作解救者。 \end{tabularx} \\ \\ \relax
7:36 & \begin{tabularx}{0.7\textwidth}{X} 這人領以色列人出來，在埃及地，在紅海，在曠野的四十年間行了奇事神蹟。 \end{tabularx} \\ \\ \relax
7:37 & \begin{tabularx}{0.7\textwidth}{X} 這人是摩西，就是那曾對以色列人說『神要從你們弟兄中給你們興起一位先知像我』的。 \end{tabularx} \\ \\ \relax
7:38 & \begin{tabularx}{0.7\textwidth}{X} 這人是那曾在曠野的會眾中和西奈山上，與那對他說話的天使同在，又與我們祖宗同在的，他領受了活潑的聖言傳給我們。 \end{tabularx} \\ \\ \relax
7:39 & \begin{tabularx}{0.7\textwidth}{X} 我們的祖宗不肯聽從，反棄絕他，他們的心轉向埃及， \end{tabularx} \\ \\ \relax
7:40 & \begin{tabularx}{0.7\textwidth}{X} 對亞倫說：『你為我們造神明，在我們前面引路，因為領我們出埃及地的這個摩西，我們不知道他遭遇了甚麼事。』 \end{tabularx} \\ \\ \relax
7:41 & \begin{tabularx}{0.7\textwidth}{X} 那時，他們造了一個牛犢，又拿祭物獻給那像，為自己手所做的工作歡躍。 \end{tabularx} \\ \\ \relax
7:42 & \begin{tabularx}{0.7\textwidth}{X} 但是神轉臉不顧，任憑他們祭拜天上的日月星辰，正如先知書上所寫的：『以色列家啊，你們四十年間在曠野，何曾將犧牲和祭物獻給我？ \end{tabularx} \\ \\ \relax
7:43 & \begin{tabularx}{0.7\textwidth}{X} 你們抬著摩洛的帳幕和理番---你們神明的星，就是你們所造為要敬拜的像。因此，我要把你們遷到巴比倫外去。』 \end{tabularx} \\ \\ \relax
7:44 & \begin{tabularx}{0.7\textwidth}{X} 「我們的祖宗在曠野，有作證的會幕，是神吩咐摩西照著他所看見的樣式做的。 \end{tabularx} \\ \\ \relax
7:45 & \begin{tabularx}{0.7\textwidth}{X} 這帳幕，我們的祖宗同約書亞相繼承受了，當神在他們面前趕走外邦人的時候，他們把這帳幕搬進承受為業之地，直存到大衛的日子。 \end{tabularx} \\ \\ \relax
7:46 & \begin{tabularx}{0.7\textwidth}{X} 大衛在神面前蒙恩，祈求為雅各的家預備居所。 \end{tabularx} \\ \\ \relax
7:47 & \begin{tabularx}{0.7\textwidth}{X} 但卻是所羅門為神造成殿宇。 \end{tabularx} \\ \\ \relax
7:48 & \begin{tabularx}{0.7\textwidth}{X} 其實，至高者並不住人手所造的，就如先知所言： \end{tabularx} \\ \\ \relax
7:49 & \begin{tabularx}{0.7\textwidth}{X} 『主說：天是我的寶座，地是我的腳凳。你們要為我造怎樣的殿宇？哪裡是我安歇的地方呢？ \end{tabularx} \\ \\ \relax
7:50 & \begin{tabularx}{0.7\textwidth}{X} 這一切不都是我手所造的嗎？』 \end{tabularx} \\ \\ \relax
7:51 & \begin{tabularx}{0.7\textwidth}{X} 「你們這硬著頸項，心與耳未受割禮的人哪，時常抗拒聖靈！你們的祖宗怎樣，你們也怎樣。 \end{tabularx} \\ \\ \relax
7:52 & \begin{tabularx}{0.7\textwidth}{X} 先知中有哪一個不是受你們祖宗的迫害呢？他們把預先宣告那義者要來的人殺了。如今你們成了那義者的出賣者和兇手了。 \end{tabularx} \\ \\ \relax
7:53 & \begin{tabularx}{0.7\textwidth}{X} 你們領受了天使所傳佈的律法，竟不遵守。」 \end{tabularx} \\ \\ \relax
7:54 & \begin{tabularx}{0.7\textwidth}{X} 眾人聽見這些話，心中極其惱怒，向司提反咬牙切齒。 \end{tabularx} \\ \\ \relax
7:55 & \begin{tabularx}{0.7\textwidth}{X} 但司提反滿有聖靈，定睛望天，看見神的榮耀，又看見耶穌站在神的右邊， \end{tabularx} \\ \\ \relax
7:56 & \begin{tabularx}{0.7\textwidth}{X} 就說：「我看見天開了，人子站在神的右邊。」 \end{tabularx} \\ \\ \relax
7:57 & \begin{tabularx}{0.7\textwidth}{X} 眾人大聲喊叫，摀著耳朵，齊心衝向他， \end{tabularx} \\ \\ \relax
7:58 & \begin{tabularx}{0.7\textwidth}{X} 把他推到城外，用石頭打他。作見證的人把他們的衣裳放在一個名叫掃羅的青年腳前。 \end{tabularx} \\ \\ \relax
7:59 & \begin{tabularx}{0.7\textwidth}{X} 他們正用石頭打司提反的時候，他呼求說：「主耶穌啊，求你接納我的靈魂！」 \end{tabularx} \\ \\ \relax
7:60 & \begin{tabularx}{0.7\textwidth}{X} 然後他跪下來，大聲喊著：「主啊，不要將這罪歸於他們！」說了這話，就長眠了。 \end{tabularx} \\ \\ \relax
8:1 & \begin{tabularx}{0.7\textwidth}{X} 掃羅也贊同處死他。從那一天開始，耶路撒冷的教會遭受到大迫害，除了使徒以外，眾門徒都分散在猶太和撒瑪利亞各處。 \end{tabularx} \\ \\ \relax
8:2 & \begin{tabularx}{0.7\textwidth}{X} 有些虔誠的人把司提反埋葬了，為他大大哀哭。 \end{tabularx} \\ \\ \relax
8:3 & \begin{tabularx}{0.7\textwidth}{X} 掃羅卻殘害教會，挨家挨戶地進去，拉著男女關在監裡。 \end{tabularx} \\ \\ \relax
8:4 & \begin{tabularx}{0.7\textwidth}{X} 那些分散的人往各地去傳福音的道。 \end{tabularx} \\ \\ \relax
8:5 & \begin{tabularx}{0.7\textwidth}{X} 腓利下撒瑪利亞城去，向當地人宣講基督。 \end{tabularx} \\ \\ \relax
8:6 & \begin{tabularx}{0.7\textwidth}{X} 眾人都聚精會神，同心合意地聽腓利所說的話，一邊聽他的話，一邊看他所行的神蹟。 \end{tabularx} \\ \\ \relax
8:7 & \begin{tabularx}{0.7\textwidth}{X} 因為有許多人被污靈附著，那些污靈大聲呼叫，從他們身上出來；還有許多癱瘓的、瘸腿的都得了醫治。 \end{tabularx} \\ \\ \relax
8:8 & \begin{tabularx}{0.7\textwidth}{X} 那城裡，有極大的喜樂。 \end{tabularx} \\ \\ \relax
8:9 & \begin{tabularx}{0.7\textwidth}{X} 有一個人名叫西門，向來在那城裡行邪術，自命為大人物，使撒瑪利亞的居民驚奇。 \end{tabularx} \\ \\ \relax
8:10 & \begin{tabularx}{0.7\textwidth}{X} 所有的人，從小到大都聽從他，說：「這個人就是神的能力，那稱為大能者的。」 \end{tabularx} \\ \\ \relax
8:11 & \begin{tabularx}{0.7\textwidth}{X} 他們聽從他，因他很久以來用邪術使他們驚奇。 \end{tabularx} \\ \\ \relax
8:12 & \begin{tabularx}{0.7\textwidth}{X} 當他們信了腓利所傳神國的福音和耶穌基督的名，連男帶女都受了洗。 \end{tabularx} \\ \\ \relax
8:13 & \begin{tabularx}{0.7\textwidth}{X} 西門自己也信了；既受了洗，就常與腓利在一處，看見他所行的神蹟和大異能，就覺得很驚奇。 \end{tabularx} \\ \\ \relax
8:14 & \begin{tabularx}{0.7\textwidth}{X} 在耶路撒冷的使徒聽見撒瑪利亞人領受了神的道，就打發彼得和約翰到他們那裡去。 \end{tabularx} \\ \\ \relax
8:15 & \begin{tabularx}{0.7\textwidth}{X} 兩個人下去，就為他們禱告，要讓他們領受聖靈， \end{tabularx} \\ \\ \relax
8:16 & \begin{tabularx}{0.7\textwidth}{X} 因為聖靈還沒有降在他們任何一個人身上，他們只奉主耶穌的名受了洗。 \end{tabularx} \\ \\ \relax
8:17 & \begin{tabularx}{0.7\textwidth}{X} 於是使徒按手在他們頭上，他們就領受了聖靈。 \end{tabularx} \\ \\ \relax
8:18 & \begin{tabularx}{0.7\textwidth}{X} 西門看見使徒一按手，就有聖靈賜下，就拿錢給使徒， \end{tabularx} \\ \\ \relax
8:19 & \begin{tabularx}{0.7\textwidth}{X} 說：「請把這權柄也給我，使我手按著誰，誰就可以領受聖靈。」 \end{tabularx} \\ \\ \relax
8:20 & \begin{tabularx}{0.7\textwidth}{X} 彼得對他說：「你的銀子和你一同滅亡吧！因為你想神的恩賜是可以用錢買的。 \end{tabularx} \\ \\ \relax
8:21 & \begin{tabularx}{0.7\textwidth}{X} 你在這道上無份無關；因為你在神面前心懷不正。 \end{tabularx} \\ \\ \relax
8:22 & \begin{tabularx}{0.7\textwidth}{X} 你要為你這樣的惡而悔改，祈求主，或者你心裡的意念可得赦免。 \end{tabularx} \\ \\ \relax
8:23 & \begin{tabularx}{0.7\textwidth}{X} 我看出你正在苦膽之中，被不義捆綁著。」 \end{tabularx} \\ \\ \relax
8:24 & \begin{tabularx}{0.7\textwidth}{X} 西門回答說：「請你們為我求主，使你們所說的，沒有一樣臨到我身上。」 \end{tabularx} \\ \\ \relax
8:25 & \begin{tabularx}{0.7\textwidth}{X} 使徒既作了見證，並且宣講了主的道，就回耶路撒冷去，一路在撒瑪利亞好些村莊傳揚福音。 \end{tabularx} \\ \\ \relax
8:26 & \begin{tabularx}{0.7\textwidth}{X} 有主的一個使者對腓利說：「起來！向南走，往那從耶路撒冷下迦薩的路上去。」那路是曠野。 \end{tabularx} \\ \\ \relax
8:27 & \begin{tabularx}{0.7\textwidth}{X} 腓利就起身去了。不料，有一個埃塞俄比亞人，是個有大權的太監，在埃塞俄比亞女王甘大基的手下總管銀庫，他上耶路撒冷去禮拜。 \end{tabularx} \\ \\ \relax
8:28 & \begin{tabularx}{0.7\textwidth}{X} 回程中，他坐在車上，正念著以賽亞先知的書， \end{tabularx} \\ \\ \relax
8:29 & \begin{tabularx}{0.7\textwidth}{X} 聖靈對腓利說：「你去！靠近那車走。」 \end{tabularx} \\ \\ \relax
8:30 & \begin{tabularx}{0.7\textwidth}{X} 腓利就跑到太監那裡，聽見他正在念以賽亞先知的書，就說：「你明白你所念的嗎？」 \end{tabularx} \\ \\ \relax
8:31 & \begin{tabularx}{0.7\textwidth}{X} 他說：「沒有人指教我，怎能明白呢？」於是他請腓利上車，與他同坐。 \end{tabularx} \\ \\ \relax
8:32 & \begin{tabularx}{0.7\textwidth}{X} 他所念的那段經文是這樣：「他像羊被牽去宰殺，又像羔羊在剪毛的人手下無聲，他也是這樣不開口。 \end{tabularx} \\ \\ \relax
8:33 & \begin{tabularx}{0.7\textwidth}{X} 他卑微的時候，得不到公義的審判，誰能述說他的身世？因為他的生命從地上被奪去。」 \end{tabularx} \\ \\ \relax
8:34 & \begin{tabularx}{0.7\textwidth}{X} 太監回答腓利說：「請問，先知說這話是指誰，是指自己，還是指別人呢？」 \end{tabularx} \\ \\ \relax
8:35 & \begin{tabularx}{0.7\textwidth}{X} 腓利就開口，從這段經文開始，對他傳講耶穌的福音。 \end{tabularx} \\ \\ \relax
8:36 & \begin{tabularx}{0.7\textwidth}{X} 二人正沿路往前走，到了有水的地方，太監說：「看哪！這裡有水，有甚麼能阻止我受洗呢？」 \end{tabularx} \\ \\ \relax
8:37 & \begin{tabularx}{0.7\textwidth}{X} 腓利說：『你若是一心相信，就可以。』他回答：『我信耶穌基督是神的兒子。』 \end{tabularx} \\ \\ \relax
8:38 & \begin{tabularx}{0.7\textwidth}{X} 於是他吩咐把車停下來，腓利和太監二人一同下到水裡，腓利就給他施洗。 \end{tabularx} \\ \\ \relax
8:39 & \begin{tabularx}{0.7\textwidth}{X} 他們從水裡上來，主的靈把腓利提了去，太監再也看不見他了，就歡歡喜喜地上路。 \end{tabularx} \\ \\ \relax
8:40 & \begin{tabularx}{0.7\textwidth}{X} 後來有人在亞鎖都遇見腓利；他走遍那地方，在各城宣揚福音，一直到凱撒利亞。 \end{tabularx} \\ \\ \relax
9:1 & \begin{tabularx}{0.7\textwidth}{X} 掃羅不斷用威嚇兇悍的口氣向主的門徒說話。他去見大祭司， \end{tabularx} \\ \\ \relax
9:2 & \begin{tabularx}{0.7\textwidth}{X} 要求發信給大馬士革的各會堂，若是找著信奉這道的人，無論男女，都准他捆綁帶到耶路撒冷。 \end{tabularx} \\ \\ \relax
9:3 & \begin{tabularx}{0.7\textwidth}{X} 掃羅在途中，將到大馬士革的時候，忽然有一道光從天上下來，四面照射著他， \end{tabularx} \\ \\ \relax
9:4 & \begin{tabularx}{0.7\textwidth}{X} 他就仆倒在地，聽見有聲音對他說：「掃羅！掃羅！你為甚麼迫害我？」 \end{tabularx} \\ \\ \relax
9:5 & \begin{tabularx}{0.7\textwidth}{X} 他說：「主啊！你是誰？」主說：「我就是你所迫害的耶穌。 \end{tabularx} \\ \\ \relax
9:6 & \begin{tabularx}{0.7\textwidth}{X} 起來！進城去，你應該做的事，必有人告訴你。」 \end{tabularx} \\ \\ \relax
9:7 & \begin{tabularx}{0.7\textwidth}{X} 同行的人站在那裡，說不出話來，因為他們聽見聲音，卻看不見人。 \end{tabularx} \\ \\ \relax
9:8 & \begin{tabularx}{0.7\textwidth}{X} 掃羅從地上起來，睜開眼睛，竟不能看見甚麼。有人拉他的手，領他進了大馬士革。 \end{tabularx} \\ \\ \relax
9:9 & \begin{tabularx}{0.7\textwidth}{X} 他三天甚麼都看不見，也不吃也不喝。 \end{tabularx} \\ \\ \relax
9:10 & \begin{tabularx}{0.7\textwidth}{X} 那時，在大馬士革有一個門徒，名叫亞拿尼亞。主在異象中對他說：「亞拿尼亞！」他說：「主啊，我在這裡。」 \end{tabularx} \\ \\ \relax
9:11 & \begin{tabularx}{0.7\textwidth}{X} 主對他說：「起來！往那叫直街的路去，在猶大的家裡，去找一個大數人，名叫掃羅；他正在禱告， \end{tabularx} \\ \\ \relax
9:12 & \begin{tabularx}{0.7\textwidth}{X} 在異象中看見了一個人，名叫亞拿尼亞，進來為他按手，讓他能再看得見。」 \end{tabularx} \\ \\ \relax
9:13 & \begin{tabularx}{0.7\textwidth}{X} 亞拿尼亞回答：「主啊，我聽見許多人講到這個人，說他怎樣在耶路撒冷多多苦待你的聖徒， \end{tabularx} \\ \\ \relax
9:14 & \begin{tabularx}{0.7\textwidth}{X} 並且他在這裡有從祭司長得來的權柄，要捆綁一切求告你名的人。」 \end{tabularx} \\ \\ \relax
9:15 & \begin{tabularx}{0.7\textwidth}{X} 主對他說：「你只管去。他是我所揀選的器皿，要在外邦人、君王和以色列人面前宣揚我的名。 \end{tabularx} \\ \\ \relax
9:16 & \begin{tabularx}{0.7\textwidth}{X} 我也要指示他，為我的名必須受許多的苦難。」 \end{tabularx} \\ \\ \relax
9:17 & \begin{tabularx}{0.7\textwidth}{X} 亞拿尼亞就去了，進入那家，把手按在掃羅身上，說：「掃羅弟兄，在你來的路上向你顯現的主，就是耶穌，打發我來，叫你能再看得見，又被聖靈充滿。」 \end{tabularx} \\ \\ \relax
9:18 & \begin{tabularx}{0.7\textwidth}{X} 掃羅的眼睛上立刻好像有鱗一般的東西掉下來，他就能再看得見，於是他起來，受了洗， \end{tabularx} \\ \\ \relax
9:19 & \begin{tabularx}{0.7\textwidth}{X} 吃過飯體力就恢復了。 \end{tabularx} \\ \\ \relax
9:20 & \begin{tabularx}{0.7\textwidth}{X} 立刻在各會堂裡傳揚耶穌，說他是神的兒子。 \end{tabularx} \\ \\ \relax
9:21 & \begin{tabularx}{0.7\textwidth}{X} 凡聽見的人都很驚奇，說：「在耶路撒冷殘害求告這名的不就是這個人嗎？他不是到這裡來要捆綁他們，帶到祭司長那裡去嗎？」 \end{tabularx} \\ \\ \relax
9:22 & \begin{tabularx}{0.7\textwidth}{X} 但掃羅越發有能力，駁倒住在大馬士革的猶太人，證明耶穌是基督。 \end{tabularx} \\ \\ \relax
9:23 & \begin{tabularx}{0.7\textwidth}{X} 過了好些日子，猶太人商議要殺掃羅， \end{tabularx} \\ \\ \relax
9:24 & \begin{tabularx}{0.7\textwidth}{X} 但他們的計謀被掃羅知道了。他們晝夜在城門守候著要殺他。 \end{tabularx} \\ \\ \relax
9:25 & \begin{tabularx}{0.7\textwidth}{X} 他的門徒就在夜間用筐子把他從城牆上縋了下去。 \end{tabularx} \\ \\ \relax
9:26 & \begin{tabularx}{0.7\textwidth}{X} 掃羅到了耶路撒冷，想與門徒結交，大家卻都怕他，不信他是門徒。 \end{tabularx} \\ \\ \relax
9:27 & \begin{tabularx}{0.7\textwidth}{X} 只有巴拿巴接待他，領他去見使徒，把他在路上怎麼看見主，主怎麼向他說話，他在大馬士革怎麼奉耶穌的名放膽傳道，都述說出來。 \end{tabularx} \\ \\ \relax
9:28 & \begin{tabularx}{0.7\textwidth}{X} 於是掃羅在耶路撒冷同門徒出入來往，奉主的名放膽傳道， \end{tabularx} \\ \\ \relax
9:29 & \begin{tabularx}{0.7\textwidth}{X} 並和說希臘話的猶太人講論辯駁，他們卻想法子要殺他。 \end{tabularx} \\ \\ \relax
9:30 & \begin{tabularx}{0.7\textwidth}{X} 弟兄們知道了，就帶他下凱撒利亞，送他往大數去。 \end{tabularx} \\ \\ \relax
9:31 & \begin{tabularx}{0.7\textwidth}{X} 那時，猶太、加利利、撒瑪利亞各處的教會都得平安，建立起來，凡事敬畏主，蒙聖靈的安慰，人數逐漸增多。 \end{tabularx} \\ \\ \relax
9:32 & \begin{tabularx}{0.7\textwidth}{X} 彼得在眾信徒中到處奔波的時候，也到了住在呂大的聖徒那裡。 \end{tabularx} \\ \\ \relax
9:33 & \begin{tabularx}{0.7\textwidth}{X} 他在那裡遇見一個人，名叫以尼雅，得了癱瘓，在褥子上躺了八年。 \end{tabularx} \\ \\ \relax
9:34 & \begin{tabularx}{0.7\textwidth}{X} 彼得對他說：「以尼雅，耶穌基督醫好你了，起來！整理你的褥子吧。」他立刻就起來了。 \end{tabularx} \\ \\ \relax
9:35 & \begin{tabularx}{0.7\textwidth}{X} 凡住呂大和沙崙的人都看見了他，就歸向主。 \end{tabularx} \\ \\ \relax
9:36 & \begin{tabularx}{0.7\textwidth}{X} 在約帕有一個女門徒，名叫大比大，翻出來的意思是多加；她廣行善事，多施賙濟。 \end{tabularx} \\ \\ \relax
9:37 & \begin{tabularx}{0.7\textwidth}{X} 當時，她患病死了，有人把她清洗後，停在樓上。 \end{tabularx} \\ \\ \relax
9:38 & \begin{tabularx}{0.7\textwidth}{X} 呂大原與約帕相近；門徒聽見彼得在那裡，就派兩個人去見他，央求他說：「請快到我們那裡去，不要耽延。」 \end{tabularx} \\ \\ \relax
9:39 & \begin{tabularx}{0.7\textwidth}{X} 彼得就起身和他們同去。他到了，就有人領他上樓。眾寡婦都站在彼得旁邊哭，拿多加與她們同在時所做的內衣外衣給他看。 \end{tabularx} \\ \\ \relax
9:40 & \begin{tabularx}{0.7\textwidth}{X} 彼得叫她們都出去，然後跪下禱告，轉身對著屍體說：「大比大，起來！」她就睜開眼睛，看見彼得，就坐了起來。 \end{tabularx} \\ \\ \relax
9:41 & \begin{tabularx}{0.7\textwidth}{X} 彼得伸手扶她起來，叫那些聖徒和寡婦都進來，把多加活活地交給他們。 \end{tabularx} \\ \\ \relax
9:42 & \begin{tabularx}{0.7\textwidth}{X} 這事傳遍了約帕，就有許多人信了主。 \end{tabularx} \\ \\ \relax
9:43 & \begin{tabularx}{0.7\textwidth}{X} 此後，彼得在約帕一個皮革匠西門的家裡住了好些日子。 \end{tabularx} \\ \\ \relax
10:1 & \begin{tabularx}{0.7\textwidth}{X} 在凱撒利亞有一個人名叫哥尼流，是意大利營的百夫長。 \end{tabularx} \\ \\ \relax
10:2 & \begin{tabularx}{0.7\textwidth}{X} 他是個虔誠人，他和全家都敬畏神。他多多賙濟百姓，常常向神禱告。 \end{tabularx} \\ \\ \relax
10:3 & \begin{tabularx}{0.7\textwidth}{X} 有一天，約在下午三點鐘，他在異象中清楚看見神的一個使者進來，到他那裡，對他說：「哥尼流。」 \end{tabularx} \\ \\ \relax
10:4 & \begin{tabularx}{0.7\textwidth}{X} 哥尼流定睛看他，驚惶地說：「主啊，甚麼事？」天使對他說：「你的禱告和你的賙濟已達到神面前，蒙記念了。 \end{tabularx} \\ \\ \relax
10:5 & \begin{tabularx}{0.7\textwidth}{X} 現在你要派人往約帕去，請一位稱為彼得的西門來。 \end{tabularx} \\ \\ \relax
10:6 & \begin{tabularx}{0.7\textwidth}{X} 他住在一個皮革匠西門的家裡，房子就在海邊。」 \end{tabularx} \\ \\ \relax
10:7 & \begin{tabularx}{0.7\textwidth}{X} 向他說話的天使離開後，哥尼流叫了兩個僕人和常伺候他的一個虔誠的兵來， \end{tabularx} \\ \\ \relax
10:8 & \begin{tabularx}{0.7\textwidth}{X} 把一切的事都講給他們聽，然後就派他們往約帕去。 \end{tabularx} \\ \\ \relax
10:9 & \begin{tabularx}{0.7\textwidth}{X} 第二天，他們走路將近那城，約在正午，彼得上房頂去禱告。 \end{tabularx} \\ \\ \relax
10:10 & \begin{tabularx}{0.7\textwidth}{X} 他覺得餓了，想要吃。那家的人正預備飯的時候，彼得魂遊象外， \end{tabularx} \\ \\ \relax
10:11 & \begin{tabularx}{0.7\textwidth}{X} 看見天開了，有一塊好像大布的東西降下，四角吊著縋在地上， \end{tabularx} \\ \\ \relax
10:12 & \begin{tabularx}{0.7\textwidth}{X} 裡面有地上各樣四腳的走獸、爬蟲和天上的飛鳥。 \end{tabularx} \\ \\ \relax
10:13 & \begin{tabularx}{0.7\textwidth}{X} 又有聲音對他說：「彼得，起來！宰了吃。」 \end{tabularx} \\ \\ \relax
10:14 & \begin{tabularx}{0.7\textwidth}{X} 彼得卻說：「主啊，絕對不可！凡污俗和不潔淨的東西，我從來沒有吃過。」 \end{tabularx} \\ \\ \relax
10:15 & \begin{tabularx}{0.7\textwidth}{X} 第二次有聲音再對他說：「神所潔淨的，你不可當作污俗的。」 \end{tabularx} \\ \\ \relax
10:16 & \begin{tabularx}{0.7\textwidth}{X} 這樣一連三次，那東西隨即收回天上去了。 \end{tabularx} \\ \\ \relax
10:17 & \begin{tabularx}{0.7\textwidth}{X} 正當彼得心裡困惑，不知所看見的異象是甚麼意思時，哥尼流所差來的人已經找到了西門的家，站在門外， \end{tabularx} \\ \\ \relax
10:18 & \begin{tabularx}{0.7\textwidth}{X} 喊著問有沒有一位稱為彼得的西門住在這裡。 \end{tabularx} \\ \\ \relax
10:19 & \begin{tabularx}{0.7\textwidth}{X} 彼得還在思考那異象的時候，聖靈對他說：「有三個人來找你。 \end{tabularx} \\ \\ \relax
10:20 & \begin{tabularx}{0.7\textwidth}{X} 起來，下去，跟他們同去，不要疑惑，因為是我差他們來的。」 \end{tabularx} \\ \\ \relax
10:21 & \begin{tabularx}{0.7\textwidth}{X} 於是彼得下去見那些人，說：「我就是你們要找的人，你們是為了甚麼緣故在這裡？」 \end{tabularx} \\ \\ \relax
10:22 & \begin{tabularx}{0.7\textwidth}{X} 他們說：「百夫長哥尼流是個義人，敬畏神，為猶太全民族所稱讚。他蒙一位聖天使指示，叫他請你到他家裡去，要聽你講話。」 \end{tabularx} \\ \\ \relax
10:23 & \begin{tabularx}{0.7\textwidth}{X} 彼得就請他們進去住宿。次日，他起身和他們同去，還有約帕的幾個弟兄跟他一起去。 \end{tabularx} \\ \\ \relax
10:24 & \begin{tabularx}{0.7\textwidth}{X} 又次日，他進入凱撒利亞，哥尼流已經請了他的親朋好友在等候他們。 \end{tabularx} \\ \\ \relax
10:25 & \begin{tabularx}{0.7\textwidth}{X} 彼得一進去，哥尼流就迎接他，俯伏在他腳前拜他。 \end{tabularx} \\ \\ \relax
10:26 & \begin{tabularx}{0.7\textwidth}{X} 但是彼得拉他起來，說：「你起來，我自己也不過是人。」 \end{tabularx} \\ \\ \relax
10:27 & \begin{tabularx}{0.7\textwidth}{X} 彼得和他一邊說話一邊進去，見有好些人聚集， \end{tabularx} \\ \\ \relax
10:28 & \begin{tabularx}{0.7\textwidth}{X} 就對他們說：「你們知道，猶太人和別國的人結交來往本是不合規矩的，但神已經指示我，無論甚麼人都不可看作污俗或不潔淨的。 \end{tabularx} \\ \\ \relax
10:29 & \begin{tabularx}{0.7\textwidth}{X} 所以，我一被邀請，沒有推辭就來了。現在請問，你們為甚麼叫我來呢？」 \end{tabularx} \\ \\ \relax
10:30 & \begin{tabularx}{0.7\textwidth}{X} 哥尼流說：「四天前，這個時候，我在家中守著下午三點鐘的禱告，忽然有一個人穿著明亮的衣裳站在我面前， \end{tabularx} \\ \\ \relax
10:31 & \begin{tabularx}{0.7\textwidth}{X} 說：『哥尼流，你的禱告已蒙垂聽，你的賙濟在神面前已蒙記念了。 \end{tabularx} \\ \\ \relax
10:32 & \begin{tabularx}{0.7\textwidth}{X} 你要派人往約帕去，請那稱為彼得的西門來，他住在海邊一個皮革匠西門的家裡。』 \end{tabularx} \\ \\ \relax
10:33 & \begin{tabularx}{0.7\textwidth}{X} 所以我立刻派人去請你。你來了真好。現在我們都在神面前，要聽主吩咐你的一切話。」 \end{tabularx} \\ \\ \relax
10:34 & \begin{tabularx}{0.7\textwidth}{X} 彼得開口說：「我真的看出神是不偏待人的。 \end{tabularx} \\ \\ \relax
10:35 & \begin{tabularx}{0.7\textwidth}{X} 不但如此，在各國中那敬畏他而行義的人都為他所悅納。 \end{tabularx} \\ \\ \relax
10:36 & \begin{tabularx}{0.7\textwidth}{X} 神藉著耶穌基督—他是萬有的主—傳和平的福音，把這道傳給以色列人。 \end{tabularx} \\ \\ \relax
10:37 & \begin{tabularx}{0.7\textwidth}{X} 這話在約翰傳揚洗禮以後，從加利利起，傳遍了猶太。 \end{tabularx} \\ \\ \relax
10:38 & \begin{tabularx}{0.7\textwidth}{X} 神怎樣以聖靈和能力膏了拿撒勒人耶穌，這都是你們知道的。他到處奔波，行善事，醫好凡被魔鬼壓制的人，因為神與他同在。 \end{tabularx} \\ \\ \relax
10:39 & \begin{tabularx}{0.7\textwidth}{X} 他在猶太人之地和耶路撒冷所行的一切事，有我們作見證人。他們竟把他掛在木頭上殺了。 \end{tabularx} \\ \\ \relax
10:40 & \begin{tabularx}{0.7\textwidth}{X} 第三天，神使他復活，使他顯現出來； \end{tabularx} \\ \\ \relax
10:41 & \begin{tabularx}{0.7\textwidth}{X} 不是顯現給所有的人看，而是顯現給神預先所揀選為他作見證的人看，就是我們這些在他從死人中復活以後和他同吃同喝的人。 \end{tabularx} \\ \\ \relax
10:42 & \begin{tabularx}{0.7\textwidth}{X} 他吩咐我們傳道給眾人，證明他是神所立定，要作審判活人、死人的審判者。 \end{tabularx} \\ \\ \relax
10:43 & \begin{tabularx}{0.7\textwidth}{X} 眾先知也為這人作見證：凡信他的人，必藉著他的名得蒙赦罪。」 \end{tabularx} \\ \\ \relax
10:44 & \begin{tabularx}{0.7\textwidth}{X} 彼得還在說這些話的時候，聖靈降在一切聽道的人身上。 \end{tabularx} \\ \\ \relax
10:45 & \begin{tabularx}{0.7\textwidth}{X} 那些奉割禮的信徒和彼得同來，見聖靈的恩賜也澆在外邦人身上，就都驚奇； \end{tabularx} \\ \\ \relax
10:46 & \begin{tabularx}{0.7\textwidth}{X} 因聽見他們說方言，稱讚神為大。於是彼得回答： \end{tabularx} \\ \\ \relax
10:47 & \begin{tabularx}{0.7\textwidth}{X} 「這些人既受了聖靈，跟我們一樣，誰能阻止用水給他們施洗呢？」 \end{tabularx} \\ \\ \relax
10:48 & \begin{tabularx}{0.7\textwidth}{X} 他就吩咐奉耶穌基督的名給他們施洗。於是他們請彼得住了幾天。 \end{tabularx} \\ \\ \relax
11:1 & \begin{tabularx}{0.7\textwidth}{X} 使徒和在猶太的眾弟兄聽到外邦人也領受了神的道。 \end{tabularx} \\ \\ \relax
11:2 & \begin{tabularx}{0.7\textwidth}{X} 等到彼得上了耶路撒冷，那些奉割禮的信徒和他爭辯， \end{tabularx} \\ \\ \relax
11:3 & \begin{tabularx}{0.7\textwidth}{X} 說：「你竟進入未受割禮之人當中，和他們一同吃飯！」 \end{tabularx} \\ \\ \relax
11:4 & \begin{tabularx}{0.7\textwidth}{X} 彼得就開始把這事逐一向他們解釋，說： \end{tabularx} \\ \\ \relax
11:5 & \begin{tabularx}{0.7\textwidth}{X} 「我在約帕城裡禱告的時候，魂遊象外，看見異象，有一塊好像大布的東西降下，四角吊著從天縋下，直來到我跟前。 \end{tabularx} \\ \\ \relax
11:6 & \begin{tabularx}{0.7\textwidth}{X} 我定睛觀看，見內中有地上四腳的牲畜、野獸、爬蟲和天上的飛鳥。 \end{tabularx} \\ \\ \relax
11:7 & \begin{tabularx}{0.7\textwidth}{X} 我還聽見有聲音對我說：『彼得，起來！宰了吃。』 \end{tabularx} \\ \\ \relax
11:8 & \begin{tabularx}{0.7\textwidth}{X} 我說：『主啊，絕對不可！凡污俗或不潔淨的東西從來沒有進過我的口。』 \end{tabularx} \\ \\ \relax
11:9 & \begin{tabularx}{0.7\textwidth}{X} 第二次，有聲音從天上回答：『神所潔淨的，你不可當作污俗的。』 \end{tabularx} \\ \\ \relax
11:10 & \begin{tabularx}{0.7\textwidth}{X} 這樣一連三次，然後一切就都收回天上去了。 \end{tabularx} \\ \\ \relax
11:11 & \begin{tabularx}{0.7\textwidth}{X} 正當那時，有三個從凱撒利亞差來見我的人，站在我們所住的屋子門前。 \end{tabularx} \\ \\ \relax
11:12 & \begin{tabularx}{0.7\textwidth}{X} 聖靈吩咐我和他們同去，不要疑惑，還有這六位弟兄也跟我一起去，我們進了那人的家。 \end{tabularx} \\ \\ \relax
11:13 & \begin{tabularx}{0.7\textwidth}{X} 那人就告訴我們，他如何看見一位天使站在他家裡，說：『你派人往約帕去，請那稱為彼得的西門來， \end{tabularx} \\ \\ \relax
11:14 & \begin{tabularx}{0.7\textwidth}{X} 他有話要告訴你，因這些話你和你的全家都可以得救。』 \end{tabularx} \\ \\ \relax
11:15 & \begin{tabularx}{0.7\textwidth}{X} 我一開始講話，聖靈就降在他們身上，正像當初降在我們身上一樣。 \end{tabularx} \\ \\ \relax
11:16 & \begin{tabularx}{0.7\textwidth}{X} 我就想起主的話如何說：『約翰用水施洗，但你們要在聖靈裡受洗。』 \end{tabularx} \\ \\ \relax
11:17 & \begin{tabularx}{0.7\textwidth}{X} 既然神給他們恩賜，像在我們信主耶穌基督的時候給了我們一樣，我是誰，能攔阻神嗎？」 \end{tabularx} \\ \\ \relax
11:18 & \begin{tabularx}{0.7\textwidth}{X} 眾人聽見這些話，就不說話了，只歸榮耀給神，說：「這樣看來，神也賜恩給外邦人，使他們悔改得生命了。」 \end{tabularx} \\ \\ \relax
11:19 & \begin{tabularx}{0.7\textwidth}{X} 那些因司提反的事遭患難而四處分散的門徒，直走到腓尼基、塞浦路斯和安提阿。他們不向別人講道，只向猶太人講。 \end{tabularx} \\ \\ \relax
11:20 & \begin{tabularx}{0.7\textwidth}{X} 但內中有塞浦路斯和古利奈人，他們到了安提阿也向希臘人傳講主耶穌的福音。 \end{tabularx} \\ \\ \relax
11:21 & \begin{tabularx}{0.7\textwidth}{X} 主的手與他們同在，信而歸主的人數很多。 \end{tabularx} \\ \\ \relax
11:22 & \begin{tabularx}{0.7\textwidth}{X} 這風聲傳到耶路撒冷教會的人耳中，他們就打發巴拿巴到安提阿去。 \end{tabularx} \\ \\ \relax
11:23 & \begin{tabularx}{0.7\textwidth}{X} 他到了那裡，看見神所賜的恩就歡喜，勸勉眾人要立定心志，恆久靠主。 \end{tabularx} \\ \\ \relax
11:24 & \begin{tabularx}{0.7\textwidth}{X} 這巴拿巴原是個好人，滿有聖靈和信心，於是有許多人歸服了主。 \end{tabularx} \\ \\ \relax
11:25 & \begin{tabularx}{0.7\textwidth}{X} 他又往大數去找掃羅， \end{tabularx} \\ \\ \relax
11:26 & \begin{tabularx}{0.7\textwidth}{X} 找著了，就帶他到安提阿去。他們足有一年和教會一同聚集，教導了許多人。門徒稱為「基督徒」是從安提阿開始的。 \end{tabularx} \\ \\ \relax
11:27 & \begin{tabularx}{0.7\textwidth}{X} 當那些日子，有幾位先知從耶路撒冷下到安提阿。 \end{tabularx} \\ \\ \relax
11:28 & \begin{tabularx}{0.7\textwidth}{X} 內中有一位，名叫亞迦布，站起來，藉著聖靈指示普天下將有大饑荒；這事在克勞第年間果然實現了。 \end{tabularx} \\ \\ \relax
11:29 & \begin{tabularx}{0.7\textwidth}{X} 於是門徒決定，照各人的力量捐錢，送去供給住在猶太的弟兄。 \end{tabularx} \\ \\ \relax
11:30 & \begin{tabularx}{0.7\textwidth}{X} 他們就這樣做了，託巴拿巴和掃羅的手送到眾長老那裡。 \end{tabularx} \\ \\ \relax
12:1 & \begin{tabularx}{0.7\textwidth}{X} 約在那時候，希律王下手苦待教會中的一些人， \end{tabularx} \\ \\ \relax
12:2 & \begin{tabularx}{0.7\textwidth}{X} 用刀殺了約翰的哥哥雅各。 \end{tabularx} \\ \\ \relax
12:3 & \begin{tabularx}{0.7\textwidth}{X} 他見猶太人喜歡這事，也去拿住彼得。那時候正是除酵節期間。 \end{tabularx} \\ \\ \relax
12:4 & \begin{tabularx}{0.7\textwidth}{X} 希律捉了彼得，押在監裡，交給四班士兵看守，每班四個人，企圖要在逾越節後把他提出來，當著百姓辦他。 \end{tabularx} \\ \\ \relax
12:5 & \begin{tabularx}{0.7\textwidth}{X} 於是彼得被囚在監裡，教會卻為他切切禱告神。 \end{tabularx} \\ \\ \relax
12:6 & \begin{tabularx}{0.7\textwidth}{X} 希律將要提他出來的前一夜，彼得被兩條鐵鏈鎖著，睡在兩個士兵當中；門前還有警衛看守。 \end{tabularx} \\ \\ \relax
12:7 & \begin{tabularx}{0.7\textwidth}{X} 忽然，有主的一個使者顯現，牢房裡有光照耀；天使拍彼得的肋旁，叫醒了他，說：「快起來！」鐵鏈就從他手上脫落下來。 \end{tabularx} \\ \\ \relax
12:8 & \begin{tabularx}{0.7\textwidth}{X} 天使對他說：「束上腰帶，穿上鞋子。」他就照著做了。天使又對他說：「披上外衣，跟我來。」 \end{tabularx} \\ \\ \relax
12:9 & \begin{tabularx}{0.7\textwidth}{X} 彼得就出來跟著他走，不知道天使所做是真的，以為見了異象。 \end{tabularx} \\ \\ \relax
12:10 & \begin{tabularx}{0.7\textwidth}{X} 他們經過了第一層和第二層監牢，就來到往城內的鐵門，那門就自動給他們開了。他們出來，走過一條街，忽然天使離開他去了。 \end{tabularx} \\ \\ \relax
12:11 & \begin{tabularx}{0.7\textwidth}{X} 彼得清醒過來，說：「現在我真知道主差遣他的使者，救我脫離希律的手，和猶太人所期待的一切。」 \end{tabularx} \\ \\ \relax
12:12 & \begin{tabularx}{0.7\textwidth}{X} 他明白了，就到那稱為馬可的約翰的母親馬利亞家去，在那裡已有好些人聚集禱告。 \end{tabularx} \\ \\ \relax
12:13 & \begin{tabularx}{0.7\textwidth}{X} 彼得敲外門時，有一個使女，名叫羅大，出來應門， \end{tabularx} \\ \\ \relax
12:14 & \begin{tabularx}{0.7\textwidth}{X} 認出是彼得的聲音，歡喜得顧不了開門，就跑進去報信，說彼得站在門外。 \end{tabularx} \\ \\ \relax
12:15 & \begin{tabularx}{0.7\textwidth}{X} 他們對她說：「你瘋了！」使女堅持真有其事。他們說：「那是他的天使。」 \end{tabularx} \\ \\ \relax
12:16 & \begin{tabularx}{0.7\textwidth}{X} 彼得不停地敲門；他們開了門，一見是他，就很驚奇。 \end{tabularx} \\ \\ \relax
12:17 & \begin{tabularx}{0.7\textwidth}{X} 彼得做個手勢，要他們不作聲，就告訴他們主怎樣領他出監；又說：「你們要把這些事告訴雅各和眾弟兄。」然後，他離開往別處去了。 \end{tabularx} \\ \\ \relax
12:18 & \begin{tabularx}{0.7\textwidth}{X} 到了天亮，士兵中起了不少騷動，不知道彼得到哪裡去了。 \end{tabularx} \\ \\ \relax
12:19 & \begin{tabularx}{0.7\textwidth}{X} 希律找他，找不著，就審問警衛，下令帶走他們處死。後來希律離開猶太，下凱撒利亞去，住在那裡。 \end{tabularx} \\ \\ \relax
12:20 & \begin{tabularx}{0.7\textwidth}{X} 希律向推羅和西頓的人發怒。他們那一帶地方是從王的土地供應糧食的，因此就託了王的內侍大臣伯拉斯都的情，一心來求和。 \end{tabularx} \\ \\ \relax
12:21 & \begin{tabularx}{0.7\textwidth}{X} 希律在所定的日子，穿上朝服，坐在位上，對他們演講。 \end{tabularx} \\ \\ \relax
12:22 & \begin{tabularx}{0.7\textwidth}{X} 民眾一直喊著：「這是神明的聲音，不是人的聲音。」 \end{tabularx} \\ \\ \relax
12:23 & \begin{tabularx}{0.7\textwidth}{X} 希律不歸榮耀給神，所以主的使者立刻擊打他，他被蟲咬，就斷了氣。 \end{tabularx} \\ \\ \relax
12:24 & \begin{tabularx}{0.7\textwidth}{X} 神的道日見興旺，越發廣傳。 \end{tabularx} \\ \\ \relax
12:25 & \begin{tabularx}{0.7\textwidth}{X} 巴拿巴和掃羅完成了供給的事，就回到耶路撒冷，帶著稱為馬可的約翰同去。 \end{tabularx} \\ \\ \relax
13:1 & \begin{tabularx}{0.7\textwidth}{X} 在安提阿的教會中，有幾位先知和教師，就是巴拿巴和稱為尼結的西面、古利奈人路求，與希律分封王一起長大的馬念，和掃羅。 \end{tabularx} \\ \\ \relax
13:2 & \begin{tabularx}{0.7\textwidth}{X} 他們在事奉主和禁食的時候，聖靈說：「要為我分派巴拿巴和掃羅去做我召他們做的工作。」 \end{tabularx} \\ \\ \relax
13:3 & \begin{tabularx}{0.7\textwidth}{X} 於是他們禁食禱告後，給巴拿巴和掃羅按手，然後派遣他們走了。 \end{tabularx} \\ \\ \relax
13:4 & \begin{tabularx}{0.7\textwidth}{X} 他們既蒙聖靈差遣，就下到西流基，從那裡坐船往塞浦路斯去， \end{tabularx} \\ \\ \relax
13:5 & \begin{tabularx}{0.7\textwidth}{X} 到了撒拉米，就在猶太人各會堂裡宣講神的道，也有約翰作他們的幫手。 \end{tabularx} \\ \\ \relax
13:6 & \begin{tabularx}{0.7\textwidth}{X} 他們走遍全島，直到帕弗，在那裡遇見一個術士—猶太人的假先知，名叫巴耶穌。 \end{tabularx} \\ \\ \relax
13:7 & \begin{tabularx}{0.7\textwidth}{X} 這人常和士求‧保羅省長在一起。士求‧保羅是個通達人，他請巴拿巴和掃羅來，要聽神的道。 \end{tabularx} \\ \\ \relax
13:8 & \begin{tabularx}{0.7\textwidth}{X} 只是術士以呂馬(他的名字翻出來就是行法術的意思)敵對使徒，設法使省長遠離這信仰。 \end{tabularx} \\ \\ \relax
13:9 & \begin{tabularx}{0.7\textwidth}{X} 掃羅，又名保羅，被聖靈充滿，定睛看他， \end{tabularx} \\ \\ \relax
13:10 & \begin{tabularx}{0.7\textwidth}{X} 說：「你這充滿各樣詭詐奸惡，魔鬼的兒子，一切正義的仇敵，你還不停止扭曲主的正道嗎？ \end{tabularx} \\ \\ \relax
13:11 & \begin{tabularx}{0.7\textwidth}{X} 現在你看，主的手臨到你身上，你會瞎眼，暫時看不見日光。」立刻迷濛和黑暗籠罩著他，他到處摸索，求人拉著手領他。 \end{tabularx} \\ \\ \relax
13:12 & \begin{tabularx}{0.7\textwidth}{X} 省長看見所發生的事就信了，因對主的教導感到驚奇。 \end{tabularx} \\ \\ \relax
13:13 & \begin{tabularx}{0.7\textwidth}{X} 保羅和他的同伴從帕弗開船，來到旁非利亞的別加，約翰卻離開他們，回耶路撒冷去了。 \end{tabularx} \\ \\ \relax
13:14 & \begin{tabularx}{0.7\textwidth}{X} 他們從別加往前行，來到彼西底的安提阿。在安息日，他們進了會堂就坐下。 \end{tabularx} \\ \\ \relax
13:15 & \begin{tabularx}{0.7\textwidth}{X} 在讀完了律法和先知的書，會堂主管們叫人過去，對他們說：「二位弟兄，你們若有甚麼勸勉眾人的話，請說。」 \end{tabularx} \\ \\ \relax
13:16 & \begin{tabularx}{0.7\textwidth}{X} 保羅就站起來，做個手勢，說：「諸位以色列人和一切敬畏神的人，請聽。 \end{tabularx} \\ \\ \relax
13:17 & \begin{tabularx}{0.7\textwidth}{X} 這以色列民的神揀選了我們的祖宗，當百姓寄居埃及的時候抬舉他們，用大能的手領他們從那地出來。 \end{tabularx} \\ \\ \relax
13:18 & \begin{tabularx}{0.7\textwidth}{X} 他在曠野容忍他們，約有四十年。 \end{tabularx} \\ \\ \relax
13:19 & \begin{tabularx}{0.7\textwidth}{X} 他消滅了迦南地七族的人後，把那地分給他們為業， \end{tabularx} \\ \\ \relax
13:20 & \begin{tabularx}{0.7\textwidth}{X} 約有四百五十年。此後，他給他們設立士師，直到撒母耳先知的時候。 \end{tabularx} \\ \\ \relax
13:21 & \begin{tabularx}{0.7\textwidth}{X} 從那時起，他們要求立一個王，神就將便雅憫支派中基士的兒子掃羅給他們作王，共四十年。 \end{tabularx} \\ \\ \relax
13:22 & \begin{tabularx}{0.7\textwidth}{X} 他廢了掃羅之後，就興起大衛作他們的王，又為他作見證說：『我尋得耶西的兒子大衛，他是合我心意的人，他要遵行我一切的旨意。』 \end{tabularx} \\ \\ \relax
13:23 & \begin{tabularx}{0.7\textwidth}{X} 從這人的後裔中，神已經照著所應許的為以色列人興起一位救主，就是耶穌。 \end{tabularx} \\ \\ \relax
13:24 & \begin{tabularx}{0.7\textwidth}{X} 在他沒有出來以前，約翰已向以色列全民宣講悔改的洗禮。 \end{tabularx} \\ \\ \relax
13:25 & \begin{tabularx}{0.7\textwidth}{X} 約翰快走完他的人生路程時，說：『你們以為我是誰？我不是；但是有一位在我以後來的，我就是解他腳上的鞋帶也不配。』 \end{tabularx} \\ \\ \relax
13:26 & \begin{tabularx}{0.7\textwidth}{X} 「諸位弟兄—亞伯拉罕的子孫和你們中間敬畏神的人哪，這救世的道是傳給我們的。 \end{tabularx} \\ \\ \relax
13:27 & \begin{tabularx}{0.7\textwidth}{X} 耶路撒冷的居民和他們的官長，因為不認識這基督，也不明白每安息日所讀的先知的書，把他定了死罪，正應驗了先知的預言。 \end{tabularx} \\ \\ \relax
13:28 & \begin{tabularx}{0.7\textwidth}{X} 雖然他們查不出他有該死的罪狀，還是要求彼拉多把他殺了。 \end{tabularx} \\ \\ \relax
13:29 & \begin{tabularx}{0.7\textwidth}{X} 他們既實現了經上指著他所記的一切話，就從木頭上把他取下來，放在墳墓裡。 \end{tabularx} \\ \\ \relax
13:30 & \begin{tabularx}{0.7\textwidth}{X} 神卻使他從死人中復活。 \end{tabularx} \\ \\ \relax
13:31 & \begin{tabularx}{0.7\textwidth}{X} 有許多日子，他向那些從加利利同他上耶路撒冷的人顯現，這些人如今在民間成為他的見證人。 \end{tabularx} \\ \\ \relax
13:32 & \begin{tabularx}{0.7\textwidth}{X} 我們報好信息給你們，就是那應許祖宗的話， \end{tabularx} \\ \\ \relax
13:33 & \begin{tabularx}{0.7\textwidth}{X} 神已經向我們這些作他們兒女的應驗，使耶穌復活了。正如《詩篇》第二篇上記著：『你是我的兒子，我今日生了你。』 \end{tabularx} \\ \\ \relax
13:34 & \begin{tabularx}{0.7\textwidth}{X} 論到神使他從死人中復活，不再歸於朽壞，他曾這樣說：『我必將所應許大衛那聖潔、可靠的恩典賜給你們。』 \end{tabularx} \\ \\ \relax
13:35 & \begin{tabularx}{0.7\textwidth}{X} 所以他也在另一篇說：『你必不讓你的聖者見朽壞。』 \end{tabularx} \\ \\ \relax
13:36 & \begin{tabularx}{0.7\textwidth}{X} 大衛在世的時候，遵行了神的旨意就長眠了，歸到他祖宗那裡，已見朽壞； \end{tabularx} \\ \\ \relax
13:37 & \begin{tabularx}{0.7\textwidth}{X} 惟獨神使他復活的那一位，他並未見朽壞。 \end{tabularx} \\ \\ \relax
13:38 & \begin{tabularx}{0.7\textwidth}{X} 所以弟兄們，你們當知道：赦罪的道是由這人傳給你們的， \end{tabularx} \\ \\ \relax
13:39 & \begin{tabularx}{0.7\textwidth}{X} 你們靠摩西的律法在不得稱義的一切事上，每一個信靠這位耶穌的都得稱義了。 \end{tabularx} \\ \\ \relax
13:40 & \begin{tabularx}{0.7\textwidth}{X} 所以，你們要小心，免得先知書上所說的臨到你們： \end{tabularx} \\ \\ \relax
13:41 & \begin{tabularx}{0.7\textwidth}{X} 『要觀看，你們這些藐視的人，要驚訝，要滅亡，因為在你們的日子，我行一件事，雖有人告訴你們，你們總是不信。』」 \end{tabularx} \\ \\ \relax
13:42 & \begin{tabularx}{0.7\textwidth}{X} 他們走出會堂的時候，眾人請他們在下一個安息日再講這些話給他們聽。 \end{tabularx} \\ \\ \relax
13:43 & \begin{tabularx}{0.7\textwidth}{X} 散會以後，有許多猶太人和敬虔的皈依猶太教的人跟從了保羅和巴拿巴。二人對他們講話，勸他們務要恆久倚靠神的恩典。 \end{tabularx} \\ \\ \relax
13:44 & \begin{tabularx}{0.7\textwidth}{X} 到下一個安息日，全城的人幾乎都聚集起來，要聽主的道。 \end{tabularx} \\ \\ \relax
13:45 & \begin{tabularx}{0.7\textwidth}{X} 但猶太人看見這麼多的人，就滿心嫉妒，辯駁保羅所說的話，並且毀謗他。 \end{tabularx} \\ \\ \relax
13:46 & \begin{tabularx}{0.7\textwidth}{X} 於是保羅和巴拿巴放膽說：「神的道本應先傳給你們；只因你們棄絕這道，斷定自己不配得永生，我們就轉向外邦人。 \end{tabularx} \\ \\ \relax
13:47 & \begin{tabularx}{0.7\textwidth}{X} 因為主曾這樣吩咐我們：『我已經立你作萬邦之光，使你施行我的救恩，直到地極。』」 \end{tabularx} \\ \\ \relax
13:48 & \begin{tabularx}{0.7\textwidth}{X} 外邦人聽見這話很歡喜，讚美主的道，凡被指定得永生的人都信了。 \end{tabularx} \\ \\ \relax
13:49 & \begin{tabularx}{0.7\textwidth}{X} 於是主的道傳遍了那一帶地方。 \end{tabularx} \\ \\ \relax
13:50 & \begin{tabularx}{0.7\textwidth}{X} 但猶太人挑唆虔敬尊貴的婦女和城內有名望的人，迫害保羅和巴拿巴，把他們趕出境外。 \end{tabularx} \\ \\ \relax
13:51 & \begin{tabularx}{0.7\textwidth}{X} 二人對著眾人跺掉腳上的塵土，然後往以哥念去了。 \end{tabularx} \\ \\ \relax
13:52 & \begin{tabularx}{0.7\textwidth}{X} 門徒滿心喜樂，又被聖靈充滿。 \end{tabularx} \\ \\ \relax
14:1 & \begin{tabularx}{0.7\textwidth}{X} 同樣的事也發生在以哥念。保羅和巴拿巴進了猶太人的會堂，在那裡講道，所以有很多猶太人和希臘人都信了。 \end{tabularx} \\ \\ \relax
14:2 & \begin{tabularx}{0.7\textwidth}{X} 但那不順從的猶太人煽動外邦人，使他們心裡仇恨弟兄。 \end{tabularx} \\ \\ \relax
14:3 & \begin{tabularx}{0.7\textwidth}{X} 二人在那裡住了好些日子，倚靠主放膽講道，主藉他們的手施行神蹟奇事，證明他恩惠的道。 \end{tabularx} \\ \\ \relax
14:4 & \begin{tabularx}{0.7\textwidth}{X} 城裡的眾人卻分裂了：有依附猶太人的，有依附使徒的。 \end{tabularx} \\ \\ \relax
14:5 & \begin{tabularx}{0.7\textwidth}{X} 那時，外邦人、猶太人和他們的官長，一齊擁上來，要凌辱使徒，用石頭打他們。 \end{tabularx} \\ \\ \relax
14:6 & \begin{tabularx}{0.7\textwidth}{X} 使徒知道了，就逃到呂高尼的路司得和特庇兩個城，以及周圍地方去， \end{tabularx} \\ \\ \relax
14:7 & \begin{tabularx}{0.7\textwidth}{X} 在那裡繼續傳福音。 \end{tabularx} \\ \\ \relax
14:8 & \begin{tabularx}{0.7\textwidth}{X} 路司得城裡有一個兩腳無力的人，他從母腹裡就是瘸腿的，老是坐著，從來沒有走過。 \end{tabularx} \\ \\ \relax
14:9 & \begin{tabularx}{0.7\textwidth}{X} 他聽保羅講道；保羅定睛看他，見他有信心，可得痊癒， \end{tabularx} \\ \\ \relax
14:10 & \begin{tabularx}{0.7\textwidth}{X} 就大聲說：「起來！兩腳站直。」那人就跳起來，開始行走。 \end{tabularx} \\ \\ \relax
14:11 & \begin{tabularx}{0.7\textwidth}{X} 眾人看見保羅所做的事，就用呂高尼話大聲說：「有神明藉著人形降臨在我們中間了。」 \end{tabularx} \\ \\ \relax
14:12 & \begin{tabularx}{0.7\textwidth}{X} 於是他們稱巴拿巴為宙斯，稱保羅為希耳米，因為他總是帶頭說話。 \end{tabularx} \\ \\ \relax
14:13 & \begin{tabularx}{0.7\textwidth}{X} 城外有宙斯廟的祭司牽著牛，拿著花環，來到門前，要同眾人一起獻祭。 \end{tabularx} \\ \\ \relax
14:14 & \begin{tabularx}{0.7\textwidth}{X} 巴拿巴和保羅二位使徒聽見，就撕開衣裳，跳進眾人中間，喊著： \end{tabularx} \\ \\ \relax
14:15 & \begin{tabularx}{0.7\textwidth}{X} 「諸位，為甚麼做這些事呢？我們也是人，性情和你們一樣。我們傳福音給你們，是要你們離棄這些虛妄的事，歸向那創造天、地、海和其中萬物的永生的神。 \end{tabularx} \\ \\ \relax
14:16 & \begin{tabularx}{0.7\textwidth}{X} 他在從前的世代，任憑萬國各行其道； \end{tabularx} \\ \\ \relax
14:17 & \begin{tabularx}{0.7\textwidth}{X} 然而他未嘗不為自己留下證據來，就如常行善事，從天降雨，賞賜豐年，使你們飲食飽足，滿心喜樂。」 \end{tabularx} \\ \\ \relax
14:18 & \begin{tabularx}{0.7\textwidth}{X} 二人說了這些話，總算攔住眾人不獻祭給他們。 \end{tabularx} \\ \\ \relax
14:19 & \begin{tabularx}{0.7\textwidth}{X} 但有些猶太人，從安提阿和以哥念來，挑唆眾人，並且用石頭打保羅，以為他死了，就把他拖到城外。 \end{tabularx} \\ \\ \relax
14:20 & \begin{tabularx}{0.7\textwidth}{X} 當門徒圍著他的時候，他站了起來，走進城去。第二天，保羅同巴拿巴往特庇去。 \end{tabularx} \\ \\ \relax
14:21 & \begin{tabularx}{0.7\textwidth}{X} 保羅和巴拿巴對那城裡的人傳了福音，使好些人成為門徒後，又回路司得、以哥念、安提阿去， \end{tabularx} \\ \\ \relax
14:22 & \begin{tabularx}{0.7\textwidth}{X} 堅固門徒的心，勸他們持守他們的信仰，說：「我們進入神的國，必須經歷許多艱難。」 \end{tabularx} \\ \\ \relax
14:23 & \begin{tabularx}{0.7\textwidth}{X} 二人在各教會中選立了長老，禁食禱告後，把他們交託給他們所信的主。 \end{tabularx} \\ \\ \relax
14:24 & \begin{tabularx}{0.7\textwidth}{X} 二人經過彼西底來到旁非利亞， \end{tabularx} \\ \\ \relax
14:25 & \begin{tabularx}{0.7\textwidth}{X} 在別加講了道，就下亞大利去， \end{tabularx} \\ \\ \relax
14:26 & \begin{tabularx}{0.7\textwidth}{X} 從那裡坐船回安提阿去。當初，眾人就在這地方，把他們交託在神的恩典中，要完成現在所做的工。 \end{tabularx} \\ \\ \relax
14:27 & \begin{tabularx}{0.7\textwidth}{X} 他們一到那裡，就聚集了會眾，述說神藉他們所行的一切事，並且神怎樣為外邦人開了信道的門。 \end{tabularx} \\ \\ \relax
14:28 & \begin{tabularx}{0.7\textwidth}{X} 二人在那裡同門徒住了一段日子。 \end{tabularx} \\ \\ \relax
15:1 & \begin{tabularx}{0.7\textwidth}{X} 有幾個人從猶太下來，教導弟兄們說：「你們若不按照摩西的規矩受割禮，不能得救。」 \end{tabularx} \\ \\ \relax
15:2 & \begin{tabularx}{0.7\textwidth}{X} 保羅和巴拿巴跟他們發生了激烈的爭執和辯論；大家就決定指派保羅、巴拿巴和本會的幾個人，為所辯論的事上耶路撒冷去見使徒和長老。 \end{tabularx} \\ \\ \relax
15:3 & \begin{tabularx}{0.7\textwidth}{X} 於是教會為他們送行。他們經過腓尼基、撒瑪利亞，沿途敘說外邦人歸主的事，使眾弟兄都非常歡喜。 \end{tabularx} \\ \\ \relax
15:4 & \begin{tabularx}{0.7\textwidth}{X} 他們到了耶路撒冷，教會、使徒和長老都接待他們，他們就述說神同他們所做的一切事。 \end{tabularx} \\ \\ \relax
15:5 & \begin{tabularx}{0.7\textwidth}{X} 惟有幾個法利賽派的信徒起來，說：「必須給外邦人行割禮，吩咐他們遵守摩西的律法。」 \end{tabularx} \\ \\ \relax
15:6 & \begin{tabularx}{0.7\textwidth}{X} 使徒和長老聚集商議這事。 \end{tabularx} \\ \\ \relax
15:7 & \begin{tabularx}{0.7\textwidth}{X} 辯論了許久後，彼得站起來，對他們說：「諸位弟兄，你們知道神早已在你們中間揀選了我，讓外邦人從我口中得聽福音之道，而且相信。 \end{tabularx} \\ \\ \relax
15:8 & \begin{tabularx}{0.7\textwidth}{X} 知道人心的神也為他們作了見證，賜聖靈給他們，正如給我們一樣； \end{tabularx} \\ \\ \relax
15:9 & \begin{tabularx}{0.7\textwidth}{X} 又藉著信潔淨了他們的心，他們和我們之間並沒有甚麼分別。 \end{tabularx} \\ \\ \relax
15:10 & \begin{tabularx}{0.7\textwidth}{X} 現在你們為甚麼試探神，要把我們祖宗和我們所不能負的軛放在門徒的頸項上呢？ \end{tabularx} \\ \\ \relax
15:11 & \begin{tabularx}{0.7\textwidth}{X} 相反地，我們相信，我們得救是因主耶穌的恩典，和他們一樣。」 \end{tabularx} \\ \\ \relax
15:12 & \begin{tabularx}{0.7\textwidth}{X} 眾人都默默無聲，聽巴拿巴和保羅述說神藉著他們在外邦人中所行的神蹟和奇事。 \end{tabularx} \\ \\ \relax
15:13 & \begin{tabularx}{0.7\textwidth}{X} 他們講完了，雅各回答說：「諸位弟兄，請聽我說。 \end{tabularx} \\ \\ \relax
15:14 & \begin{tabularx}{0.7\textwidth}{X} 剛才西門述說神當初怎樣眷顧外邦人，從他們中間選取人民歸於自己的名下； \end{tabularx} \\ \\ \relax
15:15 & \begin{tabularx}{0.7\textwidth}{X} 眾先知的話也與這意思相符合。 \end{tabularx} \\ \\ \relax
15:16 & \begin{tabularx}{0.7\textwidth}{X} 正如經上所寫的：『此後，我要回來，重新修造大衛倒塌了的帳幕，從廢墟中重新修造，把它建立起來， \end{tabularx} \\ \\ \relax
15:17 & \begin{tabularx}{0.7\textwidth}{X} 使剩餘的人，就是凡稱我名的外邦人，都尋求主。 \end{tabularx} \\ \\ \relax
15:18 & \begin{tabularx}{0.7\textwidth}{X} 這話是自古以來顯明這些事的主說的。』 \end{tabularx} \\ \\ \relax
15:19 & \begin{tabularx}{0.7\textwidth}{X} 所以，我的意見是不可難為那歸向神的外邦人； \end{tabularx} \\ \\ \relax
15:20 & \begin{tabularx}{0.7\textwidth}{X} 但是要寫信吩咐他們禁戒偶像所玷污的東西、血和勒死的牲畜，禁戒淫亂。 \end{tabularx} \\ \\ \relax
15:21 & \begin{tabularx}{0.7\textwidth}{X} 因為歷代以來，摩西的書在各城都有人宣講，每逢安息日，也在會堂裡誦讀。」 \end{tabularx} \\ \\ \relax
15:22 & \begin{tabularx}{0.7\textwidth}{X} 那時，使徒、長老和全教會認為應從他們中間揀選人，差他們和保羅、巴拿巴一同到安提阿去，所揀選的就是稱為巴撒巴的猶大和西拉。這二人在弟兄中是領袖。 \end{tabularx} \\ \\ \relax
15:23 & \begin{tabularx}{0.7\textwidth}{X} 他們帶去的信說：「使徒和作長老的弟兄們向安提阿、敘利亞、基利家外邦眾弟兄問安。 \end{tabularx} \\ \\ \relax
15:24 & \begin{tabularx}{0.7\textwidth}{X} 我們聽說，有幾個人從我們這裡出去，用一些話騷擾你們，使你們的心困惑，其實我們並沒有吩咐他們。 \end{tabularx} \\ \\ \relax
15:25 & \begin{tabularx}{0.7\textwidth}{X} 我們認為，既然我們同心定意，就揀選幾個人，派他們同我們所親愛的巴拿巴和保羅到你們那裡去。 \end{tabularx} \\ \\ \relax
15:26 & \begin{tabularx}{0.7\textwidth}{X} 這二人曾為我主耶穌基督的名不顧自己的性命。 \end{tabularx} \\ \\ \relax
15:27 & \begin{tabularx}{0.7\textwidth}{X} 所以我們派猶大和西拉去，他們也會親口述說這些事。 \end{tabularx} \\ \\ \relax
15:28 & \begin{tabularx}{0.7\textwidth}{X} 因為聖靈和我們決定除了這幾件重要的事，不將別的重擔放在你們身上， \end{tabularx} \\ \\ \relax
15:29 & \begin{tabularx}{0.7\textwidth}{X} 就是禁戒偶像所玷污的東西、血和勒死的牲畜，禁戒淫亂。這幾件你們若能自己禁戒就好了。祝你們安康！」 \end{tabularx} \\ \\ \relax
15:30 & \begin{tabularx}{0.7\textwidth}{X} 他們既奉了差遣就下安提阿去，聚集會眾，把書信交給他們。 \end{tabularx} \\ \\ \relax
15:31 & \begin{tabularx}{0.7\textwidth}{X} 眾人念了，因為信上鼓勵的話而感到欣慰。 \end{tabularx} \\ \\ \relax
15:32 & \begin{tabularx}{0.7\textwidth}{X} 猶大和西拉自己也是先知，就用許多話勸勉弟兄，堅固他們。 \end{tabularx} \\ \\ \relax
15:33 & \begin{tabularx}{0.7\textwidth}{X} 二人住了些日子，弟兄們打發他們平平安安地回到差遣他們的人那裡去。 \end{tabularx} \\ \\ \relax
15:34 & \begin{tabularx}{0.7\textwidth}{X} 惟有西拉決定仍住在那裡。 \end{tabularx} \\ \\ \relax
15:35 & \begin{tabularx}{0.7\textwidth}{X} 但保羅和巴拿巴仍留在安提阿，和許多別的人一同教導，並傳揚主的道。 \end{tabularx} \\ \\ \relax
15:36 & \begin{tabularx}{0.7\textwidth}{X} 過了些日子，保羅對巴拿巴說：「讓我們回到從前宣揚主道的各城，看看弟兄們的情況如何。」 \end{tabularx} \\ \\ \relax
15:37 & \begin{tabularx}{0.7\textwidth}{X} 巴拿巴有意要帶稱為馬可的約翰同去； \end{tabularx} \\ \\ \relax
15:38 & \begin{tabularx}{0.7\textwidth}{X} 但保羅認為不宜帶他去，因為馬可從前在旁非利亞離開他們，不和他們一起工作。 \end{tabularx} \\ \\ \relax
15:39 & \begin{tabularx}{0.7\textwidth}{X} 於是二人起了爭執，甚至彼此分手。巴拿巴帶著馬可，坐船往塞浦路斯去； \end{tabularx} \\ \\ \relax
15:40 & \begin{tabularx}{0.7\textwidth}{X} 保羅則揀選了西拉，也出發了，蒙弟兄們把他交於主的恩典中。 \end{tabularx} \\ \\ \relax
15:41 & \begin{tabularx}{0.7\textwidth}{X} 他就走遍了敘利亞、基利家，堅固眾教會。 \end{tabularx} \\ \\ \relax
16:1 & \begin{tabularx}{0.7\textwidth}{X} 後來，保羅來到特庇，又到路司得。在那裡有一個門徒，名叫提摩太，是信主的猶太婦人的兒子，他父親卻是希臘人。 \end{tabularx} \\ \\ \relax
16:2 & \begin{tabularx}{0.7\textwidth}{X} 路司得和以哥念的弟兄都稱讚他。 \end{tabularx} \\ \\ \relax
16:3 & \begin{tabularx}{0.7\textwidth}{X} 保羅要帶他同去，只因那些地方的猶太人都知道他父親是希臘人，就給他行了割禮。 \end{tabularx} \\ \\ \relax
16:4 & \begin{tabularx}{0.7\textwidth}{X} 他們經過各城，把耶路撒冷使徒和長老所決定的規條交給門徒遵守。 \end{tabularx} \\ \\ \relax
16:5 & \begin{tabularx}{0.7\textwidth}{X} 於是眾教會信心越發堅固，人數天天增加。 \end{tabularx} \\ \\ \relax
16:6 & \begin{tabularx}{0.7\textwidth}{X} 因為聖靈禁止他們在亞細亞講道，他們就經過弗呂家、加拉太一帶地方。 \end{tabularx} \\ \\ \relax
16:7 & \begin{tabularx}{0.7\textwidth}{X} 到了每西亞的邊界，他們想要往庇推尼去，耶穌的靈卻不許。 \end{tabularx} \\ \\ \relax
16:8 & \begin{tabularx}{0.7\textwidth}{X} 他們就越過每西亞，下特羅亞去。 \end{tabularx} \\ \\ \relax
16:9 & \begin{tabularx}{0.7\textwidth}{X} 夜間，有異象向保羅顯現。有一個馬其頓人站著求他說：「請你過來，到馬其頓來幫助我們！」 \end{tabularx} \\ \\ \relax
16:10 & \begin{tabularx}{0.7\textwidth}{X} 保羅既看見這異象，我們就立即設法往馬其頓去，認為神呼召我們傳福音給那裡的人。 \end{tabularx} \\ \\ \relax
16:11 & \begin{tabularx}{0.7\textwidth}{X} 我們從特羅亞開船，直行駛到撒摩特喇，第二天到了尼亞坡里； \end{tabularx} \\ \\ \relax
16:12 & \begin{tabularx}{0.7\textwidth}{X} 從那裡來到腓立比，就是馬其頓這一帶的一個重要城市，也是羅馬的駐防城。我們在這城裡住了幾天。 \end{tabularx} \\ \\ \relax
16:13 & \begin{tabularx}{0.7\textwidth}{X} 在安息日，我們出城門，到了河邊，知道那裡有一個禱告的地方，我們就坐下來對那些聚會的婦女講道。 \end{tabularx} \\ \\ \relax
16:14 & \begin{tabularx}{0.7\textwidth}{X} 有一個賣紫色布的婦人，名叫呂底亞，是推雅推喇城的人，素來敬拜神。她在聽著，主就開導她的心，使她留心聽保羅所講的話。 \end{tabularx} \\ \\ \relax
16:15 & \begin{tabularx}{0.7\textwidth}{X} 她和她一家都領了洗，就求我們說：「你們若以為我是真心信主的，請到我家裡來住。」於是她堅決請我們留下。 \end{tabularx} \\ \\ \relax
16:16 & \begin{tabularx}{0.7\textwidth}{X} 後來，我們往那禱告的地方去時，有一個被占卜的靈附身的使女迎面走來，她使用法術使她的主人們發了大財。 \end{tabularx} \\ \\ \relax
16:17 & \begin{tabularx}{0.7\textwidth}{X} 她跟隨保羅和我們，喊著說：「這些人是至高神的僕人，對你們傳講救人的道路。」 \end{tabularx} \\ \\ \relax
16:18 & \begin{tabularx}{0.7\textwidth}{X} 她一連好幾天這樣喊叫，保羅就心中厭煩，轉身對那靈說：「我奉耶穌基督的名吩咐你從她身上出來！」那靈立刻出來了。 \end{tabularx} \\ \\ \relax
16:19 & \begin{tabularx}{0.7\textwidth}{X} 使女的主人們見發財的指望沒有了，就揪住保羅和西拉，拉他們到市上去見官； \end{tabularx} \\ \\ \relax
16:20 & \begin{tabularx}{0.7\textwidth}{X} 又帶他們到行政官長們面前，說：「這些騷擾我們城的，他們是猶太人， \end{tabularx} \\ \\ \relax
16:21 & \begin{tabularx}{0.7\textwidth}{X} 竟傳佈我們羅馬人所不可接受、不可遵守的規矩。」 \end{tabularx} \\ \\ \relax
16:22 & \begin{tabularx}{0.7\textwidth}{X} 群眾就一齊起來攻擊他們。官長們吩咐撕開他們的衣裳，用棍子打； \end{tabularx} \\ \\ \relax
16:23 & \begin{tabularx}{0.7\textwidth}{X} 打了許多棍，就把他們下在監裡，囑咐獄警嚴緊看守。 \end{tabularx} \\ \\ \relax
16:24 & \begin{tabularx}{0.7\textwidth}{X} 獄警領了這樣的命令，就把他們下在內監，兩腳拴在木架上。 \end{tabularx} \\ \\ \relax
16:25 & \begin{tabularx}{0.7\textwidth}{X} 約在半夜，保羅和西拉正在禱告，唱詩讚美神，眾囚犯也側耳聽著的時候， \end{tabularx} \\ \\ \relax
16:26 & \begin{tabularx}{0.7\textwidth}{X} 忽然，地大震動，甚至監牢的地基都搖動了，監門立刻全開，眾囚犯的鎖鏈也都解開了。 \end{tabularx} \\ \\ \relax
16:27 & \begin{tabularx}{0.7\textwidth}{X} 獄警一醒，看見監門全開，以為囚犯已經逃走，就拔刀要自殺。 \end{tabularx} \\ \\ \relax
16:28 & \begin{tabularx}{0.7\textwidth}{X} 保羅大聲呼叫：「不要傷害自己！我們都在這裡。」 \end{tabularx} \\ \\ \relax
16:29 & \begin{tabularx}{0.7\textwidth}{X} 獄警叫人拿燈來，就衝進去，戰戰兢兢地俯伏在保羅和西拉面前。 \end{tabularx} \\ \\ \relax
16:30 & \begin{tabularx}{0.7\textwidth}{X} 然後獄警領他們出來，說：「二位先生，我必須做甚麼才可以得救？」 \end{tabularx} \\ \\ \relax
16:31 & \begin{tabularx}{0.7\textwidth}{X} 他們說：「當信主耶穌，你和你一家都必得救。」 \end{tabularx} \\ \\ \relax
16:32 & \begin{tabularx}{0.7\textwidth}{X} 他們就把主的道講給他和他全家的人聽。 \end{tabularx} \\ \\ \relax
16:33 & \begin{tabularx}{0.7\textwidth}{X} 當夜，就在那時候，獄警把他們帶去，洗他們的傷；他和他所有的家人立刻都受了洗。 \end{tabularx} \\ \\ \relax
16:34 & \begin{tabularx}{0.7\textwidth}{X} 於是獄警領他們上自己的家裡去，給他們擺上飯。他和全家的人，因為信了神，都滿心喜樂。 \end{tabularx} \\ \\ \relax
16:35 & \begin{tabularx}{0.7\textwidth}{X} 到了天亮，官長們打發差役來，說：「釋放那兩個人吧。」 \end{tabularx} \\ \\ \relax
16:36 & \begin{tabularx}{0.7\textwidth}{X} 獄警就把這些話告訴保羅：「官長們打發人來，要釋放你們，現在可以出監，平平安安去吧。」 \end{tabularx} \\ \\ \relax
16:37 & \begin{tabularx}{0.7\textwidth}{X} 保羅卻說：「我們是羅馬人，並沒有定罪，他們竟在公眾面前打了我們，又把我們下在監裡；現在要私下趕我們出去嗎？這不行！叫他們自己來領我們出去吧！」 \end{tabularx} \\ \\ \relax
16:38 & \begin{tabularx}{0.7\textwidth}{X} 差役把這些話回稟官長們；官長們聽見他們是羅馬人，就害怕了， \end{tabularx} \\ \\ \relax
16:39 & \begin{tabularx}{0.7\textwidth}{X} 於是來勸他們，領他們出來，請他們離開那城。 \end{tabularx} \\ \\ \relax
16:40 & \begin{tabularx}{0.7\textwidth}{X} 二人出了監牢，往呂底亞家裡去，見了弟兄們，勸慰他們一番，就離開了。 \end{tabularx} \\ \\ \relax
17:1 & \begin{tabularx}{0.7\textwidth}{X} 保羅和西拉經過暗妃坡里、亞波羅尼亞，來到帖撒羅尼迦，在那裡有猶太人的會堂。 \end{tabularx} \\ \\ \relax
17:2 & \begin{tabularx}{0.7\textwidth}{X} 保羅照他素常的規矩進去，一連三個安息日，根據聖經與他們辯論， \end{tabularx} \\ \\ \relax
17:3 & \begin{tabularx}{0.7\textwidth}{X} 講解和說明基督必須受害，從死人中復活；又說：「我所傳給你們的這位耶穌就是基督。」 \end{tabularx} \\ \\ \relax
17:4 & \begin{tabularx}{0.7\textwidth}{X} 他們中間有些人聽了勸，就跟從保羅和西拉，還有許多虔敬的希臘人，尊貴的婦女也不少。 \end{tabularx} \\ \\ \relax
17:5 & \begin{tabularx}{0.7\textwidth}{X} 但不信的猶太人心裡嫉妒，聚集了些市井流氓，搭夥成群，煽動全城的人闖進耶孫的家，要把保羅和西拉帶到民眾那裡。 \end{tabularx} \\ \\ \relax
17:6 & \begin{tabularx}{0.7\textwidth}{X} 那些人找不著他們，就把耶孫和幾個弟兄拉到地方官那裡，喊叫著：「這些攪亂天下的人也到這裡來了， \end{tabularx} \\ \\ \relax
17:7 & \begin{tabularx}{0.7\textwidth}{X} 耶孫竟收留他們。這些人都違背凱撒的命令，說另有一個王耶穌。」 \end{tabularx} \\ \\ \relax
17:8 & \begin{tabularx}{0.7\textwidth}{X} 眾人和地方官聽見這些話，就惶恐了， \end{tabularx} \\ \\ \relax
17:9 & \begin{tabularx}{0.7\textwidth}{X} 於是收了耶孫和其餘的人的保證金後，釋放了他們。 \end{tabularx} \\ \\ \relax
17:10 & \begin{tabularx}{0.7\textwidth}{X} 當夜，弟兄們立刻送保羅和西拉往庇哩亞去；二人到了，就進入猶太人的會堂。 \end{tabularx} \\ \\ \relax
17:11 & \begin{tabularx}{0.7\textwidth}{X} 這地方的猶太人比帖撒羅尼迦的人開明，熱心領受這道，天天查考聖經，要知道這道是否真實。 \end{tabularx} \\ \\ \relax
17:12 & \begin{tabularx}{0.7\textwidth}{X} 所以，他們中間有許多信了，又有希臘的尊貴婦人，男人也不少。 \end{tabularx} \\ \\ \relax
17:13 & \begin{tabularx}{0.7\textwidth}{X} 但帖撒羅尼迦的猶太人知道保羅又在庇哩亞傳神的道，就往那裡去，煽動挑撥群眾。 \end{tabularx} \\ \\ \relax
17:14 & \begin{tabularx}{0.7\textwidth}{X} 於是，弟兄們立刻送保羅到海邊去，西拉和提摩太卻仍留在庇哩亞。 \end{tabularx} \\ \\ \relax
17:15 & \begin{tabularx}{0.7\textwidth}{X} 護送保羅的人帶他到了雅典，他們領了保羅的命令，叫西拉和提摩太趕快到他那裡來，然後回去了。 \end{tabularx} \\ \\ \relax
17:16 & \begin{tabularx}{0.7\textwidth}{X} 保羅在雅典等候他們的時候，看見滿城都是偶像，就心裡非常難過。 \end{tabularx} \\ \\ \relax
17:17 & \begin{tabularx}{0.7\textwidth}{X} 於是他在會堂裡與猶太人和虔敬的人，以及每日在市場上所遇見的人辯論。 \end{tabularx} \\ \\ \relax
17:18 & \begin{tabularx}{0.7\textwidth}{X} 還有伊壁鳩魯和斯多亞兩派的哲學家也與他爭辯。有的說：「這胡言亂語的要說甚麼？」有的說：「他似乎是宣傳外邦鬼神的。」這是因保羅傳講耶穌與復活的福音。 \end{tabularx} \\ \\ \relax
17:19 & \begin{tabularx}{0.7\textwidth}{X} 他們就把他帶到亞略巴古，說：「你所講的這新學說，我們也可以知道嗎？ \end{tabularx} \\ \\ \relax
17:20 & \begin{tabularx}{0.7\textwidth}{X} 因為你有些奇怪的事傳到我們耳中，我們想知道這些事是甚麼意思。」 \end{tabularx} \\ \\ \relax
17:21 & \begin{tabularx}{0.7\textwidth}{X} 原來所有的雅典人和居住在那裡的外國人都無暇管別的事，只是談談或聽聽新聞。 \end{tabularx} \\ \\ \relax
17:22 & \begin{tabularx}{0.7\textwidth}{X} 保羅站在亞略巴古當中，說：「諸位雅典人！我看你們凡事很敬畏鬼神。 \end{tabularx} \\ \\ \relax
17:23 & \begin{tabularx}{0.7\textwidth}{X} 我到處走走的時候，仔細觀察你們所敬拜的，發現一座壇，上面寫著『獻給未識之神明』。你們所不認識而敬拜的，我現在向你們宣告： \end{tabularx} \\ \\ \relax
17:24 & \begin{tabularx}{0.7\textwidth}{X} 他是創造宇宙和其中萬物的神；他既是天地的主，就不住在人手所造的殿宇裡， \end{tabularx} \\ \\ \relax
17:25 & \begin{tabularx}{0.7\textwidth}{X} 也不用人手去服侍，好像缺少甚麼似的；自己倒將生命、氣息、萬物賜給萬人。 \end{tabularx} \\ \\ \relax
17:26 & \begin{tabularx}{0.7\textwidth}{X} 他從一人造出萬族，居住在全地面上，並且預先定準他們的年限和所住的疆界， \end{tabularx} \\ \\ \relax
17:27 & \begin{tabularx}{0.7\textwidth}{X} 為要使他們尋求神，或者可以揣摩而找到他，其實他離我們各人不遠。 \end{tabularx} \\ \\ \relax
17:28 & \begin{tabularx}{0.7\textwidth}{X} 我們生活、行動、存在都在於他。就如你們的詩人也有人說：『我們也是他所生的。』 \end{tabularx} \\ \\ \relax
17:29 & \begin{tabularx}{0.7\textwidth}{X} 既然我們是神所生的，就不應該以為神的神性像人用手藝和心思所雕刻的金、銀、石像一般。 \end{tabularx} \\ \\ \relax
17:30 & \begin{tabularx}{0.7\textwidth}{X} 世人蒙昧無知的時候，神並不追究，如今卻吩咐各處的人都要悔改。 \end{tabularx} \\ \\ \relax
17:31 & \begin{tabularx}{0.7\textwidth}{X} 因為他已經定了日子，要藉著他所設立的人按公義審判天下，並且使他從死人中復活，給萬人作可信的憑據。」 \end{tabularx} \\ \\ \relax
17:32 & \begin{tabularx}{0.7\textwidth}{X} 眾人聽見死人復活的話，就有人譏誚他；又有人說：「我們會再聽你講這事。」 \end{tabularx} \\ \\ \relax
17:33 & \begin{tabularx}{0.7\textwidth}{X} 於是保羅從他們當中出去了。 \end{tabularx} \\ \\ \relax
17:34 & \begin{tabularx}{0.7\textwidth}{X} 但有幾個人依附他，信了主，其中有亞略巴古的議員丟尼修，和一個名叫大馬哩的婦人，還有幾個與他們一起的人。 \end{tabularx} \\ \\ \relax
18:1 & \begin{tabularx}{0.7\textwidth}{X} 這些事以後，保羅離開雅典，來到哥林多。 \end{tabularx} \\ \\ \relax
18:2 & \begin{tabularx}{0.7\textwidth}{X} 他遇見一個生在本都的猶太人，名叫亞居拉。不久前，他帶著妻子百基拉從意大利來，因為克勞第命令所有的猶太人都離開羅馬。保羅去投靠他們。 \end{tabularx} \\ \\ \relax
18:3 & \begin{tabularx}{0.7\textwidth}{X} 他們本是製造帳棚為業。保羅因與他們同業，就和他們同住，一同做工。 \end{tabularx} \\ \\ \relax
18:4 & \begin{tabularx}{0.7\textwidth}{X} 每逢安息日，保羅在會堂裡辯論，勸導猶太人和希臘人。 \end{tabularx} \\ \\ \relax
18:5 & \begin{tabularx}{0.7\textwidth}{X} 西拉和提摩太從馬其頓來的時候，保羅正專心傳道，向猶太人證明耶穌是基督。 \end{tabularx} \\ \\ \relax
18:6 & \begin{tabularx}{0.7\textwidth}{X} 當他們抗拒他、毀謗他的時候，他就抖掉衣裳的灰塵，對他們說：「你們的罪歸到你們自己的頭上，與我無干。從今以後，我要往外邦人那裡去。」 \end{tabularx} \\ \\ \relax
18:7 & \begin{tabularx}{0.7\textwidth}{X} 於是他離開那裡，到了一個人的家裡，他名叫提多‧猶士都，是敬拜神的人，他的家靠近會堂。 \end{tabularx} \\ \\ \relax
18:8 & \begin{tabularx}{0.7\textwidth}{X} 會堂的主管基利司布和全家都信了主，還有許多哥林多人聽了就信，而且受了洗。 \end{tabularx} \\ \\ \relax
18:9 & \begin{tabularx}{0.7\textwidth}{X} 夜間，主在異象中對保羅說：「不要怕，只管講，不要沉默， \end{tabularx} \\ \\ \relax
18:10 & \begin{tabularx}{0.7\textwidth}{X} 有我與你同在，沒有人會下手害你，因為在這城裡有許多屬我的人。」 \end{tabularx} \\ \\ \relax
18:11 & \begin{tabularx}{0.7\textwidth}{X} 保羅在那裡住了一年六個月，將神的道教導他們。 \end{tabularx} \\ \\ \relax
18:12 & \begin{tabularx}{0.7\textwidth}{X} 到迦流作亞該亞省長的時候，猶太人齊心起來攻擊保羅，拉他到法庭， \end{tabularx} \\ \\ \relax
18:13 & \begin{tabularx}{0.7\textwidth}{X} 說：「這個人教唆人不按著律法敬拜神。」 \end{tabularx} \\ \\ \relax
18:14 & \begin{tabularx}{0.7\textwidth}{X} 保羅剛要開口，迦流對猶太人說：「你們這些猶太人哪！如果是為冤枉或奸惡的事，我理當耐性聽你們。 \end{tabularx} \\ \\ \relax
18:15 & \begin{tabularx}{0.7\textwidth}{X} 既然你們所爭論的是關乎用字、名目和你們的律法，你們自己去辦吧！這樣的事我不願意審問。」 \end{tabularx} \\ \\ \relax
18:16 & \begin{tabularx}{0.7\textwidth}{X} 於是，他把他們逐出法庭。 \end{tabularx} \\ \\ \relax
18:17 & \begin{tabularx}{0.7\textwidth}{X} 眾人就揪住會堂的主管所提尼，在法庭前打他。這些事迦流都不管。 \end{tabularx} \\ \\ \relax
18:18 & \begin{tabularx}{0.7\textwidth}{X} 保羅又住了好些日子，就辭別了弟兄，坐船到敘利亞去。百基拉、亞居拉和他同去。他因為許過願，就在堅革哩剃了頭髮。 \end{tabularx} \\ \\ \relax
18:19 & \begin{tabularx}{0.7\textwidth}{X} 到了以弗所，保羅就把他們留在那裡，自己進了會堂，和猶太人辯論。 \end{tabularx} \\ \\ \relax
18:20 & \begin{tabularx}{0.7\textwidth}{X} 眾人請他多住些日子，他沒有答應， \end{tabularx} \\ \\ \relax
18:21 & \begin{tabularx}{0.7\textwidth}{X} 就辭別他們，說：「神若許可，我還要回到你們這裡來。」於是他上船離開以弗所。 \end{tabularx} \\ \\ \relax
18:22 & \begin{tabularx}{0.7\textwidth}{X} 他在凱撒利亞下了船，上耶路撒冷去問候教會，隨後下安提阿去。 \end{tabularx} \\ \\ \relax
18:23 & \begin{tabularx}{0.7\textwidth}{X} 他在那裡住了些日子，又離開了那裡，逐一經過加拉太和弗呂家各地方，堅固眾門徒。 \end{tabularx} \\ \\ \relax
18:24 & \begin{tabularx}{0.7\textwidth}{X} 有一個生在亞歷山大的猶太人，名叫亞波羅，來到以弗所，他很有口才，很會講解聖經。 \end{tabularx} \\ \\ \relax
18:25 & \begin{tabularx}{0.7\textwidth}{X} 這人已經在主的道路上受了訓練，心裡火熱，精確地講論和教導耶穌的事；可是他只知道約翰的洗禮。 \end{tabularx} \\ \\ \relax
18:26 & \begin{tabularx}{0.7\textwidth}{X} 他開始在會堂裡放膽講道；百基拉、亞居拉聽見，就接他來，將神的道路給他更精確地講解。 \end{tabularx} \\ \\ \relax
18:27 & \begin{tabularx}{0.7\textwidth}{X} 他想要往亞該亞去，弟兄們就勉勵他，並寫信請門徒們接待他，他到了那裡，多多幫助那些蒙恩信主的人， \end{tabularx} \\ \\ \relax
18:28 & \begin{tabularx}{0.7\textwidth}{X} 因為他在公眾面前極力駁倒猶太人，引聖經證明耶穌是基督。 \end{tabularx} \\ \\ \relax
19:1 & \begin{tabularx}{0.7\textwidth}{X} 亞波羅在哥林多的時候，保羅經過了內陸地區，來到以弗所，在那裡他遇見幾個門徒， \end{tabularx} \\ \\ \relax
19:2 & \begin{tabularx}{0.7\textwidth}{X} 問他們：「你們信的時候領受了聖靈沒有？」他們說：「沒有，我們連甚麼是聖靈都沒有聽過。」 \end{tabularx} \\ \\ \relax
19:3 & \begin{tabularx}{0.7\textwidth}{X} 保羅說：「這樣，你們受的是甚麼洗呢？」他們說：「是受了約翰的洗。」 \end{tabularx} \\ \\ \relax
19:4 & \begin{tabularx}{0.7\textwidth}{X} 保羅說：「約翰所施的是悔改的洗禮，他告訴百姓當信那在他以後要來的那位，就是耶穌。」 \end{tabularx} \\ \\ \relax
19:5 & \begin{tabularx}{0.7\textwidth}{X} 他們聽見這話以後，就奉主耶穌的名受洗。 \end{tabularx} \\ \\ \relax
19:6 & \begin{tabularx}{0.7\textwidth}{X} 保羅給他們按手，聖靈就降在他們身上，他們開始說方言 和說預言。 \end{tabularx} \\ \\ \relax
19:7 & \begin{tabularx}{0.7\textwidth}{X} 他們約有十二個人。 \end{tabularx} \\ \\ \relax
19:8 & \begin{tabularx}{0.7\textwidth}{X} 保羅進會堂，一連三個月放膽講道，辯論神國的事，勸導眾人。 \end{tabularx} \\ \\ \relax
19:9 & \begin{tabularx}{0.7\textwidth}{X} 後來，有些人心裡剛硬不信，在眾人面前毀謗這道；保羅就離開他們，也叫門徒與他們分開，就在推喇奴的講堂天天辯論。 \end{tabularx} \\ \\ \relax
19:10 & \begin{tabularx}{0.7\textwidth}{X} 這樣有兩年之久，使一切住在亞細亞的，無論是猶太人是希臘人，都聽見主的道。 \end{tabularx} \\ \\ \relax
19:11 & \begin{tabularx}{0.7\textwidth}{X} 神藉保羅的手行了些奇異的神蹟， \end{tabularx} \\ \\ \relax
19:12 & \begin{tabularx}{0.7\textwidth}{X} 甚至有人從保羅身上拿走手巾或圍裙放在病人身上，病就消除了，邪靈也出去了。 \end{tabularx} \\ \\ \relax
19:13 & \begin{tabularx}{0.7\textwidth}{X} 那時，有幾個巡迴各處念咒趕鬼的猶太人，擅自利用主耶穌的名，向那些被邪靈所附的人說：「我奉保羅所傳的耶穌命令你們出來！」 \end{tabularx} \\ \\ \relax
19:14 & \begin{tabularx}{0.7\textwidth}{X} 做這事的是猶太祭司長士基瓦的七個兒子。 \end{tabularx} \\ \\ \relax
19:15 & \begin{tabularx}{0.7\textwidth}{X} 但邪靈回答他們：「耶穌我知道，保羅我也認識，你們卻是誰呢？」 \end{tabularx} \\ \\ \relax
19:16 & \begin{tabularx}{0.7\textwidth}{X} 被邪靈所附的人就撲到他們身上，制伏他們，勝過他們，使他們赤著身子，受了傷，從那房子裡逃出去了。 \end{tabularx} \\ \\ \relax
19:17 & \begin{tabularx}{0.7\textwidth}{X} 凡住在以弗所的，無論是猶太人是希臘人，都知道這件事，也都懼怕；主耶穌的名從此就更被尊為大了。 \end{tabularx} \\ \\ \relax
19:18 & \begin{tabularx}{0.7\textwidth}{X} 許多已經信的人來承認並公開自己所行的事。 \end{tabularx} \\ \\ \relax
19:19 & \begin{tabularx}{0.7\textwidth}{X} 又有許多平素行邪術的人把他們的書都拿來，堆積在眾人面前焚燒。他們計算書價，得知共值五萬塊銀錢。 \end{tabularx} \\ \\ \relax
19:20 & \begin{tabularx}{0.7\textwidth}{X} 這樣，主的道大大興旺，而且普遍傳開了。 \end{tabularx} \\ \\ \relax
19:21 & \begin{tabularx}{0.7\textwidth}{X} 這些事過後，保羅心裡決定要經過馬其頓、亞該亞，就往耶路撒冷去。他說：「我到了那裡以後，也必須到羅馬去看看。」 \end{tabularx} \\ \\ \relax
19:22 & \begin{tabularx}{0.7\textwidth}{X} 於是他差遣兩個助手提摩太和以拉都往馬其頓去，自己暫時留在亞細亞。 \end{tabularx} \\ \\ \relax
19:23 & \begin{tabularx}{0.7\textwidth}{X} 那時，因這道路而起的騷動不小。 \end{tabularx} \\ \\ \relax
19:24 & \begin{tabularx}{0.7\textwidth}{X} 有一個銀匠，名叫底米丟，是製造亞底米神銀龕的，他使從事這手藝的人生意發達。 \end{tabularx} \\ \\ \relax
19:25 & \begin{tabularx}{0.7\textwidth}{X} 他聚集他們和同行的工人，說：「諸位，你們知道我們是倚靠這生意發財的。 \end{tabularx} \\ \\ \relax
19:26 & \begin{tabularx}{0.7\textwidth}{X} 你們看到，也聽見這保羅不但在以弗所，也幾乎在亞細亞全地，引誘迷惑了許多人，說：『人手所做的不是神明。』 \end{tabularx} \\ \\ \relax
19:27 & \begin{tabularx}{0.7\textwidth}{X} 這樣，不僅我們這行業陷入被藐視的危險，就是大女神亞底米的廟也要被人輕看，連亞細亞全地和普天下所敬拜的女神的威望也受損害了。」 \end{tabularx} \\ \\ \relax
19:28 & \begin{tabularx}{0.7\textwidth}{X} 眾人聽見，就怒氣沖沖，喊著說：「大哉，以弗所人的亞底米！」 \end{tabularx} \\ \\ \relax
19:29 & \begin{tabularx}{0.7\textwidth}{X} 於是滿城都騷動起來。眾人抓住與保羅同行的馬其頓人該猶和亞里達古，齊心衝進劇場。 \end{tabularx} \\ \\ \relax
19:30 & \begin{tabularx}{0.7\textwidth}{X} 保羅想要進到民眾那裡，門徒卻不許他去。 \end{tabularx} \\ \\ \relax
19:31 & \begin{tabularx}{0.7\textwidth}{X} 連亞細亞的幾位官員，是保羅的朋友，也打發人來勸他不要冒險到劇場裡去。 \end{tabularx} \\ \\ \relax
19:32 & \begin{tabularx}{0.7\textwidth}{X} 聚集的人亂成一團，有的喊這個，有的喊那個，大半不知道為了甚麼聚集。 \end{tabularx} \\ \\ \relax
19:33 & \begin{tabularx}{0.7\textwidth}{X} 猶太人把亞歷山大推出去，人群中有人慫恿他，他就做手勢，要向民眾申訴。 \end{tabularx} \\ \\ \relax
19:34 & \begin{tabularx}{0.7\textwidth}{X} 但他們一認出他是猶太人，大家就異口同聲喊著：「大哉，以弗所人的亞底米！」約喊了兩小時。 \end{tabularx} \\ \\ \relax
19:35 & \begin{tabularx}{0.7\textwidth}{X} 城裡的書記官安撫了群眾後，說：「以弗所人哪，誰不知道以弗所人的城是看守大亞底米的廟和從宙斯那裡落下來的像的守護者呢？ \end{tabularx} \\ \\ \relax
19:36 & \begin{tabularx}{0.7\textwidth}{X} 既然這些事是駁不倒的，你們就要安靜下來，不可妄動。 \end{tabularx} \\ \\ \relax
19:37 & \begin{tabularx}{0.7\textwidth}{X} 你們把這些人帶來，他們並沒有偷竊廟中之物，也沒有褻瀆我們的女神。 \end{tabularx} \\ \\ \relax
19:38 & \begin{tabularx}{0.7\textwidth}{X} 如果底米丟和他同行的手藝人有控告的事，自有公堂，也有省長，他們可以彼此控告。 \end{tabularx} \\ \\ \relax
19:39 & \begin{tabularx}{0.7\textwidth}{X} 你們若有別的事請求，可以在合法的集會裡解決。 \end{tabularx} \\ \\ \relax
19:40 & \begin{tabularx}{0.7\textwidth}{X} 今日的擾亂本是無緣無故的，有被控告的危險。這次的騷動，我們也說不出理由來。」 \end{tabularx} \\ \\ \relax
19:41 & \begin{tabularx}{0.7\textwidth}{X} 他說完這些話，就叫眾人散會。 \end{tabularx} \\ \\ \relax
20:1 & \begin{tabularx}{0.7\textwidth}{X} 騷亂平定以後，保羅請門徒來，勸勉了他們，就辭別他們，往馬其頓去。 \end{tabularx} \\ \\ \relax
20:2 & \begin{tabularx}{0.7\textwidth}{X} 他走遍那一帶地方，用許多話勸勉門徒，然後來到希臘， \end{tabularx} \\ \\ \relax
20:3 & \begin{tabularx}{0.7\textwidth}{X} 在那裡住了三個月。他快要坐船往敘利亞去的時候，猶太人設計害他，他就決定從馬其頓回去。 \end{tabularx} \\ \\ \relax
20:4 & \begin{tabularx}{0.7\textwidth}{X} 同他到亞細亞去的，有庇哩亞人畢羅斯的兒子所巴特，帖撒羅尼迦人亞里達古和西公都，還有特庇人該猶和提摩太，又有亞細亞人推基古和特羅非摩。 \end{tabularx} \\ \\ \relax
20:5 & \begin{tabularx}{0.7\textwidth}{X} 這些人先走，在特羅亞等候我們。 \end{tabularx} \\ \\ \relax
20:6 & \begin{tabularx}{0.7\textwidth}{X} 過了除酵節的日子，我們從腓立比開船，五天以後到了特羅亞，和他們相會，在那裡住了七天。 \end{tabularx} \\ \\ \relax
20:7 & \begin{tabularx}{0.7\textwidth}{X} 七日的第一日，我們聚會擘餅的時候，保羅因次日要起行，就為他們講道，直講到半夜。 \end{tabularx} \\ \\ \relax
20:8 & \begin{tabularx}{0.7\textwidth}{X} 我們聚會的那座樓上有好些燈火。 \end{tabularx} \\ \\ \relax
20:9 & \begin{tabularx}{0.7\textwidth}{X} 有一個少年，名叫猶推古，坐在窗口上，沉沉入睡。保羅講了多時，少年睡熟了，從三層樓上掉下去，扶起來時已經死了。 \end{tabularx} \\ \\ \relax
20:10 & \begin{tabularx}{0.7\textwidth}{X} 保羅下去，伏在他身上，抱著他，說：「你們不要慌亂，他還有氣呢！」 \end{tabularx} \\ \\ \relax
20:11 & \begin{tabularx}{0.7\textwidth}{X} 保羅又上樓去，擘餅，吃了，再講了許久，直到天亮才離開。 \end{tabularx} \\ \\ \relax
20:12 & \begin{tabularx}{0.7\textwidth}{X} 他們把那活過來的孩子帶走，大家得到很大的安慰。 \end{tabularx} \\ \\ \relax
20:13 & \begin{tabularx}{0.7\textwidth}{X} 我們先上船，起航往亞朔去，想要在那裡接保羅；因為他是這樣安排的，他自己本來打算要走陸路。 \end{tabularx} \\ \\ \relax
20:14 & \begin{tabularx}{0.7\textwidth}{X} 他既在亞朔與我們相會，我們就接他上船，來到米推利尼。 \end{tabularx} \\ \\ \relax
20:15 & \begin{tabularx}{0.7\textwidth}{X} 我們從那裡開船，第二天到了基阿的對岸；再下一天，在撒摩靠岸，又過了一天，到了米利都。 \end{tabularx} \\ \\ \relax
20:16 & \begin{tabularx}{0.7\textwidth}{X} 因為保羅早已決定要越過以弗所，免得在亞細亞耽延，他急忙前行，假如可能的話，在五旬節前能趕到耶路撒冷。 \end{tabularx} \\ \\ \relax
20:17 & \begin{tabularx}{0.7\textwidth}{X} 保羅從米利都打發人往以弗所去，請教會的長老來。 \end{tabularx} \\ \\ \relax
20:18 & \begin{tabularx}{0.7\textwidth}{X} 他們來了，保羅對他們說：「你們自己知道，自從我到亞細亞的第一天，我怎樣跟你們相處， \end{tabularx} \\ \\ \relax
20:19 & \begin{tabularx}{0.7\textwidth}{X} 怎樣凡事謙卑，以眼淚服侍主，又因猶太人的謀害經歷試煉。 \end{tabularx} \\ \\ \relax
20:20 & \begin{tabularx}{0.7\textwidth}{X} 你們也知道，凡對你們有益的，我沒有一樣隱瞞不說的，或在公眾面前，或在每一個人的家裡，我都教導你們， \end{tabularx} \\ \\ \relax
20:21 & \begin{tabularx}{0.7\textwidth}{X} 不論猶太人和希臘人，我都已證明他們當在神面前悔改，信靠我們的主耶穌。 \end{tabularx} \\ \\ \relax
20:22 & \begin{tabularx}{0.7\textwidth}{X} 現在我被聖靈催迫要往耶路撒冷去，雖然不知道在那裡會遭遇甚麼事， \end{tabularx} \\ \\ \relax
20:23 & \begin{tabularx}{0.7\textwidth}{X} 但知道聖靈在各城裡向我指證，說有捆鎖與患難等著我。 \end{tabularx} \\ \\ \relax
20:24 & \begin{tabularx}{0.7\textwidth}{X} 我卻不以性命為念，只要走完我的路程，完成我從主耶穌所領受的職分，為神恩典的福音作見證。 \end{tabularx} \\ \\ \relax
20:25 & \begin{tabularx}{0.7\textwidth}{X} 「我素常在你們中間到處傳講神的國；現在我知道，你們眾人以後不會再見到我的面了。 \end{tabularx} \\ \\ \relax
20:26 & \begin{tabularx}{0.7\textwidth}{X} 所以我今日向你們作證，你們中間無論何人死亡，罪不在我。 \end{tabularx} \\ \\ \relax
20:27 & \begin{tabularx}{0.7\textwidth}{X} 因為神一切的旨意，我並沒有退縮不傳給你們的。 \end{tabularx} \\ \\ \relax
20:28 & \begin{tabularx}{0.7\textwidth}{X} 聖靈立你們作全群的監督，你們就當為自己謹慎，也為全群謹慎，牧養神的教會，就是他用自己血所買來的。 \end{tabularx} \\ \\ \relax
20:29 & \begin{tabularx}{0.7\textwidth}{X} 我知道，在我離開以後必有兇暴的豺狼進入你們中間，不顧惜羊群。 \end{tabularx} \\ \\ \relax
20:30 & \begin{tabularx}{0.7\textwidth}{X} 就是你們中間也必有人起來，說悖謬的話，要引誘門徒跟從他們。 \end{tabularx} \\ \\ \relax
20:31 & \begin{tabularx}{0.7\textwidth}{X} 所以你們要警醒，記念我三年之久，晝夜不斷地流淚勸戒你們各人。 \end{tabularx} \\ \\ \relax
20:32 & \begin{tabularx}{0.7\textwidth}{X} 現在我把你們交託給神和他恩惠的道；這道能建立你們，使你們和一切成聖的人同得基業。 \end{tabularx} \\ \\ \relax
20:33 & \begin{tabularx}{0.7\textwidth}{X} 我未曾貪圖一個人的金、銀或衣服。 \end{tabularx} \\ \\ \relax
20:34 & \begin{tabularx}{0.7\textwidth}{X} 你們自己知道，我靠兩隻手工作來供給我和同工的需用。 \end{tabularx} \\ \\ \relax
20:35 & \begin{tabularx}{0.7\textwidth}{X} 我凡事給你們作榜樣，叫你們知道應當這樣勞苦，扶助軟弱的人，又當記念主耶穌的話，說：『施比受更為有福。』」 \end{tabularx} \\ \\ \relax
20:36 & \begin{tabularx}{0.7\textwidth}{X} 保羅說完了這些話，就和大家跪下來禱告。 \end{tabularx} \\ \\ \relax
20:37 & \begin{tabularx}{0.7\textwidth}{X} 眾人痛哭，抱著保羅的頸項跟他親吻。 \end{tabularx} \\ \\ \relax
20:38 & \begin{tabularx}{0.7\textwidth}{X} 叫他們最傷心的，就是他說「以後不會再見到我的面」那句話。於是他們送他上船去了。 \end{tabularx} \\ \\ \relax
21:1 & \begin{tabularx}{0.7\textwidth}{X} 我們離別了眾人，就開船直航到哥士，第二天到了羅底，又從那裡到帕大喇。 \end{tabularx} \\ \\ \relax
21:2 & \begin{tabularx}{0.7\textwidth}{X} 我們遇見一隻船要往腓尼基去，就上船起航。 \end{tabularx} \\ \\ \relax
21:3 & \begin{tabularx}{0.7\textwidth}{X} 我們望見塞浦路斯，就從南邊行過，往敘利亞去，在推羅上岸，因為船要在那裡卸貨。 \end{tabularx} \\ \\ \relax
21:4 & \begin{tabularx}{0.7\textwidth}{X} 我們在那裡找到了一些門徒，就住了七天。他們藉著聖靈的感動，告訴保羅不要上耶路撒冷去。 \end{tabularx} \\ \\ \relax
21:5 & \begin{tabularx}{0.7\textwidth}{X} 幾天之後，我們又出發前行。他們眾人同妻子兒女都送我們到城外，我們都跪在灘上禱告，彼此辭別。 \end{tabularx} \\ \\ \relax
21:6 & \begin{tabularx}{0.7\textwidth}{X} 我們上了船，他們就回家去了。 \end{tabularx} \\ \\ \relax
21:7 & \begin{tabularx}{0.7\textwidth}{X} 我們從推羅行完航程，來到了多利買，問候那裡的弟兄，和他們同住了一天。 \end{tabularx} \\ \\ \relax
21:8 & \begin{tabularx}{0.7\textwidth}{X} 第二天，我們離開那裡，來到凱撒利亞，就進了傳福音的腓利家裡，和他同住；他是那七個執事裡的一個。 \end{tabularx} \\ \\ \relax
21:9 & \begin{tabularx}{0.7\textwidth}{X} 他有四個女兒，都是未出嫁的，都會說預言。 \end{tabularx} \\ \\ \relax
21:10 & \begin{tabularx}{0.7\textwidth}{X} 我們在那裡多住了好幾天，有一個先知，名叫亞迦布，從猶太下來。 \end{tabularx} \\ \\ \relax
21:11 & \begin{tabularx}{0.7\textwidth}{X} 他到了我們這裡，就拿保羅的腰帶，捆上自己的手腳，說：「聖靈這樣說：『猶太人在耶路撒冷要如此捆綁這腰帶的主人，把他交在外邦人手裡。』」 \end{tabularx} \\ \\ \relax
21:12 & \begin{tabularx}{0.7\textwidth}{X} 我們聽見這些話，就跟當地的人苦勸保羅不要上耶路撒冷去。 \end{tabularx} \\ \\ \relax
21:13 & \begin{tabularx}{0.7\textwidth}{X} 於是保羅回答：「你們為甚麼這樣痛哭，使我心碎呢？我為主耶穌的名，不但被人捆綁，就是死在耶路撒冷也是願意的。」 \end{tabularx} \\ \\ \relax
21:14 & \begin{tabularx}{0.7\textwidth}{X} 既然保羅不聽勸，我們就住了口，只說：「願主的旨意成就。」 \end{tabularx} \\ \\ \relax
21:15 & \begin{tabularx}{0.7\textwidth}{X} 過了這幾天，我們收拾行李上耶路撒冷去。 \end{tabularx} \\ \\ \relax
21:16 & \begin{tabularx}{0.7\textwidth}{X} 有凱撒利亞的幾個門徒和我們同去，帶我們到一個早期的門徒塞浦路斯人拿孫的家裡，請我們與他同住。 \end{tabularx} \\ \\ \relax
21:17 & \begin{tabularx}{0.7\textwidth}{X} 我們到了耶路撒冷，弟兄們歡歡喜喜地接待我們。 \end{tabularx} \\ \\ \relax
21:18 & \begin{tabularx}{0.7\textwidth}{X} 第二天，保羅同我們去見雅各；所有的長老也都在場。 \end{tabularx} \\ \\ \relax
21:19 & \begin{tabularx}{0.7\textwidth}{X} 保羅向他們問安，然後將神用他在外邦人中所做的事奉，一一述說了。 \end{tabularx} \\ \\ \relax
21:20 & \begin{tabularx}{0.7\textwidth}{X} 他們聽見了，就歸榮耀給神，對保羅說：「弟兄，你看猶太人中有數以萬計的信徒，而他們都是熱心於律法的人。 \end{tabularx} \\ \\ \relax
21:21 & \begin{tabularx}{0.7\textwidth}{X} 他們曾聽見人說，你教導所有在外邦的猶太人離棄摩西，對他們說，不要給孩子行割禮，也不要遵守規矩。 \end{tabularx} \\ \\ \relax
21:22 & \begin{tabularx}{0.7\textwidth}{X} 眾人必聽見你來了，這可怎麼辦呢？ \end{tabularx} \\ \\ \relax
21:23 & \begin{tabularx}{0.7\textwidth}{X} 你就照著我們的話做吧！我們這裡有四個人，都有願在身。 \end{tabularx} \\ \\ \relax
21:24 & \begin{tabularx}{0.7\textwidth}{X} 你帶他們去，與他們一同行潔淨的禮，替他們繳納規費，讓他們得以剃頭。這樣，眾人就會知道，先前所聽見關於你的事都是假的；而且也知道，你自己為人循規蹈矩，遵行律法。 \end{tabularx} \\ \\ \relax
21:25 & \begin{tabularx}{0.7\textwidth}{X} 至於信主的外邦人，我們已經根據我們的決議寫信，叫他們要禁戒偶像所玷污的東西、血和勒死的牲畜，禁戒淫亂。」 \end{tabularx} \\ \\ \relax
21:26 & \begin{tabularx}{0.7\textwidth}{X} 於是保羅帶著那四個人，第二天與他們一同行了潔淨禮，進了聖殿，報告潔淨期滿的日子，等候祭司為他們各人獻上祭物。 \end{tabularx} \\ \\ \relax
21:27 & \begin{tabularx}{0.7\textwidth}{X} 那七日將完，從亞細亞來的猶太人看見保羅在聖殿裡，就煽動所有的群眾，下手拿住他， \end{tabularx} \\ \\ \relax
21:28 & \begin{tabularx}{0.7\textwidth}{X} 喊著：「以色列人哪，來幫忙！這就是在各處教導眾人糟蹋我們百姓、律法和這地方的人。不但如此，他還帶了希臘人進聖殿，污穢了這聖地。」 \end{tabularx} \\ \\ \relax
21:29 & \begin{tabularx}{0.7\textwidth}{X} 這話是因他們曾看見以弗所人特羅非摩跟保羅一起在城裡，以為保羅帶他進了聖殿。 \end{tabularx} \\ \\ \relax
21:30 & \begin{tabularx}{0.7\textwidth}{X} 於是全城都騷動，百姓一齊跑來，拿住保羅，拉他出聖殿，殿門立刻都關了。 \end{tabularx} \\ \\ \relax
21:31 & \begin{tabularx}{0.7\textwidth}{X} 他們正想要殺他，有人報信給營裡的千夫長，說耶路撒冷全城都亂了。 \end{tabularx} \\ \\ \relax
21:32 & \begin{tabularx}{0.7\textwidth}{X} 千夫長立刻帶著士兵和幾個百夫長，跑下去到他們那裡。他們見了千夫長和士兵，就停下來不打保羅。 \end{tabularx} \\ \\ \relax
21:33 & \begin{tabularx}{0.7\textwidth}{X} 於是千夫長上前拿住他，吩咐用兩條鐵鏈捆鎖，又問他是甚麼人，做了甚麼事。 \end{tabularx} \\ \\ \relax
21:34 & \begin{tabularx}{0.7\textwidth}{X} 群眾中有的喊這個，有的喊那個；因為這樣亂嚷，千夫長無法知道實情，就下令將保羅帶進營樓去。 \end{tabularx} \\ \\ \relax
21:35 & \begin{tabularx}{0.7\textwidth}{X} 保羅一走上臺階，群眾擠得兇猛，士兵只得將保羅抬起來。 \end{tabularx} \\ \\ \relax
21:36 & \begin{tabularx}{0.7\textwidth}{X} 一群人跟在後面，喊著：「除掉他！」 \end{tabularx} \\ \\ \relax
21:37 & \begin{tabularx}{0.7\textwidth}{X} 保羅快要被帶進營樓時，對千夫長說：「我可以對你說句話嗎？」千夫長說：「你懂得希臘話嗎？ \end{tabularx} \\ \\ \relax
21:38 & \begin{tabularx}{0.7\textwidth}{X} 那你就不是從前作亂、帶領四千兇徒往曠野去的那埃及人了。」 \end{tabularx} \\ \\ \relax
21:39 & \begin{tabularx}{0.7\textwidth}{X} 保羅說：「我本是猶太人，生在基利家的大數，並不是無名小城的公民。求你准我對百姓說話。」 \end{tabularx} \\ \\ \relax
21:40 & \begin{tabularx}{0.7\textwidth}{X} 千夫長准了。保羅就站在臺階上，向百姓做了個手勢，要他們靜下來，保羅就用希伯來話對他們說： \end{tabularx} \\ \\ \relax
22:1 & \begin{tabularx}{0.7\textwidth}{X} 「諸位父老弟兄，請聽我現在對你們的申辯。」 \end{tabularx} \\ \\ \relax
22:2 & \begin{tabularx}{0.7\textwidth}{X} 他們聽保羅說的是希伯來話，就更加安靜了。 \end{tabularx} \\ \\ \relax
22:3 & \begin{tabularx}{0.7\textwidth}{X} 保羅說：「我原是猶太人，生在基利家的大數，但在這城裡長大，在迦瑪列門下按著我們祖宗嚴緊的律法受教，熱心事奉神，就如你們大家今日一樣。 \end{tabularx} \\ \\ \relax
22:4 & \begin{tabularx}{0.7\textwidth}{X} 我也曾迫害信奉這道路的人，置他們於死地，無論男女都捆綁，關在監裡。 \end{tabularx} \\ \\ \relax
22:5 & \begin{tabularx}{0.7\textwidth}{X} 這是大祭司和議會的眾長老都可以給我作證的。我又從他們那裡領了致弟兄們的書信，往大馬士革去，要把在那裡的信徒綁起來，帶到耶路撒冷受刑。」 \end{tabularx} \\ \\ \relax
22:6 & \begin{tabularx}{0.7\textwidth}{X} 「當我走近大馬士革的時候，約在中午，忽然有一道大光從天上下來，照射在我周圍。 \end{tabularx} \\ \\ \relax
22:7 & \begin{tabularx}{0.7\textwidth}{X} 我就仆倒在地，聽見有聲音對我說：『掃羅！掃羅！你為甚麼迫害我？』 \end{tabularx} \\ \\ \relax
22:8 & \begin{tabularx}{0.7\textwidth}{X} 我回答：『主啊！你是誰？』他對我說：『我就是你所迫害的拿撒勒人耶穌。』 \end{tabularx} \\ \\ \relax
22:9 & \begin{tabularx}{0.7\textwidth}{X} 跟我一起的人看見了那光，卻沒有聽見那位對我說話的聲音。 \end{tabularx} \\ \\ \relax
22:10 & \begin{tabularx}{0.7\textwidth}{X} 我說：『主啊，我該做甚麼？』主說：『起來，進大馬士革去，在那裡有人會把指派你做的一切事告訴你。』 \end{tabularx} \\ \\ \relax
22:11 & \begin{tabularx}{0.7\textwidth}{X} 我因那光的閃耀不能看見，跟我一起的人就拉著我的手進了大馬士革。 \end{tabularx} \\ \\ \relax
22:12 & \begin{tabularx}{0.7\textwidth}{X} 「那裡有一個人，名叫亞拿尼亞，按著律法是虔誠人，為所有住在那裡的猶太人所稱讚。 \end{tabularx} \\ \\ \relax
22:13 & \begin{tabularx}{0.7\textwidth}{X} 他來見我，站在旁邊，對我說：『掃羅弟兄，你看見吧！』就在那時，我恢復視覺，看見了他。 \end{tabularx} \\ \\ \relax
22:14 & \begin{tabularx}{0.7\textwidth}{X} 他又說：『我們祖宗的神揀選了你，讓你明白他的旨意，又看見那義者，聽見他口中所出的聲音。 \end{tabularx} \\ \\ \relax
22:15 & \begin{tabularx}{0.7\textwidth}{X} 因為你要將所看見的、所聽見的，對著萬人作他的見證人。 \end{tabularx} \\ \\ \relax
22:16 & \begin{tabularx}{0.7\textwidth}{X} 現在你為甚麼耽延呢？起來，受洗，求告他的名，洗去你的罪。』」 \end{tabularx} \\ \\ \relax
22:17 & \begin{tabularx}{0.7\textwidth}{X} 「後來，我回到耶路撒冷，在聖殿裡禱告的時候，魂遊象外， \end{tabularx} \\ \\ \relax
22:18 & \begin{tabularx}{0.7\textwidth}{X} 看見主對我說：『你趕緊離開耶路撒冷，越快越好，因為這裡的人不接受你為我作的見證。』 \end{tabularx} \\ \\ \relax
22:19 & \begin{tabularx}{0.7\textwidth}{X} 我就說：『主啊，他們都知道，我從前在各會堂裡把信你的人監禁，又鞭打他們。 \end{tabularx} \\ \\ \relax
22:20 & \begin{tabularx}{0.7\textwidth}{X} 當你的見證人司提反被害流血的時候，我也站在一旁贊同；又為打死他的人看守衣裳。』 \end{tabularx} \\ \\ \relax
22:21 & \begin{tabularx}{0.7\textwidth}{X} 主對我說：『你去吧！我要差你到遠方外邦人那裡去。』」 \end{tabularx} \\ \\ \relax
22:22 & \begin{tabularx}{0.7\textwidth}{X} 眾人聽他說到這句話，就高聲說：「這樣的人，從地上除掉他吧！他是該死的。」 \end{tabularx} \\ \\ \relax
22:23 & \begin{tabularx}{0.7\textwidth}{X} 大家一邊喧嚷一邊摔衣裳，向空中撒灰塵。 \end{tabularx} \\ \\ \relax
22:24 & \begin{tabularx}{0.7\textwidth}{X} 千夫長下令把保羅帶進營樓，叫人用鞭子拷問他，要知道他們向他這樣喧嚷是甚麼緣故。 \end{tabularx} \\ \\ \relax
22:25 & \begin{tabularx}{0.7\textwidth}{X} 他們剛用皮條把他捆上的時候，保羅對站在旁邊的百夫長說：「一個羅馬人，又未被定罪，你們就鞭打他是合法的嗎？」 \end{tabularx} \\ \\ \relax
22:26 & \begin{tabularx}{0.7\textwidth}{X} 百夫長聽見這話，就去見千夫長，報告說：「你要怎麼辦呢？這個人是羅馬人。」 \end{tabularx} \\ \\ \relax
22:27 & \begin{tabularx}{0.7\textwidth}{X} 千夫長就來問保羅：「你告訴我，你是羅馬人嗎？」保羅說：「是。」 \end{tabularx} \\ \\ \relax
22:28 & \begin{tabularx}{0.7\textwidth}{X} 千夫長回答：「我用了許多銀子才得到羅馬公民的身份。」保羅說：「我生來就是。」 \end{tabularx} \\ \\ \relax
22:29 & \begin{tabularx}{0.7\textwidth}{X} 於是那些要拷問保羅的人立刻離開他走了。千夫長一知道他是羅馬人，又因為曾捆綁了他，也害怕起來。 \end{tabularx} \\ \\ \relax
22:30 & \begin{tabularx}{0.7\textwidth}{X} 第二天，千夫長為要知道猶太人控告保羅的實情，就解開他，下令祭司長們和全議會的人都聚集，然後將保羅帶下來，叫他站在他們面前。 \end{tabularx} \\ \\ \relax
23:1 & \begin{tabularx}{0.7\textwidth}{X} 保羅定睛看著議會的人，說：「諸位弟兄，我在神面前，行事為人都是憑著清白的良心，直到今日。」 \end{tabularx} \\ \\ \relax
23:2 & \begin{tabularx}{0.7\textwidth}{X} 亞拿尼亞大祭司就吩咐旁邊站著的人打他的嘴。 \end{tabularx} \\ \\ \relax
23:3 & \begin{tabularx}{0.7\textwidth}{X} 這時，保羅對他說：「你這粉飾的牆，神要打你！你坐堂是要按律法審問我，你竟違背律法，命令人打我嗎？」 \end{tabularx} \\ \\ \relax
23:4 & \begin{tabularx}{0.7\textwidth}{X} 站在旁邊的人說：「你竟敢辱罵神的大祭司嗎？」 \end{tabularx} \\ \\ \relax
23:5 & \begin{tabularx}{0.7\textwidth}{X} 保羅說：「弟兄們，我不知道他是大祭司；因為經上記著：『不可毀謗你百姓的官長。』」 \end{tabularx} \\ \\ \relax
23:6 & \begin{tabularx}{0.7\textwidth}{X} 保羅看出他們一部分是撒都該人，一部分是法利賽人，就在議會中喊著：「諸位弟兄，我是法利賽人，也是法利賽人的子孫。我現在受審問是為有關死人復活的盼望。」 \end{tabularx} \\ \\ \relax
23:7 & \begin{tabularx}{0.7\textwidth}{X} 說了這話，法利賽人和撒都該人爭論起來，會眾分為兩派。 \end{tabularx} \\ \\ \relax
23:8 & \begin{tabularx}{0.7\textwidth}{X} 因為撒都該人一方面說沒有復活，另一方面沒有天使和鬼魂；法利賽人卻承認兩方面都有。 \end{tabularx} \\ \\ \relax
23:9 & \begin{tabularx}{0.7\textwidth}{X} 於是大大地爭吵起來；有幾個法利賽派的文士站起來爭辯說：「我們看不出這人有甚麼錯處；說不定有鬼魂或者天使對他說過話呢！」 \end{tabularx} \\ \\ \relax
23:10 & \begin{tabularx}{0.7\textwidth}{X} 那時爭辯越來越大，千夫長恐怕保羅被他們扯碎了，就命令士兵下去，把他從眾人當中搶出來，帶進營樓去。 \end{tabularx} \\ \\ \relax
23:11 & \begin{tabularx}{0.7\textwidth}{X} 當夜，主站在保羅旁邊，說：「放心吧！你怎樣在耶路撒冷為我作見證，也必怎樣在羅馬為我作見證。」 \end{tabularx} \\ \\ \relax
23:12 & \begin{tabularx}{0.7\textwidth}{X} 到了天亮，猶太人同謀起誓，說「若不先殺保羅就不吃不喝」。 \end{tabularx} \\ \\ \relax
23:13 & \begin{tabularx}{0.7\textwidth}{X} 參與這陰謀的有四十多人。 \end{tabularx} \\ \\ \relax
23:14 & \begin{tabularx}{0.7\textwidth}{X} 他們來見祭司長和長老，說：「我們已經發了重誓，若不先殺保羅就甚麼也不吃。 \end{tabularx} \\ \\ \relax
23:15 & \begin{tabularx}{0.7\textwidth}{X} 現在你們和議會要通知千夫長，叫他把保羅帶到你們這裡來，假裝要詳細調查他的事；我們已經預備好，在他來到這裡以前就殺掉他。」 \end{tabularx} \\ \\ \relax
23:16 & \begin{tabularx}{0.7\textwidth}{X} 保羅的外甥聽見他們設下埋伏，就來到營樓裡告訴保羅。 \end{tabularx} \\ \\ \relax
23:17 & \begin{tabularx}{0.7\textwidth}{X} 保羅請一個百夫長來，說：「你領這青年去見千夫長，他有事告訴他。」 \end{tabularx} \\ \\ \relax
23:18 & \begin{tabularx}{0.7\textwidth}{X} 於是百夫長把他領去見千夫長，說：「被囚的保羅請我到他那裡，求我領這青年來見你；他有事告訴你。」 \end{tabularx} \\ \\ \relax
23:19 & \begin{tabularx}{0.7\textwidth}{X} 千夫長就拉著他的手，走到一旁，私下問他：「你有甚麼事告訴我呢？」 \end{tabularx} \\ \\ \relax
23:20 & \begin{tabularx}{0.7\textwidth}{X} 他說：「猶太人已經約定，要求你明天把保羅帶到議會去，假裝要詳細查問他的事。 \end{tabularx} \\ \\ \relax
23:21 & \begin{tabularx}{0.7\textwidth}{X} 你切不要隨從他們，因為他們有四十多人埋伏，已經起誓，若不先殺掉保羅就不吃不喝。現在都預備好了，只等你的允准。」 \end{tabularx} \\ \\ \relax
23:22 & \begin{tabularx}{0.7\textwidth}{X} 於是千夫長打發那青年走，囑咐他：「不要告訴人，你已將這些事報告我了。」 \end{tabularx} \\ \\ \relax
23:23 & \begin{tabularx}{0.7\textwidth}{X} 於是，千夫長叫了兩個百夫長來，說：「預備步兵二百、騎兵七十、長槍手二百，今夜九點往凱撒利亞去； \end{tabularx} \\ \\ \relax
23:24 & \begin{tabularx}{0.7\textwidth}{X} 也要預備牲口讓保羅騎上，護送到腓力斯總督那裡去。」 \end{tabularx} \\ \\ \relax
23:25 & \begin{tabularx}{0.7\textwidth}{X} 千夫長又寫了公文，大略說： \end{tabularx} \\ \\ \relax
23:26 & \begin{tabularx}{0.7\textwidth}{X} 「克勞第‧呂西亞向腓力斯總督大人請安。 \end{tabularx} \\ \\ \relax
23:27 & \begin{tabularx}{0.7\textwidth}{X} 這個人被猶太人拿住，快被殺害時，我得知他是羅馬人，就帶士兵下去，把他救了出來。 \end{tabularx} \\ \\ \relax
23:28 & \begin{tabularx}{0.7\textwidth}{X} 因為我要知道他們告他的罪狀，就帶他下到他們的議會去。 \end{tabularx} \\ \\ \relax
23:29 & \begin{tabularx}{0.7\textwidth}{X} 我查知他被告發是因他們律法上的爭論，並沒有甚麼該死或該監禁的罪名。 \end{tabularx} \\ \\ \relax
23:30 & \begin{tabularx}{0.7\textwidth}{X} 後來有人把要害他的計謀告訴我，我立刻把他解到你那裡去，又命令告他的人在你面前告他。」 \end{tabularx} \\ \\ \relax
23:31 & \begin{tabularx}{0.7\textwidth}{X} 於是士兵照所命令他們的，連夜把保羅帶到安提帕底。 \end{tabularx} \\ \\ \relax
23:32 & \begin{tabularx}{0.7\textwidth}{X} 第二天，由騎兵護送保羅，他們就回營樓去。 \end{tabularx} \\ \\ \relax
23:33 & \begin{tabularx}{0.7\textwidth}{X} 騎兵來到凱撒利亞，把公文呈給總督，就叫保羅站在他面前。 \end{tabularx} \\ \\ \relax
23:34 & \begin{tabularx}{0.7\textwidth}{X} 總督讀了公文，問保羅是哪一省的人；一知道他是基利家人， \end{tabularx} \\ \\ \relax
23:35 & \begin{tabularx}{0.7\textwidth}{X} 就說：「等告你的人來到，我才詳細聽你。」於是他命令把保羅拘留在希律的衙門裡。 \end{tabularx} \\ \\ \relax
24:1 & \begin{tabularx}{0.7\textwidth}{X} 過了五天，亞拿尼亞大祭司、幾個長老和一個叫帖土羅的律師下來，向總督控告保羅。 \end{tabularx} \\ \\ \relax
24:2 & \begin{tabularx}{0.7\textwidth}{X} 保羅一被傳來，帖土羅就開始控告他，說：「腓力斯大人，我們因你得以享受國泰民安，並且這一國的弊病，因著你的遠見得以改革。 \end{tabularx} \\ \\ \relax
24:3 & \begin{tabularx}{0.7\textwidth}{X} 我們隨時隨地都滿心感激不盡。 \end{tabularx} \\ \\ \relax
24:4 & \begin{tabularx}{0.7\textwidth}{X} 為了不敢耽擱你太久，我只求你寬容一下，聽我們說幾句話。 \end{tabularx} \\ \\ \relax
24:5 & \begin{tabularx}{0.7\textwidth}{X} 我們看這個人如同瘟疫一般，是鼓動普天下所有的猶太人作亂的人，又是拿撒勒教派裡的一個頭目。 \end{tabularx} \\ \\ \relax
24:6 & \begin{tabularx}{0.7\textwidth}{X} 他甚至連聖殿也要污穢，我們就把他捉拿了。要按我們的律法審問， \end{tabularx} \\ \\ \relax
24:7 & \begin{tabularx}{0.7\textwidth}{X} 可是，呂西亞千夫長前來，甚是強橫，從我們手中把他奪去，下令告他的人到你這裡來。 \end{tabularx} \\ \\ \relax
24:8 & \begin{tabularx}{0.7\textwidth}{X} 你自己審問他，就可以知道我們所控告他的一切事了。」 \end{tabularx} \\ \\ \relax
24:9 & \begin{tabularx}{0.7\textwidth}{X} 眾猶太人也隨著控告他，說：「這些事情確是這樣。」 \end{tabularx} \\ \\ \relax
24:10 & \begin{tabularx}{0.7\textwidth}{X} 總督示意叫保羅說話，保羅就回答：「我知道你在本國作法官多年，所以我樂意為自己申辯。 \end{tabularx} \\ \\ \relax
24:11 & \begin{tabularx}{0.7\textwidth}{X} 你查問就可以知道，從我上耶路撒冷去禮拜到今日不過十二天。 \end{tabularx} \\ \\ \relax
24:12 & \begin{tabularx}{0.7\textwidth}{X} 他們並沒有看見我在聖殿裡跟人辯論，或在會堂裡、在城裡煽動群眾。 \end{tabularx} \\ \\ \relax
24:13 & \begin{tabularx}{0.7\textwidth}{X} 也不能對你證實他們現在所控告我的事。 \end{tabularx} \\ \\ \relax
24:14 & \begin{tabularx}{0.7\textwidth}{X} 但有一件事我向你承認，就是我正按著他們所稱為異端的道事奉我祖宗的神，又信合乎律法和先知書上所記載的一切。 \end{tabularx} \\ \\ \relax
24:15 & \begin{tabularx}{0.7\textwidth}{X} 我對神存著這些人自己也接受的盼望，就是義人和不義的人都要復活。 \end{tabularx} \\ \\ \relax
24:16 & \begin{tabularx}{0.7\textwidth}{X} 因此，我勉勵自己，對神對人，時常存著無虧的良心。 \end{tabularx} \\ \\ \relax
24:17 & \begin{tabularx}{0.7\textwidth}{X} 過了幾年，我帶著賙濟本國的捐項和供物上去。 \end{tabularx} \\ \\ \relax
24:18 & \begin{tabularx}{0.7\textwidth}{X} 正獻的時候，他們看見我在聖殿裡已經潔淨了，並沒有聚眾，也沒有吵嚷， \end{tabularx} \\ \\ \relax
24:19 & \begin{tabularx}{0.7\textwidth}{X} 惟有幾個從亞細亞來的猶太人—他們若有控告我的事，應當到你面前來告我。 \end{tabularx} \\ \\ \relax
24:20 & \begin{tabularx}{0.7\textwidth}{X} 不然，讓這些人自己說，他們看出我站在議會前的時候，有甚麼不對的地方。 \end{tabularx} \\ \\ \relax
24:21 & \begin{tabularx}{0.7\textwidth}{X} 縱然有，也不過是為了一句話，就是我站在他們中間喊說：『我今日在你們面前受審，是為了死人復活。』」 \end{tabularx} \\ \\ \relax
24:22 & \begin{tabularx}{0.7\textwidth}{X} 腓力斯本是詳細認識這道，就拖延他們，說：「且等呂西亞千夫長下來，我再審判你們的案。」 \end{tabularx} \\ \\ \relax
24:23 & \begin{tabularx}{0.7\textwidth}{X} 於是他下令百夫長看守保羅，要從寬待他，不可攔阻他的親友來供給他。 \end{tabularx} \\ \\ \relax
24:24 & \begin{tabularx}{0.7\textwidth}{X} 過了幾天，腓力斯和他夫人猶太女子土西拉一同來到，就叫保羅來，聽他講論信基督耶穌的事。 \end{tabularx} \\ \\ \relax
24:25 & \begin{tabularx}{0.7\textwidth}{X} 保羅講論公義、節制和將來的審判，腓力斯害怕起來，就回答：「你暫且去吧！等我有機會時再來叫你。」 \end{tabularx} \\ \\ \relax
24:26 & \begin{tabularx}{0.7\textwidth}{X} 腓力斯又指望保羅送他銀錢，所以屢次叫他來，和他談論。 \end{tabularx} \\ \\ \relax
24:27 & \begin{tabularx}{0.7\textwidth}{X} 過了兩年，波求‧非斯都接了腓力斯的任；腓力斯要討猶太人的喜歡，就把保羅留在監裡。 \end{tabularx} \\ \\ \relax
25:1 & \begin{tabularx}{0.7\textwidth}{X} 非斯都到省裡上任，過了三天，就從凱撒利亞上耶路撒冷去。 \end{tabularx} \\ \\ \relax
25:2 & \begin{tabularx}{0.7\textwidth}{X} 祭司長和猶太人的領袖向他控告保羅；又央求他， \end{tabularx} \\ \\ \relax
25:3 & \begin{tabularx}{0.7\textwidth}{X} 向他求情要對付保羅，把他提到耶路撒冷來，他們要在路上埋伏殺害他。 \end{tabularx} \\ \\ \relax
25:4 & \begin{tabularx}{0.7\textwidth}{X} 非斯都就回答：「保羅押在凱撒利亞，我自己快要往那裡去。」 \end{tabularx} \\ \\ \relax
25:5 & \begin{tabularx}{0.7\textwidth}{X} 他又說：「所以，你們中間有權的人與我一同下去，那人若有甚麼不是，就讓他們控告他。」 \end{tabularx} \\ \\ \relax
25:6 & \begin{tabularx}{0.7\textwidth}{X} 非斯都在他們那裡住了不超過八天或十天，就下凱撒利亞去；第二天開庭，下令把保羅提上來。 \end{tabularx} \\ \\ \relax
25:7 & \begin{tabularx}{0.7\textwidth}{X} 保羅來了，那些從耶路撒冷下來的猶太人周圍站著，提出許多嚴重而不能證實的事控告他。 \end{tabularx} \\ \\ \relax
25:8 & \begin{tabularx}{0.7\textwidth}{X} 保羅申辯說：「無論猶太人的律法，或是聖殿，或是凱撒，我都沒有干犯。」 \end{tabularx} \\ \\ \relax
25:9 & \begin{tabularx}{0.7\textwidth}{X} 但非斯都要討猶太人的喜歡，就回答保羅說：「你願意上耶路撒冷去，在那裡為這些事受我的審判嗎？」 \end{tabularx} \\ \\ \relax
25:10 & \begin{tabularx}{0.7\textwidth}{X} 保羅說：「我現在站在凱撒的審判臺前，這就是我應當受審的地方。我並沒有對猶太人做過甚麼不對的事，這也是你明明知道的。 \end{tabularx} \\ \\ \relax
25:11 & \begin{tabularx}{0.7\textwidth}{X} 我若做了不對的事，犯了甚麼該死的罪，就是死我也不辭。他們所控告我的事若都不實，就沒有人能把我交給他們。我要向凱撒上訴。」 \end{tabularx} \\ \\ \relax
25:12 & \begin{tabularx}{0.7\textwidth}{X} 非斯都和議會商量了，就回答：「既然你要向凱撒上訴，你就到凱撒那裡去吧。」 \end{tabularx} \\ \\ \relax
25:13 & \begin{tabularx}{0.7\textwidth}{X} 過了些日子，亞基帕王和百妮基來到凱撒利亞，拜訪非斯都。 \end{tabularx} \\ \\ \relax
25:14 & \begin{tabularx}{0.7\textwidth}{X} 他們在那裡住了好些日子，非斯都將保羅的案件向王陳述，說：「這裡有一個人，是腓力斯留在監裡的。 \end{tabularx} \\ \\ \relax
25:15 & \begin{tabularx}{0.7\textwidth}{X} 我在耶路撒冷的時候，祭司長和猶太的長老把他的事稟報了，要求定他的罪。 \end{tabularx} \\ \\ \relax
25:16 & \begin{tabularx}{0.7\textwidth}{X} 我回覆他們，無論甚麼人，被告還沒有和原告當面對質，沒有機會為所控告的事申辯，就先定他罪的，這不是羅馬人的規矩。 \end{tabularx} \\ \\ \relax
25:17 & \begin{tabularx}{0.7\textwidth}{X} 及至他們都來到這裡，我沒有耽誤，第二天就開庭，下令把那人提上來。 \end{tabularx} \\ \\ \relax
25:18 & \begin{tabularx}{0.7\textwidth}{X} 控告他的人站起來告他，所控告的並沒有任何我所預料的那等惡事。 \end{tabularx} \\ \\ \relax
25:19 & \begin{tabularx}{0.7\textwidth}{X} 不過，有幾樣辯論是有關他們自己敬鬼神的事，以及一個名叫耶穌的人，他已經死了，保羅卻說他是活著的。 \end{tabularx} \\ \\ \relax
25:20 & \begin{tabularx}{0.7\textwidth}{X} 我對這些事不知該怎樣處理，所以問他是否願意上耶路撒冷去，在那裡為這些事接受審判。 \end{tabularx} \\ \\ \relax
25:21 & \begin{tabularx}{0.7\textwidth}{X} 但保羅要求我留下他，要聽皇上判斷，我就下令把他留下，等我解他到凱撒那裡去。」 \end{tabularx} \\ \\ \relax
25:22 & \begin{tabularx}{0.7\textwidth}{X} 亞基帕對非斯都說：「我也願意親自聽聽這個人。」非斯都說：「明天你就可以聽他。」 \end{tabularx} \\ \\ \relax
25:23 & \begin{tabularx}{0.7\textwidth}{X} 第二天，亞基帕和百妮基大張旗鼓而來，與眾千夫長和城裡的顯要進了大廳。非斯都一聲令下，就有人將保羅帶進來。 \end{tabularx} \\ \\ \relax
25:24 & \begin{tabularx}{0.7\textwidth}{X} 非斯都說：「亞基帕王和在這裡的諸位，你們看這個人，他就是所有在耶路撒冷和這裡的猶太人曾向我懇求呼叫，說不可容他再活著的。 \end{tabularx} \\ \\ \relax
25:25 & \begin{tabularx}{0.7\textwidth}{X} 但我查明他並沒有犯甚麼該死的罪，並且他自己也已向皇帝上訴了，所以我決定把他解去。 \end{tabularx} \\ \\ \relax
25:26 & \begin{tabularx}{0.7\textwidth}{X} 論到這個人，我沒有確實的事可以奏明主上。因此，我帶他到你們面前，尤其到你亞基帕王面前，為要在查問之後有所呈奏。 \end{tabularx} \\ \\ \relax
25:27 & \begin{tabularx}{0.7\textwidth}{X} 因為據我看，解送囚犯而不指明他的罪狀是不合理的。」 \end{tabularx} \\ \\ \relax
26:1 & \begin{tabularx}{0.7\textwidth}{X} 亞基帕對保羅說：「准你為自己申訴。」於是保羅伸手辯護說： \end{tabularx} \\ \\ \relax
26:2 & \begin{tabularx}{0.7\textwidth}{X} 「亞基帕王啊，猶太人所控告我的一切事，今日得以在你面前辯護，實為萬幸。 \end{tabularx} \\ \\ \relax
26:3 & \begin{tabularx}{0.7\textwidth}{X} 更慶幸的是你熟悉猶太人的規矩和他們的爭論；所以，求你耐心聽我。 \end{tabularx} \\ \\ \relax
26:4 & \begin{tabularx}{0.7\textwidth}{X} 「我自幼為人如何，從起初在本國的同胞中，以及在耶路撒冷，所有的猶太人都知道。 \end{tabularx} \\ \\ \relax
26:5 & \begin{tabularx}{0.7\textwidth}{X} 他們若肯作見證，就知道我從起初是按著我們教中最嚴緊的教門作了法利賽人。 \end{tabularx} \\ \\ \relax
26:6 & \begin{tabularx}{0.7\textwidth}{X} 現在我站在這裡受審，是為了對神向我們祖宗的應許存著盼望。 \end{tabularx} \\ \\ \relax
26:7 & \begin{tabularx}{0.7\textwidth}{X} 這應許，我們十二個支派，晝夜切切地事奉神，都指望得著。王啊，我正是因這指望被猶太人控告。 \end{tabularx} \\ \\ \relax
26:8 & \begin{tabularx}{0.7\textwidth}{X} 神使死人復活，你們為甚麼判斷為不可信呢？ \end{tabularx} \\ \\ \relax
26:9 & \begin{tabularx}{0.7\textwidth}{X} 「從前我自己認為必須竭力反對拿撒勒人耶穌的名， \end{tabularx} \\ \\ \relax
26:10 & \begin{tabularx}{0.7\textwidth}{X} 我在耶路撒冷也曾這樣做過；我不但從祭司長得了權柄，把許多聖徒收在監裡，而且他們被殺，我也表示贊成。 \end{tabularx} \\ \\ \relax
26:11 & \begin{tabularx}{0.7\textwidth}{X} 在各會堂，我屢次用刑強迫他們說褻瀆的話，我非常厭惡他們，甚至追逼他們，直到外邦的城鎮。」 \end{tabularx} \\ \\ \relax
26:12 & \begin{tabularx}{0.7\textwidth}{X} 「那時，我帶著祭司長的權柄和命令往大馬士革去。 \end{tabularx} \\ \\ \relax
26:13 & \begin{tabularx}{0.7\textwidth}{X} 王啊！我在路上，中午的時候，看見從天上有一道光，比太陽還亮，四面照射著我和跟我同行的人。 \end{tabularx} \\ \\ \relax
26:14 & \begin{tabularx}{0.7\textwidth}{X} 我們都仆倒在地，我就聽見有聲音用希伯來話對我說：『掃羅！掃羅！你為甚麼迫害我？你用腳踢刺棒是自找苦吃的！』 \end{tabularx} \\ \\ \relax
26:15 & \begin{tabularx}{0.7\textwidth}{X} 我說：『主啊，你是誰？』主說：『我就是你所迫害的耶穌。 \end{tabularx} \\ \\ \relax
26:16 & \begin{tabularx}{0.7\textwidth}{X} 起來，站著，我向你顯現的目的是要派你作僕役，為你所看見我的事，和我將要指示你的事作見證人。 \end{tabularx} \\ \\ \relax
26:17 & \begin{tabularx}{0.7\textwidth}{X} 我也要救你脫離百姓和外邦人的手。我差你到他們那裡去， \end{tabularx} \\ \\ \relax
26:18 & \begin{tabularx}{0.7\textwidth}{X} 要開他們的眼睛，使他們從黑暗中轉向光明，從撒但權下歸向神；使他們因信我而得蒙赦罪，和一切成聖的人同得基業。』」 \end{tabularx} \\ \\ \relax
26:19 & \begin{tabularx}{0.7\textwidth}{X} 「因此，亞基帕王啊！我沒有違背那從天上來的異象； \end{tabularx} \\ \\ \relax
26:20 & \begin{tabularx}{0.7\textwidth}{X} 我先在大馬士革，後在耶路撒冷和猶太全地，以及外邦，勸勉他們應當悔改歸向神，行事與悔改的心相稱。 \end{tabularx} \\ \\ \relax
26:21 & \begin{tabularx}{0.7\textwidth}{X} 為這緣故，猶太人在聖殿裡拿住我，想要殺我。 \end{tabularx} \\ \\ \relax
26:22 & \begin{tabularx}{0.7\textwidth}{X} 然而，我蒙神的幫助，直到今日還站立得穩，向尊貴的和卑微的作見證。我所講的，並不外乎眾先知和摩西所說將來必成的事， \end{tabularx} \\ \\ \relax
26:23 & \begin{tabularx}{0.7\textwidth}{X} 就是基督必須受害，並且首先從死人中復活，把亮光傳給猶太人和外邦人。」 \end{tabularx} \\ \\ \relax
26:24 & \begin{tabularx}{0.7\textwidth}{X} 保羅這樣申訴時，非斯都大聲說：「保羅，你瘋了！你的學問太大，反使你瘋了！」 \end{tabularx} \\ \\ \relax
26:25 & \begin{tabularx}{0.7\textwidth}{X} 保羅說：「非斯都大人，我不是瘋了，我說的乃是真實和清醒的話。 \end{tabularx} \\ \\ \relax
26:26 & \begin{tabularx}{0.7\textwidth}{X} 王也知道這些事，所以對王大膽直言，我深信這些事沒有一件能向王隱瞞的，因為都不是在背地裡做的。 \end{tabularx} \\ \\ \relax
26:27 & \begin{tabularx}{0.7\textwidth}{X} 亞基帕王啊，你信先知嗎？我知道你是信的。」 \end{tabularx} \\ \\ \relax
26:28 & \begin{tabularx}{0.7\textwidth}{X} 亞基帕對保羅說：「你想稍微勸一勸就能說服我作基督徒了嗎？」 \end{tabularx} \\ \\ \relax
26:29 & \begin{tabularx}{0.7\textwidth}{X} 保羅說：「無論少勸還是多勸，我向神所求的，不但你一個人，就是今天所有聽我說話的人都要像我一樣，只是不要有這些鎖鏈。」 \end{tabularx} \\ \\ \relax
26:30 & \begin{tabularx}{0.7\textwidth}{X} 於是，王和總督以及百妮基跟同坐的人都站起來， \end{tabularx} \\ \\ \relax
26:31 & \begin{tabularx}{0.7\textwidth}{X} 退到裡面，彼此談論說：「這個人並沒有犯甚麼該死該監禁的罪。」 \end{tabularx} \\ \\ \relax
26:32 & \begin{tabularx}{0.7\textwidth}{X} 亞基帕對非斯都說：「這人若沒有向凱撒上訴，早就被釋放了。」 \end{tabularx} \\ \\ \relax
27:1 & \begin{tabularx}{0.7\textwidth}{X} 既然非斯都決定要我們坐船往意大利去，就將保羅和別的囚犯交給御營裡的一個名叫猶流的百夫長。 \end{tabularx} \\ \\ \relax
27:2 & \begin{tabularx}{0.7\textwidth}{X} 有一隻亞大米田的船要開往亞細亞沿海一帶地方去，我們上了那船，就起航了；有馬其頓的帖撒羅尼迦人亞里達古和我們同去。 \end{tabularx} \\ \\ \relax
27:3 & \begin{tabularx}{0.7\textwidth}{X} 第二天，我們到了西頓。猶流寬待保羅，准他往朋友那裡去，受他們的照應。 \end{tabularx} \\ \\ \relax
27:4 & \begin{tabularx}{0.7\textwidth}{X} 我們又從那裡開船，因為遇到逆風，就貼著塞浦路斯的背風岸航行， \end{tabularx} \\ \\ \relax
27:5 & \begin{tabularx}{0.7\textwidth}{X} 渡過了基利家、旁非利亞一帶的海面，就到了呂家的每拉。 \end{tabularx} \\ \\ \relax
27:6 & \begin{tabularx}{0.7\textwidth}{X} 在那裡，百夫長找到一隻亞歷山大的船要往意大利去，就叫我們上了那船。 \end{tabularx} \\ \\ \relax
27:7 & \begin{tabularx}{0.7\textwidth}{X} 一連多日，船行得很慢，我們好不容易才來到革尼土的對面；又因被風攔阻，我們就貼著克里特島背風岸，從撒摩尼對面航行。 \end{tabularx} \\ \\ \relax
27:8 & \begin{tabularx}{0.7\textwidth}{X} 我們沿岸前進，十分艱難，來到一個名叫佳澳的地方，離那裡不遠有拉西亞城。 \end{tabularx} \\ \\ \relax
27:9 & \begin{tabularx}{0.7\textwidth}{X} 航行的日子久了，已經過了禁食的節期，行船又危險，保羅就建議， \end{tabularx} \\ \\ \relax
27:10 & \begin{tabularx}{0.7\textwidth}{X} 對眾人說：「諸位，我看這次航行，不但貨物和船要受損傷，大遭破壞，連我們的性命也難保。」 \end{tabularx} \\ \\ \relax
27:11 & \begin{tabularx}{0.7\textwidth}{X} 但百夫長信從船長和船主，不信保羅所說的。 \end{tabularx} \\ \\ \relax
27:12 & \begin{tabularx}{0.7\textwidth}{X} 且因在這港口不適宜過冬，船上大多數的人都主張開船離開這地方，或者能到非尼基去過冬。非尼基是克里特的一個港口，一面朝西南，一面朝西北。 \end{tabularx} \\ \\ \relax
27:13 & \begin{tabularx}{0.7\textwidth}{X} 當南風微微吹起時，他們以為對目的地已有了把握，就起錨，貼近克里特開去。 \end{tabularx} \\ \\ \relax
27:14 & \begin{tabularx}{0.7\textwidth}{X} 過了不久，有一股叫「友拉革羅」的東北巨風從島上撲來， \end{tabularx} \\ \\ \relax
27:15 & \begin{tabularx}{0.7\textwidth}{X} 船被風抓住，無法頂風航行，我們只好任它漂流。 \end{tabularx} \\ \\ \relax
27:16 & \begin{tabularx}{0.7\textwidth}{X} 我們貼著一個叫高大的小島的背風岸急航，好不容易才保住了救生艇。 \end{tabularx} \\ \\ \relax
27:17 & \begin{tabularx}{0.7\textwidth}{X} 既然把救生艇拉上來，他們就用纜索捆綁船底，又恐怕在賽耳底淺灘上擱淺，就落了篷，任船漂流。 \end{tabularx} \\ \\ \relax
27:18 & \begin{tabularx}{0.7\textwidth}{X} 我們被風浪逼得很急，第二天眾人就把貨物拋在海裡。 \end{tabularx} \\ \\ \relax
27:19 & \begin{tabularx}{0.7\textwidth}{X} 第三天，他們又親手把船上的器具拋棄了。 \end{tabularx} \\ \\ \relax
27:20 & \begin{tabularx}{0.7\textwidth}{X} 許多天都沒有看到太陽和星辰，又有狂風大浪催逼，我們獲救的指望都放棄了。 \end{tabularx} \\ \\ \relax
27:21 & \begin{tabularx}{0.7\textwidth}{X} 眾人已有好幾天沒有吃東西，保羅就出來站在他們中間，說：「諸位，你們本該聽我的話不離開克里特島，就不致遭到這樣的損失和破壞。 \end{tabularx} \\ \\ \relax
27:22 & \begin{tabularx}{0.7\textwidth}{X} 現在我勸你們放心，除了損失這條船，你們中間沒有一人會喪失性命。 \end{tabularx} \\ \\ \relax
27:23 & \begin{tabularx}{0.7\textwidth}{X} 因為昨夜，我所屬所事奉的神的使者站在我旁邊， \end{tabularx} \\ \\ \relax
27:24 & \begin{tabularx}{0.7\textwidth}{X} 說：『保羅，不要害怕，你必定站在凱撒面前；並且神已把安全賜給與你同船的人了。』 \end{tabularx} \\ \\ \relax
27:25 & \begin{tabularx}{0.7\textwidth}{X} 所以，諸位可以放心，我信神怎樣對我說，事情也要怎樣成就； \end{tabularx} \\ \\ \relax
27:26 & \begin{tabularx}{0.7\textwidth}{X} 只是我們必須在一個島上擱淺。」 \end{tabularx} \\ \\ \relax
27:27 & \begin{tabularx}{0.7\textwidth}{X} 到了第十四天夜間，船在亞得里亞海漂來漂去。約在半夜，水手以為漸近旱地， \end{tabularx} \\ \\ \relax
27:28 & \begin{tabularx}{0.7\textwidth}{X} 就去探測深淺，探得有十二丈；稍往前行，又探深淺，探得有九丈。 \end{tabularx} \\ \\ \relax
27:29 & \begin{tabularx}{0.7\textwidth}{X} 恐怕我們撞到礁石，他們就從船尾拋下四個錨，盼望天亮。 \end{tabularx} \\ \\ \relax
27:30 & \begin{tabularx}{0.7\textwidth}{X} 水手想棄船逃走，把救生艇縋下海裡，假裝要從船頭拋錨的樣子。 \end{tabularx} \\ \\ \relax
27:31 & \begin{tabularx}{0.7\textwidth}{X} 保羅對百夫長和士兵說：「這些人若不留在船上，你們就不能獲救。」 \end{tabularx} \\ \\ \relax
27:32 & \begin{tabularx}{0.7\textwidth}{X} 於是士兵砍斷救生艇的繩子，由它漂去。 \end{tabularx} \\ \\ \relax
27:33 & \begin{tabularx}{0.7\textwidth}{X} 天快亮的時候，保羅勸眾人都用餐，說：「你們一直捱餓等候，不吃甚麼，已經十四天了。 \end{tabularx} \\ \\ \relax
27:34 & \begin{tabularx}{0.7\textwidth}{X} 所以我勸你們吃點東西，這是關乎你們獲救的，因為你們各人連一根頭髮也不至於掉落。」 \end{tabularx} \\ \\ \relax
27:35 & \begin{tabularx}{0.7\textwidth}{X} 保羅說了這話，就拿起餅來，在眾人面前祝謝了神，然後擘開來吃。 \end{tabularx} \\ \\ \relax
27:36 & \begin{tabularx}{0.7\textwidth}{X} 於是他們都放心，就吃了。 \end{tabularx} \\ \\ \relax
27:37 & \begin{tabularx}{0.7\textwidth}{X} 我們在船上的共有二百七十六個人。 \end{tabularx} \\ \\ \relax
27:38 & \begin{tabularx}{0.7\textwidth}{X} 他們吃飽了，為要使船輕一點，就把船上的麥子拋到海裡。 \end{tabularx} \\ \\ \relax
27:39 & \begin{tabularx}{0.7\textwidth}{X} 天亮的時候，他們不認得那地方，只見一個有岸可登的海灣，就想法子看能不能把船靠岸。 \end{tabularx} \\ \\ \relax
27:40 & \begin{tabularx}{0.7\textwidth}{X} 於是他們砍斷纜索，把錨丟到海裡，同時也鬆開舵繩，拉起頭篷，順風向著岸行去。 \end{tabularx} \\ \\ \relax
27:41 & \begin{tabularx}{0.7\textwidth}{X} 但碰到兩水夾流的地方，就擱了淺，船頭膠住不動，船尾被浪的猛力衝壞了。 \end{tabularx} \\ \\ \relax
27:42 & \begin{tabularx}{0.7\textwidth}{X} 士兵的意思要把囚犯都殺了，免得有游水脫逃的。 \end{tabularx} \\ \\ \relax
27:43 & \begin{tabularx}{0.7\textwidth}{X} 但百夫長要救保羅，不准他們任意而行，就吩咐會游水的，跳下水去，先上岸； \end{tabularx} \\ \\ \relax
27:44 & \begin{tabularx}{0.7\textwidth}{X} 其餘的人則用板子或船的碎片上岸。這樣，眾人都獲救，上了岸。 \end{tabularx} \\ \\ \relax
28:1 & \begin{tabularx}{0.7\textwidth}{X} 我們既已獲救，才知道那島名叫馬耳他。 \end{tabularx} \\ \\ \relax
28:2 & \begin{tabularx}{0.7\textwidth}{X} 當地人非常友善地接待我們；因為正在下雨，天氣又冷，他們就生了火歡迎我們眾人。 \end{tabularx} \\ \\ \relax
28:3 & \begin{tabularx}{0.7\textwidth}{X} 那時，保羅拾起一捆柴，放在火中，有一條毒蛇，因為熱的緣故鑽了出來，纏住他的手。 \end{tabularx} \\ \\ \relax
28:4 & \begin{tabularx}{0.7\textwidth}{X} 當地的人看見那毒蛇懸在他手上，就彼此說：「這人必是個兇手，雖然他從海裡獲救，天理仍不容他活著。」 \end{tabularx} \\ \\ \relax
28:5 & \begin{tabularx}{0.7\textwidth}{X} 保羅竟把那毒蛇甩在火裡，並沒有受傷。 \end{tabularx} \\ \\ \relax
28:6 & \begin{tabularx}{0.7\textwidth}{X} 當地的人想他快要腫起來，或是忽然倒下死了，但等了好久，見他沒有甚麼異樣，就轉念說他是個神明。 \end{tabularx} \\ \\ \relax
28:7 & \begin{tabularx}{0.7\textwidth}{X} 離那地方不遠有一些田產，是島長部百流的。他接納我們，盡情款待了我們三日。 \end{tabularx} \\ \\ \relax
28:8 & \begin{tabularx}{0.7\textwidth}{X} 當時，部百流的父親臥病不起，患了熱病和痢疾。保羅進去見他，為他禱告按手，治好了他。 \end{tabularx} \\ \\ \relax
28:9 & \begin{tabularx}{0.7\textwidth}{X} 從此，島上其餘的病人也都來，得了醫治。 \end{tabularx} \\ \\ \relax
28:10 & \begin{tabularx}{0.7\textwidth}{X} 他們又多方面尊敬我們，到了開船的時候，又把我們所需用的東西送到船上。 \end{tabularx} \\ \\ \relax
28:11 & \begin{tabularx}{0.7\textwidth}{X} 過了三個月，我們上了亞歷山大的船起航。這船以「宙斯雙子」為記，是在那海島過冬的。 \end{tabularx} \\ \\ \relax
28:12 & \begin{tabularx}{0.7\textwidth}{X} 我們到了敘拉古，停泊了三日； \end{tabularx} \\ \\ \relax
28:13 & \begin{tabularx}{0.7\textwidth}{X} 又從那裡起錨開船，來到利基翁。過了一天，起了南風，第二天就來到部丟利。 \end{tabularx} \\ \\ \relax
28:14 & \begin{tabularx}{0.7\textwidth}{X} 我們在那裡遇見一些弟兄，他們請我們同住了七天。就這樣，我們來到羅馬。 \end{tabularx} \\ \\ \relax
28:15 & \begin{tabularx}{0.7\textwidth}{X} 那裡的弟兄們一聽見我們的消息，就到亞比烏市和三館來迎接我們。保羅見了他們，就感謝神，越發壯膽。 \end{tabularx} \\ \\ \relax
28:16 & \begin{tabularx}{0.7\textwidth}{X} 我們進了羅馬城，保羅蒙准和那個看守他的兵另住在一處。 \end{tabularx} \\ \\ \relax
28:17 & \begin{tabularx}{0.7\textwidth}{X} 過了三天，保羅請當地猶太人的領袖來。他們來了，保羅對他們說：「諸位弟兄，雖然我沒有做甚麼事干犯本國的百姓和我們祖宗的規矩，卻在耶路撒冷被囚禁，交在羅馬人的手裡。 \end{tabularx} \\ \\ \relax
28:18 & \begin{tabularx}{0.7\textwidth}{X} 他們審問了我，有意要釋放我，因為在我身上並沒有該死的罪狀。 \end{tabularx} \\ \\ \relax
28:19 & \begin{tabularx}{0.7\textwidth}{X} 但猶太人反對，我不得已只好上訴於凱撒，並不是有甚麼事要控告我本國的百姓。 \end{tabularx} \\ \\ \relax
28:20 & \begin{tabularx}{0.7\textwidth}{X} 為這緣故，我請你們來見我當面談話，我原是為以色列人所指望的那位才被這鐵鏈捆綁的。」 \end{tabularx} \\ \\ \relax
28:21 & \begin{tabularx}{0.7\textwidth}{X} 他們對他說：「我們並沒有接到從猶太寄來有關於你的信，也沒有弟兄到這裡來向我們報告，或說你有甚麼不好的地方。 \end{tabularx} \\ \\ \relax
28:22 & \begin{tabularx}{0.7\textwidth}{X} 但我們願意聽聽你的意見，因為我們知道這教門是到處遭人反對的。」 \end{tabularx} \\ \\ \relax
28:23 & \begin{tabularx}{0.7\textwidth}{X} 他們和保羅約定了日子，就有許多人到他的住處來。保羅從早到晚向他們講解這事，為神的國作證，並引摩西的律法和先知的書勸導他們信從耶穌。 \end{tabularx} \\ \\ \relax
28:24 & \begin{tabularx}{0.7\textwidth}{X} 他所說的話，有的信，有的不信。 \end{tabularx} \\ \\ \relax
28:25 & \begin{tabularx}{0.7\textwidth}{X} 他們間彼此不合，就分散了；未散以先，保羅說了一句話：「聖靈藉以賽亞先知向你們祖宗所說的話是對的。 \end{tabularx} \\ \\ \relax
28:26 & \begin{tabularx}{0.7\textwidth}{X} 他說：『你去對這百姓說：你們聽了又聽，卻不明白；看了又看，卻看不清。 \end{tabularx} \\ \\ \relax
28:27 & \begin{tabularx}{0.7\textwidth}{X} 因為這百姓的心麻木，耳朵塞著，眼睛閉著，免得眼睛看見，耳朵聽見，心裡明白，回轉過來，我會醫治他們。』 \end{tabularx} \\ \\ \relax
28:28 & \begin{tabularx}{0.7\textwidth}{X} 所以，你們當知道，神這救恩已經傳給外邦人；他們會聽的。」 \end{tabularx} \\ \\ \relax
28:29 & \begin{tabularx}{0.7\textwidth}{X} 保羅說了這些話，猶太人議論紛紛地走了。 \end{tabularx} \\ \\ \relax
28:30 & \begin{tabularx}{0.7\textwidth}{X} 保羅在自己所租的房子裡住了足足兩年。凡來見他的人，他都接待， \end{tabularx} \\ \\ \relax
28:31 & \begin{tabularx}{0.7\textwidth}{X} 放膽傳講神的國，並教導主耶穌基督的事，沒有人禁止。 \end{tabularx} \\ \\
[1ex]
\hline
\hline
\end{longtable}
$^{1}$我們一起同心祈禱.
我們親愛的天父上帝.
我們全敬畏的心.
來到您的面前.
願您的聖靈在我們心裡工作.
藉著聖經的說話.
來教導提醒我們.
讓我們深深的明白.
我們所信的.
是真實的.
我們的神.
也是值得我們去尊敬的.
在這個時刻.
這個時代.
我們的神.
給予我們的使命和責任.
求您.
藉著聖經的說話.
堅固我們.
也指引我們當行的路.
看清楚我們前面的方向.
祈禱奉耶穌基督的名.
阿們.
很開心能夠和大家見面.
我很感恩.
在聽別人講道的時候.
我也有領受.
我今天講的講道.
是從其他同道的分享中.
得到靈感的.
有一次.
有位牧者講道的時候.
提到《使徒行傳》第28章31節.
我們可以將這段.
我主要講兩段經文.
可以播出來給大家看.
28章31節.
這節經文裡面.
就是講.
放膽傳講神國的道.

$^{41}$將主耶穌基督的事教導人.
最後一個小節.
中文有五個字.
沒有人禁止.
這位牧師告訴我們.
這五個字.
在希臘聖經裡面.
他卻只有一隻字.
很特別.
這個講法.
令我很有興趣.
我問一些問題.
如果有一個人寫一本書.
最後那一句.
最後那一隻字.
會不會其實是有.
弦外之音.
或者有伏筆呢.
很自然我就會問另一個問題.
就在《使徒行傳》.
第一章第一節裡面.
會不會有一些字眼.
有一些觀念.
是和28章第31節.
最後那一句.
最後一隻字.
或者一個小段落.
是互相對應的呢.
這樣問也很合理.
你也聽過.
有首尾呼應這回事.
我就真的回去找.
你也可以一起找.
在手機或在聖經裡面.
我找到《使徒行傳》第一章第一節.
好像不太有關係.
提阿非羅大人等等.
我就一次一次看下去.
結果我看到第八節.
我們可以播放第一章第八節.

$^{81}$他說.
但聖靈降臨在你們身上.
你們就必得著能力.
並要在耶路撒冷.
猶太傳地 撒瑪利亞.
直到極作我的見證.
我就覺得好像《使徒行傳》第一章第八節.
和第28章第31節.
就好像一頭一尾.
而這個一頭一尾.
就真的說得通.
聖靈的能力.
在這班人身上.
令他們去為耶路撒冷作見證.
然後放膽.
傳講神教的道.
沒有人禁止.
很合理的.
很清楚的.
如果這個想法.
是說得通呢.
那又要問下一個問題.
那我就問下大家.
在你的印象之中.
當你聽到《使徒行傳》的時候.
你會想起.
你會想起《使徒行傳》的內容裡面.
最令你深刻的是什麼呢?.
不要緊.
說一下也可以的.
最篇幅.
最多是關於什麼?.
看來.
大聲一點.
聽不到?.
聖靈是吧?.
也對的.
其實呢.
如果我們用篇幅來看.
大家應該很記得.

$^{121}$有個保羅三次的宣教旅程.
對不對?.
從篇幅來說.
我剛才說篇幅的量來說.
那我又問一個問題.
我們一路學聖經的時候.
就說.
在《使徒行傳》裡面.
我們是看到保羅三次的宣教旅程.
但是如果我加上.
剛才所說的想法.
一頭一尾的時候.
《使徒行傳》的作者就是路加醫生.
路加醫生是不是有一個想法.
我嘗試去說.
保羅三次宣教旅程的時候.
我是有一個框架的.
我來去說出.
保羅怎樣去體現這個訊息.
就是聖靈降臨在他身上.
他去周圍傳福音.
並且是.
放膽傳講神的道.
沒有人禁止.
會不會都講得通.
當我這樣想的時候.
我就要去追尋下去.
所以我可以和大家說.
每一次聽到別人說.
我有收穫.
我就回去馬上去工作.
成為我下一次講章的內容.
就這樣.
我講到有三部分.
這次我要改變一下.
我就什麼都不寫出來.
因為.
小玲就識笑了.
說我整個都寫出來了.
這次就想和大家看三部分.

$^{161}$周牧師都和我說.
沒那麼長.
第一部分是聖靈的能力.
第二部分就是我們要.
我們放膽傳道.
第三就是沒有人能禁止.
其實是幫大家對這兩節經文.
有一個準確的認識.
我們看回第一章第八節.
聖靈降臨的時候.
我們看到.
我們看回第一章第八節聖靈降臨的時候.
我不再讀一次這節經文.
其實最令我.
震驚的.
因為我現在有時看經文.
我看回它的原文.
最震驚的就是.
見證這兩個字中文.
其實沒什麼的.
最多是說你講見證.
不是足夠的.
你要行動 最多是這樣.
沒有一個講字.
但當我看希臘文聖經的時候.
見證這個字的字根是.
殉道.
我看完之後.
我發覺怎麼可以.
這個令我很震撼的事情.
接著我就問一個問題.
這個.
見證的字根殉道.
這個意思.
在本身的聖經.
究竟出現了多少次.
它總共出現了十次.
有兩次在路教.
福音.
有一次在《史唐人傳》.

$^{201}$有一次在《帖撒羅尼迦加前書》.
剛才我已經講了.
路加醫生就是.
路加福音和《史唐人傳》的作者.
而這個作者竟然是.
在新約聖經裡面.
有十次關於.
見證.
字根是殉道這個字.
他使用的那個人.
有九次是路加醫生講的.
那你就問一個問題.
路加醫生.
對於見證的看法是怎樣呢.
嘩 我.
一邊想的時候我都有點害怕.
不知道你有沒有感覺到我那種害怕呢.
就是說我是不是真的.
明白到那個.
見證的那個含義.
和那個力度呢.
我就讀一節經文給大家聽.
就是《史唐人傳》第三章第十五節.
這節經文.
很簡單的就是說.
你們殺了那個生命的主.
神卻是叫他從死裡復活.
我們都是為這件事.
作見證.
很簡單的一句說話.
不過你放在背景裡面就不簡單了.
我就講回.
那些聖經的歷史.
當耶穌被釘十字架之後.
那班門徒是怎樣的.
很驚慌.
是躲起來的.
你不記得有一次.
多瑪說我摸過耶穌.
我不算數的 我不信的.

$^{241}$有一次就是他們.
門徒聚集在一間屋裡面.
還關上門.
但耶穌沒有開門.
就進去了.
進去就和多瑪說你摸我的釘.
我不是要講這個部分.
而是講.
那班門徒是很.
害怕的時候.
大家要知道那班門徒很害怕.
害怕什麼.
下一個會不會抓了我.
你明不明白.
好了就要躲起來.
我們講完這番說話.
跟著有耶穌和他們吃早餐.
福音書尾段很多時候.
這些記載 叫耶穌顯現的記載.
好了就算你看到耶穌復活.
是不是你馬上可以.
由很害怕.
變成不害怕呢.
你問這個問題.
如果我們再留心看的時候.
無論馬太福音.
最後的大使命的經文.
路加福音或小行主義的經文.
都是講耶穌升天.
聖靈降臨.
這班人開始有轉變.
好了你想他轉變多極.
最多我都是走開門.
跟隔壁屋的人說.
耶穌復活了.
很小聲說 是不是很正常.
但如果你細心看小行主義的時候.
第一章的時候.
已經在講那班人.
去哪裡跟人說.

$^{281}$這班人就走去耶路撒冷的聖殿.
跟人說.
耶穌已經復活了.
大家要記得這個歷史和地理.
我想你們有人去過聖地.
是不是.
你知道耶路撒冷聖殿.
其實有一個地方.
就是猶太公會的辦事處.
有沒有印象.
如果你一想到這個地理的時候.
你就會說.
這班門徒有什麼事.
勇喔.
走去耶路撒冷聖殿都不是.
經文記載.
在所羅門郎的郎子那裡.
在聖殿的走廊.
跟人說耶穌復活.
那裡有什麼.
那裡有猶太公會的總部.
那你想一想.
你這樣去做.
有什麼後果.
是不是很高風險.
很高風險.
經文就說你們殺了生命的主.
你再想一想他在什麼處境去講.
在耶路撒冷聖殿那裡.
跟人說.
大家都會想到.
他不被人抓都算偷笑.
對不對.
我們會問為什麼他們有這麼大的轉變.
由初初耶穌被抓.
之後躲在房子裡.
到現在.
史萊姆傳第一章.
講到聖靈降臨之後.
跟著就走出去.

$^{321}$還要去耶路撒冷的聖殿.
去講耶穌復活的事.
所以如果我們看的時候.
他們走出去講耶穌復活的事.
基本上就是.
用他們的.
整個生命.
和他們的.
安全來做代價.
我這樣講.
講歷史就是真的這樣.
所以我讀完之後也有種震撼感.
就是說.
其實我們今天講見證.
就算行動作見證.
坦白一點說.
我們有很大的程度.
都不會跟得上.
耶穌基督那幫門徒.
你明白我的意思嗎.
所以.
我也不知道.
有時候有這樣的發現是好.
還是沒有這樣的發現是好.
有這樣的發現就很開心.
好像明白了聖經.
但是不好的地方就是.
不知就沒事.
知道了就大件事了.
知道了就覺得自己不行了.
那裡才是好驚.
你明白嗎.
所以那種矛盾的心理.
是很大的.
第二.
我就想跟大家講放膽傳道.
其實.
我就讀另一節經文.
在《四章31》那裡講.
禱告完了.

$^{361}$聚會的地方是震動.
他們就被聖靈充滿.
放膽傳講神的道.
在這裡我自己會覺得.
我沒有試過.
這些經驗.
我沒有試過這些經驗.
所以我不能夠去講.
禱告完了地震動.
我又吹很有能力.
但我很想講.
其實你看回保羅的人生.
裡面.
你可以看到.
他怎樣被聖靈充滿.
因為我們常常講聖靈充滿.
多數是第二章.
可以講別國的華麗.
其他的謊言等等.
其實不是那麼簡單.
因為如果我們有一點.
保羅的人生歷史的資料.
你可以看到.
保羅大概在主後.
34年是信主的.
之前他是迫基督徒.
然後.
在47到57年.
他去開始.
宣教的黃金期.
之前他要去.
在沙漠裡面被磨練.
學習等等.
然後.
到第57到62年.
是坐監期.
62到67年.
在羅馬被軟禁.
其實是準備殉道.
所以保羅他.

$^{401}$被聖靈充滿.
絕不是一次兩次的事件.
不是做一兩次轟轟烈烈的事.
而是長時間.
去面對挑戰.
如果你有機會.
再看回《史唐傳》.
你看到保羅去傳福音.
去到一個城市.
特別是大掃那邊.
過小阿實大掃是他的鄉下.
去到一個地方.
找找那裡有沒有會堂.
猶太人散居在別處.
就建不成聖殿就會建會堂.
如果城市沒有會堂.
那怎麼辦呢?就找去那些.
猶太人聚居的地方.
好了你可以想像.
保羅這種宣教的做法.
是會有什麼結果呢?.
想都想到.
要麼就有人歡迎你.
聽聽都沒有問題新的宗教.
但是有些人會怎樣.
很敵黨保羅.
所以看《史唐傳》裡面.
是有人想用石頭扔死他.
有人想.
合謀去害死他.
你看到這些事情.
絕對不是一件浪漫的事.
還有.
保羅那時候.
沒有人那麼方便.
攀山涉水的時候.
是有危險的.
有山賊有土匪.
所以當我們去看的時候.
他憑什麼去繼續.

$^{441}$這種的服事呢?.
當然你可以說.
他以前迫害基督徒.
後來他在大馬骰的路上.
遇見了主的榮光.
主對他說.
你為什麼迫害我.
從此他就改變人生.
你可以說這件事.
令他一生很震撼.
但是他明白了.
主是神的兒子就好了.
他可以不用這樣.
但是他卻要這樣.
去事奉.
如果你看保羅的書信.
他會說我不以聖名為念.
你看到.
記得這些經文.
他是這樣放膽去傳道.
甚至他在.
羅馬被軟禁的時候.
有些人想著去鼓勵他.
到回頭那些人.
被保羅去鼓勵.
好了說到這裡.
我就想下地了.
保羅這麼偉大我們學不到.
有什麼我想和大家說的呢.
其實呢.
我們都活在一個多元宗教的世界.
在香港地.
這件事和我們的信仰有什麼關係呢.
我就說一下我小小的背景.
我記得.
就是.
我曾經在.
崇基學院那裡.
我都讀過一些哲學.
宗教 神學等等.

$^{481}$讀哲學.
有西方哲學 中國哲學.
宗教.
當然有新的宗教研究方法.
當時80年代.
我有一次就上了去新亞.
修一科叫人類學.
我都不明白我為什麼會修一科這樣的東西.
不知道學分多了還是什麼.
我學了一件事.
原來雲南的少數民族.
還有一種文化人類學的思想.
他告訴我們.
每種文化都有它的文化符號.
你不要先說他迷信.
不要歧視他.
而這些說法今天很多人都懂.
我說這麼厲害.
這個說法對於宗教研究有什麼影響.
其實有影響的.
當時80年代已經開始有人研究民間宗教.
在台灣就研究媽祖.
在香港就研究王大仙.
什麼都有.
其他大學都有人做這些研究.
你可以找到這些論文.
他們說的就是我們太多先入為主的說法.
說他們迷信.
但其實有社會功能.
燒災解難的意識.
但其實有社會功能.
燒災解難的意識.
不過我有種性情.
我不是很習慣去廟裡做田野調查.
因為他們做這些要去拍.
我說不是很適合.
所以我就不走這條路.
但是這件事影響什麼.
就是說民間宗教我們都不要隨便批評.
說明我們更加不要批評.

$^{521}$你又不懂又怎樣.
當時我都讀了一些東西.
但那個影響是什麼.
我是想知道.
但我又是基督徒.
我又是想知道.
但我又是基督徒.
我們不一定要這樣走人生.
有很多不同的思考.
但我又成為基督徒.
我又要讀神學.
其實兩者是一個擺動.
好像有一個人.
現在這本書可能沒有了.
林語堂不知道有沒有人聽過.
他有一本書叫信仰之旅.
就是說他到了晚年.
他回到基督教的懷抱.
事實上有很多人.
他的人生就是.
他接觸很多不同的宗教.
很多不同的思想.
他享受其中.
或者遊走的其中.
或者你的親戚朋友可能是這種人.
都不一定.
在那個階段他開始發現.
人的宗教.
人的智慧.
和神的智慧是兩樣東西.
就是這個.
說到這件事.
和放膽傳道有什麼關係呢.
其實我們的前人.
一直以來.
都在做這件工作.
自從基督教來到中國.
一直都有.
我怕我越說越興起.
有很多東西要說.

$^{561}$唐朝有個景教.
景教.
其實是一個.
大秦景教碑.
大秦其實就是羅馬.
羅馬帝國的一個宗教.
其實就是基督教.
他只不過是基督教的一個分支.
來到中國.
所以其實一直以來.
福音都有去到中國.
再近期一點.
近期一點.
有一樣東西可能你一定都知道.
1842年南京條約.
是不是.
但你知不知道英軍是哪一年.
進入香港的.
其實是1841年.
1841年英軍進入香港之後.
澳門又派了.
幾位神父.
過來香港考察.
結果你們查證一下.
香港天主教的.
教區是哪一年成立的.
是1842年.
香港聖公會是1843年.
有跡可尋.
故事是甚麼呢.
就是說當澳門.
駐澳門的宣教士.
來香港看看牧養的可能性的時候.
他們就制定了一些計劃.
所以南京條約.
其實不是.
即是他簽之前已經知道要簽.
不是說那天才做.
所以很多人已經做了部署.
結果.

$^{601}$駐澳門的宣教士.
大部份都被調走.
要麼調來香港.
要麼調去五口通商的口岸.
所以澳門的教會的發展.
是比較緩慢.
我講這些有甚麼特別呢.
其實.
我們其實在今天.
是承繼前人的.
工業.
我們有沒有像他們那麼放膽.
去傳道 我有些很掙扎.
因為.
為甚麼呢.
當我接觸這些儒式道的人的時候.
我就抱著態度.
不要去得那麼盡.
因為我不懂.
但神給我一個機會.
因為有些牧者.
很喜歡教授中國文化.
我就陪他們.
五次我去.
答應他一次.
你想想我是哪一類人.
我是mung 車邊那種人.
但我發覺一件事.
原來在我身邊很多人.
覺得不是必須要信耶穌.
是絕對有理由.
你同不同意我這樣講嗎.
他不會覺得.
我不信耶穌我犯了甚麼大錯.
倒過來我會問.
我為甚麼可以說他不信耶穌.
不對呢 你明不明.
你一定要解的.
你不能說了就算.
所以對我來說.

$^{641}$我覺得在現在這一刻.
我們真的需要.
想想.
從甚麼去和一個.
有他的人生宗旨.
他又不覺得自己.
壞到哪裡去的人.
你要信耶穌.
你夠膽跟他說嗎.
還是你叫他.
你有空就信耶穌吧.
你明不明我的意思.
如果你這樣說.
我是那個人.
我坐在你那裡都會驚訝.
你先搞定你自己.
你那麼沒有層次的.
我再講第三件事.
沒有人能禁止.
最近因為有些同道.
其他宗教人士.
和我交流.
可能你們聽過.
我沒聽過也不一定.
他們開始講不同的文化板塊.
其實這些我以前也聽過.
西方文明.
就是海洋文明.
中國文明就是大陸文明.
海洋文明和大自然搏鬥.
所以需要神明.
我們就不是.
在大自然和四季的調和.
我們是很和諧的.
這些說法是幾十年前說的.
這些文化發展令人操纏到不得了.
但近年有一個新的說法.
大陸的板塊和海洋板塊.
互相改變了.
以前是海洋板塊.

$^{681}$厲害一點.
所以大英帝國很多船艦.
但現在不行了.
現在是陸地板塊.
我覺得很有趣.
但我從來沒見過.
有專家有專書寫這些.
隨便你說.
我就想了一件事.
現在不是講陸地和海洋.
是講太空.
美國GPS.
我們中國有北斗.
你們兩個不同系統.
在那裡可以.
那你為什麼大陸板塊.
海洋板塊.
海洋板塊不是比你還厲害.
我想問的問題.
有這些人.
他這樣說.
對我們基督教有什麼影響.
其實言下之意.
有一個叫Samuel Huntington的人.
就講文明衝突的理論.
基督教和伊斯蘭教的衝突.
然後就西方文明的沒落.
言下之意.
西方文明現在倒塌了.
基督教從西方文明來.
都倒塌了.
你明不明.
如果你要在.
全基督教的時候.
你要改頭換面化妝才行.
你明不明我的意思.
化妝.
他來教你化妝.
教你化妝教你改.
然後你看印度那個.

$^{721}$小成佛教來到中國都變成大成.
叫教好改一下.
你明不明我說什麼.
我心裡面很掙扎.
如果這樣講.
究竟我聽上帝的話.
還是聽你的話.
我聽聖經的教導.
聽神學.
我們的信仰傳統還是跟你呢.
我們承認使徒信經.
是幾千年裡面的傳統.
我們基督教有傳統.
但今天.
在我們周圍就有很多這些聲音.
有很多人.
他當這些是對的.
他就已經衝了過去.
但周圍我們.
今天一個認真的基督徒.
你想一下我們承繼著什麼.
前人做過什麼.
我們今天怎麼做.
早幾天我們神學院的.
同事就去了.
西貢崇真會.
那裡一個.
退休會.
他們找我.
搞一個遊戲.
我不懂搞遊戲.
但我都不懂.
我不懂搞遊戲.
但我就找一下崇真會的歷史.
不知道大家有沒有人.
懂得崇真會的歷史.
可能都不知道.
在1847年.
巴西警會.
即是瑞士的巴西警會.

$^{761}$就派遣了韓山明牧師.
和黎力基牧師來中國.
如果你是.
潮州人的時候.
你在潮汕教會都應該聽過.
黎力基牧師.
因為他們最先.
就去潮汕.
那裡傳道.
結果.
潮汕的知縣趕走了他們.
趕走了他們兩位之後.
他們就.
周平他們.
後來叫崇真會.
巴西警會就定義在.
客家人中間傳福音.
所以如果你是客家人.
你沒理由.
如果你是客家教會.
你不會聽韓山明和黎力基的.
黎力基.
就是德國人.
韓山明是瑞典人.
有很多小故事.
他們最初.
就跟禮賢會的宣教士.
國事牧師學廣府話.
廣東話等等.
但你想一下.
當他們要走去客家人中間.
當時.
學客家話.
是很高難度.
因為沒有錄音帶.
你聽的.
你可以說.
這班人你怎樣想像.
禮賢會的宣教士.
就去東莞.

$^{801}$去九江去新會傳福音.
這些就成為.
其中一些人物.
還有一些去到很遠很遠的地方.
去傳福音.
而這班是外國人.
你想一下.
這班外國人可以說一些地方的方言.
是我們都未必可以.
你想一下.
還有你想一下.
十九世紀的時候.
中國的生活條件.
和今天的生活條件.
你來做個比較.
你就可以想像.
他們死不少人.
搭船來香港可以在船上死.
可以被舌頭壓死.
可以生了子女.
很快腰瘠死.
在香港澳門的墓地.
有不少他們的墓碑.
但是我們有時.
就好像覺得.
這些都是屬於過去.
和我們好像.
拉不上什麼關係.
但我很想說.
無人能禁止.
黎歷基穆斯來中國.
是23歲.
很多人都是.
20多歲.
來去做一件.
我們叫轟轟烈烈的事.
但是我們.
23歲的時候在做什麼呢.
你們明白我的意思嗎.
我們23歲的時候在想什麼呢.

$^{841}$我想在這裡.
就希望藉著這兩節經文.
彼此互勉.
就是說神已經告訴我們.
祂要將福音.
向全世界去宣講的.
並且《使徒行傳》第28章第31節說.
是沒有人可以禁止的.
是沒有人可以禁止你的.
問題是你做不做.
沒有人能夠禁止下一個說法是什麼呢.
我很想挑戰大家.
就是說.
你有沒有想過你是向外國人傳福音呢.
有沒有這樣的大志呢.
可能你說我沒有.
不要緊你和你的下一代說一下.
從理論上是否可能呢.
絕對有可能.
你想想你的親戚.
很多人住在哪裡.
在加拿大,美國,英國,德國,荷蘭.
哪裡都有的.
他們的下一代都已經說過那裡的話了.
所以可能與不可能.
是我們自己的思維.
但是聖經就已經這樣命令我們了.
我們要放膽傳講神的國.
就教導主耶穌基督的事.
沒有人禁止.
但是你要知道你和我在世的日子.
其實都不會很多的.
對不對.
都不會很多.
既然不是很多.
我剛才都說了.
我以前就遷就著.
其他宗教的人.
不要得罪他,不要踩傷他.
也不要冒犯他.

$^{881}$我也很會做人.
但是到了這個階段的時候.
我又會覺得.
可能我要想一想.
我不要得罪我的神.
既然你是相信的時候.
為什麼你不敢說.
我的事比你的事更好呢.
是不是我說我的事比你的事更好呢.
就是我霸道呢.
這個問題就是.
如果我不是發現過這是真理.
我去相信.
那我相信什麼呢.
你明白我說什麼嗎.
如果我不是真的像保羅那樣.
他明白神是真實的.
他就不會奮身去傳道.
所以我想這個都值得我們反省.
你越是信得真.
你就會去得盡.
如果你不是很相信.
我就輕輕和你過兩招就算了.
我又不傷你.
你又不傷我自己.
我講了我的見證.
我向上帝有個交代.
但是慢慢你會想一想.
上帝會怎樣和你說呢.
我們一起同心祈禱.
我們真的天賦上帝.
我們知道.
我們在你的面前.
只可以傳敬畏的心.
我們都是很沒用的人.
很普通的人.
我們要學習的就是.
你藉著這麼多人在世界各地.
不斷地做你自己的工作.
剎下種子.

$^{921}$但是我們會不會在今天.
我們比較安逸.
而輕忽了你對我們的呼召.
你對我們的託付.
主 既然我們是尊敬你.
願意服侍你的時候.
究竟我們會留下些什麼.
給以後的教會.
以後的信徒.
究竟我們是不是為他們建立一個平台.
讓他們繼續起飛.
還是我們只不過是維持局面.
還是我們只是告訴自己.
我都有做到神要我們做的事.
我們的主就懇求你自己.
在我們的內心.
你向我們的說話.
是你感動我們.
不是我們自己要成就什麼大事.
求你自己的愛.
你的話語.
你的靈在我們心裡工作.
讓我們明白聖經的訓示.
讓我們真的忠於我們的信仰.
忠於我們的上帝.
禱告奉耶穌基督的名.
\newpage



\section{路加福音 10:1-12}
\label{sec:PIfMPVh5oUA}
\textbf{2023-07-16 早堂《耶穌訓練門徒》路加福音 10\_1-12 講員:曾裔貴牧師}
\newline
\newline
連結: \href{https://youtube.com/watch?v=PIfMPVh5oUA}{\texttt{https://youtube.com/watch?v=PIfMPVh5oUA}} ~~~~ 語音日期: 2023-07-23
\newline
\newline
\hyperref[sec:vtGjR9atL3U]{\small{< < < PREV SERMON < < <}}
~
\hyperref[sec:index_chronic]{\small{[返順時目]}}
~
\hyperref[sec:index_scriptual]{\small{[返順卷目]}}
~
\hyperref[sec:zajnPwak_dU]{\small{> > > NEXT SERMON > > >}}
\newline
\newline
路加福音 10:1-12
\newline
\begin{longtable}{cl}
\hline
\hline
章節 & 經文 (和合本修訂版)\\
\hline
10:1 & \begin{tabularx}{0.7\textwidth}{X} 這些事以後，主另外指定七十二個人，差遣他們兩個兩個地在他前面，往自己所要到的各城各地去。 \end{tabularx} \\ \\ \relax
10:2 & \begin{tabularx}{0.7\textwidth}{X} 他對他們說：「要收的莊稼多，做工的人少。所以，你們要求莊稼的主差遣做工的人出去收他的莊稼。 \end{tabularx} \\ \\ \relax
10:3 & \begin{tabularx}{0.7\textwidth}{X} 你們去吧！看！我差你們出去，如同羔羊進入狼群。 \end{tabularx} \\ \\ \relax
10:4 & \begin{tabularx}{0.7\textwidth}{X} 不要帶錢囊，不要帶行囊，不要帶鞋子；在路上也不要向人問安。 \end{tabularx} \\ \\ \relax
10:5 & \begin{tabularx}{0.7\textwidth}{X} 無論進哪一家，先要說：『願這一家平安。』 \end{tabularx} \\ \\ \relax
10:6 & \begin{tabularx}{0.7\textwidth}{X} 那裡若有當得平安的人，你們所求的平安就必臨到那家，不然，將歸還你們。 \end{tabularx} \\ \\ \relax
10:7 & \begin{tabularx}{0.7\textwidth}{X} 你們要住在那家，吃喝他們所供給的，因為工人得工錢是應當的；不要從這家搬到那家。 \end{tabularx} \\ \\ \relax
10:8 & \begin{tabularx}{0.7\textwidth}{X} 無論進哪一城，人若接待你們，給你們擺上甚麼食物，你們就吃甚麼。 \end{tabularx} \\ \\ \relax
10:9 & \begin{tabularx}{0.7\textwidth}{X} 要醫治那城裡的病人，對他們說：『神的國臨近你們了。』 \end{tabularx} \\ \\ \relax
10:10 & \begin{tabularx}{0.7\textwidth}{X} 無論進哪一城，人若不接待你們，你們就到大街上去，說： \end{tabularx} \\ \\ \relax
10:11 & \begin{tabularx}{0.7\textwidth}{X} 『就是你們城裡的塵土黏在我們的腳上，我們也當著你們擦去。但是，你們該知道神的國臨近了。』 \end{tabularx} \\ \\ \relax
10:12 & \begin{tabularx}{0.7\textwidth}{X} 我告訴你們，在那日子，所多瑪所受的，比那城還容易受呢！」 \end{tabularx} \\ \\
[1ex]
\hline
\hline
\end{longtable}
$^{1}$我們一起同心禱告.
感恩天父讓我們能聚集 敬拜你.
我們可以聆聽聖經的真理.
還有我們剛才用詩歌來頌揚.
我們這位三一的上主.
很盼望整個會場充滿一份平安.
充滿一份屬天的福樂.
我們帶著一份開放的心靈.
來等候我們的上主,教導,差遣.
我們這樣祈禱.
奉耶穌基督的名.
阿們.
不知道在大家的經驗裡.
有什麼選擇是你覺得很寶貴.
到今天還記得的呢.
我在五年級的時候.
五年級小時候.
幾十年前.
很開心.
有運動比賽.
原來是選拔代表學校.
當時我是讀巴富街小學.
代表學校參加香港小學的校運會.
我又很矮.
一個選拔.
拼了命跑一個籃球場的長度.
最終得到第六名.
選拔五個.
後來第五名的那個.
不知道什麼原因.
被淘汰了.
我就自動補上第五名.
可以代表學校去比賽.
很開心.
整晚睡不著.
誰知道那個是四人的接力賽.
只不過是後備.
當中的後備.
最終都沒有參賽.
我已經感覺很開心.

$^{41}$盧家福音第十章這一段.
是四福音裡唯一一段.
無端端選照七十個評訓圖.
選照他們去傳道.
和十二使徒的吩咐是完全一樣.
你就會覺得這段聖經是很特別的.
第九章作者已經交代了一個總結句.
就是說耶穌基督是面向耶路撒冷.
他一直向耶路撒冷.
就是說其實九章的五十一節.
耶穌被接上升的日子將到.
他就定義向耶路撒冷去.
接著打發使者在他的前頭走.
接著那些使者就先去到薩瑪利亞的村莊.
那些使者極有可能是十二使徒.
但現在來到第十章的時候.
耶穌基督是呼召七十個人.
但沒有說他們是什麼背景.
他們是什麼職業.
也沒有提到他們有什麼訓練.
就是這樣吩咐.
似乎他們是已經跟隨耶穌的普通信徒.
但在這個時候.
耶穌給他們一個很神聖的使命.
他們要為耶穌.
先在耶穌的前面.
好像為耶穌開路.
預備人的心靈.
去到不同的家庭.
如果我們一直看下去的時候.
表面上我們覺得有些.
為什麼要說這些.
有四個不需要,不要做.
還有就是有一樣東西必須要做.
就是第五節.
無論你去到哪一個家庭.
首先要向那個家庭說.
願這個家庭得到平安.
就是他們這次兩個兩個出去傳道的目的.
就是要將贖天的平安帶到那個家庭.

$^{81}$這裡有些技術性.
告訴他們不要想那些日常需要.
鞋,食物,乾糧,衣服.
不需要去想.
但又說去到那個家庭.
人家給你吃什麼.
你就吃什麼.
你要在那個地方住.
人家怎樣接待你.
你就要這樣去歡迎的接受那個款待.
即是說不是街頭的派單張報道.
是深度的家訪.
在家庭裡跟那些家庭.
有機會深入去探討.
什麼是贖天的平安.
但是有一個很大的訊息.
很重要的訊息.
就是一定要宣告的.
就是神的國度臨近你們了.
因為去到後面.
耶穌說的.
第十二節打後.
應該是十一節的下半節.
就說了.
如果人不接待你們.
你們就去到街上.
就說了.
就是你們城裡的塵土.
黏在我們的腳上.
我們幫著你們拆去.
好像有一種.
跟這些否定,不信,拒絕福音的人.
好像有一個劃清界線.
即是說你們不接受那個平安.
那個平安就會給了我們傳道的人.
然後留意一下耶穌說的這句話.
雖然是這樣.
即是說你們面對那個拒絕.
雖然是這樣.
你們該知道神的國臨近了.

$^{121}$叫兩個兩個報道團的人.
要記住這件事.
不要記住那些人的拒絕,反對.
或者一些白眼.
或者一些不好的,不友善的對待.
要記住的是你們之前說了的訊息.
他們拒絕.
但是最終耶穌叫他們要留意.
要深化的.
就是雖然是這樣.
雖然即是說那些人的拒絕.
雖然如此.
你們該知道神的國臨近了.
即是怎樣.
即是這些短期的傳道者.
他們更深的經歷,感受.
和唯一重要的.
就是他們傳揚神的國.
他們自己是在神的國裡的國民.
和他們感受到神的國度.
原來是沒錯.
受拒絕.
但是他們舞弊那個國的君主.
天父,神,父神.
耶穌所拒絕.
他們要更加感恩.
自己是在那個國度裡面.
去望外面的世界.
這個就很值得我們大家一起去想一想.
這段經文的整個脈絡.
還有我們很盼望弟兄姊妹.
你是國度裡面的參與者.
你是在神國度裡面服侍的人.
首先第一件事就是.
耶穌呼召了70個人之後.
幫他們分了35隊.
然後第一件事叫他們做什麼呢.
就是叫他們從神的眼光.
主耶穌的眼光去看周圍的世界.
原來整個世界是一個和場.

$^{161}$第二節就對他們說.
要收的莊家多.
作供的人少.
70個不少.
十二使徒已經有了.
現在再揀選70個.
35隊去不同的村莊.
不少了.
是耶穌設立的.
但是耶穌的眼裡面.
天父的莊家很多.
作供的人很少.
所以就叫這70個人.
求莊家的主.
打發工人出去收神的莊家.
信了主的人.
在耶穌基督裡面參與服侍的人.
第一件他發自內心會做的.
就是認同耶穌的眼光.
原來真的很多人還沒有在救恩裡面.
很多人需要救恩.
很多人是滅亡.
信了主的人才發現.
自己要真心向神求告.
求告多些的工人.
多些人看到世界的方良.
世界是每天都在運作.
高度的商業化.
科技的發展.
一日千里.
但是人心在神眼中.
是待修割的莊家.
工人要看到這個.
這樣的熟天望下來的眼光.
不是否定地上的人.
全部都是要立刻滅亡.
而是待修割的莊家.
由拜子變成墨子.
就是要有工人去禱告.
這些已經是被揀選的.

$^{201}$準備出去傳道的停信徒.
但是耶穌叫他們祈禱.
所以他們一路祈禱的時候.
他們必然會帶著耶穌的差遣.
感動就出去了.
然後耶穌第三節告訴他們.
其實不要說這句就好了.
還沒說就告訴他們.
你們是羊羔.
你們去的不是一個和場那麼簡單.
是進入狼群.
為什麼耶穌說這麼恐怖的話.
35隊進入狼群的羊羔.
你想想到底是一個怎樣的教導呢.
就是你一去就死無生.
羊是最沒有抵抗力量的動物.
但是耶穌是很清楚的說了.
告訴他們.
似乎都不希望他們靠自己血氣的力量.
去做那些工作.
似乎都不是靠他們去爭辯.
去用武力帶著一些防身的棒.
類似諸如此類.
似乎耶穌都不是要他們這樣做.
要他們經驗到他們的福音.
和外面的拒絕有一個很大很大的距離.
我預備這段經文的時候.
突然間無意中找回一些檔案.
在2001年的5月27號.
在菲律賓一個阿布沙耶夫的恐怖組織.
菲律賓的.
就綁架了很多外國人.
當然籌碼換取贖金.
去和政府軍對抗.
輾轉菲律賓的政府軍和美軍展開營救.
最終放了很多大部分的外國人質.
放剩三個人.
其中一對夫婦是宣教士.
一直在菲律賓做傳道.
幫宣教士的組織運送病人.

$^{241}$在叢林區駕飛機運送.
這個恐怖組織.
無論如何都不放這三個人.
其中一對夫婦就是宣教士.
最終四處在叢林逃亡.
因為政府軍不斷剿滅他們.
用了一年時間.
最終政府軍找到.
展開槍戰.
這位丈夫宣教士.
在槍戰當中陣亡.
死了.
都不知道是政府軍開槍.
還是恐怖分子開槍.
這個遺傷.
腳受槍傷.
輾轉和救.
一年後.
他回到美國.
將他內心的心路歷程說出來.
他那一年.
周圍跟隨恐怖分子.
在叢林不斷逃走.
他內心痛恨這些人.
很痛恨.
很差的待遇.
和對待那些人質.
很嚴重.
但是包括他的丈夫離世之後.
他回到美國.
他自己一路再沉澱他的生命.
他發現.
他居然用詩篇的23篇第5節.
來總結他內心的感受.
在我的敵人面前.
你為我擺設筵席.
你用油告了我的頭.
使我的福杯滿日.
他感受到.
雖然他內心.

$^{281}$在那一年裡面有極大的憤恨.
對這些恐怖分子的仇恨.
他的丈夫也是這樣死了.
但是最終.
他說他自己因為神的大愛.
他選擇接受和寬恕這些害他丈夫.
和令他這一年裡面.
受到很多很多不好的待遇的這些人.
他選擇接受和明白.
上主總是有他的恩典和旨意.
已經過了一段很長時間.
如果你有機會可以上他的網頁.
你會看到他的見證.
他現在都在教會裡面很忠心的服侍.
他的子女都大了.
這些就是我們告訴我們.
耶穌沒有欺騙我們.
我們作為一個傳道者.
將福音傳講的時候.
總是會有人對天國有敵黨.
我們怎樣去將這些表面的敵黨.
來消化了它.
再輸送出來的在我們的內心.
仍然是一份真愛對別人的接納呢.
你看主耶穌在這裡講的時候.
第六節.
如果有當得平安的人.
你們所求的平安.
就會臨到那個家庭.
即是說原來是有成果的.
當門徒兩個兩個去到這些村莊.
去到那些家庭.
那些家庭是欣然的接受.
聽這兩個人所講的神國度降臨.
甘信不疑接受神的平安.
就會透過這兩個人臨到那個家庭.
那個家庭往後的生命改變.
就有神會負責.
有神會在那個家庭裡面.
來搞動那個家庭.

$^{321}$他們會經歷屬天的平安.
是因為這兩個人聽到耶穌的吩咐.
他們要出去進入這些家庭裡面.
那個家庭若有當得平安.
這個是你和我都不知道的.
我們不去講.
我們不去傳.
我們不知道到底是拒絕平安的家庭.
還是欣然的接受的家庭.
我們的童宮在幼稚園裡面.
我們開始開幸福小組.
有些小朋友的家長.
真的來聽.
真的接受.
真的開始探索福音的訊息.
很快就會進入那個真理.
生命的改變.
是由聽明白屬天的神國度降臨.
而來準備迎接平安.
所以耶穌在這裡的教導是很實在的.
沒錯.
陽高代表著我們的內心.
不是靠自己的力量.
是靠神的恩典.
明白處境.
明白世情.
但仍然傳一個叫人得到平安的福音.
弟兄姊妹.
我們教會需要動員信徒.
你自己是否已經在神國度裡的事奉者.
你敢不敢受到耶穌向你的選照.
你能夠被選上是極大的福樂.
當你被選上之後.
其實你就應該.
你應該有一個屬天的眼光.
再看這個世界.
不只是商業化高度科技化的一個世界.
是一個人的心靈是虛空.
是不知道前途是怎樣的一個世代.
有錢的人.

$^{361}$沒錢的人.
同樣面對的可能就是困惑.
可能就是不知道人生是怎樣的意義.
是需要我們幫他們去解釋的.
所以耶穌叫這班人首先祈禱.
然後就要出去.
當去到的時候.
他們就會感受到那個平安的價值.
接著就說.
工人得工價是應當的.
所以就不要見義思遷.
不要見好的待遇.
就留久一點.
留到不捨得走.
成為了那個人的一個操縱者.
不要.
讓神的愛帶領著.
什麼都不需要介意.
物質永遠都是物質.
重要的是大家去經驗到神的平安.
這個是很重要的一個發言.
接著就是有一些的不接待.
耶穌也說了.
第十節.
無論進那一城.
剛剛和第八節是相對的.
無論進那一城.
第八節.
人接待你.
第十節.
無論進那一城.
人不接待你.
所以是相對的.
就是說那個人應該不要理會家庭是怎樣.
曾經和他有過節的.
說得差一點.
令到他不接待我.
我出去就說不關我的事.
不要這樣.
不要理會那個人是怎樣.

$^{401}$說的是一個原原本本的福音.
不帶著一個歧視.
不帶著個人的成見.
這個當然很難.
每個人都有自己的喜惡.
但是傳道者應該要原原本本說明那個福音.
就算那個人是你不喜歡.
你仍然用愛去說平安的福音.
所以那個不接待.
應該盡量降低到對方對福音.
不是對你的態度.
所以第十節.
無論進那一城.
人如果不接待你們.
你們就到街上去.
耶穌有教的.
所以原來福音是無價之寶.
但是福音不是賤價大平賣.
你不喜歡我求你買.
不是.
盡量說明白之後.
當人真的這樣決心.
硬心的拒絕.
我們要有一個這樣的態度.
就算你們城裡的塵土.
黏在我們的腳上.
我們也當著你們的面拆去.
做給他們看.
你們這樣拒絕神.
其實有後果的.
不是我們要咒你.
是神的審判將來會淋到你們身上.
這裡就告訴我們.
我們都需要有這種的緊迫感.
我們不是說一下不聽.
算了有他有他.
沒所謂.
而是我們有一份危機感.
那個危機感就是.
對方是滅亡.

$^{441}$我們怎樣能夠從一個滅亡的人搶救.
回到神的國度裡面.
很盼望我們的弟兄姊妹.
我們走出街.
或者去到工作的場所.
我們見到的同事朋友.
我們有一份心靈的焦急.
很希望他能夠信主.
因為神的國臨近.
你留意第十二節.
真是講審判.
我告訴你們.
當審判的日子.
所多瑪.
這是創世記十九章.
天火焚城的故事.
那個記載.
所多瑪所受的.
比你們去的那個城.
還容易受.
不過審判更加嚴重.
你就可以想像到.
原來這兩個出去的.
三十五隊的傳道者.
他們是將一個搶救人靈魂的.
這麼重要的職事.
來到出去做.
所以傳福音不是開玩笑.
是很嚴肅的.
他告訴人.
我們在講一些.
沒有關乎到你的生命永恆的事.
這是很寶貴.
很重要.
很快我們跳到第十七節.
就是整個傳道.
不知道用了多長時間.
一個星期.
三天.
聖經沒有記載.

$^{481}$但是給我們看到的.
就是百分之一百的回報.
三十五隊.
個個都是興奮地回來.
十七節.
七十個人歡歡喜喜地回來.
主啊.
因為你的名.
就是鬼都服了我們.
鬼都服.
人就更加服.
就是說那些家庭.
如果有超自然的問題.
有靈界的附身.
和醫病.
都能夠做到.
耶穌給他們短期的能力.
去做到這些超自然的工作.
分外興奮.
人接受平安.
他們已經有一份滿足感.
他們已經能夠有一個好的被接待.
可以想像.
他們在那個家庭裡面的居住.
談論道理.
一定是很暢快.
大家好開心.
很坦誠.
再加上有靈界的經歷.
他們的手.
按手.
那些鬼都走.
你想像得到.
這三十五隊的短孫隊.
非常振奮.
十八節.
耶穌對他們說.
我曾經看見撒旦.
從天上墜落.
像閃電一樣.

$^{521}$我已經給你們權柄.
即是說差遣的時候.
給了那個權.
屬天的權柄.
這三十五隊的短孫隊.
你們的能力就是可以踐踏蛇.
以及蠍子.
代表著靈界邪惡的一面.
有勝過仇敵.
一切的能力.
斷沒有什麼能夠害你們.
原來耶穌告訴他們.
你們經驗到這些.
我早就給你們權柄.
所以你們現在能夠經驗到.
一點都不意外.
耶穌不會覺得很意外.
因為耶穌告訴他們.
我就是給你們一個經驗權柄能力.
的一個時刻.
所以你們就有這種收穫.
不過留意第二十節.
耶穌基督很語重心長地.
加個言而.
不可以因鬼服了你們就歡喜.
要因你們的名字.
記錄在天上歡喜.
很特別.
35對不約而同的100\%的.
表示這個短期的福音工作.
大獲全勝.
很多的家庭.
都是因為我們說天國的降臨.
他們就深信不疑.
接受,接待我們.
而且我們還替主你去做一些.
超自然的醫治和拯救的工作.
但是耶穌居然在這裡.
跟這35對的人說.
不要只是有一個表面的興奮歡喜.

$^{561}$要因為不改變的,永恆的.
這個歡喜要經常地思考.
就是你的名字記錄在天上歡喜.
就是說你們是屬神的國度的子民.
你們是被神接納.
在神的眼中.
現在甚至乎是能夠去事奉的一些領袖.
這樣歡喜.
這個是很快人心醒.
事奉的果孝和事奉的人.
要分清楚.
要能夠常常想到.
就算沒有果孝的事奉的人.
神都是愛他.
他的名字都是記錄在天上.
這樣我們人的心態就能夠越來越憐靜.
越來越能夠明白.
不是為果孝去服事.
我們不是服事那個果孝.
是服事我們這一位愛我們.
拯救我們的神.
相信如果你是有事奉經驗的弟兄姊妹.
你會明白兩者是有很大的分別.
一個事奉主的人.
就算他事奉沒有果孝.
他都仍然會感恩讚美神.
因為他能夠有份於事奉.
所以在21節耶穌就歡樂了.
為了這一班人.
他們的工作歡樂.
但為了一些不信的人.
他就會感到痛心.
求主幫助我們能夠在這段經文裡.
主耶穌教我們.
要在神的國裡面看外面的世界.
你會很不同.
我想起有一個藝人.
鄭敬基.
他自己講見證.
他的事業最成功的時候.

$^{601}$他的太太前妻和他離婚.
他信主的.
因為三番四次背叛太太.
太太終於忍無可忍.
和他決絕離婚.
當時是他的事業.
演藝的事業最成功的時候.
其實他一段時間不成功.
很憤世嫉俗.
很多很多的埋怨.
自然就是抵受不住很多的誘惑.
就是剛才所講的對太太的背叛.
這個突如其來的打擊.
他怎麼辦呢.
他一路一路去想.
留下香港的事業.
還是立刻飛去加拿大.
來挽回這段婚姻呢.
後來他終於選擇去加拿大.
但是太太心意已決.
不接受這個復合.
他有兩年的時間.
完全沉淪在他自己的痛苦裡面.
完全回不出來.
人人都見到那個情況那個時間.
他是一個信主的信徒.
兩年後.
他的太太見到他.
這種的警方.
主動的來幫他.
和他說.
我已經訓練到他有一對子女.
子女不會仇恨你.
我這個太太.
都已經寬恕你.
接納你.
你要接納自己.
再加多一句.
我的主.
你的主.

$^{641}$也都已經接納你.
對他來說.
是一個很大的.
一個的.
好像亮光一樣.
讓鄭敬基完全的被這番說話.
來改變.
他重新的振作起來.
開始回和前妻.
這個家庭.
保持的關係.
到今天還是一個很好的關係.
他的子女對爸爸的接納和愛.
尤其是他感受到最大的接納.
就是太太帶著兩個子女.
都寬恕他這個錯誤.
和由太太的口裡面.
讓他有好心的那份感受.
主耶穌已經寬恕了你.
為什麼你不肯寬恕自己呢.
親愛的弟兄姊妹.
最寶貴的人生.
是我們深深感受到.
救恩將我們的生命拯救.
以致我們所做的一切.
是為主耶穌甘心樂意去寬恕.
去傳遞神的國度.
不是一種施捨給你.
你不要就隨你吧.
我已經得救了.
不是這樣的.
我們是發現我們被愛.
十價的大愛拯救.
所以我們是很樂意.
和人分享神的國度.
即便接受的人.
可能起初是拒絕.
對我們有很多的懷疑.
甚至乎用很多不好聽的話去否定.
但我們仍然帶著的.

$^{681}$是一份愛的包容,寬恕.
因為主耶穌寬恕我們.
因為神的國度.
我們的名字已經記錄在天上.
求主幫助我們能夠透過這段經文.
讓我們能夠.
每一位弟兄姊妹.
祖堂的.
我們都能夠一起.
接受耶穌基督的訓練.
成為一個傳遞神國度的門徒.
我們一起祈禱.
多謝主.
很盼望我們能夠在神豐厚的恩典裡面.
再去看這個和長世界的需要.
沒錯 世界需要改善失業率.
改善氣候.
和國與國之間的戰爭.
到今天都還沒有停止.
肯定不會一下子就停止.
在蘇丹.
在很多我們沒有聽過的部落.
可能每天都有屠殺.
這一切在上主的眼裡面.
你都是清清楚楚.
求主打發工人.
讓我們每一個裝備好自己.
去收裝架的裝備.
新心靈都是有這樣的準備.
被上主來差遣.
我們這樣禱告奉耶穌基督的名.
阿們.
\newpage



\section{}
\label{sec:zajnPwak_dU}
\textbf{2023-07-13 周四晚查經班《被擄歸回歷史研讀(20)》 講員:曾裔貴牧師}
\newline
\newline
連結: \href{https://youtube.com/watch?v=zajnPwak-dU}{\texttt{https://youtube.com/watch?v=zajnPwak-dU}} ~~~~ 語音日期: 2023-07-25
\newline
\newline
\hyperref[sec:PIfMPVh5oUA]{\small{< < < PREV SERMON < < <}}
~
\hyperref[sec:index_chronic]{\small{[返順時目]}}
~
\hyperref[sec:index_scriptual]{\small{[返順卷目]}}
~
\hyperref[sec:epe4ZHzeWUY]{\small{> > > NEXT SERMON > > >}}
\newline
\newline
$^{1}$我們一起有簡單的祈禱.
多謝天父再一次我們在晚上.
有一段時間來讀聖經.
都盼望我們能夠透過經文的啟示.
和所描述的當時歷史背景.
讓我們能夠掌握.
也讓我們能夠從其中一些重要的教訓.
得到一些教導.
我們在我們的教會生活.
在現在這個世代.
我們應該怎樣來跟隨主.
都是需要有神的聖靈的指引和感動的.
特別在教會的事奉.
更加需要神的恩典.
我們將這一段時間仰望,交託.
也都盼望錄音的訊息.
可以成為很多人的提醒.
彼此的勉勵,互相的扶持.
甚至乎現在跨越了那個時間.
在網路上可以更加多人有機會去聽.
盼望去祝福到很多華人的信徒.
能夠聽懂廣東話的.
都能夠得到神話語的教導.
我們這樣仰望,禱告.
奉耶穌基督的名.
阿們.
第七章的尼希米記.
就是城牆做好之後.
尼希米就發現了一個問題.
就是城牆做好之後.
城門也都安好門山.
但是耶路撒冷城是廣大的.
但是居住的人口稀少.
還有都沒有一些房屋來建造.
就是說一個空曠的聖殿範圍.
急切需要將人口遷徙.
進入聖殿聖城裡面.
所以這是一個很大的工程.
哪些人肯搬進來聖殿呢?.
會不會當時經過了百多年之後.

$^{41}$有些環境.
有些不同的滄海桑田.
一些的貿易.
有些的路線.
特別是你可以想像到.
有些大街有些新的發展.
可能對它來說.
是會有一個很大的幫助.
但是如果有些道路改了.
可能過去很繁榮的.
慢慢又會因為不方便.
人流又會減少.
這些我們都會見得到.
你可以想像到在行鄉那邊的街.
和行鄉對面的街都差很遠.
因為很多時候.
人們回來在元朗下車.
就是行鄉第一個站.
就會很多人下車.
那個舖位是特別矜貴.
特別多人流.
但相反對面街.
人們就是早上上班才會趕到那裡乘車.
通常很少人流.
不知道你有沒有發現這個問題.
所以對面街的商舖.
很多都是二三線.
或者不是很重要的大企業.
你不會見到在以前嘉麗華戲院.
嘉麗華酒樓附近的那些.
是那些金舖.
通常是賣那些生貨舖.
賣那些不是那麼貴重的東西.
你可以想像到.
會不會在今年雷月.
他們重建回去的時間.
過了一段很長的時間.
其實他們已經有很多的改變.
所以很多人都住在外圍.
特別是他們找回自己的中派.

$^{81}$自己那個支派的地方的時候.
他們就在那裡居住.
所以你會看到.
要他們用一個行政的指令.
來逼令他們回去耶路撒冷聖城裡面居住.
是一個相當難度的課題.
我們見到第八,第九章.
去到第十章.
就是這個七月初一的一個嚴肅會.
這個第八章.
一直去到十章.
就是有兩次的.
可以說是一個奮興會.
就帶來百姓有一個很大的決心.
發現自己是民族神選民的身份.
即是提升他們有自己的身份價值.
提升他們在神的旨意裡面.
他們被擄歸回.
或者在被擄的行列的下一代.
是生在一個重歸,重建的聖城的子民.
不是生在波斯帝國.
而是第一批,第二批上來的猶太人.
他們安頓了.
也都有聖殿了.
但是一直以來的國家.
尤其是以耶路撒冷做單位.
就是因為城門,城牆都是破落.
所以外邦人就很容易來買賣.
也就是說要將猶太教的嚴謹去溝淡.
令到你們不能夠有一個很純正的心靈去敬拜神.
更加就是會荒廢最直接的.
就是對神敬拜的事奉.
所以第八,第九章是一個非常大的奮興.
帶來人心的柔軟.
人的心靈對於自己的屬神的身份.
有了這個意識,強烈的意識之後.
很多事情就很容易.
因為是由自己內心的感動而發出的.
第十章就提到他們去簽名.
來去願意遵守摩西的律法書.

$^{121}$來去在一些重大的事情上.
來去有那個決心跟隨神的旨意.
例如十章的三十一節.
「這地的居民如果守安息日或什麼聖日.
帶了貨物或糧食來賣給我們.
我們就必定不賣.
簽約在神的面前起誓要守安息日」.
不是一種死守.
而是發現安息日是帶著一個敬虔的心靈去敬拜神.
能夠有這樣的一個心靈的movement.
心靈的一個感動.
就帶來一個很大的一個祝福.
所以你會看到十一章不再難.
就是這個尼欽米做了一個舉措.
就是可能跟百姓跟領袖有了一個協議.
表面上這個都是一個行政的指令.
來到去怎樣能夠讓更加多的人搬進耶路撒冷呢.
首先就是首當其衝的就是.
祭司利未人這些專職的.
本身就是要在聖殿裡面服侍的.
他們是二話不說.
根本就是他們應該要住在耶路撒冷.
來到去行上的功績.
但是都需要有其他十個支派的支持幫手.
這些人他們不是英雄.
他們是被神感動的人.
他們來到去很願意.
很樂意被在尼欽米制定的一個方法.
怎樣可以公平地能夠動員到.
其他十個支派的人.
以色列人可以住在耶路撒冷呢.
用了一個辦法.
就是將十個支派逢十抽一.
逢十個人抽一個男丁進去住.
當然涉及的是他們的家庭.
是他們的兒女.
就整個家庭這樣move去耶路撒冷.
好像現在搞豬抽居屋那樣.
他們就被很幸運地選中.
背後有一班人.

$^{161}$他們已經認同自己是神選民的身份.
很樂意照著尼欽米的動員教導.
大家一起.
其實最重要的是恢復聖殿.
行上的敬拜.
當有行上的敬拜.
就會不斷深化.
每個禮拜深化神選民的身份.
就是獨經.
就是重陽神.
就是獻祭.
就是讓摩西五經的教導.
透過利未人的研習班.
讓更加多的以色列人.
對自己的身份認同.
不單止是國民身份.
神選民的身份認同.
你可想而知.
那個重中之重.
都不是說要移民遷徙這麼簡單.
不是要搬遷居民這麼簡單.
而是要讓聖殿有真正的規模.
來活化很多百姓的心靈.
所以十一章你會看到.
第一節百姓的首領住在耶路撒冷.
其餘的百姓.
祭籤每十個人當中.
就使一個人來住在聖城耶路撒冷.
九個人就住在別的城邑.
有他們自己的居所.
即是說他們本身分了支派的地.
他們就在那裡居住.
但是會抽一個人進去耶路撒冷聖城居住.
第二節就是.
凡是甘心樂意住在耶路撒冷的.
百姓都為他們祝福.
是一個蒙福.
他們能夠住在聖城.
來扶持.
來幫助聖殿的工作.

$^{201}$可能他們未必是祭司還是美人.
但是他們可以在聖殿.
他們可以有譬如看守那些城門.
他們可以有其他的工作.
由利兄勿來去委派.
十二章就記載那些名單.
一個很詳細的記載.
接著就有一個很特別的日子.
就是整個城牆城門的奉獻禮.
就是上個星期我們說的十二章.
所以你看到由第七章一直到第十二章.
又是一個很完整的段落.
其實可能只不過是一個很短的時間.
因為建好城牆之後.
差不多去到七月了.
是六月中的時間.
但是七月初一就開始有吹國節.
開始利兄勿和以斯拉兩位的領袖.
先不去做那些行政指令.
這個都是很大的啟發.
不是疏通人心.
是讓神的道在人的心靈裡面.
先來產生感動.
因為只有你兩個人.
你自己內圍的祭司你未人.
自己感覺是神選民沒有用.
很多回去的人根本不覺得自己.
需要這樣來為神毀身.
其實很難做的.
只有很少人去做.
其實很難的.
所以要有一個嚴肅會.
動員所有人是甘心樂意的.
大家回到摩西五經的教導.
到底我們的身份是怎樣呢.
到底在這個時代.
神有什麼託付我們這個民族呢.
為什麼我們的先輩.
他們要回來重建聖殿呢.
重建聖殿又所為何事呢.

$^{241}$輾轉來到差不多有二十天.
因為去到七月二十一日的嚴肅會.
駐貧節.
可想而知在這二十多天裡.
帶來很多百姓的心靈.
得到很大的憤慨.
他們對自己的身份.
有一個很強烈的歸屬.
認定自己神一定有些東西.
一定有些東西託付我們這個民族.
要來到去有一些作為.
所以他們這樣做的時候.
真的完成了那個遷徙之後.
開始有一個很大型的奉獻禮.
就是將整個聖城的每一個城牆.
每一個城門.
來行一個奉獻禮.
由祭司來高沒這些城門.
意味著.
當然不是我們中國人的迷信.
每一個重要的位置要拜一拜.
廉章的那些葵春,紅子,辟邪.
不是這樣.
而是他們真的在整個城牆.
為整個圍繞著耶路撒冷.
來祝福.
由祭司他們在那裡抹油.
意味著整個聖城耶路撒冷.
正正式式恢復一個歸神為聖.
和分別出來之後.
每個禮拜的安息日.
他們就會有恆常的,穩定的.
敬拜的生活.
這個就是由第七章.
一直到第十二章的段落.
所以其實是大概二十日.
二十多日的嚴肅會.
如果嚴格來說.
七月初一只是朝聚.
真正的運作可能是一兩日的時間.

$^{281}$但是由一到第六章.
就是那個李希敏的個人回憶錄.
所以你看看整本書是很工整的.
一到第六章.
李希敏的前因後果.
為什麼他會有這個積分.
本來在很遠很遠的地方.
亞納什西王波斯王朝的書山城.
做一個高級的酒精.
為什麼要來到這麼勞師動眾.
上去耶路撒冷呢.
就是內心有一個passion.
內心有些不抵得.
我是神的選民.
我是猶太人.
我見到,聽到.
現在耶路撒冷的破落.
我是於心不忍.
我是很想有一些的作為.
很想有這個改變.
就是說這些觀察其實不是他的.
這些根本已經變成不是新聞了.
不是什麼突發性新聞.
是人所皆知的.
你是猶太人你就知道了.
不是什麼特殊性的重大的事情.
但是知歸知.
知道又如何呢.
現在你是一個亡國的地方.
你是比周圍的人.
你想一下書山城這麼遙遠.
長途路遠.
山高皇帝遠.
你去到河西.
裡面有他們自己的原居民.
他們的勢力一點都不少.
也某程度是效忠波斯帝國的.
你又如何有這樣的能力.
去改變現存猶太人的光景呢.
人家給你可以苟且偷生.

$^{321}$你有回自己的身份.
不用你逼你去拜偶像.
那你還想要怎樣呢.
還想要有一個很輝煌的.
所羅門時代的聖殿敬拜.
其實沒有可能的了.
也都沒有需要的了.
但是黎希敏.
給我們看到整本書的起義.
就是他有些不甘心.
覺得神選民不應該只是這樣.
所以他就一直等待祈禱.
由第一章到第二章.
他祈禱了差不多四個月.
一直等待.
等待裡面他有計劃.
因為他是一個領袖.
他是一個辦事殷勤.
來盡責的一個酒精.
所以他原來都是一個.
相當有謀略的人.
他就籌劃等待.
神的時間到了.
連亞達斯西王都發現.
他的心靈有些心事.
所以心事不一定是壞事.
你看到牧者傳道人.
時時都有心事.
未必是不好的.
有心事就是說他很記掛.
某些人事物.
大顛大廢的.
完全都沒有任何東西.
當然都是可以很好.
但是你要問一個問題.
他有沒有記掛.
他所服侍的教會的人事物呢.
有沒有為那些事情擔心.
而來到有查犯不思.
而來到有失眠.

$^{361}$或者有其他的壓力呢.
那個本來是正常的.
應該有的責任.
所以尼欽米終於被他的老闆.
發現他心情不妥.
酒精應該是歡樂的.
歡樂得來對他的內心.
有些心事.
老闆不是責難他.
而是真的知道這個.
這麼忠心的牧人.
一定是有些事情的.
那機會來了.
等了四個月.
機會來了.
就向亞達斯西王.
講述他的內心.
我列祖的墳地.
我列祖的地方.
現在這樣破落.
我上去有一些小小的作為.
貢獻我的青春.
來請假.
可能是無薪假期.
上去真的來到.
希望重建城牆.
使那些破落.
不會再繼續嚴重的惡化下去.
有了城牆就有了分別為勝.
有了分別為勝.
這個民族就開始有一些.
可以再深入一點的.
對神的認識的教導.
就是這樣.
不再需要被外邦人任意欺凌.
所以整件事就是這樣.
一到第六章.
就記載尼欽米.
他回去的過程.
那些姦生.

$^{401}$那些遭人白眼.
持份者的攻擊.
各種類型的攻擊.
使尼欽米.
只要令到他的心灰意冷.
只要傳道牧者不要心灰意冷.
就有得做了.
一心灰意冷就沒有事可以做了.
因為哀莫大於心死.
往往都是.
當然不是每個人.
有些傳道離開不是心死.
有另外的帶領.
但是說一個普遍的情況.
如果你已經沒有了那個passion.
很多事情就會像海浪一樣蓋過你.
你就會很難面對.
所以尼欽米面對這些事情.
他在每一章都有祈禱.
他倚靠神.
面對那些困難.
逐件逐件地梳理.
拆解.
使到最終第六章.
就是整場做好了.
但是還沒有獻奉獻禮.
第七章就立刻給我們看到.
發現有整場又如何呢.
都沒有人.
有人又如何呢.
都沒有一個正常健康的.
敬拜的聖殿的事奉.
帶來的只會是慢慢慢慢.
會慢慢陰乾陰乾.
那就沒有了.
所以一步一步.
首先開一個嚴肅會.
既然原來本來想著是一個奉獻禮.
但不是.
推到第十二章才奉獻禮.

$^{441}$先動員所有的人.
聽了能夠明白的.
當時耶路撒冷的人.
一齊來到水門前的寬闊處.
來有這個聚會.
聚會帶來的就是神聖靈的工作.
接著聽了能夠明白.
再加上利未人分小組的演習.
令他們更加認定.
自己神的選民身份.
和外邦人是不同的.
外邦人就是外邦人.
他們是不敬拜神的.
我們既然敬拜神.
我們就不要將我們自己.
這麼重要的信仰.
看為是平常.
看為是一個例行公事.
那就帶來一個有改變.
所以你會看到.
做好這些心靈工程之後.
再動員那些人回去.
再洗到第十二章.
就是將那些祭司利未人.
應該在不同崗位的工作.
擺好位置.
然後預備需要的材料.
物資等等.
所謂官服完職.
大家各按各職去服侍.
聖殿就有可為的地方.
所以就去到十二章.
我們最後那段經文.
就是四十三節.
那一日眾人獻大祭而歡樂.
就是奉獻城門城牆的祭.
因為神洗到他們大大的歡樂.
連婦女帶孩童也都歡樂.
甚至耶路撒冷中的歡聲.
聽到懸處.

$^{481}$這個是我們上一次說.
第十四節就是預備.
他們去敬拜神的各樣物資.
利顯米就分別為聖一個庫房.
收取百姓奉獻的舉祭.
和初熟之物.
每個支派十分之一.
來收納.
就是說讓那個敬拜.
能夠有各樣的資源.
來進行.
使整個聖殿的奉獻.
聖殿的服侍.
回到摩西時代的常規.
這個是非常之重要.
就是將那個福興帶回大家.
都來感覺到自己選民的尊貴身份.
你可以想像.
十二章的四十七節.
當所羅巴伯和利顯米的時候.
以色列眾人將歌唱首們.
每日所當得的份.
供給哈門.
又給利未人當得的份.
利未人又給亞倫的子孫當得的份.
利未人都要奉獻十分之一.
利未人是從百姓那裡收取十分之一.
以色列人的十分之一.
是奉獻給利未人.
利未人將他們收入的十分之一.
奉獻給祭司.
就是說祭司的後人.
祭司的家族.
他們是有百分之一.
就是這樣的意思.
整本書本來如果來到這裡完結.
就很漂亮了.
城門,城牆和那些門樓.
都建好了.
現在是一個很大規模的.

$^{521}$很開心的.
一個奉獻城牆城門的禮.
還要將一些庫房.
一個專屬神的地方.
分別為聖.
將那些奉獻送入庫房裡.
當然有專人來打理.
來記錄和發放.
來處理.
整件事應該由第七章.
七月初一.
一直來到這裡.
這個時間雖然沒有說到.
什麼時候做這個奉獻禮.
但是是一個很完美的結束.
但是好景不常.
十三章全章.
又打回原形.
不得已.
就令到尼熙米極其憤怒.
上星期也提過一點點.
大家如果有機會看過.
你就會感受到.
為什麼一個領袖不在.
整個民族又打回原形.
到底有什麼暗示.
在這本書裡.
叫我們留心呢.
第十三章的第一節.
當日人念摩西的律法書.
給百姓聽.
如見書上寫著說.
阿們人和摩押人.
永不可入神的會.
這是《申命記》的二十章的記載.
二十還是二十三呢.
好像是二十三.
等一下可以站一站.
因為他們沒有拿食物和水.
來迎接以色列人.

$^{561}$且顧了巴蘭就助他們.
巴蘭的故事就是.
《民數記》的二十三章.
甚至應該是二十四章.
但我們的神.
使那就助便為祝福.
這是顯著我們曾經說過巴蘭.
這個先知.
第三節.
以色列民聽見這律法.
就與一切閒雜人.
即是當地人來絕交.
第四節.
先是蒙派管理我們神殿中庫房的祭司.
以尼亞實與多比亞結親.
多比亞就是第四章說過的.
當地的土著.
土豪.
山大王的領袖.
那些人攻擊以色列人.
不准他們回來重建的.
其中一個重要的領袖.
居然現在這個大祭司.
與這個以尼亞實.
對不起.
這個以尼亞實祭司.
與多比亞來結親.
就出現一個肯定的問題.
第五節.
便為他預備一間大屋子.
就是從前修傳素祭與鄉棄名.
和照命令供給利未人歌唱.
守門五谷新走和遊的十分之一.
並歸祭司舉祭的屋子.
即是上文所說的.
就是十二章四十四節所說的.
尼欽米預備了一間庫房.
是用來收取百姓的獻祭的地方.
那個房.
那個地方.

$^{601}$是分別為姓了.
現在祭司不知道用什麼理由.
把這間屋子給了多比亞.
令到很嚴重的問題.
就是那些庫房的物件.
百姓的奉獻.
如何處理呢.
這些技術性的情況問題.
如何解決呢.
這個涉及非常之大的一個問題.
到底他會否中飽私囊呢.
這個多比亞拿了神.
奉獻給神的那些物件.
是祭司利未人才可以享用的.
而他會否在那個房間裡面.
他享用了呢.
他偷竊呢.
這個是很大的問題.
但是似乎容許他這樣做.
到底是用什麼理由可以這樣做.
而可以欺瞞到當時的領導團呢.
這個真的要了解調查一下.
暗示著利未人可能不在.
可能一段比較長的時間都已經不在.
因為以為可以自動幫.
已經搞定了整場城門.
已經有一個敬拜的制度.
已經有十分之一的猶太人.
搬遷入耶路撒冷.
已經設立了祭司利未人的這些崗位.
行常的有歌唱的.
這樣來到去大家都官服原職.
各自在不同的崗位盡心的服侍.
那黎希米呢.
好像應該功成身退.
他本身的職責是在.
波斯帝國的亞特什西王的面前.
來到去做酒精的.
那應該怎樣處理呢.
所以你看到這裡就告訴我們第六節.

$^{641}$那時我不在耶路撒冷.
因為巴比倫王亞特什西王的三十二年.
我回到王那裡.
過了多日我向王告假.
我來到耶路撒冷.
就知道以利亞實.
維多比亞在神殿的院內.
預備屋子的那件惡事.
我甚牢牢就把多比亞的一切家具.
從屋裡都拋出去.
吩咐人潔淨這間屋.
誰將神殿的器皿和掃祭乳香.
又搬進去.
果然就是那個房間.
就是分別為勝的庫房的地方.
居然那間祭司在黎希米不在的時候.
回到述職的時候.
他本來就是猶大的省長.
現在突然之間沒有了這個省長之後.
祭司為什麼會這麼大膽這麼放肆.
來找多比亞.
明知道當時是極力阻止.
黎希米去重建城牆的一個領袖.
四章提到的多比亞.
真是惡貫滿盈.
從屬靈的角度.
來攔阻工程.
為什麼祭司會和他.
這樣好像平輩為奸.
為什麼來容許.
甚至將庫房這麼神聖的地方.
給多比亞居住呢.
很難接受很難想像.
那個聖殿的敬拜.
一下子一落千丈.
一下子好像比以前還要差.
還沒恢復之前.
為什麼會這樣呢.
第十節.
我見利未人所當得的份.

$^{681}$無人供給他們.
是很自然的.
一定是一個比較長的時間.
利未人不再被尊重了.
聖殿的敬拜也是可有可無了.
利未人你們做得不好.
不應該受到這樣的對待.
不應該用最優惠的待遇.
我們不覺得你是神的選民.
神的僕人.
所以你們跟一個普通的猶太人看待.
你要就要.
不要就不要.
對於利未人來說.
他們既然不受尊重的時候.
就經營副業.
你有沒有聽過傳說有人開的士.
我的同學也開過.
他真的.
當然我想我可能是興趣.
還是他找招不得晚.
真的不是開玩笑.
不只是興趣這麼簡單.
我就沒有小巴牌.
不然開小巴專線也好.
但全職就不准.
那就不行了.
這裡告訴我們.
利未人現在要做兼職.
甚至公職的利未人與歌唱的.
驅國奔回自己的田地去了.
表面上不是做兼職.
不是被撈.
實際上就是不受尊重.
回到自己的田地耕種.
挺好的.
傳道人可以來勤奮地下田.
星期天才來港渡.
挺好的.
內地就是這樣.

$^{721}$很多的傳道人.
白天他們要經營自己的工作生意.
星期天就去事奉.
這樣挺好的.
又不用負累教會.
但是是不是這樣呢.
你看看這一段經文.
很清楚的描述.
是因為一下子.
聖殿不被猶太人尊重.
首先不尊重的是祭司.
祭司自己帶頭不尊重.
引入多比亞.
因為跟他有一個恩親的關係.
居然膽大妄為到一個地步.
就是把那些憲制的地方.
給了多比亞.
很自然的.
憲制的地方一定是已經很荒涼了.
荒涼到一個地步.
房間沒有人用.
旁邊也沒有什麼人用.
又有水又有電.
不如租出去.
分開幾個tong 房.
不就可以收入教會.
挺好的.
二樓地下.
找工程來分.
分五間.
上面就分四間.
這樣就有收入.
穩定的收入.
不用這麼多奉獻也夠了.
這樣就可以什麼呢.
其實很簡單.
加個水錶.
加個電錶.
獨立電錶.
不要收多人錢.

$^{761}$這樣就.
你想想多好.
你可以想像到.
對聖殿的輕視.
是由祭司帶頭.
祭司自己引進多比亞.
祭司見到這個地方.
沒有人住.
沒有人用.
不是因為那個地方有什麼問題.
而是根本那個憲制已經出現了一個.
可能結構性的問題的落差.
而導致到.
可能那個結果.
不是那個原因.
原因是領導的祭司利未人.
他們自己都已經受到一個.
可能很大的虧損.
就是不被當時的百姓的尊重.
所以你看到這個現象就是.
他們奔回自己的田地.
他們是有田地.
他們不可以有房屋.
利未人規定了.
他們是要散居在.
他們周圍的支派.
有地方給他們住.
但他們有自己的田地.
他們可以像下蘭那樣.
或者老了.
退休的時候.
都有一些自己的生計.
這個是當時的供給制度.
但是現在是一個不尋常的供給制度.
就是他們要自力更生.
他們沒有了聖殿的憲制.
而來的那些供應.
所以人不是喝水就可以服侍的.
是需要有當時的新瓶.
當時的供應.

$^{801}$物質的供應.
如果沒有的時候.
而領袖不幫他們想.
就只能夠自己想.
所以他們就奔回各自的田地.
有什麼後果呢.
後果就是最終的受害.
就是聖殿慢慢荒涼.
可以有可以沒有.
大家都不再明白摩西的律法書的教導.
出現的後果就是.
百姓更加對神無心.
對神無心的時候就更加.
就出現我們上一次說過的.
安息日的事情.
很多人來做買賣.
都不再care.
所不所謂禮拜日都無所謂.
敬拜神都不是最重要的事情.
這些是屢見不爽的.
一定是慢慢.
有時候很尷尬.
我經常都說.
禮拜見到頂主門打招呼.
他不是上教會.
禮拜又不是上班又不是輪班.
又不是上教會.
會出現在會所的酒樓.
或者是其他地方.
你就開想而知.
又不是在翻了晚堂.
又不是其他的.
原來對敬拜已經慢慢.
沒有了心靈的態度.
這是一個結構性的惡性循環.
所以真正要解決問題.
原來要等尼熙米回來.
所以這一章聖經.
為什麼會突然間有這麼好的.
聖殿的城牆的奉獻禮.

$^{841}$大的奉獻禮.
歡樂之聲聽到遠處.
即是大家極其的感恩.
周圍的人都聽到.
那個歡樂的.
由衷而發出的歡聲.
但現在充滿了的是敵視.
猜忌.
對神無心.
然後告訴我們.
第十一節.
我就斥責官長說.
為何離棄神的殿呢.
我便招罪利未人.
使他們照舊公職.
即是要逐個.
去勸勉他們.
你們的身份是利未人.
你們的服事.
是要在聖殿裡面專職.
專心去扶持祭司.
去做獻祭.
幫百姓獻祭贖罪.
在節期的時間.
或其他聚會.
你們要教導開的研習班.
使百姓能夠聽明.
神的摩西律法書的教訓教導.
即是主學老師.
你應該要教探索班.
初信班.
再好一些的.
再有不同的領袖.
可以教培靈的成長班.
或者團契的團長.
你應該去鼓勵團友.
一起在教會裡面同心事奉.
應該是這樣的運作.
但現在利米米回來之後.
首先就是斥責官長.

$^{881}$然後是勸勉利未人.
這是一個領袖.
他要做的很重要的工作.
第十二節.
由大眾人就把五穀新酒和油的十分之一.
送入庫房.
我派祭司施利米亞.
文士薩督和利未人比大雅.
作副官管理庫房.
副官是哈蘭.
哈蘭是薩克的兒子.
薩克是瑪他尼的兒子.
留意這裡.
第七章曾經說過.
七章第二節.
現在這裡在第十三章再說.
這班人利米米欽點的揀選出來.
在現在這麼嚴重的改革裡面.
好像再一次改革.
第二次的改革.
再一次找一些忠信的人出來.
這些人都是忠信的.
當然是對神.
他們的職份.
他們將所供給的份.
給他們的弟兄.
我的神啊.
求你因這事紀念我.
不要淘沒我為神的殿.
與其中的禮節所行的善.
這個是利米亞回來之後再看到.
而他要運用他省長的職份.
和他的屬靈的影響力.
來去撥亂反正.
你可以想像得到.
很多人歡迎喜歡利米米.
但亦都有一些可能有利益關係的.
潛在的利益衝突的.
是會憎恨他.
攻擊他.

$^{921}$所以他要祈禱.
神啊不要抹去我的善行.
這一份要在你面前蒙紀念的.
為神的殿.
為神的殿去堅守.
為神的殿去站立.
這個是這個時代.
極需要的教會的領袖.
站住在神的殿.
來為百姓.
為神來守住.
那個屬靈的氛圍.
這個是很重要.
這些經文雖然已經.
二千四百多年.
主耶穌未來這個世界之前的四百八十多年.
現在已經過了二千年.
即是二千四百八十多年的人.
一個領袖.
由亞達斯西王.
由波斯帝國的舒山城.
再一次.
那時沒有西鐵.
那時更加沒有高鐵.
要慢慢地走.
有些駱駝的.
這樣一路上去.
是千里迢迢地.
由那麼遠.
由那麼富庶的地方.
回到自己的鄉下耶路撒冷.
為著的是什麼呢.
為著的是要再一次恢復聖殿的敬拜.
所以整個尼希米基是在說.
聖殿的敬拜是需要有聖文.
聖文必須是要有靈明的更新.
你有那個身份證是猶太人.
是.
但是你有沒有那個心靈的.
感動光照.

$^{961}$你是真正的神的選民呢.
很多人都有錦繡堂的洗禮證.
那個洗禮證還要找些相架鑲起.
當然我沒有見過.
我胡說而已.
說笑.
但是那個洗禮證就放在那裡.
那個人就完全不回教會了.
那你洗禮證來幹什麼.
告訴別人.
你洗禮證了.
但你不回教會了.
現在你在哪裡見到你呢.
有時在馬場.
有時在哪個地方.
投注站.
這樣.
那你掛個洗禮證在那裡幹什麼.
所以是聖文的更新.
你回來看到的情況.
一看到.
由祭司這麼高層次的帶頭.
已經去到一種靈性的破落.
就需要來去撥亂反正.
你想想他遇到的困難是很大的.
又不知道在那裡如何興風作浪.
這個尼熙米.
又不知道拿什麼權回來.
又在搞事.
回來耶路撒冷.
最終.
他是為了聖殿能夠恢復.
那個昔日神的心意.
那裡來努力.
那裡來用心.
那裡來招致很多人對他的責難.
因為會受到嚴重的利益受損害.
因為他一撥亂反正.
很自然就會受到損害.
但是他要勉勵利未人.

$^{1001}$你們的身份是要扶持聖殿的憲制.
你們是要有責任.
終身制的.
去教導百姓明白律法書.
你要扶持祭司.
來在聖殿.
恆常的不同的節日.
安息日.
要來去公職.
所以他就是這樣來勉勵.
接著更重要的就是.
要恢復分別為聖的庫房.
要有忠信的弟兄來看守.
來管理.
來發放.
你就可以想像.
這些都是很重要的背後邏輯.
特別強調他們是忠信.
忠心.
信實.
忠心就是說.
不理別人怎麼看.
我是忠於這個信念.
忠於我的神.
這是一個很強烈的個人對神的效忠.
我們需要有弟兄姊妹.
要有忠心.
不是一次做到深夜兩點.
一次是每個人都可以.
但是經常性.
每個人見到.
都一樣是這樣來服事.
有固然好.
榮耀歸給神.
沒有的.
我都是這樣忠心去做.
因為我知道上主是看著的.
這些其實是很簡單.
小朋友都懂得的信念.
但是往往在教會我們很少找到.

$^{1041}$也證明忠信原來是需要尋找的.
不是每個人都可以肩負一些重要的崗位.
需要經過很多的鍛鍊.
很多的試驗.
使那個人真的可以來服事.
哥林多後書有同樣的經文.
保羅這樣來推薦這班人.
哥林多後書的第八章.
你就可以想像原來有些事情.
如果沒有人就做不到.
第八章的第一節.
弟兄們啊.
我把神賜給馬其頓宗教會的恩.
告訴你們.
希望哥林多教會不要再這樣了.
就使他們在患難中.
受到大肆煉的時候.
仍然有滿足的快樂.
即是對舊恩有一個很大的感動.
在患難受大肆煉.
都仍然滿足快樂.
即是主耶穌在他們的生命裡面是很大.
大到一個地步.
苦難很快就過去.
將來去見耶穌.
我得到的讚賞就更大.
這就是在極窮之間.
還格外顯出他們樂捐的厚恩.
我可以證明他們是按照著他們的力量.
甚至是過了他們的力量.
即是說不是斤斤計較.
大家有沒有聽過最近在台灣.
炒到根本是一件沒法想像的事情.
叫做白飯之亂.
一班大學生27個人.
去一間大學附近的開員會.
還是慶功宴.
在附近的那間類似小炒的地方.
去吃飯.
叫菜吃的.

$^{1081}$白飯任裝的.
菜還沒叫齊.
27個人.
兩桶白飯吃光.
吃完之後.
那時已經很晚了.
老闆就說你們叫菜吧.
這麼晚再煮就怕沒人吃浪費了.
這班大學生二話不說.
上網給付評這間餐廳.
接著一整晚的消費.
27個人是6900元台幣.
吃得少到.
這樣一件很小的事.
炒到整個台灣就擦天.
不斷的輿論.
接著大學的系主任.
帶著這班學生.
去向老闆道歉.
道歉完之後.
那些學生回頭.
我們從來沒想過道歉.
是主任抓我們來.
又再大炒大炒.
整個台灣亂到.
有這樣的質素嗎.
為什麼要跟著計較.
最近最後.
老闆說不行了.
不斷給人翻查我過去很多的事.
放上網.
揭露很多的事.
他說現在無限期結業.
整個白飯之亂.
為了白飯.
就是這樣.
你可以想像得到.
為什麼要跟著計較.
因為大家已經沒有那種愛.
沒有那種信任.

$^{1121}$但你看馬其頓宗教會.
保羅說我可以證明.
他們是按照著力量.
即是十分之一.
而且過了力量.
自己甘心樂意的捐助.
再三的求我們.
准他們在這些.
供給聖徒的恩情上有份.
即是保羅你給我們有份.
我們很想去服侍.
我們雖然有限制.
有我們的缺乏.
但我們很想在這個聖宮上有份.
有份他們覺得是一個honor.
你可以想像得到.
跟剛才那個白飯之亂.
是多麼大的一個差距.
你可以想像得到.
有神的感動的人心.
我們只是在人間水平的.
根根計較.
那個據理力爭的問題.
就很難去比較.
所以你留意一下.
保羅就是這樣勸戈林多教會.
第五節.
並且他們所作的.
不但照我們所賞望.
更照神的旨意.
先把自己獻給主.
又歸附我們.
原來一切不是因為.
保羅是我的最喜歡的講員.
你說什麼我們就做.
不是的.
更加照的是神的旨意.
他們先是來獻給主.
即是他們整個的生命.
因為救恩.

$^{1161}$因為保羅所傳的福音.
他們的生命已經改變了.
所以不是因為保羅.
他們要追星.
不是因為Mirror.
我們因為神的救恩改變了我們.
對於事物的分析我們不同了.
所以現在我們認同這個聖公.
公給聖徒.
即耶路撒冷的窮人.
他們很想有份.
他們很想知道.
我們是外邦人.
我們怎樣能夠幫助.
這些神的選民.
現在有缺乏.
先把自己獻給主.
弟兄姊妹.
投入奉獻商的錢.
不是你去施捨給教會.
不是你很慷慨.
算了.
這個月又鬆動.
多些吧.
不是的.
是自己先獻給主.
再歸附了我們.
所以第六節.
因此我勸提多.
既然在你們中間.
你們即哥林多教會.
開辦這些詞彙的事.
就當辦成.
你們既然有在信心.
口才,知識,熱心.
和待我們的愛心上.
都格外顯出滿足.
就當在這詞彙的事上.
也該格外顯出滿足.
即是說.

$^{1201}$哥林多教會.
你們一直都這麼尊重我這個職份.
這麼認同我的事工.
到處去傳道.
你們現在.
為什麼你們到現在.
都還沒有行動呢.
人家馬其頓教會.
是這麼困難.
他們都是一些.
有限的力量的教會.
他們都過了他們的份.
第八節.
我說者說.
不是吩咐你們.
而是藉著別人的熱心.
試驗你們愛心的實在.
不是吩咐.
如果要吩咐.
就沒有意思了.
如果不是自己的心靈的感動.
是沒有意思的.
所以現在就是試一下.
你們的愛心是否實在.
藉著別人的熱心.
來看看你們的愛心.
是否實在.
你們知道我們主耶穌基督的恩典.
祂本來富足.
卻為你們成了貧窮.
叫你們因他的貧窮.
可以成為富足.
我在這事上.
把我的意見告訴你們.
是和你們有益的.
因為你們下手辦這事.
而且有這個心意.
已經一年了.
如今就當辦成這件事.
就是這個辭位的事.

$^{1241}$收取獻金.
奉獻給耶路撒冷的窮人.
既有願作的心.
也當照你們所有的去辦成.
因為人如果有一個願作的心.
就是由你們自己的內心.
出現一份真誠的感動.
神會感動你去做.
就一定會做成的.
既有願作的心.
也當照你們所有的去辦成.
因為人若有願作的心.
必蒙月立.
乃是照他所有的.
並不是照他所無的.
我願不是要別人輕省.
你們受累.
這就叫君平.
就是要你們的婦孺.
現在可以補她們的不足.
使她們的婦孺.
將來也可以補你們的不足.
這就叫君平.
這就叫彼此的服事.
即是說神現在感動婦孺的教會.
去奉獻給有需要的教會.
將來有需要的教會變成婦孺.
又會將她們的婦孺.
奉獻給另外的教會的需要.
這就叫聖經所說的君平.
所以不是大鍋飯.
是誰有婦孺被感動就去做.
所以我們要有一種眼光.
時時留意身邊的人事物的需要.
要有一種敏銳.
不是敏感是敏銳.
我們怎樣將神給我們的婦孺.
隨時都可以去幫助有需要的人.
如經上所記.
多收的沒有餘.

$^{1281}$少收的沒有缺.
引的是嗎哪事件.
星期六安息日之前.
因為安息日要休息.
兩天不能收.
很特別的.
收雙倍.
平日食量大.
會預備多些.
但到第二天都不會剩下.
但那些少的又不夠食.
可想而知那是神跡.
但現在神透過這樣的公書制度.
記住是人的心.
如果不先獻給主.
就沒什麼可說.
不要稀罕這些錢.
因為這些錢一定不是他真心樂意.
出於真心樂意.
神是月立的.
所以不需要擔心.
教會自然會有心人.
來被神感動奉獻.
他會明白這個道理.
我今天的婦孺.
就是神為我預備的.
當然不會抵消他的努力.
他的恩勤.
他的健康.
他的才幹.
他的魄力.
他的機遇.
完全沒有否定.
不過總體而言.
是神的恩典賜給他生命氣息.
他來服侍.
更加有一個很強烈的使命感.
我的出現是神的恩典.
所以我的一切都是神給我的.
我是清清楚楚的一個管家.

$^{1321}$如果我是一個管家.
就不需要去論識了.
第十六節.
多謝神感動提多的心.
叫他代你們恩勤.
像我一樣.
他固然是聽了我的勸.
即是叫他去.
但自己更是熱心.
情願往你們那裡去.
我們還打發一位兄弟和他同去.
這人在福音上得到眾教會的稱讚.
我想說的就是這一句.
你可以想而知.
舒圖保羅也不是蠢的.
看人一些這麼難做的工作.
一年都沒有任何感動.
保羅就是很明白背後的道理.
找到一些忠信的弟兄.
來上心這些事工.
是神的工作.
一呼他們就明白.
而且他們自己是甘心樂意.
這樣去做這些工作.
所以重要的崗位.
如果那個人不靈性好.
那個崗位本身都被他搞亂.
就會帶來很大很大的問題.
甚至有時那些是眼光.
那個眼界.
看那個事物.
是不是屬靈的眼光.
這個是至關嚴重重要的.
所以你會看到保羅.
那些都小事.
你們奉獻都小事.
保羅要的是你們哥林多人.
有沒有像馬其頓教會那樣.
歸服主.
看到聖公的寶貴.

$^{1361}$看到現在有一個很寶貴的契機.
外邦人能夠回饋給猶太人.
很多人未信耶穌.
回饋給猶太人.
未必是那個獻金的大量.
而是那份心意的寶貴.
外邦人得到神救恩.
藉著耶穌.
你們不信的耶穌.
你們的彌賽亞.
現在我們外邦人能夠信主.
能夠得到這麼大的恩典.
現在我們將我們受感動的獻金.
奉獻給耶路撒冷的窮人.
你想想這個這樣的回饋的制度.
背後有一個很大的恩典.
很大的一個屬靈的意義.
所以哥林多人.
如果他們再不在這裡熱心.
其實就很可惜.
沒有對錯.
但是他們會失去很大的屬靈的blessing.
這個就是保羅要來做.
所以你留意一下.
得到宗教會的稱讚.
這個人.
這個誰呢.
我打發一位兄弟.
和提多同去.
是誰呢.
你想想.
是一個多麼寶貴的人.
所以盼望弟兄姊妹.
我們去看回尼希米記的時候.
尼希米在這些重要的崗位.
他要找的是忠信的弟兄.
恢復那個庫房的運作.
收取那些百姓的奉獻.
奉獻給祭司尼美人.
奉獻給聖殿的需用.

$^{1401}$讓整個聖殿的運作.
回到正常.
恢復到敬拜的熱情.
這個是尼希米回來的時候.
進行的改革.
撥亂反正.
這個就是第十五節開始.
很自然的.
靈性出現了問題.
不再care聖殿的奉獻敬拜.
很自然就留意自己的日常生活.
貿易往來.
哪裡有便宜的東西.
哪裡的東西大減價.
LV開倉五節.
怎麼都要排隊.
三點鐘都要排隊了.
第十五節.
那些日子.
我在猶大見有人在安息日榨酒.
搬運禍菌台在路上.
又把酒葡萄無花果和各樣的擔子.
在安息日擔入耶路撒冷.
我就在他們賣食物的那天警戒他們.
這些都是外邦人.
雞仔不管管馬鷹.
馬鷹不管管雞仔.
關我什麼事.
是你們有買有賣.
我拿進來賣你們可以不買.
就是說猶太人.
他們已經對於安息日.
沒有在那種昔日的改革的時候.
那種熱情.
現在需要又再改革了.
需要再一次來到去有一個憤興了.
警戒這些人.
第十六節.
又有推羅人住在耶路撒冷.
他們把魚和各樣的貨物運進來.

$^{1441}$在安息日賣給猶大人.
我就斥責猶大的貴族.
所以就是那個矛頭所在.
那個焦點.
那個根源.
你們怎麼行這惡事犯了安息日呢.
從前你們列祖豈不是這樣行.
以致我們神使一切的災禍.
臨到我們和者成嗎.
禮拜九點鐘.
好像說笑一樣.
和十一點鐘.
走上去會所.
洪經理.
為什麼你容許我的會友在這裡喝茶.
關我什麼事.
他們訂桌子的.
難道我不給他們吃嗎.
好像說笑一樣.
但現在就要警告這些貴族.
就是說黎希敏一不在.
馬上一落千丈.
馬上那些事情.
本來大家都不再計較了.
現在又回頭.
又是容許.
還要是各種的必需品.
和這些的貴族.
必需品和這些的.
很豐富的果實.
又回到安息日裡面.
擔入去賣.
大家都吹朱若霧.
有便宜的東西.
安息日最多人聚集.
那也是.
真的最多人的時間.
就是猶太人的安息日.
所以做買賣的.
當然找這些日子最多人的.

$^{1481}$難道禮拜一來教會.
人影都沒有.
只有Esther在.
那你怎麼可能會.
有事情可以做到呢.
當然是禮拜一才最多人.
所以你可以想像得到.
如果將來我們的廣場.
九點鐘和十一點鐘.
不是茶水檔.
是很多人來擺賣.
然後有其他的東西.
你可以想像得到.
猶太人就是這樣.
就是落在這樣的一個.
再一次對神的敬拜.
沒有心.
再一次對神的聖殿.
已經是看作平常.
所以這裡.
尼希米第十八節.
從前你們列祖豈不是這樣行.
以致我們神洗了一切的災禍.
臨到我們和者成嗎.
現在你們還犯安息日.
洗到憤怒越化.
臨到以色列.
在安息日的前一日.
耶路撒冷的城門有黑影的時候.
因為是星期五的黃昏.
我就吩咐人將門關鎖.
不過安息日不准開.
我又派我幾個僕人管理城門.
免得有人在安息日.
擔什麼擔子進城.
於是商人和販賣各樣貨物.
一兩次住宿在耶路撒冷城外.
我就警戒他們說.
你們為何在城外住宿.
若再這樣我必下手拿辦你們.

$^{1521}$不要忘記他是省長.
尼希米.
即是他也有一些行政權.
所以這些人蠢蠢欲動.
說說而已.
嚇嚇我們而已.
我們做著做著.
總是有顧客.
現在不是沒有.
是他在攔住.
所以如果我們的東西漂亮又便宜.
猶太人就一定喜歡.
所以他們都在城外準備入城.
既然這麼堅決的執行力.
很強的執行力.
即是我要拿辦你們.
所以從此以後.
他們在安息日就不再來了.
我吩咐利未人潔淨自己.
來守城門.
所謂潔淨自己.
不只是洗澡.
用什麼沐浴露.
而是他們要自潔.
明白自己的位份.
自己的身份.
自己是神的僕人.
利未人.
守住城門是一個重要的天職.
真的我守住的是百姓的心.
敬拜神.
真的來讓百姓.
可以心無旁騖的來敬拜.
過了安息日.
大家敬拜完神.
一樣可以去做買賣.
是沒有問題.
神不是讓我們不惜人間煙火.
而是安息日是一個很重要的指標.
是神和他的選民的立約的記號.

$^{1561}$是以色列人的摩西之約的記號.
所以耶穌基督未來這個世界之前.
他們是要守住這個身份.
這個神聖的約的記號.
是守安息日.
所以這個就是很重要.
舊約的人他們要守住.
我們新約.
我們就已經.
耶穌說再不是在耶路撒冷.
再不是在基利森山敬拜的.
是要用心靈和誠實去敬拜神.
所以我們不需要死守.
反正現在有星期二晚的.
一些教會有崇拜.
元基堂有崇拜.
星期二晚.
星期六更加有.
星期日就更加有.
所以和不同地方都有崇拜.
所以不再是局限於那個地點時間.
仍然都是那個情況.
就是心靈誠實的敬拜神.
所以我們不是要像猶太人那樣.
安息日什麼都不能做.
不是這樣的意思.
而是我們要建立一個.
隨時都讓神的聖靈帶領和感動.
這就是他們的情況.
第22節.
利未人要潔淨自己.
來守城門.
使安息日為聖.
就是這個意思.
安息日為聖.
不是那個日子.
是那班人帶著一個虔誠.
敬拜神的心靈去敬拜神.
所以小朋友我們要自小教導他.
星期日我們一起回教會敬拜神.

$^{1601}$要教導的.
要分別為聖的.
即是和其他六天是有點不同的.
他要慢慢從實踐.
從觀察.
他會慢慢去學到.
到他長大.
他自己作為當家作主.
他自己成家立室.
他自己成為成年人.
譬如他自己去外國讀書.
他也要突然有個意識.
就是禮拜媽媽不在了.
沒有人再餓我了.
但是我有感動.
我是神的兒女.
我就應該找個禮拜堂去敬拜神.
這個教導.
就是需要從小就這樣培育.
所以分別為聖.
是這個意思.
我的神求你恩謝事紀念我.
照你的大慈愛連恤我.
將一些這麼大的.
做到慣性.
又再一次的改革.
又再一次的扭轉.
其實又是很難的.
好像不是與人方便.
回來就在左頭左勢.
在這裡去唧唧噤噤.
但是黎熙敏看到的就是.
我們是神的選民.
神仍然給我們.
未有主耶穌來這個世界之前.
我們是要守安息日.
這個是鐵律.
不能更改.
我都沒有權更改.
反而就是我們應該.

$^{1641}$怎樣將這個安息日轉化.
不是死守.
耶穌的時代當然就是要.
反過來就是要叫他們不要死守.
現在我們應該怎樣將安息日.
看為是我們去敬拜朝見神.
的一個神聖的日子.
而這個神聖帶來我們的生命.
是會有不斷的進心的認識.
這個就是真正背後的重要.
所以難道真的買賣.
遲一秒鐘一分鐘.
都是神要追究嗎.
而是那班人的心.
第23節.
那些日子我也見尤大人.
娶了阿實達阿們摩訶的女子為妻.
他們的兒女說話.
一半是阿實達的話.
不會說尤大的話.
所說的是照著各族的謊言.
我就斥責他們.
就坐他們.
打了他們幾個人.
拔下他們的頭髮.
叫他們止住神氣勢.
必不將自己的女兒.
嫁給外邦人的兒子.
也不為自己和兒子.
娶他們的女兒.
我又說.
以色列王所羅門.
不是在這事上犯罪嗎.
在多國中並沒有一王將他.
且蒙他神的所愛.
神立他作以色列全國的王.
然而連他.
這麼有智慧.
這麼有才幹.
這麼得蒙神恩寵的.

$^{1681}$這個所羅門王.
都被外邦的女子引誘犯罪.
因為拜那些不同民族的神.
所以他娶了不同的.
有些人說是政治婚姻.
即是說你娶了那些公主.
自然就會天下太平.
就締結和約.
就相安無事.
你的國家又會更加壯大.
所以這是所羅門的智慧.
但是被自己的智慧影響.
就沒有想過.
原來他娶了這麼多公主.
這些公主.
人家飲食的是假神的養分.
嫁了給你.
他就不會跟你去信耶和華.
反過來就會帶所羅門去拜偶像.
所以所羅門就拜很多很多偶像.
27節.
如此我起聽你們行者大惡.
娶外邦女子干犯我們的神.
大祭司以利亞實的孫.
耶和耶大的一個兒子.
就是和倫人參巴拉的女婿.
我就從我這裡把他趕出去.
我的神求你紀念他們的罪.
因為他們沾污了祭司的職任.
違背你與祭司利未人所立的約.
這個約你要追究在摩西五經裡面.
尤其是利未記.
另外就是民數記.
就有這個榮約與祭司利未人立了.
所以現在這個大祭司的兒子.
與多比亞參巴拉結盟.
結親.
就帶來和幻.
將世俗的價值觀引進聖殿裡面.
30節.

$^{1721}$這樣我潔淨他們.
使他們離絕一切外邦人.
派定祭司和利未人的班次.
使他們各盡其職.
我又派百姓按定期獻柴.
和初熟的土產.
我的神求你紀念我.
施恩與我.
最多最多的是每一張利兄妹的祈禱.
發自內心感覺.
需要內心的力量.
內心的力量就是使他看事物.
看得清楚.
看得對是神的旨意.
他就不理人的反對.
就堅決的去做改革.
這個是你個人的野心.
意氣用事.
因為是神的旨意.
可以是很近似的.
甚至乎嚴重來說.
可以是一線之隔.
可以是你個人的一個野心.
是你個人的一個好大喜功.
但是如果這個是神的旨意.
這個民族.
這個選民.
是需要撥亂反正.
回到摩西的律法書.
神選民的那個職份.
那個規模.
就是同外邦人離絕.
利兄妹記就是想探討這個課題.
透過建立一個城牆.
分開了外邦人.
透過將聖殿恢復規模的敬拜.
一個秩序.
井井有條.
大家都是帶著一個歡樂的心靈敬拜神.
但是現在離開了一段時間.

$^{1761}$又再一次打回那個原形.
那個原形就不是聖經要求的原形.
是人要求的水平.
就是無所謂.
不要緊.
都是一樣.
敬拜可以每個禮拜都有.
但是有這些買賣.
都是重要.
那就變了慢慢慢慢.
就是無所謂.
是帶來那個世俗化進入.
神的選民的那個敬拜裡面.
所以這一張聖經.
好像突然之間.
令到那個屬靈的奉獻禮.
嚴重的失色.
要再補記這一張.
不是刻意的來到去.
想要挖苦那些猶太人.
而是很想告訴我們.
一個地方原來一直到神還沒有回來之前.
都極度需要領袖.
屬靈的領袖.
一個稱為神明下的地方.
需要有這樣的眼光的領袖.
有這樣的領袖.
自然就會有一個恢復.
自然就會有一個改革.
自然就會帶來一個宗教的復興.
所以也可以反過來.
那個領袖不在.
就會帶來可以人.
原來真的有罪性.
一下子可以再掉回一個深淵的谷底.
很快就馬上沒有了那個屬靈.
你可以想像得到.
這個就是人類共同的命運.
如果沒有神的僕人的存在.
來禱告.

$^{1801}$來去議題面明的.
無論是警告,勸勉.
提摩太后書第四章.
第四章第一節.
我在神面前.
並在將來審判活人死人的基督耶穌面前.
憑著祂的顯現和祂的國度去祝福你.
祝福是命令charge你.
務要全度.
無論得時不得時.
總要專心並用百般的忍耐.
各樣的教訓.
責備人,警戒人,勸勉人.
因為時候要到.
人必厭煩純正的道理.
以妥法養.
就隨從自己的情慾.
增添好些師父.
並且掩耳不聽真道.
偏向荒謬的言語.
你卻要凡事謹慎.
忍受苦難.
作全度的功夫.
盡你的職份.
就是全度.
要謹慎.
要忍受.
要忍耐.
忍耐什麼呢?.
有苦難的.
你盡你的職份去做.
現在我在神.
和將來審判活人死人的基督耶穌面前.
來祝福你.
還要憑著祂的顯現.
和祂的國度去祝福你.
嘩!說得這麼恐怖.
保羅.
你想我繼續去傳道.
說得這麼恐怖.

$^{1841}$提升到這麼高層次.
神和基督耶穌.
在這位將來要審判活人死人的基督耶穌面前.
還要憑著祂的顯現.
就是說祂一定會回來的耶穌.
即是你的老闆一定會回來的.
我現在憑著祂的顯現來祝福你.
祝福什麼呢?.
就是務要傳道.
和你現在要盡你的職份.
即是說神已經揀選了你.
有沒有聘書是另一個問題.
你要忠心地在你有生之年.
去做傳道.
來盡你的職份.
你要凡事謹慎.
你要來忍受苦難.
這個也有很多教導的.
神樂院也有很多提醒的.
老師也有很多經驗的.
也告訴我們.
去到工場是有很多不同的困難的.
總之總是比你在預期.
你已經有心理準備.
還要多.
你就知道.
那些保羅所說的忍受苦難.
是很實在的.
所以不需要擔心.
傳道人就是用苦難去養大的.
所以去到一個年紀.
他的道開始就真的能夠.
祝福很多人.
就是他被苦難一直養.
越養越大.
越養越成熟.
越養越覺得務要傳道.
或是反對的聲音.
總之是謹慎忠於所託.
這是傳道人很重要的.

$^{1881}$所以苦難只會令我們更加成長.
更加成熟.
更加謙卑.
更加殷勤.
來將神的道原原本本說明白.
不是為義氣之爭.
不是借題發揮.
不是指桑罵槐.
不是.
年輕的時候可能是.
但慢慢慢慢就發現.
其實不需要.
好好地將神的道說明白.
就已經足夠.
因為聽完之後.
每個人自己面對神的聖靈的感動.
所以這是很重要的.
黎軒偉就是在他的崗位裡面.
他盡忠.
離開一會兒.
整個國家又再一次.
跌回一個很低的水平.
又好像和貿易.
和日常生活的民生.
又拉上關係.
好像不能夠提升靈命.
所以一個地方.
不可以沒有了屬靈的領袖.
那個領袖就是有那個眼光.
去研判這個民族.
這個群體.
這一班選民.
他們需要的那個屬靈的餵養.
所以這是很重要的一個提醒.
十三章是兩個角度.
原來人沒有了屬靈的提醒.
是會自己慢慢地跌低.
都不為意的.
但有屬靈的提醒.
又可以再回到神期望的水平.

$^{1921}$所以今天我們有沒有寶貴.
我們是神的選民.
還是我們很寶貴.
我們自己現在地上的身份.
你隨便打開那個錢包.
哪間會所的會員.
馬會的會員.
什麼遊艇會的會員.
什麼高球會的會員.
什麼的會員.
但是神選民的身份.
是不是放在最高的位置呢.
孫家的弟兄姊妹.
或者其他教會的弟兄姊妹.
我們對神選民的身份.
有沒有這樣的榮耀感呢.
有沒有這一份的幸福感.
我們能夠因為主耶穌的救恩.
我們能夠現在在將來的永恆裡面.
我們是有份的.
所以現在我們所做的一切.
都是為神去效力.
盼望神祝福我們.
今天我就講到這裡了.
因為待會要開會.
多謝主.
讓我們都完了《尼希米記》.
盼望我們下一個禮拜.
我們看的《以斯帖記》.
又是一個新的歷史觀.
我們都能夠從不同的角度.
去看看被擄歸回的時期.
有神興起的僕人侍女.
他們的忠心.
他們那個忠信的服侍.
我們很盼望.
孫家有忠信的弟兄.
弟兄姊妹.
一同的起來服侍.
我們默默地在神裡面去做.

$^{1961}$一定是有神的祝福和紀念的.
我們這樣仰望感恩.
禱告奉耶穌基督的名.
阿們.
(音量注意).
\newpage



\section{歌羅西書 2:1-19}
\label{sec:epe4ZHzeWUY}
\textbf{2023-07-22 晚堂《打穩根基》歌羅西書 2\_1-19 講員:謝靄薇牧師}
\newline
\newline
連結: \href{https://youtube.com/watch?v=epe4ZHzeWUY}{\texttt{https://youtube.com/watch?v=epe4ZHzeWUY}} ~~~~ 語音日期: 2023-07-25
\newline
\newline
\hyperref[sec:zajnPwak_dU]{\small{< < < PREV SERMON < < <}}
~
\hyperref[sec:index_chronic]{\small{[返順時目]}}
~
\hyperref[sec:index_scriptual]{\small{[返順卷目]}}
~
\hyperref[sec:dFeC0BqGntE]{\small{> > > NEXT SERMON > > >}}
\newline
\newline
歌羅西書 2:1-19
\newline
\begin{longtable}{cl}
\hline
\hline
章節 & 經文 (和合本修訂版)\\
\hline
2:1 & \begin{tabularx}{0.7\textwidth}{X} 我要你們知道，我為你們和老底嘉人，和所有沒有與我見過面的人，是何等地勤奮； \end{tabularx} \\ \\ \relax
2:2 & \begin{tabularx}{0.7\textwidth}{X} 為要使他們的心得安慰，因愛心互相聯絡，以致有從確實了解所產生的豐盛，好深知神的奧祕，就是基督； \end{tabularx} \\ \\ \relax
2:3 & \begin{tabularx}{0.7\textwidth}{X} 在他裡面蘊藏著一切智慧和知識。 \end{tabularx} \\ \\ \relax
2:4 & \begin{tabularx}{0.7\textwidth}{X} 我說這話，免得有人用花言巧語迷惑你們。 \end{tabularx} \\ \\ \relax
2:5 & \begin{tabularx}{0.7\textwidth}{X} 雖然我身體不在你們那裡，心卻與你們同在，很高興見你們循規蹈矩，對基督的信心也堅固。 \end{tabularx} \\ \\ \relax
2:6 & \begin{tabularx}{0.7\textwidth}{X} 既然你們接受了主基督耶穌，就要靠著他而生活， \end{tabularx} \\ \\ \relax
2:7 & \begin{tabularx}{0.7\textwidth}{X} 照著你們所領受的教導，在他裡面生根建造，信心堅固，充滿著感謝的心。 \end{tabularx} \\ \\ \relax
2:8 & \begin{tabularx}{0.7\textwidth}{X} 你們要謹慎，免得有人用他的哲學和虛空的廢話，不照著基督，而是照人間的傳統和世上粗淺的學說，把你們擄去。 \end{tabularx} \\ \\ \relax
2:9 & \begin{tabularx}{0.7\textwidth}{X} 因為神本性一切的豐盛都有形有體地居住在基督裡面； \end{tabularx} \\ \\ \relax
2:10 & \begin{tabularx}{0.7\textwidth}{X} 你們在他裡面也已經成為豐盛。他是所有執政掌權者的元首。 \end{tabularx} \\ \\ \relax
2:11 & \begin{tabularx}{0.7\textwidth}{X} 你們也在他裡面受了不是人手所行的割禮，而是使你們脫去肉體情慾的基督的割禮。 \end{tabularx} \\ \\ \relax
2:12 & \begin{tabularx}{0.7\textwidth}{X} 你們既受洗與他一同埋葬，也就在此禮上，因信那使他從死人中復活的神的作為跟他一同復活。 \end{tabularx} \\ \\ \relax
2:13 & \begin{tabularx}{0.7\textwidth}{X} 你們從前在過犯和未受割禮的肉體中死了，神卻赦免了你們一切的過犯，使你們與基督一同活過來， \end{tabularx} \\ \\ \relax
2:14 & \begin{tabularx}{0.7\textwidth}{X} 塗去了在律例上所寫、敵對我們、束縛我們的字據，把它撤去，釘在十字架上。 \end{tabularx} \\ \\ \relax
2:15 & \begin{tabularx}{0.7\textwidth}{X} 基督既將一切執政者、掌權者的權勢解除了，就在凱旋的行列中，將他們公開示眾，仗著十字架誇勝。 \end{tabularx} \\ \\ \relax
2:16 & \begin{tabularx}{0.7\textwidth}{X} 所以，不要讓任何人在飲食上，或節期、初一、安息日等事上評斷你們。 \end{tabularx} \\ \\ \relax
2:17 & \begin{tabularx}{0.7\textwidth}{X} 這些原是未來的事的影子，真體卻是屬基督的。 \end{tabularx} \\ \\ \relax
2:18 & \begin{tabularx}{0.7\textwidth}{X} 不要讓人藉著故作謙虛和敬拜天使奪去你們的獎賞。這等人拘泥在所見過的幻象，隨著自己的慾望無故地自高自大， \end{tabularx} \\ \\ \relax
2:19 & \begin{tabularx}{0.7\textwidth}{X} 不緊隨元首；其實，由於他全身藉著關節筋絡才得到滋養，互相聯絡，靠神所賜的成長而成長。 \end{tabularx} \\ \\
[1ex]
\hline
\hline
\end{longtable}
$^{1}$各位大人 節哀平安.
在我們看神的懷疑之前.
我們一起祈禱.
「沒有恩典 沒有慈愛的主.
當我們來到你面前.
何等的歡喜快樂.
因為充滿了你的恩典.
以至你聖靈在當中感動我們.
若然我們有什麼是不蒙你喜悅的.
求主去挪開 剔除.
以至我們可以轉身可以看到你對我們說話.
如你在當中給我們悟性.
我們可以明白你的真道.
以至我們可以走出來.
我們這樣祈禱 教託 奉主耶穌基督的名求.
阿們.
任何的建築物都離不開根基.
根基很重要.
無論將來的建築物是怎樣.
多麼的宏偉.
或者是平時我們住的樓宇.
但是根基是不可以忽視的.
因為這個工程做得不好的話.
就會影響你.
颳颱風的時候.
或者是地震的時候.
更加很明顯.
一個很重要的考驗.
特別是看到有些地區.
曾經地震之後.
你發現原來裡面的水泥不是全部塞滿的.
原來裡面還有塞滿膠樽的.
很多其他不重要的東西也塞進去了.
你會看到偷工減料.
所以根基是很重要的.
我們看神的話語的時候.
同樣想到.
弟兄姊妹這個屬靈的生命.
都需要留意的.
是不可以輕忽的.

$^{41}$這個圖雖然很小.
但是你要留意那幾個黃色的部分.
你會看到今天我們說的.
哥羅西教會.
它在史瓦斯亞以弗所東邊的一個市鎮.
在左邊那裡.
就是以弗所.
在我右邊就是哥羅西.
然後上方就是一個老底家.
在這個老底家.
你會看到原來是當時重要的金融中心.
哥羅西和它距離只有19公里.
不是很遠.
當我們看保羅.
在一章29節的時候.
他就告訴我們.
他是很辛勞的去傳福音.
但我們又知道這個教會.
不是他去建立的.
但是他曾經在以弗所派了童工.
去到哥羅西的地方.
所以保羅他認為自己.
都應該對這個地區負上一些責任.
為什麼呢.
因為當童工去到當中的時候發現.
原來哥羅西的教會.
是出現了一些假教師.
假教師他堅持什麼呢.
你要認識上帝.
要完全的得救.
你必須要加一些東西.
去崇拜某些靈界的首領.
或者是權威.
人必須要守一些禮儀.
這些禮儀就很糟糕了.
這些禮儀.
猶太人他們有國禮.
甚至還要去嚴格的遵守一些有關食物.
或者是其他的規條.
你會看到原來這個教會裡面.

$^{81}$都出現了一些異端的教訓在裡面.
所以我們在屬靈的生命上.
我們要成長.
我們必須首先要懂得分辨.
去分辨真理.
異端很多時候就是.
讓你覺得.
他們有一些特別的智慧.
他們很注重這些.
令到人們很羨慕.
很想學他們的追求.
所以今天我們會將聖經裡面.
今天讀了很多經節.
我們今晚就會消化完它.
首先就是.
在開始的時候你會看到.
二章四字那裡說到.
免得有人用花言巧語迷惑你們.
花言巧語就是說.
有些人說話很厲害.
但是說到學術的時候.
他很動聽.
說著說著.
你都不知道誰是真誰是假.
可能你都跟著他.
那就糟糕了.
保羅提醒我們.
不要被他們迷惑了.
甚至去到.
第八節那裡.
就說到要謹慎.
那裡說到.
他們的哲學.
現在的和修本.
就是虛空的廢話.
以前我們看畫本的時候.
就看到理學和虛空的妄言.
總之一句話來說.
就是一些欺騙人的哲學.
是空虛的.

$^{121}$這些哲學.
一定不是照著耶穌基督的教訓.
去教人.
所以裡面夾著一些傳統.
是人的傳統.
是人的說話.
不是神的說話.
所以保羅就形容這些是.
世上的小學.
就是說很膚淺的.
同時.
在十六節那裡就說到.
當人們信了主之後.
猶太人以前有很多.
飲食上去留意的節期.
就好像現在我們中國人都有.
初一十五.
信了主之後.
有些人就懷疑.
我還要不要做以前的事呢.
好像很怕.
那些神靈會不會得罪了他呢.
保羅就提醒.
你不要因為你沒有做這些.
就覺得自己低於別人一等.
以致被別人評斷你.
或者審判你.
他說千萬不要有這樣的想法.
所以叫他們要很小心.
甚至是.
這個重罪.
十八節說到敬拜天使.
甚至教人要故作謙虛.
故作謙虛是什麼意思.
謙虛本來是一個美德.
但是故作的話就是假的.
不是真心的.
是一個謙虛的態度.
保羅就說.
你不要聽到他們這樣說.

$^{161}$以致.
拿不到神給你的獎賞.
將來我們都有一個.
生命的冠冕是給你的.
但是如果你跟著.
這班人去做的時候.
敬拜天使當然是.
更加不是合乎神的旨意.
只有我們的神.
我們可以敬拜.
除他以外.
其他都不是神.
所以.
這些教導一看.
就不是神在教導我們.
所以是很重要的.
還有.
二十一節.
那裡更加說到.
不可難,不可常,不可摸.
這類的規條.
更加是.
如果你一碰到.
跟著他走的時候.
原來是很危險的.
在當時.
他們很流行.
洛斯底主義.
洛斯底也是很注重知識的.
在第二,第三世紀的時候.
是很盛行的.
他覺得神並非.
掌控人命運的.
唯一.
而且基督只是一個受造物.
是最高等的.
他這樣看的.
他也覺得.
人還要去敬奉.
其他的諸靈.

$^{201}$星獸當中的諸神.
還要加一些東西.
才合乎標準.
所以保羅就在斥責他們.
這些全部是很膚淺的.
所謂的.
世上的小學.
人的傳統.
二端很喜歡用人的傳統.
去教導人.
其實人的傳統裡面.
如果沒有真理,沒有根基.
是沒有用的.
所以耶穌基督經常去責備.
法利賽人.
他說他們死守.
一些傳統.
但是.
他們違背了神的誡命.
離開了神的誡命.
所以耶穌去責備他們.
這裡還教你.
就已經不對了.
所以保羅要叫他們.
很謹慎,很小心.
他很恐怕.
信徒.
受了人.
這樣去教導的時候.
就離開了.
偏離了真道.
偏離了神的道.
以致將來.
不是屬神的.
本來神透過耶穌基督.
給我們這麼豐盛的恩典.
這麼豐盛的生命.
反而去隨從.
這些二端所教你的.
反而你失去神給你的.

$^{241}$冠冕,神給你的獎賞.
所以.
他說千萬不要.
去跟從這些.
第二方面.
他教我們要.
生根建造在哪裡.
原來建造在主的.
懷裡上.
大自然的植物.
要植物向上.
要去生長.
你一定做很多功夫.
要除蟲,施肥.
要燒罐.
有很多事情要做.
或者打風的時候.
可能你所做的努力都沒有了.
所以要花很多心機.
怎樣可以生根建造呢?.
原來.
在教會裡.
對於頂尖寺門的.
屬靈的生命.
是要教導真理的.
上次我們有福音聚會.
報道會.
或者有個人報道.
領了人歸主.
他決志了.
但之後怎樣呢?.
之後有沒有好好跟進呢?.
去栽培他呢?.
做栽培的人要花很多功夫.
我們其實不是.
只是陪他聽一聽.
報道會.
我們就贊就完了.
不是的,你要時時陪他聚會.
關心他.

$^{281}$打電話問候他.
有很多事情要跟進.
教栽培班的導師就行了.
還有很多事情要跟進.
可能他又遇到什麼事情呢?.
所以其實.
是希望他.
是好的.
在神的話語裡打好根基.
如果他只是信主.
今天舉了手.
但之後就不見了人.
他就流失了.
因為他根本就不明白.
神的話語是什麼.
不明白這群人.
為什麼我們要祈禱.
完全沒有人教他.
他不知道.
所以他就流失了.
因為我回來不回來也不要緊.
所以保羅.
去教導提摩太.
去栽培他成為一個.
教會的領袖之外.
他還教導他怎樣呢?.
他說你要教導其他人.
正如提摩太後書2:15.
要竭力在神面前.
做一個經得起考驗的.
無愧的工人.
就應該安著正義.
去講解真理的話.
自己都要去.
好好努力.
才可以去告訴別人.
所以這裡是說.
要教導的人要花點功夫.
你要實際的.
去教導他.

$^{321}$讓他明白.
明白之後他要做得出來.
才算數.
你教導他可能水過阿背.
可能那晚精神狀態.
不太好.
可能他下班了也很累了.
你那晚跟他說了一個小時.
可能下星期他都忘記了.
可能你又會很失望.
但不要緊的.
神要我們用心.
要竭力.
去教導他.
所以我們要花點.
功夫去責任.
另外就是大家很熟悉的.
馬太福音28章.
講到大使命的經文.
他教導我們.
就要去教導.
教導他.
去遵守.
不是教導他.
要他去遵守.
直到世代的終結.
是很長時間的.
一個培訓.
所以教導是很重要的.
因為教導.
弟兄姊妹就會.
有根有機.
如果他決志.
沒有英武裝.
就什麼都沒有了.
因為他都不知道.
主耶穌跟他的關係.
原來我是這麼卑微的.
原來我要去敬畏神.
如果沒有人教他就不會了.

$^{361}$是不是.
當你慢慢讀正經的時候.
發現原來我們這麼卑微的時候.
我們應該要感恩.
我們每一天對神應該有一個敬畏的態度.
神給我生命氣息.
我們應該珍惜.
神叫我們要去傳福音.
我們要去講.
不懂得講也要講.
因為神在你身上有太多恩典.
你要跟別人分享.
所以要教導.
教導之後.
你就會對基本的認識.
知道一點.
還有你發現.
房子有了框架.
有了槓根.
加上水泥.
鑽石.
砌好了.
是不是就完了.
原來不是的.
教會裡每一個信徒.
要怎樣去生根建造在整個話語裡.
不是說.
你自己好好的追求就夠了.
原來大家要一起.
因為大家都是在耶穌基督的身體裡.
所以.
生根建造.
還要合而為一.
如果我們.
每一個都追求.
當然好.
但每一個都追求的時候.
會不會每一個都去比較.
我比你厲害.
我今天去哪裡進修.

$^{401}$大家好像要拿學位去比較.
這樣就完蛋了.
原來我們讀了.
是為了怎樣去事奉.
怎樣去討神的喜悅.
於是我們這群人走在一起的時候.
我不行.
你幫我.
做功課.
大家一起討論的時候.
就會領悟多一點.
大家要.
合適才行.
教會裡也要合適.
無論你來的年數.
是早一點或是晚一點.
可能有些年紀大一點.
有些年輕一點.
可能混不來.
但聖經教我們要合而為一.
怎樣合而為一.
因為我們的基礎就在耶穌基督裡.
無論你是老還是年輕.
其實可以混在一起.
可以互相扶持.
所以合而為一.
是很重要的.
因為.
《避一前書》2章5節教導我們.
你們作為活石.
原來我們每個人.
在神裡.
在靈宮裡是一個活石.
這個活石要.
發揮你的功能.
所以我們.
每一個人都是耶穌基督.
耕種的田地.
也是.
他建造的房屋.

$^{441}$所以你不可以藉口說.
多一個不重要.
少一個也不重要.
反正我都沒什麼用.
多一個也不多.
只不過多一個人擺位置.
教會就會興旺嗎.
如果有這樣的想法就錯了.
因為在屬靈的生命裡.
因為在說屬靈的生命的建立.
所以不能少了你一個.
這些磚頭.
少了你一個.
就不完整了.
人們走路.
就掉下去.
所以少了你一個是不行的.
所以是活石.
這些磚沒有泥.
就會掉下去.
所以要黏在一起.
所以教會裡大家要合而為一.
是很重要的.
就好像你建造房子.
原來有一個角落.
要犧牲一塊磚.
要磨掉角落.
才可以黏在一起.
原來神做我們每一個人.
是剛剛好的.
在這個靈宮裡.
有一些驕傲的.
一些不明要神的.
他就磨掉你.
然後讓你可以很順利的.
是剛剛好 砌得好的.
所以我們大家就在這個.
砌得好的靈宮裡.
就發揮作用了.
成為聖潔的祭司.

$^{481}$耶穌基督獻上.
夢神如立的.
屬靈的祭物.
所以這個很重要的.
所以這個哥羅西教會裡.
出現什麼.
有破口.
有破口就不可以合而為一.
有破口.
異端就會乘虛而入.
所以破口就很重要.
風吹雨打.
怎麼去倒塌.
有時候看電視.
好像骨牌效應.
倒塌了.
原來下面的水泥可能有問題.
經過水浸之後.
全部倒塌了.
看到這個教會裡的根基.
裡面的人.
有沒有合而為一呢.
哥羅西教會裡.
出現了問題.
就是我們剛才所說的.
這些異端的問題.
就是因為他們不.
照著耶穌基督.
如果你跟了他.
你就被他俘虜了.
你不是跟耶穌 是跟了他.
聖經教我們要把這些東西.
拆掉.
這些和我們敵對的.
我們要將它釘死在十字架上.
要淘掉.
這些所教的.
是不對的字句 束縛我們.
但是我們信了神.
我們應該凌虐自由.

$^{521}$但這些束縛你的.
你要拆掉.
不要讓人.
在飲食上 規條上.
手指旗上.
去評斷你.
你不要畏畏縮縮的.
以為自己沒有做.
就好像自己做錯事.
不是的 因為你根本不需要死守這些東西.
最重要的就是.
你不要跟著他走.
做一些假裝的謙虛.
去敬拜天使.
這樣你的獎賞.
就拿不到了.
因為神本來要給你的.
第三方面就是言行一致.
首先就是.
夢照去事奉.
你看看.
今天教會裡面.
是很多事奉的崗位.
大家有沒有留意.
最近是不是.
我們開始有事奉人員招募.
你不知道.
有沒有看到.
宗教已經熟了.
神又呼召你了.
你看不看到.
很多明顯的事奉崗位.
你可以去事奉.
你有沒有留意到.
有沒有關心到.
有沒有回我們的祈禱會.
就算直播 你有沒有看.
我們願不願意.
按照神的心意.
我們去事奉.

$^{561}$我們願不願意.
按照神的旨意.
去服侍這個世代.
耶穌說.
我的食物就是進行猜外來者的旨意.
作成他的功.
如果你從來.
都沒有開始.
我希望你從現在開始.
有心不怕遲.
希望你現在.
的起心肝.
當然我們都要培訓.
每次來事奉人員都要培訓.
做主席都要培訓.
每個崗位.
都要去受訓.
你以前沒有做過.
不要緊 從現在開始.
神感動你.
你就起來 去接受培訓.
培訓是需要時間的.
當機會一來的時候.
你就可以接受職責.
當然.
很多時候我們會覺得.
怎樣呢.
輸到用時.
慌行少.
但你知道養兵千日.
只用了一時.
是不是.
養了很久.
所以要浸.
你現在不做.
我們像我們年紀大怎麼辦.
趁現在我們有氣力的時候.
我們要好好裝備.
隨時叫到你的時候.
你就可以去接受職責.

$^{601}$是很重要的.
蒙召的人多.
損上的人少.
但我們不知道誰被損上.
想想當年和你一起洗禮的人.
二十年後.
他在哪裡呢.
想一想.
你的好朋友.
但我們要記得.
我們將來都要面對工作的審判.
我們作為信徒.
我們要向神交仗.
所以這個很重要.
第二個就是要忠心.
在亞馬遜書.
第七章十節裡.
講到一個人.
叫做亞馬遜.
他是伯達利的.
一個祭司.
他就打發人.
向以色列的王.
耶路撒冷.
去控告亞馬遜.
他說亞馬遜.
在煽動以色列人.
去背叛你.
背叛王.
他控告亞馬遜.
然後他對亞馬遜說.
他說你在這個地方已經講夠了.
你不如去猶大那個地方.
去那裡繼續講寓言.
去講道.
那裡都可以.
供應到你.
你可以戶口.
靠那裡生活.
他們會供養你的.

$^{641}$你不用在這個伯達利的地方.
去講寓言.
不用在那裡再講道.
為甚麼.
他說這個地方是宮廷.
是聖殿.
是國家的.
做禮拜的中心.
不需要了.
你都去猶大那裡.
他不准他說.
這番話.
亞馬遜聽到會怎樣.
如果你是亞馬遜.
你聽到會怎樣.
是一個很大的侮辱.
因為他很清楚.
是神揀選他的.
雖然他出身是慕容.
但他知道.
神揀了他.
亞馬遜是一個祭司.
同樣都是在做.
事奉的工作.
但兩個人的態度.
是否差天共地.
今日你的事奉.
是否變質.
你想想你剛剛信主的時候.
是否很火熱.
起初的愛.
那時很熱心.
很願意事奉.
但整天整的時候.
變成老油條.
可能走了一段時間.
走了去世界.
都忘記要回教會.
其實神的工作.
不是熟悉的人.

$^{681}$可以承擔的.
亞馬遜說我原本不是先知.
我都不是先知的門徒.
根本不是你說的那種.
為了錢去講道.
我是牧人.
我是修理桑樹的農夫.
當我在牧羊的時候.
神呼召了我.
他很清楚神呼召了他.
所以他呼召他的時候.
他就沒有退縮.
神叫他去講語言.
他就去講.
當他面對攻擊的時候.
他沒有退縮.
他沒有灰心.
他很清楚這是神的旨意.
他就去了.
神叫他講他就講.
雖然有些人去恐嚇他.
但他不願意.
去恐嚇他.
但他不怕.
你有他.
他都不為之所動.
他總有放膽的.
去講神的話.
其實很多時候魔鬼都會阻止你.
去為神.
去講見證.
為神去講他的道.
魔鬼會阻止你的.
但你會妥協嗎.
還是你會.
衷心的去服事呢.
在埃及.
有一個山洞教會.
名叫默卡達姆洞穴教堂.
在埃及.

$^{721}$這個伊斯蘭教徒.
高達80\%.
以上的回教信仰國家.
竟然擁有一個.
可以容納上萬人的洞穴教會.
他就在開羅.
在埃及開羅的.
東南部的.
默卡達姆洞穴教會.
可以是.
在埃及最大的一個教堂.
每一個禮拜.
有七萬個基督徒.
在那裡讚美.
敬拜神.
這個都是世上最古老的.
一個基督教社區之一.
這個洞穴教會都被稱為.
聖西蒙修道院.
這個地方.
周邊.
住了很多.
以撿垃圾.
為生的人.
所以這個.
山.
都叫做垃圾山.
很特別.
這個教堂.
如果你有機會去的時候.
你會看到.
洞穴裡面都有一些壁畫.
寫著一個馬太福音.
28章6節.
他不在這裡.
他已經復活了的經文.
寫在那裡.
曾經在大概.
十世紀的時候.
伊斯蘭教徒.

$^{761}$統治那個地方.
佔領那個地方.
對一班基督徒是很大的逼迫.
甚至.
要殺光他們.
因為他們不可以容忍他們.
要斬首.
有一天.
這個首領.
他叫做哈利弗.
他有一天.
在圖書館裡面.
忽然間看到一個聖經的經文.
這個聖經的經文.
你很熟悉的.
就是一個介塞種.
這樣的信心.
那個金句.
如果你有一個介塞種的信心.
你就可以移山了.
他看了這個經文之後.
他就去召集.
當時幾個教皇.
他就叫他們來訪談.
跟他們約法三章.
他說如果你.
好像這裡所說的.
你叫這個山移動的話.
我就.
饒你們命.
不殺你們.
這個教皇聽了之後很害怕.
真的沒命.
我怎樣叫他移山.
他就很不信心.
他就去向神呼求.
神就派了一個.
獨眼俠.
一個皮匠.
叫做西蒙.

$^{801}$就來到教皇的身邊.
就跟他一起呼喊.
好像給他信心.
大聲的.
奉耶穌的名.
希望這個山可以移開.
很奇妙地.
他們醫藥.
是在醫療上面.
哈利法就來到.
他們兩個繼續很大聲的.
向神祈禱.
主啊求你移動這個山.
移動這個山.
他們很大聲的.
向神呼求.
當他們讚美敬拜.
在祈禱的時候.
在山上很劇烈地.
在搖晃.
有些人看到.
在山下都可以看到陽光.
根據當時的情形來看的時候.
這個統治者哈利法.
他都害怕.
他看到這樣的情形之下.
他很害怕.
還有他身邊的人都很害怕.
這個山搖動.
他說上帝是偉大的.
你證實了你的信仰是真實的.
於是那幫人就不用死了.
根據埃及的約瑟夫納薩勒.
這個牧師就說.
原來那個山大概移動了.
兩英里.
大概三公里的距離.
在2019年.
埃及的約瑟夫納薩勒.
這個牧師就去台灣.

$^{841}$去鼓勵基督徒.
他就說這個山的例子.
這個牧師就去說.
這座山.
這個山的歷史.
這個牧師就去說.
這座山.
神的應許.
他說你們都可以經歷.
當時是2019年.
他說3月到9月.
你們會經歷一個很巨大的復興.
就好像上帝去拯救他的子民.
使得埃及人.
他們在等待的時候.
在埃及的大山可以移動.
同樣你們都會經歷.
神的應許.
就是這個經文.
他就進一步說.
他說這個大山.
有兩個意思.
一個是實際的山.
一個是人的心思意念.
他說到人的心思意念是什麼呢.
他說台灣的基督徒面對最大的山.
就是你們的心思意念.
每個人都需要.
去正視自己的山.
就是你個人的心思意念.
要去正視這個問題.
他說所有舊有的成規.
和綑綁都要下去.
他說柴火夠了.
祈禱也足夠了.
但是那些在教會.
已經受過訓練的.
成熟的精兵.
你可不可以走出這個社會.
去為主 去傳福音.

$^{881}$做見證呢.
可不可以呢.
從前對於傳福音.
可能缺乏熱誠的基督徒.
在今年.
那時是2019年.
他說今年.
要重新燃點.
復興的火焰.
他呼籲當時在台灣的基督徒.
他說台灣可以作為.
亞洲傳福音的起點.
首要目的就是.
要向廣大的華人世界.
去傳講神的訊息.
然後進而影響.
周邊地域的人.
要把握著.
這一波前所未有的.
復興的浪潮.
叫他們要守住.
這個陣腳.
要去傳福音.
我們弟兄姊妹.
我們又如何呢.
今天.
我們在座的弟兄姊妹.
你願不願意在這個.
灣曲背貌的世代裡.
願意勤讀聖經.
明白真道.
做神無瑕疵的兒女.
今天.
我們的行事為人.
是否和主耶穌基督的福音.
相稱呢.
我們是否願意為.
福音的信仰.
齊心努力呢.
不怕敵人的威嚇呢.

$^{921}$我們一起祈禱.
愛你的主.
愛你的神.
每一個屬你的兒女.
每次我們來到你的面前.
何等的恩惠.
因為是你選照了我們.
我們何等的卑微.
何等的不配.
當我們來到你的面前的時候.
我們看你的說話的時候.
你的話語一解開.
就可以使愚人通達.
我們很多時候不明白的.
但我們需要你給我們悟性.
以致你對我們說話.
時時去啟動我們.
讓我們知道.
哪些是應該要去做的.
哪些是要拿走的.
求你透過聖靈去感動我們.
以致我們一起起來.
讓主興起.
我們要去傳福音.
使多人能夠得到.
你的福音的好處.
求主使用我們.
我們這樣祈禱.
教導下奉主耶穌基督的名教.
\newpage



\section{約翰福音 21:15-22}
\label{sec:dFeC0BqGntE}
\textbf{2023-07-23 加恩樓《跟隨我》約翰福音 21\_15-22 講員:陳靜娟姑娘}
\newline
\newline
連結: \href{https://youtube.com/watch?v=dFeC0BqGntE}{\texttt{https://youtube.com/watch?v=dFeC0BqGntE}} ~~~~ 語音日期: 2023-07-25
\newline
\newline
\hyperref[sec:epe4ZHzeWUY]{\small{< < < PREV SERMON < < <}}
~
\hyperref[sec:index_chronic]{\small{[返順時目]}}
~
\hyperref[sec:index_scriptual]{\small{[返順卷目]}}
~
\hyperref[sec:_FPkb_FFsjA]{\small{> > > NEXT SERMON > > >}}
\newline
\newline
約翰福音 21:15-22
\newline
\begin{longtable}{cl}
\hline
\hline
章節 & 經文 (和合本修訂版)\\
\hline
21:15 & \begin{tabularx}{0.7\textwidth}{X} 他們吃完了早飯，耶穌對西門‧彼得說：「約翰的兒子西門，你愛我比這些更深嗎？」彼得對他說：「主啊，是的，你知道我愛你。」耶穌說：「你餵養我的小羊。」 \end{tabularx} \\ \\ \relax
21:16 & \begin{tabularx}{0.7\textwidth}{X} 耶穌第二次又對他說：「約翰的兒子西門，你愛我嗎？」彼得對他說：「主啊，是的，你知道我愛你。」耶穌說：「你牧養我的羊。」 \end{tabularx} \\ \\ \relax
21:17 & \begin{tabularx}{0.7\textwidth}{X} 耶穌第三次對他說：「約翰的兒子西門，你愛我嗎？」彼得因為耶穌第三次對他說「你愛我嗎」，就憂愁，對耶穌說：「主啊，你無所不知，你知道我愛你。」耶穌說：「你餵養我的羊。 \end{tabularx} \\ \\ \relax
21:18 & \begin{tabularx}{0.7\textwidth}{X} 我實實在在地告訴你，你年輕的時候，自己束上帶子，隨意往來；但年老的時候，你要伸出手來，別人要把你束上，帶你到不願意去的地方。」 \end{tabularx} \\ \\ \relax
21:19 & \begin{tabularx}{0.7\textwidth}{X} 耶穌說這話，是指彼得會怎樣死來榮耀神。說了這話，耶穌對他說：「你跟從我吧！」 \end{tabularx} \\ \\ \relax
21:20 & \begin{tabularx}{0.7\textwidth}{X} 彼得轉過身來，看見耶穌所愛的那門徒跟著，就是在晚餐時靠著耶穌胸膛說「主啊，出賣你的是誰」的那門徒。 \end{tabularx} \\ \\ \relax
21:21 & \begin{tabularx}{0.7\textwidth}{X} 彼得看見他，就問耶穌：「主啊，這個人怎樣呢？」 \end{tabularx} \\ \\ \relax
21:22 & \begin{tabularx}{0.7\textwidth}{X} 耶穌對他說：「假如我要他等到我來的時候還在，跟你有甚麼關係呢？你跟從我吧！」 \end{tabularx} \\ \\
[1ex]
\hline
\hline
\end{longtable}
$^{1}$今早跟大家分享的講題就是.
「跟隨我」.
「Follow me」.
很熟悉的這三個字.
記載在約翰福音的21章15到22節.
如果大家有聖經在手.
看看約翰福音的21章的第一到14節是說什麼呢?.
那個場景是一個打魚的場景.
那個場景就好像在路加福音第五章記載.
當耶穌呼召彼得,雅各,約翰跟隨他們的時候.
是一個場景一樣的.
他們打了一晚通宵雨都打不到魚.
接著耶穌出現了.
在那個時候出現.
告訴他們可以打魚到哪裡.
結果他們打了很多魚.
當時路加福音五章記載.
「你要來跟隨我」.
這個是那個場景.
那個是剛剛耶穌出道.
彼得開始跟隨耶穌的片段.
但是現在來到今天的約翰福音21章.
已經是耶穌死在十字架上.
埋葬,復活,向門徒顯現了兩次之後.
這次是第三次顯現.
第三次顯現的場景就是彼得,約翰.
他們這一群門徒又再次去打魚.
又是打了一晚雨都打不到魚.
接著耶穌出現告訴他們.
「你殺網在哪一邊?在右邊你下去」.
一下去他們打了153條魚.
這是聖經記載的.
好了,雨打完了.
接著他們發覺在岸上有什麼呢?.
有炭火,有炭燒魚,有早餐.
那個場景是很熟悉的.
但一會兒會再重提.
為什麼說炭火這麼熟悉呢?.
有一個重新出現的場景.
在炭火,燒魚.

$^{41}$門徒已經發現原來是復活之後的耶穌.
開始認得他了.
接著大家就坐下來.
耶穌叫他們吃飯.
他們就吃飯.
來到我們剛剛讀的第15節.
那裡就說他們吃完早飯.
就是吃完這一餐飯.
你要記得那個場景.
當時不只是耶穌和西門彼得在場.
旁邊還有其他門徒都在場.
在這個時候.
鏡頭好像只是對焦了耶穌和彼得的對話.
但其實當時有其他門徒在場聽到他們的對話.
和大家看的這段經文.
可以分為四個部分.
關於愛主,關於服事主,榮耀主和跟隨主.
我們先看愛主.
第15到17節.
有請大家一起讀出來.
1,2,3.
第16節.
第17節.
這三節裡有句說話大家都看到.
重複的就是.
你愛我嗎?.
第15節.
你愛我比這死更深嗎?.
第16節.
你愛我嗎?.
第17節.
你愛我嗎?.
有些人.
有很多學者都會研究那個「愛」字.
希臘文說的是Agape.
或者Failure.
說的是一個真愛.
說的是愛的層次的高或低.
《孤物論》說的是愛的層次的高或低.
也有學者後來說.

$^{81}$其實這個也是冠上約翰.
寫這本《約翰福音》的作者.
約翰喜歡的一個寫作手法.
很喜歡輪流交替.
愛,真的愛.
這兩個在整個《約翰福音》中都中環地出現.
所以無論如何.
這個就是說你是否真的愛我.
你真的愛我比這些更深嗎?.
《聖經》那句.
你真的愛我比這些更深嗎?.
《易造》說.
Do you love me more than these?.
這些是重數.
但是這些是什麼就不知道.
所以有很多學者就在想.
這些「這些」指什麼呢?.
是不是指約翰出來打魚.
他的財產,他的器具.
或者是跟他有關的東西.
可能就是耶穌在提醒他.
你是不是真的愛我.
多過愛周圍的物件.
甚至是這些魚呢?.
無論這些是什麼.
有可能真的像一些比耶穌重要的東西.
可以令人不專注愛耶穌的東西.
對彼得來說.
不知道是什麼.
或者對你來說又是什麼.
但是無論如何.
就是在說神到底在人心裡的份量,比重.
夠不夠那些東西更加重呢?.
有一個人物.
他有綜合名字.
就叫做李愛銳.
或者李愛銳牧師.
Eric.
這位Eric先生.
如果大家有一定的資料.

$^{121}$應該會看過《烈火戰車》這部電影.
裡面提及的就是這個人物.
他1902年出生.
當時他1924年的時候.
22歲左右.
他參加了巴黎奧運會.
在那裡原本他應該是跑100米的.
你聽過他的故事就知道.
到要跑100米的之前那天.
他才發現原來安排在星期日要跑100米.
但是因為他是一個基督徒.
星期日是要歸給主的.
是一個安息的日子.
所以他就說他拒絕出賽.
當時是令到整個英國的人民.
包括王子等等都大衛失望.
因為100米是他的強項.
所以他一直的訓練都是集中在100米.
竟然他拒絕出賽.
拒絕出賽之後.
奧委會再幫他想想.
安排他跑200米和跑400米.
每個人都會覺得.
你一向的訓練.
如果你是做運動的.
你都知道跑100米短跑和200米400米是不同的.
要接受的訓練或者技巧.
所以他可以拿到金牌的機會.
根本是很微的.
但他竟然真的選擇跑.
為的就是他要星期日要崇拜.
其實甚至有些人說.
其實你跑100米是下午跑.
崇拜是上午回來.
是沒有衝突的.
你回完崇拜.
下午就可以去跑步.
但他都拒絕.
因為他認為星期日是全日歸給上帝的安息日.
所以他就不跑了.

$^{161}$結果大眾傳媒很多宣傳.
很多人都很生氣.
但又想不到幾天或一個星期之後.
他200米是跑到拿到銅牌.
到400米的時候是拿到冠軍金牌.
令到所有人大跌眼鏡.
當年只有22歲的一個年輕人.
他在選擇上做了一個選擇.
他知道什麼是更加重要.
你可以說當日耶穌問他.
你愛我比這些更深.
如果這些的這些就指比賽.
可以拿獎的榮譽.
或者是可以令很多人開心的一件事.
他就照樣去跑.
但他的選擇就不是這樣.
也有些學者會研究這一句說話.
你愛我比這些更深.
這個更深的意思可能是指其他人.
是指其他門徒.
還記得彼得曾經三次雞叫耳先.
有三次不認耶穌.
知道有這件事吧.
是發生在什麼時候的.
就是耶穌那一晚跟他們吃完逾越節晚餐.
已經預告給他們聽.
有人會因為他的命而跌倒.
當時彼得站出來說.
就算每個人都會跌倒.
我也不會跌倒.
就算我要坐牢和你一起死.
我也願意.
這個就是當時彼得很豪氣地說.
但耶穌已經跟他們預告一聲.
今晚雞叫耳先.
你要三次不認我.
所以這三次的失敗相當沉重.
就算到了耶穌復活顯現之後.
彼得心裡仍然有一個很大的刺.
就是耶穌還會不會接納我呢.

$^{201}$真的很難看 很羞澀.
當時我這麼豪氣地告訴別人.
每個人都跌倒.
我也不會跌倒.
每個人都失敗.
我也不會失敗.
偏偏他就是三次不認主.
這件事件還要記載在四卷科目書.
可想而知其實所有門徒都知道.
甚至你再看原來有一卷科目書.
記載到當時有一個門徒.
當時耶穌被抓了.
帶了去大祭司的原址.
當時天氣很冷.
所以有些人點了炭火取暖.
圍著炭火取暖的場景出現.
而門徒有一個.
他認識大祭司裡的僕人.
所以幫他傳話.
令彼得可以進入原址裡.
和其他僕人一起取暖.
順便他很想打探.
到底耶穌被抓到原址審查會發生什麼事.
似乎他已經很勇敢.
他跟隨著.
但在取暖的時候.
有人問一句.
你好像是他的門徒.
他就說我不認識他.
接著轉眼第二次又有另外一個人.
那個就是彼得斬了一個兵丁的親戚.
這是記載在聖經裡.
你會發現聖經有很多線索和資料告訴我們.
那個人又是這樣說.
我認得你是耶穌的門徒.
接著彼得又拒絕.
不是不是我不認識他.
第三次他仍然在其他人面前否認不認識耶穌.
這三次拒絕承認耶穌的經歷.
令他跟耶穌有一個很大的距離.

$^{241}$耶穌還會不會接納我呢.
這次他再來海邊看我們打魚.
叫我們下網.
又預備了早飯.
還點了一個炭火.
其實重複當天發生的場景.
在這一刻耶穌再問他.
你愛我比這些門徒更深嗎.
當天你說不會的.
每個人都會跌倒.
到底你愛我比其他門徒更愛我.
更愛得深呢.
耶穌透過這三次問彼得.
你愛我嗎.
其實就是給他三次機會.
公開確認他對耶穌的愛和忠誠.
來抵消他以往三次否認耶穌.
這三次彼得的回應是什麼.
他說主啊.
是的你知道我愛你.
主啊是的你知道我愛你.
主啊你無所不知.
你知道我愛你.
耶穌有沒有曾經這樣問過你.
你愛我嗎.
當時你怎樣回應耶穌呢.
你有沒有很直接的跟耶穌說.
是的我愛你.
你什麼都知道.
你知道我怎樣的.
你知道我怎樣對你.
你知道我怎樣想你.
或者對我們來說.
覺得說我愛你.
對著自己最親密的人.
都覺得未必這麼輕易說出來.
何況是對著他的老師.
他的主.
但是彼得直言說.
我真的愛你.

$^{281}$你知道的.
所以這三次令他可以公開承認對主的愛.
其實是讓他有機會完完全全.
被耶穌基督再次接納他.
耶穌基督真的接納他回轉.
你怎樣知道他接納他.
就是你看下一句.
耶穌說你餵養我的羊.
你牧養我的羊.
你餵養我的羊.
愛主呢.
其實你有沒有想過.
自己怎樣去衡量你對主的愛是有多深.
怎樣去量度呢.
如果你說你真的愛主.
很重要的是.
你要清理你和主之間的障礙.
就好像彼得這樣.
他和耶穌之間的最大障礙.
就是他曾經否認過耶穌.
但是現在的障礙.
是耶穌幫他清除了.
同樣的彼得深深明白到.
一個人需要悔改的重要性.
如果要清除和耶穌之間的障礙.
最重要的是要悔改.
《使徒行傳》二章.
記載彼得已經出來傳道.
建立教會.
他公開的講論的時候.
對著幾千人講.
他其中一個最重要的.
就是你們個人要悔改.
要在耶穌基督的名受洗.
洗了你們的罪.
得赦免.
就會領受所賜的聖靈.
如果你要和神有密切緊密的關係.
一定要無時無刻去反省.
你和耶穌之間有沒有一些障礙.

$^{321}$正如開頭領師.
司果敬拜的領事也提到.
會不會我們有一些陰暗面.
我們有一些不見光的東西.
其實只有主最清楚.
我們都知道一切.
會不會就是這些罪.
難了我們和耶穌可以有一個.
更加親密.
更加直接的愛的關係.
所以一定要清理這些障礙.
你才可以很豪氣地說.
主啊 是啊 我愛你.
你知道我愛你.
一定要悔改.
而且不只是一次.
是每天的祈禱裡.
我們都去反省.
都去回想.
我又有什麼地方不妥善.
我得罪了上帝 得罪了人.
真的要認認真真去悔改.
神應許過.
你只要是真心去悔改.
祂就願意赦免你.
洗淨你一切的不義.
和耶穌有一個很好的關係.
第二個呢.
愛主之後.
就可以去服事主了.
有一個愛的動機.
來令我們可以服事主.
又是重複這三節經文.
這三節經文呢.
十五節.
耶穌是馬上告訴你.
你餵養我的羊.
十六節.
你牧養我的羊.
十七節.

$^{361}$你餵養我的羊.
留意這個是我的羊.
是耶穌基督的羊.
是主的羊.
不是你的羊 我的羊.
那些是我的羊 我餵養牠.
那些不是我的羊 我不餵養牠.
不是的 其實全部都是主的羊.
羊又是指什麼呢.
大家都會明白.
是指一群信主的群體.
裡面的每一個.
你們每一個都是主的羊.
教會裡面每一個.
都是主的羊.
彼得就是承接了.
耶穌給他的一個使命.
就是你要去建立信仰的群體.
就是建立教會.
你的名字就是教會的盤石.
所以你就是要牧養.
餵養屬於主的羊.
無論牠是用小羊.
Lamb.
還是羊 Sheep.
又再用Sheep.
正如有學者說.
他的寫作手法就是這樣.
他用的是阿加皮的真愛.
還是Faileo 愛也好.
又或者用Lamb 或者Sheep.
小羊 羊也好.
都是一路交替地用.
就是要強調那個重要性.
和涵蓋的範圍.
所以無論是大羊還是小羊.
他是年輕的羊.
我們有很多小朋友的羊.
我們有長者的羊.
都是主的羊.

$^{401}$而需要做什麼呢?.
餵養 餵養牧羊.
或者是餵養.
這個是說無論身體 心靈.
一個全人的關懷.
全人的幫助.
牧羊令他們在靈命上成長.
在心智上都要成長.
彼得蒙神赦免他的罪.
清除障礙.
接著重新接受耶穌基督的委派.
你可以去餵養我的羊.
我的小羊.
或者是牧養我的羊.
服事主其實是什麼呢?.
就是接受上帝給你的任務.
在《耳窟所書》五章二十五字提到.
作丈夫的你們要愛自己的妻子.
正如基督愛教會為教會寫記.
留意這句 基督愛教會.
就是耶穌基督主愛教會.
可以愛到為教會連自己的性命都不要.
或者很多人我們都會說.
教會裡面真的很麻煩.
教會人又多.
是非又多.
亂轉的事又多.
也聽過有些人因為教會發生這麼多事.
就很失望.
我還是轉教會.
繼續尋索尋索去哪裡.
但你要留意.
這裡說的是基督愛教會.
愛教會就好像愛媳婦一樣.
他愛教會愛到一個地步.
就算教會有多少不足都好.
基督是為教會死的.
願意為教會死.
基督願意愛教會多於愛自己的生命.
如果連主都這麼愛和珍惜教會.

$^{441}$這麼珍惜每一個信主的羊.
每一個信主的群體的話.
我們是誰呢.
我們可以說我不喜歡這一間.
我不喜歡那間.
我不喜歡這些羊.
我不喜歡那些羊.
那些是主的羊.
沒錯 教會裡面充滿很多軟弱的人.
正因為我們每一個都是軟弱的人.
所以我們真的需要主.
我們才需要教會.
我們需要在一個信仰群體裡面.
去一起學習 一起成長.
一起如何在主的話語上更加紮根.
更加去建立.
一路更加剛強.
在這裡成長的時候.
我可以更加好好地去服侍主.
所以不只是看到某些軟弱.
不妥善.
我會很輕言說.
我不喜歡他怎樣怎樣.
其實不需要說這麼多.
反而最重要的是.
耶穌基督有沒有揀選你呢.
祂有沒有派你去服侍人呢.
這個反而你要好好反省.
如果基督可以愛教會.
去到一個地步.
為教會連自己的命都不要.
而祂又揀選你.
派你去服侍祂.
這個其實是一個非常光榮的任務.
又說回李愛蕊.
《烈火戰車》這套戲就停止了.
她可以跑到400米拿金牌.
拿到最高的榮譽.
全國的人民大衛開心.
興高采烈.

$^{481}$歡迎她回到自己的國家.
但她當時1924年拿完冠軍.
同年她又大學畢業.
竟然她做了一個決定.
又是令所有人跌眼鏡的.
她不是應該好好做好體育家.
或者做一個教育家也好.
或者繼續去參加奧運.
那時才22歲.
四年一次.
即是二十多歲.
26歲可以再來過.
其實可以拿更多的金牌.
更多的獎牌.
為國爭光.
但竟然在第二年她做了一個決定.
就是要去中國宣教.
去天津新學書院.
做一個教師.
又要說一些背景.
原來這位Eric牧師.
他的父母是宣教士.
他小時候是在天津出生.
所以他會說普通話.
他有一個中文名叫李愛睿.
他就是一個這樣的宣教家庭成長.
他對主耶穌是一個很忠誠的追隨者.
所以到他大學畢業之後.
他不理會周圍的人.
甚至大眾傳媒的很多輿論.
說他們這樣就去宣教.
但他做了這個決定.
他說要回去他出生的地方.
要服侍當地.
當時還是比較落後的一個國家.
在那裡做一個教師.
他教什麼科目呢?.
他教數學,科學.
特別是教聖經和教體育.
他就用他自己的才能.

$^{521}$因為他在體育方面有很多經驗.
他甚至可以幫助當時要在天津建一個很大的體育場館.
他還可以給很多意見.
令到當時亞洲一個很有規模.
設備很齊全的體育館.
他在任教,做老師接近20年.
他向很多學生.
不單是傳福音.
不單是教聖經.
他的新教,他的賢教.
都令到很多人很喜歡這位老師.
所以服侍主.
如果上帝用得著你.
或者你被上帝使用.
這是光榮的一件事.
是一個榮幸.
反過來.
如果耶穌說我不用你.
我一定要用你.
你反而這樣你會很害怕.
如果耶穌不用你.
你會更害怕.
耶穌再次使用彼得.
接納他.
要用他.
同樣今天.
耶穌基督已經呼籲你.
呼召你跟隨他.
你可以為他做些什麼呢.
他會在禱告裡面尋求.
尋求耶穌如何去用你.
使用你.
第三點要說的是榮耀主.
十八字.
我實實在在地告訴你.
你年輕的時候自己戳上帶子.
隨意往來.
但年老的時候.
你要伸出手來.
別人要把你戳上.

$^{561}$帶你到不願意去的地方.
耶穌說.
使子彼得會怎樣死.
來榮耀上帝.
說了者說.
耶穌對他說.
你跟從我吧.
榮耀主.
沒有人知道自己會怎樣死.
甚至不是每個門徒都知道.
耶穌會預告他們怎樣死.
譬如門徒雅各.
其實他在史堂傳記.
他是第一個殉道.
但是耶穌有沒有事先跟他預告.
如果你會用死來榮耀他呢.
不知道.
聖經沒有記載.
但是彼得.
當耶穌已經給了他任命之後.
隨即耶穌就已經馬上告訴他.
我很鄭重地告訴你.
你將來.
你會被帶到一個你不願意去的地方.
而且.
是用死來榮耀上帝.
意思是說.
他會一路服侍耶穌.
服侍到死那一天.
正如啟示錄也有記載.
他說你不要致死忠心.
我就賜給你生命的冠冕.
耶穌基督已經預告了給他聽.
你會用死來榮耀他.
又要說回李愛銳牧師.
開頭提到他.
1925年去了天津服侍.
做老師.
其實一路服侍接近20年.
直到1941年.

$^{601}$其實當時他已經結婚.
有太太.
有兩個女兒出世.
有一個女兒還在他太太懷孕當中.
但是當時因為日軍侵華.
所以很多村教士都會撤離.
包括他首先送走了他懷孕的太太和兩個女兒.
回到加拿大.
但他自己就繼續留在.
當時已經轉了.
在河北省的中南部.
繼續去宣教.
去救治傷兵.
接濟難民.
去到1943年.
當時日軍偷襲珍珠港.
所以很多西方人士.
日軍都很不喜歡.
就將他們囚禁在集中營.
而他就是其中一位.
他回不去見他太太.
但卻要在集中營裡.
在那裡足足住了兩年.
直到他病逝.
所以他都見不到他第三個女兒出世.
但這位李牧師.
他在集中營那兩年.
不只是甚麼都不做.
每天都垂頭喪氣.
覺得很慘很慘.
怎樣過日子.
又有一樣特別的地方.
他這個集中營.
又容許他們在裡面.
給小朋友.
因為裡面其實都有小朋友被囚禁.
你想一下西方人士宣教士.
特別是有一位宣教士.
可能你都聽過他的名字.
叫戴德生.

$^{641}$戴德生牧師.
他的曾孫.
戴紹增牧師.
和他的兄弟姊妹.
被困在同一個集中營裡面.
戴紹增牧師的父母.
當時在中國另外一個地方.
比較安全的地方繼續宣教.
當時他們都很想.
可不可以去探望他們的子女.
很掛念他們.
不知道他們怎樣辦.
但後來他們在祈禱裡面.
神跟他們說.
你們安心吧.
你們緊張我.
我一樣會緊張你們所緊張的人.
所以他們的幾個子女.
被困在這個集中營裡面.
之後竟然遇上了李愛睿牧師.
李愛睿牧師特別疼愛.
特別照顧小朋友.
特別疼愛這幾個小朋友.
當時戴紹增牧師.
原來他已經有14歲.
是一個青少年.
在裡面李愛睿牧師.
繼續教他們聖經.
幫他們傳遞很多正面的.
正向的思想.
很多時候告訴他們.
你們要勤力做運動.
因為他自己原本也是體育老師.
要怎樣鍛鍊好自己強壯的身體.
以備有一天你們可以離開集中營的時候.
可以再見到你們的父母.
見到你們的家人.
所以時時都說一些很有希望.
帶給他們希望的說話.
鼓勵了很多很多在集中營裡面的.

$^{681}$無論大人還是小朋友.
當然你會說.
他自己偏偏可能太辛勞.
太過份.
後來在集中營裡面服侍了兩年.
有一晚他離世了.
離世的時候.
集中營有很多人去送行.
也有些人發現.
特別有個醫生幫他做一些醫理檢查.
發現原來他腦裡面有一個很大的腫瘤.
做手術也做不到.
但他自己不知道.
因為他一直都不舒服.
但沒有說話.
無論怎樣也好.
他願意將他的一生.
就算環境是怎樣惡劣.
怎樣看下去好像沒有希望也好.
他都努力的去服侍上帝.
服侍他身邊的人.
去教導他們.
去傳遞一個美好的訊息.
我們每一個人都不知道我們生命的長短.
但是神母說.
我們一定要去到殉道去榮耀祂.
但不要說到.
就算是未死之先的我們現在的人生.
我們能不能去到榮耀主.
直至見主面呢?.
講了這句話之後.
耶穌就對彼得說.
你跟從我吧.
你跟從我.
路加福音9:23記載.
耶穌對眾人說.
如果有人要跟從我.
就要寫幾天天背起自己的十字架.
來跟從我.
自己的十字架.

$^{721}$你的十字架.
跟我的十字架.
跟彼得約翰.
跟里愛潤穆斯的十字架.
都不同.
但每個人就要為自己所做的負責任.
這句我們很明白.
但我們同樣要為我們所背的十字架.
對得起我們的主.
要負責任.
天天背起自己的十字架.
每個人有自己的十字架.
你可以說有自己的重擔.
有你自己的生活困難.
有你自己生活的挑戰.
有你的背景.
有你的強項.
有你的弱項.
但無論環境如何都好.
每一天都會有很多挑戰.
就背好那一天的十字架.
為那一天做好自己.
可以隨時隨刻.
我們可以面對上帝.
認見上帝.
天天背起自己的十字架.
來跟隨耶穌.
這是耶穌對我們很鄭重的呼籲.
你要有一個捨己的心理準備.
最後.
耶穌重覆了兩次.
跟約翰和彼得說.
你要跟隨我.
二十字.
當耶穌說完你要跟隨我.
第一次說完.
彼得立即轉身來.
就看一看.
看一看耶穌所愛的那門徒.
很多學者都認為.

$^{761}$這就是指約翰.
寫約翰福音的作者.
因為他最清楚.
彼得這個時刻.
轉個頭.
看一看那個約翰.
接著他就說.
主啊.
這個人將來會怎樣呢.
其實意思是.
我就要以死來去榮耀你.
那個人又怎樣呢.
耶穌就跟他說.
如果我真的要他等到我來的時候.
我會再回來的時候.
他還沒有死的話.
跟你有什麼關係呢.
你跟從我吧.
的確.
後來我們看到.
彼得是真的為主殉道.
但是約翰是服侍主.
直到他年老.
可能是自然的死亡.
彼得被耶穌呼召.
去做一個建立和牧養教會.
做一個牧人.
一個牧者.
一個傳講基督的人.
做一個見證人.
甚至最後.
可能是受折磨.
痛苦而死.
但是耶穌也可以呼籲約翰.
做一個守望者.
他寫約翰福音.
他寫啟示錄.
他看到很多將來.
耶穌再回來的異象.
他可以是一個作者.

$^{801}$是一個寫書的人.
他活到年老.
自然去世.
每一個人有不同的恩賜.
有不同的時候.
被上帝揀選.
使用.
但是其實我們全部人.
都是同樣.
是蒙愛.
被重視.
可以被神大大使用.
不論你什麼背景.
什麼的性情.
你對事情.
你對教會.
有什麼不同的想法.
你做事的方式是怎樣.
你的想法是怎樣.
上帝可以隨他的心意.
來去用每一個人.
用你和用我.
但是我們有同一樣的感恩.
就是我們同樣都是蒙上帝揀選的人.
《約翰福音》21章26節.
這樣說.
若有人服侍我.
就當跟從我.
我在哪裡.
服侍我的人.
也要在哪裡.
若有人服侍我.
我忽必尊重他.
所以強調的就是.
跟隨耶穌.
你跟隨他.
他去哪裡.
你就去哪裡.
我的主吩咐我要在這裡.
我就在這裡.

$^{841}$我的主吩咐我要去別的地方.
我就要去別的地方.
主吩咐你怎樣跟隨他.
你就怎樣跟隨他.
他吩咐你怎樣做.
你就怎樣做.
你服侍他的時候.
天父尊重你.
尊重你.
有時候我們會被很多周圍的環境.
人的聲音,說話影響.
好像李愛潤牧師.
他可以繼續跑他的一百米.
很有把握地拿他的金牌.
但不是的.
他不理會周圍的人怎樣說話.
怎樣大眾傳媒的輿論.
他就是選擇我要翻崇拜.
我要敬拜主.
他就是去跑.
選擇一個不是他的強項.
二百米,四百米.
他靠著上帝.
他一樣可以去榮耀主.
去拿金牌.
他大學畢業.
他可以繼續在很多人的歡呼聲之下.
繼續做好他一個體育家.
參加比賽.
為國爭光.
但在這個時候.
他又不理周圍人的聲音.
他只專心聽到一把聲音.
就是主耶穌對他的吩咐.
他要服侍他.
耶穌說要他去天津.
他就去天津.
離開自己的國家.
甚至.
他有機會和他的妻兒離開.

$^{881}$當時已經是打仗打到興嘻嘻的.
二次大戰裡.
他可以和他們一起回加拿大.
但他又選擇了留下.
當然他也想不到.
後來被關進集中營.
但他進入集中營.
他沒有放棄.
他沒有怨天罵人.
亦有些記載.
他曾經一度經歷過靈性的低潮.
他心裡也有疑惑.
為何我會去到這樣的田地呢?.
他也有掛念他妻兒的時間.
他寫了很多書信.
亦有很多人將他的書信記錄下來.
但他就做了這樣的選擇.
繼續在集中營裡幫助其他人.
鼓勵其他人繼續傳福音.
他怎會知道他所鼓勵的.
例如大紹生牧師.
他們後來也繼續服侍上帝.
他們成立了中國內地會.
亦即是現在的OMF海外基督使團.
現在這個OMF在全世界有幾千個宣教士.
你怎會想到一個十多歲的小男孩.
將來可以發展到開一個警察會.
派這麼多宣教士出去呢?.
不小看每一個小羊.
我們的小朋友.
由幾歲到十多歲到二十多歲到三十多歲.
每一個年紀都可以服侍主.
每一個環境你都可以服侍主.
我們就善用現在我們有的機會和自由.
現在仍然是可以自由傳福音的.
要把握這些機會.
真的要好好反省.
如果上帝沒有使用你.
你真的要好好反省.
主啊!為何你不使用我?.

$^{921}$是否我不行呢?.
被主使用是一個榮耀.
若有人服侍我.
我必必尊重他.
愛主是我們的一個動機.
令我們有動力去服侍.
我們服侍的對準焦點.
是事奉我們的上帝.
服侍他的羊.
我們的榮耀目的是榮耀主.
不是把榮耀歸給自己.
不是榮耀歸給教會.
不是榮耀歸給機構.
不是榮耀歸給任何人.
而是真是榮耀主自己.
最後的跟隨主.
跟隨主是你和主一個直接.
很個人的關係.
好好地跟隨主.
直至見主的面.
我們一起祈禱.
因為你愛教會.
甚至為教會寫命.
你愛我們.
甚至為我們寫命.
我們可以為你做些什麼呢.
主呀.
求你保守我們一個單純的信心.
和專一的心.
就是一條心去跟隨你.
專心去跟隨你.
專心去聽到你對我們呼召的聲音.
聽到你要我們用什麼方式.
什麼途徑.
或者什麼場合去服侍你.
去事奉你.
無論環境際遇如何.
那些都不會難了我們.
可以去到主的跟前.
不會難了我們和你的關係.

$^{961}$正如彼得三次否認你.
主呀你都親自要去扶持他.
將他三次的挽回.
給他機會.
讓他明白.
主耶穌你不是戒懷.
不是再介意他之前做了什麼.
而是他現在要重新在你面前立志.
要跟隨你.
要愛你的時候.
主耶穌就給他的任務.
就是要餵養.
去牧養.
主耶穌我們今天將嘉欣樓.
每一個弟兄姊妹.
我們的童宮.
我們的大人.
小朋友.
我們都交到你面前.
主耶穌你要使用每一個.
你要呼召每一個成為.
我們所合用的器皿.
以致到我們真是彼此相愛.
亦都互相去鼓勵提醒.
因為愛的是愛主你自己.
照顧牧養的.
牧養你的羊.
屬於你的羊.
主耶穌求你.
讓我們在你的說話上.
更多的學習.
更用心的去明白主的話語.
以致到我們更加掌握得清楚.
主你的心意是如何.
我們又能夠好地將你的.
用你的說話去教導其他人.
去牧養去餵養其他人.
主耶穌我們如果有軟弱.
失望的時候.
主耶穌求你提醒我們.

$^{1001}$就好像那位李愛潤牧師.
在這麼惡劣的環境.
但他沒有放棄.
他仍然鼓勵其他人.
何況我們今天.
可以在這個地方.
不單止是享用一個.
舒適的環境.
來親近你敬拜你.
而是我們亦都要走出去.
向我們的親朋戚友.
我們可以接觸.
或甚至接觸不到的人.
我們如何去為他們祈禱.
將福音帶給他們.
鼓勵他們.
主因為你是活著的主.
你是又真又活的主.
今天與我們同在同行.
主幫助我們.
專心去跟隨你.
祈禱交託奉耶穌基督的性命而求.
阿們.
謝謝大家.
\newpage



\section{約翰福音 17:6-19}
\label{sec:_FPkb_FFsjA}
\textbf{2023-07-23 早堂《寶貴的禱告》約翰福音 17\_6-19 講員:曾敏姑娘}
\newline
\newline
連結: \href{https://youtube.com/watch?v=_FPkb_FFsjA}{\texttt{https://youtube.com/watch?v=\_FPkb\_FFsjA}} ~~~~ 語音日期: 2023-07-25
\newline
\newline
\hyperref[sec:dFeC0BqGntE]{\small{< < < PREV SERMON < < <}}
~
\hyperref[sec:index_chronic]{\small{[返順時目]}}
~
\hyperref[sec:index_scriptual]{\small{[返順卷目]}}
~
\hyperref[sec:iuYWy9ExSgQ]{\small{> > > NEXT SERMON > > >}}
\newline
\newline
約翰福音 17:6-19
\newline
\begin{longtable}{cl}
\hline
\hline
章節 & 經文 (和合本修訂版)\\
\hline
17:6 & \begin{tabularx}{0.7\textwidth}{X} 「你從世上賜給我的人，我已把你的名顯明給他們。他們本是你的，你把他們賜給我，他們也遵守了你的道。 \end{tabularx} \\ \\ \relax
17:7 & \begin{tabularx}{0.7\textwidth}{X} 現在他們知道，你所賜給我的一切都是從你那裡來的； \end{tabularx} \\ \\ \relax
17:8 & \begin{tabularx}{0.7\textwidth}{X} 因為你所賜給我的話，我已經賜給他們，他們也領受了，又確實知道，我是從你出來的，並且信你差了我來。 \end{tabularx} \\ \\ \relax
17:9 & \begin{tabularx}{0.7\textwidth}{X} 我為他們祈求，不為世人祈求，卻為你所賜給我的人祈求，因他們本是你的。 \end{tabularx} \\ \\ \relax
17:10 & \begin{tabularx}{0.7\textwidth}{X} 凡是我的都是你的，你的也是我的，並且我因他們得了榮耀。 \end{tabularx} \\ \\ \relax
17:11 & \begin{tabularx}{0.7\textwidth}{X} 我到你那裡去；我不再留在世上，他們卻在世上。聖父啊，求你因你的名，就是你所賜給我的名，保守他們，使他們像我們一樣合而為一。 \end{tabularx} \\ \\ \relax
17:12 & \begin{tabularx}{0.7\textwidth}{X} 我與他們同在的時候，我奉你的名，就是你所賜給我的名，保守了他們，我也護衛了他們；其中除了那滅亡之子，沒有一個滅亡的，好使經上的話得以應驗。 \end{tabularx} \\ \\ \relax
17:13 & \begin{tabularx}{0.7\textwidth}{X} 現在我到你那裡去，我在世上說這些話，是要他們心裡充滿了我的喜樂。 \end{tabularx} \\ \\ \relax
17:14 & \begin{tabularx}{0.7\textwidth}{X} 我已把你的道賜給他們；世界恨他們，因為他們不屬世界，正如我不屬世界一樣。 \end{tabularx} \\ \\ \relax
17:15 & \begin{tabularx}{0.7\textwidth}{X} 我不求你把他們從世上接走，只求你保全他們，使他們脫離那惡者。 \end{tabularx} \\ \\ \relax
17:16 & \begin{tabularx}{0.7\textwidth}{X} 他們不屬世界，正如我不屬世界一樣。 \end{tabularx} \\ \\ \relax
17:17 & \begin{tabularx}{0.7\textwidth}{X} 求你用真理使他們成聖；你的道就是真理。 \end{tabularx} \\ \\ \relax
17:18 & \begin{tabularx}{0.7\textwidth}{X} 你怎樣差我到世上，我也照樣差他們到世上。 \end{tabularx} \\ \\ \relax
17:19 & \begin{tabularx}{0.7\textwidth}{X} 我為他們的緣故使自己分別為聖，為要使他們也因真理成聖。 \end{tabularx} \\ \\
[1ex]
\hline
\hline
\end{longtable}
$^{1}$弟兄姊妹平安.
大概差不多三年前.
我就去了紅欣堂服侍.
原來眨下眼時間就已經過得這麼快.
但我回來的時候仍然很有那種親切感.
在座當中都有不少是我過去所認識的弟兄姊妹.
見到大家仍然在神的家在敬拜.
是親近上帝心裡面是感到很雀躍.
亦都很開心的.
有我所不熟悉的朋友都在這裡聚會.
盼望上帝的恩典 上願你們同在.
開始的時候很想再一次一同禱告.
要設下你的華語.
成為我們腳前的燈 路上的光.
願意你進入到我們每個人的生命裡面.
對我們的心靈說話.
願意聖靈在我們心中工作.
讓我們的心是深深的被你感動.
我們就去回應你的教導.
奉耶穌基督的名字祈禱.
阿們.
今天這段經文是講耶穌的禱告.
非常豐富 亦都非常寶貴.
我深信可以成為我們每一個很大的提醒和幫助.
如果大家手上都有聖經的時候.
我們不妨翻開.
約翰福音十七章.
因為今天有些位我們都需要回到聖經的經文裡面去看.
寶貴的禱告.
耶穌基督在這裡做了一個很好的示範.
一個很好的示範.
如何去禱告.
祂為誰禱告呢.
是為屬於祂的人.
是為門徒去禱告.
這一段特別是指著為門徒去禱告.
經文的結構非常清晰.
分開幾部分.
禱告的第一部分是為何要禱告.
在這裡耶穌親自講為何要為門徒禱告.

$^{41}$有幾個原因.
第一.
在經文裡面講.
因他們本是你的.
耶穌基督說這些門徒本來是天父的.
他們本來就是屬於神的.
在神將他們賜給耶穌之前.
這班人已經是屬於神的.
這裡講的是一個非常特殊的愛的關係.
沒錯.
《約翰福音》裡面的明句是甚麼.
神愛世人.
神愛世界上每一個人.
祂不會只是愛A君B君.
而不愛其他人.
每個人都愛.
信的不信的都愛.
但對於神來說.
信的人和祂的關係是不一樣的.
是有分別的.
這個愛的關係非常特殊.
很親密.
在裡面屬於上帝的人.
領受從神而來的啟示.
屬於上帝的人會信服神.
會信靠神.
依靠神.
有從神而來的喜樂和平安.
這些對於不信的人.
是沒有這些層面的關係.
正如今天我們講一個家庭.
裡面有父母和子女.
但這個父母是心中常常充滿愛的.
他也很愛他的鄰居.
很愛所有他遇到的人.
他的同事等等.
他每個人都愛的.
但對於父母來說.
子女是不同的.
我相信這一點我們很容易明白.

$^{81}$我們不會說子女和家裡的親人.
其實沒有分別.
跟街上的人一樣.
是有分別的.
那個關係親密很多.
是有血緣的關係在裡面.
裡面很多的互動.
只有家裡的人才有.
而對於上帝來說.
屬祂的人就是有這種的.
神的家裡的人的關係.
是外面的人沒有的.
正正是因為這種關係.
耶穌基督要為這些門徒禱告.
這是第一件事.
第二件事.
這些門徒不單是屬於天父的.
亦都是屬於耶穌基督自己的.
在經文裡面.
大家會發覺很強調一件事.
凡是我的都是你的.
你的亦都是我的.
耶穌說屬於天父的.
屬於耶穌自己的.
都是屬於天父的.
屬於天父的又是屬於耶穌的.
這種的關係.
耶穌和天父的關係就是這麼親密.
父與子的關係.
所以這班門徒同時間.
既是屬於天父的.
亦都是耶穌的.
更加因為是天父將這些門徒.
賜了給耶穌.
這些人和主有很親密的關係.
更加後面我們看到.
四卷呼音書到最後的時候.
看到一件事就是.
耶穌基督走上十字架.
祂犧牲自己生命.

$^{121}$為了全世界的人.
亦都為了祂的門徒.
耶穌的捨己的大愛.
亦都為門徒帶來新的生命.
那個關係非常非常緊密.
在這裡而且不單止說到這件事.
在經文裡面說到.
並且我因他們得了榮耀.
耶穌和門徒的關係親密到什麼呢.
重要到什麼.
是耶穌要在門徒的生命裡面.
在他們的生活裡面得榮耀.
兩者是緊緊結連的.
今日我和你.
我們的生命和耶穌基督是緊緊結連.
既是因為那份救恩.
我們結連在一起.
但亦都因為.
今日我和你的生命要榮耀上帝.
耶穌基督要因為我和你而得榮耀.
所以不能夠分割.
不可以說耶穌就還耶穌.
我們還我們.
不可以的.
今日我們就是這樣.
因為這個緣故.
因為這個這麼親密的這麼重要的緣故.
所以耶穌基督要為他的門徒去禱告.
來到這裡很想問一問大家.
我們怎樣看自己的生命呢.
我們的生命有沒有榮耀主耶穌呢.
主耶穌說要透過我們的生命去榮耀他.
我們的行事為人.
我們的事奉.
我們的一切.
是不是都是榮耀主呢.
還是分割呢.
星期日的生活就是星期日.
一至五就是一至五.
一至五怎樣.

$^{161}$沒有人知道沒有人見到.
最重要是星期日回來好好看看就行了.
是不是這樣呢.
還是基本上我整個人都屬於耶穌基督.
無論在教會不在教會.
我都是榮耀他呢.
這個值得我們想.
原來這個是很重要.
其中一個耶穌要為門徒禱告的原因.
第三樣東西.
為什麼要為門徒禱告.
我到你那裡去.
我不再留在世上.
他們卻在世上.
耶穌基督好知道自己即將離開門徒.
他要踏上受苦之路.
他要被釘在十字架上.
要去受難.
他不能夠在肉身上.
親自的直接的去和門徒同在.
不能夠直接的在這個位置.
用肉身的這個位置去保護和幫助他們.
所以他要為他們禱告.
門徒將要面對很大很大的挑戰.
在這個時候.
他們仍然需要依靠上帝.
不過所依靠的不一樣了.
他們要依靠在他們裡面的聖靈.
去幫助他們.
所以在這裡耶穌要為他們禱告.
求天父去保守他們.
這三個原因就是為什麼耶穌要為他的門徒去禱告.
接下來經文很清晰地說.
耶穌禱告的內容是什麼.
十七章十一到十九節就說整個禱告的內容.
分兩部分.
第一部分是說.
十七章的十一到十六節.
門徒得蒙保護.
門徒蒙保守.

$^{201}$耶穌即將離去.
但是他很愛他的門徒.
所以他要去祈禱.
求天父去保護他們.
這裡非常豐富.
他說:求你因你的名.
就是你所賜給我的名.
保守他們.
使他們像我們一樣合而為一.
第一件事.
他是求天父怎樣保護他們.
是因他的名.
神的名不單是代表了神自己的屬性.
其實神的名也代表了他的能力.
耶穌基督求天父以他的能力去保守這班門徒.
為什麼要保守他們.
使他們合而為一.
就好像天父和耶穌基督合而為一一樣.
很重要.
是要藉著天父的.
是神的名.
神的能力去保護.
而最終是要使門徒合一.
說到這一點是非常非常重要.
所憑藉著的是神的名 神的能力.
神的能力能夠使我們合一.
但同時間我們也是為著神的名.
神自己去合一.
我們為著神的名合一.
上帝是我們合一的原因.
上帝也是我們合一能力的來源.
所以今天弟兄姊妹在座的各位.
我們在這裡合一不是因為我們有什麼筵席.
有什麼好玩的 好吃的.
有什麼活動.
所以我們合一.
也是純粹因為神自己.
我們合一.
以上帝的能力也會保守我們.
合一是非常不容易的.

$^{241}$我們每個人都不同.
我去了紅恩堂事奉之後.
紅恩堂其中一個特色.
他們很喜歡穿一件.
專屬於紅恩堂的T恤.
紫色的T恤.
無論他們是獻屍.
他們有一些大型的活動.
堂慶.
他們就很喜歡全部穿紫色的T恤.
統一的去印製.
看上去好像很一致.
一片紫色的海.
但我想說.
就在這一片紫色的海背後.
大家都非常不同.
整個身形不同.
那件衣服的尺碼已經不同.
大中小 麻加大加小.
穿上身後的效果又不一樣.
不單外表不同.
裡面更加不同.
每個人有不同的成長背景.
有不同的性格.
快慢不一.
對事情的看法不一樣.
所以走在一起.
需要很多的磨合.
要去合一.
是一點都不容易.
我深信事奉的過程.
一定有很多很多的聲音.
正如我知道.
剛剛過去這個星期.
在錦繡堂這邊.
應該有暑期遊樂園.
整個星期.
童宮或是頂梓梅.
很辛苦的.
大家心裡火熱的.

$^{281}$參與了很多裡面的事奉.
我記得我那時候在這裡.
都參與了好幾年.
暑期遊樂園的服侍.
很明白當中那種辛苦.
也有那種快樂.
但其中一樣東西.
就是要學合一個功課.
大家都很不同.
但卻有同一個目標.
就是要搞好整個的活動.
讓一群小朋友過一個.
很豐富的.
很充實的假期.
讓他們有機會認識耶穌基督.
大家的目標很清晰.
為了這個目標.
大家一同努力.
縱然大家有不同的性格.
有不同的處事方式.
都可以合一.
這是非常寶貴的.
完成之後.
大家是很珍惜的.
因為覺得很難得.
很奇妙.
我們這麼不同.
都可以完成一個.
這麼重要的活動.
一個事奉.
對於紅印堂的頂梓梅.
是一樣的.
嘩 大家是這麼不同也好.
可以一起完成那些事奉.
是何等的美好呢.
我記得過去幾年的時間.
我們試過不止一個.
大型的聚會.
例如我們的堂慶.
我們也試過出去.

$^{321}$幫醫院探訪.
有不同的服侍等.
在裡面大家不停地去經歷合一.
不停地去經歷.
那種同心服侍的美好.
這個是很重要.
也是上帝的心意.
上帝想他的門徒合一.
證明他和天父是合一.
這個就是他其中一個.
很重要的禱告內容.
除了合一.
另外很重要.
耶穌基督很想我們.
在他的喜樂裡面有份.
這是第二樣.
現在我到你那裡去.
我在世上說這些話.
是要他們心裡充滿了我的喜樂.
耶穌知道自己將要受苦.
將要被釘在十字架上.
臨離開前.
他仍然很想他的門徒.
有他的喜樂.
你說耶穌的喜樂.
和我們一般說的快樂.
有什麼不同呢.
是不同的.
一直在耶穌的觀念裡面.
那個喜樂是什麼呢.
就是因為信服天父.
因為信服天父.
遵行天父旨意.
而經歷天父的愛.
他的喜樂的意義就是.
信服天父.
遵行天父的旨意.
是經歷天父的愛.
而有的喜樂.
耶穌想他的門徒.

$^{361}$得著這件事.
想門徒都去信服天父的旨意.
去信服神.
都能夠去經歷天父.
經歷天父的愛.
而有這一份的喜樂.
這一份喜樂.
就是世人不能夠奪走的.
是很永恆的.
無論今天外在的環境怎樣.
發生什麼事.
沒有人可以拿走.
這一份的喜樂.
我記得大德生.
在他的生平事蹟裡面.
曾經這樣分享過.
有一次他和一個事奉上的夥伴.
一同經歷一個暴動的時間.
其實那一次他也有被打.
經歷了很多混亂.
他有受傷被打.
但是完了之後.
他不是滿心的苦毒.
記著那些人怎樣傷害他.
還在想為什麼這麼混亂呢.
會不會弄到我也會受傷呢.
我們會經歷這樣的處境呢.
他不是這樣.
而在過程裡面.
他笑得出來.
一種是由心裡面發出來的喜樂.
因為耶穌基督而有的喜樂.
服侍主而有的喜樂.
還可以開玩笑.
為什麼呢.
就是主耶穌的這份喜樂.
信服天父 信服神的旨意.
為了服侍上帝 甘心受苦.
所以大德生的名言是什麼呢.
如果他有多少英鎊呢.

$^{401}$他全數要給中國.
有多少性命也好.
都是全部為中國擺上.
這是他信服天父的旨意.
信服神的帶領.
因為這個信服.
他經歷神的愛.
就有喜樂.
這個可以超越環境.
接下來門徒將會面對很大的困難.
艱難.
會有逼迫生命的危險.
事實上如果大家有機會.
回看耶穌門徒最後的結果.
會知道耶穌的門徒大部分最後是殉道.
在殉道之前都面對很多苦難.
有面對信仰上的逼迫.
在過程中.
這些是避免不了的.
但是你看到耶穌的門徒.
最後是很信服主耶穌的心意.
信服天父的帶領.
信服天父的旨意.
甘心捨己犧牲.
並不覺得很苦澀.
很怨毒.
而是生命很有靈.
這個禱告很重要.
就為他們禱告心裡充滿喜樂.
而不是祈禱說.
保守我們門徒離開苦難.
保守他們不受患難.
只有喜樂.
天天只有順暢.
只有一切光明美好的.
不是這樣.
今天有時候我們在教會裡的禱告.
我們都有個傾向.
很想這樣禱告.
無論為自己為家人.

$^{441}$或者為我們所關心的幫助的肢體也好.
很想救他們離開一切的患難.
離開一切的困難.
人生最好是順利的.
天天都是陽光普照.
這樣就有喜樂了.
其實這裡不是神裡面的觀念.
喜樂不是這樣一回事.
耶穌基督的喜樂是不同的.
耶穌為他的門徒祈禱.
才會有這樣的喜樂.
第三件事很重要.
只求你保存他們.
使他們脫離那惡者.
他們不屬世界.
正如我不屬世界一樣.
在這裡我們看回經文.
第三點.
耶穌為門徒禱告的內容裡.
你會看到耶穌很想強調什麼.
十四節那裡這樣說.
我以把你的道賜給他們.
世界恨他們.
因為他們不屬世界.
正如我不屬世界一樣.
耶穌基督將神的道賜給門徒.
門徒因為神的道.
生命改變.
在光明裡面去生活.
門徒信耶穌之後.
得蒙救贖.
有了救恩.
他們整個生命不再一樣.
不在黑暗裡面生活.
在光明當中生活.
神的道是有這麼大的功效.
神的道可以照亮.
揭穿這個世界的黑暗.
所以世人很恨門徒.
因為這班人在光中生活.

$^{481}$他們生命所承載的道.
可以照亮這個世界黑暗的角落.
人在罪惡當中.
在黑暗當中.
不想光照進來.
這也是約翰福音.
其中一個理念.
一個這樣的內容.
在黑暗中的人.
不喜歡光明.
不喜歡神的道.
不喜歡神.
世界恨他們.
世界恨這些門徒.
因為他們不屬世界.
正如耶穌基督不屬世界一樣.
今天我們在世界生活.
是有一個很大的拉力.
我們的人還在生活.
有一天我們好像和其他人沒什麼分別.
一樣都要吃喝.
起居飲食.
所有東西都沒有改變.
你會發現我們和別人是一樣的.
但是我們有些東西很不一樣.
我們整個價值觀.
整個生命的質地是不一樣的.
有上帝在我們當中.
我們的行事為人.
我們的思想是被上帝掌管的.
所以我們的價值觀和世界有很大分別.
如果我們今天和一些未信主的朋友聊天.
你會發覺大家的價值觀完全一樣.
這樣要值得思考.
要值得思考.
我的生命改變了沒有.
我的生命是否真的和世界的生命不一樣.
耶穌基督的生命和世界的生命不一樣.
絕對不一樣.
記得很多年前.

$^{521}$我有機會和中學的兩個女同學見面.
真的很多年前.
那時候我已經信了耶穌.
但是兩位同學都未信耶穌.
大家聚在一起聊天.
說生活的事情.
去旅行.
不知道為什麼說旅行.
會不會兩個未結婚的異性.
會一起去旅行呢.
男女朋友.
大家一起去旅行了.
當時說到我就說不會的.
我說很大的試探.
很危險.
一男一女未結婚.
一起會發生一些上帝不喜悅的事情.
我記得我兩個同學.
以一個很奇怪的眼光看著我.
兩個互相看對方.
互相笑一笑.
感覺就覺得這個人很古板.
是這樣的想法.
什麼時代.
一男一女不可以去旅行.
很平常.
大家突然間有一個很大的隔膜.
很明顯大家看事情不一樣.
接著就會說到男女關係.
其中一個會說到跟男朋友發生關係等等.
一路說到原來大家對事情的看法是很大的分歧.
一路我又說我的看法.
他又說一下他的看法.
最後大家不能夠彼此有一個.
不會達至一個協調.
或者這個位置絕對不可以合一.
不會一起.
大家談不攏.
因為大家的價值觀實在太不同.
他們堅持他們的.

$^{561}$我又堅持我的.
但不代表大家因為這樣就翻臉.
之後我們不見面不談話.
只不過心知肚明.
大家實在很不同.
我想問在座的弟兄姊妹.
大家是否都是這樣去看見.
自己跟世界不一樣.
除了說這些行事為人.
還有你怎樣去看家庭.
你看家庭跟世界看家庭有沒有不同.
你看子女跟世界看子女有沒有不同.
世界希望你子女得些什麼.
跟你希望你子女得些什麼.
是否有分別.
你對你另一半的期望.
跟世界對你另一半的期望有沒有不同.
如果你發覺你這些位置全部都是一樣的時候.
我在想你跟世界是否很近.
或者我們跟世界.
其實應該這些位置我們跟世界很遠.
不一樣.
因為耶穌基督的緣故.
我們是應該不一樣.
所以耶穌在這裡禱告.
是很重要的一件事.
祂知道祂的門徒仍然在世界生活.
但他們卻不屬於世界.
所以世界很恨惡他們.
還有最重要的一件事.
使他們脫離那惡者.
從一開始創世記的時候開始.
有神工作的地方.
那惡者又來搞破壞.
你記得神創造天地之後.
神創造人類.
神跟人的關係本來很親密.
很甜甜蜜蜜.
上上很緊密在一起.
但接著撒旦就來工作.

$^{601}$引誘人犯罪.
人犯罪之後痛苦死亡進入世界.
跟上帝之間有一個很大的鴻溝.
撒旦就由起初的時候不斷在這裡工作.
破壞神人關係.
直到今天其實說真的.
這個世界都在撒旦的掌管下.
他仍然要去破壞神跟人的關係.
仍然要去逼迫.
要去傷害那些屬於上帝的人.
這個一直是他的詭計.
而耶穌就要為門徒禱告.
祂不求他們離開世界.
而是求神保護他們離開那惡者.
這一點是非常非常重要.
但我們不要輕看這一點.
我們很容易就落在撒旦的圈套裡.
很容易生命就軟弱就跌倒.
所以這裡絕對是一個禱告裡很重要的內容.
剛才那三點是門徒得蒙保守.
他們合一.
他們有主的喜樂.
他們能夠脫離那惡者.
接著在這個禱告的第二部分.
耶穌基督為門徒禱告.
就是為門徒得以乘勝去禱告.
在利未記上神說.
你們要成為聖傑.
因為我是聖傑的.
但原來不只是舊約.
神讓人認識他是聖傑的性情.
這個屬性.
在新約的時候.
上帝仍然沒有改變.
他仍然是聖傑的.
他想屬於他的人都是聖傑的.
所以耶穌要為他的門徒禱告.
要得以成勝.
在十七節這裡這樣說.
求你用真理使他們成勝.

$^{641}$如何令門徒成勝.
是用真理.
是用神的道.
是何等重要.
上帝的說話.
今年紅恩堂的年提.
是跟聖經有關的.
是說「學務真理 遵行主道」.
我們就在紅恩堂推讀經.
每一天就一章.
鼓勵弟兄姊妹讀.
上半年讀舊約.
下半年讀新約.
讓他們認識聖經多一點.
沒有神的話語.
其實我們就沒有了屬靈的養分.
我們很難就這樣說.
「聖長」「聖聖」.
有沒有方向.
如何成為聖傑.
上帝的期望.
祂的要求是甚麼.
我完全都不認識.
那如何成為聖傑呢.
所以在這裡很重要.
就是求你用真理使他們成勝.
用神的道 神的話語.
令門徒成為聖傑.
你的道就是真理.
是何等寶貴.
你如何猜我到世上.
我也照樣猜他們到世上.
天父猜顯耶穌基督.
來到這個世界.
去做主的工作.
而同樣地耶穌基督也照樣猜派門徒.
去到人群當中.
進入到世界當中.
去作主的工作.
為他得著人的生命.

$^{681}$但是要做到這一點.
主說我為他們的緣故.
使自己分別為聖.
耶穌基督自己是聖傑的.
刻意地將自己分別為聖.
甚麼叫分別為聖.
將自己獻給上帝.
就是一個祭物那樣獻給神.
其實分別為聖.
本來的意思就是.
和世界分出來.
你和其他是不同的.
是專屬上帝的.
今天弟兄姊妹我和你.
都是分別為聖.
上帝呼召我們.
教會的定義是甚麼.
上帝呼召一個.
屬他的群體.
一群人出來.
這個群體.
是屬於耶穌基督.
上帝呼召這群人出來.
和這個世界分別出來.
專屬於他自己.
為他自己的緣故聚集在一起.
這個群體叫教會.
所以我和你是要分別出來.
和世界是有分別的.
在這裡耶穌特別將自己分別為聖.
作為一個祭物獻情給神.
為了要使他們.
使門徒也因真理成為聖傑.
作為門徒也要將自己分別為聖.
和世界分別出來.
將自己獻給上帝.
今天說了很多分別出來.
成為聖傑.
弟兄姊妹.
你將自己獻上了嗎.

$^{721}$是怎樣獻上.
和世界不同.
你今天回去就捉住這一句問自己.
我和世界是否一樣.
還是不同呢.
我盼望我們是完全不同.
在世生活卻是與世界不一樣.
你才可以成為那道光.
別人看到你.
哦 不同的.
屬耶穌的人是不一樣的.
盼望大家都可以得到這樣的評價.
耶穌的禱告是要門徒成為聖傑.
今天這個禱告非常豐富.
我覺得絕對可以實踐在我們的生命當中.
無論是為自己的禱告.
或是為我們身邊的人.
大家在家中有兒女.
為兒女禱告.
就可以這樣為他們禱告.
為他們禱告就是求神.
讓他們能夠在整個教會群體裡合而為一.
求神讓他們能夠.
讓兒女們能夠分享主的喜樂.
也能夠分別為聖.
將自己獻上給上帝.
你可以為兒女們禱告.
讓他們得蒙保守 保護.
是很好.
在教會當中服事.
特別是服事年輕人.
都可以這樣為年輕人禱告.
更加需要.
我深信這幾年對於牧羊年輕人方面.
我相信教會有很大挑戰.
這個禱告很寶貴.
為我們的年輕人得蒙神的保守.
他們能夠在主裡面合一.
有喜樂.
離開那惡者.

$^{761}$年輕人能夠成為聖傑.
去禱告.
我們也可以彼此為對方的生命去禱告.
我深信當我們這樣去禱告的時候.
上帝都要更大的去賜福教會.
我們一起祈禱.
藉著耶穌基督的禱告.
讓我們看到你何等的愛我們.
在你眼中.
我們和世人是有很大分別的.
我們專屬於你.
和你的關係是親密的.
天父你很想賜福給我們.
很想保護我們.
所以耶穌基督的禱告裡面.
總是為了門徒的不同方面去祈求.
祈求門徒得蒙保守.
祈求門徒成為聖傑.
主也求你幫助我們眾弟兄姊妹.
能夠成為聖傑.
要和世界分別出來.
主也在世界很努力的去影響我們.
但求你給我們.
求要去堅守你的道.
要你的道天天在我們生命裡面去工作.
改變我們心思意念.
改變我們行事為人.
讓我們活得不一樣.
能夠榮耀上帝你的名.
主求你不單賜福洪恩家.
更加賜福給錦宣家.
讓弟兄姊妹的生命.
都是榮耀.
都是豐富.
是上帝你所喜悅的.
我就是這樣將眾人交上.
奉耶穌基督的名字.
祈禱.
阿們.
\newpage



\section{以斯帖記}
\label{sec:iuYWy9ExSgQ}
\textbf{2023-07-27 周四晚查經班《被擄歸回歷史研讀(21):以斯帖記》 講員:曾裔貴牧師}
\newline
\newline
連結: \href{https://youtube.com/watch?v=iuYWy9ExSgQ}{\texttt{https://youtube.com/watch?v=iuYWy9ExSgQ}} ~~~~ 語音日期: 2023-07-29
\newline
\newline
\hyperref[sec:_FPkb_FFsjA]{\small{< < < PREV SERMON < < <}}
~
\hyperref[sec:index_chronic]{\small{[返順時目]}}
~
\hyperref[sec:index_scriptual]{\small{[返順卷目]}}
~
\hyperref[sec:c_zS1WaX1q0]{\small{> > > NEXT SERMON > > >}}
\newline
\newline
$^{1}$我們同心一起祈禱.
多謝天父,我們再一次在晚上恢復茶經時間.
今晚我們開一卷新的《被老歸回》書卷.
但很特別的地方是這卷《以斯帖記》.
完全沒有提到耶路撒冷的事情.
和我們之前看的《以斯拉記》和《尼希米記》.
有一個明顯的分別.
但又是確確實實是在說被老的時期.
在波斯帝國書山城所發生的事情.
而且更重要的是關乎猶太人整個民族的命運.
所以這卷書其實都是相當重要.
幫助我們如何在一個極大的民族苦難.
看到神的看顧.
和我們能夠看到主是掌管萬有.
掌管著微小事情的主.
叫我們大家都來到去有一個敏銳的心靈.
不是一種我們普通人的疑神疑鬼.
而是敏銳,神同在的心靈.
這個很需要我們一起學習.
求主帶領我們.
也可能有一些弟兄姊妹.
未必是在現場去聽.
透過Zoom來聽到.
但我們盼望成為了一個六濟之後.
可以成為不同的弟兄姊妹.
在不同的地方.
在他們最方便的時候去聽這些訊息.
主的聖靈來工作.
叫我們明白真理.
和有一個實踐的心靈.
奉主耶穌基督的名.
阿們.
很感恩我們開始了晚上的茶經班.
因為上個星期是暑期遊樂園.
所以禮堂不方便設置.
因為音響禮堂會比較方便.
整個彌興禮記完了.
我們很簡單說過.
一到第六章.
就是由他起義要重建耶路撒冷的聖場.

$^{41}$完成一到第六章.
第七章就是.
不是重建聖場之後.
而是準備走聖門的奉獻禮.
即是整個耶路撒冷是屬神的.
但更重要的不是那些磚頭是屬神.
而是耶路撒冷的居民是屬神.
所以第七章很重要的訊息.
特別是一到第五節.
就是利米早就知道人丁稀少.
如何能讓百姓的心.
甘心樂意感到榮譽.
感到榮幸地居住在耶路撒冷.
就是第八章我們看到.
當然第七章之後就是加抱.
第八章第九章甚至延伸到第十章.
就是一個奉興會.
差不多持續了二十多天.
因為剛好是七月初一的崔國節.
一直延展到駐彭節.
駐彭節就是由七月十五.
一直到七月二十一.
如果再有一個節期.
最大一天就是二十二.
所以差不多是整個七月份.
都有這種心靈的準備.
當百姓大家都預備好.
所以這個某程度叫溝通.
其實都不是人和人的溝通.
有時如何說.
如果不是神聖靈工作.
如何說都沒有用.
都不能激動到百姓的心.
但百姓的心能夠遇見神.
很肯定就是因為聖經的真理.
明白和他們樂意遵行.
所以第十一章就是有抽籤儀式.
不是抽居屋 白表 白居二.
而是真的將百姓.
十個支派.

$^{81}$因為利微支派一定住在耶路撒冷.
因為數了加普.
合資格的按加普已經在耶路撒冷居住.
是強制性的.
他們會感覺榮幸可以事奉神.
但其他的支派就十抽一.
十個抽一個住在耶路撒冷.
十一章就是奉獻禮.
好像逐個方位.
不是拜四角.
中國人都有拜四角.
但是為了整個新的重建的聖城.
城牆城門.
這十一章說他們行奉獻禮.
是非常感恩 充滿歡樂.
然後過了一段時間.
這個時間又不知道多長.
再回到上來的時候.
故態復民.
問題又再一次發生.
即是說原來百姓不保持靈敏.
又會很容易受環境.
受周邊的所謂原居民.
就是那些外邦人的影響.
所以他們更加意識到.
神的選民身份.
和神選民以外的外邦人身份.
差不多由那個時候開始.
越來越強烈的分開.
所以去到耶穌的時代.
他們已經很明白.
外邦人那種無份無關.
這個保羅更加上心.
所以去到耶穌的年代.
又說遠了.
去到耶穌的年代.
去到保羅的時候.
他就清清楚楚的.
救恩完成了.
耶穌來到這個世界.

$^{121}$再沒有分猶太人 外邦人.
這個在羅馬書.
在迦拉太書.
和在以弗所書.
是極之清楚地說了這件事.
救恩帶來全世界的人.
只是你有信心還是沒有信心.
你是透過神的救恩.
他的愛子的寶血.
成為神的兒女.
甚至是屬靈的.
亞伯拉罕的子孫.
屬天的子孫.
要麼就沒有神.
沒有救恩.
就只有這樣的分別.
所以就不再歧視.
不再按血統去分猶太人.
非猶太人.
但是現在我們是將時光倒流.
我們回到被擄歸回的時候.
他們需要有這種.
很強烈的分別為勝的意識.
是要分別自己是神的選民.
所以守安息日.
他們看來是很嚴肅的.
所以李顯米真的不准他們.
安息日買賣.
去百佳是罪大惡極.
安息日.
就是要這樣分別.
因為救恩未完成.
耶穌未來這個世界.
新的約未完成之前.
猶太人就有這個很強烈的.
自己的選民身份的意識.
我們講完以色列李顯米記.
就發現整個焦點是環繞耶路撒冷.
環繞重建.
環繞聖文他們面對重建的困難.

$^{161}$面對很多當地的迦南人.
對重建的阻撓.
一路整個經文延展.
以色列開始.
就是古列王的頒佈的詔書.
所有人可以回到自己的原居地.
來重建.
因為上一兩首就是巴比倫帝國的時候的問題.
波斯帝國還有一個懷柔政策.
就是可以回到自己的原居地重建.
只要你拜我們的波斯.
照進貢就可以了.
我們給你有有限度的宗教上的自治.
就帶來不會這麼多國家.
來去很強烈的反抗.
這是某程度上一個很大的得正.
當然古列的子孫.
不是每個人都有這樣的手段.
所以有時會有其他的反叛,抗議等等.
但是無論如何.
如果由以色列記開始第一章.
去到尼希米記.
其實那裡已經是一百二十多年了.
所以你可以想像得到.
由536主前的536年.
去到尼希米重建城牆.
主前的大概是447或448年.
有時是頭尾的.
特別是有時是潤年等等.
你可以想像得到是一百多年了.
但是在這兩本書裡面.
仍然涵蓋他們的心路歷程.
回歸的內心世界.
和回到去之後.
他們緊張什麼呢.
當然是緊張聖約.
即是神選民和神的立約.
即是回到摩西五經的水平.
回到那個標準去討論.
才可以討論的.

$^{201}$所以他們就要帶回去.
所以祭司獻祭.
在聖殿裡面公職.
你的家譜有沒有你的身份.
你沒有身份.
不是不相信你.
不過要看清楚你的家族.
是不是祭司家族.
是不是真是利未人.
即是意味著他們很緊張摩西律法書的教導.
所以他們很看重那個家譜.
所以以色列記有家譜.
尼希米記又有家譜.
你就可以想像得到.
他們就是很緊張.
到底事奉的人合不合資格.
不是他有沒有這樣的能力.
而是他有沒有那個位份.
這個是我們在這兩卷書所見到的.
所以一直都在說耶路撒冷.
一直都在說神的選民.
一直都在說摩西的律法書.
所以他們要按照規矩守節.
所以他們守的駐朋節.
尼希米記有提過.
從來沒有試過有這樣的駐朋節.
不知道他用什麼做比較.
是用摩西的時代.
是用所羅門的時代.
還是他們在外邦的時候.
根本沒有守節.
那種氣氛都沒有了.
還是這次能夠在一個自己的地方.
一班神的選民去守駐朋節.
這些都是轟動的.
是一個很好的里程碑.
因為真真正正整個城牆城門.
是和外邦人已經沒有份的.
已經把那些人趕走了.
只是猶太人在尼希米.

$^{241}$在這個以色列.
在一些祭司尼美人的帶領下.
守的一個屬於自己的駐朋節.
這些真是猶太人的一個里程碑.
正正式式好像和外邦人.
至少關了門.
至少鎖了閘.
是自己的一個敬拜.
我們說了這些引言之後.
突然間出現一卷叫以斯帖記.
又歸納在回歸歷史裡面.
你看那個內文你都知道.
一定是在波斯帝國的書山城.
阿黑徐老王.
阿黑徐老是一個帝號名.
即是一個統稱.
波斯王的其中一個統稱的帝號.
但可能有另外的阿黑徐老.
但無論如何.
看他的行文記載的人.
有些人大部分是說是以色列.
到底他想說的整體的信息是甚麼呢.
讀完十章之後.
你就會發現到.
很重要的就是那個叫普爾節.
不是普里節.
不是雲南的普里節.
是普爾.
即是一個拯救的日子.
奇怪的地方就是.
在摩西的五經裡面.
神規定的節期是沒有普爾節的.
所以在這個時代.
在外邦人來下的書山城的猶太人.
他們自己當然就是以以色列和活底改.
作為主角來設立的一個節期.
所以很特別的.
當然以色列一直就接近被擄歸回.
對不起不是被擄歸回.
接近耶穌降生.

$^{281}$接近他們會分散在不同的地方.
後期也有一些再嚴重一點的.
對猶太人的迫害.
其中我們也有時會提過.
主前的165年或166年.
就有一個叫馬加比的事件.
就是當時有一個敘利亞王.
叫埃提亞古四世.
這個有歷史記載.
他是很討厭猶太人.
他要將猶太人所敬拜的神來羞辱.
所以當時他打贏了埃及.
所以他如日中天的強盛.
就去到耶路撒冷.
去到聖殿.
你想想一個外邦的君王.
一個拜假神的君王.
要在聖殿公然攻擊猶太人的神.
猶太人最忌諱的是甚麼呢.
就是甚麼動物呢.
就是豬.
不吃豬.
提起也覺得是污穢.
其實豬是乾淨的.
我小時候養豬就知道.
很喜歡乾淨的豬.
很喜歡洗澡的.
在聖殿裡面獻上豬.
要來羞辱神明.
所以當時有一些很敬虔的猶太人.
就是瑪加比的家族.
其中一個兒子.
他是祭司.
就叫馬提亞.
沒辦法接受.
就鼓動猶太的男丁.
就像游擊隊一樣來抗爭.
終於打敗了埃提瓦古四世.
把他們趕走.
然後他們就重光.

$^{321}$結成聖殿.
所以耶穌的時代有種光節.
因為沒有燈火.
據聞聖殿的燈一直不熄.
油一直不斷.
所以他們就重光.
所以就叫做重光節.
又叫修殿節.
因為很髒.
被外邦人獻祭.
所以他們要這樣.
所以有一個叫赫魯卡.
又是一個新的節期.
是今天的猶太人都會守的.
耶穌也有守的.
耶穌有守普爾節.
有守到這個重光節.
所以你可以想像到.
在耶穌還沒降生之前的猶太人.
都有幾次滅族的危機.
但是又見到神的保守.
所以以斯帖記在聖經的舊約.
都是幾重要的書卷.
是講到他們這個民族.
怎樣能夠避免滅絕.
滅絕就是說.
神透過這個民族.
出來的彌賽亞就沒有了.
因為沒有大衛之派.
沒有猶大的君王出現.
就是說舊約很多的寓言.
就應驗不了.
所以對於猶太人來說.
我想敬虔的人.
他們會等待那個彌賽亞降臨.
他們就會有心水清.
可能他們沒有想到.
在他們面對現實的環境的時候.
沒有想到這麼全面.
但是我們現在所謂馬後炮.

$^{361}$我們二千多年之後.
二千四百多年之後.
去看回以斯帖記.
我們當然容易說.
但是如果你在當下.
你是以斯帖記時代的人.
你就不會想得這麼遠.
你只會想到現在.
我們怎樣去面對.
這個重大的危機.
我們的神還救不救我們.
這是其中讀以斯帖記.
你可以發現到的.
另外一個更大的問題.
就是有了兩批人.
剛才也說了.
以斯拉記第一章已經說了.
是古列王的元年.
其實就是應驗耶利米蘇.
七十年的刑罰的日子滿.
就會有古列大帝頒佈詔書.
猶太人就回到自己的原居重建.
所以有六萬多人回到重建.
其實是很少的.
因為在七十年的被擄.
在波斯帝國.
當然由巴比倫到馬代.
到波斯帝國.
他們經歷了幾個朝代的更替.
但是他們一直保持著.
自己神選民的身份.
他們繼續在外邦.
一定會買賣.
一定會勤奮的工作.
一定會準備自己.
特別是敬虔的猶太人.
一定會準備自己.
回去重建.
因為神對我們這個民族的刑罰.
那個應許是七十年.

$^{401}$所以一代的人.
如果小時候.
小朋友就被人捉住.
賣豬仔.
賣豬仔被人捕捉.
他們就算七十年後.
都是一個老人家.
八十多歲.
所以第一次重建聖殿的時候.
真的有些老人家在哭.
就是被捕歸回的老人家.
所以你就可以想像得到.
有兩批人.
有一批真的帶著這份信心.
回歸回去重建.
回到上去就變成好像寄居者.
因為你自己的原居地被人搶了.
你已經沒有自己的地方.
所以是很艱難.
很辛苦才能取回自己的地方.
當然是有王的詔書.
即是說准許他們.
根據地契.
他們來取回那些地.
如果沒有王的授權.
當地的人根本不會理你.
也不會怕你.
所以最後第三次的重建.
牽連到尼希米.
甚至乎是猶大的省長.
是由亞特西西王給他的身份.
其實他已經是.
在波斯帝國裡面的一個高官.
在王身邊的一個酒政.
所以有這些政治的身份.
讓他們可以有一個法理的依據.
來重建.
來取回自己的地方.
即是神的選民的地方.
現在出現一個更大的問題.

$^{441}$就是沒有回去.
一直都有各種原因.
沒有回去的猶太人.
還被不被神眷顧呢?.
更大的就是.
更大要想的就是.
他們還是不是神的立約的子民呢?.
這個問題.
如果我們讀完《二師帖記》.
我們就會有充分的答案.
所以這本書是很重要的.
有些小小的背景.
馬丁勞德.
十六世紀的宗教改革.
加德國的一個神父.
後來成為傳道者.
他晚年.
神父是不可以結婚的.
你也知道.
晚年的時候.
他娶了一個太太.
原因.
後來我看了一些書.
我就知道了.
因為既然你自己帶領.
所有人.
宗教改革.
脫離天主教.
很多人跟隨你.
有些神父.
有些修女.
那些修女就變成了老姑婆.
沒有人理會.
沒有人照顧.
那怎麼辦呢?.
馬丁勞德就.
既然你都.
突然間你自己開竅.
聖經原來沒有禁止結婚的.
神職人員不會禁止結婚的.

$^{481}$不是他們天主教所說的.
所以他自己.
衝破了他自己.
你想想一個.
二十多歲的年輕人.
就承諾自己給神.
給天主.
這樣來事奉.
終身的事奉.
真的沒有想過要結婚的.
一心就是為神的福音.
這樣來事奉.
在晚年的時候.
看到那個實際的現況.
你自己就說.
現在我們在神裡面.
我們是自由了.
我們恩信稱義了.
我們現在是伯伯的救恩.
不用靠行為了.
很多人跟著你這樣信.
這樣離開天主教的教廷.
那怎麼辦呢?.
馬丁勞德.
我覺得很欣賞和很感動.
他就以身作則.
真的取一個.
我不知道那個是不是修女.
來作一個示範.
即是說.
連我作為一個修士.
一個神父.
雖然他還沒有去到主教.
他都是這樣.
因為信仰而帶來的生命的改變.
在神的眼中.
原來是聖經的真理.
是教導是可以的話.
他就去做了.
當然.

$^{521}$如果這些課題.
保羅都講很多.
如果我們將來有機會.
在一些專題的討論.
因為查經班就可以討論很多不同.
比較特別的題目.
用時間大家揭多些聖經.
去由聖經去醞釀那個.
即是formation那個真理.
就更加來得結實一點.
concrete一點.
所以這個你可以想像得到.
但我想講馬丁勞德.
就不是講他結婚.
馬丁勞德就說.
祂的祭司帖記就不應該放在聖經裡面.
包括這個亞國書.
因為他覺得整本書.
連一個神字都沒有提到.
是真的.
十章一句神都沒有.
無論你用什麼的稱呼.
主,神,耶和華.
這些主神.
你看看尼希米記是多的是.
每章都有祂向神祈禱.
祈求你紀念我.
怎樣怎樣.
充斥著那種祈禱.
但是以司帖記.
第四章連以司帖說要祈禱.
都是很隱晦的.
都不會說我們向神祈禱.
都沒有提到.
你開上一組.
他們只是禁食.
好像很委婉地沒有提到.
有沒有什麼難言之隱呢.
我們後代的人當然不知道.
但是有一個.

$^{561}$好像很搞笑.
一來都不到馬丁路德說了算.
他是十六世紀的人.
舊約的Canon.
所謂正典已經是在大概主前的.
我不知道他實際的時間.
那個最後的.
所謂的成為舊約的經典.
是幾時的.
因為都是一些編輯.
所以我們只能夠有一兩個重要的資料.
就是耶穌的時代.
就已經肯定有七十事易館.
因為耶穌看的舊約.
就是七十事易館.
什麼叫七十事易館呢.
就是因為去到所謂的.
主前大概三百多.
一直去到二百多的年代.
因為猶太人已經亡國了.
輾轉他們又繼續流逝.
很多猶太人已經不認識希伯來文了.
或者他們說的話就是亞蘭文那些文字.
所以就有必須性.
要來將.
因為那時候大家都知道.
希臘帝國.
雖然很快就因為亞歷山大大帝的死.
突然間三十二歲的駕崩.
就整個國家一分為四.
他的四個將軍拿了那個國家.
那個帝國.
但是希臘的文化.
就沒有因為大帝的離開.
來消失.
反而希臘的文化.
就統治整個歐洲.
整個亞洲.
所以你可以想像得到.
這個是很重要的.

$^{601}$所以有迫切性.
要將希伯來文的舊約.
翻譯做希臘文.
所以Septuagint七十事譯本.
就是舊約的希臘文的翻譯本.
當然是人手.
不會有一些電腦化.
所以這個就是耶穌時代.
猶太人他們讀的舊約.
很多人都不懂希伯來文.
所以就有這個當時流通的世界語言.
就是希臘文.
所以就有這個Septuagint.
另外也很重要的.
就是1948.
發現的死海古卷.
我們2017年去聖地.
我們去過那個地方.
有很多那些洞.
猶大的那些曠野.
很多的洞.
那些山洞.
其中就發現了.
在死海的一些地方.
有一些山洞裡面.
出現了價值連城.
就是無價的.
因為是主前的.
12世紀的首抄本的聖經.
是舊約的聖經.
是希伯來文的聖經.
不是Septuagint.
不是希臘文的翻譯.
是希伯來文的首抄本.
所以就很重要.
真是震驚全世界.
原因是這個死海古卷.
那個首抄本的希伯來文.
和我們今天在現存.
可以在書局買得到的.

$^{641}$舊約的希伯來文的聖經.
是近乎沒有那個錯誤.
你想想我們現在現存的.
舊約的希伯來文.
當然是原文的希伯來文.
不是我們現在中文的舊約.
那個希伯來文的聖經.
是13世紀的馬索拉抄本.
13世紀突然推前.
回到主前的第二世紀.
可能一百六十幾年.
那個時代.
你可以想像得到.
將現存完整的希伯來抄本的舊約.
推上去那麼早期.
是近乎沒有太大的錯誤.
首抄的手誤一定有一點點.
甚至標點符號或數字.
但整個內文等等.
是近乎沒有錯誤.
原封不動地成為了.
我們今天認為最權威的馬索拉抄本的舊約.
所以我們和合本.
都是用馬索拉的抄本.
原文直譯中文.
這個當然就是和合本的形成故事.
就是當時因為西教士去到中國傳道.
迫切需要的就是一個這麼大幅圓的國家.
你只說這樣傳.
固然有這些傳道者.
但更重要的是.
如何讓中國人能夠將來.
下一代他們能夠有大量的人信著.
必須就是將希伯來文和希臘文.
盡快翻譯成中文.
所以他們就一個成立的委員會.
找一些專家.
又請一些漢學的博學齋老師.
即精通中文的這些人.
來幫忙.

$^{681}$如果我要這樣翻譯.
哪些字眼會好一點呢?.
經過十多年.
終於完成了我們今天的和合本.
官話的和合本.
成為了當時的一個很重要.
1919年的白話文運動.
五四運動的前身.
就是這樣而產生的.
所以你可以想像到.
聖經的重要性和真確性.
我想說這件事.
在死海古卷裡面.
居然找不到二司哲記.
這個就是外人質疑.
有沒有神.
在死海古卷裡面找不到.
而塞亞書是完完整整的一卷書.
是找到的.
當然可以很簡單的解釋.
因為那個所謂牧羊人的小朋友.
因為他想找那些羊.
他怕那些山洞有鬼怪.
好像我們小時候.
我家人就經常說.
你不要上去北架山的水塘.
那裡有鬼的.
我們那裡已經傳到整條村子.
小朋友都知道.
北架山上的九龍水塘.
我們叫做猛鬼湖.
所以小朋友就怕了.
不敢上去.
晚上颳風下雨.
我們就整個人蓋上被子.
害怕.
因為真的怕鬼.
小時候.
後來長大了就知道.
原來那些大人嚇我們.

$^{721}$怕我們上去玩.
因為那個水塘很小.
但是有時候他開那個水泵.
就會將那些人.
如果游泳就會扯下去.
一定會死的.
你會游泳都會死的.
所以那些大人就嚇我們.
所以一樣.
那些牧羊的小朋友.
就扔石頭進去.
但是扔爛了那些牙罐.
就是那些塗那些做那些牙.
乾牙那些牙.
就很怕.
就找那些大人來.
終於發現很多那些卷軸.
就拿回去.
不知道是什麼來的.
又不識字.
就拿去市場賣.
賣到多少錢就多少錢.
因為卷得很漂亮.
很精緻的.
於是就被一些.
真的識聖經的考古學家發現.
就不得了.
就是說有些可能已經拿了來透火.
或者有了來不知道包什麼.
賣了.
那就沒有了.
因為當時是.
其實有依據的.
因為耶穌的時代.
前後.
都有一個族群.
在昆蘭這個地方.
他們來引世居住.
他們其中很重要的工作就是.
抄寫聖經.

$^{761}$所以這些文士.
就隱居在那些地方.
有他們自己的曲譜.
有他們自己的手側.
好像修道院那樣.
這些就是死海的這些昆蘭的人.
叫做愛色尼人.
所以.
以斯帖記.
是不是這麼巧.
拿來透火呢.
就不知道.
但是無論如何.
死海古卷就沒有以斯帖記的.
其他全部有.
是沒有以斯帖記.
所以令人質疑.
又全卷十張都沒有提到神的字.
任何一個描述都沒有.
還有就是沒有在死海古卷.
因為死海古卷不代表是一個權威.
只不過在那個地方.
剛好有一大群人.
在那裡隱居.
在那裡抄寫聖經.
以致那裡有大量的手抄本出土.
被發現.
所以就成為了一個很重要的.
近代史裡面的一個考古的重要的發現.
所以現在已經不是隨隨便便.
你可以拿到原裝的手抄本來做.
是一些.
我聽我的老師說.
是一些特別.
真的對這一方面有權威.
才會給他原稿去做研究.
因為太重要了.
這個就是.
我們簡單將以斯帖記.
這些很外圍的趣味.

$^{801}$這些花絮.
給大家先有一個了解.
又要留意.
我們新約.
我們就看不到.
以斯帖記在整個舊約的佈局.
因為新約.
我們叫四分法的舊約聖經.
摩西五經就是五卷書.
接著就是歷史書.
接著就是詩篇.
接著就是那些先知書.
所以叫四分法.
但是猶太人的希伯來文的舊約聖經.
當然猶太人只是有舊約.
因為猶太人就是不接受新約.
新約是外邦人基督徒的書卷.
所以他們只有他們的希伯來文的舊約.
但是舊約.
希伯來文的聖經.
是叫三分法.
就是他們的托拉.
就是摩西五經.
接著他們有前先知書.
後先知書.
接著就是雜類.
就是一般的書信.
一般的書卷.
是這樣分的.
所以以斯帖記.
都是歸入在著作裡面的其中一卷.
反而以斯拉和尼希米是一卷書.
不是兩卷書.
原文的希伯來文是一卷書.
歷代志上下也是一卷書.
所以舊約最後的那卷書是歷代志.
不是馬拉基書.
因為馬拉基書是在另一個類別的.
也沒有所謂十二小先知.
是前先知和後先知.

$^{841}$所謂前先知就是以被擄作為單位.
所以這些都是一些稍為大家要留意的.
因為我們手頭上一信主.
有的就是新舊約傳書的中文和合本.
一定是用了Septuagint的分法.
那典故為何會突然將它分成四份呢.
我又不是太清楚.
至少是這樣的形成情況.
我們現在看這個獨一無二的書卷.
但是歸入了被擄歸回的時期.
內文也是很肯定地說了.
的確是被擄歸回的時期.
是波斯帝國底下.
書山城的首都的一些事情.
那裡橫跨的是大概一年.
即是記載的時間發生的.
我們就可以透過這十章聖經.
看看在這一年裡.
神如何隱藏地拯救祂自己的子民.
或者我們看不到天上的平衡世界.
我們要看看地上的猶太人.
敬虔的猶太人.
如何在一個突然間的危難.
只不過是一個不向哈曼的人跪拜.
就招致滅族.
即是無妄之災.
如何來到他們將祂化解呢.
他們用什麼手段呢.
這些就是似乎聖經是想交代的.
當然成卷書的終點就是普爾節的首.
那個首是傳過所有的.
在波斯帝國底下的127個省的猶太人.
只是猶太人.
規定是猶太人要守的.
所以如果你有機會去到以色列居住.
去到普爾節.
他們會有一個小朋友的小說.
就好像我們灣仔鵝頸橋那樣打小人.
他們就不是打小人.
就是打哈曼.

$^{881}$因為哈曼就是整個普爾節的焦點人物.
所以你可以說世世代代.
由以色列開始.
沒抵改開始.
一直到今天.
猶太人都會守這個普爾節.
當然很正常.
在普爾節他們就會讀《耳私帖記》.
好,先喝口水.
有沒有什麼問題想討論,想問的.
希望將一些引言和大家簡單分享一下.
這卷書絕對不難讀的.
但是我們看看大家信主的年日和經歷.
要從那些細微裡面去.
既然現在是聖經.
無論如何歸納在聖經的經卷裡面.
這個就是真正的信仰.
我們福音派.
我強調我們是福音派.
我們是堅持著這個很重要的信仰.
已經成為了聖經一定有神的保守和默示.
一定是有一些東西.
就算我們今天讀不明白.
是我們不明白.
不是神的聖經有錯.
這些信念希望大家很牢固.
你可以說是很老土.
你可以說很不科學.
正正信仰本身就不是科學解釋.
如果科學解釋是世界的人都相信.
你說耶穌現在再顯現是人都相信.
這個就是信仰的問題.
所以我們不是大包圍.
我們真的相信有神來過這個世界.
祂已經回去預備祂的地方.
我們將來一離開世界.
就進入那個永恆的天家.
等待將來的新天新地.
可能是幾千年後幾萬年後才出現.
但那個已經不是用數字去計算.

$^{921}$因為我們將來就去到靈界.
這個就是福音派的信仰.
所以你想清楚.
如果你不想著這樣相信.
你就去自由派是沒有問題的.
不過到時你自己不根據聖經去信.
其實你沒有信仰.
就等於無神.
因為自由派是不相信有復活.
不相信有天堂.
不相信有死人復活和將來審判.
這是很嚴重的事.
回到一個問題.
其實你信什麼呢.
你都不信有這些事情.
你覺得那些是超自然的難以理解.
那何來信仰呢.
我們的信仰就是相信神來過這個世界.
祂是神又是人的這位耶穌基督.
祂又為人類釘十字架.
就是因為祂的救贖.
我們就可以拯救恩順清二.
我們將來就要面對神.
去到祂的面前.
這些是很簡單的信念.
但一定要很堅定地在我們的心靈裡.
成為我們大家的信仰.
這就是福音派和自由派的分別.
如果沒有問題我們就開始讀經了.
我不知道有沒有新朋友第一次來.
因為上個星期浸禮班有個同學.
他就很想聽.
如果他也出現在這裡.
我們都簡單講一下.
因為可能有些人會聽到.
我們查經班很主要.
我們會盡量將經文去解釋清楚.
和有些背景的.
盡量幫大家按我所認識的去解說.
希望大家能夠透過經卷去.

$^{961}$首先一定要讀完整本書.
至少你能夠簡單地分到段.
在你的腦海裡面.
你能夠大概.
收起書你都能夠講到大概的意思.
你能夠有那個小小的框架.
你自己再研讀的時候.
就能夠慢慢地豐富那個內容.
更重要的是.
有什麼肅靈的教訓提醒.
這就是我們查經班最重要的目標.
所以主要是一個半小時.
真的以讀經verse by verse來去研讀.
當然因為這些舊約的題材.
這些敘事文體.
未必需要逐節逐節地去解釋.
有時候整個plot的佈局.
特別這本書有很多的恰巧.
很多就是說有沒有這麼對.
但正正就是這麼對.
正正就是那些timing.
完全fit到漏油.
那麼對的timing.
一直累積那些劇情的佈局.
就好像我們看懸疑偵探片.
但那個不是作出來的.
是真的有歷史的根據的.
這才是聖經的奇妙之處.
所以敘事文體.
我們未必需要逐節逐節地去解釋.
和新約很不同.
譬如羅馬書那些.
真的有很多的教義.
很多的真理.
還有很多邏輯的一路的推論.
就需要我們去整本書.
需要逐節逐節地去解釋.
和要累積整本書.
上文下理.
前前後後說過的東西.

$^{1001}$再彼此提醒.
和要引述其他的經文.
就變成我們在讀經的時候.
不單純我們只是去背一下金句.
所以你可以說.
如果你不喜歡這樣長篇大論地讀.
有時候會很悶.
但盼望大家慢慢累積對聖經的認識.
建立一份對讀經的興趣.
和在其中找到裡面的生命之道.
第一章的《以斯帖記》.
阿黑隋老作王.
從印度直到古實.
即是非洲地方.
統管127省.
阿黑隋老王在書山城的宮登基.
在位第三年.
即是剛剛繼承.
即是通常是父父爸爸的王位.
在位第三年.
為他一切首領臣僕設擺賢直.
在這一章裡面有三個賢直.
為他一切首領臣僕設擺賢直.
有波斯王馬代的權貴.
因為馬代比較弱.
波斯後來完全吞併了他.
所以有這個時期.
馬代仍然出現.
有波斯王馬代的權貴.
就是各省的貴就與首領在他面前.
他把他榮耀之國的豐富.
和他美好威嚴的尊貴.
給他們看了許多日.
尤其是180日三羅十八.
即是六個月.
你可以想像得到.
這是一個非常強大.
內文已經告訴我們.
是波斯帝國早期輝煌的時候.
後期因為和希臘打仗.

$^{1041}$就一路一路衰弱.
這是很自然的.
每一個國家的興衰都是這樣.
第五節.
今日子門寮.
又為所有住書山城的大小人民.
剛才就宴請那些高官貴就.
好像國慶日.
在灣仔飲宴的都是那些在職的高官.
或者是那些解放軍的將領.
中聯辦的那些官員等等.
或者社會上的一些比較重要的.
可能行政會議成員.
立法會議員諸如此類.
現在不單止這樣.
過了這個筵席之後.
第五節.
今日子門寮.
又為所有住書山城的大小人民.
在御園的院設擺筵席七日.
有白色,綠色,藍色的杖子.
用細麻繩,紫色繩.
從銀環內繫在白玉石柱上.
有金銀的床塔.
擺在紅白黃黑玉石的鋪石地上.
用金器皿刺酒.
器皿各有不同.
御酒甚多.
足顯王的厚意.
一來是厚意.
即是王好像現在太開心了.
國度太強大了.
就來去超視他的太平盛世.
國家的勇武,力量,國度的豐富.
很合理的.
自己內圍的官員.
就一百八十日.
其他的就七日.
這樣來展示他的國力.
他的風厚.

$^{1081}$第八節.
喝酒有例.
不准勉強人.
因王吩咐宮內的一切神載.
讓人各隨己意.
第九節.
開始出現重要的轉折人物.
這個瓦實堤.
不過出現了一幾節就不見了.
因為被王廢了.
本來得到王喜歡的皇后.
就消失.
引進了主角爾斯帖.
皇后瓦實堤.
在阿黑徐老王的宮內.
也為婦女設擺了筵席.
所以是三個筵席.
即是說有男的飲宴.
普通平民百姓.
又有七天的飲宴.
接著皇后又為她內宮的.
可能那些夫人.
不同的官階的夫人.
或者其他的使節.
諸如此類.
都有一個女的.
一個婦女的筵席.
到第七天就有事發生了.
第十節.
第七天阿黑徐老王喝酒.
心中快樂.
就吩咐在他面前.
設立的七個太監.
有名字的.
米胡曼,比斯他,哈波那.
這些名字應該都是波斯名字.
比格他,阿拔他,西達,甲嘉.
請皇后雅實堤.
頭戴著皇后的冠冕.
去到王的面前.

$^{1121}$使各等臣民看她的美貌.
因為她容貌甚美.
王喝大了.
就想太太穿上後冠出來.
展示她的美貌.
如果你是一個很愛美的皇后.
這是一個可以展示的時候.
快些帶上金銀首飾珠寶.
金珠掛在身上.
好像結婚一樣.
那些叫排場.
告訴你嫁得好.
還要多少斤重.
是不是這樣呢?皇后.
第十二節.
皇后雅實堤.
卻不肯遵太監所傳的皇命而來.
所以皇臣發怒.
心如火燒.
我們讀聖經.
讀到這些宮廷的.
好像什麼什麼什麼還傳.
我沒有看我不知道.
這些宮廷鬥爭.
後宮的那些.
不懂得說.
我不懂得說名字.
皇后或者妃嬪.
爭寵和競爭.
為什麼聖經會記載這些東西呢?.
所以讀敘事文體.
又要有一些想像力.
作者用有限的文字.
描繪一些圖畫.
現在的場景就是.
在一個宮廷的飲宴.
其實人盡皆歡.
因為皇帝有皇帝飲.
平民百姓有七天的飲宴.
去到慶高彩利的第七天.

$^{1161}$那些內宮裡面的皇后妃嬪又飲.
大家都感受到.
那種國態升平.
不是打仗.
不是現在兵臨城下.
現在是展示我們這個波斯帝國.
最強大.
民心都是很一致的一個國度.
大家都很開心.
皇帝有一個.
古代的皇帝.
相信一定是.
我們都不用想.
我們都是受著明朝清朝.
這些電視劇洗腦.
皇帝就是自尊.
說一句就是.
不能更改.
什麼都要做.
差不多要摘月亮下來.
都要找到.
可以想像得到.
聖經記載.
其實是描繪一個局面.
描繪一個圖畫.
要我們將我們的思想放進去.
所以有很多空間.
我們可以在其中想到.
作者想帶出什麼呢.
是想告訴我們.
他現在是寫一個宮廷鬥爭的劇本.
還是描繪一件.
關係到猶太人的.
輕衰的事情.
所以每一個細節.
我們要留心的.
因為往後.
往後一直累積緊張.
那個時間點.
就很重要.

$^{1201}$如果時間點不對.
整件事一點都不會發生.
所以第十二節就是.
皇后雅實提.
卻不肯遵太監所傳的皇命而來.
可想而知七個太監回去回報.
皇的那種憤怒.
在眾人面前.
你想想.
怎樣可以下台呢.
你是一個皇帝.
可想而知.
現在我們在港台.
突然叫Yoyo拿咪上來.
Yoyo說我不拿.
聚會中.
不拿.
我就是不拿.
我不上來.
你可以想像到.
全崇拜的人聽著.
Yoyo不肯.
違抗唐主任命令.
她的咪沒有電.
她坐著不動.
大概.
你想想提升多少百倍.
這個時候的雅實提.
為什麼她喝大了呢.
一個皇后居然在大庭講眾.
不聽皇帝叫七個太監.
來邀請你.
要穿最漂亮的.
去展示你的美貌.
本來就應該反應.
作為女士.
說讓我穿漂亮的衣服.
可以飲宴.
現在你可以去到.
香港哪間酒店最高級.

$^{1241}$我不知道.
你去到穿最漂亮的.
有私家車載你去.
你扮到最漂亮.
然後你在眾人面前.
有很多權貴在.
你跟他握手.
手特意拍著.
那顆鑽石很漂亮.
人人都想的.
但這個雅實提.
不知道什麼原因.
她就因為這樣.
而成為了皇帝.
所以皇甚法路.
心如火燒.
所以你可以想像到.
作者現在描繪一個.
這樣的圖畫.
給我們看到.
跟雅實提有什麼關係呢.
第十四節.
那時在王左右常見皇面.
國中坐高位.
你可以想像到.
位高權重的權臣.
都在自己皇帝身邊.
皇帝怎麼可能.
忍得住這口氣呢.
你可以想像到.
那個嚴重性.
雅實提為什麼沒有人提她呢.
這個事件.
又對猶太人有什麼關係呢.
其實二十節記真的很好看.
所以大家不要等我.
逐一跟你說.
你回家今晚一次過.
整個十章讀完.
你就會覺得.

$^{1281}$很比較俗一點.
很過癮.
你會發現神的奇妙.
那個保守.
當然我們講聖經就要逐段解釋.
十四節.
那時在王左右常見皇面.
國中坐高位的尤波斯.
和馬代的七個大臣.
就是剛才提過.
這個格時拿,士達,阿瑪他.
他斯斯,米歷,馬西拿.
米武干.
都是達時無的名節人.
即是說很有那種傳盤的眼光.
為帝國,為王的無論尊嚴.
國家的大事.
來到去上方設法.
來到去思考的.
即是叫做體通全局的.
就發現這件事.
皇帝不是再關乎你個人的.
家庭的糾紛了.
你是不是曾經令太太不開心.
所以太太就公開的羞辱你.
在發女皇的皇后的脾氣.
似乎他們不是這樣看.
作者這樣描述.
即是說他們看到這件事.
背後帶來的漣漪效應就嚴重了.
這個瓦實題就沒有人提它了.
想不到背後的事關重大了.
所以有時太過衝動要小心.
有時說錯一句說話.
回不了頭.
所以瓦實題是一個這樣的舉動.
導致她殺身之禍.
一個女性,一個皇后.
被廢,被打入爛宮.
你就可以想像到那個嚴重性.

$^{1321}$這裡提到.
按皇帝常規辦事.
必須詢問知禮明法的人.
就是這七個人.
即是皇帝會問他們.
皇問他們說.
皇后瓦實題不遵太監所傳的皇命.
照例應當怎樣辦理.
她就安在皇和眾首領面前回答.
可能她就是那個領袖.
好像行政會議召集人那樣.
開始發言.
皇后瓦實題者是不但得罪皇.
並且有害於皇閣省的臣民.
因為皇后者是必定傳到眾婦人的耳中.
說阿剋隋老王吩咐皇后瓦實題.
到皇面前他卻不來.
所以當面在大庭講眾不給皇帝面子.
這件事已經不是皇帝家庭的事了.
是國家的一個大的羞辱.
即是其他127個省的平民婦女.
就有樣學樣了.
因為當時沒有YouTube也沒有直播.
更加不會有WhatsApp.
立即傳到全世界.
人們都知道瓦實題拍了照.
根本就沒有這件事.
飛甲傳書也要騎馬很長時間.
所以瓦實題的確是誇張了.
但這誇張的確是事實.
皇后也這樣做.
一個這麼大庭講眾宴請全國人民的七天的飲宴.
而導致這樣的一個羞辱.
對皇帝當面的羞辱.
帶來其他婦人原來可以不需要給面子自己的丈夫.
可以公開的羞辱.
只要有理有據就行了.
可能瓦實題就認為自己很合理的理據.
我為什麼要拋頭露面.
我為什麼要給你在這裡.

$^{1361}$是因為你喝大了就要我出來.
那我就很沒面子.
突然內心不懂得轉彎的意念.
凌駕了自己的理性.
所以情緒過大.
大家都要小心.
感恩我們到了這個年紀應該沒有什麼情緒.
應該不會你感性先行.
什麼都感性先行.
我們都應該理性.
當然過份理性就變成麻木不仁.
應該感性和理性要平衡.
和看一些事情.
所以不要過度的理性.
更加不要過度的感性.
這裡就告訴我們.
這個事件現在已經提升到.
好像國家的一個示範.
帶來皇帝說.
是喲我沒有想到這麼嚴重.
現在不是我沒面子這麼簡單.
是成為了一個錯誤的榜樣.
來讓其他人有樣學樣這樣就不好.
所以這個情況現在是去到一個不能回頭的.
米武乾就說了.
第十七節.
因為皇后者是必傳到眾婦人的耳中.
說阿恰徐老王這樣吩咐她都不肯來.
他們就藐視自己的丈夫.
這是會出現的.
是一個事實.
雖然和我們今天YouTube的傳播.
WhatsApp的傳播的快速.
我們也可以想像得到.
那個效果那個問題只會帶來的網上公審.
未求證未Fact Check就傳播天下.
這些是過速的傳播.
又是要很小心.
要小心的求證據.
不要隨便轉發.

$^{1401}$這些都是特別經過幾年的社會事件.
更加是我們要吸取這些教訓.
我們不是盲目不仁什麼都不理.
但是我們其實都是要很小心.
那些資訊的傳遞.
因為太快.
有時回不了頭.
以往是人和人需要時間.
需要大家面對面.
很容易就可以得到真正的大家溝通的共識.
但是文字有時是有它的限制.
所以有時盡量開會都不要WhatsApp.
因為有時會說的東西.
未必涵蓋到我想表達的內容.
有時很容易就會誤會.
所以真的不好.
但是現在他們已經意識到.
那個消息的傳遞開去.
加上以訛傳訛.
加鹽加醋.
就更加憤外的嚴重.
所以問題就要在它開始的時候.
就要立即有一個解決.
否則就沒有了這種意境效由.
就會帶來一個後續的惡化.
在處理國事來說是對的.
是應該的.
所以我們就看下去.
下面之後.
第十八節.
今日波斯和馬代的眾夫人聽見皇后這事.
必向皇的大臣照樣行.
從此就必定會大開藐視和憤怒之端.
透過一個這樣的事件.
可能是喝醉酒事件.
就導致到皇后.
很多波斯帝國裡面的婦女.
好像爭取自己的權益.
可能他們沒有權.
現在就變成原來可以爭取.

$^{1441}$原來皇后說一句皇帝都不能出聲.
那就變成真的不好.
出現這個情況.
他們就要立即去制止.
第十九節.
王若以為美.
就降旨.
寫在波斯和馬代人的禮中.
即是說成為一個案例.
永不更改的.
不准瓦實提再到皇的面前.
將他皇后的遺訓賜給他還好的人.
所降的旨意全片通過.
這個是有意識地去傳播.
不是那些加了鹽醋的.
那些閒言閒語的流通的.
那些叫做是非.
這個是一個官方的.
決定了的.
對國家是有利的.
一個正式的宣佈.
兩者是完全不同的.
當然人人都喜歡聽的就是那些是非.
因為加了自己的描述.
傳到另外一個人.
另外一個群體.
完全改觀了.
官方的人不喜歡聽.
也不喜歡看.
這個就是人性.
大家要留意.
所以全面天國之後.
所有的婦人.
無論丈夫貴賤都必尊敬他.
王和眾首領.
都以米武竿這個大臣的說話為美.
王就照這話去行.
他是不是.
因為還是第七日的飲宴.
是不是因為累積了歡樂.

$^{1481}$突然之間被一盤冷水淋洗之後.
皇帝自然有實權.
就有一種憤怒.
也下不了台的面子.
所以一出了金口.
就收不回了.
到了一個這樣的層次.
給你的意見.
要立為一個案例.
要頒佈天下.
現在要執行的就是.
立即廢了瓦實體.
就這麼嚴重.
不是給他一些空間.
譬如罰他什麼什麼.
而是要立即提升.
一個這麼愛的皇后.
因為一次的失誤.
而皇帝當時可能酒精上腦.
都未真正坐定去思考.
就立即聽了米武竿這個大臣的一個意見.
所以你可以想像到.
這樣的處理.
當坐定睡醒就慘了.
就開始後悔了.
但是不能後悔.
因為你已經執行了.
所以整件事背後好像有安排.
如果我們讀完這十章.
我們應該有一個這樣的信心.
就是整本書.
如果不是人為的作品.
而是憑天馬行空來創作.
而是一個事實的描述.
這就是事實的描述.
只能夠有一個結論.
就是真的有安排.
冥冥當中有一個主宰.
在主宰著那件事.
這是以色列的記想我們能讀到的.

$^{1521}$他從來沒有想過後代的人.
有外邦人.
有基督徒會讀這本書.
是他們猶太的書卷.
是他們關乎猶太人的民族命運.
但殊不知我們都是亞伯拉罕的屬靈子孫.
我們都是同神立了這個新約.
耶穌用血所立的新約.
我們是用聖餐作為我們的記號.
我們都是有神和我們立約的子民.
所以到第十章之後.
我們會說到形容.
我們新約的基督徒.
我們怎樣形容我們在神的新約底下.
怎樣展示我們的身份.
和遇到一些難處.
我們怎樣來敏銳神的同在.
所以在苦難裡面.
基督徒和其他不同的.
所謂其他就是其他未信的人.
唯一的不同不是我們刀槍不入.
是我們很敏銳有神同在.
以致我們有力量去面對今生的困難.
這是基督徒和非基督徒的分別.
非基督徒沒有就是沒有.
無論他怎樣膽正命平.
都是沒有真正的平安在他的心靈裡面.
基督徒就要操練.
在平常的日子我們操練認識神.
跟隨神.
我們在困難的日子.
我們就會敏銳神的同在.
他所賜的平安.
和我們內心裡面的盼望.
所以這裡就告訴我們.
王有沒有想錯呢.
就照者說去行.
就發詔書.
用各省的文字各族的方言.
因為有其他外族人他統管的.

$^{1561}$印度去到古實.
包括猶太人當然.
各族的方言通知各省.
使道為丈夫的.
在家中作主.
各說本地的方言.
那個結論就是.
當時男系的社會.
丈夫在家作主.
因為他廢了自己的皇后.
這個瓦實提.
第二章.
第一章有沒有問題想討論一下.
未必在可能大家不習慣這樣的氣氛.
開鏡頭討論.
我也不習慣.
所以我就很想大家善用.
在群組WhatsApp群組.
來有一些你的問題.
或者你的討論.
或者你的分享.
大家一起.
有一個言經的小組.
大家彼此分享.
當然如果你說有些問題.
你又不敢公開問的.
怕尷尬.
或者怕人笑.
不要緊的.
你就可以私下WhatsApp我.
WhatsApp我通常都會馬上回覆的.
我們再看第二章.
這時以後.
黑錘老王的憤怒指釋.
沉澱了.
慢慢的東西完結了.
那些酒精.
那些腎上腺素都慢慢平靜了.
人就才發現自己做錯事了.
似乎都是人性.

$^{1601}$很容易就是這樣.
皇帝都有他這個失誤.
即是說一次事件是嚴重.
但不應該用這麼極端.
去到最高的層次.
就這樣處理.
即是應該是否可以當下.
可能又沒有那個智慧.
去當下的判斷.
想到自己就是那個最大的受害.
最愛的那個皇后廢了.
立了例.
再找回她.
出現的就是這種失落.
黑錘老王就憤怒指釋之後.
他就想念亞實提.
和他的行為.
所行的.
即是說亞實提不只是美貌.
而是他也有智慧.
也有在皇身邊.
對皇都有益處.
有可能提點.
可能有分享.
皇帝也可以分享到.
他不只是貪慕虛榮.
不只是不斷的.
由早到晚都是去喝茶.
去逛商場.
吃好東西.
買珠寶.
最近有什麼新到的.
春季的新的裙子.
他就去買.
就不是.
他真的可以來幫助皇帝.
然後就說.
並怎樣降旨辦他.
於是皇的侍臣對皇說.
不如為皇尋找美貌的處女.

$^{1641}$皇可以派官在國中的各省.
127個省.
招聚美貌的處女.
到書山城的女院.
交給掌管女子的太監希蓋.
希蓋這個人又是重要.
給他們當用的香品.
他們要潔淨自己.
準備好.
可能要符合宮廷的禮儀.
未必是說要怎樣去讓他們.
特別的經過一個月訓練.
香港小姐現在準決賽.
入了圍去哪個地方.
我不知道.
我沒有看.
去哪個地方集訓.
然後教她儀態.
教她化妝.
諸如此類.
這個似乎就不是.
一來都要查她的身份.
她是甚麼人.
因為127個省.
有沒有一些行色.
皇帝的.
很害怕的皇帝.
所以都要查他們的身份.
皇所喜愛的女子.
可以納為皇后.
代替瓦實體.
皇以這事為美.
這句說話出現第二次了.
就如此行.
即是說皇帝.
所以這句都是.
這個《以斯帖記》.
其中一個鑰節.
皇是一個重要的人物.
以這件事為美.

$^{1681}$遲些去到第六章.
他又會有這個描述.
對皇帝來說.
他的賞念.
就讓太監可以進言.
就告訴他.
找個人代替就可以了.
他很聰明.
就不是在皇宮.
後宮再找人代替.
似乎皇帝都.
只喜歡瓦實體的.
其中一個原因.
正正是後宮那些.
他都看不上眼.
就是這個瓦實體.
所以就找一些新的.
是皇帝未見過的.
是極之美貌的.
當然去到皇帝身邊.
其實經過很多的挑選.
其實皇帝沒有機會逐個看.
初選根本不可能給皇帝.
一定是經過篩選.
所以那個欺詐.
如果能夠買通他.
古代就是這樣.
買通這個欺詐.
其實很容易.
因為美不美就沒人知道.
你把美的淘汰了.
剩下的是你買通的.
都可以的.
所以欺詐都扮演著一個重要的角色.
如果以斯帖.
得不到欺詐的歡心.
認為他不美.
就被淘汰了.
就沒有這件事發生.
你又會見到處處都有神管理.

$^{1721}$既然現在是對的時候.
就開始選美.
就看看用多少時間.
應該就不是強搶民女.
現在就貼個黃榜出來.
就說現在皇帝要徵召皇后了.
誰可以入圍呢.
現在有初選了.
就好像去TVB報名.
參加香港小姐.
填表 拍照.
你是什麼背景.
原來你是幾個學位的.
什麼什麼的輩分.
誰的姪女.
類似這些.
就很多那些八卦新聞.
娛樂圈的新聞就會出現.
誰誰什麼背景.
我想古代是一樣的.
他們一樣會有.
誰大臣.
這個是機會.
如果我的姪女入圍.
選中.
做了皇后.
那你可以想而知.
我就是國后了.
你可以想像得到.
大家就會有很多的競爭.
明爭暗鬥.
因為有一個位置.
這個位置就是皇后的位置.
你可以想像得到.
那個競爭性是多麼的大.
127個省.
難道每個那些大官.
自己的女兒.
不要啊 千萬不要選.
一入皇宮就深似海.

$^{1761}$千萬不要.
一定不會.
講者 誰都去.
只要你稍微覺得自己漂亮.
你都去.
無所謂的.
填一填表格.
不用錢的.
萬一抽中.
入選.
你都可以在那裡有一個初選.
至少去見報.
所以是很吸引的.
第五節.
舒生城有一個猶大人.
名叫沒底改.
開始進入主題了.
猶太人有關的故事.
作者很早很早.
一開始就馬上點出.
沒底改是猶太人.
只有他自己知道.
所有人都不知道.
去到第三章.
哈曼就知道.
這個沒底改是猶太人.
但是.
沒底改.
叫以斯帖千萬不可以洩露.
他是猶太人的身份.
有什麼背後的考慮.
這麼有智慧的做法呢.
這些又是作者巧妙的.
透過幾節的經文文字的演繹.
讓我們見到對文學的.
不是修辭.
那個佈局的巧妙.
和優美的地方.
我們要欣賞.
所以讀《敘事文題》.

$^{1801}$都帶著一份讀文學的心態.
去欣賞他的佈局.
那個描述的圖畫.
所以我們的腦海應該有幅圖畫.
現在又如何呢.
他是猶太人.
留意他是.
「辯雅敏茲派基氏的曾孫」.
「辯雅敏茲派」.
「素羅」就是「辯雅敏茲派」.
「基氏的兒子」.
是「侍妹的孫」.
「艾爾的兒子」.
可能是同族的那些蘇伯兄弟.
譬如我姓曾.
我在「薄羅」.
如果有「薄羅」出來的人.
又是「薄羅」的.
又是姓曾的.
咦 可能我們的太公是兄弟都不一定.
是同村的兄弟.
不足為奇.
現在作者很想告訴我們.
他是「辯雅敏茲派基氏的曾孫」.
「侍妹的孫」.
下星期我會引述.
《十五二記》上十五章.
十七章.
另外《出埃及記》十七章.
關乎到亞甲族的事件.
亞馬力人.
這些有關連.
因為哈曼正正就是亞馬力人的後人.
原來兩個一千年後的世仇.
相遇在皇宮.
這件事就好玩了.
整件事原來是一個屬靈戰爭.
原來讀《以斯特記》又會有所謂恩怨情仇.
所以哈曼就說.
我殺你這個活該.

$^{1841}$都便宜了你.
我要將你整個民族連根拔起.
跟著他就和王去求.
王就甚麼都不知道.
給了他這個權.
製造了一個極大的滅族危機.
出現了以斯特的遺訓問題.
四章二十二,二十三節.
「焉知你得了皇后的遺訓.
不是為這個時候嗎?」.
這個時候就是現在滅族的時候.
今晚我們就說到這裡.
我們到達大法九點.
我們一起祈禱.
多謝主.
我們能夠去看《以斯特記》.
幫助我們對這一卷著作.
所描繪的主前大概四百多年的.
波斯帝國的舒山城.
猶太人他們的危機.
他們是怎樣來到去.
同心地去禁食,禱告.
去倚靠神.
又見到很出奇地.
神在暗中的保護.
和神的主權.
掌管萬有的能力.
使整件事轉危為機.
為一個安全.
我們每一個基督徒.
無論在香港或不同的地方.
我們都要面對屬靈的爭戰.
幫助我們明白.
什麼是屬靈的爭戰.
什麼是我們需要去倚靠神.
讓我們能夠從生活.
我們看到的不是一般未信的人.
一個日常的生活.
現在三十七度,三十六度.
而是我們能夠在我們的生活裡.

$^{1881}$看到我們和神的同在.
神的恩典和福音的大能.
在我們的生命裡.
求主繼續透過這一卷書.
讓我們繼續深入的研讀.
我們仰望禱告.
奉耶穌基督的名.
阿們.
多謝您收睇時局新聞,再會!.
\newpage



\section{以弗所書 4:11-16}
\label{sec:c_zS1WaX1q0}
\textbf{2023-07-29 晚堂《建立基督的身體》以弗所書 4\_11-16 講員:盧思源姑娘}
\newline
\newline
連結: \href{https://youtube.com/watch?v=c-zS1WaX1q0}{\texttt{https://youtube.com/watch?v=c-zS1WaX1q0}} ~~~~ 語音日期: 2023-08-08
\newline
\newline
\hyperref[sec:iuYWy9ExSgQ]{\small{< < < PREV SERMON < < <}}
~
\hyperref[sec:index_chronic]{\small{[返順時目]}}
~
\hyperref[sec:index_scriptual]{\small{[返順卷目]}}
~
\hyperref[sec:qI_KwviwgzU]{\small{> > > NEXT SERMON > > >}}
\newline
\newline
以弗所書 4:11-16
\newline
\begin{longtable}{cl}
\hline
\hline
章節 & 經文 (和合本修訂版)\\
\hline
4:11 & \begin{tabularx}{0.7\textwidth}{X} 他所賜的有使徒，有先知，有傳福音的，有牧者和教師， \end{tabularx} \\ \\ \relax
4:12 & \begin{tabularx}{0.7\textwidth}{X} 為要裝備聖徒，做事奉的工作，建立基督的身體， \end{tabularx} \\ \\ \relax
4:13 & \begin{tabularx}{0.7\textwidth}{X} 直等到我們眾人在信仰上同歸於一，認識神的兒子，得以長大成人，達到基督完全長成的身量。 \end{tabularx} \\ \\ \relax
4:14 & \begin{tabularx}{0.7\textwidth}{X} 這樣，我們不再作小孩子，中了人的詭計和欺騙的法術，被一切邪說之風搖動，飄來飄去。 \end{tabularx} \\ \\ \relax
4:15 & \begin{tabularx}{0.7\textwidth}{X} 我們反而要用愛心說誠實話，各方面向著基督長進，連於元首基督， \end{tabularx} \\ \\ \relax
4:16 & \begin{tabularx}{0.7\textwidth}{X} 靠著他全身都連接得緊湊，百節各按各職，照著各體的功用彼此相助，使身體漸漸增長，在愛中建立自己。 \end{tabularx} \\ \\
[1ex]
\hline
\hline
\end{longtable}
$^{1}$願意我們開始之前 我們都有一個祈禱 好嗎.
抽離我們再一次來到你的面前.
一同去聽你的話語.
一息眼孩子要分享你的信息的時候.
求真理的聖靈與我們眾人同在.
預備我們的心 讓我們無論說的.
我們聽的 都是進入你自己的真理當中.
以致被主你自己的話語去得著提醒 得著幫助.
我們恭敬的 都將以下的時間去交上.
求聖靈親自與我們同在.
奉主耶穌基督明智祈求 阿們.
各位姐妹 大家不知道在三個月前.
4月27日晚上 大家有沒有參與.
教會舉行那個心系吾家再躍動的晚宴.
大家有參與的 可能可以揮揮手.
都看到有些人有參與.
在那天 我們在酒樓擺了超過50位的酒席.
在當中 教會不同的群體.
我們有外語的群體 有兒童 有少年 有金靈.
即是青年 不同的目區 我們都有參與.
也都在當中 有一些帶著自己的家人朋友去參加.
而在最後的環節 大家如果有參與.
不知道還記不記得 做了什麼.
當時我們在台上有一張寫著.
心系吾家再躍動的一個Banner.
當時最後的環節 就是將「在」字貼上一個字.
大家有沒有印象是什麼字 有參與的.
齊 聽到嗎.
當最後的時候 我們將「在」字貼上「心系吾家齊躍動」.
當我們這段時間說「心系吾家齊躍動」的時候.
這個不只是單單一個口號.
大家可能在過去幾年的疫情的影響.
坦白 教會都是面對很多不同的挑戰.
甚或可能因著環境的限制 有很多的事工都要停滯.
而現在這段時間 我們看到疫情已經慢慢開始緩和的時候.
我們作為錦宣家的一份子.
實在很需要我們每一個再次一起去躍動起來.
為著神的家 一同去擺上.
而今天分享的經文當中 都有很多的提醒.
在當中 神的話語教導我們如何建立教會.

$^{41}$建立基督的身體 建立神的家.
首先問問大家的問題 大家覺得教會是什麼呢.
在二弗所書 即是我們今天讀的這本書裡.
保羅在二弗所書一章二十三節裡說.
教會是基督的身體.
很多時候我們都說我們回教會了.
很容易我們會覺得教會就是這個地方.
就是這個建築物 這個地方就是教會.
沒錯 這個地方事實上是一個教會.
但聖經裡面告訴我們 教會不是單單說這個地方.
教會不是一件硬件 不是一個死物.
教會是你和我.
是你和我 我們一起去組成教會.
是有生命的.
所以聖經說教會是基督的身體的時候.
今天你和我就是基督的身體的一部分.
當我們去揭開今天我們看的經文的時候.
四章十一至十二節當中提到.
他所賜的有使徒 有先知 有傳福音的 有牧者和教師.
為要裝備聖徒做事奉的工作 建立基督的身體.
聖經在這裡提到 神賜給教會有不同的恩賜.
賜給不同人有不同的恩賜.
恩賜英文是gift 即是禮物.
因此是上帝給我們的禮物.
而恩賜的定義就是.
只由聖靈隨他自己的意思.
給每一個信徒有一個特殊的才能.
而恩賜的目的是要來服侍教會.
如果我們看羅馬書十二章 或者哥林多前書十二章.
當中也提到很多很多不同的屬靈恩賜.
想問問 今天在座的每一位.
大家知不知道自己有什麼恩賜.
記得之前我在讀神學的時候.
神學院的老師教屬靈恩賜的時候.
他同樣問過我們這個問題.
我當時是不懂得回答的.
因為覺得自己好像沒有什麼特別的恩賜.
我好像看不到有什麼.
有什麼恩賜在自己身上.
但是當時老師教我們說.

$^{81}$我們每一個信耶穌的人.
我們最少最少也有一個恩賜.
那怎樣去知道自己的恩賜呢.
而我們不清楚自己的恩賜的時候.
原來不是因為我們沒有我們的恩賜.
原來是你還沒有發掘自己的恩賜.
那怎樣去知道呢.
原來要知道我們的恩賜.
就是將我們自己放在不同的崗位上去嘗試.
而且當我們去運用我們的恩賜的時候.
上帝是會加給我們的.
之前聽有一個朋友分享.
他原本是教書的.
他在教會開始服事的時候.
很容易立即想到在主教學做教導的工作.
但是聽他分享的時候.
他慢慢發覺他的恩賜好像不是在教導.
雖然教導對他來說不是很困難.
在一次的機會裡面.
這位的肢體有機會出去傳福音去報道.
完結之後他心裡很雀躍很開心.
他很熱心的去傳福音.
身邊的人都被他感染了.
後來他發現原來在傳福音的過程裡.
他發掘到神給他的那份恩賜.
另外我曾經聽過一位姐妹分享.
她覺得自己平時很平庸.
好像沒有什麼特別的才能.
好像沒有什麼特別的恩賜.
有一次她有機會在教會裡.
嘗試幫助教會團契.
做一個家中的愛妍的接待.
在過程裡發現她很有服事款待的恩賜.
在那次的愛妍裡.
她打點得很正正有條.
參與的弟兄姊妹很開心.
而且在過程中這位姐妹自己也覺得很開心.
她還說很希望很樂意.
開放她自己的家裡給團友去參加聚會.
因為她發現原來上帝給了她服事的恩賜的時候.

$^{121}$當她運用她的恩賜的時候.
去服事教會.
她感到很開心.
所以原來有時候恩賜未必是說我們的專長.
我們的才剛.
當然上帝是可以使用我們的專長.
使用我們的才剛去做他要做的事.
但我想說的是.
有時候恩賜未必是我們所想的.
而當我們嘗試將自己放在教會不同的崗位上.
我們的恩賜會慢慢被發掘出來.
回到聖經裡.
在《熱忽所書》提到.
有使徒先知傳福音.
有學者教師.
有些人稱這個為五重即是.
其實這個是初代教會的事奉模式.
有些人會討論.
究竟現在還有沒有使徒有沒有先知呢.
如果我們看回《熱忽所書》.
在二章二十節和三章五節裡說.
我們知道使徒和先知是當時建立教會的根基.
但當第一世紀過後.
教會的根基已經立定了.
而且加上當時十二使徒.
連保羅著作的權威都已經可以得到認可之後.
其實現在是沒有使徒這個職份.
不過使徒的功能我相信仍然有.
就好像在一些宣教士.
他帶著使命帶著神給他的意象.
去到不同地方建立教會.
這些都是使徒的功能.
而先知的職份在舊約裡.
我們知道先知就是神的代言人.
很多時候先知去說話之前.
就是耶和華如此說.
但我們今天知道.
上帝要跟我們說的話.
已經完完全全在整本聖經裡.
所以現在沒有先知的職份.

$^{161}$但是我們如果看回新約.
哥林多前書十四章.
我們更加知道現在先知的功能.
不是說我們那些未卜先知的先知.
而是先知講道之能.
就是可以講解上帝的說話.
讓人得到安慰勸勉.
至於在經文裡提到的.
傳福音牧者教師.
我相信大家很容易在教會就知道.
上帝實在有賜予一些人.
有這些職份.
也有這些恩賜.
不過最重要的是.
當我們看到保羅在二弗所書裡說.
神為什麼賜不同的恩賜.
和職份的人在教會當中.
原因是要裝備聖徒做事奉的工作.
建立基督的身體.
聖徒.
大家一聽到會想起什麼.
可能大家馬上想到.
聖徒一定是一些特別聖潔的人.
可能是一些特別敬虔的人.
那些就是聖徒.
但如果我們看聖經.
聖潔的意思.
其實是分別出來的意思.
分別為聖的意思.
他不是說本質上的聖潔.
他說可以分別出來.
給神使用.
今天我們每一個相信主耶穌的人.
我們都是被神在萬民之中分別出來.
所以今天我們在座每一個.
我們都是聖徒.
如果這樣看的時候.
回到經文.
這段的意思就是說.
耶穌基督賜下恩賜給我們.

$^{201}$目的是要裝備我們每一個的信徒.
很叫我們一起去事奉.
去建立基督的身體.
讓基督的身體被建立起來.
大家可以想一下.
如果你的身體只有一半懂得活動的時候.
是怎樣呢.
就好像中風一樣.
一半懂得活動.
很難去活動.
所以教會.
我們應該在我們的崗位上.
在我們的恩賜上.
我們去承擔事奉的工作.
我們一起建立基督的身體.
這個我們看到是聖經裡面所教導我們.
一個很理想的教會的事奉模式.
舉例就好像一個管弦樂團.
一些大型管弦樂團.
都稱為交響樂團.
原來我看一個很大型的交響樂團.
起碼超過70個以上的演奏家.
甚至過百個演奏家去參與這個樂團.
大家不知道有沒有看過這些表演.
在當中每一位的演奏家.
他們都運用自己不同的樂器去奏出音樂.
最後可以組成一首很美妙的樂章.
如果缺少了一個.
整首樂章就變得很單調.
建立教會.
建立基督的身體都是一樣.
都是需要大家一起去努力.
接著如果我們再去看一段經文的時候.
上帝去賜下恩賜要建立基督的身體.
但是這個其實只是一個開始.
還有一些延續性的目的.
當中就是說.
「直等到我們眾人在信仰上同歸於一.
認識神的兒子得以長大成人.
達到基督完全長成的新娘」.

$^{241}$在這裡本來在說教會的建立是一個過程.
是需要時間的.
就好像一個小朋友.
他要長大成人.
是需要時間的.
基督的身體都是需要時間去成長.
但是成長與不同增長.
一個小朋友他可以長高長大.
他今年比去年大一歲.
他可能高了多少.
重了多少.
但是這些只不過是身體上的成長.
但實際內裡可以是沒有成長的.
現代有一個字叫做Kidult.
大家有沒有聽過.
就是Kid加Adult.
就是小孩加成人的合體.
當中就比喻一些心理上長不大的成年人.
就是心思意念還是像小朋友一樣.
雖然他已經是一個成年人.
所以成人與小孩的分別在哪裡呢.
小孩我們知道小朋友可能比較幼嫩.
比較幼稚.
未必很穩定.
他需要大人去餵養照顧.
特別是嬰兒.
他肚子餓,尿尿或身體不舒服的時候.
就直接哭出來表達.
小朋友,小孩子是需要被人服侍的.
而小孩子的表達都是很直接.
沒有遮掩.
也很容易被人騙.
而在追求屬靈的生命過程.
都是這樣.
都是這樣需要成長.
不過我們不可以一直停留在做小孩的層次.
聖經裡面告訴我們.
我們不要再作小孩子.
那裡是說不要再作屬靈上的小孩子.
第十四節那裡這樣說.

$^{281}$這樣我們不再作小孩子.
眾了人的詭計和欺騙的法術.
被一切邪說之風搖動.
飄來飄去.
當時保羅為何這樣說呢.
因為當時保羅所面對的.
是一個叫做諾斯底主義.
這個學說引誘當時的信徒.
偏離了神的道.
他們相信二元主義.
認為所有物質都不是好的.
唯有靈性才好.
所以他們否定了耶穌基督的位格.
也否定了耶穌道成肉身.
而今天我們所處的環境.
我們都要小心.
最近有機會和弟兄姊妹聊天分享的時候.
原來在我們的群體裡.
都有不少遇到異端向他們招手.
都是很近期發生的.
其中有些甚至在元朗Yoho Mall.
不知道大家有沒有遇過.
在元朗Yoho Mall遇到的.
有些人外表和平常人沒有分別.
他們會主動問對方.
你有沒有去教會.
你是否基督徒.
接著一起聊一些信仰的內容.
甚至想邀請對方去參與他們的教會.
其中一個肢體和我分享.
後來他問對方.
你是甚麼教會.
原來他們對方的教會叫做上帝的教會.
聽名字好像很相似.
大家都是上帝的教會.
但當他發現的時候.
原來他們的教義.
在網上看到.
他們認為上帝是有母親的形象.
而且他們更加認為.

$^{321}$唯有他們這個教會.
才可以迎接到.
新婦迎接到聖靈的教會.
所以弟兄姊妹.
聽下去的時候.
我們其實要小心.
耶穌在福音書裡面也說.
在末世裡面.
有很多的假先知.
很多的假教師會出現.
誘惑我們.
以致離開所信仰的真道.
所以我們需要求主.
讓我們有智慧去懂得.
如何去分辨.
如果我們回到經文裡面.
聖經說我們不要再做小孩子的時候.
那怎樣才能為之長大呢.
聖經裡面說.
直到我們眾人在信仰上.
同歸於一.
認識神的兒子.
新譯本在《伊治的經文》這樣說.
直到我們眾人對神的兒子.
都有一致的信仰和認識.
可以長大成人.
達到基督豐盛長成的新娘.
整個意思就是說.
我們對神的兒子.
要有一致的信仰和認識.
不過這個認識.
不是單單頭腦上的認識.
而是在生命上經歷到信仰.
以致我們真的可以長大成人.
成長到一個地步.
就是可以活出基督的樣式.
就是我們的生命.
可以流露到耶穌基督的性情.
這個就是滿有基督長成的新娘.
所以當我們說建立基督的身體.

$^{361}$有基督長成的新娘.
不是在說教會人數的增加.
而是在說教會屬靈的生命的成長和成熟.
教會屬靈的生命成長和成熟.
也就是我們屬靈的生命的成長和成熟.
所以我們要追求屬靈生命的成長.
我們要追求屬靈生命的進步.
我們不要去做屬靈的小孩子.
接著我們剛才說在經文裡.
我們已經看到神給了教會有恩賜.
而且恩賜的目的是甚麼.
當中的內容 當中的要求.
接著我們看回這段經文.
當中是繼續提到具體的實行.
在經文第56節這樣說.
我們反而要用愛心說誠實話.
各方面向著基督長進.
憐於元首基督.
靠著他全身都憐得緊湊.
百折各按各職.
照著各體的功用彼此相助.
使身體漸漸增長.
在愛中建立自己.
在這兩節的經文裡面.
提到要建立基督的身體裡面.
實際有四個原則.
第一個原則就是要用愛心說誠實話.
在這裡的誠實不是只是包括言語的層面.
即是我們有沒有說謊的誠實.
而是包括行為的層面.
也就是在言語和行為上.
要流露出愛心和真理.
如果在新譯本裡面.
他講的是要在愛中過誠實的生活.
另一個意思是更加清楚的.
即是要在愛裡面持守真理.
就是因著愛神的緣故.
不計代價地去持守真理.
真理大家怎樣去看呢.
特別我們處於現在這個世紀.

$^{401}$其實很容易被後現代的思想影響.
後現代的主義是甚麼呢.
就是認為很多東西都是相對的.
是沒有絕對的.
有很多東西他們會質疑真理的絕對性.
甚至質疑一些權威.
但在聖經裡面教導我們.
神的話就是真理.
所以當我們建立基督的身體的時候.
我們必須要在真理的基礎上.
如果離開了真理.
這也不是教會的建立.
而當我們持守真理的時候.
帶出一個結果.
也是第二個原則.
就是各方面向著基督長進.
連於元首基督.
意思是承認基督是頭.
元首的意思就是基督是頭.
就是接受基督的權柄.
頭在身體上是一個很重要的角色.
大家可以想像頭是一個帶領的角色.
當你的頭轉向哪個方向的時候.
你的身體很自然就會配合.
而且頭是有腦袋的.
也會指揮身體去決定應該怎樣運作.
所以當這裡說要連於元首基督的時候.
我們要接受耶穌基督的帶領.
接受祂的權柄.
也是真正的去信服耶穌基督.
當我們去決志相信主的時候.
我相信每一個人都是這樣去相信耶穌.
就是我們接受耶穌成為我們的救主.
祂在十字架上為我們袋贖.
救了我們.
但是耶穌除了救了我們之外.
祂也是我們生命的主.
這是很重要的.
很多時候我帶領別人去決志的時候.
我都很刻意去提.

$^{441}$耶穌除了救了我們之外.
也是我們生命的主.
祂救了我們.
我們的命是屬於祂的.
我們信主的一刻.
我們的生命是應該由耶穌去決定.
就好像你開車.
當你信主的時候.
耶穌上了你的車.
耶穌應該是車的司機.
我們最多坐在副駕位.
我們的車的方向.
是應該由耶穌來開.
但是坦白.
現實可能有時候我們的生命.
是將耶穌坐在旁邊.
或者坐在後座.
或者根本耶穌沒有上我們的車.
我們的車的方向去哪裡.
究竟是由我們自己決定.
還是真的由耶穌去決定呢.
所以在經文裡面說.
我們要憐於元首基督.
不是單說我們要接受.
耶穌基督的權柄.
而是真正的去信服耶穌基督.
真正的去跟著祂的帶領而行.
如果我們再看下去的時候.
第三個原則就是.
我們全靠著祂.
全身都連接得緊湊.
我們知道教會.
剛才也說了.
教會我們不是一個人.
教會是你和我.
是我們每一個的信徒.
我們提過教會不是.
教會的身體不是一個的肢體.
是由很多的肢體去組成.
在我們基督徒的群體.

$^{481}$聖經一向教導.
我們都不是單打獨鬥.
之前我也有提過.
可能大家也聽過.
一個關於十字架的說法.
就是一個十字架.
有中向和橫向.
中向就是我們和神建立的關係.
橫向就是我們和人的相交.
今天我們作為基督徒.
在教會.
沒錯我們一定要和神建立那份關係.
個人的關係.
但另一方面我們也不能缺少.
和弟兄姊妹團契層面的相交.
而基督的身體.
就是我們這些很多很多不同的肢體.
一起去組成.
成為一個身體.
但聖經裡面說.
你組成一個身體.
單單放在一起也不夠.
第四個原則說.
我們要百折各按各職.
照著各體的功用.
彼此相助.
就是原來當上帝.
將我們每一個.
配搭成為一個身體的時候.
我們要在不同的崗位上.
去顯出功用出來.
之前剛才也提過.
神給我們每一個.
都有不同的恩賜.
我們每一個最少也有一樣恩賜.
你的恩賜可能和旁邊那個不同.
你有的.
其他人未必有.
其他人有的.
你又可能未必有.

$^{521}$如果我們再看.
現代中文譯本.
在這一節.
也更加好.
他將百折各按各職.
譯為關節跟絡.
即是表達出.
一種血肉相連.
那種密切的關係.
就是因為我們肢體之間.
那個關係這麼緊密的時候.
唯有我們每一個的肢體.
我們各人發揮.
我們自己的恩賜.
我們不斷地互相去配搭.
教會這個群體.
不是一個人.
這樣才可以去成長.
而最後.
保羅在這裡這樣說.
當每一個的肢體.
去發揮.
他自己的功用的時候.
耶穌基督的身體.
即是教會.
不單止.
會在愛中被建立起來.
同時.
也會在愛中漸漸地去長大.
在這裡保羅沒有說.
要在律法之中.
去被建立出來.
也沒有說.
要在工作效率之中.
被建立出來.
在這裡說.
是在愛中建立出來.
所以我們看到.
當我們說整個教會建立的過程.
最後愛是一個很重要的元素.

$^{561}$愛是一切的原動力.
一切所做的原因.
是因著愛.
是因為神愛我們每一個.
而我們回應神的愛的緣故.
所以我們一同去建立基督的身體.
最後.
盼望今天的訊息.
都成為我們彼此的勉勵.
願意真是我們被神的愛去激勵.
以致我們一同的.
可以去願意去發揮.
我們各人的恩賜.
去努力去建立基督的身體.
去建立神的家.
我們一同的.
同心去禱告.
參與了我們再一次的道理面前.
你提醒我們.
你給我們今天每一個相信你的人.
都有恩賜.
而我們的恩賜.
我們是為著要建立教會.
求你幫助我們.
幫助我們在我們個人上.
我們真是知道.
發掘神給我們的禮物.
以致我們真是可以在教會裡.
好好地去運用出來.
最後我們真是盼望.
耶穌基督你的身體.
就是教會可以被建立.
主我們都是這樣去禱告.
去仰望,教託你.
因為我們知道.
耶穌你是教會的元首.
你是教會的頭的時候.
我們就是這樣去禱告,仰望,教託.
求你自己的話語成為我們的激勵.
成為我們的提醒.

$^{601}$讓我們去建立教會的時候.
我們去事奉的時候.
不是單單因著需要.
因著律法.
乃是因著愛你的緣故.
以致我們希望教會.
去被主你自己建立起來.
天父我們都是將我們每一個去交上.
求主你繼續去帶領.
願意我們每一個在你的愛中.
繼續去成長.
繼續去學習主你那份的樣式.
以上的禱告.
奉主耶穌基督明智祈求.
阿們.
\newpage



\section{以弗所書 4:11-16}
\label{sec:qI_KwviwgzU}
\textbf{2023-07-30 加恩樓《建立基督的身體》以弗所書 4\_11-16 講員:楊明麗姑娘}
\newline
\newline
連結: \href{https://youtube.com/watch?v=qI_KwviwgzU}{\texttt{https://youtube.com/watch?v=qI\_KwviwgzU}} ~~~~ 語音日期: 2023-08-08
\newline
\newline
\hyperref[sec:c_zS1WaX1q0]{\small{< < < PREV SERMON < < <}}
~
\hyperref[sec:index_chronic]{\small{[返順時目]}}
~
\hyperref[sec:index_scriptual]{\small{[返順卷目]}}
~
\hyperref[sec:i3m_mXZnkC4]{\small{> > > NEXT SERMON > > >}}
\newline
\newline
以弗所書 4:11-16
\newline
\begin{longtable}{cl}
\hline
\hline
章節 & 經文 (和合本修訂版)\\
\hline
4:11 & \begin{tabularx}{0.7\textwidth}{X} 他所賜的有使徒，有先知，有傳福音的，有牧者和教師， \end{tabularx} \\ \\ \relax
4:12 & \begin{tabularx}{0.7\textwidth}{X} 為要裝備聖徒，做事奉的工作，建立基督的身體， \end{tabularx} \\ \\ \relax
4:13 & \begin{tabularx}{0.7\textwidth}{X} 直等到我們眾人在信仰上同歸於一，認識神的兒子，得以長大成人，達到基督完全長成的身量。 \end{tabularx} \\ \\ \relax
4:14 & \begin{tabularx}{0.7\textwidth}{X} 這樣，我們不再作小孩子，中了人的詭計和欺騙的法術，被一切邪說之風搖動，飄來飄去。 \end{tabularx} \\ \\ \relax
4:15 & \begin{tabularx}{0.7\textwidth}{X} 我們反而要用愛心說誠實話，各方面向著基督長進，連於元首基督， \end{tabularx} \\ \\ \relax
4:16 & \begin{tabularx}{0.7\textwidth}{X} 靠著他全身都連接得緊湊，百節各按各職，照著各體的功用彼此相助，使身體漸漸增長，在愛中建立自己。 \end{tabularx} \\ \\
[1ex]
\hline
\hline
\end{longtable}
$^{1}$今天有一樣很特別的.
我很想你.
現在和你旁邊的那位說一聲.
今天很開心.
在你的念想我見到神的愛.
和你旁邊的那位說一聲.
很開心今天在你的念想.
我見到神的愛.
是啊,我回來後也很開心.
我見到你有我的團友在.
很開心的.
很少機會能夠見到面.
但是回來後.
又好像很久沒見.
但又好像很親切.
進來後.
過了玻璃門第一件事.
我從那邊進來.
過了玻璃門第一件事.
你知道是什麼事嗎?.
有一個團友問我一句話.
要不要喝咖啡.
哇,多溫馨啊.
很感恩.
在哪裡會有這樣的說話?.
在家裡.
在家裡會有這樣的說話.
要不要喝一杯咖啡?.
你咖啡想不想加奶?.
這是家人很親切的問候.
回來後.
有一個家的感覺.
這是我們的家.
在家裡.
很多事情是不同的.
回家後.
你也會有些事情.
有些事情你想得到.
有些事情你想付出.
最近.

$^{41}$我先生有幾天.
不舒服.
不舒服的意思是.
他不舒服的意思是.
不是病.
不舒服的意思是.
不是很舒服.
不過不是病.
所以不舒服是不看醫生的.
各位太太有沒有同感?.
不舒服是不看醫生的.
我沒有病.
我只是不舒服.
作為太太的一家人.
你知道她想做什麼.
我沒辦法說話.
我就說.
我煮粥.
我想吃粥.
不是你病了.
是我想吃粥.
然後在街市附近.
買了一條魚回來.
做魚粥.
她很開心.
她覺得.
其實我就是想吃這些東西.
為什麼你知道?.
她不知道我為什麼會知道.
她有問題.
回到一個家.
有些事情我們知道.
不說我們都知道.
有些事情是.
慢慢就知道.
慢慢就會知道.
這是我們在家裡的感覺.
進了一個家.
你就會知道.
你知不知道自己?.

$^{81}$認不認識自己?.
在一個家裡.
你會做什麼?.
有什麼事情.
你會在一個家裡參與?.
你幫忙探訪.
關心.
你幫忙為一些重要的事情.
去祈禱.
幫忙教主.
有很多不同的事情.
可能是你的參與.
幫忙.
買好的咖啡回來.
幫忙.
做好的咖啡.
幫忙.
預備好美味的奶.
加入咖啡.
幫忙看看.
哪裡有喝完咖啡的杯.
都是一個家裡.
大家要做的事情.
到底你的恩賜.
是什麼呢?.
你的恩賜是教導?.
是行政?.
財務?還是服務?.
服務一定是大家都可以的.
很多不同的事物.
都需要我們去服務.
還有一些是歌頌讚美.
大家一回來.
順拜開始的時候.
大家很想做一件事.
歌頌.
讚美主.
這是我們回來教會.
第一件很想做的.
我們都很需要有這樣的恩賜.

$^{121}$弟兄姊妹去幫我們.
但剛才我們讀了一段經文.
他不單說給我們有這些恩賜.
還有特別的.
他專注的恩賜.
首先在第七節他說.
我們每一個人.
我們都是蒙恩的.
蒙了一個恩典.
主耶穌基督.
每一個人都賜一些恩賜給我們.
這些恩賜是主耶穌基督.
自己衡量分配的.
他覺得我們哪一個.
應該有什麼恩賜.
他就給我們什麼恩賜.
這些是一個恩.
是一個禮物.
送給我們的.
不用給錢買.
神就是給了我們.
甚至有很多人.
他不用學.
他做到的.
他懂的.
但特別他又多說了幾樣東西.
他所賜的.
還有一些是做使徒.
有使徒有先知有傳福音的.
有牧師和教師.
準備聖徒做事奉的工作.
建立基督的身體.
原來神為了在頭家裡面.
有人去分配工作.
建立頭家.
做下督導.
他做到一件事.
就是分配了幾種人在這裡.
一些是做使徒先知.
一些是傳福音的.

$^{161}$一些是牧師和教師.
使徒先知傳福音的.
都是有一個「有」字在前面.
意思是有這種人在前面.
牧者和教師.
前面只有一個「有」字.
不是有牧者和有教師.
表達牧者和教師.
是一種人.
這是一個人.
他會有這樣的.
有牧養和教導的恩賜.
這是他給我們的.
我們當中.
有些人可能強項就是.
他有一個「有」字.
有些人可能強項就是.
當教會還沒成立的時候.
他幫忙去建立教會.
有些是先知講道的能力.
有些是傳福音的.
明明你跟他說了很多次福音.
他都未必肯聽.
但某人跟他說.
都未說完.
「我信了,我祈禱了,可以了」.
是這樣的.
剛剛信主的時候.
我跟他說了很多年福音.
他都不斷說.
後來有位師母.
那晚只跟他說了一句.
「你信吧,我和你祈禱」.
他不斷地說了很多次.
「你信吧,我和你祈禱」.
他那晚就真的信了.
有些人.
他特別帶領很多人信主.
那是神給他的禮物.
他就是去.

$^{201}$傳福音,叫人能夠信主.
有些人.
他是對牧羊.
關心,教導.
很有這種恩賜.
在教會當中.
他做的是幫助.
弟兄姊妹有需要的時候.
他就在他們身邊.
去關懷他們.
陪他們走過一段路.
陪他們.
可以走過很多不同的路.
開始去探學.
探索哪一班適合我學.
開始去做一下事奉.
開始.
面對一些人生.
自己的問題,自己的事情.
也需要有一些.
牧者,教師在身邊.
牧者,教師,使徒.
先知,傳福音的.
都是做一樣的事情.
是裝備聖徒.
裝備聖徒,做事奉的工作.
建立基督的身體.
誰是聖徒?.
先看看誰是神聖.
是不是頭上有一個光環?.
照聖經的教導.
每一個信徒.
都是聖徒.
每一個聖徒.
都是平信徒.
我們都是在主的恩典當中.
我們一起得著這個恩典.
就是做聖徒.
做聖徒.
我們一起走一條路.

$^{241}$就是要在信仰當中.
我們能夠.
一起去到.
更認識,更明白.
這個信仰.
更認識基督.
能夠讓整個教會,基督的身體.
基督的身體是指教會.
整個教會.
都能夠均衡地成長.
我今早.
過來的時候.
有位姊妹送我過來.
開車送我過來.
我從錦繡那邊過來.
她開車送我過來的時候.
我就知道她有一隻手.
她最近打球受傷了.
腫了.
當我們說均衡成長的時候.
我們每一個部份都這麼大.
不會某一個位特別.
發達了.
哪一個位特別小.
有的話.
都是有不舒服.
需要康復.
我們需要均衡地成長.
我們均衡地成長.
是有使徒.
有先知,有傳福音.
有牧師和教師.
帶領著的時候.
教會一起成長的時候.
整個教會.
不再是小朋友.
不是說我們當中沒有小朋友.
我們有兒童教育在.
我們有主一學.
不單止是有小朋友.

$^{281}$而是每一個人.
信心都很堅定.
有時小朋友的信心都很強.
很堅定.
最厲害是和誰傳福音.
爺爺,奶奶.
你未必能叫他們.
回教會.
孫女們都能叫他們回教會.
我們在信心上.
信仰上.
不再做一個小孩.
以至於別人可以騙我們.
哄我們.
或者用一些不太清楚的話.
引我們去信.
一些不正當.
的異教.
或者異端.
這些現在真的很多.
甚至有人會在.
你坐車的時候.
或者等車的時候.
有人跟你說.
你回教會開不開心?.
原來有些人會覺得不開心.
然後就說.
不開心你來我的教會.
你來我的教會就會很開心.
其實我不知道他那間是什麼教會.
原來要.
引人進入異端很容易.
你不知道是什麼教會.
就可以有人去了那間教會.
我們需要更加認識清楚.
認識清楚我們的信仰.
不要容易被人帶走.
不要容易被人騙了.
這是神.
祂希望教會當中.

$^{321}$有這樣的培育.
所以教會時時有不同的課程.
剛才我聽見鄧全說.
我們有一個.
淺談三一論的課程.
可能不止這些課程.
這兩三年.
在YouTube上大家聽到.
很多不同的教導.
最重要是對的.
找到一些你熟悉的,知道正確的.
教導的導師.
或者牧師.
回到教會你可以受到.
教會覺得是正確的.
已經安排好的教導.
我們可以得到.
正確的教導.
讓我們在信仰上.
可以朝著正確的方向.
一直去成長.
這個方向我們一直成長的時候.
我們出現一個警方.
大家都會信得一樣真.
這間教會是什麼教會.
保羅這樣寫信給他們.
是不是他們當中有不正確的信仰.
原來不是.
這間教會是伊弗所.
伊弗所是保羅花了很多時間.
在當中教導他們.
起碼他一開始.
就用了兩年的時間.
在這裡傳福音和建立教會.
後來他說.
他在這裡差不多三年.
在當中除了培育信徒.
他還做了一件事.
就是繼續安立.
接下來的傳道人.

$^{361}$長老和牧師.
他也在當中建立.
讓這間教會在信仰上很穩固.
這間教會.
在信仰上很穩固.
在應該是.
頭五個世紀當中.
公元這五個世紀.
頭五個世紀的時候.
基督教信仰的中心地帶.
一直移向伊弗所.
直到黑暗時期的開始.
伊弗所.
都是基督教的宗教信仰的中心點.
保羅跟他們說.
要很堅定地尋求正確的信仰.
讓牧者在當中去牧養.
不是因為他們走錯路.
而是告訴他們.
要繼續前進.
為什麼要繼續去紮好根基呢.
伊弗所在2015年.
伊弗所城.
有很多不同的地方.
發掘了出來的.
是做了什麼.
是文化保育區.
是考古出來很多文物.
是真實有這個地方.
這個地方很繁華.
它是近一個海港.
它有一個海港.
所以船隻可以進去.
有很多豐富的漁產.
也有很多不同地方來的物品.
所以商貿上很發達.
這個城市發達.
所以有不同的地方去到.
它是發達到一個地步.
它是羅馬五大城市之一.

$^{401}$它是發達到一個地步.
所以人口有很多不同的民族在這裡.
不同的民族在這裡.
於是有很多不同的宗教信仰.
在保羅去傳福音的時候.
是有很多不同宗教信仰的人.
信了主 認識了主.
開始回教會.
他們將家裡裡面.
一些不是信基督的書籍和物品.
都拿出來燒了.
他們信得很認真.
有人因此而改學.
不做他原生所做的.
拜祭偶像的行業.
這是一個很認真的教會.
在一個這麼認真的教會的時候.
其實它是.
它是存在在一個.
還有很多異教的地方.
而且在第一世紀.
一直進展的時候.
越去到第二世紀的時候.
當時的封信已經寫了.
這個方向就是.
異端在那裡很繁榮.
人多.
貿易好.
交通方便.
於是異端在當中.
發展得很好.
保羅的教導是很有神的感動.
真的要在信仰正確的時候.
扎穩根基.
信得更加堅固.
認識神 認識得更多.
更加穩陣.
保羅也跟他們說.
他在他要被捉拿的時候.
他將已婚所的長老.

$^{441}$邀請了他們去見最後一面.
這個時候是何時呢.
是耶穌基督出來宣教了三次之後.
他差不多就要被猶太教的人.
來謀害他 要殺害他的時候.
他聖靈感動他.
他知道他再去耶路撒冷的時候有危險.
但他覺得我還是要去的.
於是乎他還是要去.
但去之前他就召喚了已婚所的長老.
來跟他們說.
他會遭遇到什麼事.
跟他們道別.
告訴他們我們不是很長時間.
會見不到面.
是以後都可能見不到面.
他遇到最危險的話.
他應該會沒有了生命.
但這次他被收監.
他沒有立刻被行刑.
他反而是後來再被放監出來.
再一次被捉拿的時候才喪失生命.
但他跟已婚所有這麼密切的關係.
亦都是他很清楚.
很認識已婚所 很愛已婚所的教會.
他有跟他們這樣建立的關係.
但他仍然跟他們說.
他繼續在信仰當中盡心.
他講到信仰當中盡心的時候.
其實他一早在四章第一節.
就講了.
怎樣在信仰上盡心.
他說.
他這個時候寫已婚所書的時候.
就已經過了.
之前一次我剛才說.
他請了已婚所的長老.
跟他們道別.
說自己會被捉拿.
他真的被捉拿 坐了幾年監.

$^{481}$他寫信慰問他們的時候.
他說.
我是為主而坐監的.
我勸你們一件事.
就是什麼呢.
你信了主.
有一個新的生命.
有永恆的生命 有永恆的盼望.
我們有永恆的生命 有永恆的盼望.
要怎樣呢.
原來我們的生活.
我們活出來的生活.
要表現出這種生命.
最近在成都.
是不是有一個運動會.
開始了.
他們的跑姿.
他們出來的那種氣勢.
很能表現出他們就是運動員.
無論什麼年紀.
是導師也好 領隊也好.
或者是那幫運動員.
他們出來就有那種氣勢.
你有沒有看到他們.
拿著一支很長的棍.
上面有一個國旗.
然後他們可以拿著國旗.
不簡單的.
很重的.
我拿著.
我不能甩.
這個是.
你的生命.
那種生命力.
可以在生活當中表現出來.
你就是可以在.
無論什麼環境當中.
都將這種生命力.
可以釋放出來.
當我們有永恆的生命的時候.

$^{521}$我們這種永恆生命的生命力.
在無論什麼境況.
它都能夠表現出來.
表現出我們有永生的盼望.
表現出我們有一位.
永恆的主一直在照顧我們.
表現出我們的生命.
那種質素.
是我們有所期盼的.
我們不是失望.
不是痛苦地過日子.
雖然日子不容易.
日子一定不容易.
繁華的都市.
痛苦的日子.
我們都能夠.
在繁華的都市.
痛著會很強.
不過有主的時候.
我們知道.
如何生活在這個世代.
我們有什麼事.
我們都還可以禱告.
我們有困難.
都還可以訴苦.
我們有困難.
我們都知道我們有盼望.
我們有神.
祂不是遠在天邊那麼簡單.
祂在天上坐著為王.
祂在你我的心裡.
祂知道我們的日子.
祂陪伴我們經過.
最近有位姊妹跟我說.
她說.
我很想去那個地方旅行.
但是呢.
我又不想跟兒子說.
讓他請我去旅行.
我又不想跟丈夫說.

$^{561}$好像逼他要跟我去.
一個他不一定想去的地方.
我又不想跟丈夫說.
好像逼他要跟我去一個他不一定想去的地方.
我又不想.
特意去求神.
所以我沒有祈禱.
誰知道.
我打個電話.
問姊妹.
你有沒有空陪我去旅行.
姊妹都沒有回覆我.
但是晚上我收到兒子的電話.
他說.
媽媽.
我想什麼時候請你和爸爸去哪裡旅行.
我女朋友不會去.
是我們三個一起去.
因為如果女朋友去晚了.
她就不能陪爸爸媽媽.
姊妹覺得.
我都沒有祈禱.
我都沒有要求.
我只是想一下.
我找姊妹說一下.
神聽到嗎.
神是不是聽到.
真的嗎.
聽到嗎.
他們說是.
原來神真的聽到.
你都說是.
原來神聽到.
我一個.
這麼微小的期望.
其實又與性命無關.
的事.
神竟然都應找我.
是.
祂就在我們心裡.

$^{601}$你不說那句都知道.
祂就是這樣.
照顧幫助我們.
我們在生命當中.
我們得到永恆的主的看顧.
在生活上.
我們要表現出來.
我們是一個有盼望的人.
我們是一個有永恆生命的人.
保羅說.
我們一個這樣的人.
我們應該是怎樣呢.
我們凡事要謙虛溫柔忍耐.
用愛心互相寬容.
以和平彼此連繫.
竭力保持聖靈所賜的合一.
原來.
我們不是要.
裝扮得很謙虛.
裝扮得很溫柔.
而是我們要讓生命流露.
一件事出來.
我們要讓生命流露一件事.
就是神想我們合一.
我怎樣都要學會合一.
原來保羅告訴我們.
我們要合一.
我們有永恆的生命.
我們學習去合一.
有永恆的生命.
代表我們將來要去的地方.
在天上.
在神的樹.
有一個永恆的國度.
我們永遠都在樹上.
有人說這樣會很悶.
不是,那裡多姿多彩.
沒有什麼你不能做的.
這樣就很好.
不過你會不會有些地方不能去.

$^{641}$在天堂你會不會有些地方不能去.
我和阿邊的不太合.
我看到他後面我會轉身.
會不會有這樣的時候.
如果你在天堂.
有很多人你不想見到.
不能見到,不想和他講說話.
那去到天堂你就很沒有自由.
很多地方不能去.
一會兒你去找個地方喝咖啡.
原來見到他.
又要出來了.
讓我們在主裡面學習合一.
學習合一.
我們將來都要在同一個地方.
一起享受永生.
我們都要一起去敬拜.
一起去敬拜同一位的主.
在天堂也要一起去敬拜同一位的主.
我們未至於到.
要去到天國的時候.
跟著告訴天父.
天父,你讓我進去嗎.
我先不進去.
你幫我選擇一個位置.
他在左,我就在右.
不用我.
原來我們可以在天國的時候.
我們合一.
我覺得在佳音樓有一樣東西.
很感恩.
是大家一起去做.
咖啡.
很多不同的人一起去學習.
和炮製.
我很開心.
在程序表我就見到.
私事的群組.
有新的人.
有些名字我不認識.

$^{681}$或者我認識的.
不過當時他還不是私事.
有些人.
現在是幫忙兒童的工作.
有些不同的名字出現了.
是的.
這裡是一個合一的地方.
我們回來.
我們可以是其中的一份子.
我們一起去努力.
這是我們的家.
普羅跟他們說.
神很想我們做的就是合一.
而我們要合一.
有困難的話不用擔心.
我們盡力.
但聖靈會幫助我們.
使我們合一.
而且我們有一個很能夠合一的基礎.
很多的一.
為什麼全部一都是紅色的呢.
因為大家方便看到.
不是代表一就不合格.
紅色有很多不同的意思.
我們有同一個身體.
我們是基督做我們的頭.
我們是這個身體裡面.
我們有同一位的聖靈.
在我們心靈裡面做感動的工作.
幫助我們.
教導我們.
安慰我們.
教我們祈禱.
我們接受同一個的洗禮.
我們信同一個的主.
我們有同一個的主宰.
這是我們生命的救主.
和生命的主宰.
我們有同一個的神.
我們是得蒙這個神.

$^{721}$做我們的召喚.
認識祂 歸向祂.
這個神是一致的.
有人會問.
到底誰是神呢?天父主耶穌基督還是聖靈呢?.
如果我祈禱我只是向天父祈禱.
基督和聖靈會不會不開心呢?.
感恩啊.
天父不單叫我們合一.
聖靈幫助我們合一.
基督拯救了我們.
讓我們可以信同一個的神.
而且天父主耶穌基督聖靈.
他們是合一的.
他們做了榜樣讓我們看到.
什麼是合一.
在救恩上天父做了策劃.
主耶穌基督做了執行.
聖靈感動我們去信主.
但是主耶穌基督沒有說.
很不公平啊.
我這麼痛苦.
我這麼痛苦那幾個小時你做什麼呢?.
沒有啊.
他們真的合一的.
他們合一地.
對我們有同樣的要求.
期望我們都是合一.
沒做到呢?.
不要緊.
繼續學習.
神不是叫我們一下子轉變.
而是給我們時間.
我們一路去學習.
讓我們的生命變得更加的合一.
在合一當中的時候.
我們有人有不同的恩賜.
幫助我們在.
基督的教會當中.
基督的身體上面.

$^{761}$我們更加能夠調節得到合一.
靠著基督的時候.
我們全身可以.
拼湊得更加好.
各節.
各按各職.
全身都能夠彼此互相幫助.
使到身體.
能夠繼續去增長.
在愛當中被建立.
弟兄姊妹.
我們能不能夠調節得很合適呢?.
其實是可以的.
你買回來的拼圖.
它是能夠拼湊的.
因為它的製造.
就是要拼湊得很合一.
所有的連接位.
是能夠合得很好.
神製造我們的時候.
也是的.
我們是能夠合併得很好.
就是因為各有長短.
不同的恩賜.
所以我們合併得很好.
我們不需要要求每個人有同樣的恩賜.
同樣的強項.
同樣的思想.
不過我們有一個合一的心.
大家就能夠合併得很好.
思想不一樣.
不要緊.
我們有共同的信仰.
我們的信心會一路一致.
一路邁向一致.
能夠彼此在信心上.
互相鼓勵.
信得堅定.
其實人有很多東西都不同.
同一個家庭成長.

$^{801}$也有很多不同的東西.
家庭當中.
你有多少兄弟姐妹.
我們這個年紀應該多一點.
面對青少年.
他們一個喜,都是一個子.
我們有很多兄弟姐妹.
兄弟姐妹會不同的.
會有不同的恩賜.
每個人的強項都不一樣.
但是不同的強項.
就可以在同一個家庭當中.
才開心.
大家互相幫補.
互相鼓勵.
一個很厲害.
另一個有另一方面厲害.
大家就是這樣併在一起.
這是神給我們的恩典.
就是因為.
我們不同.
所以能夠一起建立起來.
如果一些東西.
是你認識的.
我不認識沒有問題.
有時候我找你幫忙.
如果有些東西是我認識的.
你有時候就找我幫忙.
我們在一頭家.
教會就是我們的家.
我們在一頭家當中.
我們就是這樣一起併合.
有時候遇到一些問題.
不只是你遇到.
可能大家都曾經遇到過.
出聲.
有人幫到忙.
我們不知道為什麼.
不過神預備了在我們身邊的人.
就是祂預備了.

$^{841}$做我們的天使.
神預備你在人身邊的時候.
你也是人家的天使.
我們要這樣.
在一間教會當中.
一起建立起來.
這個身體就會增長.
不只是人數多了.
是家庭更加健康.
大家的關係.
我們的信仰.
會更加鞏固.
面對困難面對問題的時候.
我們能夠一起去處理和面對.
這是神給我們的恩典.
我們可以這樣.
一起去生活.
表現出.
我們有永恆的生命.
有主的生命.
在我們裡面.
你有什麼恩賜.
恩賜由神而來.
恩賜是基督.
按著祂的意思.
分配給你的.
一定有的.
你能夠想像.
應該給多少.
給什麼.
你有的.
或者你的恩賜.
是一個建立教會的恩賜.
你在教會.
在監牢.
還是很少人的時候.
你已經在雪.
有神的預備.
你在雪.
你看著監牢一路增長.

$^{881}$你看著領師.
從每個星期.
幾乎你要去領師開始.
到現在.
有人負責鍵盤.
有人負責吉他.
有和唱.
有團隊.
有一個團契.
你看著它成長.
我們一路.
一路看著.
但是.
原來是這樣的.
教會是沒有過客的.
每個回來.
都是教會的一成人.
教會沒有過客.
教會沒有觀眾.
崇拜.
我們是做給神看.
每一個都參與.
敬拜歌頌.
我們是給神看.
我們當中沒有觀眾.
沒有過客.
我們是其中一份子.
你的東西.
神給你的東西.
你的智慧.
聰明.
學識.
你的才幹.
在教會裡.
神是想你怎樣去做.
祂想我們能夠一起合一.
建立好一間教會.
在神裡建立好.
讓我們的教會.
在神的面前.

$^{921}$能夠信仰方面.
更加強壯.
更加密切.
全福音上我們更有力量.
更有恩賜.
人數上繼續去增長.
保羅寫這幅所書的時候.
四章一到三節就告訴我們要合一.
去到第十六節.
四章十六節.
他都說我們要合一.
我們是由信主之後.
我們知道我們和神有一個和好的關係.
我們先和神合一.
然後我們和弟兄姊妹.
開始學習合一.
現在經過教會的栽培.
我們更加認識多了弟兄姊妹.
我們有不同的恩賜.
我們回來事奉.
我們回來學習有成長.
我們一路事奉.
一路成長的時候.
大家整個群體更加合一.
我們是由個人和神的.
關係的合一和好開始.
去到整個群體我們都去學習.
合一.
合在當中成長.
我們要整體一樣地成長.
弟兄姊妹.
我也想有兩個人請你們幫忙.
祈禱.
我問准了這兩個家庭.
有兩個人我很想你們為他們祈禱.
一個就是我們現在.
關心的一個家庭.
她的小朋友.
現在大概九個月至十個月大.
還在醫院.

$^{961}$還沒回家.
她一出生.
就是磅半.
大家知不知道磅半有多大.
我為了知道磅半有多大.
我特意去豬肉舖.
那裡有骨肉.
有胸肺有腦.
但一磅半.
就是這麼多.
這是生命.
是個嬰兒.
是個經過醫生.
做了很多事.
她生存到今天.
她現在有差不多二十磅.
每一磅.
每一個安士都是恩典.
在疫情期間.
生了嬰兒.
不可以接回家.
不可以見她.
這個嬰兒就是從這樣開始.
你現在坐在這裡.
可以抱她.
都有差不多大了.
我們期待她能夠在八月尾出院.
她要怎樣才能出院呢.
有兩件事一定要做到.
你猜是什麼.
第一.
呼吸.
不呼吸.
就一定要留在醫院.
第二.
能吃.
不能吃.
也要留在醫院.
醫院會用不同的管去幫助她.
兄弟姐妹.

$^{1001}$我們的生命可能就是這樣.
我們初訊的時候.
只有一磅半的淑齡生命.
很多人在我們的生命當中.
做了很多功夫.
建立我們.
幫助我們.
讓我們的生命可以成長.
我們可以從一磅半.
變成十五磅.
二十磅.
由一磅半到十五磅.
醫生告訴我們.
十五磅的時候.
已經是十倍了.
我們會覺得.
一個七八個月的小朋友.
十五磅.
好像還是輕輕的.
但是已經是十倍了.
她可以.
第一口喝到奶.
媽媽可以進去餵奶.
那個感覺.
多辛苦都可以.
在市區上班.
下班後.
馬上衝回來.
去屯門醫院.
又消毒一段時間.
穿了袍子.
進去抱她一會兒.
吃奶.
很辛苦的.
回家還有一個哥哥.
哥哥也是兩歲多三歲.
辛苦嗎.
辛苦最重要是能抱到她.
媽媽說.
是的.

$^{1041}$我們的生命.
神也是這樣看我們.
祂想我們能夠彼此相愛.
生命成長.
教會也是這樣成長.
我們信主的生命也是這樣成長.
這個小朋友.
最近做了一個測驗.
做了一個測驗.
就是看她的呼吸情況.
醫生說她.
還是需要養氣.
不過.
那個度數可以回家.
多開心.
家裡三代人.
都很開心.
最後孫子可以回家.
兒子可以回家.
回家?.
還沒馬上回來.
為什麼.
爸爸媽媽要學習.
如何讓兒子回來.
如何做運動讓他呼吸.
如何幫他做口部運動.
讓他可以吃東西.
他現在能吃一點.
不過要繼續操練.
原來.
不是每個小孩出生.
就一定會喝奶.
你身邊的弟兄姊妹.
是不是有些事情做得不太好.
是的.
不是每個小孩出生.
就會喝奶.
我們要一起去建立我們的家.
建立我們的教會的時候.
我們要有這個容量.

$^{1081}$讓每一個在當中成長.
合一地一起成長.
讓這個家更加強健.
身體更加強健.
可以好好地去增長.
我正在關心另一個人.
也請大家為他禱告.
剛才那個為他祈禱.
讓他八月底才能回家.
家裡要有設施.
父母要有訓練.
哥哥要能接受弟弟回家.
因為弟弟一定會用父母很多時間.
哥哥多少會被忽略.
另一個家庭我關心的.
也是一個七個月大的.
兩個都是七星孩子.
另一個也是一個一出生.
七個月大的.
他好一點.
兩磅二.
也有一公斤.
以後你買一公斤東西的時候.
你去想想這個小朋友.
他是三十年前出生.
所以當時的醫學.
沒有現在那麼發達.
他三十歲了.
他三十歲.
他退休了.
他不是找不到工作.
他是要退休了.
你想想.
如果我們三十歲的時候.
不是因為你找到工作.
我提早退休.
而是被告訴你.
你要退休了.
你的心情是怎樣.
我是在關心這個家庭.

$^{1121}$是的.
原來.
七星孩子健康快樂成長.
在社會工作不是必然的.
他也有讀書.
他也有出來工作.
但是.
但是到現在三十歲.
他其實退休了三年.
嘩.
是的.
教會裡面會有些人你會覺得.
為什麼他不長大.
他好像不是.
信主那麼上下時的人.
是的.
我們生命會有軟弱的時候.
人會有軟弱的時候.
但是我們為了有需要的人.
更加要去禱告.
更加要幫助他.
讓他能夠信主三十年.
都要繼續好好去成長.
讓他的生命更加強健.
這位我們幫他的.
不是幫他成長.
他已經三十歲.
但是.
他的基建不會再更加好.
他的身體的基建不會再更加好.
我們很希望他能夠早點信主.
很希望他能夠有永恆的盼望.
我們不知道他生命會有多長短.
因為他突然之間.
27,28歲的時候.
因為他突然之間.
27,28歲的時候.
醫生就告訴他.
你是時候要退休了.
我記得他是在疫情前發生的.

$^{1161}$接著就進入疫情.
所以他應該是.
27歲左右.
現在已經30歲了.
是的.
他需要永恆的生命.
他需要被建立.
我們身邊也有人需要.
他需要被建立.
在我們的身體當中.
成為一個合一的一部分.
在我們的身體當中成為一個合一的一部分.
阿兄姊妹.
你有恩賜嗎?.
有的,找到沒有?.
在教會我們.
沒有人是過客.
沒有人是來觀賞的.
我們每個人都是一份子.
你有恩賜.
用你的恩賜.
成為一個健康.
有生命力.
能夠合一.
能夠和人一起去事奉的人.
我們.
讓自己可以參與.
也讓身邊的人可以參與.
好不好.
和你身邊的那位說一句.
歡迎你加入事奉.
一起努力.
好不好.
歡迎你加入事奉.
一起努力.
我們一起禱告.
參見天父上帝.
我們求你紀念.
這兩位七星仔.
主啊,他們需要你的生命氣息.

$^{1201}$求主建立幫助他們.
求主紀念他們的家庭.
紀念他們走過的路.
求主安慰他們.
主啊,求你紀念我們.
紀念我們這個大家庭.
紀念我們走過的路.
我們需要你.
我們很需要你.
甚至主啊,嘉恩樓.
接下來要做一個很大的決定.
主啊,我們每一個.
都是嘉恩樓的一份子.
我們每一個.
都是你的兒女.
我們很想學一件事.
主啊,我們不只是放在心上.
我們要向你求.
我們要求你賜給我們合一.
我們會有不同的恩賜.
會有不同的思想.
因為我們有不同的人生經歷.
但是我們將.
這麼大的決定交在你手中.
我們很想.
知道你的心意.
很想按你的心意去決定.
求主建立我們,幫助我們.
讓我們有合一的心.
一起去分享.
一起去守望,一起去禱告.
去仰望你.
去做這麼重要的決定.
天父求主憐憫.
幫助我們.
主啊,讓嘉恩樓.
你的家在你的面前蒙恩典.
我們一起去建立.
你的家要一起去成長.
我們要一起學習.

$^{1241}$榮耀你的聖名.
願這個這麼重大的決定.
是一個榮耀你聖名的決定.
我們仰望你奉主耶穌.
基督的聖名,阿們.
謝謝大家.
\newpage



\section{馬太福音 18:1-9}
\label{sec:i3m_mXZnkC4}
\textbf{2023-08-06 加恩樓《天國的小孩》馬太福音 18\_1-9 講員:鄧泰康傳道}
\newline
\newline
連結: \href{https://youtube.com/watch?v=i3m-mXZnkC4}{\texttt{https://youtube.com/watch?v=i3m-mXZnkC4}} ~~~~ 語音日期: 2023-08-08
\newline
\newline
\hyperref[sec:qI_KwviwgzU]{\small{< < < PREV SERMON < < <}}
~
\hyperref[sec:index_chronic]{\small{[返順時目]}}
~
\hyperref[sec:index_scriptual]{\small{[返順卷目]}}
~
\hyperref[sec:1pEPdzdQ8_4]{\small{> > > NEXT SERMON > > >}}
\newline
\newline
馬太福音 18:1-9
\newline
\begin{longtable}{cl}
\hline
\hline
章節 & 經文 (和合本修訂版)\\
\hline
18:1 & \begin{tabularx}{0.7\textwidth}{X} 當時，門徒前來問耶穌：「天國裡誰是最大的？」 \end{tabularx} \\ \\ \relax
18:2 & \begin{tabularx}{0.7\textwidth}{X} 耶穌叫一個小孩子來，讓他站在他們當中， \end{tabularx} \\ \\ \relax
18:3 & \begin{tabularx}{0.7\textwidth}{X} 說：「我實在告訴你們，你們若不回轉，變成像小孩子一樣，絕不能進天國。 \end{tabularx} \\ \\ \relax
18:4 & \begin{tabularx}{0.7\textwidth}{X} 所以，凡自己謙卑像這小孩子的，他在天國裡就是最大的。 \end{tabularx} \\ \\ \relax
18:5 & \begin{tabularx}{0.7\textwidth}{X} 凡為我的名接納一個像這小孩子的，就是接納我。」 \end{tabularx} \\ \\ \relax
18:6 & \begin{tabularx}{0.7\textwidth}{X} 「凡使這些信我的小子中的一個跌倒的，倒不如把大磨石拴在這人的頸項上，沉在深海裡。 \end{tabularx} \\ \\ \relax
18:7 & \begin{tabularx}{0.7\textwidth}{X} 這世界有禍了，因為它使人跌倒；絆倒人的事是免不了的，但那絆倒人的有禍了！ \end{tabularx} \\ \\ \relax
18:8 & \begin{tabularx}{0.7\textwidth}{X} 如果你一隻手或是一隻腳使你跌倒，就把它砍下來扔掉。你缺一隻手或是一隻腳進入永生，比有兩手兩腳被扔進永火裡還好。 \end{tabularx} \\ \\ \relax
18:9 & \begin{tabularx}{0.7\textwidth}{X} 如果你一隻眼使你跌倒，就把它挖出來扔掉。你只有一隻眼進入永生，比有兩隻眼被扔進地獄的火裡還好。」 \end{tabularx} \\ \\
[1ex]
\hline
\hline
\end{longtable}
$^{1}$大家好,今天繼續跟大家分享一個我覺得是重要的訊息.
當我們一直說我們要發展,我們要向神的國前進的時候.
有些基石其實是非常重要的,要去定的.
這個時間跟大家講這段經文,就是大家一定要記得有一個前提.
為什麼這班門徒會去問耶穌天國裡面誰最大.
如果大家還記得,那些門徒曾經在聖經當中去爭辯.
將來的時間誰在耶穌的左邊和右邊,誰在裡面最大.
可能我不知道,聖經裡面沒有說,或者當時他們的爭辯就攤牌.
你看看我為耶穌的怨顧,我做了什麼,我放棄了那隻魚船.
然後有一個再大,我放棄了一把媽媽,或者是這樣子,我不知道.
總之他就說,我在神的國裡面,我應該是最大的,類似這樣.
但是你會發覺,耶穌的舉動很簡單,帶一個小朋友來到他面前.
天國裡面,誰最大?他最大.
為什麼這麼重要呢?當今天我們說我們要去建立神的國的時候.
首先大家就要搞清楚什麼叫神的國,神的國需要什麼.
很多時候我們很容易就會將屬世裡面的眼光和價值觀套在神的國裡面.
我們看到這個世界上,什麼叫做厲害?我不知道大家現在是不是還在流行,我沒有經常最新的最新的.
很多時候什麼叫做厲害的教會?就是ABC的教會叫做最厲害的.
什麼叫做ABC?attendance,很多人數,超級教會,過萬人.
B,building,我們一間很漂亮的,所以教會越建越高.
如果大家還記得,當時在溫州那裡曾經有過拆十字架事件.
然後你就會發覺,當你不看,一看,那些教堂是什麼?每一間都這麼厲害.
每一間都這麼厲害,C,cash,我們教會又很有錢.
這樣就叫做厲害的教會,所以教會應該追求ABC,是不是?聽完節目.
有一段日子,我想大家留意,連教會都有潮流.
有一段時間突然間講教會增長,然後就蜂擁的所有教會都一起去學教會增長.
看似是對的,但問題是慢慢慢慢就變成比賽.
你這間教會增長多少,他們的教會增長多少?.
很容易就陷入了這些熟悉的價值觀裡面.
耶穌在這裡的時候,為什麼他要把這個小孩子拿出來?.
因為他想告訴大家,我天國裡面所需要的是一個小孩子的樣式.
什麼叫做小孩子的樣式?不是今天你自己去套入,你認為小孩子是這樣.
然後我想大家回到聖經裡面,很習慣地,每一樣東西當聖經講,回到聖經講.
什麼叫做小孩子的樣式?就是上帝創造時候的樣式.
上帝創造人的時候的樣式,就是小孩子的樣式.
其實最初亞當夏娃被創造出來的時候,他們的生活是非常簡單的.
你留意一下,整個創世記當中,他寫得很清楚的.
全部的東西都是上帝做的,定位很清楚,上帝做,上帝定了所有的規矩.
地面就是地面,海裡就是海裡,天上就是天上.
然後創造人的時候,他告訴別人,你們要吃樹上的果子.

$^{41}$你教清楚,上帝命定了大家不是吃草的.
你們吃樹上的果子,某程度上,我想大家留意.
當上帝讓人去管理,是叫他管理動物,他沒有叫人管理植物.
意思是什麼?其實植物是代表上帝賜予,就是說人的生命是由上帝賜予.
不斷不斷的供給,我不需要你走去擔心.
大家都知道,如果植物在一個很正常的環境裡面的時候.
它會很周而復始的不斷的結果,對不對?.
它不需要你擔心,人不需要花點心力去做這件事.
然後他給了我一個工作的範圍,你幫我好好的管理動物.
大家也知道,當時的世界是否需要管理?.
我需要做一個羊欄嗎?不用的,在這個世界裡面,大家都相信.
所有的動物其實都是可以一起相處的,沒有所謂的空晚的野獸.
當時上帝的命定,原本是所有的動物都是吃草的.
就是說上帝創造了一個時候,當時的人是非常簡單.
我有上帝的供給,我滿足於上帝的供給.
我每天有上帝要我所做的事,我很有目標,每天都要做什麼事.
小孩子是一個這樣的狀態.
套回我們今天所見的小孩子.
坦白說,小孩子的滿足是什麼?.
你猜一下,在正常情況下,在你還沒跟他講錢有多重要之前.
就是說你還沒跟他賺錢去買玩具之前.
你給他一張100元,給他一包糖,他選哪樣?.
你帶他去吃一頓好的,明天我給你一千元,他選哪樣?.
我要吃一頓好的.
我要形容的就是,其實今天小孩子,我講得直接一點.
是被我們大人教壞的.
是我們的大人慢慢慢慢開始成長之後.
你開始滲入了很多人與人之間的比較.
你開始滲入了很多這個世界裡的條條框框.
你開始忘記了上帝創造你最初的原意.
你用你自己不斷去追求,不斷去追求,不斷去追求.
你認為這個世界上那些是最有價值的.
我不知道大家有沒有發覺,其實今天這些是屬靈爭戰.
我可以講得嚴重一點.
今天很多人沒有事情做,做什麼?.
香港不是最流行,但在大陸就很流行,就是叫做殺手機.
看什麼?看抖音.
一個一個不斷看,你不要看小哥那樣東西.
那裡就可以看你幾個小時.
我想大家留意,你的價值觀就是從那裡來的.

$^{81}$它慢慢就會去改變你的思想.
什麼叫做漂亮的標準?.
我想問問大家,到今天,有圖有沒有真相?.
我想告訴大家,你千萬不要再相信照片裡所謂的美女.
我絕對告訴大家,大家有沒有看到那些美顏的效果.
一個阿婆可以變成一個美女?.
大家知道吧?.
她告訴你什麼叫美的價值觀.
什麼叫有身份地位.
什麼叫賺到錢.
什麼叫吃喝玩樂,是一種享受.
你慢慢慢慢就被這個世界裡面不斷改變你.
再加上我們個人的經驗.
我們在這個世界上很想賺到別人的讚譽.
這個世界人對我的讚譽.
這個世界人對我的稱讚,對我的肯定.
我們慢慢用錢,用其他的物質,用所有東西去堆砌我們這個人出來.
然後你就教你的下一代.
所以你要好好讀書,為什麼?.
因為將來你才有出色的,你才能夠賺到錢.
你才能夠得到爸爸媽媽今天認為有價值的那個.
所以當我們的子女達到這個世界一個錢的標準的時候.
我們心滿意足.
真是聰明的兒子,真是聰明的女兒.
如果我們那個讀書不成的,那怎麼辦?.
一講你兒子的樣子,就不要提了.
其實我想大家回想起上帝所說的小孩子.
就是最初那一個完全依靠上帝滿足於上帝他自己的小孩子.
當今天我跟大家說,小朋友你看他多天真,他都不為明天憂慮.
接著大人就會回答什麼?.
當然不用了,他都不用為明天的餐,是吧?.
弟兄姊妹,今天當你忙碌生活的時候.
你信不信得過你的父母?.
你信不信得過你的天父?.
今天小朋友不擔心,因為他很相信爸爸媽媽是不會扔下我的.
爸爸媽媽是會令到我下一餐有飯吃的,對嗎?.
今天當你亡生亡死的時候,你有沒有相信你的天父.
同樣地,儘管今天我不知道會怎樣.
但是我的天父是會供給我每天所需的,你信不信?.
當你相信的時候,你就慢慢地能夠成為小孩子.

$^{121}$我可以很斬釘截鐵地告訴大家.
今天你所需要的遠遠比你認為的低.
今天只要你有一餐你喜歡吃就行了.
我經常跟別人說,鄧全道除了我現在吃farewell飯之外.
我是不會去吃自助餐的.
為什麼呢?自助餐有很多東西吃.
但是很多全部他們認為貴一點的東西是不關我的事.
那些鵝肝,我受不了.
那些魚生,我又不可以.
生蠔,我又不吃.
全部貴價不飽的,我都覺得不對.
我去夾火腿,第一輪出去的時候通常是拿西瓜.
我是這樣的,吃完一餐之後,平時那餐我幾十塊就搞定了.
但是我經常都跟別人說.
開始去吃自助餐,吃你喜歡的就行了,對嗎?.
我為什麼不喜歡的時候還要去夾羊,我不喜歡的吃,去讚自己.
這樣才能賺回來,沒有這個需要.
我媽媽永遠都有句名言,我寧願餓死都不要飽死.
飽是很辛苦的,我媽媽經常都這樣說.
所以弟兄姊妹,今天其實我們每一個人所需要的東西很少.
你晚上能夠有一個安睡,你就已經很好了,對嗎?.
你能夠有一個相對比較安穩的地方可以住,就可以了.
所以大家都聽過很多中國人的名言.
儘管你有千萬畝的田地和房產,但你需要睡的都是那張床.
所以我記得我以前曾經說過一個例子.
有一段時間,大家記不記得?.
在一個五星級的家,我還二十多歲的時候.
跟著國內賣的那些,對於香港人來說那些是神經病的.
我們住的那些只有三四百呎,但一去到之後.
整棟都給你了,那時候還很便宜,我忘記了是三十多萬還是四十多萬.
很多香港人就蜂擁去買,哇,真棒.
買完之後,旁邊的人說恭喜你,有間很大的房子.
但我很記得那時候我有個同學買了之後,我走去恭喜他.
他的樣子是燒焦的,看著我.
跟著時間說恭喜你,我說不是啊,買了房子還不開心.
開心啊,他每個星期的娛樂是什麼呢?.
就是逢星期六早上就出發,就去到國內.
之後穿好整套,然後就清潔整間房子.
清潔完之後剛剛好,第二天就坐車回香港.
我問他一年裡面住多少?.

$^{161}$他說一年回去兩次,帶朋友去威一陣玩兩次.
然後全年就靠自己走去做清潔.
弟兄姊妹,其實今天上帝由始至終都告訴大家.
你發覺今天很多有錢的人慢慢當他能夠賺到生活所需.
特別是一些很窮的地方出來的人.
他賺到很多生活所需之後,你發覺一個很奇怪的現象.
他還是寧願回到以前的農村裡面,去建一間屋子住.
你問他為什麼有時間?因為在這裡空氣好,靜.
但你要記得,你最初不用這麼努力,空氣好,你也是這麼靜.
所以你會發覺其實我們真正上帝想你的是什麼呢?.
你在他裡面的滿足和開心,因為唯有這個才不會變.
2017年有一個地方是全世界最快樂的地方,叫不丹.
經濟不好,經濟不是特別優秀,但他們的人很快樂.
成人很快樂,小朋友都很快樂.
但當你今天再搜尋不丹的時候,再不是一個快樂的地方.
因為資訊今天很發達.
他們就是從這些抖音,從這些裡面的時候,看到這個世界裡面.
原來這樣才叫好,這樣才叫好,這樣才叫好.
但每個人都很嚮往,全部衝出去這些地方賺錢.
但那個地方再不快樂,開始有很多比較.
所以我很記得我以前帶短孫的時候,帶得多的時候我就開始提大英姐妹.
以前我們香港人真的很有心,買Kitty的筆.
買老鼠的文具,然後送給他們,一個package,很漂亮的,紙很漂亮.
然後後來突然間,不要啊,你送文具給他,麻煩你送中華牌.
因為我們親眼見過有些現象,當你送的時候,你不會每個人都一樣.
送完之後,為什麼你的那支那麼漂亮,為什麼我沒有.
然後小朋友開始走出來了,我可不可以跟你換,我可不可以要這個.
開始有很多攀比,很多比較.
本來他們很快樂,本來他們有一支筆就可以了.
本來他們的文具就很開心.
但就因為他們見過.
大英姐妹,我不是說見過不應該,我想強調我們都需要有見識,我們都需要有進步.
所以我想強調的,當我今天講這篇道的時候,不是叫大家既然是這樣,我們躺平.
反而都是這樣,我們什麼都不用做,不用追求.
不是,我們要努力的,你努力是為上帝.
你的努力是為了建立上帝的國.
所以我今天早上在這篇道,我都在天崇那裡講.
我很怕那些年輕人聽完之後,說我們躺平.
我說不是,是你要為上帝的國的緣故,你要努力.
努力的不是拿好成績,努力的是學習.

$^{201}$讓你有更多的能力可以去建立上帝的國.
或者建立上帝的子民.
所以無論你是不是適合讀書,是另外一回事.
但上帝總是賜給你各方面的才能.
我經常說我就是人版,我在香港的教育政策下.
我是一個失敗者,我不是一個記性好的人.
我不是那種不斷的,可能我現在讀子民的學校.
我想說我可能會很厲害.
但我想說,到我後來讀神學的時候,其實我很沒信心.
上帝還跟人說,當上帝呼召我讀神學,我要報四間.
四間的原因就是怕他不收我.
然後上帝很感恩,三間都說收我,我竟然有得選擇.
然後我就開始讀,我一讀神學之後,就發覺原來是兩個世界.
上帝給我的才能不是死忌,上帝給我的才能是消化.
給我的才能是理解,給我的才能是寫論文.
所以我讀到最後的時候,我竟然可以拿獎學金去畢業.
上帝給我們的東西,總不是你所認為的,弟兄姊妹.
不是我厲害,而是上帝在過程中幫你放在一個應該的位置.
在你認為最沒信心,最沒自信的地方,上帝建立了你.
所以當有些人問我,我讀完四年神學之後,有什麼感想.
我只學會了一件事,原來我是這麼卑微.
原來只需要上帝放我在一個適合的位置,我就能夠發揮.
所以你就會發覺上帝要求你回轉,你回轉成一個小孩.
而回轉去做一個注釋,請你謙卑下來.
謙卑下來,你謙卑下來,你就可以成為小孩子.
什麼叫謙卑?你要記住,所有一切是各位上帝決定.
你今天所有的成績是祂給你的.
你今天所有的資產是祂給你的.
你今天所有的機遇是祂給你的.
你沒有什麼值得誇耀的,你要誇就誇上帝的大能.
你就能夠謙卑下來.
當你能夠不斷從心裡面,心甘情願地誇上帝的大能.
我可以告訴大家,這個世界上沒有什麼能夠再困住你.
因為你再不會跟別人攀比,你再不會比較.
你只會專心地在上帝的眼裡,問自己做好了沒有.
上帝你交給我的工作,我完成了沒有.
我不知道有多少大家姐妹,當我問你.
有多少大家姐妹很享受你今天的工作.
我不知道有多少人可以從心裡面回應,我說我很享受.
有些人說不用上班是最好的,早點退休是最好的.

$^{241}$大家姐妹,我想告訴大家姐妹.
當上帝給了一些恩賜給你的時候,你不要浪費.
當你能夠善用上帝給你的恩賜去做的時候.
那份再不是工作,那份是一份很有價值的人生目標.
你會做到死都不願意放手.
所以我可以肯定地告訴大家姐妹.
有一天當我65歲,我都要從這個崗位退下來.
但是教導別人的工作,沒人的工作,我是不會放下的.
只需要上帝繼續給我,可以去說話就可以了.
所以這次我不舒服的時候,我只是求上帝.
上帝你讓我繼續可以開車,開不到車我就哪裡都不用去了.
完全依賴上帝你去怎樣帶領我們.
所以他說你要回轉成為一個小孩子.
這個是你今天的定位.
接著他就告訴我們要接待小孩子.
這個小孩子裡面可以有兩個含義.
第一個含義是真實的含義.
就是這個世界上捍衛卑微.
如果大家還記得當時的片段.
曾經有些爸爸媽媽帶著小朋友來到耶穌面前.
你為他祝福,門徒是記住門徒.
他說你不要阻礙著耶穌工作.
耶穌是責備他們的.
然後有時間叫他們帶到你面前為他祝福.
為什麼這些門徒要破爛.
因為小孩子在他們的價值觀裡面是沒有價值.
沒有生產價值.
沒有更加重要的價值.
但是弟兄姊妹,在神的國裡面.
這些才是最好的價值.
但是在這個世界上很多人被人體會是沒有價值的.
體會是不入流的.
其實人是我們自己將他分為三六九等的.
所以很多時候,不知道大家有沒有見過.
我以前聽過一個最荒謬的.
當然是中學生,初中那些.
有一班女孩子在玩,不理一個玩.
接著我就問那班女孩子,為什麼不理她玩.
理由是什麼呢?因為她不漂亮.
我們這班人個個都很漂亮.

$^{281}$我們覺得她不漂亮.
所以我們覺得她拉低我們的水平.
我這樣都可以.
人是很喜歡自己走去分門別類的.
所以為什麼會有校園欺凌.
我走在一起的時間,這些就沒有價值了.
所以我踩他也可以,打他也可以,撕他也可以.
是人自己走去分類.
但是在神的裡面,上帝不是這樣分類.
特別今天我們做教會的時候.
我們更加要去接待,去愛這些在世界上看為沒有價值的.
很多時候,為什麼教會我們要帶大家一起去關心這些貧窮.
或者我們要去派飯.
原因不只是做慈善,你弄清楚,我們不是做慈善.
我們是要讓這些人知道,從我們的舉動裡面.
讓他們知道,這個世界上有人愛他們,有人關心他們.
而背後我們很希望最後他們能夠明白.
是上帝關心他們,上帝接納他們.
無論他們怎樣都好,我們都愛他們.
是這樣的意思.
所以不是我們每次在這裡拍完之後.
Oh, you know,這件事完了,做了慈善.
不是的,兄姐妹.
是你沒有真正從心裡面,當你派飯給他們的時候.
你的態度都不同.
我帶過一些人去派飯的時候.
有些人派飯就好像街邊派傳單一樣.
但我見過一個很令我感動的年輕人.
他走過去派飯的時候.
他蹲下半跪在他面前.
拿著飯,很想問他,你要不要飯.
我有機會和這個年輕人談的時候.
我剛才很欣賞你,他怎樣跟我說.
然後他就告訴我.
我爸爸也是一個街友,我找不到他.
當然他沒有再說原因.
我不知道他爸爸為什麼會成為一個街友.
但他說,剛才我給飯給這個人的時候.
我就好像面對著爸爸一樣.
他經歷過,他感受過,他知道很難.

$^{321}$我想強調的就是這種愛,這種同理心.
所以在真實的裡面.
我們教會需要去接納.
需要去接待,需要去愛這些被忽略.
被人看為沒有價值的人.
第二個在屬靈上.
屬靈上什麼叫小孩子.
就是初信的弟兄姊妹就是小孩子.
你70歲剛信耶穌,你都是小孩子.
我昨天在團體裡面才說.
我們每次有團體有新朋友.
我們今天歡迎新朋友.
但弟兄姊妹我不知道你說這句話的時候.
你知不知道代表著一樣東西.
就是說你是舊朋友.
弟兄姊妹,舊朋友的意思是什麼.
就是說你比他早信耶穌.
你比他早去經歷耶穌這份愛.
你是應該比他更成熟一點.
更懂得去怎樣走基督的路.
大家有沒有走過山.
行山都是操練.
你第一次操練可能會習慣了.
又走錯路.
慢慢你就會學會.
原來走這些隱蔽的地方.
你要怎樣開路,怎樣綁絲帶.
你慢慢學會了一些技能.
然後你就會帶著剛剛新手行山.
這裡不要走,會習慣死的.
你看著這些蝴蝶,這些絲帶.
原來就有出路.
你有經驗.
所以弟兄姊妹,當你慢慢成為一個舊朋友的時候.
你要記住,無論你想不想都好.
你都有責任帶著這些新朋友去成長.
在這些地方的時候你不需要謙卑.
你不需要跟我謙虛.
你不需要跟我說我這麼差.
不是.

$^{361}$因為你比他早經歷.
你只需要認真去信耶穌.
你認真去經歷.
你總是有更多更多的.
你可以告訴他.
是不是.
我們的生命永遠都不是.
全部充滿著那些驚濤駭浪見證的.
如果是你就大件事.
我真的恭喜你.
是不是.
每次都跟別人說我癌症好了.
我沒有工作,突然間老闆請我.
你不會的.
你的屬靈生命就是由你每一天裡面.
一點一滴去累積下來.
是上帝每天怎樣跟你走過這些路.
在一些困難的時候.
我怎樣陪伴你.
在我信你的時候.
我怎樣經歷上帝的愛.
我怎樣能夠經歷到上帝.
然後你就成為了新朋友的明燈.
你去好好去接待這些新朋友.
耶穌說當你接待他的時候.
你就等於接待我.
弟兄姊妹.
以接待耶穌的心態.
去接待這些新朋友.
但有一幅圖.
我很想給大家看.
是我一直都會保存在我自己的手機上.
我經常提醒自己.
前面是一間教會.
後面是一個懸崖.
很多時候當我們一不小心的時候.
請你留意.
我們很容易在前面帶人進入教會.
後面讓人跌進懸崖深淵.
為什麼我這樣說.

$^{401}$其實很簡單.
只需要你什麼都不做.
你做回你自己.
你做回你自己未信耶穌的時候.
你怎樣對人就怎樣對人.
你的世界價值觀是怎樣就怎樣.
你喜歡分黨分派.
你喜歡去講閒話.
你喜歡講是非.
你喜歡去埋怨.
你做回這些的時候.
然後那些人自然就會發覺.
原來你們的基督是這樣的.
原來教會是這樣的.
原來上帝是這樣的.
然後他就可以對上帝徹底失望.
原來這個世界上.
你說的上帝是假的.
所以弟兄姊妹.
為什麼這段經文後面說到.
你不要讓這些小子跌倒.
其實今天當我們一不小心的時候.
我很擔心教會成為了一個.
讓人離開上帝的地方.
當我們不斷去講這裡有上帝的愛.
我們很實際地去問.
我在哪裡見到上帝的愛.
最近我們做咖啡侍工的時候.
我們慢慢很感恩.
開始有些街坊.
我真的住了幾十年都沒有想過你.
現在這個地方很好.
有些人甚至上來visit.
看看你們在做什麼.
甚至昨天富哥和富嫂剛好在這裡.
我們開完咖啡侍工.
他們那對夫婦都很有趣.
每個男士都說.
我見你們初創.
所以來支持你們.

$^{441}$無論怎樣.
他們來到.
我經常說他們接觸到的.
可能是因為一杯咖啡.
但昨天我們的弟兄姊妹.
很友善,很親切.
他關心他們.
問他們住哪裡.
他們問你們住哪裡.
然後大家互相交流.
原來是街坊.
然後我見到富嫂和他們講養狗.
講家庭,其他.
慢慢地.
一杯咖啡他們可以坐下來.
可以開始和我們聊天.
其實這個故事在這兩個星期不斷發生.
雖然我不是經常都可以在教會.
我要出去工作.
但Fanny她在這裡留守.
她告訴我.
她說偶爾會有街坊走上來.
她有時間可以和他們聊.
有時間就沖一杯咖啡給他們.
然後和他們聊天.
我都感恩.
Fanny是一個很懂得和新朋友聊天的人.
她會和他們聊天.
講一些東西,講一些東西.
所以弟兄姊妹你會發覺.
究竟今天我們在這裡的時候.
被人看到的.
我們這杯咖啡和外面有什麼不同.
我們這裡的人和外面有什麼不同.
我不會期望第一次就.
弟兄姊妹也不要.
我記得有一次我在讀神學.
我去屯門醫院做實習.
跟著時間我就原本想學一下怎樣做軟木.
跟著那個軟木.

$^{481}$我不說是誰.
我記得他現在好像退休了.
Hallelujah他退休了.
原因是什麼呢.
因為我跟著他的時間.
做軟木很容易.
然後他就拿著.
當時我們學三幅.
大家知不知道三幅是要講九個字的呼音.
知道嗎?不知道嗎?.
其實是要講九個字的呼音.
然後他就帶著我拿著那本三幅.
我現在示範一次給你看.
他一賣到床.
第一件事.
先生你好啊.
你信耶穌了嗎?.
跟著時間那個人就開始.
轉過臉.
就是不看他了.
跟著他就跟他說.
我在這裡是軟木來的.
我現在跟你講耶穌.
跟著時間他就背著他的背.
腦後面.
跟著就在那裡跟他說了九個字.
跟著時間就問他.
那你信不信耶穌啊?.
大家都猜到了.
沒有反應.
當然沒有反應.
跟著時間.
不要緊.
下次我再來.
我只是主觀.
這是我主觀代入.
我就覺得那個人立刻臉會變青.
我不是要大家即時的.
可以讓那個人就信耶穌.
不是.

$^{521}$絕對不是.
我很相信.
他慢慢從我們的服事.
從我們的生命裡面.
他會看到.
什麼叫做被關心.
被接納.
被聆聽.
所以靈姐妹.
我很盼望我們今天在神的面前.
我們看上帝所看重.
今天這個世界上.
有很多東西是人所看重.
但上帝就是.
請你做回小孩子.
你看重我所看重.
上帝永遠都是看人為先的.
是不是?.
所以我經常說.
我不管家樓最後是怎樣.
我不管將來是怎樣.
但今天只需要上帝讓我們在這裡.
每天我們的工作要繼續.
只需要有一個機會.
我們都要繼續傳福音.
我們都要繼續去做.
因為上帝所看重人是最重要.
儘管我們在這裡很多時間.
我們這裡的人數不會很大幅增長.
但起碼曾經有人會談論.
這裡裡面有一間教會.
這裡的地方.
他們很關心人的.
他們不只是跟你喝咖啡這麼簡單.
原來他們會關心你的.
到最後的時候.
有一段經文就是.
這個世界有禍了.
因為將人叛賭.
叛賭的人的事是免不了的.

$^{561}$聖經裡面耶穌說得很清楚.
是用世界的價值觀.
這些讓人跌倒的事是免不了的.
是一定會發生的.
突然之間會.
就是說.
你不要對這個世界裡面有任何期望.
這個世界只需要你今天很認真的跟隨上帝的時候.
就必然會和它產生衝突.
所以為什麼後來他要用到那些字眼.
就是當年那些人的手和腳.
令你跌倒的時候斬掉.
你的眼睛令你跌倒的時候劃掉.
為什麼要去到這麼絕.
意思就是要和這個世界的價值觀.
完完全全清楚劃清界線.
不要和流合污.
因為你跟這個世界走的時候.
最後的結果就是叛賭這些小孩子.
所以上帝很想你搞清楚.
他就說寧願你殘缺上天堂.
都比你完完整整落地好.
弟兄姊妹只需要你很認真做基督徒.
你和這個世界的價值觀一定有衝撞.
一定有衝撞.
你一定變得殘缺.
我舉個例子.
我聽過有些弟兄姊妹跟我分享.
做會計.
做會計的時候.
其實做會計不是計數.
是做數.
接著老闆就跟你說.
你幫我做這盤數.
做這盤數出來幹什麼.
少交稅.
可以有些比較好的運作.
但倒過來去到某些時候.
你知道我不應該這樣做.
聖經教我是.

$^{601}$事就說事.
不事就說不事.
但如果你按照這個世界價值觀的時候.
你老闆就會稱讚你.
你聰明.
你做得好.
但你按照神的價值觀去做的時候.
你不行.
你不可以.
老闆會覺得你不行.
我稱之為.
這個就是在這個世界裡面的殘缺.
再沒有價值.
你明白我意思嗎.
就好像這個世界上.
很多傷殘人士被人看不起一樣.
但弟兄姊妹.
聖經裡面就是想告訴你.
如果你不跟這個世界的價值觀去完全割斷的話.
你斷不能進天國.
因為.
在神的國裡面的時候.
他容納不了這些事.
他要你一個完全的改變.
你才能夠可以去佔領天國.
在最近的時候.
我曾經看到.
在網上有些這樣.
他經常說.
基督教是什麼呢.
基督教是一個騙人的信仰.
為什麼呢.
最初你跟每一個未信仰的人說.
你只需要舉手決志.
你就上天堂了.
是不是.
跟著時間.
但不是的.
當他信了耶穌之後.
你去讀經.

$^{641}$你要事奉.
你要什麼.
弟兄姊妹.
所以當我們今天傳福音的時候.
請你小心.
不要跟別人傳一個簡單的福音.
你要告訴他.
為什麼令我們上天堂的原因.
是因為我們.
上帝改變了我們的價值觀.
我想告訴弟兄姊妹.
今天你很努力的.
去賺取這個世界所有的東西.
我想強調.
我不是叫大家不要努力.
你明白我的意思嗎.
但我想說.
你今天所賺到的任何的東西.
到最後的時候.
你都要放下在這個世界.
你都要赤裸裸的去面對上帝.
你終極的價值都在天國上.
你今天當你面對上帝的時候.
你憑什麼.
但他就說得很清楚給你聽.
當你不能夠回轉成為一個小孩子.
你斷不能進天國.
所以弟兄姊妹.
上帝要你的.
是改變你的生命.
當你不能夠改變你的生命的時候.
你所教的新朋友.
就會像你一樣.
所以坦白說.
我不知道大家的觀感是怎樣.
自從改了讓小朋友和我們一起敬拜之後.
我真的滿心歡喜.
我是很快樂的.
因為小朋友真的在裡面的時候.
看著我們大人我們去敬拜.

$^{681}$所以我經常提醒大人.
麻煩大家敬拜的時候專心一點好不好.
小朋友在看著你.
你現在.
你看到他的時候.
我敬拜的時候.
將來你看到你的兒子就是這樣.
讓他們在這裡一起成長的裡面.
他就一起將我們最好的東西告訴他.
所以弟兄姊妹.
同樣的.
你自己的子女.
今天你是一個怎樣的基督徒父母.
你的子女就是一個怎樣的基督徒.
你今天作為教會的前輩.
你是舊的朋友.
新朋友.
就是會變成怎樣.
所以求神幫助我們好不好.
讓我們能夠帶著這一種榮耀.
能夠帶著這一種上帝所喜悅.
去進入神的國.
我經常很渴望家樓這個地方.
能夠弟兄姊妹真的能夠經歷.
在地如在天.
當我這樣說的時候.
鄧全道,你真是太妄想了.
不是的,弟兄姊妹.
其實很容易的.
只需要你真的能夠改變.
你完完全全看上帝是最重要.
你相信上帝的公認.
你的生命就可以在地如在天.
你就能夠經歷在天上.
上帝與我們同在.
所以我今天跟年青人說.
你連在地上都那麼討厭上帝.
你上到天上就大件事了.
你天天都對著他.
在地上的時候.

$^{721}$如果你有多喜愛上帝.
當有一天的時候.
你就越去渴望上帝.
上帝原來如此信實.
所以我最後的時候也這麼說.
為什麼有時候上帝偶爾要給我們一些困難.
偶爾他要讓我們看為重要的東西.
我們暫時要拿走.
是因為他想你重新再去看.
什麼才是你的生命.
你的人生裡面最重要的東西.
如果你現在問鄧全道.
最重要的是什麼?.
健康.
我現在經常提醒自己.
大家看到我都提醒我.
雖然我也覺得很慘.
就是不要翹腿.
一翹腿的時候.
其實我也挺相信.
是我生活習慣令我今天這樣.
看手機的頸部.
然後我的腳的時候.
翹腿或者坐得不好.
坐得不好.
或者坐幾個小時又不起床.
又不喝水.
各樣東西.
醫生現在告訴你.
你就是這樣.
但你就會發覺一件事.
健康原來最重要.
當你面對困難的時候.
原來健康最重要.
有什麼可以換回健康?.
沒有.
我現在只可以說.
我現在的狀態是全然倚賴上帝.
有些突然之間都問我.
你現在做什麼?.

$^{761}$針灸.
做完針灸那怎麼辦?.
沒有事情可以做.
有很多突然之間會介紹其他的.
但問題就是.
其實繼續仰賴上帝.
我想大家留意.
當你這樣的時候.
儘管你在這樣的狀態當中.
你也不會覺得會失望.
因為上帝從來無變.
上帝對你的愛從來無變.
你也相信上帝今天.
放你在這個境地當中.
有祂美善的原因.
靈姐妹.
求神幫助我們.
我們一起成為天國小孩的樣式.
好嗎?.
一起同心協同.
差天父.
我們求主幫助.
帶領我們.
主人,我們真的很盼望.
我們能夠成為上帝的喜悅.
我們真的很盼望.
我們能夠按照聖經所說.
我們能夠成為小孩的樣式.
我們今天作為成年人.
我們經常想這些的時候.
我們會覺得很難.
這個世界是這樣.
或者我們已經習慣了.
但不是的,上帝.
你告訴我們.
這些才是我們出路.
你的說話從來都無變.
你從來都是叫我們活出.
屬於上帝的生命.
所以我求主幫助.

$^{801}$讓我們不要再考慮.
讓我們不要再去疑惑.
讓我們一聽到你的呼召的時候.
我們即時願意去作出一些行動悔改.
讓我們的生命能夠更加簡單.
像小孩.
以至我們能夠在神的面前.
能夠活得滿有盼望.
滿有出路.
我們多謝贊美利.
其中一方教主耶穌.
得勝名至其球.
阿們.
\newpage



\section{歌羅西書 4:5-6}
\label{sec:1pEPdzdQ8_4}
\textbf{2023-08-06 午堂《著地氣的傳福音》歌羅西書 4\_5-6 講員:曾裔貴牧師}
\newline
\newline
連結: \href{https://youtube.com/watch?v=1pEPdzdQ8_4}{\texttt{https://youtube.com/watch?v=1pEPdzdQ8\_4}} ~~~~ 語音日期: 2023-08-08
\newline
\newline
\hyperref[sec:i3m_mXZnkC4]{\small{< < < PREV SERMON < < <}}
~
\hyperref[sec:index_chronic]{\small{[返順時目]}}
~
\hyperref[sec:index_scriptual]{\small{[返順卷目]}}
~
\hyperref[sec:FlUdBHop01k]{\small{> > > NEXT SERMON > > >}}
\newline
\newline
歌羅西書 4:5-6
\newline
\begin{longtable}{cl}
\hline
\hline
章節 & 經文 (和合本修訂版)\\
\hline
4:5 & \begin{tabularx}{0.7\textwidth}{X} 你們要把握時機，用智慧與外人來往。 \end{tabularx} \\ \\ \relax
4:6 & \begin{tabularx}{0.7\textwidth}{X} 你們的言談要時常帶著溫和，好像用鹽調味，讓你們知道該怎樣應對每一個人。 \end{tabularx} \\ \\
[1ex]
\hline
\hline
\end{longtable}
$^{1}$各位弟兄姊妹 早安.
看一看 有些新朋友在我們中間.
神的殿裡敬拜 我們一起有祈禱.
我們打開聖經 盼望.
舍圖保羅對於弟兄姊妹.
哥羅西教會的勸勉.
成為我們大家.
真的能夠活出這個救恩.
帶給我們的信心.
帶給我們自己有的盼望.
可以跟其他人分享.
我們仰望等候.
奉耶穌基督的命 阿們.
保羅單純是使徒.
他過去未信主.
他是一個拉比.
在猶太教裡有相當的學習.
地位.
所以導致他要迫害.
一些跟隨耶穌的人.
提摩太前書第一章第十二節.
這一位主耶穌基督.
因他以我有信心.
派我去服侍他.
但是他說提摩太.
你知不知道我的從前是什麼.
我是一個蝕毒神.
迫不得人和沒有萬人的人.
即是說他以前.
除了他以為的宗教之外.
不會再有任何的信仰.
所以不是猶太教的人.
就是異端邪說.
就是沉淪滅亡.
即是說猶太人以外的沒有救恩.
所以他會沒有萬人.
即是說用很難聽的話.
來對這些外邦人踐踏.
用言語.
所以如果這兩節聖經.

$^{41}$保羅是一個屬靈的言語治療師.
教我們怎樣說話.
整個生命因為受著福音的影響.
而對外人要怎樣去言說呢.
第五節.
尤其是強調.
愛人這個概念.
我們就要集中在這裡去看.
因為在整個歌羅西書的上文.
由三章的十五節開始.
就已經講到教會生活的美滿一面.
十五節他說要基督的平安.
在你們的心裡面作主.
你們都是這樣去蒙召.
歸為一個群體.
在教會裡面應該怎樣彼此去服事呢.
用什麼心靈態度去服事呢.
當用各樣的智慧.
把基督的道理.
豐豐富富的存在心裡.
這個不只是背熟聖經那麼簡單.
把神的道.
價值觀.
屬天的屬靈的價值觀.
來不斷地學習.
讓自己能夠.
不只是今生世俗的價值觀.
同時間擁有對屬靈事物的視野.
那個價值觀.
所以就叫做豐豐富富.
把神的道.
存在心靈裡面.
然後用詩章.
頌詞.
靈歌.
彼此的教導.
互相的勸誡.
心被恩感.
恩典感動.
來歌頌神.

$^{81}$相當明顯.
這個就是教會崇拜裡面.
恆常彼此的服事.
敬拜神.
互相的勸勉.
互相的建立.
用.
這裡提到特別是詩歌.
特別觸動我們心靈的.
盛載.
救恩.
神的豐厚恩典的.
那些內容.
使我們能夠認識.
我們自己的生命.
第十七節.
無論作甚麼.
說話行事.
都要奉主耶穌的名.
開啟了一個很強大的視野.
就是我們在神的國度.
這個世界不是物質世界那麼簡單.
這個世界以外.
看不到的屬靈世界.
是徹徹底底的存在.
我們就是那班.
可以來接觸靈界.
屬靈世界的一班基督徒.
我們可以凡事.
說話行事.
我們要在主耶穌的同在下.
來進行.
來做.
即是說.
我們的主是與我們同在的.
我們要尋求的是祂的喜悅.
因為祂是我們生命的主.
基督就是我們的生命.
祂顯現的時候.
我們都要顯現在榮耀裡面.

$^{121}$這是歌羅西書在前面已經說了.
在這裡告訴我們.
我們是有一神的信仰的一班弟兄姊妹.
接著就開始說.
離開教會之後.
我們的生活應該是怎樣呢.
所以就開始涉及家庭的倫理.
職場的倫理.
我們應該怎樣在神安排我們不同的人生崗位.
來做好我們的本份呢.
所以就馬上說到.
作妻子的要馴服自己的丈夫.
在主裡面.
即是說家庭關係.
涉及的不只是兩個人.
是有神介入的夫婦關係.
在主裡面.
作丈夫的要愛你們的妻子.
不可以輔待他們.
當時的希臘人.
斯巴達人.
非常強悍.
你看戰狼三八就知道了.
斯巴達人是戰士.
如果在外面是戰士.
回到家裡可能會很恐怖.
信了主耶穌.
就要學習愛妻子.
你可以想像到.
保羅好像碰到你的家庭.
你的職場的倫理.
都和他有關.
因為保羅是使徒.
他要明白.
你要知道.
你信了主之後.
不只是教會生活.
我可以干涉.
我可以勉勵.
而是你的家庭.

$^{161}$你的職場.
我都應該來教導的.
這樣的一個信仰生活.
然後有兒女.
就要凡事在主裡面的喜悅下.
聽從父母.
作父親的.
就不要惹兒女的氣.
免得他們失去志氣.
然後再說.
在職場的文化下.
作僕人的.
要凡事聽從肉身的主人.
不可以只是在眼前.
巴結 奉承.
好像土人喜歡這樣.
總要全心誠實.
敬畏主.
這是一個很強大的座右銘.
任何的職場都適合用.
我凡事全心誠實.
是因為敬畏主.
我要來做的每一個工作.
很重要.
帶來的是.
你只會令到不同的同事.
對你都會慢慢.
建立了一份信任.
因為你這個人.
原來怕一個就是神.
這樣的僱員.
我們要盡快請.
他一定不會.
背後有很多不好的手腳.
因為他懼怕神.
全心誠實的敬畏主.
教一個奴隸.
要這樣做奴隸.
然後作主人的.
你就可以想像到.

$^{201}$一樣的道理.
第23節.
無論作什麼.
都要從心裡作.
像是給主作的.
不是給人作的.
明顯是給人作的.
見得到的.
但連我們那個動機.
保羅都來到去教導.
你的動機對.
你做出來的事情.
總是會對.
你可以想像到.
這裡是說一個基督徒.
因為在一神信仰底下.
對做物主的敬畏.
對神已經揭示的奧秘認識之後.
你應該怎樣來到在家庭.
在生活.
在教會.
來到有不同的表現.
而且一切的大前提.
是對主耶穌的愛.
信服和敬畏.
這樣的底下.
因為下面就說.
我們知道將來會得著基業.
作為一個賞賜.
因為你們所事奉的是主基督.
那些行不義的.
必定會受到不義的報應.
主不會偏待人.
就是說.
祂是看著的.
祂會知道哪個人.
你的人生.
如果用得不正確.
你要面對將來的報應.
這個是對於初代教會.

$^{241}$真的沒有什麼的力量.
抗衡.
猶太教的逼迫.
看你是異端.
羅馬政權.
羅馬的皇帝.
他們看自己是神.
羅馬的政權.
對於這些基督教.
當時是迫害的.
兩邊都受著壓力.
但是.
這一班基督徒.
得到的是正確的屬靈價值觀.
只需要認定主是我們的基督.
你要在每方每面.
盡力去做好.
我們不是.
任何一個人都不是屬靈的超人.
昨天周牧師說了.
我們每個人都是蒙恩典.
經常蒙恩典.
繼續蒙恩典的罪人.
明白自己的身份和內心.
我們就有一個很大的力量.
面對周遭的環境.
這個是在教會裡面的弟兄姊妹.
要這樣對自己的身份.
有這份信心.
很快就說到四章一節.
也是跟上文.
很公平的.
說了僕人之後.
就說了主人.
要公公平平的代僕人.
又再一次帶我們去看.
你在神眼中的.
怎樣做一個僱主.
因為知道你們又有一位主在天上.
不斷告訴我們.

$^{281}$主耶穌在天上.
看著每一個倫理關係.
看著我們每一個弟兄姊妹.
這樣是一個很大的信心.
我們要建立這種信心.
接著第二節.
保羅就開始說他們使徒這班人.
他們現在面對的困難.
但是仍然有機遇的.
因為他知道神的道.
是不會被任何的人.
任何的環境捆綁的.
我們要做的就是.
你們要行切的禱告.
在此警醒感恩.
也要為我們禱告.
求神給我們開傳道的門.
能夠說基督的奧秘.
保羅就強調.
我現在在這個坐監的狀態.
被捆綁.
但是仍然要祈禱.
有一個時機.
可以將神的道去講解.
這是基督的奧秘.
神的奧秘.
其實就是福音.
就是我們的主耶穌.
第二章有這樣的經文.
他說我願意你們了解.
我為你們和老弟家教會的人.
和一切我沒有見面的人.
是多麼的盡心竭力.
要叫他們的心得到安慰.
因愛心互相的聯絡.
接著留意.
他說以至豐豐足足的.
在悟性當中有充足的信心.
使他們真知道神的奧秘.
什麼是神的奧秘呢.

$^{321}$就是基督.
所積蓄的一切智慧知識.
都在他的裡面藏著.
各位這個就是保羅勸勉.
哥羅西教會.
不要受那些異端.
在當時對他們的影響.
有些人教他們要敬拜天使.
其他的一些手.
猶太教的規條.
這一切要用的時候.
是不會有功效的.
就是說要帶給哥羅西教會.
有一份在耶穌裡面.
認定我們認識了這個奧秘.
在耶穌裡面.
我們要建立這一份信心.
以至我們就能夠在教會裡面.
繼續彼此的扶持.
彼此的去勸勉.
使我們的靈命能夠成長.
以至我們在不同的人生.
不同的倫理關係上.
我們作為丈夫.
父親.
教會的僱員.
服事主的一個牧者.
和在其他的朋友當中.
我們怎樣去行動.
我們自己的品格.
這個就是保羅很緊張.
哥羅西的教會信徒.
不需要被異端邪說影響.
好好地踏實地.
對基督這個奧秘去認識.
建立這樣的信心.
所以去到第三節.
就說這個奧秘就是基督.
要我們傳遞.
當然保羅有他們的職份.

$^{361}$他們作為使徒的職份.
他們要將這個奧秘講解清楚明白.
譬如他繼續要寫書信.
提波太前後書都未寫.
他要繼續傳遞這個真理.
使更多地域的人.
可以明白神的救恩.
還有很多猶太人需要明白.
彌賽亞已經來了.
外邦人需要明白.
需要有一位對真神的敬拜.
就是耶穌基督.
最重要來到第五第六節.
就是和弟兄姊妹.
哥羅西教會的弟兄姊妹.
來去講.
要抓緊時機.
買熟那個時間.
來和外人交往.
這是一個吩咐.
是一個福氣.
我們基督徒是有的信心.
帶給其他人.
得到救恩的唯一的管道.
因為去到第七節.
一直去到尾.
保羅就立即在近尾的時候.
將他不同的年青的傳道同工.
逐個逐個來到去.
當然不是所有都列舉了.
例如列舉不同的傳道同工.
他們現在在不同的地方.
不同的遭遇.
不同的事奉.
來到去講出來.
但是第五第六節.
是唯一要求弟兄姊妹.
要來去怎樣把握你們自己的時間.
盡量令到外人.
來去得著福音恩典.

$^{401}$用什麼態度.
用什麼方式.
來展示這個福音.
所以這個涉及的就是.
我們對救恩.
基督的奧秘.
有多少的掌握.
我們的內心有多少的自信.
就視乎我們會不會和人去分享福音.
這個就是一個很實在.
反映到我們內心.
有沒有這份自信的問題.
第五第六節.
你們要愛惜光陰.
用智慧與外人交往.
你們的言語.
要常常帶著和氣.
在這個交往的過程.
我們應該用什麼態度呢.
好像用言去調和.
我們的說話.
要令到人感受到一份趣味.
感受到我們有一份熱誠.
我們的內心.
甚至感受到我們的動機.
就更加美麗了.
以致我們能夠可以.
這裡告訴我們.
就可以知道.
該怎樣去回答.
這些的國人.
即是說外人.
保羅做了一個相當清楚的分野.
就是只有在耶穌認識.
這個奧秘的基督徒.
和不是在耶穌的救恩裡面的外人.
外人是得不到救恩.
未得到救恩.
需要我們用我們的言語.
帶領他們明白這個基督的奧秘.

$^{441}$所以神已經透過保羅.
將這個寶貴的職份.
分給我們每一個人.
保羅有他自己的神安排的職份.
但每一個基督徒.
都有我們自己人生的職份.
神的吩咐.
我們要透過我們的言說.
去帶領人認識主耶穌的救恩.
經文的分段.
就是這樣將這兩節.
和上下文做了一個對比.
即是告訴我們.
教會的生活弟兄姊妹.
享受完與神與人的敬拜相交之後.
我們就是應該要走到外面.
看到有很多的外人.
未有福音.
不明白基督的奧秘.
沒有神國度的倫理價值觀的人.
他見到的就是他能夠見到的世界的現況.
看不到這個世界以外有真神.
我們基督徒就是應該看到.
同時間我們就是應該要帶領人.
認識主的救恩.
這裡說的是保羅.
好像現在告訴我們.
他彷彿變成了一個.
屬靈的言語治療師.
從前他是沒有問.
不信猶太教的人.
但他現在因為看到.
基督耶穌那份的恩典.
和對他的肯定.
他自己有很大的生命改變.
他願意和外邦人分享.
這個救恩的榮耀.
即是說保羅自己有很大的信心.
原來有時說話.
如果你咬字不是很清楚.

$^{481}$是會涉及到你可能會沒有信心.
有時被人笑.
如果你是嚴重的口吃.
需要有言語治療師.
有個同學的牙齒不太好.
在五月的時候.
我和他一起.
我們幾個同學上去.
做了那些絞刑的牙齒之後.
他說話真的變得很正.
他說現在很有信心.
說到不怕被人笑.
原來真的會有幫助.
大家有時聽廣告.
現在就比較少.
一說話就會笑.
有空就整理牙齒.
下句是甚麼.
過關都不怕檢查.
一聽就覺得很全神.
好像令到你突然很有信心.
整理牙齒之後.
保羅說的不單純是物質.
我們的真正的口齒靈力.
保羅說的.
我們有沒有這份自信.
去傳揚救恩.
帶給人生命的改變.
我剛才有說過這個例子.
在晚堂.
我們有一對夫婦.
每個星期日的早上.
因為她要上班.
每個星期早上.
她都會去新元朗中心吃早餐.
差不多每個星期都是這樣的時間.
認識了附近幾台.
好像是街坊.
其中有一對姊妹.
一個90歲一個91歲.

$^{521}$很好談.
這對夫婦.
其實很多人都認識的.
不過不說名字.
帶領他們兩位信主.
上星期邀請我去跟他們聊天.
信得很單純.
很奇妙地.
原來有些人這麼容易就信主.
但是.
你試一下回想一下.
如果弟兄姊妹只是吃東西.
跟隔壁沒有什麼可以談的.
河水井水不乾飯的.
吃完就走了.
你就沒有了這些契機.
找這些機會.
可以帶領到這對姊妹信主.
其中一位.
對不起是兩位姊妹.
是我們五堂一位姊妹的中學同學.
說的是六十多年前.
所以遲些我們在洗禮.
都會邀請這位姊妹來出席.
你可以想像得到這個世界原來這麼奇妙的引導.
有一位傳道人.
在美國.
他寫了一個分享.
他說有一天.
要去芝加哥.
來到港島.
去領會.
就上機了.
隔壁有一個婦女.
中間是沒有人坐的.
她就坐在窗口的位置.
有感動.
叫她要跟這個婦女說福音.
她就一直都不敢.
輾轉掙扎.

$^{561}$最終都沒有說.
於是就去領會.
接著都有點懊悔.
為什麼不找這個機會.
來跟同機的乘客.
這個婦人來講福音.
但是.
很奇妙地.
她回程回家的時候.
坐的班機.
這個婦人跟她同一班機.
更加奇妙地.
坐在隔壁的位置.
真的沒有可能的事.
那個婦人就說.
為什麼會這麼巧合.
傳道人就用這個巧合.
成為一個港島的題目.
跟她講.
在神裡面.
不會是巧合的.
就帶領這個婦人.
跟她在全程裡面.
好像一個贖罪一樣.
我現在不講福音就不行了.
真的要跟她講福音.
而這個婦人.
因為這個巧合.
就開放自己的心靈.
去講她自己內心.
和家庭的需要.
這樣就變成一個.
神安排的契機.
親愛的弟兄姊妹.
我們需要有這份.
屬靈的自信.
基督就是我們的生命.
他顯現的時候.
我們都會顯現在榮耀裡面.
有一個台灣.

$^{601}$其實很大的一個新聞.
應該是2003年的時間.
因為我追蹤那個新聞.
最近再看.
可能大家沒有機會留意到.
有一個姓謝的叫謝依霞.
一個年輕人.
她是一個咖啡店的店長.
為什麼成為新聞人物呢.
這個人原來很壞.
有一個常客.
時時去喝咖啡的老夫婦.
退休了.
很有錢.
這個謝依霞.
謀財害命.
將她兩個落藥.
在咖啡裡面毒死之後.
就將她棄屍在荒郊野嶺.
接著呢.
接著報警.
發現屍體.
自編自導自演.
很難想像一個女孩子.
會有這樣的惡毒的心腸.
當然輾轉揭發了.
知道她拿了別人的錢.
還要在上法庭的時候.
自己作一個故事.
說是這個陳俊福.
這個退休的男士.
經常來咖啡店勾引她.
用錢要和她發生關係.
她沒辦法.
她給這麼多錢.
忍受不了那個誘惑.
然後被太太發現.
要斷絕關係.
諸如此類.
作個故事.

$^{641}$在法庭裡面.
最終判死刑.
後來.
有傳道人.
去到監獄裡.
帶她去信主.
輾轉開道她.
說她自己成長的故事.
輾轉.
這個傳道人真的帶了她洗禮.
接著我不知道台灣是不是有這樣的法例.
因為這樣的緣故.
她的死刑現在改判.
一個25年的終身監禁.
至少要25年才可以出監.
原來是她的媽媽.
在很遠的地方搬來和她住.
她已經在監獄裡.
不斷去探望她.
和支持她.
再叫教會的牧者來探望她.
使她的生命在監獄裡有一個徹底的改變.
親愛的弟兄姊妹.
福音對人的生命改變的能力.
是很大的.
所以我們求神幫助我們.
我們的言語.
要帶著的是恩典.
好像嚴調和的.
讓人感受到我們那個熱誠.
剛才早堂我也有講一個例子.
在在有些弟兄我也見過.
我就喜歡泊車很遠的.
在元朗.
爛地近馬田路口.
很遠.
因為方便我走進市區.
可以出汗做運動.
特意走遠一點.
藉這些機會.

$^{681}$不知為何.
耕田的那個男孩.
我們叫他肥仔.
擔著一支煙.
牧師.
他知道我是牧師.
其實我沒有告訴他.
有機會聊聊天.
有時過時過節.
給他一個月餅.
他們很開心.
不知為何.
可以有機會.
分享到信仰.
弟兄姊妹我們能不能夠讓.
愛人.
記住是沒有神國度救恩的人.
是不信就注定沉淪的人.
他們能夠透過我們的言說.
我們的信心.
可以帶到他們能夠接受救恩.
保羅說.
他是一個屬靈的言語治療師.
保羅說.
你們的言語.
要常常帶著恩典.
和好像用鹽調和.
你就可以知道.
怎樣去回答各人.
很盼望我們能夠真正在神的國度裡面.
我們有份.
剛才都說了.
在教會裡面.
盡量享受與神.
與人的相交.
和敬拜.
這是一個力量.
使我們能夠走出去.
我們是神國度的人.
我們見到的是沒有救恩的人.

$^{721}$怎樣幫他們去認識主耶穌的愛呢.
求主幫助我們每一位.
我們一起祈禱.
多謝主.
讓我們能夠在神的國度裡面.
我們可以來到去透過詩歌.
一起的敬拜你.
你是屬天的主宰.
掌管萬有的真神.
你已經將你自己的救贖計劃.
藉著耶穌基督.
已經來到去揭示出來.
使我們知道.
這個神國度的奧秘.
多謝主將我們拯救.
盼望我們成為別人的祝福.
成為別人的一個治療師.
帶領他們認識神的恩典.
神的大愛.
我們這樣仰望感恩禱告.
奉耶穌基督的名.
阿們.
\newpage



\section{希伯來書 3:13-14}
\label{sec:FlUdBHop01k}
\textbf{2023-07-30 午堂《信仰的初心》希伯來書 3\_13-14 講員:陳韋安牧師}
\newline
\newline
連結: \href{https://youtube.com/watch?v=FlUdBHop01k}{\texttt{https://youtube.com/watch?v=FlUdBHop01k}} ~~~~ 語音日期: 2023-08-09
\newline
\newline
\hyperref[sec:1pEPdzdQ8_4]{\small{< < < PREV SERMON < < <}}
~
\hyperref[sec:index_chronic]{\small{[返順時目]}}
~
\hyperref[sec:index_scriptual]{\small{[返順卷目]}}
~
\hyperref[sec:RptdKZSozLI]{\small{> > > NEXT SERMON > > >}}
\newline
\newline
希伯來書 3:13-14
\newline
\begin{longtable}{cl}
\hline
\hline
章節 & 經文 (和合本修訂版)\\
\hline
3:13 & \begin{tabularx}{0.7\textwidth}{X} 總要趁著還有今日，天天彼此相勸，免得你們中間有人被罪迷惑，心腸剛硬了。 \end{tabularx} \\ \\ \relax
3:14 & \begin{tabularx}{0.7\textwidth}{X} 只要我們將起初確實的信心堅持到底，就在基督裡有份了。 \end{tabularx} \\ \\
[1ex]
\hline
\hline
\end{longtable}
$^{1}$祝求你賜給我們說話.
將聖經的說話展現.
讓我們能夠清楚明白.
我們的悟性,我們的靈性.
都能夠完全地掌握.
我們生命的實踐.
你給我們的說話.
我們去反省我們和你的關係.
我們整段信仰的路程.
求主你帶領我們.
幫助孩子能夠清楚說出你的說話.
共主明求 阿們.
我們會說有關初心的題目.
不知道大家有沒有發現.
聖經也不多說初心的題目.
我們最多找到.
耶穌叫我們做回小孩子.
讓我們做回天國裡面的小孩.
不過我們發現.
聖經裡面不少經文.
都是反對我們做小朋友.
保羅的書中也有很多說到.
我們不要再做心智上的小孩.
我們要在心智裡面做大人.
在這本書裡說.
要認識神的兒子.
得以長大成人.
沒有多少的新娘.
使我們不再做小孩.
所以我想大家都是大人.
都是一些年紀或經驗都很豐富的成年人.
甚至乎是一些長輩.
我們回到信仰的起點.
去思考信仰初心.
和我們這一刻的生命的境況和我們的關係.
經文在希伯來書第三章十四節.
和本是這樣翻譯的.
我們若將起初確實的信心堅持到底.
就是基督裡有分了.
我們要先明白經文裡面的意思.

$^{41}$和本就翻譯成起初確實的信心.
其實如果你看原文的時候.
你會發現原文裡面其實是沒有信心這個字的.
Sorry 其實是沒有信心這個字的.
這裡都不是和本的問題.
實際上這本書是一個比較難去翻譯的一段經文.
我們很難將原文的字.
去到很準確變成中文.
如果你看原文的時候.
其實那個字不是一個形容詞加名詞.
不是確實的信心.
其實只有一個字很簡單.
就叫做Hypothesis.
不如再說一次.
我玩到叫做Hypothesis.
再說一次好嗎.
1 2 3.
Hypothesis是一個非常複雜少有的字.
這個字是為什麼呢.
其實就不是信心這個字.
信心是Pistes.
希伯來書都講信心.
特別第三章都講過幾次有關信心的經文.
第二節有信字.
第五節有信字.
十二節有不信這個字.
不過十四節就不是關於信心這個字.
不過都明白的.
因為原文裡面的整段經文都是關乎於我們的信仰.
信心的事情.
所以經文裡面就翻譯成確實的信心.
什麼叫做Hypothesis呢.
這個字其實讀神學的話都大概聽過的.
這個有關三不論的字.
東方教會就會將上帝三位一體.
就叫做三個Hypothesis.
一個Usea.
即是三個位格本體.
一個的本質.
所以Hypothesis是解作上帝的本體.

$^{81}$一些Substance.
一些比較本質Substance.
這樣的意思.
所以大概可以解作什麼意思呢.
一些裡面很堅固不會改變的東西.
那些東西就叫做Hypothesis.
一些不會變的.
一些很基本很起初很基本簡單的東西.
就叫做Hypothesis.
所以文字就說我們應該要堅持住.
一些你最起初的Hypothesis.
這個就是經文裡面的意思.
所以經文就教導我們.
你要去持守住一些你起初在信仰裡面.
一些很基本不會變很簡單的東西.
這個就是經文裡面的教導.
不過說到哪裡呢.
信仰裡面這些最基本最簡單最原始最內在.
最重要的這些本質Substance是什麼呢.
說了這麼久.
經文是說什麼呢.
其實整個希伯來書第三章.
如果你看整段經文的時候.
你會發現經文是說一些關乎於.
以前舊約裡面摩西以色列人.
曠野40年有關的故事.
經文頭六節在希伯來書.
他將耶穌和摩西做比較.
帶出昔日摩西以色列人曠野40年的故事.
然後經文就開始去回憶昔日.
以色列人在曠野裡面的情景.
想當年以色列人出埃及的日子.
當初起初的時候.
上帝帶領著以色列人離開這個糟糕可怕的埃及.
我記得一走到去的時候.
那天很黑我們都很累.
我們就走了一整天的道路.
越走就越看不到前面.
後面一大堆的埃及的追兵追殺著我們.
突然之間天上出現了一大個巨型的火柱.

$^{121}$照亮整片的天空.
我記得當年就在這條火柱救了我們.
我們可以繼續走下去.
我記得當我們沒有水音的時候.
好像個個都捱不住的時候.
我們牧者摩西就用他們的木匠來擊打一盤石.
一打 突然間水就噴噴噴地飆出來.
然後就排在這裡 把水放進來喝.
我記得當我們去餓的時候.
到了晚上 突然間一隻小鳥飛過來.
整個營地都是小鳥.
我就抓著小鳥來吃.
早上露水上升之後.
看到一些很白色很像湯圓飯糰的東西.
我就拿來吃.
那些就是叫嗎哪.
我記得雲彩.
當我們在西冷山住了一年的時間.
整個的天空裡面都是雲彩.
沒有離開過我們.
描繪著滿有恩典和慈愛.
不過不知道為什麼.
這些起初這麼美好的片段都沒有了.
是不是因為時間的關係呢.
以色列人開始埋怨.
同一班人 同一個地方 同一個上帝.
同一個嗎哪 同一條雲柱 火柱.
在那裡雲彩.
但不知道為什麼再產生不了化學作用.
沒有了起初那種感覺.
沒有了起初那種力量.
很難吃 我想吃牛扒.
為什麼火柱這麼遲才到.
天都黑了.
可不可以不要只是盤石出水.
我想要喝豆漿.
對不起 我覺得摩西木曜不到我.
我想找一個更加明白我的牧者.
隨著時間的過去.
很多東西雖然是一樣的.

$^{161}$同一間教會 同一班人.
但可能很多東西都已經不同了.
40年.
最後上帝要求這班人在曠野裡足足40年的時間.
這班人留在同一片的地方裡足足40年.
40年曠野有一個很重要很特別的屬靈意義.
一個很特別的意義就是.
40年的曠野是一個不斷去兜兜轉轉.
沒有離開 也是不斷去兜圈.
由回到起點 又回到頭.
又重新再來過的時間.
你問為什麼會這樣.
這個曠野是不是有一個結界.
為什麼他們走不出去.
可能上帝正正要這班人在40年裡不斷重新去學習一些很起初很基本的東西.
上帝要他們學習的不是一些新的事情.
不是學一些新的信仰知識.
不是一些新的技能 一套新的神學.
而是上帝要他們在曠野裡重新去找回一些他們已經有的東西.
一些最基本最簡單最重要的東西.
一些不變的東西.
血泡.
希伯來書第三章第九節.
再次去提40年這個字.
在那裡你們的祖宗試驗我.
同一個曠野 同一個上帝 同一班色靈們.
這班色靈們在曠野裡40年.
學習一個最基本 回到最起初的信仰.
我想我們信裡面最大的問題不是你沒有.
而是你明明有 但好像沒有一樣.
這個可能就是信裡面我們要思考的問題.
我們不是沒有什麼東西 或者不懂什麼.
我們明明都知道 明明都有 但好像沒有一樣.
原來世界上有很多東西不是你要去學回來的.
你不需要去學習某個課程 或者看YouTube去學回來.
一直都放在你心裡面.
你要做的只不過是很簡單.
你要去堅持住去持守住它.
好好地放在你心裡面.
繼續放在心裡面.

$^{201}$這些不是一些很特別的東西.
甚至乎是一些很老土的東西.
你是簡單的 你要去簡單地去堅守住它.
持守住它.
Truth is in simplicity.
一些最好的東西 最真的東西 最美的東西.
都是很簡單的.
主耶穌很愛我.
盡心盡性盡意地愛著你的神.
不要犯罪 相信祂 感恩等等.
這些這麼基本簡單的東西.
我們信就是這樣.
沒有什麼深奧玄妙的佛家道理.
你不需要這樣誤透或者不懂.
一些很簡單基本的東西.
你好好地去持守住.
這個就是我們的信仰.
這些你聽過幾萬萬次.
我們講道都引導到很多次.
每次講到不同的方式 都講出同樣差不多的道理.
這個就是講完的困難.
所以一些最基本的東西.
就是我們信仰裡面最重要的東西.
不知道大家有沒有聽過一個台灣很有名的廚師.
他的名字叫做江振成.
你可以找他的名字.
一個很國際知名的廚師.
台灣人小時候就在士林夜市裡面長大.
幫忙幫忙在夜店幫忙.
15歲在法國學廚.
21歲成為法國米芝蓮三星級首席助理廚師.
31歲成為全球最佳150名廚師之一.
然後2011年去新加坡開了他屬於自己的餐廳.
叫做Restaurant Andree.
2017年這間餐廳成為全球排名第14名的餐廳.
全亞洲第二名的餐廳.
一個非常排隊要排一整年的餐廳.
不過非常富有的事情是.
去到2018年的時間.
他這位廚師就突然間親手結業.

$^{241}$自己創辦的這個餐廳.
大家可以看回Netflix.
Netflix裡面有一套紀錄片.
就講18年他重新開始的時候的一個紀錄片.
名字叫做初心.
可以搜尋初心這一個字.
原來這位15歲離開台灣的年輕人.
當年他已經42歲.
決定毅然要結束一切的事業.
回到自己的家鄉台灣裡面.
去尋找他的初心.
網上找到很多關於他的訪問.
一些有關初心的分享.
很值得跟大家聊聊.
他說初心給他一個很大的動力.
有兩個不同的情況.
初心是會幫助到他的.
第一個情況是什麼呢.
當你遇到很大困難的時候.
初心成為了拉你到原點.
一個很重要的決定.
讓你知道你應該做什麼決定.
什麼事情觸發你的初心.
然後答案就可以浮現.
你就知道你應該怎麼做決定.
這是第一個情況.
就是當你面對困難的時候.
初心會幫助到你.
第二個情況就是當你事業很暢順的時候.
當你遇到很多的信警,權力,金錢.
一切都不斷增加的時候.
你更加需要一個初心.
當你遇到這兩個不同的情況.
初心都可以幫你尋找自己走的道路.
因此對他來說初心不是純粹回到以前.
不見過去這麼簡單.
初心對他來說是一個不斷追求自己.
向前走一個很重要的力量.
一個追求完美的方法.
就是一個廚師不斷追求自己.

$^{281}$處於烹飪中的卓越和完美.
對他來說完美是什麼呢.
什麼叫完美呢.
完美不是當你名成利就.
完美是放下那一刻.
所謂放下就是當你做一個創作的時候.
他一定知道這個創作最後會變成什麼樣子.
沒有人會告訴你這個東西是什麼樣子.
都沒有權告訴你是什麼樣子.
只有你自己知道這個創作是一個什麼樣子.
當你最後願意把你的名字寫到完工的時候.
那一刻當你放下了就是完美.
但不知道你明白這句話的意思.
他只會為了你的初心告訴你.
你的完全.
你人生裡面的原備是一個什麼樣子.
初心是你的起點.
初心也是你的終點.
你要跟隨你自己的初心.
一直由以前跑到現在.
從這一刻跑到你將來的最後.
都是初心去帶動你自己的生命.
亭姐妹.
我們的信仰都是這樣.
所謂昔日起初的信仰.
去帶領我們走到這裡.
直到完滿的日子.
上帝叫我們好好地去持守住他.
堅持住他.
這些最基本的事情.
這份熱情是沒有辦法去培養出來的.
這份熱情在我們內心裡面.
我們回到起初信仰的熱情.
再一次去觸碰我們.
生命裡面最基本.
最簡單.
最重要的這份熱情.
它是我們一直走上路的一個很重要的力量.
最近經常問自己的問題.
究竟初信署的時候的我更加愛耶穌.

$^{321}$還是現在的我更加愛耶穌呢.
你不妨問問自己.
初信署的時候的你更加愛耶穌.
還是現在的你更加愛耶穌呢.
哪一個的我是一個更加好的基督徒呢.
可能這裡很多姐妹都已經信了超過十年,二十年.
三十年,四十年都有.
不知道四十年前的你.
還是現在的你更加愛主耶穌.
這段影片有很多不同的感慨.
我自己就開始和大家一樣.
開始有老花.
開始看不到字.
可能有些人都看著我長大.
十幾歲到現在.
開始有老花.
看那些說明書.
要拍照放大才看得到.
我都還未配眼鏡.
都不敢配著.
都覺得好像有些.
未接受到這個的事實.
我自己是九八年在錦宣家信耶穌.
那時候也是中六.
第一次回教會.
禮拜五晚在幼稚園裡回團契.
第一次回到團契裡.
第一次去祈禱.
接觸這個信仰.
開始認識.
就開始回初信班.
在錦宣家裡決志.
我記得自己第一次很誠懇地向耶穌呼求.
求主幫助我家人.
我自己和耶穌哭了很久.
去拯救我家人.
我記得第一次聽見上帝的呼召.
我覺得既然這是天父的世界.
不如我就去事奉祂.
開始走這條讀神學的道路.

$^{361}$第一次在教會裡實習.
那時候只有23歲.
穿著西裝在尖沙咀教會裡實習.
很緊張很緊張.
第一次離開神學院去德國讀書.
開始懷著一些冒險去德國進修.
六年之後.
博士考試就完成了.
第一次鬆口氣.
坐火車回柏林的時候.
開始正式事奉你.
轉眼間就到現在這一刻.
從一個18歲到現在40多歲的一個人.
前陣子我都和一些壽學生踢球.
今天都有弟兄踢完球回來教會.
原來現在這個年紀踢球.
我試過一次.
踢完之後整個人都喘了一小時.
回家之後不停地喘.
不知為何不停地喘很久.
望著自己24歲的身影.
看到自己昔日的事奉.
那種熱誠.
我自己不停地喘氣.
覺得自己跟不上這個20多歲的年輕人.
覺得眼前這個20多歲的年輕人很厲害.
很簡單很天真很有信心.
雖然有入世未深.
有些害羞.
但仍然有一份信仰的初心去伴他而行.
很羨慕這個年輕人.
雖然如此.
我知道他的初心我仍然可以有.
因為我和他擁有同一份初心.
因為我就是他.
他的起初信仰那份初心一直都可以在這裡.
就像希伯來書所說.
我們要將起初的初心堅持到底.
直到最後.
或者你可以將它叫做初心也行.

$^{401}$你叫做感動也行.
熱情也行.
所以很多時候叫做耶穌基督的愛.
更加很多時候叫做呼召.
耶穌一開始在你很多年前.
某個報道會缺字那一刻.
就已經呼召你.
呼召你去成為基督徒.
每一天再次成為基督徒.
懷著這份起初最基本最簡單的信仰.
就是這樣去追隨著這份初心.
每天成為基督徒.
記不記得自己第一次回來教會的時候.
第一次在教會裡事奉.
拿著程序表在門口排的時候的情景.
第一次唱詩班.
很緊張.
很怕自己唱錯的時候的情景.
第一次做詩琴.
第一次做教會的詩詩的時候的情景.
那份又戰驚又珍而重之的感覺.
正正是我們信仰裡很美麗的圖畫.
希伯來書三章這樣說.
總要趁著還有今天.
天天彼此相勸.
免得你們中間有人被罪迷惑.
心裡剛硬.
心裡有一個很奇怪的字眼.
叫做總要趁著還有今天.
But it's not one and another daily.
Well, it is still cold today.
聽到很奇怪.
趁著今天還是今天的時候.
其實經文是來自95篇詩篇.
為願你們今天聽他的話.
然後詩篇就重複以色列人40年曠野那段經文.
趁著還有今天.
意思是說今天不再是以前.
今天已經是今天.
聽起來好像非洲每60秒有1分鐘那樣.

$^{441}$意思是說今天不再是昨天.
現在就是今天.
昔日起初的信仰.
你今天的任務.
就是要很好的在今天裡.
趁著還有今天.
將這份起初的信仰.
繼續傳流到今天.
這就是你們今天我們要做的事情.
今天要做的就是.
將昔日這份起初的初心.
繼續流到今天.
當明天變成今天的時候.
重新再一次的.
將這份起初的信仰.
再一次成為今天.
我們要好好持守著的信仰的力量.
我們有沒有今天好好持守著這份信仰呢.
這就是我們今天要做的事情.
做基督徒.
我想這裡很多年紀都比我大.
你不單是成年人.
更加是別人的長輩.
更加是甚至是一個很有資深經驗的人.
但是這些東西.
阻礙不到我們成為信仰裡的小孩子.
懷著這份赤子之心.
來做基督徒.
無論你的頭髮有多白.
無論你的眼睛有多花.
你的牙有多少.
或者你的腳有多不行.
這些都沒有了我們那份信仰的初心.
這份初心是你以前年輕的時候.
和現在是沒有改變過的.
你只要很簡單的.
去繼續堅持下去.
持守我們基本的東西.
將這些基本的東西.
在今天好好守著.

$^{481}$這就是我們今天要做的事情.
堅持去做一個屬靈的弟弟.
主宋說讓小孩子來到我那裡.
不要禁止他們.
因為在神國的正是這樣的人.
我實在告訴你們.
反要承受神國的若不將小孩子.
斷不能進去.
所謂天國的孩子.
就是一群堅持著這份起初的信仰.
直到中老再回到最後天家的人.
雖然我們身體逐漸衰殘.
年紀越來越老.
身上有很多不同的傷痕.
無論是真實的還是比喻上的.
我們仍然可以去持守著這份信仰的初心.
懷著這份初心的人.
就是天國的小孩.
我們一起來去回到我們天國起初的起點.
我們祈禱.
再次呼求你.
因為我們很記得我們在信仰裡最起初的模樣.
雖然我們年紀或者社會.
或者很多都不同了.
但我們這份信仰可以不變.
我們能夠保持著.
我們堅守著這份信仰的起初的初心.
將這份對你的愛.
這份盼望.
這份愛心.
繼續去持守下去.
我們不需要多.
單單的來持守著今天我們的信仰.
今天我們代人都能夠充滿著愛心.
對你保持著簡單的信心.
對明天.
這個未知的明天仍然持有對你的盼望.
我們這份信仰.
從起初到今天.
再到他日.

$^{521}$我們離開世界的時候.
一直的去持守下去.
因為這是我們行走這條天國度路的必須的力量.
求主你這樣的憐憫幫助.
逢轉命求.
阿們.
\newpage



\section{使徒行傳 17:16-31}
\label{sec:RptdKZSozLI}
\textbf{2023-08-12 晚堂《進入人群的福音》使徒行傳 17\_16-31 講員:陳麗蟬牧師}
\newline
\newline
連結: \href{https://youtube.com/watch?v=RptdKZSozLI}{\texttt{https://youtube.com/watch?v=RptdKZSozLI}} ~~~~ 語音日期: 2023-08-15
\newline
\newline
\hyperref[sec:FlUdBHop01k]{\small{< < < PREV SERMON < < <}}
~
\hyperref[sec:index_chronic]{\small{[返順時目]}}
~
\hyperref[sec:index_scriptual]{\small{[返順卷目]}}
~
\hyperref[sec:kFsvDzgp4nI]{\small{> > > NEXT SERMON > > >}}
\newline
\newline
使徒行傳 17:16-31
\newline
\begin{longtable}{cl}
\hline
\hline
章節 & 經文 (和合本修訂版)\\
\hline
17:16 & \begin{tabularx}{0.7\textwidth}{X} 保羅在雅典等候他們的時候，看見滿城都是偶像，就心裡非常難過。 \end{tabularx} \\ \\ \relax
17:17 & \begin{tabularx}{0.7\textwidth}{X} 於是他在會堂裡與猶太人和虔敬的人，以及每日在市場上所遇見的人辯論。 \end{tabularx} \\ \\ \relax
17:18 & \begin{tabularx}{0.7\textwidth}{X} 還有伊壁鳩魯和斯多亞兩派的哲學家也與他爭辯。有的說：「這胡言亂語的要說甚麼？」有的說：「他似乎是宣傳外邦鬼神的。」這是因保羅傳講耶穌與復活的福音。 \end{tabularx} \\ \\ \relax
17:19 & \begin{tabularx}{0.7\textwidth}{X} 他們就把他帶到亞略巴古，說：「你所講的這新學說，我們也可以知道嗎？ \end{tabularx} \\ \\ \relax
17:20 & \begin{tabularx}{0.7\textwidth}{X} 因為你有些奇怪的事傳到我們耳中，我們想知道這些事是甚麼意思。」 \end{tabularx} \\ \\ \relax
17:21 & \begin{tabularx}{0.7\textwidth}{X} 原來所有的雅典人和居住在那裡的外國人都無暇管別的事，只是談談或聽聽新聞。 \end{tabularx} \\ \\ \relax
17:22 & \begin{tabularx}{0.7\textwidth}{X} 保羅站在亞略巴古當中，說：「諸位雅典人！我看你們凡事很敬畏鬼神。 \end{tabularx} \\ \\ \relax
17:23 & \begin{tabularx}{0.7\textwidth}{X} 我到處走走的時候，仔細觀察你們所敬拜的，發現一座壇，上面寫著『獻給未識之神明』。你們所不認識而敬拜的，我現在向你們宣告： \end{tabularx} \\ \\ \relax
17:24 & \begin{tabularx}{0.7\textwidth}{X} 他是創造宇宙和其中萬物的神；他既是天地的主，就不住在人手所造的殿宇裡， \end{tabularx} \\ \\ \relax
17:25 & \begin{tabularx}{0.7\textwidth}{X} 也不用人手去服侍，好像缺少甚麼似的；自己倒將生命、氣息、萬物賜給萬人。 \end{tabularx} \\ \\ \relax
17:26 & \begin{tabularx}{0.7\textwidth}{X} 他從一人造出萬族，居住在全地面上，並且預先定準他們的年限和所住的疆界， \end{tabularx} \\ \\ \relax
17:27 & \begin{tabularx}{0.7\textwidth}{X} 為要使他們尋求神，或者可以揣摩而找到他，其實他離我們各人不遠。 \end{tabularx} \\ \\ \relax
17:28 & \begin{tabularx}{0.7\textwidth}{X} 我們生活、行動、存在都在於他。就如你們的詩人也有人說：『我們也是他所生的。』 \end{tabularx} \\ \\ \relax
17:29 & \begin{tabularx}{0.7\textwidth}{X} 既然我們是神所生的，就不應該以為神的神性像人用手藝和心思所雕刻的金、銀、石像一般。 \end{tabularx} \\ \\ \relax
17:30 & \begin{tabularx}{0.7\textwidth}{X} 世人蒙昧無知的時候，神並不追究，如今卻吩咐各處的人都要悔改。 \end{tabularx} \\ \\ \relax
17:31 & \begin{tabularx}{0.7\textwidth}{X} 因為他已經定了日子，要藉著他所設立的人按公義審判天下，並且使他從死人中復活，給萬人作可信的憑據。」 \end{tabularx} \\ \\
[1ex]
\hline
\hline
\end{longtable}
$^{1}$各位電影姐妹晚安.
在今年年初世界各地都開始通關了.
所以每個人都可以去旅行了.
大家去旅行了沒有?.
有些人就說香港的人就是報復式去旅行.
無論長的假期又要去.
短的假期都要出去一下.
去透一下氣.
今天一次出外旅行的時候.
大家最著緊的是甚麼呢?.
去到那個地方立即找些美景拍照.
找些最美的角度.
又或者早前已經上網了.
找一找哪裡有好東西吃.
又或者去找資料.
某一個地方最好去買些甚麼.
到處去購物.
但剛才主席幫我們讀出來的經文.
保羅去到雅典這個地方.
一個很有文化的地方.
也很有藝術的地方.
他不是去欣賞裡面的宏偉建築.
也不是去注意那些雕刻是多麼的精緻.
那些神像是多麼的有神氣.
而經文告訴他.
他心裡面非常的焦急.
為甚麼會這樣呢?.
因為雅典這個地方.
在公元之前的第三,第四世紀的時候.
已經有一些很美的建築物.
這些當然來到現在已經變成了石柱.
你想想公元前的第二,第三世紀.
已經很宏偉,很美.
那時候的人的智慧.
真的不知道為甚麼會這麼厲害.
但保羅不是去慢慢欣賞.
其實這次去雅典.
是他第一次去這個地方.
他去到第二次.
和提摩泰和西拉有旅行步道.

$^{41}$他去到不同的地方.
這次他去到菲律賓,提薩羅里加和比利亞.
去到這三個城市的時候.
他會被人抓,追,鎖,打.
在我們右上角上高.
特意弄了一個紅色的.
三個城市都在那裡.
就是菲律賓.
在菲律賓那裡被人追,被人打.
然後就跑去提薩羅里加.
然後又是這樣的情況.
人們繼續追,追到比利亞.
之後就不行了.
到這三個城市非常混亂.
在迫不得已的情況下.
有些人就把保羅帶到雅典.
我用了一個黃色的.
在那附近就是雅典.
雅典是希臘的首都.
也是希臘裡面一個最大的城市.
這個城市最出名的.
就是它的一個文化中心.
哲學,藝術,戲劇都是非常出名的.
去到一個這麼有文化氣質.
這麼有藝術的地方.
但心裡面卻是非常難過.
這個難過這個字.
有些人就逆去做這個非常心急.
其實是表達他們的情緒.
是一個很高漲的.
甚至高漲到一個地方就是抽筋.
我想很重要.
當你的情緒很緊張的時候.
去到抽筋的時候.
這是一個非常非常緊要的情況.
為什麼保羅這麼有趣.
去到一個這麼有文化.
我想在那時候的建築物都仍然很漂亮.
為什麼他的情緒搞成這樣呢.
就是因為他覺得一個這麼有文化.

$^{81}$這麼有智慧的地方.
但現在整個城都是偶像.
這個雅典這麼有智慧的人民在這裡.
但他們全部的盼望.
所有的東西都寄託在這些偶像的身上.
他很心痛.
我們的上一代.
我自己的爸爸媽媽.
我不知道你們的爸爸媽媽怎麼樣.
是拜祖先 拜關公 拜王大仙.
但我爸爸媽媽他們都說.
他們不知道在拜什麼.
不過之前家人拜他們就拜了.
但經過了很多的禱告.
很多的見證.
很不容易才把我爸爸媽媽.
帶到他們耶穌那裡.
他們來去接受了耶穌基督.
但到下一代.
今年年初就和我的家人去旅行了.
去日本旅行.
發覺年輕的一代.
沒錯他們不再拜祖先.
不再拜關公.
也不再拜王大仙.
但有一個親戚.
他就玩瑜伽玩了很久.
玩到很高深那些.
原來會這樣.
他們其實和我爸爸媽媽以前差不多.
初一十五要吃齋.
還有他們不吃牛肉.
不是他們不喜歡吃牛肉.
而是不知道為什麼.
瑜伽那些都是這樣.
是不可以吃牛肉的.
另外一個年輕的親戚又怎麼樣呢.
他們去到日本.
在香港又不會.
走過什麼廟又不會進去拜拜.

$^{121}$但去到日本的時候.
經過那些廟宇.
我們自己計劃的行程.
都沒有進去.
他就說請你等一等.
走過每一處看到有廟的.
他都要進去拜一拜.
當時我也很心痛.
雖然我還沒有去到保羅.
心痛到抽筋.
我沒有.
但裡面有一種心痛.
一種慨嘆.
那些年輕人.
知道他們其實都在基督教學校讀書.
無論小學或者中學.
讀了這麼久.
聽過了.
但為什麼都不明白.
都不肯相信.
但去到這些瑜伽.
或者日本那些宗教.
他們聽過的時候.
竟然就信到十足.
我心裡面都很痛.
很痛.
又想想.
要很多很多禱告.
繼續那個工作.
因為保羅心裡面.
非常難過.
或者非常焦急的程度.
去到抽筋.
那個字是怎樣的呢.
我們看看出埃及記第20章第5節.
當以色列人.
出埃及.
摩西準備帶他們進去迦南地的時候.
耶和華已經奮不顧腹.
不可跪拜那些像.

$^{161}$因為迦南地裡面.
有很多偶像.
說明不可以拜他們.
不可以事奉他們.
因為耶和華你們的神.
是忌邪的神.
這個忌邪這個字.
用希伯來文.
對希臘文的時候.
就是這個.
和這個心急難過這個字.
是相關的.
就是說.
當耶和華和以色列這個百姓.
他們的關係.
上帝是他們的神.
他們是上帝的子民.
那個關係很密切.
就好像夫婦之間的關係.
所以耶和華和這班子民說了.
你們進到迦南地的時候.
不可以拜那些像.
不可以事奉他們.
好像夫婦關係.
你去拜他們.
事奉他們的時候.
就好像毀了婚約一樣.
有一位太太告訴我.
當她知道她的老公.
在外面有另一個女人的時候.
她告訴我.
她生氣到情緒.
她感受到.
知道那些血一直一直.
飆升到腦子裡面.
真的感覺到.
頭暈暈的.
當時她真的很想進廚房.
拿出一把刀.
就一刀插死她老公.

$^{201}$這個難過.
心急.
或者這個疾度這個字.
就說真的情緒很高漲.
保羅當時去到.
原本這個地方這麼有文化.
這麼有藝術.
沒理由他們不認識耶和華.
他很緊張很心痛.
各位電影姐妹.
想一下我們自己.
我們對我們家裡未信的家人.
親戚朋友.
我們有沒有很緊張.
很焦急呢.
這個保羅.
他就覺得.
看到整個雅典城.
都是偶像來的.
那些人這麼多的人.
他們全部走歪了.
他們永遠的路是怎樣呢.
就是步向死亡.
其實當我們想深一層.
如果當我們未信主的家人.
親戚朋友.
如果沒有耶穌的時候.
他將來的終局.
都是這樣的.
所以電影姐妹.
我們要迫切一點.
多一點為他禱告.
我相信.
終有一天.
他們都會加入耶穌基督的國度裡面.
第二方面我們就看一下.
保羅去到雅典的時候.
因為提摩太和西拉都還沒有到.
他自己在城市裡面四處逛.
他先去到猶太的會堂.

$^{241}$跟自己的同鄉.
猶太人聊天.
一些敬虔人.
他們都是愛邦人.
不過他們都尊重耶和華.
都信耶和華.
在那裡跟我們聊天.
然後他就每一天.
都在那些市集那裡.
跟別人辯論.
在市集那裡不是去購物.
買紀念品.
而是當時那些人就是這樣.
他們在市集那裡除了有買賣之外.
他們還有交易的地方.
就在這些市集.
他們甚至當你有一些.
你自己的思想.
一些想法的時候.
你就可以站在街頭一些地方.
就在那裡去講一下你的理論.
只要你講一下.
講一下.
有些粉絲認同你.
因為你的粉絲越來越多的時候.
他就會邀請你去到另外一個地方.
叫做阿勒巴古裡面.
來去演講的.
聽聞古希臘裡面.
這個很出名的哲學家.
蘇格拉底的時候.
他都是這樣.
最初在市場那裡.
在那裡講一下.
什麼叫做幸福.
什麼叫做美滿.
講著講著講到出名了.
接著就去阿勒巴古.
接著就被很多人認同了.
其實香港現在都是這樣.

$^{281}$你看到那些人.
在公園裡面唱歌.
無論早上或者黃昏的時候.
你在公園那裡.
你看看.
很多人在那裡唱歌.
唱著唱著.
唱到上電視台.
什麼中年好聲音.
接著再唱的時候.
就成為了很出名的歌星.
你聽一下.
以前的那一代.
梅艷芳那些.
去唱《星堡之夜》.
接著又成為了很多人都喜歡的.
其實古希臘的時候.
那些人就是這樣.
這個保羅.
他講自己的信仰的時候.
就講著講著.
就被那些人吸引了.
於是都叫他去了阿勒巴古.
不過有一個問題.
就是那裡的雅典人.
這個這麼有文化的地方.
每個人都覺得自己是學者.
每個人都覺得.
自己有一套的理論.
所以不要緊.
那些人就可以在那裡.
在街頭任人說.
總之有粉絲走過來聽.
這樣就可以了.
不過那些人太多人說了.
所以那些人就怎樣呢.
當新聞來聽聽.
講一下.
這樣就可以了.
他們就這樣.

$^{321}$原來是這樣的.
聽完之後.
可能就走去那裡說.
原來是這樣的.
這樣講的.
在感受就少一點了.
在我們住大廈那些.
特別一些屋村那些的管理處.
其實又有類似這樣的情況.
就是當日頭的時候.
很多退休人士.
最喜歡去哪裡呢.
就是去管理處.
和那些管理員.
在那裡打一下牙交.
特別現在天氣這麼熱.
管理處那裡又有冷氣.
所以管理員就收集了.
我們整個大廈的人的資料.
這個人的媳婦.
究竟好不好.
那個兒子孝不孝順.
然後誰什麼時候轉工.
就什麼都知道.
很多的資料.
有一位師姐她說.
她退休之後.
她就會去做管理員.
因為在大廈裡面.
那些人就在那裡說話.
有時加上了一些是是非非.
有些是無無聊聊.
這樣.
她說管理員聽了.
然後聽了A座是怎樣的.
就告訴B座C座D座那些.
她就說要在那裡聽了之後.
就將這些無無聊聊.
或者是是非非的東西.
就化為福音.

$^{361}$來跟他們說.
不過她退休幾年.
我都沒有見到她做.
繼續說這個雅典.
那些人在那裡閒談.
說話.
在裡面.
就有很多很多的偶像.
每個人都說他自己的東西.
然後每個人都可以做一些東西出來.
她在當中保羅就看到了.
這麼多的神廟.
這麼多的祭壇.
有些人說.
就這個雅典這個地方.
這個神明.
比人還多.
因為當時希臘這個地方.
就說裡面那些周圍不同的地方.
有不同的神明.
加起來都不夠雅典多.
所以可以想像到.
雅典是多麼厲害.
不過其實我又欣賞.
雅典人其實很敬畏神明.
他們就不怕萬一.
最怕一萬.
他們很害怕.
萬一拜少了一個神.
或者拜錯了就慘了.
所以他們就無所不拜.
無神不拜.
連一個神.
就怕不知道叫什麼名字.
總之就立一個碑.
叫做「未識之神」.
就全部都立了.
全部都拜了.
有拜錯.
沒有放過.

$^{401}$這個就是希臘人.
為什麼他們這麼多的思想.
這麼多的神明.
當中這些神明.
或者這些思想.
就分了兩個學派.
第一個學派.
就叫做伊碧溝魯派.
這個是無神論者.
他就會覺得.
不理這麼多了.
神明在不在都不要緊.
人生苦短.
最重要的就是活得開心.
Happy.
今朝有酒今朝醉.
明日愁來明日當.
明天發生什麼事.
不要管.
今天最重要.
Happy.
最重要.
享樂.
這一派的主義.
一路一路發展的時候.
到後期就變成了.
做一個很重慾的主義.
就搞到人際關係.
那些情慾亂搞.
另外一派.
他們的思想就是斯多瓦派.
就是另外一個極端.
就是禁慾主義.
他們就說.
世界上所有的物品.
所有的東西.
都是神來的.
所以他們就什麼都拜.
就像我們香港也是.
看到一棵樹.

$^{441}$很大一棵.
很特別.
就又拜.
又一棵許願樹.
這塊石頭.
又好像一塊高一點.
一塊矮一點.
又好像一個男孩和一個女孩.
所以又去拜這個恩怨石.
所以他們就什麼都可以找出來去拜.
就說到.
全世界所有的東西.
都是神明.
所以我們生命就要謹慎一點.
我們這個身體.
是一個臭皮囊.
所以苦待下.
虐待下.
自己不要緊.
所以他們可以不吃.
不睡.
不洗澡.
做得辛苦一點.
吃力一點.
他們就覺得.
總之要苦待自己.
以至在裡面那份靈性.
就可以提出來.
這樣才是一個真正的人生.
這兩派.
學派.
時時都有爭論的.
享樂主義.
禁慾主義.
他們又有不同的神明來支持.
保羅激動到抽筋.
就覺得.
為什麼他們有了宗教的思想.
覺得這個世界上一定有一個的神.
所以希臘裡面有男神.

$^{481}$有女神.
有戰爭的神.
有和平的神.
有智慧的神.
有藝術的神.
有農業的神.
什麼神都有.
為什麼那位創天造地.
勝過死亡的權勢的耶穌.
他們卻沒有拜到呢.
他們卻不認識呢.
所以他們心裡面很激動.
這個真神.
他們為什麼會不認識呢.
剛剛在一個培靈會裡面.
周曉輝牧師講道裡面.
講了一個屬靈方程式.
我將屬靈的方程式.
套入雅典的情況.
周牧師說.
如果沒有了神的時候.
就等於零.
你覺得什麼都是好的.
什麼都好.
但當裡面沒有了神的時候.
一切都是沒用的.
這個雅典.
它有神明.
或者有神的觀念.
所以它有很多神像.
偶像來去拜.
但裡面沒有了真正的神.
真正的上帝.
所以拜多少神都好.
是沒用的.
第二樣東西.
只要有神.
就算你沒有了.
都是一個好人生.
只要在我心靈當中.

$^{521}$有了主耶穌基督的時候.
其他東西我們得不到的時候.
我們的生命.
都可以是一個好人生.
好有意義的人生.
就好像奧古斯丁所說.
我們心靈裡面有很多很多的洞.
我們需要被人尊重.
很需要有金錢.
很需要一些朋友.
很需要很多很多的東西.
但其中一個洞.
屬於信仰.
並且你要找得對.
所以當你屬於信仰.
屬於真神的洞.
沒有填滿的時候.
你有了金錢.
有了地位.
很多人尊重你都好.
你有了.
但你生命都覺得不滿足.
都覺得欠了些什麼.
但如果我們在心靈當中.
屬於上帝真神的洞.
被填滿的時候.
我們就算沒有金錢.
沒有地位.
沒有人尊重都好.
我們都覺得生命裡面.
是有力量有意義的.
Good.
凡事都好.
沒有了上帝的時候.
就等於零.
有真神耶穌基督.
就算我們沒有金錢.
沒有地位.
沒有知識.
沒有什麼都沒有.

$^{561}$但你仍然可以過一個.
很好的人生.
當保羅在雅典那裡.
經文就說到他.
他到處走走的時候.
他到處走.
他不知道提摩太和西太會什麼時候到.
都不會立即到那麼快到.
他到處走走走.
他仔細的觀察.
看看周圍的環境.
他不是找一個很美的景.
或者很美的角度的景來拍照.
他認識一下雅典這個地方.
他發現了一個祭壇.
就是獻給未識之神.
所以保羅就說.
他們未識不識.
他們那麼多神裡面.
就是沒有了主耶穌基督.
所以下一步就是.
我現在向你們宣告.
宣告什麼.
你們未認識的上帝.
未認識的神.
是一位怎樣的神.
原來你們未認識的神.
這個才是最偉大.
最真實.
最厲害的那一個.
他就開始介紹這位神.
這位神是萬物的創造主.
他創造宇宙和其中萬物的上帝.
他既然是天上高的主.
他就不住在人手所做的奠禹裡面.
所做的神廟.
他不用住在裡面.
也不需要人手去服侍.
好像缺少什麼似的.
所以不用想著要給什麼.

$^{601}$給什麼是不用的.
相反的是.
是上帝將他的生命.
氣息.
萬物都去給過萬人.
所以不要想著.
上帝缺少什麼.
所以我們要給過上帝.
相反的是.
上帝按著我們需要什麼.
去給過我們.
我記得我曾經說過.
不知道有沒有在漫堂說過.
在琴樓錦的時候.
因為所有的雞都被宰了.
我好記得那時候我坐小巴.
聽到兩位師奶在我後面聊天.
他們說.
真是慘了.
雞被宰了.
現在拜神沒有雞.
怎麼辦呢.
不知道買鴨他喜不喜歡吃.
買鵝他喜不喜歡吃.
很怕買了那些不合他的口味.
他不喜歡吃就慘了.
另外一位師奶就聰明很多.
她說是呀.
我也很怕買錯.
我不管那麼多.
我拿個利是封.
放五百塊進去.
然後放在那裡去拜神.
由得他.
我說你喜歡買什麼吃.
就去買吧.
什麼都好買.
拜完之後就放回袋子裡.
有時候我們拜偶像就是這樣.
但是保羅很清楚地告訴我們.

$^{641}$也告訴雅典那時候的人知道.
我們的上帝是一切的源頭.
他不需要我們給他什麼.
相反.
他安著我們每一個人的境況.
去將生命.
將氣息.
將萬物.
賜給我們.
第二方面.
保羅繼續說.
他說他是去定準我們年限和所住的疆格.
他定了我們每一個人.
你有多少歲.
或者你有多少日子.
他全部定了.
我們的生活.
我們的行動.
我們的存在.
都在於他.
他全部計劃了.
他在當中.
他創造了我們.
他維護著我們.
眷顧著我們.
所以只要我們專心去依靠他的時候.
我們不需要那麼多的顧慮.
耶穌基督在我臨上十字架之前.
他也跟一幫門徒說.
他在說他們兩個彼此之間的關係.
你們若常在我裡面.
就是那幫門徒.
時時在主耶穌基督的裡面.
我的話也常在你們裡面.
只要我們專心地跟著耶穌來走的時候.
他的話就成為了我們生命裡面很多的鼓勵支持.
凡你們想要的.
你們祈求.
就給你們成全.
所以是耶穌安著你的警方給我們.

$^{681}$而不是我們要努力去做些什麼.
給我們這位創造的主.
到最後.
就是說這位雅典人還沒有認識的神.
就是一位公義的審判主.
第31節他說.
是因為他已經定了日子.
要按著他所設立的人.
按公義審判天下.
上帝一定會審判天下.
所以不要想著有時候我們看到有些人.
很兇的.
我們就覺得很不公道.
為什麼上帝不懲罰他.
中國人也有句話.
惡有惡報 善有善報.
若然未報 時辰未到.
到時候的時候.
一定會報.
所以有時候當我自己也去.
遇到一些很不公道的事情.
遇到一些覺得沒有公義的事情.
很激氣 很激氣 很生氣.
但我們堅信.
堅信上帝是那位公義的主.
是要審判大地的主.
他一定會為我們主持公道的.
所以將這個審判權.
交回給上帝.
自己不要亂來.
我們認識了.
雅典人未認識的神.
我們認識了這位創造的主.
認識了他捐購我們的.
也都拯救我們的主.
這個詩篇50篇裡面.
上帝自己說的.
我就想我們一起讀了.
詩篇第50篇.
10至15節.

$^{721}$作為我們今天整篇訊息的總結.
上帝自己說的.
我們一起讀.
1 2 3.
林中的白獸是我的.
千山的牲畜也是我的.
山中的飛鳥我都知道.
田野的酒獸也都屬我.
我若是不用告訴你.
因為世界和其中所充滿的都是我的.
我喜吃公牛的肉.
我喜喝公羊的血.
又要向至高者還你的怨.
並要在患難之日求告我.
我必答救你.
你也要榮耀我.
我們同心圖告.
因為你讓我們先嘗到救恩的滋味.
你讓我們經歷在沒有耶穌的時候.
我們生命的迷茫.
我們生命的黑暗.
但多謝你一直主耶穌基督.
我們找到方向我們找到盼望.
幫助我們在周圍很多人.
他們仍然在迷茫.
仍然在污糟的情況.
求你給他們保羅那份.
看到那些人不信不認識上帝.
那個激動的心.
給我們迫切的為他們去禱告.
求主幫助.
讓我們都將真的未信的人.
未信的家人親戚都教導主.
願你的救恩早日臨到.
多謝主耶穌基督垂聽禱告.
奉主基督的聖名.
阿們.
\newpage



\section{約翰一書彼得前書 1:1-2}
\label{sec:kFsvDzgp4nI}
\textbf{2023-08-13 加恩樓《使人蒙福的「道」》約翰一書 1\_1-2,彼得前書 1\_23-25 講員:樓恩德牧師}
\newline
\newline
連結: \href{https://youtube.com/watch?v=kFsvDzgp4nI}{\texttt{https://youtube.com/watch?v=kFsvDzgp4nI}} ~~~~ 語音日期: 2023-08-15
\newline
\newline
\hyperref[sec:RptdKZSozLI]{\small{< < < PREV SERMON < < <}}
~
\hyperref[sec:index_chronic]{\small{[返順時目]}}
~
\hyperref[sec:index_scriptual]{\small{[返順卷目]}}
~
\hyperref[sec:FIbxv36OMMk]{\small{> > > NEXT SERMON > > >}}
\newline
\newline
$^{1}$各位弟兄姊妹平安.
我們一同有個禱告.
我們親愛的天父.
也是創造萬物的主宰.
拯救我們的耶穌基督.
在我們心裡感動我們的聖靈.
我們現在一同齊集來到你面前.
求你用你的話語.
在我們面前給我們照明.
藉著聖靈在我們心裡的工作.
就在我們這個時刻.
我們所頌唱的詩歌.
我們口裡的言語.
心裡的意念.
我們的奉建都能夠蒙你的圓立.
奉耶穌名字祈求.
阿們.
和大家一同有這個時間來敬拜神.
和今天有機會和大家分享神的話.
我今天選了這個題目.
使人蒙福的道.
在這兩年裡.
是很留意去思考的.
也可以說整個恩典.
今年講九賠靈會.
那十天的聚會.
那二十八場的信息.
都是使我更加堅定相信.
在現今這個時代.
你和我活著的時代.
所需要的.
不是任何的經濟理論.
任何的政治主義.
任何數據給你的指引.
做生意人要看數據.
科研要講數據.
很多東西要講數據.
但是今天早上想和大家看的就是.
如果你要蒙福.
我用蒙福這個字.

$^{41}$待會我會和大家解釋.
你所需要的.
只有一個字.
就是道.
剛才讀了兩段經文.
我和大家再提一次.
這段經文裡面.
是約翰一書一章一到二節.
你可以留意到有幾個字.
叫生命之道.
約翰講這個生命之道的時候.
他說他是聽見的.
是看見的.
是親眼見過的.
親手摸過的.
請問一樣東西.
不知是什麼東西.
你可以看過.
看過.
聽過.
又可以摸過的東西是什麼呢?.
理論.
是你以為觸感到的東西.
第二節他說.
這個生命已經顯明出來.
他講明了這是生命.
這個生命是看過.
過去又作過見證.
現在我們能夠感覺到一個永遠的生命.
他是可以傳給你們的.
今天在座中間.
我們有青少年人.
有些成人.
在你人生裡面有些什麼東西.
到今天你可以很肯定的告訴別人.
你是親自經歷過的.
一講這段聖經的時候.
所以我馬上聯起了.
有一個聖經裡面記載的人物.
這個人物本來很富裕.

$^{81}$但後來遭遇了很慘痛的經歷.
一日之間幾乎家破人亡.
什麼都沒有了.
最後連自己身體都患了嚴重的皮膚病.
有三個朋友有四個朋友來跟他聊天.
到最後的時候.
他講了一句話.
這句話我相信我們在座中間.
成年的基督徒.
或者年輕的基督徒都曾經講過的.
不過你沒有想清楚.
我從前風聞有你.
我現在親眼見到你.
請問這位約伯他怎樣見到神?.
以利亞摩西他們想見神都見不到.
這位約伯怎樣見到神?.
難道他那句名言那句經歷.
是一句大話?.
是一句安慰人的話?.
你都講過的.
做見證.
我從前風聞有耶穌.
現在我親眼見到耶穌.
我聽過有人這樣講見證的.
我病到要死了.
結果現在好了.
我心想你見到耶穌.
你不是見到醫生嗎?.
上個星期放榜了.
考試前驚到七彩.
考完出來成績很好.
我從前聽聞有耶穌.
現在我見到耶穌.
耶穌在哪裡?.
在市場上?.
所以我今天特別用這段經文.
然後我再講下面一段.
跟大家思想一下.
到底你所想有的福.
是甚麼福?.

$^{121}$下面那段是在彼得前書第一章說的.
你無夢了重生.
不是由於能壞的種子.
而是由於不能壞的種子.
是藉著神活潑常存的道.
又講道.
因為凡有血氣的盡都如草.
他的榮美好像草上的花.
草必枯乾花必凋謝.
唯有主的道是永存的流心.
我今天跟大家所講的道.
不是一個道理.
不是合你心意.
不合人心意.
這是永遠常存的道.
所傳給你們的道.
就是這個福音.
這個道原來就是福音.
在彼得前書第一章二十五召說.
主的道是永存的.
所傳給你們的福音.
就是這個道.
所以我今天要跟大家講的就是福音.
你說牧師.
這是主人崇拜.
不是福音主人.
正正就是因為這個福音.
整個道你忽略了.
而我接著恐怕.
因為我上一堂在錦繡堂是過時了.
過了十分鐘.
我希望今天不會過時.
不過我想你們以後沒什麼.
如果你們覺得已經過時了.
你舉手我立刻停.
你不舉手我就一路講下去.
你說這樣都行.
因為很多基督徒.
特別我對年輕人少年人說.
你們這個年紀可以站在台上.

$^{161}$這麼有次序.
能夠將神.
我覺得你這樣會唱雞尾歌.
耶穌歌是合在一起的.
不會說基督歌就這樣出來的.
你們一樣有編排的.
有沒有人編排?.
誰編排?.
不要叫了 舉手吧.
少年 你可以有福的.
我雖然不是很會音樂.
但我聽得出來.
好像雞尾歌一樣.
但是沒有接到接縫.
是不是?.
我們所講的是道.
是神的話語.
你可能在座中間的聽道.
聽了十幾二十年.
最近香港十天的聚會.
有幾個人.
不是因為你沒聽不懂.
因為他忙到七彩.
三年前我們叫他來的時候.
因為疫情.
請他來的時候.
疫情.
那年我們安排的講完.
全部在海外.
我們特別安排攝影隊.
拍攝隊去他那裡和他拍攝.
他沒有時間做.
他想推了那年之後.
不用再來.
因為他不是一個普通港島人.
他以前是神學院的院長.
現在做的是.
一千間教會的傳道人.
訓練的trainer.
即是訓練傳道人的人.

$^{201}$所以當這些是題外話.
不過都讓你知道.
他想這次來.
是因為他是一個好朋友.
特別要跟他說了.
他說好來.
因為一般請海外講完.
三天七天.
已經很長時間了對他們來說.
因為他們不是沒事做.
等你請一個港島人.
所以十天.
是真的很少人肯答應.
到第二天.
不知道是第一還是第二天.
他講起一個見證.
講他太太的祖母.
差不多是一百年前.
預備來中國宣教.
結果患了病.
來不了.
但是他的家族不斷為中國來禱告.
講完之後.
他以為做一個見證這麼簡單.
我馬上有深刻的回應.
因為我的祖母都是宣教士.
來中國傳福音.
以至她病得醫治信耶穌.
得到上帝的應許.
然後安然離開世界.
所以給我一個好心的印象.
完了之後.
有機會跟他談的時候.
我就說.
你太太祖母的心願.
我問她.
有沒有成就到.
她停一停.
我說你這次來到香港這個地方.
跟她成就了.

$^{241}$因為香港已經是中國的地方.
這句話是一個閒談交通.
她是有感而發作見證.
我是有感而發給她一點點回應.
我相信神工作.
接著我禱告.
後來我們再跟她談的時候.
或許神准許她將來還會再來.
不過貿然這兩三年.
我想除非她有其他歲月.
我不知道.
所以弟兄姊妹.
當神的靈在的時候.
她工作的時候.
台上的道.
講的時候.
不是講那個人怎樣有學識.
有學問.
有多少個銜頭.
是她有沒有經歷過神的話語.
然後將她自己生命的見證.
透過神的話語的操練.
然後擺在眾人的面前.
然後你所聽到的.
不是她自己在吹水.
那是將生命之道擺在你面前.
當你接受生命之道的時候.
你一定有回應.
我舉一番聖經的例子給你聽.
有兩個門徒.
耶穌釘死之後.
就算耶穌復活之後.
他們都沮喪了.
離開耶路撒冷.
向二萬五師路上走.
當他們走的時候.
突然有一個陌生人.
他們應該認識的.
但他們太多這道問題.
遮蔽了他們的視線.

$^{281}$裡面看不到跟他們同行的是誰.
那個人跟他們一路走的時候.
就將聖經上指著基督所說的.
在舊約裡面的記載.
一樣的東西告訴他.
是將聖經上記載的東西告訴他.
聖經怎麼記載?.
這個人聽著聽著.
越聽越覺得好像有些東西有料到.
就在村莊一起去吃飯.
當吃飯的時候.
這位他們認為陌生人的人.
祝祭的時候.
這兩個人眼睛明亮.
當他明亮的時候.
不見了的時候.
他們說了一句話.
剛才他跟我們說話的時候.
我們不是心裡火熱嗎?.
每一次到你.
每一次不是這一次.
每一次當你去到神的面前.
聽到的時候.
你覺得還不說完.
還是心裡有個回應.
你裡面靈魂還沒蘇醒.
你不會有回應的.
就算回應也是假的.
但你裡面靈魂蘇醒之後.
你知道這是神的道.
是永存的道.
是福音.
就是耶穌基督自己.
我等一下不會說太多我自己的話.
會將一張一字聖經給你看.
你自己讀完之後.
你聽聽到底今天.
神有沒有藉著聖經的說話跟你說話.
如果有的.
你接受到這個道的話.

$^{321}$你就一定有蒙福.
這個福不單止今生.
直到永生.
如果沒有的話.
要說我一句衝突的話.
衝撞的話.
你浪費時間.
為什麼要來.
剛才在下雨.
我不來也可以.
我為什麼要練習.
帶著大家去敬拜神.
今天我們能夠來到神的面前.
我相信剛才那句說話.
這福音.
這個道就是神的福音.
我會問道是什麼呢.
想一想.
第二個問題.
什麼是福.
今生是福.
永生是福.
這個答案就是.
兩樣都要.
沒理由今世衰到鐵底.
將來才得永生有福.
如果今生.
被你得到了.
你所期望的福.
滿不滿足.
我可以寫一個包單給你.
我肯定告訴你.
你不會滿足的.
我剛才提到.
上個星期放榜.
我想你們中間都見過放榜的.
放榜之後.
有機會進修.
最後拿到一張文憑.
一張紙.

$^{361}$你猜那張紙.
可以讓你用多久呢.
我做了校長之後.
每年都要看人家申請的時候.
那張紙.
很重要的.
沒有那張紙就不可以了.
到我自己收拾.
我搬了幾次.
我從學校退休.
搬到教會.
教會退休.
我搬走的時候.
不是在家裡.
在辦公的地方.
兩次搬走的時候.
我打開文件一看的時候.
我心想還有沒有用呢.
到我今天.
那些文憑對我還有沒有用呢.
有.
我留給我的小孩看.
給我的孫看.
以前我都是讀過書的.
這張是那時候獎金的文憑.
還有什麼用呢.
我有一次想找.
我連文憑都不知道放在哪裡了.
今天對你是很重要的.
是不是.
所以我自己在這裡就問.
如果福不單止今生有永恆的話.
那我要追求的是什麼呢.
還有一個東西我要問的.
問什麼呢.
等我自己搞定.
到底這個福.
是不是只說信耶穌得永生的福.
會自動落的.
不要緊.

$^{401}$如果信耶穌得的福.
只有永生的.
那今生呢.
耶穌曾經說過一句話.
我來了是要叫揚得生命.
這個揚得生命就是人.
並且要活得豐盛.
什麼叫活得豐盛呢.
不是沒有艱難.
不是沒有病痛.
而是你處處能夠將這個永生之道.
透過你的生活.
讓人看到.
不是你厲害.
不是你健康.
不是你長壽.
而是看到整個恩典.
說到這裡.
很多時候我不想說.
很多人會讚.
說你的頭髮很漂亮.
是不是.
我從學校退休已經19年了.
年輕人.
在我教的學生中間.
有些事業很成功.
40多50歲就走了.
所以在這段時間裡.
我很能夠紀念那班畢業生.
有不少他們在學校的時候.
是信主的.
很熱心神的工作.
我上個月去外面負責傳道工作的時候.
遇到一個學生.
他現在在某某神學院裡.
做中文部的主任.
他中四的時候.
學校舉辦福音週.
他可以四點鐘起床為同學祈禱.
我想你們四點鐘還在做夢.

$^{441}$也有人昔日很熱心.
今時今日已經做了.
在專業上有相當地位.
但已經不回教會.
有些人到處打電話給我聽.
校長我確診.
不是確診.
我醫生判了我是第四期.
我心想你應該打電話給醫生.
你打給我幹什麼.
我跟他祈禱.
然後上帝給他.
本來判斷他一兩年要走的人.
結果活多四五六年.
後期的時候他不斷去做見證.
年輕人.
你有前途的.
是不是活得豐盛.
不是你計劃到的.
活不活得豐盛.
不是地上的東西可以給你的.
你的意志你的勤力都沒有用的.
那不用嗎?.
不是.
當你得到這個道之後.
它可以給你有力量.
我最近想一個問題.
對發生了的事件.
你不需要相信.
不是吧.
你只有接受.
你生在哪個家庭.
你有沒有信不信.
你就是接受.
你讀書成績好不好.
以前過去的那些.
你只有接受.
接受過去的歷史.
兩個態度.
一個是感恩.

$^{481}$一個是埋怨.
是不是.
請問你對過去發生的事件.
感恩多還是埋怨多.
如果埋怨多.
你前路都是不上下的.
但你感恩多的時候.
我今天早上來的時候.
幾乎找不到教會的感受堂.
因為我想.
其實我在你們有些人.
未出生的時候已經住在這裡.
你不相信.
1980年的時候.
你出生的時候.
我住在荔枝南路.
當時教會的祈禱會在我家進行.
分居祈禱會.
但我今天早上說到.
我有一個遺憾的事.
昔日這裡這麼好的地方.
可以神運跑步.
我沒有跑過.
我連單車都沒有踩過.
因為沒空.
過去的事.
我只有感恩.
我感恩在那個地方住的時候.
我在三層樓高.
我上面那層樓我自己做書房.
我每天都去靈修.
看景色.
在裡面.
還有一樣東西.
我兩個小孩和我一起.
在這個地方住.
當有宣教士來香港找.
有些機構說.
宣教士來.
家庭可以接待.

$^{521}$我和他們兩姐弟說.
有個宣教士姐姐.
兩個想來.
你們兩個人合一間房.
給她一間房行不行.
好.
這些是他們從小要學習的.
所以我鼓勵.
感受法院有房間的人.
如果將來宣教士來的時候.
接待一下宣教士.
我前兩天參加會議.
說到將來怎樣去做宣教工作.
我就和其中一個.
將來要做宣教工作的人說.
我說中國人的家庭.
如果不開放家庭.
很難讓下一代做宣教士.
中國教會經常說宣教大會.
其實動不了人很多.
但一旦基督教家庭肯開放.
讓宣教士去的時候.
最得益的是什麼人.
最得益的是小孩.
他接觸到宣教士.
聽他講宣教故事.
已經夠了.
聖靈會工作.
所以今天.
我相信神給我自己.
藉著神的道.
得到益處.
得到福氣.
我想告訴大家.
到底這個道就是福音.
這個福音是什麼福音呢.
請看下一節聖經.
福音是一個永遠的福音.
不要將全福音當作.
人有難有需要.

$^{561}$福音不是那個黑蓋子想掏錢的時候.
給他一筆錢.
福音是永存的.
你看聖經怎麼說.
到將來有一個天石空中.
還要存一次.
這是永遠的福音.
是給地上各方各族各國的人.
所以福音不是只為今生.
不是信耶穌得今生.
信耶穌得永生.
這福音就是耶穌.
所以耶穌在約翰福音十四章第六節怎麼說呢.
我就是道路真理生命.
若不藉著我.
是沒有一個人.
可以重見天父.
所以除耶穌以外.
別無拯救.
就算從基督教思想孕育出來的.
任何一種理論主義神學都好.
若不是講耶穌的話.
都不能引你進入永生.
今時今日.
所謂基督教這三個字.
我講以下說話不是詆毀.
我是有時覺得不開心.
第一次世界大戰.
第二次世界大戰.
最近打得這麼激烈的俄烏戰爭.
這些牽涉參與的國家.
都是什麼國家來的.
伊斯蘭國家.
不是.
全部都是基督教國家.
有沒有想過.
最近我聽到弟兄.
他預備講迪基督這個話題.
他跟我討論.
可能我告訴你.

$^{601}$我跟他討論的時候.
我說將來迪基督.
不可能.
我說不可能.
我不知道.
是那些非基督教國家產生的.
伊斯蘭,印度教.
或者是那些無神論的地方.
迪基督要產生出來的地方.
但可有可能.
但我自己估計.
是基督教國家裡面.
正正就是一個信仰的危機.
你不是信這個道.
你信了道理,理論.
最近我舉個例子.
可能你會實在經歷這些.
我很擔心.
我也很懷疑.
我很著緊.
信耶穌的人.
到底你的信仰是虛擬還是實體.
我來過這裡.
幾次了.
有幾次來的時候.
幾乎是網上做的.
我請問.
你前一陣子見到我在屏幕上出現說話.
和今天站在這裡說話.
有沒有不同呢.
一樣是這個人.
但你知道屏幕上出現的時候.
可以很虛擬的.
今天你信這個耶穌.
接受這個福音,這個道.
是真實的呢.
還是虛擬的呢.
什麼叫虛擬呢.
真實和假的分別在哪裡呢.
真是真,假是假.

$^{641}$但虛擬的東西.
可以半真半假.
我舉個不是你喜歡聽的.
但給你知道的例子.
如果你發現一張紙幣.
那張紙幣裡面.
那些數目字.
你喜歡那個數目字多還是少呢.
一零.
一零零.
一零零零零零.
一零零零零零零零.
你喜歡哪一個呢.
你說呀.
我當然喜歡一零零零零零那些啦.
那你拿韓幣.
你笑你知道的.
因為我最近經過韓國.
我太太說.
六七年前去過.
不是幾年前去過韓國.
說我還剩下一些紙幣在那裡.
說那些不如你去韓國機場買東西買了它.
我打出來看.
嘩這麼大銀碼的.
然後去到那裡發覺不是的.
喂.
那些.
這些銀幣你懂得分.
還有一種這些紙幣呢.
那些銀碼很厲害的.
過百萬千萬的.
不過燒得那些.
年輕人不知道什麼叫燒得那些.
我懶得說出來.
是不是.
你知道那些銀紙是燒得的.
你會不會拿手上手.
你會不會洗那張銀紙的.
不是銀紙.

$^{681}$那張燒得的紙幣.
你不會燒的.
你知道是假的.
但最糟糕的就是.
幾年前香港曾經說有人偽造.
偽造香港一千元紙幣.
不知道是2001年還是2010年.
所以說完之後.
我拿錢我不敢拿一千元.
我怕那個出銀機都搞錯了.
那我拿著寫文書.
我寧願拿五百元我不拿一千元.
是不是.
偽鈔.
偽真度越高.
你越容易上當.
是不是.
偽鈔只是一個影印機.
影印兩方面給你.
彩色給你的時候.
你分別出.
道具你分別出.
但偽真度越高的東西.
你越容易上當.
今天信仰對你一樣偽.
你自己不去查考聖經.
你不認識耶穌.
只聽人說這個說那個.
將時事論,時事資訊說在你面前.
大隊之後信耶穌.
你今生有福氣.
你信不信.
弟兄姊妹.
今天我相信.
我們所接受的福音.
若果這個福音.
你不全面接受.
你接受一部分.
不接受一部分.
或者這麼說.

$^{721}$你相信一部分.
好像世人相信一樣.
你不會相信還沒發生的.
看不到的東西.
神賜給我信心.
我怕你的信仰是虛擬的.
一有虛擬的信仰.
你不要將來.
去到天堂門口.
天國門口的時候.
大跌眼鏡.
懊悔.
聖經上有這樣的記載.
我相信有一個人物.
如果你信聖經的記載.
你都會看他很不值的.
就是羅德的太太.
是不是.
她不是自己走的.
她是被天使拖著走的.
天使拖她走.
離開了那個.
天火焚成的地方.
她最後都會.
走不掉.
要不離地地.
要建立在地上.
很多人說.
我們今天所說的.
信耶穌.
說你很離地.
我說真真正正的信仰.
一定不離地.
一定是貼地.
同時也是很離地.
為什麼.
因為歷史告訴我們.
如果你在信仰上.
是半真半假.
自己感受到.

$^{761}$有些著數.
有些考驗的時候.
懷怨上帝.
帶你出埃及.
十個災經過紅海.
去到迦底斯.
說要去的時候.
那些人說不行.
就怨言來了.
一沒水就說要喝水.
一沒肉吃.
上帝就給你吃.
真的.
結果那些不信的.
不是他們完全不信.
是信一點不信一點的人.
全部倒閉在曠野.
沒得進迦南地.
所以提醒智慧.
求主讓我們看到.
今時今日.
如果我們不是在至聖的真道上.
得著造就的話.
恐怕我們找不到耶穌.
成為我們道路真理生命.
你去不到父神那裡.
以下我和大家.
透過這些聖經的經文.
和你重溫.
重走一次.
看看你現在是.
踏在人生道路上.
還是在永生道路上.
個個都走人生路.
因為我說話.
你都在聽.
不過這條路.
聖經怎麼說.
有一條路人以為正.
結果成為死亡之路.

$^{801}$知不知道.
一個正字.
多少話.
你寫過.
多少話.
五話.
知不知道寫個Y字.
知不知道.
多少話.
在正字上加個什麼字.
就成為Y字.
很簡單.
不正就Y.
正的字只有一條.
就是耶穌基督.
你不走在永生道路上.
你一定走歪路.
走歪路的人.
去不到永生.
去不到.
這個福音.
是耶穌基督自己.
這個福音是永遠的.
羅馬書第一章16句說.
我不以福音為恥.
這個福音本是.
神的大能.
什麼叫大能.
你想見一些奇怪的東西.
見一些人做不到的東西.
但神的大能.
這條話是要救.
一切相信的人.
不是救.
一切好人.
不是救一切行善的人.
不是救一切勤力的人.
不是救一切精打細算的人.
我不是說那些不好.
福音是給那些.

$^{841}$需要的.
給那些需要.
而他們自己.
能夠謙卑在神面前的人.
這些人.
神要用大能.
不是改變他們的生活.
神要用大能改變他們的生命.
生命的改變.
才有豐盛的生活.
今天很多人不知道.
要福音很簡單.
廉價福音.
總之我相信.
我相信.
就像坐了扶手電梯.
升降機.
一樓去二樓.
由地球去到將來.
永恆的天國.
不是的.
聖經告訴你.
福音就是道.
是整個華語.
大家很熟悉的.
約翰福音一章一字.
太初有道.
道與神同在.
道就是.
神.
如果說.
耶穌將道顯明.
他將神顯明.
那麼耶穌就是道.
道就是耶穌.
有些聖經我沒有打出來.
約翰福音六章六十三節.
耶穌說.
我講的話就是.
靈就是生命.

$^{881}$耶穌講的話就是.
神自己.
這個道.
他講的話就是使人的靈魂.
蘇醒的那些話.
他講的話就是.
引導你行永生之路.
所以做基督徒.
成為神的兒女.
信耶穌.
你可以不讀聖經.
不聽道.
只喜歡聽道理的話.
我真的替你擔心.
我擔心我不敢講實.
為什麼呢?我不知道內心世界是怎樣.
你是沒有生命的.
你的靈魂還沒有蘇醒.
詩篇二十三篇.
懂得背了.
懂得背了嗎?.
不懂背很危險的.
因為第三第四節.
你背不成.
將來的日子很難過.
第三節怎麼說.
他是我的靈魂蘇醒.
為自己的命.
引導我走二路.
我雖然.
行過死刃的幽谷.
也不怕遭害.
這句話.
我會經歷的.
我不知道你會不會經歷.
有人將死刃的幽谷.
比喻成為.
這次我大病病到幾乎要死.
醫生都說沒救了,這次我活過來了.
No,那個還不是死刃.

$^{921}$那個只是大病的幽谷.
你人生一定要.
走過那條死亡的通路.
一定走的.
沒有人可以跟你同行.
你自稱的人沒辦法跟你同行.
專科醫生.
自己都走不過.
你唯一可以跟你同行的.
是誰?.
因為你與我同在.
那當時是不是耶穌.
在那些伊斯蘭人見到.
穿著白袍白衣服.
出現在你眼前.
No,他的話.
就是他自己.
當他的話在你心裡.
當他的話在你腦海裡.
你知不知道.
你一個耶穌就活在你心中.
活在你腦海中.
不知道你們有沒有經歷過.
我聽過不少見證.
包括我自己經歷過.
我睡不著的時候.
我擔心的時候.
我要睡覺.
幾乎睡不著的時候.
我就相信聖經說.
就也思想耶和華律法的.
這人便為.
什麼?.
我說神啊.
我是你所愛的.
我只不過想成為一個有福的人.
所以答案來了.
那你想一想.
一想的時候.
越想睡覺.

$^{961}$就越睡不著.
我說好啊主啊.
讓我試一試你的話語有多大功效.
菲拉比書.
四章第六第七節.
應當一無掛慮.
只要凡事.
將你所要的.
高騷神.
神賜出人意外的平安.
必在基督耶穌.
裡保守你的心懷意念.
我說主啊.
這聖經是你所說的.
我現在真的有點擔心.
有點煩惱.
睡不著.
如果你的話語是這麼真實.
我就現在向你禱告.
所要的高騷神.
我現在想睡覺.
你知道.
這個果效是什麼嗎?.
要我起床再吞兩口安眠藥.
這樣睡覺嗎?.
你起床了也醒了.
當我這樣禱告完畢之後.
不知為何.
我都不用數綿羊.
我已經進入夢鄉.
你說你怪人.
我現在經常收到學生的見證.
和他們的一些東西.
我是很受激勵.
有一對學生.
後來兩人結成夫婦.
在教會裡面都很熱心.
在事業上都有成就.
兩人都是從事金融界會計行業.
有一天.

$^{1001}$她的丈夫去參加.
聖福查經班.
那天下雨.
她的腳有點癢.
她去抓一抓.
抓傷了.
她就這樣回家.
過了一會.
發炎.
紅腫發燒.
她立即去醫院.
去醫院之後一查.
立即轉去聖天府醫院.
因為中了食肉菌.
醫生檢查完之後.
跟她的家人說.
她急症室做手術.
如果手術成功.
切斷了腳.
救回一條命.
如果不做這個手術.
恐怕命不保.
在這個混亂狀況裡面的時候.
這個同學.
這個女同學.
她的好朋友.
是我的姊妹.
就WhatsApp給我.
校長有這樣的事幹.
請你祈禱.
我之後就祈禱.
但我不單止祈禱.
我心想誰可以幫到她.
專科醫生給她的選擇.
就是要不進去切腳.
不切腳就等.
有一次聖經.
在我腦海裡面出現了.
因他受的變傷.
使我得意志.

$^{1041}$因他受的懲罰.
使我得平安.
是不是.
跟這個有什麼關係.
我禱告中間.
我真的為她祈禱.
但我不是就這樣祈禱完.
我把禱文和聖經合在一起.
發短訊給我的同學.
她現在忙得七彩.
我當時不知道.
適當的時候.
把禱文給她.
她馬上收到禱文.
WhatsApp給我女同學.
我女同學.
又要操忙.
決定幫不到忙.
結果.
在大約五,六點鐘.
就推進去手術室.
做手術.
一做起碼幾個小時.
如果能做到.
做不到就早跑出來.
所以到.
大約深夜.
她就唯有回家睡覺.
回家睡不著覺.
怎麼睡.
當她由.
十二點,一點,兩點.
到三,四點都睡不著覺的時候.
她不知道怎樣.
她就拿起了.
床頭的手機.
打開手機.
就是同學.
發的禱文給她.
我在那裡略為解釋一下.

$^{1081}$我說.
耶穌的邊廂.
已經受了.
是歷史.
耶穌不需要今天為你受邊廂.
是不是.
二千多年受過邊廂.
是不是.
耶穌的懲罰.
已經成就了歷史.
我說.
歷史的事已經發生了.
你不能改變.
但現在要改變的就是.
你能不能得到平安.
你現在不平安.
你想不想得到平安.
現在這個病是這樣的病.
你想不想得到醫治.
這是未發生的.
要將歷史上的事.
轉成你的期望.
加上能夠成就的話.
只有一個字.
就是信不信.
所以很多人祈禱.
其實你帶著不信的信心.
不信的惡心祈禱.
想上帝答應你的話.
你都很危險.
當他看到的時候.
我現在是一個平安.
只要我相信耶穌.
已經為我受了懲罰.
我就有平安.
如果我要得醫治.
我相信耶穌已經為我受了邊傷.
歷史上的.
我就得去得醫治.
所以歷史發生的事實.

$^{1121}$和成不成就在你身上.
不是其他的因素.
主要就是你.
相不相信.
接不接受神的道.
他看完之後.
他說主啊,我現在想睡覺.
我信.
當他說完這個禱文之後.
他就睡著了.
他睡著了之後.
到第二天早上六點多.
家人就說.
你老公還在醫院.
他馬上衝回去.
衝回去的時候.
手術場就出來.
出來的時候看到人在.
就安心了一半.
再看到腳還綁著.
就安心了.
頂著門.
這些是.
豐文有你.
和親眼看到你.
是基於一個人的信心.
做多少事.
這些具體的.
生命之道的見證.
不是叫我們什麼都不做.
不是.
你沒有力量做.
你唯有用.
神藉著聖靈賜給你的信心.
去成就.
神的應許.
歷史的東西.
不是你信不信的.
全世界都和他說.
耶穌已經降生.

$^{1161}$他說不信你不用理他.
你圍著祈禱就行了.
對整個元老來說.
耶穌已經受了死.
他不信就由他.
那你復活節過來.
為什麼要過復活節呢.
為什麼要過聖誕節呢.
這些是歷史.
信與不信和你.
得到豐盛的生命.
得永生是沒有什麼關係的.
如果你信的是.
耶穌的降生.
接受的耶穌的降生.
耶穌的受死.
耶穌的復活.
我解釋這個給你聽.
耶穌的升天.
聖經記載了.
怎麼可能.
起碼要穿太空衣.
講得不好.
太空衣.
耶穌的.
講得不好.
耶穌的屍體.
怎麼脫殼而出.
因為我師父還在.
不見了一個人.
不見了血肉之軀.
已經殘了.
然後他又出現了.
怎麼信得過.
你用理性.
想去解釋所有神的話語.
你不會得救的.
你想用人的方法.
來成就神的事情.
不會成功的.

$^{1201}$亞伯拉罕.
撒拉這個家庭.
蒙上帝應許.
說要輪福.
萬國也因你得福.
他們信的.
不過他們用人的方法.
想來成就神的工作.
一次都沒有成功.
第一次.
我好像跟你講過.
耶和華神.
在上世紀十六章.
對亞伯拉罕說的時候.
亞伯拉罕說.
以利爾謝.
在我家出生的.
我不知道他有沒有契他做兒子.
差不多要承繼我所有遺產的.
就是他了.
你要賜福就他了.
耶穌說不是.
是你和撒拉的.
結果.
耶穌是一個很薄的人.
你看他和.
以撒娶老婆.
那件事.
你都沒有那麼的信心.
如果我祈禱.
有個姑娘出來.
姑娘不但給我水喝.
還給我下塌喝.
下塌不是喝一兩瓶水.
是喝一卡車水.
真的被他祈禱.
但不是這個.
撒拉更聰明.
亞伯拉罕能生.
我不能生.

$^{1241}$給我下塌.
還搞我抱出來.
是不是.
直到最後他們兩個完全心裡暗笑.
耶和華神藉著使者顯明.
和他講清講楚.
到九十九歲的時候.
從七十五歲到九十九歲.
那個關鍵中間.
到二十五年時間.
以撒真的生出來之後.
他們兩個都皆大歡喜.
所以以撒的名字有歡笑的意思.
我自己往往提醒自己.
亦提醒一些.
做神工作的人.
不要用自己的頭腦.
想幫上帝.
完成他的計劃.
沒可能.
不是我們不想.
你要設計要想.
不過當你越設計越想.
之後你就越害怕.
你做大生意.
你敢不敢投資.
你投資後收穫如何.
你不知道.
你越想越大越害怕.
所以在神裡面.
做神的工作只是一個可能性.
你越害怕的時候.
正正不是要你越想辦法.
越想數據.
你越害怕的時候.
做一樣東西.
懇切禱告求主賜給你信心.
正經上是這個意思.
所以.
當我們知道福音是神的大能的時候.

$^{1281}$他救一切相信的人.
是你.
你如何相信.
所以我們從聖經看到.
道就是神.
太初有道.
道與神同在.
道就是神.
下一節聖經說.
信子的人有永生.
你信就行了.
信者不在永生.
神的真道在他身上.
好,接著下面.
我和大家看福音的內容.
約翰福音三章.
三到七節.
聖經怎麼說.
耶穌放大了下一句.
下一章就放大了.
耶穌回答說.
我實實在在告訴你.
人如果不重生就不能見神的國.
重生.
不要以為.
舉手.
你要必須生命的更生.
第四節.
尼哥提姆說人已經老了.
怎麼可以重生呢.
怎麼可以重復出來呢.
第五節耶穌說.
我實實在在告訴你.
人如果不是從水和聖靈生的.
就不能見神的國.
一講到聖靈.
不要害怕.
哎呀.
什麼派什麼派.
不用理會那些.

$^{1321}$聖靈就是神的話語.
就是神自己親自的感動.
接著第六節說.
從肉生生的就是肉生.
從靈生的就是靈.
你不從聖靈.
不從神的話語跟你更生的話.
你始終是肉聖靈.
你怎麼做好人.
都真不要拿去.
都是很虛怯的.
甚至乎可以說很虛擬的.
第七節耶穌說.
你們必須重生.
你們不要以為稀奇.
講到這裡的時候.
跟大家很快看看林多前書十五章.
有幾節聖經.
這個事態是下一章我放大了.
弟兄們.
我告訴你們說.
血肉之軀不能承受神的國.
扭壞的不能承受不扭壞的.
這個一定會扭壞的.
下面一句.
我擺一件奧秘的事.
不是每個人都明白的事.
除非聖靈感動你.
告訴你們我們不是都要睡覺.
這個睡覺這個字當時講.
不是指睡覺.
是指死的意思.
乃是都要改變.
裡面改變到外面.
好.
下一節.
就在一煞之時.
三眼之間.
浩同會此吹響的時候.
浩同要響.

$^{1361}$死人要復活.
成為不扭壞的.
我們也要改變.
下一章.
這必扭壞的.
是總要變成不扭壞的.
這必死的.
總要變成不死的.
如果要進入永生.
一定有這個過程.
沒有這個過程進不了永生.
這必死扭壞的.
既變成不扭壞的.
既變成不死的.
那時經上所記.
死被德性吞滅的話就應驗了.
神話一定要應驗的.
將來死是不可以再影響人.
因為人永遠不死.
所以在哥林多前書有講到.
死你的德性權勢在哪裡.
你的毒勾在哪裡.
然後就解釋到.
為什麼死因為人犯罪.
罪的功架是死的.
權勢是律法.
律法阻礙了我們.
所以聖經記載.
使我們直著我們的主耶穌基督德性.
沒有任何辦法.
沒有任何的主義.
沒有任何人的行為.
包括你的努力.
使你靠著德性.
唯有靠耶穌基督.
道路真理生命.
這是神的道.
所以我親愛的弟兄們.
無要堅固不可搖動上上.
極力多作主公.

$^{1401}$因為知道你們的身體才是奴婦.
心靈奴婦在主裡面不是圖源.
所以因勤作主的工作.
很多時候被性靈感動.
你要做就去做.
不做你經歷不到神話的大能.
時間到.
再看下面那段.
哥林多前書第十五章.
第十八節.
寄傳耶穌基督從死裡復活.
為什麼你們中間有人說.
沒有復活這件事情呢.
不是說現在.
初期教會也是這樣.
耶穌在生的時候.
發到的人都不信復活.
遵守律法.
不信復活.
十三節.
若沒有死人復活的事.
基督也就沒有復活了.
這句話.
說耶穌的復活和人的復活.
拉上關係.
你不信耶穌已經復活了.
有永生.
永遠存在.
你也不會這樣做.
耶穌降生是道成肉身.
神變成人.
耶穌的復活是將人.
變回成永遠的榮耀身體.
那樣東西你不信.
也不會發生在你身上.
那你有永生的得救嗎.
第十四節.
若基督會有復活.
我們所傳的便是謊言.
你們所信的也是謊言.

$^{1441}$怎麼謊言呢.
第十五節.
並且顯明我們是為神妄作見證.
你告訴別人.
我是見證的.
因為復活的生命.
不能夠今時今日在地上.
直到你將它顯明出來的時候.
將來都不會顯明.
因為我們見證神是叫基督復活了.
若有死人真不復活.
神也就沒有叫基督復活了.
那神有沒有叫基督復活呢.
如果有的話.
你要信的時候.
這個影響到你了.
第十五節.
並且明顯我們是為神.
剛才我讀了.
叫基督復活.
第十六節.
因為死人若不復活.
基督也就沒有復活了.
第十七節.
基督若沒有復活.
你們的信便是途人.
你們仍在罪中.
我罪得赦免.
不信耶穌復活.
你的罪仍然在.
我已經相信耶穌降生.
受死.
但復活是一半一半.
我不相信他會再來.
因為這個世界只要改成天堂.
便可以了.
你們仍在罪裡.
一個仍在罪裡的人.
有沒有永生?.
所以十八節.

$^{1481}$就是在基督裡.
誰了的人也滅亡了.
你不信的話.
那些已經過去死去的人.
是否也滅亡了?.
你也是將來滅亡的.
所以哥林多傳書第十二章第九節.
聖經怎麼說呢.
信心不是你與生俱來有的.
聖經說.
是蒙聖靈賜的信心.
由聖靈賜給你.
當你認罪悔改的時候.
聖靈在你心中.
就給你信心.
相信耶穌已經復活.
臣子修行四章二十一節.
一起讀好嗎?.
他要按著那能叫萬有.
歸服自己的大能.
將我們這卑賤的身體.
改變形狀.
和他自己榮耀的身體相似.
所以耶穌登山變貌的記錄.
在啟示錄約翰見到.
白魔海島上的耶穌.
見到的顯現.
這些是實在的東西.
問題就是.
我和你從生得救的人.
將來都要經歷這些東西.
最後一節一起看.
一起讀羅馬書八章十一節.
然而.
叫耶穌從死裡復活的靈.
若住在你們心裡.
那叫基督耶穌從死裡復活的.
也必藉著住在你們心裡的聖靈.
使你們必死的身體.
又活過來.

$^{1521}$這些身體必死的.
但將來能夠活過來.
是靠著基督的靈.
所以弟兄姊妹.
我將這個福音.
神的道.
我會藉著神的話.
在你們面前讀給你們聽.
你們自己已經讀了.
但問題在於.
你們信一點不信一點.
還是全信了呢?.
這關乎你們的夢幻.
你們想信一點不信一點的話.
我告訴你.
你們的人生不會豐盛.
你們信一點不信一點的話.
我很擔心你們的永生會有問題.
但當你們肯完全相信的時候.
聖經說你們一定無誤服.
願主賜福給大家.
藉著神的話語.
好好反悔反覆思想.
你們查考聖經.
因為聖經裡面有永生.
給我作見證的便是這經.
約翰福音五章三十九節.
耶穌說.
好好讀聖經.
也求聖靈幫助你.
在今生直到永遠.
都因著福音.
聽信祂的話語.
成為一個蒙福的人.
我們祈禱.
藉著你自身的真道.
讓我們知道我們的蒙福.
是和你的話語有關.
我們的蒙福.
藉著你的道.

$^{1561}$使我們成為聖潔.
求主恩待.
讓我們每個人.
因著你的話語.
聖靈的感動.
心裡火熱.
願意樂意的去事奉你.
所以聽我們的禱告.
風言壽命之機構.
\newpage



\section{哥林多後書 4:1-18}
\label{sec:FIbxv36OMMk}
\textbf{2023-08-20 午堂《傳道者的心》哥林多後書 4\_1-18 講員:鄧泰康傳道}
\newline
\newline
連結: \href{https://youtube.com/watch?v=FIbxv36OMMk}{\texttt{https://youtube.com/watch?v=FIbxv36OMMk}} ~~~~ 語音日期: 2023-08-22
\newline
\newline
\hyperref[sec:kFsvDzgp4nI]{\small{< < < PREV SERMON < < <}}
~
\hyperref[sec:index_chronic]{\small{[返順時目]}}
~
\hyperref[sec:index_scriptual]{\small{[返順卷目]}}
~
\hyperref[sec:cMmv_XzlkAc]{\small{> > > NEXT SERMON > > >}}
\newline
\newline
哥林多後書 4:1-18
\newline
\begin{longtable}{cl}
\hline
\hline
章節 & 經文 (和合本修訂版)\\
\hline
4:1 & \begin{tabularx}{0.7\textwidth}{X} 所以，既然我們蒙憐憫受了這事奉的責任，就不喪膽， \end{tabularx} \\ \\ \relax
4:2 & \begin{tabularx}{0.7\textwidth}{X} 反而把那些暗昧可恥的事棄絕了，不行詭詐，不曲解神的道，只將真理顯揚出來，好在神面前把自己推薦給各人的良心。 \end{tabularx} \\ \\ \relax
4:3 & \begin{tabularx}{0.7\textwidth}{X} 即使我們的福音被遮蔽，那只是對滅亡的人遮蔽。 \end{tabularx} \\ \\ \relax
4:4 & \begin{tabularx}{0.7\textwidth}{X} 這些不信的人被這世界的神明弄瞎了心眼，使他們看不見基督榮耀的福音。基督本是神的像。 \end{tabularx} \\ \\ \relax
4:5 & \begin{tabularx}{0.7\textwidth}{X} 我們不是傳自己，而是傳耶穌基督為主，並且自己因耶穌作你們的僕人。 \end{tabularx} \\ \\ \relax
4:6 & \begin{tabularx}{0.7\textwidth}{X} 那吩咐光從黑暗裡照出來的神已經照在我們心裡，使我們知道神榮耀的光顯在耶穌基督的臉上。 \end{tabularx} \\ \\ \relax
4:7 & \begin{tabularx}{0.7\textwidth}{X} 我們有這寶貝放在瓦器裡，為要顯明這莫大的能力是出於神，不是出於我們。 \end{tabularx} \\ \\ \relax
4:8 & \begin{tabularx}{0.7\textwidth}{X} 我們處處受困，卻不被捆住；內心困擾，卻沒有絕望； \end{tabularx} \\ \\ \relax
4:9 & \begin{tabularx}{0.7\textwidth}{X} 遭受迫害，卻不被撇棄；擊倒在地，卻不致滅亡。 \end{tabularx} \\ \\ \relax
4:10 & \begin{tabularx}{0.7\textwidth}{X} 我們身上常帶著耶穌的死，使耶穌的生也在我們身上顯明。 \end{tabularx} \\ \\ \relax
4:11 & \begin{tabularx}{0.7\textwidth}{X} 因為我們這活著的人常為耶穌被置於死地，使耶穌的生命在我們這必死的人身上顯明出來。 \end{tabularx} \\ \\ \relax
4:12 & \begin{tabularx}{0.7\textwidth}{X} 這樣看來，死是在我們身上運作，生卻在你們身上運作。 \end{tabularx} \\ \\ \relax
4:13 & \begin{tabularx}{0.7\textwidth}{X} 但我們既然有從同一位靈而來的信心，正如經上記著：「我信，故我說話」，我們也信，所以也說話； \end{tabularx} \\ \\ \relax
4:14 & \begin{tabularx}{0.7\textwidth}{X} 因為知道，那使主耶穌復活的也必使我們與耶穌一同復活，並且使我們與你們一起站在他面前。 \end{tabularx} \\ \\ \relax
4:15 & \begin{tabularx}{0.7\textwidth}{X} 凡事都是為了你們，好使恩惠既藉著更多的人而加增，感恩也格外顯多，好歸榮耀給神。 \end{tabularx} \\ \\ \relax
4:16 & \begin{tabularx}{0.7\textwidth}{X} 所以，我們不喪膽。雖然我們外在的人日漸朽壞，內在的人卻日日更新。 \end{tabularx} \\ \\ \relax
4:17 & \begin{tabularx}{0.7\textwidth}{X} 我們這短暫而輕微的苦楚要為我們成就極重、無比、永遠的榮耀。 \end{tabularx} \\ \\ \relax
4:18 & \begin{tabularx}{0.7\textwidth}{X} 因為我們不是顧念看得見的，而是顧念看不見的；原來看得見的是暫時的，看不見的才是永遠的。 \end{tabularx} \\ \\
[1ex]
\hline
\hline
\end{longtable}
$^{1}$各位姐妹早晨.
見我戴著口罩 大家就知道我有點感冒.
我昨晚的團期.
因為上星期那位姐妹去做主席.
我跟她說 你記不記得上次我們在五堂講道.
之後我就在旁邊說 待會你找我祈禱.
不過也是一個見證時刻的時候.
因為每次都咳到脫肺.
講道那段時間是不咳的.
所以每次曾牧師編我回來五堂講道我都咳的.
曾牧師考慮一下.
今天跟大家分享這段經文.
為什麼叫傳道者的心.
其實如果大家知道在哥林多後書.
如果大家都明白 大家讀過一些上過主要學.
你會知道原來哥林多的書信是有四卷的.
有1,2,3,4.
現在我們看的前書是第二卷.
我們的後書是第四卷.
所以你會發覺去到第四卷的時候.
其實保羅在之前跟我們處理了很多問題.
關心了很多事情.
其實去到第四卷的時候.
保羅其實是跟這班人.
你可以說是一個很真誠的表白.
因為當時有一個信徒懷疑他的誠信.
懷疑這個人所做的事.
所以他在裡面說了很多很多.
所以去到這一章的時候.
其實他講了很多我在傳福音裡面.
我自己真心的話和我們的心態.
當講到傳道者的時候.
我想再強調.
不是只講傳道人 不是講牧師.
我是想大家要記得.
傳道者是所有的信徒.
因為當你信了耶穌之後.
上帝讓你明白一個很重要的事實.
屬於神的人在世界上最核心的意義.
就是以你的生命去見證神.

$^{41}$讓人去認識上帝.
將人帶到上帝的面前.
這個說話很重要.
所以我再講一次.
屬於上帝的人在地上最核心的意義.
是以你的生命去見證神.
讓人認識上帝.
將人帶到上帝的面前.
所以在座裡面.
每一位弟兄姊妹你們都是傳道者.
所以和你身邊的弟兄姊妹一起鼓勵.
我們是傳道者.
我們都是傳道者.
和你身邊的弟兄姊妹一起鼓勵.
我們都是傳道者.
所以傳道不是.
我想再強調不是教會的責任.
不是牧師的責任.
不是傳道人的責任.
弟兄姊妹很流行的.
我有個什麼什麼的親戚.
我有什麼什麼的朋友.
鄧傳道你可不可以幫我和他傳福音.
其實我當然很樂意.
你明白嗎 我很樂意.
但反過來我腦裡很多時候.
第一個彈出來的就是.
為什麼你不傳.
你都可以講的.
你都可以和他分享的.
有些人就說.
是的 我都明白的.
自己家裡的先知是不受歡迎的.
所以我家裡很感恩.
我家裡有七兄弟姊妹.
我們七兄弟姊妹都信了耶穌.
但沒有一個是由兄弟姊妹傳給他們的.
因為我們各自很有性格.
全部都是朋友之間.
但你記住都是交際.

$^{81}$你幫我傳 我幫你傳.
起碼你都肯去傳.
這個是重要的.
所以你就會發覺.
今天我想和大家去看的時候.
我們一起從保羅.
他自己的心態.
去學什麼叫傳道者的心.
其實是什麼一回事.
我會將它分成.
剛才那段經文有幾個主題.
第一件事是什麼呢.
我們是對傳道的定位.
保羅一開始的時候他就說.
其實我們今天能夠得著傳道者這個身份.
完全是因為神的憐憫.
是神的憐憫.
其實從來沒有一個人在本質上.
他是配去事奉上帝.
從來都沒有.
因為我們每一個都是罪人.
我們沒有一個人去配.
所以其實這個身份.
是上帝透過他的憐憫.
而去呼召我們進到他的事奉當中.
弟兄姊妹.
這個心態是非常重要.
所以其實今天我經常都說.
我們在生活.
我們在我們所有的能力上面.
我們很努力去追求.
譬如師班.
我們要把很多時間放進去.
我們要練到最好.
我們去講說話的技巧.
廚房的技巧.
我們要花時間去練的.
大家都明白.
如果你走出來去講.
我都和大家講過.

$^{121}$我以前的傳道.
我的老師.
是神學院的老師.
他教我們講道學.
他第一句和我們講的就是.
當然我們聖經是重要的.
但是請大家學好你們怎樣講編道.
為什麼呢.
就算你編道是多好都好.
但是你講到大家都睡了.
是沒有意思的.
所以原來是要學的.
要訓練的.
但是我們去精益求精.
我們做到最好的原因.
我們不是要去增敬.
我們不是要去比較.
我們不是要去爭第一.
而是我們要學到最好的.
以最好的姿態去事奉上帝.
這個才是真實的.
正因為這個緣故.
因為我們每一個人都是蒙憐憫.
才能夠有機會去參與這個事奉.
所以我形容是很難得的.
所以很多時候.
我作為傳道人.
在教會裡面.
我們去鼓勵弟兄姊妹.
邀請弟兄姊妹一起事奉.
我同樣都會有一個問號.
弟兄姊妹好像對於事奉.
是有一種選擇.
或者我形容.
好像我很難走出那一步.
但是我不知道.
如果今天有一個人.
是你很尊敬的.
是你很愛的.
他邀請你一起來參與他的宴會.

$^{161}$或者邀請你去參與他的工作.
我不知道你會有什麼感覺.
我不知道你會不會即時答應.
這麼難得的機會.
我很記得.
我從小開始玩話劇.
我由學校的話劇團.
一直玩到地區的話劇團.
當時我們在那個年紀的時候.
我們話劇團的天花板.
就是香港話劇團.
那個是全香港公營.
我只會覺得.
我只是有份去在旁邊看他們表演.
但是我記得有一次.
我不知道什麼原因.
我收到一個電話.
如果是玩話劇的人應該認識他.
有一位叫麥秋先生.
他是話劇的一個好前輩.
突然間我收到電話.
我受寵若驚.
請問你是不是鄧太康先生.
我說是的.
我聽到某某某導演提起你的名字.
我們很想邀請你去幫忙.
在一套話劇裡面幫忙.
我突然間很震撼.
我想說我是無管動.
接著其實在整個過程裡面.
我是很開心的.
我想說.
我從來沒有想過.
原來我有機會去到一個.
我認為這麼榮譽.
這麼有質素的一個地方.
我可以有份參與.
雖然那個不是一個主角.
是一個小角色.
但是問題是在裡面的時候.

$^{201}$我覺得很快樂.
弟兄姊妹.
今天上帝就是同樣的道理.
究竟他在你的心目中.
他是不是你都認為.
因為得到憐憫.
你才有機會參與事奉.
或者甚至你是得到這個憐憫.
你才有機會去傳福音.
你才能夠做這個傳道者.
如果是的話.
你應該是與有榮焉.
你應該是要感到很光榮.
而寶萊就告訴大家.
因為我們知道這是上帝的憐憫.
所以他一回就回到他的本質.
他的性格裡面.
因此我不氣餒.
因為他說是上帝揀選我的.
所以剛才大家讀的經文.
沒有人可以阻擋我與上帝的關係.
因為是他揀選我的.
我曾經聽過一些人分享.
他說他以前回教會的時候.
當然我不知道他到什麼年紀.
有些弟兄姊妹突然跟他說了一句.
你沒有資格去事奉神.
然後他們受了很大的傷害.
這個陰影直至到今天.
但有機會我與他談的時候.
我跟他說.
我只有一個人可以跟你說.
你有沒有資格.
那個是誰呢.
是上帝.
所以你會發覺寶萊在說.
我不會氣餒的.
他說所以我是絕對會去氣絕我的舊我.
我不會耍手段的.
我不會欺騙人的.

$^{241}$我不會扭曲上帝的說話.
當我接了上帝的邀請之後.
我要成為一個正直無偽的人.
我只說上帝的說話.
很多時候今天經常與我們去爭戰的.
我可以告訴弟兄姊妹.
第一名與你爭戰的是你自己.
是我們的舊我.
我們不知道為什麼舊我裡面的性格.
很容易無端端又走出來.
很容易我們又會走回以前的特質出來.
但是在聖經裡面.
寶萊很清楚.
他說我的舊我是什麼.
已經與罪一起釘在十字架上.
弟兄姊妹是不是.
所以你就會發覺.
我們今天去面對自己的時候.
首先你作為一個傳道者.
你要成為一個正直無偽的人.
第二件事.
他說當你面對那些拒絕你的人的時候.
在聖經當中他一樣提到.
有魔鬼的參與.
請弟兄姊妹記得.
在屬靈的事件上.
有所謂屬靈爭戰.
他說什麼呢.
他說這些魔鬼會蒙蔽了那些拒絕你的人的眼睛.
讓他們看不見耶穌基督的榮光.
看到這個邏輯嗎.
所以弟兄姊妹.
你要對付的是魔鬼.
而不是拒絕你的人.
清楚嗎.
你不應該生氣那些拒絕你福音的人.
甚至攻擊你福音的人.
你不應該生氣他們.
因為背後有魔鬼.
所以你對付魔鬼最好的方法是什麼.

$^{281}$祈禱.
你要以上帝的能力.
你要以上帝的話去對付魔鬼.
耶穌在山上的時候.
同樣都是用神的話去對付魔鬼.
祂一個神蹟都沒有用過.
所以這個是我們今天作為傳道的定位.
所以你只需要做好自己.
你不需要和外面的人去爭競.
你不需要和外面的人去爭拗.
我們不要停留在一種爭辯當中.
我仍然很相信真理就是真理.
我可以告訴弟兄姊妹.
在今天這個現世當中.
有很多人是說歪理多過真理.
你講真理他都會跟你講歪理.
但你不是要和他爭辯到變回真理.
你不需要.
耶穌有一句名言.
大家有沒有看聖經經常都讀到.
有耳可聽就應當聽.
其實耶穌對於那些拒絕的人.
他都沒有那麼好氣.
我不知道大家有沒有發覺.
所以你就會發覺我們今天.
去將福音傳給一些願意去聽.
願意去接受主耶穌基督的人.
對於其他的人他聽不聽.
交給上帝.
第二件事我們要去講的是傳道的內容.
保羅他講得很清楚.
我們要傳的是基督是主.
不是我們.
是他的榮耀.
不是我們的榮耀.
當我看到這段經文的時候.
很多時候我就會馬上想起.
有很多我們的見證.
你記住一線之差.
你看看當我相信了耶穌之後.

$^{321}$上帝多麼祝福我.
我以前是一個很貧窮的人.
但上帝給了我很多機會.
所以我今天成為一個.
很豐足很有能力的人.
你看我以前是很蠢的.
但上帝幫助我.
所以我讀書後來越來越厲害.
我今天讀了多少個博士的學位.
頂子們.
一線之差.
沒錯你裡面有提過神.
但究竟你所傳的是主.
還是你自己的榮耀.
從哪裡可以看到呢.
就是當你發覺別人小看你.
認為很重要的事的時候.
你會不開心.
你會覺得反感.
甚至你會反彈.
你就知道究竟你正在傳主.
還是在傳你自己.
如果那件事是上帝賜給我.
上帝同樣可以拿走.
這是祂的權利.
所以上帝可以在那一刻.
令到我很豐足.
但同樣我也可以接受.
有一天突然間我變得很貧窮.
我不會埋怨上帝.
我們的傳是基督是主.
所以我們應該很清楚.
今天你跟別人講的每一句說話.
你是高舉上帝的榮耀.
上帝在你生命當中.
祂有甚麼作為.
我再強調一線之差.
甚至你可以騙到全部的人.
但相反上帝自己很清楚.
而接著就是因為主的緣故.

$^{361}$你要甘願去服侍人.
即是將道化成生命去活出來.
上帝曾說要有光照在黑暗裡.
所以當你說你要做傳道者的時候.
請你首先成為那個.
被上帝光照黑暗內心的人.
首先你要去完全認識和體會.
上帝的榮耀大能的人.
弟兄姊妹有很多口才.
有很多聖經的知識.
其實我們都可以學回來.
但究竟你對於這些知識.
上帝跟你講的說話.
你有多少的體會.
他是不是真的在你生命當中.
就好像聖經裡面所講.
神的說話是兩刃劍.
將你內心都否開.
去看清楚自己的內面.
而跟上帝坦白講.
當你真的去看自己內心的時候.
你唯一可以跟上帝講的就是.
求主憐憫我.
當你有這種心的時候.
你可以帶著其他的人.
去進到上帝的裡面.
所以後來他用了一個例子.
當你讓上帝照好你的內心的時候.
保羅就說.
所以我就確知自己.
只不過是一個脆弱的啞氣.
是一個這麼寶貝.
放在我這個啞氣的裡面.
這個啞氣沒有任何作用.
我們都不知道為什麼上帝.
會竟然這麼愛我.
會竟然去用這個這麼不配的人.
這麼不足的人.
弟兄姊妹.
當你真正成為傳道者的時候.

$^{401}$你所剩下來的.
就是一份對上帝的謙卑.
因為你很清楚.
所有一切的能力.
都是從上帝而來.
你自己.
你就會越來越體會到.
主怎樣去使用我們這個啞氣.
弟兄姊妹.
這種心態其實和我們今天.
在這個社會上.
很多人所追求.
是背道而馳的.
我們每一個人都很希望.
出到來的時候.
是光芒萬丈.
是閃閃星輝的.
我們很希望每一個人.
都看我們是金氣.
是銀氣的.
其實很多人的內心裡面.
我們不想自己做啞氣的.
但是我再強調.
我經常都說.
我再強調再強調.
你用不用去追求卓越.
要.
但是這個卓越.
不是為了榮耀你自己.
而是為了更好的去事奉神.
而令到這個啞氣.
我就這樣形容.
原本這個啞氣很骯髒.
這個啞氣的本質是不會變.
但是你可以將它洗到.
立立令.
你可以將這個啞氣上面.
去增加一些藝術的氣質.
去成為一件人見人愛的啞氣.
因為在裡面.

$^{441}$它找到寶貝.
所以你越是願意去誇上帝.
在你生命當中那種大能的時候.
你就越能夠可以去展現.
一種生命出來.
保羅說正因為如此.
所以他面對磨難壓迫.
困惑迷茫.
被人追逼.
甚至被打倒.
都從來沒有敗下來.
因為保羅說我在苦難中.
我們身上一直帶著的是耶穌的死.
好讓耶穌的生命.
也都彰顯在我們的身上.
弟兄姊妹原來傳道者.
最能夠去體會的.
也是回到我們福音的原點.
耶穌基督最大的能力.
就是復活的大能.
但是必須透過死.
才可以展現復活大能.
明白嗎.
所以為什麼我們作為傳道人.
經常都聲嘶力竭地去提弟兄姊妹.
弟兄姊妹去死呀.
在神裡面要死得徹底一點.
我們在神裡面.
我們受針的人.
我們在水裡面去象徵的.
我們在水裡面.
在泥土裡面.
我們已經是復活過來了.
你越死得頭徹.
你就越能夠可以展現上帝的大能.
記不記得有一句.
大家很喜歡的經文.
一粒墨紙若不落在地裡死了.
仍舊是一粒.
若是死了.

$^{481}$就結出許多紙粒來.
上帝給我們的能力就是這麼奇妙的.
你越是放下自己.
越是讓自己破碎.
越是讓自己被上帝所用.
上帝就越能夠去用你.
而展現出來的那種能力.
竟然去到一個情況.
沒有任何一個處境壓迫.
可以打倒你.
因為你有復活的大能.
這種大能可以超越你生命裡面.
所有的一切.
所以當今天如果你問我.
我們作為傳道者.
我們傳什麼道.
傳生命見證的道.
你不是跟他講一個三福.
不是跟他講一個四律.
而是你用你的生命跟他講故事.
所以有些人就說.
跟家人傳福是最難的.
我有保守 我有保留.
其實你的家人是最清楚你的壞事.
從小到大看著你的時候.
但反過來.
當你的生命轉變.
也是他見證得最清楚.
我記得我以前的太太.
她當時她被信耶穌.
我再回教會.
後來我回教會.
我重新再去努力.
後來我記得那段時間.
我們學習.
回教會很積極.
學習三福.
學習三福其中一個最難的地方.
就是要背45分鐘的內容.
然後那個人坐定聽你講45分鐘.

$^{521}$最後的時間你就要跟他決戰.
你要做實習.
最好實習的對手當然是找太太.
她還未相信.
她採取的策略就是.
你回教會吧 多好.
我不回的.
然後我就坐定.
無論如何.
你讓我實習行不行.
我就對著她講三福.
講完之後.
其實真的實習.
我就跟她說.
最後的時候我想邀請你.
你信不信耶穌.
實習完.
我打算起身走.
跟著她跟我說 我信.
跟著我很記得.
當時我跟她說 你不要玩了.
我說我實習而已.
我太太可以跟我說.
她說我是信.
我真心信.
接著我問她為什麼.
她說在過去.
我們兩夫婦.
你還未回教會之前.
我們的婚姻是怎樣 你很清楚.
我們的關係是怎樣 你很清楚.
她說我以前用盡一哭二涼三上吊.
沒辦法跟你改變.
但你回教會半年而已.
我看著你變.
我看著你改變.
她說我是完全沒辦法.
不從你身上承認這個上帝是真.
而且很厲害.
所以她就是這樣.

$^{561}$後來她信耶穌 信得很好.
所以弟兄姊妹.
我想強調的不是我們做了什麼.
你明白嗎 連我也沒有信心.
但原來上帝在不知不覺之間.
去改變我們.
而令到我們身邊的人.
去看到我們的故事.
這個就是傳道的內容.
所以弟兄姊妹.
今天我們更加要求上帝幫助我們.
在神話當中.
讓我們認清我們在傳什麼道.
我們是一個怎樣的生命.
正視自己生命的情況而去改變.
到最後的時候.
傳道是有目標的.
傳道是什麼呢.
我們要活出屬於神的生命.
為著叫未信的人去得著永生.
這個是傳道的目標.
最終你要讓人能夠得到永生.
我上次的圖.
我今天忘記給Yoyo.
那幅圖我經常都保存在電話裡.
那幅圖就是前面是一個教堂的入口.
但後面是一個懸崖掉下去.
經常提醒自己.
我們今天在教會帶領人回來的時候.
你究竟是要帶他們上天堂.
還是帶他們去地獄.
為什麼我這樣說呢.
不是啊 我們回來教會當然是上天堂.
不一定的 弟兄姊妹.
我經常都說.
當你和人說.
你記住我用這個字.
你不要哄人入局.
你回來教會吧.
教會裡面很有愛.

$^{601}$第一樣就是首先挑戰自己.
你回來教會.
你是不是自己內心真的覺得教會很有愛.
第二樣.
回來教會.
你是否很有把握這個人去感受到愛.
所以剛才我才說.
如果我們自己在我們教會.
我們信耶穌的裡面.
我們自己都有些問題.
我們都還沒有處理好的時候.
我們怎樣去將這個福音.
介紹給我們身邊的人呢.
當你帶了他回來之後.
你是等於推他去地獄.
為什麼呢.
哦 原來信耶穌是這樣的.
所以我告訴你.
他就更加有把握不回教會.
不信耶穌.
但如果你能夠帶他回來教會的時候.
哇 原來基督徒是這樣的.
那就不同了.
他真的在這裡可以找到.
聖經裡面所寫的信 望 和 愛.
所以弟兄姊妹.
求上帝幫助我們.
我們很清楚帶這些人去得著永生.
因為保羅說.
我們是憑著對上帝的信心.
第一樣東西.
他是按照著舊約裡面的先知.
他們所寫的時候去相信.
他怎樣信 我都怎樣信.
他很相信以前的上帝是怎樣.
今天同樣對我們的上帝都是怎樣.
弟兄姊妹.
同樣我想告訴你.
這個都是我們今天讀聖經的價值.
這本聖經已經超過了2000年了.

$^{641}$如果你只是當是一本2000年前的作品.
接著的時間我們繼續去看.
是沒有意義的.
但是如果你很相信.
2000年前的上帝.
也都等於是今天的上帝.
2000年前他會做的事情.
他今天同樣都會做.
他的能力從來沒有改變.
弟兄姊妹.
你的生命就不一樣了.
所以主耶穌.
他才跟我們門徒說了一句話.
你的信心只需要像一顆芥菜種.
你就能夠可以移散.
就正因為他很清楚這個上帝.
就是以前的上帝.
聖經裡面的說話就是讓你看到.
他在以色列人的生命當中.
做了什麼事.
施行了什麼大事和奇蹟.
所以如果你問我的話.
我可以告訴弟兄姊妹.
這個上帝從來沒有改變過.
你在聖經裡面所見到的神蹟.
他今天同樣都可以發生.
這個是他的主權.
所以如果你問我.
我信不信神蹟.
我信神蹟.
我在短宣裡面我經歷過神蹟.
但是你問我香港有沒有見過那些神蹟.
我沒有.
我沒有親身經歷過那些神蹟.
那些那麼誇張的神蹟.
但是當我要很理性去比較的時候.
唯一的分別就是什麼呢.
香港什麼都有.
短宣那個農村的地方是家徒四壁.
這個神蹟是怎樣的呢.

$^{681}$我記得我去短宣.
然後我們都在城市.
然後有一個姐妹.
她說鄧全道.
你可不可以和我們去農村裡面.
找一個老人家.
我說可以啊 為什麼啊.
他說他發燒.
我說發燒啊 看醫生沒有.
他沒有錢看醫生.
我們其他的人也籌不到錢去看醫生.
你可不可以去為他祈禱.
那個圖畫就真的好像聖經裡面所講.
耶穌去到彼得的岳母家.
那個情況一樣 躺在那裡.
我們去到.
我見到那個景象.
那個真的什麼叫家徒四壁.
兩張木椅搭著木板.
那個人就睡在上面.
其他什麼都沒有.
全部石具在地上 沙地.
你可以想像到沒有任何的東西.
我去到之後.
我心裡面有一份很大的感動.
我和神說.
我神說他真的什麼都沒有.
神說他真的什麼都沒有.
最後你可不可以憐憫他.
然後就為他祈禱.
最後如果你的醫治大能真的可以的話.
你給他.
我是真的見證著這個老人家.
他立刻退燒.
他真的坐起來.
他真的好像聖經裡面所講的一樣.
彼得的岳母.
他真的立刻起來的時候.
他什麼都沒有.
家裡剩下的可能是地上的一兩隻雞蛋.

$^{721}$你知道他們很客氣那些很寶貴的東西.
他立刻起來的時候就塞給你.
他說我什麼都沒有.
你當這些是醫藥費.
他說醫藥費.
我說不是的 不關我的事.
我說這個是上帝的.
但他就立刻很精神地還給你.
弟兄姊妹有沒有神蹟.
有 有.
但是弟兄姊妹.
當你認為自己仍然什麼都有的時候.
上帝怎樣在你身上去展現這種大能呢.
原來上帝他從來都沒有改變過.
所以當保羅他講完這些的時候.
他教我們一種很重要的視野和心態.
他說身體衰殘.
但是沒有辦法可以阻擋我們的內心.
身體是會一天一天地衰殘.
但是內心可以不斷更新.
弟兄姊妹.
我想大家都見過.
有一些病人.
有一些老人家.
他們雖然身體有很多病很多事情.
但是他們的內心是很有那種火熱.
我想在這裡特別提名的一位.
我真的很感恩.
其實我看著曾太太我是很感恩的.
他面對的不容易 大家都知道.
但是他的內心我不知道大家有沒有見到.
他內心裡面那種平安那種動力.
那種我要在神面前尚有一息.
我都要繼續去事奉上帝.
我很見證著我記得有一段時間他走得不是很好.
但是之後他祈禱.
後來我就見到他.
他仍然可以和大家一起列隊出來.
一起去敬拜.
弟兄姊妹在你們身邊很多的見證人.

$^{761}$他告訴大家弟兄姊妹聽聖經所說的話是真實的.
只要你真的按著聖經所說的信.
你就能夠經歷到聖經所說的保羅所說的大能.
我們的心可以不斷更新.
所以我很感恩.
我現在每次回來都不斷去慰問我的腳.
現在的情況怎樣.
但是我真的很想告訴弟兄姊妹聽.
我沒有把焦點放在腳上.
腳算什麼呢.
你是可以繼續的.
上帝沒有封你的嘴.
所以我每次講道之前.
我都說每次都是上帝告訴我.
每次都是靠我.
每次cut的時間這半個小時靠我.
每次都是靠祈禱去度過的.
每次都是靠祈禱去經歷上帝原來你可以這麼厲害.
然後我就把這種厲害告訴其他人.
上帝就是這麼厲害.
所以弟兄姊妹.
面對目前的苦難.
保羅告訴你是很短暫輕微.
相反你要得到的是永恆的價值.
你今天所面對的困難其實很快就過了.
但是你在神的裡面.
你沒有因為這個緣故而在祂當中得勝.
方法就是改變你的焦點.
不要把這個困難放在你的眼中.
而是把永恆的價值放在你的眼中.
最後我用一個結束.
有些弟兄姊妹這段時間都多人移民.
有些弟兄姊妹移民去英國.
這個不是我們教的弟兄姊妹.
是我以前曾經有一次做輔導去認識.
我去了英國.
後來他早兩天找我.
問我有沒有空.
我那時候應該是晚上十點鐘.
他說我現在英國萬城.

$^{801}$他跟我分享.
他分享了一件事.
當然詳細我不說.
但他真的很困擾他.
他覺得他不知道如何處理.
但是弟兄姊妹.
上帝就是這麼奇妙.
我預備說到.
最後就是說.
你不要把你目前的困難.
把焦點放在你的困難當中.
你要放在上帝永恆的價值裡面.
這些東西很短暫.
之後我就跟他分享.
我說其實你現在面對的困難.
其實很小.
跟他說完之後.
我記得我跟他大概通話.
應該是二十分鐘之內.
跟著他說了一句話.
Aka 你很厲害.
二十分鐘我就已經沒事了.
之後我就跟他說.
不是我厲害 是上帝厲害.
我剛才跟你說的那句話.
是聖經的說話.
不是我教你.
是聖經教你.
弟兄姊妹.
上帝的道是活潑的.
是滿有能力的.
弟兄姊妹.
放心去信.
如果我們作為基督徒.
我們都不拿著我們這把寶劍.
你如何去面對你人生裡面其他的問題呢.
你如何去成為一個傳道者.
將道傳給其他人呢.
所以求神讓我們成為一個.
真實相信上帝的傳道者.

$^{841}$用你的生命我們一起.
去將福音傳給其他人.
我們同心祈禱.
撒下天賦我們求你親戚的幫助.
主啊 你的道是又真又活又信實.
當剛才師班去唱詩.
他們獻詩去唱到.
主啊 當有主去幫助.
有你的同在的時候.
我們還要去懼怕誰呢.
還有誰可以去抵擋我們呢.
主啊 孩子滿心裡面就很感動.
我真的很心願.
我們每一個弟兄姊妹.
他們真的有這一份這樣的信.
很確信.
以至於我們有一份這麼大的影響力.
去能夠將這個祝福帶給我們身邊的人.
我們求你幫助我們.
叫我們不要小看你.
而更加我們知道.
主啊 所有一切能力都是從你而來.
我們求你叫我們將我們所有一切.
奉獻在上帝你的腳前.
求你使用.
求你繼續幫助使用我們.
當我們面對困難的時候.
我們和保羅所想的一樣.
所有的困難是又短暫又輕.
但我們可以繼續走上帝你的道路.
主你必定有你自己的能力.
幫助帶領我們去走過這些路.
將一切所有我們交託.
其他奉教主耶穌得勝 明治其空.
\newpage



\section{以賽亞書 58:3-7}
\label{sec:cMmv_XzlkAc}
\textbf{2023-08-20 早堂《禁食的意義》以賽亞書 58\_3-7 講員:陳韋安牧師}
\newline
\newline
連結: \href{https://youtube.com/watch?v=cMmv_XzlkAc}{\texttt{https://youtube.com/watch?v=cMmv\_XzlkAc}} ~~~~ 語音日期: 2023-08-22
\newline
\newline
\hyperref[sec:FIbxv36OMMk]{\small{< < < PREV SERMON < < <}}
~
\hyperref[sec:index_chronic]{\small{[返順時目]}}
~
\hyperref[sec:index_scriptual]{\small{[返順卷目]}}
~
\hyperref[sec:TEOKFaxZL4g]{\small{> > > NEXT SERMON > > >}}
\newline
\newline
以賽亞書 58:3-7
\newline
\begin{longtable}{cl}
\hline
\hline
章節 & 經文 (和合本修訂版)\\
\hline
58:3 & \begin{tabularx}{0.7\textwidth}{X} 「我們禁食，你為何不看呢？我們刻苦己心，你為何不理會呢？」看哪，你們禁食的時候仍追求私利，剝削為你們做苦工的人。 \end{tabularx} \\ \\ \relax
58:4 & \begin{tabularx}{0.7\textwidth}{X} 看哪，你們禁食，卻起紛爭興訟，以兇惡的拳頭打人。你們今日這種禁食無法使你們的聲音聽聞於高處。 \end{tabularx} \\ \\ \relax
58:5 & \begin{tabularx}{0.7\textwidth}{X} 這豈是我所要的禁食，為人所用以刻苦己心的日子嗎？我難道只是叫人如蘆葦般低頭，鋪上麻布和灰燼嗎？你能稱此為禁食，為耶和華所悅納的日子嗎？ \end{tabularx} \\ \\ \relax
58:6 & \begin{tabularx}{0.7\textwidth}{X} 我所要的禁食，豈不是要你鬆開兇惡的繩，解開軛上的索，使被欺壓的得自由，折斷一切的軛嗎？ \end{tabularx} \\ \\ \relax
58:7 & \begin{tabularx}{0.7\textwidth}{X} 豈不是要你把食物分給飢餓的人，將流浪的窮人接到家中，見赤身的給他衣服遮體，而不隱藏自己避開你的骨肉嗎？ \end{tabularx} \\ \\
[1ex]
\hline
\hline
\end{longtable}
$^{1}$(陳志全) 各位節目早安.
今天我們會和大家分享禁食的訊息.
剛才我們也有說過.
但今天也不是傳統的《識經港道》.
今天也會比較宏觀地和大家分享禁食的事件.
特別是這個暑假.
可能大家在這個暑假也會去旅行.
暑假最困難的不是花錢.
而是回來要減肥.
很多時候,特別是香港人.
其實也是很豐富.
我們對於吃也是不困難.
反而減肥是困難的.
可能對大家來說.
當然禁食不等同於減肥.
但也是有一點關係的兩者.
所以今天我們就和大家說一下.
有關禁食這件事情.
我自己也有肥過.
肥過又瘦過.
如果你認識我很久的話.
你也會覺得相對我什麼年紀的我.
如果20年前就肥了.
1年前就瘦了.
其實我想大家也知道.
我試過最肥的時候.
也不是很誇張.
也是COVID那段時間.
那時候又戴著口罩.
大家不太看到我的樣子.
就不斷地肥下去.
一來是因為運動.
二來是工作力大.
可能大家也知道.
當你工作很辛苦的時候.
你會變胖.
這是一種工傷.
你工作很辛苦的時候.
你也會嘗試多吃一點東西.
回來感受一下.

$^{41}$然後又很晚了.
又吃宵夜.
最肥的時候.
其實也試過一段時間.
有一次我在港島.
網上直播.
打開口罩.
大家就看到我身邊的目了.
第二天早上我去屯門.
教會港島.
有個學生來吃午餐.
說昨天在那邊.
昨天看到你的港島.
他突然很牧者心腸地跟我說.
其實這樣也挺好看.
現在就覺得不行.
要減肥了.
不能再這樣肥下去.
開始在網上找一些減肥的APP.
發現到一種東西.
原來這東西叫斷食.
大家有沒有聽過.
間歇性斷食.
就是你在一天裡面.
一段時間不吃東西.
原來也挺有效的.
就是這樣.
會是168,186那些.
就是16小時不吃東西.
然後有個食物窗戶.
剩下8小時.
那時候我就採取一個更激進一點的.
就是24小時.
就是20小時不吃東西.
4小時才吃東西.
其實也挺有效.
那時候也減了十多磅.
很快的.
後來發現原來.
我們所說的.

$^{81}$香港人很流行的那些斷食.
其實斷食和我們以前所說的.
這個就是160減到152磅.
我們可以稱呼這些東西.
你可以叫做斷食.
叫做減肥.
或者叫做不吃東西.
如果在一個信仰來講.
也有類似的.
我們又叫做守齋.
如果是天主教或者伊斯蘭教.
叫做守齋.
我們以前就叫做.
或者叫齋戒.
我們以前就叫做禁食.
當然兩樣東西是不同的東西.
不過其實那個的本質.
或者所做的都是差不多.
就是不吃東西.
不知道大家有沒有試過.
我們教會其實比較少提倡斷食.
通常都是一種的教會活動.
一年就做一次活動.
大家一起參加教會活動.
報名之後大家就斷食.
但這個似乎不是本身最重要的原意.
因為禁食和斷食.
重點不是一年一次的象徵式參加活動.
而是成為一個你很重要的生活習慣.
或者是信仰的方式.
所以一年一次的活動是好的.
因為可以鼓勵大家開始純粹地做.
但這不是純粹一年一次參加的算數.
而是成為一個大家在平時生活裡面.
都應該可以定期去實踐的信仰回應.
這是第一點.
禁食不是純粹像我們以前去營地一樣.
吃一次過大家很好玩的手軟就算了.
而是一個應該是成為大家在信裡面.
定期生活的一部分.

$^{121}$一個習慣要做的事情.
第二點就是禁食不是完全等同於一種禱告方法.
禁食不是純粹一種更加靈的禱告方法.
不是因為想要更加高命中率.
就要特意禁食來提高命中率.
禁食歸禁食 禱告歸禱告.
雖然很多人會黏在一起做.
但禁食不等同於禱告方法.
禁食即是不吃東西.
就是完全將食物一段時間禁止自己.
所以我們發現禁食不是一種純粹為了操控上帝的禱告結果的方法.
雖然有關係.
我們說禁食是出現了一個事件之後.
一件很特別的事情出現了.
你就開始用禁食不吃東西來回應.
或者是一個很自然的方法.
你不想吃東西了.
然後你還是會禱告.
然後出現一個結果出來.
三件事是有關係的.
但又不是完全是一個因果關係.
不是因為你禁食.
成就令你神就聽你的祈禱.
而是一件很重要的事出現了之後.
你就開始禁食並禱告.
所以三件事是有關聯.
但不是完全是因果關係的事情.
我們開始說說禁食在聖經裡面的基本基礎.
在舊約裡面.
我們發現聖經裡面出現了很多次禁食這個字.
甚至乎比靈修更多的字.
有出現了五十多次禁食這個字.
在舊約裡面你會發現.
如果將這個字分開三種不同情況的話.
舊約裡面有三種不同的情況.
是有禁食的出現的.
第一個是什麼呢.
第一個就是有關預備迎見上帝的情況.
所謂禁食其實是一種叫做克苦自我的狀態.
一個是克苦己心的狀態.

$^{161}$是讓自己經驗一種難受的.
因為禁食是不自然的.
但一種這樣的狀態.
就是這樣.
讓自己處於一種克苦己心的狀態.
為什麼這樣做呢.
第一個情況就是為了預見上帝的預備.
特別是在摩西在西太山迎見上帝之前.
當他要預備自己領受從神而來的十誡之前.
摩西就禁食四十晝夜.
來預備自己.
這個第一種情況.
就是禁食是讓我們更加準備自己.
禁食不是本身的目標.
而是一個中間的階段.
讓我們能夠透過禁食.
在身體靈性上好好地預備迎見上帝.
所以這是第一個情況出現禁食的描述.
就是一個人要去預見上帝之前.
必須要做的事情就是禁食.
第二個情況就是認罪.
每當一個人發現自己有罪的時候.
發現自己有很大罪的時候.
他都會非常痛改前非.
從而去禁食.
譬如你看到當薩姆爾記裡面記載.
薩姆爾帶領一群猶太人去禁食認罪.
立志不再拜偶像.
或者在約拿書裡面也一樣.
當約拿宣告利尼比城的審判的時候.
整個城市的人都一起去禁食.
徹底去認罪.
所以禁食往往都是和自我的反省.
發現自己有罪有關連.
因為一個非常痛切.
一個非常哀傷的情況.
承認自己有罪.
第三就是一些國難的時候.
去到一個非常大的民族危機的時候.
以色列人都會一起去禁食.

$^{201}$當一個人或者整個民族.
遇見極大的災難.
感覺到非常無助的時候.
他們都會去禁食.
這個不是一種手段.
而是一種很自然的流露.
因為太過哀傷.
太過需要去尋求上帝.
他們就會去禁食.
去這樣的表現.
所以禁食在舊約裡面有三種不同的情況.
一個就是預備遇見上帝.
一個就是認罪.
第三就是遇見一個非常大的民族災難的時候.
這個就是舊約所說的描述.
當然到了兩者之間的時候.
猶太人慢慢將會變成一種很律法式的規矩.
一種表達傾聽權的方法.
但耶穌是批判的.
耶穌是不贊成一種.
純粹為一種敬詞表現的方式.
耶穌自己也是一個禁食者.
他自己在一個曠野裡面.
都是禁食.
來去四日周夜.
禁食和禱告.
兩者是關聯.
但不一定是完全一樣的事情.
另外保羅也是一樣.
保羅當他遇見耶穌之後.
都是在大馬士革路上.
有三天的時間.
來斷食.
最後保羅也提過.
在光榮多後書裡面提過.
在臨後第六章.
他說反倒在各有的事上.
表明自己是上帝的用人.
就如在許多的忍耐患難.
貧乏困苦.

$^{241}$鞭打監禁擾亂勤勞.
警醒不食.
這個也是保羅曾經所做過的表現.
就是不吃東西.
所以整個教會的歷史裡面都是這樣.
很多的教父都是保持著這樣的習慣.
每個禮拜裡面.
每逢禮拜一,三.
都會有禁食的時間.
成為了一個很重要的.
一種的宿命的傳統.
但你會發覺.
我們環教會比較少.
或者作為一種活動參加算數.
我們好像做了就做了.
但好像體會不到那種意義在裡面.
但我們發覺原來.
當你真的去定期不吃東西的時候.
一個健康地不吃東西.
就是一直放在時間裡面.
這種飢餓的狀態.
往往是令我們有些很特別的宿命的意義.
今天就跟大家說說.
當你處於一種禁食狀態的時候.
能夠可以會遇到的一些宿命的意義.
第一個就是一種我稱之為減的靈性.
斷食其實是一種實踐一種減的靈性.
很具體的.
當你發現你不吃東西的時候.
當你身體處於一種很缺乏的狀態的時候.
你就很體現到一種減的靈性.
什麼叫減的靈性呢.
其實減是一種好事.
我們很多時候都覺得減是一種不好的事情.
減身減人減放飯時間.
都是一些被人剝奪了的事情.
原本有的被人減少了.
你就會埋怨會不忿覺得不公平.
所以我們都很傾向反方向而行.
我們就會去爭奪.

$^{281}$我們就會去backup.
我們會團職居其.
明明覺得怕自己不夠.
明明已經夠了.
都會留一些用來做後備的.
儲備backup.
無論你的hard disk.
無論你的錢.
或者你的脂肪都一樣.
所謂脂肪是什麼.
就是一些你明明不需要用.
但儲蓄來不時之需的東西.
所以我們傾向不是減的.
我們傾向是會多的.
多過我們所需要的.
怕自己將來不夠用.
這些就會成為我們的脂肪.
在屬靈裡面也有屬靈的脂肪.
我們有一些明明不需要的東西.
我們都覺得要先留著.
怕將來沒有.
這些就是我們生命裡面的脂肪.
太多反而會令我們.
遭遇太多不好的地方.
如果你跟我差不多年紀.
我們小時候.
我們小朋友玩遊戲機.
我們有一種東西.
以前我們不是玩正版的.
我們小時候是玩翻版遊戲.
我們小時候有一種東西.
有一個名字叫超音博士.
就是有一個東西.
會令你玩到三千多個免費遊戲.
所以這是一個很震撼的事情.
一個小朋友.
突然有三千多個遊戲.
是一件很誇張的事情.
但你無法應付.
一個小朋友玩不到三千多個遊戲.

$^{321}$你只會浪費了它.
每隻遊戲可能玩幾下.
你就不玩了.
因為太多了.
多到你不知道有什麼是值得玩的.
所以都是這樣.
我們生命裡面.
往往都會遇見一些不是少.
而是太多.
多到我們沒辦法專注在一些重要的事情上.
所以減去是重要的.
減去是讓我們能夠真正發現.
你生命裡面一些很重要的東西.
一些值得你將它放在高位.
放在重要頭三位的東西裡面.
當你的生命清單裡面太多東西的時候.
多到令你失望.
沒有焦點.
令你不知道哪些是你生命裡面最重要的東西.
所以這時候減.
正正是讓你能夠可以慢慢發現.
你生命裡面最重要的東西.
就是減剩三個.
減剩一個.
你就知道原來你的上帝.
你某些重要的人.
對你來說是最重要的.
所以我們太過來到.
生活裡面太過恐慌.
很怕自己的生命裡面缺少.
無論是你的金錢.
或者你的物質.
或者你在屬靈上面也是一樣.
不過知道減是分兩款.
一個叫做主動的減.
一個叫被動的減.
被動的減我們稱之為剝削.
即是強奪或者侵略.
但是什麼叫做主動的減.
主動的減就是我們自己願意的犧牲.

$^{361}$我們的捨己.
我們的愛.
所以當我們生命裡面主動減少東西的時候.
這是一個很重要的屬靈的操練.
耶穌也說.
若有人要跟從我.
就當捨己天天背十字架跟從我.
這種所謂的捨己.
就是拒絕自己.
去學習去對付減少自己生命裡面有的東西.
這個正是很重要的地方.
就是禁食.
正是一種很具體的體會.
透過你不吃東西.
這種缺乏.
來經驗到這種上帝減的靈性.
禁食是一種你自願的東西.
它不是一種上帝命令你一定要做的東西.
聖經沒有寫一個十禮拜禮佛要禁食多少次.
但是當你願意自己主動減的時候.
你就來體驗這種跟上帝很特別的關係.
不是一種律法主義的要求.
而是我們自己生命裡面自願的寫下.
第二個.
就是當你願意來到禁食的時間裡面.
你就可以去經驗到耶穌基督的真實和同在.
我經常說less is more.
少更加多.
但是多在哪裡呢.
這個多正是我們能夠經驗到這位上帝.
上帝的豐盛不是靠我們的生命累積回來的.
恩典是不能儲存回來的.
恩典只能夠在你的缺乏裡面去經驗到.
所以當你來到體驗到自己的缺乏的時候.
這種less正正是我們能夠發現到基督耶穌的豐盛.
真人這個你不是沒辦法知道的.
當然豐富豐富豐盛.
但是當你來到減少的時候.
你就能夠看到感受到基督耶穌在缺乏的時候.
給我們那種力量.

$^{401}$聖經裡面有很多經文.
一粒麥子不落在地裡死了.
仍舊是一粒.
若是死了就結出許多一粒子來.
凡要救自己生命的必喪掉生命.
凡為我和福音犧牲的必救了生命.
我們聖經裡面有一個很出名的人物.
就是聖凡西氏.
是一個十三世紀裡面的托缽修士.
他正正就是願意來實踐聖經的誡命.
就是寫了一切去跟從耶穌.
當他願意寫下一切去跟從耶穌的時候.
整個歐洲都被震撼.
幾千人跟從他一切成為一個這樣的修士.
這種自願的捨棄稱之為聖貧.
神聖貧窮.
這種聖貧是一個很有力量的經驗.
當我們自願寫起自己擁有的東西的時候.
經驗到的往往是多於我們所擁有的.
所以這是我們經驗耶穌的其中一個方法.
當你來到真日的時候.
其實是說恩典救我肉.
什麼叫做足夠.
足夠的意思是說是一種剛剛好的狀態.
我們唯有透過減法來不斷地倒水.
我經常說倒水也是一樣.
倒水怎麼倒都不會有一個準確的份量.
你只能夠裝滿它之後.
倒到剩下這麼多的地方.
這個程度之下就是我們能夠足夠的程度.
所以我們生命裡面也是一樣.
這種禁食裡面的飢餓感.
往往是讓我們發現我們足夠的水平在哪裡.
透過一種不斷的減.
來尋找我們真正經驗上帝耶穌基督恩典足夠的程度.
這個就是我們能夠專注上帝的很重要時刻.
當你處於一種飢餓狀態.
處於一種要由上帝恩典來使自己足夠的狀態的時候.
這種狀態是非常能夠讓你專注上帝.
這個真的不試過你不會知道.

$^{441}$當你處於一種飢餓狀態的時候.
你這個人其實是會越來越專注.
專注在一些你平時沒有發現的東西裡面.
這個時候你去祈禱.
是能夠更加的穿心.
不是特別靈.
而是能夠更加的專注禱告.
當你試一下一六八.
譬如說你七點鐘晚上吃完飯就不吃東西了.
到第二天中午吃東西的時候.
當你痛過的時候.
這個時候當你習慣了的時候.
你會發覺你人是更加精神的.
更加專心.
你能夠更加去發現耶穌基督的恩典在生命裡面.
這個就是我們生命裡面的狀態.
我們能夠很經驗上帝的兩個狀態.
一個就是你在跑步的時候.
我喜歡跑步領修.
因為當你不斷來到心跳流汗的時候.
你發現自己是一個很基本的受造物.
剩下心跳和流汗的一個受造物.
今時今日一樣.
當你不吃東西的時候.
你就會成為一個像塵土一樣的狀態.
你就不是一個什麼人.
只不過是一個很普通的需要食物的受造物.
這個時候你去禱告.
你會更加能夠去發現自己的本相.
所以當我們來到去實踐這種狀態的時候.
你就能夠發現我們生命裡面的上帝.
第三.
我想說的就是剛才的經文.
禁食和社會公義.
當然禁食是我們很個人的屬靈的操練.
但是我想告訴你.
禁食不單單是一種我們靈性自肥的方法.
不是純粹一種增加靈性.
或者增加和神關係的方法.
這是對的.

$^{481}$但是禁食不單單是關於自己個人靈性操練的事情.
經文問給我一次.
剛才問經文.
就是以賽亞書的經文.
好 預備.
1 2 3.
他們說:我們禁食,你為何看不見呢?.
我們刻苦記心,你為何不理會呢?.
看那!你們禁食的日子仍求利益.
勒逼人為你們做苦工.
你們禁食卻互相尊敬.
以凶惡的拳頭打人.
你們今日禁食.
不得使你們的聲音聽聞於上.
我所揀選的禁食.
不是要鬆開凶惡的繩.
解下壓上的鎖.
使被欺壓的得自由.
截斷一切的壓嗎?.
不是要把你的餅.
分給飢餓的人.
將漂流的窮人接到你家中.
見赤身的給他衣服遮體.
顧失自己的骨肉而不掩藏嗎?.
原來….
先回到座位.
原來耳塞書中提到.
禁食是一個很重要,很新的角度.
它正正是要重新去定義.
猶太人禁食是一個很重要的範疇.
禁食不單單是一種自己個人.
即是靈性操練的東西.
這是一個效果,肯定有的.
但這不是純粹我們自己個人.
操練一種自肥的方法.
當禁食失去我們對社會的貧窮人.
或其他人的關注時.
禁食只會是我們留於自己自私地.
與神關係好的狀態.
耳塞書說.

$^{521}$當我們禁食而忽略了.
社會上其他人的貧窮和飢餓時.
禁食就失去意義.
所以我們發現.
原來禁食是一個很重要.
讓我們方言世界和社會.
禁食可以稱為一種.
我們與這些人同行的一個很重要的方法.
稱之為Solidarity.
中文翻譯為休戚相關.
是一個很特別的字.
我們與他們站在同一邊.
你未必能夠有那麼多時間去服侍他們.
或與他們住在同一個地方.
畢竟錦繡花園和深水Po 是很遠的距離.
但當我們的心在那裡的時候.
我們就可以藉著禁食.
與他們一起去經歷這種苦難.
這是聖經所說的.
我們的禁食其實是一個很重要的向導.
就是讓他們在生命中實現與這些人同行的狀態.
我們未必能夠扭轉社會不平等的狀況.
貧富懸殊或許有很多不公平的現象.
你也無法透過自己改變整個世界的氣候.
但當我們的心在那裡的時候.
如果我們願意實踐耶穌基督的教導.
與這些人同行的時候.
我們可以透過禁食與他們同行.
正正這也是耶穌所做的事.
所以當我們有飢餓時.
有人口渴時.
耶穌就在那個地方.
所以當我們實行禁食的時候.
這種禁食是一種我們與這些人.
面對著世界的一個很具體的方法.
這未必是有用的問題.
而是我們願意與他們站在同一邊的問題.
當我們實踐這種禁食.
我們正正是與他們處於同樣的飢餓狀態.
同樣的缺乏狀態.

$^{561}$我們能夠經歷耶穌基督在當中.
在貧富中的豐盛狀態.
所以聖經裡面說.
禁食更加是一種探望的源頭.
耶穌說過.
新郎與陪伴之人同在的時候.
陪伴的人豈能哀動呢.
但日子將到新郎要離開他們.
那時候他們就要禁食你.
當耶穌在地上的時候.
耶穌說你不需要禁食.
因為我在這裡.
當耶穌升天之後.
當耶穌還沒回來之前.
在這段時間裡.
在大地幽暗的日子裡.
當社會仍然出現很多不公平.
和很多飢餓貧窮的時候.
我們就需要禁食.
來度過這段時間.
來與這些人同行.
來體會到他們的困苦.
離開我們一些的安恕.
一個禮拜一個月一兩次的禁食.
這是我們很重要.
實施耶穌基督的教導.
然後他說.
這是一種我們面對黑暗時候的盼望.
聖經說.
四月五月禁食的日子.
七月十月禁食的日子.
必變為由大家歡喜快樂的日子.
與和歡樂的節期.
所以你們要喜愛誠實和和平.
當我們禁食的時候.
我們是與這個世界期盼.
耶穌基督最終的現職.
當耶穌基督再回來的時候.
沒有人會再飢餓.
沒有人再需要禁食.

$^{601}$有一個永恆的現職.
一起去同享受基督耶穌的歡樂.
但是之前.
我們需要禁食懷著盼望.
見到這個世界的困苦.
我們願意一起來禁食.
懷著盼望去同行.
所以今天很鼓勵大家一起來.
實踐一種很具體的聖經教導.
不需要等教會來辦禁食的日子.
不需要報名.
你只需要定期地.
很簡單的.
可能一個星期一次.
一個月一次.
自己選日子出來.
定期地做一個重點.
不是一個活動式參與.
而是定期地不吃東西.
你叫做斷食也好.
不吃東西也好.
減肥也好.
這個順其自然可以減肥.
又可以讓你靈性上體驗到這麼多東西.
減的靈性.
經驗基督耶穌的豐富.
和與社會上的人同行.
在禁食的時間裡.
當你感覺到飢餓的時候.
這是正常的.
當然是應該的.
嘗試一個星期一個月.
經驗這樣的感受.
很少有的.
很少很餓的.
但當你自願地做的時候.
這成為我們這樣的屬靈操練.
懷著盼望.
懷著一種追求上帝的心.
實施出來.

$^{641}$我相信這是一個很好的經驗.
當你真的成為一種習慣的時候.
我猜你心靈裡會更加豐盛.
因為less is more.
愛是祈禱.
因為你自己也在生命裡做我們的榜樣.
你經常在聖經裡被記載.
你去禁食.
你去禱告.
更加去教導我們.
禁食不單單是一種自我屬靈的操練.
或者是進步.
更加是我們和社會.
在優雅的物質裡.
和一群有需要的人士來同行.
求主幫助我們.
能夠在禁食裡.
可以嘗試一個很簡單的操練.
我們從168開始.
從簡單的一些program開始.
來實踐一個這樣的習慣.
又能夠在當中.
去實踐一種這麼重要的屬靈教導.
從舊約到今天一直流傳的教導.
很多人叫做不同的名稱.
有些人叫齋戒.
有些人叫守齋戒旗.
有些人禁食.
但求主幫助我們.
又能夠有這樣的恆心.
來嘗試實踐禁食的計劃.
又能夠在當中.
去見到更多的豐富.
奉主並求.
\newpage



\section{路加福音 10:25-37}
\label{sec:TEOKFaxZL4g}
\textbf{2023-08-20 加恩樓《愛鄰如己．照樣做吧》路加福音 10\_25-37 講員:陳麗蟬牧師}
\newline
\newline
連結: \href{https://youtube.com/watch?v=TEOKFaxZL4g}{\texttt{https://youtube.com/watch?v=TEOKFaxZL4g}} ~~~~ 語音日期: 2023-08-23
\newline
\newline
\hyperref[sec:cMmv_XzlkAc]{\small{< < < PREV SERMON < < <}}
~
\hyperref[sec:index_chronic]{\small{[返順時目]}}
~
\hyperref[sec:index_scriptual]{\small{[返順卷目]}}
~
\hyperref[sec:zrsNFCgx1bQ]{\small{> > > NEXT SERMON > > >}}
\newline
\newline
路加福音 10:25-37
\newline
\begin{longtable}{cl}
\hline
\hline
章節 & 經文 (和合本修訂版)\\
\hline
10:25 & \begin{tabularx}{0.7\textwidth}{X} 有一個律法師起來試探耶穌，說：「老師！我該做甚麼才可以承受永生？」 \end{tabularx} \\ \\ \relax
10:26 & \begin{tabularx}{0.7\textwidth}{X} 耶穌對他說：「律法上寫的是甚麼？你是怎樣念的呢？」 \end{tabularx} \\ \\ \relax
10:27 & \begin{tabularx}{0.7\textwidth}{X} 他回答說：「你要盡心、盡性、盡力、盡意愛主—你的神，又要愛鄰如己。」 \end{tabularx} \\ \\ \relax
10:28 & \begin{tabularx}{0.7\textwidth}{X} 耶穌對他說：「你回答得正確，你這樣做就會得永生。」 \end{tabularx} \\ \\ \relax
10:29 & \begin{tabularx}{0.7\textwidth}{X} 那人要證明自己有理，就對耶穌說：「誰是我的鄰舍呢？」 \end{tabularx} \\ \\ \relax
10:30 & \begin{tabularx}{0.7\textwidth}{X} 耶穌回答：「有一個人從耶路撒冷下耶利哥去，落在強盜手中。他們剝去他的衣裳，把他打個半死，丟下他走了。 \end{tabularx} \\ \\ \relax
10:31 & \begin{tabularx}{0.7\textwidth}{X} 偶然有一個祭司從那條路下來，看見他就從另一邊過去了。 \end{tabularx} \\ \\ \relax
10:32 & \begin{tabularx}{0.7\textwidth}{X} 又有一個利未人來到那裡，看見他，也照樣從另一邊過去了。 \end{tabularx} \\ \\ \relax
10:33 & \begin{tabularx}{0.7\textwidth}{X} 可是，有一個撒瑪利亞人路過那裡，看見他就動了慈心， \end{tabularx} \\ \\ \relax
10:34 & \begin{tabularx}{0.7\textwidth}{X} 上前用油和酒倒在他的傷處，包裹好了，扶他騎上自己的牲口，帶他到旅店裡去，照應他。 \end{tabularx} \\ \\ \relax
10:35 & \begin{tabularx}{0.7\textwidth}{X} 第二天，他拿出兩個銀幣來，交給店主，說：『請你照應他，額外的費用，我回來時會還你。』 \end{tabularx} \\ \\ \relax
10:36 & \begin{tabularx}{0.7\textwidth}{X} 你想，這三個人哪一個是落在強盜手中那人的鄰舍呢？」 \end{tabularx} \\ \\ \relax
10:37 & \begin{tabularx}{0.7\textwidth}{X} 他說：「是憐憫他的。」耶穌對他說：「你去，照樣做吧！」 \end{tabularx} \\ \\
[1ex]
\hline
\hline
\end{longtable}
$^{1}$靜姐姐早安.
今天我又有機會來到嘉欣樓.
跟大家一起敬拜讚美.
一起用上帝的說話.
來彼此勉勵.
有一個人笑.
笑這個猶太教裡面.
很多的律法.
有人數著數著.
在聖經裡面的律法.
有多少條呢.
總共有613條.
於是他就說.
這麼多條怎麼記得.
怎麼做得完.
他就去找一個拉比.
拉比.
我現在就試一下.
用一隻腳站在那裡.
看你能不能夠.
在我用單腳站的時候.
你將你所有猶太人的律法.
全部數出來.
如果是這樣的話.
我就會加入你這個猶太教.
做猶太教徒.
這個拉比說.
行 你單腳站了.
好.
拉比就說.
一句可以說完.
愛你的論社.
好像愛自己.
其實律法裡面.
最主要的中心就是這句.
其他都是一些註解.
所以好好地學習.
剛才成熙說.
我們讀出那段經文.
是說到有一個律法師.

$^{41}$其實他心裡面.
有些不是很好的意念.
他就去問耶穌.
怎樣才能夠.
拿到這個永生.
但是這個律法師.
他就覺得他自己是.
有這個牌子的.
有這個教師牌.
是一個律法師.
也是當時.
這個宗教裡面的領袖.
但是他覺得耶穌.
他沒有正式牌子.
沒有拿牌子.
還沒有被認可.
但是他又處處感受人歡迎.
所以他就想了一些題目.
去為難他.
想試一下.
他是不是真的有.
這個教師的資格.
那耶穌也藉著這個機會.
就點出了.
雖然在摩西五經裡面.
有六百多條的律法.
但是其中的正義就是.
要盡心盡性盡力盡義.
愛主你的神.
然後怎樣去做愛主我們的神呢.
就是要愛淪舍如己.
今天我們就藉著耶穌基督.
所用這個.
好殺瑪利亞人的比喻.
我們來彼此的勉勵.
我們知道.
有一個人.
他就從耶穌的比喻裡面說到.
有一個人.
從耶路撒冷.

$^{81}$去耶利哥城.
其實耶利哥城和耶路撒冷.
相距不是太遠.
大約二十里左右.
但是耶路撒冷.
在高山上高.
在海拔二千呎以上.
而耶利哥就在海拔以下的.
一千三百呎.
相距是很長很多的.
在耶利哥城裡面.
他們的路都很險峻.
還有裡面有很多山洞.
那些賊就很喜歡.
躲在山洞裡面.
當有些人走過的時候.
他出其不意的時候.
他就會走出來.
去搶別人的財物.
甚至會把他打到半死.
在這些荒山野嶺裡面.
這些人被打到半死.
通常到最後都一定死.
在耶路撒冷和耶利哥城.
路途當中有很多危險.
但同樣在人生路程裡面.
都是充滿了很多險惡.
很多陷阱.
就像山洞裡面有很多洞洞.
我們看到.
在前段時間.
有些人被騙去泰國.
被騙去柬埔寨.
前幾天我看到一些散文.
在中國大陸裡面.
竟然有些人.
他自己甘心.
要走到柬埔寨.
去幫忙做欺詐人的工作.
但結果就不被接納.

$^{121}$知道為什麼嗎?.
因為他們電腦又不懂.
很多科技都不懂.
騙不到人.
所以就趕回中國大陸.
結果被抓了.
又要坐牢又要罰錢.
在人生裡面.
我們有時候.
當我們只尊愛我們自己.
我們就不理會其他人.
原來你偷騙這些東西.
他都不理會.
他就想為了賺一點錢.
人就是這樣.
尊為我們自己的利益.
尊為自己的地位.
尊為自己的威嚴.
有時候會做出一些盲目不仁的事.
為了滿足自己的私慾.
我們要發達.
要出位.
要讓人知道.
要有很大的面子.
總之每件事都要爭第一.
結果做出一些埋沒良心的事情.
在聖經裡保羅就說.
這些人是怎樣呢?.
他說你是母親情.
賣主賣友.
任意妄為.
自高自大.
講人是非.
自把自為.
我曾經在一間基督教小學教書.
待課一年.
第一天上學.
我已經覺得火藥味很重.
因為我是待課.
所以校長很恩代我.

$^{161}$不使我開會.
當他們開完會出來的時候.
我看到他們個個都很憤怒.
其中一位老師就大力拍桌子.
就算你用槍指著我.
我也不會相信耶穌.
原來校長就安排了他教六年級的聖經.
在基督教學校裡.
校長是教基督徒.
但那個老師是拜佛的.
因為他們兩者之間有衝突.
所以校長就想為難老師.
安排他教六年級的聖經.
我想校長為了出一口氣.
就要為難人.
但老師就全年黑口黑面.
到最後那位校長退休.
我們就合資買東西送禮給校長.
這位老師就不合資.
他自己送一份禮物.
他就在寬送會送一份禮物給校長.
他拿了一支鉛筆出來送給你.
廣告也有說.
有時我們為了自己的專業.
為了自己的喜好.
為了自己的特性.
我們跟別人吵架.
結果就算我們贏了一場架.
我們會輸了一個家.
輸了一個關係.
但作為耶穌基督的門徒.
我們領受了上帝很多很多的恩典.
我們需要學習.
在一些事情裡面.
我們需要去退後一步.
我們凡事就像在哥林多前書第十三章裡所說.
不求自己的益處.
凡事只求上帝得到榮耀.
別人得到益處.
第二方面.

$^{201}$我們要看看除了強度之外.
我們要看看祭司和利未人.
經文說到.
有一個祭司從那條路經過.
他看到有個人躺在那裡.
他從另一邊繞了一圈就走過去了.
又有一個利未人又經過了.
他看了一看.
同樣繞了一圈就走過去了.
在這個比喻裡面.
這個祭司走過.
他看到有人被打得半死的躺在地上.
但他完全不需要理會就走了.
可能他要趕著去聖殿上班.
趕著時間.
因為到了耶穌時代.
這個祭司可以說是有人門之患.
在出埃及.
上帝藉著摩西定了這個祭司的律例的時候.
亞倫的子孫.
他是可以做這個祭司的.
到了耶穌時代已經很多很多.
以至於那些祭司.
他們可能整年的裡面.
只需要去到聖殿裡面.
做一兩個月的事.
所以這次祭司經過這個路程.
看到有人被打得半死.
但他又害怕.
因為這個律法是講明的.
如果碰到他的時候.
如果他死了.
他又成為不潔淨.
他又不可以去到聖殿.
去擔這個職事.
哎呀不行不行.
等得這麼辛苦.
於是他又繞了一圈路走過去了.
我試過一次.
星期日早上回崇拜.

$^{241}$當時是七點多.
在地鐵站.
去到尖沙咀地鐵站.
看到有人在嘔吐.
我都站一站在那裡.
每逢星期日我都會穿得整齊一點.
乾淨一點.
我就想他嘔吐.
我靠近去幫他.
會的.
你覺得事奉重要還是幫人重要呢.
特別他們在荒山野嶺.
去到聖殿工作重要.
錯過了就要等下年.
還是救人重要呢.
然後那個利未人.
利未人又看到.
這個人就睡在路上.
但同樣地.
又走了.
他可能又想想.
趕時間.
還有自己夠不夠能力幫他.
一個人能不能抬他.
死了又麻煩.
又佔了一些不聖潔的東西.
他又可能想想.
四圍可能是一個騙局.
因為知道有很多山洞裡面.
有些人是假裝受傷.
靠近去幫他的時候.
其他的同黨就會走出來.
連他都會受傷.
特別現在香港.
我們灣仔的地鐵站.
很多人在爛手爛腳.
有時候我們都知道.
這些是騙局.
幫不幫呢.
有時候我都很大的掙扎.

$^{281}$我見到有些小朋友的家長.
拖實小朋友的時候.
有些照樣給錢.
但有些走來騙人.
但我覺得都要求聖靈給我們.
去幫人的時候.
都要靈巧像蛇.
順良如甲子.
怎樣做才是最好的呢.
回說一下.
耶穌基督就說.
這個祭司利未人.
他們都沒有幫到受傷的人.
我們回想一下.
在猶太人的家庭裡面.
他們每一天早上和晚上.
都會唸這兩節經文.
以色列你要聽.
耶和華我們的上帝是獨一的主.
你們要盡心盡性盡力愛耶和華你的神.
這兩節經文.
整段經文叫Shema.
意思就是聽.
用了你要聽這個字.
在希伯來文來說.
是第一個字.
聽不是說我聽到.
左耳進了然後右耳出了.
不是的.
在希伯來文這個字.
什麼叫聽呢.
我們聽了之後.
我們按照裡面那種說話的教導.
我們實行出來的時候.
這個才是叫聽.
如果左耳進右耳出.
又沒有行動.
那些就不是叫聽.
他們這個猶太人.
他們每天早上和晚上.

$^{321}$都會唸一次這兩節經文.
當他們有事.
他們要出門.
出門的時候.
他們有些經文盒掛在門口.
這是現代版.
那些人去完以色列旅行送給我的.
他們就會摸一摸這個盒.
提醒他們自己.
無論是工作還是讀書.
總之你要實行一件事.
就是要盡心盡性盡義.
愛主你獨一的上帝.
所以當律法師來問耶穌的時候.
我應該要做些什麼.
才能得到永生呢.
耶穌就對他說.
律法上寫了什麼.
你是怎樣讀的.
其實在質疑這個律法師.
我們有世之大.
每天讀起碼兩次早上和晚上.
每次出門都要摸一摸這段經文.
究竟在日常生活裡面.
你有沒有走出來呢.
你讀了什麼呢.
耶穌問這個律法師.
然後他就說了.
他就再將這句話.
加上他自己的一個說法.
給我們所有信徒一條新的命令.
他們要盡心盡性盡義盡力.
愛主你的神.
又要愛人如己.
這是耶穌在新約的時候給我的新的命令.
我們每天的靈修.
我們每個星期的崇拜.
有時候我們都會覺得.
又是說那些很悶的.
但如果按聽這個字來說.

$^{361}$是聽了然後實行出來的時候.
才叫做真正的聽的時候.
當我們聽了耶穌上帝的話語的時候.
我們要真正的落實實走出來的時候.
我們才叫做真正的明白.
真正的知道.
接著我們看這個薩瑪利亞人.
經文就說這個薩瑪利亞人.
真是很有慈心.
他上前用油用酒.
幫他敷一敷傷口包紮了.
然後又把他搭上他的傷口.
把他帶到旅店.
並且將來給錢幫他.
說如果這兩個銀幣不夠的話.
遲些我回來的時候.
我一定會補給你們.
是一個很有慈心的人.
但原來這個薩瑪利亞人和猶太人.
是世仇.
大家都很熟悉.
當耶穌基督在井旁.
和薩瑪利亞人聊天.
問他可不可以給我一些水喝.
薩瑪利亞人就說.
我們世世代代都不會相往來.
你為什麼要問我拿水喝.
為什麼猶太人和薩瑪利亞人成為世仇呢.
就是因為薩瑪利亞人.
本來是在北國裡面.
但當北國被亞述所侵佔之後.
他們亡國.
有些人就一直留在耶路撒冷.
有些人就繼續留在北國.
他們因為和周圍的迦南民族通婚.
這班猶太人就覺得很順政.
他們和外族人通婚的人是雜種.
他們誤會稱為上帝的選民.
所以世世代代猶太人和薩瑪利亞人是不會相交的.
薩瑪利亞人做了餅出來.

$^{401}$猶太人一定不會吃.
他們種的麥子也不會買.
他們有任何生畜拿到市場.
他們也不會買.
所以在井旁的時候.
薩瑪利亞人很奇怪.
為什麼耶穌會問他們要水喝和和他們聊天呢.
當時那些猶太人對薩瑪利亞人的看法.
他們可以在會堂裡面.
他們有去會堂裡面敬拜.
他們可以公開地罵他們.
或者可以大聲地祈禱.
求神千萬不要恩待他們.
千萬不要給他們什麼好處.
但在這個比喻裡面.
再一次說.
人要衝破這些仇視的關係.
在世世代代被視為敵人的薩瑪利亞人.
他看到受傷的人是猶太人.
他馬上動了持心.
幫他包紮傷口.
幫他送到酒店.
並且給了錢給他.
這個薩瑪利亞人.
在猶太人裡面覺得他們是雜種.
覺得他們不配做上帝的選民.
但這個被看低被歧視的薩瑪利亞人.
他卻真真正正地走出了上帝的十誡.
你要盡心盡性盡力愛耶和華你的上帝.
他落落實實地從受傷的猶太人裡面走出來.
又要愛倫社如己.
猶太人的倫社.
或者我們一般人的倫社.
覺得住在我們周圍的人就是倫社.
在加州的那些人.
在我們隔壁的就是倫社.
但耶穌基督在這裡跟我們說.
替我們每一個人的倫社.
不單只是住的地方的遠和近.
這個倫社也不會包括我們究竟是不是同一個種族.

$^{441}$是同一個言語.
同一個宗教信仰.
耶穌所說的倫社.
就是當其他人有急切的需要的時候.
我們就要去伸出援手.
現在疫情過後.
逢星期三星期六都會賣旗.
賣旗的時候通常都會看看機構.
看到是基督教的.
我們就會很樂意給下去.
雖然他們不是什麼大.
但有些人說如果都是其他宗教.
我給不給呢.
值得想一想.
耶穌說這個需要.
他人的需要就成為了我們的倫社.
原來愛人如己這句話都不是新的.
在耶穌基督裡面都不是新的.
在利未記第十九章三十三節這樣說.
若有愛人寄居在你們的地上和你同住.
不可欺負他.
寄居在你們那裡的外人.
你們要看他如本地人.
並要愛他如己.
因為你們在埃及地也作過寄居的.
我是耶和華你們的神.
在耶穌在舊約裡面已經吩咐了.
怎樣實行這個盡心盡性盡義.
愛主我們的上帝呢.
就是要愛人如己.
有位姐妹七十多歲了.
她身體也有很多病.
但她的兒子很不幸.
在四十歲的時候又中了風.
他又沒有結婚.
所以她照顧起來有點吃力.
這位弟兄.
他也挺肥挺健碩的.
他很少出街.
特別是在疫情之下.

$^{481}$經常困在家裡.
有位剪頭髮的弟兄.
知道了之後.
他就主動願意出來.
定期上去家裡幫她剪頭髮.
不收錢的.
另一位做物理治療的姐妹.
知道了.
這位七十多歲的姐妹.
真的不能和這麼大隻的兒子出街.
她又自動的.
一個月兩次.
星期日下午.
她就用輪椅幫這位男士.
推他出街.
讓他去吸收大自然的氣息.
又見見街外的人.
這位髮型師.
這位物理治療師.
就成為了中風男子的好鄰舍.
在耶穌眼中.
只要任何人有需要的時候.
我們不需要分任何的界限.
我們就是他的鄰舍.
我們就要伸出援手.
弟妹們.
上帝要我們在整本聖經裡.
無論有多少條律法.
上帝就要我們去.
把愛實踐出來.
上帝是喜悅.
聽祂的話.
又實行祂說話的人.
這就是一個愛鄰舍.
所以耶穌說.
愛你的鄰舍.
照樣做吧.
我們一起禱告.
因為你先愛我們.
以致我們可以學習.

$^{521}$怎樣去愛其他人.
但有時候我們都會軟弱.
生命軟弱.
心靈軟弱.
我們愛我們自己會多一點.
求你赦免.
幫助我們吧.
因為這個世界裡.
真的有很多的凶險.
有很多的欺騙.
很多的事情.
我們真的不明白.
也不知道.
幫助我們.
看到有人有需要的時候.
給我們有智慧.
給我們有認識.
以致我們在息事裡.
我們去伸出援手.
求主恩代.
祝福我們每一個人的生命.
就是主你律法.
愛的實踐.
這是上帝你所喜悅的.
求主幫助.
賜福.
讓我們每一個人.
在你面前的人的生命.
都有一個愛的生命.
多謝你水天禱告.
奉靠主基督的聖名.
阿們.
\newpage



\section{約伯記 28:1-28}
\label{sec:zrsNFCgx1bQ}
\textbf{2022-05-22 午堂《明智藏哪裡》約伯記 28\_1-28 講員:周曉暉牧師}
\newline
\newline
連結: \href{https://youtube.com/watch?v=zrsNFCgx1bQ}{\texttt{https://youtube.com/watch?v=zrsNFCgx1bQ}} ~~~~ 語音日期: 2023-08-25
\newline
\newline
\hyperref[sec:TEOKFaxZL4g]{\small{< < < PREV SERMON < < <}}
~
\hyperref[sec:index_chronic]{\small{[返順時目]}}
~
\hyperref[sec:index_scriptual]{\small{[返順卷目]}}
~
\hyperref[sec:lKgHJObwVL4]{\small{> > > NEXT SERMON > > >}}
\newline
\newline
約伯記 28:1-28
\newline
\begin{longtable}{cl}
\hline
\hline
章節 & 經文 (和合本修訂版)\\
\hline
28:1 & \begin{tabularx}{0.7\textwidth}{X} 「銀子有礦；煉金有場。 \end{tabularx} \\ \\ \relax
28:2 & \begin{tabularx}{0.7\textwidth}{X} 鐵從土裡開採，銅從礦石鎔出。 \end{tabularx} \\ \\ \relax
28:3 & \begin{tabularx}{0.7\textwidth}{X} 人探索黑暗的盡頭，查究礦石直到極處，那是幽暗和死蔭； \end{tabularx} \\ \\ \relax
28:4 & \begin{tabularx}{0.7\textwidth}{X} 他在無人居住之處開鑿礦穴，在無足跡之地被遺忘，與人遠離，懸空搖擺。 \end{tabularx} \\ \\ \relax
28:5 & \begin{tabularx}{0.7\textwidth}{X} 地出產糧食，地底翻騰如火。 \end{tabularx} \\ \\ \relax
28:6 & \begin{tabularx}{0.7\textwidth}{X} 地的石頭是藍寶石之處，那裡還有金沙。 \end{tabularx} \\ \\ \relax
28:7 & \begin{tabularx}{0.7\textwidth}{X} 鷙鳥不知那條路，鷹眼也未曾見過。 \end{tabularx} \\ \\ \relax
28:8 & \begin{tabularx}{0.7\textwidth}{X} 狂傲的野獸未曾踩踏，猛烈的獅子也未曾經過。 \end{tabularx} \\ \\ \relax
28:9 & \begin{tabularx}{0.7\textwidth}{X} 「人動手鑿開堅石，翻倒山的根基， \end{tabularx} \\ \\ \relax
28:10 & \begin{tabularx}{0.7\textwidth}{X} 在磐石中鑿出水道，親眼看見各樣寶物。 \end{tabularx} \\ \\ \relax
28:11 & \begin{tabularx}{0.7\textwidth}{X} 他封閉河川不得涓滴，使隱藏之物顯露出來。 \end{tabularx} \\ \\ \relax
28:12 & \begin{tabularx}{0.7\textwidth}{X} 「然而，智慧何處可尋？聰明之地在哪裡？ \end{tabularx} \\ \\ \relax
28:13 & \begin{tabularx}{0.7\textwidth}{X} 智慧的價值無人能知，活人之地也無處可尋。 \end{tabularx} \\ \\ \relax
28:14 & \begin{tabularx}{0.7\textwidth}{X} 深淵說：『不在我裡面。』滄海說：『不在我這裡。』 \end{tabularx} \\ \\ \relax
28:15 & \begin{tabularx}{0.7\textwidth}{X} 智慧不可用黃金換取，也不能用白銀秤她的價值。 \end{tabularx} \\ \\ \relax
28:16 & \begin{tabularx}{0.7\textwidth}{X} 俄斐的金子和貴重的紅瑪瑙，以及藍寶石，不足與她比擬； \end{tabularx} \\ \\ \relax
28:17 & \begin{tabularx}{0.7\textwidth}{X} 黃金和玻璃不足與她比較；純金的器皿不足兌換她。 \end{tabularx} \\ \\ \relax
28:18 & \begin{tabularx}{0.7\textwidth}{X} 珊瑚、水晶都不值得提；智慧的價值勝過寶石。 \end{tabularx} \\ \\ \relax
28:19 & \begin{tabularx}{0.7\textwidth}{X} 古實的紅璧璽不足與她比較；純金也不足與她比擬。 \end{tabularx} \\ \\ \relax
28:20 & \begin{tabularx}{0.7\textwidth}{X} 「智慧從何處來呢？聰明之地在哪裡？ \end{tabularx} \\ \\ \relax
28:21 & \begin{tabularx}{0.7\textwidth}{X} 她隱藏，遠離眾生的眼目，她掩蔽，遠離空中的飛鳥。 \end{tabularx} \\ \\ \relax
28:22 & \begin{tabularx}{0.7\textwidth}{X} 毀滅和死亡說：『我們風聞其名。』 \end{tabularx} \\ \\ \relax
28:23 & \begin{tabularx}{0.7\textwidth}{X} 「神明白智慧的道路，知道智慧的所在。 \end{tabularx} \\ \\ \relax
28:24 & \begin{tabularx}{0.7\textwidth}{X} 因為他鑒察直到地極，遍觀普天之下， \end{tabularx} \\ \\ \relax
28:25 & \begin{tabularx}{0.7\textwidth}{X} 要為風定輕重，又度量諸水， \end{tabularx} \\ \\ \relax
28:26 & \begin{tabularx}{0.7\textwidth}{X} 為雨定律例，為雷電定道路。 \end{tabularx} \\ \\ \relax
28:27 & \begin{tabularx}{0.7\textwidth}{X} 那時他看見智慧，就談論她，堅定她，並且查究她。 \end{tabularx} \\ \\ \relax
28:28 & \begin{tabularx}{0.7\textwidth}{X} 他對人說：『看哪，敬畏主就是智慧；遠離惡事就是聰明。』」 \end{tabularx} \\ \\
[1ex]
\hline
\hline
\end{longtable}
$^{1}$我從北角帶來向大家的問安.
大家都是主內宣道會的弟兄姊妹.
願你們在疫情裡面充滿主的恩典.
讓我今天可以和錦繡堂的各位分享神的話語.
今天我們看《約伯記》.
我相信你過去也略略看過.
整本看不容易.
所以今天我們也抽取其中一章.
去思考其中神給我們的真理.
信主年日越長.
我不知道你自己有沒有感受到.
你的熱情.
信仰的熱情.
是消退的.
遇上劇變.
甚至你會覺得這個世界都是倒塌的.
信仰對於你來說有一種.
沒什麼力量的感覺.
更大的問題.
更大的困局.
更大的苦難.
有時候人是沒力去應對.
反而我們是很希望得著的.
有沒有一些更實在.
人生的洞見.
人生的insight.
可以讓你去經歷到呢.
《約伯記》.
是一卷智慧的文學.
我們都認識了.
它是抽空這樣建構了一個極度差的處境.
一個異人.
他在生命裡面經歷.
一個很大的難關.
很多的刺在他生命當中經過.
而他這個正直的人.
因為一場天庭的辯論.
無情的苦難突然間就臨到.
約伯頓失所有.
他很不明白.

$^{41}$身邊的朋友來問他.
探他.
大家開始拿出一句一句的淑寧的金句.
甚至開始跟他理論.
但是各人都解答不了.
為什麼為什麼為什麼.
約伯會遭遇到這一切的事.
而且到了後期他們還.
blame the victim.
你是不是做了一些很壞的事啊.
很多的壞事我們不知道.
有報應啊.
他們說了二十多章.
到了二十七章.
這位約伯仍然堅持自己是清白.
正直.
然後兩章之後的二十九章.
約伯.
轉向神.
他在神面前去起誓.
他要求神去回答他為什麼.
夾在二十七章.
朋友的控訴.
和二十九章要神去回應.
這個中間.
就是我們今天要看的這一章二十八章.
這一章不是出自約伯.
或者他朋友的口.
而是來自經文的旁白所讀出的一篇詩歌.
正正就是因為之前他們的討論.
這幾個朋友和約伯的辯論.
車輪戰的樣子.
裡面很多說話道理.
此是而非.
旁白在這個位置要停一停大家.
要打破一個沉默.
好好引導當時聽這本書的神的子民.
引導今天我們讀這本書的人.
回到正軌上再去看.
再去看.

$^{81}$所以今天我們讀約伯的二十八章.
我們一起去思想我們自己的人生.
到底智慧在哪裡.
願意這首與別不同的詩歌.
一起校正我們的心態.
讓我們找回一個動力.
回到我們生命信仰裡面一個很基本.
我們怎樣面對我們所讀的聖經的態度.
我們就一起祈禱教託.
天主上願意你的話語.
讓我們心靈有振奮.
願意你的靈在我們裡面.
指教我們當明白的話.
多謝你.
祈禱是奉主耶穌基督.
得勝名自祈求.
阿們.
大概很快就說.
牧師我知道了.
一讀經文大約都知道你想說什麼.
人生的智慧在神那裡.
要認識多點上帝.
記得翻主的學.
記得讀經靈修.
敬畏主.
不要做壞事.
那我們就得著人生的智慧了.
對.
不過我們在這裡再看多一點.
再看多一點.
既然這件事這麼巧這麼重要.
為什麼我們不是很熱衷地做.
第二.
什麼叫做敬畏主.
遠離邪惡.
第三.
聖經這裡所說的智慧.
對於經歷一個連續的苦難的信徒.
那個用處在哪裡.
我們今天就透過《約伯記》這個28章.

$^{121}$一起去找人生的洞見.
要找寶藏.
藏寶圖就十分重要了.
我不知道你小時候.
我小時候是看《鮮巴力點記》的.
拿著地圖就去尋寶了.
是呀找寶藏.
要有藏寶圖.
要知道那個地方在哪裡.
然後你才找到對的人.
找到那樣的東西.
《約伯記》28章三個地方提過.
原文提過地方.
同一個字而已.
第一節那裡說.
今只有惹念的地方.
然後12節.
聰明之地.
聰明的地方在哪裡.
23節.
智慧的所在.
其實原文都是所在的地方.
中文未必能夠翻譯出來.
三個地方指向.
我們今天要掘的寶藏.
這個藏寶圖帶我們去找.
很重要的這件事.
第一段我們剛才看.
告訴我們.
其實我們人都真的OK的.
不是一開始就說.
人很不濟.
墮落.
人有罪.
這裡說人的智能很高.
你看看第三第四節.
不如我們一起讀第三節.
預備起.
人探求黑暗的盡頭.
查明在陰暗中的石礦.

$^{161}$直到極處.
第四節說在無人居住之處.
足跡罕至之地.
開坑探礦.
他們與人遠離.
原在坑中搖來搖去.
這裡是二三千年前的人的寫照.
他們當時已經懂得將.
地裡的山裡的.
很遙遠的天涯海角.
沒有人居住的地方.
那些貴金屬.
這裡有說.
其實經文繼續說.
銀啊金啊鐵啊銅啊藍寶石啊.
各樣各樣的珍寶挖出來.
這些是當時人視為至寶來的.
很有價值的.
最近的新聞.
我不知道你有沒有看到.
說.
有些科學家開始拿月球.
拿回來的那些泥土.
去種植.
為什麼.
他們要試一下.
月球的泥是不是種到一些植物.
然後如果是成功的話.
就預備在月球上面.
將來建一個基地.
那就.
以至他們可以再去遠一點的地方.
不用每次都由地球出發.
在月球出發會方便一點.
近一點.
便利一點.
預備去哪裡.
去火星.
那去火星做什麼呢.
那就不是找火星人了.

$^{201}$火星原來有很多的.
很貴的貴重的金屬.
稀土.
地球現在欠稀土.
稀土是什麼呢.
就是你弄你手機那些.
那些的晶片.
很重要的一些物品.
所有的電腦都要用稀土.
去提煉一些的礦物.
誰擁有稀土.
就擁有現在世界很多的資源.
但是大家各人去爭了.
不夠了.
於是要去更遠的地方.
去火星.
希望可以掘到這些重要的.
礦物和戰略資源.
經文特異顯出人類的與眾不同.
他以當時很厲害的動物來比較.
第七節那裡說到飼鳥和老鷹.
這些你都知道的了.
小朋友都明白.
天空的霸主.
飛鳥的霸主.
而他們的視力.
這裡說的是看東西.
他們的視力是動物裡面最頂尖的.
不過.
他們都看不透.
人類要去的地方.
地上最厲害的猛獸獅子.
王者.
勇猛.
沒有什麼地方他們不敢去的.
即使是這樣.
都有些他們未能夠踏足.
人類將要去的地方.
當然今天我們更加明白.
他們不可能上太空.

$^{241}$人很厲害.
是不是.
在這些背景.
在這些河谷底.
甚至水底下的地方.
人只要能夠建造水壩.
其實經文繼續說下去.
截斷水源.
就會露出那個河底.
水底.
他就可以挖到那些寶物.
總之這個大地有什麼寶藏.
人類就用盡一切的.
財智,技術,力量,資源.
找到為止.
其實你明白的.
當你覺得一件事很有價值.
你很重視它的時候.
你是會排除萬難地去做的.
我和小朋友都很喜歡看太空的科學節目.
其中一個片段是說到.
一個無人太空探測器.
他叫做卡西尼號.
他在地球出發經過七年之後.
去到很遠的土星.
然後這個探測器就圍繞土星.
在這裡運行十幾年.
期間不斷不斷每天去收集.
不同的資料,數據.
然後拍了很多照片.
就經過很遠的路途.
傳送回來地球.
這個卡西尼所傳送回來的數據.
讓人類揭開土星那個環.
那麼漂亮的光環的秘密.
是什麼來的.
它可以存在多久.
科學家甚至可以很準確推論到.
這個光環還可以存在一億年.
不過科學節目裡有一幕.

$^{281}$是說一個女科學家.
她說她每天最緊張的事.
每天一回到基地裡最緊張的事.
就是立刻去檢查這個無人機.
送回來的資料.
看看電腦,看看那個機.
有沒有印出什麼東西.
然後就檢查這個卡西尼號的數據.
那些機器的健康數據.
是不是運作良好.
那些訊號接收是不是沒有受阻隔.
傳送回來的資料有沒有亂碼.
她說這些事很重要.
她形容自己好像這個無人探測器的媽媽.
每天就是為了這部飛行器去掛心,操心,留心.
當中我相信你們有媽媽的.
你怎樣為孩子操心.
這個女科學家就是這樣去掛心.
那個很遠在土星的太空探測器.
當人重視一件事的時候.
是會掛心,留心,操心和賞心的.
是會排除萬難這樣去做的.
我們都知道神的話語很重要.
既然是這樣.
為什麼我們不是想知道更多.
不是在我們很忙碌的人生裡.
怎樣都找些時間去學去問去.
找出其中的奧秘呢.
約伯記的詩歌繼續去形容給我們聽.
他以採礦,煉金,煉鐵,煉銅當時最高的科技.
如果今天寫約伯記的話.
大概就是講太空科技,軍事科技.
我們的互聯網科技.
他表達人類的文明那種科技的發展.
人類的智能超卓.
今天我們的建築物已經可以穿越雲頂了.
我來到這裡就看到大家景色優美.
全部都是平房.
我曾經在將軍澳住過.
住50樓那麼高.

$^{321}$住50樓高有什麼好處呢.
差不多就是因為將軍澳很濕.
到了春天的時候就一片雲霧.
但是我就稍稍地因為我高.
就在雲霧的上面.
感覺到人類的建築就是這麼厲害.
雲上太陽那樣經常見到.
是啊,我們人類的智能科技就是這麼高卓.
但是我們的.
即使我們知識增長.
我們對生命的洞見是不是更有.
insight,有更加洞見.
是不是對人生更加知道.
哪些是可為,哪些是不可為.
是不是這樣.
其實你打開你的手機,你上IG.
你上Facebook.
你上任何的討論區.
看YouTube.
你見到的.
是美善,智慧,德行.
還是更多的惡俗,愚蠢和邪人.
弟兄姊妹有時候我們覺得.
神的話語很重要.
不過有時候真的挺悶.
會不會因為其實我們心底裡覺得.
這個世上有更多的東西值得我們花時間.
更加實際.
值得我們去不斷追尋去探索.
去得到.
就好像我們當中在座的弟兄.
不知道你是不是和我一樣.
很喜歡科技玩意.
很喜歡音樂音響.
又或者很喜歡機械汽車.
總不知道現在有消費券就開始新行.
就想買東西.
於是怎樣?找很多資料看.
不斷上網找.
你知道的,一找就找幾個小時.

$^{361}$一看就每天都在看.
為什麼?因為覺得那些東西.
很想得到,很重要.
很喜歡,很寶貴.
於是我就.
賞心,操心,留心.
花時間.
在這些事情上.
聖經呢?.
關於神的事呢?.
約伯記28章提醒我們.
是有科技,人類智能裡面.
讓我們看到.
我們的世界很吸引.
但裡面你是找不到.
真正人生的意識.
你找不到真正的智慧.
約伯記.
帶我們去想,那在哪裡找?.
第二段,十二節.
我又請你跟螢幕一起讀,預備起.
然而智慧在何處可尋.
聰明之地在哪裡呢?.
聖經裡面所說的智慧.
就不是智慧.
不是在說經驗.
不是說我們頭腦積累的知識.
也不是說那些才智.
有另一個譯本.
將第二句頭兩個字.
聰明,譯為明智.
明智,就是Understanding.
Insight.
對事物有透徹的領悟.
這個是洞見,這個是明智.
人類無疑有很高的智能.
但未必對人生有洞見.
早五六年.
香港有一個流行的詞.
高分,低能.

$^{401}$大概你明白了.
是哦.
我們發展到今天.
人類有互聯網更加厲害.
即使是這樣.
是否真的對人生有洞見呢?.
對生命看得透呢?.
十三節上半那裡說.
聖經裡面所說的智慧的價值.
人不知道.
是因為不知道所以不去找.
不想去找.
當有一天真的知道了.
知道極為重要了.
於是就用他一生所積累的財富,知識,技術.
希望盡一切去換.
十五至十九節告訴我們.
拿光所有的寶物去換.
都換不到.
因為那種價值.
是無法比較.
是買不到這種智慧的.
我們所有的都不能換.
我不知道這是否我們在疫情裡面的寫照.
有些弟兄姊妹在疫情裡面很努力.
沒有浪費時間.
網上這麼多資源.
這麼多不同的課程.
主要是學講座,講道.
拼命聽.
也都可以.
有些人真的趁這個機會去讀一些書.
去進修.
去上一些無牆課室等等.
但你慢慢發覺.
有時候即使我們以這種的做法.
用你有的金錢,時間.
在一個短時間裡面去換.
你是不是真的對人生.
神給我們的生命更有動見呢?.

$^{441}$十三節下半節這樣說.
這種的智慧在活人之地找不到.
在我們人間裡面是找不到這種真正的智慧.
大概那個應用就可以對我們今天的基督徒說.
是呀,人類的技術很厲害.
建很多很美麗的建築物.
很宏偉的建築物.
甚至.
你可以建一個很莊嚴神聖的教堂.
不過這個不等於我們真正是神聖.
我們有宏偉的堂子.
不等於我們真正擁有聖經裡面的智慧.
約伯記繼續問我們.
哪個地方找得到呢?.
他繼續在十四節列舉人.
未能夠探索的地方.
譬如.
深淵,滄海.
海洋是古代人所怕的東西.
古代人類征服不了大自然的力量.
其實今天我們也征服不了.
兩個多月前.
半夜時間.
不知道你有沒有被搖醒.
有些人點頭.
那晚睡著睡著.
好像被人搖醒了.
那就看看是不是老婆搖醒我.
不是的.
原來第二天早上起來.
一看就有新聞告訴我們.
在香港東南方.
92公里的淺層處.
發生了4.1級地震.
所以當天.
有很多市民在睡夢中感受到那種搖晃.
有八千人打電話去天文台查詢.
甚至有人報警.
搖一搖.
就令我想起三月份.

$^{481}$十一年前.
311日本東北大地震.
你看到的.
你在視像看到的.
那種海底一震的時候.
觸發的海嘯.
轉瞬間席捲了.
福島那些地方.
沖爛了他們的核電廠.
大自然的力量的大.
是令人懼怕的.
不過即使力量這麼大.
智慧都不是其中.
這個讓我們體會到.
是喔.
我們人世間裡面.
很懼怕一些力量.
看得見的.
看不見的.
在大自然的.
在我們人類社會的.
這些力量.
我們會怕.
但是裡面.
經文讓我們知道.
智慧不在其中.
若果你依從這些力量去生活.
仍然得不到人生的洞見.
經文來到第20節.
他再次敏感.
智慧在哪裡來的.
聰明的地方在哪裡.
哪裡挖到這個寶藏.
這句說了算.
跟剛才讀的12節很相似.
你看清楚一點.
有一點點的分別.
12節.
那裡說可以在哪裡找.
人可以去哪裡.

$^{521}$去到找到.
答案是無處可尋.
用了很多節跟你說了.
20節就是.
沒得去找了.
那我們可不可以知道.
它的來源在哪裡.
從哪裡來.
找不到就問它的來源.
不在這個大地.
不在我們眼見的事物當中.
於是人就轉向靈界去尋求.
22節那裡說.
滅亡和死亡.
都說我們只封聞他的事罷了.
人轉向問一些靈界的事情.
滅亡原本的讀音是Abaddon.
這個字在啟示錄都有用過.
我列出來給你看.
九章十一節.
那裡說.
他們有無底坑的使者.
作他們的王.
按著希伯來話名叫阿巴頓.
這個就是滅亡.
希伯來文滅亡的讀音.
滅亡阿巴頓.
無底坑的使者.
其實就是形容啟示錄裡說的撒旦.
撒旦和死亡的權勢.
是超自然靈界的力量.
他們似乎很厲害很厲害.
知道更多的東西.
其實古代的人都是這樣去想的.
於是就有些人去占卜去問鬼.
你記得.
素羅王去問米.
他找文米婆召喚薩姆爾的靈魂.
當時這個聲稱是薩姆爾的靈魂.
上來跟他說了些什麼.

$^{561}$在素羅王人生很混亂.
他要找到一個洞見insight的時候.
他轉向尋求靈界的事情.
靈性的事情.
他會去問米婆.
從這個靈界起得到一些新的資訊.
有沒有啊?.
沒有啊.
沒有新的東西.
沒有洞見.
沒有他見完之後還害怕.
當然今天我們不會去問米.
二十世紀末了.
在靈性上.
上個世紀.
我們的世界開始著重spirituality.
不只是基督教.
不同的世俗的科學.
人民的科學都是講靈性.
不是教會的專利.
所以有一大堆的東西出現了.
我講一些詞.
看看你有沒有聽過.
治癒.
心靈刑事.
洗練.
洗滌.
放空.
淨化.
能量.
正念.
你總是聽過一兩個的.
為什麼有這些東西出現呢?.
就是因為現代人過著一個很多物質.
享樂.
消費的生活形態.
但是在我們心底裡.
我們的確會感受到.
我們意識到我們靈性裡有缺乏.
甚至饑荒.

$^{601}$於是人就去找.
於是世界就給一些solution出來.
給一些product出來.
是啊.
各位基督徒你不去找.
其他人會去找的.
在美國的一位律師.
他叫做Angela Krummer.
她有一次上TED Talk.
講了這番話.
她的題目是.
My Failed Mission to Find God.
她說她自己是一個極重靈性的人.
她不是基督徒.
她說她相信某種程度上.
必定有神的存在.
她說她只是不確定.
這個神是誰.
她決定踏上自己找神的旅程.
她向世界不同的宗教去探索.
但是她說她下定決心.
在這個靈性之旅裡.
就好像我們基督徒的用詞.
她說她要避開那些明顯的地方.
略過那些大型宗教.
就是她不會來教會的了.
她說取而代之.
她說她要冒險.
還有一些.
不是那麼多人踏進的靈性的邊緣.
像不像.
她說去不同的地方探索.
她說她參加了一些.
由巫師舉辦的Potluck.
她去了位於紐約市的.
LGBT多元中心.
在那裡和一些僧侶.
一起喝了五加侖的火山水.
試試會不會體驗到.
神的靈性充滿.

$^{641}$她在另一個會議中心.
得到一個印度聖人的擁抱.
她說她是一個很熱.
很多煙的印第安蒸氣浴室裡面.
在那裡重念咒語幾個小時.
她又去過南美洲一個海灘上.
和一個喝著Tequila酒的靈媒.
一起召喚亡靈.
她知道自己靈性缺乏.
她說她要窮一生.
去努力找到.
可以令自己好像神那麼完美的信念.
她去找.
她會找到嗎.
當然我們今天在教堂裡.
知道她一定找不到.
她錯過了什麼.
她從人性期待美好的角度去找.
其實說到底.
她都是在找健康快樂愛情滿足.
她認為只要某一個信仰.
給她這些東西.
那個就是神.
她甚至很想努力成為.
好像神那麼完美.
不過我們看到這些事例.
起碼她去找.
比我們懶的時候好.
連神都不想找.
你又不是我回來教會的.
我不知道你今天早上出來.
有沒有內心交戰.
我今天是留在家裡看.
還是回來實體呢.
不是在家裡看不好.
而是我們總是有很多理由.
給我們去推搪.
不像別人那麼積極.
去尋求靈性的需要.
當然我們仍然有弟兄姊妹.

$^{681}$很期待有超凡的屬靈體會.
去讀不同的課程.
參加網絡的教學.
甚至有些去體驗一些特會.
不過就是不肯好好坐下.
簡簡單單拿起聖經來讀.
超越銀河系的智慧在哪裡來呢.
明智藏哪裡呢.
是否在人世間的靈性尋索裡.
是否在我們所體會到的.
各種宗教體會體驗裡.
一些很興奮的聚會.
很受感動的唱頌.
是否在那裡.
既然不是在銀河系.
是否在人類裡的靈性操練呢.
我們剛才看過也不是.
來到第三段23節.
我請你一起讀.
螢幕預備起.
神明白智慧的道路.
也曉得智慧的所在.
弟兄姊妹我們很多時候讀藥百技是被迫的.
當我們受苦我們想找答案.
就打開藥百技看看.
其實會不會是我們用錯了方法.
問錯了問題.
是喔事實上我們可能進入了一個基督徒的功利循環.
我們很想讀神的話語是很有用的.
每一次讀完最好有亮光.
有大正正切中我要找的東西.
給我指引.
給我事先的提醒.
讓我做好抉擇.
去還是留.
打不打第四針.
當你面對聖經.
你是不是當聖經是一個運程和星座.
有些弟兄姊妹會覺得勝經最好給我安慰.
最近已經這麼多事情不行.

$^{721}$環境又不好.
可不可以有些經文安慰一下我.
你以為療癒小品文嗎.
最好給我前路更加樂觀.
心靈更加正面.
你以為心靈雞湯.
最不好的就是你覺得上帝會因為你.
有時候忘記讀聖經而罰你.
你當上帝是偶像.
少做少添香油降災給你.
其實明明讀聖經只是一個很自然平淡的事.
對生命一件有益平凡的事.
但我們就是一直不想碰.
一碰就要大鑼大鼓.
特別講座,培靈會要有刺激.
有新意才聽得入耳.
我想我們搞錯了重點.
你拿起這本聖經你求什麼.
耶穌對著當代的法利賽人.
他這樣說.
你們查考這經以為內中有永生.
今天我們讀聖經覺得沒有領受.
是不是因為我們只是當聖經.
是一個挖到我要保障的地方.
我這樣說.
如果你開路你就指示我吧.
聖經為你服務.
難怪我們越讀越乾.
越讀越悶.
越來越覺得找不到答案.
因為我們真的不知道自己要找什麼.
耶穌繼續說.
給我作見證的就是這經.
然而你們不肯到我這裡來得生命.
耶穌的說話就是提醒.
你要找的不是那些道理.
而是神自己.
你可不可以用你本聖經認識天父更多.
約伯記重重複複問一個問題.
一個這麼愛神.

$^{761}$一個這麼公義的人.
當他沒有健康,財富,快樂,愛情.
你知道他老婆也埋怨他.
你還剩下什麼.
來到這裡.
他還剩下從上帝而來的智慧.
以至他開始有動見.
轉面去找上帝二十八章.
接著說下去.
真正的智慧就在神那裡.
因為神體見,描述,建立,查探.
這個智慧的價值.
智慧的來源.
是從上帝那裡來的.
當人類越探索宇宙,觀察宇宙.
我不知道當中有沒有弟兄你都熟悉的.
大爆炸理論.
有些姐妹都熟悉的.
不要這麼歧視了.
科學家仍然想找出一個起源.
於是他們在歐洲建立了一個很大的實驗設備.
叫做大型強子對撞器.
你忘記了就記得這三個英文字母.
LHC.
上網查一下.
這個是一個將物質裡面很小的質子.
將它加速.
透過一個機器加速.
接近光速.
然後兩粒質子對撞.
產生一個超高的能量.
多高呢?.
就是14萬億電子伏特.
很厲害.
希望這個撞擊的能量可以轉換到.
找到其中再小的物質.
叫做希格斯波息子.
不是波子.
是希格斯波息子.
科學家稱之為神的粒子.

$^{801}$很希望在這個物質裡面找到最最最最最微小的粒子.
去問到宇宙的起源.
不過即使他找到這粒子.
他還要再問.
它是從哪裡來的?.
是誰去指揮當時第一個的大爆炸?.
誰指引這個萬物的出現?.
正經明確說了.
一早說了是上帝自己.
是祂的創造.
智慧不是世上的智力經驗.
知識,聰明才智得出來的結論.
真正的智慧是有來源的.
這個來源不在銀河系外邊.
而是從神而來.
所以神自己的說話.
我們一起讀28節.
預備起.
他對人說敬畏主就是智慧.
遠離邪惡就是聰明.
什麼叫做敬畏上帝遠離邪惡?.
敬畏當然好像一種害怕的情緒.
但敬畏也有分兩種.
因為害怕而逃跑.
嚇到逃跑.
現在很多人要走了.
另一種是怕他走的感受.
就好像怕沒有了的感受.
就好像你坐的士,坐巴士,坐私家車.
坐飛機你會帶安全帶.
因為你怕沒有了安全帶.
萬一有意外你會出事.
你多不舒服也好.
你都不敢不帶.
不是不是.
有些人怎樣都不帶.
那我說另一樣.
你一定不敢不帶.
電話.
沒有了電話很不方便.

$^{841}$你會怕很多事做不到.
連給錢都給不到.
不能坐車.
你不敢不帶.
你怕沒有了電話.
或者眼鏡.
你一起床就戴眼鏡.
因為你不戴就看不清楚.
你不敢不帶這些物品.
因為你怕這些物品遠離你.
令你失去了幫助.
力量和安全.
你不想它遠離你.
這一種都是一種敬畏的心情.
不只是怕神.
你罰我.
不只是這樣.
怕沒有了.
而且敬畏不是慌張.
它是一種活力.
一個上鏈的音樂盒.
越上鏈越敬畏.
越上鏈越慌張.
越敬畏.
扭得越緊.
放手的時候.
音樂就更加微妙.
動聽.
清晰.
有方向.
當你敬畏.
就有動力注入你的生活當中.
以至你聽得見上帝的說話.
行出祂為你人生所譜出的樂章.
那遠離邪惡呢?.
就是你越接近光.
你越見到自己的黑暗面.
你越近神.
讓神在你身邊.
你就有個動見.

$^{881}$有光就看得見哪裡有不妥.
知道無論如何.
都要放下這些纏累你和我的罪.
丁姐妹.
你在網上很久了.
這幾個月.
你自己在家裡很久了.
你知道你跌到一個位置.
近光.
放下它們.
敬畏主.
就是智慧.
遠離邪惡.
就是聰明.
《約伯記》說這麼多的智慧.
對一個受苦的約伯.
或者今天的信徒.
的用處在哪裡?.
原來說這麼多.
《約伯記》28章.
讓我們看見面對苦難.
我們本來問一大堆27章這麼多的.
Why?.
Why me?.
但是去到這個位置.
開始懂得問.
Who?.
誰?.
Who am I?.
我們是誰?.
即使我們智能超卓.
科技發達.
擁有很多.
我們都只不過是出於塵土.
我和孩子說.
看完這些片我就和他說.
太陽系是位處於銀河系.
這個我們大家都認識.
在宇宙裡面有多少個.
好像銀河系這樣的Galaxy呢?.

$^{921}$原來在宇宙裡面.
天文學家發現.
根據他們已經放上太空望遠鏡.
你知道有個太空望遠鏡.
不是在地球看上去的.
在太空可以望得更遠.
去觀察整個宇宙.
科學家已經觀察到.
整個宇宙大約有2000億個.
銀河系這樣的星系.
你想像不到.
這麼大.
Who am I?.
所以你明白.
為什麼約伯後來.
他見神.
神要和他說那些東西.
他感受到了.
Who am I?.
我們本來出於塵土.
對於整個宇宙銀河.
我們多麼渺小.
所以我們要問.
Who?.
還要問.
Who made us?.
誰造我們?.
甚至是誰提升我們.
這些塵土竟然可以和一個.
創造宇宙銀河.
2000億個銀河系的主對話.
好像約伯那樣.
可以和他質詢.
和他祈禱.
求他給答案.
可以抗議.
問他拿.
自毀.
弟兄姊妹來到這裡.
約伯記這一章.

$^{961}$就是讓我們在苦難裡面.
從Why?.
問Who?.
願意我們都存著敬畏的心.
你就平平實實.
拿起你的聖經.
打開你的聖經APP.
下載了很久.
有些電話有個功能.
那個APP不用.
會自己uninstall.
你的聖經APP.
沒有被uninstall.
你就打開你的聖經.
去讀上帝的話語.
認識它.
自毀在哪裡.
你明知藏哪裡.
不就是在那份敬畏裡面.
不是慌張.
而是怕它遠離你.
不可以失去.
讓神的話語.
祂的生命燃點.
你注入動力.
好像上鏈一樣.
你就不會亂了.
你知道怎樣走下去.
你可以抱出.
上帝早已為你預備人生的樂章.
明知藏哪裡.
願我們每一個人都找到.
我們一起祈禱.
天外上帝.
我們多謝你.
想起我們的渺小.
你的偉大.
你卻這樣去愛我們.
讓我們有你的話語.
可以拿在手上.

$^{1001}$所以主一讓我們願意打開聖經.
我們每天去讀.
打開我們的聖經APP.
打開我們的聖經網站.
我們就去唸頌你的說話.
一點一滴地吸收明白.
得著你所賜人生的洞見.
多謝你.
我們的祈禱是.
奉主耶穌基督得勝命之祈求.
阿們.
\newpage



\section{出埃及記 23:1-13}
\label{sec:lKgHJObwVL4}
\textbf{2022-12-18 加恩樓《爭議中持守公正與良善》出埃及記 23\_1-13 講員:周曉暉牧師}
\newline
\newline
連結: \href{https://youtube.com/watch?v=lKgHJObwVL4}{\texttt{https://youtube.com/watch?v=lKgHJObwVL4}} ~~~~ 語音日期: 2023-08-25
\newline
\newline
\hyperref[sec:zrsNFCgx1bQ]{\small{< < < PREV SERMON < < <}}
~
\hyperref[sec:index_chronic]{\small{[返順時目]}}
~
\hyperref[sec:index_scriptual]{\small{[返順卷目]}}
~
\hyperref[sec:iRhu93YiYcU]{\small{> > > NEXT SERMON > > >}}
\newline
\newline
出埃及記 23:1-13
\newline
\begin{longtable}{cl}
\hline
\hline
章節 & 經文 (和合本修訂版)\\
\hline
23:1 & \begin{tabularx}{0.7\textwidth}{X} 「不可散佈謠言；不可與惡人連手作惡意的見證。 \end{tabularx} \\ \\ \relax
23:2 & \begin{tabularx}{0.7\textwidth}{X} 不可附和群眾作惡；不可在訴訟中附和群眾歪曲公正，作歪曲的見證； \end{tabularx} \\ \\ \relax
23:3 & \begin{tabularx}{0.7\textwidth}{X} 也不可在訴訟中偏袒貧寒人。 \end{tabularx} \\ \\ \relax
23:4 & \begin{tabularx}{0.7\textwidth}{X} 「若遇見你仇敵的牛或驢迷了路，務必牽回來交給他。 \end{tabularx} \\ \\ \relax
23:5 & \begin{tabularx}{0.7\textwidth}{X} 若看見恨你的人的驢被壓在重馱之下，不可走開，務要和他一同卸下驢的重馱。 \end{tabularx} \\ \\ \relax
23:6 & \begin{tabularx}{0.7\textwidth}{X} 「不可在貧窮人的訴訟中屈枉正直。 \end{tabularx} \\ \\ \relax
23:7 & \begin{tabularx}{0.7\textwidth}{X} 當遠離誣告的事。不可殺害無辜和義人，因我必不以惡人為義。 \end{tabularx} \\ \\ \relax
23:8 & \begin{tabularx}{0.7\textwidth}{X} 不可接受賄賂，因為賄賂能使明眼人變瞎，又能曲解義人的證詞。 \end{tabularx} \\ \\ \relax
23:9 & \begin{tabularx}{0.7\textwidth}{X} 「不可欺壓寄居的，因為你們在埃及地作過寄居的，知道寄居者的心情。」 \end{tabularx} \\ \\ \relax
23:10 & \begin{tabularx}{0.7\textwidth}{X} 「六年你要耕種田地，收集地的出產。 \end{tabularx} \\ \\ \relax
23:11 & \begin{tabularx}{0.7\textwidth}{X} 只是第七年你要讓地歇息，不耕不種，使你百姓中的貧窮人有吃的；他們吃剩的，野獸可以吃。你的葡萄園和橄欖園也要照樣辦理。 \end{tabularx} \\ \\ \relax
23:12 & \begin{tabularx}{0.7\textwidth}{X} 「六日你要做工，第七日要安息，使牛、驢可以歇息，也讓你使女的兒子和寄居的可以恢復精力。 \end{tabularx} \\ \\ \relax
23:13 & \begin{tabularx}{0.7\textwidth}{X} 「凡我對你們說的話，你們都要謹守。別神的名，你不可提，也不可用口說給人聽。」 \end{tabularx} \\ \\
[1ex]
\hline
\hline
\end{longtable}
$^{1}$各位電視節目早安.
我代表宣導的北角堂向大家問安.
多謝教會邀請我.
亦都今日可以來到監獄跟大家分享訊息.
在近年我們或者.
未必會再直接提很多公義公正的事情.
但誠系體重一個社會.
一個社群裡面有沒有公義和公正.
這一樣的事情不會因為時局.
或者社會的文化而有改變.
在詩篇89篇14節這樣說到.
公義和公正是你寶座的基礎.
神的寶座的基礎即是那一樣東西是不會動的.
在那裡的.
不會再改變的.
慈愛和信實走在你前面.
這一樣東西一直在神的寶座流出來.
一直在發出來.
一直在走,不會停的.
當我們提到公義和公正.
這一對同義詞.
中文都比較難去分辨.
大氣一百來文.
你聽讀音都知道是兩個字.
公義是siddakā.
公正是mishpat.
你記得一個是時音.
一個是悟音.
兩個不同的音.
有分別的,用來形容兩種我們會遇到的情景.
我神學院的老師就曾經用過這樣的例子去作出講解.
我覺得很好,所以我複述用他類似的方式去講.
他說公義siddakā.
在這樣的公義的社會裡面是甚麼呢?.
就好像我們要出街工作.
上學了,坐車.
拿著個袋,背包上車.
上車後,突然收到電話.
聽電話,一邊聽一邊聽.
到站下車了,就下車了.

$^{41}$一邊講一邊下車.
怎料下車後才發覺,忘了拿著個袋.
袋放在椅子上.
那怎麼辦呢?你當然是去追巴士.
追不到,就要去總站.
然後去找車長.
車長說你去那裡,去新聞營領處.
你去到新聞營領處那裡.
一講,是這個袋嗎?是呀.
打開來看看,裡面的錢包.
你的電腦,全部都在這裡.
拿回你的東西.
好像沒有發生過.
撿到的人.
路不拾遺.
這就是siddakā.
公義的社會.
在公義的社會裡面,每個人都做回他應該做的事.
不是他不拿的東西.
這就是公義.
公義又好像甚麼呢?.
最近有講到有公展會在維園.
不知道你們那麼遠會不會特意出去了.
不用了,你們有賣物會.
一會兒完了就去賣物會.
公展會很多人.
因為疫情後.
大家不是每個人都能夠立即飛到日本.
去公展會買東西都好.
買一元鮑魚.
要排隊入場.
排隊.
對呀,跟規矩.
有秩序.
當公展會一開了.
大家就按著秩序.
輪流入場.
這就是公義.
siddakā.
在這個秩序整形的社會裡面.

$^{81}$每個人都做回他應該做的事.
要入場就排隊.
應該的.
如果我們的社會社群天天都是這樣,多好.
每個人.
自動自覺做回應該做的事.
路不拾遺,有秩序.
但我們都知道.
逢有人的地方.
就有問題.
排隊很簡單.
但有些人不守規矩.
要打GYM.
因為要搶一元鮑魚.
入先入後是有分別的.
拿著車子衝衝衝.
看到有少少隊友.
罅位就滑了進來.
這個時候被他打了GYM後面.
那個人就說排隊.
這裡有位置我不知道.
是可以了,隨便吧.
兩個就吵了.
這個時候.
有爭議.
不是,你要排隊.
大家就口角.
附近的.
本來正在入場的.
工作人員就來問.
什麼事,幫他們解決.
誰先誰後.
他們兩個當然說.
是我先我先.
就問前面,有沒有看到.
看不到,後面幾個,是他.
他打GYM.
拖著車子的阿善.
就說麻煩你排後面.
太太,跟我來吧.

$^{121}$都買到的,放心吧.
這個.
就是公正.
公正是什麼呢.
就是.
公義其實有秩序.
有人做應該做的事.
當這個正常的狀況.
被擾亂了.
就要引入公正.
要解決這個問題.
要學原則和方法.
這個就是公正.
我陳詠英老師的講解.
其實會幫助到我們分.
什麼是公義,什麼是公正.
公義就是一般.
我們希望運作正常的社群.
公正就是.
正常運作裡擾亂了.
有爭議,有困難.
怎樣去解決.
這個就是公正的方式去處理.
其實.
有人的地方.
一個社會.
一個社群,一個公司.
一個教會,一個團契.
一個小組,一個家庭.
都會有需要公義和公正.
這兩方面在我們生命裡.
有這個定義.
我們接著就可以看經文.
但進入經文之前.
我想再和你一起祈禱.
求神幫助我們.
天上上帝願你的真理.
讓我們眼睛看得見.
讓我們靈裡都去.
明白過來.

$^{161}$以致我們作為你的兒女.
在日常生活裡.
能夠活出你這份公義和良善.
幫助我們聽到的時間.
祈禱奉主耶穌基督.
得勝名字祈求,阿們.
《出一及記》的23章1至13節.
剛才領事也幫我們讀了.
其實有兩幅圖畫.
因為你現在.
都不用聖經了,我告訴你.
這13節其實有兩幅畫面.
你應該可以想像到.
第一幅就是.
一個城門口.
等於.
AKA他們的.
法庭.
他們在城門口做法庭.
人人都經過,人人都去圍觀.
長老坐在城門口,其實就是法庭.
一個這樣的場景.
另一幅圖畫就是路邊.
田邊,日常生活的場景.
這兩幅圖畫在這13節.
經文裡是交錯出現的.
1至3節是講法庭.
4至5節是講日常.
然後又回到6至9節.
又是法庭,10至13節又是講日常.
這兩幅圖畫交錯出現.
不過有一個共通點.
就是記錄了.
我們在生活裡.
很多處境是有爭議的.
一群人一起生活.
是會有.
大家不同意見,甚至爭執.
打鬥.
是會有這樣的事出現的.

$^{201}$而爭執的對象.
或者不同人.
是會因為你的權力的分配.
你的財富的不同.
你的身份的不同.
這裡講到窮人.
仇敵,寄居的人.
有分別.
我們可以一次過.
先講第一個場景.
不如我請弟兄幫我讀這三節.
好嗎?預備起.
不可隨火.
布散謠言.
不可與惡人.
聯手妄作見證.
不可隨眾行惡.
不可在爭眾的事上.
隨眾偏行.
作見證屈枉正直.
也不可在爭眾的事上.
偏護窮人.
然後我們跳去第六至九節.
我請姐妹讀.
預備起.
不可.
這兩段.
有四段.
第一第三段.
是剛才說法庭的圖畫.
先給我們看到一個原則.
我先說原則.
在爭議裡面持守公正.
爭議即是公義被擾亂.
要公正.
就不會因為多數人而扭曲事實.
這兩組經文.
是法庭裡面有人說的.
對法庭裡面有人說的.
出現在法庭裡面的人要留心.

$^{241}$甚麼人呢?.
三個可能性在法庭裡面出現.
第一當然是被告和原告.
第二類就是證人.
不關事的.
不過可以看到.
第三類就是審判官.
主要是這三個裡面.
有人會在法庭出現.
或是判斷.
對錯的場面出現.
而這段經文主要針對後面兩組人.
證人.
審判官.
聖經有甚麼要教導這兩組人呢?.
第一至第三節是有關證人的教導.
我先回到經文.
開首的第一項.
不可除火.
布散謠言.
其實原文的意思是.
直接這樣翻譯的.
不要拿起虛假的東西.
你見到虛假的東西.
你不要拿起.
不要將你聽回來.
隨手拿來的消息.
不知道是不是真的.
未經證實的.
甚至是一些說笑的.
或者喜笑怒罵的東西.
就像是一支雜琢一樣當真相散出去.
因為這些說話很可能.
轉了幾次之後.
會變成對一些人的指控.
人家就當是事實.
這樣說的.
其實現代心理學都告訴我們.
我們人的記憶很有趣.
可以局部改變.

$^{281}$所以為甚麼.
很多有權力的人.
很希望透過一些教育.
去改變我們的記憶.
當時的記憶.
是可以受到虛構的事情影響.
有一項研究.
研究員就將一張.
印有賓尼套的迪士尼廣告.
給參與測試的人看.
用了很多的形容詞.
很多的說話不停地說.
這個有賓尼套的迪士尼廣告是怎樣的.
有甚麼特色 有甚麼特點.
但心水清楚知道.
這廣告肯定是錯的.
為甚麼.
迪士尼沒有賓尼套.
你想想.
迪士尼沒有賓尼套.
賓尼套是華納兄弟.
我看你們的樣子.
不算太小 應該還能分辨到華納兄弟和迪士尼.
OK.
賓尼套那一套.
那一堆東西是華納兄弟.
即使是.
都頗明顯的.
錯誤.
一個錯誤.
經過研究員不斷和測試對象.
去講解 去讓他看那幅圖畫.
竟然在測試裡面.
有高達16\%的人.
以為自己真的.
對呀 我小時候去迪士尼.
真的和賓尼套拍照.
言之鑿鑿.
他記憶中說.
真的有 還能形容.

$^{321}$怎樣和賓尼套抱著.
這個假的資料.
扭曲了他真實的記憶.
和印象.
心理學研究出來.
很科學化.
你看看 連這麼小的事都影響人.
所以假的說話.
假的傳聞.
一些不是真的東西.
一句都不要傳.
說著說著.
就會有人當真.
甚至轉過頭來.
是會影響到別人.
害到別人.
我們繼續看.
作為證人不單止不可以給.
假的證供.
第一節 下半.
第二個不可.
那裡說甚麼.
不可與惡人連手.
妄作見證.
新譯本.
是這樣譯的 與惡人攜手.
作假見證 惡人原本的文字.
是解作一個罪有應得的人.
即是他真的有罪.
你不可以和他.
站在同一陣線.
不可以為了幫他脫罪.
而作出假的陳述 假的口供.
當然一定不可以和他一起去害人.
但我想我們都不會.
幫壞人的 即明知道.
是壞人你不會幫他.
但如果在這裡.
一起做這件不好的事.
是一大班人.

$^{361}$整班人.
都這樣做.
你怎樣.
譬如在公司裡面.
每個同事都喜歡.
照顧某一個人.
每個都喜歡.
指他背脊.
每個都喜歡在老闆面前.
有意無意踩低他.
都覺得.
有他都不要緊 最好他離開.
不要在我們這間公司.
你身為他的同事.
你能不能公公道道地.
幫他講一句說話.
都不容易 不要講公道說話.
你不參與一份.
有人在Facebook鬧爆.
惡搞.
取笑.
如果那個流傳的東西.
不是建基於一個事實.
是有些東西很活該.
很絕核 很搞笑.
但你能不能公公道道地.
不分享 不讚好呢.
甚至在某些明顯錯誤的資訊裡面.
願意作出中肯的評論呢.
你感覺到.
不是很容易.
我們的選擇可能不是這樣選擇.
因為人人都覺得沒什麼.
聖經在第二節說.
你不可隨眾行惡.
這條題我們不可以單單.
有多少人一起做這件事.
都覺得沒所謂而做.
扭曲了事實.
影響了判斷.

$^{401}$上帝說不可以 這是行惡.
你想想為什麼上帝.
那時候要和以色列人說這句話.
這是出埃及的.
出埃及那班是什麼人.
奴隸.
沒錯 是奴隸.
每個都是弱勢社群.
每個都是窮人.
在埃及做奴隸.
當然出來的時候.
又問了鄰居拿了一些金.
最後那晚.
但基本上他們的本質.
或者他們一直過去的生活.
他們對自己的認知都是窮人.
弱勢 需要人幫的.
聖經說 神的說話.
即使是這樣.
你作為我的子民.
你不可以不分是非黑白.
幫助弱勢的人.
更加不可以單從.
有多少財富去判斷一個人.
是不是罪有應得.
現在我們常常說仇富.
有人會問 幫人都不行嗎.
鋤強扶弱.
明知他是弱勢 當然站在他那邊.
一定沒有錯的.
聖經怎麼說呢.
在這裡我們必須回到十誡裡.
和這段最有密切關係的第九誡.
一起讀這節 第九誡.
預備起.
不可作假證供.
陷害你的鄰舍.
其實原文裡沒有陷害這個字.
不可作假證供.
指證你的鄰舍.

$^{441}$不一定是陷害的 比較中立的.
即是怎樣 你有沒有想過.
假證供當然是.
好像想陷害 害人家.
但其實有人用一些.
不是正確的證供.
去幫人.
神怎麼看這件事.
在一個判斷的場景裡.
一班人要處理的爭執.
訴訟裡.
就算有一邊是弱勢.
你都不可以不擇手段的.
鋤強扶弱.
這不是聖經的價值觀.
我們都記得.
在爭議裡.
要持守公正.
就是不因多數人.
去扭曲真相.
好 跳到這裡.
第六至第九節.
文字裡的第三段.
也是說第一個的圖畫.
法庭的圖畫 第二組人.
說了證人 說審判官.
這裡有五項條文說不可.
不可在窮人.
爭重的事上屈往正直.
當遠離虛假的事.
不可殺無辜和有意的人.
不可受賄賂.
不可欺壓忌懼的.
聖經很清楚 作為有權.
缺權的人.
要小心一切虛假的指控.
不要佔立假口供 不要殺錯無辜.
不要將沒有錯事的人.
判罰 更加不要收受利益.
以致顛倒是非.

$^{481}$你說是啦 是啦 做法官.
法庭應該是這樣.
這裡也提到.
判斷的人.
不一定只是做官 你知道以前在城門口.
那些其實都是跟你一起.
生活的隔離左右.
只是他是長老 領袖.
幫忙去解決.
社群裡面的紛爭.
換言之就是說.
可以應用在我們生活裡面的場景.
可以是你在公司裡面.
你是老闆 你是主管.
你是team head.
有兩個同事來跟你說.
有些事你要去判斷.
誰對誰錯.
可以是你在學校裡面做老師的兩個學生.
幾個學生來跟你說剛才發生的事.
你要做老師幫他們去判斷.
可以是各位.
在家裡面做家長.
兩個小朋友吵架.
你幫他們說誰對誰錯.
每一個人都總會有.
一些位置 你有一個判斷權.
在那時候.
你要行使這個判斷權.
你就可以用這些原則.
不要屈窮人 這樣明白嗎.
不要聽假的東西.
不要害錯人.
不要因為你偏愛其中一個小朋友.
平時他.
孝順你一些.
受賄賂.
這裡有很多原則.
對我們依然應用到.
更加當然說大一點的環境.

$^{521}$不要受賄賂 廉潔.
這不是香港的專利.
聖經早就說了.
但這裡最明顯最顯眼的.
是我們還沒有讀到的第九節.
一起讀吧.
不可欺壓既居的.
因為你們在埃及地.
作過既居的.
知道既居的心.
這句是比有權病人.
更加要留心.
不要因為親疏有別就欺壓外族人.
少數族裔要公公道道作出判斷.
其實這些仍然是我們今天社會.
要追求的目標 成長的方向.
要達至不欺壓既居的.
在我們香港這個社會仍然有一大段路要走.
差不多三十年前.
香港開始興建新機場.
現在已經夠了.
現在還在說三跑.
當時為了要建機場.
輸入了很多外勞.
大量的南亞裔來香港做勞工.
在建築方面去參與.
到了這個機場.
真的落成了 赤駱角.
搞定了.
不少南亞裔已經在香港過了七年.
足夠時間.
可以有居留權.
所以為什麼在這十多二十年.
我們見到很多南亞裔.
在我們的社群裡面一起生活.
曾經我讀過一個新聞.
我照讀它的標題.
它說幾萬元.
不貪心.
南亞學生路不拾遺.

$^{561}$發生什麼事呢.
原來有一個南亞學生撿到一個陳先生.
掉了的錢包.
這個錢包裡面有三萬元現金.
這個南亞學生撿到之後.
就交了去警局.
結果警局聯絡失主陳先生.
陳先生很緊張.
去到警局拿回自己的錢包.
當然要查一下有沒有撿到錢.
三萬元都在裡面.
於是訪問他的人就問他.
你覺得怎樣.
陳先生就回答這個記者說.
想不到是被這間學校的學生撿到.
還要是南亞裔的學生.
我對南亞人改觀了.
這個報道.
給我們看到什麼叫壓迫寄居.
你想想.
壓迫不一定是我們去奴役他們.
讓他們做苦工.
不一定是我們敲詐他們.
壓迫可以是.
對他們一種負面的觀感.
以至認為這些族裔的人.
一定是貪心.
窮一定是會偷東西.
不窮都會出古惑.
所以上到法庭.
一個大環境.
大家有這樣的觀感.
就算是審判官都容易被這些看法.
所影響.
這是另一種.
欺壓.
其實我換一個字眼.
告訴你同一個標題.
你就明白.
如果有一天我們在英國讀到.

$^{601}$一段這樣的標題.
幾萬元不貪心.
香港學生.
路不拾遺.
你大概知道這個社會.
怎樣看待香港人.
在英國.
所以.
我們有沒有.
壓迫我們身旁的.
少數族裔呢.
我知道這裡有.
外傭的事工.
教會意外呢.
你怎樣看.
香港人.
幫你送.
Foodpanda.
都是香港人.
不過膚色跟你不同的人.
剛才看到.
貫穿兩幅圖畫的.
都是一些人物.
窮人.
跟你不對的人.
恨你的人.
寄居的人.
這些人物.
跟我們有距離的.
爭取一些東西.
這些人性已經說.
你特別要留意.
如果你是有判斷權力的人.
你不要被你一些.
社會地位你的情感.
你的一些種族家國的.
差距.
政治的立場.
而阻礙到一個公正的判決.
更加不要扭曲.

$^{641}$一個事實.
以致一些跟你很不同的人.
得不到公正的判斷.
反而被.
無形欺壓.
我們都記得.
在爭議中要持守公正.
是什麼呢.
就是不會因為對方是外人.
跟你不同不是你圈的人.
不是你這一班人.
而屈往正直.
這是第一幅.
法庭的圖畫.
跟著跟你看.
這裡經文的第二幅圖畫.
是不同的圖畫.
請你幫我讀第四第五節.
預備起.
若遇見你.
若看見.
恨你人的奴.
厄我在重駝之下.
不可走開.
勿要和奴主一同抬開重駝.
第四節.
這裡說你仇敵.
好像很激.
你氣的人.
換句話說.
若遇見你氣的人.
第五節反過來說.
恨你的人.
即是氣你的人.
無論你氣他也好.
他氣你也好.
不理什麼原因.
大家不妥大家.
如果你看見.
他的生畜出了問題.

$^{681}$你會怎麼做.
這裡說牛,驢.
這兩樣東西是重要的.
因為在他們的農業社會裡.
牛是耕田.
也很偶然會用來做食物.
其實他們不是常常吃的.
因為很貴重.
用來獻祭.
一生人可能獻了多少錢牛.
很重要的東西.
貴重的.
驢主要用作運輸.
運貨,運人.
代表著什麼.
他要賺錢的工具.
他想像是一個貨櫃車.
一輛貨車,一輛的士.
有個弟兄是的士司機.
他在路上見到一個行家.
也是的士司機.
車子不知何故停在路邊.
見到那輛車.
打開車頭蓋.
他就明白了.
當然是拋錨有問題.
於是他沒有再嚇他.
二話不說就扭tie .
停在的士的前面.
問他有沒有什麼可以幫忙.
要呀要呀.
他就立即拿出電纜.
搭橋幫他托著車.
這個弟兄會出手相助.
不會袖手旁觀的.
大家行家.
但如果換成一個假如的處境.
如果現在拋錨的.
這個的士司機.
剛剛和他有過節.

$^{721}$大家吃飯的時候.
靠在一邊.
爭位.
有時候也有很多人.
又或者.
之前因為爭客.
發生過碰撞.
和你不合.
生你的氣的人.
上次碰到他.
只是一點點.
又要賠錢.
每次見到我都在說我.
爆粗地罵我.
你看到他現在拋錨.
你停不停下來.
我不是開的士.
所以不知道.
好呀換個公司場景.
你的同事.
預定了.
幾分鐘之後.
就要去會議室和大老闆.
有個簡報.
他突然在.
位置大叫.
我不見了我的USB手指.
糟了我怎樣去會議室.
怎麼辦.
剛好你聽到他大叫的時候.
你看到自己的電腦插著.
自己的USB手指.
你借不借給他.
大家同事給他用一下.
但如果昨天你才偷聽到他在茶水間.
說你壞話.
他之前.
試過在老闆面前.
毒我.
你會不會借你的USB手指.

$^{761}$給他.
你不要害他.
你幫不幫他.
會不會.
你很小心慢動作.
將你桌面上那一疊文件.
推到電腦上.
遮住你的USB手指.
他到處去找看誰有.
可能到他搞定.
去到老遠幾個位置以外.
找到同事終於拿了手指已經過了鐘.
嘩.
看看老闆怎樣說他.
最好老闆一氣之下解僱他.
這樣說我.
你沒有害他.
對於你不喜歡的人.
你會不會有這個試探.
幸災樂禍.
甚至主動希望有災禍淋到他身上.
不是你去害他.
是上天給他的.
上帝給他的.
神在經文裡.
吩咐我們不要袖手旁觀.
不要看著他沒有了.
他的穩食工具.
他所種的財物.
而導致生活有結局.
困難生活不保.
你持守你的良善.
對你生氣的人.
你去救急扶危.
行不行.
為什麼經文講完法庭的圖畫.
又拿出這幅田邊路邊牛路的日常圖畫呢.
以前有學者認為.
這樣交錯擺位是排錯了.
但我想提出.

$^{801}$經文的鋪排用意很明顯.
其實是要你和我去理解.
當我們常常想著.
很極端的法庭.
去告人.
極端的裡面當然我們要求有公正.
但這個公正離不開.
我們日常的處境.
如果在無風無浪的日常生活裡.
面對生氣你的人.
你生氣的人有難.
我們都不能夠活出.
那一刻的良善.
幾乎可以肯定去到極端處境.
大家在爭議.
在告.
你很難.
把持到公正.
有意無意間.
你所表達的.
會有偏差.
接著我們再看最後一段經文.
是10至13節.
我們只請你讀一節.
12節.
6日你要做公.
第7日要安息.
使牛奴可以歇息.
並使你疲累的兒子.
和寄居的都可以舒暢.
在10至12節這三節經文裡.
主要提到一個我們恆常操練的問題.
就是能不能夠放下.
你自己一些基本的利益.
而信靠神.
剛才說.
法庭那些很極端.
就算走失牛羊.
賺錢有問題.
拋錨.

$^{841}$不會每天都這樣.
一年都遇不到一次.
你要表達你的愛.
不是極端.
不是偶然出現的狀況.
反而要在你日常裡.
什麼日常.
就是要延伸到和你距離很近的一班人.
你有沒有表達到這份良善.
哪些人呢.
就是窮人,疲累,寄居.
經文拉闊我們的視野.
當我們常常想著.
在爭議當中.
爭論當中.
在社會裡我們要有公正.
對仇敵要有良善.
要愛仇敵.
其實拉闊到日常.
你有沒有對你身邊的一些人.
作出美好的安排.
以致你不會不知不覺間.
好像剛才說的那個南亞人賺到錢.
去壓迫寄居者.
你看我們香港很多家庭都有這一類人.
幾乎合乎.
這節經文所說的三個角色.
窮人.
疲累,寄居.
哪些人呢?很快就想到.
姐姐.
沒錯就是你的家庭.
或者你公司裡面一些.
和你的膚色不同的人.
在這節目我們有沒有好好地.
讓你家庭的家庭.
從六天工作之後有安息.
有有有.
今天都是放假.
還是他們一齊.

$^{881}$由他們出門到回家.
他們真正休息都只有.
八九個小時,十個小時.
出去之前.
會不會還在幫你工作.
下班回來.
是不是還要洗這個洗那個.
你想想.
他們有沒有一天的休息.
換一個場景.
今天星期日.
可能昨天有些人可能還要上班.
或者星期五.
我們放兩天的更好.
老闆跟你說.
你接著放假了.
休息了.
週末.
休息沒有問題.
不過你休息那天早上.
你先回來清理東西.
然後放假.
然後晚上你八點結束.
你的放假.
我再發電郵給你.
幫公司處理好一些事.
一次趕工作,趕預劃.
OK.
但如果每個星期日都是這樣.
那一天.
是不是休息日.
當然不是.
搞錯了.
星期日就是星期日.
怎麼會叫我工作.
一天又按食.
所以各位Sir,Madam.
想不想.
一個星期.
自己洗一天碗.

$^{921}$行不行.
或者家傭放假回來.
你說明天才做.
碗有他放在那裡.
今天你放假.
東西亂一點點不要緊.
不要搞事,快點去洗澡休息睡覺.
你能不能有這樣的寬容.
或者不是寬容.
能不能做回一整天的休息.
上帝設立安息日.
其中一個目的就是為了你.
和我和其他人.
在工作之間有喘息.
歇息.
其實這個字.
就是在創世記那裡.
神吹一口氣.
進人那裡.
那個同一個字.
幫你回氣.
這一段進一步讓我們看見.
頂智慕我們說要信得過神.
以致我們恆常,直順,安息,洗人.
特別是一些弱勢的人,寄居的人.
疲累窮人有喘息,回氣.
在這種我們追崇,讚順,繼進的社群生活裡.
我們要操練我們日常的良善.
不用說得很遠.
在你隔鄰,在你家裡.
請你開始做.
不是到嚴重爭議才說我們要公正,公義.
若想在爭議裡.
持守公正.
就必須先在你的日常生活中.
活出良善.
上帝頒佈律法.
因為他知道我們的特性就是在這個秩序井然的社會.
一個涉及公義的社會裡.
不會自動自覺.

$^{961}$常常是這樣的.
有人的地方就會有困難,有問題,有爭議.
又不對.
所以我們需要公正.
一些準則,一些原則.
神給我們很多律例典章.
是從這個角度出發.
讓我們解決那時候的衝突,爭議.
用神的角度.
如何去回復一個社群的公義.
當然我們是要繼續關心社會.
無論教會,團契,小組,自發.
總要參與.
但今日經文給我們一個基本功.
你要在爭議裡持守公正.
請你在日常生活中活出良善.
不是去到極端法庭.
不是去到偶然見到別人.
別人有困難.
而是在日常生活裡.
每天你的生活裡的小事.
你能不能活出這份良善.
這一段是讓我們愛神,愛人,愛論社.
有一個基本的功力.
當有壓力,有爭執,有困難,有訴訟.
你在那些位置,日常裡.
有好的基本功.
你去到最終仍然能夠有這份公正和良善.
我們一起祈禱.
頂著如果你剛才聽到的時候.
你想起一些人和事.
我請你自己向上主禱告.
父啊,或許我們想起一些人.
想起一些我們正在面對的處境和困難.
在這些很難處理的處境裡.
求主讓我們能夠活出良善.
正如你在《初一及記》第十三節說.
上帝所吩咐的,你們都要遵守.
幫助我們能夠聽從,信服,活出良善.
多謝你.

$^{1001}$我們祈禱是奉主耶穌基督,得勝名至祈求.
阿們.
\newpage



\section{以斯帖記}
\label{sec:iRhu93YiYcU}
\textbf{2023-08-24 周四晚查經班《被擄歸回歷史研讀(22)以斯帖記》 講員:曾裔貴牧師}
\newline
\newline
連結: \href{https://youtube.com/watch?v=iRhu93YiYcU}{\texttt{https://youtube.com/watch?v=iRhu93YiYcU}} ~~~~ 語音日期: 2023-08-26
\newline
\newline
\hyperref[sec:lKgHJObwVL4]{\small{< < < PREV SERMON < < <}}
~
\hyperref[sec:index_chronic]{\small{[返順時目]}}
~
\hyperref[sec:index_scriptual]{\small{[返順卷目]}}
~
\hyperref[sec:fSxNb3R9eIk]{\small{> > > NEXT SERMON > > >}}
\newline
\newline
$^{1}$鄧英姐妹,我們可以恢復查經.
已經停了三次.
我們一起來禱告.
希望能夠慢慢勾起對二師帖記.
和當時歸回歷史的背景認識.
然後進入內文去探討作者.
我們不知道是誰.
他要帶出的當時背景.
和關乎民族的遭遇.
由張力引發的轉變.
神如何去解救民族.
這本書可以肯定是講關於.
普爾節的由來.
為什麼猶太人會多受一個節期.
轉捩點在哪裡.
我們可以一起去探討.
有個簡單祈禱先.
多謝天父我們再一次恢復.
系統地一起研讀聖經.
能夠從每一章探求內容.
我們很想去尋索.
到底神在歷史裡.
如何藉著一些危機.
來施行你的拯救.
更重要的是透過猶太人這個民族.
讓我們看到.
彌賽亞降臨之前.
與神立約的選民.
神的選民.
他們是如何忠於主.
他們如何能夠在極困難的時候.
來去持守他們的信念.
他們的信仰是活潑的.
幫助我們在現在.
面對21世紀.
很多很多不同的轉變.
我們都要提防.
哪些敵人.
好像哈曼一樣.
要來抵擋我們.

$^{41}$要來置我們於死地.
我們需要求主讓我們能夠有.
這種肅靈的眼光.
願神帶領幫助我們.
我們一起在你面前等候禱告.
奉耶穌基督的名.
阿們.
看這本書的其中.
峰迴路轉的地方.
作者很巧妙地.
將每一件事.
鋪排給我們知道.
慢慢地.
一層一層地組織.
交疊的章歷.
要一次過.
改變.
要去到第六章.
突然之間.
告訴我們.
這位統管著.
127個省的.
阿黑徐老王.
當時是一個波斯很強大的國度.
不是一個墮落到.
差不多要亡國的.
一個皇朝.
突然提到.
他一直.
很有氣勢.
很有魄力.
但突然之間.
第六章第一節就說.
那一晚皇帝睡不著.
沒什麼可以做.
就叫大臣.
將一些歷史.
讀給他聽.
希望讀到下午晚睡.
突然讀到一段.

$^{81}$第二章我們會說的經文.
所以由這裡開始.
是整個《以司帖記》的轉變.
但是去到第二.
第三章.
就是哈曼和.
這個末底改.
那個結怨.
大家同朝為臣.
但因為哈曼被抬舉.
這個末底改.
就不向他跪拜.
招惹了.
哈曼的怒氣.
你可以想像到.
當時得令到一個地步.
就是.
我殺你一個末底改.
我都覺得是少.
我要將你.
和你的民族.
都要借助皇的勢力.
來徹底瓦解.
希望求皇.
立一個法令.
讓這個民族.
猶太民族.
要消失.
在波斯帝國.
所以你可以說.
他的仇恨.
其實由他自己內心的.
自卑感.
作祟而導致.
整個民族.
要被毀滅.
可想而知.
那個情況都沒有立刻化解.
一直要等待.
這個末底改.

$^{121}$就發現.
事情越來越嚴重.
他就要冒險.
去找他的養女.
其實後來.
就是他叔叔.
他是義斯帖的叔叔.
來希望解救.
然後就說.
宮廷的禮儀.
那些手續.
一直拖.
拖了很久.
拖到根本上.
都不知道怎麼辦.
唯有等神蹟.
那個神蹟居然.
是一個皇帝睡不著.
而改變了整個局面.
你就可以想像到.
這本書的.
峰迴路轉都是作者.
一定是猶太人.
一定是拉比.
一定是他們對這件事.
在他們猶太人的.
角度去寫.
所以他就.
現在一直在營造.
這個民族其實.
岌岌可危.
怎樣能夠得到解救呢.
還有神是不是真的存在.
繼續去愛.
和他立約的.
這個民族呢.
神用什麼方法去.
解救呢.
我們又見不到一點的神蹟.
但是見到人的參與.

$^{161}$見到末底改.
仍然都是很堅信神會解救.
還有義斯帖.
無論他現在.
方不方便也好.
他也冒死去見王.
很多的時間點.
如果不是神安排.
其實是可以很大件事.
第二章開始.
第一章就是交代了.
為什麼義斯帖.
這位皇后.
會成為了.
阿赫徐老王的皇后.
就是要交代這件事.
第一章就是因為.
阿赫徐老王的太太.
瓦實緹.
因為皇后.
喝醉了之後.
很想將她的太太.
皇后展示給.
當時的飲宴.
的官員.
可能是外交的使節.
來看.
戴上皇冠.
穿上最漂亮的.
戴上珠寶.
如果你是一個.
很炫耀的.
炫富的.
你就求之不得.
可以在眾多的外國的使節.
所有的官員.
和官太.
名流.
上流社會的名媛.
來做展示.

$^{201}$你是戴最多的鑽石.
全部.
你家裡的.
首飾箱裡的東西.
都拿出來了.
但是居然.
這個瓦實緹.
一下子下不了台的那口氣.
我就是不出.
皇后面對這樣的.
當眾羞辱.
就發現一件成為了很嚴重的事.
本來就是.
喝醉了之後.
可能失言.
衝口而出.
希望炫耀一下.
自己的國力.
從太太的.
裝扮.
去做一個基本的展示.
但是.
帶來的就是一個.
很尷尬的場面.
那些在皇帝身邊.
的七個大臣.
就是那些謀士.
出謀獻策.
關乎國體的.
就說不行.
他一個人不做.
他的身份是皇后.
127個省的.
其他官員的官太.
有樣學樣的時候.
所有的官員.
可能是男系的社會.
沒有地位了.
自己的太太.
都照樣煮碗.

$^{241}$那時候沒有直播.
又不是直播帶貨.
又不是有zoom.
又不是有什麼影像.
其實要傳都要快馬.
那些傳書.
但是都會令到.
後果.
這樣的一個.
不聽命令.
會帶來.
真是有辱國體.
其實關以斯帖什麼事呢.
但是你就可以.
很見到人生.
一直都說.
當下你不知道.
神有什麼計劃.
在你的生命裡面.
神有什麼計劃.
我們都不知道.
其實好看的地方.
就是他將一個人.
發生的事.
有些事情.
未去到最後.
其實都不知道.
神怎樣去解救.
神怎樣去出手.
神怎樣去使到.
一些我們不想發生的事情.
是會離我們而去.
這個是一個.
立約的子民.
堅持守神.
的誡命的猶太人.
他們面對這樣的一個危機.
是怎樣算呢.
神是否當下.
馬上出手.

$^{281}$來叫王.
或者叫其他人.
來對付哈曼呢.
又不是反而我們見到那個情況.
是越來越一直這樣嚴重.
一直去到第五章.
不單止不能解救.
甚至乎立了一個禁令.
去到十二月.
所有在一百二十七省的.
波斯人都可以.
向猶太人.
逼迫殺害.
所以.
後來才改了.
這個普爾節.
就是當日那個禁令.
改不了的了.
但是猶太人就可以起來反抗.
你就可以想到.
神那個奇妙的作為.
第二章我們讀一讀.
這一段經文.
以後.
因為開始的時候講了.
不能回頭了.
(陳健波議員).
合神,對不起,合王的眼緣.
當然要過五關斬六將.
首先你就要買通這個熹柑.
如果你的熹柑他認為不行.
那麼你就被淘汰了.
那你就永遠都見不到王的了.
這個大家都明白了.
有很多事情不會直接去到王那裡的.
所以這個太監熹柑一定是有手握大權的了.
就是你中謠好聲音.
有些評判團.
不過這麼多的傳媒報道.
很難說有幕後的黑幕或者是賄賂.

$^{321}$很難的.
但是古代就不是了.
一定是會讓這個熹柑有極大的權.
可以來去操縱那個選舉.
第五節了.
蘇珊成有一個猶大人.
或者叫猶大人.
猶大之派.
名叫沒抵改.
是汴雅門人.
姬氏的曾孫.
侍妹的孫子.
艾爾的兒子.
記住他是汴雅門人.
就是這個掃羅王的支派.
其實很近.
都是那個族的人.
還有那些叔伯兄弟.
都挺像是掃羅那條線的普系的.
他是艾爾的兒子.
從前巴比倫王尼布格尼撒.
將猶大王耶哥尼雅和百姓.
從耶路撒冷撈去的時候.
是第一批被撈的人.
就是耶哥尼雅的時候了.
沒抵改也在其內.
沒抵改撫養他叔叔的女兒哈大沙.
即是說那個時間不是太久.
現在還是被撈的期間.
即是說這個是第一代被撈的人.
所以是早過尼熙米很多.
我們要記得這些歷史.
第七節.
沒抵改撫養他叔叔的女兒哈大沙.
後來就叫做以斯帖.
當然就是猶大的名和波斯的名.
因為他沒有父母.
即是他的弟弟可能早死.
就收養他為女兒.
這女子又容貌俊美.

$^{361}$她父母死了.
沒抵改就收她為自己的女兒.
即是本來是叔叔.
現在變成了女兒.
即是爸爸.
王的御子傳出.
就招聚許多女子.
去到書山城.
交給掌管女子的希蓋.
以斯帖也送入王宮.
交付希蓋.
希蓋喜悅以斯帖.
就恩待她.
給王給她需用的香品.
和她所當得的份.
又派所當得的七個宮女服侍她.
使她和她的宮女.
搬入女院上好的房屋.
以斯帖未曾將直觀.
中俗告訴人.
因為沒抵改祝福她.
不可叫人知道因為亡國.
也都因為可能當時.
對猶太人未必那麼友善.
所以暫時先不說.
這個先不說.
原來都有很重要的背景因素.
因為不說就帶來哈曼.
不知道以斯帖的身份.
因為第三章就是.
哈曼知道沒抵改是猶太人之後.
要將整個民族都要完全滅絕.
然後用一些手段.
令到王都簽了議協書.
就答應了哈曼由他做執行人.
就要去到12月.
將這個民族完全在波斯帝國鏟除.
但是唯一一個人.
他不知道的就是以斯帖.
是猶太人.

$^{401}$所以很奇妙地沒抵改.
叫以斯帖暫時不要說.
他的直貫他的宗族.
真的聽爸爸說就沒有說.
第十一節.
沒抵改天天在女院前邊行走.
要知道以斯帖平安不平安.
並後事如何.
很擔心女兒.
因為始終後宮很多的鬥爭.
什麼恩還傳.
我不知道我沒看.
你就知道背後幕後的鬥爭.
要爭做王后.
現在是選拔階段.
未定案.
未定案就是說我都有機會.
你都有機會.
現在我們又見到.
欺誣後代以斯帖.
可能純粹是從美貌.
又或者可能是人格人緣.
這些很難說.
但聖經現在鋪排給我們見到的.
這一個序士.
就是以斯帖.
現在可能有機會.
會擠生在王后的身份.
如果他是王后的身份.
那就是第四章.
大家很熟悉的經文.
焉知道你現在得了王后的位分.
正正就是為現在這個危機的出現嗎.
沒抵改就很有那種的mindset.
就是有一種危機意識.
一方面沒有強迫她入宮.
為了榮華富貴.
可能是全國要招募.
每家每戶18歲以上就要參選.
就要進去被召見.

$^{441}$這些就沒得選擇.
不是自己貪慕虛榮去做.
但是沒抵改又很有sense地.
女兒做了王后之後.
沒因為這樣而失了位分.
變了自己.
應該叫國寇諸如此類的.
這樣來到外府.
就會變成突然之間.
飛黃騰達之後.
得意忘形.
就忘記了自己猶大民族的身份.
但是都不講那一段了.
現在反而沒抵改.
自己招惹了.
其實不關她事的.
她招惹了哈曼.
就出現了一個滅族的危機.
十二節開始.
眾女子照例先潔淨身體12個月.
就是要一路來到去.
似乎可能將那些鄉音.
你的鄉下話.
你的不同的文化背景的東西.
都經過訓練.
你想想香港小姐的.
初選之後的復選准決賽.
都要去哪裡扎訓.
我又不知道.
我只是普通的新聞.
經過不知道多久的時間.
嚴格的篩選.
和你怎麼走那些戲劇.
怎麼來到去.
在鏡頭面前.
和你的談吐.
那些化妝.
一年都差不多了.
因為這麼大的帝國.
由不同的省來到去.

$^{481}$將這些女子送入女院.
六個月就用活藥油.
六個月就用香料和潔身之物.
換了日期.
然後哀次進去見阿合徐老王.
女子進去見王是這樣的.
從女院到王宮的時候.
凡他所要的都必給他.
晚上進去.
次日回到女子的第二院.
交給掌管妃嬪的太監沙甲.
這個是再低階一點.
除非王喜愛她.
再提名召她.
就不再進去見王.
或者可能是王身邊的太監.
第十五節.
末底改叔叔阿比孩的女兒.
這次又說她的弟弟的名字.
就是阿比孩的女兒.
阿比孩.
就是末底改收為自己女兒的議師帖.
按次序當進去見王的時候.
除了掌管女子的太監.
希該所派定給他的.
他別無所求.
就是說他有信心.
足夠的需用.
大家都是基本的香料.
他都不會要多.
不需要額外的化妝.
不需要有另外的一些秘密武器.
或者其他的東西.
他只需要基本的就已經足夠.
也讓我們看到他有自信心.
接著就是.
凡是看見議師帖的都喜悅他.
第十六節.
阿合徐老王第七年的十月.
就是提別月.

$^{521}$議師帖被引入恭見王.
王愛議師帖過於愛眾女.
他在王眼前蒙寵愛.
比眾處女更甚.
王就把王后的冠冕戴在他頭上.
立他為王后.
代替瓦實體.
王因議師帖的緣故.
吸眾首領和臣服.
設擺大言席.
又豁免各省的租稅.
並召王后議大頒賞賜.
就是說再一次令到這位阿合徐老王.
有回那份妃室.
有回那份歡欣.
令到全國都很興高采烈.
所以在這種情況下.
他又大頒賞賜.
給那些重要的內圍官員等等.
又令到國家.
現在選了王后.
不同的省就可以減稅.
都是一些福利.
有些國家甚至特赦.
第十九節.
第二次焦聚處女的時候.
沒抵改坐在朝門.
議師帖照著沒抵改所祝福的.
還沒有將直冠宗族告訴人.
因為議師帖遵沒抵改的命.
如撫養他的時候一樣.
當那時候沒抵改坐在朝門.
王的太監中有兩個守門的.
辟探和提獵名.
老衡阿合徐老王.
就是那些內侍臣.
禁軍那些.
就是那些御林軍.
就想要下手害他.
沒抵改知道了.

$^{561}$因為他是個閒人.
他是等待自己的女兒.
議師帖的消息.
所以不為意.
聽到這兩個在朝門的.
守門的人的密議.
一些秘密的.
一些陰謀.
所以他知道之後.
沒抵改知道了.
就告訴王后議師帖.
議師帖奉沒抵改的命.
將那個honor.
是給回議師帖.
對不起是給回沒抵改.
是沒抵改來提名.
不關我議師帖事.
是我的爸爸.
他發現了這個陰謀.
要行食皇帝.
所以接著就調查.
報告於王.
拆救者事果然事實.
就把義人掛在木頭上.
將這事在王面前.
寫於歷史上.
就停在這裡了.
將這一筆.
本來要即時賞賜的.
沒有.
就記錄在案.
記在歷史裡面.
似乎告訴我們.
當時他們很多事都很嚴謹.
這一次宮廷的.
差不多是暗殺事件.
其實處理了就.
不需要公告天下.
但是王就希望.
記在歷史裡面.

$^{601}$即是第六章.
王睡不著覺.
要起身看歷史.
看歷史每天都有記載.
為什麼突然間會揭露那一段.
就是這一句.
就是活底改.
這個發言.
而未有賞賜.
那你就可以想象到.
原來.
神的工作.
我們都不知道.
我們正在做的事.
原來將來是神.
來改變.
那個事情的.
那個的種子.
我們當下做的.
活底改當下發言.
這個的陰謀.
就挺身而出.
告訴以斯帖聽.
以斯帖就.
因為事態嚴重.
就告訴.
這個皇帝知道.
但是要具名.
即不是好像廉政公署那些匿名投訴.
是具名的.
是要奉活底改之名.
來報告這件事.
發言果然屬實.
接著皇帝就要求.
要寫在他面前.
寫在歷史裡面.
為什麼會煞有介事.
為什麼要這樣拖著尾巴呢.
現在不是寫小說.
現在是真的人生的一個發展.

$^{641}$一個時間的進行.
真是一件歷史的事.
但是當時就不是歷史.
是正在發生.
所以這本書.
值得我們看完之後.
要問的就是.
其實我們每一秒鐘.
都是在製造歷史.
一定是.
其實全個地球的人都是.
前一秒鐘之前.
就是歷史.
但是當下你參與的.
決策的每一件事.
必須要很.
心靈的單純.
依靠神.
因為我們不知道.
神怎麼用我們做過的決定.
來行事.
成卷書就是告訴我們.
就是這樣.
我們現在看當然容易.
當時沒抵改.
這樣幫了皇帝.
又沒有論功行賞.
接著第三章.
那個哈曼還要針對我.
和我的民族.
很不公平的.
很冤屈的.
第三第四.
第五章都還沒有出路.
還要是他們迫切的.
禁食祈禱.
不停的去呼天搶地.
但是到了第六章.
就是那個轉捩點.
就是皇帝記起這一節的.

$^{681}$歷史的記載.
所以我們讀聖經.
和我們在人生.
特別是因為一定是基督徒.
如果不是基督徒.
他們就不會想.
他們當下.
只會將來是遺憾.
或者是有眼光.
或者不枉此生.
沒有信仰的人.
但我們有信仰就不同了.
當下我們每一個決定.
大家有沒有印象.
羅馬書八章的二十八節.
正正就是這個意思.
我們曉得萬事都互相效力.
叫愛神的人得益處.
就是他所召的人.
他就會將我們帶來.
所召來的人.
就叫我們來得到榮耀.
當然還有就是兒子的身份等等.
所以你就會看到.
造物主有掌管萬有的能力.
但他往往都不是當下.
每一秒鐘立刻報答.
報應.
不是這樣的.
有時過了可能好幾十年.
有時可能是要很長的時間.
我們才會見到轉機.
所以這個當下是很緊要.
不是一般人所說的活在當下.
而是這一卷書是說.
這個民族.
現在他們都是無可奈何地.
因為被老.
被老就生活在波斯帝國.
是要等待七十年.

$^{721}$那個被老歸回.
但現在他們又不是黎希米的年代.
他們要歸回回去重建.
有神的拯救事實根據.
現在是一路等待那個復興.
被老的刑罰過了.
就可以回歸.
但現在在回歸之前.
這段時間.
他們如何認定自己仍然猶太人身份.
還有含金量呢.
含金量的意思是.
神還看管我們這個民族呢.
還是繼續憤怒地發怒氣.
來拒絕我們這個民族呢.
因為很快就是滅族的危機.
是由末弟自己引發的.
他引發的原因.
就是第三章說的.
他不跪不拜.
以猶太人的身份.
就是說.
我根據神要求我們這樣做.
現在得不到神給我們的恩典和出路.
我曾經幫過這個阿黑除老王.
現在他恩將仇報.
還要支持哈曼這個惡計.
那你就可以想象到原來.
在末弟的內心.
到底阿黑除老是一個昏君.
還是喜歡選后妃.
和縱容哈曼.
會不會第六章之前.
末弟是什麼樣的內心呢.
又是值得我們去重置旅行間呢.
我們去思想.
但我們見到他就不是.
他就是一個不相信宿命論的人.
他相信的是神保守的猶太敬虔的領袖.
第三章.

$^{761}$對不起.
這事以後.
阿黑除老王抬舉阿甲族.
哈米大他的兒子哈曼那阿甲族.
就是有些歷史的指引.
就是出埃及的十七章.
當時的民族.
來到去逼迫出埃及的這群以色列人.
就導致到神要透過摩西來宣告.
會世世代代.
和這個亞馬力人為仇.
這個亞甲族.
其實是亞馬力的後代.
所以原來哈曼.
原來他的祖宗.
可能還有這些深仇大恨.
一直落到.
射鳥英雄傳.
契丹人和漢人.
上一代的恩怨.
變成了元瀾洪烈.
還有那個.
很久沒看了.
忘記了.
和郭靖.
你就可以想象到.
那個問題.
但這個是一個.
屬靈的歷史的問題.
不揭開就揭開十七章.
出埃及記.
十七章.
讀回這些歷史.
大家就有印象了.
既然講到這個亞甲.
講到這個問題的時間.
第八節.
十七章的第八節.
出埃及記.
那時亞馬力人來到尼非定和以色列人爭戰.

$^{801}$摩西對約書亞說.
你為我們選出人來.
出去和亞馬力人爭戰.
明天我手裡要拿著神的像.
站在山頂上.
於是約書亞照著摩西對他所說的話.
就去行了.
和亞馬力人爭戰.
摩西阿倫和胡以.
就是兩位輔助摩西的領袖.
當時的父老.
都上了山頂了.
摩西何時舉手.
以色列人得勝.
何時垂手.
亞馬力人得勝.
但摩西的手發沈.
他們就搬石頭來.
放在他耳下.
他就坐在上面.
阿倫與胡以就扶著他的手.
一個在這邊.
一個在那邊.
即是說他的手.
不用太高舉.
但都是舉起的姿勢.
他的手就穩住了.
直到日落的時候.
約書亞就用刀殺了亞馬力王.
和他的百姓.
耶和華對摩西說.
我要將亞馬力的名號.
從天下傳言塗抹.
你要將這話寫在書上作紀念.
又念給約書亞聽.
他在山下打印象.
摩西就捉了一座壇.
起名叫耶和華離西.
這個意思就是.
這個名字反映著神某些的屬性.

$^{841}$就是耶和華是我們的升旗.
即是說是我們的升旗.
代表著國家.
民族的戰爭.
升起旗就是古代出戰.
是很重要.
所以耶和華離西.
神是我們的.
不是保護.
是我們的徵戰.
又說耶和華已經起了誓.
必世世代代和亞馬力人徵戰.
即是耶和華是徵戰的神.
是這個意思.
所以舉起升旗就是戰爭.
所以不是針對人.
是針對亞馬力這些形態的人.
如果用蜀靈的眼光.
就是一種這樣的形態.
所以我們稍後看到哈曼.
他就具備這個.
永遠與神為敵這樣的人格.
就代表著亞馬力是神要.
一定和這些人徵戰.
因為耶和華離西是徵戰的神.
亞馬力就是不得到神喜悅的人.
那些後代.
或者應該說先祖.
大家如果沒有忘記.
蘇羅沒有殺到阿蓋.
就是亞馬力派的阿蓋.
就招致神的憤怒.
我立蘇羅為王.
我後悔了.
和三母飲說的.
三母飲上的十五章.
我們曾經說過的.
正正就是這樣.
要滅絕正盡.
他就留下殺光了.

$^{881}$留下一個.
就是太子.
沒有殺到他.
他連戲他.
這個阿蓋.
就招致後患無窮.
蘇羅好像也有一些這樣的.
表面就好像.
神啊你不要那麼過分.
這個無辜的.
你也要殺了他.
我們就其實.
如果這麼解釋.
用人性的解釋.
有時我們就很難去明白.
神為什麼要趕盡殺絕.
沒有殺到太子.
還遷怒於蘇羅.
就後悔要立了這個蘇羅王.
不聽你的說話.
連戲不是很好嗎.
以和為貴.
但是神說不是.
他就是阿蓋的後人.
他就是阿蓋.
阿瑪力的這些品格的後裔.
好,回到第三章.
所以哈曼就代表著.
敵對神的阿瑪力.
(四).
這裡就很明顯凸顯了沒底改.
他以自己猶太人.
神選民的身份為榮.
才敢這麼做.
每個人都適時務者為盡傑.
適時務者就是隨時都轉態.
隨時都見到情勢不對路.
就立刻隱瞞隱藏.
很正常的.
但是不正常的就是沒底改.

$^{921}$不正常的就是他以自己是神選民為榮.
沒底改.
如果用我們廣東話.
阿哥,現在你在哪裡.
現在你不在以色列.
你現在不在耶路撒冷.
你以自己的身份當然是為榮.
阿哥,現在你在波斯的蘇珊城.
現在是什麼情況.
是什麼局面.
現在突然間有一個哈曼.
他被抬舉.
王吩咐所有的.
他之下的.
一人之下.
萬人之上的宰相.
我們積觀叫宰相.
的哈曼.
是要受到禮拜的待遇.
那個級數.
就是說拜王就肯定了.
但是連帶哈曼.
其他的神僕都要拜.
跪拜當然.
那就是一個命令.
不道理.
拜不拜.
真的信仰.
真金不怕紅爐火.
那麼純一的信仰.
沒底改就展示給我們看.
這些不是一些愚忠.
不是一些不懂得變.
做人那麼笨.
那麼死蠢.
不是,他知道.
他明白了那些事情.
他完全瞭解.
現在當下他不拜的後果.
但是.

$^{961}$他一定是以自己猶太選民.
其實亡國.
他是被撈.
他是以自己這個身份為榮.
所以.
錦書家我們在這兩天.
我們一起去.
再去談新一年的事宮.
其實重中之重.
由去年已經提出.
就是如何讓錦書家.
你受浸加入成為我們的會友.
你以這個身份為榮.
你感覺有含金量.
你感覺你是很有福.
你感覺你是很有幸福感和榮耀感.
很難在這個社會後現代.
很難個人主義抬頭.
洗禮證.
有什麼價值呢.
你進入哪間大學.
或者哪間子女的中學名校.
你就真的以那裡為榮.
你要買他們的債券.
都為榮.
都是的.
正正就是他們有實力.
但是神兒女的身份.
神選民的身份.
你能不能夠好像沒底蓋那樣.
不跪不拜.
好像主耶穌那樣.
撒旦,退我後邊去.
當拜主你的神.
但要事奉他.
你可以想象得到.
這些是我們生命裡面最重要的.
對造物主的敬拜.
沒有東西可以交換.
沒底蓋就做到了.

$^{1001}$之前.
公出了那些陰謀.
皇帝沒有賞賜.
現在雖然是大衛.
不是針對沒底蓋.
但沒底蓋知道有一個意識.
我不能向人下跪.
跪拜.
我只能夠向神下拜.
下跪.
我初期信主.
我們教會都教.
針查都不准跪.
爸爸媽媽.
那些邏輯和價值觀.
都是從那裡來的.
只能夠敬拜神.
不能向人跪拜.
當然你說可以斟酌.
是不是過分了.
可能有些可以斟酌.
不一定是錯了.
或者完全對了.
但這就肯定了.
沒底蓋真的面對這樣的危機.
不是由他自己犯錯製造的.
但聖經的描述就是.
唯獨沒底蓋不跪不拜.
你可以想象到.
如果是哈曼.
再加上他好像回到舊象.
我的老祖宗阿瑪力.
除非不知道.
知道是猶太人就不得了.
開始就發現.
《成卷書》的峰迴路轉.
《成卷書》的張力就在這裡開始.
你本來想著.
這個民族為什麼會有普耳節呢.
當然是因為以斯帖得寵.

$^{1041}$叫皇帝設一個普耳節.
讓猶太人可以多一個節日.
在亡國不要那麼痛苦思鄉.
原來不是.
原來普耳是一個.
轉危為機的一個節日.
好像我們的中秋節.
反清復明的一個抗爭.
這裡告訴我們.
在朝門的神僕就問沒底蓋.
你為何違背王的命令呢.
他們天天勸他.
他還是不聽.
他們就告訴哈曼.
要看沒底蓋的事.
站得住還是站不住呢.
因他已經告訴他們.
自己是猶太人.
就是說本來不關哈曼事.
本來不關皇帝事.
更加完全不是針對擺架子.
來去奚落這個.
得到皇來去信任的哈曼.
雖然你這樣理解.
當時得令的哈曼.
記住啊.
亞馬力在他內心.
是他的內心的價值觀.
就是說他裡面一定是有一些.
極端的驕傲.
在他裡面才成為了.
這個所謂哈曼.
成為了所謂亞馬力的.
神所爭護.
如果我們看箴言.
看一些其他的經卷.
比如《已不所書》.
神有告訴我們.
信了主耶穌之後.
舊的品格就不可以再有了.

$^{1081}$這裡就告訴我們.
哈曼就開始.
他內心的那些.
惡的不好的價值就出現了.
第五節.
哈曼見末底改.
不貴不敗.
他就怒氣填空了.
他們已經將末底改的.
本族告訴哈曼.
他就由一個小小的傲慢.
小小的一個憤怒.
來到他自己提升力度.
就是要將這個人.
和他們這個民族.
所以為拜的真神.
就不向任何人下跪下拜.
提升到一個這麼嚴肅的.
屠殺事件.
就是說比希特拉的屠殺猶太人.
還要來得早.
就是說希特拉又是.
哈曼的化身.
又是阿瑪力.
當然是屬靈的解釋.
阿瑪力的內心.
一個小小的.
當時德國的一個小小的官員.
但是可以因為時勢做英雄.
就變身成為一個殺人狂魔.
就要將在日耳曼民族.
最大的威脅的猶太人.
因為很精英的.
要將他們完全滅絕.
將他們所有的資產.
全部來到成為國家.
德國的財富.
你可以想象得到.
哈曼是在任何一個時代.
都會再出現的.

$^{1121}$他以為下手害沒抵改.
一個人是小事.
就要滅絕厄弗魯王.
就是中國所有的猶太人.
就是沒抵改的本族.
用什麼方法.
當然是透過王.
因為最好用.
所以皇帝身邊的大臣.
完全信任.
也有危機.
因為他的眼光.
就等於你自己的眼光.
所以有時都沒辦法避免.
如果多些有不同角度的.
中間的人.
就會比較平心靜氣.
所以一個領袖千萬不要忽視.
有同工有下屬.
向你提反對的意見.
就是說平衡你的眼光.
你的單一角度.
看這件事.
有時先入為主.
會主觀.
有時會有自己的驕傲.
一定是這樣.
但別人可能平心靜氣.
可以不是這樣.
如果你能夠謙卑.
你就停一停.
聽一聽別人的意見.
就會帶來.
沒那麼危險.
你就可以看到.
阿合徐老王很輕易.
就將他的介子.
給哈曼去做一個蓋印.
一蓋了印之後.
是誰都沒得更改的.

$^{1161}$皇帝自己都沒得更改的.
第七節.
阿合徐老王十二年正月.
就是離散月.
猶太人的月份.
人在哈曼面前.
按日日月月祭普爾.
就是祭遷.
所以普爾就是遷的意思.
希伯來文.
要定何月何日為吉日.
扎定十二月就是阿達月.
又是猶太人的節日.
那個月份.
哈曼對阿合徐老王說.
有一種民.
散居在王國各省的民中.
他們的律例與萬民的律例不同.
表面是對的.
沒錯的.
也都不守王的律例.
就是說.
關乎到敬拜真神的.
如果和波斯的律例有抵觸.
當然是不會跟波斯的律例.
比如一定要拜.
獨一耶和華神.
這個很合理的.
但是敵人就會將這些事情扭曲.
就要來到在王的面前.
來到他說有那麼壞.
得那麼壞的壞話.
如果王沒有做一個平衡的.
去停一停.
再找其他的人再去咨詢一下.
一方面這樣聽哈曼的一面之詞.
就出事了.
改變不了.
所以.
王聽到那麼大的威脅.

$^{1201}$這個民族不聽波斯帝國的律例.
有自己一套的律例.
如果那個領袖呼籲他們的神民.
推翻阿克徐努或者波斯帝國.
那就很大件事了.
收服到你.
臣服到.
當然沒問題了.
其實現在是做懷柔政策.
但是懷柔政策的危機就是.
他可以靜靜起革命.
他可以私底下.
自己的民族.
在這裡做聯絡.
裡應外合.
時機到了.
就可以推翻你的王位.
高度鎮壓就少這些問題.
沒那麼快出現.
但是如果懷柔.
有些不好的.
那就會有這個危機了.
所以王都是會擔心的.
很合理的擔心.
道弟都不守王的律例.
所以容留他們是與王無益.
這些是做慣官.
懂得將點到即止.
製造危機感.
給王有一個空間去擔心.
去懼怕.
現在還沒有一個有勢力形成.
現在要殺他一個措手不及.
在他現在的凝聚初期.
就要立刻擒賊先擒王.
就是這個沒底改.
要制止他.
果然其他的神文都見到.
這個沒底改真的有自己一套的律例.
是不跪拜你哈曼.

$^{1241}$意味著哈曼說的是真話.
他們真的有自己一套.
就會帶來對王是無益的.
這就是說這個危機表面正公成立.
王帝都聽到.
是真的.
那些神僕都說.
是他真的沒跪拜.
很不給面子哈曼.
王你安排了的.
吩咐了的.
一直說的時候.
第九節.
王若以為美.
請你下旨意滅絕他們.
我就捐一萬他連得銀子.
交給掌管國堂的人.
納入王的府庫.
不是得到那個獎賞.
我反過來.
我主動捐一萬他連得.
就是那個道量行的單位的重量.
的銀子.
來給國父.
就好像令到王.
你不做不行了.
因為現在這件事.
一方面好像天秤一樣.
一方面有國家危機.
另一方面就是.
這個忠心的手下.
主動的請紅去做.
因為你要去這樣做.
因為滅絕是不得到民心的.
現在哈曼是將自己的官位.
都壓住上去.
甚至還要將自己的.
有的那個俸祿.
拿一萬他連得出來.
來給國父.

$^{1281}$令到皇帝.
其實你不可能不接受這個offer.
如果是關乎到國家的命運.
第十節.
於是王從自己手上扎下戒指.
就是這裡.
就是波斯帝國的一個.
他們的習慣的一個命令.
王一出了金口.
就沒得修改.
包括王自己都沒得修改.
就吸猶大人的仇敵.
阿甲族哈米大他的兒子哈曼.
作者很巧妙地在這裡.
詳細的來說述.
他不是一個波斯帝國的哈曼.
他是阿甲族.
這個哈米大他的兒子.
這個哈曼.
所以現在將兩個.
好像屬靈的仇敵.
猶太民族屬神的.
阿甲族.
阿瑪力之派的.
兩個是為仇的.
是神好像令到他們為仇的.
好像在伊甸園的蛇.
和這個夏娃為仇.
一樣的.
那你就可以想象到.
這裡就是王就怎麼處理呢.
王對哈曼說.
這些銀子仍然賜給你.
這個民族也都交給你.
你可以隨意待他們.
那你可以想象到.
這個的哈曼略施小計.
就可以使得整個猶大的民族.
都沒等到神的審判員.
就要在神的國度以外.

$^{1321}$來到去被滅絕.
這個是不是神的旨意呢.
肯定不是.
因為神只是刑罰以色列人70年.
他們的民族.
他們的後裔.
是要回去重建.
為了什麼.
為了將來的尼賽亞降臨.
神的救恩的歷史.
神的救恩的進程.
是要一路繼續的.
而不是耶穌怎麼來.
那些民族都沒了.
所以這個哈曼他不識歷史.
更加不敬畏神.
更加不知道神有一個永恆計劃.
他現在在他面前的.
就是一個仇恨.
就是這個民族.
猶太民族.
所以他以為自己可以.
代表著黑暗的魔鬼勢力.
其實是不行的.
我們就可以想象得到.
所以哈曼真的代表著.
黑暗邪惡的勢力.
如果我們有哈曼的內心.
就要求主憐憫了.
因為後代.
比如猶大就是哈曼的內心.
他由於貪婪.
就去到自毀自己的人生.
就走不回頭.
希特勒諸如此類.
所以你要和神的旨意計劃去為敵.
就一定是沒好下場的.
所以這裡就讓我們看到.
來到第三章這裡.
作者一路描繪.

$^{1361}$這個民族現在到底.
是不是在神的掌管以外.
被滅絕呢.
神還掌不掌管著.
這個民族的歷史呢.
神是不是還在照管保護這個民族.
因為是亞伯拉罕的子孫.
神和亞伯拉罕立了這個約.
這些全部我們看回舊約.
看回神一路的作為的時候.
到底神是不是信實的呢.
他是不是應許了.
他自己都沒能力去改變呢.
現在就有一個當下的危機了.
怎麼辦呢.
第十二節.
正義十三日就照了黃的書記來.
因為是正月的時候離散月.
就祭簽.
現在去到十三日.
就是說這件事件.
發酵了十二日.
現在正式來到去成為一個文書.
改變不了的了.
這個命運.
不過唯一的就是.
還有時間.
現在是正月.
還有十二個月.
剛剛新星頭.
去到聖誕節.
差不多.
就是說有十二個月.
給抹底改去奔走.
去拯救自己的民族.
就是說這十二個月.
正正就是他們切切依靠神.
禁食祈禱的時間.
包括連這個耳詩帖.
所以皇帝看了歷史.

$^{1401}$可能很後.
所以有時神好像將我們放在.
急到一定要神.
一定要出手.
神就是不出聲.
神就是不出手.
就要看我們當下的人.
如何去仍然信靠神.
仍然不要認輸.
仍然依靠神.
當然我們說因為.
未有事件.
未有其他的事情.
我們很難具體化那些事情.
未發生更加我們不知道.
將來回望才知道.
抹底改一樣.
現在都不知道.
當下他就.
如果他看到這個禁令.
你就可以想象到那個危機.
自己的同胞.
會不會怪責抹底改呢.
你一個人的不跪不拜.
就害死我們整個.
姦淫之帝國裡面的猶大人.
這個罪名.
還有亡國.
神是否還看顧我們這個民族呢.
這是很大的歷史問題.
正如十三日.
就趙了皇的書記.
趙鑄哈曼一切所吩咐的.
用各省的文字.
各族的方言.
奉阿錘老皇的明寫旨意.
傳與總督和各省的省長.
因為有不同的民族.
因為波斯帝國是非常強大的.
在這個時期.

$^{1441}$並各族的首領.
又用皇的介旨.
丐人交給逆卒.
傳到皇的各省.
吩咐將猶大人.
無論老少婦女孩子.
在一日之間.
十二月.
就是阿達月的十三日.
全言剪除.
殺戮滅絕.
並奪他們的財為掠物.
抄錄這旨意.
頒行各省.
宣告各族.
使他們預備等候那日.
由正月十三.
現在頒這個正式的文書.
十二月十三.
剛剛足足有.
如果不計前後.
其實十一個月.
你就可以想象到.
那個危機如何化解.
誰可以化解.
這裡是第四第五章未交代.
第四第五章當然是高潮所在.
就是末底改.
去confirm以斯帖.
已知道你今天.
得了皇后的位分.
豈不是為了現在.
這個民族的危機.
其實以斯帖可以得自己的爸爸養父.
是因為你不跪不拜.
才有這個危機.
關我什麼事.
但是以斯帖不會這樣說.
他們都知道這是一個.
由阿瑪力.

$^{1481}$由阿甲.
由哈曼.
所引發的一個屬靈戰爭.
是一個民族.
面對一個敵.
好像魔鬼一樣的敵對.
要針對這個民族.
所以如果不是在神許可.
誰有權將神的旨意.
神的選民去滅絕.
你可以想象到哈曼.
可以那麼驕傲到一個地步.
借用壓垂老.
都不是很英明的皇帝.
其實只是製造了其他民族的擔憂.
隨隨便便就拿禁令.
就要來滅絕一個民族.
你的懷柔政策原來都是只會說.
原來小小事就要將人殺盡.
對自己的國家都會帶來動蕩.
但是不管了.
因為在這個時候.
哈曼已經是當時得令的一個宰相.
第十五節.
「亦遵奉皇命急忙起行.
子已也全遍了書山城.
王和哈曼坐下飲酒.
書山城的文卷都慌亂」.
兩個強烈的對比.
作者真是很巧妙地.
好像生花描筆.
王和哈曼完成了.
聖旨吸印了.
那些皇帝的聖旨就全遍各省.
即是127個省.
沒理由一個人全篇那麼多.
真是分散了.
即是要製造很多很多的聖旨.
即是要吸印都吸127份.
不單止.

$^{1521}$有省有市有其他鄉村.
要頒布了.
都需要一段時間.
對的,去到12月才執行.
因為有些地方是執行不了的.
就好像你投票.
有些海外投票都要先行.
你可以想象得到.
書山城就充滿了混亂.
懼怕,驚慌.
這個就是尤大人的寫照.
這個時候.
在驚慌的情況下.
最需要的就是那個屬靈的領袖.
他怎樣依靠神呢?.
他是用什麼策略依靠神呢?.
我稍微喝一口水.
尤其是在一些不是那麼容易的時代.
屬靈的領袖.
就是帶給那個群體.
有一個眼光.
就是要用屬靈去看那件事.
那才能夠渡過一個危機.
現在肯定是一個危機了.
有11個月.
書山城就立即知道了.
因為是首都.
但是其他的省市.
沒那麼快知道.
知道之後可能都是4,5月.
甚至到6月.
但是書山城現在就有得振沒得睡了.
哈曼就喝酒了.
和黃喝酒了.
就好像凸顯了黑暗掌權.
神的選民遲早被滅絕.
當然.
寫在聖經裡面.
一來不是猶太人在自圓其說.
而是真有其事.

$^{1561}$真有這個歷史發生過.
寶貴的就是.
當下活底改作為屬靈的領袖.
他怎樣去回應這件事呢?.
活底改知道所作的這一切事.
就撕裂衣服穿麻衣蒙灰塵.
在城中行走.
痛痛哀號.
到了朝門前停住腳步.
因為穿麻衣的不可以進朝門.
黃的御子所到的各省各處.
猶太人大大的悲哀.
禁食哭泣哀號.
穿麻衣躺在灰中的神多.
全個民族.
在這個時候感受到禁食的重要.
感受到禁食可以聲聞於天.
感受到自己的民族.
就快要面對滅絕了.
他們不會妥協.
他們不會覺得是宿命論.
宿命論就會怪責神.
你都不拯救我們.
算了 捆破了.
誰可以改變這個形態呢?.
耶路撒冷根本沒有居民.
現在是被捕的期間.
根本沒有黎希米的出現.
還沒出生.
那怎麼辦呢?.
反過來.
你怪我們自己以前得罪了神.
但已經在刑罰.
你想想這是一個這麼大的危機.
但我們看到整個民族.
在這個時候意識到嚴重性.
大家都懂得向神去禁食祈禱.
哭泣穿麻衣.
整個首都.
整個國家.

$^{1601}$都不是悲哀的皇帝駕崩的日子.
每個猶太人都穿麻衣.
很不老套.
但這正是民族的特性.
就是他們會守神的誡命.
不只是一般地上的律例誡命.
就是說現在我們這個民族.
面對這些事情.
神有要求.
神有規定.
就是只要.
我們又站起這些經文.
歷代至下.
大家可能都熟悉.
歷代至下的七章.
就是所羅門的祈禱.
所羅門造成了耶和華殿.
七章的十一節.
歷代至下.
所羅門造成了耶和華殿和皇宮.
在耶和華殿和皇宮.
凡他心中所照著的.
都順利造成.
耶和華向所羅門顯現.
對他說.
我已經聽到你的禱告了.
也都選擇這個聖殿.
作為祭祀我的殿宇了.
我如果使得天輩失不下雨.
使得蝗蟲吃著地的出產.
或者使得瘟疫流行在我民中.
這一班稱為我名下的子民.
如果自卑.
禱告.
尋求我的面.
轉離他們的惡行.
我必從天上隨聽.
赦免他們的罪.
醫治他們的地.
我必睜眼看.

$^{1641}$則而聽在子處所獻的禱告.
現在我已經選擇這個殿分別為聖了.
使得我的名永在其中.
我的眼我的心也必常在那裡.
你如果在我面前效法.
理父大位所行的.
遵行我一切所吩咐的.
謹守我的律例典章.
我就必兼顧你的國位了.
這個就是.
你可以想象到.
神同所羅門.
第一次他的正式就職典禮.
當晚.
就是散了那些人.
百姓回去各歸各家的時候.
所羅門靜下來.
其實應該當時.
他整個人沈甸.
準備睡覺的時候.
神就向他顯現.
來告訴他.
我已經定了這個聖殿.
是我月立.
我同在.
我的眼我的心都會在那裡.
所以這一群稱為我名下的人.
假如有一天.
當你感覺到.
無論天災.
無論一些事情.
你們來向我.
在聖殿祈禱.
我就會隨聽.
這是一個應許.
所以猶太人.
現在面對一個這麼大的滅族之災.
他們就醒起了.
當然就是有活體改.
他們似乎都不能夠再有.

$^{1681}$他們的聖殿敬拜了.
也可以因為這樣的緣故.
開始早期會堂.
未是這個時間.
會堂的初營.
有他們一些家族.
自己來守節旗.
會讀摩西的律法書等等.
但是去到真正的會堂.
就未是這個時段.
但無論如何.
他們現在就要面對.
這個滅族之災的時候.
他們就向神來禁食祈禱.
所以第四節.
皇后以斯帖的宮女和太監.
來將這件事告訴以斯帖.
她尋時憂愁.
就送衣服給活體改村.
要她脫下麻衣.
她卻不受.
似乎意識到.
如果這個民族再這麼高調.
其實等同於.
和皇的命令對抗.
你無事無干.
全個猶大民族.
穿上麻衣.
很禮貌.
等於你試想想.
感受堂為了抗議.
叫所有會友.
一起穿黑衣黑褲.
在湖邊舉起牌.
抗議.
亂說.
抗議什麼的.
維修之類.
很嚴重.
等於結集一個.

$^{1721}$強大的勢力.
去對抗阿克徐老.
這個命令.
所以以斯帖也感受到.
不要這麼高調.
帶來再嚴重的反感.
可能加快.
滅族的速度.
現在以斯帖.
還未表態自己是猶太人.
所以.
在皇宮裡面.
倖免於難之前.
他聽到.
供女告訴他.
以斯帖也很關心.
沒抵改.
他平常有沒有當官.
在官場又如何.
諸如此類.
也很掛心.
爸爸的情況.
但聽到的消息.
就不是那麼容易.
所以.
這個.
以斯帖送衣服.
接著就是我們待會會看到的.
那節金句.
第四章那節金句.
就是.
十四節那段.
第五節.
以斯帖就把.
王所派侍候他的一個太監.
名叫哈他甲.
召來.
吩咐他去見沒抵改.
要知道這是什麼事.
是什麼緣故.

$^{1761}$為什麼一群猶太人.
天天都穿上麻衣.
是什麼事.
是否想引起.
國家的.
處理.
是要瞭解的.
第八節於是哈他甲.
出到朝門前的寬闊處.
見沒抵改.
沒抵改將自己所遇的事.
並哈萬惠滅絕.
尤大人應許捐入.
王府的銀數.
都告訴了哈他甲.
又將所抄寫的.
全篇書生成的.
要滅絕尤大人的旨意.
即是那些副本.
photocopy交給哈他甲.
要他轉交.
給以斯帖.
要告訴他原因.
是什麼.
為什麼現在.
一群人.
禁食哭泣披麻.
蒙灰禱告.
是一個引起社會.
爭論的事件.
是會要調查的.
並祝福他.
進去見王.
為本族的人在王.
面前懇切祈求.
表面.
這個.
沒抵改.
以為借助.
以斯帖.

$^{1801}$可以直接.
去到王的面前.
但沒抵改可能不知道.
他立了的例.
連他自己都不能更改.
而且是否一個最適合的時機.
後宮不能幹政.
是否以斯帖可以隨便說兩句.
就可以改變.
王的缺址.
和王和哈曼.
之間的關係.
是否一個女子.
王后就可以改變.
現在關乎到整個民族.
和我們這個.
波斯帝國會有.
危害.
難道以斯帖一兩句話.
就可以改變嗎.
所以時機不對.
就會帶來反效果.
但沒抵改.
他只能盡他能做.
他現在唯一有的生機.
出路.
就是他的女兒.
感恩的是他沒有想過.
安插自己的女兒.
進入王宮.
要坐到不擇手段.
什麼毒計.
宮廷鬥爭.
無論如何都要坐到王后的位置.
你怎樣來競爭都好.
完全沒有這些心思.
完全沒有這些部署.
完全只不過是.
因為有一個.
廢後事件.

$^{1841}$而導致王要在全國.
頒布.
要選立新的王后.
所以似乎.
十八歲以上就要被.
召入宮.
所以沒得選擇.
不是現在貪慕虛榮.
或者沒抵改有些心計.
要安插自己的女兒.
來到王宮.
所以越是.
少人為參與.
越是一些計謀.
處心積慮的心計.
人的心.
和神的心.
有時差很遠.
有時人的心.
都不夠的時候.
就會用手段.
心計本來都是一個分析.
但心計有時.
往往會容易偏向於.
一些不好的念頭.
和想快些達成.
那件事就會用一些.
手段.
所以心計往往是負面.
是對的,所以越是想得太多.
就會.
帶來反效果.
現在沒抵改.
一初的初心.
完全不是因為他.
要安插一個.
自己的女兒來到.
在皇帝身邊.
更加不是密謀.
要推翻波斯.

$^{1881}$更加不是想飛黃騰達.
貴為.
皇親國戚.
仍然守住.
自己猶太民族.
什麼意思呢.
等候神的刑罰.
過了七十年.
他們就可以被撈歸回.
但這個問題就是.
我們見到.
這個情況是多嚴重.
對不起.
我可能搞錯了歷史的時間.
鴉鎚樓的時候.
已經是古烈之後的第三代.
所以.
沒抵改是當初被撈.
所以又有些.
奇怪的地方.
為什麼他又被撈.
現在歸回了.
在那邊已經有.
索羅巴伯上去重建.
這群人是留下在.
蘇珊城的一群猶太人.
沒有上去重建聖殿.
我要更正.
錯誤了歷史.
那個倒置了.
但無論如何.
沒抵改他現在.
要面對的這個情況.
他就求自己的愛女.
以斯帖.
你要為民族出頭.
祝福他進去見王.
第九節.
哈他甲回來將沒抵改的話.
告訴以斯帖.

$^{1921}$以斯帖就吩咐哈他甲去見沒抵改.
在那裡傳話.
說了一番說話.
解明聖旨.
和.
為什麼猶太人要披麻.
要禁食要哭泣.
要蒙灰.
以斯帖.
思考了一段時間之後.
就再叫哈他甲傳話.
因為他不能隨便出宮.
沒抵改.
更加不可能入宮.
他只能在朝門外面行走.
但他也不能進去.
因為他披麻.
就不可以進去.
就和他說了第十一節.
王的一切神祿.
和各省的人民都知道.
有一個定例.
如果不蒙召.
禪入內院見王.
無論男女.
都必定會被致死.
除非王.
向他伸出金杖.
不得存活.
現在我沒有蒙召進去見王.
已經三十日.
即是說.
王帝已經沒見我.
兩個原因.
已經玩厭了.
或者公務很繁忙.
王帝都沒空.
去夜夜笙歌.
但現實就是.
加增.

$^{1961}$不是原意那麼簡單.
是加增危機.
原來最後一個指望.
都幻滅了.
如果大家聽到這番說話.
就知道多嚴重.
第十二節.
人就哈他甲.
把以斯帖者話告訴.
王的一切神祿.
就把以斯帖者話.
告訴沒底改.
沒底改托人回復.
好像對話來對話去.
因為原因就是.
王后女兒不可以出宮.
沒底改更加.
沒權沒身份.
入宮.
帶來這個傳話.
沒有WhatsApp.
如果有WhatsApp就簡單.
自己對電話.
現在就是要靠這個哈他甲傳話.
感恩的中神.
忠於兩邊講真話.
這句話.
就是十三節.
就是以斯帖大家都很熟的金句.
沒底改托人回復.
以斯帖說.
你莫想在皇宮內.
強過一切猶大人.
得免者禍.
既然你是猶大人.
始終你和我們這個民族.
掛勾.
你不可以倖免於難.
你都不要明哲保身.
你都會面對這個情況.

$^{2001}$不過這個意思當然不是.
責怪他.
反而就是叫他.
憑信心.
走出你.
現在唯一你可以.
拯救我們這個民族.
的那個.
盼望.
剛才說了.
三十日沒見你.
你如果毅然進去.
皇帝向你伸金杖.
就沒事了.
而且皇帝肯聽你的陳詞.
講述背後的原委.
皇帝就改變這件事.
怎麼改變呢.
這個不知道了.
沒人可以改變.
因為皇帝自己都不可以再改變.
可以不改變.
可以立一條新例.
去推翻舊的.
立一條新例.
令到猶大人可以在當日.
十三日的普耳日.
來對抗.
那些敵人.
那就不同了.
皇帝現在頒下了另外一個法例.
就是准自衛.
就可以改變.
不是將舊的推翻.
而是立一條新的例.
原來這樣出路.
就是需要.
沒底改.
不要信宿命.
信自己是神選民.

$^{2041}$信做物主.
是創造萬有的神.
這些信念如何灌輸下一代.
灌輸我們自己.
就是我們有生之年.
在重大事件.
我們要揀選神.
認定神.
我們下一代才能夠學到那些教訓.
不然我信什麼呢.
我的父母.
有事就.
媽你不要浪費時間想神.
當然想辦法.
解決了不用找神.
下一代就知道原來不用找神.
沒問題.
他學到什麼功課呢.
現在沒底改.
他要求.
所以這一節13,14節.
就.
如果在我們一直看的上下文.
那個張力.
就不是一個金句那麼簡單.
而是一個民族.
不信.
命運的宿命.
信創造.
宇宙的神.
信我們這個民族.
是神所閱立.
是神所拯救的民族.
所以基督徒我們應該要.
有這一種很強烈的信仰.
我們是神所愛的.
神尚且.
不愛惜他的兒子.
為我們眾人捨了.
豈不也把萬物.

$^{2081}$和他一同.
次級我們.
保羅就說誰能.
控告神所揀選的人.
有基督耶穌已經死了.
並且復活了.
現在就坐在神的右邊.
誰能定我們的罪呢.
你看看保羅.
一直這樣說的時候.
誰能誰能誰能誰能.
我們要有一個很強烈的.
信仰的信念.
我們是已經屬神.
只不過還沒斷氣.
離開這個世界.
所以我們不能飛到天堂.
我們的靈魂.
保羅很想告訴我們.
雖然還受著物質世界.
地球的吸力.
但其實我們應該要.
在那邊生活.
這樣想東西.
用一些類比.
下個月你都要去英國.
你現在所有東西.
你都是想著英國.
想著那邊要買一件乾濕衣服.
好看一點的,經常下雨的.
等諸如此類.
你應該想好那邊.
基督徒就要這樣.
我們很少很少人會.
突然間想起.
我什麼時候會復活.
耶穌回來.
你興高采烈,很少.
我們根本都好像不懂這些真理.
我們根本不敢.

$^{2121}$這樣去想.
一有事的時候.
方章就死了.
改變不了,就好像.
認輸了.
但末底改這裡寶貴的地方.
就是十三節.
人就把以司帖者說.
告訴末底改,末底改就托人.
回復以司帖.
這些話都很重.
哈他甲就不知道怎麼演繹.
你爸爸告訴你.
焉知你得到皇帝的遺憤.
不是為現在的機會嗎.
這個金句.
是由哈他甲告訴以司帖.
如果你是以司帖.
聽到這個太監這樣說.
都挺有趣.
末底改托人回復.
以司帖說.
你莫想在皇宮里強過一切猶大人.
得免這禍.
此時你若閉口不言.
猶大人必從別處得解救.
可以奇妙地.
末底改原來不是.
最後一張皇牌沒有就沒有.
你若堅持不出聲.
明哲保身.
現在你仰尊處優.
現在你要和自己的民族割席.
你可以的.
但是猶大人.
必從別處得到解脫.
咦,這是信心來的.
這就是對神的一種.
還沒發生那件事的信心.
希伯來書十一章.

$^{2161}$信就是所望之事.
信就是所望之事的實底.
是未見之事的確具.
古人因為這樣的信.
就知道諸世界都是神創造的.
你沒有見到那個創造過程.
但是見到那些現在的東西.
就知道是造物主創造的.
你就可以想象得到.
有造物主存在.
是猶大人很強烈的信念.
所以這個末底改.
在這裡確認他的女兒.
依然是在這一方面.
在這一方面.
確認他的女兒.
焉知你得了皇后的位分.
不是為了現在的機會嗎.
這個機會就是你現在.
代表猶大民族.
直接進入皇宮.
向皇帝.
你的丈夫.
來表白.
我們是含冤.
我們是受到一個迫害.
看皇帝怎麼定奪.
你要表態.
你要站在猶大民族.
這個身份.
這一邊.
所以求主讓我們今天.
沒有迫害.
沒有任何形式的迫害之前.
基督徒保羅自己說的.
凡是敬前度日.
難免受迫害.
盡臣的國.
必須經歷很多艱難.
所以我們如果真的.

$^{2201}$尋求將來的天國.
人生.
我們不是為自己自招的麻煩.
自己製造一些迫害.
而是如果真的有迫害.
我們是要站得住.
我們是神的選民的身份.
我們是要有一個這樣的信仰.
成為一個強烈的信念.
我們要為真理來站穩.
這個世界是有.
造物主.
是有神.
是有.
他的界限的.
這個世界是要面對審判的.
我們要有這個.
很強烈的.
意識.
所以現在以斯帖.
就要面對這個取捨.
第十五節.
以斯帖就吩咐人.
回報梅底蓋.
即是哈達格.
這次是第三輪的對話.
你當去焦聚書山城.
所有的猶大人.
為我禁食三晝三夜.
不吃不喝.
宮女也要這樣禁食.
然後我遺離.
進去見王.
我若死就死罷.
於是梅底蓋照以斯帖.
一切所吩咐去行.
這裡就彰顯了.
神同在的證據.
雖然沒有提一個字.
關於神,但是這個民族.

$^{2241}$守神的誡命.
即是說他們立刻聯想到.
歷代志下七章.
十四節.
所羅門的應許.
即是說,沒錯.
不是一定回到聖殿.
猶太人,特別是在被擄.
之後.
他們相信聖殿是特別.
強WIFI,直上天庭.
所以在聖殿祈禱.
是靈赦靈厭.
但是現在.
總之你是奉神的命.
向神祈禱,猶太人.
神都會隨聽.
尤其是關乎到.
民族的滅絕.
所以以斯帖聽到這番說話.
這樣證明他.
焉知這個位份.
就是現在.
你就要代表這個民族.
當初沒有隨便說,現在更加要說.
即是說.
原來都是智慧來的.
我們不是隨便說.
基督徒,但是到關鍵的時候.
基督徒是我們的王牌.
我們是信奉耶穌基督.
是我們的honor.
是我們的幸福感.
是我們的榮耀.
我們感宣家有沒有這樣的弟兄姊妹.
我死.
我都以我是神而屢為榮呢.
這個我們就.
不需要懼怕了.
有亡的牧者.

$^{2281}$存不存在都不要緊.
因為我們有一班.
很成熟的弟兄姊妹.
站在這個位置.
去為神作見證.
是很需要這樣.
所以我們是需要教導弟兄姊妹.
對真理的認識.
然後你內心.
最強烈的就是.
你以成為神而屢.
這個身份.
為你最緊要最寶貴.
平時當然你不會.
隨便被人見到,當然透過.
你的潛移默化.
但是到真正要表態的時候.
是呀,我就是為.
耶穌基督.
而來到這樣的.
行徑.
所以這裡就是.
以斯帖真是很好.
他的自小的信仰.
沒有因為做了皇后.
就來到去.
改變了.
或者甚至乎已經.
刻意.
避免成為.
受牽連.
那.
似乎大有人在在歷史里.
但是真正的基督徒.
就是在這些位置.
真金不怕紅爐火.
就是這樣展示信仰.
赤裸裸的單純.
死就死罷了.
但是.

$^{2321}$都要三日三夜.
大家禁食祈禱.
不吃不喝.
講到明,不是168.
是真正的禁食.
三日三夜.
其實四十周夜.
都真是不出奇.
其實大家如果沒試過.
你試試你都知道.
譬如我們要照搶鏡.
我第一次.
在法國醫院.
一整晚.
射完之後.
第二天早上.
其實我很少有這個感覺.
我整個肚子空了.
沒有東西.
全部出來的是水.
是白色的清水.
整天.
我以往.
七點多就不吃東西.
胃就開始不舒服.
感覺很辛苦.
我由整天.
到晚上因為很多檢查.
是完全沒有吃的.
那個需要和感覺.
餓的感覺.
原來禁食去到一些階段.
其實只需要喝水.
就行了,三日三夜OK.
三日三夜都不喝水.
不吃東西.
但水就已經可以維持.
原來人體的結構.
是這麼奇妙.
所以現在去到關乎滅族.

$^{2361}$就真是需要.
這樣來依靠神.
就是智律行家.
作者不是故意不想寫神的字.
不是要避忌.
而是好像凸顯了.
這個民族已經完全.
對他們的信仰是深化了.
亡國又如何呢.
我們是神選民.
明成利就.
好景的時候.
仍然是神選民.
沒有取締.
但到有哈曼出現.
來成為敵人.
令到.
沒底改要面對一個危機.
仍然是猶太人.
不跪不拜.
不是針對你.
是針對我是猶太人.
我不隨便向人下跪下拜.
皇帝可能另外另計.
可能有些宮廷的禮儀.
要這麼做.
但我不會這樣做.
因為我是猶太人.
但現在以斯帖清楚了.
他知道.
他意識到危機是真實存在.
所以他就挺身而出.
以自己是猶太人的身份.
不是皇后身份.
是猶太人的身份.
大人身份.
是猶太人的身份.
大家一齊和他的宮女.
都一齊禁食三日三夜.
這件事.

$^{2401}$更加.
不是要聚集人氣.
不是聚集人氣.
集氣.
不是.
是要大家一齊.
很有意識地向神祈禱.
所以很有意識地.
一個群體祈禱.
是很重要的.
一個信息.
向神.
去懇求.
所以建立成為一個習慣.
成為我們群體的一個.
看重.
我們看重祈禱.
覺得人的力量.
有限.
我們不是要哀.
不是.
而是我們看重.
如果這些事情是這麼嚴肅的.
為什麼我們沒有人會看重.
這本身是一個問題.
是凸顯了一個問題.
我們靠什麼呢.
靠我們自己的力量.
靠我們的計劃.
靠我們有管理學的才華.
靠什麼呢.
我們有沒有依靠神呢.
依靠完神之後.
再用我們的才華是沒有問題的.
但如果單純什麼都依靠才華.
不依靠神.
那是問題所在.
不如我們不要叫基督徒.
我們普通一個群體就行了.
任何一個區議會.

$^{2441}$立法會.
大家一起祈禱.
禁食.
沒有這樣的.
人家做事很順暢.
但為什麼基督徒.
神的教會需要祈禱呢.
這不是一個儀式.
不是一個整個鐘頭的禮序.
其中一個.
而是我們看這件事.
真是我們唯有依靠神.
才叫我們這樣做.
所以這就凸顯了.
這個民族.
這群人.
當然他們已是代表.
其他散居在其他省的.
都是這樣敬虔.
有這群人.
他們作為屬靈的領袖.
當然是封眼.
因為慕底哥直接得罪了阿曼.
但是.
帶來的就是轉機.
就是因為他們.
依靠神的禁食祈禱.
這個是.
在作者這四章的描繪裡面.
就將那個.
風波.
有一個危機.
透過禱告尋求神.
不是要喚醒神的記憶.
而是神一路看著整件事.
問題就是.
發生在當下的群體.
如何將這件事.
帶進去神的計劃.
和在神裡面.

$^{2481}$尋求出路.
這個很重要.
所以我自己在事奉多年.
很渴慕的就是.
建立弟兄姊妹.
對於祈禱是真心.
去依靠神.
和成為了大家一個.
很重要的彼此相交.
其中不可或缺的部分.
和對聖經的認識.
福音.
全福音我就不是很強.
但我就很有負擔.
去查經.
祈禱.
很希望建立到一群人.
無論自己或家庭.
或自己的事業.
或教會的聖公.
都是真的以祈禱.
作為很重要的力量來源.
今天就來到這裡的時間.
我們就祈禱了.
多謝主.
我們去到第四章.
在第四章發生.
哈曼代表著邪惡.
對神的敵對.
對猶太人的滅絕.
我們看到整個歷史.
在一個發展里.
以為黑暗可以得勝.
但我們看到神.
一路的管理.
現在我們還沒有看到神的出手.
但是我們看到.
有一群人.
雖然表面是由他們而顯生.
但我們看到.

$^{2521}$以改和以斯帖.
以他們為中心核心的猶太人.
他們有真實的信仰.
有一個敬虔的內心.
他們看到.
這個需要.
就是要向神禱告.
唯有禱告.
禱告之後會有智慧.
會有力量.
會有心靈的膽量.
就是單憑一個女子.
以斯帖.
誤會地去進建.
進入內室.
進去見王.
好像很危險.
但是我們感受到.
當他們祈完禱之後.
有心裡面的決心和力量.
求主幫助.
教會.
在現在這個時代.
其實我們都需要把脈.
有移民潮.
有經濟的很多轉變.
很多人心.
其實經過了疫情.
幾年.
很多人其實很困擾.
很多事情可能表面.
是服尚.
是開心.
是很多旅遊消費.
但很多人其實都很多的.
內心的困擾.
或者是疲倦.
甚至有些內心的.
抑鬱都不出奇.
很盼望神都幫助我們.

$^{2561}$作為教會我們能夠.
真是將神給我們的.
恩典,力量.
和那份幸福.
去傳遞,去傳揚.
就是耶穌基督的.
救人的福音.
我們仰望著奉耶穌基督的名.
阿們.
孫悟空.
破功.
孫悟空.
我貪了.
滾.
不要.
翱翔.
林俊雄.
\newpage



\section{彼得前書 5:1-11}
\label{sec:fSxNb3R9eIk}
\textbf{2023-08-27 加恩樓《父慈子孝的神家》彼得前書 5\_1-11 講員:鄧泰康傳道}
\newline
\newline
連結: \href{https://youtube.com/watch?v=fSxNb3R9eIk}{\texttt{https://youtube.com/watch?v=fSxNb3R9eIk}} ~~~~ 語音日期: 2023-08-29
\newline
\newline
\hyperref[sec:iRhu93YiYcU]{\small{< < < PREV SERMON < < <}}
~
\hyperref[sec:index_chronic]{\small{[返順時目]}}
~
\hyperref[sec:index_scriptual]{\small{[返順卷目]}}
~
\hyperref[sec:xZ5COJOEvOY]{\small{> > > NEXT SERMON > > >}}
\newline
\newline
彼得前書 5:1-11
\newline
\begin{longtable}{cl}
\hline
\hline
章節 & 經文 (和合本修訂版)\\
\hline
5:1 & \begin{tabularx}{0.7\textwidth}{X} 所以，我這同作長老，作基督受苦的證人和分享將來所要顯現的榮耀的人，勉勵在你們中間的長老們： \end{tabularx} \\ \\ \relax
5:2 & \begin{tabularx}{0.7\textwidth}{X} 務要牧養在你們當中神的群羊，按著神的旨意照顧他們，不是出於勉強，而是出於甘心；也不是因為貪財，而是出於樂意。 \end{tabularx} \\ \\ \relax
5:3 & \begin{tabularx}{0.7\textwidth}{X} 不要轄制所託付你們的群羊，而是要作他們的榜樣。 \end{tabularx} \\ \\ \relax
5:4 & \begin{tabularx}{0.7\textwidth}{X} 到了大牧人顯現的時候，你們必得到那永不衰殘、榮耀的冠冕。 \end{tabularx} \\ \\ \relax
5:5 & \begin{tabularx}{0.7\textwidth}{X} 同樣，你們年輕的，要順服年長的。你們大家都要以謙卑當衣服穿上，彼此順服，因為「神抵擋驕傲的人，但賜恩給謙卑的人。」 \end{tabularx} \\ \\ \relax
5:6 & \begin{tabularx}{0.7\textwidth}{X} 所以，你們要謙卑服在神大能的手下，這樣，到了適當的時候，他必使你們升高。 \end{tabularx} \\ \\ \relax
5:7 & \begin{tabularx}{0.7\textwidth}{X} 你們要將一切的憂慮卸給神，因為他顧念你們。 \end{tabularx} \\ \\ \relax
5:8 & \begin{tabularx}{0.7\textwidth}{X} 務要謹慎，要警醒。因為你們的仇敵魔鬼，如同咆哮的獅子，走來走去，尋找可吞吃的人。 \end{tabularx} \\ \\ \relax
5:9 & \begin{tabularx}{0.7\textwidth}{X} 你們要用堅固的信心抵擋他，因為知道你們在世上的眾弟兄也正在經歷這樣的苦難。 \end{tabularx} \\ \\ \relax
5:10 & \begin{tabularx}{0.7\textwidth}{X} 那賜一切恩典的神曾在基督裡召了你們，得享他永遠的榮耀，在你們暫受苦難之後，必要親自成全你們，堅固你們，賜力量給你們，建立你們。 \end{tabularx} \\ \\ \relax
5:11 & \begin{tabularx}{0.7\textwidth}{X} 願權能歸給他，直到永永遠遠。阿們！ \end{tabularx} \\ \\
[1ex]
\hline
\hline
\end{longtable}
$^{1}$好,弟兄姊妹早安.
今天很感恩,剛才敬拜隊和我們唱的一首詩歌.
上帝最喜悅的敬拜.
上帝喜悅的敬拜就是將我們全人奉上,是嗎?.
所以今天和大家分享的時候.
為什麼我停了題目就叫做「父慈子孝神的家」呢?.
如果大家,特別是中國人,聽到「父慈子孝」,令你想起什麼呢?.
應該是要令你想起一個很幸福的家庭.
家庭裡面能夠做到「父慈子孝」的話.
你就會發覺這個家是幸福的.
你就會發覺其實今天彼得在這本書裡面.
他和我們分享了很多很多,到了第五章,他要去結尾的時候.
他講了一番心底的說話.
我在任何經文還沒講之前,我可以告訴大家一個很重要的進路是什麼呢?.
你就會發覺原來他告訴我們,每一個人在教會裡面都有自己的角色.
如果你能夠堅守你的角色,做好你的角色.
然後這個地方自然就會有幸福隨著你而來.
亦都是上帝所喜悅.
所以就像剛才詩歌裡面所講的一樣.
上帝喜悅的是你將你自己全然奉上.
我想借給大家,你們自己都有自己的家庭.
很多時候你們自己都會感覺到.
什麼叫做家不成家呢?.
我會將它細分到家不成家的意思就是.
每一個人你們都不是做回自己的角色.
爸爸不像爸爸,媽媽不像媽媽,子女不像子女.
舉個例子,子女永遠都是頂心杉.
然後你們說什麼的時候,你們就一言九鼎.
什麼都覺得你不對,跟你死去活來.
做爸爸媽媽永遠都不是愛你的,永遠都是壓榨你.
你跟他說道理,他就跟你說什麼?.
總之你閉嘴,我是你爸爸.
我想大家留意,這樣的家庭就會出事.
所以你會發覺當彼得在這裡說的時候.
他鼓勵當時的教會的時候.
我想大家留意,一直給我們的鋪陳都是什麼呢?.
全部都是以自身出發.
弟兄姊妹,我想你首先記住這個要訣.
你想幸福,你想你的生命,你想你的家庭.
去到上帝面前,是從自己出發.

$^{41}$人是很喜歡推卸責任的.
發生任何事都好,都是不關我的事.
我是最好的,都是別人錯.
為什麼家裡這麼不好,兒子不乖.
老公不好,老婆不好,永遠都是這樣的.
但你會發覺在聖經裡面都是,先回到自己那裡.
你就會發覺彼得一開始的時候說.
我這個作長老的,作基督受苦的見證.
在將來得到榮耀的時候才再勸你們,是什麼意思呢?.
他說其實彼得在這裡跟我們說.
他是以自身的見證去鼓勵你的.
就是他今天鼓勵大家.
你在神裡面,你要面對苦難,你要在上帝面前面對苦難.
首先彼得他自己是親身經歷苦難.
親身經歷困難.
而他在上帝裡面,他找到出路.
他用他的生命經歷的見證去跟你分享.
這樣很不同,你明白我的意思嗎?.
我想告訴大家,其實今天站在台上.
我要跟你在聖經裡面說一些很離地的大道理.
其實很容易的.
很容易的.
因為其實那些只是知識.
我要用原文幫你識經,幫你解經.
其實不難的,我坦白說.
我讀書不厲害,對我來說當然是難的.
但在在裡面,我都很喜歡讀書.
我所認識的肥沙,他讀到四十歲都繼續讀.
很喜歡讀書,很喜歡去研究.
所以你會發覺,在裡面我要把知識傳遞給你,很容易.
之前我跟大家說過,接下來大家不可以不面對AI.
check GDP.
很簡單的,其實它厲害到一個地步.
就是很強大的資料庫,你問他什麼,他會馬上回答你.
所以我都說,你不是避開他,你是善用他.
所以我現在查經,我用的.
我和大家一路在查經的時候.
有些問題大家突然間問我.
你真的當鄧全道真什麼都知道嗎?當然不是.
上次我們查雅歌.

$^{81}$雅歌裡面的百合花代表什麼?.
你明白嗎?這些是以前要找很多時間找資料的.
找一找查之後才知道在卷書裡面代表什麼.
我就不是了,我現在拿著電腦.
大家還在聊的時候,我都在問.
我想問一下,在聖經雅歌裡面.
百合花有什麼含義?.
哇,它真的很厲害,很快就彈出來了.
其實它幫到我們很多忙.
有人就說,他在那裡演的時候.
其實大家崇拜他很簡單.
這裡弄了一個虛擬的傳道人.
比Aka帥.
這個我可以肯定告訴你.
一定是你看著他的樣子,簡直是奧巴.
然後高大衛猛.
出來的approach又漂亮.
說話又可以去create.
然後我跟你說,有一篇道,很容易的.
但接著我就會想.
那我跟AI有什麼分別?.
AI是不是真的可以取代一個牧者?.
但我就告訴大家.
這個正好就是彼得所說的.
AI他沒有辦法去經歷人生.
他跟你說知識.
但他沒有經驗.
我們作為牧者.
彼得他作為他們的前輩.
他是用他的生命經驗過出來.
之後再去鼓勵接下來的人.
弟兄姊妹,這樣才真正叫做成全.
其實今天彼得他鼓勵的時候.
就是說我們每一個人今天.
在後面就在說.
有人作為牧者,有作為屬靈長輩,有作為師父.
有人是會作為羊仔,做晚輩,做徒弟.
我們每一個人都需要不斷去將上帝這些救恩去承傳下去.
但承傳不是靠嘴巴,而是你的生命.
今天我們傳出去的東西.

$^{121}$特別是基督徒,我們說傳福音傳福音.
我經常都說很簡單.
你傳給別人,你是不是真的覺得是福音來的.
如果你本身信耶穌裡面.
我自己都滿懷裡面很多的掙扎.
做基督徒是很吃虧的.
又不可以這樣,又不可以這樣,又不可以這樣.
接著我做基督徒的時候就做到很累.
其實做基督徒是很累的.
接著你做基督徒,給不到你那種生命的力量.
接著我就問你傳什麼給別人.
最後剩下的是什麼呢.
信耶穌,可以上天堂.
因為這個你都是信還是不信的.
但我經常都問的問題是.
你是否堅信你可以上天堂.
你是否堅信這個信仰在你今天的生活當中.
已經可以帶領你在地如在天.
是不是呢.
靈姐妹,彼得在這裡落了一個很重要的開頭.
就是我以自己曾經為基督去受苦而去作見證.
去鼓勵你們.
所以他確信這個信仰是真實的.
有力量的.
而到最後的時候是可以得到榮耀的冠冕.
所以你會發覺在聖經當中有很多我們很看重的人物.
這些長者其實是非常厲害的.
彼得相傳是倒釘十字架的殉道.
你不要說十字架.
今天我們面對的病我們都很不安.
不知道該怎麼辦.
所以我才說我相信上帝在裡面是滿有恩典的.
所以我們今天儘管面對疾病面對困難.
我仍然相信耶穌基督在這個困難疾病當中.
是在跟我們說一些東西.
教我們一些東西.
最近我都跟大家分享.
現在還沒有正式跟大家說.
接下來請大家為我祈禱.
我應該會做手術.

$^{161}$會做一個頸椎的手術.
剛才當我一說的時候.
有些人跟我說.
鄧炫東,我聽說做這個手術是寫不回來的.
但我可以告訴你.
剛才我坐在這裡.
我預備面對這個手術的時候.
我心裡不是很懼怕.
但我剛才去敬拜的時候.
我多了一份感動.
我會覺得原來我們今天能夠坐在這裡.
我們去敬拜上帝.
剛才我腦裡突然出現了很多圖畫.
平時你知道鄧炫東像螞蟻.
我在崇拜的時候走來走去.
坐不定.
原來我們能夠很安心很專心在裡面.
去享受在上帝當中的崇拜.
原來是這麼寶貴.
我也在想.
如果有一天上帝這麼恩待我.
我真的醒不來.
那我最多就上天堂.
但又反過來我又重新想.
我今天事奉的意義是什麼.
突然之間原來上帝給我們的每一個難處.
每一個艱難.
原來是美好的.
當你說你要將這個福音傳給別人的時候.
你有沒有這份確信.
我也相信這是給我作為一個很重要的鼓勵.
我今天在傳什麼給你們.
我傳給你們的東西,我信不信.
我是否經歷過.
我是否和彼得一樣.
我是否很有信心地告訴你們.
你們不需要擔心.
上帝就在裡面.
我很相信這些都是生命編織出來的.
接著下去的時候,彼得鼓勵他們.

$^{201}$他就分開了兩個群體.
一個群體我稱之為作牧者.
我們這些.
作長輩的,屬靈長輩.
我上一篇道,我都和大家說過.
你想不想也好,有一天你都會成為屬靈長輩.
因為你被人信耶穌,信得早.
另外一個就是師父.
很實在的.
你願意去承擔這個做師父.
帶著徒弟一起成長.
你可能做導師.
你可能幫忙教聖經.
你做這個身份.
在這裡面的時候,他直接和我們說.
第一句就是.
「務要目仰在你們中間神的群揚」.
你要目揚.
弟兄姊妹.
這是彼得給那些牧者.
那些領袖一個很重要的.
我們今天所做的事情.
你是不是目仰著弟兄姊妹.
他們是不是得到目揚.
剛剛報告說接下來會轉會.
有些朋友都問我.
其實我轉不轉會很掙扎.
我不會回以前的教會.
但來到這個時間,我又好像不知怎樣.
通常我會給兩個指標.
第一.
你在這間教會裡面,你找不找到神.
第二.
你在這間教會裡面.
你是不是得到目揚.
這是很重要的.
弟兄姊妹,今天我們在這裡.
你也知道我們家樓最喜歡的就是新教作.
我們有很多不同的東西.
但最重要的是這些東西出來的時候.

$^{241}$是不是真的目揚到弟兄姊妹.
目揚的意思是領路.
你為這些羊去尋路.
你要去幫牠尋找水源.
尋找食物.
這些朋友來到當中.
這些羊仔有沒有能量可以得到.
第三,你要確保牠的安全.
上次在哪裡,我跟大家說.
請大家小心.
千萬不要教會前門入.
後面就是懸崖.
讓牠掉進地獄.
我們今天每一個人有沒有這種觀念.
每一個來到的朋友.
其實我們就是前輩.
我們是需要目揚牠的.
我們是需要目揚.
第二,在剛才的經文.
另外一個版本是這麼譯的.
他在安神的旨意管教.
要目揚上帝託付給你們的羊群.
要甘心樂意地照顧牠們.
不是出於勉強.
不是為了要從中得到好處.
而是因為你們渴望事奉上帝.
我們要目揚牠們.
我們要按著神的旨意.
裡面包括了管教.
包括了教導,包括了愛.
但更重要的是出於甘心樂意.
而不是勉強.
這個在家裡是非常重要.
大家試過的.
當你煮飯.
無論是爸爸還是媽媽煮飯.
突然間子女走進來跟你說.
媽媽,我幫你洗碗.
我幫你洗點東西.
你開心嗎.

$^{281}$你開心的嘛,對不對.
但我以前媽媽是很含蓄的.
我一進來就說幫她.
她就說出去吧,妨礙她.
想清頂,我會馬上戒掉.
千萬不要有這些.
你一潑了一隻冷水,下次就不理了.
但反過來.
你在做到一頭煙.
然後你看到兒子在外面打遊戲機.
女兒在外面玩.
然後你插個頭出去.
我進來幫忙啦.
然後她的臉都燒焦了.
把東西丟下,然後進來幫她.
是兩個完全不同的圖畫.
是不是,弟兄姊妹.
甘心樂意.
你今天為上帝去做.
牧羊羊群.
你是不是甘心樂意.
還是你是勉強.
你做事奉.
你一勉強.
你就會出現埋怨.
在神的家裡面.
最忌諱的就是埋怨.
所以我經常都顧慮你,弟兄姊妹.
最近我們做的咖啡侍工.
我經常都說,第一個要訣是什麼.
不是有需要.
不是因為我們要人.
所以你來參加,不是.
而是你喜不喜歡這件事.
我想說喜歡.
你最初可能什麼都不懂.
但喜歡可以很多.
你可以直接喜歡咖啡.
你可以喜歡過程服務人.
但有一個更重要.

$^{321}$你喜不喜歡.
你明不明白在這杯咖啡的背後.
我們對人的那種服侍.
我們對人的那種安慰.
你明不明,你喜不喜歡.
如果你喜歡,你加入吧.
但如果你只是因為.
需要而幫忙的話.
我可以告訴你.
通常到最後都失敗.
大家有沒有試過參加教會.
經常去舉行那些叫讀經運動.
你去參加.
你去做.
你去做那些經常去舉行的讀經運動.
有沒有.
一開始的時候.
我們今年要立志.
我們要讀什麼什麼.
有些人就很乖的.
怎樣都好,跟著表照讀.
但我想問.
通常是成功還是失敗.
通常是失敗的.
為什麼.
問題是你喜不喜歡神的話語.
你喜歡嗎.
你愛慕他的話語嗎.
你覺得他的話語重要嗎.
如果你愛慕.
你覺得他重要.
我可以告訴你.
你怎樣都會把時間用出來看.
是不是,弟兄姊妹.
甘心樂意是很重要的.
我們今天作為牧者.
我們去牧養弟兄姊妹的時候.
是需要甘心樂意的.
但如果我們是勉強的.
誰叫我是牧者.

$^{361}$接著的時間.
我今天應該休息的.
你叫我的時間.
在這裡跟你說.
你以為弟兄姊妹聽不出來嗎.
她聽得出來的.
你想跟她談還是不想跟她談.
她很清楚的.
弟兄姊妹在這裡的時候.
提醒我們.
我們在裡面是要從心出發.
去牧養這些群羊.
而另外更加重要的是什麼呢.
你不是為了在裡面得到好處.
舊的版本亦為前才.
但不只是前才.
當然當時有當時的背景.
但在這裡說好處.
其實當你接觸人的時候.
要小心你的心有沒有歪念.
你想不想在建立的過程中.
想得到一些好處.
她說你要小心.
我們不是為了這些好處.
我們不是為了他們的回饋.
我們不是為了得到一些什麼利益.
而反而是因為我們渴望.
去事奉上帝.
因為你很感受.
你給我們這個機會.
我們能夠去做.
你在整個出來的時候.
你才能夠令羊仔感受到愛.
感受到你裡面的誠懇認真.
另外一個層次.
你不要克制他們.
另外一個版本亦為.
你不要專橫地對待他們.
即是高高在上.
有你說沒有他說.

$^{401}$我叫你這樣你就這樣.
這些是人性原本的裡面.
因為人性的原本就是.
我越高就越有power.
越有權柄.
我就可以指使你去做什麼.
命令你去做什麼.
其實這種權柄是人很追求的.
我想告訴你.
人很喜歡特權.
我舉個例子.
我是感恩的.
我不是覺得是特權.
我剛剛去看醫生.
原本醫生都很貴.
有一位弟兄介紹員工.
我帶了八折給你.
你心裡很開心.
看醫生都有八折.
但問題是反過來.
你是不是追求這一樣東西.
那個人介紹我的.
我多有面子.
不是的.
在裡面是感恩.
而反過來反而是叫你做榜樣.
就回到彼得剛才所說.
你要做榜樣.
今天你要去讓人的時候.
你要成為一個榜樣.
你要做到給他看.
什麼叫將難.
當一個擔子交給人.
是連你自己都難的時候.
你不要叫人做.
有時候我經常說講大道理.
用什麼呢.
人們走到你的時間.
犯罪犯罪不對.
我經常都說一樣東西.

$^{441}$先問自己.
你犯不犯罪.
你都有犯罪.
我不是說犯罪對.
不要一去就去到鐘擺在另一邊.
我是說的時候.
每一個人面對罪很不容易.
他裡面有很多心理的掙扎.
很多原因.
我們是和他同行.
你做目的就是和他同行.
你不是只是指責他.
你是去幫助他和你一起渡過.
正因為我們明白.
不容易.
所以到最後的時候.
我們是渴望得到神最後的讚賞.
這是給這些人的總結.
第二部分的印.
是關於小羊.
或者叫做宿命晚輩.
或者叫做徒弟.
他說第一個.
他說你要去純服年長.
我想大家留意.
這個是基本的規則.
上帝創造的時候.
有一些角色.
剛才我說寫得很清楚.
爸爸就是爸爸 媽媽就是媽媽.
子女就是子女.
每一個角色裡面.
上帝賦予了他一些權柄.
所以管教只可以.
爸爸媽媽管教子女.
不是子女管教爸爸媽媽.
有很多不同的.
有很多責任.
所以我和我的兒子分得很清楚.
意思是什麼呢.

$^{481}$我要賺錢 我要養你.
我要供書教學.
你的責任是什麼呢.
你的責任就是讀好你的書.
這是很基本的.
在我們角色的不同.
所以你會發覺.
這是上帝告訴我們.
首先你要尊重上帝的權柄.
因為這個權柄.
從神而來.
你是尊重上帝的權柄.
今天是沒錯.
做爸爸媽媽的我們是人.
我們是軟弱的.
我們做得不好的地方.
我經常提醒自己.
我媽媽很硬頸.
有時反過來.
不舒服的時候.
看醫生.
那些就會出來.
我只可以在旁邊提醒自己.
到某一天我會變成這個年紀.
請不要這麼硬頸.
你永遠都要記得.
在我做子女的時候.
爸爸媽媽原來這樣做會帶來壓力.
我希望我不會這樣.
我只可以這樣.
但反過來.
會不會因為你爸爸媽媽的硬頸.
他做得不好.
你就不需要尊重他.
不是,聖經裡說得很清楚.
你都要尊重他.
因為他就是你爸爸媽媽.
因為你是馴服上帝.
他自己的權柄.
但很有趣,經文一落的時候.

$^{521}$一轉就轉到什麼呢.
就是你們眾人.
要謙卑.
要彼此馴服.
因為神阻擋驕傲的人.
一轉頭.
就說你要彼此馴服.
很有趣.
所以你會發覺這才是真正的核心.
就是眾人.
都要謙卑.
彼此的服侍.
彼此的馴服.
為著的是要服侍對方.
我用一個.
做一個例子.
鄧全道.
我不喜歡吃魚.
如果你問我,我無法解釋給你聽.
我又覺得有腥味.
我又很怕那些滑溜溜的口感.
我很怕.
所以我從來不吃魚生.
我從來不吃三文魚.
我是不吃的.
所以上次他們.
找我去吃飯的時候.
一見到我在日本放題.
我就.
這樣.
無論如何,我都不吃.
所以我很怕魚.
很滑溜溜的口感.
我吃那些很粗糙的.
但我媽媽蒸魚.
你怎樣.
你會怎樣.
我不吃的.
我都說我不吃.
這樣嗎.

$^{561}$不是.
彼此馴服.
我會.
借著媽媽.
她的辛勞.
她開心的緣故,我會吃一些.
我不會批評.
你要發覺在這裡的時候.
我們有一個很基本的架構.
不過在裡面.
什麼叫父慈子孝.
就是這個意思.
爸爸本身有沒有這個權柄,有權柄.
但他寧願成為一個.
慈父,他寧願.
讓他的子女感受到什麼叫愛.
他寧願放下.
他自己的一些權利.
一些利益,而寧願滿足.
這個子女.
做子女的,什麼叫子孝.
他真的會懂得孝順.
令到爸爸媽媽他們真的會感覺到開心.
他知道爸爸媽媽的需要.
所以有時候.
有些事不用這樣,我見到有些事很粗暴.
爸爸媽媽很不喜歡.
子女玩遊戲機.
有時會衝去房間拔電掣.
我不知道是真是假.
我見到網上有一個.
拿一個棒球棍.
打爛電腦.
我見過有這樣的短片.
很誇張.
但弟兄姊妹,你自己也很清楚.
這個不是方法.
方法就是.
做子女的,他真的明白.
爸爸媽媽在這一刻,有些事是為你好.

$^{601}$你做點睡,為什麼.
你做點睡,他會多生一塊肉嗎.
他不會的,你沒有為了爸爸媽媽.
他擔心的緣故.
而你願意去做點睡.
你願意去.
多穿一件衣服.
弟兄姊妹,彼此.
馴服,慰著的是能夠.
祝福,造就到對方.
所以他說,不馴服.
是什麼呢?不馴服是來自驕傲.
你挑戰他的權威.
你認為自己比他厲害.
背後我認為自己比上帝厲害.
你就自然不會馴服.
一切的馴服是源於什麼呢.
你服在神的大能手下.
你馴服於他的設定.
你馴服於他的詞語.
你馴服於他的規則.
一切的驕傲是源於什麼呢.
後面說到,突然間他又說到.
你很憂慮,卻急神.
原來一切的驕傲是源於你的憂慮.
為什麼我這樣說呢.
你很憂慮別人看你的眼光.
你就會在別人面前膨脹自己.
你就會瓦起一些鋼.
你就會發現原來是.
上帝在裡面很走趣.
常常叫我們將我們的憂慮卸給急神.
因為上帝會照顧一切.
接下來你心裡就會平安.
你就不需要再瓦起鋼.
你會相處得更好.
你不再為未來的擔心.
這就是作為晚輩.
作為羊仔.
我們需要去做的.

$^{641}$父慈子孝屬靈的教會.
為了什麼呢.
後面給了我答案.
你要謹慎警醒.
因為你們的仇敵魔鬼.
如同孝叫的獅子遍地遊行.
尋找可吞吃的人.
原來大家留意.
我們今天所有的東西.
是有屬靈爭戰在裡面.
一間教會裡面.
為什麼我們要.
一定要維持一個美好的關係.
因為混亂的教會.
是魔鬼最好的機會.
一間教會裡面.
充滿著.
互相的挑剔.
互相的批評.
一間教會裡面充滿著.
我不服你,你不服我.
自然這間教會.
就會散了.
這間教會就不會有一個一致的方向.
所以你就會發覺.
在聖經裡面.
很有趣的.
為什麼他會跟大家說這麼多權柄的東西呢.
我只是想.
我是一個傳道人.
傳道人是不是很聰明呢.
是不是一定很聰明呢.
不是的.
你教清楚.
我是不是所有東西.
都是最厲害的.
不是的.
但是反過來.
傳道人在教會裡面.
卻是有一個很重要的責任.

$^{681}$要帶領大家的屬靈方向.
所以我想問一下.
特別是現代.
現代裡面多了很多的機會.
大家有機會進修神學.
甚至有些大家在神學的領域裡面.
你拿了PhD.
我只不過是一個小蝦子.
我是一個學士.
可能你比我優秀.
但我就是在神裡面.
我有這個權柄.
上帝給了我權柄,我作為牧者.
我要去帶領教會的屬靈的方向.
頂子梅.
我告訴頂子梅,這個世界上.
沒有最好的.
在教會裡面,我們只是.
順從上帝自己的.
規則裡面.
我很清楚,我作為牧者.
我要盡我自己最大的能力.
去.
小心求證.
去小心的去尋找.
教會.
屬靈的方向是應該怎樣.
所以很多時候,我應該要小心這個位置.
這是我的責任.
而頂子梅是需要馴服在.
這些位置裡面.
我想告訴頂子梅.
一間教會這樣的運作.
你才會繼續向前走.
你才會繼續有一個很清楚的方向.
當然裡面有沒有.
共識,當然有,但我永遠都是那一句.
所以我就很推崇.
一件事,就算現在行政上.
我們有一個叫管治委員會.

$^{721}$管治委員會裡面.
每一個頂子梅,每一個同工.
他們都有一個角色負責的範圍.
每一個範圍.
不是一個人做完,一個人做死.
下面有人幫忙.
但我永遠都提倡.
就是到最後.
如果有些方向.
大家不一致,由誰決定?.
由組長決定.
你這樣才可以.
繼續向前走.
如果你裡面出現很多不服.
我不服,我不服你.
我不認你這個,你那個部分就散了.
但不是.
我們很信服,因為你是組長.
你有這個權柄.
所以你決定了,我們就一起向前走.
是吧.
而裡面的時候,我們每一個人都多了一份謙卑.
我們不會知道.
全局.
我們只可以盡力而為.
所以在上帝.
裡面的時候,他就告訴我們.
當你這樣做,魔鬼就沒有機會進來.
魔鬼就沒有機會.
魔鬼就沒有機會在裡面的時候.
去抵擋我們.
反而反過來,彼得說.
我們要同心以對神的信心.
去抵擋他.
什麼叫信心?聽從神.
跟從神,是必成的.
這個就是我們的信心.
我們相信上帝的大能.
所以你就會發覺,他告訴你不用怕.
為什麼?你身邊有很多的見證人.

$^{761}$跟你一起告訴你.
到最後的時候,我們會得勝.
所以你就會發覺彼得在這裡的時候.
其實他是幫一間教會.
去打好一個很重要的基石.
弟兄姊妹.
我們今天每一個人.
有你自己的角色.
弟兄姊妹,做好你的角色.
在裡面的時候,剛才我所說.
當你有機會.
幫忙成為咖啡師.
一起幫忙沖咖啡的時候.
請你沖好杯咖啡.
去.
做好它.
讓別人感覺到.
你裡面是有一份誠意.
其實裡面很有趣.
最近我們都說.
昨天才剛開完會.
我們有Eddie.
他幫我們不斷地.
去監察這個品質.
這個情況之下.
他會一直弄.
告訴我們怎樣才算是最好.
說真的,我要沖杯咖啡.
跟他沖杯好咖啡,其實是一線之差.
他弄完之後.
原來是有分別的.
那你就要去做.
你去做好它.
當你今天.
上帝放你在這裡的一個位置.
請你盡力做好它.
最簡單,今天我們有很多弟兄姊妹.
一起幫忙參與師事的工作.
師事的工作.
很多瑣碎.

$^{801}$很多雜務的東西.
我們要派東西,我們可能要收奉獻.
很多其他,不過弟兄姊妹.
其實你是第一個防線.
朋友,特別是未信任的朋友.
來到,看到那個是你.
你給他的.
是一個很歡迎.
還是你給他的有一種.
看不到.
很忙.
還是我們以人為先.
我們能夠給一個.
印象,一個形象給人.
我們今天.
弟兄姊妹他們參與敬拜.
我很欣賞他們.
他們不斷擠時間出來練習.
為什麼要練習.
練習是我們要呈現一個.
最好的敬拜.
準備好的我們去.
給上帝帶領弟兄姊妹一起敬拜.
弟兄姊妹,當每一個部份.
你都盡心盡意盡力去做好它的時候.
這間教會.
會不得了.
所以其實很在乎的,究竟你的心裡面.
你是按著什麼去進行那件事.
可以很簡單.
可以很隨便.
到最後我想用一個例子.
去做結束.
我上次和一位姊妹去聊天.
姊妹她.
都是在她的機構.
裡面的時候去做一些管理的工作.
我在教會我都要做一點點管理.
起碼我下面也有朋友可以幫我.
那時候大家在交流.

$^{841}$其實我很開心.
那次的對話.
平時她剛剛轉了部門.
然後她有新的環境.
她要去.
做這個管理.
那些人個個都不懂.
她說了一件事是令我驚醒的.
她說其實.
我去帶領.
人的時候.
有時候會被人用一個感覺.
她覺得我很mean.
每一樣東西要求很高.
但她說.
其實原因是什麼呢.
因為我.
管理裡面我有一個很重要的理念.
我很想.
跟過我的人.
因為有一天我們都會分開的.
我很希望這個人.
去到每一個.
其他人的面前.
人家都喜歡他.
我做了那麼多事.
他在那一刻可能會覺得很討厭.
為什麼你要那麼嚴格.
但問題是.
我很想告訴他到有一天的時候.
他能夠可以.
讓其他人都喜歡他.
在那一刻你不明白的.
但反過來.
到後來的時間.
就很有意思.
她說我教我自己的子女.
我們都是這樣教的.
但當她突然說的時候.
我驚醒了.

$^{881}$原來我們今天回來.
我們跟同工合作.
我們不只是工作上的關係.
原來是有牧養的關係.
原來我們都很盼望.
有一天的時候.
她也會喜歡他.
喜歡他不是什麼.
大家有沒有合作過.
有些人跟你合作的時候.
遲交遲到.
沒要求.
所有的事的時候.
是很討厭的.
但如果你跟他合作的時候.
準時.
準時快手.
令你很開心.
所以我可以告訴鄧英正.
吳全離開了.
當然是很可惜.
但我還沒有這種感覺.
但在九月尾.
我們將會有一位同工離開.
我就會覺得是世界末日.
就是我們有一個幹事同工.
他專門幫我們設計海報.
之前大家見過他.
Elsi.
他很快手.
他完全.
稱得上是.
秘書.
他現在回到感受那裡.
他一直都是成為.
曾牧一個秘書的角色.
曾牧昨天.
早兩天開會他也說.
經常催我交什麼.
沒有人敢提我.

$^{921}$是他.
做得很好.
所以弟兄姊妹.
原來我們每一個人當我們做了一件事.
我們背後有一個很清楚的理念的時候.
你就會發覺.
不知不覺之間我們建立了那個地方.
上帝很常.
所以這個姐妹提醒我一件事.
當你今天做主管.
請你盡力去做好你的主管.
你背後的時候.
是要祝福到你的下屬.
所以求神幫助.
我希望這個例子能夠提醒我們.
今天大家回去的時候.
求上帝你做爸爸媽媽.
請你做好一個爸爸媽媽的角色.
你做子女請你做好一個子女的角色.
你是神.
祂的家裡面的養子.
請你做好這個養子.
你今天做大嶺請你做好大嶺.
我們一起去建立.
這個屬於上帝的地方.
求神祝福你.
我們一起去禱告.
將我求主你自己去大嶺.
願意我們在這裡成為一個.
我們各就各位.
在自己各自的各位當中.
他能夠得著恩典的地方.
願意你自己在這裡看顧.
帶領我們.
我願意在這裡成為一個美善.
成為一個每一個人來到.
都能夠經歷到喜悅.
經歷到愛.
經歷到更生同在的地方.
我求主你自己首先更生我們在這裡的.

$^{961}$在座每一個弟兄姊妹.
願意我們的生命先被改變.
願意我們的生命當中.
不停留在我們自己的難處裡面.
我求主你自己親自去保守大嶺.
祈禱奉教主耶穌得勝明志祈求.
更多.
\newpage



\section{哥林多前書 16:15-16}
\label{sec:xZ5COJOEvOY}
\textbf{2023-08-27 午堂《為主同工,同勞的人》哥林多前書 16\_15-16 講員:曾裔貴牧師}
\newline
\newline
連結: \href{https://youtube.com/watch?v=xZ5COJOEvOY}{\texttt{https://youtube.com/watch?v=xZ5COJOEvOY}} ~~~~ 語音日期: 2023-08-29
\newline
\newline
\hyperref[sec:fSxNb3R9eIk]{\small{< < < PREV SERMON < < <}}
~
\hyperref[sec:index_chronic]{\small{[返順時目]}}
~
\hyperref[sec:index_scriptual]{\small{[返順卷目]}}
~
\hyperref[sec:fjBMO0wHlN4]{\small{> > > NEXT SERMON > > >}}
\newline
\newline
哥林多前書 16:15-16
\newline
\begin{longtable}{cl}
\hline
\hline
章節 & 經文 (和合本修訂版)\\
\hline
16:15 & \begin{tabularx}{0.7\textwidth}{X} 弟兄們，你們知道司提法那一家，是亞該亞初結的果子；他們專以服事聖徒為念。 \end{tabularx} \\ \\ \relax
16:16 & \begin{tabularx}{0.7\textwidth}{X} 我勸你們順服這樣的人，和一切與他同工同勞的人。 \end{tabularx} \\ \\
[1ex]
\hline
\hline
\end{longtable}
$^{1}$頂止每日主女便平安.
今早的正道不是一個詳細的解經.
因為更重要的是今天我們邀請了岑偉龍牧師.
其實已經好幾年沒有回來了.
疫情和過去的事奉.
稍後他就會再詳細將他在這麼多年.
在英國的南倫敦的事奉.
來跟大家分享.
所以我這個是一個引言.
是一個我們在主裡面.
必須要大家一起同工.
保羅更加寶貴的一個描述.
對於Steven.
不知道有沒有人叫Steven呢.
有的 其實已經舉了手了 好幾個.
他是阿該亞省 即哥林多教會.
很早期就信主的家庭.
保羅幫他們施洗.
現在保羅面對哥林多教會內部的一些紛爭的時候.
去到結尾.
他引述了一個他在內心裡面.
肅然起敬的家庭.
為他們來感恩.
又要求哥林多教會要馴服這樣的行徑.
在教會裡面來事奉的弟兄姊妹 家庭.
我再讀多一次.
弟兄們 你們都曉得.
斯提反一家.
是阿該亞初結的果子.
保羅曾經在哥林多有一年半的時間.
來向 主要是在猶太人的會堂.
來講道 有猶太人信主.
但更多是對他的迫害.
他就離開了.
去到一個叫提多猶時到的會堂的地方.
在那個家裡 那個是管會堂的猶太人.
在那裡聚會 開始向外邦人傳道.
有很多哥林多人聽到來信主.
現在我們看回這裡.
這個斯提反很早期就受到保羅的教導.

$^{41}$信主 他們的生命有改變.
全家歸主.
而且他們有一個很寶貴的.
保羅對他們的讚賞.
專心以服侍聖徒為念.
專心服侍聖徒.
這個斯提反在保羅的眼中.
對這些聖徒的內心的肅然起敬.
因為他們在教會裡.
不單止不是製造麻煩.
是專心服侍其他人.
專心留意哪些聖徒有需要他們服侍.
這樣一位斯提反.
和他們全家這樣歸主.
使徒保羅就告訴我們.
我勸你們哥林多人要信服這樣的人.
並且一切同工同勞苦的人.
同工為神.
與保羅同工的人.
當時是有困難的.
教會的存在有逼迫的.
去到一年半之後.
保羅就面對當時哥林多的一些地方官.
就對他的迫害.
他要走了.
你可以想像得到這個同工.
是要面對一定的壓力的.
但是保羅都不會介意.
很寶貴的地方就是.
連當時哥林多教會.
都有弟兄姊妹.
看到自己以服侍神.
服侍使徒.
服侍其他的聖徒.
作為他們的主要在教會裡.
用心的去留意不同人的需要.
盼望弟兄姊妹我們能夠在神的家裡.
來彼此的留意.
彼此的關心.
哪些人有需要.

$^{81}$可能先是禱告裡去關心.
可能去觀察.
可能去留心.
去留意.
這是主所喜悅的.
這個時候.
真的要介紹岑牧師了.
1998年.
我都還沒有來錦繡堂.
已經和錦繡有一個差遣.
去英國.
來看到當地有很多華人的需要.
就是這樣來到一路一路直堂.
總共有沒有二十多年.
1989年.
我記錯了.
34年.
在英國的南倫敦.
等一下岑牧師會講詳細一點.
有多少間堂他直了等等.
這是很重要的.
我們錦宣家在過去的前輩.
一路的帶領和感動.
不單純是社區的福音需要.
還有很多地方的福音需要.
所以今天真的有我們的前輩.
在我們中間.
來講述神怎樣感動岑牧師和太太一家人.
來去同宮同牢府.
在神的國度裡面.
所以今天很簡短的正道.
重點是岑牧師等一下會有比較詳細的見證和分享.
還有一個影片.
我們現在把時間交給岑牧師.
各位弟兄姊妹.
早安.
很高興今天又可以回來和大家.
分享.
應該是.
聽岑牧師講.

$^{121}$應該是兩年前.
倒回來的.
因為疫情的緣故.
所以延遲到現在.
其實前輩應該是.
陳一華牧師.
就是前輩.
我在這裡工作的時候.
他已經是我的牧師.
所以他應該是前輩.
我是84至87年在這裡聚會.
那時候在幼稚園.
我那時候是五重節的老師.
原來我有些學生認得我.
我不認得你們了.
不好意思.
在教會見到是很好的.
在教會裡面事奉是很好的.
想不到有些學生.
在主的裡面.
仍然很愛主.
很感恩.
也有些在幼稚園的時候.
他說我已經在幼稚園為你祈禱了.
也很感恩.
幼稚園應該是三十多年了.
為我祈禱.
最近有一次.
陳偉安牧師.
他去英國.
去我們南倫敦宣道會講道.
他一進來.
岑牧師.
幼稚園就為你祈禱了.
陳偉安牧師.
有個牧師.
幼稚園為我祈禱了.
看到神很祝福錦宣家.
我相信.
神祝福錦宣家.

$^{161}$就是要祝福更多的人.
不論在這個地方.
在元朗.
洪水橋.
甚至以外.
更多的人得到這個福氣.
聽到有位.
我以前的學生.
我做短孫.
在土耳其.
又是感恩.
我們應該多些.
去將福音傳開.
我相信我們的盼望.
在於主.
我們主.
在主的裡面.
必定是有福的.
也多謝曾牧師.
給我有機會去分享.
在過去這麼多年.
為什麼要分享過去這麼多年.
可能因為我已經退休.
我已經過了退休年齡.
所以就分享過去三十多年的.
的福事.
因為你們支持了我.
三十多年.
多謝大家.
如果沒有你們的支持.
沒有你們的待討.
沒有你們的關心.
我相信我都不會站在這裡.
多謝大家.
三十多年.
很長了.
其實我是建立了五間堂會.
以前那些.
就是獨立教會.
為什麼呢?因為是其他的機構.

$^{201}$因為以前宣道會.
在英國.
是沒有.
聯會.
所以也沒有差會.
所以宣道差會差了我們過去.
就和其他的機構.
合作.
所以就變成了合作的.
在當地的福事.
就建立了四間教會.
那時候.
後來開始有聯會.
宣道會有聯會在英國.
所以.
黎家麗牧師.
香港宣道差會的總幹事.
就邀請我們回來宣道會.
所以我和我太太就回來宣道會.
在英國的地方.
建立宣道會.
所以最後這間.
就叫做南倫敦宣道會.
以前那些就是華人教會.
因為.
以前的機構.
是有不同宗派的人在當中.
所以他們建立的教會.
不可以用宗派的形式.
去建立教會.
希望大家明白.
所以現在的情況就是這樣.
在南倫敦.
這個地方.
已經有十幾年了.
總共有三十四年.
剛才說了.
在南倫敦已經有十幾年了.
在這個地方.
我們很感恩.

$^{241}$神仙差我過去.
香港的弟兄姊妹.
支持我們.
現在我們又服侍香港.
黎的弟兄姊妹.
我們想請大家先看看.
一個影片.
我們就讀了一個影片.
自從2018年回港述職之後.
轉眼間就五年了.
五年日子.
不算太長.
但面對這個世事的變化.
真是風雲萬變.
甚至從南宣的角度來看.
這五年.
也都讓南宣翻天覆地的改變了.
所以我們就定了.
這次的報告主題.
就是以不變應萬變.
不變是古舊十八的救恩.
耶穌基督活主.
生命之道.
萬變就是回應萬變的世代.
所用的工具和方法.
影響南宣最大的變化.
就是香港新移民.
在2021年初.
南宣的移民.
就逐漸多起來了.
當時.
南宣在2021年8月.
開始實體崇拜的時候.
崇拜人數就驟升.
溫布頓所租用的地方.
都不夠用.
特別大人小朋友.
都要找地方住.
於是我們只有四處去找地方.
南宣就為.

$^{281}$神就為南宣.
預備了一個殘障人士.
托管的中心.
由於這個殘障人士.
托管中心.
只是星期一到星期五來工作.
我們能夠租用.
星期一的一整天.
甚至星期六的半天.
在2021年的10月.
我們就開始.
在這個地方實體崇拜.
大壇將中心.
沒有足夠的座位.
讓我們去崇拜.
我們只有自己動手.
去裝配買回來的椅子.
大人小朋友一起去坐.
一下子.
下午一百張的椅子.
就出現在我們眼前.
另一個問題出現了.
殘障中心.
不能夠提供給我們.
擺放椅子和教會雜物的地方.
但允許我們在花園.
建一間小屋來擺放.
於是弟兄們動手了.
建了一間儲物的小屋.
和兩個儲物大櫃.
就好像這樣的.
這些都是新移民.
他們一起合作做的.
一百張椅子.
一下子就坐滿人了.
現在我們第一堂崇拜.
大約有120人.
第二堂崇拜有100人.
我們第二方面的改變.
就是南宣的飲食文化.

$^{321}$我可以從私房菜.
講到群宴.
私房菜是什麼東西呢?.
就是在我們家裡開始.
因為移民來到英國很多.
在星期天和他們在教會.
一起聊天是不容易的.
而且也很難長談.
於是我們就想到.
請他們回家吃飯.
於是我們就定了每星期五.
邀請一兩個家庭.
來我們家裡吃飯.
認識我們.
晚晚弟兄姊妹.
有一個交談的題目之一.
就是你去過牧師的飯沒有?.
你說牧師沒有做什麼菜呢?.
是挺好吃的.
我們租用了殘障中心.
是有廚房的.
而崇拜後也是午膳的時間.
為免到街外買食物.
於是就開展了膳食關注小組.
弟姐妹自發地預備送菜.
讓大家看看南宣每星期的菜單.
在菜單上還寫了教會四個目標.
我們也不少得.
在英國傳統的聖誕大餐.
看看這個菜單.
在團年的時候.
我們做了團年的盆菜.
做了24個盆菜.
雖然飲飲食食.
我們也不會忘記傳福音.
在世界盃決賽那一天.
崇拜完結後就借用禮堂.
用大銀幕看世界盃.
邀請微信主的朋友.
看完之後還有吃火鍋.

$^{361}$還有微信主.
有些微信主出席.
我們還有些留下來.
讓大家看看南宣星期日的飯菜.
真的可以餵飽一百人.
飯堂的情況就是這樣.
我們在裡面.
大家能不能找到我們呢?.
第三方面就是從無到有的教會中心.
南宣擁有教會中心.
真是神莫大的恩典.
在神開門關門當中.
讓弟兄姊妹體會到神的引領.
認定目標.
沉著應戰.
靠著神終於在2021年11月5日.
拿到教會中心的鑰匙.
但另一個問題出現了.
就是我們沒有錢裝修.
登高一副.
弟兄姊妹就起來呼應.
一起動手去為教會中心裝修.
連小朋友都一起出席.
來參加遊圍欄.
將圍欄遊遊和種花.
今天我們能夠享受弟兄姊妹辛勞的成果.
在教會的花園.
我們有聚會.
在教會的室內.
不同的房間.
我們都有不同的聚會.
另一方面想跟大家分享.
教會洗禮的人數增加了.
隨著教會人數增加.
洗禮人數也都上升.
不單是這樣.
我們更有自己的浸池.
不用再去租用其他教會.
想甚麼時候洗禮都可以安排.
南宣的浸池就是這樣.

$^{401}$有自動吹氣.
又可以將熱水加熱.
洗禮弟兄姊妹曾反映這樣說.
水有時還太熱.
這個就是南宣的浸池.
第五方面想跟大家分享的.
就是南宣的智力推動宣教.
鼓勵弟兄姊妹參加段宣.
疫情之前我們去了柬埔寨.
但由於疫情關係我們停辦了兩年.
到2022年再出發去曼城.
今年我們去了馬島.
今年去馬島是很特別的.
為甚麼特別呢?.
就是因為科技更加進步.
我們段宣隊去馬島的時候.
每天他們都將他們的行程.
傳送到教會裡面.
教會負責的弟兄姊妹.
就將它傳發到教會的每一個不同的小組.
於是整個教會上上下下.
就立時知道段宣隊當天做了甚麼.
還有怎樣為段宣隊祈禱.
有些弟兄姊妹就更加反映說.
好像在等待段宣隊的消息.
以致自己覺得自己是在做段宣一樣.
第六方面我們分享.
教會的牧者和領袖人數都是增加的.
我們兩個在2010年來到南宣.
每處崇拜的人數只有雙位數字.
怎麼可能支持聘請傳道人呢?.
但是在2020年的5月.
當時移民潮還沒有開始.
我們為了發展普通話事工.
我們參迦南宣10年之後.
就從中國請了嚴傳道.
來南宣和我們一起的事奉.
發展普通話的事工.
又因為新移民的需要.
在兩年之後2022年的9月.

$^{441}$就聘請了張牧師負責牧養這班新移民.
新移民帶著他們的子女來到教會.
一下子青少年和兒童的數目又大量的增加.
所以我們在辦主學上.
我們都需要很多人手.
主要學現在的學生大約有60多人.
而青少年也有50多人.
在今年我們少年人就參加了.
英國辦的Spring Harvest這個營會.
在營會當中少年人有5個都會決志順主.
而我們決定在2023年今年的10月.
我們就開始第一次少年人的英語崇拜.
最後想和大家分享我們家庭成員的增加.
2020年7月我們第二個孫子在疫情裡面就出世了.
現在給大家看看我們現在的家庭照.
現在南倫敦宣道會有兩堂崇拜.
第一和第二堂.
當我們回去之後.
明年將會增加為四堂崇拜.
因為第二堂是翻譯的.
廣東話翻譯成普通話.
現在不是了.
現在第二堂都是廣東話.
然後是英語崇拜給青少年.
然後是國語崇拜或者普通話崇拜.
給那些在那裡定居的中國來的弟兄姊妹.
所以明年將會是四堂的崇拜.
這個都是神的恩典.
都是為我們需要聘請一位青少年的牧者禱告.
因為要牧養這班青少年.
要中文和英文都要熟悉.
因為他們在那裡長大將會講英文的.
這個牧者英語應該要很流利.
請為我們這個請牧者禱告.
我們也都講一下.
我們宣道會在英國有宣道聯會.
宣道聯會就是說不同的宣道會.
大家聚在一起成為一個聯會.
聯會的牧者彼此都有一個支持幫助.
就好像香港區聯會.

$^{481}$那個就是英國區聯會.
初初是幾對宣教士.
初初是幾對宣教士和一兩對牧者.
自從英國的新移民很多來到之後.
大家都看到宣道會已經增加了.
現在由上至下.
由北至南大概有13間宣道會在英國.
牧養很多都是香港的新移民的教會.
大家看到打星那些.
我們仍然是希望在這些地方再建立教會.
去牧養香港來的新移民.
因為香港來的新移民他們去到有華人教會的.
就參加華人教會.
如果沒有華人教會的他們就參加英國人的教會.
但英國人教會聽都可能沒問題.
但是團契分享就很大問題了.
還有就是有些英國人的教會不是很福音性的.
所以對某些頂子妹他們就是翻一下不翻一下.
又對他們的兒童因為很多都帶有兒童去的.
那些不是很興旺的英國人教會.
也都沒有兒童主教學校那些.
所以他們在這些地方.
所以我們都為這些地方祈禱.
都希望能夠建立一些教會適合他們的需要.
當然其他地方都有的.
但這幾個地方就因為有宣道會的頂子妹.
要求我們在那裡建立教會.
因為他們那裡沒有適合的華人教會.
所以都為這些祈禱.
我們都是聘請一些傳道牧者.
去在這些地方開荒建立教會.
打聲那些的地方.
這個就是近這幾年.
就是2021年之後.
這個2022年的退休會.
我們就多了這麼多童工了.
初初是幾對宣教士加兩位牧者.
就是八個童工.
現在有15位了.
其實2023年我們多了4位.

$^{521}$就是有19位童工了.
很多都是在英國.
大家都知道很多都去了英國了.
在英國加入我們宣道會.
去牧養我們剛才所說的教會.
我們都說一下這個回教徒.
或者穆斯林的時宮.
我們的宣道會.
對當地的穆斯林很有負擔.
首先在曼徹斯特這個地方.
曼徹斯特這個地方.
向這個穆民傳福音.
為什麼呢.
因為那一位童工在曼徹斯特那裡開始.
所以就在那裡開始.
因為也有很多的難民.
和很多新來的穆斯林的人.
也都在曼徹斯特這個地方定居.
所以就在那裡開始.
這個就是一個婦女的小組.
其實婦女最好由女士去接觸他們.
所以我們那一位是男的宣教士.
所以現在我們請了一位女的宣教士.
是香港去的.
一位女的宣教士去接觸.
這班是巴基斯坦的.
也有一些是其他的民族.
女士接觸女士.
向他們傳揚福音.
這個就是查經班.
查經班是家庭式的.
你看到有些是義工.
有些華人是義工.
我們都希望我們的宣導會.
華人宣導會去參與.
在當地的跨文化的福音工作.
所以在不同的查經班裡面.
不同的小組裡面.
他們都有參與.
我們當地的查經.

$^{561}$因為有不同的地方的難民.
所以他們都很清楚.
敘利亞的就還敘利亞.
伊朗的就還伊朗.
伊拉克的就還伊拉克.
庫爾德的就還庫爾德.
大家都明白的.
土耳其的就還土耳其.
因為他們去到.
當然在別人的地方.
但是他們的國家裡面的關係問題.
都會帶到當地英國裡面.
所以我們不可以合得來.
不過信了主的.
有受了洗的.
受了洗的就會帶他們進教會裡面.
就是曼徹斯特宣導會.
就帶他們進教會裡面.
那裡有個講英語的.
叫做International Church.
他們信了主.
很奇怪.
只有神才將他們的.
家壽國限甚麼都好.
才能夠拜得低.
他們不同地方來的.
穆斯林背景的基督徒.
是可以一切的敬拜.
一切的領受聖餐.
所以我們就用這個方式.
然後就將這個福音帶到他們.
福音是和平的福音.
就將這個福音帶到他們當中.
讓他們信了主之後.
有一些信了主之後.
就向自己人傳福音.
所以現在在英國.
很多的穆斯林難民朋友.
仍然需要福音.
這個就是最新的消息.

$^{601}$除了當地的華僑之外.
在這幾年多了這麼多人.
就是新移民有20萬.
我想不止的了.
華人留學生有21萬.
很多都是中國去的.
中國仍然派很多的人.
或者有很多中國的人.
仍然是去英國留學的.
福建人的勞工是27萬.
當地有很多福建人的勞工.
特別是做餐館外賣.
現在很多是給了福建人去做了.
給大家看一下.
我們去傳福音.
仍然有一些異端.
除了耶和華見證人.
摩門教之外.
那些就是英國本地的異端.
這些就是中國帶去的異端.
原來東方閃電.
也趁機向這些新移民.
傳他們的信仰.
這一張就是他們的海報.
教英文.
吸引一些在網上教英文.
吸引一些在網上的.
去參加他們的英文班.
然後再做一個廣告.
告訴他們我們有一個茶經班.
請他們去那個茶經班.
為什麼我們會知道呢.
因為我們有一些人上了英文班.
接著聽到茶經班.
去了茶經班.
發覺說的話不一樣.
所以我們就留意得到.
所以給大家看到.
那邊的實在有很多.
有40\%的新移民是信主的.

$^{641}$是基督徒.
有60\%的是未信主的.
所以對.
有很多是回教會.
也有很多是未回教會.
我們剛才看到.
我們南宣也有很多人因此而信主洗禮.
其實這個和長也是很大的.
最後多謝大家的支持關心.
請繼續為英國的華人.
穆民的福音工作.
禱告和支持.
我們也看到你們在加拿大也有.
我們也是為你們在加拿大的事工禱告.
我們也為Alice的事工禱告.
我的好朋友Alice.
多謝大家.
\newpage



\section{彼得前書 5:1-5}
\label{sec:fjBMO0wHlN4}
\textbf{2023-08-27 早堂《彼此順服的教會》彼得前書 5\_1-5 講員:盧思源姑娘}
\newline
\newline
連結: \href{https://youtube.com/watch?v=fjBMO0wHlN4}{\texttt{https://youtube.com/watch?v=fjBMO0wHlN4}} ~~~~ 語音日期: 2023-08-29
\newline
\newline
\hyperref[sec:xZ5COJOEvOY]{\small{< < < PREV SERMON < < <}}
~
\hyperref[sec:index_chronic]{\small{[返順時目]}}
~
\hyperref[sec:index_scriptual]{\small{[返順卷目]}}
~
\hyperref[sec:9XvZPOyWZdE]{\small{> > > NEXT SERMON > > >}}
\newline
\newline
彼得前書 5:1-5
\newline
\begin{longtable}{cl}
\hline
\hline
章節 & 經文 (和合本修訂版)\\
\hline
5:1 & \begin{tabularx}{0.7\textwidth}{X} 所以，我這同作長老，作基督受苦的證人和分享將來所要顯現的榮耀的人，勉勵在你們中間的長老們： \end{tabularx} \\ \\ \relax
5:2 & \begin{tabularx}{0.7\textwidth}{X} 務要牧養在你們當中神的群羊，按著神的旨意照顧他們，不是出於勉強，而是出於甘心；也不是因為貪財，而是出於樂意。 \end{tabularx} \\ \\ \relax
5:3 & \begin{tabularx}{0.7\textwidth}{X} 不要轄制所託付你們的群羊，而是要作他們的榜樣。 \end{tabularx} \\ \\ \relax
5:4 & \begin{tabularx}{0.7\textwidth}{X} 到了大牧人顯現的時候，你們必得到那永不衰殘、榮耀的冠冕。 \end{tabularx} \\ \\ \relax
5:5 & \begin{tabularx}{0.7\textwidth}{X} 同樣，你們年輕的，要順服年長的。你們大家都要以謙卑當衣服穿上，彼此順服，因為「神抵擋驕傲的人，但賜恩給謙卑的人。」 \end{tabularx} \\ \\
[1ex]
\hline
\hline
\end{longtable}
$^{1}$大英姐妹 早安.
在開始之前 我們也有一個祈禱.
參與了我們再一次來到你的面前.
可以讓我們一同回到殿中 一同去敬拜.
一同去聽你的話語.
一時間只要分享你自己的話語的時候.
求真理的聖靈.
都與我們眾人同在.
無論說的 聽的.
我們都被聖靈你自己親自去感動.
帶領.
願意我們都進入你自己的真理當中.
被你自己的話語去得著提醒.
得著幫助 得著造就.
我們恭敬的 將以下時間交上.
奉主耶穌基督明治祈求.
阿們.
大英姐妹 平安.
在這段時間.
大家如果有留意.
教會其實都有很多不同的活動.
大家有沒有看到外面.
在門外面都貼了一張Banner.
當中都預告了.
教會將會有很多不同的聚會.
我們的主題.
都是圍繞著.
心系吾家齊躍動.
教會就是很希望藉著不同的聚會活動.
可以給過去幾年的時間.
因為疫情的影響.
其實教會很多的事工都停頓了.
今天希望可以重新再躍動起來.
不過當中真的需要我們每一位弟兄姊妹.
在裡面的參與和投入.
不過坦白.
特別是一些在教會比較長時間的弟兄姊妹.
都會明白.
我們在教會裡面.
都有很多不同的人.

$^{41}$我們有不同的背景.
我們有不同的性格.
我們有不同的年齡.
我們有不同的特質.
而不同的人.
對不同的事物.
都有不同的反應.
不過教會.
就是我們這些不同的人.
走在一起.
聖經告訴我們.
教會.
不只是單單這個地方.
這個建築物.
不是一件硬件.
而是由你.
和我.
一起去組成一個教會.
是一個有機體.
如果我們看聖經之前的一些講道.
很多時候我們都會分享到.
當中特別一些經文.
特別是保羅書信都提到.
怎樣去建立.
這個教會的身體.
而今天我們看的.
就是彼得前書.
都鼓勵大家可以翻開聖經去看.
彼得前書當然顧名思義.
就是由彼得所寫的.
他就是耶穌十二門徒.
當中的大弟子.
而彼得寫這本書的時候.
其實他當時.
正值在羅馬的時候.
那時候如果看回一些歷史.
就是尼祿王開始.
大規模地去迫使.
一些基督徒.
而整個羅馬帝國裡面的基督徒.

$^{81}$他們的信仰.
都是受到一些迫使.
甚至折磨.
有些甚至是被殺害.
所以當時.
耶路撒冷教會.
就開始分散到地中海一帶.
而彼得.
他知道這個時.
當時信徒的警方的時候.
他就寫這封信.
就派西拉去安慰他們.
在書信裡面.
如果我們看彼得前書.
捲書是一封信來的.
如果我們由一開始.
去看到第四章的時候.
其實看到彼得.
在當中都.
勸勉教會怎樣.
在神面前可以忠誠地去生活.
也都在裡面說了很多.
從逼迫的角度.
去提出一些勸勉.
直到我們看彼得前書的.
第四章的十七節尾段.
彼得在當中提到.
將來末世神的審判.
是由神的家開始.
是神對他子民的那種磨練.
所以彼得在這裡提醒信徒.
就是要早作預備.
在磨練之前.
都要站立穩.
以至在將來永生的時候.
可以得到那份獎賞.
所以回到我們今天那段經文.
彼得前書第五章第一節.
第一句提到所疑.
這個所疑是在回應.

$^{121}$剛才上文所說.
就是說.
既然審判是從神的家開始的時候.
所以今天在教會裡面的弟兄姊妹.
我們要憤外的去留意.
但是當我們讀今天的經文.
特別是剛才我們主席讀的時候.
一至四節.
大家不知道讀的時候有沒有一種感覺.
覺得好像這段經文.
不是跟普遍大部分的會眾去說.
好像就是跟我們這些.
牧者或者傳道人去說.
好像跟會眾沒有什麼關係.
不知道大家有沒有這種感覺.
坦白說我起初預備的時候.
我都有一點這種感覺.
因為這段經文其實是我們按照教會.
我們每一個星期有一個講道表去講的.
當我去預備的時候.
我起初都有一種感覺.
但是當我慢慢再去細讀.
去預備的過程裡面.
我發現這段經文.
不只是單單跟我們這些傳道人.
長老做牧羊的工作的人去說.
而神的話.
是對我們每一個都有提醒的.
相信我們在當中都是可以有學習.
經文在第五章第一節那裡這樣說.
所以我這同作長老.
作基督受苦的證人.
和分享將來所要顯現榮耀的人.
勉勵在你們中間的長老們.
在這裡我們見到彼得稱呼自己為什麼.
同作長老.
但是剛才我也說了.
彼得其實他的身份.
他是一個使徒.
他這樣稱呼自己的時候.

$^{161}$我們見到彼得他不是.
好像以一個上級和下屬的身份去說.
相反他勸勉當時教會的弟兄姊妹.
就是他站在同作長老的地位去說.
大家去想想長老.
究竟什麼是長老呢.
長老是不是就是年紀大的就是叫長老呢.
原本是這個意思就是年長.
但是我們看到在聖經裡.
在猶太的社會.
自從摩西我們讀民數記的時候.
摩西他立了70位長老去幫忙管理事務開始.
當時猶太的社會開始出現了長老.
沒錯的.
當中的人可能會比較年長一點.
或者德高望重一點的人.
當時這些人很多時候都是成為猶太會堂的領袖.
但是如果我們讀剛才說彼得這個時間的時候.
他是一個初期教會的時間.
我們知道長老這個職事.
其實是按恩賜去選出來和按立的.
是可以成為當中那時候教會的領袖.
所以其實這裡說的長老.
就直接的說不是和年齡有直接的關係.
如果我們引申下來我們現在的處境.
長老當然剛才說的可能是我們這些牧者傳道人.
做牧羊工作的弟兄姊妹.
但是也不限於這些的職事.
而是關於那些在教會裡.
受託照顧教會弟兄姊妹的人.
一些有份去做帶領的人.
例如在我們教會當中.
是怎樣呢.
譬如堂董執事部長.
可能團契的團長 導師.
甚或主一學的師班的導師.
或者你在兒童幼兒青少年牧區的導師.
甚或可能是一些小組上的小組長.
因為我們知道昔日保羅的情況.
他是沒有這些分工的.

$^{201}$不會像我們現在這樣分得那麼仔細.
不是保羅 sorry.
彼得在昔日是沒有這些分工的.
但是彼得在這段話.
是和教會有份做帶領角色的人去說的.
在這裡彼得對長老這樣說.
務要牧養在你們當中神的群羊.
按著神的旨意照管他們.
不是出於勉強而是出於甘心.
也不是因為貪財而是出於樂意.
不要合際所託付你們的群羊.
而是要作他們的榜樣.
在《聖經》裡我們很多時候都會看到.
都會將一些領袖作為一個牧人.
或者將一些人作為羊.
如果我們記得在《約翰福音》21章的時候.
昔日主耶穌復活之後.
其實耶穌都在加里里海和彼得去顯現.
也都三次確定彼得要牧養主的羊.
這個是上帝給他的即使.
所以當我們今天去看彼得前書的時候.
彼得在這裡似乎又重覆神對他.
主耶穌對他的呼召對他的命令.
不過現在他就和一些童作長老都是這樣說.
他說要以牧養的即是為念.
但是這個牧養是有特定的要求.
就是要按著神的旨意.
就是要按著神的要求原則標準去看管羊群.
而不是按著自己的意思.
不是按著自己的喜好去照顧.
接著我們下去看這段經文的時候.
有三個對比.
他說不是出於勉強而是出於甘心.
勉強的意思就好像一個強迫被迫的意思.
好像和甘心相反.
剛才說過昔日的長老.
他剛才說是被教會推舉出來或者選立出來.
或者當時他自己個人的意願會被忽略.
但是給提醒這些的事奉不是出於勉強.
如果引用我們現在的處境.

$^{241}$今日我們弟兄姊妹.
我們都不是應該因為一些環境.
或者因為別人勉強的情況下.
我們去承擔教會的一些的職事.
而是我們需要甘心樂意地去事奉.
有時我間中和弟兄姊妹去聊天.
間中都會聽到他們都參與教會的一些的事奉.
不過在他分享當中.
有時分享到心底裡面好像有少少不是真的那種樂意.
很多時間中會聽到.
你會參與那個事奉.
那為什麼呢.
他就說可能原因是某某可能要我做.
那我不可以私拒絕.
那就做吧.
沒錯.
在表面我們見到弟兄姊妹參與事奉是一件好的事.
但是坦白.
如果在心底裡面不是真正去樂意的時候.
其實是不健康的.
在聖經裡面我們見到.
就是我們在神的一切事工裡面.
我們應該是甘心樂意地去做.
因為我們知道我們有份去參與神國的事的時候.
這個是我們的福氣.
我們應該是有那份積極.
有那份主動的甘心.
而不是好像一種被迫或者被勉強的心態.
而接著下去.
彼得繼續在說.
不是因為貪財而是出於樂意.
貪財呢.
這裡不是說貪錢的意思.
意思是指.
喜歡以一些不誠實的方法去賺取一些利益.
記得在以前曾經聽過.
不是在這間教會的.
在我以前曾經聽過一些教會.
他說有些爸爸媽媽.
他們很想他們的兒女可以入讀他們心儀的一些基督教學校.

$^{281}$但是原來如果父母或者家人有回教會.
或者有參與教會的活動的時候.
或者他們在基督教學校.
原來對於他們入讀的學校是有份的.
所以就加入了教會.
我不肯定這件事是不是真確.
但是如果.
如果我們參與教會的事奉.
或者我們的動機是為了自己利益的時候.
這種心態其實都是在聖經裡面告訴我們.
原來是不對的.
而在第三個的對比.
其實都是有一個警惕的.
他就說不要合制所託付你們的群若.
而是要作他們的榜樣.
合制的意思.
很明顯就是以一些權柄控制人.
很多時候都是在說一些君王.
他怎樣去治理他的百姓.
有一份絕對的權威.
但是如果我們在福音書裡面看.
我們知道耶穌的教導都不是這樣的.
耶穌的教導就是說我們千萬不要像世上的君王那樣.
以合制為治理之道.
反而要成為僕人謙卑的服侍.
記得之前我還沒有讀神學.
我仍然在工作的時候.
公司有一位的director.
我當時是做商場管理.
簡單來說就是商場所有的東西都要打理整個商場.
裡面所有東西都要負責.
偶爾你知道要巡一下場.
就是走一下商場.
看看有什麼不妥.
其實大家走在街上去商場最明顯看到不妥的就是清潔.
清潔不妥的時候.
大家怎樣去看一個商場覺得它乾不乾淨呢.
我相信就是洗手間.
大家都可能覺得那個商場乾不乾淨.
用他的洗手間就可以評定這個商場乾不乾淨.

$^{321}$當然剛才說那個director.
他作為一個公司的director當然不用去巡洗手間.
他在head office坐在辦公室工作.
不過如果他有機會來到商場裡面.
又或者聽他自己放假的時間.
他都會去巡一下不同的商場.
因為我們有很多不同的商場.
就巡一下不同的商場.
其實作為他這個身份.
作為一個公司的director.
他絕對是可以去叫下屬去做.
或者叫下屬做完告訴他究竟有什麼問題.
他不需要自己去做.
但是就是因為他這樣做的時候.
他給一些員工去明白.
什麼是真正管理一個商場.
所以當時他在公司裡面很多下屬都很服他.
甚至他的工作態度.
他的影響都很影響我自己的工作態度.
我在想一個非基督徒的人.
當他以身作則的時候.
尚且可以影響這麼多的人.
所以當我去讀這段經文的時候.
我體會到為什麼彼得這樣去提醒我們.
他提醒我們.
當神給我們可以做帶領的角色.
給我們權柄的時候.
不是給我們這些權柄去控制.
去控制人.
相反我們知道這些權柄是神給我們.
所以要學習耶穌的榜樣.
效法耶穌基督.
而耶穌就是那位很謙卑成為僕人的樣式.
如果我們看到第四節的時候.
他說到了大目人顯現的時候.
你們必得到那永不需查榮耀的冠冕.
大目人.
我們剛才說了我們都知道是在說耶穌基督.
所以今天無論我們在哪一個崗位當中.
最後我們每一個其實都是屬於主耶穌的若.

$^{361}$祂是我們那位大目人.
祂是我們那位大牧者.
當他讀這句經文的時候.
給我們其實是一個很大的鼓勵和提醒.
提醒就是今天上帝給我們的羊或者人.
給我們的事奉.
其實都是託管的.
重要的是我們怎樣將這些人.
這些事奉帶到上帝的身邊.
而在鼓勵.
這句話給我們一個鼓勵的時候.
就是當我們可能事奉到很累.
或者真的牧羊到很累的時候.
就可以跟神說.
其實這些都是屬於上帝.
上帝自己會去負責.
引申我們剛才說在教會可能我們不同的人.
我們有不同的崗位.
但是都是求主給我們一份審查.
其實今天我們所事奉的.
我們所做的.
都是屬於主.
盼望當主耶穌回來.
就是大目人去顯現的時候.
盼望我們每一個都可以被主耶穌稱讚.
我們是那位又良善又忠心的僕人.
但是可能坐在這裡的大兄姐妹.
都會有些心裡會想.
我在教會都仍然沒有一些帶領的角色.
又或者剛才盧姑娘你提的崗位.
或者事奉.
我都沒有.
剛才那段經文.
好像不是很關自己事.
但是知道.
雖然可能今天.
在教會你沒有這個崗位.
沒有這個角色.
但是或許未來.
上帝有機會會放你在這些崗位當中.

$^{401}$這些我們都不知道.
又或者就算是.
我真的沒有了.
我真的未來都看不到自己有這個崗位.
我只是一個信徒來的.
但是這段經文.
都不是對你沒有用的.
我們很重要.
很重要就是我們可以真的為著.
教會的領袖我們去禱告.
求主幫助.
幫助他們教會的領袖.
可以以聖經的教導去帶領教會.
大家不要小看這些禱告.
我相信.
如果我們每一個的.
弟兄姊妹.
我們都為教會裡面的帶領者去禱告的時候.
我們的教會.
是更加.
更加可以.
行在神的心意當中.
回到聖經.
彼得和長老講完話.
剛才一至四節就是對象.
就是一班長老.
去到第五節的上半節.
他這樣說.
同樣.
你們年輕的.
要信服年長的.
聖經這些連接詞大家去看的時候.
其實很有趣的.
同樣.
的意思是什麼.
同樣就是表示.
現在以下所講的東西.
就是和上面所講的東西.
有關係.
但是上面講的東西明明在講.

$^{441}$長老的職責.
兩者都好像沒有邏輯的關聯.
所以我們知道這個同樣.
其實是在回應剛才說在第四章尾尾那裡關於.
末日的審判.
就是末日.
那個審判是由神的家開始.
所以剛才說一至四節.
就是提醒教會的領袖.
要時刻的警醒.
去盡牧養的責任.
而在這裡.
就是提醒一些其他的會眾.
他說你們年輕的.
剛才說.
年老年長的不代表.
和年齡有直接關係.
所以這裡.
年輕的也不是在講實際的年齡的年輕.
可能大家已經說我已經不年輕了.
所以這裡不是在講.
年齡上的東西.
剛才說長老就是可能.
當時是按恩賜而被選立的人.
大家想像按恩賜被選立的.
可能他們在相信相對地.
靈命上都是比較成熟.
所以這裡說的年輕.
很大機會就是可能在說一些.
靈命上比較年輕的.
也是在說教會裡面除了一些長老之外.
其他的人.
可以這樣說.
要信服年長.
信服.
可能在教會很久大家都聽到.
信服教導很多.
但是有沒有想過.
信服.
其實和屈服是不同的.

$^{481}$信服和屈服有時很相似.
雖然他們都是在描述.
在某種壓力或者一些權威面前的一些行為和態度.
但是他們含義上有少許的分別.
當我們說我們信服的時候.
我們通常就是在說.
我們信服服從某種的權威.
指令命令規則.
但是這些很多時候都是出於.
自願和你自己合理認同的一種行為.
而這種服從.
是你基於對那件事的尊重信任忠誠.
所以當我們說我們一個人信服的時候.
通常會照著指示要求方式去行動.
而且在過程裡面.
你是會感到滿足.
也都覺得有目的感.
不過屈服是什麼呢.
屈服很多時候就是在說面對外面一些的壓力.
困難那種的退讓.
那種的屈從.
去暗示就好像被迫或者無法抗拒的情況下.
因為權制壓力威脅或者其他的因素所導致.
當一個人屈服的時候.
他通常就要放棄自己的意願.
自己的立場自己的權力.
接受一些他們不可以改變不可以反抗的情況.
簡單總括來說就是.
信服就強調一個很主動和很自願的一個行為.
而屈服就是暗示一種被迫和一種被動的情況.
但聖經在這裡說的是信服.
在這裡說的是你們年輕的不是因為那些年長的.
所以迫於無奈的去屈服.
有時候我們在工作的環境在社會的環境很多時候都會聽到.
你做大的.
你說是.
你是老闆.
你說什麼都要聽的.
沒錯有時候當我們說這些話的時候.
心底裡我們可能會覺得.

$^{521}$因為他高級一點.
因為他高過我大過我.
厲害一點所以就要聽.
但是心底裡這個是不是真的服呢.
可能這些只不過是屈服而不是信服.
不過彼得在這裡說的信服.
就不是那種麻木的去信服.
而是說當有前設的就是當如果.
基於如果那些長老那些領袖他是盡忠職守.
按神的旨意去帶領的時候.
而他又不是核製羊群而是做廣羊的時候.
我們其實可以假定到也可以想到.
這個信服其實不是一個很沉重的擔子.
對於被帶的其實就是很自然.
很應該這樣全心去信服去接受.
而去聽從那份的領導.
如果我們繼續看這段經文的時候.
第五節下面這樣說.
你們大家都要謙卑.
當衣服穿上彼此信服.
因為神抵擋驕傲的人.
但次因給謙卑的人.
這裡很快彼得又轉了對象.
剛才就跟長老們說.
然後就跟年輕的說.
現在突然間就跟大家說.
對象又轉了.
這次是對所有人說.
他說你們大家包括什麼.
無論你是年長的年幼的.
在教會的每一個人.
我們都要以謙卑蓄夭.
蓄夭是什麼意思呢.
蓄夭原來當時的意思就好像你工作的時候.
穿一條圍裙.
然後你穿一條圍裙是什麼.
就是防止弄髒衣服.
所以蓄夭的意思有一個意思.
就好像穿上的意思.
穿上的意思.

$^{561}$其實如果我們在聖經裡面.
很多時候都會用穿衣服去形容.
我們信徒要過新的生活.
譬如在耳窟所書.
就是我們要脫去舊人.
穿上新人.
在加拉太書就是我們要披戴主耶穌基督.
很多時候聖經都會這樣去形容.
我們要穿上一些東西.
而彼得在這裡說.
我們要穿上什麼.
穿上謙卑.
我們每個人都要放下自恃的心態.
學習別人比自己強.
我們學習以謙卑彼此的去相待.
之後彼得說了原因為什麼要這樣做.
他很直接.
就是因為神抵擋驕傲的人.
但是次因比謙卑的人.
這裡彼得他引用的是.
《箴言》七十四譯本三章三十四節的內容.
這個也是一個很重要的真理.
阻擋即是抵擋的意思.
即是敵對.
這裡說的是神敵對這些人.
驕傲的人.
即是這些是神視為敵人.
即是聖經說驕傲的人.
上帝去敵對他的時候.
他的道路是不亨通.
相反神是次因比謙卑的人.
所以彼得為什麼在這裡去說.
弟兄姊妹我們為什麼要以彼此去相待.
因為這個很清楚說明了.
原來神做事工作的原則.
就是他抵擋驕傲的人.
但是次因比謙卑的人.
弟兄姊妹一開始我們說.
我們教會都有很多不同的人.
我們有不同的背景喜好.

$^{601}$坦白有時分歧是在所難免的.
有時我們不得不承認.
我們人很多時都是喜歡別人聽自己說話.
這個情況其實不一定在教會裡面的群體出現.
大家可能想起在工作的場景.
即是上司對下屬.
同事之間.
又或者在家裡面.
父母對子女.
又或者配偶丈夫和妻子的關係.
甚或可能男女朋友的關係.
又或者朋友之間的關係.
其實這些情況很多時都會出現.
有時我們很容易都會覺得.
我們自己看的東西才是對的.
我們很多時都會覺得.
希望別人認同自己.
但是當我們如果堅持我們自己的東西是對的時候.
有時有時.
有時可能紛爭和分歧就會出現.
不過彼得在這裡給了我們一個很重要的原則.
就是無論發生什麼事.
如果我們每一個人都是穿上謙卑的衣服.
我們都是穿上謙卑的衣服的時候.
即是我們願意放下我們自己的時候.
當我們願意這樣做.
有時那些分歧.
那些紛爭.
那些問題.
可能無形中就會得到解決.
最後今天聖經的教導教導我們.
無論我們在什麼崗位.
我們都學習彼此的信服.
或許可能現在剛才說在那些崗位當中.
你是有份做帶領的.
是年長的.
但是可能你雖然年長.
但是對於某些人來說你都是年輕的.
但又或者你現在真的很年輕的.
很年幼的.

$^{641}$但始終有一天你都成為年長的.
願意我們都學習彼此信服的功課.
我們真的同心合一地去同行.
教會才真正服從神的真理.
我們真正的成為一間傳言對主信服的教會.
亦都是一間彼此信服的教會.
願意上帝的說話.
都成為我們彼此的提醒和勉勵.
我們一同去祈禱.
主讓你藉著昔日彼得所教導.
你提醒主耶穌.
是呀.
你這位大牧人.
你有一天會再回來.
當你再回來的時候.
你提醒神.
你會審判.
而審判是由神的家開始.
求你幫助我們.
幫助我們今天藉著聖經.
你自己的話語的時候.
無論我們今天在哪一個崗位當中.
或許我們在教會裡面.
我們是有份做帶領的.
或許我們在教會裡面.
我們是一個信徒.
但無論我們在哪一個崗位當中.
我們都學習彼此的去信服.
我們學習彼此的去幫助.
因為神你說.
你敵黨驕傲的人.
相反你恃恩比謙卑的人.
求你幫助我們.
幫助我們每一個弟兄姊妹.
真的我們好像正如剛才聖經裡面所說.
我們穿上謙卑的衣服.
以至無論我們在教會群體.
我們和弟兄姊妹相處.
甚或無論在我們的工作家庭裡面.
我們真的成為那個榜樣.

$^{681}$成為主耶穌的樣式.
讓更多人認識主你.
所以主也求你幫助我們.
願意我們在你的話語裡面.
我們真的去實踐出來.
以致我們成為一個合你心意的兒女.
天父多謝你聽我們祈禱.
奉主耶穌基督名字祈求.
阿們.
\newpage



\section{以斯帖記}
\label{sec:9XvZPOyWZdE}
\textbf{2023-08-31 周四晚查經班《被擄歸回歷史研讀(23)以斯帖記》 講員:曾裔貴牧師}
\newline
\newline
連結: \href{https://youtube.com/watch?v=9XvZPOyWZdE}{\texttt{https://youtube.com/watch?v=9XvZPOyWZdE}} ~~~~ 語音日期: 2023-09-05
\newline
\newline
\hyperref[sec:fjBMO0wHlN4]{\small{< < < PREV SERMON < < <}}
~
\hyperref[sec:index_chronic]{\small{[返順時目]}}
~
\hyperref[sec:index_scriptual]{\small{[返順卷目]}}
~
\hyperref[sec:ay_HsjJTAHo]{\small{> > > NEXT SERMON > > >}}
\newline
\newline
$^{1}$今晚我們又再一次能聚集,一起有一個簡單的祈禱.
多謝天父讓我們今天晚上,雖然很快到凌晨天氣也會非常惡劣.
但是我們能夠在禮拜堂這裡去查考聖經.
幫助我們能夠明白經文,當然有很多字面我們都明白了.
但是讓我們看到作者這個敘事的歷史的.
可能是猶太人,可能是以色列.
都幫助我們能夠從其中去明白經文背後.
想要承載的那個時代的一個需要.
也都讓我們看到在主耶穌的救恩裡面.
我們都是德性的,我們都是屬神的子民.
幫助我們對這個身份的一個認同.
還有一份很強的榮耀感.
我們等候在你的面前,禱告奉耶穌基督的名,阿們.
上個禮拜我們就說了有一個三天的禁食禱告.
意味著大家都感覺到那個事態嚴重.
必須要無論是個人的安危.
或者是整個猶大民族的生死存亡.
都是需要當下去倚靠神.
雖然沒有提到向誰祈禱.
假設已經知道猶大人一定是向神禱告.
一定是向著耶路撒冷的主去祈禱.
也都很肯定他們是知道.
耶路撒冷已經被重建有聖殿.
因為已經過了一段長時間的被擄歸回.
在這個時候是不是孤軍作戰.
127個省的猶大人.
在12月13日就要面對滅族的危機.
現在是正月,還有12個月去處理這個問題.
現在就沒抵改,一直禁食祈禱.
一直尋求唯一的辦法.
就是他的繼女,以斯帖來拯救自己的民族.
以斯帖也都輾轉透過太監的傳譯.
來跟叔父說.
他已經30天沒有去皇宮被召見.
可能皇帝公務繁忙.
也都可能暗示皇帝不再像起初那樣.
對以斯帖有一個愛.
也都不出奇的.
但是無論如何就是三日三夜的禁食.
全個民族.

$^{41}$當然未必去到其他省份的猶大人.
因為那個通訊不是今天的通訊.
不是今天的WhatsApp或者Zoom.
所以是有困難的.
但是書山城裡面的猶大人.
他們得到那個消息來祈禱.
這個就肯定了.
因為四章的17節.
於是沒抵改.
照以斯帖一切所吩咐的就去行.
即是通報招聚書山城所有的猶大人.
為他準備要冒險見王.
來到去禁食三晝三夜.
不吃不喝.
即是連飲水都不可以.
是極度的危急.
極度的一個需要.
來到去放下最基本的飲食.
這些偶然間人生都有的.
但是當然不可以變成一個常規.
但是我們需要有這個操練.
第五章.
第三日.
即是真的禁食三晝三夜.
第三日以斯帖穿上朝服.
進王宮的內院對殿站.
王在殿裡坐在寶座上.
對著殿門.
王見皇后以斯帖站在院內.
就施恩於她.
向她伸出手中的金像.
以斯帖便向前摸像頭.
王對她說.
皇后以斯帖啊.
你要什麼你求什麼.
就是國的一半也必賜給你.
這句話出現了三次.
即是說原來.
克錘老王並未對以斯帖.
來到他的愛慕冷淡了.

$^{81}$反而可能真是公務的原因.
令到他有三十日都沒有召見.
這一位皇后.
接著第四節.
以斯帖說.
王若以為美.
就請王戴著哈曼金日.
赴我所預備的筵席.
王說.
叫哈曼速速照以斯帖的話去行.
於是王戴著哈曼.
即是兩個人.
赴以斯帖所預備的筵席.
當然是指大臣唯有哈曼.
皇帝的隨從一定有.
那些太監.
或者那些.
現在叫保鏢.
一定會隨行.
接著就是說.
在酒席的言前.
王又再問以斯帖.
這個時候哈曼在場.
你要什麼.
我必賜給你.
你求什麼.
就是國的一半.
也必為你成就.
可以愛這位太太皇后.
愛到這樣.
即是開金口.
甚至乎是有這樣的offer給皇后.
為什麼會令到皇得到這麼大的信任愛慕.
到底是以斯帖用美貌.
還是用他的智慧.
還是用他其他的一些.
譬如一些法術.
你知道看那些語場片都有.
通常那些西宮都會法術.
但是令到皇可以這樣對自己.

$^{121}$這個都是一個很大的成就.
相信以斯帖一定不是用一些古靈精怪的手段.
是真的用他的美德.
用他對神的敬畏.
來換取到皇帝丈夫的歡心.
接著就是第七節.
以斯帖回答說.
我有所要.
我有所求.
我若在皇眼前蒙恩.
皇若願意賜我所要的.
准我所求的.
就請皇戴著哈曼.
再付我所要預備的筵席.
好像吊著皇帝的胃口.
再付即是第二天.
再有一個筵席.
都不把握機會.
已經說了兩次.
亦都是走席的歡暢的時間.
這個時候為何不直接了當地說呢.
其實明天說都是在哈曼的在場.
哈曼不是說現在退席.
現在要避席.
要他走開.
不是的.
是當面.
戳口戳面地說.
就是這個仇人哈曼.
要滅絕我們這個民族和我.
為何今晚和明晚又有甚麼不同呢.
這個真的沒有辦法去明白.
到底是否這個.
是以斯帖他自己的一些智慧.
令到皇有一個這樣的懸念呢.
這些真的不知道.
唯一我們知道的就是第六章.
皇喝完之後.
可能飲保利飲得太多.
睡不著.

$^{161}$或者飲咖啡.
當然說笑.
那時候哪有呢.
可能就是太開心.
太過歡暢.
睡不著.
而他要拿歷史來看.
那就很奇怪.
所以如果當晚說了.
那個後果.
可能皇帝會調查.
又或者是否以斯帖.
在這裡來重傷.
他現在這一位得寵的大臣哈曼呢.
兩個局勢都很難說的.
所以我們真的很難想像.
以斯帖這麼危急.
這麼關乎到自己民族的命運.
現在已經抓緊了機會可以見皇.
皇亦都伸出金像.
皇亦都現在來開金口.
你有什麼所要.
就是國的一半我都給你.
說了兩次了.
即是說皇帝現在其實很開心.
那個狀態不是一個很憤怒.
很六神無主的情況.
那為什麼不說呢.
接著就是第九節.
那日哈曼心中快樂.
因為很正常.
你不說.
皇帝又帶我去筵席.
明天.
我還沒讀完下面那一段.
就是第八節的下半節.
就是請皇帶著哈曼.
再赴我所要預備的筵席.
明日我必定照皇所問的說明.
即是說這樣的懸念.

$^{201}$就令到哈曼有事要做.
因為她真是興高采烈.
因為她真是感覺自己.
不單止被皇帝賞識.
連皇后都賞識她.
只是她和皇帝.
兩個人赴宮廷的飲宴.
而且還是私人的宮廷飲宴.
即是說是很特選的身份.
才可以出席.
不是一般性的.
好像國慶的例行酒會.
每個人都可以去到.
而是只有一個就是哈曼.
所以有這一段的描述.
一點都不奇怪.
因為她的人本身就是這樣的.
見高拜見低踩.
見到這個沒底蓋.
不跪不拜.
她就很沒有意思.
然後就很想滅絕他.
滅絕他原來都是很小事.
即是當他是一隻蟻.
一踩就死.
但是如果能夠將她這個民族.
都連根拔起.
就令到哈曼更加感覺自己.
那種權傾朝野.
這就是哈曼的計謀.
所以當以斯帖.
又是沒有讓人知道.
他是猶大人的身份.
哈曼都不知道.
哈曼要滅絕的.
原來就是以斯帖的民族.
很多的時間點.
到底是巧合.
還是他們這一班猶大人.
在認真的禱告裡面.

$^{241}$來守住自己的身份.
不是要特意隱瞞.
而是非到必須的時候.
就不公開自己的身份.
某程度上這些都是智慧.
這些都是一些.
雖然不能叫他很深的誠腐.
或者內心有一些心術等等.
心計.
而是有些成熟的人.
不隨便.
什麼都告訴別人.
這些都是我在人生裡面.
相信在座或者在網上的弟兄姊妹.
你做人這麼久.
真正成熟的人.
都不是隨便這樣.
有碗話碗 有碟話碟.
記住不是那些負面的.
很深的誠腐.
心計很重.
不是那些 是很正面的一個人.
有時很謙厚的.
不隨便說自己.
很張揚自己.
我們今天香港人說的.
很show off自己.
什麼都告訴別人.
有多少個頭銜.
有什麼功績.
其實那些不需要說.
不是假謙卑.
而是真是一個成熟的人.
不是刻意要包裝.
而是有一個.
真是好像一種.
個人內心的一個修養.
真是不知道以斯帖.
為什麼有這樣的成熟.
有這樣的智慧的眼光.

$^{281}$停留一晚.
不說.
哈曼就做事了.
哈曼第九節.
那日哈曼心中快樂.
歡歡喜喜的出來.
但見沒底改.
在朝門不站起來.
連身也不動.
就滿心牢牢沒底改.
這麼奇妙.
出朝門也見到沒底改.
為什麼呢.
睜眼睜鼻.
沒底改.
其實也有官職.
因為以斯帖已經給他做官.
但為什麼仍然要.
這樣在哈曼面前出現.
刻意挑起哈曼的憤怒.
這些我們不太知道.
當時的場景是怎樣.
但唯一我們知道的.
就是激怒哈曼.
讓他.
去到一種.
完全的失控.
這種失控.
就是將他極其的驕傲.
即是說.
他內心的品性.
完全暴露出來.
使他好像不斷.
在積蓄神的憤怒.
使到神.
如果要報應他的時候.
其實真是報應.
第二晚.
他就死於非命.
自己所做的刑具.

$^{321}$行刑的刑具.
五杖.
可能是中國人的杖.
不知道是否尺多.
很高的一個木架.
要來掛這個.
沒底改.
成為他自己.
這樣離開世界的下場.
這就告訴我們.
他見到這個時候.
滿心惱怒沒底改.
哈曼暫且忍耐回家.
叫人請他朋友.
和他妻子.
塞利斯來.
作者就很巧妙地.
將前後兩次的對比.
這次他就很開心.
很興奮.
回去見人就說.
見到太太塞利斯.
就告訴他.
皇后只是指名道姓.
只是叫我.
唯一的大臣.
去負這個賢職.
第二晚.
他又回家.
不過垂頭喪氣.
又見到太太塞利斯.
和內閣的.
他自己也有些下人.
謀事.
又說回他的心情.
這次說的是興高采烈的心情.
明晚.
是另一個心情.
你可以想像.
作者也很巧妙地.

$^{361}$捕捉到.
每個人不同處境.
帶來的.
內心世界.
第十一節.
哈曼將他父后的榮耀.
眾多的兒女.
和皇台舉他.
使他超乎首領.
臣僕之上.
都述說給他們聽.
有些小偷.
在他面前.
因為現在很開心.
開心到一個地步.
今晚都睡不著.
不過不是看歷史.
可能更加瘋狂地.
自己的圍內.
吃火鍋,飲宴.
一定是這樣.
因為實在太興奮.
那些小偷的臣子.
沒有警告他.
反而是巴結他.
讓他們知道.
第十二節.
哈曼又說.
不單止是這樣.
我現在權傾一時.
權傾朝野.
不單止是這樣.
明晚還有一次的筵席.
哈曼又說.
皇后已私帖預備筵席.
除了我之外.
不許別人隨皇附席.
明日皇后又請我.
隨皇附席.
改坐在朝門.

$^{401}$須有這一切榮耀.
也與我無益.
這樣的對比.
明明.
那個榮耀.
大到一個地步.
好像發光一樣.
一個人.
怎會令到你完全無益.
怎會令到你的心情.
突然之間冷靜下來.
沒有理由.
你很開心很興奮.
為何你會這樣.
假設.
又是開玩笑.
鄧全就確認了.
按目通過了.
已經開始籌備.
但突然之間.
走過.
感受有個職工.
不經意掃地.
潑水潑到他.
他很不開心.
對職工很生氣.
按目通過了.
與我無關.
一點也不開心.
根本不會.
是兩回事.
級數差很遠.
但為何哈曼.
可以小氣到一個地步.
要整個心胸.
窄到一支針也進不去.
為何人可以這麼恐怖.
你是一個這麼有見識.
國家的國度.
這麼大.

$^{441}$127個省.
從印度去到很遠的地方.
都是波斯帝國.
最如日中天的時候.
你是那個大神.
你不是那些.
地中海.
或者是.
那些小的島.
那些不知道什麼島.
什麼名字的.
什麼巴布亞.
新加萊亞.
那些島的島主.
是美國的國家.
元首的重要.
謀士.
你想想差多遠.
但哈曼居然.
一個人都容不下.
他的內心.
要置他於死地.
才能消除他的.
仇恨.
頂尖會仇恨.
會令我們跌鑊.
一定會降歸.
你心裡面.
有仇恨.
你肯定不會再見到神的面.
這是很嚴重的罪.
人很容易.
會有這種.
心結.
很容易會.
內心裡面.
積蓄這些憤怒.
所以.
新約聖經差不多每一卷書.
講完那些教義的真理.

$^{481}$之後都勸我們.
要在一個新生命裡面.
釋放自己.
在靈性上.
要追求.
與眾人和睦.
寬恕.
解除.
那些冤仇.
這些.
是很重要的真理.
最終是幫我們自己.
因為.
箴言說的.
疾度就是骨中的扭爛.
那些骨都枯.
古代的人.
他們不太懂.
哪有X光.
哪有MRI.
哪會診斷到你的骨的密度有多高.
如果有些人年紀大.
骨枯就會知道.
你的密度有多高.
古代哪有這些.
但他們說的不是物質的骨.
是心靈裡面.
所以哈曼.
他有.
但這樣描述不是過份.
那又真的很貼切.
還有萬貫家財.
還有權經朝野的權力.
但是.
這個人的內心.
不配得到這些尊榮.
因為他內心有極大的.
對這個末底改的成見.
仇恨.
而遷怒於.

$^{521}$他是猶大族.
再說就是.
原來他的太公太公的祖先.
是亞馬力.
這個支派.
是與神為敵.
所以我們就完全明白.
那個原因.
所以與神為敵.
真是.
我們應該要用神給我們的力量去抵抗.
所以.
《已不鎖書》六章的全副軍裝.
抵擋魔鬼.
就不是靠血氣之勇.
要穿起神所賜的全副軍裝.
好在磨爛的日子.
抵擋仇敵.
所以要站穩了.
以真理.
當作帶子蓄腰.
接著就有其他的裝備.
其實都是在說.
真理.
人愛.
福音.
和神的聖靈的寶劍.
就是神的道.
來去抵抗魔鬼.
抵抗魔鬼.
所以都是護衛著.
我們自己的新心靈.
用救恩.
使我們自己成為一個新造的人.
所以這個哈曼.
是代表著.
一個與神為敵的人.
一個極致的人.
在他之下.
無論是沒有他那麼強的權力.

$^{561}$或者一般平民百姓.
一樣有仇恨.
那就要.
來到去.
求神幫助.
每個人都有不同.
自己人生的跑道.
總有些人你喜歡.
有些人你不喜歡.
但聖經就是要我們.
盡力與眾人和睦.
若是能行.
總要盡力與人和睦.
羅馬書十二章.
不記得.
二十多節.
總要這樣做.
這就是保羅教導.
也是主耶穌教.
你的仇敵.
你就要和他和好.
這是相當重要.
哈曼.
從來沒有尋找.
和睦.
寬恕的機會.
也不是真的得罪他.
只不過.
沒有在他面前.
下跪.
所以.
茶餐就出現.
如果你的心.
喜歡聽一些.
公為你的說話.
這些說話.
在你現在成功.
這些說話就出現.
在你失敗的時候.
也不會告訴你.

$^{601}$有危機.
在你成功的時候.
你是皇帝身邊的紅人.
當然有這麼高的帽子.
有這麼惡毒的.
計謀獻給你.
後果是你負責.
不是我們這些賊仔負責.
所以.
第十四節.
他的妻西莉絲和他一切的朋友.
之前說是.
那些謀士.
原來是一些蠶蟲師爺.
因為你是米飯班主.
自然說你.
喜歡聽的說話.
所以如果有些下屬.
有些同工.
真的說一些.
你不喜歡聽的說話.
而那些說話又真的有.
不同的角度.
其實這些是很好的朋友.
因為他真的將.
你沒有想過的.
看不到的角度.
告訴你.
而不是就著你自己喜歡的.
來公為你.
這些是很難得的.
真正的同工.
真正的朋友.
例如立一個五丈高的木架.
明早求王將.
末底改掛在旗上.
然後你可以.
歡歡喜喜的去.
隨王赴宴.
哈曼爾者說為尾.

$^{641}$就叫人.
作了木架.
飲完.
出門口見到末底改.
回去將.
他的那個最近.
飛黃騰達.
連升幾級.
那個興高采烈.
告訴他們.
但是.
唯有末底改.
不向我這樣來跪拜.
我心裡就不喜歡.
那一切的尊榮.
都和我無益.
好像總有一條刺在內心.
有刺就拔掉.
用最好的辦法.
最快的辦法.
就是拔刺.
置他於死地.
他只不過是一隻木塞.
有什麼所謂呢.
你現在已經是皇帝身邊.
最重要的臣子.
你說什麼.
皇帝一定會聽.
哈曼爾說是.
那我就可以更加.
可以來到去.
事無忌憚享受這些榮華富貴.
不再需要有人.
刺眼刺鼻.
也彰顯他真是有實權.
整個故事去到第五章高潮.
因為.
三日三夜的禱告.
末底改.
一路.

$^{681}$事而不捨的來告誡自己的女兒.
你不要以為.
你在皇宮就可以倖免於難.
這個時候.
如果你不伸手.
來救援.
猶大人我們.
都能夠從別處.
得到拯救.
不過.
既然你現在有皇后的位份.
你是直接.
那個當事人.
你是代表著猶大族.
代表著你自己.
你更加代表著你現在有皇后這個位份.
你就應該要來到去.
向皇帝.
陳情表達.
你的擔心.
和這個危機.
是由哈曼.
所誣衊和製造出來的.
這個民族.
我們這個民族從來沒有想過.
要推翻.
和製造.
混亂和紛爭.
所以.
那個位份很重要.
大家就要學習.
你在那個位裡.
在那個位份裡.
你就要講你在那個位份裡.
應該講的話.
所以這個涉及到承擔.
那個承擔是很重要的.
而且是那個時機.
是你必須要講.
你不講就不行.

$^{721}$當然.
你可以說我不祈禱了.
我沒有信心.
我現在辭退皇后的位份.
和我的同胞.
一同面對那個國難.
我和你一起.
做回平民百姓.
大家一同禁食祈禱.
大家一同有苦同當.
有難同當.
就算了.
都可以是一個位份.
不過.
似乎以斯帖不是要.
戀棧他的榮華富貴.
而是他現在知道.
這個時機.
是神為他預備.
這樣的人.
在那個時代.
能夠站在破口.
這些真是時代的英雄.
所以以斯帖.
他的主一學教育.
他的猶大人的身份.
的榮耀感.
他現在得到的是.
全個波斯帝國.
最重要的身份.
但是都不夠他.
成為猶大人的身份.
來得重要.
和尊榮.
他以自己是猶大人.
的身份感到榮耀.
現在受到威脅.
他決心和他的同胞.
共渡那個難關.
他就是那個唯一.

$^{761}$可以挽回.
扭轉局面的.
那個神.
擺下來.
當然他們當時不是這樣想.
後來事後應該是這樣想的.
因為他們都是.
信奉神的.
身份.
神擺在.
皇宮裡面.
一個這麼重要的人.
來解除.
整個猶大滅族的危機.
也令到.
黑錘老王.
有一個.
不要因為一個政治錯誤的研判.
導致一些不必要的.
聖靈塗炭.
這是一個很大的.
轉變.
所以.
普爾節.
就是一個轉危.
危機的一個.
重要的節日.
所以我們人生都應該有.
我們自己的普爾節.
雖然我們不是猶太人.
我們都應該有我們自己的.
普爾.
我們有我們自己的生命的改變.
這裡就是.
以自己猶大身份.
為榮.
今日香港.
以至於一直擴散這個訊息.
全世界.
現在在地球的基督徒.

$^{801}$有多少的基督徒.
以自己是神選民的身份.
為榮.
這就是浩大的基督徵兵.
這就是.
很大的力量.
因為我們有的是.
做物主給我們的.
那個應許.
這就是很重要.
我們去七除八扣.
我們就發現其實是極少.
極少.
今日我們的基督教.
說自己.
聲稱自己是基督徒.
是有分的時候.
有利益的時候.
希望我不會這樣研判錯.
希望就是.
不是的 曾牧師.
何時我要背負十字架.
何時我就願意走出來.
如果有這樣的弟兄姊妹.
基督教又怎會被人嘲笑.
基督教又怎會.
成為了可有可無.
基督教.
怎會被人越來越.
覺得是一個好像沒什麼實力.
的基督教.
如果我們能夠像以茲帖那樣.
我們會背的.
焉知你得了皇后的位份.
豈不是為現今的機會.
其實不是這個機會.
是這個危機.
是這個危機站出來.
所以我們平日的裝備.
令到我們.

$^{841}$自身是一個機會.
甚至乎一個危機.
我們要站出來.
我們能不能夠有這個勇氣.
我們能不能夠有經過.
三日三夜的祈禱.
而有那個智慧.
有那個內心的平靜.
這些是很重要.
我們是否已經.
自亂陣腳.
譬如大家都一定會知道.
這兩年的移民潮.
有多少人認認真真的禱告.
我清楚神的差遣.
帶領.
去到這個地方.
我要的就是將福音傳到更遠.
這個.
是否我們最高的.
最崇高的理想.
因為不是.
現在這個地方是沒有平安.
我一定要走.
你個個都走了.
走不到的那些怎樣.
關我什麼事.
這些不關我的事.
我只顧及我自己.
我的家人就算了.
那裡有沒有聚會.
我也不知道.
不是一個真正的聖經的教導.
的那個遺訓.
真正的遺訓就是.
在神要安排你的一個危機裡面.
你來到去挺身而出.
不是匹夫之勇.
而是.
過去的一切的裝備.

$^{881}$就是為著.
你在這個危機裡面.
成為那個中樓底柱.
以斯帖就扮演.
著這個角色.
第六章.
王睡不著覺.
哈曼都睡不著覺.
不過.
就是開心到一個地步.
他第二天早上就要.
致末弟國於死.
似乎他繼續會暢飲.
繼續的.
來到去載歌載舞.
但是.
這個時候很平靜.
神的工作.
默默的去做.
如果沒有王的.
這個睡不著覺.
就不會有往後第六第七章.
那個轉危.
危機的一個.
改變.
神怎樣令到皇帝.
睡不著覺.
剛好那些歷史官.
揭那些歷史.
又揭中這個事件.
記住.
以斯帖被從來都不是.
一份小說.
是一個作者.
道讚出來的一個小說.
一個懸疑.
偵探小說.
未去到結尾.
你都不知道結局.
譬如國內現在有一套電影.

$^{921}$很賣得.
但是那個.
口碑.
是兩極的.
有些人說很好看.
因為去到最後.
你才知道原來是這樣的結局.
你想都想不到.
但有些人就說不好.
是好收得.
我不說名字了,免得賣告白.
但是這裡告訴我們.
神令到皇.
睡不著覺.
有一個戲劇性的.
當夜的扭轉.
而且去到天亮.
哈曼剛起床.
進入皇宮.
就遇見皇帝的決定.
所以這裡開始說的.
那個時間點.
真的讓我們想像得到.
神掌管著這個民族.
其實他們祈禱了很久.
現在去到這個時間.
其實已經是三月.
到第七章我們就知道.
那夜皇睡不著覺.
就吩咐人取歷史來.
念給他聽.
正如見書上寫著說.
皇的太監中有兩個守門的.
闢探和提烈.
兩個太監的名.
想要下手害阿克徐老王.
沒抵改.
將這事告訴皇后.
因為他不可以直接進去見皇帝.
就告訴他的女兒.

$^{961}$女兒就轉告.
但是那個honor是要給.
沒抵改.
不是皇后自己要的.
所以那個歷史觀.
聽皇帝這樣說.
就要記住是給沒抵改.
但是皇通常都有.
說完就說完.
忘記了賞賜.
通常有時皇帝就因為日理萬機.
寫完就忘記了.
那件事原來.
是未完結的一個.
的.
論功行賞的.
這樣的事件.
就變成.
那個時間.
差不多臨近天亮.
皇帝就說.
這樣的事不賞賜怎麼行.
就當我好像對沒抵改.
好像不公平.
第三節.
王說沒抵改行了這事.
其他什麼尊榮.
著為沒有.
事後皇的臣僕回答說.
什麼?王說.
誰在院子裡?那時哈曼.
正進皇宮的外院.
要求王將沒抵改.
掛在他所.
預備的木架上.
當然是在想著跟皇帝說.
是他的心在想著.
或者是作者告訴讀者.
就不是皇帝知道.
因為他根本沒有機會再說這件事.

$^{1001}$臣僕就說.
哈曼站在院內.
王說叫他進來.
哈曼就進去.
王問他說.
王所喜悅尊榮的人.
當如何待他呢?.
哈曼心裡說.
他還沒有馬上回答.
因為他以為對號入座就是他自己.
因為現在皇帝.
一路興高采烈.
第二天都邀請他.
去赴皇后的.
私人宴會.
所以現在皇帝突然間說.
哪個人.
如果是我喜歡的.
是我所器重的.
要可以得到什麼的尊榮呢?.
哈曼當然有.
再高級的.
當然是去四季酒店.
當然是去坐羅斯萊斯.
遊整個香港島.
中環,金鐘.
開篷的諸如此類.
然後還要上遊艇.
諸如此類.
當然說是自己.
因為以為是他自己.
所以過分的對號入座.
哈曼心裡說.
王所喜悅尊榮的.
不是我還是誰呢?.
哈曼回答說.
王所喜悅尊榮的人.
當將皇上穿的朝服.
和戴冠的御馬.
都交給王極尊貴的一個大臣.

$^{1041}$名他.
不是那個被戴的人.
是再有尊貴的.
可能是宰相.
可能是其他的監國.
來幫忙去.
加官進爵那個人.
你可以想像得到.
哈曼盡他能夠想到的.
他.
或者可能他內心認為.
這樣的排場.
就是人人都見到.
我穿著王的朝服.
是王派給我的.
清朝的那些王馬褂.
穿著.
然後騎著御馬.
即是那些馬本身都有裝飾.
裝飾了.
這樣走出去.
大街大巷的示眾.
讓人看到.
這個就是王所尊榮的大臣.
果然是.
走遍城裡的街市.
在他面前宣告說.
王所喜悅尊榮的人.
就如此待他.
還要這樣到處去說.
即是打鑼打鼓.
好像那些考到狀元.
就要上京.
然後有些人.
一路吹著.
打鼓送他進京城.
這樣.
這些相信大家都會看到.
大家都會明白.
什麼是真正的尊榮.

$^{1081}$不同的時代一定有不同的方式.
重點是那個Honor.
現在哈曼.
認為他自己就應該得到.
之前.
好像.
不知道他在王的面前.
有什麼功勞.
而王這麼器重他.
而他得到一個.
高位.
而他又來到去.
見到末底改求王.
王又聽他講.
這個民族因為他強調.
上綱上線.
是會危害.
國家的安全的這個民族.
所以你給我.
我就奉獻給.
庫房多少多少的錢.
那王帝覺得.
這件事我這個.
這麼信任的大臣.
幫我發言.
因為之前都有人想殺王帝.
所以很合理地.
王帝當然擔心.
這個國家和自己的安全.
所以就准許了.
這個哈曼.
可以來到去.
有一個這樣的法例.
12月13日.
因為其實是制簽的.
制到這個亞達月.
12月13日.
就是吉日.
就可以滅盡.
國內127省.

$^{1121}$包括書山城的所有的猶大人.
你想想.
可以想像得到.
對王帝來說.
幫自己去掃平.
那些異見人士.
反對的.
攻擊國家的.
都是一個辦法.
也都借哈曼的口.
有可能引起當時.
波斯人的一個妒忌.
這些我們不是太知道.
那個細節.
但是我們知道那個行動.
就是這個哈曼.
所引發出來的.
所以現在哈曼以為.
更大的權.
更大的就是.
人都見得到的尊榮.
就來到去講給.
王帝知道.
就是要穿著朝服.
要騎著御馬.
王帝的馬.
而且還要有些人拖著這匹馬.
到處跟人說.
告訴人知道.
這個就是王所尊榮的人.
要這樣對待他.
第十節.
王對哈曼說.
你束束的將這衣服和馬.
照你所說的.
向坐在朝門的尤大人.
沒抵改去行.
凡你所說的.
一樣不可缺.
沒抵改.

$^{1161}$坐在朝門.
可能真的有他的官職.
如果不是怎麼會.
王帝都會知道呢.
那個官職可能都不是低的.
因為經過皇后.
居中的一個建議.
推薦等等.
所以沒抵改.
都是忠心辦事的人.
不是要製造紛爭.
不是要引起皇帝的注意.
希望巴結皇帝.
我們曾經講過.
聖經人物.
大家還有沒有印象.
大衛第三個兒子.
這個.
忘記了他的名字.
有沒有人可以提醒我.
大衛的兒子.
最終都是死於非命的.
稍微等一等.
大家完全忘記了.
要殺爸爸的.
大衛.
是亞沙隆.
未做皇帝之前.
早上.
神那麼早就站在十字路口.
等百姓上朝.
求問皇帝.
大衛.
你這件案.
有理.
你這個冤情.
我深深的感受到.
聖經的作者.
其實有很多暗示據.
他暗中.

$^{1201}$得到百姓的心.
亞沙隆.
人家.
做太子.
當然是睡到日上三竿.
不用怎麼做.
但他不是,很勤力.
很勤力.
天光就已經在.
十字路口等那些百姓上朝.
去求問大衛皇帝.
所以.
沒抵改.
似乎不是這樣的人.
他不是一個真正有心計的人.
所以.
有心計也是求神要憐憫.
聰明的人.
才會有心計.
我不是要論斷聰明.
聰明是正面的.
但人心.
再加上聰明.
一擺錯位.
就變成心計.
要小心.
心歪了.
看事物也會歪了.
急功近利.
又會變成.
趨炎附勢.
又會變成.
有時那種明察秋毫.
加上了.
自己的眼光.
就變成.
明察秋毫.
帶著一個.
背後的手段.
和利益.

$^{1241}$所以聰明的弟兄姊妹.
真的要求主憐憫.
也不是祈禱要求神.
讓你蠢一點.
聰明是好的美德.
是正面的.
不過聖經說的真正的聰明.
是認識和敬畏耶和華.
這是真的很重要.
哈曼.
看到的沒抵改.
仍然是這樣的時候.
但現在皇帝有一個大的反轉.
就是說.
祈禱開始生效.
神終於出手.
當然.
如果你不相信耶穌.
那麼巧合.
這些只是小說家的橋段.
不是怎樣製造.
出人意表的懸疑.
令票房可以收得.
如果這麼平鋪直敘.
哪有人看電影.
沒錯.
小說一定要這樣.
加上各種角色.
或者是加一些竄改的劇情.
但這裡是原汁原味.
作者是將歷史.
說出來.
再敘述.
猶大人.
為什麼要守普爾節.
為什麼以自己的民族為榮.
這是整本書.
最重要的中心思想.
在這個時候.
沒抵改仍然扮演著.

$^{1281}$一個亡國.
和忠心的.
在波斯帝國.
來當官.
和忠於自己的民族.
和信仰.
猶太人的一神教.
對耶和華的敬拜.
對十誡.
對自己的安息日.
遵守.
但招致的.
就是無妄之災.
就是哈曼.
作者已經交代了.
哈曼是亞馬力族.
所以是與神為仇的民族.
現在哈曼就面對.
突然間的逆轉.
昨晚興高采烈.
當夜就做了一個.
五丈高的木架.
是怎樣的形具呢?.
如果我們那時候有手提電話就好了.
讓我們看到是怎樣的形具.
怎樣殘忍.
怎樣要將沒抵改.
剝皮.
拆骨.
我們如果見到就好了.
但可惜沒有了.
那個時代都不會有.
也沒有人畫下它.
變成我們都不能夠在聖經裡面.
找到任何的形具.
變成我們都不能夠在聖經裡面.
找到任何的蛛絲馬跡.
但無論如何.
沒抵改就面對.
本來他想置於死地的猶大人.

$^{1321}$沒抵改的對頭哈曼.
現在還要在他面前.
他就是那個扮演.
這個人就是得到.
王的尊榮的人.
就要這樣對待他.
是王吩咐他去說.
所以第十節就是李卓宗要去.
第十一節.
於是哈曼將朝服給沒抵改穿上.
史達奇上馬.
走遍城裡的街市.
在他面前宣告說.
王所喜悅尊榮的人.
就如此待他.
沒抵改迎回到朝門.
沒將這些擺上心.
因為整件事都未完結.
現在.
王的尊榮的人.
就如此對待他.
未完結.
現在突然間得來的尊榮.
他都不擺在心裡.
都不會.
據為是自己一個.
一個人可以免於難.
仍然回來站在朝門.
回到朝門.
即是繼續在朝門裡面.
去等待.
整件事.
等待皇后以司帖的回覆.
事態的發展.
所以作者.
用一些很小的說話.
就已經讓我們看到猶大人.
不會貪生怕死.
他們.
真是很優質的民族.

$^{1361}$真是以自己.
雖然現在是面對神的.
刑罰.
的一個處境.
但是仍然不會放棄.
自己神選民的身份.
如果真是要滅族.
以司帖也盡了力.
那就滅.
是神要我們毀滅.
是神懲罰我們.
但如果不是仍然.
未去到最後一刻.
他們都不放棄.
這亦都是在智女行間.
隱隱透視到.
就是我們不要肅明論.
我們真是不知道.
明天如何.
但是不知道明天如何.
其實知道明天會爆風球.
很厲害.
特首都要宣佈叫人出來.
第一次穿上.
一樣的衣服.
外套.
通常台灣的領導人才這樣.
香港很少這樣.
個個穿自己的西裝.
很嚴肅.
個個穿著不知在哪裡買的套.
個個一樣的.
在裂隊.
天文台台長都要走出來.
明天可以預計到很嚴重的風.
但是猶大人.
雖然他們.
不知道神.
如何解救他們.
但他們未去到最後一秒.

$^{1401}$仍然認定.
神可以在別處.
解救他們.
這個信念是很重要.
基督徒要有.
主耶穌已經在.
十誡上得勝的.
信心.
誰能使我們與基督的愛.
隔絕.
難道是患難魔.
是逼迫魔.
是飢餓魔.
是赤身怒體魔.
是危險魔.
是刀劍魔.
諸如此類.
保羅一直這樣.
他能夠數的.
其實是他自己經歷過.
經歷過這些痛苦.
這一切都不能夠使我們與神的愛.
隔絕.
因為這一份的愛.
是在我們的主.
宣告.
雖然他還未見到主耶穌.
他仍然面對地上.
很多猶太人對他的.
白眼,侮辱.
嘲笑.
和外邦人.
羅馬政權對他的迫害.
諸如此類.
但他的內心是很堅強.
他內心的信念.
信仰是很.
實淨.
可以感染他身邊的年輕的同工.
所以.

$^{1441}$97年的時候.
我剛剛信主.
還未信主.
不太清楚.
不是,是87年,88年.
97是.
我不是見很多.
因為那時候信了不久.
剛出來傳道.
很多教目.
因為.
中英宣言.
諸如此類.
基本法.
很多大的教目.
低調地離開.
離開香港.
不告訴別人.
引發香港的教會.
突然間清黃不接.
很多信徒很失望.
教目低調地.
自己離開.
(問:信仰的目標是怎樣?).
(還是要相信神?).
(還是要學習?).
(還是要信仰?).
這些都是.
信心的問題.
如是者.
過了這麼多年.
每個時代都有危機.
每個時代.
神都會透過不同的事情.
考驗我們.
信仰.
忠貞.
讀以斯帖記.
不可能不面對.
要面對的問題.

$^{1481}$就是我們那個忠貞.
是否以神國度的選民為榮.
以斯帖.
或者是活體改.
他們是第一個知道的.
可以馬上遠走高飛.
走上耶路撒冷.
都可以逃避這個問題.
或者.
低一低頭.
跪拜.
解決了問題.
但是似乎.
告訴我們神有他的時間.
人可以做的.
就是.
不是要去探知.
那個隱蔽的.
神的事情.
一些奧秘.
或者用什麼方法.
來尋求.
那些.
可能是很靈恩.
不是.
應該做的就做.
應該做的就是站穩自己的身份.
應該做的就是禁食.
祈禱.
應該做的就是.
誰有那個位份的就要來.
代表去說.
如果有.
咨詢的我們要表達.
什麼是基督教.
什麼是我們的信仰.
但是最終不行.
不容許有這樣的宗教.
不要緊.
因為我的神不會因為你不容許.

$^{1521}$就不存在.
但是怎樣能夠令到我的神.
讓人見到他的存在.
就是我們甘心樂意.
放下我們的權利.
接受那個.
制裁甚至乎刑罰.
初期教會.
任何一個使徒都是這樣.
面對的.
要坐牢不要緊.
我是為基督耶穌被囚的保羅.
我是為基督耶穌來受苦的.
的獨人.
或者見證.
這些是充斥著初期教會.
他們的信念.
感恩的就是他們去到很後期.
他們已經很寶貴.
自己這個.
神獨人的身份.
很寶貴.
基督教已經不再是一個.
很弱雞的宗教.
已經是可以.
抗衡任何的.
勢力的一個信仰.
因為他們認識.
神的力量.
所以這裡.
就讓我們見到.
沒底改他們的祈禱.
現在真的有那個答應.
接著又有一個.
很戲劇性的一個.
和剛才昨晚的哈曼.
回家見到.
西利斯太太.
又見到她的.
謀士再說.

$^{1561}$相同的話.
就說她是受容.
第十二節的中間.
哈曼卻悠悠悶悶的蒙著頭.
急忙回家去了.
即是說.
做了一天的車夫.
說了周圍去.
告訴人知道.
沒底改就是王所尊貴.
所尊榮的人.
大家要這樣對他.
很醜.
非常之.
應該是悠悶悶.
急忙回家.
將所遇的一切事.
詳細說給他的妻西利斯.
和他的眾朋友聽.
這班賊子.
他的智慧人和他的妻子.
西利斯對他說.
你在沒底改面前.
此而敗落.
他如果是猶大人.
你必不能勝他.
終必在他面前敗落.
之前又不說.
現在這個主子.
落難了.
說出真話.
之前意氣風發的時候.
又不告訴他.
警告哈曼.
不要和猶大人.
爭一日的長短.
因為他們敬拜的神.
可能是真神.
一直都不說.
現在落難了就說這些風涼話.

$^{1601}$令哈曼.
心更加痛.
當場心臟病.
發死.
他們還與哈曼說話的時候.
王的太監就來.
催哈曼快去赴.
以師帖所預備的筵席.
對 晚上.
做完車夫馬夫之後.
就要回家.
因為沒有地方傾訴.
他內心的.
完全沒有面.
被落面的那種尷尬.
他就要回家.
被落面的那種尷尬.
那種侮辱.
你想想一個驕傲的人.
如果當眾被人侮辱.
其實那種痛.
很難想像.
一個人很謙卑的.
從來都不介意別人怎麼說他.
因為他不是這樣的人.
完全不會介意.
其實內心的平安.
是沒有人可以奪去.
但一個驕傲的人.
不斷聽的.
不斷的敬畏.
恭維的說話.
突然之間見到真相.
是自己好像完全沒有穿衣服的真相.
就有很大的羞恥.
所以他一定要.
立即衝回家.
告訴他們.
有沒有什麼出謀獻策.
有什麼可以挽回.

$^{1641}$沒有想過自己要面對生死.
那些人就說.
你在尤大人面前.
這樣恥而敗落.
終必.
你一定會.
面對.
在他面前.
敗落.
如果你不能勝過他的時候.
第七章.
我沒有帶水.
不要緊.
王帶著哈曼來赴.
皇后以斯帖的筵席.
因為哈曼要回去.
面聖.
因為催他.
皇后的時間到了.
皇的太監要來催他.
和皇帝一起.
再赴席.
所以你可以想像到.
一晚的扭轉.
是皇帝睡不著覺.
就是成為.
神聽祈禱的一個.
這麼重要的開始.
一方面.
是作者.
給我們有這個懸念.
這個事情的發展.
但更加看到的就是.
造物主.
掌管著萬有.
掌管著皇的心.
我們以為皇一定是.
喜怒無常.
是的.
多數都有權.

$^{1681}$是的.
但我們又想到神可以管理住一個皇.
突然之間這麼關心.
朝政的事.
突然間聽到.
末底改的未賞賜.
而立即.
來做.
剛好哈曼準備.
來求皇.
將末底改.
給他.
因為上次都可以了.
這次為什麼不行呢.
給末底改他.
真正幫皇帝去行刑.
即是說他的死期到了.
誰知當晚第二晚.
就是他自己哈曼的死期到.
第七章的第二節.
在第二次在.
走席言前.
皇又問以斯帖.
第三次了.
皇后以斯帖說.
你要什麼我必賜給你.
你求什麼.
就是國的一半也必.
為你成就.
皇后以斯帖回答說.
我若在皇眼前蒙恩.
皇若以為美.
我所願的是.
願皇將我的性命賜給我.
我所求的是求皇將我的本族賜給我.
因我和我的本族被賣了.
要剪除殺戮滅絕我們.
我們若被賣為奴為婢.
我也閉口不言.
但皇的損失.

$^{1721}$敵人萬不能補足.
現在開始將自己置於猶大人這個民族裡面.
讓皇突然發現.
最近要滅絕的民族.
豈不是猶大.
可能皇帝未知這麼詳細也不出奇.
他只知道要准許哈曼.
來制簽.
制倒12月13日.
和要滅絕.
那個與皇.
與波斯為敵的民族.
可能沒有詳細說是甚麼民族.
但無論如何.
現在以斯帖將自己擺在自己的民族裡面.
這是很強大的.
對信仰的認信.
才會這樣做.
因為一般人都會怕死.
一般人都不希望有任何自己已經得到的利益.
一下子就沒有了.
但信仰可以超越了皇后的身份.
她的民族.
不只是一種普通的民族性.
而是她真是以自己神選民的身份.
與她的同胞.
來到一起的共度.
好像摩西.
醒覺自己是以色列人.
她的媽媽.
她的姐姐米尼暗.
她的媽媽叫約基別.
她意識到自己.
人就生在皇宮.
但心靈是以色列.
即是說小時候聽過很多聖經故事.
羅亞方舟.
阿當夏娃.
這些口傳的歷史.
以色列人一定是聽過.

$^{1761}$所以摩西長大之後.
不想再做埃及人.
希望與她的同胞一起面對無理的欺壓.
如果她再熟悉一點歷史.
以前不是這樣的.
埃及和以色列是很和好的.
我們的老祖宗約瑟.
是帶來埃及人對我們很優待的.
給了我們最好的地方.
歌山地去牧羊.
但好景不常.
有不認識約瑟的新王.
可能是外族入侵.
但又有類似血統.
所以又稱自己是埃及的法老.
名號一樣但可能民族不同.
這些不足為奇.
蒙古也有很多不同的支派.
摩西的民族性和信仰.
令她要將自己放在民族安危.
這些民族性和信仰.
若然如此強烈.
從小就在她的血液流通.
成為日後大用的人.
也不是無跡可尋.
從小就有穩固的信仰.
所以奉勸.
現在在座的可能是公公婆婆.
爺爺媽媽.
未必是父母.
對自己的下一代,下兩代.
基督教的教育.
深化在他們的血液裡.
不要忘記信仰是最大的遺產.
你給她信仰.
將神給她.
她自己去經歷神.
她自己一生去倚靠神.
勝過千萬家財.
不是不好的一定是好.

$^{1801}$但那個來得更加永恆穩固.
因為神是永遠不變.
財物是可以變.
會貶值.
但神在她的生命裡.
如果能夠越來越去寶貴自己的信仰.
一路去裝備去追求.
日後被神大用的通常都是這些人.
這個真是很重要.
譬如現在的白玫君牧師.
就是白會員.
播道神學院的院長的兒子.
非常敬虔的父母.
報院長的太太.
報她.
我忘記了她的名字.
是幾代都是基督徒世家.
很敬虔的報師母.
我聽過她的講道.
分享就是很敬虔的一位父女.
她的兒子就是白玫君.
現在是全球知名的華人的牧者.
你可以想像到.
小時候他們的DNA已經滲入他們的血液.
他們的信仰.
父母的敬虔.
對他們的影響.
偶然間都聽到他們的講道.
就這樣講自己的父母的事情.
這些都是真的很重要的.
所以我們看到哈曼就面對她自己所製造出來的惡果.
所以行惡永遠都不會有好的下場.
聖經又告訴我們.
一個這麼極惡極驕傲的人.
他的下場就來得更加快.
更加慘痛.
哈曼的猶大.
一直都是偷錢的.
偷錢黨.
一直都是混入在12個的使徒當中.

$^{1841}$以為耶穌是一個這麼仁慈的老師.
祂不會察覺.
其實耶穌就明察秋毫.
但是就沒有將祂供出來.
相信有很多不同的暗示等等.
很多教訓裡面都有講.
就不是縱容祂.
買祂怕.
而是希望祂能夠有一個悔改的機會.
最後的和祂親聚.
都是一個暗號.
祂都知道的.
不過把心一橫.
因為老師再沒有利用價值.
不如就決心背棄他.
就是黏在法利賽人和祭司獎的行列.
希望不單自己不受牽連.
更加可以出賣.
都帶來一個可以換取到一些少量的利益.
這些就是內心一直都很驕傲.
而沒有改變自己內心的罪惡.
導致的下場就是猶大.
他會自己上吊去吊死自己.
後來小譚村就說那條繩斷了.
他就跌在石頭上.
腹部都穿了.
長穿肚爛就是這個意思.
就是猶大的下場.
還沒有計算.
聖經對他的描述就是.
這個人往陰間痛苦的地方.
往他自己的地方去.
就是說一定不是上天堂.
將來在天堂是見不到猶大的.
這就是哈曼他的下場.
現在以斯帖在這裡就說.
第五節.
「黑錘老王問皇后以斯帖說.
善感喜意如此行的是誰?.
這人在哪裡?」.

$^{1881}$就是說暗示他真的不知道.
實際具體的那個操作.
哪個民族這樣來想與波斯為敵.
這個是哈曼的無中生有.
製造那個敵對.
引起王的擔心.
而聽他講立了那個命令.
以斯帖就說.
仇人敵人就是這個惡人哈曼.
所以第二晚.
整個氣氛不同了.
當然沒有飲宴.
因為皇帝飲宴之前.
再一次問.
第三次問以斯帖.
你到底要求什麼呢?.
為什麼你要設兩晚的筵席呢?.
一定是你內心有心事.
你是想要那顆綠寶石.
哪個地方呢?.
要建造一個行宮還是怎樣呢?.
要買一艘遊艇還是怎樣呢?.
不是的.
現在以斯帖.
真的在找第二晚的機會.
又是只有哈曼和皇帝.
但是哈曼因為經過了早上的事件.
其實他面帶了愁容.
沒有昨晚那種興高采烈.
意氣風發.
鬆毛鬆翼.
但是在這個時候.
就是殺生之禍的開始.
哈曼在王和皇后面前.
就沉驚惶了.
王便大怒.
起來離開酒席.
往御苑去了.
哈曼見王定意要加罪於他.
就起來求王后以斯帖救命.

$^{1921}$王從御苑回到酒席之處.
見哈曼伏在以斯帖所靠的塔上.
王說他竟敢在宮內.
在我面前凌辱王后嗎?.
這話一出王口.
人就蒙了哈曼的臉.
事後王的一個太監.
名叫哈波那.
之前出現過的.
哈曼為那救王.
由宮的沒底改.
作了一個五丈高的木架.
現今立在哈曼的家裡.
誰打小報告的?.
好像那些講者一選完之後.
就有什麼消息.
就有什麼的好外.
那些消息來源從哪裡來的?.
選的時候又不說.
選到才說.
就是說.
他如果拿不到冠亞季.
講完也沒意思了.
都淘汰了.
就是他拿到才值錢.
那些消息.
那誰.
昨晚又不見他講的.
誰來的消息.
為什麼會洩露了呢?.
因為是怎麼樣呢?.
他興高采烈一陣子.
到處跟人說.
所以哈波那.
都知道有這個木架.
這次就不是推舉他.
是再加多一腳.
踩他進去那個深淵裡面.
已經黃大怒了.
還要說.

$^{1961}$這個哈曼為了.
黃你所賞賜的那個中神.
墓底改.
預備了一個行刑的木架.
今晚就要置他於死地了.
你不說.
沒人說你啞的.
你這樣對哈曼.
這樣來落井下石.
一來當然是對黃的忠心.
但是好像屋漏兼逢年夜雨.
哈曼就面對接連的打擊.
相信現在.
當時馬上要面對的就是那個行刑.
第十節.
對不起.
上面先.
黃說.
把哈曼掛在旗上.
於是人將哈曼掛在他.
為墓底改所預備的木架上.
黃的憤怒.
這才止息.
這個就是來到第七章.
整件事的扭轉.
是因為黃帝睡不著覺.
黃帝要聽歷史.
黃帝聽到墓底改的功德.
未賞賜.
時機.
也是另外一邊的鏡頭.
以撕帖.
決心以自己猶大民族的身份.
來去進見黃帝.
因為他有這個位份.
就是皇后的位份.
就是說捨他其誰.
沒有第二個人了.
所以如果我們能夠找到.
十個一百個.

$^{2001}$知道自己的位份.
在適當的時候.
站在神的真理.
為福音為神來作見證.
那個力量就非常之大了.
很大的感染力.
和那個震撼力.
連死都不懼怕的人.
就是殉道士為神的福音.
所以不是亂做的.
是裝備好.
是很成熟.
是很知道真理.
很知道現在這一刻不說就不行的.
這些人就是一路是神所預備.
好像以利亞的時代.
以利亞就說我為神大發熱心.
第二天就被阿爺洗別追殺.
然後就走了去羅騰樹下.
神就恢復了他的心靈.
恢復了他的內心鬥志.
也仍然用他.
叫他回去.
要告立三個人.
告立以利沙.
告立一個靈士的兒子.
這個我不知道是否叫做彼哈達.
和耶護那個君王.
神就跟以利亞說.
但是我在以色列留下了七千人.
未曾向巴力屈膝.
未曾向巴力親嘴.
七千個是無名小卒.
散居在以色列都不知道是誰.
但是他們有一個很大的決心.
現在我們的民族.
我們的宗教面對外族.
耶洗別是西頓.
西頓就是這個.
突然間忘記了名字.

$^{2041}$阿斯塔祿巴力的這些假神.
但是不向巴力親嘴的.
七千個尼希米.
尼希米,對不起.
以利亞,七千個.
不只是你一個,是有七千個.
那是從哪裡來的呢?.
以色列很小的,我為何不知道.
有這麼多的,好像間諜一樣.
你不要管,不關你的事.
是我供養他們的.
七千個散居在以色列.
以利亞其實很醜.
只有他一個大發熱心.
所以如果我們有這些眼光.
看看不同的宗派.
浸信會,播道會.
不同的宗派.
有些人這樣為神.
這樣來持守真道.
這樣來為主擺上.
其實很多美麗的故事.
所以有時很多不是很喜眼的弟兄姊妹.
他們很敬虔.
有時那些婆婆不是很識字.
但他們很單純的敬虔.
真的好像閃耀的七千人.
所以我們要追求做那個.
沒有人知道的七千人.
當然神要興起以利亞是真的.
但其實我們要追求做.
是人都做到的.
只要你內心對神有堅定的信念.
信仰就行了.
不需要甚麼資格.
是人都可以做到的那七千人.
就是給神嚴重你的內心.
所以這裡就讓我們看到.
在這七章裡面.
就是神怎樣來到去.

$^{2081}$透過沒底改.
透過以斯帖.
透過那個民族.
透過他們的禱告.
即是意味著.
原來他們的信仰是非常真實的.
他們的信仰是經得起考驗的.
他們的祈禱真是直上天庭.
是神聽的.
是神會解救的.
是神會在適當的時候.
所以是適當的時候.
神的時候.
不是我們的時候.
不是我們的方便.
今天的媒體世界.
甚麼都是要方便.
要立即按個鍵就要有東西看.
三秒都開不到網頁.
你就走了.
我們已經發展到.
訓練到我們是沒有耐性.
這個是祈禱等候神的大忌.
基督教甚至乎都不用祈禱.
我們想完覺得OK就行了.
沒有了一種等候神.
所以有時神都要透過一些環境.
要我們一起去等候.
去禱告.
有些事情大到一個地步.
都不是你的普通常識就可以過到的.
真的需要整個民族這樣祈禱.
我一直的心願就是.
真的看到祈禱寶貴.
然後願意大家一齊回來.
在公土會裡面去禱告.
不是要製造一些只是氣氛.
而是真的大家看到祈禱的寶貴重要.
現在都有網絡直播.
但情願按著TVB看無聊電視.

$^{2121}$都不情願去祈禱會.
這個是現實.
是社會告訴我們的現實.
不要怪那些不信的人.
看不起我們基督教.
你們這班人說到祈禱這麼重要.
但你們都不祈禱.
我們也不需要像打書仗那樣.
而是我們要不斷去鼓勵.
無論新一代或下一代.
都讓他們看到真的祈禱的重要.
不是一個例行公事.
不是掛在口唇邊的一種咒語.
是真的大家放下自己去祈禱.
今晚就說到這裡.
我們下一次就由第八章.
說到第十章完結.
之後我們整本書做一個簡單的總結.
和那個時代.
我們又如何在以司帖的時代.
如何做一個研判等等.
我們一起禱告.
再一次多謝主.
我們能夠逐節去看.
看到他們的張力.
看到這個民族.
他們有敬虔的沒底改和以司帖.
當然是以他們兩個做一個代表.
有更多沒有提名的.
我們也是這樣在禱告裡面去祈求.
去化解這一次的災難.
放下在聖經裡面.
不是要我們去追求一些超自然的神蹟.
或者好像很奇怪的靈因.
而是要我們建立.
我們在我們的時代.
基督教的榮耀感.
對神的榮耀感.
對信耶穌.
對成為神兒女.

$^{2161}$這個身份的寶貴和尊榮.
幫助我們在越是不適合講福音的環境.
讓我們更加以福音為榮.
而不是以福音為恥.
願神幫助我們.
也讓神的話語.
在很多人的心靈裡面.
得到那個建立幫助.
再一次我們為到明天後天.
嚴重的颱風.
香港我們其實也挺有信心.
不會有很大的災難.
好像溫代的年代那樣.
因為很多的樓宇結構.
都有非常高的防風的能力.
但是我們擔心的就是.
如果轉入內地的農村.
帶來的那個破壞.
可以是非常嚴重的人命傷亡.
所以我們也盼望求主憐憫.
很多人的心靈.
身心靈的脆弱.
盼望這個颱風.
雖然它是沒有衛甲.
只不過是自然界的人.
釀成的風暴.
但是我們盼望神都管理.
真的叫這一切都停止.
願主來到去施恩.
我們這樣禱告奉耶穌基督的名.
阿們.
\newpage



\section{提摩太後書 4:1-8}
\label{sec:ay_HsjJTAHo}
\textbf{2023-09-03 加恩樓《保羅最後的勸勉》提摩太後書 4\_1-8 講員:鄧泰康傳道}
\newline
\newline
連結: \href{https://youtube.com/watch?v=ay-HsjJTAHo}{\texttt{https://youtube.com/watch?v=ay-HsjJTAHo}} ~~~~ 語音日期: 2023-09-05
\newline
\newline
\hyperref[sec:9XvZPOyWZdE]{\small{< < < PREV SERMON < < <}}
~
\hyperref[sec:index_chronic]{\small{[返順時目]}}
~
\hyperref[sec:index_scriptual]{\small{[返順卷目]}}
~
\hyperref[sec:Bres3qEZ6Yk]{\small{> > > NEXT SERMON > > >}}
\newline
\newline
提摩太後書 4:1-8
\newline
\begin{longtable}{cl}
\hline
\hline
章節 & 經文 (和合本修訂版)\\
\hline
4:1 & \begin{tabularx}{0.7\textwidth}{X} 我在神面前，並在將來審判活人死人的基督耶穌面前，憑著他的顯現和他的國度鄭重地勸戒你： \end{tabularx} \\ \\ \relax
4:2 & \begin{tabularx}{0.7\textwidth}{X} 務要傳道；無論得時不得時總要專心，並以百般的忍耐和各樣的教導責備人，警戒人，勸勉人。 \end{tabularx} \\ \\ \relax
4:3 & \begin{tabularx}{0.7\textwidth}{X} 因為時候將到，那時人會厭煩健全的教導，耳朵發癢，就隨心所欲地增添好些教師， \end{tabularx} \\ \\ \relax
4:4 & \begin{tabularx}{0.7\textwidth}{X} 並且掩耳不聽真理，偏向無稽的傳說。 \end{tabularx} \\ \\ \relax
4:5 & \begin{tabularx}{0.7\textwidth}{X} 至於你，凡事要謹慎，忍受苦難，做傳福音的工作，盡你的職分。 \end{tabularx} \\ \\ \relax
4:6 & \begin{tabularx}{0.7\textwidth}{X} 至於我，我已經被澆獻，離世的時候到了。 \end{tabularx} \\ \\ \relax
4:7 & \begin{tabularx}{0.7\textwidth}{X} 那美好的仗我已經打過了，當跑的路我已經跑盡了，該信的道我已經守住了。 \end{tabularx} \\ \\ \relax
4:8 & \begin{tabularx}{0.7\textwidth}{X} 從此以後，有公義的冠冕為我存留，就是按著公義審判的主到了那日要賜給我的；不但賜給我，也賜給凡愛慕他顯現的人。 \end{tabularx} \\ \\
[1ex]
\hline
\hline
\end{longtable}
$^{1}$頂尖的早安.
今天就預備這篇導.
稍後我講導我會短一點.
不然時間太長.
我今天預備這篇導的時候.
我相信大家看這篇導的時候會覺得.
咦 今天這麼特別.
原因是我想大家都知道.
接下來我要做一個手術.
其實這個手術裡面.
難又不是難.
你說有沒有危險.
又有危險.
所以在這種情況下.
當我思考的時候.
這篇導是從心裡浮出來的.
但當我講這篇導的時候.
我想強調.
當我看這篇導的時候.
坦白說我連保羅的十分之一都沒有.
他的勤奮.
他的盡力.
他的毀身.
坦白說我十分之一都沒有.
裡面他講的很多.
譬如他講他自己.
他說他打過了當拍的路.
我跑進了.
坦白說我不敢去這樣去講.
反而我裡面加了很多問號.
美好的象.
我打過了.
當跑的路我跑進了.
這個道我守住了.
但是.
我沒有保羅這一種.
這麼.
我形容是一種.
很強的素人生命.
我覺得我不是.

$^{41}$但我有和他同一個心情.
我都會去想.
坦白說如果問我面對手術的時候.
我不懼怕.
我經常和電影姐妹說.
坦白說在裡面如果上帝.
這麼恩待手術.
如果不成功.
他就帶我回天教.
對於我來說我一點都不懼怕.
但反過來其實我問了一個問題.
我作為一個傳道人.
我可以留下什麼.
我會重新再思考.
一個傳道人.
一個傳道人原來可以留下的.
是一個生命的見證.
和保羅在這裡一樣.
對電影姐妹的祝福.
保羅就算你多強.
但到有一天.
他也說他的生命到一個時間.
他就要完結.
他要離開.
但留下來的是什麼呢.
那我也會重新思考.
在這段日子裡面.
他在聊.
究竟這個地方是怎樣.
我們一班的回眾.
愛慕養的羊仔.
你們是怎樣.
我心裡面.
我應該留下什麼給大家.
這是我在片導裡面不斷思考.
所以我會想.
其實我在想.
我有沒有這麼恩典可以安排到.
我在做手術之前的最後片導.
但我發覺我不是這麼如意.

$^{81}$為什麼呢 我明天才見醫生.
我明天才知道究竟我是什麼時候做.
我不知道是九月做還是十月做.
我還不清楚.
我下次再站在這裡.
跟大家說是十月一號.
如果主約許可.
十月一號再跟大家繼續說.
到時我會說另一篇.
但如果九月的時候我做完那篇導.
究竟之後是怎樣 我不知道.
雅各書曾經有一段經文.
其實明天如何我們還不知道.
你們的生命是什麼呢.
你們原是一片雲霧.
出現笑事就不見了.
你們只當說.
主約願意.
我就可以活著.
也可以做這事.
做那事 弟兄姊妹.
所以我們每一個人.
去求的都是主約.
如果你願意.
所以為什麼我作為傳道人的時候.
我一直都很強調.
為什麼我經常都鼓勵弟兄姊妹.
在地上你今天所追求的東西.
是不是真的有永恆價值.
說到嘴爛.
我都是這樣說.
我昨晚在團結裡面才說.
上帝賜給你的東西.
你儘管去享受.
我是很喜歡傳道書.
傳道書裡面.
講一件事 這個世界就是這樣.
所以你當.
趁著還有機會的時候.
享受上帝所賜給你的份.

$^{121}$甚至乎再嚴格來說.
你不去享受上帝所給你的份的話.
是一種.
上帝對你的責備.
你要去享.
你要趁著還有時間.
你和你所愛的人一起去享樂.
是不是.
傳道書教的.
但同樣的我們是不是給這些享樂.
給這些東西我們去佔據了我們的生命呢.
所以今天我講的.
這篇道的心情是什麼呢.
之前有一套電影.
其實我是很想.
我和我一個最愛的人.
我曾經和他說過.
其實很多時候我自己作為傳道人.
我都做很多的安息禮拜.
我舉個例子.
我爸爸離開的時候.
我都很感恩很多大英姐妹.
來送別安慰我們.
是很好的.
但這套電影裡我覺得是更好的.
他就是在這個主角.
他要知道自己癌症.
他要走的時候.
他提前開了一個.
告別會.
他在裡面.
他可以親耳的聽到.
他的朋友和他說的話.
他可以親身和他的朋友.
去說一番說話.
我和我一個.
我所愛的人.
我都很盼望將來有一天.
我的離開我都希望是這樣.
所以我今天帶著這種心情.

$^{161}$和大家講這篇道.
你會發覺這篇道的裡面.
當阿保羅.
他要將他這個.
的事情傳承下去.
給提摩太太他所非常愛的.
一位同工的時候.
他知道他有很多軟弱.
他知道他裡面會有很多不容易.
但是在裡面的時候我想大家留意.
第一句他已經說.
我在神的面前並在將來審判活人死人的.
基督面前憑著.
他的顯現他的角度祝福你.
什麼意思呢.
其實他在講的就是.
其實上帝是一直都得勝.
我想大家留意.
我們今天在地上為上帝打仗.
我們去面對很多不同的挑戰.
但你記住.
保羅之所以.
能夠一直挑戰到底.
直到他被囂電的時候.
是因為他一直帶著一種.
就是上帝已經得勝.
弟兄姊妹.
今天其實當你說.
你要和上帝同工的時候.
請你要記得上帝已經得勝.
你不是在打一場.
不知道贏還是輸的仗.
我想大家留意.
保羅是帶著一份得勝的信心.
和盼望去打這場仗.
你會發覺.
其實整條線.
是按照著上帝的心意去進行的.
大家有沒有留意.
其實在聖經裡面.

$^{201}$一早已經和你說了.
在末後的時候.
啟示也說了的結局.
已經說完了.
就是說.
我們今天在神面前的時候.
我們去繼續為上帝去工作.
我很想告訴弟兄姊妹.
你帶著一種已經贏的心態.
去繼續你的工作.
什麼叫已經贏的心態.
接下來的監牢不知道怎麼辦.
其實我們不會知道.
這個地方究竟是不是繼續的.
但是問題就是.
我都問了.
是不是沒有了監牢這個地方.
就代表會輸.
真正的輸不是沒有了監牢.
真正的輸.
不是沒有了監牢.
真正的輸.
是大家的心裡.
沒有了上帝.
我們這個群體裡.
沒有了上帝.
才是真正的輸.
但反過來.
監牢這個地方重不重要.
很重要.
因為我們是從這裡開始.
我們是在這裡.
我們經歷了很多很多.
我們的回憶.
這裡有很多我們一起的努力.
我們大家的關係.
是在這裡形成的.
所以我想強調.
為什麼我一直強調.
我們要求問上帝的心意.

$^{241}$因為唯有在神的心意當中.
你才會贏.
所以到這裡的時候.
當我們一直向前走的時候.
不是可有可無.
而是如果上帝.
很清楚地叫我們離開.
我們甘心離開.
但如果.
我們還不是很清楚.
我們不是很確定.
上帝要我們離開.
我們應該要去到最後一刻.
我去面對這個手術.
我坦白說我不知道上帝的心意.
到現在我都不知道.
坦白說.
可能是到我要做完這個手術.
我才知道上帝的心意.
如果上帝要我安然離開.
我安然離開.
但如果上帝主若許可.
祂要你繼續.
你就不是再好像.
還有很多的時間.
我的心態一樣.
如果上帝今次經歷了這個手術.
祂讓我可以繼續去事奉祂.
繼續和大家一起去努力.
我可以告訴大家.
我已經很開心.
大家一起去努力.
我可以告訴大家不要再浪費時間.
這個手術告訴我.
原來人的時間.
和機會是有限的.
今天我們欠上帝的債.
我們有沒有繼續努力地還給上帝.
我們今天所做的每一件事.
我們有沒有繼續回饋上帝.

$^{281}$所給我們的生命.
我們有沒有浪費我們的生命.
我們有沒有為了一些不需要糾纏的事.
我們在糾纏.
我們有沒有全心全意.
將自己去獻上給上帝.
所以不是上帝有沒有能力.
而是人是不是跟得上祂的步伐.
大兄姐妹.
當你跟上帝.
你已經贏了.
大家有沒有看聖經.
每一個聖經人物.
在他的過程當中.
其實都是.
考驗你的信心.
亞伯拉罕.
被告訴他有個兒子.
是考驗到他100歲的時候.
才有個兒子.
但最後.
仍然是按照上帝的旨意成就.
大兄姐妹.
過程不是上帝贏和輸.
而是你有沒有上帝.
你信不信上帝.
這是非常重要的.
第二件事在他的角度.
唯一要我們做的是什麼呢.
是傳道.
和牧養.
牧要傳道.
無論得時不得時.
總要專心遍用百般人來.
各樣的教訓責備人警戒人勸勉人.
因為時候要到.
人必厭煩純正的道理.
以多發養就隨從自己的情慾.
天天抱惜師父並且掩耳不停.
爭道偏向荒謬的言語.

$^{321}$你卻要凡事謹慎.
忍受苦難.
做傳道的工作盡你的職份.
無論順.
或逆.
大兄姐妹.
我想告訴你家欣樓有一種很重要.
我們核心的精神.
大家知不知道嗎.
其實我們在每一次分享和講道德都不斷告訴大家.
我們現在做的工作並不會去看家欣樓是否繼續的.
我想告訴大家.
如果我們看著的話,咖啡店不會出現的.
都預備倒閉.
你還整來幹什麼.
浪費了錢.
不是.
是我們很清楚要做的就做.
所以在疫情期間.
我告訴大家一件很感動的事.
當全香港的教會.
還在籌算.
不知如何是好的時候.
我們在核心開會的時候.
第一件事我們知道.
不要讓時空停下來.
因為這一刻.
弟兄姊妹更需要牧羊.
人的心靈更需要.
所以如果你記得的話.
當時我們開了很多網上的課程.
多了很多網上的講座.
我們即時聯絡了很多不同的講員.
在裡面的時候.
叫弟兄姊妹不要停.
要繼續.
效果是怎樣.
我可以告訴弟兄姊妹.
效果我們是不知道的.
因為其實大家也明白網上的世界.

$^{361}$點擊率有100.
是不是有100個人坐下來全程聽你講.
我們是不知道的.
但我們終於做了.
但有一件很奇妙的事發生.
就是當時.
錦繡堂還是經濟緊張的時候.
還在慰藥.
大家也知道.
疫情加上移民.
全香港都受到衝擊.
我們所有的東西都不夠.
我們奉獻收入的時候.
突然間有一個不是我們教會的弟兄.
間中就會回來錦繡堂.
去問曾牧師.
曾牧師我想給一個大額的奉獻.
給教會可以嗎.
當然可以.
為什麼你會這樣做.
他說了一番話.
我今天也記住.
是上帝給我們的回應.
我在外面的教會.
我的舞會.
在疫情的裡面.
什麼都沒有做.
停止了.
但我看到你們很努力.
我看到你們.
繼續將這些牧羊推到網上.
你們的支持.
突然間有一筆很大的奉獻.
就來到.
弟兄姊妹.
聖經說的.
一生的果效是隨著你.
去出現的.
你對上帝有多少信心.
你對上帝有多少的時候.

$^{401}$你就承諾去做.
所以保羅說得很清楚.
保羅在當時是疫還是信.
疫.
在人們眼中的時候.
他是在輸的.
被人抓的.
預備死的.
但他對提摩太的提醒是什麼.
無論得時不得時.
你都要傳道.
弟兄姊妹.
我們有這份傳承的精神.
弟兄姊妹.
你們有沒有.
我們有沒有在我們生命裡面.
都很確定.
無論信還是疑.
上帝的教導.
我們都要傳達.
無論信還是疑.
上帝都在.
所以我們就憑著這份信心.
繼續向前行.
第二,專心.
不要理會外在的環境.
你不要理會別人.
聽不聽,爛不爛.
而是你傳不傳.
傳福音在歷世.
歷代都難的了.
我好記得我來到這裡的時候.
第一次和弟兄姊妹交流的時候.
我第一個收到的訊息就是.
這裡好難.
然後我可以數到100樣的難.
是不是.
這裡不同.
錦繡堂的時候有個湖.
人們走出來,有個商場.

$^{441}$人流都不同.
這裡怎樣,大家住在這裡.
晚上鬼影都沒有.
時間點對點,由家去到.
百佳那裡買完東西就立刻回家.
人都不停留下.
人流的,怎樣做.
弟兄姊妹.
你記不記得在疫情之前.
我們有這個情況.
如果我錯了.
請一些資深弟兄姊妹跟進我.
如果我來到的時候.
如果我沒有印象,沒有記錯.
當時的崇拜人數,一般普遍來說.
60多到70.
是當時崇拜.
但是.
大家記不記得世界盃.
這些沙發是世界盃搬進來的.
後面的桌子是世界盃搬進來的.
是慢慢因為世界盃這個時宮.
而形成我們今天有這種文化.
世界盃裡面.
曾經試過總決賽那晚.
超過200人.
所以世界盃的時宮裡面.
帶動了我們這間樓.
曾經在崇拜裡面.
衝到130人.
大家記不記得.
其實我想告訴弟兄姊妹.
從來因為這個核心價值.
我們都不是看難與不難.
而是要做就要做.
如果大家當日有份參與.
的時候,世界盃是我剛剛入職.
我7月1日入職.
我入職的那個星期裡面.
我記得那晚通宵.

$^{481}$我記得有兩晚我睡在祈禱房裡.
因為世界盃結束了.
已經凌晨4點.
我好記得,Carol那時我第一次認識.
她煮東西很厲害.
我走過來時,她煮東西很厲害.
我才知道,當時我才相認.
原來你是肥沙的老婆.
當時如果你看回那些片.
當時我們拍了一段片.
大家是同心齊心合力.
一起去做的.
我們整條走廊擺滿了食物.
弟兄姊妹完全不吝嗇.
將她所有的.
我們就奉獻出來.
弟兄姊妹,這個叫做神的價.
所以你就會發覺,在這裡看的時候.
不是看外在環境.
而是用盡你一切的方法.
這裡寫的就是.
教導,警告,勸戒.
我們要去做的是牧羊.
和做全方位的工作.
弟兄姊妹,佳音樓.
最核心,我們做的每一樣東西.
什麼方法都不重要.
你可以打高球.
你可以咖啡.
你可以帶去旅行.
你可以玩音樂.
你可以教唱歌.
你可以去教書法.
所有的方法都無所謂.
目的都要回到原一.
傳福音和牧羊.
就是說,那些朋友來到的時候.
他能不能夠從你所教你的人的生命裡.
感受到福音.
弟兄姊妹參加這個課程.

$^{521}$他能不能夠被牧羊.
所以我一直在和大家宣傳.
我們要的咖啡大師.
不只是會沖咖啡.
而在裡面你要知道.
一杯咖啡的含義.
不只是咖啡.
而在背後是給人的安慰.
溫暖和鼓勵.
你可以藉著一杯咖啡.
坐下來和他慢慢聊天.
你可以坐下來.
和他慢慢聊天.
大家記不記得.
在世界盃之後.
很多教會是不容許的.
我記得在座堂的教會.
不要在禮堂裡吃東西.
大家都覺得很興奮.
我不是說不對.
我想說每個地方的特色.
但我們在世界盃之後.
有一班朋友是未信耶穌的.
他們不會崇拜耶穌.
但我印象很深.
當時我們後面的場地.
全部坐滿了.
每個人手上都有一杯咖啡.
是一邊喝咖啡一邊崇拜.
這樣都可以?.
弟兄姊妹,不要停留在那一點.
我想告訴大家.
我見到有很多弟兄姊妹.
從那張桌子坐進來.
從他們做一個局外人.
變成我們的弟兄姊妹.
變成我們的會眾.
弟兄姊妹.
將你的眼光放遠.
但不要偏離.

$^{561}$當你很清楚有這種核心價值的時候.
我們的心可以很廣,可以很闊.
上帝說得很清楚.
要將你的障目延伸進去.
另外一樣.
保羅是用他的信心去見證.
他說.
我現在被囂張.
我離世的時候.
到那裡美好將我打過.
當跑過路我已經跑盡.
所信的道我已經守住.
從此以後由公義的冠冕為我存留.
就算按照公義審判的主.
到那一天要賜給我們.
不但賜給我.
也要賜給凡外無他顯現的人.
我們在互相勸勉.
剛才我都提過.
就算我做傳道人.
我還有很多不足的地方.
你問我我的模樣做得好不好.
我坦白說.
我做得不夠好.
上次我們在團體裡面的交流.
有些大影子會提到.
我們這裡的模樣不夠.
其實我心裡面一層.
我都在想.
我會不會做了另外一些事.
而去耽誤了我去做模樣.
但我們在這裡.
我們要給予一次勸勉.
是甚麼呢.
就是到有一天.
當你要離開的時候.
你離開之前.
保羅說.
你要打美好的仗.
你要跑盡.

$^{601}$所能跑的路.
你要守住.
你要守的道.
你要.
保羅用他的生命.
我們也見證著.
他的生命.
他能夠在很多人面前說.
我已經做到.
弟兄姊妹.
我們還有機會去做.
以前做得不好.
不要緊,我們還有時間.
我們還有機會.
弟兄姊妹.
要打美好的仗,要跑盡你的路.
要守住你的道.
第二件事,對你的未來要有盼望.
因為上帝.
會有永恆的獎賞.
等著我們.
我們每一個人.
要去得到的獎賞.
如果你是確信.
你自己離開這個世界.
之後回到上帝那裡.
你就一定要等候上帝這個美好的獎賞.
但昨天在團期.
我和弟兄姊妹分享.
為甚麼今天上帝在地上.
仍然給我們時間.
為甚麼他仍然給我們時間.
透過這些去塑造我們.
去塑造我們認識他.
塑造我們能夠愛慕他.
你可以想像一下.
如果你現在都這麼討厭.
教會裡面.
無論唱詩歌.
無論你覺得.

$^{641}$我又要這樣說.
我想說這個不是家樓的現象.
這是香港普遍的現象.
一聽到講到的時間.
不要講這麼長.
45分鐘了還講.
講短一點.
我想告訴弟兄姊妹.
我之前經常去宣教.
只有香港的弟兄姊妹叫我講短一點.
我以前去國內宣教.
我們原本的計劃.
這麼短時間當然要濃縮.
我們就要去港島.
我們想去港勝景.
當時的環境還可以在台上.
我可以告訴你.
鄧全都曾經試過台下有500人.
國內的弟兄姊妹很熱心.
拿著自己的椅子回來教會.
坐在椅子上聽.
我們原本的計劃是什麼.
每兩個小時拿一個break.
之後才吃飯.
但我的經歷是我想到國內的時候.
跟他們講了兩個小時.
主持人說現在我們休息了.
弟兄姊妹走過來說不要停.
不要停我想聽.
接著OK繼續講.
都吃飯了.
不要緊拿著飯來吃.
你吃飯跟我們講.
我最高的紀錄.
是在上面連續講了11個小時.
弟兄姊妹我不是責備我們.
我不是責備.
我是想講的是我們今天有沒有.
渴慕上帝永恆他自己的掌上.
我再強調.

$^{681}$有一天.
你現在所賺到的地上的東西.
你要放下.
你想到神裡面.
你得到什麼.
是看你今天.
你為上帝做什麼.
還有你是不是屬於他.
如果你在地上.
你已經這麼討厭或不喜歡上帝的東西.
你上到天堂.
我想告訴你.
你就大件事.
誇張一點.
弟兄姊妹我不是說賭錢.
到今天仍然很喜歡四方城.
玩玩.
衛生麻雀.
和朋友.
上到天堂沒有麻雀.
很大件事.
你今天你現在都不渴慕上帝的地方.
你上到去.
你渴慕上帝的天家.
你不要玩.
所以弟兄姊妹.
上帝在地上.
給你經歷很多東西.
重新給你去調教你的生命.
我很看這次我的病.
這個手術是上帝對我一個很大的調教.
我都問我做了什麼.
剛剛紮正這個地方.
這個地方很有趣.
醫生都說.
很少會剛紮正的位置.
左邊影響我的運動神經.
右邊影響我的感官神經.
很奇怪.
但就是在那裡.

$^{721}$但問題是.
上帝讓不讓許可.
上帝的許可.
讓我好起來.
繼續去做的話.
那你想怎樣過啊鄧全堂.
你想怎樣繼續你的生活.
所以弟兄姊妹到最後的時候.
保羅說.
他是次級那愛無他顯然的人.
請大家記得這句說話.
聖經說過一個圖畫.
我們不要那麼低層次.
究竟我信耶穌將來我得不得救.
聖經的圖畫就是.
有些人因為信耶穌的緣故.
你得救的.
不過像什麼.
像一個在火堆裡搶出來的炭.
剛剛好救到你.
上到去之後.
那個樣子.
你想嘆的樣子.
弟兄姊妹.
我們想上到上帝的面前.
我們成為什麼.
所以.
我很盼望.
我也求上帝.
我也不斷祈禱.
請大家也為我祈禱.
主若許可.
我可以做這事.
或做那事.
我也很盼望.
我能夠繼續為上帝工作.
所以請大家繼續祈禱.
為我收網.
但這是上帝的主權.
交給上帝.

$^{761}$但我盼望.
如果.
你就當我說最後一句.
如果上帝接我回去.
弟兄姊妹.
我很盼望你傳承.
這一份家因流的精神.
讓家因流.
繼續.
行神的路.
不要讓家內的貨.
關了.
繼續去做神的工作.
不要浪費時間.
我一直同心祈禱.
七天復活求主.
你親自保守和幫助.
我們的生命.
在你的裡面.
是你掌管的.
所以我們全部交給你.
我們知道有你帶領之下.
我們也不舉辦.
但我盼望神.
保守著家樂和每個弟兄姊妹.
願意我們一起.
很努力地繼續走你的路.
守住你的道.
打美好的仗.
直至到有一天.
我們在這裡.
我們沒有悔.
我們沒有遺憾.
我們能夠完成上帝的工作.
將一切所有交代給你.
求你親自保守.
請你奉教主耶穌.
\newpage



\section{彼得後書 1:3-8}
\label{sec:Bres3qEZ6Yk}
\textbf{2023-09-03 午堂《在信仰上加倍努力》彼得後書 1\_3-8 講員:曾裔貴牧師}
\newline
\newline
連結: \href{https://youtube.com/watch?v=Bres3qEZ6Yk}{\texttt{https://youtube.com/watch?v=Bres3qEZ6Yk}} ~~~~ 語音日期: 2023-09-05
\newline
\newline
\hyperref[sec:ay_HsjJTAHo]{\small{< < < PREV SERMON < < <}}
~
\hyperref[sec:index_chronic]{\small{[返順時目]}}
~
\hyperref[sec:index_scriptual]{\small{[返順卷目]}}
~
\hyperref[sec:naPXdHer1J8]{\small{> > > NEXT SERMON > > >}}
\newline
\newline
彼得後書 1:3-8
\newline
\begin{longtable}{cl}
\hline
\hline
章節 & 經文 (和合本修訂版)\\
\hline
1:3 & \begin{tabularx}{0.7\textwidth}{X} 神的神能已把一切關乎生命和虔敬的事賜給我們，因我們認識那用自己榮耀和美德召我們的神。 \end{tabularx} \\ \\ \relax
1:4 & \begin{tabularx}{0.7\textwidth}{X} 因此，他已把又寶貴又極大的應許賜給我們，使我們既脫離世上從情慾來的敗壞，就得分享神的本性。 \end{tabularx} \\ \\ \relax
1:5 & \begin{tabularx}{0.7\textwidth}{X} 正因這緣故，你們要分外地努力。有了信心，又要加上德行；有了德行，又要加上知識； \end{tabularx} \\ \\ \relax
1:6 & \begin{tabularx}{0.7\textwidth}{X} 有了知識，又要加上節制；有了節制，又要加上忍耐；有了忍耐，又要加上虔敬； \end{tabularx} \\ \\ \relax
1:7 & \begin{tabularx}{0.7\textwidth}{X} 有了虔敬，又要加上愛弟兄的心；有了愛弟兄的心，又要加上愛眾人的心。 \end{tabularx} \\ \\ \relax
1:8 & \begin{tabularx}{0.7\textwidth}{X} 你們有了這幾樣，再繼續增長，就必使你們在認識我們的主耶穌基督上，不至於懶散和不結果子了。 \end{tabularx} \\ \\
[1ex]
\hline
\hline
\end{longtable}
$^{1}$我們一同祈禱.
感恩天父讓我們有很快就有美好的天氣.
我們能夠帶著一個愉快的心靈.
來到神的殿敬拜.
剛才的詩歌要我們發光.
當時的背景是一些船遇到海難.
因為燈塔沒有了光.
沒有了紫路的明燈.
我們人生都是向著永恆.
我們能夠有主耶穌在我們的生命裡.
幫助我們能夠在整個人生裡.
我們可以活出神的聖情.
我們一同在你面前等候.
這段經文都相當握要告訴我們.
基督徒的靈命成長的階梯和步驟.
幫助我們不要小看這些經文.
帶給我們的生命改變.
將來我們可以豐豐富富.
因為我們的生命就是始我們的主.
進到永恆裡面.
我們這樣禱告奉耶穌基督的名.
阿們.
這個月份九月.
因為十月就是堂慶.
九月我們牧者我們童宮.
都真的很想弟兄姊妹.
那個靈命門徒生活.
尤其是我們怎樣去裝備自己.
可以在人生被神使用.
所以這個第一講門徒的生命.
相信這裡剛才主席所讀了的.
尤其是第五到第八節.
總共講了基督徒有八樣的生命成長特質.
信心.
然後就要有一個德行.
然後你要有知識.
無論地上的知識.
屬靈的知識.
你能夠掌握運用知識的智慧.
然後你又有節制.

$^{41}$然後你又有忍耐.
對人對其他人的忍耐.
然後你內心裡面有一個很大的虔誠.
再大一點的就是愛弟兄和愛眾人.
相信斯圖比德將這八樣東西寫出來.
他自己充分地在他的人生.
由他放下漁船漁網.
跟隨耶穌之後.
他就是這樣走上這條神聖的道路.
有三年的時間跟隨著耶穌.
往後因為主耶穌升天.
他就要獨當一面在耶路撒冷.
輾轉去到不同的地方.
後來主要主力都是在耶路撒冷.
來作當時的教會領袖事奉.
帶領牧羊教會.
他的內心一定是因為這個身份.
而經歷很多很多的磨難.
很多的挑戰.
但領袖他這一刻寫的書信.
告訴我們原來.
成聖就需要這樣去按部就班.
有料要加上有料要加上.
中文的翻譯都是很想大進這種.
彼此這些不單不是互為排斥.
而且是整個生命互為調和.
其實就是一個神的性情有份.
你和我有份於神的性情.
還有最終十一節.
每個人豐豐富富進入主救主榮耀的國度.
今早又看新聞你都知道.
都是香港的光榮.
梁朝偉先生.
李安導演說.
每一部戲只要有他的出現.
整個電影就會生息不少.
我想大家都同意.
如數家珍你可以數到很多很多.
你想想葉問已經有固定的形象.
他都可以將新的形象.

$^{81}$給我們打破舊有的觀念.
可想而知李安導演說.
他一個眼神就等於很多很多.
其他的演員用了很多對白.
都未必能夠將那個角色.
深度地在觀眾面前呈現.
41年的敬業樂業.
終身成就獎.
彼得說我快離世了.
第十五節.
並且我要盡心竭力.
使你們在我去世以後.
時常紀念這些事情.
紀念什麼事情.
循序漸進的基督徒的生命.
這八方面的特質.
有沒有越來越互為調和.
越來越深化增長.
彼得說我就快要去世了.
我離開這個世界了.
他告訴我們.
他是屬靈的一個富有的人.
豐豐富富的進入神的國度.
其實他是殉道被斬頭.
對不起.
他是最終要被釘十字架的一位使徒.
他告訴我們.
不是什麼神秘.
不需要有什麼特別的秘笈.
更加不需要有些奧秘的事情.
你才可以來到有神的聖情.
有份於神的聖情.
其實是有的.
十八節.
我們和他在聖山的時候.
親自聽見這些聲音.
就是神和他的愛子說.
這是我的愛子.
我所喜悅的.
他這樣來說.

$^{121}$是我們親自聽見.
從天上出來的.
我們並有先知更加確的預言.
你可以想像得到.
彼得有他個人.
因為使徒.
因為神特別的一些工作.
在他的身上.
有很特殊的神跡.
他見過五餅二魚神跡.
吃過那些餅和魚.
他也看著在海裡面.
本來是沉的.
但耶穌說你走過來吧.
他真的走過去.
耶穌就捉住他上船.
經歷過這些神跡.
是他在那三年裡面的經歷.
但他現在就快離世之前.
告訴我們基督徒.
只需要認定你的起點在哪裡.
和你的終點要帶什麼走.
起點在哪裡.
就是三節.
神的神能已經將一切.
關乎生命前景的事.
賜給我們了.
皆因我們認識.
這位用自己榮耀和美德.
召我們的主.
就是這個恩典的選召.
就是我們信耶穌.
打開我們的心靈.
這就是那個起點.
每個人都是一樣的.
每個基督徒.
每個弟兄姊妹.
你都是這樣.
將你的心靈打開.
得到這個選召.

$^{161}$耶穌的救恩.
但是美德告訴我們.
有一個很長的生命階梯.
我們要進入.
你要最後得到的就是.
第十一節.
這就是第十節所說的.
只要我們肯這樣做的時候.
就使我們所蒙的恩召和揀選.
就會堅定不移.
頭尾都說到恩典的揀選和呼召.
然後這個揀選.
原來是有一個深度.
就是被揀選.
蒙召信主的弟兄姊妹.
就是使我們如何去到將來的終極.
有進入永恆的豐富.
原來端視是第十節的重點.
接著就是這樣說.
你們如果行這幾樣.
就是剛才十五節說過的.
你們行這幾樣.
就是上文所數的八個靈性的特質.
你們是認真每天.
在我們不同的處境去學習的時候.
接著就說永遠不失腳.
這樣就是這個.
.
說了的起點.
大家是完全一樣.
但是仕途彼得說.
有一個的終點.
他期望每個都是終身成就.
每個弟兄姊妹.
有你自己的終身成就.
你如何去看待生命的這個課題.
有沒有在我們的生命裡面.
彼得已經列舉了.
不是你見過甚麼異象.
不是你經驗過甚麼神蹟.

$^{201}$甚至不是你自己能夠行神蹟.
好像上天下地那樣.
不是.
就是生命我們裡面.
因為我們已經有份於神的性情.
如何讓這個性情不斷的去發展.
成為我們自己的成熟基督徒生命呢.
信心加上德行.
信心是在你的內心裡面.
未必一下子看到你的內心有這個認信.
但是一個信了主的弟兄姊妹.
就開始在德行上會有不同的表現.
每個人都有不同的快慢.
我很快接近信主的那半個月.
突然間有一個念頭.
不要吸煙了.
一天吸兩包煙.
突然間不吸.
從此就甚麼都沒有了.
聞到味道我都會有些反感和有些俗的.
這麼奇妙的神跡性的一種醫治.
在我的內心.
完全沒有煙癮.
完全沒有困難.
很奇妙的工作.
神可以這樣做.
但每個人都不同.
有些人面對不同的事情.
他可能要很漫長地對付.
但是這裡告訴我們.
信心德行是一對的出現.
基督徒很重要就是對神.
對這個宇宙萬物的認知.
特別是第二章第三章.
主在來的應許是確實的.
有人以為他擔言.
其實不是擔言.
是寬容我們.
主漢一日如千年.
千年如一日.

$^{241}$即是說一直要有一種這樣的眼光.
來貼近神的心牆.
我們才能敏銳對永恆事物的真確.
原來知識是需要不斷地與時俱進.
有了知識的人.
原來在基督徒的生命裡面.
是會有自我節制的.
兩者一條和.
就是一個很成熟的基督徒.
有知識的人.
可以游刃有餘地對不同的人.
你不會很驕傲.
你會很能夠有一種同理心.
對別人.
尤其是一些沒有學問的人.
你會有一種尊重,欣賞等等.
如果你沒有了聖經所說的.
這一種知識.
而你的內心沒有真正的.
在主的生命裡面成長.
那就變成很沙塵,很炫耀.
和不懂駕馭自己的知識.
真正的知識是會帶來節制.
然後還會有忍耐.
和加上有虔誠.
你想想你去揣摩一下.
默想一下.
忍耐加虔誠.
是怎樣的基督徒呢?.
那種成熟的品味是怎樣的呢?.
還有就是.
在內心發自的愛弟兄,愛眾人.
你想想一個基督徒.
彼得就是一個出生基層的漁夫.
耶穌在加利利海邊.
來去呼召他.
他就撇下跟隨耶穌.
就開展他的人生.
仍然是一個猶太人.
但是他能夠成為一個愛愛邦人.

$^{281}$哥尼留事件.
舌頭行轉的第九章.
本來在那裡全度探訪.
有一個燒皮匠西門的家.
準備等午飯.
突然間見到異象.
天打開.
就有個聲音跟他說.
有很多地上的走獸.
叫他起來撞了一下.
是一個異象.
異象他都會說話,對話.
煮呀不行的.
這些逐而不結的食物我不可以吃的.
因為我是猶太人.
我的血統留著的是利未記的經文.
告訴我小時候不敢吃這些東西.
起來撞了一下.
神所結證的.
你不可以再視為不結.
這些異象這樣跟他說三次.
接著就有人來訪問.
就是哥尼留的那些丁丁.
有沒有一個人叫做彼得的.
我們要找他.
你可以想像得到.
彼得深深的去懊悔.
為什麼我慢過神那麼多.
神愛世人.
我作為使徒傳道一段這麼長時間.
我看見一個外邦人.
我會有一種格格不入.
會有歧視.
我正不是傳道要告訴人.
神的救恩嗎.
你去到那個地方.
哥尼留和他的親屬密友.
和那些丁丁完全聚在一起.
有天使跟他說.
會有人來跟你說基督教的福音.

$^{321}$哥尼留原來是一個很虔誠的羅馬的白夫長.
比彼得更加的渴慕認識神.
這一次的教訓.
神就要來點醒彼得.
神所潔淨的.
你不可以再捍衛.
逐而不潔.
晚年的時候.
彼得能夠真是伸手.
讓那些羅馬的丁丁.
因為當時都是尼祿的年代.
來處他死刑.
第三章有說.
我要走的日子到了.
正如主指示我的日子.
這是約翰福音21章.
彼得看不過眼.
約翰是怎樣的呢.
耶穌說.
我要他等到我回來.
與你何干.
你跟隨我吧.
晚年的時候.
耶穌預言彼得願意伸出他的手.
不抗爭.
你要叫我不信耶穌.
不見證耶穌是不行的.
但你要抓.
你要殺是無所謂的.
因為我是要跟隨耶穌的.
這樣的屬靈生命的極致.
這就是屬靈的風風負負.
走在神的旨意裡的使徒彼得.
所以弟兄姊妹.
今早.
穆者與大家彼此勉勵.
我們有沒有循序漸進.
由信主的開始.
裝備這八樣神聖情的特質呢.
信心要加上德行.

$^{361}$然後要有知識和節制.
忍耐要加上虔誠.
愛是要讓人看得到的.
你可以想像得到.
如果我們能夠真的站在一個地方.
那個地方突然間.
人們不敢再亂說話.
因為我們的出現.
帶來一種潔淨和虔誠.
這也是很重要的基督徒的品格.
讓人能夠看得到.
所以很盼望神將恩典給我們.
重要的是彼得在這本書.
講了三次很重的三次.
要恩勤.
第五節.
一章的第五節.
正恩者緣故.
正恩者緣故的意思就是上文說.
我們已經因為這個恩典的選照.
脫離了世界從情慾來的敗壞.
即是說我們已經恩信稱義.
得到永生.
既然我們已經不再是在敗壞行列裡.
已經與神的聖情有份.
然後就說正因為這樣的緣故.
你們就要分外的恩勤.
第十節.
所以弟兄們.
應當更加的恩勤.
使你們所蒙的恩照和揀選堅定不移.
原來信了主耶穌之後.
我們如何運用我們的人生.
這個恩勤加倍的努力.
正正就是使我們在這些生命的特質裡.
彼此不是互為衝突.
是慢慢慢慢互為調整.
互補.
和越來越有一種和諧.
活出神的聖情.

$^{401}$到三章十四節.
親愛的弟兄.
你們既然盼望這些事情.
上文說的就是新天新地.
必然會在地上出現.
使自己沒有玷污無可指責.
然後我們安然的見主.
很有平安.
很坦然的面對我們人生的終局.
或者是主的再回來.
為什麼呢.
因為我們已經是豐豐富富.
我們已經是終身的成就.
是大家都能夠認同的.
盼望頂尖會能夠在這些循序漸進裡面學習.
我記得就是.
因為我們當時入學.
我們有19個同學.
有些讀過大專.
但大部分的同學和我這個年紀的.
我們只是讀過中學.
可能中六我們都沒有讀過.
突然間要學寫論文.
簡直沒辦法想象.
我第一次上系統神學.
出了名的老師是Killer.
給分數是非常緊的.
但也要求很高的.
我給大家知道就是廖炳棠牧師.
其實是人都知道的.
很緊的分數.
上他的系統神學.
可以肥你.
你就要重考的了.
第一年.
因為連續三季都是他教.
第一年對我們這些從來沒有寫過論文的.
是什麼叫論文.
什麼格式.
完全不懂的.

$^{441}$我蠢到一個地步就是.
以為那些footnote下面的線是用手畫的.
原來按個按鈕就會出現的了.
在office裡面的work.
完全不懂.
但我想說這件事很緊要.
他知道我們的程度.
第一堂就說了.
未來的十堂.
是會教什麼的.
這一季.
接著很快就會講功課的了.
會給三至四條的.
大的論文題目.
你們就去做了.
對我們那個的幫助.
就是第三堂.
就要立即交論文的大綱.
和你的書目.
和你為什麼要這樣做.
你有什麼論點等等.
逐個同學去見.
當然寫慣論文的同學.
或者他之前已經讀了夜校課程的.
就沒問題了.
但對於我們這些是從來沒有寫過的.
就要很仔細的教我們.
那個格式和正反對比的論述.
和最後你有什麼自己的論點.
這樣去講.
這樣去做.
好像抓住你的手.
循循善誘那樣.
很辛苦.
很困難.
第三年的系統神學.
我居然是全班拿A的.
簡直沒法想像.
我連那個footnote按哪個按鈕.
原來可以按按鈕就有個footnote出來.

$^{481}$我都不懂.
但可以能夠慢慢再看那些功課的論文.
真的連那些碩士班的同學都覺得寫得很好.
這個正正就是我體驗到.
循序漸進的學習.
一個好的老師.
在你的生命裡面.
沿途的那個指引.
使徒彼得.
不跟我們講什麼高賢大智.
我們見過主的威榮.
親耳聽過天上來的聲音.
他不是跟我們講這些特殊奧秘.
講的就是我們能不能夠每日.
以我們基督徒的恩典的選照這個身份為榮.
有這份榮耀感.
幸福感.
然後在生活裡面.
你看福音書.
你看《少道行傳》.
多數是困難才孕育我們的信心.
如果沒有任何的困難.
我們不會發展我們的靈命的成長.
其實任何都是.
所以盼望.
敢孫家的弟兄姊妹.
我們能夠在一個這麼優越的環境.
不是每個地方的人都可以這樣.
我跟林牧師去印尼.
來到他短孫的時候.
他開車載我去一個教堂.
他說這個教堂就是我第一次出來傳道.
在那裡事奉的.
他說他剛剛出來傳道的時候.
附近周邊有些穆斯林.
去焚燒這個教堂.
這個是新重建的教堂.
這樣面對你在那個事奉裡面的困難.
你可以想像得到.
不是每一個環境都歡迎基督教.

$^{521}$但是基督徒就是有這個很清楚的成長階梯.
信心加上你的品格 德行.
你有知識.
你更加有節制.
你能夠有忍耐.
無論對環境 對別人.
因為你內心有一個很大的虔誠.
愛充滿了你的內心.
因為你已經擁有神的性情 生命.
很盼望神祝福我們每一位.
最後我想用一個例子.
我都講過一次.
有一位叫黎彥平的牧師.
其實我不認識他的.
不過在我們中間婚禮的正道.
我坐在台上聽到他說.
很深的感受.
他訓練一對新人.
是他教會的弟兄姊妹.
他說什麼叫愛.
他說了他媽媽.
黎彥平牧師的爸爸很早就離開世界.
但這位母親對丈夫的愛.
要他們的兒女都要銘記.
每年的爸爸的壽辰.
或者是死忌.
都要全家來去墳墓 掃墓.
來去紀念.
這位牧師說他很深刻.
媽媽已經很大年紀了.
走不了了.
但堅持要去.
走上很多級樓梯.
去為著子紀念丈夫.
他就很感受到媽媽對丈夫的愛.
那種見證.
哇 我很深刻.
我不知道那對新人是不是很深刻.
愛就應該是這樣.
親愛的弟兄姊妹.

$^{561}$神對我們的愛.
就是這麼堅貞不移.
就是希望我們每一個弟兄姊妹.
好像彼得恩晚年的時候.
要殷勤 要殷勤成長.
他回來了.
要殷勤裝備我們自己.
成為屬靈的富有的人.
進入神的國度.
我們在他面前拿的是終生的成就獎.
我們一起祈禱.
多謝主讓我們能夠透過這段簡單的經文.
彼此的提醒和勸勉.
盼望主的聖靈感動我們.
奉主耶穌基督的名.
阿們.
\newpage



\section{撒母耳記下 23:13-19}
\label{sec:naPXdHer1J8}
\textbf{2023-09-10 加恩樓《大衛與三個沒有名字的勇士》撒母耳記下 23\_13-19 講員:陳韋安牧師}
\newline
\newline
連結: \href{https://youtube.com/watch?v=naPXdHer1J8}{\texttt{https://youtube.com/watch?v=naPXdHer1J8}} ~~~~ 語音日期: 2023-09-12
\newline
\newline
\hyperref[sec:Bres3qEZ6Yk]{\small{< < < PREV SERMON < < <}}
~
\hyperref[sec:index_chronic]{\small{[返順時目]}}
~
\hyperref[sec:index_scriptual]{\small{[返順卷目]}}
~
\hyperref[sec:TXT2__xkqM8]{\small{> > > NEXT SERMON > > >}}
\newline
\newline
撒母耳記下 23:13-19
\newline
\begin{longtable}{cl}
\hline
\hline
章節 & 經文 (和合本修訂版)\\
\hline
23:13 & \begin{tabularx}{0.7\textwidth}{X} 開始收割的時候，三個侍衛下到亞杜蘭洞，到大衛那裡。非利士的軍兵在利乏音谷安營。 \end{tabularx} \\ \\ \relax
23:14 & \begin{tabularx}{0.7\textwidth}{X} 那時大衛在山寨，非利士人的駐軍在伯利恆。 \end{tabularx} \\ \\ \relax
23:15 & \begin{tabularx}{0.7\textwidth}{X} 大衛渴想著說：「但願有人從伯利恆城門旁的井裡打水來給我喝！」 \end{tabularx} \\ \\ \relax
23:16 & \begin{tabularx}{0.7\textwidth}{X} 這三個勇士就闖過非利士人的軍營，從伯利恆城門旁的井裡打水，拿來給大衛喝。他卻不肯喝，將水澆在耶和華面前， \end{tabularx} \\ \\ \relax
23:17 & \begin{tabularx}{0.7\textwidth}{X} 說：「耶和華啊，我絕不做這事！這三個人冒生命的危險，這不是他們的血嗎？」大衛不肯喝這水。這是三個勇士所做的事。 \end{tabularx} \\ \\ \relax
23:18 & \begin{tabularx}{0.7\textwidth}{X} 洗魯雅的兒子，約押的兄弟亞比篩是這三個勇士的領袖；他曾舉槍殺了三百人，就在三個勇士中得了名。 \end{tabularx} \\ \\ \relax
23:19 & \begin{tabularx}{0.7\textwidth}{X} 他在這三個勇士中是最有名望的，所以作他們的領袖，只是不及前三個勇士。 \end{tabularx} \\ \\
[1ex]
\hline
\hline
\end{longtable}
$^{1}$午安,鄧小妹.
很開心我們能在一起敬拜.
今天的講道是《撒武爾記》下的經文.
這篇道理我上星期已經講過了.
如果不幸上星期打風聽過的話.
沒問題的.
因為這篇道理雖然跟我息經.
但我覺得這篇道理更為加恩堂.
加恩樓頂智報而設.
所以今天講的向導是跟大家一起想想我們的處境.
所以今天我們會講《撒武爾記》下的大衛.
和三個幫他打水的人的故事.
基本上這個名稱都講了整個故事.
我們就一起思考一下.
對我們的信仰,對我們今天加恩樓面對的情況.
我們一起思考.
多謝祈禱.
最後祈求你幫助我們.
今天你自己的聖言,你的說話.
來到我們這個群體裡面.
你讓我們去恭敬的.
來聆聽你自己對我們每一個弟兄姊妹的說話.
你自己的旨意.
你自己對加恩樓的旨意.
你自己來成就.
讓我們能夠去聆聽.
願意去康復在你面前.
我們一起來尋求.
讓你自己更加讓我們這個群體.
能夠去實踐你對教會的要求.
你對我們每一個人生命裡面的要求.
求主你這樣的幫助.
奉主名求,阿們.
先講講經文的背景.
十母耳記下第23章是十母耳記的最後一部分.
如果你看十母耳記的話.
你會發現整個十母耳記是順著的.
即是從哈拿生子.
到十母耳記的成長.
到素羅王被高納.

$^{41}$到大衛興起.
到素羅和大衛之間的爭戰.
然後大衛統一天下.
然後大衛作王等等.
然後就別氏巴,阿撒隆等等.
整個十母耳記是順著的.
按著猶太人的歷史一步一步的去記載.
但是來到十母耳記的最後部分.
你會發現十母耳記下的第21到24章.
卻不是這樣的.
你會發現最後的三章經文.
其實就好像一本書的附錄.
Appendix那樣.
是一個不跟之前的附錄內容.
很多不同的內容都編輯在最後的部分.
說說這個附錄包括什麼呢.
他無端端會加插了詩篇18篇.
大衛讚美耶和華的詩篇.
還有一段大衛臨死之前.
耶和華上帝曉語大衛的一段說話.
還有大衛數文的故事.
都在後面.
今天我們所說的.
大衛和他的勇士們的名單.
都是在這個附錄內.
為什麼這些勇士會有這麼多的遺訓.
竟然能夠記載在聖經十母耳記下的最後部分呢.
一個字就是勁.
這30個勇士的名字能夠被記載.
因為他們是很厲害的人.
他們是大衛王朝期下最能打的英雄.
如果是三個字的遊戲.
每個人都是85以上的武將.
每個人都是非常厲害的.
名副其實是戰狼30.
非常勇猛的30個勇士的名字.
今天要說的正正就是戰狼30其中的三個.
今天我們來看看.
13節.
修國的時候有30個勇士的三個人.

$^{81}$下到雅杜蘭洞見大衛.
非利士的軍兵在理凡欽谷安營.
那時大衛在山寨.
非利士人的防營在伯里群.
大衛學想說甚願有人將伯里群城門旁井裡的水打給我喝.
這三個勇士就闖過非利士人的營盤.
從伯里群城門旁的井裡打水拿來奉給大衛.
總括來說.
如果你用一句說話來總括整個故事的話.
這個故事基本上是一個foodpanda的故事.
是一個送外賣的故事.
三個勇士接了一張order.
然後就進入伯里群的井旁.
叫他外賣一杯井水給雅杜蘭洞的大衛.
沒有提示的.
基本上經文記載了這個故事.
當時如果留意文件的話.
寫著理凡欽谷.
這個在十武記下第五章記載過.
那時正是大衛快要統一天下的時候.
你看經文就看到.
基本上是埋齋.
能夠攻破這個伯里群.
這個非利士人的伯里群.
就能夠統一天下.
所以你看到.
這個正正就是大衛將要統一天下.
一個非常如日方中的大衛.
他用手機拿了一張order的故事.
我們今天嘗試用比較現代化的說法.
所以你看到這個外賣的故事裡面.
外賣的餐廳在哪裡.
就是這個伯里群城門的井旁.
外賣的地址就是這個阿杜蘭洞.
阿杜蘭洞和伯里群距離多遠呢.
如果你google map就能夠google到.
大概就是.
結果找到是28公里.
就是說你們步行的話.
大概要走6小時.

$^{121}$由這個阿杜蘭洞走過去伯里群.
去去6小時.
倒一倒水回來.
總共就12小時.
如果你不休息的話.
不過ok的.
這班是勇士來的.
他們是不用休息的.
所以你看到.
如是者這三個人.
真的就是南亞車手.
他們是中東人來的.
即刻由這個阿杜蘭洞出發.
前往伯里群.
中間是經過這個利凡欽谷.
又稱為個人谷.
敵方的陣營.
千山萬水去到這個伯里群裡面.
打水.
拿個水壺裝好.
用保溫袋裝好.
立刻就上車.
沿途再經過敵方的陣營.
再回到阿杜蘭洞大衛那裡.
就將這個外賣送給大衛簽收了.
一件非常之厲害的事情.
不過最令人覺得震撼的是.
當大衛接過這三位勇士.
千山萬水這張order的.
這壺冰涼新鮮.
在保溫袋拿出來的井水.
真真正正的千山萬水之後.
他竟然將整壺水.
倒在地上.
今天的重點不是來查究.
大衛為什麼要這樣做.
今天的重點不是這裡.
不過也姑且說一下.
一些醒學者對於大衛的行動.
或者整件事.

$^{161}$大衛有什麼想法和動機.
這些都是一些學者的推測.
第一個我們可以排除的.
就是大衛口渴.
因為基本上大衛不是口渴.
不是出於生理的需要.
所以要人幫他倒水喝.
基本上你看到雖然大委身處在阿杜蘭洞.
但這個山洞不淒涼的.
不是那些躲起來淒涼的山洞.
基本上是有水的.
起碼不需要跑這麼遠的路去拿水.
所以我們可以排除原因.
不是因為口渴的原因.
第二個可能的原因就是大衛思鄉.
原來大衛打發人去伯利恆拿水是什麼原因.
因為伯利恆正正是大衛出生的故鄉.
當然大家覺得伯利恆就是耶穌出生的地方.
但你知道伯利恆本身在舊里面.
被稱之為大衛之城.
所以正正就是大衛從小長大的地方.
所以大衛特別要拿一張order去伯利恆井水的話.
是因為想品嚐一下小時候古早味的井水味道.
就是這樣.
這是其中一個可能性.
第三 都有人認為比較陰謀論的推測.
大衛其實純粹是為了宣洩自己的權力.
當時大衛如日方中將要統天下的皇帝.
所以為了要展示自己的權力.
隨意地叫人冒生命危險拿水.
我們以前看色神都這樣說.
喜歡怎樣就怎樣.
這些就叫拋下辣.
這個就是楊雷思的情況.
這是一個學者的推測.
當然也有學者認為.
大衛這個舉動其實是暗藏一個軍事的戰略部分.
表面上是叫人去打水.
其實是為了試探敵方陣營.
來當一下菲律賓人的情況.

$^{201}$這是一份paper寫出來的看法.
不過這也不是我們今天要關注的地方.
今天我們嘗試從這三位勇士的角度去思考這件事.
無論大衛的動機是怎樣也好.
一個千真萬確無可否定的事實就是.
這三位勇士千山萬水將自己的生命豁出去.
幾經辛苦將井水打回來之後.
被大衛一下子倒在地上.
那一刻其實是不知道你會怎樣想.
如果你是那三位勇士.
你用盡你的努力將你的生命和心血.
冒著生命危險闖入敵方陣營.
打打水再回來.
倒給大衛喝的時候.
被人倒了你會有什麼感受.
基本上根據聖經記載.
聖經是沒有記載的.
聖經沒有記載這三位勇士的反應.
就停在那裡.
他們沒有拍到這三位勇士的表情.
不過如果是你的話.
當你花盡了你120\%的心血.
卻被人一下子倒在地上.
你是不知道怎樣去理解這件事.
百做百費心機.
徒然無為.
如果真的是foodpanda也OK.
基本上我也賺了外圍的錢.
吃不吃也不關我的事.
基本上我賺了就可以走了.
但當你真的將自己的心血.
好好地奉上.
然後那個人倒在地上的時候.
你會有什麼感受.
不過你會發現.
我們人越大.
我們面對這個世界.
你會發現原來這個世界.
這個社會裡面.
其實有很多類似的事情.

$^{241}$很多類似的情況.
幾經辛苦滿流血汗.
百百付出.
然後的結果.
是只有一個.
努力了.
付出了.
卻沒有帶來任何的結果.
運動員為了備戰奧運.
花了四年時間.
在比賽前突然傷了.
不能出賽.
在公司在學校裡面.
一個預備了半年.
很多晚通宵的項目.
突然上面說給預算.
就不做了.
五百多頁的計劃書全部扔掉了.
一個地方.
你在這裡建立了很多的心血.
付上了很多的.
每個星期的心血.
突然間說要停.
本身這篇道我在Full Church說.
因為我們Full Church的閱題叫做「上莊」.
不知道你有沒有上過莊.
大學的時候有沒有上過莊.
對我來說上莊是什麼呢.
上莊就是大學生一班人.
一起很努力很認真.
花很多晚通宵去做一些沒什麼結果的事情.
不知道大家有沒有這樣看.
很努力很努力.
但做出來其實沒什麼特別改變.
對世界裡面.
我這樣去定義我之前上莊的回憶.
很多晚通宵.
但結果未必很多人會看到.
甚至不太在意.
一起不斷地不斷地搶自己.

$^{281}$然後不求回報.
最後其實沒什麼特別做出來.
很大件事很風光偉大.
那不是的.
這不是我自己的看法.
我上網也試過.
搜尋一些上網的人對於上莊的原因.
網上有一個叫做大學生上莊七大原因.
看看網上有什麼原因.
第一就是認識莊園.
就是認識人.
第二就是學習如何運用軟技能.
就是學習.
第三就是拓展人脈.
就是認識人.
第四就是展示威風.
不知道為什麼.
可能拍那些照片.
現在數相很有型.
第五就是自取.
第六就是拿贖分.
第七就是出鋪.
原來七個原因裡面沒有一個關上莊的事.
當你做了一件事.
你那個原因是不關那件事.
那你想想那件事是什麼事.
不是負面的.
不是說上莊的人很敷衍.
或者是完全不懷企圖.
相反上莊的人是很熱血的.
很多晚上大家一起做一個part time.
然後就好好的去搶.
但其實大家也不是特別重視.
或者覺得ok的.
做不到什麼出來.
其實大家也不是很介懷做什麼.
重點是那個過程.
重點是大家這麼多年之後那班友誼.
就是你不斷在這裡消時間.
消GPA也好.

$^{321}$重點是大家曾經一起努力過奮鬥過.
這個就是很多人對於上莊的回憶.
暑假的時候我和一班教會的牧者.
就是流堂裡面的牧者一起去台灣.
去那裡退休.
我們經過台灣的夜市.
我們經過台灣夜市裡面的小時候彈珠堂.
一個彈珠店.
這個彈珠店基本上都是標榜.
是兒童回憶.
很特別很多部小時候玩的彈珠機.
玩的彈珠會是領獎品.
基本上你唱基本台幣.
然後你拿了很多彈珠.
然後你將彈珠放進彈珠機裡彈.
贏了就會拿更多彈珠.
然後你拿了更多彈珠就玩下去.
整件事是很無聊的.
就這樣玩很多彈珠贏去.
我們就大家一起享受了一個小時.
大家玩彈珠機.
不過我特別欣賞.
雖然無聊說無聊.
但特別欣賞的是什麼呢.
就是店裡面台灣文青味道的標語.
這個店裡面寫著.
很值得跟大家分享.
你不要看我們好像很無聊.
一群人在玩彈珠好像贏不了很多東西.
他就寫了幾句文青feel的標語.
他說很多事情都是這樣.
你認為值得就是值得.
就像玩彈珠一樣.
就這樣寫出來.
就這樣推銷這件事.
不要理我做什麼.
總之我覺得值就值.
第二句這樣寫.
就算很多事情沒有發生也沒有關係.
最重要的是我們一起玩彈珠.

$^{361}$度過美好的一天.
明明店裡面賣什麼.
就是賣無聊但我喜歡.
我覺得這件事對我來說.
有一定程度很重要的.
一些屬靈的提醒.
我覺得一間教會.
教會這件事情就是這樣.
做教會也好.
玩彈珠也好.
甚至上妝也好.
有一個很類似的地方.
教會是什麼呢.
教會就是一群人.
大家聚在一起.
大家去做一些事情.
做完之後可能對這個世界.
沒有什麼特別的結果.
但是OK.
這個對我來說是一個很重要的屬靈的反省.
那次台灣退休我們有參觀The Hope.
可能大家聽過一間台灣的教會.
基本上Full Church和The Hope都類似.
都是一些比較活躍的教會.
但是我發現.
我們都是有些不同的地方.
The Hope是一間很大的教會.
但其實他們又未必和我們這個地方的間流完全一樣.
不需要比較.
因為大家正正都是一群人聚在一起去做一些事情.
如果我們純粹是為了一個很商業的.
或者是一些像外面公司那樣去想.
那件事有沒有效果.
這樣做得不行就取消.
這個不是教會要去做的事情.
如果我們任何一間教會.
只是去單單計較那件事有沒有效.
純粹用效益去做的話.
純粹去想有沒有用.
有沒有後果去想的話.

$^{401}$這就不是一間教會.
這只是一個宗教的生意.
不是說結果不重要.
不是說我們不去想.
而是教會不是單單用結果去衡量它本身的存在.
牧養或者我們要做的使命.
耶穌基督的命令.
這些全部都不是和我們的結果有關係.
沒有結果都要去做.
這件事不是要去衡量的.
我們可以回到最基本.
不如看網上宋算了.
這基本上是最高性價比的.
拍個影片播個歌.
然後看完就行.
很便宜.
但這不是教會.
我們是一群人聚在一起.
一起去做一些事情.
這才是我們最重要的地方.
回到經文裡面.
我們發現故事並沒有停在那裡.
當大衛將三個勇士帶來的水.
墊在地上的時候.
大衛為了這個舉動.
說了一句非常非常重要的一句話.
大衛強調他倒了的水.
不是純粹將水倒在地上.
而是將水墊在耶和華的面前.
並且倒過來說.
耶和華這三個人冒死去打水.
水是他們的血一半.
原來一群人一起去做一些事情.
不一定是講求那個結果.
大衛將整件事的本質去改變.
大衛重新定義他們三個人所做的事的本質.
他們做了一件信仰的事情.
他們不是純粹做生意.
這三個人冒死去在百里行那裡.
付上諸多的心血.

$^{441}$這是在耶和華面前所囂顫的一件事.
這場生命的敬拜.
我們不能不承認.
信仰從來都是一件沒什麼用的事情.
大家都是在外面工作.
教會作為一個NGO.
或者一個很邊緣很微小的NGO.
其實對於社會沒有什麼特別大的改變.
或者經濟貢獻.
我送外賣反而可以改變世界.
送一碗麵給一個人可以吃到.
但是一群教會大家聚在一起.
有什麼特別的貢獻呢.
其實沒有什麼特別的幫助.
但是又如何呢.
一群基督徒.
我們一群人能夠聚在一起.
大家一起做一些事情的時候.
可能沒有什麼特別的結果.
沒有什麼特別的改變.
但是OK.
因為這是一件信仰的事情.
這是在耶和華面前所擺上的一件事情.
羅馬書十二章.
大家很熟悉的經文.
正正就是講這件事.
所以你們要將身體獻上當作活祭.
最後說什麼.
是一個理所當然的敬拜.
原文所謂的事奉.
其實就是敬拜這個字.
羅馬十二章所說的.
你們將身體獻上當作活祭.
這是一個理所當然的敬拜.
我們基督徒的生命.
大家聚在一起.
我們的生活就是一場敬拜.
這個敬拜是理所當然的.
什麼叫敬拜.
敬拜沒有說贏或輸.

$^{481}$敬拜沒有說有效或沒有效.
敬拜不是看後果.
敬拜就是一場在上帝面前呈現的敬拜.
無論目標是什麼.
結果什麼.
成效什麼都好.
它就是一場在上帝面前的敬拜.
所以當我們一群弟兄姊妹.
一起做一些事情的時候.
無論是沖咖啡.
或者是預備好這個地方.
或者是任何事情.
這場本身就是一場在上帝面前的敬拜.
你不需要去看結果是怎麼樣.
結果就是在神面前.
在地面上的敬拜.
所以弟兄姊妹我們一起去想一想.
我們今天在做什麼.
坦白說不單單是這裡.
或者是感受.
或者是什麼.
你都知道.
我自己在今年年初.
在建道裡面.
舉辦了一個香港教會前景會議.
我負責整個會議的籌備.
基本上這兩三年裡.
我不斷有很多跟中派教會領袖的訪問.
很多不同的交談.
基本上這幾年香港教會.
是一個下降軌的裡面.
政治運動.
疫情.
移民潮.
基本上每一間教會.
大大話話有些是跌了兩成三成都有.
弟兄姊妹離開.
牧者又離開.
很多很多時候非常艱難.
在這個大艱難裡面的時候.

$^{521}$我們問我們做什麼好.
基本上是沒有一個解決方案.
沒有一個能夠扭轉局勢的方法.
我想不是單單人數問題.
而是教會裡面的士氣問題.
基本上是跌到覺得回來做什麼.
或者不知道可以做什麼好.
但我發現一個最基本能夠做的事情就是.
你不要管.
基本上我們一班的教會.
仍然留在香港的人.
我們一起去做一些事情.
你不要問那個結果是怎麼樣.
只要大家仍然肯去做的時候.
這個就是一個好的開始.
搞一場福音報道會也好.
開一個新的咖啡師公也好.
搞什麼小組也好.
任何的事情只要你仍然肯去做.
做一些教會應該做的事情.
做一些教會要做的事情的時候.
不用問那個結果是怎麼樣.
你肯去做的話.
這個就是一個好的開始.
對於我們佳欣樓頂師妹來說.
我覺得這個對我們特別重要.
我記得佳欣樓幾年前開始的時候.
看著這個地方.
跟著大家一起每個星期聚會.
這班人.
雖然有很多不同的改變.
但相識面很容易看到.
每個星期見到的人.
大家都是回來佳欣樓裡面.
其實這個就是一個群體.
你也知道我是錦繡堂信主.
我中六的時候是錦繡堂信主.
基本上作為一個傳道人.
基本上我離開錦繡堂.
跟著去孫德那裡做方剛的事.

$^{561}$然後就到建道.
然後去德國.
從錦繡堂到孫德堂.
到柏林宣道會.
到回到香港開流堂教會.
基本上我看到很多不同教會群體的成長.
錦繡堂仍然是我的母會.
但我知道錦繡堂不是我每個星期都會回來的地方.
但我知道教會就是這樣.
群體的重要性就是大家一班人.
一起去做一些事.
一個熟悉的面孔.
這裡其實很像早期教會.
早期教會其實就是這樣.
一班人在一個地方聚在一起.
大家去為主去做一些事.
不要問結果是怎樣.
仍然的來做.
2000年前沒有那麼多信仰品牌.
沒有恩福還是感受.
還是全聖教的東西.
教會就是教會.
一班人在一個地方聚在一起去做一些事.
這個就是上帝眼中的群體.
大家可以去.
重要的不是問我們.
更加重要的問題是.
究竟我們這班人.
每個星期在一個地方聚在一起.
做些什麼呢.
我們是不是真的為主去做一些事呢.
這個才是我們真正思考問題的重點.
教會不是純粹一個建築物.
不是純粹一些資源.
不是純粹一個名字.
教會是一班人聚在一起去為主做一些事.
如果這班人願意一起去努力去為主做一些事的話.
據我這十幾二十年經驗.
開了那麼多間教會的經驗.
這些問題其實是次要問題.

$^{601}$仍然是可以解決的.
仍然是次要的.
重點是這班人是不是真的為主做一些事呢.
如果不是的話.
那麼多資源其實是沒有什麼意思的.
如果這班人只是圍圍圍地聚在一起.
沒有什麼特別的使命.
沒有什麼特別的一起去做一些事的話.
這個未必一定是我們最想見到的教會情況.
所以很值得我們一起去問一問.
究竟這個群體是怎樣.
我們大家在這個地方聚在一起.
願不願意一起去為主做一些事.
這個是我們最重要要思考的問題.
我一直在祈禱.
祝你求你幫助我們.
因為你在地上設立群體.
設立教會.
正正是要我們一起聚集.
並且為你去做一些事.
我們不問結果是多少.
我們有多少資源.
有多少的後果.
我們單單要知道你是否當中呼召我們.
一起去聚集並且為你做事情.
求主讓我們能夠去問清自己.
我們怎樣來思考我們在這個地方.
怎樣為你去作供.
讓這個教會真真正正成為你所喜悅的教會.
奉基督耶穌名字祈求.
\newpage



\section{提摩太後書 4:1-8}
\label{sec:TXT2__xkqM8}
\textbf{2023-09-10 早堂《保羅最後的勸勉》提摩太後書 4\_1-8 講員:鄧泰康傳道}
\newline
\newline
連結: \href{https://youtube.com/watch?v=TXT2_-xkqM8}{\texttt{https://youtube.com/watch?v=TXT2\_-xkqM8}} ~~~~ 語音日期: 2023-09-12
\newline
\newline
\hyperref[sec:naPXdHer1J8]{\small{< < < PREV SERMON < < <}}
~
\hyperref[sec:index_chronic]{\small{[返順時目]}}
~
\hyperref[sec:index_scriptual]{\small{[返順卷目]}}
~
\hyperref[sec:LQDQ7CliZ7U]{\small{> > > NEXT SERMON > > >}}
\newline
\newline
提摩太後書 4:1-8
\newline
\begin{longtable}{cl}
\hline
\hline
章節 & 經文 (和合本修訂版)\\
\hline
4:1 & \begin{tabularx}{0.7\textwidth}{X} 我在神面前，並在將來審判活人死人的基督耶穌面前，憑著他的顯現和他的國度鄭重地勸戒你： \end{tabularx} \\ \\ \relax
4:2 & \begin{tabularx}{0.7\textwidth}{X} 務要傳道；無論得時不得時總要專心，並以百般的忍耐和各樣的教導責備人，警戒人，勸勉人。 \end{tabularx} \\ \\ \relax
4:3 & \begin{tabularx}{0.7\textwidth}{X} 因為時候將到，那時人會厭煩健全的教導，耳朵發癢，就隨心所欲地增添好些教師， \end{tabularx} \\ \\ \relax
4:4 & \begin{tabularx}{0.7\textwidth}{X} 並且掩耳不聽真理，偏向無稽的傳說。 \end{tabularx} \\ \\ \relax
4:5 & \begin{tabularx}{0.7\textwidth}{X} 至於你，凡事要謹慎，忍受苦難，做傳福音的工作，盡你的職分。 \end{tabularx} \\ \\ \relax
4:6 & \begin{tabularx}{0.7\textwidth}{X} 至於我，我已經被澆獻，離世的時候到了。 \end{tabularx} \\ \\ \relax
4:7 & \begin{tabularx}{0.7\textwidth}{X} 那美好的仗我已經打過了，當跑的路我已經跑盡了，該信的道我已經守住了。 \end{tabularx} \\ \\ \relax
4:8 & \begin{tabularx}{0.7\textwidth}{X} 從此以後，有公義的冠冕為我存留，就是按著公義審判的主到了那日要賜給我的；不但賜給我，也賜給凡愛慕他顯現的人。 \end{tabularx} \\ \\
[1ex]
\hline
\hline
\end{longtable}
$^{1}$好 弟兄姊妹早安.
我看看今次可不可以破這個魔咒.
每次講到我都咳嗽.
但每次我在講到又不咳嗽.
這是要見證著上帝的恩典.
見證奇蹟的時候.
我今天和大家分享的這段經文.
在這段日子裡思考了很多 想了很多.
我想如果有部份的弟兄姊妹.
或者大家上星期的報告知道.
我接下來要做一個手術.
最近大家都看到我的腳走得一般般.
最後醫生發覺在頸椎的地方.
有骨刺壓著中樞神經.
所以接下來會做一個有風險的手術.
醫生說因為骨刺只離開中樞神經兩毫米.
所以在過程裡其實有很多思考.
但當然我想強調上帝有很多恩典.
這次幫我做的醫生是非常有經驗的.
在做脊椎這方面其實很有權威.
在過程裡其實我就會不斷思考.
當然醫生很直接.
你知道醫生一定會跟你講所有風險.
他說你有機會在病床上做完.
不小心碰到中樞神經.
你會癱瘓.
你也有機會就這樣走了.
但我一直都說不定.
當我一直思考的時候.
我只會想如果上帝在這次手術裡將我帶走.
那我會怎樣呢.
我去想的時候.
一個傳道人其實可以唯一留下的是什麼呢.
後來我一直思考的時候.
我就覺得就好像保羅那樣說.
原來保羅儘管他在我們心裡有多偉大.
但他唯一可以留下的就是他生命的見證.
和對弟兄姊妹的祝福.
因為有一天的時候.
當上帝帶他走的時候.

$^{41}$他遺留下來的時候我們很感恩.
上帝選照他.
所以我們今天才會有這麼多的新約書信在裡面.
他繼續教導我們.
這個是上帝的恩典.
當然我今天講這篇道的時候.
我想強調我連保羅的十分之一都沒有.
無論是他的毀身.
無論是他的熱心.
無論是他對上帝的付出.
他那種的熱情.
坦白講我自己真的十分之一都沒有.
我沒有他的能力.
我沒有他這種毀身的生命.
但我有他的心情.
我今天去講的時候.
突然間有一段經文在我腦裡浮現.
在亞國書那裡他曾經說.
其實明天如何你們還不知道.
你們的生命是什麼呢.
你們的完事一片雲霧.
出現笑時就不見了.
你們只當說主若願意.
我們就可以活著.
也可以做這事或做那事.
在這段經文裡面最觸動我的就是什麼呢.
主若願意.
頂尖的原來我們今天每一個人.
我們繼續去努力去過生活的時候.
最關鍵的條件是主若願意.
其實我們人真的很差.
剛剛星期五大家都知道.
現在這裡變成了水塘.
大家知道嗎.
現在大家很安坐的地方變成了水塘.
一場雨而已.
我們已經倒水下鞋.
我也覺得很好笑.
身體裡面我也聽過很多人說你這個骨刺.
一支骨刺而已.

$^{81}$一個小到都不知道怎樣形容的骨刺.
但就對人生裡面產生了一個.
很大很大的影響.
昨天剛剛奸流也播放了那套電影.
愛潘B.
如果大家有機會看的時候就會知道.
無論你有多少計劃.
無論你有多少的鴻圖大志.
但是一粒病毒而已.
就將一個家庭裡面全面的反轉.
所以在這個過程裡面.
我一邊去思考的時候.
我發覺無論怎樣也好.
我將我自己的生命交給神.
有些頂尖問我.
你面對著這個疾病.
這次的手術你心情如何.
我可以很坦然很真實地告訴頂尖們.
我真的不擔心.
因為我只是在過程裡面.
我跟上帝去爭辯.
每逢去到一個關鍵.
我就會跟上帝很真誠地去爭辯.
我說我的性格我很清楚.
我不喜歡去連累人.
所以我說上帝.
很簡單的.
要不你就在病床上帶我走.
要不你就令我好起來.
令我可以繼續去事奉你.
如果你認為我事奉你還不夠.
讓我有機會繼續去改進的話.
你繼續將生命賜給我.
這是我跟上帝的禱告.
所以我今天帶著一個什麼心情.
跟大家說呢.
我不知道大家有沒有看過一套電影.
那套電影叫《非誠勿擾二》.
這套電影裡面.
其中有一個是我今天很深刻的印象.

$^{121}$裡面其中一個角色去還了一個絕症.
這個絕症是一定會走的.
但是他的選擇是什麼呢.
在他還沒有走的時候.
之後他就開了一場道別會.
就是告別會.
那個主角可以坐在那裡.
自己去聽每一個人對他的說話.
懷念.
他也可以很真心地跟他們說一聲.
謝謝和告別.
我今天帶著這樣的心情.
去跟大家說這篇道.
因為我經常都會覺得.
仍然有機會去珍惜每一個的一刻.
是上帝給我們的恩典.
所以每一個相聚每一個相處.
在什麼都還沒說之前.
我鼓勵弟兄姊妹.
好好地珍惜每一次我們的相聚.
每一次的相處.
因為下一刻.
究竟你身邊的這個人還在不在.
我們不知道.
在這段經文裡面.
其實你發覺保羅.
他最厲害的地方是什麼呢.
他說我在神面前.
並在將來審判沒人死人的基督面前.
憑著他的顯現和他的國度去祝福你.
保羅去勸勉他們的前提是什麼呢.
是這個上帝是一直都得勝.
就是說他今天告訴我們.
我們所信的這個上帝.
我們今天在地上很多很多的一場努力.
並不是在等候一個結果.
而是一切都按照著他的旨意去進行.
所以你就會發覺我們的教導裡面.
我們所有的門訓課程都告訴大家.
我們要很努力去尋求神的旨意.

$^{161}$因為今天其實我們每一個走上上帝的路的時候.
首先我們要知道本身就是得勝.
你不是去等待結果.
而是你在打一個得勝的仗.
你是要打仗.
但我們在打一個已經贏了的仗.
如果大家有機會看聖經.
你要知道啟示錄已經告訴你.
這個世界最後的末後是怎樣的.
其實我們是贏的.
我們今天面對任何的事情.
是帶著一種贏的心態去走.
儘管我今天帶著這個病.
我不知道結果.
因為掌權在上帝那裡.
但儘管我假設我走.
不是代表輸.
大家教清楚.
聖經有一句很寶貴的話.
就是什麼.
一粒墨紙若不落在地裡死了.
仍舊是一粒.
但落在地裡就結出許多紙粒.
我會覺得上帝無論讓這件事怎樣進展下去.
只要是他的心意.
這件事就會成為祝福.
所以是贏的.
今天我們每一個人.
我們要帶著上帝的能力.
我們要跟上祂的步伐.
這才是我們真正需要去做的事情.
因為上帝已經在贏.
但你一直去墮後.
很多時候我們說我們信這個信仰.
很多次在台上我都不斷去鼓勵弟兄姊妹.
我們這個信仰是很有能力的.
是可以改變生命的.
我曾經說過很多個見證.
我之前的教會是做吸毒者的.
我們是看著他的生命改變的.

$^{201}$上星期我參加了一個按目禮.
這個童工是由以前的吸毒者.
他後來信了耶穌.
到後來他回到戶外機構去事奉.
之後很多年他讀完神學.
成為傳道人.
到上星期去按目.
他在台上唯一可以分享的是什麼.
他都說 全是恩典.
弟兄姊妹我們這個信仰.
是可以衝破你人生裡面所面對的一切疾病.
所有你所面對的困難.
所有面對的關係.
是可以的.
其實上帝是得勝的.
但是你跟得上嗎.
你跟到上帝的步伐嗎.
這個才是關鍵.
所以保羅一開始跟提摩太太說的時候.
我想大家都知道.
在聖經裡面所形容的提摩太太.
是一個相對軟弱的童工.
但是其實他就告訴她.
你不要看著你自己的軟弱.
因為上帝是得勝的.
你所信的上帝是得勝的.
我們今天為什麼說我們是基督的精兵.
是因為我們的元帥是贏的.
所以你才是精兵.
所以弟兄姊妹.
當我們今天信了耶穌之後.
我們不是好像很可憐.
即是沒有什麼能力.
好像很軟弱.
不是的 弟兄姊妹.
保羅說什麼.
保羅說在他的軟弱上.
看出基督的完全.
第二件事.
務要全道無論得時不得時.

$^{241}$總要專心.
並用百般的忍耐各樣的教訓.
責備人 警戒人 勸勉人.
因為時候要到.
人必厭煩純正的道理.
以多發揚.
就隨從自己的情慾.
增添好色師父.
並且掩耳不聽真道.
偏向荒謬的言語.
你卻要凡事謹慎.
忍受苦難.
做全道的功夫.
盡你的積分.
第二件事你會發覺.
保羅提醒在他的國度裡面.
我們唯一要做的只有兩件事.
就是全道和五目仰.
無論得時不得時.
無論順或逆.
上帝的工作從來都不應該停下來.
我可以告訴大家聽.
現在是不是教會.
我講的不只是感受堂.
是不是所有教會最差的時間.
今天很多人都說.
因為我們經歷過好.
現在我們相對.
由退休潮 又移民潮.
又奉獻降低.
之後全部出來的數字都不好.
如果大家知道教新每五年的時候.
全部數字都是跌的.
現在很差.
靈姐妹.
知不知道教會增長最迅速的時候是何時.
如果你有機會看回歷史.
就是教會最困難的時候.
問題是究竟你做與否.
保羅在這裡很清楚地說.

$^{281}$你不要看著順或逆.
大家記不記得.
如果有我這麼老.
我經常都說我喜歡扮老.
二十多年前我在這裡做神學生.
那時候我們還是只有幼稚園.
我們沒有禮堂的.
只有幼稚園.
和我這麼老的那些.
還記得有一場大火嗎.
慘不慘啊.
幼稚園啊 著火啊 是不是啊.
學校怎麼搞啊.
教會崇拜怎麼搞啊.
弟兄姊妹.
最美的圖畫就出現了.
我好記得弟兄姊妹齊心努力.
當時沒有那麼多童工的.
齊心努力回到幼稚園那裡.
一起還原.
一起做清潔.
一起想方法.
怎樣令到那裡盡速去完成.
早幾天這裡變成水塘.
慘不慘啊.
很多時候.
可能大家還在想.
如果不是這樣.
就不會這樣.
人很喜歡的 是不是啊.
不是啊.
我那天看到一幅很美的圖畫.
我真的覺得那天是在家裡.
弟兄姊妹穿什麼衣服回來的.
T恤 短褲 踢拖鞋.
大家在這裡做什麼啊.
拔水.
大家抹地 拖地.
幫忙還原這個地方.
弟兄姊妹 這裡就是家.

$^{321}$當家有事的時候.
我們只能埋怨嗎.
不是.
有沒有想過.
為什麼上帝讓這裡成為水塘.
平時弟兄姊妹.
在什麼情況之下才回來這裡.
可能只剩下聚會.
但是那天在聚會以外的時間.
上帝給了大家機會一起回來.
一起下手.
我記得我在耕林的時候.
耕林裡也有一班很愛教會的弟兄姊妹.
有一次我記得好像是打完風還是什麼.
我忘記了.
總之回去我們要做清潔.
然後有一個弟兄很好.
他跑來跟我說.
說下次不要只我們做.
叫弟兄姊妹回來做.
聽起來好像是埋怨.
不是的 不是的.
弟兄姊妹告訴我.
下次叫弟兄姊妹回來做.
因為他們這樣才會有歸屬感.
是真的 弟兄姊妹.
我還記得幼稚園的時候.
我們很有歸屬感.
為什麼.
因為招行晚策.
每個星期早上都要弟兄姊妹一起回去幫忙.
設椅子 設音響.
是大家下手下腳去做的.
弟兄姊妹保羅提醒.
無論順或逆.
我們就要去做.
原來我們今天在上帝的帶領當中.
順或逆.
都是在他的計劃掌管當中.
所以剛才我才說.

$^{361}$不要看著你身邊的問題.
你看著問題 問題只會越多.
你看著的是因電.
為什麼我今次面對的疾病我這麼坦然.
我講一些見證給大家聽.
我最初知道自己是頸椎問題的時候.
說真的 我心裡都在想怎麼搞呢.
有很多人跟我說.
頸椎最好不要動就不要動.
特別現在我知道這個情況更加.
就不要亂動.
我也試過針灸很多其他的東西.
之後就去到不行了.
情況真的變差了.
無論如何 我就先去了大陸.
在廣大深圳醫院去做檢查.
我最初拍的是腦的問題.
然後看腦外科.
醫生就幫我做MRI.
我給大家一個參考.
在國內做MRI便宜很多.
之後他就幫我照了腦.
照了頸椎的MRI.
看完之後他說你的腦沒事.
但你的頸椎有骨刺.
當時他已經叫我做手術.
但我對於醫療的很謹慎.
然後我就沒有答應.
後來我回來.
我就想等到香港的醫生幫我.
因為我正在排期.
我就拿了MRI給政府醫生幫我看.
看完政府醫生後.
政府醫生看完後.
都說做手術了.
但坦白說我是擔心的.
我要很坦誠.
因為在現在香港的醫院裡.
因為太多人流失了.
我形容是.

$^{401}$太多新的人.
我不知道誰幫我.
我不知道會發生什麼事.
後來我認識了一位弟兄.
他以前是拿打素骨科的主管.
他移民了.
後來我發給他.
問他可不可以幫我看MRI片.
他說不行.
因為我不可以只看片.
我要看臨床.
所以他就介紹了一個醫生給我.
是一個非常好的朋友.
合作了很多年的朋友.
然後我就去看.
看完後.
醫生很直接了當地告訴我.
幫我檢查完.
亦都很清楚解釋了我為什麼會有這樣的情況.
因為以前醫生一直都解釋不到.
為什麼我左右的情況會這樣.
他解釋給我聽.
他說我現在窄轉的位置.
左邊影響運動神經.
右邊影響感覺神經.
所以剛剛在一個交叉位.
然後他說你這個情況.
物理治療沒有用.
吃藥沒有用.
你唯一就是做手術.
然後他跟我說.
如果不是必要.
他說我是做脊椎的.
不是必要.
我不會叫你做手術.
然後時間.
我接受了.
後來我跟回澳洲那邊的弟兄商量.
他說其實分別就是.
你做私家的話.

$^{441}$其實你是用金錢去換康復和安全.
因為這位醫生很有經驗.
但我想告訴你.
我就發覺恩典原來是這麼豐富.
他收了我的手術費是多少呢.
十萬元.
我問過其他人.
做這個手術.
如果你找一個有經驗的醫生.
你沒有二十多三十萬.
你不用想.
這個已經是恩典.
第二件事.
他這次.
平時他會由他再加一個.
普通的醫生去幫他.
他這次是由他.
再加一個他一個好朋友.
都是脊椎的轉科醫生.
跟我一起做.
但他沒有收我多錢.
他很直接告訴我.
因為多一個醫生.
將你的風險再減低.
另外一個在病的裡面.
我跟每個電影節目都說.
我最感恩的地方.
就是我從來都不是痛.
現在令到我的腳不舒服.
令到我行路不方便.
但我從來都不是痛.
我知道痛是很辛苦.
但上帝很保守.
弟兄姊妹.
其實在信義當中的時候.
上帝的恩典從來都沒有離開過.
上帝只要的是.
你去做.
無論得時不得時去做.
所以弟兄姊妹.

$^{481}$我們不要停下來.
我們要做神的工作.
因為當你看到上帝的恩典的時候.
你就可以繼續向前.
而不是去看著這個困難.
停下來.
我記得之前疫情的時候.
大家記得這裡全部聚會.
因為政府說一定要停下來.
其實我知道很多香港的教會是停下來的.
是失去措手不及的.
有些教會連直播都未搞定.
都未知道應該怎樣做.
但我們很感恩.
我們當時.
其實即時開了很多在網上的講座.
網上的課程.
我們繼續去做牧羊.
繼續去做傳福音.
在那段日子裡我們從來沒有停過.
而上帝亦都很感恩.
在疫情之前.
祂就已經給我們設置好這些直播系統.
我們可以繼續透過直播.
去和弟兄姊妹繼續崇拜.
弟兄姊妹所以.
我們要繼續向前.
而接著祂提的是什麼.
你要專心.
你不要理會外在的環境.
不是人聽不聽.
不是難不難.
問題是你傳不傳.
所以你會發覺我在荃灣宣道會出身.
如果大家知道最初的牧者裡有幾位.
藤斤溫牧師.
我們荃灣宣道會是都乾初牧師.
都乾初牧師是外省人.
撈鬆你.
不會說廣東話.

$^{521}$但是他就這樣在香港開始傳道.
拿著風琴.
手風琴.
走出去的時候.
所以我小時候跟著他們報道的時候.
我不知道大家有沒有印象.
還會揭著紙那些.
揭著紙一頁一頁反過去那些.
我們去做街頭報道.
我們的前輩在最困難的時候去做工作.
領姐妹.
所以我很鼓勵你們.
不要因為環境而停下來.
不要看著自身的困難.
上帝說過.
你先求我的國和我的義.
我將你所需用的一切.
我要加給你.
我們越是淪惜.
我們越是無.
所以保羅耶德提.
用盡一切的方法去教導.
警告.
勸勉.
弟兄姊妹.
我們很需要讓弟兄姊妹得到牧養.
真正的牧養不是只是.
我想強調.
不是不做.
不是只是和大家喝下茶.
探下大家.
真正的牧養裡面.
我們真是需要將大家帶回上帝的面前.
所以求神幫助我們.
好嗎.
我相信我們作為傳達人的時候.
我們不斷要調整我們的牧養.
我們不斷要去調教.
亦都鼓勵弟兄姊妹.
我在家人那邊.

$^{561}$我就會直接和他們說.
你們想到在教會裡面.
有什麼不足.
你們覺得不夠.
你們要出聲.
不要觀望.
不要吃花生.
告訴我們.
而更加漂亮的童話是什麼.
是你們自己想到的一些辦法.
然後和傳達人分享.
很多的事工.
是可以百花齊放.
是可以一起去做.
因為你們每一個人都有上帝對你們的感動.
和恩賜.
而最後的時候.
保祿是什麼呢.
保祿用了一段說話.
就是大家經常都聽到.
他說我現在被囂點.
我離世的時候到了.
那美好的將我已經打過了.
當跑過路.
我已經跑盡了.
所信的道我已經守住.
從此以後有公義的冠冕為我存留.
就算按照公義審判主.
到了那一天.
要賜給我的.
不但賜給我.
也要賜給那凡愛慕我顯現的人.
保祿到最後的時候.
其實我們見到.
是用他的見證.
用他信心的見證.
去成為他們的鼓勵.
這個給我自己很大的提醒.
說上帝的話.
說哪裡很容易.

$^{601}$說真的.
我記得上次我講到都講.
現在很流行查GDP.
整個AI人工智能的傳達人用什麼.
你知道現在科技.
打個靚仔的我.
其實都還有.
整個明星樣子的.
跟大家講到.
幾爽.
但問題分別的就是.
究竟我們今天做傳達人.
或者弟兄姊妹.
我們今天作為信徒.
我們有沒有一個見證.
讓大家見到.
保祿他夠膽講.
他說美好將我帶過.
當跑過路.
我跑盡了.
所信的道我已經守住.
其實我在這裡我只是配.
我配和大家去分享的是.
互相勉勵.
剛才我都提過.
我連保祿的十佰特都沒有.
我的障礙未能完的.
我的路未跑盡的.
我的道守得好不好.
可能還有缺陷.
所以我一直求上帝.
上帝你給我機會.
如果可以的話.
讓我繼續.
繼續去改善.
但保祿在這裡的時候.
他提醒在他離開之前.
每一個人在他離開之前.
要做的三樣東西.
要打美好的仗.

$^{641}$弟兄姊妹.
要跑盡你能夠跑的路.
要守住你的道.
這個是我們信徒必須要做.
如果這個信仰是我們最高的標準.
弟兄姊妹.
請你放在你整個生命當中.
整個家庭裡面.
讓到我們一代一代的時候.
都能夠繼續打美好的仗.
怎樣去打.
怎樣去跑.
怎樣去守.
最簡單.
當我們在最困難的日子裡面.
你身邊有很多受苦的人.
弟兄姊妹.
你有沒有為他.
代表著上帝.
去向他付出.
有沒有他成為他的祝福.
有沒有他成為他的安慰.
有沒有將他帶到神的面前.
弟兄姊妹.
讓我們一起的彼此勉勵.
另外.
他有對未來的盼望.
因為他說有官面為他存留.
弟兄姊妹.
我們每一個人都渴望上帝永恆的掌上.
等待著我們.
很坦白說.
想到這一點的時候.
其實我心裡面是心驚膽跳的.
我都想過.
如果我回到神的面前.
我會怎樣呢.
我很坦白說.
我從心裡面.
我只是求神憐憫.

$^{681}$我真的覺得自己未合格.
我不是說我不行.
就是不要教錯.
但我覺得我未合格.
我去到這一刻.
我仍然會覺得.
我不夠膽說.
我去到神的面前.
啊 上帝.
你頒個獎給我.
我不夠膽說.
因為原來還有很多的事.
我們可以去做.
所以我們需要去仰望上帝.
盼望著這一份獎賞.
去繼續走我們人生的路.
因為保羅說.
這些獎賞是賜給所有愛慕他顯現的人.
弟兄姊妹.
你渴慕他的顯現.
你期望他回來.
上星期我在教會跟他們說.
很真實的 弟兄姊妹.
如果你連地上都不享受教會生活.
你都不享受神的話語.
你都不享受他的詩歌.
你上到去.
你很大件事.
舉個例子.
有些弟兄姊妹說.
不是啊 我不是賭錢.
我喜歡打四方城而已.
打下衛生麻將而已.
很享受這個過程.
弟兄姊妹.
你上到天家沒有麻將機.
不會有人燒給你的.
你搞清楚.
我們今天在地上.
究竟我享受什麼.

$^{721}$我享受什麼.
我的滿足是什麼.
所以上帝的獎賞.
只是賜給那些凡愛慕他顯現的人.
弟兄姊妹.
這是一個問號.
究竟我們今天顯現期望什麼.
所以求上帝幫助我們好嗎.
我不知道上帝在未來.
這兩個月會怎樣帶我.
我不知道.
但是我只是作為一個傳道人.
我真的很想告訴大家.
儘管我不在了.
我也希望弟兄姊妹.
繼續走神的路.
繼續去打這場仗.
繼續去守住這個道.
可以嗎.
盼望這是一個傳道人.
對於他一群無恙的羊仔.
一個最後的期盼.
希望敢守堂的弟兄姊妹.
能夠活出上帝的生命.
我會在此說.
我多麼相信上帝會帶我回來.
所以我們繼續一起努力.
求神一起幫助.
讓我們每一個人都不要枉過.
因為時間有限.
同心祈禱.
秋天苦我們求主幫助.
帶領我們.
因為你總是將最好的來賜給我們.
主恩但是我們所面對的疾病.
或者其他都告訴我們.
其實我們每一個人的生命.
我們的時間都有限.
求你讓每一次.
當有這些生命出來的時候.

$^{761}$我們都去明白.
我們都能夠提醒我們.
去知道原來我們每一個人.
都要善用我們的時間.
就正如孩子在每一次安息禮拜.
我都說每一個離開的弟兄姊妹.
其實都去提醒我們參加的人.
原來有一天我們都要走.
所以我求你幫助.
我們多謝你透過保羅他的提醒,勸勉.
去告訴我們.
我們還有很多很多你的工作要去做.
我們還有很多很多可以去為你而去活的機會.
所以我求你幫助我們.
願意我們帶著信心.
帶著恩典我們走這條路.
我們多謝你.
祈禱奉靠主耶穌.
得成名字祈求.
好 再見 拜拜.
\newpage



\section{詩篇 8:1-9}
\label{sec:LQDQ7CliZ7U}
\textbf{2023-09-09 晚堂《神你是誰?而我又算甚麼?》詩篇 8\_1-9 講員:張美薇牧師}
\newline
\newline
連結: \href{https://youtube.com/watch?v=LQDQ7CliZ7U}{\texttt{https://youtube.com/watch?v=LQDQ7CliZ7U}} ~~~~ 語音日期: 2023-09-13
\newline
\newline
\hyperref[sec:TXT2__xkqM8]{\small{< < < PREV SERMON < < <}}
~
\hyperref[sec:index_chronic]{\small{[返順時目]}}
~
\hyperref[sec:index_scriptual]{\small{[返順卷目]}}
~
\hyperref[sec:eRch8741F34]{\small{> > > NEXT SERMON > > >}}
\newline
\newline
詩篇 8:1-9
\newline
\begin{longtable}{cl}
\hline
\hline
章節 & 經文 (和合本修訂版)\\
\hline
8:1 & \begin{tabularx}{0.7\textwidth}{X} 耶和華－我們的主啊，你的名在全地何其美！你將你的榮耀彰顯於天。 \end{tabularx} \\ \\ \relax
8:2 & \begin{tabularx}{0.7\textwidth}{X} 你因敵人的緣故，從孩童和吃奶的口中建立了能力，使仇敵和報仇的閉口無言。 \end{tabularx} \\ \\ \relax
8:3 & \begin{tabularx}{0.7\textwidth}{X} 我觀看你手指所造的天，並你所陳設的月亮星宿。 \end{tabularx} \\ \\ \relax
8:4 & \begin{tabularx}{0.7\textwidth}{X} 人算甚麼，你竟顧念他！世人算甚麼，你竟眷顧他！ \end{tabularx} \\ \\ \relax
8:5 & \begin{tabularx}{0.7\textwidth}{X} 你使他比神微小一點，賜他榮耀尊貴為冠冕。 \end{tabularx} \\ \\ \relax
8:6 & \begin{tabularx}{0.7\textwidth}{X} 你派他管理你手所造的，使萬物， \end{tabularx} \\ \\ \relax
8:7 & \begin{tabularx}{0.7\textwidth}{X} 就是一切的牛羊、田野的牲畜、空中的鳥、海裡的魚，凡游在水裡的， \end{tabularx} \\ \\ \relax
8:8 & \begin{tabularx}{0.7\textwidth}{X} 都服在他的腳下。 \end{tabularx} \\ \\ \relax
8:9 & \begin{tabularx}{0.7\textwidth}{X} 耶和華－我們的主啊，你的名在全地何其美！ \end{tabularx} \\ \\
[1ex]
\hline
\hline
\end{longtable}
$^{1}$請姐妹平安.
很高興今天又有機會回來大家的當中.
跟大家分享神的話.
我們開始的時候.
我們一起低頭祈禱.
天父上帝求你使用孩子.
幫助孩子.
給孩子所說的.
是你要孩子說的.
更加幫助會中.
他們有能聽的耳朵.
有能明白你話語的頭腦.
更加有柔軟的心.
能夠將你的說話.
藏在他們的心裡面.
幫助大家.
有能行道的手和腳.
去行你的道.
聽我們的祈禱.
奉耶穌基督的名給.
阿們.
當我預備這篇講章的時候.
因為早一個星期已經開始寫.
我就寫了.
9月1號的大風.
誰知今天再講的時候.
9月1號已經是濕濕碎碎的.
今天就是昨天.
昨天是黑雨.
大家都知道影響很大.
我星期五早上才坐船坐車.
回來感受.
去看感受的情況.
於是我就WhatsApp一個姐妹.
問問她.
感受平安嗎.
她就給了一系列的影片給我.
看到原來你們這裡.
很厲害.
水流得很厲害.

$^{41}$然後今天.
因為我今天很忙.
早上6號到下午5點多.
我都在外面出出入入.
原來謝牧師今天下午3點多的時候.
就給了我WhatsApp.
張牧師你今天來港島嗎.
原來他們又給了這一堆的.
我想是一樣的.
因為我都沒空看.
他就說你一定要鼓勵一下弟兄姊妹.
因為是靠他們的努力.
我們今天才可以有一個這麼整齊.
這麼清潔的地方.
可以一起在這裡崇拜.
如果你們不知道什麼事的話.
那你們就去問問.
去看看.
究竟教會發生什麼事.
以致謝牧師又是這樣說.
另外那位姐妹都說.
很難得很難得很難得.
本來就想和大家講一下.
其實8月31號.
是一個很特別的日子.
我們說那是超級藍月亮.
所謂的藍月亮.
並不是說月光變了藍色.
而是說在同一個月裡面.
有第二次的滿月.
那就是叫藍月亮.
那什麼叫超級月亮呢.
我想大家都知道.
其實月亮圍繞著走的那條軌道.
其實是橢圓形的.
但有些時候滿月的時候.
就近一些的地球.
有些時候就遠一些.
所以當我們去看的時候.
我們就會看到.

$^{81}$大家都是拍月亮.
大家都是同樣的放大.
是不是看到.
左手邊那個和右手邊那個.
好像距離都遠的.
右手邊那個就是8月31號.
那晚的藍月亮.
那晚雖然是16號.
不是15號.
但比2月那次的.
不是藍月.
那次的月圓.
就是8月31號的月圓.
比2月那次的月圓.
那個距離.
看到那個差別都大.
所以.
還有我們發覺.
當有超級藍月亮的時候.
就會引起天文大潮.
所以沿海低窪的區域.
就會海水倒灌.
局部淹水.
所以為什麼我們會看到.
今次蘇拉來的時候.
政府就很大緊張.
很緊張很緊張要去應付.
因為那時候.
剛剛是8月31號.
是超級藍月亮.
那時候就知道.
那個潮漲潮積之類.
是很厲害的.
所以很嚴陣以待.
誰知道這次就沒什麼事.
大家都很慶幸很開心的時候.
誰知道第二個星期.
接著來那個黑雨.
另外一個颱風.
殘餘的橫流.

$^{121}$就影響到.
我們今年這麼大的難處.
這次其實已經不是滿月的時間.
但當我們去說到月光的時候.
我想大家都有讀過這首詩.
床前明月光而是地上霜.
舉頭望明月低頭思故鄉.
我記得我還沒讀幼稚園.
我媽媽就已經教我.
後來到我讀小學的時候.
都有讀這首詩.
其實詩當中描繪.
在午夜的時候月光滿地.
但雖然眼前的月光是這麼美好.
但叫人追難忘的是什麼.
就是我們的故鄉.
意思就是說.
這首詩是一個遊子寫的.
他寫自己的心情.
我們會發覺原來人的憂傷.
都會被月亮的影響.
我又和大家去看看.
原來我們這次31號16號是滿月.
今天是多少.
大家知不知道今天是多少.
今天是25號.
如果大家今晚看到有月光的話.
應該是看過這樣的月光.
其實月光是我們每晚看到的.
都不一樣的.
是隨著日子轉變的.
在短短的一兩個星期.
天上的月光看來就被人改變了很多.
而當我們看到月光的時候.
也都會影響人的作為.
有人說月黑風高是什麼.
殺人夜.
什麼意思.
可以做很多為非索大的事.
人看到滿月的時候會怎樣.

$^{161}$就會很開心很快樂.
所以人是被月光所影響的.
當我再給大家看這三個.
大家認不認得是什麼太空人.
可能大家都未必認得.
但其實他們當中在1969年7月21號的時間.
他們當中最左邊的那個.
他就是岩士堂.
中間就是柯林斯.
然後再旁邊的就是艾德尼.
岩士堂是什麼.
他是第一個登陸月球.
他的腳是第一個踏落月球那裡.
他的腳印.
他就會說.
一個人的一小步是什麼.
是人類的一大步.
雖然他們是這麼高的科學水平.
這麼高的先進.
但是他們三個太空人在月亮軌跡上的時候.
他們就讀到第三第四節.
我們剛才所讀的第三第四節.
就是說.
並你所陳設的月亮星宿.
人算什麼.
你竟顧念他.
世人算什麼.
你竟眷顧他.
一個好像很科學.
很前眾.
很前面的一個太空人.
他看到月球這麼美妙的一面.
他更加踏足在月球上.
他竟然會用這一句來讀出來.
我們看天文學家從前.
以前時時說月球是由地球而生的.
但是當人從月球回來的時候.
他帶了很多什麼.
帶了很多月球石回來.
所以對科學有很多的幫助.

$^{201}$經過很多化驗之後.
發覺地球是沒有月球石上的物質的.
意思是什麼.
意思是月球不是地球spin出來的一塊東西.
不是由地球而生的.
因為其實我們看詩篇第八篇都說.
他說什麼.
原來月球並非需要地球.
那是什麼.
神陳設在地球附近的一個星宿.
在這裡已經很清楚地說.
原來月球並不是地球來的.
我們又再看清楚一點.
詩篇第八篇.
然後他寫了一首詩篇.
和第十九篇是相映成趣的.
如果大家有時間可以回去看第十九篇.
第十九篇是寫太陽的.
第八篇是寫月光.
我們知道其實黑夜的詩歌.
很多時候都蘊藏著很深沉的悲戚.
或者百般的辛酸.
好像李伯剛才所看的靜夜詩.
他會寫到柔子詩香等等.
但當我們看詩篇第八篇的時候.
我們卻看到大衛這個詩篇.
是沒有悲哀憂戚的感覺.
一個月裡面.
月亮我們看到的.
真的是不一樣的.
但是大衛當他看到月亮的時候.
他可以覺得做人是很幸福,很喜樂的.
其實為什麼呢?.
為什麼大衛他寫的月光.
是這麼正面,這麼陽光.
我們本來都很難說陽光.
但其實他說的話就是這麼陽光.
因為他看到月光之下背後.
是有一個創造的主在.
所以他在這裡的時候.

$^{241}$他會和神有相交.
所以我們可以在這裡看到.
原來他從月光這裡.
可以看到一些什麼樣的東西.
所以讓我們可以學習.
從他的詩篇第八篇裡.
所以今天我說了三點.
第一點就是說.
神啊!你是誰呢?.
其實第一節就說.
耶和華我們的主啊!.
你的名在全地何其美.
你將你的榮耀彰顯於天.
而第九節也說.
耶和華我們的主啊!.
你的名在全地何其美.
大衛在月夜湧出頌讚.
所以我們會說這首歌是一首讚美詩.
詩人或者我們聖經學者會覺得.
它是一首讚美詩.
當他看到月亮的時候.
他就會看到原來神的榮耀.
是彰顯於天.
我們活在一個覆蓋在主的恩愛之下.
美麗的世界.
因為裡面說到神的名.
他說你的名在全地何其美.
你的名在全地何其美.
神的名是什麼呢?.
我們知道神會當祂向摩西透露.
祂說什麼呢?.
祂說我是.
I am who I am.
我是自有永有的.
神是至尊無比的.
祂是自我存在的.
祂是是因為祂自己的「是」.
祂不是其他的因素叫祂這樣的「是」.
祂有是因為祂自己的「有」.
不是因為其他的因素叫祂「有」.

$^{281}$祂是永不改變的神.
祂是 祂永是 祂在 祂永在.
祂有 祂永有 祂是永不改變.
祂是一切事物的源頭.
祂所做的世界是其大無比.
當中是有無限多的星體.
天上星體有多少?.
大家都數不盡.
有人說地上的沙.
算起來也可能沒有天上的星獸那麼多.
誰可以數呢?.
神將祂的榮耀彰顯於天.
不僅如此.
神更加是什麼呢?.
神是我們的主.
我們這麼偉大的主.
這麼偉大的神.
祂與我們人是有關係的.
祂是主宰.
祂是統治者.
祂是審判者.
祂是我們的救世主.
祂是我們的供應者.
祂是我們生活的維持者.
祂是我們的後盾.
神並不是創造了我們就不理我們.
神創造了我們.
繼續與我們在一起.
祂與我們有立約的關係.
因為神與我們有立約.
所以祂一定是以祂立約的愛.
來愛我們.
所以當我們再去看神是誰的時候.
我們就可以看到第二節.
第二節說.
「你因敵人的緣故從孩童和赫拉的口中建立了能力.
使仇敵和報仇的閉口無言.
並你所陳設的月亮星獸」.
這裡讓我們看到.
其實當大衛去看夜空的時候.

$^{321}$他第三節讓我們看到.
並你所陳設的月亮星獸」.
就好像太空人一樣.
他去看到那個天.
他立到腳落月球.
他仍然會讀這一節.
因為他看到神的大.
我們可以說.
可以有三種不同的人.
一種人很迷信.
他們很怕.
他們會凡物都拜.
凡物都信.
因為他們很怕.
另一種人是怎樣.
是無神論.
看到天的話.
怎樣.
和尚打散.
無法無天.
我是最大的.
我們知道一定不是.
但是大衛就是第三種.
他是一個敬畏者.
他思想的根基.
他的真理和知識.
他去了解到.
他有信心去接受.
神啟示的真理.
他看到神無限的崇高.
誰的手能夠觸及天上高的星.
誰有能力去左右天體的運行.
神不需要甚麼.
神也不需要人去做他的幫助.
神也不需要人去為他設立很多東西.
因為我們知道.
神是最大,最有能力.
他的智慧是無法測度.
他創造了宇宙.
他創造了窮倉.

$^{361}$他創造了深海大洋.
他創造了星宿.
也創造了這麼小的沙粒.
但我們可以看到.
從大衛在這裡來說.
原來他是有敵人的.
他有仇敵和報仇的.
因為他說:你因敵人的緣故.
從嬰孩和赫拉的口中建立了能力.
使仇敵和報仇的閉口無言.
神原來是有仇敵.
神是有敵人.
在靈界當中我們都知道.
我們知道有撒旦.
撒旦其實是一個墮落的天使.
在靈界當中也有戰爭.
但我們知道最後誰是得勝者.
最後神是得勝者.
但神如何去處理仇敵的張狂攻擊呢?.
仇敵在興風作浪.
挑撥離間等等.
神是靠什麼去和仇敵抵擋呢?.
這裡說:從孩童和赫拉的口中建立了能力.
孩童和赫拉是形容什麼呢?.
形容一些無助的孩童.
不單止是孩童.
還是赫拉的孩童.
是一個很無助的孩童.
但我們卻知道.
原來當耶穌在聖殿裡結證聖殿的時候.
孩童在那裡說:阿和珊,名是應當稱頌的.
有人就說:你要去斥責他們,不應該這樣.
耶穌也引用這一句話.
從孩童和赫拉的口中.
神是建立了能力.
我們更加.
大家都知道.
當我們去讀聖經或者教兒童主義學.
都知道大衛和哥利亞去抵擋的時候.
哥利亞是什麼呢?.

$^{401}$是一個大巨人.
大衛是什麼呢?.
大衛是一個小孩子.
對起哥利亞.
哥利亞有什麼呢?.
他有盔甲.
他有劍.
有棘.
大衛有什麼呢?.
大衛沒有.
大衛只有石頭.
大衛有他的丫叉.
但他卻可以用一個丫叉.
他的石頭.
去打中哥利亞.
哥利亞是跌倒.
最後大衛用哥利亞的刀去割了哥利亞的頭.
所以你會發覺.
原來大衛是很厲害的.
孩童或者赫拉都可以很厲害.
所以神不需要用和仇敵一樣強勁威猛的東西.
去抵擋仇敵.
神是用一些孩童和赫拉去抵擋仇敵.
所以當我們看到這裡的時候.
我們就會問.
神真的厲害嗎?.
厲害到無可比擬.
當我們看到神厲害的時候.
我們可能要去看看.
究竟人是算什麼呢?.
所以我今天第二點就是說.
我人又算什麼呢?.
其實我是誰.
是每一個人都在問的問題.
Eric Erickson.
他是發展心理學的專家和心理分析學者.
他是用心理在社會裡發展的理論去註清.
他將人生分為八個階段.
有嬰兒期,幼年期,學前期,學齡期.
青年期,成年期,中年期和晚年期等等.

$^{441}$每一個階段都有很多特別的東西.
另外也有一些學者將它發揚,發表.
以致有人說一個青年.
我們今天有一些年輕人坐在我們當中.
青年人會問什麼問題呢?.
就會問三個問題.
一般來說我們可以將問題分.
第一個就是我在尋找我是誰.
其他人怎樣期望我呢?.
我的形象是怎樣呢?.
我是否多過別人對我的標籤呢?.
這是青年人第一個會問的問題.
我是誰.
第二個問題就會問.
我適合在哪裡呢?.
How do I fit in?.
在什麼環境我會感到安全?.
我怎樣和其他人相戀?.
我在哪裡會感覺到被需要呢?.
這是青年人會問的第二個問題.
當青年人繼續成長的時候.
他會問第三個問題.
我能夠提供什麼貢獻呢?.
What difference can I make?.
我能夠為這個群體,為這個社會.
為這個世界有什麼貢獻呢?.
在什麼範疇我可以幫到別人?.
當我在前面打算的時候.
我怎樣去選擇呢?.
很多時候教會裡面的青年人的導師.
就會按照這些需要去幫他們解答.
幫他們去設計一些課程.
令他們可以對自己有信心.
亦都可以靠著和別人的相連.
去fit in 到.
亦都可以在群體裡面有貢獻.
蘇格拉底有一天在白天的時候.
他拿著燈籠在尋找東西.
有人問他你在尋找什麼?.
蘇格拉底說我在尋找我.

$^{481}$很多人都不知道自己是誰.
不單止青年人.
當我們繼續成長的時候.
我們很可能就找不到我是誰.
我們又看看大衛.
大衛他在這篇詩篇裡面.
他又說人是誰呢?.
在第四節 第五節.
他說人算什麼 理竟顧念他.
西人算什麼 理竟眷顧他.
你使他比神微少一點.
賜他榮耀 尊貴為冠冕.
這是和修本.
但如果我們看和合本的翻譯.
第五節就是說.
你使他比天使微少一點.
你可以用天使或者用神.
我們在看.
當我們去看人是誰的時候.
我們會看這首詩篇第八篇裡面.
有很多個「你」字.
你 你 你.
總共有十二個「你」字.
這十二個「你」字裡面.
有七個和神是誰有關.
而另外那五個和「我」.
即是人 是誰有關.
所以我們會看到.
人算什麼 理竟顧念他.
神顧念我們.
人有什麼可誇的呢?.
其實隨時都可能遭到各樣的天災人禍.
好像昨天一樣.
一場的黑雨.
終於有兩個人離世.
有四個人重傷.
有百多個人受傷入院.
我們沒有預算.
我們很多人都好像錯過了.
有些天文台的前副主任就會說.

$^{521}$我們現在人的工具追不上.
所以我們沒有可能很正確.
可以預測到.
接下來的風是多大.
人其實真的遭遇各樣的天災.
各樣的人禍.
最後怎樣呢?.
入了墳墓.
身體就歸於塵土.
不論富貴貧賤都是這樣的.
我剛剛上星期幫我先生的弟父殺灰.
真的 很年輕的一個 62歲.
他可能讀書很厲害.
在工作上都很厲害.
他是稅務局裡面的一個很好的同工.
但是他也敵不過一些疾病.
無論人富貴貧賤都是這樣的.
雖然我們這麼無力.
我們這麼無能.
但是神卻愛我們.
保養了我們.
顧念我們.
顧惜我們.
他時時想到人.
時時將人放在他的心中.
替人做了很多很多美好的計劃和安排.
答應我們的需要.
賜福給我們.
並且萬事為人的好處是效力.
我們施恩 發憐憫給我們.
所以我們說 神顧念我們.
神眷顧我們.
神更加給我們.
比神微少一點.
或者比天使微少一點.
從居所而言.
我們是地上高.
神或者天使在哪裡.
在天上高.
從能力來說.

$^{561}$我們的能力是軟弱到像蟲一樣.
但是神或者天使卻有可畏的大能.
他好像風 好像火一樣.
從實質而言.
我們人是被造的一個.
是由塵土而造.
但是神或者天使.
在實質而言.
天使是靈體.
好像神一樣.
人雖然有神的形象.
但是人卻是被造的.
很軟弱的.
但是神卻為人類.
預備了十架的救恩.
讓我們可以到末後的時候.
我們可以和神一起享受這個榮生.
這個就是天使所無的.
原來當初神造人的時候.
是照著神的形象和樣式創造的.
同時神也給我們人的榮耀 尊貴和權柄.
執行管理萬物的任務.
神造人的目的是要人帶一個冠冕.
這個就是說.
他賜人榮耀 尊貴為冠冕.
握著神所賜的權柄去做管理者.
治理神所造的大地和生物.
今天是九月九日.
大家記不記得2022年九月八日發生什麼事?.
2022年九月八日.
那天是英女皇逝世的一天.
英女皇逝世一年.
我們看到在英女皇逝世的國葬裡面.
很多人都在她的國葬裡面一起參與.
在她的棺頂有她的冠冕.
我們知道其實英女皇是神所設立.
神也讓她有這個榮耀.
我們也聽到不同的子民.
很遠很遠都要去到她出殯的地方.
去趕她去瞻仰最後 向她致敬.

$^{601}$我們看到其實可以看到.
他們是很想向英女皇做最後的道別和致敬.
一個好的國 一個好的冠冕.
給人一個很大的緬懷.
詩人在詩篇第八篇裡面指出神的名是偉大的.
他驚嘆神的無限 奇妙和大量.
另一方面他驚訝人的脆弱 渺小和卑賤.
是在這兩種情緒裡面交織.
詩人更大的驚訝就是偉大的神.
竟然願意主動去保護和看顧人.
將敗壞的人提升到一個極之尊貴的地位.
並且神更加給人一個尊貴的身份和使命.
派人去管治他所做的萬物.
包括加速地上的走獸.
天上高的飛鳥和海裡各類的魚類.
神就是這樣恩待人.
第三點我想說.
在說第三點之前.
我和大家說說這篇詩篇.
有些聖經學者很喜歡分析一些聖經的經節.
由第一節到第九節.
上下包著的 A 和 A'.
A 就是神威嚴的聲明.
我們看到耶和華的名.
我們的主的名在全地何其美.
最後一節也是.
耶和華你的名在全地何其美.
然後B就是說神在全地掌權.
然後他就會說.
你因敵人的怨顧.
從孩童和赫拉的口中建立了能力.
對你受敵和報仇的閉口無言.
並你所陳設的月亮星宿.
神在全地掌權.
然後他就說人的卑微.
然後他說人算什麼.
你竟怨顧他.
世人算什麼.
你竟怨顧他.
這裡是人的卑微.

$^{641}$但我們又看到神又給我們人的尊貴.
你使他比天使微小一點.
或者你使他比神微小一點.
更加賜他榮耀尊貴為他的冠冕.
然後B'是什麼呢.
我們現在看第六到第八節.
就是神委派人去掌權.
所以很多時候學者會這樣去研經的時候.
就將聖經的經文這樣去對比.
更加去明白聖經在說什麼.
所以我們去到最後第三段的時候.
就是神是委派人去掌權.
我就是用神對我的恩惠來說.
第一就是你派他管理你手所做的.
使萬物就是一切的牛羊.
田野的牲畜.
空中的鳥.
海裡的魚.
凡游在水裡的.
都伏在他的腳下.
所以看到原來神派他管理神所做的.
使萬物都伏在他的腳下.
所以這個就是治理的意思.
神派他的.
就好像創世記第一章這樣說.
我們要照我們的形象.
按我們的樣式做人.
使他們管理海裡的魚.
空中的鳥.
地上的牲畜.
全地和地上爬的一切昆蟲.
所以神就照他的形象做了我們.
照他的形象做男做女.
做我們每一個.
我們每一個都有不同的恩賜.
有不同的身份.
有不同的背景.
但我們每一個都是神所做.
我們每一個神都派我們管理他手所做.
然後我們再看神給我們的管理範圍.

$^{681}$第七節就會說到.
一切的牛羊.
田野的牲畜.
牛羊是家養的.
田野的獸.
我們人是沒有辦法控制的.
我們人可能會懼怕的.
但這些其實都是神給我們管理的.
是我們管理的範圍裡面.
所以我們不要像迷信那些人.
就要去拜這件獸.
拜那件獸.
其實無論是圈養的家畜.
或者在田野裡面的.
都是我們所管理.
然後更加說.
凡空中的鳥.
海裡的魚.
凡游在水裡的.
驚行海盜的.
都有我們所做.
所以才說人有保護環境的任務.
我們不可以任意暴殺.
我們不可以破壞生態.
神和人建立關係.
這個關係不僅止我們人和人.
更加包括我們的大自然.
在我們的秩序當中.
都應當保持和諧.
月亮是神所做的.
宇宙間好像月亮這麼大的行星.
其實千千百百萬粒.
我們人比月亮渺小很多.
所以詩人說人算什麼.
神竟眷顧.
今天當我們抬頭看月亮的時候.
我們會不會像大衛一樣.
想到神.
我們會不會也和大衛一樣.
想到人.

$^{721}$人是多麼的幸福.
人是有神眷顧的.
當我們有神眷顧.
我們還怕什麼呢.
為什麼我們輕看自己呢.
月亮雖然很大.
但是只是為人安排設備的.
飾物罷了.
是神擺設在太陽和月球旁邊.
人才是神的精心傑作.
宇宙萬物之中.
唯有人才稱得上是神心中所愛的.
所以神派人管理.
他所創造的一切.
我們人也要彼此的尊重.
我們要尊重其他的人.
因為其他的人.
都一樣是神所管理的.
神被派管理.
神所做的萬物.
不只是加速.
也包括野獸,飛禽和一切魚類.
人參與神的管理.
是神賜的榮耀尊貴.
但不是絕對的權利.
不可以獨立於神之外.
同樣我們人是管理其他的人.
我們也不可以獨立於神之外.
我們的管理.
是要降伏在神的管理之中.
當我們看見神創造奇妙偉大.
我們會奇怪.
為何他會眷顧他時時失望的人呢?.
我們知道我們人時時都叫神失望.
但神比我們比祂自己卑微一些.
當我們懷疑自己的價值.
或感到自卑的時候.
記住神看你是有很高很高的價值.
我們不要單單看別人怎樣看你.
我們要看神怎樣看你.

$^{761}$我們是極尊貴的.
因為我們是主的形象.
主是我們的創造主.
我們因著神.
已經宣告了我們對祂是何等有價值.
我們就當除去那些自覺沒有價值的自卑感.
我們人是要彼此的尊重.
我們人更加要知道.
其實神是尊重我們的.
我們要認真承受托管的管家的身份和任務.
這是高尚無比的身份.
我們可以跟神的關係裡面.
得到自我的身份.
我時時說到.
都是用這裡來結束的.
「於我何過」.
今天聽了我說了很久.
我又應該是搭三.
說了七個字.
八個字.
今天說了八個字.
可能過了我的時間.
但大家聽到這篇道之後.
你會怎樣回應?.
你會不會被月亮影響你的思緒?.
你會不會被月亮影響你的心情?.
或者你會被人影響你的心情?.
我們知道我們是神所造的.
神看我們是尊貴的.
神肯託付我們要做的事.
我們要努力去做好.
無論是生畜或者是野獸.
我們都要做好我們有的半份.
我們低頭祈禱.
天父上帝.
當我們來到你面前的時候.
我們要向你承認.
我們真的沒有想過.
因為你讓我們知道.
我們是尊貴的.

$^{801}$我們就願我們每一個.
都有尊貴的身份.
我們都要好好地.
用我們尊貴的身份.
來榮耀你.
來稱頌你.
來管理你賜給我們.
管理的一切.
聽我們的祈禱.
奉耶穌基督的名堂.
阿們.
\newpage



\section{提摩太後書路加福音 9:57-62}
\label{sec:eRch8741F34}
\textbf{2023-09-10 午堂《若要跟從主》路加福音 9\_57-62,提摩太後書 1\_1-12 講員:郭熹瑜宣教士}
\newline
\newline
連結: \href{https://youtube.com/watch?v=eRch8741F34}{\texttt{https://youtube.com/watch?v=eRch8741F34}} ~~~~ 語音日期: 2023-09-13
\newline
\newline
\hyperref[sec:LQDQ7CliZ7U]{\small{< < < PREV SERMON < < <}}
~
\hyperref[sec:index_chronic]{\small{[返順時目]}}
~
\hyperref[sec:index_scriptual]{\small{[返順卷目]}}
~
\hyperref[sec:Kpbv61GAjrM]{\small{> > > NEXT SERMON > > >}}
\newline
\newline
$^{1}$各位弟兄姊妹 早安.
早安.
很開心今天能夠在這裡跟大家分享.
也很開心在過去的十年裡.
其實也見證了神自己在我們工場裡面很多不同的工作.
在錦書家裡面一直在不同方面支持我們.
與我們同行.
還記得在2008年的時候.
我們在錦書家裡面實習.
2007年.
2007年的時候在這裡實習.
然後也在我們2014年出工場的時候.
教會也有支持我們.
而過去我們在2018,2019年回來述職期間.
其實我兒子爾信也有定期跟小芬學習性交.
所以今天也可以一起獻醜.
也見到教會在不同方面如何與我們同行.
以及支持我們.
非常感恩.
我們上一次在2018,2019年的時候.
回到香港有一年時間的述職.
2019年我們就回到工場.
繼續我們第二期的事奉.
這裡有一段經文其實就是亞比斯的禱告.
這個是我上一次述職期間其中一天的靈修.
當時我很記得我祈禱.
讀完這個經文我就祈禱.
神啊 究竟我第二期要我為你做些什麼呢?.
你要怎樣用我們的一家呢?.
我就向神大膽去求.
求神擴張我的境界.
想與我同在.
然後神就好像經文一樣.
祂真的去回應.
祂讓我們在第二期有不同經驗的服侍.
還記得2019年我們回到香港.
最初的三個月我們要重溫一些日文.
我不知道大家有沒有學過外語.
不過我們回來一年.
回到香港的日文都跌窩跌得很厲害.

$^{41}$所以回到香港就要讀三個月.
之後我們在2019年.
由5月開始到2020年3月.
我們就參與了一間叫做平安傳基督教會.
這間教會裡面的開方直堂報道的服侍.
你見到有一個表.
一系列列了的工作.
就是當時我們在那大概八個月的時間.
我們去做的.
當時我們和另外三位宣教士組成了一個團隊.
就在這個教會在這個地區裡面做報道的工作.
我們做無論是料理班.
還是小朋友的活動班.
興趣班.
或者是親子的那些未入學的小朋友.
親子的活動班.
都有的.
另外我們都有搞ukulele的小組.
讓到有興趣的弟兄姊妹.
和附近的居民可以來參加.
我們其中一張相片.
就是我們去用ukulele.
在冬天站的地鐵站外面去報街音.
不過你看到就好像很好.
其實周圍圍著都是雪.
冷到手指都變硬了.
但這個都是神秘愛.
有機會在那段時間參與的服侍.
也都有機會和弟兄姊妹一起經歷了教會.
當年是他們開始了十年的時候的一個慶祝.
特別感恩的是.
雖然我們只是短短在那裡八個月的時間.
但神讓我們有機會去接觸.
不同的新來教會的朋友.
特別是一些年長的新朋友.
你看到相片裡面.
其中一個太太.
當時我和她有機會.
她回來我邀請她.
做一些初訊的查經,栽培.

$^{81}$經過大概三個月的查經之後.
這個太太是決志信主.
而且在查完經之後半年.
她再上浸禮班.
她就在教會受浸.
當時這個太太是71歲.
在日本我們常常說是福音硬土.
這個都是事實.
都仍然是.
不過神也讓我們見證.
也見到神自己在當中的工作.
因為祂從來都沒有放棄日本人.
也讓我們繼續經驗.
上帝使用我們的恩典.
所以我們非常感恩.
一回到去.
神就給我們有機會結果子.
剛才那個禱告都說到.
求神擴張我們的境界.
在想不到的情況之下.
2020年的3月.
我們就開始承擔了.
除了一些報道工作之外.
更加承擔了.
差會在北海道.
我們有自己的語文中心.
而我成為了當中的語文學習顧問.
語文學習顧問是什麼呢.
其實主要是因為我們會有新來的宣教士.
他們需要學日文.
我的工作就是安排他們怎樣學日文.
和跟進他們每一個人學習的進度.
我們的課堂全部都是一對一的.
所以我就要跟進每一個人.
學得怎樣.
學得快 學得慢.
也要陪伴每一個新來的宣教士.
和他們的家庭.
去到經歷來到第一年工場.
第一年 第二年.

$^{121}$在工場的所有適應.
所以神都給我們有機會.
做多一點.
不單只是服侍日本人.
更加服侍我們自己差會的宣教士.
很感恩.
雖然全世界那幾年都在疫情裡面.
很多的不容易.
但神就為我們預備了.
有很多對日本有心的宣教士.
來到我們的工場.
我們有一段時間封關.
大家都來不了日本旅行.
但到了2022年3月開關.
其實到現在.
我們已經有接近30個宣教士.
帶著25個小朋友來到我們工場.
厲不厲害.
而這一堆的宣教士和他們的小朋友.
都是今天我們去事奉的對象.
你看到相片裡面.
就是第一批.
一開關.
第一批已經等了一年多的宣教士.
來到的時候.
我們一起拍了一張照片.
除了去幫助宣教士之外.
我們更加有14個.
是日本的童工.
是日本人.
我們聘請他們來教宣教士日文.
我們的工作包括管理.
我們所有的日文老師.
定期跟他們開會.
和定期跟他們做評估.
或者提供培訓.
所以對我來說.
這也是我們在第一期宣教.
沒有經驗過的事.
沒有做過的事.

$^{161}$都是神給我們一些很新的學習和嘗試.
除了這個語文學習顧問之外.
神更加在2020年7月的時候.
又讓我們有機會承擔多一個崗位.
就是成為了北海道地區的事工主任.
差會我們分了三個地區的事工.
有最北的北海道.
然後有東北地區.
然後有近東京的我們叫關東地區.
我們承擔的就是最北.
北海道那個島的地區的事工主任.
當中我們的工作除了剛才說的.
更加包括我們要管理和協調.
所有在北海道我們差會的宣教士.
他們的工作.
他們和教會的合作.
無論是派哪些宣教士去教會工作.
派去哪些教會.
都是我們其中一部分的工作.
另外我們也有機會.
因為這個職責的緣故.
特別是希如多一些機會.
會飛到東京和所有地區的主管.
一起去開會,溝通.
也有機會飛到不同的地方.
去到港島.
在相片裡面你們看到有一座建築物.
就是我們北海道的辦公室.
這座大樓其實只有三層.
最底的那層就是我們的辦公室.
中間那層就是我們的語文學校.
最高那層就是我們的禮堂.
用來聚會的地方.
我們的工作之一包括管理這個大廈.
管理大廈剛才弟兄也說了.
前兩天水浸.
所有同工都要出動去清潔.
我們也是同樣的.
就是管理這座大廈的所有東西.
很奇妙地過去一年.

$^{201}$不知為什麼不斷爆水管.
廁所又壞了.
我們的工作包括去處理.
抹地,掃水和弄廁所.
這些都是我們過去兩年多的經驗.
這一批是什麼呢?.
這一堆就是我們有機會.
可以跟日本很多不同的教會一起合作.
特別是有機會跟一批.
在日本比較年輕的牧師去同行.
我們上次回來塑膠的時候有分享過.
日本的牧師平均年齡已經到65歲以上.
但是沒有年輕人也不是.
很感恩神都興起了一批40多歲的牧師.
他們剛剛讀完神學.
出來開始承擔教會牧羊的工作.
我們因著我們的角色緣故.
我們有機會跟他們同行.
一起商討他們教會的福音事工發展.
然後派出宣教士跟他們合作.
這裡就是其中一部分的童工.
在分享裡面.
剛才Karen也特別分享到我們的一些工作.
為什麼也想分享一下我們的事工呢?.
很想讓弟兄姊妹都知道.
當你們支持我們的時候.
我們其實是在參與一些什麼類型的工作.
有直接在前線的報道工作.
同時也有一些是領導.
或者要去管理宣教士的團隊.
去幫助新的宣教士去適應的工作.
特別我們是海外基督使團的宣教士.
我們很重視的.
都是能夠去幫助每一個童工.
跟他們去同行.
而我們其中一個特色就是.
會四年在工場.
一年才回來.
有一年就回來述職.
一期就五年.

$^{241}$所以可能在十年裡面.
只能見到弟兄姊妹.
第二次再回來述職去見大家.
但是在工場上也很深感受.
原來在香港.
每一個弟兄姊妹教會.
你們對我們的支持,代討.
都是很重要.
都是去幫助我們能夠在前線去服侍神.
今天我們講的主題.
特別是經文的內容都講到.
若要跟從主.
在這段經文裡面.
其實都有三個人講到要去跟從神.
回想起最初.
我們為什麼會踏上宣教的路.
都是一個很簡單的立志.
為什麼會成為基督徒.
都是一個很簡單的回應.
就是主耶穌說我很想跟從你.
主耶穌說我知道你愛我.
我想去回應你.
正正在這段經文裡面.
就有三個人.
這三個人分別都去到耶穌面前.
有兩個是主動跟耶穌說的.
主動跟耶穌說主耶我想去跟從你.
而有一個是耶穌向他發出邀請.
跟從我來.
當這三個人在一段很短的經文裡面出現的時候.
我們就會想為什麼聖經裡面.
將這三個人放在同一段這麼小的段落裡面.
原來整個的大段落.
在路加福音9章23節.
耶穌曾經這樣說.
聖經裡面這樣說.
耶穌有對眾人說.
有人要跟從我.
就當捨己天天背起他的十字架來跟從我.
這段經文是我們很熟悉的經文.

$^{281}$耶穌向我們眾人發出邀請.
我們是不是願意去跟從耶穌.
然後要作的回應是.
就當捨己背起我們的十字架去跟從耶穌.
這三個人.
我們看回這三個人.
這三個人裡面其中兩個都很清楚說過.
你無論往哪裡去.
我要跟從你.
曾幾何時我們有沒有這樣去回應神呢.
曾幾何時我們在唱詩歌的時候.
在我們靈修的時候.
有沒有都有同樣的一種情緒.
一種心意思意念.
神啊無論如何我都會跟從你.
神啊我會獻上我的一切去給你.
但是過後會不會都有這幾個人的一些想法.
主啊容我.
但是主啊容我去做這些那些.
容我去完成這些和那些.
我就可以放心.
放下一切去跟從你.
三個的人.
三個不同的處境.
今天不會有機會我們可以深入每一個處境去想.
究竟這些處境和我們有沒有關係.
我們相信這些經文我們都熟悉.
耶穌對這些人說的處境的回應.
都是叫他專心去仰望神.
專心去看清楚.
你既然今天願意去跟從主的時候.
你願不願意都真是定睛望著主耶穌.
放下我們覺得這一刻對我們來說是重要的事情.
不過當我們說放下.
當我們說撇下的時候.
又要再看回整個路加福音的大段落.
當耶穌作出邀請.
如果有人要跟從我的時候.
除了有這三個人的出現.
去到第十四章.

$^{321}$耶穌同樣去說到.
凡不背著自己十字架跟從我的.
也不能作我的門徒.
但是在這句經文之後.
祂說出的就是.
如果我們要建一個房子.
應該先計算一下代價.
不計算好代價.
建不建得成才開工的話.
其實可能不是最有智慧.
接著另外一個比喻說到.
如果我們去打仗.
不夠打的時候.
可能先應該去派使者求和.
因為最後三十三節都是這樣說.
你們無論什麼人.
若不撇下一切所有的.
就不能作我的門徒.
開始是說要撇下一切.
最後的經文說撇一切.
不過中間提醒的是我們要計算代價.
我們有沒有好好地計算代價.
在我們計算代價的同時.
其實也要想一想.
我們有時候也會鼓勵.
要憑信心去做.
但當我們說憑信心的時候.
其實我們有沒有看回現實.
究竟我們是否知道面對著一個什麼樣的挑戰.
我們在說要付上一個怎樣的代價.
剛才Karen也分享到.
在疫情裡面特別Karen有機會和很多新來的宣教士同行.
曾經有一個宣教士.
就這樣跟他說.
我一直都知道要做宣教士.
做宣教的工作真的很不容易.
不過我沒有想過原來真的難得這麼厲害.
我沒有想過原來真的這麼艱難.
在他的處境.
和大部分宣教士一樣.

$^{361}$面對著什麼艱難呢.
語言上的挑戰.
我們很多宣教士來到.
其實連去超級市場買藝都覺得不容易.
在疫情裡面他們自己要去醫院.
中了COVID之後.
我們是要陪他們去醫院看醫生.
其中一個原因就是他們連填自己的地址.
連填所有的表.
醫生問他所有的東西他都回答不到.
語言是一個最基本最開始的挑戰.
生活上的挑戰.
我們自己集訪.
集訪冬天下雪的時候.
除了冰天雪地很美.
大家去旅行覺得很開心之外.
就是每天生活的艱難.
我們有宣教士在下雪之後.
結冰了.
我們有宣教士在雨汶中心下雪.
斷了一條肋骨.
我們要教宣教士.
怎樣一步一步慢慢去學行路.
我們有童工.
在心靈裡面.
在他情緒上面面對不同的艱難.
因為在一個過渡時期.
離開家裡又來愁別井的情緒.
又掛念家裡的情緒.
疫情裡面更加多這些心靈的需要.
文化上.
團隊上.
作為一個國際差會.
我們不只是適應日本的文化.
還有很多來自世界各地.
不同的童工他們背後的文化.
在這些當中都是一些挑戰.
我們在工場裡搬了五次屋.
雨汶是由我們出發十年前.
到現在總共轉了七次學校.

$^{401}$這些全部都是每一個.
宣教士將會.
或者已經在面對的挑戰.
但是神.
他跟我們同行.
神必定會引導.
我們清楚見到那個代價.
願意回應.
然後就憑信心去到踏上.
我想這個就是神給我們.
很清楚的呼召.
回到香港這段時間.
我們有幾個月.
無論去到不同的教會裡面去推動.
關於宣教或者都在.
海外教士團辦公室裡面有服侍.
不過同時都有一些機會去多些閱讀.
看書和分享的時候.
有機會看了.
隋掌管明天這本書.
是何以思宣教士.
他在1910年代的時候.
來到中國去宣教的很多經歷.
在這本書的記載裡面.
他自己說道.
他有一天都問自己為什麼今天要搞到這樣的田地.
為什麼要自討苦吃.
其實他在討什麼苦吃.
什麼事情已經在發生.
其實他只是在學廣東話.
學極都學不懂.
講極都講不正.
經常都被人笑.
在他覺得很鬱悶的時候.
他心裡面就響起這段聲音.
不是你自己要來的.
是我帶領你來到這裡.
因為我與你同在.
我在你身邊.
這是神對何以思宣教士的回應.

$^{441}$我想我們每一個人.
去跟從神的時候.
多了說我們願意跟從神的時候.
神的應許同時是伴隨著.
在路加夫音整個大段落.
去到第十八章.
是那個段落的總結.
這個段落去到總結的部分的時候.
先是有一個少年的官.
那個少年的官就想去跟從耶穌.
問耶穌.
我們都要變為一切去跟從的時候.
他就悠悠瘦瘦地走了.
然後彼得就去到回應.
我們都變成一切了.
我們有什麼能夠去得到.
然後耶穌就這樣去回答他.
我實在告訴你們.
人為神的國.
放下房屋或是妻子弟兄父母兒女.
沒有在今世不得百倍.
在來世不得永生的.
這張照片是我們一個宣教士.
在愛爾蘭來的.
單身的女宣教士.
在我們大學的時候.
我們所有同工和小朋友都跟她拍照.
她就用這段經文去鼓勵我們.
她說她撇下了她的家人朋友.
在工場裡面三十多年.
但是神按著她的應許.
去祝福她.
她在我們工場裡面接觸過的宣教士.
宣教士的小朋友.
她特別去支援很多宣教士的小朋友.
的孫子.
她說我接觸過超過一百個孫子.
與他們同行.
而正正在她離開工場的時候.
使團在北海道的同工有六十位.

$^{481}$有四十多位小朋友.
神按著她自己的應許.
去回應這個宣教士.
神會親自去祝福.
親自去幫助.
當我們說若要跟從主.
第二段想看的經文.
是在提摩太后書一章一至十二節.
今天沒有讀到這段經文.
但當我們一方面說我們要撇下.
但一方面我們要看神怎樣去祝福每一個跟從主的人.
我們認清究竟神怎樣去祝福我們願意跟從他的時候.
或者這才是我們回應神呼召最重要最重要的基礎.
在一章一節這樣說到.
照著基督耶穌女生命的應許.
作基督耶穌的使徒.
然後第二節說到連文.
第六第七節說到恩賜.
最後到第九至十二節說到神自己的保存.
鼓勵你今天有機會可以再讀這段經文.
但想用四點去說一下.
若我們要跟從主.
我們需要認定的是主自己的應許.
若我們要跟從主.
我們要思想活在主自己的連文裡.
若我們要跟從主.
我們要知道其實神已經將各樣的恩賜賜給我們.
賜給教會.
若我們要跟從主.
我們必須要認定主自己是託付我們的那位主.
但同時也是那位去保守成全的主.
提摩太後書.
特別第一章第二章.
裡面說的是保羅提醒提摩太.
要聽從主的命令.
不是從自己的意思.
或者是一個自討苦吃.
很不容易或者我們說的是他甘願面對苦難痛苦的狀況.
刻意進入一個很艱難的狀況裡面.
選擇在艱難裡面堅持.

$^{521}$甚至看到的不是現況裡面的艱難.
是超越現況看到神給我們的盼望.
這是提摩太後書第一第二章主要的內容.
而第一節.
剛才也讀到.
保羅提到自己事奉神的身份.
是奉神自己的旨意.
是神自己的呼召.
是神自己的應許.
而他因為照著耶穌基督女生命的應許.
所以可以作試圖.
可以回應神.
我還記得我們去回應神自己的帶領.
參與了宣教的服事.
去出工場之前.
神就讓約翰福音十五章十六字的經文成為我們的鼓勵和提醒.
不是我們揀選了要去回應.
是神揀選了我們.
是神讓我們能夠出去為主去到結果子.
在第一期四年裡.
其實很多不同的工作.
在職堂的報道裡.
都沒有機會和人很深入地讀聖經.
甚至幫助人去到做初信栽培.
帶領人信主.
但在過去四年.
是神讓我們見證到八位頂尖妹受浸.
是神讓我們經歷到結果子的喜樂.
這八位頂尖妹有四位是我們剛才說的.
在平安傳教會裡.
我們有機會個別和他們做初信栽培.
有部分是我們的同工和他們做初信栽培.
然後我們加入教會.
大部分都是六十歲七十歲弟兄姊妹妹.
另外有四位頂尖妹是比較年輕弟兄姊妹妹.
在一間日本教會裡和教會同行.
第一期的時候已經見到有部分的教會頂尖妹.
他們的小朋友在五六年級.
教會的牧師我們就一起有青年的聚會.
到過去兩三年.

$^{561}$縱使在疫情裡.
這班頂尖妹已經成長.
已經去到高中讀大學.
他們繼續留在教會.
教會這六七年裡和我們同行栽培他們.
最後他們都願意去受浸加入教會.
甚至有一位的姐妹.
她也在大學裡帶了一位未信的朋友回來教會.
這位弟兄後來也受浸加入教會.
神自己的應許是沒有落空.
是神自己的應許讓我們能夠結果治.
是祂揀選我們每一個熟悉祂的人.
能夠在我們的處境裡結果治.
可能時間不是即時.
但是神自己的應許必定會去到實現.
第二點很想說的是.
要活在主的憐憫裡.
提摩太前書後書.
保羅的閩安語也這樣說.
願因為憐憫平安.
從父神和我們主耶穌基督歸於你.
是很獨特的.
保羅大部分書信都是說願因為平安.
不過去到和提摩太說的書信裡.
特別是加上憐憫在中間.
相信保羅很清楚知道.
作為一個事奉神的人.
很需要主的憐憫.
而主的憐憫是什麼呢?.
無論舊約聖經和新約聖經都提醒我們.
神就是憐憫的神.
神也叫我們用祂的憐憫.
去憐憫我們身邊的人.
我們要學讓主去憐憫我們身邊的人.
當我們再看看憐憫的意思的時候.
其實憐憫就是指主和以色列人之間.
立約的關係裡會有的相應行動.
然後有了那份關係就產生了憐憫.
因此和憐憫都是神賜給我們的.
在日本我有機會去到不同的教會.

$^{601}$有港島的服侍.
一方面因為我們自己的工作崗位.
很多教會都會來邀請我們.
曾經有一個教會打電話來.
和我們分享.
教會的牧師就特意打來.
他說我有三年裡沒有放過一天假.
三年裡每個星期的主日都是在港島.
我未來一個星期希望放假.
你們可不可以派孫教士來港島幫我們.
日本很多教會都只有一位牧師一位牧者.
三年裡他沒有一天休息.
是一個73歲的牧者.
他退休了後來看到教會有需要.
就繼續出來去承擔.
於是我們就去幫助.
另外一個教會也因著他們的傳道人.
在很短時間裡有抑鬱的情況.
就需要停工後來離職.
教會也有執事會的主席打來問我們.
可不可以來幫忙.
這間教會是我們之前在ONMAP建立.
已經有十多年是由日本的牧者他們去帶領.
在這個處境下我們就嘗試去回應.
有一年半的時間.
我和Karen都會每個月去到那個教會.
去到港島.
和他們有帶領主餐的時候.
之後我們就請回一個牧師.
然後我們就可以和他們的教會同行.
很多時候去到教會港島的時候.
有很多的機會我都分享過一個見證.
很想和大家分享.
特別是港島主的年敏.
是一位台灣的姐妹.
她在她成長的過程中經歷了神自己.
經歷了家庭裡的愛.
然後她就立志在舊約裡做更多去推動的工作.
在一個大學的畢業禮裡.
她分享了一個故事.

$^{641}$這個故事說到在她十一歲的時候.
她的媽媽就患上了癌症.
然後要接受化療.
看到媽媽身體很虛弱.
心裡第一次發現可能媽媽會離開世界.
心裡有很多驚恐和擔憂.
但是她回應這些驚恐和擔憂.
她做了一個決定.
就是她要更加努力讀書.
她要考第一.
她要做全班最好的一個學生.
只有一個原因.
原因就是很想她做得更好的時候.
媽媽會不捨得她.
不會離開世界.
於是有一天她考到第一.
拿著她那張成績表.
打算回到家裡給媽媽看.
她媽媽就叫她.
她媽媽說:你來到我床邊吧.
我有話要跟你說.
她說:女兒,我有件事很想你知道.
很想你知道.
當媽媽有一天真的離開了世界.
你要記住.
無論你是做得多麼不好.
甚至被人抓去坐牢.
你考試考到最後.
都不重要.
媽媽是從起初就愛你.
媽媽從起初就愛你.
這個從起初就愛你的聲音.
在她心裡不斷響起.
成為她自己很大的提醒.
不是因為她做得多麼好.
是因為那份關係.
她媽媽從起初就愛她.
我們跟天父的關係何嘗不是這樣呢.
在我們還作做人的時候.
我們天父已經愛我們.

$^{681}$耶穌基督就為我們釘上十字架.
不是因為我們做了什麼.
我們配得這份愛.
是因為主自己的憐憫.
從起初就在這份關係裡.
淋到我們當中.
弟子們,如果我們不經歷主這份憐憫.
我們怎樣能夠將這憐憫帶給身邊的人.
當我們不經歷原來神的憐憫.
是口口地淋到我們心裡的時候.
我們又怎樣能夠宣講這份憐憫.
讓我們身邊還沒經歷神的愛的人.
能夠經歷這份愛.
我們時常都記住.
若果要跟從主.
我們必須活在主自己的憐憫裡.
第三點.
很想說的是.
神其實給了我們各樣的恩賜.
不過在提摩太后書裡說.
恩賜如火挑旺之前.
先說到提摩太太承繼了.
她能夠有的.
是祂祖母和母親無畏的信.
我們在承傳著一個怎樣的信心.
我們在教會裡承傳著一個怎樣的傳統.
其實整個教會不是一年兩年.
是在說過去十年二十年三十年四十年五十年.
很多弟兄姊妹一起有一個信仰的傳統在當中.
這個信仰的傳統是神給我們的恩賜.
我們怎樣去擁抱.
在創新裡繼續繼往開來.
但同時.
當然保羅都提醒提摩太太.
他得到的恩賜.
保羅去按手藉著.
賜給提摩太太的恩賜.
會再一次如火挑旺起來.
第七節是這樣說的.
因為神賜給我們的不是膽怯的心.

$^{721}$而是剛強人愛謹守的心.
恩賜可以說是才幹能力.
恩賜也可以說是剛強人愛謹守的心.
有時候我們看到很多事情都很多懼怕.
但懼怕的靈不是出於神.
神給我們的是不膽怯的心.
剛強人愛謹守的心.
按著我們的屬靈成全.
信仰成全.
按著神給我們的教會裡各樣的恩賜.
每個弟兄姊妹你們各有的能力才幹.
我們可以去服事主.
可以去回應主.
在北海道.
我們冬天都有機會出去戶外的活動.
這張照片其實有一次我們去露營的時候.
就把地下的冰挖出來.
砌了一個磚的牆.
冰的磚要砌一面牆.
都很不容易的.
一來我挖了出來的時候.
不是正正方方的.
很多硬硬的.
硬硬的砌的時候.
就覺得一推就倒.
第一個辦法.
可以令它不倒.
而且變得很堅固.
就在零下十度的處境下.
我們只要淋一點水.
在五至十分鐘之內就結冰了.
結冰之後就很穩固.
我想我們每個人都硬硬的.
我們每個人有我們的恩賜.
可能都不是砌得很完美.
發揮得很好.
不過當聖靈淋在當中.
就好像那些水淋在冰上面的時候.
就將我們連繫.
叫我們能夠牢不可破.

$^{761}$這個是神自己給我們有信仰的承傳.
有各樣的恩賜的美好.
我們在這間房間有六十位同工.
整個OMF在日本工場裡面有140位同工.
我們都來自不同的背景.
有不同的信仰承傳.
但我們都見證到神給我們能夠配合我們的恩賜.
活出主的愛.
然後去傳揚福音.
在我們的侍工裡面.
有很多同工會想用很多創意的方法去傳福音.
包括運動,cafe,教育樂器.
甚至一些藝術.
我們都試過在我們自己的玻璃女中心外面.
有一個雪場建立了出來.
然後用一些聖誕的裝飾.
有馬槽的圖像.
希望幫人去認識信仰.
神讓我們每個人都有不同的恩賜.
彼此配搭去服侍主.
最後很想說說日本教會的情況.
我們在北海道都和十幾間教會有合作.
每間教會都有不同的處境.
但都看得見.
其實神用我們的宣學士不同的恩賜.
去配合幫助不同的教會.
頂尖的我們的恩賜預備好.
一起彼此配搭去服侍主了沒有呢?.
最後很想說下去.
第四點 主的託付和保守.
主其實託付了很多使命給我們.
按著祂的旨意和恩典.
特別在十二節這樣說道.
「為者緣故,我也受者些苦難,然而我不以為此,.
因為知道我所信的是哪一個,.
也深信他能保存我所交付他的,直到那日.」.
同時有一個日本說道.
「也深信他能保存我所交付他的,直到那日.」.
相信兩個翻譯都對.
一方面是神託付給我們.

$^{801}$神把責任交給我們.
但神交給我們的責任.
我們同樣透過禱告.
同樣透過獻情去讓神去到成就.
我們只是一個中間的媒介.
我們在過去幾年有很多不同的承擔.
很多的承擔是遠超過我們覺得自己有的能力和經驗.
不過神給我們有以下幾個學習.
特別有兩點很想提.
第一點是聆聽的功課.
我想作為信徒,作為服事主的人.
作為跟從主的人.
我們要聆聽神的聲音.
但神都叫我們要聆聽我們內心的聲音.
很多艱難不容易的時候.
我們有沒有正面去處理.
我們要聆聽我們身邊的人的聲音.
認真去聽,而且去回應.
最後聆聽社會,世界的需要.
另一點我們覺得很重要.
神提醒我們的是.
神自己是掌權作王的主.
我們不是憑自己的意思.
是神自己的主權.
我們需要越時常用禱告去主人面前.
讓主自己的主權去成就.
這張照片是我們後院種了西瓜出來.
除了宣教工作,也有我們生活.
種了這個西瓜很開心.
七公斤那麼重.
但不是一年就種得到.
我們種了五年才種到一個西瓜.
西瓜每年都死的.
每年都種新一棵下去.
不過第一年種完,開花沒有果子.
第二年都是這樣.
第三年大了一點.
但是都大不了.
第四年好像西瓜球那麼大.
覺得可以試試.

$^{841}$切開它,白白的,不好吃.
終於到第五年才結果子.
才能有一個能吃的果子.
但第六年行不行我們都不知道.
不過每年都改進.
每年都去改良.
每年都嘗試去到種.
在我們參與日本一間教會裡面.
它有這個報道的策略.
說要先鬆肚.
部分教會的聚會其實是不講福音的.
等人進來教會.
是鬆肚的工作.
然後殺種,淋水,收割.
日本的狀況可能要二十年.
由鬆肚到那個人真的信主.
是要二十年的時間.
我想香港其實都可能是這樣.
因為很多基督教的幼稚園,中小學.
其實幫了很多鬆肚的工作.
以致人相對地容易去接受福音.
我們見到這個時間表會不會膽怯呢.
還是其實神都給我們有剛強的心.
去接觸很多未信的人.
這個圖有一千個人.
只有四個藍色的是基督徒.
可能幾乎見不到.
日本的處境就是這樣.
只有0.4\%是基督徒.
見到這個圖你會覺得很沒有盼望.
沒有辦法去幫助那九百多人去信耶穌.
但我深信是神交託給我們.
要繼續去傳揚主的福音.
雖然不容易.
但神同時會去到承擔.
在過去幾年.
我們很被鼓勵的.
聽到這首詩歌很被鼓勵.
這首詩歌特別在副歌說到.
為我心之所信的是誰.

$^{881}$並且也深信.
祂必定能夠保守我所交付他的都傳避.
直到那一日.
等於我們有沒有一些東西.
是神已經交託給我們.
叫我們要去承擔.
而我們是不是覺得很不容易.
很艱難.
但神的應許是.
我們要知道我們信的是哪一個.
我們跟從的是哪一個.
祂必定去到成就.
很想請以信.
今年十四歲了.
我們出工場的時候五歲.
他有一個見證的分享.
真的很感恩.
可以五年前在這裡.
有機會學習性言語.
然後可以五年後回來這裡.
可以有機會在崇拜獻奏.
不知道大家看到這張相片.
會想些什麼呢.
這張相片是在2020年的夏天.
拍攝的一張相片.
這塊白板寫著什麼呢.
有一個圖.
有很多很小的字.
這是一個心臟的圖.
當年我在這裡讀書.
去準備考入一間日本的中學.
可能大家看過日劇.
日本的漫畫都會知道.
日本其實考試的競爭是很大的.
要在之後中學才會讀的內容.
就要在中學考入一間中學的考試.
就會出的了.
我覺得讀書是好像行山一樣.
行山就是要離開自己的一個舒適區.
要做好準備才可以離開這個地區.

$^{921}$行山有很多很難行的地方.
好像雪溪,河邊,石澗.
不同地方有不同的難處.
固然會有容易行的地方.
但是應該有很多地方都是很難行的.
當年我考試讀書的時候.
都有想過放棄不去讀,不考試.
就好像行山放棄不再繼續行.
最終行一個山會怎樣呢?.
就是會到達目的地.
去到山頂.
但這個不是靠自己的能力.
而是有神的幫助.
有很多不同的弟兄姊妹的代賭.
我們才可以去到山上一個最高的地方.
最終我考到這個試.
考到入這間學校.
這間學校是一間天主教的學校.
是日本難找的一種類別的學校.
在正門都有一個很大的十字架.
在這間學校的校訓就是作炎作光.
我覺得這個校訓是很好.
然後也想用這個校訓.
希望可以帶多些朋友回教會.
最後想用多一個比喻.
這張是冬天在波萊島行山的一張照片.
在波萊島行山是需要一雙很特別的鞋.
叫做Snow Shoe.
沒有這雙鞋我們就會沉入樹林裡.
然後很難走.
我覺得這雙鞋就像神的幫助.
在一個很難走的地方.
好像人生的路.
神就會像這雙鞋托著我們.
讓我們每一步都可以容易走.
雖然有神的幫助.
我們也需要自己去走.
我們要一步一步靠著神.
還有很多不同人的代祖的幫助.
然後去走人生這條路.

$^{961}$要過山過谷.
謝謝大家.
若要跟從主.
我們必須認定主的應許.
若要跟從主.
記住神是愛我們.
神的憐憫環繞我們.
我們必須要跟從主.
我們必須要認定主的應許.
若要跟從主.
記住神是愛我們.
神的憐憫環繞我們.
我們必須要跟從主.
我們必須要認定主的應許.
若要跟從主.
記住神是愛我們.
神的憐憫環繞我們.
我們必須要活在主的憐憫當中.
神已經賜給我們各人有恩賜.
而且幫我們教會成為一個教會.
彼此恩賜配搭去服事主.
在社區裡面為主作炎作光.
主的託付.
主的保守是很真實.
若要跟從主.
我們要認定主是保守.
同時聽清楚主已經託付我們.
路加福音9章23節這樣說.
耶穌有對眾人說.
若有人要跟從我.
就當捨己天天背起他的十字架來跟從我.
若一齊祈禱.
天父上帝我們向你禱告.
多謝你.
因為是你自己呼召我們.
因為是你應許.
你自己要去揀選我們.
願你幫助我們.
繼續在我們生命的不同場景裡面.
去服事主.

$^{1001}$主亦都多謝你.
因為我們雖然不配.
但你將很多恩賜,能力,才幹,信仰的傳統賜予給我們.
讓我們能夠承擔.
你給我們不同的積分.
願你幫助我們去回應.
願你幫助我們去認清.
你必定保守,必定大力.
主讓我們天天背起我們的十字架.
去跟從你.
我們這樣禱告奉耶穌的名.
\newpage



\section{撒迦利亞書}
\label{sec:Kpbv61GAjrM}
\textbf{2023-09-14 周四晚查經班《歸回先知書研讀(1)撒迦利亞書》 講員:曾裔貴牧師}
\newline
\newline
連結: \href{https://youtube.com/watch?v=Kpbv61GAjrM}{\texttt{https://youtube.com/watch?v=Kpbv61GAjrM}} ~~~~ 語音日期: 2023-09-16
\newline
\newline
\hyperref[sec:eRch8741F34]{\small{< < < PREV SERMON < < <}}
~
\hyperref[sec:index_chronic]{\small{[返順時目]}}
~
\hyperref[sec:index_scriptual]{\small{[返順卷目]}}
~
\hyperref[sec:BaO8MW3rJk0]{\small{> > > NEXT SERMON > > >}}
\newline
\newline
$^{1}$各位弟兄姊妹晚安,我們一起祈禱,感恩我們又能夠在禮堂或線上家中透過Zoom.
可以一齊去研讀神的話語,我們一直都在歸回的歷史.
有三卷書卷我們來研讀,我們可以回頭去看.
在被老歸回的時候,才知道怎樣配合建造.
和要發出神秘的訊息,預言等等.
我們很盼望這些我們比較少去讀的經文.
崇拜可能比較少會講,可能會有少少陌生.
但我們都希望能夠透過一起的研讀,有少少掌握.
雖然有些未必容易解,但都幫助我們能夠透過不同的體裁的經文.
讓我們看到上主祂的作為和多樣性的啟示.
使我們能夠有一個更加豐富對神的作為的認識.
求主帶領我們一起在你面前等候禱告.
奉耶穌基督的名,阿們.
上星期因為有些技術性的問題,上一堂課錄不到.
所以這次我們要保險一點,Zoom也做了後備錄影.
方便如果音響錄不到,都可以有一個補救.
否則會有些浪費,有些人想聽,聽不到就有點可惜.
為何會講撒迦利亞書?.
因為之前我們看《以斯拉記》的時候.
因為它只有兩章,很快我們看完《哈該書》.
撒迦利亞和哈該都是同時期的先知.
到最後一卷的《馬拉記》,是另一個時代.
但是因為我們又要看《尼軒米記》.
我們先放下撒迦利亞書.
所以撒迦利亞書今天我們會開始講.
要幫助大家的記憶.
很早期的,在大利烏王的時候.
即是說是剛剛古列王吩咐他們上去重建.
過了大概十五年,那個建造工程.
因為當地的原居民,耶路撒冷的原居民.
阻撓那個建造,所以一直停頓.
停頓到大利烏王的時候.
在那個時候,首先是由先知出來.
有神的密使去吩咐.
當然是和當時的領袖所羅巴百.
和大祭司約書亞.
來到去有一個雙方面的.
好像大家彼此的勉勵.
來到去再鼓動百姓.
所以先知哈該,他總共有三篇講道.

$^{41}$我們以往提過,不詳細再回頭看.
其實都重要就是在那個建造裡面.
最重要就是既然已經有一班人.
其實五萬多人,第一次回去重建聖殿的時候.
但是遇到一些阻力.
連所羅巴百,連大祭司約書亞.
都招架不住那些阻礙.
所以就停了十五年.
所以都是很荒廢,很荒涼的爛尾樓.
耶路撒冷原來都有爛尾聖殿.
只是建了一個根基,然後就沒有辦法再做了.
所以當時的百姓又不會回頭回到波斯帝國.
所以就開始私算自己的重建.
自己的居所,自己的耕種,自己的買賣.
因為群龍無首,沒有人出來領導百姓.
使大家的心再次向神.
所以就有時間到了.
神就感動哈該和易多的孫.
就是這個撒迦利亞.
所以就分了兩卷書.
各自有各自的領袖.
各自有神秘的密士的寓言.
所以我們這次可以回頭看歷史比較早期.
第一批上去重建的時候先知的寓言.
這些經文表面上相當沉悶.
因為很多的寓言,尤其是第一到第七章.
其實你可以很簡單找到有八個的儀仗.
這些儀仗有些很難解的.
有些很天馬行空的,好像科幻顯.
但是你就可以發現到.
在這個時期,大概是主前的五百多年.
神都已經透過先知.
好像有些類似神密教材.
有些好像是叫他們有些幻想.
這樣來告訴他們.
其中一個只是一個重點.
就是耶路撒冷要有敬拜.
其實就是這麼多的.
有了敬拜之後,就是第八章鋪路.
第八到十四章.

$^{81}$就是講到猶太這個群體日後的光景.
所以很快到第九章就預言彌賽亞降臨.
所以這本書其實就是彌賽亞的一個先知書.
譬如九章,大家可能比較熟的.
有時有些牧者都會引述的.
九章第九節.
其實福音書有引述.
「錫安的民應當大大喜樂.
耶路撒冷的民應當歡呼.
看!你的王來到你這裡.
他是公義的,並且私行拯救.
謙謙和和的騎著驢.
就是騎著驢的驅使」.
哪裡出現過?.
騎驢進城嘛,耶穌.
就是直接引述這段經文.
他就是那位彌賽亞.
所以你可以想像得到.
一到第七章是有八個的儀仗.
在撒迦利亞書第八章開始.
就講到這個國度耶路撒冷聖殿.
將來是會怎樣的呢?.
某程度上第八到第十四章.
就是一個預言將來.
這個民族這個國家的光景.
當然是一個亮麗的光景.
又牽涉到和列國的關係.
所以先知有時都超越了地域.
不只是局限於耶路撒冷.
外面的國家的興衰.
都和他有關.
所謂和他有關.
是和耶路撒冷有關.
和神的選民有關.
所以又會有這樣的啟示.
大概有這個框.
所以有些解經家都說.
可能是兩個撒迦利亞.
一到第七章這樣的體裁.
第八到十四章又是另外一種的體裁.

$^{121}$會不會是兩個作者呢?.
這些就超越了我們大家.
這個階段的能力.
比較複雜的討論.
但無論如何在舊約的聖經.
撒迦利亞書是一卷書.
不是兩卷的.
我們都以往說過.
先知書沒有分所謂大先知,小先知.
大小先知其實是後人.
特別是七十事易本的分段.
將比較大的篇章.
譬如以賽亞書是66章.
就叫所謂大先知.
其實在猶太人他們的傳統.
其實就沒有大小之分.
只有一個分法.
就是前先知書和後先知書這樣分.
所以我們要有些印象.
所以後人.
我忘記了是否再後一點的人.
就會這樣理解.
所以我們中文翻譯的和合本.
很明顯就是大先知,小先知.
所以十二小先知的書卷.
其實就不是這樣分.
給大家有一個這樣的印象.
我們可以進入經文的世界.
記得當時是波斯帝國.
古列王准許天下各國的人.
因為他們吸納了巴比倫和亞述的高壓政策.
不得到民心.
所以當波斯帝國打敗列強之後.
他就採取懷柔政策.
准許所有被擄,亡國的國家的人.
可以回到自己的原居地.
來建造你們的神廟.
所以猶太耶路撒冷.
只不過是其中一個民族.
不過當然這個就是神早就預言了.

$^{161}$所以原來神都可以使用一個國家的政策.
來達成被擄歸回,重建聖殿.
七十年的刑罰戒滿.
你可以想像到.
連古列王在《二次華夏》書的四十五章.
是神的僕人.
他沒有想過現在他做的政策是要配合神的計劃.
他不會讀聖經.
他不會知道耶利米的預言.
他不會知道達爾利曾經這樣說過.
所以你可以想像到.
神的作為有時真的超越了地上的國度.
這就讓我們有少少立體.
對於當時的政治.
當時的國家的興衰.
對於整個重點是到底神的選民.
現在在神的救恩計劃裡面去到哪一點呢?.
現在是要被擄了.
誰都難不倒的.
所以當時末世的先知.
譬如耶利米,以西結.
是主催一定要投降巴比倫的.
你不投降.
你就是等於違背神的旨意.
所以耶利米的下場是很慘的.
是被眾叛親離的.
每個人都說他賣國.
但是神的旨意就是說.
在這個時候.
整個民族一定要被擄到巴比倫.
是不能回頭的.
但是到了被擄歸回的差不多時間.
譬如達爾利他在晚年的時候祈禱.
或者有被擄歸回的時間.
那些人就想去重建.
所以你看到哈該,撒迦利亞在這個時候出現.
就是他們停頓了不重建.
反而就被先知去責罰他們.
所以你就可以想像到.
原來不同階段的先知.

$^{201}$都要有很敏銳的觸覺.
就是神在這個時候.
對這個民族有什麼的應許.
有什麼的吩咐.
所以先知就要配合神的時間.
現在是責罰.
還是安慰.
還是鼓勵,加油.
先知就要按時間去說神要先知說的工作.
所以我們就要留心這些經文.
所以第一章頭六節.
就是舊事重提.
為什麼你們不能夠得到神的祝福呢?.
就是因為不順服.
現在在這個階段.
就是他們已經停頓了.
聖殿已經爛尾了.
十幾年了.
現在怎樣去從根基再去建造呢?.
一度就是六節.
「大理烏王第二年八月.
耶和華的說臨到亦多的孫子.
比利加的兒子先知撒迦利亞說.
耶和華曾向你們列祖大大的發怒」.
所以你要對以色列人說.
你就是先知撒迦利亞.
撒迦利亞的希伯來的意思.
這個字就是神眷顧的意思.
如果哈該.
希伯來文的意思就是關於節期.
所以哈該很多時候.
是在重要的節期說話.
好像出來做一個代言.
我們看到神就和撒迦利亞說.
所以第三節.
「你要對以色列人說.
萬君之耶和華如此說.
你們要轉向我.
我就轉向你們.
這是萬君之耶和華說的.

$^{241}$不要效法你們列祖.
從前的先知呼叫他們說.
萬君之耶和華如此說.
你們要回頭離開你們的惡道惡行.
他們卻不聽也不順從我.
這是耶和華說的.
你們的列祖在那裡呢.
那些先知能不能夠永遠存活嗎.
只是我的言語和律例.
就是我吩咐我僕人眾先知的.
豈不臨到你們列祖嗎.
他們就回頭說.
萬君之耶和華定義按我們的行動作為.
向我們怎樣行.
他已照樣行了」.
好像幫他們做一些歷史的功課.
過去為什麼這個民族.
被擄到這麼遠的巴比倫.
這個米索波大米的平原地方.
即是現在大概是伊拉克那一帶.
這樣來到整個國家沒有了.
70年沒有了.
現在回頭來那五萬多人.
其實是很少數.
回頭來重建.
但是面對一些困難的時候.
他們又開始手腳放遠了.
好像沒有了那個很強烈的.
類似猶太的復國主義.
不能夠一鼓作氣了.
因為那個氣已經被人打沉了.
現在神就興起先知.
一方面為什麼要這麼迂迴呢.
因為先知的系統.
至少在這個民族.
在猶太人.
其實當時的近東.
對於一些能夠講語言的人.
都會很尊崇的.
都會很敬重的.

$^{281}$比如非洲那些巫師.
能夠行法術的那些.
那些很落後的村子.
整個村子最重要的那個人.
就是那個巫師.
那些祭典就要那個巫師出來.
但是這個就是神已經在這個民族.
一直興起的積分.
所以猶太人他們始終.
對先知講的話.
他們都是會有尊重的.
他們都是會來聆聽的.
所以神就在這個時候.
就不是向所羅巴伯.
不是向大祭司約書亞.
來去顯現.
反而是興起兩位先知.
其實他們是什麼背景的人呢.
是不是跟隨上去重建的人呢.
聖經又沒有太多的交代.
只是告訴我們.
所羅巴伯就是在巴比倫.
被神的靈感動.
來到去上去重建的一個領袖.
是省長來的.
這個切巴薩.
所羅巴伯的另外一個名字就是切巴薩.
你就可以想像得到.
當時連一些領袖.
是當年回歸.
回到去.
都有一個很大的軟弱.
但是現在呢.
透過先知的安慰.
相信肯定是安慰.
因為你們轉向我.
我就一定會轉向你們.
這是一個很大的安慰.
就是說既往不咎.
不會再追究他們的問題.

$^{321}$在這個時候.
是一個安慰的時候.
在這個時候.
是一個要鼓動大家的心.
來看到重建聖殿.
是最重要的工作.
所以一切的其他的事務.
暫時都要分配時間.
和優先.
哪些是優次.
那個次序就要重新排列.
所以先知就扮演著這個.
這個時候就不是警告.
是提出過去的歷史.
我們的祖宗.
受過的遭遇.
但是現在不應該再重蹈覆轍.
現在是新一代.
而且這是一個重建的時候.
所以神就透過先知.
來發出安慰的說話.
其中第六節的重點就是.
人物全非.
就是人面全非.
但是物都沒有了.
但是有的仍然停留的.
就是有神的神諭.
即是先知的預言.
和先知曾經在適當的時候.
說適當的話.
那個重要.
你過了之後.
馬後炮 事後孔明.
是人都會說的.
但是先知最重要的身份.
那個價值.
就是在適當的時候.
未發生之前.
提出那個.
你可以說是危言聳聽.

$^{361}$其實重要的是.
將神要說的話.
重要的說出來.
這個就是先知最重要的.
因為最重要的警告.
你不出聲.
就不知道.
怕被人說.
怕受牽連.
怕得失人.
那你就很容易做先知了.
但是最重要的時候.
你能夠說了.
聽不聽就不關你的事了.
就是他自己面對神了.
這個時候.
當那些事情一直演變.
原來真的是.
真的是這樣發生.
所以耶利米後來.
被很多猶太人景仰他.
因為他在最重要的時候.
他不計較自己的名聲.
不計較自己生命的安危.
都要來勸百姓放手.
投降.
去到巴比倫.
做買賣.
蓋房子.
生孩子.
準備等候70年的回歸.
他甚至乎自己都買了一些地.
你可以想像得到.
簡直是難以想像.
那些事情還沒發生.
耶路撒冷的先知.
那些就是賊孩先知.
因為當時西底加王.
是非常之極力.
想展示信心.

$^{401}$我們怎會輸給巴比倫.
我們的國家有神.
那麼強大.
怎會輸給巴比倫呢.
所以他當然希望身邊的先知.
告訴他們.
我們一定打勝仗的.
神一定會站在我們這一邊的.
我們不需要懼怕的.
雖然巴比倫的軍隊.
尼布薩拉旦已經在冰林城下.
但是不需要怕的.
你想想有兩批.
一批就是站在西底加王身邊的.
但是一些是很孤單的.
被王遺棄的.
將他放在一個枯井裡面.
任由他沒得吃.
沒得喝水.
當然有一些領袖去接濟耶利米.
在耶利米書裡面充斥著說.
你可以想像得到.
先知的說話.
在最重要適當的時候.
他必須要說神要透過他所說的話.
所以在這一段.
這個序言.
就是神透過撒迦利亞.
當然還有哈戈.
哈戈就特別強調.
要做了.
特別強調那個行動.
真的這樣上山取木料.
真的去強調輕功建造.
但是撒迦利亞.
似乎是一個精神領袖多一些.
是說一些很抽象的預言.
但是整個重點.
都是說耶路撒冷將要發生的事.
所以這一大段.

$^{441}$我們先來一個小小歷史的回顧.
過去神藉著先知說話.
列祖聽的就會得到教導.
但是大部分不聽.
就會得到報應.
神的刑罰.
審判.
所以現在什麼都沒有了.
人都過去了.
因為是歷史.
過去了.
那些先祖都過去了.
過去的先知都過去了.
耶利米都死了.
都已經在主耶穌的懷裡了.
但是現在去到這個時代.
是神要復興以色列.
猶太選民被擄歸回.
可能第三代.
甚至可能第一代.
都還有些老人家.
小時候被擄的.
現在去到七八十歲都有的.
哈戈書第二章就是這樣說.
所以你可以想像得到.
在這樣的離愁別緒底下.
現在能夠在聖地耶路撒冷.
根本就是自己一起去做.
這個聖殿的重建.
但是需要的可能就是.
這一份的挑戰.
這一份的安慰.
所以先知要拿捏那個時代.
百姓的需要是很重要的.
要拿捏到他們現在.
好像有一種屬靈的把脈.
現在這個時代.
百姓面對著什麼呢.
那個困難來自哪裡呢.
是否能夠克服呢.

$^{481}$還有是否有神的時間呢.
先知就是很重要的時間.
神的時間到了.
所以先知就會站出來.
好像說一些還沒有發生.
但是他好像已經見到那個結果了.
所以應該先知是第一個.
捉夢的人.
建築捉夢的人.
來給予一個遠景.
那個群體.
向他們預言的那幫人.
所以這裡告訴我們.
現在他們是可以得到神的迴轉.
只需要他們一幫人來迴轉.
這個迴轉相信肯定就是.
一起要建造聖殿.
在這個時候如果有人說.
不是呀不可以做呀.
神沒有要我們這樣做.
這些肯定是神所禁止的.
因為這個時候.
無論被擄的人.
已經實實在在站在耶路撒冷.
只不過是當地人的阻撓.
這些阻撓原來是可以排除的.
那個阻力絕對不夠神的幫助來得大.
所以有神在我們這一邊.
這個真是最重要.
所以如果我們還沒有看下面的異象之前.
我們在這一段裡面有什麼教訓呢.
我們要把脈一個時代.
我們看到這個時代現在是在什麼處境.
還有我們一個群體應該怎樣去.
在神永恆不變的福音裡面.
福音是一個很大的能力.
可以改變生命.
我們怎樣能夠踏上順風車.
將福音來給人呢.
這個一定是對的.

$^{521}$永恆是對的.
但是不同的時代.
不同的群體.
對福音都有不同的身心靈的阻隔.
怎樣排除這些障礙呢.
這個就是每一個時代每一個教會.
都必須來小心把脈.
小心去幫它分解它們的困擾.
每個人的困擾都不同的.
中環的人的困擾.
和在天水圍的街坊.
一些的主婦的困擾都不同的.
你要將一個福音.
永恆不變的耶穌.
怎樣可以傳遞得準確.
在這些人的困擾心靈裡面呢.
這個就是那個牧者.
那個先知.
那個傳道者.
就要很準確的拿捏.
和那個群體怎樣去承載.
福音說了.
怎樣承載那些人內心的困擾呢.
就是那個後續.
怎樣能夠配合到福音工作.
來做那個繼續的播種.
可能就是消水.
這些就是需要不斷的把脈.
更加重要就是.
我們要留意到.
到底現在這個整個大時代.
福音怎樣來向這些.
未信的朋友來傳遞呢.
所以這一段聖經.
是令到百姓.
配合哈該先知.
不是彼此衝突.
是彼此的調和.
帶出一個很和諧的訊息.
就是要做了.

$^{561}$哈該就和他們一起.
落手落腳去做了.
真是很實在的.
去取木料去畫圖跡.
很實在的一齊去建造.
但是撒迦利亞.
相信是更多.
在那個精神上.
比他們有一些更重要的寓言.
所以這裡有一個序言.
一到第六節.
就是撒迦利亞書的一個.
先是一個序言.
接著就會突然間有八個義將.
一併這樣說出來.
去到第七章.
所以現在我們就逐個義將去看了.
每一個義將其實.
多或少都要告訴我們.
有一位非比尋常的人.
其實我們新約的人就很幽福.
新約的人一點都不覺得是神秘.
就是耶穌.
但是你覺得很奇怪.
為什麼當時明明是.
實實在在要建造聖殿.
重建.
接著有自己的憲制.
要有自己的城牆.
這些我們都讀過了.
尼希米基我們都讀過了.
但是為什麼在這個時候.
撒迦利亞這麼不現實.
需要人手幫助.
為什麼這麼不現實.
去說一些將來很遠的.
五百多年後.
現在是主前說這個寓言的人.
撒迦利亞是主前的.
四百多年.

$^{601}$還沒到五百多年.
也不是.
五百四十多年.
也有五百多年後.
耶穌才出現在伯尼行.
這麼快就說了.
好像神也巴不得.
他現在看到的是一個爛尾的聖殿.
但是神看到的不是爛尾的聖殿.
神看到的是五百多年後.
神根本就沒有時間.
連撒迦利亞都看到.
他所謂的看到.
就是很模糊的.
我們就很清楚了.
他看到的是五百多年後.
有一個人.
騎著白馬的不知道是什麼人.
很特別的.
每個騎馬的都要向他下拜.
我們就發現.
原來舊約的寓言.
都挺好看的.
都挺特別的.
第七節開始.
大利烏第二年的十一月.
又過了.
有時間的.
有神秘的時間.
好像有個時間表.
好像告訴你.
是哪段時間我講的道.
就是西巴特月.
二十四日.
十一月的二十四日.
耶和華的華臨到易多的孫子.
比利加的兒子.
先知撒迦利亞說.
見一人騎著熊馬.
站在.

$^{641}$蛙地.
番石樓樹中間.
在他身後又有熊馬.
黃馬.
和白馬.
我對與我說話的天使說.
明顯是異象.
因為撒迦利亞.
是見到天使的.
但是.
他見到的那個情景.
是一個番石樓樹的地方.
一個蛙地.
就說主啊.
這是什麼意思呢?.
他說我要指示你.
這是什麼意思?.
第十節.
那站在番石樓樹中間的人說.
這是奉耶和華差遣.
在遍地走來走去的.
那些騎馬的.
對站在番石樓樹中間.
耶和華的使者說.
我們已在遍地走來走去.
見全地.
都安息平靜.
我們先有這個圖畫.
異象其中一個的.
作用.
就好像畫了一個圖畫.
給我們.
所以啟示錄.
這是我很深刻的.
就是.
我忘記了.
什麼博士?.
中神的老師.
他說啟示錄.
其中一點.

$^{681}$很重要.
啟示錄是天啟文學.
有很多天上的圖像.
透過好像一幅圖畫.
讓我們見到.
所以.
現在撒加尼亞書.
他用的.
體材.
都有天啟文學的含義.
在.
其實.
兩若中間.
就會有很多.
天啟文學的.
作品出現.
所以撒加尼亞.
雖然未到那個時間.
那麼後期.
但在他這個時候.
都開始有天啟文學的.
這些體材.
出現在聖經裡面.
所以民間應該更加多.
但聖經.
有一個指向.
一定是說將來末後.
尤其是那位彌賽亞.
有相關連.
所以我們就有這個.
類似是詩情畫意.
類似是有一種.
好像欣賞藝術.
看天啟文學的.
這些敘述.
十二節.
於是耶和華的使者說.
萬君知耶和華.
李魯恨耶路撒冷.
和猶大的成音.

$^{721}$已經七十年.
舊事重提.
被擄的原因.
就是神的憤怒.
那麼現在這個時代.
是不是還要懼怕神的憤怒呢.
似乎不是.
似乎是歸向.
轉回.
就是說現在不再是憤怒的神.
現在是安慰.
和施恩的.
這一位耶和華.
所以你看到下面.
他就會這樣說.
你不施憐憫要到何時.
第十三節.
耶和華就用美善的安慰.
回答那與我.
說話的天使.
這裡有好幾位天使.
所以很混亂.
到底誰才是真正的.
真命天子呢.
第十四節.
與我說話的天使對我說.
你要宣告說.
萬君知耶和華如此說.
我為耶路撒冷.
為石安.
其實只是一樣的稱呼.
石安城石安山.
就是耶路撒冷.
所以就是雙重的稱呼.
特別的親切.
感受.
如果再說元朗.
感受.
就是更加凸顯了感受的重要.
元朗有好多地方.

$^{761}$新田也是.
竹苑也是.
元朗的感受.
就是石安的耶路撒冷.
是心裡火熱.
神心裡.
不是憤怒.
這次是火熱.
神的心裡火熱.
神自己說.
與天使說.
我為耶路撒冷.
心裡火熱.
是極其的火熱.
我甚擄擄那安逸的烈國.
因我從前稍微擄擄我民.
他們就加害過分.
他現在的擄擄.
即是憤怒.
是向烈國.
是向那些安逸.
即是以為扮演著.
可以將耶路撒冷猶太人.
趕走的烈國的人.
這些外邦人.
拜假神的.
以為可以替神.
去懲罰了.
這個.
以色列猶太民族.
神說沒錯.
那七十年是的.
神有一個.
稍微的擄行.
對猶太人.
所以時代的把脈.
很重要.
這個時候.
撒迦利亞的時候.
就是神心裡火熱.

$^{801}$為誰火熱?.
為耶路撒冷.
為猶太人.
為這個聖殿的重建.
為這個地方有憲制.
將來五百多年後的彌賽亞.
按照著神的.
應許降臨.
這個加勒頓大書四章.
七節應該是.
及此時候滿足.
神就差遣他的兒子.
為女子.
所生.
且生在.
耶路撒冷.
降臨.
來拯救這個民族.
和世界上的人.
所以這裡告訴我們.
神現在要怎樣?第十六節.
所以耶和華如此說.
然今我回到.
耶路撒冷,仍思憐憫.
我的殿必重建.
在其中,準繩.
必拉在耶路撒冷之上.
就是那些.
類似那些蜜豆.
以前.
建房子都是靠這些.
那時候沒有那些.
高水平的電子.
的量度.
現在用電子.
就可以很直.
以前是靠繩子彈一些蜜.
我都玩過.
小時候錦田建房子的時候.
見到那些工人.

$^{841}$彈那些蜜.
挺有趣的,就很直.
沿著那條線.
建那個磚牆等等.
這些一般建房子都是這樣.
所以這裡說.
必拉準繩在耶路撒冷之上.
這是萬軍之耶和華說的.
你要再宣告說.
萬軍之耶和華如此說.
我的誠邑.
必在豐盛發達.
耶和華.
必在安慰石安.
簡選耶路撒冷.
你想想.
簡直是lum 到爆炸.
不斷說神要安慰.
神心裡火熱.
神要來拉準繩.
在耶路撒冷.
這個時候.
他們勉為其難.
開始半信半疑.
停頓了.
爛了沒十幾年.
你現在要我們再重建.
誰說的.
哈該說神給我有說話.
撒迦利亞說.
神給我有這個預言.
甚至有異象.
那個騎紅馬.
他好像有些小孩.
其實是天使.
在這裡被吩咐.
在全地走來走去.
看到全地都安靜.
撒冷.
然後再詳細解解.

$^{881}$第一個異象.
就是馬的異象.
解釋就很清楚.
神透過先知.
已經看清楚了.
四處都看到.
沒事.
很平安.
很安靜.
現在是一個可以重建的時候.
所以不要小看這些預言.
對百姓來說.
是一個很大的安慰.
和支持.
如果沒有神站在他們那邊.
多麼賣力都做不到.
但是如果有神的安慰.
有神的應許.
多難都會做得成功.
而且那個過程.
更重要就是.
神的話語成為.
整個群眾.
整個群體的凝聚的原因.
如果不是.
我聽誰說.
你是領袖.
你講這一句.
你就是領袖.
這才是最重要.
所以我們看到這個異象.
就是神透過.
撒迦利亞.
他在晚上就看到.
這個人騎著紅馬.
站在.
那個窪地.
可能是濕地.
翻石榴樹中間.
翻石榴在猶太人.

$^{921}$他們有些意思.
多種子.
就是.
生命的復興.
可以想像得到.
猶太人會這樣想.
我們想的就是.
翻石榴.
有什麼特別.
但是猶太人他們會理解.
這是一個復興的.
一個預表.
接著在他身後.
有紅馬黃馬黃白馬.
也就是說不只是那些.
人騎的馬.
一定是有人騎.
這裡所說的人其實是天使.
因為第九節.
我對與我說話的天使說.
主啊這是什麼意思呢.
那位天使就說.
我要指示你這是什麼意思.
第十節就開始說.
那些騎馬的.
在異象裡面.
騎馬的開始有些動作.
當然那個重點.
是放在騎紅馬的.
翻石榴樹中間的那一位.
那站在翻石榴樹中間的人說.
這是奉耶和華差遣.
在遍地走來走去的.
那些騎馬的.
對站在.
所以真的有人騎.
紅馬黃馬白馬.
都有人騎的.
對站在翻石榴樹中間.
耶和華的使者.

$^{961}$現在就再詳細一點.
那個騎紅馬.
站在翻石榴樹中間的這一位.
原來是耶和華的使者.
所以就來到去說.
是那些騎馬的.
當然我們看到的.
在異象裡面.
就是三隻馬.
是不是真的這麼少呢.
我不知道.
這個時候被吩咐.
到處走的這幾個.
原來是天使來的.
又匯報給這個騎紅馬的.
在翻石榴樹中間的.
耶和華的使者知道.
所以你看到.
已經有不同階級.
騎紅嗎哪個.
原來是高級的.
原來是耶和華的使者.
但是其他的似乎是天使.
我們在遍地走來走去.
見全地都安靜.
安息平靜.
好像向他匯報.
我們已經巡過一遍.
這樣的意思.
撒迦利亞就很有福.
他就看到這個異象.
所以先知在舊約.
特別在撒母耳記開始.
撒母耳記常已經開始.
先知其實叫先見.
叫seer.
即是看到.
其實先見是.
比別人見得更早.
先知就是.

$^{1001}$沒錯.
比其他人再早一點知.
因為先知是.
先見的.
先見的意思.
是先見的.
因為先見是.
你再去報那個訊息.
所以這裡就告訴我們.
撒迦利亞就看到.
有一個耶和華的使者出現.
是尊貴的.
非比尋常的.
因為其他騎馬的.
是向他匯報.
他們走來走去的.
那個結果.
耶路撒冷一定是那個中心焦點.
耶路撒冷現在.
安息平靜.
接著.
於是.
耶和華的使者說.
似乎是剛才騎紅嗎哪個.
耶和華的使者.
除非有另外一個使者.
萬君之耶和華就說.
你老恨耶路撒冷.
和猶大的成音.
你不思憐憫.
要到何時.
這是一個祈禱.
這個使者向神.
發出呼籲.
好像撒迦利亞.
看到一個舞台劇.
首先出場的.
就是有一個報警.
一個翻石樓樹的中間.
所以大家真的要.

$^{1041}$我從來沒有看過舞台劇.
人們說很好看.
我就不懂欣賞.
當我們幻想.
現在我們就去到那個舞台劇.
我們坐在下面.
看到報警了.
就看到那個騎紅馬的.
在翻石樓樹中間.
原來他在指點.
一些其他騎馬的.
有黃色馬.
我們沒有見過什麼叫黃色馬.
就不是應該蓋布.
黃色馬會那些就是了.
那些賽馬就是了.
有不同的騎師有不同的制服.
黃色馬.
又有白色馬我們見過.
紅馬我們都見過.
中式的.
當時通常都有.
鮮紅的就不會.
但是在異象裡面.
撒迦利亞就見到了.
騎紅馬的.
一直都站在翻石樓樹中間的.
那位使者.
就一直等待.
跟著那些去周圍走的.
部下.
就騎著馬回來了.
就報告說現在.
我們視察過.
都是安息平靜的.
跟著轉那個畫面.
就是這位耶和華的使者.
突然之間走到來.
就大聲.
不知道什麼形態.

$^{1081}$所以舞台劇就很鮮明.
可能仰天長嘆.
向上天祈禱.
你到什麼時候當然不是一種質問.
是一種呼求.
就是呼求神.
怎樣已經70年了.
你老恨耶路撒冷.
似乎為耶路撒冷祈禱.
這個耶和華的使者.
所以你可以想像到.
主耶穌在臨離開世界之前.
跟大家揭一揭.
應該是《馬太福音》.
23章.
我要找一找.
主耶穌在這裡.
如果大家可以幫到我的忙.
想一想在哪裡.
耶穌基督.
祂在這裡向神祈禱.
呼喊.
耶路撒冷啊.
你嘗殺害.
才知.
由.
打死那些來到.
我暫時找不到.
否則今晚給大家.
如果大家搜尋到.
應該比我快.
22章.
對不起.
23章的37節.
我的記性都差不多了.
沒錯了.
主耶穌在這裡明顯.
23章是說什麼的.
就是.
法理賽人的七禍.

$^{1121}$文氏法理賽人的七禍.
接著27節.
就是耶路撒冷.
對不起.
耶路撒冷啊.
你嘗殺害才知.
又用石頭打死.
那奉差遣到你這裡來的人.
我多次願意聚集.
你的兒女.
好像母雞把小雞聚集在.
翅膀底下.
所以你有時見過一些故事.
那些地方火燭.
那些母雞燒死.
接著拿開.
死了的母雞.
那些小雞就在母雞下面.
遮蓋著它.
其實地震也是.
那些媽媽都怎樣壓死.
都要保護自己的小朋友.
這裡就告訴我們.
主耶穌好像是母雞.
保護祂的子民.
指示你們不願意.
看那你們的家成為.
方牆樓給你們.
這個就是先知的悲哀.
自己悲哀.
因為看到那個現實.
一定會遭遇神的懲罰.
所以以色列.
猶太.
70年.
羅馬人.
來毀壞聖殿.
他們就會被擄到天下.
即是分散到天下.
你可以想像得到.

$^{1161}$這個就是耶穌的呼喊.
祂差不多就快要.
在這個受難州裡面.
祂要入城.
你可以想像得到.
主耶穌就是這樣的一個祈禱.
所以原來.
在舊約.
在這些寓言裡面.
都已經有耶和華的使者.
向神呼喊.
不過這個就不是一個悲哀的呼喊.
而是好像要喚醒.
神的恩典.
原來這個時候.
耶和華使者的呼喊.
原來這個時候.
就是求神制止.
其實不需要.
因為這個時候.
所以是一個非常之.
復興的時候.
要有一個興高采烈的心.
為什麼呢?.
因為神為他自己的.
這個地方耶路撒冷.
他要大發熱心.
你就可以想像得到.
如果神重新眷顧這個民族.
這個恩典.
只會是加倍的.
湧流到.
這個民族當中.
所以十三節.
耶和華就用.
美善的安慰.
對這個耶和華的使者.
對這個.
站在.
番石樓樹中間.

$^{1201}$騎紅馬的人.
就是撒迦利亞.
報警又轉另外一幕.
然後天上.
用安慰美善的.
說話安慰.
這個使者.
這個使者是為耶路撒冷求.
平安的.
所以你有時留意.
聖經福音書.
主耶穌那種仰天長嘆.
其實你.
靈修.
你會很感受到耶穌的愛.
對這個民族.
在這個民族七十年後.
主後七十年.
主耶穌的時候.
就是主後的大概三十多年.
現在主後七十年.
即是還有三十多年.
整個耶路撒冷就被羅馬人.
攻陷了.
大家去到那個馬薩大.
就是那些羅馬的軍隊.
提多將軍.
因為他們在山上.
他們就抓那些土裡.
一個螺旋形的上山.
上到去最終他們都是要.
自盡.
不可以給外邦人殺自己.
情願大家就被人來殺.
給自己的同胞去殺.
那你可以想像到.
耶穌是遇見到.
三十多年後的耶路撒冷.
血流成河.
所以他來到去.

$^{1241}$呼喊耶路撒冷.
耶路撒冷.
你嘗殺害先知.
又打死那些奉差遣而來的.
你們都不肯.
願意悔改.
耶穌說我好似母雞.
來人肉.
但你們都不願意.
這個使者.
在這個時候就不是一種.
呼喊一種悲哀.
是一種為這個.
民族去求更.
大的恩典.
所以神用美善安慰的說話.
回答.
那與我說話的天使.
接著又說天使.
很特別.
當然這是我.
或我們新約的人.
這樣理解.
可能真是天使.
可能未必是主耶穌.
但你可以想像到.
這個人是有些特別的.
比起其他使者.
是有特別不同的.
如果我們按照著.
三一論.
主耶穌的職份就是.
在救恩上.
祂成為聖子.
來到這個世界.
所以祂在地上.
祂就有為人求平安的職份.
最終就是在.
十誡上為人類求平安.
就是祂的死.

$^{1281}$祂的流血.
將人類的罪拯救.
所以.
如果我理解這位耶和的使者.
這位天使是主耶穌.
其實都不是太過分.
因為祂是.
特別的使者.
但更重要的是.
神的宣告.
我為耶路撒冷.
為石安.
心裡極其火熱.
不再擄擄民眾.
而且.
祂是極其火熱.
即是說.
像李家超特首.
什麼都按結果行先.
石澳.
我的同學在石澳堂.
現在還在淋雨.
回不去怎麼辦呢.
上個星期就很緊張.
要有崇拜.
李家超說.
九點前一定要通.
鐵律.
務求一定要通車.
類似這樣.
神現在是.
心裡火熱.
為民族.
為耶路撒冷.
心裡火熱.
這時候是興功建造的時候.
不要再爛尾的時候.
所以你可以想像到.
這本書其實充滿了盼望.
撒迦利亞書.

$^{1321}$很少人讀.
因為舊約我們好像看不上眼.
變了我們很多時候.
很少透過這些異象.
去看神.
在那些事情還沒發生.
祂已經有這樣的心情.
我們的造物主.
和祂的民族.
和猶太人.
因為他們的關係已經幾千年了.
所以可以說一些很深的說話.
很窩心的說話.
這個民族受了這麼多的.
蹂躪.
七十年.
很慘的.
沒有了自己的國家.
很悲哀.
但是現在在這個時候.
重建的時候.
更加透過兩位先知.
去講述一些安慰的說話.
使到神的選民.
得到建立.
所以你可以想像到.
對於當時重建的人.
特別是所羅巴伯.
大祭司約書亞.
第四章又會再出現約書亞.
特別是第四章的異象.
是提到.
在聖殿裡面.
很特別的就是.
源源不絕的.
有一些油燈.
需要油.
我的燈需要油.
就是由那裡的靈感來的.
就是在聖殿的一個異象.

$^{1361}$就是聖殿.
怎會有一棵橄欖樹呢?.
橄欖樹的油.
不斷流向那些燈芯.
即是說.
燃燒的那些油.
即是說.
源源不絕的那個供應.
在那個聖殿生生不息.
所以你想想.
整個撒迦利亞書是充滿了盼望的.
所以其實.
這些先知書.
雖然我們從來都不會揭.
但其實讀下.
都讓我們回頭看回那些事情未發生.
耶穌未降生.
但都讓我們能夠看到.
原來神已經有祂的計劃.
原來神已經有祂的心意.
尤其是我們透過這些先知書.
更加可以觸摸到.
神的心腸.
不是鐵石心腸.
是祂不同的情況.
祂就會有不同的心情.
如果我們讀耶利米書.
你會感受到一位憤怒的神.
對自己的民族.
子民的定義的刑罰.
那個是.
無法解釋的.
那個定義的刑罰.
那個是.
沒有人可以再祈禱.
神改變他的心意.
但來到這裡.
神現在要.
祝福這個民族.
這個祝福都不是一種.

$^{1401}$說一下就算的安慰.
而是用心裡火焰.
這樣來講述.
祂的內心的心情.
當然更重要就是.
因為那個救恩.
我們真的需要有.
這個民族.
慢慢人肉彌賽亞的出現.
神對人類.
的永恆救贖.
所以現在過了.
二千年.
基督教的走向是怎樣呢.
現在伊斯蘭世界的.
生育率.
佔據.
是越來越多了.
如果基督教一直都不再生BB.
基督教裡面的弟兄姊妹.
不是說基督教.
越來越少.
歐美就少了.
現在連華人都少了.
如果是這樣的時候.
將來再過多一百年後.
伊斯蘭的人口.
會越來越.
統治整個地球.
這個都是一個.
大家都要正視的問題.
如果沒有了敬虔的後裔.
整個世界將來會怎樣呢.
神的計劃.
是由下一代世代的人去想.
現在我們不是這樣的時候.
但是我們是看到.
那個趨勢的發生.
所以我們其實都要求主.
讓我們能夠.

$^{1441}$擔心.
為下一代去禱告.
他們不單純是.
買不買到房子.
工作越來越艱難.
而是在說那個作見證的時機.
那個更加來得迫切.
十五節.
剛才我們簡單讀過.
神曾經.
擄奴耶路撒冷.
其實是.
稍微的擄行.
但是現在.
我甚擄奴那些安逸的列國.
就是說.
幸災樂禍.
落井下石.
看到這些民族受到神的刑罰.
就以為.
做物主只是泥菩薩.
幫不了這個民族.
反過來我們的神.
我們所拜的神.
反而令到我們很好.
就等於.
譬如有個好朋友.
他拜那些.
很迷信於某些神.
隨便一個.
有些人拜那些很厲害的.
四面佛.
很成功的.
生意做到很大.
但是你信耶穌.
你就有很多困難.
又有不同的病.
就這樣看.
這樣比較.
他就真的是安逸的列國.

$^{1481}$你的神都不幫你的.
你的神好像令到你.
這麼多的困難.
這個就是安逸.
和表面上好像神都不祝福.
但是.
當然新約我們就不是.
這麼簡單的.
從那個結果來評論.
因為.
心靈的問題.
所以基督徒有另外一個.
很反映到信仰的真實.
就是.
就算那個環境.
表面是很困難.
但是基督徒內心.
深處.
是有真正的平安.
是有神比的盼望.
因為真正的永恆.
才是那個.
獎賞的地方.
暫時.
你多麼的風山水起.
成功.
當然我們不是吃不到的葡萄.
我們就用一種這樣的心態.
而是.
他有沒有真正與神.
結立的平安呢.
有沒有真正的永恆的盼望呢.
當那個生命去到結束.
他所擁有的東西.
能不能夠帶走呢.
這些就是我們作為信主的人.
一定要去想.
所以這裡告訴我們.
現在我們站在神這邊的民族.
神現在為這個民族.

$^{1521}$心裡火熱.
神現在真正的老奴.
是那些安逸的列國.
就是說神要施行.
那個審判.
在這些其他的國家.
的角度上.
例如在.
摩尼哈該書或者.
以色列記.
都出現過的參巴拉.
多比亞.
這些領袖.
這些阿拉伯人基善.
這些人都是代表著.
他們那個民族.
他們拜的假神.
所以現在.
你會想像到神.
為自己的子民.
來到申冤的時候.
不是局部性.
不是真正永恆的申冤.
真正永恆的申冤是要去到.
新天新地.
是要去到那個白色的大寶座.
所以大家要記住.
在今生我們不要有.
一種過分的樂觀.
就是.
在今生.
必定會有一個絕對公平的.
審判,是沒有的.
真是.
我們要發洩耐心.
去向神效忠.
因為最終.
都是神去定奪.
誰是對誰是錯.
又想起羅馬書的.

$^{1561}$十四章.
又和大家揭一揭.
這些都是重要的經文.
保羅這樣講的時候.
其實他都是語重心長的提醒.
我們不要隨便去.
小看其他人.
不可以論斷其他人.
十四章的第一節.
是沒有什麼對錯.
關係的事情.
就是文化.
倫理的一些事情.
你信心軟弱的.
你們要接納他.
但是不要辯論所疑惑的事.
接著他就舉例子.
有人信百物都可吃.
就是豬紅又吃得.
什麼都吃得.
但是有些猶太人是不行的.
有人信百物都可吃.
但是軟弱的.
吃不吃蔬菜.
吃的人就不可以輕看不吃的人.
不吃的人也不可以論斷.
所以媽生兄弟.
看不起人.
人家又會論斷你.
接著你又會論斷他.
這樣互相排斥.
所以你可以想像到.
保羅用的字眼.
很將人性講出來.
你做到的.
你又看不起其他人.
人家不敢這樣做.
但是就會論斷你做的人.
保羅說千萬不可以這樣.
不可以論斷吃的人.

$^{1601}$因為神已經收納他了.
第四節很重.
你是誰.
竟論斷別人的僕人.
他或站住或跌倒.
自有他的主人在.
所以不是論斷.
是善意的.
大家的溝通.
和打開心靈的彼此接納.
每個人都有不同的自己.
每個人都有不同的.
自己的文化.
背景.
每個人都有不同.
如果不關乎真理.
其實互相欣賞.
互相接納.
互相溝通.
重要的是彼此能夠有.
在主裡面的一種溝通.
彼此來欣賞.
你是誰.
竟論斷別人的僕人.
他或站住或跌倒.
而且他也必要站住.
因為主能夠使他站住.
即是說.
他也要為自己.
隨隨便便地吃東西.
來向神負責.
不是要向你負責.
所以你就可以想像到.
這樣就帶來一個很通達.
大家都會彼此.
有一個空間.
互相會拉近大家的距離.
如果不吃的人.
原來後來就沒有關乎真理.
原來是可以的.

$^{1641}$他信了主之後.
原來有些事情.
是不需要再執著.
他慢慢消除了內心的罪疚感.
軟弱.
他可以慢慢成為.
有被接納的一個群體.
的一個成員.
就不會永遠是一個對立面.
所以保羅在這裡.
是關乎到教會的模樣.
即是說.
要有這樣的胸襟.
這樣的一個觸覺.
對人.
不是你我行我素.
我是這樣的.
我這樣做沒問題.
都要想想別人的感受.
第六節.
又說另外一個例子.
有人看這日比那日強.
有人看日日都一樣.
只是各人心裡.
要意見堅定.
保羅是守安息日.
猶太人.
去到星期五的黃昏.
墳外感覺自己神聖.
我們.
希望文化慢慢都會改變.
星期日的早上.
你令到你的整間屋.
如果你是屋主.
令到整間屋.
我們準備好.
盡快吃完早餐.
大家穿好一點.
準備一起回教會.
希望是我們家庭的文化.

$^{1681}$這樣的時候.
你就會感覺到聖日的重要.
猶太人就.
經年累月.
世世代代就安息日.
由亞伯拉罕開始.
由摩西定了.
第四日是安息日.
第四屆是安息日.
就已經是.
後代一直遵守.
所以保羅.
你可以想像一下保羅.
他對信仰的透徹的理解.
他不會再叫.
其他人堅持.
而他自己.
很虔誠地.
如果去到安息日.
他自己都會進去會堂.
去講道.
他都不做事.
你可以想像一下保羅.
現在他就教好了.
有人看這日比那日強.
有人看日日都是一樣.
只是各人心裡要意見堅定.
首日.
安息日.
人要為主守.
吃的人要為主吃.
不吃的人.
原來大家都是神的兒女.
原來大家.
做和不做都不是那麼重要.
關係才是最重要.
和神的關係.
和人的關係.
才是最重要.
所以保羅就告訴我們.

$^{1721}$我們沒有一個人是為自己活的.
信了主耶穌之後.
因為十二章已經將身體當作活祭獻給神.
我們沒有一個人為自己活.
也沒有一個人為自己死.
我們若活著.
是為主而活.
若死了.
是為主而死.
所以我們或活或死.
總是主的人.
你可以想像到.
保羅講關於論斷.
他都有一個那麼大的圖畫.
用一些神的教材.
他自己是國內人.
這樣去苦口婆心.
去勸勉教會.
一定是來自五湖四海.
一定是來自不同文化背景.
一定是來自不同的國家民族.
但是有一個更大的身份.
改變了.
就是大家是神永恆.
國度的兒女.
所以我們要帶著.
這些神的教材.
去勸勉教會.
所以我們要帶著.
這樣的身份.
要很強烈.
有這個意識是神兒女.
這樣回到.
所以有時弟兄姊妹.
回來做私事.
原來他們平日.
是很重要的.
CEO.
當然你會覺得.
大材小用.

$^{1761}$這是另一個話題.
因此可以再用得好一點.
但是如果.
他以神兒女的身份.
我現在星期日.
來到神的家.
在教會裡面.
牧者安排我們做這個私事.
派情罪表.
我不單止不會覺得.
是卑下.
我覺得我能夠服侍神.
服侍弟兄姊妹.
就是很尊貴.
這樣想就對了.
每個弟兄姊妹都是這樣想.
大家都是神的兒女.
來一起服侍.
所以這裡告訴我們.
在這裡.
神要來施憐憫.
對這些安逸的列國.
來施以懲罰.
所以我們求神.
讓我們能夠.
把握著神的心腸.
透過這一段的寓言.
這個意象.
本來就是主角.
是耶和華的使者.
騎著紅馬.
在這個翻石樓地的.
這個使者.
但問題就是.
他向神呼天呼喊.
說你何時安慰耶路撒冷呢?.
神就在天上回答了.
所以如果這個舞台劇.
撒迦利亞看到的.
就是有一把.

$^{1801}$很雄厚的聲音.
向這個使者說.
不是向撒迦利亞.
向耶路撒冷說.
不是向撒迦利亞.
是向這個使者的回覆.
我為耶路撒冷心裡火熱.
現在我會將列國.
會將列國.
來懲罰.
他會擄勞列國.
使神的工程.
能夠恢復.
所以就是第十六節.
所以耶和華如此說.
現今我回到耶路撒冷.
凝思憐憫.
原來很重要.
所羅巴巴.
約書亞他們上去重建.
更重要的.
不是帶那些錢.
金銀財寶.
有的,有那些工人.
有那些瀑被.
有那些牛奴.
有那些在聖殿.
巴比倫王.
擄掠了的那些.
器皿.
他們帶上去.
但原來更重要.
神同在耶路撒冷.
這是最重要的.
如果沒有神同在.
聖殿只是一個.
金碧輝煌的空洞的地方.
有神同在.
就有很重要的.
神聖的工作.

$^{1841}$所以這裡告訴我們.
神說.
我的殿必重建在其中.
準繩必拉在耶路撒冷之上.
這是萬軍之耶和華說的.
第一章.
如果大家回去翻一翻.
數一數萬軍之耶和華.
在哈該書和撒迦利亞書.
是出現很多次的.
比其他的先知書.
還要多.
特別強調這個身份.
萬軍.
即是說.
他是.
我回去查英文.
和那個.
試一下查一下那個原文.
第十七節.
你要再宣告.
萬軍之耶和華如此說.
我的誠邑.
必在豐盛發達.
耶和華必在安慰.
錫安揀選耶路撒冷.
第一章就.
停在這裡.
就是第一個意象.
就是先是翻石樓樹.
紅馬的耶和華使者.
和他.
發現了.
全地安息平靜之後.
就向神.
來去呼求.
神就回答.
神的回答就是.
他要重建耶路撒冷.
他要再在耶路撒冷.

$^{1881}$顯現.
和要賜福給這個民族.
所以這個重建.
在這個時候.
神的重建.
就來得很有意義.
也很有神的幫助.
這個真是事半功倍.
兩位先知的安慰.
相信當時的.
索羅巴伯.
和大祭司約書亞.
一定是如有神助.
這樣來去.
再勉勵.
和那些百姓.
一齊去同工.
所以神的預言.
神的意象.
是很重要的一個群體.
需要有的.
保證.
所以.
趙少福音派.
今天.
大家有沒有留意.
港台是放在中間的.
大家有沒有留意.
沒有考究.
因為禮儀教會.
港台是要靠邊站的.
港台在中間.
是福音派的.
強調的.
因為馬丁路德的宗教改革.
發現.
唯獨聖經.
當群體回到聖經.
回到聖經.
那個宣講.

$^{1921}$就會令人心振奮.
所以就會帶來.
復興.
這個的迎信.
福音派到今天.
都是這樣迎信.
所以為什麼我們要特別崇拜.
一定要有起碼.
二十分鐘至三十分鐘的講道.
正道的原因.
不是某些人特別口才.
而是我們仍然相信.
神會使用.
他的僕人使女.
向我們講.
神的心意的說話.
所以你可以.
透過這些禮儀的設計.
都是這樣的.
心意.
所以第一個意象.
已經告訴我們.
是一個很大的.
安慰.
因為神要回到石安.
如果你看.
以西結書.
神是一步一步由聖殿走.
一路走.
逐步走.
用耶和華的榮光.
做一個標題.
我忘記了是第八章.
以西結書的經文.
一路這樣走.
走到離開耶路撒冷.
你可以想像那種悲慘.
如果我們.
看到靈界.
神一路走.

$^{1961}$你可以想像這個民族.
有七十年的苦難.
但現在不是了.
神現在回來了.
不只是人回來了.
是神自己回來了.
去眷顧石安.
還有一點時間.
我們也看完.
第十八節.
第二個意象.
又是神要保護.
他的子民.
見有四角.
我就問.
神與我說話的天使說.
好像是一個解話的天使.
又是什麼意思?.
撒迦利亞.
常常好奇地問.
這其實是一個.
很好的問題.
我們看到一些事情.
我們內心也要有些思想.
所以盼望我們.
帶著我們的靈魂.
其實每個人都帶著.
沒理由不帶著.
帶著我們的靈魂來崇拜.
敬拜神.
有時候每件事.
你都可以觀察.
看在眼裡.
然後問這是什麼意思?.
這樣會令到你的靈魂.
上心.
譬如你看到一些人.
很隨便地.
你問這是什麼意思?.
你的內心就會有些觸動.

$^{2001}$你就會有些感想.
會化做你內心的禱告.
譬如你.
還讚神.
看到不同的.
事奉的人.
很用心.
很謙卑.
很單純.
聚會充滿了大家的歡樂.
開心.
你就會感受到.
這是什麼意思?.
雖然這裡是借題發揮.
本來的意思不是這樣的.
但是這個可以引申到.
我們去到聚會.
你也可以問自己這個問題.
還有.
神在我身上.
有什麼我可以貢獻?.
所以你就不是觀眾.
你自己真是參與者.
所以.
撒迦利亞就是參與者.
看到異象.
你是欣賞你是好奇.
你問這是什麼意思?.
他很想了解.
和我們現在的民族.
重建有什麼關係?.
你這樣看這些異象.
其實我們回到我們自己.
其實我們都是神.
今天讓我們來崇拜.
有什麼可以看到?.
其實都是這是什麼意思?.
這個問題在我們心裡.
每一個星期.
你帶著這個問題來.

$^{2041}$神會給你解答.
他就回答.
這是打散猶大.
以色列和耶路撒冷的國.
四方的仇敵.
耶和華又指.
四個匠人給我看.
我說.
他們來作什麼呢?.
又再問.
所以才知道都要有好奇心.
給你一個異象.
你又什麼都不問.
不問就變成.
你根本已經心如止水.
但如果你還是有那個熱誠.
你就要問.
又是什麼意思?.
四個的公將.
公將當然是興建.
你想想.
撒迦利亞日頭.
這個異象是晚上看到的.
日頭就看到很多人搬東西來.
去重建.
撒迦利亞都會去工地看.
你可以想像.
現在晚上就看到有公將.
接著就說.
他說.
這是打散猶大.
以色列和耶路撒冷的國.
這是打散猶大的國.
使人不敢抬頭.
但這些匠人來威嚇列國.
打掉他們的國.
就是舉起.
打散猶大地的國.
這個異象就很簡單.
但背後的含義就很深.

$^{2081}$讓撒迦利亞看到的.
不是日頭看到的.
那些竇木老.
整聖殿的人.
猶太人的同胞.
現在在異象裡面.
神給他首先看到.
原來耶路撒冷過去的雞毛鴨血.
原來有四國.
打散猶大.
打散以色列.
打散耶路撒冷的四國.
即是四方的仇敵.
敵黨神子民.
但神又給他看到.
四個匠人.
就是公匠.
他們有一個.
四個匠人.
就是他們有能力.
去打散四國.
以致神的地方.
不會再受威嚇.
而且可以來去重建.
異象就是這麼簡單.
但後面的就會再複雜一些.
就不是這麼簡單.
所以你可以看到.
又是一個主題.
全部是關乎耶路撒冷的重建.
甚至關乎到這個民族.
再被神留意.
是留意的祝福.
以及講述給撒迦利亞知道.
過去.
是有一些表面的.
國與國之間的競爭.
因為耶路撒冷.
以色列是剛剛在.
幾個強大的帝國的中間.

$^{2121}$歐亞非的中間.
歐洲亞洲.
非洲的中間.
所以由埃及.
是可以經過以色列.
去到亞洲.
但是那是地上的.
國土而已.
歷史而已.
撒迦利亞看到的是.
因為他是精神的領袖.
他看到的就是那個異象裡面的.
靈界.
所以他看到的就是.
原來過去一路的歷史.
其實將來都是.
四國是敵對神子民.
但是有四個將人.
當然沒有再具體講.
這四個將人是哪些呢.
具體的顯明.
只是講了.
拆毀和建造.
那個角.
就好像大錘一樣.
有沒有見過拆樓的大錘.
大家可能沒有見過.
一些很大的錘.
這樣打在樓上就倒了.
可能這些角.
都是當時他們用來.
拆毀的.
但是現在神就.
被撒迦利亞看到.
就是有四個將人.
是重建的.
是興建的.
是眷顧耶路撒冷的.
所以我們看到.
兩個異象.

$^{2161}$第一張已經不斷告訴我們.
那個焦點就是.
在神再一次.
眷顧耶路撒冷.
眷顧他自己的子民.
所以這個時候.
這個再重建.
爛了尾的.
聖殿.
15年.
一直到大理烏王的第二年.
這樣的一個.
停頓.
神現在那個時間到了.
他要來興起先知.
哈該.
撒迦利亞.
所以撒迦利亞的預警.
那個密使是豐富過哈該.
因為還有很多.
其他的異象.
都好像一脈相承.
是指向耶路撒冷.
所以你可以想像到.
神重新.
眷顧耶路撒冷.
是這個時代的主調.
是.
先知撒迦利亞.
的時代的主調.
所以沒有.
其他的外族.
可以來去抵擋.
神的指紋.
所以當時有.
波斯的人.
上書.
去到.
大理烏王的時候.
他們又再一次.

$^{2201}$尋查那些典籍.
發現原來真是當年.
古烈的時候.
是准許他們重建的.
所以又.
給他們有一個.
可以興功建造的.
理由.
所以你會看到.
真是神的話語.
是先行的.
撒迦利亞首先.
自己要被復興.
他講的那些語言.
就能夠安慰.
那個時代.
重建的百姓.
所以這裡都求神讓我們能夠.
對於神自己的.
話語.
我們自己.
都有一個裝備.
所以如果大家從來.
都不讀這些舊約的書卷.
其實我們就會.
少了對神.
他那種有血.
有肉的感情.
會少了一點.
因為去到新約神已經透過耶穌.
來顯現.
做物主很少.
直接對話.
耶穌當然已經反映神的性情.
但是某程度.
都不是很足夠.
所以我們都能夠.
透過大家.
很興趣地讀讀.
這些舊約的先知書.

$^{2241}$讓我們能夠稍微掌握.
神自己的性情.
我們今天一起祈禱.
多謝天父給我們今天.
開始第一章的撒迦利亞書.
多謝你給我們.
看到一個很重要的訊息.
就是撒迦利亞的年代.
神你重新.
眷顧這個民族.
要將恩典帶給這個民族.
是從你眷顧石安.
留心的私恩.
還有你心裡面.
火熱.
來祝福這個民族.
如果不是真的有神自己的憐憫.
眷顧.
其實沒有一個人類.
沒有一個民族可以繼續的存活.
多謝你給我們有救恩.
多謝你藉著耶穌基督.
將我們拯救.
使我們能夠明白.
我們的人生因為主耶穌而有意義.
是真正的意義.
求主讓我們能夠有生之年.
把握.
將神的恩典.
去跟其他人分享.
使更多人得到.
神你自己的救恩.
多謝你聽我們禱告.
奉耶穌基督的名.
阿們.
\newpage



\section{詩篇 128:1-6}
\label{sec:BaO8MW3rJk0}
\textbf{2023-09-16 晚堂《福氣:敬畏主+遵主路》詩篇 128\_1-6 講員:陳麗蟬牧師}
\newline
\newline
連結: \href{https://youtube.com/watch?v=BaO8MW3rJk0}{\texttt{https://youtube.com/watch?v=BaO8MW3rJk0}} ~~~~ 語音日期: 2023-09-19
\newline
\newline
\hyperref[sec:Kpbv61GAjrM]{\small{< < < PREV SERMON < < <}}
~
\hyperref[sec:index_chronic]{\small{[返順時目]}}
~
\hyperref[sec:index_scriptual]{\small{[返順卷目]}}
~
\hyperref[sec:Ah5GT26Za2E]{\small{> > > NEXT SERMON > > >}}
\newline
\newline
詩篇 128:1-6
\newline
\begin{longtable}{cl}
\hline
\hline
章節 & 經文 (和合本修訂版)\\
\hline
128:1 & \begin{tabularx}{0.7\textwidth}{X} 凡敬畏耶和華、遵行他道的人有福了！ \end{tabularx} \\ \\ \relax
128:2 & \begin{tabularx}{0.7\textwidth}{X} 你要吃勞碌得來的；你要享福，凡事順利。 \end{tabularx} \\ \\ \relax
128:3 & \begin{tabularx}{0.7\textwidth}{X} 你妻子在你內室，好像多結果子的葡萄樹；你兒女圍繞你的桌子，如同橄欖樹苗。 \end{tabularx} \\ \\ \relax
128:4 & \begin{tabularx}{0.7\textwidth}{X} 看哪，敬畏耶和華的人必要這樣蒙福！ \end{tabularx} \\ \\ \relax
128:5 & \begin{tabularx}{0.7\textwidth}{X} 願耶和華從錫安賜福給你！願你一生一世看見耶路撒冷興旺！ \end{tabularx} \\ \\ \relax
128:6 & \begin{tabularx}{0.7\textwidth}{X} 願你看見你的子子孫孫！願平安歸於以色列！ \end{tabularx} \\ \\
[1ex]
\hline
\hline
\end{longtable}
$^{1}$弟兄姊妹晚安.
今天又有機會跟大家一起用上帝的說話.
彼此的勉勵.
快要中秋節了.
給大家一個燈謎猜猜.
謎底就是這樣.
謎面是這樣的.
一口填衣六全.
今年的中秋晚會不知道會不會出這題.
知道的 不過你們不必回答我.
一口填衣六全.
就是這個.
Yoyo怎麼不出給我.
出了出了.
就是這個福字.
在我們中國人來說.
會覺得有這個三著.
有這個土地.
又能夠每一頓飯都吃得飽.
就是這個福氣.
中國人就是這樣很著眼於一些物質上的.
我們要有什麼.
譬如我們要有房子.
有車 有老婆 有子女.
有很多的東西.
就是一種福氣.
所以我們過年的時候一定說.
恭喜發財.
有很多的錢.
就可以買到更多的物質.
但是是不是真的擁有這些物質.
房子 車 房子這些東西.
有子女 老公 老婆呢.
就等於幸福.
一個福.
大家都過來人.
大家心裡面都有不同的滋味.
不同的經歷.
但是在這個屬靈的層面裡面.
又怎麼看這個福字呢.

$^{41}$這個屬靈的福字.
它就是看旁邊的左手邊那裡.
它看一個 當是一個「士」字旁.
這個「士」字就指的是這個神.
然後旁邊就一口田.
就是說先有這個神.
和這個人在這裡.
是神和它所做的人在一起.
然後又賜給它有這個田地.
就是說它會將一些物質.
就賜給我們.
這個就是這個福氣.
這個在屬靈層面的福氣.
這個就和我們中國人很不同.
中國人總之是物質就可以了.
但是在這個屬靈的層面.
用「士」字來看的時候.
是先有神.
然後它就和它所做的人是一口.
一齊 聯合在一起的時候.
它就在物質上.
它會供應給我們.
這個以色列人.
他們每年都要有三次.
上耶路撒冷那裡去獻祭的.
這個耶路撒冷就在山上高的.
所以他們在不同的地方.
一路去耶路撒冷城的時候.
就越走越高 走得很辛苦的.
當他們就去到這個耶路撒冷城的時候.
他們就會唱歌的.
在詩篇裡面.
有好幾篇有這個叫做「上行詩」的.
這個詩篇有128篇.
就是其中的一篇.
他們一路唱詩的時候.
在那裡大家彼此的鼓勵.
他們越來越近聖殿.
所以他們心情又很開心 很激動.
但另一方面.

$^{81}$因為越走越高 越走越斜.
他們就很辛苦.
當他們一路唱詩的時候.
特別是在128篇裡面.
其實描繪著一幅很美的圖畫.
就是在說一個敬畏人的生命.
他會是怎樣的呢?.
我們今天就一起來看看這個.
敬畏神的人的生命是怎樣的.
當摩西帶了一幫以色列人出了埃及.
然後將會進入迦南地.
當他們在這個曠野的時候.
耶和華就兩次三次四次這樣.
去祝福這一幫的以色列人.
將來有一天.
你們要進入迦南地的時候.
你們要怎樣呢?.
在《新命紀》第十章裡面其中一章.
它就說到.
你們要敬畏耶和華你的神.
遵行他的道.
愛他.
盡心盡性事奉他.
進入迦南地的時候.
一定要敬畏他.
還有要忠心的事奉他.
遵守他的誡命律例.
這是我今天所吩咐你的.
所以當他遵守他.
敬畏他.
還有遵守他的律例的時候.
上帝就會賜福給他.
上帝就會賜福給這個以色列人.
甚麼叫做敬畏呢?.
敬畏很簡單的說.
就是聽話.
聽話.
也就是一個尊敬.
以上帝為上帝.
有時候我們聽到那些人吵架.

$^{121}$那些夫婦吵架.
你究竟有沒有聽我說話的?.
或者有時候吵得激烈一點.
你究竟有沒有當我是老公的?.
當我們講到這些說話的時候.
我們就對人沒有了那份尊重.
上帝要我們那個敬畏.
就聽他的話.
他要我們怎樣做的時候.
我們知道他所講的是對的.
也知道我們的好處不會在他以外的時候.
我們聽他的話.
這個敬畏是一個在心態上高的.
所以很難講要做些甚麼的.
但是他很清楚地講.
就是要去進行他的路.
是要走耶和華的路.
在第一至三節.
很清楚地講到.
是不從惡人的計謀.
不站罪人的道路.
不助褻瀆萬人的作為.
講到那些惡人.
在做無理取鬧的事情.
千萬不要跟著他一起無理取鬧.
那些惡人的思想.
淫邪的思想.
想著不好的事情.
你千萬不要跟著他一起去做.
不要跟著他一起去想.
這個是上帝在消極的層面裡面.
千萬不要跟這個罪惡完完全全分開.
然後他就說.
為喜愛耶和華的律法.
晝夜思想他.
這個就說我們要專注在上帝的律法裡.
晝夜思想.
就是說為了我們去時時去想他的說話.
讀他的說話的時候.
這樣我們那個屬靈根基.

$^{161}$才能夠打得深入.
這些人這樣才是一個有福的人.
他要像一棵樹栽在溪水旁.
按時候結果指.
葉子也不枯幹.
凡他所做的盡都順利.
這個詩篇裡面其實很多時候都會講到.
敬畏他.
走他的路.
這樣我們就會怎樣呢?.
凡事順利.
就有福了 就有夢福了.
上帝要我們先跟他建立一個緊密的關係.
我們看到.
我們香港每一次打風.
又或者剛剛上個禮拜的黑雨.
我們看到香港很多的樹木.
都很容易塌下來.
甚至連根拔起.
那次黑雨的時候.
我就看到新聞那裡就講到.
為什麼香港的樹木這麼容易塌下來呢?.
他說香港的樹木很多時候.
在不知道什麼地方種到差不多.
然後就拿到街邊的路那裡去種的.
以致根本它不是長得很深.
是很淺的抓住一些的地.
以致大風大雨的就很容易塌下來.
同樣都是.
如果當我們那個屬靈根基打得不深.
不穩定.
我們根本聖經就放在桌面時時都不看的.
時時都不讀的話.
我們這個屬靈的根基就很膚淺.
如果一遇到人生裡面有什麼風浪的時候.
就很容易會塌下來.
很容易離開教會.
所以在聖經裡面不斷地提醒我們.
我們一定要打穩我們屬靈的根基.
要去敬畏耶和華.

$^{201}$去走祂喜悅的道路.
如果我們敬畏祂.
又走祂的道路的時候.
又有什麼好呢?.
這個蒙福的人.
蒙福的人凡事都順利.
不過你要記住.
這個凡事順利.
就是上帝給過我們的.
有時我很怕的.
很怕聽到頂尖的說.
哎呀真是好彩.
今天搭到村巴不用等很久.
哎呀今天真是好彩.
去到街市裡面有菜買.
又是好彩.
當我們說好彩.
當我們說.
哎呀這次真是幸運.
這些全部都是和一些幸運.
一些有關的.
但我們知道.
我們屬於耶穌基督的人.
不單止屬於耶穌基督的人.
應該世界上所發生的所有事情.
都有上帝掌管著的.
所以當我們遇到一些事情.
你發覺.
嘩真是好順利.
以後不要再說好彩了.
一切都有祂掌管的.
所以當我們遇到一些.
好順利的事情.
我們知道.
這是上帝給過我們特別的祝福.
在《施比人版》28篇裡面.
有三種上帝給過我們的.
第一種是.
吃勞碌得來的.
即是說.

$^{241}$當我們辛勤工作了之後.
我們能夠享受成果.
能夠坐在這裡.
嘩坐在這裡.
吃飯.
好開心 好舒服.
但現在的香港人.
很多時候.
在工作裡面.
為了拿到成就.
或者為了有一份穩定的工作.
不斷地去賺錢.
不斷地工作.
不斷地工作 做 做 做.
我有一個親戚.
其實不單止一個.
數落了三,四個.
三十多歲.
其中一個.
他在某一間醫學院裡面.
做一些行政工作.
他們整隊team很好的.
想搞好一些行政.
以致教授或者醫學生.
可以在上課的時候.
有些設備.
妥當一些.
做的事會好一些.
於是整隊team很好的.
很合作的.
但是慢慢一直做下去的時候.
你知道一直做事的時候.
自然會有些碰撞.
上司就會越要求越來越多.
你會發覺到.
他每天九點鐘就要回到公司.
但是通常下班.
是晚上十一點多鐘.
如果在開學前.
或者到考試的時候.

$^{281}$他們下班的時間.
是凌晨兩,三點.
他們星期六,日.
除了星期日早上不用上班之外.
星期六,日.
星期日下午都要上班.
星期日不是讓他們上崇拜.
而是星期日.
讓他們睡覺,休息.
下午繼續上班.
有些什麼,有些力量.
我經常都說.
你們變成了工作的奴隸.
結果我看到這班年青人.
做到死氣沉沉.
都不跟我們有什麼家庭聚會.
經常說胃痛.
說他的同事更慘.
開始有些幻覺出來.
有時被上司罵得厲害的時候.
有些幻覺.
進入了精神病院.
他的親戚也開始有抑鬱.
吃抑鬱藥.
所以你會發覺.
這個能夠吃奴隸得來的.
不是必然的.
當中三十多歲的人.
薪金很高.
星度也很高級.
但是他們能不能吃到.
從奴隸得來的呢?.
我知道他的親戚.
有時他們不會出去買飯.
總之就放一些餅乾.
放一些蛋糕在那裡吃.
所以,弟兄姊妹.
這個能夠去做事.
去辛勤的做事.
然後去吃一餐.

$^{321}$能夠去享受奴隸得來的.
是一種福氣.
上帝叫我們.
工作有事.
休息也要有事.
除了工作狂之外.
有些人可能因為疾病.
所以他真的吃不到東西.
但是上帝說.
吃奴隸得來的.
是一種福氣.
第二種.
就是在家庭裡面.
那份歡愉.
在第三節裡面.
他講到妻子好像多結果子的葡萄樹.
看到這些葡萄子.
是不是很開心?.
是否很想吃?.
是吧?.
這個葡萄樹.
在中東一帶.
是一個很時事堅的農作物.
我們知道葡萄的用處很多.
他們可以用來慢慢釀酒.
可以做葡萄乾.
葡萄餅.
這個葡萄餅.
是中東主要的糧食.
葡萄樹在這裡成長.
在一個很安定的環境.
是一些爛藤的植物.
不能夠掛在這裡.
掛得好像一般好看.
又掛在這裡.
希望它吸收更多陽光.
當你掛來掛去的時候.
它就出不到這個葡萄.
它一定要在一些很安定的環境裡面.
才能成長.

$^{361}$所以在竹林層面裡面.
這個葡萄樹.
是象徵著平安.
所以是妻子在內室.
好像多結果子的葡萄樹.
不只是說它生養眾多.
而是可以帶給整個家庭裡面.
來和諧的相處.
還有我們看到葡萄樹.
果實纍纍.
剛才我說.
我在網上看到這幅.
我覺得很開心.
整個心就好像.
總之很開心.
你覺得它很美很美.
很吸引.
如果真的時候.
就很想立即撕一些下來吃.
就在那裡說到.
這個美.
是說到人的美.
那麼吸引的時候.
不是因為它的外貌.
這個比喻到這個屬性層面裡面.
不是說到它的外貌.
而是裡面的內在美.
弟兄姊妹.
我想你們都試過這樣的經歷.
就是有時.
你見到.
走著走著街.
見到一些很美的女孩.
我也會的.
有時見到一些美的女孩.
或者她穿的裙子.
多美的時候.
我也會看一看她.
多看一兩眼.
但是.

$^{401}$當你發覺到她突然間.
拿起一支煙來.
去吸一吸的時候.
或者你去隔離的時候.
不經意聽到她說幾句粗口.
你會覺得她的美都沒有了.
甚至你會想反胃.
想作嘔.
所以.
她說到這個葡萄樹.
很美.
外表很美 很吸引.
但是說到.
這個人的層面的時候.
因為敬畏耶和華.
她是尊衡她的道理的人.
那個家庭.
所以說到裡面的內在美.
我有一次跟一位老太太.
陪她去博愛醫院那裡.
探望她的丈夫.
一進到病房.
就跟病床的病人.
還有家屬打招呼.
你丈夫進去兩天.
三天.
為甚麼病房那麼熟.
原來.
當她拿食物出來.
給她的丈夫吃的時候.
就自己吃了.
我看到這個老太太.
她真的很美.
她就周圍.
這張床.
你今天好不好一點.
今天有沒有胃口.
嘩 家裡人的食物那麼好吃.
一會兒又走去.
跟其他病人的家屬聊天.

$^{441}$我站在那裡看的時候.
嘩 我覺得她還像牧師多於我.
一個很美.
充滿了愛的一位老太太.
一個敬畏耶和華的人.
進行她所吩咐.
以愛來待人的時候.
在日常生活當中.
很容易流出一個很和平.
一個很美的生命.
並且處處都為人著想.
所以說到這個夢幅的人.
他的家庭是歡愉.
因為他的妻子也是敬畏耶和華.
他表達出來的.
給這個家庭平安.
他處處有愛的表現.
接著這個夢幅的家庭.
就是家庭歡愉.
第三節下半節講到.
兒女圍繞著他的卓子.
如同橄欖樹苗.
原來橄欖樹是全世界最暢壽的樹.
它結很多很多的果.
據說現在有千多二千年的橄欖樹.
仍然存在在意大利裡面.
這個橄欖樹老了.
你說 嘩 千多二千歲.
所以有些橄欖樹很醜.
不斷扭曲.
當它老了之後.
它一直出那些原本出的樹幹.
它已經不能夠再結果子.
不過上帝很特別.
你給它上千多二千年的命.
它又不可以再出這個果子.
但上帝卻在樹的底部.
會出一些新的葉子.
你會看到在樹底下.
會出一些新的樹苗.

$^{481}$一路出一路出.
看上去就像很多新葉圍住了這棵老樹.
詩人就說這棵這樣的情景.
就像兒女圍著爸爸媽媽.
一起坐在飯桌裡面吃飯.
這個橄欖樹就象徵著家庭的和諧 溫馨.
你的兒女圍繞著你的卓子.
如同橄欖樹苗.
我爸爸媽媽先後離開了世界.
回天父家已經十幾年了.
在以前我們六兄弟姐妹.
或多或少都會有些爭拗.
特別有子女那些.
同年齡的才慘.
爭取成績 爭取入哪間學校.
出來工作的時候.
又爭取誰比較有成就.
但我覺得當我作為上帝的兒女的時候.
我就需要學習做一個和平的使者.
特別當我爸爸媽媽都先後回天父家了.
我開始聯絡每一個家庭的哥哥姐姐.
起初很難的.
這次就叫這三個家庭.
這個不要來 這個和這個吵架.
有時叫這個的時候.
他來不來 他來我就不來了.
一個家庭都是這樣.
爸爸媽媽在的時候勉強都可以.
但是一走了的時候.
維繫需要真的要很多力量.
但現在經過了十幾年的努力.
我們很感恩 很感恩.
大時大節我們全家大小.
除了那些要上班的那些.
大時大小我們一起吃飯.
我們現在大約每兩個月左右.
我們就會輪流.
不只是我們這些年長的姑姐姑媽那些.
就是我們六兄弟姐妹.
我們就會輪流去到不同的.

$^{521}$姐姐,哥哥,弟弟那裡的地方吃飯.
上個星期他們就回到我家裡了.
聊聊天 說說話.
吃完飯就切一個月餅.
我們那天出席都九個人.
連同哥哥,嫂子,弟夫那些.
切了一個月餅而已.
每人一口袋子 喝一杯茶.
之後我和我姐姐說.
我覺得在那天我們覺得很幸福.
她都說是.
覺得那份幸福是很滿足.
是說不出那份歡愉.
所以就好像詩人所說.
當家人圍在一起的時候.
兒女圍繞的桌子.
好像橄欖樹苗.
在當中那份和諧,溫馨.
所以這個橄欖樹是和平.
我們在當中作為上帝的兒女.
在家庭裡面努力維持和平.
還有就是橄欖樹苗.
原來橄欖樹除了那些.
橄欖的果子很有用.
可以榨橄欖油 很漂亮的.
它老了之後.
樹底下那些樹苗.
它們慢慢長大的時候.
又會同樣出很多橄欖果實.
它們出來的果實更漂亮.
榨出來的油更香.
還有這棵橄欖樹.
它的木材很堅硬.
可以用來做家具.
還有在古代的時候.
這個橄欖油.
就是當高末君王的時候.
用橄欖油.
當我們教會去做末尤祈禱會的時候.
我們不會用花生油.

$^{561}$也不會用芥花籽油.
我們會用橄欖油.
這個在《根翻聖經》裡面.
是一個祝福的油.
所以這棵橄欖樹.
象徵著團結,生命力,有用,代代相傳.
一個蒙福的人.
他的家庭就是上帝這樣賜福.
不過也要大家一起去努力.
當這個家人一起吃飯的時候.
真的很開心.
閒話家常.
不過你看到現在很多的家庭.
是怎樣吃飯的.
各自修行.
他不拿自己的飯碗.
夾好菜.
然後坐在電視機旁.
或者回到房間吃飯.
已經偷笑了.
可能勉勉強強坐在一起.
在桌子上.
然後各人目口目面.
然後各人自己拿著電視機.
自己戳戳戳.
爬兩檔.
又戳戳戳.
老人家就在外面看著電視.
或者有很多人的家庭是這樣的.
家不成家.
或者我們已經是一段時間是這樣.
好像之前我們自己家也是這樣.
我們挽回不了這些冷漠的境況.
不過弟兄姊妹.
從禱告開始.
還有我們需要有實際的行動.
這個蒙福的人除了敬畏,尊敬,聽上帝的話的人之外.
下一步就是怎樣?.
進行他的教訓.
就是要實質地走出來.

$^{601}$所以最初走的時候是辛苦的.
要聽很多的話.
但是在今天裡面.
我看到上帝就是這樣賜福我們一家.
雖然可能背後都會有些是是非非.
是吱吱喳喳的.
但是圍在一起的時候.
我們感到歡愉.
我們感到溫馨.
所以弟兄姊妹第一步是祈禱.
第二步就是我們先做一個和平使者.
我們要走出來.
我們終有一天可以見到.
這個以女為妖是卓子.
如同橄欖樹苗的幸福.
好了 講完這個家庭之後.
就講最後了.
就是這個蒙福的人.
他可以令到這個社會國泰民安.
有什麼資格?.
我又不是做官.
在第五節裡面.
聖經這樣講.
「願你一生一世看見耶路撒冷的興旺.
願你看見你子子孫孫.
願平安歸於以色列」.
敬畏耶和華的人.
遵行他道路的人.
除了個人有那個福樂.
吃勞碌所得來的.
家庭也都會蒙福.
其實施人就說.
接著可以去延伸出去.
所處的社會.
各國家都可以能夠平安.
可以國泰民安.
我們知道家庭是一個塑造人格.
塑造價值觀的地方.
所以一個健康的家庭.
就能夠塑造下一代.

$^{641}$他們有一個健康的生命.
有一個正確的價值觀.
我想在好幾年之前.
國內已經流行一句話.
多一個基督徒.
少一個罪犯.
多一間教會.
少一座的監獄.
一個敬畏耶和華的父母.
他遵行他道的.
他也都能夠去培養他的子女.
成為一個有良好品格的子女.
所以弟兄姊妹.
第一步祈禱.
第二步就是你去走出來.
這個美國汽車大王亨利福特.
在公司大變之後.
他就回家了.
雖然經歷了大變.
他不是變成窮光蛋.
但是他在家鄉裡面.
建了一間很小很簡陋的房屋.
那些人說.
你都不是窮到沒錢了.
為什麼建一間這麼小這麼簡陋的房屋.
他就回答.
他的家庭不是一座建築物.
不是一座建築物.
是一個家.
他這間房屋雖然小.
並且很簡單.
很樸實.
但如果在家裡面.
充滿了耶穌基督的愛的時候.
這間就是豪宅了.
這間是豪宅了.
所以我們要令到我們的房屋.
變成豪宅.
除了努力工作之外.
更加需要在屬靈層面裡面.

$^{681}$將耶穌基督的愛在當中實行出來.
所以弟兄姊妹.
這裡說到一個夢幸的人.
赫魯魯德來.
你妻子在你的內室.
好像多結果子的葡萄樹.
兒女圍繞著你的卓子.
兒童橄欖樹苗.
所以我們要珍惜這些場面.
努力地去建立.
有人說.
幸福是上帝賜福給我們的.
家庭是上帝建造.
在地上建造的天國.
家庭是上帝建造在地面上的天國.
不是地獄.
幸福的家庭.
是一個天國的實現.
盼望我們在家庭裡面.
我們大家一起努力.
我們同心禱告.
親愛的天父上帝.
我們多謝你.
因為你定義賜福給你所創造的人.
我們多謝你.
我們知道無論個人.
或者家庭.
或者全地.
你都很想我們.
去領受你給我們的各種福氣.
幫助我們.
幫助我們在生活的層面.
在工作的層面裡.
我們敬畏你.
聽你的話.
遵行你的道.
以致無論個人.
我們的家庭.
整個社會.
都能夠成為大夢.

$^{721}$祝福的人.
禱告感恩.
奉主基督的聖名.
阿們.
\newpage



\section{約翰福音 14:1-14}
\label{sec:Ah5GT26Za2E}
\textbf{2023-09-17 早堂《你使我堅壯》約翰福音 14\_1-14 講員:蕭建輝牧師}
\newline
\newline
連結: \href{https://youtube.com/watch?v=Ah5GT26Za2E}{\texttt{https://youtube.com/watch?v=Ah5GT26Za2E}} ~~~~ 語音日期: 2023-09-19
\newline
\newline
\hyperref[sec:BaO8MW3rJk0]{\small{< < < PREV SERMON < < <}}
~
\hyperref[sec:index_chronic]{\small{[返順時目]}}
~
\hyperref[sec:index_scriptual]{\small{[返順卷目]}}
~
\hyperref[sec:Hdjnt297x1A]{\small{> > > NEXT SERMON > > >}}
\newline
\newline
約翰福音 14:1-14
\newline
\begin{longtable}{cl}
\hline
\hline
章節 & 經文 (和合本修訂版)\\
\hline
14:1 & \begin{tabularx}{0.7\textwidth}{X} 「你們心裡不要憂愁；你們信神，也當信我。 \end{tabularx} \\ \\ \relax
14:2 & \begin{tabularx}{0.7\textwidth}{X} 在我父的家裡有許多住處；若是沒有，我就早已告訴你們了。我去原是為你們預備地方去。 \end{tabularx} \\ \\ \relax
14:3 & \begin{tabularx}{0.7\textwidth}{X} 我若去為你們預備了地方，就必再來接你們到我那裡去，我在哪裡，叫你們也在哪裡。 \end{tabularx} \\ \\ \relax
14:4 & \begin{tabularx}{0.7\textwidth}{X} 我往哪裡去，你們知道那條路。」 \end{tabularx} \\ \\ \relax
14:5 & \begin{tabularx}{0.7\textwidth}{X} 多馬對他說：「主啊，我們不知道你去哪裡，怎麼能知道那條路呢？」 \end{tabularx} \\ \\ \relax
14:6 & \begin{tabularx}{0.7\textwidth}{X} 耶穌對他說：「我就是道路、真理、生命；若不藉著我，沒有人能到父那裡去。 \end{tabularx} \\ \\ \relax
14:7 & \begin{tabularx}{0.7\textwidth}{X} 既然你們認識了我，也會認識我的父。從今以後，你們就認識他，並且已經看見他了。」 \end{tabularx} \\ \\ \relax
14:8 & \begin{tabularx}{0.7\textwidth}{X} 腓力對他說：「主啊，將父顯給我們看，我們就知足了。」 \end{tabularx} \\ \\ \relax
14:9 & \begin{tabularx}{0.7\textwidth}{X} 耶穌對他說：「腓力，我與你們在一起這麼久了，你還不認識我嗎？看見我的就是看見了父，你怎麼還說『將父顯給我們看』呢？ \end{tabularx} \\ \\ \relax
14:10 & \begin{tabularx}{0.7\textwidth}{X} 我在父裡面，父在我裡面，你不信嗎？我對你們所說的話不是憑著自己說的，而是住在我裡面的父在做他的工作。 \end{tabularx} \\ \\ \relax
14:11 & \begin{tabularx}{0.7\textwidth}{X} 你們要信我，我在父裡面，父在我裡面；即使不信，也要因我所做的工作信我。 \end{tabularx} \\ \\ \relax
14:12 & \begin{tabularx}{0.7\textwidth}{X} 我實實在在地告訴你們，我所做的工作，信我的人也要做，並且要做得比這些更大，因為我到父那裡去。 \end{tabularx} \\ \\ \relax
14:13 & \begin{tabularx}{0.7\textwidth}{X} 你們奉我的名無論求甚麼，我必成全，為了使父因兒子得榮耀。 \end{tabularx} \\ \\ \relax
14:14 & \begin{tabularx}{0.7\textwidth}{X} 你們若奉我的名向我求甚麼，我必成全。」 \end{tabularx} \\ \\
[1ex]
\hline
\hline
\end{longtable}
$^{1}$好 各位錦宣家弟兄姊妹平安.
代表華鬼堂送上祝福.
願主的恩惠平安與你們同在.
今天來到這裡 我自己有一種親切感.
不單止是唱詩歌 因為我也是唱生命聖詩 紀哲的詩歌長大.
不過當然到了近年可能就多一些短詩.
所以有一種親切感.
另外親切感就是因為牧羊臥城長的許金金牧師.
這裡是他第一個牧會的地方.
所以今天來到這裡的時候有一種親切感.
今天可以在當中和大家一起分享主的話.
剛才讀經的時候 相信大家最耳熟能詳的兩段經文.
分別是第二至三節和第六節.
我們先和大家看看這兩段經文說什麼.
二至三節 我去完事為你們預備地方去.
我去為你們預備地方 就必再接你們到我那裡去.
我在那裡叫你們也在那裡.
經文指出耶穌先離開到父那裡為門徒預備地方.
而按我們今天理解這段經文.
耶穌是為了跟隨他的人.
先為我們預備一個永恆的家鄉.
其實耶穌是神子 根本用不著自己去預備.
執頭執尾 對不對.
所以耶穌在這裡用了當時門徒能夠明白的說法.
為人預備地方 家主要去招呼人客的時候.
他要好好打點地方 親自來我去為你預備地方.
是表達很歡迎貴賓的到來.
所以耶穌親自為人預備地方.
是指耶穌很歡迎門徒到父那裡去.
至於第六節 更加是我們人人都會背的經文.
耶穌說 我就是道路真理生命.
若不藉著我 沒有人能到父那裡去.
按照第六節的下半節說 到父那裡去.
就是說耶穌自己就是引導人去到父那裡的道路.
耶穌是真理 他是道成肉身來到世間.
將神國的真理去指示世人.
而他自己也是真理本身.
耶穌是生命 他是天上降下來的嗎哪.
吃耶穌這個生命的糧 就將要得著生命.
第二至三節 第六節.

$^{41}$給我們很安慰 很有盼望 很感動的感覺.
而上述這兩處的經文.
在耶穌和門徒對話的時候.
他是不是想表達這些安慰 盼望 感動呢.
約翰福音十四章 今天我們看的這段經文.
是在學者稱為耶穌賞十價錢最後講論的篇章裡面.
整個的講論是很容易看到分界線的.
是由第十三章的三十一節開始.
一直到十七章的結尾.
十七章的二十六節 就是那個結尾.
這個講論的開頭是十三章三十一節.
他寄出去 耶穌就說 最頂那句.
這個他是在說猶大.
吃完晚餐之後 猶大先行離開.
做什麼呢 要找官兵 法利 塞人 祭司來捉拿耶穌.
當猶大離開之後 耶穌就開始和餘下的門徒.
講這一段的最後講論.
這裡是開始 到十八章的一節.
耶穌說了這話 就和門徒出去 最下面那句.
說了這話 原來十八章一節以上的經文.
就是最後講論的末端.
十八章一節之前的就是這話的內容.
我們可以看到整個講論的範圍.
而我們今天看的十四章 一至十四節.
先和大家看第一節.
你們心裡不要憂愁.
到底有什麼愁呢.
當時耶穌說話的氣氛是怎樣.
我們和大家很快去看一段短片.
是耶穌傳 你也會看過的.
當中大概分半鐘 記載關於最後晚餐的內容.
大家看的時候 特別留意門徒的神情.
我們看一看短片.
(門徒們的晚餐).
我一直都希望在受難之前 和你們吃逾越節晚餐.
我告訴你們 在上帝的國降臨之前.
我不會再吃這個晚餐.
使到葡萄樹結出果子.
你們把這杯葡萄酒分來喝.
我告訴你們 在上帝的國降臨之前.

$^{81}$我不會再喝這些酒了.
你賜食物給地上的人.
這是我的身體為你們而寫的.
你們要這樣做來紀念我.
這杯是我用血所流的新約 為你們所流的.
你們看 出賣我的人 和我坐在同一張桌子上.
當然 人會按照上帝的安排受死.
但是出賣人質的人就有禍了.
不是我啊 沒可能的.
其實我加快了一點.
所以你們看到門徒的動作.
為什麼轉得這麼厲害.
加快了一點.
從片段可見 當耶穌說要受難.
在神國降臨之前.
他不會再吃這個晚餐.
然後出賣我的人 和我同坐一桌的時候.
門徒感到困惑 迷茫 憂愁.
也就是說 第十四章的背景.
氣氛是相當沉重.
但是不知道去了哪裡.
耶穌又說有人要賣你.
又說要離開大家.
整段的講論是以一個沉重的氣氛開始.
這個氣氛要一直去到整個講論的哪個位置.
門徒才開始有一點點正面的表達.
快速地我們看一個數據.
整個最後講論的篇幅一共有125節.
門徒說的話和耶穌說的話佔的比例有多少呢.
我們看看.
十三章那裡門徒說了兩節.
十四章三節.
十五章十七章門徒沒有說話.
耶穌自己一個說完.
接著十六章有四節.
門徒從這個段落當中.
這個圖片我們看到.
門徒只是說了九節的說話的內容.
就是說整個的講論九成的篇幅.
是夫子耶穌說完.

$^{121}$而餘下的一成.
我和大家再去看看門徒其實說了些什麼.
我們看看紅色的字.
主往哪裡去.
我為什麼現在不能跟你去.
怎麼知道哪條路呢.
求主將父顯給我們看.
主啊為什麼要向我們顯現.
不向世人顯現呢.
這是什麼意思呢.
到底是什麼意思呢.
我們不明白他所說的話.
還有最後.
二十九至三十節.
門徒說如今你是明說.
並不用比喻了.
夫子去教導的時候喜歡用比喻.
現在我們曉得你凡事都知道.
也不用人問你.
因此我們信你是從神出來的.
可以看到在這九節當中.
門徒說話的這九節當中.
是有六個問題在裡面.
都是在表達他們不明白耶穌說什麼.
是去到十六章的二十九至三十節.
這裡才看到有我們信你是從神出來的.
故此整個耶穌的離別講論.
絕大部分都是在一種很困惑迷茫的氣氛當中.
是要去到十六章尾這裡.
才看到門徒有一點正面的宣告.
這個數據告訴我們.
如果我們只是閱讀最後講論比較前面的部分的話.
我們基本上只會看到門徒迷茫和困惑的一面.
而門徒的信心宣告是要去到比較後段的十六章尾才出現.
今天和大家看的這一段十四章一至十四節.
這個段落一個還沒有盼望可言的段落.
可以給我什麼學習呢.
十四章一節你們心裡不要憂愁.
一節四節這裡耶穌在回應彼得前面十三章結尾的一個提問.
小姐們我還有不多的時候與你們同在.

$^{161}$後來你們要找我.
但我所去的地方你們不能到.
誓門彼得就問耶穌主往哪裡去.
耶穌回答說我所去的地方你現在不能跟我去.
後來卻要跟我去.
彼得他很想知道耶穌會去哪裡.
耶穌就回答了這個一至四節的內容.
我們停一停想一想.
彼得他想知道地點.
按常理如果耶穌要回答彼得的話.
就應該回答一個有實質地名的地方.
今天我來這裡講道.
我起碼知道是錦繡堂.
有個實質的地名.
地名也有個地址在背後.
然而耶穌說去傅家裡.
請問傅家裡是哪裡.
經文有沒有解釋傅家裡在哪裡.
沒有.
事實上門徒當時已經很迷茫.
耶穌回答了一個對門徒來說是答非所問的答案.
可想而知門徒是非常不解.
應該回答一個有地方的地名地址什麼都好.
應該是這樣.
但夫子回答傅家裡.
門徒聽完之後真的很不明白.
而的確耶穌在第四節對彼得的回應.
我往哪裡去你們知道哪條路你們也知道.
引來了另一位門徒多瑪在第五節的疑問.
多瑪對他說主啊我們不知道你往哪裡去.
又怎麼知道哪條路呢.
哪條路?哪條路?我們怎麼會知道.
多瑪說耶穌啊我們都不知道地點又怎麼知道路線.
多瑪打開手機Google Map.
他搜尋傅家裡找不到.
即是哪裡呢.
再者耶穌的回應是真的在回答多瑪的問題.
耶穌說我就是道路真理生命.
若不借著我沒有人能到傅那裡去.
你們若認識我也就認識我的傅.

$^{201}$從今以後你們認識他並且已經看見他.
這是耶穌回答多瑪的.
多瑪說主啊我都不知道你去哪裡.
我們又怎麼知道那條路呢.
但主耶穌就回答了這一句給他們.
耶穌你回答的是什麼.
為什麼你不直接把地方告訴我們呢.
我們一群門徒都已經非常迷茫.
你又說要走又不知道去了哪裡.
你又說不再吃這個晚餐.
你還回答了一個我們真的很不明白的答案.
你給個地方名我們知道就行了.
你為什麼要說認識我就認識我的傅.
從今以後又認識他並且已經看見他.
這些是什麼來的.
不明白啊門徒真的不明白.
對於多瑪來說.
他都是覺得耶穌沒有直接去回答他的問題.
而一波未平一波再起.
耶穌給多瑪的回應.
引來另一位門徒肥力的要求.
求主將傅獻給我們看.
我們就知足了.
這麼多疑問不要再問了.
直接一點吧.
獻個傅給我們看吧.
很直接的要求.
對於肥力來說.
他的訴求是很直接的.
獻傅給我們看.
各位這些是思想.
這些是肥力.
你等著耶穌獻傅給你們看.
獻到的.
耶穌你說的是真實的.
你有大能.
而一切如果主耶穌.
不直接將傅獻給我們看的回應.
我們都可能會視之為.
你在這裡兜圈.

$^{241}$轉彎抹角.
不正面回應問題.
你直接獻就可以了.
訴求很直接.
當主耶穌沒有直接去獻的時候.
我們開始心裡面.
就可能會有些疑惑.
有一些潛台詞.
耶穌你直接獻傅給我們看就行了.
不要轉彎抹角.
你獻不到.
那你是不是真的有能力.
耶穌你是不是真的.
是不是.
訴求很直接.
我們要見傅.
但主耶穌你當刻.
如果你不獻的時候.
提問者就好像會有很多疑問在當中.
我們認知的上帝是有大能.
有神蹟奇事.
覺得他講一句話可以創天造地.
那為什麼不滿足我們一點點的要求呢.
我們其實都有這種心態.
我們有沒有試過這樣想.
耶穌啊.
你不醫好我.
你不是神.
因為在你有神蹟奇事.
那就拜託神蹟發生在我身上.
耶穌我們夫妻關係不好.
你又說愛我.
又說會幫我.
那拜託你馬上幫我處理好.
我不管.
你不再相信我的話.
我不相信你了.
耶穌啊.
我和這個男或女朋友的關係.
這段關係你知道我很愛他.

$^{281}$那為什麼他要走呢.
你幫我留住他.
耶穌啊.
我生活很拮据.
事態越來越差.
我都沒有說大魚大肉.
那你為什麼不看顧我.
耶穌啊.
我兒女有很多成長問題.
學業的問題.
你幫我處理好.
耶穌啊.
我父母對我有很多過度的操控.
你幫我讓他們明白.
我也需要空間.
我們還可以有很多.
很大的事情.
放在我們的主耶穌面前.
呼求耶穌.
耶穌.
等於我們要留意.
這種心態.
其實有一定程度的危險.
甚至是高度的危險.
因為我們將一些又大又難的事.
直接抽空了處境.
放在相對上帝大能的角度.
去要求上帝直接滿足我們.
我們再看一次這一句.
我們將一些又大又難的事.
抽空了處境.
有些事可能搞錯了.
可能是你做錯了.
抽空了處境.
上帝你就立刻來幫助我救我.
你抽空了處境.
你再將問題放在相對上帝大能的角度.
這樣一定對的.
我們的上帝這麼厲害.
你這樣相對去求上帝的時候.

$^{321}$要求上帝這樣就直接滿足我們.
我們要這樣向上帝求的時候.
站在上帝的角度.
我們這樣說一定對的.
一定是在上帝沒有大能的事.
沒有難成的事.
一定是大山可以挪開.
小山可以遷移.
一定是曠野開道路.
沙漠開江河.
不過我們要細心留意.
我們的上帝是不是這樣行事.
你看聖經裡面.
聖經給我們認識的上帝.
他是不是這樣行事.
舉一個我們都很熟的例子.
加密山上面.
以利亞先知大戰巴力和亞瑟拉先知.
你想一下.
上帝要對付巴力和亞瑟拉先知的話.
上帝直接擊殺他們就行了.
直接擊殺.
為什麼要差派一個以利亞去做呢.
為什麼上帝不直接自己去做.
要叫一個人去做.
我們的上帝是不是真的會隨時走出來.
抽空我們生活的處境.
然後直接按個鍵.
全部事情都搞定了.
幫你搞定了.
頭髮都換面了.
你看上帝在聖經裡面啟示他自己的時候.
他是差派過多少人去為他完成很多不同的事.
我們只要想一下.
其實上帝是你有大能.
你自己做就行了.
為什麼要叫人去做.
我們的上帝是不是隨時就這樣出來.
一次過按個鍵就搞定了.
總結三位門徒.

$^{361}$他們的問題.
他們的訴求.
彼得說主往哪裡去.
多瑪說主啊我們不知道你往哪裡去.
怎麼知道哪條路呢.
肥力華求主將父獻給我們看.
我們就知足了.
三位門徒的問題有一個共通點.
就是耶穌拜託你這一秒就告訴我們我想知道的答案.
我想知道的你告訴我吧.
大家有沒有真的這樣想過.
耶穌你這麼有大能力.
你這一秒就告訴我.
這一秒就可以幫我解決全部的問題.
神你說一句話可以創天造地.
我又跟隨了你這麼久.
這個小小的要求.
為什麼你不馬上應允我.
這個小小的要求不算過分吧.
暫時我們看1-11節這個講論.
我們可以總括一句.
如果門徒在這個講論的前段.
就已經明白了耶穌所說的話.
從而令他們的憂愁困惑減退的話.
可以說是沒什麼可能.
同樣道理弟兄姊妹.
如果耶穌真的對我們當前面對的艱困.
解釋你聽.
真的這一秒就給你答案.
你猜你的憂愁是不是可以除去呢.
這一秒真的回答你.
是不是真的可以除去呢.
未必的.
因為解釋可能跟你所想的會有出入.
既然都未必是你想聽的東西.
你的憂愁又怎麼會除去呢.
而很可能我們的主已經回應了我們.
只不過是我們不太接受那個答案.
耶穌他回答了門徒.
只不過他們不明白亦難以接受.

$^{401}$又或者我補手一些說.
以他們當時生命的階段.
他們未能夠明白.
的確都好像我們的寫照.
在我們生命當前的階段.
其實有很多東西用心去問的時候.
我們都不太明白.
1至11節原來門徒.
他當刻有很多疑惑不明白.
對門徒來說.
夫子跟他們說這一段話.
到底意義何在呢.
真的要聽夫子說一大段不明白的事情.
整個講論.
耶穌最後的講論.
去到16章尾才在氣氛上有少少好轉.
好正面宣告.
我們就信你了.
在14章仍然一個這麼沉重的段落.
沉重當中.
門徒他們可以做些什麼.
弟兄姊妹在我們人生沉重的路段.
我們還有很多不明白.
未見出路未見曙光的時候.
我們又可以做些什麼.
以下我們就一起去看看.
可以做些什麼.
我們看看我們在人生沉重的階段.
我們可以做些什麼.
而這個亦是我今天分享的最後部分.
我看兩處的經文.
一節的下半節.
你們心裡不要憂愁.
你們信神也當信我.
十一節你們當信我.
我在父裡面.
父在我裡面.
即或不信也當因我所做的事.
信我.
當耶穌回答門徒的時候.

$^{441}$其實他們知道門徒當刻不能夠明白.
就算在過去的日子.
曾經多番解釋給他們.
聽人子必須受害這件事.
子與父緣為一.
這是真理.
但他們未能夠明白過來.
請我們看一看.
這一段記載著這麼多不明白的經文.
經文的開頭和後面.
我們留意一下.
黃色的.
有個信字.
前和後.
耶穌是以信心作為勉勵門徒的一個很大的重點.
耶穌語重心長地跟他們說.
你們信神也當信我.
原來耶穌想表達的意思是.
你們不明白的時候.
一定要對我有信心.
要信我.
要信靠我.
從而去面對你們心中的迷惘困惑.
此刻你們未明白.
但你們一定要信我.
一定要信我.
就算未明白.
都要信.
這一句.
給未信的朋友聽到.
他們可能就會這樣回應.
o下?不明白你都要信?.
你都幾盲目.
請大家留意.
第十一節耶穌是怎麼說的.
你們當信我.
我在父裡面.
父在我裡面.
即或不信.
也當因我所做的事.

$^{481}$信我.
即或不信.
這四個字.
將全句.
分出了兩種信心的狀況.
上半句.
你們留意一下.
你們當信我.
我在父裡面.
父在我裡面.
意思是.
門徒對於子與父原為一.
或者再前一點.
耶穌說.
他要為他們預備地方.
這些說話.
他們未明白.
未能掌握.
未能完全接受.
但他們知道.
說話既出自夫子的口.
一定是正確的.
但卻仍然未能了解.
我既知道夫子說的話一定是正確的.
但我未能明白.
未能完全了解.
他們相信夫子.
見過夫子施行過的神蹟奇事.
聽過他們的教導.
是帶有權柄.
所以出自夫子的說話.
一定是對的.
但我又未明白.
前半句信心的狀況是.
雖然未明白.
但仍要相信.
而下半句.
也當因我所做的事.
信我.
意思是.

$^{521}$過去曾經發生過的事.
鐵證如山的事.
你經歷過的事.
你要因為這些事.
所做過的事.
信我.
下半句顯出信心的狀況是.
有著很多真實的經歷.
到我們否認說沒有發生過的經歷.
上半句.
我們知道是真的.
但未能明白.
下半句就是.
很多很多發生過的事.
要因為這些事.
信我.
夫妻關係愛的心口.
是有過多少經歷.
去建立出來.
心交朋友的關係.
有經歷過多少赴湯蹈火.
孫家弟兄姊妹的關係.
有走過多少同行同心的歲月.
愛同主的關係.
有經歷過多少黑夜前行.
堅心相信去熬煉出來.
真實發生過的事.
建立出來的.
信心是一個互動的旅程.
即是過去神曾經做過的工作.
這些工作可以在我們身上.
或者我們觀察到弟兄姊妹的成長.
或者其他的事情.
很多上帝成就了的事情.
藉此去增添我們的信心.
因為那些事真的發生了.
真的做過了.
這些事的成就.
也都增添了我們的信心.
推動我們繼續有信心.

$^{561}$有耐性.
有耐性迎向那些我們還未能夠明白的事.
亦越能夠去明白一些更加深入難明的事.
我們過往想信心.
通常是一種信心的跳躍.
那樣東西不知道怎麼搞.
沒有信心的時候.
完全倚靠上帝.
信心跳躍全部交託給上帝.
這個你會發覺有點似上半句.
既知道是對的.
但又未能夠完全明白.
你不能夠掌握.
有點似上半句.
信心的跳躍.
但今日這段經文主耶穌說的時候.
讓我們去看見.
除了一種信心的跳躍之外.
第二就是也當因我所做的事.
信我.
建立信心很互動.
前面有些事情你仍然未明白.
沒有信心很缺欠的時候.
你要倚靠上帝.
上半句.
但下半句上帝帶你走過很多路.
完成過很多事.
真實的經歷發生過.
也當因我所做的事信我.
我們的信心已經被主建立了.
過去很多事情已經建立了.
與主走過同行歲月.
這些經歷我們去看.
有沒有走過的路.
你發覺都建立了你的信心.
你有沒有發現.
你的生命已經被主去見證.
原來主耶穌這個段落.
主要想表達的.
沒錯.

$^{601}$祂是想說.
是為我們預備地方去.
沒錯.
是想說.
祂是道路真理生命.
沒錯.
祂是想說.
祂與父園為一這個真理.
帶動門徒去明白.
這些事情的約時.
是信心.
一節實一節.
這裡的關鍵是信心.
整個禮別講論的開始.
受運慘澳的開始.
耶穌是跟門徒說信心.
弟兄姊妹不知道我們有什麼困擾我們.
我們生命這一刻不明白.
在我們此刻的年紀.
我們的身份.
我們的責任.
我們成長的際遇.
一些你面對的難關.
過去發生過了.
但是無疾而終的一些事情.
你仍然未明白的事情.
如果你覺得上述種種.
很困擾著你.
你心中有很多的疑問.
一切上代揭曉之前.
你未明白的.
但請按照耶穌教導所說.
也當因我所做的事信我.
我們人是這樣的.
很不安.
對於明天的盼望不知道.
掌握不到的時候.
信很弱.
上半句我們未明白的事情.
很害怕很擔心.

$^{641}$但不要緊.
信心另一方面是什麼.
也當因我所做的事信我.
走過了很多經歷的時候.
那些真實發生過.
你不能否認.
你不能說那些是假的.
你經歷過.
你被感動過.
從當中你可以找到.
原來你的生命.
你的信心是被建立過.
被建立了.
你一定還有很多事情是不明白的.
但我們是真的不需要一下子.
就覺得上帝我真的不信你了.
不行的.
不需要一下子就推倒建立過的東西.
因為上帝.
他真的按照我們每一個人生命的境況.
去逐步帶領我們.
建立了我們.
又給多一點點難關我們走一走.
再建立我們多一點.
上帝他逐步建立我們.
我們要好像耶穌這個教導這樣說.
越翻所有走過的經歷.
弟兄姊妹我們要細心去數.
越走過一些艱難的路.
除了真的很辛辣.
真的很艱難痛苦之外.
你有沒有發現.
你走過這些路之後.
你堅強了很多.
你成熟了很多.
你長大了很多.
你明白事理多了.
明白人情世故多了.
明白一切是恩典.
明白原來我沒有什麼值得誇口.

$^{681}$我想說當我們細心去數的時候.
我們一定可以數到.
有很多很多是很多太多太多.
上帝做在我們每一個人生命裡面的工作.
你越能夠數到.
按耶穌的教導所說.
我們就越能夠去信靠耶穌.
也當因我所做的事信我.
今天的講道.
我的講題是.
你使我見狀.
耶穌領航我們的人生.
但我們走過恩典的路.
這些路徑.
願我們都能夠從心裡去發現.
原來是父神在我們的生命.
使我們的生命更見狀.
我們都知道我們現在.
可能在迎向一些艱難的日子.
你的信心夠不夠.
前面那些還沒明白.
迎向這些日子的時候.
我們的信心夠不夠.
我們要好好去數.
上帝是愛我們.
他一定是已經建立了我們.
裝備我們.
使我們見狀.
頂字幕再說.
這個段落.
1至14.
是一個沉重的段落.
滿途仍然在迷惘當中.
而或許我們最近.
甚至近一年半載.
光景都是有些暗淡無光.
我們都可能迷惘.
或者是真心對自己的能力有疑惑.
沒什麼信心.
我想我們都看重一個事實.

$^{721}$就是一個在這麼不好的光景之中.
主耶穌仍然對滿途有極大的信任.
縱然耶穌知道彼得要不認他.
猶大要賣他.
滿途將要驚惶失措四散.
但耶穌仍然對他們有極大的信心.
時間關係我們最後看12節.
我實實在在告訴你們.
我所做的事.
信我的人也要做.
並且要做比這更大的事.
因為我往乎那裡去.
沒錯我們都明白.
主耶穌為什麼我們是不明白的.
不明白.
你看主耶穌這裡的講.
這裡宣告.
信我的人是要做一些比我更大的事.
我們不明白主耶穌.
為什麼仍然對滿途有這麼大的信心.
並要做比這更大的事.
更大的事.
我們可以嗎.
做比主耶穌更大的事.
我們的信心這麼微小.
我們可以嗎.
我們現在一臉迷茫.
我們可以嗎.
主耶穌這一番話是極之寶貴的.
你試想一個小朋友.
他可能能力不高.
很多方面都沒有其他人這麼厲害.
好像被人比較的時候都被人比下去.
小朋友自己都覺得可能傷心難過.
這個時候有什麼可以令到小朋友繼續加油努力不放棄.
是父母對他的支持.
是父母對他的欣賞.
是父母稱讚他願意繼續用心嘗試.
是父母對他繼續有信任.
因為有愛我們的主耶穌.

$^{761}$願意這樣去相信我們每一個.
我們的生命才得以被建立上來.
我們是在主的愛裡面被建立.
你使我建壯.
主耶穌對當刻迷茫的滿途說.
你們要做比我更大的事.
同樣我們仍然在軟弱當中懷疑自己的價值.
懷疑自己的能力的時候.
連我們都不太相信自己的時候.
主耶穌相信我們.
父母對子女說.
無論世界怎樣對你呼呼喝喝.
無論世界怎樣標籤你沒用.
但是爸爸媽媽相信你.
你繼續成長的生命.
必定是建立在爸爸媽媽對你的愛和信任之上.
你使我建壯.
弟兄姊妹願我和你都能夠靠主建壯.
一切可能我們感到迷茫困惑的事.
我們的主心知道.
我們繼續正如剛剛上次的經文所說.
有種信心繼續仰望他.
亦數算生命當中走過的道路.
我們的信心是已經被上帝建立.
我們同心祈禱.
因為我們都是一群軟弱的人.
就好像我們今天所看的經文當中.
滿途很多的不明迷茫軟弱.
沒錯我們都有很多軟弱.
很多不明白.
甚至生命當中對你有質疑的時候.
但求你幫助我們.
我們不單止像今天主耶穌的教導.
也當因你曾經在我們生命當中成就的事.
你建立了我們的信心.
你加添了我們的力量.
因你所做的事我們繼續去信靠你.
從經文當中我們看到一個很寶貴.
我們都覺得自己可能很軟弱的時候.
主耶穌你竟然對一群迷茫的門徒.

$^{801}$有這麼大的信任.
要做更大的事.
你對我們有很大的信任.
你的信任今天又再一次激勵我們每一個.
我們可能每一個已經嘗試過.
但已經放棄了不再努力的事.
今天再一次去激勵我們.
繼續在你的愛裡面繼續去成長.
不要輕易放棄.
繼續有一個對你依靠最大的信心去仰望你.
沒錯我們不知道我們生命前面的道路.
將要走的路讓人走.
我們知道你掌管了明天.
你帶領著我們每一個人的生命.
你按我們的需要繼續去指引我們去明白.
求你幫助我們.
我們看見的是.
在走過的歲月裡.
你叫你每一個人都被你去見蹤.
我們為此獻上我們滿心的近人讚美.
我們禱告奉主的名.
阿們.
\newpage



\section{路加福音 11:5-13}
\label{sec:Hdjnt297x1A}
\textbf{2023-09-17 加恩樓《好吧!天父!我信靠你!》路加福音 11\_5-13 講員:蕭建輝牧師}
\newline
\newline
連結: \href{https://youtube.com/watch?v=Hdjnt297x1A}{\texttt{https://youtube.com/watch?v=Hdjnt297x1A}} ~~~~ 語音日期: 2023-09-20
\newline
\newline
\hyperref[sec:Ah5GT26Za2E]{\small{< < < PREV SERMON < < <}}
~
\hyperref[sec:index_chronic]{\small{[返順時目]}}
~
\hyperref[sec:index_scriptual]{\small{[返順卷目]}}
~
\hyperref[sec:tDVwCjoCqSI]{\small{> > > NEXT SERMON > > >}}
\newline
\newline
路加福音 11:5-13
\newline
\begin{longtable}{cl}
\hline
\hline
章節 & 經文 (和合本修訂版)\\
\hline
11:5 & \begin{tabularx}{0.7\textwidth}{X} 耶穌又對他們說：「你們中間誰有一個朋友半夜到他那裡去，對他說：『朋友！請借給我三個餅； \end{tabularx} \\ \\ \relax
11:6 & \begin{tabularx}{0.7\textwidth}{X} 因為我有一個朋友旅途中來到我這裡，我沒有東西招待他。』 \end{tabularx} \\ \\ \relax
11:7 & \begin{tabularx}{0.7\textwidth}{X} 那人在裡面回答：『不要打擾我，門已經關了，孩子們也同我在床上了，我不能起來給你。』 \end{tabularx} \\ \\ \relax
11:8 & \begin{tabularx}{0.7\textwidth}{X} 我告訴你們，雖不因他是朋友起來給他，也會因他不顧面子地直求，起來照他所需要的給他。 \end{tabularx} \\ \\ \relax
11:9 & \begin{tabularx}{0.7\textwidth}{X} 我又告訴你們，祈求，就給你們；尋找，就找到；叩門，就給你們開門。 \end{tabularx} \\ \\ \relax
11:10 & \begin{tabularx}{0.7\textwidth}{X} 因為凡祈求的，就得著；尋找的，就找到；叩門的，就給他開門。 \end{tabularx} \\ \\ \relax
11:11 & \begin{tabularx}{0.7\textwidth}{X} 你們中間作父親的，誰有兒子求魚，反拿蛇當魚給他呢？ \end{tabularx} \\ \\ \relax
11:12 & \begin{tabularx}{0.7\textwidth}{X} 求雞蛋，反給他蠍子呢？ \end{tabularx} \\ \\ \relax
11:13 & \begin{tabularx}{0.7\textwidth}{X} 你們雖然不好，尚且知道拿好東西給兒女，何況天父，他豈不更要把聖靈賜給求他的人嗎？」 \end{tabularx} \\ \\
[1ex]
\hline
\hline
\end{longtable}
$^{1}$Hello各位嘉欣樓的弟兄姊妹平安.
帶來來自南區一小時十五分鐘的車程距離.
在大家當中分享願我們一起被主的說話去激勵.
今天這個講題.
天父我信靠你.
我抄在IG上有一個小朋友對著天空大叫.
It's alright God, it's okay, I trust you.
就是這樣來的講題.
好吧天父我信靠你.
今天想藉著這段經文.
想和大家思考祈禱.
講道也需要大家的聯想和想像.
你要想很多.
你會在講道當中.
我希望我們可以更深的去發現.
我們的禱告一字一句原來都是很重要.
重要到會嚇死你的.
這個講道也希望大家去發現.
我們人生在世.
其實我們和神的關係和父神的關係.
才是重中之重.
其實禱告的內容.
有時候當你發現了和神的關係是怎樣的時候.
你會發覺那些東西不是很重要.
今天和大家思考.
藉著禱告去思考和神的關係.
有個小熱身.
為什麼我說我們禱告的內容.
某程度上當你更加認識和父神的關係的時候.
沒有那麼重要呢.
熱身一,這是主耶穌的教導.
馬太福音.
你們禱告的時候.
不可像那些外邦人用許多重複的話.
他們就以為話多了就必蒙垂聽.
接著再來.
你們不可效法貪玩.
因為你們沒有祈求以先你們所需用的.
你們的父祖而知道.
這句相信是我們認識我們的父的時候.

$^{41}$一個很重要的教導.
我們認識我們的父.
我們的神.
全知的.
祂心知我們.
認識我們的肺腑.
我們母胎未成形的時候.
已經認識我們.
什麼都知道.
所以有時候我們祈禱.
祈禱內容的時候.
我們都一定會有這樣的回馬槍.
刺一刺自己.
是的.
明明父都知道.
那我說是做什麼呢.
會不會的.
有這種小小的.
明明祂都知道.
祂應該自動撥給.
經文所說.
主耶穌的教導說.
那些外邦人以為話多了必蒙垂聽.
你會看到一條公式.
原來外邦人的理解.
總之重重複複地說.
總之我一定要得到.
主耶穌教導我們不要學他們.
為什麼.
因為其實你的父真的知道.
是真的知道.
是不是.
這是第一個一.
第二個熱身.
都是主耶穌的教導.
同樣是馬太福音.
隨後的經文.
所以不要憂慮吃什麼喝什麼.
串什麼.
這都是外邦人所求的.

$^{81}$你們所需用的一切.
你們天父是知道的.
再一次重複說.
你會看到.
吃什麼喝什麼串什麼.
你都不用擔心.
因為父知道的.
是不是.
所以當然我會覺得.
都是主耶穌教導.
主禱文你都會背.
起碼有一句.
是我們日用的飲食.
都賜給我們.
都幸運有一句.
以致我們平時求什麼.
明明我們的父.
這麼豪大.
這麼豐富.
什麼都可以.
賜福給我們的時候.
吃什麼喝什麼串什麼.
不用求了.
或者我們都會說話.
主禱文都有教.
是不是.
日用飲食.
然後又叫我們.
千萬要謹慎.
不單止是.
其實外邦人祈禱的危險性.
都不是重複的話語.
而是他們認為.
必定會這樣.
可以控制上帝去應允他們.
你要捕捉這個位置.
不然有時候.
重複祈禱.
每天我們都重複.
祈禱謝飯祈禱.

$^{121}$那算不算重複呢.
不是這回事.
外邦人的情況就是.
他覺得可以控制上帝.
可以控制那個應允.
重複重複重複.
必蒙誰聽.
這是熱身.
先打個底.
關於祈禱.
主耶穌教導我們.
要避免這些.
也打了個強心針.
我們的上帝.
我們的父.
真的這麼知道我們.
生命裡面的需要.
今天這段經文.
我想進一步和大家細細地思考.
我一開始所說.
我們和父神的關係是怎樣.
《老家方》十一章.
這裡的分段很明顯.
今天想只和大家看五節十三節.
因為一至四節主禱文即背.
你上網聽到.
有些講員也會分享.
其實主禱文逐句也可以說一篇.
所以不可能在這裡掩飾.
所以就看下面.
五至十三節.
五至八節中間是半夜借餅的例子.
然後九至十三節說.
祈求尋找求門.
我們可以看到.
在這三個段落裡面.
有一樣東西很明顯出現.
就是關係.
頭尾包著的.
主禱文說我們在天上的父.

$^{161}$下面九至十三節說.
你們中間作父親的.
是用關係去思考禱告.
有了關係才進入禱告.
中間那段是朋友的關係.
所以三個段落會看到.
在關係裡面展開對話.
頭尾是父.
中間的例子是朋友.
要了解剛才五至八節的例子.
我們要認識一世紀當時的文化背景.
簡單說三個.
第一個就是他們高度重視款待客人.
如果大家有小牧師想去你家坐坐.
我就會說很開心.
你一定要來太好了.
可以招呼你.
我真是幾生手度.
當然我們沒有幾生.
我們一生而已.
高度重視款待客人.
是光榮.
是權利.
可以招呼你.
我就真的非常開心.
第二.
每家每天烤的包.
只是公當天需用.
很正常.
他們沒有雪櫃.
也不是一個很富庶的年代.
所以食物很珍貴.
不會做太多.
如果家裡媽媽負責管理飯食.
你知道家裡人的胃口.
知道兒子不喜歡吃這些.
你不會煮太多.
會浪費.
沒有雪櫃的話.
又會變壞.

$^{201}$所以不會多.
最後就是每個住戶.
家裡一般都不是太大.
明白了這三個背景之後.
我們進入五至八節.
半夜借餅的例子.
讀經我們已經發現了.
總之有一個朋友.
去打擾一個朋友.
誰知道這個朋友沒有東西招呼他.
他又去打擾第三個朋友.
就問他去借東西.
半夜借三個餅.
然後那個朋友就在門打開回應他.
回應有三個不想借的理由.
分別就是半夜沒有餅.
為什麼.
剛才我說了.
每家每戶不會只吃食物.
三更半夜就吃光了.
基本上都已經沒有了.
如果真的這麼幸運有一個.
就先借給你.
第二門已關閉.
剛才說他們家不大.
他們開門不像現代門.
開門沒有聲音.
靜音門.
你推開除非教學位生鏽.
如果不是就沒有聲音.
他們不是的.
開門就很吵.
總之簡單說全屋都會醒了.
不單止醒了.
甚至隔壁屋都聽到.
第三.
孩子已睡.
這個精彩了.
我想做父母的非常清楚.
如果當天晚上你難得能夠安心小朋友睡覺.

$^{241}$你就不要吵醒他.
讓他自己半夜肚餓自然醒.
你不會吵醒他.
醒了之後你再安心他就要搞很久.
又大戰三百回合.
孩子已睡.
所以在門打開的朋友就說.
不要了.
你不要搞我.
有三個理由.
真的幫不了你.
不過.
經文說.
但因他情馳迫切地直求.
必起來照他所需用的給他.
情馳迫切.
當我見到這四個字的時候.
我就想起多年前的一條網上短片.
我們先看一看片段.
請看片段.
(我吃魚翅).
(為何會沒有魚翅).
(我都不明白).
(我都不吃魚翅).
(我要吃魚翅).
(為何會沒有魚翅).
(這麼多魚).
(怎會沒有魚翅).
(你告訴我魚翅在哪裡).
(你又說藏起來).
(我要吃魚翅).
(我特地來吃魚翅).
(我要吃魚翅).
(今次才吃完).
(不讓人吃).
(我要吃魚翅).
(我特地來吃魚翅).
(我要吃魚翅).
好.
大家覺得她夠迫切嗎.

$^{281}$你還能做到這麼情馳迫切嗎.
再來一次我們加強.
是吧.
行不行啊各位.
夠不夠情馳迫切.
Dee王老師你可以表演.
你可以再迫切一點.
那台交給你吧.
是不是.
我馬上想起這條片.
先跟大家做一個小小的調查.
你們是海鮮檔的老闆.
現在有位客人來到.
這麼情馳迫切.
(我要吃魚翅).
各位老闆.
你會不會真的去幫幫他.
幫他舉手.
喂情馳迫切.
情馳迫切.
是客人來的.
你這樣服務態度.
今時今日服務態度不行.
有沒有人可以幫幫忙.
有啊太好了.
如果時間許可.
我可以訪問你.
行不行訪問一下.
為什麼你會想幫他.
因為你是待客服務之道.
什麼.
(我不是服務態度).
(我是廚師).
你是廚師.
但他沒有魚翅.
所以你尊重自己的職業.
所以客人有要求也幫幫他.
非常好.
如果用剛才的說法.
關係.

$^{321}$廚師好像和客人的關係.
也幫他.
但基本上九成九.
這麼多位老闆.
都好像不太幫忙.
OK.
這是一個很實際的想法.
情馳迫切.
我們也要想.
是不是情馳迫切.
那個人在你面前.
情馳迫切.
好像有求必應.
當然現在.
他幸運.
經民港的要求是三個餅.
是不是.
如果.
最近出了iPhone15.
你的子女情馳迫切一下.
各位父母.
是不是.
我要15.
iPhone15.
要不要.
這麼多位父母.
iPhone15.
可能父母就說.
我15個14個給你.
這樣就可以談了.
是不是.
我們想.
很現實地想.
情馳迫切是不是最大.
如果永遠.
我們身邊朋友也好.
或者和上帝的關係.
我們是不是每一樣東西求的時候.
加情馳迫切.
四個字.

$^{361}$再加上一會看了一段經文.
又尋找求見.
叩門.
給你開門.
這樣就發達.
是不是.
當然我們有時候反過來.
我就怕.
是不是我不夠迫切.
是不是要做多點東西再迫切一點.
什麼叫做多點東西再迫切一點.
我們又沒有教我們這些招.
你要想清楚.
到底是否加情馳迫切.
就是那個王牌.
這個字情馳迫切.
在原文裡面.
有那個.
大膽.
勇敢.
刁念.
呂鎮忠就說不顧一切.
和修就說不顧面子.
新易就更厲害.
厚顏.
我們的面皮就薄一點.
這麼厚顏去求好像不太行.
情馳迫切.
是否.
真是我們每次.
去求上帝的時候.
的那個.
格言座右銘.
沒錯.
我覺得裡面有一些很積極的鼓勵.
就是你總之.
大膽.
放心去求.
去向上帝禱告.
但是我們又回馬槍.

$^{401}$剛才馬太福音說.
你的父又知道你要什麼.
哎呀.
有時候.
那怎麼求.
求什麼.
我應該.
我既知道祂又知道我.
然後情馳迫切好像.
好像我應該迫切一下.
但是好像.
你的心裡面又覺得過不了自己.
因為身邊有個人迫切你的時候.
你都不會回應他.
那為什麼上帝又是必要回應我呢.
是不是.
情馳迫切這個字.
在翻譯上面.
這一節經文很精彩.
因為它可以有另外一個譯法.
而那個譯法.
是可以令整句話.
反過來的.
我們看看.
A本身.
這個我用和修.
我告訴你們.
雖不因祂是朋友起來.
但因祂情馳迫切.
然後直求就必起來.
給他.
就是剛才一開始我理解的那個.
B.
我告訴你們.
那人即使不念朋友之情而起來給他.
也會因為不想丟臉而起來.
把他所需要的東西給他.
你看到反轉了.
A.
是求助者情馳迫切.

$^{441}$厚顏去求屋主給餅.
B是屋主不想丟臉.
所以一定要幫你.
你看到嗎.
喂.
那人生選擇了.
大家又想想它.
給三秒大家想想.
到底這個翻譯上面.
嘩玩這麼大.
到底A還是B呢.
A就是那個求助者厚顏.
B就是反過來.
屋主不想丟臉.
所以我也要幫你了.
認為是A的.
請大家舉舉手好不好.
好.
OK.
謝謝.
那認為是B的呢.
我看到手是少一點的.
是少一點的.
剛才是A多一點的.
沒錯了.
這一句到底我們那個人生交叉點.
怎樣抉擇好呢.
到底是A是求助者不想丟臉.
還是那個幫助者.
屋主不想丟臉呢.
關鍵其實在於半夜.
半夜.
你半夜去找一個朋友.
譬如今天晚上.
你晚上打電話給你身邊的大英姐妹.
凌晨兩三點.
我相信你不是只想跟她說晚安.
你有事找她.
你試一下你晚上12點多一點.
有人叮噹按你門.

$^{481}$你立刻在家裡嚇一跳.
看看房東眼.
看看誰在玩什麼.
半夜這個時間已經代表什麼.
代表你去找一個人的時候.
無論.
你那件事一定是大件事.
因為我們做人很正常.
你不會這麼晚去鬼混.
你不會這麼晚去找人.
你這麼晚去找人是很迫切的緣故.
你這麼晚去找人是因為你真的搞不定.
半夜這是一個很關鍵的時間.
的點來的.
真的.
你半夜叮噹.
心裡很害怕很害怕.
打開門 幸好關了.
下一秒你就開始.
什麼事啊 為什麼這麼晚.
這麼晚找我.
你是不是有什麼事.
除非他shock了.
一定有大事找你.
所以.
我看的時候.
翻譯.
是多一點傾向.
B那邊.
屋主.
不單止.
求助者.
他基本上不用太過.
不用那種情緒迫切.
因為你晚上找隔壁的朋友.
已經在說.
那件事是很迫切的事.
這是第一.
第二 為什麼我剛才說要認識.
一世紀文化.

$^{521}$就是那種高度重視款待客人.
的文化.
所以直接說.
人家這麼晚.
都來找你.
你沒有原因.
你沒有理由可以不幫他.
他們很高度重視款待客人.
有人這麼晚來找你.
一定是有很大件事.
所以你事必要幫他.
這是和文化有很強烈的掛勾的情況.
所以我們在翻譯上.
我會多一點傾向.
B那邊.
屋主是很有.
甚至乎那種.
義務一定要幫.
我們生活都有.
譬如你坐公共交通工具.
隔壁可能有個小朋友.
雙腳就插了.
過去對面的椅子.
你心裡也會覺得.
不如也提醒他.
小朋友腳要回來.
乖乖的.
你也會覺得有種義務.
在文化裡.
有種義務要提醒他.
去幫他.
和你在一世紀.
高度重視款待客人的文化裡.
有一種.
有朋友.
還有經文推出.
一個更重要的極致.
是「雖不因朋友」.
就算不是朋友也幫.
意思是.

$^{561}$「雖不因朋友」也會起來.
給他.
我們覺得是朋友就幫.
經文推出一個極致.
不是朋友.
倒幫.
很厲害.
形成文化的元素.
是很強.
甚至家主.
真的要幫他.
本來不想幫.
因為半夜.
可能沒有餅.
小朋友都睡了.
但因為文化裡的.
規範.
文化裡的必要性.
所以最終給予.
求助者.
幫助.
在這一段.
我們可以看到.
耶穌講完一至四節主禱文.
引用這個.
例子想說明什麼.
主禱文裡.
其實求了很多.
他引用這個例子.
說明的時候.
讓我們看到.
「雖不因朋友」都幫.
更何況.
我和大家說.
一開始我們和父的關係.
一開始我們在天上的父.
最下面那段.
中間有做父親.
不是朋友都幫.
更何況我們和父之間.

$^{601}$不是他應該.
更加會回應你嗎.
文化裡不是朋友都幫.
我們和父有關係.
他更加會回應我.
繼續下去.
你覺得發達了.
五至八節.
「雖不因朋友都幫」.
來吧.
祈求就給你們.
尋找就尋見.
叩門就給你們.
開門.
上面上文說的.
「雖不因朋友都幫」.
更何況.
他是我們的父.
還不幫.
還不給.
我們簡單看看.
他所說的中間做父親.
有兒子求餅.
反給他石頭求魚.
反給他蛇求蛋.
反給他蠍子.
看看圖片.
石頭和餅.
外表很相似.
有些像蛇的魚.
又有些像魚的蛇.
很相似.
最後.
我已經盡力尋找.
蛋和蠍子.
哪裡相似.
我還是看《釋經書》.
總之相似.
外貌.
你驟眼看.

$^{641}$很相似.
當然我覺得這個最不相似.
但拼合前面.
石頭和餅.
蛇和魚.
他都想說兩樣東西很相似.
但兩樣東西會有完全相反的效果.
是不是.
中間做父親.
有兒子向你求餅.
反給他石頭求魚.
反給他蛇求蛋.
反給他蠍子.
這一段在承接上.
求就給.
然後他再說.
簡單說藍色字.
我們的父.
是不會玩你的.
你怎會.
他給你餅.
你求餅他給你石頭.
經文很直接.
我們的父不會去玩你的.
你心裡知道.
他知道我們需要.
他知道你緊張什麼.
他知道你祈求什麼.
你不要說出口他都知道.
我們的父.
不會玩你的.
求的東西.
是我們很正常.
也是我們可能.
生命裡面.
一些很正常的東西.
甚至乎.
我們的祈求.
很多時候.
是非常正確的.

$^{681}$是很需要的.
我們有沒有求過以下這些.
我想有一個和諧有愛的家庭關係.
很強健的體魄.
很堅壯的心靈.
不要被人說兩句.
塌了.
我想有很好的弟兄姊妹.
在我身邊.
我想有兩脅插刀的朋友.
我想有聰明智慧.
我想有好的成績.
我想自己是一個.
成熟成長.
有情有義的人.
想不想.
這些很正常很好的東西.
是不是很重要的東西.
我想我們去留意的就是.
很多時候我們祈禱.
求了這些東西.
我們可能會比較少.
去思想.
我們少去思想的是.
那件事由我們祈禱.
到蒙英雲.
會發生什麼事.
由你祈求.
去到蒙英雲的時候.
會發生什麼事.
我們看今天第二條的短片.
好,謝謝,Eddie.
我們不如再看多一次.
再看多一次.
我們的想像,我們思考我們每一句的祈禱.
不單我們要發現,原來我們的上帝真是超級信實的.
祂真是聽你祈禱,如剛才避難所那首歌.
每個祈求,每個禱告,祂要應運你,祂有聽的.
不過,我們要去想多一點,就是我們想祈求.
蒙英韻的事情是要怎樣發生.

$^{721}$你看完這段片,我也是在IG上看的時候.
那一刻真的會心打針.
真的要好好,清醒地去思考自己所獻的每一個祈禱.
我們常常都求上帝給我們聰明智慧.
你會發覺,聰明智慧到底是怎樣得來的.
沒錯,書本上會教你.
但如果謙奉實戰經驗,還是紙上談兵.
原來上帝可能會給你一些難題去經歷.
大家現在家人留在這裡,即將要續租.
可能中秋晚會,你想有個很開心的晚會,怎樣出現的.
求一下主,你給我吧.
那件事,大家都覺得很滿足的罪會是怎樣發生的.
是怎樣發生的.
是不是每個弟兄姊妹都參與在當中.
都很用心在裡面,很想維繫愛裡面的關係.
又祝福到這裡的街坊朋友.
我們祈禱,很想有出現的事情,到底是怎樣發生的.
你很想堅壯一點嗎.
上帝給我們難題,去讓我們堅強.
是的,你覺得人生經過一些風浪之後,你會堅強很多.
我讀神學之前,我都上班了.
都試過被人解僱.
第一次被解僱的時候,我心想被解僱很不開心.
當你轉了兩三次工的時候,你會覺得很小事.
搞錯了,小事拿來說.
是不是,人生經過很多不同的大大小小的難關的時候.
你覺得自己明白了很多東西.
你覺得有些事情,其實不是天塌下來的.
不用去到這麼大.
你走過的路,你祈求很多東西.
我想說,你試一下用這個角度想的時候.
你會發覺上帝不斷在應穩你.
但有時候,可能我們被經文那種,好像這麼直接的.
求餅就應該給餅,直接的事情.
是不是,但我們生命裡求很多東西,不是這麼直接的.
是不是,看回剛才的.
我們要按下去.
這個,和諧有愛的家庭關係.
誰去有那個很溫柔,很願意去愛對方的那種口吻.
那種關心,去相向的對待.

$^{761}$誰去做的.
是不是大家一起.
還是有人喜歡,要先走一步,總有一個.
誰去做的.
你強勁的看拍.
我真的經常看影片,那些舉啞鈴.
做HIIT,看就看很多了.
做就做很少的.
強勁的看拍是怎麼來的.
祈禱就可以有,是不是.
真的.
你要經歷過一點點的歷程才可以有.
堅壯的心靈.
沒錯,可能過往很受別人的話影響.
自己也要去審視,到底我的價值.
在上帝面前.
上帝怎麼看我.
身邊的人,他說的話.
我應該用一個什麼態度,什麼角度去聽.
聽要聽多少.
我自己是否可以.
也肯定自己,但也嘗試聽別人.
那種堅壯的心靈是怎麼來的.
你是不斷反覆思考,又經歷一些對你有傷害的事.
你在當中整整合合地出現.
好的弟兄姊妹.
那種關心,那種愛.
那種彼此的付出.
那種願意大家一起花時間.
去到佳音樓或教會當中.
我們一起建立,我們夢寐以求的教會生活.
也是一定要花時間的.
大家一定要一起結在當中.
否則我們很想有那種教會的熱鬧.
那種愛,那種人情味是不會出現的.
我們祈禱,我們一定是祈這些的.
只不過是祈完之後,那件事又會怎麼發生.
朋友,聰明智慧,剛剛也說了.
馬上害怕,馬上「啊,樹一球,你給我聰明智慧」.
原來是會有困難的.

$^{801}$原來有問題的.
原來有些事情經歷過.
但你真的,很真實的.
你處理過一些傑達的事情.
大問題的事情.
你犯錯過的事情.
你醒過來的人.
你覺得自己在神面前.
真的,警醒了.
有些事情你明白了的時候.
不要重複.
上帝不喜悅.
我試過了.
我盡量不斷不斷要改變,是不是?.
你會學到聰明智慧.
成績.
成熟,成長.
有情有義的人.
有情有義.
都真的是經年累月.
風吹雨打地練回來.
是不能隨便說一句話就可以有.
我們求這些.
我們要可以想很多,思想很多.
還有,你真的要小心一點.
我們的上帝真的很信實的.
很信實.
祂聽到你心裡說什麼.
是不是?.
你的祈求.
祂很信實的.
祂就帶你完成那件事.
應允你.
是不是?.
所以我們看到主耶穌在五字八字說.
父會回應我們.
雖不忍朋友都回應.
更何況有關係.
再去到後面.
就是說這種祈求就會給你們.

$^{841}$那個成就的過程是怎麼樣?.
我們再最後一點點碰到主禱文.
裡面祈求的每一樣.
我們就要想.
願人都尊你的名為聖.
為什麼每一個人都可以尊上帝為聖?.
是因為你尊神為聖.
你將上帝擺在你生活的首位.
你每一件事都不斷操練操練.
可能現在想自己多一點.
但是我開始越來越願意禱告多一點.
交託多一點.
看上帝想我怎樣多一點.
尊上帝為聖.
那些人才知道.
你們基督徒真是尊神為聖.
是不是?.
主耶穌教我們.
如果主禱文之後就開一個新的篇章.
沒有其他東西說.
沒有5-13節這裡說的時候.
那這個經文就是這麼多.
下面我們要思考的就是.
你求上帝會給.
甚至我們的父.
他會帶我們去發現這件事.
蒙英韻的整個過程.
我們要很清楚知道.
我們和父的關係是怎樣.
我們可不可以信得過我們的父.
13節這裡說.
你們尚且不好.
地上的父親可能尚且都未是最完美.
尚且不好.
都知道拿好東西給兒女.
何況天父豈不將聖靈求給他的人嗎?.
這裡是說.
當我們遇上.
就是現在你想求餅.
為什麼好像是石頭.

$^{881}$求魚.
為什麼好像是蛇.
求蛋.
為什麼好像是蠍子.
這個過程.
我們真的需要.
有一個忍耐.
有一個等待.
聖靈.
在我們每一個人生命裡面.
祂要給我們力量.
去走過這個.
好像不是很對應.
我們的祈禱的這些日子.
我們很需要.
很深的.
信靠我們的天父.
對祂有最大的信任.
這個就是為什麼我想說.
當我們發現和神.
父神的關係.
有多少的時候.
我們就開始可以有一點對於.
所求的事情.
當然我們都會照祈禱去求.
但你要發覺那個關係是最重要的.
你沒有了和神的關係.
是什麼意思呢?.
你會看到全部的都是.
石頭.
蛇.
和蠍子.
是的.
你信不過上帝的.
人生一定不是一帆風順的.
有事情初之.
迎面而來的時候.
上帝你做什麼.
為什麼會這樣.
是不是?.

$^{921}$我們人會這樣想的.
但如果我們信不過上帝.
信不過我們的父神的時候.
我們看到的東西.
真的會看到蠍子.
看到石頭.
看到是蛇.
但如果.
我們真的能夠這樣信過上帝.
知道祂一定有祂的帶領.
有祂的安排.
我此刻經過的事情.
求聖靈給我們.
有力量.
有那種心裡面的.
打開的眼睛.
知道上帝.
帶領著我們.
走著一段怎樣的路程.
我們會看到的是.
反而是反過來.
仿似蠍子.
蛋和石頭.
你會看到.
不是.
蠍子.
就是蛇和魚.
蠍子,蛇和石頭.
你仿似.
好像看到一些不喜歡的東西.
但有上帝.
你信得過祂的時候.
你仍然可以經歷.
可能還沒有這麼快有蛋的事情.
但你不至於失卻.
你不至於對上帝不信任.
所以我們很需要祈求.
聖靈不單止讓我們明白很多事情.
我們心裡面.
我們也要建立和我們的父神.

$^{961}$有那種很深的關係.
我再一次說.
這段經文.
讓我們看到的關係.
父.
頭尾包著.
中間這個所謂的朋友.
經文也說.
「雖不恩,朋友都幫」.
我們的父.
聽我們每個祈求.
每個禱告.
求主幫助我們.
我們在祂的神面前.
去禱告的時候.
我們細心留意.
所祈求的每一樣事情.
到底我們的神.
會怎樣帶領我們.
去看到這件事.
蒙英雲.
我們有一點思想的時間.
半夜.
我們的半夜.
你放心.
主是不會不知道的.
祂一定知道.
那種迫切.
我們生命裡面有很多迫切的事情.
是會有很多迫切的事情.
信得過父.
祂知道.
那你可以安心.
第二.
我們的需要.
你看到嗎.
經文所提出的.
都是推進一些很極端的表達.
只是微小如三個餅.
主都理會.

$^{1001}$「雖不恩,朋友都幫」.
這麼微小的三個餅.
更何況.
我們那種半夜.
我們那種迫切.
我們那種覺得.
生命裡面需要.
上帝去帶領我們走過的事情.
父又怎會不理.
五至八節.
以朋友作對比.
剛才已經再說了.
「雖不恩,朋友」.
但都會回應.
更何況我們的天父.
最後我們每一個人都要檢視.
我們對我們的父神.
有多大的信心.
在我們的生活裡面.
我們有多少不斷實踐.
去信靠祂.
無論是讀經.
真的不少.
生命裡面的靈量塑造.
不少的.
生命裡面.
生活裡面.
你有多少東西.
在你交託父.
你真的看他的教導行先.
以致你的行事為人.
是真的按上帝的教導而去行.
我們對父神操練的信心.
是每一天操練的.
所以求主幫助我們.
今天我們藉著這段經文.
願我們都好像一個小孩子.
可以在我們的父神面前.
這樣跟祂說.
「好吧,天父」.

$^{1041}$「我信靠你」.
我們一起祈禱.
愛我們的父.
我們知道.
你就是不會給蠍子,不會給蛇.
不會給石頭我們的父.
你知道我們生命的需要.
你按我們每一個人生命的階段.
去帶領我們走過.
你幫助我們建立.
和你有最深的關係.
以致我們走在每一段不同的路上.
特別是一些艱難的路.
我們仍然可以滿心信靠你.
我們知道.
你是聽到我們生命裡面.
很多有聲無聲的需要.
我們的禱告.
你不會不理.
因為你愛我們的父.
所以我們真的很需要.
對你有那份全面的信靠.
以致我們可能在走一些艱難的路段.
但甚至我們真的可以發現.
原來當中.
都有你給我們的蛋,魚和麵包.
祂再一次去幫助我們.
讓我們發現我們每一句禱告.
你是信實的父神.
你是聽我們祈禱.
我們這樣求.
你就真的帶領我們去應允我們的時候.
我們要去發現的是.
到底你是怎樣帶領我們走過.
你讓我們從聖靈裡面有那種智慧.
生活裡面對你帶領的發言.
以致當我們看見你每一個.
非常微小的帶領.
無微不至的安排的時候.
我們就向你發出感恩.

$^{1081}$啊,你真是愛我們的父.
你沒有不理我們.
你是知道我們生命各種的需要.
我們為此就要向你獻上感恩.
我們多謝你願你的說話.
紮根在我們的生命當中.
我們常常繼續這樣去依靠你.
我們禱告奉償命寇.
阿們.
\newpage



\section{以賽亞書 42:1-9}
\label{sec:tDVwCjoCqSI}
\textbf{2023-09-17 午堂《以賽亞先知的生命質素與使命》以賽亞書 42\_1-9 講員:郭鴻標牧師}
\newline
\newline
連結: \href{https://youtube.com/watch?v=tDVwCjoCqSI}{\texttt{https://youtube.com/watch?v=tDVwCjoCqSI}} ~~~~ 語音日期: 2023-09-20
\newline
\newline
\hyperref[sec:Hdjnt297x1A]{\small{< < < PREV SERMON < < <}}
~
\hyperref[sec:index_chronic]{\small{[返順時目]}}
~
\hyperref[sec:index_scriptual]{\small{[返順卷目]}}
~
\hyperref[sec:_MeYYEhxYZE]{\small{> > > NEXT SERMON > > >}}
\newline
\newline
以賽亞書 42:1-9
\newline
\begin{longtable}{cl}
\hline
\hline
章節 & 經文 (和合本修訂版)\\
\hline
42:1 & \begin{tabularx}{0.7\textwidth}{X} 看哪，我的僕人，我所扶持、所揀選、心裡所喜悅的！我已將我的靈賜給他，他必將公理傳給萬邦。 \end{tabularx} \\ \\ \relax
42:2 & \begin{tabularx}{0.7\textwidth}{X} 他不喧嚷，不揚聲，也不使街上聽見他的聲音。 \end{tabularx} \\ \\ \relax
42:3 & \begin{tabularx}{0.7\textwidth}{X} 壓傷的蘆葦，他不折斷；將殘的燈火，他不吹滅。他憑信實將公理傳開。 \end{tabularx} \\ \\ \relax
42:4 & \begin{tabularx}{0.7\textwidth}{X} 他不灰心，也不喪膽，直到他在地上設立公理；眾海島都等候他的訓誨。 \end{tabularx} \\ \\ \relax
42:5 & \begin{tabularx}{0.7\textwidth}{X} 那創造諸天，鋪張穹蒼，鋪開地與地的出產，賜氣息給地上眾人，賜生命給行走其上之人的神耶和華如此說： \end{tabularx} \\ \\ \relax
42:6 & \begin{tabularx}{0.7\textwidth}{X} 「我－耶和華憑公義召你，要攙扶你的手，保護你，要藉著你與百姓立約，使你成為萬邦之光， \end{tabularx} \\ \\ \relax
42:7 & \begin{tabularx}{0.7\textwidth}{X} 開盲人的眼，領囚犯出監獄，領坐在黑暗中的出地牢。 \end{tabularx} \\ \\ \relax
42:8 & \begin{tabularx}{0.7\textwidth}{X} 我是耶和華，這是我的名；我必不將我的榮耀歸給別神，也不將我所得的頌讚歸給雕刻的偶像。 \end{tabularx} \\ \\ \relax
42:9 & \begin{tabularx}{0.7\textwidth}{X} 看哪，先前的事已經成就，現在我要指明新事，告訴你們尚未發生的事。 \end{tabularx} \\ \\
[1ex]
\hline
\hline
\end{longtable}
$^{1}$我們一起同心祈禱.
我們親愛的天父上帝.
因為你是天地的主.
你留給我們有寶貴的聖經.
你讓我們經常的每個星期去崇拜.
去對你歌頌讚美.
也可以聽到聖經的訊息.
我們就懇求你.
透過這些聖經的話語.
去牧養我們的內心.
讓我們的靈魂得到健壯.
讓我們走在地上做基督徒.
這條人生的路的時候.
我們真的體會到.
你是信實的主.
願你今天也接受我們的敬拜.
你的聖靈在我們當中.
幫助我們去明白聖經的教訓.
祈禱奉耶穌基督的名.
阿們.
我在思想和大家分享.
什麼經文的時候.
我決定用以賽亞書裡面.
有四首的《獨人之歌》.
其中的第一首.
來和大家去作為講道的經文.
這一段的經文就是.
以賽亞書的四十二章一到九節.
我就給它一個題目.
原本不是這個題目的.
就是說.
以賽亞先知的生命質素和使命.
這個是根據這幾節經文的結構.
來發展出來的.
我講道通常都會分段.
這次就可以分段了.
第一個部分.
我要幫大家去明白.
以賽亞書第四十二章一到九節的背景.
裡面是會比較多一些歷史的資料.

$^{41}$第二部分就是要講.
神僕人的生命質素.
以賽亞先知.
他的生命特質是怎麼樣.
第三部分就是要講.
神僕人的使命.
究竟神要他做些什麼.
這樣就可以幫大家.
明白到這段經文.
他的教訓是怎麼樣.
首先我們去看一下.
以賽亞書第四十二章一到九節的背景.
在舊約聖經裡面.
每個人物或者每個地方的名字.
是有含義的.
以賽亞的意思就是耶和華是拯救.
這個表示.
如果你放在以色列的歷史裡面.
你就會發覺是非常之寶貴.
因為牽涉到.
北國我經常都會提.
就是主前的七百二十二年.
亞述是打敗了北國以色列.
主前五百八十六年.
巴比倫打敗了南國的猶大.
以至到王國分裂.
並且猶太有些人精英被擄去巴比倫.
我們經常講.
但以利就是這班很年輕的時候.
已經被擄到巴比倫.
後來波斯皇帝打敗了巴比倫.
就容許這些猶太人.
回去耶路撒冷重建聖殿.
所以這一段歷史我自己比較有興趣.
因為以色列人經歷了起起跌跌.
然後回去重建聖殿的時候.
是引申了很多的教導給我們去學習.
今天我就不會講這麼多.
但是就集中回去講.
以塞俠書四十二章一到九節的經文是怎樣.

$^{81}$我們看看以塞俠書第一章第一節.
是講一些人名的.
就是當烏西亞約坦亞哈斯希西加作猶大王.
剛才講了猶大就是南國的王的時候.
有什麼人出現呢.
就是阿摩斯的兒子以賽亞.
去得到默示.
論到猶大和耶路撒冷.
所以每一個先知都不是講自己的事.
是因為神給了一些訊息他.
所以他得到默示.
他默示的內容是什麼呢.
在以塞俠書一章一節就說.
他是講到南國猶大和耶路撒冷的事.
在這節經文提過幾個王.
這都是和他生存的年份有關.
有些是比較短命的.
烏西亞主前792到740.
約坦是750到732.
亞哈斯735到715.
希西加主前715到686.
剛才我講過.
722年這是重要的年份.
就是亞述打敗北國的時候.
如果是這樣的時候.
他們還沒有去到南國猶大.
被巴比倫打敗.
但是你可以想像.
唇亡齒寒的道理.
北面的兄弟被人打敗了.
什麼時候到南面呢.
在這個處境下.
神就感動先知提出訊息.
以賽亞先知.
主要不是講給北國知道.
是講給南國知道一些訊息.
當時南國還沒有去到很凶險的時候.
但是你可以這樣理解的時候.
神是猜險他去做一些事.
要他去給南國的猶太人.

$^{121}$能夠明白他們不能夠怠慢.
在以賽亞先第一章第八節.
其實是反映主前725年到780年之間.
當時敬拜群體消失.
其實耶路撒冷都失去光彩.
在同時代八世紀.
有其他三位先知.
阿摩斯,尼加和西亞.
你可能都聽過.
他們都不是在耶路撒冷做先知.
他們在以色列和猶大其他地方.
傳達上帝的話.
而以賽亞卻是在耶路撒冷蒙召.
在烏西亞王主前740年死.
所以在那段時間.
以賽亞向南國的人.
傳達神的訊息.
以賽亞是主前八世紀.
一位很重要的先知.
他極可能在主前的739到700年之間.
大概40年去做先知.
你可以想像.
其實我們的上一代的牧者.
常常都有機會慶祝他.
事奉上帝50年.
這個對我們來說是一個夢想.
你可以數一下他幾多歲出來事奉.
所以我們就不同了.
但先知以賽亞他很蒙福.
有40年為神去做事.
接下來我們就看看.
以賽亞書的結構是怎樣.
我要說這部分比較多背景資料.
以賽亞書其實可以分三個部分.
第一部分是第1到39章.
就是主前739到701的背景.
第40到55章反映605到539大被擄的年代.
56到66章反映主前的539到400年.
即是回歸的時代.
特別是主前539到500的歸回年代.

$^{161}$你可以看到.
以賽亞書牽涉的年代是很廣闊的.
不知道你們有沒有讀過以賽亞書.
我們很流行叫人快快將整本聖經讀完.
不過通常讀完都不懂這些東西.
但不要緊.
起碼你有印象知道他講什麼.
但到了我現在這個階段.
我就很想讀得明白讀得通聖經所講的東西.
以賽亞如果他是八世紀前的一位先知.
以賽亞書的內容由主前七世紀到主前六世紀都有.
所以他的前瞻性是非常強.
所以如果你真的很好奇.
很想讀聖經.
我都會鼓勵大家去讀.
以西結書 耶利米書 以西結書等等.
這些大本的先知書是幾十章.
要很有耐性才看得到.
不過其實大家都會看小說.
小說都會不睡覺去看.
一氣呵成.
但看聖經我們何時有這種熱誠呢.
以賽亞書今天我要講這段經文.
我又要讓大家知道上文下理前後的關係.
以賽亞書四十二章一到九節.
其實是擺一個大小小的單元.
每一段的經文都有一個大小小的單元.
還有一個大的單元.
所以逐步來.
它是四十一章二十一到四十二章第九節的總結.
所以你一定要掌握上文下理.
而重點有兩方面.
你一明白它就能入到這段經文.
第一是偶像不能做.
神的僕人就能成就.
第一點是這樣.
第二點很重要是.
以色列抱怨神的僕人安慰.
因為你剛才知道背景就是.
北角已經是打敗了.

$^{201}$亞述已經打敗了北角.
其實很多人是抱怨.
以色列人抱怨.
為什麼神的子民我們都被外族打敗了.
在這個背景下.
以賽亞就去安慰.
所以你大概知道結構是怎樣.
以賽亞書第四十二章一到九節.
它還有一個特點.
它是承先啟後的.
前面是四十一章二十一到二十九節.
但是它開後面的.
就是四十二章十到十二節.
所以有時候我們要明白經文前前後後的東西.
你要有連貫性.
不是斬件去讀.
以賽亞書四十一章二十一到二十九節的重點是什麼呢.
在法庭的情景裡面.
來看外邦偶像的結構.
你可能覺得很奇怪.
為什麼會這樣.
但是你想想.
聖經真的有這些.
譬如約伯記也有.
神說你跟我辯論一下.
看看你對還是我對.
法庭的情景就是用來說一件事.
究竟誰對誰不對.
所以這個是寫作手法.
以賽亞書第四十二章一到九節的重點.
就是神的僕人.
他是給世界充滿痛苦情況的答案.
這個很重要.
第一就是外邦偶像沒什麼用.
第二就是神的僕人帶出來的訊息.
正在回應充滿痛苦情況的世界的答案.
你明白這個很寶貴.
以賽亞書第四十二章十到十二節的重點.
是痛苦的世界因為神僕人的訊息來歌唱.
如果你這樣去看的時候.

$^{241}$你會很明白.
先知他講的訊息不是無聊的訊息.
不是可有可無的訊息.
而是他針對當時的時代一些很重要的問題.
神就給他一些訊息去回應.
所以我們做傳達人去看這些舊日聖經.
都有很多得著的.
舉個例子就是很多先知夢照.
如果你是研究心理學或宗教心理學.
你都可以探討一個課題.
先知夢照有沒有一個共同的pattern.
有沒有一個同樣的式樣.
每個先知摩西被夢照.
約書等等.
你可以一大堆的人物去做研究.
究竟他們有沒有類似的地方.
有沒有不同的地方.
你會知道原來這些叫做宗教經驗.
有沒有共同的東西.
其實這些很有趣的.
如果你是有興趣.
你可以這樣讀聖經.
你看看不同的人.
他的屬靈經歷.
然後做什麼.
看我自己.
我的屬靈經歷是否跟他們可以對比.
又或者如果我要去做.
現在很流行.
我做人家的mentor.
我做人家的生命師傅的時候.
我根據什麼去做.
如果我的認知面不夠闊的時候.
我只知道我的屬靈經歷.
那我就把我的屬靈經歷放大套在你那裡.
但如果我看到聖經裡面不同的人物.
他的屬靈經歷的時候.
我聽你講.
我就知道你這麼暗示彼得.
我就知道.

$^{281}$現在很流行性格分析.
看你幾號仔.
或者是究竟是哪些MPTI.
或者是怎樣的東西.
我覺得讀聖經是很有趣.
我只是這樣說.
希望你們有興趣.
我就覺得很有趣.
這次二賽亞書第41章一到二十一節的背景.
在法庭上控告外邦的偶像.
是一章七節.
一到七節和二十一節到二十九節的重點.
很特別.
就是神有絕對的主權.
但是大家請記住.
以色列人建立了這個神.
是大過所有世界民族的神.
或者神有絕對的主權.
這些信念.
不只是在大衛所羅門的時候.
他最榮華最光輝的時候體會到的.
是當他們被人打.
打到扁了的時候他們體會到的.
大家明不明白我講這句話.
很奇怪的.
神被他們學習了.
外邦的偶像不是什麼.
但是剛剛才講完.
亞述打低了北角.
你還說外邦的偶像不是什麼.
外邦的民族打敗了以色列人.
他們偶像打贏了以色列人的神.
你明不明白.
古代是這樣想的.
但是你反過來講.
外邦的偶像不是什麼.
這個真是很大的信心.
還有他們體會到.
不單止我們的神是有能力.
他還有絕對的主權.

$^{321}$他是大過其他民族的神.
還有他是創天造地的神.
以塞爾蘇第41章的8到20節的重點.
就是神還有一個很大的使命.
他要猜險以色列的聖者.
來救贖以色列.
大家有些短詩.
以色列的聖者.
你明不明白.
你要回到這個觀念的源頭.
在以塞爾蘇裡面.
你就會發覺.
這個很重要的就是說.
那群以色列人是很迷惘的.
而神就是透過先知.
將一些悉切的訊息給他們.
以塞爾蘇第42章1到9節.
是以塞爾蘇第42章10到44章22節的前奏.
所以我又將這個幅度拉大了.
以塞爾蘇第42章1到9節.
現在就真的進入這段經文解釋給大家聽.
剛才給大家一個闊的圖畫.
這段經文我怎樣分段呢.
其實一般聖經學者都是這樣分的.
就是第1到第4節.
第5到第9節.
你一分得對那個節數.
你就完全掌握那個部分的主題是什麼.
1到4節就是介紹神的僕人.
5到9節就是介紹神僕人的工作.
所以我剛才說.
第二第三部分就要處理了.
僕人的生命質素.
其實就是1到4節.
以塞爾蘇第42章1節說.
漢娜我的僕人.
我所扶持所揀選的.
心裡所喜悅的.
我已經將我的靈賜給他.
他必將功利傳給外邦.

$^{361}$漢娜這個詞.
在以塞爾蘇裡面四首的僕人之歌.
都是共同的開場白.
漢娜.
然後就介紹神僕人的工作.
他的獨特性是什麼.
這種介紹的方式.
就好像介紹偉人.
君王 先知很相似.
有一個模式在這裡.
創世記裡面介紹亞伯拉罕.
出埃及介紹摩西.
薩姆爾記下介紹大衛.
都是一樣.
都是有類似的方法在這裡.
神的僕人是神所扶持.
神所揀選的.
是神心裡所喜悅的.
神已經將他的靈賜給他的僕人.
以致到有個結果.
這個僕人不是按自己的能力去事奉.
大家明不明白.
以塞爾蘇僕人之歌第一首裡面.
說這麼多東西.
是要帶出一個訊息.
以賽亞不是憑自己的能力去事奉.
其實他也不少東西.
他是皇室裡面宮廷裡面的一個有份量的人.
他在聖殿裡面蒙召.
你記得禍災禍災那段經文第六章裡面.
他在聖殿裡面的可能是祭司.
可能是聖職人員等等.
其實他也是宮廷裡面有份的.
已經是不簡單.
不是在郊外看養那些.
不是那些的牧羊人.
他是一個不簡單的人.
但是他也會感到自己很不配.
而以塞爾蘇第四十二章二到四節.
是用了七次的筆去形容這個僕人怎樣事奉.

$^{401}$你會問為什麼是這樣.
頭三個不就在四十二章的第二節.
他不喧揚不揚聲.
也不使街上聽見他的聲音.
當然我不會在這裡說這麼多希伯來文.
我們現在很感恩.
我們香港神給我們很多的屬靈的資源.
我們有很多人教舊約聖經教新約聖經.
教舊約聖經的老師一定是懂得希伯來文.
一定可以原文識經.
新約的也是一樣可以用希利文識經.
所以這幾十年.
香港的神學訓練訓練到我們全都要捱過這個痛苦的經歷.
讀完希利文第一期第二期.
我們同學第一年一年級帶著非常衝勁進來神學院.
誰知道我們一個星期不只有一兩次測驗.
他們個個捱了一個學期就有點麻煩.
但是之後他們就會重新開始.
因為他們從來沒有想過原來讀神學要這樣讀法.
但是我們有這個條件我們就要進入原文的世界.
你看看有些故事.
唐三藏都取西經.
他都是從印度繁文翻譯到漢文.
你意志到那個宗教才可以在文化裡面生根.
今天我們不做這些事.
你有什麼結果.
你講一大輪你去別處都可以.
但是你沒有經文的準確性.
所以我不是要特意帶一些希伯來文出來嚇你.
我真的講之前做完功夫.
軒陽的意思其實就是向神呼求.
陽星其實有另外一個含義.
就是向神呼求的時候抗議.
神啊你為什麼這樣.
為什麼會這樣提呢.
因為他們經歷了很大的挫敗.
所以不在街上發聲.
不是說不要在公共地方講話.
可能解釋有兩個.
第一個就是你不要在街上講這麼多哀聲.

$^{441}$他們剛才講了北角以色列經歷了失敗.
第二就是神的僕人不要再以悲哀.
在舊式聖經有很多哀歌.
你不要再以一種受害者的心態去繼續事奉.
你不要這樣.
你明白我的意思嗎.
在當時實在是這樣的氣氛.
所以在《二賽亞書》第42章裡面.
我們可以這樣理解.
這幾個筆開頭是這樣的意思.
《二賽亞書》第42章第3節說.
阿商的爐位他不截斷.
張壇的燈火他不催眠.
他憑真實將功利傳開.
這裡你會看到神的僕人有什麼.
他經歷了痛苦和被抵擋.
如果你去受難節通常是用《二賽亞書》第53章.
他的僕人是沒有街營美容.
即是不好看的.
你會發覺這個僕人是經歷了很大的艱難.
你可以委屈也行.
或者他的失敗也行.
在這段經文裡面沒有形容神的僕人是屬靈精英.
不是精銳部隊.
他不是因此過人.
反而他好像一個弱者.
他不再年輕.
不過神卻願意扶持他的僕人.
大家看到《二賽亞書》第42章是說這些.
而阿商的爐位他不截斷.
張壇的燈火他不催眠.
神的僕人可能是飽力滄桑的人.
或者逆流大叔那些.
不過他首先是經歷了神的更新復興.
不管你有多少的沮喪.
在人生裡面有多少的挫敗.
神是會將你整理好.
《二賽亞書》第42章第4句說.
他不灰心.
另一個筆.

$^{481}$他不喪膽.
直到他在地上設立公理.
海島都等候他的吩咐.
我逐句理解給大家聽.
雖然面對著困難.
但神的僕人他不會灰心.
不會喪膽.
一個捱過痛苦的人.
他其實有承受力增加了.
他可以負重量增加.
不過一個受過苦捱過苦的人.
如果他的生命沒有更新的時候.
他很可能是將苦毒帶給其他人.
有時我們經常聽.
這個人的內心很多bitterness.
為什麼.
因為他沒有辦法去處理.
他受傷的經驗.
正如一個受過傷的人.
可以成為一個受傷的治療者.
Henry Lowen 路運神父.
寫了一本書叫受傷的治療者.
都很值得大家看.
受過傷的人.
你怎樣可以成為有同理心.
你走過你自己的幽谷.
可以幫人體諒人.
不過如果你沒有經歷過生命的轉變.
你可以成為.
我自己弄出來的觀念.
受傷的殺手.
Wounded Killer.
即是怎樣.
我都捱過苦.
你都應該捱下.
這個故事是甚麼.
我們看乳腸片.
希望大家不要覺得對女性有甚麼偏見.
即是一些惡家姑的故事.
即是我年輕的時候嫁入豪門.

$^{521}$我五,六時起床已經在拆地板.
到她做奶奶的時候.
就會和未來媳婦說.
我當年其實是五時起床做家務.
不過現在我沒有這樣要求你.
你明不明我這樣.
即是這些就是故事.
我都等了這麼多十年.
我才可以出頭.
你明不明.
即是我不捽你.
捽得久一點.
怎樣過癮.
你明不明.
很多人的內心就是.
他沒有辦法處理他自己.
人生的鬱結.
被人捽.
他現在怎樣可以寬容對人.
你明不明我意思.
所以這個是很重要.
你看《二賽亞書》裡面一個很重要的課題.
這個先知他捱過苦.
但他並沒有讓這種苦毒留在他內心裡面.
亦都不會用這種態度對人.
我覺得很有啟發性.
希望我們回去可以默想一下.
在《二賽亞書》裡面.
42章三到四節有文學上的對比.
A,B和Bπ和Aπ.
和我這麼上一年級的人.
看到A,B,B,A又有另一個聯想.
我就不說那些了.
結構是這樣的.
壓傷這個字的字根.
Razus.
張殘的字根.
Kahar.
不灰心的字根.
Kahar.

$^{561}$不喪膽的字根.
Razus.
你發覺原來很有趣.
如果你知道這些東西.
你讀希伯來文的時候.
這麼有文學的韻律.
是呀.
我再舉多一個例子給你聽.
人在希伯來文叫什麼.
Adam.
土地是什麼.
Adamah.
你再想想.
經文就說上帝用泥土中的塵土做人.
我先把塵土抽起.
上帝用Adamah做Adam.
多有押韻.
我只想告訴大家.
如果你有機會讀一下.
聖經的希伯來文.
希臘文.
你會很享受.
很簡單.
大家知不知道伯特利.
這裡有間伯特利中學.
伯特利的希伯來文是什麼.
或者是學校的名字.
Beth L.
Beth就是家.
L就是最高的神.
神的家.
永生神的家.
我鼓勵你們.
名字是有解的.
你們為子女改名的時候.
就看看希伯來文說什麼.
神僕人其實有四種很特別的性格.
很有趣.
你留心看.
他既慈愛又堅定.

$^{601}$既謙和又堅強.
你想想這四樣東西合得來.
很不容易.
為什麼他可以培養到這些性格特質.
我回答這個問題之前.
我們要搞清楚這個僕人是誰.
這樣才可以明白得多些.
大部分聖經學者說.
以賽亞書的40到48章裡面所說的僕人.
就是以色列.
究竟以色列怎樣被培養到這些性格.
你看回經文就是.
經歷苦難.
被神正念.
大家消不消化到這幾個字.
當人經歷苦難.
他回到神的面前的時候.
他就有機會看到自己的盲點.
心靈的黑暗面.
他可以嘗試靠著神的恩典.
來面對自己的失敗.
以色列不獲亡了.
知衰了.
就回到神那裡.
當我們回到神的面前的時候.
是怎樣.
我們就好像由浪子.
變成神的兒子.
以前流浪的.
現在不是了.
以賽亞書第42章1到4節.
就反映了一個皇室的處境.
同時是有先知被差遣的內容.
不過這個在皇室處境.
受先知差遣的神的僕人.
竟然就是曾經失敗了的神的子民.
神是差遣曾經失敗的人.
去為他做事.
你看這段經文多獨特.
當以色列民經歷失敗之後.

$^{641}$他回到神的面前.
神就重新接納他們.
以色列民現在成為神的僕人.
所以僕人之歌其實是在說整個以色列.
神的僕人能夠被神扶持.
被神去管理.
他才可以去事奉神.
去配合神的計劃.
奇妙的就是說.
還有另外一件事.
海島都等候他的訓誨.
訓誨是什麼意思.
就好像指律法.
神就預備列國來期待學習神的律法心.
神的僕人就要配合神的計劃.
我剛才說給你聽.
亞述打敗了北角.
我打贏你我為什麼要聽你的話.
你有多少東西要學.
巴比倫打贏.
遲些會打贏南國猶大.
我為什麼要聽你的話.
但是神就猜險他的僕人.
要預備有這個心態.
你要將他的律法.
向以色列甚至外面的國家宣講.
所以你看到這裡的時候.
你發覺是很奇妙的.
以賽斯第42章一節有三次提到公理.
希伯來文.
42章一節.
他必將公理傳給外邦.
第三節.
他憑真實將公理傳開.
第四節.
直到他在地上設立公理.
很特別這件事.
而以賽斯第42章一節和三節.
是用傳開和建立來說的.
整個意思就是.

$^{681}$神的僕人要傳開建立公理.
在全地.
最重要是全地.
以賽斯被呼召.
他不只是為以色列人說話.
或者猶大民族說話.
不只是為耶路撒冷說話.
他要有一個使命.
從今日我們新約的角度來說.
他有普世宣教的使命.
你明不明.
在當時已經有的了.
神的僕人被呼召在全地.
傳開和建立公理.
神已經預備了列國.
去期待去學習神的律法.
這個才是很重要的地方.
以賽斯第42章一到二十一節的背景.
就是法庭裡面控告外邦的偶像.
面對外邦力量的控制的以色列.
以賽亞領受的呼召不是去報復.
不是去仇愛.
不是去排斥其他民族.
而是讓神的公理在全地傳開和建立.
如果你看到這一點的時候.
你就會很明白.
就是說.
神的僕人的眼光.
不只是集中在以色列.
而是普世的民族.
很獨特的.
以賽亞需要獨特的地方就是說.
預言猶大會被懲罰.
猶大會亡的.
被擄的.
但被擄之後.
他會得到回歸神的懷抱.
然後被更新.
領受一個很重要的使命.
你要成為萬國的光.

$^{721}$大家看不看到.
以賽亞書裡面很重要的.
以色列成為萬國的光.
但你要記住.
他是衰了之後才可以做到萬國的光.
你看多奇妙.
以賽亞先知的信息不是排外.
不是不斷宣傳給敵人攻擊.
所以我要國仇家恨.
而是說.
猶大會亡國.
亡國之後會再生.
他就會成為萬國的光.
你如果在當時聽這些信息.
你都不是很容易接受.
人類的歷史經常出現.
我當年是受害者.
現在我一朝得志.
我就要加害你.
推廣報仇.
繼續敵視.
這些的信息是社會上有的.
但你看到以賽亞書不是說這些.
如果我們明白世界歷史.
起碼知道世界歷史.
你就應該明白.
聖經裡面以賽亞先知的信息.
絕對不是一般道人向善的宗教道理.
絕對不是.
是神自己的說話.
我很希望要說這麼多.
解清楚它.
就讓你明白你信什麼.
第三.
我真的會過鐘.
不好意思.
曾牧師兄.
神僕人的使命.
剛才他說.
你可以說四個鐘.

$^{761}$看人會不會走.
我說好就試試.
當然我不會這樣做.
第三個就是神僕人的使命.
第五節.
你看到筆調是不同了.
創造諸天.
鋪張窮倉.
將地和地所出的一拼.
鋪開.
賜氣息給地上眾人.
又賜靈性給行在其上的人.
那位神爺和華.
他這樣說.
說了一大堆東西去描述我們神.
但最主要是說後面在說什麼.
二賽亞書42章1到4節.
就是說神僕人的東西.
第5到9節就是說.
神對著神僕人說什麼.
以至他執行什麼使命.
他要有命令.
不是自己作的.
第5節說.
創天造地賜人肉身生命.
賜人靈性生命.
那位神.
神爺和華這個詞.
只在二賽亞書出現一次.
成本教學聖經都不見.
有其他的.
第6節就說.
我爺和華憑公義去呼召你.
是攙扶你的手.
保守你.
使你作眾民的中寶.
作外邦人的光.
這裡我看完都很感動.
眾民和外邦人可以譯做萬國.
這兩句話是平行句.

$^{801}$眾民和萬國是對照的.
神的僕人.
以至要成為萬國的祝福.
你看到很重要.
好像以阿伯利罕成為萬國的祝福.
在舊約聖經裡面.
公理.
mispart.
和公義.
ceterqua.
經常一起出現.
二賽亞書42章1到4節提到公理.
但去到第6節的時候.
你看到出現公義.
所以這些很重要的基本常識.
神按公義去呼召神的僕人.
在全地傳開和建立公理.
這個就是舊約裡面先知要做的事.
第6節.
使你作眾民的忠報.
外邦人的光.
神的僕人的使命.
是要做眾民的忠報.
和做外邦人的光.
但是你要知道背景就是.
在二賽亞書第6章第10節.
是怎樣描述以色列這個民族.
這班人是眼睛昏迷的.
現在.
即是眼蒙蒙.
是黑眼的.
現在他要成為外邦人的光.
是否有很大的對比.
即是我以前是一個不信主的人.
我都是烏下烏下的.
現在我要做人家的光.
你明不明.
自己都覺得自己做不到.
二賽亞就要傳達這個消息.
神就將這番話向二賽亞說.

$^{841}$你就做這個角色.
而這個角色亦都是代表.
整個以色列民族要做的角色.
當神的計劃開始的時候.
人就要被更新.
如果他不被更新是做不到的.
二賽亞書第42章第7節.
他就說什麼.
要開核子的眼.
領被囚的出牢獄.
領坐黑暗的出監牢.
你想想.
曾經核眼的人.
以色列民.
現在去帶領核眼的人走出黑暗.
你是不是很諷刺.
但不是很稀奇.
不稀奇.
很多人他以前是做了錯誤.
他悔改之後.
他就在這班人當中.
來去做一些工作.
我年輕的時候在戶外團契實習.
那些過來人就戒了之後.
就信了耶穌.
然後就回去這班人中間.
來去服事.
而二賽亞書這裡是說.
曾經核眼的人.
現在走去帶領核眼的人出黑暗.
是不是很有趣.
二賽亞書第42章8節.
我是耶和華.
這個是我的名字.
我必不會將我的榮耀給其他的假神.
也不會將我的稱讚給雕刻的偶像.
你看看說完這些之後.
就後面包住了.
誰說這番話.
是耶和華說這番話.

$^{881}$我是不會將這些偶像當成什麼.
面對著外邦偶像敬拜的挑戰.
臣他自己去榮耀他的名.
不過他需要以賽亞和以色列民族.
所以第9節說.
漢啊 先前的事已經成就了.
現在我將新的事告訴你.
這些事是未發生的.
我就告訴你.
所以你看到.
以賽亞才是說一些有希望的東西出來.
以前的事成就了.
前面還有事.
所以他才有這個任務.
去告訴那班以色列人.
當然今天我們有什麼新的事跟大家說.
這個就請大家為我們每一個.
事奉神的人去祈禱.
我們怎樣去讀聖經.
怎樣去祈禱.
怎樣領受神的訊息.
這些經文是強調偶像不知道未來.
只有誰知道.
只有神知道.
這些經文是對應以賽亞書第41章21-29節.
偶像不知道未來.
只有神.
所以當我們去思考現在的處境的時候.
你可以思想誰知道未來.
只有神知道未來.
所以以賽亞書第42章1-9節.
是以賽亞書第41章21-42章9節的總結.
有兩個重點.
我已經走向結束的部分.
第一就是偶像做不了什麼.
只有神的僕人做到.
第二就是以賽亞書第41章1-42章9節.
那些以色列人抱怨.
但是神的僕人是安慰.
神的僕人的使命.

$^{921}$是向全地傳開和建立神的公理.
今天或者我多說一句.
你和我有沒有珍惜時間.
你將把握時候去做更多的事.
在以賽亞書第42章1-9節的經文.
是沒有很具體講到神僕人的工作.
沒有一個job description給我們.
沒有的.
神僕人的使命是向全地傳開建立神的公理.
但是具體的東西.
神的公理具體的內容是什麼.
經文沒有講.
全地的範圍有多大.
經文也沒有講.
究竟從哪裡開始.
也沒有講.
傳開的方法是什麼.
也沒有講.
這麼多東西都不講.
是給了什麼我們.
給了一個方向我們.
我耶和華是憑公義召你.
必參附你的手.
保守你.
使你作眾民的中保.
作外邦人的光.
這一節經文.
以賽亞書第42章第6節是什麼.
是神的應許.
這段經文沒有講得很細緻.
這個job list.
以賽亞做什麼沒有很清楚.
但是他告訴我們.
這個是神的應許.
你就跟著做.
以賽亞才知道蒙召是要怎樣.
他要抓緊神的應許.
神僕人蒙召.
有一樣很重要.
他抓緊神的應許.

$^{961}$神的訊息.
他就去做.
神僕人經歷壓傷的蘆葦.
他不截斷張禪的燈火.
他不催眠.
他學會謙卑.
但是有進取心.
要為神改變局面.
所以你這樣看的時候.
你會對以賽亞.
是有一個.
或者希望你對他有多一點.
或者深入一點的理解.
他不灰心不喪膽.
因為他知道困難重重.
但是仍然不放棄.
所以這個是很值得我們學習的.
總結那裡.
我都不說那麼多了.
因為已經差不多了.
我們有一個簡單的祈禱.
我好希望.
大家.
你不記得我說什麼都無所謂.
我已經去到這個境界了.
最重要你回去看聖經.
這樣就行了.
你記得這個世界裡面.
在舊世紀裡面有個先知以賽亞.
很值得你去學.
這樣我就已經很成功了.
你明不明白.
人的記憶力很短.
但最重要就是榮耀歸給上帝.
我們一起同心祈禱.
我們親愛的天父上帝.
因為你是歷世歷代的主.
以色列這個民族.
在他的興衰裡面.
是經歷到你的同在.

$^{1001}$你是呼召你的先知.
你的祭司.
你的領袖.
去做你的事.
去安慰去教導當時你的子民.
我們的神.
我們看到.
他們都不一定是屬靈的精英.
他們不一定是偉人.
他們可能曾經失敗過.
但是你是應許過.
你是願意使用他們.
是你自己改變這一班人的生命.
是你自己去榮耀你自己的命.
我們今日活在一個比較複雜.
比較混亂的世代.
我們求你.
首先讓我們內心肯定.
讓我們看到.
我們不要浪費時間.
我們要把握機會.
將耶穌基督的福音.
向周圍去傳播.
不單止是香港.
甚至全世界.
因為這個是你的應許.
求你讓我們恢復這種勇敢和志氣.
讓我們深信.
我們的神是全能的神.
我們恭敬的祈禱.
是奉耶穌基督的命.
阿們.
\newpage



\section{撒迦利亞書}
\label{sec:_MeYYEhxYZE}
\textbf{2023-09-21 周四晚查經班《歸回先知書研讀(2)撒迦利亞書》 講員:曾裔貴牧師}
\newline
\newline
連結: \href{https://youtube.com/watch?v=-MeYYEhxYZE}{\texttt{https://youtube.com/watch?v=-MeYYEhxYZE}} ~~~~ 語音日期: 2023-09-23
\newline
\newline
\hyperref[sec:tDVwCjoCqSI]{\small{< < < PREV SERMON < < <}}
~
\hyperref[sec:index_chronic]{\small{[返順時目]}}
~
\hyperref[sec:index_scriptual]{\small{[返順卷目]}}
~
\hyperref[sec:B_Q_tzD4YvY]{\small{> > > NEXT SERMON > > >}}
\newline
\newline
$^{1}$好了,我想我們差不多可以一起簡單有一個祈禱先.
再一次多謝天父都給我們能夠可以又順利地過了一個星期.
又來到星期四晚,我們可以一起在家中,在教會.
來到去聚集,去聆聽一同的查考聖經的話語.
我們特別有機會去看被擄的歸回的時代.
先知的預言,很盼望藉著撒迦利亞八個的預言.
和往後有的經文都給我們對當時的世代.
當時的神的默示都有認識.
特別我們是在一個回望的時間,我們已經有耶穌的救恩.
我們已經是新約後代的人,盼望我們去追述這些歷史的時候.
都給我們有啟發,有學習,求主帶領我們奉耶穌基督的名,阿們.
再一次多謝各位,也有新朋友加入我們的查經班.
因為大家的程度不同,所以我們講到歸回的歷史書和歸回的先知.
相當複雜,因為大家能夠看的機會很少,未必很熟.
所以大家不要介意不厭其煩,將一些背景都跟大家重溫.
我們看看在以斯拉記第四章的經文.
這就是第一次他們重建聖殿的時候.
慶高采烈,所羅巴伯和大祭司約書亞帶著那些被老歸回的百姓.
其實總共只有五萬多人,所以不是多數,其實是很少數.
這樣來到上去,其實已經是某程度排除萬難.
也是一個新的時代,就是波斯王鼓勵大帝允許所有被老的人.
可以回到自己的原居地重建.
所以猶太人就會被神的靈激動,領取了這個聖旨.
上去耶路撒冷,輾轉在以斯拉記記載了他們上去的.
無論是人口,家譜,帶的物資等等.
但是今晚我們要看的是撒迦利亞書,特別是我們看第三個遺像.
就更加要回去看這些歷史,有些背景,大家有一個觸覺.
第四章的以斯拉記提到他們開始面對相當大的.
那些已經駐紮在耶路撒冷猶太遍地的原居民,即是那些外邦人.
就會阻撓你們上來,說這個地是你們自己的,是你們猶太人的.
已經不可以了,因為由亞述,由巴比倫已經是亡了國.
被露了去外邦,那怎麼會還是你的呢?.
所以這就是主權的問題,第四章讀多一點給大家有個印象.
第四章第一節,以斯拉記,由大和變眼敏的敵人聽說.
被擄歸回的人為耶和華以色列的神建造殿宇.
就去見所羅巴伯和以色列的族長,對他們說.
請容我們與你們一同建造,因為我們尋求你們的神,與你們一樣.
自從亞述王以撒哈頓帶我們上這地以來,我們常常都是祭祀神的.
但是所羅巴伯就意識到,這些是外邦人,他們根本那個心都不是歸向神的.
只是想保住自己的地圖,利益等等,所以他們只是假意的來奉承.

$^{41}$所以他們就,所羅巴伯就很有這個警覺性.
但所羅巴伯耶穌就是大祭司,等下我們第三章會看.
和其餘的以色列族長對他們說,我們建造神的殿與你們無關.
我們自己為耶和華以色列的神協力建造,是照波斯王古列所吩咐.
就是那一位古列大帝,第四節.
那地的民就在猶大人建造的時候,使他們的手法軟,擾亂他們.
可能是危言聳聽,可能是威逼利誘,或者告官.
或者在那裡製造輿論等等,令到他們的軍心和民心不能夠合一.
第五節就這樣描述,從波斯王古列年間,直到波斯王大利烏登基的時候.
大利烏就是允許他們再一次重建,即是說有大概十四至十五年.
聖殿就爛了尾,就不能再重建.
即是有一些外力令到猶太人,所羅巴伯他們都不能夠重建.
甚至乎他們在這個期間,就賄賣一些謀士,即是可能是一些辯護.
法庭的一些律師等等,要敗壞他們的謀算,即是告官等等.
所以就令到他們不能夠再建造,所以那個工程就停頓了.
四章的廿四節,於是耶路撒冷神殿的工程就停止了.
一直停到波斯王大利烏的第二年.
這個就是那個歷史給大家有一個小小的了解.
即是說如果去到大利烏王的時間,第二年,大概主前的519或者518.
視乎是年頭還是年尾,都重要的.
因為去到我們看這些先知書的時候.
哈該和撒迦利亞正正就是在大利烏王的第二年.
開始在百姓當中來講預言.
所以就勉勵他們要來奮勇地重建神的殿.
即是說將一個已經爛了十幾年的聖殿已經荒廢.
十幾年的沒有人出入,又不能夠每個禮拜來到去有崇拜敬拜神.
不可以燒香獻祭,十幾年荒廢.
其實那些樓宇本身可能都會有結構問題的.
但當然可能他們只是立了根基,其他什麼都做不了.
好像那些鋼筋只是外露,什麼都做不了.
這裡有一個很重要的時代,我們要拿捏的.
哈該書我們之前都講了有四篇的信息.
其實是講道來的,都是要他們重新在神的面前.
透過哈該先知給他們的勉勵.
要認真地上山去取木料去建造,很興高采烈.
所以當那個工程完成,他們就好像當時耶路撒冷都有很大的歡樂之聲.
一直聽到遠處.
我們今天看的撒迦利亞書,上個禮拜我們很快地看了兩個異象.
頭七章有八個異象,所以都不難去分辨的.
第八章是另外一個信息,是向列國的信息.

$^{81}$所以先知的角色就是在一些很關鍵,很非常的時候.
有神給的預言,他們要麼就沒有,有時很多.
有時停很久都沒有,這個是很視乎神給他們的時間.
和預言的時間很重要.
所以在這個時候,特別是他們要重建需要的是信心,需要的是鼓勵.
需要再一次認定耶路撒冷到底在神的眼中定位是怎樣呢?.
他們都需要去有一個了解.
我們上個禮拜有兩個異象,就是一個騎紅馬的.
有一隊騎兵在他身邊好像聽候差遣.
叫他們在全地走來走去,當然是一些預言,是一些象徵.
即是說這個騎紅馬的身份是非比尋常的.
他可以去求問神,為耶路撒冷求平安.
神會安慰他,這些我們上個禮拜講過了.
這個人我們都提過了,就是主耶穌曾經在馬太福音提過.
祂求平安,為耶路撒冷.
所以有時好像一些很一脈相承的,舊約有這樣的預言.
但是主耶穌真是本著祂是神的兒子.
看到這個城很快就要成為荒場,荒廢的一個廢墟.
來去求,即是在未發生之前已經有這種不只是惻隱之心.
而是一種預見耶路撒冷因為敵黨和釘這位彌賽亞在十字架.
就要面對神的公義的審判.
所以我們看到第一個異象就是好像凸顯了這位騎熊馬的天使.
其實如果我們新約,我們就知道是主耶穌.
所以撒迦利亞書其實有很多彌賽亞的預言.
第四章,第三章都有,大衛的描銳已經出現了,三章八節.
我們等一下應該有機會可以看到.
第二個異象就是講四國和有四個象人.
似乎四國是指著全地,即是那個方位東南西北.
一直操控著耶路撒冷,來打散對於猶大和以色列的這些國.
似乎是因為古代的人都會相信不同的地方都會有守護天使.
所以就是守護著整個神選民的地方以色列.
但是現在給我們看到有四個象人出現.
象人特別強調他們打散了之後會重建那個城.
這個好像給我們看到原來耶路撒冷有那個保護.
有那個重建的工程,很快就會恢復昔日的榮耀的工作.
或者第二章的哈該書,翻前一頁就有了.
第二章的哈該書,我就給大家看一下.
當他們開始這個工程的時候,那個帶來的應驗的興奮.
第二章的哈該書,尤其是第四節開始.
就是神跟索羅巴伯說,他是那個公頭,是那個總工程師.

$^{121}$是他策劃回歸帶領百姓帶領國民從波斯回到耶路撒冷.
他是一個省長,很自然就是他回來的目的就是要召集所有的猶太人.
一起來建造聖殿,但是好景不常.
面對的就是當地人的阻撓,又結合暗中的搞擾.
又明明的來告官,所以令到他們的建造突然間要停頓.
但是現在經過哈該的勉勵和預言,他們在這個時候就正式要建造這個破爛了的聖殿.
所以這個時候最需要信心鼓勵的就是這位領袖索羅巴伯.
神就跟他說,耶和華說索羅巴伯,雖然如此.
你當剛強,若撒達的兒子大祭司約書亞,這個在剩下的撒迦利亞書再會出現的.
你也當剛強,這地的百姓你們都當剛強作供,說了三次剛強.
因為我與你們同在,有神的同在,當然是藉著先知的預言.
因為當時他們都沒有聖殿,他們都沒有獻祭,更加沒有大祭司來有聖殿裡面的功績.
無論是宰殺那些的祭僧,燒香,或者有其他的祈禱等等.
所以現在需要的就是這個政治上的領袖索羅巴伯和還沒有真正履行職位的大祭司約書亞.
現在有先知更加超然的積分預言,是直接由神那裡而來的.
來安慰幾位的領袖,要他們在這個時候首先要剛強.
百姓未剛強,你首先要剛強,因為你的剛強就會帶來其他的百姓都會有一個勇氣.
一齊來去建造工程,跟著就告訴我們神這樣跟他們講.
第五節,這是照著你們出埃及,我與你們立約的話,那時我的靈住在你們中間,你們不要懼怕.
萬君之耶和華如此說,過不多時我必再一次震動天地.
這個是引出埃及記的經文.
滄海與旱地,我必震動萬國,萬國的珍寶必都運來.
我使到這殿滿了榮耀,這是萬君之耶和華說的.
很重要就是神能夠震動萬國,需要的材料,需要的榮耀.
神都可以一下子叫不同的人,或者不同的時代,不同的演變,這樣來運來.
這個殿之前曾經因為受神的刑罰,變成方牆,要被擄到巴比倫.
甚至乎神聖殿的器皿,公職的,獻祭給神的器皿,全部都被擄到巴比倫.
但是現在輾轉,所羅巴伯就帶領回來了耶路撒冷.
但是現在當他們要重建的時候,他們一定會有百感交集.
因為所羅門的第一個聖殿是非常宏偉的.
現在他們又要面對這麼多的難處,人心也都渙散.
還有其實回來的是五萬多人,所以其實是一個很小的工程.
你可以想像得到那個建造是很艱難.
出來的效果都未必好像跟過去那麼大的規模.
一定是會有了被比下去的感覺.
但是神在這個時候安慰他們,好像不再在乎那個大小.
在乎的是這一班人是從被擄的地方回來.
這一班人是帶著感動回來重建.
還有更重要的是神應許再住在這個民族當中.
很快來看之前,和大家揭一揭去以西結書.

$^{161}$一段很標誌性的經文,說到神要離開這個地方的奇怪之處.
第十章的以西結書,這些經文大家真的比較少去看.
其實很形象化,以西結書就是被擄的時期.
他們是要去巴比倫,在一個叫加巴魯河邊.
以西結在這個外邦的地方見到異象.
即是你走百萬行,可能那時候是我走過一次.
在九龍塘走去石泥背水塘.
那時候是剛剛開的那條高和路街.
一直走到大窩坪,一直再下去.
然後走到石泥背水塘.
你走到一些很荒蕪的地方,突然間見到異象.
以西結就是這樣,其實整個行列是被擄的.
他見到一些特別的異象,你見到你都會覺得很心痛.
所以你讀一下這些經文,如果你是猶太人.
顯出藍寶石的形狀,彷彿寶座的形象.
主對那穿世麻衣的人說:你進去.
在旋轉的輪內,基魯柏以下.
從基魯柏中間將火炭取滿兩手,撒在城上.
我就見他進去,那人進去的時候.
基魯柏站在殿的右邊,雲彩充滿了內院.
耶和華的榮耀從基魯柏那裡上升.
停在門槳或者門欄以上.
殿內滿了雲彩,院與也被耶和華的榮耀光輝充滿.
尤其留意整張講耶和華的榮耀.
一路一路由聖殿,由基魯柏天使的光輝.
即是約衛的光輝,一路向外延伸.
即是去到人人都能見到的地方.
第五節:基魯柏的翅膀響聲聽到外院.
好像全能神說話的聲音.
他吩咐那穿世麻衣的人說.
要從旋轉的輪內基魯柏中間取火.
那人就進去站在一個輪子旁邊.
由一個基魯柏從基魯柏身手到基魯柏中間的火那裡.
取些火放在那穿世麻衣的人兩手中.
兩隻手拿著炭火.
你明顯見到這是一個儀仗.
不是真是物質的捧著火.
那人就拿出去了.
在基魯柏翅膀之下顯出有人手的樣式.
見基魯柏旁邊有四個輪子.

$^{201}$這基魯柏旁有一個輪子.
那基魯柏旁又有一個輪子.
每基魯柏都是如此.
輪子的顏色彷彿水倉玉.
至於四輪的形狀都是一個樣式.
彷彿輪中套輪.
即是說那是一個移動的聖殿的躍跪.
那個基魯柏是移動的.
可以行的.
有四個輪子.
輪中再套輪.
第十一節:輪行走的時候.
向四方都能直行.
並不調轉.
頭向何方?.
他們也隨向何方?.
行走的時候並不調轉.
他們全身連背帶手和翅膀並輪.
周圍都滿了眼睛.
我們用物質去想就是怪物.
但是告訴我們.
他不單止是直行橫行.
都是可以這樣來去的.
但是他每一個方位.
任何的角度都不會看兩眼.
現在我看著大家.
我就看不到後面.
因為我的眼睛不會關在後面.
但是他所有的方位.
任何的角度都能看到.
即是遍察全地.
這個基魯柏.
他不是想告訴我們一個怪物.
在聖殿裡面有自己的異形.
不是.
是說造物主.
他的存在.
他的臨在.
透過基魯柏我們能看見.
其實看不到.

$^{241}$因為聖殿裡面是不准進入的.
大祭司都要裡面點滿香.
整間聖殿的致聖所.
充滿了煙雲.
你才可以進去彈血.
彈完就要走.
那裡只能夠是約會.
在那裡.
這就告訴我們.
這樣的描述.
即是說展現一個悲哀.
那個被擄去到一個形象化.
就是聖殿最重要的基魯柏約會.
他自己一路走.
那個走的意味是神的榮耀.
離開以息列.
再讀下去.
第十二節就說.
「換了眼睛.
這四個基魯柏的輪子都是如此.
至於這些輪子.
我耳中聽見說是旋轉的.
基魯柏各有四個簾.
即是四活物.
第一個是基魯柏的簾是人的簾.
第二就是人的簾.
第三就是獅子.
第四是鷹.
基魯柏升上去了.
這是我在加巴魯河邊所見的活物.
基魯柏行走.
輪也在旁邊行走.
基魯柏展開翅膀.
離地上升.
輪也不轉離他們旁邊.
那些站著.
這些也站著.
那些上升.
這些也同上升.
因為活物的靈在輪中.

$^{281}$十八節.
耶和華的榮耀.
從殿的門欄那裡出去.
停在基魯柏耳上.
基魯柏出去的時候就展開翅膀.
在我眼前離地上升.
見得到.
連先知本來不可以進入聖殿的.
他也見到這個異象.
當然可能只是他一個人見到.
大家都在避露.
他見到基魯柏.
原本應該在聖殿.
原本應該沒有人可以見到.
原本是有那個萬紫遮蔽的.
但是現在.
他還是這樣.
一路離開.
還要是眾目睽睽地見到他.
一路這樣徐徐上升.
一路強調一路見到.
接著十八節就說.
停在基魯柏耳上.
基魯柏出去的時候就展開翅膀.
在我眼前離地上升.
倫也在他們的旁邊.
都停在耶和華殿的東門口.
東門口.
聖殿的門口.
就是那個入口.
東面的.
朝東的.
就是那個入口.
一進去就見到一個大的祭壇.
現在.
神的聖殿最重要的那個藥櫃.
那個基魯柏.
其實最重要就是那個榮耀.
神的榮光.
一路一路.

$^{321}$連二世傑在加巴魯河邊都看到.
那你可想而知.
糟糕了.
等於.
任何一個拜偶像的地方.
那些觀音.
黃大仙.
在那個垃圾池.
那是很大的事情.
你想想.
是多麼的羞辱.
多麼的對猶太人的侮辱.
所以這裡是.
告訴我們.
第十九節.
下半節.
在他們耳上有以色列神的榮耀.
強調強調.
是神的榮耀.
在基魯柏耳上.
基魯柏原本是安放在藥櫃之上的.
現在已經是去到.
大家都看得到.
如果你能夠可以看到這個異象.
你就可以看得到.
那二世傑就是看到這個異象.
那就是.
好像告訴他們.
事出有因.
不是巴比倫.
不是那些軍事強大.
令到猶太人現在.
行衰運.
被擄.
其實不是.
根本的原因是.
神的離開.
不再和這個聖殿的選民同在.
祂要走了.
神的離開.

$^{361}$第二十節.
這是我在加巴魯河邊所見.
以色列神榮耀以下的活物.
我就知道他們是基魯柏.
各有四個臉面.
四個翅膀.
翅膀以下有人手的樣式.
至於他們臉的模樣.
並身體的形象.
是我從前在加巴魯河邊所看見的.
他們驅國直往前行.
成章聖經.
只是很想說.
耶和的榮耀在基魯柏耳上.
基魯柏隨著自己離開.
意味著一個這樣的審判.
已經沒有什麼可以說了.
神都要離開這個聖殿.
所以已經沒有了靈魂的聖殿.
很快巴比倫拆毀整個聖殿.
甚至猶太人.
全部被擄到巴比倫.
當然留下國內的老弱.
在那裡都沒有人要.
這就是被擄的時候的光景.
這樣的一個突顯.
現在回到撒迦利亞書.
就很不同了.
現在是神的榮耀.
神的同在.
神的安慰.
要所羅巴伯.
要大祭司約書亞.
要剛強.
百姓要剛強.
現在要重建.
為什麼呢.
因為我與你們同在了.
這是很重要.
對猶太人來說.

$^{401}$這是最大的安慰和保障.
所以如果將這些先知書.
不同時代的先知書.
一併來看的時候.
你就會看到最重要的地方.
我們又回到第二章.
我們有了一點背景之後.
這個四象人.
大家要留意.
即是神都會有他的保守在其中.
讓耶路撒冷得到重建.
很形象化.
即是我已經預備有工人.
來協助重建耶路撒冷.
去到第二章第三個異象.
更加特別.
有那些Sawade.
那些測量師.
現在是用電子的.
一射就很遠.
都可以在某個位置.
多少呎都可以數到.
以前就要靠拉繩.
古代肯定沒有這些儀器.
一定不會那麼準.
但他們一定會找到那個辦法.
我們看看這個異象.
是想帶出什麼訊息呢.
似乎一直都在說.
從外圍去看耶路撒冷第二章.
見一人手拿準繩.
Sawade去測量.
我說你往那裡去.
很特別.
在異象裡面的先知.
可以直接去問天使.
他對我說.
要去量耶路撒冷.
看又多寬多長.
說笑嗎.

$^{441}$一個人怎麼去量這麼大的耶路撒冷呢.
所以一定是一個異象.
就是說要來留意重見的意思.
與我說話的天使.
果然是天使來的.
去的時候.
又有一位天使迎著他來.
對他說.
就是說另外一個天使.
B天使.
A天使就回答了.
撒迦利亞我去哪裡.
我去量耶路撒冷.
多寬多長.
準備興功建造了.
但是有另外一個天使迎著他.
跟他說.
第四節對他說.
你跑去告訴那少年人說.
到底這個少年人是不是撒迦利亞呢.
因為有另外一個人.
耶路撒冷必有人居住.
好像無城牆的鄉村.
因為人民和生畜甚多.
耶和華說.
我要作耶路撒冷四圍的火城.
並要作其中的榮耀.
好特別.
火和榮光都在這裡出現.
一方面就是不用量.
量不到.
因為將會很多人居住.
不要說笑.
現在只有五萬多人.
在這個時候.
這個異象的時間.
都是大理烏王的第二年.
九月的時候.
所以這麼少人量度有城牆.
是一個保護.

$^{481}$你說不需要有城牆.
這個當然是說將來.
就好像告訴我們.
耶路撒冷會恢復它的榮耀.
恢復它的繁榮.
因為多人多生畜在其中.
你又可以想像到.
這個量省的異象.
好像很搞笑.
很形象化.
又好像給我們一個圖畫.
好像舞台劇.
有一個人很刻意地.
準備去量度耶路撒冷.
是天使.
然後有另外一個天使迎著他.
量度什麼呢?.
將來都會很多人.
然後就說神的語言.
在第五節.
耶和華這樣說.
所以我們就要留心.
在舊約.
這些都是一些形象.
這些都是特別提到.
神是烈火.
神在荊棘火焰中.
向摩西顯現.
出埃及的第三章.
來呼召他.
給他有一個信心.
他當年的那種.
對自己的同胞的關愛.
原來是一份感動.
不過那份感動.
未夠時候.
未夠成熟.
未夠火候.
就失敗了.
但是.

$^{521}$在四十年後.
原來他的內心.
好像心如止水.
成為一個牧羊人之後.
神就向他顯現.
好像一個已經.
沒有什麼雄圖大志的牧羊人.
一直都是在他的外父.
葉剔羅的手下.
來去安分守己.
做一個牧羊人.
以前的雄心.
以前的訓練.
全部都用不到了.
看羊就不用這麼複雜.
但是突然間見到.
經濟燒著了.
但是燒不壞.
很好奇.
一個在四十年.
做看羊的.
帶著羊到處吃草.
見到野外.
一點都不是.
特別尋常.
但是這次的尋常就是.
為什麼那些經濟.
正常的經濟.
為什麼會燒不壞呢.
一直燒 但燒不壞呢.
走過去看.
突然間有聲音跟他說.
在荊棘火叢當中.
有聲音跟他說.
摩西 摩西.
他就很害怕 跪下.
你不要過來 因為你所站的地方.
是聖地.
神的同在就是聖地.
那是曠野.

$^{561}$那是一個很荒蕪的地方.
燒著的火.
是很平常的.
山火.
秋天就一定會有.
冬天就更加會有.
重陽節 因為放孔明燈.
尤其是那些.
祖墳.
上山.
不想斬.
不想放火燒.
這些是很正常的事.
但是摩西感覺到的.
就是神在火裡面.
向他顯現.
所以這個火.
都是很特別的一個.
儀仗.
所以現在告訴我們.
神成為四圍的火城.
神在其中.
有榮耀.
我們讀一讀.
《出埃及記》的24章.
都是特別提到火.
所以來到這裡.
神這樣講對猶太人來說.
一點都不陌生.
因為他們的祖宗.
本身已經是在火裡面.
被神呼召.
24章的17節.
15節.
摩西上山.
有雲彩把山遮蓋.
就是上西瀨山.
耶和華的榮耀.
亭於西瀨山.
雲彩遮蓋山六天.

$^{601}$第七天.
他從雲中照摩西.
耶和華的榮耀在山頂上.
在以色列人眼前.
形象如烈火.
所以你看到.
火和榮耀.
是掛勾的.
舊約很早期已經.
這樣描述神自己的.
顯現.
那個榮耀.
當然.
火就不是神.
我們不是在拜一個火神.
我們不是.
造物主是烈火.
造物主是榮耀.
造物主是光.
造物主是愛.
但是愛不是神.
神是愛.
不可以反過來.
神就是光.
這個約翰一書告訴我們.
一章的五節.
你就可以想像到.
猶太人對於.
神的榮耀.
他們有這些.
形象化的概念.
所以他們一聽到.
撒迦利亞這樣說.
神這樣預言.
即是神再一次在我們的民族當中.
再一次來呼召我們.
因為這個時間.
二十四章的《出埃及記》.
正正就是摩西.
和西乃山.

$^{641}$神和民族立約.
十誡的頒佈.
和之後不同的誡命.
所以現在.
撒迦利亞再有.
這樣的火和榮耀.
你就可以想像到.
神再眷顧.
這個民族.
是一個肯定的訊息.
不是一個.
先知的一廂情願.
他自己.
很有民族的感情.
不單止是這樣.
是在神的.
旨意下.
來講述.
神的預言.
所以舊約.
耶穌還沒來之前.
好像要勾起猶太人.
對於神的出現.
有什麼.
所謂有跡可尋.
是什麼形式呢.
火是其中一個形式.
榮光.
就是其中一個形式.
到了新約.
也有講到神是烈火的.
有沒有想過.
在哪裡.
希伯來書的十一章.
對不起.
應該是十二章才對.
十二章就是.
勉勵在苦難裡面.
希伯來的信徒.
要忘記.

$^{681}$背後.
等待主耶穌的再回來.
不要因為.
現在的小小苦楚.
就去氣餒.
甚至乎.
神的管教.
是對.
每一個兒女的一個.
很重要的訓練.
所以.
下垂的手.
發酸的腿就要挺起來.
也要為自己的腳.
把道路修直.
使道缺止.
不至於歪腳.
反得全愚.
很特別的.
就這樣解就解不了.
為自己的腳修直道路.
為了自己的腳不要踩得太過.
我們明白了.
修好道路.
有時候有時候.
暴雨沖壞了.
很多坑洞.
你走下去會扭傷.
很合理.
修直道路.
關那些瘸腿的什麼事.
瘸腿就不會因為你的洞.
不瘸.
瘸腿就是先天或後天.
已經成為事實.
怎樣瘸子都能夠.
得到全愚.
瘸子就不至歪腳.
瘸了腳歪就歪.
不會直.

$^{721}$但作者說不是.
你為自己的腳.
修直道路.
你自己潔身自愛.
你自己循規蹈矩.
行神的道.
成為一條道.
成為一條路.
你行真理.
其實我們終於明白.
其實你在家裡好過見證.
信主.
他們見到父母那麼敬虔.
他們慢慢.
不喜歡回教會.
一路長大.
發現原來這個信仰是很真.
他們的歪腳.
小時候都百艷.
慢慢被分逃.
他們都會有一個.
對信仰的堅持和認識.
原來是這樣理解.
原來為自己的腳.
要修直道路.
是我們的責任.
但修直之後.
不單止你走路.
不會扭傷.
是令到那些強子.
扭傷的腳.
都可以.
痊癒.
原來聖經是那麼奇妙.
所以一個人信主.
全家得救.
某程度都是真理.
因為慢慢就會改變到.
不喜歡回教會.
但這個需要時間.

$^{761}$需要耐心.
需要忍耐.
所以不是一個必然的.
好像一定是這樣.
而是過程有很多的眼淚.
很多的堅持.
很多的放棄.
放棄你的自我中心.
願意來到去.
事事都以別人.
做優先.
好像回教會.
禮拜.
好像有一個簡單的協議.
除了禮拜我要敬拜神.
其餘時間.
你想我怎樣都可以.
這樣就很不同了.
對那個未信的家人.
他就會明白.
原來禮拜你不是故意不陪我.
你更加重要.
就是你敬拜神.
這樣的原則.
其實正正就是把道路修直.
使那些強子.
不至於歪腳.
他有時都未明白.
甚麼叫價值觀 甚麼叫信神.
如果你都給不到他.
很明顯的真理.
和堅持的時候.
他就不太知道.
到底甚麼叫信神.
但是.
如果他慢慢地看到的時候.
他們就會有心靈的改變.
我重點現場.
是看下面.
第25節這一大段.

$^{801}$25到29節.
你們總要謹慎.
不可棄絕那向你們說話的.
因為那些棄絕在地上.
警戒他們的.
尚且不能逃罪.
何況我們違背.
那從天上警戒我們呢.
當時他的聲音震動了地.
但如今他應許說.
再一次我不但要震動地.
還要震動天.
這是黑哥書2章.
剛才我們讀過的經文.
再一次的說.
指明被震動的.
就是受造之物.
都要羅去.
使那些不被震動的.
存在 常存在.
所以我們既然.
得了一個不能震動的國.
就是將來的永恆的天家.
就要當感恩.
照神所喜悅的.
用虔誠敬畏的心.
去事奉神.
為甚麼呢.
因為我們的神.
是烈火.
再一次告訴我們.
是有一個將來.
永恆的審判.
我們是.
不被震動的神.
的選民.
因為救恩已經將我們.
帶進去將來的天家.
所以我們是不被震動.
但在有生之年.

$^{841}$我們.
是要常存這個信仰.
所以真的要記住.
要懂得背.
一次相信.
是要一生一直的相信.
一直的堅持行神的真理.
你才可以得到永生.
不是一次得救.
就是永遠得救.
一次得救 永遠得救的中間.
是有一生的相信.
一直的堅持.
你才可以必然的得救.
所以我們既然.
得了這個不能震動的國.
就當感恩.
這是每一個基督徒.
信了主之後.
即時的一生的表現.
就是一個感恩的心靈.
讓這份感恩.
不斷的滋生.
不斷的散播.
然後全著虔誠.
敬畏神的心.
這樣來生活.
事奉.
為甚麼呢.
因為我們的神是烈火.
回到剛才的廿四章.
出埃及記的經文.
那個簡選.
那個呼召.
神的榮耀.
即是說.
這個烈火是包含著.
將來的審判.
那個潔淨.
在我們今生.

$^{881}$我們就要帶著這份虔誠敬畏.
對神.
所以回看這裡.
那個意思就更加的立體.
更加的玲瓏浮凸.
告訴我們.
原來神.
他現在是做耶路撒冷的火城.
他的榮耀.
充滿四維.
即是說.
又是和那個.
建造耶路撒冷有關.
所以這八個意象.
其實很簡單.
就是全部圍繞著.
耶路撒冷來說.
耶路撒冷在這個時候.
大尼烏皇帝二年.
正式的.
爛尾聖殿重建的開始.
所以.
很多這些意象.
都是針對著.
這個耶路撒冷的.
大小的工程.
現在神應許.
他自己成為那個火城.
即是說這一位的主.
再簡選.
也都願意.
他的同在.
在耶路撒冷.
他的榮耀會充滿.
這個聖殿.
再被建立的聖殿.
第二個聖殿.
就一直去到主耶穌的時代.
主耶穌離世之後.
主後的七十年.

$^{921}$那個年份就很準確了.
就是羅馬人.
完全打敗了猶太人.
他們就要分散.
去到天下.
聖殿就被焚燒了.
整個聖殿就會夷為平地.
這個就是.
聖殿今天的光景.
大家去到耶路撒冷.
那個哭場.
應該是西場.
就是所羅門那個聖殿的遺址.
唯一未塌的地方.
那真是很大塊的石.
所以很多那些石瀨.
猶太人就在那裡祈禱.
一直在聖殿.
這樣叩頭.
就是將那些祈禱的經文.
他們的願望.
塞進那些瀨裡.
希望神去垂聽.
還有他們一直這樣叩頭.
是等候彌賽亞降臨.
他們等候的是.
他們的彌賽亞.
他們所以不能夠接受.
我們基督徒說的主耶穌.
是他們的彌賽亞.
他們一直這樣等待.
所以真是所託非人.
主耶穌已經離過這個世界.
以及已經回到天上.
來為人類預備居所.
所以猶太人到今天.
大部分都是不信主.
是很難的一件事.
但是我們回頭看的時候.
當這個聖殿重建.

$^{961}$神就應許.
它成為耶路撒冷四圍的火城.
不是去燒這個地方.
而是保護這個地方.
它的榮耀.
是會同在.
火和榮耀.
永遠都是一起出現的在舊約.
並要作其中的榮耀.
有神的火.
有神的管教.
有神的榮耀.
即是意味著.
有神的看顧和同在.
所以教會.
其實是真是需要有愛.
更加要有真理.
不可以兩者.
只是有真理.
公義又沒有了愛.
是不可以的.
這就是第三個意象.
接著有些說話說.
耶和說.
我從前分散你們.
去到天的四方.
即是任何的方位.
東南西北.
現在你們要從北方之地逃回.
特別強調北方.
其實巴比倫.
是東面的.
大家知道.
稍微有些地理背景.
東面是一個大沙漠.
所以他們不會走過大沙漠.
橫行過來.
未到達都死了.
所以一定是繞個圈.
在上面.

$^{1001}$因為是沿著河流.
有法拉底河.
和底格里斯河.
就是兩條大河.
所以一定是沿著河.
那些海岸線.
河岸線.
這樣進軍.
北邊的陸地.
騎馬下來.
所以說北方.
是對的.
不是東方.
因為不會走過大沙漠.
什麼飛機坦克車都不行.
你可以想像得到.
這個就是.
北方.
到下面.
這是耶和說的.
與巴比倫人同住的石安民.
應當逃脫.
萬軍之耶和說.
在顯出榮耀之後.
差遣我去.
懲罰那擄掠你們的列國.
摸你們的.
就是摸他眼中的銅人.
那個Apple.
就是那個.
我們的銅孔.
就是說.
是非常敏感的.
就是誰都碰不到的.
不可以碰的.
現在告訴我們.
神看顧這個民族.
看耶路撒冷.
是他眼中的銅人.
這些經文在哪裡出現過呢.

$^{1041}$例如.
《雅歌》也有.
當然我們不看那裡.
第九節.
我要向他們.
他們就必作服事.
他們之人的老民.
你們便知道.
萬軍之耶和說.
差遣我了.
撒迦利亞特別強調.
他這個差遣.
第十節.
錫安城啊.
應當歡樂歌唱.
因為我來要住在你們中間.
這是耶和說的.
那時必有許多國.
歸附耶和說.
作他的子民.
他要住在你們中間.
你就知道萬軍之耶和說.
差遣我到你那裡去了.
耶和說.
必收回猶大.
作他聖地的份.
也必再揀選耶路撒冷.
特別強調.
猶大是聖地.
舊約不是很多.
這樣講.
這樣的概念.
聖地似乎.
這裡明顯的一次.
說到猶大.
這個地方.
為什麼叫聖地呢.
就是因為神的再揀選.
和神的.
同在在其中.

$^{1081}$凡是有血氣的.
都當在耶和說面前.
靜默無聲.
因為他興起.
從聖所出來了.
這一段聖經是.
接著那個.
測量員.
那個天使.
去測量耶路撒冷.
那個的.
長闊.
意味著準備大興土木.
但是有另一位天使.
就好像叫他.
不需要了.
多此一舉.
因為將來會有很多很多人居住.
甚至乎.
不需要城牆.
為什麼呢.
因為神就是那個城牆.
而且是火城.
而且是神的榮耀.
一提到城牆.
提到這個聖殿.
所以我們都.
有些感覺這一段經文.
不只是當時的.
那個的.
是的.
初步的應驗.
就是他們會開始一個重建.
都是有一個.
很奮勇地去興建.
但是說著說著.
似乎不只是.
當時代.
那個的處境.
可以涵蓋了.

$^{1121}$可能我們看看啟示錄.
我們有些印象.
22章.
相信是啟示錄裡面最寶貴.
最重要的地方.
或者我們先看了.
21章的經文.
尤其是去到梁鐸的時候.
15節.
21章的15節.
就是講將來的新天新地.
一直描述.
那個的聖城.
然後講它的尺寸.
和那個的.
不同的物料.
每一道門.
那個城牆的根基.
是怎樣的.
諸如此類.
重要就是去到22節.
我未見城內有殿.
因主神全能者.
和羔羊.
為城的殿.
那城內.
猶不用日月光照.
因有神的榮耀光照.
神就是那個殿.
造物主和羔羊.
是城的殿.
殿即是猶太人.
他們.
任何一個地方.
都需要有.
當然他們不敢任何一個地方有聖殿.
他們可以有會堂.
敬拜神的地方.
聚集的地方.
來去中毒那些.

$^{1161}$Torah的地方.
但是既然去到耶路撒冷.
才知道.
這個就是約翰.
他要找的就是.
這個城這麼宏偉.
這麼奇特.
那些珍寶.
但是一切都不及.
整個聖城.
最重要的.
就是有聖殿.
所以他就要找聖殿.
找不到.
但是原來神就是那個殿.
我們就在神的裡面去敬拜.
在羔羊的裡面去敬拜.
就是這個意思.
不需要再有一個物質的殿.
這些將來.
我們離開這個世界.
進入那個永恆.
就不解自明.
因為到時就會完全你會認識.
所以現在我們只能夠.
很象徵性地了解.
將來那個大概的概念.
和那個規模.
但是一定是一個永恆.
裡面我們就不需要.
再在這裡左問右問.
這一段.
為什麼就不需要了.
那你就可以想像得到.
這個就是.
我們回到看撒迦利亞書的時候.
我們就要有這個觸覺.
因為我們是新約.
當時他們一定是沒有這個啟示錄.
我們就.

$^{1201}$讓他們可以知道.
更多神.
陸續啟示他自己的.
這個人類的.
救恩.
那個進程.
我們在2000年後.
耶穌降生離世升天的2000年後.
我們再看回這些舊約經文.
我們就憤外感覺到.
那個局部應驗.
和全面的.
應驗的那個可貴之處.
是有一種這樣的.
局部局部.
所以當我們再看的時候.
首先就是.
在巴比倫同住的石安文.
這裡就好像.
有一些語帶相關.
固然是.
其實巴比倫已經亡了很久.
亞述.
巴比倫.
接著就馬代波斯.
接著就到波斯帝國.
所以波斯帝國是滅了巴比倫的.
已經不存在了.
那為什麼還要講這個概念呢.
你住在.
北方的巴比倫的石安文.
強調他們是石安文.
即是神的選民.
但是住在巴比倫.
那意味著不是.
物質的地理的描述.
是靈性的描述.
即是說.
你們是神的選民.
現在耶路撒冷有一個.

$^{1241}$很大的榮耀.
快要開始了.
為什麼你們還要在北方.
在巴比倫.
這樣來去同流合污呢.
那個城是要滅亡的.
很多這些經文.
譬如.
以賽亞書又有這些經文.
提到在巴比倫.
那些要快點走.
第52章.
以賽亞書.
希望大家不要介意.
揭這麼多經文.
但是一提到巴比倫.
巴比倫其實就不是在講.
那個物質的巴比倫.
其實已經亡了國了.
等一下我們來看看.
這個啟示的十字架.
等一下我們來看看.
啟示的十八章.
又提到那個大巴比倫.
第52章的第十一節.
你們離開吧.
離開吧.
從巴比倫出來.
不要尖不結靜的物.
要從其中出來.
你們剛才耶和華.
戲名的人啊.
勿要自結.
必不自急亡.
也不自奔逃.
因為耶和華必在你們前頭行.
以色列的神.
必作你們的後盾.
又在前又在後.
到底是前還是後.

$^{1281}$又是前又是後.
就是說完全的保護.
從巴比倫回來.
如果你知道在古代.
是要走差不多五個月.
剛才也說了.
是沿著河岸.
這樣走的.
不是橫行沙漠.
一有風沙就死了.
所以是很長時間.
既然那些地方.
是這麼好走.
很容易就會有土匪.
所以以前是很危險的.
所以有些王道.
是很寬闊.
可以讓馬車過.
某程度都是軍事的.
重要的國度.
所以.
很多時候都提到.
是屬靈的含義.
多一些.
不只是地理的巴比倫.
是巴比倫代表著的.
是離開神.
邪惡.
不潔.
要離開這些不潔.
回到上來.
所以這裡又出現過了.
剛才我們看的撒迦利亞書.
強調住在北方的.
巴比倫地方的錫安民.
強調.
你是住在那些.
外邦人當中.
不理你是什麼情況.
但是你在神的眼中.

$^{1321}$是錫安的子民.
這個很重要.
有了這個.
對自己身份的意識.
其實就是一個復興的開始.
我是基督徒.
為什麼我今天會坐在馬會.
和其他人一起跑馬.
當你自己都發現.
其實不需要這樣的.
其實應該坐一星期.
唱聖詩.
敬拜神.
聽道.
很開心地彼此分享.
這樣才是基督徒的真正的享受.
你有這個發現.
其實你就開始意識到.
你自己是巴比倫的錫安民.
所以這裡是.
你可以想像.
是一個靈異上的.
那個錫安.
和那個巴比倫.
也是剛才.
二創書52章.
所說的那個情況.
離開吧.
離開這個巴比倫.
不要沾染他們的不潔.
我就會收納你們.
你們離開.
必不辭急忙.
因為以色列的神.
會成為你們的前路.
成為你們的後盾.
這樣來安慰.
就是說一個有心.
要復興的.
那個心靈.

$^{1361}$就會有神的靈的感動.
和那個看守.
所以寶貴我們.
內心對神的單純.
寶貴我們內心一份.
初心的純潔.
神.
就會淋在你的內心.
來和你講說話.
給你有安慰.
所以來到這裡的時候.
我和大家一起揭開18章.
我們以往看過啟示錄.
再揭一揭回去.
給大家看看巴比倫.
這個神所不喜歡的大城.
(看著大城).
18章.
此後我看見另有一位.
大權柄的天使.
從天降下.
地就因為他的榮耀發光.
他大聲喊著說.
巴比倫大城傾倒了傾倒了.
成了鬼魔的住處.
和各樣污穢之靈的巢穴.
並各樣污穢可憎之雀鳥的巢穴.
因為列國都被他.
邪淫大怒的酒傾倒了.
地上的君王與他行淫.
地上的客商因他奢華太過.
就發了財.
我又聽見從天上有聲音說.
我的民啊.
你們要從那城出來.
免得與他一同有罪.
受他所受的災殃.
因他的罪惡滔天.
罪惡滔天.
即是天都聽到.

$^{1401}$他的不義.
神已經賞起來了.
這樣的描述.
不是神現在才突然間醒.
從來都見到.
不過時間到了.
神的容忍.
也要來施行審判了.
神已經賞起他的不義了.
他怎樣待人.
也要怎樣待他.
按他所行的家賠報應他.
用他調酒的杯.
家賠的調給他喝.
他怎樣榮耀自己.
怎樣奢華.
也當叫他照樣痛苦悲哀.
因他心裡說.
我坐了皇后的位.
並不是寡婦.
卻不至悲哀.
所以在一天之內.
他的災殃要一齊來到.
就是死亡.
悲哀.
饑荒.
他又要被火燒盡了.
因為審判他的主.
神大有能力.
你就可以想像到.
巴比倫城的那個傾倒.
將來到底這是一個什麼樣的體系呢.
是一個政治.
商業.
還是宗教的集成體.
還是怎樣呢.
我們很難對號入座.
但是.
也很肯定的就是.
那個敵黨神的一個地方.

$^{1441}$可能是一個集團.
一定是魔鬼操控的.
背叛神.
離開神.
是沒有神的.
最終是要被傾倒.
所以.
意識到這樣的情況.
趕快的投靠真神.
回到神的錫安.
當然是靈界.
不是要你去.
是屬靈的層面.
不是要你去錫安城.
去耶路撒冷.
每年都去朝聖.
就等於.
你就是錫安民.
現在是告訴我們.
我們每一個在主耶穌裡面的神.
都看我們是錫安的子民.
但是要離開.
離開巴比倫.
三次了.
以賽亞書的52章.
可能都是個被擄的時候.
撒迦利亞書第二章.
又說要在北方巴比倫的錫安民.
要歸回了.
不要和他們一起了.
去到將來的.
最終極的.
那個的集團式.
都是要被神來去傾覆了.
神想起他的不義.
所以我們求主讓我們能夠.
在有限的人生.
我們在這個.
蒙愛的群體裡面去成長.
所以一生都要回教會.

$^{1481}$要有教會成為你的Umbrella.
來保護你.
守護你.
還有你在那個地方.
你遇見神.
一起事奉神.
這樣的人生.
所以就算不同的教會.
總之他是講真理.
就有神的同在.
不要離開了教會.
而去到這個巴比倫.
所以不是物質地域地理的巴比倫.
是靈性的巴比倫.
是沒有了神的地方.
我們要記得這件事是很重要的.
一生的.
所以信主.
你小時候會背.
這時天賦世界.
但是你二十歲.
三十歲.
四十歲.
你的事業有成了.
你還回不回教會.
唱不唱這些兒歌.
你還是不是相信耶穌是你的救主.
很多人都覺得已經是老土和過時.
他自己都不知道多麼名成利就.
不需要信耶穌.
大有人在的.
我們這個社區都大有人在的.
有時見到面.
見到牧師都好似好尷尬.
但是要回到教會.
離開巴比倫.
要回到.
因為你是錫安的子民.
這個身份是神給的.
這個是很重要.

$^{1521}$所以我們更多要為那些不再回教會的人祈禱.
回教會的其實不用為他祈禱.
因為他已經很穩定.
應該為那些不回教會的祈禱.
不再出現在教會的.
有時真的要打開名冊.
逐個為名冊裡面的不再回教會的.
去流淚祈禱.
作為堂主任.
作為牧師.
真的不是開玩笑的.
要為這些弟兄姊妹祈禱.
有些可能他們在外國很好的聚會.
那沒問題.
有其他的牧者去牧羊.
但是離開羊群的.
已經在巴比倫的錫安民.
其實要歸回.
這個就是撒迦利亞書第二章的訊息.
我們再看多一點.
本來打算看第三章.
關於大祭司的約書亞.
可能未必夠時間看.
我們只是這一段還未解得完.
第七節.
與巴比倫同住.
身在巴比倫.
但實際你的生命是屬神.
很多香港的掛名基督徒.
甚至不再回教會的基督徒.
其實就是這句聖經涵蓋了.
應當逃脫離開巴比倫.
萬君子耶和華說.
在顯出榮耀之後.
我去懲罰那擄掠你們的列國.
摸你們的就是摸他眼中的同仁.
要顯出榮耀之後.
所以就不是現在重建這個階段.
往後神一定會保護他自己的選民.
一定會照顧.

$^{1561}$一定會顯出他的榮耀.
和公義的審判.
所以這是一個很大的安慰.
摸你們就等於摸神的.
眼中的同仁.
大家試試摸摸自己的眼皮.
你馬上縮了.
因為這是很自然的.
神給我們的保護機制.
因為我們的眼是非常複雜的.
所以我們自己都不可以摸我們自己.
如果我們要做眼的手術.
都要有其他的很先進的科技.
或者是怎樣的情況.
你可以想像得到摸以色列.
就是摸神眼中的同仁.
你可以想像得到是多麼的尊貴.
我們的身份是多麼的神的保護.
又很快突然間想起約翰福音十章.
十章的應該是十八節.
主耶穌說的.
我們是雙重確認.
不知道我們的信仰有沒有這個雙重確認.
雙重的保障.
耶穌是被那些法利賽人追問他.
十章的.
十八節沒錯了.
第十六節開始.
好牧人這個比喻.
第十六節開始一直到第十八節.
我另外有陽.
不是這圈裡.
是不是在說外邦人呢.
可能是也不一定的.
我必須領他們來.
他們也要聽我的聲音.
並且要合成一群.
歸一個牧人.
就是說我們是一主一信一洗一神.
第十七節.

$^{1601}$我父愛我.
因我將命捨去.
好在取回來.
沒有人奪我的命去.
是我自己捨的.
我有權並捨了.
也有權並取回來.
這是我父所受的命令.
我想說的不是這一段.
下面.
對不起.
應該是一直讀下去下面就對了.
猶太人就為這些說話起了爭論.
內宗有好些的人就說.
這個人是被鬼附的.
而且發瘋了.
為什麼這樣說呢.
有些人就說.
鬼附的人怎麼會這樣說話呢.
鬼怎麼能夠叫之前的核子眼睛開呢.
大家都很疑惑.
為什麼耶穌不斷稱呼.
這位神是他的父呢.
這樣的意思.
第廿二節.
這一段我想說的就是這一段.
很重要的一個雙重確認我們的信仰.
在耶路撒冷有修殿節.
就是那個重光節.
聖殿被毀了.
其實那個被毀是污穢了那個被毀.
是當時在主前的164 165年.
有一個敘利亞的王叫埃提亞古四世.
很憎恨猶太人.
打敗了.
去到聖殿.
他就在聖殿牽一隻豬.
在那裡屠污聖殿.
令到猶太人.
當時的一些.

$^{1641}$有一個家族叫瑪加比的家族.
其中一個兒子叫瑪他提亞.
就來奮勇地.
招聚自己的同胞.
趕走了埃提亞古四世.
重光聖殿.
就是潔淨聖殿.
在聖殿裡面重新點燈.
於是這個就叫修殿節.
現在耶穌就在修殿節.
在聖殿裡面行走.
在冬天的時候.
耶穌在殿裡所羅門的廊下行走.
當然不是耶穌入聖殿.
因為他是猶太人.
他不是祭司.
二十四節.
猶太人圍著他說.
你叫我們猶疑不定到什麼時候.
如果你真是基督.
你就明明告訴我們.
耶穌回答說.
我已經告訴你們.
你們不信.
我奉我父之名所行的事.
可以為我作見證.
只是你們不信.
因為你們不是我的羊.
這個好像是一個揀選的問題.
他們一直都疑惑.
不是他們不夠聰明.
不明白.
是他們從來都不是神的兒女.
他們是猶太人.
他們是范尼徐人.
是拉比.
就是說生命上.
他們就算有一套的知識.
那些知識沒有令到他們的生命轉化.
成為耶穌的門徒.

$^{1681}$就是被耶穌這個身份卡住了.
所以他們那些舊約的猶太教的知識.
是不會令到他們得救.
是必須認信耶穌.
才成為神的兒女.
或者就是耶穌這裡所說的.
就是他的羊.
接著下面.
二十七節.
這裡如果大家會背就好了.
十章二十八應該是.
我的羊聽我的聲音.
我也是認識他們.
他們也跟著我.
留意二十八節.
我又賜給他們永生.
他們永不滅亡.
誰也不能從我手裡把他們奪去.
這是第一個confirm.
就是賜永生.
永不滅亡.
沒有人可以在我的手裡.
耶穌你是誰.
你是木匠的兒子.
人人都可以在你手上搶.
彼得也被人搶去否認你.
所以這裡說的是.
耶穌得到榮耀之後.
就是說.
那個邏輯就是這樣.
如果你真的是耶穌的羊.
首先就是.
你是聽他的聲音.
你是願意跟隨他.
那就是意味著.
是用另一種說法.
你是有聖靈.
你是有決志.
還有你是有了神的生命.
你有聖靈在你裡面.

$^{1721}$你是已經有了一種生命的重生.
你是願意跟隨神的.
想這樣跟隨.
耶穌就說.
我就賜給他們永生.
是永不滅亡.
所以永生和永不滅亡.
是聽從耶穌的教訓.
跟隨耶穌.
所以是一次相信.
一生相信.
必定得救.
那個意思就是這樣.
還沒完.
二十九節.
我父把羊賜給我.
他比萬有都大.
誰也不能從我父手裡.
把他們奪去.
說了兩次不能奪去.
那就是說.
這是一個雙重確認的救恩.
我與父緣為一.
其實就是一位神.
這樣的強調又強調.
其實正正就是要駁斥猶太人.
他們是被拒絕.
在救恩的門外.
不是他們不可以信.
而是他們自己.
將耶穌這個榮耀的彌賽亞身份.
來去否認.
某程度上就是等於褻瀆聖靈.
所以他們是被注定要滅亡.
但是每一個信主的人.
我們要有這份永恆的確信.
永恆的保障.
所以這些經文大家要背熟它.
耶穌基督這樣說.
再不夠.

$^{1761}$再加上我父比萬有都大.
是父神將羊賜給我.
現在由他去照管.
沒有人可以在他的手上.
把羊奪去.
那就可以想像得到.
其實就是說.
這個救恩是可以肯定的.
是因信而稱義.
而得到這個永恆的天家.
只需要我們聽耶穌的聲音.
跟隨耶穌.
那就是他的羊.
這個是有講究的.
所以有些人就是.
信了主耶穌.
洗完禮.
他覺得可以不需要再回教會了.
他覺得回不回那是無所謂的.
其實根本符合不到這裡.
整個上下文是說.
我的羊認識我.
我又認識他們.
他們也跟著我.
他們又認得我的聲音.
我又認識他們.
這個關係的建立.
是生命的關係.
所以我們不是來聽道理.
我們是神的兒女.
我們在神的聖殿裡面.
一起敬拜真神.
所以這個是.
真的很重要的一個雙重確認.
正正就是撒迦利亞所講的.
神說將來如果.
來攻擊你的.
就等於摸我眼中的瞳人.
其實動物主哪裡有形象.
動物主哪裡會有一雙眼.

$^{1801}$連基督教都說.
上下片滿眼睛.
但是神這樣看他的選民.
神這樣來愛他的兒女.
所以大家要背下.
這些很重要的經文.
第二章的第八節.
摸你們的就是摸神眼中的瞳人.
所以如果我們在神的愛裡面.
假如將來我們要受苦.
受逼迫.
我們不需要擔心.
如果神不容許.
我們是不會受額外過多的苦難.
如果神容許.
即是說一定是有一個鍛鍊.
是有一個旨意.
在我們的生命裡面.
所以原來舊約已經有.
這些這麼寶貴的應許.
所以神的選民.
真的可以來到他有一個安穩.
所以有一首聖詩就是.
安穩在耶穌手中.
不需要再懼怕.
或者說我心靈得安寧.
那首聖詩.
就是那個作者的爸爸.
在大西洋.
他全家人三個女兒.
被淹死了.
沉船.
太太就寫了一個電報.
被救回上岸之後.
就寫了一個電報給丈夫.
就說只有他一個得救.
Saved Alone.
接著那個作者.
當然當時的一些牧者.
坐船去到大西洋.

$^{1841}$船長就告訴他們.
大概這個位置.
就是你三個女兒淹死的地方.
他就衝回去船艙.
就作了這一首聖詩.
我心靈得安寧.
那首聖詩是很出名的.
很好聽的.
很多的東西他不能夠解釋.
但是神有他的美意等等.
就是幾個女兒全死了.
我忘記了那個作者的名字.
我今晚在群組告訴大家.
那首聖詩突然想起.
就是第八節這裡.
摸你們的就是摸我眼中的銅人.
所以我自己.
剛剛星期六晚在講福堂講道.
一直講.
講少少見證.
我信耶穌.
1987年的12月24日平安夜.
平生第一次去聽報道會.
第一次聽了十分鐘就舉手決志.
都不知道是什麼來的.
就開展了我36年到今天.
那個信仰的人生.
是有很多的回望.
很多恩典.
其實是很多十字路口.
但是回望那些十字路口.
你不知道為什麼會選擇這樣走.
你現在回望.
將它們一個一個連接起來.
你就會發現.
不是喔 是可以連接的.
是有一個故事在那裡.
那就很奇妙.
很恩典.
每一次的回望都是.

$^{1881}$哎呀 冤孽.
慘了.
糟了.
不是喔 是恩典.
是真的能夠解釋到.
這樣是最好的選擇.
當下你完全不知道.
如果能夠這樣回望你的信仰.
其實正正就是摸你們的.
就是摸我眼中的同人.
神體我們是比同人更加珍貴.
這個信仰你想想多麼寶貴.
值得我們去跟隨.
為什麼我會和大家這樣聯想.
因為我講的經文就是提摩塔前書.
保羅說.
因他而我又忠心.
派我去服侍他.
其實我從前是褻瀆神的.
我從前是逼迫人的.
我從前是舞脈人的.
但是呢.
我是不信.
不明白的時候而作的.
但是我主的因是格外的豐盛.
恩典去改造保羅.
使我能夠在基督耶穌裡面.
就是這裡.
有信心.
有愛心.
所以保羅就是經歷到基督耶穌的恩典.
使他能夠有信心.
也就是對自己的成長和服侍.
他又開始有自信.
和有愛心去幫助其他人.
所以他能夠說他自己.
基督耶穌降世.
就是要拯救罪人.
他在罪人當中他自己承認.
他是一個罪魁.

$^{1921}$因為不是那個量.
是那個質.
他居然逼迫耶穌.
來侮辱,嘲笑和褻瀆這位耶穌.
說他不是彌賽亞.
本來是永恆的死罪.
很多猶太的拉比.
剛才我們看的十章.
《若行福音》就是這樣.
真的注定不能夠救.
但是他居然蒙神的憐憫.
保羅還要大大地用他.
要顯明在罪魁身上一切的忍耐.
給後來信耶穌的人做榜樣.
所以保羅是一個人版.
就是神的忍耐.
來去孕育他的生命.
成為別人的祝福.
所以這裡很重要.
這個同仁.
所以大家有沒有默想.
你自己是神的同仁呢?.
是神眼中的同仁呢?.
比神眼中的同仁還要來得珍貴呢?.
我們的滿足感,榮耀感是從哪裡來呢?.
每個人都會的.
大家看到這幾天的新聞就知道.
那些元父的KOL.
說我又出糧了,七十多萬了.
我們是從哪裡而來的滿足和榮耀呢?.
我們信主的應該是在神的裡面.
有這一份的滿足.
因為祂說我在祂的手.
沒有人可以在祂的手將我們奪去.
說的是我們的生命.
說的是我們的信仰.
所以求主讓我們能夠在我們的生命裡面.
都有這樣的堅持,這樣的相信.
今晚我們就說到這裡了.
我們一起禱告.

$^{1961}$多謝主讓我們稍微看了第二章八節的經文.
很盼望我們能夠繼續深化我們的信仰.
原來舊約都有這麼多的安慰.
這麼寶貴的應許.
我們的主你是四圍的火城.
更加是百姓的榮耀.
更加是我們每一個人的保護和安慰.
我們比同人Apple更加的寶貴.
願主繼續使我們的信心堅固.
能夠在主的愛裡面成長.
多謝您聽我們禱告.
奉耶穌基督的名.
阿們.
\newpage



\section{使徒行傳 21:17-26}
\label{sec:B_Q_tzD4YvY}
\textbf{2023-09-24 加恩樓《入會四則》使徒行傳 21\_17-26 講員:陳麗蟬牧師}
\newline
\newline
連結: \href{https://youtube.com/watch?v=B-Q-tzD4YvY}{\texttt{https://youtube.com/watch?v=B-Q-tzD4YvY}} ~~~~ 語音日期: 2023-09-26
\newline
\newline
\hyperref[sec:_MeYYEhxYZE]{\small{< < < PREV SERMON < < <}}
~
\hyperref[sec:index_chronic]{\small{[返順時目]}}
~
\hyperref[sec:index_scriptual]{\small{[返順卷目]}}
~
\hyperref[sec:CTfZJ96vxzI]{\small{> > > NEXT SERMON > > >}}
\newline
\newline
使徒行傳 21:17-26
\newline
\begin{longtable}{cl}
\hline
\hline
章節 & 經文 (和合本修訂版)\\
\hline
21:17 & \begin{tabularx}{0.7\textwidth}{X} 我們到了耶路撒冷，弟兄們歡歡喜喜地接待我們。 \end{tabularx} \\ \\ \relax
21:18 & \begin{tabularx}{0.7\textwidth}{X} 第二天，保羅同我們去見雅各；所有的長老也都在場。 \end{tabularx} \\ \\ \relax
21:19 & \begin{tabularx}{0.7\textwidth}{X} 保羅向他們問安，然後將神用他在外邦人中所做的事奉，一一述說了。 \end{tabularx} \\ \\ \relax
21:20 & \begin{tabularx}{0.7\textwidth}{X} 他們聽見了，就歸榮耀給神，對保羅說：「弟兄，你看猶太人中有數以萬計的信徒，而他們都是熱心於律法的人。 \end{tabularx} \\ \\ \relax
21:21 & \begin{tabularx}{0.7\textwidth}{X} 他們曾聽見人說，你教導所有在外邦的猶太人離棄摩西，對他們說，不要給孩子行割禮，也不要遵守規矩。 \end{tabularx} \\ \\ \relax
21:22 & \begin{tabularx}{0.7\textwidth}{X} 眾人必聽見你來了，這可怎麼辦呢？ \end{tabularx} \\ \\ \relax
21:23 & \begin{tabularx}{0.7\textwidth}{X} 你就照著我們的話做吧！我們這裡有四個人，都有願在身。 \end{tabularx} \\ \\ \relax
21:24 & \begin{tabularx}{0.7\textwidth}{X} 你帶他們去，與他們一同行潔淨的禮，替他們繳納規費，讓他們得以剃頭。這樣，眾人就會知道，先前所聽見關於你的事都是假的；而且也知道，你自己為人循規蹈矩，遵行律法。 \end{tabularx} \\ \\ \relax
21:25 & \begin{tabularx}{0.7\textwidth}{X} 至於信主的外邦人，我們已經根據我們的決議寫信，叫他們要禁戒偶像所玷污的東西、血和勒死的牲畜，禁戒淫亂。」 \end{tabularx} \\ \\ \relax
21:26 & \begin{tabularx}{0.7\textwidth}{X} 於是保羅帶著那四個人，第二天與他們一同行了潔淨禮，進了聖殿，報告潔淨期滿的日子，等候祭司為他們各人獻上祭物。 \end{tabularx} \\ \\
[1ex]
\hline
\hline
\end{longtable}
$^{1}$(英語)姐妹們,早安.
(英語)1863年,歐洲很多地方發生戰事.
(英語)當中意大利和奧地利正在打仗.
(英語)有一天,他們停止開戰.
(英語)各國將自己的傷兵帶回自己的地方.
(英語)來醫治他們.
(英語)在意大利的救傷隊裡面.
(英語)將戰場上的傷兵帶回他們.
(英語)將他們帶到教堂裡面.
(英語)來進行醫治.
(英語)但在醫治當中,突然發覺到.
(英語)不知道是否心急還是混亂.
(英語)其中竟然有四個敵軍的傷兵.
(英語)所以大家就吵起來.
(英語)這是奧地利的,不是我們意大利的.
(英語)於是有些人就很激動.
(英語)快點將他們拉到街上.
(英語)讓他們自生自滅,不要醫治他們.
(英語)當時有一個人,他穿著一件白袍.
(英語)他就走出來了.
(英語)不可以這樣.
(英語)不可以將他們拉到街上,讓他們自生自滅.
(英語)不可以的.
(英語)他是因為我們都是弟兄.
(英語)我們都是弟兄.
(英語)當他講出這句話的時候.
(英語)這群原本氣沖沖的人.
(英語)全部安靜下來.
(英語)然後又將這四個奧地利的傷兵.
(英語)帶回教會裡面,來醫治他們.
(英語)其實在那天,這個穿白袍的人不是醫生.
(英語)他只是一個商人.
(英語)不過他一直都在推行人道主義.
(英語)一直都在推行人道主義.
(英語)所以他覺得.
(英語)就算你抓了這四個奧地利的傷兵回來.
(英語)你都不可以任由他們自生自滅.
(英語)都要醫治.
(英語)就在那天.
(英語)就有一個德國的醫生.

$^{41}$(英語)有三個意大利的學生.
(英語)有一個法國的新聞的採訪員.
(英語)有兩個英國的青年.
(英語)他們就在禮拜堂裡面.
(英語)從早到晚.
(英語)一直將所有的傷兵.
(英語)不管他們是什麼國籍.
(英語)他們都來醫治他們.
(英語)就因著這個傷人.
(英語)他說我們都是弟兄.
(英語)這個穿白袍的人.
(英語)他就是紅十字會的創辦人.
(英語)Henry Donald.
(英語)他就在那天裡面.
(英語)就因為他一句話.
(英語)就開始在這個世界上.
(英語)就開始了這個跨國的救治.
(英語)醫治.
(英語)在戰爭的時候.
(英語)只要紅十字會一出的時候.
(英語)就不分國籍都要來去醫.
(英語)我想在我們處理了這些戰事戰爭.
(英語)在我們成長的裡面.
(英語)在我們事奉的裡面.
(英語)我相信我們每一個人.
(英語)都會有一些心靈受傷的經歷.
(英語)我們都是受傷的人.
(英語)特別在這個教會裡面.
(英語)一起成長.
(英語)一個本有一個充滿了愛的家的裡面.
(英語)但我們又發覺到.
(英語)在這個團契.
(英語)在小組的裡面.
(英語)沒錯,我們都是要來去.
(英語)不要分背景.
(英語)不要分喜好.
(英語)不要分階級.
(英語)我們都要來去有主耶穌基督的愛.
(英語)在當中成為我們的榜樣.
(英語)一起來去生活和成長.

$^{81}$(英語)但往往就是因著我們很多時候.
(英語)一些不同.
(英語)因著一些不同.
(英語)也都會同樣出現一些傷痕.
(英語)出現一些令人難過的地方.
(英語)今天我們就用這個《使徒行傳》第21章.
(英語)在裡面其實就是講到保羅.
(英語)和耶路撒冷教會裡面.
(英語)一班的舊派人.
(英語)我當保羅是新派人.
(英語)他們彼此之間.
(英語)其實有很大的衝突.
(英語)但你怎樣能夠彼此化解這些衝突.
(英語)彼此的勉勵.
(英語)以致到免去了歷史上.
(英語)基督教的新派和舊派的衝突.
(英語)這個就是保羅.
(英語)保羅在聖經裡面我們看到.
(英語)他有三次的宣教旅程.
(英語)有些說三次半,有些說四次.
(英語)那怎樣看我們第三次的旅程呢.
(英語)第一次的旅程.
(英語)保羅就在這個安提柯教會開始.
(英語)就去到不同的亞細亞的地方.
(英語)帶了不少的外邦人信主.
(英語)這個以前猶太人的觀念就是.
(英語)這個救恩只是臨到猶太人的裡面.
(英語)但現在這個保羅.
(英語)竟然去到不同的外邦的地方.
(英語)帶了很多的外邦人信主.
(英語)於是在耶路撒冷的會議裡面.
(英語)他們就開了會.
(英語)也因為福音要傳到外邦地.
(英語)上帝也很好的.
(英語)他的頭腦也很精密.
(英語)他就先給這個救派人彼得.
(英語)就讓他看到意象.
(英語)這個福音要傳到外邦地.
(英語)也讓他看到意象.
(英語)他要先行一步.

$^{121}$(英語)去到這個外邦人.
(英語)哥尼留那裡去講這個福音.
(英語)所以在第一次宣教旅程那裡.
(英語)就定了這個外邦人.
(英語)不需要像猶太人那樣受割禮.
(英語)也要注意外邦人.
(英語)不可以去拜偶像.
(英語)不可以去吃勒死的生畜和血.
(英語)為什麼要特別提出這兩件事呢.
(英語)就是因為當時這些外邦拜祭偶像的神廟.
(英語)他們在拜祭裡面.
(英語)我們也很清楚知道.
(英語)常常在拜祭完之後.
(英語)就有很多淫亂的事情發生.
(英語)覺得這樣是偶像歡喜的.
(英語)他們吃的食物.
(英語)最喜歡吃一些生畜的血.
(英語)他們殺死生畜的時候.
(英語)就拿繩子勒死他們.
(英語)所以為什麼耶路撒冷會議.
(英語)要這樣去規定他們.
(英語)不單止因為殘忍.
(英語)也是因為當他們這樣做的時候.
(英語)就會想起來.
(英語)很容易想起他們在神廟拜祭的時候.
(英語)這些的情況.
(英語)所以就叫他們不要這樣做了.
(英語)到第二次的宣教旅程.
(英語)上帝給保羅他們有很多的恩典.
(英語)建立了不同的教會.
(英語)到第三次旅程了.
(英語)他們在公元59年的時候.
(英語)去到尾程的時候.
(英語)就回到耶路撒冷教會.
(英語)就是我們剛才所讀的經文.
(英語)剛才所說的經文就是.
(英語)保羅和路加醫生和一些人.
(英語)去到耶路撒冷那裡.
(英語)我們就說他們很開心.
(英語)一群長老和當時的領袖雅各.

$^{161}$(英語)很歡喜地接待保羅,路加醫生和其他的人.
(英語)當時耶路撒冷教會的領袖就是雅各.
(英語)雅各就是耶穌基督的弟弟.
(英語)當時的教會就把行政的事務.
(英語)交給了雅各來負責.
(賴麗櫻)大家一般來寒暄問候一下.
(賴麗櫻)保羅就分享他在這個外邦的地方.
(賴麗櫻)在這個亞細亞的地方.
(賴麗櫻)他做了的事情.
(賴麗櫻)去過的地方.
(賴麗櫻)原來在過往第一次宣教旅程.
(賴麗櫻)一直到第三次.
(賴麗櫻)成立了很多的教會.
(賴麗櫻)有加拉提教會,菲律賓教會,哥林多教會.
(賴麗櫻)帖撒羅尼迦教會,以弗所教會,哥羅西教會.
(賴麗櫻)我們讀的書信裡面都有一間教會.
(賴麗櫻)在過往保羅一直來設立的.
(賴麗櫻)他說了很多的教會.
(賴麗櫻)很多的上帝成就的事情.
(賴麗櫻)這班耶路撒冷.
(賴麗櫻)就當他們是舊派的領袖.
(賴麗櫻)聽到保羅說了這麼多福音的果子.
(賴麗櫻)接著他們說到他們很歡喜.
(賴麗櫻)很欣賞將榮耀歸給上帝.
(賴麗櫻)將榮耀歸給上帝.
(賴麗櫻)有時候我們跟其他人一起合作.
(賴麗櫻)無論耶路撒冷教會或保羅.
(賴麗櫻)去到不同的亞細亞地方.
(賴麗櫻)一起合作在教會或公司都會.
(賴麗櫻)我們有時候會看到別人的成就.
(賴麗櫻)多了這麼多這麼好.
(賴麗櫻)有時候我們會心酸.
(賴麗櫻)所以當聽到別人的成就.
(賴麗櫻)我們最多就是眼瞄一下.
(賴麗櫻)不會說任何的東西.
(賴麗櫻)但其實我們看到這裡.
(賴麗櫻)看到耶路撒冷教會的領袖.
(賴麗櫻)雅各和長老.
(賴麗櫻)他們是歡喜的,欣賞的.
(賴麗櫻)並且將榮耀歸給上帝.

$^{201}$(賴麗櫻)在教會整個群體裡.
(賴麗櫻)第一樣我們要學習的.
(賴麗櫻)就是互相的欣賞.
(賴麗櫻)在一起合作,事奉的裡面.
(賴麗櫻)多說欣賞,鼓勵,建立人的說話.
(賴麗櫻)這樣會合作得愉快.
(賴麗櫻)合作得順暢一點.
(賴麗櫻)另外,雅各又繼續說.
(賴麗櫻)他說我們在耶路撒冷教會.
(賴麗櫻)我們聽到一些關於你的說話.
(賴麗櫻)他說你在外邦的地方.
(賴麗櫻)你叫那些外邦地方的猶太人.
(賴麗櫻)叫他們不用守摩西的律法.
(賴麗櫻)不需要幫你的小孩行割禮.
(賴麗櫻)不需要守猶太人的規矩.
(賴麗櫻)我們在耶路撒冷教會聽到這樣的說話.
(賴麗櫻)你覺得怎樣辦好呢?.
(賴麗櫻)雅各和那些長老.
(賴麗櫻)真的聽到這些說話.
(賴麗櫻)說到保羅在亞細亞的地方.
(賴麗櫻)一方面做外邦人工作.
(賴麗櫻)另一方面又教那些住在外邦的猶太人.
(賴麗櫻)不再需要守摩西的律法.
(賴麗櫻)不需要守規矩.
(賴麗櫻)因為摩西的律法對猶太人來說.
(賴麗櫻)是很重要的.
(賴麗櫻)很尊重摩西的.
(賴麗櫻)你覺得教他們不需要守規矩.
(賴麗櫻)很多猶太人都會很憤怒.
(賴麗櫻)但我們看到在第22節.
(賴麗櫻)雅各和那些長老的語氣是怎樣呢?.
(賴麗櫻)那些人聽到你們來了.
(賴麗櫻)聽到你們說教人不需要守規矩.
(賴麗櫻)怎樣呢?.
(賴麗櫻)和他們的語氣和用詞.
(賴麗櫻)其實這些領袖都不相信這些謠言.
(賴麗櫻)他們很想和保羅一起想辦法.
(賴麗櫻)我想大家一起合作的最重要的.
(賴麗櫻)最重要的就是互相信任.
(賴麗櫻)在這個信任之中.

$^{241}$(賴麗櫻)大家才可以一起繼續共事.
(賴麗櫻)否則我們一起工作的時候.
(賴麗櫻)我們經常懷疑對方有沒有偷懶.
(賴麗櫻)有沒有不忠心.
(賴麗櫻)又覺得對方可能做事做得不夠自己厲害.
(賴麗櫻)於是就每次都管著.
(賴麗櫻)是最慘的.
(賴麗櫻)我認識一位姐妹.
(賴麗櫻)她在琴行教琴.
(賴麗櫻)她說她的老闆.
(賴麗櫻)不懂彈琴和音樂.
(賴麗櫻)但總覺得她這樣教不好.
(賴麗櫻)於是就教她怎樣教琴.
(賴麗櫻)有時來不及.
(賴麗櫻)都會聽到一些弟兄姊妹在工作中的分享.
(賴麗櫻)她說行外的人.
(賴麗櫻)管行內的人.
(賴麗櫻)很辛苦.
(賴麗櫻)她甚麼都不懂.
(賴麗櫻)但教你要怎樣做.
(賴麗櫻)這些工作氣氛.
(賴麗櫻)一定是僵硬的.
(賴麗櫻)一定是不愉快的.
(賴麗櫻)所以在這個群體.
(賴麗櫻)或是教會裡面的.
(賴麗櫻)除了互相欣賞之外.
(賴麗櫻)第二方面就是要互相信任.
(賴麗櫻)互相信任之後.
(賴麗櫻)但真的有聽到這些謠言.
(賴麗櫻)說你們不規矩.
(賴麗櫻)她就說.
(賴麗櫻)她說弟兄.
(賴麗櫻)我們在耶路撒冷.
(賴麗櫻)這個基督教.
(賴麗櫻)和猶太教的發源地.
(賴麗櫻)我們裡面有數以萬計的猶太信徒.
(賴麗櫻)當然有些是信了耶穌的.
(賴麗櫻)就算這些信了耶穌的人.
(賴麗櫻)他們仍然熱心於摩西律法.
(賴麗櫻)那怎麼辦呢?.

$^{281}$(賴麗櫻)還有當中有不少是激進派的分子.
(賴麗櫻)雅各就擔心.
(賴麗櫻)擔心保羅他們可能會遇到危機.
(賴麗櫻)可能會抓他們.
(賴麗櫻)其實在21章以後.
(賴麗櫻)保羅就遇到很大的危險.
(賴麗櫻)所以那怎麼辦呢?.
(賴麗櫻)雅各就給予一些提議.
(賴麗櫻)不是說怎麼辦.
(賴麗櫻)怎麼辦就停留在那裡.
(賴麗櫻)他說既然那些人說你們.
(賴麗櫻)在外面的地方不守猶太的律法.
(賴麗櫻)不守摩西的律法.
(賴麗櫻)不如這樣吧.
(賴麗櫻)現在我們有四個.
(賴麗櫻)剛剛完成了那塞爾人的許願的人.
(賴麗櫻)那這個那塞爾人的許願.
(賴麗櫻)就是他們有些人.
(賴麗櫻)他們想在某一段時間.
(賴麗櫻)來更加的親近神.
(賴麗櫻)領受上帝更多.
(賴麗櫻)就好像我們現在說的退休會.
(賴麗櫻)或者現在嘉欣很好.
(賴麗櫻)現在逢星期一又有些退休日.
(賴麗櫻)很想在某一段時間裡.
(賴麗櫻)親近上帝更多.
(賴麗櫻)那時候他們就會有很長的.
(賴麗櫻)不是一天的.
(賴麗櫻)七天至到有三十天.
(賴麗櫻)在那段時間他們要剃光頭髮.
(賴麗櫻)不准喝酒,不准吃肉.
(賴麗櫻)他們自己許願要跟上帝有親密交往的日子.
(賴麗櫻)他們要在那天完結後.
(賴麗櫻)他們就會有一個叫做潔淨禮.
(賴麗櫻)又要出回外面的世界.
(賴麗櫻)來工作去事奉.
(賴麗櫻)阿葛就提議.
(賴麗櫻)他說剛剛有四個許願者.
(賴麗櫻)他們已經完成了.
(賴麗櫻)不如這樣吧.

$^{321}$(賴麗櫻)你跟他們四個人一起舉行潔淨禮.
(賴麗櫻)還有.
(賴麗櫻)你幫他們付潔淨禮的費用.
(賴麗櫻)給聖殿.
(賴麗櫻)然後給祭司看看.
(賴麗櫻)你們完成了日期.
(賴麗櫻)還有獻祭.
(賴麗櫻)證實給他們看.
(賴麗櫻)你不是不守規矩的人.
(賴麗櫻)你是很尊重猶太的律法.
(賴麗櫻)並且你跟人一起去守這些律法.
(賴麗櫻)阿葛和長老有這樣的建議.
(賴麗櫻)如果我是保羅我會怎樣做呢.
(賴麗櫻)可能我就會說.
(賴麗櫻)有沒有搞錯.
(賴麗櫻)第一我沒有不守規矩.
(賴麗櫻)我沒有叫外邦猶太人不需要受國禮.
(賴麗櫻)我只是按照耶路撒冷的規定.
(賴麗櫻)外邦人信了耶穌.
(賴麗櫻)他們都不需要受國禮.
(賴麗櫻)我沒有叫住在外邦的小朋友.
(賴麗櫻)不需要跟隨摩西的律法.
(賴麗櫻)我沒有教過.
(賴麗櫻)你們在誣陷我.
(賴麗櫻)可能保羅心裡是這樣想的.
(賴麗櫻)還要我跟那四個那西議員的人.
(賴麗櫻)行結正禮.
(賴麗櫻)還要我幫他付錢.
(賴麗櫻)我沒有錯.
(賴麗櫻)為什麼我要這樣做.
(賴麗櫻)但保羅為了耶路撒冷教會.
(賴麗櫻)和他自己起碼和諧和平的時候.
(賴麗櫻)他願意照著做.
(賴麗櫻)有時我們一起去合作.
(賴麗櫻)一起的事奉就是這樣.
(賴麗櫻)保羅和耶路撒冷教會的時候.
(賴麗櫻)最初是客客氣氣.
(賴麗櫻)彼此的歡迎.
(賴麗櫻)彼此的欣賞.
(賴麗櫻)跟著第二步深入一點.

$^{361}$(賴麗櫻)就是互相提醒.
(賴麗櫻)無論友誼或是合作.
(賴麗櫻)要一直盡心.
(賴麗櫻)唯有一直有友誼.
(賴麗櫻)或者盡心的時候.
(賴麗櫻)才能脫離表面的關係.
(賴麗櫻)表面的關係是沒有關係的.
(賴麗櫻)就說你好嗎.
(賴麗櫻)但在盡心的時候.
(賴麗櫻)大家如何在愛的裡面.
(賴麗櫻)用愛心來說誠實話.
(賴麗櫻)雅各和長老.
(賴麗櫻)就說有人這樣.
(賴麗櫻)不如你這樣做吧.
(賴麗櫻)保羅就在同樣的和諧.
(賴麗櫻)和愛心的大前提下.
(賴麗櫻)他就吞了.
(賴麗櫻)行結證禮.
(賴麗櫻)他OK.
(賴麗櫻)行結證禮.
(賴麗櫻)他OK.
(賴麗櫻)為什麼保羅願意這樣做.
(賴麗櫻)為什麼他能夠互相遷就呢.
(賴麗櫻)我們來看看.
(賴麗櫻)為什麼保羅要行結證禮.
(賴麗櫻)他要獻一隻公羊做贖獻祭.
(賴麗櫻)又要獻一隻公綿羊做平安祭.
(賴麗櫻)又要獻素祭就是一些麵條.
(賴麗櫻)斷祭就是一些酒.
(賴麗櫻)所以費用是不少的.
(賴麗櫻)但是亞國建議保羅幫這些人付了費用.
(賴麗櫻)當時其實也是這樣.
(賴麗櫻)在他們的一些建議裡面.
(賴麗櫻)就說有些比較窮的人.
(賴麗櫻)可能經過了納西爾縣的期限.
(賴麗櫻)他要行這些祭.
(賴麗櫻)然後結證完之後.
(賴麗櫻)被祭司簽了名.
(賴麗櫻)OK可以出到世界了.
(賴麗櫻)可能他不夠錢.

$^{401}$(賴麗櫻)有些富有的人可以幫他們付.
(賴麗櫻)他們這樣做的時候就表明.
(賴麗櫻)我都支持尊重猶太人過往的規定.
(賴麗櫻)所以保羅也跟著去做了這件事.
(賴麗櫻)保羅其實是抱著一個心態.
(賴麗櫻)他在真理以外的事情一定不守.
(賴麗櫻)但在真理裡面無雙大愛的時候.
(賴麗櫻)他就是竭力保守聖靈所賜.
(賴麗櫻)合而為一的心.
(賴麗櫻)這個妥協不是等於軟弱.
(賴麗櫻)反而在真理範圍以內.
(賴麗櫻)這個妥協就是一個力量.
(賴麗櫻)大家有沒有試過.
(賴麗櫻)有時候覺得我不喜歡這樣做.
(賴麗櫻)但為著福音的緣故.
(賴麗櫻)為著和諧不分裂的緣故.
(賴麗櫻)我要去做的時候.
(賴麗櫻)心裡的力量更加要多.
(賴麗櫻)我們來看看.
(賴麗櫻)在《使徒行傳》裡面.
(賴麗櫻)很多時候都出同心合意的字.
(賴麗櫻)從《使徒行傳》第一章開始.
(賴麗櫻)就提到一群婦女和門徒.
(賴麗櫻)他們在這個樓房裡.
(賴麗櫻)同心合意地去祈禱.
(賴麗櫻)接著在第三章一直下去.
(賴麗櫻)就提到一群使徒.
(賴麗櫻)一起同心合意.
(賴麗櫻)來敬拜守主的餐盒.
(賴麗櫻)這個字在整個新約聖經裡面.
(賴麗櫻)就出現了12次.
(賴麗櫻)在《使徒行傳》裡面.
(賴麗櫻)又佔了10次.
(賴麗櫻)看到路加芬非常重視.
(賴麗櫻)在神的家裡.
(賴麗櫻)我們要如何同心合意.
(賴麗櫻)我們再一次看回.
(賴麗櫻)整個《使徒行傳》.
(賴麗櫻)在這個時間裡.
(賴麗櫻)我說他們正在面對一個很大的危機.

$^{441}$(賴麗櫻)因為在過往.
(賴麗櫻)在《使徒行傳》第15章之前.
(賴麗櫻)整個宗教的中心.
(賴麗櫻)都在耶路撒冷.
(賴麗櫻)但是在第15章之後.
(賴麗櫻)就開始轉了.
(賴麗櫻)轉了去安提奧.
(賴麗櫻)然後就一直發展在.
(賴麗櫻)剛才我說過的教會.
(賴麗櫻)保羅建立的那些.
(賴麗櫻)加拿大,菲律賓,伊夫索.
(賴麗櫻)帖撒羅尼迦,哥林多.
(賴麗櫻)轉移到那裡.
(賴麗櫻)所以其實.
(賴麗櫻)耶路撒冷教會就心裡有點害怕.
(賴麗櫻)還有在人物上高.
(賴麗櫻)在這個十三章.
(賴麗櫻)這個《使徒行傳》.
(賴麗櫻)十三章之前.
(賴麗櫻)那個中心人物都是彼得.
(賴麗櫻)然後就開始轉到保羅.
(賴麗櫻)所以可能彼得都很害怕.
(賴麗櫻)過往自己是宗教界裡的領袖.
(賴麗櫻)但現在竟然去了保羅的裡面.
(賴麗櫻)所以在初期教會裡面.
(賴麗櫻)我想大家在面對.
(賴麗櫻)除了在外面面對逼迫的壓力之外.
(賴麗櫻)更加要面對在內心裡.
(賴麗櫻)很多張力,很多比較.
(賴麗櫻)但我們看到耶路撒冷教會.
(賴麗櫻)阿哥和一群長老.
(賴麗櫻)他們看到保羅的成就.
(賴麗櫻)他們將榮耀歸給神.
(賴麗櫻)我們也看到在路加芬裡面.
(賴麗櫻)不斷出現同心合意這個字.
(賴麗櫻)同心合意這個字其實就是團契.
(賴麗櫻)就是團契.
(賴麗櫻)在商業裡面.
(賴麗櫻)大家彼此訂立一些協議.
(賴麗櫻)訂立合同.

$^{481}$(賴麗櫻)大家訂立了之後.
(賴麗櫻)大家建立了一份關係.
(賴麗櫻)溫慧瑤博士說.
(賴麗櫻)好像買股票一樣.
(賴麗櫻)你買了股票就進了那間公司.
(賴麗櫻)成為了他的股東.
(賴麗櫻)如果大家有買過股票.
(賴麗櫻)你就知道了.
(賴麗櫻)你就會經常留意電視機.
(賴麗櫻)那裡在說股市行情的時候.
(賴麗櫻)我爸爸以前在生的時候.
(賴麗櫻)他不懂得.
(賴麗櫻)不過一看見綠色箭咀.
(賴麗櫻)他就很高興.
(賴麗櫻)哇升了升了升了.
(賴麗櫻)看到很多紅色箭咀的時候.
(賴麗櫻)他說死了跌了跌了.
(賴麗櫻)於是就打電話給我大哥.
(賴麗櫻)告訴他死了.
(賴麗櫻)今天股市大跌了.
(賴麗櫻)關係很密切.
(賴麗櫻)其實原來這個同心合意.
(賴麗櫻)這個字的時候.
(賴麗櫻)都可以用到在教會.
(賴麗櫻)盧家醫生就用了.
(賴麗櫻)在當時的教會.
(賴麗櫻)不單是在初初.
(賴麗櫻)以耶路撒冷女為中心的教會.
(賴麗櫻)也在之後.
(賴麗櫻)安提柯和其他不同的外邦教會裡面.
(賴麗櫻)那一份大家怎樣.
(賴麗櫻)同心合意.
(賴麗櫻)慶旺復任.
(賴麗櫻)就是我們.
(賴麗櫻)大家是一個家.
(賴麗櫻)所以剛才那首詩歌很好.
(賴麗櫻)家.
(賴麗櫻)這個不同的風浪.
(賴麗櫻)在這個家裡面.
(賴麗櫻)有主耶穌基督的愛.

$^{521}$(賴麗櫻)就在裡面.
(賴麗櫻)充滿了愛.
(賴麗櫻)充滿了那份的連結.
(賴麗櫻)充滿了那份的合作.
(賴麗櫻)在這個家裡面.
(賴麗櫻)每一個信徒.
(賴麗櫻)不再是自己.
(賴麗櫻)不再是單單的我自己.
(賴麗櫻)所以你回來.
(賴麗櫻)你個個都不認識.
(賴麗櫻)不肯去跟其他信徒.
(賴麗櫻)去聯繫在一起的時候.
(賴麗櫻)這個不是真正進入了.
(賴麗櫻)這個屬靈的家.
(賴麗櫻)在這個家裡面.
(賴麗櫻)無論是佳音樓.
(賴麗櫻)無論是錦繡堂.
(賴麗櫻)我們凡事都以天國.
(賴麗櫻)上帝的國度的事.
(賴麗櫻)為中心.
(賴麗櫻)所以這個家.
(賴麗櫻)在這個天國.
(賴麗櫻)這個真理的大前提下.
(賴麗櫻)縱然有時我們.
(賴麗櫻)會有不同的意見.
(賴麗櫻)有不同的觀點.
(賴麗櫻)有不同的喜好.
(賴麗櫻)但我們只有一個目的.
(賴麗櫻)就是竭力.
(賴麗櫻)來慶旺主的福音.
(賴麗櫻)近來無論是佳音樓.
(賴麗櫻)還是錦繡堂.
(賴麗櫻)我們最熱門的話題.
(賴麗櫻)就是佳音樓的續租.
(賴麗櫻)在前兩個星期.
(賴麗櫻)有一個很溫馨.
(賴麗櫻)很好的彼此分享.
(賴麗櫻)但在不同團體裡.
(賴麗櫻)我說這不只是佳音樓的事情.
(賴麗櫻)可能錦繡堂說不關我們事.

$^{561}$(賴麗櫻)不只是我們佳音樓的事情.
(賴麗櫻)同樣是我們錦繡家的事情.
(賴麗櫻)所以前兩個星期.
(賴麗櫻)在佳音樓有集思會.
(賴麗櫻)有分享會.
(賴麗櫻)在10月1日就輪到我們錦繡堂.
(賴麗櫻)我不知道錦繡堂會有多少人來.
(賴麗櫻)可能很多人覺得不關我們事.
(賴麗櫻)但我說不是.
(賴麗櫻)我們是同一個屬靈的家.
(賴麗櫻)大德生的.
(賴麗櫻)不知道是曾孫還是其他.
(賴麗櫻)大少曾牧師.
(賴麗櫻)他是外國人.
(賴麗櫻)但他講道的時候.
(賴麗櫻)你知道他經常說什麼嗎.
(賴麗櫻)我們中國人.
(賴麗櫻)由他的大德生牧師.
(賴麗櫻)他已經投入了中國的時風.
(賴麗櫻)竟然去到曾孫的時候.
(賴麗櫻)他講道不是你們中國人.
(賴麗櫻)而是我們中國人.
(賴麗櫻)所以我們覺得在這裡.
(賴麗櫻)我們無論嘉欣樓.
(賴麗櫻)無論錦繡堂.
(賴麗櫻)我們一起為了一個續租的事情.
(賴麗櫻)我們來禱告.
(賴麗櫻)我們去尋求上帝的心意.
(賴麗櫻)有人會從金錢來看.
(賴麗櫻)有人會從.
(賴麗櫻)嘩!這裡的菲姐姐,欣月姐姐.
(賴麗櫻)我們嘉欣樓的弟兄姊妹.
(賴麗櫻)我們一起很努力.
(賴麗櫻)不斷做這個工作.
(賴麗櫻)大家從不同的角度來看.
(賴麗櫻)但最重要的是.
(賴麗櫻)從上帝的心意.
(賴麗櫻)所以弟兄姊妹.
(賴麗櫻)我們在這裡一起.
(賴麗櫻)我們去尋求上帝的心意.

$^{601}$(賴麗櫻)我們很盼望.
(賴麗櫻)我們在和諧和平的氣氛下.
(賴麗櫻)上帝必定為我們開出一條最美好,最好的路.
(賴麗櫻)我們最重要的是互相的欣賞.
(賴麗櫻)互相的信任.
(賴麗櫻)互相的提醒.
(賴麗櫻)互相的遷就.
(賴麗櫻)上帝一定給我們一條最美好的出路.
(賴麗櫻)我們同心禱告.
(賴麗櫻)因為你給我們每一個屬於你的人.
(賴麗櫻)我們都要同心合意,興旺福音.
(賴麗櫻)我們看到佳恩的弟兄姊妹.
(賴麗櫻)他們在這個福音的事情.
(賴麗櫻)在大令敬拜的事情裡面.
(賴麗櫻)他們盡心盡意來服侍你.
(賴麗櫻)主要是在前面.
(賴麗櫻)我們正在面對的這個題目.
(賴麗櫻)我們求你給我們有指引.
(賴麗櫻)也幫助感孫家的弟兄姊妹.
(賴麗櫻)我們不覺得這個不關我們事.
(賴麗櫻)相反,我們同樣緊張處理國度的事情.
(賴麗櫻)主要知道金錢經濟不能夠破壞任何事情.
(賴麗櫻)只要是上帝你心意的時候.
(賴麗櫻)主你必定成就.
(賴麗櫻)主就求你給我們.
(賴麗櫻)我們每一個人有這份信心.
(賴麗櫻)上帝必定為我們安排最美好的.
(賴麗櫻)來興旺福音.
(賴麗櫻)我們奉靠主基督成名.
(賴麗櫻)阿們.
\newpage



\section{哥林多後書 5:17-18}
\label{sec:CTfZJ96vxzI}
\textbf{2023-09-24 早堂《重生》哥林多後書 5\_17-18 講員:黃頌言姑娘}
\newline
\newline
連結: \href{https://youtube.com/watch?v=CTfZJ96vxzI}{\texttt{https://youtube.com/watch?v=CTfZJ96vxzI}} ~~~~ 語音日期: 2023-09-26
\newline
\newline
\hyperref[sec:B_Q_tzD4YvY]{\small{< < < PREV SERMON < < <}}
~
\hyperref[sec:index_chronic]{\small{[返順時目]}}
~
\hyperref[sec:index_scriptual]{\small{[返順卷目]}}
~
\hyperref[sec:QkQu6jc1ZPg]{\small{> > > NEXT SERMON > > >}}
\newline
\newline
哥林多後書 5:17-18
\newline
\begin{longtable}{cl}
\hline
\hline
章節 & 經文 (和合本修訂版)\\
\hline
5:17 & \begin{tabularx}{0.7\textwidth}{X} 所以，若有人在基督裡，他就是新造的人，舊事已過，都變成新的了。 \end{tabularx} \\ \\ \relax
5:18 & \begin{tabularx}{0.7\textwidth}{X} 一切都是出於神；他藉著基督使我們與他和好，又將勸人與他和好的使命賜給我們。 \end{tabularx} \\ \\
[1ex]
\hline
\hline
\end{longtable}
$^{1}$各位弟兄姊妹 各位家長 各位同學.
小朋友 平安.
不知道大家有沒有聽過一句話.
叫做 我是不是有一個powerpoint.
對 就是這句話.
西方是沒有垃圾的 只有放錯地方的資源.
世界上是沒有垃圾的 只有放錯地方的資源.
這句話很有智慧.
我們眼睛看它好像垃圾的東西.
原來只要你有創意 只要你有眼光.
就可以找到它的用處.
重點其實是我們的眼睛怎樣看事物.
這次的花燈作品.
品質非常高.
我們的設計比賽.
你們看看旁邊 看到大家的作品很厲害.
父母和小朋友合作無間.
我特別欣賞一些用到環保物料特質的作品.
讓我們欣賞一下.
我第一次看到這條魚就覺得不少東西.
因為它是利用了水樽的形狀.
其實它是切開了水樽.
用了水樽的形狀和食物網袋的特質.
做了金魚尾巴的感覺.
很輕的感覺.
好像在水裡.
如果加上燈就更加好.
其實真的不容易.
我們再看.
我看到這隻水母的時候畫得很厲害.
上半圓的透明蓋.
可能是用不要的膠蓋.
加上下面的絲帶.
無論是形狀還是透明的部分.
都能夠完全用到物料的特質.
呈現出水母的感覺.
真的很不容易.
很有創意.
原來世上是沒有垃圾的.
只有放錯地方的資源.

$^{41}$如果我們有這個眼界.
所有的物品都可以物盡其用.
神創造這個世界.
每一樣東西都很美.
原來神創造每一天之後.
祂都說.
這東西是好的.
這是好的.
祂每次創造完一樣東西.
祂都會有一個評價.
這東西是好的.
祂看一切都是好的.
所以無論大自然裡.
每一樣東西.
我們都要去學習欣賞.
要珍惜愛護大自然.
神創造第六天.
是創造男造女.
看到圖畫裡面有男有女嗎.
還是只看到動物.
看到兩個人.
看到嗎.
有兩個人.
神做男做女之後.
都有一個評價.
祂說是非常好的.
在祂眼中.
祂看我們每一個都是好的.
但有時候.
可能我們會因為和別人比較.
慢慢就會覺得自己不夠好.
我學東西好像沒有人這麼厲害.
沒有人這麼快.
我數學計算沒有人這麼厲害.
語文又沒有這麼厲害.
覺得自己未必有別人這麼強.
也沒有人這麼有創意.
就好像這次的花燈比賽.
別人的作品真的很美.
會不會不懂得去欣賞自己呢.

$^{81}$又或者會不會.
我們曾經做錯的事.
原諒不了自己.
接受不了自己曾經犯錯.
或者接納不了自己的軟弱呢.
我們每一個都不是完美的.
但是我們在神眼中.
都是好的.
我們做一個好.
給自己一個叻.
每一個都是好的.
給旁邊那個叻.
我們每一個都是好的.
是的.
沒錯.
今天我們的經文是這一句.
所以若有人是基督女.
他就是新造的人.
舊事已過都變成新的了.
聖經告訴我們.
或許我們過往有很多不好的事.
我們都有軟弱的地方.
但原來我們可以重新開始.
只要主耶穌基督.
只要我們在主耶穌基督.
我們的拯救者裡面.
我們軟弱都可以變成剛強的.
Odi姑娘小時候.
讀書成績其實不是很好.
很努力很努力讀書.
但成績才算是平平.
英文更加差.
小時候我很怕墨書.
有沒有人很怕墨書的.
都有些人不怕墨書.
真是好了.
我不孤單了.
以下是Odi姑娘的黑歷史.
小朋友記住不要學.
其實我從幼稚園到小學的時候.

$^{121}$都有一個不好的行為.
小時候沒有人教我.
我已經懂得出貓.
即是偷看隔壁的同學.
然後偷看了答案.
寫在自己的墨書本.
這個是不好的行為.
我們先給一個大交叉.
不好的跌跌.
千萬不要學各位小朋友.
但是呢.
因為原來呢.
我有這個出貓的行為.
我後來發現原來是因為我接受不到自己.
接受不到自己只有70分.
當自己有70分的時候.
接受不到自己只有80分.
接受不到.
當我有80分又接受不到自己沒有100分.
那怎麼辦呢.
我才知道原來我在教會裡面.
我聽到原來天父不喜歡人不誠實.
聖經裡面這樣說.
耶穌教我們.
你們的話是就說是.
不是就說不是.
再多說就是出於那惡者.
即是魔鬼.
原來聖經教我們.
你是的就說是.
不是的就說不是.
不可以再繼續去狡辯.
我就發現原來主耶穌不喜歡我不誠實.
於是呢.
我從此呢.
就.
我去反思自己.
那為什麼我要這樣做呢.
那我就不誠實了.
當我靜靜看別人的答案.

$^{161}$我不是靠自己的時候.
於是呢.
我慢慢就不再出貓了.
因為我不想天父不喜歡.
我慢慢接納了自己.
盡了力溫習.
依然.
可能我盡了力.
不懂就不懂.
我的成績多好都不要緊.
因為這是神創造的我.
從前可能我只是看自己的分數.
現在呢.
我是重視神怎樣去看我這個人.
所以呢.
若有人在基督裡面.
他就是新造的人.
舊事已過都變成新的.
我呢.
將自己放在基督的裡面.
用基督的眼睛去看自己.
他看原來我們每一個都是好的.
耶穌欣賞我們盡力做好每一件事.
我的分數就是我的分數.
要好好地接納我們自己這個自己.
一切都是出於上帝.
他藉著基督使我們與他和好.
又將勸人與他和好的使命賜給我們.
神創造一切都是好的.
但是當我們犯罪的時候.
原本我們是與神一起的.
當我們犯罪的時候.
原來罪就將我們與神分開了.
剛才我講的出貓.
即是不誠實.
原來讓我心裡面很沒有平安.
這個罪讓我與天父完全地分開了.
但是聖經說主耶穌的眼界.
幫助我接納自己.
使我有勇氣做對的事情.

$^{201}$讓我與天父和好.
雖然這件事看似是小學生的一件小事.
但這件事其實是我人生裡面.
第一次被神的話語去改變.
這是我人生第一次的經歷.
成為新做的人.
使我可以與天父和好.
做對的事情.
所以若有人在基督的裡面.
他就是新做的人.
舊的事已經過去.
可以變成新的.
重點是我們要常常.
將自己放在基督的裡面.
要記得天父的話.
要用天父的眼界去看自己.
就可以成為新做的人.
無論我們現在幾多歲.
可能我們三歲兩歲.
還是我們已經長大了.
無論幾多歲.
我們都可以成為新做的人.
只要我們在基督的裡面.
我們怎樣才可以不變舊呢?.
就是不停地被神的話語去更新.
我們回來教會.
回來主日學.
去學習聖經.
一切都是出於上帝.
祂藉著基督.
使我們與神和好.
有勸人與他和好的使命賜給我們.
學業當然是很重要的.
但我們成為怎樣的人.
是更加重要的.
我們需要與神和好.
若然我們可以坦然地.
活在神的面前.
我們更加有使命.
使人與神和好的.

$^{241}$在學校裡.
我們怎樣與人相處呢?.
我們要將自己.
生命的改變.
經歷.
與人分享.
做父母的.
我們先自己去經歷生命的改變.
我們要學習接納.
做父母都有自己的限制.
我們要將子女帶到天父的面前.
像剛才凱晴姨姨分享.
她也覺得自己很有限.
但只能將小朋友帶到天父的面前.
我很感恩.
天父從小就將我.
我的父母將我放在教會裡成長.
讓我從小就被神的話語去教導.
今天我們透過花燈比賽.
將廢棄的資源.
可以學習重新再做.
用新的眼光去看這些資源.
讓他們重生.
今天我們也用神的眼光去學習.
怎樣用神的眼光去看自己.
接納自己.
將自己放在基督的裡面.
記得勝經的話.
生命改變與天父去和好.
讓我們有力量與別人去分享.
使別人也能與天父和好.
誰也想經歷在天父裡面的重生.
請舉手來看看.
誰也想經歷在天父裡面的重生.
我們在天父的面前納旨.
讓我們也成為新做的人.
在天父面前.
我們一起祈禱吧.
我們合手閉眼睛祈禱.
親愛的天父.

$^{281}$我們很想經歷在你裡面的重生.
你剛才看到我們有舉起手.
沒有舉起手.
心裡面所想的.
你也知道.
天父謝謝你.
創造這個世界一切都是美好的.
求你在我們的生命當中.
更新我們的眼界.
使我們能夠用你的眼光去欣賞自己.
接納自己.
好像我們欣賞自己花燈作品一樣.
也能夠將你的話語放在心裡面.
使我們成為合你心意的人.
我們祝福每一個家庭.
能夠認識你.
能夠經歷在你裡面的重生.
多謝你.
奉主耶穌基督明教.
阿們.
謝謝大家.
\newpage



\section{希伯來書 4:1-11}
\label{sec:QkQu6jc1ZPg}
\textbf{2023-09-24 午堂《兩個安息》希伯來書 4\_1-11 講員:楊慶球牧師}
\newline
\newline
連結: \href{https://youtube.com/watch?v=QkQu6jc1ZPg}{\texttt{https://youtube.com/watch?v=QkQu6jc1ZPg}} ~~~~ 語音日期: 2023-09-27
\newline
\newline
\hyperref[sec:CTfZJ96vxzI]{\small{< < < PREV SERMON < < <}}
~
\hyperref[sec:index_chronic]{\small{[返順時目]}}
~
\hyperref[sec:index_scriptual]{\small{[返順卷目]}}
~
\hyperref[sec:t6pbod4Qek8]{\small{> > > NEXT SERMON > > >}}
\newline
\newline
希伯來書 4:1-11
\newline
\begin{longtable}{cl}
\hline
\hline
章節 & 經文 (和合本修訂版)\\
\hline
4:1 & \begin{tabularx}{0.7\textwidth}{X} 所以，既然進入他安息的應許依舊存在，我們就該存畏懼的心，免得我們 中間有人似乎沒有得到安息。 \end{tabularx} \\ \\ \relax
4:2 & \begin{tabularx}{0.7\textwidth}{X} 因為的確有福音傳給我們像傳給他們一樣；只是所聽見的道對他們無益，因為他們沒有以信心與所聽見的道配合。 \end{tabularx} \\ \\ \relax
4:3 & \begin{tabularx}{0.7\textwidth}{X} 但我們已經信的人進入安息，正如神所說：「我在怒中起誓：他們斷不可進入我的安息！」其實造物之工，從創世以來已經完成了。 \end{tabularx} \\ \\ \relax
4:4 & \begin{tabularx}{0.7\textwidth}{X} 論到第七日，有一處說：「到第七日，神就歇了他一切工作。」 \end{tabularx} \\ \\ \relax
4:5 & \begin{tabularx}{0.7\textwidth}{X} 又有一處說：「他們斷不可進入我的安息！」 \end{tabularx} \\ \\ \relax
4:6 & \begin{tabularx}{0.7\textwidth}{X} 既有這安息保留著讓一些人進入，那些先前聽見福音的人，因不信從而不得進去， \end{tabularx} \\ \\ \relax
4:7 & \begin{tabularx}{0.7\textwidth}{X} 所以神多年後藉著大衛的書，又定了一天—「今日」，如以上所引的說：「今日，你們若聽他的話，就不可硬著心。」 \end{tabularx} \\ \\ \relax
4:8 & \begin{tabularx}{0.7\textwidth}{X} 若是約書亞已使他們享了安息，後來神就不會再提別的日子了。 \end{tabularx} \\ \\ \relax
4:9 & \begin{tabularx}{0.7\textwidth}{X} 這樣看來，另有一安息日的安息為神的子民保留著。 \end{tabularx} \\ \\ \relax
4:10 & \begin{tabularx}{0.7\textwidth}{X} 因為那些進入安息的，也是歇了自己的工作，正如神歇了他的工作一樣。 \end{tabularx} \\ \\ \relax
4:11 & \begin{tabularx}{0.7\textwidth}{X} 所以，我們務必竭力進入那安息，免得有人學了不順從而跌倒了。 \end{tabularx} \\ \\
[1ex]
\hline
\hline
\end{longtable}
$^{1}$大姐妹早晨.
每次回來宣導會講到我都feel at home.
我穆少驊還戴著宣導會領帶.
今天能夠在主日裡跟大家一起思考 成立話語.
剛才那段經文其實是很難解的.
我相信你看完之後或許不是太明白.
所以我必須要用一些時間解釋.
這段經文跟第三章所講很有關係.
我們沒時間再看第三章.
我將這段經文的主要原因.
為何會寫這段經文.
這段經文提到有兩個安息.
第一個安息是以色列人沒有去到 沒有進入.
然後上帝再提另一個安息.
這點剛才你讀經文的時候大致上可以掌握得到.
我們現在就問.
為何上帝在路中起誓不讓他們進入安息.
這段經文我們不再讀了.
這段經文提到原來以色列人在埃及圍爐的時候很慘.
他們沒有身份 他們是廉價勞工.
他們被埃及人虐待 被法律當他們是牛是馬.
結果他們向上帝呼求 上帝聽到他們的哀聲 看到他們的辛勞.
就差遣摩西帶領他們離開埃及.
就應許說有一個流乃與物之地.
你進入那裡 離開埃及 去到那裡 你就可以得到安息.
這段歷史大家很熟悉.
於是他們就離開埃及 經歷很多神跡奇事.
但是還沒有進入流乃與物之地.
上帝就已經發怒了 為什麼呢.
因為在曠野的時候 他們發現很艱苦.
沒有東西吃 神給他們嗎哪.
沒有肉吃 神給他們鵪鶉.
但是他們吃到嘴都歪了.
他們覺得在曠野的生活不好.
後來上帝樣樣要什麼給什麼 去到哪一界.
於是就跟他們說 你派探子進入這個地方.
看看是不是流乃與物之地.
探子看完之後 回來就說 果然是流乃與物之地.
但是他們的報告 全部都不是說流乃與物之地.
是說那些人又大又厲害 科技又進步.

$^{41}$因為迦南人已經懂得煉鐵.
以色列人不懂 他們最厲害的是野煉銅.
熨熨鐵 沒得打的.
他們城有高 有鐵車 他們進去死的.
然後他們就說 我們為什麼要出來呢.
在埃及就好 我們出來當炮灰.
於是就埋怨了.
他們每一次遇到困難 他們都有一個結論.
我們留在埃及好 不要出來.
你記不記得 在曠野他們沒東西吃的時候.
他們怎麼跟摩西說 在埃及就好.
埃及我們吃火鍋 有黃瓜 有和牛 有帶子.
多好 他們出來沒東西吃 要回去.
現在遇到困難了 更加說我們要回去.
我們不去了 我們自己找個領導人帶我們回埃及.
上帝真的很憤怒 上帝有很多神蹟奇事.
帶他們出來 出來臨門一腳 去到邊境都說不去.
於是希伯來書三章第十九節就說.
他們不能進入安息 是因為不信的緣故.
不相信 上帝很憤怒.
於是神就說 在曠野漂流四十年.
這一代的人死了之後 新一代才可以進去.
於是他們進入去四十年之後.
進了迦南地 應該是享安息.
但結果有沒有享安息呢 沒有.
進了迦南地都沒有享安息 為什麼.
他們進了迦南地 第一件事就是他們發現.
那個地方和他們想像的不同.
他們很害怕 他們的生產不夠.
很快就會看到 他們就遷就迦南人.
或者追隨迦南人去拜偶像.
他們拜巴力 巴力管理農業的神.
拜亞瑟拉 亞瑟拉是掌管生育的神.
他搖滑一個牧羊的神 說我們來到這個新的地方.
他害怕他們地不夠生產.
他害怕他們畜牧 畜牲不夠強盛.
於是很快就拜偶像 整個舊約是說什麼.
一路就拜偶像 所羅門 我最近看列王記.
所羅門在他最興盛的時候 上帝興起很多東西.
起了聖殿 但所羅門做什麼呢.

$^{81}$娶了外邦的女子 因為有政治的利益.
更加就是他娶了外邦的女子 連外邦的神都引過來.
他以為多一些的神可以保護自己的國家.
為什麼 不可以安息 不平安.
他覺得神不夠能力 拜偶像.
整個舊約 你會看到 由列王記直到被擄的時候.
上帝都是責備這件事.
他就犯奸淫 犯奸淫就是拜其他的神.
所以他們不可以安息.
為什麼會有這麼大的落差呢.
為什麼進了迦南地之後 他們仍然懷念埃及呢.
仍然不相信迦南地可以讓他們安生立命呢.
其實摩西是沒有騙他們的.
他們進迦南地之前 摩西跟他們說了一番話.
這番話很清楚 什麼叫做流乃與物之地.
流乃與物之地 我們就想著 哇 太好了.
去到躺平不用做 哇 真是飯來張口.
樣樣東西就神供應了 比以前更好.
摩西說不是 你看摩西在申命記十一章.
字可能細些 遠些未必看到 我都開始看不到.
雖然我驗眼 驗眼師說我100分.
申命記那裡 這樣不行.
所以 都不要緊.
申命記十一章有聖經你可以看.
所以你們要守我今日吩咐.
摩西一句話很有趣.
你進去的地方是流乃與物之地第十節.
你要進去得遺孽的那地.
本不像你出來的埃及地.
你在那裡殺種 用腳挑冠.
像挑冠菜園一樣 你們要過去得遺孽的那地.
乃是有山有谷 雨水滋潤之地.
是耶和華你神所眷顧的.
從歲首到年終 耶和華你神的眼目.
時常看顧那地 你想一想 摩西這番說話很有玄機.
他們進迦南地之前 他們想去流乃與物之地.
聽起來流乃與物很豐富.
很好的 摩西說.
小心 你去的地方.
和你出埃及的地方不一樣.

$^{121}$你去埃及的地方 用腳挑冠.
用腳挑冠 像挑冠菜園一樣.
他們住在埃及尼羅旁的歌山地.
很多水 像杭州 蘇州一樣.
河流遍滿 耕種根本不用想.
用腳挑冠 水就來了 滾滾而來 種什麼有什麼.
但去到迦南地 摩西說有山有谷.
我問你 耕田比在山上好.
還是在河邊好 你自己選.
當然河邊好 有山有谷.
還要仰望天下雨才有水.
所以摩西沒有騙他 你看看.
這就是住的歌山地 很多水.
以前都是這樣 樹木清脆.
因為尼羅河水源豐富.
生育更多 人多得不得了 連法老都害怕.
流乃與物之地 有沒有去過以色列旅行.
去過吧 黃昏昏的 沒有雨 乾乾的.
這個就是流乃與物之地.
摩西沒有騙他 他們這個地方有山有谷.
這個是今天拍的 二千年四千年都沒有變.
所以摩西說什麼呢 他說流乃與物之地.
是上帝看顧 按時下雨去支援的地方.
但如果你不聽話 不聽上帝的話 上帝收回沒有雨.
那你就什麼都沒有 在歌山地不用祈禱.
不用敬拜上帝 樣樣都有 靠自己搞定.
進到迦南地一定要敬畏上帝 倚靠神.
那裡就是流乃與物 流乃與物不是因為地好.
是因為上帝同在 上帝應許.
在今天你會看到 以色列的地沒有改變過.
埃及的地都沒有改變過 但今天以色列和埃及.
誰進步 誰先進 以色列.
經過亡國二千年之後 他們悔改.
上帝讓他們重新立國.
可以荒涼的地 成為一個很富足的.
以色列很乾燥 但他們有科技去耕種.
我們不說這些 我們說回原來在整流乃與物之地.
他們不能安息 他們看地 不是看神.
他們看到地這樣 進去就很害怕.
拜巴黎 拜亞瑟拉 拜很多神.

$^{161}$他們覺得上帝保佑不了他們.
上帝沒有辦法使他們很豐足.
由所羅門開始 歷代的君王不斷的拜候像.
因為他們信不過上帝 心不安息.
所以第一個安息 進入了迦南美地.
去到流乃與物之地 沒了上帝 就不是流乃與物.
然後 摩西很清楚說.
他說 你聽我的誡命 我就可以使這些地方成為流乃與物.
按時降秋月 但他們沒有信心.
結果那個地方就成為了上帝懲罰他們的地方.
剛才我們讀的經文就提到另一個安息.
就是說 進入迦南地 他們不能安息.
於是上帝就說 我給另一個安息.
這個安息是要憑信心進入的.
如果約書當時帶你進來.
如果你們相信上帝 你們就已經得到安息.
但因為你不得安息 所以上帝就應許.
上帝很好 所以他不聽話 敬拜偶像.
離開他 甚至後來國家都亡了.
被巴比倫人俘虜了他.
他悔改回來 再重建 神就說 我給另一個安息你.
這個安息是為神的子民 為相信神的人而設.
這個安息是什麼呢 已經不是一個地方的問題.
不是空間問題 而是什麼呢.
這個安息就是在上帝的懷裡.
因為那個地之所以是流乃與物 不是那個地的原因.
是因為上帝祝福 是上帝與你同在.
所以那個地方就是流乃與物.
他們之所以得到安息 因為他們知道他們安穩在神的手裡.
無論在哪裡都不要緊 不一定是加拿大 不一定是埃及.
去到哪裡 上帝的同在就是流乃與物.
耶穌基督在一個風浪裡 為什麼能夠一點恐懼都沒有.
因為他知道生命在神的手裡.
悶到很害怕 風來了 浪來了 我們死了.
去找耶穌 耶穌睡覺.
安息是什麼呢 剛才聖經說 上帝揭了他手的弓.
上帝在他創造完之後 他享受自己所做.
心滿意足 一切在他掌管下.
這個就是安息.
我們揭了弓 不是說我們什麼都不做 躺平.

$^{201}$而是說我們不再憂慮 我們不再擔心.
因為在上帝的帶領下.
所以安息不是一個空間的安息 而是心靈的安息.
希伯來書裡很清楚說.
現在神的子民就是要憑信心 努力進入這種心靈的安息.
這種心靈的安息就是透過基督的救贖.
約書亞的那個字其實就是新約的耶穌.
當我們回到香港.
在香港 過去很多年.
這個成為我們出生成長的地方.
60年代的暴動 很多人離開了香港.
香港是一個很美麗的地方.
我們在裡面享受神給我們生活.
但是一直不穩定.
所以不斷有人離開香港 有人回來.
其實不重要.
因為聖經的教導最主要不是在哪個地方.
而是上帝與你同在.
所以當我讀希伯來書的時候.
特別是在去年前年.
我一直思想上帝有什麼訊息給我.
當我回到香港的時候.
我發現很多地方的人心中很不安.
很想離開香港.
不要緊 離開香港 神祝福你.
但最重要的是你知道要做什麼.
留在這裡的我們也一樣心中平安.
上帝掌管.
我們去哪裡也好.
去英國也好 去加拿大也好 去哪裡也好.
祝福每個弟兄姊妹.
但重要的是我們選擇留在這裡.
或者我們被迫留在這裡.
我們心中要平安.
因為這個地屬於上帝.
流乃與物之地.
事實上在今天.
當我說到兩個安息的時候.
正正就是今天香港的信徒.
我們需要的.

$^{241}$如果我們心中不平安.
就算在這裡生活 我們不快樂.
但如果我們知道.
我們在這裡 我看到上帝同在.
過去六零年代 八九年.
還有很多很多的.
都在上帝的歷史引領下.
香港是一個很奇妙的地方.
香港人口有七百多萬.
但信眾 基督徒和天主教徒.
加起來大概有八十多萬.
佛教有一百萬 道教有一百萬.
其實超過四百萬人沒有信仰.
基督徒和非基督徒的比例是一比八.
也就是說八個不信耶穌.
一個信耶穌.
所以今天我們在香港.
神給了我們很多機會.
我們不是單單用一句話.
說沒辦法 要在這裡生活.
不是 原來這是神給我們的地方.
祂的美好.
因為我們認定.
神與教會同在 與我們同在.
2020年.
香港有十八萬人離開香港.
按照英國的統計數字.
大概有十八萬人.
不過離開香港的比例.
基督徒是特別多的.
基督徒是25\%.
在英國他們的統計.
那代表什麼呢.
代表就是.
這個我一會兒再說.
這個就代表了.
原來香港有很多機會.
有很多地方基督徒離開了.
我們剩下的基督徒.
其實這個就是一個機會給我們.

$^{281}$我們要為主去填補這個數字.
我們能夠在這裡生活.
掌權的在上帝手中.
君王在上帝的手中.
如果我們認定這個地方是神給我們的.
我們就好好生活下去.
我們繼續去見證.
我們繼續的去福音上見證.
有人說香港好像沒有什麼自由.
但我想跟大家說一個事實.
其實香港的公共空間.
是有很多福音機會.
我以下所說.
你不要誤會.
我不是反美人士.
我不是反外國.
我在外國生活超過20年.
我明白.
不過在外國有很多自由.
但有一樣東西大家未必知道.
我舉個例子.
如果今天香港政府說.
所有中學不准傳福音.
不准牧師講道.
如果今天香港政府說.
你會覺得香港沒有宗教自由.
你會覺得死了匪夷所思.
但是我想跟你說.
你自己fact check.
原來美國在20世紀.
有一個叫ACLU.
American Civil Liberty Union.
它推行一個極端的.
叫政教分離.
這個政教分離就是說.
所有政府裡面是不可以有宗教.
包括在政府裡面不可以慶祝聖誕節.
只有season greeting.
加拿大已經是了.
美國也是.

$^{321}$我不是說它沒有自由.
就是政教分離這個政策.
這個政策就是.
所有政府下面的教育機構.
大中小學.
是不可以有宗教.
即是不可以有基督教.
原因就是說.
因為如果有基督教.
就會offend了.
是會開罪了其他宗教人士.
我舉個例子.
我自己的例子.
我在屋泰華牧會.
我參加chaplain.
醫院的chaplain.
就是院務.
要有個訓練的.
有個神父帶領.
要上課.
我一路參加.
到參加完了.
我就問一個問題.
為什麼在整個課程裡面.
完全沒有提到.
就一個牧師.
在醫院裡面.
在一個臨危的病人.
給他傳福音.
你以為負責.
負責的是神父.
他很激動.
Jason don't impose your fate to the clan.
很惡.
知道我說什麼嗎.
他說你不能將信仰.
強加於其他人身上.
This is law.
法律來的.
那時候很惡.

$^{361}$我就說.
作為牧師.
為什麼不可以向病人傳福音.
不可以.
他問你可以.
你不可以.
你自己主動說.
我要傳福音給他聽.
我帶你祈禱.
我帶你信耶穌.
不可以.
當時我很生氣.
我說我不參加了.
我會去傳福音.
Show me.
你抓我吧.
不過我傳福音用廣東話說.
他不懂得聽.
我繼續.
我探訪一些華人.
我繼續傳福音.
我想告訴大家.
其實香港.
暫時來說.
起碼有自由.
將來我不知道.
我沒有水晶球.
能夠有機會.
神給我們留了很多機會.
教會很多事情可以做.
我們可以向姐姐傳福音.
多難得.
在其他地方不可以.
在阿拉伯國家不可以.
我剛剛去杜拜旅行.
探訪一位孫教師.
是我教會.
在那裡超過二十年.
他是做物理治療在醫院那裡.
他很想向阿拉伯伊斯蘭教傳道.

$^{401}$我問他二十年.
你有多少果子.
我只問一個問題.
他跟我說.
Zero.
零.
為什麼.
因為他進去工作.
阿拉伯人當他是菲傭.
當他是工人.
沒有機會說.
他只能夠跟其他的外僱傳福音.
真的要跟阿拉伯人說.
是沒有機會.
二十年過去.
沒有一個阿拉伯人.
他可以向傳福音.
香港可以.
阿拉伯人這裡可以向傳福音.
這些數字.
這些你fact check.
有一次我說.
有人說我抹黑美國.
我這些不fact check就說我抹黑美國.
這個law來的.
美國的supreme court.
1962年.
supreme court通過了.
禁止全美公立學校念誦主土文.
有一次我查過.
原來是被一個猶太人告的.
猶太人說.
我兒子一聽到耶穌就頭痛.
念主土文他就頭痛.
於是他說.
這個是公立學校.
為什麼你要念主土文.
是應該沒有宗教的.
他打官司打了幾年.
打到最高法院.

$^{441}$到了1962年.
最高法院下令.
全國的學校不准念主土文.
要尊重其他族裔.
1963年.
最高法院禁止學校任何形式宣讀聖經.
因為有其他宗教的人.
有不信的人.
不可以在周圍讀聖經.
任何地方.
1980年.
最高法院禁止學校張貼十誡.
連道德上的十誡都不准.
因為這是利治聖經.
1981年.
這件事就大了.
美國原來很多大學是教會出展起.
但後來就成為了公立大學.
好像香港浸會大學.
崇基學院成為了政府大學.
這些大學之前是有禮拜堂.
比如崇基禮拜堂.
大家去過吧.
很漂亮.
現在最高法院就說.
所有周立大學.
原來是基督教大學.
現在成為周立大學.
所有的宗教設施要開放.
所以有一個和尚說.
要在崇基禮拜堂.
要打齋.
你要給他.
因為這個不是禮拜堂.
是一個宗教設施.
任何宗教人士是可以使用的.
1992年.
公立學校不能夠邀請基督教牧者.
在畢業禮祈禱和祝福.
不准講道.

$^{481}$不准祈禱.
不准祝福.
2000年.
最高法院禁止校際足球比賽.
公開祈禱.
這些全部是法律.
是全國性.
有一個小朋友.
在今年.
他的祈禱文.
這個祈禱文引起很多美國人的共鳴.
特別是基督徒.
美國怎麼會這樣呢.
為什麼美國號稱基督教國家.
政府去到一種所謂的.
政教分離去到這個地步呢.
其實我在加拿大的時候.
有時候也很生氣.
不過今天不夠時間說.
沒有說那麼多.
我再說.
其實美加有很多好的地方.
我那時候說溫哥華是福地.
起碼生活氣候很自由.
不過在宗教上.
是有很多限制的.
這個小朋友就說.
他的祈禱文.
他那本是英文.
你可以查看.
他說當你祈禱.
在學校是犯規的.
還有讀聖經是違法的.
他說學生什麼都可以.
染髮也可以.
婚前性行為什麼都可以.
但是信耶穌不可以.
讀聖經不可以.
祈禱不可以.
尾二那行.

$^{521}$他引用聖書賦刑責.
是刑事責任.
犯法.
未禁窗械禁聖經.
窗械禁不了.
但是聖經全禁了.
這個就是那個情況.
第一句我要解釋的.
他說物教孩兒分對錯.
判斷行為不妥當.
因為在美國和加拿大.
有一條條例.
叫non-judgmental basis.
就是說一個人向你求助.
你不能夠說他不對.
就是說一個小朋友走來.
我離家出走.
你不能夠說他不對.
你只能夠幫他.
例如他沒東西吃.
你給他餐卷.
沒地方住.
你給他地方住.
你不能夠說他不對.
這個叫自由.
這個叫non-judgmental basis.
你要幫助他.
不是責罰他.
這個就是自由.
這個就是人權.
其他的大家自己看.
我們今天無論在哪裡生活.
最主要就是耶穌基督在我們心裡.
這個是我們所得以安穩.
可以興翰世上.
我想講兩個例子鼓勵大家.
第一個就是.
這個故事都很爭議性.
因為一般在網上.
很難找到整個過程.

$^{561}$不過我看到一個記者.
一個英文的報道.
大概是這樣.
有一個紐西蘭的爬山專家.
叫希拉蕊.
她在1953年5月9日.
全世界她是第一個爬上珠穆羅馬峰.
但這件事下來之後.
就成為了很大爭議.
為什麼呢?.
因為下山之後.
因為當時的人在中途已經死了.
很艱難爬上去.
她就恭喜他.
他就說對不起.
爬到這個高峰不是我.
第一個登上山是我的助手.
就是這個尼泊爾人.
於是人們就問丹珍.
你是第一個爬上.
因為新聞是要知道.
誰是第一個踏足最高.
希拉蕊紐西蘭的爬山專家就說.
不是我.
我助手.
於是就說恭喜你丹珍.
你是第一個爬上最高峰.
她就說對不起.
不是我.
記者就說.
你說不是你.
你兩個搞什麼.
於是繼續問.
希拉蕊說不是我.
丹珍說不是我.
不過這個登山的過程.
是希拉蕊設計的.
因為她是爬山專家.
丹珍只是一個尼泊爾的導遊.
她是幫助她.

$^{601}$所以後來的功勞歸給希拉蕊.
不過希拉蕊一直說.
不是我第一個上去.
過了很多年之後.
就有記者再問丹珍.
你可不可以誠實一點說.
當時的情況是怎樣.
其實當時的情況.
已經有很多報紙都登過.
她就說.
當時的情況是這樣.
上到山的最高.
就是經歷死亡.
差很少.
就有十幾公尺.
戴著氧氣罩.
希拉蕊就是她的主人.
因為她是爬山專家.
就跟她說.
你走到那一點.
我幫你拍一張照片.
於是她一直聽她的話.
所以丹珍就聽紐西蘭的爬山專家說.
就走到最高點.
她幫你拍一張照片.
所以今天是有那張照片.
下來的時候.
她就跟我說.
恭喜你.
你是第一個爬上最高峰.
丹珍就說.
她沒有理由.
她說你用盡你的家產.
用盡你所有的錢.
用盡你所有的精力.
你都是爬上最高峰.
第一個人踏上去.
她怎會是我.
她說對不起.
這是歷史.

$^{641}$沒得改.
她說不行.
我不可以拿這個榮譽.
當時.
所以她下來.
她就說.
我不是第一個.
之後希拉蕊就上去.
在最高峰那裡.
就挖了一個洞.
就放了一個十字架下去.
證明她是上去的.
所以上面有一個十字架.
塞了下去.
下來的時候.
丹珍後來她說.
當時希拉蕊跟她說了一句話.
她聽不明白.
她不明白為什麼.
她只說.
因為她是尼泊爾人.
這個榮譽要給尼泊爾.
她只說了這句話.
但是她還說了一句話.
當時她聽不明白.
她說現在過了十幾年之後.
她明白希拉蕊說什麼.
你猜她說什麼.
原來她下來的時候.
她就跟她說.
她說我恭喜你.
你是第一個.
她說因為這個山是屬於尼泊爾.
所以榮譽是給你的.
最好的我已經有了.
所以這個榮譽歸給你.
當時她聽不明白.
什麼叫做最好的已經有了.
你最好不是要爬山嗎.
她說不是.

$^{681}$最好的我已經有了.
所以這個榮譽.
我給尼泊爾人.
她當時不明白.
所以她沒有說這番話.
到後來她明白.
最好的我已經有了.
還有什麼恐懼.
還有什麼爭.
還有什麼我們不能安息.
面對一個這麼大的誘惑.
見歷史第一個人踏上去.
她是信理成真她應該要有.
但是她說最好的我已經有了.
這個榮譽歸給尼泊爾.
心中平安.
安息.
因為我們得到神最大的救贖.
在這個救贖裡面.
我們憑著信心.
過每一天.
最好的我們已經擁有了.
安息.
不在乎地.
在乎我們的主.
我們同心祈禱.
因為你是我們的救主.
無論去到哪裡.
最好的.
你已經給了我們.
讓我們心中平安.
繼續為主你的福音.
為你的工作.
我們努力.
願你教會復興.
願多人歸順主.
祝福我們.
愛我們家人.
愛教會.
愛主你.

$^{721}$帶領教會.
繼續建聖.
祈禱.
服主名求.
阿們.
\newpage



\section{撒迦利亞書}
\label{sec:t6pbod4Qek8}
\textbf{2023-09-28 周四晚查經班《歸回先知書研讀(3)撒迦利亞書》 講員:曾裔貴牧師}
\newline
\newline
連結: \href{https://youtube.com/watch?v=t6pbod4Qek8}{\texttt{https://youtube.com/watch?v=t6pbod4Qek8}} ~~~~ 語音日期: 2023-09-29
\newline
\newline
\hyperref[sec:QkQu6jc1ZPg]{\small{< < < PREV SERMON < < <}}
~
\hyperref[sec:index_chronic]{\small{[返順時目]}}
~
\hyperref[sec:index_scriptual]{\small{[返順卷目]}}
~
\hyperref[sec:ZbYxDMl2Zbg]{\small{> > > NEXT SERMON > > >}}
\newline
\newline
$^{1}$現在我們一起祈禱.
感恩天父讓我們在晚上再次有茶經的時間.
幫助我們對於撒迦利亞書有明白.
有一個掌握.
特別是對整卷書一個大脈絡的掌握.
願神帶領我們.
讓我們有足夠的身心靈的力量去瞭解.
去明白.
奉主耶穌基督的名.
阿們.
好,我們恢復.
撒迦利亞書一到第七章.
是有八個異象.
是很工整的.
先用異象來講解.
撒迦利亞先知所得到的啟示.
然後第八章一直到十四章.
都是有另外的.
他得到的預言.
和一些教導.
所以現在就穿插在這八個異象裡面.
但是原來又是一個信息.
重點都是神重新來看顧.
或者說保守耶路撒冷.
得到的重建.
和耶路撒冷會有昔日的興旺.
我們可以透過這八個異象.
再看看神如何去關顧.
重新去賜福給祂的子民.
我們看過第一個異象.
就是有一隊騎馬的騎兵.
當然是天使.
因為是異象.
尤其是強調騎紅馬.
這樣來向神懇求.
你何時再眷顧耶路撒冷呢.
神就搭這一個騎紅馬.
另外第二個異象就是.
那個四角.
其實重點是四個像人.

$^{41}$就是說耶路撒冷.
是真真正正得到重建的.
神要來保守這個工程.
其實已經開始了.
這個本來是豆腐渣.
本來是一個爛尾的工程.
現在是重新建造的.
這是大理污黃的第二年的時間.
現在撒迦利亞就是第二年的九月.
他得到的默示.
第三個異象就是.
他見到有一個人在量度.
轉而引出神.
應許整個耶路撒冷會得到保護.
因為神就是那個火城.
就是那個城牆.
而且還是用神作為他自己的火來保護.
所以不是焚燒.
是一個保護.
是神的屬性.
這是第三個異象給我們見到的.
接著就是有那個呼喚來求神.
神就說我會成為耶路撒冷眼中的瞳人.
他是那個Apple.
就是他一定會留意管理保守這個城.
所以你會見到.
一直都是從那個環境.
從那個物質讓人去看到.
就是讀者.
看到神現在重新眷顧耶路撒冷.
就在外圍.
有些騎兵在那裡看著.
周圍的人來到走來走去.
接著有四個像人的那個回標重建.
接著又有神自己作為那個火城.
因為那個地方已經涼度了.
不需要有城牆.
因為足夠的大.
神來到去成為那個地方的保護.
一直都見到是這樣描述.

$^{81}$在那個好像讓我們能夠從那個異象裡面.
好像一幅一幅的圖畫.
一幅你能夠想象得到的那個景象.
來到第四個異象.
很特別的就是不再是物質的那個物理的環境.
就是不是外圍遠遠地看耶路撒冷.
現在是進入聖殿.
因為大祭司約書亞就在聖殿裡面.
接著第四張就是所羅巴伯.
就是說祭司和領袖首領來到去.
都被擺在異象裡面.
而且都是正面地來扶持他們.
能夠肩負起那個聖跡的工作.
第三張天使又子給我看.
所以是一連串的異象.
是一個緊接另外一個的.
所以如果之前那三個.
就是從一個很遙遠的望去耶路撒冷.
這次就不再只是普通的耶路撒冷.
而是耶路撒冷裡面最重要的事奉人員.
就是大祭司.
天使又子給我看.
大祭司約書亞站在耶和華的使者面前.
撒旦也站在約書亞的右邊與他作對.
大家想想那個畫面.
約書亞知不知道撒旦在他旁邊.
在他的右邊.
在他的左邊就是耶和華的使者.
天使是子給撒迦利亞看的.
撒迦利亞當然一邊看一邊聽.
因為第二節.
耶和華向撒旦說.
撒旦啊 耶和華責備你.
就是揀選耶路撒冷的耶和華責備你.
這不是從火中抽出來的一根柴嗎?.
到底這句話.
撒迦利亞明不明白呢?.
或者我們今天.
現在大家在現場或者在網上聽的明不明白呢?.
到底這個在火中抽出來的柴.

$^{121}$是指約書亞.
是指耶路撒冷.
還是指撒旦呢?.
我們要瞭解一下.
第三節.
約書亞穿著污穢的衣服.
在神的使者面前.
在撒旦的旁邊.
約書亞現在最吃虧的就是.
他穿著的是污穢的衣服.
污穢當然不是指那件衣服.
掉到泥土里.
弄污穢了.
或者吃意大利粉的汁.
弄到整件衣服一塌.
不是這些污穢.
這裡說的.
一定是說屬靈的層面的.
那個不潔.
那個現在暴露在神的使者.
和魔鬼撒旦的面前的那個不潔.
這就大件事了.
那就分外顯出問題所在了.
第四節.
使者吩咐站在.
面前的說.
你們要脫去.
他污穢的衣服.
又對約書亞說.
我使你脫離罪孽.
很明顯了.
就是這樣對號入座了.
污穢的衣服.
不是那件衣服不乾淨.
是那個人本身不乾淨.
現在就.
我使你脫離罪孽.
要給你穿上華美的衣服.
我說.
我說就是撒迦利亞.

$^{161}$要將潔淨的冠冕戴在他頭上.
對不起.
應該是繼續使者這樣說的.
他們就把潔淨的冠冕.
戴在他頭上.
給他穿上華美的衣服.
耶和華的使者在旁邊站立.
好像一個加冕的儀式.
相信肯定是.
整件事一定是在說.
在獻祭祭司就職.
和聖殿的那個場景.
不是現在休假.
大祭司在很遠的地方.
一定是在一個很神聖的時刻.
很神聖的地方.
大祭司他現在.
撒迦利亞見到的這個儀仗.
第六節.
就有一些教導了.
這個儀仗之後.
耶和華的使者告誡若舒雅說.
萬君知耶和華如此說.
你若遵行我的道.
謹守我的命令.
你就可以管理我的家.
看守我的園園.
我也要使你在這些站立的人中間來往.
大祭司若舒雅.
你和坐在你面前的同伴.
可能是其他的先知.
是誰也沒有說.
之前也沒有提.
可能就是說一般性的對他的告誡.
坐在你面前的同伴都當聽.
他們是作預兆.
沒錯是那些先知.
不是那些天使.
我必使到我僕人大慰的苗若發出.
這裡又很特別了.

$^{201}$說著說著明明是告誡若舒雅大祭司的.
叫他醒醒定定.
潔身自愛.
現在已經幫你搞定了.
好像那些廣告.
那些好像很簡單.
欠債.
take a take就一筆勾銷.
其實是騙人的假的.
還要多給一些.
無論是罰款或者其他的行政費.
就是說將那些債項.
轉移到另一個債項里.
根本就沒有一筆清.
現在這裡就已經說了.
就說了給若舒雅知道.
現在已經將你污穢的衣服.
脫開了.
已經使你潔淨了.
甚至給你有那個榮耀了.
官面肯定是代表榮耀.
代表那個積分.
代表他稱職的那個服侍.
這樣來幫他加冕了.
不過若舒雅要做的.
就是要看到他自己要.
這個是第六節.
第六第七節.
其實任何人都對的了.
任何一個事奉神的人.
都需要來到去留心了.
要謹守神的道.
和他的命令.
就可以管理神的家.
看守神的園圓.
神就會使得他能夠.
站在眾人當中.
來到去行走來往.
就是說在那個積分上面.
帶來他自己都會得到.

$^{241}$很大的一個祝福和益處.
會得到別人的尊重等等.
這個是神給事奉者的榮耀.
這些真是大家其實是要.
學無.
不過我們再講下去.
第八節就很重要.
突然間出了一個預言.
突然間在意象裡面.
有對將來的啟示.
相信肯定是要拉上關係.
就是大魏王朝最緊要的後裔.
那個大魏的苗裔.
也可以這麼說.
我們一讀《馬太福音》就已經知道了.
阿巴拉罕的後裔.
大魏的子孫.
耶穌基督的家譜.
所以三章八節.
撒迦利亞已經有這樣的一個意見.
神已經講講若順亞.
就講到關於大魏的苗裔.
所以這一卷書是充滿了安慰.
因為是重建耶路撒冷.
神再眷顧曾經被祂罰的神的子民.
現在神要再成為國家的火神.
眼中的瞳人等等.
現在不僅如此.
充斥著這一卷書.
是會有耶和華的使者.
這一位大魏的苗裔.
這一位我們上個星期講過.
就是他騎著駱驅進城的這位尼賽亞.
所以這一卷書是特別多尼賽亞的預言.
三章八節就已經有了.
我必使到我僕人大魏的苗裔發出.
看那我在若順亞面前所立的石頭.
在一塊石頭上有七眼.
可能與大相關.
因為當時他們就是重建.

$^{281}$開始有很多的建築.
每天可能他們就去工地看一下那些工程.
就不用他真摯做的.
但是真的看到那些石頭慢慢搭高了.
開始來到閩魁了.
諸如此類.
但是神的真正工程是在聖殿裡面.
是在大祭司若順亞他得到潔淨.
所謂官服完職之後.
他那個服侍帶來的百姓的潔淨.
因為大家遵行神的道.
有道可聽.
有道可傳.
令到大家都有一個歸向神的心.
當然需要慢慢累積.
但是至少開始有神的默示.
有神的道.
就帶來.
這裡告訴我們.
原來有真正的石頭.
不是普通的石頭.
這些新約出現過好幾次.
彼得自己說.
你們來到主的面前.
乃是活石.
被建造成為靈工.
但是在這一句之前.
是在說耶穌基督是活石.
祂是被檢選.
被驗證過的石頭.
這是彼得前書.
二章第四到第七節的經文.
等一下我們可以揭的.
接著就告訴我們.
下面我要親自雕刻這個石頭.
並要在一天之間.
除掉這地的罪孽.
第十節.
這些經文.
猶太人就很振奮.

$^{321}$等於我們很渴望中秋節.
明晚一家人.
拿著燈籠.
小朋友.
最好是舊款那些.
我們小時候那些.
要點蠟燭的.
要拿著它很光的.
現在這些就沒有了氣氛了.
雖然是電的環保一點.
但是氣氛會更加強烈一點.
有個白兔仔.
有個輪子可以走的.
你可以感受到.
那個是真正的太平.
真正的家家戶戶.
晚上出街.
大家很開心.
你見到那些人都很開心.
這個是真正反映社會的平安.
第十節就反映這件事.
當那日.
記住.
短期的預言.
就是他們重建.
當聖殿真的弄好之後.
他們就很興高采烈.
這個是以色列記.
我們講過的第六章.
但是真真正正的那一天.
相信是千禧年了.
相信是耶穌基督回來.
這就告訴我們.
你們各人就要請輪捨.
坐在葡萄樹和無花果樹下.
這是萬君之耶和華說的.
你可以想象到整個第三章.
結合上文的兩章.
三個的異象.
無論是整個耶路撒冷.

$^{361}$有四圍的重建.
有神的保護.
神是城的火城.
就是那些牆.
四圍的牆環.
神親自保守.
看他們的子民.
是他們眼中的同仁.
就是沒有人夠膽.
來再去搞擾以色列.
接著我們就見到.
很重要的就進入到人心了.
進入到聖殿的潔淨了.
進入到大祭司.
代表著全民族的一個憲制.
一個領袖.
第三章是宗教領袖.
第四章就是民間的領袖首領.
當時是以省長.
因為沒有君王.
亡了國.
所以所羅巴伯.
第四章又是全章講他.
還有就講聖殿裡面很重要.
就是有源源不絕的金燈台.
來去炎亮.
所以你會看到.
第三第四章.
特別那個範圍是.
聖殿的那個場景.
和事奉神的職位.
這樣來去描述的.
我們很簡單的詳細的去看一看.
這個第三章先.
緊接的.
就是好像透過一些動作.
一些的形容.
撒加尼亞是一路參與著.
在整個的密士裡面.
他是一路留意著.

$^{401}$所以到底這些異象.
是怎樣來去發佈呢.
我們新約的人就比較少.
能夠這樣看到.
比如彼得能夠在哥尼樓的.
這個小徒行轉是十章.
他真的在.
因為當時他們的樓房有天台.
就是現在我們村屋的天台.
在那裡等.
是正午的開飯.
他就上天台在那裡.
脈想一下.
突然間的異象就出現.
中午的時間有異象來到.
讓他看到.
他看到的異象就是.
天上有很多的活物降臨.
就有聲音和他說.
叫他起來.
摘了這些食物.
其實是那些不潔食物.
特意讓他看到這些.
他從來都不敢吃.
不是他不喜歡吃.
是不可以吃的.
猶太人.
但是要他吃.
就違反了他的神學.
就是他的自小猶太教的教導.
那個利尾記第十章的經文的教導.
是不准吃的.
但是呢.
神就說了兩次.
神所潔剩的.
你不可以再看為是俗而不潔.
這樣來告誡他.
就是.
打破彼得.
固有的神學的框架.

$^{441}$和那個眼界.
因為我們可以想象到彼得.
他們這群人是很特殊的.
因為見著耶穌在這個世界.
見著耶穌釘十字架.
現在就是代表耶穌在地上傳道.
所以有很多舊有的觀念.
和十字架帶來的一個徹底的改變.
神如何透過十字架的耶穌.
去拯救世人.
這群使徒.
還沒有一個真正的站身的神學.
也沒有神學院再教這些東西.
所以神就要直接給一些異象.
這些使徒.
他們就好像舊約的先知那樣.
來得到一些第一身的.
從神那裡直接而來的一些異象.
當他們消化了.
瞭解了.
明白了之後.
他們就慢慢來讀書立說.
將他們所得到的.
那個原來.
如果在耶穌的救恩之後的人.
就再沒有分什麼階層.
再沒有分什麼是潔淨不潔淨了.
唯有在救恩裡面.
才是真正的潔淨了.
彼得他們.
或者保羅他們.
就會將這些他們所領受了.
其實是他們已經.
所謂茅塞頓開的那個眼光.
來成為著作.
告訴後世的門徒.
不要學我們.
我們有我們的盲點.
有我們舊有的背景.
因為我們是猶太人.

$^{481}$耶穌一天都不來.
我們一天都不會知道.
原來有這樣的眼光.
在耶穌的救恩裡面.
去看不同的人.
原來不同了.
但是彼得他們.
真的慢慢地經歷.
他們那一代的人都還沒死.
是剛剛出來傳道的階段.
所以他們的眼光.
就非比尋常的重要了.
所以當彼得能夠.
一錘定音告訴我們.
就是神潔淨的人.
你不可以再用.
舊有的觀念去看他了.
這些其實是很難的.
人的成見.
有時牢固到一個地步.
為了捍衛那個牢固.
捍衛那個固執.
情願一錯再錯.
都可以 都有的.
一生人不肯放下自己的.
固有的固執.
都帶有人在的.
但是他們就是很感恩.
因為有神給的釋放.
除去他舊有的觀念.
然後用一個.
嶄新的救恩的眼光.
去看其他的民族.
就帶來.
這個就是福音了.
帶來他們寫的書信.
就成為一個很重要的.
教導的指標了.
因為他們正正就是經歷過.
那個神學的洗禮.

$^{521}$其實很痛苦的.
因為你舊有的東西是你的安全感.
是你的安書區來的.
是你的那個保險來的.
但是一下子打破了你.
其實你猶太人.
和外邦人是沒有分別的.
大家都在救恩裡面是沒有分別的.
你不再可以說自己是阿巴拉克的子孫.
就有一個高高在上了.
是沒有的.
救恩就是令到大家都平等的.
大家都是靠恩典得救的.
這樣就完全不同了.
所以我們看回這裡的時候.
我們進入回這段經文.
我們收窄的範圍了.
就是先知撒迦利亞.
現在被帶進去.
看一些極其深入的靈性的問題.
是神才可以解決的問題.
很特別就是.
大祭司自己可能都不知道.
原來被人在這裡平頭稟祝著.
在這個異象裡面.
他居然自身在神的使者.
和魔鬼撒旦的面前.
真的好像拍電影一樣.
有個魔鬼的出現.
有天使的出現.
當然那些是將我們人幻想的東西.
來拍成電影,電視劇.
但是這裡是真的.
最重要的那個積分.
可能白天他在聖殿指揮那些人怎麼做.
要小心搬台.
這個地方是很神聖要小心.
原來他自己成為了主角.
被人.
其實是神的使者.

$^{561}$帶先知來看聖殿裡面.
最重要的就是這位大祭司.
「天使又指給我看.
大祭司約書婭站在耶和華的使者面前.
撒旦也站在約書婭的右邊與她作對」.
很特別到底大祭司約書婭自己知不知道呢.
都不重要因為這個重點是想給撒迦利亞知道.
這個重點是透過撒迦利亞.
可能他看完這個異像之後.
第二天早上見到大祭司約書婭.
就有很多的語重心長的教導.
作為一個先知.
先知雖然不可以說凌駕了祭司和君王.
這三個職份在舊約一直出現.
其實最早是先知.
其實不是去到撒武爾的時候才有先知.
其實亞伯拉罕.
神是稱呼他是先知來的.
很特別的.
應該是詩篇的107篇.
亞伯拉罕是先知來的.
但是在歷史過程.
當然是先有祭司.
因為是神親自設立的.
在《出埃及記》的經文就已經設立.
這個阿倫作為大祭司.
一定是要阿倫的後人.
即是這個利美之派.
阿倫暗戀的兒子.
阿倫才可以成為大祭司.
其他的支派就算是利美之派的其他的族弟兄.
都不可以成為大祭司.
這個規定了.
接著輾轉到王國.
就有掃羅做第一任的皇帝.
接著就是大衛王.
你會見到是神親自來告立.
設立的.
這些某程度都是叫受高者.
都是算是一定程度的.

$^{601}$好像有少少尼賽亞的身份.
但當然不是.
但是後來與君王同時期.
就是撒母耳開始.
就正正式式由先知的制度開始.
就是代表著神.
直接得到神的啟示.
他們有一個身份就是.
向百姓傳講神的默事.
這個制度會令先知.
因為君王祭司都要聆聽先知的預言.
所以某程度其實大家都是平起平坐的.
無所謂誰高一些的職份.
但是在宣講的過程.
先知就那一刻代表著神直接宣講.
所以其中一個很重要的信息就是.
耶和華如此說.
不是我說.
是耶和華要我這樣說.
那就真的凌駕一切了.
所有人都要聽.
這個就是舊約先知的那個積分.
在這個時候.
神也都仍然透過先知.
是神自己興起的哈戈和撒迦利亞.
哈戈主要是感動他們.
和督則他們來要殷勤勤奮的作功.
只是有四個很短的信息.
但是撒迦利亞很重的是安慰.
很重的是一個的醫治.
先是那個外圍的耶路撒冷.
實質的那個物質的耶路撒冷.
但是現在第三第四章.
涉及的是最重要的兩個重要的人物.
都有在那個異象裡面出現.
你可以想象得到.
整個宏觀去看撒迦利亞書.
你就可以想象得到.
不單止是神緊張那個聖殿.
不單止是緊張他的百姓.

$^{641}$他的百姓是他眼中的同人.
更加重要是事奉神的人.
神更加的體重和緊張.
所以就有這個異象給撒迦利亞看到.
所以其實第三第四章是很好看.
和很重要的對每一個事奉者的勸勉和告誡.
所以如果我們事奉神的人.
其實應該都要看一下這些經文.
所以我們很多人其實連揭都不懂揭.
其實有可惜的.
現代的基督徒已經不讀經.
就變成了連很重要的經文.
其實這些是預言.
雖然已經在耶穌基督身上應驗了.
但是還有更遙遠的千禧年要來.
我們其實已經是.
我們雖然不是猶太人.
但是我們在神的救恩裡面.
我們是屬天的亞伯拉罕的子孫.
所以我們都是分授將來神給我們的產業.
所以某程度我們都屬於神國度的子民.
所以我們有份於這個國度的時候.
我們應該要看一下.
神如何來做他的工作.
他的性質.
他如何去運作.
這裡又很詳細地告訴我們.
重點就是有地上的敵黨.
有撒旦.
撒旦就是敵黨的意思.
encounter.
就是是都要來反對的.
真正的為反對而反對就是撒旦.
明明神是公義聖潔.
他都是不接受的.
明明他是低於神的.
明明他只不過是受造之物.
他都要來抵擋神.
這個就叫撒旦.
所以又是一切抵擋神的人的祖先.

$^{681}$雖然《憲經》沒有這麼說.
你想想撒旦站在若須雅旁邊.
來控告和抵擋若須雅.
使神要來看到若須雅的過失.
一方面是好.
就是暴露在神的面前.
那解救的辦法是什麼呢.
所以救恩就在這裡出現了.
穿著污糟的衣服.
顯出自己的不潔.
內心的罪孽.
但是神立刻賜下的.
是那個潔淨.
脫掉那件衣服的一個動作.
和再給予那個榮耀.
所以一個信了主.
認識耶穌的大愛.
救恩的弟兄姊妹.
如果只是以為是一個必然的.
是你得到恩典的.
而你忘記了帶著謙卑的心靈.
在一生去服侍神.
服侍人.
那你就根本不明白.
第三章這裡.
他脫掉了污糟的衣服之後.
有潔淨了.
有神宣告的潔淨了.
還有.
給了他一個皇冠.
就是說在服侍的過程.
他會有安拿.
這個是事奉才會感受到的一份榮耀.
所以很多人不參與教會的事奉.
其實是很可惜.
在我們人生.
其實最重要就是我們服侍神.
首先最重要就是你要找到哪位是神.
做物主.
創造這個宇宙萬物.

$^{721}$我們找到了.
就是我們這一位.
獨一的三一真神.
找到之後我們就應該要等候.
他有什麼差遣.
然後事奉.
你這個找到的過程.
雙重的找到.
然後參與事奉.
這個就是安拿了.
所以很多人不去瞭解這些預言.
這些異象.
其實就錯失了他人生.
其實最大的榮耀就是服侍神.
當然那個服侍是可以.
在很多不同的性質.
不同的場景.
不只是單純在教堂.
教會裡面的.
但是教會裡面.
更加是應該我們.
義無反顧的來服侍神.
透過服侍人而我們去服侍神.
這裡告訴我們.
撒旦的敵黨.
然後神就向撒旦說.
其實不是不知道.
若舒雅的污蔚.
不過先來抵擋撒旦.
好像是勢不兩立的.
雖然我們今天不是說天使論.
撒旦,鬼魔學這些神學.
在這個系統神學的第三.
是會說關於天使論.
關於鬼魔的一些神學.
如果大家有興趣.
將來我們可以一起去看一下.
逼我能夠再看回那些書.
找回那些經文.
和大家重組一下.

$^{761}$都是一個樂趣.
但是這裡告訴我們.
神第一時間不是針對若舒雅.
處理她的罪.
而是首先來抵擋魔鬼撒旦.
這裡告訴我們.
耶和華責備你.
就是揀選耶路撒冷的耶和華責備你.
這不是從火中抽出來的一根柴.
可能說耶路撒冷.
可能是說.
站在她面前的若舒雅.
在火中被抽出來.
就是說一條很漂亮的木.
放在火堆里燒著.
但是被抽出來.
變成了一個炭.
這些我們小時候.
用柴燒火.
我們就會明白.
現在這個世界.
香港已經不會有人用木燒火了.
除非你去露營那些.
不然一般都是電子.
石油氣.
這裡告訴我們.
很重要的就是.
在意象裡面.
神是首先讓撒迦利亞先知.
他見到的一個情景.
撒旦似乎無處不在.
在聖殿他都會存在.
我都相信的.
我們都要每次聚會祈禱.
當然不是那些疑神疑鬼.
起了雞皮疙瘩.
是不需要的.
因為是一個神聖的地方.
不過既然撒旦還是一個靈界的一個物體.
他會在很多人的心裡面.

$^{801}$會挑起一些仇恨.
報道會會令到那些人沈睡.
會令到那些人很反感.
諸如此類.
但是我們要祈禱去抵擋.
在這個時候.
耶和華的使者就代表著神.
耶和華去責備你.
我想和大家看看猶大書.
有一節經文很特別.
猶大書.
即是聖經最尾二的卷書.
因為之後就去到這個啟示錄.
猶大書這樣講很特別.
第五節.
猶大書只有一卷.
據聞就是耶穌基督的弟弟猶大寫的.
他自己也這麼說的.
耶穌基督的僕人.
雅各的弟兄猶大.
教會歷史就說是耶穌的弟弟.
當然後來信了主.
生命有改變.
又是教會的領袖.
也得到默示.
神要他寫這一卷書.
全卷書大部分都是責備.
那些存異端的人.
和他們的行徑.
他們的作風.
他們的詭跡等等.
即是有跡可尋的意思.
第八節.
這些作夢的人.
也像他們污蔽身體.
興曼主治的.
即是說.
講一些很輕佻.
傲慢的說話.
偉旁在尊為.

$^{841}$這些我們大家要很小心.
我們可以評論人.
評論一些政策.
但是我們記得.
如果我們已經變成一個.
偉旁人的人.
連尊為.
古代的人更加不敢.
但是現代.
所謂人人平等.
我們慢慢就沒有了.
這樣的人生的操守.
要很小心的.
因為下面就講了.
第九節.
天使長米加勒.
為摩西的施首.
與魔鬼爭辯的時候.
尚且不敢用偉旁的說話.
罪責他.
他就是魔鬼.
不是摩西.
天使長米加勒.
在神的極近的前面.
即是等於在李家超.
我們的特首的身邊.
永遠都是.
政務司司長.
一定是他.
財政司司長.
律政司司長等等.
因為是最重要的職位.
一定是在他身邊.
這樣問他.
現在天使長米加勒.
要幫摩西的施首.
不知什麼原因.
為什麼要搶屍體.
這些很特別的.
曾經出現在《賜經》.

$^{881}$到底是什麼事.
但是重點不是這裡.
重點是第八節裡面所說的.
要我們不要學這些異端的人.
他們那種.
那把口不收.
那把口充滿了傲慢.
充滿了污蔑邪惡.
甚至乎隨便可以偉旁人.
這就告訴我們.
天使長.
就是說.
引用一個例子.
天使長這麼高職位.
面對撒旦.
是神咒座的一個靈體.
永遠是敵對神的.
就是說.
無可救藥的了.
注定是要去永火的了.
但是.
很特別的這一節《賜經》就告訴我們.
尚且不敢用偉旁的說話去罪責他.
只是這麼說.
主責備你罷.
接著第十節就說那些全異端的人了.
而且是在教會裡面行走.
不然根本不需要說.
不是告訴你.
在外面.
哪個黑社會的大哥.
你千萬不要跟他說一些.
負氣的話.
怎麼會這麼說呢.
一定是在教會裡面.
第十節.
但這些人.
偉旁他們所不知道的.
本來不知道你就不應該論斷.
而且是用一種偉旁的說話.

$^{921}$去說一些他們自己不知道的事情.
他們本性所知道的事.
與那沒有靈性的畜類一樣.
在這世上竟敗壞了自己.
好像那些沒有靈性的畜類.
畜生一樣.
作者這麼重的說話.
說這些偉旁人的人.
好像畜生一樣.
沒有靈性的畜生.
這些是很大的侮辱.
可以想象得到.
作者他很緊張.
弟兄姊妹.
平日你的說話.
你所聽的東西.
要很小心.
這個其實我們在現代的社會.
資訊流通.
和今天真的隨隨便便.
少男少女都會說粗口.
當是很平常.
覺得是很有型.
很自由.
很奔放.
我們要很小心.
有時甚至乎教會曾經一段時間.
不是說我們的教會.
是說一般性.
有時都覺得.
那個粗口是要表達那個憤怒.
沒辦法不發洩.
但是我們要留意.
經文有這樣的教導我們.
告誡我們.
現在這裡有一個很特別的例子.
主責備你罷.
這樣的意思.
他不敢多說.
不敢說多過神.

$^{961}$神責備你.
所以現在我們代表著神.
去責備你.
撒旦.
為著那個摩西的屍首.
去爭辯.
爭辯的過程.
爭辯就最容易失控.
最容易會說多了一些.
很激氣的話.
很負氣而來到去.
又會一些言辭上.
收不回的說話.
我們都是要很小心的.
因為如果你說錯了.
說多了一些說話.
令到可能你的朋友.
可能你的同事.
如果他很受傷.
那個關係永遠都收補不回.
這個我們要很留意.
包括我們的兒女也是.
我們比較少中國人向上一代的人.
說了很重的.
很粗鄙的話.
很少的了.
那些真的是五億仔.
但是同下一代.
我們就很容易.
會有這一種失控.
所以我們要很小心.
我們回來看這裡的時候.
那個責備.
都是耶和華責備你.
不是天使代表著神.
自己說多了自己的.
加鹽加醋.
這個是我們在舊約.
包括新約都見到的靈界.
他們都有這樣的秩序.

$^{1001}$這樣的守則.
這樣來到去.
在說話上面.
不會用多了.
來到去譴責.
甚至說的是魔鬼撒旦.
你怎麼罵他.
他都罪有應得的.
但不是.
那我們就很奇怪.
為什麼靈界都會有這樣的事情呢.
也都意味著.
為什麼新約.
教導信徒這麼多關於言語呢.
為什麼總要用言調和呢.
就可以知道怎麼回答國人呢.
《牙國書》第三章就說.
舌頭是還沒有人能夠自服的一個器皿.
好像生命的輪一樣.
可以說一些很大的說話.
也可以說一些安慰平安的說話.
那就是說.
所謂一口兩舌.
新約是很緊張信徒信得主之後的言語的.
你可以想象到.
這個是.
中國人就說上游深生.
但是言語都是反映我們內心的改變.
那個價值觀.
所以內心的改變.
是反映在我們的說話.
總要找到一些機會.
去對別人有造就.
一些的說話.
就算別人真的錯了.
你也找到機會去輓回.
去安慰.
去令到他知道怎麼錯.
知道怎麼來到有一個回轉.
和不是用一句話去致死一個人.

$^{1041}$這些是我們大家都要一生學習.
真的每一句說話.
是會令到一個人得到很大的建立.
如果我們學得好.
其實是一個很大的學問.
也是我們自己一個很大的祝福.
所以很感恩就是我們大家.
至少我現在看見的.
都是一些已經成熟的弟兄姊妹.
有心地來讀經.
有系統地讀經.
正正就是我們內心.
很渴望對神有認識.
不只是聽一些好笑的例子.
因為根本都沒有例子.
只是讀經.
但是去到一些很深度的東西.
我們就要停下.
要思想.
我的說話是不是真的誠信.
真實.
和對別人真的有一個安慰.
建立和鼓勵呢.
這些是真的很重要的.
所以這裡告訴我們.
耶和華責備你.
就是揀選耶路撒冷的耶和華責備你.
再一次提醒.
耶路撒冷是屬神的.
神已經在眷顧這個城了.
撒旦是沒有你的位置的了.
沒有你的位置.
你可以隨時出現.
這些靈界我們真的不是很清楚.
隨時可以去到天庭.
控告約伯.
又可以隨時站在約書婭的旁邊.
但是最終我們見到的結果.
永遠都是被神來責備.
永遠都不可以理直氣壯的.

$^{1081}$站在神的面前控告.
祂要控告的人.
因為後面就是說那個原因了.
也都為什麼.
其實就不是約書婭.
好像有刀槍不入那樣.
而是大家都見到.
這個當然是靈界了.
神見到.
天使見到.
甚至現在連到撒加尼亞先知都見到.
約書婭自己可能見不見到呢.
原來她自己穿著污糟的衣服.
在聖殿.
在神的使者.
在撒旦的面前.
這個真的我們要.
所以聖餐.
不知道你們有沒有留意的了.
真的未到聖餐之前.
上到台都要再蜜土.
再尋求主耶穌的寶血.
洗乾淨自己.
然後去派餐.
因為代表著神.
和弟兄姊妹一起領受.
神再一次的恩典.
不是好像天主教說的離撒.
但是再一次去領受十誡的恩情的恩典.
所以自己都是要再一次結證.
這裡告訴我們很真實的.
所以先知真的見到一些人見不到的東西.
可能約書婭自己都懵然不知的.
所以每一個事奉者自己做完事.
回家臨睡之前.
都是將自己今天做過的事.
謙卑的交托.
神已經盡力做好了.
有什麼不行的求你赦免.
不是很在那裡邀功.

$^{1121}$哇今天多厲害.
那個報道會.
看著那個人舉手走出來.
哇多開心.
開心固然是應該的.
但是絕對不是我們人.
就是事奉的人的功勞.
其實都是不斷在恩典裡面可以服侍.
這裡就已經講了一個事奉者.
很重要的態度.
和他內心的真正的光景.
約書婭原來是穿著污穢的衣服.
站在事奉者的面前.
事奉者吩咐站在面前的.
就說了.
你們去脫去他污穢的衣服.
就是說原來在他周圍.
都有其他的事奉者.
就是說這個耶和的事奉者就再高一點.
高級一點.
很特別的.
所以永遠我們又見到.
似乎有主耶穌的影子在那裡當中.
有一個特別高級一點的.
比如第一張那個紅馬.
又是吩咐騎馬周圍去走.
周圍在那些窪地四處游走.
接著是這個紅馬向神祈禱.
你什麼時候再眷顧耶路撒冷.
要到什麼時候.
還沒眷顧.
這樣求禱的祈禱.
現在這裡就是吩咐身邊的天使.
就幫他脫了這個污穢的衣服.
象徵.
舊約永遠都是象徵.
但是已經是足夠的那個含義.
又對約書亞說.
我使道你脫離罪孽.
要給你穿上華美的衣服.

$^{1161}$再一次穿上公義的衣袍.
得到那個潔淨.
44章的《異賽亞書》22節.
大家應該要識背.
這些都是很重要的經文.
有歌仔唱的.
有些短歌的.
我塗抹了你的過犯.
好像白雲緬默.
我都忘記了.
44章的22節.
這個是《異賽亞書》.
我塗抹了你的過犯.
像口雲消散.
我塗抹了你的罪惡.
如薄雲緬默.
你當歸向我.
因我救贖了你.
那些兒校的兒歌.
可以唱的.
就是這樣.
以色列再被神紀念.
神令到他們的罪孽.
好像口雲消散.
好像薄雲緬默.
你看看那些雲.
一直飄開的時候.
那種美麗的景象.
神看我們的罪孽.
當然了.
歸向神.
就在耶穌的救恩裡面.
得到這個塗抹.
塗抹就是不再紀念.
是沒有你的記錄.
我們如果犯過事.
良民證都拿不到.
但是神的良民證.
是不會再有記錄的.
在主耶穌的寶血裡面.

$^{1201}$這個大家要記住.
所以原來救約.
也有救恩.
幫他除了污穢的衣服.
他的罪孽.
得到那個潔淨.
再沒有了.
但是記住.
除了那件污穢的衣服.
你要穿回潔淨的華美的衣服.
所以耶穌和那些人說.
我已經洗乾淨你的罪了.
從此以後不要再犯罪.
在約翰福音第八章.
和那個婦人說.
你可以想象到.
原來舊的除去.
是要穿上新人.
新約也是說很多的.
保祿.
要穿上的.
如果不是.
我們人的心靈.
如果沒有神.
沒有神的道.
就慢慢就會荒廢.
變成一個荒蕪的地方.
又會雜草叢生.
所以我們是要經常的.
來到去.
透過神的道.
和彼此的血緬.
去護衛著我們的心靈.
這個動作就是.
天使將乾淨華美的衣服.
比約書亞.
一方面是.
他要去舊迎新.
更加重要就是.
他現在是真正的.

$^{1241}$匹配的去事奉神.
或者說去服侍人.
大祭司就是要在聖殿裡面服侍.
其實很孤單的大祭司.
是他一個人一年.
進去聖殿一次.
其餘時間.
他就每天.
可能他們會輪班.
因為多人.
就進聖所裡面.
去做那些憲祭等等.
其餘的時間.
除非他下班或者是他休息.
如果不是.
經常都只能夠停留在聖殿的範圍.
所以.
某程度真的是孤單的服侍者.
但是.
他有的就是神的陪伴.
有神的同在.
所以這裡告訴我們.
他除去了污穢的衣服.
就要穿上這個新人.
這個新人.
是照著神的形象做的.
有真理的人而和聖潔.
《爾弗所書》的四章.
應該是三十節.
要穿上新人.
舊約都已經.
這麼形象化.
他.
那件污穢的衣服.
可能是大祭司的袍.
可能在事奉裡面.
不知不覺有驕傲.
有時會有一種怨天尤人.
有時會顧影自憐.
事奉就會有這些心態.

$^{1281}$所以有時你會見到一些鼎姐妹.
做著做著就不想再做.
因為.
吃力不討好.
但是.
記住.
那個是你自己的.
可能用你自己舊有的價值觀.
去看事奉.
不是人.
你就要很緊張別人.
見不到.
甚至你事奉.
他都不太在意.
你就真的很辛苦.
我做了這麼多事.
這麼真心做.
不單止沒有人贊賞.
還給人很多批評.
那下就是一個顧影自憐.
很容易不知不覺.
就思前想後.
為什麼我這麼辛苦.
為什麼我要這樣做.
為什麼我要這樣來服侍你們.
那就不做了.
其實.
不知不覺他忘記了.
其實.
是不是用著自己舊有的價值觀去服侍.
只是很想得到別人的贊賞.
不是真心為神去做.
所以他如果發現.
這個污穢的衣服.
原來是他自己本身舊有的問題.
舊有的性情.
脫開了.
穿著的是神秘的華美衣服.
那他的事奉就完全是另一個境界了.
不再介意人怎麼看他.

$^{1321}$只是介意我有沒有盡心服侍神和服侍人.
這個就是所謂的初心.
其實原因就是.
舊的衣服已經脫開了.
穿著的是神給他的.
親自給他的華美衣服.
因為他發現.
原來真正重要的是.
令到神透過我去事奉.
讓人得到神的恩典.
當他這樣理解之後.
他的事奉進入另外一個境界.
另外一個層次.
那個時候.
不知不覺人的贊賞又會來.
那時候他也不再那麼在意贊賞.
因為真正他知道.
在人心裡得到工作的是神自己.
在他的生命裡面的工作.
所以這裡是告訴我們.
很重要就是.
我們已經除去了舊人味.
所以教會如果有很多舊人在服侍.
那不怪得那麼多碰撞.
個個都帶著自己的價值觀.
自己的agenda.
一得不到自己想要得到的.
就在衝撞.
其實還一直穿著舊友的衣服.
但是一個穿著神似的.
華美的衣服.
因為我們事奉的是主基督.
保祿告訴我們.
我們大家一起都是事奉這位主基督.
大家都是要尋求上主的喜悅.
所以穿了這件衣服之後.
自然就會緊接而來的就是那個honor.
我要將潔淨的冠冕戴在他頭上.
他們就把潔淨的冠冕.
戴在若順雅的頭上.

$^{1361}$給他穿上華美的衣服.
耶和的使者就在旁邊贊賞.
就讓我們看到.
現在是進到聖殿.
進到人心.
進到大祭司的心靈.
來處理耶路撒冷的復興.
不只是那個物質的聖殿.
四圍需要有守護.
現在是最重要的那個事奉的人.
大祭司.
所以神是眷顧每一個教會.
眷顧每一個事奉的人.
因為那些人就反映著十字架救恩的榮耀.
個個都是因為神改變.
才會用心的去事奉神的.
這個是每一個教會一樣的.
你看到那些老穆者.
就是他們一生就是這樣來服侍.
他們的生命一直隨著事奉.
有不同的層次.
有不同的境界.
有不同的心境.
原因就是不斷去舊迎新.
不斷發現是恩典奉庇著你的事奉.
每人會再有一個怨言.
原因就是我們在神的愛裡面去服侍.
所以不知不覺也得到弟兄姊妹的尊敬.
和欣賞.
這就是那個冠冕.
所以你做一年就走.
做兩年就說要走.
接著這份工作你不喜歡做.
那個崗位你又覺得怎麼樣.
很難立刻有更新.
很難有你的生命去到另一個境界的那種心境.
沒有經過挫折.
何來會有一個成長呢?.
不斷應該有很多人對你的誤解.
來撐大你的心胸,你的心靈.

$^{1401}$一件事你立刻玻璃心.
那你怎麼做呢?.
做不到事奉.
因為正正是對人的工作.
當你一有一個後路要退.
你就不能夠進入另一個境界.
但實際上神給我們的就是.
有華美的衣服.
加上那個華冠.
這個是帶來.
在本來不配事奉的人.
成為了在神眼中.
可以被神重用的一個工人.
這些是一個很超級.
不是每個人都羨慕的honor.
其實真正的honor是由神那裡而來.
所以盼望我們大家都能夠.
真的參與事奉.
無論事奉是高還是低.
是看得到還是看不到.
是神安排的.
你就忠心地去做.
那個過程.
你自己都會有靈明的成長.
好 第六節了.
這個就是實實在在對若舒婭了.
她回頭以後.
她要兼顧她的弟兄了.
跟大家揭一揭.
彼得的回轉.
《路加坊》第23章.
應該是 如果我沒記錯.
第22章.
不是23章 是22章.
第31節.
和若舒婭那個都有異曲同工.
就是被主.
除去了那些污穢的衣服.
穿著華美的衣服.
再戴上那個華冠之後.

$^{1441}$第31節.
主又說.
西門 西門.
撒旦想要得著你門.
好西你門 像西麥子一樣.
就是那些坑要飛走.
用筲箕揚起那些菊.
那些坑因為很輕.
風一吹.
那些麥子就會落回筲箕.
那些坑就會吹走.
所以現在彼得.
之前就剛剛和那些門徒爭大.
他這個人在這個時候.
他是很穩重的一個大人.
但是心靈在神.
在耶穌的眼中.
就輕飄飄的 哼哼的.
就會吹.
耶穌就告誡他.
撒旦要西你們.
好像西麥子一樣.
不過第32節就很特別了.
但我已經為你祈禱了.
叫你不至於失去信心.
當你回頭之後.
就要堅固你的弟兄.
哇 對號入座.
耶穌對號入座彼得.
他現在是那些坑坯.
輕飄飄.
就是飄飄然.
耶穌說我已經為你祈禱了.
就是說到我受害之後.
我為你祈禱.
你就能夠可以回頭.
不會失去信心.
完全的失落.
哇 你想想彼得怎麼可能接受到.
這樣的對號入座.

$^{1481}$我是差到要否認你.
我是差到沒信心.
我是差到失腳.
我是差到要來到去.
好像你所描述的.
我幾年的跟隨你.
好像白費了.
正正就是這樣.
正正就是不知道什麼原因.
他可以在一個.
妹仔的面前.
三次否認耶穌和他的關係.
耶穌就說我已經為你祈禱了.
叫你不至於失去信心.
但是感恩的地方就是.
你回頭之後.
原來耶穌的祈禱.
是可以輓回彼得.
沒有為猶大祈禱.
沒有和猶大說.
你現在比彼得更加危險.
我已經為你祈禱了.
但是約翰福音第六章告訴我們.
耶穌說那個滅亡之子.
這樣的描述猶大.
所以好像他是不可能悔改的.
當然耶穌的循循善誘是一視同仁.
每個十二個門徒.
他都是這樣說.
甚至給更加多暗示.
更加多機會猶大.
不過彼得就很不同了.
他原來得到耶穌的祈禱.
他被建立.
他有一個將來的責任就是.
你回頭之後.
就是將來你發覺毒世.
否認耶穌的.
其實你應該是要面對.
那個褻毒性靈的.

$^{1521}$但是你當下的那種失腳.
其實是被神輓回了.
你回頭之後.
是要兼顧你的弟兄.
就是說有後來者.
好像你一樣的.
你就有了失敗的經驗.
被神輓回的經驗.
你就要幫助那些有需要的人.
你不要奚落.
你不要去看死別人.
因為你比別人更加死.
你跟過耶穌三年.
你就號稱自己是最愛耶穌的門徒.
但是你比其他的門徒.
一樣這麼軟弱.
所以通常回頭之後的弟兄姊妹.
對別人就不會有一種定罪的審判.
我年輕的時候初出來傳道.
有時候有些人都看不過眼.
我真的有一種這樣的心態.
雖然不敢說出口.
但是有一種這樣的心態.
什麼信耶穌信到這樣的.
幾十年過去了.
真的幾十年了.
由我信耶穌到今天是三十六年.
由我正式全時間事奉到今天是二十七年.
真的幾十年過去了.
很明白這句話.
回頭之後你要兼顧你的弟兄.
所以永遠那些老穆者.
我都算是的了.
因為登錄了.
那些老穆者講道都是充滿了恩典的.
不會再很重的一些.
好像很重的罵.
因為深深的發現.
唯有恩典去改變人.
所以回頭之後是有事要做的.

$^{1561}$是要兼顧弟兄的.
是要在教會裡面愛神所愛的.
不然將你改變有什麼用呢.
改變完你之後.
你原來繼續過你自己的生活.
又是大有人在的.
我不是只講感宣.
是一般的教會都是這樣.
很多人信了主洗禮.
不再緊張自己的靈明.
不再緊張參與教會的服事.
回頭之後.
正正就是你的生命一個很大的改變.
就是要將這個改變的能力.
神怎麼改變你.
去告訴其他人知道.
和要和神一樣.
為那些不肯改變的人流淚.
所以彼得不是每次都成功的.
就是這些流淚.
又是會改變到其他人的生命.
所以回頭看這裡的時候.
這個榮耀之後.
就是要告誡他.
真正可以在神的電語出入來往.
是需要遵行神的道.
謹守神的命令.
你就可以管理我的家.
看守我的園圓.
所以真正的事奉.
都不是外表上的人見到的那一層事奉.
是你內在內心.
對神的忠誠的跟隨.
就是走在真理裡面.
不妥協.
就是神要你說的話.
你就要忠心的說.
這是很重要的回轉的一個信徒.
需要不是去濫用神的恩典.
保羅第六章的羅馬書.

$^{1601}$一路上文講.
罪在哪裡顯多.
五章的廿四節.
恩典也顯多.
就如罪作王叫人死.
恩典也借著疑作王.
叫我們在主耶穌基督里得永生.
接著第六章就說.
我們既然在恩典裡面.
是不是可以犯罪呢.
斷乎不可.
我們在罪上死了的人.
怎麼可以再在罪中活呢.
豈不知我們這受洗.
歸入基督耶穌的.
是歸入他的死.
所以我們借著洗禮.
和他同埋葬.
完事叫我們一舉一動.
有新生命的樣式.
好像基督從死裡面復活一樣.
這是羅馬書第六章一至三節.
一樣是一樣.
就是回頭之後.
你要堅固弟兄.
回頭之後藥宣說.
你要謹守我的道.
遵行我的命令.
就是說要讓.
為什麼耶路撒冷會失落呢.
因為就是百姓不行神的道.
所以大祭司現在就要行神的道.
讓國民回到.
所謂back to the bible.
所以有時好像都很誇張.
其實真的.
一個國家回到聖經.
就自然國泰民安.
那個國家尊重聖經.
國民實踐聖經.

$^{1641}$所以大家不要再說美國是.
用基督教立國.
是.
是那些創辦人那些founder.
是很好的基督徒.
但從來都不是用聖經立國.
所以我們有時這些誤解.
但是你會見到就是.
那個大原則的方向就是這樣.
藥宣他真的需要這樣來到去.
讓國民有一個普及的聖經教育.
他自己要督信聖經.
實踐聖經.
他自己要去尊崇神.
自然他就可以在百姓當中來往行走.
這個所謂的就是.
成為了國家的很重要的.
不是精神支柱.
是那個屬靈的領導.
所以屬靈的領導是很重要的.
以致真的讓百姓能夠看到.
神真的在我們中間.
接著就是第八節.
開始出現這些預言.
講著講著其實可以收的了.
不需要了.
再給我們看到就是.
不只是約書亞.
大衛的描繪.
才是真正將來令到整個國家.
得到那個建立.
所以都不是維繫在人.
人始終是有他的限制.
大祭司就算多好的大祭司.
他都會過去的.
他都會死的.
所以希伯來書就在這個概念裡面.
就講很多的.
永遠的大祭司耶穌.
是沒有死更替的.

$^{1681}$他就可以永遠的活著.
而且他現在就在天上.
為每一個人去祈禱.
所以這裡告訴我們知道.
真正的釋放和拯救.
是仍然在將來.
就是這個大衛的描繪.
我必使我僕人大衛的描繪.
發出看.
我在若舒雅面前所立的石頭.
在一塊石頭上面有七眼.
又是那個形象.
就是說沒有一個方位.
他是看不到的.
啟示錄就更加多了.
以西潔書都有了.
那個四活物.
都已經有這些形象出現.
就是說神是見到每一個的環境.
每一個人的內心.
需要有這個聖靈.
需要有主耶穌這一位的描繪.
這樣來到去立著了.
是神親自雕刻的.
也和大家揭開了.
這個《以蔡雅書》的28章.
這些是重要經文來的.
講到來蔡雅的.
彼得都印的.
《以蔡雅書》的28章.
第14節開始.
所以你們這些涉盲的人.
就是核管住在耶路撒冷者百姓的.
似乎講那些領袖.
要聽耶和華的話.
28章的14節.
接著第15節.
你們曾經講.
我們與死亡立約.
與陰間結盟.

$^{1721}$敵軍如水.
障瞞經過的時候.
必不臨到我們.
就是用這麼重的.
的氣勢.
就是說我們一定沒事的.
一定OK的.
可以大步跨過的.
敵軍都不會來衝擊到我們的.
因我們以謊言為秘鎖.
在虛假以下藏身.
第16節.
一個很強烈的反差.
就是神定一個石頭.
所以那個基石.
那個穩固的根基.
訴諸於當時的領袖.
所謂黎菩薩過河都幫不到他.
但現在神就說我自己立一個根基.
在耶路撒冷.
第16節.
所以主耶和華如此說.
看啊.
我在石庵放一塊石頭.
作為根基.
是試驗過的石頭.
是穩固的根基.
寶貴的防角石.
信靠的人必不著急.
我以公平為準誠.
以公義為善堂.
冰雹必衝去謊言的秘鎖.
大水必漫過藏身之處.
你們所謂與死亡納的約必然廢掉.
與陰間所結的盟必立不住.
當然.
取笑他們.
哪裡可能與死亡立約.
就是告訴我們.
神自己親自設立他的愛子.

$^{1761}$成為這個穩固的根基.
就是那個寶貴的防角石.
我們翻到彼得前書.
正是引用這段經文的.
所以剛才我們說的神.
讓撒迦利亞知道.
神自己雕刻一個石頭.
第二章彼得前書.
第三節開始.
不要看太長.
你們如果嘗過主因的滋味.
就必定如此.
第四節開始說.
主乃活石.
是一個living stone.
是活的.
不是過去.
不是死的.
固然是被人所棄.
釘十字架.
被當時代的人嘲笑.
被自己的同胞拒絕.
說他是拿撒達人.
說他是木匠的兒子.
說他根本不可能說自己是彌賽亞.
卻是被神所揀選.
所寶貴.
你們來到主面前.
也就像活石.
原來我們有一個活潑的生命.
我們有性靈.
我們已經是神的子民.
所以我們被建造成為靈工.
一個spiritual house.
我們這群人的生命.
就是神的靈工.
是我們這群人.
所以永遠是建立基督徒的生命.
我們要發現自己是活石.
是任何一個聖殿的材料.

$^{1801}$就是我們這群人.
你們來到主面前.
也就像活石.
被建造成為靈工.
作聖潔的祭司.
借著耶穌基督奉獻.
神所閱立的靈祭.
所以我們是去奉獻靈祭給神.
我們是一個祭司.
將人帶到神的面前.
祭司是為兩類人祈禱的.
獻祭的.
犯罪的獻贖罪祭.
感恩的獻平安祭.
簡單來說就是這兩類.
所以我們作為祭司.
我們知道有些人是硬心.
我們為他祈禱.
求神去改變他.
有些人蒙恩.
為他來到感恩.
有些人的生命改變.
我們為他感恩.
所以我們是一個屬靈的祭司.
聖潔的屬靈的祭司.
獻屬靈的祭.
所以原來我們有這麼尊貴的位分.
我們又可能未必知道.
如果我們不讀比特前書.
比特自己講的.
看我把所揀損.
所寶貴的房角石.
安放在石安.
起不就是我們剛才讀的.
《異象雅書》的二十八章的十六節嗎?.
信靠的人必不至於羞愧.
所以他在你們信的人就為寶貴.
原來那個時代.
很多人就算見到耶穌.
都是不信.

$^{1841}$因為完全用自己的價值觀.
自己的想法.
舊約的尼賽亞.
預言一定是榮耀的.
怎麼會是這麼衰殘的.
這麼面色.
這麼難看的.
那個身份這麼寒酸的.
就說自己是尼賽亞呢?.
又不能夠行神跡的.
所以原來我們一定要用屬靈的眼光.
去看耶穌.
他在你們信的人就為寶貴.
在那不信的人就有話說.
在《引聖經》.
詩篇一百一十八篇.
大家很熟的.
這是耶和華所定的日子.
我們在其中要高興歡喜.
很多的分離都引用的.
其實就是這一段經文.
其實說的是.
「將人所棄的石頭.
已作了房國的頭塊石頭」.
是說那個日子.
當然在這樣的歡樂下.
在神的愛子的撮合.
又成為一個重要的日子.
沒有什麼錯的.
不過原本詩篇的作者是想說.
那個石頭.
被棄的石頭.
現在成為房國的頭塊石頭.
那個令人振奮的日子.
意思是.
大家如果回到原典.
至少原文這樣說過.
大家知道一下會好一點.
不然就變成斷章取義.
隨便合用就用來證明.

$^{1881}$這是耶和華所定的日子.
還不是婚禮.
還不是我的婚禮.
是對的不是錯的.
不過這個日子是說.
那個杖人棄的石頭.
即是那些猶太人.
處死耶穌的愛子.
在那個「各國他」.
「這是耶和華所定的日子」.
然後耶穌的復活.
我們就在其中高興歡喜.
所以現在是作了房國的頭塊石頭.
所以這些經文.
已經是在舊約的時候有預言.
神已經告訴撒迦利亞.
其實真正的改變.
是因為這塊石頭被設立.
是神親自雕刻.
所以你想想耶穌一生.
祂有一個名字.
以上阿順的53章第5或第6節.
祂的名字就叫作「憂患之子」.
Man of Sorrow.
即是祂一生三十多年.
三十多年是沒有好日子過.
充滿憂傷.
Man of Sorrow.
祂是憂患之子.
所以神雕刻的耶穌.
是要這樣的犧牲.
可以想象得到.
我們設法令到我們的下一代.
或者我們自己.
都要盡量過好日子.
這個是常俊心是沒有問題的.
但是人生我們知道了.
真的知道了.
不到你選擇的.
譬如我一生下來.

$^{1921}$就在九龍塘的摸區.
不到我選擇的.
沒得選擇的.
接著成長.
接著發現.
去到四五歲發現.
哇家裡這麼多人.
哇有十三兄弟姐妹.
沒得選擇的.
不到我選擇的.
我發現的時候.
就開始要面對那些弱肉強食了.
冰箱永遠空的.
家裡又窮.
冷藏了的je 喱.
還沒硬.
就已經被人喝光了.
自己都是那個參與者.
原來最快就是.
等那些人不在的時候.
喝了那些je 喱水.
因為又甜又好吃小朋友.
那你沒得選擇的.
耶穌有得選擇的.
他來這個世界降生.
但是我和大家聊聊天.
大家可能未必是很熟的.
《耳塞爾書》其實是53章.
所以如果你在人生.
信了主之後有一些小小的苦楚.
其實信不信耶穌都有的了.
那你根本不需要那麼介意.
耶穌都是這樣的.
第四節.
可能整章大家更加感動一點.
我們所傳的有誰會信呢.
耶和華的蚌壁向誰顯露呢.
他在耶和華面前生長如嫩芽.
一路解讀《耳塞爾書》.
表面就是說.

$^{1961}$這個受苦的僕人就是以色列這個民族.
但是解解下你會發現.
其實不是的.
是說一個人將來就要受這樣的苦.
這位肯定就是彌賽亞的了.
甚至乎新約已經不跟你再講了.
就是對號入座.
引這些經文就是耶穌的了.
應驗《耳塞爾書》這樣講的了.
你看看福音書就是這樣直接引用的.
他在耶和華面前生長如嫩芽.
杖根出於乾地.
就是那個出生.
兩歲就要走難去埃及.
因為希律王要殺兩歲以來的男孩.
就是說他都要這樣跟隨父母.
走去埃及避難.
他無皆營美容.
一生.
那時候又沒有什麼護膚品.
無皆營美容.
這樣描述他.
我們看見他的時候.
也沒有美貌使得我們去羨慕他.
他被藐視.
被人厭棄.
多受痛苦.
這個「常經憂患」.
英文原文就很清楚了.
它是「憂患之子」.
就是說他為什麼而生呢?.
是為憂患而生.
他整個人生就是憂患.
他被藐視.
好像被人掩面不看的一樣.
你有沒有試過被人這樣藐視你?.
藐你.
你想想很難受的.
他是被藐視的.
我們都會不尊重他的.

$^{2001}$承認擔當我們的憂患.
背負我們的痛苦.
我們卻以為他受責罰.
是被神擊打苦待.
哪知他為我們的過犯受害.
為我們的罪孽壓傷.
因他受的刑罰.
我們得平安.
因他受的鞭傷.
我們得醫治.
彼得就是引這一節.
在異章的二十四節.
就是說我們得到心靈的醫治.
其實是新心靈的醫治.
我們想想耶穌的鞭傷.
我們會得到心靈的醫治.
所以雖然是殘酷的凝聚.
但是邊邊這樣鞭落耶穌.
我們得到的是醫治.
我們好像羊一樣走迷路.
各人都偏行自己的路.
但是耶和華使我們眾人的罪孽.
記住是眾人.
全世界的罪人.
都歸在他身上.
你可以想象到.
一個人背負著全人類的罪.
天地都昏暗的.
在十二點鐘正午的時候.
這個就是主耶穌.
我想說的就是.
祂是為憂患而生的憂患之子.
這個就是神雕刻的石頭.
不是用最美的人生.
我不是仇富的.
每個人都有他的人生.
有他的舞台.
有他要走的跑道.
神是將來會來追究審問的.
神給你這樣的平台.

$^{2041}$給你這樣的人生.
你有沒有用好你的平台呢?.
你有沒有用好你的人生.
去榮耀神呢?.
所以如果很有錢.
很有社會地位.
但是很拒絕神的人.
將來的審判就很嚴重.
因為他很有資本.
很有人生的資源.
至少他聰明.
至少他有很多空閒的時間.
很多人根本找招不得晚.
很辛苦的基層.
你要他信耶穌.
他很難.
告一告假就沒有錢了.
有時也要體諒這樣的情況.
要有一種同理心.
但是神就雕刻這個石頭.
在一塊石頭上面有七眼.
萬君子耶和華說.
我要親自雕刻這個石頭.
並要在一天之間.
除掉這地的罪孽.
記住是因為耶穌的犧牲.
是神定義將他壓傷.
壓傷.
有沒有見過很大的千斤頂.
壓在鐵上.
一路壓到扁.
我們小時候家裡開豬油廠.
很多豬肉炸完變成豬油渣.
但是因為散燒.
有些機器就像千斤頂.
放在裡面.
變成一餅一餅.
一按鍵就壓下去.
耶穌是一個人.
是他的愛子.

$^{2081}$被罪孽壓傷.
被神的憤怒壓碎.
這就是神的雕刻.
所以如果神在我們大家的生命裡面.
他有主權的.
神在我們的生命裡面.
他賜下.
他收取.
我們要盡快學這些功課.
收取你的健康.
收取你開心的.
你自己的興趣.
收取一些甚至是你自己最愛的人.
我們要學這些功課.
以致我們要明白.
他是親自的雕刻.
他一定會用更好的來補償.
所以我們求神幫助我們.
我們是管家.
一定要盡快學.
我們是管家.
我們只不過是受託管.
去管理.
如果是這樣的時候.
我們盡力去愛.
但是到要離開的時候.
我們就不要死線濫打.
對神永遠懷恨.
這樣就不好.
我們要有這樣的主權.
神是他有他的主權.
所以耶穌基督.
倘若可行.
求你叫這個苦杯撤去.
但是你不要照我的意思.
照你的意思.
就是說你繼續壓傷我.
我是受得住的.
你可以想象到主耶穌.
所以這些不是拍戲.

$^{2121}$是真是聖經的真理.
告訴我們主耶穌.
就是這樣順服至死.
所以他能夠成為我們人類的救主.
所以他就能夠除掉.
一天之內除掉地上的罪孽.
所以當信主耶穌.
一家都得救.
所以你要考慮認耶穌為主.
心裡面信神.
叫他從死裡復活.
就必定得救.
這個就是口應心信.
所以我們要將耶穌這個救恩.
講清楚給人知道.
對不起,過了時間.
我們一起禱告.
多謝主.
我們能夠在第三章.
看到這個這麼寶貴的儀像.
是聖殿裡面大祭司的儀像.
多謝你施行拯救.
將人的罪惡除掉.
使人可以在你的面前.
得著新生命.
和我們是一個聖潔的身份來事奉.
幫助我們剛才.
剛才揭了很多很多經文.
也讓弟兄姊妹能夠慢慢消化.
對聖經.
對自己所得到的救恩.
更加牢固的掌握.
以至我們在事奉上.
充滿了一份.
甘心樂意的服事.
對神我們去效忠.
這樣感恩禱告.
奉耶穌基督的名.
阿們.
\newpage



\section{彌迦書 6:6-8}
\label{sec:ZbYxDMl2Zbg}
\textbf{2023-10-01 加恩樓《選擇站在上帝那邊》彌迦書 6\_6-8 講員:鄧泰康傳道}
\newline
\newline
連結: \href{https://youtube.com/watch?v=ZbYxDMl2Zbg}{\texttt{https://youtube.com/watch?v=ZbYxDMl2Zbg}} ~~~~ 語音日期: 2023-10-03
\newline
\newline
\hyperref[sec:t6pbod4Qek8]{\small{< < < PREV SERMON < < <}}
~
\hyperref[sec:index_chronic]{\small{[返順時目]}}
~
\hyperref[sec:index_scriptual]{\small{[返順卷目]}}
~
\hyperref[sec:EvC3MODORZI]{\small{> > > NEXT SERMON > > >}}
\newline
\newline
彌迦書 6:6-8
\newline
\begin{longtable}{cl}
\hline
\hline
章節 & 經文 (和合本修訂版)\\
\hline
6:6 & \begin{tabularx}{0.7\textwidth}{X} 「我朝見耶和華，在至高神面前跪拜，當獻上甚麼呢？難道獻一歲的牛犢為燔祭來朝見他嗎？ \end{tabularx} \\ \\ \relax
6:7 & \begin{tabularx}{0.7\textwidth}{X} 耶和華豈喜悅千千的公羊，或是萬萬的油河嗎？我豈可為自己的過犯獻我的長子，為自己的罪惡獻我所親生的嗎？」 \end{tabularx} \\ \\ \relax
6:8 & \begin{tabularx}{0.7\textwidth}{X} 世人哪，耶和華已指示你何為善。他向你所要的是甚麼呢？只要你行公義，好憐憫，存謙卑的心與你的神同行。 \end{tabularx} \\ \\
[1ex]
\hline
\hline
\end{longtable}
$^{1}$好, 電影節目早安, 早安, 夠不夠我的精神?.
今天應該是豆粉的樣子都是很累的, 但我訪問過他們, 今天是累得開心.
因為昨晚我們在見證著上帝很多的恩典, 上帝把人數加添下來的時間塞爆了這個地方.
所以我們很感恩, 待會我們其實也會有些跟大家分享的.
昨晚大家猜不猜得到來到這裡的人有多少人? 昨晚那個最熱的聚會.
超過500人, 500人是很保守的, 超過500人.
當我們仍然自己都小信, 令我們自己都覺得沒有可能的時候, 上帝就做給你看.
所以為什麼今天跟大家分享這個講題很重要?.
其實我們人生有很多的選擇, 但最重要的選擇對於我們來說是站在上帝那邊.
我經常說這個世界上太多價值觀, 太多不同的做法.
我們可以用自己人的方法, 我們可以用自己的計算, 用自己的各樣東西去成就一件事.
但我想告訴大家, 實實在在, 如果我們這樣做, 我們是最蠢的.
因為上帝才是所有事情的負責人.
我們只需要選擇站在上帝那邊.
尼加書是一個很重要的背景, 是當時的以色列民已經腐敗.
這班人選擇要走世界的路, 去做了很多不好的事情.
其實到了第六章的時候, 耶和華一開始就跟他們說, 我要跟我的子民爭辯.
我要跟他們爭辯, 我要跟他們去討論, 我要跟他們一起去談清楚這件事.
他們第一個說話就是, 究竟我做了什麼讓你們厭煩我?.
當時他們去選擇拜其他的偶像, 他們雖然都會在聖殿有敬拜, 但他們同時去拜其他的偶像.
他們去做他們自己喜歡的事情, 而拜其他的偶像帶來的結果就是他們的生活和內心有很多的腐敗.
然後跟他們說我做了什麼, 展示給他們看的是, 我所做的只有兩件事.
第一, 是將你們從維羅之家的埃及地領了出來, 是拯救.
第二, 你記不記得你們打過什麼仗的時候, 是我去拯救你們.
上帝很強調的就是, 我給大家的只是拯救, 我只是做了拯救.
有什麼難道我的拯救令你們很厭煩?.
然後這班子民是怎麼回應?.
第六第七節的時候, 他們的回應就是, 第一個想法就是.
上帝既然這樣說, 你想我們獻什麼給你? 我獻什麼給你?.
弟兄姊妹, 當你一聽這句話的時候, 其實是很笑話的.
我們今天一路去講奉獻奉獻奉獻, 是不是?.
奉獻的很重要的基礎是, 所有的東西都是上帝的.
所有的東西是上帝賜給你的.
他今天叫你所當立的十分之一, 是你應該的, 弟兄姊妹, 是你應該的.
如果反過來, 今天我們跟上帝去爭辯的時候, 哇, 上帝, 你要我拿這麼多的?.
我不夠, 你還要拿我? 你搞錯了焦點.
這班子民就是為什麼? 我想告訴你, 為什麼他們敢去說, 你要我獻什麼給你?.
因為他們在生活的上面, 他們因為這些鬼賊的緣故, 他們賺到很多的物質和錢財.
所以他們反過來的時候, 他們的價值觀就是這些物質和錢財是最重要的.
他列舉了一大堆, 其實是很可笑的事.

$^{41}$第一, 你記住, 全部都是以神聖之名而去做惡事, 是什麼事呢?.
難道你要我們去獻初生的繁祭嗎?.
就是說, 是不是我們獻給你的祭你不夠?.
你是不是要我們獻最好的祭給你?.
好笑的, 本來你根本就應該獻最好的祭, 是不是?.
但是他問上帝, 你是不是要我們獻最好的祭給你?.
你是不是要成千的羊? 你是不是要成萬的油糕?.
到最後, 他竟然搬了自己的祖先出來.
難道你要我們像亞伯拉罕那樣, 獻我的長子贖罪嗎?.
頂枝妹, 我們要小心.
當我們不是站在上帝那邊的時候, 我們站在自己的位置, 站在惡那邊的時候.
我可以跟上帝的, 就是這種攤牌式的方法.
覺得上帝, 你現在還不夠嗎?.
你還嫌我們不夠嗎?.
但是上帝, 他最關心的不是獻什麼禮物.
甚至乎不是你們認為的成就.
他們覺得他們很有成就.
你想想, 我沒有這樣的價值, 我沒有這樣的錢.
我哪裡能夠獻上千千萬萬的羊?.
我哪裡能夠獻上這麼多的東西給你?.
我不是要你們的獻物, 我不是要你們的成就.
上帝關心的是人內心和生活裡面充滿著很多的不公義.
上帝最關心的那個是你.
你作為他的子民, 他最關心的那個是你.
今天你的內心是怎樣?.
下面的經文就是上帝篤穿他們.
惡人家中不仍有非義之才和可惡的小聲竇嗎?.
我若用不公道的天秤和狼中鬼寨的法馬, 豈可算為清潔呢?.
城裡的富戶滿行強暴, 其中的居民也說謊言, 口中的舌頭是鬼扎的.
這班人完全不覺得自己有什麼問題.
因為他被世界價值觀已經同流合污了.
他們覺得這是正常的.
就像我以前跟大家說過, 我在國內做了很長時間宣教工作.
在那些環境裡最不容易的是什麼呢?.
這是事實.
就是他們要讀一間好的學校的時候, 他們有一個文化, 就可能要封利是.
當我和鄧英哲會談的時候, 我們這裡是這樣的.
我想說, 他們連弟兄姊妹也已經習慣了.
覺得這件事是很合理的, 是很正常的.
你要讀書嗎? 給錢吧.

$^{81}$最可怕的地方就是當有一天我們被世界價值觀同流到一個地步.
我們覺得完全沒有問題, 這是最恐怖的.
而倒過來你還去質疑上帝, 上帝是不是你要求太多?.
所以在這段經文裡, 上帝在那裡, 你就看看誰去派定了懲罰的形象.
然後他所說的懲罰是什麼呢?.
擊打他們, 所以他們會滿身傷痕.
第二件事, 因為罪的緣故, 變得荒涼.
鄧英哲會, 教會裡面荒涼不是因為沒有錢, 搞清楚.
教會裡面荒涼是因為我們不走上帝的路.
我們要很清楚一件事, 你越願意為上帝去做事.
上帝就不斷把資源賜給我們.
所以你就會發覺, 在這裡的時候, 他們說你們會變成怎樣?.
你要吃卻吃不飽, 你在虛弱中必顯在你們中間.
你必羅去卻不得獲救, 你所救護的我必交給刀劍.
你必殺種卻不得收割, 端橄欖卻不得油抹身.
端葡萄卻不得走殼, 因為你守暗你的惡規.
行阿蝦家一切所行的, 順從他們的計謀.
因此我必使你荒涼, 使你的居民令人癡笑.
你們也必擔當我們的羞辱.
上帝在這裡說得很清楚, 他們不斷在高舉我們有什麼物質.
我們有千千萬萬的羊, 千千萬萬的油.
但現在上帝告訴你, 我全部收回.
所以弟兄姊妹, 上帝在裡面就是剛才跟大家讀的那段很重要的經文.
你們聽清楚, 上帝要什麼?.
什麼是為之善?.
弟兄姊妹, 其實嘉因樓, 我經常說.
嘉因樓我們很需要找清楚我們的核心.
你會發覺我很感恩嘉因樓, 我們願意用很多的方法.
我們很開放, 用很多的方法.
去令到其他人能夠來到當中.
我們很願意讓一杯咖啡去吸引人.
我們昨晚在這裡做這個活動的時候, 有很多時候.
的而且確有的人因為遊戲的緣故而上來的.
因為氣氛而上來的, 是事實來的.
但我們的目的就是要讓他們上來.
但我們用這些方法不是我們裡面的變惡, 我們要搞清楚.
我們裡面的價值觀從來沒有變過.
我們在這裡不是要做一間會所.
我們不是要成為一間很成功的企業, 不是的.
我們所做的任何事都是為了傳福音.

$^{121}$我們所做的任何事都是為了讓更多人認識我們背後的十字架.
所以, 頂英姐妹上來告訴你什麼是重要的.
第一件事, 秉行公義.
在當時就有太多的惡, 太多的詭詐.
我可以告訴頂英姐妹, 我作為嘉因樓的傳道人.
這裡我要跟大家說得很清楚.
任何方法我都很接受, 任何方式我都很接受.
但唯一的一件事, 我不接受.
平時, 我相信鄧傳道給大家的感覺, 平時都算是一個很溫柔的人.
但當去到一點的時候, 我不會任由他.
如果大家內心裡有惡, 有一個不好的動機.
有一些傷害其他人的想法.
我不會容許.
如果你發覺有一天鄧傳道, 我也是這樣的人, 請你指出我, 請你罵我, 請你罵醒我.
而同樣, 當頂英姐妹發覺自己內心有惡, 請你告訴我.
教會裡我們所做的任何事, 都會有甩漏.
做任何事都不會完美, 一定會有錯失.
但最重要的是我們內心不是有惡.
我們一直在走上帝的公義.
我們所做的每一件事, 都要對得起上帝.
這個才是最重要的.
而第二件事, 好施憐憫頂英姐妹.
對其他人, 我們充滿著的是憐憫.
憐憫裡面, 包括更大的, 是欣賞.
我都說, 其實說真的, 剛才我們都很興奮.
昨晚有這麼多人能夠來, 當中我也很興奮.
但我想告訴你, 就算昨晚沒有人來, 我一樣這麼興奮.
為什麼?.
我真的很欣賞我們這班年青人, 在下班之後.
忙忙碌碌, 回來之後又要開會, 又要訂東西.
短短的兩個星期, 搞了昨天這麼大的事.
接著之後, 我們這班叔叔, 爸爸媽媽.
我們在旁邊的時候, 我們有什麼能幫忙, 就幫忙.
是不是?.
全部人, 我們都是出於一個心的時候, 只想做好這件事.
我們裡面是充滿著欣賞, 而憐憫是什麼?.
儘管他有些錯失, 你見不見到, 他已經很盡力了.
盡力去做, 沒什麼值得去批評.
但有沒有錯失? 有.
可不可以檢討? 可以.

$^{161}$下次可以做得更加好, 是不是?.
最重要是心裡面帶著對其他人的憐憫.
我真的很盼望, 嘉欣家這個地方, 是每個人回來得到醫治的地方.
不是我們只掛在口裡, 所說的愛.
而是我們在行動的裡面, 可以讓人感受到這份愛.
所以當弟兄姊妹, 如果很簡單, 我經常間中會去做咖啡師.
見到街坊上來, 說真一句.
他奉不奉意喝一杯咖啡, 我真的完全不介意.
但反而有時候他走過, 我會問他要不要喝咖啡.
其實我是帶著一份真正關心的.
我見到有些人的樣子其實是有受用的.
我見到有些家庭是帶著家裡, 可能平時他的兒女.
不是那麼方便在平時的時間, 但現在沒什麼人.
他可以帶她上來輕鬆一下, 玩一下.
曾經有些我不認識, 我相信也不是信徒.
跟我談了一個多兩個小時.
很細心地拿著一杯咖啡分享家裡的事, 我很感動.
以前是做不到的.
弟兄姊妹, 我們需要是一個好使憐憫的地方.
而最後帶著謙卑的心與神同行.
這句是一句來的.
帶著謙卑的心.
當我們每次見到上帝, 你也不會比上帝厲害.
弟兄姊妹.
你不會比上帝厲害.
我們在上帝的面前其實是差到不得了.
所以我覺得我們成功的秘訣是什麼?.
在神面前凝碎.
所以我覺得很感恩的圖畫是什麼?.
剛才我說過, 我們做的事有很多失誤.
我們做的事會很敏感.
但最厲害的地方是當我們將最好的呈現上了.
上帝去完成這件事.
所以我跟那班年青人一起祈禱和鼓勵和聊天.
我都說, 放心吧, 我們的上帝很厲害.
所以昨天趁Allen不在, 今早Allen去上班.
昨天我跟Allen說, 應該是去到大概八點半.
當時我們計算到已經超過三百人.
我第一時間當然是鼓勵那班年青人去找他們.
我跟Allen說, Allen, 你猜是多少人?.

$^{201}$他說, 三百人, 應該超過三百人.
他呆了一看, 看著我.
我想強調的眼神是對上帝的驚訝.
對上帝的肯定.
弟兄姊妹, 所以我只鼓勵弟兄姊妹.
我們每個人都需要站在上帝那邊.
選擇站在上帝那邊.
我們一定在前面, 上帝會無限的使用我們.
所以我很鼓勵, 好嗎?.
最後我還是這樣說.
最近我們跟很多NGO和一些人來聊天.
我跟他們說, 每個人來到, 未認識, 但參觀完之後.
我說這個地方你就當是迪士尼.
你有什麼夢想帶進來.
我們這裡是一個可以創造夢想.
可以去發生夢想的地方.
所以弟兄姊妹, 求神繼續使用我們.
秉行公義, 好施憐憫.
全謙卑的心與神同行.
一起同心祈禱.
最後我們要多謝上帝的恩典.
我求主, 您親自帶領我們.
我們真的很盼望, 我們不是要去到.
彌加先知當時見到的一群子民的光景.
我們真的求你親自保守帶領我們.
我們真的要知道.
我們不是要這些所有世界所看見的物質.
這些東西我們要為榮.
我們要知道我們走在上帝的道路上.
看上帝的工作.
我求主, 讓我們同樣的.
當今天我們去敬拜你的時候.
我們帶著謙卑的心.
最後如果我們心裡面.
被一些惡念觸動的時候.
我求主, 你的聖靈保護我們.
求你去卸去我們心裡面的惡念.
我們求你最重要.
賜給我們一個很平凡.
一個很單純.

$^{241}$我們可以跟隨你的心.
就將一切我們交在你面前.
將整個間流.
每一個弟兄姊妹.
我們都交在你面前.
其利奉教主耶穌,得勝名至其球.
阿們.
\newpage



\section{以弗所書 4:11-16}
\label{sec:EvC3MODORZI}
\textbf{2023-10-01 早堂《天降恩賜齊共享》以弗所書 4\_11-16 講員:陳麗蟬牧師}
\newline
\newline
連結: \href{https://youtube.com/watch?v=EvC3MODORZI}{\texttt{https://youtube.com/watch?v=EvC3MODORZI}} ~~~~ 語音日期: 2023-10-03
\newline
\newline
\hyperref[sec:ZbYxDMl2Zbg]{\small{< < < PREV SERMON < < <}}
~
\hyperref[sec:index_chronic]{\small{[返順時目]}}
~
\hyperref[sec:index_scriptual]{\small{[返順卷目]}}
~
\hyperref[sec:x5xyTkGVbQ0]{\small{> > > NEXT SERMON > > >}}
\newline
\newline
以弗所書 4:11-16
\newline
\begin{longtable}{cl}
\hline
\hline
章節 & 經文 (和合本修訂版)\\
\hline
4:11 & \begin{tabularx}{0.7\textwidth}{X} 他所賜的有使徒，有先知，有傳福音的，有牧者和教師， \end{tabularx} \\ \\ \relax
4:12 & \begin{tabularx}{0.7\textwidth}{X} 為要裝備聖徒，做事奉的工作，建立基督的身體， \end{tabularx} \\ \\ \relax
4:13 & \begin{tabularx}{0.7\textwidth}{X} 直等到我們眾人在信仰上同歸於一，認識神的兒子，得以長大成人，達到基督完全長成的身量。 \end{tabularx} \\ \\ \relax
4:14 & \begin{tabularx}{0.7\textwidth}{X} 這樣，我們不再作小孩子，中了人的詭計和欺騙的法術，被一切邪說之風搖動，飄來飄去。 \end{tabularx} \\ \\ \relax
4:15 & \begin{tabularx}{0.7\textwidth}{X} 我們反而要用愛心說誠實話，各方面向著基督長進，連於元首基督， \end{tabularx} \\ \\ \relax
4:16 & \begin{tabularx}{0.7\textwidth}{X} 靠著他全身都連接得緊湊，百節各按各職，照著各體的功用彼此相助，使身體漸漸增長，在愛中建立自己。 \end{tabularx} \\ \\
[1ex]
\hline
\hline
\end{longtable}
$^{1}$有四個人被困在一間細小的房間裡面.
四邊都是牆壁.
在房間最高最高的地方有一隻小窗.
他們在這個房間裡面.
他們很想去看看上帝是否在牆外.
因為覺得困在這個房間裡面.
他們好像甚麼都沒有經歷.
於是他們想到一個辦法.
就是一個一個疊上去.
最高那個就可以爬出去.
爬到窗外面就看到外面發生甚麼事.
看看能否找到上帝.
就想想誰站在最低那個.
等了一段時間之後.
A就說好 我站在最低.
沒多久B就說我疊第二個.
剩下C和D.
C就說我來第三個.
到最後D就最幸運.
騎上去他們三個人上高.
就站到頂.
經過一番努力之後.
他們很辛苦.
特別是下低的三個.
很辛苦.
支撐了十幾二十秒之後.
就全部倒下來了.
A,B,C都很想知道.
究竟D能否看到外面.
能否看到上帝.
但D透了一口氣之後.
我看到外面很多東西.
但我看不到上帝.
四個人就覺得很難過.
用了這麼多的力.
竟然看不到上帝.
沉靜了一段時間之後.
在最低的A就說.
在剛才的經歷裡.
我經歷到上帝.

$^{41}$說我是最下低的那個.
被你們踩得很辛苦.
頂不住的時候.
我發覺到上帝給我力量.
B也是這樣說.
他說我也是經歷到它.
因為當我下低有A.
上高又有C,D的時候.
我都覺得很辛苦.
我都撐得很辛苦.
我很想腳踢了A.
我就躺下去.
但我看到我經歷到上帝.
叫我支持到底.
到C了.
C就說.
其實我是最後一個.
我只要踢一踢高腳.
我就可以看到窗口.
但我知道踢一踢腳.
下低的那兩個很辛苦.
所以上帝告訴我.
我不可以這樣做.
我經歷到祂的提醒.
D說.
是呀.
我在窗外.
我想找上帝.
我找不到.
但我跟剛才你們.
個個很主動做A,B,C.
給機會.
我不用被人踩緊.
我經歷到上帝的愛.
看到你們這麼合作.
這樣去支持著我.
去到最高最高的位置.
我在你們的身上高.
我看到上帝.
上帝是個靈.

$^{81}$是看不到的.
人是看不到的.
看到就會死的.
不過.
從我們在這個肢體之間.
從人的生命裡面.
我們很多時候.
我們都看到上帝.
經歷到上帝.
剛才淑卿為我們讀那段經文裡面.
其實就是在說.
上帝給我們每一個人.
他們有不同的恩賜.
不同的才幹.
但是因為我們同一個身體.
同一個教會.
在當中我們怎樣彼此的事奉.
配搭的裡面.
我們都可以經歷到上帝.
首先我們知道.
這個恩賜的來源在哪裡.
在這個《已忽所書》第四章.
就是前幾節第八節是這樣說.
「他」這個「他」指的是耶穌基督.
祂升上高天的時候.
擄掠了仇敵.
將各樣的恩賜賜給人.
是說到耶穌基督.
祂在死裡復活.
之後我們看到.
從死裡復活.
去勝過了死亡的權勢.
魔鬼就成為了.
打敗了祂.
耶穌和那些門徒在地上.
高四十天之後就升天.
第五十天之後就是聖靈降臨.
其實在詩篇裡面.
這一節.
祂升高天的時候.

$^{121}$擄掠了仇敵.
這裡是詩篇六十八篇.
在古代的詩人已經預言到.
當耶穌基督升天.
差派這個聖靈來到.
這個聖靈就是一個最大最大的禮物.
給我們每一個信徒.
聖靈也都按著祂的心意.
將各樣的恩賜賜給我們.
以致我們得著了.
上帝而來不同的恩賜.
不同的能力的時候.
我們就在教會裡面.
彼此分擔不同的職份.
接著我們看看.
這個上帝給我們的恩賜.
究竟目的是在哪裡呢?.
在十二十三節.
祂這樣說.
是為要裝備聖徒.
做事奉的工作.
建立基督的身體.
直到我們眾人在信仰上同歸於一.
認識上帝的兒子.
得以長大成人.
達成基督完全長成的新娘.
上帝給我們有不同的恩賜.
剛才搭車進來的時候.
就遇到我們的同工 Dave.
說到中秋節那天很熱鬧.
應該是整場最辛苦的是他.
所有搬搬抬抬.
又爬高爬低.
那些全部他做了.
但我們這些.
就只是走走站站聊聊天.
大家不同的恩賜.
但目的都是要裝備信徒.
裝備這個字.
原來在原文裡面.

$^{161}$是一些外科手術的用詞.
原來當人裡面的骨頭移位.
有時有沒有試過.
不知道各有一個動作扭一扭.
扭傷了.
或者有些人因為意外.
滑了一滑.
骨頭就不知如何撐了出來.
每逢這些骨頭移位之後.
就需要做手術.
將它合起來.
我們中國人很厲害.
那些鐵打.
我媽媽那時在生的時候.
有一次不知怎樣被東西夾.
整條腳都腫了.
裡面的骨頭不知怎樣.
鐵打師傅一路跟他聊天.
就打一聲.
媽媽好痛呀.
就痊癒了.
但在希臘人的時候就說.
要經過手術.
然後就要慢慢固定位置.
休息.
之後就重新回到原位.
整個身體.
骨骼的地方就恢復正常.
所以這個裝備就是一個手術.
將一些錯位的東西.
去更正.
讓它去穩固之後.
就出來工作.
這個裝備信徒.
裝備信徒.
其實就是這個意思.
我們未信主之前.
我們人生就走錯了路.
沒有方向.
不認識主耶穌基督.

$^{201}$但當我們信了主之後.
耶穌基督將我們從這個罪惡.
黑暗的地方拯救出來.
當我們信了耶穌.
回到教會.
為甚麼又有浸禮班.
又有盤石班.
又有主日學.
有不同的課程呢.
就是要我們去接受聖經的教導.
將我們過往在這個世界裡面的一些思想.
習慣做事.
未信主之前通常是怎樣.
喜歡貪小便宜.
喜歡沙沙塵塵.
又喜歡有時說話加兩句粗口.
就是藉著這些聖經的教導.
將我們帶回耶穌基督裡面.
這應該要怎樣.
以至生命得著訓練之後.
我們就可以出來去事奉神.
我們大家彼此都是經過訓練裝備的時候.
我們就可以建立基督的身體.
建立教會.
所以我們大家弟兄姊妹.
我們需要各自得著耶穌基督的救恩.
得著適當的訓練.
我們就要在不同的位置上.
不同的崗位裡面.
互相的配搭.
發揮我們個別的功用.
我講道的.
我帶茶經的.
我絕對不會去爬高爬低.
去舔重的.
所以我們要按著恩賜彼此的服侍.
裝備信徒.
我們經過訓練.
建立基督的身體.
基督的身體原本的文字又是甚麼呢.

$^{241}$是一些關節.
一些筋膜.
表達出我們整個教會.
是耶穌基督的身體.
我們是血肉相連的.
有一個很密切很密切的關係.
亦都因為有筋骨.
這些關節.
任何一處出現問題.
我們都會痛得很辛苦.
但是我們就要大家一齊.
去怎樣彼此的配搭.
千萬不要讓任何一處受傷.
不要一處斷了.
否則我們就做不到事.
除了裝備信徒.
還有恩賜要我們怎樣呢.
在信仰上高同貴於日.
在舊的和合本就說到.
在真道上高.
我們要同貴於日.
我們先說在信仰這個字.
在真道這個字.
真道這個字.
就是上帝的說話.
信仰或真道都是.
在上帝的說話.
就是生命的糧食.
剛才我說到.
在基督的身體.
關節的那些.
黏黏的.
在當中.
這些血管.
去聯絡每一個身體.
然後給營養.
上帝的說話.
就在我們屬靈的生命當中.
去餵養著我們.
這個上帝的說話.

$^{281}$在詩篇說.
上帝說話.
我們腳前的燈.
是我們路上的光.
是我們人生旅途上高.
指引我們應該要怎樣走.
在這個希伯來書.
更加厲害.
上帝的說話.
就好像一把利劍.
可以刺透.
剖開人裡面的心思意念.
有時我們心裡面.
想了一些東西.
不是很正確.
不是很對的.
以為沒有人知道.
但當我們去看聖經的時候.
或者去聽道的時候.
我們發覺到上帝.
就會藉著他的說話.
來給我們提醒.
你好改變了.
在《雅各書》裡面.
更加說到.
上帝的說話.
就好像一面鏡子.
在我們的本相.
有時我們都會以為.
我們自己是很好的人.
我記得很多年前.
當我決定.
相信耶穌之前那一刻.
牧師說道.
世上所有都是罪人.
我但心裡面卻說.
我是好人.
因為周圍的同學.
都說我是.
很好的小姐.

$^{321}$我以為我自己是.
很好的人.
但很特別的那一次經歷.
上帝就讓我.
把不好的東西.
別人不知道的.
我說謊話.
有些什麼很好的筆記.
自己藏起來不讓別人看.
別人的就要.
我那些就一定不拿出來.
上帝就給了一幅一幅圖畫.
你點不好的圖畫.
拍了出來.
牧師繼續說.
我們人都是罪人.
到需要耶穌的時候.
就不斷地哭.
我不是真正的好好小姐.
這個上帝的說話.
就一面鏡子.
世上所有的罪人.
我覺得我是好好小姐.
一照就.
全部的本相都出來了.
所以.
祂說我們在這個真道上.
在這個聖經的上高.
同歸於一.
上帝就要我們藉著.
祂的說話.
我們來彼此合一.
為什麼在這兒.
《耳窟所書》裡面突然間會這樣說呢.
就是因為當時.
耳窟所教會.
他們正在面對一個.
洛斯底主義.
這個洛斯底主義和我們.
那個信仰是有一點點不同的.

$^{361}$他們都信耶穌.
但是他們覺得一連串的.
神有不知道七等還是九等的.
耶穌是最低等的.
因為覺得.
每逢神接觸到.
這個物質的時候.
都是最低等的.
所以耶穌基督是最差的那個神.
最聖潔最聖潔.
耶和華是最高最高那個的.
祂不會接觸人的.
但是我們一路的信仰.
都是說耶穌基督.
是道成了肉身.
最令我們感動的.
就是耶穌基督.
來到世間.
和我們體上生活裡面的苦與樂.
所以保羅說.
一定要我們的真理.
我們的信仰.
一定要見基於.
在當時試圖的教導.
試圖的教訓.
在信仰上高.
在真理上高.
要同歸於人.
以弗所書.
第一至第三節.
做一些理論的教導.
第四節到最後.
就是一些日常生活的實踐.
踏入第四章.
其實就是一些生活上高的實踐.
這個一節.
在裡面.
都是一個中心詞.
你看到在第三節.
裡面說.

$^{401}$「聖靈所賜合而為一的心」.
第四節.
「身體只有一個,聖靈只有一個,一主一信一洗一神」.
說到這個一字.
就是大家都是一致的.
大家是聯絡的.
彼此聯合的.
這裡就說到.
整個的身體.
是同一個的身體.
就好像我剛才說到.
基督的身體.
說到關節.
血管.
說到我們整個身體裡面.
有心肝脾肺腎.
大家每個人.
都是不同的功能.
但是怎樣連在一起呢?.
就是靠這些血管.
一直連在一起.
同樣都是.
當耶穌基督呼召我們.
去成為了.
這個人.
呼召我們加入了.
他的教會.
在同一個身體裡面.
他就成為了.
我們彼此之間的聯絡.
在當中.
每一個人.
我們有不同的背景.
有不同的喜好.
有人喜歡吃東西.
有人喜歡睡覺.
性格更加不同.
在合作的時候.
有些人就急驚風的.
很快就做事的.

$^{441}$有些人就慢吞吞的.
一起合作就會激死人的.
但是怎樣在這些.
很多元化不同的東西.
合在一起.
裡面一定是很多的張力的.
我記得那時候有一個.
看到兩位姊妹一起合作.
一個就很快.
一個就慢慢慢慢吞吞.
慢吞吞很生氣.
亂衝.
衝在垃圾桶裡.
很多的張力.
但也都一起不同的.
性格合作的時候.
很多的矛盾.
怎樣去追求.
這個合一呢.
就是當我們.
有些不同的意念.
出來的時候.
我們不斷的去指.
最壞就是他.
衝得這麼快.
誰不看他衝在垃圾桶裡.
又不是怪他慢吞吞.
誰不看他又怎樣怎樣.
但應該在這麼多不同的當中.
我們就要去為自己.
去做一個反思.
我是慢吞吞的人.
我可不可以快一點呢.
以致距離不要拉那麼遠呢.
所以就在這裡講.
怎樣在真理上.
來去同歸於一.
就是我們在同一身體裡.
我們怎樣用真理.
用上帝的說話.

$^{481}$我們可以去聯繫在一起呢.
在事奉當中.
真正的考驗.
不是事奉的成果.
是否達成得到.
我是否達到這個目標.
在中秋晚會裡.
我都覺得很成功.
但不是的.
最重要不是那個成果.
而是在過程裡.
我們有沒有.
我們一起事奉的肢體.
我們有沒有去認定.
耶穌基督.
就是我們同一個身體.
就是同一間教會.
以致在事奉的過程裡.
我能不能夠謙虛.
不自高自大.
覺得自己最了不起.
我能不能夠以謙虛.
溫柔.
忍耐.
寬容.
來跟其他的弟兄姊妹.
一起合作.
一起相處呢.
尊重其他人的不同.
就是活出了.
合一的信仰.
有位牧者說.
這個和諧.
不是說100個人.
說的東西.
都是同一句話.
這個不是合一.
這個合一是100個人.
他們每人說自己的一句話.
可能有100幾十句話.

$^{521}$但是當這麼多不同的說話.
都出來的時候.
我們可不可以.
來彼此的尊重.
一起的討論.
去想出一個方向.
然後大家一起.
向著這個目標直走.
有時我們覺得.
有些很固執的人.
就是.
我提出了這個.
每個人都和我不同.
你敢不敢和我不同.
我們來兇你.
來惡你的時候.
於是就逼大家.
來說出同一個方案.
逼出來的同一個方案.
不能夠使人.
真心的.
來各按恩賜.
完成整個的事工.
所以這個.
本來就是說.
凡是我們用上帝的說話.
如何用這個謙虛.
溫柔 忍耐 寬容.
我們活出合一.
還有.
在第十五第十六節裡面說到.
用愛心說誠實話.
很多時候我們都會說.
我也說過.
以後就要小心說了.
要用愛心說誠實話.
多面向著基督長進.
連於元首基督.
靠著他全身都連接得緊湊.
百折各按其足跡.

$^{561}$照著各體的功用彼此相助.
使身體漸漸增長.
在愛中建立自己.
保羅很喜歡.
就將身體裡面的血管.
黏緊我們.
將心肝脾肺腎黏緊.
上帝給我們的恩賜.
是要我們彼此的配搭.
不是用來彼此的炫耀.
我昨天放假就看到中信.
好像是金奇的.
說到有一個很有學問的人.
他放假了.
坐小艇出去散心.
漁夫撐船.
一個很有學問的學者.
就跟漁夫說.
你讀過多少書.
漁夫說.
我都沒有讀過書.
我甚麼都不懂.
沒多久.
很有學問的人又問漁夫.
你會不會唱歌.
漁夫說.
我都不會唱歌.
你真是一點人生樂趣都沒有.
你都享受不到人生的野景.
然後漁夫在游.
游來游去.
發現小艇的船底.
有些水一直滲進來.
漁夫問那個學者.
你會不會游水.
他說他不懂.
那你就甚麼命都沒有了.
各看因此.
那個學者最初就想著.
大人地驕傲.

$^{601}$笑人家沒有學識.
結果實際的東西他都不懂.
甚麼命都沒有.
上帝給我們的因此.
是互相的配搭.
一起在這個真道上.
追求這個成長.
保羅說.
為用愛心說誠實話.
但用愛心說誠實話.
不是說我想說甚麼就說甚麼.
因為有時我們會說了一些傷害人的說話.
然後我為你好.
所以我用愛心去說.
我覺得你真實的話.
提醒你.
用愛心說誠實話.
這個真正的意思.
誠實這個字.
其實是真理的意思.
是指到真理的意思.
再從譯出來.
它的譯法就是.
愛裡面的真理.
在愛的真理裡面.
我們去說說話.
我們在上帝的寬容裡面.
我們所說的每一句話.
我們無錯地要為對方的好處去說.
但是我們要估計一下.
對方可不可以承受得到.
我們要考慮到每一樣的情況.
如果當對方未必在這個時間承受得到的時候.
可能就不是我們說的時候.
這個用愛心來說誠實話.
是指在上帝愛的真理下.
我們衡量過對方.
衡量過當時周圍的情況.
我們來去說話.
所以這個然後跟著.

$^{641}$是這個用愛心來去說誠實話.
各方面都是向著基督來去長進.
連於耶穌基督.
所以每一樣東西都是連於耶穌基督.
我們是同一個的身體.
最後就回過頭來說.
在四章十一節的時候.
保羅就說到.
在這個以弗所教會裡面.
有不同的恩賜的人.
他們需要有不同恩賜的人.
其實保羅在不同的書信裡面.
他都指出了很多的恩賜.
有關顧的恩賜.
有說完的恩賜.
有說完各種不同的恩賜.
但他對著這個以弗所教會.
就是指出了這五種的恩賜.
使徒,先知,傳福音,牧者和教師.
使徒不單止指當時那十二位的使徒.
也說到一些將上帝的信息傳遞.
說到在教會裡面的一些去看方向.
先知舊約就說一些預言.
但到現在就說到一些在教會裡面.
用上帝的預言去適切教會.
或者現時社會的需要.
來去說這個信息.
傳福音我們看到有些人.
傳福音的人很容易相信.
有些人說來說去.
我這些人說來說去都很難相信.
不同的恩賜.
牧師和教師都是教導的恩賜.
我們看到保羅說來說去.
這裡有這麼多種的恩賜.
都是一些教導式的.
是否其他不需要呢?.
其他恩賜不需要呢?.
不是.
因為剛才我說了.

$^{681}$他們正在面對的就是異端的迷惑.
這個洛斯底主義.
所以保羅就要說.
他們這間教會裡面的人.
所有人都要在基督的真理上.
立穩根基.
所以這些全部在真理上.
教導的人就說到他們要出來.
不過當這些人出來的時候.
他們需要謙卑,溫柔,忍耐.
不能覺得自己高高在上.
來去做這個教導.
而體恤周圍民眾遇到的情況.
來去做教導.
保羅也要以弗所的信徒.
經過了很多的真理教導之後.
脫離了過往那些未信之前的.
那些的習慣各方面.
讓整個生命充滿了耶穌基督的生命.
流露出耶穌基督的聖情.
一個很秦道衛理的接班人.
他叫做喬治·懷德菲.
George Witterfield.
他在1740年的時候.
他是一個很出名的報道家.
無論在英國,歐洲,美國裡面.
他的報道是很有效的.
最出名的是甚麼呢?.
是在郊外來報道.
為甚麼要去郊外報道呢?.
因為很多人來聽他講道.
教會都放不下.
他說當時沒有廣告宣傳過.
又沒有這些Banner.
但是只要去到哪裡.
人們都湧進去聽他講道.
每次報道會之後.
有幾百人信主.
我們一次報道會.
幾個人信主都很難.

$^{721}$但上帝昔日就是聖靈的恩高領導.
有一次他收到一位朋友寫信給他.
這位朋友也是以前教會.
和他一起在祈禱會裡面的.
他就和他講了.
他說喬治·阿George.
我很欣賞你.
報道很成功.
不過我覺得在你的生命裡面.
沒有完全被聖靈充滿.
因為我從你的報道會裡面.
我沒有看到你有很多神蹟醫治.
只有人信主不夠.
喬治·阿George想了一段時間之後.
他就回覆了.
其實現在有些人去了報道會.
有些人都回來和我說一些掙扎.
為甚麼他們有神蹟醫治你們沒有呢?.
喬治·阿George就說.
他每一次去講道的時候.
看到上帝的話語的大能臨到.
臨到.
我看到在靈性上高黑暗的人.
被主的話語開了眼睛.
信耶穌.
我有看到一些靈性死亡的人.
在以前已經落到人生的谷底.
他因著上帝的話語.
信了耶穌.
生命醒了起來.
我都看到福音傳到很多生命貧窮的人.
所以上帝沒有將醫治人疾病的恩賜給過我.
但上帝就用聖經的真理.
去醫治人的靈魂的失散來給過我.
喬治·阿George就將.
將在上帝的各方面.
有神蹟,有醫治.
有講道真理的.
這些不同重點的工作.
他就這個重點的工作.

$^{761}$就是叫人的生命.
因著耶穌基督的話語.
他可以得著重生.
所以在聚會裡面.
單有醫治.
沒有上帝的話語.
這個不是屬於上帝的聚會.
單有上帝的話語.
沒有這些疾病的醫治.
同樣上帝都可以發揮.
祂的拯救人的工作.
我們一起祈禱吧.
因為聖靈按著你的心意.
將不同的恩賜給過我們.
幫助我們.
就是在得到上帝.
你給我們這麼大的禮物的時候.
我們怎樣去藉著你的話語.
裝備好我們自己.
讓我們的生命.
生命可以去更加的成長.
貼近耶穌.
生命流露出耶穌基督的氣象.
以致我們可以將你給我們的恩賜.
可以在教會運用.
叫主你的名字.
得著榮耀.
得著詛咒.
求主幫助我們.
以致我們在座的每一位弟兄姊妹.
都能夠覺按奇蹟.
建立主基督的身體.
我們多謝主.
多謝主瑞聽禱告.
奉靠主基督的聖名.
阿們.
\newpage



\section{馬太福音 9:35-10:4}
\label{sec:x5xyTkGVbQ0}
\textbf{2023-09-30 晚堂《流散的羊群》馬太福音 9\_35-10\_4 講員:褚永華牧師}
\newline
\newline
連結: \href{https://youtube.com/watch?v=x5xyTkGVbQ0}{\texttt{https://youtube.com/watch?v=x5xyTkGVbQ0}} ~~~~ 語音日期: 2023-10-04
\newline
\newline
\hyperref[sec:EvC3MODORZI]{\small{< < < PREV SERMON < < <}}
~
\hyperref[sec:index_chronic]{\small{[返順時目]}}
~
\hyperref[sec:index_scriptual]{\small{[返順卷目]}}
~
\hyperref[sec:JHYqHgNM3GQ]{\small{> > > NEXT SERMON > > >}}
\newline
\newline
馬太福音 9:35-10:4
\newline
\begin{longtable}{cl}
\hline
\hline
章節 & 經文 (和合本修訂版)\\
\hline
9:35 & \begin{tabularx}{0.7\textwidth}{X} 耶穌走遍各城各鄉，在他們的會堂裡教導人，宣講天國的福音，又醫治各樣的病症。 \end{tabularx} \\ \\ \relax
9:36 & \begin{tabularx}{0.7\textwidth}{X} 他看見一大群人，就憐憫他們；因為他們困苦無助，如同羊沒有牧人一樣。 \end{tabularx} \\ \\ \relax
9:37 & \begin{tabularx}{0.7\textwidth}{X} 於是他對門徒說：「要收的莊稼多，做工的人少。 \end{tabularx} \\ \\ \relax
9:38 & \begin{tabularx}{0.7\textwidth}{X} 所以，你們要求莊稼的主差遣做工的人出去收他的莊稼。」 \end{tabularx} \\ \\ \relax
10:1 & \begin{tabularx}{0.7\textwidth}{X} 耶穌叫了十二個門徒來，給他們權柄，能驅趕污靈和醫治各樣的疾病。 \end{tabularx} \\ \\ \relax
10:2 & \begin{tabularx}{0.7\textwidth}{X} 這十二使徒的名字如下：頭一個叫西門（又稱彼得），還有他弟弟安得烈，西庇太的兒子雅各和雅各的弟弟約翰， \end{tabularx} \\ \\ \relax
10:3 & \begin{tabularx}{0.7\textwidth}{X} 腓力和巴多羅買，多馬和稅吏馬太，亞勒腓的兒子雅各，和達太， \end{tabularx} \\ \\ \relax
10:4 & \begin{tabularx}{0.7\textwidth}{X} 激進黨的西門，還有出賣耶穌的加略人猶大。 \end{tabularx} \\ \\
[1ex]
\hline
\hline
\end{longtable}
$^{1}$各位電影節目的朋友們平安.
剛才我進來的時候.
盧醫生說我今晚很精神.
我經常覺得自己很精神.
我想我們前面坐的老人家.
每個人都覺得自己很精神.
只是覺得而已.
但是別人看到我們.
就說你今天好像不太精神.
但是我們覺得我們自己很精神.
我們剛剛聽到印尼來的牧師.
跟我們做過見證.
也知道我們今晚有些老人家.
他信了主.
今天第一次來我們的教會聚會.
所有這些事能夠發生.
就是經文裡面那句話.
耶穌說.
我們要求我們的天父.
將人猜出去.
講傳他的福音.
我們這邊是接受那個傳福音的人.
所結的果子.
在我們的生命裡面.
一定有一些人.
某一個人.
某一些人.
遵從那個吩咐.
進到我們的生命裡面.
跟我們講幾句關於耶穌基督的話.
使我們有感動的心.
讓我們可以聽到那些說話.
內心柔軟到.
可以信耶穌.
我們知道在這個世界裡面.
有很多人聽一波不信.
聽兩波不信.
聽十波不信.
聽一百波都不信.
我們今天在我們年紀老大的時候.

$^{41}$能夠有人將福音傳到我們的生命裡面來.
是上帝給我們一個很特別的恩典.
真是一個很特別的恩典.
讓我們在年紀老大的時候.
似乎我們的身體力量健康.
不如我們年輕的時候.
但是在這個年紀裡面.
仍然讓我們信得到耶穌.
信得到那些人不肯接納的.
或者拒絕的耶穌.
當聖經告訴我們.
為什麼我們要讀馬太福音第九章.
一直到第十章.
因為這兩段聖經是不可以分開的.
我們一分開的話固然有自己的信息.
但是最重要的信息就是.
主耶穌說.
我們要求天父.
派人出去收割莊稼.
然後他就選擇門道.
訓練他們.
是相輔相成的.
我們今天一起來看看.
在這兩段經文裡面.
或者整段經文裡面.
上帝會給我們什麼信息呢.
經文第九章三十五節.
耶穌走遍各城各鄉.
在會堂裡教訓人.
宣講天國的福音.
有醫治各樣的病症.
他出來傳講福音.
他出來醫治各樣的病症.
醫病趕鬼.
耶穌說.
帶著一班人.
可能很多人跟著他.
可能不是太多人跟著他.
經文裡面說.
耶穌走遍.

$^{81}$然後他吩咐呼召門道.
讓我們讀經的人.
在這裡想像的話.
到底耶穌在這裡.
走遍各城各鄉.
是不是孤零零的一個人.
去到不同的鄉鎮.
去到不同的村莊裡面.
進入猶太人的會堂裡面.
宣講天國的福音.
或者他們沒有辦法去到會堂裡面.
在風災露宿野地裡面.
有人群聚集的地方.
他就過去跟人說話.
我們不是很知道.
當時跟耶穌.
周圍走的有多少人.
因為第十章他才呼召門道.
第十章他才呼召門道.
有一班人跟隨他.
可能還沒有呼召他們的時候.
他們都在群眾裡面.
跟隨耶穌.
都未到.
所以第九章第三十五節.
耶穌走遍各城各鄉.
你可以想像.
他可能是孤零零的一個人.
有可能是一班人.
無論是孤零零的一個人.
無論是一班人跟隨他.
他忠心地想做完.
上帝的交代給他的工作.
他走遍各城各鄉.
有很多人就說.
我們基督徒.
或者我們中文人.
我們中國人.
喜歡把事情煲大.
其實走了一條圈.

$^{121}$他就走遍各城各鄉.
可能我們讀中文的人.
可能很多時候就說.
說大了一點.
這麼多條鄉村.
他走得完嗎?.
但是從經文.
馬太,馬可,盧迦,約翰.
四本福音書.
和其他的經文裡面.
讓我們看到.
一班人.
奉了上帝.
自己的使命.
心裡面有所感動.
無論藉著他們的旅行.
無論藉著他們.
由南到北去耶路撒冷過節.
或者他們要出去外.
做生意,謀生.
無論用什麼環境.
他們總會將.
他們所聽到的.
關於耶穌基督的事.
關於福音的事.
關於耶穌基督怎樣出生.
關於施洗約翰怎樣出生.
這些都是當時他們耳熟能詳的故事.
我們相信.
在耶穌基督成長的過程裡面.
他的父母不斷將他怎樣出生的事告訴他.
然後將施洗約翰.
他的堂兄弟.
怎樣出世的事.
告訴他.
兩件事都像神蹟一般的事.
很難發生.
經文的記載就告訴我們.
一個就是.
童貞女瑪利亞生子.

$^{161}$超乎我們想像之外.
當瑪利亞.
當約瑟將這個故事.
在他們的鄉下.
不斷流傳的時候.
你猜那些人會怎樣想呢.
你還說.
我們不找石頭打死你.
已經算你夠運了.
這是行人的婦人.
童貞女.
未曾嫁人.
怎可以和一個男人生子.
但是聖經就告訴我們.
這是上帝的工作.
是聖靈的工作.
但是瑪利亞.
她的兒子.
是聖靈的.
但是瑪利亞我相信.
在耶穌長大的過程中間.
在她生活的境況裡面.
這個故事一直不離開她的口.
不能夠離開她自己.
關於這位兒子.
這位爺子.
因為一出世的時候.
天使就跟她說.
你要給他一個名字.
不能隨便給他起名.
不能回家問爸爸.
這個名字應該怎樣.
現在的子女.
很多時候就把名字起了.
你選擇過了.
不要回家問你.
你給他什麼名字.
但是這個耶穌一出世的時候.
天使就跟她說.
你要給他起名.

$^{201}$叫什麼.
耶穌.
沒有別的名字.
只有這個名字.
因為他要把自己的百姓.
從罪惡裡救出來.
在她成長的過程裡.
她就不斷說這個兒子.
這個上帝預定.
要把她自己的百姓.
從罪惡裡拯救出來的.
這個兒子.
這個耶穌.
我相信她的故事是不斷說的.
不斷說的.
所以我相信一定有很多人.
住在納沙勒.
這個小小的鄉村裡的人.
已經耳熟能詳.
這些故事.
如果我們再詳細一點.
說施洗約翰.
他怎麼出生.
又是奇蹟般的事.
如果你問薩加利亞.
如果你問以利沙白.
喂 你們年紀長老大.
好像.
已經不能生了.
但是上帝竟然讓你生了一個兒子.
到底是一回什麼事呢.
在施洗約翰.
長大的過程裡.
他如何出生.
神蹟般的故事.
又不斷地說.
就說.
這樣的故事.
流傳到人的心裡.
不能夠抹殺.

$^{241}$因為歸根結底.
這是上帝的故事.
這是上帝自己的工作.
不是人的工作.
是上帝的工作.
基督徒的工作.
在你的生命裡.
你一定要尋找.
上帝在我們生命裡工作的痕跡.
軌跡.
我鼓勵大家.
你自己真的細心想一想.
你坐下來的時間.
多了的話 你想一想.
想一想福音在你生命裡的痕跡.
想一想上帝在你生命裡的工作.
想一想上帝.
在你身邊的人.
一些基督徒.
在你生命裡的工作.
如何將福音.
帶到你身邊.
你又肯相信.
直到現在.
你去想一想.
你不會很難.
就會看到原來.
上帝的手不斷在這裡工作.
一步步來到工作.
一步步帶領.
一步步做一些難了.
一步步做一些開門的工作.
讓我們今天仍然能夠走在.
上帝的身邊.
你去想一想.
瘟疫擴張.
疫症猖狂的時候.
我們夢想著.
上帝保守.
或者是.

$^{281}$現在流感.
都非常厲害的時候.
我們又夢想著.
上帝保守.
我們現在都習慣了.
不戴口罩.
但是我們知道.
上帝在我們的生命裡.
有一些工作.
祂真的保守我們.
讓我們無論去到哪個地方.
遇到任何人.
都能夠述說上帝在我生命裡的故事.
我們每講述一件事.
後面就一定有一個故事.
一定有一個故事.
你又將這個故事.
如何和上帝的故事.
串連起來.
我相信我們今天所有的基督徒.
在生命裡你慢慢去想.
慢慢去禱告.
尋找你生命裡的軌跡.
你不能夠忘記一些事.
不能夠忘記一些人.
將這些人 將這些事串連在一起的時候.
你就看到.
上帝的故事.
將這個故事和你生命裡的故事.
合併起來的時候.
你就更加看到.
上帝在我們生命裡的工作.
或者我將我自己的故事講.
我媽媽生了12個小孩.
但是那11個.
很小的時候就死了.
有些是六七歲死.
有些一出生沒多久就死了.
只有這個人養得大.
在我從小到成長的過程中間.

$^{321}$我很少去問媽媽.
覺得她告訴你.
到了神時.
某時.
中秋節.
端午節.
新年的時候.
我們家裡會吃媽媽自己做的湯圓.
因為我是揚州人.
吃鹹湯圓.
鹹湯圓的拳頭這麼大.
第一次吃的時候.
湯圓裡面當然有紅豆沙 芝麻.
為什麼突然間.
弄得好像菜肉包這麼大.
一到這些時候.
她就說.
如果我還在的話 今年60歲了.
我們一面快樂地.
吃著湯圓.
她就一面很痛苦地想起.
60年前的事.
當我這樣想的時候.
我就說 上帝啊.
你將我12個小孩.
只有一個人.
流傳於世.
到底你有什麼故事呢.
你的什麼故事呢.
上帝將恩典.
給了一個12個小孩.
給了這個小孩.
這個小孩到底有什麼特別的地方呢.
成長的過程裡.
不斷在問.
不斷在尋找.
不斷說 上帝啊.
為什麼會這樣呢.
我們每逢快樂的時候.
就是媽媽最痛苦的時候.

$^{361}$總是會說.
你的大家姐.
我只知道到那時那刻.
我的大家姐就60歲.
我那時候還是.
丁氏那樣大.
可能六七歲.
只知道.
原來她第一個出生的是一個女孩.
就這麼多了.
其他的知不知道呢 不知道.
只是名字都不知道.
男女都不知道.
只知道這些故事的點點滴滴.
將這些點滴.
湊在一起的時候.
就發覺到原來上帝有他自己的工作.
能夠在神樂院裡.
服侍三十多年.
能夠退休 上帝的工作.
能夠去到外國讀書.
對一個窮家子弟來說.
是想像不到的東西.
一切.
都是上帝的故事.
是上帝的恩典.
就是真正有一班人.
聽從上帝的吩咐.
懇求上帝.
差派人出去收莊價.
差派人出去傳揚福音.
有印尼一班牧師.
有牧養的.
一千多二千人.
有一班未來.
可能中間還有更偉大的人物.
在那些小孩子裡.
在那些中學的學生裡出現.
各位.
寓言這兩段經文.

$^{401}$再一次提醒我們.
不要忽略上帝.
在我們生命裡所做的任何工作.
尋找他在我們生命裡的軌跡.
無論你去到世界.
任何一個地方.
你都不要忘記.
你千萬不要忘記.
你一定不能夠忘記.
上帝仍然工作.
他不是在宇宙萬物裡.
掌管天地宇宙萬物的運行.
讓星宿.
讓各個星球.
讓各個整個宇宙.
能夠運作有序.
他還在我們人間.
在我們微不足道的人間.
在我們未來的世界.
在我們微不足道的人間.
在我們細小的生命裡.
也有工作.
如果我們細心去想.
那些工作的時候.
因為有一群人不辭勞苦.
不計較收入.
不計較金錢.
不計較名譽.
不計較地位.
只是計較我能夠為上帝付出.
造就了我們今天能夠信耶穌.
耶穌出來的時候.
有人.
走遍各城各疆.
在會堂裡教訓人.
宣講天國的福音.
醫治各樣的病症.
傳福音是一件艱難的事.
講解福音永遠都是一件艱難的事.
有時候我們人間的語言.

$^{441}$不足夠來解釋.
來演繹上帝話語的偉大.
我們要去學習.
來演繹上帝話語的偉大.
有時候人心過硬.
拒絕上帝自己的話語.
不想聽.
世間充滿了這樣的人.
但我們今天所領受的責任.
我們今天所領受的使命.
上帝在我們生命裡的痕跡.
就是你不要計較那麼多.
你的責任就是.
忠心去傳揚.
你去傳.
你的責任就是傳.
你去講.
無論有人聽不聽.
你的責任就是傳.
去講.
我們不知道什麼時候.
我們所講的東西.
能夠落到人的生命裡.
產生一些漣漪.
讓他能夠產生一些感動.
讓他能夠相信耶穌.
我們不知道.
但我們只能夠忠心地.
做上帝教導給我們的工作.
就像以色列第六章第八節那樣.
耶穌說:誰因為我問去.
我可以差遣誰.
那段經文我們都很熟悉.
以色列我再請差遣我.
上帝就跟他說.
喂!你要講.
講多久都沒有人聽.
但以色列不是說.
那你叫我去做什麼.
以色列問了一個很神奇的問題.

$^{481}$沒有人聽.
要講到什麼時候.
上帝說.
要講到什麼時候.
我現在就揭露先機.
天機告訴你.
直到他們家破人亡.
他們就會聽了.
直到他們生命裡遇到危機的時候.
他們就會聽了.
直到他們被擄外邦.
受外邦人絕大的苦難.
痛苦的時候.
他們就會想我對他們的好.
直到那天.
你去講.
最近看到一篇文字.
我相信那篇文字.
對我們今天所有的牧師.
傳福音的人都有很大的提醒.
或者幫助.
在文章裡面.
有一位外國的牧師這樣問.
他說在現代的世界裡面.
他稱呼這個叫後基督教的時代.
什麼叫後基督教呢.
基督教已經過去了.
他說現在在教會裡面講道理.
因為太多人不是很讀聖經.
因為有些人信主信得很遲.
又沒有上過主教學.
到了信主的年代裡面.
主教學又不是那麼流行.
到了我們這個年代裡面.
每個人要把主教學的元素.
在家庭小組的聚會裡面講.
小組聚會裡面講一些不講一些.
不是刻意講一些.
是不夠時間講那麼多.
所以講一段又停一段.

$^{521}$講一段又補充一段.
所以對整個聖經的理解是很少的.
傳統的理解是不少的.
知不知道俄巴底亞是誰.
什麼俄巴底亞.
知不知道猶大書在哪.
斷古猶大都在舊約.
誰知就在新約.
很多人是不明白的.
一點都不明白.
這個牧師就說.
我們去到一個這樣的年代裡面.
所以很多牧師都不明白.
在這樣的年代裡面.
所以很多牧師講到的時候.
不能夠假設你的會眾.
是明白你講的話的.
有時候覺得.
你說約拿斯之後那是什麼水.
你給大魚吃了那個.
吞出來那個.
大鯨魚.
我聽過.
但是他怎樣和上帝對話.
和上帝說.
我死了我死了.
我死了他更好.
沒聽過.
我們中亞可能.
我希望沒有這樣的人.
但是這個牧師就說.
他發覺到因為他講完之後.
很多人.
就是.
沒什麼心機聽廣道.
我們自己讀神學.
幾十年前已經說.
有些人覺得廣道是沒用的.
是沒用的.
為什麼沒用呢.

$^{561}$沒人肯坐在那裡三十個小時.
三十分鐘聽你講話.
你一站出來講完.
他已經.
好像關了按鈕.
就關了按鈕.
就是說我現在可以睡覺了.
我現在可以休息了.
我們禮拜堂裡面.
沒有一條柱子.
那些柱子後面的位置是黃金的.
每個人都拿著那些位置坐.
因為睡覺的時候沒人知道.
所以那牧師就說.
因應這樣的情況.
他就設計了幾個問題.
什麼問題呢.
他說第一個問題就是說.
我這個禮拜講完道.
下個禮拜回來.
抓一些會友來問.
會友A你過來過來.
記不記得牧師上個禮拜講了什麼.
這個問題分開三條.
牧師上個禮拜.
講了哪一段經文.
講了什麼道理.
上個禮拜裡面.
有沒有想那段經文.
翻查一下聖經.
看看這段經文對你的生命.
有什麼影響.
有什麼作用.
如果這三條都答中.
有啊.
就給自己十分.
第二條.
如果你問他.
過來過來.
記不記得上個禮拜牧師講了什麼.

$^{601}$哇.
什麼都不記得.
只記得.
你講了哪一段聖經.
給自己五分.
問第三條.
過來過來.
記不記得牧師上個禮拜講了什麼.
不記得.
不過記得他的笑話.
倒扣十分.
當我想看完之後.
詳細想想.
其實給牧師的分.
實際上只是給你自己的分.
我不用你舉手.
上個禮拜牧師講了什麼道理.
用哪一段聖經.
在上個禮拜.
直到今天為止.
有沒有想那段聖經.
那段聖經對你有什麼幫助.
或者沒什麼幫助.
或者對你有衝擊.
或者沒衝擊.
這些關係.
有沒有呢.
你自己回答自己.
當我們細心思考這些問題的時候.
我們就發覺到.
講到還有什麼用.
反正你都不記得.
我問過你自己.
這麼晚回來幹什麼.
不如在家裡看電視.
弄點水果.
喝杯咖啡.
在這裡享受一下.
反正回來聽又不記得.
整個禮拜裡面.

$^{641}$聽了等於沒聽.
聽了的道理和自己的生命毫無關係.
這樣才有出現.
這樣才神奇.
牧師全道講道.
為了什麼事.
到底為了什麼.
這些就是我們每一個信徒的人.
要在上帝的面前思考.
到底我們信仰.
所為何事.
我願意你回去思考.
你所聽的道理.
和你的生命有什麼連結.
做全港的人.
就要盡量將聖經裡面的道理.
和我們的生命.
做一些結連.
做一些提醒.
讓你生活到.
那個禮拜.
開始的時候.
你能夠仍然想起.
主的說話.
對你生命裡的結連.
當主耶穌.
周圍去講訊息.
一定有人.
歡言接納.
一定有些人水過鴨背.
一樣有些人.
擁抱的.
不肯的接納.
無論什麼人都好.
但是他看到一群什麼人.
他就看到.
原來你不接納.
原來你這麼喜歡聽.
原來你是.
水過鴨背.

$^{681}$他不是用這樣來分別人.
他看到的是什麼.
他們看見很多人.
憐憫他們.
因為他們將困苦.
流離.
沒有牧人的羊.
這些說話對我們今天做宣講.
似乎宣講不達果效.
聽完.
等於沒聽.
不知道怎樣.
是不是要繼續堅持講下去.
這些說話就提醒我們.
他要你看的.
不是說他接納.
他不接納.
他多有心聽.
他沒心聽.
不是叫我們做這樣的分別.
經文的教導就說.
我們能不能夠.
透過主耶穌基督的眼睛.
看到他們是困苦.
流離.
失散的羊.
他們就很需要.
你不斷這樣去講.
在不斷傳講的過程裡.
起碼有一句.
兩句半句.
能夠進入他們的心.
讓他們在生命裡.
遇到艱難的時候.
就是那一句半句.
在他們生命裡產生.
極大的動力.
讓他們不放棄.
我們今天.
所期盼的.

$^{721}$我們想做的就是這樣.
儘管沒有人願意聽.
儘管有人覺得.
很浪費時間.
但是所有的傳道牧者.
都不能夠興奮這個即時.
就算半句都好.
只要你記得半句都好.
我都不放棄這個機會.
就為了這個半句.
能夠進入你的心裡.
讓聖靈可以工作.
就為了這個半句.
因為.
我們透過主耶穌基督的眼睛.
所看到的.
是流離.
失散.
需要上帝的話語.
來幫助.
滋潤鼓勵的人.
所以他對門徒說.
要收的莊稼多.
做工的人少.
所以你們當舊莊稼的主.
帶法人出去.
收他的莊稼.
在這裡我們不能夠.
不想在這個世界裡.
有很多的人.
有很多很出色的人.
有很多不是那麼出色的人.
或者我應該說.
有很多很聰明的人.
有些不是那麼聰明的人.
有些很絕頂聰明的人.
響應上帝的呼召.
出去做這個工作.
上帝說在這裡請你收起我.
我願意去到最艱難的地方.

$^{761}$最沒有人願意去的地方.
最痛苦的地方.
去到那個地方.
傳揚你自己的道理.
我知道那裡艱難.
但我不計較.
我只知道你叫我去就行了.
只要你清晰地發師號令.
告訴我.
我要你去.
就去那個地方.
我就會去.
因為我看到的就是那.
沒有牧人的羊.
我聽到的那是一個猜派我的上帝.
手上有很多的見證.
我隨處拿來.
都有一些見證.
讓我們看到上帝自己的工作.
不要輕忽在家庭裡.
家庭主婦的責任.
我們中間都有一些家庭主婦.
一想到家庭主婦.
我就想到奧古斯丁的母親.
不知為何.
我都想到我母親.
我母親不識字.
一隻字都不識.
連自己的名字都不識寫.
簽紙要吸手指毛.
吸手指毛.
但在我長大的過程裡.
我就經常看到她和一班姐妹.
在教會裡面也好.
去到個人家裡也好.
有幾個姐妹.
就在大清則庫開聲.
跪在那裡祈禱.
那時都五六十歲.
她八十多歲離開世界.

$^{801}$那時都五六十歲.
難道那時沒有.
骨質疏鬆.
沒有關節痛.
看到她跪在地上半小時.
還跪在那裡祈禱.
那時不懂得想這些.
那時只說.
哇 他們祈禱這麼大聲.
好像呼天搶地.
又不識字.
未讀過書.
在揚州鄉下養鴨.
養鴨人家.
只懂得祈禱.
然後在她安息禮拜的時候.
有一個和她一起.
在那裡祈禱的姐妹說.
在她上來講一些見證的時候.
她就說.
儒牧師 你知不知道你媽媽.
整天都為你祈禱.
她不想你做律師.
不想你做會計師.
不想你做醫師.
只想你做牧師.
我說我媽媽沒有告訴我.
那個姐妹就說.
她祈禱有效.
一個養鴨人家的.
沒什麼知識的.
一個女孩.
嫁了來給揚州的.
在銀行裡面做學師.
實際上做學徒.
什麼都做的.
學師.
留下一個小孩.
不斷為她祈禱.
希望她有朝一日.

$^{841}$成為一個全道的人.
最近看到在奧古斯丁所寫的.
懺悔錄裡面.
第九卷Book 9.
有一段講她媽媽的說話.
你可能不知道奧古斯丁的媽媽叫什麼名字.
莫妮卡.
不過你一定會記得.
張國明唱的那首歌叫什麼.
那首歌的名字叫什麼.
叫莫妮卡.
可能.
如果你說我知道那首歌.
就是我們同年代.
莫妮卡.
我們看奧古斯丁的懺悔錄裡面.
所寫她媽媽的事跡.
好像一輩子.
她自做了兩件事.
嚴格講起來.
只是做了一件事.
因為莫妮卡.
是一個很虔誠的基督徒.
但是她嫁了給一個異教徒.
她老公不信耶穌.
不知道信什麼宗教.
奇奇怪怪的.
第三第四世紀的宗教.
而且家長脾氣甚壞.
被脾氣.
不得了.
但是莫妮卡.
總是她發完脾氣之後.
人都走了.
她就請她老公.
坐下在沒有人的地方.
就說你剛才這樣發的.
是不對的.
你說.
這樣的生活過了多少年呢.

$^{881}$因為奧古斯丁.
是北非人.
非洲人.
在羅馬帝國管理下的非洲.
在羅馬帝國那時候.
通常那些人.
少男少女什麼時候結婚呢.
天理炎熱.
早熟.
所以十三四歲.
十四五歲就嫁人了.
我們不要說她那麼老.
十五歲.
通常十四五歲就嫁人了.
莫妮卡.
她老公離開世界那一年.
莫妮卡剛剛四十歲.
說多說少.
她嫁給這個那麼大脾氣的男人.
起碼都超過二十五年.
在二十幾年裡面.
她就不斷地為她老公祈禱.
希望她老公有一天.
能夠信到耶穌.
終於在莫妮卡四十歲那一年.
她老公臨死之前.
信了耶穌.
生命到了最後關頭.
她信了耶穌.
莫妮卡很開心.
但是讓莫妮卡更加開心的是什麼呢.
她的兒子信了耶穌.
她老公死那一年.
莫妮卡四十歲.
到五十六歲.
那一年 莫妮卡就離開世界了.
也是那一年.
奧古斯丁信了耶穌.
奧古斯丁信了耶穌之後.
我們知道奧古斯丁在教會歷史.

$^{921}$到現在為止都是一個非常出色的神學家.
他留下的著作非常多.
用拉丁文來寫著作.
寫得非常多.
所以我們今天仍然要思考一些神學問題的話.
還是要看奧古斯丁.
但是在那本《布拉第九卷》裡面.
他記述他媽媽在生的時候.
跟他講的東西.
下面這一段是這樣說的.
他說兒子 兒子.
莫妮卡跟奧古斯丁說.
兒子 對於我來說.
我此生已毫不.
毫無留戀.
我不知道還有什麼.
我可以做.
為什麼我還要留在這個世界上.
因為我的預謀已經.
已經滿足了.
過去所以要暫留此世.
不過希望.
你在我未睡世之前.
能夠信到耶穌.
而上帝的恩典超越我本來的盼望.
使我看見.
你竟然能夠輕視人世的幸福.
成為上主的僕人.
我還有什麼可以做呢.
我覺得我的工作已經完畢了.
這是奧古斯丁在生的時候.
他媽媽在生的時候.
聽他媽媽所說的話.
不久後.
莫妮卡就離開了世界.
由這些戰篇直接裡面.
我們就可以看見.
莫妮卡一世人就是盼望.
一世人就是祈禱.
這個師奶就是祈禱.

$^{961}$祈禱到兒子能夠信到耶穌.
我們中間有很多師奶.
我希望我們都能夠.
忠心地.
腳踏實地地.
誠心地.
全心全意地.
為你的兒孫祈禱.
為你身邊的應該信到耶穌.
但是現在還沒有神信耶穌的人.
實實地為他們祈禱.
直到有一天.
你能夠說出像是莫妮卡那樣的話.
我存在世界上的目的.
就是祈禱到.
他信耶穌.
然後.
在馬達福音第十章.
耶穌建了十二個門徒.
一定有一班人跟著他.
但是現在在一班人跟著他的人群中間.
他選了十二個人.
這班人的組成.
是令我們很.
為什麼上帝要選人.
選得很出奇.
選一個賣耶穌的猶大.
經文這樣寫出來.
誰知道賣耶穌的人叫做猶大.
我們不知道猶大有沒有結婚.
他的兒子也挺難受的.
他的太太也挺難受的.
一見到猶大.
就不用說了.
賣耶穌那個.
還有就是.
瑞利馬達.
瑞利也是瑞人.
專門幫羅馬人來工作.
還有很多資料說.

$^{1001}$一些瑞利.
那個官職.
是買回來的.
價高者得 必回來的.
如果你是必回來的話.
斷股你做了那個位置.
就拼命賺錢.
賺夠本才行.
但是這樣的人.
跟另外一個人.
由粉瑞黨的西門.
第四章粉瑞黨的西門.
粉瑞黨的西門.
如果用現在的俗話來說.
我們叫做刀仔黨.
為什麼呢.
他隨身在咒口裡面.
說你們擺了些刀仔.
專門會殺什麼人.
暗殺什麼人.
瑞利馬達這些人.
專門幫羅馬人來工作.
那些猶大人.
那些漢奸.
叫做尤奸.
尤奸在後面砍他兩刀.
為什麼耶穌會.
放一個這樣的人在那裡.
還有阿肯和阿肥.
他們兩個兒子亞國約翰.
記不記得.
他們兩個人來求耶穌.
可能他們自己不好意思說.
叫媽媽來說.
媽媽就見到耶穌說.
將來你德國降臨的時候.
天國降臨的時候.
你最重要留下兩個位置.
留下兩個位置.
一個給我這個.

$^{1041}$一個給我那個亞國約翰.
讓他們一個人坐在你左一個人坐在你右.
他們兩個兒子一聲都不出.
馬上被人打.
多好多好.
一群這樣的組合.
如果我們看回第九章.
懇求莊家的主.
差派人出去.
原來差派一群.
這樣的人.
我們馬上就想到.
我們今天承擔責任的人.
承擔責任的人.
有好多過他嗎.
我們甚至.
比他還差.
但是上帝竟然將責任.
給過我們.
當然也有很多人.
條件比他們更好.
對上帝來說.
忠心耿耿.
無論在任何情況.
至死忠心絕不退後.
有這樣的人.
只要我們能夠按照上帝.
自己的心意而行.
我鼓勵大家在這段時間.
如果你有時間.
你可以讀一讀一些宣教士的記錄.
或者宣教士的生病.
看他們怎樣信耶穌.
怎樣按照上帝的心意.
去到世界不同的地方.
來做上主交代給他們的工作.
最近回到.
內地會.
戴德生所創辦的內地會.
他有些規條.

$^{1081}$所有根據去內地.
中國內地做宣教士的.
那些外國人.
全部都要穿長衫馬褂.
穿中國服飾.
吃中國的東西.
好像中國人流蘇.
那些西人說很反感.
為什麼要這樣呢.
他說不要第一眼就被人拒絕你.
你穿和他們一樣的衣服.
穿中國的服飾.
吃中國的東西.
學中國的話.
將福音容易傳給他們.
中間有很多很出色的人.
有班劍橋.
在劍橋大學畢業的.
我們稱呼叫劍橋七子.
全部都是劍橋大學畢業的.
二話不說.
劍橋大學畢業.
高頭就去中國傳福音.
有些牛津畢業的.
很多醫生.
很多傳教士.
有很厲害的.
有些不是很厲害的.
但他們同一個呼召.
我要將福音帶到中國內地.
讓那些人能夠信到耶穌.
直到現在.
我們都能看到.
這班厲害與不厲害.
出色與不出色.
全部都不是問題.
主耶穌說.
我只看到一班.
流離失所的羊.
他們需要牧羊.

$^{1121}$只要他們能夠信到耶穌.
體會到上帝在他們生命中的工作.
你就可以知道.
上帝照顧流離的羊群.
因為祂要使用我們.
使用所有的人.
願主幫助我們.
憐憫我們.
我希望.
這個星期裡.
你繼續思考.
如果在這段經文裡.
你不記得所有的事.
但希望你能記住.
一句半句聖經裡的說話.
但希望你能記得.
一句半句聖經裡的說話.
無論是要收的莊家多.
作供的人少.
肯求主.
差派人出去.
收割他的莊家.
又或者說.
我完全透過主耶穌基督的眼光.
所看到的.
不是分門別類.
不是叻與不叻.
不是聽到留心.
水過鴨背.
我不是看到這些東西.
願我透過主耶穌基督的眼睛.
看到那些流離失所.
困苦無助.
我好像羊.
沒有牧人一般的羊群一樣.
願主幫助我們.
在接下來這個星期裡.
上你的話語.
不離開你的心.
也不離開你的思想.

$^{1161}$更加不要離開你的生活.
多謝主.
用你的話語提醒我們.
我們都是流離失所.
困苦無助的羊群.
願主接著你的話語.
接著你的僕人使律.
繼續不斷的慰養我們.
讓我們的生命成長.
承接上帝的交託.
給我們的工作.
多謝主.
誰聽我們的禱告.
奉耶穌基督的聖名.
\newpage



\section{撒母耳記下 23:13-19}
\label{sec:JHYqHgNM3GQ}
\textbf{2023-10-01 午堂《大衛與三個沒有名字的勇士》撒母耳記下 23\_13-19 講員:陳韋安牧師}
\newline
\newline
連結: \href{https://youtube.com/watch?v=JHYqHgNM3GQ}{\texttt{https://youtube.com/watch?v=JHYqHgNM3GQ}} ~~~~ 語音日期: 2023-10-04
\newline
\newline
\hyperref[sec:x5xyTkGVbQ0]{\small{< < < PREV SERMON < < <}}
~
\hyperref[sec:index_chronic]{\small{[返順時目]}}
~
\hyperref[sec:index_scriptual]{\small{[返順卷目]}}
~
\hyperref[sec:CXJyuYm4_pE]{\small{> > > NEXT SERMON > > >}}
\newline
\newline
撒母耳記下 23:13-19
\newline
\begin{longtable}{cl}
\hline
\hline
章節 & 經文 (和合本修訂版)\\
\hline
23:13 & \begin{tabularx}{0.7\textwidth}{X} 開始收割的時候，三個侍衛下到亞杜蘭洞，到大衛那裡。非利士的軍兵在利乏音谷安營。 \end{tabularx} \\ \\ \relax
23:14 & \begin{tabularx}{0.7\textwidth}{X} 那時大衛在山寨，非利士人的駐軍在伯利恆。 \end{tabularx} \\ \\ \relax
23:15 & \begin{tabularx}{0.7\textwidth}{X} 大衛渴想著說：「但願有人從伯利恆城門旁的井裡打水來給我喝！」 \end{tabularx} \\ \\ \relax
23:16 & \begin{tabularx}{0.7\textwidth}{X} 這三個勇士就闖過非利士人的軍營，從伯利恆城門旁的井裡打水，拿來給大衛喝。他卻不肯喝，將水澆在耶和華面前， \end{tabularx} \\ \\ \relax
23:17 & \begin{tabularx}{0.7\textwidth}{X} 說：「耶和華啊，我絕不做這事！這三個人冒生命的危險，這不是他們的血嗎？」大衛不肯喝這水。這是三個勇士所做的事。 \end{tabularx} \\ \\ \relax
23:18 & \begin{tabularx}{0.7\textwidth}{X} 洗魯雅的兒子，約押的兄弟亞比篩是這三個勇士的領袖；他曾舉槍殺了三百人，就在三個勇士中得了名。 \end{tabularx} \\ \\ \relax
23:19 & \begin{tabularx}{0.7\textwidth}{X} 他在這三個勇士中是最有名望的，所以作他們的領袖，只是不及前三個勇士。 \end{tabularx} \\ \\
[1ex]
\hline
\hline
\end{longtable}
$^{1}$.
今天我們會說《殺母耳記》下一段大衛和其下三個勇士的故事.
基文記載在《殺母耳記》下的第23章.
23章其實是整個《殺母耳記》最後的一部分.
如果你看《殺母耳記》的話.
《殺母耳記》是一本順著的書卷.
整本書是按著歷史的年份一段一段地說下去.
從哈拿生仔開始.
《殺母耳記》上經文.
到殺母耳小朋友的事蹟.
蘇羅王被以色列人高納成為皇帝.
再說到大衛興起.
然後蘇羅和大衛之家兩家的鬥爭.
到了大衛王朝的建立.
統一了整個以色列.
再說到大衛成為皇帝之後的故事.
拔士花 阿沙隆等等.
整個書卷上下都是按著年份一段一段地說下去.
你會發現在以色列文歷史的記載.
不過來到《殺母耳記》下的最後一部分.
你會發現從21章到24章最後那幾章書卷.
其實是不跟時序去記載的.
你會發現21章到24章就好像一本書的附錄.
appendix一樣.
就是夾雜了很多不同種類的題材.
放在整本書的最後一章.
這些附錄包括什麼呢.
你會發現appendix無端端會加插一篇詩篇.
詩篇18篇.
就是大衛自己曾經親手寫的一篇詩篇.
讚美上帝的詩篇.
還加了一段.
大衛臨死前又說上帝去向他曉喻一段說話.
一段臨告生命最後一段上帝的說話.
還有一段就是24災.
一段獨立的故事.
就是大衛數民的故事.
然後上帝懲罰他的故事.
最後還有一段.
就是今天所說的經文.

$^{41}$就是一段清單.
一個清單.
就是大衛的勇士們的清單.
我們問一下為什麼這批勇士.
能夠這麼好.
這麼有份量.
能夠記載在薩姆爾記的附錄裡面呢.
原因只有一個.
就是厲害.
因為這班人.
這30個勇士.
是名副其實都是大衛王朝旗下.
一個最能打得最能抽的英雄.
如果薩姆爾記是一隻三國志遊戲的話.
基本上每個人都有數據.
你可以用他的武將.
非常厲害.
名副其實是戰狼三十.
說的正正就是戰狼三十裡面.
其中的三個.
30個裡面的三個.
經文裡面寫.
修國的時候有30個勇士.
當中的三個.
去到亞都蘭洞見大衛.
非利士的軍兵在利凡欽安營.
那時候大衛在山寨.
非利士人安營在伯利倫.
大衛又渴想說.
甚願有人將伯利倫城門旁井裡的水打給我.
這三個勇士就闖過非利士人的陣營.
那個營盤.
去到伯利倫城門旁的井裡.
就打水.
拿水去奉給大衛.
基本上總括來說.
整個故事.
用一句我們現代人的說話來說.
整個故事基本上是一個Foodpanda的故事.
是一個外賣的故事.

$^{81}$三個勇士就一起.
就千里迢迢地去接一個外賣order.
從伯利倫的井旁.
外賣一杯.
不是真奶.
而是一杯井水.
來給亞都蘭洞的大衛.
沒有貼士.
經文記載.
在大衛還沒有去統一天下的時候.
留意.
經文裡面所說的利凡欽谷.
這個地方.
如果打開《深海記》第五章.
同樣出現了利凡欽谷這個字眼.
利凡欽谷又名巨人谷.
大衛將要去到統一天下之前.
最後.
最後擊破伯利倫.
這個菲律賓人的陣營.
就能夠統一天下.
所以你發現.
整個故事是說什麼.
快要統一天下.
如日方中的大衛.
用手機下了order.
用生活法律來說.
外賣的餐廳就在伯利倫門旁.
井旁的井那裡.
外賣的地址就在阿都蘭洞.
阿都蘭洞和伯利倫相距有多遠呢.
如果你真的來的話.
可以回去google map.
得出結果.
你搜尋兩個地方的名字.
是28公里.
步行的話大概需要六個小時.
一來一回.
你去一去.
浸浸水.

$^{121}$再回來.
大概需要12小時.
假設你不休息的話.
不過ok的.
這三個是勇士.
速度是很快的.
他們不需要休息.
所以這三位南亞車手.
中東的外賣車手.
由阿都蘭洞出發.
前往伯利倫.
途中經過利凡欽谷.
這個巨人谷.
敵方的陣營.
千山萬水.
去到伯利倫門旁的井.
浸浸水.
放在保溫檔裡面.
再回頭.
再經過敵方的陣營.
穿過去.
回到阿都蘭洞.
將外賣遞給客人.
簽收.
一件非常之禮貌不起.
非常之震撼的事情.
最令人覺得驚訝的.
就是大衛當他接過外賣的時候.
隨手將這壺新鮮的冰涼水.
從保溫檔拿出來.
真真正正的千山萬水.
他怎樣做呢.
他就將整壺.
放在地上.
一下子就倒在地上.
我們所討論的不是大衛為何要這樣做.
我們不用大衛的角度去思考這個故事.
當然大衛的動機.
有很多性經學者都可以去研究.
去猜今天大衛為何要這樣做.

$^{161}$雖然今天不是重點.
但不妨都講一下.
為何大衛會將這些水.
或者說為何會打發人.
要千里迢迢.
去拿水喝呢.
有幾個可能的意見.
第一個我們首先要排除的.
就是大衛不是口寒.
大衛不是出於生理的需求.
很辛苦沒有水喝.
就要求大衛幫他打水.
三個勇士都不是因為這樣緣故.
幫他打水.
雖然你看到阿杜蘭洞這個字.
但沒有你想像中那麼淒涼.
阿杜蘭洞不是躲藏在山洞.
而是一個軍事的陣營.
最少是有水的.
有一個水源在那裡.
所以大衛不是因為口渴.
所以就打發人去拿水喝.
我們先排除這個.
第二個可能的原因是什麼呢.
可能就是大衛私鄉.
大衛打發人去伯利恆那裡打水.
因為伯利恆其實就是大衛的故鄉.
雖然你今天比較熟悉的就是伯利恆.
就是耶穌出世的地方.
但伯利恆就是叫做大衛的城.
所以其實伯利恆本來就是大衛從小長大的地方.
就好像大家的元朗一樣.
從小就長大在這個地方裡面.
所以當大衛特別要拿一張order.
要喝這個伯利恆城門旁的井水.
可能就是想品嘗一下小時候那些鄉下的井水的味道.
第三個可能的原因.
有人認為充滿了一個惡意的陰謀論.
大衛其實純粹是要宣洩自己的權力.
當時大衛如日方中將要統治天下.

$^{201}$所以他為了去展示自己的權力.
就隨便叫人冒著生命危險.
叫他的手下幫他打水喝.
就好像我們以前看色神.
喜歡怎樣就怎樣就叫做power.
第四個學者有人都認為寫了一篇論文.
關乎於這個原因.
大衛這樣做可能是出於一個戰略的考慮.
大衛知道這是關乎於第五章的對戰.
所以表面上大衛叫人去拿水喝.
可能有些人認為是戰略上的原因.
為了打發勇士去試探敵方的軍情.
這是學者的推測.
不過這也不是我們今天要講到關注的地方.
我們要關注的地方.
我們嘗試站在這三個勇士的角度去想這件事.
無論大衛的動機是怎樣也好.
一個千真萬確無可否認鐵砲的事實是什麼.
這三個勇士山長水遠將自己的性命都豁出去.
幾經辛苦就是要將水打給大衛.
大衛卻一下子就倒在地上.
糟蹋了.
不知道你會怎樣想.
如果你是那三位勇士的話.
你本著一條命跑了十二小時.
拿回來的東西一下子被大衛倒在地上.
你有什麼感受呢.
基本上根據聖經記載.
聖經是沒有記載這三位勇士的反應.
不過如果是你.
你真的付出120\%的心血.
幾經很好的去付出努力.
甚至本著一條命.
卻被人家一杯倒在地上.
不知道你會怎樣去理解這件事情.
白做.
浪費心血.
徒然.
很無謂.
如果真的是foodpanda我也覺得OK.

$^{241}$如果我送外賣的話.
我當然知道你吃不吃.
基本上我都賺了錢.
你馬上吃一口也好.
不關我的事.
基本上我都是收了錢就走人.
但當你連你自己的生命都擺上.
得出來的結果.
卻被人家糟蹋在地上.
你說說感受.
有什麼感受呢.
不過是真的.
當我們年紀越大.
你發現這個世界上有很多很多的事情.
本質可能都是類似的.
幾經辛苦.
滿流血汗.
白白的付出.
得出來的結果.
卻是沒有什麼結果的.
努力付出了.
但最後一番心血都是白做的.
一個運動員為了參加奧運會練了四年.
突然出賽前受傷.
一個建築師或是一個公司的項目.
一班同事一起花了半年的心機.
做了很多晚通宵.
做了五百多個計劃.
突然間公司說沒有資金.
項目就取消.
很多時候你會發覺.
原來世界上有很多這樣的情況.
你付出的心血和結果其實是不成正比的.
太多這樣的例子.
事實上我們不知道大家如何去理解這些事情.
暑假的時候.
我和我的教會.
一班同工就去了台灣退休.
我們經過一個夜市.
夜市有一個彈珠店.

$^{281}$那間彈珠店叫小時候彈珠店.
那間就是賣古早味道的彈珠店.
即是懷舊.
裡面有很多小朋友小時候玩的彈珠機.
他們買的是童年回憶.
我們一班人就.
基本上玩什麼呢.
基本上你唱幾百元台幣.
然後就可以得到一堆彈珠.
然後你就可以彈珠機彈.
贏了就得到更多彈珠.
然後你彈珠玩下去就得到更多彈珠.
其實就是這樣.
沒有什麼特別意思.
整件事就是這樣.
你就可以玩彈珠機.
可以玩一段時間.
我們就一班牧者就一起去玩.
不過無聊說無聊.
我欣賞這個彈珠店裡面的標語.
台灣很多很文青的口號.
我覺得很精彩.
他如何推銷他的彈珠店呢.
玩彈珠機這些東西.
有兩句很值得大家去聽一下.
這樣的口號.
他說很多事情都是這樣.
你認為值得就是值得.
就像打彈珠一樣.
你覺得值就行了.
你不要理會有什麼用.
你不要理會有什麼玩.
總之一種情懷.
你覺得值就值.
你就可以去玩.
另一句就叫.
就算很多事情沒有發生.
有沒有關係.
最重要的是我們一起打彈珠.
度過美好的一天.

$^{321}$基本上他就是推銷這件事.
其實沒有什麼貢獻.
沒有什麼特別的東西.
你玩彈珠開心就可以了.
我覺得自己在這個彈珠店裡面.
我自己有一種很特別的反省.
我覺得作為教會.
作為一班基督徒.
無論是教會也好.
打彈珠也好.
或者你小時候.
讀大學的時候的上裝也好.
基本上一班人.
大家聚在一起.
很熱血地做一些事情.
往往未必有很大的結果和改變.
不過其實OK的.
就算好像沒有什麼事情發生了.
你的經驗的經歷.
其實是OK的.
對我來說這是一個很重要的信仰反省.
那次台灣退修.
我們也有機會去探訪台灣的教會.
探訪了The Hook.
大家也聽過.
這間跟我們差不多類型.
比較媒體多一點的.
年輕人多一點的教會.
發覺那些教會.
其實也是很多這樣的方法去運作.
基本上它也是台北靈良堂.
這樣一個有關係的教會.
所以也是很多這樣的手法去做.
但我們發覺我們教會並不是這樣.
我們也花了很多心思.
放在牧羊那裡.
牧者希望能夠跟得更貼近人.
或者希望能夠更加關心弟兄姊妹的生命.
雖然做了很多這些事情.
其實未必有什麼改變.

$^{361}$你放在很多心血的弟兄姊妹那裡.
可能未必有這麼多即時看到的東西.
不過我覺得牧羊可能就是這樣.
基本上我們未必做很多即時有結果的事情.
教會不是外面那些世俗的機構.
更加不要說那些JPEGS攪碗店那些.
基本上教會是一個對社會上沒有什麼經濟很大的效益.
你來到不會有身家多了.
對經濟貢獻也不會多了.
教會如果相反地只是做一些有效的事情.
有回報的事情.
即是會升的事情.
如果只是去做這些為了效益而去做的話.
純粹去思考有沒有用.
有沒有後果去做這件事呢.
這個其實不一定是叫教會.
這是一盤生意.
不是說結果不重要.
教會也是要去看看情況會有什麼結果.
但是我發覺很多事情都不只是去計較結果.
最少你不能夠看即時結果.
牧羊,上帝的使命,聖經給我們的命令.
這些從來都不是叫你想後果而去做.
而是單單你要去實踐出來.
所以教會不是一個這樣的地方.
教會就是一班人大家一起很努力.
付出很多心血.
但往往未必會看到即時很多的果效的地方.
如果好像外面公司那樣.
一不行就cut.
一不行就要裁員.
教會不是這樣去做.
回到經文裡面.
十五頁的經文我們發現故事並沒有停留在那裡.
大衛那三個勇士將水帶下來給大衛那裡.
然後大衛將它們囂電在地上.
大衛沒有停止他所說的話.
大衛繼續說非常重要的話.
大衛強調他所倒下的水不是純粹倒在地上.
而是他特別稱呼為將水電在耶和華的面前.

$^{401}$並且禱告耶和華.
這三個人冒死去打水.
這水就像他們的血一般.
原來一班人一班基督徒一起去做一些事.
其實不一定是講求那個結果.
大衛將整件事的本質去改變.
或者說大衛重新去定義這件事的本質.
將它變成一件信仰的事情.
不是純粹送外賣.
而是一件信仰的事情.
當他們做一件信仰的事情的時候.
整件事我們思考這件事的標準和思維就很不同.
這三位勇士冒死要從伯利恆那裡打水過來.
他們的心血代表著他們的生命.
根本上是一個將血獻上的祭.
所以大衛去定義這件事是向耶和華面前獻上的祭.
簡單來說這是一個生命的敬拜.
真人如果我們純粹從一個世俗的角度去看.
即是從外面的人眼光去看一間教會的時候.
信仰其實很多事情都沒有什麼特別回報.
沒有什麼特別具體的利益.
整件事是不會一直升或升多少.
什麼時候止蝕這些事情.
信仰往往都是去做一件你應該要做的事情.
敬拜本來就是一件這樣的事情.
敬拜是沒有輸贏的.
敬拜就是敬拜.
敬拜就是一件你需要做的事情.
一班基督徒聚在一起.
好像在元朗錦繡花園的角落裡.
大家每個星期聚在一起.
好像做完之後對這個世界沒有什麼特別的改變.
不過是可以的.
《羅馬書》十二章.
大家都熟經文.
正正是說這個事情.
本來說將身體獻上.
當作活祭.
這個是理所當然的敬拜.
我們基督徒的生命.

$^{441}$我們基督徒的生活本來就是一場敬拜.
你根本敬拜是應該要做的.
不是問敬拜之後有什麼好或不好.
有些是贏有些是輸.
敬拜就是你不需要太理會.
它的成敗得失.
有沒有後果.
有沒有用.
因為無論這件事情的目標是怎樣.
結果是怎樣.
成效是怎樣.
這個本來就是在耶和華面前的敬拜.
這個就是本身你要做這件事情的原因.
我想教會也是一樣.
教會屹立在錦繡花園這個地方裡.
本來就是一場在上帝面前的敬拜.
我們就去做這些事情.
每件信仰的事情最後都逃不過最ultimate的上帝.
我們就在祂面前敬拜祂.
屬於祂.
教會沒有一個行動不是敬拜.
我們聚在一起.
一起去做一些事情.
正正就是一種敬拜.
大家看到教會掛了很多燈籠.
昨天教會有中秋節聚會.
我們每年都做的.
或者每年都做很多這樣的街坊活動.
你問這些活動最後有多少成果呢.
有什麼人數帶回來呢.
可能你會覺得其實做出來之後都不知道有什麼用.
我都不知道有什麼特別的.
看到有什麼數字.
不過這個就是我們基督徒要做的事情.
我們大家聚在一起.
作為基督徒.
正正就是去服侍這些街坊.
告訴他們耶穌很好.
我們一起去做一些事情.
無論是在監牢也好.

$^{481}$在錦繡堂也好.
這邊也好.
我們就是不斷去做基督徒應該做的事情.
教會應該做的事情.
我自己在今年四月.
代表建度負責全港的香港教會前景會議.
在這幾年有很多我們叫做COVID移民潮之後.
第一個全港教會的大型會議.
去一起集氣.
大家聚在一起.
香港教會可以怎樣走下去呢.
你知道很多教會人數下跌了很多.
無論是COVID移民.
基本上很多教會都真的.
不要說頂智慧.
牧者可能都不在了.
我發覺重新來到去.
給一個答案是不容易的.
有些什麼方法.
有些什麼方法有結果呢.
能夠立刻說出來.
上台上說完之後.
大家就跟.
就能夠扭轉這趨勢.
你發覺是沒有這麼容易的答案.
不過我自己參加完這個會議之後.
見過很多教會的不同領袖.
很多不同教會的傾談之後.
發覺最起碼一件事情就是.
大家起碼仍然堅持.
士氣很重要.
起碼不要覺得已經沒有事情可以做.
單單願意繼續去做一些事情.
單單去繼續去做一些教會應該要做的事情.
這個就是我們的意義所在.
無論你說是搞一下教會福音大旅行.
搞一下幸福小組.
搞一下中秋節聚會.
什麼都好.
重點不是那些聚會.

$^{521}$那個地方是不是大到什麼用.
而是你只要願意不要洩氣.
仍然繼續做.
做回上帝交託給我們做的事情的時候.
這個就是一個最重要.
我覺得情境會裡面最有意思的事情.
所以在教會我們感受堂.
我們一起在疫情之後.
我們仍然留在香港的弟兄姊妹.
仍然繼續回來敬拜的弟兄姊妹.
只要我們仍然一起去齊心付出.
我們100\%的努力去做一些事情的時候.
這個正正就是我們的敬拜.
不需要問那件事有沒有效.
有沒有用.
最少這個不是我們負責去問的問題.
我們要做的就是付出我們的生命.
放在耶和華面前.
將它成為我們應當的敬拜.
我一直祈禱.
再次請你繼續使用我們的教會.
讓我們仍然去堅持.
發揮我們教會的功能去見證你.
去切實的彼此相愛.
每個禮拜仍然在敬拜你.
我們不知道這些敬拜.
這些事情.
這些活動帶來什麼具體的果效.
我們雖然會去思考.
但我們知道這個不是我們唯一的判斷理由.
我們知道我們要盡我們的本份.
我們要做的就是獻上我們的生命.
用出120\%的努力.
就好像當日三位勇士一樣.
將我們的生命獻上.
成為一個理所當然的敬拜.
求神你給我們有智慧.
給我們仍然有一份的信心.
繼續的來為你做好事.
一起去做一件死所欺人的事情.

$^{561}$奉尊名求.
\newpage



\section{撒迦利亞書}
\label{sec:CXJyuYm4_pE}
\textbf{2023-10-05 周四晚查經班《歸回先知書研讀(4)撒迦利亞書》 講員:曾裔貴牧師}
\newline
\newline
連結: \href{https://youtube.com/watch?v=CXJyuYm4-pE}{\texttt{https://youtube.com/watch?v=CXJyuYm4-pE}} ~~~~ 語音日期: 2023-10-06
\newline
\newline
\hyperref[sec:JHYqHgNM3GQ]{\small{< < < PREV SERMON < < <}}
~
\hyperref[sec:index_chronic]{\small{[返順時目]}}
~
\hyperref[sec:index_scriptual]{\small{[返順卷目]}}
~
\hyperref[sec:G0FSq209Kac]{\small{> > > NEXT SERMON > > >}}
\newline
\newline
$^{1}$我們可以在晚上再思想撒迦利亞書的經文.
雖然都是一些很長時間的寓言.
但是那個持久和對於日後耶穌基督的降臨.
和關於以色列這個民族.
其實都是有那個研究的價值.
都幫助我們能夠明白.
如果有些是深的.
都讓我們能夠得到一些教導.
使我們能夠慢慢掌握經文的意義.
我們今天晚上等候在你的面前.
多謝你讓我們能夠在教會.
在網上我們都可以來學習聖經的真理.
我們這樣禱告奉耶穌基督的名.
阿們.
我們簡介上一個禮拜.
特別是提到第四個異象.
就是這個大祭司.
他在聖殿裡面.
這個撒迦利亞先知他所見到的.
尤其是強調大祭司的積焚.
和當他見到這個異象之後.
尤其是髒衣服脫下.
然後又有潔淨的衣服和戴冠冕.
然後實實際際叫約書亞.
當然是叫這個先知撒迦利亞.
來跟約書亞說.
大祭司約書亞說.
這個是第三章的第六節那段經文.
然後就引述引進了彌賽亞大衛的子孫.
和這一塊石頭.
神親自雕刻的石頭.
神是怎樣來栽培他的愛子.
將來成為穩固的根基.
防國的頭塊石頭.
不是用很漂亮的裝飾.
而是他沒有階型美容的.
我們受苦的主.
這一位憂患之子.
是為憂患而生的.
一點都不是寓意詳片那些的苦情.

$^{41}$是他要注定來這個人世間.
是為人類的罪受苦的.
我們的耶穌基督.
所以他的人生就是這樣.
常經憂患.
常常經歷憂患.
多受痛苦的.
今天去到第五個義將.
總共有八個義將.
第五章就有三個義將.
一併地說完.
其實應該去到第五第六章.
今天第五個義將金燈台.
很特別.
金燈台然後就會說到所羅巴百.
所以一個是大祭司.
一個是那個起聖殿的領袖.
即是鼓動百姓回歸的所羅巴百.
兩位的領袖.
第四章的第一節開始.
那與我說話的天使.
又來叫醒我.
好像人睡覺被喚醒一樣.
突然間被人叫醒.
他問我說.
你看見了什麼.
我說.
我看見了一個純金的燈台.
頂上有燈盞.
燈台上有七盞燈.
每盞有七個管.
即是每一盞燈有條管.
有一條喉.
來輸送油.
去燈盞.
即是說自動的.
不斷有油.
不用添油.
有別於聖殿的.
本身的擺設.

$^{81}$因為大祭司.
或者祭司也可以.
他可以進去聖殿.
或者進去會幕.
他就是做那個點燈.
即是更換那些燈芯.
去到燒到差不多沒有了.
有沒有燒過油燈.
你們小時候.
即是大光燈都是要用那些燈芯.
如果燒完.
你就要弄高一點那支燈芯.
一路這樣燒著.
少少油這樣碰著它.
它就一路這樣燒.
大概燒一晚.
大光燈.
它自己會熄的.
底座就是裝著火水.
這裡就告訴我們.
有條管直接去到燈盞.
剛剛睡醒.
你看人家對於那個儀仗.
有時我們如果越仔細讀經.
就越能夠對那個經文.
有一個透徹一點.
深入少少的默想.
和那個發言.
人家剛剛被翻醒.
撒加令啊.
在儀仗裡面被天使拍醒了.
好像有人睡醒了.
醒了.
拍醒.
你見到什麼.
見到燈台.
一句說完.
不是.
逐樣去描述.
你就知道原來.

$^{121}$這個和我們讀經的態度都有關連.
如果我們去認真讀經.
是會有很多的發言.
因為他等一下就要問天使.
是什麼意思呢.
為什麼會這樣呢.
接著說第三節.
旁邊還有兩棵的橄欖樹.
在聖殿裡面有棵樹.
一棵是在燈台的右邊.
一棵在燈台的左邊.
我問與我說話的天使說.
主啊.
這是什麼意思呢.
與我說話的天使就回答我說.
你不知道這是什麼意思嗎.
其實是很噁心的.
在聖殿裡面.
為什麼會突然間有兩棵橄欖樹呢.
一棵在左右擺設.
你現在不是像我們過年.
擺兩盆桔子.
擺大門口擺兩盆桔子.
是對的我們中國人.
但是你在聖殿裡面.
明明是金燈台.
金燈台是沒有東西的.
是有另外有的.
是有五桌裝陳設餅的.
另外一個香壇.
是三樣東西.
聖殿怎麼會突然間多了兩棵樹呢.
所以這個一定是一個異象.
有別於正常的聖殿.
其實索羅巴伯.
對不起.
撒迦利亞.
如果他進到聖殿.
他都要死.
因為他不是祭司.

$^{161}$他是先知.
第六節.
他對我說.
天使就對他說.
這是耶和華指示索羅巴伯的.
上文第三章就指示大祭司約書亞.
現在先知撒迦利亞見了兩件事.
第三章就是他見到.
關於大祭司約書亞的情景.
第四章又關於省長索羅巴伯的情景.
就是耶和華指示索羅巴伯.
萬君之耶和華說.
不是依靠勢力.
不是依靠才能.
乃是依靠我的靈方能成事.
我們很喜歡引這節聖經.
四章六節很多人引的.
其實很好用的.
真的好像萬能key.
一到祈禱會.
或者是一些類似是報道會.
就會引這節聖經.
但是上下文說什麼呢.
這裡說的不是依靠勢力.
不是依靠才能.
是要依靠神的靈方能成事.
現在他們在做什麼呢.
在先知白天的時候.
現在這個異象.
白天的時候.
很明顯就是他們現在被鼓勵.
重新建造荒廢了.
來破壞了美的聖殿.
建了根基之後.
就被人來干預.
那些原居民.
阿拉伯人干預.
就來告官.
然後就要逼令他們停工.
因為你都沒有開工紙.

$^{201}$那怎麼讓你開工呢.
然後就停了.
一停就停到大利烏王的第二年.
然後哈該先知和撒迦利亞先知.
就有神的密使.
來去抗訓那些百姓.
你們還住在天花板的房屋.
神的殿方樑.
你們都不理嗎.
這是哈該說的.
然後兩位先知就來勸勉.
其中一定是以這個省長做單位.
他一定是要現在焦聚公頭.
焦聚不同的公將.
還有上山要取一些木.
因為要木造的.
要最好的木.
然後來去畫跡.
然後來去平頂.
很多很多工程.
所以是要依靠勢力的.
是要依靠才能的.
是要依靠人的.
是要振奮軍心的.
還有一些人可能要做飯吃的.
那些婦女又要煮飯.
有時唱下歌.
令到那些百姓那些男丁.
就很振奮.
大家一窩蜂一條心.
去興建重建神的聖殿.
因為他們上來回歸.
都是這個目的.
由巴比倫.
由波斯帝國.
回來就是想建聖殿.
所以這個時候.
他們在這樣的情況下.
很肯定是在用人的力.
很肯定是省長運用他的影響力.

$^{241}$去推動整件事的.
但是才見到一個很重要的儀仗.
是關乎到這個所羅巴巴.
大家要留意.
記住那個圖像.
在聖殿裡面.
他特別見到的就是金燈台.
然後有七盞燈.
這些不特別的.
金燈台本身就是七個燈枝.
即是點七個位置.
特別的地方就是.
那些燈盞裡面是有條管.
每盞燈的燈盞都有條管.
那條管.
待會我們告訴大家.
是搭去那些橄欖油.
所以橄欖樹在這裡.
就是那個原因.
即是直接有源頭.
由源頭供應那些油.
不用那些祭司每天去添油.
你可以想像得到.
第七節.
大山你算甚麼呢.
在所羅巴巴的面前.
你必成為平地.
好像神蹟.
耶穌講過的.
你一粒芥菜種.
這麼小的信心.
就是叫這座山移開.
它就馬上移開了.
耶穌這樣教導的.
這是盧伽芬第十章.
典故原來在這裡.
原來由這裡而來.
接著就是.
他必搬出一塊石頭.
安在殿頂上.

$^{281}$人且大聲歡呼說.
願因為歸與這殿.
接著第八節.
耶和華的話又臨到我說.
所羅巴巴的手納了這殿的根基.
他的手也必完成這功.
你就知道萬軍之耶和華.
猜顯我到你們這裡來了.
誰藐視這日的事為少呢.
這七眼乃是耶和華的眼睛.
遍察全地.
見所羅巴巴手拿善陀.
就歡喜.
原來這裡有一個暗示.
是讓先知撒迦利亞知道.
現在在外面在白天.
所羅巴巴和他的族弟兄.
在這裡重建聖殿.
到底能否完成呢.
其實用了四年多.
接著先知在異象.
在聖殿裡面.
他聽到神自己親口說.
他納了根基.
一定可以完工.
因為不是依靠勢力.
是依靠神的靈.
所以原來很重要是.
有神的同在.
有神的力加在所羅巴巴身上.
由誰去輸送這個力呢.
一定是先知他領受了這個異象.
神向他去講解這個異象.
他是比所羅巴巴.
更加知道這個工程的始末.
更加知道所羅巴巴.
他可以儘管日間拼命地做.
但原來真正的力量.
是來自神在他內心的聖靈工作.
和先知的勸勉,安慰和加力.

$^{321}$所以這裡真的很重要.
我們見到舊約.
已經有聖靈工作這個部份.
所以引這支聖經是對的.
一般的工程.
有一些重大的福音工作.
我們引這支聖經是對的.
不過整個典故.
都是告訴我們.
先知就要為這一方面去禱告.
先知是應該在暗室.
為日間的工作.
來去更多的禱告.
以至那個省長.
他不只是依靠他的力量.
他的智慧聰明.
因為依靠的勢力.
就是他的人脈,才幹,才華.
人見得到的影響力,親和力.
但如果裡面沒有聖靈.
一切都是人的工作.
先知就是要為那個領袖.
來去禱告.
以至他能夠感受到.
一切都是聖靈工作.
他能夠將榮耀歸給神.
一切都是.
如果不是神給我有這樣的智慧.
有這樣的機遇.
有這樣的人脈.
我怎能夠做得成呢.
所以這句說話嚴格來說.
是向那個領袖所羅巴伯說的.
所以一切日後有重大要帶領的.
就是需要去學習這個道理.
先知就是更加第一時間知道.
神要他去傳話給所羅巴伯.
讓他知道其實神和你同在.
神在你的工程裡面.
因為這個是神的工程.

$^{361}$你去做一定有神的幫手.
和加的力量.
而且這裡告訴我們.
誰去藐視這一天的事為少呢.
這七眼乃是耶和華的眼睛.
不是怪物.
不是異形.
是造物主遍察全地.
即是說他看得到.
見所羅巴伯的手拿起善陀就歡喜.
一個領袖拿個善陀.
意味著他也是那個.
落手落腳去做的人.
他不是高高在上.
他真的是落手落腳去做.
當然一個的暗示.
不是叫他每天搬東西搬柴.
他應該做的策劃他就不做.
不是這個意思.
而是告訴我們.
神的眼目是七眼遍察全地.
所以當然不代表我們回家的晚上.
可以大無私一樣不關門.
就不怕有賊.
打開門都不怕的感受.
就不是這樣的.
因為我們是信主.
那些賊不是信耶穌的.
他不會覺得有神的七靈.
七眼.
我們要記住.
我們做了那些事之後.
但其實我們內心要明白.
神是七眼.
他是遍察全地.
即是說原來有些角度.
連我們想都想不到的恐慌.
其實神已經幫我們想了.
幫我們保護了.
我們想到的就盡量去做.

$^{401}$但有時這些是我們想不到的.
超出了我們的預算.
你怎會想到黃大仙廣場會水浸.
沒有人會想到的.
香港還不是在很遠的.
內地或非洲.
這些就是我們想不到的東西.
但記住我們信的神.
他講到明的.
神的七眼遍察全地.
不要藐視這一天為少.
繼續讀下去.
我又問天使.
這登台左右兩棵的橄欖樹是什麼意思.
讀經仔細就有這個好處.
看完之後.
很奇怪.
猶太人我們小時候就知道.
上利未記.
讀出埃及記.
聖殿裡面.
那時還沒有聖殿.
是會幕.
會幕裡面有什麼擺設.
雖然我進不去.
但我知道的.
我們上主一學.
我們知道的.
小時候.
即是說猶太人.
裡面只有三樣東西.
只有三個重要的神器名.
金燈台.
多少爪.
三爪最高.
香台.
即是燒香的香壇.
不是香台.
香壇.
是三爪最高.

$^{441}$另外就是一個餅桌.
有十二個聖餅.
所以大衛逃難的時候.
問了祭司.
我忘記了名字.
拿一些陳設餅.
祭司說.
這些陳設餅.
不是可以隨便吃的.
除非有特殊原因.
大衛說我有皇命的緊急.
你給我吧.
那祭司就不如有詐.
給了他.
耶穌有引那個典故.
大衛吃陳設餅.
都尚且無罪.
因為很特殊.
所以神是有恩典的.
有餅桌.
有香壇.
還有金燈台.
只有金燈台.
沒有其他.
怎麼會突然多了兩棵橄欖樹.
現在不是過年.
有兩盆桔.
是有兩棵橄欖樹.
所以讀經仔細.
他一見到那個圖像.
其實剛被人拍醒.
見到那個圖像.
很奇怪.
他就很想問.
為什麼會有兩棵橄欖樹.
所以很顯然.
就是這個意象裡面.
最重要的啟示.
這燈台左右的兩棵橄欖樹.
是什麼意思.

$^{481}$好奇.
又是讀經的一個很重要的態度.
不是一種為批判而批判.
好奇就不同批判.
好奇就是.
我很想明白.
所以其實.
大家上茶經班要多問問題.
多看有什麼.
覺得很好奇.
就要問.
現在天使就告訴他.
第十二節.
第二次問他說.
這兩斤橄欖枝.
在兩個流出金色油的金嘴旁邊.
是什麼意思.
越看就越仔細.
在那棵樹.
有一個金枝.
橄欖枝.
然後流出金色油的金嘴旁邊.
即是那盞燈.
一定是不會熄滅.
因為有橄欖樹的汁液.
而且有一個管.
直接流去燈盞.
意思就是這樣.
那個圖像.
第十三節.
他對我說.
你不知道這是什麼意思嗎?.
好像又再一次引起他的好奇.
更加大的好奇.
就是天使說.
什麼你不知道.
應該知道的.
見圖識字的.
連那個為什麼要有橄欖樹.
都已經是一個好奇.

$^{521}$原來橄欖樹不是擺設.
是一個功用.
是樹的汁液直接去點燈.
我說煮一下我不知道.
第十四節.
他說這是兩個受高者.
站在普天下主的旁邊.
不說油字可.
說到就非同小可.
就是一節聖經解完整個意象.
原來整個意象是想說這兩棵橄欖樹.
是在普天下的主面前站立的兩位使者.
即時的應驗.
一定就是大祭司若孫亞和所羅伯伯.
一定是這樣.
即時的應驗.
當然你可以說.
撒迦利亞和哈該都可以.
但似乎就不是.
因為那個聖殿建好之後.
那個禮儀.
關於獻祭就一定是大祭司若孫亞.
現在建聖殿的工程.
一定是責任在省長所羅伯伯身上.
所以現在神透過天使.
給先知撒迦利亞看到這個意象.
這個意象沒什麼特別.
但唯一特別的地方就是有兩棵樹.
不過要說明是橄欖樹.
我們看到可能是盆桔.
可能是一些其他的東西.
可能是過年沒有菊花.
那些是劍蘭類似.
你都要分到那個種類.
他不單見到兩盆花.
還見到是橄欖的樹.
有橄欖的枝.
流出的汁液.
那你就可以想像到兩個見證人.
在新約出現過.

$^{561}$在啟示錄的我和大家揭一揭.
我們以前在啟示錄都讀過.
有兩個見證人出現的.
我忘記了是十二還是十三章.
對不起 十一章才對.
有一根位置.
這個啟示錄的十一章.
有一根位置似給我.
當作糧度的帳.
且有話說.
起來將神的殿和祭壇.
並在殿中禮拜的人都糧一糧.
只是殿外的院就要留下不用糧.
因為這是給了外邦人的.
他們要遣答聖城四十二個月.
三年半.
我要使到我那兩個見證人.
穿著毛衣傳道一千二百六十日.
你計算一下.
其實就是三年半.
他們就是那.
當然是三百六十.
隨得盡的猶太人三百六十.
我們三百六十五.
又四分一日.
所以也不是人家不準確.
不過人家是用那個.
曆法不同的意思.
第四節他們就是那兩棵橄欖樹.
有沒有留意到.
兩個燈台.
立在世界之主面前的.
就是我們剛才讀的四張儀仗.
第五個儀仗.
金燈台儀仗.
金燈台有兩棵橄欖樹.
有生命的光.
當然就是那些油.
去到那個燈盞.
那個燈盞就是源源不絕.

$^{601}$有燈有油.
就可以點亮.
照亮這個世界.
所以原來傳道人.
是照亮這個世界.
照亮這個地上.
兩個見證人.
在這個時候.
其實已經是大災難的時候.
都有兩個見證人.
在這裡來到去.
怎樣為神.
立在世界之主的面前.
是兩棵橄欖樹.
是兩個燈台.
一對的.
可能彼此有照應.
又不知怎樣.
所以很奇妙.
將來怎樣那兩個人.
真的只有他們兩個.
怎樣去傳道呢.
其實很孤單.
很容易的.
有三年半他要這樣來傳道.
接著我們讀下場.
第四節.
對不起.
第四節我讀了.
第五節.
若有人想要害他們.
就有火從他們口中出來.
消滅仇敵.
這些我們解不了.
是不是真的好像.
他的口有噴火呢.
還是用其他的方法呢.
我們不知道.
凡想要害他們的.
都必這樣被殺.

$^{641}$這兩個人有權並.
在他們傳道的日子.
叫天閉塞不下雨.
又有權並叫水變為血.
叫天閉塞不下雨.
你想起誰才知道.
以利亞.
叫水變血.
想起誰.
摩西.
並且能夠隨時隨意.
用各樣的災殃.
攻擊世界.
他們作完見證的時候.
那從無底坑上來的獸.
必與他們這兩個人交戰.
而且那個獸贏了.
這個就是地上來的地獸.
我們過去說過.
如果我們真的這樣解釋的時候.
那個地獸.
其實就是宗教邪惡的勢力.
並且德勝把他們殺了.
他們的屍首就倒在大城裡的街上.
這城按照著靈異.
就是索多瑪又叫埃及.
與其不是就是.
以利亞敵黨耶誓別.
阿哈摩西敵黨法魯.
就是埃及.
雖然.
以利亞是和這個.
西頓的女王耶誓別.
但無論如何.
按靈異.
就是說敵黨神的.
就會是這樣.
原來將來有兩個這麼有power.
傳道又傳又用神蹟.
但是最終都被地獸贏了.

$^{681}$殺了.
接著留意第九節.
從各民各族各方各國中.
是不是zoom呢.
可以同時間馬上看到.
全世界.
你們傳道三年.
三年半.
現在就給你們.
三天半不准埋葬屍體.
又不可以埋葬放在墳墓.
如果對於某些信仰來說.
如果你不馬上埋葬.
就是上不到天堂.
拉登都要馬上24小時就要埋葬.
不知道是海葬還是怎樣.
不然你激發起那些阿拉伯人的反抗.
就不行.
所以那些海豹部隊馬上幫他埋葬.
第十節.
地上的人為他們的死歡喜快樂.
互相攜送禮物.
因這兩位先知.
曾經叫住地上的人受痛苦.
過了三天半.
有生氣從神那裡進入他們裡面.
他們站起來復活.
看見他們的人甚是懼怕.
兩位先知.
又是見證人又是傳道者.
更加是先知.
兩位先知聽見有大聲音從天上來說.
對他們說.
上到這裡來.
即是被提.
他們就駕著魂上了天.
他們的仇敵也看見了.
正在那時候地大震動.
城就倒塌了十分之一.
因地震而死的有七千人.

$^{721}$其餘的都恐懼.
歸榮耀給天上的神.
這個是第一個災禍.
第六之號的中段插曲.
突然間加插了兩個見證人.
仍然留意.
他們的描述和我們今天看的完全一樣.
他們就是那兩棵橄欖樹.
兩個燈台.
立在世界之主面前.
所以很特別.
大衛起聖殿.
當然不是他起的.
是他的兒子起的.
所羅門起聖殿.
真的在他的聖殿門口.
有兩條銅柱.
一個叫牙根.
另一個我忘記了名字了.
要看列王記上起聖殿那裡就會看到.
大家可以回家看一看.
或者在網上幫我找一找.
就有兩條銅柱.
就好像一個見證人.
這樣來到去立在神的面前.
所以無論舊約.
先知撒迦利亞.
被天使指引他.
看到這些異象.
他看到的是金燈台旁邊的橄欖樹.
左右的橄欖樹.
和橄欖枝有油流出.
來讓燈盞可以生生不息的點燃燈.
意味著神的工程已經開始了.
在當時的情景.
那個場景.
工程開始了.
而且一定這次不會爛尾了.
因為神照管著.
所以先知要傳遞給所羅八百的.

$^{761}$就是不要倚靠勢力.
不要倚靠才能.
雖然是的.
是要倚靠.
但是更要倚靠的就是神的靈.
方能夠成事.
因為這件事本身就是一個神的工程.
所以教會的聖公.
如果沒有了祈禱.
沒有了大家靜心的.
同心的去倚靠.
其實多麼漂亮.
多麼人間的.
多麼人為的.
如果神不認可.
其實都是沒有意思的.
真的沒有意思的.
所以這是一個很大的學問.
一個自會的態度.
一個讓一些工程能夠.
真的有質量保證.
不是豆腐渣.
就要是真的經得起時間的考驗.
所以不是一些煙花式的興奮.
是很踏實的.
大家是很認真的.
不需要大多打鑼打鼓出聲的.
默默地做.
默默地去承擔教會的聖公.
就是這樣.
很多的教會就是這樣建立的.
所以有些人都說.
教會去到第三第四代就會沒落.
因為沒有第一代的人.
那種艱苦.
那種同心看到一個異象.
要開始這裡有一個教會.
大家一起出心出錢出力.
這樣去做.
所以很多年長的弟兄姊妹.

$^{801}$還很懷念這個教堂的起好.
是因為真的由零.
朝陽晚策的幼稚園.
真的每個星期都要搬來搬去.
地方又不夠.
又隔不到聲.
很多的很簡陋.
到一個教堂.
一次性的來到.
大家奉獻.
那個心是很強烈的.
這是第一代.
這是第一代的聖徒.
那個內心.
你打開他的內心.
就是單純因為福音的愛改變他.
到第二代了.
很明白上一代的困苦.
他們的擺上.
很多的故事.
他們是口耳相傳.
他們是很明白父母輩那份愛主的心靈.
所以他們成長.
都是有這些的記憶.
所以他們又可以很殷勤的讀書.
來到去有社會地位.
有成就.
他們會很有那個心.
就是為神去發展他們的事業.
他們的工餘的時間.
都是為神去做.
第三代了.
危機就出現了.
沒有第一代的祖父母輩那些故事.
也都沒有那個經歷.
因為一出來就是一個教堂.
一出來自小就在幼稚園.
這樣來上主一學.
慢慢慢慢.
根本不是很明白很多的故事.

$^{841}$也都沒有經歷.
都不需要經歷神蹟.
因為太順暢了.
危機就開始了.
第三第四代.
似乎都好像是一個模式.
不是說感宣家.
絕對不是就.
絕對沒有這個意思.
而是一般的教會.
去到第三第四代.
開始出事就是這樣.
就變成.
將那些變成一些教產.
就看為是必然的.
應份的.
是應得的.
是屬於我們第三第四代的.
這個時候.
從來都沒有落手落腳.
平手抵觸.
來奉獻.
又當今年.
又怎樣來到去.
殷勤.
怎樣不吃.
我都要為教會建教堂.
來奉獻.
第四代沒有這個東西.
沒有這個問題出現.
危機就在那裡出現.
因為沒有信心的功課需要去學.
這個告訴我們.
我們求主要幫助.
我們要留心.
我們已經進入第五代都有了.
我們其實真的要小心.
第四第五代他們成長之後.
他們肯定會成為某些人的.
成為執事.

$^{881}$成為信徒領袖.
他們的眼光有沒有第一代的聖徒.
我們的前輩們的眼光呢.
有沒有那些謙卑.
有沒有那些信心呢.
所以這個正正就是.
這個牧養的問題所在.
所以這兩個見證人.
和這個.
撒加尼亞.
見到的兩位領袖.
其實那兩棵的橄欖樹.
在聖殿裡面.
其實意味著.
所羅巴伯和大祭司約書亞.
原來在神的眼中.
其實他們是聖殿裡面的橄欖樹.
是見證人來的.
他們的工作就是要見證神的榮耀.
和能力.
所以這個也可以這樣說.
真是讓後世.
到了耶穌基督出來了.
真是不斷揀選的是那些見證人.
他們那三年多跟隨耶穌.
其實學不到什麼技巧.
學到的就是耶穌經常罵他們.
指責他們這樣對這樣不對.
阿彼得就比較多一點.
因為很多時候他就會.
出了他自己個人的性格.
但是原來這些.
真是生命的功課了.
真是門徒的訓練了.
真是讓他們將來沒有耶穌之後.
他們就真是可以來到去.
用他們的品格.
耶穌訓練的生命去做服事了.
所以他們能夠帶來.
整個世界的震動.

$^{921}$就是有見證人.
立在世界之主面前的見證人.
所以兩個已經夠數了.
不是加計會的口號.
兩位見證人.
是兩個燈台.
所以你會看到橄欖樹.
是有那個執液燈台.
就是說有源源不絕.
他們有神的聖靈.
有神的生命.
有神的話語.
來供應他們那個時代.
譬如起聖殿的時代.
到後來以色列上來.
都是興起以色列.
再後來就興起我們看過的尼希米.
都是這樣的一脈相承.
他們都是被神的靈感動.
所以原來更重要就是.
兩個見證人.
是兩位先知.
是兩棵橄欖樹.
是兩個燈台.
還要是立在世界之主的面前.
所以傳道者的眼界.
不要那麼狹小.
是全地的.
眾物主是七眼.
他是遍察全地的.
所以我們要有一個胸襟.
看到別人的工作.
我們要感恩.
我們要來去欣賞.
我們要為對方做得不好的去祈禱.
就不是批判.
就不是見高就拜.
見低就踩.
做不到傳道.
一般基督徒都不應該.

$^{961}$何況傳道人.
我們應該有這樣的胸襟和眼光.
所以這個是很重要的.
所以傳道者我們應該.
無論是多留意新聞之餘.
更加要看到.
整個世界是立在眾物主的面前.
所以我們要為一些.
未發生的事要祈禱.
守望.
去制止那些黑暗.
罪惡蔓延.
這些是我們的責任.
所以我們回來這裡的時候.
你就可以見到.
這個是很重要的一個.
被撒迦利亞先知.
看到的異象.
他看到兩位領袖看不到的.
就是神已經指示撒迦利亞.
看到關於大祭司.
原來大祭司都需要脫下那件污糟的衣服.
真正能夠事奉的人.
是穿上神的衣袍.
戴的是神給的冠冕.
所以不要乞求人對你的讚賞.
哪需要的.
牧者傳道就是有一個超然的地位.
你不需要讚我.
因為我做了我應該做的.
神就喜悅就行了.
不需要的.
你看不起我.
貶低我.
都不介意我盡了我的力.
我只有這麼多的料.
我為什麼要扮呢.
不需要扮的.
做好我自己.
那就是神喜悅就夠了.

$^{1001}$所以我們要記住.
就算基督徒都一樣.
我們都要追求在神眼中的那個你.
是蒙他的喜悅.
耶穌都是這樣教我們的.
所以我們要盼望.
好像撒迦利亞先知所見到的異象.
見到大祭司穿污糟的衣服.
在聖殿裡面穿污糟的衣服.
不得了.
大家又不熟悉了.
大祭司去公職的時候.
首先要去到門口的洗浴盆.
有盆水在那裡.
要洗身的.
然後要換衣服的.
換那套制服.
很講究的.
穿金戴銀的.
又有十二顆寶石的.
然後就進去.
還要帶著香爐進去.
進去之後.
在致聖所裡面.
要把香爐擺好.
讓整個致聖所.
就是那個聖所裡面.
還有一個莫萬遮住的.
就叫致聖所.
致聖所有什麼呢.
只有一個藥櫃.
什麼都沒有.
是沒有燈的.
你不要帶著光管.
帶著燈.
電話開著眼燈就死了.
進去之後其實已經黑黑的了.
要把煙遮滿整個致聖所.
為什麼呢.
讓神的憤怒公義.

$^{1041}$看不到你這個大祭司.
真的是這樣教導的.
這樣才可以進去彈血.
其實大概的.
因為什麼都看不到.
拿些血彈七次.
這樣就走了.
還要退後走.
出來的聖所有燈的.
因為有金燈台.
出來聖所之後又再出去.
就是外面的回望門口.
就是這樣了.
再更換衣服.
就可以把陳設餅更換.
就可以吃.
和獻祭員那些就可以吃.
當然全年要獻的.
那個贖罪祭就不准吃的.
全隻贖罪祭生都要燒掉.
這個就是大祭司進聖殿.
但是現在很搞笑.
撒迦利亞先知.
看到那個大祭司.
穿著髒衣服.
進入聖殿裡面.
但是原來.
令到他潔淨的是神.
幫他更換那件衣服.
和給他一個華觀.
所以我們回來.
查經班的弟兄姊妹.
記得我們在主耶穌裡面.
你已經有你的冠冕.
你已經有你的乾淨衣服.
你已經有你的位份去服侍.
從祈禱開始.
從讀明白聖經開始.
你會開始有神的心腸.
一路這樣擴大.

$^{1081}$你看其他人的眼光都不同.
你會坐在這裡.
為今天的講者,領唱,主席.
會去祈禱.
你不會有那麼多平頭品足.
而是你有一個心腸.
這個時刻是神聖的.
神同在的時刻.
我就應該要參與在這個敬拜裡面.
這樣的心情.
你日積月累.
其實你慢慢就變成虔誠的聖徒.
就是這樣建立上來的.
所以不是有一個.
今天很興奮.
那些歌很棒.
他們講話很好笑.
那些全部可以是一個協助.
但是重點仍然是我們的心來到神面前.
第四章見到的所羅巴伯.
原來他的勢力是來自.
必須是依靠聖靈.
接著用一個就是.
兩棵橄欖樹.
在聖殿裡面的出現.
代表著,暗示著.
原來那個見證在神的面前.
就是說我們如果多在神的面前.
我們就不怕在人的面前.
明不明白那個道理.
因為我們多在神的裡面.
明白神的心意.
我們就不怕在人的面前.
為神作見證.
就不怕.
不知道初信主謝法你怕不怕呢.
試過了.
馬上彈起.
怕死人知道你在那裡閉上眼.
然後眨眼睛.

$^{1121}$其實你是正在祈禱的.
初信的時候衝紅燈.
腳都癢了.
都想跟人過.
但又好像不好意思.
有這個情況.
但是慢慢那個生命成長.
你要記住.
我們是立在聖殿.
立在兩個金燈台的側邊.
接著在啟示錄.
甚至乎就是說.
他們在整個世界敵對神的人的面前.
只剩下他們兩個.
他們都是那兩棵橄欖樹.
兩個燈台.
而且是站立在世界之主的面前.
好像將聖殿裡面的秘密.
在外面的世界去展現出來.
所以我們說那個福音.
有時要說到很有力.
讓人感受到.
我聽完.
我如果不信.
我可能會落地獄.
這個人說得很真實.
如果他這樣說是真實.
而我是.
那個時候應該不是在那裡.
嬉笑怒罵.
你說什麼.
怕什麼.
地什麼獄.
不應該是這樣的.
如果那個力救.
不是大聲.
是那個聖靈的力救.
那個人是怕自己落地獄.
他會信.
一些人很懼怕神的.

$^{1161}$他會感受到你說的那個赦罪.
他會有平安.
這樣離開那個會場.
這些才是那個力.
所以兩個見證人就是這樣.
所以撒迦利亞才知道.
他見到兩位領袖見不到的.
見到大祭司有他的瑕疵.
見到所羅巴伯那個力量來源.
原來是聖靈.
所以這些公開了的秘密.
今天我們都識背的四章六節.
其實更加就是我們要實踐.
更加正正就是.
是人都知道的真理.
但是有多少人回到祈禱會.
我絕對沒有意思說.
不回到祈禱會就沒有力量.
不回到祈禱會就是不屬靈.
絕對不是這樣比較.
而是教會是需要更多的禱告.
不是為我們自己.
是為整個世界.
是為整個見證.
很多地方受到迫迫.
我們要關注.
雖然我們可能一生都不會去那些地方.
我們不會在那些地方.
神要我們去傳道.
但是既然有這些資訊.
去到我們的腦袋.
我們是不是應該為他們祈禱呢.
很多事情需要祈禱.
不只是能不能進入哪間學校.
為我們的兒女去祈禱.
而是我們的心要一直這樣再長大.
就是祈禱的心靈的容量.
是要增大.
就會令到我們對神的事情.
我們的關注就會更加多.

$^{1201}$這就是兩棵燈台.
兩棵橄欖樹的遺像.
第五張第六個遺像.
不如我和大家揭一揭.
剛才我很想讓大家知道的.
就是這個所羅門.
我找到了.
第六張就是那兩棵銅柱.
我看看哪裡沒有那麼多.
或者我們由這裡開始看起.
沒有了第六張是太多了.
還是第七張這裡開始會好些.
七章十五節.
他所羅門製造兩根銅柱.
每根高十八爪.
一爪就是尺半.
為十二爪.
為即是那個圓周.
那個Circumference.
那個圓周.
即是圓形的邊界.
又用銅鑄了兩根柱頂.
安在柱上.
各高五爪.
你想想那個圖畫.
很宏偉.
柱頂上有裝修的網子.
和擰成的鏈索.
即是可能用一些銅.
每頂七個.
即是挑花.
網子周圍有兩行石榴.
遮蓋柱頂.
即是滲下來.
一些可能是花.
類似我們有些膠花.
兩個柱頂都是這樣的.
狼子的柱頂.
它的徑就有四爪.
刻著百合花.

$^{1241}$兩柱頂的那些古圖上.
挨著網子.
各有兩行石榴環繞.
兩行共有二百.
他將兩根柱.
立在殿廊前頭.
右邊一根.
起名叫瓦根.
對了沒錯了.
左邊一根.
就是叫做波瓦斯.
有名的.
在柱頂上刻著百合花.
這樣做柱的工作就完畢了.
做了兩根銅柱.
即是說.
他好像都有一些的感動.
要有兩根銅的柱.
當然又挑花很美.
按古代的美感.
我們今天都不知道是什麼.
有兩根銅柱.
立在聖殿的門口.
讓人感受到威嚴.
銅代表舊約.
代表公義.
因為很多都是用銅去做的.
銅的祭壇.
在回望門口.
是用銅的祭壇.
香壇就是金香壇.
有些當然是包了金的.
另外就是陳設餅也是金的.
另外那個燈台也是精金做的.
至聖所的藥櫃更加不用說.
一塌連得的精金去打造.
沒有燒焊.
用人手去打造出來.
所以是很精巧的.
很幼細的功夫.

$^{1281}$是打造出來的.
所羅門說得很清楚.
建聖殿在聖殿的宮地.
不准有鑿的聲音.
真的很誇張.
你可以想像得到.
在很遠很遠的山上.
你要切割好石頭.
運到來.
一運來就凹凹陷.
凹凹陷這樣來到去.
一放下去就要對位.
很難的.
給我鑿一鑿.
都不准.
你可以想像得到.
是很過份的暴君.
但是他就認為.
那個宮地是那麼神聖.
不應該有這些鑿聲.
真的有些誇大誇張.
但無論如何.
也代表著這個聖殿.
是多麼的隆重金碧輝煌.
但是如果裡面的人.
獻祭的人他們不跟隨神.
神說得很清楚.
摩西律法已經說了.
神都不會在這個聖殿裡面.
這個聖殿就會被外邦人踐踏.
他們就會被擄.
所以這個就是.
那個舊恩的歷史.
猶太人的歷史的進程.
在這個時候大衛要為神有居所.
神說你要為我找居所.
我要為你找居所.
所以大衛的投資.
那個回報是很可觀.
他想著.

$^{1321}$我自己住在香柏木的房屋.
神就住在回木的侍羅.
怎麼可以呢.
他就要為神大興土木.
神就跟先知拿丹說.
你告訴大衛.
你要為我建聖殿就不可以.
因為你的手沾了很多血.
你的兒子就可以幫我建聖殿.
但是神反而祝福大衛.
所以這個投資那個回報.
他世世代代大衛的子孫.
就成為了君王.
一直得到神的Blessing.
所以我們一代人為我們的子女去求福.
真正不是王大仙求福.
是為他們的靈性去有好的愛主.
其實是很大的投資.
你不要忘記大衛就是這樣.
所以是很重要的.
為我們的子女去祈禱.
譬如上個星期.
儲牧師說的那個Monica.
就是奧古斯丁的媽媽.
就是這樣一生為她的兒子.
其實她的兒子奧古斯丁.
早年在大學教書.
其實他的生活是很放蕩的.
有很多未婚的太太.
但是終於有一天.
他突然間看到羅馬書.
然後就悔改信主.
這個十三章的羅馬書.
奧古斯丁是這樣改變.
所以我們一代人有意識地知道.
神的恩典會臨到人.
我們的子女就更加應該是受恩典的人.
其實大家調查一下.
我們一起事奉神的人.
我們的教會領袖.

$^{1361}$他們的子女現在在哪裡呢?.
很多人甚至乎都未必回教會.
這些是需要祈禱的.
當然不是拿出來像鞭屍那樣.
不是這樣的.
而是要明白.
其實真正的祝福是說.
大家一起很敬虔地一代一代傳下去.
讓他們能夠感受到神的恩典.
所以這裡告訴我們.
這兩根柱子.
所有人都知道要有一個見證.
所以銅的見證.
都是代表著那個公義.
來到聖殿.
要感受到神的威嚴,公義.
所以這個都是值得我們去想的.
或者我們看看尼加書.
站在這裡的尼加書.
可能都有少許暗示.
尼加書.
所以就是小先知書.
《以摩訛那尼》.
尼加書的第四章.
揭不出來就不要緊了.
我讀給大家聽.
尼加書.
上次我引述過這段.
我弄錯了.
這裡我們已經看過了.
不是說這一方面.
是說神的公義方面.
應該另外一段.
就是不如看《以賽亞書》.
《以賽亞書》的十一章.
十一章就會有這裡.
關於公義的概念.
十一章.
這裡說從耶西的本.
必發一條.

$^{1401}$那個本的意思原文是個「等」.
那個「殊等」.
就是那個木字沒有了上面那一動.
就是叫「等」.
真是叫「等」.
「打等」那個「等」.
「殊等」.
是有這個字的.
必發一條.
即是說在那個樹的底座.
會有一個枝條發出.
從他更新的枝子.
必結果實.
留意一下.
描述這位耶西的本.
這位大衛的子孫.
這位尼賽亞.
耶和華的靈必住在他身上.
就是使他有智慧和聰明的靈.
謀略和能力的靈.
知識和敬畏耶和華的靈.
他必以敬畏耶和華為樂.
這樣描述一個尼賽亞的內心.
有智慧,有聰明,有謀略,有能力,有知識.
重要就是敬畏神.
甚至敬畏神是他的樂事.
回教會是否大家開心的事呢?.
應該是的.
回來就敬拜神.
我們的教堂也是很莊嚴敬虔.
是的.
但是會不會是我們生命裡面.
一個越來越高的位置呢?.
這就是敬畏耶和華為樂.
行審判就不憑眼見.
就是不看表面.
斷是非也都不憑耳聞.
就是說不是聽到就當是.
要求證.
尋求那個證據.

$^{1441}$不會聽完就算.
卻要以公義審判貧窮人.
所以不是要偏幫哪一個人.
不是要助紂為虐.
什麼雞蛋高牆.
而是要誰是公義就要找那個公義出來.
不要管他是富人還是窮人.
高等法院那個.
那個什麼.
正義女神.
遮住眼睛的.
看不到東西的.
意思就是這樣.
憑公義不憑眼見.
斷是非不是憑耳聞.
就是這個意思.
就是說公正言明.
不會因為你是有錢.
或者因為你是窮.
這些正正就是任何的國家.
任何的社會.
一直都解決不了的問題.
我爸爸是李剛.
解決了問題.
我們知道.
就是不是這樣.
就是審判就看人面.
斷是非就是靠.
斷股.
這樣就不行.
那個社會就會亂.
但彌賽亞的日子就不是這樣.
所以其實我們期望.
渴慕彌賽亞降臨.
真正的和平.
真正的公義.
正直判斷世上的謙卑人.
以口中的象擊打世界.
以嘴裡的氣殺六惡人.
留意下第五節.

$^{1481}$公義必當他的腰帶.
腰帶就是束起腰.
一方面是古代的人行動.
準備做事.
就會束腰.
新約羅馬的人都是這樣的.
所以要束起腰.
來準備去爭戰.
這裡告訴我們.
公義是彌賽亞的腰帶.
信實是他脅下的帶.
你可以想像得到.
有公義有誠實.
就是彌賽亞.
這樣描述這位耶西之根.
這一天第十節.
到那一天.
耶西的根立作萬民的大旗.
他就是元首.
他是那個遂令者.
外邦人必尋求他.
就是我們.
他安息之所大有榮耀.
當那日主必二次伸手.
救回自己的百姓所餘剩.
接著就是一直說那個.
下面的我們不再描述.
說到其他外國的外邦地方.
但重點就是.
這個彌賽亞.
這位將來要來的主.
他就是這樣的.
有神的靈在他內心.
他就會有這種生命的特質.
以聖靈所結的果子.
豈不是人愛喜樂和平.
忍耐恩慈良善.
信實溫柔節制.
豈不是這些特質.
和彌賽亞給我們的生命.

$^{1521}$一樣是一脈相承.
所以這些經文.
是比較重要的金句.
大家都要追求和學習.
盡量我們人生就是這樣.
來到有這種.
屬靈的眼光和屬靈的生命.
所以兩個燈台.
兩個見證人.
和所羅門的聖殿.
兩條同處.
相信就是反映公義信實.
讓人見得到.
所以教會是應該讓人見得到.
我們有愛更加有公義.
更加有信實.
所以其實真的.
有些婆婆.
信傳道人多於信自己的兒子.
很多年前.
我在我的舞會做傳道.
有一位婆婆.
有一天叫我和另一個傳道人.
那個傳道人已經不在我們的教會.
去了另一間教會.
就是納維足.
其實他有個兒子.
所以兒子很生氣.
納維足就是我的執行人.
將所有的現金.
我是那個執行人.
要奉獻給博愛的靈養院.
因為他在這個靈養院.
幫到他很多.
他覺得很有需要.
將他的錢給了那些靈養.
作為基金服務.
服務更加多人.
不少錢.
雖然不是過千萬.

$^{1561}$但是兒子就分到一層樓.
兒子就很生氣.
媽媽給了錢人家.
我就很難做.
但是那個納維足執行人.
他講了一句話.
當然他的兒子不在現場.
信教會牧師多於相信兒子.
因為兒子未信耶穌.
對信仰很反感.
所以有時候我都很尷尬.
某程度就是.
很多年老的聖徒弟兄姊妹.
人家心裡面是很清楚.
神的僕人一定會有.
某程度的公義和信實.
為神去服侍.
所以信你多於信兒子.
就在醫生和另外一個傳道的面前.
就去納維足.
我就要執行.
當然後來就很煩.
因為他剛離世.
博愛的靈養院.
就和屯門合併.
沒有了這個名字.
納維足是有這個名字的.
寫到明博愛醫院靈養中心.
沒有了這個名字.
律師說不行,我執行不了.
沒有了這間,就和博愛說.
你吃了它吧,給你們的.
不行,我們不要的.
兒子會告我們.
搞了很久.
後來屯門就說.
不關我的事,不是寫給我的.
我怎會要你呢.
搞了很久,搞了幾年.
要找大律師尋求法律意見.

$^{1601}$又花錢.
最後屯門肯要.
還要靠區議員協助.
公務員就是這樣.
那裡沒有我的名字,我要了.
到時追我怎麼辦.
過了這麼久,兒子也不理了.
因為知道不是他的.
那層樓他都賣了.
類似這樣.
所以其實不是說.
我有什麼特別的東西.
教會是真的應該.
要有公義.
要有信實.
這是彌賽亞的作風.
耶穌基督來到這個世界就是這樣.
人們知道他是.
又公義又誠實.
是誰說的呢?是那些法利賽人說的.
控告他的時候.
拉比,我們知道.
你是不會看人的情面的.
這樣這樣,怎樣做.
對不對.
其實教會追求在這個地上.
就是應該要這樣.
不爭名不爭利.
不是為了.
服侍哪一個群體.
我們是要將人帶進救恩裡面.
這就是教會存在.
唯一這件事如果被人禁止.
我們只能夠用愛.
去接納那些難阻的人.
也不是用武力.
也是用愛.
唯有牧者.
帶頭接受.
對方認為的處罰.

$^{1641}$你不准我說福音.
告訴別人天堂有地獄.
有永生.
有神.
我只能夠繼續用我的生命去宣講.
這個訊息.
我不會告訴別人.
政權現在不准我說.
所以現在這個世界就沒有神.
因為這個是真理.
我們要為真理.
所以就是公義和信實.
我們就要為這個真理付代價.
暫時我從做傳道到今天.
都沒有這些明顯的逼迫.
來受苦.
是不是有生之年要來受呢.
預了.
因為你做傳道就預了這件事.
我們看第五章.
用一點時間.
還有一點時間.
我們看看第五章的第六.
第七個義將.
開始是另一個角度.
我們說的頭五個義將.
頭三個義將.
是在外面看耶路撒冷的聖殿.
聖城.
最後這兩個義將.
是帶領先知進到聖殿裡面.
看到裡面的人.
更加完全看到所有的需要.
這次就不是.
而是看外面的世界.
第五章.
這些連接詞.
告訴我們一個接一個的義將.
有時甚至是睡了.
很搞笑的.

$^{1681}$說明是義將.
就像現在的VR.
可以看到電影.
是不是他們已經古代有這些VR.
甚至是AI遙距.
真的很誇張.
可以看到一些義將.
天可以打開.
看到很多東西.
譬如保羅自己都看到.
在大馬骰的路上.
是日光日白.
有個雲.
有個聲音和他說話.
周圍的人聽不到.
是他聽到.
周圍的人聽到是打雷的聲音.
他聽到是很清楚.
素羅素羅.
為什麼你要逼迫我.
他可以回答.
主啊你是誰.
很害怕的狀態.
這樣說的主.
我就是你所逼迫的耶穌.
有對話的義將.
很奇妙.
第五章.
這次看外面的世界.
見有一幅飛行的書卷.
他問我說.
你看見什麼.
我回答說.
我看見一飛行的書卷.
長二十爪寬十爪.
所以讀經靈修.
真的要細心.
是屬於你自己的時刻.
安靜細心.
你看到什麼.

$^{1721}$要他留意.
天使.
要他先發現了.
然後他就會發問.
所以由發現到發問.
到尋求.
那個好奇心.
得到那個解答.
這個是靈修.
他見到有飛行的書卷.
書卷是在飛行的.
不是UFO.
是長二十爪.
寬十爪.
很大.
真的好像一個.
現在的無人機.
他對我說.
這是發出行在遍地上的咒詛.
凡是偷竊.
必按卷上這面的話除滅.
凡起假誓.
必按卷上那面的話除滅.
即是說.
譬如說.
書卷裡面一定有字.
不是無字天書.
不是那些.
九音真經.
不是那些乾坤大羅.
而是那些繁文.
說的書卷對猶太人來說.
一定是舊約的摩西五經.
或者舊約的重要的律法書.
等等.
甚至說到的就是.
凡是干犯這一樣.
很明顯是關乎公義.
這次不是說.
屬神的人的公義.

$^{1761}$是說神的公義.
是對外圍的人的審判.
即是說.
外面的世界一直充斥著罪惡.
誰是標準?.
一定是聖經的真理是標準.
這個標準就是可以去審判其他的人.
這就告訴我們.
是發出遍行傳遞的.
在遍地上發出的是什麼呢?.
這個咒詛.
上文一到第二章.
頭三個的二章.
都是Blessing.
都是對耶路撒冷的安慰.
有沒有記得?.
這個城就是我眼中的銅人.
我就是那個火牆.
火城,不好意思.
來保護這個民族.
這個聖城.
我要來成為穩固的石.
諸如此類.
接著還有大祭司.
是可以公職.
因為他穿乾淨衣服.
和戴冠鍊.
所羅巴伯有聖靈.
他就有力量.
可以來為這個聖城的最重要工程.
一直建造.
但是現在.
二章就對外了.
真真正正神的審判.
要發出對外.
整個書卷.
說的是咒詛.
所以這段經文.
我們比較少看.
但是你可以說.

$^{1801}$如果我們說人就人.
不信耶穌落地獄.
某程度是真的.
當然我們不是看到人就說.
你要落地獄.
那又誇張了.
因為每個人的處境都不同.
但是無論如何.
我們就看到.
這裡特別說的二章.
這次才看到的.
是一個毫無恩典的審判.
因為你看到的.
就是發出走在遍地上的咒詛.
今天為什麼還是有恩典呢?.
因為主耶穌還沒有回來.
宣道會.
我們說四重福音.
浸禮巴一定要上.
四重福音最重要的主再來.
就是那個皇冠.
主再來的含義.
不是聖徒被提這麼簡單.
也可以暗示著.
是不信的人.
要面對真正的咒詛審判.
這是雙重的意義.
所以我們尋求主再來.
其中信徒要做的.
就是宣教.
要告訴未得之民.
主耶穌很快要回來.
所以這是很重要的.
雙重的意義.
所以現在這個書卷.
在撒迦利亞看到.
其實很恐怖.
因為書卷上所記載的.
例如.
偷竊.

$^{1841}$按卷上的說話.
除滅.
起假釋的.
按卷上的說話除滅.
沒有恩典可言.
這就是對將來的世界的審判的預言.
才看到這一點.
第四節.
萬君之耶和華說.
我必使道這書卷出去.
進入偷竊人的家.
和指著我明.
起假釋的人的家.
指著我明就是猶太人.
即是他們一直都解決不了.
他們都解決不了.
外邦人就更加不用解決.
他們都是信神.
說到明不可以起假釋.
他們都是這樣的時候.
何況其他人呢.
這正是告訴我們.
這三個意象.
第六第七個意象.
都是說到罪惡要被審判.
和怎樣可以真正知識罪惡.
就是第七個意象.
那個婦人.
在一個良器裡面.
帶到無人之地.
是一個很特別的意象.
所以第六個意象是審判.
第七個意象是指著罪惡.
即是說彌賽亞的再回來.
又是這兩樣事情的發生.
你想一下將來我們去到永恆.
就不會再犯罪.
亦都不會再有罪.
你可以想像得到那個美麗.
那個永恆.

$^{1881}$是多麼值得我們去羨慕.
地上的幾十年很快就會過去.
進入的就真的再沒有罪惡.
和永遠的公義和良善.
在那個新天新地裡面.
這告訴我們.
這個書卷要出去進入這些人的家.
必常在他家裡.
你可以想像得到.
家有一老餘有一寶.
或者雖然我們中國人都說.
家家有一本難唸的經.
不夠這本難唸.
必常在他家裡.
連房屋大木石都毀滅.
即是說不信神的人.
不敬畏神的人.
面對的就是再沒有恩典的審判.
用一些這麼形象化.
這個書卷會飛入他家裡.
誰偷東西飛去你家裡.
然後就停留在你家裡.
必常在他家裡.
其實我們讀聖經.
就是可以令到我們不再犯罪.
真的有這個潔淨作用.
讀聖經是真的有這個潔淨作用.
所以一個地方幾暴戾的人.
如果讓他接觸到聖經.
他又真誠地讀.
他的生命會改變.
所以改變我們不犯罪.
是聖經的奇妙能力.
所以大家來聽查經班.
已經是進入一個聖經的世界.
我們有一個意識和神有一個相交.
所以你看多恐怖.
你說到處飛當然沒有問題.
不關我的事.
我家裡是重門深鎖的.

$^{1921}$我們沒有犯罪.
不是的,這個書卷好像有導航.
誰偷竊就飛去他家裡.
誰起家世就飛去他家裡.
這個常在的書.
不是難唸.
是你知道罪名.
然後連房屋大木石都毀滅.
你可以想像到.
不是抄家.
抄家都還有一間屋.
是連房屋大木石都被毀滅.
即住在這裡的人更加不用說.
這個就是飛書卷的異象.
即是神現在開始眷顧.
眷顧這個visit可以是審判.
可以是私恩.
所以我們不要單方面永遠都是恩典.
到第七個異象.
與我說話的天使出來對我說.
見所出來的是甚麼.
我說這是甚麼來的.
他說.
所出來的是一個良器.
即是一個可能是一些罐或元柱的東西.
他又說.
這是惡人在遍地的形狀.
所以是一對的.
這兩個異象是一對的.
在遍地的形狀就是圓的.
我見有一片圓的圓被舉起來.
在良器中的是一個婦人.
啟稟也有一個婦人.
稱她為大淫婦.
啟稟的十七章.
我應該沒有記錯.
我們看過了.
當然因為罪惡的容器是陰性.
所以會有這個婦人希伯來文來說.
當然也有一個婦人在這個異象中.

$^{1961}$第八節.
天使說這是罪惡.
他就把婦人拎在良器中.
良器本身已經是罪惡.
罪上加罪.
良器本身是罪惡.
把一個婦人是罪惡的.
放在良器中.
即是一擔擔.
那怎麼辦呢.
留意一下.
重點反而不只是描述她.
而是要拎到去視那地.
告訴我們.
將那片圓的圓拎在良器的口上.
見有兩個婦人出來.
她們翅膀中有風.
飛得甚快.
似乎是神的使者.
翅膀如同灌料的翅膀.
她們將良器抬起來.
圓在天地中間.
我問與我話的天使說.
所以有對話.
她們要將良器抬到那裡去呢.
因為她都見到.
兩個有翅膀的.
即是應該是天使.
是婦人的形狀.
就拎起了良器.
力大無窮.
即是在靈界.
不會像我們物質那樣想重量.
好像很輕盈地拎起.
就飛到那裡去呢.
先知就好奇了.
飛去那裡呢.
就問那個天使.
所以我們靈修有時都要與神對話.
不是有思覺失調.

$^{2001}$是真的祈禱與神對話.
將你讀到的聖經.
告訴神.
希望如果是要求的.
我們就去進行.
如果是安慰的.
我們就去領受.
這樣對話的.
所以靈修是對話式的.
與神傾談的.
所以你見到天使.
與撒迦利亞有話題的.
他真的問的.
所以他就知道.
如果不問就不知道.
飛了走就飛走了.
他就問.
原來他就答了.
要往士拿地去.
為他蓋造房屋.
等房屋齊備.
就把他安置在自己的地方.
士拿地是他自己的地方.
應該要去的地方.
似乎那個含義.
我們未必需要了解.
為何一定要去士拿.
這些就不應該是好奇的問題.
而是他自己應該去的地方.
剛才一路上文很清楚.
那個容器是罪惡.
那個婦人是罪惡.
要帶到士拿地.
即是遠離神.
遠離耶路撒冷.
當然是以耶路撒冷做那個方位.
遠離神的地方.
去到士拿地.
在那裡屬於他自己的地方.
記不記得.

$^{2041}$即是說罪惡再不會再有.
這時間.
其實就是十七章到十九章的啟示錄.
就是講這件事.
當魔鬼最終被拎入火壺.
第十九章首先就是.
那個獸和那個迪基督.
地獸和海獸被拎入火壺之後.
最終這個邪惡的鐵三角.
邪惡的三惟一體.
撒旦和那個海獸和那個地獸.
賈先知和那個迪基督.
三個邪惡的三惟一體.
就被拎入永恆的火壺.
所以就真的意味著.
罪惡真正的止息.
所以今生我們就要為這些事情祈禱.
求神叫我們遠離罪惡.
將來就不用了.
因為罪惡這個意象告訴我們.
將來這個罪惡就有兩個使者.
拎他去到士拿地.
回到他自己的地方.
即是回家鄉.
他們就是應該去的.
即是神是公然.
所以在天堂就不會再有罪惡了.
我們就進入永恆的天家.
當然我們的生命會改變.
我們不會再有帶著我們的Greedy.
我們的貪婪.
我們的嫉妒.
我們的憤怒.
那你怎樣進入天堂呢.
應該不需要了.
所以進入永恆.
我們一定是有一個新的生命.
所以這裡告訴我們兩個意象.
一個就是飛行的書卷審判罪惡.
所以都盼望我們都要在神的面前.

$^{2081}$多一些的安靜.
反省.
另外就是這個良氣的婦人.
給我們感受到.
是會有真正的神的審判時間.
罪惡再不會來到去猖狂的.
和大家揭一揭.
最後揭這個啟示錄17章.
溫習這些經文就會OK了.
17章的啟示錄.
拿著七個碗是憤怒的碗.
七位的天使.
有一位前來對我說.
你到這裡來.
我將坐在縱水上的大淫婦.
即是泉地.
所要受的刑罰只給你看.
地上的君王與他行淫.
可能是宗教上.
他們都不敬拜真神.
住在地上的人喝醉了他淫亂的酒.
我被聖靈感動.
天使帶我去到曠野.
我就看見一個女人.
騎在朱紅色的獸上面.
這個獸有七頭十角.
遍體有褻瀆的名號.
遍體不是七眼.
是褻瀆神的名字.
遍體.
從哪個角度.
她都是咒罵神.
很恐怖的怪物.
神的七眼就是七靈.
遍察全地.
甚麼都看見.
那是安慰對信徒來說.
但這個獸.
騎著的這樣褻瀆的獸.
遍體都是充滿了褻瀆神的名號.

$^{2121}$即是你想得出.
想不到的.
不可以用文字的.
它都有.
這樣來褻瀆神.
所以基督徒最怕.
一個有聖靈的基督徒.
最怕就是褻瀆神.
千萬不要.
不可以當神來玩.
很嚴肅的.
真的不可以.
有時連不信的人.
都知道不可以玩你們的神.
即是不可以說笑.
如果基督徒自己拿神來說笑.
就最佳一等.
這告訴我們第四節.
哪個女人是穿著紫色和珠紅色衣服.
所以表面看不到她是這麼.
敵黨神.
得罪神.
穿著的是紫色.
珠紅色.
最漂亮的.
那個時代的.
最華麗的衣服.
但她的內心.
不是神喜悅.
用金紙寶石珍珠為裝飾.
手裡拿著的是金杯.
杯中盛滿了可憎之物.
就是她淫亂的污穢.
在她額上有明寫著說.
奧被災大巴比倫.
作世上的淫婦.
和一切可憎之物的母.
所以是那個Mother.
Mother of Crime.
Mother of 邪惡的意思.

$^{2161}$就是敵黨神的.
Adversary 的 Mother.
一切敵黨的之母.
所以真是描繪那個罪惡的本質.
所以她騎著的那隻怪獸.
不覺得是怪.
是同類.
騎著的是怪獸.
遍體都是褻瀆神的名號.
你可以想像得多恐怖.
這個世界對神的譴責.
接著.
我又看見這個女人.
喝醉了聖徒的血.
和為耶穌作見證之人的血.
以敵黨神的僕人兒女為樂為榮.
為她最大的滿足.
有些人真是很喜歡.
去挑戰.
嘲笑基督徒.
這些人其實根本不需要太多對話.
因為他們是為反對.
為喜樂而喜樂.
不是真正想尋求.
是想嘲笑你.
說到你不夠他說.
這樣就算了.
我看見她就大大的稀奇.
天使對我說.
你為什麼稀奇.
將這個女人和騎著她.
七頭十角的獸的奧秘告訴你.
你所看見的獸.
先前有如今沒有.
將要從無底坑裡面上來.
無底坑是靈界還是怎樣呢.
似乎都是靈界.
又要歸於沉淪.
凡是住在地上.
明治從創世以來.

$^{2201}$沒有記載生命冊上的.
見先前有如今沒有的.
之後再有的獸.
就必稀奇了.
智慧之心就可以在這裡思想.
一直在描述一座城.
是一個七山的城.
然後就說到那些王等等.
所以你可以想像到這個淫婦.
我不說這麼多了.
十五節.
天使又對我說.
你所看見那淫婦坐的重水.
就是多民多人多國多方.
即是全世界.
你所看見的那十個獸.
必定會恨這個淫婦.
使她被冷落.
又要吃她的肉.
用火將她燒盡.
因為神使到諸王同心合意.
遵行神的旨意.
把自己的國級獸.
直等到神的話都應驗.
你所看見那個女人.
其實就是管轄地上眾王的大城.
將來有一個大城.
有這個女人的代表.
所謂大淫婦.
即是對真理徹底的抵擋.
這樣來招聚那些人.
然後還有一個在地裡面出來的怪獸.
有十隻角七個頭的.
很恐怖.
都不夠恐怖.
她的身體遍體都是充滿了.
不是眼睛.
是褻瀆神的名字.
那些名號.
我都不敢隨便說.

$^{2241}$諸如此類.
你可以想像得到.
這樣的裝備.
這樣的佈局.
告訴我們將來在地上會有.
不知道什麼時候.
是會有這個日子.
神的榮耀好像暫時在地上完全消失.
因為他們得勝.
又殺基督徒.
但是無論如何.
當神的時間到.
那個地獸是由地無底坑上來.
但始終歸於沉淪.
大淫婦.
神想起罪的時候.
同樣都會被審判激怒.
所以這個就是我們需要.
透過今天第五.
對不起第六和第七個議帳.
關於那個罪惡的審判.
和最終神要將罪惡指釋.
成為我們大家的留心.
但是第四章.
我們應該很熟悉的第六節.
不是依靠勢力.
不是依靠才能.
其實是要.
不過不是單單依靠勢力.
不是單單依靠才能.
是要看到.
是要那個當事人看到.
是真的需要依靠神的靈.
方能夠成事.
我們一起禱告.
多謝主賜給我們又有一個半小時.
我們大家可以很安靜.
很靜心的去聆聽.
第五第六的議帳.
我們很盼望神給我們能夠繼續.

$^{2281}$透過這八個議帳.
我們可以去認識神的心意.
你對耶路撒冷的愛.
你重新看顧眷戀這個民族.
但是你對於其他敵黨神和真理的.
你是會有飛行的書卷.
去到他們的家中去審判.
更加這個良氣的婦人.
是會被拐到她自己應該要的地方.
就是這個士拿地.
很盼望神都給我們看到.
我們去追求聖潔.
非聖潔是沒有人可以見主的.
願我們能夠在你的面前.
蒙你的恩典.
多謝主聽我們的感恩禱告.
奉耶穌基督的名.
阿們.
\newpage



\section{馬可福音 1:40-44}
\label{sec:G0FSq209Kac}
\textbf{2023-10-07 晚堂《醫治為要得潔淨》馬可福音 1\_40-44 講員:黃國維院長}
\newline
\newline
連結: \href{https://youtube.com/watch?v=G0FSq209Kac}{\texttt{https://youtube.com/watch?v=G0FSq209Kac}} ~~~~ 語音日期: 2023-10-11
\newline
\newline
\hyperref[sec:CXJyuYm4_pE]{\small{< < < PREV SERMON < < <}}
~
\hyperref[sec:index_chronic]{\small{[返順時目]}}
~
\hyperref[sec:index_scriptual]{\small{[返順卷目]}}
~
\hyperref[sec:iO1qZX8Cmq4]{\small{> > > NEXT SERMON > > >}}
\newline
\newline
馬可福音 1:40-44
\newline
\begin{longtable}{cl}
\hline
\hline
章節 & 經文 (和合本修訂版)\\
\hline
1:40 & \begin{tabularx}{0.7\textwidth}{X} 有一個痲瘋病人來求耶穌，向他跪下，說：「你若肯，你能使我潔淨。」 \end{tabularx} \\ \\ \relax
1:41 & \begin{tabularx}{0.7\textwidth}{X} 耶穌動了慈心，就伸手摸他，說：「我肯，你潔淨了吧！」 \end{tabularx} \\ \\ \relax
1:42 & \begin{tabularx}{0.7\textwidth}{X} 痲瘋病立刻離開他，他就潔淨了。 \end{tabularx} \\ \\ \relax
1:43 & \begin{tabularx}{0.7\textwidth}{X} 耶穌嚴嚴地叮囑他，立刻打發他走， \end{tabularx} \\ \\ \relax
1:44 & \begin{tabularx}{0.7\textwidth}{X} 對他說：「你要注意，千萬不可告訴任何人，只要去，讓祭司為你檢查，又因為你已經潔淨，獻上摩西所吩咐的祭物，作為證據給眾人看。」 \end{tabularx} \\ \\
[1ex]
\hline
\hline
\end{longtable}
$^{1}$各位弟兄姊妹晚安.
我們一起祈禱.
將現在的時間交託.
你曾經來到我們當中.
工作,醫治,講道.
讓我們明白.
天父是一位怎樣的上帝.
今日求你活現.
在我們當中,透過你自己的說話.
讓我們能夠認識你更多.
更新我們的生命.
好讓我們懂得怎樣去跟從你.
祈禱耶穌基督的名祈求.
阿們.
耶穌來到我們當中的時候.
也做了很多事情.
行神蹟,醫病,講鬼.
在聖經裡面.
當我們認真去看聖經記載.
不同的神蹟.
其實不只是說.
耶穌醫好我們.
或者耶穌很大能.
其實每一個的記載.
都有自己一些更深的心意.
剛才主席幫你讀了.
《馬學福音》第一章的經文.
幾節的經文.
是在說耶穌.
遇見一個大麻瘋的病人.
不過當我們細心去看的時候.
發覺原來耶穌.
不只是醫了他身體這麼簡單.
也都當我們細心看.
會發覺這段經文.
向我們一班.
經歷過幾年疫情.
這樣的狀況的信徒.
都有很多的提醒.
不過你要了解.

$^{41}$耶穌遇見這個麻瘋病人.
其實做了什麼.
我們就會明白.
原來在舊約裡面.
以色列人已經有很多.
關於麻瘋病.
這個傳染病的條例.
我們明白了那個背景.
我們才知道耶穌.
其實做了什麼.
在這樣的狀況裡.
我和大家看了一點經文.
在曠野以色列人.
領受律法的時候.
利美金已經有提到.
一些如何去到.
當以色列民裡面.
發覺有麻瘋病.
要如何去處理.
利美金第十三章就說.
耶和華就和.
摩西亞倫就說.
如果你們發覺.
在你們當中.
有些人.
肉裡面生了節子.
長了扇子之類.
你們就要去.
把節子長成扇子之類.
發覺有些跡象.
如果要去到.
覺得他去患麻瘋.
第二節後面就說.
就要把這個人帶到.
祭司亞倫.
或者亞倫作祭司一個子孫.
的面前.
做什麼好呢 第三節.
祭司要看一下他.
皮上的災病.

$^{81}$教他怎麼看的 聖經有教.
如果災病那裡的毛.
變白.
現象在心的皮上.
肉上的皮之類.
原來那個人就.
患了大麻瘋.
我們都會明白.
麻瘋是一個傳染病.
尤其是在以前的.
世代裡面.
沒有養醫.
他們要有一個方法去辨認.
今天其實對我們來說.
社會裡面都有傳染病.
過往幾年的疫情.
過去之後.
現在都很多流感.
這些走不掉的東西.
今天我們如果有這些病.
那怎麼知道.
去看一下醫生 做一下檢測.
過往的時代沒有.
那怎麼好呢.
以色列人就.
他們沒有醫生 沒有醫療系統.
那怎麼做好呢.
在祭司的傳統裡面.
發生這件事.
為什麼用祭司的傳統呢.
因為如果你要將這些.
傳染病控制得.
有效.
是需要有一個.
大家都聽從的方法.
如果不是的話.
是會控制不到.
而在以色列的人裡面.
最有效的方法.
就是那個祭司的.

$^{121}$傳統 所以那時候.
上帝就說.
當你懷疑有人有麻瘋病的時候.
你就要帶到祭司那裡.
祭司就要決定.
究竟他是不是真的.
患了大麻瘋.
如果他發覺真的有大麻瘋.
所謂叫做確診了.
那怎麼好呢.
同樣地.
利未記第十三章比較後.
那時候就教怎麼做.
就說新上長有大麻瘋.
災病的 他要怎樣.
衣服要撕裂 蓬頭散髮.
蒙著上唇.
要大聲叫不潔淨.
不潔淨.
很明顯為什麼要做這些.
因為老遠就被人知道你有病.
老遠知道你有傳染病.
其他見到你的人.
就要怎樣 要避開.
不要受傳染.
四十六節就說這個狀況.
災病在他身上的日子.
他就是不潔淨.
他既是不潔淨.
就要獨居.
型外.
那種不要傳染人.
所謂今日的說話就是隔離.
這樣的狀況.
因為沒得醫.
不過我們留意有一個字眼.
叫做不潔淨.
今日我們說叫確診.
不過原來.
在舊約裡面.

$^{161}$就叫做不潔淨.
這樣的東西.
他就不能夠去到.
跟其他人一起相處.
我們就.
看到類似經文裡面.
舊約的律法.
不單止是麻風病.
文數第五章.
那裡都有說其他的病.
例如第二節.
麻風病患流症.
和因屍體.
不潔淨都要送到型外.
所以我們發覺.
原來不只是麻風有其他.
他們那時候.
恐怕會傳染的.
狀況都一律.
要送到型外.
第三節就無論男女.
都要送 將他們送到型外.
後面那句很厲害.
免得他們.
玷污了他們的影.
這是我.
是誰 耶和華.
住在他們中間的地方.
這個說話其實對以色列人.
來說是很有效的.
他知道自己說麻風病.
是所謂不潔淨.
而不潔淨就不能夠.
玷污那個型.
那個型是誰在裡面的.
耶和華.
所以對他們來說很有效.
因為他覺得自己不能夠玷污這位聖潔的上帝.
他們就感受這些.
律法.

$^{201}$是獨居型外.
不過同時我們可以.
理解他們身體有病.
已經很多人拒絕.
很多不舒服.
也都覺得.
原來在上帝的面前.
都不潔淨.
聖經來到這裡我們覺得很慘.
不過舊約又不是很絕情.
有沒有機會.
好得回.
又有.
我們繼續看下去.
第十四章就有說.
萬一那個人.
自己在外面住.
真的好得回那怎麼辦.
第十四章就討論這件事.
這個是萬麻風病.
的人得到潔淨.
要做什麼.
既然確診找祭司.
好得回都是找祭司.
祭司要怎樣.
先不要叫人回來.
出去檢查他.
看看他的皮膚.
好得回.
第四節就要.
帶一些祭物去獻祭.
這裡彷彿跟我們說.
其實醫治都是在上帝的手裡.
當你好得回.
你要獻祭.
去到感恩.
不過除了獻祭這些宗教禮儀之外.
在第四節.
有說一些.
要做的事.

$^{241}$要洗衣服.
剃去所有的毛髮用水洗澡.
他就潔淨.
很明顯這個效用.
不只是宗教禮儀.
這麼簡單.
他們可能憑經驗知道.
都洗過澡洗過衣服.
都是好.
很容易.
開個洗衣機.
不過大家知道.
那時候以色列人在曠野.
要洗衣服洗澡.
都不是很簡單.
不過最後一句很特別.
當他洗完澡洗完衣服之後.
他就可以回營.
不過營要在.
自己的帳篷之外.
住七天.
他都恐怕那些傳染病.
會帶回自己的帳篷.
不要進去.
回營.
住多七天.
七這個數字真的很特別.
因為經過幾千年.
今日都是用七天去做隔離.
聖經裡面都很有智慧.
就算今日都好.
經過幾千年.
我們仍然有些病是治不到.
仍然有這些方法去隔離.
好了.
洗完澡洗完衣服.
終於搞掂.
到第八天.
這個病的人.
就要拿兩隻沒有殘疾的小公羊.

$^{281}$其他的祭物.
要去獻祭.
但十一節提到什麼.
原來有一個叫做.
宣佈潔淨的祭司.
他要找祭司.
祭司看到他.
做完所有的程序之後.
接納了祭物之後.
要宣佈.
那個人潔淨.
然後連同祭物.
安置在耶和華面前.
會幕的門口.
為什麼.
因為會幕的門口.
是全部人都看到的地方.
大家記得當一個人有病.
他就要.
縫頭散髮,撕裂衣服.
大聲叫,所有人都要知道你病.
但他痊癒了.
那時候同樣.
要和所有人.
宣佈他沒事.
為什麼.
因為他要回復一個群體.
正常的生活.
我們看看在舊約.
這個制度裡面.
我們發覺原來.
我想在人類很多文化.
都是的,遇見一些.
不懂得醫治的傳染病.
好像蜜蜂一樣,必須要去.
不要讓那個病.
去到散開.
不然很大件事.
有很多是分開的方法.
如斯利文.

$^{321}$用了祭司的制度.
不過我們會發覺.
在舊約的律法裡面.
其實經歷這些.
患麻風要被.
去獨居營外不潔症的人.
其實都很慘.
都有說到身體病已經不舒服.
不舒服很需要人照顧.
但他們不能夠被照顧.
要一個人.
還有他也感覺到.
我的不潔症.
好像上帝也不要我一樣.
都是一件很悲慘的事情.
不過在舊約的律法.
都留有一些.
恩典,為甚麼?.
因為讓你可以痊癒,可以回頭.
其實.
想清楚一點,是否一定可以讓你回頭.
其實都未必.
既然你患了這麼.
嚴重的傳染病.
但可以趕你走.
以後不要回來.
自生自滅都可以.
但無論如何.
這些災病被人帶來.
很多痛苦.
很多困難.
在這樣的背景下.
我們會明白.
究竟耶穌.
當他遇見一個.
帶麻風的病人,究竟他做了甚麼.
這樣就開始記載.
在.
《馬可福音》第一章.
40節.

$^{361}$我們一起讀.
這三節.
一二三.
有一個麻風病人.
來求耶穌.
向他跪下說.
你若肯,你能使我.
潔淨.
耶穌動了慈心.
就伸手摸他.
說:我肯,你潔淨了.
麻風病.
立刻離開他.
他就潔淨了.
我想讓大家看看.
這個麻風病人.
見到耶穌跟他說了一句話.
叫做:你若肯.
你能使我潔淨.
這句話其實問得很奇怪.
當我們有些.
很難醫的病.
去找醫生.
你會問個甚麼問題呢.
那個醫生.
行不行的.
他有沒有能力.
他有沒有所謂的track record.
他是不是一個好人.
我們會問能力.
行不行的問題.
但麻風病人不是問.
耶穌你行不行.
是問你肯不肯.
為甚麼呢.
其實你想想.
我們墮入麻風病人裡面.
他對於耶穌有能力.
醫他,他完全沒有懷疑.
可能他聽到耶穌的大能.

$^{401}$但是.
他知道這位能夠醫治.
麻風病人一定是一個.
從很聖潔的上帝.
那裡來的一個人.
而當他患麻風的過程裡面.
其實他覺得自己不潔淨.
所有人.
見到他都避開.
我一個這樣不潔淨的人.
遇見一個這麼厲害.
從上帝而來的耶穌.
我走過去.
都不敢.
所以這句.
你若肯這句話.
是反映在他心裡面.
他對自己形象的低落.
他覺得其實.
沒有人肯理我.
這個從神而來的耶穌.
你肯不肯呢.
所以他問你若肯.
你能否使我潔淨.
他不是問能夠的問題.
耶穌知道.
耶穌聽到,所以你見到耶穌對他.
答案是甚麼,耶穌後面就.
第一句跟他說的話.
用回他的字.
我肯.
我不單止有能力,我是很願意.
肯這個字.
廣東話聽起來就像被人求.
求你肯不肯.
求你,你肯.
但是在原文裡面.
這個肯的字是很正面的.
是在說他的渴望.
英文甚至叫.

$^{441}$I desire,我很想去.
依著你,耶穌說.
我不是只給你求,我肯.
而是我見到你的時候.
令你得到意志,是我的渴望.
為甚麼耶穌會這樣呢.
因為聖經記載.
當耶穌見到這個.
盲風病人的時候,他第一個反應.
是叫做動了.
慈心.
動了慈心是甚麼意思呢.
在新漢語譯本裡面.
是這樣翻譯,是憐憫的心.
憐憫這個字.
在聖經裡面.
都多出現,通常是形容.
上帝對於罪人.
或者對於病人.
的一個狀況.
憐憫.
憐憫這個字的意思是甚麼呢.
其實憐憫的意思.
就是當.
上帝遇見.
一些本來應該是.
很尊貴.
是神兒女.
有神的形象.
尊貴的樣式的人.
在不幸的情況下.
落入一些很悲慘的處境裡面.
的反應.
即是.
如果大家要分辨清楚一點.
憐憫這個字.
和很可憐.
有一點點不同.
很可憐覺得你.
搞到這樣,很骯髒.

$^{481}$我這麼有能力.
我就幫你一下.
平常我們覺得很可憐的人.
好像很低等.
我幫你一下.
不然我就很厲害.
但在聖經裡面往往強調.
憐憫這個字不是這樣.
原來.
一個人表面有病.
有罪.
很老到.
對人來說.
我們會覺得他很骯髒.
我們很害怕.
但耶穌往往.
是看穿那個人的表面.
看到每個人.
其實是很尊貴.
但是一個這麼尊貴的人.
竟然被病.
被罪.
去到折磨不像一個人.
他心裡面.
犯出那種感覺.
就是一種憐憫的感覺.
這裡給我們很多提醒.
其實我們人很自然.
都是看到人的表面.
或者災病裡面.
或者傳染病裡面.
我們有時候都會很害怕.
那些有病的人.
或者我們看到.
那些人病了很久.
樣子都已經很難接受.
害怕他.
不過原來耶穌看人.
不是看表面.
他是看到表面背後.

$^{521}$那份尊貴.
這裡給我們很多提醒.
有時候在疾病傳染病裡面.
我們都會覺得.
無論病的人.
或者很可憐.
或者有些害怕.
或者甚至有些.
傳染病的人知道自己病.
會很怕.
連累到其他人.
過去幾年可能都有這樣的經驗.
吃完飯之後.
他病了他很害怕.
很擔心自己會害到人.
第一節不是說.
我們要掉以輕心.
對於傳染病.
不過其實我們有沒有看到.
那個狀況就是.
我們看到病的人去吃飯.
不是首先去怪責他連累自己.
而是看到.
其實那個人自己本身.
都是不應該.
被病去纏繞.
反而多了一份.
所謂主耶穌基督.
一種憐憫的心腸.
所以.
其實我們看到原來耶穌看人.
和我們真的很不同.
他會看到人的尊貴.
產生那份憐憫的心.
而有所行動.
很渴望.
去醫治他.
接著耶穌做了一個.
很特別的行動.
耶穌動了慈心.

$^{561}$未說話之前.
就伸手摸他.
伸手摸他.
對於耶穌去醫治人來說.
是一個很多餘的動作.
在其他.
醫治的記載裡面.
耶穌只是想去碰人.
甚至有些人隔空.
去都不用去就能醫治.
所以很明顯耶穌不是因為摸他.
而醫治他的麻瘋病.
還有.
在路加福音.
同樣的記載裡面.
對於這個.
麻瘋病人多了一點.
的資料.
在路加福音.
五章十二節.
是說這個人滿身長了麻瘋.
滿身長了麻瘋.
滿身長了麻瘋的意思是什麼.
就是不是剛剛.
昨天病了.
你知道皮膚病是要慢慢來的.
慢慢發散.
如果大家有經驗.
生過濕疹就知道.
你不會突然一天起來全身都是.
慢慢來.
所以我們知道這個病人.
其實是病了很久.
有麻瘋病人.
即是同時說他被隔離了很久.
又或者說.
原來他被隔離.
不要說被人碰.
走近一點.
都沒有人敢.

$^{601}$所以耶穌伸手摸他.
這個動作.
是給他一個很大的安慰.
和接納.
一個人被隔離.
幾天都這麼辛苦.
今天我們都可以.
講一下電話.
FaceTime 一下.
聊聊天.
這個年代是沒有的.
不要說被人捉摸.
我們開始看到.
耶穌在做什麼.
耶穌去醫治.
不只是說你的身體有病.
我醫治.
我醫好你就行了.
他發覺我們人在疾病裡面.
可能最大的障礙.
最大的困難.
未必是身體的.
健不健康.
而是當有病.
尤其是傳染病的時候.
他落入一個.
和其他人要去到分隔.
很覺得被拒絕.
很覺得自己.
污穢的那種感受.
其實耶穌.
伸手摸他.
醫治他.
心靈裡面.
那種渴望.
還有.
其實對耶穌來說.
摸他也是冒險的行動.
在傳染病裡面大家都知道.
走得很近.

$^{641}$伸手摸.
自己都危險的.
不過.
我們會很快覺得.
耶穌.
他醫人都可以了.
自己不會病的.
好啊 耶穌不會病.
不過大家知道在舊約的條例裡面.
你不病.
你不潔淨的人自己都很麻煩.
都是的.
在今天.
其實都會是.
傳染病裡面.
如果我們要對人那種憐憫.
那種接納關懷.
可能都會讓我們.
落入某種的冒險.
可能挑戰的是.
我們願不願意.
沒有這個險去到.
下一發耶穌呢.
好.
接著耶穌動了慈心.
憐憫的心.
伸手摸他.
說了我肯你潔淨了.
實際的記載就是.
麻風病立刻離開他.
一個很特別的記載.
可以說他的病好了.
痊癒就行了.
為什麼要用.
立刻離開他.
去形容呢.
我想對於那個.
病人來說.
在他患病的.
經歷裡面.

$^{681}$是其他人離開他.
上帝離開他.
健康離開他.
不過聖經記載.
當他遇上耶穌的時候.
是麻風病.
離開他.
我想這句話彷彿.
都提醒我們.
聖經裡面有說.
疾病,孤單,刀劍.
都不能夠將我們和.
上帝的愛去到.
隔絕.
很形象化地去到描繪.
耶穌.
我們認識耶穌.
在生命裡面離開我們.
最終.
都是疾病.
那些孤單.
這裡來到第42節.
麻風病.
立刻離開他.
他就潔淨了.
來到這裡其實故事.
已經很圓滿了.
病又醫好了.
耶穌又接納他.
心靈都醫治了.
不過聖經.
繼續說多兩節.
才完結.
當然有那個意思.
還有43.
和44兩節.
不如我們一起讀這兩節.
耶穌嚴嚴地叮囑他.
立刻打發他走.
對他說.

$^{721}$你要注意.
千萬不可告訴任何人.
只要去.
讓祭司為你檢查.
由因為你已經潔淨.
獻上摩西所吩咐的祭物.
作為證據.
給眾人學.
耶穌醫好他.
吩咐他的說話.
44節跟他說.
你要注意千萬不可告訴任何人.
這句就和麻風病人.
關係不大.
這句是因為.
耶穌初時出來工作的時候.
不想這麼快張揚.
是這樣的背景.
所以說.
後面這句是和麻風病人說的.
只要去讓祭司為你檢查.
我們讀了.
你未記就明白.
那個是以色列人.
他就說.
那個是以色列人.
一個所謂社會上.
去處理.
傳染病的.
制度.
這個真的很有意思.
耶穌怎樣說.
原來耶穌是肯定.
社會一些的.
防止傳染病的制度.
你見不見到.
耶穌說.
這個社會上一天.
還有這些傳染病.
社會上去面對.

$^{761}$這些傳染病的制度.
其實是需要存在的.
他沒有否定.
不過耶穌來.
他不是只是去.
跟著那個制度去工作.
而是在.
接觸人的裡面.
是.
關懷,接納,理解.
我想這裡.
給我們一些提醒.
是甚麼呢.
有時我們做信徒.
或者教會群體.
今天對這個社會.
有沒有傳染病.
有時我們都無能為力.
我們做得甚麼.
如果等於你是醫護.
可能多些事.
做到.
但一般群體,教會群體.
彷彿覺得.
在疫情裡面.
我們唯有跟著.
那些制度.
去做事.
不過我們發覺.
耶穌做的事提醒我們.
社會上的制度.
當然有.
它的作用.
但是我們.
跟從耶穌的人.
原來在這些困難的環境裡.
就應該做的事.
就是耶穌那種.
的憐憫.
那種去到人的身邊.

$^{801}$去關懷.
去冒險.
即使有很多困難.
我們都願意去.
跟從耶穌.
效法耶穌所做的事.
在這裡給我們另一種的看法.
究竟.
耶穌的門徒.
可以做甚麼.
接著下去.
最後一句.
已經潔淨了.
獻上摩西豐富的祭物.
作為證據吸眾人看.
你得了潔淨.
你還需要回到社群裡.
過你的正常.
群體生活.
這句說話.
是告訴你.
原來疾病.
不只是身體的問題.
是更加.
跟群體有關.
我想我們很深刻地體會.
過去的幾年.
其實表面.
是一種病毒.
表面是很多東西.
要停下來.
但其實最深刻影響的.
是群體.
很多很多.
我們過往.
能夠做到的事.
能夠怎樣去建立群體.
原來.
是被這些疾病.
去到影響了.

$^{841}$我們知道.
原來教會裡面是一些.
很學習.
真誠去彼此相愛的群體.
我們發覺.
原來在災病裡面.
基督徒最大的.
見證.
就是我們願意活出一個.
怎樣的群體.
來到這裡.
可能要做一點總結.
我為了這段經文.
定了一個題目.
叫做「醫治.
為了得潔淨」.
醫治是甚麼?.
身體的健康.
但如果大家記得.
潔淨或不潔淨.
在舊約裡面是在說.
要獨居營外.
是在說.
要和群體.
去到隔絕.
是在說被拒絕.
這個問題.
所以身體上的健康.
往往.
是會影響.
他們的社群裡面.
怎樣去生活.
我們發覺.
耶穌是很渴望.
醫治人的.
見到麻瘋病人.
是醫好他們的身體.
不過原來.
有健康不是停在健康.
健康的目標.

$^{881}$是要讓我們.
在有一個群體裡面.
有彼此.
相愛的機會.
的機會.
原來我們.
經常渴望得到的健康.
的目的不是只是健康.
是我們怎樣.
運用我們健康的身體.
在群體裡面.
彼此相愛.
在這樣的狀況裡面.
我們發覺有些反省.
第一個反省是.
健不健康都好.
我們都要彼此相愛.
耶穌給我們最大的.
誡命是甚麼.
愛神愛人.
最大的誡命.
就是愛上帝愛人.
甚麼叫誡命.
誡命就是.
你甚麼狀況下都要做.
就是誡命.
不過你有沒有想過.
為何耶穌給我們的誡命.
不是要我們健健康康的生活.
因為如果耶穌給我們的誡命.
是健康身體健康的話.
沒有甚麼人可以守到誡命.
是嗎.
健不健康.
不到我們決定.
因為這樣的緣故.
有健康沒有健康都好.
我們都需要.
見證.
耶穌彼此相愛.

$^{921}$甚至我們會發覺.
在信徒裡面.
有些弟兄姊妹.
即使身體不是健康.
都好.
他都見證著.
其他人.
有些長期病患.
患了癌症的弟兄姊妹.
即使身體軟弱.
有機會.
有力量他都去服侍人.
去見證上帝.
這是一個很美麗的狀況.
我想教會群體都一樣.
不是等到這個社會健健康康.
才去愛人.
可能是在當社會裡面.
很多疾病.
這些愛才來得.
更加重要.
這些見證.
才是更加漂亮.
第二.
再反省下去.
健康是很好的東西.
耶穌都醫人.
不過我們發覺.
原來健康.
身體健康不是.
最重要的目標.
我們發覺原來.
健康不好.
成為我們的偶像.
我們很少這樣想.
不過如果我用.
另一個比喻大家會明白多些.
就是我們.
經常說錢很容易成為偶像.
金錢.

$^{961}$金錢是不是好東西.
當然是好東西.
當我們發覺原來金錢.
是比其他東西重要的時候.
金錢就成為我們的偶像.
健康.
都可以成為我們的偶像.
當我們發覺.
最緊張的東西.
我花很多心機.
花很多時間.
去保護我的健康.
過於.
我要愛我們的上帝.
過於.
我要愛身邊的人.
我們可能已經將健康.
成為一個偶像.
上帝.
是很渴望.
我們健康.
不過我們要問.
究竟我花在.
我身體健康的精力.
其實是不是已經.
漸漸成為我們生命裡面的偶像.
我想都是很需要.
我們去到反省.
最後.
我們又會記得.
耶穌討論健不健康.
有一段經文.
我想這段經文.
給我們一些提醒.
馬太福音第九章十二節.
耶穌說了一句話.
健康的人就不用醫生.
有病的人.
才需要醫生.
這句話是在什麼情形裡面去說的.

$^{1001}$就是當一些.
的理財人.
他發覺耶穌和那些.
罪人混在一起.
很不滿耶穌.
這樣去做.
耶穌就這樣說.
耶穌的意思是.
其實全世界所有人.
都有病.
而反理財人.
最大的問題是.
以為自己沒有病.
而當一個人以為自己沒有病的時候.
他就會對那些.
所謂有病的人.
沒有憐憫.
拒絕.
甚至傷害.
去隔絕.
將自己提起.
去隔絕.
將自己提高.
我想基督徒.
其實每一個都是.
承認自己有病.
當我們承認自己有病的時候.
我們就.
不會將一些.
真是定為有病的人.
看為和我們不同.
我想今天的社會裡面.
一個很大的問題.
就是人和人.
或者人群和人群之間的.
那種相隔.
那些矛盾.
那些紛爭.
我們看到.
原來耶穌透過.

$^{1041}$醫治麻瘋病人的事情.
去跟我們說.
人的社會.
可以因為疾病.
分隔人.
可以因為地位分隔.
可以因為種族.
立場.
屬於誰.
彼此.
分隔.
但是主耶穌基督來到.
其實跟我們說.
每一個人都有神的形象.
每一個人.
如果.
當落入一些悲慘的狀況裡面.
都應該憐憫.
都是我們.
愛的對象.
所以今天.
我想這段經文.
去提醒我們一件事.
主耶穌基督來.
救贖我們的生命.
他同時救贖我們的群體.
他知道人和人之間.
最大的困難.
就是拒絕.
就是隔菜.
而當他來到的時候.
他願意去接納我們這一班.
其實本身.
不應該被他接納的人.
我們有應當效法.
我們的耶穌.
在我們的社會裡面.
在冷漠,在隔膜裡面.
接待身邊的人.
讓所有的人.

$^{1081}$都能夠得到上帝的憐憫.
我們願意這樣去遵行.
我們一切有祈禱.
因為其實我們每一個.
都是污穢的罪人.
你沒有嫌棄我們.
你來到.
你看穿我們.
罪的肉身.
看到我們每一個都是.
寶貴的神兒女.
所以你選擇我們.
你愛我們.
讓我們能夠得到.
上帝很尊貴的.
兒女的名份.
所以求你幫助我們.
我們都效法你.
用同一樣的眼光.
去接待身邊的人.
在世界裡面真的有很多分隔.
可能疾病的分隔.
可能很多.
階級的分隔.
很多族裔之間的分隔.
但天父.
其實世界上每一個.
都是你寶貴的創造.
求你幫助我們.
在今天.
這麼多分隔的世代裡面.
我們成為你和平之子.
將到你那份無比的.
接納.
憐憫和愛.
能夠透過我們的生命.
和我們的淪陷.
去分享.
求你這樣去幫助我們.
更新我們.

$^{1121}$也聽我們的祈禱.
禱告奉耶穌基督命其求.
\newpage



\section{約翰福音 4:1-42}
\label{sec:iO1qZX8Cmq4}
\textbf{2023-10-08 早堂《生命轉化之旅》約翰福音 4\_1-42 講員:陳淑娟牧師}
\newline
\newline
連結: \href{https://youtube.com/watch?v=iO1qZX8Cmq4}{\texttt{https://youtube.com/watch?v=iO1qZX8Cmq4}} ~~~~ 語音日期: 2023-10-12
\newline
\newline
\hyperref[sec:G0FSq209Kac]{\small{< < < PREV SERMON < < <}}
~
\hyperref[sec:index_chronic]{\small{[返順時目]}}
~
\hyperref[sec:index_scriptual]{\small{[返順卷目]}}
~
\hyperref[sec:wwGm5m9DNT8]{\small{> > > NEXT SERMON > > >}}
\newline
\newline
約翰福音 4:1-42
\newline
\begin{longtable}{cl}
\hline
\hline
章節 & 經文 (和合本修訂版)\\
\hline
4:1 & \begin{tabularx}{0.7\textwidth}{X} 耶穌知道法利賽人聽見他收門徒和施洗比約翰還多， （ \end{tabularx} \\ \\ \relax
4:2 & \begin{tabularx}{0.7\textwidth}{X} 其實不是耶穌親自施洗，而是他的門徒施洗，） \end{tabularx} \\ \\ \relax
4:3 & \begin{tabularx}{0.7\textwidth}{X} 他就離開猶太，又回加利利去。 \end{tabularx} \\ \\ \relax
4:4 & \begin{tabularx}{0.7\textwidth}{X} 他必須經過撒瑪利亞， \end{tabularx} \\ \\ \relax
4:5 & \begin{tabularx}{0.7\textwidth}{X} 於是到了撒瑪利亞的一座城，名叫敘加，靠近雅各給他兒子約瑟的那塊地。 \end{tabularx} \\ \\ \relax
4:6 & \begin{tabularx}{0.7\textwidth}{X} 雅各井就在那裡；耶穌因旅途疲乏，坐在井旁。那時約是正午。 \end{tabularx} \\ \\ \relax
4:7 & \begin{tabularx}{0.7\textwidth}{X} 有一個撒瑪利亞婦人來打水。耶穌對她說：「請給我水喝。」 \end{tabularx} \\ \\ \relax
4:8 & \begin{tabularx}{0.7\textwidth}{X} 因為那時門徒進城買食物去了。 \end{tabularx} \\ \\ \relax
4:9 & \begin{tabularx}{0.7\textwidth}{X} 撒瑪利亞婦人對他說：「你是猶太人，怎麼向我一個撒瑪利亞女人要水喝呢？」因為猶太人和撒瑪利亞人沒有來往。 \end{tabularx} \\ \\ \relax
4:10 & \begin{tabularx}{0.7\textwidth}{X} 耶穌回答她說：「你若知道神的恩賜，和對你說『請給我水喝』的是誰，你早就會求他，他也早就會給了你活水。」 \end{tabularx} \\ \\ \relax
4:11 & \begin{tabularx}{0.7\textwidth}{X} 婦人對耶穌說：「先生，你沒有打水的器具，井又深，哪裡去取活水呢？ \end{tabularx} \\ \\ \relax
4:12 & \begin{tabularx}{0.7\textwidth}{X} 我們的祖宗雅各把這井留給我們，他自己和兒女以及牲畜都喝這井裡的水，難道你比他還大嗎？」 \end{tabularx} \\ \\ \relax
4:13 & \begin{tabularx}{0.7\textwidth}{X} 耶穌回答，對她說：「凡喝這水的，還要再渴； \end{tabularx} \\ \\ \relax
4:14 & \begin{tabularx}{0.7\textwidth}{X} 誰喝我所賜的水，就永遠不渴。我所賜的水要在他裡面成為泉源，直湧到永生。」 \end{tabularx} \\ \\ \relax
4:15 & \begin{tabularx}{0.7\textwidth}{X} 婦人對他說：「先生，請把這水賜給我，使我不渴，也不用到這裡來打水。」 \end{tabularx} \\ \\ \relax
4:16 & \begin{tabularx}{0.7\textwidth}{X} 耶穌對她說：「你去，叫你的丈夫，再到這裡來。」 \end{tabularx} \\ \\ \relax
4:17 & \begin{tabularx}{0.7\textwidth}{X} 婦人回答，對耶穌說：「我沒有丈夫。」耶穌說：「你說沒有丈夫是對的。 \end{tabularx} \\ \\ \relax
4:18 & \begin{tabularx}{0.7\textwidth}{X} 你已經有過五個丈夫，你現在有的並不是你的丈夫。你這話是真的。」 \end{tabularx} \\ \\ \relax
4:19 & \begin{tabularx}{0.7\textwidth}{X} 婦人對他說：「先生，我看你是一位先知。 \end{tabularx} \\ \\ \relax
4:20 & \begin{tabularx}{0.7\textwidth}{X} 我們的祖宗在這山上敬拜神，你們倒說，應當敬拜的地方是在耶路撒冷。」 \end{tabularx} \\ \\ \relax
4:21 & \begin{tabularx}{0.7\textwidth}{X} 耶穌對她說：「婦人，你要信我。時候將到，你們敬拜父，既不在這山上，也不在耶路撒冷。 \end{tabularx} \\ \\ \relax
4:22 & \begin{tabularx}{0.7\textwidth}{X} 你們所敬拜的，你們不知道；我們所敬拜的，我們知道，因為救恩是從猶太人出來的。 \end{tabularx} \\ \\ \relax
4:23 & \begin{tabularx}{0.7\textwidth}{X} 時候將到，現在就是了，那真正敬拜父的，要用心靈和誠實敬拜他，因為父要這樣的人敬拜他。 \end{tabularx} \\ \\ \relax
4:24 & \begin{tabularx}{0.7\textwidth}{X} 神是靈，所以敬拜他的必須用心靈和誠實敬拜他。」 \end{tabularx} \\ \\ \relax
4:25 & \begin{tabularx}{0.7\textwidth}{X} 婦人對他說：「我知道彌賽亞—就是那稱為基督的—要來；他來了，會把一切的事都告訴我們。」 \end{tabularx} \\ \\ \relax
4:26 & \begin{tabularx}{0.7\textwidth}{X} 耶穌對她說：「我就是，正在跟你說話呢！」 \end{tabularx} \\ \\ \relax
4:27 & \begin{tabularx}{0.7\textwidth}{X} 正在這時，門徒回來了。他們對耶穌正在和一個婦人說話感到驚訝，可是沒有人說：「你要甚麼？」或說：「你為甚麼和她說話？」 \end{tabularx} \\ \\ \relax
4:28 & \begin{tabularx}{0.7\textwidth}{X} 那婦人留下水罐，往城裡去，對眾人說： \end{tabularx} \\ \\ \relax
4:29 & \begin{tabularx}{0.7\textwidth}{X} 「你們來看！有一個人把我素來所做的一切事都說了出來，難道這個人就是基督嗎？」 \end{tabularx} \\ \\ \relax
4:30 & \begin{tabularx}{0.7\textwidth}{X} 他們就出城，來到耶穌那裡。 \end{tabularx} \\ \\ \relax
4:31 & \begin{tabularx}{0.7\textwidth}{X} 就在這個時候，門徒求耶穌說：「拉比，請吃吧。」 \end{tabularx} \\ \\ \relax
4:32 & \begin{tabularx}{0.7\textwidth}{X} 耶穌對他們說：「我有食物吃，是你們不知道的。」 \end{tabularx} \\ \\ \relax
4:33 & \begin{tabularx}{0.7\textwidth}{X} 門徒就彼此說：「難道有人拿甚麼給他吃了嗎？」 \end{tabularx} \\ \\ \relax
4:34 & \begin{tabularx}{0.7\textwidth}{X} 耶穌對他們說：「我的食物就是要遵行差我來那位的旨意，完成他的工作。 \end{tabularx} \\ \\ \relax
4:35 & \begin{tabularx}{0.7\textwidth}{X} 你們不是說『到收割的時候還有四個月』嗎？我告訴你們，舉目向田觀看，莊稼熟了，可以收割了。 \end{tabularx} \\ \\ \relax
4:36 & \begin{tabularx}{0.7\textwidth}{X} 收割的人已經得工錢，為永生儲存五穀，使撒種的和收割的一同快樂。 \end{tabularx} \\ \\ \relax
4:37 & \begin{tabularx}{0.7\textwidth}{X} 『那人撒種，這人收割』，這話可見是真的。 \end{tabularx} \\ \\ \relax
4:38 & \begin{tabularx}{0.7\textwidth}{X} 我差你們去收你們所沒有辛勞的；別人辛勞，你們享受他們辛勞的成果。」 \end{tabularx} \\ \\ \relax
4:39 & \begin{tabularx}{0.7\textwidth}{X} 那城裡有好些撒瑪利亞人信了耶穌，因為那婦人作見證，說：「他把我素來所做的一切事都說了出來。」 \end{tabularx} \\ \\ \relax
4:40 & \begin{tabularx}{0.7\textwidth}{X} 於是撒瑪利亞人來見耶穌，求他在他們那裡住下，他就在那裡住了兩天。 \end{tabularx} \\ \\ \relax
4:41 & \begin{tabularx}{0.7\textwidth}{X} 因為耶穌的話，信的人就更多了。 \end{tabularx} \\ \\ \relax
4:42 & \begin{tabularx}{0.7\textwidth}{X} 他們對那婦人說：「現在我們信，不再是因為你的話，而是我們親自聽見了，知道這人真是世界的救主。」 \end{tabularx} \\ \\
[1ex]
\hline
\hline
\end{longtable}
$^{1}$感受堂的弟兄姊妹 早安 主內平安.
很感恩能夠再一次來到這裡跟大家分享神的話語.
多謝主席剛才讀出一段很長的經文.
理論上這段經文可以分開五六篇講一講.
但我知道我只有半小時.
所以我從另一個角度去跟大家分享這段經文.
我的講題是生命轉化之旅.
這麼長的經文如果逐節看是非常困難的.
所以今天我想用另一個形式去呈現.
不知道大家有沒有看過舞台劇.
今天我嘗試帶大家進入一個好像舞台劇的環境當中.
去感受一下當時這個婦人怎樣經歷一個人生的生命轉化之旅.
如果大家有去看過舞台劇的時候.
你會發現你的座位通常都有一本長刊.
舞台劇的長刊通常講什麼呢.
講人啊 大致的內容啊 或者背景啊等等.
今天我先給大家一本長刊.
首先講一下人 講一下時間.
什麼時間 如果根據經文講.
其實這個場景發生是什麼呢.
就是猶太人的第六時.
就是正午 非常曬的.
你想像一下 現在要進入一個猶太人的時間.
正午 跟外面很大距離 外面現在狂風雷暴.
但不是 很曬的 去到一個地方.
那個地方在哪裡呢.
經文所提的就是耶穌要去加利尼.
當時講必須要經過撒瑪利亞.
於是他就到了撒瑪利亞一座城叫做聚迦.
那裡有一個井.
其實我想大家先了解一件事.
就是一般的猶太人.
如果他要來往猶大地和加利尼.
通常是不會走這條路的.
雖然事實上這條路是捷徑 是最快的.
但是敬虔的猶太人是絕對不會走的.
原因就是他必須要經過這個撒瑪利亞的山區.
當這座地方住滿了.
他們認為一些不敬虔或者不純正的撒瑪利亞人.
在猶大人和撒瑪利亞人之間.

$^{41}$自古以來一直都是爭鬥不斷的.
在耶穌的時候.
這兩族的居民更加積怨甚深.
關係非常緊張.
其實關係源於他們的種族和宗教的問題.
所以一般敬虔的猶太人是不會走這條路的.
但是耶穌這次很例外.
他卻走這條路.
為什麼呢?.
一會兒看下去你就知道了.
他來到一個地方.
有一個井.
雅各井.
大約你就可以想像.
在猛烈的陽光下.
耶穌來到雅各井前面.
人物有誰呢?.
剛才主席讀過的.
有耶穌.
一個撒瑪利亞的婦人.
還有一群門徒.
門徒出現很少.
因為門徒一進城.
接著他去哪裡?.
買東西吃.
差不多是完場他才回來.
所以他沒有看到中間那段發生什麼事.
還有城內裡面的撒瑪利亞人.
大約有這些人物.
好了.
宗教關係.
剛才說過了.
猶太人和撒瑪利亞人從來是不來往的.
如果剛才讀經的時候第九節說得很清楚.
所以到了二十七節.
回來見到耶穌和撒瑪利亞婦人說話的時候.
他們這麼稀奇.
因為剛才說他們的積怨是極深的.
雙方都不對的.
原因就是猶太人認為撒瑪利亞人是不敬虔的.

$^{81}$撒瑪利亞人是北國亡國的時候.
亞述將外族人移居到以色列地.
而亞述引入的外族人.
就和當時的以色列人和外族人一起所生的後裔.
叫做撒瑪利亞人.
簡單來說我們叫做不順政.
或者叫做集了.
血統上.
根據聶王記下.
亞述曾經在以色列這個地方.
最初這些外族人進來的時候.
他都曾經猜派以色列的祭司教導他們宗教的禮儀.
不過當時以色列宗教已經非常混亂.
再加上撒瑪利亞人接受了教導之後.
又混入了他們自己外族.
拜外邦偶像的敬拜.
所以猶太人更加認為他們不敬虔.
而且撒瑪利亞人一直不承認先知書的作品.
也不承認智慧書.
主要他們是看五經.
和猶太人很不同.
還有一個最大的爭拗點就是.
撒瑪利亞人認為敬拜的地方.
就是在事件境內的基尼心山上的新聖殿.
他們拒絕承認耶路撒冷為朝聖之地.
這也是他們的爭拗之一.
所以待會兒你會看到經文.
為什麼婦人突然問耶穌.
究竟哪裡敬拜才對.
這些是小說的背景.
介紹了人 事 地.
我們開始進入整個的場景.
第一幕在第一到九節.
我叫它做相遇.
耶穌獨自在雅各井旁.
遇到一個撒瑪利亞的婦人.
她主動和他傾談.
當時是正午.
大家留意一下.
正午的時候天氣就很熱.

$^{121}$很曬.
一般正常人都不會正午出來工作.
事實上也是.
因為當時去水打井的粗活.
是婦人做的.
不是像我們現在這樣家教自來水.
一般水打井的婦人.
會什麼時候去呢.
早上一早.
因為沒那麼熱.
或者黃昏.
其實沒什麼人會在中午去.
你看到這個婦人你都猜到.
她一定有些事情.
才會中午選擇.
最沒人的時間去水打.
在這一幕裡.
我們先明白到一件事.
這個婦人.
她是被社會.
或者社交上孤立的一位.
可能她自我孤立.
又或者她被孤立.
因為她在最沒人的時候.
才出來.
當然你後面看你會明白.
因為她有些明星上的纏繞.
她做了一些事情.
你想想.
你想像就想像到.
如果一班婦人.
一起出來打水.
通常有什麼做的.
有什麼做的.
當然是說是非.
有什麼做.
一班婦人通常都是.
我經常都說笑.
我們教會最熱鬧就是廚房.
因為最多婦人在.

$^{161}$那做什麼.
我沒這樣說.
都是閒話家常.
你想想有一個明星敗壞的人.
在當時的社會.
明星敗壞很嚴重.
她一定是成為被說的對象.
或者被人指指點點.
所以這個婦人選擇.
在中午來到這裡.
但你一路看的時候.
你發覺很特別.
這個婦人.
她很怕見到人.
耶穌卻在這裡.
好像在等她.
好像不像一個偶遇.
剛才我說.
其實耶穌是不應該.
走這條路去加利利的.
他卻選擇走這條路.
去到加利利.
主耶穌還要主動留意.
主動和這個婦人聊天.
主耶穌跨越了一些界限.
因為在那個時代.
第一男子絕對很少.
在公開場合裡.
和婦人聊天.
就算是一對夫婦.
兩夫婦也好.
他們也不會在公開場合裡.
輕輕鬆鬆.
像我們現在這樣.
男子如果單身的話.
他更加不會單獨.
和一個婦人聊天.
或者接觸她.
大家想想.
耶穌當時是一位拉比老師.

$^{201}$他是嚴格遵守這些規條.
以作眾人的模範.
但是他卻主動地.
去和這個被社會所排斥.
被人視為污穢有罪的人.
去主動和他交談.
使得這個生命之旅.
踏入了一開始的部分.
主耶穌親自地去接觸.
和他猶太信仰.
完全不同的婦人.
和他身份地位天淵之別的婦人.
和他背景完全不同的婦人.
一個被孤立被排斥的.
撒瑪利亞婦人.
而且這個婦人在最終掙扎.
耶穌就將真道傳給她.
主與她的相遇.
使得這個婦人從黑暗裡見到光.
使得生命得著轉化.
所以第一點.
我想提醒姐妹思想一件事.
我們的生命其實都經歷了一個轉化.
我們在認識耶穌之前.
都在罪裡面兜兜轉轉.
耶穌主動來尋找我們每一個人.
當我們成為主的門徒之後.
我們接著又如何呢.
主耶穌親自去尋找這個婦人.
走進不一樣的地方.
去找了一個不一樣的人.
弟兄姊妹.
今天我們願不願意走進人群當中.
去找那些和我們身份.
背景 信仰 種族都不一樣的人.
去將福音傳給他.
使得他認識主.
生命得著轉化.
主動走出你的安書區.
走到不同的領域裡.

$^{241}$很多時候我們要和親友傳福音.
已經不容易了.
還說要叫你和一些陌生人.
甚至和你身份地位.
或者甚至種族都不同的時候.
其實不容易的.
但是耶穌親自做了一個示範給我們看.
他就是走進這些地方.
去找當中的撒瑪利亞人.
你願不願意呢.
你願不願意走進社區裡.
去找這些撒瑪利亞人呢.
我教會有一個童工.
他主要做外展的工作.
外展的工作包括做露宿者的工作.
包括去到球場找一些邊青.
他每一年其實差不多四五次.
都做一些大型的教會的活動.
那個活動是什麼呢.
他會招聚我們的學生們.
去做一個叫拿世議人計劃.
就是帶我們的學生走進社區.
比如可能是一個復活節.
一個聖誕節.
或者一個暑假.
每天走到社區裡.
去接觸和我們身份背景不一樣的人.
很多教會曾經都問我.
牧師呀 面對移民潮.
面對年青人願意回教會.
教會可以做些什麼.
我很深信每一件事都不會偶然的.
但我們要知道當這些事情發生的時候.
我們不是坐在這裡等.
而是我們願意走出去.
我很多次都見到我的年青人.
回來分享的時候.
當中不少信義黛.
他說他的生命不一樣.
因為當他願意走出去.

$^{281}$接觸那些和他完全不同的人.
他大開眼界.
他看見就動了慈心.
我們是不是只是關上門.
在教會裡崇拜呢.
崇拜當然重要.
但是生命和生命的接觸和碰撞.
也都很重要.
耶穌當日親自去尋找你.
今天你願不願意因著神的緣故.
走出去尋找更多失喪的靈魂呢.
所以第一點想大家思考的.
就是你願不願意主動的.
走去人群當中尋找那些撒瑪利亞人.
你會說我不知道怎樣尋找牧師.
其實很多的.
我經常都舉一些例子.
我剛剛在神學院教完一科.
當中有一個70歲的.
應該剛剛生日.
之前上課是69歲.
他上完我的課之後.
有一天他走來找我.
上到中間.
他說牧師牧師.
我很想去服侍人.
但我不知道怎樣去服侍一些.
在水深火熱的人.
我說你祈禱吧.
雖然你今天不知道怎樣做.
課堂完結他又走來找我.
他說牧師牧師我正在祈禱.
我還不知道.
我說你試試走出去吧.
他說嗯.
我也繼續祈禱試試.
隔了一個月左右.
他走來找我.
他說牧師.
我有一個患了情緒病很嚴重的親友.

$^{321}$我知道你教會有輔導.
你可不可以幫幫忙.
我說好啊你帶他來吧.
帶了一個親友來輔導.
來輔導期間.
因為他親友輔導了一個小時.
他進我房間聊天.
他說牧師.
我現在走出去了.
我說你怎樣走出去.
他說當我祈禱的時候.
神給了四個.
我很久沒有聯絡的親友.
四個家庭.
總共二十多人.
這四個家庭都是.
我們整個家族都不會再聯絡的.
因為有好賭的.
有情緒病的.
有借錢借到整個親.
所有宗族都跑路了.
怕了他們.
有那些叫什麼.
吸毒的.
總共四個家庭.
我說你現在怎樣做.
他說我現在每一天.
陪著這四個家庭去社福律署.
去戒毒.
去見心理輔導.
帶他去找社工.
我聽完之後很感動.
當他走出我寫字樓的時候.
我和我同工就坐在寫字樓.
他的背影走出去.
他動著一支拐杖.
一個六十九歲.
大手術的一個姐妹.
他竟然能夠活出這種生命.
去尋找撒瑪利亞人.

$^{361}$他說所有的家族成員.
都對他們真的備至質.
他卻走過去.
一個一個希望能夠做到.
上一個月他有一天邀請我.
他說牧師我七十歲大壽.
我會在我教會搞感恩會.
我會請這二十多位家族成員.
四個家庭來教會.
我不知道行不行.
但我想他們認識神.
原本我答應了他去.
但那天我剛好有些重要事.
我去不了.
完了我的會議之後.
我WhatsApp他問他怎樣.
他說很感恩.
他說他二十多個家族成員.
幾乎都來了.
有六個人信了主.
最近找他.
他說有一個已經上了慕道班.
預備洗禮.
弟兄姊妹.
很多時候要我們去接觸一些.
跟我們完全不同的人.
是不容易的.
但記住.
神沒有做基督徒.
神是做人的.
每一個人都是尊貴的.
我們的責任就是將福音遍傳.
好第二幕.
我稱之為對話.
在第十到第二十六節.
一段很長的對話.
分開兩段.
第一輪的對話.
主耶穌提到關於水的問題.
因為旁邊有個井.

$^{401}$一開始主耶穌問.
這個婦人拿水喝.
耶穌將屬世的需要.
先說起.
慢慢轉移屬天的事實.
水是人生基本的需要.
每一個人的身體都需要水.
主耶穌以婦女生存的基本需要.
作為引子.
在這個世界.
她需要水.
但在永恆裡.
她需要活水.
在第十節.
主耶穌提到活水.
說喝這個活水.
永遠就不喝了.
因為耶穌就是這個活水的源頭.
主耶穌實在是我們人生當中.
能夠使我們完全滿足的活水.
當時這個婦人聽完.
其實都很糊塗.
你留意到.
不是很明白耶穌的意思.
所以她還問耶穌.
我又沒有打水的器皿.
井又深.
你怎樣拿水呢.
怎樣有活水呢.
不是很明白.
另外這個婦人在認知上.
撒瑪利亞地雅各井.
只有這個地方提供水.
但耶穌告訴她.
喝井水的人.
其實永遠都會渴.
但如果喝了生命泉源的活水.
永不渴.
雖然這個婦人好像有點迷茫.
你是不是那一刻她馬上知道你是神呢.

$^{441}$其實不知道的.
有點迷茫.
但她都渴求這個活水.
當主耶穌將福音.
一步一步傳給這個婦人的時候.
讓她慢慢明白真理.
但你看到一個傳福音的策略.
什麼策略呢.
就是從她切身的問題去談起.
從人生的基本需要.
就是水.
我們有沒有想過.
當我們做福音的工作的時候.
其實我們是不是每次都一開始.
或者我們都不理了.
總之一見到人就問.
你有沒有聽過人生的好消息.
你想耶穌沒有.
是不是這樣呢.
其實一個真正的整全的福音.
是包括一個人的身心舍靈的健康和飽足.
其實耶穌也是.
耶穌在他行道的三年半.
你留意一下.
他走遍各城各鄉.
他在做什麼.
他不單單是宣告悔改的福音.
他更加是趕鬼治病.
是不是.
他將有需要的人帶到他面前.
他身體的需要.
他先幫他去處理和解決.
今天我們經常說我們要傳福音.
但我們傳的福音.
會不會有一些只是過於簡單.
或者幸福的福音呢.
整一個大型報道會.
叫人信主.
包保你平安.
很多的成功的見證.

$^{481}$是不是只是這樣呢.
真正整全的福音.
就是要將整個人的身心舍靈重新建立.
讓他認識神.
讓他重回自己的尊貴和價值.
這個才是整全的福音.
我經常舉一個例子.
當有一個人肚餓的.
走進感受堂.
他說我很肚餓.
你可不可以給我一杯水和一包.
難道我們就跟他說.
你聽了人生的好消息沒有呢.
你不會的.
你不會拿一本四律三福跟他說完.
才給他一包和一杯水吧.
你會說不會的牧師.
不過我會請他吃完一包和一杯水.
然後就馬上跟他說.
但是我們會不會真的問他.
還有什麼需要呢.
一個人的整全是很重要的.
很多時候有人說傳福音很難.
其實福音的工作.
不是靠我們單單告訴別人.
主耶穌有多好.
福音的工作是應該這樣的.
我們以上帝的愛.
基督的愛.
用我們的行動.
去關懷那些真是有需要的人.
無論他是心靈貧乏.
實質物質貧乏.
然後讓人見到上帝的Being.
神就是愛.
為什麼有時候人很不願意信主.
當然有很多原因.
其中一個原因就是.
他在我們身上見不到基督的愛.
今天快要八號風球.

$^{521}$不知道大家記不記得上次十號風球.
你們家有沒有漏水.
我剛剛跟你們一位姊妹說我家漏水.
我們教會除了旺角之外.
我們還有一個在上水.
今天八號風球我最掛心的是.
因為我們很主要上水我們做的是.
我們的街坊全部住寮屋.
當時我是很擔心因為還排洪.
我猜他們水浸的機會很大.
所以那天晚上我立刻打給我的傳道同工.
我說無論如何今晚你了解了情況.
明天一早什麼都不要做.
就是發動所有弟兄姊妹去清淤泥.
還有立刻幫他們買回來.
如果真的浸.
我都猜到他們每次浸.
他們床都浸了.
那些電器.
雪櫃都是在水裡飄來飄去.
但那晚他們還要睡覺.
雖然開了那些庇護中心.
其實我們那晚都立刻打給街坊.
你需要進去庇護中心嗎.
十個街坊都拒絕.
為什麼.
因為進了庇護中心可能間屋都吹到沒有.
他們要守衛家園.
在那天早上.
我的同工和弟兄姊妹.
真的覺得他們很厲害.
是即時.
因為十月封是第二天.
其實那些床.
他們發給我的照片.
有些屋已經浸過腳踝.
有些是浸過上下格床的第一格.
已經浸到這麼高.
浸了整間屋.
全間屋都是淤泥.

$^{561}$但那晚他們還要睡覺.
因為接著還在下大雨.
他們也不想去庇護中心.
那天我們發動很多弟兄姊妹.
包括旺角堂也好.
上水堂也好.
即時做什麼呢.
第一幫他們買回所有床.
床褥.
接著呢.
幫他們買了一些緊急電器.
因為我看到有一個雪櫃.
真的很誇張.
飄在水上.
已經完全不能用了.
買回雪櫃.
有一個小朋友哭得很厲害.
我以為他在做什麼.
發短訊給我.
原來他唯一一部.
三年前二手的ipad也浸爛了.
他很怕不能上網購物.
他不是不怕不能上網打機.
他說我每天一篇都還沒做.
他真的是這樣.
我們很感動.
我們馬上幫他們買了所有東西.
在買的過程中.
我們很感受到.
第一就是街坊們很好.
一邊買一邊講價.
其實我的同工也不講價.
因為快點幫你買了.
其實當時我們還沒籌到款.
我同工問牧師有沒有多少預算.
我說不知道.
因為還沒籌款.
因為這麼急.
先買完.
不要理.

$^{601}$還沒籌到款.
但是我的同工就去買.
一邊買的時候.
我們做了五十多戶.
有一半都還沒信主.
他們都.
不因為我們買給他們.
而貪心要多一點.
不斷在店鋪裡磨.
一百塊都講價.
五十塊都講.
不斷講不斷講.
我同工發給我看.
我覺得人和人之間的那種愛.
就是這樣建立.
其中一個同工分享.
我很感動.
有一個同工就這樣跟我說.
當中有一個婦女.
她有一位先生一直很反對她回來教會.
這位婦女不敢告訴她信了耶穌.
只是說帶小朋友回來補習.
但這位先生一直買電器的時候.
他自己暗暗沉沉.
他說原來信耶穌那些.
不是我想像中那麼可怕.
我以為他們只是叫人奉獻.
他誤解.
我以為進了教會是負責奉獻.
但原來不是.
原來信耶穌的人比我想像中更有愛心.
然後拍一拍他老婆.
你帶我回去.
結果兩個星期後.
我們見到這位男士回來教會.
我們最初沒有辦法跟他傳福音.
你想想他老婆回來教會都要瞞著.
但弟兄姊妹.
什麼叫做讓基督的愛彰顯在我們身上.
我想大家思想.

$^{641}$今天我們每一言每一行.
當社會面對這麼大的風浪的時候.
十號風球很多人流離失所.
我們在做什麼呢.
我們是哇!正!當日假期.
今天生日千萬不要祈禱.
想著風球打到明天.
明天不用上班.
其實如果風打到明天.
有很多人其實又沒有了一天的工資.
有很多家庭面對的困難更大.
今天我們作為基督徒.
我們應該怎樣將上帝的愛.
在我們裡面透過行動可以活出來.
我想大家去思想.
第二輪的對話.
耶穌開始討論完水的問題.
接著開始討論到生命的問題.
耶穌將婦人的罪.
生命的罪輕輕帶出來.
不過這個婦人極度聰明.
馬上轉話風.
當婦人說我想要活水.
耶穌就話題一轉.
叫你丈夫來.
接著她就說我沒有丈夫.
為什麼她說她沒有丈夫呢.
耶穌說她都答對的.
沒錯的.
因為她離了婚或者是寡婦.
現在那些男士都不是她的丈夫.
主耶穌為什麼明明說要活水.
突然間說這件事.
耶穌就要將她的罪輕輕帶出來.
讓她看到自己生命中的實況.
主耶穌是神當然知曉一切.
但他都讓這個婦人重新提醒自己.
她現在生命的狀況是怎樣.
所以大家很清楚.
為什麼她中午來打水.

$^{681}$因為明星一直纏繞著她最大的問題.
婦人見到自己的狀況被披露.
罪被人說出來很醜.
所以她不想耶穌再追問她道德方面的問題.
所以她急轉彎.
在20節她立刻搬出一個問題.
就是猶太人和撒瑪利亞人.
由久以來的宗教分歧.
就是敬拜的地方.
所以你明白為什麼她突然提敬拜的地方.
因為一開始說了背景.
所以她問究竟在基里森山的新聖殿敬拜.
還是在耶路撒冷的聖殿敬拜.
主耶穌沒有逃避婦人這個屬靈的話題.
反而把這個屬靈的話題.
教導婦人什麼叫做真正的敬拜.
就是用心靈誠實去敬拜.
然後耶穌告訴她.
我就是這位主.
在第26節主宣告她就是彌賽亞.
他將這個婦人帶進去整個光的當中.
人接觸到光就知道自己黑暗.
耶穌就是真光.
當人願意接受耶穌成為光明之子.
最得赦免.
這個婦人遇見光.
她的生命就開始不同的更生和轉化.
每一個人都需要的.
所以我們今天屬主的人已經是光.
我們有沒有將光帶進黑暗裡.
使人都可以得著心意更生轉化.
轉化這個字其實是來自這段經文的.
婦人為什麼能夠心意更生而變化.
生命轉化是因為耶穌進到她的裡邊.
當人遇見神接納真理.
就是生命轉化的開始.
今天我們每一個人都應該有生命轉化.
我們開始了我們的旅程裡邊.
但我們是不是真的被轉化了呢.
剛才施班讀了一段經文.

$^{721}$很清晰地告訴我們.
如果我們作為基督徒.
我們的使命是盡心盡意盡性愛你的神.
並且愛倫如己.
真正的愛就是關顧到身邊的人.
身心舍靈的需要.
看見人貧窮缺乏.
給他一杯涼水.
看見未信者的靈魂.
我們願意將他帶到神的面前.
我們有沒有盡諸般的智慧.
將人帶到神的面前呢.
主耶穌他親自做了一個傳福音的示範.
從屬地的需要去了解.
開始了一個話題.
接著慢慢去講屬天的需要.
又將這個人的罪慢慢顯出來.
顯他的罪出來.
不是要當刻定他的罪.
而是要拯救他.
將他帶進真理的裡面.
去到第三幕.
收割.
門徒出場了.
門徒突然出來.
買完東西回來.
很奇怪.
看見耶穌竟然跟一個婦人說話.
奇怪就沒有人敢問.
但無論如何.
他問耶穌還吃不吃.
他買東西回來.
耶穌就說我有食物了.
耶穌吃了什麼.
就是進行猜我來者的旨意.
作他的功.
今天你有沒有食物.
你有沒有遵猜你來者的旨意.
作他的功.
你有沒有做上帝給你的工作.

$^{761}$有沒有好好將福音遍傳.
使人和神和好.
耶穌跟門徒說了一個莊稼的比喻.
莊稼已經屬鳥了.
說明了門徒就應該信服行事.
將福音遍傳.
耶穌先做了一個示範.
其實他做了一個完美的示範.
他在井旁殺種.
接著做收割.
不是收割了嗎.
不是那個婦人就信了嗎.
不單止.
他在井旁殺種.
還有些收割.
在罪中章.
在39到42節.
收割什麼.
在39節.
這個撒瑪利亞再次出場.
收割就在眼前.
因為基於撒瑪利亞婦人作見證.
很多人來找耶穌.
接著他們再經歷耶穌的講道.
他們很多人信了主.
這個就是收割.
直到眾人都宣告耶穌是救世主.
我覺得大家也可以看看這個婦人的生命經歷.
她原本是一個被社會排斥的.
你想像一下在這個中午之前.
她連接觸人群都不敢.
婦人都不敢走去城裡面.
靜靜的日間去打水.
千萬不要接觸到人.
但是當真道進入她裡面.
她成為光的時候.
她竟然可以走回城裡面.
跟所有人說.
其實你有沒有想過.
她怕不怕別人不相信她.

$^{801}$你有沒有想過.
她會不會覺得別人說.
你什麼人啊.
你跟我說這些話.
你是罪人.
你的生活很惱教.
她沒有這樣想過.
因為她生命轉化了.
她就走回合城裡面.
告訴別人她看到的所有東西.
弟兄姊妹.
今天我們呢.
當我們領受了真道的時候.
我們有沒有走到合城裡面.
去見證神呢.
還是我覺得可以.
我有穿這個福音.
我自己可以.
其他人的事.
當然我們傳福音.
就好像這個婦人一樣.
信不信呢.
合城信不信呢.
其實不是我們的能力.
我只是介紹給她.
這些人呢.
她自己來找耶穌.
經歷一定的對談.
然後她們自己經歷神.
我們都是.
我是學完全道會成長的.
學完全道會有句說話.
就是我們單單傳講基督.
並將結果教給神就行了.
靠著聖靈行事.
因為不是依靠勢力.
不是依靠才能.
那是世界上的靈方能成事.
但問題是我們有沒有走出去.
有沒有去做呢.

$^{841}$這個很重要.
是吧.
最後我想鼓勵大家.
教會和信徒.
要走進冒險之旅.
耶穌進入撒瑪利亞地.
超越界限.
走出一般人會走的路.
去找一些完全不一樣的人.
去將福音傳給他.
走進不一樣人的內心.
關顧他身心社靈.
教會能不能夠呢.
教會和信徒.
能不能走進未信的群體.
用基督的愛.
透過你們的行動.
讓人見到上主的being.
然後將榮耀歸給主.
還是今天.
別人只見到一群.
關上門Hallelujah的信徒.
願聖靈提醒我們.
什麼是盡心盡義.
盡性去愛我們的神.
並且愛倫如己.
願我們都能夠.
謹守遵行大使命和大誡命.
願主祝福大家.
\newpage



\section{撒迦利亞書}
\label{sec:wwGm5m9DNT8}
\textbf{2023-10-12 周四晚查經班《歸回先知書研讀(5)撒迦利亞書》 講員:曾裔貴牧師}
\newline
\newline
連結: \href{https://youtube.com/watch?v=wwGm5m9DNT8}{\texttt{https://youtube.com/watch?v=wwGm5m9DNT8}} ~~~~ 語音日期: 2023-10-14
\newline
\newline
\hyperref[sec:iO1qZX8Cmq4]{\small{< < < PREV SERMON < < <}}
~
\hyperref[sec:index_chronic]{\small{[返順時目]}}
~
\hyperref[sec:index_scriptual]{\small{[返順卷目]}}
~
\hyperref[sec:HWMu4dld7_M]{\small{> > > NEXT SERMON > > >}}
\newline
\newline
$^{1}$我們記掛現在在二巴的戰爭,一些都是我們不想見到,而又的確發生,一些彼此的仇恨,一直積聚下來的慘痛..
我們很盼望神來制止,特別我們一直在看舊約的時候,都是和耶路撒冷的平安有關係,我們很盼望我們的主來施恩..
雖然現在我們很難透過聖經去預測那些局勢,但是我們盼望的都真是人與人之間,因為福音帶來的和睦,因為福音帶來的可以成為神兒女,這個最重要的身份,願神來幫助..
我們再一次來到今天晚上,我們可以再打開撒迦利亞書,很盼望我們真是能夠看到這位耶和華公義的苗裔主耶穌,你要掌權,求主讓我們能夠有這份信心,這樣來等候..
我們仰望等候禱告,奉耶穌基督的名,阿們..
我們來到第六章的經文,簡單之前說過一點,但是我們再認真仔細去看第六章,其實就是八個義將裡面最後一個,這個第八個義將和第一個義將是有一點關聯的,因為都是提到一些戰馬,一些騎馬的,當然有不同的,不同的就是這次這些馬是有馬車的..
如果在古代,馬加上車可能是戰車,等於今天的現代,如果打仗就是坦克..
這裡就告訴我們第八個義將,尤其是提到北方,北方是來自最強大的軍隊,為什麼不是東方,因為是大的沙漠,所以就不會由沙漠過來,一定是沿著河岸線,所以北方是最強大的以色列,以色列做單位,.
又忘記了,第二節,第一輛車套著紅馬,第二輛車套著黑馬,第三輛車套著白馬,第四輛車套著有斑點的撞馬,特別強調是有馬車的..
(主持人翻譯).
接著就是有斑點的馬往南方去,撞馬出來要在遍地走來走去,天使說,你們只管在遍地走來走去,他們就照樣行了..
(主持人翻譯).
後來就有白馬回來問,你什麼時候要眷顧耶路撒冷,但是這裡四輛馬車或者戰車,就是奉差遣要出去的,出去了就沒有回來了,而且是在遍地走來走去,不過特別強調,.
雖然我們很難分析到白色,紅色,黑色有斑點有沒有什麼特別的意義,我們很難去做一個這樣的評估,可能有也說不定,不過特別強調,其中有兩輛戰車往北方去,特別強調天使要安慰撒迦利亞先知,.
告訴他們,因為北方的戰車已經安慰了天使的心靈,就是說令到北方得到太平才會有安寧,所以似乎這個最後的異象,.
都是再一次告訴我們,不單止耶路撒冷神會再眷顧,成為這個地方的火城,保護這個地方,而且上個星期我們說了,在桶裡面的婦人,在那個糧氣被運到她應該屬於她自己的地方,.
就是那個叫士拿地,就是說罪惡遠離耶路撒冷,當然這個不是目前可以應驗,因為我們看到的基本上解決不了,但是將來神的時間到,.
以色列不單止他們有保護,神再次眷顧,罪惡遠離他們,甚至乎周圍的列國用第八個異象來給我們看到,四圍都是受保護,.
北方的列強過去都是巴比倫,一直攻打耶路撒冷,這一切都會過去,所以八個異象都是好的兆頭,全部都是在現在他們重建聖殿的時候,.
重建之後的第二年,大理烏王的第二年的時間,這樣重建,但是一而再的這個疊一個的異象,都是想告訴我們神再眷顧他的子民,.
和將榮耀再給這個民族,去到第九節,特別提到有一個祭司,他要背起兩個職任,大祭司本身已經是全民族最崇高的,.
關乎到憲制的,另外有君王的職份,但是在這裡有一個預言,當然立刻應驗就是約書亞,但是有一個預言就是有一個大衛的苗裔,.
他會在這兩個職任裡面,他一個人背起,基本上我們在新約大部分的福音書和保羅的書信,都是提到主耶穌既是祭司,尤其是君王,.
很奇妙就是這些舊約的經文,在主前的大概五百多年,現在去到耶穌的年代,就立刻應驗了,雖然不是今天我們在教會的時代,可以見到耶穌回來有這兩個職份,.
但是我們的主在這段經文裡面已經應驗了他自己既是君王又是祭司,第九節,耶和華的話臨到我說,你要從被擄之人中取黑代多比亞耶大亞的金銀,.
為什麼特別從這幾個人,這幾位領袖拿金銀財寶,可能他們是一些很敬虔的人,這三人是從巴比倫來到西班牙移居約西亞的家裡,特別提到他們是被擄歸回的人,.
可能有一定的社會階層地位,也可能特別敬虔,所以神就叫撒迦利亞去他們這三個人的家裡拿金銀,為什麼呢?.
因為當日你要盡他的家,取這金銀作冠冕,戴在約撒達的兒子大祭司約書亞的頭上,約書亞現在忙著重建,鼓勵著百姓,當然重建的project,那些工程就由索羅巴伯帶領著做,.
但約書亞是在做那些祭祀禮儀,和什麼時候可能要有什麼的讀經訓練,約書亞正在做,現在很特別,神就要叫撒迦利亞做好這個冠冕之後,要戴在大祭司的頭上,.
又不是戴在索羅巴伯,索羅巴伯是民間的領袖,是他帶領第一批的人由巴比倫來響應這個古烈王的呼召,要回去自己的原居地來建造耶路撒冷的聖殿,.
現在兩個,一個是索羅巴伯,是民間領袖,省長,另外就是大祭司,大祭司本身不是做君王的,官面很忌諱,我是做大祭司,我不是做君王的,為什麼要戴官面在我的頭上呢?.
這裡就是第十二節,對他說,萬君之耶和華如此說,看那名稱為大衛苗裔的,他要在本處長起來,那個苗裔,並要建造耶和華的殿..
本來大祭司約書亞,現在的而且確實實在在協助索羅巴伯來建造神的殿,或者重建,荒廢了十幾年,因為有本地人來阻撓,.
現在很特別的是,將來有一個真正的苗裔,他會建造耶和華的殿,所以好像與大相關,現在的大祭司約書亞的而且確是參與重建實實在在,爛了尾的聖殿,.
但是現在告訴我們,透過這個戴官面的具體行動,告訴我們將來有一個真正的加冕,真正的既是大祭司,又是君王,在舊約已經這麼玲瓏浮凸,.
告訴我們將來會有這一位大衛苗裔,即是所說的彌賽亞,來去坐這個位置,十三節,他要建造耶和華的殿,並擔負尊榮,留意這裡,坐在位上掌王權,又必在位上作祭司,使兩職之間,籌定和平..
即是很和諧,即是說祭司和君王一個人擔當兩個職位,而且是做得非常的貼切,是很恰如其位的,這樣來到將來,又既掌王權,又可以在祭司這個職任上,.
不是的,你來獻贖聚祭,很憂愁,很困苦,那祭司就要很了解你的內心,聽完你說了,就幫你獻贖聚祭,跟著代表神來跟你說,這個祭生已經擔待了你的罪,你平平安安回家了,.
這就是離聖殿獻祭的心靈的狀態,所以這裡告訴我們,原來將來有一個人是大衛苗裔,說得很清楚很白的,是大衛的苗裔,.
就是說一定是大衛的子孫,將來在這條根裡面,在這個根源裡面長大,他就可以有這樣的榮耀,而且做君王,君王就是帶來國泰民安和平,君王帶來的最重要都是公義的審判,.

$^{41}$祭司就帶來百姓的心靈得到無論是祝福,禱告,帶給百姓有心靈的真正滿足,其實就是將百姓帶到神的面前,因為他看不到神,他看到祭司,祭司幫他獻祭,.
但是祭司代表神說的話,就是等於神向他作出的應許,或者是安慰,所以這裡給我們一個很大的福氣,就是我們本來舊約不是屬於我們的,.
但是我們在耶穌的裡面,我們就能夠有份於這位彌賽亞帶給世人,這裡本來就只是猶太人,但是帶給世人有這麼大的一個普遍的祝福,.
就是君王的治理,牧養和審判,祭司帶來的安慰和祝福,帶給那些憲制者..
十四節,這官臣要歸欺憐,就是這個黑袋,多比亞,耶大亞和西班牙的兒子,他叫做賢,或者就是約書亞,放在耶和華的殿裡為紀念,在神的面前被紀念這個官臣,.
所以我們要在這裡的時候稍微停一停,看看大衛的描繪這個概念,其實由以賽亞書已經開始有的了,以賽亞書的第四章,本來一直都在責備當時的耶路撒冷,.
都是充斥著很多假冒為善,很多的罪得罪神,無論男男女女,或者不同的階層,因為第四章都挺多描述當時的一些耶路撒冷的婦女,他們那些的情況,似乎當時他們的拜偶像都沒有停止過,.
第四章有一個這樣的預言,在那日七個女人必拉著一個男人說,我們吃自己的食物,穿自己的衣服,但求你許我們歸你名下,求你除掉我們的羞恥,.
到那日,耶和華發生的妙,必華美尊榮,地的出產必為以色列逃脫的人,顯為榮華茂盛,主必以公義的靈和焚燒的靈,將錫安女子的污穢洗去,.
亦都是第三章上文所描述的他們的種種惡行,這裡告訴我們,將錫安女子的污穢洗去,又將耶路撒冷中殺人的血除乾淨,那時勝在錫安留在耶路撒冷的,.
就是一切住在耶路撒冷,在生命冊上記名的,必稱為聖,即是有些所謂愚民,不是打魚的魚,聖魚的魚,就是愚民,在二賽亞書很多出現這個愚民這個字,這個概念,.
即是說有些表面是死靜,其實就是被揀選,就是chosen people,chosen這個概念,在二賽亞書又會出現,經常提這些愚民,那聖魚的將來就必稱為聖了..
第五節,耶和華也必在錫安泉山並各會眾以上,使到白日有煙雲,黑夜有火焰的光,不是現在的導彈,即是神將來會親自降臨,白日你會見到煙雲,黑夜會見到火焰,都合理,雲處火處,因為白天你見不到火,因為太陽太光,就會有些煙雲,.
晚上就不會有光,所以會有些火焰,所以雲處火處,帶我們去看神顯現在西瀨山和他們過紅海的時間的景象,.
即是說神再次會降臨,會紀念耶路撒冷這個民族,因為在全榮耀之上必有遮蔽,必有停止,即是有個亭,避雨亭,你可以想像得到如果你在野外,.
在亭上或山洞上,白日沒有事,完全不會有問題,這裡說必有停止,白日就可以得蔭避暑,如果你行山就知道,有個亭可以避到熱,又可以作為藏身之處,躲避狂風暴雨,.
所以我認識一個弟兄叫蔭亭,就是這個字,蔭就是蔭避,亭就是亭止,相當適合聖經,將來神會成為我們的蔭避,成為一個亭止,.
你可以很疲乏的旅客,可以在那裡休息,在那裡避暑,突然間狂風暴雨,你記住你在山上,當堂山洪暴發,如果你有個地方可以避到這些山洪,那真的是那些狂風暴雨,.
所以我們人生都需要記住,神永遠是我們的蔭避,神永遠是我們的蔭亭,很重要的,這裡告訴我們,是因為大衛發生的苗裔,所以就會有這一份尊榮,一直一直來到,再次臨到耶路撒冷,.
所以神自己親自成為我們的蔭避,所以如果改名,男孩子改蔭亭都不錯,他的弟弟叫蔭泉,泉水的泉,又適合聖經,.
我們再揭下去,耶利米書23章,還有他的大哥叫蔭盤,蔭避的盤石,又適合聖經,二次元書28章,是說有一個盤石,就是我們的主耶穌的活石,.
所以用蔭字暗示父親相當敬虔,他的下一代要有上主的蔭避,這個都真的很好,23章整章聖經不可以讀完,因為相當長,但是我們都要留意,由第一節讀起,.
耶利米書23章,雖然我不夠時間跟大家讀,這章聖經說了一些預言之後,就會說一些當時的耶利米時代的先知,因為耶利米被人label了是假先知,為什麼假呢?假在哪裡呢?.
就是他的預言,說神給我的預言,因為巴比倫已經在攻打耶路撒冷,所以耶利米就說,神告訴我們,你們去投降,跟他們去巴比倫,去那裡自買田地,安居樂業,七十年之後你們就會回來,這樣的信息叫做預言,.
所以其他的同行,同工,一定是不可能贊成,因為當時的西底加王主催要打到輸為止,不可能我們會投降巴比倫,所以在西底加王身邊的先知,適時無,.
當然就皇帝想怎樣就怎樣,就說神告訴我們,我們一定可以抵抗巴比倫,是不會失敗的,所以耶利米現在是出賣我們的國家,要我們投降,那你很難分辨到,誰是真誰是假,.
有些人就更加要來迫害耶利米,就將他關起來,不給他吃東西,就將他放在一個枯乾的井裡面,希望讓他自己自生自滅,耶利米就很有骨氣,你們什麼時候戰敗的時間,就是神的時間應驗了,.
這樣的言論有個別的領袖,不知道同不同意他這樣說,就找一些餅,接濟他,每天叫人吊一個餅給他,就是放在井下面給他,你可以想像到耶利米你只需要講一般人講的就已經足夠了,.
為什麼要這樣去陷自己在一個這麼痛苦的地方,所以如果你有興趣去看他對那些先知的警告,其實從第九節開始就有了,.
23章我們主要看的第五節,.
耶和華說:那些殘害趕散我草場之羊的牧人有禍了..
耶和華以色列的神斥責那些牧養他百姓的牧人,如此說:你們趕散我的羊群,並沒有看顧他們,我必討你們這行惡的罪..
這是耶和華說的,我要將我羊群中所餘剩的,從我趕他們到的各國內,焦聚出來,令他們歸回本圈,他們也必生養眾多,我必設立照管他們的牧人,牧養他們,他們不再懼怕,不再驚惶,也不缺少一個..
這是耶和華說的,一開始就說那些趕散群羊的牧人,那就是說那些不是好的牧人..
第五節,我們留意這裡,耶和華說,日子將到,我要給大衛興起一個公義的苗裔,他必掌王權,行事有智慧,在地上施行公平和公義..
這就是苗裔本身的定義,神的苗裔是由大衛王朝的子孫出的,苗裔將來做什麼呢?有什麼帶來地上的福氣呢?重點是公義,提了很多次了,他的行事有智慧,在地上一定是施行公平和公義..
使得在地上的人有平安,因為你會有一個標準,標準是他這個公義的苗裔所定的,大家都活在那個標準底下,就很有平安了..
因為他不會收賄賂,他不會偏幫人,他是按真理,公義來做審裁的..
所以第六節,在他的日子,猶大必得救,以色列也安然居住,他的名稱為耶和華,我們的義..
耶和華的另一個名字就是公義的日頭,在比特前書和後書提及過,可以想像到主耶穌來這個世界,他如何能夠做到這麼重要的職份呢?.

$^{81}$所以任何一個人,任何一個制度,任何一個君王,我們在目前也看不到,連大衛王我們也看不到,他真的這樣在他的任內以公平公義來施政和來執掌王權..
因為一執掌王權,就慢慢就沒有公平,就沒有公義,這個很可惜,人性就是這樣..
這個就告訴我們,到了耶利米蘇,這個描繪又再次強調,在這個時候,其實整個耶利米蘇是國破家亡的時候,正正就是沒有一個君王是好的..
由約西亞之後的全部君王都不好,約西亞的兒子,約哈斯,約瓦敬,約瓦根,西底加,西底加其實不是約西亞的兒子,是他的兒子,不過就不是約瓦敬的兒子,不是子傳給另外的兒子..
西底加其實是約哈斯的兄弟,所以你可以想像得到整個猶大都是注定一定要被腦被神來懲罰,巴比倫來將整個猶大結束..
但是在這樣的情勢下,神透過耶利米,他其實很慘,沒有人聽他講,受盡凌辱,西底加王在國情的會議裡面,就召耶利米入宮,就問他現在應該怎樣怎樣..
其實也不是耶利米,其實是耶利米的書記,或者是他的一些先知門徒,就叫巴祿,就來進見,就說耶利米有一些東西給皇帝你看,皇帝看完之後,當時是冬天,就有些炭火燒,.
馬上將書卷放下去燒了,真的先知,得到的就是這樣的結局,皇帝又怕他,又想聽聽他講什麼,但是不喜歡聽,不喜歡自己的喜好,就不理他..
這個就是輸打贏要,一般都是這樣的,你不是遷就我的想法,說明我們是要振奮軍心,打贏巴比倫,讓我的名聲可以被提高,你說叫我們投降,投降最容易的,就馬上可以解除那些威脅了..
但是成為亡國奴,聖殿裡面的器皿都要被帶去成為俘虜的,擄掠了的物品,好像火燒圓明園,八國聯軍,所以這裡就告訴我們,耶利米不灰心,耶利米看到的是將來的耶路撒冷,是看到將來的彌賽亞..
他看不到耶穌,因為他已經死了,在這個時候是主前的六百多年,但是我們看到,我們現在住在我們心裡面的就是主耶穌,和他現在在天上幫我們預備居所,他要再回來,就是這位耶和華,我們的兒,將來就要執掌皇權了..
就是以公平公義來施行他的治理,他的皇權,所以你可以想像到任何一個國家,古今中外,都不會有一個國家,又公平又公義,都不會有一個君王可以一直以公平公義去治理他的國,令到那些百姓很全神的,.
在農業社會安然居住,夜不閉戶,可以這麼安殊的,就是說五穀豐收,有新酒盈日,百姓就可以在葡萄樹下,無花果樹下,來到去暢所欲言,大家來到去很安然的,這樣沒有人驚嚇的,這樣的描述..
你可以想像到,現在我們在香港都很太平,沒有疫情了,我們有那份內心的,出街我們又不需要戴著口罩,很開心的到處去,大的節期,譬如中秋節,小朋友都出街就證明很安全了,證明很安然居住了,開開心心的拿著燈籠,這些都是中國人我們明白的安然居住..
但是神給我們的救恩,那個規模是要這樣的水平,所以我們要為全球祈禱,就算在我們一生祈完了,回天家了,地上仍然打生打死,但是都不代表神的旨意是不會成就,因為將來主耶穌再回來的時候,就是這個應驗..
所以我們要記住,這個大衛的描述在耶利米斯又再出現了,以賽亞書是離開了耶利米很早的,以賽亞書是主前的大概八百多年,因為所羅門大概一千年..
還有一二百多年之後,耶利米斯出現了,再次印證大衛的描述是神給耶利米的一個預言,然後到了第十五節,再有出現了..
我都不讀了,因為下面就說了假先知,怎樣去判斷假先知..
另外再揭一段,三十三章,同樣都是耶利米斯的,三十三章,其實耶利米斯是很好看的自傳,有預言,但是更多有耶利米自己的獨白..
難看的就是他會交錯在不同的君王的時期,是那裡難看,但是其他很好看的,他的內心的痛苦的一個陳述,尤其是第十九章,說他自己..
三十三章,耶利米還囚在護衛兵的院內,被人監禁了.耶和華的話第二次臨到他說,成就的是耶和華,造作為要建立的也是耶和華..
耶和華是他的名,他如此說,你求告我,我就應運你,並將你所不知道又大又難的事指示你..
說笑地說,這就是三三三,耶利米斯的求救電話,我們香港是999,耶利米斯就是333,神說你求告我,我會將又大又難的事指示你..
因為他得到的預言,沒有人聽,沒有人覺得是神的預言,只會覺得這個耶利米斯是出賣國家,其實很難..
第四節,論到這城中的房屋,和猶大王的宮室,就是拆毀為黨敵人高磊和刀劍的.耶和華以色列的神如此說,人要與迦勒底人(即是巴比倫)征戰,正是那死屍充滿這房屋..
就是我在怒氣和憤怒中所殺的人,因他們的一切惡,我就掩面不顧這城.看!我要使到這城得以存於安書,使城中的人得依治,又將封城的平安和誠實顯明於他們..
我也要使到猶大被老和以色列被老的歸回,並建立他們和起初一樣,我要除正他們的一切罪,就是向我所犯的罪,又要赦免他們的一切罪,就是干犯我違背我的罪..
這城要在地上萬國人面前,使我得頌讚,得榮耀,名為可喜可樂之城.萬國人因聽見我向這城所賜的福樂,所施的恩惠平安,就懼怕戰驚..
耶和華如此說,你們論這地方說,是荒廢無人民無生畜之地,但在這荒涼無人民無生畜的猶大城邑和耶路撒冷的街上,必在聽見有歡喜和快樂的聲音,新郎和新婦的聲音,並聽見有人說:.
要稱謝萬君之耶和華,因耶和華本為善,他的慈愛永遠長存..
而你不說這些預言,就是被擄,荒廢,七十年後,這個聖殿會重新建造,會有歡樂,會有規模和神的祝福眷顧,但是都未發生,聖殿還活著,敵人現在只不過是攻打我們,.
我們只要守住我們的城牆,就可以解除這個危難,但是耶利米說不要浪費時間,鬆開手,跟隨巴比倫去,所以就帶來的就是他自己要面對這些預言,帶給他自己的仇父..
第十二節,萬君之耶和華如此說,在這荒廢無人民無生畜之地,並其中所有的城邑,必再有牧人的住處,他們要使到羊群,躺臥在那裡,在山地的城邑,高原的城邑,南地的城邑,遍雅敏地,耶路撒冷四處的各處,.
和猶大的城邑,必再有羊群從數點的人手下經過,這是耶和華說的,遍山都是羊,整個以色列佈滿了神的恩典,他們是一個牧羊的民族,.
有經濟,有福樂,你可以想像得到是一個很好的預言,但是問題就是現在有聖殿,現在有國度,怎麼可能這個國家無端端會沒有了,成為荒廢無人民無生畜的地方呢?.
大家看不到這個遠景,看不到前面有一個危機,就是巴比倫危機,看不到,所以怎麼可能會有國破家亡,七十年後重建,回歸這樣的片段呢?.
所以耶利米的說話是不合時宜,只是一個報憂不報喜的先知,十四節,耶和華說,日子將到,我應許以色列家和猶大家的恩言必然成就,當那日子,那時候,我必使大衛公義的苗裔長起來..
再一次在耶利米書裡面再預言大衛的苗裔,所以這個概念對猶太人是很牢固的,他們等待的彌賽亞就是大衛的子孫,所以馬太福音大家懂不懂得背?.
第一句,亞伯拉罕的子孫,大衛的家譜,耶穌基督的家譜記在下面,然後就說那個家譜的甚麼了..
即是告訴我們大衛公義的苗裔會長起來,他必在地上施行公平和公義,在那日子猶大必得救,耶路撒冷必安然居住,他的名稱為耶和華,我們的義..

$^{121}$第二次再一次去肯定,這樣說了這個苗裔,所以很特別的地方就是當我們回到撒迦利亞書,我們要看這段經文的時候,或者我們都揭一揭馬太福音的第一章,.
可能大家會覺得很奇怪,為甚麼馬太這位瑞麗,她會一開始記載福音書的時候,首先記載一個家譜,十四代的家譜,就是亞伯拉罕的後裔,大衛的子孫,耶穌基督的家譜,是由亞伯拉罕開始,不是由阿當開始..
蒙應許的子孫,是由亞伯拉罕繼起,是神應許的,在創世記十二章,就說你的子孫將來就會蒙福,諸如此類,所以一再一再去兼顧神自己的預言,.
那個應許,那個祝福,在大衛的苗裔身上,所以很特別,如果你懂得舊約的概念,關於大衛苗裔,所以如果你是猶太人,如果你是很平心靜氣去讀馬太福音的福音書,你就會很驚訝,.
他講的那個苗裔,如果耶穌真是這樣的家譜,耶穌就是那個苗裔,耶穌就是那位施行公平公義,稱為耶和華我們的兒的那一位大衛的苗裔,彌賽亞,所以馬太福音是給猶太人讀的,.
嚴格來說是寫給不肯信耶穌的猶太群體的,是要讓他們知道,沒錯,耶穌就是我們的苗裔,神的應許的苗裔,你可以想像得到,這個是多麼震撼的一個信息,.
所以我們基督徒要有這樣的處境,要有這樣的裝備,對舊約有一定的認識和掌握,我們就會有這種震撼感,如果不是小朋友叫他看馬太福音,或者剛剛信耶穌的,家譜很悶,.
因為跟我們沒有關係,但是我們越明白了舊約所承載的那個應許,居然現在兌現在耶穌的身上,對猶太人來說,如果他不是戴著一個有色眼鏡,其實是一件很震撼的事,居然那個木匠的兒子,原來就是神的愛子,原來就是神大衛的苗裔,.
將來帶領我們猶太人得著永恆福樂的這位救主,更加是我們外邦人的光,這個又是另外在二次元書,耶利米書都有記載的那些預言,.
所以主耶穌就是真的扮演著這樣的角色,在他的淫奈,就是在地上的事奉,職訓,就是這樣帶領人進入公義,.
所以我們信耶穌,保羅告訴我們,嚴格來說,不是因為你可愛我可愛,神愛我們,神只是有愛,其實不是的,保羅一講到救恩,是從那個公義那裡開始,是神的義藉著耶穌基督去成就出來的,.
所以救恩是彰顯的,不單純是神的愛,神的愛去到第八章的羅馬書才講的,真正講救恩是神的義,幫大家重溫,我們以往講過羅馬書,第一章,.
一章的羅馬書,所以我們講這個描繪,很重要是講的是公義和平的描繪,施行的是公平和公義,所以你會看到保羅的很重點,很緊張,他不是講愛先行,一定有,.
但是更重點是由第一章的十八節講到三章甚至四章尾,都是講公義這個課題,所以我們有時要全面掌握福音,不只是覺得很開心,有了耶穌,而是我們要從很重要的義開始,.
我長話短說,不讀這麼多,怕拆開了我們的主題,因為第七章都很想今晚有一點可以講解的,.
由十八節開始,這裡有一個很完整的段落,原來神的憤怒從天上顯明在一切不乾不易的人身上,這樣開波講救恩,原來一到第十八節,.
是講他自己為什麼要寫信給羅馬教會,他有什麼身份,他為什麼突然間會成為神的特選使徒,開始稱讚羅馬教會,.
然後進入他要講明福音的原因,因為他是一個欠債者,欠神福音的債,所以他要盡力將福音講明白,.
最好的機會就是讓羅馬教會有對救恩的規模的認識,然後更重要就是第十五章,最重要的動機就是羅馬教會,.
當保羅去到的時候,認識他,和他一起同工,將福音傳給歐洲大陸更多的未得之民,原因就是這樣寫羅馬書,.
所以由十九節開始,福音不是說神很愛地上的人,地上的人很可愛,神要愛,其實不是,在救恩裡面我們才成為可愛,.
一開始就告訴我們,原來在保羅的眼中,任何地上的人,任何一個信耶穌的人,很少,因為他很怕這個世界有神,很怕這個世界有審判,.
快點去找誰是神,我相信他,從來都很少,甚至報道會變成給你面子,說得這麼好笑,儘管舉手,這樣的信耶穌其實很不穩定,.
那可以是一個入門,可以是一個開始,打開了救恩的門,他繼續去探索,他會進入到這個階段,但是這樣不穩定,.
其實真正的穩定就是,保羅說,其實我們不要想多了,我們人類,我們現在在造物主的眼中,其實神的憤怒,.
造物主,他是自有永有,他的憤怒,他可以隨時叫一個宇宙不見了,就行了,你想想他創造的有黑洞,所有的東西,連光都被他吞噬,.
我們都不知道是什麼現象,還未能夠近距離觀察,只是天文的發現,黑洞之後是什麼呢?所有東西都被他吞噬,.
有些人說可能有白洞,去那裡再出來,這些已經超過了我們人類,可以感知的了,但是現在告訴我們,保羅很斬釘截鐵告訴我們,.
救恩其實是什麼呢?其實就是如果沒有一位公義的主,擔當我們的罪,這是三章二十三節,我們就要面對神的憤怒,.
他在天上已經顯明一切不填不二的人身上,就是哪些人呢?就是那些行不義,阻擋真理的人,.
神的事情,人本來應該知道的,但是已經失去了這種功能,我們一出生,我們是不知道做物主的存在的,沒有一個民族會知道的,.
這個就是救恩的開始,所以一直說,其實由十九節一直說到三章二十節,都是說罪人面對憤怒的神,.
我們沒有公義,我們去哪裡找公義來成為我們的庇蔭呢?成為我們的那個蔭被的亭呢?成為那個蔭被的盤石呢?.
第二十節,所以凡有血氣的,說的是地球的人類,沒有一個因行律法,因為第二章說猶太人,神的選民,.
因為可以行律法,可以在神的面前稱義,因為律法只不過是叫人知道自己是一個怎樣的人,律法不會將你拯救的,.
所以第二十一節開始,有沒有留意,救恩原來就是神的義彰顯,救恩不是神的愛彰顯,第二十一節,.
但如今神的義在律法以外已經顯明出來,有律法和先知為證,就是神的義因信耶穌基督,加給一切相信的人,並沒有分別,.
原來我們信耶穌,我們被加一個很重要的概念,就是神的義現在加在你和我的身上,原來我們一信耶穌,我們就立即獲得神的義,.
即是神已經宣告我們沒有罪,我們甚至被祂的義去遮蓋,這個義是怎樣來的,怎樣運作呢?相信大家已經完全信了主,明白了,.
但保羅在這裡要說救恩最重要的部分,是先從神的憤怒和神的義怎樣賜給信耶穌的人,.

$^{161}$那個原因就是二十三節,因為世人都是犯罪,規缺了神的榮耀,如今卻蒙神的恩典,因基督耶穌的救贖,就白白的稱義,.
神設立耶穌作挽回制,是憑著耶穌的血,藉著人的信,所以我們見到耶穌流血,我們信這件事,祂的流血就是擔待我的罪,.
這個就是信,這個信就立即帶來神的義,加在我們這些不義的人身上,所以你可以想像得到,這個就是要顯明神的義,因為祂用忍耐的心寬容人先時所犯的罪,這個忍耐的心其實就是神的愛,.
等待歷世歷代的人去決志信主,聽明白救恩,我就二十多歲信主,由我出生一直忍耐的心,忍到我二十多歲,其實很難的,.
我犯過很多罵神,得罪神,得罪人的事,連老師也罵,媽媽也惹怒,這樣忍耐,有些人再大一點,甚至臨離世才信主,一生遠離神,得不到神的,或者不覺得需要神的恩典,.
臨在病床才信主,其實失去了很多很多本來屬於祂的福氣,沒有,但是神都是這樣用愛去忍耐去寬容,所以你可以想像得到,神多想把祂的義加給不義的人,.
所以我們信主就是有一位公義的苗裔擔待了我們的罪,所以我們這樣去看猶太人的彌賽亞,我們就會很明白,一方面我們同時間也要明白猶太人為什麼這麼激氣你們外邦人,說耶穌就是你們的彌賽亞,對他們來說是很難接受,.
但是我們透過保羅寫這封信我們明白了,原來神愛的是全世界的人,不是只是愛猶太人,是每一個眾生,每一個人類有血氣的,只要我們透過神現在定的救恩的辦法,我們只需要打開我們的心靈,就立即可以得到神的義,.
所以我們要記得我們不只是得到神的愛,那個也被神的義賜要,當然不應該這樣比較的,因為神是一個全面的神,我們一擁有這位神,祂的愛,祂的公義,祂的榮耀,祂的美德都是屬於我們的,我們要宣揚神的美德,.
所以這裡告訴我們,這就是神的義,這就是信耶穌一個很重要的原位,接著其實去到第八章才真正說神的愛,當然有些偶然間也會說的,譬如第五章,說到信主有什麼的福樂也有說的,.
例如第五章的第五節,聖靈會將神的愛囂貫在我們的心裡,在患難當中我們也會歡歡喜喜,為什麼呢?基督徒如果你有了聖靈,你就可以在患難當中歡歡喜喜,這些要學的,在平安日子我們會歡喜,.
我們吃完一個自助餐,很美味,我們也會歡喜,和一班朋友我們暢談,很開心,也會歡喜,但是在患難當中,基督徒就要追求這樣的層次,這樣的一個quality的life,就是在患難中,不是我們想,但是人生你面對不同的處境,在患難當中我們都要學習歡歡喜喜,為什麼呢?.
第五章說了,第五節,盼望不至於羞恥,因為所賜給我們的聖靈將神的愛,這裡才第一次說神的愛,羅馬書說救恩,說到第五章才說一個基督徒,所以一個未信的基督徒某程度患難令他不會感受到神的愛,.
基督徒就不同了,因為已經有了聖靈,因為已經將那個心靈打開了,所以神就可以進入他的內心,其實你想一下我們基督徒應該是很高檔的,我們不需要天天都求平安,保住我,出街不准下雨,什麼都不准,裁員都不關我的事,我是被解僱的那個,.
這樣你始終有一次,如果你的人生去到一個低谷,那就很慘,你從來都沒有病過的,一病就很嚴重,但是基督徒就不是,反而我們要學習就是,原來我們在禮拜堂很開心,很平安的敬拜,是為著我們經歷神的那個牧養恩典,.
以致我們去到一些很困難的人生境遇,都不再懼怕,為什麼呢?不會覺得羞恥,這裡就告訴我們,保羅告訴我們,那個秘訣就在這裡,因為所賜給我們的聖靈,他會將什麼在我們的心靈裡面囂冠呢?.
就是用神的愛,在患難裡面就不是用神的義,是用神的愛,保羅真的分得很清楚,所以他將一位很全面的神,在整個救恩裡面,他會來講解得很清楚,所以又有神的榮耀,是我們要歡歡喜喜去盼望的,所以盼望是榮耀,.
盼望不是神的愛,盼望不是神的義,盼望是神的榮耀,所以你看到保羅講的這位做物主的神,很全面的,所以盼望我們有時間真的自己坐下,放杯咖啡在你旁邊,一口氣自己讀完整卷的羅馬書,至少由第一章十九節,一直讀到八章的最後一節,應該是三十一節,.
整個救恩明白了,因為第九到十一章是講猶太人為什麼不能夠救,十二章開始是講倫理的生活教導,所以由第一章的十九節,一直讀到八章的最後,就是救恩到底所為何事,救恩帶給我們信主的人,有什麼人生的改變,.
還有在我們的人生,我們要用什麼眼光去等候那個永恆,那個榮耀,我們是要追求的,你想一下,已經是很重要的人生功課,所以這些有時候講開,就要給大家認識,其實我們今晚講的是一個主題,是大衛的描繪,.
所以羅馬書,保羅,根本他已經完全消化了他自己是一個拉比,他對舊約的純熟認識,當他打通了任督二脈之後,就是耶穌向他顯現,他通了,他就將舊約那些的寶藏,本來就是和猶太人有關的,.
現在深入淺出,告訴外邦人,救恩就真的將舊約的描繪講得清清楚楚,主耶穌就是神的兒,主耶穌就是擔待我們罪惡的那個贖罪的制,所以我們只需要信這位耶穌,我們就得到神的兒,.
你有了神的兒之後,你就是神的子民,你就有這個身份,我們就和猶太人一樣,歡歡喜喜等候神的國度來了,接著就開始講我們在患難裡面,你經歷的,一般人經歷不到的,就是神的愛,.
透過聖靈會囂冠,囂冠在你的心裡面,淋水,會淋水,不斷的淋,所以越是經歷困難的弟兄姊妹,他如果祈禱依靠神,他的靈命是會茁壯成長的,所以我們不希望有患難出現,.
但是人生,神容許人生的不同的情況來塑造我們的人生,那就很精彩了,那些不信的人可以在我們的生命裡面來添上不同的色彩,他添上色彩,如果他是你的敵人,他沒有想過幫你塗漂亮一點的,.
他只是想幫你弄得這麼醜,但是原來這些不漂亮的東西,在神的聖靈的改變裡面,成為我們可以經歷神的愛的很重要的不可少的元素,所以那些本來想攻擊我們的人,其實是我們的恩人,如果你這樣想,這樣看,其實你的靈命又成熟了,.
因為我們根本就不會有敵人,我們的敵人就是魔鬼,所以我們要面對的不是別人對我們的攻擊,而是讓我們反而可以經歷神的愛,透過賜給我們的聖靈來去囂觀在我們的心裡,因為有了神的愛,其實有了愛,自然就有盼望,.
你在這麼患難裡面,你都感受到神的愛,感受到人間有盼望,你就會有積極的力量面對人生,所以基督徒是應該要這樣去看人生的處境,.
我們回到最後這一段,第七章還有時間的,可以講多一點的,我們要留意這個描繪去到撒加尼亞書再次有預言,再次告訴我們約書亞他現在要背起兩個職份,就是祭司和君王,.
這個描繪首先就是透過大祭司約書亞去顯明,所以對於猶太人來說,我看過一些記載,一些書,就說在耶穌的時代,我們去過昆蘭,猶大南面的山洞,有些叫死海古卷,.
其實就是當時耶穌的時代,有一個派別,猶太教都有些分支,這些某程度上都是一些異端,當時未至於這麼嚴重,這些叫愛色尼人,他們看不過眼耶路撒冷的宗教腐化,.
他們就自己遠離人世,遠離人群,就在猶大的曠野的地方,就自己組成一個退休避世的一個社團,有很嚴格的社團的條例,就要遵守的,就是黑虎,所以他們抄聖經是帶給我們後世有很大的祝福的,.
他們不斷抄寫聖經,很嚴謹的生活,這些叫愛色尼人,如果你讀教會的歷史,或者一些舊約的背景,都會有需要讀的,因為出土了很多他們的文獻,特別死海古卷有一些都是提到他們,.
所以就有這個情況,但是原來愛色尼人正正就是不接受為什麼尼賽亞會又是大祭司又是君王,說明大衛的子孫只能是君王,為什麼君王又做了祭司呢?.
所以他們就不是很明白,他們就公然的反對這個寓言的概念,我有些記載就看到,但是對於我們今天就沒有這個困難,因為主耶穌,如果你讀希伯來書,你就會很清楚主耶穌就是大祭司,他更加是先知,而且是君王,.
所以原來一個人可以背起三個職任,不過這裡就告訴我們在這裡的家面,就有兩職之間,他籌定和平,一個人,大衛的苗裔,表面就是將官面帶在約撒達的兒子,大祭司約書亞的頭上,.
就是將他加冕成為君王和祭司,好像一個實物教材,但是他當然都很奇怪,為什麼會幫我加冕呢?大祭司應該很尊榮的,但是他不明白的,包括撒迦利亞都不明白的,.
但是無論如何我們又看到這個寓言在將來,其實四百多年之後,這個時候就是五百多年,就應驗在耶穌基督的身上,耶穌基督擔之無愧的,他是我們的大祭司,是永遠的大祭司,不如又再揭一揭,我都怕大家不是很熟,.
希伯來書的第八章,先說了祭司這個概念,用了三章聖經說大祭司耶穌,你可以想像那個分佈是多麼的懸殊,第八章希伯來書,我們所說的事,.
其中第一要緊,就是我們有這樣的大祭司,已經坐在天上至大者寶座的右邊,在聖所,就是真的帳幕裡面作執事,這帳幕是主所知的,不是人所知的,凡大祭司都是為了獻禮物和祭物設立的,.

$^{201}$所以這位大祭司也必須有所獻,他如果在地上,必定不可以為祭司,所以這個就是問題了,大祭司明明是亞倫的後人,必須是利美的子孫,亞倫的後人,亞倫支派的後人,才可以成為大祭司,.
所以我們說過的薩姆爾,其實他不是大位,他不是利美的子孫,薩姆爾不能夠成為大祭司,因為他不是亞倫的後人,.
所以薩姆爾是另外一個體系的超級祭司,他肯定是先知,先知的門徒也是他開始有的,但是獻祭不是他獻,是他吩咐的祭司獻,.
所以蘇羅自己私下獻祭,不是要等到薩姆爾來,薩姆爾獻,是薩姆爾來吩咐祭司獻,薩姆爾是超級祭司,所以我們要留心這一點,.
但是這裡告訴我們主耶穌是這位永遠的大祭司,一路去討論這位大祭司,他首先如果在地上就不合資格了,.
因為舊約講明必須是利美的子孫,亞倫之派的後人才可以做大祭司,所以耶穌講明他是大位的子孫,.
他是由猶大之派出來的,一定不可以在地上做大祭司,但是現在作者告訴我們,我們在講的是在天上永遠神的真正帳幕的大祭司,.
永恆的大祭司,由第八章一直講,講到第十章,講到主耶穌不單止是祭司,他是完美的祭物,他的獻祭,他的受死,.
所以主耶穌就是這樣將我們拯救,這個就是祭司的概念,如果講君王,其實一章和第二章都有的,這個是希伯來書都有提到的,.
神記載古代的時候,藉著眾先知多次的多方面的不同的方式,來去曉喻列祖,曉喻就是預言,耶和華如此說,不是先知自己的個人的意見,.
就在這個末世,藉著他的兒子,曉喻我們,所以耶穌就是神最大最後的曉喻,那個啟示,所以如果忽略這麼大的啟示,我們就不能夠陶醉了,.
由早以納他為承受萬有,也曾藉著他創造諸世界,他就是神榮耀所發的光輝,所以是由神的本體而出的主耶穌,這個尼西亞信經的一個概念,.
從神而出的光輝,他不是受造,特別要強調他不是受造,他是神榮耀所發的光輝,他就是神本體的真相,見到耶穌就是見到神,很特別的,我們都不太明白這些概念,.
一個人,現在作者說到他超越到這樣,他是神本體的真相,是神榮耀所發的光輝,他創造諸世界,甚至納了他承受萬有,所以將來主耶穌就是承受一切的萬有,我們就和主耶穌一同作王,所以這個概念,這個就是君王開始的描述,.
然後告訴我們,他常用他的權能,權能的命令托住萬有,即是一切都是為耶穌而做的,一切都是歸於他的,倚靠他的,一切都是他來托住的,.
你可以想像到,造物主,我們的耶穌基督,他的榮耀,他的能力,他洗乾淨了人的罪,坐在高天至大者的右邊,他所承受的命,既比天使的命更加尊貴,就遠超過天使了,這樣的四字的描述,已經是一開始的序言,作者已經告訴我們,主耶穌是超越一切,.
希伯來斯也是和猶太人說的,和那些快要沒有聖殿敬拜的猶太人說,對一些認識這個世界是由神做的,我們是亞當的後裔,我們是承受亞伯拉罕的祝福,應許的,我們是大衛的子孫,我們等待彌賽亞的,.
我們在祭司的任下,我們獻祭的,小時候就要讀利未記,你們又要開始,由上主一學開始,就不是帶一元來奉獻,是要帶那些獻祭的祭僧和父母一齊上來獻祭,這樣的一個成長的,但是希伯來斯說,不可以再獻祭去聖殿,.
因為一開始四節已經將耶穌,不再說他地上的家譜,不再說他是怎樣降生伯利恆,不再說他在馬槽,不再說他跟爸爸要到埃及逃難,因為希律要殺光兩歲的小朋友,不說那些了,一開始說他那個本源,他那個本身的榮耀的身份,這樣去說耶穌,.
我們很少在教會這樣對主耶穌的敬拜,很少人去說希伯來斯,這裡已經很難解釋了,他是神榮耀所發的光輝,.
做物主的榮耀就在耶穌身上見到,誰會見到呢?見到的就是一個木匠的兒子,見到的就是一個沒有階型美容的一個出身貧寒的窮家小子,人生就很坎坷,還要是生在那剎那,.
你想一下,小時候你是住在九龍塘還是住在大坑東,如果你認識的石劍美,朱古仔,你可以想像得到你住在哪裡呢?你住在慈雲山嗎?你住在深茂坪嗎?住在藍田嗎?.
我們一聽就知道那些地方龍蛇混雜,不是呀,我住在西半山,已經不同了,已經知道那個身份的出處,所以主耶穌甘心樂意稱自己為拿撒勒人,.
所以大家不要介意自己的出身,那個是由神定的,不到我們定的,因為你生在哪個人家呢?就是神定的,後天你的努力怎樣爭奪你的原生家庭對你的影響,但是救恩告訴我們,其實不需要這樣爭奪,.
反而你應該擁抱,就是神安排你的人生,你的種子在哪裡出生,成長,將來發芽,將來來榮耀神,這個是很美麗的故事,但是正正因為我們的主耶穌原來是這樣的榮耀來到這個世界,原來是這樣選擇來拯救世人,這個描繪就很有意思,.
所以任何一個人感懷身世,你一拼在耶穌身上,你馬上就變得很有安慰,真正的安慰,因為我們的主是神榮耀所發的光輝,是神本體的真相,他來到這個世界洗乾淨人的罪,是用什麼方法呢?.
就是他自己由童貞女懷孕降生,時間對他的限制,要逐秒逐秒,那時候沒有錶的,逐秒逐秒這樣來長大,限制來的,他是沒有任何的東西可以限制他的,.
但是他要逐秒這樣成長,要去到十二歲才上聖殿,跟父母上聖殿朝拜,三十歲神要定了他三十歲出來,他十二到三十歲,這十八年就不知道去了哪裡,一點都沒有記載,.
所以有些刺經就說他去了印度,好像學魔術,有空很悶就變一些東西出來,自己在那裡玩,是那些刺經的描述,因為聖經沒有記載,那些人就在那裡幻想,他既然是神,自然就會這樣,.
就是說耶穌真的會被人笑,會被人嘲笑侮辱,真的會看不起他,你可以想像得到他要忍受十八年沉潛在拿撒勒做一個木匠的兒子,沒有東西可以改變這個環境,因為神的時間還沒有到,他就要等待,.
等十八年,什麼青春都過了,才出來做傳道三年多,就要馬上離世,不可以展露他的才華給世界的人,就只是一小撮的人知道他的身份,要很漫長的二千年後才來到我們香港,知道有耶穌這個救恩,.
這樣迂迴的,不需要用一個Zoom的直播,全世界的人就可以馬上看到,是要經過很多很多的時間,這個就是公義的描繪,所以在舊約已經有那個寓言,國破家亡的時候,那個寓言都沒有收回,.
給耶利米就是說,沒錯,我們看到的是國家因為人民的離開神,是必須不可以再改變的了,哪個禱告都不可以改變的了,是必定要受懲罰去到這個巴比倫,但是七十年就看你有沒有信心了,.
七十年你在那裡買田買地,殷勤的準備,囤積,致富,七十年之後就可以重歸聖城重建,很少的百姓是有這樣的信心,大部分的都是帶著眼淚離開,被逼離開,很多都是不想離開,.
耶利米在那個離開的行列,不過是被人抓去一個叫做,突然間忘記了那個領袖的名字了,抓了他去,後來就好像說在埃及死了,耶利米,他都沒有去到巴比倫了,.
因為又是那些當時的將軍,就不信了,又要逃難,就抓了耶利米做那個人質,這樣帶他,後來那個將軍都被巴比倫殺了,巴比倫人就很敬重耶利米,這些是後期耶利米書的記載,好像是五十二章的記載,這樣來告訴我們那些事情,.
去到被擄歸回,再一次在第八個二章,就有這一番的預言,再確定那個苗裔擔起的就是君王和大祭司這個身份,.
他來到去可以在兩職之間,兩職之間,籌定和平,和就是,我翻到這個馬拉基書的第,對不起,在撒迦利亞書先,第六章這一段,.
兩職之間籌定和平,這個冠冕要歸希臘,多比亞,耶大亞和西班牙的兒子,放在耶和華的殿中為紀念,就是說這個榮耀是放在神的殿裡面,作為一個紀念,這個殿還沒有建好,很快就有那個殿了,.
那個冠冕真是,相信撒迦利亞真是做了一個冠冕,就將這個冠冕放在神的殿裡面,作為一個紀念,就是說神會對他自己的應許,這個預言會來對言,.

$^{241}$所以真正敬虔的猶太人,他們讀明白這些舊約的應許,所以很多人是等待那個彌賽亞,還更加奇妙,就是差不多去到主耶穌的時候,原來很多的猶太人是等候那個誰是那個幸運兒,什麼幸運兒,.
同女懷孕降生,所以他們就一直渴望誰的大姐會懷孕這個彌賽亞,猶太人是這樣想的,真是敬虔的,他們知道彌賽亞的時間到了,.
所以我們就是一個回望式,因為彌賽亞已經來了,救恩已經傳到我們中國,來到我們香港,2000年了,就是說現在我們要等的是彌賽亞的第二次降臨,.
我們要這樣的心情,猶太人還是第一次等待彌賽亞的降臨,所以他們現在還活在他們困難的處境,所以周圍的國家都是對以色列有這樣的敵視,.
再加上你自己的地方有一幫人是你在高壓的統治著,真是一個很大的困局,所以我們等待的應該是用猶太人的等候彌賽亞的心情去活我們地上的人生,.
我們真是有一個等待的,就是彌賽亞第二次降臨,這個是新約的聖經應許的,不是我們以為我們扮自己是猶太人,其實不是的,.
真是新約告訴我們主耶穌是第二次會回來的,這個第七章我們下個星期再開始,因為開始了有些東西要教的,.
我們最後看這一段,就是這個帖撒羅尼加前書的第四章,說到彌賽亞第二次降臨,也是這個新約的應許,帖撒羅尼加第四章,.
(主禮物).
耶和華如此說,保羅這麼肯定地說,.
我們者活著,還存留,到主降臨的人,斷不能夠在那已經睡了的人之先,因為主必親自從天降臨,.
有呼叫的聲音,和天使掌的聲音,又有神的號吹響,那在基督裡面死了的人,就是睡了的人,就首先復活,以後有一個身體,有一個形體,是怎樣的就真的不知道了,.
有一個可以去到天堂的形體,或者榮耀的身體,這樣說我們在地球暫時未必經歷到,我們所謂的有時見到鬼魂,鬼怪,那些所謂靈界,靈體,那些都一定不是將來的榮耀,.
所以將來不會有鬼的,放心,不會有靈界的,因為必先復活,一定有一個榮耀的身體,復活的身體,所以那些死了的人,其實本身已經沒有身體的,.
說明死了,不知道變灰還是變什麼,你不要說五百年,你過了三十年,你的名字都沒有人記得你,何況是幾千年,根本不會有人認識你的,但是每個人都會按照自己的呼叫,就會復活,.
所以這是一個很大的安慰,對人類,對我們基督徒一個很大的安慰,那些死了的人是首先復活,然後以後就是那一代人,耶穌回來的時候還沒死的那一代人,.
我們者活著還存留的人,必和他們,就是那些死了的,一同被提到雲裡,在空中與主相遇,這樣我們就要和主和那些人永遠的同在,.
所以你們當用這些話彼此勸諱,說明彼此勸諱,就是說這個是預言,這個是真實的啟示,這個真是人類一定按字面,這樣來按次序的復活,.
還有最大前提是這位復活的耶穌會第二次再回來,就是我們每一個基督徒的盼望和等候,所以這個是很重要,對於我們看這位公義的描繪,.
這位彌賽亞第一次降臨,他已經開展公平公義的治理,就是說他的救恩帶來的地上可以得到神的義,可以得到平安,因為是與神相和,第五章羅馬書,恩順清義帶來的福氣,.
第五章大家會不會背,很早期我初信的時候就背羅馬書,我們的恩順清義就直著我們的主耶穌基督得以與神相和,就是Find Peace in Our God,恩順清義的人立即是與神相和,.
當然就是神的義,直著耶穌加給我們,我們都不知道,我們知道的就是第一時間你在神裡面,因為這個和好,我們就能夠有這份平安,我們又直著他,恩順得以進入現在所站的恩典之中,其實是恩寵,.
我們在一個Favor裡面,我們Find Favor in God,直著耶穌基督得以進入現在所站的恩典之中,並且歡歡喜喜,盼望神的榮耀,整個人生原來就是一個榮耀的盼望,就是救恩帶給我們,不單止是這樣,就算在患難當中,我們都是歡歡喜喜,.
因為患難生忍耐,忍耐生老練,老練生盼望,盼望不至於羞恥,為什麼不需要羞恥呢?因為剛才所說的第五節,所賜給我們的聖靈,將神的愛囂貫在我們的身上,.
所以有時候這些小時候背好了經文,有時候坐車,默想,有時候人生有些逆境,你背一下,那些經文就會安慰你的心靈,所以如果不讀經,你自己儲藏了神的話太少,你就會靠自己的記憶,靠自己的經驗,每個人都有經驗,.
但是那個經驗是人間的經驗,我們沒有神而來的心靈的安慰和平安,盼望這個也是我們去查經班,大家一起仔細去讀經的原因,將我們不同的經驗彼此分享,.
我們祈禱了,感恩我們在今晚能夠很有系統的去看我們這位公義的苗裔,原來很早神已經預言,藉著大衛的子孫,這位苗裔要來去成為以色列的拯救,.
阿們..
\newpage



\section{馬可福音 5:21-43}
\label{sec:HWMu4dld7_M}
\textbf{2023-10-15 加恩樓《這是一個怎樣的信心》馬可福音 5\_21-43 講員:老耀雄牧師}
\newline
\newline
連結: \href{https://youtube.com/watch?v=HWMu4dld7_M}{\texttt{https://youtube.com/watch?v=HWMu4dld7\_M}} ~~~~ 語音日期: 2023-10-18
\newline
\newline
\hyperref[sec:wwGm5m9DNT8]{\small{< < < PREV SERMON < < <}}
~
\hyperref[sec:index_chronic]{\small{[返順時目]}}
~
\hyperref[sec:index_scriptual]{\small{[返順卷目]}}
~
\hyperref[sec:cM9RUvSdlG8]{\small{> > > NEXT SERMON > > >}}
\newline
\newline
馬可福音 5:21-43
\newline
\begin{longtable}{cl}
\hline
\hline
章節 & 經文 (和合本修訂版)\\
\hline
5:21 & \begin{tabularx}{0.7\textwidth}{X} 耶穌又坐船渡到對岸，有一大群人聚集到他身邊；他正在海邊。 \end{tabularx} \\ \\ \relax
5:22 & \begin{tabularx}{0.7\textwidth}{X} 有一個會堂主管，名叫葉魯，也來了，一見到耶穌，就俯伏在他腳前， \end{tabularx} \\ \\ \relax
5:23 & \begin{tabularx}{0.7\textwidth}{X} 再三求他，說：「我的小女兒快要死了，求你去為她按手，使她痊癒，可以活下去。」 \end{tabularx} \\ \\ \relax
5:24 & \begin{tabularx}{0.7\textwidth}{X} 耶穌就和他同去。有一大群人跟隨他，擁擠著他。 \end{tabularx} \\ \\ \relax
5:25 & \begin{tabularx}{0.7\textwidth}{X} 有一個女人，患了經血不止的病有十二年， \end{tabularx} \\ \\ \relax
5:26 & \begin{tabularx}{0.7\textwidth}{X} 在好多醫生手裡受了許多苦，又花盡了她所有的，一點也不見好，反而更重了。 \end{tabularx} \\ \\ \relax
5:27 & \begin{tabularx}{0.7\textwidth}{X} 她聽見耶穌的事，就夾在眾人中間，從後面來摸耶穌的衣裳， \end{tabularx} \\ \\ \relax
5:28 & \begin{tabularx}{0.7\textwidth}{X} 因她想：「我只摸到他的衣裳，就會痊癒。」 \end{tabularx} \\ \\ \relax
5:29 & \begin{tabularx}{0.7\textwidth}{X} 於是她的流血立刻止住，她覺得身上的疾病好了。 \end{tabularx} \\ \\ \relax
5:30 & \begin{tabularx}{0.7\textwidth}{X} 耶穌頓時心裡覺得有能力從自己身上出去，就在眾人中間轉過來，說：「誰摸我的衣裳？」 \end{tabularx} \\ \\ \relax
5:31 & \begin{tabularx}{0.7\textwidth}{X} 門徒對他說：「你看眾人擁擠著你，還說『誰摸我』呢？」 \end{tabularx} \\ \\ \relax
5:32 & \begin{tabularx}{0.7\textwidth}{X} 耶穌周圍觀看，要見做這事的女人。 \end{tabularx} \\ \\ \relax
5:33 & \begin{tabularx}{0.7\textwidth}{X} 那女人知道在自己身上所成的事，就恐懼戰兢，來俯伏在耶穌跟前，將實情全告訴他。 \end{tabularx} \\ \\ \relax
5:34 & \begin{tabularx}{0.7\textwidth}{X} 耶穌對她說：「女兒，你的信救了你，平安地回去吧！你的疾病痊癒了。」 \end{tabularx} \\ \\ \relax
5:35 & \begin{tabularx}{0.7\textwidth}{X} 耶穌還在說話的時候，有人從會堂主管的家裡來，說：「你的女兒死了，何必還勞駕老師呢？」 \end{tabularx} \\ \\ \relax
5:36 & \begin{tabularx}{0.7\textwidth}{X} 耶穌不理會他們所說的話，就對會堂主管說：「不要怕，只要信！」 \end{tabularx} \\ \\ \relax
5:37 & \begin{tabularx}{0.7\textwidth}{X} 於是他帶著彼得、雅各和雅各的弟弟約翰同去，不許別人跟著他。 \end{tabularx} \\ \\ \relax
5:38 & \begin{tabularx}{0.7\textwidth}{X} 他們來到會堂主管的家裡，耶穌看到一片吵鬧，並有人大聲哭泣哀號， \end{tabularx} \\ \\ \relax
5:39 & \begin{tabularx}{0.7\textwidth}{X} 就進到裡面，對他們說：「為甚麼大吵大哭呢？孩子不是死了，是睡著了。」 \end{tabularx} \\ \\ \relax
5:40 & \begin{tabularx}{0.7\textwidth}{X} 他們就嘲笑耶穌。耶穌把他們都趕出去，帶著孩子的父母和跟隨的人進了孩子所在的地方， \end{tabularx} \\ \\ \relax
5:41 & \begin{tabularx}{0.7\textwidth}{X} 就拉著孩子的手，對她說：「大利大，古米！」翻出來就是說：「女孩，我吩咐你，起來！」 \end{tabularx} \\ \\ \relax
5:42 & \begin{tabularx}{0.7\textwidth}{X} 那女孩子立刻起來走動—她已經十二歲了；他們就非常驚奇。 \end{tabularx} \\ \\ \relax
5:43 & \begin{tabularx}{0.7\textwidth}{X} 耶穌切切地囑咐他們，不要讓人知道這事，又吩咐給她東西吃。 \end{tabularx} \\ \\
[1ex]
\hline
\hline
\end{longtable}
$^{1}$大英姐妹平安.
根據我的統計和資料顯示.
我已經有五年沒有在家欣樓分享了.
所以很高興這次再來家欣樓和大家分享.
特別記掛著家欣樓的大英姐妹.
為什麼呢?.
在我三月患病的時候.
我記得家欣樓都為我拍了一條片.
為我打氣.
證明大英姐妹都是關心我的.
我都深繫家欣樓.
稍後再和大家分享我的康復之路.
以及我在這個病中的神學反省.
開始的時候我們有一個祈禱.
親愛的恩主.
讓我們能夠從他們當中一朝你的話.
可以學習不同的肅靈教導.
從你的教導裡可以明白你的心意更多.
好叫我們能夠按著心意而行.
讓你名被高舉.
我們祈禱.
去奉教主耶穌基督得性名字祈求.
阿們.
我們先了解一下經文的背景.
剛才已經有人為我讀出了23節經文.
血柳婦人的故事其實只有10節.
由25節開始到34節就結束了.
另外血柳婦人的故事.
也是由兩組經文包著.
分別是21節到24節的上文.
以及35節到43節的下文.
所以血柳婦人本身不是一個獨立的故事.
而是依附著上下文的一連串的事件.
《色經學者》就將一種寫作題材.
叫做三文治式的結構.
所以看血柳婦人的故事.
不應該單獨去看.
還要看上下文的故事.
否則就會錯過了作者三文治式的結構心思.
故事上是講什麼呢?.

$^{41}$在一個管會堂的業老.
他再三請求耶穌去救他的女兒.
因為他的女兒快要死了.
在耶穌趕去業老家的途中.
就發生了血柳婦人的事件.
最後血柳婦人被治好了.
下文就說.
婦人被治好的同時.
業老的女兒就死了.
耶穌應該不需要去.
不過耶穌沒有理會.
繼續前去.
他進入了業老的家.
趕出所有人.
最後將業老的女兒從死裡拯救出來.
整個故事涉及了兩組人物.
兩件事件.
兩件事件都涉及信心和行為.
無論是婦人的血柳病被治好.
還是女孩死而復生.
很明顯地讓讀者知道.
兩件事都是神職醫治.
業老的行為很明顯地.
他在眾人面前.
向耶穌再三請求.
顯示他對耶穌醫治的信心.
婦人的行為是隱藏的.
他靜靜地走向耶穌後面.
沒有人知道婦人的企圖和信心.
兩件事都指向不同的性別.
不同的身份.
業老是一個男性.
有姓名.
管會堂的.
是一個有身份.
受人尊重的人物.
婦人是女性.
沒有被記載.
更加沒有身份和地位.
兩人的身份和地位都截然不同.

$^{81}$可以這樣說.
一個是有身份.
有地位.
有名聲.
受人尊敬.
另一個沒有身份.
沒有地位.
根本是一個被社會拋棄的弱勢社群.
兩件神職被捆綁在一個沒有間段的故事裡.
究竟作者想我們明白一個甚麼訊息呢.
我們開始進入經文.
管會堂的業老.
他的女兒病了.
雖然經文沒有說這個女孩.
她是生了甚麼病.
或病了多久.
但經文說她快要死了.
女孩的生命與死亡很接近.
一個有身份有地位的業老.
他應該找過很多醫生去醫治他的女兒.
但始終都是徒勞無功.
女兒現在病了.
一點喜悉都沒有.
在這一刻.
業老有一份絕望的感覺.
他不知道可以為女兒做些甚麼.
當他知道耶穌在海邊的時候.
他就走過去.
要見耶穌.
雖然經文沒有交代到.
業老為甚麼會對耶穌這麼有信心.
但當他看了三文治結構裡.
中間那一層的經文.
我們就知道了.
他一見到耶穌.
就俯伏在耶穌的腳前.
再三求祂.
他說我的小女兒快要死了.
求你去為她按手.
使她痊癒.

$^{121}$可以活下去.
業老對女兒病情的無助感.
推動他去見耶穌.
當他見到耶穌的時候.
他就請求耶穌的幫助.
耶穌成為他在絕望裡的唯一盼望.
耶穌對於業老的請求.
亦都欣然接受.
業老的女兒雖然現在仍然未醫好.
但業老這一刻.
應該是充滿了信心和盼望.
於是我們一起在情.
回到業老的家裡.
想醫治業老的女兒.
不過故事的發展.
被另一件事打斷了.
原來一個患了血漏的婦人.
在人群當中.
她患了血漏病已經12年了.
12年看了很多醫生.
開始醫不好.
花盡了她所有的.
不但沒有好轉.
病情反而更重了.
這12年.
婦人曾經有個盼望.
盼望有醫生可以醫好她的病.
但失色了.
一次又一次的失望.
財富沒有了.
病情卻越來越嚴重.
不單止是這樣.
血漏是一個不潔的象徵.
所有和婦人接觸過的人.
或者物件.
都需要潔淨.
因此沒有人願意接近婦人.
婦人在日常的社交.
宗教的生活裡.
都是被排擠的.

$^{161}$可以這樣說.
是一個被社會遺棄的弱勢社群.
婦人在這個情況下.
應該陷於絕望中.
經文說她聽見耶穌的事.
馬可福音第一章.
記載了耶穌醫好患麻風的人.
在第二章.
耶穌醫好一個癱子.
在第三章.
耶穌醫好一個萎縮手的人.
在第四章.
耶穌平靜風和浪.
在第五章.
耶穌在加拉森.
趕出了二千隻污鬼.
因此耶穌的名聲已經被傳開了.
婦人聽見耶穌的事.
孽瘤也同樣聽見耶穌的事.
他們都是從別人裡面所傳來的訊息.
或者故事裡面.
知道耶穌這個人.
他們知道耶穌是大有能力的.
也曾經醫治過很多絕望的人.
他們都希望藉著耶穌的能力.
病患可以得到醫治.
就像那些患麻風的.
癱子的.
被鬼附的人一樣.
能夠得到耶穌的憐憫.
以及醫治.
這一刻.
婦人和耶穌都在人群裡面.
兩人的距離是這麼近.
基督能對婦人所想的.
作出了一個很具體的描述.
她說:我只要摸到他的衣裳.
就會痊癒.
這是一個什麼信心呢?.
這個信心不需要向耶穌請求.

$^{201}$這個信心不需要經過耶穌的同意.
這個信心只需要摸到耶穌的衣裳就成就了.
於是.
婦人真的從後面摸到耶穌的衣裳.
神蹟就發生了.
經文說.
婦人的血流止住了.
疾病就好了.
困擾了12年的病.
也困擾了12年的心靈壓抑.
這一刻.
都得到釋放.
對於婦人來說.
這是件多麼興奮的事.
這個信心超越了一般人的想像.
婦人認為.
摸到衣裳就可以了.
她沒有經過耶穌的同意.
婦人以這種信心.
去支持她摸耶穌衣裳的行為.
而這個信心卻成就了婦人所期望的.
這個信心也印證了耶穌的大能.
縱使耶穌處於不知情的情況下.
耶穌在事情開始的時候完全是被動的.
婦人未經過耶穌同意.
就主動去摸耶穌的衣裳.
不過.
耶穌縱然是被動.
但無礙於她欣點的施慾.
當耶穌發現有能力從自己身上出來的時候.
耶穌就轉為主動.
去向眾人問.
誰摸我的衣裳.
門徒認為在這個擠擁的場合.
被人捉摸得到.
是一件很自然不過的事.
為什麼耶穌還在意誰摸我.
不過耶穌定義要知道.
誰做這件事.
婦人沒有逃跑.

$^{241}$她知道逃避不了.
於是就俯伏在耶穌面前.
經文形容婦人是恐懼戰驚.
她恐懼戰驚地將事實告訴了耶穌.
婦人感到恐懼戰驚.
是不是因為很怕被耶穌責罵呢.
事實上她的確沒有得到耶穌的同意.
就是摸她的衣裳.
她是不是又擔心耶穌會收回她的能力.
讓她的病不能夠康復呢.
耶穌的回應不是對婦人的斥責.
而是讚賞和安慰.
他說:女兒,你的信救了你.
平安地回去吧.
你的疾病痊癒了.
婦人本來是無形無性的.
是被社會所不略的.
但現在卻有一個尊貴的女兒命份.
她以前是隱藏地去到耶穌後面.
不想被人發覺.
怕被人接觸.
但現在她光明正大.
在耶穌和眾人面前.
從容地離開.
她的病不單單對耶穌的信心被治好了.
她的信心也令她的罪獲得赦免.
得到拯救.
因為耶穌不是單單對婦人說.
你的病好了.
而是加多了一句.
你的信救了你.
婦人的信心讓她獲得拯救.
這個拯救不單單是血流病.
也是生命的拯救.
在猶太人的傳統.
身體上的殘缺.
都是犯罪的懲罰.
唯有罪得赦免.
身體才能夠康復.
因此耶穌對婦人的赦罪.

$^{281}$也讓婦人的身心靈都得到康復.
婦人的身心靈被醫治.
應該不會再感到恐懼.
她的人生應該充滿了感恩和讚美.
她在絕望中帶著隱藏的信心而來.
然後匍匐在耶穌面前.
恐懼戰驚地向耶穌說了她的故事.
最後帶著神的恩典和歡欣的心離開.
血漏婦人事件的結束.
卻是另一個信心故事的開始.
有人從接女兒的家來到.
跟著她說:你的女兒死了.
何必還勞駕老師呢?.
有人指出女兒已經死了.
耶穌也無需來.
因為一切都完結了.
死亡將一切的盼望都終結了.
在這一刻熱路有什麼反應呢?.
因為有身份,有地位,有名聲.
負責管會堂的人.
此刻阿爵無聲.
一言不發.
都不知道該說什麼好.
甚至在這個三文治式的結構裡.
熱路在下半部的故事裡.
連說一句台詞的角色都沒有.
好像完全消失了.
在這個信仰故事裡.
耶穌卻主動跟熱路說.
不要怕 只要信.
耶穌這句話正正反映.
熱路當時的心態.
就是非常之怕.
以及失去信心.
女兒一死.
令他感到害怕.
女兒死去的消息.
令他失去一切的信心.
眼前的耶穌.
是否能夠令他的女兒死而復生呢?.

$^{321}$他沒有信心.
這就是他在下半部的故事裡.
完全沒有角色的原因.
耶穌沒有理會熱路的反應.
他帶著三個親密的媒徒.
就是彼得,雅各,洪逸翰.
繼續前往熱路的家.
當時家裡的人為了女兒的死而難過.
眾人都在哭.
耶穌的反應就是.
為什麼哭呢?為什麼大吵大哭呢?.
孩子不死了 誰著了.
眾人都嘲笑耶穌.
耶穌就把他們趕出去.
進入了女孩的房間.
然後對女孩說.
女孩 我吩咐你起來.
這句話充滿了權柄.
叫人不得不聽從.
不單人要聽從.
連死亡都要聽從.
女孩就立即起來.
恢復了生命的氣息.
正如耶穌所說.
她只是睡著而已.
現在睡醒了.
馬上可以起來行走.
好像沒有病一樣.
經文透露.
他們非常驚奇.
這個女孩已經十二歲了.
他們驚奇當然是因為女孩能夠走動.
沒有死去.
也因為耶穌的權柄和大能.
竟然可以令死去的人.
恢復生命氣息.
他們因為見證到一個神蹟而驚奇.
不過他們除了驚奇之外.
是不是應該對耶穌有一份感恩之情呢?.
是不是對耶穌身邊的身份.

$^{361}$有一份敬畏之心呢?.
對於業佬來說.
會不會因為自己失去了對耶穌的信心.
而感到慚愧或是懊悔呢?.
血路婦人患了病十二年.
女孩也是十二歲.
這是巧合還是作者另有心意呢?.
經文沒有解釋.
不過無論是婦人十二年患病的苦痛日子.
還是女孩十二年的無憂無慮的青春歲月.
她們也在絕望中遇到耶穌.
在她們遇到耶穌之後.
生命也經歷了極大的變化.
婦人十二年的病患被釋放.
罪得赦免.
女孩的病被治好.
生命得以延續.
有美好的人生.
她們不但同樣被耶穌視為女兒.
有著尊貴的身份.
同樣在未來的世界中.
也有一段難忘和不一樣的人生.
在兩個故事中.
葉老主動尋找耶穌.
耶穌主動回應.
就算旁人認為女孩已經死了.
事件也告一段落.
就算葉老知道女兒死了之後.
沒有進一步的請求.
甚至信心全失.
就算被葉老家人嘲笑.
被人認為這是風言風語.
人死了還當作睡覺.
但耶穌對女孩的憐憫.
沒有因別人的誤會或無知.
而阻礙了祂對女孩的請求.
也不會因著葉老的信心而放棄.
血肉婦人也是主動尋找耶穌.
但卻隱藏著信心的行為.
令耶穌被動地回應.

$^{401}$就算婦人沒有得到耶穌的同意和許可.
就算她沒有耶穌的依償.
偷偷摸摸地進行.
但對於需要耶穌拯救的人.
耶穌對人的憐憫.
超越了人的一般信心.
因此無論是主耶穌主動的.
或被動地回應人的請求.
無論人是否失去了信心.
或信心的行為沒有得到耶穌的同意.
但對於需要被拯救的人.
對於耶穌所憐憫的人.
耶穌仍然樂以賜予豐富的恩典.
這就是我們所相信的.
以及所爭取的三一真神的屬性.
弟兄姊妹們.
從上述兩件神蹟的依字.
我們看到一個有身份有地位的葉老.
在他感到絕望的時候.
他主動尋找耶穌.
所以耶穌為他的女兒醫病.
他呼服在耶穌面前.
只要耶穌願意為他的女兒按手.
女兒就可以痊癒.
葉老向眾人顯示出他信心和誠意.
不過當有人說他的女兒已經死了.
葉老就信心全失.
但無礙耶穌拯救的恩典.
葉老縱然從有信心到失去信心.
但神的恩典仍然臨到.
我們看到一個沒有身份沒有地位.
連名字都沒有的婦人.
婦人也很絕望.
她用盡她所有的.
但情緒更加嚴重.
她主動尋找耶穌.
所以耶穌醫好她的血流病.
她隱藏著的動機.
靜靜跟著耶穌後面.
她的信心很特別.

$^{441}$她不需要耶穌的同意.
也不需要耶穌的應許.
她認為只要摸到耶穌的衣裳就可以了.
最後她的信心得以成就.
婦人信心的行為雖然未被耶穌認可.
但神的恩典仍然臨到.
兩種信心.
是否可以作出一些比較呢?.
婦人被醫治的過程中.
信心好像未被考驗過.
未被挑戰過.
起碼沒有被死亡打擊過.
相反,葉老的信心.
卻被女兒的死亡挑戰過.
以致信心全失.
普遍人都認為.
讓死人復活的信心.
應該比病得醫治的信心更大.
這樣才合理.
事實上,耶穌醫治神蹟.
施行神蹟醫治.
不在乎人的信心大小.
只在乎主耶穌的憐憫.
如何應用這段經文呢?.
是非常困難的.
教會的教導就是.
我們一定要對神有信心.
神蹟是會幫助我們的.
在葉老的事件中.
在女兒死後.
神真的喪失了信心.
我們是否可以在人死後.
仍然很有信心地宣告.
這個死去的人可以再次站起來呢?.
現實上我們很難做到.
婦人的信心行為.
未經過耶穌的同意.
是否我們的信心.
不需要主耶穌的同意.
我們都可以成就呢?.

$^{481}$兩種信心.
在現實生活中.
都遇上困難.
或者我們可以這樣去形容.
如果我們真的一點信心都沒有.
神就一定幫不了我們.
不過信心是必須的.
多大的信心才能成就呢?.
這個卻不是必然的條件.
信心越大.
神當然會越立.
但信心微小.
神的憐憫和應許.
卻可以補足我們信心的不足.
今天講題是.
這是一個怎樣的信心?.
今天我們如果要尋求耶穌.
不論我們的身份,地位或性別.
耶穌同樣會接納我們.
當我們要尋求耶穌的幫助的時候.
我們都需要對耶穌有信心.
不過耶穌為我們所成就的東西.
與我們的信心大小沒有必然的關係.
最重要的是主耶穌的憐憫.
只要是神所揀選的.
是神要定義請求的.
神所做的功必定會成就.
我們知道神是一個完美的神.
但神卻沒有要求信徒同樣的完美.
神只要求我們信徒.
盡心盡性盡義盡力.
每個信徒所能盡的能力和效果都不一樣.
所以不能以成果作標準.
而是以信徒所盡的心思和行為.
信徒是否盡了心思與實踐呢?.
這個主是一定會知道的.
業老以及血流婦人.
在我們絕望的時候.
他們已經盡了信心和能力.
或者業老最後喪失了信心.

$^{521}$白海婦人的行為未經主耶穌的同意.
但主耶穌看到他們已盡的心.
耶穌就憐憫他們.
補足了他們的信心.
因此今天我們不是要比較信心的大小.
而是要審查自己對主的愛.
是否已經盡了一切的心思意念.
我今年與神的關係有一個很大的突破.
大家知道我在3月的時候.
因為嚴重的腹膜炎.
大腸破裂了.
需要做緊急手術.
切除了大腸的破裂地方.
根據當時的醫護界評估.
死亡率70\%.
生存率只有30\%.
當然最後是神仙救了我.
否則我都不會站在這裡.
我之後就有一個領袖.
每個人都會受到死亡的威脅.
尤其是我做牧師.
很多時候都要和別人做床邊灑泥.
灑完泥之後我當然要安慰他們.
不用怕 主耶穌與你同在.
我問問自己.
別人面對死亡的時候你這樣安慰別人.
你自己面對死亡的時候會怎樣?.
所以我自己覺得.
我這個病就是要印證.
究竟我在死亡的時候.
我的心裡面是不是真的不怕.
事實上我的心裡面很平靜.
在過程中都沒有怕過.
我也感受到神與我同在.
所以我覺得.
這次手術是一個恩典之旅.
是神以這個經歷回應我的提問.
也是我和神的關係的一種肯定.
我感受到這次患病所得到的屬靈果子.
比我失去的更多.

$^{561}$我初以為整件事已經告一段落.
安心養病.
養完病就繼續事奉.
誰知在一次組織化驗之後.
又確診患上了淚癌和血癌.
是一種蛋白質不能分解的病.
會引致內出血和產生血泡.
這些血泡會根據身體不同的器官.
造成嚴重的傷害.
經過醫生的判斷.
盡快接受標靶治療.
我現在已經在接受標靶治療.
進入第二期.
我再和問神說.
這次我要學一個什麼功課.
我領受的就是.
與主同在的經歷.
不只是手術前後.
不只是一晚.
而是整個信仰的生命.
正正是這個病的不確定性.
危險性.
讓人感到憂慮,牽掛.
以及死亡的溫血.
導致我們可能會埋怨神.
或者對神產生懷疑.
而我領受的就是.
不要讓這些不確定性的東西.
影響我和神的關係.
影響我事奉的生活.
要我經歷更多.
更依靠神和神的關係.
更親密的肅靈之旅.
在我患病期間.
繼續可以讓我講道.
讓我帶團契.
讓我有事奉的機會.
我現在正在做標靶和化療.
副作用不是很大.
很多人如果在治療期間.

$^{601}$副作用很大的話.
基本上做不到什麼事奉.
你看到我今天都精神奕奕.
雖然有時候舌頭的關係.
說話不是很清晰.
但起碼可以講到一篇道.
站了這麼久.
所以我整個肅靈生命.
所需要的.
不單單是要對神有信心.
更加需要的是.
神對我們憐憫.
我們在肅靈生命.
固然要對神有信心.
要靠和仰望上主.
同樣都需要神的看顧和憐憫.
而獲得神眷顧的唯一方法.
就是主和主人保持一個親密的關係.
這個關係就像葡萄樹與枝子一樣.
能夠在主身上.
支取源不絕的動力.
才可以在逆境中.
持守信仰.
否則信仰所淺.
都是空談.
或者只是一個口號.
我們一起祈禱.
三位上主我們要稱頌您.
因著您的慈愛憐憫充滿我們.
您的愛彌補我們信心的不足.
您的愛遮蓋著我們所犯的罪.
您的愛讓我們生命全然成聖.
您的愛使我們可以與您復和.
也可以與別人復和.
您的愛讓我們進入神的致勝所.
與您相遇.
在您面前.
我們生命可以有順忘愛.
順便我們每一個屬您的.
都在您的愛中脫鑿成長.

$^{641}$成為一個合您心意的僕人.
被您使用.
讓您的名被廣傳.
被高舉.
讓我們能夠一生去跟從您.
我們祈禱.
是奉教主耶穌基督得勝名字祈求.
阿們.
願主祝福大家.
\newpage



\section{以弗所書 1:17-23-4:7-16}
\label{sec:cM9RUvSdlG8}
\textbf{2023-10-15 早堂《尊基督為元首的門徒相交》以弗所書 1\_17-23,4\_7-16 講員:曾裔貴牧師}
\newline
\newline
連結: \href{https://youtube.com/watch?v=cM9RUvSdlG8}{\texttt{https://youtube.com/watch?v=cM9RUvSdlG8}} ~~~~ 語音日期: 2023-10-19
\newline
\newline
\hyperref[sec:HWMu4dld7_M]{\small{< < < PREV SERMON < < <}}
~
\hyperref[sec:index_chronic]{\small{[返順時目]}}
~
\hyperref[sec:index_scriptual]{\small{[返順卷目]}}
~
\hyperref[sec:c_6waYrNYvI]{\small{> > > NEXT SERMON > > >}}
\newline
\newline
$^{1}$弟兄姊妹早晨.
一起同心祈禱.
感恩我們再一次聚集.
有剛才的詩歌敬拜.
亦有我們獻上今早的頌禱.
我們等候在神的面前.
讓我們一班弟兄姊妹.
認識我們自己的身份.
認識主耶穌的救恩.
帶給我們的拯救.
以致我們能夠在教會.
同心去事奉.
我們要長大成人.
滿有基督長成的生亮.
我們這樣禱告.
奉耶穌基督的名.
阿們.
為什麼要尊基督為元首的相交.
門徒的生活.
在這段聖經.
很清楚地說了.
小徒保羅對於以弗所教會的要求.
其實這是神的旨意.
任何一個稱為神名下的教會.
都要朝這個目標設立教會.
就是要成熟.
長大成人.
你不會生個嬰兒.
你去到五歲就不要再長大了.
我沒錢養你了.
不會的.
他會自然地生長.
長大之後.
他各方面的身段.
和他整個人的標準的配合.
這些全部要均衡健康地成長.
教會也是一樣.
誰是頭的元首呢.
就是第一章第三章.
保羅就用了兩個很長的祈禱.

$^{41}$去帶進一個很重要.
給教會每一個弟兄姊妹認識的.
一個最大的真理.
就是原來耶穌.
他不單只是神的兒子.
他的受死.
他的埋葬.
和他的復活升天之後.
他成為了地上教會每一個弟兄姊妹.
那個元首.
那個頭.
就是說原來教會就是他的身體.
在地上讓人見得到的身體.
我們是基督的身體.
我先和大家看看第一章的經文.
有一個祈禱.
第十五節開始.
很快地將這個祈禱讀給大家聽.
因此我既聽見你們是順從主耶穌的人.
又親愛眾聖徒.
禱告的時候.
上上提到你們.
祈禱開始了.
求我們主耶穌基督的神.
榮耀的父.
將賜人智慧和啟示的靈.
賜給人有屬靈的智慧.
賜給人對啟示有明白的這位聖靈.
來賞賜給你們.
使你們能夠真知道他.
他就是上文所說的主耶穌基督的父神.
第十八節.
並且照明你們心中的眼睛.
心靈的眼睛.
信了主耶穌的弟兄姊妹.
有聖靈在內心.
現在要照明.
Enlighten你的內心.
讓你能夠照明.
你要對一些事情有認識.

$^{81}$什麼事情呢?就是永恆.
並且照明你們心中的眼睛.
使你們知道他的恩兆.
有何等的指望.
恩典的選兆.
有指望的選兆.
每一個弟兄姊妹.
是有指望的.
他在聖徒當中所得到的基業.
有何等豐盛的榮耀.
你和我在神的眼中.
在這位造物主的眼中.
因為救恩.
帶給我們成為一份榮耀.
將來他看到這份產業.
他看為是榮耀.
不是頹垣敗瓦.
是一個榮耀的基業.
造物主的眼界.
肯定比我們闊很多.
如果你住過一些很美的地方.
突然間你住一些很小的地方.
那你就可能心暗了.
因為你住過很美的地方.
那些才叫真正的基業.
但是現在造物主說.
他在聖徒當中所得的基業.
是有何等豐盛的榮耀.
並知道他向我們信的人.
所顯的能力.
是何等的浩大.
不斷提何等何等.
去提升到極致的描述.
何等的浩大的能力.
是甚麼浩大能力呢.
小犬.
吹到人很麻煩.
不能乘車回市區.
還是怎樣呢.
他說這個浩大的能力.

$^{121}$就是二十節.
就是照他在基督身上所運行的大能大力.
曾經運行過甚麼大力.
在耶穌身上呢.
似乎我們看不到的十字架的時候.
有大力令到他可以不用死.
群眾的呼喊聲.
是令到他直接的致命.
釘十字架斷氣.
見不到有任何的神蹟.
那些人嘲笑地離開那個刑場.
但是三天之後.
基督在墳墓出現.
不是見鬼.
不是萬聖節.
是基督在墳墓有一個大力.
造物主的復活生命的能力.
在我們現在所信的耶穌身上.
他復活過來.
他說這份大能大力.
現在就運行在我們身上.
使到他能夠從死裡復活.
叫他在天上坐在自己的右邊.
主耶穌現在不是在這個赫西瑪利園.
他現在不是在那個葛國他的墳墓.
他現在是在天上至大者榮耀的主的右邊.
在神的寶座右邊.
我們敬拜的是這位主.
這是一個祈禱.
想我們能夠明白.
而且他現在的身份.
是遠超過一切執政掌權有能主治.
和一切有名.
不單止是今世的有名.
是連來世的有名都超過了.
然後又將萬有伏在他的腳下.
使到他為教會作為萬有之首.
這裡就開始提到他是教會的元首.
接著第23節.
就是和我們今天要說的道有關連.

$^{161}$教會就是他的身體.
充滿萬有者所充滿的.
充滿萬有的造物主.
充滿著教會.
這是一個首先的概念.
教會和教堂是不同的.
我們每個星期回來教堂.
有時在市中心.
街坊就問你在哪裡工作.
我在教堂建築物裡面.
街坊會問的.
因為知道我們是牧師.
很好奇在做甚麼.
教堂是這個建築物.
這個硬件.
但這裡說的是我們每一個.
信耶穌得救的弟兄姊妹.
認識耶穌的救恩.
你認同他的死和你有關.
他的復活和你有關.
你和我就是教會.
我們就是神的兒女.
我們也是門徒.
甚至四章第七節說的.
我們是聖徒.
這是不同的概念.
但最重要的是.
我們這班人的聚集.
才叫教會.
我們在這個禮拜堂.
教堂這個建築物聚會.
但我們這班人就是教會.
這裡要留意.
因為教會是他的身體.
是這位充滿萬有者.
所充滿的.
我們就要問一個很嚴肅的問題.
我們要怎樣成長.
身體如果是一個有機的身體.
不是一個洋娃娃.

$^{201}$他是會長大的.
怎樣的長大.
去到甚麼的標準.
才是合乎神的旨意.
四章七到第十六節.
就有三方面的教導.
第三章我們再看這個祈禱.
保羅又再一次為爾忽所教會禱告.
三章可能是第十四節開始.
因此我在父的面前屈膝.
他在坐牢.
在爾忽所.
他現在寫這封信的時候.
他在羅馬半遠近的坐牢.
因為福音的緣故.
天上地上的國家.
都是從這位的父神而得到名的.
即是說一切的眾生.
都是由他而來的.
是他命立就立的.
第十六節.
求他按著他豐盛的榮耀.
藉著他的靈.
叫你們心裡的力量剛強起來.
再一次提到我們心裡的力量.
第一章的祈禱就是要我們裡面的眼睛明亮.
看到恩典的選照的盼望.
這裡告訴我們.
我們要有力量.
這個力量用來做甚麼呢.
用來彼此相愛.
這裡告訴我們.
第十七節.
使基督因你們的信.
住在你們心裡.
叫你們的愛心.
是有根有基.
能以和眾聖徒.
一同明白基督的愛.
是何等長闊高深.

$^{241}$並知道這愛.
是過於人所能夠測度的.
便叫神一切所充滿的.
充滿了你們.
能夠照著運行在我們心裡的大力.
再一次說裡面的大力.
這個大力是令我們有愛的根基.
可以和眾聖徒.
不是你所喜愛的.
不是你自己所喜歡的小徹.
能以和眾聖徒.
所以必須是神的大能大力.
所以必須是基督藉著我們的信.
祂要首先住在我們心裡.
是有層次的.
所以我們不能夠盡力的愛眾聖徒.
是因為我們根本從來沒有基督住在我們心裡.
我們從來都沒有求.
這個大能大力.
令到我們心裡的愛有根基.
以致我們才能夠開始學習.
怎樣愛眾聖徒.
這裡告訴我們.
這個愛是過於人所能夠測度的.
想一想基督的犧牲.
這份大愛.
無條件為你,為我的犧牲.
這就是帶來我們的愛.
能夠體現.
這位主的愛.
是過於人可以去測度.
接著就說教會.
第20節.
神能夠照著運行在我們心裡的大力.
充充足足的成就一切.
超過我們所求所想的.
第21節.
但願他在教會中.
不是教堂.
是教會.

$^{281}$一群懂得救恩耶穌大愛的人.
一群已經分別了自己.
不再屬於這個世界.
屬於神的兒女.
神可以在教會當中.
在我們大家的生命裡.
得著榮耀.
直到世世代代.
永永遠遠.
阿們.
這就是在說教會.
所以我們今天集中的主題.
本來就是這裡有一個教導.
我們看第7節.
到第16節.
四章這段經文.
很重要有三方面.
讓教會來長大成人.
重點是成熟.
身心靈的健康.
錦繡堂42年了.
不是這個教堂.
是我們這一群人聚集.
由我們起義要在錦繡花園.
直堂有聚會開始.
第42年了.
如果用一個現在一個人的標準.
42歲正值壯年.
應該就不需要每半年去驗身.
看看你有沒有三高.
看看你有沒有其他的疾病.
當然除非你有其他的.
已經顯示出來的症狀.
那就要了.
如果正常有運動.
有健康飲食.
42歲是很健壯.
應該是可以是最重要的階段.
無論發展事業各方面.
但現在我們說的是教會.

$^{321}$我們就要問.
怎樣令到教會能夠朝向長大成人.
有這個目標.
大家能夠達至呢.
有些教會是超過100年.
我的同學在長洲浸信會.
超過100年.
百年老店就有很多問題.
因為又是長洲.
很多人搬走了.
出了市區.
所以很多不同的情況.
但我們要問.
怎樣叫做健康呢.
神的教會.
怎樣才是健康呢.
嘗試一下第七節到第十六節.
有三點和大家彼此的勉勵.
第七節開始.
我們各人蒙恩.
都是照基督所量給我們各人.
每個人的恩賜.
第七節已經告訴我們.
很重要新的修訂版.
真的非常之好.
強調每個人說了兩次.
舊的和話本只說了一次.
每個人.
我們每個人.
都是照基督將他的恩典.
來量給我們每一個人.
所以第一個我們要明白.
怎樣令到教會成熟.
必須每個人要有這個意識.
你信主.
主耶穌就會將恩賜.
恩典來給你.
這個恩賜是要來讓教會成長.
讓教會健康的發展長大.
這是最重要的.

$^{361}$第一個我們要認識的.
來到教會.
你回來有什麼想法呢.
你信耶穌入門.
開始的時候.
我們不是太熟悉.
你來到想聽崇拜.
想安靜.
想敬拜神.
但這是入門的開始.
你有沒有盡快意識到.
你在主的恩典裡.
獲得你自己那一份.
是量給你的恩賜.
為著的不是你自己個人的一種享受.
為著的是服事.
因為第十二節再說.
這是第一方面.
大家需要有這樣的認識.
所以進入第二個重點.
要有人幫你去發掘.
去發現和發展你的恩賜.
你就可以在教會裡.
來有那個服事.
記得不是在教堂.
是我們一群人.
如何能夠來到去.
在神的兒子.
他是教會的頭.
我們是他的身體.
如何來成長成熟.
所以第一方面就必須.
你要認識自己的恩賜.
你要去發掘.
要有人幫你去發掘.
要和我們有相交.
要來到不是你自己一個人.
每個禮拜坐在那裡孤零零.
坐完就走.
那個還不是一個成熟的想法.

$^{401}$聖經是告訴我們.
你必須要認識.
舊恩帶給你的生命改變.
所以第二方面.
感恩神賜給教會有目的的人.
第七節可以再補充一點.
因為突然之間.
保羅說了一個德姓的海船.
保羅的年代.
哪些國家是德姓的海船.
是有目共睹的呢.
你看過電影《奔驅》.
你就知道那些羅馬的將軍打勝仗.
回來的時候進入海船門.
就有兩邊的女孩子.
有些鮮花來審給將士.
德姓的海船的香氣.
但也隨著進城的列隊的時候.
都會有一些俘虜.
用一些鉤鉤著.
準備完成這個海船之後.
就要行刑來刺死.
所以是一些香氣為那些德姓.
那些不香的氣.
就是為那些要準備面對死亡的俘虜.
這是當時經常有的場景.
保羅就用這個海船.
告訴我們哪個海船.
在第八節.
他首先引述詩篇的第68篇第八節說.
所以經上這樣說.
他升上高天的時候.
擄掠了仇敵.
將各樣的恩賜賞賜給人.
這個豈不是論功行賞.
羅馬的將領將軍.
海船回來之後.
就會論功行賞的那個場景.
第九節.
既然說升上.

$^{441}$豈不是先降在地下.
那降下的就是遠升諸天之上.
要充滿萬有的.
特別這樣提.
主耶穌是首先降到地上.
成為一個人.
那撒勒耶穌.
然後在地上傳道.
最終被人來嫌棄.
被人來藐視.
被人來誤解.
就釘十字架.
這是我們大家所知道的.
但是世人所不知道.
而我們今天認順的.
就是我們的主耶穌已經復活升天.
保羅特別強調.
他升上去之後.
擄掠了仇敵.
相信不是那些得罪他.
逼迫他的那些祭司長.
或者彼拉多.
這裡說的是屬靈界的黑暗勢力.
死亡 魔鬼.
來到他已經得勝了.
他擄掠仇敵之後.
就將各樣的恩賜.
來賜給人.
所以這是一個海船的恩賜.
弟兄姊妹.
教會不是求你回來幫忙做義工.
表面是.
實際不是 保羅說.
是這位得勝的耶穌.
復活之後.
將恩賜賞賜給每一個.
信他的弟兄姊妹.
所以這是無上的光榮去服侍神.
不是廉價的義工.
你的mindset要改變.

$^{481}$不要搞錯.
即是說你能夠有恩賜.
是神賜給你.
是糧給你的.
然後你能夠參與去事奉.
為著讓身體成長.
成長之後.
你就是帶著一份無上的光榮感.
參與服侍的.
是要有這樣的心態.
能夠事奉.
是因為他將恩典給過我.
我可以參與服侍.
然後這是得勝的服侍.
現在我們不是服侍一位.
是呀 二千年前他死了.
不過我們都懷念他.
所以我們聚集在一起.
每個星期就紀念他.
不是呀.
我們紀念的是這位主耶穌.
已經復活.
在天上坐在父神寶座右邊.
所以我們的基督教.
是亮麗的宗教.
我們的身份.
我們現在是在得勝的裡面.
去發展我們地上的傳道的事工.
我們不是fight for一個virtually.
我們是fight from一個virtually.
由那個得勝開始.
我們去服侍.
所以教會.
是一直一直有成長的.
二千年沒有任何的政權.
沒有任何的黑暗勢力.
包括人的軟弱.
可以打倒教會.
當然教堂可以沒有.
教堂失修就沒有.

$^{521}$大家不肯奉獻就沒有.
就難.
但那個教會是那班人.
所以我們可以擴散去到任何地方.
為神的福音.
來領受神給我們的恩賜.
我們就可以去服侍.
這種力量就是來自神.
很大的力量.
這個第一方面大家要記得.
怎樣讓第一方面.
能夠真是有效地運作呢.
感恩就是十一節.
神賜給教會.
有五個積分.
他所賜的有使徒.
有先知.
有傳福音的.
有牧者和教師.
五個積分.
為著什麼呢.
這五個積分.
怎樣存在呢.
是由神所賜.
神賜給教會.
五個特別積分的人.
來做什麼呢.
就是十二節.
裝備.
使每一個的聖徒.
就是第七節所說的.
你會慢慢去發掘到你自己的恩賜.
你會發現.
原來你服侍神是一份恩典.
來在教會裡面.
又參與每一個的事奉工作.
這個就是第十二節.
我們這五個積分的存在.
在教會裡面的目的.
所以又是生生不息.

$^{561}$教會需要牧者.
昨晚我也說了一個.
很特別的討論.
我的兒子在車上回家時問我.
最終的結果是怎樣.
因為昨晚我們說過.
讀神學的期間.
我們神學生討論了一個問題.
當然有些老師在場.
其實教會需不需要請傳道人呢.
全職的傳道人需不需要呢.
因為現在神學的普及教育.
不用做傳道.
也可以讀神學.
甚至有學問.
有學識.
有好好的工作.
不需要做傳道.
繼續在教會熱心服侍.
朝向這個方向.
其實不用全時間受新的傳道人.
我們就一路討論.
我沒有說昨晚的答案.
我兒子很好奇.
突然問我.
到底怎樣.
最終的結論是怎樣.
我當然回答了.
梁院長和我們一起討論.
他說根本沒有可能.
絕對不是因為我們要找傳道人.
要請他回來工作.
而正正就是第十一節所說.
是由神賜予教會特別的職份.
這些職份.
不單純是我們在這裡做打工.
沒錯.
我們每個童工都有一個合約.
我們要跟那個合約去工作.
但是我們每一位的牧者.

$^{601}$都很清楚自己的身份.
我們是蒙召.
是他所賜予教會.
來出現的.
就有這五個的職份.
而且很特別.
在這段聖經說的這些職份.
只是說真道建立的職份.
其實保羅在哥林多前書十二章.
去到第十四章.
他舉了很多不同的恩賜.
還有一些是行神蹟的.
可以講方言的.
可以翻譯方言的.
等等的恩賜.
但是這裡是隻字不提.
特別強調是建立信徒對真理認識的恩賜.
的專職的牧者.
這樣就可以想像得到.
保羅是刻意這樣去營造.
教會能夠健康的建立.
其中必不可少的.
就是聚會的那班人.
對救恩的傳遞的認識.
是否具備呢.
這個就很重要了.
至少我們在初信班到浸禮班.
每一個準備洗禮的.
清楚自己跟神的關係.
清楚自己承認自己是一個罪人.
透過主耶穌的寶血.
得到救贖潔淨.
這個至少是我們到今天.
都堅持我們必須牧者.
教回初信班浸禮班的原因.
就是我們要讓每一個信徒.
都對自己的救恩.
是要有很清楚的認識.
這是很重要的教會成長的真理.
第二個就是第十四到第十六節.

$^{641}$出現在愛裡面建立.
就是說教會必須有真道.
教會必須在愛裡面去建立真道.
在信徒的心靈裡面.
所以第十一節.
教會的領袖.
教會的弟兄姊妹.
有沒有將這些牧者的積分.
擺在他自己應該發揮的位置呢.
那個教會就健康了.
那個教會就能夠一路的成長了.
如果我們做一些不是神賜給我們去做.
而我們將應該做的我們不做.
那就是一個很大的問題了.
所以使徒是一班.
沒錯固然十二使徒.
是有他們很特殊.
因為跟隨耶穌.
所以變成一提到使徒.
就是提到這十二個人.
但是真正的意思.
是新約的意思.
不是出人高人一等.
不是特別的有什麼經歷.
是他們有一個使命.
有使命的門徒.
所以我們應該藉著牧者的.
終身蒙召.
毀身的事奉.
會感動你.
他們是有使命.
這個使命應該感動我們.
我們都要參與在教會的服事.
都要帶著一個使命感.
大家明白使徒的意思就是這樣.
先知不是說微卜先知.
不是說真牧師.
到底現在以巴那個戰爭怎樣走向.
我們不是做一個微卜先知.
先知在保羅的理解是講道.

$^{681}$因為舊約是有一些先知講預言.
但是在保羅的用法上.
哥林多前書十三章第七.
可能不用了.
我若有先知講道之能力.
又有傳遞的信心.
卻沒有愛.
我就好像明的羅響的筆.
先知是講道.
將神的真理深入淺出.
融會貫通的講解給回眾知.
去建立去餵養.
讓弟兄姊妹能夠成長.
接著有傳福音的.
你就可以想像到.
那個真是一個很特別的恩賜.
有些童宮.
我其中講一個.
梁月華傳道.
如果你認識他.
講幾句他就會進入到福音.
特別的恩賜.
不是那麼容易.
你就可以想像到.
神給教會有這些職份的牧者.
就是要怎樣.
就是要去培訓.
裝備每一個聖徒.
去準備做服事的工作.
親愛的弟兄姊妹.
教會的牧者.
我們有沒有這樣去看待.
你的出現.
在大家當中的牧養教導.
然後你的心靈得到的成長.
最終是要服事的.
最終每一位弟兄姊妹.
無論你新來的.
或者是很多年服事的.
你都是要參與教會的服事.

$^{721}$因為十三節.
保羅講的一個人的成長.
十三節說直至到我們眾人.
即是每一個信徒.
在真道上.
所以為何上文講的五個職份.
都是和真理有關係的.
他不是要將整個的恩賜的調查表.
每一項都列出來.
他特別強調的.
教會每一個信徒.
是需要對真理.
有一個全被的掌握認識.
所以是要直至到.
這個生命成長.
是要在真道上同歸於一.
認識神的兒子.
以致我們長大成熟.
這個就是目標.
長大成熟.
直至到.
要等.
身體的成長是要等的.
我記得很小的時侯.
和他晚上吃完飯.
去惠康買米.
因為沒有米了.
那時好像沒有五公斤的.
因為現在大家越來越不想吃那麼多米.
就越來越細包.
以前是十公斤.
十五公斤.
就買一包十公斤的.
那時他還是三年班.
一個三年班.
他怎可以抽一個十公斤.
你可以想像得到.
出力的是我.
不是他.
表面上好像是他抽著.

$^{761}$但其實是我出力.
要等.
一個生命的成長.
現在他還快過我.
開始又要練馬拉松.
有時我追不到他.
生命原來是真的會成長的.
敢宣家四十二年.
你每個星期來採教會.
你有沒有感覺到.
你自己的生命需要成長呢?.
恩賜已經是藉著耶穌的升天賜了給你.
牧者行常的循循善誘.
就是要讓你去發現.
去發掘.
和去發展你的恩賜.
直至到你參與事奉的工作.
所以這是教會的成長.
健康的發展.
如果不是呢?.
十四節就出現了.
會搖來搖去.
不是坐海盜船那麼簡單.
是真理會被人挑戰.
十四節.
使我們不會再做小孩子.
中了人的詭計和欺騙的法術.
被一切的異教之風來搖動.
飄來飄去.
就隨從各樣的異端.
不單止是不進則退那麼簡單.
是會令你失去你對永恆的盼望和救恩.
所以為什麼要均衡的成長.
要健康.
不是過度的膨脹.
是要健康.
現在很明白了.
不知道在座有沒有參加一些訓練.
一些健身班.
特別有專人教的.

$^{801}$那個吸引就是很多時候那些男教練.
又大隻又帥.
每個禮拜見他都有點動力.
逼你去健身.
每個禮拜回來你見到的牧者.
就是你的屬靈的健身教練.
是神設立的.
是讓你能夠成長.
是挑戰你人生為誰而活.
是告訴你.
是真的有永恆你和我要進入.
將來我們要向神交征.
這個就是告訴我們.
我們不要做小孩子.
被異教之風搖動.
你可以想像得到多慘.
飄來飄去.
我們應該是越來越對真理穩固.
我們的愛心是有根有機.
在第15到16節講了兩次.
愛心在教會裡面講誠實的話.
在愛裡面16節.
在愛裡面彼此的建立.
在這個地方就是我們一班.
這個就是我們屬靈的家.
一起去成長.
所以弟兄姊妹.
今天就是這三方面.
我們需要來到彼此的勸勉.
在教會裡面的相交.
所以如果有些人只是回來團契.
我不回崇拜的.
我不回主義學的.
我只是回團契.
團契不要查經.
我們只是聊聊天吃飯.
這個不叫聖經的團契.
可以叫團契.
可以叫fellowship.
但是沒有真理.

$^{841}$也不是和眾聖徒的相交.
講不上是神期望的教會生活.
我們一班人.
真是在不同的場景裡面.
去彼此的學習.
彼此的成長.
很盼望.
這個月我們選取的經文.
就是因為我們是堂慶42週年.
42歲是精壯.
我今年60了.
跑步還沒有索氣.
還可以踩單車上大堂.
原來現在神給我們身體.
在各方面的科技飲食.
相對可能幾百年前的人.
我們是健康了.
但我們在說的是.
教會基督的身體是否健康呢.
在座的.
如果隨便有人問你.
羅馬書到底講什麼的.
你讀都沒有讀過.
你講什麼呢.
我們在座的弟兄姊妹.
是很有追求的.
隨時可以裝備好.
去教聖經給其他弟兄姊妹.
這個才是健康的教會.
求主幫助我們能夠.
真的看到保羅的羽眾心腸.
對於我們基督的身體.
這個教導.
大家一起的學習.
所以這兩方面是很重要.
在真道上.
同歸於一.
認識神的兒子.
得以長大成人.
和在愛中彼此的建立.

$^{881}$我們一起祈禱.
多謝主.
今早我們能夠一起.
將已不所書這段聖經.
來彼此的學習教導.
懇求主來感動弟兄姊妹.
一同的發掘.
發展.
和在教會裡.
同心的服侍.
奉耶穌基督的名.
阿們.
\newpage



\section{撒迦利亞書}
\label{sec:c_6waYrNYvI}
\textbf{2023-10-19 周四晚查經班《歸回先知書研讀(6)撒迦利亞書》 講員:曾裔貴牧師}
\newline
\newline
連結: \href{https://youtube.com/watch?v=c_6waYrNYvI}{\texttt{https://youtube.com/watch?v=c\_6waYrNYvI}} ~~~~ 語音日期: 2023-10-21
\newline
\newline
\hyperref[sec:cM9RUvSdlG8]{\small{< < < PREV SERMON < < <}}
~
\hyperref[sec:index_chronic]{\small{[返順時目]}}
~
\hyperref[sec:index_scriptual]{\small{[返順卷目]}}
~
\hyperref[sec:qZ9s1IR331E]{\small{> > > NEXT SERMON > > >}}
\newline
\newline
$^{1}$我們一起祈禱.
感恩天父給我們在晚上再一次有機會來研讀.
先知的著作《撒迦利亞書》.
我們會由第七章開始看.
離開了八個的異象之後.
我們來看神如何處理當時耶路撒冷的居民.
和一些被撈歸回的猶太人.
他們的問題和審判列國.
我們都很盼望.
這些可能已經是很多年前的一些事情.
未必和我們有直接的生活上的關係.
但是對將來耶路撒冷的預言.
對神的選民.
對於聖地.
特別是將來的局勢.
我們很盼望神給我們能夠有一個留心.
為神的選民.
為以色列去禱告.
但是我們看到現在的情況.
我們知道我們不能夠.
簡單幾句話就可以分析.
就可以來做一個很清楚的研判.
但是我們都盼望神給我們能夠.
透過舊約的先知的預言.
給我們先有一個概念.
能夠掌握到到底神如何去處理.
你自己的選民.
他們和你的關係.
求神帶領我們明白經文的內容.
我們這樣禱告奉耶穌基督的名.
阿們.
我們講了頭六章.
也都處理了八個異象.
最後我們看到.
大祭司約孫亞.
又成為大祭司.
又成為了加官面的君王.
當然是一個象徵性.
引發了到底為什麼會.
無端端給皇官約孫亞.

$^{41}$是很難想象.
但是如果我們明白了.
將來的尼賽亞.
他不單止是先知.
是祭司.
更加是大魏的子孫.
他是君王.
所以我們的主耶穌.
他是那一位和平之君.
是很重要的.
你可以說幾個的積分.
都在他的身上.
這個就很特別.
所以在舊約都已經有這些預言.
將來的尼賽亞.
他身兼兩職.
可以在他的積分里.
籌訂和平.
帶來真正的.
整個世界的和睦.
這些是很.
令我們很感到振奮的事.
其實一直讀下去.
開始的審判.
其實對猶太人.
以色列是一個很大的安慰和祝福.
因為其中一個概念就是.
神要再回到耶路撒冷.
神要再親自眷顧.
他自己的選民和地圖.
這個對猶太人來說.
是夢寐以求的事.
所以這個就是先知書的預言.
其中寶貴的地方.
特別因為現在的時代.
已經是經過了70年的被撈.
他們現在在這個時候.
是回去第一批的人.
重建耶路撒冷的聖殿.
後來因為有當地人的反對.

$^{81}$就一路建造延後了十幾年.
去到現在大理污王的第二年.
就正正式式.
透過哈該先知和撒迦利亞先知.
再勉勵百姓.
他們就重新去建造.
是由第二年.
大理污王第二年.
由六月開始建造.
現在撒迦利亞先知.
其實我們之前說的.
這八個異象.
這六張聖經.
其實在第二年的六月.
現在就去到第二年的九月.
撒迦利亞.
現在就過了兩年.
第七章開始.
就是大理污王的第四年的九月.
他又特別提到.
是猶太的月份.
他們有他們自己的名稱.
基斯留約大概是五至六月.
這裡告訴我們.
基斯留約的初四日.
耶和華的話臨到撒迦利亞.
那時伯特利人已經打發沙利息和利堅米勒.
並跟從他們的人去懇求耶和華的恩.
並問萬君子耶和華殿中的祭司.
和先知說.
我歷年以來在五月間哭泣齋戒.
現在還當這樣行嗎?.
五月就不是很大的節期.
猶太人他們七月就會比較多節期.
會堆在一起有幾個重要的節期.
七月初一是摧國節.
七月初十很重要.
七月初十就是贖罪日.
七月十五一路.
總共八天就是駐彭節.

$^{121}$上一次的恐怖襲擊.
就是在駐彭節週末的時間.
因為家家戶戶都會在野外.
搭那些營來駐彭.
紀念他們出埃及的時候.
駐的帳幕等等.
要招行猛策.
他們每年的駐彭節都會這樣做.
找些棕樹枝搭一個簡陋的棚.
來紀念.
是一個歡樂的節期.
所以五月其實就不是太大的禁食的需要的日子.
但是現在他們在尋求釋法.
我們做了這麼多年.
要不要再做呢?.
就來到去懇求神的恩.
還有找了三個代表.
去尋求大祭司.
和先知的解釋.
到底我們現在怎麼做呢?.
其實他們緊張的是.
神會不會我們做漏了.
做得不夠呢?.
我們要怎麼做呢?.
你告訴我們吧.
現在很多人開始很有神心.
很想來尋求神的指引.
我們讀下面.
神是有解答的.
不過去到第八章才解答.
反而是神用一些過去的歷史來責備他們.
首先要解釋.
到底你們的哭泣,齋戒.
你們的禁食是所為何事呢?.
慢慢是否已經慣性.
總之是去到那個日期就要禁食.
就變成你所做的都不知道做什麼.
其實就不好了.
第五節讀回來.
你們要宣告.

$^{161}$第四節好一點.
萬君子耶和華的話就臨到我說.
就是撒迦利亞.
你要宣告國內的眾民和祭司說.
你們這七十年.
就是被擄的七十年.
在五月,七月禁食悲哀.
豈是絲毫向我禁食嗎?.
七十年是幾代人了.
可能一至兩代.
有些可能是三代人了.
被擄他們一直等待回歸.
在被擄之地.
就來到繼續在他們的首節里.
可能是約定俗成.
可能是一些老人家.
吩咐他們這樣做.
慢慢去到之後.
就因為在外邦的波斯地方.
當然是巴比倫.
跟著就是馬代波斯.
這樣來到交替的政權.
一路這樣來到更替之後.
他們一路等待日期.
就是七十年的期滿.
所以個個家家戶戶的猶太人.
都不敢不守這些節日.
不敢去齋戒禁食.
因為很想得到神在外邦的眷顧.
他們也有這種慢慢的.
人在異鄉.
你去到三藩市.
你見到一間雲吞面鋪.
你的眼睛都發光.
你都很想.
雖然很難吃.
你都覺得很好吃.
因為你吃過香港的.
吃過自己小時候.
媽媽帶你去吃的雲吞面.

$^{201}$你就思鄉的情.
所以他們是很緊張.
每年都會在五月七月.
來到齋戒禁食.
問題就是現在回來.
我們是不是要這樣做呢.
他們就很想去求釋法.
但是這段經文.
重點是神帶他們去看.
到底禁食齋戒這些節日.
到底真正背後的動機是什麼呢.
這個很重要.
這個就是將我們的宗教性.
所以有時又說遠了.
但其實都不是說遠的.
是有關的.
復活節.
尤其是平安夜.
很多閒集人來教會崇拜.
晚上的.
如果晚上有聖惡.
有平安夜的聚會.
有很多閒集人來.
他們可能未必信耶穌.
但他們很想有這種宗教氣氛.
他們會來教堂.
如果外面的中環更加多.
你們可以想象到.
那個宗教性都會不知不覺.
滲透在一般人的內心.
但是平日他們就不會覺得需要神.
但突然間有一種莫名的空虛.
很想去教堂.
去看什麼彌撒.
去聽你們的聖詩.
可能都是唱平安夜.
他們憤怒地感受到.
這種親切感.
這個宗教性.
原來每一個時代都會有.

$^{241}$今晚神就透過他們的求問.
正正式式說了.
到底宗教性要放在哪個合適的位置.
才可以和神有一個結連.
這一段聖經第七第八章.
就有很清楚的解答.
所以其實很寶貴.
因為猶太人自己問自己重要的節日.
我們這樣做.
所謂何事呢?.
做了七十年了.
第五節.
你要宣告國內的眾民和祭司說.
你們這七十年在五月七月禁食悲哀.
豈是絲毫向我禁食?.
你們吃喝不是為自己吃.
為自己喝嗎?.
當耶路撒冷和四圍的城邑有居民.
正興盛南地高原有人居住的時候.
耶和華直從前的先知.
所宣告的話.
你們不當聽嗎?.
你就可以大概感受到.
原來焦點不是在神那裡.
所以禁食變成很中性.
齋戒首節旗是很中性的一件事.
但是你對準的是誰呢?.
這個就要做的人就要問了.
第八節.
耶和華的話又臨到撒迦利亞說.
萬君子耶和華曾對你們的列祖如此說.
要按治理判斷.
各人以慈愛憐憫弟兄.
不可欺壓寡婦孤兒寄居的和貧窮人.
誰都不可心裡謀害弟兄.
這些豈不是他們被神來懲罰.
要被擄去到外邦的種種的惡行嗎?.
就是這樣一路積累.
積蓄神的憤怒去到極點了.
去到極致了.

$^{281}$神就要將這個民族來懲罰.
透過外邦的人巴比倫.
來將他們擄到遠方.
所以這裡神就再次透過撒迦利亞先知.
重提舊事.
好像要翻案.
來提到他們列祖的問題.
第十一節.
他們卻不肯聽從.
扭轉肩頭.
塞耳不聽.
使心硬如金剛石.
不聽律法.
和萬君之耶和華用靈直從前的先知所說的話.
故此萬君之耶和華大法烈怒.
萬君之耶和華說.
我曾呼喚他們.
是一路說過去的歷史.
過去的先祖.
過去的先知.
來到一路苦口婆心的.
來到去勸勉.
來到去預言.
來到去警告.
以色列人要歸向神.
但是呢.
都是一路都不聽.
將來他們呼求我.
我也不聽.
第十四節.
我必以旋風吹散他們.
到數不認識的萬國中.
就是被擄了.
這樣他們的地就荒涼.
甚至無人來往經過.
因為他們使到美好之地荒涼了.
這個是沒有答案.
到底怎樣才是真正神.
合乎神心意的齋戒禁食呢.
五月七月.

$^{321}$在這些節期的日子.
還應不應該做呢.
所以這個是.
回過頭來.
去到.
必然要問的問題.
猶太人.
他們一路到底.
要怎樣將那個心.
那個信心.
擺在這些實際的.
參與的這些宗教行為裡面呢.
要怎樣做才是.
合乎神的心意呢.
這些是一個有宗教的人.
才會問的問題.
因為正正就是周而復始.
這樣做這些宗教的禮儀.
一些不信神的人.
不回教堂的人.
他們沒有這些問題.
沒有這些掙扎的.
他們不需要去問的.
因為我根本不信.
你們這些禮儀.
可以承載一份的恩典.
可以承載帶來的那個信心.
所以是有信仰的人.
才會問的.
但是現在就是了.
他們就要問了.
我們這樣做.
是不是做完.
都是得個吉呢.
應該怎樣才是.
合乎神的心意的.
一個敬拜.
或者是齋戒呢.
第八章就有答案了.
神一會兒會在第.

$^{361}$好像第九還是第十節.
是會有解答.
我們一路讀下去.
這兩章聖經.
很多時候提到過去的.
先知的預言.
而且是由聖靈去執先知講的.
所以先知不是講自己的說話.
較為貼切的就是.
譬如你每晚看新聞.
每個國家有他們那個.
外交部的發言人.
現在就叫做毛靈.
就是這位.
他不可以講自己的東西.
他只能照稿讀.
是外交部要他講的官方的說話.
多一點的他都不能夠回答.
有時候笑笑口地避開.
有時候問他.
那個秦剛在哪裡.
沒有補充.
沒有東西回答.
他不可以自己在這裡爆肚.
不可以.
因為他是外交部的發言人.
一樣.
先知是神直住.
先知向百姓講的.
所以他有責任就是.
講神要他講的.
這個是舊約很重要的一個職份.
和那個功能.
就是要將神.
向眾百姓講的說話.
來講出來.
所以如果加多減少.
就變成是假先知.
就算你很有見地.
你講你自己的東西.

$^{401}$所以舊約是有規範的.
我們看第八章.
萬君之耶和華的話臨到我說.
萬君之耶和華如此說.
我為石安心裡極其火熱.
我為他火熱.
向他的仇敵發烈怒.
耶和華如此說.
我現在回到石安.
要住在耶路撒冷中.
耶路撒冷必稱為誠實的城.
萬君之耶和華的山.
必稱為聖山.
所以聖城就是這樣的意思.
就是這樣來的.
萬君之耶和華如此說.
將來必有年老的男女.
坐在耶路撒冷街上.
因為年紀老邁就手拿拐杖.
城中街上必滿有男孩女孩還灑.
萬君之耶和華如此說.
到那日這是在廬城的民眼中.
看為稀奇.
在我眼中也看為稀奇摸.
這是萬君之耶和華說的.
萬君之耶和華如此說.
我要從東方從西方救回我的民.
我要領他們來.
使他們住在耶路撒冷中.
他們要作我的子民.
我要作他們的神.
都憑誠實和公義.
萬君之耶和華如此說.
當建造萬君之耶和華的殿.
立根基之日的先知.
所說的話.
現在你們聽見.
應當手裡強壯.
那日以先人得不著顧家.
生畜也是如此.

$^{441}$且因敵人的緣故.
出入之人不得平安.
乃因我使道眾人互相攻擊.
但如今我待者如城的民.
必不像從前.
從前就是不斷的去刑罰.
現在就是立意的去施恩.
和他會重新眷顧這個城.
這是萬君之耶和華說的第十二節.
因為他們必平安撒種.
葡萄樹必結果子.
地上必有出產.
天也必降金爐.
我要使道者如城的民.
享受這一切的福.
所以保羅都這樣說的.
在羅馬書的十一章提到.
是有些特選的選民.
就是這些所謂的愚民.
所以這個選民的概念.
保羅也有這樣提到.
我等下再和大家揭羅馬書那一段.
我們稍微讀完這一大段.
你會見到一路說.
是神現在將他過去.
這個民族受盡神的責罰.
現在好像是完全的完結.
將來呢.
到底這個是什麼時候的將來呢.
是短暫他們現在重建聖殿的復工的開始.
過了兩年.
因為第二年的六月開始恢復.
現在第四年的九月.
就是這個基斯留月.
就有預言.
因為這個伯特利人.
來尋求神的釋法.
就變成了.
他們很想瞭解一下.
到底神現在是否在祝福我們的禁食.

$^{481}$這個節旗和這些宗教的行為呢.
但是現在我們見到.
神好像很充分地將他的恩言.
來講述.
你想想國泰民安.
人口.
連老人家.
又不是人口老化的慘況.
是充滿人口.
你對於一個荒蕪了的.
荒涼了的一個被老的地方.
現在充滿人口.
是一個很大的祝福和歡欣.
這些我們很明白.
因為有個疫情的幾年.
所有東西停擺停頓.
任何人都不可以聚集超過兩個人.
你覺得簡直難以想象.
但是現在神再眷顧耶路撒冷這個聖城.
所以這裡給我們有一個很大的.
對將來這個民族.
真的說這個猶太民族.
是很有盼望的.
但是怎樣能夠達成將來的應驗呢.
真的在當時的人的眼中是看為稀奇.
就是沒有可能.
所以這些預言擺在這裡.
什麼時候會應驗呢.
什麼時候是落實的日子呢.
這些都不是我們能夠預測到.
第十三節.
猶太家和以色列家.
你們從前在列國中.
怎樣成為可就坐的.
照樣我要拯救你們.
使人們稱你們是有福的.
你們不要懼怕.
手要強壯.
萬君子耶和華如此說.
你們列祖惹我發怒的時候.

$^{521}$我怎樣定義降和.
並不後悔.
現在我照樣定義.
施恩與耶路撒冷和猶太家.
你們不要懼怕.
你們所當行的是什麼呢.
這裡說了.
十六節開始說了.
各人與淪捨說話誠實.
在城門口按治理判斷.
使人和睦.
誰都不可心裡謀害淪捨.
也不可喜愛喜假釋.
因為這些事都為我所恨惡.
這是耶和華說的.
接著十八節就回應剛才的問題.
萬君子耶和華的話臨到我說.
萬君子耶和華如此說.
四月五月禁食的日子.
七月十月禁食的日子.
必變為由大家.
歡喜快樂的日子和歡樂的節期.
所以你們要喜愛誠實與和平.
其實很白說.
很清楚已經說到.
本來那些節期本身是悲哀的.
是需要你認真的.
贖罪日就是.
神.
等下我又和大家揭一揭.
什麼叫贖罪日呢.
那個儀式.
那個節期.
是如何守法呢.
這個是你沒記23章.
耶和華的七節期.
其實大家都要熟的.
那些是說明要悲哀的.
但是現在神說將來.
不需要了.

$^{561}$連那些悲哀.
禁食.
守齋戒的日子.
只要你個人本身.
是合神心意.
那些節期.
就會更加奮愛.
來達到.
他預期的功效.
功能.
但是如果反過來.
你個人本身不是神所喜悅.
到時你.
怎麼哭到死.
都變成虛假.
變成一種儀式.
變成一種宗教.
那就是.
這兩章聖經.
的焦點.
因為就是由這些.
有虔誠信仰的猶太人.
他們很想求個明白.
安心.
就是說我們做了70年.
阿爺教我們.
阿爸叫我們這樣做.
又有不同的人叫我們這樣做.
現在我們在外邦做.
做到回來了.
回到耶路撒冷了.
那現在我們五月七月.
是不是都還要禁食齋戒.
還要這樣來悲哀呢.
這是一個很好的問題.
一個很牽連到.
到底一個有信仰的人.
他應該.
怎麼可以得知到.
自己和神的關係.

$^{601}$是一種和諧呢.
這個是.
很值得我們去思想一下.
討論一下.
第20節.
萬君子耶和華如此說.
將來必有列國的人.
外邦人.
和多城的居民來到.
這城的居民必到那城.
說我們要快去.
懇求耶和華的恩.
尋求萬君子耶和華.
我也要去.
必有列邦的人.
和強國的民.
來到耶路撒冷.
尋求萬君子耶和華.
懇求耶和華的恩.
萬君子耶和華如此說.
在那些日子.
必有十個人.
從列國諸族中出來.
不同的民族.
拉著一個猶大人的衣cum .
說我們要與你們同去.
因為我們聽見.
神與你們同在了.
以馬來尼.
神與人同在.
這兩章聖經.
源於當時的伯特利人.
有幾個被派遣的.
一些使節.
來到耶路撒冷.
其實他們已經在耶路撒冷.
他們去到聖殿.
因為聖殿正在建造.
其實是一個工地.
一個工程.

$^{641}$大祭司和所羅巴伯省長.
日間督工.
可能他們去到一些安息日.
又會教導百姓.
又會讀律法書.
又會有一些節旗來守.
但是可能還是一個帳篷.
未是一個完整的聖殿.
但是可能那些重要的功能.
都開始會恢復.
但是他們就很想問一個.
在他們內心積累了很久的.
到底我們現在在聖城.
在聖殿的環境.
我們被老歸回.
我們還應不應該.
來守這些齋戒節旗呢.
這個就是第三節.
我們去找祭司和先知.
來問的問題.
這個問題.
等下八章的十八節.
十九節就解答了.
第三節就是說.
我們歷年以來.
在五月之間哭泣齋戒.
現在還當這樣行嗎.
要怎麼做.
現在不同地方了.
我們被老歸回了.
現在整個世界局面都轉變了.
我們不需要再寄人籬下了.
當然也有些寄人籬下的.
沒有回去重建聖殿的.
也沒有回歸的.
一直都住在波斯帝國的.
例如我們知道的以斯帖.
末底改.
他們都是住在書山城的.
他們一直會守住這些齋戒節旗的.

$^{681}$但是問題就是.
當時也不是那麼方便.
用資訊軟件.
資訊流通不會那麼快.
是靠用碼的驛站.
來傳遞信息.
其實快碼也要一頭半個月.
來傳遞信息.
但是現在真的有一個.
比較官方的.
一個使節.
一個訪問團.
去求問釋法.
到底要不要繼續這樣做呢.
就算繼續做.
應該用什麼心態去做呢.
這些是很好的問題.
應該是每一個.
信了主耶穌的弟兄姊妹.
都應該要問的問題.
我們會不會慢慢慣性了呢.
洗禮是不是高峰.
一路一路慢慢走下坡呢.
慢慢信仰就變成很個人化.
自己喜歡怎樣就怎樣.
回聚會喜歡就回.
有空就回.
沒空就不回.
這樣是不是一個正式的.
真正的宗教的行為.
那個生活呢.
這些是息息相關的.
是很貼近我們的心靈的.
有沒有一個.
就算外邦人信了主耶穌.
一樣是要問的.
因為大家雖然是.
有不同的時代.
但是敬拜著.
獨一的真神.

$^{721}$還有到底我所做的.
這個宗教的行為.
是不是其實是他喜悅的呢.
如果不是做了幾十年之後.
原來他到時一不嘔銷.
不認數.
你們這樣的宗教.
是你們自己喜歡的.
那就很大件事了.
那我們自己一廂情願.
去做了.
每個星期坐定定在那裡.
原來神說我不收納你們.
因為真的不同的教會.
有不同的形態.
如果你去到黑人的教會.
美國比較多.
特別是在南方.
黑人教會.
他們偏向會靈恩.
他們偏向會在歌在舞.
很多很多的.
他們會叫喊.
很有和應的.
那些講完那句.
他們全場會和應.
他們差不多是整個聚會.
很有那種氣氛的.
你去到一些可能白人的教會.
就很沈靜的.
尤其是一些.
封建派的教會更加沈靜的.
來到去經歷.
這樣的一個.
類似安靜.
很多一定是有背後的神學.
但是神學.
慢慢慢慢又變成一個習慣.
變成一個固定了的.
約定俗成的.

$^{761}$這樣的一個模式.
有時大家都習慣之後.
你就很抗拒突然間的改變.
突然間加插一些東西.
你會覺得很惡毒的.
譬如偶然間.
在崇拜.
突然間有個年輕人站起來.
叫回應.
其他人就覺得很不舒服.
甚至有沒有搞錯.
是不是有鬼符.
有靈界.
想很多東西.
連那個講完可能都有很多東西想.
什麼事呢.
一般可能我們當他就是.
他情緒太高漲.
或者他有可能.
情緒的問題.
我們比較少.
會從另一個角度去看.
當然一般都沒錯的.
一般都是這樣的.
正正因為我們已經慣性了.
這個敬拜的模式.
我們就不接納.
其他的模式.
這個某程度都是我們的限制.
或者應該是每個人都有這種的限制.
慣性了自己的模式.
一下子要改.
其實很難的.
那個本來.
其實不是一個絕對的真理.
那個模式只是我們習慣了.
所以有時我們就要.
開放我們自己的心靈.
在那個開放裡面.
又不要失去了那個真理.

$^{801}$這個就真的很難.
又要守住真理.
但又不可以很固步自封.
這個怎麼能夠來到.
有一個協調呢.
那就是很高的學術.
的藝術.
所以這裡又透過這一大段經文.
真的給我們有一些解答.
開啟了我們一些.
可能我們平日都少留意的.
那個關注.
原來神關注的是.
我們那個人的內心.
到底是怎樣的一個人.
夠因神的真理.
神的那個屬性.
在我們的生命裡面.
到底有多少的改變呢.
我們是不是越來越貼近.
神的那個誠實和和平呢.
這個是後面第十九節.
所講的.
所以那個解答的問題.
答案就是.
所以你們要喜愛誠實和和平.
這個就變成了.
你們守不守節.
不是那麼重要了.
甚至一個本來是要我們.
哭泣齋戒的節期.
我們可以使用來變成.
歡樂的節期.
就是說原來.
神不是要我們.
只有一個外表的表情.
是我們裡面的人的態度.
這個正正就是信仰.
對一個人最重要的教導.
因為我們可以向人.

$^{841}$有不同的面具.
但是我們在神的面前.
是沒有任何的面具的.
所以這個就是.
很值得我們去透今晚的經文.
當年他們這樣問.
就是說問的人真的很想.
得到神的喜悅.
他們所做的這些.
宗教上的.
恆常的.
這樣的守節期.
和他們的祖宗教他們要這樣做.
就會得到神的恩典.
沒有人敢反對不做.
做了慢慢做.
可能就有各自的版本的演繹.
就會帶來很懸殊.
所以他們現在在這個時候.
直接這樣問.
真的給我們有一個很好的啟發.
所以你會看到神就透過過去.
用了差不多大半張聖經.
來講回歷史.
講回為什麼神要來到去.
循循善誘的透過先知.
教導百姓.
一些真理和民族之間.
彼此相處的生活.
如何能夠和睦共處.
在那個憲制過程.
其實就是透過.
其實舊日已經有信心了.
透過這樣的憲制守節期.
其實是一個民族慢慢對真神.
有了認識.
這樣來慢慢展現到神的真理.
在一些未認識神的外邦民族.
或者迦南的其他民族.
但是我們看到的就是.

$^{881}$人性在進入迦南之後.
慢慢慣性了埃及做奴隸的時候.
那種困苦那種受欺壓已經忘記了.
因為已經安居樂業在迦南地了.
也慢慢得到神眷顧.
來到去在生活里.
但是慢慢人性就是.
慢慢又是弱肉強食了.
慢慢又是大家來到你娛我詐的.
一些商業行為了.
慢慢慢慢又是建立自己的勢力了.
慢慢又是強取豪奪了.
慢慢又會用手段來去欺壓了.
來去強搶了.
對其他人剝削了.
慢慢又去人性那種.
不知不覺在沒有神底下的生活了.
建立自己的王國了.
有些人就越來越有錢.
有些人就越來越貧窮.
扶貧的系統又不根據神的教導.
因為聖經有教的.
有錢的你就要刻意扎加衣.
將你自己有剩的要給那些窮人的.
這個世界不會斷絕窮人的.
但是這個世界是可以有人有錢的.
有錢的你怎麼能夠建立一個扶貧連門的系統呢.
這些在舊約摩西的五經有教導的.
真的這樣來到要有錢的.
要顧惜那些窮乏的人.
但是現在我們看到的就是.
他們問這個問題.
神就打壓著這個機會.
藉著先知來到去有這個啟示.
就是將過去以色列人遭遇的神的刑罰.
說出來.
是不是神是一位喜怒無常的神呢.
喜歡虐待自己的子民呢.
一定不是.
一定是他們是不聽話的百姓.

$^{921}$還有不斷神就透過先知來到去警告.
來到去預言.
來到去叫他們離開自己的罪孽.
都是人際關係的問題.
人和人交往的問題.
也都是息息相關的.
就是說任何一個時代都是這樣的.
你怎麼能夠可以讓整個民族.
每個人一定會有不同的命運.
三衰六旺一定有的.
有些人天生是殘缺的.
聖經也有說的在舊約.
所以那些做祭司的兒子如果有殘缺.
就不可以做祭司的.
他可以在爸爸家裡吃那些祭物.
但是他不可以做祭司.
因為他身體有殘缺.
有教導的.
就是意味著神沒有做任何的人.
就像複製人一樣.
每個都是完美的.
每個都是智商150的.
每個都是很帥的.
是沒有這一樣的.
有些人真的天生殘缺的.
是不容易的.
正正是這樣.
正正就是人生.
正正是因為罪.
帶來整個世界受了審判.
所以來到這個世界.
每個人的命運都不同的.
也都正正是這樣.
才可以發揮到那種守望相助.
互相的關愛.
彼此的扶持.
在有的蒙恩的.
感覺自己是幸福的.
如何去幫助那些所謂的less fortunate.
沒有這麼幸運的人.

$^{961}$很好聽的.
其實就是那些很慘的人.
有些人真的很慘的.
你可以想象到.
這個就是那個.
精潔的問題所在.
他們現在很想問.
我做這些宗教行為.
是不是就是相應得到神的祝福和恩典呢.
神就好像暫時不說那些憲制.
齋戒 節旗.
是說什麼呢.
就是說舊約的時候.
你們的先祖.
有沒有聽從先知的教導.
來到有生命的成長呢.
譬如我青少年的時期.
未信耶穌.
小時候就在木屋區長大.
幾乎是入黑社會.
入不到.
家裡很嚴.
但是.
根本上就是你那時候是一個暴風少年.
又對家庭的背景是很奮勢執著.
你很自然就是很多東西.
在讀書又沒有遇到很好的老師.
基督徒的老師.
在那些Band 5的學校.
你喜歡就聽.
不喜歡就請你出去.
不要妨礙.
這樣的學校.
怎麼會有好的老師為我流眼淚.
放了學和我補習.
沒有這些的.
我就是這樣.
那個階段.
就這樣過了人生.
你可以想象得到.

$^{1001}$一定是有很多很多的問題.
但是感恩的地方就是.
不知道為什麼神會遇到我.
尋找我.
讓我能夠得到祂的恩典.
而生命一路一路地改變.
這個一路一路地改變.
是要慢慢將你自己過去的成長.
你的原生家庭對你的影響.
要慢慢在恩典裡面去忘記.
去寬恕.
和去接納.
這樣來到去.
好像殺出一條新的血路.
就成為今天的我.
每個人都一樣.
所以現在神就和舊約的故事.
來到某種程度.
先是責備你現在去提問題的這三個人.
這三個人其實很無辜.
我們去問而已.
沙尼息,尼堅米勒.
並跟從他們的人.
去懇求耶和華的恩.
第七章和第八章.
提了三次去求神的恩.
懇求神的恩.
現在猶太人被撈回去.
就是很想去尋求釋法.
我們現在做的宗教的行為.
是不是神喜悅的呢.
要不要再做呢.
因為已經回到耶路撒冷了.
其實做一定是要的.
因為你是猶太人.
你根本不用問.
一定要做的.
但是他們似乎不只是這麼簡單.
真正背後的意思.
就是說.

$^{1041}$怎麼做才是真正蒙神的喜悅呢.
得到神最重要的blessing呢.
但是他們忘記了.
神其實在這個時候.
在被撈歸回的這個時代.
所以你讀舊約一定要記住.
不同的時代.
你讀的先知書是不同的.
因為你讀的是審判的時代.
就全部是罵的了.
但是現在你見到是被撈歸回的時候.
神就說那個安慰的說話.
和神是說將來的預言的.
所以你就可以想象得到.
整個氣氛.
整個tune是完全不同的.
現在在這個時候.
他們就很想問.
神才是來到去舊事重提.
其實問題就是.
想說回.
怎麼可以得到神的恩.
得到一個地步.
第八章說了兩次.
外邦人會拉住一個猶太人.
我們要去求你們的神.
尋求你們的恩典.
即是說猶太人將來.
不是現在.
將來.
他們得到的恩典.
大到外邦人看到.
見到這群猶太人.
很有blessing.
十個人拉住一個猶太人.
要求你帶我們去.
見一下神的恩典.
有沒有記得耶穌在.
約翰福音第十二章.
有幾個希臘人.

$^{1081}$很想見耶穌.
他們找了白菜大人.
拿彈葉.
很想去見一下耶穌.
可以想象到.
當時耶穌的門徒.
不知不覺.
他們都流露一些氣質.
有耶穌的少少風範.
風采和blessing.
他們跟過耶穌.
有時都沒有那麼暴慮.
沒有那麼激動.
外邦人都見到.
這些門徒和一般的猶太人.
是有些不同.
其實我們回教會.
你們自己都知道.
你回教會久了.
慢慢有一種氣質.
其實外邦人.
即是未信的人.
其實是會觀察到的.
是會知道的.
所以我們慢慢浸淫在.
神的愛里.
我們看事物的角度.
我們不會再那麼執著.
我們不會再像以前那樣.
見到人就捏住他的頸.
你要盡快要怎樣怎樣.
就不會.
所以其實那個精潔.
新約其實已經說得很透徹.
主耶穌已經透過不同的比喻.
說得很透徹.
不過在舊約.
在這個被擄歸回的時候.
他們這樣去尋問.
神就和他們舊事重提.

$^{1121}$告訴這幾個使節.
應該怎樣可以得到.
神的blessing.
以至你們去守節.
那個節可以變成.
歡樂的節.
那個節因為你們的出現.
因為你們的生命.
你們的內涵.
可以是一個很哀傷的節.
可以變為歡樂的節.
這不就是正正那群人的內心.
是充滿了神的恩典.
所以這個舊事重提.
值得我們再細心去想.
首先就是神提到他們.
不聽從前的先知所宣告的說話.
你們是不是應該聽.
第八節.
耶和華又說又臨到撒迦利亞說.
這是第七章的第八節.
這個是對他們的一個直接解釋.
如果你們不是這樣做.
間接地說.
你們的守節期.
齋戒.
哭泣.
都變成是裝模作樣.
變成是一種外表的行為.
本來外表的行為.
不可以說錯的.
因為是那個節本身賦予.
需要的表情.
或者那個態度.
但是.
更深的就是.
那個參與守節的人.
那個朝拜者.
那個朝聖者.
他的內心.

$^{1161}$才是神體重.
所以如果我們六天.
在我們的生意.
在我們的生活.
在我們和我們最親的人相處.
我們是出現了第九節.
你們的列祖曾經就是這樣.
要按治理判斷.
各人以慈愛憐憫弟兄.
壓寡婦孤兒寄居.
全部是那些社會上.
你想想寡婦.
沒有了丈夫.
舊約的時代.
是等於沒有了生計.
是不會有什麼扶貧系統.
是不會有優先上公屋.
你可以想象到.
你的整個社會保障是沒有了.
丈夫是對外的.
賺錢的人.
寡婦何來有工作能力呢.
是靠其他人接濟的.
如果你是寡婦.
更加有母兒女.
就更加慘了.
那你孤兒.
就雙親都沒有了.
寄居的.
就是說你在外地來的.
沒有國籍.
沒有戶籍.
那你可以想象到.
那種的情況.
我們香港太小了.
你想想如果你是在大西北的.
走來深圳打工.
可能沒有戶籍.
很多事.
人海茫茫.

$^{1201}$除非你很聰明.
這樣你能夠成長.
不然很困難.
你可以想象到.
神要我們有一種這樣的.
對別人的憐憫的心腸.
那就是說神是一位憐憫的神.
要按照治理去判斷.
那些貧窮人.
誰都不可以心裡謀害弟兄.
就是有一個奸詐的心靈.
就是存心有些惡毒.
或者心計很重.
對人有些動機是不懷好意.
這些是信仰最大忌.
所以信了主就可以.
將自己的心靈擺在神面前.
讓神的愛去充滿.
去填補.
甚至那些不能改變的冤仇.
都可以化解.
你就會變成有神的那種恩典.
有那種公義.
你是一個和睦的人.
所以這樣提.
這些比那個憲制更加重要.
第十一節.
他們就是舊事重提.
先祖未被神懲罰之前.
居住在聖地.
他們是入迦南.
一路由摩西帶領之後.
若孫一路他們在那裡殖民.
在那裡開始他們的王國.
開始有士司時代.
接著去到大魏的王國時代.
一路一路去到最後的.
黃西帝家的時候.
所以一段相當長的.
如果大魏王到被擄.

$^{1241}$大概有千一年左右.
或者一千年.
差不多一千至一千一年.
因為數字不太準確.
你可以想象到.
是一個很長的歲月.
經過很多的世代.
但慢慢地.
他們不是因為在埃及學壞.
其實不學也壞.
人性就是這樣.
所以他們入到迦南.
其實最重要就是.
學習摩西的律法書.
守十芥.
定時定後去敬拜神.
那個心靈就會在神的.
無論律法 聖靈底下.
先知的教導底下.
會慢慢有了這種神的心腸.
神的性情.
舊約很少這樣說的.
新約就經常這樣說.
我們要將我們的心志.
改換一生.
穿上新人.
這個已忽所需四章.
二十九節左右.
這個新人是照著神的形象做的.
是有什麼?.
真理的人意和聖潔.
很清楚的.
保羅說要我們像阿當一樣.
新的阿當.
阿當就是照著神的形象做的.
現在保羅說.
你們信了耶穌.
已忽所.
我們十一月就會去已忽所.
你會看到那個戴安娜女神.

$^{1281}$那個神廟等等.
保羅就教.
信了主耶穌.
我們的心志要改換一生.
不是換心手術.
是因為有聖靈進入.
我們就要改換我們的價值觀.
我們整個人的想法.
要穿上一個新人.
這個新人不是穿著新衣服那麼簡單.
是要有神的形象.
就是真理的人意和聖潔.
你就可以想象到.
新約就教很多.
但舊約就很實在.
和一群以色列人說.
一說就說回到歷史.
回到從前.
回到看猶太人自己.
如何來到.
一直每放月下的離開神.
但是有趣的地方就是.
他們一直到西帝家的年代.
他們都是不斷固定的守節旗.
不敢不守.
你想想他們不敢剃梳.
因為舊約禁止.
他們一定要守安息日.
誰不守就很大件事.
慢慢被那些只是外在的宗教的行為.
牢牢困住自己的心.
以為做了這些就心安理得.
但實際真正的敬拜.
原來是以靈以成.
正正就是那個問題的癥結所在.
所以原來舊約已經反映新約的教導.
新約的教導其實在舊約的時候.
根本就是那個心態的問題.
所以他們如果不聽律法書.
第十一節就說.

$^{1321}$他們卻不肯聽從.
扭轉肩頭.
失而不聽.
如何扭轉脖子.
凌轉面.
失而不聽.
使心硬如金剛石.
不聽律法和萬君之耶和話.
用靈直從前的先知所說的話.
即是說原來每一次去到節旗.
當先知利未人要讀經.
要來教導律法的時候.
要來彼此苦心自問.
過去一個星期.
大家生活的時候.
大家就立刻扭轉肩頭.
立刻塞住耳朵.
其實現在不用戴耳機都可以了.
基本上一說到某些東西你不喜歡聽.
你自然聽不到.
人可以發明到.
訓練到自己.
可以是聽不到的.
所以很多時候報告.
人們聽不到.
跟他說有報告過嗎.
其實每個星期都報告過.
但很奇怪就是這樣.
所以重點仍然是.
你想想.
其實是頭腦.
耳朵.
其實是將聲音波段進入去.
變成有意思的文字.
你聽得明白廣東話才聽得到.
其實重點原來是心.
令自己的心很剛硬.
心已經很麻木.
這是信耶穌.
信主最大忌就是.

$^{1361}$心已經硬了.
很難.
比一些乾硬化還嚴重.
因為你的心已經硬到.
金剛石.
你怎樣弄都不會爛.
就是不聽神的律法書.
所以最終神要用.
極大的烈怒.
使他們面對荒涼.
才能夠重新反思.
要用一個這麼大的手術.
令硬的心慢慢變得柔軟.
所以你會想象到.
做物主使選民要被捕.
聖殿要被拆毀.
聖殿裡面的縣制器皿.
是分別為聖的.
都要被帶到外邦.
帶到巴比倫的廟里.
這樣供奉.
他們都很清楚.
這些是神聖的器皿.
那些巴比倫人都不敢去動.
人家都有些神心.
不敢去動器皿.
他們都不敢這麼做.
但是猶太人就慢慢.
對自己的宗教掉以輕心.
所以我們都感受到.
神的內心對民族.
那種一路一路的忍耐.
一路一路的.
在接觸神的憤怒.
但是神就來到.
仍然給他們機會.
所以第八章的頭段經文.
就很特別.
因為現在的烈怒.
轉變成為.

$^{1401}$神大發熱心.
為石安心裡面.
極其的火熱.
我為他火熱.
向他的仇敵發熱露.
這是最緊要的地方.
所以如果我們將它化做.
最簡單的理解.
首節旗最重要就是.
知道神他在這裡.
我們做的這一切.
是因為我們和他的關係.
希望是做在他的喜悅裡面.
這是最緊要.
我們現在有的節旗就是.
我們的聖餐.
我們每一次每一個月.
假如你定了是晚堂的.
你來到聖餐的時候.
你差不多好像旁若無人.
不要理會其他人.
你自己帶著一個虔誠的心.
為什麼我要吃這些餅.
和要喝這個小小的杯.
為什麼要這麼做.
為什麼要這麼安靜.
敬虔的來做.
原因就是有一個歷史的事件.
是和我這個人有關的.
如果周圍都沒有人需要信耶穌.
只有我這個人是罪人.
耶穌都會為我釘十字架.
這個時候.
這個十字架的救恩就很重要.
所以我來到首節的時候.
其他的節就不重要了.
我的心要完全集中.
在這個短短的時間.
來思念我的主耶穌.
為我的犧牲.

$^{1441}$他是甘心樂意.
他為我的救恩來到大發熱心.
他心裡面極其的火熱.
所以這裡就是.
這幾個時節.
聽完這番話.
一方面不能不認同.
因為歷史已經發生了.
我們的先祖就是這樣.
受到神的刑罰.
現在他們就是活在一個.
神現在是心裡.
為他的聖神火熱的一個時代.
所以現在他們一定是.
做得很甘心樂意.
當然.
慢慢地.
當他們又故太復文的時候.
又會慢慢地改變.
所以我們之前看的以色列記.
正正就是這樣.
第二批過了二十多年之後.
再上來的.
再去到尼希米再上來的.
又是會有這樣的情況出現.
尼希米再過久一點.
過了一百二十多年.
你可以想象到.
每一個時代.
都有我們對神的背影.
有我們慢慢地成長.
有我們慢慢地明白神的心腸.
明白神對我們的極其的火熱.
我們要來到去.
對神.
我們就不只是一種.
外在外表的宗教行為.
而是我們真的有一個.
很單純的心靈.
不為什麼.

$^{1481}$就是要來到去.
尋求和敬拜我們這一位主.
所以焦點就放在神那裡.
所以第八章就告訴我們.
現在是神自己重新回到石安.
這是第三節.
要住在耶路撒冷.
耶路撒冷到時.
就稱為誠實的誠.
所以是因為神.
而帶來生命改變.
所以只是靠人.
怎麼可能做到這些慈愛.
憐恩.
任何的福利系統.
都有它的瑕疵.
有它的問題.
任何的.
無論民主極權.
或者是共產.
什麼的政權.
都是會有它的瑕疵.
而且人.
始終是人.
所以就算整個國家是基督教.
都不會有好的基督徒.
是沒可能的.
因為涉及的是很多公共的政策.
每個人都有不同的尺.
所以說自己是基督教的國家.
其實做的事情一點都不基督教.
所以一定會是這樣的.
所以我們都不需要有些憧憬.
或者有些幻想.
基督教國家會好一點.
有時拆開來說就是.
基督教的機構.
有時更加mean.
對下屬有時更加剝削.
當然我不敢說那麼多實際例子.

$^{1521}$免得.
但是的確有人真的這樣經驗到.
感覺到.
所以其實正正就是因為.
那個是人性.
和有它的限制.
但是最重要就是.
為什麼耶路撒冷將來.
可以成為誠實的誠呢.
因為神.
祂再眷顧石安.
祂要住在耶路撒冷.
就是說神再回來.
這些就是我們對將來.
在這個新天新地來到之前的.
千禧年國度.
耶穌基督再回來.
就可以改變人心了.
再眷顧耶路撒冷了.
所以現在我們千萬不要太逆差.
這個是不是就是.
現在耶穌回來的時候.
第三次世界大戰.
哈曼吉多頓.
現在是不是猶太人.
最終要被滅國來圍困.
然後他就要呼求耶穌.
可能是.
但是我不敢逆差.
因為我說任何的東西.
都是對雙方面的死者.
已經過萬人.
是非常不敬.
對於現在這樣的一個.
這麼嚴峻的危機.
戰爭的危機.
和現在人道的災難.
盡量不要說我們的評論.
盡量現在大家要真是認真祈禱.
一定是有不對.

$^{1561}$有雙方面的錯.
有積累歷史的錯.
各方面的了.
但是現在我們要的.
不是去做評論家.
而是我們每一個基督徒.
真的大家一起祈禱.
昨晚的祈禱會.
回來有19個.
有同工告訴我.
在線有98位.
一直在線.
就是一直從頭到尾.
都有27位在線的登錄.
當然我們不是要求數字.
而是我們能夠做的就是.
至少你在一個這麼安全的地方.
神的電.
燈火通明.
我們這樣來祈禱.
是我們最低的付出.
你想想.
在一個現在完全沒有水沒有電.
沒有食物沒有燃料.
不斷的轟炸.
那種恐懼.
生活的.
誰錯誰對.
根本不要再說了.
要人道救亡.
所以現在這個時候.
我們就要祈禱.
很特別的就是.
剛好我們讀到撒迦利亞書.
不是特意為它來講.
而是真的讀到原來神.
給這幾個時節要知道.
時代不同了.
過去你們有一個大錯誤.
就是你們是神的選民.

$^{1601}$你們是住在聖城.
但是你們的心.
沒有一直一直成長.
以致慢慢那些舊的人性.
那些罪性又再出來.
慢慢安居樂業之後.
就不滿足自己的那份.
就會踩人家的街.
去搶奪別人的所有的.
由欺壓那些孤兒寡婦.
慢慢就變成.
社會上的結構性的罪.
有錢的就越來越富有.
沒有的就越來越慘.
你可以想象到.
如果你在香港.
你有五六套房.
都不知道住哪裡好.
每年要交的.
差一樣交很多.
你忙到那些在tong 房住.
很可憐很可憐.
說完之後.
你又繼續繼續.
這個問題就是.
如果你又是基督徒.
你有沒有想過.
可以多走一步.
記住你有的不是罪.
是神給你的恩典.
你有這樣的資源.
但是你有沒有想過.
可以多走一步.
去幫到有真正需要的人呢.
這些就是.
他們就是沒有了.
而且變本加厲的.
去欺壓了.
所以神就透過.
先知已經提出.

$^{1641}$你至少建立一些.
社會的扶貧制度.
至少你們每一個.
去到首節的猶太人.
那個是你的同胞.
是你的弟兄.
你不單止沒有憐憫.
還欺壓他.
你想想神.
看到這樣的心情.
他會怎樣呢.
這個就是.
舊約過去的歷史.
但是現在神在這個時候.
被老歸回.
他就提出了一個很緊要的.
一個預言的真理.
就是這個城.
其實最終會成為.
誠實之城.
不是罪惡之城.
有些人巴西就是.
特別里約熱內努.
那些貧民區.
就是罪惡之城.
差不多好像以前的城寨.
但是原來耶路撒冷.
真正的誠實之城.
是耶路撒冷.
你很難想象.
現在我們看不到結局.
看不到將來的預言.
但是現在.
聖經就已經告訴我們.
這個城會成為誠實之城.
原因就是神住在這裡.
所以基督是我們一家之主.
不要只是那個牌.
而是我們認真的落實.
就是我們的兒女.

$^{1681}$看到我們的改變.
就很不同了.
譬如我們看到.
有小朋友跟父母回來祈禱會.
回來一起敬拜神.
其實他們耳喧目艷.
他們長大.
當他們自己獨立.
成家立室的時候.
他們會感受到.
這一份父母輩的blessing.
他們會有這一種.
對神的愛.
對人的愛.
這些一定會有的.
所以現在很重要的就是.
神在眷顧石安.
是最大的祝福.
所以教會我們要求神的恩典.
祂的同在在聚會當中.
在我們中間.
我們慢慢.
如果大家都是心存謙卑.
大家焦點都是對準神.
教會就很順暢.
因為根本不需要爭拗.
哪個意見都無所謂.
總之教會一直帶領著.
我們跟真理.
這樣走就行了.
總有一些對.
有一些沒那麼好.
沒絕對.
但是如果我們要找那個絕對.
叫餐叫死.
根本就不是神的喜悅.
所以這裡告訴我們.
最重要的焦點就是.
到底神是不是在我們中間.
眷顧我們這個城.

$^{1721}$所以那個光景就是.
滿了老人家.
城中街上滿有男孩女孩在這兒玩耍.
你可以想象那個景象.
是多麼美麗.
人是最重要的動物.
整個耶路撒冷城.
曾經的荒涼.
現在是充滿人口.
而且是歡樂的.
不是驚嚇的.
不是在這裡走那個恐雜的.
是很開心.
是沒有人來侵害的.
你現在製造的浪漫.
製造那些蔓延著的恐怖主義.
你越是這樣來到去狂轟濫炸.
現在連阿根廷.
美國的領事館.
以色列的領事館.
都收到炸彈恐嚇襲擊.
你可以想象到.
一直在蔓延著.
整個中東的國家.
都開始來到不斷上街示威.
燒美國國旗.
以色列的國旗.
你可以想象到.
現在以色列充滿的是那個懼怕.
充滿的是那種擔心.
可能有些人都會移走.
但是似乎就很少會這樣.
因為他們的民族性也很強.
但是這些是神才可以做到.
一個誠實的誠.
描述裡面的安居樂業.
國泰民安的歡樂.
所以來到教會.
弟兄姊妹.
要立刻將自己擺作.

$^{1761}$即是立刻圖傳自己.
我是神的兒女.
這些要是一些很受壓迫的民族.
會有這種形態浮現流露.
我剛才提到黑人的教會.
因為他們在美國一直受壓制.
他們是很多因素.
令他們是社會的最底層.
有些根本你不可以想象.
他有明天的.
18歲都未必夠.
就已經是坐監.
要麼就槍死在黑幫的仇殺.
搶奪地盤等等.
要麼就在監獄裡面.
那些小朋友出世是不會有爸爸的.
你可以想象到整個黑人的社區.
充斥的就是毒品.
很多的仇殺.
黑幫的橫行.
你就可以想象到.
他們根本讀不到書.
一代一代的結構性問題.
所以黑人的教會.
分外他們的感受就是.
當牧者根據聖經告訴他們.
你們在外面的生活受壓制.
但是你們來到教堂.
你們是神的兒女.
所以這個不是要麻醉他們.
而是這個真的給他們一個生命的出路.
無論我在日常的生活.
我們的結構性受到剝削.
但是我來到教堂.
神告訴我.
我真是祂的兒女.
我是有盼望的明天.
這個就是.
可能某程度的癥結所在.
如果我們都不稀罕.

$^{1801}$這個抽象的身份.
那你日常的生活.
那個身份就緊要了.
你就很緊張你的身份.
你就很緊張人怎麼看你.
所以你如果不好.
你的卡片不夠亮麗.
那你就覺得很自卑.
很憤世嫉俗.
很內向的避開.
很扭曲.
那你就可以想象到.
這個就正是為什麼有些社群.
有些教會.
會分外有一種活潑.
和興旺的原因.
如果大家有機會.
根本你可以隨時去體驗.
你去一下一些印尼教會.
譬如我們有的.
或者你禮拜.
你下午去教育路的基福堂.
很開心的.
他們敬拜神.
又跳舞呈獻給神.
他們又唱詩歌.
跟著會有人每個禮拜.
都會出來分享他們的見證.
他們的生活見證.
有些可能是很小的事.
但他們大家彼此是很投入.
會一起講阿們.
那你就可以想象到.
他們的氣氛.
分外有一種力量.
其實某程度上是一種醫治.
心靈的醫治.
所以我們回來.
這些又不可以人為製造.
我們要提醒.

$^{1841}$勸勉.
但如果整個教會.
我們很精英化.
個個都是叻人.
個個都不需要像神的眷顧.
我都很好的.
有時未必那麼看重.
那個永恆身份.
當然我們絕對不是要.
落低到社會基層.
才會這樣感受.
也不需要那麼極端.
仍然是那個心態.
到底什麼才是永恆.
什麼身份才是傳到永遠.
什麼身份才是神.
會真正喜悅的.
這個就是問題的關鍵.
所以終極的關懷.
應該不是地上的短暫的得失.
我們重要的.
我們要忠心的.
但真正永恆的才是永恆.
所以這裡告訴我們.
將來這個民族的良厲.
不是因為他們有高的科技.
個個都是很叻的精英.
而是因為神再回到耶路撒冷.
神再眷顧石安.
我讀這一段之前.
和大家再補充一點.
剛才都說想和大家讀的.
就是羅馬書第二章的經文.
其實可以去到下面第八章.
再講都可以.
不過現在先讀了.
其實大家都已經有過.
上過羅馬書你就會知道.
保羅其實第二章就試探討.
到底猶太人他自己的同胞.

$^{1881}$有什麼要悔改呢.
愛邦人的罪相對比較明顯.
拜偶像.
但是猶太人不拜偶像.
但又不得到神的恩典的原因.
就是第二章的十七節這一段.
算是挺重的.
好像他有一個對手.
就是有些猶太人在他面前.
你稱為猶太人.
又依靠律法.
且指著神去誇口.
既從律法中受了教訓.
就曉得神的旨意了.
又能夠分別是非.
又深信自己是吸核子領路.
就是愛邦人.
那些人是沒有神的.
沒有宗教的.
沒有永生的.
那他就是核子領路.
他是開眼的.
核子是那些愛邦人.
是黑暗中人的光.
是蠢笨人的師傅.
是小孩子的先生.
保羅這樣去描述.
好像荃方偉那樣.
你全都厲害了 猶太人.
你是有律法 有神的旨意.
你都知道了.
不過第21節.
你既是猶太人.
你既是教導別人.
還不教導自己嗎?.
你講說人不可偷竊.
自己還偷竊嗎?.
你說人不可姦淫.
自己還姦淫嗎?.
你厭惡偶像.

$^{1921}$自己還偷竊妙中之物嗎?.
你指著律法誇口.
自己倒犯律法.
佔欲神嗎?.
神的名在愛邦人中.
因你們受了褻瀆.
其實是很大的罪.
很大的事.
愛邦人認識不到神的原因.
是因為猶太人.
不是一個漂亮的示範單位.
他們表現出的.
可能比愛邦人更加的.
來到去犯律法.
可能做生意更加的.
對別人剝削.
就變成了原來.
他們就是好像剛才第七章所講的.
他們的內心.
沒有憐憫.
沒有顧念其他人.
內心也充滿了詭詐.
要來剝奪.
自己的同胞.
的弟兄.
這樣的一種強搶.
的一種強搶.
你可以想象到.
愛邦人因為猶太人.
而來到去褻瀆神.
正如經上所記.
你如果行律法.
國禮固然於你有益.
開始講那些禮儀.
講那些宗教的行為.
如果犯律法.
你的國禮就算不得國禮.
國禮對猶太人來說是很重要.
第八天行國禮是很敬虔的.
保羅是第八天行國禮.

$^{1961}$耶穌也是第八天行國禮.
所以第26節.
所以那些未受國禮的愛邦人.
如果遵守律法的條例.
他雖然未受國禮.
豈不是算是有國禮.
本來講的那個精神.
而且那本來未受國禮的.
如果能夠全守律法.
豈不是要審判.
你這些有疑問.
和國禮.
警犯律法的人.
有疑問.
塗炬外表.
塗炬儀式.
甚至說得比較俗一點的廣東話.
就是只有一把口去說.
談論宗教.
談論律法.
什麼節日.
就可以隨隨便便說出來.
所以.
那幾個博士.
在東方上來.
去哪裡.
一定是去皇宮.
找王子當然是在皇宮.
沒理由去貧民區.
一定是去皇宮.
一定是找猶太人的王.
所以就去找希律王.
大希律.
希律王就很害怕.
生下來做猶太人的王.
我已經是.
是羅馬的分封的王.
他就去找那些文士.
叫那些文士來釋法.
究竟在哪裡說.

$^{2001}$他自己不讀經.
就問那些文士.
基督在哪裡.
那些文士.
揭都不用揭.
就說尼加書五章二節.
大衛的城伯利行.
在那裡會有救主出生.
多快就回答得到.
但他們自己又不去伯利行.
那幾個博士就去.
去到的時候.
路加芬記載.
去到的時候就見到瑪利亞.
接著他就拜瑪利亞.
和那個小朋友.
就是耶穌.
你可以想象到.
保羅在這裡.
他很明白.
他自己是很認真守律法.
但他只是明白了那個律法的表面.
他明白不了背後的深度.
將來律法所.
將來要來的人.
那個尼賽亞就是耶穌.
他是沒有罪.
他才可以讓這個新生命.
可以在我們每一個信的人.
內心去滋生.
所以猶太人.
根本連保羅都沒有到那個水平.
這些猶太人.
連基本的守律法都沒有.
他們就口口聲聲說.
我們是有律法的一個民族.
所以保羅在這裡.
是28節.
很重要的各位.
因為外面作猶太人.

$^{2041}$不是真猶太人.
外面肉身的割離.
也不是真割離.
唯有裡面作.
才是真猶太人.
真割離也是心裡.
在乎靈不在乎疑問.
這人的稱贊.
不是從人而來.
是從神而來.
只要心裡做.
做給神看.
所以我們所謂的感宣人.
如果我們真的要這樣說.
我們不是想和其他教會分開.
感宣人.
我們是否真的從心裡是感宣人.
感宣是怎樣的呢.
我們是否很著重傳福音.
彼此相愛.
很著重去猜傳.
很著重去來到大家在聚會里.
彼此服事.
這是定義.
所以如果在外面做感宣人.
就不是真的感宣人.
真的感宣人.
是要做內涵的事.
所以這是很重要的.
生命的價值.
保羅其實好像將.
剛才那段經文.
做了一個他個人的反省.
他自己的信仰反省.
他自己其實.
是認真守律法.
他也感受到.
他自己守律法.
其實也做不到神律法的要求.
但是他一直等待.

$^{2081}$以為拼命做好就行了.
誰知道原來.
不需要.
原來.
這個義.
是憑信心就行了.
義人是因信得生.
保羅自己這麼說.
有那個信.
因為已經有一個義.
是成就了.
你只要信這個義.
你就可以得生.
那個義就是耶穌基督的救贖.
所以我們回來.
這個撒迦利亞書.
其實是一脈相承.
救約神就已經這樣說了.
要記住.
如何能夠讓.
這個生命能夠貼近神.
就是要.
認同神的同在.
帶來的真正的祝福.
神的同在.
是會慢慢改變.
我們的生命.
順便讀二弗所書四章這一段.
也是重要經文.
就是說.
新舊生命的轉變.
第25節.
不是,原來是17節.
這一段.
四章的17節.
所以我說.
且在主裡面確實的說.
上文剛剛.
上個禮拜的講道就說了.
我們要發現.

$^{2121}$每個信主的弟兄姊妹.
神都會讓你.
有恩典.
的恩賜去服事.
你要如何發展.
發揮你的恩賜.
需要有專職的牧者.
來讓你去發展.
然後.
你要來到在愛裡面.
一同的.
百折各按各跡.
在愛裡面.
建立身體.
就是基督的生命.
就是我們每一個基督徒.
接著17節.
很具體的教導.
如何來用恩賜.
來彼此在愛裡面建立.
所以我說.
且在主裡面確實的說.
你們行事.
即我們平常的生活.
不要再像外邦人.
我們以前是外邦人.
其實他們一直都沒有改變.
就是外邦.
不是猶太人.
但心態不再是外邦人.
即不再是.
未信耶穌的那種狀態.
不要再像外邦人.
傳一個虛妄的心行事.
即天不怕地不怕.
叫虛妄的心.
他們心地昏迷.
內心完全沒有良知.
於是.
他們就會.

$^{2161}$失去.
良知.
與神所賜的生命隔絕.
所以這是.
心冤.
信和不信是一個心冤.
所以如果不信.
就是與神的生命隔絕.
是一線.
表面是一線.
其實是一個心冤.
都因為自己的無知.
心裡面剛硬.
良心既然喪盡.
心地昏迷.
放縱私慾.
貪行種種污蔑.
你們學了基督之後.
就不是這樣了.
如果你們聽過他的道.
領了他的教.
學了他的真理.
也是上文所說的.
有專職.
牧者的教導.
所以這裡.
是不斷告訴我們.
教會很重要.
就是真理.
所以真理是教會的.
儲石和根基.
所以他們.
有這些生命的教導之後.
就要脫去.
你們從前行為上的.
舊人.
好像用一個衣服的.
概念的比喻.
這個舊人.
是因為私慾的迷惑.

$^{2201}$漸漸變壞了.
又要將你們的心智.
改換一新.
是心智裡面的.
心靈改換.
並且穿上新人.
這個新人.
是照著神的形象.
做的.
神賜給我們一個新的形象.
就是說我們是新的阿當.
阿當就是神的形象.
做的.
現在我們有一個新的形象.
所以很多信了主耶穌的人.
上也變了.
是.
和諧很多.
和藹可親很多.
笑容是很親切的.
信了主耶穌的.
真的會慢慢.
你的內心.
是改換了一新.
有之於內.
就盈之於外.
真是這樣的.
所以是照著神的形象做的.
有真理的人意.
所以你是對真理.
是有一個這樣的.
渴求.
和聖潔.
和你會來到.
有神的心.
的心腸.
所以在行事為人.
有些事你自己問自己就行.
有些人挑戰.
牧者可不可以吃煙.

$^{2241}$沒有教導說不可以吃.
可不可以打牌.
問這些好像很實在的問題.
其實不用問.
神喜悅你你就做.
如果不喜悅.
你是否符合真理.
的人意和聖潔.
是保羅很明顯地教導.
不明顯.
更加明顯.
你以前是偷東西的.
你是小偷.
你不要再偷了.
就要勤力做事.
甚至有餘錢.
可以去糾纏其他人.
說得這麼白.
真的有些人是小偷.
在信主之前.
是這樣的.
你可以想象到.
都是提到內心.
容易發憤.
有些人可能.
有各種因素他要偷錢.
他要去偷東西.
你可以想象到.
有些人很容易就有.
苦毒很容易就會惱人.
很容易就會憤怒.
很容易就會讓罵.
讓罵就是說是說非.
搬弄是非.
在這裡來到去.
就是暗箭傷人.
挑撥離間.
又毀謗.
就是這樣.
明明不是的.

$^{2281}$講得加鹽加醋.
並一切的惡毒.
都要在你們中間除掉.
跟著要以恩慈相待.
傳一個憐憫的心.
彼此的堯恕.
為什麼呢.
正如神在基督裡面.
堯恕你們一樣.
多麼實在的生活.
新的守則.
好像信了耶穌.
你就要根據這些守則來生活.
重點就是因為裡面有了.
一個新的心智.
裡面穿上了一個新人.
是照著神的形象.
所以我們是反映著神的性情.
所以難怪第八章.
他們如果活得好.
當然這是預言將來的.
以色列.
他們有令到外邦人.
要懇求神的恩.
所以這兩章.
講了三次懇求神的恩典.
第七章的開始.
就是他們有.
使節的訪問團.
派了沙利息.
和利堅米勒.
去懇求耶和華的恩.
很想得到神的恩典.
怎樣得到呢?.
我們現在這樣做的守節.
五月七月的齋戒.
對不對呢?.
怎樣做得再好一點呢?.
神說不是外表儀式.
是內心的更換.

$^{2321}$改換你的心.
所以這個是新約.
其實保羅都已經教了給我們知道.
今晚我們就講到這裡.
我們下次再補充一點.
就進入第九章.
去到第十四章.
是有對列國的審判的.
多謝主耶穌.
讓我們今晚.
可以透過第七第八章.
我們可以來看到.
神一路在每一個時代.
對每一個民族.
都是很想我們去學習.
你自己的生命.
性情.
多謝你讓我們有救恩.
因為主耶穌的救贖.
我們得到重生.
我們可以成為神的選民.
我們能夠有一個.
改換了的心智.
幫助我們能夠明白.
不是外表的一個儀式.
是內心的態度.
以致當我們的內心.
是有公義.
有誠實.
有平安.
這樣貼近神的心腸的時候.
和感受到神的同在.
帶給我們的信息.
和神的同在.
帶給我們的福樂的時候.
我們自然就會成為別人.
懇求恩典的一個.
管道一個媒介.
願主幫助我們.
讓我們能夠繼續.

$^{2361}$可以將我們所得到的.
美好的恩典救恩.
去告訴其他人.
多謝你讓我們今天晚上一起.
學習真理奉耶穌基督的命.
阿們.
\newpage



\section{馬太福音 10:34-42}
\label{sec:qZ9s1IR331E}
\textbf{2023-10-21 晚堂《作主門徒》馬太福音 10\_34-42 講員:陳麗蟬牧師}
\newline
\newline
連結: \href{https://youtube.com/watch?v=qZ9s1IR331E}{\texttt{https://youtube.com/watch?v=qZ9s1IR331E}} ~~~~ 語音日期: 2023-10-24
\newline
\newline
\hyperref[sec:c_6waYrNYvI]{\small{< < < PREV SERMON < < <}}
~
\hyperref[sec:index_chronic]{\small{[返順時目]}}
~
\hyperref[sec:index_scriptual]{\small{[返順卷目]}}
~
\hyperref[sec:ZmEDOfaqivc]{\small{> > > NEXT SERMON > > >}}
\newline
\newline
馬太福音 10:34-42
\newline
\begin{longtable}{cl}
\hline
\hline
章節 & 經文 (和合本修訂版)\\
\hline
10:34 & \begin{tabularx}{0.7\textwidth}{X} 「你們不要以為我來是帶給地上和平，我來並不是帶來和平，而是刀劍。 \end{tabularx} \\ \\ \relax
10:35 & \begin{tabularx}{0.7\textwidth}{X} 因為我來是要叫『人與父親對立，女兒與母親對立，媳婦與婆婆對立。 \end{tabularx} \\ \\ \relax
10:36 & \begin{tabularx}{0.7\textwidth}{X} 人的仇敵就是自己家裡的人。』 \end{tabularx} \\ \\ \relax
10:37 & \begin{tabularx}{0.7\textwidth}{X} 愛父母勝過愛我的，不配作我的門徒；愛兒女勝過愛我的，不配作我的門徒。 \end{tabularx} \\ \\ \relax
10:38 & \begin{tabularx}{0.7\textwidth}{X} 不背自己的十字架跟從我的，不配作我的門徒。 \end{tabularx} \\ \\ \relax
10:39 & \begin{tabularx}{0.7\textwidth}{X} 得著性命的，要喪失性命；為我喪失性命的，要得著性命。」 \end{tabularx} \\ \\ \relax
10:40 & \begin{tabularx}{0.7\textwidth}{X} 「接納你們的就是接納我；接納我的就是接納差遣我來的那位。 \end{tabularx} \\ \\ \relax
10:41 & \begin{tabularx}{0.7\textwidth}{X} 把先知當作先知接納的，必得先知的賞賜；把義人當作義人接納的，必得義人的賞賜。 \end{tabularx} \\ \\ \relax
10:42 & \begin{tabularx}{0.7\textwidth}{X} 無論誰，只因門徒的名，就算把一杯涼水給這些小子中的一個喝，我實在告訴你們，他一定會得到賞賜。」 \end{tabularx} \\ \\
[1ex]
\hline
\hline
\end{longtable}
$^{1}$各位姐妹平安.
有一位英國的記者.
他曾經得過攝影獎.
在2021年的時候.
他在網上發布了一段短片.
這裡是其中的一張照片.
講到一個美國的軍人.
他拿著一瓶礦泉水或者一瓶水.
他蹲在地上.
他就去問一些小朋友.
這些小朋友都是阿富汗難民的小朋友.
他們在機場等著.
將要飛去不同的地方.
這些小朋友和他們的家長.
其實在那裡的地方已經等了起碼兩天了.
他們不僅沒有地方去洗手間.
也沒有水喝.
整個地方就烏煙瘴氣 很辛苦.
其實水對於當時的美國的年輕軍兵來說.
都是很珍貴的.
因為有限量的去配給的.
但是對於這幫逃難的人.
無論大人還是小朋友都好.
他們更加是難能可貴的.
當年輕的軍兵逐個小朋友問.
你想不想喝一口水給他們喝的時候.
小朋友的回應很可愛.
不同很天真地說謝謝.
就帶來了一些父母周圍的人的笑聲.
令到整個場景原本是憂心忡忡的.
原本是煩躁不安的.
是焦急的 不知道路向的.
但是在那個情況下.
就有一些歡笑聲.
點點的歡笑聲就發了出來.
剛才讀的那段經文是在馬太福音的.
馬太福音就在強調這個門徒要出去.
為主你去作見證.
先將這個福音傳給以色列家.
就是猶太人.

$^{41}$然後最後一直傳到去萬民.
給全世界的人都知道.
這個馬太福音的主題是「以馬來利」.
當在第一章講到耶穌基督出生的時候.
祂就說「必有同女懷孕生子.
人要稱祂為以馬來利」.
以馬來利發明出來的就是上帝與我們同在.
到了第十八章.
當耶穌準備升天的時候.
耶穌就頒佈這個大使命.
祂就說「凡我所吩咐你們的.
都教導他們遵守」.
教導這幫門徒.
「看啊 我天天與你們同在」.
耶穌與門徒同在.
以馬來利直到世界的終極.
這個整卷書裡面.
都含有這個「以馬來利」.
神與我們同在.
而剛才第十章.
就是講到耶穌要差遣十二個門徒.
去到不同的地方.
去講這個福音的 講這個真理.
而在這個十章的第34到42節.
祂就對這個門徒有一些要求.
就是要他們將生命的主權交給上帝.
一生要尊上帝為首.
於是上帝的同在.
上帝的恩典就會臨到他們了.
我們看看第一方面.
就是這個「作主一門徒」的時候.
我們要將這個主權交給上帝.
在第34到39節 祂這樣講.
祂說「你們不要以為我」.
就是耶穌.
「你們不要以為耶穌來到地上高.
會帶給地上高有和平.
我來到不是帶給地上和平.
而是刀劍.
因為我來是要叫人和父親對立.

$^{81}$女兒和母親對立.
婆媳和婆婆對立.
人的仇敵就是自己家裡的人」.
聽到這段經文的時候.
是不是有些人害怕 是吧?.
信了耶穌之後要和家人對敵.
沒什麼可能的.
叫人信耶穌都很難 是吧?.
明明耶穌就是和平之君來的.
祂應該帶來地上高和平的.
為什麼反而會帶來戰爭呢?.
原來為地上帶來這個刀劍.
這個「刀劍」這個字.
就是與《耳窟所書》第六章十七節.
它這樣說的.
「拿著聖靈的寶劍」.
這個就是同一個字.
這個「刀劍」和「聖靈的寶劍」是同一個字.
但這個「聖靈的寶劍」是指什麼呢?.
就是上帝的道 上帝的說話.
而這個《希伯來書》第四章十二節.
就說到上帝的道是活潑的.
上帝的說話很有能力的 很有功效的.
比一切兩刃的劍更加的鋒利.
又是這個「劍」.
又是指同一個字.
「刀劍」都是同一個字.
不過它正正式式.
不是真的用來打仗殺人的那些「刀劍」.
而是指上帝說話.
上帝說話那份的活潑.
能夠進入到魂與靈.
還有骨節和骨髓之間.
能夠否入否開.
連心中的思想意念 主都知道.
這個「劍」不是用來讓我們去打架.
而是這個「劍」 這個就是上帝的道.
成為了我們信徒生活.
來判斷事情的準則.
原來在猶太人的文化裡面.

$^{121}$如果有兩個人在為一件事爭執.
無論吵架也好 講道理也好.
旁邊的人都可能有不同的看法.
於是他們就會把一把刀放在地上.
就不用來劈炮 不用來打.
以後我們做一個表決.
贊成A的就站在這邊.
贊成B的就站在另外一邊.
就是說它是用來表明立場.
你究竟是贊成A還是贊成B呢?.
其實有時去做一些事情.
來表明立場是很難的.
誰是誰非 很難的.
我很怕去跟一個親戚過馬路.
他整天不跟 不單不跟紅綠燈.
最糟糕的就是 這裡明明有紅綠燈.
那裡也有的 但當他一路走路.
看到對面的那個店鋪就想去的時候.
他就在這裡直接的走過去.
整天都綁著他 喂 前面有紅綠燈.
他說 這裡直走就到了.
我說犯法的 警察抓的.
他說 人人都是這樣的了 抓了再說.
我說 人人都是這樣的了.
我又看一下 周圍的人個個都站在那裡準備過馬路.
但是一句 人人都是這樣的了.
是不是就對了呢?.
村貴桶底 人人都是這樣的了.
未婚同居 未婚生子為例 泊車 說謊話.
人人都是這樣的了 是不是?.
但是不是人人都是這樣的話呢?.
就是一件對的事情了.
這個刀 象徵著上帝的說話 上帝的真理.
當它放在這個中間的時候.
就是說我們作為信徒的 作為信徒的.
我們所有判斷任何的事情.
都要以上帝的說話來成為我們的準則.
因為當我們在信仰上高.
我們需要去表明立場.
我們不能夠一腳踏兩船.

$^{161}$耶穌說不能夠一方面事奉上帝.
另外一方面又事奉媽媽.
我們需要有一個決定 表明我們的立場.
耶穌說到了.
他說你不能夠愛自己的父母.
你去多於愛我.
當我們有些事物 地上的那些事物.
它吸引了我們.
但是我們竟然愛這些事物.
多於去耶穌的時候.
這個就是我們離開了上帝.
遠離了上帝.
當我信耶穌的時候.
其中一個最大的難處是什麼呢?.
因為過往我們逢星期天都會走過來一起去喝茶的.
但是信了耶穌星期天就要回教會崇拜.
最初他們覺得我很沒有人情味.
不再和他們一起喝茶.
但是後來慢慢就改變了.
我完了崇拜 下午練師.
唱完詩班之後.
臨著四點多回去和媽媽買菜煮飯.
用另外一個方式和他們一起.
特別當我很清楚上帝呼召我去讀神學的時候.
我就和我媽媽說.
媽 我要讀神學.
可能我也不知道什麼叫讀神學.
她的反應就是說.
哦 什麼來的?.
哦 你自己決定吧 去吧.
然後當我知道神學院收了我.
我就告訴我媽和我家人.
我要進長洲讀神學了.
我不教書了.
於是很多雜鼻的聲音就來了.
哇 這麼好一份工作.
人工又高 鐵飯碗來的.
為什麼要去讀神學 搞什麼呢?.
現在大哥就去英國讀書.
我們正在供大哥去讀書.

$^{201}$你不做事 給不給家用?.
哇 很多很多很多雜鼻的東西就出來了.
當時我就和上帝說.
我不想因為讀神學的緣故.
以致和我家人破裂了關係.
我要求我上帝.
可能上帝如果媽媽不答應給我去讀神學的時候.
我不會去交學費的.
到了交學費那天.
我媽一直都不出聲 我都很心急.
但是到了交學費那天 我仍然等.
我出門上學的時候.
我媽就問我.
你今天是不是要去長洲交學費?.
我說是啊.
我就想著等她說 你去吧.
她又沒有出聲.
我搭地鐵的時候突然間就想到了.
咦 她是不是想和我進去呢?.
所以我一回到學校的時候就打給她了.
媽 你是不是想和我一起進長洲?.
她說 不知道檔口生意怎麼樣 看看吧.
如果沒有什麼人吃 我就和你去.
我說不要緊 我放學了 我回來找你.
那天我上課的時候 經常有空我就會去祈禱.
天父上帝求你.
叫我今天爸爸那個檔口最好沒有什麼生意.
當我那天下午去到的時候 我很開心.
真的只有小貓三,四隻在那裡吃東西.
我相信上帝是那位信實的神.
在祂裡面沒有難成就的事.
所以祂一點多的裡面.
整個大排檔只有三,四個人.
但是祂都沒有說話的 都坐船的時候了.
因為肚子餓了 下去最底層.
我們就去買點公仔麵吃.
一邊吃的時候 祂就說.
以後上學你就不要坐在樓下那層了.
因為那麼多粗魯的男人 你不要在那裡 危險啊.
原來祂間接就已經答應了我.

$^{241}$我覺得原來真的堅持到底.
跟隨耶穌的時候 祂一定會開出路的.
要看用什麼方法 你會為我們開出路的.
所以耶穌說 不要以為祂來到地上高.
是帶給地上高和平.
祂來不是的 而是要帶給地上高的刀劍.
因為祂的真理一下來的時候.
很多不是真理的事情 都要顯明出來.
祂就要糾正這些不合乎真理的事情.
在路加福音裡面 同樣記載了這兩節的經文.
在這兩節經文之前 耶穌又做了另外一個比喻.
祂就說到一個有錢人請客的比喻.
祂就請了一些人來 ABCD(ABCD)請了人來.
但是到時間開飯了 竟然都沒有人來.
於是就打發那些工人.
就去快點再提醒那些客人要來了.
大家趕快去叫他 喂喂喂 阿A 阿A.
主人請你去吃飯了 忘記了嗎?.
阿A就說了 我買了塊田了.
我現在要快點去看看塊田怎麼樣.
那他就去找阿B了 阿B就說.
我買了一隻牛 我要趕快去看看那隻牛怎麼樣.
要試一下那隻牛行不行.
阿B就去找阿C了 阿C 你忘記了去吃飯了嗎?.
阿C就說 我剛剛結婚 要陪我的老婆.
沒空去吃飯了.
結果呢 這三個人都沒有辦法.
去享受上帝的筵席 上帝的福分.
這個比喻就告訴我們.
今天攔了我們去親近神.
去事奉神的東西.
未必一定是不好的東西.
相反 可能就是上帝給了我們一些恩典的事情.
所以人生的遺憾就是.
神給我們有不同的恩典 不同的恩賜.
但是有時候我們因為太喜歡這些恩典.
太喜歡這些恩賜的緣故.
以致攔了我們去跟神有一個更深的溝通.
跟神有一個更深的事奉.
耶穌藉著這一段話.

$^{281}$給我們每一個人去反省.
在我們生命當中 最看重的是什麼呢?.
主耶穌不是真的教導我們.
不疼愛我們自己的爸爸媽媽.
不喜歡我們的子女.
不喜歡我們的兄弟姐妹.
是真正的 要人將主權交給耶穌.
祂這樣說.
你們無論什麼人.
若不撇下一切所有的.
就不能作我的門徒.
若不撇下一切所有的.
就不能作我的門徒.
祂就要看我們.
看我們那個生命所看重的是什麼.
第二方面.
作主的門徒 我們會得到這個賞賜的.
第42節.
「無論誰 只因門徒的名.
就算把一杯涼水給這些小子 中的一個喝.
我實在告訴你們 他們一定得到賞賜」.
主耶穌教導完那幫門徒.
要將生命主權交給上帝.
然後就答應他們.
告訴他們 凡跟隨祂的人.
和其他人分享上帝恩典的人.
一定得到賞賜.
上帝很熱鬧.
我們去幫助周圍一些軟弱的人.
祂都很樂意.
很樂意我們將從祂而來的恩賜,恩典.
去跟一些弱小的人,貧窮的人去分享.
當我們這樣不斷去學習這個分享的路的時候.
不單止是得到上帝的賞賜.
最大的就是我們的生命可以脫離那份自私.
貪婪,小器,習道這些敗壞的性情.
可以慢慢的脫離開它.
以致我們生命更加活出耶穌基督的美善.
上星期一 我和一位姐妹去醫院探望她的叔叔.
她的叔叔無親無故在香港.

$^{321}$也要靠綜援來過日子.
他躺在床上 都很淒涼.
插滿了喉管.
他很口渴 因為不能喝水.
我看姐妹用一些很粗的棉花.
蘸一些水 讓他的唇邊變硬一點.
讓他喝一下水 濕一下他的嘴唇.
我就看到叔叔.
他閉上眼睛 一直都閉上眼睛.
他輕輕點一點頭.
好像感到很滿足.
一個粗一點的棉花棒.
一點點的水 蘸在嘴唇裡面.
他已經感到很滿足了.
我走了之後 第二天姐妹告訴我.
她叔叔跟她說了一聲「I love you」.
就這樣在病床上 蘸一點點這個水.
對她來說 真的是很感激.
當姐妹聽完之後.
她說她自己心裡面的眼淚就一直流.
就好像將她的生命清洗.
也都給她生命一份.
真是一個不懂得怎樣形容那份的感動.
這個就是上帝給她那份的禮物.
她說那份的甘甜 那份的甘甜.
有一天 門徒問耶穌.
「我怎樣才可以去服侍這些小子呢?」.
「我怎樣才可以服侍到你呢?」.
她說 給一幫門徒說.
「我們什麼時候見你餓了 或者渴了呢?」.
「我們什麼時候見到主你赤身露體」.
「什麼時候見到你病」.
「什麼時候見到你在監獄裡面」.
「什麼時候知道我們要去侍候你呢?」.
主耶穌就說了 這個王就說了.
「這些的事 你見人餓 給人東西吃」.
「見到人口渴 你給人水喝」.
「你見到人赤身露體 給人衣服穿」.
「你見到人有需要 你就去服侍」.
「你見到人在監獄裡面 你就去關懷」.

$^{361}$「只要坐在弟兄當中 其中一個最小的身上」.
「就是坐在我身上」.
「坐在一個小子的身上」.
「就是坐在耶穌基督的身上」.
另外 在馬太方第25章40節那裡.
它就說「你坐在弟兄當中最小的人的身上」.
「就是坐在耶穌的身上」.
這個小子 說最小的人 最小的弟兄.
就是說一些地位很微小 很卑微 貧窮的人.
原來 怎麼說我們服侍到耶穌呢?.
就是我們以愛心去照顧 去看顧周圍的鄰舍.
用愛心來服侍他們 就等於服侍了耶穌.
原則上 人人都能夠做到的 是不是?.
送上一杯涼水 可能那杯涼水對你來說很珍貴的.
在這個沙漠地帶 或者有限制的供水底下.
但是你送上去 你都有損失 你可能變成沒有的.
但是通常的情況底下 是一件不難做的事.
不難做的事 只是舉手之勞.
但原來這樣做的時候 給人很大的鼓勵.
很大的支持 耶穌說.
當我們這樣 在其它一些有需要的群體.
有需要的鄰舍底下 我們這樣做的時候.
就是服侍了主耶穌.
高銘獻牧師在一篇靈修分享中這樣說.
在牧師的時候 每一個人都要接受這個審判.
其中一幫人的審判是怎樣的呢?.
就是他沒有服侍到鄰舍 沒有服侍到弱小的弟兄.
所以就不符合耶穌基督的標準.
耶穌基督的標準就是要做在弱小的身上告.
做在弱小的身上告 如果沒有做的時候.
他們就會受到刑罰.
並且這個刑罰是直至永恆的.
高牧師給我們一個提醒.
就是有時候我們在做教會的侍工裡面.
我們重視了這件事究竟是否成功呢?.
或者你看一下這件事有多少個數字呢?.
我們很留心 很重視了它.
如果是這樣的時候 我們會失去了事奉的焦點.
他說 人不應該只是為了侍工而存活.
特別是教會裡面所有的侍工.

$^{401}$都需要為服侍的人而存在.
所以教會裡面所有的侍工.
都應該以人為中心 求主幫助我們.
最後我們就用一篇的土文.
叫做《只能活一次的禱告》.
不知道作者是誰 但我覺得裡面的內容很好.
我們就一起讀這一段的土文.
成為我們的祈禱.
1,2,3.
我來世上只能活一次.
若有任何能做的好事.
或任何能向人表達的善行.
讓我立誓去做.
讓我不辭焉 不卸責.
因為我將不會再從這條路經過.
\newpage



\section{加拉太書 6:9-10}
\label{sec:ZmEDOfaqivc}
\textbf{2023-10-22 加恩樓《及時行善》加拉太書 6\_9-10 講員:何智雄牧師}
\newline
\newline
連結: \href{https://youtube.com/watch?v=ZmEDOfaqivc}{\texttt{https://youtube.com/watch?v=ZmEDOfaqivc}} ~~~~ 語音日期: 2023-10-24
\newline
\newline
\hyperref[sec:qZ9s1IR331E]{\small{< < < PREV SERMON < < <}}
~
\hyperref[sec:index_chronic]{\small{[返順時目]}}
~
\hyperref[sec:index_scriptual]{\small{[返順卷目]}}
~
\hyperref[sec:LM01sAVsGrw]{\small{> > > NEXT SERMON > > >}}
\newline
\newline
加拉太書 6:9-10
\newline
\begin{longtable}{cl}
\hline
\hline
章節 & 經文 (和合本修訂版)\\
\hline
6:9 & \begin{tabularx}{0.7\textwidth}{X} 我們行善不可喪志，因為若不灰心，到了適當的時候就有收成。 \end{tabularx} \\ \\ \relax
6:10 & \begin{tabularx}{0.7\textwidth}{X} 所以，一有機會就要向眾人行善，向信徒一家的人更要這樣。 \end{tabularx} \\ \\
[1ex]
\hline
\hline
\end{longtable}
$^{1}$很開心能夠再一次回到.
嘉欣家這個地方跟大家一起見法上帝.
我以前也是這樣孫的童宮.
不過N年那麼長.
主要是讓我可以去孫一堂.
去牧養弟兄姊妹.
弟兄姊妹.
平安.
今天說平安其實很難.
是不是.
每一天打開電視.
打開你的新聞.
打開你的手機.
太多的事.
是太多.
你覺得.
嘩.
怎麼辦.
無能為力這個感覺.
我不知道大家有沒有.
我就有.
大家不要以為牧師很厲害.
嘩 行了.
牧師你靠著上帝祈禱.
可以搞定.
這個牧師不行的.
其他的牧師我不知道.
這麼近.
這麼深.
怎麼辦.
今天我盼望用加拿大書跟大家一起分享.
透過加拿大書跟大家講一個訊息.
加拿大書是保羅他去了第一次的宣教旅程.
加拿大在哪裡呢.
就是今天的土耳其.
不知道大家有沒有去過這個地方.
不要緊 大家看看圈住的地方就是.
接下來的選舉都會去土耳其.
就是這些地方.
當保羅去了這些地方.

$^{41}$建立了教會之後.
他就離開.
之後他收到消息.
裡面的弟兄姊妹.
面對很大的困難.
很大的挑戰.
為什麼.
是因為有很多不同的人.
包括一些說自己信教.
信猶太教.
信舊約聖經的人.
去到教會當中跟我們說.
信耶穌.
不只是信而已.
大家知不知道什麼叫信耶穌.
你是不是信耶穌的.
什麼叫信耶穌.
裡面有些人說.
你要受割禮的.
你做得好不好.
不要做好事.
有沒有犯罪.
怎麼行.
這樣都行嗎.
不行.
但耶穌的福音是什麼.
不因為我們做得多好.
也不因為我們做得多差.
也不因為我們怎樣.
只要我們相信耶穌基督.
道成肉身來到這個世界上.
為我們死.
我們相信祂流出的寶血.
我們生命就能夠得到潔淨.
因信聖義.
就是我們信仰的核心.
弟兄姊妹.
你是不是因信聖義.
但當時的弟兄姊妹說.
不是啊.

$^{81}$他告訴別人.
你不行.
這樣還行嗎.
你有沒有割禮.
有很多困難.
很多疑問在裡面.
於是保羅就寫了這封信.
所以整本加勒太書讀完之後.
你猜一下.
保羅想你知道什麼.
讀完整本書.
可能你還是.
這樣.
不要緊.
你記住一句金句.
你就知道加勒太書說什麼.
就是加勒太書第五章.
第25節.
大家看見嗎.
一起讀一次.
預備起.
「我們既靠聖靈得生.
就當靠聖靈行事」.
我們既靠聖靈得生.
就當靠聖靈行事.
如果我們每一個人.
都是因為聖靈的感動.
聖靈的帶領.
接受耶穌基督.
為我們個人的救主.
我們得到新的生命.
我們就要什麼.
要靠聖靈行事.
這是理所當然的.
所以整本加勒太書.
可以分成三部分.
第一部分是講到.
保羅自己的身邊.
因為有太多人.
太多人過來跟人說.

$^{121}$信耶穌怎樣怎樣.
保羅就說.
我的信息.
是從哪裡來的.
不是從人來的.
是從誰來的.
耶穌基督.
親自跟我說的.
所以這個身份是很重要.
第一部分.
保羅就說.
我的身份是這樣.
所以我說的.
不是我自己想說.
我是聽得清楚.
是主耶穌要我說的.
第二部分.
就跟一群加勒太人說.
你們所面對的困難挑戰.
知道明白.
有人跟你說要守這個.
有一個律例等等.
都知道.
但這不是信耶穌.
信耶穌是什麼?.
恩,信,清,義.
這是第二部分.
第三部分就是.
如果是這樣的時候.
我們就有自由了.
大家有沒有感覺到.
信耶穌是有自由的?.
不是說你做基督徒.
你一樣不行.
這樣不行.
什麼都不行.
很慘的.
全部被綁住.
但在耶穌基督裡面是什麼?.
是有自由的.

$^{161}$但這個自由要小心.
不是亂用的.
某些人說自由是什麼?.
我想做什麼就做什麼.
我要怎樣做就怎樣做.
我覺得應該怎樣就怎樣.
這個叫自由.
這個是不是叫自由?.
這個不是叫自由.
大家知道嗎?.
這個叫做我想怎樣就怎樣.
做基督徒不是這樣.
在耶穌基督裡面有自由.
它的表現是怎樣呢?.
《加勒太書》裡面有說.
加勒太書說到.
當我們這樣的時候.
我們不單只顧著自己的需要.
我們要看別人的需要.
都是很重要的.
當我們要和別人一起.
分擔一些重擔.
自己也需要承擔自己的重擔.
要在供應給僑童工一些生活上的需要.
不單止這樣.
在第九節和第十節裡面.
今天我想說的.
當我們得到這份自由.
在我們內在內注的時候.
我們要行善.
行善是什麼意思?.
不是滿東華那些捐錢.
不是.
那些也可以說是.
大聖經上說的.
行善是什麼?.
為其他人著想.
做一些其他人所需要的事情.
但是.
現在這個世代.

$^{201}$做好事是怎樣的?.
有沒有好報?.
最近大家聽到經常有騙案.
什麼什麼什麼.
很多人都是.
好心未必有好報.
第五節是這樣的.
你的經歷是不是這樣?.
我的經歷就是這樣.
但是.
聖經提醒我們.
這個世代是這樣的.
現在是黑暗的世代.
這個世代當中有很多很難的事情.
有很多很黑的事情.
做好事未必有好結果.
但是聖經就說.
那怎麼辦?.
就不要做了.
現在有個詞叫躺平.
有沒有聽過?.
做都是這樣.
不如什麼?.
不做了.
躺平在這裡.
睡覺不好.
是不是?.
不是.
聖經不是這樣的教導.
聖經教導的是什麼?.
繼續做好事.
繼續要行善.
因為.
雖然這個世代是這樣的.
但是不要灰心.
不要喪志.
為什麼?.
因為到了時間.
你們就會修成.
這一點是重要的.

$^{241}$所以在《迦太書》裡面.
保羅用了一個比喻.
一個眼光給我們看.
行善是什麼?.
另外在《迦太書》第五章第十六節.
跟我們說了一個很重要的真理.
當我們信了耶穌之後.
在我們裡面.
其實有兩個律.
大家知不知道?.
一個是聖靈.
聖靈在我們當中.
跟我們說做什麼.
第二個就是自己的自律.
是我想做什麼.
如果我們不按照聖靈的帶動.
我們就會.
按照自己的行善.
問題就出現了.
這就是《迦太書》教導我們.
一個非常重要的信息.
所以我們需要什麼?.
常常都要聽.
聖靈跟你說了什麼.
有沒有跟你說.
什麼時候跟你說.
他對你怎麼說.
你有沒有信心.
有沒有眼光.
按照著聖靈的帶動.
去做一些事情.
這個很重要.
怎麼做呢?.
做成什麼樣呢?.
我有一首詩歌送給大家.
我請我們聽聽.
我覺得這個歌很好聽.
很優美.
大家聽下去就知道了.
《迦太書》 詞曲:李宗盛.

$^{281}$是我送給我們的知識.
在怨恨之中是有悲悲哀.
在感傷之中傳送的故事.
在懷疑之中獻上的信心.
服從我們要讓大道發光.
能願意留你教會我們所有.
能求你要來達到的大任.
至深愛人但願意付上光.
是我送給我們的知識.
很幽悄的環境.
祂說什麼呢?.
祂說在憎恨的地方播下愛.
在傷痕當中播下寬恕.
在懷疑當中播下信心.
在失望當中播下相信.
聖靈感動我們的時候.
要我們做什麼.
就是一個和平之節.
在《迦太書》裡面再給我們一個訊息.
一個圖畫.
應該要做什麼.
夠了,謝謝.
因為太長了.
如果大家喜歡可以上youtube聽就行了.
很好聽.
在《迦太書》裡面給我們一個圖畫.
讓我們清楚知道.
我們行善不要灰心.
不要喪志.
為什麼?.
不知道大家有沒有種菜.
種瓜種豆.
有沒有?.
我知道大家家裡有些人會種這些東西.
當你種的時候.
你會知道種瓜做什麼?.
種豆呢?.
非常好.
但問題是否馬上就有?.
是否馬上就有?.

$^{321}$當然不是.
你也知道.
要等.
等什麼?.
等它慢慢成長.
等什麼?.
等時機.
不是要有就有.
想就有.
不是這樣的.
在我們的生命當中都是一樣.
所以保羅在《迦太書》的六章第九節.
告訴我們.
我們行善不可喪志.
若不灰心.
是什麼?.
到了時候.
就要修成.
原來有時候我們灰心喪志是因為什麼?.
因為我們沒有耐性.
我們覺得應該有.
對了.
但上帝說什麼?.
還不是時候.
還不是時候應該怎樣?.
不要灰心.
還不是時候怎樣?.
不要喪志.
為什麼?.
來到的時候就行了.
你的生命是否有這些經驗?.
有沒有?.
在我的錢包裡.
經常都有一張票.
讓我看看.
就是這張.
這張什麼票呢?.
我兒子小時候.
我帶他去泰國曼谷.
我當了宣教士.

$^{361}$去了三年.
去到那邊人生路不熟.
只能坐車.
但這裡是去地方的.
坐車是一件不容易的事.
大家有沒有去過泰國曼谷?.
我想有些.
有沒有坐過巴士在曼谷?.
我想很少.
但作為宣教士.
你每次都要坐的士.
還有的士貴.
有一天.
我一樣這樣.
我打算去某些地方.
在巴士站等巴士.
那邊的巴士不是很快.
等一等.
看一看.
傻乎乎的.
突然有一個人開電單車.
來到我面前.
然後他給我.
我看一看.
不知道是什麼.
因為是泰文.
不懂.
還未懂得看.
之後再看一看.
咦.
是一個月票.
看看日期.
是不是今天都可以用.
咦 是的.
於是我怎樣.
車來的時候.
我給他看.
試試是否真的.
你知道有很多假的東西.
咦 可以.

$^{401}$這張票.
跟我說什麼.
這張票.
就像上帝和我一樣.
我在這裡.
去泰國一點都不容易.
特別是那時候.
我介紹一下自己.
我有三個兒子.
三個都是兒子.
大兒子和一對媽兒子.
那時候他們是幾歲.
我大兒子差不多六歲.
我和太太膽子粗粗.
這樣的東西.
都肯去.
就這樣.
去到之後.
有很多困難挑戰.
大大小小.
逛街吃飯都是一個挑戰.
上帝和我說什麼.
我在這裡.
你所需要的我給你.
不用錢.
我從來都沒有想過.
我竟然可以在曼谷.
不用錢.
帶一程車.
所以我學會了之後.
我就開始學一下.
行線.
做一些好事給別人.
有一次我去曼谷.
之後再去曼谷.
大家有沒有去曼谷搭BTS.
有沒有去過.
BTS在上面.
不是下面.
有些叫日票.

$^{441}$日票是什麼意思.
你買完之後可以無限搭.
之後我買了一張票.
我用了.
差不多黃昏的時候.
就想.
還可以用.
我不用用了.
那張票怎麼辦.
我就想起這裡.
於是我就膽粗粗.
就問人.
Jeffrey.
你知道什麼叫Jeffrey嗎.
不用錢免費的.
你問他有什麼反應.
你自己做一些善事.
別人一定要很回報.
拍手說.
不是的.
之後我再給.
Jeffrey.
接著就是那個人.
男人.
背後那個拿了之後.
進去.
對於我來說.
我覺得很感恩很開心.
為什麼.
因為聖靈跟我說.
我不知道這個可以幫到多少.
但是我在做什麼.
行善.
做一些事情.
令人開心.
讓人感受到一點點的盼望.
讓人感受到原來這個世界上.
仍然有人會這樣做.
為什麼.
因為我受過.

$^{481}$我自己就懂得怎樣去付出.
突然間.
我要說的就是.
時機.
這裡特別提到.
第九章第六章第九節.
到了時候就要修成.
第十節說到.
如果時機到的時候.
你就要去做.
在信主的家當中.
更加要這樣做.
什麼叫時機.
對於我來說.
時機就是適當的時間.
適當的地點.
適當的人做適當的事情.
說適當的話.
這個就是時機.
我再說一次.
是適當的時間.
適當的地點.
適當的人做適當的事.
說適當的話.
這個就是時機.
難不難.
難 當然難.
但行不行.
行.
我住在元朗.
剛才說了.
我和太太.
有時候是單人無事.
平日起床.
大家會做什麼.
起床之後.
非常好.
我都擦.
之後我就會去.
喝茶.

$^{521}$大家去過元朗喝茶.
有很多酒樓.
之後我就會去.
我會選一間.
不如去另一間.
我們就去另一間.
去哪一間.
大家都可能有這些困難.
隨便吧.
於是我就去了那間.
去了那間上面.
很多位.
大家選哪個位.
無所謂.
都沒有了.
於是我就選了角落.
就在那裡.
我們就過去坐.
大家都知道喝茶是怎樣.
喝茶很窄.
桌子和桌子之間很近.
經常都會聽到旁邊的人說什麼.
你不想聽.
但你沒有辦法.
我們坐下來.
我們點了東西吃.
然後開始喝茶.
突然間我們聽到旁邊有一個人.
媽媽和小朋友.
為什麼這樣.
你吃東西吧.
我叫了所有.
你還.
你知道我很辛苦嗎.
叫了你那麼多次.
你都不吃.
你看你的手.
人人見到你的時候都說你很瘦.
你吃東西吧.
這些對白.

$^{561}$有沒有聽過.
有沒有.
有沒有發生在你家.
你都不敢說.
沒有感恩.
但當你聽到的時候.
你會有什麼反應.
我太太.
師母聽完之後.
她就過去.
跟媽媽聊天.
跟媽媽說.
你的小朋友可能會有.
吞煙困難.
如果有的話.
你可以跟老師聊天.
或者他可以.
叫你去看一些.
聊完就回來.
當時我就在.
祈禱.
我可以做些什麼.
我有一點點感動.
最後都.
沒有了.
師母回來.
我打算結帳離開.
準備離開的時候突然間有感動.
我走近.
坐在小朋友旁邊.
我跟他說.
你很特別.
他看著我.
你真的很特別.
你把你的食物.
放在你的口裡.
咀嚼咀嚼咀嚼很久.
你知道是很好的嗎.
慢慢.
到你的胃的時候.

$^{601}$你很容易消化.
你真的很厲害.
小朋友有什麼反應.
哭了出來.
之後媽媽就很緊張.
我師母就跟他說.
不要緊,讓他哭吧.
媽媽就立刻拿了紙巾出來告訴我.
不要緊,男孩可以哭.
我跟小朋友說.
不要緊,你哭就哭吧.
之後我就跟他說.
現在媽媽.
幫你點了這麼多食物.
你選一種就好.
不要所有食物.
只選一種,其他給媽媽.
你知道他選了什麼嗎.
蝦餃.
他還要拿起筷子.
夾蝦餃在自己的碗裡.
我看了之後.
我再問.
因為我知道.
他的年紀.
我看了之後就知道他是小學生.
之後就跟媽媽.
希望多聊聊天.
我問他.
你讀哪間學校的.
你為什麼會知道.
這麼厲害.
我聽到之後簡直嚇一跳.
上帝.
是不是這麼巧.
世界上.
這麼巧的時候.
他竟然說是徐澤林小學.
我問他,你知道我是誰嗎.
我是何牧師.

$^{641}$他當然不認得.
我懶得嚇他.
我跟他說不要緊.
我剛剛去了香港.
我剛剛回來.
我跟他說你叫什麼名字.
你以後在學校見到我.
你要叫我.
因為我不認得你,我不記得你叫什麼名字.
你告訴我你叫什麼名字.
我想說的是什麼.
你知道我想說什麼嗎.
適當的時間.
適當的地點.
適當的人.
做適當的事.
說適當的話.
你知道我跟他說什麼嗎.
我說.
你知不知道你很帥.
我真的很帥.
我和你自拍一張.
我和你自拍了一張.
兄弟姐妹.
我想想.
如果這番話.
如果不是我說.
當然可以.
但如果是我說.
就和你說不一樣.
明白這個道理嗎.
什麼叫時機.
兄弟姐妹.
這個世代當中.
很黑暗.
很多困難,很多挑戰.
但作為基督徒.
我們不是躺平.
我們不只是祈禱.
我們仍然有事可以做.

$^{681}$但我們需要捉緊時機.
如果時機來.
是你的時候.
要捉緊.
大家知不知道什麼叫.
事半功倍.
但很多時候.
我們是事倍功半.
大家明白這個意思嗎.
為什麼.
因為時機.
後面的燈關了.
我想大家.
如果突然之間.
大家去到一個黑暗的地方.
這裡行不行.
不行,還有些光.
大家閉上眼.
請大家全部都閉上眼.
合上眼睛.
突然之間.
你進入黑暗的裡面.
你第一個反應.
你要做什麼.
你第一個反應.
你要做什麼.
突然之間.
我們進入一個黑暗的環境當中.
你第一個反應.
你想起的.
你會做什麼.
請大家打開眼睛.
我會拿起我的手提電話.
打開燈.
為什麼.
全黑的時候.
你需要什麼.
就是光.
當然現在的光和那些比.
當然.

$^{721}$沒得比.
但試想一下.
全黑的環境.
這個光.
是什麼.
這個光很重要.
是不是.
突然之間.
我想說的是什麼.
突然之間,你知道你是光嗎.
耶穌說.
我們是世上的光.
是用來照亮其他人的.
在黑暗的世代當中.
我們沒錯.
我們當中有很多困難.
有很多挑戰.
請你拿起手提電話.
有沒有.
有的,請你拿起手提電話.
打開行不行.
行不行.
難不難.
我們需要的就是這樣.
如果你和我都願意行善.
你和我也願意不放棄.
你和我也不願意灰心.
你和我也願意有需要的時候.
拿起手提電話.
打開燈的時候.
這個世界就可以不一樣.
幾個月前我去了以色列.
很特別.
我不知道為什麼我會去以色列.
之前我也不太知道.
但之後.
為什麼我會去以色列.
其實主要都是因為泰國人.
那時候我聽到在泰國.
有幾萬個泰國人在那裡打工.

$^{761}$沒有人理會他們.
沒有人關心他們.
於是我拍了心口.
沒有人去,我就去.
我會不會說泰文呢.
十幾年沒有說了.
不過還可以說一點.
(泰文).
如果我去到之後.
起碼和他們說(泰文).
他們也會覺得近一點.
於是我去了以色列.
兩個星期前.
當我聽到新聞.
看到新聞.
哈馬斯突然間.
五千個飛彈飛進去.
然後我就看看在哪裡發生.
大家知道在哪個位置嗎.
什麼位置.
大家認熟吧.
那裡叫加沙.
加沙走廊在哪裡.
在圖上.
我很記得那天.
是人生當中第一次.
開左tie 車.
很害怕.
平時我們是右tie .
那次是左tie .
嚇得我七彩.
開到哪裡呢.
就是加沙走廊.
我很記得那天.
開車去到這個地方.
我的右手邊.
就是加沙.
眼都看到.
接著再想起.
今天想起的時候.

$^{801}$在那段時間.
我認識了一些猶太人.
認識了一些以色列人.
在那邊.
我好傷痛.
我就會問.
他們怎樣.
很特別.
在那段時間.
我去了以色列.
在那段時間.
我認識了一些猶太人.
在今天.
我們知道這件事.
所以對於我來說.
我是好痛.
你知不知道.
在那次的襲擊當中.
死了多少個泰國人.
28個.
8個受傷.
還有16個失蹤.
我聽到這些消息的時候.
我心裡是好痛.
因為我.
見過這些人.
我可以做些什麼.
我好相信.
上帝告訴我.
既然我去到.
我知道這些消息的時候.
我就為他們祈禱.
因為大家都有手機.
大家知道手機有WhatsApp.
有Line.
大家知道什麼叫Line.
泰國人是不用WhatsApp.
於是就拉在一起.
在不同的地方.
我認識了很多泰國人.

$^{841}$就知道他們的第一手資料.
他們是否安全.
我懂得怎樣為他們祈禱.
我想說的那些是什麼.
我不是想播放.
我想說的是.
各位.
弟兄姊妹.
拿起你的手機.
打開你的燈.
照亮一下.
你身邊的地方.
有很多人很需要的.
可能就在你旁邊.
可能今晚.
你就會看到他.
在抓緊機會.
祈禱.
我做的不一定是你做的.
我去過以色列.
我去過以色列才做的.
我見過很多徐南雄的家長.
不等同你也會做.
不是,是每個人都有.
上帝在每個人的心目當中.
都安排了你的時機.
都有的.
問題在哪裡.
問題是你知不知道.
你沒有抓緊這個時機.
你沒有去做.
其實也不難的.
今天我來到這裡.
再看這個地方.
我記得有一次.
原本有很多問題.
很多地方.
我心想怎麼辦.
哇.
接著怎樣.

$^{881}$時機到的時候會怎樣.
適當的時間.
適當的地點.
適當的人.
做適當的事.
說適當的話.
今天就行.
原來.
原來大家都經歷過.
是剛剛經歷過這件事.
不是新的事.
不只是我才能經歷到.
是大家都可以.
問題就是.
問題是什麼.
你有沒有把握.
你有沒有上上祈禱.
你的行事為人是按照聖靈的引導.
我覺得.
我想怎樣.
我覺得應該.
是不是這樣.
我有沒有想過.
我應該去那間酒樓喝茶.
我有沒有想過.
在那間酒樓我會見到這個小朋友.
完全不是我的安排下.
沒有這麼巧合的.
我不相信.
世界上有這麼巧合的東西.
在我裡面是什麼.
是上帝的安排.
上帝讓我捉緊這個時機.
我做的事不是很多.
但是這個家庭.
有很大的變化.
我最近11號.
我們孫一棠.
我們去了浸會員.
有一個日營.

$^{921}$我們邀請了徐三小學的老師.
校長也來.
副校也來.
老師也來.
很開心.
不單這樣.
家長也來.
學生也來.
你知道那天見到誰嗎.
何牧師.
他媽媽跟我說.
牧師我真的很想講見證.
他講了很多.
我好心記得牧師母跟我說的話.
牧師母跟他說了什麼話.
他說小朋友成長不要太緊張.
要求健康就行了.
是健康就行了.
不要說他這麼瘦.
你為什麼這麼瘦.
人家都說你瘦.
健康就行了.
這位家長好心記住師母這句話.
她改變了.
她跟她的小朋友的關係都改變了.
一家人的關係都改變了.
難不難.
改變大不大.
誰可以做得到.
你跟我都可以.
適當的時間.
適當的地點.
適當的人.
適當的做事.
適當的話.
這些就行了.
鄧志武我很想你都能夠見到.
為什麼.
事件不是最重要.
人人都有困難.

$^{961}$人都有挑戰.
你有你的我有我的.
我的困難不比你的小.
你的困難不比我的大.
重要是什麼.
在事件當中.
你能夠經歷到你的上帝.
是一個怎樣的上帝.
是不離不棄的上帝.
是He will make a way的上帝.
在困難當中.
一路堅持一路等候.
時間到的時候.
路就開了.
你和我就能夠一起度過.
可以嗎.
我們一起建上祈禱.
在準備祈禱的時候.
我再一次邀請大家.
如果你願意去拿起你的手機.
打開你的燈.
請你現在就拿起你的手機.
亮起你的燈.
這個世代很需要我們的燈.
這個世代很需要的是光.
是基督的光.
我們一起祈禱.
多謝你.
因為你從來都沒有離開過我們.
你知道我們正在面對什麼.
挑戰有多大.
我們心情有多多.
我們要求的是什麼.
不只是解決這個問題.
求的是你在我們裡面.
求的是讓我們真的很深的經歷到.
我們會因為這些事.
危機永遠可以變成一個契機.
契機是我們可以更深的經歷到.
你的同在.

$^{1001}$你的保守.
和你的帶領.
這個是最重要的.
面前我們仍然有很多困難.
很多的挑戰.
一件又一件.
一堆又一堆.
我們需要的不是去解決所有的問題.
我們需要的是耶穌基督.
我們一次有個經驗的時候.
我們就經驗在我們心裡面.
常常帶著一個感恩的心.
去活每一天.
這就是基督徒.
看見我們身邊有人需要的時候.
不要灰心.
不要喪志.
不要覺得自己沒用.
伸出我們的雙手.
一個電話.
一個問候.
可能已經改變了他的生命.
可能改變了他的家庭.
可能改變了他的關係.
主啊!幫助我們.
我們不需要做很多事.
我們需要做的是對的事.
是聖靈感動帶領的事.
多謝你!讚美你!.
我們同心合意在你面前一起禱告.
仰望奉耶穌基督的名字祈求.
阿們!.
\newpage



\section{約拿書 1:1-17-4:1-3}
\label{sec:LM01sAVsGrw}
\textbf{2023-10-22 午堂《經歷有恩典憐憫的神》約拿書 1\_1-17,4\_1-3 講員:陳燕雯姑娘}
\newline
\newline
連結: \href{https://youtube.com/watch?v=LM01sAVsGrw}{\texttt{https://youtube.com/watch?v=LM01sAVsGrw}} ~~~~ 語音日期: 2023-10-25
\newline
\newline
\hyperref[sec:ZmEDOfaqivc]{\small{< < < PREV SERMON < < <}}
~
\hyperref[sec:index_chronic]{\small{[返順時目]}}
~
\hyperref[sec:index_scriptual]{\small{[返順卷目]}}
~
\hyperref[sec:xr_m3_qAqFs]{\small{> > > NEXT SERMON > > >}}
\newline
\newline
$^{1}$大家早晨 各位弟兄姊妹平安.
很開心可以來到錦繡堂跟大家宣講神話語.
其實在五月的時候.
我和Z7的歐洲及亞太區執行董事Kirk都來過.
在當中分享Z7的事公.
不知道今天有多少個在當天在場呢.
都很多 其實你們對Z7都不陌生.
我是Z7的事公拓展主任.
我在2017年的時候回歸香港.
然後入讀神學進修.
2023年在Z7香港服事.
在我未回來香港之前.
我曾經去過阿爾巴尼亞 奧比利亞服事了三年.
在一間宣教士的子女學校教宣教士子女.
我也認識在當地服事的宣教士.
聽到那些宣教士分享.
他們去到的頭一兩年主要都是要學習語言和文化.
甚至有些警察會要求這些宣教士.
要有一段長的時間和當地人一起住.
目的就是希望他們可以體驗他們的生活和文化.
使這些宣教士不單止在頭腦上知道.
也真是親身的體會.
為什麼我這樣分享我小小的背景和這些事呢.
一會兒我就會揭曉.
讓大家一方面知道我小小.
另一方面留心聽為什麼我會說這件事.
我們就一起來祈禱.
親愛的天父 多謝你的恩典.
你帶領我們朝聚.
你也透過宣教士去到不同地方傳講福音.
以致到就算在香港我們也得著蒙受.
以前西教士傳福音來到香港的福氣.
以致我們也能得著這個救恩.
神愛的求你也給我們有一個淌開的心.
在約拿書裡這件事給我們有學習.
也給我們能夠領受.
主耶穌求你聖靈幫助我們.
求主也幫助我們說話 聽話.
讓說話能夠清晰 聽話能夠聽得明白.
我們同心的禱告.

$^{41}$盼望我們一同走這條天路.
祈禱交託 奉主耶穌基督的名祈求 阿們.
今天我的講題就是.
經歷有恩典憐憫的神.
剛才主席幫我讀了約拿書的第一章和四章的頭三節.
但其實我們今天會整本約拿書去看.
不單看一章這麼簡單 而是看完四章.
所以我會簡單地講一講.
約拿書裡面說了些什麼.
或者大家都很熟悉的.
就在很久很久以前.
大概是主前的八世紀.
在一個很遠的地方.
距離以色列大概八百至一千公里的亞述.
的一個大城 尼尼微.
而這個國家的人就.
行了很多勁惡事.
做了很多 譬如拜偶像 邪淫 這些事.
到一個地步 耶和華要親自介入.
他差派了若拿來到當中.
和他們宣告一個審判的信息.
大家可以看看這個地圖.
一方面我們看到在地圖的右手邊.
即是東面 看到沿海 地中海就有以色列.
當中有猶大 耶路撒冷那個地區.
而神其實呼召若拿去到.
以色列的東北亞述國利尼尼城那裡.
而若拿他就去了約柏近海岸那裡.
打算搭船去西邊.
就去了這個塔斯.
塔斯其實就是西班牙的地方.
在若拿還沒有到達的時候.
神就做了一些事.
若拿以為這樣就可以解決了 逃避了.
神叫他去東 他就去西.
他就想好了.
在他上船沒多久 神就掘起大風浪.
以致到船差不多要爛 要沉.
船上的水手和船上的人都不知道怎麼做.
他們就大家都求神拜佛 求他們的神.

$^{81}$希望可以平安.
但若拿在這個時候反而就睡在船底.
船長就問他 你快點來求你的神吧.
然後若拿就跟他們說.
船員和水手也知道這件事是因若拿而起.
之後那些水手把若拿丟下海的時候.
風浪就停止了.
若拿會不會因為這樣而葬身在大海呢.
不是 耶和華有恩典有憐憫.
耶和華派一條大魚吞了若拿.
若拿在雨的覆裡三日三夜.
在這段時間 若拿禱告向神祈求.
若拿憐憫他 恩待他.
就在一個乾地放他出去.
撿回一條小命 若拿應該感恩了.
神第二次跟他說.
我要你去尼尼尾城宣告.
若拿這次就去.
宣告說四十日後這個城就要毀滅了.
那些人聽到他們的反應是怎樣呢.
他們相信信服.
他們立刻從大到小都悔改認罪 哀傷.
披麻蒙灰.
其實若拿早就知道.
他一方面知道.
這個亞述帝國去攻擊侵略以色列.
所以他很不喜歡去到他們當中.
另一方面若拿又知道.
這個耶和華是有恩典有憐憫的神.
不輕易發怒 而且有豐盛的慈愛.
他知道當人願意悔改回轉的時候.
耶和華會不幹栽和比尼尼尾城.
所以耶和華叫他這樣做.
他不想有這樣的結果.
他就唯有不聽逃避耶和華.
但是那些人真的認罪悔改.
若拿就因為這樣而不高興.
甚至發怒.
神真的有恩典有憐憫.
不輕易發怒.

$^{121}$他內心的來教導若拿.
讓若拿有反思.
約拿書這個故事.
大家可能聽了很多次.
其實裡面有很多教導教訓可以分享.
我選了三樣跟大家分享.
第一個 全港福音是基於神的恩典和憐憫.
如果我們認定全港福音是基於神的恩典和憐憫.
我們就不需要怕不夠口才.
不懂說話 或擔心別人會拒絕我們.
我們看看若拿說了什麼.
其實若拿所說的很簡單.
他和那些水手和在船上的人.
他們不認識耶和華.
若拿就跟他說.
他是希伯來人.
我是敬畏耶和華的.
天上的神 祂創造了滄海和陸地.
祂就是把耶和華的本性.
祂的屬性 祂是創造主.
我們來跟這班水手說.
接著我們看看若拿去到尼里味城又說了什麼.
很簡單 一句話.
若拿去到尼里味城的時候.
那個城很大 要走三天這麼久.
但他真的走了一天.
就是說再過四十天你們就要滅亡了.
耶和華吩咐若拿說這個審判的信息.
若拿就照著說了一個這樣的神審判的信息.
其實不單止若拿去宣講.
我們看到保羅都很多時候去到不同的路程裡去宣講.
甚至乎在土耳其 中東的地區.
都有很多保羅的足跡.
保羅在哥林多前書裡這樣說.
他說弟兄們 我從前到你們那裡去.
並沒有用高賢大智向你們傳講神的奧秘.
因為我立定主義.
在你們中間什麼都不想知道.
只知道耶穌基督和他釘十字架的事.
我在你們那裡的時候又軟弱又懼怕.

$^{161}$而且戰戰兢兢.
所以我所說的話 說的道.
都不是用智慧的話去說服人.
而是用聖靈無能力去證明.
使你們的信不是憑著人的智慧.
而是憑著神的能力.
保羅其實是很有口才 很有智慧的.
在《使徒行傳》當中.
可能大家都會聽到保羅怎樣去答辯.
但保羅居然在哥林多前書裡.
他跟會眾說.
他不是靠他的高賢大智.
他也不是靠他的智慧的話去說服人.
他靠的是聖靈.
他靠的是耶穌基督的恩典.
他所宣講的就是耶穌基督釘十字架事跡.
這個就是神的恩典和憐憫.
傳福音是基於神的恩典和憐憫.
我們看看約拿所宣講的結局是什麼.
那群水手和船上的眾人.
他們聽到知道.
原來這是創造天地萬物的主宰.
這個就是獨一真神的時候.
他們就大大敬畏耶和華.
有些譯本將這個字譯為敬畏.
英文就是fear.
他不只是怕這麼簡單.
還有一種尊敬的意思.
敬畏耶和華.
其實敬畏這個字就跟之前約拿跟水手說.
我是希伯來人.
我是敬畏耶和華.
那個字是同一個字.
這裡就說到那群原本拜假神拜偶像的水手.
當他們認識耶和華之後.
他們就敬畏耶和華.
而且還獻祭和許願.
我們看看尼里尼城的結局又是怎樣.
當約拿宣講審判的信息的時候.
尼里尼城的人.

$^{201}$大致他們的王小至的平民百姓.
他們從大到小都穿上麻衣.
不理人畜都披上麻衣.
而且各人懇切呼求悔改離開惡行.
放棄手中的強暴.
他們回轉歸來.
神不願一人沉淪.
乃願人人悔改.
這個並不只是新約的教導.
其實神是一個宣教的神.
他的心腸一直都沒有變.
甚至在舊約的時候.
你都看到他的心腸.
神可以藉著約拿.
來叫那班水手和船員認識真神.
去向耶和華獻祭許願.
神也可以使用約拿.
去使這個滿有惡行的尼里尼城.
回轉悔改.
歸向耶和華.
傳福音是基於神的恩典和憐憫.
在中東和北非廿多個國家當中.
有超過九成的人口是穆斯林.
如果不是基於神的恩典和憐憫.
Z7不可能也不會成立.
是一件不可能的事.
一方面基於神的恩典和憐憫.
弟兄姊妹願意將身體獻上.
當作活祭.
他們甚至冒生命危險.
願意出鏡讓人知道.
公開承認他們是基督徒.
講他們的見證.
做這些基督教節目.
不單是這樣.
如果不是神的恩典和憐憫.
在一個這樣的未得之民當中.
真的不可能有人會回轉歸順.
釘十字架的耶穌.
所以我很想在此鼓勵弟兄姊妹.

$^{241}$如果我們堅信認定.
我們無論是傳福音 宣教.
甚至我們在日常生活中做見證.
我們的根基是基於耶穌基督.
祂的恩典和憐憫.
我很深信在錦宣街的弟兄姊妹.
他們無論做印尼的侍工.
或者泰國的侍工.
或者是建堂.
或者傳福音報道的侍工.
你們一定會經歷到.
是神的恩典和憐憫.
接下來你們的堂慶.
是否四十二年了.
經歷了四十二年的光陰歲月.
我很深信錦宣街也經歷到.
是神的恩典和憐憫.
所以我在此鼓勵弟兄姊妹.
不要灰心 不要放棄.
可能你為未信的家人朋友去祈禱.
希望他們信耶穌.
又或者你們猜派的宣教士.
在宣教工場裡去服侍.
好像沒什麼果效.
聽不到有什麼果子.
記住 我們傳福音.
是基於神的恩典和憐憫.
在抓緊這件事的時候.
我們就這個力量繼續去做.
繼續忠心地去走.
第二個我分享的是.
經歷有恩典憐憫的神的改變.
其實剛才也有提到.
我們比較一下約拿和水手.
和利利美的人.
約拿她自己說的.
她說我是希伯來人.
我敬畏耶和華.
這個耶和華就是天上的神.
祂創造了滄海和陸地.

$^{281}$祂知道誰是耶和華.
祂認識我.
祂甚至說自己是希伯來人.
是什麼意思呢.
其實祂告訴那些外國人.
祂是希伯來人的意思就是說.
其實你知不知道.
我們這個民族是被神揀選的.
我們是有神的揀選.
有那個的恩典的.
是一個特別的民族.
但是她的回應是什麼呢.
約拿卻起來.
要逃往他思去.
躲避耶和華.
她居然去躲避耶和華.
要逃走.
她說她自己認識耶和華.
甚至敬畏.
為什麼她會這樣做呢.
我小時候很怕我爸爸.
所以我就會很乖.
不敢頑皮.
因為怕他會罰我.
用藤條打我.
但是約拿如果她怕敬畏的時候.
為什麼她不應該聽說.
怕被神罰.
反而神叫她去東.
她就去了西呢.
我們看看這班水手.
他們的改變.
那班水手聽到約拿所宣告.
這是創造天地萬物的主宰的時候.
他們就敬畏耶和華.
同樣都是敬畏耶和華.
而他們的行動是什麼呢.
他們的行動就是.
向神獻祭和許願.
許願有些英文的翻譯就是make vows.

$^{321}$即是立誓盟.
就好像在婚姻結婚裡.
新娘新郎.
來去立誓盟.
至於那班水手和船上的人.
和耶和華立什麼誓盟我就不清楚了.
不過我猜猜吧.
不知道會不會說.
我以前不認識耶和華.
所以就拜這些假神這些偶像.
現在我認識誰是獨一的真神.
我就離開這些偶像.
我就專一去獻祭奉獻和敬畏.
這個獨一真神.
我就和你去許這個願.
立這個誓.
不知道.
但是你看到.
他們的生命有改變.
又再看看尼尼未成.
約拿就去到尼尼未成.
他就去宣講.
看看尼尼未成的反應.
尼尼未成的人就信了神.
宣告禁食.
從大到小穿上麻衣.
各人懇切呼求神悔改.
離開惡行.
摒棄手中的強暴.
同樣原本那些尼尼未成的居民.
是作惡的.
窮凶極惡的.
以致約拿討厭他們.
不想他們能夠得到救恩.
也因為這個緣故.
要去到一個地步.
要宣判審判的訊息.
但是當這班居民聽到這個訊息的時候.
他們的反應就是相信神.
而且他們都有行動轉變.

$^{361}$他們悔改.
痛悔.
披抹.
蒙悔.
之前做惡事的.
有很多殘暴的事情.
聖經沒有仔細地說.
但是可能你看一些色經書.
就會看到他們的殘暴是很恐怖的.
但是他們居然因為這個緣故.
他們悔改.
回轉.
不再做這些事情.
或者之前是拜那些偶像.
那些偶像是掩飾很多邪淫的.
他們現在不做.
不跟他們.
悔改.
回轉.
我們看到原來.
當水手.
或者尼尼微的人.
經歷到耶和華的恩典憐憫之後.
他們的生命有改變.
那些水手遇到困難.
大風暴的時候.
他經歷耶和華的拯救.
他生命改變.
那些被罪纏繞的人.
他們聽到神的寬恕.
那個赦免.
赦罪之恩的時候.
他們生命有改變.
弟兄姊妹.
我不知道在你的生命裡頭.
當你無論是經歷.
你的生命的大風浪.
困境.
甚至.
差不多去死的時候.

$^{401}$會不會有經歷到.
神的恩典和憐憫.
讓你有這個信心.
跟神一起度過呢?.
又或者.
有沒有經歷過.
神向你的寬恕.
赦罪.
這個恩典.
以致你回轉.
悔改.
不再救我就除去.
不再走回以前.
纏繞我們赦罪的時候.
那些光景.
那些歲月呢?.
第三樣.
耶和華以恩典憐憫待人.
其實呢.
在這件約拿書的事當中.
不單止水手.
得著這個救恩.
李李未成.
得著耶和華的救恩.
其實.
約拿自己.
都很需要.
神的恩典和憐憫.
約拿是認識耶和華的.
我們在經文裡看到.
他從第一章已經說.
他說他是希伯來人.
就是說他應該是熟悉.
摩西的律法.
知道耶和華帶領他們.
出埃及.
過紅海.
這些神蹟奇事.
他甚至說.
他自己是敬畏耶和華.

$^{441}$他清楚知道.
這個耶和華是創造天地.
創造滄海陸地的神.
而在他在大禹的.
三日三夜裡.
他的祈禱.
裡面你看到.
有很多.
他都是引用詩篇.
詩人裡的禱告和呼求.
他知道.
他要仰望耶和華的聖殿.
他知道他在急難當中.
求告耶和華.
耶和華會應允他.
他在陰間.
來去呼求耶和華.
耶和華會聽他的聲音.
他二章六節那裡說到.
他的祈禱.
他說耶和華.
我的神啊.
你會將我的性命.
在坑裡頭拉上來.
約拿又知道.
他的禱告.
可以去到神的面前.
去到聖殿當中.
約拿也都很清楚.
知道救恩是出於耶和華.
其實約拿很清楚的.
他知道的.
他甚至在第四章那裡說.
他說我知道你是有恩典.
有憐憫的神.
不輕而發怒.
並且有豐盛的慈愛.
他說的沒錯.
因為這個正正是耶和華.
他自己這樣去介紹他自己.

$^{481}$當耶和華向摩西顯現的時候.
在荊棘火裡頭.
摩西問你是誰.
我怎樣去跟別人說.
耶和華在自我介紹裡.
他就有這樣去說.
他說.
他就是那個.
耶和華 耶和華.
是有憐憫 有恩典的神.
不輕而發怒.
並且有豐盛的慈愛和誠實.
這個是耶和華自己所說的.
他向摩西介紹他自己.
就是這樣.
這是他的屬性.
本質就是這樣.
不單止是這樣.
摩西聽到耶和華這樣說他自己之後.
當摩西有需要向神祈禱的時候.
摩西很聰明.
他引用耶和華他自己的屬性.
來向神祈禱.
在《民數記》的時候.
摩西向耶和華祈禱.
引述他所說的話.
用回.
你說你是這樣的神.
拜託你憐憫我們這群.
都頗不聽話的以色列人.
在曠野裡走來走去.
摩西用回這番說話來向神祈禱.
不單止這樣.
還有.
在詩篇裡.
詩人大衛他也清楚耶和華.
他也是用神的說話來向神祈禱.
所以其實如果大家讀聖經的時候.
你會發覺這句話.
在很多篇幅裡都出現過.

$^{521}$從不同的人口中裡出現過.
除了約拿這個先知.
還有約爾先知.
也同樣引用耶和華這句話.
他自己的屬性.
向那群南國的猶太人說.
你們要張開眼睛.
看清楚.
這個神是這樣的神.
你們快點回轉.
否則我們就會滅亡.
在這樣的境況.
約拿他頭腦上知道耶和華.
其實他有沒有真正經歷到.
耶和華的恩典和憐憫呢.
我們來看看.
其實整本約拿書.
耶和華所說的話不多.
只是五句.
我們一會兒看看是哪五句.
現在我們想在這裡看看.
約拿和耶和華的對比.
在第四章當記載到.
約拿知道神的屬性.
有恩典有憐憫.
不輕而發怒.
有豐盛的慈愛.
約拿因為這樣而發怒.
經文說了幾次他發怒.
甚至第九節說.
他說我發怒以致死是對的.
這樣發悔氣和耶和華說話.
我死了都可以.
這樣發悔氣和神對話.
他就發怒.
我們對比一下.
耶和華.
這亦是四章二節說.
你是有恩典有憐憫的神.
不輕而發怒.

$^{561}$且有豐盛的慈愛.
不輕而發怒.
有英文翻譯是slow to anger.
耶和華真的做出他所說的.
他真的慢很慢都不發怒.
我不知道大家有沒有試過.
有些人可能有些東西踩中你.
激動你.
你就會火滾發怒.
但你看看約拿一次又一次這樣對神.
神叫他去東.
他不聽話去西.
神有這個屬性.
他悔悔氣氣說了.
看到人們回轉.
他又不開心.
甚至向神發怒.
神對約拿沒有發怒.
神怎樣呢.
神很有耐性.
他用他的方法來教導.
引導約拿.
他想約拿.
不是只是頭腦上知道.
我是有恩典有憐憫.
他要約拿他真的經歷體驗.
這個有恩典有憐憫的神.
剛才我說了.
耶和華在這個約拿書裡面.
只說了五句話.
頭兩句就是叫約拿.
你要去尼尼尾城去宣講.
就是吩咐約拿的事.
之後的三句就在第四章.
其實這三句.
耶和華都是一個問題.
第四章的第四節.
約拿發怒.
耶和華對約拿說.
你這樣發怒對不對.

$^{601}$耶和華很有智慧.
他也很有耐性.
他用一個問題來引導約拿去反思.
第九節.
你因為這棵披麻樹這樣發怒.
對不對.
第二次問他.
然後第三次.
耶和華說.
這棵披麻不是你種的.
又不是你使它長大.
一夜成長一夜死去.
你尚且愛惜它.
何況在尼尼尾城.
其中不能分辨左右手的.
有十二萬那麼多.
而且還有很多畜生.
我豈不愛惜呢.
整本書就結束了.
沒有一個大團圓結局.
好像一個開放式結局.
究竟約拿的反應是怎樣呢.
究竟約拿之後是怎樣呢.
他的答案是怎樣呢.
他怎樣去回答神呢.
還是其實他沒有回答呢.
不知道.
聖經沒有記載.
但是聖經在這裡就讓我們看到.
神是要以恩惠慈愛.
來讓約拿去認識他.
他用恩惠慈愛去對待約拿.
如果是這樣說的話.
其實可能神想要去處理的是約拿這個人.
神是想得著約拿這個人的生命.
所以才差派他去尼尼尾城.
傳講這個訊息.
原來真是可能有時候.
我們遇到的事.
是神想得著我們的生命.

$^{641}$想我們認識經歷他的時候.
給我們的一些遭遇.
其實目的是要建立我們.
讓我們不單止是頭腦上認識.
這個恩典,憐憫的神.
也都在我們的心裡去接受.
初初的時候我提到.
我在阿爾巴尼亞短孫的時候.
認識在當地做長孫工場的一些孫教士.
他們提到他們頭一兩年.
其實都不是做什麼士工.
是要去學語言.
認識當地的文化.
甚至乎他們是要住在當地人的家裡.
一段長的時間.
我們有沒有住在神的裡面呢?.
這些孫教士和那些人一起住的時候.
他們就會認識那些人的文化.
原來你這樣吃飯的.
原來你這樣跟人打招呼的.
原來你們這樣去接待人的.
透過一起住,一起生活認識.
我們有沒有邀請神和我們一起生活呢?.
在我們的生活裡有沒有神在當中呢?.
其實不單在聚會裡我們有神.
其實可能在我們上班.
或者甚至我們娛樂的時候.
跟家人吃飯.
或者甚至去看電影.
我都可以邀請神在我的生命裡.
這樣的時候你就能夠.
不單在頭腦上.
在聖經裡知道耶和華的屬性.
在這些微小的事情裡.
你都可以經歷,跟祂同行.
認識耶和華是一個怎樣的神的時候.
我們無論在服侍,傳福音.
我們生活面對困難的裡頭.
我們都更加有力.
因為我們不是靠己力.

$^{681}$我們是依靠神.
有人說世界上最長的距離.
其實就是這裡.
就是腦子.
到這裡就是心.
我盼望姐妹們.
不要像若娜一樣.
在頭腦上知道耶和華很多的事.
而不是真真正正去認識,體驗,經歷這個神.
我們一起祈禱.
主啊,多謝你.
你真是那個有恩典,有憐憫.
輕易發怒,且有豐盛的慈愛的神.
這是你的屬性.
所以你對待若娜.
對待我們這裡在座每一個.
你都是這樣.
而且你讓我們能夠經歷到.
這些恩典,憐憫的時候.
我們生命得著改變.
我們生命得著力量.
我們無論去傳福音,去見證,去做事奉.
我們知道我們是基於耶和華的恩典和憐憫.
來成事.
我們因為耶和華的恩典,憐憫.
我們能夠有一個聖靈帶領的生命.
可以不去走在罪的當中.
也有力活出活潑的生命.
神啊,你真是不輕易發怒.
有豐盛的慈愛.
你對待我們每一個.
其實我們自己心知.
你對我們多有耐性.
一次又一次的給我們機會.
一次又一次的用你的慈情,愛愴.
牽引我們回來.
我們知道.
主,我們希望不是只是停留在頭腦上的知道.
我們盼望我們的心一心一意跟隨主.
和我們弟兄姊妹一起彼此鼓勵.

$^{721}$走這條天路.
我們這樣去祈禱,感恩,交托.
是奉主耶穌基督得勝的名字祈求.
阿們.
願主祝福大家.
\newpage



\section{以弗所書馬可福音 4:14-16}
\label{sec:xr_m3_qAqFs}
\textbf{2023-10-28 晚堂《在愛中成長》以弗所書 4\_14-16,馬可福音 12\_29-30 講員:楊明麗姑娘}
\newline
\newline
連結: \href{https://youtube.com/watch?v=xr_m3-qAqFs}{\texttt{https://youtube.com/watch?v=xr\_m3-qAqFs}} ~~~~ 語音日期: 2023-10-31
\newline
\newline
\hyperref[sec:LM01sAVsGrw]{\small{< < < PREV SERMON < < <}}
~
\hyperref[sec:index_chronic]{\small{[返順時目]}}
~
\hyperref[sec:index_scriptual]{\small{[返順卷目]}}
~
\hyperref[sec:ROGX9cAg8Cg]{\small{> > > NEXT SERMON > > >}}
\newline
\newline
$^{1}$各位晚安.
這張照片大家都很熟悉的.
這裡是哪裡呢?這裡是對面.
有一天我在這裡拍照.
拍到對面的情況.
我發覺天空很特別.
天空的顏色很特別.
天空是很有線條的.
所以我拍下了它.
到要講今天的講道的時候.
我就覺得這個很好用.
今夕何夕呢?明天是我們的唐慶.
記得是幾多週年嗎?.
四十二.
明天是我們四十二週年的唐慶.
其實今晚的講完不是我.
我只是因為發生了一些事故的緣故.
所以我今晚在這裡跟大家分享這個訊息.
我們很感恩有這個偶然的機會.
今夕何夕?.
四十二年前大家在哪裡呢?.
你說明天是我們四十二週年的唐慶.
四十二年前我們在哪裡呢?.
我想四十二週年.
在四十二年前.
感受唐.
是一個只是旅遊的教會.
我去的年日是不多的.
聽聞一些弟兄姊妹的分享.
就是當時他們買了一個地方在這裡住.
不是.
買了一個地方在這裡周末住.
周末的時候就會帶兒女過來.
好像一個周末的旅行.
就會在這裡度假.
所以這個是一個度假教會.
這個是感受的第一身.
是一個度假教會.
也有傳道人.
他們很體重.

$^{41}$雖然是周末的旅遊時期的時候的教會.
但是他們很認真地去策劃和牧養.
讓到教會可以漸漸地成長.
令到教會一直有成長的時候.
就有居民住在這裡.
不是周末進來旅遊的.
就在這裡繼續聚會.
在這40年當中.
這12年當中.
做了很多的事.
我聽聞過.
當年的關故部是怎樣開始的呢.
當年的關故部是去到護老院當中開始的.
去到護老院當中的時候.
他們去探長者的時候.
見到原來長者不單是缺乏關心要被探訪.
而且是長者的能力衰退.
大家怎樣吃橙的.
原來吃橙是很不容易的.
他們就發覺原來吃橙.
老人家是自己不能夠將橙去皮.
原來吃橙是有很好的纖維.
我們吃橙都很需要纖維.
但他們原來是吞不到纖維.
會再嘔吐出來.
所以去探訪的人可能就是撕了橙.
餵他吃的時候.
發覺他吃不到一塊.
然後吐到他全身都是.
那個探訪者.
當然被探訪的人也吐到全身都是.
於是他們是要即場清理.
然後繼續半解一點.
讓他可以吃到.
原來他們是穿著被嘔吐物.
玷污了的衣服.
繼續在試縫探訪.
不是說好 我做完這次.
就這次.
下個禮拜你看我來不來.

$^{81}$原來不是.
他們繼續這樣做關顧.
不是做一個禮拜.
不是做一年.
做了很多年.
然後他們發覺.
不單是橙他們吃不到.
吃粥也很有困難.
他們的愛心去到每一個地方.
當時40年前的教會在哪裡呢.
是在對面的幼稚園.
我相信在座有人.
你是當年在那裡幫忙招巷午茶的.
早上去到.
早上去到幫忙砌牆.
砌主人殼的房.
砌崇拜的房.
然後崇拜了.
可能沒有空去崇拜.
去了哪裡呢.
我相信Eric當時也在那裡.
去了兒童主人殼.
不過是在教你.
不是現在你教人.
是 當時是這樣.
然後快快又快又拆東西.
變了有桌子.
做甚麼.
一起吃午飯.
再繼續有活動.
當年原來是這樣的.
是這麼溫情.
明天有沒有聚餐.
明天你有沒有報名聚餐.
你猜明天聚餐會有多少人呢.
我不知道.
不過報名的超過300.
當時好像是旅遊的教會.
之後就可以有禮拜堂.
有主人殼.

$^{121}$有兒童主人殼.
也有成人的崇拜.
慢慢地一步一步.
為甚麼會來到這裡呢.
因為不夠位子坐.
因為很擠.
不是不想去昭和午茶.
而是真的很擠.
不夠位.
不夠位子給兒童.
也不夠位子給成年人.
於是就要很努力去禱告.
看怎樣可以有開展.
後來我們有這個地方.
這個地方是怎樣來的呢.
這個地方原來是弟兄姊妹.
一元一元奉獻出來的.
我們沒有寫哪一張椅子是哪一個奉獻的.
沒有寫哪一道簾.
哪一塊玻璃是哪一個奉獻的.
但是每一樣東西.
是弟兄姊妹奉獻回來的.
是當年你們有份奉獻回來的.
今天我們繼續.
有兒童主人殼.
都是在那邊.
幼稚園那邊.
有青少年的聚會.
有天崇.
都是在幼稚園那邊.
那邊我們沒有荒廢.
不過換了人了.
換了什麼人呢.
當年在那裡崇拜的人.
他們的子女.
當時在那裡讀著兒童主人殼的.
現在在那裡教著兒童主人殼.
你太太.
接著呢.
當年在那裡崇拜的他們的孫子.

$^{161}$就是你們的子女.
現在在那裡兒童崇拜.
三代人.
最少三代人.
如果當時.
回來的是青年.
他帶著自己父母現在回來的.
可能是四代或者五代人去說.
四十二年.
我們不是繼續留在一個.
是旅遊的人周末過來崇拜.
有沒有這樣的人呢現在.
有 還有的.
有人單車.
過來元朗.
進來感受.
都會回來崇拜.
有人過來香港一日遊.
都會過來香港崇拜.
來到感受.
但是這個不是我們的主體.
我們的主體.
就是四十二年來.
神一直帶領你們.
一個一個回來.
一個一個踏進來.
一個一個愛恕.
幫忙去建立這間教會.
不單止是你的一元.
一元 一元 加一元.
是你的一分力.
加一分力 加一分力 加一分力.
建立出來的教會.
我們預備.
明天四十二週年.
這個是神給我們的恩典.
明天.
明天我們是一個聯合的崇拜.
明天是我們一個聯合的崇拜.
我們有家恩 有雄恩在我們當中.

$^{201}$我們有青少年.
在我們當中.
我們有兒童師班在獻師.
我們有非語 印語的在我們當中.
四十二年.
神在我們當中做了很多事.
做了很多事.
是我們想不到的.
甚至我們不是很知道的.
但原來.
祂一直沒有放棄.
一個當時是旅遊的人.
進來過週末的時候回教會.
我們是有今天的.
我們今天是有好幾個聚會的地方.
是神賜予給我們的.
我們的今天是為了什麼.
我們今晚是為了預備明天.
預備明天我們有週年的堂慶.
我們也預備明天之後.
我們進入的是第四十三週年的開始.
我和大家一起看一段聖經.
為什麼會想起這一段經文呢.
這一段經文是說.
教會是神所建立的.
教會是神在基督裡面建立的.
基督是我們的頭.
我們是基督的身體.
教會是基督的身體.
在這個身體當中.
神一直建立我們的時候.
有些說話我們可以護免的.
他說.
我們不再作小孩子.
眾了人的詭計和欺騙的法術.
被一切邪說之風搖動.
飄來飄去.
我們在座只有兩個小孩子.
我們一向都已經不再是小孩子了.
是的.

$^{241}$但是他的說話說到.
小孩子的時候不是說我們的年紀.
或者我們的身形.
他是說到在信仰上面.
我們要不飄蕩.
我們要穩定.
才不是小孩子.
什麼原恩惠飄來飄去呢.
其中一樣東西是會被人欺騙.
信仰有沒有欺騙的.
大家有看新聞.
信仰有沒有欺詐.
是不是要交費才有人為他祈禱呢.
在欺詐當中.
還記得嗎.
每一次祈禱要收多少錢呢.
那些書是不用的.
兩個星期前我們有活游祈禱會.
一直都有活游祈禱會的.
這麼多年.
有沒有收費的呢.
沒有.
都沒有叮囑一句祈願圖.
出去你自由奉獻.
沒有.
祈禱嘛.
被祝福需要奉獻嗎.
不用.
這個是恩典.
神給我們的恩典.
我們可以禱告.
我們可以彼此代禱.
《雅各書》第五章更加講到.
有什麼需要就真的要找弟兄姊妹.
還有教母去為他禱告.
是的.
原來社會當中有很多詐騙.
甚至用信仰的名堂都可以詐騙.
就算不是詐騙在金錢上面.
會詐騙在信仰上面.

$^{281}$更加嚴重.
詐騙在金錢上面.
你有少少損失.
錢還可以再賺.
但在信仰上面.
沒有一個正確的信念的時候.
被人騙了的時候.
我們會進入異端.
會被欺騙.
更加嚴重.
我們的生命所需要的.
被欺騙.
所需要的信仰.
所需要的信念.
所需要的生命資料.
我們被欺騙的時候.
這個更嚴重.
還有一些東西.
我們可能會自己飄來飄去.
弟兄姊妹為什麼你會坐在這裡呢?.
因為你回來感受堂崇拜.
還有一些人.
他不是每個星期都回到一間教會.
他還會飄去不同的堂會.
他不是飄去某幾間堂會.
可能是飄去不同的堂會.
讓自己能夠屬於一間教會.
都是我們的需要.
讓我們能夠安定.
在教導上面能夠有系統.
被造就.
神都想我們在一個穩定的地方.
穩定的聚會.
這一書.
我們反而要用愛心說誠實話.
各方面向著基督長進.
連於元首基督.
這一節的金句.
很多人都背了.
也有些可能成為了團契的團訓.

$^{321}$我們很重視這一節的聖經.
都是很好的.
我們會認真去看.
但是我又給多兩個不同的譯本.
大家看一下第十五節.
第十五節說.
相反我們要以愛心真誠地生活.
在各方面朝向他長進.
他就是頭.
是基督.
這裡說的他.
是在說耶穌基督.
我們各方面朝著耶穌基督長進.
我們不只是在言語上面.
我們是不是要說誠實話.
他在說的誠實.
原來不是在說話而已.
是在說我們的生活.
生活的態度.
我們是不是用愛心.
以誠實的態度來生活.
這個是新漢語譯本.
還有一個譯本.
是環球聖經譯本.
他這樣說.
他說.
而可以在愛中誠實做人.
在各方面長進.
達到基督的新量.
他是教會的頭.
他這一節不是說話.
我們只是說誠實話.
說誠實話.
我聽過很多不同的人問過的問題.
我們不是說要說欺騙人的話.
不過這說話不只是在言語上.
而是我們要做一個誠實的態度.
去真誠的態度去做人.
所以我問你戶口的密碼的時候.
你不需要給我的.

$^{361}$那個是你真誠對自己的時候.
和真誠對我的時候.
都不應該將你戶口的號碼給我.
但是呢.
以真誠的態度去生活.
我們以真誠的態度去生活.
無論這個是我們要做的事情.
你的工作.
你待人.
或者教會當中的相處.
還是我們對神.
我們在事奉上面.
我們原來需要一樣東西.
就是要真誠.
是真誠去處理.
真誠去生活.
真誠不一定我們時時做的事情都很對.
但是在態度上面.
我們就是用真誠來面對.
當我們發覺自己有不正確的時候.
我們都要真誠去面對這個不正確.
真誠來承認自己的錯.
這樣才可以有新的開始.
他說我們需要真誠.
原因是我們不再是孩子.
我們也不受欺騙.
我們不飄來飄去.
我們以真誠的態度去生活.
不過其實14至16節.
14節和15節都不是這一段的中心點.
這一段的重點.
應該在第16節.
對不起.
錯過了第16節.
第16節是說什麼呢.
全身靠著他.
藉著每一個關節的支持.
照著每一個部分的功能.
得以配合 聯繫.
讓身體漸漸成長.

$^{401}$在愛中建立自己.
這個自己是在說.
不是你和我.
是在說教會.
是在說我們大家一起.
甚至是在說耶穌基督整個普世的教會.
我們是在裡面靠著基督.
我們是各個關節的支持.
然後每一個部分能夠發揮功用.
這個是神給我們的.
原來我們的功能.
我們每一個都有我們的功能.
不過在教會的事奉上.
和你外面的工作都有點相似.
就是不是一個人做完所有的事.
不是只是做了自己的那部分.
我們原來有一些關節位.
是彼此聯繫著.
而怎樣能夠聯繫得合適.
是靠著耶穌基督.
我們去聯繫得合適.
這些關節位重要嗎.
原來醫學上有一科.
不是骨科.
有一科是關節科.
原來關節很重要.
走路能不能沒有膝蓋呢.
就這樣大腿接著小腿.
手臂能不能沒有手肘呢.
就這樣一隻手.
很多事我們做不到.
沒有關節位.
沒有轉折位.
我們是動不了的.
如果轉折位受傷.
不能夠好好運作的時候會怎樣呢.
你問問有膝蓋痛的人.
他可以解釋得很清楚給你聽.
他以前可以怎樣走.
現在又要怎樣走.

$^{441}$昨天我和一位談電話.
他就說我搬了出去元朗住.
因為我不想走那麼多樓梯.
我的膝蓋不能夠支持我走那麼多樓梯.
原來關節是很重要.
關節是要來支持的.
關節不一定要來出力.
但沒有關節很多事是做不到的.
我們全身有很多重要的位.
有很多是很有能力的位.
但有一些位是關節.
是他要去支持著.
讓其他的位能夠去運作.
我記得有一年.
有一位弟兄他下小巴的時候.
因為小巴錯了一錯.
他就很快地掉了下來.
他不是跌倒.
他只是從小巴掉了下來.
他不是掉在地上.
但他站穩在那裡以為沒事.
過了一兩天他發覺.
他的尾指又腫又痛.
原來就是這樣動一動.
尾指的骨斷了.
於是就要包著他.
有一塊小石膏.
讓他癒合.
當然短期內又不能穿鞋.
短期內又不是那麼帥.
最嚴重原來不是這樣.
嚴重的是以後那個小小的腳趾尾.
是影響他的平衡.
只是一個很小的一件.
弟兄姊妹回來教會.
我們沒有很小的一個人.
沒有.
每一個都很重要.
我們覺得他無關痛癢.
無關重要.

$^{481}$什麼都不用找他做.
不是 原來他不在.
我們運作不了.
原來人平衡.
沒有了腳趾尾是平衡不了.
當我們有腳趾尾的人.
我們是不明白的.
我們看了就覺得.
很多問號.
為什麼不行.
你還有四隻腳趾.
原來醫生真的說是不行的.
是沒有以前平衡得那麼好.
弟兄姊妹我們回來.
每一個都很重要.
我們在這裡.
每一位都很重要.
我們以真誠的心.
一起結連.
以真誠的心.
在教會當中.
彼此關心.
一起成長.
才能夠建立到這間教會.
當年大家.
在幼稚園那邊聚會的人.
很多都回來這裡.
不過過了這二十多年.
不是每個人都回來繼續崇拜.
有些已經去了天國.
有些可能移民了.
搬到別處.
有些可能要在家中崇拜.
有很多不同的因由.
但是.
每一個都是這間教會的一份子.
每一個都是聯合起來.
一起建立起來的教會.
每一個都有份.
於這個成長.

$^{521}$曾經教導兒童主義學的.
他的學生.
現在是兒童主義學的老師.
甚至是崇拜的主席.
是青少年的導師.
還有青年部的執事.
甚至是堂董.
當年所做的.
每個星期回來教的主義學.
並不是一件小事.
你是在建立今天的教會.
今天我們在這一書.
我們學習地預備去建立.
明天的教會.
明天的感宣.
最近我一直在思想這一節的經文.
其中一個原因就是很多字不懂.
大家都很懂得背.
這個就是第一要緊的.
耶穌這樣說.
第一要緊的就是.
以色列你要聽.
主我們神是獨一的主.
你要盡心盡性盡義盡力愛主你的神.
這些是我們每一個字都讀得到.
又不是很明白的.
為什麼一開始.
耶穌回答他的時候.
一開始就要用.
以色列你要聽.
神是拯救以色列的.
主耶穌基督也是拯救我們的.
但如果有一天.
神或主耶穌基督跟我們說.
你聽著.
我們是不是會很認真地去聽呢.
我想是的.
不過其實我們聽得明白多少呢.
聽得到多少呢.
我想其實聽這個字.

$^{561}$我們真的要很認真地去學.
大家是不是很聽得懂別人說的話.
我有一個經驗就是.
當我過了兩個星期很忙的時候.
過了之後.
我先生就會說.
我那天跟你說了什麼.
你都不記得了.
是的.
你說了什麼.
嗯了你就不會再繼續說.
我說yes了你還說.
聽不到你說什麼.
聽了我都會不記得了.
我會放在一個我不需要想起的位置.
不需要被它騷擾我的腦中的任何一個位置.
不會再找回它出來的檔案.
是不是.
原來是的.
聽.
聽是我們要慢慢學的.
我們會聽.
聽了之後不知道放在哪裡.
還有一些東西呢.
是當年我爸爸說的.
說了你都不明白.
我現在回想一些當年的事.
我就會覺得.
是的我當年真的不明白.
現在我明白多少呢我不知道.
是真的不明白的.
當你教小朋友的時候他不明白.
是很正常的.
他真的不明白.
他不是不聽.
不過是聽不完.
天父要跟我們說話的時候.
我們認不認真呢.
我們認真的.
我相信弟兄姊妹.

$^{601}$他們對這本聖經都很認真.
不過我們這麼認真的時候.
從聖經當中.
我們真的聽了多少.
留在心裡呢.
我們真的還要去聽.
這個是我覺得很重要的.
我們需要去聽.
我們需要聽.
我們需要放在心裡.
因為祂是我們的神.
我們再沒有第二個主.
第二個神.
再沒有其他的拯救.
我們聽祂說的時候.
我們需要做一件事.
這裡出現得最多就是四個「盡」字.
要怎樣呢.
盡心盡性盡義盡力.
愛主你的神.
大家都識背.
寫都沒有問題.
你可以默它出來.
不過這四個字.
心 性 義 力.
一向都是我們的問題.
心性義力是什麼呢.
原來心呢.
不是說心.
是說腦.
心是說思想的東西.
我們大家知道.
心的功能是用來泵血的.
跟血的運行有關.
腦是用來思想的.
心原來是盡心.
我們要好好去想.
我們要盡我們的心思意念.
去愛我們的神.
嘩.

$^{641}$我們一天有多久.
有多多.
有多隆重地去.
用我們的心去想愛神呢.
我們不會說.
今天沒有什麼祈禱.
很掛念天父.
我們不會說.
是喔.
天父喔.
我今天應該要怎樣去.
更加表達今天愛祂呢.
我讀哪一道的聖經呢.
我繼續怎樣去讀聖經.
或者讀多少呢.
出了很多本.
跟聖經有關的書.
最好就不要買.
因為買了要負責任去擺它.
擺完它之後.
你每天都會受到影響.
我們會想很多.
跟愛神不是很關的東西.
我們擺下它.
我們的心思用了多少去愛神呢.
我記得是.
是.
大少曾牧師離世前幾個月.
我有一個很意外的機會.
跟他同席吃飯.
那個時候.
其實他身體一般般.
我已經穿衣服穿得多.
因為我很怕冷.
他起碼比我多穿一件大褸.
是真的很厚的大褸.
坐在那裡吃飯.
可想而知他身體.
當時已經不是很行.
但是他仍然要赴宴.

$^{681}$原因就是那一次.
是一個猜拳的聚會.
他有些事情.
他要去談 要去說.
一個過了幾個月就離世的人.
他用盡他的生命.
在那個時候他仍然掛念著.
宣教的事.
仍然掛念著內地會的發展.
不是在中國大陸 是在外地.
在歐洲 非洲方面.
內地會的發展.
他還在想什麼呢.
他還在想.
怎樣可以去支持這個發展.
其實他當時說話的聲音已經很小.
盡他一生去事奉神.
他的心思所牽掛的.
去到最後的年日.
仍然是宣教的事.
我們心靈裡面.
可以有多重視呢.
有一個人.
我曾經.
很怕他比我早回教會.
他坐在那裡.
我想他也知道我在說他.
因為我回來他一定已經在那裡.
他負責在那裡 差不多開門.
他未到九點就到了.
他不是工作人員.
但是他每天都差不多在那裡.
為什麼呢.
他很多事情回來做.
我真的不明白他為什麼要回來.
我覺得他很多事情都懂.
他為什麼要回來做呢.
後來我才知道原來很多事情他還在學.
在音樂上他的造詣很深.
很多他都懂.

$^{721}$但原來他還是回來學不同的樂器.
他的心思仍然是要回來做什麼.
去學什麼.
他的心思也是.
今天回來教會的時候.
教會哪裡有什麼.
他要去幫忙整理.
整一整的呢.
因為我後來發現原來他回來.
他真的帶了一個心來.
弟兄姊妹我們的心.
我們的心思.
我們的心意.
向著神嗎.
當你有心向著神的時候.
你的力就會向著神.
力是什麼呢.
原來力.
前面後面加上去的東西都跟他有關.
力講的是力氣.
你做事很快.
跑步很快.
你的力氣.
力講的是你的能力.
原來大家可能有很多不同的能力.
游泳可能是你的能力.
跑步是你的能力.
爬山是你的能力.
唱歌是你的能力.
音樂是你的能力.
藝術是你的能力.
收拾東西是你的能力.
煮東西是你的能力.
購物是你的能力.
你可以有很多不同的能力.
原來這些能力神給過你的時候.
你可以用這些能力來事奉神.
用這些能力去表達你對神的愛.
不是.
不是用這些能力去表達你對神的愛.

$^{761}$是用盡這些能力去表達你對神的愛.
Fanny Crosby.
一個盲眼的人.
竟然能夠作出美麗的詩歌.
Nick Woodward.
大家認識他.
無手無腳的報導家.
他們是用盡他們的心.
用盡他們的力氣.
去愛這位神.
我們每個人都有我們的心和力量.
你有你醫學上的能力.
你有你電子上的能力.
甚至有人有他寫字的能力.
是不是可以愛神.
原來可以的.
每個人年初一都有人來幫你寫寫揮春.
什麼是你的能力呢.
什麼是神給你的.
想你要來愛他的能力呢.
畫畫 設計.
和諧粉彩告訴我們.
填顏色也是.
明日的愛言.
不是在外面訂了什麼回來.
大家在廚房拿出來的.
你可以煮.
你可以撿.
原來都是神給你愛教會.
愛神的能力.
我漏了.
我煮不了.
怎樣呢.
今天有三位姐妹回來.
洗了三百多隻碟.
三百多支叉.
三百多支羹.
三百多支.
我們可以用神給我們的.
我們去愛神.

$^{801}$這是神給我們的.
我們盡力.
會不會做得太多.
三百多個我都覺得多.
但他們說不洗不安心.
不洗還有些塵.
突然姐妹這麼好.
這麼有愛心回來聚餐.
接著吃了些髒東西.
下肚.
然後每個禮拜她不舒服.
於心不忍.
是不是做得太多.
我說不是.
他說四十二週年有多少次.
一次.
為誰做.
為了要來的人做.
為了教會做.
為了神做.
三樣都是.
盡力愛主.
我們的神.
我們愛神.
我們需要祂幫我們.
我們才能更加愛祂.
我們需要祂幫助我們.
能夠關節與骨.
能夠連絡得合適.
然後每一個部份.
才能夠發揮到功用.
突然姐妹有時自己一個人做事容易些.
是不是.
好像是.
但自己一個人做到多少.
自己一個人洗三百多隻碟.
三百多支叉.
三百多支羹之後.
你就覺得以後都不想再洗東西了.
但是.

$^{841}$有第一次姐妹一起的時候.
好像時間過得快些.
除了水一直流之外.
大家彼此的分享.
彼此的鼓勵.
都一直流暢.
我們需要一起去合作.
原來合作不容易.
合作需要關節.
關節不是外事.
關節是支持.
是神的設計.
讓我們整體的身體有支持.
神的設計和引導.
讓我們每一個每一個我們進來了.
我們彼此要支持.
一起以行動去愛我們的神.
我們彼此支持.
以行動愛我們的神的時候.
我們才能夠.
大家做到那種愛.
盡力去愛.
我們明天做一件事.
是為教會的生日.
今天我們還是在做準備.
為明天做準備.
弟兄姊妹我更加願意大家.
我們真的是為明天.
為明天的教會做準備.
神使萬有都要歸於基督的腳下.
並且去賜教會做萬有的元首.
教會是基督的身體.
是各位充滿萬有者所充滿的.
我們能夠在教會當中.
一起彼此連繫.
一起事奉.
都是神的恩典.
神所預備的.
讓我們都一起為明天.
我們付上努力.

$^{881}$為明天付上我們的一分力.
我們一起低頭祈禱好嗎.
我們一起禱告.
親愛的天父上帝.
我們來到你的面前.
我們很希望.
我們能夠有好的預備.
去慶祝明天的週年紀念.
我們更加願意是你使用我們.
讓我們為著明天及以後的教會.
付出努力.
求主你加力量給我們.
我們願意將我們的力量.
我們的能力.
我們的心力.
我們的體力.
奉獻給你.
求主你紀念.
幫助我們.
主耶穌基督是我們的元首.
祈求你讓我們有智慧的眼力.
可以跟隨基督.
時時跟隨祂.
不轉離.
求主讓基督充滿我們和我們的教會.
求主讓我們每一位.
牧者和導師.
都得著你的恩典.
能夠在牧養上.
教導上.
都好好地去盡責.
也讓我們能夠在宣教和拓展上.
常常有你給我們的眼光.
有進一步的進展.
求主幫助我們.
讓我們年少的一代.
能夠起來.
將來一起事奉你.
祈禱奉主耶穌基督的聖名.
\newpage



\section{腓利門書 1:4-20}
\label{sec:ROGX9cAg8Cg}
\textbf{2023-10-29  42周年堂慶聯合崇拜《友伴．前行》腓利門書 1\_4-20 講員:周曉暉牧師}
\newline
\newline
連結: \href{https://youtube.com/watch?v=ROGX9cAg8Cg}{\texttt{https://youtube.com/watch?v=ROGX9cAg8Cg}} ~~~~ 語音日期: 2023-10-31
\newline
\newline
\hyperref[sec:xr_m3_qAqFs]{\small{< < < PREV SERMON < < <}}
~
\hyperref[sec:index_chronic]{\small{[返順時目]}}
~
\hyperref[sec:index_scriptual]{\small{[返順卷目]}}
~
\hyperref[sec:j9CuOv5EhDA]{\small{> > > NEXT SERMON > > >}}
\newline
\newline
腓利門書 1:4-20
\newline
\begin{longtable}{cl}
\hline
\hline
章節 & 經文 (和合本修訂版)\\
\hline
1:4 & \begin{tabularx}{0.7\textwidth}{X} 我在禱告中記念你的時候，常為你感謝我的神， \end{tabularx} \\ \\ \relax
1:5 & \begin{tabularx}{0.7\textwidth}{X} 因聽說你對眾聖徒的愛心，和你對主耶穌的信心。 \end{tabularx} \\ \\ \relax
1:6 & \begin{tabularx}{0.7\textwidth}{X} 願你與人分享信心的時候，能產生功效，讓人知道我們所行的各樣善事都是為基督做的。 \end{tabularx} \\ \\ \relax
1:7 & \begin{tabularx}{0.7\textwidth}{X} 弟兄啊，由於你的愛心，我得到極大的快樂和安慰，因為眾聖徒的心從你得到舒暢。 \end{tabularx} \\ \\ \relax
1:8 & \begin{tabularx}{0.7\textwidth}{X} 雖然我靠著基督能放膽吩咐你做該做的事， \end{tabularx} \\ \\ \relax
1:9 & \begin{tabularx}{0.7\textwidth}{X} 可是像我這上了年紀的保羅，現在又是為基督耶穌被囚的，寧可憑著愛心求你， \end{tabularx} \\ \\ \relax
1:10 & \begin{tabularx}{0.7\textwidth}{X} 就是為我在捆鎖中所生的兒子阿尼西謀求你。 \end{tabularx} \\ \\ \relax
1:11 & \begin{tabularx}{0.7\textwidth}{X} 從前他與你沒有益處，但如今與你我都有益處。 \end{tabularx} \\ \\ \relax
1:12 & \begin{tabularx}{0.7\textwidth}{X} 我現在打發他回到你那裡去，他是我心肝。 \end{tabularx} \\ \\ \relax
1:13 & \begin{tabularx}{0.7\textwidth}{X} 我本來有意將他留下，在我為福音所受的捆鎖中替你伺候我。 \end{tabularx} \\ \\ \relax
1:14 & \begin{tabularx}{0.7\textwidth}{X} 但不知道你的意見，我不願意這樣做，好使你的善行不是出於勉強，而是出於自願。 \end{tabularx} \\ \\ \relax
1:15 & \begin{tabularx}{0.7\textwidth}{X} 他暫時離開你，也許是要讓你永遠得著他， \end{tabularx} \\ \\ \relax
1:16 & \begin{tabularx}{0.7\textwidth}{X} 不再是奴隸，而是高過奴隸，是親愛的弟兄；對我確實如此，何況對你呢！無論在肉身或在主裡更是如此。 \end{tabularx} \\ \\ \relax
1:17 & \begin{tabularx}{0.7\textwidth}{X} 所以，你若以我為同伴，就接納他，如同接納我一樣。 \end{tabularx} \\ \\ \relax
1:18 & \begin{tabularx}{0.7\textwidth}{X} 他若虧負你，或欠你甚麼，都算在我的賬上吧， \end{tabularx} \\ \\ \relax
1:19 & \begin{tabularx}{0.7\textwidth}{X} 我必償還。這是我—保羅親筆寫的。我並不用對你說，甚至你自己也虧欠我呢！ \end{tabularx} \\ \\ \relax
1:20 & \begin{tabularx}{0.7\textwidth}{X} 弟兄啊，希望你使我在主裡因你得益處，讓我的心在基督裡得到舒暢。 \end{tabularx} \\ \\
[1ex]
\hline
\hline
\end{longtable}
$^{1}$各位感宣家弟兄姊妹 主女平安.
很想在這個堂慶日子裡面.
送兩句話給大家.
災於活水江河 靈明日進.
蘊納基督盤石 誠摯宣道.
當中是蘊含著我們宣道會.
一個很強調的精神.
我們不單止自己的靈明要盡心.
我們也繼續承擔主所託付給我們的使命.
坐在這裡也謹代表自己的教會.
北宣去恭賀感宣家42周年的堂慶蒙恩.
也代表宣道會 香港區聯會去恭賀大家.
剛才主席幫我們讀了這個肥理門書.
今天和你看這段經文.
我們到底在這裡當中可以領受到一個怎樣的訊息呢.
聖經裡面又是怎樣看一個教會.
一個群體的關係是怎樣呢.
當肥理門書這樣寫的時候.
這裡提到.
我聽見你對主耶穌和眾聖徒有愛心和信心.
你看到那個愛的字.
你就知道牧師你是說彼此相愛.
不單止這樣 在這裡更加提出.
在這段經文裡面很強調一個中心思想.
就是這班彼此相愛的人是有團契關係的.
你說我知道 我也會回團契 回小組.
或者我們一起查經.
不單止是那個合意的小組群體.
而是到底我們這個家.
這個整個教會是不是有彼此團契的關係.
這個更加重要.
團契這個字的原文在新約聖經出現了有20次.
其中有一次出現在肥理門書.
剛才我聽經文也不是很看得懂.
待會會告訴你在哪裡.
肥理門書是保羅在監獄裡面所寫的信件.
信當中保羅為一個從肥理門家裡逃走出來的奴僕.
或者奴隸歐尼西姆去求情.
歐尼西姆逃跑之後輾轉遇到保羅.
也因此信了耶穌.

$^{41}$表面上這封信是處理一個個別的事件.
修補兩個人的關係.
但實際保羅是行文之中讓大家體會到.
他做了三個的示範.
我們走在這條天路裡面.
我們說的群體 我們說的教會.
在這條天路裡面怎樣可以得到友伴.
一個信仰群體又可以怎樣一起前行.
我們會透過這封信一起去學習一下.
我們自己的群體 我們的教會生活.
是否也可以向著這個目標.
友伴 前行 繼續邁進.
我們一起將這段聽到的時間再交代給主.
剛才我們在敬拜的當中.
我們很感受到這個家是由很多人去組成.
亦都彼此願意同心.
我們知道我們的目標是對準你.
亦領受你給我們的使命.
現在我們打開你的話語.
願聖靈指教我們當明白的真理.
多謝你.
我們的祈禱是奉主耶穌基督得勝名字祈求.
阿們.
在剛才的四至六節那裡.
我們再一起看看.
不如我請大家一起和我讀這三節經文.
預備起.
我聽見你對主耶穌和眾聖徒有愛心和信心.
我每逢禱告提到你的時候.
願你與眾人分享你的信心的時候.
會產生功效.
使我們可以知道在我們中間的一切善事.
都是為基督作的.
這裡提到愛心.
在第四五節提到.
他說:肥利們你有的愛心不是掛在口邊的.
根據聖經學者的分析.
在肥利們的年代曾經發生過一些天災.
極之可能的就是肥利們見到教會裡面.
因災難而生活缺乏的人.

$^{81}$他們缺少生活的物資,錢糧.
於是肥利們就用自己所有的去幫助這些弟兄姊妹.
其中第五節更加貼近原文的翻譯是這樣的.
因為我聽說你對眾聖徒有愛心.
對主耶穌有信心.
原來在這封信當中不斷提到的愛.
對象是弟兄姊妹.
是主內的弟兄姊妹.
是基督徒之間.
是聖徒的群體有的彼此相愛.
這個愛心不只是我們對我們見到的人.
在街外的人有連續恩賜.
而這種愛心在我們的群體都要常常有.
或者我們都稍為在想一想.
當這個時刻問你.
你在腦海裡面想起哪一兩位敢宣家的弟兄姊妹呢?.
你應該想像到 馬上想到.
現在常常和你WhatsApp通訊 微信的弟兄姊妹是哪些呢?.
可能你有時候想到的人不是這些很關切的人.
可能是你要不斷學習以愛心去對待的人.
會不會有些時候自己其實對教會人的愛心.
是比對自己家裡的人多.
在我們真實的家庭裡有時候是這樣的.
在教會的家裡會不會都是這樣呢?.
對外面的人會寬容.
但卻是對本應親近的人.
有時候會多了少少計較.
在第六節裡說.
願你與眾人分享你的信心的時候.
這個分享其實就是我一開始說的團契的字.
在聖經的原文就是團契的字.
當然這個團契不是限制一個小組一個群體.
而是可以說是一個教會一班人.
聚集一起敬拜常常見面的一班人.
在這個其中 分享是否這樣呢?.
是有分享的 我見到外面擺了很多食物攤位.
一會兒大家會分享食物.
不單止這樣.
對於肥利們 保羅是否想跟他們說.
記住了 肥利們.

$^{121}$你這麼有愛心 你就要有愛心到底.
剛才一開始說過.
這封信表面上是說兩個人要修補關係.
是不是保羅想說.
你去原諒那個得罪你的人.
不是.
保羅不是這樣去教肥利們.
或者命令肥利們.
你要這樣這樣做.
保羅是首先說我去分享我怎樣實踐.
因著信耶穌基督而對人能夠有的愛心.
對我身邊的聖徒有的愛心.
在第八第九節他這樣說.
他本應可以吩咐肥利們.
因為教會的群體裡面他的位份是很高的.
但是保羅沒有說一人少一句.
保羅也沒有計較自己的使徒地位.
你就聽我說 給我面子.
他不是這樣.
他也沒有計較自己的面子 個人榮辱.
保羅說我是憑著.
我是透過.
我是通過愛去請求你.
去求你肥利們.
我求你.
你可以像我一樣這樣去做.
我在其他教會聽過一件這樣的事.
但不是這樣算 不是這樣算.
有兩位弟兄.
他們一起去開教會執事會.
在過去一段很長的年日裡面.
領袖弟兄A也得到領袖弟兄B的幫助.
每一次開完會都送他回家 順道.
每次在路上開車的時候.
弟兄A和弟兄B會談剛才開會的事.
說一下教會的事情.
大家繼續分享.
有一次就是這樣開完執事會之後.
回家的途中.
弟兄A說開教會的崇拜.

$^{161}$就和弟兄B說.
現在崇拜這樣這樣.
我覺得應該可以這樣.
於是他就把自己的意見.
這樣這樣這樣.
不斷地和弟兄B說.
在那一次車程的分享之後.
弟兄A不知為何總是覺得.
弟兄B接著就有意無意嘗試避開他.
甚至開完會之後.
總是他去過廁所.
轉頭他就開車走了.
沒有接著他一起走.
弟兄A思前想後.
是不是正因為弟兄B.
上次和他說崇拜的事.
他是有份負責的.
是不是他有點不高興呢.
但是他說大家都是領袖.
大家都是在教會裡面去管理的人.
我當時已經很客氣和他.
去說這件事很中立的.
去分享一些看法.
所以我不是很明白.
為什麼他會這樣氣我.
再過多一段時間.
弟兄A再三反省.
他發現是不是我沒有換位思考.
我沒有在我這位弟兄的角度去看這件事.
當我去分享我的意見.
很中肯的意見的時候.
我作為執事的資格的時候.
我是用權利和地位的出發去反映.
對方B弟兄所收到的是一些誠實話.
B弟兄他反省.
是 我是說了誠實話.
但是我沒有對弟兄有愛心.
保羅為肥利門去感恩.
因為他真的可以展現到愛心.
在有災難的時候.

$^{201}$怎樣去幫助身邊人 社區的人.
但是保羅他繼續去教導肥利門.
或者去求肥利門.
你可以像我這樣.
嘗試超越我們的位份 我們的權利.
用這一份很單純的愛心.
去對待我們身邊的聖徒.
你看到嗎.
教會群體裡面.
我們因著這個大團契的關係.
我們是著重能夠很實在的彼此相愛.
說的話的人自己實踐出來.
保羅去分享.
他也是這樣做.
這個經文對我也有很大的提醒.
回想過去這幾年我們所面對的人際那種.
我形容很intense 很緊張.
很容易擦槍走火.
一句說話 一句的貼文.
可以變成很多爭拗的一年.
超越權利.
用愛心去對待我身邊的人.
這個是主要給我的提醒.
現代人的關係的確是很麻煩.
因為我們常常被這個社會去教導.
我們要用功利 得失 效益 好處去計算.
對的 因為我們是這樣生存的.
我們這樣生活的.
我們去賺得我們每天所需用.
是用這些方式.
所以這是很中立的事情.
但是如果我們所有的人際關係都是這樣的話.
就很可怕了.
例如一個人轉了工.
他自然會繼續和舊同事有一定的聯繫.
不過慢慢就會越來越少.
因為大家不再圍繞同一個目標.
為了同一間公司去賺到一些利益.
以前的關係是會越來越淡的.
你轉了工 你明白的.

$^{241}$還有一些聯絡的.
大概在工作上有需要.
又或者今天留一線.
他日還有生意可以談 還有合作的機會.
不過這些狀況如果進來教會群體.
特別在疫情裡面 我們見少了.
你就會發覺有時候我們仍然是.
跌落這個社會 教導我們.
以功利 得失 效益 好處.
去計算我們和別人的來往.
費利門和歐尼西姆的關係.
起初都是一種功利的關係.
用今天的說法 歐尼西姆是一個工具人.
為什麼呢 因為他的存在.
是幫費利門去做事的奴隸.
幫他做不同的事務.
所以在十一節裡說.
保羅對費利門說.
他從前對你沒有什麼好處.
現在對你對我都有好處.
這個好處 這個致詞.
其實是用回耶穌在馬達芬說過.
五千 二千 一千的比喻裡面.
那個一千的 最後主人說.
他是沒用的僕人.
那個沒用的兩個字.
用就是好處或益處.
在原文是同一個字.
歐尼西姆對費利門沒有好處.
因為他逃跑了.
不再在我身邊工作 為我做事.
所以對我來說是沒用 沒益處 沒好處.
打個比喻.
今天在你公司裡面有一個超級不負責任的同事.
不盡責 不做事.
你作為主管 你可以怎樣做.
見到他回來就是吃早餐.
很快就訂午餐.
然後訂完午餐就想喝下午茶.
喝不喝珍珠奶茶.

$^{281}$當然是等下班.
如果是這樣的話.
你見到一個他都不是很productive.
很能夠為公司的好處著想的人的時候.
你會覺得你不想留他.
當然現在一個這樣的人.
你都要留 因為不夠人.
但回想費利門.
回想歐尼西姆.
他這樣逃跑了.
或者他不再幫自己做事.
費利門是真的很生氣.
當時的律例是可以.
你抓到這個奴隸 你可以處死他.
羅馬的法例可以這樣做.
但保羅這個時候竟然說.
這些都是以前 過去了.
接著還說 以前或者這個階段.
對你沒有用的歐尼西姆.
如今對你對我都有好處.
其實保羅在說什麼.
我們很熟悉這一節的經文.
提摩太后書.
人若自傑 脫離卑賤的事.
就必作貴重的器皿.
成為聖傑 合乎主用.
你們見到的用字.
這個就是好處的字.
同一個字詞.
歐尼西姆的用.
和今天我們一起成為聖傑.
合乎主用的用是同一個字.
你想想你家裡的碗碟.
一隻碗 一隻碟.
作為裝東西的器具.
如果可以用.
基本上都是裝食物 裝湯.
最重要是什麼.
看它是金造 銀造 還是木造.
不是 我想無論怎樣用都好.

$^{321}$最重要是乾淨.
如果你在碗櫃裡拿碗出來.
看到上面還有菜渣.
還搭著蟑螂腳.
就算相金給你幾顆閃鑽.
你都不會用來裝飯.
不合用.
一個器皿要給主人使用.
不一定是新的.
不一定是漂亮的.
不一定是名貴的.
但是這個器皿能夠給主使用.
是要清清潔潔 乾乾淨淨.
洗好之後 抹乾了.
放回碗櫃裡.
隨時主人可以拿來用.
古羅華這個以前沒有好處.
沒有用的歐尼西姆.
現在變成怎樣.
十六節 一起讀.
不再是奴隸 而是高過奴隸.
是親愛的弟兄.
對我固然是這樣.
對你來說.
不論按肉身或在主內的關係.
更是這樣.
按剛才說的羅馬社會的制度.
看目標功利得失.
效益好處的計算.
一個逃走的奴隸沒有用的.
是有害的.
但現在保羅180度大反轉.
告訴肥靈門.
歐尼西姆對你對我都有好處.
不僅如此.
他還將奴隸奴僕.
一下提升到弟兄的地位.
保羅的思想在當時非常超前.
甚至可以說是荒謬的.
當時沒有人會這樣想.

$^{361}$但經文說有好處和沒有好處.
其實是在說什麼.
是提出在我們團契關係的群體.
在教會當中.
我們可以學的就是.
我們待人不求.
這個世界說的好處.
各位 主內的關係的確不是功利的計算.
早幾個星期.
亞運.
不知道你們有沒有機會了解.
香港這次很厲害.
很多金牌.
有水有勁.
很多劍擊有勁.
很多項目都拿到很多獎牌.
再早十幾年的時候.
香港隊參加亞運.
比較突出的項目.
只有兩三樣.
其中一項就是單車.
你都知道我們單車隊很厲害.
在2006年的亞運.
香港單車隊拿到六面金牌.
非常厲害.
不過當時拿到金牌的人.
你會想起哪些人物呢.
當年.
可能你會想.
是不是黃金寶.
這個名字我們基本上香港人都認識.
但在那一年亞運.
不是黃金寶拿到金牌.
而是他的師弟張勁衛.
在比賽之前.
其實港隊的教練已經計劃好了.
接著的個人賽是以黃金寶為搶分的主力.
他的師弟張勁衛.
這個21歲的年輕人.
就是做副手.

$^{401}$負責幫師兄在車隊前面開路擋駕.
你知道單車要贏就要破風.
如果有人在前面幫你破風.
後面的人就踩得輕鬆些.
以至到時他要為搶分的時候.
他的力量可以爆發出來.
不過在比賽過程當中.
一定有很多變數.
令到阿寶由最初去帶開.
這個組合帶開.
到中段阿寶開始發覺自己虛耗過量.
後勁不計.
這個時候他馬上改變想法.
他示意自己的師弟張勁衛.
他爬到他前面.
他跟他說.
我和你換.
你做主力.
自己就運用經驗速度去牽制其他有威脅的對手.
去協助這個師弟衝前.
最終這個張勁衛.
在衝線之前壓倒對手奪得金牌.
原本身為人人都期待他拿金牌的主將.
亦都是他應該去負責搶分拿牌的阿寶.
竟然放棄了原本應該是他的獎牌.
他以身擋敵.
造就師弟一舉擲金.
後來有人訪問阿寶.
你這樣很可惜.
大家都期望你的實力應該是拿金牌.
阿寶毫不介意地說.
以往都是眾師弟為我開路特惠.
將機會給我.
成就我拿獎牌.
今次我只是去協助自己的師弟上前.
證明他們都有奪金的實力.
而且香港單車隊不只是只有我黃金寶.
人家贏等於自己贏.
這是一個團隊的勝利.
因為這不只是計算自己個人得失功利好處.

$^{441}$本來都是這樣想.
而且經文更加說.
現在的團隊是誰是主人.
誰是僕人.
肥理門以往你是主人.
歐連西武是僕人.
但是在十六節的結尾說.
我們都只是在主內的關係.
恩著信同一位主.
大家是親愛的弟兄姊妹.
所以待人不求自己的好處.
這是保羅藉著這一節的經文.
再一次去求肥理門去明白這件事.
他自己也親身示範.
在前幾節他這樣說.
歐連西武對於保羅來說是很寶貴的.
好像契仔一樣.
因為在他坐牢的時候.
歐連西武信了耶穌.
他就留在保羅身邊去服侍他.
你會覺得很奇怪.
為什麼坐牢有其他人進去呢.
是的 在羅馬時代坐牢是沒有食物的.
不會給你飯吃的.
沒有皇家飯的.
你要找你外面認識的人帶食物給你.
去照顧你.
你有病你要外面找個人來.
幫你看醫生,療傷.
所有東西都是由外面供應.
甚至有時候監獄可以用一些你認識的人留下來.
幫你,陪你一些東西.
歐連西武就留在保羅身邊.
幫他這個上了年紀的保羅.
身體不好的保羅.
但是他說我為了大家的好處.
我很想保留這個像契仔一樣的歐連西武.
但是保羅知道我不會求自己.
求自己的好處.
真正對歐連西武對肥你們有好處有用的.

$^{481}$不是你們抹走過去不了了之.
算了算了我寫封信你們就這樣.
遙距原諒他就算了.
而是真正對歐連西武對肥你們有好處的.
就是他們可以面對面一起去看這件事.
最終在主裡面有復和,開展.
因著同信一位主.
而有這個團契.
彼此相愛的關係.
領事館這個道理是很簡單.
我們常常聽的.
但應用在我們自己教會生活裡面.
倘若在過去這幾年我們有大環境,小處境.
所引致各樣各樣的相處碰撞就很不容易了.
保羅示範了什麼是不求自己的益處.
是的.
他說代人就是不求自己的益處.
願意我們都一起去學習.
我們這個群體,我們這個家.
甚至我們這個宣道會.
誰是主場.
是你還是他呢.
是牧師傳道還是執事領袖.
是導師還是團友.
是這個地點的人還是那個地點的人.
都不是我們所信的主.
是我們所信的主耶穌基督是主場.
所以在這樣的旗號下.
我們的確可以學待人不求自己的益處.
繼續我們一起看17節.
我請大家一起讀這節.
所以你要是把我看作同伴.
就接納他好像接納我一樣.
這個同伴在希臘文是有一個common.
或者公共,共同的意思.
是用這個公共,共同這個字變化出來的.
所以團契相交原來都是和這個同伴這個字詞.
那個公共那個字是相同的來源.
這樣的話帶給我們這個字詞起碼有三重的意思.
第一 同伴 中文已經告訴我們了.

$^{521}$是companionship 是一同,是共同,是一起.
你都明白了 總之是一齊.
就是這個意思 都比較普通的.
未講到那種很緊密的 原文那種很緊密的關係.
所以第二種在中間那裡就是一種夥伴partnership.
有些事今天我和你拍檔.
今天我和你一起做這個事奉.
那你明白了 一起拍著相.
有多一點默契的.
在這個空間時間上我們聚在一起.
我們有相同的追求 相同的目標.
那拍檔是什麼呢 當然我們明白.
保羅上上提到的夥伴是在福音上一齊努力的.
就好像今年的大家的念題 全陽福音.
很明顯的 這個很清晰的目標.
一起去做這件事.
而第三種這個原文同伴這個字還有一個.
有福同享 有難同當的意味.
所以就在最上面那個字詞.
也是我們今天最要掌握的一個友伴.
一個friendship.
不單是companionship.
不單是partnership 還是friendship.
可以彼此分擔 彼此分享.
保羅和他的同工其實就是這樣.
一齊的 有一個目標.
一起全陽基督同時在他們的生活裡.
有很多的結連 很多的千絲萬縷的相處的時刻.
一起去面對不同的事情 彼此去扶持 傾訴心事.
這個是友伴.
我們可以想一想.
疫情讓我們去跟別人很疏離.
今天錦宣家可以一起聚首一堂.
連同孔恩堂 連同嘉欣樓的弟兄姊妹.
連同不同語言的群體 我們有菲律賓的姊妹 一齊.
今天在這裡我們聚在一起.
在這麼大的群體裡 有沒有你的同伴呢.
有啊 一起敬拜嘛 是不是.
但是否只是停留於回來教會說聲Hi.
離開教會說聲Bye.

$^{561}$有時候真的是這樣.
我有的 我有一起事奉的.
今天我們一起唱詩班.
一起在後面做影音 做招待 做預備食物.
我們一起是伙伴 在事工上是伙伴.
大家是不是友伴呢.
能不能說得心事呢.
當然我們不可能跟所有人是友伴.
但是你身旁有沒有這樣的弟兄姊妹.
你可以去不單止事奉 還有生命的交流.
還能夠做真正的對工朋友 友伴呢.
我們怎樣可以找到這樣的對工呢.
我記得有一年自己在團契選職員.
當時負責選舉的弟兄姊妹就公佈.
未來這一年就有這個這個.
我在其中一個名單.
我看名單一出 哇.
未來一年怎樣事奉啊.
他都想選職員.
他這麼古怪.
這個還經常吹水 怎樣事奉啊.
那個 剛剛大家面左左.
經文都提醒我們.
你都讓我去理解到一件事.
有時候一切的友伴是伙伴 無法選擇.
是神放我們在這裡的.
保羅跟腓利門是團友 是友伴.
如果是這樣的話 請你也接納歐尼西姆.
因為歐尼西姆都是我的友伴.
也可以是你的友伴.
而17字的下半節說請你接納他.
好像接納我一樣.
這個團契關係是主安排 不是隨自己喜好.
保羅說那你接納他吧.
這個接納其實是一個很有動感的字詞.
在紹行傳28章形容保羅一次船難之後.
他們沉船 游到岸邊.
有一句經文是這樣說.
當時在島上有一班土人.
看到這班船難游上來的人.

$^{601}$就說看待他們有非常的情分.
當時下雨天氣又冷 就生火接待我們種人.
你可以想像一下這種情景.
我們幻想一下.
沉船很辛苦游啊游啊游到岸邊.
抱著船的舷木 終於去到岸邊 沒氣了.
在船難之前 其實他們沒有吃到什麼東西.
一直被風浪打擊.
上岸之後 一直趴在海灘上的時候.
前面有一班園居民在生火.
氣 火 氣 火 吃人族.
他們很熱情地拉著這些在海灘上趴著的人.
拉他們去火邊.
不是燒他們.
拉他們去火邊讓他們取暖.
所以接納這個字詞很有動感.
不是只有一張嘴.
是歡迎你來這裡.
這是最基本的.
但不止這樣.
不只是站在這裡.
接納不只是提一個歡迎.
而是延伸一個歡迎.
伸出你的手 拉他們來.
向大家這個群體中心走多一步.
你是主動做這件事的.
其實中文字的接納納字都很精彩.
喬斯邊加一個內字.
我們可以想像.
很全神地告訴你.
你用一條帶章的那個人牽引你.
你這個內圈 接納.
接納一個圈外的人.
你 安書區.
你 朋友圈以外的其他人.
進來裡面.
可能你都會像我當年選職員的時候有抗議.
要我接納一個這樣的人不容易.
一個身家清白好人好者都沒有難度.
但要接納一個和我面左左.

$^{641}$和我有很多爭執的人.
還要我做主動.
他不是很便宜嗎.
明明他得罪我的.
接納這個字在另一處也用過.
在羅馬書十五章七節.
那裡說.
所以你們要彼此接納.
如同基督接納你們一樣.
基督要我們想起.
這個接納不單是我們用我們的能力.
其實我們很多時候是沒有能力甚至不想.
我們的心不舒服.
所以我們接納的動力或者起點.
不是在我們能否解決那些事件.
而是更重要的是.
我們是先被主接納了.
先是由耶穌來到這個世上.
將我們拉入他的國裡面.
神extend a welcome.
將我們牽扯到他的愛當中.
如果是這樣的話.
我們是信這位主.
我們是可以有能力靠著主變得慷慨.
可以接納我們從前要擋在我們圈外的弟兄姊妹.
剛才我說那個團體職員選舉的事.
好了,選了,真的選了那些人.
一年之後.
我的看法完全改變了.
我真的很感受到.
每一個人是主擺在這個團隊當中.
每一個人都是屬於主.
也可以是我的友伴.
在那一年裡.
對我最多意見的團友.
我覺得他經常不聽我的說話.
經常幫我包頸的人.
當我學,我只是停下來聽多一點.
我發覺那個關係不同了.
甚至在事奉以外.

$^{681}$我們可以一起去街上看電影.
一起去打球.
一起玩.
不只那間由夥伴變成友伴.
神就使用他.
我發覺在一年之後.
我覺得上帝是派他來收止我的生命.
今天我們仍然是很好的團友.
友伴.
說回頭.
那些我以為相處不到的弟兄或者姊妹.
他們何嘗不是多走一步.
延長歡迎.
將我接納在他們的內圈裡.
弟兄姊妹沒有主耶穌基督慷慨的接納行動.
我們也沒有這樣的.
但我們是信主的群體.
我們高舉耶穌基督的十字架.
我們這個信主的群體.
是能夠靠著主耶穌基督的愛.
去給予我們的接納.
我們也是這樣去接納一些.
跟我們很不同的人.
去接納在我們安書圈以外.
主愛親愛的弟兄姊妹.
這個友伴的功課.
在服償的階段.
可以是我們最新的目標.
放下過去疫情裡.
放下過去社會很多的事情.
是我們還要去面對.
但是我們的起點.
是因為主接納了我們.
我們信主.
我們就能夠去接納身旁不同的人.
我們不需要靠自己.
是靠主.
在主裡.
你做得到.
我做得到.

$^{721}$保羅知道肥里門一定會聽他的意見.
甚至超出在信裡的要求.
所以弟兄姊妹.
今天是錦宣家四十二週年的堂慶.
十二年是一個很長的歷史.
有歷史就有故事.
有故事就有美麗的故事.
有艱難的故事.
但是無論如何.
我們在這一點可以慶祝.
我們依然在主這個群體裡.
我們願意走在一起.
我們就可以成為友伴.
展望前行.
疫情過後社會的變幻.
大家靠主能夠維繫這個關係.
無論如何.
經文鼓勵我們.
我們學習超越自己的權利地位.
我們用愛心去對我們家裡的弟兄姊妹.
我們去待人.
不要像社會那樣去計算功利.
我們學習不求自己的好處.
更加我們學習.
我們因為同有一位主.
同信一位主.
我們領受的慷慨.
我們都是這樣去對我們身邊.
每一個親愛的弟兄姊妹.
求主幫助我們恩著信.
不斷去建立這個真正大團契的家.
最後很快說一說.
你想不想知道肥利門他們最後怎樣.
《聖經》沒有寫.
歐尼西姆他們怎樣.
教會歷史是這樣說的.
後來肥利門.
看到歐尼西姆親自帶著這封保羅的書信來.
他看了之後.
他原諒了歐尼西姆.

$^{761}$接納了他.
拉回到他家裡.
這個教會裡.
還給他自由.
不再看他為奴隸.
看他為弟兄.
後來歐尼西姆在主後一百年左右.
成為了以弗所的主教.
這個奴隸成為了帶領教會的主教.
他還負責收集保羅的書信.
所以為什麼這封信.
這麼個人的信.
匯入了保羅的十三封書信.
也在我們今天新約的聖經裡.
你看看.
因著真正友伴的關係.
這個教會.
這個團體.
這個對公.
這兩個很不同的人.
很有過去歷史的人.
他們成為友伴.
一切前行.
為主作了美好的事.
全陽福音.
同樣我們可以因著主的教導.
一切得著更新.
我們祈禱.
主求你大大賜福給金孫家.
當我們每一個人願意.
在你的十字架底下.
成為一群.
能夠彼此相愛.
互相視為友伴.
這個群體就能夠被你帶領前行.
願主繼續使這個地方.
去成為作炎作光.
一個美好的見證.
多謝你.
我們祈禱.

$^{801}$在奉主耶穌基督的聖名之祈求.
阿們.
\newpage



\section{撒迦利亞書}
\label{sec:j9CuOv5EhDA}
\textbf{2023-11-02 周四晚查經班《歸回先知書研讀(7)撒迦利亞書》 講員:曾裔貴牧師}
\newline
\newline
連結: \href{https://youtube.com/watch?v=j9CuOv5EhDA}{\texttt{https://youtube.com/watch?v=j9CuOv5EhDA}} ~~~~ 語音日期: 2023-11-03
\newline
\newline
\hyperref[sec:ROGX9cAg8Cg]{\small{< < < PREV SERMON < < <}}
~
\hyperref[sec:index_chronic]{\small{[返順時目]}}
~
\hyperref[sec:index_scriptual]{\small{[返順卷目]}}
~
\hyperref[sec:URxw02niNWY]{\small{> > > NEXT SERMON > > >}}
\newline
\newline
$^{1}$我們停了一個星期.
我們恢復撒迦利亞書的研讀.
和大家一起祈禱.
多謝天父再一次讓我們能夠在晚上.
一同學習聖經的真理.
多謝您給我們有機會仔細去看.
撒迦利亞書這本歸回的歷史書.
和先知書的系列.
幫助我們能夠明白.
尤其是對於將來的預言.
聖經真是安定在天.
因為是神的默示.
使我們能夠了解以色列神的選民.
將來在神的手中的光景.
那個的福樂,那個的復興.
因為我們知道不是靠人的.
真的不是人可以速食這件事.
願神幫助我們.
尤其是對於彌賽亞的預言.
那個準確性是令我們吃驚的.
幫助我們能夠從其中明白.
我們仰望等候在您的面前.
禱告奉耶穌基督的名.
阿們.
我們回到第七,第八章.
有少許上一個星期.
不夠時間講得詳盡的.
我們接著就會進入第九章到十四章.
真是一個很遙遠的預言.
差不多你讀撒迦利亞書.
你有那個概念一直讀.
就算你不是太.
在新約裡面去了解.
其實已經浮現了.
彌賽亞的那個.
模糊的片段,樣貌.
你就可以甚至乎.
如果我們是新約.
回頭讀舊約.
其實有主耶穌祂的出現.

$^{41}$成就舊恩.
我們就可以對號入座.
用新約的那個世界的眼光.
來去看舊約的預言.
正是彌賽亞就是耶穌基督.
很奇妙的.
裡面的那些預言.
你想都沒有想到.
福音書是直接這樣引用撒迦利亞書的.
所以我們.
其實我們很熟悉那些經文.
但是當我們看回那個原典的.
撒迦利亞書的時候.
其實我們就要問一個問題.
當然我們要完整.
整個九至十四章這些預言.
我們就要問這個問題.
耶穌基督在祂地上的成長.
祂看到撒迦利亞書.
祂有什麼感受呢.
我們帶著這個心情再看舊約的時候.
可能會有很不同的一番滋味.
仍然都很慨嘆就是.
今天不只是香港.
很多的基督徒.
舊約根本連揭都不揭.
怎樣揭都不懂揭.
不要說不揭.
這個是一個很大的悲哀.
基督徒是不讀經.
來到去.
自己有機會去研讀.
我們回來看回第七章.
就是有一個訪問團去尋求釋法.
我們在被擄的時候.
一直都在五月,六月來到齋戒禁食.
現在有好日子過了.
耶路撒冷重建聖殿了.
現在也有一班人被擄歸回的.
五萬多人回來了.

$^{81}$現在應該怎樣呢.
問這個問題的當然是伯特利人.
第七章的第二節.
那時伯特利人已經打發沙利息和利堅米勒.
並跟從他們的人.
去懇求耶和華的恩.
就是在這樣的情況下.
就有一個小小的訪問團.
似乎不是秘密地去訪問.
而是很想去尋求釋法.
到底我們應該怎樣呢.
答案就在八章的十六節.
我們看過了.
但是因為十六,十七節.
有些比較重要的經文.
跟大家再簡單地回顧一下.
然後我們就會進入第九章.
我們可以由第十四節開始看.
就是那個問題.
我們應不應該繼續齋戒禁食.
我們應該怎樣做才得到神的恩.
因為他們是要懇求耶和華的恩典.
現在很特別在這段聖經提了兩次.
將來這個民族.
是會有很多人拉著猶太人的衫尾.
我們要透過你們去得到神的恩典.
如果我們現在在這個以巴的戰爭底下.
你是猶太人你可能會很害怕.
拉著你的衫尾.
會不會是自殺式的炸彈.
會不會是一些恐怖襲擊.
但是撒迦利亞書是充滿了平安的寓言.
對這個民族的將來美好的願景.
是很淋漓盡致地說的.
其實現在是大理烏王第二年的九月.
撒迦利亞先知得到寓言.
在這個時候其實是當年的六月.
兩位先知有神的密使.
就來勉勵當時的領袖省長.
就是所羅巴伯.

$^{121}$大祭司就是約書亞.
來勉勵他們興功重建.
破壞了美的聖殿.
他們正式現在開工了.
在這個時候九月.
撒迦利亞就憤外得道.
首先就是一到第六章.
有八組的一些異象.
我們都說完了.
到第七章.
回到這裡就有些人事的問題.
就是有些可能被擄歸回上來的人.
他們就很想問.
到底我們在當年.
這麼多年在被擄的時候.
在巴比倫.
在波斯.
我們一直做的禁食的守節的事情.
現在應該怎麼做呢.
怎麼做才會更加合成心意呢.
這裡有解答了.
我想跟大家再將這一段聖經.
認真地再讀一次.
八章的十四節.
萬君之耶和華如此說.
你們列祖惹我發怒的時候.
我怎樣定義降和並不後悔.
現在我照樣定義.
施恩與耶路撒冷和由大家.
你們不要懼怕.
你們所當行的是怎樣呢.
正正就是第三節他們所問的.
我們要去尋求.
去懇求耶和華的恩.
我們應該怎麼做.
現在神就說了.
你們所當行的就是這樣.
所以十六,十七節是金句.
各人與鄰舍說話誠實.
在城門口按治理判斷.

$^{161}$使人和睦.
城門口就是說當時他們是.
德高望重.
尋求民間訴訟的一個地方.
那裡一定是城門口有一些受人尊敬的長老.
來到去.
大家都信服的.
大家都信任的.
這樣來到去幫忙.
做一個仲裁.
重點都不是誰對誰錯.
公說公有理.
而是在那件事情上.
尋求最好的和解.
最好的大家的和睦共處.
所以需要有一個這樣的.
一些德高望重.
大家都信任的長老.
長者.
來到去民間的領袖.
來到去在城門口.
所以如果你讀路得記.
有沒有留意.
波亞斯就去到城門口.
去尋求那些長老.
當時的父老.
來到去怎樣解說.
到底路德是不是可以嫁給我.
波亞斯和那個人.
那個人是最近的親屬.
後來就說我不要.
我放棄這個權利.
波亞斯就可以娶路德.
是要有審裁的.
是要有長老判斷的.
有一個儀式.
一個拋鞋的儀式.
所以這個就是城門口.
重點就是按自公自理判斷.
重點是洗人和睦.

$^{201}$所以教會很需要有牧者.
來到去幫助弟兄姊妹.
來到去排除一些的誤會.
一些的糾紛.
一些的衝突.
一些的嚴規.
牧者是要扮演這個角色.
如果一個教會能夠對牧者有信任和尊重.
很多問題就可以解脫.
化繁為簡.
不會將小事變大事.
不會暗中又拉攏不同的勢力.
這樣就沒有了.
教會就是.
可能一件事整個教會就反轉了.
是很可惜.
是很浪費大家的心力.
教會是應該集中一起.
和睦地去傳揚福音.
集中力量.
對外.
對人的心靈.
將它帶進去救恩裡面.
所以這裡是再一次提醒我們.
舊約新約都是這樣.
告訴我們.
一定要有一個和睦共處.
誰都不可心裡謀害淪陷.
也不可喜愛喜假釋.
因為這些事都為我所恨惡.
這是耶和華說的.
講得很清楚了.
公正和追求公義.
追求和睦.
追求與眾人和睦.
若是可行.
《羅馬書·保羅》十二章說.
總要盡力與眾人和睦.
你可以想像得到.
這樣投射我們的做物主.

$^{241}$祂就是真理.
公義.
聖潔.
和信實.
和滿有慈愛的神.
如果我們人一生.
我們可以在教會裡面學習.
在聖靈的幫助下學習.
我們的生命.
我們的屬靈.
又會一直成長.
否則.
我們只是口號式地掛在嘴唇邊.
而我們自己一生都不去實踐實行.
救恩只會遠離我們.
我們不會改變到我們自己的生命.
所以這是很重要的.
雖然是兩節短短的聖經.
已經講述了.
神不是那麼緊張你的儀式.
第三節就是說.
我們應不應該繼續做這些事.
神就說.
那些不是不重要.
但是相對於人的內心的品性.
那些儀式就真的不是那麼重要了.
十八節.
萬君之耶和華的話臨到我說.
萬君之耶和華如此說.
四月五月禁食的日子.
七月十月禁食的日子.
必定會變為由大家歡喜快樂的日子.
和歡樂的節期.
所以你們要喜愛誠實與和平.
所以是相輔相成的.
是兩條腿的.
你很敬虔的去守那些節期.
是因為你內心本來就有誠實.
是有和平.
那你就去做這些事情.

$^{281}$就相得益彰了.
不然就變成了宗教的虛偽.
所以你會看到主耶穌.
為什麼那麼生氣那些法利賽人呢.
就是太虛偽了.
他們做到嘩+01.
看到耶穌的門徒.
在安息日.
去摘那些墨水來吃.
因為肚子餓.
就指著耶穌說.
你看看你的門徒.
安息日做了一些不可以做的事情.
太過矯枉過正了.
完全沒有了一個憐憫的心靈.
什麼都要齊齊齊.
什麼都要準確.
什麼都不可以有錯.
一絲不苟.
但實際上.
對人就有一個很大的排擠.
那就變成沒有了那種通容.
沒有了那種彈性.
沒有了那種最重要的.
人和人之間的相處.
當然不是叫你包庇.
當然不是叫你隻眼開隻眼閉.
那種所謂的通容.
是真正的.
是看到一些情況的需要.
來做一個彈性的處理.
所以這裡告訴我們.
那些齋戒.
都可以因為那個人.
而變成是由大家歡樂的日子.
齋戒本來就應該是.
嚴肅.
禁食.
悲哀.
反省.

$^{321}$但是.
神可以令到這些事情.
變成為一個歡樂.
所以這裡又再一次告訴我們.
猶太人將來他們要面對.
有一個徹底的改變.
就是來塞亞降臨的時間.
但是朝著這一路的發展.
我們就看不到人可以.
這樣來改變整個局面.
也不會看到他們.
又得到眾民的歡心.
因為第20到23節.
就講到將來的猶太人.
很矜貴.
因為每個人要拉著他的衣服尾.
你要讓我去見神.
我們要去求神的恩.
去懇求神的恩.
出現了兩次.
第20節.
萬君之耶和華如此說.
將來必有列國的人.
和多城的居民來到.
這城的居民必到那城說.
我們要快去懇求耶和華的恩.
外邦人.
列國.
再不是上文講的伯特利的人.
來懇求神的恩.
現在是外邦萬國的居民.
就是說全世界將來.
焦點都在耶路撒冷.
都在以色列.
在神的選民.
原因是什麼呢.
彌賽亞回來了.
絕對不是今天這樣.
人見人憎.
第21節的中間.

$^{361}$尋求萬君之耶和華.
我也要去.
必有列邦的人和強國的民.
來到耶路撒冷.
尋求萬君之耶和華.
懇求耶和華的恩.
兩次出現了.
巴不得我能夠成為你們的同胞.
進入你們的國籍.
那個當然是地上.
但現在說的是屬靈的那個Blessing.
是臨離盡至這樣將來會出現.
所以這些是一個很美麗的願景.
怎樣可以讓這件事發生呢.
我們真的沒辦法看到.
不能夠.
我們在我們腦海裡面自己幻想.
但是新約耶穌已經說了.
祂一定會再回來.
會有號角聲.
有天使掌的聲音等等.
帖撒羅尼加前書四章就已經說了.
我們應該是13號還是14號.
我們就會在帖撒羅尼加那個地方.
那個舊城.
當然現在是一個很現代的城市.
但是我們會去舊城看帖撒羅尼加.
就會看回保羅的那個宣告.
這就告訴我們.
怎樣可能會發生呢.
但是神可以令到這個日子一定會來臨.
所以這個是一個將來的預言.
沒錯現在是在重建那個聖殿.
剛剛重新開工.
因為停了十幾年了.
現在還有大概四年左右.
他們做好了.
但是這裡說的重建.
是一個短暫的應驗.
但是真正永恆的應驗.

$^{401}$是要來塞亞的降臨.
才會萬國.
不是攻打耶路撒冷.
是去懇求耶和華的恩.
所以如果你是當時的猶太人.
你就會感受到那種榮耀.
那種幸福.
第23節.
萬君子耶和華如此說.
在那些日子.
必有十個人從列國諸族中出來.
拉著一個猶太人的衣襟.
蚊衫尾.
說我們要與你們同去.
因為我們聽見神與你們同在.
原因就是這樣.
所以我們要切心期望彌賽亞降臨.
對整個以色列的改變.
對整個世界的徹底的扭轉.
不要再身臨淘汰.
好像主耶穌降臨之前.
是會有這些悲哀.
是會有這些戰爭.
但是與我們基督徒也有關.
我們都要看到這個日子來臨.
所以我們是要祈禱.
彌賽亞快點降臨.
所以我們不可以去估計.
那個日子那個時辰.
但是我們是要有一個懇求主的在來.
改變這個世界的各方面的命運.
因為實在是太過令到人與人之間.
處於一種很繃緊的戰爭狀態.
就是越來越多國家有核彈.
你怎麼辦呢.
一個都已經搞不定了.
說幾十支幾百支這樣一起射的時候.
整個世界就滅亡了.
你可以想像得到.
人類製造那種緊張.

$^{441}$來自己嚇自己一代一代過去.
現在就實際在我們面前見到的.
人道的災難.
在巴沙.
在這些地方.
你就可以想像得到.
最無辜受害的都是平民百姓.
我們完全見不到以色列是受歡迎.
拉著它的衣cum .
我們很想去尋求你的神.
現在哪裡有呢.
但是我們一定要相信聖經的華語.
那個應驗.
那個權威.
因為是神自己默示.
將來一定會按字面來應驗的.
好了對九章.
這裡去到十四章.
都是一些對將來的局面的預言.
都提過了.
真的可以對號入座.
在大家的腦海裡面.
真的不會有錯的.
一直看著彌賽亞的模糊的身影.
出現在這個預言裡面.
我們從新約回頭讀舊約.
二千年過了.
耶穌基督如果離開這個世界.
回到天家.
祂會再回來.
我們用這個概念去看舊約.
就不同舊約當時的人.
第一手的撒迦利亞得到的預言.
很模糊看到有一個彌賽亞.
但是是誰來的呢?不知道.
但是我們新約就肯定對號入座.
就是我們的主耶穌.
第九章.
耶和華的默示應驗在哈德拉地大馬色.
世人和以色列各支派的眼目.

$^{481}$都仰望耶和華.
和靠近的哈瑪推羅西頓.
因為這以成的人大有智慧.
有趣的地方就是.
就是現在以巴的地方.
加沙地大.
就是推羅西頓這一帶.
推羅為自己收足保障.
因為他們是那些肥尼基人.
航海.
積蓄銀子如塵沙.
堆起精金如街上的泥土.
主必趕出他.
打敗他海上的權力.
他必被火消滅.
亞實基倫看見必懼怕.
就是菲律士一帶地方.
現在還是那些名字.
加沙.
就是加沙地大.
加沙看見甚痛苦.
以及倫因失了盼望蒙羞.
加沙必不再有君王.
亞實基倫也不再有居民.
師生子必住在亞實特.
這裡是菲律士的五個城.
就是那一帶地方.
就是沿海推羅西頓的那一帶地方.
敘利亞下面西面近地中海的.
就是推羅西頓和菲律士沿海的地方.
「我必除滅菲律士人的驕傲」.
「我必除去他口中大血之肉」.
「和牙齒內可爭之物」.
就是一直說那些典故.
外邦人是拜偶像.
吃那些是神厭惡的食物.
猶太人是不敢不提的食物.
他們就可以隨便來當為可以吃.
但在神的眼中是可憎的食物.
接著就是第七節.

$^{521}$「他必作為如成的人歸與我們的神」.
「必在猶大象族長」.
「以甲鄰人必如耶布斯人」.
「我必在我家的四圍安營」.
「使到敵軍不得任意往來」.
「暴虐的人也不再經過」.
「因為我親眼看顧我的家」.
「錫安的民應當大大歡樂」.
「耶路撒冷的民應當歡呼」.
「漢啊,你的王來到你這裡」.
「他是公義的,並且施行拯救」.
「謙謙和和的騎著路,就是騎著路的驅使」.
福音書出現了兩次.
很快和大家揭一揭.
萬太的二十一章.
耶穌騎路進城.
萬太福音是在撒迦利亞書之後.
所以不用揭太多.
第二十一章.
「耶穌和門徒將近耶路撒冷」.
「到了伯伐旗,在橄欖山那裡」.
「耶穌就打發兩個門徒,對他們說」.
「你們往對面村子裡去」.
「必看見一匹驢山在那裡」.
「還有驢驢仔同在一處」.
「你們解開,牽到我這裡來」.
「如果有人對你們說什麼,你們就說」.
「主要用他,乃人必立時讓你們牽來」.
「這是成就」.
留意一下.
「是要應嚴先知的話說」.
「要對錫安的居民說」.
「看啊,你的王來到你這裡」.
「是溫柔的又騎著路」.
「就是騎著路驅使」.
我們可以接下第12章的約翰福音.
這就是耶穌基督的.
受難州的騎驢進城.
很有名.
有詩歌都將這件事.

$^{561}$成為一首相當動聽的詩歌.
《王騎驢進城》.
第12章的約翰福音.
第12節.
第12章的12節.
第二天有許多人.
這裡說第二天.
即是上文.
我們曾經說過約翰福音了.
拉薩路一家人擺酒.
慶祝拉薩路復活.
耶穌也坐直.
一個由死復活.
從墳墓出來的拉薩路.
本來病死了.
那些人生氣到一個地步.
不單止要除滅耶穌.
甚至要將拉薩路再殺一次.
殺了他.
你可以想像得到.
當時的人已經是多麼的硬心.
接著第12節就說.
第二天有許多人上來過節.
聽見耶穌將到耶路撒冷.
拿著鐘樹枝出去迎接他.
即是說多了少許補充.
在這個橄欖山沿著.
走過去對面的耶路撒冷.
當然經過吸輪溪.
一個小的山丘.
那些人就識貨.
知道這樣騎驢.
不是普通的.
就是王.
正正就是撒迦利亞書的寓言.
王騎著驢驅進城.
所以很特別的地方就是.
那些人真的有準備.
就拿著鐘樹枝出去迎接他.
就是耶穌.

$^{601}$喊著說.
和珊娜奉主名來的.
以色列王是應當稱頌的.
耶穌得了一個驢驢就騎上.
如經上所記的說.
「錫安的民,不要懼怕.
你的王騎著驢驢來了.
這些事門徒豈先不明白.
等到耶穌得了榮耀以後.
才想起這話是指著他寫的.
並且眾人果然向他這樣行了」.
有沒有留意到作者再加上一些事後的孔明.
事後的補述.
原來門徒當時都沒有這樣對號入座.
「這麼巧而已.
騎驢仔個個都有.
大把人騎.
這麼巧的時間.
這麼巧有隻驢仔.
那就騎吧.
不要想多了」.
但後來想著想著.
不是的.
不對勁.
這真是應驗撒迦利亞書的預言.
所以我們要回來看撒迦利亞書.
回到原典的出處.
就是《福音書》的作者.
很清楚地告訴我們.
應驗不是哪個王.
耶穌就是那一位舊約預言的君王.
也是他們等待的彌賽亞.
耶穌真的照著.
其實為什麼可以這麼巧.
又叫那些門徒去找那個地方.
找到那個人.
就問他拿隻驢仔給耶穌騎入城.
就完全應驗撒迦利亞書九章.
是不是主耶穌見到這些預言.
他都要這樣來跟隨去應驗呢.

$^{641}$耶穌的內心.
小時候讀到撒迦利亞書的時候.
因為這些經典.
他們已經有了.
在會堂他們可以聽拉比教導的.
舊約就一定有了.
新約就肯定沒有了.
所以每一個猶太人.
自小就在會堂學聖經.
家裡就一定不會有.
因為很珍貴.
所以他們是要去到會堂去聽.
讀到這些預言的時候.
耶穌的內心又如何呢.
知道自己是來撒亞.
這樣的心情又如何呢.
我們可以看回經文的內容.
不過第十節就很特別.
福音書的作者是沒有引用的.
所以第九第十節是一段來的.
是一個段落來的.
第十節就是.
我必除滅以法連的戰車.
和耶路撒冷的戰車.
征戰的弓也必除滅.
他必向列國講和平.
他的權柄必從這海管到那海.
從大河管到地極.
這一位和平之君.
這一位騎著驢驅.
正常就應該騎著戰馬.
所向披斐的君王.
才是世界的人認為權力眾望所歸.
認為這樣才是勝利.
但是主耶穌是以和平進入耶路撒冷.
最終我們知道.
他是要在十誡上犧牲.
完全手無寸鐵.
奄奄一息的一個人質.
但是原來第十節.

$^{681}$是這樣的格局.
這樣的君王.
是這樣的一個治理.
他的統合權管治權.
是去到列國.
去到縱海島.
去到地極.
所以我們如果不讀聖經.
我們就不知道原來有兩個階段.
原來福音書的作者都很醒目.
他們不忍這一節.
忍了就一口笑.
你講的你的彌賽亞.
你讓我看看.
你跳下來.
我們就信你了.
那些當時的群眾.
耶穌就是跳不了.
其實只是幾眼釘.
用力就可以跳下來.
但是不行.
他就是不跳.
不跳的力量才大.
不跳含冤離世才是大的力量.
中國人都說死剩一口.
臨死都咒下面的人.
他是避害不說威嚇的話.
完全默默的接受自己是一個死囚.
要擔待世人的罪孽.
世人的激動他.
他不為所動.
他是有那個力量.
但是如果他一用那個力量.
他就不可以拯救世人了.
他真的要這樣來到去.
在眾人的面前.
其實好像甘心情願認低威.
你想想.
那個你在你身上.
你認低威.

$^{721}$其實很難.
更何況我們的主耶穌.
所以這個是一個很大的力量.
這個力量.
是可以帶來整個宇宙改變.
神的救恩.
藉著耶穌可以傳到世界每一個人.
只要相信耶穌就可以得到了.
但是原來我們的耶穌.
第十節才是那個見真章.
原來真正的彌賽亞.
是要將來掌管萬國萬民.
所以福音書的作者不引述.
不是他們不懂事.
是未應驗.
是知道是有兩個階段的再來.
所以你可以想像得到.
今天我們的基督教就被人嘲笑.
或者我們根本都不是很緊張.
我們自己所信的.
但是在神的預言裡面.
彌賽亞是先受苦後得榮耀.
所以真正的榮耀是在將來主的再來.
第二次的再來.
所以第十節說.
我必除滅以法蓮的戰車.
和耶路撒冷的戰馬.
徵戰的弓也必除滅.
再無戰爭.
真正的管治是和平.
在愛和平裡面去管治.
彌賽亞真的不是浪得虛名.
他的權柄必從這海管到那海.
從大河管到地極.
第十一節.
由新的段落.
錫安娜.
我因與你納若的血.
將你中間被擄而囚的人.
從無水的坑中釋放出來.

$^{761}$你們被囚而有指望的人.
都要轉回保障.
我今日說明.
我必加倍賜福給你們.
我乃猶大作賞賢的功.
我乃以法蓮為將功的戰.
兩個大的支派.
即是等於代表著.
北面的以色列和南面的猶大.
即是整個以色列的意思.
神體是神的器具.
戰爭的器具.
即是錫安娜我要激發你的眾子.
攻擊希臘的眾子.
使你如勇士的刀.
耶和華必顯現在他們耳上.
他的箭必射出像閃電.
主耶和華必吹割.
乘南方的旋風而行.
萬軍之耶和華必保護他們.
他們必吞滅仇敵.
踐踏彈石.
他們必喝血吶喊.
猶如飲酒.
他們必將盛滿血的碗.
猶像彈的四角滿了血.
當那日.
耶和華他們的神.
必看他的民如群羊.
拯救他們.
因為他們必將冠冕上的寶石.
高舉在他的地上.
他的恩慈何等大.
他的榮美何其盛.
五谷健壯少男.
伸酒培養處女.
好像是有一場的戰爭.
這個戰爭是會帶來列國受審判.
但是那個工具.
就是神會興起他自己的選民.

$^{801}$就是神的上了賢的公賤的賤.
那個賤上了賢的.
就是以法蓮就是張公的賤.
猶大就是那個賢.
那個公.
以法蓮就是賤.
就是說神會興起他自己的選民.
來對付列國.
將來這一場大戰.
以色列會得勝.
描述那個情景.
這個古代.
因為是舊約的年代.
所以他們都是用弓箭,馬.
殺戮敵人.
來復興之前的一場戰爭.
是不是啟示佬所說的哈米吉多頓大戰.
很難對號入座.
但是無論如何.
神要再在他自己的選民當中.
來作供.
使得他們能夠.
特別那個描述就是因為神保護他們.
所以他們能夠吞滅仇敵.
他們能夠像喝酒一樣喝敵人的血.
盛滿了血的碗.
又像彈四角滿了血.
就是好像電掣一樣.
來殺戮敵人.
所以到了將來.
當那日出現很多次.
在這五章聖經裡面.
當那日.
就是說到底什麼時候呢.
我們不要估計.
因為沒有人可以估計.
是不是公元三千年呢.
還是公元四千年呢.
沒有人可以估計的.
至少我們不在這個世界了.

$^{841}$除非主耶穌這個時代回來.
我們都不知道.
我們有機會.
但是我們不能夠猜度.
不可以去計算.
但是當那日必然會出現.
是撒迦利亞書.
在第九到十四章.
是出現很多次的.
所以大家回去有機會.
你數一數出現多少次.
現在這裡就已經說了.
那個日子.
猶大會得勝.
以色列會復興.
整個地上.
神在看他的子民.
是他的群羊.
牧養他自己的子民.
拯救他們.
再加上就是.
這樣的描述.
冠冕上面的寶石.
冠冕已經很榮耀了.
冠冕上面的寶石就是猶太人了.
神好像對他們有一種.
好像很特別的偏愛.
其實就是神一直的管理.
管教他們.
讓他們能夠慢慢的.
回復這份選民的氣質.
那個身份.
使得可以想像得到.
冠冕上面的寶石.
高舉在神的耳上.
整個地上將來的焦點就是以色列了.
所以神的出現.
帶來的就是和睦.
那個平安.
和榮耀.

$^{881}$在猶太人身上.
他們當然就不是他們自己本身.
而是彌賽亞賦予的生命的改變.
所以這個第九章.
有好幾段經文.
就是頭八節.
是講到神要對付列國.
接著中間最正宗.
就是第九第十節.
彌賽亞的進入聖城.
所以短暫的應驗.
已經應驗了.
耶穌第一次來.
祂就是騎著驢軀.
進入耶路撒冷.
真是猶太人看祂是那一位的君王.
不過後來不對勁.
為什麼不是幫我們復興.
對付羅馬.
趕走羅馬人.
巡撫彼拉多.
分封的王希律.
還有就是肥力.
這個加利利分封的王.
為什麼不是將他們全部趕走.
使得我們能夠重建我們的以色列.
所以後來發現不對勁.
就要致耶穌於死地.
原來是假貨.
不是真正的彌賽亞.
但是聖經已經預言了.
第九第十節.
第十節很明顯就是.
君王治理的管治權.
可以延伸到列邦地極.
還要用和平.
但是很難想像.
在這樣的預言下.
福音書的作者真的能夠.
你可以說他們真的.

$^{921}$耶穌升天之後.
聖靈的工作改變他們.
他們的言經的準繩度.
可以知道原來.
是的.
九和第十節是有兩個時間的.
彌賽亞耶穌基督降臨是有兩次的.
所以我們第一次見不到耶穌降臨.
但是我們等待的是第二次主耶穌的再來.
從天上再來這個世界.
其實是祂自己應許了的.
是新約祂親口答應使徒的.
祂一定會再回來的.
而且是駕雲再來的.
譬如我們站在使徒行傳第一章.
應該是第八節.
使徒行傳第一章.
是天使這樣說的.
我們從第六節開始看.
使徒行傳第一章.
他們聚集的時候就問耶穌.
主啊,你復興以色列國就在這時候嗎?.
他們很關心的問題是真的.
彌賽亞既然這樣成就了救恩.
按道理應該跟隨就復興以色列了.
耶穌對他們說.
「…賦憑著自己的權並所定的時候日期.
不是你們可以知道的.
但聖靈降臨在你們身上.
你們就必得著能力.
並要在耶路撒冷,猶太,全地.
和撒瑪利亞直到地極都作我的見證」.
也就是說有一段時期.
彌賽亞好像懸空了.
是需要地上的人.
信主的弟兄姊妹去傳揚福音.
去到地極.
第一批人就是使徒.
跟隨耶穌的門徒.
由耶路撒冷開始去傳.

$^{961}$所以來到二千年後的香港.
我們今天就接續著這個使命去傳.
也就是說彌賽亞的降臨.
是端扶我們外邦的教會.
來傳揚救恩.
讓不同的民族得到福音.
接著第九節.
「說了者說,他們正看的時候.
他就被取上星.
有一朵雲彩把他接去.
他就是耶穌.
便看不見他了.
當他往上去.
他們定睜望天的時候.
忽然有兩個人身穿白衣站在旁邊說.
加利利人啊.
你們為什麼站著望天呢?.
這離開你們被接星天的耶穌.
你們見他怎樣往天上去.
他還要怎樣來?」.
這是一個應許.
他將來都會駕雲再回來這個地上.
所以我仍然都是認為.
以色列一定是將來那個世界的焦點.
耶穌基督會腳踏橄欖山.
來拯救他的選民.
改變整個世界的格局命運.
所以第九,第十節.
是很明顯說有兩個時間.
彌賽亞第二次以君王.
真正的君王的身份.
來到再降臨地上.
所以如果這樣的分析.
一千年的千禧年.
又不可以說完全錯的.
真的地上是有一個千禧年.
「禧」就是說是好的和平的.
不會再有戰爭的.
所以去到十一節就再說.
將來石安被復興.

$^{1001}$神再回來眷顧他的子民.
以他們為弓.
以法連為箭.
這樣來做一個比喻.
當然不是真的射那個人出去.
不是炮彈射那個人出去.
而是真的再復興這個民族.
使他們能夠有無論尊嚴.
無論力量.
使列國都要降服.
這裡是有戰爭的暗示的.
也有很血腥的.
怎樣能夠取得平衡的公義呢.
這些就超出我們能夠認知的範圍了.
但現在肯定不是這場戰爭的規模和手段.
因為我不認為現在以色列這樣做.
就是在打聖戰.
我就不認為.
但是我們就不要從新聞讀回聖經.
我們應該在聖經裡面找到這個線索.
將來神的時間到.
彌賽亞第二次降臨.
前後都夾雜著戰爭.
都夾雜著神要對抗.
對付那些敵黨祂的人.
要抬舉祂自己的子民.
來成為弓箭.
所以我們都應該要羨慕.
我們在教會裡面.
我們期望自己在神的眼中.
我們是不是神牧養的群羊呢.
我們是越來越瘦.
又生枝.
又有很多暗瘡.
你有沒有見過這些狗.
你戒那些瘦猛猛的.
牠的生命力應該很強的.
不用打針.
又不用做美容.
但是很瘦弱.

$^{1041}$但是這裡告訴我們.
當那日耶和華他們的神.
必看他的民如群羊.
拯救他們.
因為他們必像冠冕上的寶石.
我們有沒有羨慕成為.
不是要追求最厲害.
而是我們在神眼中.
是祂冠冕的寶石.
我們是神牧養的群羊.
我們看看《約翰福音》十章.
耶穌自己說的.
十章的第十節.
或者第九節開始.
可能再長一點點第七節.
所以耶穌又對他們說.
「我實實在在的告訴你們.
我就是羊的門.
凡在我以先來的都是賊.
是強盜.
羊卻不聽他們.
我就是門.
凡從我進來的必然得救.
並且出入得草吃.
到賊來.
無非要偷竊殺害毀壞.
我來了.
是要叫羊得生命.
並且得的更豐盛.
是要叫羊得到生命.
得到豐盛.
圓文的意思.
不是更豐盛.
是得生命得豐盛.
我是好牧人.
好牧人為羊捨命.
若是故宮不是牧人.
羊都不是牠自己的.
牠看見狼來.
就是那些故宮.

$^{1081}$就撇下羊逃走.
狼抓住羊.
趕散了羊群.
故宮逃走.
因牠是故宮.
並不顧念羊.
我是好牧人.
我認識我的羊.
我的羊也認識我.
正如父認識我.
我也認識父一樣.
並且我為羊捨命.
我另外有羊.
不是這圈裡.
我必須領牠們來.
牠們也要聽我的聲音.
並且要合成一群.
歸一個牧人了.
我父愛我.
因我將命捨去.
好在取回來.
沒有人奪我的命去.
是我自己捨的.
我有權並捨了.
也有權並取回來.
這是我從我父所授的命令.
你可以想像到.
當時那些人.
對主耶穌的論述.
身份.
產生很大的好奇.
而主耶穌是避而不答.
但是現在.
在這個時候.
在這個所羅門的廊下.
在這個修奠節的時候.
那些人咄咄逼人.
逼耶穌一定要表態.
耶穌就說一些猶太人.
一聽就會明白的.

$^{1121}$就是群羊的牧者.
那個好的牧者.
真正的牧人.
撒迦利亞書之後.
很多關於牧人的比喻.
十章開始已經有了.
你可以想像到.
主耶穌祂說自己這個身份.
猶太人就很羨慕很想.
好牧人就應該是這樣的.
一定是對群羊有保護.
又讓他們能夠.
可以越來越.
有一種所謂豐盛.
就是越養越肥.
多毛.
不會有生的暗病.
不會有寄生蟲.
有瘡.
你可以想像得到.
主耶穌說我是好的牧人.
從我進來的必然會得救.
並且出入得草吃.
耶穌基督來.
就是要使羊得到生命.
得到豐盛.
就是這裡的意思.
冠冕上面的寶石.
我們要追求有主耶穌的喜悅.
我們要追求我們的生命成長.
成熟.
基督徒不成熟.
不因為我們神有偏心.
這個普遍的恩典.
是已經給每一個人的了.
有哪個人緊張.
珍惜和渴慕.
這個是很大的問題所在.
根源的問題所在.
下個月12月了.

$^{1161}$很多外國人.
經濟好的時候.
現在不好當然是.
他們的Blood Friday.
感恩節之後.
就有一個願望清單.
我最想要在聖誕節收到什麼禮物.
我們列一下這些清單也好.
我們是不是全部是物質的清單呢.
還是我們心靈都想.
我的平安能夠一直增加.
我可以成為平安之子.
我的盼望越來越有那個確具.
我對主的愛是與日俱增.
真的有一個很熱切的渴慕.
將來可以見到主耶穌和眾聖徒.
在教會裡面.
見到很多的後輩.
會默默地為這些小朋友祈禱.
雖然不能看著他們長大.
按正常的情況.
都會慢慢地漸漸老去.
但是這個漸漸老去的過程.
你有沒有在教會裡面更多的.
因為你已經很難再有很大力的活動.
那個活動能力一直減退.
但是那個腦袋那個心靈.
是應該與日俱大的.
那個心胸越來越闊.
看事物越來越有一種包容.
你能不能夠這樣來為很多的後輩.
尤其是新的傳道人.
更多為他屈膝祈禱.
默默地為他祈禱.
使他將來給神大用.
這些就是寶石.
官面上的寶石.
那個光輝.
我們很少人懂得這樣來到去.
覺得教會好像沒有什麼給他.

$^{1201}$有福氣,福利.
其實不需要的.
你坐下來就已經是福氣了.
你坐下來來敬拜神.
和親近人.
就已經是福氣了.
你的內心就會有一份的Blessing.
這樣離開.
回家.
繼續延續你的生命力.
所以這裡告訴我們.
當那日.
其實新約已經應驗了.
耶穌說.
我來就是使人得生命.
得豐盛了.
所以我們基督徒的人生.
不要只是希望.
有生活順利.
兒孫進入好的學校.
他們將來有好的工作.
這些一定是對的.
因為世界的人都是這樣期望的.
但是.
多走一步就是.
希望他們的靈性.
被神大用.
將來他們的人生.
這個就是我們基督徒的眼光.
這個眼界要這樣開闊.
否則就很難了.
最近都知道有年青人.
在外國讀書回來.
很關心教會的事務.
哪個牧者有什麼需要.
他都可以說到.
你可以想像得到.
愛主,愛教會的心.
日於言表.
離開了這麼遠.

$^{1241}$幾千里外.
都這麼了解教會的需要.
運作和祈禱.
這樣的靈性.
正正就是教會需要的下一代.
正正就是需要我們不同的長輩.
我們要為這個教會祈禱.
不斷有新一代的年青人.
有屬靈的眼光.
在世界的洪流.
可以不受影響.
每個人的內心的願望清單.
不是要裝模作樣.
寫我要怎樣屬靈.
我要看完多少次聖經.
不是要這樣做給人看.
而是內心真正的渴慕.
所以我自己很開心.
在一個忙裡面.
逢星期四晚都可以很透徹地.
跟大家講聖經.
將一些我們比較少留意.
或者沒有機會去讀的經文.
認真大家一起.
苦口婆心地認真地去討論.
真的這樣來看到.
這些才是真正基督徒.
我們要在神的裡面去追求.
我們要成為祂冠冕上面的寶石.
將來神就要令到祂的選民.
就是要這樣.
但現在未能夠這樣.
但我們基督徒就不應該了.
因為我們真的眼巴巴看到.
耶穌就是這位彌賽亞.
就是這樣死在耶路撒冷.
成為了我們今天信主教會的元首.
所以我們就應該要追求.
有這樣的成長.
這種成熟.

$^{1281}$所以保羅有很多的勸勉.
就是期望我們能夠.
比如和大家企立羅馬書.
就講到這個榮耀.
八章.
第31節.
羅馬書八章的31節.
既然如此.
還有什麼可以講.
就是這個意思.
因為他講到明.
那些被稱議的人.
是叫他們「得榮耀」.
八章的30節.
「如先所定下的人.
又召他們」.
就是我們信耶穌的人.
召回來的人.
就稱他們為「義」.
就是說沒有罪的了.
神已經赦免了他們的罪.
成為他們的兒女了.
所稱為「義」的人.
是要叫他們「得榮耀」.
神體為的榮耀.
「得榮耀」.
基督徒是要「得榮耀」.
所以你不要因為你的歲數.
因為你的各方面的限制.
這是神很清楚的旨意.
是保羅自己親口講的.
他是和外邦人講的.
和羅馬的教會這樣講都很貼切.
因為羅馬人在意大利首都.
見過世面.
多麼的虛幻的場景他們都見過.
羅馬的海船經常都有.
那些皇帝打勝仗.
海船們勝利的歸回.
拿了很多努力的戰利品.

$^{1321}$很多戰俘用一些魚鉤牽著.
就快要行刑了.
兩邊的甲道有很多的婦女.
撒那些香花.
所以有香氣.
和有死的氣味.
就是那些戰俘等等.
所以保羅講的榮耀.
是講述天神眼中的榮耀.
基督徒由被神選召.
到你信耶穌.
到你得救.
到你來到被稱義.
神說透過保羅.
就是說我們要在進入永恆之前.
我們應該是靈明的成熟.
得到榮耀.
這是保羅告訴我們.
所以他就這樣說.
既然如此.
還可以說什麼.
神這樣幫助我們.
誰可以抵擋我們.
神既然不愛惜祂自己的兒子.
為我們眾人都捨去.
豈不是也把萬物和他.
一同白白的賜給我們.
誰能控告神所揀選的人.
有神稱他們為義.
誰能定他們的罪.
有基督耶穌已經死了.
而且從死裡復活.
現今在神的右邊.
也替我們祈求.
誰能使我們與基督的愛隔絕.
難道是患難.
第一件就提患難.
因為保羅正正面對著.
因為傳道帶來的各種苦難.
是困苦.

$^{1361}$因為他要面對很多不同的處境.
困苦是說生活上的壓迫.
是迫不.
是飢餓.
是赤身怒體.
是危險.
是刀劍.
如經上所記.
我們為你的緣故終日被殺.
人看我們如張載的羊.
詩篇的44篇.
第22節.
第37節.
然而.
剛才說的是.
假如真的有患難.
有困苦.
假如真的有迫不.
假如真的有這些.
但是然而.
Therefore.
然而.
靠著愛我們的主.
在這一切的事上已經得勝有餘.
就是說.
上文所說.
假如的事情.
每個人的人生都會或多或少會面對的.
假如有一些橫疾.
是你一生要面對.
第二個人都不會.
但是沒辦法.
你的命就是要這樣受那些苦.
但是.
如果沒有了然而.
那就很慘了.
然而靠著陪伴我們.
愛我們的主.
在這一切上文的假如的事上.
已經是得勝有餘.

$^{1401}$就是說.
這個病纏繞我.
不纏繞任何一個人.
不要緊.
因為有主耶穌和我一起面對.
所以不再是假如.
是事實.
事實又如何呢?.
因為愛我的主和我一起面對.
那就可以跨性這個疾病了.
就是這樣的意思.
得勝有餘就是.
我不是沒有了這個病.
是我跨性這個病.
我陪伴這個病.
這個病成為了我可以榮耀神的一個場景.
這樣的意思.
譬如暗失之後.
不可以見一點光.
但是他的靈性.
很多傳道人去找他的時候.
是得到幫助.
得到安慰地離開.
是得到力量鼓勵地離開的.
其實不認識他.
但是那個暗失之後.
我已經為你祈禱了.
這樣的.
一見到那個傳道人.
就說我已經為你祈禱了.
其實不認識的.
這些長輩.
你可想而知.
原來得勝有餘就是這樣的意思.
不是沒有那樣東西.
所以神容許地上有疾病.
患難.
不同的人生處境.
都是磨練我們學習.
在主耶穌的愛的同在底下.

$^{1441}$去跨性我們的困難.
38節.
因為我深信無論是死是生.
是天是掌權.
是有能是現在.
是將來的事.
是高處是低處.
是別的受造之物.
都不能夠叫我們與神的愛隔絕.
這愛是在我們的主基督耶穌裡.
神的愛是藉著耶穌.
彰顯在我們的心靈裡面.
這東西是永遠屬於基督徒的.
所以如果我們不去追求主的同在.
不會感受到我們有好目人的模樣.
我們不會感受到我們得生命.
得祂的豐盛.
我們就面對困難.
我們就不能夠變成.
那個冠冕的寶石.
我們散發出的可能就是很多的埋怨.
很多的怨天尤人.
很多的反抗.
很多的苦毒.
基督徒都有的.
是很苦毒的都有的.
有時牧者我們都改變不了他.
人生不要停留在你自己那個囚禁的監獄.
就是你的心靈.
可以很多事做的.
但是.
這裡告訴我們.
冠冕上的寶石.
是我們一生要去追求的.
好了,稍微停一停,喝一口水.
幫大家再.
抓緊一點那個君王的概念.
先跟大家揭幾段《牛若經文》.
《創世記》的49章.
因為第九章這裡已經很清楚告訴我們.

$^{1481}$主耶穌就是那位彌賽亞的君王.
和平之君.
49章其實就是雅各的臨終的預言.
49章.
可能看快一點.
第十節.
也是很出名的.
這個就是.
雅各對十二個兒子的祝福.
有些就是沒有福的.
不會有福給他的.
例如.
第五節.
西冕和利未.
他們是弟兄.
記不記得他們做過什麼.
就是殺了士劍全家全族.
為了他的妹妹.
這個底娜.
來報仇雪恨.
說你想娶我們的妹妹.
行.
我們猶太人是要.
走過來的.
那你們就走吧.
士劍就說.
他愛底娜.
就鼓勵.
勉強因為他是王子.
是他族長的兒子.
就叫自己的叔伯兄弟一起.
全部行割禮.
其實等他們割完禮之後.
那個傷口還沒有癒合.
西冕和利未.
這裡就說.
他們的刀劍是殘忍的器具.
我的靈不要與他們同謀.
我的心不要與他們聯絡.
因為他們趁勞殺害人命.

$^{1521}$淫意砍斷牛腿大根.
他們的怒氣暴烈可咒.
他們的憤恨殘忍可助.
我要使他們分居在雅各迦裡.
散住在以色列地中.
又說是祝福的.
對這兩個兒子.
全年咒詛來的.
一點福都沒有.
真是咒詛來的.
舊事重提.
翻舊帳.
現在才臨死之前的秋後算帳.
爸爸為什麼這樣.
臨死應該.
你不分多分少一點.
一點祝福都沒有.
還說是屬靈的祝福.
你可以想像得到.
西冕和利未就很大件事了.
有一個預言就是.
我要使他們分居在雅各迦.
即是分散他們這兩個支派.
果然利未支派將來就分散了.
他們十二支派的原因是什麼.
因為他們成為了利未人.
祭司的家族.
所以後來他們就可以分散.
當然這個後話.
因為這個咒詛改變成為祝福.
就是神人摩西在他臨終之前.
替利未支派有一個很大的平反.
這些後話有機會再看.
但是我們想看的第八節.
就是猶大.
很特別.
猶大支派將來會有個君王出現.
第八節.
猶大,你弟兄們必讚美你.
猶大的意思就是讚美.

$^{1561}$你手必撼住仇敵的頸鏖.
你父親的兒子們必向你下拜.
猶大是個小獅子.
我兒你找了食便上去.
你屈下身去餓如公獅.
遵如母獅.
誰敢惹你呢.
那些象必不離猶大.
象必不離他兩腳之間.
直等細羅,就是平安者.
來到萬民都必歸順.
有一個這樣的預言.
在猶大會有一個領導者.
這個細羅.
就是會掌王權.
他會來到去.
萬民會歸順.
有一個預言.
在雅各很早期.
這個民族都未出埃及.
已經有這個預言.
將來猶大之派會有這樣的君王.
我們很快又揭尼加書.
應該大家都熟的.
第五章.
尼加書的第五章.
即大先知之後.
以摩訶那尼.
即約拿書之後就是尼加書.
第五章的第一節.
成群的民啊.
現在你要聚集成隊.
因為仇敵圍攻我們.
要用象擊打以色列審判者的臉.
百里行以發他.
你在猶大諸城中為少.
將來必有一位從你那裡出來.
在以色列中為我作掌權的.
他的根源從耿古.
從太初就有.

$^{1601}$尼加書第五章第二節.
就是那幾個博士.
上來找希律王.
大希律.
就問.
生在猶太的君王在哪裡.
希律王就很害怕.
生下來做猶太人的王.
豈不是就是我.
我是分封的猶太的王.
但是他很害怕有人取締他.
是什麼原因呢.
就找那些以色列的博士,文士來.
就問這些文士.
基督生在哪裡.
他們一引就引.
五章第二節.
就是百里行以發他.
就是那些文士.
所以他們是熟讀聖經的.
到底他們有沒有一起去百里行.
去找那個基督呢.
就不知道了.
聖經就沒有記載.
突然間鏡頭就轉到去.
在百里行的野外.
就有天使,天君來降臨.
在那些牧羊人當中.
告訴他們.
基督生了在百里行.
他們去到就見到瑪利亞抱著那個嬰孩.
他們就拜瑪利亞.
你可以想像到.
這個是君王的寓言.
再看詩篇一段.
72篇.
這些是比較明顯的.
彌賽亞的寓言的經文.
和我們剛才看的第九章.
那個君王都是非常貼切.

$^{1641}$就是第九章的九節十節.
所以盼望大家都回家.
將九章撒迦利亞書九至十節都背熟它.
很多人只認識第九節.
不知道有第十節.
不知道原來將來彌賽亞.
統管萬國地極.
是以他的和平去統管.
你可以想像到.
即是說你認識了多過很多人.
72篇.
是所羅門的詩.
可能讀快一點.
可能第一節開始.
給大家有少許概念.
神啊求你將判斷的權柄賜給王.
將公義賜給王的兒子.
他要按公義審判你的民.
按公平審判你的困苦人.
大山小山都要因公義使民得享平安.
他必為民中的困苦人伸冤.
拯救窮乏之輩.
壓碎那欺壓人的.
太陽還存月亮還在.
人要敬畏你直到萬代.
他必降臨.
像雨降在已割的草地上.
如金林滋潤田地.
在他的日子.
二人要發旺大有平安.
好像月亮長存.
第八節留意了.
他要執掌權柄.
從這海直到那海.
從大河直到地極.
接著下面不讀了.
你可以想像到.
所羅門他有那個影子.
因為他都是和平之君.
所羅門有另外一個名是神賜給他的.

$^{1681}$叫耶底底亞.
可能是一個譯音.
希伯來文就是平安的君王.
所以這裡一路說真正的治理.
是要用和平不是用手段不是用強權.
是要用愛是要用和睦.
這樣來治理是用公義的.
但是你看到所羅門有生之年都不行.
他自己都不是這樣.
但是他描述的那個彌賽亞的治理.
又是一樣了.
去到第八節就是.
去到地極.
都是這位君王那個治權的範圍.
你可以想像到.
這個一定是彌賽亞的掌權的日子.
才會有這樣的可能性.
整個世界趨向了公義和和平.
那個美麗的光景.
就會是人人都羡慕.
你可以想像到秋天.
你下午走過湖邊這樣很舒服.
你見到小朋友一堆一堆在大草地上.
這樣玩這樣野餐.
這些景象你就可以看到.
星平和睦家家戶戶都這麼開心.
就是那個真正的管治.
這個是聖經的期望.
是聖經的預言.
但是你可以想像到現在怎麼可能有呢.
你如果有機會看到那些難民.
你就知道那種慘況.
我們還有一點時間進入第十章.
特別提到那個牧羊.
開始提到這位浮出水面的牧者.
大牧人.
我們的主耶穌.
九章剛才已經提到.
他要來親自去牧羊.
他自己的群羊.

$^{1721}$九章十六節.
會使他們能夠成為他的冠冕的寶石.
現在第十章這裡.
就告訴我們.
他要將那些假的.
使人離開神的偶像的敬拜.
來去懲罰.
更加將那個怒氣.
向那些牧者來去發作.
第一節開始.
當春雨的時候.
春雨即是叫晚雨.
Later rain.
地中海氣候.
十月就是秋雨.
就是叫初雨.
即是第一陣的雨.
去到大概三四月.
第二陣就叫秋雨.
或者叫晚雨.
聖經很多講的.
一提到這些.
就是那些好日子.
第一節.
當春雨的時候.
他們要向發閃電的耶和華求雨.
他必為眾人降下金林.
使到田園生長菜蔬.
因為家臣所言.
似乎那些巴力.
又是人沒有什麼訴求.
你在迦南地.
因為那些氣候那麼乾燥.
那麼多變.
如果能夠有四季的雨水.
滋潤田地.
你是一個農耕的社會.
自然就可以豐收.
人都是需要想豐收.
想有好日子過.

$^{1761}$所以他們拜巴力.
都不是沒有原因.
因為他們就是要保障自己的農作業.
所以家臣就成為了.
拜耶和華以外的選擇.
因為家臣所言的是虛空.
占卜所見的是虛假.
作夢者所說的是假夢.
他們白白的安慰人.
只有一個「講」字不單止.
是毅然聳聽.
是帶你離開真理的.
無論是預言.
占卜.
或者是那些疑夢.
因為人都是想有些靈界的經驗.
靈界告訴你怎樣怎樣.
你就很想得到.
但是神就說他要摧毀這些.
因為.
所以眾人就如羊流離.
因無牧人就受苦.
不單止是這樣.
再加上就是.
真正本來扮演著牧羊角色的領袖.
扮演不了.
失職之後.
那些羊,即是百姓.
就流離了.
經常都描述這樣.
沒有人帶領就流離了.
大家都不知道什麼正確的方向.
然後就受苦了.
所以現在的光景就是.
神要百姓去祈禱.
祂就會降下金林.
就是說.
你最需要的.
你尋求我.
我就會賜下給你.

$^{1801}$當時的世代一定是對雨水最緊張.
因為不可以生存.
是一定會面對.
一乾旱.
一失修.
就會帶來很大的生活.
甚至乎國家面臨的危機.
所以他們為什麼一定要拜巴力.
但是在.
這個.
以利亞的年代.
正正拜巴力就是沒有雨水.
所以就要鬥法.
原來原因是.
以利亞祈禱叫天不要下雨.
所以那些人就很生氣.
以利亞.
以利亞就叫齊那八百五十個巴力的先知.
上加密山.
神要來做一個判斷.
但是現在不用了.
神就說叫百姓.
主動的來去.
在那個季節.
應該是有晚雨.
春雨的時候.
就要向發閃電的耶和華求雨.
就是說真正給到福氣你的是神.
今天當然不是我們再這樣祈禱下雨.
今天我們就要看到.
我們的生命的呼吸.
我們的生命的一切.
其實都是神賜給我們的.
我們要回到這個賜恩的源頭.
去仰望祂.
去尋求祂.
去認識祂.
我們就會得到.
祂所賜給我們的甘霖.
田園生長的菜蔬.

$^{1841}$不要再求問那些家臣.
不知道基督徒還會不會去看相.
還會不會去灣仔鵝頸橋打小人.
還會不會去黃大仙.
自己走過去,撕一撕紋埋.
走過去裝朱鄉.
好像王夏惠.
一定要裝頭柱香.
是不是基督徒我也不知道.
那些人說他是.
但是就又扮鬼扮馬.
剛剛的對比.
假神,家神.
家神一定是和你很熟的.
你很有感情的.
你小時候就契了他.
猶太人就是這樣.
他們就是進到迦南之後.
發現拜神,真神很麻煩.
又要去聖殿.
又要很多手續.
又有些事不能做.
但是拜巴力就什麼都不用.
你給了他.
他就會賜福給你.
就會下雨.
這樣就很靈.
人就是需要那個靈驗.
其他事你不要碰我.
總之靈驗就可以了.
但是神就說.
其實真正的福樂是神那裡而來.
要回到那個源頭.
要向這位發閃電的耶和華求雨.
當然他不是閃電神.
是那個源頭.
生命的源頭.
所以我們信了主.
其實我們就真正意識到.
我們生存的價值所在.

$^{1881}$原來是出於我們的造物主.
我們才有我們的生命氣息.
我們的生活.
我們的疆界.
都是由他所定.
所以我們要回到這個源頭.
去看我們的人生.
這樣就會很不同.
我們作為牧者都見到不同的人.
真的有不同的際遇.
有些人真的很慘.
但是我們不能出聲.
因為神給每個人的際遇.
我們不能出聲.
不關我們的事.
我們不能去指指點點.
神啊你這樣不公道.
他就這麼慘.
我們不可以這樣說.
但是我們是要牧羊去陪伴.
和要讓他知道.
你現在的慘.
沒錯.
在那個相對裡面是很辛苦的.
因為有比較.
如果全世界每個人都是.
跛手跛腳就沒問題了.
好像我們小時候住在武安區.
每個人都這麼窮沒問題.
但是當你長大之後.
你有很多比較.
你就會思前想後.
正正不要這樣思前想後.
回到那個源頭.
牧者就是要做這件事.
回到那個源頭.
我解釋不了.
我沒得解釋.
我更加不可以代神解釋.
但是告訴你.

$^{1921}$回到這位像閃電的神.
他就會降下那個春雨.
他就會使到眾人.
得到那個生長的菜蔬.
就是有回你的尊嚴.
有回你的那個幸福感.
有回你自己那個生命的意義力量.
每個人都有自己的跑道.
有些人是跑一百公尺.
有些人是跑長跑的.
每個人都有不同自己的人生.
是神賦予的.
牧者就是告訴你.
無論你怎樣的人生.
神都愛你的.
在基督耶穌裡面.
祂已經收納了你.
所以你是在你那個人生去榮耀神的.
所以這裡就是.
加神就不是了.
不要管那麼多.
你賄賂我我就會給你.
所以很多人就情願拜那些四面佛.
拜那些人.
不知那個甚麼泰國那個伯甚麼.
伯甚麼.
對了 伯龍王.
即是很靈的.
可能真的很靈的.
但是基督徒我們不需要.
我們要回到那個次因的主的源頭.
祂要我怎樣.
那我就怎樣.
這種不是一種的叫做無奈.
而是在神安排的構造的世界裡面.
去活好我的角色.
將來就會得到我那個份的賞賜.
每個人有他的角色.
至於苦難.
這些超過了我們人類能夠理解的.

$^{1961}$我們最好就不要太問.
那個背後的因果那些東西.
而是我既然肯定知道.
耶穌基督愛我這個罪人.
為我釘十字架.
祂給我的人生一定是祂認為最好的.
所以應該這樣理解.
我都提過.
我在八歲九歲.
那個時候已經有意識了.
住木屋區.
住木屋區都有些人很有錢的.
做生意.
但是我家裡就特別窮.
十三兄弟姐妹.
很多的吵鬧.
很不想在那個家裡.
知道有個有錢人想收養我.
最慘的是長大之後才知道.
當時沒有人告訴我.
就帶我去加拿大.
媽媽當然說不行.
雖然說十三個.
沒有了一個是無所謂.
當然不行.
始終是骨肉.
但是我長大之後.
再去想.
如果我當時被帶去加拿大多倫多.
現在是什麼世界.
我已經可能是當時很出名的人.
或者可能是很壞的人也不奇怪.
但是當我第二次去探望我姐姐的時候.
她告訴我.
收養你的娘太就住在附近.
氣得我快死了.
當年如果她收養了我多好.
但是你在基督裡面再回望.
總有神管你住.
正正就是要在這樣的環境.

$^{2001}$來到一路長大.
一路很多很多青少年階段的憤世嫉俗.
又不肯讀書等等.
才能夠在那個階段順處的生命改變.
所以我不會再有一種遺憾.
或者是可惜.
媽媽為什麼你這麼殘忍.
不把我給了人.
不需要.
反而就應該從另一個角度.
回看一切有神管你.
一切如果不是神掌管著.
不容許.
自然就不會發生.
我父母都未信耶穌.
他怎麼會知道有神的計劃.
但是我們信了主就應該這樣理解.
有神的管理.
所以我才會在我姓曾爸爸的家長大.
現在我們的兄弟姐妹的關係非常密切.
非常好.
各樣東西都是有那個掌管的源頭.
所以我們求主讓我們能夠.
向這位發閃電的耶和華求雨.
不要向家臣.
現在那些人就說被那些KOL.
佔便宜騙幾十萬.
太多這些家臣.
現在的YouTube頻道更加多.
每個人都告訴你.
我教你這樣就一定可以.
當然我說完了.
但是這段聖經告訴我們.
回到那個切恩典的源頭.
就是這位愛我們的耶和華.
舊約已經預言了.
他們要回去.
一來不是舊事重提.
去責備他們.
而是他們要回去.

$^{2041}$你想尋求春雨.
你想尋求真正的Blessing.
回到切恩典源頭的主.
這位能夠發閃電的耶和華.
所以.
夠不夠時間呢.
好了不看下面了.
下次我們第三節開始.
我們下個星期四再繼續.
留意一下那個當那日.
大家回去做一下功課.
由第九章到十四章.
出現多少次.
當那日.
出現很多次.
一起祈禱.
多謝天父的恩典.
讓我們能夠透過第九章.
其實第八章我們就見到一個預言.
將來萬國的人.
都要去懇求耶和華的恩.
都會拉住猶太人的衣襟.
來希望認識神.
幫助我們能夠看到.
我們在主耶穌裡面.
是很有那份榮耀.
我們要將這個救人靈魂的福音.
很有勇氣地向人宣講.
來讓更多人明白.
這個世界原來是有真理.
是有真神的愛.
求主幫助.
多謝你們讓我們有這個半小時的時間.
一起認真的學習聖經的真理.
奉主耶穌基督的名.
阿們.
多謝各位.
阿們.
\newpage



\section{詩篇 116:1-19}
\label{sec:URxw02niNWY}
\textbf{2023-11-04 晚堂《以感恩為祭》詩篇 116\_1-19 講員:曾裔貴牧師}
\newline
\newline
連結: \href{https://youtube.com/watch?v=URxw02niNWY}{\texttt{https://youtube.com/watch?v=URxw02niNWY}} ~~~~ 語音日期: 2023-11-07
\newline
\newline
\hyperref[sec:j9CuOv5EhDA]{\small{< < < PREV SERMON < < <}}
~
\hyperref[sec:index_chronic]{\small{[返順時目]}}
~
\hyperref[sec:index_scriptual]{\small{[返順卷目]}}
~
\hyperref[sec:EpQdyi4IyoI]{\small{> > > NEXT SERMON > > >}}
\newline
\newline
詩篇 116:1-19
\newline
\begin{longtable}{cl}
\hline
\hline
章節 & 經文 (和合本修訂版)\\
\hline
116:1 & \begin{tabularx}{0.7\textwidth}{X} 我愛耶和華，因為他聽了我的聲音和我的懇求。 \end{tabularx} \\ \\ \relax
116:2 & \begin{tabularx}{0.7\textwidth}{X} 他既向我側耳，我一生要求告他。 \end{tabularx} \\ \\ \relax
116:3 & \begin{tabularx}{0.7\textwidth}{X} 死亡的繩索勒住我，陰間的痛苦抓住我，我遭遇患難愁苦。 \end{tabularx} \\ \\ \relax
116:4 & \begin{tabularx}{0.7\textwidth}{X} 那時，我求告耶和華的名：「耶和華啊，求你救我！」 \end{tabularx} \\ \\ \relax
116:5 & \begin{tabularx}{0.7\textwidth}{X} 耶和華有恩惠，有公義，我們的神有憐憫。 \end{tabularx} \\ \\ \relax
116:6 & \begin{tabularx}{0.7\textwidth}{X} 耶和華保護愚蒙的人；我落到卑微的地步，他救了我。 \end{tabularx} \\ \\ \relax
116:7 & \begin{tabularx}{0.7\textwidth}{X} 我的心哪！你要復歸安寧，因為耶和華用厚恩待你。 \end{tabularx} \\ \\ \relax
116:8 & \begin{tabularx}{0.7\textwidth}{X} 主啊，你救我的命脫離死亡，使我的眼不再流淚，使我的腳不致跌倒。 \end{tabularx} \\ \\ \relax
116:9 & \begin{tabularx}{0.7\textwidth}{X} 我行在耶和華面前，走在活人之地。 \end{tabularx} \\ \\ \relax
116:10 & \begin{tabularx}{0.7\textwidth}{X} 我信，儘管我說：「我受了極大的困苦。」 \end{tabularx} \\ \\ \relax
116:11 & \begin{tabularx}{0.7\textwidth}{X} 我曾驚惶地說：「人都是說謊的！」 \end{tabularx} \\ \\ \relax
116:12 & \begin{tabularx}{0.7\textwidth}{X} 耶和華向我賞賜一切厚恩，我拿甚麼來報答他呢？ \end{tabularx} \\ \\ \relax
116:13 & \begin{tabularx}{0.7\textwidth}{X} 我要舉起救恩的杯，稱揚耶和華的名。 \end{tabularx} \\ \\ \relax
116:14 & \begin{tabularx}{0.7\textwidth}{X} 我要在他的全體百姓面前向耶和華還我所許的願。 \end{tabularx} \\ \\ \relax
116:15 & \begin{tabularx}{0.7\textwidth}{X} 在耶和華眼中，聖民之死極為寶貴。 \end{tabularx} \\ \\ \relax
116:16 & \begin{tabularx}{0.7\textwidth}{X} 耶和華啊，哦，我是你的僕人；我是你的僕人，是你使女的兒子。你已經解開我的捆索。 \end{tabularx} \\ \\ \relax
116:17 & \begin{tabularx}{0.7\textwidth}{X} 我要以感謝為祭獻給你，又要求告耶和華的名。 \end{tabularx} \\ \\ \relax
116:18 & \begin{tabularx}{0.7\textwidth}{X} 我要在耶路撒冷當中，在耶和華殿的院內，在他的全體百姓面前，向耶和華還我所許的願。哈利路亞！ \end{tabularx} \\ \\ \relax
116:19 & \begin{tabularx}{0.7\textwidth}{X} 見上節 \end{tabularx} \\ \\
[1ex]
\hline
\hline
\end{longtable}
$^{1}$我們一起祈禱.
感恩天父讓我們能夠在今天晚上.
再一次我們可以在聖餐裡面.
去思念耶穌基督和我們的生命關係.
使我們的心靈可以有這份平靜.
有內心的福樂.
真是屬靈的心靈福樂.
是錢都買不到的.
有時甚至困難.
都不會阻礙我們這份平安.
這份安息.
我們今天晚上會看這篇詩篇.
我們真的很想用感恩作為我們的祭.
我們的內心充滿一份對上主的重陽.
更加我們能夠在生活裡面.
在我們每天大家都不同的崗位.
不同的生活裡面.
我們去學習重陽我們的主.
我們這樣等候禱告.
奉耶穌基督的名.
阿們.
很快就是聖誕節了.
在歐美的地方.
中國人就比較少.
他們很快就會有一個願望的list.
我聖誕希望收到什麼禮物.
感恩節有什麼禮物.
我們在這幾年.
我們都定了.
我們的傳道同工都定了.
11月我們要學習感恩.
我們過去很少這樣想的.
我們就覺得不是歐美的人的專利.
他們都可能與舊恩和信仰無關的了.
總之就是感恩節.
他們就很重視的.
但是我們是要學習這個感恩的文化.
怎樣能夠承載我們每一個人的生命呢.
這一篇詩篇.
其實不是獨立的.

$^{41}$113到118篇.
有一個系列.
叫做Hallelujah的詩篇.
是猶太人他們在逾越節.
他們會唱頌的詩篇.
其中116篇.
就是這一篇詩篇.
就是逾越節最高潮的時候.
他們就會唱這一首的詩歌.
來頌揚這一位主的拯救.
這一篇詩篇.
其中重要的主題.
就是這個詩人.
他示範他自己.
怎樣面對很大的困難裡面.
他一生要學習求告神.
這個功課.
我們一生要學.
相信我們每一個人.
如果你是信了主耶穌.
你的入門.
一定是因為你首先.
你求告主名.
你就得救.
一定是這樣.
當然你是怎樣能夠信.
就是由聽道而來.
聽道就是由基督的說話而來.
這是羅馬書十章.
告訴我們的步驟.
就是有人跟你說.
關於救恩.
你聽到這個道.
多特別.
聽下有點感動.
你開始更加.
生命的比對.
有個對照的時候.
你越來越你的心靈打開了.
到有一天.

$^{81}$或者一個時候.
你在你的決志者的面前.
你有一個很簡單的祈禱.
主耶穌.
我願意相信你.
這是每一個人必須經過的步驟.
就是你要求告主名.
不過詩篇116篇.
作者詩歌的人.
不知道是誰.
他有一個很大的生命功課.
是要我們大家一起學習.
這裡說的是.
他要一生的求告這位主.
而且更加重要的是.
他告訴我們.
他對耶穌基督.
對這位真神.
他有一個愛.
第一節已經說了.
為什麼他會愛這位神.
而且我們看下去.
這個愛.
不是只有一種感情上的愛.
是要有一個行動.
是整個意志的投入.
對神有一個這樣的效忠.
你留意下.
其中就很清楚.
第二節就已經說了.
他為什麼要愛這位耶和華.
這裡告訴我們.
因為他聽了我的聲音.
和我的懇求.
他既然向我質疑.
我一生就要求告他.
一生要求告神.
不是停留在信耶穌.
決志的那個祈禱.
而是他開始和神建立一個生命的關係.

$^{121}$到底他有什麼遭遇.
讓我們可以來效法.
或者有相近的遭遇.
他就不說得太仔細了.
他只是一些.
可能是一些描述.
但不是具體的某一些的情景.
當然下面他會提到.
人都是虛荒的.
第十節.
十一節.
第十節就是保羅在哥林多後書有引述.
他經歷這個困難.
憤外感受到耶穌基督的復活大能.
在這裡作詩歌的人告訴我們.
他曾經描述.
眼流淚.
腳幾乎跌倒.
在他面前有神的同在扶持.
他才可以走活人的路.
所以他就發現.
人都是說謊的.
虛荒的.
那就是說意味著.
他要在禱告裡面.
避開了那些困難.
和那些會令到他失腳的人.
使他能夠深深的感受到.
神在他生命裡面的同在.
聽他的祈禱.
這篇詩很強調.
神聽他祈禱.
所以他要愛這位神.
這個經歷.
這個信仰.
就很深了.
就是我們人生.
要隨著我們人生的際遇.
你的每日生活.
就要來提升禱告的階梯層次.

$^{161}$因為不只是入門的.
信主的決志祈禱.
而是在你的生活裡面.
你的禱告.
你感受到神聽你祈禱.
他是側耳聽.
他會來成就他的旨意.
所以他一生要求告.
我們這位神.
帶來的結果.
這個互動就真的很美麗.
因為我面對困難.
因為我認識這位神.
他是聽祈禱的.
他是有恩惠有公義.
是以憐憫為懷.
他這樣在禱告過程.
對這位神不是遠遠的.
陌生的一個高高在上.
而是他愈來愈感受到.
他們猶太人.
在出埃及記的十誡裡.
這位神的屬性.
這位神是他感受到.
不是摩西的神.
是這個詩人的神.
這樣在他的生活裡.
他能夠和神的屬性拉上關係.
和他更重要的.
作為一個對其他人的教導.
就是他要一生愛這位神.
當他在很大的困難裡的時候.
他開始描述第三節.
就說到他一些遭遇.
所以有遭遇.
有困難.
人才會認真祈禱.
沒有困難.
你期望你的基督徒人生.
是一帆風順.

$^{201}$是不會令到你的祈禱生活.
祈禱的生命.
對神的認識.
是會增加的.
是會有一個相輔相成.
似乎如果我們沒有經驗到.
這些地上的困難.
或者我們逃避一些.
神要我們學習的功課.
我們不會成長.
似乎這裡這樣說.
是幾合聖經作者的經驗.
那個教導.
他遇到什麼事情.
相信不是在無的放矢的.
來胡亂的構想.
是真有其事的心靈經歷.
死亡的繩索纏住我.
陰間的痛苦將我抓住.
我遭遇的是患難受苦.
那時他又祈禱.
那時我就求告耶和華的命.
耶和華求你救我的靈魂.
心靈 靈魂 內心.
所以不是一按個鍵.
就可以解除的困難.
可能是持續的.
可能是因為他的背景.
他的際遇.
是必須要面對的.
這樣就無法逃避.
唯一的方法.
最好的方法.
最重要的教導就是.
回到神的裡面.
向這位有恩典有憐憫.
有公義的神去求告.
四節聖經已經說了三次求告.
一生的求告.
在困難很大的困境的時候求告.

$^{241}$我們有多少弟兄姊妹.
在這樣的困難裡面.
說說自己的一些經歷.
可能我以前有機會都說過.
我讀建道第二年神學的時候.
開始要上原文.
很痛苦 很難.
每個星期測驗.
第二年因為我們四年.
開始跟同班的同學.
就開始又擦出火花.
來自不同背景.
來自不同的神學.
開始有很多不同的看法.
因為學校要求每一班的同學.
必須有自己的班會.
有自己的祈禱會.
我們隔個禮拜五的中午.
十幾年前就一起有祈禱.
其中一個同學就很靈恩.
我就很飢餓.
我頂不住他.
說唱的歌又手舞足蹈.
我那時候大家都是神學生.
但我頂不住他.
他就被委派做主席.
每次都是他大詩歌.
很簡單.
我又經常要走讀.
就不去祈禱會.
於是就令到那班同學分了兩班.
因為我這樣一個頂不住他的神學.
就開始出現一些張力.
有些同學就很尷尬.
怎麼辦呢.
想挽回又不行.
大家都是神學生.
令到我自己每次見到他就避開.
於是很深刻.
有一次上完教會歷史.

$^{281}$在長洲坐船回元朗.
兩個小時四點鐘放學.
趕四點三船.
衝下山要走.
回到家元朗是六點.
一路拖著行李箱一路走回家.
突然間電話響.
原來剛剛上完課的老師.
他很細微的觀察到我.
近來很不開心.
其實我不知道他為什麼會見到.
隱藏都沒有人知道.
是同班同學就知道.
因為那些事情.
有些神學看法.
大家面左左.
平時是沒有事的.
但這位老師很細微的主動打給我.
我還沒有回家.
還沒有放下行李.
就聽到電話.
他說: 瑞貴 近來我觀察到你很不開心.
有為你祈禱 什麼事.
下星期上學你來找我.
我和你聊聊天.
我當然說: 沒事沒事.
我回到家之後有很深的反省.
很深的自己在神面前去祈禱.
其實是令到全班的氣氛不好.
其實祈禱會是好的.
大家一起為讀書的過程.
大家在實習的教會.
大家在不同的家庭生活.
大家彼此扶持.
這只是因為不同的做法.
唱歌跳舞好些 活潑些.
我就受不了.
就因為這樣受不了 令到大家都很尷尬.
經過老師這樣打電話給我.
其實他完全不知道什麼事.

$^{321}$我整個的生命突然間有一個改變.
當然不是第二個星期就有行動.
我回去多謝這位老師.
也過了短短一段時間之後.
我向這個同學認錯.
我看到自己的陰暗.
看到自己的胸襟這麼有限.
我來到這裡只是因為別人的做法.
他很有心想帶動大家的氣氛.
一起敬拜 一起祈禱.
居然在我這麼有限的神學裡面.
就認為這麼靈恩不行.
就變成了不關靈恩是否靈恩的事.
只不過個人的心胸太狹隘.
導致班裡很緊張.
老師一個電話扭轉我.
令到我正視自己的內心 祈禱.
令到我能夠從死亡的繩索.
纏繞著我的都可以解套.
感恩啊.
男士說對不起是很大件事.
彼此完全能夠釋懷.
所以我們這麼多年已經畢業了.
我們是很好的同學.
有時放假我們就去玩等等.
可以是這樣的一個扭轉.
神差派這位老師.
來打這個電話給我.
好像讓我纏繞著的繩索.
陰間的痛苦不要再抓住我.
需要祈禱.
所以作者深深感受到.
神沒有離開他.
到底什麼抓住他.
什麼纏著他.
什麼讓他面對周遭的患難困苦.
唯一的解脫 唯一的辦法.
就是有人帶你跳到神裡面去禱告.
所以教會其實是一個很重要的場景.
就是大家安靜聆聽上主的道.

$^{361}$讓你可能舉一反三.
你有一個足淚旁通.
是喔 曾牧師你這個經歷.
其實可能我正在面對.
我都可以去向神尋求解套.
這樣的祈求.
正正就是私人感受到神愛他.
所以他願意愛這位神.
當然第七節更加寶貴.
就是當他面對這些事情.
他自己來發現.
自己原來不是那麼聰明.
如果用自己的想法.
其實很愚昧.
當他落到一個謙卑的境況的時候.
他就發現他可以在神裡面迎歸安樂.
心靈迎歸安樂.
我看到一個例子.
都很感動我自己的心靈.
有一對夫婦 很好的夫婦.
他是台灣陳雲子.
雲就是雲端的雲.
子就是地的子.
你可以搜尋一下.
當然散會之後.
搜尋一下.
他是台灣恆春一個醫院.
台灣最南面的一間醫院.
很窮的地方.
那個醫院的院長.
他做了很多年了.
但他其實在台北開一間.
叫信望愛的醫院診所.
眼科的專科.
是很賺錢的.
是很多人的.
因為他本身是很專業的.
眼科的專科醫生.
但是太忙了.
他整個人burn out.

$^{401}$要放下他的醫務所.
回去美國.
去休養.
當一切都停頓了之後.
他突然之間.
又想到自己的生命有什麼意義呢.
他突然間想到.
他有這個專業.
是神給他的.
為什麼他不可以來貢獻.
幫助其他人呢.
他就決心.
那個病稍微好轉了.
他就回到台灣.
當時他就走遍整個台灣.
去尋找什麼是神.
可以令到他貢獻他的人生.
來去服侍神的呢.
他一直去觀察的時候.
終於讓他發現.
恆春這個地方.
台灣最南面的一個地方.
是非常的窮乏.
方圓幾百里.
一個眼科醫生都沒有.
所以當地的人.
很多白內障.
而且那些少數民族.
山區的少數民族.
那些原住民.
是很迷信的.
認為要做眼的手術.
是會得罪神明的.
就是這些想法.
所以情願盲.
都不會做手術.
所以真的有很多眼的疾病.
是很無辜的.
這樣去浪費了時間.
這一位陳醫生.

$^{441}$是很好的基督徒.
他看到神.
要他在他的人生.
要在這個地方來貢獻.
所以他用了很多的時間.
和他的太太.
來到去那個地方.
去視診.
幫助裡面的這些人.
慢慢改變他們的想法.
使到這個恆春這個地方.
慢慢那個眼的疾病.
一路一路減少.
甚至去到一個地步.
一些原住民.
都來和這位陳醫生.
建立了那個關係.
信任他.
由醫療慢慢進入福音.
因為他發現.
就算救到那個人的眼睛.
救不到他的心靈的眼睛.
其實都是沒有用的.
重要的是救到那個心.
所以這個陳醫生.
就是這樣將他自己的生命.
在一個窮乏的地方燃燒.
這個豈不是後面作者.
所告訴我們的十二節.
他看到地上的人.
如果沒有神.
都是虛荒.
他面對的是極大的困苦.
但當他要來去看到.
神在他生命裡面的工作.
那個改變的時候.
他說我拿什麼去報答.
神在我的生命裡面.
祂賜給我的一切的厚恩.
他就是要舉起這個救恩的杯.

$^{481}$去清揚耶和華的名.
他有一句名句.
他說與其一個人.
這樣荒廢自己的人生.
慢慢的生鏽.
倒不如為神用盡他.
使別人得到建立.
生命的改變.
所以他用了他的人生.
出錢出力.
來幫助這些原住民的時候.
他的身體反而越來越好.
原來忙碌.
原來撲心撲命的幫助這些窮人.
自己反而得到更大的平安.
喜樂滿足.
這樣的方程式.
相信基督徒一個事奉神的人.
才會明白.
有時候很忙碌很忙碌的服事.
但完了之後.
內心裡面是有一份很大的滿足.
因為你是無私地服事這位神.
舉起的是救恩的杯.
去清揚耶和華的名.
所以這些告訴我們.
在第十四節.
我要在眾民的面前.
向耶和華還我的怨.
作者許過什麼怨呢.
他還什麼怨.
當然就是獻感恩祭.
猶太人他們會有一個平安祭.
就是感恩.
向神還怨.
保羅他這麼忙碌的宣教.
他上耶路撒冷的時候.
他就去找祭司還怨.
他是一個拿世爾人.
他剃頭.

$^{521}$他有一個怨.
是和神之間的.
弟兄姊妹我們在信主的人生.
我們有沒有向神禱告呢.
一生的求告神呢.
禱告神為了我們大家.
無論有事無事.
都一定不是例行公事.
是我們越來越渴望.
和神有禱告的生活.
更加重要就是.
我們人生有沒有一個向神許怨.
神啊你要成就我所期盼的.
如果成就我會還怨的.
這是很寶貴的.
不是交易.
是很寶貴的生命學習.
我們對神有所求.
神會成就在我們生命裡面的工作.
你就看到神在你生命裡面的改變.
你就會來向神還所許的怨.
接著十五到十八節.
在神的眼中.
看一個聖人的死.
是很昂貴.
不是寶貴那麼簡單.
很昂貴.
耶和華我真是你的僕人.
我是你的僕人.
是你婢女的兒子.
你已經解開了我的綁索.
又要求告耶和華的命.
耶和華的眼中看聖文的死.
是極其的昂貴寶貴.
到底是一個正面的意思.
還是負面的意思.
聖文怎樣死.
為甚麼會死.
作者好像不是講得很詳細.
但他立即說.

$^{561}$我就是你的兒女.
我就是你的僕人.
我就是屬於你的.
我曾經被纏綁.
但現在我已經被解套.
即是說我現在的內心是自由的.
很寶貴.
發現自己與神有這樣的關係.
神看聖文的死.
是極其的昂貴.
但作者告訴我們.
他就是聖文.
他就是神的僕人.
這樣將自己放在神的眼光裡.
這樣看自己的人生.
是一個自由的人生.
心靈自由的人生.
做的事是為神而做.
不再是為利益.
為別人.
可能過程是為別人.
但心裡是為神而做.
因為神已經解開了他的綁索.
親愛的弟兄姊妹.
一個教會最寶貴的地方.
就是我們看到每一個都是神的聖文.
每一個在神眼中都是寶貴昂貴.
每一個人的死.
每一個人的生.
都是在神的裡面.
即是說就算不知道什麼原因.
我們面對死亡.
在神的眼中極其的寶貴昂貴.
但我們在有生之年.
一起以這個身份來服侍神.
我們就是聖經新約所說的.
我們就是一隊好像軍隊一樣.
是福音的軍隊.
我們是有一個使命.
我們是有一個目標.

$^{601}$我們是有一個差遣.
是神要我們去做.
他安排要做的.
因為我們已經解開了這個綁索.
來呈獻給神.
因為我們就是屬於他的人.
所以你看到他有一個示範.
他說我要在耶路撒冷.
在他的殿裡面來還願.
到底是不是真的上了聖殿.
來獻祭一路頌唱呢.
我們就沒有這個背景.
不是太知道.
但是這裡很感受到.
作者他在神的選民當中.
來去重陽這一位的神.
在他的生命裡面的拯救.
和一而再提到這個求告裡面.
包含的是有一個還願.
神聽他祈禱.
神抬舉他的生命.
使他能夠深深發現.
自己是在神的選民的聖民的行列裡面.
這樣的一個猶太人.
一個這樣的心境.
親愛的兄弟姐妹.
我們是不是以感宣家,感宣人.
作為我們的一個誇耀.
我們能夠信主.
我們能夠求告神.
我們可以在我們的生活裡面去經歷.
一般人不求告神的人.
是不會經歷的.
那種平安.
那種內心的福樂.
那種心靈可以迎歸安樂.
這樣的一個心境.
這個就是.
牧者很需要不斷提醒.
不斷去教導弟兄姊妹.

$^{641}$我們真的要學習.
在我們的生活裡面.
感受到神同在.
我們一起去向神有更多的禱告.
這個禱告不單是停留在謝飯的祈禱.
而是在我們一個行動的.
我們的觸覺.
我們的思想.
我們所留意到的事情.
因為我們是有智慧的動物.
我們是應該看到世情.
化作我們的禱告.
不是應該看到二巴的戰爭.
只要罵完一句就算了.
是化作祈禱的.
是要有憤怒.
是要有悲哀.
是要有內心對那些人面對的困難.
那種的憂戚與共.
這樣的基督徒.
就是這篇詩篇告訴我們.
我們在這個月.
我們要學習感恩.
每天能不能夠有幾件事情.
我們可以學習.
自己學習感恩.
能不能夠常常在感恩裡面.
開始我們的生活.
這些都是我們大家一起.
需要學習.
很盼望基督徒不要再停留在.
只是謝飯的祈禱.
希望我們能夠有一個很大的觸覺.
心靈的觸覺.
我們在不同的事情.
隨時的多方的禱告.
和經歷到神聽祈禱.
所以我們要愛這位神.
我們一起祈禱.
讓我們今天晚上.

$^{681}$可以一同來到你的面前.
多謝你聽我們的禱告.
也讓我們能夠在真理詩篇裡面學習.
詩人真的能夠經驗到.
神聽祈禱.
他要一生的去求告.
他要在活人之地.
讓其他人看到.
他是一個有恩典經歷.
去還願的人生.
願主幫助我們.
我們這樣祈禱.
奉耶穌基督的名.
阿們.
\newpage



\section{哥林多後書使徒行傳 9:30-13:2-27:22-24}
\label{sec:EpQdyi4IyoI}
\textbf{2023-11-05 午堂《不同階段有不同學習:神永遠有奇妙恩典》使徒行傳 9\_30,13\_2,27\_22-24,哥林多後書 12\_9 講員:蔡少琪院長}
\newline
\newline
連結: \href{https://youtube.com/watch?v=EpQdyi4IyoI}{\texttt{https://youtube.com/watch?v=EpQdyi4IyoI}} ~~~~ 語音日期: 2023-11-08
\newline
\newline
\hyperref[sec:URxw02niNWY]{\small{< < < PREV SERMON < < <}}
~
\hyperref[sec:index_chronic]{\small{[返順時目]}}
~
\hyperref[sec:index_scriptual]{\small{[返順卷目]}}
~
\hyperref[sec:z09lg9RcA3o]{\small{> > > NEXT SERMON > > >}}
\newline
\newline
$^{1}$大家好, 感恩在錦索堂一齊敬拜神.
最近有兩個主題比較敏銳.
一個是發覺人在不同階段面對的問題.
要學的功課很不同.
我和一些人一樣.
這幾年開始可以坐兩元機.
兩元機發覺可以長周抵.
快船有二十多元一次.
兩元就抵.
但我只想香港庫房要扣多少錢.
不關我的事.
另一個是人生不同階段之外.
因為香港世界都面對很多逆風.
階段是一個順境還是逆境.
這個世代都有這個感受.
第二個觀念就是.
原來神在不同階段的時候.
祂都給我們不同的挑戰磨難.
但提醒就是神永遠有奇妙恩典.
這是一個主題.
世界很動盪.
大家知道有打仗.
二巴的情況.
兩邊都有人幫忙說話.
這些打仗我都不會參與討論.
是很複雜的事.
但總之都很血腥.
很多張力.
世界動盪.
有時在香港亞洲來說.
都不知道會不會戰火.
有時某日在香港發生或亞洲發生.
希望不會.
很多時代末世觀的感覺.
最近開始多個留意.
就是疫情後.
發覺其中一個擔心.
和一些做生意的第一次會談話.
似乎有部份都有感受.
就是經濟的浪都很深.

$^{41}$我自己就想起一句金句.
一起讀藍色字.
現在這地的饑荒已經二年了.
還有五年不能耕種不能收成.
曾經在前幾年都讀過這個經文.
已經覺得很難.
但發覺原來還有五年.
可能再有五年都不奇怪.
人生的事.
這個動盪期.
COVID後似乎都很動盪.
我自己比較關心中國的經濟.
似乎都聽聞很多公務員減薪.
甚至有些小地區.
連巴士都要出不了工資.
很複雜.
歐洲都是一個情況.
香港很多人移民去英國.
但英國找便宜工就容易.
我自己一代.
我兒女都剛剛讀Master的時期.
我女兒是COVID期間讀大學的.
COVID讀大學有一個特色.
就是很多時候都在家.
讀完之後.
我女兒讀書都不錯.
但讀完之後都沒有朋友.
還有老師都不認識她.
記得那時候去英國申請Master.
要老師幫她寫Reference.
老師就拖很久.
因為老師可能在Zoom裡面.
見了很多人都不認識她.
到寫Reference的時候.
都不知道寫什麼.
所以這個時代不同時代.
有很多不同類型的挑戰.
今天第一段就用了三段保羅的經歷.
讓我們看到人生.
其實他的遭遇的處境.

$^{81}$都略有不同.
保羅第一段經文.
可能大家很少詳細讀.
不過我們再讀一次.
第30次.
保羅就送他下蓋撒尼亞.
打發他往大掃去.
我溫溫保羅的人生.
保羅很小的時候.
在猶太人讀宗教信仰.
是很聰明的.
後來可能耶穌傳道的時候.
可能回了鄉下.
有機會耶穌升天後.
回來猶太人附近地區.
就迫不得教會.
大家熟悉吧.
在大馬塞與耶穌.
在大馬塞那三年很勤奮傳道.
這一段就是.
他在大馬塞被迫迫迫逃離.
去到耶路撒冷.
他想在耶路撒冷這麼大基地的地方.
去很長的時期.
誰知神說.
這個地方有人接納你.
你回鄉下吧.
大掃去鄉下.
這個時間.
大家沒有留意.
我們熟悉的保羅是從安提阿開始.
老統時份.
信主零年到三年.
在大馬塞.
三到十三年在他鄉下.
十三年在安提阿.
大概是這樣.
所以我們所熟悉的保羅.
多數是在安提阿之後.
大概這一十年.

$^{121}$保羅是寂寂無名.
講到這些挑戰的時候.
就像這裡說.
大概有信主後.
三到十三年.
是一個寂寂無名艱難期.
如果本來三十多歲信耶穌.
如果他是一個聰明人.
例如三十多歲聰明人.
像大律師.
什麼不同行業.
可能是哈佛.
清華 劍橋畢業.
如果你問你.
三十多歲到四十多歲.
要你很忙還是沒有工作好呢?.
如果聰明人.
要你十年沒有工作.
要你十年忙還是十年沒工作.
容易嗎?.
我懷疑十年有空.
是很艱難的.
這一代的風浪.
就是工作市場變了很多.
社會變了很多.
所以神給我們的挑戰.
有時候我們在想.
某個時代的挑戰.
但是聖經人物.
有些人物很厲害的人物.
摩西他有王子的學問.
但神讓他抗疫一年.
抗疫了多少年?.
四十年之後才容許.
問旁邊的人.
如果神讓你癱瘓十年到四十年.
你願意嗎?.
可能比忘十年更早.
對保羅來說.
可能沒有準備.

$^{161}$還有一個觀念.
保羅當時的猶太人.
很反對基督教.
有時候用木孫.
我喜歡用回教的字眼.
熟悉的字眼.
猶長的兒子信耶穌.
回到自己的鄉下.
回教的地方.
見鄉親.
你覺得鄉親是歡迎他們.
還是想把他們扔死?.
可能把他們扔死.
所以這十年.
其實是保羅很不容易的時期.
不是簡單的時期.
但聖經沒有很多記載.
有一個記載是中間這一段.
下面那段就說到.
巴拿巴去大掃找保羅.
已經到了.
大概信主十三年後的遠.
中間一段就是.
在臨後十一章.
二十四節第四行.
讀.
被猶太人鞭打五次.
每次四十.
減去一下.
即是.
狂打了五次.
保羅試過被官打過.
但這裡說猶太人打.
就好像中文傳統一樣.
你不長進.
或者鄉親不喜歡你.
在祠堂打你.
你明白嗎?.
所以學者認為.
可能多數是.

$^{201}$他回到鄉下大掃期.
被鄉里打.
或者父親.
所以這裡似乎.
都記載.
八至十年是一個不容易的時候.
很艱難的時候.
人生就是這樣.
我們想的人生.
是一路向上的人生.
還是平的人生.
還是向下的人生.
香港以前的社會.
都是很多時候向上.
香港六六六七大暴動.
暴動幾年又向上.
那時賣樓的.
那些臭小子.
我自己讀書.
我是港大八一到八四.
我讀師父的.
是師父的武功.
八三八幾就帶翠姨去北京.
第一次.
香港就倒了一陣子.
後來才有聯儲.
聯繫匯率.
我畢業那年是師父武功.
過了一年才有功.
所以後來就做了貿易.
後來做中國貿易.
又有另一個機遇.
所以不同時代.
那個順逆.
香港是一逆風就上去.
有些觀念.
一逆風就上去.
但很多國家不是.
有些非洲國家一借就借十年二十年.
有些借四五十年.

$^{241}$中國歷史很厲害.
中國歷史以前出名的.
王朝首都.
開封.
今天開封厲害嗎.
可能普通.
中國王朝曾經是西安.
西安厲害嗎.
普通.
但每個地方都向上.
香港未來當然.
我們不討論向上還是橫.
人生就是這樣.
不一定向上.
有時是低.
每個時代看東西都不同.
舉例.
我自己經歷八年代.
那時對中國貿易很負擔.
後來就白五入中國貿易.
那時中國貿易騰飛.
兩三年裡搭飛機搭了一百次.
那時很年輕.
有個美容師請我們幾個經理.
那時香港.
如果記得八年代.
利息很高的時候.
那時他請我們經理.
我們做了兩年生意很好.
他就加我人工.
每人加十幾\%.
我們四個就決定集體辭職.
最後公司就玩我們.
我不知道他加其他人又如何.
就加我25\%.
那時我都不知.
我都不知.
那年代是加薪很猛的時候.
上海十年前.
加薪很猛現在就減人工.

$^{281}$人生裡都有高評落的時候.
我們基督徒經歷神.
向上的時候.
有什麼?.
打橫的時候有什麼?.
向落的時候去經歷神.
我做教育.
美國都進入一個.
某些大學.
基督教大學倒閉潮.
因為疫症下.
小孩少等等.
香港最近有些幼稚園倒閉潮.
因為移民潮.
香港很多Band 1的學校.
收不到以前Band 1的學生.
收到Band 2的學生.
因為有些中產移民了.
很多動盪.
在宣道會為例.
宣道會以前.
一個很出名宣教的大學.
叫Lyon.
NYACK在紐約.
但過去十年經營不算很好.
最近十年經濟不好.
最後倒閉了.
所以提醒一件事.
就是時代有時很多浪.
不是說你負責的機構.
負責事業.
負責教會.
一定是很風光.
我們熟悉的是.
一百多年前的宣教士.
大部分在八五年.
二百年前的都是英國.
但英國今天的西洋教會.
人多還是人少.
人少.

$^{321}$所以我們有時時代.
有不同的挑戰.
前陣子政府想香港夜繽紛.
但我發覺.
香港晚上夜宵潮多.
因為前陣子.
我的熟悉的親戚在加拿大.
回不來.
我陪他去掃墓.
剛剛沒事做.
去了伯爵林.
我家在那邊.
約了幾個親戚.
那裡我見到在尖沙咀吃午餐.
我很少在那些日子.
去尖沙咀.
早到.
我去了尖沙咀一個商場.
叫I SQUARE.
我早到半個小時.
帶他們去I SQUARE.
我發覺十二時多的時候.
都沒有人去買東西.
買東西的地方都沒有人.
我都做經營.
想想老闆買東西都沒有人.
都有點心寒.
因為我有經歷.
因為我那區在疫症前.
開了一個比較漂亮的西餐廳.
都頗大.
坐位好像一半的禮堂.
但每次去.
我都一個人坐.
一個起司蛋糕加一個好咖啡.
五十多元加一.
我最喜歡坐在那裡.
因為我怕去到靜.
有人認到我 坐牧師叫來叫去.
希望見到人.

$^{361}$那個地方特別好 經常都見不到人.
每次都一兩個.
但其實我很擔心老闆會走下.
捱了幾年.
都覺得不走下都很奇怪.
怎樣捱呢.
最後一次我再去.
又坐下 因為沒有人招呼我.
看見廚房黑漆漆.
奇怪了 老闆走過來.
我說今天不開業.
我又不開業 又關門.
我心裡說.
我都是算熟客.
都很忙.
算了 我又不喜歡爭拗.
就出去.
過了幾天再去 吉特里封舖.
原來欠了一年租.
差不多.
其實我預計的事.
真的發生了 根本是很艱難.
疫情下辦這些餐廳.
夜宵調.
真的不容易.
我上星期.
去了一個教會港島.
剛好他們有些弟兄姊妹做玩具業.
我跟他們談了兩句.
我心想中美吵架.
跟電子業有影響.
玩具可能少一點.
他說都不是.
很多美國商都說.
反正吵架不如搬到墨西哥.
少了在中國樂丹.
我們行業都少了幾成.
我其實不太懂.
我只是談一個.
我沒有預計.

$^{401}$總之一個觀念就是.
人生裡面有些階段.
可能好像約瑟說的.
方年可以是兩年三年五年.
你控制不到.
教會都是因為移民潮.
教會有些老化.
年輕人其實這一代.
找便宜公義 找好功能.
很多的挑戰.
所以回到保羅.
保羅其實是沒有想過.
會面對十年的沉默期.
因為保羅是一個行動型的保羅.
在大馬錫信主.
他就在國會堂.
講耶穌.
說他是神的兒子.
很特別在大馬錫被逼迫的時候.
特別講一句說話.
紅色字的四個字是道.
他的門徒.
即是說保羅在大馬錫.
三年已經訓練了自己的團隊.
有保羅的門徒出現了.
他做師父了 徒弟了.
如果一個人在教會事後三年.
有一群人和他合作.
那些人算不算做得勤力.
這個就是保羅.
所以保羅沒有想過.
回鄉下.
但聖經記載在.
四行尊二十二章.
重溫那段經文.
主就說要你回鄉下.
保羅就講一段很長的說話.
我很多見證.
在耶路撒冷才做才好.
耶穌說我不回去了.

$^{441}$主就很硬的說一句.
下面一段的藍色字.
你去吧 一起讀吧.
你去吧.
我要差你遠遠地.
往外幫人拿女兒去.
所以很明顯.
保羅是不想去的.
這十年不是保羅爭回來的.
你問一個比較行動型的.
小時候讀書好的人.
叫他十年安靜.
還是十年勤力.
我懷疑他都想十年勤力.
今天中國那些年青人.
要躺平.
其實找不到工作才躺平.
躺平是很艱難的經歷.
所以.
聖經就在這十年.
記載的不多.
但其中有意思的就是.
在哥林多後書記載.
聖經說他前十四年.
被提到.
第三層天去.
記得這個經歷嗎.
保羅死在臨後.
大概信主後22年.
算回.
十四年前.
22減14.
就是幾年.
18年.
再說一次 0到3在大馬色.
大概3到13.
就在鄉下大掃.
所以.
第8年就在大掃期.
所以有一次.

$^{481}$我默想之後就開始明白道理.
第三層天.
可能很榮耀的經歷.
神不是當保羅順境的時候.
給他的.
可能他倒下.
埋怨神的時候給他的.
神呼召我去.
我離開家.
有很多前途都不去.
接著去事奉你.
接著放我十年不工作.
十年沒發揮.
那些人會埋怨嗎.
會埋怨.
保羅大公的光景裡.
所以聖經提醒.
原來神給我們.
很特別的經歷.
未必是我們最成功的時候.
有時神給我們很特別的經歷.
是我們最無能無助的時候.
這是一個安慰.
神特別.
看你的時候.
可能是你最不行的時候.
神看你.
希望你不用那麼不行.
神看你.
不一定是你最不行的時候.
神的恩典.
不是我們順境很厲害.
捐一千萬捐一百萬.
給教會的時候才看你.
是你最不行的時候去看你.
聖經講的特別.
我自己小時候.
在黃大仙寨子居.
七層樓長大.
那時爸爸在我.

$^{521}$零到三四歲的時候.
事業比較順利.
買了自己的私宅有個工廠.
但三到四歲的時候.
被人騙了一點錢.
周轉不靈就破了產.
所以家裡由私宅搬到.
寫字區.
因為我很小就沒有私宅的印象.
因為我四五歲就在寫字區.
所以我性格就像寫字區.
風格.
我喜歡穿拖鞋短褲走街.
因為好像小時候那種氣質.
寫字區如果大家經歷過.
那段階段.
是很骯髒的.
很多老鼠走來走去.
所以我那時候曾經小便兩條蟲那麼長.
杜蟲.
花塔餅那些.
你聽過沒有.
我那區那時候.
治安沒有那麼好.
試過婦女被人箍頸.
我有個親戚女士被人箍頸.
其實我小時候都不知道她箍成怎樣.
總之見到紅了.
我媽媽就命令我做一件事.
那時候我大約八歲.
當我姐姐去中間公眾浴室洗澡的時候.
叫我站在門口.
其實我只有八歲.
那一代就是這樣走過.
我們家裡連奶奶十個.
只有奶奶早過世.
就九個.
所以一間長屋就九個人一起睡.
飯煲床木板床.
所以那一代的人去神后院讀書.

$^{561}$由我的為例.
九人房搬了四人房.
但今天不是了.
今天有個弟兄姊妹.
有時我們見到有兩人房就要祈禱很久才進去.
因為可能自小都有人房.
有些住屋村都可能有人房.
所以有時我們年長一點都明白.
大家心態不同.
很難說誰捱苦不捱苦.
只是處境不同.
有時會收了去.
等等的情況.
其實說真的今天到我.
我出差都不喜歡兩人房.
因為我喜歡自己靜.
那個是我幾十年前的我.
我選了那一次.
要我回到那個我都未必可以.
那一次高興一點.
說那十次未必可以.
總之不同時代有不同的遭遇.
甚至是用遭遇去磨練我們.
好像摩西四十年抗疫.
我自己在讀PhD的時候.
在美國留學的時候.
神給我一個例子提醒我.
無論我們是謙卑卑微都好.
我們都可以與神同行.
那時我們在一個小的教會.
國語教會.
有些東南亞人.
有些台灣人.
有些大陸人.
香港人.
有四十人的教會.
我們那時教會有少少增長.
我們就希望接觸當地的華人.
吸引他們回教會.
借了一個教會.

$^{601}$新人教會很靚.
他們一個地毯的hall.
大概大家這個hall多一格.
即是三格都很大.
我們借完之後.
全部地毯.
所以我們很怕他們.
不多來教會的人.
弄髒地下.
所以我和執事會主席有個責任.
就是整場都不做台上的工作.
就拿著垃圾袋.
扔在這裡.
因為怕弄髒我們借人.
所以我那天就.
自己站著就拿著垃圾.
開始.
有一個家庭是大陸人來的.
他就來了教會兩三次.
他以為教會甚麼都可以送給他.
那個小朋友五六歲.
就見到台上抽獎禮物.
有個木公仔 大獎 細獎.
他就不留意大獎.
留意細獎 木公仔.
Kevin 我問Kevin.
木公仔.
平時教會有些小禮物.
給你五歲 但這個抽獎.
那時候我順口說一句.
你祈禱一下.
說完之後我有些後悔.
因為半跌小孩子是很重罪的.
後來我才發現原來他沒有買抽獎券.
那時候很忙 要處理.
接待客人 撿垃圾.
很忙 進了廚房.
走來走去.
那九個字裡.
突然在台上就叫Kevin.

$^{641}$原來我抽到獎.
不是大獎.
是個木公仔.
其實在廚房.
我看不到台.
我都猜到多數抽到獎.
真的抽到獎.
上台拿了獎.
我就看都不看 放下給小孩子.
就進廚房開心地撿垃圾.
因為我學到一個神學.
就是撿垃圾都與神同工.
這神秘有時候.
未必是順風的學習.
有時候是逆風裡.
學習與神同行的故事.
所以保羅都是.
神預備他的時候.
在他軟弱裡學到這個功課.
所以保羅.
有一個名句.
下面那個藍色字是什麼.
我的恩典救你.
很多人都很喜歡這句.
恩典救庸.
特別當我們開始.
年長了.
年老了.
慢了.
你會發覺用力去示奉神都好.
但慢慢都沒什麼好力.
我記得有一次.
去一個較大的教會聚會.
我經常去.
他們有一個老弟兄很熱心.
你來到他就經常握手.
做私事.
他沒有在門口接待.
坐在後排.
因為認到他.

$^{681}$就走過去.
他很笑容的.
他說不知道.
因為我前兩個月.
做了腳的手術.
站不久.
所以就坐下.
我覺得他有一種不開心.
人看他.
好像一個.
簡單不深的示奉.
他喜歡這樣服侍神.
但.
當不行的時候.
可能這個示奉都拿走了.
所以人生的事.
有時要知道.
不同光景都經歷神.
神的恩典是救庸.
所以我也是教.
教神學生.
今年的口號叫.
1899.
是命不留.
1899年成立.
明年就是我們125週年.
宣傳會在中國歷史.
都差不多那麼長.
所以1899之後.
中國有沒有經歷.
軍閥割據.
清朝變中華民國.
有沒有.
第一次世界大戰.
有沒有國共內戰.
二戰.
後來中國.
又有文革.
又有其他.
我想提同學.

$^{721}$我們經歷的遭遇.
我們真的不知道前景是順與逆.
因為香港習慣順.
但世界歷史.
真的不知道順與逆.
但無論順與逆.
神都永遠與他同在.
無論順與逆.
神永遠有恩典.
我們勇敢去.
示奉神.
這樣提他們.
在香港時代有個出名的.
演講叫劉福群牧師.
他當時.
經歷日本侵華的時候.
那時日本在北方先打.
打到南方的時候.
廣州淪陷.
當時見到在廣西.
吳州.
被戰火攻擊.
他逃難去到.
越南邊境.
當時戰火無情.
其中一個牧師的太太.
太師母就死了.
他們逃到山區.
叫復興村.
那些.
少數民族山區.
那些好像苗寨.
那些山區.
我看那幅相.
很特別的是當時有軍訓.
如果你今天問我.
見到有沒有軍訓.
如果沒有這個歷史.
我一定沒有.
處境就是這樣.

$^{761}$很深刻的是我重溫這個歷史的時候.
就是當他.
劉福群牧師逃難去到.
近越南的時候.
他覺得走投無路.
曾經有決定將建道.
帶去轉入越南.
因為當時由北方南方.
已經是日本侵華.
怎知準備去的時候.
就是39年.
9月希特拉.
在歐洲攻入波蘭.
世界大戰.
發生.
法國怕在這邊和日本開戰.
封了邊境.
不讓中國人過境.
結果這個建道小故事被迫.
進不了越南.
停辦了幾年後來不知不覺.
來了香港.
如果當年去越南.
就沒有了蔡牧師.
曾牧師的校友.
你明白嗎.
或者沒有了宣導會.
所以整個主權就很有趣.
上一代其實都在風浪裡面.
嘗試找路.
最後神封了門.
結果.
開了另一個門.
其實香港的門.
或者我們個人的門.
家庭的門我們都不知道.
上一代都可能會祈禱.
越南那時有宣導會.
有校友在.
去越南就應該順利.

$^{801}$世是這樣的.
神封了門就封了門.
人生有很多這樣的故事.
第二段和大家說.
就是亡祿奇.
我試讀回經文.
十三章第二節.
他們事奉主.
禁食的時候.
聖靈說要為我.
分派巴拿巴和蘇洛.
去做我照他們.
所做.
最近就默想這一段.
大家最熟悉的經文.
後來去馬其頓異象.
哥林多教會 雅典.
也就是這一段.
但其實這段時間.
只有八九年.
我自己就.
八九蒙召.
是很想去中國事奉.
現在敏感時代.
我就做交流.
因為以前用的字眼比較敏感.
有一次想交流.
有人想我去分享異採書.
那時候我又怕過關敏感.
很早期 九零年代.
我就把我自己.
看開的聖經.
就撕剩個異採書.
就好像很薄的.
過關 希望沒有人留意.
誰知我過完關之後.
他說那個地方被人封了.
結果我就在那裡.
戇居居幾天.
我就撕爛了我的聖經.

$^{841}$又一事無成.
上去都做不到什麼.
記得有一年.
有人帶我去探訪.
去到一個溫州的山頂位.
開車九個小時.
去到郊外.
走了半小時.
上去一個小山頂.
我說你們為什麼在這裡聚會.
這麼僻.
他說原來有典故.
因為當文革的時候.
我們前輩就走到這個山頂位.
就祈禱.
為什麼這個祈禱好呢.
因為通常.
關心他們的人就不肯行山.
所以多數是安全一點.
你無異議就明白了.
走到僻的.
我就明白他們的傳統.
為什麼要走到僻的地方祈禱.
或者聚會就是一個屬靈傳統.
你明白嗎.
但是幾年前.
拆十字架.
香港事件.
其實現在探訪中國都比較敏感.
就是福音工作.
那幾年有機會去探訪.
你們都會.
大陸交流很多.
但現在都相對少了.
所以有些時間.
好像這段很忙的期.
但未必永恆.
保羅這些很忙的三次行程.
大概八年九年.
突然靜了.

$^{881}$所以有時一個提醒.
就是.
當你很忙的時候.
要學習感恩.
忙的時候都很辛苦.
讀到這段經文黑色字.
除了者外面的事.
還有眾教會關心的事.
天天壓在我身上.
有誰軟弱.
我不.
有很多底子會熱心.
這裡做生意很忙.
忙完之後教會的事奉很忙.
辦個機構很忙.
有些時候真的很忙.
但有些時候神不理.
太靜可能你懷緬.
很忙的時候.
好像這位很出名的宣道會領袖.
見到錢長澤富文.
如果你看網上.
教會在越南開始.
就是福音由他帶到越南.
在很忙的時候.
他也遭遇人生一個.
隱藏的困難.
就是他有一段時間.
太太憂鬱症.
有時人生好像這幾年.
香港人憂鬱症都多了.
壓力多了.
有些人會用名句.
讀一讀黑藍色字.
從前他是神給我一隻右手.
現在他是神給我一個十字架.
人生裡面有時真的有順境期.
有逆境期.
所以西方形容.
有些艱難期叫Happy Problem.

$^{921}$即是你做生意還是賺到錢的.
頭髮白了點.
忙一點.
你想說頭髮白點.
忙一點還是每天.
等倒閉的時候好一點.
可能做生意的人覺得.
都是忙一點.
比等破產.
等倒閉的時候好一點.
你明白嗎.
跟別人說.
有得忙應該高興.
有得事奉有得忙是感恩.
有個教會.
有一次我遇到一個老姐妹.
她家裡算富裕.
都捐錢給教會.
她教會的師班樓梯比較多.
一級級上.
比較窄.
有一年她腳軟.
站不住.
要坐輪椅.
最後她決定不唱師班.
對一位唱了幾十年的師班的姐妹.
對她都很大打擊.
她很喜歡事奉.
很喜樂的事奉.
但身體不佳就做不到.
所以人生裡面.
無論我們自己家人相聚.
或者生意教會事奉.
有些時候當你很忙.
仍然有服侍的能力.
我們應該學習感恩.
或者明天我站不住.
我就要坐後面.
人生裡面.
有些事情.

$^{961}$我們應該坐後面.
人生就是這樣.
忙裡面.
有忙的掙扎.
這幅相是我以前在城鎮主任學師時.
拍的.
那時很忙.
一年講五六十篇.
有很多會議.
人事處理.
個人真的很累.
講完都知道.
下星期做兩三篇都不奇怪.
所以我習慣逢星期一.
在九龍城區經常遇到人.
有時想靜下來.
就拿著這麼僻的石澳.
希望遇不到人.
那段時間我喜歡石澳在海灘.
走的安靜.
有一次近聖誕.
很忙的.
我去靜下來.
就像穿短褲拖鞋.
走水邊.
散步默想.
換了衣服就留下了.
荷包在公寓室.
回來還淋了.
神啊我還求你講信息.
連荷包都不見了.
簽證報失後身份證還要處理.
我煩到嘔.
你幫我找回來.
我頂不住了.
我這樣吐了四天都沒有效果.
我打算再吐一兩天都認命.
但那時很不高興.
神啊我都覺得很不妙.
你還加個難處給我.

$^{1001}$吐了四天.
吐的方法未必很合理.
但我還是很開心.
吐到第四五天.
旺角太子警官打電話給我.
你是否蔡少奇.
你的錢包在旺角警署.
我到今天都不明白.
為什麼錢包可以從石澳運到旺角.
可能那個人拿到之後.
要祈禱很久才放到警署.
但那時對一個.
很不高興的人.
我都覺得很不妙.
但那時對一個很忙的我.
是一種安慰.
我這麼不高興.
我都勉強來掌摑神.
神都不離不棄.
這不是正確的祈禱觀.
但這亦是我們人生裡面.
就是這樣.
有些事真的不行.
你就去找神.
所以我想神的憐憫就是這樣.
祂不同階段都可以有不同階段.
去幫我們.
第三階段就是.
今日讀的經文就是沉船期.
保羅有兩年.
無端幫人捐錢.
被人都要坐監.
上訴該殺.
誰知.
黑運重重.
坐船都沉船.
沉完船還要被毒蛇咬.
你見到嗎.
就是一路黑下去.
這個經文裡面.

$^{1041}$最近我是很感動.
一起讀讀這個紅色字.
保羅不要害怕.
你必定站在該殺面前.
並且和你同船的人.
神都賜給你.
這個時候我特別喜歡這一段.
為什麼呢.
因為我開始有兩元坐車的時段.
都有些開始.
六五歲又退休.
準備就繼續.
有些計算倒數.
某些日子.
我小時候都很喜歡快.
和認叻.
我以前其中一個叻.
是數學算叻.
因為我去一個名校讀書.
我那年中四的名校讀書.
我們那個數學組.
是拿全港數學冠軍.
那時候我都算是招職數學的人.
中六的時候.
有個老師教的數是差一點.
那時候我讀書.
教的數是差一點.
有一次他自己弄一條很深的數.
又不懂得解.
不懂得解就很普通.
有時候我們都不懂得解.
最蠢的就是不懂得處理.
他不懂得解.
應該說專門給你一條深數.
你逐個做下.
他又有面子.
不懂得解又不懂得轉頭.
面壁半個小時.
然後都解不到.
那時候整堂.

$^{1081}$那些男生在後面.
老是這樣.
那時候我和另一個同學算叻.
現在另一個就做了.
新加坡大學數學教授.
那時候我心想.
其實很容易.
你找我或他出去就搞定.
那時候我是招職時代.
但現在不是.
現在自己都頓了.
都老了.
皺皮了.
講什麼都勾.
年輕人興什麼都不懂.
潮語都不懂.
現在我最近就喜歡這個祈禱.
什麼祈禱呢.
就是當你又坐牢.
又失敗.
又絕望又沉船.
都可以救一整船人.
慢一點.
當你什麼都沒有.
什麼都失敗.
仍然可以幫一整船人.
這個厲害一點還是你很叻的時候幫一整船人.
我起碼在這個階段.
我喜歡這個.
即是說.
我就算坐輪椅.
我在病床裡面.
或者我腦裡開始有些痴呆症.
原來神都可以用我的見證.
幫到一整船人.
所以我就想到兩句說話.
一句說話就是.
死定都有出路.
第二就是.
無能的時候都可以幫很多人.

$^{1121}$跟別人說.
死定都有出路.
不行都可以幫很多人.
我們經常以為很行才幫人.
這個不是聖神號.
我們這個雅氣.
我們幫到人的時候.
顯出了能力.
是出於什麼.
神不是出於我們.
保羅就說.
恩典救用.
在軟弱裡面顯出大能.
所以我很喜歡默想這一句.
死定都有出路.
死定都有幫人.
所以沉船.
在經歷裡面.
我們得救的指望都絕了.
所以一個做囚犯的人.
參加死囚犯.
沉船.
被毒蛇咬.
經歷連串的.
黑運,厄運.
竟然救到很多人.
就見證我們有的是真神.
我們所存的真是福音.
所以最近很喜歡默想這一段.
年輕的時候.
靠自己去衝.
努力,智慧.
捐錢去事奉神.
這是很美麗的.
但有些階段是一事無成.
或者有很多限制.
但原來那個限制的時候.
神都能夠用我們.
都與我們同行.
同歸於盡.

$^{1161}$我們都能夠用我們.
都與我們同行.
同供.
我覺得是更美麗.
宣導會有一個歷史.
就是第一個來中國宣教士.
就是賈錫德.
有一幅宣導會.
香港宣導會的照片.
是把不同時代的僕人.
弄成一張.
他坐船.
積極傳福音.
但染了病.
未到中國.
日本死了.
所以這個歷史.
就像中國的名句.
出師未至.
身先死.
但當我們再追故事.
來說就很感動.
就是他太太坐第二班船.
繼續帶七歲.
和四歲的女兒去中國.
這三個姐妹.
將一生埋在中國.
好像他大女兒.
和他太太都撞在中國.
而他二女老年.
回美國就撞在美國.
後來我再追故事.
就找到.
他在美國的墓碑.
我感動的就是.
他的墓碑寫了一句話.
.
我其實不太認識這位姐妹.
但爸爸.
去傳福音.

$^{1201}$去宣教.
竟然他們家不放啟示奉.
容易嗎?.
屬靈來說.
厲害嗎?.
還一生去事奉.
.
如果我們的毛志銘.
有些句都很感恩.
(笑).
幾十年裡面.
他們經歷不同的風浪.
但他一生.
事奉著.
(笑).
見到.
剛才講的例子就是本來去越南的.
後來再上任的院長.
他在印尼集中死了.
就是習富民.
他留在長洲.
上面的別墅.
就成為香港的故事.
開始.
神有奇妙的預備.
我們計不到的.
神封了這個門.
就可以另一個門.
我們真的計不到.
後來很多畢業生都不講了.
我最後結束用了.
四行傳最後一個捷經文.
粒粒的中文字一起讀.
帶著.
所有勇氣.
無人能.
禁止.
四行傳28章是很特別的.
他沒有講保羅有沒有出獄.
後來我們知道出獄了.

$^{1241}$後來再殉道.
但那裡沒有講.
我覺得路加用了一種表達.
就好像華懷仁.
的一種講法叫留白.
留一些空間.
給你去想事情.
所以有些西人的牧師就說.
就好像神提醒我們.
四行傳寫了28章.
第29章是你參與寫的.
28章是你望人寫.
第29章是你寫的.
所以我在見道和學生講.
其實百多年裡面.
香港見道世界都很多故事.
神的恩典不離不棄.
我們經歷順或逆.
但關鍵順與逆都可以與神同工.
順與逆都可以經歷神.
第二就齊了.
你是見道的29章.
你是神的29章.
如果沒有29章.
這個機構教會就完蛋了.
如果感受堂沒有29章.
那就沒有未來.
家庭家族沒有29章.
以前多風光都玩完.
所以我就提他們成為29章.
那你提你旁邊的.
你要寫你們29章.
不要退縮.
不要閉關.
我們都是29章.
你就是29章.
其實聖經就是這樣.
聖經特別在上傳尾.
用了很長篇幅講保羅怎樣沉船.
其實回頭看不是偶然.

$^{1281}$好像預備一個訊息的總結.
就算我們經歷沉船.
政治的逼迫.
我們都要承受.
神的恩典仍然是不離不棄.
甚至我們沉船.
原來還可以救船人.
所以提醒基督徒.
帶著所有勇氣去試夢神.
沒有人能夠禁止神的過渡.
所以同教會說.
帶著所有勇氣.
無人能夠禁止你是第29章.
三句好不好.
帶著所有勇氣.
無人能夠禁止你是第29章.
這個就是今天的訊息.
我們一起重新祈禱.
天父我們人生或長或短.
可能都經歷順與逆.
但今日香港世界.
逆的風比較猛烈.
我們學習不同逆風的時候.
仍然不會灰心.
仍然珍惜你給我們空間去服侍你.
我們或許會疲倦.
或許有憂傷.
但仍然我們是你的僕人.
我們要抱捨你的29章.
或許我們遭遇風浪.
我們遭遇沉船.
遭遇病患.
但我們仍然可以與你同工.
最絕望的時候.
原來都可以幫助很多人.
求你使用我們每一個.
亦使用錦繡堂的弟兄姊妹.
祈禱加多過萬次耶穌.
得勝冥求.
\newpage



\section{撒迦利亞書}
\label{sec:z09lg9RcA3o}
\textbf{2023-11-09 周四晚查經班《歸回先知書研讀(8)撒迦利亞書》 講員:曾裔貴牧師}
\newline
\newline
連結: \href{https://youtube.com/watch?v=z09lg9RcA3o}{\texttt{https://youtube.com/watch?v=z09lg9RcA3o}} ~~~~ 語音日期: 2023-11-11
\newline
\newline
\hyperref[sec:EpQdyi4IyoI]{\small{< < < PREV SERMON < < <}}
~
\hyperref[sec:index_chronic]{\small{[返順時目]}}
~
\hyperref[sec:index_scriptual]{\small{[返順卷目]}}
~
\hyperref[sec:Jspoigh_Nw0]{\small{> > > NEXT SERMON > > >}}
\newline
\newline
$^{1}$各位晚安,我們開始今晚簡短的茶經班.
多謝天父的恩典,我們一齊來到第八課的撒加尼亞書.
我們很盼望,我們一直看這些無論是異象.
或者是對於將來以色列局勢的預言.
一方面我們不是為了好奇,而是神你自己.
有對你的民族,對你的子民的管教,審判.
和牧羊,尤其是今晚我們看到.
好幾段經文都說到神是那一位好的牧者.
和你會興起在以色列當中有好的牧羊人.
好像撒加尼亞,本來祂是一個先知.
但是祂要扮演著一個牧羊人的身份,角色.
來去預表將來耶穌基督的出現.
所以我們都盼望藉著這些經文.
讓我們明白將來的世界.
也讓我們看到神的作為.
在你自己的子民當中,求主帶領我們.
我們這樣仰望禱告,奉耶穌基督的名.
阿們.
開始第十章,其實是第十四章.
有比較難解的經文.
當你去到裡面,你就會發現.
有些好像突然間加插了一些戲劇.
或者是一些可能是寓言.
所以我們要,是有困難的.
但是總的來說.
牧羊人這個主題都是很清楚的.
變成以色列這個民族.
不是短暫的目前,是很遙遠的將來.
神是會有那個工作.
特別是去到第十四章,就天下太平了.
是一個很榮耀的景況.
所以有時候我們讀一下舊約的先知書.
都給我們有很大的振奮.
我們上一次講到第十章的第二節.
我們由頭再讀起,十章開始.
希望今天能夠講到兩章至三章.
因為十一,十二章都有一些關聯.
第十章開始.
「當春雨的時候,你們要向發閃電的耶和華求雨.
他必為眾人降下甘霖,使到田園生長才疏.

$^{41}$因為家臣所預言的是虛空,卜士所見的是虛假.
作夢者所說的是假夢,他們白白的安慰人.
所以眾人如羊流離」.
這是群羊的寫照.
因沒有了牧人就受苦了.
「我的怒氣向牧人發作,我必懲罰公山羊.
因我萬軍之耶和華,眷顧自己的羊群.
就是由大家,必使他們如進馬在陣上.
防割石,釘子,征戰的弓和一切的掌權都從他而出.
他們必如勇士在陣上,將仇敵踐踏在街上的泥土中.
他們必征戰,因為耶和華與他們同在.
騎馬的也必羞愧.
我要堅固由大家,拯救弱失家,要令他們歸回」.
其實已經歸回了,這本書寫作的時間.
就是大理烏王的第二年九月.
一連串八個的異象.
然後由第九章開始.
一直說到將來以色列這個民族.
他們的興衰,他們的福祉等等.
所以其實已經歸回了.
這些猶太人已經在所羅巴伯的帶領下歸回.
所以這裡肯定是在說一個很遙遠的.
真真正正的全國以色列的復興的時候.
「我要連戌他們,他們必像未曾棄住的一樣.
都因我是耶和華他們的神.
我必應允他們的禱告.
以法連人,必如勇士」.
又說弱失家,又說以法連,或者嗎哪西.
因為弱失兩個兒子就是嗎哪西和以法連.
以法連是相對在北面的支派來說是比較大的支派.
所以很多時候一個統稱就是猶大家.
就是南面,北面就是以法連家.
「他們心中暢快如同喝酒.
他們的兒女必漢見而快活.
他們的心必因耶和華喜樂.
我要化師聲聚集他們」.
好像那些獅子叫的聲音.
「因我已經救贖他們.
他們的人數必加增.
如從前加增一樣.

$^{81}$我雖然播散他們在列國中.
他們必在遠方紀念我.
他們與兒女都必存活,且得歸回.
我必再領他們出埃及地」.
這是再次的出埃及,即是第二次的拯救.
第一次的拯救就是以色列人成為奴隸.
在埃及,神就揀選呼召摩西.
回去埃及拯救他的同胞.
用十災,現在這個是新的出埃及.
其實在《易卓雅書》也有提過新的出埃及.
現在這裡就更加明顯一點.
告訴我們一個再站身的神的拯救.
使這個民族好像釣魚一樣.
去到不同的地方,將他們釣回來.
使整個以色列真正的復歸於一.
全家歸主.
第十節的中間.
「焦聚他們出亞述.
令他們到基列和黎巴嫩.
這地尚且不夠他們居住.
耶和華必經過苦海,擊打海浪.
使到尼羅河的深處都枯乾.
亞述的驕傲必至卑微.
埃及的權柄必然滅沒.
我必使到他們倚靠我.
得以堅固.
一舉一動,必奉我的名.
這是耶和華說的」.
或者我們可能讀完十一章一點點.
因為後面說到有些牧人的.
到底是一個比喻.
還是一個好像醫學語言的語言呢?.
還是一個戲劇呢?.
有不同的角色呢?.
因為有買賣雙方的.
十一章提到.
有些賣家買羊的.
也有牧羊人的.
有用不同的飼料餵養群羊的.
撒迦利亞就要扮演那個角色.

$^{121}$就是好牧人的角色.
所以是有些難解釋.
到底是想暗示什麼呢?.
十一章的第一節.
「黎巴嫩啊,開開你的門.
任火消滅你的香柏樹」.
黎巴嫩最出名就是香柏樹.
高大的,木質又漂亮.
聖殿就是用黎巴嫩的香柏樹去做的.
即是裡面的內部裝修.
「松樹啊,應當哀號.
因為香柏樹傾倒.
街尾的樹毀壞.
巴山的橡樹應當哀號.
因為茂盛的樹林已經倒了.
聽啊,有牧人哀號的聲音.
因他們榮華的草場毀壞了.
有哨狀獅子咆哮的聲音.
因約旦河旁的叢林荒廢了」.
到底是在說一個嚴重的自然災害.
還是周邊的國家.
是他們以他們自己的出產.
以他們所誇耀的過去誇耀.
現在因神的介入.
這些周邊的國家.
他們似這些自然的資源.
因失收,似乎有一場很大的火.
甚麼都沒有了.
即是誇耀的全部都化為烏有.
這些在目前都很難看得到.
應驗是按照著這樣的描述.
但如果這樣描述就真是很奇妙了.
因周邊的國家正正就是以色列的敵人.
正正就是一直都在他們周邊成為以色列的敵人.
或者我們可能因為說到關於牧人這個主題.
都給大家有一個機會認識一下.
其實很難解的,這一段第四節開始.
一直到十七節.
「耶和說:我的神如此說:你撒迦利亞要牧養者張載的群羊」.
即是說牠是那個牧羊人.

$^{161}$如果當這個是一個戲劇.
現在那個角色就有個劇本.
撒迦利亞就是那個牧羊人.
他要牧養一些快要拿去屠宰去賣的羊.
即是說他要養到他肥肥白白.
要他有好的出產等等.
第五節就提到有些買家.
「買他們的宰了他們」.
即是說是那些買羊.
即是說他們是做開肉檔.
去這些羊圈買羊.
到底誰是那個買家呢.
即是說是那些羊群的僱主.
這裡說「以自己為無罪,賣他們的雪,耶和說是應當稱重的,因我成為富足」.
很特別就是買家就說我沒有罪的.
我是買來殺的.
我是真金白銀這樣買賣的交易的.
我不是去偷不是去搶那些羊的.
因為他有價值.
即是說賣到很高的利潤.
其實想說什麼呢.
相信都很肯定是在說以色列這個民族.
即是百姓就是群羊.
其實一直每一本書都是這樣說的.
神真是以他的子民就是群羊.
所以現在就在這個賣賣的協議底下.
到底想帶出什麼訊息呢.
和撒迦利亞又扮演著什麼角色呢.
到底他是否吃力不討好呢.
那個牧羊人很盡心去做好了.
到底百姓的群羊領不領情呢.
對撒迦利亞這個牧羊人會不會怪責他.
甚至乎遠離他呢.
因為下面說到三十塊錢耶穌被出賣.
耶穌被自己的民族.
被自己所有的當時代的同胞.
完全棄絕.
連他的門徒都棄絕他.
你就可以想像到那個棄絕是很大的.
但是現在撒迦利亞在主前的五百多年的時間.

$^{201}$首先來到預先演出一套劇集.
就是牧羊人被出賣.
最終是三十塊錢被出賣被棄絕的一個戲劇.
如果是戲劇的話.
如果是一個預言.
又是給我們有一個這樣的.
將來要發生的事.
就是按足字面應驗在彌賽亞身上.
所以這段聖經有時候那種有趣性和難解.
就是在這些地方.
我們看看那個戲劇.
他想描述的那個大概的內容.
有了買賣雙方.
有了群羊.
也有牧羊人.
有四個角色.
第六節.
耶和華說.
我不再連戍這地的居民.
必將這民交給各人的鄰舍和他們王的手中.
他們必毀滅這地.
我也不救這民脫離他們的手.
沒錯的確那些買家.
正正就是神興起的那些列國的君王.
周邊的國家.
來到去要對付以色列.
神交付給這些列國的君王.
所以那些買家就說.
我給足錢的.
我不是來偷搶的.
我要這些羊買來屠宰的.
就是說如果神將那些群羊.
就是他的子民交付敵人.
那麼那個審判是很大.
或者說那個管教.
在聖經的世界.
聖經的角度就是.
用蜀靈的角度.
那個叫做管教.
就是神交付群羊給敵人.

$^{241}$敵人的眼中就是說.
我用錢買的.
那麼戲劇的元素就在這裡出現.
第七節.
於是我牧羊者將宰的群羊.
將宰就是快要給神來刑罰的百姓.
就是群中最困苦的羊.
我拿著兩斤杖.
一斤我稱為榮美.
一斤我稱為聯索.
聯索就是那個聯合的意思.
英文就是Union.
Union.
這樣我牧養了群羊.
一月之內我除滅三個牧人.
因為我的心厭煩他們.
他們的心也憎嫌我.
我就說:我不牧養你們.
要死的由他死.
要喪亡的由他喪亡.
如剩的由他們彼此相食.
我截斷那稱為榮美的杖.
表明我廢棄與萬民所立的約.
當那日就廢棄了.
那些仰望我的困苦羊.
就知道所說的是耶和華的話.
我對他們說:你們若以為美.
就給我公價.
不然就罷了.
於是他們給了三十塊錢.
作為我的公價.
耶和華吩咐我說.
要把眾人所估定美好的價值.
雕給瑤戶.
那些出賣耶穌的那些祭司長.
那個當時的議會.
與猶大談好了三十塊錢.
後來那些錢猶大要不到.
因為猶大自己走了去上吊自殺.
那些人就說很不老禮.

$^{281}$又不敢要那些錢.
你不要沒人要就省下.
他們不敢.
因為說好了是猶大出賣他們自己的老師.
所以他們讀了這段經文.
似乎沒有這樣對號入座.
那些祭司長將那三十塊錢.
就雕給瑤戶.
其實是幫那些沒人去告老歸田.
還是說嚇死異鄉的人.
安葬那些異鄉人.
去買了一塊田叫血田.
我們等一下看看《四道行傳》第一章有記載.
接著第十四節.
撒迦利亞不是扮演著猶大的角色.
是他不被這些群羊接納.
是被人排擠.
撒迦利亞說你也要給我工錢.
好,我們就給你三十塊錢.
說好了.
接著我便將這三十塊錢.
在耶和華的殿中雕給瑤戶.
雕給瑤戶,瑤戶即是做那些Potter.
Potter即是那些叫做甚麼.
那些塗土.
即是有個東西轉.
用泥來搓那些做一些公仔.
那時他們不是做哆啦A夢.
不是做Hello Kitty.
是做那些偶像,那些公仔.
用來拜.
所以他們就這樣來到.
撒迦利亞就扮演著一個好像.
一個失職的牧羊人.
其實就不是他失職.
是那些百姓群羊不要他.
其實彼此有一個爭論.
即是好像變成神憎鬼厭那樣的撒迦利亞.
很特別的一個經文.
所以我們當然最特別就是.

$^{321}$主耶穌讀到這段聖經.
將來他要被全民族出賣.
和嫌棄這位好牧人.
要來到去被三十塊錢的公價.
來去置於死地.
自己最愛的門徒猶大出賣他.
這要告訴我們第十四節.
我猶節斷稱為聯索的那根杖.
表明了我廢棄猶大與以色列弟兄的情誼.
十五節,耶和華又吩咐我說.
你再取愚昧牧人所用的器具.
因我要在這地興起一個牧人.
他不看顧喪亡的,不尋找分散的,.
不醫治受傷的,也不牧養強壯的,.
卻要削肥羊的肉,撕裂牠的蹄子.
沒用的牧人,凋棄羊群有禍了.
刀必淋到牠的蚌壁和幼眼上.
牠的蚌壁必全然枯乾.
牠的幼眼必昏暗失明.
這裡又更加難解釋.
撒迦利亞不單扮演一個好牧人.
幫買賣商方去牧養羊.
希望強壯的可以出欄.
以前的豬和牛可以出欄,可以賣.
現在撒迦利亞養好了百姓,百姓不要牠.
現在神要牠扮演一個不好的牧人.
相信很難接受.
壞事做盡,一個牧羊人.
不好的地方就是不理群羊的生死.
這是最大的諷刺.
這是十一章,好像有些戲劇成份.
所以你看到舊約的《先知書》.
都幾多元的體材,可以說神是很創意的神.
透過不同的戲劇讓事情在將來發生.
真是在主耶穌的年代,耶穌鞠躬盡瘁.
為自己的民族,由出世到祂出來傳道.
都是一心一意去幫人,走遍各城各鄉.
來到去宣講神的道,見到很多人困苦流離.
祂就來去憐憫他們等等.
最終沒有人接納祂,甚至嘲笑祂,侮辱祂.

$^{361}$就是應驗了十一章的這個好牧人.
被三十塊錢出賣,被人來厭棄.
所以這是很特別的寓言.
所以《先知書》原來都幾好看.
雖然有些我們很難解釋.
但有些儘管我們解釋不到都不要緊.
因為整個去到新約.
就是這樣來按舊約的字面應驗.
所以有時看新約,有時幾特別的地方就是.
原來神這樣保守祂自己的寓言的說話.
有些是真的要按字面應驗.
譬如三十塊錢,真的有這麼巧.
還是他們真的這樣知道.
似乎那些祭司長又不會這麼敬虔.
要跟足那個劇本去做.
那不就是他們自己理虧.
他們就是那個賣家,賣自己的百姓賺錢.
那些賣家就是將來的猶太的領袖.
那你就可以想像到神的說話.
有時就算人怎樣想干預都干預不到.
神的話語自己會應驗.
有時我們讀舊約,如果未有新約.
當然我們不知道在說什麼.
但是有了新約再回望舊約,就很不同了.
這是一個很重要的信號.
聖經是神的默示.
不是普通人寫,不是撒迦利亞.
他有突然間的感動,自己作一些作品出來.
是他自己天馬行空地作出來.
是有神的管理,是有神的默示.
這樣來到,要這樣記載,要這樣發生.
將來就是這樣應驗,這樣就很特別了.
這樣理解聖經的權威.
就是在這裡不斷讓我們看到這些事情.
連主耶穌自己在十誡上.
十誡七言的其中一句就是.
「我學了」.
但是有人呼應告訴我們.
他想起經上的說話,未應驗.
然後說「我學了」.

$^{401}$這是詩篇的二十二篇.
你可以想像到耶穌基督.
是真的有這樣的對於舊約的認識.
對於神的話語.
關於祂的預言的彌賽亞.
主耶穌要來滿足.
神在舊約裡面的預言的成就,要求.
你去理解,有時一拼我們看舊約新約.
這樣才有東西看的.
因為猶太人他們不相信新約.
其實沒有東西看的.
因為舊約就是說了一些.
好像是不是會發生的事情.
到底所指何人呢?.
彌賽亞是誰呢?.
如果沒有了耶穌,什麼都是假的了.
什麼都是沒有意思的.
但是當我們基督徒,基督教.
我們有新約,我們就很不同了.
回望,看回舊約的經文.
因為根本就是使徒他們寫出來.
告訴我們,這樣就是應驗舊約的經文.
他們不是特意要對號入座.
而是有教會我們的傳統.
尤其是福音派的傳統.
我們仍然相信使徒們.
他們是有特別的靈感,特別的啟示.
跟我們後世的信徒是有些不同的.
因為他們是剛剛橫跨耶穌.
他們見過耶穌.
也知道耶穌死而復活升天.
他們去傳道.
神是給他們這一班人.
特別是電機的人.
有特別的靈感來寫聖經的.
因為他們親眼見過耶穌.
所以他們不可以只是.
當時如果有相片當然好.
如果有電子化的網路當然好.
但是他們是要將那個真理.

$^{441}$來著書立說.
來成為後世的信徒的典範.
和那個真理依據.
不可以人云亦云的憑感覺.
憑一個自己的主觀的感受.
所以信仰不是建立在那些感覺上.
是建立在聖經的真理上.
所以一拼這樣看的時候.
看寓言都要有這樣的心靈的前設.
去理解今天我們基督教.
為什麼這麼捍衛聖經的無誤.
神的默示是不會改變的原因的理由.
十二章先不揭了.
我們回來十章.
簡單解一解這段經文.
全章是神再抬舉.
再使到以色列出埃及.
他們可以再被在列國當中被焦聚.
這裡要和大家看一段經文.
申命記.
申命記的二十八章.
這章聖經是摩西之約的一個很重要的.
你可以說是條款.
可以說是全卷的.
不單止申命記.
其實摩西五經.
全個申命記的摩西之約.
創出利民生這五卷書.
最重要的精華就是二十八章.
所以如果能夠背熟這一章.
就很不簡單.
因為神很清楚.
既然那些以色列人.
你們在西乃山和我納了約.
在寫申命記的時候.
其實已經離開了埃及四十年.
摩西就快要死了.
在摩亞平原那個曠野的地方.
摩西最後的一篇講道.
就是申命記了.

$^{481}$叫他們重新記住.
神和我們這個民族納了約.
某程度上摩西代表以色列人.
等於簽了紙.
即是約已經生效了.
二十八章是說什麼呢.
就是說生死和福.
二十八章.
即是說簽了約之後.
現在才發現那些條款很難啃.
原來這些條款是.
即是說他們簽了成為神的子民之後.
你不要再下排輕輕.
神現在和我們立約.
雖然有十戒.
但是現在如果我們成為他的子民.
有什麼福氣.
如果不是他的子民會怎樣呢.
這個叫做申命記的摩西之約.
是二十八章.
是很重要的一章.
可能我要讀多一點.
你如果留意聽從耶和華理神的話.
謹守遵行他的一切誡命.
就是我今天所吩咐你的.
他必使到你超乎天下萬民之上.
你若聽從耶和華理神的話.
這以下的福必追隨你.
臨到你身上.
你在城裡必蒙福.
在田間也必蒙福.
你身所生的.
地所產的.
生畜所下的.
以及牛毒羊羔都必蒙福.
你的康子和你的磚面盤都必蒙福.
即是你的恆生戶口.
你的匯豐戶口.
你的投資戶口.
你的定期戶口.

$^{521}$你的內地人民幣戶口.
你的金的戶口.
你的各方面的.
全部都蒙福.
即是這樣說.
你出也蒙福.
入也蒙福.
你想想.
出入平安.
出入也蒙福.
在城門口也蒙福.
在田間耕作的豐收蒙福.
什麼都很順利.
這個世界有沒有這樣的人呢?.
有的.
不過可能未必靠神.
但這裡是說以色列這個民族.
他們現在在摩亞平原.
快要進入迦南.
摩西一離世他們就進入迦南.
進入了應許之地.
流乃與物之地.
現在神就這樣最後.
透過摩西跟他們說.
這個十誡之約.
摩西的西乃山之約.
或者叫摩西之約.
因為摩西簽的.
十誡.
他領受法版.
現在就說詳細的條款.
那個內容.
那就真的沒辦法了.
以色列人難怪這麼蒙福.
什麼都是福.
什麼出入都是福.
第七節.
仇敵起來攻擊你.
耶和華必使他們在你面前.
被你殺敗.

$^{561}$他們從一條路來攻擊你.
必從七條路逃跑.
在你倉房裡.
並你手所辦的一切事上.
耶和華所命的福必臨到你.
耶和華你神也要在所給你的地上賜福於你.
你若謹守耶和華你神的誡命.
遵行他的道.
他必照著向你所起的誓.
立你作自己的聖文.
天下萬民見你歸在耶和華的名下.
就要懼怕你.
你在耶和華向你列祖起誓.
應許賜你的地上.
他必使到你身所生的.
生畜所下的.
地所產的.
都綽綽有餘.
耶和華必為你開天上的虎腹.
按時降雨.
按時是很重要的.
失驚無神.
突然秋天下黑雨.
是很大件事.
是按時的.
按時降雨在你地上.
在你手裡所辦的一切事上.
賜福於你.
你必謝給許多國民.
卻不至向他們謝太.
即是家家戶戶都有餘的.
那些存款有很多個零的.
你可以想像得到.
是多麼的蒙福.
你若聽從耶和華理神的誡命.
就是我今日所吩咐你的.
謹守遵行.
不偏左右.
也不隨從.
事奉別神.

$^{601}$耶和華就必使到你.
作首不作尾.
但居上不居下.
成為人上人.
十四節的祝福.
節節都是說.
福杯滿日.
只需要是一個動作.
聽從神的誡命.
即是說要不斷的揣摩.
什麼是神喜歡我們要做的.
這個民族.
什麼是神不想我們做的.
就是這樣叫聽從.
小時候就沒得去想了.
要這樣做.
一早起床一定要刷牙.
沒得不刷.
媽媽叫到七點就要起床.
吃早餐.
吃完早餐就洗臉.
背書包上學.
一定要這樣做.
但你一路長大.
進到大學.
七點還不起床.
都還沒轉身.
昨晚剛剛開完會.
三四點才睡.
在宿舍住.
還八點吃早餐.
當然不可能.
我自己會調整的.
午餐十一點鐘.
剛剛醒.
洗完臉就和同學一起.
一路長大.
聽從就不同了.
就不是七點鐘要做這件事.
八點鐘要做這件事.

$^{641}$是小朋友要這樣做的.
但長大之後.
他就要去思想.
去揣摩.
什麼是父母的教導的精髓.
精髓就是.
吃早餐原來很重要.
守時很重要.
有規律的生活很重要.
但因時制宜.
已經讀大學.
已經住宿舍.
你不可能.
每天早上七點醒.
八點一定要坐好.
穿好衣服吃早餐.
不是不行.
有規律不是錯的.
但有時需要調整.
甚至要打破這些規律.
因為已經有自己生活的節奏了.
所以以色列人.
這十四節是一個.
很大的一個方向.
聽從無福.
什麼的福.
神說各樣你想到的都是福.
不過你看看這張聖經有多少節.
六十八節.
還有五十四節的奏坐.
二十八章其中出名的地方就是.
十四節的祝福.
五十四節的奏坐.
第十五節開始.
你如果不聽從耶和華.
理神的話.
不謹守遵行他的一切誡命律例.
就是我今天所吩咐你的.
這以下的奏坐.
都必追隨你.

$^{681}$臨到你身上.
剛剛和上文所說的.
一個對立的.
你在城裡必受奏坐.
在田間也必受奏坐.
你的糠子和你的磚面盤.
都必受奏坐.
你身所生的.
地所產的.
以及牛獨羊羔.
都必受奏坐.
你出也受奏坐.
入也受奏坐.
耶和華因你行惡離棄他.
必在你手裡所辦的一切事上.
使得奏坐擾亂.
責罰臨到你.
直到你被毀滅.
遠速速的滅亡.
耶和華必使到瘟疫貼在你身上.
直到他將你從所進去.
得為業的地上滅絕.
耶和華要用勞病.
肺癆.
熱病.
火症.
藥疾.
刀劍.
可能藏屍.
旱風.
霉難.
攻擊你.
也都要追趕你.
直到你滅亡.
你頭上的天要變為銅.
腳下的地.
要變為鐵.
耶和華要使到你.
那降在你地上的雨.
變為塵沙.

$^{721}$即是變成沙漠化.
你可以想像.
如果你在北京.
突然間大風沙.
那種情景都很痛苦.
又出不到街.
出街吹到整個臉痛得不得了.
從天臨到你身上.
直到你滅亡.
耶和華必使到你.
敗在仇敵面前.
你從一條路去攻擊他們.
必從七條路逃跑.
是你逃跑.
你必在天下萬國中拋來拋去.
你的屍首.
必給空中的飛鳥.
和地上的走獸.
作食物.
並無人鬧趕.
控趕.
耶和華必用埃及人的倉.
病治倉.
牛皮鱔.
乳戒.
攻擊你.
即是這些橫疾.
這些是叫做.
不知道什麼的.
古靈精怪的疾病.
都是一次過發在你身上.
使你不能夠醫治.
耶和華必用癲狂.
眼核.
心驚.
攻擊你.
心驚膽戰.
你可以想像得到.
那種痛苦.
你必在五間摸索.

$^{761}$好像核子在暗中摸索一樣.
你所行的必不亨通.
時常遭遇欺壓搶奪.
無人答救.
你聘定了妻.
別人必與他同房.
你建造房屋.
不得住在其內.
你栽種葡萄園.
也不得用其中的果子.
你的牛在你眼前宰了.
你必不得吃他的肉.
你的奴在你眼前被搶奪.
不得歸還.
即是說充滿了強暴.
充滿了不公義的社會.
你的羊歸了仇敵.
無人答救.
即是周圍的搶奪.
好像以前那些土匪.
誰是山寨王.
要搶就搶.
沒人敢出聲.
你的兒女必歸與別國的民.
你的眼目終日切忘.
甚至失明.
即是強搶民女.
必被你所不認識的國民吃盡.
對不起.
你的土產和你的奴鹿得來.
你時常被欺負.
受壓制.
甚至你一眼中所看見的.
必至瘋狂.
即是人都瘋了.
甚麼的衰運.
你都日日這樣走.
日日走衰運.
你可想而知那種慘情.
耶護華必攻擊你.

$^{801}$使你膝上腿上.
從腳掌到頭頂.
掌毒瘡無法醫治.
身外物沒了就沒了.
自己身體健康是最重要.
都不放過你.
你的當事人要生毒瘡.
由腳底到頭頂.
三十六節.
耶護華必將你和你所立的王.
領到你和你們早不認識的國去.
在那裡你必是木頭石頭的神.
你在耶護華領你到的各國中.
要令人驚喜.
笑談.
譏笑.
被人嘲笑.
剛剛上個星期我們不是說.
以色列人將來.
有十個外國人.
抓住一個猶太人.
抓住他們的衣襟.
即是mung 衫尾.
我們要見耶護華.
我們要去求那個恩典.
我們去尋求神的恩典.
這樣的將來的光景.
以色列人又是神的冠冕的寶石.
但是現在摩西的預言.
摩西的警告.
我就平平安安回天家了.
你們就要受苦很久了.
如果你們進到迦南之後.
那個約是正式生效的.
其實生效了.
你們就知道那個味道了.
這樣的危言聳聽.
摩西.
還有真的很諷刺.
十四節.

$^{841}$一邊聽一邊開心拍手掌.
但是由十五節開始.
不斷說奏座.
淋漓盡致的奏座.
方方面面.
不只是由那些物質這麼簡單.
心靈變成瘋狂.
你可以想像到.
我們走了.
不在這個國家了.
去別的國家了.
成為譏笑笑談.
別人笑這個民族.
你說你的神這麼厲害.
為什麼你們搞成這樣.
你記得路德跟隨他的婆婆拿訛米.
回來的時候.
他們那些以色列人就.
這個不是顯微拿訛米.
他們不認識路德.
其實他們在外國.
都是受盡這些的凌辱.
這樣回來.
你可以想像到.
這裡告訴我們.
我不再讀下去了.
因為都挺多的奏座.
不過你可以想像到.
這些的奏座.
一路這樣說的時候.
其實不是沒有出路的.
是有出路的.
我們看一看這個.
應該是歷代志.
在這黃記之下的歷代志上.
歷代志下應該是第七章.
應該第六章會對.
就是他獻殿的時間.
第六章的歷代志下.
也是很出名的經文.

$^{881}$就是有史以來.
第一個聖殿.
猶太人第一個聖殿.
供奉耶和華的.
以前就是會幕.
一路朝行晚策.
其實入了迦南之後.
就不需要策了.
所以就在士羅那個地方.
但是在士羅一路都沒有人去想.
神的殿應該怎樣的呢.
那個規模啊.
那個建築啊.
直到當然.
因為大家都是一個時思時代.
後來真的穩定了.
去到大衛了.
真的做了王了.
國家復歸統一了.
不會有很多的諸侯割據之後.
大衛去想一下這件事.
有神的心去想一下.
都是得到神的喜悅.
當然大衛就不可以建造聖殿.
不可以說不可以.
不可以之中.
大衛是有生之年.
預備很多很多建造聖殿的材料.
他要來為國家的庫房.
積累很多很多的金銀寶石.
來為他的兒子將來建造聖殿.
來慶功的那些材料.
所以雖然神不需要你做.
但是你仍然都有那個心.
為神去建造.
這個是很難得的.
我們其實都要這樣.
這個不是叫賄賂.
這個是告訴我們.
神有他的原因不用你做.

$^{921}$但是你也可以參與.
在其他方面的建造.
這個其實我們就要學大衛了.
你可以說他是緊張他自己的家族.
因為如果他建好了聖殿.
或者他的兒子將來建好聖殿.
那就是說那個父蔭.
就是神的Blessing的蔭秘.
就延續他的後代.
你可以說他有這樣的私心.
但是最重要的就是.
他真的有一個心.
為神的家來建造聖殿.
所以神就說.
你為我建造聖殿.
就是你為我建造家.
我為你建造你的家.
這樣的交易是很著數的.
大衛.
所以他真的是一個很醒目的人.
他在找神的Blessing.
所以你可以說猶大支派.
一直的君王蒙福.
都是因為大衛王這個事件.
和他一生對神的忠誠.
我們說的是第一個聖殿.
是很重要的.
因為是一個里程碑.
這個聖殿的起好.
所羅門真的是一個很有智慧的人.
他第一件事要求的.
就是為在這個聖殿裡面敬拜的人.
他的子民.
一來就是如果這個聖殿.
是大家公認的一個聚集地.
是一個宗教的聚集地.
那所羅門的國度就會穩固.
他就可以延續千秋萬代.
所以這個聖殿是一個很重要的凝聚.
他就很會祈禱.

$^{961}$他就著《申命記》28章去祈禱.
我們看看第十三節.
六章的十二節.
所羅門當著以色列回眾.
站在耶和華的壇前.
舉起手來.
所羅門曾經做一個銅臺.
長五爪,寬五爪,高三爪.
那個比例是按那個比例.
是按繪幕的那個比例.
長闊高比例.
所以大家不要以為是錯的.
就放在院中.
因為是銅臺.
就站在臺上.
當著以色列的回眾跪下.
向天舉手.
說耶和華以色列的神啊.
天上地下沒有神可比你的.
你向那盡心行在你面前的僕人.
守約斯前外.
向你僕人我傅大衛所應許的話.
現在應驗了.
你親口應許親手成就.
正如今天一樣.
耶和華以色列的神啊.
你所應許你僕人我傅大衛的話說.
你的子孫若謹慎自己的行為.
遵守我的律法.
像你在我面前所行的一樣.
就不斷人錯.
以色列的國為現在求你.
應驗這話.
耶和華以色列的神啊.
求你成就.
向你僕人大衛所應許的話.
神果真與世人同住在地上嗎?.
看那天和天上的天.
他很懂得擦鞋.
他真的很懂得神的心.

$^{1001}$你是造物主.
整個宇宙萬物都是你創造的.
怎麼會你要被區限在一個地方呢?.
但其實就是他想神在這裡.
需要有神的同在.
尚且都不足夠你居住.
何況我所見的這殿.
就算多大.
神都不足夠你住.
為求耶和華我的神.
誰顧僕人的禱告祈求.
撫聽僕人在你面前的祈禱呼籲.
願你晝夜看顧這殿.
即是想神住在這裡.
雖然不夠你住.
不應該限制你.
但求你限制你自己.
你眷顧我們這個民族.
你因為我爸爸的緣故.
你和我爸爸納了約.
所以求你看顧.
即是他想神留意.
裡面敬拜的人.
那個的福樂,平安和盼望等等.
所以這些祈禱.
其實我們又要學.
我們信主之後.
其實要學習禱告.
這些就是一些很重要的典範的祈禱.
他很明白神的屬性.
他深知道我做多少都沒有用.
不是這個硬件多美,多有規模.
首先就是自己自認無有.
但是神啊.
你答應了我爸爸.
你說要看顧我們這個民族.
真的要兌現.
現在我們會照著你的教導.
在這個地方去有敬拜.
不是一年兩次.

$^{1041}$而是我們每個禮拜行常.
這樣去聚集敬拜你.
你要看顧才行.
這是一個很重要的祈禱.
你可以說所羅門真的很有智慧.
他還要求的就是神的同在.
其他的他不需要.
第20節.
願你晝夜看顧這殿.
就是你應許立為你名的居所.
求你垂聽僕人向此處禱告的話.
你僕人和李文,以色列.
向此處祈禱的時候.
求你從天上你的居所垂聽.
垂聽而赦免.
人如果得罪淪赦.
有人叫他起誓.
他來到這殿.
在你的壇前起誓.
求你從天上垂聽.
判斷你的僕人定惡人有罪.
照他所行的報應在他頭上.
定惡人有理.
照他的而賞賜他.
所以有時拆開這樣說.
猶太人他們很緊張聖殿.
正正就是因為他們知道.
聖殿的地理環境.
和那個屬靈的氣氛.
屬靈的性質.
是特別強WIFI.
直上天庭.
去到神的殿.
神垂聽.
所以他們一定要去到聖殿祈禱.
所以所羅門現在這樣祈禱.
限制了神在這個地方.
祝福和垂聽和審判.
按公義.
很重要就是告訴我們.

$^{1081}$以色列人就是這樣緊張那個殿.
但是神說.
沒錯.
我會同在.
但是有一天.
如果你們不聽話.
這個殿都要被拆毀.
所以所羅門的聖殿就是這樣被壓述.
被巴比倫拆毀.
所以這裡就是一個建殿的時候.
你可以說正式的開幕典禮.
所羅門有這個祈禱.
真是很重要.
一個很有智慧的君王.
不需要祈禱.
那些百姓聽聽話話.
只需要神在這個地方.
在這個聖殿.
間接.
他的王位就很穩固.
因為百姓其實想什麼.
就是想神祝福.
想神有恩典.
想我們自己這個國家.
大家安居樂業.
大家四景平安.
夜不閉戶.
大家出入平安.
平安你的國度就穩固.
就不需要擔心.
所以所羅門的祈禱.
真是很寶貴.
其實作為父母.
你想你的子女將來長大.
敬虔.
又不會入黑社會.
又不會坐牢.
又不會成為害群之馬.
其實很簡單.
你為他祈禱.

$^{1121}$從小就帶他回教會.
他從小在教會裡長大.
有神透過不同的場景.
不同的人.
不同的事.
來慢慢塗鑿他的生命.
他將來是怎樣都不會離開神.
這是最大最好的投資.
什麼教育都不需要.
帶他回教會.
還有一起敬拜神.
這正是所羅門的智慧所在.
不會求我的軍隊.
我的幕僚.
他求的是有這個聖殿.
最重要是有神你的同在.
你的眼目的看顧.
國家就會很穩固.
再讀一點.
因為神聽了他的祈禱之後.
來應運了一件事.
二十四節.
你的文以色列如果得罪你.
敗在仇敵面前.
又回心轉意承認你的名.
在這殿裡向你祈求禱告.
求你從天上垂聽.
赦免你文以色列的罪.
使他們歸回你賜給他們和他們列祖之地.
他們常常都提到.
所羅門提到.
必須要在這個殿裡發生.
你的文因得罪你.
你懲罰他們.
使得天閉塞不下雨.
他們若向此處禱告.
承認你的名.
離開他們的罪.
求你在天上垂聽.
赦免你僕人和你文以色列的罪.

$^{1161}$將當行的善道指教他們.
且降雨在你的地.
就是你賜給你文遺孽之地.
國中若有饑荒.
等等等等.
我不讀這麼多了.
第三十八節.
其實一直在說的這些奏坐.
正正就是《申命記》二十八章.
有些我沒有再讀的那些奏坐.
就是照足《申命記》二十八章.
所羅門來到祈禱.
將那個奏坐逆轉的唯一的途徑.
辦法就是回到這個殿.
回不去的.
在遠方被擄的.
就朝向耶路撒冷去禱告.
其實就是朝向耶和華祈禱.
這樣也對的.
將來耶穌說.
耶穌的救恩成就之後.
將來的人敬拜.
不需要再在基里森山.
不需要再在耶路撒冷.
因為神是一個靈.
敬拜的是用心靈成神去敬拜神.
所以神已經將聖靈給了我們.
何時我們去仰望神呢?.
我們其實就去到天上的耶路撒冷了.
這裡很重要就是三十八節.
他們若在老道之地盡心盡性歸復你.
又向自己的地.
就是你賜給他們列祖之地.
和你所選擇的城.
並我為你明所建造的殿禱告.
求你從天上你的居所.
隨聽你民的禱告祈求.
為他們伸冤.
赦免他們的過犯.
我的神啊.

$^{1201}$現在求你睜眼看.
則而聽在此處所獻的禱告.
耶和華神啊.
求你起來.
和你有能力的約桂.
同入安息之所.
這個是詩篇的一三二篇大衛的詩.
所羅門引他爸爸的這篇詩.
就是說怨聖殿最重要的那個器皿.
就是約桂去到安居之所.
當時那篇詩篇是沒有找到的.
因為約桂被擄了.
後來回來之後.
就去到一個叫做俄別以東的地方.
在這個基列耶林.
失散了.
大衛就去迎接那個約桂.
所羅門在這裡.
其實可能已經安放好了.
他就引述這篇詩.
來讓約桂.
就是說神同在的記號.
就在現在我這個聖殿了.
我是那個建築的人了.
所以等同於就是神.
祝福我們這個家庭.
我們這個猶大之派.
我的家族.
要有那個一路繼承的皇位的人.
大概這樣.
當然他不是這麼有限的.
他都是為整個國家.
在這個時候.
他有他的身份.
因為他是所羅門王.
他有這個身份.
為全國來禱告.
其實應該是那些大祭司祈禱的.
所羅門就扮演著這個角色.
他自己在這裡行開光之禮.

$^{1241}$他自己在這個開幕裡面.
他帶領祈禱.
神就收穫.
第七章.
所羅門祈禱以筆.
就有火從天上降下來.
燒盡凡祭和別的祭.
耶和華的榮光充滿了電.
因耶和華的榮光充滿了耶和華的電.
所以祭司不能進電.
那火降下.
耶和華的榮光在電上的時候.
以色列眾人看見.
就在鋪石地俯伏叩拜.
稱謝耶和華說.
耶和華本為善.
他的慈愛永遠長存.
就是說神公開的閱立.
所羅門的禱告.
所羅門所禱告的所有的東西.
神照單傳授.
他一定會看顧.
只要所羅門這樣說.
這一句神的話語.
這個民族聽取神的教導.
走在神的公義下.
他們就會蒙福.
所羅門當然是大事的慶祝.
敬拜完當然一起吃.
所以就將祭物.
將他預備好的禮物.
每個人有多少東西呢.
後面有說的,不過不讀了.
將這些東西給百姓.
有平安祭,有掃祭,有滋遊等等.
第十節.
7月23日,剛剛完結了駐彭節.
王遣散眾民.
他們因見耶和華向大衛.
和所羅門與他們以色列所施的恩惠.

$^{1281}$都心中喜樂,各歸各家去了.
帶著平安喜樂回家.
我們現在這個民族有了一個聖殿.
很重要的,其實真的是里程碑.
對於這個民族有自己的國家.
現在有自己的君王.
有自己的一個四景平安.
現在更加有的就是.
最重要的就是每個人的宗教的中心.
就是聖殿.
就有神的祝福.
這是很重要的一個指標.
又拆開說.
天主教他們的傳教策略.
他們會去一些比較貧窮的地方.
去到之後一定是找那個城的中心點.
來建教堂.
你可以再留意一下香港的教堂.
有哪些是在那些區域.
由它建造的時候.
那些地方是重要的中心點.
就一直擴散開去.
就不會一開始找很遠很偏遠.
找去韶關,去什麼地方.
不會的,找一個中心點.
你可以說坐落在那裡.
讓那些周圍的人來到這個教堂.
來朝見他們的神.
或者望彌撒.
這是他們傳教的策略.
和基督教很不同的.
基督教曾經在七十八十年代.
是騙傳.
周圍去直堂.
幾十人去直堂.
又去開教會.
開很多很多教會.
我們宣道會今年到現在.
已經超過一百一十間.
超過一百一十間.

$^{1321}$一百一十幾我就不記得了.
不是每間都幾百人.
過千人.
有更多,有超過百成.
一百人都不夠.
幾十人的都不少.
到社會一直轉變.
那個經營.
你可以說那個叫傳道的策略.
因為樓價高企.
地少人多.
很困難的經營.
這個當然拆開了說.
所羅門的聖殿.
就有這個凝聚四方的子民.
的那個重要的.
你可以說那個凝聚力.
這個都很重要的.
感恩我們在這裡.
終於建了這個教堂.
對這個社區都是一個很大的祝福.
盼望我們作為.
神的子民.
我們都應該將這份Blessing.
來到去.
透過我們的言說.
我們人與人之間的待人接物.
流露出我們是屬神的兒女.
有時現在都和那些.
餐廳的職員都熟了.
有時都慢慢有機會可以.
給他們認識一下.
你做什麼的.
我在教會做牧師的.
他們很好奇是什麼來的.
原來有些人真的不是很知道.
盼望我們都能夠.
透過所羅門這個建殿.
帶出這一種的Blessing.
去到其他地方.

$^{1361}$但是還沒說完.
真正重要的是這裡.
第十一節.
所羅門造成了這一切之後.
第十二節.
耶和華向所羅門顯現.
對他說.
我已聽了你的禱告.
也選擇這地方.
作為祭祀我的奠宇.
我若使到天閉塞不下雨.
或使到蝗蟲吃這地的出產.
或使到瘟疫流行在我民中.
這稱為我名下的子民.
若是自卑禱告尋求我的面.
轉離他們的惡行.
我必從天上垂聽.
赦免他們的罪.
醫治他們的地.
我必睜眼看.
則而聽在此處所獻的禱告.
現在我已選擇這殿.
分別為聖.
使我名永在其中.
我的眼我的心也必常在那裡.
你若在我面前.
效法你父大位所行的.
遵行我一切所吩咐你的.
謹守我的律例典章.
我就必兼顧你的國為.
正如我與你父大位所立的約說.
你的子孫.
必不斷人作以色列的王.
倘若你們轉去.
丟棄我指示你們的律例誡命.
去事奉敬拜別神.
我就必將以色列人.
從我賜給他們的地上拔出根來.
並且我以我為己民.
所分別為聖的殿.

$^{1401}$也必捨棄不顧.
使他在萬民中作笑談.
被譏笑.
這殿雖然甚高.
將來經過的人必驚訝說.
耶和華為何向這地和這殿如此行.
人必回答說.
是因此地的人離棄耶和華.
他們列祖的神.
就是領他們出埃及地的神.
去親近別神.
敬拜事奉他.
所以耶和華.
使這一切災禍臨到他們.
講完了.
這就是神聽所羅門祈禱.
但是帶著一個應許的.
就是說只要你.
令到這個國家的國民.
可能在祭司,先知的領導下.
不只是政治.
而是宗教上等等.
來去跟足神.
透過摩西大衛的爸爸.
那個律例典章來進行.
就可以得到神的任避.
神就會看顧這個殿.
假如什麼時候不聽.
轉去敬拜別神.
神的咒詛就臨到.
其中的咒詛就是.
神會轉來.
不再在這個聖殿裡面.
最重要的是那個靈魂.
沒有了神的同在.
這是一個聖殿的維護.
最重要的地方.
所以來到神的家.
神的殿.
我們要有一種敏銳的心靈.

$^{1441}$神在這裡.
我們要學習.
這樣去敬畏神.
這樣去敬拜神.
我們的靈命就會蒙福.
我們在不同的地方.
都來得到對神的敬重.
跟我們剛才看的撒迦利亞書.
第十章有什麼關連呢?.
你可以想像得到.
那個時代已經不同了.
神現在的刑罰時間已經過去了.
以色列人已經從.
被擄的巴比倫.
波斯回歸.
回歸上來了.
就是說現在這個時代.
不再是耶利米審判.
滅亡的時代了.
神要離棄聖殿的時代了.
神現在要重新.
回到這個聖殿.
有神的同在.
所以神是.
透過兩個的領袖.
撒迦利亞.
和哈該.
來勉勵當時的省長.
這個索羅巴伯.
和當時的大祭司約書亞.
勉勵他們來重建這個聖殿.
重建好之後.
神的同在就會出現了.
所以這個聖殿的重要.
就是慢慢慢慢讓四百多.
當然沒有五百年那麼久.
將來的彌賽亞耶穌.
就在這個聖殿裡面.
來去履行.
他在神的救恩.

$^{1481}$所扮演的角色.
他成為贖罪祭.
這樣來去.
在以色列人當中.
來去完成神的救恩的託付.
所以在現在這個重建聖殿.
是有這個歷史的任務的.
所以這個聖殿一直在建的過程.
這是第二年的大利烏王.
九月講的預言.
但是現在去到第十章.
神對於將來的以色列.
是有一個這樣的寄望.
是一個這樣的美麗的.
將來的復興.
第十章就講到將來的復興.
可以去到什麼的層次.
可以回到出埃及立約的時候.
申命記二十八章的光景.
那個水平.
你可以想像到.
是多麼奇妙的地方.
一個民族這樣滅亡了.
可以回到原地.
所以第十章相對比較容易去理解.
神要懲罰那些帶領百姓.
走錯路的假先知.
特別是我們剛才第十章講的.
就是那些家臣的預言.
卜士即是占卜的.
作夢的即是作夷夢.
他們這樣白白的安慰人.
即是說話好像那些風水佬.
騙你十年八年那樣.
什麼水晶啊.
紫微斗數啊.
什麼守文啊.
占卜啊.
什麼塔羅牌啊.
什麼運程啊.

$^{1521}$諸如此類.
又去黃大仙求籤啊.
又那些什麼伏羈啊.
道教那些啊.
諸如此類.
不過是不同的方式.
人都很好奇.
對於前途.
對於那個運程.
我的運是怎樣呢.
很好奇.
很多人都想知道.
但是這告訴我們.
神要懲罰這些曾經帶領.
不知道什麼方式.
帶領他的子民.
走錯了路的這些人.
使到在第十章的第二節.
所以眾人因為這樣的迷信.
或者是這些假的先知.
都是歸類去假先知的.
即是你假借神的名.
假借你是有些靈驗.
可以來幫人趨吉避凶.
這些都是屬於假先知的.
使到神的百姓如羊群流離.
因為沒有牧人.
沒有真正的牧者來帶領.
所以就受苦了.
但是神就會懲罰這些牧羊人.
將來他的民族.
第四節就很寶貴了.
以色列人好像防國石釘子.
征戰的工.
和一切的掌權.
都從他而出.
在尤大加那裡會興起這樣的領袖.
可以使到防國石.
即是有那個水的穩固根基.
釘子都是這樣的意思.

$^{1561}$釘穩了在牆上.
和那個征戰的工.
即是射箭的工.
神會再興起.
從尤大加那裡興起.
重要的是他現在重新眷顧這個民族.
所以第十章一直看下去.
全部都是說.
將來神要再紀念尤大加和約瑟加.
即是說全個以色列.
都是蒙神再度的眷顧.
去到什麼水平呢.
回到好像出埃及的時候.
這樣的一個光景.
即是最美好的日子.
由一個奴隸的一班人.
神現在眷顧台舉.
透過摩西帶領國民浩浩蕩蕩的.
去到這個西奈半島.
跟著和這個民族立約.
要賜福各種各樣的福氣.
和說你們要入迦南地.
是我亞伯拉罕要給他的應許之地.
是說地圖的應許.
真的全部應驗了.
但是現在神就透過撒加尼亞才知道.
現在重建聖殿.
將來的好日子.
要去到哪一個水平呢.
去到出埃及的時候.
這個是第九節.
十章的第九節.
我雖然播散他們在列國中.
他們必在遠方紀念我.
他們與兒女都必存活且得歸回.
我必再領他們出埃及地.
招罪他們出亞述.
領他們到基列和尼巴嫩.
這地方尚且不夠他們居住.
你可以想像得到.

$^{1601}$這個是新的出埃及.
是說那個救贖.
不是說現在大家一起回去埃及.
準備過紅海.
就不是這個意思.
而是在天下萬國.
他們好像有一個突然的救贖.
那是怎樣呢.
我們不能夠用我們現在的眼光去益測去判斷.
但保羅是有這個信念的.
當愛邦人的數目天滿.
揭一揭羅馬書十一章.
羅馬書的十一章.
可能不可以讀這麼多.
第二十五節.
十一章的羅馬書二十五節.
這三章聖經第九第十第十一章.
是在說一個問題.
保羅就是要和羅馬的教會.
來說述一個情況.
當然與大相關.
你們愛邦人.
你們有這麼多的資源.
你們應該要很殷勤的.
準備好你們的內心教會.
接待我和傳道者.
因為我們來到的時候.
我們要把福音傳到更遠的地方.
所以你們不要把福音.
扣留在你們自己的教會內部.
要一不斷的去傳福音.
為什麼呢.
因為這裡就告訴我們.
我還沒說這裡.
就是三章的羅馬書九十十一章.
試探討一個很重要的問題.
是和愛邦人說.
不是和猶太人說.
就是為什麼你們眼睜睜.
看到這麼多猶太人拒絕福音.

$^{1641}$而你們愛邦人.
包括羅馬教會.
在保羅的年代.
保羅的時期.
為什麼這麼多愛邦人這麼容易信主.
為什麼猶太人這麼難去信主.
好像有一個跨越不了的難.
跳不過的難.
為什麼呢.
他就解釋.
我沒時間和大家說.
但是二十五節.
就有一個很重要的啟示.
保羅有這個很強的信念.
就是說.
就是剛才第十章的撒迦利亞書.
所描述的.
將來有一個.
不知道什麼原因的新的出埃及.
新的拯救.
整個民族得到解救.
是不是一夜之間得到解救呢.
這些我們都很難去講述了.
二十五節羅馬書十一章.
弟兄們.
我不願意你們不知道者傲備.
恐怕你們自以為聰明.
就是以色列人有幾分是硬心.
等到愛邦人的數目添滿了.
於是以色列全家都要得救.
如經上所記.
必有一位救主從錫安出來.
要消除那國家一切的罪惡.
他是引的是以暢亞書五十九章的保羅.
又說.
我除去他們罪的時候.
這就是我與他們所立的約.
這是二十七章的以暢亞書.
第二十八節.
就著福音說.

$^{1681}$這裡說的以色列全家得救.
是說全家歸主.
不是那個物質.
生命的得救.
是說整個民族.
全家的歸神.
所以是怎麼可能做到呢.
就唯有在神的時間到.
在一個時間.
愛邦人的數目添滿了.
保羅說.
以色列就全家歸主了.
就給了一個這樣的offer給我們.
保羅.
接著第二十八節.
他解釋.
就著福音說.
他們為你們的緣故是仇敵.
我們蒙福因為福音.
因為耶穌拯救我們.
白白無條件犧牲拯救我們.
我們信.
這是福音.
我們就得到這個救恩.
但是以色列人不信耶穌.
就著我們來說.
耶穌是我們的救主.
是我們的恩人.
是我們的牧人.
但是猶太人就看耶穌是他們的敵人.
所以就著我們的緣故.
以色列人就是仇敵.
神福音的仇敵.
記住不是我們的仇敵.
我們是不應該有敵人.
我們要為猶太人祈禱.
我們要為萬族萬民祈禱.
他們是福音的仇敵.
是耶穌的仇敵.
就著簡選.

$^{1721}$簡選就是說.
神和亞伯拉罕所立的約.
和摩西所立的約.
和大衛所立的約.
和神自己在舊約的應許.
對這個民族.
將來他們有一個新的出埃及的拯救.
這個時候.
這個就是簡選.
我們就不是簡選.
我們是得到神的福音.
當然在新約的角度.
我們都是被簡選的.
我們都是被召的.
這裡告訴我們.
就著簡選來說.
他們為列祖的緣故.
是蒙礙的.
現在是神的敵人.
現在是福音的敵人.
我說的是猶太人.
但是因著亞伯拉罕的信心.
原來可以在神的裡面一直有效.
所以神的簡選是不會廢棄的.
神應許了亞伯拉罕.
他的後裔會蒙福.
神是會算的.
算數的.
他會繼續延續.
所以現在猶太人就硬心了.
但是福音就傳到外邦了.
但是原來有一個日子.
當那個簡選的時間時機到了.
即是滿足了.
成熟了.
神就會出手.
所以就著列祖的緣故.
猶太人是蒙礙的.
我們要有這些眼光.
保羅有這些眼光.

$^{1761}$他絕對沒有偏袒任何一個民族.
他是以事論事.
很難得的.
他自己是有很大的包袱的.
他自己本身是一個拉比.
他在未信耶穌之前.
是真的不明白.
為什麼救恩要臨到外邦人.
為什麼神要憐憫這些外族人.
但是當他在耶穌裡面再看人類.
看神的救恩的時候.
他就知道.
這個是無條件的簡選和福音.
是任何人都可以的.
所以猶太人你有紅眼睛.
看不起外邦人.
其實是你搞錯了.
將自己當作是主人.
其實那個恩典是神的.
神是創造萬物的主.
他怎麼會容許猶太人.
好像自己成為了當家作主.
所以猶太人就想多了.
以為自己是唯一的神的選民.
這些都是在背後的價值觀.
我們都要小心.
慢慢我們信了主內.
當你常常做領袖.
你就慢慢不知不覺會有一種權力慾.
你要小心了.
你慢慢可能會有一種猶太人的迷思.
那種我們是選民.
我們是蒙福的.
這樣其實是不健康的.
本來一個領袖.
本來就是要運用神給他的.
或者匯眾給他的某些權.
去做事去辦事.
這個無可厚非是對的.
但是辦完事之後.

$^{1801}$你一定要看自己是一個普通的匯眾.
回不了頭.
其實真的很容易就會失腳.
失腳的意思就是慢慢不知不覺.
有一種.
又不是叫做Label.
不知不覺地分開了.
我們是特別一點.
你是普通一點.
這個是信主某程度的一個小心的大忌.
猶太人就是這樣落在整個民族.
面對咒詛被撇棄.
所以保羅自己是很痛心的.
二十九節.
因為神的恩賜和選照.
是沒有後悔.
所以正正就是.
福音和揀選.
都是在神的愛裡面去運行的.
所以應許了列祖.
和他現在要大師恩典給外邦人.
全部是神的.
神有什麼理由.
我要問清楚你猶太人.
我可不可以師恩給外邦人.
是沒有這件事的.
所以猶太人他們慢慢.
以為自己是唯一的選民.
這個價值觀和這個心結.
其實很多後世的教會是需要警醒.
需要小心.
是需要慢慢不要變成家族式經營.
這個是很大的一個要留意的地方.
跟著就告訴我們.
第三十節.
你們從前不信服神.
如今因他們猶太人的不信服.
你們倒蒙了連戌.
這樣他們也是.
即是現在不信服.

$^{1841}$叫他們因著司給你們的連戌.
現在也就蒙連戌.
因為神將眾人都圈在不信服之中.
特意要連戌眾人.
跟著保羅就引述.
這一段經文.
心災神豐富的智慧和知識.
他的判斷何其難測.
他的蹤跡何其難尋.
誰知道他的心.
誰作過他的謀事.
誰是先給他.
使他後來償還.
因為萬有都是本於他.
依靠他歸於他.
願榮耀歸給他直到永遠.
阿們.
說到結尾.
這樣來說阿們.
這樣來說去稱頌這位的眾物主.
即是說猶太人是有盼望的意思.
即是說猶太人將來.
你不要理會.
將來他們全家會歸主.
是會得救.
但其實舊約.
我們現在看的撒迦利亞書十章.
將來就有一個真真正正的.
按照字面應驗的.
新出埃及.
周圍的列國都要神服於以色列.
他們到時就齊齊整整.
完全符合得到神的恩典的行為.
這樣的一個改變.
但現在不是.
現在還未見到這個情況.
現在似乎是慢慢朝野.
全世界對他的限制.
因為越做越離譜.
那個智慧變成一個.

$^{1881}$藉口的一個入侵.
和一個聖靈淘汰.
這個絕對不是神這樣做的.
所以將來那個全家歸主.
我個人認為.
就一定不是用戰爭.
而是以色列人不斷地輸.
輸到一個地步.
他們就要呼求神.
即是回到申命記二十八章.
回到歷代志下的第六章.
他們會向天上的神.
真誠的禱告祈求.
十二章有一個叫洗罪的傳言.
很特別.
十一章就是在說.
以色列人大獲全勝.
但十二章說的就是.
真正的勝利.
是國民的心歸向神.
所以神就開了一個傳言.
是洗罪的傳言.
是很特別的結構.
今晚就不夠時間說第十一章開始.
要和大家說一聲.
因為接下來的星期日.
一直到二十四日.
我都會和一班弟兄姊妹.
去到希臘和七教會.
即是土耳其和保羅的行蹤.
所以有兩個星期.
星期四晚都要暫停.
我是第二個星期四的.
下午黃昏就回到香港.
但可能都會有時差.
所以要停兩次的查經班.
盼望大家回去有機會去看一下.
剛才我在下午就將.
上星期問的問題.
到那天在九到十四章.

$^{1921}$出現了的次數.
我幫大家列出來.
給大家這兩個星期的功課.
就是回去看經文.
到哪天.
其實在先知書.
先知書比較明顯.
有人計算過超過六十次.
關於到哪天的.
大家如果有興趣.
可以逐卷去看一看.
關於將來的尼塞爾瓦降臨之前.
以色列人要面對的情況.
但因為第九到十四章.
出現了十多次.
所以我先給大家一個概念.
我們由十一章開始.
我們會再說.
今晚就停在這裡.
我們有兩個星期停一停.
大家其實可以回聽那些講道.
那些查經.
有時我都聽回.
我才記得我自己講到哪裡.
都不錯的.
偶然有些頂展會留言.
但因為我很少會上去看.
所以我一般都不會.
去看那些留言.
所以如果你有些問題.
你在群組裡問.
我就盡快回答.
但如果你不是在群組裡.
而你是經常聽的.
你不要介意你給個電郵.
或者你WhatsApp我.
我就加你入群組.
我們有甚麼材料呢.
通常我會在群組裡發佈.
更好就是大家有問題討論.

$^{1961}$關於聖經的.
我盡量一定會解答.
希望大家一起學習聖經.
我想預告一下.
講了舊約一段長時間.
歸回的歷史書和先知書.
當然我們還有一卷馬拉基書未講.
一定要交代整個歸回的歷史和先知.
但之後我想講一卷新約.
大家可以去想想.
我給大家幾個選擇.
我們講過羅馬書.
不過那時剛開始.
沒有錄音.
全部一張都沒有.
完全沒有錄音.
很可惜.
我不介意重新講一次羅馬書.
因為能夠講都是一個福氣.
看看大家想不想.
不然就推而求其次.
我真的很想講的就是希伯來書.
很深但又很重要.
我又很想講.
但如果都不行.
大家覺得不好.
這些是可以投票可以民主的.
啟示錄我們有講過.
約翰福音又詳細講過.
這些全部在YouTube的記錄裡.
大家聽回就可以了.
羅馬書和希伯來書都沒有講過.
羅馬書是沒有記錄.
是很重要的書卷.
有系統將舊因講明.
第三卷我想講的.
可能是帖撒羅尼迦加前後書.
加上末世論.
變成是神學加上聖經.
這三卷我暫時在這段時間.

$^{2001}$等候裡面有這樣的感動.
當然大家可以提議.
我希望是新約.
你可以提議哪一卷.
新約或者是這三卷.
大家可以投票.
都幾好玩的.
大家可以不要介意.
按你自己的感動.
你喜歡哪一卷你自己投票就可以了.
我嘗試今晚整理好WhatsApp的投票.
大家可以去嘗試.
盼望大家一直在學習裡面.
都一起去成長.
一起祈禱.
多謝天父給我們今晚.
在第十章的撒迦利亞書.
看到將來有一個以色列的復興.
但我們更加看到.
申命記的二十八章.
和歷代治下的第六章.
神對於這個弱的.
特別是立約的人.
那個生命.
那個靈性的體重.
我們想到.
神你藉著撒迦利亞.
來勸勉當時.
在重建聖殿的一班人.
這班是歸回.
回到耶路撒冷.
來居住的國民.
某程度他們就是.
響應神的那個祝福.
重新離開自己的家園.
回到自己的故土.
來重建.
有他們很重要的價值.
但我們都很盼望神.
透過我們今晚所看到的.

$^{2041}$綜合了的訊息.
讓我們真的能夠在神的面前.
看到你看顧神的家.
你看顧我們的人生.
但是我們需要履行這個新約.
我們和主耶穌所立的.
我們要一生敬愛你.
和進行你的旨意.
更加將這個福音和人分享.
使我們自己的生命都能夠活潑.
願神幫助我們.
在未來的兩個星期.
求主保守弟兄姊妹.
無論是自己的溫習.
自己的讀經.
或者我們一班三十多人的弟兄姊妹.
我們在希臘追隨保羅的腳蹤的時候.
都給我們有很深的體會.
我們會去到七教會的一些廢墟.
爾忽所,露底加,善美娜.
這些地方的時候.
都給我們有很大的觸動.
因為我們真的腳踏著.
當時的教會.
保羅在那裡傳道的地方.
讓我們都能夠從其中.
有一個很大的靈性的思考和反思.
我們仰望著奉耶穌基督的名.
阿們.
\newpage



\section{馬可福音 16:9-20}
\label{sec:Jspoigh_Nw0}
\textbf{2023-11-12 午堂《傳揚復活的基督》馬可福音 16\_9-20 講員:曾裔貴牧師}
\newline
\newline
連結: \href{https://youtube.com/watch?v=Jspoigh-Nw0}{\texttt{https://youtube.com/watch?v=Jspoigh-Nw0}} ~~~~ 語音日期: 2023-11-14
\newline
\newline
\hyperref[sec:z09lg9RcA3o]{\small{< < < PREV SERMON < < <}}
~
\hyperref[sec:index_chronic]{\small{[返順時目]}}
~
\hyperref[sec:index_scriptual]{\small{[返順卷目]}}
~
\hyperref[sec:lMH9_3_EQe8]{\small{> > > NEXT SERMON > > >}}
\newline
\newline
馬可福音 16:9-20
\newline
\begin{longtable}{cl}
\hline
\hline
章節 & 經文 (和合本修訂版)\\
\hline
16:9 & \begin{tabularx}{0.7\textwidth}{X} 〔 凡耶穌所吩咐的，她們簡潔地告訴彼得和他周圍的人。這些事以後，耶穌親自藉著他的門徒，從東到西，把那神聖、不朽、永遠拯救的福音傳出去。阿們！〕〔 在七日的第一日清早，耶穌復活了，先向抹大拉的馬利亞顯現；耶穌曾從她身上趕出七個鬼。 \end{tabularx} \\ \\ \relax
16:10 & \begin{tabularx}{0.7\textwidth}{X} 她去告訴那向來跟隨耶穌的人；那時他們正哀慟哭泣。 \end{tabularx} \\ \\ \relax
16:11 & \begin{tabularx}{0.7\textwidth}{X} 他們聽見耶穌活了，被馬利亞看見，可是不信。〕 \end{tabularx} \\ \\ \relax
16:12 & \begin{tabularx}{0.7\textwidth}{X} 〔 這些事以後，門徒中有兩個人往鄉下去；正走著的時候，耶穌以另一種形像向他們顯現。 \end{tabularx} \\ \\ \relax
16:13 & \begin{tabularx}{0.7\textwidth}{X} 他們去告訴其餘的門徒，那些門徒還是不信。〕 \end{tabularx} \\ \\ \relax
16:14 & \begin{tabularx}{0.7\textwidth}{X} 〔 後來十一使徒坐席的時候，耶穌向他們顯現，責備他們不信，心裡剛硬，因為他們不信那些在他復活以後看見他的人。 \end{tabularx} \\ \\ \relax
16:15 & \begin{tabularx}{0.7\textwidth}{X} 他又對他們說：「你們往普天下去，傳福音給萬民聽。 \end{tabularx} \\ \\ \relax
16:16 & \begin{tabularx}{0.7\textwidth}{X} 信而受洗的必然得救，不信的必被定罪。 \end{tabularx} \\ \\ \relax
16:17 & \begin{tabularx}{0.7\textwidth}{X} 信的人將有神蹟隨著他們：就是奉我的名趕鬼；說新方言； \end{tabularx} \\ \\ \relax
16:18 & \begin{tabularx}{0.7\textwidth}{X} 手能拿蛇；若喝了甚麼毒物，也不會受害；手按病人，病人就好了。」〕 \end{tabularx} \\ \\ \relax
16:19 & \begin{tabularx}{0.7\textwidth}{X} 〔 主耶穌和他們說完了話以後，被接到天上，坐在神的右邊。 \end{tabularx} \\ \\ \relax
16:20 & \begin{tabularx}{0.7\textwidth}{X} 門徒出去，到處傳福音。主和他們同工，藉著伴隨的神蹟證實所傳的道。〕 \end{tabularx} \\ \\
[1ex]
\hline
\hline
\end{longtable}
$^{1}$請每位主人平安.
早安 我們一起祈禱.
感恩天父讓我們能夠今早再一次聚集.
我們在十月有堂慶.
上星期雄恩也有十二週年的堂慶.
我們昨天晚上有一個很大型的福音報道會.
音樂的報道會.
我們又見到很多很多的家長.
帶著很小的小朋友.
都一起在整個晚上感受到聖靈的工作.
這個地方是能夠傳遞這位復活的救主的福音.
這個道是神你自己賜給我們人.
我們要去宣講.
你就會與我們同在.
今早我們以這個題目再一次彼此來思想的時候.
願你在十一月稱為一個感恩節的月份.
以及我們在錦宣家.
我們有猜拳的年會等等的事情.
我們一起來思想.
復活帶給我們的生命意義.
我們一起等候禱告奉耶穌基督的名.
阿們.
剛剛剛才祈禱的時候也有提到.
上個星期洪恩的十二週年堂慶.
當然我代表教會和一班弟兄姊妹領袖去崇拜.
以及去有一些的協助.
講道的正正就是老耀洪牧師.
他坐在我旁邊.
上台講道.
見到他的身軀很瘦削.
他的西裝已經不合身了.
太鬆了.
感受到神的僕人那份盡心盡力去服侍.
講道非常有震撼力.
感受到有一份很大的力量驅拆著他.
繼續盡他一口氣去服侍.
過兩天.
因為我是第一個收到消息.
轉發給不同的群體.
特別是領袖的群體.

$^{41}$同工的群體.
一直每天跟著他的進度.
又看到神很奇妙的拯救.
透過醫生.
透過藥物.
透過手術.
透過香港的醫療體系.
使老牧師現在他能夠自己去發短訊.
這是很大的生命力.
弟兄姊妹.
我們這個月裡面.
這是感恩節.
西方社會很重視感恩.
不過去到十一月.
要完結他們的感恩節之後.
他們很緊張他們的願望清單.
我在聖誕節的時候收到什麼禮物呢.
列了很多很多的清單.
而很少能夠再回頭想想.
生命的感恩.
生命的意義.
對於上主賜給我們每一天的生命.
我們能不能夠.
就算好的日子.
不好的日子.
我們都能夠充滿一份感恩呢.
很想今天藉這段聖經.
來到大家一起彼此的勉勵.
這段聖經有很特別的地方.
去到最後.
整個馬可福音其實就是很想告訴我們.
耶穌想告訴他的門徒.
跟隨的門徒知道.
他是哪一位.
他就是那一位的彌賽亞.
但是來到第十六章.
作者仍然告訴我們一個.
當然是當下的結果.
就是耶穌死而復活了.
門徒一概的不信.

$^{81}$這個不信出現了好幾次.
就令到我們覺得很困惑.
為什麼不是寫得好看一點呢.
其中在這段聖經.
剛才主席讀了給我們知道的.
在比較傳統的重要的抄本裡面.
有些經文是沒有的.
但是在《使徒行傳》裡面.
的而且確裡面所提到的每一個神蹟.
神或者耶穌伴隨著使徒.
他們去傳道的時候.
所來到去施行的神蹟.
在《使徒行傳》你可以找得到.
但是這一卷書的最後這一段.
比較傳統的重要的抄本是沒有的.
由九節到第二十節.
正正我們今天說的這段聖經.
在傳統的抄本是沒有的.
當然有其他的抄本.
不是那麼權威的就有的.
所以和合本選取認為是有.
所以就放在這裡裡面.
所以成為我們今天中文的群體.
就是中國人的群體的華語.
我們是有這段聖經.
作為我們的權威來了解的.
相信最重要的就是第二十節.
門徒最終他們能夠明白了.
主耶穌的死而復活和升天.
然後他們就到處宣傳福音.
留意這裡是主耶穌和他們同宮.
然後用神蹟追隨著他們.
證實所傳的道.
不是要誇讚他們行神蹟的能力.
有哪些是驅魔人.
有哪些能夠化腐朽為神奇.
不是那些 那些是次要.
那些是伴隨著道本身.
是證實所傳的道.
是那個道改變人生命.

$^{121}$這個是我們要留意.
要一起著重教會是要有道.
所以這段聖經.
我們要回頭看看.
有好幾個驚訝.
作者特別想highlight給我們知道.
父女的驚訝 門徒的驚訝.
再加上最後主耶穌對他們的責備.
使他們終於有一個很大的.
慢慢的整個心靈改變.
但是大家有沒有留意.
其中有一個人.
應該是最大的驚訝.
就是在主耶穌復活之後.
因為這個人被點名.
天使點名說一定要見他.
要他在加利利出現.
你會猜到是誰.
就是這位彼得.
很重要告訴我們.
我們要留心這件事情.
所以我們來求神幫助我們.
在我們的人生復活.
我們有沒有錯過呢.
你錯過復活節.
不要緊 那個是節期.
但是如果我們的生命錯過了.
這位復活的主.
向我們的心靈顯現.
那這個就是一個很大的錯失.
是一個很大的人生的損失.
剛剛昨天看到一個新聞.
都很有趣.
但又很反映到人性的貪婪.
在南非開普敦.
有一個女生大學生.
18歲.
她就不知道什麼緣故.
可以拿到一些政府提供的.
一些救濟的食物.

$^{161}$每個月有1400個.
克里南非的錢幣.
那個幣值.
1400當400元.
1400元.
但突然有一天.
在2017年的6月1號.
她去打保看看戶口.
那個月入了1400萬.
1400元變成1400萬.
她當然不相信.
當然去按機.
按到.
每個月1400元.
突然間那個戶口有1400萬.
你猜一下她會怎樣做呢.
相信一個沒有神生命的人.
是政府的程序搞錯.
她瘋狂的購物.
瘋狂的用.
用到當政府真正發現的時候.
是73日後.
她用了多少錢呢.
1400萬.
結合港幣是600萬.
她在73日裡面.
未被發現.
她感覺很驚訝.
為什麼會這麼有錢呢.
突然間成為了富婆.
用了53萬.
每天買名牌.
盡量去揮霍.
終於第73日.
被政府發現了.
還不是政府突然間找到的錯誤.
是超級市場裡面.
她無意中將那個收據放下.
人家看到.
為什麼這個女孩會這麼有錢呢.

$^{201}$才告發她.
法官來到她說.
她那種貪婪.
是有意識地去做這些事情.
每天每天這樣做.
跟她的同學朋友.
她在這個驚訝裡面.
她所表現出來的.
法官對她很無情的譴責.
你作為一個這麼有知識的大學生.
你本來是拿著救濟金的.
那個是基本給你的食物的.
1400元.
但是入多了很多零之後.
你居然間這樣據為己有.
還要來盡量去揮霍.
判了她五年監禁.
後來輾轉.
因為她的稅數.
就可以來去守行為.
香港應該不行了.
可以不用坐牢.
還有就是.
不知道什麼原因.
那個辯護的律師幫她辯護.
可以令到她那53萬.
不用還給政府.
那就是說她賺了五十多萬.
這個就是.
我們想一想這個驚訝.
帶給她的一個.
可以以為是濫用的一個.
這樣的行為.
來去將本來納稅人.
要資助這些有困難的大學生.
的一些生活費.
她拿來這樣濫用.
各位你看一看這段聖經.
在之前那段.
有一個叫做茉達拉的瑪利亞.

$^{241}$和另外兩個婦人.
剛才我們沒有讀的.
第一到第八節.
給我們有很大的感動.
這幾位的婦女.
曾經主耶穌幫她.
她們是帶著一份的等待.
因為安息日什麼都不准做.
安息日一完.
她們就第一時間.
很早很早.
她們三個人.
這個第二節.
就跑到墳墓.
然後令她們很驚訝.
很驚訝的地方.
你可以預想得到.
構想得到.
她們在墳墓裡面.
預計會發生的就是.
那個大的石頭.
其實是那個墓門.
因為是用自然的.
造成的一個山洞.
用一個大很大的石頭.
是有羅馬的兵丁.
來蓋了一個封條.
基本上是沒有辦法打開的.
她們想著的只是.
到時如果有兵丁在的時候.
能不能夠打開墓門呢.
因為裡面的人已經死了.
也過了安息日.
我們想去告末.
其實不是再用什麼香料.
其實是防腐.
防臭.
她們是想著這樣.
但相信令她們很大的驚訝.
就是她們在清早見到的.

$^{281}$就是第四節.
那石頭原來很大.
她們抬頭一看.
卻見到石頭已經滾開.
各位你有沒有想過.
你去尋找耶穌的屍體.
你會有一個這麼大的驚訝呢.
復活的主.
給這三個婦人.
有一個很大的意外驚喜.
就是她們想著用什麼辦法.
來接觸到耶穌的屍體.
就用她們預備好的香料.
幫助耶穌防腐防臭.
但是第二個驚訝.
大家都可以見到.
第五節.
沒有了屍體.
見到的是一個天使.
一個少年人.
天使變成一個少年人的樣子.
可以見到.
可以對話.
你想一下這個驚訝.
沒有了石頭.
沒有了屍體.
但有一個少年人坐在墳墓裡面.
天未亮.
你可以想像到.
如果去和合石.
見到這樣的情景.
你會害怕.
這個驚訝.
這個意外的收穫.
更加重要的是.
當這個少年人和他們說的時候.
這個少年人穿著白袍.
就甚驚恐.
很正常.
三個婦人.

$^{321}$他們來到這個墳墓的時候.
那個少年人就和他們說.
不要驚恐.
你們尋找的這位釘十字架的拿撒勒人耶穌.
他已經復活了.
不在這裡了.
這個更加大的驚恐.
死了.
大家都預期.
一定是死了.
人不能復生.
死了.
在他們有限的眼光.
能夠預想的.
耶穌多偉大都好.
死亡吞噬了他.
整個十字架的刑罰.
已經將耶穌完全吞噬了.
大家接受的現實.
包括門徒.
何況幾個婦女呢.
能夠做的就是.
盡量讓耶穌.
不是繼續被人這樣.
不光彩的埋葬.
可以是這樣.
但是.
各位你覺得很奇妙的一件事.
就是少年人向這三個婦女的宣告.
一心充滿了愛對主耶穌的敬重.
換來的是.
一而再的驚訝.
換來的就是.
這個驚訝到底是一個幻覺.
還是一個事實呢.
是不是生命的主.
真的向我們三個婦人顯現呢.
半信半疑.
似乎是當時的狀態.
是當時他們的內心.

$^{361}$已經復活.
不在這裡了.
在哪裡呢.
接著少年人就說.
你們來看他的地方.
就是那個空的墳墓.
什麼都沒有.
第七節.
你們可以去告訴他的門徒.
剛才我說的.
最重要令人驚訝的就是彼得.
你們可以去告訴他的門徒.
和彼得.
為什麼要特別強調彼得.
一定要去加利利.
去見復活的耶穌呢.
其實回頭看.
你就會看到.
在第十五章.
已經告訴我們.
彼得是這樣來否認耶穌.
這是我們大家一早已經知道的事實.
在這裡的時候.
這個少年人.
這個天使.
和這三個婦人.
來講的時候.
他們帶著的是充滿恐懼的驚訝.
來聽到.
見到眼前的事實.
然後他們就來到去.
再聽到一個.
他們要帶的訊息.
回去耶路撒冷的.
似乎就是馬可樓.
來報這個訊息.
三個婦女帶著這個訊息.
回去告訴門徒.
然後就告訴我們.
第八節.

$^{401}$他們就出來.
從墳墓那裡逃跑.
又發抖又驚奇.
什麼也不告訴人.
因為他們害怕.
這個驚訝.
這個驚恐.
一直都沒有消退.
到底他們的內心.
當時的心情.
到底對這個訊息.
是深信不疑.
還是他們在恐懼裡面.
以為只是一個幻覺.
發生在他們的內心.
是他們一個心靈的投射.
太愛耶穌了.
太掛念耶穌了.
所以有這個幻覺.
但似乎又不是.
這個人真的邀請他們去看.
和來到去.
真的給他們一個訊息.
帶回馬可樓上.
現在聚會的十一個門徒.
因為猶太已經死了.
就是我們開始看的.
第九節這段聖經.
提到七天的第一天清早.
主耶穌復活的這個情景.
他再一次提到.
向這個末大拉的瑪利亞.
來顯現.
來告訴我們.
馬可特別強調.
末大拉的瑪利亞.
他在他的人生.
怎樣遇過耶穌.
耶穌怎樣改變他的生命.
耶穌未死之前.

$^{441}$在他傳道的日子.
怎樣來讓這個末大拉的瑪利亞.
得到一個生命的改變.
留意聖經怎樣描述.
特別現在回頭要描述.
末大拉的瑪利亞.
他的生命.
這裡告訴我們.
耶穌從他身上.
曾經趕出七個鬼.
他就去告訴.
那些向來跟隨耶穌的人.
向來.
即是一直跟耶穌傳道的.
十二使徒.
或者有些沒有提名的.
其他的門徒.
向來跟耶穌的.
為何現在不在墳墓裡面.
跟他一起.
在安息日過了之後.
就在這個主日.
來找耶穌基督.
和來告末他的身體.
安息日不能做.
為何這些事情.
會變成這幾個婦女.
不是那些向來跟隨耶穌的門徒.
來做領耶穌的屍體嗎.
這裡告訴我們.
婦女的生命.
對於耶穌基督的感動.
那個復活.
帶給他們的改變.
是很大的一個觸動.
婦女很容易就會受到.
主耶穌的愛的幫助.
而他們有生命的一個回應.
末大拉的瑪利亞.
我們可以想一想.

$^{481}$他遇見耶穌的時候.
雖然只是一輪半找.
不是很多的描述.
但作者好像故意想告訴我們.
這一位.
在清晨來找耶穌的.
在他人生曾經.
是受控制.
受邪靈的驅控.
沒有他的自由.
沒有他的心靈的自主.
一直都是在那個.
不知道什麼靈界.
有什麼背景.
令到他要受靈界的驅控.
以至他內裡有鬼.
還有七隻鬼.
在他裡面.
但是耶穌釋放他.
讓他得到自由.
往後.
他就是這樣跟隨.
和其他敬虔的婦女.
服侍耶穌.
路加福音第八章.
路加醫生有記載.
耶穌出去傳道.
有好些婦女用自己的財物.
供應耶穌和他的門徒.
相信末大拉的瑪利亞.
特別強調.
他不是其他的瑪利亞.
是這個末大拉.
這個地方的瑪利亞.
可能他以前.
很多人都認識他.
他是有名有姓.
是有鬼父很長時間的.
但現在他的生命.
交付給主耶穌.

$^{521}$成為愛主.
服侍主的其中一個婦女.
當然上文提到的.
有撒羅米.
似乎是瑪可.
他們那個家族的一個.
不知道是不是媽媽還是怎樣.
但這裡告訴我們.
你看看第十節這裡.
他就去告訴向來.
跟隨耶穌的人.
那時他們正哀動哭泣.
男士的感情世界.
有時不是這麼豐富表達.
如果連續哭幾天.
是很難得的.
通常半天哭完就不再哭.
當然不是一概而論.
一般的理解.
女性的感情世界豐富.
有時一想起就受觸動.
又會哭.
但男士就很難.
所以這個時候很近.
因為剛剛安息日.
可能還有另一個大的安息日.
因為是逾月節.
所以可能有重疊的兩個安息日.
接著去到七日的第一日.
即是主日.
所以門徒還在房間裡面.
十一個門徒還哀動哭泣.
那件事還是很新鮮.
歷歷在目.
耶穌被人這樣釘死.
很慘.
很不公義.
但很快就會忘記.
不過十一節.
這樣去描述門徒當下的狀態.

$^{561}$這段聖經其實很想.
讓我們大家一起反思.
如果我們遇見復活的主.
復活的主的真理訊息.
如果我們視而不見.
我們好像沒有經歷過.
其實人生就這樣流走.
一次,兩次,三次.
這樣來向我們傳遞.
第一次就是瑪利亞.
趕去跟門徒說.
正正在哀動哭泣.
歷歷在目的我們的主.
這樣死於非命.
男士到現在還在哭泣.
這件事一定是很大的事.
還有十一個門徒一起哭泣.
還有其他人.
你可以想像得到.
但是落差實在太大.
聽見耶穌復活.
即是茉蒂拉的瑪利亞.
趕去樓房跟這些門徒說.
耶穌復活了.
聽見耶穌活了.
被瑪利亞看見.
卻是不信.
正正反映我們男士的分析力.
眼見為真.
我們是用理性來行事.
哭歸哭.
死了就是死了.
復活是你想多了.
沒可能的.
很理性地不應該這樣想.
作者是很如實反映人性.
那種對復活的驚訝.
復活真的會帶給我們的生命.
有一個由驚訝而來的生命力.
所以這個不信.

$^{601}$是非常合理的人性.
合理到一個地步.
第二次了.
第十二節.
路加會比較詳細一點.
叫做以馬五思的事件.
這是以後門徒中間有兩個人往鄉下去.
走路的時候.
耶穌變了形象.
向他們顯現.
他們就去告訴其餘的門徒.
其餘的門徒也是不信.
這兩個不是使徒.
不是最內圍的十一個耶穌.
選擇的三年跟隨耶穌的使徒.
這兩個是普通的門徒.
都是跟耶穌.
但是沒有跟得這麼內圍.
這兩個人.
其中一個叫格劉巴.
其實是完全不知道是誰.
因為路加福音二十四章.
偶然間提到.
有兩個人去鄉下做一些事情.
接著耶穌就改變形象.
向他們兩個顯現.
跟他們一路走一路聊天.
進到屋裡面.
一直都不發現耶穌.
耶穌跟他們講解聖經.
一路講解.
接著準備進屋.
他們要去吃餅的時候.
用餐用飯的時候.
見到耶穌舉起手.
他們認得這對手.
認得是耶穌.
趕緊跑回去耶路撒冷.
由以馬五師去耶路撒冷.
有二十五里路.

$^{641}$跑上去告訴門徒.
不過馬可福音很簡潔地用了一句.
告訴我們.
耶穌向我們顯現.
他們就去告訴其餘的門徒.
跑了二十五里回去耶路撒冷.
其餘的門徒不信.
這個驚訝第二次臨到.
但是門徒不為所動.
弟兄姊妹如果復活的主.
每個禮拜向你顯現.
希望你的生命交付給他.
讓他來掌管.
當然不保證你的人生.
從此一帆風順.
因為每個人有神安排的跑道.
每個人有不同的人生.
是神安排的.
我在查經班也有提過.
有機會透過聖經.
說一些自己人生成長的見證.
我在八歲那年.
很深刻.
但是當時沒有人告訴我.
如果告訴我.
我的人生不會再這樣.
我家是很大的家族.
有十三個兄弟姊妹同住.
我是排第九.
我有四個哥哥四個姐姐.
兩個弟弟兩個妹妹.
全部現在還在.
所以那個家族一直很大.
不是那個多的問題.
就是在很窮的母區長大.
現在的友一城.
那個時候叫做小西湖村.
不知道什麼原因.
我爸爸認識一個很有錢.
在公務局做到很高級的.

$^{681}$當時的政府官員.
不知道什麼原因.
他的太太姓梁的梁太.
很喜歡我.
要收養我.
當時八歲.
沒有人跟我說這個事情.
沒有人叫我去選擇.
去還是不去.
他要帶我移民去加拿大.
當時七十年代.
手指那時候可能是七一年左右.
1971年.
完全沒有人告訴我.
到我長大了.
很大了.
我媽媽跟我說.
小時候本來有人收養你的.
我當然不肯.
後來想想.
十三個給一個.
沒所謂.
阿神怎樣去管理我的未信的父母呢.
我早幾年去加拿大多倫多.
我姐姐那裡.
我姐姐原來跟梁太很熟.
我姐姐跟我說.
梁太當年說收養你的那個.
就住在我們附近.
剛剛離世.
我說你們為什麼這麼遲告訴我.
你早點跟我說我的人生都不同了.
我不再姓曾了.
你想想我七一年在加拿大成長.
我不是在武區長大的一個街童.
中學階段.
每天要站在主任室門口的那個就是我.
但是現在我回想這麼多年了.
神的安排是最好的.
如果不是這樣的成長.

$^{721}$怎會有日後我的信主.
信得這麼玲瓏浮凸.
第二年信主.
1989年.
應該是1988年.
1987年的12月24號.
平生第一次去所謂教堂.
其實不是教堂.
是一間小學的天台禮堂.
那個聚會是差不多完結呼召.
我才衝上去.
全身大汗就問那個牧師叫李寧牧師.
他說現在我要呼召各位家長.
是一個家長的平安夜晚會.
不關我的事.
是我的弟弟在那個地方聚會.
就告訴我現在呼召這個世界有神.
我們要信耶穌.
就問哪個家長.
當時我是一個二十多歲的年輕人.
在最後排.
最後去到全身汗都還沒乾.
誰要信耶穌舉手.
我不知道為什麼我第一個舉手.
我這個舉手就到今天.
我沒辦法相信.
這位主耶穌的真實在我的生命裡面.
復活的主居然向我顯現.
讓我來到生命有一個很大的改變.
我們在團契.
晚上的團契我們都有分享過.
一個本來二十多歲.
每天吃開兩包煙.
一天的一個年輕人.
一夜之間.
其實不是夜.
是一日之間.
上午突然間說我不吃煙了.
就不再吃了.
以後不再吃.

$^{761}$聞到煙味肺部會不舒服.
吃了這麼多年煙.
由十三歲開始偷偷地在家裡吃.
在學校吃.
到十三,十四,十五歲.
家裡一起吃.
每個人都吃男生.
不需要再遮遮掩掩.
但可以一下子完全不吃煙.
就是信主的那一下.
生命的改變.
很感恩就是.
我們回頭去看.
門徒不信.
但是天使傳遞的訊息給這三個婦女.
一定要叫彼得去加利利.
耶穌約定他要見他.
要再使用他.
相信這個驚訝.
是讓所有的門徒.
尤其是彼得他人生的生命.
一個最大最大的改變.
親愛的弟兄姊妹.
你看看第十四節.
耶穌怎樣在責備完之後.
大大的使用他們.
你就可以想像到.
復活的生命.
就是想我們每個人都能夠遇見和改變.
後來十一個門徒坐直的時候.
耶穌向他們顯現.
責備他們不信.
心裡面剛硬.
原來原因不是真理難以接受.
原來心靈不夠柔軟.
原來是心靈未改變.
這是我們每一個人要解決的問題.
但是心裡面的剛硬.
不是靠你吃什麼補心的丸.
是遇見這一位真實愛我們每一個的耶穌.

$^{801}$心靈才可以柔軟.
因為他們不信那些在他復活以後.
看見他的人.
就是那幾個婦女.
就是兩個二萬五仙的門徒.
還有其他的人.
第十五節.
他又對他們說.
你們往普天下去.
傳福音給萬民聽.
信而受贈的必然得救.
不信的必被定罪.
信的人就會經歷.
神同在帶來的生命改變.
你自己改變.
可以改變其他人.
而且是一個印證.
是有印證.
你可以想像得到.
十四和十五之間兩節的聖經裡面.
將神的愛.
將神的寬容接納.
是多麼的淋漓盡致.
責備他們的不信.
心裡面剛硬.
然後十五節.
又對他們說.
你們去普天下.
傳福音給萬民聽.
因為唯有你們.
這十一個門徒.
才真正經歷生命的蛻變.
經歷這位復活的主.
那個大愛大能.
這個生命力.
才可以改變其他人.
他們就是這樣出去.
十九節.
主耶穌和他們說完了話.
後來被接到天上.

$^{841}$坐在神的右邊.
二十節就告訴我們.
門徒出去.
到處宣傳福音.
主和他們同供.
用神跡脫著.
證實所傳的道.
可能是一個很慨略的.
將門徒.
在耶穌復活之後的一大段時間.
所做的事.
來去概括地講述.
但是重點在這裡.
門徒是經過心靈的改變的手術.
由剛硬變得柔軟.
弟兄姊妹要拿捏這個部分.
每個星期.
是否帶著一個期望.
來遇見耶穌呢.
有沒有驚訝讓你發現呢.
昨天晚上有一個.
分享見證的.
烏干達的一個女士.
小時候就在瓦都圖長大.
她的村裡長大.
原來她成長的時候.
有兩年時間.
被自己的堂哥.
性侵犯.
足足兩年.
後來輾轉淘來到瓦都圖.
一路成長.
給她的有新的身份.
新的盼望.
在主耶穌裡面有新的一個開始的人生.
所以她長大了.
這次大學畢業了.
來到香港.
我們昨晚她講見證.
我們有沒有遇見耶穌.

$^{881}$在你的生命.
在你的生活裡面.
讓你感到驚訝.
祂真是愛你.
祂真是認識你.
祂真是在你的生命裡面.
有計劃呢.
讓你的生命有那個的意義.
和改變.
每個人都不同的.
每個人的場景.
神安排都不同的.
剛剛開始的時候講.
Gary牧師.
但其實我們有另外一位牧者.
他在倫敦生存.
那個康復進度非常非常的理想.
可以和他繼續WhatsApp.
很精神.
15號就可以上班了.
是很不同的大家的人生際遇.
記住我們沒有辦法.
知道我們下一步.
明天.
但掌管的在.
做物主和愛我們的.
耶穌的生命.
的手裡面.
所以記住.
天使要帶這個訊息給三個婦女.
記得叫彼得去加利利.
耶穌沒有忘記.
否認她的.
發毒誓.
否認她的彼得.
這個是很大的一份的愛.
教會歷史告訴我們.
馬可.
就是巴拿巴的表弟.
後來是跟彼得.

$^{921}$傳道的.
所以他記得這件事.
是其他福音書沒有記載.
一定要叫彼得.
所以.
我們每一個人的人生.
其實一定要回望.
過去很多的日子.
你在這個地方.
有沒有遇見過耶穌.
遇見神.
你有沒有將你自己的片段.
一點一點的連在一條線.
這是喬布斯.
告訴我們.
他在他的人生.
2005年的時候.
他講的speech.
在史丹福大學.
他終於發現原來他的人生.
一點一點.
連在一條線.
他就發現人生.
其實是成長不是成功.
一個成長成熟的人生.
每一個基督徒.
應該都是成長.
不是成功.
我們是成長.
在舊恩在復活的主裡面.
我們認識.
我們的自己.
在神安排的人生.
去服侍這個世代.
然後慢慢.
帶著尊嚴來老去.
這是我們人生.
應該這樣去面對.
求主.
祝福我們在這個月.

$^{961}$我們要來慶祝感恩節的時候.
我們的內心.
真的要來帶著這一份.
對復活的主的期盼.
有沒有驚訝呢.
復活的主.
有沒有令到你的生命.
有驚訝呢.
因為你的信仰.
已經淡如水.
沒有任何可以來觸動你的.
你的內心.
是不是已經去到剛硬呢.
我們一起禱告.
多謝主.
今早我們很盼望.
我們每一位弟兄姊妹.
再一次遇見你.
你用不同的方式.
來向我們顯現.
講說話.
使我們再一次被你的愛.
被你的真理.
去感動.
因為你的道.
就是可以改變人生命.
多謝你聽禱告.
以耶穌基督的名.
阿們.
\newpage



\section{哥林多前書 9:24-27}
\label{sec:lMH9_3_EQe8}
\textbf{2023-11-12 加恩樓《為主奔跑》哥林多前書 9\_24-27 講員:陳一華牧師}
\newline
\newline
連結: \href{https://youtube.com/watch?v=lMH9_3-EQe8}{\texttt{https://youtube.com/watch?v=lMH9\_3-EQe8}} ~~~~ 語音日期: 2023-11-15
\newline
\newline
\hyperref[sec:Jspoigh_Nw0]{\small{< < < PREV SERMON < < <}}
~
\hyperref[sec:index_chronic]{\small{[返順時目]}}
~
\hyperref[sec:index_scriptual]{\small{[返順卷目]}}
~
\hyperref[sec:zCZXBv4_hT8]{\small{> > > NEXT SERMON > > >}}
\newline
\newline
哥林多前書 9:24-27
\newline
\begin{longtable}{cl}
\hline
\hline
章節 & 經文 (和合本修訂版)\\
\hline
9:24 & \begin{tabularx}{0.7\textwidth}{X} 你們不知道在運動場上賽跑的，大家都跑，但得獎賞的只有一人？你們也要這樣跑，好使你們得著獎賞。 \end{tabularx} \\ \\ \relax
9:25 & \begin{tabularx}{0.7\textwidth}{X} 凡參加競賽的，在各方面都要有節制，他們不過是要得會朽壞的冠冕；我們卻是要得不會朽壞的冠冕。 \end{tabularx} \\ \\ \relax
9:26 & \begin{tabularx}{0.7\textwidth}{X} 所以，我奔跑，不像無目標的；我鬥拳，不像打空氣的。 \end{tabularx} \\ \\ \relax
9:27 & \begin{tabularx}{0.7\textwidth}{X} 我克制己身，使它完全順服，免得我傳福音給別人，自己反而被淘汰了。 \end{tabularx} \\ \\
[1ex]
\hline
\hline
\end{longtable}
$^{1}$今天在崇拜當中和你們一起唱詩歌.
感覺到很雀躍.
還有大家唱歌的氣氛是很興奮.
我想可能會不會受到Watutu影響呢?.
昨晚我們在錦繡堂.
Watutu來表演唱歌.
講見證,不知有多興奮.
特別是小朋友們在歌在舞.
令我們沒試過晚堂這麼興奮和雀躍.
嘉欣樓這裡.
上星期你們有多少人參加Watutu的聚會?.
是不是很精彩?.
非常精彩,很感恩.
天父給我們的恩典是非常足夠.
不論我們在嘉欣,在錦繡.
上帝也給我們很多恩典.
牧師今早和大家分享的題目是.
為主奔跑.
剛才已經有機會讀了整段聖經.
為何保羅會說到24節.
開始用運動員的比方來說呢?.
因為前面的上文.
第九章的第一節到23節.
原來保羅在說一件事.
這件事是保羅說我怎樣去傳福音.
我的經歷如何.
我的心態如何.
我的態度如何.
在1到23節就表露無遺.
不知大家有多少人讀過臨前的第九章.
如果你讀過,你應該有些印象.
就是保羅說什麼?.
保羅說我傳福音,願沒有可誇的.
因我是不得而的.
若不傳福音,我就有禍了.
有沒有印象?.
有些印象吧.
保羅又說,無論對著猶太人.
對著律法的人.
對著沒有律法的人.

$^{41}$對著軟弱的人.
我就是做那些人.
有沒有留意到?.
為了什麼?.
為了得謝人.
所以你會發覺保羅傳福音很有智慧.
他對著不同的對象.
他首先要了解他們.
知道他們的文化.
認同他們.
幫他們有一個共通點.
建立之後,又有信任.
結果保羅就傳福音給他們.
所以你會看到保羅傳福音的模式.
給他們很好的啟發.
還有保羅說.
凡我所作的,都是為福音的緣故.
為要使人同得這福音的好處.
記不記得?.
這個也是保羅的心願.
他一生人的使命.
是要傳福音.
因為他在大馬識路上.
有沒有印象?.
主向他顯現.
光照他.
他就整個倒下了.
趴下了.
眼睛瞎了.
那次的經歷.
是保羅一生人第一次經歷到.
遇見主耶穌之後.
他看到自己的問題.
看到自己的毛病.
看到自己的污穢.
所以他就認罪悔改.
因為他以前迫不及待的.
接著耶穌跟他說.
從此之後.
你要起來.

$^{81}$去到你的同胞.
去到外邦人.
去到君王當中.
傳我的福音.
就是說.
保羅一生人.
領受一個很特別的任務.
就是要傳福音.
所以不怪第九章.
一到二十三節就已經鋪陳了.
很詳細的.
保羅說自己怎樣去傳福音.
這個經歷.
但到了二十四節.
筆鋒一轉.
轉到什麼呢?.
原來保羅想帶出一個.
運動員的形象.
為什麼呢?.
因為他覺得傳福音的人.
事奉上帝的人.
都好像一個運動員.
所以他在二十四節開始.
就拉到這個運動員的角色當中.
可能我們會問.
為什麼他會用運動員.
比方到基督徒呢?.
原來有原因的.
因為有三個原因.
第一個原因是什麼呢?.
第一個原因.
是因為運動員有賞賜的.
有賞賜.
有獎賞的.
因為大家都知道.
出來比賽的運動員.
拼命的.
為什麼這麼拼搏呢?.
因為每個人都想拿獎.
不過很可惜.

$^{121}$獎牌只有三個.
比賽的人有幾十人.
所以每個人拼命.
拼命去拼搏.
希望拿到獎牌.
回來.
第二個原因.
運動員的心態.
他去到盡.
拼到盡.
都不計較.
他目的是什麼?.
希望拿到獎賞.
這種拼搏的精神.
應用在我們信耶穌的人身上.
就很貼切.
為什麼呢?.
因為我們今天做基督徒.
我們事奉上帝去傳福音.
同樣也要有這種態度.
這種精神.
就是全力以赴.
認真的.
拼命的.
為了天國福音的緣故.
來擺上自己.
要有這種精神.
來撥出去.
這是很重要的.
如果你明白這一點.
你就會明白為什麼百多年前.
有些人千里迢迢.
遠道重陽.
來到中國.
這個大地.
去做宣教士.
傳福音.
大家有聽過吧?.
不論是馬勒信也好.
不論是大德生也好.

$^{161}$這些愛中國人的宣教士.
不計一切.
都要來到中國大陸.
這個荒蕪的地方.
來傳福音.
為什麼呢?.
因為他們知道.
將來在上面.
有沒有獎賞呢?.
有獎賞的.
而這個獎賞.
是一個永不扭壞.
永不衰殘的.
生命的冠冕.
所以現在明白.
為什麼這些人.
會這麼認真.
這麼拼命.
不怕艱難.
來傳福音.
因為每一個這樣做的人.
都得到上前.
可能你也會.
發覺你身邊.
有些人.
好像我的朋友.
我遇到的朋友.
本來有一份高薪厚職.
做得不錯.
如果他一直做下去.
他的工作.
真的可以養他一輩子.
真的好說好聽.
打什麼都不用擔心.
你明白我的意思嗎?.
但是不明白為什麼.
做到中途他又不做.
打信,打紙,辭職呢?.
跑去.
報讀神學.

$^{201}$然後全時間.
事奉上帝.
有沒有這樣的人?.
有啊!.
我們會問這些人.
是不是這裡有問題呢?.
不是!.
是因為主光照他們.
所以他覺得.
人一世物一世.
幾十年裡.
原來現在主呼召我.
我要打拼.
拼命.
認真.
去踏上.
蒙召的道路.
所以你明白這些人.
他們這樣做.
是因為他們知道將來有沒有賞賜.
有獎賞.
我又開始明白.
我以前認識一些人.
為什麼無緣無故.
提早五年,提早十年退休呢?.
有時候我都很不明白.
你知道如果提早五年,提早十年.
那他的祥奉就什麼?.
沒有了缺的.
被人剁了缺.
但是他願意嗎?他願意.
為什麼?因為他覺得.
他要利用餘下的.
下半場.
人生下半場.
上帝要他做的工作.
所以他提早退休.
現在我們會問.
這些人.
為什麼會這樣做呢?.

$^{241}$因為他們都是運動員.
那種精神.
為了追求有賞賜.
有上帝給他的獎賞.
他甘願放下.
目前所有的福利.
他的好處.
都不在話下了.
願意走上.
事奉上帝的道路.
我們知道保羅說.
我們好像運動員.
是因為我們有獎賞.
第二個原因.
為什麼這樣描繪呢?.
原來運動員.
保羅怎麼說呢?.
他說既然我們要賽跑的話.
你們也要這樣跑.
意思就是我們也要去跑.
也要去操練.
大家都知道.
如果你要去賽跑.
參加馬拉松.
要不要操練?.
要操練的.
我曾經見過一個人.
拿了別人的號碼.
跑馬拉松.
原來他也沒有什麼操練.
拿了別人的號碼.
結果在青魚隧道.
跑了不久.
跌倒了.
整個大字型的.
然後有人來搶救他.
大家知道什麼原因嗎?.
是因為他沒有操練.
沒有操練就去跑.
拿了別人的號碼去跑.

$^{281}$結果在隧道.
整個躺下了.
各位頂尖的朋友.
如果喜歡運動員.
要不要操練?.
要的.
無論你喜歡什麼運動.
要不要操練?.
要的.
台上十分鐘.
台下十年功.
一定要操練.
不操練不行.
你記不記得.
講起運動.
我們香港人.
面有光彩.
運動員在亞運.
成績亮不亮麗?.
亮不亮麗?.
亮麗.
不論是一般的運動員.
傷殘的運動員.
他們的成績都拿得不錯.
在亞運會.
所以我們香港人.
走在街上挺胸昂首都可以.
都有些光彩.
為什麼拿這麼多獎牌回來?.
頂尖的朋友.
誰是拿第一金?.
今年的亞運.
誰為香港拿第一金?.
我沒有這麼快.
不是去到這個.
知道是誰嗎?.
今年為香港亞運拿第一金.
是兩個男人.
做什麼?.
划艇.

$^{321}$記不記得划艇?.
你看這兩個運動員.
真的瘦猛猛.
一點都不是健碩那些.
不是健碩那些.
隨時我和你都比他健碩.
你明白我的意思嗎?.
但這兩個運動員.
竟然在划艇那裡.
拿到金牌.
還是第一金.
真的很金.
為什麼?.
各位頂尖的朋友.
原來這兩個運動員.
有沒有操練?.
不單有操練.
還有秘操.
什麼叫秘操?.
秘密操練.
他們這兩個人.
用了一年時間.
去到西藏高原的地方.
操練自己.
一年之後.
下山去到平原.
這兩個人夠不夠氣?.
夠氣了.
你明白嗎?.
無論氧氣什麼氣都去.
都夠氣了.
我就開始明白.
為什麼旁邊的艇.
烏茲別克.
又大又健碩.
都不夠他撐得快.
為什麼?.
因為要秘操.
放棄了一年時間.
去操練.

$^{361}$目的是要什麼?.
拿到金牌.
為什麼保羅用運動員.
給我們基督徒?.
是因為每一個基督徒.
都需要操練.
我們為什麼要操練.
要裝備自己?.
是因為上帝.
不想你和我.
空窗上陣.
大家明白我的意思嗎?.
如果你去事奉上帝.
去傳福音.
你沒有去裝備.
沒有去接受訓練.
心口掛個勇字.
撞板多過什麼?.
多過吃飯.
我曾經見到一個弟兄.
參加什麼課程呢?.
名字都很長.
叫做顏關大使訓練課程.
有沒有聽過?.
這個課程上八課.
我問這個弟兄.
為什麼你會參加這個課程.
上八課這麼多?.
還要晚上來上課.
這個弟兄就這樣回答我.
他說因為我最近.
不知為什麼.
無論在教會也好.
在公司也好.
在我的好朋友圈子也好.
總是遇到一些人.
很不幸.
患了癌症.
我的心很難過.
但我又不知道怎樣幫他們.

$^{401}$我又不知道說什麼.
我又不知道應該怎樣做.
越忙.
一個吳國宇.
就傷害到對方.
心更不想.
所以我決定報這個課程.
我希望自己能夠.
裝備自己.
訓練好自己.
以致我幫我的朋友.
幫我的弟兄姊妹的時候.
我懂得怎樣去幫他們.
所以各位弟兄姊妹們.
基督徒是要有裝備.
我們現在明白為什麼教會.
有時候辦這麼多專題講座.
有時候開這麼多主日學課程.
目的是什麼?.
目的是要讓我們人人都有機會.
接受訓練.
接受裝備.
所以如果你發覺有些課程適合.
你就快點去報名.
不要浪費這些機會.
以致你更加被上帝使用的時候.
你發揮得淋漓盡致.
這是第二個原因.
我們看第三個原因.
為什麼保羅用運動員比方基督徒呢?.
這個人真是人見人愛.
真是車見車載.
你喜不喜歡他?.
喜歡他.
他很有笑容.
很有陽光.
這些都是一個很正面的形象.
不只這樣.
原來每一次他贏了錦標.
記者採訪他的時候.

$^{441}$他都說了一句話.
我下一次再去游的時候.
可能再游得好一點.
就是說他有沒有進步的空間?.
有啊.
其實他已經拿到金牌.
為什麼他還這樣說呢?.
因為每一個運動員都知道.
他自己的紀錄不能夠永遠的保留下去.
你相不相信?.
頂級的運動員.
無論他的紀錄五年沒有人打破.
十年沒有人打破.
但他保證不到20年之後被人打破紀錄.
有沒有這樣的情況?.
有啊.
所以你刷新紀錄是運動員的目標.
為什麼呢?.
因為他很想不斷的進步.
這就是我們今天做基督徒.
好像運動員的原因.
是因為你和我都有進步的空間.
我不知道你今天信了耶穌多久.
可能十年,三年,五年,一年.
不要緊各位頂級的姐妹們.
你記得信了耶穌之後.
原來我們靈命的操練.
靈命的成長.
是有進步的空間的.
我們可以來到什麼?.
令到我們屬靈的生命.
不斷的透過我們的裝備.
透過我們的訓練.
令到我們的靈命更加成長.
記不記得保羅在《迦拉泰書》.
告訴我們聖靈有九樣的果子.
都不背出來.
我想你應該懂得背.
其實我們有很多要去追求的.
一些很重要的靈命操練.

$^{481}$在彼得後書有講過.
有了信心又要加上德行.
有了德行又要加上知識.
有了知識又要加上忍耐.
有沒有印象?.
一路加上去.
原來我們的靈命成長.
是不斷的有進步的空間.
所以各位頂級的姐妹們.
如果你信了耶穌之後.
你會發覺自己活得這麼開心.
這麼精彩.
這麼有幹勁.
這麼有衝勁.
是因為你懂得進步.
因為你有進步.
如果你原地踏步.
你就不進就側退.
所以我們發覺到.
運動員,比方教徒是很貼切的.
所以有些人信了耶穌很久.
但他的靈命可能比他更低更差.
為什麼?.
因為他沒有追求成長.
但當你有追求成長的時候.
你會發覺你不斷地體會到主的心意.
經歷到主的幫助.
既然寶萊提到我們做教徒.
是運動員有這三個原因的時候.
在這段經文裡面.
寶萊也說到做運動員怎樣可以成功.
怎樣可以達標.
接著就跟大家分享三個秘訣.
做運動員可以成功.
可以達標.
有哪三個秘訣呢?.
第一個秘訣是什麼呢?.
就是要有清晰的目標.
所以寶萊怎麼說呢?.
寶萊說我奔跑不像什麼?.

$^{521}$不像無目的.
或者Kitty再拿一張.
原來我們能夠成功能夠達標.
是要有方向有目標的.
因為寶萊告訴我們.
他就是有目標方向.
他打拳不是打空氣.
他奔跑不是萬無目標.
是有目標有方向.
所以今天做教徒.
你會發覺到我們真的需要有目標有方向.
我曾經見過一個人去參加馬拉松比賽.
為什麼他跑得很索氣.
好像氣喘如牛一樣.
再看清楚一點.
臉青口純白.
為什麼跑得這麼辛苦呢?.
原來這個運動員這個跑手是跑錯線.
跑錯了線.
因為他原本是跑半馬的.
他一個不經意就跑到全嗎哪裡.
結果你會發覺他真的跑到氣喘.
是因為迷失了方向.
別人轉彎的時候他不轉.
他繼續向前跑.
變成跑到全嗎哪邊.
說起頂尖的門.
如果我們沒有目標沒有方向的時候.
你會發覺到我們過的日子是一個很混亂的日子.
很明顯的是什麼呢?.
很明顯的是當你相信了耶穌.
你今年也很想有些事情去做.
那你要祈禱了.
求上帝指引你今年要做些什麼呢?.
有些人覺得不如參加兒童事工.
因為我也很喜歡小朋友.
所以我就參加兒童事工.
有些人覺得他的目標方向是什麼呢?.
我很想做老人家的工作.
長者事工.

$^{561}$為什麼呢?.
因為我覺得老人家都很淒涼.
因為老年沒有什麼人關心.
很孤單.
所以我想做長者事工.
其實是需要有什麼呢?.
目標和方向.
現在已經到了2023年.
還有一個多月就要到2024年了.
我不知道你今年有沒有為自己訂立了一個目標.
是怎樣去事奉呢?.
如果你發覺上帝已經給你有清晰的目標的話.
你就專心集中去做.
如果你沒有呢.
沒有就要回去好好祈禱.
悔改.
上帝我今年真的淪盡了.
沒有目標沒有方向.
我希望你寬恕我.
這個多月好好想一想.
靜下,祈禱.
然後我希望2024年.
上帝你給我一個目標.
一個方向.
讓我知道新一年.
我應該怎樣集中我的精神力量.
恩賜時間.
擺在什麼呢?.
擺在要去做的事奉的崗位.
又或者上帝你想我新一年.
向多少個人傳福音.
能夠帶領他們信主呢?.
可能上帝給你的數字目標都未定.
可能是五個,可能四個,可能十個都未定.
如果你有領受的時候.
你就集中你的精神注意力.
專注去做.
所以我們發覺運動員之所以成功.
是因為有方向有目標.
今天我們都需要什麼?.

$^{601}$有目標有方向.
第二個秘訣是要有什麼呢?.
原來保羅說.
做一個運動員.
凡事都有節制.
剛才讀了這篇經文有沒有印象?.
為什麼要有節制?.
因為如果你沒有節制.
肯定拿不到牌.
你沒有節制.
肯定什麼?.
你會大失水準.
所以你看到這幅圖.
這個運動員.
舊底都很帥.
手瓜氣濺.
不過現在他發了達.
有錢就胖了.
這個體操王子.
那時候他拿金牌的時候.
沒有一件事難倒他.
無論是吊環也好.
木馬也好.
單槓也好.
他能不能犯都可以.
你明白嗎?.
你知不知道為什麼可以?.
因為他那時候有節制.
現在可能沒有節制了.
他那時候有節制.
為了要拿牌.
為了要拿金牌.
他就限制自己.
30天裡面.
一粒飯都不吃.
只喝什麼?.
喝粥水.
為什麼?.
因為他想保持自己的體重.
不要超磅.

$^{641}$不要過磅.
保持自己的身形水準.
重量.
來面對比賽的挑戰.
所以底子沒有問題.
原來節制的意思是什麼?.
節制的意思就是說.
你懂得取捨.
凡事都有節制.
今天很多人.
吃東西沒有節制.
放縱自己.
說得好聽叫放題.
說得不好聽叫放縱.
不斷地吃.
吃到自己沒有三高.
都變成三高.
這就是因為我們不懂得取捨.
作為一個運動員.
他能夠成功達標.
因為他懂得取捨.
有些事情他很想去做.
是他的興趣.
例如他很喜歡打球.
很喜歡玩這些運動.
但想起如果這段時間.
家人需要我陪他.
需要我陪家人參加報道會.
參加什麼聚會.
我願不願意放下我的興趣.
陪家人去呢?.
大家明白我的意思吧?.
這就是取捨.
凡事都有節制.
當你知道上帝要你做一件事的時候.
你可能真的要取捨.
有些你很心愛的事情.
很有興趣要做的事情.
可能你要暫時放下.
完成了上帝要你做的功夫.

$^{681}$這就是凡事都有節制.
我們再看最後一個秘訣.
他如何達標如何成功呢?.
保羅說:我是功克己身,叫身服我.
這個功克在原文的意思是.
向著敵人致命的地方發出攻擊.
是這個意思.
例如大家看見打拳的那些人.
你知道打拳的拳手.
最致命的地方在哪裡嗎?.
哪個部份最致命呢?.
這個最致命.
所以為什麼拳手每次都要打對方的頭.
十拳打不中,打到一拳中就夠了.
一拳都打到對方暈頭轉向.
所以功克就是這個意思.
要向著敵人致命的地方發出攻擊.
保羅為什麼要這樣寫呢?.
他說這個攻擊不是打對方.
是打誰呢?.
是打自己.
為什麼呢?.
功克己身的意思是.
我們今天做基督徒.
如果能夠拿到獎牌.
能夠成功.
能夠被上帝所稱讚.
千萬不要忽略我們自己的問題.
因為我們信了耶穌之後.
我們這個人裡面.
時時都有一種老我.
一種舊我.
來控制著我們.
或者有我們的情慾.
有我們的慾念.
七情六慾來引誘我們.
所以我們很多時候都容易中了魔鬼撒旦的試探.
掉進了罪惡的網羅當中.
你說是不是很危險呢?.
所以保羅說.

$^{721}$你做一個運動員.
你要記得要什麼?.
要功克己身.
就好像箴言告訴我們.
自伏己心的.
是強如取勝.
如果一個成功的將軍.
他打仗很厲害.
攻打了一個城市.
一個城被他打下來.
好像很厲害.
但都不及一個人能夠控制自己的內心.
那種慾念.
所以各位頂尖的姐妹們.
我們就明白到保羅說要功克己身.
當我們功克己身的時候.
你就發覺到.
我們這個人為主作供.
為主事奉的時候.
就能夠達標.
就能夠成功.
否則呢?.
接著保羅說什麼?.
免得我全福音級別人.
自己反倒被淘汰.
這裡告訴我們.
自己反倒被棄絕.
大家有沒有留意到.
棄絕這個意思.
在原文是不合格的意思.
叫做棄絕.
或者說得白一點.
棄絕了.
這裡好像自己很熱心.
為主作供.
最後晚節不保.
有沒有這樣的情況?.
有的.
所以我們要很小心.
很謹慎.

$^{761}$各位頂尖的姐妹們.
今天和大家分享這個題目.
為主奔跑.
你要記得.
除了我們要掌握保羅告訴我們的秘訣之外.
時時都要謹記.
在我們佳音樓的弟兄姊妹要謹記.
上帝給我們有機會續約.
在這個地方繼續來傳福音.
繼續來關心我們的左鄰右里.
去關心四周的人.
我們要好好裝備自己.
為主奔跑.
因為擺在前面的路.
是有很多的挑戰.
你知道整個北區都會區.
會有很大的變化.
所以我們今天要努力.
要為主奔跑的時候.
我們不要每個人都沒有氣沒有力.
沒有神沒有力.
我們要做一個健康的運動員.
這個健康的運動員.
就是保羅剛才所說的.
有他的秘訣.
有他的要求.
使我們能夠達標.
被上帝大大的使用.
好不好?.
好,我們一起祈禱.
我們看到保羅真的樹立了一個很好的典範.
讓我們看到傳福音.
事奉上帝不是徒然的.
是有恩典,有獎賞的.
將來你有一天為我們所預備的冠冕.
是一個不扭壞,不衰殘.
是一個生命的冠冕.
我們Hallelujah多謝主.
但願我們嘉恩的弟兄姊妹.
更加努力被主訓練,裝備.

$^{801}$繼續好好地親近你,依靠你.
上上的熟讀上帝你的話語.
能夠拿起聖靈的寶劍.
來抵擋魔鬼的引誘試探.
讓我們能夠站立穩.
又為主作供.
給你大大的使用.
以致我們的奔跑不疲倦.
行走不乏力.
多謝耶穌.
誰聽我們的禱告.
奉耶穌名求.
阿們.
\newpage



\section{詩篇 150:1-6}
\label{sec:zCZXBv4_hT8}
\textbf{2023-11-18 晚堂《當來讚美神》詩篇 150\_1-6 講員:張美薇牧師}
\newline
\newline
連結: \href{https://youtube.com/watch?v=zCZXBv4_hT8}{\texttt{https://youtube.com/watch?v=zCZXBv4\_hT8}} ~~~~ 語音日期: 2023-11-21
\newline
\newline
\hyperref[sec:lMH9_3_EQe8]{\small{< < < PREV SERMON < < <}}
~
\hyperref[sec:index_chronic]{\small{[返順時目]}}
~
\hyperref[sec:index_scriptual]{\small{[返順卷目]}}
~
\hyperref[sec:uKxjJHU8_7c]{\small{> > > NEXT SERMON > > >}}
\newline
\newline
詩篇 150:1-6
\newline
\begin{longtable}{cl}
\hline
\hline
章節 & 經文 (和合本修訂版)\\
\hline
150:1 & \begin{tabularx}{0.7\textwidth}{X} 哈利路亞！你們要在神的聖所讚美他！在他顯能力的穹蒼讚美他！ \end{tabularx} \\ \\ \relax
150:2 & \begin{tabularx}{0.7\textwidth}{X} 要因他大能的作為讚美他，因他極其偉大讚美他！ \end{tabularx} \\ \\ \relax
150:3 & \begin{tabularx}{0.7\textwidth}{X} 要用角聲讚美他，鼓瑟彈琴讚美他！ \end{tabularx} \\ \\ \relax
150:4 & \begin{tabularx}{0.7\textwidth}{X} 擊鼓跳舞讚美他！用絲弦的樂器和簫的聲音讚美他！ \end{tabularx} \\ \\ \relax
150:5 & \begin{tabularx}{0.7\textwidth}{X} 用大響的鈸讚美他！用高聲的鈸讚美他！ \end{tabularx} \\ \\ \relax
150:6 & \begin{tabularx}{0.7\textwidth}{X} 凡有生命的都要讚美耶和華！哈利路亞！ \end{tabularx} \\ \\
[1ex]
\hline
\hline
\end{longtable}
$^{1}$弟兄姊妹平安.
很高興今天有機會再一次回來.
跟大家分享有關神的話.
讓我們開始的時候一起低頭祈禱.
天母上帝願你使用孩子.
也幫助弟兄姊妹.
讓他們有一個能聽到的耳朵.
能將你的話語藏在他們的心裡.
更加有明白你的話語的頭腦.
也有手和腳去回應.
去行出你的道.
聽我們的祈禱.
奉耶穌基督的名求.
阿們.
不知道大家有沒有去看過這齣戲.
叫做《反轉老朋友》.
主角叫做Riley.
不知道有沒有看過.
這齣戲是比較舊的.
我去找就發現.
原來他腦裡的五個朋友.
是很精彩的.
他有什麼呢.
他有一個叫做阿曹.
他經常發怒.
有一個叫做阿曾.
做的事很多時候都看不起.
或者做各樣的事.
另外當然有阿樂.
很開心的.
也有阿驚.
很容易的.
時時都會發怕的.
當然還有阿岑.
在五個情緒裡.
其實最喜歡的人.
就是阿樂.
很多人都很喜歡阿樂.
我們常常都很喜歡.
要常常喜樂.

$^{41}$但很多時候發覺.
原來在那齣戲裡.
主角除了阿樂之外.
還有阿愁.
因為如果我們沒有悲傷.
又怎麼能夠配合喜樂呢.
每一個情緒都有它的用途.
樂觀的人不一定每秒都快樂.
有時候勉強將所有回憶美化成快樂.
只是自欺欺人.
很多時候阿joy阿樂都會傷心.
阿愁都會笑.
找到平衡點.
才能遇到最好的.
所以我們發覺.
在戲結尾的時候.
阿樂和阿愁竟然會一起.
去解決他們很多問題.
其實每個不同的民族.
都有不同的民族性.
民族性都有可能.
會隨著時間改變.
我不知道你們現在做父母.
跟我小時候.
父母對我們有沒有什麼不一樣呢.
我記得小時候.
很少聽到父母稱讚.
我不知道你們是不是.
很多時候父母都是責備的多.
讚的少.
有時候做得很好都覺得是英分.
考試考到98分.
考到第二名.
父母還會說.
為什麼不拿第一.
為什麼不拿100分.
我們年紀大的就會這樣.
其實現在來說.
不是這樣的.
我們發覺在美國也是.

$^{81}$小朋友做到六七十分.
父母已經說.
Good, good, well done.
所以現在的父母都學了.
覺得是要這樣.
很多時候我們發覺.
在我們的文化底下.
很多時候我們對外人.
稱子女做什麼.
他很蠢.
他是我的乞丐.
或者乞丐女.
小朋友的自尊得不到見賴.
其實人人都喜歡別人讚.
什麼時候小朋友會很開心呢.
讓阿樂在他裡面充滿呢.
是怎樣呢.
當有人讚他的時候.
其實不單止小朋友喜歡別人讚.
其實政治人物.
都喜歡被人讚.
我想李家超.
特首都很喜歡被人讚.
讚美別人成為了人與人之間.
關係持續良好的必需品.
而合理又恰當的讚美.
原來是一件很美很美的事.
但我們看報紙.
大多數都是彈的.
很多東西都是彈來彈去.
同樣的我們對神.
我們又懂不懂去讚神呢.
今天我們會用詩篇第150篇.
用一個綠河的查經法.
你有沒有聽過.
綠河查經法就是問.
為什麼,怎樣,什麼時候等等.
去看一看我們怎樣去讚美神.
當我們去看主席以下.
幫我們讀出這段經文的時候.

$^{121}$我們會發覺.
裡面紅色的是什麼.
紅色字全部都是.
Hallelujah 讚美神.
讚美神.
讚美耶和華.
Hallelujah 或者是讚美他.
總共有多少次.
總共出現了13次.
如果大家心水清很快地數.
裡面有45個字是紅色的.
而全篇詩篇.
知不知道有多少個字.
是120個字.
所以全篇詩篇有三分一.
其實超過三分一是說.
我們要讚美神.
我們的經歷其實會影響.
我們怎樣去讚美神.
當我們去問.
我們要讚美誰.
很自然.
當然是讚美神.
我們來到與神的相交上高.
我們很多時候都不太懂得讚美神.
所以其實讀詩篇是好的.
因為詩篇很多時候都會是讚美神.
幫助我們學習怎樣的讚美神.
當我們會發覺.
我們來到神面前.
我們很少去讚美.
我們多數去求神.
我們很多時候都會祈求.
去求神.
而不記得或不懂得怎樣去讚美神.
我想大家都聽過.
很多人都說過這個故事.
有兩個天使.
一人揹著一個袋.
一個天使去收集人們對神的讚美.

$^{161}$另一個天使去收集人們對神的祈求.
結果誰是滿得很厲害.
當然是收集祈求的那個.
但當我們去看剛才讀出的經文.
他說要讚美神.
其實是命令.
你們要.
你們要讚美神.
他第一句就說.
Hallelujah.
我們要讚美神.
去到最後那一句.
都仍然會說.
我們要讚美神.
Hallelujah.
還有記住.
這首詩篇是150首詩篇最後的一首.
所以去到最後那一首.
最後那一句.
都仍然說要讚美神.
所以我們知道.
讚美神是私人命令我們做的.
是重要的.
是重複的.
不是可有可無.
是一定要做的.
在什麼情境我們都要學讚美神.
在世間.
人們對受讚美的神.
不願意讚美.
不懂得讚美.
但其實到我們去到天家的時候.
大家知不知道在天上我們是做什麼的.
有人說我們可能是彈樹琴.
有人說我們可能都是讚美神.
所以我們在地上高.
我們就要學習.
所以去看詩篇.
就會幫助我們怎樣去讚美神.
我們今天就去看第一個問題.

$^{201}$為什麼要讚美神.
在詩篇第150篇第二節.
那裡這樣說.
他說要因他大能的作為.
讚美他.
因他極其偉大.
讚美他.
神的大能.
神的作為.
所以我們看到.
為什麼要讚美神.
第一個是因為神大能的作為.
其實任何人物要得到讚美.
絕對不應該單憑他的身份.
或者他的地位而得到.
必須因為他具有美好的作為.
美好的表現.
讚美是贏回來的.
我們讚美神.
是因為神贏得我們的讚美.
並不是天然.
並不是天生的.
更加不是強迫出來的.
當我們說到神的大能.
神有什麼大能.
我們要讚美他.
我們很自然就想到什麼.
就想到神的創造.
天和地都是神做的.
是不是.
地和其中的一切.
都是他做的.
天地顯出神的大能.
神也做出了人類.
做出了你.
做出了我.
我們為神的創造.
我們要讚美神.
我們更加要為神在歷史當中的大作為.
讚美神.

$^{241}$當我們去看舊約的歷史.
新約的歷史.
我們會看到神揀選亞伯拉罕.
神將以色列人領出埃及.
神將以色列人帶進迦南地.
為他們征戰.
神也當以色列人被擄.
在當中顯出他保護的大能.
當以色列亡國之後.
以色列人到今.
竟然可以歸回耶路撒冷.
其實我們都要讚美神.
雖然亡國二千多年.
仍然在二戰之後復國.
至今仍然存在.
所以神的大能是很大很大.
是奇大無比.
當然我們今天也會面對.
以巴戰爭.
究竟誰是誰非呢.
但我們除了因為神的大能.
神的大作為.
要讚美祂之外.
我們要因為神極美的大德.
去讚美祂.
記住耶穌基督在十字架上.
成就了普世的救贖工作.
不單單只救贖以色列人.
也延伸到外邦.
延伸到萬國萬民.
更延伸到哪裡呢.
延伸到我們當中.
我們今天能夠坐在聖所裡.
我們可以一起讚美神.
其實是因為神極美的大德.
祂對祂所做的人的慈愛.
祂降雨給好人.
神也降雨給不好的人.
祂叫白天照日.
也同樣叫白天照惡日.

$^{281}$神的大德是極美.
神的醫治.
神對有需要的人的醫治.
對我們身心靈的醫治.
神對我們的公認.
神對我們的看顧.
神不單單看顧那些竭力盡心.
去事奉祂或去生活的人.
神也看顧一半被人忽略的人.
孤兒寡婦.
神也同樣看顧.
其實當我們大家坐在這裡.
講起感受的歷史的時候.
大家也會知道.
神有很多大能.
很多大德.
在我們當中彰顯出來.
好像我們起初的時候.
我們發覺當我們四十多年前.
想在這裡植堂的時候.
大業主不租給我們.
他怕租給基督教會之後.
回教,佛教,道教又要租的時候.
那怎麼辦呢?.
於是就拒絕了弟兄姊妹.
你們知不知道?.
為什麼我們終於可以在這裡植堂呢?.
可以建立感受堂呢?.
因為當時的弟兄姊妹很聰明.
租不到地方.
他們怎麼辦呢?.
於是他們就說.
知道這裡是新的社區.
很有需要.
一定很需要一間幼稚園.
所以就在這裡建立了幼稚園.
我相信那一年幼稚園就在九月開課.
感受堂就在十月在這裡建立.
教會就因著神的大德.
我們可以在這裡建立.

$^{321}$神也不單止讓感受在這裡建立.
神也叫感受可以植出孫德堂.
去彰顯神的大能.
孫德堂在五周年,十周年.
十五周年,二十周年.
都出了見證集.
在見證集裡.
很多弟兄姊妹都不斷地寫文章.
寫厚厚的一本見證集.
來歌頌神,讚美神.
就算是年輕的弟兄姊妹.
少年部的弟兄姊妹.
都有份去讚美神.
我記得我初信的時候.
當時我參加一個團契.
團契裡的團長就跟我們說.
他說因為我們那時候.
每個星期都會一起開會.
雖然人少少.
但是我們一個星期就會開團契的會.
另一個星期就會開跟進的會.
每個星期都要開會.
團長就跟我們說.
我們每個星期回來都要有分享.
要分享什麼呢?.
要分享我們的靈修心得.
同時也要分享.
我們看到神在我們附近做了什麼.
起初的時候是很困難.
好像看來看去都看不到.
神有什麼做.
但是整下整下.
當眼睛開了,習慣了之後.
就會發覺很多時候.
很早就看到神在工作.
於是就說為什麼都不到星期六.
為什麼都不到星期六.
很想很想去讚美神.
是的.
神祂的大能的作為.

$^{361}$祂極美的大德.
所以我們是需要去讚美祂.
另外我們再要問.
我們應該在哪裡去讚美神呢?.
詩人在詩篇第150篇第一節裡說.
Hallelujah!你們要在神的聖所讚美他.
在他顯能力的窮倉讚美他.
聖所是人做的地方.
是人用來敬拜神的地方.
我們知道其實不同的宗教.
不同的人物.
都做不同的聖所.
如果我們去看的話.
在他的聖所.
究竟有什麼的聖所呢?.
我們可以說這是一個.
會教寺.
這是一個印度廟.
這是一個埃及的聖所.
這是一些比較.
部落的圖騰.
可能是加拿大那邊的.
他們就是用這些來敬拜他們的神.
其實中國我們也有天壇.
也是用來敬拜我們的神.
雖然他們未必是我們所了解.
是敬拜一個真的神.
但是他們每一個的民族.
每一個的宗教.
他們都會去建立聖所.
來去讚美神.
去敬拜神.
如果我們看.
其實聖所在詩篇裡面來說.
其實是指聖殿.
後來以色列人被擄之後.
聖殿被毀.
神的子民也在各處的地方.
設立會堂.
以致我們今天在各處的地方.

$^{401}$我們都會有禮拜堂.
或者我們所謂的教堂.
教堂可以是很金碧輝煌.
可以是很漂亮.
但也可以是很簡單.
很簡潔.
因為當人有什麼.
人就會在那裡建立一些地方.
來敬拜神.
如果我們是很有錢的.
很富有的.
他們會建立一些很大的教堂.
如果他們比較簡單.
或者比較欠缺的.
他們會建立一些簡潔的教堂.
甚或一些的貿寮.
我想弟兄姊妹都會.
在當中有弟兄姊妹記得.
你們去泰國去短孫的時候.
其實那時候.
就在這些四邊都沒有牆的地方.
但卻可以在裡面.
一起去敬拜神.
教會是由信徒組成的.
這個是重要的.
敬拜的地方並不是最重要.
但教堂不是最重要.
但教堂仍然是用來敬拜神的建築物.
對敬拜神的地方.
我們是要很重視.
很記得有一次.
藤近輝牧師跟我們說.
他起初在北宣聚會的時候.
他那時候是北宣的傳道人.
是北宣的牧師.
他說有一個星期六的早上.
他在教會裡面工作的時候.
他就想去教堂裡面找一些東西.
然後他就見了幾個上海來的老婆婆.
他們其實是很有錢的老婆婆.

$^{441}$但他們卻願意.
星期六一早回到教堂裡面.
去打掃乾淨.
讓弟兄姊妹在星期日的時候.
可以有一個很乾淨很整齊的地方.
去敬拜神.
所以藤牧師說.
見到他們那麼專心.
那麼事奉神的時候.
他就很安靜地退出去.
恐怕打擾到他們那麼神聖的服務.
很簡單的清潔.
都是很神聖的服務.
為了基督的緣故.
很多時候在聚會當中.
以心靈的歌頌.
以口讚美.
以及用手勢表達對神的重量.
這一切是詩篇希望教導我們的.
這句話是潘博華所說的.
我們知道到末後的時候.
我們會有黃金街.
會有碧玉城.
會有玻璃海.
在三層天上高的聖所裡面讚美神.
是面對面的讚美神.
所以詩人說.
我們要在聖所裡面讚美神.
詩人更加說.
我們不單止在神的聖所裡面讚美神.
我們更加要在祂顯能力的窮倉裡面讚美神.
神造的世界.
彰顯神榮耀的地方.
我們都要讚美神.
我們在整個宇宙的任何地方.
都同樣要向神獻上讚美.
在天地當中讚美神.
我們有時唱些詩歌都有.
我們不單止在禮拜堂裡面讚美神.
我們更要在天地當中讚美神.

$^{481}$因為天地彰顯神的榮美.
太陽,日月,星宿都讚美神.
都將神顯明出來.
所以保羅說.
自從做天地以來.
神的永能和神性是明明可見的.
雖是眼不能見.
但藉著所做之物.
就可以了解.
叫人無可推諉.
在大如天空的星體.
小如電子,原子等等的粒子.
規律其實都是一樣的.
叫人驚訝神的奇妙.
在人體的細密.
亦不是科學家所完全能夠明白的.
連人的眼眉毛.
腳的小指頭.
都有一定的用處.
窮瘡不斷地在讚美神.
怪不得勝經說.
如果人不開口讚美神.
連石頭都要發聲讚美神.
其實整個宇宙.
是一個很大很大的聖所.
所以人要在神顯能力的窮瘡.
讚美他.
我們看看.
這是一個很美的景.
這張照片其實沒有做過任何手腳.
是同事的太太拍的.
是很美的.
窮瘡是很美麗.
但亦有時是烏雲密佈.
是行雷閃電.
更加會有龍捲風之類的風在影響著.
但我們卻知道雨後卻可以有天空.
所以無論天色是蔚藍的.
或是暴風雨當中.
我們同樣都可以讚美神.

$^{521}$其實人處於順境或逆境當中.
我們都要學習開口讚美神.
順境讚美神是容易的.
但逆境我們都要讚美神.
所以有人說.
在逆境當中說一句讚美神.
比在得意的時候說一千次讚美神.
是更加有價值.
甚至我們看到保羅的生平.
他在監獄當中都一樣可以讚美神.
當他讚美神的時候.
我們看到在菲律賓的監獄裡.
地大震動 監門大開.
弟兄姊妹.
你們現在處於順境當中嗎.
你們要讚美神.
還是處於逆境當中呢.
或是人生低潮當中呢.
你仍然要開口讚美神.
神是不改變的.
無論何境況.
祂仍然愛你.
祂不會丟下你.
祂會幫助你.
有人會問我們怎樣去讚美神呢.
在第三至第五節詩人說.
要用覺聲讚美他.
鼓瑟彈琴讚美他.
擊鼓跳舞讚美他.
用詩言的樂器和燒的聲音讚美他.
用大響的筆讚美他.
用高聲的筆讚美他.
詩人用了三節這麼長的篇幅來說.
我們可以怎樣讚美神.
這裡有各類的樂器.
我知道今天你們.
聖樂的弟兄姊妹去了國內交流.
所以前面的樂器沒有用.
所以今天我們只有鋼琴和風琴.
但我們知道詩篇裡有說.

$^{561}$有各類的樂器.
包括了弦樂.
有琴,樹琴,室樂.
有詩言的樂器.
有管樂.
有覺,有燒.
有擊打的樂器.
有筆,鼓等等.
各種樂器.
各有各特殊美妙的聲音.
當同時演奏合來一起的時候.
是一首多姿多彩的奏曲.
也包括不同種族的人.
用不同種族很特別的樂器.
都同樣可以去讚美神.
同時有各種姿勢.
各種不同樂器的演奏方法.
有些用口吹.
有些用手撥.
有些是擊打的.
有些是用手敲.
有些是用手彈.
又或者有時是全身來跳舞.
又或者用口來歌唱.
我們的口演一個很好的樂器.
一舉一動各種的方式.
我們都可以一起去讚美神.
敬拜神.
還有用大小不同的聲浪.
有高調的,有低調的.
有聲浪大的,有聲浪小的.
去活出讚美的人生.
我們知道其實.
當中也有詩班.
也有福音樂曲.
不同的民族.
也有不同文化的音樂.
同樣去敬拜.
創天造地的獨一真神.
用我們各樣的東西.

$^{601}$各種的姿勢.
各種不同的聲音.
一起和起來.
一起去讚美神.
接著要問誰應當讚美神呢?.
凡有生命的都要讚美耶和華.
Hallelujah.
第一個要讚美神的是誰?.
是你們.
是信徒.
神的子民屬神的人.
信徒一定要先起來帶領敬拜,讚美.
所以我們當中有領唱.
他帶領我們一起去唱.
Hallelujah.
這是信徒所熟悉的.
也願意敬拜和讚美.
不單是和其他信徒.
也要和其他人一起去讚美.
所以很特別.
今天聖樂部的弟兄姊妹.
會到國內一起去讚美神.
你們要在神的聖所讚美神.
屬神的人每週的崇拜.
都在這裡讚美神.
因為讚美神是一件重要的事.
其實不單是公開的讚美.
我們個人的靈修.
我們個人的讚美.
也是一個很好的機會.
每天都要讚美神.
每天都要親近神.
神是很重視我們的敬拜.
很重視我們的讚美.
和謝恩.
其實不是因為祂自大.
乃是因為我們的益處為念.
提升為以基督為中心.
使我們的心智.
以主耶穌基督為焦點.

$^{641}$而且越來越似祂.
所以信徒們.
我們一定要學好,學會.
每天,每週,每個時間.
我們在教會裡去讚美神.
我們更加要在天地之間讚美神.
但其實不單是講信徒.
私人是講什麼呢.
凡有氣息.
什麼叫凡有氣息呢.
凡有氣息即是神所創造的.
神的子民是屬神的人.
神的創造每一個人都是屬神.
神都很希望每一個人都能夠讚美他.
所以不單是已經認識神的人.
神是希望每一個有氣息的人.
都要來讚美神.
凡是活人都要來讚美神.
其實剛才大家看過.
我給大家看不同民族.
他們都需要神.
他們都會讚美神.
就好像我們知道.
印尼未開發的獵頭族.
美洲的印第安人.
非洲的土人.
會教徒,佛教徒,印度教徒.
他們都了解到人不能夠控制的事.
他們會敬畏天,敬畏宇宙.
敬畏大自然.
敬畏他們的神.
他們都會做聖所來敬拜他們的神.
誰曉得那創造天地的神.
誰能曉得耶穌基督降世為人.
背負世人的罪而走上十甲的救恩.
誰又曉得勝靈奉派來到世界上高.
叫人為罪,為義,為審判,自己,責備自己呢?.
羅馬書說.
哪有奉差出去宣講神的大能和大德的人.
神將全福音的使命交給我們.

$^{681}$神要我們來完成祂的使命.
要我們叫凡有氣息的都來讚美神.
我們要在這裡讚美神.
我們更加要出去向世上我們所認識的人.
我們不認識的人去宣講神的大能,大德.
羅馬書第十章十四到十五節說.
然而人還沒有信他,怎能求告他呢?.
沒有聽見他,怎能信他呢?.
沒有人傳揚,怎能聽見呢?.
如果沒有奉差顯,怎能傳揚呢?.
如經上所記.
那些傳美事報喜訊的人.
他們的腳中何等皆美.
弟兄姊妹.
你們就是那些腳中皆美的人.
你們要去傳福音.
大德生曾經說過.
甚麼時候神的兒女才會覺察到.
神要我們往天下傳福音給萬民聽.
這個使命並不是被我們塵封在國際上.
不是讓我們束之高閣.
神給我們這個使命.
我們每一個人都要去.
我剛才進來的時候.
有同工跟我說.
原來十一月是猜全月.
大家要出去紀念神給我們的使命.
所以你們不單止在這裡要讚美神.
你們更加要出去.
讓其他人都明白,都了解.
然後去讚美神.
這是我們的使命.
我們要讚美神.
我們要奉差出去宣揚神.
將凡有氣息的都來讚美神.
我最近講道經常用這個結束.
大家如果有聽過都記得.
於你何干.
講了大概八個字.
你會怎樣去讚美神.

$^{721}$你又會怎樣去讓普天下的人.
凡有氣息的人都讚美神.
讓我們低頭祈禱.
天父上帝.
願每一個弟兄姊妹都珍惜.
亦都去努力.
亦都喜悅去讚美你.
更加幫助每一個弟兄姊妹.
都能夠去完成你給我們的大使命.
叫凡有氣息的都來讚美你.
我們這次禱告.
是奉耶穌基督的名禱.
阿們.
\newpage



\section{約書亞記 14:6-15}
\label{sec:uKxjJHU8_7c}
\textbf{2023-11-19 加恩樓《生命的投資》約書亞記 14\_6-15 講員:蔣文忠傳道}
\newline
\newline
連結: \href{https://youtube.com/watch?v=uKxjJHU8-7c}{\texttt{https://youtube.com/watch?v=uKxjJHU8-7c}} ~~~~ 語音日期: 2023-11-21
\newline
\newline
\hyperref[sec:zCZXBv4_hT8]{\small{< < < PREV SERMON < < <}}
~
\hyperref[sec:index_chronic]{\small{[返順時目]}}
~
\hyperref[sec:index_scriptual]{\small{[返順卷目]}}
~
\hyperref[sec:z_EDsyn25oM]{\small{> > > NEXT SERMON > > >}}
\newline
\newline
約書亞記 14:6-15
\newline
\begin{longtable}{cl}
\hline
\hline
章節 & 經文 (和合本修訂版)\\
\hline
14:6 & \begin{tabularx}{0.7\textwidth}{X} 猶大人來到吉甲，約書亞那裡，基尼洗族耶孚尼的兒子迦勒對約書亞說：「耶和華在加低斯‧巴尼亞指著我和你對神人摩西所說的話，你都知道。 \end{tabularx} \\ \\ \relax
14:7 & \begin{tabularx}{0.7\textwidth}{X} 耶和華的僕人摩西從加低斯‧巴尼亞差派我窺探這地的時候，我剛四十歲。我把心裡的話向他報告。 \end{tabularx} \\ \\ \relax
14:8 & \begin{tabularx}{0.7\textwidth}{X} 雖然同我上去的眾弟兄使百姓膽戰心驚，我仍然專心跟從耶和華－我的神。 \end{tabularx} \\ \\ \relax
14:9 & \begin{tabularx}{0.7\textwidth}{X} 那日，摩西起誓說：『你腳所踏之地必要歸你和你的子孫永遠為業，因為你專心跟從耶和華－我的神。』 \end{tabularx} \\ \\ \relax
14:10 & \begin{tabularx}{0.7\textwidth}{X} 現在，看哪，耶和華照他所說的使我活了這四十五年。當以色列人在曠野飄流的時候，耶和華曾對摩西說了這話。現在，看哪，我已經八十五歲了。 \end{tabularx} \\ \\ \relax
14:11 & \begin{tabularx}{0.7\textwidth}{X} 現今我還很健壯，像摩西差派我去的那天一樣；無論是戰爭，是出入，我現在的力量和那時的力量一樣。 \end{tabularx} \\ \\ \relax
14:12 & \begin{tabularx}{0.7\textwidth}{X} 請你將耶和華那日所說的這山區給我。那日你也曾聽說，這裡有亞衲族人，以及寬大堅固的城，或許耶和華會照他所說的與我同在，我就把他們趕出去。」 \end{tabularx} \\ \\ \relax
14:13 & \begin{tabularx}{0.7\textwidth}{X} 於是約書亞為耶孚尼的兒子迦勒祝福，把希伯崙給他為業。 \end{tabularx} \\ \\ \relax
14:14 & \begin{tabularx}{0.7\textwidth}{X} 所以希伯崙成了基尼洗族耶孚尼的兒子迦勒的產業，直到今日，因為他專心跟從耶和華－以色列的神。 \end{tabularx} \\ \\ \relax
14:15 & \begin{tabularx}{0.7\textwidth}{X} 希伯崙從前名叫基列‧亞巴；亞巴是亞衲族最尊貴的人。於是國中太平，沒有戰爭了。 \end{tabularx} \\ \\
[1ex]
\hline
\hline
\end{longtable}
$^{1}$.
好像一種否白自己人生.
是因為他說的話 所以使我今天用這個題目跟大家分享.
我們的生命投資投放在什麼上面呢?.
這是我希望《帕望直柱加勒》的生命故事.
刺激到我們 讓我們更多思考我們自己.
又如何去規劃我們自己的人生呢?.
無論大小老幼都沒有問題的 其實都是在一個人生階段當中.
我們都要再往前看 去看人生的路如何再走下去.
不知道你會如何投資你的生命呢?.
剛才我這裡是再突顯加勒剛才說的經文.
他將自己人生分了幾個階段 有時我們都會的.
我們經常說人生下半場 過了某個階段我就下半場了.
有時都未必是 但自己當自己是.
然後就預備在下一屆如何再走下去.
或者我們退休之後如何過 還有一個金陵的時間.
又如何活下去呢?.
有時我們都會為自己這樣去界限 去分段 去看自己的人生.
然後就為自己有一些目標再去建立.
然後就再讓自己向前再邁進.
加勒裡面也說了他40歲的時候.
就是摩西派他和約書亞和十個兄弟.
當時如果看民數記第13,14章.
我不會詳細讀那一段 但有時我都會引用那段經文.
裡面說到這12個人 別看低他們.
其實是12個支派的領袖.
所以就各有一個領袖做代表 進入加勒去窺探那個地方是怎樣.
當時加勒就說我是40歲.
一個很有衝勁的領袖.
想帶領伊斯蘭人進入一個新的地方 新的方向的一個年紀.
去到後面 剛才也讀了.
去到經文剛才第10節的時候.
接著 大家都很熟了.
因為正正因為這個經歷.
其實上帝作了一個很重要的決定.
就是令到當時的第一代在曠野猶太人進入不了迦林地.
大家都很熟悉的曠野漂流40年.
有一段這樣的時間.
加勒就說了在這45年裡面.
就是當伊斯蘭人在曠野漂流這個時間.

$^{41}$原來加勒在當中經過這個曠野日子.
然後耶路撒冷曾經跟我講過.
曾對摩西說過這番話.
就是說他由耶路撒冷對摩西說關於加勒的說話.
到他今天現在跟約書亞說這番話的時候.
相差45年.
到現在我85歲了.
也是他預備進入一個人生的另一個階段.
如果大家記得當時的故事.
是剛才講了一點.
因為我會再詳細再講多一點.
就是關於伊斯蘭人的命運.
也是關乎加勒自己.
他能否什麼時候得到上帝給他的地方.
頂智婆我不知道大家.
當加勒這樣去將自己人生.
40歲的時候.
接著45年的曠野.
然後到85歲再進入人生的階段.
對於你來說你又會怎樣規劃去看自己人生呢.
譬如你想起一些自己聽到的生命故事.
他會這樣分享.
譬如一些20多歲的年輕人.
20歲或者20多歲.
有一個年輕人跟我說.
我畢業了.
我很想看看自己有沒有宣教心智.
所以用了一年時間去了非洲.
他說沒得輸.
我人生年輕.
儘管試一下.
所以未想得很清楚.
就決定了將一年時間放到非洲.
也有些同學.
他毅然在讀一個主科的時候.
他中間會決定轉科.
他覺得自己的主修不是他最有興趣的科目.
或者先拿一年gap year.
再探索一下.
然後再決定.

$^{81}$年輕人可以這樣去決定人生.
或者未思考得很清楚.
但他覺得我有青春.
我有時間.
就這樣儘管試一下.
這樣做下去.
但到了40歲的時候.
大概40歲的朋友跟我說.
他覺得到了人生很重要的轉捩點.
究竟他繼續這份工作.
高薪厚職.
但好像沒什麼挑戰.
還是轉去一個新的環境.
再給自己新的挑戰.
再給自己新的突破.
他就在這裡思前想後.
在一個這樣的人生的階段.
他覺得自己還有拼搏的精神.
但多了東西要關心.
多了東西要去負擔.
所以不是這麼輕鬆.
好像20歲的時候.
說去就去這樣決定.
到了60歲又如何呢.
即是今日香港人都60歲.
60多歲.
或者提出退休.
都會想在工作裡面退下來.
我有朋友就是他退休了.
他跟我說.
年少時一直有個心願想開咖啡店.
會不會現在從一份做了30多年的工作.
可以退下來.
有自己時間.
就去做回年少未完成的夢想呢.
但到了60歲的時候.
原來發覺不容易去冒這個風險.
不容易做一些冒險的決定.
覺得自己的負擔.
自己的顧慮更加大.

$^{121}$所以思前想後之後.
都不容易或者不敢於行出一個決定.
就算我不知道你在哪個階段.
什麼狀態.
你如何去再預備自己.
再向前走.
或者你如何回顧.
你過去經過一段人生的階段也好.
我們今日從迦勒的這個身歷.
他自己為自己去回顧.
或者看前面的路的時候.
那種階段性的時候.
我們可以讓我們的生命更多.
去思想自己.
又如何再向前走下去.
在經文裡面.
其中一個都重複了幾次的一句話.
去專心跟從耶和華.
如果是譯得比較好一點.
其實那個字是一個.
就是完全很圓滿地跟從.
英文就是fully的意思.
是一種.
我譯了做全心.
全心.
跟專心有一點不同.
專心就好像你專注.
將你的專注力放在一些東西上.
全心就是在說.
我們沒有被其他東西佔據.
去全心全意跟從耶和華.
其實在《聖經》裡面.
都多說這些概念.
或者在全心的相反.
如果在新約裡面就說一種叫做三心兩意.
double-minded.
或者double-heart的意思.
就是說你有很多東西兼顧.
今天提到.
你想想你做一份工作的時候.

$^{161}$你主宰的工作又兼領一些職務.
那你不能夠全心做好你本身最主要的工作.
我們很多東西兼顧.
很多東西去尖前顧後.
這樣的狀態.
就是一種不能夠全心.
三心兩意的狀態.
當時在十二個探子去窺探這個地方的時候.
那十個的朋友.
那十個弟兄.
你見到第八節.
他都這樣說.
雖然和我上去的眾弟兄.
使百姓膽戰心驚.
但是我專心跟從耶和華.
他的譯法應該是.
但是我專心跟從耶和華.
那十個弟兄.
大家記得那個故事.
是不是.
文書的十三章.
是派了他們去.
其實要窺探那個地方.
是不是一個流拿與物資地.
其實他們已經做足功夫了.
拿了那些葡萄回去.
給摩西給以色列人.
真的正的.
如果進去住.
應該是一個農產品.
各樣畜牧等等.
可能是一個很美好的地方.
那個地方真的正.
但是他們也都說.
那裡的人很高大.
然後誠有堅固.
說是一個吞吃居民之地.
因為文書的十三章.
第三十二節.
所以既是確認這個地方很好.

$^{201}$但是也都很憂慮.
那些環境的因素.
這個就是三心兩意.
他不能夠專注在.
放全心在上帝.
叫他看一件事.
他給了另一些東西.
都令到他都分心.
將心思放在那裡.
他就在那裡搖擺不定.
甚至比那個環境因素.
更加蓋過了他自己.
原先要全心跟從耶和華的那條路.
但是迦勒他自己的路.
專心跟從耶和華.
去到曠野之間.
你不要忘記.
其實當以色列子民被懲罰.
在曠野飄流四十年的時候.
迦勒也在當中.
接著經文也都說他.
是專心跟從耶和華.
你當然不要忘記.
那些是一班搏弈之民.
很多怨言.
很多投訴.
還要想回到埃及地.
你記得當時在整個曠野的路程上.
上帝又擊打盤石.
接著很多這些.
懲罰他們的行為出現了.
但是迦勒在中間.
跟這班人一起的時候.
仍然能夠專心跟從耶和華.
到他最後經文說到.
第十二節的時候.
十三節的時候.
和十四節.
當時摩西按照這個說話.
約書亞給了他的時候.

$^{241}$他能夠得到希伯倫的地方.
因為他也是全心跟從耶和華.
所以迦勒的一生.
就讓我看到是一個.
全心跟從耶和華的一生.
無論他在四十歲的時候.
在曠野日子的時候.
到他接著要回應上帝的應許的時候.
他一生裡面都沒有受到的環境.
沒有受到自己的一些變化.
而影響得到.
突然之間會盼望我們.
都能夠學習迦勒在中的洋食.
我們是否可以將人生.
最精壯的年歲.
大家可以想一下.
你放了最長的時間.
放在哪裡呢.
在你現在活到今天的時候.
你回顧的時候.
會不會是.
譬如你的生意.
有一份工作.
或者放在讀書那裡.
或者是有時候.
要放在照顧家人的上面等等.
可能有不同的事情.
在我們人生當中要面對.
要去應付的時候.
我們將我們最多的時間.
最精力.
最好的光陰.
放在什麼地方上面呢.
而在那些事情上面.
是否也是全心地跟從上帝.
作出一些重要的決定和投放呢.
我們能不能夠去.
好像迦勒一樣.
不顧一切.
其他的因素影響.

$^{281}$去為上帝而活呢.
這是很盼望我們今天.
去思想我們的生命.
接著我會從.
迦勒的生命當中三個階段.
跟大家分享.
去看迦勒怎樣可以.
每一個階段當中.
都是全心地跟從.
我們的上帝.
第一個階段是剛才說了.
40歲的時候.
背景就是《民數記》的.
13章14章.
波西派了12個探子.
去窺探迦南地.
而10個探子回來.
向以色列人報惡信.
英文就是一個bad report.
就是.
我們有時候工作寫很多報告.
你怎樣寫報告都視乎你.
怎樣理解那些資料.
他將那些情況.
大家看到同一個環境.
同一個.
看到同一個現象.
但是回來寫下來的時候.
就是一個bad report.
令到以色列人當中心裡融化.
很恐懼.
有一番混亂在當中.
他看到.
是不是.
好像經文13章13次這樣說.
他看到那些人.
是巨人.
巨人中的亞納族人.
我們看自己眼中.
他們乍看一樣.

$^{321}$而他們眼中看我們也是一樣.
是不是.
都很仔細的.
他連.
不只是自己看自己.
連那些亞納族人怎樣看他.
都看到.
這個是他自己的照形象.
是看低了自己.
這個就是那個.
回報.
令到他們恐慌的原因.
這個環境的因素.
令到他們.
不敢投放自己生命在上帝的路上.
看看迦勒的反應.
迦勒撕裂衣服.
然後說必會領我們進這地.
在文述的14章後面就繼續詳細說.
迦勒當時這個行動很想.
停止到這班十個兄弟.
所對於以色列人的反應.
所以他走出來.
想告訴你一個上帝的心意如何.
就發生了剛才那件事.
上帝就決定了.
第一代的意思是不能夠進迦林地.
而只有約書和迦林地.
你問上帝如何評估.
如何評價迦勒的生命.
他說.
為我的僕人.
對我的僕人迦勒.
因他有另一個心智.
專心.
剛才說的就是全心跟從我.
我要領他進入他要去的地方.
他後裔都要進入那個地方為業.
他有另一個心智.
這個是.

$^{361}$another.
一個與別不同的心智.
甚至是一個獨排眾議的心智.
所以有時候.
我們不是說眾人所做的事情.
一定是美事.
當時就是一個例子.
迦勒就看到.
無論是十個兄弟再加上以色列人的想法傾向.
其實是遠離上帝.
上帝就說.
迦勒有一個與別不同的心智.
那個心思.
那個靈.
是跟從耶和華的.
他能夠看到事情的真相.
他能夠有不同的報告打出來.
他解讀當時的情況有不同的理解.
他看到上帝的心意.
他不受這個環境的因素所決定.
等於我們要看到別人看不到的東西.
不要誤會我看到一些奇怪的東西.
而是看到同一件事情.
但是你從上帝的角度來看.
那個意義.
那個能夠給你.
帶給你的決定是不同的.
以致我們能夠跳出一步.
去做一些一般人不會做的事情.
這個就是迦勒當時所做的事情.
大家想想.
其實我們今天坐在這裡很舒服.
敬拜吧.
大家都記得.
或者留意過.
其實錦繡堂也是由西教士進來.
開幼稚園.
開始有教會.
以後我們能夠發展到今天有監牢的聚會.
這個整個環境.

$^{401}$我們能夠有這樣的資源,地方.
其實是由一群來到香港的傳教士.
去傳福音給當地的人.
給他們香港的人.
這個整個宣教運動是從十九世紀開始的.
譬如大家聽過建橋七傑.
可能聽過戴德新的很感動人的故事.
有當時的學生的自由運動.
其實在十九世紀末至二十世紀中之內.
這幾十年.
是有很多很多傳教士.
由英國,美國等等.
來到不單止是內地,香港,日本.
東南亞金三角等等的地方.
是將福音傳給這些未得之民.
你想想其實他們這個行動.
是很不可思議的.
他們在當時的西方或者歐洲地方.
其實無論是教育,環境,生活,文化等等.
都是他們很享受的地方.
但是來到東南亞一帶.
其實是有很多東西適應的.
可能大家聽過很多宣教士故事.
其實來到一分鐘.
未來到小的地方.
在船中已經過身了.
來到的時候可能很短時間.
適應不了水土,飲食習慣等等的時候.
接著有很多的疾病.
那些藥疾,傷寒.
各種的風土病,濕熱等等.
是取去了很多很多宣教士的生命.
他們可以將壞的報告帶回去.
亦都可以將好的報告帶回去.
但是大家事實上看到.
其實有更多更多宣教士來到這個地方.
以致建立了很多的教會.
以致今天我們才能夠有這個地方.
可以去敬拜上帝.
不只是看環境去決定生命的投資.

$^{441}$這是很多的先人都留下了很多優美的腳蹤.
讓我們去看得到,回顧得到.
不受環境所限.
不受環境的因素去影響我們.
是在迦納或者我剛才所說的.
很多來到東南亞的宣教士.
給了我們一個美好的提醒.
接著是一個曠野的生活.
迦納有45年.
跟這班以色列人一起.
在曠野漂流的日子.
迦納沒有一條短路.
上帝你不要那麼辛苦.
直接由曠野調升,提升到.
降落到他要去的基伯倫的地方.
他跟這班以色列人一起經過這45年.
曠野是一個什麼樣的地方呢?.
象徵著人生的什麼樣的階段呢?.
我們有兩個方向大家去思想.
曠野是一個健步,閒步的歷程.
我們前途未明的人生的路途.
也沒有一些工作的成就.
也沒有一些別人的讚賞.
去掩蓋著我們生命的真實.
那種沒有東西可以遮蓋的實況.
我想問大家有沒有試過.
經歷過這些階段.
我有一個朋友.
以前開了一間叫曠野書店.
不知道大家有沒有機會去.
那時候在油麻地幾年時間.
就已經關門了.
他說開這間書店.
就是在反映他人生的狀態.
好像給自己一些放逐的時間.
給自己放任一下的時間.
放下任何的工作,任何的服侍.
去經過那幾年.
然後再上路,再走自己人生的路.
這種閒步健步的狀態.

$^{481}$不知道大家有沒有經歷過呢?.
另一些朋友經過過.
我想起他在工作上.
經歷過一些待業或者失業的時間.
他在雷曼金融風暴的時候.
本來之前是一個.
都算是在銀行界一個很有成績的人.
突然之間他的銀行就倒閉了.
大量裁員.
他本來是一個很高級的崗位.
就突然間什麼都失去了.
於是就很茫然.
不知道下一步怎麼走下去.
然後就約一些朋友聊天.
平時我都沒那麼空閒.
平時現在就很多時間.
跟朋友聊天,喝下午茶,聊聊天.
過一些比較休閒的生活.
等候著上帝下一步怎麼走下去.
不能夠去決定自己.
但卻是一種心智和信心的磨練.
不知道大家最近有沒有看.
一套都很出名的戲叫《白日之下》.
一套真的很值得看的戲.
除了很多觸動的位置.
對我反思的其中一個位置.
其實是其中一兩個配角.
他選擇了謀生.
就把家人放在院舍被照顧.
然後他就感到後悔.
有時候我們的抗議就是這樣.
究竟我們在這個路口的時候.
是否要放下一切.
應該照顧家人.
以家人的相處,需要為先.
但你什麼都做不到.
或者你覺得在這段時間.
對自己,對社會,對很多事情.
都好像沒有貢獻.
但只是把時間.

$^{521}$陪伴一些有需要的人.
可能這種狀態.
也是那種曠野,生涯的狀態.
「堅步行步之下」同時間也是一種信心磨練.
提醒自己要謹守上帝吩咐,要操練.
上帝叫我們追求良善.
叫我們怎樣去對人.
要我們怎樣去持守一些價值觀.
在這個狀態當中.
好像很被動,一無所有的時候.
其實需要繼續去謹守上帝給我們的吩咐.
要仰賴上帝每天的供應.
每天上帝給我們足夠的恩典.
去過這段日子.
不知道大家有沒有經歷過.
這些曠野的人生.
事實上我們的人生也不是.
一生裡面只有風調雨順,只有順境.
有時候高低起跌.
那些低谷,那些逆境.
其實就是可能在表示.
我們進入了一個曠野人生的階段.
需要我們好好地持守.
好好地等候上帝幾個階段.
接著嘉立.
預備進入人生的85歲之後的一個新的階段.
我很喜歡她一句話.
她說現在的力量和那時的一樣.
所以我特地在這一段裡面作一個標題.
人生的下半場.
不知道兄弟姐妹你有什麼目標呢?.
有沒有一些東西一直存在心裡面.
你很想完成.
這個時候是時候了.
你要踏出一步去完成呢?.
嘉立大概在她的人生當中.
等了45年.
當然不是白等.
這個時候可能是一個上帝安排的時間.
所以她和約書亞說.

$^{561}$我要上路了.
我要進到上帝應許給我的地方了.
我相信上帝要完成.
祂在我45年前.
應許過給我的一切的話.
要應驗了.
這個相信是.
嘉立到了這個狀態的時候.
覺得是一個新的階段的開始.
在第十節.
我再讀給大家聽.
我現在看.
耶和華照她所說的話.
小我活了45年.
當以色列人在抗日飄來的時候.
耶和華曾對摩西說過這話.
現在看啊我已經85歲了.
現今我還很健壯.
像摩西派我出去的那一天一樣.
無論是戰爭是出入.
我現在的理論和那時都是一樣.
請你將耶和華那天所說的那地.
山地分給我.
就是有時候我們看自己.
我沒問題啊很棒啊.
如果大家有機會知道.
人生只會越來越摧殘.
如果你跌過傷過那些.
你是沒辦法回到最.
完全百分之一百的那種狀態.
人生這種身體的磨損.
身體實質的體能下降一定會發生.
但是嘉立很厲害.
我85歲和我40歲一樣.
還可以出生入死.
還可以打仗.
可以進入那個地方對付那些亞納族人.
但是你看清楚一點.
其實他沒有像40歲的時候.
那麼一聲號令.

$^{601}$那麼有信心.
就是說或許耶和華會照他所吩咐的.
與我同在.
我就可以將他們趕出去了.
這個或許很有意思.
是不是.
如果40歲的時候.
或許當時.
這些以色列人沒有背誦上帝的時候.
真的要約書亞和嘉立.
去做統帥帶領著以色列人.
那時候他應該馬上衝了出去.
但是到了85歲就不是這樣了.
他說雖然我還可以走得走得.
但是如果耶和華真的要與我同在.
我就應該進入這個地方.
保守很多.
可能他要想的東西都多了很多.
他都未必有在自己身上.
有那麼有十足的信心同伴.
很多東西都要計算.
或者我們進入一個階段.
都不同了.
你要想想不只是想自己.
你想想你的家人.
想想你自己擁有過一切的東西.
你怎樣去處理呢.
這個就是我們隨著年紀長大的時候.
就不像年輕的時候那種.
可以那麼意氣風發地說.
我要一生都為上帝.
很轟烈地去活的時候.
嘉立這樣說的時候.
他不是將他的目光放在自己力量上面.
甚至環境沒有變的.
那些亞納族人一樣那麼高大.
他們的城一樣那麼堅固.
但是沒有變的就是.
好像我剛才說的.
經文都說他是全心地去跟從耶和華.

$^{641}$是上帝的呼召沒有變.
是上帝的應許說話沒有變.
他固定了這件事.
時間到了.
他就去回應了.
或者上帝仍然今天對他繼續使用.
這個就是一個人.
有時候隨著自己的智慧年歲成長的時候.
人都會謙卑下來.
他用一個比較謙厚的心去回應上帝.
到最後結語的時候.
特別是到約書亞記最後那幾節.
是不是?.
剛才我也說過了.
十四節就說.
所以希伯倫承了基尼洗族.
耶夫尼的兒子嘉立的產業.
直到今天.
因為.
不是因為他一樣那麼強壯.
不是因為那些人變了.
因為他全心去跟從耶和華.
以色列的上帝.
希伯倫從前叫做.
十五節再說.
希伯倫從前叫做基列阿巴.
基列的意思是一個城市的意思.
阿巴是亞立族最尊貴的領袖.
以前那個地方就叫做阿巴城.
就是用他的人名.
建造那個城市的名稱.
大家就叫做希伯倫了.
於是那個地圖裡面.
太平.
從此沒有戰爭.
以色列人正式進入了.
按著上帝的應許.
去佔有了他們這塊應許之地.
弟子們今天.
了問自己的時候.

$^{681}$上帝又給我們每一個.
擁有了什麼呢?.
我們的能力.
如果這個.
我經常想這個好像一個表.
給自己大家.
你試試填一下.
你有什麼能力呢?.
我有音樂的才華.
我有電腦的能力.
我讀了很多書.
讀了很多學位.
你可以寫下很多你在.
直到今天的時間.
你拿了的一些的qualification.
你寫下自己的履歷.
寫下很多經驗.
能力.
才幹.
你有很多的資源.
建立出來.
你有很多.
自己能夠判斷到很多的.
人生起起跌跌的經歷之下.
你的判斷的能力.
當你知道自己擁有這麼多這些東西的時候.
接著再問的問題就是.
你願意為上帝.
投資自己這些擁有的東西.
有多少呢?.
我剛才不是問大家要全心.
好像迦勒那樣.
有時候我們都未必做了全心全意.
我們願意有多少.
是給上帝呢?.
就從我們願意給那裡開始.
然後求主幫助我們.
可以更多更多.
以致全心都歸給我們的上帝.
全身根從來講一個.

$^{721}$不為心裡面感到憂慮.
和懼怕的事情東西.
所去搖動我們.
環境因素繼續是.
很多不明朗.
甚至很多暗湧.
可能你自己的.
身邊的人.
好像迦勒那樣很多那些第一代的.
以色列人.
他們不是跟隨著耶和華.
他會在這裡.
可能會牽動到你.
那種存在價值.
我是不是要跟隨他們一樣.
這樣去回應呢?.
我們是不是要發怨言呢?.
這是當時以色列人喜歡做的.
是不是?.
常常這樣做.
但是迦勒他堅持.
專心全心在上帝身上.
然後接著就是他應許實現.
上帝必定會在我們生命當中.
成就他的作為.
是不是?.
他拿著一件事.
四十五年前講過一句話.
他真的會發生的.
可能他很心急.
是不是?.
但是上帝給他抗夜的日子.
是一個準備磨練的時間.
在上帝的時間當中.
我們可以回應他.
迦勒他經過了四十五年抗夜.
沒有忘記上帝的應許.
去回應上帝.
以致在上帝的人生的規劃裡面.
那個時間裡面.

$^{761}$他投放自己.
去回應上帝.
可能你覺得.
為什麼不是四十歲呢?.
我沒有人知道.
為什麼要去到八十五歲呢?.
為什麼不是早一點好一點呢?.
是不是更多時間.
去服侍上帝更年輕力壯呢?.
可能這些是我們的角度會這樣想.
上帝要他經過四十五年的抗夜日子.
當我叫大家這樣去.
規劃自己的時候.
當然我讀經文的時候.
都一著不著地反省自己.
或者回顧自己.
所以我可以跟大家說一點.
我自己一些.
將自己分了階段的一些體會.
跟大家分享.
我不是四十五歲.
我三十五歲就.
舉家離開香港的時候.
就是.
都是離開錦宣家的服侍.
去了北美讀書.
當時就.
帶著兩個兒子.
大兒子是六歲.
小兒子就未夠三歲.
當時有很多的聲音.
我記得.
想起朋友怎麼說.
有些是唱反調.
就是.
你記住當時的香港.
和環境和今天完全不同.
他說你帶兩個兒子去.
一起捱嗎?.
一起受苦嗎?.

$^{801}$就是用一個這樣的角度.
來說我這個決定.
另外就.
更親近一點.
就是我要舉家去.
帶著太太一起去.
外父也有很多聲音.
他沒有直接跟我說的.
跟我太太說.
接著也有很多壓力.
為什麼什麼都.
就是放下了不工作.
還要去讀書.
這些親人的壓力的聲音.
然後當我自己預備讀一個.
PhD讀書的時候.
想著研究方向的時候.
找一些接近讀過書的.
過量人分享.
就是離開之前.
接著有一盆冷水.
你還做這些研究.
接著.
是不是不應該做呢?.
一個很大的信心的考驗.
當然我不敢說自己獨排眾議.
但是堅持.
其實是上帝帶的.
就應該要你走下去.
到了六年之後.
回到香港.
接著在FS服事.
到現在十一年的時間.
有時候跟太太分享的時候.
很多微妙的事情.
不能夠詳細跟大家分享.
有機會去見面的時候.
分享故事.
深入一點再跟大家說.
那個感受就是我們覺得.

$^{841}$原來對小朋友的祝福.
是比對我祝福還要大.
這個是想都沒想過.
當時按照環境裡面的意見.
可能覺得對家人不是很好.
但原來上帝的心意是很微妙.
很奇妙.
在他的時間.
除了我有我的時間.
可能在我們一起去的一家四口當中.
都有上帝的時間.
以致讓他們經歷最好.
上帝最大的祝福.
而我們特別服事當中.
都是一個自己能夠盡力去做的崗位.
繼續去服事我們的主.
弟兄姊妹.
我不知道對於你們來說.
下一步的時候要怎樣走下去呢.
究竟我們的人生當中.
也都有沒有好像.
按照像迦勒這樣說.
我們願意按照上帝的心意.
排除我們身邊的就算很多的困難.
或者很多的意見也好.
只要你覺得這是上帝安排你的時候.
你堅持下去.
你就可以看得見上帝給你的祝福.
或者好像對迦勒那樣.
最終的時候.
這個應許之地.
都能夠給到他.
按他的時間當中去踏上.
拍摩我們的人生當中.
不會停滯不前.
也都繼續有上帝的呼召在心中.
以致我們繼續去領受祝福.
成就他的應許.
我們一切都有禱告.
藉著迦勒的故事.

$^{881}$讓我們都去護過自己的生命.
可能我們會是像迦勒.
四十歲的時候那種衝勁.
或者進入了四十五年的曠野漂流日子.
又或者想不到有人生的再一步突破.
新的階段的開始.
但無論如何都好.
我們體會到.
迦勒的生命是全心跟從你.
按你的旨意而行.
按你的時間的預備去回應他.
主祈求你.
無論我們落在什麼的光景.
我們都做好準備.
特別當我們落在曠野日子裡面的時候.
我們可能會覺得消沉.
但同時間是讓我們學習.
等候你和信心磨練的功課.
以致到你再呼喚我們.
帶領我們的時候.
我們能夠有信心.
有一份準備好的心.
有一份明白你的旨意那種觸覺的心.
去回應你.
願你親自帶領我們頌弟兄姊妹.
讓我們的生命成為一個輕鄉的兒子.
我們的生命的投資是放在你的身上.
多謝你天父的恩典.
同心禱告奉主耶穌的名禱.
\newpage



\section{以賽亞書 49:1-7}
\label{sec:z_EDsyn25oM}
\textbf{2023-11-19 午堂《成為萬邦之光》以賽亞書 49\_1-7 講員:陳麗蟬牧師}
\newline
\newline
連結: \href{https://youtube.com/watch?v=z-EDsyn25oM}{\texttt{https://youtube.com/watch?v=z-EDsyn25oM}} ~~~~ 語音日期: 2023-11-22
\newline
\newline
\hyperref[sec:uKxjJHU8_7c]{\small{< < < PREV SERMON < < <}}
~
\hyperref[sec:index_chronic]{\small{[返順時目]}}
~
\hyperref[sec:index_scriptual]{\small{[返順卷目]}}
~
\hyperref[sec:Th6Drt1kZ1s]{\small{> > > NEXT SERMON > > >}}
\newline
\newline
以賽亞書 49:1-7
\newline
\begin{longtable}{cl}
\hline
\hline
章節 & 經文 (和合本修訂版)\\
\hline
49:1 & \begin{tabularx}{0.7\textwidth}{X} 眾海島啊，當聽從我！遠方的眾民哪，要留心聽！自出母胎，耶和華就選召我；自出母腹，他就稱呼我的名。 \end{tabularx} \\ \\ \relax
49:2 & \begin{tabularx}{0.7\textwidth}{X} 他使我的口如快刀，把我藏在他手蔭之下；又使我成為磨利的箭，把我藏在他箭袋之中； \end{tabularx} \\ \\ \relax
49:3 & \begin{tabularx}{0.7\textwidth}{X} 對我說：「你是我的僕人以色列；我必因你得榮耀。」 \end{tabularx} \\ \\ \relax
49:4 & \begin{tabularx}{0.7\textwidth}{X} 我卻說：「我勞碌是徒然，我盡力是虛無虛空。耶和華誠然以公平待我，我的賞賜在我的神那裡。」 \end{tabularx} \\ \\ \relax
49:5 & \begin{tabularx}{0.7\textwidth}{X} 現在耶和華說話，他從我出母胎，就造我作他的僕人，要使雅各歸向他，使以色列聚集在他那裡。耶和華看我為尊貴，我的神是我的力量。 \end{tabularx} \\ \\ \relax
49:6 & \begin{tabularx}{0.7\textwidth}{X} 他說：「你作我的僕人，使雅各眾支派復興，使以色列中蒙保存的人歸回；然而此事尚小，我還要使你作萬邦之光，使你施行我的救恩，直到地極。」 \end{tabularx} \\ \\ \relax
49:7 & \begin{tabularx}{0.7\textwidth}{X} 救贖主－以色列的聖者耶和華對那被人藐視、本國憎惡、統治者奴役的如此說：「君王看見就站起來，領袖也要下拜；這都是因信實的耶和華，因揀選你的以色列的聖者。」 \end{tabularx} \\ \\
[1ex]
\hline
\hline
\end{longtable}
$^{1}$鄭英姐妹平安.
平安.
上星期有四天.
我上了杜豐山退休.
其中兩天我和幾位同學.
去到杜豐山大十字架看日出.
這幅相是同學教我.
怎樣拍日出相.
是我拍的.
是很漂亮的.
不過後期加了工才有這麼漂亮的光彩.
有一天我們遇到一位老師.
他說杜豐山的童工.
最近在割草.
有些草越長越高.
發現了古蹟.
他叫我們一起去看古蹟.
於是跟著這位老師去看.
他說是最近才發現到的.
可能有幾十年.
甚至一百年跟著杜豐山一起建的時候.
在一個像塔的地方.
有四個面.
寫著一些鼓勵人祈禱的經文.
我們看完後很開心.
發現了新大陸.
於是就走了.
但走的時候.
我們發現到.
我們每一位同學的褲子.
都變成這樣.
很多的褒委.
我不知道叫什麼.
都黏在褲子上.
我更加是重災區.
無論是褲子還是鞋子.
還有鞋帶還有外套.
那件T恤.
全部黏滿了.
所以我們就要去做.

$^{41}$去撕掉這些針刺的行動.
有些同學說.
不行的一定要撕到全部清除.
否則會很痛.
如果撕不清.
有時撕不清根底.
那怎麼辦.
他說皮膚會敏感.
他看著我的褲子.
他說你的褲子.
可能要報銷了.
我說真的心痛.
但真的沒有辦法.
撕得完.
但又要去吃早餐.
又要去參加聚會.
所以都要走.
到下午的時候.
我都忘記了我的褲子.
有很多針刺的行動.
也不覺得我的腳有什麼敏感.
原來我發覺到.
到下午的時候.
這些褒委已經乾死了.
所以原本針刺人的行動.
就不再刺人.
也穿了褲子大半天.
我都不覺得有什麼.
皮膚不舒服.
有什麼問題.
我都不覺得有什麼問題.
我都不覺得有什麼皮膚敏感.
但又讓我們在整件事裡面.
我們不同的同學.
有不同的發現.
我自己發現了一樣東西.
當我們活在這個罪的權勢底下.
生命是痛的.
是腐敗的.
所以被這些蒙子.

$^{81}$扎實是很痛的.
也很擔心會有皮膚敏感.
但當我們.
藉著主耶穌基督的寶血.
我們可以能夠.
離開了罪的權勢的時候.
生命就不會再痛.
也不會再壞.
所以當陽光一路曬下去.
當時間一路過的時候.
這些蒙草就死了.
它不再影響到我的身體.
可以繼續平平安安地過這幾天的營會.
所以靠著主耶穌基督的寶血.
我們脫離罪的權勢.
上帝就說要去.
我們要去.
使萬邦能夠得到光明.
在這個以賽亞書裡面.
有四首我們稱為「獨人之歌」.
「獨人之歌」是預言到.
尼賽亞的身份.
尼賽亞的工作.
尼賽亞的身份.
尼賽亞的工作.
也藉著這四首詩歌.
教導我們信徒.
我們要去做一些事情.
剛才Patrick為我們讀到49章那裡.
是第二首「獨人之歌」.
是說到這個獨人.
他來到這個世界的使命.
是要將這個福音.
一路地去展現開.
而這個獨人呢.
這個獨人其實不只是指耶穌基督.
因為在聖經裡面.
他講這個獨人的時候.
一方面就真的指那些.
很卑微地位很低的奴隸.

$^{121}$但另外一方面.
是指到一些位高權重的大臣.
所以也是指到尼賽亞.
不過他那份謙卑.
在集體性來說.
就是指到以色列家.
一講到以色列家的時候.
就包括了你和我.
但也在個別性.
是專指到尼賽亞.
他謙卑.
甘願擺上生命.
完成了上帝.
差派他來到地上.
這個救贖的使命.
你看第一方面.
上帝選召他的獨人.
成為了這個萬邦之光.
在第一節.
在第一節是這樣說.
在中海島.
當聽從我遠方的眾民.
要留心聽.
上帝要這個以色列家.
留心聽他講說話.
他更加要向全世界的人.
宣告一件事.
就是尼賽亞.
他自己的獨人.
他將要來救贖.
所以第一節的後半節.
他說自出母胎.
也就是說選召我.
自出母腹.
他就稱呼我的名.
以賽亞宣之.
就預言七百多年之後.
尼賽亞.
耶穌基督.
他就來到世間降生成為人.

$^{161}$但很容易就說.
他這個萬邦之光.
我們不行.
上帝沒有給我們這樣的呼召.
但在我對宣教的認識裡.
我知道他整個宣教.
分開三個層面.
第一方面叫M1.
Mission 1.
按照《使徒行傳》第一章第八節說.
將福音傳到耶路撒冷和猶太傳遞.
這個很容易.
就是和文化的人傳福音.
和我們在香港周圍的鄰居.
傳福音.
第二就是將福音傳到薩馬利亞.
一些近文化的地區.
可能去到國內.
或者台灣.
太北.
因為太北很多人都是說普通話.
這些都應該可以.
雖然我們普通話是刨冬瓜.
但都可以刨到一點.
到第三類.
就是要將福音傳到地極.
是一個完全跨文化.
跨語言的使命.
這些是上帝特別的選照.
去到非洲.
去到很遠很落後的地方.
那些語言要全部重新再學.
那些飲食文化完全都不同.
這個上帝.
上帝早知道自己創造的世人.
我們這班世人.
很容易就陷入一個罪惡.
失敗.
苦困的光景.
但上帝就一早選照了這個僕人.

$^{201}$他就專指到尼賽亞.
但也進一步指到我們每一個信徒.
在我們沒有福的時候.
他已經提醒我們的.
因為說到耶穌基督.
在沒有福的時候.
他是聖靈感恩的.
並且耶穌基督這個名字.
是上帝藉著天使指明.
祂要叫耶穌的.
本來以色列人.
是上帝特別挑選的民族.
希望他們在這個世界上.
成為祭司成為先知.
成為宣教的民族.
但以色列一次又一次的背道.
一次又一次的犯罪.
其實我們也是.
很多時候犯了罪.
做錯了事.
認罪.
悔改.
沒多久又再犯.
但以色列人他們太嚴重了.
所以上帝就興起周圍的列邦.
來去管教他們.
也在他們第一步.
被擄到巴比倫之後.
七十年之後.
上帝又興起波斯.
這個波斯王採取懷柔的政策.
所以他對被擄的民族很好.
他又出錢又出力.
讓那些以色列人.
回到耶路撒冷重建聖殿.
這些以色列人當時很感恩.
但回頭看.
回到聖殿的時候.
很快就念念不忘巴比倫.
上帝就覺得.

$^{241}$這些民族搞不定.
於是就在這個.
希伯來民族裡面.
興起了一個中堡.
就是耶穌基督.
就是神的兒子.
他就將百姓從罪惡當中.
拯救出來.
如果再進一步講到.
上帝其實都同樣選召你和我.
成為了萬邦之國.
雖然我們未必一定.
每個人都去到第三重.
去到非洲.
去到很遠的地方.
但起碼在我們周圍身邊的人.
我們都需要和他們去講.
耶穌基督的救恩.
將這個被人盼望的救恩.
傳過給他們.
他在這裡說.
上帝要使我們成為萬邦之國.
我們作為信徒.
我們的生命應該是光明的.
有時我都會講到一個比喻.
假設這個房間裡面.
裡面是沒有窗.
沒有燈的.
黑漆漆的.
當你關上它.
半年 一年 兩年之後.
裡面全部都是蜘蛛網.
全部都是老鼠 蟑螂.
很骯髒.
沒有人喜歡進去.
但如果將.
當我們有一天.
我們拿著一個大電筒.
開了一道門.
一照著大電筒進去的時候.

$^{281}$接著會有什麼發生呢.
裡面的老鼠 蟑螂 蜘蛛.
全部都會飛出來.
所以上帝要我們作萬邦之光.
在地上高位去發光.
就是要在世人.
在一個暗淡.
在生命幽暗.
沒有盼望 沒有方向的裡面.
我們藉著.
我們燃點耶穌基督的光的時候.
叫世人得到盼望.
生命可以得到豐盛.
但上帝選擇你和我.
成為了無論你是做第一階段.
第二階段 第三階段都好.
不是因為我們有什麼長處.
也不是因為我們有什麼優點.
有什麼好行為.
而是上帝憐憫我們.
祂的主權淋到我們的生命當中.
祂要差派我們出去.
在我們差派出去的時候.
真是奉主命來工作的時候.
我們就經歷到了.
我記得在神學院的時候.
有兩位姐妹.
她去到在國內.
在很北的地方去短孫.
但有一次很危險.
她們真的被人抓到.
就抓她們去跟她們聊天.
審問她們.
就好像我們剛才讀的啟應經文.
你被人抓去見官的時候.
不用想著說什麼.
不用想的.
聖靈會引導的.
在那次裡面.
她們就在那裡.

$^{321}$前面就是審問她們的人.
怎麼辦呢?.
怎麼回答呢?.
怎麼說呢?.
後來就想到了.
一個懂普通話.
懂得聽普通話.
另一個又不懂得聽又不懂得說.
所以當那些人問話的時候.
問到她的時候.
她就不懂得聽又不懂得說.
所以就要這個人.
來翻譯廣東話.
其實上帝就給她們機會.
大家在那裡互通.
想想什麼的辦法.
想想什麼的答案.
是可以最好的.
上帝就是這麼奇妙.
以致我們為她們去工作的時候.
上帝就讓我們經歷到.
原來聖靈的工作.
可以讓我們在不同的時間裡面.
經歷祂奇妙的作為.
我們原本都是不配的.
我們也覺得我們沒有用.
所以能夠被神選召.
這兩位姐妹回到.
在我們神學院講見證的時候.
她說「全是恩典」.
所以弟兄姊妹.
上帝都呼召了你和我.
成為了他們的僕人.
我們要好好地珍惜.
但我們也會容易說.
「哎呀我很害羞我很害怕」.
「我很怕被拒絕」.
「我不懂得說」.
人人都怕被拒絕.
但我們想到.

$^{361}$他們拒絕的時候不是拒絕我們.
而是拒絕耶穌基督的時候.
我們就將這個不好的感覺.
交給耶穌.
叫耶穌給我們.
在使命上我們有份.
在聖經裡面說到.
使命給了我們.
若果我們不存的時候.
我們就有禍了.
第二方面.
上帝要我們成為了萬邦之光.
祂不會扔下我們不管的.
祂會給我們栽培.
在第二節.
祂使我的口如塊刀.
把我藏在祂手刃之下.
又使我成為無力的箭.
把我藏在祂的箭袋之中.
這個刀和劍.
都是非常鋒利.
是說到什麼呢.
就是上帝的話語.
上帝的話語.
就像希伯來書所說.
可以去刺透人心.
將骨折和骨折之間.
都可以去否開.
所以當我們宣講耶穌基督的時候.
講耶穌基督的話的時候.
我們不需要怕被人拒絕.
不需要別人覺得不中聽.
上帝的話一出的時候.
這個能力就出來了.
他說.
祂使我們的口快如刀.
然後把刀.
藏在祂的手刃之下.
就是拿緊了刀.
當你把刀磨利之後.

$^{401}$你拿緊在手上.
即是怎樣.
隨時可以用.
隨時就可以用.
然後這個箭.
箭頭磨利之後.
就放在箭袋之中.
隨時可以拿出來射.
如果你把刀不利.
你出去打仗的時候.
你不會拿一把鋒利的刀.
殺人獸的刀出去.
不會的.
如果箭頭不尖.
你也不會拿出來打仗.
用來射人.
就是說.
上帝親自到來.
每一個屬於祂的僕人.
所以.
當我們說我們沒有膽量.
我們很害怕.
我們不懂得說的時候.
接受訓練.
在上帝的話語裡面.
多去讀.
多去接受上帝的訓練.
但另外.
上帝也會在我們的生活.
我們周圍發生的事裡面.
會訓練我們.
我記得應該是七八年前.
我就帶著這個福音樂曲隊.
去到澳門.
有四日三夜的短訓.
最近也跟一些團契分享.
當時黎光偉傳道.
就在澳門.
他就訂了一輛小巴.
二十四人的小巴.

$^{441}$跟我們在澳門.
的公園裡面.
我們就下去唱一下方月曲.
跟人家派一下單張.
在我們出門之前.
去三町四足.
就不斷地強調.
去到這個三角公園的時候.
千萬不要下車.
因為那裡有三山五嶽的人在那裡.
那些人有自己的地盤.
他們也有拉耳鼓.
也有打洋琴.
整隊樂隊在那裡.
唱他們自己的樂曲.
不過因為那是他們的地盤.
任何人都不可以進去.
不可以去參與.
所以你們千萬不要下車.
千萬不要跟他們有任何的交往.
否則不知道會發生什麼事.
我們就記住記住.
所以當司機先生告訴我們.
快要到這個三角公園了.
每個人都好像有點害怕.
不知道會發生什麼事.
到了那個時候.
司機先生就在那裡繞一繞.
來去走.
你們看一下好了.
看下去那些人黑鼓甩骨.
又有一點兇神惡煞.
遠看都看到他們周圍的紋身.
現在就真的在地鐵周圍看到男男女女的紋身.
現在就不怕了.
但那時候很害怕.
他們就說停在那裡讓你們看清楚一點.
突然間有人在後面拍一拍我.
下去了.
那我就說下去了.

$^{481}$司機先生又開了門.
那我就一起下去.
我心想後面的頂層人員拍我.
叫我下去.
我作為傳道人不下去.
這麼沒膽.
哪裡有面子.
那就下去了.
但很感恩.
最初都害怕.
他們聽著聽著看著.
接著就開始大膽了.
就拿著他們的歌詞.
走上去問一下.
我們很想和你們唱一下.
最初很黑臉的.
唱唱唱.
要一百元三首歌.
不是一百二十元三首歌.
我們給我們給的.
不過這些歌詞是不同的.
但我們知道他們的樂曲.
有什麼春花秋月.
就像那些調名.
何唐乙式.
我全部都作的.
他們就懂得了.
總之說到調名.
他們就懂得彈懂得拉.
大家越唱越開心.
我們說唱三首給一百二十元.
很開心.
然後我們說這麼興奮.
我們就開始派福音擔將.
就跟周圍的人說.
結果上高的師傅說.
這麼熱鬧.
沒試過唱得這麼熱鬧這麼開心.
不收錢.
我們說不要緊.

$^{521}$我們預了.
我心想我們也預了給他們.
我們也給了.
他們說你們也要給.
給你們八折一百元.
但結果我們唱了五首歌.
到晚上我們檢討.
檢討的時候我就說.
誰這麼大膽在後面拍我.
叫我下去.
結果他說沒有人.
沒有人拍我.
也沒有人叫我.
他們每個人都很乖.
記住黎光遺傳道的吩咐.
我相信是聖靈的工作.
我相信在很多短宣宣教的經歷裡.
我們是不斷經歷的.
當我們不踏出去的時候.
經歷不到上帝的大能.
上帝在一些艱苦或危險的地方.
祂會給我們幫助.
是聖靈的工作.
我們到今天為止.
我們有時在樂曲隊.
我們來講話的時候.
我們都興奮.
雖然那天每個人都不停地稱讚我.
很勇敢.
我實在很羞愧.
我下台的時候才罵誰這麼壞.
有時我們覺得.
一些很危險.
很艱難.
沒什麼可能的事.
但其實在上帝的眼裡.
只是一些很小很小的事.
上帝給我們一個很大的使命.
就是要將祂的救恩.
給人豐盛生命的福音.

$^{561}$要不斷地傳開去.
我們更加要作愛邦人的光.
在上帝的裡面.
凡事都能.
上帝是那位行其事的上帝.
在祂的裡面.
是無所不能.
最後一方面.
上帝就應許.
我們出去為祂們去事奉的時候.
是會有成果的.
在第六節.
向我還要使你作萬邦之光.
使你施行我的救恩直到地極.
第七節.
向君王看見就站起來.
領袖也要下拜.
第七節就講到彌賽亞.
他被釘十字架的時候.
很多的羞辱.
但他成就了救恩.
於是上帝就給他.
這個超乎萬明之上的名.
使無論在天上的地上的地底下.
一切受造之物.
都無不口稱耶穌基督為主.
使榮耀歸於父上大.
所以就預言君王都要起來.
領袖都要下拜.
就講到我們每一個人.
肯去事奉的時候.
就有這個成就.
上帝給我們有成就的.
在這個事奉的過程裡面.
我們或者會覺得.
有時覺得沒有用去徒然.
跟整隊的短孫子一起出去的時候.
一開始的時候.
我記得在澳門碼頭.
就有些是是非非.

$^{601}$又有些人有些人爭執.
又有些人覺得死了.
不知怎樣怎樣.
很多的事情.
經文都告訴我.
有時我們會被藐視的.
有時我們會被批評.
被誤會的.
但當人去看低我們的時候.
不要緊.
在第六節裡面.
耶和華告訴我們.
去看我們為尊貴.
因為耶和華就是我們的力量.
在第六節講到.
所以我們知道.
我們只要忠心的去事奉.
上帝自己去做判斷.
人說什麼.
我們不需要去理會.
我們最主要是忠心去為主工作.
在羅斯福做總統的那個年代.
有一對老孫教士.
他們在非洲做了很多年.
退休後回到美國.
做孫教士.
其實他們真的很艱辛.
他們沒有退休金.
亦因為在非洲這麼落後的地方.
他們染了很多病.
所以他們回到美國的時候.
或多或少都對他們退休後的生活.
有些擔憂.
他們坐船回到美國的時候.
剛好羅斯福總統放大假.
去了非洲打獵.
就回到美國.
他說當埋岸回到美國紐約的時候.
很多政要和傳媒歡迎總統.
放完大假回來了.

$^{641}$但這對老孫教士夫婦.
心裡有些心疼.
這個總統放完大假.
那麼多人歡迎他們.
但我們在非洲幾十年.
回來是那麼冷清.
為什麼我犧牲了自己那麼多東西.
回來的時候沒有人理會呢.
這個老孫教士有點不濟.
太太就跟他說.
你自己問上帝吧.
所以這個孫教士回到家之後.
很不開心就進了房間裡去祈禱.
幾個小時之後.
他就很開心地走出來了.
太太問他.
哇!上帝跟你說了什麼.
這麼開心.
他說我一路祈禱.
我跟上帝甚怨.
我覺得很不值.
為什麼羅斯福那麼受歡迎.
我這麼冷清.
這麼慘淡.
上帝在我禱告當中出現.
拍一下肩膀.
他說你還沒回家.
那怎麼會有人歡迎你呢.
這個老孫教士突然間去醒覺.
上帝告訴他.
你的賞賜不是地上的賞賜.
你在將來回到真正的天家的時候.
自然有耶穌基督.
有天使,天君一起來歡迎你.
所以,靈姐妹.
我們今天在地上高.
這只是我們暫時的家鄉.
我們真正永恆的家鄉.
是在天上高.
所以將在地上高一切的金錢.

$^{681}$一切的人物.
一切的地位.
我們都不需要看得那麼重要.
我們最重要的.
就是我們在靈魂回到天家的時候.
我們肯定主耶穌基督.
和天使,天君和眾聖徒.
一同列隊來歡迎我們.
並且主耶穌基督.
一定會和我們每一個忠心事奉祂的人.
歡迎忠心耿耿的僕人士女.
回到天家.
所以求主幫助我們.
我們的焦點要在耶穌基督的身上高.
我們千萬不要定睛在地上高的人和事物.
最後,上帝呼召我們每一個人.
成為祂的信徒.
成為祂的僕人.
要我們作祂萬邦之光.
要我們離去這個.
將舊恩傳到.
由耶路撒冷傳到地極.
所以在全方位的事上高.
在宣教的事情上高.
我們對上帝有什麼回應呢?.
我們同心禱告.
因為我們先得著耶穌基督給我們豐盛的生命.
使用我們.
將我們生命裡面的一切.
或者禱告.
或者時間,金錢.
甚至差派我們.
要我們去到很遠的地方.
主求你使用我們.
將主你的呼召.
臨到在你要選照的人身上高.
被主使用是蒙福的.
我們的禱告是奉主基督的聖名.
\newpage



\section{歌羅西書 3:15-17}
\label{sec:Th6Drt1kZ1s}
\textbf{2023-11-25 晚堂《心存感恩》歌羅西書 3\_15-17 講員:盧思源姑娘}
\newline
\newline
連結: \href{https://youtube.com/watch?v=Th6Drt1kZ1s}{\texttt{https://youtube.com/watch?v=Th6Drt1kZ1s}} ~~~~ 語音日期: 2023-11-28
\newline
\newline
\hyperref[sec:z_EDsyn25oM]{\small{< < < PREV SERMON < < <}}
~
\hyperref[sec:index_chronic]{\small{[返順時目]}}
~
\hyperref[sec:index_scriptual]{\small{[返順卷目]}}
~
\hyperref[sec:Q0yYJrC6Zoc]{\small{> > > NEXT SERMON > > >}}
\newline
\newline
歌羅西書 3:15-17
\newline
\begin{longtable}{cl}
\hline
\hline
章節 & 經文 (和合本修訂版)\\
\hline
3:15 & \begin{tabularx}{0.7\textwidth}{X} 你們要讓基督所賜的和平在你們心裡作主，也為此蒙召，歸為一體。你們還要存感謝的心。 \end{tabularx} \\ \\ \relax
3:16 & \begin{tabularx}{0.7\textwidth}{X} 當用各樣的智慧，把基督的道豐豐富富的存在心裡，用詩篇、讚美詩、靈歌，彼此教導，互相勸戒，以感恩的心歌頌神。 \end{tabularx} \\ \\ \relax
3:17 & \begin{tabularx}{0.7\textwidth}{X} 你們無論做甚麼，或說話或行事，都要奉主耶穌的名，藉著他感謝父神。 \end{tabularx} \\ \\
[1ex]
\hline
\hline
\end{longtable}
$^{1}$好 弟兄姊妹平安.
願意我們開始之前.
我們都有個祈禱.
你的殿當中.
一同去敬拜.
一同去聽你自己寶貴的話語.
一會兒孩子要分享.
你自己的話語的時候.
求真理的聖靈.
與我們眾人同在.
無論我們說的.
聽的.
我們都進入.
你自己的真理當中.
願以我們被主你自己的話語去得著.
天父我們恭敬的.
將以下聖餐交上.
奉主耶穌基督名字而求.
阿們.
弟兄姊妹不知道大家.
知不知道剛剛過去這個星期.
其實有一個節日.
大家知不知道.
有些點頭疙瘩.
其實就是剛剛過去的星期四.
就是感恩節.
如果在美國的話.
每年進入了十一月.
他們多數都已經稱為感恩月.
每年十一月的第四個星期四.
就叫做感恩節.
其實我們教會.
我們在十一月除了.
我們都是一個猜拳月份之外.
其實我們都是一個感恩的月份.
如果說回感恩節的時候.
在美國是一個很全國性的節日.
如果說感恩節.
其實要追溯回原來在當時.
清教徒去到美洲大陸的歷史.

$^{41}$原來在公元大約十六世紀末.
十七世紀初的時候.
當時英國.
他就發起基督新教徒的.
宗教改革運動.
因為他們當時.
一班的清教徒.
他反對當時羅馬天主教的操控規矩.
所以當時英國的國王.
亨利八世.
他在1534年的時候.
就宣告英國的國王.
成為教會的領袖.
造就了當時英國國教的獨立.
而且削弱了當時羅馬天主教的勢力.
不過很可惜.
英國當時的教會.
慢慢地.
開始淪為腐敗.
於是就有一班的清教徒.
他很想改變教會.
但是因為這樣.
就當時被政府的官員去迫害.
導致他們要去其他國家避難.
直至到1620年的9月.
有102位的清教徒.
他就搭一架叫做五月花號的木船.
橫到大西洋.
去美洲的大陸.
就在11月21日.
抵達了美國.
現在叫做普利茅斯的一個地方.
不過很可惜.
當時因為是冬天.
天氣很寒冷.
而且有很多傳染病.
使得那102人.
有一半的人.
是死了.
剩下50人生存.

$^{81}$到了第二年的春天.
因為當地有一些土著印第安人.
很同情這一班新移民.
所以就送了很多生活的必需品.
又教他們捕魚.
狩獵.
種粟米 種南瓜.
養火雞等等.
這一班新移民.
蒙上帝的賜福.
就風調雨順.
也因為印第安人的幫助.
獲得了一個大豐收.
終於就渡過了生活的難關.
過著一些比較安定的生活.
於是那一班的清教徒.
就邀請這班印第安人.
一起參加一個感恩的慶典.
厚賜了這些恩賜.
就成為了今天感恩節的由來.
顧名思義.
感恩節.
就是向神發出感恩.
剛才說如果在美國或加拿大這些國家.
一家大小.
是會去圍爐述說.
一年裡面神給予的那份恩典.
那份感恩.
回到聖經.
其實感恩.
我們很多時候在聖經裡面.
都知道是一個聖經的教導.
都是神喜悅我們去做的.
今天很想和大家一起去看看聖經.
究竟怎樣可以學習.
我們真是有一份心存感恩的心.
鼓勵大家可以打開聖經.
我們今天看的經文.
就是在歌羅西書.
如果我們在歌羅西書的時候.

$^{121}$我們知道.
這個是使徒保羅所寫的一封信.
當時他在監獄裡面寫的.
他是寫給哥羅西教會的基督徒.
今天我們看三章十五至十七節.
不過如果我們看.
早一點在三章十二節那裡.
我們看到.
保羅提到.
基督徒是神的選民.
是聖潔無礙.
所以他在開初的時候提到.
我們要脫去舊人和舊人的行為.
要穿上新人.
保羅就在當中勸勉信徒.
要彼此寬容.
體諒.
彼此饒恕.
要在基督的愛裡面合一.
而接著我們今天剛才讀的三節經文.
其中其實都有一個主題.
在短短這三節的經文裡面.
很強調感恩.
在十五節那裡看到.
他說.
在第十六節最後那裡也是這樣說.
他說.
「以感恩的心歌頌神」.
在第十七節那裡.
我們也看到.
我們看到單單這三節的經文裡面.
保羅都很強調感恩.
保羅在當中提醒.
作為神的中心的子民.
作為神所揀選的選民的時候.
有時我們現在很容易被視為禮貌.
或者禮節.
例如別人對你好.
你就說一聲謝謝.
或者為我們有的東西去感恩.

$^{161}$不過如果我們看保羅在這裡所說.
他的感恩其實是涉及很重要的.
在他們內心裡面一個信念的核心.
甚至是一個神學信念的核心.
就是在這段段落裡面.
我們看到保羅和其他的書信都強調.
他提到人感恩的元素.
首先第一是.
所以如果這樣的話.
而且不單單是維持人際關係那麼簡單.
而且第二保羅在這裡很強調.
在保羅來說.
他主要是回應神.
藉著他兒子耶穌基督成就的恩典作為.
所以保羅在這裡說的感恩.
是對神恩典的回應.
如果回過頭來.
我們剛才說感恩是作為我們基督徒.
一個新人身份的標記的話.
我們一起看看保羅對舊人的描述.
是一個怎樣呢.
我們看羅馬書一章21節裡說.
因為他們雖然知道神.
卻不把他當作神榮耀他.
原來不信的外族人.
當時經歷神的憤怒.
原因是因為他們不尊神為神.
這裡不單單說不信的人.
也當中說不當作神榮耀神.
即是不尊重神.
我們在這裡看到.
保羅把感恩連於.
信服基督的主權.
承認祂是創造主.
承認祂是一切美善的源頭.
如果我們讀聖經的時候.
很多時候聖經都說.
如果我們過分自我中心的時候.
堅持只靠自己的時候.
就好像等同拜偶像的行為.

$^{201}$相反感恩就是承認.
我們是需要創造主.
我們需要將我們關注的焦點.
回到神的身上.
有些心理學家這樣說.
不懂得感恩的人.
原來在人際關係裡.
是一種自戀的標記.
他說原來你需要感恩的時候.
是需要謙卑承認自己的不足.
是承認我們原來是需要人幫的.
心理學家這樣說.
他認為如果一個人只是關心自己.
一個自戀者.
他雖然可能都會遇到一些值得感恩的事.
不過可能他經常拒絕.
就好像承認自己不足.
侮辱了他那個自我.
所以一個自戀的人.
很多時候都不懂得怎樣去表達感恩.
相反相反就是說.
懂得感恩的人.
原來是需要有一份謙卑的心.
回到聖經裡面.
剛才保羅說感恩.
其實是信服基督的主權.
就是承認祂是我們的創造主.
祂是一切美善的源頭.
因為耶穌基督的時候.
我們才可以成為剛才說的新人.
在基督裡面一個新造的人.
保羅在歌羅西書第15節這樣說.
就是要讓基督所賜的和平.
在你們的心裡作主.
基督的和平.
也就是解作基督的平安的意思.
基督的平安.
其實這裡可以有兩種解釋.
一個解釋就是那份平安.
就是我們基督徒彼此之間那份平安.

$^{241}$那份和睦相處的平安.
另一個解釋就是認為.
平安就是神賜給我們基督徒內心的平安.
特別當我們遇到困難.
在壓力底下.
我們在神面前那份不會失調的平安.
也在耶穌說.
祂已經賜下平安給我們.
這份平安不是世人所賜.
所以這裡的平安是說我們在神面前那份平安.
而這句裡面特別提到.
要讓基督的平安在我們心裡作主.
作主原來有個意思.
就是在運動場裡面做評判或裁判員的意思.
這句說話就是說原來要讓基督的平安在我們心裡作主.
就是今日我們基督徒在我們的生活裡面.
無論遇到甚麼困難.
或者一些衝突.
或者一些紛爭.
或者在基督徒彼此之間的困難也好.
我們要怎樣.
我們就是要讓基督的平安作為我們決定的原則.
讓祂的平安成為我們一切決定的元素.
或者我反過來說.
我們想想如果失去了這份基督的平安的時候.
當你遇到事情失去了基督的平安的時候.
大家去想想.
當然遇到一些很順利的事情的時候.
當然容易.
但相反遇到一些逆境.
或者遇到一些你不想遇到的事情的時候.
因為可能我們當遇到一些逆境.
或者遇到一些我們不想遇到的事情的時候.
第一個反應可能我們就會問.
為什麼.
為什麼會發生.
甚至我們會抱怨.
為什麼這些事情會發生在我們身上.
但是當我們去抱怨的時候.
其實就令我們失去了那份平安.

$^{281}$而抱怨也偷走了我們那份平安.
抱怨和感恩其實是對立的.
因為當我們抱怨的時候.
問為什麼這些事情會發生在我們身上.
或者問為什麼的時候.
我們很多時候的專注力就在困難上面.
我們看不到那個問題解決的方案.
甚至很多時候當我們去抱怨的時候.
我們伴隨的情緒是什麼呢.
可能是不開心,憤怒,沮喪,不滿.
這些事情都讓我們失去了內心那份平安.
所以為什麼保羅他在這裡.
他是用一個命令的語氣.
他命令信徒要讓基督所賜的平安.
在我們的心裡面作主呢.
因為當我們這樣的時候.
我們有神給我們的那份平安的時候.
我們用神的平安掌管我們的心的時候.
以致當我們無論遇到任何的情況.
接著第十六節他繼續這樣說.
當用各樣的智慧把基督的道豐豐富富地存在心裡.
用詩篇贊美詩靈歌彼此教導.
互相勸誡.
以感恩的心歌頌神.
如果在《新漢語》譯本這一段是這樣說.
要讓基督的道豐豐富富地住在你們裡面.
以各樣的智慧彼此教導.
互相勸誡.
以詩章聖詩靈歌以感恩的心歌頌神.
基督的道其實是一個很特別的.
因為我們很多時候看到都是說.
神的道 主的道.
這裡很特別強調是基督的道.
顯然當然是說耶穌基督那份教訓.
也是說神的真理.
而這個基督的道要怎樣呢.
這裡說形容基督的道要住在我們的心裡.
住的意思.
好像家裡居住的意思.
是長久地居住在我們裡面.

$^{321}$這裡補充提醒說.
我們要讓神的道 基督的道長久地住在我們裡面.
怎樣住呢 就是豐豐富富的方式.
即是我們要以基督的道在我們的生命裡去作主.
而且這裡繼續說.
要讓各樣的智慧彼此教導.
懷著感恩的心去歌頌神.
剛才我們也提到.
其實保羅一直都在說.
今天基督徒蒙照.
不只是單單是說.
我們為著我們的東西是感恩之餘.
而敬拜也是人回應神恩典的責任.
在這篇經文裡.
提到要用詩篇 贊美詩 靈歌.
被視為教導和勸戒的途徑.
可能很難明白.
即是怎樣呢.
即是唱歌這樣去互相勸戒 教導.
雖然我們知道我們唱歌.
例如剛才我們唱詩 歌頌的時候.
我們的對象是神.
但原來我們藉著一同去敬拜.
一同唱詩的時候.
我們是彼此教導的一個工具.
因為藉著當我們一同集體敬拜的時候.
我們可以彼此提醒神的作為.
向神發出一個真誠的敬拜.
藉著在敬拜裡.
我們宗鄉和神的關係.
我們橫向和弟兄姊妹會眾之間的層面.
敬拜其實添上了一個很重要的意義.
原來我們回來教會.
或者我們說我們敬拜的時候.
不只是單純表達一個個人的感受.
原來在當中 在敬拜過程裡.
我們紀念上帝為我們所作的事.
我們紀念神是那位掌管萬有的神.
我們和神有一份愛的關係.
有一份立約的關係.

$^{361}$所以弟兄姊妹.
為什麼我們經常說我們要回來教會.
我們要一起去敬拜神呢.
因為原來當我們集體一同去敬拜的時候.
我們有彼此提醒 彼此教導的作用.
我們唱詩敬拜的時候.
我們不只是單單去唱一首詩歌.
而是我們在感恩當中.
我們在詩歌裡.
我們去紀念.
我們去稱頌神在我們身上的恩典作為.
我們真的用口唱心和.
用讚美去回應神的恩典.
記得之前有一次短宣的經歷.
當時是同一班的弟兄姊妹.
參與在德國當地教會的崇拜聚會.
有機會和當地的弟兄姊妹一起去敬拜.
剛好在聚會裡.
大家不約而同都選了一首歌.
剛才我們也有唱的.
就是歡欣這首歌.
不約而同大家一起選了.
我記得當時.
德國的弟兄姊妹.
他們用德文去唱歡欣這首歌.
之後我們一班的弟兄姊妹.
我們用粵語去唱歡欣這首歌.
之後大家就用英文去唱這首詩歌.
在過程中.
我不懂得德文.
或者他們也不懂得粵語.
可能我們也不太明白.
大家對方的語言.
但是藉著同一天去敬拜的時候.
我們很感受到.
我們在主裡面的合一.
主內一體的感覺.
還有在敬拜裡.
我們除了唱詩歌.
感到很歡欣和感恩之外.

$^{401}$藉著一同敬拜的過程.
神在當中讓我們看到.
祂是萬國的主.
祂是各國各族各方的主.
在敬拜的當中.
我們彼此也有一份提醒和教導.
就像剛才經文所說.
我們藉著詩歌裡.
有互相的勸誡和教導.
回到聖經裡.
接著在第17節.
那裡這樣說.
你們無論做什麼.
說話或行事.
都要奉主耶穌的名.
如果在新譯本.
這句就是這樣說.
凡你們所作的.
無論是言語或行為.
都要奉主耶穌的名.
按照我們剛才一直說的理解.
我們知道感恩.
其實是一種認信.
也是一種屬於敬拜的範疇.
所以我們會明白.
為什麼保羅提醒信徒.
凡事都要去感恩.
特別我們再看.
在帖撒羅尼加前書.
有兩節經文.
帖撒羅尼加前書.
都是保羅所寫的書信.
當中更加指出.
原來信服基督的主權.
是跟聖潔的生活.
和感恩的生活有關係.
在帖撒羅尼加前書四章說.
神的旨意就是要你們聖潔.
在帖撒羅尼加前書五章十八節.
寫著凡事謝因.

$^{441}$因為這是神在基督.
向你們所定的旨意.
如果你看這兩節經文的關係很清楚.
保羅在說.
這就是神的旨意的公式.
來定義什麼是神的旨意.
我們一併去理解的話.
保羅其實在這裡說得很清楚.
原來感恩就好像.
那一個聖潔的生活.
回過頭來.
也解釋了今天我們在歌羅西書.
剛才說三章十二節那裡說.
神所揀選聖潔蒙愛的人.
之後的段落.
這一段的段落.
之後保羅就細說.
關於感恩的主題.
我們明白為什麼保羅會這樣說.
所以對於今天我們.
我們這群基督徒來說.
感恩不只是局限在.
人際關係的層面.
也不只是局限在.
感恩節的假期裡.
我們提到感恩.
相反我們作為基督徒.
作為神所揀選聖潔蒙愛的子民的時候.
我們感恩是要承認.
神在我們身上所做的一切是主權.
因為唯有神才是一切美善的源頭.
而且他所做的.
是為了要完成他對我們所定的旨意.
回到今天.
剛才說感恩就是承認.
我們所做的事.
有上帝一切的主權在當中.
但原來有研究發現.
在實際的研究發現.
原來當人去表示感恩或感激的時候.

$^{481}$是可以改變人的腦神經.
可以活得更自信更快樂.
原來感恩也是有一份力量.
有些科學家真的這樣研究.
我看到一些記錄.
他找了一群人.
在日記上寫感恩事情的人.
發現他們的心理狀態和大腦也有些不同.
比起普通人他們更加自愛.
不會受負面情緒影響.
原來感恩是可以帶來改變.
當中也有科學根據.
我分享一點給大家聽.
原來當我們去感恩的時候.
我們會減低釋放負面的情緒.
原來當他接受讚美或釋出善意的時候.
大腦負責情緒.
大家可能有些在座會知道.
原來在大腦裡面.
有個叫扁桃核和海馬迴的地方會被啟動.
這個有助於我們處理情緒.
這種處理可以幫助我們更加輕鬆和平靜.
除此之外原來感恩也可以減少痛楚.
原來當科學家研究.
當他經常表達感恩的時候.
身體原來會調節多巴胺.
多巴胺是一種神經的遞質.
它可以提升我們的情緒和快樂感.
也減少感覺上的痛楚.
而且在研究裡面.
科學家發現可以改善睡眠.
除了情緒處理.
大腦剛才說的海馬迴原來也會負責睡眠.
他說當我們感恩的態度的時候.
可以促進海馬迴的正常運作.
使得會有深度和質量好一點的睡眠.
而且感恩也可以調節壓力.
當面臨壓力的時候.
原來感恩在科學裡面研究.
是減少體內的皮質醇.

$^{521}$原來皮質醇是和壓力有相關的荷爾蒙.
當皮質醇降低的時候.
我們的身體就會很容易去面對負面的情緒.
提升心理的韌度.
自然你很從容地面對那份壓力.
而且感恩也可以減少焦慮和抑鬱的情緒.
因為感恩原來可以增強前額葉的腦神經活動.
前額葉就是大腦負責一些內疚,愧疚感和負面情緒的區域.
總括來說.
原來感恩可以幫助我們變得更加積極.
也更有活力.
感恩對我們心理健康和人際關係.
都有正面積極的影響.
怎樣去活出.
我們知道感恩原來是這樣.
我們怎樣可以將感恩的習慣融入我們的生活裡.
坊間如果大家看其實有很多的建議.
首先可能大家寫下.
這個可能覺得很老套.
但原來寫下一些感恩的事項的時候.
幫助我們去專注眼前美好的事物.
我有一個朋友他真的有一本感恩日記.
他每天都嘗試寫下當日三件他值得感恩的事和原因.
當他回顧一些感恩的日記的時候.
當中他會被提醒神過去對他的慈愛和信實.
或者可能大家即使沒有時間去寫下也好.
就盼望嘗試留一些時間去想想.
你今天有沒有一些值得感恩的事情.
可以和你的家人朋友去分享.
也幫助我們培養一份感恩的態度.
我曾經有一些朋友都是基督徒朋友.
在他們的家庭他們會有家庭崇拜.
或者夫婦二人之間他們會有祈禱.
他們可能每晚在祈禱的時候.
他們彼此都會分享當日有什麼值得感恩的事情.
藉此去培養看到日常生活中感恩的事.
或者透過默想.
我們都很少和家人去聊天.
或者透過自己默想.
我不知道大家會不會有這個習慣.

$^{561}$可能趁著自己洗澡的時間.
又或者是茶的時間.
去思想一下一天當中有沒有什麼是值得感恩的呢.
有沒有什麼是值得我們向神發出的讚美呢.
這些習慣都可以讓我們將感恩帶進我們的生活中.
最後 大家會怎樣去拍望呢.
今天我們剛才三節很短的經文.
但當中保羅都提醒我們.
我們要有一份感恩的心.
我們要常存感恩的心.
而感恩不是單單是為了我們生活中得到滿足的事情的回應.
相反 我們縱然在困難的當中.
縱然在艱難的當中.
都讓感恩的事提醒我們.
神對我們那份恩惠.
也提醒神是我們那位創造主.
祂正在掌管一切.
我們不是一個人去面對我們的事情.
拍望我們不要失去神基督的平安在我們心中.
願意我們活出那份心存感恩的生命.
我們一同去祈禱.
昔日你藉著保羅對我們的教導.
你提醒我們.
我們要向你凡事謝恩.
因為你是那位愛我們的主.
你是那位配得我們讚美的主.
或許在我們的生活中.
有時我們遇到順境的時候.
我們很容易懂得去感恩.
但當特別遇到一些逆境的時候.
我們很容易會失去那份平安.
求主幫助我們.
幫助我們不要失去基督那份平安在我們心中.
正如保羅所提醒.
願意基督的平安在我們心中作主.
以致當我們無論遇到什麼事情的時候.
我們都可以有一份感恩的心.
因為你呼召我們.
因為我們一切都是你所賜.
你是那位美善的源頭.

$^{601}$求主幫助我們.
幫助我們常常上傳感恩.
無論我們對身邊的人.
我們對你.
我們都是活出感恩的生命.
願我們將一切榮耀,頌讚歸給你.
奉主耶穌基督名字祈求.
阿們.
\newpage



\section{提摩太後書 1:1-5}
\label{sec:Q0yYJrC6Zoc}
\textbf{2023-11-26 加恩樓《代代相傳的信仰》提摩太後書 1\_1-5 講員:黃頌言姑娘}
\newline
\newline
連結: \href{https://youtube.com/watch?v=Q0yYJrC6Zoc}{\texttt{https://youtube.com/watch?v=Q0yYJrC6Zoc}} ~~~~ 語音日期: 2023-11-28
\newline
\newline
\hyperref[sec:Th6Drt1kZ1s]{\small{< < < PREV SERMON < < <}}
~
\hyperref[sec:index_chronic]{\small{[返順時目]}}
~
\hyperref[sec:index_scriptual]{\small{[返順卷目]}}
~
\hyperref[sec:nQU0VbnDk9Y]{\small{> > > NEXT SERMON > > >}}
\newline
\newline
提摩太後書 1:1-5
\newline
\begin{longtable}{cl}
\hline
\hline
章節 & 經文 (和合本修訂版)\\
\hline
1:1 & \begin{tabularx}{0.7\textwidth}{X} 奉神旨意，按照基督耶穌裡所應許的生命，作基督耶穌使徒的保羅， \end{tabularx} \\ \\ \relax
1:2 & \begin{tabularx}{0.7\textwidth}{X} 寫信給我親愛的兒子提摩太。願恩惠、憐憫、平安從父神和我們的主基督耶穌歸給你！ \end{tabularx} \\ \\ \relax
1:3 & \begin{tabularx}{0.7\textwidth}{X} 我感謝神，就是我接續祖先用純潔的良心所事奉的神，在祈禱中晝夜不停地想念你。 \end{tabularx} \\ \\ \relax
1:4 & \begin{tabularx}{0.7\textwidth}{X} 我一想起你的眼淚，就急切想見你，好讓我滿心快樂。 \end{tabularx} \\ \\ \relax
1:5 & \begin{tabularx}{0.7\textwidth}{X} 我記得你無偽的信心，這信心先存在你外祖母羅以和你母親友妮基的心裡，我深信也存在你的心裡。 \end{tabularx} \\ \\
[1ex]
\hline
\hline
\end{longtable}
$^{1}$各位弟兄姊妹平安.
我是第一次在佳欣樓跟大家見面.
你好 我是Odi.
我平時為什麼會沒有來這邊.
因為我11點至12點15分.
在錦宣那邊就是兒校的開始.
就撞了大家的時間.
今天很開心可以跟大家一起分享訊息.
我們一起有個祈禱吧.
差異的主多謝你.
我們很渴慕你的話語.
願你的聖靈來我們心裡.
讓我們得以更生得以滿足.
多謝你 奉耶穌名求 阿們.
好 想問大家一個問題.
大家可以想想.
最影響你生命的一個人.
究竟是誰呢.
媽媽 是.
還有沒有其他答案.
想到嗎.
父親 是的.
爸爸媽媽.
在預這篇講章的時候.
為了我的例子可以客觀一點.
所以我也做了一些功課.
最近我在教會以外.
認識了一班音樂創作的年青基督徒.
我跟他們討論過這些問題.
究竟他們的家庭和信仰有沒有什麼關係呢.
這班人來自不同的教會.
年齡大約23至28歲初職.
剛好我跟他們聊天的時候.
發現原來他們全部都是信義代.
即是父母在教會裡.
其實可能是牧者 傳道人.
或者一些資深的會友.
我跟他們討論完之後.
大家都認同有一個結論.
就是大概我們都不知不覺.

$^{41}$最影響我們的.
原來是我們的家庭和我們的父母.
最近我正在上一個心理分析輔導的課程.
在課程裡我了解到一個大人.
無論大至已經退休60多歲.
還是40多30多歲.
都要面對自己的生命問題.
歸根究底要回到自己的童年.
你的照顧者對你的影響是很深遠的.
今天我們一起來看一位年青人的生命.
提摩泰.
在保羅13封的書信裡.
有4封是有特定修信人的.
有2封原來是寫給提摩泰的.
即是提摩泰前後書.
而提摩泰後書更是保羅臨終前最後一封的書信.
我今天不會像平時的《釋經講道》那樣.
我們今天會穿插在不同的書卷.
看看提摩泰和保羅以及提摩泰的家庭.
對提摩泰的影響.
讓我們反思一下我們應該怎樣栽培我們的下一代.
保羅第二次宣教旅程裡.
大概是主後50至52年.
他再遇到提摩泰.
所以他第一次宣教旅程時.
已經把福音傳給提摩泰.
他第二次宣教旅程就再遇到提摩泰.
《使徒行傳》第16章1-5節這樣說.
後來保羅來到特比.
又來到路斯德.
在那裡有一個門徒名叫提摩泰.
是信主猶太婦人的兒子.
他的父親是希臘人.
路斯德和以哥廉的弟兄都稱讚他.
保羅要帶他同去.
只因那地方的猶太人都知道他父親是希臘人.
就給他行了國禮.
大家都知道.
猶太人其實熟讀摩西五經.
而提摩泰是信主猶太婦人的兒子.

$^{81}$聖經學者認為.
保羅是第一次宣教旅程裡經過路斯德.
而帶提摩泰一家去信耶穌.
信主之前的信仰基礎.
就是來自他的外祖母羅爾和母親尤尼基.
剛才我們也讀過這段經文.
在提摩泰後書1-5節.
保羅這樣寫.
我記得你無偽的信心.
這信心先存在你外祖母羅爾和你母親尤尼基的心裡.
我深信也存在你的心裡.
由此可見外祖母羅爾和母親尤尼基.
是用摩西五經去教導提摩泰.
以致他認識創天造地的耶和華.
相信神的應許.
更加容易明白彌賽亞的拯救.
因而相信耶穌.
當保羅傳福音給他的時候.
就很容易相信了.
提摩泰後書3章14-16節說.
至於你要持守所學習的和確信的.
因為你知道是跟誰學的.
並且你知道你從小明白聖經.
這聖經能使你因在基督耶穌裡的信.
有得救的智慧.
外祖母羅爾將所信的傳給母親尤尼基.
再傳給提摩泰.
因此其實培育敬虔的後代.
需要用神的話語去教導.
這裡提摩泰從小就熟讀聖經.
不單是教導小朋友.
而是首先家長要相信.
你要將那份信心傳給下一代.
你要自己先有那份信仰.
所以第一點要和大家分享的是.
表裡如一的榜樣.
當我訪問那一批年輕人信義代的時候.
他們分享了一個共通點.
他們很欣賞父母的身教.
有幾位父親都是很有原則的嚴父.

$^{121}$而媽媽是溫柔的.
不一定每個家庭都是這樣.
但剛好就是這樣.
其中一位說他在父親身上學到原來有些原則.
有些事情是不能做的.
而他爸爸教導他的事情.
他爸爸自己是真的做得到.
例如爸爸教導他要截劍.
他爸爸平時是不會坐的士的.
他有金錢管理金錢官身教.
而他在母親身上學到母親的同理心.
而他們的教導行事為人.
他的父母都是和他所教導的.
和他自己的人是相稱的.
他們也有分享到一些反面的.
有一個姐妹分享.
她有相反的經歷.
她對父母很不滿.
因為她父母是人前人後兩個人.
在別人面前她的父母是很好的.
但在家裡是另一個樣子.
而她更加失望的是.
因為她的父母是牧者.
因此她覺得牧者也不是很神聖.
也是人來的.
她有些失望.
有一位其實是青少年的導師.
他分享他自己在教會裡牧養青少年.
但那些青少年.
其中一位是父母在家裡打開聖經的時候.
是一個樣子.
合上聖經的時候又變成另一個人.
因此提出只是做一些禮儀上的事情.
可能父母也會做齊.
和他一起讀經祈禱.
但只是做禮儀上的事情.
父母內裡沒有改變.
沒有用新教去教導.
這樣的教導對孩子來說.
可能是一個更加難的難阻.

$^{161}$或者是半腳石.
因此父母要先做好自己.
先做好自己的時候已經是一個很好的榜樣.
很好的教導.
另一位也是青少年導師的分享.
他說青少年人很想回教會.
但他的父母早一個半月已經不讓他回教會.
因為要溫習.
很重要的要溫習.
但如果我們教導孩子.
每個星期都要到神面前.
相信安息日的重要性的時候.
讓我們不要成為阻擋孩子到神面前的人.
讓我們好好教導孩子分配時間溫習.
而不是剝奪孩子安息日到神面前的時間.
在這裡也想和大家分享.
我們可以分享自己的軟弱.
因為我們不是完美的.
我們和他們聊天的時候發現了一件事.
我自己也是信三代.
我也有發現.
信二代門也有發現父母有形象崩壞的一天.
即父母不是完美.
因為父母也是人.
因此我們更加需要和孩子分享自己的軟弱.
和讓孩子一同去經歷.
為什麼這樣說呢?.
有一位姐妹分享我覺得很好.
她很深刻.
因為她的導師在中學時期.
她還是一個學生的時候.
她的教會導師和她分享她的小產.
竟然導師在中學團體中和她分享她的軟弱.
讓這班學生一起經歷和為她祈禱.
她很深刻.
因為這個經歷讓她知道.
原來當她遇到困難的時候.
她也可以這樣去祈禱.
後來她見證著這位導師.
由小產到有一個小朋友.

$^{201}$再有另一個小朋友.
一共有兩個小朋友.
而她自己就親自成為這兩個小朋友的導師.
所以這個經歷影響了她的生命.
各位弟兄姊妹.
我們在別人面前分享成功好像很容易.
但原來分享我們的軟弱.
我們的軟弱如何變成剛強.
原來對別人更加有益.
保羅自己也說.
他說他有一根刺架在我身上.
就是撒旦的差役來折磨我.
免得我過於抬高自己.
為了這事我曾三次求主讓這根刺離開我.
他對我說.
我的恩典夠你用的.
因為我的能力.
視在人的軟弱上顯得完全.
所以我更喜歡誇自己的軟弱.
好使基督的能力浮遍我.
為基督的緣故.
我以軟弱,凌辱,艱難,迫害,困苦為學習樂的事.
因為我什麼時候軟弱.
什麼時候就剛強了.
保羅也跟別人分享他生命的那一根刺.
他不怕讓人知道他軟弱.
因為他何時軟弱.
何時就剛強.
同時在主的內.
當我們在主面前展現我們的軟弱的時候.
神的愛就像一個覆蓋.
覆蓋著我們.
讓我們顯現祂的大能在當中.
有一位信義代發現父親形象崩壞的那一天.
因為他以前覺得爸爸是正確的.
但他發現原來爸爸也有軟弱.
他覺得失望.
原來爸爸也有做錯事的時候.
他分享很感恩他的信仰很堅固.
當他知道爸爸的軟弱的時候.

$^{241}$他更加可以抓緊神.
因為他知道人原來是脆弱.
每一個人都有他的軟弱.
所以爸爸的軟弱令他的信仰沒有瓦解.
因為他可以更加抓緊神.
但避免孩子因為信仰,因為人而瓦解的話.
我們要做的就是將他們帶到神的面前.
更加需要讓孩子可以參與自己的軟弱.
就好像那位導師.
他和孩子一起祈禱.
他們和你一起去經歷神的大能.
一起去經歷你如何從軟弱變成剛強的時候.
這是他們生命的經歷.
成為他們的經歷.
第二點想和大家分享的是.
教會會一起培育孩子.
除了外祖母羅爾和母親尤尼基之外.
保羅也是提摩泰的屬靈父親.
第一次宣教旅程之後.
提摩泰信了耶穌.
五年之後保羅第二次宣教旅程.
再次遇到提摩泰.
覺得他已經夠成熟了.
因為路斯德和以哥念的弟兄都去稱讚他.
說他這個年輕人真的很好.
保羅決定要帶他出去宣教.
他更加在書信上稱提摩泰為他的真兒子.
無論是提摩泰前書還是後書.
提摩泰前書的紅色字說.
「因信主作我真兒子的提摩泰」.
他這樣形容提摩泰.
提摩泰後書說.
「親愛的兒子提摩泰」.
即是當他是兒子一樣.
保羅書信上稱提摩泰為自己的真兒子.
意思是屬靈的兒子.
因為是他帶提摩泰信主.
後來因為提摩泰的父親是希臘人.
學者猜測提摩泰應該是爸爸身上學到很流利的希臘文.
因此他可以幫助保羅寫信.

$^{281}$原來有四封書信.
哥林多後書,菲勒比書,歌羅西書和費利門書.
都是提摩泰和保羅聯名寫的.
我們看一個例子.
哥林多後書.
「奉上帝的旨意,作基督耶穌使徒的保羅和他的弟兄提摩泰」.
他寫這封書信.
「保羅認提摩泰屬靈兒子,除了是兒子之外,更是他親密的童工」.
「在保羅眼睛不好,軟弱時,他過一根刺,提摩泰成為他的幫助」.
保羅也在信中關心提摩泰的身體.
鼓勵和支持他.
在提摩泰前書523節.
「為了你的胃,又患病,不要只學水,要學眉,學點酒」.
他很關心他的身體.
還有我們經常聽到的一句.
「不可叫人小看你年輕,要在言語,行為,愛心,信心,清潔上作信徒的榜樣」.
提摩泰要堅持所信的事和所做的事.
要相稱,要作信徒的榜樣.
之後他更加鼓勵他.
因為他要和提摩泰分開.
在第一章7節說.
「因為上帝賜給我們的不是膽怯的心,是剛強仁愛自制的心」.
他鼓勵提摩泰.
我們未必有一個和我們出生入死的屬靈父親.
好像提摩泰.
但在教會裡,我們有很多屬靈長輩.
今天早上在祖堂崇拜.
感受祖堂崇拜.
其實是一個家庭崇拜.
同時也是一個孩童奉獻禮.
所以我今天穿了這件兒校T恤和大家見面.
孩童奉獻禮在乎的不是家長送小朋友去教會.
而是父母有立志成為榜樣的心.
這是家長的心,是很重要的.
用神的話語去教導孩子.
同時教會也承諾一起去培育小朋友.
為這些家庭去守望.
前一陣子有一位家長.
因為懷孕了,他很苦惱.
因為他已經有兩個小朋友.

$^{321}$也好像不太能應付.
沒有能力再多撫養一個.
他很想放棄這個小朋友.
當他真的要放棄這個小朋友那一天.
他很混亂,他來找我.
我為他祈禱.
在祈禱裡,我很心痛.
我也感受到聖靈與我們同在.
所以我求主賜給這位家長.
賜給他信心,也憐憫這個家庭.
最後我也告訴他.
其實我們是為你開心.
因為你有這個小朋友.
因為孩子是神賜給我們的產業.
他不是自己一個人撫養這個小朋友.
是教會的弟兄姊妹一起看著這個小朋友成長.
後來神改變了這位家長.
上個月他跟孩子一起受了咒.
是的,他拯救了這個小朋友.
你的小朋友在教會成長.
其實是大家一起去培育這個小朋友.
我自己是一個信三代.
跟提摩太一樣,婆婆已經信主.
我媽媽也信主.
我在教會成長裡.
生命裡也有很多的屬靈母親.
因為教會有很多不同的姨姨.
在不同的時期,有不同需要的時候.
他們就會出現幫助我.
有一位信二代的分享跟我很相似.
他說青春期的時候不喜歡聽父母說話.
可能我們很乖的小朋友也是.
青春期不知道為什麼會這樣.
但是他說姨姨們他們無條件的關懷.
讓他感受到那份愛.
所以姨姨叔叔的一句話.
不知道為什麼會勝過父母的一百句話.
在教會的群體裡.
我們可以有很多父母的模特兒.
因為我們的父母其實也很有限.

$^{361}$但是教會有很多姨姨叔叔.
有時候他們的身上可以補足到我們父母的不足.
我的姨媽是我舞會裡的長老.
有一次的短宣裡.
很深刻的是我跟她同房.
晚上她教我跪在床邊祈禱.
拿著一本叫做「祈禱的手冊」在那裡祈禱.
很深刻.
我也吸收了這個經歷.
成為了我現在也會教小朋友.
我們要跪在那裡祈禱.
另外,我媽媽是一個很不喜歡擁抱的人.
我需要擁抱.
有一個姨姨知道這件事的時候.
她自己也很喜歡擁抱.
所以她每次看到我都會抱著我.
雖然好像只是一件小事.
但是對當時的我來說.
青春期的我來說.
一個擁抱其實拯救了我的世界.
甚至是我讀神學的時候.
他們也會記得我.
約我聊天,為我祈禱.
他們會想念我.
間中我也會想念他們的湯水.
各位弟兄姊妹.
教會的群體.
我們跟弟兄姊妹的關係.
其實影響著我們的小朋友.
他們會看到我們怎樣相處.
從中去學習.
甚至他們在彭培之間.
他們會聽對方分享.
他們有沒有類似的經歷.
他們會看到對方怎樣去經歷這些困難的事情.
例如升學,找工作.
他們看到原來有導師,牧者一起同行的時候.
他們就會知道原來我自己.
將來我不用怕.
我自己可以這樣去經過.

$^{401}$第三點想跟大家分享的是.
屬靈的婆婆,公公,長輩.
我想問問這裡有沒有.
已經成為了公公,婆婆,爺爺,奶奶的.
舉手來看看.
咦,也有.
今天我們的主題其實是.
代代相傳的信仰.
孩童奉獻禮.
在乎父母奉獻的心.
但這個心是延續的.
當我們由父母.
榮升為祖父母,外祖父母.
不要覺得自己沒有用.
很多老人家都覺得已經沒有用了.
我們不要這樣想.
因為你依然是一個好的榜樣.
你也要繼續用神的話語去提點你的兒女.
就算你的兒女已經成家立室.
已經有孫子,孫女.
甚至你已經退休了.
三年前我進來感孫的時候.
在兒校有一半的兒童導師都離開了.
而我又不是太認識教會的人.
於是Aka傳道.
就呼籲不同的肢體來幫忙.
當時有幾位.
奶奶,婆婆級的.
他們幾十年前.
已經是忠心在兒童部服務的姐妹.
他們回來幫忙.
他們說:我都不懂得用電腦.
於是他們就預備了自己最熟悉的方法.
就是用大字簿.
有沒有人知道大字簿是什麼呢?.
舉手來看看.
你小時候也是這樣嗎?.
有很多大大張的海報.
就用手寫.
以前沒有這麼幸福.

$^{441}$有這麼好的電視,螢幕.
他們真的很用心地預備大字簿.
教小朋友唱歌.
我很感動.
神會用我們.
只要我們有願意的心.
他們不停地幫忙.
幫了兩個月.
他們是很疲勞的.
因為他們已經退休了.
慢慢地.
我就看到有不同的弟兄姊妹.
願意加入來幫忙.
在這個導師感恩分享會裡.
奶奶,婆婆級的導師們.
他們感動到流淚.
當他們分享的時候.
很感動.
因為他們覺得神依然用自己.
而且小朋友真的很可愛.
下一輩的導師.
就是爸爸,媽媽級.
他們看到自己的.
原來他們的小時候主一學老師.
奶奶,婆婆級的那幾位都回來.
他們都覺得自己.
還是要回來兒校幫忙.
覺得自己要回應神.
出來事奉.
所以我們不要小看自己的年長.
你要忠心.
你的付出.
依然會成為大家的好榜樣.
所以請大家紀念兒校.
最後鼓勵大家.
不要因為怕環境而不去生育.
因為生養眾多.
和培育敬虔的後代.
是我們要做的事情.
神是信實的.

$^{481}$祂是超越一切環境的神.
提摩太當時其實正在經歷尼祿王.
要捉基督徒.
甚至保羅是在監獄裡面殉道.
他們的環境比我們現在更差.
所以神的恩典是會伴隨著愛祂的人.
《申命記》這樣說.
我們一起讀.
1,2,3.
你知道耶和華,你的神,他是神,是信實的神.
他向愛他,守他誡命的人.
守約,思慈愛,直到千代.
神一直都很想祝福我們.
只要我們遵循祂的話語.
愛祂,守祂的誡命.
今天我在這裡都很想奉耶穌的名去呼召.
如果今天神感動你.
你願意成為教會孩子的好榜樣.
一起牧養,培育教會的下一代成長.
請你在神面前舉手.
很有信心的舉手.
大家先.
求主加力.
讓我們都有信心.
讓我們一起祈禱.
差的天父,多謝你.
今天我們在提摩太的家庭裡.
我們都學習到.
我們身為父母的,長輩的.
神體重我們奉獻的心.
我們剛剛有舉手的,你都看見.
或者沒有舉手的,你都知道.
讓我們盡力去學習成為孩子表裡一如一的好榜樣.
不怕讓孩子參與自己的軟弱.
讓他們一同去經歷神的大能.
我們在神的教會裡.
也要成為孩子們的屬靈的父母.
在教會當中活出美好的團契.
求主祝福感孫家.
家恩樓的孩子們.

$^{521}$在主的愛裡成長.
被神的愛包圍.
奉主耶穌基督的名求.
阿們.
\newpage



\section{使徒行傳 11:19-30}
\label{sec:nQU0VbnDk9Y}
\textbf{2023-11-26 午堂《要作榮耀主的教會》使徒行傳 11\_19-30 講員:趙家碩牧師}
\newline
\newline
連結: \href{https://youtube.com/watch?v=nQU0VbnDk9Y}{\texttt{https://youtube.com/watch?v=nQU0VbnDk9Y}} ~~~~ 語音日期: 2023-11-28
\newline
\newline
\hyperref[sec:Q0yYJrC6Zoc]{\small{< < < PREV SERMON < < <}}
~
\hyperref[sec:index_chronic]{\small{[返順時目]}}
~
\hyperref[sec:index_scriptual]{\small{[返順卷目]}}
~
\hyperref[sec:gy4cK2fsHdc]{\small{> > > NEXT SERMON > > >}}
\newline
\newline
使徒行傳 11:19-30
\newline
\begin{longtable}{cl}
\hline
\hline
章節 & 經文 (和合本修訂版)\\
\hline
11:19 & \begin{tabularx}{0.7\textwidth}{X} 那些因司提反的事遭患難而四處分散的門徒，直走到腓尼基、塞浦路斯和安提阿。他們不向別人講道，只向猶太人講。 \end{tabularx} \\ \\ \relax
11:20 & \begin{tabularx}{0.7\textwidth}{X} 但內中有塞浦路斯和古利奈人，他們到了安提阿也向希臘人傳講主耶穌的福音。 \end{tabularx} \\ \\ \relax
11:21 & \begin{tabularx}{0.7\textwidth}{X} 主的手與他們同在，信而歸主的人數很多。 \end{tabularx} \\ \\ \relax
11:22 & \begin{tabularx}{0.7\textwidth}{X} 這風聲傳到耶路撒冷教會的人耳中，他們就打發巴拿巴到安提阿去。 \end{tabularx} \\ \\ \relax
11:23 & \begin{tabularx}{0.7\textwidth}{X} 他到了那裡，看見神所賜的恩就歡喜，勸勉眾人要立定心志，恆久靠主。 \end{tabularx} \\ \\ \relax
11:24 & \begin{tabularx}{0.7\textwidth}{X} 這巴拿巴原是個好人，滿有聖靈和信心，於是有許多人歸服了主。 \end{tabularx} \\ \\ \relax
11:25 & \begin{tabularx}{0.7\textwidth}{X} 他又往大數去找掃羅， \end{tabularx} \\ \\ \relax
11:26 & \begin{tabularx}{0.7\textwidth}{X} 找著了，就帶他到安提阿去。他們足有一年和教會一同聚集，教導了許多人。門徒稱為「基督徒」是從安提阿開始的。 \end{tabularx} \\ \\ \relax
11:27 & \begin{tabularx}{0.7\textwidth}{X} 當那些日子，有幾位先知從耶路撒冷下到安提阿。 \end{tabularx} \\ \\ \relax
11:28 & \begin{tabularx}{0.7\textwidth}{X} 內中有一位，名叫亞迦布，站起來，藉著聖靈指示普天下將有大饑荒；這事在克勞第年間果然實現了。 \end{tabularx} \\ \\ \relax
11:29 & \begin{tabularx}{0.7\textwidth}{X} 於是門徒決定，照各人的力量捐錢，送去供給住在猶太的弟兄。 \end{tabularx} \\ \\ \relax
11:30 & \begin{tabularx}{0.7\textwidth}{X} 他們就這樣做了，託巴拿巴和掃羅的手送到眾長老那裡。 \end{tabularx} \\ \\
[1ex]
\hline
\hline
\end{longtable}
$^{1}$大家好 很感恩有機會跟大家分享.
亦代表大埔堂向大家問安.
大家有沒有去過大埔 都有的 其實我們很近.
今天我穿這件風褸 其實今年是大埔堂60週年堂慶的時間.
大埔堂60年了 不知道錦繡堂42.
我戴了真木絲一點 雖然頭髮我想大家都相當白.
在這裡能夠可以一起去敬拜.
剛才我們看到有不同國籍的裡面.
猜拳主日的裡面 體會到我們在地若天萬國萬民的去敬拜神.
是多麼美好的一個景象.
不過可能大家會覺得 今天我們在香港這個地方.
經歷了好幾年 無論疫情 全球經濟的下滑的時候.
香港又好像很多人移民.
究竟今天神擺我們教會在這個地上.
特別在香港這個地方 面對著這些困境的時候.
我們有什麼應該要做的呢.
在《士同行傳》第一章八節裡面.
原來初期教會給我們看到的就是.
當時他們仍然面對著很多很多不同的衝擊.
但耶穌在升天之前就跟他們說好.
他耶穌基督的計劃是怎樣呢.
但聖靈降臨在你們身上 你們就必得著能力.
並要在耶路撒冷猶太傳地和撒瑪利亞直到地極作我的見證.
原來這個福音的故事 弟兄姊妹.
是要由我們每一個繼續延續下去.
我們今天就透過看安提阿的教會.
讓我們看看對我們有什麼激勵 有什麼提醒.
先介紹一個宣教士給大家去認識.
這個宣教士叫賈錫德醫生.
紅色圈那個 有沒有見過這幅圖畫.
這幅圖畫其實是我們宣教士歷代裡面.
宣道會歷代以來來到中國宣教的宣教士.
中間裡面圈著的賈錫德醫生.
他未曾踏足中國.
不過他卻稱為宣道會附和的第一人.
為什麼會這樣呢.
賈錫德醫生本身是加拿大人.
我是不是應該要說慢一點呢 讓他翻譯好一點.
他很年輕的時候就有這個宣教的心智.
1887年的時候 也是宣道會成立的時間.

$^{41}$他回應著神的呼召.
所以兩夫婦去讀神學.
然後就加入了宣道會.
成為了首批來到中國的宣教士.
不過這個賈醫生的太太因為要繼續他的學業.
所以就先留在美國的地方.
賈醫生就說自己買船票先去中國.
心裡面想著打點好一切.
然後才直接接家人來到中國的地方.
賈錫德醫生和他的太太.
臨離別之前說了一句這樣的說話.
他說我已經準備好為主而死.
我已經準備好為主作供.
不過誰知這一句說話.
就成為了賈醫生的道別說話.
當時要由美國一路坐船去到中國的地方裡面.
要經過大西洋太平洋.
他們要先去到日本.
才轉船去到中國.
這麼長的行程裡面.
賈醫生特意買那些下等艙的船票.
為什麼呢.
因為下等艙的船票就能夠有機會和那些華人可以共處.
學習他們的語言 傳福音給他們聽.
不過在這個船程裡面.
不幸地賈醫生患上了天花.
所以幾個月之後.
他終於到達日本了.
不過也在日本環濱這個地方病逝.
我們中國人有一句.
看看大家是否認識.
出師未捷身先死.
我不知道翻譯的時候怎樣可以翻得到.
這位賈醫生一心想著來到中國宣教.
不過他未曾踏足中國這片土地的時候.
就已經死了.
為什麼呢.
究竟有什麼意義呢.
其實神為什麼不讓他先來到中國.
起碼做一點事 是不是好像有價值一點呢.

$^{81}$老實說 我也不懂得回答.
不過我深信這位賈醫生.
他沒有後悔.
而他的犧牲也不是白白的.
因為他的故事激勵了很多宣教士.
也讓很多人因此而相信耶穌.
也在他去世一年之後.
賈醫生的遺孀和他兩個女兒.
加上一些宣教士.
終於再由美國的地方輾轉來到中國.
建立了宣道會在中國裡面首個宣教站.
就讓我們看到.
其實弟兄姊妹今天.
或者我們都會面對著很多環境的轉變.
很多的艱難.
不過我們要知道不是去問為什麼.
而是要問我們怎樣可以繼續為主向前去走.
今天我們看安提柯教會和這段經文.
就讓我們看到原來他們仍然面對著很艱難的環境.
但是神讓他們怎樣建立成長以至轉化.
經文裡面大概可以分三個段落.
19至21節.
我們看到他們怎樣從北向去到亞向.
成為了一個福音廣傳的地方.
22至26節就看到他們怎樣繼續在神的恩典裡面成長.
27至30節就是神怎樣轉化他們.
以至他們成為了一個普世宣教的基地.
首先我們要的就是突破我們自己的框框.
從北向去到亞向.
我們要說安提柯的教會.
先說少少地理上的.
看到安提柯在哪個位置嗎.
敘利亞的安提柯.
其實他就在耶路撒冷以北大概500公里的地方.
都頗遠的.
是一個沿海的城市.
當時羅馬帝國底下一個相當重要的城市.
很多很多外邦人在這裡聚集.
而經文裡面11章19節.
一開始就寫了.

$^{121}$那些因司提反的事遭患難四散的門徒.
其實就告訴我們.
因為在使徒行傳第八章.
司提反因著神的緣故去殉道.
當時這班門徒大遭逼迫而四散.
不過他們四散的時候不只是逃跑.
他們去到撒瑪利亞巴勒斯坦的一帶.
不過他們隨走隨存.
亦都有些門徒走到更遠的地方.
更北的地方.
即是塞浦路斯.
菲利基這一帶的地方.
一直走到安提柯.
很遠很遠.
一個外邦人的城市.
他們很忠心的.
不只是逃難.
願意的隨走隨存.
不過在經文裡面19節的下.
就提到他們有一個這樣的情況.
經文裡面說.
他們不向別人講道.
只向什麼人講道.
只向猶太人講道.
這班猶太人他們信了主.
雖然要四散逃難的時候.
但他們仍然拿著福音的使命.
不過他們當時這些猶太人的移民.
他好像我們中國人一樣.
如果我們選擇要移民的話.
通常會選擇什麼城市.
我想這段時間最多人移民的.
就應該是英國.
應該我們都會選擇一些大城市.
無他的.
因為多一些同鄉.
亦都可能容易些謀生.
當時這個安提拉.
是羅馬帝國底下第三大的城市.
交通很發達.

$^{161}$亦都商業很繁盛.
所以當時很多的猶太人的移民.
都會來到這個城市.
要知道這班猶太人.
他們總有一個自己的框框.
覺得自己的民族比較優越.
所以他們初到貴境的時候.
不向其他人去傳福音.
只是向他們自己的同胞傳福音.
你覺得合理嗎.
其實都合理的.
試想想.
如果要你去和其他民族去傳福音.
用你不是你自己的語言.
不是你自己的文化去傳福音難不難.
其實有一定的難度.
就這樣說.
我估計可能你們容易些.
要你們用英文來講福音.
可以嗎.
都沒有人會點頭的.
我想英文都可以.
印尼文講福音可以嗎.
你們當然可以.
我們就不可以.
我只會講Salamat Bagi.
有時很難的.
不止難.
有時是我們不肯.
但經文第20節就和我們說了.
內宗有一些賽普路斯和古里內人.
這班信徒來到安提阿的時候.
也向希臘人去傳講主耶穌.
由於當時他們是在羅馬帝國統治下.
希臘文是他們普遍會用的國際語言.
就好像我們今天用英文.
因著這班人他們成長很特別的背景.
他們是來自賽普路斯.
來自一些本身已經在外地成長的.
猶太的信徒.

$^{201}$所以他們很大機會去到耶路撒冷.
認識了耶穌 相信了耶穌之後.
當他們逃難的時候.
就回到他們本身熟悉的外邦人的城市.
來到安提阿.
這些賽普路斯成長的人.
就好像我們今天那些叫ABC或BBC.
大家知道是什麼吧.
就是在外國成長那些叫黃皮白心的人.
他們就因著能夠粗一口流利的希臘文.
也都能夠有國際性的視野.
以致他們就突破自己那種民族優越性的框框.
願意向其他的種族.
願意向那班安提阿的人去傳講福音.
這裡讓我們看到什麼呢.
這班的信徒因著迫不及待困難.
他們要四散.
不過這些限制沒有讓他們停止了福音的工作.
他們民族的優越性.
同樣有人不會因為這些環境.
或者自己內心裡面的一些限制.
而局限了福音的工作.
局限了他們自己.
他們願意突破自己.
以致他們由不向願意去到也向其他人去傳福音的門徒.
其實我們要突破我們自己難不難呢.
就好像我們上崇拜那樣.
大家是不是通常上五堂崇拜的.
這裡是禮堂的照片.
大家是不是通常坐的位置都是附近的位置.
有沒有試過回來坐的時候.
弟兄姊妹 這個位置是我的.
大家知道這個位置不是你的.
不過我們坐慣了.
總是會覺得好像別人坐了我的位置.
又或者我們完了崇拜的時候.
通常都會那幾位熟的弟兄姊妹.
聊聊天 去哪裡吃飯.
當然有些習慣其實無傷大雅.
不過會不會作為一個信徒.

$^{241}$我們信主耐了.
真的有一些固有的想法.
有一些限制了我們.
要阻攔了我們去實踐福音使命的情況.
弟兄姊妹 你願不願意放下呢.
你願不願意看到神要挑戰你.
突破你自己這些固有想法的一些地方.
以至為主去行多一步呢.
沒錯 今天的確在香港的教會裡面.
有很多很多不同的弟兄姊妹移民到外國.
又或者你們心裡面都會很不捨得.
有一些弟兄姊妹去到其他地方.
我們為她祈禱.
同時盼望她們帶著福音的使命.
去到當地裡面 繼續實踐使命.
弟兄姊妹 今天我們留在這裡的地方.
我們不只是懷緬過去.
我們過往過往多好多好.
而停滯不前.
今天我們留在這個地方.
留在香港.
其實有沒有看到原來神已經將很多很多.
不同民族的地方.
不同民族的人 由遠方帶到我們身邊.
甚至帶到我們家裡面.
我為了感受堂去感恩.
我們有菲律賓的團體.
有印尼的群體.
還有非洲來的姊妹在裡面.
其實在香港有幾十萬的非傭人傭.
我們想 教會正在做.
不關我事.
不是的 弟兄姊妹.
都關我們事.
神將他們擺在我們身邊.
你願不願意關心他們呢.
今天不單止有不同語言的群體.
今天香港的政策帶了很多很多.
我們叫做新移民 新香港人來到裡面.
究竟我們會不會願意接納他們呢.

$^{281}$還是他們非我族女.
弟兄姊妹 沒錯.
有時要去傳福音給不同的群體是難的.
不過難不在於語言上.
難在於我們心裡面.
有沒有關切到不同群體的心.
當然我們也向的意思.
不是只顧著其他人.
我們自己本身這個群體.
我們自己本身身邊的朋友.
都是我們需要關心的.
不過聖經提醒我們.
不要忽略到這些其他群體.
又或者很容易被人歧視.
這些群體的需要.
以致我們願意站起來.
按照我的教會.
他們優勢不在於他們在一個很繁榮的都市裡面.
而是在於他們有一個很有包容性.
很接納程度很高的群體.
這班信徒所建立的教會.
所以當他們願意突破自己的限制.
由很內在的不幸.
以致成為一個願意外展也向外邦人傳福音的教會.
神就自然祝福他們.
神就與他們同在.
以致叫信主歸主的人.
每一天都加珍給他們.
我深願錦繡堂同樣弟兄姊妹.
你都願意突破你們自己.
當我們願意突破的時候.
神又會讓其他的人.
看到怎樣從門徒去到基督徒.
經文裡面第22節裡面去到第二個段落.
講及到這班門徒他們願意突破他們自己的界限.
當時很多很多安提拉的人願意去加入教會.
相信耶穌發展得很快.
以致風聲就傳到耶路撒冷教會.
好像舞會一樣.
他們就猜派著教會的代表去到當中.

$^{321}$看一看他們的環境.
了解一下他們的情況.
不過他們就不是派使徒約翰或使徒彼得.
派了誰去呢.
派了巴拿巴去.
很大原因派巴拿巴去.
就是巴拿巴本身是塞浦路斯人.
在安提拉旁邊大島裡面成長的人.
當時巴拿巴去到那裡的時候.
就見到感恩果然福音傳得很快.
這個教會也發展得很快.
不過你要知道突然間一間教會多了這麼多人.
混不混亂.
都會有很多混亂有很多事情需要處理.
巴拿巴來到這個地方.
他不是在裡面指指點點.
而是見到他們這樣的情況.
願意留下.
立定心智.
放上自己留下去做好他們.
勸勉他們栽培他們的工作.
建立這班初信者.
也幫助他們學習怎樣幸狗的去靠主.
經文裡面除了去形容巴拿巴很忠心服事之外.
都特別讚賞他的品格.
說他是一個怎樣的人.
說他是一個好人.
大家是不是好人.
不承認的 大家不是好人.
我們都是好人.
聖經同樣去形容巴拿巴是被聖靈充滿大有信心.
沒錯 巴拿巴是一個很謙柔.
盡心去造就別人的一個領袖.
而他的品格就影響整個教會.
巴拿巴他的名字其實叫做勸慰之子的意思.
因為他其實來到耶路撒冷教會的時候.
他已經是第一批賣了自己的家當奉獻給教會.
要甘心樂意供應給教會的人.
當其時巴拿巴面對著人見人怕的素羅的時候.
他是特意引薦素羅給使徒去認識的.

$^{361}$聖經裡面去形容巴拿巴.
很明顯和保羅 即是素羅是很不同的領袖.
你不會見到聖經形容巴拿巴四處去講道.
或者和人去辯論.
因為巴拿巴他的恩賜是在於個人的栽培 牧養上.
巴拿巴去到安提阿的教會牧養了一段時間.
教會仍然發展得很好 發展得很快.
不過他就看到一班會眾的需要.
所以他就特意走了230公里那麼遠的地方.
就由安提阿特意去到大數那裡.
找那個本身很可怕的保羅過來.
和他一起去帶領著教會.
要知道巴拿巴在安提阿教會牧養了一段時間.
有一定的位置.
現在他要找另外一個人.
無論在教道 在講道上都比他強的一個人去配搭.
其實很難的 胸襟要很闊才做得到.
不過巴拿巴他心裡面的就是記掛著這一班弟兄姊妹的需要.
他也不只是要看自己的位置.
以至他能夠好好地去栽培弟兄姊妹之外.
也栽培保羅成為一個教會很重要的領袖.
兩個領袖一起去服侍.
中間一定有很多不同的困難.
二十六節裡面好像很輕描淡寫地寫.
他們足有一年的功夫去和教會一同聚集 教導教會.
就好像我們講的童話故事.
他們兩人結婚之後就快快樂樂地生活下去.
結了婚的是不是這樣的.
快快樂樂地生活下去.
是的 有時.
不過結婚久了的都知道.
中間過程裡面一定有不同的艱難.
有不同的事情發生.
只不過聖經裡面很簡單地去描述.
其實去到迦拉太書.
保羅在迦拉太書裡面提及過.
因著基伐 即是彼得和巴拿巴.
在一個和外邦人吃飯的場景裡面.
因為保羅見到巴拿巴有不適之處.
所以就當面去責備巴拿巴.

$^{401}$沒錯 保羅因著神學.
因著一個正確的觀念的情況.
他做出這個動作.
不過我們更加看到的就是.
巴拿巴他在這個事情裡面.
沒有說要去反駁.
沒有說要去爭論甚麼.
他仍然保持著謙卑的態度.
才會目仰著整間教會.
事實上 當我們回到教會的日子.
又或者我們會覺得.
看到教會有很多很多不同的問題.
不同的艱難.
不過有時問題的出現.
究竟是教會的問題.
還是我們自己胸襟的問題呢.
會不會有時.
其實心底裡面真的是我們自己小器的呢.
有時會覺得我提出的意見.
為甚麼沒有人理的.
為甚麼那些人不重視我的意見呢.
人與人之間可能在事奉的過程裡面.
有時出現隔膜.
他說你做得這麼好 你就做完吧.
我退下來 我不做了.
甚至你們用你們的腳表達你們的不滿.
原來教會的發展.
不是講求我們自己達到自己的目的.
達到自己的成就.
而是要將焦點放在神的工作裡面.
當我們每一個都願意放下自己.
盡心為主去作工.
去造就別人的時候.
神就自然祝福這間教會.
安提拉教會裡面.
我們看到的.
有一班很忠心願意培育他人的領袖.
亦都整間教會裡面願意放下自己.
以致成為榜樣.
以致神祝福他們能夠自然而然地成長.

$^{441}$為甚麼我這樣說.
因為很特別.
經文裡面第26節.
就在這裡裡面用了一個小結.
他說門徒稱為基督徒是由哪裡開始.
由安提拉開始.
今日我們信耶穌的人.
你會稱呼自己是甚麼.
基督徒好像很簡單的一個稱呼.
沒甚麼特別.
不過原來聖經裡面稱呼為門徒的.
一般都是信主的群體裡面.
我們彼此一個稱呼.
門徒.
而這個基督徒的名稱.
說由安提拉開始.
即是說由安提拉的教會.
他們這班信耶穌的人.
是首先由他們稱為基督徒.
為甚麼會是他們稱為基督徒.
其實正是一班外面不信耶穌的人.
不過見到這班教會裡面.
一班信耶穌的人.
他們很有見證.
他們很有愛心.
或者翻譯成廣東話.
就是耶穌仔 耶穌女.
表達著外面群體的人.
見到這班教會弟兄姊妹的見證.
一個很大的群體.
願意祝福社區的群體.
就稱呼他們為基督徒.
所以門徒稱為基督徒.
從安提拉起手.
雖然好像是一個.
今天我們很普通的稱呼.
但是卻是別人對他們認定.
認定他們見證.
認定他們榜樣的一個稱呼.
所以別具意義.

$^{481}$究竟弟兄姊妹.
今天錦繡堂.
存在在錦繡花園.
元朗區這個地方的時候.
我們的存在有沒有被其他人看到.
我們是基督徒.
我們的存在有沒有被其他人看到.
我們是祝福這個社區的呢.
我心願神繼續帶領錦繡堂.
為主的作美好的見證.
這樣我們就要做到.
從一間蒙恩的教會.
去到一間施恩的教會.
第三個段落裡面.
就說當那些日子.
即是巴拿巴和保羅.
在安提拉教會裡面.
牧養了一段時間之後.
當時就有幾位先知.
由耶穌撒冷來到安提拉這個地方.
聖靈就感動這幾位先知.
祂說指明天下.
即是當時巴勒斯坦地.
地中海一帶的地方.
將會有什麼事情發生.
有什麼發生.
有大饑荒發生.
按著歷史考古的顯示.
大概一年之後.
的確有饑荒發生在那個地方.
不過要留意.
當這幾位先知說有饑荒發生的時候.
饑荒發生了沒有.
沒有.
就算饑荒已經發生.
也不代表著.
這班人會動起他們憐憫的心.
因為這個饑荒發生在哪裡.
當時地中海一帶.
即是說除了耶路撒冷的教會.

$^{521}$會面對饑荒之外.
安提拉教會會不會面對饑荒.
同樣會面對.
不過安提拉教會就看到別人的需要.
以至他們有包容有關心的DNA.
以至他們願意為了對方的需要.
主動地去籌集一些捐款.
送到耶路撒冷的教會.
他們還派了兩個很重要的領袖.
將他們的心意派到耶路撒冷的教會.
讓我們看到的就是.
這間願意突破界限的教會.
連接著這班屬靈美好傳統品格的領袖.
他們有著福音看到別人的需要.
以至他們本身是剛剛發展著.
一間蒙恩的教會.
不過他今天已經轉化了.
成為一間私恩的教會.
成為了原本他只不過是一個宣教的工場.
但轉眼間就成為了一個有宣教心智.
成為了普世宣教基地的安提訛教會.
怎樣才能夠成為一個.
沒有恩典又願意私恩的教會呢.
其實在歷史裡面.
往往其實我們願意踏出來的時候.
神都讓我們看到不同的境況.
這些的境況可能是一些艱難.
一些困難.
卻喚醒我們心裡面的心智.
來到這裡.
小小的去分享.
在大埔堂的一些小小的歷史.
我剛才說大埔堂今年是多少歲.
60歲.
不過我們作為宣道會的堂會.
宣教都是我們很重要的一環.
亦都是我們必然要做的地方.
我們大埔堂起步猜全宣教上的事.
都不是很長的日子.
我們上年是猜全部30週年的紀念時間.

$^{561}$我們有一本的特刊.
我們在1992年的時候.
當時有幾個弟兄姊妹.
有同工 有一個心智.
幾位的 只有七位.
成立了一個叫猜關小組.
嘗試關心猜全的事情.
我們去到1994年.
才正式成立我們的猜全部.
亦都在那時候我們再去計劃思想.
我們猜全的工作是怎樣做.
當時在我們猜全部成立的時候.
我們就定了.
我們以一個猜全基金的方式.
去支援 去支持我們宣教的工作.
定了計劃.
每一年都有一個猜全基金的目標.
今年應該這樣說.
2024年我們猜全的基金.
即是說要求或者期望弟兄姊妹.
在十一奉獻 常費奉獻以外.
我們特別在猜全上面去奉獻.
我們2024年猜全奉獻的目標是200萬.
我們亦都在初頭完全不知道怎樣做猜全事工.
去到今天 過去的幾年裡面.
我們有12個宣教士的單位.
是我們所支持關心.
亦都有一些弟兄姊妹.
是由我們自己的堂會所猜出去的.
在這裡面想和弟兄姊妹分享的.
不是大報堂做得猜全的工作有多好.
我深信感受堂可能有更好的一些情況.
不過我就很體會到.
當我們願意為神擺上的時候.
神同時間又祝福不單止遠方的工作.
亦都祝福教會裡面不斷地去成長.
大報堂在過去的日子裡面.
我們不單單關心遠方的工作.
亦都關心本地的工作.
以致叫大報堂不斷地去成長.

$^{601}$我們起步的時間.
我們都有做印尼事工.
起步的時間比大家可能遲.
我們在2011年的時候.
才開始做印尼的事工.
亦都在當時在印尼裡面聘請了.
印尼的叫做童工.
來到自己的堂會裡面去牧羊.
過去的十年裡面.
其實我們印尼的事工的發展.
是最快或者叫做最好的.
初初只有幾個的姐姐來到.
聽得明白吧.
你認識我們印尼那邊的姐姐嗎.
認識.
有時有些見過.
Ransom那些就見過了.
由初初幾個的姐姐.
到今天我們上星期的印尼崇拜.
有50個姐姐實體聚會.
還有20多個姐姐在網上聚會.
因為姐姐不是一定星期日放假.
很體會到的就是.
神告訴我們.
當我們願意為神擺上.
不只是看自己有多少的艱難.
多少的缺陷.
而是當你為神擺上的時候.
神就有不同的方式去祝福我們.
我們經歷了好幾年.
在疫情裡面的困境.
其實不止香港.
全球的地方裡面我們都在艱難.
經濟沒錯是下滑.
我們又面對著移民的地方.
教會一定有影響.
但是為了弟兄姊妹你們感恩.
為了錦繡堂去感恩.
我們不只是停在這裡.
我們仍然要學習.

$^{641}$怎樣去關心社區.
做好不同民族不同社群的工作.
我們要做宣教的工作.
要做福音的工作.
我們要成為一個榮耀神的教會.
原來最重要的一環.
就是我們願意突破自己.
承認我們不足.
承認我們有困難.
卻是願意同時間.
甘心樂意為神擺上.
有沒有想過.
神盼望你成為一個怎樣的基督徒呢.
深信神對弟兄姊妹.
你每一個都有祂特別的心意.
我們能夠從耶路撒冷的教會.
這一幅是我們恩利姐姐.
和我們自己的敬拜隊.
一起去敬拜崇拜的時候.
一個圖片.
我想說的是.
神從耶路撒冷安提柯教會.
讓我們看到的.
原來基督徒必然.
必然會面對環境的轉變.
甚至有逼迫挑戰.
昔日的信徒.
他們大遭逼迫.
不過他們不只是著眼於.
自己的安全.
自己的生活.
而是願意為神走多一步.
願意將福音傳到外邦.
我們看到的是安提柯教會.
是一群包容性很高.
又謙和的一群信徒.
以致助紂別人.
幫助別人.
有廣闊的視野.
神就祝福他們成為一個普世宣教的基地.

$^{681}$今天.
我們縱然好像面對著內憂外患.
不過我們不至於說叫做大遭逼迫.
幸好我們總會面對很多的限制.
但不是我們灰心失意的時間.
而是我們嘗試在這個環境.
這個地方.
怎樣可以突破我們的思維.
讓我們活在這個地方.
別人看到我們有基督徒的生命.
有主的生命.
能夠榮耀主.
能夠榮耀神的教會.
願主祝福我們每一個.
\newpage



\section{馬太福音 2:1-12}
\label{sec:gy4cK2fsHdc}
\textbf{2023-12-02 晚堂《聖誕禮物》馬太福音 2\_1-12 講員:楊明麗姑娘}
\newline
\newline
連結: \href{https://youtube.com/watch?v=gy4cK2fsHdc}{\texttt{https://youtube.com/watch?v=gy4cK2fsHdc}} ~~~~ 語音日期: 2023-12-06
\newline
\newline
\hyperref[sec:nQU0VbnDk9Y]{\small{< < < PREV SERMON < < <}}
~
\hyperref[sec:index_chronic]{\small{[返順時目]}}
~
\hyperref[sec:index_scriptual]{\small{[返順卷目]}}
~
\hyperref[sec:7XxklgVuie4]{\small{> > > NEXT SERMON > > >}}
\newline
\newline
馬太福音 2:1-12
\newline
\begin{longtable}{cl}
\hline
\hline
章節 & 經文 (和合本修訂版)\\
\hline
2:1 & \begin{tabularx}{0.7\textwidth}{X} 在希律作王的時候，耶穌生在猶太的伯利恆。有幾個博學之士從東方來到耶路撒冷，說： \end{tabularx} \\ \\ \relax
2:2 & \begin{tabularx}{0.7\textwidth}{X} 「那生下來作猶太人之王的在哪裡？我們在東方看見他的星，特來拜他。」 \end{tabularx} \\ \\ \relax
2:3 & \begin{tabularx}{0.7\textwidth}{X} 希律王聽見了，就心裡不安；耶路撒冷全城的人也都不安。 \end{tabularx} \\ \\ \relax
2:4 & \begin{tabularx}{0.7\textwidth}{X} 他就召集了祭司長和民間的文士，問他們：「基督該生在哪裡？」 \end{tabularx} \\ \\ \relax
2:5 & \begin{tabularx}{0.7\textwidth}{X} 他們說：「在猶太的伯利恆。因為有先知記著： \end{tabularx} \\ \\ \relax
2:6 & \begin{tabularx}{0.7\textwidth}{X} 『猶大地的伯利恆啊，你在猶大諸城中並不是最小的；因為將來有一位統治者要從你那裡出來，牧養我以色列民。』」 \end{tabularx} \\ \\ \relax
2:7 & \begin{tabularx}{0.7\textwidth}{X} 於是，希律暗地裡召了博學之士來，查問那星是甚麼時候出現的， \end{tabularx} \\ \\ \relax
2:8 & \begin{tabularx}{0.7\textwidth}{X} 就派他們往伯利恆去，說：「你們去仔細尋訪那小孩子，找到了就來報信，我也好去拜他。」 \end{tabularx} \\ \\ \relax
2:9 & \begin{tabularx}{0.7\textwidth}{X} 他們聽了王的話就去了。忽然，在東方所看到的那顆星在前面引領他們，一直行到小孩子所在地方的上方就停住了。 \end{tabularx} \\ \\ \relax
2:10 & \begin{tabularx}{0.7\textwidth}{X} 他們看見那星，就非常歡喜； \end{tabularx} \\ \\ \relax
2:11 & \begin{tabularx}{0.7\textwidth}{X} 進了房子，看見小孩子和他母親馬利亞，就俯伏拜那小孩子，揭開寶盒，拿出黃金、乳香、沒藥，作為禮物獻給他。 \end{tabularx} \\ \\ \relax
2:12 & \begin{tabularx}{0.7\textwidth}{X} 因為在夢中得到主的指示，不要回去見希律，他們就從別的路回自己的家鄉去了。 \end{tabularx} \\ \\
[1ex]
\hline
\hline
\end{longtable}
$^{1}$聽眾晚安.
是不是很久沒有送聖誕禮物了?.
好像很久都沒有大家有聖誕的祝賀.
聖誕的派對.
希望今年大家有一個熱鬧的聖誕.
在我們慶祝聖誕的時候.
大家記不記得今天是什麼日子?.
剛才陳牧師跟我們說了.
今天是章林期.
明天是章林期第一主日.
我們會有聖誕圖片登出.
這是我們每年都有的節日.
不過過往三年.
可能大家慶祝得沒有那麼熱鬧.
相信今年會更加熱鬧.
在一起預備的時候.
章林期的意思是預備.
預備主耶穌基督的來臨.
應該是猶太人預備理賽亞尼的日子.
我們也預備自己的心.
等待聖誕的來臨.
我們去慶祝之前的一個預備.
第一個禮拜 第一個主日.
章林節的第一個主日的預備.
應該是給我們一些醒察.
我們就看看這段經文.
馬太福音的第二章.
剛才很好.
馬弟兄幫我們讀得很清楚.
我想大家都很熟悉.
在希律作王的時候.
耶穌在猶太的伯尼行降生.
那時候有幾個.
不知道有多少個.
有幾個從東方來的博士.
來到耶路撒冷來找耶穌.
幾個不是三個.
我們從小看聖誕卡都是三個.
為什麼呢.
可能因為禮物是三樣的禮物.

$^{41}$其實很多人都可以送這三樣的禮物.
黃金 乳香和末陽.
因為這三份的禮物.
有代表尊重皇帝.
有代表對皇帝的尊崇.
和一種朝拜.
有一些很獻祭的心意.
所以可能不同的博士.
都會帶這幾份禮物來.
他們來到就先去了耶路撒冷.
因為他們跟著一顆星去.
去耶路撒冷.
為什麼呢.
因為耶路撒冷和伯尼行很近.
之間的距離是九公里.
真的非常近.
一顆星來說.
他們就跟著一顆星去.
去到那裡.
那顆星也沒有再出現.
於是他們就朝著一個方向去.
他們去的目的.
他說了在這裡.
我們在東方看見他的星.
特來拜他.
他們是要來拜他的.
這群就是一群.
從東方來的博學之士.
他們去到耶路撒冷的時候.
當然是四圍問.
其中一個他們會去問的地方.
就是去到那裡最漂亮的屋去問.
因為他們知道是猶太人的王要降生.
猶太人的王要降生.
應該是降生在一個很漂亮的屋裡.
預備做王.
所以他們就去找到希律的宮殿.
希律王聽見了.
有人從東方來.
要拜猶太人的王.

$^{81}$其實他很害怕很緊張.
因為當時他就在羅馬帝國的政治底下.
派去猶太省.
在猶太省當中做分封王.
他在這裡做王的時候.
突然間這裡多了一個王.
代表什麼呢.
對他來說是一個危機.
所以他就不安.
他不安的時候.
耶路撒冷全城的人都不安.
為什麼呢.
我想現在在以巴.
全城的人都不安.
以巴的戰爭底下.
為什麼呢.
因為巴沙不是很大.
有個什麼的戰爭.
很容易臨到很多不同的人.
百姓期望國泰民安.
全城都沒事.
他可以安居樂業.
可以安享.
突然間說有另一個王.
對他們來說.
原來他們覺得是害怕.
並不覺得是一個喜訊.
害怕的時候.
覺得不安的時候.
希律就找了一些祭司.
和民間的文士來問他們.
祭司長是猶太教當中.
代表百姓向神獻祭的.
文士是他們博學於聖經的.
舊約的聖經.
和一些律法典章.
要教導百姓如何守律例.
這班人就被希律召了去.
希律就問了他一件事.
基督要在哪裡降生.

$^{121}$意思就是說彌賽亞.
基督就是彌賽亞.
猶太人一直在等的彌賽亞.
要在哪裡降生.
似乎這班祭司長和民士.
不是要花很多時間去找資料.
他們很快就可以回覆希律.
他在伯利恆降生.
很近的 九公里.
就算當時不是一條直線飛過去.
也很方便.
可以有當時的交通工具.
來到過去.
不是遠的地方.
應該很容易處理.
他們就說到為什麼會是伯利恆.
大家如果有尼加書.
有聖經.
尼加書 舊約十二小鮮之書當中的尼加書.
五章二節就說到.
在猶太地的伯利恆不是最小的城.
為什麼呢.
其實猶太的伯利恆.
是真的很小的一個城.
當時可能只有一條街.
一條大街.
不是很多住戶.
但是又怎樣呢.
他為什麼不是最小呢.
他不是最小.
因為有一位最偉大的.
又是一個說降生.
就是主基督.
這位主基督要降生.
做一個統治者.
不單只是做一個統治者.
還要牧養以色列民.
他不只是要管治以色列民.
他來要牧 要養以色列民.
這個就是祭司長和民事回復希律.

$^{161}$按道理也不錯.
這個訊息應該百姓會很安心.
但似乎不是.
希律是暗地裡.
他是暗地裡去召見這班博學之士.
這班博士.
他不是直接宣他們.
進去希律的皇宮見他們.
他是暗地裡召見他們.
然後就問他們.
這個彌賽亞會在伯利恆出生.
你去伯利恆找他.
我就比較忙.
你去找他.
找到他了.
你就回來告訴我他在哪裡.
我也想去朝拜他.
是不是很好.
東方的博士也相信是.
就去了.
他們覺得有個指引.
起碼有個方向.
伯利恆在哪裡.
就去那個方向.
伯利恆和耶路撒冷很近.
大家看到一個紅色的圓點.
那裡就是伯利恆.
上面有一個紅色的橫線.
那個位置就是耶路撒冷.
其實是很近的位置.
這個就是看到在死海的旁邊.
兩個都是一個很近的地方.
猶太人覺得這兩個地方都很重要.
耶路撒冷是猶太教信仰的中心.
伯利恆有什麼重要.
原來在以色列人的地方當中.
包括在猶太省.
只有一個城被稱之為大衛之城.
大衛是大衛王.
大衛在以色列人眼中.

$^{201}$他是神應他的子孫一路作王的那個王.
那個大衛王.
這座伯利恆城.
就是稱之為大衛之城.
所以彌賽亞在那裡降生.
表明了彌賽亞降生.
是真的要做以色列的王.
所以大家是很重視的.
伯利恆的地位就在這一處.
雖然它是一個最小的城.
真的是一個小鎮.
但是它有它很淑寧的意義.
是說的.
尼加書就說到伯利恆這樣說.
伯利恆的爾法他.
到底是伯利恆還是爾法他呢.
有些地方就會有早期的名字和後期的名字.
伯利恆這個小城.
以前就叫做爾法他.
因為它很小.
後來就叫做伯利恆.
在猶太諸城當中雖然是最小.
但將來必定有一位從那裡出來.
在以色列中為我作掌權者.
他的根源自耿古.
從太初就有.
尼加書就說明了.
這個要出來做掌權者的.
不只是大衛的後裔.
他的根源自耿古.
從太初就有.
其實他是神的地位.
是形容神.
才說他是源自耿古.
源自太初.
這個也是約翰福音第一章第一節就說到.
太初有道.
就是在說耶穌基督.
這個就是在說伯利恆要降生.
這個真的是彌賽亞.

$^{241}$是耶穌基督.
在說這些的時候.
我們來看看怎樣.
希律王聽見這件事本來就心裡不安.
希律不是猶太人.
但耶路撒冷傳承的人是猶太人.
而且都是猶太教的人.
都是信耶和華的人.
為什麼知道一位從神那裡要來.
要牧養他們的神.
是神來的.
不只是王.
為什麼他們聽了之後心裡不安呢.
這個值得我們想一想.
他們真的是猶太人.
是猶太教徒.
是信耶和華的.
他們是在等待彌賽亞的來臨.
但是他們聽到彌賽亞降生的消息.
他們的反應是不安.
這個是他們當時的情況.
這個是希律.
在書本當中找到他.
希律的像.
是不是真的這樣呢.
我真的沒見過他.
但是在書本上找到的就是這樣.
希律是什麼人呢.
希律是羅馬帝國.
他去猶太省作王.
管治猶太省.
他是一個統治者.
他聽到有人要來做統治者.
他緊張也是正常的.
他的反應是怎樣呢.
他的反應就是.
當這些東方的博士去朝拜完耶穌之後.
他就想著得到這個消息.
知道他住在哪裡了.
就可以去殺了耶穌.

$^{281}$這件事他的計劃是瞞不過天父的.
天父知道就派使者.
跟這幾位的博士說.
你們從第二條路走.
不要去跟希律說.
彌賽亞是什麼地方.
免得他來殺彌賽亞.
他不是想要來朝拜.
他是要來殺滅耶塞亞.
要剿滅他將來的禍患.
對於希律來說.
耶穌基督帶給他的不是安寧.
我們經常說信了主之後.
你是安寧的.
心裡面有平安.
有神的看顧照料.
有什麼事你可以向祂祈禱.
但是對希律來說.
耶穌的來臨是令到.
他的王位受到威脅.
如果耶穌基督來.
做了君王.
那就是令到羅嗎哪邊.
覺得你統治不了.
那我換一個人.
或者你們兩個打完贏了之後.
那個做吧.
那他的一生事業就沒有了.
耶穌基督的降臨.
令到希律覺得.
我的事業受到威脅.
這個希律是什麼人呢.
這個叫大希律.
這個意思是.
他是senior.
他是第一代希律王朝的君王.
還有幾個跟著他的.
是他的子女.
是他的兒子.
是後來的希律王.

$^{321}$這個希律王是生於主前74年.
在主前37年開始.
羅馬政府派了他去猶太省.
你都要做猶太人的君王.
在那裡的時候.
他就一直統治到主前四年.
主前四年耶穌基督出世了.
是的.
你記住公元多少年.
是後期退回去數的.
不能夠數得那麼清楚.
應該照計主耶穌基督.
大概在主前四至六年.
很多學者都認為.
是主前六年出世的.
為什麼呢.
因為在同一章的聖經當中.
13至15節就說到.
後來希律為了自己的王權.
他就寄言這班博士.
不告訴我耶穌在哪裡.
我就殺了整個百里行城.
不是 是百里行城和百姓行城.
四位不是初生的嬰兒.
是兩歲以下的孩童.
所以估計他應該那個時候.
耶穌基督應該是大約一歲多兩歲.
而希律是有殺錯沒有放過的人.
可能耶穌一歲.
他連兩歲的人都殺了.
這是當時的希律.
他真的有殺錯沒有放過.
是 他對所有人都懷疑.
他的妻兒子女.
誰令他覺得不忠心.
有機會想謀朝篡位.
他就殺了他.
他的兒子也被他殺了.
不要說他的朝臣.
對他來說.

$^{361}$耶穌基督會威脅到他的王權.
不是他一個人的王權.
他得了這個王權之後.
是可以世襲的.
下一代都可以做的時候.
耶穌來威脅到他的時候.
他就真的要殺了他.
所以他就殺了百里行和百里行四位.
山區 郊區的孩童.
令到發生了一件很多家庭受到創傷的事.
為什麼呢.
因為威脅到他的事業.
弟兄姊妹.
我們今天是張臨期第一主日.
我們是一個反省的時候.
有沒有一些事物.
或者是信仰的事.
令到我們感覺到.
我們的事業.
我們的苦幹.
受到威脅的時候.
我們會很抗拒.
很大的抗拒.
不單止是不信.
甚至是做出一些.
很不應該做的行動.
這個值得我們想一想.
當時除了希律之外.
希律是這樣.
除了希律之外.
還有一班人沒有出聲的.
這班人是什麼人呢.
是百姓.
百姓不知道他的目的.
但這班百姓都是不想有危險.
你塞阿萊林嗎.
他們去了百里行嗎.
他們去吧.
我沒空 是不是.
大家都要工作.

$^{401}$我沒空.
我要照顧家庭.
他們去吧.
他們去朝見耶穌.
看看將來怎樣.
千萬不要打到你.
是不是.
這班是耶路撒冷的百姓.
耶路撒冷是當時猶太教的中心點.
宗教中心點.
現在我們稱之為聖地.
是不是.
在那裡居住的人.
在那裡居住的人應該很敬虔的.
他們住在聖殿附近.
聖殿山在那裡.
他們時時都得到很多提醒.
很多祭司在那裡.
很多文士在那裡.
時時得到教導.
他們的反應是.
你塞阿萊林.
你自己看著自己.
我也要看著自己.
在整件事當中沒有發聲.
沒有任何他們的對白.
沒有說下他們有什麼反應.
是不是信徒也有時會這樣.
有危機的時候.
我們怎樣面對呢.
求主幫助我們.
在新一年當中.
我們能夠和神的心意更加貼近.
神對我們有差遣的時候.
我們能夠積極去參與.
今天我收到一些相片.
我們有弟兄.
他有份帶領去土耳其.
一些救濟的工作.
和傳呼音的工作.

$^{441}$他傳了一些相片回來.
和當地的人為他們做一些醫療.
除了醫療服侍之外.
和他們講福音.
和他們一起祈禱.
有很多事情原來我們可以參與.
在別人危機的時候.
我們可以幫助.
求主也幫助我們.
我們做一個.
在一個自由地可以敬拜的人.
能夠關心到.
不能夠有自由敬拜的人.
能夠將神的福音.
向別人的孫子傳揚.
這件事不是一個故事.
這件事能夠記載下來.
因為有一班博士.
這班博士找了希律去問了這些問題.
所以知道彌賽亞的降生.
這班博士是什麼人呢.
他們不是某間大學的研究生.
這班博士意思就是說.
他們專門研究天文和一些星象學.
也就是說他們是其他國家.
外國 即猶太人的外國.
或者是羅馬帝國的其他地方的.
一些省份的什麼呢.
的祭祀.
他們也敬拜耶和華.
不是.
他們簡單來說是異教的祭祀.
信第二種宗教的人.
信第二種宗教的祭祀.
應該是敬拜他的宗教.
原來這一班的祭祀.
這一班博學的祭祀.
他很留意.
到底除了自己所拜的那位神之外.
還有一位是要降生的.

$^{481}$於是乎他們不遠千里而來.
來到耶路撒冷要尋找耶穌.
這班人是很特別的.
這班人不是猶太人.
不是猶太教徒.
不是明白猶太教是什麼的.
不是信奉耶和華的.
但是來找彌賽亞.
是我們應該慚愧的.
是嗎.
是猶太人很應該慚愧的.
來找彌賽亞的.
去到伯利恆.
去朝拜耶塞亞的.
去獻上黃金與香活躍的.
是一班異教.
不是異教徒.
是異教的祭祀.
這一班異教的祭祀.
竟然在尋求.
弟兄姊妹.
你有沒有遇過有異教徒.
尋求信仰向你求問.
如果是應該怎樣做.
我們第一個反應就是.
哇 他是異教徒.
做到很高位的.
我怎會懂得怎樣跟他說.
對嗎.
是喔 怎算.
算了 他是不是異教徒.
我們都覺得不懂得怎樣跟他說.
有什麼分別.
我們應該做的就是.
好好的祈禱.
深度祈禱.
既然他問就答.
叫他信主的永遠都不是我們.
無論他是異教徒.
他是無神論者.

$^{521}$或者他是一個.
曾經在基督教學校讀書的人.
他問到我們.
我們都不能夠叫他信主.
我們只能夠按我們所明白的.
向他解答一點點.
回答他.
叫他信主的是神.
唯有聖靈能夠在他心中.
感動他 令他信主.
這個才是.
令人可以得到救恩.
是神.
祂計劃了這個救恩.
耶穌基督為我們完成了救恩.
聖靈在每個人的心靈當中.
去感動我們.
讓我們能夠明白.
其實這一班祭司.
這班異教的祭司.
他們受到神的靈的感動.
他們是看到那顆星就過來.
他們不是住在附近的.
住在附近的那些.
看到那顆星都沒有反應.
是嗎.
是喔.
然後他們朝拜完之後.
他們是有神差遣.
使者去跟他們說.
叫他們不要走回原路走.
不要去耶路撒冷告訴希律.
你走另一條路回家鄉.
他們就聽話地.
走另一條路回家鄉.
沒有去找希律.
他們是很認真的.
他們是很認真地.
尋求信仰的人.
是很認真地.

$^{561}$按著他們所得到的領受.
按著神.
藉著天使去感動他們的說話.
他們去做的.
好像聽話過我們一樣.
是喔.
原來他們是這樣.
如果有機會.
真的有機會有人問你.
無論他是無神論者.
他是很久沒有回教會的人.
或者他正在回其他宗教的人.
主都會給你力量.
他來找你.
神會不知道的.
可能是神預備的機會.
神要使用你.
讓你可以讓他知道.
原來有天國.
原來信主可以得永生.
原來附近有教堂.
原來我可以回教會.
是可以的.
求主幫助我們.
他們要獻上禮物.
唯有他們可以獻上禮物.
還有兩個人.
這幅照片那個不是嬰兒.
那個是小孩.
還有兩個人陪著小孩耶穌.
就是瑪利亞和約瑟.
這兩個人一直很近親地.
在服侍耶穌.
我們叫事奉.
他們原來很低調.
他們在大理行生活了一兩年.
都很低調地沒有出聲.
為什麼要這麼低調呢.
不是應該宣揚一下嗎.
原來他們除了宣揚之外.

$^{601}$可以做一些低調的事奉.
為什麼這麼低調呢.
原來耶穌基督的降生是很低調的.
我沒有去過聖地.
但是聽聞一些從聖地回來的人.
都有些照片 有些分享.
大理行有一座殿叫主誕堂.
主誕堂是門口不高的.
如果大家是比較高大.
就要蹲下身進去.
堂就應該是在比較地下的位置.
原來耶穌基督未必在那間主誕堂出世.
那是一個紀念.
不過耶穌基督出世的地方.
應該都類似這樣.
因為大理行人住的地方.
是分開兩層的.
一層在地下 一層在地上.
大家以前聽過前舖後居.
大理行人就是上面那層居住.
下面那層就要養運輸的動物.
所以耶穌基督說是一個養馬的地方.
出世的.
然後他的BB廠就是馬槽.
應該是一個低落一層的地方.
因為當時耶穌基督出世的時候.
所有的客店都住滿了人.
當時的客店應該是B and B.
不是真的Hotel.
因為Hotel真的很貴.
記載在哪裡呢.
《路加夫音》第二章.
大家可以看看.
耶穌基督出生.
神的兒子出生.
他真的不是希律家.
他是君王家.
他是一個近乎是地底下的位置.
是一個養馬或是養驢的地方.
他的BB床.

$^{641}$如果說是一個馬槽.
即是裝馬的食物的東西.
其實需要好好地潔淨一下.
都不是乾淨.
他很低調地出生.
那晚有誰來朝拜呢.
沒有.
如果朝拜就不會讓他住在那個地方.
誰都會給他騰房間.
不過有天使向野地的牧羊人傳播這個消息.
有牧羊人來朝拜過.
不過他們走了之後.
這個消息沒有傳出去.
不知道他們還記不記得呢.
但是瑪利亞和約瑟一直都照顧著BB.
當時一直都照顧著牠成長.
原來都有一些一直都在事奉.
未必很多人知道.
但是他們繼續努力事奉.
而有一些真的聽了都不以為意.
當時很開心.
過完之後就忘記了.
都沒有將這件事放在心.
每年都去找一下耶穌在哪裡.
祝祂生日快樂.
沒有.
也沒有去耶路撒冷當時的宗教地點去講.
是的.
弟兄姊妹.
我們可能有時都是.
聽了很多.
聽完我們繼續走我們人生的道路.
每天我們都要走的路.
我們繼續走.
我們都會忘記了.
不說聽了會忘記了.
你讀經都會忘記了.
問一下你.
上個星期六你為了什麼事祈禱.
我們就是這樣.

$^{681}$連自己祈禱的事都忘記了.
不是神有沒有回覆我們.
回覆了我們都不知道我們祈禱了.
求主幫助我們.
我們忙碌的腳步下.
我們留些空間給自己.
留些空間給人.
在忙碌的腳步下我們留意一下神.
跟我們有什麼說.
有什麼事等著我們去做.
在你的認識的朋友當中.
你的親戚當中.
你的隔離鄰舍.
你隔離兩條街的人.
還有一些都是在忙碌當中.
今年大家都很忙碌.
工作的忙工作.
旅遊的忙旅遊.
只是想深圳吃飯都很忙.
是不是.
在大家的忙碌當中.
有多少位認識神.
還沒有回教會.
有多少位還沒有認識神.
有多少位等著你跟他分享.
你祈禱神聽了你.
按著你的祈禱賜福你.
剛才說到.
可能有一些不是信耶穌的人.
怎麼跟他分享呢.
我們身邊可能很多異教徒.
剛才說到有些異教的祭司.
我也有機會跟幾個不同的異教的.
是不是叫祭司呢.
異教的宗教人物.
有一個是扶乩.
有兩個應該是幫人算命.
還有一個是看風水.
但是他們有需要的時候.
經過弟兄姊妹的介紹.

$^{721}$有機會跟他們傳福音.
都看到他們.
原來是真的到自己有需要的時候.
他們知道自己所信的.
能不能夠給他們平安.
能不能夠給他們永生.
其實他們知道的比我們可能還多.
因為我想他們屬於祭司類的.
這樣的情況下的時候.
他們怎樣呢.
他們去到自己有需要的時候.
他們覺得自己幫不到自己.
自己所敬拜的幫不到自己.
於是他們怎樣做呢.
於是他們需要尋求信仰上的幫助.
原來真的跟他們分享的時候.
他們是有機會會信主的.
他們是有機會會很認真地去祈禱的.
他們當看祈禱這件事.
是比我們更加認真.
我看到他們原來帶到一個異教的宗教人物.
信主的時候.
是影響到他一家的.
不會是他一個人信的.
是影響到他一家的.
求主幫助我們.
當他真的要找我們做的時候.
我們不要告訴他.
我很忙啊.
下個小時就已經有約了 天父.
我們不要這樣跟他說.
我們就說 天父 我不是很懂.
不過你願意教我的話.
你怎樣教我 我就怎樣做 是不是.
我們求主幫助我們.
有神去幫助我們 我們做得到的.
我們真的很不懂嗎.
那我們就學.
第一樣我們要學的.
今天我們自省的.

$^{761}$第一樣我們要學的.
上一次讀聖經是什麼時候呢.
這件事要想一下.
我們要看下.
我們跟著這一年.
我們要怎樣看聖經.
上一次很認真地祈禱是什麼時候呢.
有沒有一個恆常我們祈禱的時間.
有沒有一個恆常.
我們跟家人一起祈禱的時間呢.
除了真的謝飯.
坐在一起謝飯之外.
我們一直學的時候.
教會都有一些預備的.
有一些課程.
我們下一個12月.
我們可以開始參加去學習的.
也有一個報道會.
什麼時候的報道會.
12月24號.
應該大部分人都不用上班的.
約定你的朋友.
請他回來報道會.
為了約他.
你不要一開口就打電話約你報道會.
你陪陪他.
找找他聊聊天.
約他喝茶.
慢慢地邀請他回來.
約他回來報道會.
讓他有機會聽福音.
有需要的時候介紹給教務式.
讓教務又可以認識他.
跟他分享多一點.
跟著這個季度的主日學.
有很多是適合大家去學習的.
也都適合你的朋友去學習的.
可以帶他們回來.
求主幫助我們.
能夠有教會給我們一些資源的時候.

$^{801}$我們就好好一起去努力.
讓我們是什麼.
我們是信徒.
我們不只是平信徒.
平信徒我們都不是躺平的信徒.
我們是認真的信徒.
求主幫助我們.
接下來這一年成為一個更認真的信徒.
我們一起低頭禱告.
真是天父上帝我們來到你的根前.
是你的恩典.
我們何許人.
你竟然願意讓基督為我們而死.
拯救我們.
叫我們可以因著你的緣故.
因著基督的犧牲.
因著聖靈的感動.
我們可以得著救恩.
得著永生.
求主幫助我們.
幫助我們能夠將我們得到的恩典.
跟我們身邊的人分享.
讓他們都能夠得著.
讓他們都能夠被你來拯救.
都能夠稱呼你為天父.
成為你的兒女.
求主幫助.
主啊願你使用我們每一個.
建立我們每一個.
成為你的僕人.
願望你奉主耶穌基督的聖名.
阿們.
\newpage



\section{使徒行傳 5:12-32}
\label{sec:7XxklgVuie4}
\textbf{2023-12-03 午堂《傳講生命之道》使徒行傳 5\_12-32 講員:曾裔貴牧師}
\newline
\newline
連結: \href{https://youtube.com/watch?v=7XxklgVuie4}{\texttt{https://youtube.com/watch?v=7XxklgVuie4}} ~~~~ 語音日期: 2023-12-06
\newline
\newline
\hyperref[sec:gy4cK2fsHdc]{\small{< < < PREV SERMON < < <}}
~
\hyperref[sec:index_chronic]{\small{[返順時目]}}
~
\hyperref[sec:index_scriptual]{\small{[返順卷目]}}
~
\hyperref[sec:Qx_35UZ1A44]{\small{> > > NEXT SERMON > > >}}
\newline
\newline
使徒行傳 5:12-32
\newline
\begin{longtable}{cl}
\hline
\hline
章節 & 經文 (和合本修訂版)\\
\hline
5:12 & \begin{tabularx}{0.7\textwidth}{X} 主藉使徒的手在民間行了許多神蹟奇事；他們都同心合意地聚集在所羅門的廊下。 \end{tabularx} \\ \\ \relax
5:13 & \begin{tabularx}{0.7\textwidth}{X} 其餘的人沒有一個敢接近他們，百姓卻尊重他們。 \end{tabularx} \\ \\ \relax
5:14 & \begin{tabularx}{0.7\textwidth}{X} 信主的人越發增添，連男帶女都很多， \end{tabularx} \\ \\ \relax
5:15 & \begin{tabularx}{0.7\textwidth}{X} 甚至有人將病人抬到街上，放在床上或褥子上，好讓彼得走過來的時候，或者影子投在一些人身上。 \end{tabularx} \\ \\ \relax
5:16 & \begin{tabularx}{0.7\textwidth}{X} 還有許多人帶著病人和被污靈纏磨的，從耶路撒冷四圍的城鎮來，他們全都得了醫治。 \end{tabularx} \\ \\ \relax
5:17 & \begin{tabularx}{0.7\textwidth}{X} 於是，大祭司採取行動，他和他所有一起的人，就是撒都該派的人，滿心忌恨， \end{tabularx} \\ \\ \relax
5:18 & \begin{tabularx}{0.7\textwidth}{X} 就下手拿住使徒，把他們押在公共拘留所內。 \end{tabularx} \\ \\ \relax
5:19 & \begin{tabularx}{0.7\textwidth}{X} 但在夜間主的使者開了監門，領他們出來，說： \end{tabularx} \\ \\ \relax
5:20 & \begin{tabularx}{0.7\textwidth}{X} 「你們去，站在聖殿裡，把這生命的一切話講給百姓聽。」 \end{tabularx} \\ \\ \relax
5:21 & \begin{tabularx}{0.7\textwidth}{X} 使徒聽了這話，天將亮的時候就進聖殿裡去教導人。大祭司和他一起的人來了，叫齊議會的人和以色列人的眾長老，然後派人到監牢裡去把使徒提出來。 \end{tabularx} \\ \\ \relax
5:22 & \begin{tabularx}{0.7\textwidth}{X} 但差役到了，不見他們在監裡，就回來稟報， \end{tabularx} \\ \\ \relax
5:23 & \begin{tabularx}{0.7\textwidth}{X} 說：「我們看見監牢關得很緊，警衛也站在門外，但打開門來，裡面一個人都不見。」 \end{tabularx} \\ \\ \relax
5:24 & \begin{tabularx}{0.7\textwidth}{X} 守殿官和祭司長聽了這些話，心裡困惑，不知這事將來如何。 \end{tabularx} \\ \\ \relax
5:25 & \begin{tabularx}{0.7\textwidth}{X} 有一個人來稟報說：「你們押在監裡的人，現在站在聖殿裡教導百姓。」 \end{tabularx} \\ \\ \relax
5:26 & \begin{tabularx}{0.7\textwidth}{X} 於是守殿官和差役去帶使徒來，並沒有用暴力，因為怕百姓用石頭打他們。 \end{tabularx} \\ \\ \relax
5:27 & \begin{tabularx}{0.7\textwidth}{X} 他們把使徒帶來了，就叫他們站在議會前。大祭司問他們， \end{tabularx} \\ \\ \relax
5:28 & \begin{tabularx}{0.7\textwidth}{X} 說：「我們不是嚴嚴地禁止你們，不可奉這名教導人嗎？看，你們倒把你們的道理充滿了耶路撒冷，想要叫這人的血歸到我們身上！」 \end{tabularx} \\ \\ \relax
5:29 & \begin{tabularx}{0.7\textwidth}{X} 彼得和眾使徒回答：「我們必須順從神，勝於順從人。 \end{tabularx} \\ \\ \relax
5:30 & \begin{tabularx}{0.7\textwidth}{X} 你們掛在木頭上殺害的耶穌，我們祖宗的神已經使他復活了。 \end{tabularx} \\ \\ \relax
5:31 & \begin{tabularx}{0.7\textwidth}{X} 神把他高舉在自己的右邊，使他作元帥，作救主，使以色列人得以悔改，並且罪得赦免。 \end{tabularx} \\ \\ \relax
5:32 & \begin{tabularx}{0.7\textwidth}{X} 我們是這些事的見證人；神賜給順從的人的聖靈也為這些事作見證。」 \end{tabularx} \\ \\
[1ex]
\hline
\hline
\end{longtable}
$^{1}$請姐妹主女們平安.
我們一起再祈禱.
感恩天父給我們今早我們能夠再一次聚集.
很盼望教會我們能夠將主耶穌救贖世人的恩典.
赦罪的權柄.
我們能夠充分了解明白.
大家一起齊心努力.
為福音效力.
今天是第一個星期的章臨期.
我們很盼望在過去一些先賢.
一些過去在教會的歷史裡面.
他們一直建構四個星期的等待耶穌基督的復活降臨.
也是我們今天等待主耶穌第二次回來的時間.
教會需要有這種警醒,醒悟的心靈.
來仰望主第二次從高天降臨的時間.
這個萬民得到救贖的時候.
幫助我們明白我們處身在現在我們的社會.
應該怎樣為主作見證.
願神帶領我們.
今早我們等候在你的面前.
求主藉著小行傳這段發生的歷史.
給我們能夠有啟發.
一起的護臉.
奉主耶穌基督的名.
阿們.
剛才也提過.
我們也有一個.
當然不是很隆重.
好像禮儀教會一樣.
我們會有章臨期的慶祝.
今天第一個星期的主日.
就是傳統禮儀教會.
他們會燃點紫色的蠟燭.
象徵著四個星期的章臨期的開始.
就好像等於教會年的新年.
新的開始.
紫色也好像代表著.
一種對於救恩的了解.
尤其是那個醒悟.
有一種這樣的含義.

$^{41}$慢慢已經成為了.
很多年教會的歷史裡面.
他們都會這樣做.
特別是禮儀教會.
無論天主教或聖公會.
他們就會隆重其事.
這樣去做這些事情.
我們今天在這段經文裡面.
主要集中想說一說.
也都相當明顯.
是天使使道使徒.
可以繼續站在聖殿裡面.
來去說.
特別留意第21節.
使徒聽了這番說話之後.
天將亮的時候.
就進聖殿裡面去教導人.
當然就是全港耶穌基督.
是怎樣的一位主.
特別是第29節.
他重申.
向當時審問他的祭司長.
和當時的民間長老.
來去重申我們全港的.
我們教導的是甚麼的內容呢.
很清楚.
彼得和眾使徒回答.
我們必須順從神.
勝於順從人.
即是你們不准我說.
但是你可以禁止我這個人.
但不能夠讓這個訊息被阻礙住.
這個訊息就是第30節.
你們掛在木頭上殺害的耶穌.
我們的祖宗的神.
就是這位猶太人的神.
已經使到他復活了.
神也都將他高舉在自己的右邊.
在神寶座.
永恆的寶座的右邊.

$^{81}$使到他成為元首.
使到他成為人類的救主.
使到以色列人得到悔改.
並且罪得到赦免.
留意第32節.
一個總結句很重要.
不僅僅是我們膽大不怕死.
為這件事作見證.
民間所有人都說.
耶穌死了.
屍體都不見了.
但是他們好像力排眾議.
要將這個見證繼續傳遞.
不僅僅是我們試圖.
留意一下這裡很重要.
第32節.
下半節.
彼得告訴我們.
神賜給順從的人的聖靈.
聖靈在哪些人的生命裡面.
特別的顯現呢.
就是願意順從神的人.
聖靈會在他們的生命裡面.
作見證.
為這些事情作見證.
其實如果我們就這樣看.
那個戲劇性不是很大.
因為第五章.
作者路加醫生給我們有幾件事情.
來襯托這個見證.
在聖殿這一次的見證.
因為四章有一個這樣的復興.
就是當時有一個叫做巴拿巴的人.
他是塞浦路斯人.
其實他自己的原名是約瑟.
他是利未人.
當然他後來在塞浦路斯住.
他就上到耶路撒冷.
當時是有一個小型的復興.
那個復興就是家家信主的.

$^{121}$都賣了自己的田產.
來交在使徒的腳前.
讓使徒去調配.
運用州制不同的人.
按當時的需要來分配.
當時形成了一股氣氛.
大家很想耶穌快點回來.
很想整個耶路撒冷.
很多人能夠信主明白救恩.
所以就有這樣的一個.
你可以說復興帶來的.
大家不介意自己的財物.
這裡是非常之美麗的.
第四章三十二節.
許多的信徒一心一意.
沒有一個人說自己.
任何的東西是自己的.
都是大家公用.
使徒就有大能見證主耶穌復活.
這樣來描述.
當時的小小背景.
因為這幾件事.
待會我們將它解釋完之後.
再重新看三十二十九節.
你就會發現.
裡面有很強的生命力.
就是生命之道帶來的改變.
這裡告訴我們.
凡物大家公用.
因為有很多人將田產房屋.
都賣了.
將所賣的錢拿出來.
放在使徒的腳前.
照每個人所需要的.
分給每個人.
然後就如描述剛才提過.
就是這個巴拿巴.
三十六節.
有一個利未人.
名叫約瑟.

$^{161}$使徒稱他為巴拿巴.
很明顯這個就是一個花名.
巴拿巴翻出來.
就是一個安慰之子.
很喜歡安慰人.
他就生在塞浦路斯.
三十七節說他也有田地.
也賣了.
又將錢拿來放在使徒的腳前.
為什麼這樣描述.
教會當時的光景呢?.
就是說大家好想促成.
耶穌基督快點回來.
快點來完成.
主耶穌在升天之前提到的大使命.
大家願意付出自己的部分.
甚至有些人是全部.
在他基本的生活可以解決之後.
他就將自己擁有的.
全部拿出來.
任由教會去分配.
你說多麼美麗.
很特別.
張醫生再另外舉多一個例子.
就是第五章.
有對夫婦.
他又照做.
不過那個動機是相當邪惡.
在一個大復興的時期.
這樣的行徑.
突然間死於非命.
還要死在教會.
他帶來非常震撼.
這樣描述.
而引發下面的29節.
那個生命之道.
很值得我們今天大家一起.
在章林期的第一個禮拜.
去想想我們的生命.
有沒有一個醒悟.

$^{201}$對神的道.
對耶穌基督的復活.
我們有沒有一個醒悟.
以致我們的生命.
很踏實地按部就班的成長呢.
因為有一對夫婦.
她不按部就班.
她不想成為屬靈.
她想人看到她以為屬靈.
原因就是.
有一個人名叫阿拿尼亞.
和他的太太叫撒菲拉.
賣了田產.
因為當時大家都是這樣做.
大家都認為這樣是為神.
為神的國度.
為福音.
為主耶穌.
為復活的耶穌.
我出一分力.
賣了田產.
那你說有什麼問題呢.
第二節.
盧家很快就揭穿那個故事.
他想要鋪排的.
那個生命之道.
在有些人身上原來沒有.
沒有生命之道在他內心.
這個非常嚴肅.
弟兄姊妹.
第二節.
她將錢私自留下一部分.
更加諷刺.
更加有趣的就是.
再加一句.
她的妻子也知道.
不是被迫的.
非情緣的造假.
來配合劇情.
而是知道丈夫所做.

$^{241}$因為之後就有事發生.
即是說兩夫婦處心積慮.
希望人家看她是屬靈.
而不是她真正的成為屬靈.
成為屬靈要很多時間.
很多人生.
很多經歷.
之後就有很多經歷.
試圖.
很多的挫敗.
監禁.
審問.
恐嚇.
很多這些事情.
但是生命之道在內心.
就是無論環境順和逆.
我都要繼續傳揚.
這一位我無法否認的.
復活了的耶穌.
所以他的妻子知道之後.
其餘的部分就拿來.
放在使徒腳前.
第三次出現.
放在使徒腳前.
個個都拿來放在使徒腳前.
心甘情願.
按你自己的需用.
按你自己的能力.
可以做多少.
是沒有任何的壓力的.
在這裡就出現了.
成為屬靈.
即是我對屬靈的事.
我真是很明白.
主耶穌的愛實在太大了.
我願意在我必須的需用之外.
那些不太需要的.
我願意奉獻出來.
讓神的家可以怎樣去運用.
這個叫做成為屬靈.

$^{281}$但願意人看他是屬靈.
這個假裝出事了.
第三節.
彼得為什麼會知道呢.
彼得就說.
阿拿尼亞.
為什麼撒旦充滿了你的心.
使你欺騙聖靈呢.
把賣田地的錢.
撕至留下一部分呢.
田地還沒賣.
不是你自己的嗎.
即是說是你的.
沒問題的.
你可以完全不賣的.
賣了.
錢都是你作主的.
你給多少是沒問題的.
你說彼得.
我也要想想我們夫婦的生活.
我只能奉獻一部分.
是沒問題的.
邏輯很清楚.
接著.
這句就大件事了.
你怎麼心裡會這樣做呢.
你不是欺騙人.
是欺騙神.
這個就是.
很想人看他是屬靈.
而他不是成為屬靈.
弟兄姊妹.
我昨天代表教會去澳門.
黎剛慧牧師.
我們以前的童工.
他在那裡成長.
初中.
一路長大.
後來這班大學生.
有感動.

$^{321}$有好幾個去建道神學院.
讀神學.
講到明.
他們心裡的照明.
學成之後.
當然有一段時間.
在我們中間實習.
來到去成為童工.
他們很清楚.
幾位的神學生.
是要回去澳門.
來到去.
在這個硬土地.
他們所成長的地方.
來到去傳福音.
昨天我有機會見證.
這個教會.
由一個幾十人.
已經是三十六年.
到現在買了一個地方.
可以坐三百人.
這一班的執事.
是小時候成長的.
在老牧師.
這位蔡牧師.
這麼多年的心血.
這樣來到去建立底下.
這一班的執事.
現在是個個都在社會上.
算是澳門的社會精英.
有牙醫.
有政府的公務員.
有教師等等.
他們由初中.
和黎光偉牧師一起長大.
有些甚至乎是幼稚園的同學.
這樣的幾十年.
積累的生命力.
現在神好像讓種子.
萌芽成長.

$^{361}$弟兄姊妹.
成為屬靈.
要很長時間.
每日自己來到去面對神.
讀經 禱告.
不斷讓神的道.
去連正自己的心靈.
現在出現了一個這樣的風潮.
大家蜂擁認為.
怎樣是屬靈.
賣田產是屬靈.
因為有很好的明星.
被人看到.
阿娜琳 薩菲拉.
他們都這樣做.
如果不是這樣做.
可能會構成.
你們有沒有做 有沒有賣.
其實問問而已.
沒有惡意.
沒有若有所指.
他們就好像有些膽怯.
不賣好像很沒面子.
賣了.
一來可以令到別人不會出聲.
又可以帶來.
有一個明星.
突然間那個屬靈的身份.
在很多人的眼中.
會被抬高.
這樣就做了.
這件事.
第五節.
阿娜琳一聽見彼得這樣說.
就仆低 斷氣.
所有聽見的人.
都非常之懼怕.
神在一個大復興時期.
讓教會的罪惡.
一些很想急速成為屬靈人的.

$^{401}$突然間倒地斷氣.
你可以想像得到.
對教會的震撼有多大.
即是彼得很有屬靈的權柄.
明察秋毫.
看到你內心有一個.
這麼嚴重的裝假扮屬靈.
本來是完全不需要的.
沒有規定.
你一定要每個人來到教會.
來到耶路撒冷的教會.
都要賣掉你的田地.
從來沒有這樣的要求.
但是阿娜琳這樣死了.
你可以想像得到.
處心積慮.
兩夫婦.
太太不認同.
是很合理的.
丈夫當時是一個家庭.
社會的頭.
如果丈夫又很暴躁.
很專橫.
太太要就範.
是沒有問題.
一定是這樣.
但是恐怖的地方就是.
第七節.
過了三個小時.
路加要這麼詳細記載這件事.
他的太太回來.
不知道丈夫死了.
送去醫院都不知道.
還不知道所發生的事.
第八節.
彼得好像給一個最後機會.
彼得對她說.
你告訴我.
你們賣田地的錢.
是不是就這麼多.

$^{441}$如果我們抽離上下文.
有沒有搞錯.
彼得你這麼緊張人家奉獻的錢.
奉獻不夠.
你就好像不放過他.
但是你有上下文.
你就知道不是這樣的事.
她說.
這句話是最後的悔改機會.
其實不關我的事.
是我的丈夫阿拿尼亞.
逼我要這樣去裝假.
其實我不想的.
我都說過一次.
我在屯門醫院工作的時候.
我的上司.
剛剛聽完他大老闆的電話.
大家談完一些事工.
那時候下午三點半.
出來跟我們說.
我現在去做你治療.
如果大老闆打電話來.
你說我離開了.
本來不需要我這個崗位.
進去我上司的房間聽電話.
剛剛我對上的上司.
就放大家.
變成到我那個階級.
走了沒多久.
房間真的響電話.
個個都說快點去聽電話.
不知道是不是大老闆打來.
在總部的大老闆.
果然是.
那一刻的試探就很大了.
感恩啊.
那時候剛剛初訊.
鎮靜了兩秒鐘.
是啊.
我上司剛剛離開了.

$^{481}$他說他去了做你治療.
大老闆就很激氣.
輾轉就說不關我的事.
我說了我應該說的部分.
是他說你叫我這樣說.
但是薩菲拉不是這樣.
薩菲拉說就是這些.
有意識地配合丈夫的劇情.
又不知道丈夫已經死於神的擊打.
而給他一個最後機會.
不揭穿他.
很想他有一個醒悟.
親愛的電影姐妹.
張琳琪是給我們醒悟.
上個星期一.
有機會跟神樂園的同學.
因為他們有些講車不上.
上去珠海.
去朝聖的山墓.
第一次去.
哇 真的很大.
我看到一樣東西很特別.
很多張琳琪的擺設佈置.
連內地無神的.
都擺設這些聖誕裝飾.
很大的櫃檯.
擺滿了聖誕裝飾.
但是背後的生命之道.
相信他們不會留意.
相信如果沒有人幫他們解釋.
他們只是當作一個聖誕點綴.
而不是真正達至他們內心.
對主耶穌的復活.
這種敬畏和期望自己的內心.
能夠成為那個成聖.
所以撒菲拉這樣說之後.
第九節.
彼得對她說.
你們為什麼同謀去試探主的靈呢?.
一而再說不是得罪我.

$^{521}$是得罪聖靈.
你看埋葬你丈夫之人的腳.
已經到門口了.
他們也要把你抬出去.
第十節.
她立刻撲倒在彼得腳前.
斷了氣.
那些年輕人進來見她已經死了.
就把她抬出去.
埋在她丈夫的旁邊.
第十一節就是一個總結.
全教會和所有聽見這些事情的人.
都非常之懼怕.
這個就是罪惡.
因為要大進的就是二十九節這一段.
彼得和使徒們被監禁之後.
就要在公會的面前.
來講述他們到底見證什麼事情.
就是這個生命之道.
接著還有一點點的.
路加的描述.
第十二節提到.
主耶穌藉著使徒.
在民間行了許多的神蹟奇事.
他們同心合意聚集在所羅門的廊下.
有沒有印象.
約翰福音十章.
在這個修奠節的時候.
主耶穌就在這個所羅門的廊下.
被那些法利賽人.
齊心的來逼迫他.
你快點講吧.
很招認.
到底你是不是基督.
耶穌這樣講.
接著耶穌就講到.
我的羊會聽我的聲音.
你們不是我的羊.
我的羊是會跟隨我.
我會賜給他們永生.

$^{561}$他們永不滅亡.
誰也不能從我的手裡面.
將他們奪去.
是在講那些門徒.
你可以想像得到.
現在這一班的使徒.
耶穌回到天家了.
他們現在聚集在所羅門的廊下.
就是當年.
其實是去年.
耶穌在那個廊下.
來到去講道.
和那些法利賽人對質.
證明耶穌是不是基督復活的彌賽亞.
但是現在他們就是見證.
耶穌就是復活的彌賽亞了.
親愛的兄弟姐妹.
這裡非常值得我們大家一起.
在章林期.
大家一起去醒悟.
大家一起去反思.
到底我們對舊恩.
是不是真的深切的明白.
和我們真的有那個經歷.
在我們的人生際遇裡面.
我們都能夠有這個經歷.
以致我們不是去論述一個很偉大的舊恩.
而是我們正在經驗舊恩.
帶給我們的生命的改變.
我們不是要追求一個成功的神學.
我們是要追求一個成長的生命.
是慢慢地來到與日俱長的屬靈的生命.
尤其是先從內心裡面.
對神的聖靈的那個順從.
這個是三十二節那裡所描述的經文.
接著下去.
路加稍微描述一下.
主耶穌藉著使徒做過什麼呢?.
第十五節.
甚至有人將病人抬到街上.

$^{601}$放在床上和浴紙上.
好讓彼得走過來的時候.
他的影投在這些人身上.
這些真是非比尋常.
影都可以救人.
可以醫病.
還有許多人帶著病人.
和被鬼附的錢磨的人.
錢磨即是法作.
由耶路撒冷四圍的聖真理.
他們全部得到醫治.
這些是頭條號外的新聞.
周邊的人都蜂擁來.
而且是神跡性的立即醫好.
這個就是作者要給我們看到的.
當時的情況.
復興可以去到這樣的情況.
在教會內部.
亞拿尼亞似乎這樣描述.
不是一個普通的.
剛剛初信的猶太人.
可能有相當在教會內的資歷.
初信的不會想人覺得他屬靈.
祈禱都不敢.
說兩句就收聲.
因為沒有那麼多詞彙.
去承載他腦袋內對神的思念.
但信主內了.
就會有很多很多的口吻.
很多很多重複聽而學習的.
那種屬靈的套語.
但仍然是告訴我們.
我們要醒悟的就是.
基督教主耶穌在我們的生命裡.
到底帶來甚麼生命的改變.
我們的屬靈生命有沒有與日俱長.
是這個問題.
試圖有一些超自然的神蹟.
發生在他們的生命裡.
在傳道的過程.

$^{641}$到底維持了多少日.
這裡沒有詳細交代.
因為是一個很概括的描述.
即是當時有這樣的一個光景.
教會是一個好像很大的輕盛的時期.
這樣傳道.
你說多麼容易.
每個人自己來教會.
自己來到耶路撒冷.
在四處的城鎮.
自己來.
有人抬過來.
我走不了了.
不要理你.
你抬我去吧.
我見到彼得就行了.
甚至有些謠傳.
我見到他的影子.
我都會好的了.
果然可以.
你可以想像到.
神藉著使徒們.
做了這些事情.
但這些未必能夠.
一直這樣來長傳.
因為如果使徒們.
他們的內心的屬靈生命.
配合不了的時候.
其實這些可能是談一言.
掩眼法.
弟兄姊妹.
在人生信了主之後.
到底我們追求的是.
神給我們的恩賜.
甚至越超自然的恩賜越好.
還是我們追求這位愛我們的神.
他在我們的生命裡面的改變.
為我們內心裡面最重要的一部分.
我們追求的是.
很多很多能夠讓人看到我們屬靈的恩賜.

$^{681}$例如你的說話是很有才華的.
很有說服力的.
你的唱詩歌是非常動聽的.
可以令人對你有很大的那種欣賞.
這些本身不是錯的.
是恩賜來的.
但是我們在這些追求的恩賜裡面.
有沒有忘記我們自己內心的裝備呢.
我期望敢宣會一位.
參與事奉的弟兄姊妹.
我們是不是自己先屬靈吃飽.
然後去服侍呢.
我們有多少天沒有靈修呢.
我們沒有調查這些東西的.
我們沒有問.
你今天回來領唱嗎.
你有沒有靈修今早.
今早沒有不要緊.
因為太忙太趕.
早堂都要.
但是昨天上個禮拜上個月都沒有.
這就是問題了.
即是說我們內心原來是想靠超自然的屬靈.
那就不耐煩了.
不能夠內心是真正的心悅誠服.
被救恩被聖靈來自服.
沒有了一個對聖靈的馴服的心.
因為下面就出現大事了.
事奉神的人有時是可以.
很精彩的.
這麼大的神蹟.
我的影子可以救人.
不是吧.
我是彼得來的.
我不是什麼神醫.
但是第十七節.
你不要忘記這個世界有撒旦.
這個世界是有妒忌的人.
我們要預了事奉神的人.
你要預了你是加入主耶穌的軍隊裡面.

$^{721}$有人會批評我的嗎.
我不做了.
其實早就有了.
於是大祭司採取行動.
他和他所有一起的人.
就是撒都該派的人.
撒都該是不信復活.
這個世界是有人蓄意敵對神.
你全陽有神.
有耶穌復活.
你直接挑戰他的信念.
滿心忌恨教會.
千萬不要落在這樣的分派.
有這樣的對神工作的忌恨.
你看到別人被神大用.
你應該心裡面拍手掌.
很感恩.
千萬不要有一絲的忌恨.
我們要學習.
當然這些是敵對勢力.
更加嚴重.
下手拿著使徒.
把他們押在公共拘留所.
但在夜間主的使者開監門.
領他們出來說.
你們去站在聖殿裡.
把這生命的道.
一切的話講給百姓聽.
首先就是天是夜媽媽開了監門.
放了這班使徒.
一晚的時間.
到底有多少個鐘呢.
是十二點鐘凌晨才釋放.
還是很早的晚上就釋放呢.
還要是公共的外監.
不是內監.
很多人看到.
那些兵兵都嚴嚴看守住.
但是使徒們有一晚時間去想.
哇 傳揚耶穌.

$^{761}$我們又能夠行神蹟.
我們做那麼多好事.
幫那麼多耶路撒冷人.
我們都要遭遇和耶穌一樣的受逼迫.
會不會很快就要被釘十字架呢.
相信他們有思想.
有祈禱.
一晚時間.
他們要怎樣呢.
因為天使給他們的吩咐就很清楚.
告訴他們你們去在聖殿站著.
將你們傳講的生命之道.
去告訴百姓.
不是叫他們去行神蹟.
而是叫他們傳遞生命之道.
這裡有幾樣東西.
天使為什麼不自己去講呢.
能夠可以如入無人之境.
解除他們的武裝.
讓他們可以釋放.
那些鎖好像完全不需要開的.
都可以就這樣出去.
為什麼天使不自己向百姓顯現呢.
意味著一而再新約告訴我們.
天使只是報信.
是我們這些經驗過救恩的弟兄姊妹.
不同的民族.
我們要負起傳福音的責任.
天使是沒有這個責任.
他是沒有被授權去直接向人顯現.
他向使徒顯現就可以.
向其他的人.
未信的人.
他沒有權授給他去顯現.
如果顯現就很多人信.
不是的.
另外就是神會負責.
因為使徒們不是遠走高飛.
說去英國還是去加拿大好呢.
還是去澳洲還是去荷蘭呢.

$^{801}$荷蘭原來沒有人移民.
因為荷蘭話很難學的.
還是怎樣呢.
那個問題很清楚.
在耶路撒冷為我作見證.
他們在聖殿.
是不是在所羅門的廊下.
來傳講這個道呢.
說這個是要很大的勇氣.
但是神會負責任.
你只要順從聖靈的指引.
去講解這個真理的道.
就是能夠救人的生命的福音.
所以後面很諷刺的就是.
那些祭司長.
其中一個指控是什麼呢.
你們令到耶路撒冷.
充滿了你們的聲音.
很特別的一個指控.
這個豈不是耶穌在升天之前.
要十一個使徒要做的事情嗎.
這裡就告訴我們.
應該是第幾節呢.
第28節.
看啊 你們倒把你們的道理.
充滿了耶路撒冷了.
巴不得這樣宣家.
我們現在在加州花園.
洪水橋.
在元朗的宣德堂.
再加上其他的宣道會.
我們有110多間的宣道會.
巴不得這個道理.
是充滿了整個新界西區.
每個人都知道.
這位主耶穌.
你相信祂的生命會得到改變.
這是一個救人的福音.
盼望我們帶著這個訊息.
在我們的棚背.

$^{841}$在我們的工作的地方.
來傳揚.
以致真的可以拯救很多人的生命.
親愛的弟兄姊妹.
來到第21節的時候.
試圖聽了這些說話.
天將亮的時候.
他們就跟隨.
內心聖靈的引導.
讓他們有一個.
戰戰兢兢的勇氣.
去傳遞這個真理.
這個需要很大的.
屬靈生命去承載.
不是一下子.
一瞬間的.
一剎那的虛榮.
那些的扮屬靈.
是需要勇氣的.
盼望我們的弟兄姊妹.
信主.
無論你從什麼入門都好.
你要記住.
是福音的真理.
耶穌的生命.
承載你和我的生命.
這個是很重要.
所以他們二話不說.
其實整晚應該睡不著了.
給我一定睡不著了.
要再入聖殿.
再要面對大祭司的逼迫.
其實是要想清楚.
這個代價的.
他們就繼續這樣做.
果然.
祭司長打算去入監裡面.
去將他們提出監.
準備審問.
準備行刑.

$^{881}$你可以想像得到.
就是這樣的結局發生.
不過.
聖經告訴我們.
這些人見到監控完之後.
就回來回報.
然後再去的人.
因為他們已經在聖殿.
再去的人就叫他們.
去到審判的地方.
是邀請他們去.
不是捉拿他們去.
因為當時的百姓.
是來認同這些使徒.
他們的生命力.
他們的內心有一個很強的.
生命之道的力量.
第26節說.
於是守殿官和差役.
去帶使徒來.
並沒有用暴力.
因為怕百姓用石頭打他們.
輾轉最後就是這一段的聖經.
告訴我們.
我們豈不是嚴嚴的禁止嗎.
接著就是彼得這番話.
彼得和眾使徒就回答.
我們是必須順從神.
勝於順從人.
第32節.
我們這些事情的見證人.
和神賜給順從的人的聖靈.
也都會為這些事情作見證.
弟兄姊妹.
這些事情.
就是你們掛在木頭上.
殺害的耶穌.
我們祖宗的神.
已經使祂復活了.
神將祂高舉在自己的右邊.

$^{921}$所以今天我們傳揚的是.
這一位復活.
將來第二次要回來這個地球.
拯救世人.
拯救教會的元首.
基督.
我們一起禱告.
親愛的主.
盼望教會.
能夠在章林期的第一個禮拜.
也是教會年.
一個新年的開始.
我們很盼望.
我們每一位弟兄姊妹.
都能夠心靈明亮.
對真理有認識.
有降伏.
而且是我們的內心.
在神的真道的練靜.
那個栽培.
來成長.
我們做一個真心實意的基督徒.
一起在神的家裡面.
同心的服事.
願主幫助.
也都除去我們一些錯誤的神學觀念.
一些可能一夜之間.
就可以變屬靈的成功神學.
求主幫助我們除開.
使我們很踏實的.
一起.
孫家的弟兄姊妹.
來天天的靈修.
在神的話語裡面裝備.
我們這樣禱告.
奉耶穌基督的命.
阿們.
\newpage



\section{撒迦利亞書}
\label{sec:Qx_35UZ1A44}
\textbf{2023-12-07 周四晚查經班《歸回先知書研讀(9)撒迦利亞書》 講員:曾裔貴牧師}
\newline
\newline
連結: \href{https://youtube.com/watch?v=Qx-35UZ1A44}{\texttt{https://youtube.com/watch?v=Qx-35UZ1A44}} ~~~~ 語音日期: 2023-12-08
\newline
\newline
\hyperref[sec:7XxklgVuie4]{\small{< < < PREV SERMON < < <}}
~
\hyperref[sec:index_chronic]{\small{[返順時目]}}
~
\hyperref[sec:index_scriptual]{\small{[返順卷目]}}
~
\hyperref[sec:_n0nRPZl6sE]{\small{> > > NEXT SERMON > > >}}
\newline
\newline
$^{1}$好,弟兄姊妹,我們停了幾次.
今晚能夠有身體健康,可以恢復.
跟大家一起讀聖經.
研究歸回的先知書.
一個簡單祈禱.
再一次多謝天父.
因為身體健康不是必然的.
有時可能很少的疾病.
都會令我們說不出話.
或者令我們不舒服.
多謝你讓我們有身體健康.
也讓我們能夠有一個渴求.
對聖經話語的明白.
我們知道舊約的先知書.
去到撒迦利亞書.
後面這六章有很多對將來.
以色列國度的預言.
有些我們都很難去理解.
還有那個時間到底是何時會發生.
這些都超越了我們能夠預知的能力.
但我們盼望從這六章裡面.
給我們一個準確的大方向和原則.
就是上主尤其是彌賽亞降臨.
是非常重要的一個主題的時候.
都讓我們認真學悟.
舊主的降臨.
再一次回來這個地上.
來拯救世界.
求主帶領我們今天晚上的研讀.
我們一起禱告奉耶穌基督的名.
阿們.
這樣我們能夠去到最後這六章.
我們間間竭竭地說過第九第十章.
十一章有少少輪廓都有說過的.
有一個比較特別.
這六章聖經是指著將來的預言.
剛才在祈禱裡面都提到.
無論第九章第十一章.
都有提到彌賽亞.
和新約的福音書有引述.

$^{41}$直接說到彌賽亞的受苦.
甚至乎是字面的應驗.
例如十二章提到.
他們會仰望這位被砸的人.
其實就是約翰福音第十九章就引述.
主耶穌就是被人刺穿他的肋旁.
有血水流出來.
就是那一節聖經的引述.
所以你可以想像到這是十二章的第十節.
就是他們會仰望這位他們所砸的.
必定為我悲哀等等.
你可以想像到這六章最後的預言.
有別於之前的那些比喻.
是說到將來以色列國度的預言.
第九章很快我們曾經說過.
就是神會再回來.
使列邦.
當然是周邊的國家.
例如約旦,敘利亞,托羅西頓等等的周邊國家.
就會來受到懲罰.
神就要出手了.
拯救祂自己的子民.
這個主要是第九章.
其中就夾雜了.
有一位謙卑的君王再出現在這個地上.
就是第九章九節.
是很熟悉的.
就是福音書引述的主耶穌騎著驢.
驅進城.
這個就是受難周一開始的時間.
「錫安的民,九章九節,應當大大的喜樂.
耶路撒冷的民,應當歡呼.
看!你的王來到你這裡.
他是公義的.
並且施行拯救.
謙謙和和的騎著驢.
就是騎著驢的驅子」.
原原本本的按字面應驗.
福音書.
真是耶穌很特別地叫那些門徒.

$^{81}$你去到那個地方.
跟那些人說.
要用這個驢仔.
那個人就會給他.
好像不是很尋常.
但是就按字面應驗.
耶穌得到驢驢之後.
就騎著驢驢仔進城.
那些周圍的百姓.
因為是首節.
準備上首逾越節.
就以為是大軍進城.
就是軍王.
大家鋪了一些中樹枝.
一些衣服.
來到高聲唱和珊娜和珊娜.
至高之處榮耀歸給神的那首詩歌.
來歡迎彌賽亞.
他們認為的彌賽亞進城.
但是後來他們發現.
主耶穌不是這樣榮耀的德性.
反而是很羞辱的落敗.
釘十字架.
他們就徹底的掉眼鏡.
但是這裡是在說彌賽亞降臨.
所以我們上一次.
其實上兩次還是什麼.
第十節是沒有引導的.
福音書的作者知道.
是有兩個階段.
所以第一次的再來.
騎著驢的進城.
就是講到這位謙和的軍王.
第十節就是.
我必除滅以法連的戰車.
和耶路撒冷的戰馬.
徵戰的公也必除滅.
他必向列國講和平.
他的權柄必從這海管到那海.
從大河管到地極.

$^{121}$所以有一種這樣的.
漸進式的.
局部式的顯現.
的應驗.
所以有時候讀這些的寓言.
有它的難的地方就是這樣.
因為原來兩節之間.
是有一個逢曲的.
可能這個逢曲.
是多少千年是沒有人知道的.
耶穌他第一次來.
他們知道了.
門徒他們後來寫信的時候.
就能夠找到耶穌.
那個進城.
原來就是應驗撒迦利亞書九章.
但是.
福音書的作者就很巧妙地.
避開第十節未應驗的.
因為耶穌不是用這樣的方式.
來去.
好像是.
冰林城下.
徹底的得勝的.
所以.
福音書是全部沒有引到第十節.
但是舊約的寓言.
就會是比較.
原原本本的說了.
但是他自己.
寓言的那些先知.
他們自己都不知道.
到底那個時序.
是要怎麼發生的.
所以有時候我們讀這些寓言.
就要有這個觸覺.
感恩的就是新約的作者.
幫我們釐清了.
這些問題.
他們幫我們.

$^{161}$做了這些詮釋.
所以我們就知道.
原來福音書的作者.
都不是亂抄的.
他們抄下來.
他們記錄那些事件.
他們又會知道.
在哪裡應驗了.
那個應驗原來只不過是一個.
局部性的應驗.
就不是全面的了.
所以.
我們看往後的那五章聖經.
因為第十二章又會有重複.
又是說神怎樣來拯救.
他自己的子民.
不過又會加多了一些元素.
因為純粹是軍事上.
政治上的勝利.
其實沒有真正改變到.
百姓的心靈.
所以.
需要有另外一個傳言.
一個洗罪的傳言.
十二章.
這些經文又是很出名的.
十二章第十節這裡.
又是有引述的.
第十章.
九章我們當說完了.
對列國的審判.
和允許這位謙和的君王.
進城.
來施行一個拯救.
和一個審判.
當然那個審判.
還不是現在馬上應驗.
就是他再回來的時候.
就會以這個形態.
來做一個全地的審判.

$^{201}$當然新約也有這樣的預言.
主耶穌再回來的時候.
就是改變一切.
第十章.
也告訴我們.
神會藉著.
他自己.
來興起的那個牧者.
來將以色列.
親自的去牧養.
可能讀回一點點.
當春雨的時候.
你們要向發閃電的耶和華求雨.
他必為眾人降下金林.
使道田園生長菜蔬.
因為家臣所言的虛空.
卜士即是占卜那些.
所見的是虛假.
作夢者所說的是假夢.
他們白白的安慰人.
所以眾人如羊流離.
因為沒有神的道.
沒有方向.
沒有了那個安慰.
沒有了那個指引.
自然就會流離.
如牧人沒有牧人群羊的苦楚.
這樣描述.
因為當時他們都是尋求.
似乎是在說被擄之前.
以色列一貫的形態.
有先知但不聽先知的講道.
情願傾向跟隨那些迦南人.
拜偶像.
拜不同的假神.
甚至興起一些.
幫到他們的.
可以滿足百姓內心的訴求.
例如你想我要求雨.
我就找一些假先知來占卜.

$^{241}$或者用一些不同的法事.
使你們可以得如所願.
但這些不是神的先知.
更加不是百姓的一些牧者.
所以最終會令群羊流離.
其實屢見不鮮.
如果教會沒有了神的道.
弟兄姊妹自然就會感覺到.
缺乏慰養.
沒有神的道的指引安慰.
和督責教訓.
自然就會慢慢對自己的生命.
屬靈生命會有一種停滯.
和失去方向.
失去方向仍然要生活.
就會專注在生活.
多於屬靈的生命追求.
所以這些都可以完全對號入座.
任何一個教會任何一個時代.
如果沒有道.
沒有牧者的牧羊.
就一定會有一種流離感.
就會有一種困苦.
因為沒有牧人而困苦.
所以很自然.
往往先知書都是神透過先知.
來譴責,攻擊那個時代的領袖.
因為某程度你是負起了屬靈的帶領.
你要作為領袖.
你就有帶領的責任.
如果沒有神就會譴責.
所以第十章的三節.
我的怒氣向牧人發作.
我必定懲罰公山羊.
不是綿羊 是公山羊.
我們也說過.
馬太坊的二十五章.
是說綿羊,山羊.
山羊是不好的.
不知道為什麼在主耶穌.

$^{281}$說比喻也要說.
山羊就是那些不得救的.
沉淪的.
綿羊是好的.
進入永生的.
是不是在他們猶太人的傳統.
山羊就代表不好.
綿羊就代表好的.
其實兩者都是不同的品種.
所以似乎大家對這些典故都要認識.
否則新約的比喻就不夠立體.
你就覺得很奇怪.
什麼是山羊 什麼是綿羊.
如果我們讀完.
可以揭一揭二十五章的比喻.
讓大家可以看看.
這裡說神的怒氣向牧人發作.
然後懲罰公山羊.
因我萬軍之耶和華.
眷顧自己的羊群.
就是由大家必使到他們如駿馬在陣上.
然後提到房國石釘子增薦的弓.
無論是穩固的.
房國石是穩固的.
穩固的根基.
釘子就是釘穩某些東西.
那個增薦的弓是對外的.
好像有一個攻擊性的.
不只是穩固的.
這樣來描述.
在他神自己親自的牧羊底下.
由大家就會有這種形態.
即是由屬靈根基穩固.
又能夠釘穩像釘子一樣.
而且穩固在神的應許.
那個藥 摩西之藥.
又能夠好像是神手中的弓.
這樣射出箭.
而且一切掌權都是由祂而出.
他們必如勇士在陣上.

$^{321}$將仇敵踐踏在街上的泥土中.
他們必爭戰.
因為耶和華與他們同在.
騎馬的也必定羞愧.
這樣描述在神的牧羊底下.
和上文一至到第二節.
沒有牧人的情況底下.
大家傾向.
用我們現在的說法.
看看什麼星座.
什麼雲層紫微斗數.
什麼掌紋.
甚至現在有什麼手指的紋.
我不太清楚.
很多千奇百怪的.
水晶球.
其他的甚至有些再嚴重的.
真的去拜一下那些廟.
我們那裡有個天后廟.
大堂.
會不會有弟兄姊妹去拜呢.
不知道.
沒有見過.
會不會有其他的東西呢.
這些就是所謂的占卜家臣.
家臣就是自己家裡有自己的公仔去拜.
有些作夢.
有些異夢等等.
但是原來在神的牧羊裡面.
就會帶來由大家有那種爭戰的能力.
和有一個穩固的根基信仰.
所以神親自的牧羊.
神興起的牧者.
如果能夠恰如其分.
在那個牧羊裡面.
大家就會專注在神的道.
和神應許的將來的.
種種的屬靈的福氣.
大家就會有一個羨慕.
例如我們會一起去傳福音.

$^{361}$會讓人看到福音的大能.
改變生命的能力.
和我們會經驗到.
因為看到人改變.
生命改變.
這些就是在牧羊裡面帶來的.
原來你不知不覺.
你是房角石.
你是釘子.
你是爭戰的弓.
有時候拆開講一些例子.
偶然的.
好像是上個星期.
突然間收到很長的WhatsApp.
他就把一些星期四晚的茶經班.
給我看.
我其實不認識他.
他說來過我們錦宣家聚會.
接著就回到美國.
他的妹妹就患重病.
就很虛弱.
他就不知道什麼原因.
就給那些星期四晚的茶經班的連結.
他去看.
去聽.
應該是聽.
應該是沒得看.
一邊聽一邊聽.
他說他的心靈有很大的改變.
還有他本來是很嚴重的情緒.
因為那個病很重.
聽了一整年.
就說他的妹妹現在很開心很積極.
還要給他的家人聽.
是收到WhatsApp.
其實我真的不認識這位姐妹.
是他的妹妹在美國.
就很感恩.
原來其實我完全不知道.
我說的話其實他可能聽不明白.

$^{401}$很深的舊約.
他說聽不明白.
但他一邊聽的時候.
他內心是得到安慰.
甚至現在聽了主日講道.
又把主日講道的連結給我看.
說哪一分鐘叫家人聽.
比我更加仔細.
這個很重要.
神親自去牧羊的時候.
當他去眷顧自己的羊群.
就去遊大家.
其實鬧到傻的上文.
一到第二節.
但原來當神要再眷顧自己的子民.
其實這個時候寫作動機就是在重建聖殿.
本來荒廢了爛了尾的.
現在因為撒迦利亞和哈該.
兩位先知被神呼召去勸勉他們.
其中當然是所羅巴伯和大祭司約書亞.
在這個時候就是大利烏王的第二年9月.
所以這個時候正正就是發禦言的時候.
就是日頭大興土木的時候.
重建聖殿的時候.
所以原來建好的硬件沒有用.
原來那班人將來會有一個這樣的亮麗的前景.
就是神親自牧羊.
他們就會成為了防國石.
成為了釘子.
成為了增建的弓.
在神手中很有用.
你想想在古代.
你拿著一支弓.
有整袋箭.
你就很有攻擊性.
你就不是挨打了.
你想想有些箭可以射敵人.
古代不是用槍.
你可以想像到神親自會眷顧自己的子民.
所以這些預言其實對撒迦利亞.

$^{441}$對所羅巴伯.
對任何的領袖.
都是一個很大的安慰.
所以盼望我們能夠在神的話語上.
默默地去研讀.
我們自己內心得到溫暖.
得到建立.
得到安慰.
我們就能夠成為了神增建的弓.
當然不是去殺人.
是將人帶進基督的愛裡面.
我們自己感受到愛.
我們自己有這份在主耶穌裡面的幸福榮耀.
我們就願意和別人分享.
不再害怕.
不再拍人笑.
而是很理直氣壯.
我可以很嚴肅地告訴你.
你真的不相信嗎.
你要面對神.
所以這裡很值得我們去再看.
所以第九章是對列國的審判.
其中就是這位謙和的彌賽亞第一次降臨.
所以任何一個領袖必須要學習主耶穌.
我們是用謙和去得到人的心靈.
不是用武力.
所以我絕對不贊成現在的以巴戰爭.
那種一面倒的狂軍猛炸.
無論是不是神的選民也好.
這肯定不是神的心意要這樣做.
當然我們不要拆開講了太多的政治.
接著就是第六節.
我要堅固猶大家.
拯救弱失家.
要令他們歸回.
這已經歸回了.
現在將來是有一個全面的歸回.
我要連戌他們.
他們必像未曾棄住的一樣.
從來好像沒有罰過.

$^{481}$你可以想像神那種既往不再追究.
那種愛.
那個公義.
不會有任何的陰影.
你曾經說過這句話.
我永遠都記得.
神是不會的.
好像令到這個民族.
這個神的選民.
好像未曾棄絕.
其實被擄歸回就是70年.
70年要流亡去到巴比倫.
後來就要去到馬代波斯.
接著波斯王的古獵才回歸.
現在是回來了十幾年.
所以說怎麼可能未曾棄絕呢.
真是神要使自己的聖殿.
所羅門的聖殿.
完全被拆毀.
那些聖殿裡面的器皿.
施奉神.
用來敬拜.
用來祭祀神的器皿.
全部偷走了.
去到巴比倫.
你可以想像到.
其實是一個很大的羞辱.
但是神說.
那個愛.
那個救恩.
大到一個地步.
蓋過了從前的責罰.
讓你可以感受到.
神的大愛.
對一個罪人的寬恕.
我們時常都喜歡用浪子的比喻.
那個爸爸.
看到浪子回來.
根本都不會介意他很髒.
他之前在豬場裡面工作.

$^{521}$一股臭味.
衣衫藍柳.
都完全不介意.
因為那份負的愛.
大到一個地步.
足以彌補.
將一件袍子.
加在他身上.
又給他戴戒指.
你可以想像到.
神有這樣的應許.
我想我們每個人.
在主耶穌的愛的更新裡面.
在祂的眼中.
好像我們.
每一層氣訣一樣.
我們都應該要.
這樣去經驗神的愛.
挽回的恩典.
都因我是耶和華他們的神.
我必應允他們的禱告.
以法連人.
必如勇士.
他們心中暢快如同喝酒.
他們的兒女.
必漢見而快活.
他們的心.
必因耶和華喜樂.
我們可以想像到.
當神親自牧羊.
神的時間到.
祂要全面地使他們歸回.
我們不知道.
因為這個寓言.
沒有特別.
若有所指.
那個時序.
時間的次序.
哪年哪月哪日應驗.
我們不可以這樣對號入座.

$^{561}$但是將來的國度.
以色列就要這樣.
得到復興.
復興就是.
神對子民的愛.
如同將來的光景.
如同未曾氣絕一樣.
好像.
大家都記得.
這個弱失.
未來到.
後來被法老召見入宮.
解了法老的夢.
告訴法老.
這個夢.
原來是有解釋的.
首先會有.
七個豐年.
在這七個豐年裡.
要盡量囤積.
要盡量.
收購.
因為有緊接著的.
七個荒年.
寸草不生.
所以.
後來真是這樣描述.
去到荒年的時候.
那些人荒到一個地步.
完全忘記了豐年.
其實我們都會.
大概明白.
如果今天很冷.
冷到只有十度.
我們完全在內心.
忘記了夏天三十四度的時候.
那種情況.
人都會上升.
何況神告訴我們.
將來會有.

$^{601}$這樣的一個景況.
就是當他們完全.
回到地老歸回的時候.
或者完全回到地老的時候.
神的連率會令到他們.
好像未曾.
被神棄絕一樣.
這是很大的福氣.
接著描述那些光景.
以法蓮.
大部分北方的十國之派.
統稱.
他們就好像勇士.
他們的內心暢快.
好像喝完很純的酒.
當然不是.
鼓勵喝酒.
古代他們沒有那麼多娛樂.
一定是飲食.
如果能夠喝到.
很純的酒.
就心中暢快.
我們自己想我們的心中暢快.
我們認為什麼是最令你開心的.
最令你心滿意足的.
那個意思.
講的是整個民族.
不是個別一個人.
有些人很容易滿足.
有碗糖水就已經很滿足.
有些人就不是.
所以就不是個人.
是整個民族.
因為萬君之主.
將他們全面歸回.
將來有第二次真正的大歸回.
但現在這本書本身就是歸回的先知書.
在這個時候正在重建聖殿的時候.
所以這樣講的寓言.
其實很有壓力.

$^{641}$其實這樣講的寓言.
其實就很合時機.
對建造的所羅巴伯和約書亞.
有很大的推動和鼓舞的作用.
不是要來到.
好像只是拍馬屁.
不是要來到.
好像只是拍馬屁.
做吧做吧.
而是有神的應許在裡面.
當然不是當時應驗.
全面的應驗.
只是很局部的應驗.
所以你可以想像.
經過二千年.
到今天以色列都未真正享受到這種平安.
你會想像得到.
這是必須要有那位彌賽亞.
第二次回來.
才能夠成就.
後面會有些講到.
真的講到關於彌賽亞的二次降臨.
輾轉提到.
這個民族得到的復興.
當然相對其他民族.
會受到審判.
或者我們簡單讀一讀.
就會進入第十一章.
就是我們很簡單.
上一次講過.
好像有一個戲劇.
撒迦利亞.
要扮演一個好牧人.
和一個壞牧人的角色.
到底是一個寓言的講法.
還是一個戲劇.
因為有買賣雙方.
那些買家.
當然是列國的人.
想對付神的子民.

$^{681}$那些賣家.
那個場景.
我們等一下看十一章.
你就會知道.
我們讀第八節這一段.
我神.
要發師聲.
聚集他們.
因我已經救贖他們.
他們的人數必加增.
一直在講以色列的歸回.
如從前加增一樣.
我雖然播散他們在列國中.
他們必在遠方紀念我.
這個紀念就很考功夫了.
就是所羅門的祈禱.
幾時幾日.
當那些子民走投無路.
向聖殿祈禱的時候.
神啊.
求你在天上垂聽.
拯救他們.
路加十五章.
浪子是為什麼會回家的.
就是他突然有一天醒悟.
我得罪了天.
我得罪了地.
我得罪了你.
我本來是一個這麼有錢的人家.
為什麼我要在這裡.
做一個乞丐呢.
養豬呢.
他就醒悟過來.
決心想起台詞.
回去就跟爸爸講.
就是源於他的醒悟.
所以這個從遠方紀念神.
就很重要了.
這個遠方可以是.
心靈的深度.

$^{721}$有些人沉到.
沉到他不再回教會.
甚至聖經以前很敬重的.
很珍而重之的.
靈性不好.
聖經就說.
我被馬經蓋住也不怕.
也不介意.
拿來煎煲底也無所謂.
因為最好夠厚.
但是呢.
當他從遠方.
由他心靈的遠方.
紀念神.
這個動作很重要.
這個紀念.
是發自他內心深處的動機.
還有聖靈藉著不同的.
可能是他的子女.
令他突然間.
有一個醒悟.
或者用不同的方法.
尤其是信了主.
曾經經歷過耶穌基督的愛的弟兄姊妹.
這個遠方.
就憤外離得真實.
一個未信的人.
就沒有這種掙扎.
因為他根本沒有聖靈.
但是有聖靈的弟兄姊妹.
他會在遠方紀念神.
譬如嫁了一個未信的丈夫.
丈夫是很差.
攻擊信仰.
他沒辦法.
因為已經成為婚姻.
你都要忍氣吞聲.
甚至因為原因都不能回教會.
等等.
但不知什麼原因.

$^{761}$在他的遠方.
在他的心靈的遠方.
突然間什麼促使他.
這個一定是聖靈工作.
他來到紀念神.
這個動作是很強大的力量.
當他們在遠方紀念神.
當然這裡一定是在很遙遠的地理遙遠.
你想想天涯海角.
以色列人都可以歸回.
因為神自己會使他們歸回.
且得歸回.
沒錯.
我必在領他們出埃及地.
所以第二次出埃及.
就是說.
人不可能做到.
出埃及是十災.
是紅海分開.
是一個超自然的大能的拯救.
所以不是靠人力可以做到真正的歸回.
所以將來以色列一定不是靠軍事實力.
十二章又再提到這一點.
所以十二章洗罪的全完就很重要.
神會開一個全完.
這裡就有一個很亮麗的未來.
以色列全面歸回.
甚至神會再給他們更多的土地.
令他們激烈和泥巴亂.
這地尚且不夠他們居住.
耶和華必經過苦海擊打海浪.
駛到尼羅河的深處到附近.
第二次出埃及的描述.
亞述的驕傲必至卑微.
埃及的權柄必然滅沒.
沒有任何一個民族一個國家.
可以攔阻這一次真正的歸回.
甚至不只是人的歸回.
而是他們這一班的以色列人.
在心靈裡面想念神.

$^{801}$即是說有神的靈在他們裡面的感動.
這個就更加重要.
因為你個人去而他的內心都沒有變過.
根本不敬畏神就沒有意思了.
所以我必使到他們倚靠我得以堅固.
一舉一動必奉我的名.
這是耶和華說的.
其實很快和大家揭一揭.
耶利米蘇的三十二章.
有這些的預言.
其實就在.
希伯來書就已經有引述過了.
耶利米蘇.
我看一看是三十二還是三十三.
三十三才對不好意思.
三十三章.
耶和華說.
三十三章的十四節.
耶和華說日子將到.
我應許以色列家和猶大家的恩賢必然成就.
當那日子那時候.
我必使到大衛公義的苗裔長起來.
他必在地上施行公平和公義.
在那日子猶大必然得救.
耶路撒冷必安然居住.
他的名稱為耶和華我們的義.
其實還可以再揭多一點.
三十四章有沒有呢.
當然我引述希伯來書這段更加直接.
因為我暫時找不到.
剛才這個耶利米蘇那裡.
希伯來書引述的這段經文.
就是第八章有這樣的教導.
八章.
我揭什麼了.
因為這裡有提示.
原來是三十一章才對.
剛才三十三章是另一個彌賽亞的工作.
三十一章這一段.
這段聖經是重要的.

$^{841}$應該是三十一章三十一節.
耶利米蘇是很忌諱的.
一定要忌諱的.
因為希伯來書直接引述這段經文.
提到以色列或者神的選民.
將來會有一個這麼大的復興.
他說耶和華說日子將到.
我要與以色列家和猶太家另立一個新約.
即是說我們今天耶穌的救恩就是新約.
摩西和以色列人在西乃山立的就是救約.
所以我們是新約.
希伯來書是這樣肯定的.
所以那個新約的成就就是耶穌的救恩.
所以我們今天是以守聖餐來紀念這個新約.
有別於猶太人他們在西乃山的救約.
救約他們要立約的條款.
就是要行國禮和守安息日.
行國禮就是亞伯拉罕的子孫.
安息日就是摩西規定西乃山的猶太人他們要守.
但是我們新約就不是這樣了.
我們就是透過主耶穌的救恩.
憑信心就可以得到這個永遠的救贖.
留意一下這個新約在耶利米蘇的時候的預言.
因為裡面的內容是很重要的原則.
他說我要和他們立的約.
第32節.
不像我拉著他們祖宗的手.
就是西乃山.
領他們出埃及地的時候.
與他們所立的就是救約.
我需作他們的丈夫.
就是說神體.
祂自己是以色列文的丈夫.
就是說有一個恩典的約在其中.
中國人都說.
一夜夫妻.
就不是百二蚊.
是百夜恩.
你嗲一下都有一種很深的味道.
這個是丈夫.

$^{881}$神是以色列人的丈夫.
就是有一個不會改變的約.
對這個民族.
將他們興起.
揀選.
藉著摩西立了這個西乃山的約.
設立十誡.
我雖然而已.
但是他們卻背了我的約.
就是說原來十誡.
無論你寫得多漂亮.
那個約板怎樣精美.
那些人守不到.
沒有心去守.
沒有能力去守.
就變成背約了.
這個就是耶和華說的.
這次留意了.
33節.
耶和華說.
那些日子以後.
當然就是主耶穌來成立救恩的情況.
以色列所立的約乃是這樣.
新約的條款.
我要將我的律法放在他們裡面.
寫在他們心上.
那就是說.
以前就是寫在石板.
以前就是有一些手抄的卷軸.
但是.
白紙黑字都沒有用.
守不到.
但是當主耶穌成就救恩.
聖靈進入了我們每一個信徒的內心.
包括猶太的信徒.
神的道就寫在心板.
和在裡面.
那是什麼意思呢.
就是說我們能夠透過聖靈.
我們明白真理.

$^{921}$和我們能夠願意.
和有能力去實踐.
與神相和相交的真理.
以致我們和神有一個很密切的關係.
就不是一種抽離的.
是呀有十戒.
不過守不到.
所以這是新約很重要的精神.
聖靈在我們的心裡面.
才讓那些律法寫在我們的心板.
所以大家來聽道.
要帶著心靈.
尤其是查經班最好.
願者才來.
你不願你不用來.
崇拜就不是.
你願不願你都要來.
你不來更加顯得你不願意.
但不是三百人坐在那裡都是願的.
所以寫在心板的意思就是說.
整個的心靈是降伏在聖靈之下.
這裡告訴我們.
神說我要作他們的神.
他們要作我的子民.
就是說大家的關係.
是有一種通.
或者一種順.
或者是一種和諧.
這是很重要的.
你感受到和諧.
你才會很願意不斷去親近.
即是你和你的孫子.
和你的兒子.
你的女兒.
你都沒有話題.
最好是他自己搞定.
不用我煩.
否則你多想有時間陪他.
見他.
這就是和諧關係.

$^{961}$所以新約可以帶來.
我要作他們的神.
原來是雙向的.
他們要作我的子民.
是雙向的.
所以為什麼新約基督教.
很強調是雙向的.
一個神人相交.
不是一個單向的.
我們去黃大仙就是單向.
你靈驗.
我就給你大支香.
我就給多些香油錢.
那些幫手的人.
我就去還願.
不靈驗我就找別人.
但是神人相交.
新約就不是這樣.
他要作我們的神.
我們很主動要作他的子民.
你就可以想像到.
基督教那種其中美麗.
那個大的力量來源.
就是每一個弟兄姊妹.
有很強烈的意識.
我要成為神喜悅的兒女.
這個就是新約帶給我們的榮耀.
還沒完.
他們各人不再教導自己的香輪.
自己的弟兄.
你要認識神.
因為他們從最小到最大.
都必定會認識我.
認識就帶來關係的復和.
認識是一種不是頭腦上.
是有經歷的認識.
如果感受到親近神.
就會得到罪的赦免.
得到我們不會再記住自己的罪惡.
因為神都已經寬免了.

$^{1001}$這個就是新約.
其實剛才數手指都已經有四個條款.
新約舊約的限制.
其實就不是神的律法有什麼瑕疵.
是人做不到.
這個是羅馬書第七章.
這樣來解說.
但是新約就不同了.
主耶穌的救恩完成.
就將聖靈 神的靈.
進入我們每一個人的內心.
所以就令到我們與神復和.
與神和諧關係.
神人的相交.
是一直進行到永恆的天家.
這個就是為什麼這麼重要.
這個新約和有別於舊約的關係.
所以在撒迦利亞書的預言.
都提到將來就會有這個光景.
但是其實其他的先知書.
都有涵蓋將來新約帶來的復興.
所以耶利米蘇就更加詳盡.
來描繪這個新約.
你想想耶利米蘇.
在國破家亡的時候.
來寫的一卷書.
他自己也是被人迫殺的光景.
被人放進去那個枯井裡面.
來自生自滅.
就是半監禁的.
這樣的狀態.
但是神給他的預言就是說.
將來是會有新約的出現.
而新約的那個執行人就是耶穌.
所以我們回到第十章.
就告訴我們.
就有這個光景.
神就是要這樣來將他們的子民.
來拯救.
他們就得到這個復興.

$^{1041}$先休息一下.
我們就進入第十一章.
如果我們整章去看.
就很簡單.
但是逐句看就有點複雜.
好.
似乎有一個戲劇.
一個drama.
第十一章撒加尼亞書.
黎巴嫩啊.
開開你的門.
任火消滅你的香柏樹.
松樹啊應當哀號.
因為香柏樹.
香柏樹代表著那些高傲.
因為那些樹本身的體態就是這麼硬朗.
又皆美又高.
橡樹就象徵著它的堅韌.
其實描繪著不同的民族.
他們的特性.
都是不約而同地被毀壞.
因為神的審判度.
皆美的樹毀壞.
巴山的橡樹應當哀號.
因為茂盛的樹林已經倒了.
聽啊有牧人哀號的聲音.
因為榮華的草場毀壞.
有哨狀獅子咆哮的聲音.
因為約旦河旁的叢林荒廢.
就是說那些用樹.
不同的樹來描述周邊的國家.
面對神的審判.
他們的高傲.
他們的勇力.
像樹.
你都可以想像到.
那些橡木.
那些安木是很堅固的.
那些老鬼以前都是用這些木.
現在是用石.

$^{1081}$魚製的石.
一來就沒那麼多木.
你可以想像到.
古代如果能夠有這些樹林.
都是一個資產.
也可以誇耀.
但是.
這些全部要哀號.
全部要哭泣.
因為都是神的審判能力.
連帶那些動物.
因為要棲身.
你沒有那些樹林.
你沒有那些保護.
自然就會受到牽連.
所以神的審判.
有時就會令到不同的持份者.
都要面對審判帶來的後果.
後來下面提到.
突然間又加插了牧人.
因為其實在《先知書》.
牧人本身已經是整個民族.
百姓的一些屬靈領袖.
木樣的.
大家已經很清楚.
職份提到的屬靈領袖.
所以很多時候都會用這些做一個比喻.
或者甚至乎是一些直接的譴責.
讀完這裡會看到.
有一個很特別的.
因為和撒迦利亞有關.
他是好牧者.
後來要扮演一個角色.
做一個壞的牧者.
不稱職之餘.
是要被人譴責的.
大家聽一聽.
第四節.
耶和華我的神如此說.
撒迦利亞要木樣.

$^{1121}$這些將要宰的群羊.
即是養大要去市場賣的.
不是全部羊都是.
有些羊是要養大.
來繼續繁殖的.
那些不會賣的.
也不會拿去宰的.
但是一般的就會拿去宰的.
其實我們那時候養豬也是.
有些某些豬是留下來.
將來要做配種的.
就不賣的.
因為它本身那些品種可能好.
這裡說撒迦利亞要做木樣.
是不是真的有做呢.
其實又不知道.
是一個戲劇去扮演角色.
還是實際有做呢.
賣他們的宰了他們.
以自己為母贅.
賣他們的說.
耶和華是應當稱重的.
因我成為父足.
因為他們販賣那些羊給那些買家.
因為我們風雨了.
販賣了那些羊.
表面看這個戲劇.
就好像買賣雙方有什麼問題呢.
買家盡力去養大.
等到那些張宰的羊我們買了.
如果用那個場景上下文.
其實買家就不是普通人.
是列國的那些惡勢力.
要把百姓當羊去宰殺.
買家就說自己沒有罪.
就呼之欲出了.
真真有詞.
還看到自己的豐富.
富足.
不理民間的疾苦困境.

$^{1161}$而他們覺得賣了這些百姓.
就是自己富有.
反而是神的祝福.
是恩典.
那就不是普通的買賣羊了.
因為如果羊比喻回去是百姓.
那就買賣雙方都有罪.
買的動機不良.
是要對付以色列人.
以為神放棄這個民族.
所以他們可以來到示意的攻擊.
但是賣家其實就是那些民間的領袖.
因為如果牧羊人是撒迦利亞.
其實看著雙方.
這樣殘忍對待神的子民.
如果你作為牧羊人.
其實你就很不開心.
你就會很憤怒.
也很無助.
所以繼續下去的時候.
他們就應當稱頌.
因為他們富足.
接著牧羊他們的並不連戍他們.
接著第六節.
就是說有其他的牧羊人.
都是隻眼開隻眼閉.
見到都出不到聲.
第六節.
耶和華說.
我不再連戍這地的居民.
必將這民交給各人的鄰舍.
和他們王的手中.
他們必毀滅這地.
就是那些買家.
我也不救這民脫離他們的手.
於是我牧養者將宰的群羊.
到底他真的有做.
還是在戲中的角色呢.
就是群中最困苦的羊.
我拿著兩斤杖.

$^{1201}$一斤我稱為榮美.
一斤我稱為聯索.
這樣我牧養了群羊.
一月之內我除滅了三個牧人.
就是不稱職的.
只看著羊面對困苦.
面對買賣雙方.
買就是敵人.
賣的就是民間的以色列領袖.
都是對視群羊如草.
如地底泥.
根本就不會珍惜.
這些牧羊人都被撒迦利亞.
像炒魷魚般解除.
我除滅三個牧人.
因為我的心厭煩他們.
他們的心也憎嫌我.
我就說我不牧養你們.
要死的由他死.
要喪亡的由他喪亡.
餘剩的由他們彼此相食.
我截斷那稱為榮美的杖.
表明我廢棄與萬民所立的約.
當日就廢棄了.
這樣那些仰望我的困苦羊.
就知道所說的是ye 和華的話.
我對他們說.
你們若以為美就給我功架.
因為他是牧羊人.
好牧人為羊捨命.
有多少功架呢.
遣散費.
就給三十塊錢你.
你們如果以為美的就給我功架.
不然就罷了.
於是他們給了三十塊錢.
作為我的功架.
這是約翰福音十九章.
那個的.
猶大出賣耶穌.

$^{1241}$接著ye 和華吩咐我說.
我要把眾人所估定美好的價值.
雕給瑤瑚.
我便將這三十塊錢.
在ye 和華的殿中雕給瑤瑚了.
我又截斷稱為聯索的那根杖.
表明我廢棄猶大與以色列弟兄的情誼.
ye 和華又吩咐我說.
這是下一個扮演的角色.
這裡如果整段去看.
一方面神如果要來對付周邊的列國.
周邊的列國就不能再依靠自己的.
可能是聰明智慧.
可能是地理的險要.
可能是有很多的謀事.
可能是有很多的資源.
一切都會化為烏有.
因為神的攻擊只會連根拔起.
會令到攻擊以色列的.
就要面對神的審判.
但是後來神就好像透過撒迦利亞.
做一個話劇.
就是現在有一班群羊.
好像現在在牧場裡面.
有一些群羊有買家來買.
你要養肥養白牠.
讓牠們可以出欄去屠宰.
賣給那些買家.
那些買家神就說他們是不會顧惜那些百姓的.
神都會將這些百姓交給這些買家.
就是侵奪他們的姓名的.
這些列邦的外族的人.
所以買家其實就是不懷好意.
接著有賣家.
賣家就更加.
你可以說是他更加無恥.
原來我賣了我有價值.
就是說我根本不理那些百姓的死活.
我賣了我就有價值.
你可以想像得到.

$^{1281}$賣了那些百姓.
可能真的有這樣的情況也說不定.
所以作為領袖.
其實要面對我們那個中心.
是不是牧羊神的群羊.
因為這個是很反映到群羊的那個肅靈光景.
這裡就告訴我們.
撒迦利亞就兩難.
那些買家不是好人.
賣家都不理那些百姓的死活.
自己自肥就行了.
但是如果你是撒迦利亞.
神就說你要扮演一個好牧人的角色.
你要牧羊.
但是牧羊面對很多的阻力.
其他的牧人其他的童工阻撓他.
他就遣散除滅了這三個牧人.
他就厭煩那些牧人.
那些牧人也都憎恨他.
撒迦利亞是不是這樣就會得到百姓的歡心呢.
或者百姓對他的牧羊的認同呢.
不是的.
他帶來的就是.
他最終要說這些好像劈炮不撈的說話.
我不再牧養你們了.
要死的由他死.
要喪亡的由他喪亡.
餘姓的由他彼此相食.
一個好牧人說這樣的說話.
記住是一個戲劇.
是一個話劇.
甚至乎應該都是一些預言.
所以不是撒迦利亞真的這樣.
而是他的內心.
好像已經去到一種沒什麼可以做到.
百姓不介意被賣.
百姓甚至乎落到賣家的手.
只會埋怨神.
埋怨神沒有拯救.
那你就可以想像得到.

$^{1321}$撒迦利亞始終都要面對這種灰心喪志.
還好的就是一個話劇.
這樣做給以色列人看.
就是說原來百姓可以硬心到.
就算交付給敵人.
他們都好像仍然硬心腸.
是不會去思念好牧人的牧羊.
去向神有一個迴轉等等.
這裡就告訴我們.
撒迦利亞去到一個動作.
就是他截斷兩斤的象.
即是說那些象是幫他去打那些白艷的羊.
或者勾回那些掉進坑裡的羊.
這些象成為他的得力助手.
也是一個標誌.
一個牧羊人應該是拿著牧羊棍.
有名的.
聯索和榮美.
他也要截斷.
有解釋.
第十一節是這樣說.
十節應該說.
我截斷稱為榮美的象.
表明我廢棄與萬民所立的約.
當日就廢棄了.
是什麼意思其實很難解釋.
因為突然間說萬民的約.
其實是說以色列百姓被賣.
怎樣理解這件事.
可能是似乎在這個時刻.
好牧人.
撒迦利亞被自己的同胞百姓棄絕.
又要面對買賣雙方對他的排擠.
就是說他盡心盡力的牧羊.
得不到任何的回報.
就是說如果百姓沒有真正的心靈的悔改.
耶穌就算走幾十倍的神蹟.
也不會改變百姓的硬心.
果然是耶穌真的這樣走了這麼多神蹟.
他們明明知道.

$^{1361}$但是只要你不符合我們認為的彌賽亞的標準.
你不是這個標準.
你就不是彌賽亞了.
就變成他們自己製造他們自己心裡的彌賽亞.
所以耶穌就不會再是其他人了.
只不過是木匠的兒子.
那就應該要釘十字架了.
因為你說自己是神的兒子.
所以你就可以想像得到.
撒迦利亞最終神安排他在這裡做一個話劇.
就是我要離職了.
你們給我一個公價吧.
三十塊錢.
很特別就是真的按字面這樣應驗.
我們揭一揭去.
應該有人方便有的了.
(揭曉).
由大出賣耶穌那裡.
(揭曉).
我稍微揭一段.
(揭曉).
十二章就是說到他離開.
進入黑暗裡面.
可能其他的福音書會比較有描述一點.
《小徒行傳》一章就有說到.
由大出賣的價錢.
然後他就去他自己的道路上了.
一章的十八節就提到由大了.
(揭曉).
這個就是他的公價了.
如果我們可以準確一點.
馬太的.
我看看是哪一章.
我們揭開由大出賣耶穌這一段.
如果大家找到可以幫幫忙.
看看誰可以快一點.
(揭曉).
十四節當下十二門徒.
有一個稱為加略人由大.
去見濟斯長說.

$^{1401}$我把他交給你們.
你們願意給我多少錢.
他們就給了他三十塊錢.
從那時候他就找機會.
要把耶穌交給他們.
是這麼巧合.
還是要真的這樣對號入座.
彌賽亞現在撒迦利亞就扮演著.
這個好牧人的角色.
而受到自己百姓集體的排擠攻擊.
就是說在這個話劇裡面的這一幕這一場.
就是主耶穌被自己的百姓氣絕.
而用三十塊錢的公價將耶穌出賣.
但是撒迦利亞得到的只是三十塊錢的公價.
神就說你要將這些公價.
似乎是不宜的公價.
其實撒迦利亞沒有不宜的.
但是他就將這些公價三十塊錢.
神就說你要雕刻那個瑤瑚.
瑤瑚就是做那些泥土.
似乎給他們去做.
可能會做一些偶像.
一些公仔.
泥菩薩來去拜.
所以很特別就是神叫他雕刻瑤瑚.
猶大真的將他的錢.
剛才小杜行傳說.
那三十塊錢他自己畏罪潛逃.
就將那些錢扔下了.
但是那些人就將那些錢.
用來買了一個異鄉人的地.
來埋葬那些死了沒有人幫忙埋葬的人.
這樣來用.
所以雕刻瑤瑚.
也是撒迦利亞書提到的.
所以在這個畫劇裡面.
好牧人是不得好死的.
好牧人象徵的就是.
他做多少都得不到百姓任何的認可和接納.
你可以想像到將來真正的大牧者來到這個世界.

$^{1441}$在自己的百姓當中.
約翰福音很清楚說.
他到自己的地方來.
自己的人都不接待他.
你可以想像到他面對的就是.
所有的人對他的誤會.
你想像一下這麼熟悉.
跟他幾年的門徒.
都尚且會對他的身份有很多的半信半疑.
猶太甚至感覺到自己好像受騙.
幾年跟隨的耶穌.
原來是一路退縮.
一路沒有了那種彌賽亞的光彩和能力.
所以最終猶太都選擇背棄耶穌.
也就是說.
都是好像撒迦利亞的心境.
一個牧者要面對這樣的困難.
死的猶他死.
喪亡的猶他喪亡.
餘剩的猶他彼此相食.
一個牧者要說這番話.
他服侍的群體是極其的硬心.
絕對不是一種很負氣的說話.
所以這個話劇將來就應驗在耶穌身上.
他就是好牧人.
他就不是要索要那個功架.
而是猶太用三十塊錢出賣耶穌.
這樣的彌賽亞的預言.
都在這些話劇裡面.
在新約是按字面應驗.
所以有趣的地方就是.
新約的作者一來不是作出來.
真是猶大問那些祭司長.
祭司長你給多少錢我出賣耶穌.
三十塊錢.
在他們這些祭司長的眼中.
或者是猶大的角色.
只是值三十塊錢.
可能三十塊錢幣是很少的.
不是金幣.

$^{1481}$可想而知猶大徹底的硬心.
去到多大的程度.
如果給一千億,一千萬.
你就想一下值得.
但是三十塊錢你都要出賣自己的老師.
可想而知一來他感受到那個威脅.
越來越政治氣候越來越嚴峻.
因為愚節.
因為似乎收到風聲.
可能羅馬人都會藉機會鎮壓.
如果你猶大是和耶穌一起的.
你都會受牽連.
現在不單止不需要受牽連.
可能就會得到大祭司和祭司長.
一些官職都不奇怪.
所以猶大都是一個有理想的人.
他是一個懂得計算的人.
所以他真的計算過.
認為值得三十塊錢.
不過他應該有讀舊約.
應驗了撒迦利亞得到的工錢.
而那些是不義之財.
撒迦利亞就聽從神的說話.
要將這些價值.
第十三節.
十一章的十三節.
耶和華吩咐我說.
要把眾人所估定美好的價值.
雕給瑤護.
我便將這三十塊錢.
在耶和華的殿中雕給瑤護了.
跟著就跟手再截斷了.
連索這一根杖.
表明了廢棄.
剛才的榮美.
就是廢棄和萬民的約.
現在就是廢棄猶大和以色列弟兄的情誼.
即是徹底的民族和萬民.
都來到與神有一個疏離.
在這個第一部分的話劇裡.

$^{1521}$撒迦利亞作為好牧人.
其實內心很難受.
很想去木洋.
很想去警告群洋.
不要落在這些賣家的手.
很想給這些賣家.
即是那些屬靈的民間領袖.
要好好的秉公行義.
去木洋群洋.
但是奈何那個故事是一個悲劇.
撒迦利亞面對眾叛親離.
他就要取回自己的價值.
這個三十塊錢.
所以大家記住.
在新約裡面.
對於主耶穌的彌賽亞的預言.
越來越是對號入座的.
按字面應驗.
接著就很特別.
另外一個場景.
撒迦利亞就要扮演一個不好的先知.
對不起,不好的木洋人.
我口水有點乾,對不起.
或者都是遠一點說.
十五節開始.
耶和華又吩咐我說.
你再取愚昧牧人所用的器具.
因我要在這地興起一個牧人.
他不看顧喪亡的.
不尋找分散的.
不醫治受傷的.
也不牧養強壯的.
卻要吃肥羊的肉.
撕裂他的蹄子.
無用的牧人.
凋棄羊群有禍了.
刀必淋到他的傍臂和右眼上.
他的傍臂必全然的枯乾.
他的右眼也必昏暗失明.
只是告訴撒迦利亞.

$^{1561}$他要扮演一個不好的.
不稱職的牧人.
所表現出來的那種不稱職.
帶來的群羊是另外一種傷害.
所以是一個愚昧的牧人.
叫他做一個愚昧角色.
是怎樣的形態呢?.
要興起一個愚昧.
他就會.
因為最主要的牧養就是.
令到羊能夠.
例如不要有寄生蟲.
毛長得一路順滑.
肉一路長大.
不會偏行己路.
走去迷路.
或者被野獸攻擊.
但是一個愚昧的牧羊人.
不看顧喪亡.
不尋找分散.
也不醫治.
因為最需要去纏果.
最需要用藥酒去包裹.
不醫治受傷.
也不牧羊強壯.
但是他自己要吃肥羊的肉.
撕裂牠們的蹄子.
所以這裡就告訴我們.
神眼中的愚昧.
對一個牧羊人.
就是這樣的形態.
所以這個沒用的牧人.
有面對的禍就是.
神會興起刀兵.
來攻擊.
所以這裡就是.
在這一章裡面.
特別凸顯了.
這個民族他們本身的.
無奈的光景.

$^{1601}$和他的悲哀.
相信悲哀源於.
不是沒有好牧人.
而是好牧人不能夠.
牧養這些群羊.
是群羊自己硬心.
但是另一個場景.
第二個就是愚昧的牧羊人.
是沒用的牧羊人.
好和不好.
都是帶來傷害.
你可以想像到.
以色列這個民族.
以這個做一個單位.
神興起好牧人.
你們又不聽.
硬心.
興起不好的牧人.
就更加令到群羊受到虧損.
所以歸根結底.
在這一章裡面.
都好像預示.
將來的彌賽亞尼.
其實真的不是這麼多人.
會有這種屬靈的眼光.
對主耶穌的身份.
地位.
他要成就的救恩.
是會能夠明白.
是會跟隨.
所以初期的教會都是少數的猶太人.
來信主.
祭司長法利賽人.
其實不多.
是去到使徒的時代.
開始會多一些.
你可以想像到.
耶穌就是好牧人.
但是不好的牧人.
就一直都充斥.

$^{1641}$就是百姓.
耶穌自己在十章的.
約翰福音的講道.
都有提到.
好牧人.
和不好的.
分別.
詳細就不用說了.
所以用這麼大段的篇幅.
來說一個好牧人.
遇到的困難.
其實都給我們.
有一些啟發.
我們都要.
忠心的牧羊.
講好神要我們講的.
不要來避重就輕.
因為.
你遇了群羊.
不是每個人都能夠.
預備好心靈去接受.
神的道.
很多自己被賣.
都自己是蠢蠢不知.
你就可以想像到.
每個人都要面對.
將來神的審判.
但是.
當下到一個光景.
如果那個地區.
或者那個民族.
那個時間.
那個年代.
只是有.
不好的愚昧牧人.
那些百姓就更慘了.
因為根本就沒有人理你.
也都不會講那些真理.
給你聽.
那你就更慘了.

$^{1681}$就是.
不會看顧.
不會尋找.
不會醫治.
那你就可以想像到.
我們寧願有一個好的牧人.
但是好牧人往往就是.
神的道是不會令到他們有改變.
永恆福音持續我們曾經查過.
不斷的總結.
他雖然在他們行了很多的神蹟.
他們最終.
都是不相信耶穌.
那你可以想像到.
那個.
挫敗感是非常之大的.
但是很奇妙的就是.
主耶穌是.
心裡面柔和謙卑.
他不會.
灰心.
直到真理.
能夠得到那個彰顯.
那這個就是.
我們看了這三張聖經.
九十.
十一章.
其實都是環繞著.
以色列將來的寓言.
還有他們.
將來的復興.
和歸回.
夾雜的又提到.
牧羊.
好牧人的.
比喻甚至這個可能是.
一個戲劇.
撒迦利亞要扮演.
好牧人.
要扮演愚昧的.

$^{1721}$好用的牧人.
這樣的角色.
如果對他來說.
都不容易受的.
他要這樣去拿了.
愚昧牧人所用的器具.
即是那個棍.
那些裝石仔的.
背囊.
來去.
帶領羊群.
出去飲水.
吃草.
但是他是愚昧的牧人.
這個角色.
對撒迦利亞來說.
其實很難受.
好牧人又難受.
愚昧的牧人又難受.
十二章.
很特別.
有兩段.
第一段.
以色列有很大的德性.
第二段.
以色列人.
有一個心靈的德性.
就是.
家家戶戶.
男男女女老少.
都來到.
自己獨在一處.
什麼叫獨在一處呢?.
一定是牽連著.
上文說的.
第十節的應許.
神會將私恩.
叫人懇求的靈.
來挑灌大衛家.
耶路撒冷的居民.

$^{1761}$接著他們會.
仰望神.
就是他們所擇的.
這個字.
正正就是約翰福音十九章.
所出現的那個字.
真的.
篤進耶穌的.
肋旁.
仰望這位被擇的人.
這裡所指的人.
一定是在說神.
或者.
彌賽亞.
為我悲哀.
所以上文的.
一到第九節.
是大獲勝利的.
歡欣.
十到十四節.
是各人仰望.
被擇的.
不是歡欣.
是悲哀.
悲哀是歡欣的前奏.
真正在神.
面前.
每個人.
每個以色列人.
那個就是.
從遠處.
思想神.
想念神的.
異曲同工.
所以這裡告訴我們.
原來.
以色列靠他的政治.
靠他的武力.
那個德性.
一方面不會耐.

$^{1801}$你打贏人就代打你.
而是真正內心的.
那個尋求神.
每一個.
以色列人.
很仔細的描述.
這麼詳盡.
這麼沉醉.
可想而知.
每個人都有強烈的意識.
仰望這位被擇的.
就是耶穌.
我們揭一揭.
楊仁寬引述的這一段.
所以力量真的很大.
所以.
有些人是看耶穌傳.
信主.
其實是對的.
因為拍攝是非常嚴謹.
忠於原著.
來拍攝.
真是有個主角.
插進去.
十九章.
耶穌的死.
喉嚨有點癢.
所以可能我們說完這一段.
就要完了.
第三十一節.
十九章的三十一節.
約翰福音.
猶太人因為這一日.
是預備日.
猶欣拿安息日.
是個大日.
似乎是如意節之後.
有一個剛好的安息日.
所以就變了雙重的假期.
就求彼拉多.

$^{1841}$叫人打斷他們的腿.
把他們拿去.
免得失守.
當安息日留在十字架上.
於是兵丁來.
把頭一個人就是那個強盜.
的腿.
和耶穌銅釘的.
另外一個人的腿.
都打斷了.
來到耶穌那裡.
見他已經死了.
那兩個強盜.
其實可能是還沒死的.
因為他們沒有受那麼多的刑罰.
但是耶穌.
因為受之前的.
釘十架之前的摧殘.
逼迫.
所以已經是.
斷氣了,沒力氣再.
呼吸了.
就窒息而死了.
見他已經死了.
就不打斷他的腿.
唯有一個兵丁拿槍.
砸他的肋旁.
砸即是刺進去.
隨即有血和水.
分開了.
即是心臟停頓了一段時間.
血和水分開了.
這些有醫學根據.
流出來.
看見這件事的.
那個人就作見證.
那個人是兵丁.
還是什麼呢.
他的見證也是真的.
並且他知道自己.

$^{1881}$所說的是真的.
叫你們也可以信.
第36節就是.
作者約翰.
幫我們做一個總結.
告訴我們.
舊約出現過.
這些是成就.
為要應驗經常的.
話說.
他的骨頭一斤也不可.
截斷.
其實這是出埃及的關於.
逾越節的羊羔.
你們吃的時候.
那些骨頭是不可以截斷.
接著留意第37節.
經上又有一句說.
他們要仰望.
自己所摘的人.
撒迦利亞書的12章10節.
原來舊約的寓言.
是這樣的方式.
真的對號入座.
按字面應驗.
所以有些寓言.
你不可以當為一種.
象徵說法就算了.
是真的按字面應驗.
這些經文.
所以那個仰望.
在新約翰福音.
就是說.
這件事成就.
將來他們就會.
仰望這一位.
他們親手摘的.
耶穌彌賽亞.
對於我們.
外邦人也是沒有這些背景.

$^{1921}$我們一信就得到救恩.
但猶太人.
如果真正信主的猶太人.
他們能夠讀到新約.
其實真的非同小可.
舊約我們只是有撒迦利亞書.
已經有三段.
是直接按字面應驗.
在耶穌身上.
你可以想像得到.
如果猶太人自己讀新約.
他本身對舊約很熟悉.
他們真的有很大的感動.
所以希伯來文的新約.
其實很重要.
就面大兄就印了一大批.
送去.
以色列的聖經公會.
希望有猶太人.
拿起一本新約.
希伯來文的.
他有聖靈感動他.
他去看.
因為猶太人他們是沒有新約的.
他們是不信不接受新約的.
所以印了希伯來文.
你們就不能.
說你不懂希臘文.
你不懂中文.
用希伯來文.
去印新約.
很有意義.
我們華人去印.
但這裡告訴我們.
可能下一次再說.
就是這段聖經.
所以就算多強大的勝利.
如果只是靠武力.
改變不了民族的命運.
但如果是.

$^{1961}$仰望神.
帶來的悔改.
那個力量就很大.
所以其實教會.
真的需要.
一些奮卿會.
需要大家對神有心.
認識神.
我們一起祈禱.
感恩我們能夠.
用大約.
一個半小時.
可以一起來去研讀.
第九第十第十一章.
盼望我們還有餘下的.
三章的撒迦利亞書.
上主都給我們有.
那個眼光.
給我們有一個很強的心靈.
看到彌賽亞.
就是耶穌.
就是我們今天已經認識的.
主耶穌.
就是這位舊約寓言的.
這位.
謙和的.
救主.
騎驢進城的.
更加是在十字架.
被羅馬的兵丁.
用槍去截穿.
立旁的.
希望大家都能夠對.
舊恩有認識.
有明白.
更加將人帶進.
舊恩的那個.
泉源裡面.
多謝你聽我們的感恩禱告.
奉耶穌基督的名.

$^{2001}$阿們.
\newpage



\section{路加福音 3:1-9}
\label{sec:_n0nRPZl6sE}
\textbf{2023-12-09 晚堂《君王來臨的平安喜樂》路加福音 3\_1-9 講員:曾裔貴牧師}
\newline
\newline
連結: \href{https://youtube.com/watch?v=-n0nRPZl6sE}{\texttt{https://youtube.com/watch?v=-n0nRPZl6sE}} ~~~~ 語音日期: 2023-12-12
\newline
\newline
\hyperref[sec:Qx_35UZ1A44]{\small{< < < PREV SERMON < < <}}
~
\hyperref[sec:index_chronic]{\small{[返順時目]}}
~
\hyperref[sec:index_scriptual]{\small{[返順卷目]}}
~
\hyperref[sec:KJG_idsYK_w]{\small{> > > NEXT SERMON > > >}}
\newline
\newline
路加福音 3:1-9
\newline
\begin{longtable}{cl}
\hline
\hline
章節 & 經文 (和合本修訂版)\\
\hline
3:1 & \begin{tabularx}{0.7\textwidth}{X} 凱撒提庇留在位第十五年，本丟‧彼拉多作猶太總督，希律作加利利分封的王，他兄弟腓力作以土利亞和特拉可尼地區分封的王，呂撒聶作亞比利尼分封的王， \end{tabularx} \\ \\ \relax
3:2 & \begin{tabularx}{0.7\textwidth}{X} 亞那和該亞法作大祭司。那時，撒迦利亞的兒子約翰在曠野裡，神的話臨到他。 \end{tabularx} \\ \\ \relax
3:3 & \begin{tabularx}{0.7\textwidth}{X} 他就走遍約旦河一帶地方，宣講悔改的洗禮，使罪得赦。 \end{tabularx} \\ \\ \relax
3:4 & \begin{tabularx}{0.7\textwidth}{X} 正如以賽亞先知書上所記的話：「在曠野有聲音呼喊著：預備主的道，修直他的路！ \end{tabularx} \\ \\ \relax
3:5 & \begin{tabularx}{0.7\textwidth}{X} 一切山窪都要填滿；大小山岡都要削平！彎彎曲曲的地方要改為筆直；高高低低的道路要改為平坦！ \end{tabularx} \\ \\ \relax
3:6 & \begin{tabularx}{0.7\textwidth}{X} 凡有血肉之軀的，都要看見神的救恩！」 \end{tabularx} \\ \\ \relax
3:7 & \begin{tabularx}{0.7\textwidth}{X} 約翰對那出來要受他洗的眾人說：「毒蛇的孽種啊，誰指示你們逃避那將要來的憤怒呢？ \end{tabularx} \\ \\ \relax
3:8 & \begin{tabularx}{0.7\textwidth}{X} 你們要結出果子來，和悔改的心相稱。不要自己心裡說：『我們有亞伯拉罕為祖宗。』我告訴你們，神能從這些石頭中給亞伯拉罕興起子孫來。 \end{tabularx} \\ \\ \relax
3:9 & \begin{tabularx}{0.7\textwidth}{X} 現在斧子已經放在樹根上，凡不結好果子的樹就砍下來，丟在火裡。」 \end{tabularx} \\ \\
[1ex]
\hline
\hline
\end{longtable}
$^{1}$各位弟兄姊妹晚安.
請我們同心祈禱.
再一次我們恭敬的來到您的面前.
求主打開我們的心靈.
在將臨期的第二個禮拜.
我們不單止好像第一個星期.
我們要醒悟.
這一位君王真的隨時會再回來.
不單止是我們每年.
行常慣性去預備聖誕節.
而是我們大家的新心靈.
在主耶穌第二次回來之前.
我們真的要有一個警醒.
有一個悔改.
來迎接.
以致真正的屬靈的平安喜樂.
是會臨到我們大家的心靈.
使我們能夠在主和聖靈的引導下.
我們來結出好的果子.
與所蒙的恩典相稱.
今晚我們等候在您的面前.
求主施恩奉耶穌基督的名.
阿們.
剛才祈禱也提到.
這個月的講題.
我們會以將臨期四個禮拜.
準備進入聖誕節.
所以每一個禮拜都有特定的題目.
嚴格來說如果禮儀教會.
第一個禮拜就會點起紫色的蠟燭.
來預表著要有一個這樣的準備.
首先就是一個醒悟.
對自己有一個醒覺.
即是說主隨時會回來.
我們應該怎樣準備自己呢.
今天第二個禮拜.
我們就要選取這段經文.
這段經文本身.
當時的場景.
就已經告訴我們.

$^{41}$約翰有召喚去到曠野.
在曠野那裡.
走遍那一帶約旦河的地方.
來宣示一個嶄新或時代的開始.
就是君王要降臨.
這段經文.
路加醫生用了很多當時的敘述.
來襯托約翰所得到的.
最重要的就是神的呼召.
在這段經文裡面.
提到當時有最高的政治領袖.
就是羅馬的君王.
蓋薩提比劉.
就是緊接著蓋薩阿古士都.
第二位的羅馬皇帝.
他的之下當然就有猶太的分封的王.
就是這個我們熟悉的彼拉多.
接著也有另外的行省.
加利利的分封的王.
這個就不是大希律.
就不是殺主耶穌的大希律.
是希律王的兒子.
這個將來會處死約翰.
所以大家是有瓜葛的.
因為約翰在之後一段時間.
就直斥希律娶了他太太.
就是他弟弟的太太.
來犯了當時的罪.
所以希律王就懷恨在心.
希律王的太太就是希羅底.
另外還有其他的分封的王.
這個呂撒涅.
就作亞比利尼的分封的王.
接著還有不單止政治.
還有宗教上.
就是當時猶太最高的祭司.
大祭司.
即是在耶路撒冷.
在聖殿那裡來公職事奉.
有同時間阿拿和該亞發.

$^{81}$他們兩個是有親戚關係的.
因為我忘記了.
阿拿還是該亞發是他的外父.
所以兩個人同時間作大祭司.
當然一個可能是離任了.
另外一個就繼續接續了大祭司的位置.
所以這裡就已經有兩節聖經.
作者盧迦很想告訴我們.
當時的政治氣候.
宗教上的情況.
要凸顯了一個.
第三節就告訴我們.
撒迦利亞的兒子約翰在曠野.
如果我們不知道誰是撒迦利亞.
其實很多基督徒都應該知道.
就是他在聖殿裡面.
按照著班次去公職.
而在聖殿裡面見到神蹟.
見到天使向他顯現.
告訴他.
你的太太以利沙白將要生一個兒子.
你一定要幫她改名叫做約翰.
撒迦利亞就說.
我這麼大年紀怎麼可能生孩子.
這次他不相信.
天使就說.
只因為你不相信.
你就要啞巴直到你的兒子出世.
所以第一章就已經交代了.
約翰原來是祭司家族.
約翰原來是可以接續撒迦利亞.
在聖殿裡面事奉的一個祭司.
但是某個程度.
這是第二章經文的背景.
這個小朋友一路長大.
第二章的近尾時間.
就告訴我們約翰的事奉.
一路就在曠野那裡.
來到第八十節.
一章的八十節.

$^{121}$那孩子漸漸長大.
心靈強健.
住在曠野.
直到他顯明在以色列人面前的日子.
這樣描述約翰.
好像被人遺棄在曠野那裡長大.
心靈強健.
直至到他有一個時間.
就是神的時間到的時候.
他就要顯明在眾人面前的日子.
而且在預言他出世之前.
作者特別強調.
天使和撒迦利亞說.
你的兒子將來要使很多以色列人迴轉.
歸於神.
這個是一章近尾的經文.
我們做一點分析的時間.
我們就發現有少少作者.
用了當時歷史君王.
分封的王.
和宗教的.
大家都人所共知的.
相當有位份的阿拿該亞法.
這樣來襯托一個住在曠野.
一直等待一個日子.
要顯明在以色列人當中的約翰.
而他本身的出身是祭司家族.
是撒迦利亞的兒子.
很特別.
到底他在等什麼呢.
到底他為什麼自己會走到曠野那裡住呢.
這些都是我們大家應該感到好奇.
有這樣的身份 有這樣的出身 有這樣的地位.
為什麼要住在曠野呢.
曠野有什麼可以令到人會有羨慕呢.
雖然沒有描述他在曠野生活的一輪半找.
因為馬太馬賀是有描述的.
提到他吃的是蝗蟲野蜜.
他在那裡很簡樸的生活.
穿的是駱駝毛的衣服等等.

$^{161}$就是這樣來描述約翰等候的時期.
不過第三章第三節.
當描述了這些君王和當時歷史背景的時間.
然後有句話很重要.
神的話臨到他.
神的話不是臨到阿拿該亞法.
神的工作不是藉著該撒提比留.
因為當時是非常顯赫.
非常之大的帝國.
整個以色列都要完全成為一個附屬的國.
已經沒有自己的國家了.
以色列是完全亡國了.
但是在這樣的政治氣候下.
神一直訓練撒迦利亞的兒子.
撒迦利亞認不認同媽媽.
以利沙伯認不認同他的兒子.
要在曠野那裡等候.
而不是在耶路撒冷裡面受訓練.
受最好的祭司訓練.
將來接續做祭司呢.
似乎整條路完全不是這樣.
不過路加醫生給我們一句很重要的話.
就打開了這個章林的序幕.
因為還要是引述.
二賽亞書的四十章.
講到那個先鋒來的時候.
會帶來的那個改變.
這裡告訴我們.
當神的話臨到他.
第三節.
他就來到約旦河一帶的地方.
宣講悔改的洗禮.
洗到罪得到赦免.
有些解經家說.
這個洗禮為什麼令到猶太人這麼厭煩.
感覺無法呢.
因為其實本身不是給猶太人的.
是給那些要進猶太教的外邦人.
你要接受這個洗禮.
你才可以成為我們猶太教的一份子.

$^{201}$如果真的這樣的時候.
他走到約旦河一帶地方.
來到洗禮的時候.
很特別就是.
按道理你不再是祭司.
你又在曠野.
你又用這個洗禮這個儀式.
按道理就不受歡迎.
按道理就沒有什麼人會出去.
準備接受你的洗禮.
但不知道什麼原因.
會有什麼特別的吸引力.
他聽到神的話.
知道以為是神差遣他.
在那裡洗禮.
準備人心來悔改.
然後他就去到這裡的時候.
當他宣講悔改的洗禮.
洗到罪得赦免的時候.
然後一直都是作者盧家醫生的描述.
正如先知以賽亞書上.
這就是以賽亞書的四十章.
第三節.
在曠野有人聲喊著說.
預備主的道.
修直他的路.
不過留意.
第五第六節.
馬太和馬可都沒有引述.
只是盧家引述整段以賽亞書的經文.
第五第六節.
是非常有考究.
因為第七節有很多猶太人.
不同階層的猶太人.
就出到約翰那裡.
第五第六節就這樣描述.
一切的山窪都要填滿.
大小山崗都要削平.
彎彎曲曲的地方要改為正直.
高高低低的道路要改為平坦.

$^{241}$相信這樣壓穩的舊約的經文.
有他們希伯來的修辭.
相信肯定我們不要按字面去解釋.
你用什麼鏟泥車鏟開山崗.
削平填滿那些低窪.
相信一定是在說高傲.
卑微的可以來到去拉平.
拉平什麼呢.
預備一位君王到來.
就是第六節這裡.
凡有血氣的都要見神的救恩.
原本以賽亞書四十章.
是說耶和華的榮耀必然顯現.
凡有血氣的都會見到神的榮耀.
作者就加上他自己的一個詮釋.
似乎希利盟約他這樣來到去再解釋.
都不可以說錯.
凡是有血氣的會見到神的救恩.
這樣來到去描述.
以賽亞提到的為君王準備的.
這一位的先頭部隊的先鋒.
當時的世代.
如果有些很重要的人物出巡.
都會有這樣的準備.
相信我們都不覺得稀奇.
國家元首如果要去美國.
你想想那些只是廚師的準備.
都是很講究.
要有自己的御廚.
剛剛看了一篇報導.
很多年前的.
就是二零一幾年的時間.
習近平主席去英國.
有一個國事訪問.
兩個多月前.
已經秘密找了文華東方的主廚.
不告訴他你要接待哪個元首.
這是一個秘密的任務.
你要立即找文華東方兩個.
你認為最出色的廚師.

$^{281}$中西的.
就三個人去英國準備.
和國內的國家元首的御廚.
準備那些菜單.
是要高度機密.
不可以讓人知道.
你想想這麼多的複雜的.
所以我們都不會難理解.
一個先鋒為一個君王的準備.
這裡告訴我們.
很重要的就是.
在張臨期的第二個星期.
怎樣可以得到心靈的平安喜樂呢.
怎樣能夠可以經歷神的救恩呢.
怎樣能夠可以使我們.
來見到這位主.
這裡後面第七節.
真真正正回應.
高高低低 大小山崗.
無論是高低都可以拉平.
原來有一個秘訣.
這個秘訣就是.
千萬不要看自己的出身.
和倚靠自己的出身.
而忘記了對神恩典的倚靠.
因為第七節.
約翰對那些出來要受洗的眾人說.
讀蛇的種類.
誰指示你們逃避將來的憤怒呢.
你們要結出果子來.
與悔改的心相稱.
即是說要有一個回轉.
一個心靈的轉向.
才能夠看到自己有什麼需要.
要改變 要悔改.
因為自己心裡面.
說我們是亞伯拉罕的子孫.
即是說我們有亞伯拉罕是我們的祖宗.
這個好像一個免死金牌.
對當時代的猶太人.

$^{321}$即是意味著主耶穌的出現.
救恩的開始.
神的榮耀要顯出的時候.
一個劃時代的新約的開始.
約翰就是先頭部隊.
為這位君王準備.
即是說約翰之前的所有的猶太人.
都要立即明白.
亞伯拉罕的確是肉身的祖宗.
的確肉身是亞伯拉罕的子孫.
但是當耶穌要來的時候.
約翰就斬釘截鐵.
代表神的說話.
你們要結出悔改的果子.
和心裡面的回轉相稱.
不要自己心裡面.
即是未去找約翰.
已經有十足的把握.
約翰一定會歡迎我們.
看你是甚麼級數的歡迎.
如果是平民百姓的一般的接待.
如果你是很高階的一個領袖.
例如宗教的領袖.
例如政治的領袖.
約翰就應該要好好的接待我們.
但是一視同仁.
約翰真是有一種的冒犯.
見到每一個出來的人.
他好像二話不說.
都沒有任何的審查.
每個人出來都立即告訴他們.
似乎這個不是特意的.
要來這樣為罵而罵.
而是好像要點出一件事.
一個新的時代開始了.
每一個人其實在神造物主的眼中.
嚴格來說是魔鬼,撒旦,毒蛇的種類.
是因為有救恩耶穌的出現.
他成為我們的救贖主.
我們因為信他.

$^{361}$就可以成為真真正正的亞伯拉罕的子孫.
即是說原來猶太人不可以因為自己是猶太人.
就可以免疫於這個改變.
外邦人亦都不需要因為自己的身份出身是外邦人.
沒有亞伯拉罕的血統.
就好像不能夠得到這個救贖.
這樣的拉平就是救恩精彩的地方.
親愛的兄弟姐妹.
因為下面第八節就提到怎樣悔改,怎樣回轉.
就意味著和以塞俠書40章的寓言是完全很吻合.
怎樣能夠將高高低低,大大小小的山崗削平.
那些低窪的填滿.
原來就是需要約翰的警告.
我再讀下面第八節這一段.
這裡告訴我們結出一果子和悔改的深相稱.
然後第八節的中間都是.
很奇妙.
神能夠製造另外的亞伯拉罕的子孫.
是約翰所說的.
所以雖然這裡就這樣按字面是有難以理解.
但我們有救恩,有聖靈的弟兄姊妹.
我們就明白.
神就是興起每一個信的人.
不論是近處,遠處.
不論是什麼人.
只要在主耶穌裡面.
就成為了屬靈的亞伯拉罕的子孫.
所以這裡是神的工作.
神能夠從這些石頭.
礦物.
真的是字面的石頭.
是毫無用處.
是何等的卑微.
甚至乎卑賤的.
完全沒有生命的石頭.
可以成為亞伯拉罕的子孫.
彼得自己也是這樣說的.
彼得前書的二章第四節.
你們來到主的面前.
也就像活石.

$^{401}$被建造成為靈宮.
作聖傑的祭司.
奉獻神所閱立的靈祭.
我們本來是被廢棄的石頭.
但是在主的手中.
救恩裡面.
我們就是活石.
我們可以大家聯絡.
在永恆裡面.
我們就成為屬靈的宮殿.
那裡獻祭給這位永生神.
你想想我們多麼的尊貴.
這裡就有兩個分別.
讀石的種類.
和亞伯拉罕的子孫.
約翰不是想為了罵人而罵人.
而是他要揭示一個新時代的開始.
耶穌基督還沒來到這個世界.
救恩還沒正式開始.
一路的預言.
都是神和亞伯拉罕立了這個永約.
接著從摩西建立了這個西乃山的舊約.
但是當主耶穌來到這個世界.
當很快耶穌就出現.
來到去傳道.
來到去告訴人.
在祂裡面就有這個永恆生命了.
而且耶穌自己也親口說過.
在約翰分第八章.
你們不聽我講.
你們不能夠信.
因為你們不是亞伯拉罕的子孫.
耶穌這樣說的.
說那些猶太人.
猶太人就說.
你明明告訴我們你是否基督.
我們就相信你了.
耶穌說你們不相信.
因為你們不是亞伯拉罕的子孫.
你們是出於你們的父親撒旦.

$^{441}$八章的五十四節.
約翰福音.
這樣來描述.
就是說原來主耶穌的出現.
約翰的宣告.
不是無的放矢.
而是每一個人.
無論你本身是猶太人.
希伯來人.
你還沒有在耶穌的救恩裡面.
意味著今天的猶太人.
沒有耶穌的救恩.
他們都是撒旦毒蛇的種類.
我們因為有主耶穌.
有聖靈.
我們就是神的兒女.
更加我們可以稱為亞伯拉罕的子孫.
享受得著神的應許.
給予我們福氣.
在這裡繼續說的時間.
第九節不單止提到要有一個迴轉.
不要以自己肉身的身份.
來作為可以免去神的憤怒的一個藉口.
或者一個擋箭牌.
因為第九節.
現在苦子已經放在樹根上.
凡不結好果子的樹就砍下來.
這些砍樹在舊約經常都有.
以上書第五章.
那些不結果的葡萄.
就是這樣被砍下.
路加也有很多.
例如十三章有一棵無花果樹.
栽在葡萄園.
就砍下.
這裡再一次告訴我們.
約翰很迫切.
他感受到這個訊息.
神的說話臨到他.
就是說時間到了.

$^{481}$一個新時代的開始.
就由約翰宣告傳道開始.
而且這個開始就是末世.
因為苦子已經放在樹根上.
凡不結好果子的樹就砍下來.
雕在火裡.
這些描述.
一而再在耶穌的講論裡.
都會再出現.
這個雕在火裡是永恆的滅亡.
這是一個非常嚴肅的問題.
所以第十,十一和十二節.
雖然我們沒有引述.
但是就是印證約翰所說的.
是有小部分的人.
害怕和聽從和願意有一個迴轉.
也都印證原來這個預言所承載的.
無論你有多麼的高傲.
多麼的卑微.
身份上或者人體裡的身份上.
都可以改變.
只需要一個回轉.
心靈的迴轉.
第十,十一,十二節.
其中好像有實在的一個例子,案例.
這樣去問,這樣去答.
其中那個迴轉就是.
整個心靈開始對其他人有一種憐憫.
留意第十節.
中人就問了.
這樣就是說.
如果這麼嚴肅,這麼恐怖.
又審判要來.
有把斧子.
如果握著那個斧子的不是一個神經的人.
不是一個有問題的人.
而是神.
那我應該怎樣呢?.
約翰的回答就是.
有兩件衣裳.

$^{521}$就分給那些沒有的.
有食物的也當這樣行.
這麼簡單.
分啊.
相信這裡告訴我們.
對別人的那個憐憫.
相信是知足.
在神的裡面.
你的人生.
是一種恩典.
而你能夠白白的願意去捨棄.
尤其是有意識地去留意.
其他人的需要.
我們基督教裡面有很多很好的弟兄姊妹.
沈祖堯教授.
他做校長的時候.
當然要面對很多的壓力.
很多很多的事情.
但其中一樣東西真的很感動.
他的分享.
他說他有一天.
有一個很強烈的意識.
他看到正生書院.
在梅窩要建新的地方.
村民.
原居民.
極力的反對.
明明是自己買了那個地方要建的.
因為怕那些年輕人犯事.
法庭判去那個地方.
接近開放式的勞教中心.
很簡陋的.
在大嶼山的霞澗那個地方.
那時候我們在報道團.
我都參加了一年.
每個星期五我們就去正生書院.
這樣去跟那些年輕人去報道.
很多都未信耶穌.
個個十幾歲就紋身.
有些是吸毒進來.

$^{561}$有些是濫藥等等.
不同的情況.
很多都是家庭破碎.
沈祖堯教授見到這樣的情況.
看新聞報紙.
他就跟他的同事.
想到一個辦法.
我們能不能夠盡我們一點點的力.
可以讓這些年輕人.
有一個向上的機會.
在不減低大學的修生的標準底下.
能不能夠容許至少可能一個兩個.
稍為標青稍為肯讀書的正生裡面的年輕人.
可以進到中文大學呢.
他真的這樣去說服他的同事.
就跟正生的校長陳肇卓博士去談.
一談就合作.
就開始這個MI的計劃.
Mentorship的計劃.
就真的揀選了大概六個.
不單止給他們機會.
因為你可以想像到.
他們成長全部是破碎家庭.
讀書不行.
因為根本已經學壞了.
可想而知他們的英文程度.
是根本不達到的.
沈教授身體力行以身作則.
邀請這些學生被挑選的.
來幫他們補英文.
用了好幾個月時間.
每個星期校長幫他們補英文.
希望能夠讓他們的程度.
一路一路跟得上.
其實是給他們一個機會.
一來希望打破那個隔膜.
和那個歧視.
他們後來因為都熟了.
在他家院長的宿舍.
去學習補習.

$^{601}$太晚了.
容許有四個留宿.
第二天早上.
沈教授如常的.
一早就去跑步.
神魂.
跑跑下.
死了.
家裡有四個男人.
紋身的.
還有我太太.
還有其他的工人.
怎麼辦呢.
是不是應該馬上跑回家呢.
不知道發生什麼事.
但他想想.
如果他這樣做.
其實在他內心一樣有歧視.
有那個成見.
他就忍住.
按照他自己的時間.
可能半小時.
可能一小時.
他才回家.
沒事發生.
很正常.
其中一個成為了他的乾兒子.
後來進了中大畢業.
他講回這個經驗.
他說我們是不是應該.
要學習聖經的教導.
要來對別人有一份的憐憫.
他們明顯是社會已經標籤了.
你給他們一份工都沒有用.
因為根本上.
社會對他們都有一個戒心.
所以當然我們做這些事情.
我們需要有智慧.
大時有智慧.
正正就是一個標籤.

$^{641}$這裡約翰告訴我們.
跟這些猶太人說.
如果你們真的有一個迴轉的心靈.
你裡面應該會有一份神的憐憫.
迴轉有神的心腸.
來看到有兩件衣裳就分級窮人.
有食物也這樣走.
然後有稅來.
很實在的.
我收稅的.
我不能夠不收的.
我不能夠來到現在.
特意不收的.
我應該怎麼做啊夫子.
約翰就說.
對啊.
除了例定的數目.
就不要再多取了.
原來多取是行規的.
多取是沒有人會說你的.
連那個羅馬的你的上級.
都不會說你的.
因為你不多取.
你就不能夠有你自己的額外的多.
這個多取就有很大的空間.
有些人就非常非常的嚴厲.
所以你看看.
盧家鋒有一個叫薩蓋.
盧家十九章.
說這個薩蓋的信主.
他是那個地方.
耶利哥的瑞利長.
他曾經訛詐人.
用他的官職來逼令很多人.
給很重的稅.
這裡告訴我們.
例定的數目.
就是說你的迴轉.
仍然可以在你自己的崗位上.
改變那個崗位的.

$^{681}$或者那個行業的那些的陋習.
那些的罪惡.
不同流合污.
這樣很大件事的.
很難的.
要站穩.
是很不容易的.
相信大家在不同自己的行業裡面.
必定有你自己的掙扎.
有你自己面對的困難.
但這裡告訴我們.
除了例定的數目.
就不要多取了.
不是你的.
你就不能夠據為己有.
就是說你要示範.
樹立一個行業的基督教的倫理.
這個就是近年.
國航編牧師和廖炳彤牧師.
在我們課堂裡面.
他們建立的.
希望不只是個別.
你在你的工作崗位.
傳福音這麼簡單.
而是你能夠代表著.
你自己的專業.
建立到有自己的同行.
來有一個基督教的標準.
讓人看到.
我們這些的專業.
是有倫理的守則的.
基督教是這樣的.
這樣的影響力就更加大.
所以原來約翰當年.
已經有這樣的教導.
他聽到這個說理.
他聽到他如果沒有那個回轉.
聖靈工作.
其實就難之又難.
但是有聖靈.

$^{721}$又有那個心思的回轉.
那就很不同了.
他就知道應該怎樣做了.
所以你看看路加提到的那個殺該.
主啊.
我將我一半的分給窮人.
我訛詐過誰的.
我要還他四倍.
因為救恩今天來到我家.
耶穌願意和他吃飯.
他立即這樣宣告.
你可想而知.
聖靈的工作.
在這位將臨的君王裡面.
帶給人的是真正的平安.
和喜樂.
第十三節.
下半節對不起.
第十四節.
悠悠冰冰就問了.
我們應該做什麼.
約翰就說.
不要用強暴待人.
也不要訛詐人.
自己有錢糧就當知足.
這裡也反映當時的行業.
當時的作為兵丁.
會有這種的武力.
可以行駛.
所以仍然都是帶著一份對別人的憐憫.
你做對你自己的工作就夠了.
不可以越了你自己的崗位的權力.
權力根本就是無影的.
用多少根本有時很難量度.
完全在乎那個人的心靈.
有沒有這種界線.
有了權力真的很難駕馭.
最大的危險就是權力慾.
很難駕馭其實是自己的內心.
所以約翰這裡真的非常可圈可點.

$^{761}$整個敘述整個描述.
都是想告訴我們.
要準備君王的到來.
每一個行業每一個階層.
都可以見到神的救恩.
準備好自己.
去讓神進入我們大家的心靈裡面.
這樣我們就能夠享受到.
神賜給我們的平安和喜樂.
所以這個迴轉是很重要.
這個心靈的轉向.
剛好又看到一個例子.
真人真事.
有一個銀行家.
他當時的金融海嘯.
不是現在.
總公司要他放無薪假.
他就很不開心.
賺了很多錢.
但是突然之間因為收縮的經濟.
要他無薪放假.
差不多等於就快要遣散.
他很不開心.
又不知道怎麼辦.
沒有事做.
他的好朋友.
來跟他說.
你當時不是說自己很忙.
什麼都做不到.
為什麼現在不用這些時間.
陪你家人.
做一些自己有興趣.
喜歡做的事.
他想了很久.
他現在又不是手頭上欠錢.
他現在可以很空閒.
他真的用時間陪他太太.
陪他家人.
還有很多時間.
他開始發展自己的興趣.

$^{801}$他學人做餅.
越做就越有興趣.
後來自己開一間很小的餅舖.
親手去做這些餅.
後來發展到他失去工作.
因為自己辭職.
後來專心去打理這間餅舖.
他找到的收入和興趣.
那個滿足.
比他做銀行家的時候還大.
那種基本上每天對著電腦.
和目標.
要去做那個數.
他感受到很大的心靈滿足.
就是一個迴轉.
當然這是每個人的際遇.
親姐妹.
我們每一年等待聖誕節.
雖然我們不是禮儀教會.
真的有些蠟燭給你.
有一個醒悟.
要來了.
24號就是了.
每年我們都是這樣去提醒.
但很盼望我們在年底.
我們真的有一個在這個將臨期.
要經驗神給我們的平安喜樂.
很重要就是.
我們的心靈要有一個柔軟.
要有一個迴轉.
願主祝福大家.
也讓我們一起等候.
我們的救主降臨.
我們一起禱告.
多謝主讓我們今天晚上.
能夠靜心的聆聽你的真理.
多謝你讓我們能夠看到.
施洗約翰.
他在當年在神的引導.
在神的話的呼召下.

$^{841}$他成為那個時代的先知.
將神的心意命令傳遞給當時的猶太人.
來預備耶穌的降臨.
和耶穌的出現.
使到有一些部分的猶太人.
來相信和跟隨.
但我們知道有更多.
仍然是硬心拒絕耶穌的.
我們很盼望今天.
我們大家都是已經屬於主的弟兄姊妹.
這個不是一個福音性的聚會.
但是我們盼望.
主給我們大家有一份柔軟的心靈.
我們來思想.
怎樣可以在聖靈的感動引領下.
去做神要我們做.
能夠滿足神的心意.
可以幫助更多的人.
尤其是帶人進到救恩裡面.
得著耶穌的拯救.
我們這樣感恩禱告.
奉耶穌基督的名.
阿們.
\newpage



\section{馬可福音 1:1-8}
\label{sec:KJG_idsYK_w}
\textbf{2023-12-10 加恩樓《修直祂的道》馬可福音 1\_1-8 講員:楊明麗姑娘}
\newline
\newline
連結: \href{https://youtube.com/watch?v=KJG_idsYK-w}{\texttt{https://youtube.com/watch?v=KJG\_idsYK-w}} ~~~~ 語音日期: 2023-12-12
\newline
\newline
\hyperref[sec:_n0nRPZl6sE]{\small{< < < PREV SERMON < < <}}
~
\hyperref[sec:index_chronic]{\small{[返順時目]}}
~
\hyperref[sec:index_scriptual]{\small{[返順卷目]}}
~
\hyperref[sec:m2hnqzcprEM]{\small{> > > NEXT SERMON > > >}}
\newline
\newline
馬可福音 1:1-8
\newline
\begin{longtable}{cl}
\hline
\hline
章節 & 經文 (和合本修訂版)\\
\hline
1:1 & \begin{tabularx}{0.7\textwidth}{X} 神的兒子，耶穌基督福音的起頭。 \end{tabularx} \\ \\ \relax
1:2 & \begin{tabularx}{0.7\textwidth}{X} 正如以賽亞先知書上記著：「看哪，我要差遣我的使者在你面前，他要為你預備道路。 \end{tabularx} \\ \\ \relax
1:3 & \begin{tabularx}{0.7\textwidth}{X} 在曠野有聲音呼喊著：預備主的道，修直他的路。」 \end{tabularx} \\ \\ \relax
1:4 & \begin{tabularx}{0.7\textwidth}{X} 照這話，施洗約翰來到曠野，宣講悔改的洗禮，使罪得赦。 \end{tabularx} \\ \\ \relax
1:5 & \begin{tabularx}{0.7\textwidth}{X} 猶太全地和全耶路撒冷的人都出去，到約翰那裡，承認他們的罪，在約旦河裡受他的洗。 \end{tabularx} \\ \\ \relax
1:6 & \begin{tabularx}{0.7\textwidth}{X} 約翰穿駱駝毛的衣服，腰束皮帶，吃的是蝗蟲和野蜜。 \end{tabularx} \\ \\ \relax
1:7 & \begin{tabularx}{0.7\textwidth}{X} 他宣講，說：「有一位在我以後來的，能力比我更大，我就是彎腰給他解鞋帶也不配。 \end{tabularx} \\ \\ \relax
1:8 & \begin{tabularx}{0.7\textwidth}{X} 我用水給你們施洗，他卻要用聖靈給你們施洗。」 \end{tabularx} \\ \\
[1ex]
\hline
\hline
\end{longtable}
$^{1}$各位大英姐妹 早安.
很感恩今天有機會和大家一起崇拜.
今天我們是 大家很清楚.
今天是我們張琳琦的第二主日.
我們為了預備聖誕的時候.
我們也重溫一下聖誕的訊息.
也讓我們預備好心靈.
剛才我們說過.
這個禮拜我們是悔改.
上個禮拜我們是反思.
這個禮拜我們要進入悔改期.
我們又看一段經文.
看這段經文對我們有什麼提醒.
讓我們在第二主日可以做好準備.
我想張琳琦四個主日大家很熟悉.
第一個禮拜是醒悟.
我們點了紫色的蠟燭.
第二個禮拜是悔改.
我們又點了紫色的蠟燭.
我們期待下個禮拜.
我們進入一個喜樂期和希望期.
我們不是迎接耶穌基督的降生.
不過我們是預備好心靈.
在迎接聖誕的同時.
我們是預備心靈迎接主耶穌基督的再臨.
我們讀了幾段經文.
第一節《萬學福音》第一章第一節.
它說:神的兒子耶穌基督福音的起頭.
福音的起頭是怎樣呢?.
如果有人問你福音.
你可能會想福音的起頭是怎樣.
但原來萬學福音一說福音的起頭.
它不是說耶穌基督的降生.
它是說新約和舊約是連貫的.
不是新約在舊約完結後才開始.
也不是新約來否定舊約.
它是說新約來完成舊約的預言.
在舊約的書當中.
《先知書》.
大家可能很熟悉的《以賽亞書》.

$^{41}$不過未必能讀完.
《以賽亞書》有66章.
鼓勵大家早點讀完.
現在有很多不同的方式.
你可以讀完聖經.
除了自己拿著紙去看.
或者在網上看之外.
網上有很多聖經.
它是讀給你聽.
聽應該容易些聽得完.
而且要留心去聽.
不要開得很快.
留心去聽.
這樣你會聽得很有領受.
大家可以聽一下.
或者自己去讀一下.
《以賽亞書》.
《以賽亞書》當中有一段預言.
「漢娜,我要參與我的使者在你面前」.
「他要為你預備道路」.
這個預備道路是預備誰的道路呢?.
是預備主耶穌基督的道路.
要來預備主耶穌基督的道路的人.
就叫做使者.
是開路的先鋒.
馬可福音就由這個先鋒開始去講.
他講到這個先鋒的時候.
他就介紹了這個叫做「施洗約翰」.
他說:召者說.
照著這個預言所講.
施洗的約翰就去到曠野宣講悔改的洗禮.
使罪得赦.
由太傳地和全耶路撒冷的人都出來.
到約翰那裡承認他們的罪.
在約旦河裡面受他的洗.
約翰就穿著駱駝毛的衣服.
腰束皮帶.
吃的是蝗蟲野物.
我不知道大家對這段熟悉的經文.
有些什麼東西是你很有印象的呢?.

$^{81}$我第一樣有印象的就是聖經當中有幾個約翰.
所以就給他一個分別.
這個是施洗的約翰.
第二樣就是他有特別癖好.
他穿著駱駝毛的衣服.
腰束皮帶.
吃的就是我不敢接觸的蝗蟲.
蜜糖也會吃的.
這是他特別的地方.
但這段聖經一講到施洗約翰.
是他去到曠野的地方.
在《說尼路書》講到悔改的洗禮.
他是什麼出身呢?.
我們今天的崇拜主席姓莫.
雖然他有一個大家更加熟悉的名字.
但他姓莫的 莫弟兄.
施洗約翰他姓什麼呢?.
原來猶太人和以色列人以前是沒有姓氏的.
哪一個約翰呢?.
他出名的就是講他的那樣.
據聞日本的姓氏也是這樣來的.
他住在哪一個地方.
那個地方就是他的姓氏.
例如井上什麼什麼.
他姓井上.
這位約翰最出名的就是施洗.
所以分別為.
他的名字出來讓大家知道是講誰呢?.
施洗的約翰.
施洗的約翰去到曠野.
無情白事他為什麼去到曠野呢?.
話說施洗約翰當時的年紀.
大概是二十多三十歲.
我們當中有沒有二十多三十歲的青年人呢?.
原來他當年是青年人.
是青年人.
他是青年.
這個青年才俊他做什麼呢?.
大家二十多歲嘛.
讀完書畢業就要找工作.

$^{121}$他是比較特別的.
其實他不用去找工作.
他不是失業.
他不是失戀.
他不是沒有工作.
也不是呢.
叫做什麼呢?.
Working Holiday.
他不是.
他不是因為Working Holiday所以去到曠野.
Working Holiday他去了澳洲.
不是.
他是有一份.
是不是叫世襲?.
他爸爸做的那一行他可以跟著做.
猶太人本來就是這樣.
而他那一行.
他爸爸做的那一行很威風.
威風到他不用科目去做.
他爸爸是祭司.
他媽媽是利未人.
他一出生BB的時候.
他就是吃祭司的家飯.
就是祭司的.
他們利未人的食物.
祭司的食物.
不是每個人都可以去吃的.
但他一出生他就可以在.
祭司的家屬當中.
有資格去吃祭司的飯吃.
就是說他不需要有飯吃.
他可以不用太過怎樣找工作.
怎樣他都是一個祭司.
還要他是哪個地方的祭司?.
他爸爸是耶路撒冷的祭司.
他其實就是.
繼續可以在耶路撒冷做祭司的了.
但是這個青年才俊呢.
他就沒有留在耶路撒冷.
他去了曠野.

$^{161}$他不是懶惰.
他不是想找一份休閒一點的.
時時可以看一下什麼呢?.
看一下這個日出與海灣.
不是.
他是覺得神對他的領導.
帶領他去到曠野的地方.
去那裡做一個更有意義的工作.
祭司的工作是什麼呢?.
祭司的工作是代表人.
向神認罪.
求悔改.
也將神的意思告訴人民.
讓人可以與神建立美好的關係.
祭司是幫助人與神建立美好的關係.
約翰去到約旦河邊曠野地做什麼呢?.
他不是去那裡和羊,牛,駱駝,驢子說話.
他真的去過樹.
他對著人去講道.
他去的地方.
我們覺得不是富庶的地方.
耶路撒冷人多.
耶路撒冷是猶大省的省都.
省會的地方.
應該是最富庶的地方.
最多人的地方.
最密集的.
我們要做報佳音的工作.
傳福音的時候.
我們都會去找人多的地方.
為什麼這個青年人.
好像躲在曠野呢?.
聖經中說到曠野的時候.
說到什麼呢?.
耶穌基督去了曠野四十天做什麼呢?.
退休.
他去那裡做什麼呢?.
青年才俊去到那裡.
他去傳講悔改的福音.
他去傳講.

$^{201}$叫人知道自己有罪.
需要與神建立關係的福音.
這些他去對著一群猶太人.
一群認識耶和華.
信耶和華的人.
已經很多代了.
那時候還沒有傳耶穌基督復活的佳音.
耶穌基督還沒有出來傳道.
他去那裡傳講悔改的福音做什麼呢?.
全部都是信徒.
還要是全部都聽著這裡.
聽著大的.
他們懂得背舊日聖經.
從小聽到大的道.
你去到一個曠野地方講.
誰來聽呢?.
其實從小聽到大的道.
很難講的.
怎樣講才會有新意呢?.
但是他就是預備好了.
神對他有這樣的呼召.
預備他有一個道路.
他就是要去到曠野的地方.
宣講人要與神建立關係的道.
他去到曠野沒有什麼人的地方.
其實約旦河邊應該是比較富庶.
起碼有很多靚草的地方.
種植也有的地方.
但是他去到比較曠野的地方.
應該是沒有什麼人煙的地方.
在那裡傳道.
誰來呢?.
他很特別.
他就在那裡傳道.
那些人就說耶路撒冷來.
不是吧.
你去過一兩天.
當作是旅行休息一下.
不是.
是耶路撒冷的人去他那裡聽道.

$^{241}$去他那裡受洗.
他們受洗的意思是.
我承認我是一個罪人.
我需要與神重建關係.
這是耶路撒冷的人.
要去到曠野找約旦的原因.
但是他們其實在耶路撒冷.
聖殿就在那裡的地方.
為什麼要去到曠野找約旦呢?.
施洗約旦.
他有這個使命.
叫人與神重建關係.
叫他去到一個曠野的地方.
去傳講福音.
這個使命對於一個青年人來說.
他應該說.
不是吧 天父上帝.
不是在那裡的.
我領受錯誤.
我回去再多祈禱兩個星期.
應該另外有一個地方要我去的.
應該傳福音去人多的地方.
不是.
神真的叫他去那裡.
他真的聽話去那裡.
他去那裡聽講道的時候.
真的有人來聽.
不是一個.
不是兩個.
是越來越多人來聽道.
越來越多人來聽道.
而且每次講到.
都不是很安慰性的訊息.
不是鼓勵性的訊息.
不是說我們很有愛.
不是.
一直是一種告訴他.
你這樣做是有錯的.
你要知錯知罪.
你要悔改.

$^{281}$但是就有很多人.
越來越多人.
可能是因為有人去聽完之後.
很有領受.
回去他跟隔離論社講.
於是更加多人.
從耶路撒冷過去聽.
不單止.
耶路撒冷的人聽完之後.
跟隔離附近的人講.
可能伯利恆的人也去.
伯納利的人也去.
於是乎.
猶大全地的人.
都去抗耶.
那我就知道.
為甚麼約翰要去抗耶.
要給那麼多人.
都能夠一起去聽道的地方.
耶路撒冷雖然是大.
沒有那麼多空間.
沒有那麼多人.
可以集合在一起.
去聽道.
也不能夠那麼多人.
可以在那裡.
暢所欲言地悔改.
去求問.
請約翰指引他們.
一條路.
我應該怎樣悔改.
但是去到抗耶的地方.
他們有空間.
他們有空間可以聽道.
他們有空間.
可以講述自己的需要.
他們有空間.
讓約翰可以教他們.
提醒他們應該怎樣悔改.
怎樣才能夠跟神重建關係.

$^{321}$於是乎.
不單止是他爸爸服侍的地方.
耶路撒冷.
有人去聽道.
還要是隔離附近的城鎮.
再遠一點的城鎮.
都有人來他那裡聽道.
整個猶太省的人.
都來他那裡聽道.
聽什麼道呢?.
悔改的道.
來他那裡聽道.
走得遠一點.
去得遠一點的地方.
不好意思.
不像嘉欣樓.
聽完道之後沒有茶點.
高山不在那裡.
於是乎他們沒有茶點.
沒有美食.
為什麼要去聽道?.
那只是他們生命真正的需要.
他們找到他們生命的需要.
在曠野.
在一個沒有人煙的地方.
在一個沒有富樹出產的地方.
去到那裡.
他們找到生命之道.
約翰.
施洗的約翰.
就在那裡傳講他的道.
讓那裡的地方的人.
能夠有領受.
這班人.
聽到這個道.
悔改的道.
他們真的.
心靈有感動.
他們心靈有感動.
願意認罪.

$^{361}$求一件事.
求神去赦免他們.
讓他們能夠.
收命得到拯救.
其實猶太人來了.
他們有應許的.
他們是有得救的.
他們是得救的.
為什麼還需要求得救呢?.
你們這些姐妹生活得開心嗎?.
我相信那個時候.
知道自己身上有罪的人.
即使知道自己有永生.
他們也生活得不開心.
不是社會給他們的壓力.
不是工作給他們的壓力.
但是在心靈裡.
有罪咎的壓力下.
他們生活得不開心.
他們需要有這個把握.
我今天所犯的罪.
雖然在永恆當中.
神寬恕我.
我會有永恆的生命.
不過我的生命.
能夠獵家得到赦免.
我可以安心生活.
是人的需要.
施洗約翰所傳說的.
叫他們認罪悔改.
不是給他們壓力.
你不認罪嗎?.
看你今晚回去能不能睡得著.
他讓他們知道.
你認罪悔改嗎?.
你有一個新的出路.
神是赦免的.
不過不要逗留在罪惡當中.
不要為魔鬼留地步.
魔鬼時時都張著飢渴的口.

$^{401}$想要吞食我們.
但是神時時都有恩典.
我們向祂認罪悔改.
我們的生命就可以生活得喜樂.
可以生活得有平安.
施洗約翰所傳說的時候.
他們認罪悔改.
不但如此.
他們還要向約翰問.
我以後要怎麼做?.
你猜約翰怎麼跟他們說?.
稅利不要緊.
你繼續收稅吧.
神赦免你.
不是的.
約翰跟他們說.
你做稅利的.
做稅利的.
你應該收的就收.
不應該收的就不要收.
讓人民可以有安定的日子過.
你自己不需要這麼富有.
生活到就可以.
他讓稅利不再需要為自己冠上.
我是罪人稅利的身份.
他可以做一個為人民工作的公務員.
稅利去到他的屬.
可以得到出路.
這是約翰所做的工作.
讓人去到他的屬的時候.
心靈都能夠得到幫助.
除此之外.
還有什麼恩典呢?.
他告訴別人.
我今天為你施洗.
但將來還有一個人為你施洗.
在為他施洗的時候.
我們看一看.
他說將來有一位要來.
這一位要來的是要用聖靈為人民施洗.

$^{441}$他是指向主耶穌基督的來臨.
主耶穌基督來臨的時候要為他們施洗.
是聖靈的洗.
是指向一個永恆的國度.
這只是他們最重要要領受的洗禮.
弟兄姊妹.
今天我們傳揚福音.
不是傳揚我們的口才.
我們傳揚福音.
不只是傳揚教會的堂旨.
我們是在傳揚耶穌基督賜給我們的福音.
聖經當中所說的.
以賽亞書所說的預言.
有一個要來為耶穌基督預備道路的.
這個就是施洗約翰.
他要來怡平山.
他要來怡平高山低谷.
削平是「包肖」.
兩字很識讀.
削平所突出來的地方 以為平地.
也讓低谷的地方可以填滿成為平地.
在「包削」的地方.
他可以來怡平.
低谷的地方又可以填滿.
弟兄姊妹.
我們傳福音的時候.
有些人會有少許驕傲.
有些人會有少許自卑.
甚麼都是我們向他們傳福音的對象.
我們要預備他們的心來領受福音.
讓他們的心靈預備他們.
可以接受主基督的救恩.
是我們的工作.
讓基督的福音能夠去到他們當中.
就好像以賽亞書五十四章所說的預言.
讓人的心靈能夠痛悔.
謙卑的人能夠得著神的恩典.
聖靈的同住.
使得謙卑的心靈甦醒.
痛悔的心能夠多醒覺.

$^{481}$這是神想我們預備的.
我們不是預備主耶穌基督的來臨.
我們等待祂的再來.
但在聖誕節之前.
我們傳福音的同時.
我們傳福音之前做少許的事.
做一個少許的先鋒.
預備人的心靈.
讓神的恩典能夠臨到他們當中.
祂的醫治能夠來到人的心靈當中.
讓人的嘴唇能夠結出讚美歌頌的果子.
在曠野有人聲.
喊著說.
預備祂的道路.
我很想大家看一看.
我覺得是令我的日子很燦爛.
看到照片就很開心的照片.
對你們來說並不陌生.
是過往幾個月大家的努力.
你們不是星期五晚沒事做.
不是星期六晚很空閒.
而是要預備主的道路.
其實我們清楚一件事.
我們的方向是要預備去哪裡.
我們的方向是預備去哪裡.
你預備聚會的原因是什麼呢?.
不是因為大家想唱歌而已.
回家有很美麗的飯菜等待你.
不需要你清理場地.
但來到這裡.
你想預備一個可以移平大山小山.
可以填滿低谷的地方.
讓心靈有需要的人能夠得著.
弟兄姊妹我們需要繼續努力.
這個工作還沒完.
我們要朝著我們真實的目標方向.
我們不是要做一晚的聚會.
我們是要預備人心.
領受耶穌基督的救恩.
十二月是一個很好的時候.

$^{521}$不過不是一個時候的終結.
應該是一個時候的開始.
讓人能夠領受福音.
有一個新的生命.
新的心靈.
來面對新的一年.
嘉欣樓也是步向一個新的開始.
我們要有一個更加同心合意的心.
大家回來一起朝著一個共同的目標方向.
為著一個更美麗的地方.
這算不好嗎?.
很好.
但福音不單止在耶路撒冷.
施洗約翰是去到曠野地.
曠野也是一個很重要的地方.
我們要去的地方.
不一定是曠野.
但可能是福音的曠野.
可能是一個很需要福音的地方.
是很需要我們將神的愛帶到那裡.
我們是有神的使命的.
是神差遣我們的.
這段時間我們在這裡預備.
將來我們要去一個更美麗的場景.
更美麗的地方.
將整個福音帶給那裡的人.
那裡的人從哪裡來呢?.
我們不需要擔心那個地方在哪裡.
神預備到的.
神差遣到的.
人就由耶路撒冷.
從大全地到曠野地.
去領受救恩.
領受洗禮.
只要是神預備.
神差遣.
其他什麼都不需要我們去負責.
我們就是做我們的使命.
施洗約翰做什麼呢?.
他只是傳講.

$^{561}$悔改的福音.
悔改的洗禮.
甚至他不知道耶穌基督是不是真的那位要來呢?.
當他被收監的時候.
他請門徒去找耶穌問他.
你是不是真的就是那位啊?.
沒錯 耶穌基督就是那位.
有時候我們會有些時候覺得迷茫.
但不要緊 我們認定方向.
神要做的.
祂繼續會完成.
我們只是在中途完成我們一部分的使命.
整個工程.
天父去處理.
求天父幫助我們.
一起低頭 叩叩.
親愛的天父.
上帝我們來到你的庚前.
我們祈求你看顧.
主啊 這個是你的家.
家恩家.
這個是你所差遣的.
主啊 願你幫助我們.
讓我們每一位.
我們滿懷使命.
知道上帝你的指引.
我們朝著你引導的目標方向前進.
最後讓你所預備的民.
來到我們當中.
讓你的福音廣傳.
讓你所預備的人.
能夠在這裡找到家.
找到愛.
找到兄弟姐妹.
找到被你牧養的地方.
找到安寧的清水.
找到你給他們的草.
仰望你奉主耶穌基督的聖名.
\newpage



\section{以賽亞書 11:1-9}
\label{sec:m2hnqzcprEM}
\textbf{2023-12-10 午堂《共融的新天新地》以賽亞書 11\_1-9 講員:陳麗蟬牧師}
\newline
\newline
連結: \href{https://youtube.com/watch?v=m2hnqzcprEM}{\texttt{https://youtube.com/watch?v=m2hnqzcprEM}} ~~~~ 語音日期: 2023-12-12
\newline
\newline
\hyperref[sec:KJG_idsYK_w]{\small{< < < PREV SERMON < < <}}
~
\hyperref[sec:index_chronic]{\small{[返順時目]}}
~
\hyperref[sec:index_scriptual]{\small{[返順卷目]}}
~
\hyperref[sec:jzt3Ohpru70]{\small{> > > NEXT SERMON > > >}}
\newline
\newline
以賽亞書 11:1-9
\newline
\begin{longtable}{cl}
\hline
\hline
章節 & 經文 (和合本修訂版)\\
\hline
11:1 & \begin{tabularx}{0.7\textwidth}{X} 從耶西的殘幹必長出嫩枝，他的根所抽的枝子必結果實。 \end{tabularx} \\ \\ \relax
11:2 & \begin{tabularx}{0.7\textwidth}{X} 耶和華的靈必住在他身上，就是智慧和聰明的靈，謀略和能力的靈，知識和敬畏耶和華的靈。 \end{tabularx} \\ \\ \relax
11:3 & \begin{tabularx}{0.7\textwidth}{X} 他必以敬畏耶和華為樂；行審判不憑眼見，斷是非也不憑耳聞； \end{tabularx} \\ \\ \relax
11:4 & \begin{tabularx}{0.7\textwidth}{X} 卻要以公義審判貧寒人，以正直判斷地上的困苦人，以口中的棍擊打全地，以嘴裡的氣殺戮惡人。 \end{tabularx} \\ \\ \relax
11:5 & \begin{tabularx}{0.7\textwidth}{X} 公義必當他的腰帶，信實必作他脅下的帶子。 \end{tabularx} \\ \\ \relax
11:6 & \begin{tabularx}{0.7\textwidth}{X} 野狼必與小綿羊同住，豹子與小山羊同臥；少壯獅子、牛犢和肥畜同群；孩童要牽引牠們。 \end{tabularx} \\ \\ \relax
11:7 & \begin{tabularx}{0.7\textwidth}{X} 牛必與熊同食，牛犢與小熊同臥；獅子與牛一樣吃草。 \end{tabularx} \\ \\ \relax
11:8 & \begin{tabularx}{0.7\textwidth}{X} 吃奶的嬰孩在虺蛇的洞口玩耍，斷奶的幼兒必按手在毒蛇的穴上。 \end{tabularx} \\ \\ \relax
11:9 & \begin{tabularx}{0.7\textwidth}{X} 在我聖山各處，牠們都不傷人，不害物；因為認識耶和華的知識要遍滿全地，好像水充滿海洋一般。 \end{tabularx} \\ \\
[1ex]
\hline
\hline
\end{longtable}
$^{1}$弟兄姊妹早晨.
在這兩個星期大家回來教會.
發覺到在禮堂裡面開始有些聖誕的佈置了.
上星期開始有些很漂亮的天使.
今天又有耶穌出生的聖經掛了上牆.
在禮堂的門口我們又有聖誕樹了.
這一切一切都告訴我們聖誕將要臨到了.
很多時候踏入十二月的時候.
我們就會聽到一些領唱或者主席說.
今天是張臨期的第一主日.
今天正正是張臨期的第二主日.
什麼叫做張臨期呢?.
張臨期是一些禮儀教會.
他們就在這個用聖誕節12月25號.
就向前推前四個星期.
連續四個星期就叫做張臨期.
這個張臨期他們會有四個不同的主題.
每一個主題他都幫助信徒.
幫助教會在生命當中的一些反思.
在哪一個地方可以能夠更加像耶穌基督.
預備心靈迎接主耶穌基督的降生.
但是耶穌基督的降生.
就在二千多年之前已經發生了.
現在已經過了.
現在我們就要預備我們的生命.
迎接主耶穌基督的再來.
因為主耶穌基督再來的時候.
祂所建立的國度是一個新的國度.
我們經常都說這是一個新天新地的國度.
在當中是和平和諧.
彼此接納是共融的一個世界.
剛才朱太我們讀出的《以賽亞書》11章1至9節.
就形容這個新天新地的情景.
我們今天藉著這段經文.
我們一起學習 一起彼此勉勵.
先看一下當中的背景.
我們知道以色列國.
時時都不遵守耶和華的律例.
一方面時時拜偶像.
沒有尊耶和華是獨一的真神.

$^{41}$另一方面耶和華在申命記載裡.
教導很多社會當中.
要怎樣和周圍的人.
怎樣對一些貧寒的人.
但是以色列這幫人們.
這幫上帝的子民.
仍然在這個社會當中.
沒有行上帝教他們的公義.
也都在這個社會當中.
沒有行出上帝的人賜.
來照顧周圍有需要的人.
所以耶和華就要興起不同的外邦.
來去管教以色列.
當時興起的外邦民族就是亞述.
亞述當時是一個剛剛新出的國家.
不過它興起得很快.
沒多久就成為一個強國.
來去威脅著以色列.
剛才讀十一章.
在《以剎書》第九章和第十章那裡.
上帝已經不斷發出警告的信息.
他四次這樣說.
他說耶和華的怒氣還未轉燒.
他的手亦繼續伸出.
沒有收回.
即是說那幫以色列人.
無論先知他講了多少的教導.
但是他們都沒有回轉的想法.
所以上帝說他們一天不回轉的話.
他的怒氣都不會消去.
他的手 他的管教審判都不會收的.
耶和華是要一幫子民.
屬於他的子民.
那個生命是要徹徹底底的改變.
要回到創世的時候.
他所給人的生命.
大家彼此相愛.
彼此和諧地生活.
所以整個《以剎書》的主題就是.
審判 管教和救贖.

$^{81}$這個審判不是將人置諸死地.
這個管教不是刻意罰到他痛 要他死.
而是目的是要將他們拯救.
更生他們整個的生命.
在第十一章 就是這個很寶貝的應許.
就是某一天 以色列人.
他們一定會完全的回轉回到上帝那裡.
到時上帝一定的拯救.
這個和諧和平的國度.
就會臨到這個地上高了.
我們先看一下.
這個新的天地.
新天新地的管治者.
是一個怎樣的神.
在第一至五節講到.
他說是一位有智慧有聰明的靈.
這個智慧和聰明的靈.
是講到這個新天新地的掌管者.
他的判斷力非常的準確.
他不會去搞錯事的.
還有他的洞察力非常敏銳的.
可以看得透人的眼眉.
也看得透人的心靈.
他知曉世界所發生的事情.
包括了人的心靈.
另外他也是一個很有謀略很有能力的靈.
這個謀略的靈.
是講到他策劃這個事情.
計劃這個事情.
很有條理.
有計劃.
很周詳的.
只要有他在.
只要有這個管治者.
在這個新天新地裡面.
一切的事情都可以解決.
一切的事情都可以迎刃而解.
最後就是他一位有知識的靈.
也是驚威耶和華的靈.
就是講到他對上帝的心意是很清楚的.

$^{121}$他和父神也都有一個很緊密的關係.
所以他知道上帝的心意.
耶和華的心意.
是要人行正直.
行公義.
所以他就照著上帝的吩咐.
行上帝要求人的心意的生活.
從這兩三次的經文裡面.
我們看到了.
這位管治者.
新天新地的管治者.
是一個很厲害的靈.
他可以說是一個萬事通.
他也很清楚上帝的心意.
所以行一切的事情.
都是按照上帝在一路的教導.
在《新命紀》裡面的教導裡面.
來行的系行.
他知曉人的心靈.
在年輕的時候我也會寫日記的.
長大了就沒有寫了.
但是在手機裡面.
現在很容易.
在手機裡面.
很多時候我都會將一些特別的事情.
來記下來.
在幾年之前.
有一些令我很激氣.
很憂心的事情.
我就覺得很激氣.
我就覺得.
不知道為甚麼其他我都沒有.
我就要將整件的事情.
來記下來.
所以有些甚麼任何.
我覺得很重要的對話.
我都將它打粒星.
你知道我有時.
這些東西都不會不見了.
我希望整件的事情.

$^{161}$今天發生過甚麼事.
第二天又有甚麼話說過.
第三天或者一個月之後有甚麼事.
我全部都整理得很整齊.
在三年前.
我又上了杜豐山裡面.
讀了一個靈修指導的課程.
基本上整個課程.
都教我們怎樣處理生命的.
牧師就經常說.
在第一年裡面不斷地說.
我們生命裡面有很多雜質.
很多垃圾.
要將這些雜質.
這些垃圾全部清走.
唯有全部清走這些雜質的時候.
我們的生命才可以改變.
才可以更生.
倘若在裡面仍然有很多垃圾.
無謂的東西的時候.
我們是更生不了的.
我們的生命是比不上耶穌的.
當整年的都是這樣教導的時候.
曾經有一刻鐘.
腦海就閃了出來.
你在手機裡面儲存了很多.
你認為很憤怒的事情.
你沒有刪除了.
你繼續在這裡.
當時我就回應了.
我說 這些事情很不公義.
沒有公理.
不僅對我 對其他人都是.
我一定要記得它.
我一定要記住它.
以致將來有些什麼事情的時候.
我要拿出來做證據.
其實我的心思不是很好.
我想用這些證據來審判人.
每一年開學的時候.

$^{201}$我們就要去杜豐山四天的營地.
所以第二年的時候我們又去營地了.
第二天起床.
一開機.
那麼多的機都開不上.
昨晚臨睡前十點多都開完機.
回覆了很多WhatsApp的.
為什麼今天開不了呢?.
就知道 哎呀 慘了.
一定壞了.
但是呢 因為去退休了.
就不准下山的.
你還要下山修理手機.
是不行的 是吧.
所以 哎呀 慘了.
要忍三天沒有手機用.
都已經借了同學的手機來打電話.
找梁銓道救命.
梁銓道 怎麼辦.
我的手機有沒有救呢?.
他就說 看看吧.
通常都有得救的.
開一下手機 看看拿回裡面的資料.
可能其中一些都沒有了.
所以一完了營地就快點走下山了.
但是結果呢.
整個手機都燒了.
梁銓道問我 你做過些什麼事?.
我說我沒有做過什麼.
我沒有做過什麼.
那麼就給了60元還是120元開機費.
試一下可以找一些資料出來.
裡面的照片拿出來 只有半張.
還有其他東西 什麼都沒有了.
哇 什麼都沒有了 全燒了.
然後聖靈就提醒我了.
祂說 你在課程裡面.
你又說要更新生命.
又說要致死路我.
又說要將生命的雜質全扔掉.

$^{241}$你在每個星期一的課堂裡面.
跟同學又為小組祈禱.
在營裡面這麼多天裡面.
又由你去祈禱 不斷地說.
但是你卻靜靜地.
將你裡面的那些憤怒.
覺得不公義 不公平的事情.
埋在手機裡面.
你埋在手機裡面.
就是埋在你心靈的當中.
並且你有個小小的思想.
等候時機來報復.
我覺得上帝這個靈這麼厲害.
祂是智慧和聰明的靈.
祂識破了人心靈裡面所有的事情.
我想自己好像是一個人格分裂.
說了投資三年的時間.
每逢星期一 隔個星期一.
就想去學習整理生命.
但在暗地裡靜靜地將很多的東西.
髒髒的那些垃圾不肯扔掉.
扔掉了很多的.
不過總是有些.
我自己覺得不順氣的東西不肯扔掉.
就自己擠埋了.
上帝其實是允許我禱告的.
是幫助我的.
但還有一些我不肯交出來的時候.
就好像我老師可能他真的.
不過他不知道我的.
他說我經常質疑.
有一個寶箱裡放了一些東西.
不肯拿出來處理的.
就是放在寶箱裡.
我不肯處理.
上帝就出手來處理.
整架機全部銷毀.
什麼都沒有了.
到今天.
你問我記不記得整件的事情呢.

$^{281}$我告訴你 我記得.
不過呢 是不是這麼清楚呢.
有很多東西忘記了.
就那些次序都忘記了.
忘記了.
但是呢 又感恩.
當我想起來的時候.
我沒有再憤怒了.
我沒有再覺得激氣.
我都沒有再覺得.
我需要來做審判.
我需要到時候.
我就要交出來做這個證據.
人生裡面有很多很多東西.
為什麼要自己.
令自己的生命這麼辛苦呢.
聖靈對我的管教.
最初很難過的.
不過難過的是.
裡面不見了很多很美的相片.
只剩下半張紅茶茶的.
但是聖靈對我的管教就是.
祂要你去將我的生命.
救贖我的生命.
給一個乾淨的生命.
給一個更加像耶穌基督的生命.
還有 這個靈.
這個管治者的靈是怎樣呢.
祂是肯審判 不用這個眼見.
斷這個是非.
不會憑這個耳聞.
祂處理所有的事情.
都是按祂的公義.
按祂的正直.
來去處理 去審判的.
人就是這樣的.
我們眼看到一些事情的時候.
我們就斷定.
這個不是好人來的.
這個不是我的朋友來的.

$^{321}$有一天.
我和幾個姐妹.
我們去吃早餐的自助餐.
有位姐妹出去.
在倒咖啡的時候就問我.
「喂 順便幫你倒杯吧」.
我說「好啊 我要黑咖啡」.
她回來之後.
就倒了一杯這樣的東西給我.
我說「我要黑咖啡」.
「為什麼會有奶的?」.
我說「我是按了機器 按了黑咖啡」.
「不過不知道為什麼出來是這樣的呢」.
我心裡面有些不高興.
我覺得.
一去就夢察察.
不知道按了什麼按鈕.
另一方面.
就可能她覺得.
喝黑咖啡苦不好.
給你奶喝吧.
我自己判斷的.
這個都是我自己判斷的.
當時我心裡想.
「唉 我不想喝甜的東西」.
「今天早上才喝的」.
旁邊的一位姐妹就說.
「這個是黑咖啡來的」.
「你用匙羹去攪一攪它吧」.
一攪它出來的時候.
原來真的是黑咖啡來的.
攪完之後再喝一口.
真的是苦的.
一點甜味都沒有.
我就去反省.
人就看.
「哇 為什麼會有啡色的東西」.
「好像奶的在表面的」.
「就不是黑咖啡了」.
但我不懂.

$^{361}$我不懂那些咖啡機.
有些會按了出來.
原來有些東西是按在表面的.
人看外貌.
它就不是黑咖啡.
誰知道原來下面的裡面.
就真的是苦苦澀澀的黑咖啡.
人是看外貌的.
耶和華是看內心的.
在那次裡面.
上帝又給我一個很好的教訓.
最後他就說到.
這個本賜者的靈.
它是一個知識和敬畏耶和華的靈.
它一定會以敬畏耶和華為榮.
這個就說到彌賽亞.
就是他向上帝指示的.
那個按公正公義.
來審判本賜整個的世界.
我們知道在《腓立比書》裡面.
就說到.
當耶穌基督信服父神的旨意的時候.
他已至於死.
且死在十字架上.
然後上帝就將那個.
超乎萬名之上高的名.
都給了耶穌基督.
然後就無論在天上的.
地上的.
地底下.
一切受造之物.
都無不口稱耶穌基督是主.
然後下一句是什麼呢?.
使榮耀歸於父神.
所以耶穌基督.
祂知道是聖子的身份.
使到榮耀歸於父神.
祂所作的一切.
都是以敬畏耶和華為樂.
弟兄姊妹.

$^{401}$我們的生命.
我們今天的生命.
我們以什麼為樂呢?.
以跟隨耶穌基督的教導.
活出祂對我們的教訓為樂.
還是以這個.
西族的裡面.
以這個物質.
以這個金錢.
以這個地位.
來為樂呢?.
第二方面.
我們就要看這個.
新天新地的新生活裡面.
是怎麼樣的生活.
我們看到這個圖畫.
是不是很開心啊?.
這個迪士尼啊.
你看到在這幅圖裡面.
有老虎.
有紅人.
也有豺狼.
他說.
這些全部都是很凶殘的動物.
但是卻是很活潑.
很可愛.
又可以跟老鼠.
跟鴨子.
跟狗一起.
這麼生活.
這樣在這裡生活.
原來在這個新天新地裡面的新生活.
是和諧共用的.
在當中.
一個很兇猛的野獸.
可以跟一些很純良的動物.
一起生活.
這個野菜仙子就說.
他看到的異象就是.
哇 那些很狡猾.

$^{441}$很兇的豺狼.
可以跟綿羊一起住的.
那些豹.
吃人的豹.
可以跟山羊一起睡覺的.
那些獅子.
又可以跟牛毒.
那些肥畜.
可能那些有鴨.
肥鴨.
肥雞.
一起玩的.
那些小朋友又可以牽著蛇.
一起走路的.
剛才完了總堂之後.
又在弟兄那裡說.
他說陳牧師.
你說錯了一些東西.
他說蛇是不會咬人的.
我說是嗎.
就算咬不咬人也好.
我也不會跟牠玩的.
他說蛇只會吞人吞東西的.
他說.
那好吧.
不會咬人.
但是會吞人.
這樣.
但是你看到.
很開心的一個景象.
就好像你看迪士尼片.
這麼開心的.
我記得有一次行山.
我和我媽媽一起行山.
下山的時候.
就是看到兩條青竹蛇.
就是站在那裡.
然後舌頭伸出來.
哇.
如果我那時候知道.

$^{481}$蛇不會咬人的時候.
可能沒有那麼害怕.
但是牠們真的會腳軟的.
我和我媽媽說.
糟了.
向前走.
還是怎樣.
我們就決定.
不要向前.
慢慢走.
向後.
但是發覺那兩個都走不動的.
是腳軟的.
唯有慢慢熬.
就向後走.
見到任何人.
都叫牠不要再下去了.
但是.
這些這麼恐怖的情況.
當上帝的公義來到.
新天新地降臨的時候.
一切的事物都會安定下來.
那些很凶殘的猛獸.
都會順樑下來.
牠可以一起和那些順樑的動物.
一起生活的了.
在現實的世界裡面.
的確有很多很凶殘的野獸.
牠們不可以和這些牛毒,羊膏一起的.
否則就會變成很香噴噴的牛扒和羊扒的了.
但是上帝在創造人的時候.
祂的心意.
是叫萬物都能夠和諧,和平.
融洽地相處.
祂不希望那些野獸.
或者任何人去傷害其他的野獸.
或者傷害其他的人.
不過人的本相就是這樣.
總是喜歡自己最厲害的東西.
總是喜歡威猛.

$^{521}$總是喜歡自己出名.
所以看到其他同樣威猛的人.
我們就會容易大家去比拼.
大家去批評.
大家要來將對方踢出局.
看到一些沒有那麼聰明的人的時候.
就會覺得那個人不懂事.
什麼都不懂的.
是低人一級 是低B來的.
整個羊菇一樣.
有時候我們就會不作聲.
心裡面也會這樣去批評.
在過去五天.
我就和兩位姐妹一起去旅行.
我都跟她說.
今天我講到很感恩.
在五天的旅行裡面.
我就帶了我的經文.
帶了我早些找到資料.
就預備我的講章.
我就說上帝給我預備了這篇講章.
我就跟她說.
我會將這個經歷跟她們.
跟弟兄姊妹分享的.
她們OK的.
這三個人的背景.
A 阿A就在美國長大的.
自由獨立慣的.
他也周遊列國.
去過很多很多的地方.
他不跟旅行團的.
自己走 自己坐巴士走.
他去到日本已經五六次這麼多了.
自己去的.
B 就是在香港長大的.
他整個都沒有什麼所謂的人.
總之不是得罪神.
總之在真理範圍之內的事情.
他都覺得沒什麼所謂的.
可以做的.

$^{561}$C 就是在印尼長大的.
來到香港工作.
我很體會.
這次不是猛龍就不會過江的.
這次去日本.
是第一次去長假期.
五天來說.
去這麼美麗的地方.
他很興奮.
第一天我們就用了很多時間.
到第二天就正式去旅行.
他看不同的景點.
阿A一去到的時候.
他就看這裡很多紅葉.
不過山上更多.
走快一點上山的地方.
那些紅葉更美麗.
更多景色看.
C他還沒看過紅葉.
他在門口已經覺得.
哇 很美麗.
所以又要拍很多照片.
還要拍影片.
留給他的姐妹看.
於是一個就大聲說.
快點向前走.
上高很美麗.
下面這些只是垃圾.
不用看那麼多 快點走.
C就越逼我 我就越不走.
B他沒所謂的.
上山還是在這裡.
就不知道怎麼算了.
接著說一下搭車.
A和C大家都很有意見的.
但總是有不同的意見.
在搭車的裡面.
A就說.
這些線 這些線.
這樣又去到那個地方了.

$^{601}$C就說.
這樣 這樣 這樣都可以去到.
但大家都是拿著同一份的地圖.
大家拿著不同的手機.
於是不同的手機有不同的路線.
告訴你可以去到那個地方.
所以A就說要這樣去才對.
C說這樣去才對.
我的手機是這樣說的.
我就想一下.
其實就好像香港的鐵路線一樣.
在元朗去灣仔.
很多條線可以走的.
由元朗可以搭西鐵.
在美孚轉車也可以.
在尖沙咀轉車也可以.
去到紅磡轉車也可以.
在紅磡轉車.
朱小樓忘記了.
你可以向上高走.
就轉火車.
還是向下低.
兩樣都可以 看你回不回去.
然後轉了火車.
就過到港島.
去到會展站下車也可以.
又或者你再搭多一個.
去到金鐘站下車也可以.
所以就看你跟哪條線而已.
是不是.
所以這樣之下.
怎樣搭哪條線.
都有很多的討論.
我發現.
大家的臉色.
就開始收了一些笑容了.
心想 慘了.
第二天的下午.
怎麼熬啊 還有四天.
我們大家都是主裡面的弟兄姊妹.

$^{641}$有耶穌基督在我們生命的當中.
就有新天新地.
所以那天的晚上.
我們在酒店裡面.
彼此接納的場景就出現了.
A首先發言.
他說.
「不好意思」.
A說「不好意思」.
「我不可以將我的要求」.
「擠到你們的要求裡面」.
「我沒有體諒」.
「你第一次來日本」.
「你什麼都沒看過這些紅頁」.
「所以你想將裡面的記憶」.
「全部擠在裡面」.
C也說.
「是啊」.
「不好意思」.
「我只顧著拍照」.
我覺得大家彼此去學習一件事.
我不會看我自己比另外一個人強.
我覺得這條路線是好一點的.
C也不會再看A那條線.
覺得他那條線不行.
而大家彼此走哪一條路呢?.
對吧?.
我說好 第二天.
你繼續向山上的高處.
C 你就要學習一件事.
繼續向前走.
山上的風景是真的更加的美麗.
不就只是停留在閘口.
走上半山看到.
原來真的比原來更美麗.
之後的幾天.
大家就能夠去協調得到.
然後更加美麗的就是.
今天我想去購物.
好啊 你去吧.

$^{681}$那就繼續上去自己的景點.
大家都會看手機去那條路線的 是吧?.
就由一天彼此分開.
到吃飯坐在一起.
到最後在機場裡面.
因為A要回美國了.
他就跟C說對不起.
他說對不起.
我真是在那天第二天那裡.
我有點兒臉色.
A就說對不起.
我經常只顧著拍照.
C就說對不起.
我經常只顧著拍照.
A就說人美麗.
當然要多拍些照片.
我不美麗.
所以我只要把風景放進腦袋就可以了.
然後我們三個人就一起擁抱.
我們一起的祈禱.
我覺得在當中我們很感恩.
我們在第二天的酒店裡面.
我們自己已經彼此講了自己的一些要求.
然後我們彼此接納對方.
我們彼此怎樣去協調.
免得這麼多天都臉色.
到最後那天才離開去認錯就慘了.
但上帝讓我們在第二天發現的時候.
我們就做了.
以致接著那三天我們很開心.
第四天分開的時候.
我們再一次的彼此希望對方接納.
一幅很美麗的圖畫.
然後我們還相約.
明年我們在印尼再見.
這就是一幅很美麗的新天新地.
這個共融和諧的場景.
這個就是凶惡的猛獸和順量的動物.
可以一起的生活.
這個共融.

$^{721}$這個新天新地裡面的共融是甚麼意思呢.
就是一些人.
大家有不同能力.
有不同背景的人聚在一起.
在當中我們會有共通的地方.
我們喜歡旅行.
我們喜歡美麗的東西.
但也有很多的差異.
有些人喜歡上山.
有些人就停留在那裡就行了.
但我們怎樣一起的生活.
一起的工作.
事奉的時候.
我們怎樣在很多不同的差異裡面.
我們怎樣努力去衝破這些障礙.
我們首要的是要互相的接納.
接納大家的不同.
我們都不將自己的要求.
希望對方都一樣.
我們互相的欣賞.
互相的尊重.
新天新地.
不是要去到將來.
耶穌基督再臨這個時候才有.
耶穌基督就是要我們在今天的裡面.
將新天新地裡面.
那個的共融.
那個的和平.
就要在這個裡面裡活出來.
在我們的家人.
在我們的同事朋友當中.
裡活現出來.
中國人有一句話.
是要「修身齊家治國平天下」.
即是說.
當我們在外面.
在這個工作.
在這個教會裡面.
我們做得多麼好.
但我們心靈裡面.

$^{761}$用一個寶盒.
裝了我們一些的老鵝.
裝了一些的垃圾.
在這裡的時候.
這個外面的和諧共融是假的.
如果當我們和我們的家人.
都沒有很好的關係.
在外面的共融.
都是假的.
所以我們首先要想一下.
反省一下.
我們和我們自己內心的世界.
我們有沒有真的.
將我們一些不好的.
一些的質地.
生命裡面的質地.
甚至自己都不肯去面對的東西.
我們肯不肯交出來.
由耶穌基督來管理.
由耶穌基督來更生.
然後下一步.
就是看一下.
我們在當中.
我們和我們家人.
我們對我們的父母.
和我們的父母的關係.
有沒有和諧呢.
有沒有尊重呢.
我們夫婦之間.
和妻子和丈夫之間.
我們有沒有彼此的相愛.
彼此的支持呢.
我們和子女.
和我們的親戚又怎麼樣呢.
修身齊家.
然後我們才平過.
我們才會和別人說.
在教會.
在工作的範圍裡面.
我怎樣可以做得到.

$^{801}$所以求主耶穌基督幫助我們.
以致我們在彼此之間.
我們不再有仇恨.
不再有恐懼.
我們不再有批評.
不再有鬥爭.
彼此之間我們就去接納其他人.
和我們一起是不同的.
我們怎樣彼此的要恕.
彼此的寬恕.
唯有這樣的時候.
這個新天新地就會出現.
在這個家庭.
在教會工作的範圍裡面.
這個新天新地就會出現.
所以這個章林奇.
說要整理生命的意思.
就是這樣.
還有.
在這個新天新地裡面.
吃奶的嬰孩.
在這個毀蛇的.
這個毒蛇的洞口玩.
這個斷奶的嬰兒.
也都按守在毒蛇的穴上.
我們說 嘩.
如果當我們作媽媽.
作婆婆公公爺爺的.
我們說 嘩.
為什麼有個嬰兒要和毒蛇玩.
一定大叫救命.
看到新天新地就行了.
新天新地就行了.
我們看到這裡說到嬰孩.
說到幼兒.
在這個二賽書裡面.
經常說到有這些的景象出現.
這個嬰兒.
這個小朋友.
給人的感覺.

$^{841}$就是一些軟弱.
沒有能力.
要人照顧的嬰兒.
是不是?.
但是呢.
原來在二賽亞先知裡面.
他經常用這個嬰兒.
小孩子.
去指向彌賽亞.
在這個二賽書九章六至七節裡面.
他更加用這個小孩.
去指向將來的彌賽亞.
是怎樣的上帝.
我們一起讀這兩節的經文.
1,2,3.
因有一嬰孩為我們而生.
有一子賜給我們.
政權必擔在他的肩頭上.
他名稱為奇妙賊士.
全能的上帝.
永在的父.
和平的君.
他的政權平安必加增無窮.
他必在大位的寶座上.
治理他的國.
以公平公義.
使國堅定穩固.
從今直到永遠.
他還有一句就是.
萬君之爺和華的熱心.
必成就這事.
這個小孩.
給人的感覺是軟弱.
沒有能力.
不懂照顧自己的.
但是呢.
原來這個指向彌賽亞.
到這個降臨的時候.
他就是那位.
滿有智慧聰明的靈.

$^{881}$滿有謀略能力的靈的彌賽亞.
我們只要將生命.
給他們去管治.
肯給他們去整理的時候.
於是耶穌基督的人愛.
耶穌基督的和平.
耶穌基督新天新地的共融.
就會在我們生命裡面.
立即發揮光輝.
保羅教導我們.
不可陷路到日落.
下一句是什麼呢.
在《伊佛所說》.
不可陷路到日落.
下一句就是.
不可給魔鬼留地步.
所以原來當我們.
不肯學習耶穌基督.
對世人的饒恕.
對世人的接納.
我們都不肯在我們自己生命裡面.
裡活出來的時候.
就是給魔鬼留地步.
損害了耶穌基督的名字.
所以弟兄姊妹.
求主幫助我們.
在這個.
就說回這個嬰孩那裡.
就是這個.
第一節.
他說從耶西的殘幹.
必長出這個聯誌.
這個耶西.
都不知道是誰人來的.
原來就是大衛的爸爸.
為什麼以賽亞才知道.
他不說從大衛的後代裡面.
必長出這個聯誌.
這個聯誌一定會有果實.
其實就是在說.

$^{921}$從大衛的後裡面.
後裔裡面.
有尼賽亞的誕生.
但為什麼上帝要感動以賽亞才知道.
要寫這個從耶西的殘幹呢.
耶西是一個寂寂無名的人.
沒有人認識的.
不過只要上帝的權能臨到的時候.
這個寂寂無名的人.
他都可以生出一個.
這麼有能力的.
大衛成為了以色列第二任的皇帝.
也都在之後.
之後能夠生出尼賽亞.
上帝就要告訴我們.
不要小看我們自己軟弱.
不要整天覺得我們軟弱.
覺得我們沒力.
只要上帝的權能臨到的時候.
我們也都可以為上帝去成就大事.
只要我們願意的時候.
上帝會工作.
他就說到.
在我聖山各處.
這一切都不傷人 不傷 不害民.
因為認識耶和華的知識.
要充滿遍地.
好像水充滿海洋一半.
在將來新天新地.
在今天我們活現這個新天新地的時候.
這個和平共融.
就會臨到在這個地上高了.
我們一起唱這個.
「以馬來利懇求降臨」.
我們唱第二第三節.
「耶西之象懇求降臨」.
我們一起來唱吧.
「耶西之象懇求降臨」.
(耶西之象懇求降臨).
(撒旦心中釋放自問).

$^{961}$(地獄心處清皎自問).
(始終得食食沒自憤).
(困人困人 已死亦亡).
(以馬來利懇求降臨).
(情神若看客求降臨).
(只知江南困壞眾生).
(燥火長夜又暗又慘).
(來來光明皆散後還).
(困人困人 已死亦亡).
(以馬來利懇求降臨).
我們同心土局.
我們來去各自的安靜.
想一想在我們生命當中.
有沒有一些家人.
或者我們的同事.
我們的弟兄姊妹.
我們覺得我們很憂心.
不可以接納得到的呢?.
我們願不願意將他們都交給耶穌?.
你都願不願意將我們自己的生命都交給耶穌.
由耶穌來去管治的呢?.
只要我們肯的時候.
耶穌基督的新天新地.
就會在我們的生命當中.
彰顯出來.
親愛的主耶穌基督.
我們多謝你.
原來在你自己的政權管治底下.
是這麼美好的.
是可以這麼和諧.
這麼共融的.
我們多謝你.
主就求你.
賜給我們勇氣.
賜給我們力量.
去整頓我們的生命.
去治死我們裡面的老鵝.
一些不肯交出來給你們去管理的老鵝.
我們交出來吧.
幫助我們.

$^{1001}$主要就是在我們生命的裡面.
我們真的活出主你的和平.
主你的和諧.
求主工作.
以致我們真的可以用一個喜樂的生命.
去迎接主耶穌基督的降臨.
每一年我們去迎接主你的誕降.
紀念你的拯救.
直到將來.
我們要迎接主你的臨在.
我們多謝主耶穌基督.
多謝主的工作.
禱告 祈求 感恩.
奉主基督的聖名.
阿們.
\newpage



\section{約書亞記 14:6-15}
\label{sec:jzt3Ohpru70}
\textbf{2023-12-16 晚堂《生命的投資》約書亞記 14\_6-15 講員:蔣文忠傳道}
\newline
\newline
連結: \href{https://youtube.com/watch?v=jzt3Ohpru70}{\texttt{https://youtube.com/watch?v=jzt3Ohpru70}} ~~~~ 語音日期: 2023-12-19
\newline
\newline
\hyperref[sec:m2hnqzcprEM]{\small{< < < PREV SERMON < < <}}
~
\hyperref[sec:index_chronic]{\small{[返順時目]}}
~
\hyperref[sec:index_scriptual]{\small{[返順卷目]}}
~
\hyperref[sec:oC6cD_CKYi8]{\small{> > > NEXT SERMON > > >}}
\newline
\newline
約書亞記 14:6-15
\newline
\begin{longtable}{cl}
\hline
\hline
章節 & 經文 (和合本修訂版)\\
\hline
14:6 & \begin{tabularx}{0.7\textwidth}{X} 猶大人來到吉甲，約書亞那裡，基尼洗族耶孚尼的兒子迦勒對約書亞說：「耶和華在加低斯‧巴尼亞指著我和你對神人摩西所說的話，你都知道。 \end{tabularx} \\ \\ \relax
14:7 & \begin{tabularx}{0.7\textwidth}{X} 耶和華的僕人摩西從加低斯‧巴尼亞差派我窺探這地的時候，我剛四十歲。我把心裡的話向他報告。 \end{tabularx} \\ \\ \relax
14:8 & \begin{tabularx}{0.7\textwidth}{X} 雖然同我上去的眾弟兄使百姓膽戰心驚，我仍然專心跟從耶和華－我的神。 \end{tabularx} \\ \\ \relax
14:9 & \begin{tabularx}{0.7\textwidth}{X} 那日，摩西起誓說：『你腳所踏之地必要歸你和你的子孫永遠為業，因為你專心跟從耶和華－我的神。』 \end{tabularx} \\ \\ \relax
14:10 & \begin{tabularx}{0.7\textwidth}{X} 現在，看哪，耶和華照他所說的使我活了這四十五年。當以色列人在曠野飄流的時候，耶和華曾對摩西說了這話。現在，看哪，我已經八十五歲了。 \end{tabularx} \\ \\ \relax
14:11 & \begin{tabularx}{0.7\textwidth}{X} 現今我還很健壯，像摩西差派我去的那天一樣；無論是戰爭，是出入，我現在的力量和那時的力量一樣。 \end{tabularx} \\ \\ \relax
14:12 & \begin{tabularx}{0.7\textwidth}{X} 請你將耶和華那日所說的這山區給我。那日你也曾聽說，這裡有亞衲族人，以及寬大堅固的城，或許耶和華會照他所說的與我同在，我就把他們趕出去。」 \end{tabularx} \\ \\ \relax
14:13 & \begin{tabularx}{0.7\textwidth}{X} 於是約書亞為耶孚尼的兒子迦勒祝福，把希伯崙給他為業。 \end{tabularx} \\ \\ \relax
14:14 & \begin{tabularx}{0.7\textwidth}{X} 所以希伯崙成了基尼洗族耶孚尼的兒子迦勒的產業，直到今日，因為他專心跟從耶和華－以色列的神。 \end{tabularx} \\ \\ \relax
14:15 & \begin{tabularx}{0.7\textwidth}{X} 希伯崙從前名叫基列‧亞巴；亞巴是亞衲族最尊貴的人。於是國中太平，沒有戰爭了。 \end{tabularx} \\ \\
[1ex]
\hline
\hline
\end{longtable}
$^{1}$領子妹很感恩在章林期的日子.
和大家一起在這個晚上.
一起去敬拜主.
也一起去聆聽上帝的道.
去讓上帝去提示我們的生命.
可能領子妹都知道.
有一個消息是我們敬愛的老爺牧師.
在周三在主懷安息.
求主都可以繼續去安慰.
與師母和他的兒女同在.
在他們不捨和哀傷日子.
都能夠得到主的恩典.
得到主的憐憫.
相信老牧師一生都是將他自己.
投放在主的服事.
將整個生命投放在回應主的這個職事上.
相信都很對應到我今天和大家.
想分享例如迦勒這個生命的一個故事.
我們會問就是我們自己又是.
怎樣去投放自己的生命.
或者我用的詞就是投資自己的生命.
並不是在說一些金錢的投資.
而是在說我們在屬靈上.
去整個生命怎樣投資自己.
怎樣投放自己下去主.
希望我們投入的事情中間.
剛才我們讀的經文裡面.
其中一個特色就是你見到.
迦勒去向著約書亞.
好像在做一個將整個人生去回顧.
或者一個自白一樣.
分享自己的人生.
有幾個位置是在說他自己的一些歲數.
譬如他說到自己.
這個經文是翻到他民數記13,14章.
我就不會和大家詳細回到那段經文.
但我都會有些時候和大家點題.
去提出那段經文的對理解.
這段經文裡面一些重要的地方.
當時他是40歲.

$^{41}$其實約書是同場的.
如果大家記得是12個探子.
去窺探一個迦南的美地.
如果大家留心在民數記13章的記載.
其實這12個探子是12個支派的領袖.
應該是他們看重.
會很尊重的一些當時的領袖.
派了12位去那裡.
一起去確立上帝是否像一塊迦南美地.
給以色列人將要進入去居住.
當時迦南說自己是40歲.
和其他的兄弟一起去窺探這個地方.
其實當時這段經文都是一個很重要的.
可以說是以色列人的分水嶺.
就是因為他們這班以色列人.
不去跟從耶和華的心意.
以致耶和華在那次的經歷裡面.
去和以色列人宣告.
他們要漂流曠野40年.
他們那一代的以色列人.
不能夠進入迦南地.
除了約書亞和迦納.
所以今天這段經文.
在約書亞記裡面.
就是兩個之間的對話.
過了多久呢.
如果按照迦納的字法裡面.
是過了45年.
然後按照他所說的.
讓他活了45年.
40再加45.
就是說現在已經85歲了.
所以他說到了今天.
他就覺得.
然後說曾經說過的這個應許.
就應該看今天就要藉著我.
去成就對我迦納的這個應許.
所以說如果這樣說的時候.
在這段經文裡面.
你會看到迦納將他人生規劃.

$^{81}$好像是分了幾個階段去看.
0至40歲可能是一個階段.
一個青年的階段.
40至85歲.
一個曠野漂流的階段.
然後去到85歲.
才進入人生的高峰.
進入人生正在成就.
上帝要成就生命的應許的階段.
再者頂智妹不知道你又覺得.
你會怎樣規劃你的人生.
或者會不會多不多好像迦納這樣.
你有時候會回顧一下自己的生命.
有時候想起一些聽過的故事.
譬如我自己一個學生機構的工作.
有時候二十多歲的年輕人都會說的.
很多事情譬如他站在一個十字路口.
他有些掙扎.
有時候都還沒有想得很清楚.
但是他會說我都是年輕力壯.
尚有很多可以嘗試的空間.
可以錯的空間.
毅然決定做一些事情.
我認識一個青年.
他在大學裡面本來是在讀化學的.
但一直心裡面都有一種很想去宣教的夢.
特別是去非洲.
於是他特意附修非洲研究在大學裡面.
完了之後毅然就放一年時間.
去參與一個非洲的宣教一年.
他說沒有什麼可以輸人生年輕的階段.
就去試一試這個想法.
但譬如我們去到中年的時候.
譬如大約四十歲或者四十多歲的時候.
可能我們就多了不同的考慮.
有時候都是站在一個路口當中.
究竟我繼續累積在我前半生裡面.
累積的功夫特別是在工作裡面.
還是毅然去透過這些經驗.
我們去試一些新的崗位.

$^{121}$新的挑戰.
我認識一個母親就是.
她以前是做社工的.
但就因為要教養兒女.
就毅然放下工作.
但到四十多歲的時候.
她放下十多年.
但覺得她要孩子大了.
就毅然出來再次進入那個場景當中.
去做社工.
她心裡面最初投入的時候.
都很賤勁.
但她覺得如果這一下不投放下去的時候.
就不知道什麼時候可以放下去了.
所以她覺得應該要把握時機.
去為自己的工作事業去闖一闖.
但譬如我們去到六十歲.
屈落幾歲的時候.
即將退休或者退休年紀.
可能我們有些人.
我們心裡面可能年少當中.
都還有一些未完成的夢.
我認識一個朋友就是.
剛剛退休.
他跟我說.
其實他小時候裡面很想.
其實不是預了這份工作.
做到退休的.
即是他說.
一直都有一個夢想.
很想開一間咖啡店.
好像很自由自在地.
做一些喜歡做的生意.
喜歡有空去跟客人聊天.
說一下信仰.
譬如一種很寫意.
但營運得到的生活.
但到了退休的時候.
很多上面有父母.
下面有兒女.

$^{161}$很多東西就已經不願意.
這樣冒風險了.
所以思前想後的時候.
都遲遲都未能夠去做這些決定.
去毅然踏上.
去做一個這麼大的投資.
這麼大的轉變.
不知道你們又在想什麼.
或者在人生當中.
是一個什麼狀況.
是不是一個要轉變.
轉折或者掙扎的當中.
迦納告訴我.
他無論幾個階段都好.
他的人生裡面.
都有一句在經文當中.
重複了的一個描述.
經文裡面分別是第八節.
第九節和十四節.
聖經都說他是專心跟從耶和華.
這個專心跟從耶和華.
我覺得譯得更加好.
就是我用了一個全心去跟從耶和華.
因為那個字本身是一個副詞.
其實是Fully Follow.
就是完全這樣去跟隨那個意思.
全心的意思就是說.
他沒有被其他的事情.
去佔據了他的心思意念.
這個都是不容易的.
譬如幾個經文的描述.
正正就是說他不同的階段.
無論他四十歲的時候.
當時都.
大家如果熟悉文述的記載.
就是那十個兄弟.
是回來報惡信.
衛獨,迦納和約書亞.
就站出來為耶和華說話.
那十個兄弟令到百姓都戰驚了.

$^{201}$就是膽戰心驚.
但是我專心或者全心跟從耶和華.
其實意思就是這樣.
但是我去專心全心去跟從耶和華.
到了去抗耶的日子裡面.
其實是更不容易.
你記住當上帝懲罰以色列人.
抗耶飄流.
如果大家讀文述記都很記得.
常常都反叛.
常常都後悔.
又說要回到埃及地.
又說這裡不夠吃.
那些東西又沒有肉吃等等.
你見到很多這些激怒摩西激怒上帝的經歷.
在文述當中記載.
你記住迦納一直和這班活役的子民一起.
有沒有受感染呢?.
有沒有被那些人渲染了迦納呢?.
這個有時對於迦納來說.
他想全心跟從耶和華.
其實身邊的人全部都沒有這樣做.
對他來說是一個很大的挑戰.
自己很專注走自己應該走的路.
稍後我會說說抗耶的路是一條怎樣的路.
但這個時候聖經都說他是全心跟從耶和華.
到了最後一節十四節的時候.
到了一個結局.
當時剛才說了迦納向約書亞說.
這個時候會不會是上帝要承受應許呢?.
即是四十五年前說的話會不會今天要承受呢?.
十四節結局就是希伯來成為了迦納的產業.
直到今天即是約書亞寫的時候的日子.
有一個原因的.
因為他全心跟從耶和華.
即是說無論他打贏仗得到希伯來.
或者是跟著維持產權.
或者是沒有遺棄或沒有人再次反擊攻佔他的地方.
一直維持著都是屬於迦納的地方.
整個過程裡面因為他全心跟從耶和華.

$^{241}$如果大家心水清的時候都記得.
這裡我都會說一點點那個經歷.
他全心跟從耶和華.
所以他知道那個地方.
當時報告回來的時候.
如果記得就在文書裡記載.
他們要拿一些初熟的葡萄回來.
因為他們要查一查那個地方是不是流乃與物資地.
其實是說那個地方是不是好好居住的地方.
因為是上帝給他的.
所以要檢查一下.
拿那些當地的收成回來.
其實拿回來的時候全部都是很好的.
全部都是我們吃的最大顆的葡萄.
全部都很甜美.
但是另外同時間他們見到的那些亞納族人.
也都是很巨大.
那些城也都是很堅固.
所以同時間那十個兄弟.
譬如在文書13章32節說.
那個地方是吞吃居民之地.
當他體會到上帝的心意.
原來果然是流乃與物資地.
同時間又被那些環境被那些人嚇怕了.
這個地方是吃人的.
在新約聖經就說是一個心懷異議.
一個雙重心.
他不能夠只是在維護上帝的心意.
另外他繼續被一些環境因素去嚇怕了他們.
以至最終他不能夠全心的跟從而來.
向子民報了一個惡信.
令到子民很驚動 很驚嚇.
以至整班子民去背負上帝.
這個就是當時的背景.
弟子們我們人生裡面最精壯的年歲的時候.
也有沒有全心去跟從上帝.
作出最重要的決定和投放呢.
大家都是想一想.
我們人生當中你將你的年歲投放得最久最多的時候.
在什麼事上呢.

$^{281}$有些人可能是他放在工作裡面.
或者是他自己放在自己一直追求不同的讀書學位裡面.
或者是不同的事上 或者是你的兒女上.
無論如何我們有沒有在事上.
你投放的精力 投放你的心思意念.
都是按著上帝的心意而行呢.
有沒有好像迦勒一陣我們分享的故事裡面.
他不顧一切為上帝而活呢.
這個就是我相信我們今天都要問自己的問題.
我們接著去看一看迦勒幾個人生的階段.
他怎樣都是一樣的全心的跟從耶和華.
第一個階段我剛才說了.
在他四十歲的時候.
他進入迦南地作為一個探子.
在文宿13章14章裡面都說了摩西派12個探子窺探迦南地.
其中10個就向以色列人報惡信.
其中好像文宿13章33節.
他說那10個人 所謂10個兄弟這樣說.
我們去窺探那裡看見巨人.
就是巨人中的那些亞納族人.
我們自己眼中看自己好像渣曼一樣脆弱渺小.
而在他們眼中都是一樣的形象.
即是他們連亞納族人都去窺探.
他們只是一個探子.
其實基本上他們不會自己大搖大擺在亞納族人面前.
但連他們怎樣看以色列人.
他都猜到 都了解到他們的心思.
是一種很自卑 看自己看得很低的一種價值.
連覺得別人看他都是自己很低的一種價值.
所以在這個環境 被當地的人嚇怕了他們.
他看不到上帝要建立他們.
他只見到人比人的時候.
他在人面前是一個很渺小 很脆弱的一種形象.
加納見到這樣的情況.
他施理衣服.
他告訴他們 必須領我們進入這個地方.
當然單憑加納和以色列的力量.
在當時都不能改變這些以色列人的背景.
但上帝在審判以色列的時候.
他都為了加納給了一個評語.

$^{321}$就好像文宋十四章二十四節.
唯獨我的僕人加納 因為他另有一個心志.
專心 其實都是同一個字 全心跟從耶和華.
我要領他進入他所去的那地.
他的後裔必得那地為業.
所以其實就是憑著這句話.
文宋十四章二十四節.
去到約書亞記 剛剛我們今天讀的經文的時候.
就向約書亞講回這番話.
這番話除了講了耶和華應許他.
那塊地是一定得到的.
他會唯獨加納.
全部當時的以色列人都不能進入加納.
唯獨加納他會進入這塊亞納族人住的地方.
他有個條件的 因為他說加納是另有一個心志.
是與眾不同的 與別不同的.
所以鄧小平說有時眾人捍衛美的事.
要留心去做 就是這個意思.
全部人都是活在恐慌之中.
當時沒有人夠膽覺得要進入一個迦南之地.
唯獨加納他走出來.
他領受了上帝一樣的心志.
去到上帝廁果然是迦南美地.
是一個美好之地 那就可以了.
那應該我們就進去吧 不需要懼怕那些人.
那些堅固的城 上帝必定同在 必定幫助.
這個就是另一個心志.
去看到事情的真相.
事情的真相不是環境來決定.
事情的真相是由上帝來決定.
上帝的應許要成就 雖然未成就.
但那個地方終會成為他們的地方.
因為上帝說了一定會成就這件事.
所以真相是上帝要成就在我們中間的所有事情.
當讀到這裡的時候 大家不知道有沒有留意到.
其實我們今天很多香港教會.
其實都有很多西方來的宣教士的粉絨.
以致我們今天有這麼好的教堂.
如果大家都記得 就算是感受堂的背景.
都是有西教士在感受堂裡開幼稚園開始.

$^{361}$而成立了一間宣道會.
當時的宣教運動可以說是大概由十九世紀末到二十世紀中期.
可能大家都聽過大得新 一個很出名的宣教士.
有一千條的生命都會獻上給中國內地.
可能聽過劍橋七傑 七位劍橋大學的精英.
他們毅然一起搭上船去到內地.
去將福音傳給國內的同胞.
在十九世紀最後十年的時間裡.
但如果大家都有留心讀這些宣教士歷史.
你會知道他們不是一個很浪漫的事情.
在整個宣教運動當中 真是死傷無數.
很多宣教士不要說去到一個地方.
他想去到哪裡宣教.
在船的途中都經不起那些風浪 可能已經輕生.
或者去到一個東南亞所有的地方.
基本上是一種很難適應的對地.
習慣住在歐洲或者美國的地方的時候.
來到那些風土 那些習俗.
特別很多地 香港很多的熱病 傷寒 內那些利疾那些.
其實是殺害了很多宣教士來到的生命.
如果大家去看看那些在帕馬迪的墳場.
你就會留心很多宣教士的命.
就留在那些墳地裡 你會發覺很多的生命都很短暫.
但是如果不是第一代宣教士來到香港.
或者內地或者其他東南亞地方.
他們堅持這裡果然是美德之民之地.
那樣的時候第二代第三代就不會繼續來了.
如果將一個惡信告訴你 回過頭來.
向他們的猜拳的總部說.
那裡真的是吃人之地來的.
很多疾病很多的不知什麼的.
奇難雜症很多的水路不服.
其實殺害了無數生命的時候.
他們就不會再來這個地方.
將福音帶給我們可能上一代很多的信主的前輩.
其實就是他們堅持全心跟從耶和華.
以致福音能夠傳到給很多香港,內地.
東南亞所有很多地方的信徒建立教會.
以致有自治 有自己本色本土化的教會.
能夠跟著陸續出現.

$^{401}$今天我們就不需要有西教士在當中主持教會.
這是有一個歷史的進程的恩典在中間.
看見上帝的心意.
不被環境因素去嚇怕的生命投資投放.
是我們需要去學習的.
亦都是在迦納或者我剛才所說的.
很多的來華宣教士的生命當中我們見得到.
接著迦納進入了一個曠野的人生.
有45年曠野的人生.
以前我有一個朋友.
也是我神學院的同學.
當時我不是在香港.
他開了一間曠野書店.
在油麻地區 不知道大家有沒有機會.
上過這間二樓書店.
因為開了時間不是很長.
好像開了一年多兩年左右.
就已經營運不到 就結業了.
我記得他說過.
他開這間書店為何叫曠野書店.
其實是反映他自己當時體會人生的狀態.
曠野是一個甚麼地方呢.
我有兩點要和大家分享.
第一點就是曠野是一個見步行步的歷程.
一片荒蕪.
其實你沒有東西可以見到方向.
見到一條清晰的路去走.
為何說曠野兜來兜去.
兜兜轉轉地走呢.
就是一種這樣的意思.
可能說用在人生裡面.
就是一種人生前景未明的狀態.
亦都沒有一些工作的成效.
沒有別人讚賞這些外表上.
去遮蓋我們內心裡面一些無助.
一些軟弱的狀態.
不知道大家有沒有看到一套戲叫白日之下.
是一套好看的.
很值得看的電影.
我不是想說戲裡面的內容.

$^{441}$我只是說其中一幕令我都有觸動.
就是有一個角色.
他落了一種掙扎裡面.
究竟他應該要將一生的時間.
投放在工作需要謀生的裡面.
還是在照顧他的家人的裡面呢.
我都記得一個我的朋友.
他說他本來都是在做傳道工作.
他進入了一個很熟練.
亦都很有成效的時候.
他毅然要放下.
因為他要放下時間照顧他的母親.
他不知道照顧多少年.
但他覺得這個時候他要珍惜這些時間.
或者是他自己習無旁貸.
都需要做一件事.
所以他毅然放下工作.
投入了一個照顧母親.
每天都是一路那種程序.
那種時間表.
這樣陪伴他母親.
可能進出醫院.
要預備各樣的生活.
飲食 起居等等.
需要幫助他的事情.
這個不是他自己本身應該.
按照人生發展裡面要做的事.
但是因為照顧家人的緣故.
他將自己放下了.
停頓了好像在一種這樣的狀態裡面.
不住重複做一件事.
他不知道要做多久.
唯有讓他知道一定要繼續做下去.
繼續重複地去盡他自己的責任.
去回應他家人的需要.
抗疫有時就是一種這樣的狀態.
未知道下一步要怎樣走.
但你只能夠把握每一步.
做回手上應該做的事.
又或者可以說是一個信心磨練的地方.

$^{481}$因為什麼都沒有.
你不是可以有你自己.
剛才說了很多時候.
平時累積得來的工作成果 資源.
在你身邊包圍著你.
可能你只是赤裸裸地要緊守.
上帝吩咐我們的誡命的操練.
你也要每天仰望上帝的公認.
你也記得以色列人在抗疫的時候.
要靠嗎哪 要靠鵪鶉幹為的食物.
你不能留到第二天.
所以你每天都要相信上帝.
適當時候會公認我們所需.
在信心磨練的時候.
不是一個停頓的時候.
也不是一個被動的時候.
其實正正反而是一個預備期的時候.
準備我們的生命.
倘若上帝呼喚我們走下一步.
我們準備好自己.
我們不會自怨自艾.
我們不會挑低了.
或者是躲遠了自己.
或者是放棄了自己去回應上帝.
而是隨時在信心之中.
在尋問上帝之中.
在圍著上帝的恩典之中.
這樣去過每一天.
所以到了迦勒經過了這45年.
抗疫漂流的時候.
和這班活躍的子民在一起.
每天生活.
對於他來說不知道難不難受.
他就是不斷地操練自己的信心.
操練自己對著上帝的全心去依靠他.
到了上帝的時候.
他仍然是準備好自己去回應上帝.
進入了第三段的經文人生的階段.
他的人生下半場.
竟然是85歲才發生.

$^{521}$我們不知道大家.
你的人生下半場進入了沒有呢?.
你預計正在進入自己人生下半場.
你有什麼目標想去進發呢?.
大家看看迦勒的回應.
在約書下記14章11和12節.
我們一起讀.
螢幕上的經文.
預備1,2,3.
「現今我還懇見狀.
像摩西差派我去的那天一樣.
無論是戰爭是出入.
我現在的力量和那時的力量一樣.
請你將耶和華那天所說的這山地給我」.
就是剛才民數記14章那句話說的那件事.
他一直存在心裡.
就好像剛才說有些人將夢想.
小時候夢想還存在.
不過我們好像越來越遠了.
都沒有辦法去成就夢想.
但是約書有一個不只是一個夢想這麼簡單.
是上帝給他的說話要成就的一件事.
他一直記住.
一直問上帝是何時呢.
我準備好自己.
你何時要成就這件事呢.
這個就是約書那種等候.
卻不是一種被動的等候.
而是一直維持住自己的信心.
直到覺得上帝時間到了.
他就能夠毅然去回應上帝.
當然如果40歲的時候.
你想像一下.
當然這是一個FF的想像.
不是真實發生.
如果當時加勒頓的書.
舉起這支旗.
全部以色列人都悔改的時候.
他們不需要走入抗日漂流的時候.
這樣都不錯.

$^{561}$但是當時上帝懲罰了這班以色列人.
但他就跟加勒講.
你會進入這個美好之地.
加勒漢都會問 何時呢.
然後你又懲罰那班以色列人.
要40年漂流抗日.
那我呢.
你是否立即用直升機.
我立即調我進入這個艱難地.
不用經過.
陪伴那班以色列人一起搏弈.
一起掙扎呢.
你都會想.
都會問 上帝你的計劃是如何發生呢.
跟他們漂流抗日40年.
不知道漂流到何時.
那上帝會忘記了.
你跟我說過的話呢.
有時我們人性上都會有難免的掙扎.
過了45年.
是真的不同.
你看看加勒在14章12節的說話.
我讀給大家聽.
那一日你也曾經這樣說.
這裡有阿納族人.
以及寬大堅固的城.
或許 而又說會照他所說的.
與我同在.
我就把他們趕出去.
用了或許這個字.
我想大家留心.
可能40歲的加勒就覺得.
我一定可以的.
但是到了85歲的加勒就說.
或者上帝都繼續和我一起.
那我就可以了.
穩妥一點.
還有可能他考慮的東西都多一點.
這個很難免.
人生的經歷會是這樣.

$^{601}$但是關鍵是上帝與他同在.
或者如果上帝都繼續.
上帝忘記了那個說話.
45年來都是一樣的時候.
一樣上帝這樣與他同在的時候.
他就可以了.
雖然加勒之前都說.
我現在力量都一樣.
但是真正的最大的安穩.
最大的那種保障.
不是看人的力量.
而是上帝的同在.
而人的回應.
就是是否都是全心的去跟從耶和華.
耶和華的呼召仍然有效.
問題是我們那個心智對著他.
有沒有因為年歲的增長而減退了呢?.
還是一樣堅定地去全心去跟從耶和華呢?.
弟兄姊妹.
我不知道上帝給我們每一個人擁有些什麼呢?.
不能夠給大家一張紙填寫.
如果有表格.
我問你有什麼能力.
有什麼才幹.
或者過了這麼久的人生.
你覺得自己有什麼判斷的力量.
有什麼經驗.
有什麼資源.
是你擁有呢?.
可能你都會獵到很多很多.
你可能覺得我有自己的履歷去寫出來.
我做過什麼.
不同的人生.
也拿過不同的獎項.
也拿過不同的人給我們的一些確認.
一些成績等等.
可能都會獵出很多很多.
但接著我們要問自己的就是.
我們有願意為上帝投放自己擁有的東西.
有多少願意放在上帝的手上呢?.

$^{641}$但問題不是問.
要將下屎全部都交給上帝.
而是我們問清楚自己.
真實問自己.
既然我擁有這麼多的時候.
我願意給多少給上帝呢?.
就由那個部分開始.
一半都不要緊.
我就願意將一半給上帝.
我一半留給自己先.
都不要緊.
由我們願意放上的份量開始.
繼續將更多更多投放給上帝.
因為真正的人生的祝福.
就是在乎將整個人生都投放給他.
都歸向給他.
是一種全心的跟從.
不為我們心裡的憂慮.
或者環境帶來的懼怕.
而令到我們有所動搖.
我相信我們人最容易受環境影響.
亦都受心裡的一些意念影響.
而那些意念卻不是來自於上帝.
我們學習全心去更多更多的.
去服從我們的上帝.
跟從我們的上帝.
而且我們相信的是.
上帝的應許是會實現的.
好像加拿大一樣.
45年來.
最終看下上帝要在這個時候.
將這塊地方交給加納的首相所擁有.
上帝都必定在我們生命當中.
要成就他的作為.
所以去到約書記最後14節.
14章的14節.
15節.
他特別是15節說.
加納希伯來這個地方.
本來叫基列阿巴.

$^{681}$阿巴就是加納特朗普最厲害的一個.
基列是一個城的意思.
本來希伯來不是按著希伯來話叫的.
是叫阿巴城.
即是阿納族人的大佬的名字來命名的城市.
你想像其實世界上都有很多這些城市.
你想想不同國家都有用人的名字.
當那個地方那個國家最厲害的人.
就以那個地方去立名作為城市.
改名了.
現在歸入正式是以色列人十二支派擁有的地方.
改名為希伯來.
直到約書記的時候.
一直國中太平.
沒有戰爭.
一直都是在加納族的民族裡.
所擁有這個地方.
等之母我們有沒有準備好自己呢.
剛才我說了.
在加納裡邊.
由他自己作探子.
到進入去曠野的時候.
到他真的回應上帝的時候.
可能我們人生當中.
可能都好像經去那樣.
不會說整個生命都是很順利.
中間還有些曠野的時候.
要我們等候.
等待上帝到的時候去回應.
我們有沒有準備好自己呢.
這樣去回應上帝呢.
盼望我們都能夠好好的去養生命.
無論在什麼階段當中.
都是這樣去全心去跟從耶和華.
當我自己去準備訊息.
或者是都去.
好像加納一樣.
都做一些回顧的時候.
我都會和大家分享一些.
我自己人生的經歷.

$^{721}$其實都好像加納一樣.
有些不同的階段.
我自己是35歲的時候.
其實那時候還是錦宣的傳道人.
離開了香港.
就帶著一家四口去北美讀書.
就是讀我的神學博士.
當時一去就去了六年的時間.
那六年都好像大概捍衛.
在加納人生當中.
也不是真的抗夜.
但當中是你停頓了你的事奉.
你就在這裡學習.
你在這裡裝備自己.
和去體驗一些不同的.
上帝給自己.
要怎樣領受將來.
事奉的路怎樣走下去.
當時都有趣.
當我自己要去走的時候.
那個年頭其實.
很多人還是很願意.
譬如去到內地工作.
或者讀書.
很多時候都是對於.
在香港仍然很有信心.
是很值得投放在當中.
去發展的時候.
所以當時我自己.
帶著一家去外國讀書的時候.
我記得有朋友跟我說.
你這樣去走的時候.
你自己一個人受就好了.
不要帶著兩個兒子.
一起去受.
當時兩個兒子跟著我去的時候.
一個是六歲 一個是三歲.
另外聲音就來自.
再緊密一些的人.
譬如我家裡的人.

$^{761}$譬如我岳父.
對我們一家人走的時候都有聲音.
他覺得怎會放下工作.
而去讀書沒有任何收入呢.
當然他的聲音通常都是我太太聽的.
他幫我擋掉.
那些不同的.
對我們微言的聲音.
都做研究的時候.
想著去跟一些.
學成歸來的時候的一些學者分享.
我想做這個研究.
那時候也記得.
那個聚會裡面.
當我說想研究哪個神學家的時候.
都給他一盤冷水撥過來.
你還研究這個人.
這個人其實很多人都研究了.
沒有什麼可以研究了.
其實是很多的聲音.
或者很多的環境.
或者人的挑戰在中間.
很感恩.
那六年時間其實都是一個信心的操練.
到了我讀完書回來的時候.
結果投入了現在.
事奉的機構到現在11年了.
一直在當中.
一直的事奉 一直的服侍香港的年青人.
有時候我跟太太都會說.
等回頭.
最初好像有朋友這樣說.
對於孩子的經歷未必是好的.
但現在看著他長大之後.
其實我們有時候都說笑.
其實對於他們的祝福是最大的.
對於我反而是最小的.
因為我要改變的東西是最少的.
也沒想過我太太會進入神學院教書.
可能他們三個的經歷.

$^{801}$得到的恩惠或者上帝帶給的祝福.
反而我們想都沒想過.
所以頂尖的盼望.
無論我們在一個什麼階段都好.
有時候都會是一個信心的挑戰.
或者有很多的聲音.
或者很多的惡信.
都不是對於我們來說.
是一些很有利的消息.
但盼望大家都去.
專心仰望我們的神.
無論我們落在什麼階段.
都去盡力去.
跟從上帝的心意而行.
當我們好像迦勒那樣回顧人生.
可能我們都是一個看見人生當中.
只有祝福 只有上帝的應許.
要成就 要去發生.
好 我們一起低頭禱告.
在一個張臨期的日子.
藉著迦勒的人生.
都給我們一個勉勵.
給我們一個提醒.
無論我們怎樣去規劃自己.
或者看自己.
在一個什麼人生的階段都好.
無論我們在壯年.
要進入一個.
無論是退休或者有時.
一些事業的高峰.
或者過去的時候都好.
或者甚至是一些曠野.
迷失 迷惘的日子都好.
相信你一直都和我們同在.
就好像迦勒所相信的一樣.
而我相信.
你要在我們生命當中應許的事情.
一定會最終按你的時間.
按你的旨意成就.
就盼望上帝你幫助眾弟兄姊妹.

$^{841}$我們的心都是學習.
全心的去跟從你.
不被環境的.
或者其他別人的聲音.
去嚇怕.
或者去減退我們對你的信心.
變得心懷異議.
願你幫助 保守.
願我們學習投放更多.
投入你的國當中.
在你的旨意裡面.
聽我們同心的禱告.
奉救主耶穌.
得勝於自然.
\newpage



\section{路加福音 2:25-35}
\label{sec:oC6cD_CKYi8}
\textbf{2023-12-17 加恩樓《必看見所立的基督》路加福音 2\_25-35 講員:張漢強牧師}
\newline
\newline
連結: \href{https://youtube.com/watch?v=oC6cD-CKYi8}{\texttt{https://youtube.com/watch?v=oC6cD-CKYi8}} ~~~~ 語音日期: 2023-12-19
\newline
\newline
\hyperref[sec:jzt3Ohpru70]{\small{< < < PREV SERMON < < <}}
~
\hyperref[sec:index_chronic]{\small{[返順時目]}}
~
\hyperref[sec:index_scriptual]{\small{[返順卷目]}}
~
\hyperref[sec:EJ51POSz_gM]{\small{> > > NEXT SERMON > > >}}
\newline
\newline
路加福音 2:25-35
\newline
\begin{longtable}{cl}
\hline
\hline
章節 & 經文 (和合本修訂版)\\
\hline
2:25 & \begin{tabularx}{0.7\textwidth}{X} 那時，在耶路撒冷有一個人，名叫西面；這人又公義又虔誠，素常盼望以色列的安慰者來到，又有聖靈在他身上。 \end{tabularx} \\ \\ \relax
2:26 & \begin{tabularx}{0.7\textwidth}{X} 他得了聖靈的啟示，知道自己未死以前必看見主所立的基督。 \end{tabularx} \\ \\ \relax
2:27 & \begin{tabularx}{0.7\textwidth}{X} 他受了聖靈的感動，進入聖殿，正遇見耶穌的父母抱著孩子進來，要照律法的規矩而行。 \end{tabularx} \\ \\ \relax
2:28 & \begin{tabularx}{0.7\textwidth}{X} 西面就把他抱過來，稱頌神說： \end{tabularx} \\ \\ \relax
2:29 & \begin{tabularx}{0.7\textwidth}{X} 「主啊，如今可以照你的話，容你的僕人安然去世； \end{tabularx} \\ \\ \relax
2:30 & \begin{tabularx}{0.7\textwidth}{X} 因為我的眼睛已經看見你的救恩， \end{tabularx} \\ \\ \relax
2:31 & \begin{tabularx}{0.7\textwidth}{X} 就是你在萬民面前所預備的： \end{tabularx} \\ \\ \relax
2:32 & \begin{tabularx}{0.7\textwidth}{X} 是啟示外邦人的光，是你民以色列的榮耀。」 \end{tabularx} \\ \\ \relax
2:33 & \begin{tabularx}{0.7\textwidth}{X} 孩子的父母因論耶穌的這些話就驚訝。 \end{tabularx} \\ \\ \relax
2:34 & \begin{tabularx}{0.7\textwidth}{X} 西面給他們祝福，又對孩子的母親馬利亞說：「這孩子被立，是要叫以色列中許多人跌倒，許多人興起；又要成為毀謗的對象， \end{tabularx} \\ \\ \relax
2:35 & \begin{tabularx}{0.7\textwidth}{X} 叫許多人心裡的意念顯露出來；你自己的心也要被劍刺透。」 \end{tabularx} \\ \\
[1ex]
\hline
\hline
\end{longtable}
$^{1}$各位早安.
給點早安的感覺.
我姓張的 張翰強.
是中華基督教會旺角堂的牧師.
中華基督教會旺角堂可能大家都經過過.
去到旺角有一棟基督教大樓.
通常你們去基督教大樓都會買本書.
去到基督教大樓.
我們教會就在樓下.
如果大家有機會來到旺角.
歡迎大家.
我們教會的咖啡機沒那麼漂亮.
不過我們都有咖啡喝.
歡迎大家一齊喝咖啡.
我感覺到今天好像是一個合一主一.
原因就是上個星期.
使徒永富先生來到我教會講道.
今天輪到我來.
也是屬於他的教會來講道.
我覺得是一個很好的感受.
亦代表我教會.
代表我太太在後面.
你不知道是誰就不要緊.
我最重要是知道.
向各位問好.
能夠在不同地方一起敬拜上帝.
是一件很美好的事.
今天的主題就是.
必看見所立基督.
到聖誕節.
我們重新紀念耶穌基督的降生.
通常我們在幼稚園看的聖景.
有什麼場合.
聖誕老人?.
馬槽?.
馬槽是錯的.
那時候是沒有馬的.
普通居民不准養馬.
因為是打仗用的.
在動物的地方.

$^{41}$有機會我再講.
但我不知道有沒有機會.
然後還有什麼?.
三博士.
但我們不知道是不是三個博士.
聖經也沒有說三個.
還有.
剛才那兩個我們不確定.
但有一個很肯定經常出現.
就是有個老人家出來抱個嬰兒.
那個是誰?.
剛才讀聖經.
沒有留意.
西勉.
大家打開聖經.
我很想和大家看聖經.
有手機有聖經都可以看.
主要是講兩三節.
西勉的事蹟是很奇特的.
奇特到一個地步.
他本身講了耶穌整個生平.
那個寓言.
那個寓言是一個很精彩的部分.
其實我們研究聖經的人知道.
路加紀刻意地去寫這個寓言.
是直接指向路加福音的結尾.
耶穌釘十字架那裡.
是很能夠息息相關的.
如果用我們現在寫文的方式.
我們就叫做起承轉合.
一開始就講了提綱契令.
然後結局一直下去.
我們就看到那個過程是很精彩的.
今天我們要講西勉.
當然整個場景.
就是在耶路撒冷.
我們不知道為什麼西勉會一直在聖殿.
又不知道他的生平狀況.
甚至連多少歲我們都不知道.
一般我們講西勉.

$^{81}$如果回教會聽過主一學故事.
你會覺得西勉一定是一個老人家.
似乎是.
因為他說見了耶穌之後就怎樣.
就死了.
但凡離開世界的人.
是不是一定老.
不知道.
但在聖經沒有說.
傳統有說.
傳統教會的教父.
一直告訴我們.
他是高齡離開的.
足足113歲.
長壽.
香港也是最長壽.
大家知不知道.
但長壽.
只有長壽.
就差一點.
長壽得來有盼望是很重要的.
西勉這麼長壽是不是很好呢.
可能不是.
為什麼呢.
我講實際的狀況給大家知道.
耶路撒冷這個地方.
今天我想大家都很關注.
特別這兩個月.
足足打了一個多兩個月.
是一個很特別的地方.
如果今天有機會去耶路撒冷.
整個耶路撒冷城其實分了四楷.
有不同的宗教.
其實佔領了在不同的地方.
背後有不同的政治勢力.
這個地方.
自古以來到今天.
都是一個.
你叫做兵家必爭之地也好.
你叫他做本身是一個很奇特.

$^{121}$不知道為什麼大家都.
要爭取這個地方也好.
總之這個地方是很有趣的.
包括整個猶太人的傳統.
從亞伯拉罕開始.
就經而慢慢形成了.
對於後來的耶路撒冷.
一個重視.
而這個重視一直到猶太人.
認為是他們的所謂聖神.
這個場景.
在耶路撒冷.
耶路撒冷是一個什麼地方呢.
在耶穌出生之前.
大概是主前63年.
有一個將軍.
這個將軍可能你聽過他的名字.
叫做龐貝.
龐貝古城聽過.
有個誠意去命名.
你知道這個人是多麼了得.
他的功勳在哪裡呢.
就是因為他是羅馬人.
裡面的曾經.
逐六個.
做皇帝的其中一個.
當時有所謂元老院.
然後後來就.
跟著有些增敬.
然後有些人開始走出來.
然後去跟他們對抗.
然後慢慢去爭取到皇位.
最後得到的皇位大家知道就是屋大維.
等一下我會講一下這個歷史.
在主前的63年.
龐貝將軍帶了大批.
羅馬軍.
當時猶太人.
就在他們自己.
一個小小的國土裡面.

$^{161}$就叫做.
有些自治權的.
但是他們是不肯聽令任何的政權.
而羅馬經已如日方中.
一個小小的政權.
在巴勒斯坦地.
就很自然不能容得下.
於是龐貝就帶著.
大批的人馬.
就去剷平了整個耶路撒冷.
在那次的.
殺戮裡面.
那個殺戮的數字是相當驚人的.
如果用今天的角度講.
都是很驚人的.
就是有萬多個猶太人被殺.
類似是屠城.
你想一下.
這個地方能夠擺到多少人呢.
其實當時都不多.
都是幾萬到十萬之內.
在這樣的狀況沒有了萬多人.
你想一下如果你是猶太地的.
猶太人.
你目睹這些事件發生.
你是倖存者.
你知道你的城邑裡面.
有萬多個居民.
就這樣犧牲了.
沒有了姓名.
甚至珠寶財物都拿走了.
如果你是倖存者你的感覺是怎樣.
很深刻.
應該是吧.
並且可能會有仇恨.
因為那萬多人裡面.
總會有些是你的親人.
總是這樣.
你看家生那裡打仗.
看到小朋友在叫.

$^{201}$我爸不見了,我媽不見了,我哥不見了.
我弟弟現在被壓在石頭堆.
那個人繼續生存會怎樣.
你記住這個印象.
這個印象就是我們待會要講西面的狀況.
當這件事發生之後.
其實我們.
起碼知道兩件事.
第一就是他們帶走.
殺了很多人.
除了帶走很多金銀財物.
也帶走了很多當時的猶太人.
然後.
在當地的耶路撒冷.
這個地方其實居民急速下降.
因為變成了一個廢城.
然後就在那一代的年青人.
或者少年人.
他們倖存的.
就一直生存下來.
就是.
去到耶穌的年代.
因為是主前63年.
去到耶穌的年代這些少年人青年人.
大概是70歲.
80歲.
因為大家知道耶穌的出生.
就是在什麼時候.
00年.
通常我們有講法.
但是一般的賢經.
我們知道應該是主前的4至6年.
因為有幾個因素.
令我們很確定地知道.
那群年青人.
他們留住了.
這種深刻的印象.
幾十年來.
我想是很難撕掉的.
以前不流行輔導的.

$^{241}$甚至有輔導的,你不知道去哪裡看.
所以在那群倖存者裡面.
其中一個是西面.
他們是帶著一種.
一種無奈或者甚至憤恨.
甚至可能很希望.
人間是有和平的.
希望是有公道的.
於是我們帶著這種感覺.
待會我們去看西面.
龐貝殺戮之後.
耶路撒冷又怎樣呢.
他們留住了很長時間.
直至主前的37年.
這是一個重要的日子.
原因就是猶太人不喜歡.
被外國人統治.
羅馬人也知道很難統治他們.
於是找了一個中介人.
這個中介人其實大家讀聖經就知道.
叫做大希律.
大希律就是耶穌出世不久.
下令要殺那些BB.
兩歲以下的.
就是這個大希律.
結果這個大希律上了場.
羅馬政權很大力支持他.
也因為有一個原因.
就是這個大希律.
本身背後是有人脈的.
他的人脈就是.
屋大維,屋大維就是第一代.
羅馬的皇帝.
有些.
也不是叫疏堂.
總之有些很密切的關係.
有機會再說.
於是這個大希律一上場.
他第一件事要做.
就不是要君臨天下的姿態.

$^{281}$因為他知道這些猶太人.
是不容易去管治的.
於是他就做了一件很好的事.
對猶太人就是重建聖殿.
重建整個耶路撒冷城.
但這個人.
他是雙贏的.
原因是他不斷地重建.
但其實他背後.
如果用現代的角度看.
其實他是一個建築師.
是很出名的.
甚至他現在的建築.
是在世界十大奇景之一.
他是一個很擅長.
很喜歡建築.
所以他要求的建築.
是精準度很高.
所以他去重建聖殿.
你知不知道重建聖殿.
重建了多少年.
去到耶穌的年代.
都還沒有建好.
直到第二次聖殿被毀的時候.
其實都還沒有建好整個聖殿.
所以他不惜工本.
討好猶太人.
也討好自己.
於是大肆重建了整個耶路撒冷城.
所以當時的耶路撒冷城.
雖然經歷了在63年.
那個年代被剷平.
甚至變成.
好像廢墟一樣.
不過在大希律的手下.
慢慢這個城音重新恢復.
回到當日或再之前的光輝.
並且那些.
陸陸續續在外地.
已經逃亡的居民.

$^{321}$又回到耶路撒冷.
所以在耶穌出生前.
我們估計.
當時的耶路撒冷的居民.
是可以去到幾萬至十幾萬.
都可能有的.
但實際我們不知道.
因為我們實際沒有統計.
於是整個耶路撒冷.
已經變得很繁盛.
而當日倖存的那一代人.
留下來的西面.
他的心是怎樣呢.
是很複雜的.
看著那個環境.
曾經變成廢墟.
但一下子就變成了.
一個很繁榮的地方.
他只有一個心願.
就是希望.
看到上帝的公義和愛.
再次臨到這裡.
我相信除了西面.
實際上在那一代倖存者.
每一個人心上都有類似的狀況.
而西面在.
所謂幾萬至十幾萬.
的人裡面.
他實際上是一個很普通.
普通不過的人.
但如果記載的話.
我們是不會知道有一個人叫西面的.
事實上.
當時見過.
見過BB耶穌的人.
是很顯然.
不止西面.
如果根據聖經還有那些人.
有牧羊人.
還有就是.

$^{361}$接下來要說的那個女仙.
哈娜.
是不是只有這麼少呢.
耶穌在哪裡出生.
先說.
巴黎行.
今天好像主要是學問答比賽.
沒錯是巴黎行.
但現在說耶路撒冷.
巴黎行在哪裡.
大家地理.
南面.
他要由南面.
巴黎行一路北上.
因為要做很多的.
包括潔淨禮等等.
一路行水.
有沒有人見到他們.
我想他不會把他困在屋簷裡.
慢慢抽.
不會的.
很自然應該在這段長的旅程裡.
我相信一定很多人.
見過他.
出生的時候.
有沒有接生婆.
應該都有.
除了接生那些.
去他地方住.
應該都有些人.
知道BB出生.
其實你數下去.
一路由出生去到.
由巴黎行去到耶路撒冷.
其實基本上是一個常識.
很多人見過耶穌.
是不是.
即使他不知道他叫耶穌.
起碼知道BB在裡面.
為什麼.

$^{401}$路加只記載.
牧羊人西面和呂先知阿那.
如果我們說牧羊人.
要記載.
因為他遇見天使.
這些場面不是每個人都見到.
記載.
另一個就是阿那.
阿那是先知.
先知是偉大.
大家知道在耶路撒冷.
一定要記載,並且是呂先知.
路加一定記載.
因為如果你看路加福音.
你會發現路加是很重視女性.
和紀錦福音書不同.
這是很特別.
西面呢.
他沒有天君天使場面.
又不是先知.
起碼聖經沒有說他是先知.
從來沒有說過.
西面幾十萬人.
只是一個小人物.
但為什麼他會被記載.
其實.
很耐人尋味.
是不是因為他見過耶穌.
所以他就有理由.
但剛才我們也說.
由伯利恆出生.
一直由南面.
去到北面,去到耶路撒冷.
途中都不知道有多少人見過耶穌.
為什麼不記載那些人.
為什麼要記載.
這個西面呢.
這就是我們之前所說.
西面是作為一個老人家.
教會傳統認為.

$^{441}$他113歲.
作為一個老人家.
他曾經經歷過旁輩.
所帶來的屠殺.
對於和平.
猶太人.
要被拯救的願望.
在他身上一直揮之不去.
他心底裡.
一直都期盼著.
見到上帝的公義.
見到耶路撒冷.
重新榮耀上帝.
見到那群猶太人.
可以重新被拯救.
離開羅馬的政權.
這個一向都是他們那一代人.
一直下來的.
在他的人生裡,在死亡.
要能夠入土.
為安之前,他實在希望.
自己不要帶著.
遺憾離開這個世界.
我想我們每一個人都是.
如果我們有遺憾.
我們真的好難.
當我們離開世界的時候.
如果用中國人的說法.
死都不難.
於是你會發現.
他實在希望自己不要帶著遺憾離世.
他實在希望猶太人.
得到拯救,人間有一個.
永久的和平,不要再打仗了.
殺戮太恐怖了.
人命太過賤了.
人賤應該是屬於上帝每一條人命.
豈能再發生呢.
將來會怎樣,不知道.
經歷過戰爭的人是好擔心.

$^{481}$時常會再有戰爭的,這是事實.
簡單來說.
西面的心下一直都連繫著.
整個猶太的群體.
他的心是一直連繫著.
上帝的應許.
因為先知說下來.
一定是有一個彌賽亞要來.
他覺得一定要來.
而且一定在他離開世界之前.
故此當西面.
在聖殿一出來.
見到耶穌的嬰兒.
你看看第29節.
他一開口是說什麼.
這個嬰兒好像爸爸.
會不會啊.
如果我們去探訪.
就會了,是吧.
但也很奇怪.
一開口就是.
主啊,如今可以照你的話.
容你的僕人安然去世.
奇怪嗎,這個對話.
你想想他抱著嬰兒.
然後對著他爸媽說.
又不是像爸爸,不是像媽媽.
然後跟神說.
跟著跟神說.
我快要死了.
你抱著一個新生嬰兒在爸媽面前.
他對於嬰兒充滿了希望.
然後你告訴爸媽.
我快要死了,你是不是覺得很奇怪.
所以你讀聖經.
你重新回到處境.
你會覺得,嘩,整件事很反差.
按常識不應該這樣說.
但如果你剛才.
明白了歷史背景.

$^{521}$你明白了龐貝本身.
帶來那種令人痛恨.
看到被殺戮.
那種心痛.
看到生命如草蜢.
的時候.
你看到以色列國,包括你是猶太人.
你看到以色列國.
你會覺得國家再沒有希望.
的時候,你帶著這種心情.
你看到一個.
忠於你期盼的嬰兒.
是很自然,會像西面一樣說.
是不是,終於見到了.
不過你有沒有想過.
那個是嬰兒.
他不是已經行了很多神蹟.
是不是,我都見到耶穌.
我對他有.
用現代的商業術語.
可以投資.
信得過,是不是.
是嬰兒.
你看看你旁邊左右.
他會不會是將來Elon Musk的繼承者.
影響世界.
不知道.
我怎麼知道,可能會.
是不是.
可能將來.
我們將來有幾千個一起.
會眾參與崇拜,行不行.
行,但我們不知道.
不過有沒有這麼奇怪.
你見到嬰兒而已.
他這麼肯定.
會不會認錯了.
可能會的.
但當然我們講到最後,我們知道不會.
就因為聖靈大令到他.

$^{561}$不見了.
於是在這個過程,他講完這番說話.
令我們很明白.
西冕之所以被記載在.
路加福音.
除了因為他講了對耶穌的.
預言之外,也都因為.
就是.
他的人生.
是一直都心繫著上帝.
和上帝的.
應許.
他整個生命就在這裡.
所以當他在生命完結的時候.
這個時候他才覺得.
起碼他覺得自己是可以.
畫上句號.
西冕之後有沒有死呢.
其實你和我不知道.
這裡沒有講過,不過西冕覺得.
人生可以畫上句號,就是這個時刻.
起碼他覺得再沒有遺憾.
如果我要離開的話.
好了,我終於見到.
而他這一份終於見到.
是因為他對神那份信靠.
是因為他很深信.
有一天,彌賽亞要來.
告訴他,在他在生之年.
一定要來,所以他一定要見到.
而他見到了,在昔日.
即使多麼困難,令到他心裡.
有條刺藏著.
他沒有辦法拔掉.
但終於見到耶穌了.
所以他忍不住和神說.
他不是抱著寶寶和爸媽說.
他忍不住一抱在手上.
他就說,是啊,上帝.
這個就是了,就是你讓我.

$^{601}$見到的那個耶穌基督.
這個是寶寶.
他的信心是很大.
他很相信這個一定可以的.
是怎樣可以他不知道.
不過一會兒的預言他都說了.
今天我不說預言,只是說.
關於這個這麼普通的.
西面一個人.
一個普通人在上帝的眼中.
變得這麼重要,是因為.
他肯將他的人生.
或者他一生人的心思.
都是連繫著.
在上帝那裡.
都是將整個信心投放在.
上帝那裡.
所以弟兄姊妹我們想想我們自己.
我們在座有沒有人覺得自己不普通.
如果你說不是,張牧師我一點都.
不普通的,你舉手給我看.
我祝福你.
沒有人敢不舉手.
好,事實上.
我們真的很普通的,是不是.
有沒有這個.
人大在我們中間.
政協有沒有.
特首不要在這裡.
千萬不要在這裡.
還有沒有.
有沒有專業人士.
可能有的.
不好意思舉手.
其實我們.
就算去到什麼位置都好.
我們都是很普通的人.
普通到一個可能就是.
我們連一些感人的.
生命故事都沒有.

$^{641}$我們日出而作,日入而息.
每天.
每天都是這樣過.
很普通,普通到.
可能我們離開這個世界.
人家在我們的墓志銘想寫一些東西.
都不知道寫什麼.
就說不如寫.
他是一個很勤奮上班的人.
又或者是.
很懂得享受生活的人.
可能就是這麼多.
事實上我們人生,或者很多人的人生.
都是很普通.
不過在西面.
的經驗給我們看.
我們是.
可以不普通.
可以是不一樣的.
可以是因為.
我們和神的那份連繫.
將我們的心思和生命.
投入了和.
投放在上帝的懷抱裡面.
而我們才可以.
生命變得不一樣.
不一樣可能他依然是很普通.
不過在神的角度.
裡面和神的眼中.
是不一樣的.
我們每一個人其實都是很尊貴的.
包括這個西面.
而當我們有這份信心.
如西面一樣.
只見到一個嬰兒,很確定這個就是耶穌.
當我們有信心.
能夠見到我們身邊的人.
和事.
我們有這份信心,我們都很確定地知道.
我們的神.

$^{681}$是在這裡的,我們的神.
是在祝福我們,當我們坐在這裡.
我們都可以很知道.
如果我們有信心.
上帝正在祝福我們每一個人.
包括我們的教會.
這個是很重要的.
一個普通的人,即使很普通.
但他可以不一樣.
因為他有這份信心.
投入在上帝那裡,投入在耶穌那裡.
以至到耶穌令我們的生命.
變得不一樣.
而西面也一樣.
他能夠被記載在路加福音.
令我們知道他的不一樣.
是基於他心繫著.
他所相信的神.
我再說一個普通人的西面.
得到神肯定的祝福.
是因為他的心一直.
心繫著神.
有這份信心,心繫著上帝的.
說話和應許,是一定會發生的.
我們今天讀聖經一樣.
我們知道.
聖經裡面說過很多.
的神跡其實,對不對?.
你信不信的各位?.
信不信的?.
是真的相信?.
你信不信紅海會打開?.
信不信.
死的人可以復活?.
可能我們都相信.
不過當我們面對.
一些突如其來的衝擊.
或者一些艱難的時候.
我們再問自己.
信不信上帝呢?.

$^{721}$那個時候才變得更真實.
是不是?.
西面在這裡,他經歷了.
幾十年的心裡面的.
沉痛.
他一直都堅持這份相信.
是很真實的.
他是信上帝的.
其實我們一樣都可以.
所以我們在西面的人生裡面.
我們看到的這個.
是我們人生都值得學習的功課.
讓我們的心投放在上帝那裡.
當然西面是有不一樣的地方.
他不單止讓心.
投放在上帝那裡.
並且你看經文.
西面為什麼能夠心是上帝?.
因為他認為.
上帝是最好的.
上帝認為好的.
他都認為是好的.
在經文哪裡看到?.
你看看二十五節就知道了.
二十五節.
和合本是這麼說的.
就是這個人是什麼人?.
西面.
又公義又虔誠.
是不是?.
這兩樣東西是同情的兩面.
虔誠是什麼意思呢?.
虔誠如果你看聖經.
你打中文.
你查本聖經裡面有哪裡譯過虔誠呢?.
很多地方都有虔誠這個字.
很多.
不過這裡用的虔誠這個字是有一點點不同的.
整本聖經這裡用的.
虔誠的原文.

$^{761}$是用過四次.
雖然他都是譯做虔誠.
和其他不同.
這四次在哪裡呢?.
三次在士道行傳.
當然路加福音和士道行傳.
其實都是路加寫的.
所以我們讀路加福音.
必須要看士道行傳.
路加才會用這個字.
這個字有什麼意思呢?.
如果我們查原文.
其實他本身有一個.
所謂虔誠.
是拿一些好的東西.
意味著.
只要上帝在哪裡.
我們才可以在那裡.
拿到一些好的東西.
如果用一些.
我們廣東.
能夠明白一點的語言.
就是我們認為.
上帝好的東西.
於是我們就去上帝拿.
上帝認為好的東西.
上帝不認為好的我們拿不拿?.
不拿了,就是說上帝認為好的.
我們就認為好的.
這個就是整個所謂虔誠的價值觀.
的一個系統.
也都意味著.
大家有沒有去過街市買菜?.
應該有吧?.
不是街市,應該在百家圍巷買過.
你看到一些東西.
譬如有個橙一堆在那裡.
你會選擇一下.
你怎麼都不讓他選擇.
是我自己親手選擇.

$^{801}$選什麼?怎麼選擇?.
大隻一點,重一點.
重一點就是有水分.
有些發霉,要不要?.
不要吧,那些癌癌蠶蠶.
都不像一個橙,整個四方形.
那些可能都變了異.
都不要了,我們要怎麼樣?.
去買東西當然選擇一些.
選擇一些好的東西.
這是人之常情.
所以當說乾淨這個字.
其實對我們來說不陌生.
我們當然選擇一些好的東西.
上帝說好的.
那就好.
正常.
不過提醒姐妹.
我們是不是真的.
去街市我們是這樣.
和上帝的關係是不是這樣.
我們會不會以為.
上帝認為好的我們一定說好的.
當然我們不敢說不要.
不過我們心底裡.
可能真的覺得不好.
會不會?每個星期回來崇拜.
好不好?.
好,真的每個星期回來嗎?.
差不多.
我們再數下去.
我不想再數.
可能很困難的數下去.
所以我們想想.
上帝說好的我們說好的.
是不是這樣?.
我們每天要靈修,上帝說好的.
是不是?.
我們行動上是不是真的覺得好?.
上帝叫我們禱告.

$^{841}$很好,和上帝經常聊天.
我們行動上.
是不是真的覺得很好?.
我們有沒有做?.
我們的神說.
事奉好.
我當然知道好.
但我們行動上.
是不是真的覺得好?我們有沒有做?.
頂智慕在我們生活裡.
行為,職業.
和家人相處.
種種的生活環節.
是不是都是以神認為好的.
然後我們去做?.
這個就是.
經文之所以.
這麼肯定這一位.
的西面.
其中一個就是他是虔誠的.
並且這個虔誠的意義.
就是他以上帝的好.
為自己的好.
理論上一定可以的.
正如剛才說.
奉獻好,事奉好.
愛神好,祈禱.
全部都好,不過就是不知道.
有沒有做到.
於是在經文說.
又虔誠又什麼?.
又公義.
公義原文.
大家想想,直譯應該做什麼?.
正直.
不會兩面三刀.
不會說一套做一套.
正直,所以我說.
公義,虔誠本身兩個是同情的兩面.
正直是什麼意思?.

$^{881}$就是.
如果當代日本.
是譯做正直的人.
正直的西面.
就是扣著虔誠說.
一個正直人.
就是他怎樣相信.
然後就怎樣?.
就怎樣做.
我們叫正直.
他不會說一套.
就做一套.
而這個所謂正直的西面.
就是.
他說一套.
做也是那一套.
沒有變過.
很直覺的.
就是這樣做.
他不會說的和做的不同.
他的信念就是他的行動.
但很可惜我們見到很多人.
很多人即使有多美好的信念.
他們讀過書.
明白過聖經.
知道什麼叫好.
不過正正就是做不到正直的行動.
西面是一個有虔誠的人.
是一個以上帝的好.
作為自己的好.
他是這樣想.
有這樣的信念.
他會去做.
所以經文去肯定他.
其實是值得肯定的.
經文描述他是一個正直的人.
他深知上帝的好.
他就自己按著上帝的好.
來做上帝的好.
等知會有時候我們也是這樣.

$^{921}$我們生活上見證好不好?.
傳福音好不好?.
帶人回教會好不好?.
愛人如己好不好?.
其實我們全部都知道.
通通都好.
不是我們不敢說不好.
是我們真的覺得美好的.
但當我們認為這些都是神.
都是聖經所說的好.
我們要走出來.
然後我們才是正正直直.
可以原原本本告訴上帝.
我們是這樣相信.
我們同時是這樣做.
我們一方面相信.
另一方面卻是不理會.
不實踐的說話.
我們就是那些說.
不實踐的說.
但是那些說.
說得多麼好聽都好.
就變成那些表裡不一的人.
如果難聽一點.
我們叫這些人虛偽的人.
如果用聖經說.
這些就是假冒偽善的人.
就是耶穌經常罵他們那班人.
西面.
聖經說.
他是一個又公義又虔誠的人.
一個以神好.
作為自己好的人.
對自己從神而來的信念.
敢於轉化為行動.
用行動用生命.
來見證上帝的正直人.
這是西面.
經文對西面的描述.
是值得我們深思的.

$^{961}$同樣是值得我們深思的.
是接著經文所說.
西面除了是正直虔誠之外.
西面也是一個.
你看看25-27節.
是一個被聖靈感動的人.
他願意開放自己的心.
被聖靈感動.
和帶領他.
你看看25-27節.
三節都有說聖靈.
很有趣.
25節的結尾.
有聖靈在西面的身上.
26節開始就說.
西面他得了聖靈的啟示.
27節.
那裡說.
西面他受了聖靈的感動.
短短三節經文.
用不用這麼煩.
不停地說聖靈.
路加福音.
路加醫生不厭其煩地告訴你.
在這個西面.
在這個西面一次兩次三次.
都是有聖靈在他身上.
然後是什麼.
他很頻繁地.
去經歷聖靈.
他願意打開他的內心.
讓聖靈來到他身上.
我們都知道.
聖靈臨到一個人的身上.
或者聖靈去啟示和感動一個人.
其實通通一定是上帝的主權.
我們不可以拿錢出來.
和他交換 上帝今天給我一個聖靈.
不可以的.
沒有這些事發生的 我們知道聖經.

$^{1001}$是聖靈的臨在.
一定是上帝的主權.
不可以改變的.
我們想用什麼.
來改變都不可以.
大家記得士徒行傳出現的那件事.
有個人想去趕鬼.
於是給錢看看能不能買到一些方法和能力.
沒有 這些事.
是不會發生的.
我們都知道聖靈是這樣的.
既然是上帝的主權.
為什麼.
聖靈會臨在和啟示西面.
開放.
他的心當然是重要.
不如想想.
為什麼聖靈在旁邊.
而不在自己身上.
什麼原因呢.
當然我們自己可能知道很多原因.
不聽話 又不理會上帝.
或者聽聽地 諸如此類.
但有一點.
我們是不可以忽視.
如果我們一直看西面的狀況.
我們發現.
有聖靈臨到.
我們知道.
其中一個 剛才說了.
就是西面和上帝.
是連繫著的.
他有這份信心.
他的心連繫著.
上帝.
這個當然是很重要.
也意味著 這是一份什麼關係.
你現在坐在這裡.
男士會不會想起你老婆.
女士會不會想起你老公.

$^{1041}$不會.
現在想起上帝好了.
但你一出了教堂.
會不會想起他們.
不會 想起待會去吃飯.
但如果你告訴我.
我一出了門.
我就想起我老婆.
那是什麼意思.
關係很密切.
你只要走進去.
你都會想起他.
是因為關係很密切.
所以你走進去.
你都會想起他.
這就是西面.
他就是心繫著上帝.
於是他有這份密切的關係.
往往我們和上帝有密切的關係.
其實就是促使.
聖靈臨到我們身上.
一個重要的元素.
反觀這段經文.
以及看其他經文.
其實都有這些暗示.
你看大約26字.
所說的啟示.
那個原文的字根.
其實是叫或指示.
叫或指示.
這些意思.
無情情 為什麼會聽到.
聖靈在叫我.
其實關鍵在於.
關鍵在於.
我們很熟悉的禱告.
禱告原意就是溝通.
現代的術語就是對話.
禱告原意就是對話.
現代的術語就是對話.

$^{1081}$你說 神都說.
不只是你一說就說了.
你不用說了 我說完.
就完了.
本身禱告就是一種溝通.
有一次有人問.
德蘭修女.
大家都很關心.
德蘭修女為什麼好像聖人一樣.
見到她的臉當然有很多問題.
特別是我們有信仰的人.
我們都很想像她這樣.
於是人們一見到德蘭修女.
就想像她和常大關一定很好.
於是人們就問.
修女 禱告的時候.
你和神會說什麼.
聽聽她說什麼 我們學一下.
修女怎麼回答.
她說.
我什麼都不說.
然後她說了一句.
我只是在聆聽.
我只是聽 我不是想說.
這是顛覆了.
我們日常禱告的習慣.
我們習慣怎樣禱告.
說.
但我們忘記了什麼.
聽.
禱告的原文本身就是溝通.
用英文翻譯.
通常我們叫做communication.
有聽 有說 有說.
就是有聽.
所以西面就聽到.
為什麼聽到 她是有禱告的人.
弟兄姊妹為什麼會聽到.
上帝叫你.
神會和你說.

$^{1121}$是因為你和神說話.
上帝回答你.
叫你 和指示你.
換句話說.
西面有聖靈臨在.
和啟示.
有神的靈和他說話.
是因為.
簡單來說 是因為.
西面就是一個禱告的人.
他禱告 神回應.
他和神說 神又和他說.
事實上禱告和聖靈的臨在.
從聖經所見.
很常很常.
都是相生的.
你想一想 舊約.
但以你禱告之後 神的使者就來了.
在新約耶穌在.
馬利沿禱告 禱告到沒多久.
天使就來了.
我們都見到禱告和.
神是息息相關的.
他使者會臨到的.
而禱告和聖靈的臨在.
和啟示都是一樣.
是息息相關的.
你看《士徒行傳》第十三章.
有機會回去看.
安提柯教會的領袖.
當時還沒叫保羅.
還叫素羅 他直直無名.
他和其他人 包括巴拿巴.
包括西面的另一個.
包括路球 包括馬念.
包括保羅自己.
一起祈禱 經文說.
他們正在崇拜.
和禁食祈禱的時候.
聖靈就向他們.

$^{1161}$顯明他的旨意.
在禱告的時候.
聖靈就和他們說話.
在禱告的時候 神就和他們說話.
在他們說話的時候.
神就回應他們的說話.
故此當我們說.
我們倚靠聖靈的人.
我們必須是一個什麼樣的人.
是一個禱告的人.
我們不要.
羨慕其他人有聖靈.
我們都有.
我們不要羨慕那些人才聽到聖靈的聲音.
我們都有.
因為我們要學習做一個禱告的人.
故此以不所需.
告訴我們 靠著聖靈.
隨時多方禱告.
唯有禱告.
我們才能經歷到聖靈.
在我教會裡面有一個.
特別關注.
癌症的病人.
我教會很有趣的 很喜歡關心病人.
有一個團契是被中風的.
我是參加那個團契.
我中過風.
所以我參加.
團契是有癌症的.
有癌症病人.
或者有家屬一起參加.
最近我們請了一個.
國內人來分享.
他叫做菲烈.
最初我聽到他病的情況.
是很慘烈的.
在腦裡有癌症.
是搞不定.
不知道怎麼辦.

$^{1201}$於是我經常想他是不是搞定了才來我們那裡.
講見證.
我經常在幻想.
他一定是身體.
各樣東西很虛弱.
我一見到他.
我真的有落差.
就是他的臉色.
很紅潤.
甚至我發現他比我更正常.
我懷疑.
我是病人多於他.
然後一邊聽他講見證.
很奇妙 詳細我不講.
我大概講他的狀況.
就是他醫生判斷了有腦癌.
而腦癌的位置.
牽連了很多不同的神經.
試過化療.
電療 因為醫生說不要開刀.
那個位置太危險.
一開刀就只剩一成的機會.
當然是做其他方法.
於是就做了 做了所有東西.
他已經很虛弱.
幫不了 腦依然在.
只是沒有縮小.
但都在壓著神經.
結果無藥可醫.
之後是想領養.
止痛等過了那些日子.
他不甘心.
他禱告.
他說上帝你醫我吧.
將來我如何去奉獻給上帝.
當然一輪跟上帝講了很多.
上帝聽他祈禱.
有一天他.
好像做夢一樣.
他經歷到上帝跟他說.

$^{1241}$我會醫你的.
甚至他在那個類似夢境的東西.
他的見證講法.
他說我見到上帝.
我有信心得到醫治.
信心在哪裡.
都是相信醫生.
然後跟醫生講.
我做手術.
於是就做了那次手術.
而且手術是很成功.
即使醫生說其實只有百分之十的機會.
你真的試嗎.
試 因為他信上帝.
他聽到上帝的說話.
他祈願到得到上帝告訴他.
於是他帶著信心來醫.
結果就是他來到我們面前的時候.
是我很詫異地覺得.
上帝真的很厲害.
不單止醫好病.
連身體的健康都比我更健康.
一定是我平日的禱告不夠.
所以我應該更健康.
丁姐妹你見不見.
就是那個禱告.
是跟神連繫了起來.
就是那個禱告.
開放自己的心.
給神跟他講.
要知道他聽得到 於是他又相信.
有時聽到你未必相信.
但他聽到他相信.
因為他自己去求上帝.
而上帝回答了他.
都不是上帝回答你 是上帝回答他.
於是他如常人一樣繼續生活.
現在他很好.
不同地方不停地見證.
我們的生命.

$^{1281}$其實丁姐妹都是一樣.
是呀 我們是要禱告.
唯有禱告我們才能經歷聖靈.
唯有禱告我們才能夠.
經歷神那份真實和信實.
不再在紙本聖經裡面.
那麼簡單.
而是一個活生生的神跟你一起.
丁姐妹我們禱告.
在今天我們自己的禱告.
我們的禱告是不是只是.
在吃飯之前我們才做.
我們的禱告是不是.
只是在臨危.
臨急的時候我們才禱告.
我們的禱告.
到底有多麼真誠.
還是行禮如儀.
我們的禱告.
你看看你的手機.
是不是只是.
在whatsapp裡面.
弄一隻祈禱手.
就是禱告員.
沒錯.
吃飯祈禱也好.
我們要繼續.
臨危臨急的禱告.
我們都需要 亦都是人之常情.
whatsapp給一隻祈禱手.
沒有壞的 你給十隻都可以.
不過.
可不可以我們多走一步.
跟你關心的人.
開聲.
為他禱告.
甚至你跟他一起.
去禱告.
是有血有肉的人.
我們禱告的那個.

$^{1321}$亦都是很真實而具體的神.
他不是抽象.
上帝聽禱告.
西面在第29節.
到32節去禱告.
禱告完之後.
聖靈就啟示了西面.
說了很精彩的那兩節經文.
就是第34和35節的預言.
是預言一個眼前的BB.
他未來的一生.
這裡我不解釋了 留給大家慢慢去看.
弟兄姊妹.
事實上西面和你和我都一樣.
都是一個普通的人.
但是上帝對於他.
他是加倍的厚重.
給他莫大的祝福.
從他的身上我們有幾點.
很值得我們重新去反思和學習.
他對於神的那份虔敬.
以神的好.
作為自己的好.
他以從上帝的好.
作為他行事為人的準則.
是一個正直的人.
他的禱告能夠深系上帝.
和上帝連繫.
去聆聽神的回應和啟示.
弟兄姊妹願主都幫我們.
幫助我們成為一個虔敬.
正直禱告和行動.
去見證我們上帝的人.
讓上帝祝福我們.
讓我們低頭禱告.
在耶穌的周邊.
有很多人很多的事.
今天我們藉著西面.
這個人物求主幫我們.
讓我們每一個在座的弟兄.

$^{1361}$和姐妹.
我們都立願成為虔敬的人.
成為公義的人.
成為禱告與你連繫.
去經歷聖誕.
去祝福我們.
去連繫去經歷聖靈.
的同在的人.
求主幫助我們.
多謝耶穌.
祈禱靠耶穌基督聖靈.
\newpage



\section{約伯記 42:1-6}
\label{sec:EJ51POSz_gM}
\textbf{2023-12-17 午堂《從前風聞有你,現在親眼看見你》約伯記 42\_1-6 講員:鄧泰康傳道}
\newline
\newline
連結: \href{https://youtube.com/watch?v=EJ51POSz_gM}{\texttt{https://youtube.com/watch?v=EJ51POSz\_gM}} ~~~~ 語音日期: 2023-12-19
\newline
\newline
\hyperref[sec:oC6cD_CKYi8]{\small{< < < PREV SERMON < < <}}
~
\hyperref[sec:index_chronic]{\small{[返順時目]}}
~
\hyperref[sec:index_scriptual]{\small{[返順卷目]}}
~
\hyperref[sec:JmKumENSOqs]{\small{> > > NEXT SERMON > > >}}
\newline
\newline
約伯記 42:1-6
\newline
\begin{longtable}{cl}
\hline
\hline
章節 & 經文 (和合本修訂版)\\
\hline
42:1 & \begin{tabularx}{0.7\textwidth}{X} 約伯回答耶和華說： \end{tabularx} \\ \\ \relax
42:2 & \begin{tabularx}{0.7\textwidth}{X} 「我知道，你萬事都能做；你的計劃不能攔阻。 \end{tabularx} \\ \\ \relax
42:3 & \begin{tabularx}{0.7\textwidth}{X} 誰無知使你的旨意隱藏呢？因此我說的，我不明白；這些事太奇妙，是我不知道的。 \end{tabularx} \\ \\ \relax
42:4 & \begin{tabularx}{0.7\textwidth}{X} 求你聽我，我要說話；我問你，求你讓我知道。 \end{tabularx} \\ \\ \relax
42:5 & \begin{tabularx}{0.7\textwidth}{X} 我從前風聞有你，現在親眼看見你。 \end{tabularx} \\ \\ \relax
42:6 & \begin{tabularx}{0.7\textwidth}{X} 因此我撤回，在塵土和爐灰中懊悔。」 \end{tabularx} \\ \\
[1ex]
\hline
\hline
\end{longtable}
$^{1}$靜慧早安.
很感恩能夠平安的.
再回到大家當中.
一切可以敬拜神.
可以參與事奉.
這個絕對是上帝的恩典.
今天很想透過這段經文.
和大家去分享.
用整個約伯記的脈絡.
用他的主題.
再加上這一次.
我在手術當中經歷的.
和大家去分享.
究竟約伯記裡面.
最出名的這句話.
我從前風聞有理.
我現在親眼見你.
究竟是什麼意思呢?.
我也很想讓大家知道.
或許我要先去假設.
可能在座裡面.
有些弟兄姊妹都沒有看過約伯記.
所以我很快地去講一講.
約伯記的整個脈絡.
首先在講之前.
我要先定位.
看聖經一定要知道.
究竟他的文體是什麼.
約伯記是一個戲劇.
記住約伯記是一個戲劇.
所以你可以放在.
詩篇,箴言裡面.
什麼意思呢?.
即是當你看的時候.
重點並不是.
究竟約伯記裡面.
所呈現的事情.
有沒有發生.
而是他裡面的訊息.
清楚嗎?.

$^{41}$即是你今天看電影.
你不會去問.
究竟電影裡面的某某某某角色.
是否真實存在.
其實不是.
是這個劇作者很想透過這個角色.
去跟你講一些訊息.
所以大家看約伯記的時候.
要回到訊息而不是重點.
我想舉個例子.
就好像當你看詩篇的時候.
他說上帝有大能的旁碑.
如果你去不斷去研究.
究竟上帝的手有多粗.
那你就出事了.
又或者.
他有些話.
詩篇第十七篇八節.
求你保護我如同保護眼中的銅人.
將我隱藏在你翅膀的陰下.
究竟上帝有沒有逆呢?.
你明白嗎?.
所以我想搞清楚.
看約伯記.
重點是這一點.
為什麼呢?.
因為約伯記常常引來爭議的是什麼呢?.
是他在天庭的一幕.
約伯記開始的時候.
他就說.
電影開場.
有一個天庭的一幕.
上帝在當中.
眾子或天使都在.
然後撒旦也在.
然後上帝主動的問撒旦.
你有沒有留意到.
我的僕人約伯.
然後撒旦就說.
有 我有見到他.

$^{81}$然後你看.
他對我有多忠心.
他對我多敬畏敬虔.
然後你會發覺.
魔鬼的回應是什麼呢?.
他對你當然敬虔.
你賜了那麼多東西給他.
你給了那麼多祝福他.
然後上帝的回應就是.
OK.
既然你這樣說的時候.
我容許你將他所有的祝福.
將他拿走.
在一下之間.
在一個人身上.
一下之間.
所有的子女.
所有的.
他的資產.
所有的資產.
一下之間都沒有了.
然後.
他的老婆.
沒有死.
然後就跟他說.
然後他就說.
這些是上帝賜給我的.
他不可以收回嗎?.
他沒有得罪神.
又回來了.
上帝再問撒旦.
你見到我的僕人嗎?.
他現在沒有了東西.
但你看他仍然對我那麼敬畏.
然後魔鬼第二次.
他就跟上帝說.
當然了 你都沒有傷害他.
他不見了的只是身外之物.
然後上帝就容許.
既然是這樣的時候.

$^{121}$我容許你傷害他.
但你要留下他的性命.
然後約伯身上就掌滿毒瘡.
然後.
再下去的時間.
他老婆就出聲了.
她說你到現在.
你還相信你的上帝嗎?.
你不如不要他了.
你看他怎樣對你.
然後約伯就叫他收口.
其實經歷了兩次之後.
約伯仍然很堅持.
要去敬畏這個上帝.
他很堅持要跟隨這個上帝.
然後下去的時候.
約伯就回到地上.
那個場景.
約伯很辛苦.
他有三個朋友就來到他身邊.
陪了他.
這個很重要.
有七天陪了他身邊.
七天沒有跟他說過話.
這個坦白說.
你身邊有沒有這樣的朋友.
只是陪你 但不敢跟你說話.
純粹陪伴.
三個非常好的朋友.
然後約伯開始.
說話的時候.
這三個朋友就回應了.
重點主要.
三個人都差不多.
你一定是犯了罪.
我們相信上帝很清楚.
如果你不犯罪 上帝是不會給你災難.
所以你認罪吧.
大概裡面就是這樣.
但約伯很清楚.

$^{161}$他自己不斷去澄清.
我沒有犯罪.
但我始終沒有得罪上帝.
然後.
到另外一個朋友.
再出現的時候 他就澄清一點.
他說在上帝裡面.
不一定犯罪.
不一定犯罪才會發生這樣的事.
不過.
上帝一定有些事想你學.
所以你就在裡面.
學上帝要你學的東西.
到最後.
約伯自己去開聲.
其實他裡面是.
他想澄清 訴苦.
我沒有犯罪.
甚至他想澄清.
自己的情況 他要和上帝.
有些甚至爭辯.
上帝.
他開始懷疑上帝你的不公平.
上帝為什麼你要這樣對我.
而去到最後.
的一幕時間.
就是上帝和約伯很多的對話.
上帝沒有.
直接回應他的提問.
沒有直接回應他的說話.
上帝只是不斷問他一樣東西.
你以為你是誰.
我今天我自己.
做了很多很多的事 難道.
你能夠去了解嗎 難道.
你能夠知道嗎 我想做的.
任何的創造之物的時候.
難道你能夠控制嗎.
說了一輪.
之後 接著就來到.

$^{201}$第42章這一幕.
就是約伯聽完.
之後 啞口無言.
接著他就和.
上帝說.
我知道你是萬能的.
我知道你的旨意是不能夠.
攔阻的 誰能夠用.
無知的言語 其實我在說.
自己 誰能夠用.
無知的言語去令到.
你的旨意隱藏呢.
我所說的是不明白的 即是說.
他和上帝所說的.
我不明白 這些事太奇妙了.
所以求你.
聽我 求你.
我問你的時候 求你.
告訴我 所以到最後的.
時刻 他說我從前風聞有理.
我現在親眼見你.
接著最後一句很重要.
所以我要撤回.
剛才我的說話.
我要在塵土.
和爐灰之中懊悔.
即是說到最後約伯去到.
一個地步 他很謙卑.
下來 他不再和.
上帝去爭拗或者.
疑問今天他所面對.
苦難的原因 他不再問.
他來到神的面前.
而整個約伯記的.
重點就是.
到最後上帝賜福給他.
而給以前的.
是加倍.
這個是整個約伯記的文.
難看的地方是在裡面.

$^{241}$有很多其實很.
違反我們今天人性.
為什麼呢.
當人面對困難的時候.
就很想你告訴我.
為什麼.
我們做傳達人就會經常.
經歷到.
突然之間遇到一些情況.
他們走來.
一登傳道 為什麼我發生這件事.
其實說真一句.
這件事我們做傳達人.
都很危險.
因為問題是我們只可以.
用聖經我們所有有限的知識.
去告訴你 嘗試解答.
給你聽 但是.
坦白說 你問我那個.
是不是最真正的原因.
我以為我真的不知道.
因為我不是上帝.
今天我想和大家去講的時間.
因為約伯記.
他不是.
為了要去.
解釋上帝指紋.
為什麼受苦.
很多人看約伯記是搞錯了.
如果我想告訴你.
這個才是約伯記的重點.
如果你想解釋為什麼受苦.
你就糟糕了.
你就會發覺.
很多人會有一個錯覺.
就是上帝就是在耍我.
我在地上好好的.
你在上面和魔鬼在爭論的時候.
你將我擺到.
這個地步.

$^{281}$這樣就錯了.
其實約伯記整個重點.
是要告訴我們.
神聖智慧的重要.
我想告訴大家.
剛才在約伯記當中.
有一段是很重要的.
就是上帝和魔鬼.
主動去提約伯的時候.
其實反過來.
是他對約伯.
非常之有信心.
你看不看到我的僕人.
他今天的時間.
他行在我的道路上.
上帝的眼光.
是非常之肯定.
他的僕人對他這份信.
所以他是預視到.
到最後的時候.
約伯他這種回應.
他是預視到的.
而你就會發覺.
整個核心主題.
我剛才說的就是神聖智慧的價值.
什麼叫做神聖智慧的價值呢.
簡單而言就是敬畏神.
降伏在上帝的主權當中.
約伯記就是要告訴你.
我們要馴服上帝.
就算去到一個很極端的情況下.
我們仍然信.
這個才是約伯記.
裡面一個很核心的主題.
你會發覺.
引發了二人約伯.
他活出敬虔的生活.
究竟是什麼原因.
這個才是約伯記的問你的問題.
是不是上帝賜給他的財富.

$^{321}$和安全呢.
魔鬼的挑戰是.
人只要得到上帝的祝福.
就能夠.
得到上帝的祝福.
你會發覺.
只要得到上帝的祝福.
即是好處.
所以就會對上帝敬畏.
和過敬虔生活.
這個是魔鬼的挑戰.
而上帝的回應.
是什麼呢.
不會.
聽得清楚嗎.
魔鬼的角度就是.
你是因為得到上帝的好處.
所以你對上帝就敬畏.
你對上帝就馴服.
但當上帝把這些東西拿走後.
你就不會再對上帝敬畏.
但上帝很清楚的回應.
他會.
所以你會發覺.
一個真正的智者.
其實會一直堅持.
活出敬虔的生活.
甚至面對著很多物質的流失.
或者是痛苦的生活.
其實是很難的.
因為你會發覺.
一個真正的智者.
面對著很多物質的流失.
或者是痛苦的時候.
他仍然會堅持.
他不會計較.
上帝給了多少報酬給我.
這個才是真正智慧的生活.
弟兄姊妹.
我很想你們知道.

$^{361}$有時候我們的福音.
是很容易會令到你有錯覺.
很多時候我們在台上.
講見證的時候.
往往都是我得到一些東西.
無論那個是實質的東西.
或者那個是.
心靈的東西.
所以我就信上帝.
你明白我的意思嗎.
很容易給你一個錯覺.
但其實.
我們的信仰真真正正.
最大的意義是什麼.
就是我們本來認清.
這個真相.
我們根本在地上.
所有的東西.
不是最有價值的.
我們在地上的東西.
本身是怎樣的.
我想形容的是自然而生的.
很簡單的一件事.
你上班就自然會有報酬.
對不對.
只不過我們.
控制不了多和少.
而多和少.
大部分就是上帝賜給大家.
對不對.
所以其實我想告訴大家.
嚴格來說無論上帝賜多給你.
賜少給你.
這個上帝都沒有改變過.
他都一樣那麼值得.
你去敬畏他.
和愛他.
再退後一萬步.
我們最終的一傳一福.
講的一個例子.

$^{401}$今天大家.
呼吸的空氣.
上帝有沒有收錢.
由你一出生開始.
空氣就自然.
供應.
你可以說.
鄧全道空氣值什麼錢.
大家有沒有入過醫院.
有沒有試過要吸氧氣.
知不知道氧氣機一個月.
要多少錢租.
弟兄姊妹.
原來有很多不知不覺之間.
我們和上帝原來.
是無法去比的.
在這裡.
約伯松當中他更深的.
學習到和認識.
上帝和他之間的關係.
他指向了的就是這個主題.
原來當人的心裡面.
一有少少偏差.
我們就會和上帝的關係.
就會拉開.
我們就會失去了.
真正的平安.
在約伯的裡面有三個地方.
叫做錯誤點.
三個朋友.
他一起這麼有心.
來陪約伯.
他一定是愛約伯.
但是三個朋友.
他的錯誤在哪裡呢.
你留意一下.
他所講的全部裡面的.
神學觀點.
其實是正確的.
他所講的神學觀點.

$^{441}$是正確的.
亦都是聖經.
經常教導.
告訴你聽.
因為人犯了罪.
所以上帝透過這些疾病.
透過這些東西.
去提醒你 挽回你.
是有的.
但是這三個朋友的錯.
在哪裡呢.
他就將一般的原則.
放在整個人生.
所有的面上.
這個就是錯.
即是說.
你將一些聖經的原則.
代替了上帝.
聽不聽得明白.
聖經本來.
是上帝給我們的.
但是.
你就將一些聖經的原則.
用在.
每一個場合裡面.
而你奪取了.
上帝的主權.
我想講 上帝是一個.
有靈的神.
他是可以.
隨他的意思.
去做任何事情.
即是說 如果你將一個.
聖經的通則 將我們自己神的通則.
應用在人們身上.
這個是非常危險.
舉個例子.
鄧全道.
可以因為罪的緣故.
所以上帝讓我.

$^{481}$在疾病當中重新再反省.
但不等於.
舉個例子 戴卓爾.
Roger病了 之後就是因為他.
犯了罪.
所以上帝就要透過這個病.
去叫他重新再思考.
No 不是.
上帝和.
每一個人.
他身上的互動.
都不同的.
我經常都形容上帝.
令我們最奇妙的地方.
對我們每一個人的救恩.
是度身訂做.
上帝最清楚.
你裡面欠缺些什麼.
上帝最清楚.
你裡面的心靈是些什麼.
所以我經常鼓勵弟兄姊妹.
你越愛那樣東西.
上帝.
就越是要拿走.
因為那樣東西.
成為了你生命當中的上帝.
所以弟兄姊妹.
不要特別 我想鼓勵弟兄姊妹.
我們今天做很多同行者.
做其他弟兄姊妹的陪伴.
請你千萬不要.
變成了這三個朋友.
你千萬不要.
用你自己所擁有的一套神學.
就不斷地跟他們說.
不斷地跟他們說.
叫他們這樣這樣.
不是.
反而我可以告訴弟兄姊妹.
開頭那七天是很有意義的.

$^{521}$是真正的朋友.
就是陪伴.
所以我亦都很相信.
當你曾經讀過一些輔導課程.
首先那個老師一定會教你一樣東西.
是什麼.
大家知不知道.
是閉嘴.
聽多過說.
第二個錯誤就是.
以利戶的出現.
其實是比起三個朋友好.
所以你會發覺到最後.
上帝去責備的時候.
責備了三個朋友沒有責備以利戶.
即是第四個朋友出現.
但是以利戶同樣一個錯誤.
他雖然解釋了.
不是每一個的罪.
都一定.
每一個病一定會和罪有關.
但是他仍然去.
堅持一樣東西.
上帝有功課要你學.
所以你一定要學到這個功課.
這個都是錯.
這個是很難了解.
不是啊 上帝.
如果不是為了任何原因.
你為什麼要我經歷那樣東西.
如果你想知道真正答案.
有一天你上到天間.
你問回上帝.
而是上帝有.
他這個很絕對的主權.
你明白我意思嗎.
有時候我們自己在地上.
我們很難用我們有限的智慧.
去知道.
我們只知道一件事.

$^{561}$就是上帝做的事.
一定不會無端端我們只知道這麼多.
但是究竟.
實際是不是要你學東西.
或者其他 不知道 我舉個例子.
一會兒我都會說.
我這次.
進去醫院裡面.
原來其中一個很重要的任務.
就是讓我旁邊的伯伯信了耶穌.
上帝可以做一件事.
不是為了你.
可以為其他人 明白我意思嗎.
但是總之上帝在過程當中.
有沒有讓你去受苦.
這個.
答你不到.
有沒有讓你去經歷一些事情.
我答你不到.
但是總之上帝有他自己的主權.
而第三個就是約伯他自己的錯誤.
沒錯 聖經裡面說得很清楚.
約伯的確是一個非常之有.
非常之做得好的義人.
他也都是一個.
真是.
他連在聖經裡面都說.
連他的子女.
你可以說是一個父母.
子女他們一起聚集.
他爸爸都不斷提醒他們.
你不要犯罪 你們要敬畏神.
你們不要做任何事得罪上帝.
所以整個.
這套戲劇裡面所形容的約伯.
是一個非常非常之聖潔的人.
但是約伯在裡面.
他做了一件事.
到後來當你朋友再去.
加強他這件事的時候.

$^{601}$他就不斷想去澄清.
我沒有犯罪 上帝我沒有犯罪.
他想不斷地澄清.
而接著就去到一個地步.
就是我不知道大家有沒有經歷過.
這件事.
你要證明自己是正確.
那你做什麼.
大家知不知道.
你要證明你自己正確.
你就是要證明對方錯誤.
明白嗎.
其實就是約伯推到極致.
他完全很清楚我沒有犯罪.
接著他開始.
他的言談之間.
開始去質疑上帝.
上帝你的公平公義出了問題.
是不是你看不到東西.
是不是你.
對我有什麼問題.
開始埋怨.
開始怨恨.
開始有怨憤.
就出來.
這個原來.
才是真正令我們.
當陷在一個.
困難的處境.
甚至是苦難當中.
真正令我們萬劫不復的點.
今天我們有機會做.
輔導的時候.
很多時候弟兄姊妹.
都困在以前.
一些的境況當中.
從成長的裡面.
或許有 從經歷的裡面也有.
他困在那個境況當中.
出不到來.

$^{641}$這個才是真正.
令我們苦的地方.
所以你就會發覺.
約伯他在裡面.
他.
上帝在裡面.
他和約伯的這個對話.
他完完全全沒有正面回答他這些問題.
他一直上帝很清楚.
約伯的主題就是.
我以一個創造者的身份.
去問了你很多很多很多的問題.
你這個受造的人.
在創造者的面前.
其實.
你見到什麼呢.
你的能力有多大呢.
你會發覺約伯他學到一樣東西.
他認識到.
自己公義的生活.
對上帝敬畏.
只是受造者.
應該有的自然回應.
大家記不記得.
當耶穌進城的時候.
有很多人.
一起歡迎他.
跟著的時間.
有些人叫那些人閉嘴.
跟著耶穌的回應是什麼呢.
就算那些人不去說.
不去讚美.
石頭都會讚美.
所以弟兄姊妹.
你就會發覺.
其實聖經告訴大家.
原來今天我們作為一個受造者.
當你去認識到.
上帝奇妙的作為的時候.
你本應就要.

$^{681}$活出一個公義的生活.
一個敬畏上帝的生活.
本身就應該是這樣.
剛才我在外面.
跟弟兄姊妹分享.
我兒子讀徐澤林小學.
他是很早期的畢業學生.
以前陳學長.
陳瑞如校長有一句金句.
我現在問現在的學生.
他們沒有聽過.
我自己見到他們都這樣說.
他跟學生的.
一面其中一句說話.
你做得好是應該的.
我以前剛聽的時候.
我很忍耳.
我說不是啊 小朋友需要讚的.
做得好就讚他.
但陳校長說.
你做得好是應該的.
到今天我終於明白為什麼.
你作為一個學生.
你在學校裡面.
你就應該好好學習.
你本應就要好好學習.
你本身就應該要盡力去爭取.
我們今天的學習.
不是為了要去贊取.
其他人的讚賞.
這個是你本身應該的.
其實就是回到傳道書裡面.
這個核心主題.
我們今天每一個被上帝創造出來的人.
我們本身就應該好好地按著上帝的旨意.
好好地去活的.
是不是啊.
如果你知道原來上帝創造你是這麼寶貴的話.
根本就不需要上帝給你什麼.
你都應該要這樣活.

$^{721}$而我想說.
當我們去經歷很多高高低低的時候.
正因為我們相信上帝的主權.
而最重要這個上帝是愛我們的.
這個上帝是會為我們預備最好的.
所以我們走每一步的時候.
我們都能夠可以很放心.
因為無論你今天所經歷的任何事情.
上帝都清楚.
上帝一定在這裡.
是這樣令到我們安心下來.
我們不會再去控訴這個世界.
不會再去控訴上帝.
是不是上帝你有問題.
是這一份的安穩平靜.
能夠令到我們帶著信心去過每天的生活.
所以你會發覺整個結尾到最後.
上帝約伯的回應就是剛才那段經文.
上帝我明白了.
他說我以前封民有你 現在親眼見你.
就是這個親眼見.
原來我們心裡面一有一點點對上帝質疑.
我們所有的平安過去的祝福.
會一次過全部都沒有了.
他唯一可以做的就是回到神面前的時候.
上帝原來我是這麼微小.
原來我是不知道整個圖畫.
但我只是求你和我一起走過.
所以他說我撤回.
我要懊悔.
為什麼我要和上帝說這樣的說話.
而整個的結尾就是約伯重新得到上帝的更大的祝福.
臨到他身上.
這個大概就是整個約伯記裡面.
他很想和大家分享的重點和核心.
我就說回我的手術.
我要很真誠的告訴大家.
鄧全道面對這個手術的時候.
我不是裝出來的.
我也不是因為全道人.

$^{761}$所以我要很堅強的站在大家面前.
因為全道人我不可以這麼弱.
所以你突然問我OK.
不是 絕對不是.
而是在過程裡面.
我自己其實是真的在經歷同一樣東西.
上帝其實我的健康是你賜給我的.
約伯說了一句話.
他可以賜予 他也可以收回.
健康不是必然的.
所以上帝就將我整個的健康.
我就交回給你.
我不知道為什麼我在這一刻裡面.
我會經歷了這件事.
我不知道為什麼會這麼嚴重.
整個過程當中我沒有問過上帝為什麼.
我只是告訴上帝.
我現在這一刻就將我完完全全的交在上帝的面前.
坦白說 在這個星期.
我們親愛的老牧師安息住.
其實養安裡面重新再想了很多.
老牧師他是非常勤力.
非常熱心.
做了很多很多很多的事情.
特別現在下了紅恩堂.
他經常記掛著紅恩堂.
弟弟姐妹和他們的一些事工發展.
有同工在他臨走之前去探望他的時候.
他都仍然記掛著這個星期.
他講的那篇講章.
所以今天我收到一個消息.
今天使母會在紅恩堂講道.
代表老牧師將這篇講章和大家分享.
很多時候我們面對這個處境.
我們都可能會問 上帝為什麼.
既然這麼愛你的僕人 你賜健康給他吧.
我們會問很多這些問題.
我們很多都不明白的.
但是弟弟姐妹 其實在這裡的時候.
正好就反映了上帝的主權.

$^{801}$我們不知道為什麼.
我們不捨得老牧師.
但是 但是.
老牧師的的而且確.
上帝宣告告訴我們.
他地上要跑的路 他跑了.
他守的道 他守了.
他覺得是適當的時候 他回來了.
所以他現在就是回到神的面前.
去領取這個榮耀的冠冕.
我們沒有一個人可以自己去決定多事奉一刻.
所以當我再回來能夠去事奉.
站在大家面前講道.
我很感恩每個弟兄姊妹來慰問我的時候.
我都是這樣說.
再熟一點 就像小芬這些.
我跟他再熟一點.
我就會跟他口不擇言地跟他說.
上帝認為我還沒做夠.
上帝有絕對的絕對.
所以無論上帝給我的是什麼情況下都好.
我都覺得帶著感恩的心.
現在這個疾病很感恩 為什麼呢.
因為原來我做了原本三個小時變成四個小時.
因為裡面出了一些問題.
原來有兩支釘插在脊椎上的骨.
撐開來做.
做完之後 拔出來的時候那些骨流血.
所以後來都做了一些治療.
醫生最初的時候最怕就是內出血.
因為那個位置剛剛壓著器官.
如果一旦內出血就很麻煩.
所以這個是恩典.
我現在其實清除了.
做完之後醫生才跟我說.
他說因為你這個手術如果不做的話.
我會怕你癱.
因為你壓的神經線壓得很厲害.
還有你壓得很扁.
所以現在做了 起碼清除了.

$^{841}$你不會再癱下去.
有些弟兄姊妹誤會了.
鄧全都你現在應該走得很好.
不是.
我說我主觀的感覺告訴大家.
我的走路 我的麻痺跟之前沒什麼分別.
是差不多的.
但醫生的回應就是.
因為神經線我不知道究竟是要半年還是一年.
之後才慢慢康復.
這個也要看你自己的身體.
這個我幫不了你.
但弟兄姊妹.
重點是不是要完全康復.
不是.
重點是.
我很快我就能夠繼續去事奉.
我繼續能夠去服侍.
這個才是重點.
在這個過程裡面.
其實我跟上帝說.
我就將命交給你了.
但我就有些要求.
什麼要求呢.
第一件事.
我真的很怕痛.
如果上帝你欣點我.
我欣點的話.
你給我不要痛.
第二件事.
鄧全都是靠口運藉.
醫生說過.
因為我做在前面.
我進去的時候會動到氣管 聲帶.
有機會葉狀患會沙.
甚至會失聲.
會有這些情況.
第二件事.
上帝我靠口運藉.
你讓我繼續說話.

$^{881}$所以我做完手術之後.
首先一個人一醒了.
我就很大聲叫 姑娘.
然後姑娘就問你做什麼事.
然後我就跟她說.
沒有 我只是試一下.
完全正常.
繼續洪量.
繼續可以和大家去角落.
繼續在過程裡面上.
給了我一些欣點.
我這個傷口.
直到裡面的地方.
到今天.
由發病到今天.
做完手術到現在.
我都沒有痛過.
是完全不同.
當然你可以說.
醫生手術好.
是的 醫生手術真的好.
但我相信連醫生.
因為都是一個介紹.
所以比較熟的醫生.
我都問他為什麼我不痛.
我解釋不到給你聽.
為什麼你不痛.
上帝.
我想告訴大家.
我今天和大家分享的原因.
不是我經歷了這些.
之後有這份平安.
不是 你明不明白.
是我之前的裡面.
我很清楚上帝.
你有絕對的主權.
你可以做任何的事.
所以很安然.
上帝在後面.
給了很多很多的欣點給我.

$^{921}$好了 另外.
我原本之前.
如果大家有一次見過.
我戴著頸箍.
我原本戴頸箍.
醫生最初和我說.
要戴三個月至四個月.
但我後來.
我回去的時候.
我和醫生覆診.
我早問他.
我什麼時候可以脫下頸箍.
然後他說.
不要緊 我幫你檢查.
然後做了幾次來來回回的檢查.
我大概一個月.
我就已經不用再戴頸箍.
另外剛才一個很奇妙的地方.
我要和大家分享.
我康復得很快.
我不知道大家有沒有試過.
我真的很奇妙.
做完手術之後.
應該下了麻醉藥.
是要等到麻醉藥過.
所以等腸再重新動.
你才可以吃東西.
好了 我做完手術之後.
應該是兩個小時左右.
我就開始咕嚕.
即是腸裡面的結.
然後時間就是.
姑娘說不要吃東西.
我就不要吃東西.
我忍了四個小時.
醫生就過來看我.
醫生就說.
我問他.
我問他什麼時候可以吃東西.
不要那麼快 先吃東西.

$^{961}$你的腸要動才可以.
我幫你聽聽.
一聽.
你的腸已經重新動了.
好了 吃吧.
我想強調.
我康復得很快.
很快.
其實我做完手術之後.
第二天我就已經可以下床.
然後我就戴著頸箍.
坐在床上不行.
我就覺得我要多做運動.
在病房裡.
就好像在走馬.
在走沙圈.
在病房裡這樣走.
好了.
我走的時候.
然後就隔壁.
有個伯伯.
他做菠蘿蓋的手術.
他做完之後.
其實他有點辛苦.
他有點痛.
然後我每天都聽到他有點叫.
於是就做一下物理治療.
我走.
然後那個伯伯.
不知道是不是太悶.
看到我在走的時候.
他說.
喂 你做什麼手術的.
然後我就跟他.
我跟他頸椎.
我給他看頸椎.
頸椎這麼嚴重.
你這麼快可以走.
我就說我不知道.
他說你不痛嗎.

$^{1001}$我不痛.
然後時間.
後來說開說開的時候.
他就知道.
他應該聽到我說話的語氣.
你上教會嗎.
然後我說.
我是傳道人.
哦 傳道人.
然後我說.
為什麼你這麼熟.
你好像很熟悉.
我的子女.
每個都信了耶穌.
然後就開始跟他聊天.
沒有事情做.
就坐在那邊聊天.
我走到這裡跟他聊天.
後來就知道.
原來他的子女信了很多年.
但是他幾十年都不肯信.
但是真的很有趣.
老人家 我不知道是不是老人家.
我就跟他說.
後來我突然間.
聖靈感動我.
你想不想不痛.
你想不想像我一樣不痛.
然後他說我當然想.
信耶穌.
是不是很敢開了這樣的笑話.
但是老人家說 好啊.
他真的跟我說 好啊.
那我就真的很認真地跟他說福音.
說完之後.
我真的跟他決志.
他真的決志.
我就跟他說祈禱.
祈完禱之後.
弟兄姊妹.

$^{1041}$不要不相信上帝對初信者的特別憐憫.
祈完禱之後.
他真的告訴我.
他不同.
而最奇妙的是什麼呢.
就是當我在這裡跟他聊天的時候.
那些子女來探望他.
然後他爸爸突然間對著他的子女說一句.
Hallelujah.
我沒有跟他說過Hallelujah.
然後我就知道.
爸爸是在耍那些子女.
明白我意思嗎.
原來那些子女是信了很多年.
然後就說 爸爸你沒事吧.
然後他就走來問.
他就問他.
為什麼你會說Hallelujah.
我相信耶穌了.
然後他說.
旁邊那個人他不懂得分辨.
我說我是傳道人.
旁邊那個牧師帶我來信耶穌.
後來那些子女就走來跟我說.
我才覺得很感動.
原因是什麼呢.
這個伯伯真的幾十年.
他的子女用盡方法.
他都不肯信耶穌.
他說他真的沒有想過.
在這次的情況之下.
他能夠這麼簡單.
順利地信了耶穌.
然後我就跟他的子女開了一個玩笑.
我說或許這次是上帝讓我得了這個病.
特意傳給你爸爸的呢.
弟兄姊妹.
當你對上帝能夠有這一份肯定的時候.
災難都可以變成祝福.
當然還有更多更多更多的恩典.

$^{1081}$我很難在這裡這麼長時間跟大家分享.
但在這裡的時候.
我知道這次教會為我籌款.
去面對這個手術.
我很感恩.
因為會計和公告告訴我.
報告之後一個星期.
我原本的需要.
當時的手術費是23萬.
但是到一個星期已經有28萬.
到現在已經超過30萬.
我很感恩.
因為過程裡面.
其實我所需要的絕對是少.
後來我註冊了一些保險之類的.
我只需要向教會申請十多萬.
但是我很感恩.
能夠上帝透過這次的機會.
將剩下的錢能夠放在慈惠基金.
繼續去幫其他弟兄姊妹.
所以弟兄姊妹.
其實當你去相信上帝.
有這一份這樣的恩典的時候.
上帝的恩典就會伴隨你一生.
你放心吧.
你放心去相信上帝.
無論你今天遇到任何的情況.
你只需要肯定上帝是愛你的.
你仍然堅持去敬畏祂.
你不要質疑祂.
祂一定會在最後.
就好像約伯記一樣.
祂會將祝福更大的賜給大家.
而最重要的不是物質的祝福.
而是這份平安會跟隨你一生一世.
所以在這裡我也要去說.
弟兄姊妹.
或者今天一些未信主的朋友在當中.
我很盼望你們見證告訴你們.
人世間我們總是會經歷苦難.

$^{1121}$但是我們作為基督徒.
最大的祝福.
因為有上帝.
所以如果你願意去相信上帝.
你找我們任何一位傳道人.
我們很樂意跟你分享.
和跟你一起決戰祈禱.
我們是同心祈禱.
所以這份我們獻上感恩.
我們多謝贊美你.
主啊 我們求主你去讓我們每一個.
能夠在神的面前.
我們全然的絕對相信你.
我們完全全然知道上帝你的恩典和祝福.
我們求你親自的保守幫助我們.
很願意我們能夠因為相信你的緣故.
我們的平安.
我們能夠伴隨你一生.
我們很盼望.
我們能夠親眼看見你.
祈禱奉教主耶穌.
得性明智其金.
\newpage



\section{詩篇 90:1-17}
\label{sec:JmKumENSOqs}
\textbf{2023-12-17 早堂《赤子之心》詩篇 90\_1-17 講員:陳韋安牧師}
\newline
\newline
連結: \href{https://youtube.com/watch?v=JmKumENSOqs}{\texttt{https://youtube.com/watch?v=JmKumENSOqs}} ~~~~ 語音日期: 2023-12-19
\newline
\newline
\hyperref[sec:EJ51POSz_gM]{\small{< < < PREV SERMON < < <}}
~
\hyperref[sec:index_chronic]{\small{[返順時目]}}
~
\hyperref[sec:index_scriptual]{\small{[返順卷目]}}
~
\hyperref[sec:I2kkhTmf8Lk]{\small{> > > NEXT SERMON > > >}}
\newline
\newline
詩篇 90:1-17
\newline
\begin{longtable}{cl}
\hline
\hline
章節 & 經文 (和合本修訂版)\\
\hline
90:1 & \begin{tabularx}{0.7\textwidth}{X} 主啊，你世世代代作我們的居所。 \end{tabularx} \\ \\ \relax
90:2 & \begin{tabularx}{0.7\textwidth}{X} 諸山未曾生出，地與世界你未曾造成，從亙古到永遠，你是神。 \end{tabularx} \\ \\ \relax
90:3 & \begin{tabularx}{0.7\textwidth}{X} 你使人歸於塵土，說：「世人哪，你們要歸回。」 \end{tabularx} \\ \\ \relax
90:4 & \begin{tabularx}{0.7\textwidth}{X} 在你看來，千年如已過的昨日，又如夜間的一更。 \end{tabularx} \\ \\ \relax
90:5 & \begin{tabularx}{0.7\textwidth}{X} 你叫他們如水沖去，他們如睡一覺。早晨，他們如生長的草； \end{tabularx} \\ \\ \relax
90:6 & \begin{tabularx}{0.7\textwidth}{X} 早晨發芽生長，晚上割下枯乾。 \end{tabularx} \\ \\ \relax
90:7 & \begin{tabularx}{0.7\textwidth}{X} 我們因你的怒氣而消滅，因你的憤怒而驚惶。 \end{tabularx} \\ \\ \relax
90:8 & \begin{tabularx}{0.7\textwidth}{X} 你將我們的罪孽擺在你面前，將我們的隱惡擺在你面光之中。 \end{tabularx} \\ \\ \relax
90:9 & \begin{tabularx}{0.7\textwidth}{X} 我們經過的日子，都在你震怒之下，我們度盡的年歲，好像一聲嘆息。 \end{tabularx} \\ \\ \relax
90:10 & \begin{tabularx}{0.7\textwidth}{X} 我們一生的年日是七十歲，若是強壯可到八十歲；但其中所矜誇的不過是勞苦愁煩，轉眼即逝，我們便如飛而去。 \end{tabularx} \\ \\ \relax
90:11 & \begin{tabularx}{0.7\textwidth}{X} 誰曉得你怒氣的權勢？誰因著敬畏你而曉得你的憤怒呢？ \end{tabularx} \\ \\ \relax
90:12 & \begin{tabularx}{0.7\textwidth}{X} 求你指教我們怎樣數算自己的日子，好叫我們得著智慧的心。 \end{tabularx} \\ \\ \relax
90:13 & \begin{tabularx}{0.7\textwidth}{X} 耶和華啊，我們要等到幾時呢？求你轉回，憐憫你的僕人們。 \end{tabularx} \\ \\ \relax
90:14 & \begin{tabularx}{0.7\textwidth}{X} 求你使我們早早飽得你的慈愛，好叫我們一生一世歡呼喜樂。 \end{tabularx} \\ \\ \relax
90:15 & \begin{tabularx}{0.7\textwidth}{X} 求你照著你使我們受苦的日子，和我們遭難的年歲，使我們喜樂。 \end{tabularx} \\ \\ \relax
90:16 & \begin{tabularx}{0.7\textwidth}{X} 願你的作為向你僕人們顯現，願你的榮耀向他們子孫顯明。 \end{tabularx} \\ \\ \relax
90:17 & \begin{tabularx}{0.7\textwidth}{X} 願主－我們神的恩寵歸於我們身上。願你堅立我們手所做的工，我們手所做的工，願你堅立。 \end{tabularx} \\ \\
[1ex]
\hline
\hline
\end{longtable}
$^{1}$剛才在開始等節目進來的時候.
我就在外面打招呼.
再見到我們在那裡掛了2001年建堂的根基照片.
不知道大家平時走過有沒有再看一下.
如果你在的時候.
我發覺那時候都已經有二十多年了.
那時候我都已經二十歲了.
當中大家就將自己年齡減二十.
剛才黃恩提見到我的時候.
都說這個和我兒子一起做小丑的.
我們一起做小丑.
學他小丑班 學他福音小丑.
轉眼間來到聖誕節.
大家又會開始發覺.
2023年將要過去了.
大家年紀又會加震.
今天我們一起來思想.
有關我們的年歲 我們的生命.
特別今天安排在祖堂是很有意思的.
這裡各位的年紀.
雖然我講到十幾年.
但應該都是比我大.
大部分來說都是年紀長過我的.
我們一起來在聖經裡面.
思想我們的成長 我們自己的生命.
不知道你什麼時候覺得大過.
這麼多年裡面.
哪一刻你覺得自己是一個大人.
我自己覺得自己是大人是什麼時候.
我都不記得什麼時候.
我自己覺得不是大人是什麼時候.
我是很記得的.
我記得我二十四歲的時候.
我讀神學.
一年級就去了一間教會實習.
第一次去到在占卜堂.
占卜沙嘴那邊有一間教會實習.
那時候在長洲穿著西裝.
背著行李箱去教會實習.
也是一間中產的教會.

$^{41}$占卜堂也是一間.
去到的時候.
我自己在廁所上照鏡.
沒有人看都覺得自己不像大人.
覺得自己是一個小朋友.
背著西裝的樣子.
對著一群成年人.
覺得自己還是很小.
什麼時候覺得自己大過呢.
當然是當爸爸的時候.
當我自己當爸爸的時候.
都覺得自己好像已經大過了.
不過這都不是一個最深刻.
覺得自己最深刻的一刻.
開始白頭髮都未必是.
最近我開始有一個.
不能夠不說我自己大過.
我自己開始 上次也說過.
我自己有老花.
都有老花.
看東西真的看不到.
有些很小的都看不到.
後來我就配了一副漸進鏡.
不過我原來很心散光.
其實都不是很行.
很心散光原來配老花鏡.
其實都不是很行.
大家不明白我說什麼.
以後就知道什麼事了.
所以發覺自己的生命.
不得不承認已經是大過.
甚至和年老扯上關係.
已經再不是一個年輕人.
極其量只是一個年輕人.
一個扮年輕人的大人.
今天我們會看詩篇90篇這段經文.
我們思想我們的成長和年歲.
我讀一次.
因為我們主要看頭十二節.
主在你世世代代作我們的居所.

$^{81}$珠山未曾生出.
地與世界你未曾造成.
從耿古到永遠你是神.
你使人歸於塵土說.
你們世人要歸回.
在你看來千年如過的昨天.
如已夜過的現在這一更.
你叫他們如水沖去.
他們如睡一覺.
早晨他們如生長的草.
早晨發芽生長.
晚上閣下枯乾.
我們因你的怒氣而消滅.
因你的憤怒而驚惶.
你將我們的罪孽擺在你面前.
將我們的隱惡擺在你面光之中.
我們經過的日子都在你震怒之下.
我們年盡的歲月好像一聲嘆息.
我們一生的年日是七十歲.
若是強壯可到八十歲.
但其中所經過的不過是勞苦愁煩.
轉眼成空.
我們便如飛而去.
誰替你怒氣的權勢.
誰安著你歌手的敬畏.
替你的憤怒呢.
求你指教我們怎樣所算節的日子.
好讓我們得著智慧的心.
多謝祈禱.
最後求主你多來說話.
我們知道無論是甚麼年日.
你仍然對我們說話.
求主你讓我們能夠在.
藉著詩篇九十篇裡面.
來細聽,昔日摩西一番的說話.
求主你這樣來幫助我們.
奉主名求 阿們.
今天我們用一個特別的視野.
來閱讀詩篇九十篇.
先說說詩篇的背景.

$^{121}$詩篇九十篇是屬於.
詩篇卷四裡面的第一篇.
其實一百五十首詩篇.
是分了五大卷.
九十篇是屬於第四卷的部分.
雖然這首詩的題目.
叫做神人摩西的祈禱.
但這首詩的作者.
不一定是摩西本人.
或多數都不是摩西本人.
這首詩或以色列人.
以摩西的名義來說述.
臨終時的摩西對生命的反思.
可能大家最熟悉的就是這句.
我們一生的年日是七十歲.
若是強壯可到八十歲.
大家不少基督徒都用.
這個Biblical的標準.
來計算自己的衡量.
自己是到未還是Over了呢.
是否已經賺多了呢.
基本上來說.
如果你聽過這篇詩篇的講道的話.
大概都可以告訴你.
人的一生裡面的年歲.
我們的生命和死亡.
人的渺小.
上帝的永恆等的題目.
正如C.P.Socra的作者.
頭四節這樣說.
主啊!你的世世代代.
將我們的居所.
珠三未曾生出.
地與世界未曾造成.
從敢古道永遠你是神.
使人歸於塵土.
你們世人要歸回.
在你看來千年如過昨日.
如夜間的一更.
都是講一些永恆的角度.

$^{161}$來看我們生命的短暫.
不過如果你細心閱讀這篇詩篇的話.
其實你會發現.
詩篇不單單是講我們的生命很短暫.
時間過得很快.
人生將要過去.
而它講述我們人一生的年歲裡面.
是怎樣渡過我們的年歲.
它的年歲是怎樣流逝.
怎樣過去.
我的意思是.
如果你讀這篇詩篇的時候.
你細心發現.
你會看到這篇詩篇不單單是講上帝.
看千年如過一日.
我們的生命過得很快.
而是我們的生命不單過得很快.
而且過得不太好.
在神眼中.
真的,弟兄姊妹.
如果你看這篇詩篇的時候.
如果用年老的摩西來講述這篇詩篇的時候.
詩篇概嘆的不單單是人生短暫.
而是我們人年紀成長成熟.
一種少許嘆息和控訴.
甚至可以說.
《詩本十篇》是一篇去unify我們人生年歲一首詩篇.
找不到一個更加貼切的字去講這個unify的字.
或者說.
《詩本十篇》是對於我們人生裡面所謂成年人.
長大的人.
一些大人.
有一種控訴.
《詩本十篇》不單單是說我們時間過得很快.
人生的世上日子有限.
而是說就算讓你過到幾十歲.
五六十歲.
七八十歲.
其實都是很沒什麼的.
不知道當中大家有沒有一些比較年輕的弟兄姊妹.

$^{201}$或者你曾經以前讀中學的時候.
有沒有一種感覺.
對於一些比你年長的人.
那些長輩.
有一種很敬畏.
有一種甚至是恐懼的感覺.
這些原來叫權威恐懼症.
對於一些年紀大的人.
或者你的老闆.
有一種權威恐懼症.
以前會覺得少許.
當我讀中學的時候.
或者讀神學的時候.
對於一些四五十歲的老師校長.
我是特別很緊張.
對於一些大人很緊張.
他們是大人來的.
他們是很有智慧.
要尊敬他們.
我和他們好像參考院的感覺.
當我自己成為一個四十幾歲的人的時候.
發覺原來四十幾歲的人都是一些沒什麼的.
都是一些很普通的.
一坨東西都不奇怪.
都是一些很普通的一個人.
只不過看起來好像很大個.
年紀大.
其實真的只不過是年紀大.
大家不知道有沒有聽過.
陳昇有一首歌.
陳昇是台灣的歌手.
他到現在都口腔癌就康復了.
年紀都應該六十幾歲.
都快七歲的年紀.
六十幾歲的時候寫首歌曲.
叫做《像我們這樣的人》.
MV就是一班廢老.
彈結他和一個美女一起.
大家一起唱歌.
首詩歌正正是講什麼呢?.

$^{241}$就是講一個人年紀大就越是敗壞一種的無奈.
歌詞這樣說.
我在天上了多少網.
我在天用了多少塑膠袋.
在天走了多少路.
我在天為了多少的夢.
我在天說了多少謊.
在天撒敗了多少的諾言.
一千個宿醉的痛苦.
一半不了真心的懊悔.
我願若去跟好人吐槽.
卻跟噁心的人低頭.
像我們這樣的人.
是不是應該在上世紀就死去.
當然是一種很浪漫的嘆息.
詩篇九十篇大概都是感嘆這件事.
如果這篇詩人的作者.
是以摩西的名義去寫的時候.
一個年紀極度老邁的摩西.
一個八十歲被枯燥的老頭.
在曠野裡面經歷了滄桑四十年的人.
一百二十歲離開不了曠野的老人家.
當他去到他的生命終結的時候.
他發現其實一個很真實的經驗.
就是年紀大不代表什麼.
我們一起去看看.
了諒一下.
整篇九十篇裡面如何去講年紀.
第五六節 經文這樣說.
早晨他們如生長的草.
早晨發芽生長 晚上閣下枯幹.
經文告訴他人的一生就像草一般的短暫.
不過如果我們細味經文的時候.
經文不單單是講生長的草生命短暫.
他不斷去強調.
生長的草其實不斷地來生長.
很快的.
早晨發芽生長 晚上閣下枯幹.
早上仍然是小朋友 中午變成成年人.
晚上就步入晚年而死去.

$^{281}$這篇九十篇不斷去描述各式各樣的生長.
不斷去描述歲月這回事情.
究竟這篇九十篇如何去理解歲月呢.
我們再看看第七至第九節.
詩人這樣說.
我們因你的怒氣而消滅.
因你的憤怒而驚惶.
你將我們的罪孽擺在你面前.
將我們的隱惡擺在你面光之中.
我們經過的日子都在你震怒之下.
我們度盡了年歲好像一聲嘆息.
對詩人來說 人的一生在上帝的震怒之下.
在上帝的怒氣之中.
人類的罪孽和隱惡都擺在神的面前.
這段所說的罪孽不是一般人的罪孽.
而是甚麼呢.
詩人強調一個人是隨著他的年歲增長積累的罪孽.
詩人強調你的年紀越是增長.
你積累的罪孽和震怒.
上帝都是計算在內.
尤其是第九節這樣說.
我們經過的日子都在你的震怒之下.
我們度盡了年歲好像一聲嘆息.
詩人說我們一生經過的日子都在上帝的震怒之下.
真的 你年紀越大.
你經過的日子越多.
你走過顛覆到做過的錯事.
或者是你曾經碰過的罪孽.
都會隨著你走過的路徑增加.
而隨之而去增加.
就好像福音書裡面.
耶穌見到一個被人用石頭打的婦人的故事一樣.
耶穌就問你的年長.
你的年紀越是增長.
你的罪孽就會越重.
這是大家可能經過的事情.
曾經做過的.
知道自己做過的.
總比一個小孩子做得多.
我以前是.

$^{321}$上教會之前.
都有做過社會服務.
叫做特能童軍.
在朗平村裡面.
幫一班輕度弱智的朋友.
來到去大力做童軍.
這班人大概都二十多歲.
因為他們的心智上.
應該是適合小朋友童軍.
即是那些叫做小童軍.
初時我們童軍.
因為他們穿短褲.
如果是小朋友的話.
就安排他們穿一條短褲的截帽.
有個童軍制服這樣穿.
後來發覺因為他們太大了.
穿短褲就好像很醜怪.
後來就安排他們穿長褲.
和穿一個比較大人的服裝.
我發現原來.
不要以為他們二十多歲的小童軍.
就好像很可愛.
他們其實仍然會做很多犯罪的事情.
他們會買色情漫畫看.
他們那時候會做很多錯事.
會罵人的.
會大家吵架.
他們只不過是不像大人.
他們不會將自己醜怪的事情去隱藏.
他們只不過和小朋友一樣.
將自己有的不好的想法.
就讓人看到.
所以年紀越大.
其實相比起來.
絕對不是有任何的好.
所以詩篇最後一篇第十節.
我們一生的年日是七十歲.
若是強壯海盜八十.
但其中所驚誇不過的.
是老虎受煩轉眼成凶.

$^{361}$我們並如飛而去.
詩人說你以為你自己很叻.
做了那麼多年人.
所驚誇的只不過是老虎受煩.
大家不要誤會我不是不敬老.
事實上這裡所說的不是說老人家.
而是要我們去成長.
去想想究竟怎樣看我們的成長.
怎樣看我們的成熟.
或者我們的年歲.
詩篇正正是一首去控訴年紀.
年月長大的一首詩篇.
你覺得自己已經長大了嗎?.
或者你覺得怎樣才叫成長?.
怎樣才叫成熟?.
我們經常說我們要做回小朋友.
經常說我們要有赤子之心.
其實當我們想的時候.
做小朋友是否一件好事呢?.
其實成都很好.
我們都要成長 需要成熟.
難道我們不是要成熟嗎?.
聖經不是說我們不要做小孩子嗎?.
作為一個人我們不是應該要成長嗎?.
我們做一個大人.
還是做一個如耶穌所說的小朋友呢?.
我想問題關鍵是.
我們怎樣去看成長.
怎樣看長大.
怎樣看大人.
推翻百步去想這個話題.
當然長大是好的.
上帝創造我們.
我們有上帝形象.
理論上我們越是長大成人.
我們就越像上帝.
我們就越是成熟.
我們就越是做一個更加好的人.
我們就遠離小孩子的階段.
越來越是一個上帝所欺厭的人.

$^{401}$理論上成長是一件好事.
成長當然是一件好事.
問題是當我們在人生裡面.
我們往往未必所經歷的是一種叫做成長.
我們只不過是勉勉強強地去長大.
滾過了.
就好像那些激素雞,激素豬.
明明是沒有足夠空間去長大.
就鞠躬打一針.
肉多了,肉爆出來.
外表是大隻了.
樣子是成熟了.
但其實裡面是沒有真實的實質在當中.
稱之為強迫長大.
面對一個問題.
你可以用一個捷徑,某個方法.
繞過它,滾過它,快速過了它.
你是長大了的.
你成為了一個大人.
但可能有些東西你還未真正在神面前去長大.
面對著這個世界.
面對著這個年頭,這個社會.
很多時候我們是一個一個逼出一個一個大人出來.
你沒有辦法不成長.
你出來工作了.
沒有辦法不去用一個大人的眼光和方式去面對這個世界.
所以面對著這個世界的時候.
我們很多時候都是被逼出來.
成為了一個大人.
真人,你叫一個小朋友長大,有很多方法.
但不是每一個方法都是健康的方法.
有很多很容易的方法.
很多很快捷的方法.
很多很能夠回應問題的方法.
但在神的眼中.
這些不一定每一個都是一個好的方法.
這些反而是在上帝的震怒之下的成長.
真人一個人年紀大了.
做事了得了.
經驗多了.

$^{441}$不代表一個真正的成長.
如果你上班的話.
你看看你公司office裡面的那些OC food.
很多時候你就知道.
所謂OC food是甚麼.
就是那些經驗豐富,老練的大人,前輩.
不過in one way.
在一個錯誤的方法裡面.
厲害了,長大了,經驗老到了.
他們厲害在哪裡呢?.
可能就是一些走精面,射波,politics,賺錢的方法裡面.
所以真正關鍵是你是怎樣的長大.
你是怎樣的一個大人.
最近聽一個曾經做保險的弟兄這樣分享.
他說保險裡面行家有這樣的說法.
他說一年保險三年人.
你看見保險人你都知道.
一做保險.
人就立即經歷很多人生百態.
做人處事,研探技巧的事.
豐富了,厲害了.
人的成績深了.
做人都圓滑了.
但這個是不是真正的成長呢?.
一年保險三年人.
他做了保險兩年.
已經是六年人了.
不過他最後都沒有做到.
後來都看不過眼.
很多事情都覺得自己不適合自己的性格.
不能夠做到自己這樣.
當然保險是忠誠的.
不是批評這個行業.
起碼不適合他去做.
我覺得這件事情.
聽說你做了這麼多年人.
成為了一個大人.
比起二十年前的你.
或者是剛出來工作,讀書的你.
有什麼不同呢?.

$^{481}$很好的來到去想一想.
說話自圓其說了.
懂得隱藏自己多了.
公關笑容變漂亮了.
推銷自己懂得了得了.
吹水吹得大了.
究竟你是圓滑還是虛偽.
是自愛還是自私.
懂得疼愛自己.
還是過分放縱自己呢?.
作為大人.
我經常都有一句很有霸氣的話.
小朋友才可以做選擇.
大人不可以做選擇.
但是決定不做選擇的你.
你怎樣去選擇做一個大人呢?.
我自己呢.
我圓圓我女兒.
今年都九歲.
快十歲了.
如果很久沒有見她的話.
她已經不是一個小朋友的樣子了.
已經是一個少女的樣子.
一年前還是買衣服.
喜歡買粉紅色.
那些Unicorn的衣服.
今年已經買了黑色衣服.
灰色衣服那些.
隨著年紀越大.
她當然都面對一些少女的衝擊.
就是一些憤怒.
有些情緒.
有些陰謀論.
變成一個.
她回家後就說.
今天老師騙我們.
原來今天那個魔術表演取消了.
老師針對他們.
騙他們.
很多同班女生的陰謀論都出來了.

$^{521}$我自己呢.
為她祈禱.
叫她不要長大得太快.
寧願她慢慢長大.
寧願她保持小朋友的天真.
在一些正確的事情上長大.
慢慢地.
好好地長大.
赤字之心.
這個字是我們今天要講的題目.
這句話其實是來自於萬子的一句話.
全句叫做.
大人者不失其赤字之心者也.
一句來自萬子尼留篇的一句話.
萬子這句話是什麼意思呢.
他說大人者不失其赤字之心者也.
所謂大人.
意思未必是我們所理解所講的大人.
不是普通長大的人.
所謂大人.
其實就是那些能夠成就大事的人.
一些偉大的人.
一些所謂中國所講的君子.
一個作為君子.
或者對於一個我們信仰的角度來說.
一個上帝所喜悅的人.
一個好的人.
大人.
就是一些能夠不丟去赤字之心的人.
這個就稱之為一個大人.
一個好的大人.
我們知道.
世上有好的大人.
也有些不好的大人.
不好的大人會很厲害地做不好的事.
好的大人.
是會幫助這個世界做一些好的事情.
不是每一個大人都是真正的大人.
不是每一個真正長大了.
年紀大的人.

$^{561}$都是上帝所喜悅的人.
《詩篇》正正教導我們.
怎樣去當一個大人.
雖然《詩篇》不斷地對歲月,年紀和成長作出控訴.
不過《詩篇》最後的最後一篇.
這樣教導我們.
怎樣去做大人.
第十二節這樣說.
求你指教我們怎樣數算自己的日子.
好叫我們得著智慧的心.
求你指教我們怎樣數算自己的日子.
好叫我們得著智慧的心.
究竟怎樣去數算自己的日子呢?.
怎樣來到去.
為何詩人要教導我們去數算日子.
怎樣去數這些日子呢?.
一個不爭的事實.
我們每天都在長大.
無論你今日的年紀如何.
你每天都仍然在長大.
你今日度過了一日.
你就長大了一日.
你的生命過多了一日.
你就長大了一日.
這句說話比非洲一分鐘有六十秒過去更加真實.
一天的度過.
你就一天的來到去長大.
所以詩人就說.
在這個冥冥中的道理裡面.
我們要好好地去禱告祈求上帝.
叫我們好好地去審查我們每一日.
審查我們怎樣來到去長大.
我們要有好有智慧來到去選擇.
怎樣來到去長大.
不要隨便長大.
一個未能夠真正長大的地方.
不要隨便去長大.
寧願保留一些簡單愚笨的地方.
那些愚笨的地方.
正正就是那些仍然是善良.

$^{601}$正直正義的地方.
我們寧願愚笨地善良.
都不想成熟地去失去自己.
每個人都有些愚笨地方.
但不要緊.
寧願保持這些好像未能夠長大的地方.
寧願正確地成長.
都好過隨便地滾過去繞過去.
我想這個正正就是赤子之心的意思.
我自己都不是一個很會說話的人.
不單止口齒不伶俐.
更加不是很會說話.
有時說話和人交談就接不下去.
或者突然間會有空隙.
說話是我很愚笨的地方.
不過就算是沉默寡言也好.
擔心也好.
最少我都選擇有時候去愚笨地.
去誠實地去說話.
不要言過其實.
不要純粹討人喜歡.
赤子之心正正是說這種東西.
好好地去數據我們的日子.
好好選擇每一天你怎樣來長大.
我們仍然有成長的地方.
我們作為一個大人.
甚至一個老年人.
仍然有很多成長的地方.
每天好好地去求問上帝.
我怎樣地數算我的日子.
我今天有些什麼比昨天更加做得好.
在神眼中做得好.
這個小小慢慢地進步.
總比表面上的進步好.
我經常覺得早堂崇拜是一個很有趣的崇拜.
很多時候見到我這麼多年的弟兄姊妹.
看著大家在這裡.
我覺得這個崇拜可以成為一個兒童崇拜.
一個早堂兒童崇拜.
一個大家可以在神面前.

$^{641}$一群小朋友.
我們去做一些對的事情.
一些美善的事情.
在神面前.
可能你整個星期都是人家的公公婆婆.
夜夜麻麻.
但回到星期日八點半的崇拜裡面.
做一個在神面前的小朋友.
因為上帝去監察我們的生命.
數算我們的日子.
我們期望我們能夠回到天家的時候.
能夠可以帶著這份赤子之心.
去見我們天父上帝的面.
我們祈禱.
因為你每天視下你的恩典和慈愛.
但我們知道我們人的生命.
每天都在你的震怒之下.
我們成長.
我們年紀大了.
但求主你這樣來念穩我們.
讓我們能夠在你面前過得正直.
每天在你面前.
能夠可以慢慢地去數算我們的日子.
我們謹慎小心地度過我們的年日.
求主你赦免我們.
更加念穩我們.
求求你讓我們能夠可以.
每天越來越年老的時候.
在你面前更加被你所喜悅.
我們單單地來追求這件事.
我們祈求你幫助.
奉主名求.
阿們.
\newpage



\section{撒迦利亞書}
\label{sec:I2kkhTmf8Lk}
\textbf{2023-12-21 周四晚查經班《歸回先知書研讀(10)撒迦利亞書》 講員:曾裔貴牧師}
\newline
\newline
連結: \href{https://youtube.com/watch?v=I2kkhTmf8Lk}{\texttt{https://youtube.com/watch?v=I2kkhTmf8Lk}} ~~~~ 語音日期: 2023-12-23
\newline
\newline
\hyperref[sec:JmKumENSOqs]{\small{< < < PREV SERMON < < <}}
~
\hyperref[sec:index_chronic]{\small{[返順時目]}}
~
\hyperref[sec:index_scriptual]{\small{[返順卷目]}}
~
\hyperref[sec:pZkchuW18Ro]{\small{> > > NEXT SERMON > > >}}
\newline
\newline
$^{1}$好,弟兄姊妹,我們天氣雖然冷,但是室內也很溫暖,.
我們也能夠繼續去讀聖經,.
和大家先有一個簡單禱告,.
天父求禰讓我們能夠有一個明白的心靈,.
當我們去看舊約的撒迦利亞書,.
特別是關乎到在最後這六章,.
是有一個很明確指向將來的耶路撒冷,.
將來的神的國度,其實不單止是地上的國度,.
去到第十四章,有天上的形態出現,.
所以我們都盼望當我們去研讀這些經文的時候,.
有些是和當時代的處境有關的,.
尤其是對於治理以色列的墓者的勸告,警告,.
都需要我們大家都能夠在這個時代,.
我們來明白,.
我們將今晚我們要看的經文,.
十二,十三章仰望在神的手中,.
讓我們大家都能夠在過去片段,.
斷斷續續,盼望我們都能夠重拾思緒,.
又或者透過錄音我們可以再重新整理,.
使我們聽得更加具體,明白,.
我們仰望在你的手中,奉耶穌基督的名,阿們..
再一次感恩我們又能夠恢復到曾經班,.
我們一路這樣看,其實工程都很好大,.
能夠一路看回歸的歷史,.
有以色拉記,有尼欽米記,.
也有以史帖記,.
我們看到神在不同的時代的工作,.
藉著重建無論是聖殿,聖城,.
又或者是就算以色列人沒有回歸,.
但是都可以藉著在波斯帝國的書山城,.
一件我們無法想像的事情,.
一個叫哈曼的人,本來要滅絕猶太人,.
但是可以使到那個就坐變成一個祝福,.
有了這個普爾節,.
使到以色列人能夠成為每年都去歡慶的一個日子,.
這些都是歷史裡面主前耶穌未來這個世界之前,.
四百多年的時間出現的歷史的場景和故事,.
這些都是我們需要去留心,.
也感恩有回歸的先知書,.
貼近了和回歸的領袖,.

$^{41}$來到去很平衡地去講那個時代,.
先知一定是會聚焦在那個時代對神的心,.
百姓是怎樣跟隨神,.
和他們要做什麼可以帶來那個改變,.
所以不只是拼命去建工程,建聖殿,.
或者重建那個城牆,.
而是需要有那個心靈上面的重建,.
所以先知的預言無論是透過一些比喻,.
很天馬行空的比喻,.
來去講述一些具體以色列人要反省的一些可能是很沈重的包袱的問題,.
例如他們一直亡國的原因,.
要他們有一種痛定思痛,.
也有些關乎到預言將來的,.
尤其是第九到第十四章的預言,.
關乎到末後,.
我們提過就是到那一天出現很多很多次,.
十三章也是大規模地出現,.
好像指向一個日子,.
是遙遠的將來,.
但是我們不能準確知道什麼時候的遙遠,.
這個就是我們看先知書,.
看預言的先知書,.
必須要銘記,.
我們不可能要對號入座的,.
因為到底是哪一個時代的那一天,.
是沒有告訴我們後世的人知道的,.
我們只能夠透過預言和講述來掌握,.
十一章上一次,前一次提到,.
撒迦利亞要扮演一些戲劇的角色,.
他要扮演一個忠心的牧羊人,.
見到那些賣羊的和買羊的,.
好像一個牧場的舞台劇,.
那些賣的不顧惜那些羊,.
買的又以為自己是怎樣來有勢力等等,.
就變成了那個中間人,撒迦利亞,.
後來他又將那些不稱職的牧羊人解僱,.
但最終他還是得不到百姓的喜歡,.
所以他只能夠有三十塊錢的公價,.
所以當主耶穌看回這些經文的時候,.
主前四百多年,撒迦利亞書,.

$^{81}$你會感受到主耶穌看到自己將來,.
即是他成長的時候,.
他就是那個好牧人,.
他的下場就是被人出賣了,.
被猶大自己的門徒出賣了,.
三十塊錢,很少的價值,.
來到他出賣了耶穌,.
使得主耶穌要被交在這些祭司長,.
法利賽人的手中,.
所以很實物教材,很戲劇性,.
撒迦利亞就扮演著這個好牧人的角色,.
被出賣,被百姓所憎惡的一個結局,.
後來他又扮演另一個角色,.
就是不好的牧人,.
那個描述就是,.
那個不好就會帶來更大的破壞,.
就是群羊沒有人醫治,.
沒有牧者來幫助他們,.
來餵養他們,.
就變成沒用的牧人,.
丟棄了羊群,.
羊群就失去了那個依靠等等,.
這些其實對於撒迦利亞來說,.
他其實自己都不是很明白的,.
為什麼要這樣扮演,.
神為什麼要安排我做一個.
這樣的愚昧牧者的角色呢?.
我們看到第十一章,十五節,.
我們之前上一次提過的,.
但是去到十二章,.
仍然都有牧羊人這個身份的責備,.
十三章就再明顯一點,.
所以我們看看第十二章,.
上一次我們都有提到,.
有兩種的勝利,.
第一種是戰爭的勝利,.
以色列不知什麼時候,.
它非常強大,.
周圍的列國都被它打敗了,.
但是如果用第十節,.

$^{121}$去到十四,.
甚至去到十三章的第一節,.
那個分段你就會看到,.
那種的勝利就不是硬件,.
不是要征服,.
不是要去侵略,.
而是裡面的那個潔淨,.
所謂裡面不是以色列的內部,.
是人心靈的內部的那個潔淨,.
即是人人願意將自己,.
放在神的洗罪泉源裡面,.
來得到那個活水,.
這樣的勝利,.
就是神使得百姓的心靈改變,.
所以這一章聖經又是很戲劇性的,.
這些都是在新約有引述的經文,.
十二章第一節開始,.
關乎到默示,.
英文都有些翻譯,.
不是一個vision,.
有些是vision的,.
即是有一個視野,.
有一個啟示,.
但是很多數在舊約的,.
關乎到默示,.
是burden,.
burden即是重擔的字,.
即是關乎到以色列的重擔,.
即是得到默示的才知道,.
得到預言的才知道,.
他看這個訊息是很沈重的,.
所以叫burden的意思,.
即是要釋放一個訊息,.
是神要你說,.
其實某程度上如果那個訊息不是一個皆大歡喜,.
是關乎到一些的,.
可能是審判,可能是刑罰,.
說的人就有一種重擔的感覺,.
這個大家要留意,.
十二章第一節,.

$^{161}$耶和華論以色列的默示,.
鋪張朱天,建立地基,.
做人裡面之靈的耶和華說,.
我必使耶路撒冷被圍困的時候,.
向四圍列國的民,.
成為令人昏醉的杯,.
這默示也論到猶大,.
那日我必使耶路撒冷向聚集攻擊他的萬民,.
當作一塊重的石頭,.
凡是舉起的必受重傷,.
耶和華說,.
到那日我必使一切馬匹驚惶,.
使騎馬的癲狂,.
我必看顧猶大家,.
使到列國的一切馬匹黑眼,.
當然用現代的,.
就是一切的軍事,.
在神的時間,.
在耶路撒冷,.
都完全失去效力,.
好像沒有電,.
或者完全被黑客入侵,.
電腦的程式壞了,.
完全都插不著,.
類似這樣,.
如果在古代,.
他們沒有這些先進的飛機大炮,.
他們能夠想到的就是,.
在戰場那些馬都瞎了,.
到處亂跑,.
上來攻擊耶路撒冷的,.
他們全部好像喝醉了,.
你可以想像到,.
其實沒有可能,.
在戰爭的時候是最精銳的,.
但是去到有一個時候,.
神要使到列國攻擊,.
是完全失效,.
失靈,.
這些其實我們都知道,.

$^{201}$稍為有些黑客入侵,.
現在所有的東西都是高度的,.
中央的電腦化,.
只要癱瘓了某些東西,.
其實什麼都失靈了,.
這個是我們都能夠明白理解的,.
在這裡告訴我們,.
一切的馬匹乞懶,.
猶大的族長必心裡說,.
耶路撒冷的居民,.
倚靠萬軍之耶和華他們的神,.
就作我們的能力,.
那日,出現很多次,.
我必使猶大的族長,.
如火盤在木柴中,.
猶如火把在禍菌裡,.
他們必左右消滅四圍列國的民,.
耶路撒冷的人,.
必凝住本處,.
就是耶路撒冷,.
耶和華必先拯救猶大的帳篷,.
免得大衛家的榮耀,.
和耶路撒冷居民的榮耀勝過猶大,.
那日,耶和華必保護耶路撒冷的居民,.
他們中間軟弱的必如大衛,.
大衛的家必如神,.
如行在他們前面,.
知耶和華的使者,.
那日,我必定義滅絕來攻擊耶路撒冷各國的民,.
你可以想像到,.
那個保護,那個德性,那個榮耀,.
在那天,我們是不知道哪一天,.
什麼時期,我們都不知道,.
但是有一個預言,.
就是神定義要釋放,.
拯救祂自己的選民,.
這個絕對不是所謂偏心,.
所謂助紂為虐,絕對不是的,.
因為神有祂自己的工作,.
在這個民族的心靈裡面,.

$^{241}$所以第十節,.
有些人就說兩段經文好像不是很協調,.
有些不是和上下文有拉上關係,.
其實正正就是那個關係,.
我們表面看就是神要出手,.
使耶路撒冷,以色列,猶大,.
得到那個應該得到的榮耀,.
因為他們強大起來了,.
列國都不可以攻擊他們,.
但是現在不是說只是單方面給他們一切武器,.
而是最重要,.
反而是第十節,.
一直去到十三章的第一節的經文,.
這裡很特別的,.
這次又是提以色列的國家國護,.
不過這次提他們不是有什麼強大的武力,.
四圍的保護,.
或者是有神怎樣來令列國,.
他們的軍備全部瓦解,.
完全不是這樣,.
而是內部有一個泉源,.
有一個能夠解決他們內心問題的泉源,.
這裡真的很寶貴,.
如果我們要背今句,.
這些真是今句了,.
整段都是今句,.
「我必將那斯恩叫人懇求的靈,.
囂冠大維加和耶路撒冷的居民,.
他們必仰望我,.
就是他們所擇的」,.
這個擇字已經是一個很形象化的比喻,.
好像來象徵著這位耶和華,.
這位的牧者被擇,.
即是受傷,.
這個被擇令你想起誰呢?.
很明顯就是主耶穌被羅馬的兵丁.
用長矛去刺入他的胸膛,.
看看他死了沒有,.
流出血和水,.
約翰福音就是引這段,.

$^{281}$他們要仰望自己所擇的人,.
我們如果還沒有讀完這段經文之前,.
和大家揭一揭約翰福音十九章,.
就是耶穌離世斷氣的情景,.
所以很奇妙,.
撒迦利亞書這些舊約的經文很少人讀,.
但是新約的作者,.
福音書的作者,.
是引很多撒迦利亞書的,.
十九章,.
可能我們讀多一點經文,.
由第十七節開始,.
約翰福音十九章十七節,.
他們就把耶穌帶了去,.
耶穌背著自己的十字架出來,.
到了一個地方名叫獨留地,.
那個獨字不是夫,.
獨留,獨白那個獨,.
希伯來話就叫各各他,.
他們就在那裡釘他在十字架上,.
還有兩個人,.
當然是那兩個強盜,.
一同釘著,.
一邊一個,.
耶穌在中間,.
彼拉多又用牌子寫了一個名號,.
安在十字架上,.
寫的就是叫猶太人的王,.
拿撒勒人耶穌,.
有許多猶太人念這名號,.
因為耶穌被釘十字架的地方,.
與城相近,.
並且是用希伯來,.
羅馬,羅馬就是拉丁文,.
希利尼,就是希臘文,.
三個文字寫,.
猶太人的王,.
那個罪名,.
變相就是,.
一讀經文的時候,.

$^{321}$就以為是那個,.
開會,.
就以為是他的身份,.
其實是罪名,.
猶太人的王的罪名,.
猶太人的祭司長對彼拉多說,.
不要寫猶太人的王,.
要寫他自己說,.
我是猶太人的王,.
要讓路過的人看到,.
不要誤會了,.
他是猶太人的王,.
是他自己說自己是猶太人的王,.
我們是不承認的,.
他們要這樣拗文擲字,.
其實那個人已經釘在十字架上,.
他們都要這樣來據理力爭,.
一來就是,.
那些羅馬兵兵,.
管你那麼多,.
都不認識你誰來的,.
只是當其中一個死囚,.
當耶穌和那兩個強盜一樣,.
是罪大惡極,.
或者是私通外地,.
叛國等等,.
所以就要將他置於死地,.
但是二十二節,.
彼拉多說,.
我所寫的我已經寫上了,.
繼續留意,.
兵兵既然將耶穌釘在十字架上,.
就拿他的衣服分為四份,.
每兵一份,.
又拿他的女衣,.
這件女衣原來沒有縫而,.
是上下一片織成的,.
即是織出來的,.
他們就彼此說,.
我們不要撕開,.

$^{361}$只要拈溝,.
即是抽籤,.
看誰得著,.
這要應嚴經常的話說,.
他們分了我的外衣,.
為我的女衣,.
即是內衣,.
來拈溝,.
這個詩篇二十二篇,.
彌賽亞的仇虎詩,.
二十三篇我們很熟,.
耶和華是我們的牧者,.
但二十二篇是仇虎的僕人的詩篇,.
所以意味著作者寫得很隱晦,.
其實外衣被人拆了,.
女衣也被人拆了,.
即是赤裸裸的耶穌,.
暴露在所有的一個大街大巷的小山丘上,.
即是人們走過,.
見到這個罪名,.
猶太人的王,.
但見到這個被人打得血淋淋,.
又完全是赤裸裸的耶穌,.
是極大的羞辱,.
對於這位萬王之王的一個很隱晦的描述,.
即是大家自己可以聯想,.
是多羞恥,.
是多麼痛苦,.
是多麼值得那些人,.
即是不認識他的人,.
來詛咒他,.
即是朱般的咒罵,.
在那一刻,.
那一段時間,.
好像那裡特別黑暗,.
所以到十二點鐘,.
耶穌釘十字架的時候,.
整個天都黑了,.
你想想周圍的人,.
所有不認識耶穌的人,.

$^{401}$誤會耶穌的人,.
都盡量將他的咒罵,.
來咒這個人,.
即是他一個人擔當我們眾人的罪,.
我們就不會滅亡,.
很多人是這樣想的,.
即是避免羅馬人借機會,.
來再欺壓我們,.
即是說將這個本來賣國的,.
本來出賣我們的,.
現在找到了,.
賜死他了,.
那我們就沒事了,.
是大祭司該亞發這樣想的,.
你們其實不知道了,.
讓我們來說吧,.
之前上文有提到的,.
好,接著下面,.
釘釘果然就作了這樣的事,.
又是應驗聖經寓言,.
詩篇二十二篇的經文,.
應該是十八節,.
第二十五節,.
站在耶穌十字架旁邊,.
與他母親的姐妹,.
並格羅巴的妻子瑪利亞,.
和莫大拉的瑪利亞,.
即是在十字架下面,.
其實就是最大嫌疑,.
居然間門徒,.
他們懼怕離開,.
但是忠心的婦女,.
好像沒什麼可以輸的了,.
陪伴這位老師,.
這位主,.
最後一程,.
所以女性對信仰的堅忍,.
其實真的,.
你看看,.
二千年前都已經有了,.

$^{441}$你都可以想像得到,.
雖然婦女未必對很多真理,.
她能夠去辯論,.
可以去思辯,.
但是對於忠誠,.
對於她堅持的信仰,.
那個力度,.
有時比男性,.
比一些所謂神學家,.
更加強,.
所以往往很多的姐妹,.
支撐起教會很多的聖公,.
有些默默地做,.
都不需要讓人知道,.
不需要求明,.
不需要人知道她做過,.
你可以想像得到,.
這些姐妹就是這樣,.
在耶穌的腳下,.
當然更加有她的媽媽,.
第二十六節,.
耶穌見母親和她所愛的那門徒站在旁邊,.
當然歷史,.
或者教會的歷史都說是約翰,.
就是這位作者,.
這個漁夫約翰,.
不是施洗約翰,.
施洗約翰已經死了很久,.
就是這位約翰,.
站在瑪利亞的旁邊,.
跟他的媽媽說,.
母親看你的兒子,.
又對那門徒說,.
看你的母親,.
從此那門徒就接他到自己家裡去了,.
好像約翰就說回這件事,.
他好像是有一個很明白耶穌這個說話的托父,.
將自己的母親交託給他的門徒約翰,.
是主耶穌所愛的門徒,.
這個交託,.

$^{481}$他也記一筆,.
在今天,.
在現在他寫這本書的時候,.
他的媽媽如果還沒有死,.
就一直是約翰去供養,.
代替耶穌在肉身的身份,.
去供養這位媽媽瑪利亞,.
接著留意,.
第28節,.
「這是以後耶穌知道各樣的事已經成了,.
為要洗到經上的話應驗就說,.
我學了有一個器皿盛滿了醋,.
放在那裡,.
他們就拿海蛹摘滿了醋,.
綁在牛膝草上,.
送到他口,.
這個又是詩篇的另外一篇的壽苦詩,.
69篇的壽苦詩,.
就是講到彌賽亞他的壽苦,.
69篇其實是大衛的苦,.
接著第30節,.
「耶穌嘗了那醋就說成了,.
便低下頭將靈魂交付神了.」.
就是說他已經沒有辦法再忍受,.
他就已經死了,.
留意,.
約翰福音沒有提到天地昏暗,.
沒有提到那兩個強盜來和耶穌對話,.
那些是盧伽福音的記載,.
約翰福音特別強調,.
他是君王的身份的描述,.
接著第31節,.
「猶太人因這一日是預備日,.
似乎是一個特別的安息日,.
因為他們預節第二天就已經是安息日了,.
由因那安息日是一個大日,.
所以可能有一個雙倍的安息日,.
所以他們必須要強制安息,.
就求彼拉多叫人打斷他們的腿,.
把他們拿去,.

$^{521}$原因就是如果他們那些強盜未死,.
安息日就要馬上拔下來,.
安息日不可以掛在木架上,.
所以要快些收工,.
年三十晚要快些走,.
即是打養,冬至四點半要收工,.
是這個意思,.
不知道他們死了沒有,.
要求屍體,.
所以快的就是打斷他們的腿,.
呼吸不到就死了,.
就是這個意思,.
就求彼拉多叫人打斷他們的腿,.
把他們拿去,.
免得屍首,.
當安息日留在十字架上,.
無論是申命記十六章,.
那些死了在十架上的人,.
或者掛在木頭上的人,.
是受咒詛,.
但是有律法教導,.
就是當晚要拿下來,.
不要過安息日,.
即是對遺體的尊重,.
即是你時中遊街,.
都要給尊嚴屍體,.
所以某程度他們真是懼怕,.
如果過了就沒有問題,.
但是如果過了安息日還沒有死,.
是很不吉利,.
這個是他們的麻煩之處,.
羅馬人就沒有這些問題,.
你喜歡掛多少日就多少日,.
有些記載說,.
起碼五六日都不死,.
不斷肝煎,.
那幾眼釘不會令到他完全透不出氣,.
你想想人是不喝水三日四日才死,.
你想想這樣的煎熬,.
每一個透氣還要給他一個腳踏,.

$^{561}$他可以踩著腳踏,.
提高透氣,.
你想想每一次就痛到入心入肺,.
即是這樣的折磨,.
一個死囚,.
應該是槍斃或者注射毒藥就立刻死,.
是無痛,.
但是羅馬人就想出這樣的懲戒,.
還要每一個顯著的公路位置釘十字架,.
你可以想像到的確有殺一儆百的作用,.
你可以想像所謂羅馬的承平,.
是靠殘酷的統治令到其他城邦臣服,.
所以當時當時猶大也是亡國,.
所以如果要有這麼嚴重的,.
差不多是干犯國家的罪,.
就要這樣釘十字架死刑,.
所以他們兩個強盜未死,.
所以就打斷他們兩個的骨頭,.
不能呼吸,不能提高透氣,.
那就斷氣死了,.
即是好死,不用再煎熬,.
所以你會看到這個描述是很真實的,.
三十二節,.
於是丁丁來把頭一個人的腿,.
並與耶穌同釘第二個人的腿,.
即是在他左右兩邊的強盜,.
都打斷,即是未斷氣,.
打斷他們的腳,.
那就一定很快就馬上死了,.
「以此事來到耶穌那裡,.
見他已經死了,.
就不打斷他的腿,.
唯有一個丁丁拿槍,.
砸他的肋旁.」.
就是這個字,砸他的肋旁,.
即是刺他的肋旁,.
不是輕輕用手術刀,.
找對位置這樣切下去,.
是隨便用一根矛插進去,.
你可以想像到,.

$^{601}$沒錯,耶穌已經死了,.
是不會有肉身的痛楚,.
因為已經斷了氣,.
他的靈魂都已經離開他的身體,.
你可以想像到,.
然後有一個很客觀的,.
好像幾科學,.
如果你是懂醫學的常識,.
「隨即有血和水流出來.」.
即是已經分開了,.
即是心臟完全停頓,.
那個功能已經完全沒有了,.
即是很證實耶穌真的斷了氣,.
死了,這樣來描述,.
所以漢見這件事的那個人,.
就作見證,.
他的見證都是真的,.
並且他知道自己所說的是真的,.
叫你們也可以相信,.
「一事成了,惟要應驗經常的話說,.
他的骨頭一根也不可折斷.」.
沒錯,是如意折羔羊,.
世世代代都記住,.
吃的如意折羔羊,羊羔,.
骨頭是一根都不可以折斷的,.
即是如果跛的,.
上來獻祭,.
走路扭傷,.
就不可以獻祭了,.
就變成殘疾了,.
所以每一隻是驗過的,.
骨頭不可以折斷的,.
奇妙就是居然連這件事.
都要應驗在耶穌的身上,.
他的骨頭,你想想是千萬分之一,.
真的有機會是未死,.
那些丁丁安息日,.
那些祭司長求彼拉多.
要打斷那些死囚的腿,.
讓他們不要呼吸,.

$^{641}$你想想打斷很容易很快,.
但是來到耶穌那裡,.
免得煩,.
都死了不打了,.
就應驗了舊約的出埃及記十二章,.
兩千多年前的一個預言,.
其實那個都不是預言,.
只不過你們吃這些如意折羊羔,.
要吃得貴重一點,.
吃得刁鑽一點,.
要吃哪些,.
就好像有些人吃菜膽鮑魚,.
不要吃那些黃牙白外面那些很粗糙的,.
要吃裡面那些很細微的,.
這樣真的很刁鑽,.
要這樣,.
耶穌的腳不打斷,.
完全是因為當時的情勢,.
但是原來應驗的是出埃及記十二章的那個.
關乎到以色列人守逾越節的,.
所以耶穌在新約,.
你知不知道祂叫什麼稱呼?.
保羅說,.
我們逾越節的羔羊耶穌,.
這樣稱呼耶穌,.
保羅這樣叫,.
所以新約的人是看得很明白,.
十字架的死和耶穌的關係,.
所以你想想,.
只是約翰福音這一章,.
都已經引述了幾段詩篇的經文,.
出埃及的經文,.
還有這裡告訴我們,.
其實還有詩篇的三十四篇的經文,.
就是這樣的,.
接著第三十七節,.
這裡我讀了這麼多,.
才是在說我們撒迦利亞書這一節,.
經上又有一句說,.
他們要仰望自己所擇的人,.

$^{681}$其實是那些丁丁擇,.
擇是很聽的一個很委婉的說話,.
其實是用那支長矛,.
他們在十字架下面,.
耶穌是掛在十字架上面,.
那個十字架都未必是這麼低,.
所以他們用那支長矛這樣大力插進去,.
插入耶穌的胸膛裡面,.
這樣來證實耶穌是否真正死了,.
不是只是休克,昏迷了,.
而是要證實祂是否死了,.
所以這個擇,.
原來這個撒迦利亞書十二章十節,.
我們可以回來看回這一段經文,.
撒迦利亞書十二章,.
這裡就很不同了,.
因為他們要仰望自己所擇的人,.
他們如果是從十二章的第十節,.
到十三章的第一節,.
就是在說存已息烈的百姓,.
上文就說沒有任何一個烈國,.
在那天,就是神命定的那天,.
可以來攻擊耶路撒冷,.
因為神會保護祂自己的子民,.
他們會很強大,.
他們會很得勝,.
所有烈國的武裝都會被廢棄,.
那些馬自己都盲了,.
所有的軍事都失效,.
但是你看看第十節這一段,.
我們還未正式開始讀,.
就已經感覺到很奇妙,.
因為這個是真正人心靈的迴轉,.
仰望一位受傷的墓者,.
第十節,.
「我必將那斯恩叫人懇求的靈,.
囂冠大衛加和耶路撒冷的居民.」.
不就是上文說的大衛加,.
以色列加,猶大家,.
其實overall所有的以色列人的意思,.

$^{721}$「就是他們所摘的,.
必為我悲哀如喪毒生子,.
又為我愁苦如喪長子,.
那日耶路撒冷必有大大的悲哀,.
如米格多平原之哈達臨門的悲哀.」.
引述這個典故,.
聖經沒有出現過,.
到底是什麼事件,.
是什麼大的天災,.
還是地震,.
還是一些戰爭的傷亡,.
沒有什麼記載就不知道,.
但是對猶太人來說,.
他們會很深刻,.
我們一想起汶川大地震,.
唐山大地震,.
228,.
是不是228,.
雞寮的山泥傾瀉,.
如果我們想起這些溫代的年代,.
62年的颱風,.
我們老香港一代就會記憶猶新,.
他們一想起就是,.
這個悲哀,.
不是普通的無病呻吟,.
是真真正正全國的悲哀,.
全國悲哀,.
原來悲哀在神眼中是很厲害的,.
會帶來一個很大的復興,.
第十二節,.
一直描述的悲哀,.
境內一家一家都必悲哀,.
大衛家男的獨在一處,.
女的獨在一處,.
拿丹家,.
拿丹就是先知的家族,.
男的獨在一處,.
女的獨在一處,.
利美家,.
祭司家族,.

$^{761}$男的獨在一處,.
女的獨在一處,.
侍媒家,.
為什麼要提侍媒,.
那個是辯雁問之派,.
曾經咒罵大衛王的那個呢,.
侍媒家,.
男的獨在一處,.
女的獨在一處,.
其餘的國家,.
男的獨在一處,.
女的獨在一處,.
為什麼要這樣整天說獨處,獨處,獨處呢,.
即是好像表明著,.
無論是男的,.
對外,.
無論是工作,.
或者是一家之主,.
或者對內,.
或者是在社會上,.
未必有這麼大的社會地位的,.
總之男男女女,.
都來到去感受到,.
這一份的私恩,.
懇求的靈,.
帶來的那個的感動,.
因為你想一下,.
很正常的,.
譬如考試季節,.
爸爸回家,.
媽媽就不回家了,.
媽媽要陪著兒子考試,.
看他溫書,.
很正常的,.
男的聽完課後,.
很有感動,.
回家後,.
媽媽沒有感動,.
我都沒有聽,.
我聽回錄音帶,.

$^{801}$都未必這麼感動,.
男的獨在一處,.
女的就不是獨在一處,.
她可能是在看電視,.
在看韓劇,.
但現在這裡說,.
不斷說男的女的,.
都感受到這個私恩懇求的靈,.
那個的內心的激動,.
那個的大家都不理對方,.
大家互相都一齊懇求,.
你想一下,.
這個是全世界的大復興,.
這個是神的私恩懇求的靈,.
令到人很自發地向神禱告,.
所以其實,.
星期三晚祈禱會,.
大家多回來安靜,.
在神面前暫時放下一些燭霧,.
其實是一個很好的訓練,.
讓我們感受到神同在,.
至少你的意識,.
你要認為這個地方是神同在,.
當然不是說神不可以在大家樂,.
他不可以在牛奶公司,.
不是牛奶冰室,.
不是這個意思,.
他是無處不在的,.
但是當有一班人很認真,.
為神的國度去祈禱,.
這件事本身是嚴肅的,.
這件事本身就是屬靈,.
這件事本身就是很驚人的,.
因為大家都願意很安靜地等候神,.
所以其實這些情景是我們要去孕育,.
去鼓勵弟兄姊妹的,.
是呀,.
這個社會太忙碌了,.
正正應該說太忙碌了,.
我整個星期七日的晚上都這麼忙碌,.

$^{841}$我應該迫我自己要預留一個晚上,.
一個小時來到禮拜堂,.
和弟兄姊妹一同安靜等候,.
你越是忙的,.
你越是高級的,.
你越是真是日理萬機的,.
如果你能夠將自己擺在禮拜堂那裡,.
那一個小時,.
其實就憤外感受到,.
這種獨在一處,獨在一處,.
說了很多次,.
所以聖經不是在說什麼,.
什麼獨在一處,.
接著十三章就有了,.
又是那個那日,.
上下文都出現了很多那日,.
所以不是那些軍事的強大的德性,.
是神的私恩叫人懇求神的靈,.
懇求那種悲哀,.
你想像得到,.
這個悲哀帶來的就是那個潔淨,.
所以十三章就是說,.
那一日必給大衛家和耶路撒冷的居民,.
開一個傳言,.
是什麼傳言呢?.
是洗除罪惡與污穢的傳言,.
你可以想像得到,.
這個就是那個那日,.
所以是平衡的,.
上文一至九節,.
就是為什麼耶路撒冷這麼厲害的,.
全世界的攻擊都可以完全打不低他,.
所以列國都要來失敗呢?.
我們不知道這一日到底是怎樣發生,.
怎樣出現,.
現在不是了,.
但是我們知道的就是,.
我們要祈禱的就是,.
十節去到十三章的一節,.
是神已經給了一個很重要的指示,.

$^{881}$世界的列國教會,.
在教會裡面要來呼籲,.
我們要向神懇求,.
求這一位施恩的神,.
讓我們可以去向他懇求的靈,.
是那個靈的心靈的蘇醒,.
自然來到去囂觀我們的時候,.
我們就會來到去有一個向這一位受傷的牧者,.
仰望這一位被責的這一位主,.
這一位主肯定就是在說耶和華,.
舊約的時候,.
撒迦利亞的時候,.
是在說耶和華就是那個受傷的主,.
我們新約,.
我們是後世,.
再是四千年了,.
如果計算這個時候,.
二千年的耶穌的救恩的時代,.
這一位被責的就是主耶穌了,.
仰望主耶穌,.
所以羅馬書十章九節,.
我們一仰望就可以得救,.
因為當人的心表示相信就可以得救,.
口裡面承認就可以清義了,.
所以就是一定要仰望,.
所以報道會,.
他覺得很好聽,.
覺得很好笑,.
覺得很感動,.
全部是對的,.
那個是一個準備,.
但是最重要就是,.
或者奇異恩典,.
令到他感動,.
全部是準備他仰望這一位被責的人,.
這個是最重要的,.
他如果不打開心靈,.
他如果不離開之前,.
他沒有那個心靈的感動,.
真正他自己,.

$^{921}$我應該怎樣祈禱,.
你教我祈禱,.
我向這一位主仰望,.
他一仰望就可以得救,.
尤其是如果他明白,.
那個被責的人和他有關,.
這個就很重要了,.
所以施恩懇求的靈,.
囂冠當時的眾生,.
所以這個就是教會復興一個很重要的關鍵,.
已經告訴我們了,.
你可以想像得到,.
不是無病新淫,.
為悲哀而悲哀,.
是真誠地自己感受到神的同在,.
神的愛,.
來向神懇求,.
求祂的恩典,.
求祂的安慰和憐憫,.
這些不能造作的,.
這些只能夠我們一路去等候,.
一路去勉勵弟兄姊妹,.
去看到神的恩典,.
所以這段經文是真正的金句的金句,.
弟兄姊妹,.
我們要時時都要看一看,.
真的我們要祈禱,.
教會的復興,.
當然人數多不代表不復興,.
肯定是一個指標,.
肯定是一個量度,.
但是真正的復興,.
是有一班人感受到神的愛,.
神的恩典,.
神的威嚴等等,.
來向祂懇求,.
和為其他人懇求,.
教會需要有代禱者,.
教會需要有認識祈禱的人,.
去為教會禱告,.

$^{961}$這些是很重要的,.
教會不缺乏一些活動,.
很多很多的項目,.
忙到不得了,.
但是真正大家懇求神的靈的聚會,.
或者這樣的氣氛的氛圍,.
其實更加來得重要,.
不是說那些活動不重要,.
而是這些才真正可以去到神的心靈裡面,.
也透過這樣去到我們自己的心靈裡面,.
你看看舊約的描述,.
是多麼反映,.
大家根本不再介意有沒有活動,.
介意的是神怎樣在我的內心,.
使我得到那個復興,.
得到那個祂的恩惠,.
所以這個洗罪的泉源,.
是神要為我們而開,.
又揭曉《約翰福音》第六章,.
耶穌就是說這些經文,.
或者應該是第七章,.
那個是鑄盤節的時候,.
第七章的十四節,.
《約翰福音》,.
七章,.
或者由第十節開始看,.
耶穌過這個鑄盤節,.
但他弟兄上去以後,他也上去過節,.
不是明去,似乎是暗去,.
正在節期,猶太人尋找耶穌說,.
他在哪裡呢?眾人為他紛紛議論,.
就說他是一個好人,.
有些人就說不是,.
他是迷惑眾人的,.
但是沒有人明明的去講論他,.
因為怕猶太人去到節期的時候,.
耶穌上殿裡面去教訓人,.
猶太人就稀奇說,.
這個人都沒有學過怎樣能夠明白書呢?.
所謂書就是當時的卷軸,.

$^{1001}$應該是摩西五經那些,.
耶穌說,我的教訓不是我自己的,.
乃是那差我來的,.
人如果立志遵著他的旨意行,.
就必曉得這教訓,.
或是出於神,.
或是我憑著自己說的,.
輾轉一路講到摩西,.
講到那些真道,.
然後到第37節,.
去到節期的末日,.
即是主棚節的第八日,.
那天又是一個聖日,.
就是最大之日,.
耶穌站著高聲說,.
人若學了,可以到我這裡來學,.
信我的人,就如經上所說,.
從他腹中要流出活水的江河來,.
這個可以說是二創亞書那段經文,.
58章,.
耶穌這話是指著信他之人要受聖靈說的,.
那時,即是耶穌的時代,.
還沒有賜下聖靈的,.
因為耶穌尚未得著榮耀,.
未得著榮耀,即是說他未離開這個世界,.
未死,.
所以保衛司就未能夠代表耶穌成為另外一個的忠報,.
另外一個的保衛司來到地上,.
所以是有一個先後次序的,.
耶穌要得了榮耀,即是說他要離世升天,.
保衛司聖靈才會被差遣來到世間,.
所以當時代的耶穌的門徒,.
他們就介乎兩個鹹淡水,.
所以他們是有雙重的受聖靈的,.
我們就不是了,.
我們後世的人一信耶穌,.
就會有聖靈在我們裡面內住,.
因為聖靈是由始徒行傳的時候,.
五旬節就普遍降臨在地上,.
不需要再第二次降臨,.

$^{1041}$那次是一個很特殊的,.
他們能夠見到好像一些百甲,.
好像一些火焰,.
這樣降落他們的頭上,.
但是我們就不需要這樣做,.
這裡就已經預示了那個聖靈的全面的工作,.
舊約就告訴我們,.
施恩懇求的靈,.
就是什麼時候我們仰望受傷的耶和華,.
被壓的受傷的聖者,.
我們就可以得到神的施恩,.
這個靈在我們的內心,.
而且還有,.
再開一個洗罪的傳言,.
信信肯定就是主耶穌的寶血,.
可以洗乾淨我們的罪,.
這個是一個很雙重的意義,.
既有聖靈,.
同時間有那個赦罪的平安,.
所以我們信耶穌,.
你應該要一信耶穌就感受到,.
你內心裡面有聖靈帶來的生命,.
帶來的悔改,.
帶來的生命平安,.
這是一個信道位的門徒的特徵,.
所以我們報道會,.
我們陪談,.
我們不要很開心,.
你就先審慎地收貨,.
因為你不知道他有沒有懇求聖靈,.
他有沒有仰望這位被壓的耶穌,.
你要教清楚他,.
你不是信那個喜樂,.
你不是信那個說完的,.
很好笑的笑話,.
很感動的故事,.
你要信的是那一位被壓的耶穌,.
就是你的主,.
我們陪談完就要有這樣的眼光,.
所以不要這麼快開心,.

$^{1081}$他走出去,.
去到大家樂就已經忘記了,.
第二天就已經忘記了那些興奮,.
但是如果他沒有忘記,.
他真的仰望這位耶穌,.
和聖靈已經住在他的內心,.
那就很耐心了,.
他慢慢就會有一種裡面的感動,.
他會尋求真理,.
當然陪談完最好再做多一步,.
可能過了一個星期再打電話邀請他,.
上個星期怎樣,.
你再加固這一種的感動,.
和對那個信心的堅持,.
那個信者就容易再得到那個建立,.
如果不是很快太多的續務,.
他就會忘記了他的那個決志,.
所以這個其實是很重要,.
所以回來撒迦利亞書十二章,.
很奇妙就是,.
沒錯,神會透過他的作為,.
使到以色列蒙保守,.
去到那一天,將來的不知道哪一個時候,.
神要在全地工作,.
整個以色列蒙保守,.
但是有更寶貴的就是,.
真正的復興不是靠刀槍,.
是施恩懇求的靈,.
在百姓當中來工作,.
男的女的家家戶戶,.
無論你是什麼支派,貴族,.
皇室的大維加,利美加,.
即是祭司家族,.
你是拿丹仙姿家族,.
你是侍妹被咒詛的家族,.
侍妹是被咒詛,.
這樣即是說,.
無論你是住半山區,.
還是住瘦茂坪,.
或者你是住雞寮,.

$^{1121}$或者你是住在藍田,.
或者你住在沙頭角等等,.
你住在那些屋村,.
男的獨在一處,.
女的獨在一處,.
你可以想像得到,.
你看那些娛樂版,.
哪個明星又嫁給哪個富三代,.
過百億的,.
這些看來都沒有用,.
又不關你的事,.
報紙就是有這麼多賣這麼多,.
讓那些人仇富,羨慕,.
快點嫁給有錢人,.
將來如果這些人,.
男的獨在一處,.
女的獨在一處,.
半山區不愁衣食,.
完全不會懼怕,.
只是怕每天怎樣去玩,.
想辦法怎樣玩,.
不需要愁衣食,.
他都獨在一處,.
女的獨在一處,.
在很窮乏的,住tong 房的,.
他又內心又有感動,.
又獨在一處,.
這個就是將來的情景,.
聖經的工作,.
施恩懇求的靈,.
所以我們要有這個眼界,.
看到神的工作,.
是要令到人心進入改變,.
這個是很寶貴,.
很難得,.
十三章的第一節,.
會為大衛加和耶路撒冷的一切居民,.
開一個傳言,.
所以這個預言,.
新約已經開始了,.

$^{1161}$主耶穌告訴我們,.
當他升天得到榮耀,.
就可以將這個聖靈賜下,.
使我們能夠有活水江河在我們的內心,.
所以我們需要有這個懇求,.
第十三章第二節,.
萬軍之耶和華說,.
那一日,.
我必從地上除滅偶像的名,.
所以我再說一次,.
十二章的第一到第九節,.
是軍事上以色列的強大,.
但是第十到十三章第一節,.
是下半段的以色列真正的復興,.
是說心靈歸向神的真正的大復興,.
什麼時候會見到呢?.
是不是我們已經在天堂才見到呢?.
這個就沒有人知道了,.
但是我們在地上的教會,.
應該就是每一個時代的牧者,.
都要為這件事懇求,.
每一個家族都要這樣的時候,.
就很不同了,.
雖然不應該對號入座,.
那些指名道姓,.
譬如某些家族,.
譬如阮微家族,.
很多人的,.
又有Allen,Joyce,.
有其他的人,.
他們獨在一處,.
女的獨在一處,.
小玲家族,.
很多,.
小芬,小恩,.
還有Eric,等等,.
又獨在一處,.
女的獨在一處,.
這樣對號入座,.
神也是這樣對號入座,.

$^{1201}$那丹家族,.
利美家族,.
大衛家族,.
侍媒家族,.
這樣來到去,.
不論你是什麼貴族也好,.
被人嘲笑咒詛的侍媒家族,.
都這樣一起享受這個.
施恩懇求的靈,.
所以第十到十三章一節,.
就是一個復興的,.
很重要的指示,.
撒迦利亞收到這個信息,.
他應該是非常之振奮,.
因為他就是參與在他那個時代的懇求,.
記住,.
這個時候是撒迦利亞的大理烏王第二年,.
九月的時候,.
他這個信息,.
所以是剛剛在興建工程的時候,.
那個聖殿,.
這個爛尾的聖殿,.
現在再復修,.
所以這個時間是一個挺興奮的時候,.
但是以色列真正的榮耀,.
是每一個家族,.
每一個的支派,.
都來到去向神有一個懇求,.
十三章第二節,.
就不是一些好東西,.
就開始說那些當時的情景,.
和那些牧者的問題,.
萬君之耶和華說,.
那一日,.
我必從地上除滅偶像的名,.
不再被人紀念,.
也必駛到這地,.
不再有假先知與污穢的靈,.
若再有人說預言,.
即是說那些文迷,.

$^{1241}$或者是膠鬼,.
或者是占卜等等,.
或者是拜那些偶像巴力等等,.
如果有人再說預言,.
生他的父母必對他說,.
你不得存活,.
因為你託耶和華的名說假的預言,.
生他的父母,.
在他說預言的時候,.
要將他刺頭,.
那一日,.
又再出現這些描述,.
凡是作先知預言的,.
說預言的,.
必因他所論的異象來收窺,.
不再穿毛衣哄騙人,.
你看到非洲的巫術,.
都是這樣穿著毛衣,.
可能舊約,.
你看到以利亞,.
那些巴力先知都是這樣的,.
似乎他好像披帶著某些特殊的裝扮,.
來去,.
好像現在那些靈體附我身,.
我開始說預言,.
類似這樣的情景,.
所以有時那些鬼氣,.
某程度上都好像反映到,.
這些古代的迷信的風俗,.
所以這裡都幾傳神,.
穿著毛衣去哄騙那些人,.
因為你正正是怕這些巫術,.
他就,.
現在我有預言,.
那些人去問迷,.
就是講什麼都信,.
因為你已經整個心理,.
有了這一種先入為主,.
所以這裡告訴我們,.
將來不會再有這些虛假的預言,.

$^{1281}$教鬼的這些巫術,.
和這些所謂的異象,.
他們一定會蒙羞的,.
第五節,.
第十三章的第五節,.
他必說,.
我不是先知,.
我是耕地的,.
我從幼年作人的奴僕,.
必有人問他說,.
你兩臂中間是什麼傷呢?.
他必回答說,.
這是我在親友家中所受的傷,.
其實不知道什麼意思,.
突然之間有句這樣的說話,.
如果我們放在上文,.
就是那些假先知,.
來講那些預言,.
正常就是他們會用刀刺自己,.
巴尼先知就是這樣的,.
以利亞就取笑他們,.
都沒有什麼,.
你刺大力一點,.
他們真的這樣刺,.
但是都沒有用,.
後來以利亞一祈禱,.
就天火焚下來,.
燒了那些水溝的祭物,.
所以特別神在那一刻,.
顯靈去幫以利亞,.
以利亞和這八百個巴尼先知,.
在加密山上鬥法,.
其實就是羞辱那些假先知,.
其實他們不是假的,.
他們是巴尼的先知,.
所謂假先知,.
以色列叫假先知,.
這裡就好像描述那些人,.
又說自己是先知,.
又來講預言,.

$^{1321}$但是當有壓力,.
又說我不是先知,.
我是耕種,.
我是做一個普通的人,.
然後有人就戳穿他,.
你兩手臂中間是什麼傷,.
其實是交軌的時候,.
有些人問米,.
神打,.
都會戳自己,.
我小時候中學的時候,.
我跟同學去看人玩神打,.
那些叫彪彤,.
他真的,.
我們是三個我這麼大隻的男孩,.
捉住他,.
他上了身,.
捉住他,.
是捉不到,.
還有他的眼睛發光,.
我們其實很害怕,.
他講到明他還沒上身之前,.
就叫我們無論如何,.
如果他失控,.
我們就用大力地斬他的頸,.
讓他醒過來,.
真的很大力地斬下去,.
他才沒事,.
原來當時是什麼書悟空,.
齊天大聖上他身,.
真的好像喝醉酒一樣在動,.
你想想交軌都有這樣的東西,.
所以你可以說這個咒座,.
是一個刑罰,.
但是去到那天,.
一切與魔鬼,.
與神的真理無關的,.
都要來被排除,.
有一個這樣的復興,.
一切的交軌都被人發現,.

$^{1361}$他的父母都跟他割蓆,.
要將他刺透,.
有些在舊約的經文有教導,.
如果你發現你家鄰居,.
他回錦宣,.
他在交軌,.
你要揭發他,.
要告訴教會,.
《新靈記》有教,.
不然你一定會牽連,.
你就可以想像到,.
將來就會這樣,.
對這些交軌,.
或者講假的預言,.
用託神的名,.
來講預言,.
即是說這些人,.
其實每一個時代都有,.
又要講預言的榮耀,.
但又不要講預言之後,.
自己要為所講的負責任,.
好像我們比較俗的,.
戴口罩又要戴頭盔,.
這些先知,.
將來就不會再有了,.
因為會被揭發,.
會被刑罰,.
會被排擠,.
你可以想像得到,.
又是給這個時代,.
即是撒迦利亞的時代,.
百姓有一個反思,.
其實過去他們的失敗,.
就是沒有聽真先知的警告,.
沒有講口婆心的預言,.
聽的是一些賊系先知,.
在皇室,.
皇室的先知就一定是,.
你想聽什麼我就講什麼,.
西帝家說應不應該打仗,.

$^{1401}$周圍的先知都說,.
一定可以打贏的,.
其中一個叫米該雅,.
皇帝不喜歡她,.
那些大臣說不要,.
聽一下,.
你聽了這麼多個,.
米該雅就被召來了,.
在監獄裡面,.
未召來之前,.
那些大臣就跟她說,.
你懂得了,.
那些皇室的先知已經講了,.
米該雅說,.
神要我講什麼我就講什麼,.
去到皇宮的時候,.
皇帝就問她,.
其實應該是阿哈王,.
不是西帝家,.
西帝家是另外耶利米,.
這個米該雅就來到她說,.
那你講吧,.
米該雅說,.
我見到以色列人沒事了,.
回到自己的國土,.
很平安,.
然後那個阿哈王就在罵她,.
平時你不斷講咒詛的說話,.
現在你講這麼好聽的,.
米該雅說,.
不是,.
我見到的就是那些百姓能夠回來,.
是因為那個皇帝陣亡,.
阿哈聽了之後,.
氣得爆,.
然後就罵她,.
然後就馬上叫周圍的人掌嘴,.
這個是聶王紀上有記載的,.
聶王紀上應該有21章記載,.
真的才知道,.

$^{1441}$百姓一般都不聽的,.
這個是以色列一個很沈重的包袱,.
所以去到將來,.
所以這個那一天又是什麼時候,.
又是沒有人知道,.
神都沒有為我們揭曉,.
但是真正的虛假,.
真正的假冒神的名字的寓言,.
和這些交軌的,.
巫術的,.
這些這樣的謊話,.
全部都會消失,.
神要來興起他自己的真先知,.
這些又是非常之寶貴,.
所以一個的時代,.
如果有些很重要的講解聖經的牧者,.
聖經學者,.
傳道人,.
是一個時代的很大的福氣,.
所以很盼望為多些不同的傳道者祈禱,.
能夠他們領受神的話語,.
然後重要就是忠心地,.
無畏懼地去講解出來,.
因為一定是有這個落差,.
這就是這一段經文的描述,.
所以第五第六節就比較合理了,.
就是這些講完假預言的假先知,.
就隱藏了,.
想置身於事外,.
就說我不是先知,.
我只不過是耕地做農夫,.
我一直在這個民族都是循規蹈矩,.
大概意思,.
然後人們就揭發他,.
你的兩臂中間為什麼會受傷,.
一個又是很委婉,.
好像俄景橋那些阿姨,.
打小人的,.
為什麼你的手特別粗,.
又起疹,.

$^{1481}$你平時有什麼的,.
其實他是打小人的,.
大概是這個意思,.
就是說他所謂兩臂中間受傷,.
這些膠軌,.
這些是儀式,.
長期下來可能有些疤痕,.
有些疤痕等等,.
還有那些位置其實膠軌的人就會知道,.
你就可以想像到這樣的情況,.
將來到那一天是不會再出現,.
就是說歷史歷代我們都要面對.
真的先知假先知,.
我們要有一個聖經說的神思明辨,.
我們去過士徒行傳的地方,.
去過哥林多,.
去過帖撒羅尼迦加,.
我們沒有去到比利亞,.
比利亞那裡的人比帖撒羅尼迦加人更加優秀,.
不是他們特別厲害,.
而是他們聽到保羅講道,.
雖然你是保羅,.
你是神的使徒,.
我們都信任,.
但是你說的道,.
我們要對照一下舊約有沒有錯,.
所以他們是要將保羅的講道,.
來對照一下聖經,.
是不是真是聖經的一貫真理,.
所以他們信得更加堅固,.
這個就是比利亞人,.
他們是神思明辨,.
所以都是我們求主讓我們學習,.
今天我們很難看到一個人的額頭是假才知道,.
是他兩臂有受傷,.
有些疤痕,.
膠軌,.
我們不知道,.
但是盼望我們都能夠我們的胃口來崇拜,.
聽道,.

$^{1521}$不只是聽笑話,.
覺得很悶就閉上眼,.
塞起耳朵,.
這樣就很難的,.
其實要學的,.
不是每個人都很容易就能夠,.
你喜歡好聽的就用心聽,.
覺得不好聽,很悶的,.
你都用心去聽,.
那些悶的你吹走它,.
它都有些道,.
它都有一兩句是我有幫助,.
這樣聽道,.
你就真是所向無敵了,.
真是要這樣學的,.
真是要這樣去,.
每個牧者都有他,.
總有一兩句,.
他是在神那裡領受的,.
我有時說笑,.
高明軒牧師說,.
我載他回去出元朗,.
港元道出去元朗,.
他跟我說,.
曾牧師我有聽你查經班的,.
我說這個是全香港最大的笑話,.
你聽我查經班,.
我差很遠,.
我只是讀一下經,.
他說不是,.
我都有領受,.
我不知道他是說笑還是說真,.
我想都是真的,.
但是我們要留心,.
當然好像我現在說笑,.
其實不是的,.
將來是人都分辨到的時候,.
再沒有假才知道,.
但是今天我們很容易就我們自己,.
有我們自己的判斷,.

$^{1561}$誰說道今晚不要聽了,.
明天才來,.
明天是誰,.
不是,.
明天是李思敬牧師,.
院長,.
明天幾大都要準備好來聽他講道,.
有時是你那個準備過程,.
你自己已經復興了,.
人家都未講你已經復興了,.
或者那個普通的港元,.
又不可以起名,.
自己同工當然不要起名,.
曾銳貴,.
都預了是這樣講的,.
不用怎麼聽的,.
你未坐下你自己都已經不能復興了,.
因為你根本都不希望神給你有說話,.
其實我們就不要這樣,.
將來到那天,.
人人都看見哪些是假才知道,.
但是現在神沒有給我們這個恩典和恩賜,.
現在我們就什麼人我們都要聽,.
但是重要就是我們聽到他裡面的道,.
或者他裡面的那個對神的愛,.
這樣才叫聽道的人,.
這樣才是我們學習到,.
因為很難的,.
因為一年阿里斯堅尼兩次,.
高明軒尼可能三次,.
你一年只有四次,五次聽道,.
這樣不夠的,.
這樣沒有可能的,.
所以你應該是不同的牧者,.
都有他的有限的領受,.
他有他自己的不同的領會,.
他的演繹,.
你不要那麼介意他的演繹,.
你介意的就是他裡面,.
認真地準備好,.

$^{1601}$來領受了,.
可能是一兩句,.
有可能是他一句,.
有可能是一個笑話的故事,.
幫到你,.
令你產生一個共鳴,.
一個反思,.
這些就是對於假先知的判斷,.
雖然這個是預言將來的,.
但是對我們今天都是有一定的啟發,.
和鼓舞,.
我們可以肯定,.
神會給我們每一個時代,.
都有他的說話,.
都有他的牧者,.
來釋放神的道,.
所以我們都和牧者說,.
我們不是在講台分享給你知,.
分享就是你可以聽,可以不聽,.
其實不應該這樣理解,.
我們既然站在台上,.
我們是代表神去證道,.
見證那個道,.
這樣的意思,.
所以講道和分享是兩樣東西,.
分享就是請你不要站在台上,.
你站在下面,.
這樣叫分享,.
雖然我們不應該這樣拗文嚼字,.
但是嚴格來說是真的,.
不然講台為什麼要放在中間呢?.
福音派會這麼緊張,.
講台要放在中間呢?.
就是我們都真的認為,.
神有道給我們這個時代去聽,.
藉著他的僕人,.
他的侍女去傳道,.
這才是重點,.
不然講口才,.
講有叻的智慧的人,.

$^{1641}$大把,.
不需要我們牧者,.
這樣你就會明白了,.
到第七節了,.
我們看完這一段,.
我們可能下個禮拜就完結了,.
整個撒迦利亞書,.
最後的四章的馬拉基書,.
也都完了整個歸回的歷史書和先知書,.
大家的投票就是,.
帖撒羅尼迦前後書和末世論,.
有關聯的一併去講,.
我們又有機會講一些系統神學的課題,.
末世論是一個很大的題目,.
第七節了,.
萬軍之耶和華說,.
刀劍馬應當興起,.
攻擊我的牧人和我的同伴,.
擊打牧人,羊就分散,.
出現過的,.
馬太福音,.
馬漢福音,.
就是這段經文,.
我必反手加在微小者的身上,.
耶和華說,.
這全地的人三分之二必剪除而死,.
三分之一仍必存留,.
我要使到這三分之一驚火,.
熬煉他們如熬煉銀紙,.
好像是那個揀選的問題,.
這個揀選在這裡就不是很強,.
但是在整個舊約的先知書,.
尤其是以賽亞書,.
那個愚民,.
成愚之民,.
Chosen People,.
就很強的意識,.
這裡又好像告訴我們,.
有其中的三分之一,.
是經過熬煉的,.

$^{1681}$試煉他們如試煉金紙,.
他們必求告我的名,.
我必應允他們,.
我要說,.
這是我的子民,.
他們也要說,.
耶和華是我們的神,.
又是一脈相承,.
又是一個求告,.
他們有意識地知道自己是神的選民,.
神是他們的神,.
這個又是將來的預言,.
這裡有一個短暫在耶穌身上應驗,.
又是再一次說那個牧者,.
好牧者,.
被擊打楊勳山,.
我們站一站去馬太,.
第26章的馬太福音,.
主耶穌說的經文,.
只是撒迦利亞書已經出現三番四次,.
在福音書裡面的預言,.
是應驗,.
應驗在主耶穌,.
他這位彌賽亞身上,.
所以撒迦利亞書雖然很難讀,.
但是很多關於彌賽亞的事,.
馬太福音第26章,.
不看這麼多了,前一點,.
31節這一段,.
或者我們看上文的設立聖餐這一段,.
26節,馬太福音的26章,.
他們吃的時候,.
耶穌拿起餅來祝福,.
就張開第七門徒說,.
你們拿著吃,這是我的身體,.
又拿起杯來祝謝了,.
第七他們說,你們都喝這個,.
因為這是我納藥的血,.
為多人流出來,.
洗到醉得到赦免,.

$^{1721}$但我告訴你們,.
從今以後我不再喝這葡萄汁,.
因為第二天就要死了,.
直到我在我父的國裡,.
和你們喝新的日子,.
新筵席,將來新天新地,.
一個永恆的大筵席,.
我們又不需要用物質,.
我們不燃開千萬席,.
我們不要用物質的概念,.
因為天堂是靈界,.
不能用物質的概念,.
但一定有一個新的,.
古往今來的所有的靈體,.
聖徒,我們應該有一個靈體,.
我個人認為,.
就會在這個新天新地裡面,.
肯定有這個羔羊的筵席,.
你可以說這個聖餐,.
將來的聖餐,.
主耶穌做主席,去主領了,.
接著就說那個實在,.
第30節,.
他們唱了詩,.
就由橄欖山走去葛國他,.
當然走進去耶路撒冷,.
第31節,.
那時耶穌對他們說,.
今夜你們為我的緣故都要跌倒,.
留意一下,.
因為經上記著說,.
我要擊打牧人,羊就分散了,.
撒迦利亞書13章7節,.
是耶穌自己引的,.
所以耶穌對撒迦利亞書,.
是非常非常熟悉的,.
隨便可以引述,.
去到那個時刻的情景,.
他知道那些群羊會分散,.
原因就是擒賊先擒王,.

$^{1761}$耶穌給那些祭司長,.
和這個猶大,.
出賣耶穌的猶大,.
在這個赫西瑪利園,.
當耶穌行完那個橄欖山,.
吸輪溪,.
就走去赫西瑪利園,.
就擊打牧人,.
羊分散,.
是說他自己,.
就是撒迦利亞書這一段經文的應驗,.
所以再一次提那個好牧人被擊打,.
好牧人被砸,.
好牧人被出賣,.
哇,11,12,13章,.
都是說這個好牧人,.
30塊錢的公價被出賣,.
百姓是不認同他所做,.
就是給個良心當狗吠,.
那些百姓,.
自己情願被人出賣,.
都不介意,.
就是用羊養大之後要出欄,.
那些賣的其實就是那些外邦人,.
賣的就是那些祭司長,.
就是那些富戶,.
但是那個好牧人就被人真言了,.
用30塊錢就打發他走了,.
然後這裡就告訴我們,.
主耶穌被擊打,.
他讀到這些經文,.
他引述這些經文的時候,.
相信撒迦利亞的形象,.
一直說的那個好牧人,.
怎樣被擊打,.
怎樣被砸,.
都在主耶穌的內心,.
這個時候他正在去喀西馬利縣,.
他還沒有上十字架,.
但是他都知道自己終有一天,.

$^{1801}$會被一些很殘忍的兵丁,.
用一些長矛這樣刺他,.
刺死他,.
不是刺死他,.
他已經死了,.
這樣刺他的胸膛,.
流出血和水,.
這些情景,.
這些描述,.
對於這個好牧人,.
當時舊約一定是在說耶和華,.
自己這樣被擊打,.
他的羊分散,.
但是這位牧者將來會焦醉,.
他自己的百姓,.
會經過那個火煉,.
會得到那個重生,.
最終就是你是我的神,.
我是你的子民,.
所以我們其實有生之年,.
弟兄姊妹我們一定要,.
對哪一樣東西才是最重要,.
傳到永恆的冠冕才是最重要,.
地上的很快沒有人會記得你,.
你家裡有多少個Hermes袋,.
那些名人,.
劉鑾雄的太太,.
你家裡有多少,.
整個櫃打開,.
誰會記得你,.
當你再老一點,.
你住在病房,.
住在老人院,.
誰會看你那些Hermes袋,.
你拿來在老人院周轉走,.
誰會來看你,.
可能一個半個,.
有不同的場合,.
就已經足夠了,.
哪些是傳到永恆呢,.

$^{1841}$再一次告訴我們,.
到那一天,.
我要說,.
這是我的子民,.
神這樣稱呼,.
這個是我的子民,.
神這樣稱呼你,.
跟著他們也要說,.
耶和華就是我們的神,.
所以一些敬虔的長者,.
你看到他在老人院,.
我認識一位,.
我們以前的舞會的一個師母,.
九十幾歲,.
她在老人院,.
在湖平老人院,.
很得到整個院捨的人的愛戴,.
很低調,.
很謙卑,.
她還在做詩琴,.
在老人院的聚會,.
詩琴,.
人一見到她,.
就會感受到那種平安,.
見到老人家,.
你感受到一份屬天的平安,.
你又知道,.
這不是操練,修練,.
而是她日積月累,.
對自己身份的認識,.
耶和華你是我的神,.
我是你的子民,.
人生去到一些你已經沒有什麼活動空間能力的時候,.
這是最大最榮耀的宣言,.
宣告,.
耶和華你是我的神,.
所以求主讓我們在有生之年,.
現在是很活動,.
很多你可以服侍的時間,.
有時候我們未必那麼在意這個身份,.

$^{1881}$你很在意別人怎麼看你,.
看你做得好不好,.
一說做得不好,.
你就說很不開心,.
很激氣,.
浪費力氣,.
不需要,.
我們事奉神,.
神喜悅就OK,.
何必介意,.
不同人有時候有不同的角度看事物,.
不要那麼介意別人怎麼看,.
盡力做好自己就行,.
但是這件事不是開玩笑,.
去到斷氣之前,.
你都要認定最寶貴最榮耀就是,.
耶和華你是我的神,.
我是你的子民,.
這是最重要的,.
這就是到哪一天,.
神的兒女,.
以色列百姓,.
他們就需要有這樣的生命改變,.
當然原因就是神的熬年,.
所以這是一生的,.
我們每一個弟兄姊妹,.
其實未去到見主,.
都是一個熬年的階段,.
是試煉的階段,.
所以我們都需要不停有不同的學習,.
很簡單,.
小小分享,.
完了我們這兩張聖經,.
其實我們昨天,.
祈禱會都有簡單報告過,.
接下來的幾堂崇拜,.
都會認真報告給大家知道,.
堂董會又決定,.
因應現在紅恩堂有一個很大的需要,.
魯木斯突然間的離開,.

$^{1921}$又激發我自己在過去,.
由一月到十二月,.
我參與在紅恩堂的服事幫忙,.
我自己都有很大的負擔,.
所以我都和教會的領袖商討,.
我很想去紅恩堂幫忙,.
所以明年一月開始,.
我會全時間在紅恩堂幫忙,.
做全副恩的工作,.
但我不會減少,.
星期四晚我仍然會回來帶茶經班,.
所以大家繼續為我祈禱,.
我希望也有一些弟兄姊妹,.
如果有感動,.
可以用不同的方式去紅恩堂幫忙,.
一起去全副恩,.
那裡基層的地方的需要都很大,.
有可能有不同的方式去全副恩,.
所以盼望大家都為教會,.
為我祈禱,.
麗成木斯她會留守在錦宣,.
她會延後她明年六月的退休,.
所以我就可以有一個無後顧之憂,.
所以我會全力在星期二至星期日,.
我會在紅恩來服侍,.
盼望大家彼此禱告,.
為教會現在這個轉變禱告,.
更加其實應該為魯木斯的家人,.
他們現在的哀傷,.
或者紅恩堂的弟兄姊妹,.
他們突然間失去了一位忠心的牧者,.
都需要來禱告,.
魯木斯的安息禮拜,.
就是在十一日,.
最尾那天的黃昏時間,.
在萬國殯儀館,.
大家如果有時間,.
都一起出去對牧者最後的一個致敬,.
我們一起禱告..
再一次多謝天父給我們看了十二,十三章,.

$^{1961}$很重要我們看到,.
雖然以色列將來有一個很大的德性,.
而且他們的軍事,.
他們的勇力帶來周圍的列國都要懼怕,.
甚至乎都沒有一個國家可以攻打,.
但是我們看到有更寶貴的就是那個懇求的靈,.
賜下給以色列,.
使得整個國家,.
整個神的選民,.
都可以來去仰望這位被責的主,.
使得他們得到心靈的復興,.
和有一個洗罪的傳言,.
來給以色列可以洗除罪惡,.
我們更加看到十三章所描述的,.
這一位被擊打的牧人,.
真是這樣一個人承受神的罪惡,.
使得流散的羊可以得到歸回,.
可以得到澡煉和嗜煉,.
我們盼望我們的人生都能夠真是有這樣的機會去學習,.
更加教會都是在神的煉證裡面,.
成為明亮的燈台去照耀,.
保守我們能夠,.
不單止我們追求的是一些活動,.
而是我們真是在神的裡面,.
不斷的有神聖靈的感動,.
教會是一間追求神國度的教會,.
我們認定我們是屬神的,.
我們是神的子民,.
我們這樣仰望禱告,.
奉耶穌基督的名,.
阿們..
大師兄 這不是我自己做的嗎.
我自己做的 沒有人說是我自己做的.
我自己做的 沒有人要說誰做的.
你幫幫我.
我自己做的 沒有人要說誰做的.
\newpage



\section{路加福音 2:8-20}
\label{sec:pZkchuW18Ro}
\textbf{2023-12-24 加恩樓三語聯合崇拜《Good News of Great Joy》路加福音 2\_8-20 講員:Rev. Carol Ilidan}
\newline
\newline
連結: \href{https://youtube.com/watch?v=pZkchuW18Ro}{\texttt{https://youtube.com/watch?v=pZkchuW18Ro}} ~~~~ 語音日期: 2023-12-27
\newline
\newline
\hyperref[sec:I2kkhTmf8Lk]{\small{< < < PREV SERMON < < <}}
~
\hyperref[sec:index_chronic]{\small{[返順時目]}}
~
\hyperref[sec:index_scriptual]{\small{[返順卷目]}}
~
\hyperref[sec:7Vdee0ksMqc]{\small{> > > NEXT SERMON > > >}}
\newline
\newline
路加福音 2:8-20
\newline
\begin{longtable}{cl}
\hline
\hline
章節 & 經文 (和合本修訂版)\\
\hline
2:8 & \begin{tabularx}{0.7\textwidth}{X} 在伯利恆的野外有牧羊人，夜間值班看守羊群。 \end{tabularx} \\ \\ \relax
2:9 & \begin{tabularx}{0.7\textwidth}{X} 有主的一個使者站在他們旁邊，主的榮光四面照著他們，牧羊人就很懼怕。 \end{tabularx} \\ \\ \relax
2:10 & \begin{tabularx}{0.7\textwidth}{X} 那天使對他們說：「不要懼怕！看哪！因為我報給你們大喜的信息，是關乎萬民的： \end{tabularx} \\ \\ \relax
2:11 & \begin{tabularx}{0.7\textwidth}{X} 因今天在大衛的城裡，為你們生了救主，就是主基督。 \end{tabularx} \\ \\ \relax
2:12 & \begin{tabularx}{0.7\textwidth}{X} 你們要看見一個嬰孩，包著布，臥在馬槽裡，那就是給你們的記號。」 \end{tabularx} \\ \\ \relax
2:13 & \begin{tabularx}{0.7\textwidth}{X} 忽然，有一大隊天兵同那天使讚美神說： \end{tabularx} \\ \\ \relax
2:14 & \begin{tabularx}{0.7\textwidth}{X} 「在至高之處榮耀歸與神！在地上平安歸與他所喜悅的人！」 \end{tabularx} \\ \\ \relax
2:15 & \begin{tabularx}{0.7\textwidth}{X} 眾天使離開他們，升天去了。牧羊人彼此說：「我們往伯利恆去，看看所成的事，就是主所告訴我們的。」 \end{tabularx} \\ \\ \relax
2:16 & \begin{tabularx}{0.7\textwidth}{X} 他們急忙去了，找到馬利亞和約瑟，還有那嬰孩臥在馬槽裡。 \end{tabularx} \\ \\ \relax
2:17 & \begin{tabularx}{0.7\textwidth}{X} 他們看見，就把天使論這孩子的話傳開了。 \end{tabularx} \\ \\ \relax
2:18 & \begin{tabularx}{0.7\textwidth}{X} 聽見的人都詫異牧羊人對他們所說的話。 \end{tabularx} \\ \\ \relax
2:19 & \begin{tabularx}{0.7\textwidth}{X} 馬利亞卻把這一切的事存在心裡，反覆思考。 \end{tabularx} \\ \\ \relax
2:20 & \begin{tabularx}{0.7\textwidth}{X} 牧羊人回去了，因所聽見所看見的一切事，正如天使向他們所說的，就歸榮耀於神，讚美他。 \end{tabularx} \\ \\
[1ex]
\hline
\hline
\end{longtable}
$^{1}$謝謝.
早安.
大家早安.
我已經知道三種語言了.
無論如何.
今天我們要做的事情.
就是要觀察神的言語.
主題是.
「大喜」的好消息.
今天.
今天我們的題目是「喜樂」的好消息.
我好奇.
如果耶穌在女王 Elizabeth 的時期出生.
或者是英女皇 喬治布希.
或者是美國總統.
而不是海薩的皇帝.
會有什麼不同呢?.
我經常在想.
如果耶穌在女王 Elizabeth 的時期出生.
或者是英女皇 喬治布希.
或者是美國總統.
而不是海薩的皇帝.
會有什麼不同呢?.
有一件事可以肯定.
耶穌的出生.
也不會成為頭條新聞.
耶穌的出生.
也不會成為頭條新聞.
新聞報道.
是關於世界的事件.
而我們的主人降生.
是在黑暗中.
新聞報道.
當然是世界性的事件.
但是我們的救主降生.
其實是默默無聞.
留意這三個特點.
第八節意外的驚喜.
第十二節有不尋常的徵兆.
第十七節有不尋常的徵兆.

$^{41}$第十八節有不尋常的徵兆.
第十七節有不尋常的徵兆.
第十八節有不尋常的徵兆.
第十七節有不尋常的徵兆.
第十八節有不尋常的徵兆.
第十七節有不尋常的徵兆.
上帝就是揀選一些愚蠢的,去顯出他們的聰明.
回想去年的聖誕節,你們當中可能也有些人感到驚喜,為何自己今天會在這裡.
有沒有想過神會帶你進入他們的家.
有些人可能仍然處於黑暗中,他們做的事情很少會想起上帝的榮耀.
神很想闖進我們的生命.
這個聖誕節對你來說可能是一個很巨大的驚喜,因為你第一次明白這些意味著什麼.
第二點是不尋常的徵兆.
.
伯利恆可能是一個小鎮,但要找到一個嬰兒也不容易,所以牧羊人要得到一些記號才可以找到那個嬰兒在馬槽裡.
.
所以那個記號就是一個在馬槽裡的嬰兒.
.
這個嬰兒的第一張床就在動物的馬槽裡.
.
這個馬槽不單只是一個記號,它本身在新約裡也很有意思.
.
馬槽的意義是什麼呢?.
.
它顯示了耶穌基督願意在外面毀身.
.
在哥林多後書八章九節說.
.
經文說因為這位嬰兒願意取奴僕的樣式,降卑而致顯出上帝的榮耀.
.
所以這個馬槽的意思就是耶穌的人生是由一個憂患的人生開始.
.
它也標誌了他在地上做事的方式.
.
他以救主的身份在馬槽出世,也以救主的身份在十格上死亡.
.
馬槽的意思也意味著牧羊人可以很自由地踏到耶穌的面前.
.
有一位學者曾經說過,我們可能在到達寶座前會很害怕,但我們到達馬槽前不需要害怕.
.

$^{81}$第十七對十八節是一個難以置信的故事.
.
牧羊人就像第一批基督教傳教士一樣.
.
他們見過基督,就告訴別人基督在他們身上發生的事.
.
耶穌的好消息是很值得分享的,而且最適合分享的人就是有這些經驗的普通人.
.
你可能不知道所有事情,但當你知道神闖進你的生命時,你就有些事情可以分享.
.
我們可以跟別人分享馬槽裡的寶寶,十架上的基督,也是一個將會回歸的君王.
.
這首詩歌就是《O Little Town of Bethlehem》.
.
這首詩歌就是講到伯利恆的聖旨降臨到我們身上,我們去祈禱,去趕走我們的罪,不要讓罪進入我們身上.
今天這個嬰兒誕生在我們身上,我們聽到天使敘述這個偉大的好消息.
我們來到這裡,我們跟主在一起,以馬來尼.
.
願我們都像牧羊人那種舊的方式,但我們傳的是新的好消息.
今天我們的救主誕生了,願神賜福給大家,聖誕節快樂.
\newpage



\section{}
\label{sec:7Vdee0ksMqc}
\textbf{2023-12-24 平安夜佈道會《不一樣的聖誕》 講員:李炳光牧師}
\newline
\newline
連結: \href{https://youtube.com/watch?v=7Vdee0ksMqc}{\texttt{https://youtube.com/watch?v=7Vdee0ksMqc}} ~~~~ 語音日期: 2023-12-27
\newline
\newline
\hyperref[sec:pZkchuW18Ro]{\small{< < < PREV SERMON < < <}}
~
\hyperref[sec:index_chronic]{\small{[返順時目]}}
~
\hyperref[sec:index_scriptual]{\small{[返順卷目]}}
~
\hyperref[sec:nEkocKnhlmk]{\small{> > > NEXT SERMON > > >}}
\newline
\newline
$^{1}$各位朋友 各位親愛的兄弟姐妹.
非常難得在平安夜的晚上能夠見到大家.
今天有人跟我說 整個香港靜了.
很多食肆都有空位.
怎麼以為不是滿了.
原來有人說.
今天有44萬的香港人.
離開了香港.
去享受他們聖誕的快樂.
吃喝玩樂是人生的一個享受.
今晚在座的每一位.
能夠放下.
很多吃喝玩樂的消遣.
來到這裡跟我們一起敬拜天父.
並且參加這個報道會.
有人說平安夜開報道會.
有沒有搞錯.
今晚應該要開心一點玩一下.
為什麼要坐在這裡聽別人說話呢.
但是我又覺得.
其實平安夜.
其實是一個很好的報道時間.
因為.
今天我們要見證一個人的誕生.
剛才Kathy說到.
他快要死了.
但是.
其實每個人都經歷過一個很特別的時刻.
我們都有出生的時候.
他來到這個世界.
但是這個人來到這個世界很特別.
他幫助無數的人.
直到現在.
都是做拯救的工作.
這個人是誰呢.
我相信大家都知道.
這個就是我所信的主耶穌.
今晚.
讓我有機會來到這裡做見證.
周牧師打電話給我.

$^{41}$本來今晚.
我自己的教會.
有聚會的.
今晚有音樂晚會.
那麼.
我太太施母.
現在也在唱詩.
她不能夠來.
但是我覺得有機會.
在這個時間見證主.
和分享福音.
是一個福氣.
正如Kathy所說.
我們不知道日子有多少.
但是我們把握機會.
能夠見證一個很重要的訊息.
分享天父的恩典.
是很難得的.
大家知不知道人生有很多突發的事.
正如Kathy所說.
不知道為什麼有車會撞過來.
今晚.
可惜她介紹得太短.
就是.
有一位弟兄.
是在這貴堂的教友.
他今晚和我一起吃飯.
很難得的一個見證.
他真是所遭受的經歷.
好像約伯一樣.
但是現在他活潑的.
生活在我們當中.
而且是行動自如.
雖然他的家人離開了.
但是他沒有離開上帝.
稍後如果有機會見到他.
大家可以詳細問問他的經歷.
但是我不想代替了他說見證.
也不想加添Kathy的見證.
但是無論如何讓大家知道.

$^{81}$原來一件事的發生.
不是偶然的.
原來是有很多很多的教導給我們.
大家知不知道聖誕節有兩段經文.
是很不同的.
一段經文.
就是路加夫音大家很熟的.
在百里行的野地裡.
有牧羊的人.
按著羹字看守羊群.
大家都很清楚.
天使天君來到報喜.
不要懼怕.
今日在大衛的城裡.
為你生了救主.
就是主基督.
非常感人.
而且興奮.
歡樂的一個場面.
但是另外一段經文.
是完全相反的.
這段經文是什麼呢?.
是馬太夫音第二章.
第二章說到什麼呢?.
說到原來耶穌出生的時候.
三博士獻禮.
但是原來希律的心裡.
是有一些特別的計劃.
如果你找到這個救主.
告訴我.
幸好.
三個博士沒有聽他說.
找到之後.
獻禮之後.
走了.
但是你知不知道.
那一次的第一個聖誕節.
風聲鶴唳.
很多小朋友都被殺.
我覺得今晚.

$^{121}$我們都可以說是.
經歷一個這樣的場面.
我們今晚很舒服的.
坐這麼溫暖的.
一個這麼漂亮的禮拜堂.
剛才我來到的時候.
看見湖的時候.
我覺得你們真的有眼光.
找到這麼漂亮的地方.
來做禮拜堂.
這麼安靜的地方來做禮拜堂.
現在我們真的很舒服.
而且又聽到這麼動聽的詩歌.
而且我們有機會一起慶祝.
平安夜.
但是你知不知道.
現在在加沙的地方.
又是風聲鶴唳.
很多小朋友.
每十分鐘有一個死.
而且那個地方.
在這幾個月裡面.
死了二萬多人.
而且你看見這個世界.
充滿著憎恨.
充滿著仇殺.
充滿著各種的問題.
加上這麼冷的天氣.
甘肅突然間又地震.
你看一下.
好像哭哭流近兼逢夜雨.
真的樣樣不好的事情都出現.
真的令我們覺得很難過.
我們怎樣能夠高呼的.
普世歡鵬救主下降.
我們怎樣能夠歡呼的告訴別人.
恭祝你聖誕快樂.
是不是.
很難說.
在今年的裡面.

$^{161}$我們的心情是兩樣的.
好像聖經裡面一樣.
兩樣的.
這兩個場面是截然不同的.
第一幅圖畫.
是充滿著平安和盼望.
在黑暗當中行走的百姓.
可以見到大光.
第二幅圖畫.
就是兵荒馬亂.
有殺戮.
有逃亡.
雖然博士朝拜聖嬰.
但是他們不敢將信息.
帶起信告訴希律.
落荒而逃.
雖然有人說.
聖誕節的晚上.
如果你在太空上.
如果你是太空人的話.
你看回地球.
是全宇宙當中最美的星球.
信不信.
因為到處都開了燈.
到處都弄得這麼美的燈光.
地球真的很美.
但是.
第一個聖誕夜.
絕不是這樣.
地球的人.
雖然說慶祝聖誕.
但其實他只是慶祝一個節日.
一個節日.
不是慶祝救主的誕生.
我們互相慶祝聖誕快樂.
但是這種快樂.
只是表面的.
短暫的.
雖然聖誕節.
最早是慶祝一個人的誕生.

$^{201}$但是這個壽星宮.
直到現在.
仍然沒有人理會它.
好像你開派對.
那個壽星宮在.
但是你們只有歡樂.
只有交換禮物.
只有飲飲食食.
沒有人理會這個壽星宮.
這個壽星宮.
直到現在.
都是被忽略的.
大家都知道是誰.
各位.
它沒有禮物.
只有我們交換禮物.
它沒有生日的祝詞.
只有我們吃喝玩樂.
什麼都沒有去分.
難怪有些國家.
認為我們不應該去慶祝聖誕.
不應該叫做Merry Christmas.
應該改做什麼呢.
是Season's Greeting.
換句話來說.
今天我們過著一個.
沒有耶穌的聖誕.
一個沒有耶穌的聖誕.
不是真的聖誕.
是不是.
各位.
你心裡沒有耶穌呢.
如果你沒有耶穌的話.
你不是那個真的聖誕.
今天的聖誕.
更加不容易過的.
大家都知道.
因為世界各處都有.
不平安事件發生.
不是平安夜.

$^{241}$在世界另外一邊.
正是主戰爭.
百里行.
我去過很多次.
告訴大家.
我的經驗不是那麼理想.
我去百里行.
帶隊去百里行的時候.
我遇到三次炸彈.
而且進去的時候.
不能正當進去.
要在橫門那裡.
一條很小的門.
很小的路才進去.
帶多少個紀念品.
各位.
百里行.
今年宣佈什麼呢.
不舉行聖誕.
原來有一間主誕堂.
那個神父在那裡說.
今年我們取消聖誕.
因為我們的弟兄正在死亡.
他說一個沒有真正平安的聖誕.
是不是真的聖誕呢.
其實一個沒有真正喜樂的聖誕.
是不是真的聖誕呢.
各位.
我們只不過是想念我們自己.
忘記了這個世界.
是處在水深火熱的當中.
當年天使的報喜.
大家都很熟悉.
在至高之處.
榮耀歸於神.
在地上平安.
歸於祂所喜悅的人.
就是說.
聖誕節最重要是講平安.
剛才的見證是不平安的見證.

$^{281}$但是在主的帶領之下.
他得到平安.
但是今天這個世界上.
有很多人再沒有機會了.
各位.
我很欣賞有一個聖誕卡.
近來很少收到聖誕卡.
間中只有一兩張.
聖誕卡寫了幾個字.
是很精準的.
就是說聖誕的精神就是平安.
聖誕的喜樂就是希望.
聖誕的忠心就是至愛.
沒錯.
聖誕就應該這三樣東西.
就是聖誕的精神就是平安.
各位我們有沒有平安.
你說我沒有痛.
沒有頭痛.
睡一覺.
吃得睡得.
我平安.
不是的.
其實我想問一下.
你的心靈有沒有平安呢.
還有.
我們喜樂.
待會吃大餐.
但是這個喜樂.
是在乎我們有沒有希望.
今天這個世界上.
很多人活在沒有希望的當中.
你看看.
打開報紙.
其實.
失望灰心的消息.
比希望平安的消息多得多.
其實聖誕的忠心就是愛.
但是我們感不感受到愛呢.
各位今晚.

$^{321}$我們真的有幾種的心情.
我們有沒有這種的質素呢.
各位.
我的題目就是.
我們要過不一樣的聖誕.
這樣.
可能你說牧師.
過不一樣的聖誕.
是不是表示我們取消了慶祝呢.
是不是我們取消了一切的活動呢.
我們睡在家裡.
安安靜靜就算了.
不是的.
其實我們能不能說.
是聖誕快樂呢.
我們能不能夠大聲說普世歡騰呢.
各位.
今晚.
我就說一些很特別的.
配合今晚我的講題.
和大家分享.
我逢禮拜四.
都有查經班的.
這個查經班.
是開了40年了.
40年我在不同的教會當中.
我逢禮拜四晚上都有查經.
一方面逼自己研究聖經.
一方面和大家分享.
在今年的年底.
我的查經班.
每個禮拜都有幾十人.
而且在網上的直播.
看到的有幾百人.
我覺得在這個時候.
我決定不再教了.
因為我已經很大年紀了.
暫時不要說我幾歲.
大家都猜到的.
我只是說.

$^{361}$我事奉了主60年了.
但是我說我暫時不教了.
要停一停才行.
萬事都有一個結束的時候.
當然很多人是很捨不得.
我自己也捨不得.
但是我在這個查經班.
有一晚.
我特別放映了一個短片.
這個短片是介紹一首詩歌.
這首詩歌.
很聞名.
一會兒大家就知道是什麼詩歌了.
這首詩歌.
竟然有個很富裕的弟兄.
他竟然請了.
世界上第三位.
第三位頂級的男高音.
特別唱這首歌.
做了一個短片.
我那晚找了一個短片.
放給弟兄.
大家不記得我說什麼.
只是看了短片.
各位,很多人看了之後流眼淚.
而且.
有人說不如這樣吧.
用這首歌來開個報道會.
找一班師班來.
我不知道今晚大家唱得這麼好聽.
就找你們唱就好了.
我聽得這麼好聽.
就找一班師班來唱.
本來計劃著.
但是太過匆忙了.
不能夠邀請大家.
而且.
他們拍拍心口.
說今晚我來唱.
來唱的時候.

$^{401}$大家看這首歌.
大家先聽了.
我先介紹這首歌給大家知道.
這首歌的背景.
很感人.
而且在平安夜的時候.
讓大家知道.
怎樣獲得真正的平安.
現在我請.
幾位弟兄姊妹.
來唱這首歌給大家聽.
而且有字幕打出來.
大家留心聽.
這首歌.
我播出來的時候.
很多人流眼淚.
而且很感動.
但是.
不是容易.
人人都可以得到這種經驗.
現在我請這幾位.
我們轉來的弟兄姊妹.
跟大家分享.
有時想不到.
天上何必高崗.
有時也不想到.
天上何必高崗.
我愛我自己.
我為何在平淵.
我心念我愛我自己.
我為何在平淵.
我心念我愛我自己.
我為何在平淵.
我心念我愛我自己.
我為何在平淵.
在我心裡.
我為何在平淵.
我心念我愛我自己.
我為何在平淵.
在我心裡.

$^{441}$我為何在平淵.
我心念我愛我自己.
我為何在平淵.
我心念我愛我自己.
祝福我太陽.
是誰在遠方出現.
明天再看天晴再出現.
我心念我愛我自己.
多謝幾位弟兄姊妹.
分享一首令人家寧靜的詩歌.
特別在平安夜的晚上.
各位,為何在平安夜的晚上.
大家應該高興唱Christmas carols.
但為何我特別介紹這首歌.
我覺得這個世界很需要這首歌.
這首歌原本的英文名叫.
It is well with my soul.
即是你不用擔心.
我的show,我的靈魂是好的,安定的.
翻譯得很好.
它說我的心靈得到安寧.
各位,為何我選這首歌呢.
原來我聽這幾個高音唱的時候.
我特別查考作者的來龍去脈.
原來發現他生在1828年.
死於1888年.
他是一名律師.
而且是一名律師機構的合夥人.
特別是做地產買賣.
這位叫做 Horatio Spetford.
這個 Spetford 本來是一名很成功的商人.
但同樣他是一名很熱心的基督徒.
他與當年慕迪博士在報道的工作上有份.
他生了五個新生兒子.
大兒子是兒子.
但很可惜患了重病離開世界.
在一個非常傷心的心情裡面.
面對人生各種的問題.
本來在1873年.
他為了要令到內心有多少的安慰.

$^{481}$帶著四個女兒和太太.
預備去巴黎參加慕迪博士的報道會.
準備好買了飛機票.
但突然間他的公司說.
不行,你不在這裡就不行.
因為現在有些地產問題.
各樣的問題需要你在這裡解決.
他沒有辦法.
只好讓太太和四個新生兒子先去.
他以後再來參加他們.
但是就在那一年.
他送別了太太後犯他的事的時候.
過了九天.
在1873年11月15日.
他聽到一個噩耗傳來.
真是可說是心碎難形容.
原來他太太坐了一艘很小的遊船.
大概五千多噸.
大概二三百個房間.
大概有三百一十三人在船上.
誰知走到第九天.
突然間有另外一艘船攔腰撞下去.
整艘船沉了.
以為全部人都喪生.
但是經過了一段時間.
他聽到他的太太還存在.
現在在美國國會博物館裡面.
有他的太太和四個新生兒子的照片.
還有他太太寄來的電報.
只有兩個字.
以前不是隨便打電報的.
現在不是隨便打電話的.
電報裡面只有兩個字.
Save Alone.
這幾個字現在還在博物館裡面.
拍了出來放大.
很感動人.
原來他的太太沒有死.
只留下四個新生兒子就先離開了.
非常難過在這段時間的情況下.

$^{521}$但是他就作了這首詩.
這首詩在人生各種災難來臨的時候.
他竟然說你不用擔心.
It is well with my soul.
我的內心有平安.
而裡面的曲調是由著名的作曲家.
Philip Breese寫了.
成為一首直到現在最感人安慰人心的詩歌.
在青年性歌裡面有的.
不過現在很多人唱.
但是這首歌仍然成為一番名曲.
各位,十九世紀曾經將這首歌.
成為最感人的靈魂的詩歌.
我講到這個故事的時候.
讓大家都知道.
原來這位律師.
竟然在這個環境當中.
講出一個很重要的人生的經歷.
It is well with my soul.
就是說我的心靈得到安寧.
各位,剛才我說.
在聖誕節裡面.
天使的報誡音.
大家知道裡面有幾樣重要的質素.
就是不要懼怕.
安慰.
大喜的訊息是歡樂.
關乎萬民的是屬於所有人.
生了救主是來拯救的.
包著布臥在馬槽裡是一個窮人的福音.
這個就是記號.
這個就是希望的福音.
榮耀歸於上帝.
地上平安歸於人.
原來聖誕節應該有這幾樣質素.
但是現在我們過聖誕節.
有沒有這些質素在其中呢?.
各位,剛才我提到.
聖誕節的特質就是有平安,有喜樂,有至愛.
但是我們內心裡面.

$^{561}$有沒有這種質素呢?.
聖經裡面提到.
有什麼能夠使我們與上帝愛隔絕呢?.
難道是患難嗎?.
是各種困擾嗎?不安嗎?.
沒有什麼能夠使我們與上帝愛隔絕.
所以今天在這個聖誕夜,平安夜,聖善夜裡面.
我要宣佈一個很重要的訊息.
就是真正的平安夜.
就是你的內心有安寧的一種表現出來.
在任何環境當中.
It is well with my soul.
各位,我有一次和一個合唱團.
去到成都開一個報道會.
剛剛是完結了四川的大地震.
這間教會可以容納八百多人.
但是那一晚坐滿了.
我負責講話.
但在講之前.
師班唱詩.
我又叫師班打出歌詞.
誰知我聽師班唱的時候.
特別去到一首歌.
奇怪,我看到指揮流眼淚出來.
大家知道流眼淚不止眼淚,還有鼻涕.
我在台上,我真的很心急.
我怎樣幫助他呢?.
被紙巾又不接.
怎樣幫助他呢?.
我正在緊張的時候.
我發現師班都是個個都在發呆.
為什麼師班都忍不住呢?.
原來我在台下,很多人跟著一起唱.
而且很多人都是流著眼淚,抹著眼睛.
我在這時候覺得奇怪.
這首是什麼歌呢?.
原來我留心看著裡面的字.
真的很配合我們人生所遭受到的一切.
他說,在壓力之下,我仍然平靜.
在冷漠之中,我有關懷.

$^{601}$在絕望的當中,我有盼望.
在忙碌的當中,我能夠悠閒.
在失敗的當中,我仍然努力.
在受苦的當中,我仍然歡笑.
在困境的當中,我有出路.
在逼得的當中,我仍然屹立.
在死亡的當中,我有永遠的生命.
接著他的結論是什麼呢?.
他說,做主的兒女,何等的有福.
這首歌我相信師班都會唱.
做主的兒女,何等的有福.
原來壓力可以平靜.
冷漠可以有關懷.
在盼望當中,我仍然有希望.
在失望的當中,我有希望.
在忙碌的當中,我仍然悠閒.
在失敗的當中,我仍然努力.
在受苦的當中,我仍然歡笑.
在困境的當中,我有出路.
在逼得的當中,我仍然屹立.
各位朋友,這就是我們的福音.
原來耶穌基督來到世界上.
是給我們一種信念.
在任何環境當中,我仍然可以屹立.
可以面對.
剛才我和一位弟兄載我來.
我跟他說,近來我聽到很多不好聽的消息.
一個進了醫院,一個又跌倒.
一個就是,我特別為他祈禱.
他患了癌症.
我前晚開祈禱會的時候.
他在Zoom露面的時候.
帶著貓鬚.
我看到之後,我難過得不得了.
為什麼這個世界上有這麼多我不開心的事呢?.
特別是聖誕節的時候,為什麼會來到這樣呢?.
各位,剛才阮木說到.
其實原來人生有很多事情不是我們想像得那麼順利.
各位,我們今晚這麼平安地在這裡.
但是你知不知道離我們沒有多遠的地方.

$^{641}$或者甚至我們所認識的人當中.
現在是流淚,痛苦,失望,難過的當中.
你面對這個平安夜.
各位,讓我們都能夠相信基督教的信仰.
是讓我們可以獲得真正的平安,真正的喜樂.
兩受真正的愛.
各位,耶穌基督親自說.
你們不要憂愁.
你們相信我,應該是相信上帝.
因為在我裡面有平安.
我所指的平安,沒有人能夠阻止的.
你們心裡面不要憂愁,不要膽怯.
各位朋友,真正的平安夜.
就是在任何環境當中,仍然保持寧靜.
很難得,阮木竟然撞車了,還跟上帝說話.
這麼靜?這麼好?這麼快?.
但是,其實真真正正的人生.
是要面對這件事.
各位,我不是叫大家說.
我們在平安夜應該要遇到的事情.
但是我們有認識的人.
各位,讓我們都能夠知道.
原來基督教的信仰.
就是讓我們能夠在壓力之下有平靜.
在冷漠當中有關懷.
在盼望,在絕望當中有希望.
在忙碌的當中可以悠閒.
在失敗的當中可以努力.
在受苦的當中可以歡笑.
在困境當中有出路.
在迫不得已的當中可以屹立.
在死亡的當中,我們有永生的盼望.
各位親愛的朋友.
我盼望大家要過一個不一樣的聖誕.
剛才我說一個沒有平安的聖誕.
不是真的聖誕.
一個沒有喜樂的聖誕.
不是一個真的聖誕.
一個沒有主耶穌基督的聖誕.
不是一個真的聖誕,是不是?.

$^{681}$各位,我們是不是過一個沒有耶穌的聖誕?.
人云亦云,人食我食,人玩我玩.
這樣的聖誕,各位在座當中的朋友.
我剛才沒有問今晚有多少人是第一次來.
但是無論如何,我希望大家都能夠思想一下.
你的內心裡面有沒有真正的平安?.
你的內心裡面有沒有真正的喜樂?.
你的內心裡面有沒有真正的盼望?.
有沒有真正的喜樂,真正的愛呢?.
各位,讓我們都能夠在神的面前坦白.
各位,不要過一個沒有耶穌的聖誕,好不好?.
各位,在聖誕節慶祝的時候.
我們不要將耶穌放在旁邊.
不要像一個開生日會.
或者是他沒有份的壽星宮.
各位,今晚耶穌在你的生命當中佔什麼地位呢?.
我經常說,如果你將耶穌放在次要地位.
你等於不將祂放在任何地位.
今晚,我特別求主開啟我們的心靈.
讓我們看見真正的平安.
是從耶穌那裡來的.
真正的喜樂是從耶穌那裡來的.
真正的愛也是從神那裡來的.
各位,你肯不肯接受一個跟你有關.
而帶給你人生在任何環境當中.
仍然平靜的一個神呢?.
所以,我希望大家都接受這個挑戰.
各位,今年的聖誕節.
跟普通的聖誕節沒有什麼分別,是嗎?.
又是唱詩,又是唱聖誕歌.
又是吃那麼多餐.
那些餐都是千篇一律的.
但是,那是不是叫做Merry Christmas呢?.
我希望大家說一個真正的Merry Christmas.
好不好?好不好?.
如果是的話,我希望大家知道.
關鍵就是,你有沒有讓耶穌在你的生命當中呢?.
我希望大家思想,如果你接受這個耶穌.
那個有耶穌的聖誕,是不一樣的聖誕,是不是?.
如果說,我希望大家都能夠接受這個主.

$^{721}$做你生命的救主.
我們不能夠知道將來會怎樣.
我們不能夠人生的際遇怎樣.
我們可以預測.
但是無論怎樣,在任何環境當中.
有主的生命,是不一樣的生命.
我希望大家接受這個主耶穌.
做你生命的救主,好不好呢?.
我想請剛才的幾位弟兄姊妹.
再和大家一起唱,好不好?.
我們不用唱這麼多節.
我們唱兩節.
然後我希望大家和我一起祈禱,好不好?.
(唱).
各位朋友,當你唱的時候.
你的生命裡面有共鳴的.
請你在座當中起來.
一起唱,好不好?.
(唱).
有哪位請在座當中起來,一起唱.
(唱).
各位朋友,請在座當中起來,一起唱.
(唱).
我希望大家在這個時候安靜.
如果你希望能夠有主和你一起.
那個有耶穌的聖誕.
你接受主做你生命的救主.
請你在座當中舉高手來.
舉高手來,好不好?.
在座當中舉高手來.
我和你一起祈禱,好不好?.
看到了,看到了.
還有哪位呢?.
還有哪位?我看到了.
好,你放下就好了.
還有哪位?.
還有哪位呢?.
我想大家有一個真正具體的行動,好不好?.
你想我和你一起祈禱.
這麼多人和你一起祈禱.

$^{761}$請你走到前面,好不好?.
走到前面來,我們一起祈禱,好不好?.
走到前面來.
請基督徒的,可以帶領微信主的.
一起走到前面,好不好?.
(唱).
還有哪位?請你走到前面來.
不要客氣,我們一起祈禱.
走到前面來,走到前面來.
請在座當中是基督徒的.
請你幫忙隔離.
或者你陌生面孔的.
你鼓勵他.
使他能夠接受主耶穌.
做你個人的救主.
(唱).
還有哪位?.
(唱).
在座當中已經是基督徒.
議員在你面前作見證.
我要過一個有耶穌的聖誕.
請你一起出來.
(唱).
是基督徒的.
在人的面前作見證.
我要過一個有耶穌的聖誕.
不一樣的聖誕.
請你走到前面來,好不好?.
走到前面來.
還有哪位?.
(唱).
不要看別人是你自己的決定.
還有哪位?.
我再呼籲在座當中.
是基督徒的.
請你是主動的.
帶領微信主的朋友.
和你自己作一個廟的見證.
請你走到前面來,好不好?.
(唱).

$^{801}$還有哪位?.
(唱).
好,現在我們一起祈禱.
在祈禱的時候.
我講一句.
大家跟我講一句,好不好?.
是我們每個人的禱告.
不只是我領大家.
我們從心裡做一個真正的.
在神的面前呼求的禱告.
親愛的主耶穌.
我們多謝你.
我們知道.
在你的裡面.
我們有平安.
求你幫助我們.
在這個聖誕節當中.
我們決定.
不過一個沒有耶穌的聖誕.
有耶穌的聖誕.
才是真的聖誕.
主啊幫助我們.
讓我們從今年開始.
過一個不一樣的聖誕.
多謝主.
因為你的降生.
帶給我們有盼望.
帶給我們有平安.
帶給我們有喜樂.
多謝天父.
願今晚我們每一個人.
在這裡參加聚會的.
從你那裡.
支出更大的平安.
更大的喜樂.
願你與我們同在.
因為今年的聖誕節.
是不一樣的聖誕.
多謝你.
願一切榮耀頌讚歸給主.

$^{841}$奉耶穌基督的聖名祈禱.
阿們.
Hallelujah.
好.
各位.
請坐.
主特別幫助你.
缺知的朋友.
請牧師安排一下.
\newpage



\section{以賽亞書 7:10-17}
\label{sec:nEkocKnhlmk}
\textbf{2023-12-24 聯合崇拜《以馬內利的再思》以賽亞書 7\_10-17 講員:盧思源姑娘}
\newline
\newline
連結: \href{https://youtube.com/watch?v=nEkocKnhlmk}{\texttt{https://youtube.com/watch?v=nEkocKnhlmk}} ~~~~ 語音日期: 2023-12-27
\newline
\newline
\hyperref[sec:7Vdee0ksMqc]{\small{< < < PREV SERMON < < <}}
~
\hyperref[sec:index_chronic]{\small{[返順時目]}}
~
\hyperref[sec:index_scriptual]{\small{[返順卷目]}}
~
\hyperref[sec:bpii0IGPrEA]{\small{> > > NEXT SERMON > > >}}
\newline
\newline
以賽亞書 7:10-17
\newline
\begin{longtable}{cl}
\hline
\hline
章節 & 經文 (和合本修訂版)\\
\hline
7:10 & \begin{tabularx}{0.7\textwidth}{X} 耶和華又吩咐亞哈斯： \end{tabularx} \\ \\ \relax
7:11 & \begin{tabularx}{0.7\textwidth}{X} 「你向耶和華－你的神求一個預兆：在陰間的深淵，或往上的高處。」 \end{tabularx} \\ \\ \relax
7:12 & \begin{tabularx}{0.7\textwidth}{X} 但亞哈斯說：「我不求；我不試探耶和華。」 \end{tabularx} \\ \\ \relax
7:13 & \begin{tabularx}{0.7\textwidth}{X} 以賽亞說：「聽啊，大衛家！你們使人厭煩豈算小事，還要使我的神厭煩嗎？ \end{tabularx} \\ \\ \relax
7:14 & \begin{tabularx}{0.7\textwidth}{X} 因此，主自己要給你們一個預兆，看哪，必有童女懷孕生子，給他起名叫以馬內利。 \end{tabularx} \\ \\ \relax
7:15 & \begin{tabularx}{0.7\textwidth}{X} 到他曉得棄惡擇善的時候，他必吃乳酪與蜂蜜。 \end{tabularx} \\ \\ \relax
7:16 & \begin{tabularx}{0.7\textwidth}{X} 因為在這孩子還不曉得棄惡擇善之先，你所憎惡的那兩個王的土地必被撇棄。 \end{tabularx} \\ \\ \relax
7:17 & \begin{tabularx}{0.7\textwidth}{X} 耶和華必使亞述王臨到你和你的百姓，並你的父家，自從以法蓮脫離猶大的時候，未曾有過這樣的日子。 \end{tabularx} \\ \\
[1ex]
\hline
\hline
\end{longtable}
$^{1}$大英姐妹平安.
願意我們開始之前.
我們都有個祈禱好嗎.
是平安夜的日子.
我們可以回到你的殿當中.
一同去敬拜.
一同去聽主你自己的話語.
一會兒孩子要分享.
你自己信息的時候.
求真理的聖靈.
與我們眾人同在.
我們無論說的.
聽的.
我們一同去進入.
你自己的真理當中.
願意我們每一個.
都給主你自己的話語.
去得著.
幫助.
得著提醒.
我們恭敬的.
將以下時間交上.
奉主耶穌基督明智祈求.
阿們.
大英姐妹.
今天是平安夜.
大家有沒有想到.
今天.
打算怎樣過呢.
是不是都想著.
可能今晚我們教會.
有平安夜報道會.
或者今晚會不會都會約了一些.
朋友或者家人一起過呢.
有時我們在一些節日的時間.
我們都很希望可以.
跟一些朋友.
有人一起陪伴.
一起過.
而在聖誕節.

$^{41}$其中一個很重要的訊息.
我想在無數人頭上.
我們的詩歌或者我們今天講題.
很多時候都會提到.
二馬來尼.
我相信呢.
大家都知道二馬來尼的意思.
二馬來尼的意思.
就是.
神與我們同在.
在聖誕節呢.
我們都會去紀念.
耶穌基督的降生.
很多時候我們都會用.
二馬來尼聽到.
會說二馬來尼.
或者我們引用.
剛才讀經都有讀過.
就是《馬太福音》一章28至23節那段經文.
當中記載耶穌基督的出生.
但是大家知不知道原來.
二馬來尼這個字呢.
原來最開初.
不是在《馬太福音》裡面出現.
在聖經裡面.
最開初的時候呢.
就是在我們讀的剛才.
二賽亞書當中.
原來呢.
二馬來尼呢.
有它的歷史背景.
今天呢.
所以很想和大家一起去分享.
亦都去思想.
二馬來尼對我們的意義.
和我們的關係.
剛才提過.
二馬來尼呢.
有歷史.
原來呢.

$^{81}$它是耶和華賜給亞哈斯王的一個兆頭.
當時它是發生在.
猶太國亞哈斯在位的時候.
如果大家去讀這個.
以色列的歷史.
有一段很長的時間.
以色列是分為南國.
猶大和北國以色列.
而今天我們讀的經文呢.
就是在亞述的管治之下.
當時呢.
先說一點歷史背景.
就是亞蘭.
在上面的亞蘭和北面的以色列.
受到威脅.
所以當時亞蘭和以色列.
就希望可以聯盟去反抗亞述.
亦都逼猶大國.
這個亞哈斯.
剛才提過.
去加盟.
不過呢.
這個猶大.
這個亞哈斯不肯加盟.
於是呢.
亞蘭和以色列的聯兵.
就想攻擊猶大的耶路撒冷.
甚至呢.
想威脅猶大.
廢除亞哈斯的皇帝.
另立一個皇.
如果大家可以看聖經.
因為我們剛才就說.
看二賽亞曰十章至十七節.
但都有一點背景.
原來呢.
當時呢.
剛才提過.
在這個威脅裡面.
當時的人呢.

$^{121}$其實是很惶恐的.
聖經在二賽亞曰七章二節.
這樣形容.
有人當時告訴大衛家的亞哈斯王.
他說亞蘭和以法連已經聯盟.
王的心呢.
就和百姓的心一起顫動.
好像林中的樹林一樣.
當中呢.
提到猶大的國王.
即是亞哈斯王.
和當時的百姓.
都很不安.
就好像林中的樹林一樣.
大家幻想一下當時的處境.
我想大家在感受上.
應該會容易一點理解.
當很大風吹的時候.
吹落樹的時候.
樹就會震動.
震動很搖動.
當時呢.
聖經形容亞哈斯王.
和當時的百姓.
就是這麼緊張和害怕.
好像震動一樣.
所以呢.
在這個時間的時候.
耶和華興起一位先知.
叫做以賽亞.
去提醒亞哈斯.
在第七章三節.
耶和華對以賽亞說.
就說.
要你和你的兒子.
施亞述出去.
去上池的水溝盡頭.
往漂泊的大路上去迎見亞哈斯.
即是要他們去找亞哈斯.
但是呢.

$^{161}$原來大家看到聖經.
這個是有些意思的.
原來呢.
為什麼要和他的兒子去的時候呢.
以賽亞的名字.
本身有個意思.
就是耶和華的救贖的意思.
而施亞述的意思.
就是愚民要返回.
原來當神叫以賽亞和他的兒子去的時候.
當中的意義其實是.
接著說以色列人將來發生的歷史.
如果我們再看經文.
大家去看的時候.
講一點點之前的背景.
以賽亞和他的兒子.
真的去到亞哈斯的面前.
就和亞哈斯說.
叫他們不用怕.
鎮定.
不要因為利遜和亞蘭.
和以瑪利的兒子冒起的火頭.
而去擔心.
主要就是說.
耶和華和他們說.
這些事是不會成就的.
那些人想去攻擊你嗎.
不過是不會成就的.
甚至乎神應許.
在65年之內.
以法年.
這個國家是一定會滅亡的.
所以就叫亞哈斯不需要怕.
簡單來說.
以賽亞宣告.
就是告訴亞哈斯.
不需要擔心.
縱然現在那個環境.
好像有兩個國家要去攻擊你.
但是他們的計劃.

$^{201}$是不會達成的.
就好像兩個沒有煙的火頭.
即是火把頭.
大家不要以為沒有煙的火把頭.
就是很大火.
這裡說的.
其實他們兩個火把頭.
已經燒盡了.
那些煙.
是白煙.
即是神告訴.
藉著以賽亞告訴亞哈斯.
不需要怕.
雖然他們同盟.
你很擔心.
但是神知道這件事.
去保證亞蘭和以法年.
接下來會瓦解的.
而應付這個危機.
最好的方法是什麼呢.
就是相信神.
相信上帝的大能.
亞哈斯如果要被建立的話.
他就一定要.
必須相信.
所以上帝.
祂叫亞哈斯尋求一個機會.
第11節.
祂這樣說.
你向耶和華你的神求一個預兆.
在陰間的一心願.
或往上的高處.
這裡的意思是.
無論在哪裡.
無論是什麼.
無論有多難.
只要亞哈斯.
祂肯求.
只要祂開口的話.
神都會答應祂.

$^{241}$都會滿足祂.
如果換一句說話.
耶和華給了一張.
簽了名的支票給亞哈斯.
不過沒有說銀碼.
只要亞哈斯說得出.
神就做得到.
也都一定會為亞哈斯和猶大的百姓去成就.
大家聽到這裡的時候.
有沒有留意.
其實整件事.
不是因為亞哈斯去求而成的.
相反.
是耶和華去主動.
去積極的.
去主動給人.
甚至要求亞哈斯去接受這個訊息.
有時候.
我都會想.
我們作為基督徒.
我們都很幸福.
因為神.
允許祂說.
祂常與我們同在.
不過.
有時候我們又很容易.
去輕忽了這份同在.
又或者我們對神的同在.
好像不算是很積極.
有時候我都會想.
有些會反思.
有時候我身邊都有很多.
拜偶像.
或者信其他宗教的人.
當他們有危難的時間.
他們馬上就去求神拜佛.
他們就去全心全意去相信他們的假神.
甚至大家在新年的時候.
農曆年初一.
看到人們爭著去上.

$^{281}$投著香的.
就是迫爆的.
爭著去上投著香.
真的會看到.
可能我們會看.
他們所信的是偶像.
不過他們真的全心全意去相信他們所求的.
他們為著.
希望他們所信的.
可以保護他們.
又或者我都會想.
身邊有一些穆斯林的朋友.
當他們見到的時候.
不知道大家身邊有沒有.
他們真的會定時去祈禱.
一天真的定時定候.
他們就會去祈禱.
有時候我會反思.
今天我們作為基督徒的時候.
有時候好像我們對神.
都沒有那份主動.
幫問及一些弟兄姊妹.
你最近的禱告靈修怎麼樣.
可能大部分都會回答.
嗯 有的.
不過不是經常.
當可能有需要的時間.
或者我記得的時間.
才會做.
曾經我看過一本書.
那本書提過.
聖經原來.
它說有七千多個.
上帝給人的英雄.
它統計過.
舉個例子.
可能當中舉例.
上帝說不會撇下我們.
祂是隨時幫助我們.
換句話說.

$^{321}$聖經裡面有七千多張.
上帝給我們未兌現的支票.
這一切都是幫助我們.
不過可能有時候.
我們對於上帝的幫助的時候.
可能我們都未必是很積極.
縱然剛才我們知道.
上帝是與我們同在.
但很可能.
很可能當我們遇到問題的時候.
我們很習慣.
習慣不是即時主動去尋求神.
可能我們很習慣.
即時想.
我們自己有什麼方法去解決呢.
回到經文裡面.
我們看看阿克斯的反應.
阿克斯怎麼說呢.
他說.
我不求.
我不試探耶和華.
即是對於上帝.
耶和華很積極地主動.
去給阿克斯一個機會的時候.
阿克斯說.
我不求.
他用「不」字來形容.
其實看到這個回應的時候.
我們覺得好像很屬靈.
好像很虔誠.
因為他說.
我不試探耶和華.
但什麼是試探呢.
簡單來說.
假如人向上帝求一件事.
他的動機是要求上帝去證明他的能力的話.
這個是試探.
又或者.
如果人刻意去做一些事.
去看看上帝的反應.

$^{361}$舉個例子.
可能他刻意去犯錯.
看看上帝是否真的會懲罰.
這個也是屬於試探.
但如果我們回到這裡去看的時候.
這個情況.
如果阿克斯好像以賽亞這樣說.
向上帝求一個兆頭.
當然這個不是構成試探耶和華.
因為是上帝採取主動.
上帝已經準備好.
是在等待人的回應.
但可能我們又會想.
為什麼神會有這個要求呢.
倒不如直接將這個兆頭給阿克斯就行了.
為什麼要這麼轉折呢.
如果這個兆頭是這麼重要.
為什麼要阿克斯求先給他呢.
好像聽起來不是人試探神.
好像是上帝在試驗人.
原來我們回到這裡去看的時候.
問題的重點是.
原來阿克斯的求和不求.
是有很大的意義和重要.
我們這樣去理解的時候.
如果阿克斯求.
他表示他對上帝的信服.
也對上帝的信求.
更重要的是.
他表示他願意對上帝有一份開放的心.
如果他願意去接受這個兆頭的時候.
他是在給上帝一個空間.
給上帝一個機會去工作.
上帝願意在這個時間和他的百姓同行.
他要阿克斯向他求兆頭.
而不是直接給他的時候.
是希望阿克斯可以主動.
可以積極去開放.
去自己去信靠上帝.
去給上帝工作.

$^{401}$回到今天.
其實上帝也是一樣.
上帝是不會逼我們的.
我們有時都會問.
很多人都會問.
為什麼上帝不把我們創造成一個好像機械人一樣.
或者當他叫我們的時候.
我們就全部都做.
都說yes.
這樣不就好了嗎.
每個人都聽話.
全世界我們都相信神.
每個人都相信神.
這樣不是更好嗎.
但當我思想的時候.
我會發現上帝真的很愛我們.
他給我們一份自由意志.
給我們可以去選擇.
我經常都會去舉例.
當你很喜歡一個人的時候.
你是希望對方是甘心樂意.
去喜歡你 去愛你.
而不是強迫對方.
因為強迫的那份不是愛.
唯有你甘心樂意才是真正的愛.
所以在這裡也是一樣.
上帝很希望阿克斯可以主動.
他可以開放.
去給機會 去給上帝在當中去工作.
不過很可惜.
阿克斯他是用一個筆來回應上帝.
上帝很想去幫助人.
不過有時得到的回應是很難忘.
當我們剛才說「以馬內利」這個字.
「以馬內利」這個應許.
在歷史裡面第一次出現的時候.
竟然遇到這樣的對待.
當上帝很主動去尋找人的時候.
反而是被拒於千里.
大家有沒有試過對上帝say no?.

$^{441}$可能有時我們都很不自覺.
會對上帝say no.
我自己曾經都試過對上帝say no.
記得有一次曾經和朋友出現一些誤會.
可能心裡覺得我都不太想找回他.
不過心裡上帝就提醒.
我不應該這樣.
我應該要和他復和.
但我一直不想.
這個所謂的冷戰關係就持續了一段時間.
但心裡知道要主動和.
所以聖靈不時都會去提醒.
最後都是順服在聖靈的提醒下.
去找回這個朋友.
很感恩的.
找回的時候大家可以去聊天.
說回之前中間的誤會.
大家有機會去復和.
有時我都在想.
如果當初不是因為自己不想聽神說.
和神say no.
不願意聽神說的時候.
可能這段關係可以早些的好起來.
原來有時神是很想幫我們.
而我們其實都明知應該怎樣做.
不過有時我們都會選擇向神say no.
可能只不過在心底裡.
未必是真的很想上帝幫我們.
又或者我們覺得神的方法未必是最好.
弟兄姊妹.
今天在你心裡面.
有沒有一些事情.
上帝是想你去做的.
或者你還在猶豫當中.
甚至好像我之前那樣.
很想和上帝say no.
願意聖靈都去提醒我們.
我們都是去信服聖靈.
給我們那份感動.
回到經文裡面.

$^{481}$當阿克斯向神say no之後.
先知爾賽亞責備阿克斯.
他這樣說.
聽啊 大衛家.
你們使人厭煩喜訊小使.
還要使我的神厭煩嗎.
因此主自己要給你們一個預兆.
看啊 必有同類懷孕生子.
給他起名叫爾瑪萊利.
到他曉得棄惡擇善的時候.
他必吃與樂和風物.
因為在這孩子還不曉得棄惡擇善之先.
你所珍惡的那兩個皇的土地必被撇棄.
神積珠.
爾瑪萊利這個兆頭向阿克斯.
啟示了他的心意.
我們會想 什麼叫兆頭.
兆頭其實是一個記號.
指向一個真實的存在.
當有一個真實的存在的時候.
才有這個兆頭.
並不是因為有這個兆頭.
而有這個真實的存在.
簡單來說.
好像我們見到一本書.
你見到一本書的時候.
你會知道有一個作者存在一樣.
所以回到這裡的時候.
兆頭的重要不是在說本身.
而是這個兆頭.
在傳遞一個什麼樣的訊息給我們.
到底爾瑪萊利的應許是什麼呢.
他的內容這裡說得很清楚.
第一 會有同類懷孕生子.
他的名字是爾瑪萊利.
第二 就是當孩子還很小的時候.
那兩個來勢洶洶要消滅大衛家的國.
必被撇棄.
所以在這個兆頭.
當時對阿克斯來說.

$^{521}$其實是在說什麼呢.
爾瑪萊利的應許.
重點是在說歷史上將要發生的事.
而這件發生的事.
上帝要表明的就是.
不是偶然.
不是人為的因素.
這個和人的智慧沒有關係.
是上帝自己的工作.
是照上帝的計劃來進行.
所以爾瑪萊利這個兆頭.
就是上帝為了人的緣故.
將他自己的計劃向人去顯示.
揭示出來.
讓人知道上帝在掌權.
讓人可以安心.
讓人不需要像隨風擺的樹一樣.
隨風飄搖.
但是當阿克斯去拒絕這個兆頭的時候.
他不接受這個安排的時候.
他不認同這個計劃是好的時候.
他自己有更好更理想的安排.
其實無論是在說阿克斯那個時代.
或者是今天.
接受這個兆頭其實只有一個意思.
就是我們要相信.
相信上帝的能力.
相信上帝在掌管.
相信上帝的拯救.
所以在這裡.
耳塞又說.
你們若是不信.
必站立不穩.
信是可以叫我們在風雨的時間.
可以屹立不倒.
信就是對於上帝給我們的應許裡面.
唯一的回應.
信就是在我們處境上面.
我們學習放棄我們自己的計劃.
放棄我們自己的打算.

$^{561}$我們走在上帝的路上.
接受他那份安排.
記得之前有個見證.
因為我做青少年工作.
曾經有個見證.
有一個肢體.
大學生來跟我分享.
他在海外留學.
他選大學的時候.
他很想選海外一間大學.
因為他覺得那間大學出路好一點.
科目好一點.
排名好一點.
他就為這件事去祈禱.
交給盼望上帝為他安排.
結果的時候.
他一路祈禱.
結果的時候.
出乎他意料之外.
不是他所預期.
他去了另一個學校.
讀了另一科.
他說收到的時候.
他覺得很不開心.
他覺得上帝為何我為這件事去祈禱.
為何上帝好像不應運我的祈禱.
沒有回應.
不過當他時間過了的時候.
他跟我分享.
他說他明白.
他明白為何神當善事不應運他.
原來他原先很想讀那間心儀的大學.
他有些朋友在那裡.
原來很忙.
不單是很忙.
根本沒有時間可以讓他回到教會.
相反他去了自己那間大學.
讀了幾年的時候.
他在那間大學裡.
雖然漢始的科目未必是最好.

$^{601}$但是他裡面的團契.
大學生的團契.
是非常之好.
他說他這幾年.
他在大學生團契裡.
很多很多的成長.
他很感恩上帝帶領他去現在的大學.
大兄姐妹.
今天你是不是真的相信上帝.
在我們生命裡.
每一個計劃.
每一個安排.
我們信不信得過.
上帝為你預備的是最好.
縱然這一刻是看不見.
以馬來理.
意思是神與我們同在.
上帝應許他好像母雞一樣.
將熟悉他的人.
焦聚在他翅膀.
膀背的底下.
去保護他們.
但不過有時可能我們不願意.
有時可能我們會靈覓另一種方法.
可能我們有時會向上帝say no.
這樣以馬來理有沒有意思呢.
好像沒什麼意思.
以馬來理有沒有作用.
但試問如果是這樣的時候.
怎樣發揮作用呢.
就好像人.
我們假如手裡一個很新的工具.
我們用的時候可以事半功倍.
不過如果不用的時候.
就算這件工具有多好.
我們不給機會的時候.
有什麼作用呢.
當然以馬來理.
我不是說工具.
也不是視為工具.

$^{641}$但如果我們不接受的時候.
上帝這份好意.
只會是白費.
回到聖經裡.
我們看看阿哈斯最後是怎樣呢.
在《列王記》下記載.
阿哈斯怎樣去回應上帝.
怎樣去回應 怎樣去做.
《列王記》下十六章說.
阿哈斯派使者去亞述王提勒拉彼列色.
說我是你的僕人 你的兒子.
現在亞納王和以色列王攻擊我.
求你上來 救我脫離他們的手.
這裡很特別.
剛才我們說我們見到阿哈斯.
向耶和華說NO.
但是他走去亞述王那裡.
去說YES.
去求阿哈斯.
上帝求阿哈斯.
不是 上帝要求阿哈斯給個機會.
給個空間 給上帝工作.
但是阿哈斯說他不願意.
他不求.
但相反去求亞述王.
去求他脫離敵人的手.
阿哈斯他相信亞述.
向亞述投信任.
但是他向耶和華去關門.
最後我們看聖經我們見到.
阿哈斯甚至將聖殿和皇宮所有的金銀.
送給亞述王.
最後阿哈斯讓猶大國成為亞述的附庸國.
在亞述的逼迫之下.
正如第17節這樣說.
耶和華必是亞述王臨到你和你的百姓.
令你的父家自從二百年脫離猶大的日子.
未曾有過這樣的日子.
如果我們看這代記下.
這代記下28章這樣說.

$^{681}$亞述不但去求亞述王.
亞述不單沒有幫助他.
反而讓阿哈斯他們受苦.
最後阿哈斯失去了國家獨立的自主.
甚至失去了當時猶大國很多的財富.
弟兄姊妹.
以馬內利.
上帝與我們同在.
這是來自上帝給我們的應許.
這個應許是神向我們承諾.
上帝向我們表明.
無論在任何的處境.
任何的環境.
遇到任何的挑戰的時候.
祂依然是那位可信靠的神.
祂同在臨在在我們當中.
以馬內利不單單是聖誕節裡面的一個主題.
而是在我們一生一世的主題.
是關於上帝救贖我們的主題.
盼望在平安夜的日子.
我們在這個聖誕節假期裡面.
我們在思想以馬內利的時候.
我們向神闖開我們的心.
向神進到我們的內心當中.
因為祂是那位愛我們的神.
是那位與我們同在的神.
我們一同低頭祈禱.
你是那位以馬內利的神.
你與我們同在.
無論在我們生活裡面.
我們遇到任何的處境.
我們畫高畫低.
神你應許你時刻的.
都是陪伴我們.
與我們經過.
主要求你幫助我們.
幫助我們在平安夜的日子.
在我們紀念耶穌基督降生的日子的時間.
我們記得神你是時刻與我們一起.
盼望當我們遇到任何事情的時候.

$^{721}$我們向你闖開我們的心.
我們真的信服上帝你給我們的那份大靈.
信服你給我們那份感動.
以致我們走在你心意的當中.
成為你所喜悅的兒女.
天父求你繼續幫助我們各人.
願我們各人在聖誕的日子裡面.
我們更深經歷你對我們的那份愛.
天父我們多謝你.
聽我們祈禱.
奉主耶穌基督名字祈求.
阿們.
\newpage



\section{}
\label{sec:bpii0IGPrEA}
\textbf{2023-12-25 聖誕節崇拜《諸天萬物同歡唱》 講員:陳麗蟬牧師}
\newline
\newline
連結: \href{https://youtube.com/watch?v=bpii0IGPrEA}{\texttt{https://youtube.com/watch?v=bpii0IGPrEA}} ~~~~ 語音日期: 2023-12-27
\newline
\newline
\hyperref[sec:nEkocKnhlmk]{\small{< < < PREV SERMON < < <}}
~
\hyperref[sec:index_chronic]{\small{[返順時目]}}
~
\hyperref[sec:index_scriptual]{\small{[返順卷目]}}
~
\hyperref[sec:MkLoNpPCJk4]{\small{> > > NEXT SERMON > > >}}
\newline
\newline
$^{1}$有主的一個使者站在他們旁邊.
主的榮光四面照著他們.
牧羊人就很懼怕.
忽然有一大隊天兵和那天使讚美上帝說.
在至高之處榮耀歸於上帝.
在地上平安歸於他所喜悅的人.
大家怕不怕黑呢.
如果在黑漆漆的晚上高.
突然有事發生的時候.
或者有人大叫.
或者有大光的時候.
是不是會更害怕呢.
當耶穌基督降生的那天.
這班牧羊人就在郊野裡面.
對於這班牧羊人來說.
晚上他們不是那麼害怕的.
因為他們習慣在郊外的地方.
晚上可以看著羊.
不然羊就會被人偷走.
或者被猛獸搶走.
但是在那晚黑漆漆的晚上.
有一道強光出現了.
嚇得這班牧羊人半死.
這班牧羊人其實當時他們的處境.
不單單每一天都要在黑漆漆的晚上上來工作.
他們還要生活在黑漆漆的社會裡面.
當時管治他們這個地方的.
就是這個希律王安提柏.
他是一個荒淫無道的君王.
他以一個不法的手段.
搶了他同父異母的哥哥的太太希羅底.
當施洗約翰來指證他的錯誤的時候.
他就在他自己生日的會裡面.
藉著太太設計.
殺死了施洗約翰.
還有當耶穌基督被釘十字架的時候.
這個希律王安提柏.
是主刀走去誘訓,走去審耶穌的.
所以當時這個管治的王.
他的罪性也都帶來整個社會的人很多的不安.

$^{41}$牧羊人就住在這麼黑的環境裡去生活.
晚上來工作.
在這個黑社會裡面去生活.
當這個大光一進入了這個牧羊人的世界的時候.
最初這個牧羊人是很害怕的.
因為不知道發生什麼事.
但後來他定一定了.
他竟然可以將救世主耶穌基督降生的信息.
一路來去傳開去.
將這個榮耀歸給上帝讚美他.
人生有時都會陷在一個很黑暗,很恐懼.
也都覺得很迷茫.
不知道前面的路程是怎樣的時候.
可能是有這個重病.
又可能是一些感情的問題.
又是工作的問題.
但只要有耶穌基督的大光介入在裡面的時候.
就會經歷到耶穌基督很奇妙的帶領和指引.
生命就會改變了.
我第一次拍拖的時候.
其實我掙扎了一年才投入到那段感情.
但拍拖了一年之後.
我發覺這個弟兄他開始對我冷淡.
我就開始一路去經歷一下.
看看什麼事.
我發覺原來在教會裡面又有另外一個女朋友.
很慘的.
因為平時逛街的日子,聽電話的日子.
現在就沒有人陪我逛街,沒有人跟我談電話.
我就呆呆的坐在家裡.
初時我都會祈禱的.
不過一路我都會祈禱.
不過這個祈禱有很多埋怨上帝,怪責上帝.
為什麼會這樣?.
很多的不忿,很多的不喜歡.
但當我經過了兩個星期.
不斷地怪責,不斷地祈禱的時候.
我發覺我的心靈一路一路有轉變.
我的心靈開始看到.
人是很容易變的.

$^{81}$唯一是上帝是不會變.
他是完全可以信靠的.
所以我就第一次向上帝納志.
要一生一世的來事奉祂.
我教了初五年之後.
我又落在另一個人生的大低潮.
哇!很黑啊!.
其實我工作很順利的.
校長也很喜歡我.
學生也很喜歡我.
家長也很喜歡我.
我的同事相處得很融洽的.
教會也參加了師班的事奉.
我也是傳道部的部長.
不過不知為何覺得人生很苦悶.
淡而無味.
不知做什麼好.
我找了一個下午.
我自己一個人和耶穌講掃.
不是祈禱,是講掃.
我就告訴耶穌.
我在教會裡盡心盡性的事奉祂.
學校也有這麼盡心的工作.
為什麼呢?.
在感情的事上處處來碰釘.
我覺得弟兄很壞.
當我說這句話的時候.
我就記得.
我在第一次感情的挫敗的時候.
我就說我要一生一世的來事奉上帝.
因為人很容易變.
所以在那個時候.
上帝開始和我說話.
祂說你立了志第一次.
接著在不同的聚會裡.
我已經三次,四次呼召你起行.
但是你一步都無欲.
我就告訴上帝.
我怕不教書.
將來沒有這麼多收入.

$^{121}$我又怕我不會做.
我更加怕孤單.
上帝只回應我一句.
我與你同在.
我一聽到這句話的時候.
心裡的迷茫,孤單,無力.
全部不知道為什麼就出來了.
接著在腦海裡就出現了一幅圖畫.
就是原本很多垃圾的河.
裡面有樹枝,樹葉,紅色的背心袋.
很骯髒的河流.
突然那些垃圾.
不知道為什麼全部飛走了.
變成一條很清澈,很清澈的河流.
也經過了下午之後.
我就開始去報讀神學院.
人生有時候會遇到一些不如意的事.
很迷茫,很不安的日子.
但當耶穌基督介入在裡面的時候.
生命就不再一樣.
有光彩,有希望,有意義.
「無人人彼此說.
我們往伯利行去.
看看所成的事.
就是主所告訴我們的.
他們急忙去了.
找到瑪利亞和藥室.
還有那英鞋臥在馬槽裡.
他們看見就把天使論者孩子的話傳開了.
聽見的人都詫異無人人對他們所說的話」.
在耶穌時代.
這些牧羊人其實是很低下的人.
人們覺得他們的工作很骯髒.
也是很卑賤的工作.
所以很多人都看不起這些牧羊人.
他們的收入又不高.
還有那些羊不是他們自己擁有的.
多數那些有錢人請他們回來幫忙.
幫他們去看羊.
這些牧羊人地位不高.

$^{161}$沒有尊嚴.
而且時時都要在曠野裡看羊.
所以他們都沒有家庭.
當時更慘的是.
他們不讓進入猶太會堂.
因為覺得他們是不乾淨的人.
所以他們都是被忽略的.
有時我們也會.
我們也會在不知不覺之間.
喜歡高舉一些有錢人.
出名的人.
或者一些達官貴人.
我們覺得跟他們在一起的時候.
或者認識他們的時候.
我們會覺得聰明一點.
或者站得直一點.
是一種光榮.
不過耶穌不是這樣看人的.
耶穌是看我們每一個人的生命.
都是尊貴.
都是寶貝.
所以當主耶穌基督降生世上的時候.
這是一個很大的訊息.
祂就選了當時在社會裡.
被認為是一些很低下層.
別人看不起的這班牧羊人.
就由他們將這個大喜的訊息傳開去.
當這個牧羊人.
自己知道了元代救世主降生.
他知道了.
他都很開心.
就將這麼重要的訊息.
一直告訴周圍的人.
這班牧羊人的生命開始轉變了.
他由一個原本被人看不起的.
被人拒絕的.
被人遺棄的這些人.
但到現在.
他就成為了一個傳福音報好訊息的見證人.
雖然他們的生命轉變了.

$^{201}$但不過他所處的環境沒有改變.
他生活仍然是這麼困苦.
仍然都是這麼窮.
他仍然都要在郊外.
漆黑這樣去工作.
不過他知道.
他知道.
他的生命在耶穌基督裡面是尊貴的.
上帝仍然會重用他的.
所以弟兄姊妹.
無論今天我是一個博士級.
還是一個沒怎麼讀書的人.
又或者我是公司的總裁.
又或者我只是一個家庭裡面的主婦.
千萬不要覺得自己很有用.
或者覺得自己很沒用.
上帝看他所創造的每一個人.
都是尊貴的.
他都會使用的.
所以找著每一個時間.
上帝給我們機會去服侍他的時候.
就珍惜.
主耶穌因也擺在前面的喜樂.
興旱羞辱.
忍受了十字架的苦難.
如今已坐在神寶座的右邊.
你們要仔細想想.
這位忍受了罪人如此頂撞的耶穌.
你們就不至心灰意冷了.
這個頂撞這個字.
原來那個意思就是反對你.
他不聽你的話.
那麼耶穌基督來到這個世界.
他原本不受時空的限制.
但來到這個世界的時候.
他就受了這個時空的限制.
他天上原本是滿有尊貴榮耀的.
但來到世界的時候.
他處處被人頂撞.
當他被彼拉多尼去審判的時候.

$^{241}$全部的人都大叫什麼.
釘他十字架.
耶穌基督為了完成救贖世人的工作.
他默默的忍受.
忍受這些不公平的對待.
忍受所有人給他的凌辱.
這個就是耶穌的美.
他忍受人一切給他的凌辱.
和這個不公平.
為了救贖世人.
他都給我們寬恕.
在十二月初的時候.
我就問准我姐姐.
我就請我的印尼姐姐.
去我家做了十七年.
照顧了我爸爸十幾年.
請她去日本旅行.
過程有些艱辛.
申請簽證那裡.
繁複到不得了.
也有一些繁複的手續.
好了 搞定所有的難關.
就開始去了.
很開心.
到最後一天.
他竟然口黑臉又黑.
不知道他做什麼.
那天又不肯吃東西.
進到機場過了關.
我和他兩個人的時候.
他就說他很生我的氣.
因為我沒有尊重他.
他說你看見一盒士多啤梨.
為港幣七十多元.
他問我可不可以買.
我說這麼貴不要買.
回香港再吃.
然後他又看到一個.
日本的蛋撻四十多元.
他又問我買不買.

$^{281}$我說這麼貴不要買.
先回家.
但其實我給了錢他有自由.
他不來問我什麼都可以.
我聽到的時候.
我都有點激氣.
我沒有出聲.
我只是心想.
是真的 我給了錢他.
你自己喜歡做什麼就做什麼.
我覺得.
我幫你排除這麼多的萬難.
然後出錢 出力 出時間.
其實十二月我很忙.
我也不想去.
帶他去日本玩.
目的就是完成了他的心願.
他想下年完約的時候.
他就回印尼.
他就能夠去一次日本.
我覺得.
我出了這麼多的心機.
這麼多的東西.
但結果是這樣的收場.
他說我不尊重他.
他喜歡吃士多啤梨.
你不能吃.
我覺得.
我很心淡.
所以當我進了車站.
開始一路走去搭飛機.
不知道為什麼我們兩個.
因為大家都不開心.
又沒有說話.
我們又不懂得走去搭機場巴士.
我們直走.
走去那裡走十幾二十分鐘.
路上我一直覺得很難過.
心也很淡很淡.
我覺得.

$^{321}$我為了完成他的心願.
但結果去完旅行.
我就被他那些黑口黑面的去埋怨.
死啦.
還有大半年時間.
怎麼捱.
一路走的時候.
我真的在心裡哭.
我覺得為什麼會這樣.
付出了一些東西.
得回來的是另外一樣的東西呢.
但那時候突然間有個醒覺.
因為在這個機場播放聖誕的音樂.
就是平安夜.
我就突然間記起.
耶穌基督來到這個世界上.
都是.
他捨棄了天上高的尊貴和榮耀.
來到這個世間上高.
來拯救人.
他何嘗不是來去忍受很多人的頂撞.
很多人的辱罵嗎.
耶穌基督在地上高.
他所承受的比我更慘更淒涼.
心想.
回到香港.
又要捱幾天的黑口黑面了.
真是慘.
不過上帝很奇妙的.
在第二天的早上.
他醒的時候很開心.
他就跟我說.
姑姐早安.
哇.
為什麼這麼開心.
耶穌昨晚提醒了他.
因為你無端端有六天的假期.
因為你看到很美的風景.
有紅葉有黃葉有竹林.
還有.

$^{361}$你也吃了日本很多很多.
很特別的食物.
你不可以.
耶穌告訴他.
你不可以因為沒有士多啤梨.
沒有日本的蛋撻.
吃就這麼生氣.
他說耶穌會幫我的.
你只幫他又不是幫我.
他說如果當他忽略了上帝給他這麼多的恩典.
他只看著自己沒有那些東西的時候.
他說上帝會幫他.
如果這番話從我的口出的時候.
一定會罵多幾場架.
但是上帝就是這麼奇妙的.
在他的裡面沒有難成就的事.
他的說話滿有能力.
能夠改變人心.
所以我問准他.
我可不可以講這些經歷.
他說可以.
所以就要跟大家講.
上帝我們的耶穌.
是那位奇妙的上帝.
去承受人的頂撞.
但祂本來是尊貴榮耀的.
(主席翻譯).
在末後的日子.
萬民都要留歸這山.
這山就是上帝的山.
上帝的山.
即是有上帝的同在.
有上帝的管治.
不過我相信.
這個不需要一定去到末後的時候.
才去發生.
在今天的裡面.
只要我們肯.
在教會的裡面.
在家庭的裡面.

$^{401}$在我們工作.
在我們生活的地方裡面.
行出主耶穌基督的謙卑.
行出主耶穌基督的捨己.
饒恕的話.
這樣上帝這個和平的國度.
就會臨到這個地上.
在這一刻.
就願意我們每一個人都立志.
要活出主耶穌基督降生的意義.
這樣全世界都一起歡騰.
我們一起起來.
同唱普世歡騰.
普世歡騰教主降臨.
全民接同為我.
安心慰住.
願被提防.
宇宙萬物國長.
愛的源頭.
全自我.
不問高聖中的天地薄弱.
平英三國.
向英國聖表亮.
世上一切罪惡有上.
從此不再生疆.
普世造作變壞人意.
全愛澤及四方.
耶穌是以公義統治世界.
以信實對待萬民.
耶穌起人次晚安.
要在萬國安葬.
將天上傳.
公義榮光.
全愛其妙方式.
普天同慶君王降臨.
萬國萬民都要敬拜他.
阿們.
\newpage



\section{撒迦利亞書}
\label{sec:MkLoNpPCJk4}
\textbf{2023-12-28 周四晚查經班《歸回先知書研讀(11)撒迦利亞書》 講員:曾裔貴牧師}
\newline
\newline
連結: \href{https://youtube.com/watch?v=MkLoNpPCJk4}{\texttt{https://youtube.com/watch?v=MkLoNpPCJk4}} ~~~~ 語音日期: 2023-12-29
\newline
\newline
\hyperref[sec:bpii0IGPrEA]{\small{< < < PREV SERMON < < <}}
~
\hyperref[sec:index_chronic]{\small{[返順時目]}}
~
\hyperref[sec:index_scriptual]{\small{[返順卷目]}}
~
\hyperref[sec:lIj_mrEVPGg]{\small{> > > NEXT SERMON > > >}}
\newline
\newline
$^{1}$我們一起禱告.
感恩天父讓我們一直看歸回的歷史書.
和先知書.
撒迦利亞書.
都是平常教會很難在港台崇拜去講的.
我們都算是有系統逐段去看.
一至八章.
八個的一些預言,比喻.
和由第九到十四章.
有很重要的對將來.
對以色列,對列國的預言.
尤其是主再回來的第二次降臨.
今天晚上我們會看最後一章.
和稍後我們會看馬拉基書.
很盼望對於十二小先知.
特別是歸回的歷史.
我們都能夠有一些明白.
雖然和我們的時代.
相差已經是二千五百多年了.
但我都盼望我們在新約的時代.
我們其實更加準確對於彌賽亞的降臨.
主耶穌將要第二次再回來.
這個的應許.
都是我們要等候.
我們要有盼望.
求主帶領我們今天晚上的茶經.
禱告奉耶穌基督的名.
阿們.
很感恩一直和大家.
每個星期都有機會去讀聖經.
來看.
都認真去看.
雖然我們不是用原文識經.
一來我都不是在這一方面.
有這樣的裝備.
所以都是透過經卷.
透過一些歷史.
透過一些經文的上下文的識經.
和以色列人的歷史背景.
幫大家去組織了.

$^{41}$對於當時代盡量我們能夠還原歷史.
盼望大家能夠在往後的日子.
有不同的學習.
不同的追求.
如果大家時間有許可.
有些晚間的課程.
如果是再有深度一點.
例如一些可能真的會碰到原文等等.
那些其實都值得去報讀的.
當然要看大家的時間.
和大家的研習興趣.
這個就真的不勉強了.
不過至少你一定要讀完文之前.
你自己要對於和合本.
你都要很熟悉.
甚至某些關於舊約的神學.
猶太人的歷史.
這些可能都有基本的基礎.
如果能夠都有.
是會更加相得益彰.
所以盼望大家我們彼此一同去鞭策.
有什麼提問.
真的可以繼續在學習上延伸.
這個群組我們會繼續下去.
如果你認識有些人.
他很想對聖經有多些認識.
你不妨告訴我.
我又可以加他進入這個群組.
使得有時我們討論某些問題.
有些人又能夠知道.
我又不是太多在YouTube上.
聽過了一兩堂.
我要知道大概講到哪裡.
有時我都會按一下.
很偶然會見到有人留言.
因為已經很多了.
已經儲了我想有過百堂了.
所以如果你們有機會重溫.
又見到有人留言.
不妨又告訴我.

$^{81}$我看看那些留言是否重要要回覆.
因為我很少會整段再聽.
所以大家幫幫忙.
如果有些聽到突然.
特別是晚上.
你可能某個時段.
我自己都會.
聽一聽會睡著了.
如果你再回去聽.
其實可能又另一種明白和認識.
所以有時一閃就過了.
講的人就過了.
但聽的可能你會疏忽了.
會遺漏了.
所以現在就好了.
有了紀錄.
你聽回來可能對你看那段經文有幫助.
因為舊約始終不可能逐字去看.
所以只能夠整段去看.
那個大概的輪廓.
新約就很不同了.
因為新約是書信的形式.
就真的要逐字去了解.
甚至乎譬如保羅的書信.
他用的字和那個時代有什麼關係.
譬如很快我們講過.
《肥勒比書》四章六節.
「出人意外的平安」.
「必在基督耶穌裡面保守你們的心懷意念」.
「保守」那個字本身就是駐軍.
這個Garrison那個字.
本身就是羅馬的駐軍.
所以那個「保守」原來有軍事性.
原來神的平安是會保守我們的心懷意念.
就是說神的駐軍.
神的平安.
神的同在在我們的內心.
如果你懂得原來保羅用的希臘文那個字.
他刻意這樣用的時候.
你就會知道真的很不同.

$^{121}$一來他就是坐牢.
見到羅馬的兵兵.
守衛著他.
不讓他到處走.
保羅可以用他自己見到的情景,遭遇.
來講到基督耶穌的平安.
保守我們.
你就會很不同.
這些就是Work Study.
字詞的研究.
在那個時代.
保羅的時代.
他們是怎樣用那些字呢?.
因為保羅一定是在希臘文的世界.
當時的羅馬帝國的時代.
他是一個精通希臘文的專家.
一個拉比.
就是說他運用那些字.
一定有他自己的意思,目的.
所以你就會看到原來在舊約,新約.
有時你對原文有些認識.
你就能夠看到他們背後寫作的動機.
以及寫的優美等等.
很快又想起一個.
《約翰福音》十七章.
耶穌的禱告.
那個「叫你們」出現很多次.
就是說如果作者刻意運用那些.
「沉醉」的詞彙.
你就不要覺得他是.
「無不識丁」.
又或者是他.
為甚麼會那麼「枯燥乏味」.
原來可能變成.
就是神學上的一個用意.
不是他不懂運用那些詞彙.
而是他刻意在神學上.
強調再強調那個真理.
這個就變成我們有機會.
再碰到這些經文的時候.

$^{161}$我們就能夠.
所以大家可能開始.
從和合本開始.
再看看英文聖經.
英文比較好的就是.
新美國際版.
即是New American Standard Bible.
即是叫NASB.
或者是.
如果你看到一些比較古老的英文.
就是King James.
英皇音定本.
這兩個比較權威一點.
那個翻譯比較嚴謹.
有些就是異譯的.
不是直譯的.
大家就可以去想一下.
當然如果有機會.
真的可以讀原文就更加好.
這個就題外話.
讓大家知道一下.
我們說了.
9到14章.
不再是那些比喻了.
是一些預言.
有幾個大題目的.
耶和華是那一位的牧者.
撒迦利亞.
曾經要扮演那個牧羊人.
他要去牧養那些群羊.
但是得不到那些羊的滿意.
甚至乎是會被排擠.
覺得他沒有盡力.
還有他的工作不被討好.
甚至乎要被人遣散.
30塊錢.
所以你看到.
撒迦利亞書原來在福音書裡面.
是引述最多的舊約經文.
在福音書引舊約.

$^{201}$如果引先知書.
撒迦利亞書引很多.
你就可以看到.
彌賽亞的形象.
去到福音書的時候.
將撒迦利亞所描繪的那個工人.
那個牧羊人.
呼之欲出.
那個君王騎著驢驅進城.
九章九節.
是撒迦利亞書.
你可以想像得到.
在九章到十四章.
是有一個近乎很清楚的形象.
就是彌賽亞的兩次降臨.
所以十四章.
最後這一章.
你可以說是整個啟示錄.
都會作為那個藍本.
你可以說是那個題材.
來將這些天啟文學.
好像真的照搬一樣.
我們要看第十四章.
所以再說.
撒迦利亞書其中難讀的.
很多經文很難讀的.
但是其實分段就很簡單.
一到八章大概是用比喻.
八個的比喻.
來講述當時的情況.
和主回來之前.
怎樣來到某程度叫收拾以色列.
其實就是使他們重新有一個歸回.
一個復興.
所以撒迦利亞書是應該讀的.
因為是很正面的.
是很亮麗的將來.
但是那個最後的將來之前.
就有很多的爭戰.
很多的審判等等.

$^{241}$去到第九到十四章.
這六章又很對稱.
就是講到更遠的將來.
更遠的將來耶路撒冷的德性.
那個光景.
和將來神.
這些又是福音書和新約都出現的.
要開一個洗罪的泉源.
使到以色列人.
是從心底裡面悔改.
我們上個星期講過.
不同的家.
以色列家.
猶大家.
變雅敏.
即是名門望族.
祭司家族.
接著還有那個侍妹家族.
即是永遠都被唾棄的侍妹家族.
都能夠得到那個悔改.
個個都是男的獨在一處.
女的獨在一處.
大家都是不會再介懷自己.
那個生活的好與壞.
很緊張就是.
希望能夠去到神那裡.
得到那個的恩典.
其中就是那個收納.
即是有神的那個赦罪的恩典.
所以十三章一節.
就立即講.
神將來會開一個洗罪的泉源.
所以約翰福音十九章.
我們上個星期引述過.
耶穌基督被用長矛.
這樣去戳穿他的胸膛.
是十三章裡面提到的.
是那個叫做.
讀回那個字.
窒.

$^{281}$完全十九章的約翰福音.
就是講這個字.
窒.
即是代表著一個受傷的牧者.
在撒迦利亞書已經有的.
給我們見到.
原來主耶穌將來.
真的要好像對號入座那樣.
是被那些冰釘來刺穿他的胸膛.
所以有詩歌告訴我們.
流出的血與水.
就好像一個泉源那樣.
由主的胸膛流出.
洗乾淨我們的罪惡.
這個洗罪泉源.
原來舊約已經有了.
撒迦利亞書已經預言.
將來神會為我們開.
所以原來就是他的愛子耶穌.
十四章就更加震撼了.
因為就真的是世界末日了.
世界末日之後.
那個新的天地了.
所以你可以想像得到.
撒迦利亞書所涵蓋的那個將來的終極.
真的是最後的終極.
就是神作王在地上.
這些就是啟示錄的意境.
十四章第一節.
耶和華的日子臨近.
你的財物必被搶掠.
在你中間分散.
因為我必聚集.
萬國與耶路撒冷爭戰.
城必被攻取.
房屋被搶奪.
婦女被玷污.
城中的民一半被擄去.
剩下的民仍在城中.
不自剪除.

$^{321}$那時耶和華必出去.
與那些國爭戰.
好像從前爭戰一樣.
那日他的腳必站在耶路撒冷前面.
朝東的橄欖山上.
這些有提過的.
耶穌回來應許的.
祂說將來我由哪裡升天.
就降臨哪裡.
如果我們真的按字面去理解.
就是撒迦利亞書的那個意境.
所以如果我們.
我們下一次就要講末世論.
講貼薩諾尼加前後書.
加上講末世論.
有幾派的.
如果我們根據字面的經文應驗.
耶穌就是要在橄欖山再回來.
這裡告訴我們.
這個就是這樣的意思.
「這山必從中間分裂.
自東至西成為極大的谷.
山的一半向北來移.
一半向南來移」.
即是開一條好像黃道一樣.
使到神的選民可以經過.
「你們要從我山的谷中逃跑.
因為山谷必延到阿薩.
你們逃跑.
必如猶大王烏西瓦年間的人.
逃避大地震一樣.
而和我的神必降臨.
有一切聖者同來」.
這些貼薩諾尼加前書四章.
就是有這樣的意境.
所以舊約的這些預言.
成為將來新約他們.
某程度上是他們的題材.
是他們所領受的默示.
因為將來主耶穌會和他的聖者降臨.

$^{361}$第六節.
那一天.
一直提到那一天.
完全我們不可能預測.
所以妄想不要預測.
不需要預測.
是沒有人可以預測的.
何時主再來是不能預測的.
我們只能夠警醒.
只能夠做好我們自己的人生.
盡力去傳福音.
我們自己成為一個聰明.
有見識的管家.
就是等候或者計算.
耶穌回來的時間.
因為耶穌從來都告訴我們.
透過比喻.
要我們做一個良善.
忠心.
有見識的管家.
就是那個等候了.
如果不是.
就變成一個又惡又懶的僕人.
所以這裡告訴我們.
那個大地震.
神的聖者來.
第六節.
那一天必沒有光.
三光必退縮.
那一天星星月亮太陽.
我們人類.
特別是舊約古代.
能夠看到的就是.
不會有街燈的.
晚上就是晚上.
是陽燭而已.
陽燭都沒有.
是油燈而已.
所以.
真的能夠看到的就是.

$^{401}$星星月亮太陽.
是沒有其他東西的.
不會有任何人造的.
發光的物體.
所以三光退縮.
就是說.
去到一個天地昏暗的日子.
那一天.
出現了很多次了.
不過這一段.
是講到那個審判.
主最終極回來之前.
那些的預兆.
那些的情況.
地上的混亂等等.
大自然都要辨識.
大自然都要改變.
所以這些是和末世論的經文有關的.
譬如《馬太福音》二十四章.
是主耶穌自己親自講.
祂回來之前.
那個世界的.
讀完這段我和大家稍微揭一揭.
因為一講到末世論.
就要幫大家組織那些經文.
在新約出現過的.
給大家知道一下.
或者坊間神學界.
有哪幾個理論.
對於末世論.
到底是災前,災中,災後的被提.
有沒有大災難.
諸如此類.
這些都是大家都需要知道的.
第七節.
那日必是耶和華所知道的.
不是白晝也不是黑夜.
到了晚上才有光明.
那日必有活水從耶路撒冷出來.
記住.

$^{441}$去到啟示錄二十二章.
就有生命河的水.
會結十二回的果.
即是神的寶座那裡流出.
這裡已經有這些概念.
一半往東海.
如果這樣理解.
東海可能是死海.
如果西海再西一點.
有可能就是地中海.
當然準確未必是這個意思.
要看怎樣理解.
東下都是如此.
即是不會結冰.
又不會乾旱.
那就更加預示那個生命河.
生命的活水.
其實由耶穌升天.
聖靈降臨.
耶穌已經說.
那個生命就由祂那裡開始.
我們一信主耶穌.
領受聖靈.
就在那個人的生命裡.
從祂腹中流出活水江河.
是說到人得到聖靈.
所以這是很寶貴的救恩.
第九節.
所以過了這些審判.
過了這些災難之後.
連自然界.
我們以為是生生不息的.
我們以為是沒可能是沒有的.
你有沒有想過8月15日是沒有太陽.
沒有月光.
你有沒有想過.
沒有的,從小出生到現在.
最多是有雲遮住.
一定有的.
就算下雨,它也會有.

$^{481}$會有緣間出來的.
如果有一天8月15日是全世界都沒有月光.
將來就是這樣了.
三光都退縮.
你就可以想像到.
整個天地神賦予我們這個.
能夠計算.
能夠預測.
能夠給我們.
來去有那個.
不同的節期.
和人生有不同的意義.
所以你看看人類.
寫了多少的詩.
來紀念這些大自然.
床前明月光.
多少這些詩詞.
讓我們人生都不會那麼枯燥乏味.
自然界都是這樣的.
都不是的.
你可以想像到.
我們要記住.
我希望我們大家都不在這裡.
在天堂.
在樂園看到這些地上的事.
如果將來你在這樣的情況.
你可以想像到.
那個太陽是會變成血紅色的.
還有那個熱力.
是會像燒烤一樣.
烤地上的人.
你可以想像到.
將來的那種痛苦.
都未計神的烈怒.
那個還是自然界的痛苦.
所以這一切.
所說的轉變.
那個劇變.
都未講到最終.
神要在地上掌權.

$^{521}$這個就是第八節那裡說的.
第九節那裡說的.
耶和華必作傳遞的王.
那日耶和華必為獨一無二的.
他的名也是獨一無二的.
就是再無偶像了.
現在就.
你說你有神.
你說你的耶穌.
我都不信他.
我都沒有東西.
將來就不會了.
現在就.
我們叫人信耶穌.
這個都挺好的.
都挺導人向善的.
不過我暫時都沒有需要.
還是很客氣對宗教.
至於未至於說罵.
嘲笑.
有時有些人會的.
刻意找基督徒的.
一些不好的.
事實真的找到.
那就不好.
但是問題就是.
將來獨一無二.
即是說再無偶像.
人只能夠面對.
一是這一位.
滿有恩典慈愛的主.
一是面對這一位.
憤怒審判的神.
不能再說我拜觀音.
對不起.
我不用.
我是伊斯蘭教的.
我不需要.
沒有了.
我是無神的.

$^{561}$不會再有這樣.
所以這個是.
將來每個人都要面對這一位.
做物主.
第十節.
傳遞從加巴直到耶路撒冷.
南方的臨門.
要變為阿拉巴.
耶路撒冷凝居高位.
因為本身就是一個山丘.
大家去過耶路撒冷就會知道.
本身就是一個.
石安山的圍繞.
然後就是一個高處.
大概好像都是.
由哪裡計算呢.
那個的.
突然間又想不起了.
我.
西羅瓦茨.
西羅瓦茨就是那個水平.
上去其實很粗的.
很斜的.
這樣上去.
就是那個耶路撒冷上面的構成.
所以你可以想像到.
現在將來.
它是高處.
即是說周圍都會被掃平.
而且所謂凝居高位.
就是因為很正常.
神在耶路撒冷.
自然就是那裡.
每個人都想留歸那個山.
每個人都想去.
尋求神的恩典.
我們曾經都講過.
應該是十一章.
還是十二章的經文.
告訴我們.

$^{601}$他們要去.
懇求神的恩典.
這個是.
將來就會出現.
這樣的一個情景.
你可以想像到.
將來.
耶路撒冷的顯耀.
就是說.
就是從汴雅門門.
到第一門之處.
又到閣門.
一直講那個舊城.
現在是在重建.
重建聖殿.
很多舊的城門.
其實都荒廢.
和破落.
你可以想像到.
撒迦利亞預言的.
不是現在的短暫.
進行的工程.
本來破壞了.
末期的聖殿.
現在重建.
撒迦利亞就是其中一個.
呼喚百姓起來重建.
的先知.
現在是.
大利烏王的第二年.
九月.
他們開始重建.
然後撒迦利亞得到預言.
得到默示.
但是撒迦利亞現在講的將來.
的聖城.
是神降臨.
掌權的聖城.
所以你可以想像到.
如果你是神的選民.

$^{641}$就很榮耀.
很大的福氣.
十一節.
人必住在其中.
不在有就坐.
耶路撒冷人必安然居住.
安然居住.
在舊約是很重要的.
是在講.
告訴你.
是神祝福的明證.
人人都坐在.
葡萄樹下.
無花果樹下.
來聊天.
來享樂.
等等.
真正的國泰民安.
在這裡告訴我們.
是真真正正的.
神同在帶來的.
安然居住.
都未夠我們說.
聖誕平安夜我們去看煙花.
看倒數.
那些都已經是反映了一種平安.
但是那些.
還是很密集.
有時可能都會在地鐵.
都會有人吵架.
那不是叫平安.
但是這些描述.
是真正講神同在帶來.
以色列的.
真正的傳聞.
整個國家.
的平安.
和福樂.
所以安然居住.
是很重要的祝福.

$^{681}$在舊約.
新約也是.
很多時候都是求平安.
神同在.
的福氣.
就是平安.
所以薩林是祝福.
願神的.
恩典與你同在.
就是薩林.
我們中國人不懂說希伯來文.
薩林就變成薩林.
所以薩林就不是.
這樣叫薩林.
是真正的平安.
願你得到神的平安.
如果用保羅講的.
就是保羅更加.
誇張神的.
出人意外.
一定是你想不到的.
出人意外的平安.
保守就是.
Garrison 駐軍的字.
進駐在你的.
心靈裡面.
所以我們要多些去祈禱.
多些去仰望神.
我們就會得到心靈的平安.
如果不是就會忐忑不安.
所以這個安然居住.
大家回去做功課.
其實你.
稍為上網找一些聖經的.
Apps.
就可以做到.
舊約出現多少次.
安然居住.
你就會覺得很燦爛.
很美麗.

$^{721}$不要小看.
你這樣懂得這句說話.
其實你看多了很多經文.
多過很多只是崇拜的.
弟兄姊妹.
你看看每一段說什麼.
每一段的應許是怎樣.
這些肯定.
去到新約.
神都是這樣的祝福.
所以一個基督徒.
如果他還是怕怕怕.
其實他還沒有.
享受到神.
的同在的平安.
所以我們要留意.
十二節.
耶和華用災殃攻擊.
那與耶路撒冷爭戰的.
列國人.
必是這樣.
他們兩腳站立的時候.
肉必消沒.
你就可以想像到你站在哪個陣營.
現在神要.
那個for節.
神要for以色列.
against以色列.
現在是.
for將來的.
預言.
你想想七十年的被擄.
就是神against.
以色列.
他們要被擄去到巴比倫.
現在將來神就是for.
以色列.
誰和以色列為敵.
你就知道.
你想想那個描述.

$^{761}$眼裡面那個框.
都乾別.
錯別字那個別.
洩在口中.
潰爛.
你生過妃子你就知道.
痛到你都.
痛不欲生.
那日耶和華必使.
到他們大大擾亂.
他們各人彼此.
抽著舉手攻擊.
猶大也必在耶路撒冷.
爭戰.
那時四圍各國的財物.
到食許多金銀衣服.
必被收罪.
那臨到馬匹雷子落駝爐.
和營中一切生畜.
亦災殃,事與那災殃一般.
所有來攻擊耶路撒冷.
列國中聖下的人.
必連連上來.
敬拜大君王萬君之耶和華.
並守住彭節.
某程度是.
最開心.
最豐收的那個節期.
就是住彭節.
經過了贖罪日.
贖罪日就是猶太人的節期.
7月初十.
到7月15日就是他們的住彭節.
守八日.
到7月21日就是最大.
最大那一天.
住彭節.
他們要在西羅瓦茨.
打水.
上去聖殿.

$^{801}$倒水.
他們要這樣來.
代表住.
代表住活水.
西羅瓦茨你們知道.
在所羅門.
那個隧道.
那個引水道.
出來就是.
西羅瓦茨.
那個水tum .
因為已經缺水.
所以那個西羅瓦茨現在已經.
很缺水.
但是在耶穌的年代.
西羅瓦茨都是很重要的.
他們祭司.
在住彭節就會打水.
上去耶路撒冷.
一路唱著詩篇.
上去耶路撒冷.
耶穌很諷刺.
在住彭節的第八日.
約翰福音第七章.
人如果學了.
可以到我者裡來.
從他福中就可以流出.
活水江河.
然後那些門徒就想起.
耶穌說這番說話.
其實是說到聖靈.
在榮耀.
所以聖靈就未赤下下來.
所以是很有趣.
很諷刺.
那些祭司很開心.
打了水之後.
好像淡跳.
淡上去耶路撒冷.
倒水到聖殿.

$^{841}$他們有些儲水的地方.
代表著這個國家.
有一個風聖.
但是耶穌說的.
真正的風聖.
所以這就告訴我們.
到時萬國都要守住彭節.
但其實我們已經守住.
我們已經在福樂裡.
可以和大家再揭這些經文.
新約的約翰福音.
耶穌自己說的.
所以如果將來萬國要守住彭節.
其實我們有沒有享受.
耶穌在我們生命裡.
給予我們的這個節期呢.
第七章第十節.
這個節期大部分人都認為是住彭節.
第十節.
七章的十節約翰福音.
但他弟兄上去之後.
他也上去過節.
「不是明去,似乎是暗去」.
正在節期.
猶太人尋找耶穌說.
他在那裡.
眾人為他紛紛議論.
有些人說他是好人.
有些人說不是.
他是迷惑眾人的.
只是沒有人明明去講論.
因為怕猶太人.
因為已經引起了很多人的關注.
到十四節.
去到節期的時候.
耶穌上殿裡去教訓人.
猶太人就稀奇說.
這個人沒有學過.
怎麼明白書呢.
耶穌說.

$^{881}$我的教訓不是我自己的.
乃是那一位差我來的.
人如果立志遵著他的旨意行.
就必曉得.
這教訓或是出於神.
或是我憑著自己說的.
所以人要立志遵著他的旨意.
你聽完.
當故事聽完.
你不去遵著去行.
其實救恩是由你的知識過了.
就過了.
因為你停留不到在你的心靈裡.
因為你還沒打開你的心門.
去認信和接受.
接受是有一個行動.
就是要遵行的.
就是說.
你如果不是這樣的時候.
你聽很多.
那些道理可能會妨礙你的.
所以這裡告訴我們.
如果人有這樣的態度.
就會曉得.
耶穌就是神所差遣而來的救主.
和會對他有一個另外的眼光.
來接受他.
我和你的救主.
這裡耶穌自己說了.
人憑著自己說.
是求自己的榮耀.
唯有求那差他來者的榮耀.
就是耶穌.
這人是真的.
在他心裡沒有不二.
摩西豈不是全律法給你們嗎?.
你們卻沒有一個人守律法.
為什麼想要殺我呢?.
這些全部我們之前說過了.
說《約翰福音》的時候.

$^{921}$大家如果有機會.
你回去再重新聽.
但是重點是在37節這裡.
節期的末日.
就是第八天的鑄坪節.
就是最大之日.
耶穌站著高聲說.
人若喝了.
可以到我者裡來喝.
兩個字是不同的.
三點水是名詞.
口字邊喝下去就是動詞.
所以可以到我者裡來喝.
所以喝是一個行動.
首先你感覺乾喝.
即是你覺得自己有罪.
覺得自己沒有平安.
覺得自己怕將來的死亡.
有沒有死亡.
有沒有永生.
你開始問這些問題.
你開始對心靈對靈魂有一個關注.
不是人一世物一世.
而是有來生有永生.
你就要問這些問題.
有一個饑渴的時候.
你看完戲玩完.
為什麼你內心都有點不滿足.
你有一些空虛.
你有一些其他的問題解決不了.
這個往往就是尋找真神真理的開始.
所以人若喝了.
你看到那個步驟.
是福音.
所以耶穌說的.
你發現自己內心有問題.
這個就是福音的開始.
如果說你的耶穌很好.
我也很贊成你們說的.
道人向善很好.

$^{961}$不過我沒有需要.
和我沒有關.
那你就是拒絕福音.
他看不到自己是渴.
所以教會是要令到人渴的.
因為你不渴.
你就不會去到耶穌那裡.
來喝祂的活水.
所以耶穌說了幾個步驟.
幾個層次.
可以到我者裡來喝.
信我的人.
喝和信是同義詞.
口渴要喝水.
心靈的枯竭.
是需要靈魂的醫生.
是需要靈魂的食物.
就是需要聖靈.
需要救主.
需要救恩.
所以我們記得.
我們信了耶穌.
繼續要在主耶穌裡面.
祂才是主角.
我們不要慢慢信了耶穌.
享受的就是人和人的團契.
就沒有了耶穌的救恩.
其實是大問題.
我們圍著的是.
耶穌給的恩典.
對這位賜恩典的神.
毫無興趣.
慢慢地疏遠祂.
繼續有團契.
我們有教會團契.
真是搞笑到不得了.
絕對禁止.
不應該.
我們不可能沒有了這位神.
而在祂圍繞著祂給的禮物.

$^{1001}$我們追逐那些禮物.
說出來好像很好笑.
但其實是很悲哀的.
因為永遠是在說.
信這位耶穌.
這裡告訴我們.
信我的人就如經上所說.
經上所說就是.
《二創亞書》58章的應許.
所以揭下揭下.
有些已經解經.
就是這樣的重點.
從他福中要流出活水的江河來.
耶穌者說.
是指著信他之人.
要受聖靈說的.
他們是馬後炮.
事後孔明.
門徒在耶穌復活之後寫的.
當然會知道.
那個時候耶穌說這段經文的時候.
是未有聖靈的.
是未釘十字架的.
所以當他這樣說的時候.
肯定了一件事.
所有信耶穌的人.
在耶穌升天之後.
信耶穌的人.
歷世歷代.
古往今來.
包括什麼民族.
一信耶穌.
立即獲得的.
不是洗禮證.
是聖靈住在你的心裡.
你就得到這個活水江河.
這個是永遠和我們同在的.
所以你和我都有聖靈.
所以我們要讓聖靈在我們內心.
一路引導我們.

$^{1041}$我們要遵行神的旨意等等.
這些告訴我們.
是受聖靈所說的.
那時還沒有賜下聖靈.
因為耶穌尚未得著榮耀.
所以我們看釘十字架.
原來是耶穌得榮耀.
《肉眼福音》就強調這個主題.
我們說過了.
所以這裡告訴我們.
很重要的就是.
這個鑄棚節.
本來猶太人守鑄棚節.
就是享受那個付祝.
一來是秋收.
封城.
鑄棚.
紀念他們出埃及做奴隸.
搭棚.
但是更重要就是.
如是封城.
所以耶穌相當搞笑.
你們去打活水.
西羅馬祠打水.
我就是那個活水.
我才能夠提供那個活水給你.
我就是可以令到你.
透過相信我的救恩.
得到聖靈.
所以這裡.
回到撒迦利亞書.
將來全世界守鑄棚節.
當然全世界是願意歸服神.
站在神那邊的人.
如果你是反對神的.
那你就是那些仇敵.
就是要面對上文所說的審判.
所以是有一個分水嶺.
你想想.
如果在新約.

$^{1081}$保羅說得很盡的一件事.
就是說.
神若幫助我們.
誰能抵擋我們呢?.
羅馬書八章.
應該是三十節還是三十一節.
那個英文你已經看到.
不是神成為你的幫手.
不是.
英文就是說.
God is for you.
祂是站在你這一邊.
祂不是against you.
就是說.
我們要記住保羅的教導.
都揭曉一下.
大家可能未必很熟悉.
羅馬書八章那段的頌歌.
保羅說得很盡的.
進到一個地步.
你好像覺得他有沒有誇張了.
其實使徒他們真的有這種.
內心的那種凝訊.
那種眼光.
由十八節開始看.
是先從那個目前.
整個宇宙.
整個世界.
落在一種虛空底下.
因為什麼叫虛空底下呢?.
就是創世記的咒詛.
萬物受咒詛.
因為始祖犯罪.
本來神的創造是完美的.
是沒有罪的.
所以保羅就說回這個.
沒有辦法的.
你一出生來到這個世界.
你不能再怨天尤人.
因為我們的祖宗阿當.

$^{1121}$夏娃做了這件事.
吃了分別善惡樹果.
所以萬物受咒詛.
是不關你和我的事.
但是我們就生下來.
一個受咒詛的世界.
要等候什麼呢?.
就是等候那個真正的.
神的那個時間.
為什麼要引這些經文呢?.
因為十四章後半段.
正正就是講這個新天新地.
真正沒有咒詛的天地.
十八節.
八章的十八節羅馬書.
這些經文真的要懂得背.
羅馬書我也講過.
第三章到第八章.
要反覆反覆地讀.
因為保羅在爭論一些真理.
他要將一個恩順清義.
反反覆覆講明白.
給羅馬的信徒知道.
將整個恩順清義的原理.
怎樣去解釋恩順清義.
所以如果做福音工作的.
初信班的導師.
其實一定要熟.
要很熟.
隨時可以引導那些經文.
知道保羅在爭論什麼.
他好像有一個對手.
和他爭論一樣.
在那邊辯論.
他要將那個真理.
徹底地帶出來.
去到第八章.
就是那個最高峰.
十八節.
我想.

$^{1161}$他在那裡.
來表達他內心那種.
對真理的共鳴.
我想.
現在的苦楚.
不是做事被老闆罵這麼簡單.
那是其中一種.
很少的一種.
他在講的是萬物在那個就坐底下.
如果神沒有改變.
一直人生都在一個就坐底下.
即是傳道書那個境界的虛空.
若比將來要顯於我們的榮耀.
不足介意了.
你想想他這樣的心境.
不是去做剃度一年.
修行一年.
而是他真的說.
你如果真的認識神的那個就坐.
人類在這種虛空底下.
千千幾千年這樣來到去.
不斷.
那些人是會過去.
很小的時候又會離開世界.
不同的方式離開世界.
你會感受到那種虛空那種絕望.
人類是感受到的.
所以人類是不敢再提死亡的.
很避諱的.
但是保羅很有那種豁達的眼光.
因為他有救恩.
有耶穌.
有神.
他說如果比起將來要顯於我們基督徒的榮耀.
現在的苦楚不足介意.
他不是說他自己被人打.
受審那些那麼簡單.
是說眾生在那個就坐底下.
在神的就坐底下.
留意第十節已經說了這件事.

$^{1201}$對不起,十九節.
受造之物,切忘等候神的眾子顯出來.
因為受造之物伏在虛空之下.
根本就是在說萬物在這個虛空.
就是神的就坐.
當夏娃被趕出伊甸園的開始.
就是那個就坐.
不是自己願意,乃是因哪叫他如此.
就是神令到萬物受就坐.
地球,地球的人.
是神創造的人類.
但十一節.
但受造之物仍然指望脫離敗壞的核際.
得享神而屢自由的榮耀.
似乎說的是天使,靈界.
那種自由,那種榮耀.
將來我們去到將來天堂.
沒有人再犯罪,不會再犯罪.
那句說話是真的假的.
你不用再問了,每個都是真話.
將來天堂是沒有人犯罪的.
我都想見識一下.
將來去到天堂不用牧羊.
根本不需要傳道人.
每個人都信主.
你可以想像將來有這樣的世界.
試過有些地方.
整條村信主.
甚至乎要成為他們自己的法律.
但始終都不行,沒用的.
因為地球上的人.
就算信主都是有罪性.
就算你製造了那些規律.
慢慢都會有人打破.
你就可以想像到將來的榮耀.
神而屢自由的榮耀.
第22節.
我們知道一切受造之物.
一同嘆息勞苦.
直到如今.

$^{1241}$將一個很現實.
很真實的人類的光景說出來.
完全不是負面.
而是很正面帶進他的榮耀的理論.
第23節.
不但如此.
就是我們者有聖靈初見果子.
信了主耶穌的基督徒.
包括保羅自己.
也是自己心裡嘆息.
等候得著兒子的名分.
乃是我們的身體得贖.
你可以想像到.
將來那個榮耀的身體.
是怎樣奇妙的呢.
我們火化都不怕了.
將來會有一個榮耀的身體.
你可以想像到.
那個人類的得救贖.
是真正完全的救贖.
身體得贖.
保羅這裡說得很清楚.
他確信我們將來會復活的.
第24節.
我們得救是在乎盼望.
只是所見的盼望不是盼望.
你如果見到的就不是了.
就是說還沒發生.
還沒完全的.
Fully的Fulfill的.
那個叫將來的盼望.
所見的.
見到的就不是盼望了.
見到就是見到了.
誰會盼望所見到的呢.
第25節.
但我們若盼望那所不見的.
就必忍耐等候.
這個又是說一個態度.
所以我們要渴慕主再來.

$^{1281}$我們的渴慕.
是應該化成為忍耐等候.
忍耐等候.
就是說有個對象.
你是期望他來的.
所以我們要有這個心靈的空間.
偶然間想一下.
他回來.
你要見他.
那個裝扮.
那個忍耐.
那個等候.
就是說我們的地上生活.
是應該有要求對自己.
他回來.
我是不是裝飾整齊呢.
福音書有這些比喻.
忍耐等候是保羅自己說的.
所以既然我們盼望那個日子.
我們就應該忍耐等候.
等候就是說.
你是不斷有準備.
也是回應耶穌所說的.
不是闊綽濫理.
是真的做一個良善忠心.
還要是有見識的管家.
我們要準備我們自己.
營建我們的主耶穌.
所以是必忍耐等候.
所以如果我們完全沒有一些.
occupy你的人生.
有一些的事奉.
有一些的擔心.
有一些的參與.
有對神的國度.
對地上現在越來越嚴重的罪惡.
你都完全沒有任何的介意.
你在你的內心完全是沒有的.
我都很懷疑你是否信耶穌.
甚至你有沒有神的心腸.

$^{1321}$這裡說要忍耐.
就是說教會可能受逼迫.
很多困難.
我們要一起面對.
這就是忍耐.
主耶穌還沒回來.
無論現在穆斯林的群體.
越來越多人信穆斯林.
不信耶穌了.
或者歐洲越來越少人信耶穌了.
教堂已經不再是敬拜神了.
你都完全沒有任何的感覺.
我不知道現在歐洲人他們有什麼感覺.
所有教堂已經不再是.
拿來做禮拜堂.
做咖啡室.
如果將來錦繡堂.
又是亂說的.
沒有的.
拿來做麻將館.
抽水的.
用來維持的.
令到風琴不要生鏽.
入場要抽水.
打四圈就多少錢.
你想想.
然後你看到.
是好的.
可以維持教堂.
裝修不用這麼多錢.
這個是悲哀.
不是很好.
我們要靠這些收入.
神都不喜悅的收入.
來維持.
而信徒根本不再回教會.
不再來奉獻.
這個是悲哀.
不是應該高興.
所以忍耐等候.

$^{1361}$是說我們參與在神的國度裡面.
你自己的內心.
你的觸覺.
所以如果你的觸覺已經很不敏銳.
求神給我們眼淚.
那個心靈.
能夠有那份敏銳的心靈.
所以這裡就是.
保羅告訴我們.
如果我們等候那個未見的盼望.
就必定會忍耐等候.
然後還要說下面.
第26節.
況且我們的軟弱.
即是說他現在包底.
不是這麼硬.
原來有聖靈幫我們.
在現在這個忍耐等候的過程.
我們有力量去面對.
不是一種灰心絕望.
不是被環境打沉.
而是我們是有盼望的.
所以如果歐洲有些老的信徒.
仍然為自己的國家祈禱.
這些人就是忍耐等候.
雖然越來越少.
但是這些敬虔的信徒.
我不理周圍的環境.
我仍然去敬拜神.
這些就是一個忍耐等候.
所以我們求主給我們能夠.
祈禱在我們的內心.
越來越真誠.
越來越敏銳神的心意.
這就告訴我們.
然後就說.
況且.
所以不要灰心.
況且我們的軟弱.
都有聖靈幫助.

$^{1401}$我們本不曉得.
當怎樣禱告.
一說到忍耐等候.
就說到我們要閉上眼.
仰望神.
即是要有神的眼光.
我們如果不閉上眼.
怎樣有神的眼光呢?.
教會如果全族人都不祈禱.
教會的領袖都不祈禱.
不懂祈禱.
沒有興趣祈禱.
那就困難了.
因為我們不覺得我們有軟弱.
保羅說.
我們的軟弱就是不懂祈禱.
我們有聖靈教導我們.
這裡說我們本不曉得.
當怎樣禱告.
保羅不是吧?.
How to play你都說不懂.
有沒有搞錯?.
你傳道牧師.
你是使徒.
那個應該說.
在神的心意裡禱告.
不是那些字句.
不是在裝飾那些很漂亮的字.
而是神在這件事.
有什麼心意.
你要這樣揣摩.
神的心意.
神的眼光.
這裡告訴我們.
只是聖靈親自用說不出來的嘆息.
替我們禱告.
即是當我們貼近神的裡.
聖靈就會用他說不出來的嘆息.
說不出來即是我們都不知道.
保羅都不知道.

$^{1441}$就為了我們祈禱.
使到原來神的心意就可以成就.
27節.
監察人心的.
即是神.
曉得勝靈的意思.
因為聖靈是照著神的旨意.
去替聖徒祈求.
真的很玄妙.
超過了連保羅都未必會明白.
但他講得出來.
即是說聖靈用人說不出來的嘆息.
人是不可以講得出來的.
但神就監察人心的.
會曉得勝靈的祈禱.
因為聖靈就是照著神的旨意.
照著的是神的旨意.
不是我們的想法.
這樣就成就.
所以第28節.
又是金句裡的金句.
因為是我生日.
我們曉得萬事互相效力.
叫愛神的人得益處.
就是按他旨意被召的人.
這句說話如果就這樣讀就很難解.
因為上下文有一個張力.
一直在講不懂得禱告.
聖靈會禱告.
神會因為聖靈的嘆息來成就.
因為聖靈是照著神的旨意祈禱.
但28節告訴我們.
萬事在神旨意成就的大前提.
有些是我們根本不知道的環境.
因為是萬事.
不是一件事.
是所有東西互相效力.
這些就超過我們人能夠想的.
你想想如果我們只是一隻螞蟻.
我們怎會知道人類世界的複雜呢.

$^{1481}$我們怎會知道.
譬如你會駕車.
你去到很複雜的高架天橋.
很複雜的.
你背後想想.
那些工程師真的很厲害.
他們想到那麼多的設計.
那些還是很簡單的東西.
現在是人類萬人萬事.
在互相為神的旨意效力.
這樣來達至.
這超過我們每一個人的想法.
所以我們即管祈禱.
因為我們沒有人百分之一百知道神的旨意.
我們即管祈禱.
但神成就一定會過於我們所求.
這就是28節的大概意思.
我們曉得萬事都互相效力.
最終是甚麼呢.
叫愛神的人得到益處.
這亦是聖靈的禱告裡面的.
神旨意成就的那件事.
使得這些人.
就是按他旨意被召的人.
因為他預先所知道的人.
就預先定下效法他兒子的模樣.
使他兒子在許多的弟兄中作長子.
主耶穌就是我們的長子.
在神的救恩裡面.
我們今天慢慢的.
是一個模樣.
我們應該是越來越像主耶穌.
當然是說生命.
說屬靈的特質.
來有一個這樣的安排.
所以30節.
預先所定下的人.
又召他們來.
我們是被召得救恩的人.
所召來的人又稱他們為義.

$^{1521}$我們在神的眼中.
是已經被算為公義無罪.
神的義已經加在我們身上.
是藉著耶穌的寶血.
祂的救恩.
我們已經得到神的義.
接著所稱為義的人.
又召他們得榮耀.
即是說我們一生要榮耀神.
然後最終我們的生命.
我們整個人生.
是在神眼中是一個榮耀.
所以這個就是萬事互相效力.
這個就是聖靈在我們生命.
有生之年為我們祈禱.
是照著神在我們人生的旨意禱告.
不是照著你的喜好禱告.
是神的聖靈照著神的旨意.
神要在我們的生命裡面.
什麼的生命.
就是你被召.
你得救.
你被稱義.
你要將來得榮耀.
這樣的照明人生.
這些也是很重要的一個領袖訓練.
必定要教導的.
就是你有沒有想過你的人生.
你可以在神裡面得榮耀.
這個得榮耀.
就不是每個人都很有錢.
很健康.
很開心.
不是這樣.
每個人的人生都不同.
有些我們認為是悲哀的人生.
可以是一個在神眼中榮耀的人生.
接著是三十一節.
其實我想講這一句.
既是這樣.

$^{1561}$還有什麼要講呢.
既是這樣.
就是說.
你是被召的.
你是得救的.
你是神稱義的.
你將來要得榮耀的.
這個是神在你生命裡面的旨意.
保羅的意思就是說.
既是這樣.
本來我們每個人都是墮落.
在咒詛底下.
在虛空底下的.
但是現在主耶穌的救恩.
給了我們一個這麼宏大的永恆.
人生的歸宿.
是肯定可以.
你想想.
如果你新年.
舊曆年.
你已經訂了日本的酒店.
去到你只是給一個護照去看.
你就肯定可以入房了.
跟你完全沒有的.
到時是沒有客店給你的.
你是有兩樣的心情的.
你可以想像一下.
現在保羅告訴我們.
我們一定有永恆.
一定回到天堂.
接著告訴我們.
既然現在是這樣的時候.
我們還可以講什麼呢.
就是這裡.
神若幫助我們.
誰還能夠抵擋我們呢.
這裡的意思不是幫助.
中文翻譯不了那個味道.
God is for you.
神站在你這一邊.

$^{1601}$完全站在我們這一邊.
這樣的意思.
誰還能夠抵擋我們呢.
他現在說的不只是該殺尼祿王.
他說神不愛惜他自己的兒子.
為我們眾人捨去.
豈不是把萬物.
和他一同百百的賜給我們.
我們很喜歡萬物.
不喜歡耶穌.
最終連萬物都會沒有.
記住是神把萬物和他和他的愛子.
百百賜給我們.
所以一直說的是.
有主耶穌就有一切.
如果沒有.
你所擁有的.
慢慢都會變成枯燥乏味.
會變成過去.
所以我們不要搞錯了.
要了那個妹仔.
不要那個主人婆.
有時變成我們本末倒置了.
慢慢迅速轉變就是這樣.
翻崇拜.
享受那些音樂.
享受那個風琴.
享受那個講道好笑.
但原來不享受和耶穌的同在.
帶來的心靈滿足.
對他的敬拜.
對他的降服.
對他又多一種的認識.
如果沒有了這個最重要的.
根本我們不是朝見神.
不是去敬拜神.
我們是敬拜崇拜氣氛.
這樣就很大件事了.
不知不覺就落了一個偶像敬拜.
所以一定要追求.

$^{1641}$稍為彈錯了.
罵到那個人死.
慢慢不覺得自己是.
參加崇拜敬拜神的敬拜者.
弟兄姊妹.
以為自己來消費.
不如你去看一場好聽的音樂會.
找個最出色的.
全世界的高音去聽.
到時你付了錢.
你去罵他.
你不是剛剛好唱成這樣嗎.
這樣就OK.
但教會不是這樣.
教會是弟兄姊妹.
我們是重聚人.
蒙了恩典.
一起在這裡事奉神.
所以那個態度價值觀.
和我們對崇拜的主角.
是有密切的關係.
作為我們辦事的人.
每件事要很well prepared.
是必須的.
我們不斷會每個星期檢討.
有怎樣的安排.
但是參加者.
記住我們不是來去欣賞崇拜.
我們不是來去買服務.
我們是敬拜神.
我們是神的兒女.
我們有福氣.
來到神的家敬拜祂.
所以很多事.
你不需要再斤斤計較.
這些態度是真的不要得的.
如果我們一直帶著這樣的心情.
無論現在在現場的.
在網上的.
將來聽到這些經文的.

$^{1681}$都真的要彼此學習.
彼此勉勵.
我們不是去購買服務.
我們是來朝見神.
所以你一定要.
無論你走路來.
開車來.
差不多你準備來到教堂.
你要課自己.
我現在是來朝拜.
敬拜真神.
重點是這一點.
當然你有你的意見.
某些事可以做得好一點.
完全沒有問題的.
對的.
因為我們正正是.
大家是神的家裡的弟兄姊妹.
是可以彼此去扶持.
彼此去勸勉的.
但是你用一種挑剔.
這樣就很大件事了.
你如果日積月累地挑剔.
你有沒有搞錯.
神都會挑剔你.
這樣就大件事了.
你挑剔人.
神挑剔你.
是很嚴重的.
你和神計較.
神不和弟兄姊妹計較.
因為他已經盡了他的力.
盡了他的所能.
來去服侍弟兄姊妹.
然後你和那個弟兄姊妹計較.
神就和你計較了.
這些是很大的罪.
當然我不是拆開說.
而是這裡告訴我們.
神既然都不將他自己的愛子來去顧惜.

$^{1721}$豈不是他會將萬物和他.
一同白白賜給我們嗎.
誰可以控告神已經所揀選的人.
就是我們.
有神稱他們為義.
再一次確認這件事.
我們被揀選.
我們被稱義的.
誰能定他們的罪.
由基督耶穌已經死了.
而且從死裡復活.
這段聖經.
馬丁路德就曾經有這樣的體驗.
他宗教改革的時候.
很多的君王.
教宗.
即是那些神父.
要追殺他.
他周圍歐洲流亡.
有些當時的城邦的王.
就包庇他.
讓他避免教廷對他的追殺.
下的追捕令.
他做夢.
想起.
有一個夢.
就是撒旦控告他.
說他自己這麼多罪.
你說什麼宗教改革.
什麼.
冷清高.
他當時.
突然間聖經.
聖靈的工作在他內心.
誰能控告神所揀選的人.
他突然間.
整個心靈.
有一段時間就不開心.
很困苦.
很困擾.

$^{1761}$因為周圍的壓力.
逼得他很緊.
他突然間這段經文.
釋放他.
讓他知道.
是的.
我是在神眼中.
沒錯.
我是一個罪人.
但是.
沒有人再可以控告我了.
因為基督耶穌已經為我死了.
這樣.
如果你去參觀馬迪路德的故居.
據聞有一個地方.
他當時拿起墨水.
他們用毛筆.
沾墨水.
來寫字.
他就拿起墨水瓶.
扔到魔鬼身上.
真的弄髒了.
這是一個當時的故事.
其實就是這段經文.
誰能控告神所揀選的人.
有基督耶穌已經死了.
而且從死裡復活.
所以這個就是說.
神站在我們這一邊.
我說得太遠.
我們回到撒加尼亞書.
十四章這段經文.
將來就是萬國萬族的人守住彭捷.
你可以想像得到.
將來好像天生地之前的.
如果啟示了說的千禧年.
就是這樣的光景.
萬族的人要懇求神的恩.
要拉住一個以色列人.
我剛才翻不出來.

$^{1801}$應該是第幾章.
我再看一看.
他們要懇求神的恩典.
在發起之前的.
不是這樣的.
第八章提過了.
有兩段是說懇求神的恩.
八章的二十一節.
這城的居民必到那城說.
我們要去懇求耶和華的恩.
尋求萬君之耶和華.
我也要去.
必有列邦的人.
這個二十二節.
八章撒加尼亞書.
和強國的民來到耶路撒冷.
尋求萬君之耶和華.
懇求耶和華的恩.
不斷提到.
然後說萬君之耶和華如此說.
在那些日子.
必有十個人從列國諸族中出來.
拉住一個猶太人的伊甘說.
我們要與你們同去.
因為我們聽見神與你們同在.
一而再說將來那個.
很榮耀的救贖.
以色列人得到很大的尊榮.
萬族都要來到去.
去到耶路撒冷.
來尋求神的恩典.
看到神已經在.
這個.
鑒欖山降臨.
然後坐著.
然後不會再有偶像.
你可以想像得到.
撒加尼亞書給我們的將來的情景.
是非常寶貴.
和真正的.

$^{1841}$神再回來的日子.
用不同的方式.
這裡再守住彭捷.
這裡又再告訴我們.
萬族都要來到去耶路撒冷.
十四章的十七節.
地上萬族中.
凡不上耶路撒冷敬拜大君王萬君之耶和華的.
必無雨降在他們的地上.
埃及族若不上來.
雨也不降在他們的地上.
凡不上來守住彭捷的列國人.
耶和華也必用者栽攻擊他們.
這就是埃及的刑罰.
和那不上來守住彭捷之列國的刑罰.
你可以想像得到.
這樣的描述.
只能夠是站在哪一邊.
你站在神這一邊.
你就享受神給的恩典.
那個福樂.
那個對神的認識.
和平安.
如果不是的話.
就好像我們說的.
「蠶豆都沒一粒給你吃」.
我們上一代是這樣說的.
類似這樣.
你想沾一點雨水都沒你的份.
要這樣的奏作去到這麼嚴重.
就是說將來就不會這樣了.
也就是說.
都有一些人.
很堅決的否定神.
拒絕神.
那就要面對神的審判了.
好像用啟示錄的經文就是說.
「唇一不擇」的憤怒.
「唇一不擇」.
就是不會加上恩典.

$^{1881}$現在我們有很多恩典.
很多機會.
很多不同的方式賜福給我們.
但是將來.
那個時候.
「唇一不擇」的審判的憤怒.
這裡就告訴我們.
將來列國.
不在這個「住鵬節」.
就面對很大的痛楚.
所以其實我們每一次能夠守聖餐.
真的要感恩.
真的每一次都要珍惜.
當然不應該這樣恐嚇大家.
我自己這樣學習.
每一次可能是最後一次.
你不知道的.
突然間真的要回天家.
每一次是最後一次.
其實是幫自己建立那份珍惜.
那份敏銳.
絕對沒有說.
人生去到很灰很絕望.
不是這個意思.
而是珍惜每一次的相聚.
每一次都是.
如果是最後一次.
可以是這個教會最後一次.
可以是我的最後一次.
可以是大家的.
Fellowship的最後一次.
因為有些人可能移民.
他一生都未必再見到.
諸如此類.
我們珍惜一起相聚的時間.
的聖餐.
的聚集.
其實每樣東西都是.
每一個聚集.
都大家珍惜.

$^{1921}$其實會激發起我們那份對神的敬虔.
如果我們都不珍惜.
其實每一個聚會.
對我們來說都是.
沉悶不過.
因為你根本都不珍惜.
如果珍惜.
你就不會覺得.
回到聚會是一件苦差.
應該是一件很開心.
很快樂的事.
尤其是我們去到一個年紀.
見到有下一代.
再下一代.
我們更加應該.
感覺是一份開心.
滿足.
見下小朋友.
為他們祈禱.
拍下他們的肩膀.
拍下他們的頭.
是一件很大的祝福.
所以我們應該要有這些心境.
我們應該要建立這些.
屬靈的胃口.
其實回到教會就很忙.
因為你周圍.
看哪裡有需要.
你就默默祈禱.
其實你的心靈就很活躍.
沒有人見到.
但你的心靈和神很活躍.
在運作中.
你就會很不同.
你整個人的投入.
那些詩歌特別好聽.
講到又特別明白.
因為你整個人投入.
那個敬拜裡面.
聖靈就會在你的內心.

$^{1961}$有感動.
有工作.
第二十節.
當那日.
馬的鈴鐺上.
必有歸耶和華為聖.
的這句話.
剛剛.
今天看新聞.
有隻貓.
綁住條帶.
是什麼Gucci.
二千多元.
名人就有很多名牌.
連隻貓都要有名牌.
在那裡繩.
好像很搞笑.
錢多到要這樣用.
但將來.
這裡就不是因為.
皇家什麼馬.
馬的鈴鐺.
馬都會響.
有些鈴鈴噹噹.
都必有歸耶和華為聖.
這句話.
耶和華殿內的鍋.
必如祭壇前的碗一樣.
凡耶路撒冷和猶大的鍋.
必歸萬君之耶和華為聖.
本身那些器皿.
就是裝聖物.
裝那些獻祭的玉.
本身就已經是要乾淨.
和要聖潔.
但更加歸耶和華為聖.
即是神的聖潔.
已經是臨離盡至.
在不同的景況裡.
即是說.

$^{2001}$一些我們覺得沒什麼顯眼的東西.
一些琴簇.
都有這種歸耶和華為聖.
鈴鐺不斷提醒你.
馬的鈴鐺.
跑馬地的馬就不會鈴鐺.
因為那麼吵.
影響跑馬.
會落鞭.
或者他們會穿什麼制服.
但神的馬.
在耶路撒冷將來.
那個鈴鐺提醒你.
歸耶和華為聖.
所以我們都要求主.
給我們有這種.
我們是不是神的僕人.
我們有這些鈴鐺.
人們見到我們就感受到一種.
神同在的平安.
神同在的榮耀和聖潔.
和歸神為聖.
歸神為聖.
我們這個人是專心跟隨神的.
他的心會真的為神的國度想.
集中去怎樣讓神的國度.
福音可以傳遞.
有鈴鐺.
所以其實真的.
有時你聞到一些人.
他有沒有那種信耶穌的趣.
你看李炳光牧師.
八十多歲.
他一句說話已經感動我了.
他說平安夜.
對,沒問題.
去開心吃大餐.
教會有崇拜.
教會有聚會.
師母就在他的教會.

$^{2041}$有那些聖學的崇拜.
他說來報道會講福音.
更加是他的快樂.
更加感動.
你會看到這個老人家.
他很實地將福音講出來.
完全不會覺得他很避忌.
怕誰得失誰.
他講明這個世界真的有救恩.
唯有耶穌才帶給你平安.
什麼樣的聖誕.
如果沒有就要追求得到.
其實是很重要.
你會感受到他那種對福音的熱誠.
就是這樣.
有鈴鐺.
弟兄姊妹我們有沒有那個鈴鐺呢.
走過就感受到.
這個人信耶穌.
很溫柔.
因為耶穌帶給他溫柔.
這個人有一份平安.
是主耶穌帶給他.
這個人有一個屬靈的智慧.
是耶穌帶給他的.
如果你繼續是這麼粗魯.
又很硬的人.
其實耶穌在你的生命裡面.
有一些雕刻.
因為歸耶和華為聖.
馬的鈴鐺.
所以憲制的必來取這個鍋.
主肉在其中.
當那日在萬軍之耶和華的殿裡面.
不會再有迦南人.
即外族.
即是沒有神的生命的外邦人.
即是說.
整段聖經是說一個.
你是屬神還是不屬神.

$^{2081}$因為將來那個新天新地一定會來.
主耶穌一定會再腳踏.
橄欖山分開.
是不是真是物質.
我就不知道.
但是至少有一派的人這樣理解.
至少將來主耶穌一定會回來.
是不是這樣的字面應驗.
那個未必是最重要.
但是我們今生是不是課神呢.
或者神是不是站在我們這一邊呢.
還有我們人生.
有沒有一種讓人感受到.
我們是有救恩.
有神的呢.
所以這些描述.
我們覺得很搞笑.
但是這是對將來的預言.
這個是他們猶太人古代的時候.
記住是古代.
主前的五百多年的時間.
撒迦利亞所領受的那個意象.
但是去到我們現在二千多年了.
我們的香港.
我們信主了.
但是有些大原則是不會改變的.
就是我們信主.
是不是我們有聖靈在我們裡面.
來引導我們去禱告呢.
我們對自己的夢召有沒有感動呢.
我們自己是被神揀選.
我們是被清二不再追究我們的罪.
還有我們追求的就是神的榮耀.
我們一生有沒有一個榮耀神的這個心呢.
很多信耶穌的人未必有.
未必懂得要這樣有.
原來我們信了主耶穌.
最大的人類最大的目標.
目的.
你的存在.

$^{2121}$就是要榮耀神.
這個是宗教改革第一個最重要的教條.
這個是嘉爾文的教理問答第一條.
就是問那些小朋友.
瑞士日內瓦的小朋友.
人類存在最重要的是什麼.
賺錢.
要有太太.
有房子.
不是.
是榮耀神.
你的人生.
中世紀的人是這樣想的.
宗教改革的時候.
要榮耀神.
所以這裡就告訴我們.
來到最後了.
十四章.
舊約已經預言.
所有東西都要榮耀神.
馬的叮噹都有個歸榮耀給神.
所以如果你家裡的很乖.
你很喜歡的寵物.
你都做個叮噹提醒自己.
不是提醒他.
他是不懂的.
提醒我們自己.
是要榮耀神.
所以都求神幫助我們.
我們很快的.
第九到十四章.
就是對將來列國.
對以色列將來的量力復興.
因為神的無論第一次降臨.
和第二次降臨.
都包在九到十四章裡面.
你覺得很奇妙.
為什麼舊約會這麼奇妙.
玲瓏浮突.
講主在來呢.

$^{2161}$福音書為何引這麼多.
原因就是這樣.
不過引的是第一次主降臨.
甚至是真是對號入座.
這位主是騎著驢驅進城的.
這位主是被猶大出賣.
是三十塊錢的.
這位主他自己是好牧人.
但不為他自己的國民所接納的.
好像這個撒迦利亞這樣.
他是那個的牧者.
但被人拒絕的.
提到主耶穌是被那個的矛去刺穿的.
就是這樣來將他.
一個受傷的牧者.
主耶穌.
這全部是第一次的.
十四章就真是預言.
彌賽亞第二次降臨.
就是這位主將來他要來到腳踏橄欖山.
所以盼望我們完這個撒迦利亞書.
大家有時間再回去重溫.
再看.
有些經文就很難解的.
但如果我們整體的段落.
其實就很清楚的.
本來今天都預備了馬拉基書.
但我們可以下一個星期再看.
提一提.
因為接下來的一月.
天崇就需要用禮堂.
他們星期四晚就要來去練師.
因為將來天崇就在這裡做崇拜.
星期六的下午.
所以我們逢星期四晚的查經班.
跟網上的沒有關係的.
因為到時一樣是直播的.
但如果回現場的.
我們就在203那裡去查經班.
和會直播.

$^{2201}$都是用Zoom的直播.
所以對大家就沒有影響的.
都紀念.
因為雄恩那邊開始有很多事要跟.
都比較忙的.
盼望我都有兩邊的心力去做.
大家都為雄恩堂祈禱.
很多事工要開展.
和有些事都需要一些更新.
盼望全個教會.
無論是錦宣.
佳恩樓.
一月就叫佳恩堂.
和雄恩堂都一起.
在福音上面大家一起.
來去共同的努力.
分享資源.
讓神的工作可以做得好.
我們下個星期就回去看馬拉基書.
最後的一卷歸回的先子書.
我們一起禱告.
多謝主給我們今天晚上.
可以將第十四章.
都算是完整的.
亦都是有一些的.
教導性的來彼此分享.
願你保守我們對聖經的明白.
更加讓我們真的在.
我們的頸項.
代表著我們的信服.
能夠有一個金鈴當.
我們能夠真的提醒周圍的人.
包括我們自己.
我們是歸神為聖的.
盼望我們能夠.
真的整個人生.
有一個對神的榮耀的追求.
以致我們將來其實.
我們就是得到神給我們的榮耀.
這樣感恩禱告.

$^{2241}$奉耶穌基督的名.
阿們.
\newpage



\section{哥林多前書 9:24-27}
\label{sec:lIj_mrEVPGg}
\textbf{2023-12-30 晚堂《為主奔跑》哥林多前書 9\_24-27 講員:陳一華牧師}
\newline
\newline
連結: \href{https://youtube.com/watch?v=lIj-mrEVPGg}{\texttt{https://youtube.com/watch?v=lIj-mrEVPGg}} ~~~~ 語音日期: 2024-01-03
\newline
\newline
\hyperref[sec:MkLoNpPCJk4]{\small{< < < PREV SERMON < < <}}
~
\hyperref[sec:index_chronic]{\small{[返順時目]}}
~
\hyperref[sec:index_scriptual]{\small{[返順卷目]}}
~
\hyperref[sec:c3xZAeZuwD0]{\small{> > > NEXT SERMON > > >}}
\newline
\newline
哥林多前書 9:24-27
\newline
\begin{longtable}{cl}
\hline
\hline
章節 & 經文 (和合本修訂版)\\
\hline
9:24 & \begin{tabularx}{0.7\textwidth}{X} 你們不知道在運動場上賽跑的，大家都跑，但得獎賞的只有一人？你們也要這樣跑，好使你們得著獎賞。 \end{tabularx} \\ \\ \relax
9:25 & \begin{tabularx}{0.7\textwidth}{X} 凡參加競賽的，在各方面都要有節制，他們不過是要得會朽壞的冠冕；我們卻是要得不會朽壞的冠冕。 \end{tabularx} \\ \\ \relax
9:26 & \begin{tabularx}{0.7\textwidth}{X} 所以，我奔跑，不像無目標的；我鬥拳，不像打空氣的。 \end{tabularx} \\ \\ \relax
9:27 & \begin{tabularx}{0.7\textwidth}{X} 我克制己身，使它完全順服，免得我傳福音給別人，自己反而被淘汰了。 \end{tabularx} \\ \\
[1ex]
\hline
\hline
\end{longtable}
$^{1}$弟兄姊妹晚安.
時間真的過得很快.
眨眼間我們今晚已經是.
2023年的最後一個晚堂的崇拜了.
所以我相信大家來之前.
都很珍惜今晚的晚堂的崇拜.
正如剛才主席祈禱的時候說到.
上帝從歲守到年終.
祂都是看顧我們.
所以正如領唱的姐妹.
一唱那首詩歌.
應該是第三首.
叫做《數算主因》.
我就立刻擔高頭.
看一看大家.
看看有多少人懂得唱這首詩歌.
因為這首歌是古老的歌.
是很久很久之前.
特別那時候宣導出版社出的《青年聖歌》.
裡面的其中一首詩歌.
很古老的.
近年都多了唱一些其他的短的詩歌.
我想有沒有人懂得唱這首《數算主因》呢?.
我看一看大家的時候.
發覺形勢非常好.
不是不妙.
原來很多人都懂得唱這首詩歌.
很投入.
朗朗上口.
換句話來說.
這首詩歌提醒我們.
《數算主因》是非常重要的.
特別我們來到年尾的時候.
我們更加回頭看一看.
上帝給我們的恩典.
給我們的保守祝福.
真的有很多時候超出我們所求所想的.
特別過去一年.
你有很多的活動.
不論在香港.

$^{41}$不論在外地.
不論在任何的地方.
你發覺來到年尾的時候.
人都很健在.
人都夢神的保守.
一切其實都是上帝的恩典.
上帝的祝福.
因為來到年尾.
我想一想.
很快就是年初了.
過兩天就是新一年開始了.
我就想一想.
哪個題目會好一點呢?.
我就想起今晚這個題目.
就是「為主奔跑」.
剛才讀的那段經文.
就比較冷門一點.
不是我們那麼熟悉.
不過經卷大家都很熟悉.
林多前書.
但是那幾節聖經就比較陌生一點.
保羅其實在這一章第九章.
一開始.
由第一節到二十三節.
其實它的主題只有一個.
就是講到保羅的全福音的心志.
和他的行動是非常清晰的.
他在這二十幾節的經文當中.
其實弟兄姊妹你都很熟悉.
因為有些經文你很有印象.
你可能靈修,查經或者聽道都聽過.
保羅怎樣講他自己全福音的心志呢?.
他說我全福音是怎樣的?.
保羅說是不得已的.
記不記得?.
若不全福音我就有禍了.
保羅亦都告訴我們.
如果我甘心做這件事情.
就得賞賜.
如果我不甘心去做.

$^{81}$責任就託付在我身上了.
所以保羅亦都在這一章裡面講到.
他說我向著猶太人.
我就做猶太人.
說要得猶太人.
大家有沒有印象?.
好像沒有什麼反應.
他說我向著律法以下的人.
我就做律法以下的人.
向著沒有律法的人.
我就做沒有律法的人.
說要得著什麼?.
沒有律法的人.
保羅說我向著軟弱的人.
我就做什麼?.
軟弱的人.
無論如何總要救赦人.
大家聽我這樣讀這幾經文.
這幾經是不是都很熟悉?.
保羅沒有停在這裡.
保羅繼續說.
他怎樣說?.
他說凡我所作的.
都是為福音的緣故.
為要叫人與我同得什麼?.
這福音的好處.
所以大家就知道.
保羅由他在大馬色路上遇見主.
被主光照之後.
他回轉悔改.
接著主就將任務給他.
主怎樣說呢?.
他說從現在開始.
將來你要在同胞.
即是猶太人.
外邦人.
和君王的面前.
做我的見證.
所以保羅信主不久.
領受的使命任務.

$^{121}$亦都非常清晰.
所以我們讀到這章經文.
一到廿三節的時候.
就不足為奇.
因為我們都發覺.
真是和他的心智夢照是很相稱.
不過去到廿四節的時候.
忽然間保羅的筆鋒一轉.
就轉到一個冷門一點的課題.
是說什麼呢?.
原來他告訴我們.
我們每一個傳福音的人.
每一個事奉上帝的人.
其實我們都是一個運動員.
都是一個運動員.
我都是經常顧著說.
忘記了按給大家看.
好了.
為什麼我們都是運動員呢?.
有三樣很相似的地方.
第一樣.
原來運動員.
保羅怎樣說呢?.
他說運動員在場上賽跑.
個個都拼了老命.
用盡他們的方法.
在比賽場上來賽跑,比賽.
希望的無非是拿到一個獎牌.
你都知道.
任何的比賽.
頂多都是三個獎牌.
金銀銅三個獎牌.
但是運動員就因著這些獎牌.
他們就發奮拼勁拼搏.
都要來到什麼呢?.
希望拿到獎牌.
這個就是運動員他們最終極的目標.
希望能夠有一個獎賞在當中.
保羅就說.
原來我們每一個信了耶穌的人.

$^{161}$我們每一個傳福音的人.
原來我們都有獎賞的.
而我們的獎賞.
更加特別是不止一個人.
不止三塊獎牌.
原來每一個基督徒.
只要他是忠心的事奉上帝.
到將來有人.
他都會得到獎賞.
得到賞賜.
大家都知道.
聖經記載基督徒將來見主面的時候.
得到的賞賜.
起碼有三頂的冠冕.
我不知道你有沒有印象.
一頂就是生命的冠冕.
另一頂榮耀的冠冕.
第三頂就是公義的冠冕.
大家有印象嗎?.
原來每一個屬主的人.
每一個忠心事奉上帝的人.
我們都能夠得到賞賜.
所以我現在就開始明白.
為什麼有些人.
他們會在百多年之前.
或者二百年之前.
他們遠道重陽.
千里迢迢.
來到中國大陸的地方.
來去傳福音.
甚至將他們的性命都擺上.
不顧一切.
來到中國大陸一個荒蕪的地方.
來去傳福音.
我們明白.
因為這班宣教士知道.
將來在天上.
上帝會為他們預備什麼?.
榮耀的冠冕.
我也認識一位弟兄.

$^{201}$他的工作實際上是相當不錯的.
是公務員的工作.
鐵飯碗大家都知道.
他的工作前途是非常之好的.
也可以說是高薪厚職.
但我們不明白.
為什麼這個弟兄做得沒多久.
接著就打信辭職.
為什麼呢?.
原來他有一個感召.
要進神學院接受訓練.
將來他知道.
他要全時間去事奉上帝.
做一個傳道人.
弟兄姊妹們.
我們有時候都想.
這個弟兄是不是很傻呢?.
這麼好的工作.
這麼好的幽差.
那些人真的流口水.
你竟然自動辭職.
走去讀神學.
有時候我們都會有一個問號.
為什麼他這樣做呢?.
可能真的很多人都替他很可惜.
覺得他很傻.
但我們讀到這段經文.
知道答案了.
為什麼呢?.
因為他知道他這樣去信奉上帝的呼召.
將來必定在上帝的面前領受什麼呢?.
領受寶貴的賞賜.
我在警察團契都認識一位弟兄.
他給我的印象都覺得很特別.
實際上大家都知道.
如果他一路上班工作.
做到退休.
是非常之理想.
他的祥奉.
他的心是很可觀的.

$^{241}$但這位弟兄不明白.
為什麼未到五十歲.
他就打算辭職.
早點退休.
為什麼這樣做呢?.
原來他覺得自己上半生.
所過的日子.
很少時間能夠專注地去事奉上帝.
所以他想到人生都快過半百的時候.
不如我就想一想.
我應該下半場怎樣走呢?.
這位弟兄清楚上帝的心意.
毅然決然就辭工.
當時他身邊的同事都覺得他很傻.
為什麼這麼快就辭職呢?.
做多幾年 做到退休.
大家都知道他的祥奉是非常之何觀的.
為什麼會這樣做呢?.
今晚我們都有答案.
聖經說什麼呢?.
原來當一個基督徒信服上帝的心意.
走在上帝的旨意當中.
忠心地去事奉上帝的時候.
將來他必定得到賞賜.
今天我們明白.
保羅為什麼用運動員比方做基督徒.
第一個類同是因為什麼呢?.
是有獎賞.
我們看第二個原因.
原來基督徒跟運動員有相同的地方是什麼呢?.
就是有操練.
保羅怎麼說呢?.
既然這班人在場上賽跑的時候.
跟著保羅.
所以你們也當這樣跑.
有沒有印象?.
就是說原來每一個基督徒.
我們都要什麼呢?.
都要繼續不斷地向前邁進.
要跑.

$^{281}$就是說要操練的意思.
所以我們明白.
運動員最重要是什麼呢?.
就是要操練得很健美才行.
尤其是要比賽的時候.
開講也有說.
台上十分鐘.
台下就多少呢?.
十年功.
各位.
原來我們不能走捷徑的.
如果你說我事奉上帝.
我服侍上帝.
可以了.
我不用什麼操練.
不用什麼訓練.
我就這樣去的時候.
等一下 姐妹們.
如果我們沒有操練.
沒有訓練的時候.
我們真的會很危險的.
有一年.
當我在渣馬.
跑過青魚隧道的時候.
一天隧道沒多久.
就看到前面有個跑手.
就跌在地上.
然後大字型就躺在地上.
我跑過的時候.
我的心都很擔心.
希望這個年輕人沒事.
希望他沒什麼事.
但我看到他已經跌在地上.
又吐白沫.
我第二天早上看新聞.
就知道原來這個年輕人.
他拿了別人的號碼去跑.
就是說什麼呢?.
就是說應該跑的那個.
沒有跑到.

$^{321}$不知道是不是不舒服或者有事.
不應該跑的那個.
就跑去跑.
結果有沒有操練到?.
沒有操練到.
就在青魚隧道躺下了.
後來這個年輕人是失救.
救不回來.
各位弟兄姊妹們.
原來我們每一個教徒.
我們事奉上帝.
服侍上帝.
我們需要把握不同的機會.
來接受操練.
接受訓練.
癌症關懷事工聯會.
簡稱癌聯.
這個機構這個組織.
他們有一個特別的地方.
就是很強調提供訓練.
所以一年裡面有三次到四次.
有一個叫「顏關大使」的訓練.
我們當中和同工一齊來到.
大家事奉的時候.
每一次顏關大使的訓練的時候.
那個反應都相當好的.
不少人讀這個課程.
如果一讀就要讀八堂.
早期還十堂.
後來八堂.
各位弟兄姊妹們.
有時候我都會想.
為什麼弟兄姊妹願意抽八個星期出來.
去上堂去裝備自己呢?.
後來在我們課堂之餘.
我們有時候都會小組來分享.
聽到一些組員的心聲.
大部分的人都覺得.
他們很緊張一些患癌的親友同事.
但是他們不知道怎樣去幫他們.

$^{361}$不知道應該怎樣去做.
不知道怎樣去說應該說的話.
所以他們覺得很無助.
但是一聽到有癌聯這個課程.
他們就馬上去報名.
弟兄姊妹們.
我很欣賞這一批弟兄姊妹.
因為他們知道.
原來事奉上帝是必須有裝備.
有訓練.
就正如運動員是要有操練一樣.
我們再看第三.
為什麼運動員和基督徒是這麼雷同的呢?.
就是因為大家都有進步.
大家喜不喜歡這個運動員呢?.
很喜歡他吧?.
我想很多香港人都很喜歡他吧?.
他的名字不用說你們都知道是什麼名字.
這個運動員真的可以說是什麼呢?.
人見人愛.
下面那句我不用說了.
為什麼呢?.
每一次他贏了比賽之後.
你留意一下記者訪問他的時候.
你知不知道他講了什麼話?.
這句話重複又重複.
我聽了之後.
我得到很大的一個啟發.
他說什麼呢?.
他說如果下一次再有得比賽的時候.
我相信我會有得更加好一點.
大家明白這個意思嗎?.
就是說雖然他已經拿了冠軍了.
但是他覺得自己還有什麼呢?.
進步的空間.
哇!.
這些運動員.
正不正啊?.
厲害嗎?.
很厲害.

$^{401}$頂尖舞門.
有一個這樣的心態的運動員.
難怪他每次游泳都能夠拿到最佳的成績.
頂尖舞門.
你知道嗎?.
這個世界上任何最頂尖的運動員.
他怎樣打破紀錄.
怎樣能夠贏得比賽.
我想告訴大家.
他所得到的這個擦身的紀錄.
其實可能風塵三年,五年,十年.
始終有一天.
他的紀錄都會被人打破.
大家相不相信?.
相不相信?.
因為什麼呢?.
每一項的運動都有進步的空間.
所以我們不要想著.
這個紀錄永遠可以保持下去.
這樣不是太好.
為什麼?.
這樣就是沒有進步.
如果有進步的空間的時候.
你一定會繼續努力.
各位.
我們今天事奉著傳福音.
就需要有這種心態.
就是有進步.
如果你說我跟別人傳福音.
我也很想我的家人信耶穌.
但是你發覺自己傳了一次,兩次都不太行.
都不太有成績.
都被人拒絕.
可能你再想想.
或者我檢討之後.
發覺自己傳福音的訓練比較少.
如果是教會有一些三福報道的訓練.
我下一次都要去參加.
為什麼?.
因為他發覺原來他跟別人傳福音的時候.

$^{441}$一味都是自己說了.
別人沒得說.
那怎麼去跟別人傳福音呢?.
大家明白這個意思吧?.
所以各位弟兄姊妹妹.
原來我們傳福音的時候.
是需要什麼?.
繼續不斷地進步.
我們事奉上帝也是.
有很多人以為事奉很簡單.
不是的.
當你在教會事奉.
或者在職場裡面事奉.
你要記得.
當你參與事奉上帝的時候.
永遠都不會是一個one man band.
一個人.
一人隊長.
全部做完.
不是的.
當你去事奉的時候.
你一定會跟不同的人去磨合.
去配搭.
大家認不認同?.
當你跟其他人去磨合配搭的時候.
各位.
你知道每一個人都有性格的.
每一個人都有不同的作風.
每個人的觀點角度都不同.
這一刻我們怎樣跟弟兄姊妹.
有一個美好的配搭.
一起去事奉呢?.
需不需要裝備?.
需不需要在失敗當中.
來檢討去進步?.
弟兄姊妹們.
我們需要的.
所以當你記得.
我們事奉上帝的時候.
為什麼跟運動員那麼相似呢?.

$^{481}$因為我們永遠都有進步的空間.
如果有一天你發覺自己什麼都行.
什麼都懂.
事奉什麼崗位.
一個人殺光.
前無古人.
後無來者的話.
那就淪盡.
弟兄姊妹們.
我們就沒有進步的空間.
所以.
這個就是保羅.
以運動員比方基督徒的原因.
接著第二個部分.
跟大家看看怎樣去達標.
怎樣成為一個達標的運動員呢?.
如果我們要成功的話.
保羅都提醒我們.
有三個條件很重要的.
第一個條件.
就是我們需要方向.
而這個方向.
是一個很正確的方向.
千萬不要混亂.
千萬不要事奉.
把傳福音當成等閒事.
隨便隨便去做.
沒有方向.
沒有目標.
沒有計劃.
不行的.
弟兄姊妹們.
原來我們見到運動員.
他們很有清晰的.
目標方向.
他玩哪一種運動.
他一定是用那種運動的特質.
去找一些方法去裝備自己.
所以運動員是很需要方向的.
我亦都記得有一年.

$^{521}$當我去跑步的時候.
就發覺有一個跑手.
跑步的時候.
發覺不對勁.
為什麼呢?.
原來發覺他跑到差不多.
氣喘如牛.
接著.
多看他兩眼.
臉青口唇白.
什麼原因呢?.
原來這個跑手.
他應該跑的方向.
跑到某一個轉折點.
他就要轉彎回到終點.
因為他步的時候.
是步半馬.
但不知為什麼.
他跑到一個轉折點.
沒有轉彎.
繼續向前跑.
大家都知道繼續向前.
是什麼呢?.
是全馬.
所以我看到這個跑手的時候.
替他很擔心.
因為他跑錯了方向.
各位頂尖的朋友們.
跑錯方向.
是一個很大的.
一個挫敗.
浪費了時間.
浪費了自己的操練.
正所謂一子錯.
滿盤皆落索.
所以我們要達標.
要為主奔跑的時候.
我們要有一個清晰的方向.
頂尖的朋友們.
今天我們已經來到.

$^{561}$2023年最後的一兩天.
實在是有點累了.
如果能夠早一兩個星期.
或者早幾個星期.
你可能聽到牧師講這些堂道.
可能更加受惠也未定.
我會鼓勵大家.
真的拿些時間安靜.
給自己休息一下.
靜一靜.
想一想2024新一年.
到底上帝要我去事奉的方向.
是甚麼呢?.
2024年有甚麼使命工作.
上帝是要我去做的呢?.
各位頂尖的朋友們.
有沒有想過.
不是說所有眼前的工作.
我都全部納下去做.
那只會是甚麼呢?.
越做就越耗盡.
很快就很不濟了.
所以我們需要有方向.
可能2024年你發覺.
上帝感動你要在教會裡面.
參與多些的事奉.
又或者可能你發覺2024年.
上帝要你在職場裡面.
做多些的事奉.
又或者可能在社區裡面.
上帝讓你看到這個社區的需要.
所以我要用多些時間.
在這個社區去事奉上帝.
又或者可能你發覺.
這個家庭的崗位.
是上帝給你的一個很清晰的方向.
可能你要用多些時間.
在這個家庭崗位上.
來事奉神.
當然有些頂尖的姐妹.

$^{601}$她發覺2024年.
上帝感動她.
去醫院做羽毛義工.
這樣更加清晰.
怎樣來好好善用她的時間.
因此才幹智慧.
來去發揮得淋漓盡致.
第二個條件是甚麼呢?.
第二個條件就是需要節制.
需要節制.
所以剛才讀的聖經.
保羅說.
運動員凡事諸事都要有甚麼?.
有節制.
大家有沒有印象?.
很重要的.
這位運動員.
幾十年前他很出名的.
大家都知道.
幾十年前.
他可以說是一個很標稱的運動員.
無論是單槓.
無論是吊環.
或者是木馬.
沒有一項難得到體操王子.
大家對他都有印象.
當然他現在不是這樣的身形.
現在又肥又年紀大了.
不過他做生意就發達了.
但你知不知道當年他代表國家出賽的時候.
為了要在比賽場上保持他的身形體重.
他有30天時間.
一粒飯都不吃.
他只吃一些粥水.
為甚麼呢?.
因為他要維持自己的體形.
狀態.
所以運動員.
他要有節制.
他要犧牲.

$^{641}$不可以吃那麼多甜品.
不可以吃那麼多朱古力.
不可以吃那麼多你喜歡吃的東西.
你要甚麼?.
有節制.
然後你才能夠發揮到你最佳的狀態去比賽.
各位聽眾.
這個節制對我來說.
今天有另一個意思是甚麼呢?.
就是如果我們今天事奉上帝的時候.
這個節制.
另一個意思就是.
我們懂不懂得取捨.
取捨.
你有沒有將你每天做的工作.
我想你都會有做的.
十幾二十項工作.
你一定有優先次序.
有時有些事你要覺得你很喜歡的.
很有興趣的.
不過因為事奉的緣故.
你就很樂意.
甘心的擺下.
我認識一位醫生.
他跟我分享.
他說牧師我雖然很喜歡打高球.
這裡有很多高球的.
他說如果我某一個週末.
有我的同行的人.
同行的人.
未信耶穌.
有約到他們喝茶吃飯的話.
可以傳到福音.
他說那天那個週末.
他寧願不去打高球了.
各位頂次妹妹.
不是說打高球不好.
不要誤會這一點.
而是他懂得取捨.
他知道什麼是重要.

$^{681}$如果他要去人傳福音的時候.
這是難得的機會.
雖然他有很大的興趣.
但他都願意擺下.
今天你對我來說.
成為一個達標的基督徒.
我們懂得要去取捨.
操練我們的取捨.
最後.
需要的是什麼呢?.
恭克.
保羅說.
我是恭克己身.
叫身服我.
有沒有印象?.
免得我傳福音給別人.
自己反倒而已.
被棄絕.
恭克原文的意思就是.
向敵人要害的地方.
全力出擊.
這叫恭克.
換句話來說.
保羅的恭克就是針對他自己.
信了耶穌之後.
原來我們這個人有肉體.
有血肉 有血氣.
我們有我們的情慾.
原來這些情慾令到我們的人.
很多時候表現得很自我.
很舊我又出來了.
我們這個人是老我又出來了.
各位聽眾.
因為這些都是什麼呢?.
我們不好的行為表現.
有時候不知不覺間就流露出來.
所以保羅說我要恭克己身.
然後運動員就需要.
知道自己的弱點在哪裡.
他要特別針對弱點.

$^{721}$致以一個致命的攻擊.
什麼是最好恭克呢?.
我相信莫過於二弗所書提醒我們.
二弗所書說什麼呢?.
原來一個基督徒能夠戰勝敵人.
戰勝魔鬼撒旦的隱有試探.
就是拿著聖靈的寶劍.
就是神的道.
各位.
如果我們能夠多讀聖經.
多留意神的話語.
上帝的應許的時候.
藏在我們的心中.
我們就明白.
詩人說我將你的話藏在心中.
免得我得罪你.
弟兄姊妹們.
原來我們能夠拿到神的話.
聖經就是聖靈的寶劍.
我們就能夠戰勝老我.
救我和治我.
若果不是的話.
原來保羅說.
免得我傳福音給別人.
自己反倒就是被棄權.
這是和合本的翻譯.
現在很好.
新譯本和修本就翻譯成.
免得我被淘汰.
這個棄絕原文的意思就是.
不合格.
肥了老.
所以就被人篩了出來.
肥老了.
被人淘汰了出來.
我很熱心傳福音.
我很熱心事奉上.
但是想不到我最後.
我竟然得不到上帝的掌上.
而反而被淘汰.

$^{761}$為什麼呢?.
因為我們不懂得.
功克己身.
要身服我.
所以為什麼我們很多時候.
信了耶穌.
無論你多少年也好.
幾十年也好.
都不能夠放鬆自己.
屬靈的操練和進步.
若果不是的話.
我們很容易.
慢慢地.
跟主的關係疏離了.
跟主的關係疏離疏遠的時候.
我們其實很容易跌倒.
大家不要誤會.
難怪有時候你看到.
有些消息.
這個人信耶穌這麼久.
為什麼也會跌倒.
無他.
因為他跟神的關係疏遠了.
所以現在我們明白.
保羅說要功克己身.
然後我們才能夠達標.
最後各位感孫的弟兄姊妹.
我們來到新一年就法開始的時候.
求主幫助我們.
能夠在2024年當中.
能夠全力以赴.
全神貫注地.
為主奔跑.
所以我希望我們每一位.
都將2024年的時間.
用得最美麗.
用得最美好.
用得最有價值.
以致你接下來這一年.
你就好像一個很fit的運動員.

$^{801}$發揮得淋漓盡致.
我們低頭祈禱.
親愛的主多謝你.
因為保羅給我們的教訓.
是提醒我們.
真的好像一個運動員.
我們要全力以赴.
我們為主奔跑.
為主作功的時候.
我們真的懂得.
來到是要有正確的方向.
要有很好的節制.
也都要功克己身.
所以求主加難賜禮給我們眾弟兄姊妹.
把握新一年開始的時候.
我們可以更加成為主你手中.
合用的器皿.
能夠祝福多人.
榮神益人.
奉主名求.
阿們.
\newpage



\section{路加福音 2:1-7}
\label{sec:c3xZAeZuwD0}
\textbf{2023-12-31 加恩樓《在你心裡有空處給主嗎?》路加福音 2\_1-7 講員:謝靄薇牧師}
\newline
\newline
連結: \href{https://youtube.com/watch?v=c3xZAeZuwD0}{\texttt{https://youtube.com/watch?v=c3xZAeZuwD0}} ~~~~ 語音日期: 2024-01-03
\newline
\newline
\hyperref[sec:lIj_mrEVPGg]{\small{< < < PREV SERMON < < <}}
~
\hyperref[sec:index_chronic]{\small{[返順時目]}}
~
\hyperref[sec:index_scriptual]{\small{[返順卷目]}}
~
\hyperref[sec:p_xp8W8_F_Q]{\small{> > > NEXT SERMON > > >}}
\newline
\newline
路加福音 2:1-7
\newline
\begin{longtable}{cl}
\hline
\hline
章節 & 經文 (和合本修訂版)\\
\hline
2:1 & \begin{tabularx}{0.7\textwidth}{X} 在那些日子，凱撒奧古斯都降旨，叫全國人民都登記戶籍。 \end{tabularx} \\ \\ \relax
2:2 & \begin{tabularx}{0.7\textwidth}{X} 這第一次登記戶籍是在居里扭作敘利亞總督的時候行的。 \end{tabularx} \\ \\ \relax
2:3 & \begin{tabularx}{0.7\textwidth}{X} 眾人各歸各城，辦理登記。 \end{tabularx} \\ \\ \relax
2:4 & \begin{tabularx}{0.7\textwidth}{X} 約瑟也從加利利的拿撒勒城上猶太去，到了大衛的城名叫伯利恆，因為他是大衛家族的人， \end{tabularx} \\ \\ \relax
2:5 & \begin{tabularx}{0.7\textwidth}{X} 要和他所聘之妻馬利亞一同登記戶籍。那時馬利亞已經懷孕。 \end{tabularx} \\ \\ \relax
2:6 & \begin{tabularx}{0.7\textwidth}{X} 他們在那裡的時候，馬利亞的產期到了， \end{tabularx} \\ \\ \relax
2:7 & \begin{tabularx}{0.7\textwidth}{X} 就生了頭胎的兒子，用布包起來，放在馬槽裡，因為客店裡沒有地方。 \end{tabularx} \\ \\
[1ex]
\hline
\hline
\end{longtable}
$^{1}$各位弟兄姊妹平安.
讓我們看神的話.
支持我們一起有望主.
每有恩典 每有慈愛的主.
從歲守到年終.
你保守我們.
你帶領我們 你焦聚我們.
可以再次的來到你面前.
因為這個是耶和華的日子.
你所定的聖日.
你讓我們能夠歡呼喜樂的.
來到你面前.
與你的靈與我們同在.
由於我們心裡有什麼.
得罪你的.
與主你羅開.
以致我們能夠.
高舉你.
能夠.
不見一人 只見耶穌基督.
與你對我們說話.
當我們要去承擔.
我們將要成為.
這個.
佳人堂的時候.
有很重要的使命.
你要給我們去做的時候.
求你幫助我們.
求你讓我們有使命的時候.
讓我們有信心.
去依靠你.
能夠跨步向前.
你祝福我們每一個.
讓孩子心裡的意念.
口中的言語.
是無理如納.
奉主耶穌基督的名求.
當我們去.
看今天這個經文的時候.
我們知道.

$^{41}$就是.
當時.
當政的是奧古士督.
這個王.
他在當政的時候.
他剿平一切的勢力之後.
當時.
被羅馬帝國帶來一個.
昇平的景象.
你會看到.
當時被譽為.
羅馬和平的時期.
在公元前一世紀.
在伊弗所一塊.
碑文裡面.
奧古士督.
是世界的救主.
在小亞細亞.
連他的總督.
都寫在公銜上.
指奧古士督.
海撒奧古士督.
他誕生的這一天.
被世界帶來福音.
所以你可以想像.
路加.
他寫路加福音的時候.
他特別說到.
天使去宣告.
很重要.
宣告什麼呢?.
真正的拯救者.
就是耶穌基督.
他要拯救萬民.
全世界的人.
所以.
救主耶穌基督的降生.
是一個好的消息.
所以我們看到.
剛才的經文裡面.

$^{81}$我們看到為什麼.
在耶穌基督的時代.
那些客店.
會住滿呢?.
難道像我們今天去旅遊.
假期.
沒有地方住?不是的.
當你們剛才聽經文.
一開始就知道.
原來是奧古士督.
他要在當中.
叫人去報戶口.
難怪沒有地方.
你會看到.
就是.
那些人去登記.
沒有戶口之外.
到底.
這個約翰福音.
一章十二節.
跟你說什麼呢?.
這個他指的是耶穌基督.
他在這個世界.
這個世界是他做的.
世界卻不認識他.
他到自己的地方來.
自己的人.
並不接納他.
所以你在路加福音開始的時候.
你會看到這個嬰兒出生的時候.
已經沒有人.
去接待他.
沒有地方住.
而我們看客店的時候.
這個客店.
這個字有兩個.
一個字.
你會看到.
頭一個字.
就是指.

$^{121}$有老闆去主持.
有食物供應.
有水給你.
可以住.
第二個字.
原來.
好像有些.
給畜生.
隔開的圍欄.
或者.
在當中.
旅客也可以住.
不過就沒有老闆.
去負責.
去管理你.
沒有食物供應.
只是給你水.
我們看到.
原來那時候.
沒有那麼多人去.
每個人都要有戶口.
所以連這些地方都沒有.
你可以想像.
原來瑪麗亞和約瑟.
他們都覺得嬰兒快出生.
很緊張的情況下.
可能他們只是去一些民居.
一些很簡單的.
有些圍欄.
隔著的.
就這樣.
或者每個人都去.
很簡單的居民.
一個很簡陋的地方.
每個人都要有戶口.
每個人都去用.
難怪連嬰兒出生的空間.
都沒有.
所以你會看到.
嬰兒只能放在一個.

$^{161}$馬槽的地方.
在那麼緊張的.
關頭之下.
你會看到.
耶穌基督到自己的地方.
但是.
並沒有接納他的地方.
因為那麼榮耀的君王.
他來到地上.
神用一個很奇妙的方式.
他來到地上.
雖然沒有人接待他.
但是他就誕生了一個.
那麼卑微的地方.
雖然不是皇宮.
大家.
你會想像你的出生.
都比他乾淨很多.
起碼家人都給你一個乾淨的床.
有衣服準備給你.
但是他沒有.
他只是用布包著.
第二個景象.
你會看到.
希律.
心裡不安.
當你看回.
《馬太福音》的時候.
為什麼不安呢?.
原來有幾位博士.
他們.
看到東方的星.
他們知道.
他問希律.
他問什麼呢?.
他生下來要做猶太人的王.
在哪裡呢?.
當他問之下.
這個王就休壽了.
心裡不安.

$^{201}$於是他問誰呢?.
他就召祭司和問師.
這個王會在哪裡誕生呢?.
他就告訴他.
原來會在伯尼咸.
於是他暗中.
就去.
召這些博士.
問他.
這個嬰兒.
你大約什麼時候.
見到這顆星呢?.
他大約摸索一下.
到底什麼時候出現呢?.
然後他就像好人.
他說.
你去查探之後.
告訴我.
我會去敬拜他.
表面上.
他做得很好.
但是他的心是否這樣呢?.
他心裡不安.
但是回頭看.
你見到這個博士.
他們跟著星望.
就找到耶穌基督了.
他們去敬拜.
耶穌基督之後.
本來就應該要去.
報告給王.
但是有主的使者去指示他.
原來你不要去.
不要去了.
於是他們走另一條路.
然後有主的使者.
跟約瑟去顯現.
告訴他們.
你們快點離開.
於是他們要逃亡了.

$^{241}$離開這個地方.
去埃及.
你見到希律很生氣.
我叫你告訴我.
你不告訴我.
所以他摸索一下.
這個嬰兒有多大呢?.
於是他就將兩歲以下的嬰兒.
全部殺了.
除了耶路撒冷那裡.
百里行那裡.
四景的全部殺了.
你見到就是.
原來希律.
是這麼生氣的.
因為他很害怕.
他的王位不保.
然後你看到.
耶穌的人生裡面.
又再見到他去到.
格拉森的地方.
大家都很清楚這個格拉森人是怎樣的.
他很大力的.
又不穿衣服的.
在那裡出沒的時候.
那些人用鎖鏈.
鎖住他都不行的.
但是.
當耶穌基督拯救他.
他變成一個正常人的時候.
聖經就看到.
就是.
別人看到他改變.
只是害怕.
但是沒有人.
為耶穌基督去改變他.
行他的歧視.
去感恩.
去讚美神.
完全沒有這個舉動.

$^{281}$你看這個人.
穿得多好.
他正常了.
穿衣服.
耶穌接納他.
但是另外一幫人.
因為這些鬼.
被這個人.
扶在他身上之後.
耶穌基督將這個鬼趕出來之後.
這個鬼就要求.
去豬那裡.
這個鬼的名字叫群.
很多不是只有一隻.
那些鬼就去了豬那裡.
結果豬就掉到山崖裡.
全部死了.
這些人.
做生意的.
利益受損.
當然利當頭.
利當頭的時候.
人是給錢.
給勢力.
去充分頭的時候.
根本就看不到.
神的榮耀.
他們怎麼做呢.
你看到.
他們要耶穌離開.
但是被鬼附於這個人.
卻是怎樣呢.
耶穌請求他之後.
他穿回衣服.
他想和耶穌在一起.
但是耶穌叫他回家.
你要傳揚神在你身上的作為.
你要告訴別人.
但這個人真的做到了.
在經文的39節.

$^{321}$你會看到這個人走遍全城.
去傳揚耶穌基督.
為他做了何等大的事.
他真的照著耶穌的吩咐去做.
但其他的人.
他們好像世俗的人.
連門神都趕走了.
財神.
過年的時候.
人們會唱財神.
其實耶穌基督.
有耶穌就有一切.
但他們的耶穌都趕走了.
因為你在這裡.
阻礙我做生意.
把我的豬.
趕到山崖.
死了我都沒有生意做了.
所以他們不歡迎耶穌.
反而要趕走耶穌.
今天.
主為你寫命.
你又怎樣對待主呢.
在你心裡.
你有多少空間給主呢.
沒錯.
我們可能每個星期都會勉勵其難.
我們都會聚集.
因為我們要拿耶穌給我們去祝福.
但原來.
我們要變成.
是不是嘉恩堂.
我們有使命.
你的責任.
不是在我們身上.
而是在我們身上.
我們有使命.
你的責任.
不是聚集那麼簡單.
你只改了一個字.

$^{361}$你的肩膀就重一點.
記住.
耶穌說要付代價.
原來我們這班.
我們承擔.
耶穌給我們使命的人.
我們要跟從他.
原來要付代價.
要寫記.
要天天背在十字架.
來跟從耶穌.
不容易的.
當我們看到.
要付代價的時候.
看看瑪利亞.
神已經給了她一個使命.
要做耶穌基督的媽媽.
當她承擔這個使命的時候.
你看到她怎樣.
還要輕輕的.
從那裡.
來到哪裡.
百里行.
她不是普通的人.
她是一個.
她不是普通的人.
那時已經有身孕.
她肚子很大.
還要這樣.
這麼遠的地方.
來到這個地方.
是不是很辛苦.
她已經知道.
她要成為耶穌基督的媽媽.
但她還要這麼辛苦.
來到這個地方.
沒有地方招待她.
嬰兒還要.
躺在馬槽裡.
最後還要怎樣.

$^{401}$被人追殺.
還要逃去埃及.
大家有沒有看到.
原來我們做一個基督徒.
我們的人生不是很浮躁.
不是我們今天信了主.
洗了禮.
等著花轎.
接你去天堂.
不是這麼簡單的.
原來我們有使命的.
我們要背在十字架跟從耶穌.
原來這條路是不容易的.
你會看到耶穌.
來到的時候.
好像沒有空間.
給主.
好像兒童.
幼稚園.
他們做聖景的時候.
有一幕.
就是瑪利亞.
和約瑟要找住的地方.
每次都說沒有位置.
但在表演當中.
有個人說.
我有.
這個小朋友一說.
每個人都笑.
好像台詞改了.
怎樣接下去.
他說他有.
其實今天我們有沒有.
有沒有空間給主.
瑪利亞雖然.
已經知道.
自己是成為.
耶穌基督的媽媽.
我們要遵從.
聽耶穌的話.

$^{441}$遵從耶穌的話.
遵從天父的旨意.
但沒有保證.
你的人生是很舒服的.
不簡單.
在這條路上走路不簡單.
但神卻.
影響我們.
雖然在不簡單的路程中.
你未必是很舒服的生活中.
神有祂的美意.
我們儘管聽祂的話.
去做,雖然是不舒服的.
但在耶穌基督.
在神的計劃中.
你的人生是有意義的.
儘管是.
不容易走的.
很痛苦的.
但為何.
明明已經是真命天子.
我已經認證了.
選擇了我.
但原來這條路.
不是一帆風順的.
大家知道成為基督徒這條路.
不是一帆風順的.
你不要以為今天我信了主.
我就會發達.
我們要搞事一定是最好.
但未必的.
我們要忠心做好自己應該做的本分.
然後等著.
主的安排.
你的人生.
交給神.
但我們要做好.
我們要好好裝備自己.
所以你會看到.
原來.

$^{481}$走這條路.
要預備一個受苦的心智.
這條路不簡單.
要預備一個受苦的心智.
當神呼召你的時候.
當時神都在呼召人.
有人很大聲地說.
無論你拍我去哪裡.
我都會跟隨你的.
但耶穌隨即說.
狐狸有洞.
天空的飛鳥有窩.
人子卻沒有枕頭的地方.
原來.
你看耶穌基督.
他來了地上.
很簡陋的.
放在馬槽裡.
他死了.
也沒有.
墳墓.
他死是死在別人.
願意把安葬的地方.
給他.
這個墳墓不是他的.
是一個有錢的藥室.
他買了,沒用過的.
他讓耶穌安葬.
他連死也沒有地方.
原來我們這個心智.
就是每一個跟從耶穌的人.
原來要預備.
一個受苦的心智.
跟從耶穌不簡單.
你看到耶穌基督.
由他一生裡.
他無論是.
出生.
或是到死.
原來都不是屬於他的.

$^{521}$你願不願意跟從他呢.
在文士.
原來在《馬太福音》八章十九節.
說到這個人.
無論你去哪個地方.
我都要跟從你這個人.
原來在《馬太福音》看就是文士.
文士.
都願意跟從主.
文士他們是很熟悉.
聖經裡面是律法書.
很熟悉教人.
但他竟然被耶穌基督的權柄.
吸引了.
他願意去跟耶穌.
在當時期.
一些希臘的哲學家都會巡迴.
講說.
耶穌基督是有徒弟的.
但耶穌他竟然公然的.
去召門徒.
他召門徒.
在這裡這個人回應他.
但耶穌叫他考慮清楚.
人子是沒有枕頭的地方.
飛鳥都有他的鬥.
狐狸都有他的鬥.
但耶穌沒有枕頭的地方.
耶穌基督.
正是告訴我們要付代價的.
你不要因為.
他很有權柄.
當然會有很多人.
去敬重他.
會聽你的.
會敬重你幾分的時候.
你物質生活.
都應該是沒有問題.
但原來.
跟從耶穌的人要留意.

$^{561}$耶穌說你不要將他.
跟那些有權柄的人同等.
你不要期望跟隨他.
就可以得到.
世上的福氣.
你會得到高的職位.
或者你的生活會富裕.
其實是剛剛相反.
他要我們預備一個.
受苦的心智.
然後耶穌基督.
又呼籲人去跟從他.
這個人.
跟耶穌說什麼呢.
他說容許我.
先回去埋葬我的父親.
耶穌說讓死人埋葬.
他不就是死人.
你只管去.
全港神的國.
這句話.
其實當然埋葬.
是表示我們.
對我們的先人.
一個孝敬的表示.
孝敬父母在律法書.
都是要我們孝敬的.
但這句話不是表示.
他的父親剛剛去世.
而是可以理解到.
要求一些年老的父親.
真的希望.
理解到他們.
等到父親離世之後.
再去服侍你.
是這個意思.
在阿拉伯人家裡.
家裡有老人家.
他們都不會遠遊.
他們.

$^{601}$或者在經文會說.
這個人.
為父親進行第二次哀悼期.
可能老人家.
他已經死了一年.
我們為他哀悼之後.
就去舉行一個儀式.
就是去埋葬.
孩骨.
究竟是怎樣呢.
聖經就在經文裡.
沒有交代.
無論如何.
由耶穌的回應裡.
你會看到.
死了一個人.
如何埋葬死人呢.
這句話好像很誇張.
耶穌基督不是教我們.
不可以孝敬父母.
因為耶穌基督上十字架的時候.
祂交給祂最親愛的約翰.
去照顧祂的媽媽.
所以在這裡看到.
耶穌也是很孝順父母.
也很孝敬.
但是這裡好像很誇張.
去現實.
讓我們看到.
怎麼可能會發生呢.
只不過提醒門徒.
其實你有其他方法.
去做重要的事情.
有其他方法可以處理.
但是你要很關注的.
就是.
最重要的事情.
我們要有優先的次序.
怎麼去做呢.
曾經有一個姐妹.

$^{641}$她去做.
現在人們很喜歡旅遊.
坐船.
有一個姐妹.
她去坐船.
在旅遊途中.
醫生跟她說.
什麼都不能吃.
只能吃柑.
因為市議員每天都拿柑給她.
但是船差不多走到.
離開市鎮的時候.
市議員就說.
柑已經吃光了.
現在船已經走到一半.
已經走到一半.
都不能買了.
這個姐妹怎麼做呢.
這個姐妹就去祈禱.
她安靜地去神面前祈禱.
她不去疑惑.
神的應許.
她向神祈禱.
祈禱的時候.
就聽到遠方有些船的.
汽笛的聲音.
她聽得很清楚.
大概過了一個小時.
市議員就抱了很多柑給她.
發生什麼事呢.
市議員就說.
現在有很多柑給你.
原來那艘船迷路了.
大霧的緣故.
迷路了,他們沒有水.
水都用光了.
我們就去供給他們.
他們就報答我們.
把很多柑給我們.
於是就把柑給我們.

$^{681}$於是我就有柑可以給你.
大家看到就覺得很奇怪.
但這個姐妹做了什麼呢.
雖然現實就沒有柑吃.
但她專心的去仰望天父.
去祈禱.
她去祈禱的時候.
原來神有不同的方式.
去供應給她.
是不是很奇妙.
但我們在人生裡.
很多事情不是照你的版本去做的.
很多時候會有些事情出現了.
但我們不可以說.
我的劇本是這樣寫的.
為什麼不是這樣做的.
我們不可以橫行霸道.
神給我們過的人生.
是很有人性的.
是很特別的.
但當你遇到一些特別的時候.
你是更加有可能.
去遇到一些特別的事情.
你會更加有可能.
去留心神有什麼吩咐給你.
你要有信心去踏上.
你要去祈禱.
我們只可以靠祂.
當你什麼都幫不到你的時候.
只有神在你身邊.
成為你隨時的幫助.
你就要依靠祂.
你不要再想其他事情.
其他事情是沒有辦法的.
柑是現實沒有的.
但很奇妙.
神就為她預備.
祂就在你身邊.
神就為她預備.
你見到就是.

$^{721}$原來主耶穌.
祂要走十字架道路的時候.
祂終於被殺.
祂知道祂將要死.
但祂找接班人.
誰是祂的接班人呢.
祂在找.
很緊張的時刻.
所以叫到你的時候.
就是你有資格.
但你還有很多推搪.
很多藉口的時候.
是不應該的.
你對祂有多少信賴呢.
有人.
耶穌呼召他.
叫他跟從祂.
祂說容許我先去.
辭別我家人.
耶穌跟他說什麼.
守父就利.
向後看的人不配見.
進神的國.
你見不見到.
原來守父就利是什麼意思呢.
原來.
以前的人.
現在就說.
破拉機.
但以前不是.
利是用來耕田的.
但不是很容易做的事.
因為你要.
很專心的向前看.
還有你的雙手要緊握著利.
不可以分心.
你耕出來的路.
才會不直.
很直的路.
是很重要的.

$^{761}$主耶穌用利.
去耕田.
有個意思.
告訴你.
你跟從祂的人.
當你接受耶穌基督的時候.
你就是.
拿著利.
你要信服祂.
你要聽祂的吩咐.
如果我們.
時常回頭望的時候.
你的路就不直了.
我們不要回頭望.
我信了主之後.
我們又要.
去崇拜.
又要讀經.
又要靈修.
又要教主學.
又要回主學.
這麼多的事.
又要去做私事.
沒有了我們自己的私人空間.
你會想很多事.
或者.
不如意的事出現的時候.
走回以前.
走回頭路.
好像以色列人.
神帶他離開埃及.
因為那時候被人壓迫.
離開埃及.
一個遺老的地方.
去到曠野.
出現問題的時候.
他又懷念.
在埃及.
又有西瓜吃.
又有黃瓜又有肉吃.

$^{801}$又不用花錢就有魚.
我們今天會不會.
顧影自憐呢?.
我們會不會走回頭路呢?.
好像他們.
類似的想法.
耶穌警告我們.
守乎特利.
向後看的人.
不配進神的國.
是不是真的不可以回頭呢?.
其實可以回頭的.
不過叫你不要.
留戀你過去的罪.
記得不要留戀你過去的罪.
以前在罪終之落.
罪終的生活你不可以留戀.
反而在現今.
我們想想主怎樣.
帶我們脫離.
這個罪惡.
我們要想想.
我們的生命.
今天被主改變了.
我們有沒有感恩.
我們有沒有更加積極的願意.
去跟從祂呢?.
今天我們每個人很想得到神的祝福.
到底.
今天有什麼.
佔據了我們的心呢?.
上班的一族.
因為大部分人都還沒退休.
你們很年輕.
無論你是什麼行業.
你可能都會遇到問題.
那個時候你還會不會回教會呢?.
或者.
你是家庭主婦.
很多家務要做.

$^{841}$又要照顧孩子.
最基本的.
柴米油煙.
煎醋茶.
基本的.
這些都要去招呼周到.
我哪有空去崇拜.
學生呢?.
考試.
考不好.
所以很多人.
看不透.
就會解決了自己.
不懂做功課.
很多煩惱.
做學生有學生的困難.
不論你是小朋友還是大學生.
都有自己的困難.
做學生.
你遇到你的成績.
這麼差的時候.
我哪有空去崇拜.
可能連段階都不見你.
或者.
你是病人.
我記得.
我那時候.
還沒開始做電療.
那時候只是畫線.
我心裡就有些.
心理的問題.
開始會覺得.
去到那個地方.
會覺得想吐.
但去到洗手間.
真的會吐.
想想,糟了.
將來呢?.
怎樣去探病呢?.
原來每次來到這裡.

$^{881}$我會作嘔的時候.
我怎樣去事奉呢?.
原來事實呢.
做一做神.
再去那個地方.
真正去.
接受治療.
或者治療之後.
繼續去覆診.
原來神都可以讓你勝過.
或者你失業.
失業有很大的難關.
一家人都等著你.
失業了.
心裡很難過.
每個人都上班.
很無聊的,不知道做什麼好.
甚至每天都數錢.
現在經濟這麼差.
連夜市.
都沒有什麼人去光顧的時候.
真是.
怎樣度日子呢?.
種種的問題.
你不知道怎樣去面對.
我哪有時間.
去想主的事呢?.
我哪有時間去事奉呢?.
今天你會不會退後呢?.
當你有很多埋怨的時候.
你願不願意跟從主呢?.
你相不相信.
神會保守你呢?.
你相不相信.
祂在你身邊.
帶領你.
你仍然有出路呢?.
當你有很多埋怨的時候.
我們是不是應該.
將一些妨礙我們前進的東西.

$^{921}$我們要.
趕走它.
任何事情影響我們.
對主忠誠的時候.
我們必須.
放棄目前在.
阻礙我們的東西.
我們應該放棄它.
你願不願意呢?.
求主幫助我們.
當我們要自立了.
我們要成為.
嘉欣什麼?.
嘉欣堂.
大一點的.
一個名堂.
成為嘉欣堂的時候.
你的肩膀重一點的時候.
你要不要對主忠誠呢?.
你願不願意呢?.
要付出一個代價.
要預備一個受苦的心志.
你預備好了什麼呢?.
我們一起祈禱.
親愛的主耶穌.
你給我們這麼重的責任.
你必定給我們足夠的能力.
可以承擔的.
在預備的時候.
加添我們的信心.
讓我們勇敢的.
向前跨步.
我們不要回頭望.
我們要向前望.
因為你就在前面.
你會帶領我.
你會幫助我們.
你是我們腳前的燈路上的光.
你會指引我們的腳步.
讓我們不要看左右.

$^{961}$我們就專心一意的.
跟著你.
達到你要我們走的方向.
能夠去榮耀你的名.
求你幫助我們.
我們一起的.
齊心的.
告訴主,主我們願意跟從你.
求你保守我們的心志.
從現在開始.
直到我們見主面的時候.
我們都能夠榮耀.
歸給你.
奉主耶穌的名求,阿們.
\newpage



\section{以賽亞書 61:1}
\label{sec:p_xp8W8_F_Q}
\textbf{2023-12-31 午堂《禧年執行者》以賽亞書 61\_1 講員:楊明麗姑娘}
\newline
\newline
連結: \href{https://youtube.com/watch?v=p_xp8W8-F_Q}{\texttt{https://youtube.com/watch?v=p\_xp8W8-F\_Q}} ~~~~ 語音日期: 2024-01-03
\newline
\newline
\hyperref[sec:c3xZAeZuwD0]{\small{< < < PREV SERMON < < <}}
~
\hyperref[sec:index_chronic]{\small{[返順時目]}}
~
\hyperref[sec:index_scriptual]{\small{[返順卷目]}}
~
\hyperref[sec:R1nryPpcggU]{\small{> > > NEXT SERMON > > >}}
\newline
\newline
以賽亞書 61:1
\newline
\begin{longtable}{cl}
\hline
\hline
章節 & 經文 (和合本修訂版)\\
\hline
61:1 & \begin{tabularx}{0.7\textwidth}{X} 主耶和華的靈在我身上，因為耶和華用膏膏我，叫我報好信息給貧窮的人，差遣我醫好傷心的人，報告被擄的得釋放，被捆綁的得自由； \end{tabularx} \\ \\
[1ex]
\hline
\hline
\end{longtable}
$^{1}$各位早晨.
我們一起祈禱.
我們一起禱告.
我們來到你的面前.
我們好奇求你一件事.
求你幫助我們.
我們每一位無論講的.
聽的都得蒙你的恩典.
最後我們可以一起.
領受你的話語 一起學習.
我們仰望你 奉主耶穌.
基督的聖名 阿們.
第二樣我想請大家做的.
他們不單只是在我們當中.
陪我們一起崇拜.
我還看到有位.
其實不是在香港成長的.
他也是小朋友 幾歲大.
他從外地回來.
竟然可以和我們一起.
讀經 讀中文.
我們欣賞一下我們身邊的小朋友.
今天在我們身邊的小朋友.
可能將來.
你和他的關係會更加密切.
我不知道.
你有沒有哪些小朋友.
你特別多接觸.
我有教過.
幾堂兒童主義學.
我在攝堂.
我一直都是成年區多.
有一位.
在.
兒童主義學.
其實不是一兩個.
其實到他一路成長.
我應該是教過他兩三堂.
但是有一天.
我在地鐵看到他的時候.

$^{41}$他是一個中學生.
我不認得他.
他認得我 他和我打招呼.
到現在.
我是他的粉絲.
他每次報告天氣.
我一定會看.
為什麼呢.
因為有一個關係在.
他是我兒童主義學的學生.
我探訪過他家.
他竟然成長了之後.
他會認得我.
他會很記得這件事.
現在早上.
報告新聞.
有天氣報告的時候 我一定會看.
在我們當中.
會有一些不同的人.
這些不同的人 我們會接觸到.
也有機會會繼續接觸到.
或者在不同的時刻.
我們都會接觸到.
而今天這段經文.
我是祈禱了才找這段經文的.
我甚少說.
特別倚賽亞書.
一本很困難的聖經.
這本書.
有一節.
有一章是61章.
一節到62章.
三節 通常都是在.
聖誕後.
聖誕之後的禮拜.
會說的 是說神的.
恩典臨到.
文字是.
第一生的.
第一生就說.

$^{81}$主耶和華的靈在我身上.
要宣告一件事.
就是神高立了他.
這段聖經為什麼會.
放在聖誕節之後.
其實是在表達.
這段聖經是在說.
耶穌基督.
就是被高立的那一位.
這段聖經.
其實有幾重的意義.
第一.
他提醒猶太人.
告訴他們.
希臘你們忘記了嗎.
為什麼你們.
不去遵守希臘呢.
第二 他是說到.
以色列人 猶太人.
快要被釋放了.
因為他們被擄.
但是以賽亞書的時候.
其實他們還沒有亡國.
是一個預言來的.
這個預言是指向.
是.
古列王.
波斯王古列.
要釋放猶太人回去.
那個時候.
他們好像回到自己的地方.
重建聖城 重建聖殿.
那個希臘.
但是最重要.
他是指向主耶穌基督的降臨.
主耶穌基督是我們的拯救.
他將恩典.
對次序給我們.
所以我們好像有一個恩典.
但是這個恩典不是一年.

$^{121}$是長遠的 是一個希臘來的.
給我們很多的恩典.
這個希臘在舊約當中.
是有這樣的指引的.
在利未記.
他說到希臘是五十年一次.
在五十年一次的希臘裡面.
你要讓.
人能夠得到.
恩典得到自由.
意思就是說 如果有人.
欠了你什麼.
到了希臘那一年.
就不要再跟他計較了.
如果有人 是猶太人.
是住在你當中的.
已經虧損了的.
做了猶太教徒的.
你猶太人你要做一件事.
就是釋放他回去.
他賣了兒女給你嗎?.
賣了妻子給你嗎?.
到了希年的時候 你要釋放他回去.
讓他一家團聚.
他沒有錢.
抵押了房子給你嗎?.
你要在希年的時候.
還給他.
他不是.
他窮到什麼都沒有.
他只是抵押了.
一張被子給你.
他應該晚上就還給他.
給他晚上有被子.
還沒給嗎?.
你希年一定要完全不再跟他計較.
希年是有這個意思.
釋放.
次恩典.
讓他能夠.

$^{161}$得到他所需要的.
可以自由喜樂地.
去生活 可以感受到.
神的恩典淋到他身上.
主耶穌基督的來臨.
不單止是這樣.
主耶穌基督的來臨.
是特別一點的.
主耶穌基督的來臨.
是讓我們能夠完全地.
感受到他是從神.
父神那裡差遣來.
讓他傷心地得到.
醫治 是報告被擄的.
無論是被習慣.
擄了的.
或者是被人的惡言.
所俘虜了的.
都能夠得到釋放 得到自由.
在以賽亞書的書.
他還說到.
這一年是一個私恩的年份.
這一年是.
大家不需要工作.
都有神的賜福.
甚至乎是生活當中.
是有公義的一年.
但是還有一節的經文.
這一節的經文.
大家應該都很熟悉的.
跟剛才那節是很相似的.
是主耶穌基督說的.
在《路加福音》第四章.
不單止是說到耶穌基督遇到試探.
之後還說到.
他回到自己鄉下.
他成長的地方.
那撒勒.
在他成長的地方.
他進了那撒勒的會堂.

$^{201}$他可能就在那裡成長.
我們當中也有.
我們當中在一起崇拜的.
有些兒時就在那裡成長.
耶穌基督回到他成長的地方.
他去事奉.
人們給了他一本經書.
他就拿了一本.
以賽亞書的經書.
然後就讀了這一節的經文.
第61章的第一節.
他讀了這一節的經文.
他說神的靈.
在我身上.
因為祂用.
高告我.
叫我傳福音給貧窮的人.
差遣我.
宣告被擄的得釋放.
生命的得看見.
受壓迫的得自由.
宣告神.
決立人的欺憐.
然後.
他就開始和.
他在會堂裡的人.
來講述.
其實我們看聖經.
十九節之後.
還有兩節.
這兩節就是說.
他讀完經之後.
就把聖經放好.
就坐下來.
但原來坐下來在猶太人當中.
讀完經是坐下來的.
身邊其他人.
可能有些人站著.
不過讀經的人讀完經就坐下來.
這是猶太人會堂裡的規矩.

$^{241}$然後他就開始講道.
他坐下來講道.
這就和我們很不同.
找天試試.
牧師回來講道的時候.
牧師就坐下來講道.
我們站著聽.
他們猶太人讀經的位置.
是坐下來的.
他們坐下來的時候就開始講道.
政教教會.
都是這樣.
東正教教會都是這樣.
他們有一個位置.
叫做講道的位置.
是坐下來的.
象徵著這一點.
耶穌基督就坐下來講道.
他應該講了一段.
沒有記載下來的經文.
但是聽了之後.
這個講道的亞撒族人.
他們不是很接受.
你說你就是要來的彌賽亞.
我們很明白這段聖經.
是真的在說.
彌賽亞要來臨.
他要在我們當中.
讓我們能夠得到自由.
還要復興我們的國家.
這是猶太人的信仰.
他說你是誰呢.
你就是我們當中.
所認識的.
弱勢的兒子.
耶穌我們看著你長大的.
所以他們接受不了.
但是耶穌基督.
開始去宣講的時候.
他去告訴人家.

$^{281}$他是受高者.
他就去了自己的會堂.
受高者.
其實猶太人一直在等待.
但是當真正的猶太人.
當中的受高者.
去到他們當中.
他們竟然認不到.
也都接受不了.
甚至最後將耶穌基督釘十字架.
但是耶穌基督.
他本身就是這位彌賽亞.
我們叫他Jesus Christ.
耶穌基督.
基督這個字.
就是彌賽亞.
彌賽亞.
是主要來到.
讓我們能夠得自由.
聖經當中說.
主的靈在我身上.
因為他要用膏.
告我.
是耶穌基督.
是被神高納的那位.
是預言當中那位.
他就是受高者.
他的確來到.
讓人能夠得到自由.
讓人能夠從罪惡當中.
得到自由.
弟兄姊妹你們有沒有得到自由.
我今早和一位姊妹.
傾談到血壓.
其實血壓可能都.
提醒了我們我們是不是得了自由.
當你那天的心思.
意念不得自由的時候.
原來血壓是不一樣的.
我們星期四早上有.

$^{321}$量血壓 行程天地的.
服侍.
我也有被我們.
當中的護士姑娘們幫我量.
如果我掛著和人傾談.
原來都會高二三十度.
其實我.
血壓是很高的.
很高到很少點到一百二十.
所以一到一百二十我已經知道.
那天高了.
但是當我和人傾談的時候.
可以去到一百四十.
如果你心裡面.
掛著一件事.
前一晚睡不著.
那天早上血壓真的會高了.
我們是不是真的能夠.
得釋放.
可能看看你的血壓.
看看你早上起來的心情.
看看你晚上睡得好不好.
看看你對著面前的.
食物你知不知道自己在吃什麼.
我已經看到.
我們是不是得到釋放.
我們將我們的事交給神.
讓祂釋放我們.
在你的生活習慣當中有沒有得到釋放.
可能你現在困鎖我們的電話.
我們沒有電話就不行了.
有時可能你不帶身份證.
都可以.
不帶電話已經不行了.
是不是.
車的鎖匙在電話裡.
屋的鎖匙在電話裡.
什麼都在電話裡.
沒有帶電話一定要回家拿.
是啊.

$^{361}$我們被不同的東西.
困鎖住了.
還有一些東西我們被困鎖住了.
有沒有某些食物困鎖了你.
我們不說.
有沒有某些習慣困鎖了你.
求主幫助我們.
不要讓一些語言去困鎖了我們.
我記得有弟兄.
他信主之後.
他頭那一兩年很辛苦.
說一句話要停一停.
現在他不用了.
他現在可以很自由釋放地繼續說.
不用特別留意著.
會有什麼字眼會走出來.
因為他已經從這些字眼當中.
得到釋放.
求主幫助我們.
主耶穌基督來.
釋放了我們.
我們不要放棄了被釋放.
自由.
是不可放棄.
他不單止來讓我們.
從.
綑綁當中得到自由.
他還做了很多事.
他醫好了人.
他不單止醫好猶太人.
他特意去到.
一個地方.
醫好一個迦南人.
不是猶太人.
的女兒.
她女兒被.
邪靈附身了.
他去醫好她.
他去醫好盲眼的.
他不單止用說話.

$^{401}$你好了就走吧.
不是.
他用口水磨了泥.
塗在她眼睛上.
他用不著做那麼多動作.
但是.
他就是跟她有這樣的接觸.
甚至.
麻風病人來到他身邊.
他都去撫摸她.
告訴她.
我願意.
你得醫治了.
我剛才和大家說.
我特別喜歡聽某一個科學家.
報告天氣.
每一早.
我都等今天這個科學主任.
是不是我那位.
前主任學生.
其實現在變成我是他粉絲.
因為大家有一個.
關係.
有一個關係是不同的.
我特別喜歡聽梁麗麗說.
另外還有一個.
就是楊國忠.
楊國忠 為什麼呢.
楊國忠就不是.
我主任學生.
不過在我帶中學生.
主任學的時候.
邀請過他來.
和我們分享一些天文的東西.
這個是有一個.
關係.
你和他有關係的時候.
是不同的.
我家人說.
你那麼喜歡看.

$^{441}$你看三十一台有AI報告.
一定不會說錯的.
我和AI.
沒有關係.
他說得再好.
他的表情再怎樣.
都好.
我和他沒有什麼關係.
我不單聽那個報告.
我看到一個熟人.
弟兄姊妹.
耶穌基督來到這個世界.
他不單止.
完成了.
欺凌他作為受告者.
釋放我們得自由.
他做了一件很重要的事.
他和我們建立關係.
我們是什麼.
如果你喜歡做.
泥池 塗膠.
你都可以製作很多.
產品在家.
擺一下點一下.
你怎會對著他一路說話.
不是.
我們是受造物.
但是耶穌基督覺得.
我們是親切的受造物.
他愛我們 為我們付出感情.
為我們死於十字架.
他覺得我們和他有關係.
不單止我們覺得和他有關係.
他覺得和我們有關係.
多好啊.
我們變得多尊貴.
我們和這位真神有關係.
他賜給我們恩典.
給我們自由.
給我們釋放 還體貼我們.

$^{481}$我們有什麼.
真的要去問他.
有什麼我們真的要.
和他談一下.
他有什麼對我們的帶領.
我們真的可以和他談多一點.
我們懂的東西 告訴他.
他幫到我們.
我讀神學的時候.
要上聲樂.
聲樂的時候.
我唱歌是一個絕大的娛樂.
當中會有一位知道.
是怎樣呢.
我試過.
每個學期都要考試.
明知我不懂唱.
我都唱不到.
但總要在同學當中唱.
你怎樣都要畢業.
都要做.
那就做吧.
我記得第一次唱的時候.
那位同學是男同學.
是曾牧師認識的男同學.
他笑得.
比soprano還要高音.
差點掉在地上.
是啊 沒錯.
我是不懂唱歌的.
但有得求.
我事奉需要.
原來有得求.
弟兄姊妹 主耶穌基督很愛我們.
祂視你是親人.
親到一個地步.
祂可以為我們犧牲.
是十字架.
有什麼人和你親到一個地步.
當他需要一個神的時候.

$^{521}$你會驗血捐給他呢.
主耶穌基督就是知道我們有犯罪.
祂願意為我們捐出生命.
你願意為祂做些什麼嗎.
你願意嗎.
你說我不懂嗎.
放心吧.
你不懂的東西我一定不懂.
不過我們原來可以求.
原來可以求的.
可以求神會幫你.
讓你有機會學.
讓你有機會明白.
怎樣去做.
祂愛我們.
祂願意每一樣東西都幫我們.
讓我們都達成我們所願意做的.
祂是我們的主.
我們的神.
還願意做我們的父.
我們真的很執到.
主耶穌基督.
祂也不計較什麼日子.
去聽我們禱告.
在我們心裡.
祂也不計較什麼日子.
去聽我們禱告.
祂也不計較什麼日子.
去聽我們禱告.
在安息日的時候.
祂醫好有需要的人.
我想這個人的腳.
手肌肉萎縮了.
不是一天的.
可能都不止一年.
其實他萎縮多一天.
都可以的.
但耶穌基督接受不到他病多一天.
他就施恩典.
醫好他.

$^{561}$從病患當中釋放了他.
他在著緊我們每一個.
他著緊到一個地步.
就是他和門徒.
在安息日的時候.
走過山野的地方.
那裡有些不是山野.
那裡還有一些是玉米.
麥.
那些門徒就拿去.
麥來吃.
他就被別人質疑.
耶穌你是否拉比.
你的門徒在安息日.
竟然預備.
食物去工作.
飯安息日.
其實可能走多一段路.
就有食物.
但耶穌基督就是體貼我們.
不單止他的門徒.
他體貼我們每一個.
他為我們預備.
有什麼事.
我們就跟他講.
向他祈求.
讓他為我們預備.
無論白天還是晚上.
他都為我們預備.
讓我們能夠得到.
誰是欺凌的執行者.
當然是主耶穌基督.
有沒有錯.
是呀,他來讓我們能夠.
得到恩典.
讓我們能夠.
在生活當中得到釋放.
不單止是.
從病患當中得到釋放.
最重要是從罪惡當中.

$^{601}$釋放了我們.
讓我們能夠自由,不去犯罪.
對著一些以往.
我們會做的錯事.
我們會不斷地跟他說.
我們監控受你控制.
但他不單止這樣做.
主耶穌基督還是這樣說.
「你們去.
把所看見.
所聽見的.
告訴若翰.
就是盲人看見.
瘸子行走.
麻風病人得見證.
聾子聽見.
死人服沒.
窮人聽到福音.
咦?.
施洗若翰是耶穌基督的前驅.
但原來.
他有一個時刻.
他懷疑耶穌.
表弟,是不是你.
真是那個受告者.
我們有沒有.
信心會有動盪.
的一天.
有時遇到困難可能會.
那個時候施洗若翰.
被收監.
在困難當中.
主都看顧我們.
在患難當中他不離開我們.
我們要記得.
他在我們身上所做過的.
他在我們身邊的人所做過的.
今天下午有什麼事.
今天下午兩點鐘.
有見證分享會.

$^{641}$如果你覺得自己有.
心裡面有時會.
軟弱的時候真的要來聽聽.
無論什麼人.
遇到什麼事他都幫助我們.
他都聽我們禱告.
賜福給我們.
照料我們.
他不在乎我們是什麼身份.
窮人有福音.
死人有福活.
聾子才能夠聽見.
麻瘋病人得到潔淨.
瘸子行走.
盲人看見.
在你最困難的.
那一樣.
神最幫到你.
在你覺得.
最絕望.
那一樣.
就告訴他.
不要將絕望.
自己擁抱著.
讓主耶穌基督幫助我們.
他是.
欺凌的執行者.
不單止這樣 他還跟門徒說.
他說天上地下所有的.
權柄都賜給我了.
所以你們要去.
使萬民在我的門徒.
奉父子聖靈的名.
給他們施洗.
凡我所吩咐你們的.
都教導他們遵守.
看那我天天與你們同在.
直到世代的終結.
主耶穌基督.
都差遣我們.

$^{681}$你說欺凌的執行者.
他已經升了天.
那我們怎麼辦?不是.
主耶穌基督還留下我們.
每一個做欺凌的執行者.
他預了我們.
一起為他來做.
他所做的事情.
他也讓我們能夠.
得到聖靈的高滿.
聖靈在.
主耶穌基督.
升天之後.
又降臨.
在復活後的四十天.
主耶穌升天.
再過多十天.
聖靈降臨.
所以有聖靈降臨節.
就是復活節之後的七個禮拜日.
第七個禮拜日.
那天.
門徒.
聚在一起的時候.
他們一起領受到聖靈.
這個是.
十二位.
門徒.
當然不是猶大.
是另外有馬提亞.
來補了他的位置.
聖靈降臨.
不單只是當時的使徒.
而是我們每一位.
我們洗禮之後.
就好像主耶穌基督那樣.
洗禮之後聖靈也與我們同在.
一路與我們同在.
讓我們能夠一直.
得著他的恩典.

$^{721}$甚至乎.
你信主那一刻你也禱告.
祈求聖靈進入你的心中.
他們有聖靈降臨之後.
這一班的使徒.
耶穌基督的門徒.
在耶穌基督死了之後.
本來就很害怕的.
躲起來了.
領受聖靈之後他們就出去傳福音.
不單止出去傳福音.
是一次過有三千多人信了主.
三千多人表達.
這些全部是男士.
表達的是三千多個家庭.
是家庭.
是三千個家庭.
三千個家庭信了主.
然後他們就.
一起團契.
一起祈禱.
一起守聖餐.
弟兄姊妹.
主耶穌基督都差遣我們.
他都差遣我們.
他差遣我們.
讓我們又能夠得著他的恩典.
生命可以依靠他.
仰望他.
讓我們又能夠得著恩典.
我們得到聖靈在我們當中.
感動我們.
他也照顧我們.
讓我們每一個都能夠恩著他的緣故.
我們可以繼續好好去事奉.
而且我們還有一個使命.
我們一路事奉.
我們一路帶.
比我們更年輕的人一起去事奉.
你帶你的兒女.

$^{761}$你的孫回來.
是很對的.
你讓他們繼續在這裡.
來受栽培.
得著神的恩典.
你回來.
你讓你的團隊.
繼續學習在當中事奉.
你就是在做這個.
欺凌的執行者.
讓大家都能夠得著這個恩典.
我們在恩典當中.
可以同心合力一起去事奉.
弟兄姊妹.
耶穌基督還這樣說.
他說:我實實在在地報訴你們.
我所做的工作.
信我的人也要做.
並且要做得比這些更大.
因為我到乎那裡去.
耶穌基督.
你現在到乎哪裡去.
我們要做什麼大事呢.
我們要將耶穌基督所做的事.
那個範圍.
我們需要擴得更寬.
讓更多人能夠領受到.
讓更多人能夠得著.
讓到你身邊.
和你身邊你不認識的.
都能夠得著他的恩典.
這是他賜予給我們的使命.
他告訴我們.
你要靠著那住在我們裡面的聖靈.
牢牢守住.
所交託給你那美好的事.
讓人得著美好.
得著平安.
得著自由.
這是耶穌基督的使命.

$^{801}$交給了我們.
我們都成為希臘的執行者.
不好意思.
我想再退一步.
我們都成為希臘的執行者.
對了.
在這裡有一個心心.
特別送給大家.
為什麼呢.
今天我們會有很特別的宣佈.
大家都知道我們來到年尾.
明年又有一個新的開始.
我們要面對很多新的事物.
我們也會有很多不同的呼聲.
呼喚大家.
一起去事奉.
對呀.
我們應該一起去事奉.
耶穌基督讓我們成為希臘的執行者.
那應該去哪裡事奉呢.
無論哪一個地方.
哪一個部門.
哪裡找你.
或者感受.
或者雄恩.
或者嘉恩人找你.
我們都要做一件事.
我們都要將這件事交託給神.
我們要將這件事交託給神.
讓神去.
差派我們去做.
耶穌基督.
他差派我們.
也要我們牢牢守住這個本份.
讓我們繼續去做.
希臘的執行者.
你去哪裡做.
你真的要禱告.
神會說你知道嗎.
神有說我知道.

$^{841}$我記得有一次.
我應徵一間教會.
是一個很特別的經歷.
我寄信的時候.
我都會祈禱.
我都為這件事去禱告.
也都面試了一次.
我也去祈禱.
天父面試了.
大家都談了.
問了.
我問神說什麼.
是不是你要我去這間教會呢.
我想很清晰.
誰知道.
就是這樣.
祈禱出事.
應該大概就是.
那一兩天裡.
那位堂會的牧師.
也是宣道會的牧師.
現在退休了.
那位堂會的牧師.
就不見了我的信.
不見了我的電話.
我問神.
神說你怎麼會不見了.
我問神.
神說你怎麼會不見了.
我問神.
神說你怎麼會不見了.
我問神.
神說你怎麼會不見了.
我問神.
神說你怎麼會不見了.
我問神.
神說你怎麼會不見了.
我問神.
神說你怎麼會不見了.
我問神.

$^{881}$神說你怎麼會不見了.
我問神.
神說你怎麼會不見了.
我問神.
神說你怎麼會不見了.
我問神.
神說你怎麼會不見了.
我問神.
神說你怎麼會不見了.
我問神.
神說你怎麼會不見了.
我問神.
神說你怎麼會不見了.
我問神.
神說你怎麼會不見了.
我問神.
神說你怎麼會不見了.
我問神.
神說你怎麼會不見了.
我問神.
神說你怎麼會不見了.
我問神.
神說你怎麼會不見了.
我問神.
神說你怎麼會不見了.
我問神.
神說你怎麼會不見了.
我問神.
神說你怎麼會不見了.
我問神.
神說你怎麼會不見了.
我問神.
神說你怎麼會不見了.
我問神.
神說你怎麼會不見了.
我問神.
神說你怎麼會不見了.
我問神.
神說你怎麼會不見了.
我問神.

$^{921}$神說你怎麼會不見了.
我問神.
神說你怎麼會不見了.
我問神.
神說你怎麼會不見了.
我問神.
神說你怎麼會不見了.
我問神.
神說你怎麼會不見了.
我問神.
神說你怎麼會不見了.
我問神.
神說你怎麼會不見了.
我問神.
神說你怎麼會不見了.
我問神.
神說你怎麼會不見了.
我問神.
神說你怎麼會不見了.
我問神.
神說你怎麼會不見了.
我問神.
神說你怎麼會不見了.
我問神.
神說你怎麼會不見了.
我問神.
神說你怎麼會不見了.
我問神.
神說你怎麼會不見了.
我問神.
神說你怎麼會不見了.
我問神.
神說你怎麼會不見了.
我問神.
神說你怎麼會不見了.
我問神.
神說你怎麼會不見了.
我問神.
神說你怎麼會不見了.
我問神.

$^{961}$神說你怎麼會不見了.
我問神.
神說你怎麼會不見了.
我問神.
神說你怎麼會不見了.
我問神.
神說你怎麼會不見了.
我問神.
神說你怎麼會不見了.
我問神.
神說你怎麼會不見了.
我問神.
神說你怎麼會不見了.
我問神.
神說你怎麼會不見了.
我問神.
神說你怎麼會不見了.
所以,基督的信仰,阿們!.
\newpage



\section{撒母耳記上 15:17-23}
\label{sec:R1nryPpcggU}
\textbf{2023-12-31 早堂《作主聽命的兒女》撒母耳記上 15\_17-23 講員:曾敏姑娘}
\newline
\newline
連結: \href{https://youtube.com/watch?v=R1nryPpcggU}{\texttt{https://youtube.com/watch?v=R1nryPpcggU}} ~~~~ 語音日期: 2024-01-03
\newline
\newline
\hyperref[sec:p_xp8W8_F_Q]{\small{< < < PREV SERMON < < <}}
~
\hyperref[sec:index_chronic]{\small{[返順時目]}}
~
\hyperref[sec:index_scriptual]{\small{[返順卷目]}}
~
\hyperref[sec:qkHV9uoLu6Y]{\small{> > > NEXT SERMON > > >}}
\newline
\newline
撒母耳記上 15:17-23
\newline
\begin{longtable}{cl}
\hline
\hline
章節 & 經文 (和合本修訂版)\\
\hline
15:17 & \begin{tabularx}{0.7\textwidth}{X} 撒母耳說：「你雖然看自己為小，你豈不是作了以色列諸支派的元首嗎？耶和華膏你作了以色列的王。 \end{tabularx} \\ \\ \relax
15:18 & \begin{tabularx}{0.7\textwidth}{X} 耶和華差遣你，吩咐你說：『你去除滅那些犯罪的亞瑪力人，攻打他們，直到把他們完全滅盡。』 \end{tabularx} \\ \\ \relax
15:19 & \begin{tabularx}{0.7\textwidth}{X} 你為何沒有聽從耶和華的話呢？你為何急著撲向掠物，行耶和華眼中看為惡的事呢？」 \end{tabularx} \\ \\ \relax
15:20 & \begin{tabularx}{0.7\textwidth}{X} 掃羅對撒母耳說：「我聽從了耶和華的話，行了耶和華派我行的路，擒了亞瑪力王亞甲來，滅盡了亞瑪力人。 \end{tabularx} \\ \\ \relax
15:21 & \begin{tabularx}{0.7\textwidth}{X} 百姓卻從掠物中取了牛羊，是當滅之物中最好的，要在吉甲獻給耶和華－你的神。」 \end{tabularx} \\ \\ \relax
15:22 & \begin{tabularx}{0.7\textwidth}{X} 撒母耳說：「耶和華喜愛燔祭和祭物，豈如喜愛人聽從他的話呢？看哪，聽命勝於獻祭，順從勝於公羊的脂肪。 \end{tabularx} \\ \\ \relax
15:23 & \begin{tabularx}{0.7\textwidth}{X} 悖逆與占卜的罪相等，頑梗與拜偶像的罪孽相同。因為你厭棄耶和華的命令，耶和華也厭棄你作王。」 \end{tabularx} \\ \\
[1ex]
\hline
\hline
\end{longtable}
$^{1}$弟兄姊妹平安.
每一次回來錦繡堂.
都很有那種回娘家的感覺.
覺得非常親切.
雖然在座當中.
我未必每一個都有機會和大家聊天.
有機會建立關係.
但是大家都是同樣在神的家裡面.
很重要的一員.
首先在這裡將洪恩家那份問候和祝福.
帶給大家.
願大家繼續在神的國度裡面.
一同去打拼.
我們開始的時候.
繼續一同去禱告.
讓我們回到你的家裡面去敬拜讚美你.
是一個何等大的美好的福氣.
我們卑微的人.
竟然可以來到你的面前.
這樣去親近你.
稱呼你為我們的爸爸.
在我們生命裡面的每一天都可以經歷你.
你有次在你的話語成為我們腳前的燈.
路上的光.
願意你對我們眾人說話.
讓對友誼聽的就應當聽.
願聖靈在我們眾人生命裡面去動功.
叫到我們有信服的心.
有回應的心去回應神的教導.
奉耶穌基督的名字祈禱.
阿們.
其實今天要講的經文.
遠遠比剛才要讀的經文要長.
我不知道大家前面還有沒有一本聖經.
在你的座位前面.
如果有的時候.
我想請大家去翻開.
三號議機上十五章.
待會我們有機會去看看.
剛才那段經文上下文所講的內容.

$^{41}$如果不是的話.
就這樣聽我講.
也是可以的.
今天來到本年度最後的一個星期日.
也是最後的一天.
我們真的心願上帝的說話帶領我們.
不單止是完結這一年.
更加是對新一年給我們一個很大的挑戰.
在十五章裡.
其實剛才那段經文上文是在講一個很重要的故事背景.
故事背景是在講.
有一件事發生了.
有一場戰役.
可以說首先是來自一個命令.
上帝的吩咐.
然後一場戰爭發生了.
那個戰爭是有結果的.
在裡面我們卻看到.
上帝所揀選的君王.
在裡面是失敗了.
這就是頭一段十五章.
是一去到第九節的地方所講的.
簡單地去看的時候.
那裡是這樣說.
撒母耳去到蘇羅.
蘇羅是以色列第一任的君王.
上帝揀選他作王.
在這裡撒母耳跟他說.
是上帝差遣他去高納蘇羅王.
所以去高納作王這件事.
不是撒母耳這個先知個人的意思.
是神的揀選.
神自己的意思.
神親自揀選他的君王.
去治理他的百姓.
上帝賦予這個君王有權柄.
很大的權.
去管理整個以色列國.
但在裡面上帝有他的期望.
有他的吩咐.

$^{81}$上帝的期望是很重要.
就是先知所講的這一句.
現在你要聽從耶和華的話.
在這裡神並沒有說.
到時你管理以色列民.
你要怎樣做.
首先第一步就要做這些.
第二步做那些.
第三步就這樣做.
你要跟你的軍隊.
要怎樣開會商量.
要搞些什麼.
或者跟你的百姓.
要計劃些什麼事情.
上帝都沒有說這些.
上帝只講一件事.
一件事 唯一的一件事.
最重要的那件事.
這是上帝最期望.
他所揀選的君王要做到的.
就是要做什麼.
要聽從耶和華的話.
最重要是這樣.
今天我和你都是屬於.
上帝所揀選的人.
我們可能有不同的崗位.
有不同的服侍.
我們有不同的成長背景.
但都是屬於上帝所揀選的.
是很尊貴的 很榮耀 很寶貴的.
身體我們很珍貴.
我深信同樣的原理都是.
上帝不會在說.
今天你就要12345678.
有不同的事情你要進行.
我要和你說.
你就要和這個人說這些事.
和那個人做些什麼事.
說 說 說 說 說.
說一大堆事.

$^{121}$不是 上帝對我和你都是重要的.
有一件事就是這件事.
要聽從耶和華的話.
你要聽從耶和華的話.
最重要這件事.
聽啊 聽神聲音啊.
放下我們自己的計劃.
放下我自己個人的看法.
我今天看下去在座.
我們都是很有經驗的弟兄姊妹.
無論在我們人生的工作 家庭.
各方面的經驗.
非常豐富.
深信大家對教會都可以有很多不同的貢獻.
上帝可以大大的去用大家.
但原來神第一件事說的不是.
用你的經驗.
使用你一直對不同事情的用法.
你怎樣善用在這些事工 不同的活動裡面.
這不是上帝對我們第一件事要我們注意的事.
第一件事都是.
聽他的話.
他的話是指什麼.
什麼是上帝的話.
上帝的話可以指的是整本聖經的教導.
神就在聖經裡面跟我們說話.
所以我們經常祈禱都說.
神的話語是我們腳前的燈 路上的光.
怎樣行事為人 怎樣準則.
上帝喜悅的是什麼 不喜悅的是什麼.
怎樣去跟隨 神有一個帶領在這裡.
有準則的.
但不只是限於在這裡.
上帝對每一個人都有自己的帶領.
對我 對你.
在不同的時間.
在神的時間 神會有帶領.
我今天一路去思考這篇信息的時候.
我就想起我過了紅恩堂已經滿了三年的時間.
原來時光飛逝 這麼快就過了幾年.

$^{161}$幾年裡面有很多經歷.
感覺那幾年好像很長時間.
由一開始到現在.
其實都經歷了很多的挑戰.
很多的風浪.
真的很多.
在裡面 神的話就成為我很大的提醒.
怎樣跟人相處.
怎樣面對那些風浪.
這個是其中一個帶領.
但另一個神亦都親自對我說話.
我很感恩這幾年在紅恩家裡面.
多了很多聽到神的聲音.
神真的在裡面對我說話.
每一步去到那個位置.
我應該用一個什麼態度.
對待不同的人 不同的事.
神都會說.
在過程中 神想我去聽.
什麼叫聽.
聽了就要去跟.
我真的發覺人生是很不同的.
事奉是很不同.
神是說話的.
每一個處境 神都知道怎樣走下去是最好.
而且神有自己的旨意.
對每個人都有自己的旨意.
在裡面祂想我聽 跟著走.
每一步祂自有祂的鋪排.
所以我感到很豐富.
心裡面很雀躍.
我不知道我見到頂音姐妹.
剛才都有些頂音姐妹跟我打招呼.
你有沒有感受到我有些事不同了.
我知道自己有一個很大的轉變.
因為上帝的說話實在太重要.
上帝的命令實在太重要.
我也見證到一件事.
不是自己經常拿著經驗一套套地套下去.
這個就是A的處境.

$^{201}$套下去就這樣做.
這個是B的處境.
做開就這樣做.
就這樣套下去.
不是 每一次都不同.
神說在這個處境你要這樣做.
那我就要這樣去.
其實不用這麼辛苦.
這麼用自己的力量.
神有命令 神也有恩典.
神自然會幫助我去做到祂要我做的事.
看回這段經文.
都是很能夠看到這件事.
上帝就是通過撒母耳去吩咐蘇羅.
要聽從耶和華的話.
最重要是這件事.
接下來很重要的命令.
就是要去面對亞馬利人.
上帝的審判臨到亞馬利人身上.
他在說 以色列人從埃及上來的時候.
在路上亞馬利人怎樣待他們的話.
這裡怎樣去抵擋他們.
我都要懲罰.
神的審判在這一刻臨到亞馬利人身上.
上帝藉著以色列人去施行審判.
就是要去攻打他們.
神命令他的百姓去攻打他們.
而且不單止這樣.
上帝的命令是怎樣的.
是滅盡他們所有的.
不可憐惜他們.
將男女孩童吃奶的.
以及牛羊駱駝和牢全部殺死.
聽清楚這個命令.
不是只聽某一部份.
只聽攻打他們.
其他就不聽了.
不是的 所有事情要滅盡.
女老幼大大小小不同性別.
包括所有的財產.

$^{241}$牛羊所有的全部滅盡.
是一丁點而不正.
這是神完整的命令.
在吩咐這個命令之前.
上帝透過撒母耳說.
要聽他的話.
上帝的吩咐是很清晰.
蘇羅不可能聽不清楚.
他不可能自己聽力有問題.
或者理解力有問題.
一定聽得清楚.
在這裡聽到這個命令.
所以蘇羅是有回應的.
第四節說蘇羅開始召集百姓.
準備去打這場仗.
然後你看到很多人.
召集了二十萬步兵.
還有一萬的猶大人.
一路進去打仗.
還有策略地打仗.
在這裡當中.
蘇羅想避免無辜的人受牽連.
所以他叫一些基尼人.
離開亞馬利人.
讓他們不用受牽連.
他們就專心地對付亞馬利人.
其實亞馬利人是什麼人呢.
他們是以素的後裔.
以色列人是雅各的後裔.
雅各和以素是兩兄弟.
大家還記得在創世記的時候.
說他們之間的衝突.
是非常嚴峻的衝突.
去到他們的後人.
這些衝突都沒有停止過.
亞馬利人究竟是誰呢.
就是以素的兒子.
以理法的妾婷納所生.
他們主要居住的地方.
就是在南地和西奈的曠野.

$^{281}$這一班人是會傷害以色列民.
傷害神的選民.
在《申命記》第25章這個地方.
大家聽我分享就行了.
17到19節這個地方.
上帝要以色列民.
總有一天要去對付亞馬利人.
大家聽我讀出來.
你要記得你們出埃及的時候.
亞馬利在路上怎樣對待你.
在路上迎擊你.
趁你疲乏困原時.
擊殺所有在你後面軟弱的人.
並不敬畏神.
所以當耶和華你的神.
使你不避四圍一切仇敵擾亂.
在耶和華你神賜你遺孽的地上.
得享平靜的時候.
你要把亞馬利的名從天下逃去.
你不可忘記這事.
上帝吩咐過.
要以色列民有一天.
要執行上帝的審判.
要去對付亞馬利.
現在就是這個時候.
要去攻打他們.
亞馬利人就是這樣行惡.
傷害上帝的選民.
上帝很愛自己的選民百姓.
誰去詛咒傷害他的子民百姓.
上帝都不會放過他們.
誰去祝福他的百姓.
上帝都會祝福.
你記得這個昔日上帝對亞巴拉罕.
有一個很大的應許.
誰去祝福你的.
我都會祝福他.
誰詛咒你的.
我都會對他不好.
神是這麼愛他的百姓.

$^{321}$這裡展現這一點.
這裡說的是.
素羅召集兵隊.
軍隊去打仗.
在這個過程.
很快看到戰果.
是如期地順利.
他去攻打亞馬利.
完全可以將他們打敗.
在第八節.
生擒了亞馬利王亞格.
用刀殺盡了亞馬利的眾百姓.
這個戰果很快出現.
打贏.
一點難度都沒有.
為什麼.
為什麼一點難度都沒有.
是不是以色列民很能打.
不是.
你記得在出埃及記.
有一場戰爭.
約書大令所揀選的.
以色列民出去和亞馬利爭戰.
摩西在山上禱告.
記不記得這一幕.
摩西何時舉手.
那場仗何時贏.
何時垂手就輸了.
證明一件事.
神真的要以色列民倚靠上帝.
戰敗在於.
成敗關鍵在於神自己.
亞馬利不是想像的.
不是我們一般覺得他那麼軟弱.
其實他是一個很好鬥的民族.
很好戰的民族.
打仗很厲害.
在這裡說真的.
面對以色列民.
我不相信亞馬利是沒有辦法打贏.

$^{361}$但是這裡這麼快贏了是因為什麼.
上帝將敵人交在以色列民手裡.
是神令他們得勝.
因為神吩咐的.
神吩咐他們要去攻打亞馬利人.
上帝已經在當中.
會給以色列民有力量.
上帝掌管這場仗最後輸贏.
所以以色列民贏了.
蘇羅王贏了.
直到這一刻都還好的.
但是接著很快就看到.
這一位上帝所揀選的君王.
很快就看到他是怎樣去回應神的吩咐.
直到那一刻打勝仗還是很好.
接著你看到第九節.
蘇羅和百姓卻連飾亞格.
就是亞馬利的王.
愛惜上好的牛羊牛羊.
高揚以及一切美物.
不肯滅絕.
神說要什麼?.
毀滅所有人.
男女老幼大大小小.
不同性別.
不同階層.
是所有人一個都不留.
他們在這一刻卻愛惜亞格.
在經文沒有說.
什麼原恩惠愛惜亞格呢?.
亞格有些魅力.
是什麼原因沒有說.
你只看到他們在這裡卻是手軟.
不只是眾百姓.
經文說得很清楚.
蘇羅自己也是這樣.
君王和百姓一起.
不聽上帝的話.
你記住之前他聽得清清楚楚.
知道神的吩咐.

$^{401}$但這一刻馬上轉了態.
他愛惜他不單止這樣.
還有其他東西.
還要愛惜上好的牛羊牛羊.
高揚以及一切美物.
上帝說明所有東西都要滅絕.
一件事都不可以拿.
那些有生命的全部要殺掉.
在這裡他馬上轉了.
不只是眾百姓.
蘇羅自己也有份.
都留住了那些很好的牛羊.
那些不好的.
軟弱的牛羊自然就殺滅.
但這些不是.
全漏了.
這裡說他不肯滅絕.
在知情的情況下不肯滅絕.
在這裡我們看到.
為什麼整件事忽然風回路轉.
本來有一個很好的勝利.
歡喜快樂的.
靠著上帝得勝了.
信仰他們得勝.
本來在這裡要一同慶祝.
在這裡忽然間整件事不同了.
專為很黑暗.
這就是我們之前那段經文.
一個開始的背景.
是一個失敗.
上帝有期望.
是他所揀選的人聽從他的命令.
但是他所揀選的人卻失敗了.
失敗了之後.
其實神仍然在給人機會.
有一個正確的回應.
中間這一段.
10到23節.
這裡就說一個.
不是那麼理想的回應.

$^{441}$10到16節這裡就說到.
上帝的聲音臨到撒母耳.
不是撒母耳發現了.
才發現.
這個素羅王失敗了.
和百姓一起得罪上帝.
不是.
是神的聲音臨到.
我們跌倒.
我們軟弱的時候.
不要以為只是等某一天.
某個人.
某個弟兄.
某個姐妹.
某個傳道.
某個人發現了我跌倒了.
我失敗了.
不是的.
首先發現的是誰.
上帝自己.
上帝知道我們失敗.
上帝知道素羅失敗.
所以他對撒母耳說.
我立素羅為王.
我感到遺憾.
其實這個遺憾是指.
神很大的失望.
你說神不是全能嗎.
你選擇的時候.
你就知道他軟弱.
有一天他會跌倒.
那你就不要選擇他.
為什麼你會覺得遺憾呢.
你也知道他會跌倒.
但你看見神對人很大的憐憫.
我真的可以說.
神對人很大的憐憫.
仍然是憐憫素羅.
仍然去揀選他.
仍然期望他聽自己的吩咐.

$^{481}$走在他所定的道路上.
很可惜.
人錯過了上帝的揀選.
神所要帶來的福氣.
反而背叛上帝.
人不能得著神所賜給人的福氣.
被神揀選是很蒙福的.
我想說.
之所以為揀選.
不是每個人都會坐那個位.
不是每個人都會做那個角色.
能夠有這樣的使命照命.
被揀選是何等的蒙福.
但素羅他卻錯過了這個福氣.
神在此感到很難過.
很難過 很大的失望.
神自己去評價他.
因為他轉去不跟從我.
不遵守我的命令.
所以你看看神最重要.
最看重的是否聽他的命令.
最重要是聽他的命令.
上帝一路看著.
今天我和你在教會裡服事.
或者我們回到家裡.
去到工作環境裡.
在鄰居當中見證他的時候.
上帝每一下都看著我們.
看著我們是否聽他的命令.
是否跟從他.
還是跟從自己的那一套.
我們很懂得說很屬靈的話.
信主越來.
那些屬靈的話越好聽.
朗朗上口.
彼此說很屬靈.
這個弟兄 這個姐妹.
這位領袖.
很屬靈.
最重要不是說那些話.

$^{521}$最重要是我們走出來.
上帝給我的評價.
是否我聽他的話.
最重要是這一點.
我可以說很屬靈的話.
但我的行事為人.
完全違背上帝的教導.
我可以做一些看似很好的事.
但上帝沒有吩咐我.
上帝吩咐我另一樣事.
我偏偏不跟從.
我就按照自己的經驗.
我覺得好就去走.
都是好的.
怎樣都好.
神就由得我.
我喜歡這樣跟就這樣跟.
上帝一定讚我.
不是的.
神有神的安排.
神才是我們的主.
我們的王.
記不記得我們剛剛信耶穌的時候.
我們做過決志禱告怎樣說.
所有弟兄姊妹記不記得.
我邀請耶穌基督.
進入我的生命成為甚麼.
我的救主.
和生命的主.
甚麼叫生命的主.
上帝你說你的.
我說我的.
有空我向你匯報.
我做些甚麼.
我相信我做些事好.
你都讚我的.
這樣是不是叫生命的主.
不是的.
生命的主不是我作主.
是祂作主.

$^{561}$大小事情.
你想想到今天.
真的聽祂的話.
是不是真的聽祂的話.
真的值得我們去思考這個問題.
在這裡神很大的失望.
在這個訴羅轉去.
不跟從祂.
不遵守祂的命令.
撒母耳很生氣.
很生氣.
很生氣訴羅為甚麼要違背上帝的命令.
為甚麼不聽上帝的話.
撒母耳站在神那裡.
他的心很難過.
很激氣.
他中夜哀求耶和華.
這個哀求是在說他哭.
他哭.
他哀求耶和華甚麼.
他憐憫訴羅.
是他背叛上帝.
不聽上帝的話.
在這裡他清早起來.
就要作上帝的僕人.
走去跟訴羅說話.
訴羅去到加密的時候.
看到這裡有人告訴訴羅.
告訴撒母耳.
訴羅去到加密這個地方.
他在那裡為自己立了紀念碑.
於是就下吉甲.
撒母耳一定是下吉甲去找訴羅.
撒母耳知道訴羅為自己做了甚麼事.
為甚麼要立紀念碑.
因為神經常記住他.
多偉大的功績.
我打敗亞馬利亞.
亞馬利亞人搞擾我們以色列民.
我現在是我.

$^{601}$訴羅打敗他們.
我立紀念碑就是你們要記得我.
做一件這麼偉大這麼好的事.
其實說真的.
我剛才一開始就說.
打敗亞馬利亞是上帝的工作.
沒錯上帝派訴羅去打這場仗.
訴羅有付出代價.
軍隊有付出代價.
但是是神令這場仗得勝.
要紀念的是誰的榮耀.
上帝自己的榮耀.
將榮耀歸給神.
不是歸給自己.
經常歌功頌德自己.
我是怎樣 我是這麼好.
所以在過程中你看到一個生命.
一個真正的生命.
是不是一個榮耀上帝的生命.
還是每一刻為自己.
其實我們的生命.
由信耶穌那一刻.
我想跟弟兄姊妹說.
我相信大家都很知道這件事.
不是再屬於自己了.
我們的生命活著都不是為自己.
不只是傳道人生命.
才是為上帝而活.
是弟兄姊妹的生命.
都是為上帝而活.
活著就是要榮耀上帝.
今天能夠做任何事.
是慎自己.
是施恩憐民 賜下能力.
不是我們這個人可以做些甚麼.
所以要將榮耀歸給他.
所以你看到蘇羅的生命.
一直在跌倒 一直在下滑.
聖經留下了這個註腳.
就是讓我們看到.

$^{641}$這個生命要讓我們引以為鑒的地方.
好了 撒母耳去見他的時候.
你看到很有趣.
一去到蘇羅那裡.
蘇羅就真的不覺得自己有甚麼問題.
他就說 願耶和華賜福給你.
耶和華的命令我已經遵守了.
向人稱讚自己.
我遵守了神的命令.
你之前叫我遵守.
我遵守了 我做得很好.
很希望別人稱讚自己.
他完全沒有自知之明.
自欺又欺人.
在這個過程裡.
撒母耳就說.
我耳中聽見有羊叫 有牛鳴的聲音.
又是甚麼呢.
是甚麼來的.
他還很想蘇羅自己知道自己失敗了甚麼.
就好像在創世記的時候.
亞當夏娃犯罪之後.
他們躲避上帝.
上帝去呼喚他們 去找他們的時候.
其實神還在等待人主動跟上帝說.
我犯罪得罪了你.
上帝求你寬恕我.
他想等人主動說.
但人還是不說的.
直到神了解人們的一些事.
人們才會透露一些事.
最後就死死地氣.
被神發現了.
就知錯了.
被神發現了.
但神不是自己主動承認.
蘇羅這裡不覺得自己有問題.
當撒母耳給他一個機會.
你承認吧.
聽聽羊叫牛鳴是甚麼意思呢.

$^{681}$還在說.
這個時候開始推卸責任.
這是百姓從阿瑪利人那裡帶來的.
因為他們愛惜上好的牛羊.
他都知道自己不對.
但他要先解釋一下.
我有理由的.
我留下他們.
要獻給耶和華你的神.
其餘的我們都已經滅盡了.
要做就做吧.
但為甚麼會留下呢.
我要解釋一下.
就是要怎樣呢.
獻祭.
我就是為上帝你好.
所以我做的事是好的.
不是為自己的.
神你知不知道.
我們事奉裡面都可以這樣說.
神有一個吩咐.
有一個大令.
我未必會跟隨.
我另外覺得那件事.
我覺得是好的.
上帝為你做.
你應該懂得欣賞.
如果不是為你.
我為甚麼會這樣做.
這是好的.
是為了這樣我才做這件事.
為自己去解釋.
為甚麼我不聽上帝的命令.
上帝是否覺得這樣就很好呢.
整件事.
上帝要施行公義的審判.
在每個人的身上.
祂要祂的子民百姓做的.
就是要徹底的.
執行上帝的命令.

$^{721}$施行審判.
這個命令是很清晰的.
但掃羅為首的整個百姓群體.
都沒有真正徹底地聽上帝的吩咐.
在這裡還有很多的解釋.
是掃羅都借百姓那裡去過橋.
然後撒末爾真的忍不住.
叫他住口吧.
你閉嘴吧.
閉嘴吧.
讓我把耶和華昨晚對我說的話.
告訴你.
知不知道上帝怎樣看你.
今天剛才那段經文.
17到23節.
掃羅看自己很小.
看自己很小.
甚揀選他做以色列的君王.
仍然覺得自己很小.
但其實這個何等大的榮耀.
一點都不小.
是一個很大的榮耀.
上帝揀選他.
給他一個很重要很尊貴的位份.
是不是.
他說你不是做了以色列諸支派的元首嗎.
又說高尼作了以色列的王.
這麼重要.
神這麼體重你.
然後神給你一個吩咐.
要去除滅那些犯罪的亞馬利人.
所以這絕對是一個審判.
要攻打他們.
要完全滅盡.
你為什麼沒有聽從耶和華的話呢.
你為何急著撲向那些略物呢.
你很緊張.
說得好聽.
為上帝為憲制.
為了憲制.

$^{761}$有多迫切.
在舊約裡面.
的確上帝設立不同的憲制.
那個制度做什麼.
是教導他的百姓.
怎樣去和上帝自己相處.
去相交.
人和神之間有一個很大的鴻溝.
我們不可以直接地.
和上帝有一個相交的生活.
神就教我們透過憲制.
向他有不同的表達.
神對於憲制有很詳細的教導.
很仔細.
怎樣獻法.
獻什麼.
神的要求是什麼.
很仔細.
如果大家看利未記的時候.
看到神的要求.
我們也會看到.
神是很體重憲制.
但是在這裡.
撒慕爾很清晰地說.
所有這麼體重的東西.
都不及聽從上帝的命令.
不及.
聽從上帝的命令是很重要.
是超越一切.
今天神也很喜悅我和你.
在教會裡面服事.
上帝召集我們在一起.
敬拜讚美他.
我們各自忠心於自己的崗位服事.
在這裡一起建立神的家.
是何等美好.
是很重要的.
上帝對於我們服事.
祂也有不同的期望.
但是我想說一件事.

$^{801}$在這個之上.
是上帝要我們聽他的命令.
所以這裡說得很重要.
神在責備裡面說了一件事.
要說喜愛繁祭和祭物.
喜如喜愛人.
聽從他的話.
看那聽命勝於憲制.
沒錯.
神很喜歡我們憲制.
但是神很想我們聽命令.
聽從勝於公羊的知方.
是何等重要.
我可以參與很多的事奉.
很多的活動.
我這個人很熱心.
很忙碌.
很多外在的事我都做足了.
還不夠你這個心.
真真正正聽上帝命令嗎.
整個人真正的生命是屬於上帝嗎.
如果我只是有外在的這一切.
很多的工作.
事奉 會議 活動.
但是我沒有聽上帝的命令.
上帝也不喜願.
而且看到好蟲.
神看到那種順蟲是甚麼.
這裡說的搏抑和占卜的罪相等.
占卜的罪和邪靈占卜上的關係.
和拜偶像占卜上的關係.
神是何等的憎惡.
這些是在屬靈上何等的污穢.
他說搏抑和這些罪相等.
是這麼大的一件事.
還敢和拜偶像的罪相同.
神是多麼憎恨這些東西.
所以弟兄姊妹.
我們來到本田最後一個主日的時候.
回顧我們過去一年的生命.

$^{841}$我們的生命是否一個聽從上帝命令的生命.
真真正正整個人屬於上帝的生命.
是否一個軀殼.
一個可以有很多活動的.
很多事奉的.
很多工作的.
很多外在的事.
我做得很足夠.
別人都覺得我很屬靈的.
但我裡面是一個軀殼.
這樣的生命要小心.
上帝不會喜悅我們所拜上.
而且上帝看這件事很嚴厲.
很嚴重.
神為甚麼要在聖經這裡放下這段經文.
如果大家有機會再看新約的教導.
神真的很看重我們.
是否聽從祂的命令.
不只是單一的事件.
單一的教導.
神是很看重我和你的生命.
在這裡神宣判.
神要透過訴羅 以色列民審判亞馬利人.
事件發展到一個地步.
是神要審判訴羅.
因為你厭棄耶和華的命令.
耶和華也厭棄你作王.
之後我們知道神揀選另一位君王.
讓他興起.
在未來的日子.
代替了訴羅王.
那位是大衛王.
上帝要去審判訴羅王.
但你看到這樣之後.
訴羅雖然有兩次認罪.
他也不徹底.
悔改也不徹底.
最後還是薩姆爾埋尾.
是薩姆爾去殺死阿格王.
所以很特別.

$^{881}$如果訴羅真的知罪.
認罪不只是神寬恕我.
我得罪你.
我現在馬上做.
把牛羊全部滅絕.
我馬上做.
我把阿格拿過來.
我要殺死他.
他沒有.
沒有說過有任何行動.
這麼有趣.
只是嘴說認罪.
所以也很危險.
我們信耶穌之後.
真的知道我們生命有得罪上帝.
我們認罪.
上帝寬恕我.
寬恕我.
這個不足夠.
在嘴裡.
悔改.
徹底的悔改.
我在什麼地方得罪上帝.
就真的悔改.
不要再犯那些罪.
本年到最後一天.
不單止是感恩.
也要問神.
靠神光照我們.
是不是有罪.
未曾徹底認罪悔改.
那些罪也會難了我們.
去領受上帝更進一步的祝福.
難了我們更好的伺服神.
小心我們做的事.
上帝不會紀念.
會的.
因為上帝有上帝的期望.
上帝有上帝的要求.
在這裡.

$^{921}$這是很重要.
訴諾王.
結果成為被上帝撇棄的一個王.
被上帝偏心.
是訴諾王.
沒有去真正跟從主.
聽命令.
卻不去跟從神.
神很認識人.
看穿,看透我們的人.
神看他是厭棄他的命令.
用這個字眼.
很輕看上帝.
看自己的榮耀.
比上帝的榮耀大.
今天弟兄姊妹問我們自己的服事.
是否真的為上帝的榮耀.
我也試過問自己.
我也試過去回應上帝.
對我的提醒.
在紅印堂這幾年的事奉裡.
我剛才說到有很多風浪.
每一步神都會跟我說話.
我發現一件事.
聽命是重要的.
每一步我聽了他的聲音回應.
神就會開路.
而且我看到那條路開得很直.
很舒暢.
下一步是很順暢.
去到下一步再來.
又是聽他的命令.
他又要去開路.
我知道如果我不信服.
後果是很嚴重.
我知道.
我的生命一樣要被上帝雕琢.
弟兄姊妹想跟你說.
我今天跟你說的堂道.
不只是鼓勵大家認罪.

$^{961}$遲些離開罪惡.
我也跟自己說.
我今年裡也有不同的時候.
向上帝認罪.
求神寬恕我.
想想我自己的動機.
是否有動機為了什麼.
榮耀誰.
是否真的榮耀神.
我都問自己.
我是否每一件大小事情.
去尋求上帝心意.
現在說聽主的命令.
是否真的聽到.
很多時候我們連尋求都不尋求.
我都不知道神對我的心意是什麼.
聽都沒有聽過.
我就按照自己.
覺得一直是這樣就去.
上帝應該可以的.
聖經有些共同原則.
跟那些原則就可以.
其實不是的.
只是看了這裡.
也要問神.
對我個人.
我事奉我生活.
我家庭的帶領是什麼.
是否每個位置尋求神心意.
我都問我自己.
所以在新的一年.
我自己都要悔改成長.
更多尋求神的心意.
更多的真的要放下自己.
聽上帝的命令.
我挑戰我自己.
同樣的我挑戰你.
做聽命的兒女最重要.
其他事情都不擔心.
所有事情神都會預備.

$^{1001}$不是這麼勞累的.
很累,很有壓力.
不是這樣的.
所以神跟我說.
都是跟大家說.
在新的一年可以活得不一樣.
做主聽命的兒女.
不是這麼辛苦,這麼淒涼.
反而是喜樂的,有力的.
精彩的,豐富的.
在這裡祝福大家.
也都心願大家同樣祝福我.
在新的一年.
活得精彩,活得豐盛.
兔住喜悅.
我們一起禱告.
在這裡憐憫我們.
我們不可以說自己比素羅好很多.
回顧這一年.
可能有很多時候我們都很自我.
很自私.
永要自己.
甚至沒有想過去求問你的心意.
沒有想過要跟你的命令走.
我們做基督徒這麼多年.
很習慣就這樣生活下去.
但主其實你對我們.
是充滿了你的期盼.
你有你的心意.
天父求你寬恕我們的生命.
是一個這樣的光景.
求你叫我們每一個都悔改.
活著是為榮耀你.
活著是為你而活.
活著是為了成就你的旨意.
讓我們在新一年活得不一樣.
我們的生命是徹底地向你開放.
讓你掌王權.
你讓我們聽到你的聲音.
明白你的教導.

$^{1041}$你的帶領.
讓我們順服.
讓我們每一個都活得精彩.
討你喜悅.
深圓有一天.
我們有這個福氣.
去見你面的日子.
我們不會羞愧.
而且上帝你會讚我們.
又良善又忠心的僕人.
好心願我們在世的每一天.
因為聽從你的命令.
我們祝福很多很多人的生命.
祝福那些未信耶穌的人的生命.
祝福我們《格尼論社》的生命.
祝福我們家人的生命.
我要讓上帝來得榮耀.
求主幫助我們.
奉耶穌基督的名字祈禱.
阿們.
\newpage



\section{瑪拉基書}
\label{sec:qkHV9uoLu6Y}
\textbf{2024-01-04 周四晚查經班《瑪拉基書(1)》 講員:曾裔貴牧師}
\newline
\newline
連結: \href{https://youtube.com/watch?v=qkHV9uoLu6Y}{\texttt{https://youtube.com/watch?v=qkHV9uoLu6Y}} ~~~~ 語音日期: 2024-01-05
\newline
\newline
\hyperref[sec:R1nryPpcggU]{\small{< < < PREV SERMON < < <}}
~
\hyperref[sec:index_chronic]{\small{[返順時目]}}
~
\hyperref[sec:index_scriptual]{\small{[返順卷目]}}
~
\hyperref[sec:SWoA7YUZj_4]{\small{> > > NEXT SERMON > > >}}
\newline
\newline
$^{1}$我們一起禱告,感恩我們可以再一次打開聖經去研讀.
多謝您讓我們能夠看完哈該書,撒加尼亞書.
都是對將來的預言和當時重建聖殿的光景和百姓的心情.
也藉著先知對將來末世的預言.
雖然很遙遠,但也指向終極.
就是主耶穌再回來的時候腳踏橄欖山的情景.
我們很盼望將來以色列人會得救.
有一些剩餘的以色列人.
是心靈改變來歸向主的.
到時真是很多人歸耶和華為聖.
和會投奔神的山.
所以我們也盼望為將來去準備.
雖然我們不知道哪一個年代耶穌基督會再回來.
但既然這是一個肯定的事實的時候.
就叫我們在我們自己的時代有限的人生.
來為神的國度做準備,求主帶領.
我們今天晚上開始最後這一本書.
《馬拉基書》.
盼望裡面的訊息有些都很淺白.
但很有提醒.
因為敬拜不是一兩天的事.
不是一次的聚會,是一生的敬拜.
所以我們需要不斷有更新.
盼望這一本書成為我們大家很多的提醒.
我們仰望在你的面前.
感恩禱告奉耶穌基督的名.
阿們.
雖然我們完全不知道.
在作者馬拉基寫這本書的時候.
到底是幾時的那個年代.
因為每一個年代都有它相關的歷史背景.
所以越是準確.
就越是能夠知道當時所默示的.
預言的,所討論的.
到底是發生什麼事呢.
大部分的解經家都認為.
馬拉基書既然擺在那個正典的最後.
又或者是一個歸回的歷史的那個階段.
就是在以色列人.
過去我們講過有幾次的回歸.

$^{41}$第一次就是這個索羅巴伯和大祭司約書亞.
還有這個.
以斯拉第二次.
尼希米的第三次.
那麼夾雜第一次的回歸.
索羅巴伯回到去沒多久又遇到困難.
過了15年左右.
哈該先知和撒迦利亞先知.
來到去有神的默示.
來預言,安慰,勉勵索羅巴伯.
來重建這個聖殿.
輾轉也都完成了.
後來就有以斯拉的回歸.
他就不是去重建聖殿.
因為已經建好了.
他也都不是建造那個城牆.
他回去其實就是一個祭司去教導律法書.
所以以斯拉的大部分是.
尤其是第九章之後.
就是他們在那個水門事件.
來到去復興.
接著輾轉他們就要進行摩西的律法書.
這個就是以斯拉第二次的回歸.
那麼夾雜著的就是第三次.
尼希米.
尼希米相隔和這個索羅巴伯.
和哈該的歷史.
相隔有拉比就認為超過120年了.
就是說過了一代甚至兩代.
這個就是尼希米第三次的回歸.
這些都是我們間歇地說過了.
最重要的就是要有自己的城牆.
就是說要有自己的surrounding的時候.
就是有門有戶了.
就是說是很徹底的在舊約來說.
就是和外邦人當地的迦南人.
正正式式有界線了.
正正式式我們在耶路撒冷.
才是聖城才是屬神的子民.
所以你們在城外的就是外邦人了.

$^{81}$所以尼希米他是做了這個很大的工程.
就是要將聖潔的子民和非猶太人的外邦人.
做了一個切割分開.
當然遇到的困難就是.
正正就是你沒有這個分別的時候.
猶太人甚至那些祭司和外邦人.
又會有很多很多的聯合.
一向是這樣的.
習慣了的.
安息日方便.
就是這頭手安息日一完就開門.
就叫那些人進來了.
有些人甚至乎都漏夜偷偷地進來了.
就變成了只是為了生活的方便.
就沒有想到自己是神選民的身份.
來去敬拜.
所以就夾雜了.
尼希米就很果斷地.
就禁止那些外邦人安息日之前進來.
好像潛進來那樣.
來去引誘那些心都不專心.
來去在安息日都想著去買東西.
這個是尼希米所做的.
所以去到這個《馬拉基書》.
我們看到是有那個敬拜.
而且是很多經文去責備那個敬拜的生活.
當然他們所謂的敬拜.
因為還有聖殿.
所以他們一定是在說憲制的.
因為去到關於耶穌的年代.
沒錯有聖殿.
但是都有它一定的限制.
但是在現在這個《馬拉基書》的年代.
他們沒有那個政治的迫害.
他們是恢復了聖殿的各種儀式.
和節期的敬拜.
所以在整個《馬拉基書》.
感覺到政治方面是沒有壓力的.
敬拜也是自由運作的.
這個大前提要記住.

$^{121}$因為正正是恆常的日以繼夜的這樣的憲制.
也就是說日久見人心.
日久見敬拜者的心情.
所以在行文裡面.
作者沒有因為一些困難,逼迫.
就去勉勵信徒,猶太人.
不需要懼怕,不需要膽怯.
這些在哈該書有的.
在撒迦利亞書也有的.
因為明顯是有逼迫.
可能是現在的.
但是更多是由第九章到十四章我們說過了.
更多是將來的災難.
你回去聽回那些錄音的信息.
但是在整卷的《馬拉基書》.
完全沒有外在的威脅.
和政治的困難.
令到你不能夠敬拜.
令到你敬拜有很大的成本.
令到你需要排除萬難.
甚至生命威脅來去敬拜.
因為猶太人敬拜神,憲制是天公地道的.
是他們理所當然的.
就是說他們恢復了聖殿敬拜之後.
就是一直一直他們不敢遺忘的.
不過正正就是這樣的時候.
有一些事情已經恢復了.
但是作者這位先知.
在這個恢復的過程裡面.
就開始有神的感動,默示.
來對於自己的同胞.
尤其是他們的敬拜生活.
來提出一些問題.
來點出他們的錯誤.
用什麼方式去點出呢?.
這本書的特色就是.
似乎有五至六段的對話.
是神用第三身.
以馬拉基作為第一身.
來代表神去詢問以色列人.

$^{161}$不同的崗位的以色列人.
開始是很廣泛的.
就是一般性的以色列人.
但是慢慢就很入獄了.
對號入座了.
一開始第六節就講到祭司了.
他們那個獻祭.
代表神去收納百姓的獻祭.
出了什麼問題呢?.
所以馬拉基書.
雖然我們不應該這樣.
從一個主題去讀馬拉基書.
但是如果從一個敬拜者.
現代的敬拜者.
都是一個朝聖獻祭的行動.
我們就可以從馬拉基書.
可以吸取到一些屬靈的原則.
到底神喜悅什麼的敬拜呢?.
什麼心態的敬拜呢?.
在這本書又會浮現了很多.
神喜悅,神不喜悅的對比,對立.
所以這裡又可以讓我們看到.
在先知的對話裡.
讓我們看到原來神很緊張哪些事情呢?.
我數過大概五到六個.
可能我今晚或明天時間鬆動一點.
我就將那個大綱放在群組裡面.
讓大家先有這個概念.
是一個很奇怪的行文的對話.
就是說神這樣理解以色列人.
以色列人就說我不覺得是這樣.
有什麼理據說我們是這樣呢?.
是一個很有趣的對話.
就好像揭示以色列人.
你們做得不對了.
你們的態度不正確了.
你們對我,對神有不公敬了.
這是一個點出了問題.
你要有證據的.
以色列人就說證據何在?.

$^{201}$是這樣的對話.
所以如果你用一個大綱這樣看它.
是挺有趣的.
我們進入經文.
四章很容易讀的.
這本書裡面也有些是灰暗書引.
尤其是第四章說的那個以利亞.
主耶穌說.
你所說的那個以利亞.
其實就是施洗約翰.
《馬可福音》裡面提過的.
所以我明天就給大家一些知道的資料.
所以其實如果原來你不看.
這些十二小先知.
或者歸回的先知書.
其實新約的作者是引很多的.
很多的預言.
很快讀第四章.
第三章先.
萬君之耶和華說.
我要差遣我的使者.
在我面前預備道路.
你們所尋求的主.
必忽然進入他的殿.
立約的使者.
就是你們所仰慕的.
快要來到.
這裡耶穌就定性了.
就是在說施洗約翰.
他就是那一位的先鋒.
來去說的那個以利亞.
或者是第四章.
第四章的第六節.
看!耶和華大而可畏之日未到以前.
我必差遣先知以利亞.
到你們那裡去.
他必使父親的心轉向兒女.
兒女的心轉向父親.
免得我來就坐遍地.
直接就是路加福音第一章.

$^{241}$就已經引述了.
就已經預言在施洗約翰身上.
直接引這一段聖經.
你可以看到.
其實先知書都充斥著.
彌賽亞的那個影像.
不過未全面的實現出來.
福音書就當他們看到.
耶穌的種種的生活行徑.
和他的講論.
很多年之後當耶穌升天了.
他們就有聖靈工作了.
他們就真的知道.
原來耶穌基督就是應驗.
先知哪一卷書哪一句話.
所以他們引述那些經文.
連路加醫生他不是猶太人.
他都懂得引述.
馬拉基書的四章六節.
說到大家如果不熟的話.
有時候就變得不夠立體.
一章的第十七節.
可能讀多一點.
反正今天是一個序言.
一章的第五節開始.
撒加尼亞按照著班次.
即是說他是大祭司.
按照著抽籤的班次.
即是輪班去聖殿燒香.
燒香即是一定是大祭司.
大祭司可以進去聖殿裡面.
即是說原來施洗約翰.
是出身自名門望族.
是亞倫的後人.
連他太太即是他媽媽.
這個以利沙白都是亞倫的後人.
即是根正苗紅的意思.
即是我們今天說的富二代.
施洗約翰就是祭司的第二代.
撒加尼亞按照著班次.

$^{281}$在神的面前供祭司的職份.
照祭司的規矩祭籤.
得以進到主的殿燒香.
所以大祭司燒香的時候.
眾百姓在外面禱告.
似乎是一個很嚴肅的節期.
會不會是一個贖罪日呢.
祭司一個人在裡面.
大祭司只能夠一個人在裡面燒香.
猶太人的規矩傳統.
就是有些靈當.
我們上星期說過.
那隻馬的靈當.
歸耶和華為聖.
大祭司他們進入聖殿.
那件袍子是有些金靈當的.
鈴鈴鈴鈴的.
為什麼呢.
因為他進去聖殿.
如果整個小時都沒有動靜.
即是說他已經死了.
沒有聲音.
因為睡了覺.
但裡面是沒有床沒有椅子的.
即是聖殿裡面有聖所.
有致聖所.
唯一沒有的就是沒有椅子.
因為是要暫著事奉的.
這是聖殿裡面的特色.
不會讓你休息一下.
在裡面休息一下.
是沒有的.
還有猶太人的傳統.
祭司進去的時候.
先綁著繩子.
因為其他人是不能進去的.
所以如果他很久都沒有動靜.
就將他拉出來.
拉出來.
所以是很有趣的.

$^{321}$這樣事奉的.
這裡告訴我們.
第十一節.
有主的使者站在香壇的右邊.
向他顯現.
很特別.
在聖殿裡面.
其實在香壇裡面.
如果在聖所.
都還有光.
因為有金燈台.
但是香壇在希伯來書說.
已經在致聖所裡面.
致聖所裡面是沒有光的.
是黑漆漆的.
是不會有光的.
所以在裡面燒香.
但是在燒香的時候.
這麼神聖的.
叫做Holy of Holies.
致聖所的英文.
神的使者.
天使.
在香壇裡面顯現.
怎樣進去的呢.
因為進去是要經過門簾的.
有一個莫萬的.
是遮掩著聖所和致聖所的.
這個就是一個意象.
一個很大的神蹟.
顯給撒迦利亞看.
接著留意.
撒迦利亞看見.
驚慌害怕.
因為你在致聖所裡面.
見到天使.
即是你突然間在黑房裡面.
有些人嚇一嚇你.
你都會害怕.
何況這個是天使.

$^{361}$驚慌害怕.
第十三節.
天使對他說.
撒迦利亞不要害怕.
因為你的祈禱.
已經被聽見了.
祈甚麼禱呢.
為誰祈禱呢.
他的責任當然是為百姓祈禱.
但似乎不是.
似乎是他個人的祈禱.
為家人祈禱.
因為你的妻子以利沙白.
要給你生一個兒子.
你要給她起名叫約翰.
即是司祭約翰.
你必歡喜快樂.
有許多人因她出世.
也必喜樂.
老年得子.
撒迦利亞.
他在主面前將要為大.
探酒,濃酒都不喝.
探酒,濃酒不喝.
是代表甚麼.
歸臣為聖.
拿勢以仁.
就是民數記第六章.
有些人就是.
那些叫做甚麼.
逐家弟子.
剃度那些.
要commit可能兩個星期.
我出家.
去保蓮寺吃齋.
剃了少少頭.
然後就還俗.
即是那些明星.
原來聖經也有.
你不是祭司就不可以入聖殿.

$^{401}$但你可以歸臣為聖.
你就要分別自己.
你和誰起了誓.
這個月我要做拿勢以仁.
你有手續的.
要找祭司的.
何時何時開始.
你就清酒濃酒都不可以喝.
即是代表著.
你對於俗世的事務.
暫時放下.
所以原來有時我們去夏令會.
去夏令營.
一個星期又不看電視又不看手機.
天主教很多的.
很多這些退休的.
真的整個星期都不准看手機.
或者可能某一段時間准你看.
然後你整個人就.
沉醉在那個安靜.
默想敬拜.
聽聖經一起祈禱.
這樣的氛圍.
有根據的.
不是沒有根據的.
所以將來的施洗若翰.
他整個人生就是拿勢以仁.
他整個人生就是.
原來分別為聖給神.
所以他本來是.
祭司大祭司的家族.
即是一定會繼承.
在耶路撒冷最崇高的職份裡面服侍.
但是當那個召喚來的時候.
第三章路加夫一味講.
神的話臨到他.
他就去到那個曠野.
約旦河外.
在那裡施洗.
預備人的心歸向神.

$^{441}$所以他一生做一個拿勢以仁.
有祭司不做,做拿勢以仁.
這裡告訴我們.
他的出現當然會令到很多人開心.
從他的母腹裡面.
這是第十五節.
從他的母腹裡面.
就已經被聖靈充滿.
在媽媽的肚裡.
就已經有聖靈充滿.
施洗若翰.
不是媽媽有聖靈充滿.
是兒子有聖靈充滿.
這些很奇妙的事.
所以伊底莎白去找瑪利亞的時候.
瑪利亞的肚子在動.
很開心.
因為瑪利亞的肚子裡面.
那個更加重要.
就是耶穌.
所以第一章.
一個醫生記載這樣的胎動.
神跡的胎動.
聖靈充滿的胎動.
很特別.
但是因為這個人是神.
已經在舊約.
正正就是我們將來讀的四章六節.
告訴我們.
這個人.
他的出現.
就有以利亞的心智能力.
以利亞什麼的心智呢.
我們有沒有說過.
以利亞那個聖經人物.
說了是吧.
以利亞的心智.
就是說他獨排眾議.
整個阿哈的王朝.
都是拜巴力.

$^{481}$我堅決不拜.
基列寄居的提斯比人.
以利亞.
就和這個阿哈說.
這幾年如果我不祈禱.
就不會下雨了.
說完就自己走去那個基納溪.
那個地方來住.
直至溪水都完了.
神就說你去這個撒勒法.
撒勒法就是阿哈的太太.
耶洗別的家鄉.
就是你搞笑去踩場.
去人家那個外邦的地方.
香火那麼鼎盛的.
拜巴力的.
你就要在那裡.
還要住在寡婦那裡.
就是神好像用一些很特殊.
沒可能的.
那些烏鴉.
啖那些肉給以利亞.
饑荒喔.
饑荒.
鳥都不會和你這麼生性的.
就算不是饑荒.
都不會這麼生性.
她很生性.
以利亞的烏鴉.
有時不請自來.
自己每天準時擔肉給她.
令到她可以三年.
饑荒都不會餓死.
接著那個饑荒真的遍地了.
就不再住在基納溪了.
就去到這個撒勒法了.
又是饑荒的.
以利亞說不要緊.
你弄給我先.
接著就沒事了.

$^{521}$吃完那些最後一次.
弄個餅給她.
那個地方.
那個寡婦.
因為有這個神人在那裡.
又可以繼續來生活.
這個就是以利亞的心智能力.
當然就是神給她在那個時代.
要百姓撥亂反正.
即是普遍拜巴力.
拜耶和華是不夠力的.
因為拜了耶和華.
都沒有保證.
保障我們有雨水.
有那個農地的豐收.
所以拜耶和華不夠力.
要拜巴力這樣的意思.
但是以利亞有那個能力.
就是神使到她.
一個人祈禱就可以不下雨.
所以我們今天靈恩派.
就很崇尚以利亞.
祈一祈禱就不下雨.
我再祈禱就下雨.
這樣和阿哈挑戰.
所以還要等阿哈.
其實阿哈已經坐著那些車.
可能馬車.
趕入城.
趕入皇宮.
以利亞就快過她入宮.
因為那些雨開始來.
如果她報了一個假案.
原來是巴力令我們下雨.
就令到那些以色列人又拜巴力.
所以以利亞快過她.
不知道是跑馬拉松.
還是有對風火輪.
不知道.
快過阿哈入到城.

$^{561}$這個就是以利亞很特別的能力.
我們讀到下面.
第十六節.
她要使許多的以色列人迴轉.
歸於主他們的神.
這個誰會喜歡呢.
你要以色列人迴轉歸誰呢.
歸以色列神.
即是迴轉這個transformation.
就是心靈的迴轉.
誰願意歸向神又不用付代價.
每個人都願意.
但是歸向神迴轉.
即是要付代價.
即是耶穌說的你們要背十字架.
走窄路.
人到我者裡來.
若不愛自己愛父母.
過於愛我.
就不可以跟隨我了.
這樣十四章說的路加福音.
所以這裡告訴我們.
她要使許多的以色列人迴轉.
你可以想像得到.
不是叫就行.
是要帶來那個.
可能是審判.
可能是眼淚.
可能就是受到那個逼迫.
施洗約翰警告希羅底和希律.
遭遇那個殺頭.
那個頭真是送給他的太太.
這個就是以利亞.
如果這樣事奉.
必然會招致那個後果.
所以十七節.
她必定有以利亞的心智能力.
強調是心智能力.
似乎不可以倒轉.
要有心智才有能力.

$^{601}$有時要有一個立志.
我要決心跟隨神.
你就會自然有一種.
得到神給你的恩典.
可能你裡面有很深的愛.
對別人的關心.
那個福音的緊張.
心智和能力不可以分開.
要不要就不行.
你要能力.
但不要的心智.
那就是你炫耀你自己的神蹟.
那你變成神人.
即是你變成神醫.
心智和能力是分割不到的.
所以她有以利亞的心智.
即是說以利亞自己也要分別為性.
好像一個怪人一樣.
一個人住在基納溪.
然後走到外邦.
走到巴黎的根據地.
這樣去生活.
說些話做些話怪怪淡淡.
但這個是神給她的心智.
才有能力.
施西漢就是這樣.
他住在哪裡.
你是耶路撒冷的祭司.
為什麼你搞成這樣.
約翰.
你住在曠野.
吃的是蝗蟲野蜜.
第三章路加福音記載他.
然後就在曠野施洗.
曠野是很荒涼的地方.
猶大曠野.
真的很荒涼.
你想想這種的心智能力.
即是說他已經將自己分別為性給神.
他自己預備好.

$^{641}$神的時間到.
我就要傳這個.
可以叫人回轉歸向神的福音.
當然這個回轉歸向神.
因為耶穌還沒有來.
耶穌還沒有成就救恩.
所以是將人帶去耶穌那裡.
所以他是告訴人.
我是施洗的.
我不是基督.
我不是那個主角.
我是一個配角.
你們不要搞錯了.
我是將人帶去耶穌那裡.
去到耶穌那裡受洗.
這個就是施洗約翰.
他自己恰如其分.
對自己的生命的認識.
你可以想像到.
下面這裡就很重要了.
就是我要講的這一段.
他行在主的面前.
叫為父的心轉向兒女.
即是說父母跟子女信主.
當時當然不是叫信主.
歸向神.
叫背逆的人轉從二人的智慧.
背逆神的人.
跟隨轉變來得到二人的智慧.
即是說原來信耶穌才是真正的智慧.
背逆的人.
即是說一直都是敵黨神的.
即是自己就是那個主人.
但現在有一個智慧.
就是二人的智慧.
就是你們這一班信耶穌的.
似乎表面上什麼都好像不准做.
煙又不准吃.
走過頭轉站又不可以買六合彩.
或者又不可以進麻將館打幾圈.

$^{681}$各方面都不行.
不是喔.
你們有一種智慧是二人的智慧.
背逆的人要轉向二人的智慧.
所以你想想施洗約翰是多麼的powerful.
來抗論,施洗.
有很多人去他那裡.
接受他的洗禮.
這個就是叫背逆的人.
轉從跟隨二人的智慧.
所以其實我們信耶穌.
我們成為基督徒.
我們要追求對真理的全面.
當然你的全面.
按照神給你的不同崗位.
不同的恩賜.
我們要盡量.
我們要比較準確地去解釋聖經.
對時代背景.
因為我們要看很多背景的書.
才能夠了解當時.
他們寫這本書的時候的背景.
所以其實你要將很多東西.
來平衡地比較.
這個很重要的.
因為你要講解明白給那些背逆的人聽.
才能夠轉向二人的智慧.
是要令到他內心有一個轉向.
有一個他自己裡面的transformation.
或者他自己裡面有一個reform.
有一個改革.
這個就是轉向二人的智慧.
接著說又為主預備合用的百姓.
這個就是施洗約翰.
就是預備人心.
神可以用的百姓.
施洗約翰都帶著很多門生.
很多跟隨他的門徒.
有施洗約翰的門徒.
一直到《士徒行傳》十九章.

$^{721}$保羅去到以弗所一帶傳道.
就見到有十幾個人.
看他又像舊人.
講的東西似曾相識.
又是講得有文有路.
但是問他對於耶穌的洗禮.
從來沒有聽過.
也都不知道是什麼.
我們好像以前都揭過下.
原來是施洗約翰的門徒.
只知道約翰的洗禮.
不知道耶穌的洗禮.
很快和揭一揭.
和這節聖經有關連的東西.
給大家多一點概念.
有時你揭得多揭得多.
慢慢就開始在腦袋裡面.
都有一些概念.
幫大家讀經可以有更加豐富的立體.
十九章.
《士徒行傳》的十九章.
因為這個是施洗約翰預備的百姓.
他的門徒.
不是他們一知半解.
不是他們信得一般般.
而是他們根本從來沒有聽過耶穌的洗禮.
即是說他們一直都在等待彌賽亞.
原來他們一直都還未得救.
但是不關他們的事.
他們一直以為只有一個洗禮.
什麼原因呢?.
失散了.
因為約翰死了之後.
可能他們就怕就要逃難.
走了去亞細米亞的地方.
去了以弗所的一帶.
但是接著就沒有再有追求.
沒有再有傳道人來.
沒有再講關於進一步的救恩的全面的真理.
所以他們就變成孤兒仔.

$^{761}$十九章的《士徒行傳》.
阿波羅在哥林多的時候.
保羅經過了上邊一帶地方.
我們去過哥林多.
就來到以弗所.
保羅的以弗所就是我們去年十一月去的地方.
就是一個很大的劇院的地方.
在那裡遇見幾個門徒.
記住.
盧家醫生的記載.
首先肯定他們的行徑是像門徒.
但是保羅對於人的靈魂很敏銳.
感覺有點不妥.
我給你描述.
就問他們.
你們信的時候受了聖靈沒有.
他們就說沒有.
沒有聽見有聖靈賜下.
保羅就說.
這樣你們受的是什麼洗.
問的說話第二句就已經到肉了.
就已經中point了.
就是你們受什麼洗.
沒有聽過.
已經傳到很普遍了.
已經是沒有可能你不知道.
你洗禮就是歸入耶穌的名下.
就會得到聖靈.
那就是意味著當時.
以弗所書未寫的.
保羅還是自由的傳道.
以弗所書是他第一次坐牢的時候寫的.
你們既聽見真理的道.
就是叫你們得救的福音.
既然信了他.
這是一章的21-22節.
就受了所應許的聖靈為印記.
這聖靈是我們得到將來天上基業的憑據.
憑據就是那個訂金.
那個首期.

$^{801}$直等到神的子民被贖.
被贖就是我們被提.
叫他的恩典得著稱讚.
我們是神的訂金.
對不起.
我們是神買的產業.
因為神現在下了訂金.
下什麼訂.
大訂還是小訂.
不是.
聖靈.
神給我們每一個信的人聖靈.
那裡就很清楚.
保羅很清楚來解釋那個真理.
但是在寫行傳的時候.
保羅路過.
他未寫這本二佛所書.
但是就發現了.
所以可能這十二個人.
觸發保羅一定要寫二佛所書.
解釋信耶穌可以得到什麼.
得到天上各樣屬靈的福氣.
其中就是聖靈.
作為我們每一個信徒的訂金.
原來聖靈是訂金.
神下了訂.
訂了你.
將來你在天上得到基業.
你不信我.
我現在下訂.
下訂就是聖靈.
這樣的意思.
這裡告訴我們.
我們繼續讀下去.
保羅就是很敏銳.
所以原來約翰所預備合容的百姓.
約翰都有他自己的時代的限制.
因為他最多是叫人去跟隨耶穌.
他自己不是來塞亞.
所以如果你聽約翰.

$^{841}$所以這個就衍生了一個問題.
你跟一個人.
跟到一個地步.
讓他影響你的心靈.
那個靈魂.
那個想法.
他的限制就成為你的限制.
所以神很公平的.
始終有時會將一些領袖帶走.
讓你跟隨另外的人.
讓你能夠全面.
這個很重要的.
所以施主約翰他很好.
但是他現在說的.
令到這12個人是不得救的.
是未信耶穌.
沒有聖靈.
沒有聖靈就是不得救.
在保羅的裡面.
就是說你們是門徒.
是沒有用的.
因為你根本沒有聖靈.
你說自己是門徒都沒有用.
你繼續好像門徒一樣生活都沒有用.
因為你根本沒有聖靈.
你根本未接聖靈.
怎樣接呢.
不知道.
保羅現在.
我告訴你怎樣接聖靈.
保羅就第三次這樣說.
他首先問了你們受的是什麼洗呢.
他們說是約翰的洗.
就是施洗約翰的洗.
第四節.
保羅說.
約翰所行的是悔改的洗.
是準備人心.
來跟隨那一位將要來的基督.
真正的彌賽亞.

$^{881}$所以甚至一半都沒有.
都不要說一半.
不可以這樣分.
根本沒有.
這樣就很大件事了.
預備好心靈.
但是最終約翰死了.
那幫人就走了.
然後就瓦解了.
然後就不知道怎麼辦.
後來當然神就會憐憫.
這些人既然已經預備好自己的心靈.
要悔改要跟隨.
不過是不知道為什麼要信什麼.
所以神都差遣保羅.
似乎都不是一個偶遇.
真的讓這幫門徒.
可能只是其中一個代表.
可能還有一些分散了.
這幫人如果信了.
他們就說一個雙重洗禮的福音.
即是說既有約翰的祝福.
再加上有耶穌的救恩.
這12個人就是為主所預備合用的百姓.
就是約翰當年所做的.
到這裡.
反正我們都說了悔改洗禮.
保羅就說.
告訴百姓.
當信那在他.
他就是約翰.
以後要來的就是耶穌.
所以我們很開心.
這個報道會很好聽.
很好笑.
很感動.
很好什麼什麼.
你們這些弟兄姊妹很有善.
很有愛心.
全部是幫他準備悔改的.

$^{921}$跟著他走出去之後.
就未信耶穌.
你先不要興奮.
其實他未信耶穌.
可能他一回到家.
沒有了剛才的東西.
很感動.
流淚.
都沒有了.
他一定要.
首先就是保羅這裡所說.
告訴百姓.
當信那在他.
以後要來的就是耶穌.
所以最終我們是要踢他進耶穌的救恩裡面.
如果不是.
多好都沒有用.
一百樣好都沒有用.
踢完都沒有用.
如果他的心不信耶穌就不行了.
所以一定要他自己.
陪談的弟兄姊妹就是.
你現在跟隨我們.
和你決志.
我們說一句你說一句.
其實就是他的心歸向耶穌.
接受耶穌.
這個就是很重要.
第五節.
他們聽見者說.
就奉主耶穌的名受贈.
你可以想像.
現在才信耶穌.
搞了很久.
不知道多少年.
因為走了老去異不所了.
散居了.
幸好遇到保羅.
這個幸運當然有神的恩典.
有神的撮合.

$^{961}$因為如果他們繼續傳這個就變成異端了.
變成一個和音了.
什麼都不用想了.
你跟我的老師約翰吧.
這樣就原來根本不得救.
這樣就害死人了.
歷史歷代總有一些教派.
一些教主.
一定要人跟隨自己.
譬如文仙明.
統一教.
很多這些異端.
譬如摩門教的約瑟史密夫.
總是有人很想自己做救主.
其實他是一個教主.
你是教主.
你令人滅亡.
但你將人帶去救主那裡.
他就可以得救了.
所以一定要搞清楚.
所以他們奉耶穌的名受浸.
第六節.
很重要了.
因為靈因派會將這一節濫用.
他不要看上下文.
保羅按手在他們頭上.
按手.
所以我們有時間有機會.
可以說說按手這個在新約出現.
一般就是認同.
猜險.
所以保羅的按手.
是有一個對他們得救的認同.
所以不是保羅按手就有神蹟.
靈因派就是追求這件事.
按手就會有聖靈的激導.
聖靈的能力.
透過神醫來顯出.
你的長短腳立刻改變.
他們這樣聲稱.

$^{1001}$但其實很多時候都不行.
很多人排著隊去.
我什麼病你幫我按手.
這裡就不是這個意思.
這裡的按手就是正正告訴我們.
這十二個人其實之前.
未決志信耶穌.
洗禮是未得救的.
現在保羅按手.
留意一下.
聖靈降在他們身上.
接著他們就講謊言.
更加說預言.
所以靈因派一節聖經.
十九章六節.
什麼都有.
好像占陽三寶一樣.
又有一按手.
接著就有聖靈.
接著就講謊言.
接著就講預言.
所以他們要的就是有.
保羅的先知使徒的名分.
所以第四波的靈因運動.
就是恢復使徒的身份.
所以他們可能覺得.
有些人就是使徒.
新使徒運動.
這是第四波的靈因運動.
有機會講.
太多的要有機會講.
所以根本不夠時間講.
這裡告訴我們.
其實不是.
其實由頭到尾上下文很簡單.
他們十二個人是未信耶穌.
很虔誠.
跟隨施洗約翰.
他們一直的心都是悔改.
但是未接到耶穌.

$^{1041}$即是未得救.
所以保羅就認同他們現在肯悔改.
所謂悔改就是.
他們現在打開他們的心.
接受耶穌.
保羅幫他們按手.
接著他們就得到.
這個聖靈的降臨.
所以這個聖靈降臨.
是真實的.
真的有聖靈降臨.
不過這個是特殊案例.
不是一個普遍的案例.
是特殊的原因.
因為《使徒行傳》.
是記載了三次特殊個案.
哪三次呢.
第二章五旬節降臨.
有聖靈的普遍降臨.
以後就不需要再這樣了.
也可以這樣說.
這一次的特殊案例.
就是代表著.
有些很真誠悔改的人.
但是他們不知道.
原來耶穌才是那個救主.
所以如果和約翰.
這個背景的人.
就是當時的時代.
分散了在亞馬遜的這些人.
他們如果來聽明白了.
還有這幫門徒.
神所預備的百姓.
施洗約翰所預備的百姓.
他們現在就負起了責任.
他們就要將這個真正的.
圓滿的得救的福音.
告訴其他約翰散居的門徒.
他們就有這個責任了.
但是那些人就不需要受聖靈了.

$^{1081}$因為一信耶穌.
任何人都會立刻得到聖靈.
就不需要再按手.
聖靈就要立刻降臨.
就不需要這樣了.
所以這個特殊案例.
就不可以成為普遍的個案.
所以我們今天.
為什麼我們牧師幫我按手.
聖靈不降臨.
我又不講方言.
那就恭喜你了.
因為你不是靈因派.
靈因派就是會追求這些極端.
因為這個特殊個案.
不需要也不應該再翻做.
不需要這樣去複製.
這樣做不需要了.
因為只需要我們明白了.
聖靈是會在信耶穌的人的心裡.
就會居住了.
這個是很重要的確立.
所以如果我們追求那個現象.
找一個勁人按手.
第三波八十年代.
就是靈因派的John Wimber.
就來紅館開報道會.
神醫報道會.
這樣去讓信徒上台按手.
有聖靈的激導等等.
現在當然不流行.
其實香港沒有人buy.
那就不會有了.
但是這裡告訴我們.
這些是特殊案例.
不要再重複了.
不需要重複了.
任何人信耶穌.
都可以得到聖靈在他的心裡.
所以這個按手講方言.

$^{1121}$只不過是一個印證.
是一個肯定.
是一個認同.
這十二個人是微信到信.
而他的微信是因為思洗若寒.
其實嚴格就不關思洗若寒的事.
所以某程度思洗若寒.
當年所做的為主預備合用的百姓.
就是這十二個人.
他們現在成為傳道者.
因為他們是經歷過.
原來未有聖靈.
都有聖靈.
他們就要告訴這些人.
一直都有宗教虔誠的人.
一直都沒有聖靈.
一直都以為自己有宗教虔誠.
等於有聖靈.
那就不行了.
其實很簡單.
他們已經有那個虔誠.
但是未有那個決志歸向耶穌.
所以從小教會學校長大的人.
他有那個概念.
但是從來都未決志信耶穌.
就未有那個救恩.
亦都不會有救恩的經歷.
所以這裡盧家呼音就告訴我們.
其實講了這幾節.
是很重要的.
將來我們看馬拉基斯第四章的預言.
很重要.
因為施洗約翰就是這個預言的應驗者.
馬拉基斯所預言的.
大而可畏的日子.
未到以前.
必定有這個以利亞的降臨.
即是說原來耶穌的出現.
就是大而可畏的日子的開始.
末世的開始就是耶穌的出現.

$^{1161}$但是什麼時候是真正的終極.
當然我們講過撒迦利亞書是沒有人知道.
但是我們其實已經是末世了.
我們要記得這樣的概念.
所以我們就應該要自己都有這樣的追求.
所以我們回到第一章這裡.
是盧家醫生引述馬拉基斯所講的.
所以這裡告訴我們.
有為主預備合容的百姓.
接著第十八節.
撒迦利亞對天使說.
我憑什麼可以知道呢.
我都這麼老了.
即是說他原來有祈禱.
希望晚年都能夠有馬利亞.
對不起 以利沙白可以生孩子.
可能生女兒.
生孩子.
這麼老都祈禱.
繼續希望有孩子生.
可能有兩個原因.
因為猶太人他們會算數的.
彌賽亞可能差不多來了.
同女懷孕生子.
所以當然那個就是馬利亞.
即是一個大姐仔.
猶太的女孩.
他們會祈禱自己是不是那個同女.
在耶穌差不多來的時候.
當然馬利亞.
對不起 經常混淆他們兩個.
以利沙白當然不會這樣祈禱.
因為她已經嫁了.
但是因為一個問題.
猶太人認為你沒有孩子生.
就是羞恥.
甚至嚴重來說是咒詛.
記住這是猶太人的傳統.
不是我們基督徒.
所以你不要說人家沒有孩子生.

$^{1201}$就咒詛.
這樣就不對.
每個人都有不同.
不過也有一個很現實.
就是如果基督徒.
譬如歐美越來越富庶.
越來越追求自由.
那個小朋友一出生就十八歲.
那就不用理他.
不是的 十八年你什麼都不做.
不可以做.
又要背著他.
一秋二清.
大家經歷過那些三年零八個月.
隨時要換尿片.
歐洲人就不喜歡這麼麻煩.
生得越來越少.
基督教的世界就越來越少人.
回教徒就不是了.
鼓勵你大量生.
慢慢長大之後.
人家就洗牌.
人家就將你整個種族自然地除掉了.
將整個宗教都沒有了.
所以鼓勵頂尖的生兒女.
某程度屬靈.
也不是不屬靈.
因為很重要的就是.
你要為主預備合容的百姓.
所以他說我都這麼老了.
我的妻子都年紀老了.
我可以回答了.
我是站在神面前的嘉伯列.
原來是天使長.
原來天使長向撒加尼亞顯現.
你可以想像得到.
這件事真的很非比尋常.
因為最高級的那個來.
來向你顯現.
告訴你有兒子生.

$^{1241}$而且那個兒子有以利亞的心智能力.
就是《馬拉基書》四章六節所預言.
你想一下.
為什麼這麼嚴重的這件事.
所以如果我們不看舊約的先知書.
有時候我們錯過了很多重要的經文.
猶太人尤其是祭司.
他自己肯定是很熟經文.
舊約的經文.
他們根本沒有新約.
我們才有新約.
猶太人哪有新約.
耶穌的年代哪有新約.
都還沒寫.
但是撒加尼亞他們是很熟舊約.
天使長加伯列告訴你.
你的兒子.
快要出生的兒子.
其實現在都還沒懷孕.
他就是以利亞.
就是聖經預言的以利亞.
就是四章六節的以利亞.
你想一下.
真的嚇死撒加尼亞.
你可以想像得到.
所以我們要回來再讀第五第六節.
看!耶和華大而可畏之日未到以前.
我必猜顯先知以利亞.
其實就是施洗約翰.
耶穌已經定性了.
就是施洗約翰.
到你們那裡去.
他必定使父親的心轉向兒女.
盧家就是引用這句.
兒女的心要轉向父親.
免得我來就坐遍地.
就是說要有一個很大的reform.
一個心靈的改變.
施洗約翰就是做這件事.
他的一生很短暫.

$^{1281}$因為好像是談花一言.
一出現很多人就去洗禮.
接受他.
很快他就沒有市場了.
因為耶穌來了.
人人去耶穌那裡就不去他那裡.
他繼續傳一個叫人悔改的transformation.
得罪了希律.
殺了他的頭.
所以他很快就出現在這個世界.
也很快三十多歲就已經離開世界.
就斬頭升天了.
這個就是施洗約翰一生的貢獻.
就是叫人的心回轉.
其實一個人一生就有這個message.
很重要.
叫人的心回轉.
他一生就做到這件事.
其實很偉大.
雖然人也不多.
因為慢慢的人就跟耶穌不跟他.
不如我們又揭一揭約翰福音.
勉勵一下事奉神的弟兄姊妹.
如果你是已經在事奉.
你不要有錯誤的心態.
我們一定要.
其實事奉的第一天.
你就要有這個準備.
你是應該要被人忘記的.
第三章三十一到三十六節.
約翰福音的三章三十一到三十六節.
再前一點.
不然那個背景大家不熟悉.
第二十二節這一段.
今天我們先說一些背景.
所以先給大家鬆一鬆肚子.
可能下個星期我們就會進入經文.
那個對話.
神透過先知馬拉基來對話.
你們做得不對.

$^{1321}$我們怎樣做得不對.
這樣對話.
其實對質.
我們看這裡二十二節.
這裡都很寶貴的.
每一個事奉的人都要有這個態度.
這是以後耶穌和門徒到猶太地.
在那裡居住施洗.
約翰在靠近撒冷的哀戀也施洗.
即是有兩檔.
因為那裡水多.
眾人都去受洗.
接著第二十四節.
那時約翰還沒有下在監裡.
特別強調他還沒坐監.
即是還沒得罪希律.
約翰的門徒和一個猶太人辯論潔淨的禮.
即是說有些猶太人.
他們很有興趣約翰的一些誡條.
一些論述.
對話.
來到去有些好像宗教對談.
可能問一些很尖銳的問題.
那個人不懂回答.
就來問約翰.
即是那個門徒.
就來見約翰.
說拉比從前和你在約旦河外.
你所見證的那一位.
即是耶穌.
現在施洗.
眾人都往他那裡去了.
可能那個猶太人就不懂抬舉.
你說你是正派.
為什麼那裡又有另一個正派.
為什麼你又說自己正派.
他又說自己正派.
什麼邀諷.
還是什麼.
其實是兩黨.

$^{1361}$是兩兄弟還是什麼.
然後是恩諷.
即是你可以想像得到.
你約翰又說自己正宗.
耶穌現在又是你的門徒.
你幫他洗禮.
他現在又說自己正宗.
那誰才正宗.
那個門徒真的不懂回答.
就來問問約翰.
其實很尖銳.
你所見證的那一位.
如果我們真是在現場聽到他的對話.
這個是質問.
還是替自己的老師不值.
你所見證的那一位.
現在施洗.
這句後面真是帶著酸酸的.
眾人都往他那裡去了.
其實我想約翰都知道.
如果不是從天上賜.
人就不能得到什麼.
我曾經說.
我不是基督.
是奉差遣在他前面.
你們跟我的門徒.
你們自己可以給我作見證.
娶新婦的是新郎.
這些是當時的比喻.
大家一聽就明白的比喻.
新郎的朋友站著.
聽見新郎的聲音.
我朕喜樂.
故此我者喜樂滿足了.
他必興旺.
我必衰微.
整段經文是說那個洗禮的那個.
那個現在洗禮的那個.
人們都去他那裡的那個.
那個就是耶穌.

$^{1401}$肯定對號入座.
那個門徒就是針對耶穌.
你所見證的那個.
現在施洗.
人們都去他那裡了.
約翰就告訴我們.
首先他明白.
他的事奉是天上賜的.
如果天上要收回就收回了.
若不是從天上賜.
人就不能夠得到什麼.
大家事奉記得這個很重要的原則.
我們是因為神給我們一個事奉的機會.
去服侍的.
所以我們要感恩.
不要慢慢就變成了喧賓奪主.
你拿著那個事奉.
誰都不可以拿走.
那不應該的.
若不是從天上賜.
就沒有人能夠得到什麼.
你有這個心態.
其實這裡有幾點的.
約翰的內心的事奉態度.
今天我想很多事奉的人.
連這段經文都沒讀過都不奇怪.
很多聖經都沒讀過.
就去事奉其實有些危險的.
因為不知不覺.
你會拿著這個事奉以為是你的.
那是不健康的.
忠心去做好你的本份是應份的.
忠心做完之後.
你要告訴自己.
若不是從天上賜.
我是不能夠得到什麼.
這個不是謙虛的說法.
是真實的說法.
而事奉一來的牧者.
越會明白這個真理.

$^{1441}$接著他不單只是空泛地說.
天上賜的我現在事奉.
接著他跟他的門徒說.
是針對自己.
我曾經說過.
即是我一直這樣信奉這件事實.
我不是那個主角.
我不是基督.
即是那個是基督.
我只不過是為他準備的那個人.
所以約翰是一個很特殊的人.
因為他真是先鋒.
又是先鋒又是以利亞.
又是預備百姓的心.
所以他有好幾重的存在價值.
我們肯定沒有約翰.
那些特殊的歷史價值.
因為耶穌未來之前的約翰.
是有特殊的工作的.
是他出現耶穌才能來到的.
但我們就不是了.
我們有約翰事奉的崗位.
因此我們就不是要回頭.
我們帶耶穌來這個世界.
就不可以了.
所以約翰是很特別的.
所以他現在說第二個重要的.
內心的價值觀.
事奉的價值觀.
我曾經說.
我一直認定自己.
恰如其分的身份.
我不是基督.
你們可以為我作見證.
即是說他的門徒一直.
即是說這個門徒開始有少少歪了.
那個心態歪了.
開始為自己的老師不值了.
開始有些覺得.
老師好像為什麼要這樣呢.

$^{1481}$約翰說你不要搞錯.
我一直認定自己是.
那個準備基督來的人.
你們自己可以給我作見證.
是奉差遣在他前面.
準備的先鋒.
接著說你們自己可以作見證.
所以盼望我們的事奉.
是要讓人看到.
我一直都是這樣.
去看到自己的身份是僕人.
這個是很重要的.
慢慢就是你是主人.
調一調你的崗位.
你就化爛了.
這樣不對.
你要讓人看到.
你們自己可以為我作見證.
約翰這樣理解自己.
你想想那個價值觀多麼純正.
那個事奉態度是多麼好.
有我就開心事奉.
人們都去他那裡.
就對了.
我的工作就是要叫人們都去他那裡.
原來事奉的人是這樣想.
原來約翰是這樣思考.
我的工作就是叫人們快點去他那裡.
這樣的見證.
所以那個門徒聽到老師這樣說.
終於又明白了.
所以原來都要幫那些門徒.
來開解他的內心.
因為始終門徒未到那個境界.
肅靈境界.
一定有一種情意結.
對老師那種擔心.
但是現在老師開解他之後.
很寶貴的一個肅靈功課.
老師自己.

$^{1521}$你們自己可以為我作見證.
不是我為自己作見證.
是你們可以為我作見證.
我一直都說我不是基督.
第三個就是.
接著是一個比喻.
娶新婦是一件很開心的事情.
我們留意的.
一定是新郎的婚姻.
那個婚宴.
等待新郎的出現.
至於當時的背景是怎樣.
怎樣娶.
是不是和我們今天香港一樣.
這些不是最重要.
而是那個主角的出現.
其他的人就應該開心了.
你可以想像得到.
新郎的朋友站著.
聽見新郎的聲音.
就岑喜樂.
就是說原來新郎一直都未出現.
大家一直等待.
大家都好像以為自己是那個主角.
但是那個開心.
不要緊.
但是如果新郎來了.
而你不開心.
那就是問題了.
你就是問題了.
新郎就是主角.
整個婚宴就是他是那個主角.
所以基督得到榮耀.
主耶穌被高舉.
那你還想怎樣.
你還想有什麼不開心.
不舒服.
所以阿約翰就和那個門徒說.
三點關於事奉的態度.
多麼重要.

$^{1561}$任何一個事奉的人都要讀一讀這段經文.
新郎的朋友站著.
聽見新郎的聲音就岑喜樂.
故此,therefore.
我的喜樂.
我者喜樂.
那些人走去他那裡.
我滿足了.
喜樂滿足.
一個事奉神的人.
越做越沒有人.
他的喜樂滿足.
這個心境要很成熟的牧者才能夠經歷到.
否則真的很難.
是盡了力.
盡了力就沒有人.
這樣的心靈滿足.
故此.
那些人去到基督那裡.
又認出他是基督.
而認出我不是基督.
那我的喜樂就滿足了.
所以我們一定要告訴別人我們不是基督.
所以這個世界有很多教主.
真是很大的試探.
人人都想自己做教主.
每個人都要跟隨.
你說什麼就一錘定音.
誰敢說的話就這樣.
就好像有票房保證.
每個人都不敢出聲一定要聽.
還要是洗腦的聽.
這樣就很大罪.
我們寧願做一個時時都喜樂.
在基督旁邊喜樂的人.
見到人信主你又會喜樂.
見到人離開神你會很憂愁.
這樣的心境.
所以他必定會興旺.
我必定會衰微.

$^{1601}$接著三十一節.
本來就不想再讀下去了.
因為怕太長了.
但其實有關連的.
為什麼他能夠有這個心境呢.
你就可以看到.
到底是作者約翰說.
還是施洗約翰繼續這樣說呢.
似乎是作者約翰說.
他說其實怎樣比較.
一位是天上屬天而來的.
你和我包括作者自己本人.
我們是屬地的.
我們是阿當的後人是罪人來的.
怎樣跟耶穌比.
他是萬有之上.
從地上來的.
我們這些事奉的.
是屬乎地的.
他說的都是屬乎地.
我們有我們的限制.
說的是恩賜.
你叫我唱歌我也走音.
每個人的恩賜都有不同.
盡力唱錯.
很盡力了不過也走音.
這個就是沒得比較了.
但是你可以想像得到.
更嚴重的就是.
我們只不過是屬乎地.
所以作者約翰就在這裡加個注釋.
為什麼施洗約翰會這麼偉大呢.
就是因為他看到自己是屬乎地.
不會在這裡裝模作樣.
說我就是那個彌賽亞.
起初的時候.
每個人都以為他是彌賽亞.
但是他知道自己不是.
下面那一節.
從天上來的.

$^{1641}$是在萬有之上.
說了第二次了.
這樣跟著再說.
天上來的更加重要.
天上來的就是他帶天上的訊息給我們.
人間.
天上來到凡間.
他將所見所聞見證出來.
只是沒有人領受他的見證.
這個就是約翰福音第一章.
他到自己的地方來.
自己的人.
就是猶太人.
道不接待他.
接著一章十二節.
就是我們信耶穌的.
凡接待他的.
就是哪些人.
就是信他明的人.
他就賜他們權柄.
所以我們是有權柄的.
作臣的兒女.
這些人不是從血肉生的.
不是從人而生的.
是從神而生的.
就是聖靈而生的.
所以你想想多麼寶貴.
所以這裡又再一次回應.
說主耶穌.
天上來的.
將所見所聞.
就是神賜給人類的永生.
來見證出來.
用他自己的生命.
用他自己的救贖.
來見證出來.
指示.
作者真是可圈可點的.
來去很慨歎.
就是說他寫約翰福音的時候.

$^{1681}$真是很多猶太人不得救.
他自己是猶太人.
他自己是由漁夫成為傳道人.
跟隨耶穌.
被耶穌欽點成為他的使徒.
一路跟隨.
他現在寫約翰福音年紀很大了.
指示沒有人來領受他的見證.
所以你記住你在教會內.
你看到有些人.
真是信不到耶穌.
你就很慨歎.
很悲哀.
為他信不到耶穌流淚.
希望他能夠接受耶穌的見證.
他看不到天上.
這裡我們一直說.
約翰福音說.
有一個平衡世界.
天堂和地上.
約翰經常這樣告訴我們.
他將天上和地上.
時時做一個平衡的比較.
信不到耶穌的人.
不是他不夠我們厲害.
是他看不到天上.
因為他未打開他的心靈.
是沒有天上.
沒有天上的每個星期來崇拜.
甚至有奉獻的.
是沒有得救的.
他沒有天上的見證就沒有了.
這是很絕對的.
33節再一次再說.
回應一章十二節.
那領受他見證的就人常人.
性靈.
人常人在他的心靈裡面.
證明神是真的.
即是他開啟了他屬天的境界.

$^{1721}$我信耶穌.
1987年12月24日.
我接著去到一月.
我突然之間.
是一月還是12月底.
突然之間我回家.
整個人輕了.
很輕散.
是你沒想像.
你開啟了你對天堂的認識.
不是你見到天堂.
是你心靈裡面有天堂.
有靈界.
不是見到靈界.
是你裡面有靈界.
你就可以想像得到.
這些經文.
現在你幾十年後再讀.
原來很真實.
原來一個信耶穌的人.
打開了天堂給你.
在你的心靈裡面.
你可以想像到.
這些經文有多重要.
約翰苦口婆心地說.
用一些例子來說.
彌哥底母.
這樣來說.
接著最終他想說的真理.
主耶穌打開了天上的門.
將天上告訴我們.
他的所見所聞.
我們很好奇.
一些人從哪裡回來.
告訴我們是怎樣.
特別那時幾十年前.
沒有飛機那麼先進的時候.
你去過坐飛機就很厲害.
是怎樣的.
美國是怎樣的.

$^{1761}$如果你去過非洲.
非洲是怎樣的.
我們就很好奇.
但是在天堂來的人.
我們不好奇.
這個不奇怪.
因為我們是罪人.
約翰說了兩次.
一章十二節說了.
三章又再說了.
所以原來信耶穌的人.
就是領受他的見證的人.
即是認同和接受主耶穌.
你是我的救主.
和你是由天差遣而來的神.
這裡說.
神所差來的就說神的話.
即是將神的真理告訴我們.
所以我們有真理.
因為神賜聖靈給他.
是沒有限量的.
父愛子.
已將萬有交在他手裡.
所以我問為何這麼決絕.
信的人.
你就打開了整個永生之門.
不信的人就得不到永生.
而且神的震怒常在他身上.
神的震怒不是令他天天震.
是令他只有這個世界地上的見識.
很恐怖.
你整個人生就是地上的見識.
你沒有天堂的見識.
即是連一個婆仔坐在這裡.
黑眼訓崇拜的.
他都有天堂.
你一個很厲害的.
是拒絕神拒絕耶穌的.
你沒有天堂.
你只有地上.

$^{1801}$神的震怒.
即是說注定你是不得救的震怒.
這就真是很嚴重的一件事.
我們又越講越遠.
其實我們想講的就是.
《馬拉基書》第四章.
關於盧迦講的這個以利亞.
其實就是施洗約翰.
我們也順便講一下.
施洗約翰的事奉心境.
事奉態度.
講一下又不是拆開的.
不如順便講一下.
約翰很慨嘆在他的那個時代.
或者相類的.
原來是有兩個平衡世界.
在天堂的人是有永生.
只有地上見識的人是沒有永生.
他又想講這件事.
他要慨嘆的就是.
其實很簡單.
接待就可以了.
都是不肯.
因為不肯接待他.
就是藐視他的身份.
他的講論.
所以多大的神蹟呢.
都改變不了這些人的心靈.
所以用我們的心靈遊園.
其實就是一個救恩.
很重要的開始.
我們信主內.
如果我們的心靈遊園.
讓神的道在我們的心裡面.
又會帶來另一種階段的改變.
所以我們還有幾小時.
我們先看看經文.
放心.
去到四章六節.
我們會再詳細講上下文.

$^{1841}$我們來看第一章.
耶和華籍馬拉基全級以色列的默示.
默示可以是那個.
sign的字.
也可以是burden的字.
所以很多時候默示.
是一個重擔.
有時是那個.
我們講過耶利米蘇.
以西結書.
先知食書卷.
食了就是口中甜如蜜.
肚裡面發苦.
當然是一個異象的食.
代表著你作為那個代言人.
就會有可能是講安慰.
你就要講安慰的話.
令到人得到安慰.
但有時是要講審判的話.
令到那些人.
不要感覺.
你講什麼.
我都不覺得有什麼問題.
那就不是審判了.
就變得不夠力度了.
好像隔鄉騷養.
所以那種食小書卷.
真的暗示著一個先知.
你要扮演你自己.
刻如其分的角色.
應該要審判的話.
你就要講.
不要怕人的情面.
你是神的代言.
所以講真話的先知.
是不容易的.
第二節.
耶和華說.
我曾經愛你們.
你們卻說.

$^{1881}$你在何時上愛我們.
很有趣的對話.
馬拉基就用.
代表著神的.
好像第三身一樣.
來講論點.
你們得到神的愛.
你們以為神不愛你.
耶和華說.
以數不是雅各的哥哥嗎.
神要舉的例子.
居然是兩兄弟.
為什麼要舉兩兄弟呢.
因為這個涉及到.
簡選論.
耶和華說.
以數不是雅各的哥哥嗎.
我卻愛雅各.
「戶」就是動詞.
戶以數.
使他的山嶺荒涼.
把他的地葉交給曠野的野狗.
以東人說.
我們現在雖被毀壞.
卻要重建荒廢之處.
以東就是皮特拉古城那裡.
Petra那裡.
萬君子耶和華如此說.
任他們建造.
我必拆毀.
人必稱他們的地為罪惡之境.
稱他們的民為耶和華永遠老奴之民.
即是說.
在他們那個時代.
一定知道以東怎樣建造.
怎樣想復興.
都不行的.
背後未必是地土荒涼.
未必是貧瘠.
未必是他們沒有智慧人.

$^{1921}$原來最大的問題.
是神一直咒詛那個地方的人.
是神定義我要咒詛.
這個耳塑的後人.
這樣的一個邏輯.
定義要咒詛.
咒詛到你們真的變成罪惡之境.
變成神永遠老奴之民.
那為什麼?.
雅各.
為什麼你會愛雅各.
你會恨惡耳塑?.
這個大問題.
本來下一次要詳細講.
還有多少時間?.
還有十分鐘.
我們要揭幾段經文.
跟大家先來看看.
第五節.
你們必親眼看見.
也必說.
願耶和華在以色列境界之外.
被尊為大.
因為你們不肯認為我愛你們.
愛包括保護.
包括使你們有末後的興盛和祝福.
但是在這個時候的以色列人.
心如止水.
不認為自己的神和他們有愛的關係.
去到淡如水的敬拜生活.
沒有了愛的敬拜生活.
自然就未同酌納.
行禮如儀.
暗示有幾樣東西牽連.
所以第二段對話對質.
集中在聖殿裡.
參與事奉的祭司.
原來祭司有令百姓心如止水的原因.
所以你們會看到.
首先這裡就很奇怪.

$^{1961}$一說到我曾經愛你們這個民族以色列.
他們就立刻說.
是呀,三百年前是,現在不是.
亞伯拉罕的時代就是,現在不是.
你什麼時候愛我們.
他們現在已經落在對神可有可無的關係光景.
我們就要問這個問題.
為什麼神要定義對耳掃的懷恨這麼深呢?.
永遠的牢牢.
即是說一直的奏坐是不會離開.
這樣都很慘.
我們先看幾段.
希伯來書和新約有三段.
很有代表性的三段.
希伯來書兩段,羅馬書一段.
來講到為什麼神會恨戶耳掃.
恨他也好,還要恨他的後人.
這樣凸顯什麼呢?.
凸顯其實以色列一直沒有放棄你們.
一直就愛你們.
這個邏輯是這樣.
這個和我們今天在香港富庶的社會有關連.
我們從來都沒有經驗過要回教會.
突然間人家用機關槍掃你.
用水潑你.
吐口水在你的頭上.
你一散會,我們就沒有.
沒有這些危機.
這樣我們的信仰就慢慢變成.
有時會熱心.
有時會冰冷.
現在作者告訴我們.
代表著神.
這個對質就是揭示為什麼你們會冰冷.
冰冷的彰顯情況會是怎樣呢?.
原來就是不覺得神愛我們.
這是一個很大的思考問題.
我先揭了幾段新約經文.
下星期我們就詳細講這個揀選論.
揭了給大家知道.

$^{2001}$這是聖經很重大的課題.
至少保羅是認為揀選的.
剛才那裡某程度是揀選的.
因為那些人不肯接受耶穌的見證.
就真的一路都不行.
只見不到天堂.
只是他們不接受耶穌的見證.
某程度上反過來.
他們信不到耶穌的原因.
是他們硬心.
不是他們沒有聽過福音.
我們先看羅馬書這段.
羅馬書第九章.
這章聖經就講揀選論.
真的用耳塑,瓦角做一個例子.
第九章.
我在基督裡面說真話.
並不謊言.
可能是動詞.
有我的良心被聖靈感動.
給我作見證.
我是大有憂愁.
心裡面時常傷痛.
為我弟兄我骨肉之親.
就是自己被咒坐與基督分離.
我也是原意.
他們是以色列人.
那兒子的名分榮耀諸約.
律法禮儀應許.
都是他們的.
即是說保羅自己也是這樣.
有亞伯拉罕的同神的立約.
因為他們是以色列人.
即是第八日行國禮.
有兒子的名分.
榮耀諸約.
榮耀即是神的Glory.
諸約即是摩西的律法.
西乃莎的約.
大衛又有大衛的約.

$^{2041}$即是很多神與以色列人.
訂立了的祝福.
以約做這個單位.
即是說神一定會信守他的諾言.
律法十誡.
禮儀即是說很多不同的節期.
不同的那些.
要讓人感受到神同在.
那些禮儀.
譬如吹國節.
這些禮儀.
呼籲人來想想.
自己要來到神面前.
吹國節.
所以歐洲的教堂.
鄉村教堂.
或者一些社區的教堂.
崇拜會敲鐘.
告訴人現在來崇拜了.
今天是星期日.
要穿得整齊.
平時耕田的農夫.
每個人都要穿西裝.
去教會敬拜神.
所以我去到雄恩.
都要鼓勵弟兄姊妹.
不要那麼吵架.
基層就不等於吵架.
基層一樣可以很有端莊.
清潔去敬拜神.
這個要慢慢灌輸.
讓你感受到我是不同的.
我今天是星期日.
禮拜即是我去朝見神.
這個是要教導的.
所以禮儀是他們的.
應許對亞伯拉罕.
和對眾先知等等的應許.
都是他們的.
列祖就是他們的祖宗.

$^{2081}$即是亞伯拉罕.
以撒雅各那些.
按肉體來說.
按照肉身而言.
Bodily speaking.
基督都是猶太人所生的.
基督是猶太人.
不過保羅是很保險的.
一說到他是肉身.
他馬上再加一句.
他是萬有之上.
永遠可稱頌的神.
阿們.
這樣寫信的.
突然間馬上停一停.
耶穌?.
他是萬有的主宰.
永遠可稱頌的神.
阿們.
保羅這樣的敬虔.
我們要學的.
我們一說到.
主耶穌都比較對.
剛剛我不是說那個女士.
問這些問題.
耶穌還是主耶穌.
還是這樣.
他開始懂得問這些問題.
這個就是保羅他的內心.
一說到基督是地上而來.
馬上就說.
馬上補充.
他是在萬有之上的主宰.
所以我們要有這種.
兩個平行世界的眼光.
一說到他是.
按照肉身來說.
基督都是你們那裡所生的.
然後馬上停一停.
不可以就這樣說這件事.

$^{2121}$立即說.
他是在萬有之上.
所以你看看保羅這些行文的敬虔.
第六節.
這裡開始說簡選.
「這不是說神的話落了空.
因為從以色列生的.
不都是以色列人.
也不因為是亞伯拉罕的後裔.
就都作他的兒女.
唯獨從以撒辛的.
才要稱為你的後裔.
這就是說肉身所生的兒女.
不是神的兒女.
唯獨那應許的兒女.
才算是後裔」.
因為所應許的話是這樣說的.
「到明年這個時候.
我要來撒拉必生一個兒子」.
當然就是以撒.
創世記的十八章.
「不但如此.
還有利伯加.
既從一個人.
就是從我們的祖宗.
以撒懷了孕」.
利伯加就是他的僕人.
在巴旦亞蘭所娶回來.
嫁了給以撒.
以撒六十歲的時候.
生了一對孿生的孿生子.
孿生子就不一起出嫡.
一定是一個出一個後.
一個先一個後.
以素先出.
然後就到雅各.
留意一下這裡.
「既從一個人.
就是從我們的祖宗.
以撒懷了孕」第十一節.

$^{2161}$這個真是揀選的.
更加重要的一個指標性的經文.
雙子都未生下來.
接著保羅還要加一個.
他自己的注釋.
善惡都未作出來.
其實你不用說了.
我們讀過創世記二十五章.
兩個都一啖啖.
兩個都是惡的.
兩個都沒有善.
兩個如果在天平稱一稱.
其實雅各是更加詭詐.
雅各是更加深沉.
以素就發了出來.
但是雅各就更加陰險.
所以善惡未發出來.
保羅的意思是.
他們好像三歲都未定八十.
因為未出世.
但是神已經揀選.
保羅這樣說.
善惡都未作出來.
只因要顯明神揀選人的旨意.
不是在乎人的行為.
乃在乎召人的主.
就是這一句.
揀選論.
所以今天我們能夠一直信耶穌.
就是真的非常大的恩典.
是神揀選我們.
為什麼揀選呢.
沒得解釋了.
這個就是神的事了.
我們回到天堂就知道了.
但是我們真的看到天堂.
因為從萬有之上的主.
已經來到將祂所見證的.
我們相信了.
一信的人就打開天堂的門.

$^{2201}$我們的心靈就有平衡世界了.
所以我們要為那些沒有平衡世界的人.
就是沒有天堂的人.
來祈禱.
因為我們不知道他是否被神揀選.
但是他一信耶穌.
我們就知道他是被神揀選.
所以我們是不知道先後次序的.
我們只能夠知道.
誰信耶穌呢.
他一直信耶穌呢.
他就是被神揀選了.
因為從來不關我們的善惡的.
不是因為你和我特別乖.
特別是回基督教學校.
從來都不是的.
是因為善惡都未作出來.
只因要顯明神的旨意是不更改.
就是這個揀選的旨意.
我們一起禱告了.
多謝主給我們開始了這個背景的講解.
盼望我們一直看這些對質.
神透過瑪拉基和以色列人的那個時代的對質對話.
有時好像是一些很挖苦性.
但是其實是一種寫作的修飾.
一種修辭.
來到很苦口婆心的.
點出以色列人可能正在面對的問題.
可能是一個不太多又逼迫.
因為已經過了波斯帝國的時候的重建.
也都未到羅馬帝國的迫害的時候.
所以可能他們在這一段時間.
真是有一種慢慢對信仰可能只是一種行禮如儀.
很盼望神都給我們從其中吸取了一些屬靈的教導.
我們這樣禱告奉耶穌基督的名.
阿們.
\newpage



\section{箴言 29:18}
\label{sec:SWoA7YUZj_4}
\textbf{2024-01-07 午堂《2024錦宣家．齊躍動》箴言 29\_18 講員:陳麗蟬牧師}
\newline
\newline
連結: \href{https://youtube.com/watch?v=SWoA7YUZj_4}{\texttt{https://youtube.com/watch?v=SWoA7YUZj\_4}} ~~~~ 語音日期: 2024-01-09
\newline
\newline
\hyperref[sec:qkHV9uoLu6Y]{\small{< < < PREV SERMON < < <}}
~
\hyperref[sec:index_chronic]{\small{[返順時目]}}
~
\hyperref[sec:index_scriptual]{\small{[返順卷目]}}
~
\hyperref[sec:Y_5AXHxN_yY]{\small{> > > NEXT SERMON > > >}}
\newline
\newline
箴言 29:18
\newline
\begin{longtable}{cl}
\hline
\hline
章節 & 經文 (和合本修訂版)\\
\hline
29:18 & \begin{tabularx}{0.7\textwidth}{X} 沒有異象，民就放肆；惟遵守律法的，便為有福。 \end{tabularx} \\ \\
[1ex]
\hline
\hline
\end{longtable}
$^{1}$弟兄姊妹平安.
今天是2024年.
第一次的五堂崇拜.
在這裡祝福大家新的一年裡面.
身體心靈都壯健.
我們要齊心的.
我們為主的功力去努力.
其實本來今年的六月我就會退休了.
將會轉為半職.
所以我已經安排了自己很多的.
不是身後事.
不過就是退休之後的事情.
我又知道65歲之後買保險是很難的.
所以那些人就叫我.
你在做事的時候.
65歲之前就要買定了.
結果又去買了.
真是很心痛.
然後又去計劃去旅行.
已經有人說跟我去美國.
還有去新疆.
也有人邀請我去短孫.
我就問定了他.
去短孫的時候可不可以順便旅行.
他說可以可以.
我就答應了.
我又想一下.
我又報了接下來這一年.
報了有四次的宿營.
就是一些安靜脈賞.
又有Day Camp那些.
又交了一大筆錢.
安排好了.
在12月初.
當我放大假去日本旅行之前的四天.
幾位堂董就找我聊天.
我想著半職太多了.
你少一點吧.
三分一職好了.
想著是這樣的.

$^{41}$誰知道他說不是.
因為教會現在的情況.
他們就邀請我.
延遲退休的日期.
所以就叫我考慮.
我就經過了幾天的禱告.
我就感到平安的.
我覺得我現在都精神的.
我現在都走得不錯的.
於是我就接受了這個任命.
在這個今年裡面.
教會將會有很大的變化.
就說曾牧師.
由今天開始.
他就會去到這個紅恩堂.
帶領著紅恩堂的發展.
推動這個福音的工作.
因為我在那裡都難的.
所以都多為祈禱.
也都在嘉恩樓.
在今年開始就轉為這個叫做嘉恩堂.
也就是說他們將會.
就來去獨立成為一間分堂的.
不過現在還沒有獨立之前是一個自立.
鄧全道主力就會過去.
這個嘉恩堂那裡來去事奉.
我就留在這個感受堂.
和大家一起來去努力.
去將主耶穌基督給我們那個大使命.
我們一起來去實行出來.
今天過新年.
所以我只是選了一節的經文.
和大家一起來去勉勵.
我們一起去讀.
《箴言》第29章第18節.
我們一起讀.
1,2,3.
沒有異象民就放肆.
為遵守律法的便為有福.
先看一下異象這個字.

$^{81}$什麼叫異象呢.
就不是說我們見到有什麼奇特的東西.
有些人見到異象.
又見到有個人穿白衣服發光.
不是那些.
也不是我們晚上睡覺的時候.
做夢的那些異象.
更加不是白天坐著.
做白天夢的那些.
更加不是那些.
異象就是指道上帝對人的說話.
上帝對人的啟示.
直接可以說上帝對人的指引.
上帝給人的使命.
上帝給人的目標.
上帝就會藉著聖經.
藉著祈禱.
來表明在我們每一個人的生命裡面.
有不同的指引.
每一個人的生命都是獨特的.
我們從上帝裡面去領受了獨特的異象.
或者領受了獨特的使命的時候.
我們的人生就變得有目標.
有方向.
當人有目標.
有方向的時候.
就帶給我們有動力.
也都有合作.
我們看一看幾幅漫畫.
這是蒲公英那本《小冊子》裡面的.
剛剛今期出的.
姐姐就跟弟弟說.
「考試到了,快點讀書了」.
弟弟就說.
「不用了,我已經學會了很多東西,夠了夠了」.
姐姐說.
「你現在只是讀小學而已」.
「你懂得甚麼,將來可以做甚麼」.
弟弟怎麼回答呢.
「我比幼稚園的人懂得多東西」.

$^{121}$「所以將來我就要教幼稚園」.
我不知道這是不是他的一個意象.
今早我很感動.
因為在廣場那裡.
知不知道朝聖茶室是幾點開檔的.
朝聖茶室是早課完結後才開檔的.
但今天早上突然來了二十多位.
信義中神學院的同學.
他們來參觀一下感受.
現在就過了嘉欣堂那邊.
突然間很熱鬧.
我們在101那裡.
跟他們歡迎完,講解完之後.
富哥就說帶他們遊一下感受.
富哥就說.
外面有個茶水檔,可以光顧一下.
當時已經開了檔.
但還未擺好東西.
也有三位姐妹弟兄在那裡.
還未開檔.
因為早課完才開檔.
那時是九點三四個字左右.
突然間二十多個人.
來到這個愛咖啡那個.
愛薑茶那個,愛花茶那個.
然後那個又愛餅乾.
我看著他們個個快手快腳.
個個五十多歲的姐妹.
跑上樓梯不知道要拿些甚麼.
我看見那幾位姐妹.
後期又再加入一些弟兄.
我看到他們為了招呼一些新來賓.
其實他們都不知道他們是甚麼人.
總之就突然間來了二十多個人.
他們開這個朝聖茶室.
只有一個目的.
就是給每一個來到的人.
對教會一個很溫馨的感覺.
所以他們都不理這些二十多個人.
在哪裡而來.

$^{161}$也都不理.
即是說未夠鐘開檔.
可能外面的茶餐廳就會.
未夠鐘開檔.
十點鐘才來.
他們不會.
雙方快手快腳.
拿咖啡,沖茶,開餅乾.
富哥本來說.
你們沒有麵包的.
我說要不要去麵包檔買.
但你看到其實一班.
在朝聖茶室事奉的弟兄姊妹.
他們有同一個目標.
就是叫每一個人來到的時候.
去感受到教會那份溫馨.
那份對他們的接待.
那份溫暖.
在當中我很感動.
其實我告訴你.
你們的事奉.
也都感動了一位國內來的姐妹.
他們問我們.
奉獻箱在哪裡.
拿了奉獻封.
然後我帶了她進來.
親手帶進奉獻箱裡.
我想當人有那份目標.
去擺定我們的朝聖茶室.
目標就是這樣.
不是叫她十點鐘之後來.
而是她何時來.
只要我們可以的時候.
我們都努力去招呼其他人.
所以相信.
我們在面對這個異象的時候.
我們面對上帝給我們使命的時候.
我們就要做到最好.
在2024年.
今天是第一個主日.

$^{201}$我問一下大家.
有沒有立了今年的目標.
沒有.
包括我之前都沒有的.
覺得很老土.
這些是小學生作文才寫的.
但我發覺不是的.
其實有目標.
有使命是好的.
我是負責常恩師班的.
他們晚上會獻師.
我們在過去兩年左右.
我們一個月有兩次的星期五下午.
我們一起查經.
有時我帶.
也都訓練那些姐妹.
自己來去大廈查經.
在今年裡面.
我想到三,四月.
我突然間有感動就鼓勵他們.
不如我們今年就一起訂讀經計劃.
於是有些人就訂了.
一年要看完整本聖經一次.
有些人只是看完舊約.
有些人就訂了新約就夠了.
有些人就說.
我看字很慢.
我只看四福音就夠了.
到了十一月左右.
我們就來一個總結.
看看將來會怎樣.
真的有人可以看完整本聖經一次.
也都有人可以達成.
他之前訂的目標.
更加厲害的.
有人是超過他的目標.
我是其中一個達不成目標的.
很差勁.
但不過我自己的檢討就是.
有目標比沒有目標好.

$^{241}$當我不訂的時候.
我想我不會立志從創世記開始看下去.
但我訂了的時候.
忽然間看到WhatsApp.
他們說他們看到哪個時候.
我就會努力了.
雖然我只看了一大半.
但總起碼沒有目標比好.
所以他說沒有異象.
民就放肆.
如果當我沒有為自己去訂的時候.
可能我每天都會有靈修.
但我就不會去決心去速讀聖經.
又有一次.
有一位弟兄來跟我聊天.
問到孫家現在有多少位童工.
他說曾牧師將會轉去紅恩堂.
鄧全道就主力在嘉恩堂.
我們這邊就剩下謝牧師.
他是半職的.
楊姑娘 盧姑娘 Odi姑娘 和我.
我一聽完之後.
那位弟兄說.
哦 就這麼多了.
沒有了.
我說哦 就是這麼多了.
接著他說.
不要緊.
我們有很多弟兄姊妹可以動員.
是不是.
他說真的.
每一間教會都不會有很多的童工.
不過最重要的.
就是一班弟兄姊妹.
一班弟兄姊妹.
肯起來一起同心事奉.
這個就是最重要的.
在疫情過了之後.
在去年.
去年我們教會有一個全教會的晚宴.

$^{281}$很虛潭在沙田.
那次427.
一個全教會的晚宴.
負責人告訴我.
那次動用了幾百位弟兄姊妹來事奉.
幾百位.
好像我沒有記錯.
是四百多位弟兄姊妹.
在晚宴之後.
全部人一起宣告一件事.
錦宣家齊躍動.
在那次晚宴之後.
我就看到.
大家 弟兄姊妹.
開始動了.
朝聖茶室.
以前有時冷清清的.
是不是.
但之後我發覺到.
這個去飲用的多了.
更加重要的是.
我們看到.
參與事奉的弟兄姊妹.
越來越多了.
所以弟兄姊妹.
當我們享受的時候.
想一想.
我可不可以抽一些時間來幫忙呢.
我都看到.
去推動這個教會大旅行.
雖然最後都不是太多人參加.
但我很欣賞.
那些負責的弟兄姊妹.
每一個堂的之後.
我就看到他們.
拿著單張派給人.
拿著筆.
拿著報名表.
就叫人簽名鼓勵人去.
我看到很感動.

$^{321}$弟兄姊妹都很希望能夠.
將一班弟兄姊妹.
錦宣家的人.
圍在一起.
到今年.
一會兒.
1月13日.
那個下午茶.
在俱樂部.
嘩.
很快一推出.
就很快就滿了.
我相信不是因為只有20元.
就可以在俱樂部喝下午茶.
而我相信.
弟兄姊妹.
我們說得很明白.
那次是整個的.
全部的團契.
一起彼此的分享.
但我覺得.
都不夠的.
個個都參與這些活動.
不夠的.
在2024年的時候.
我一直都邀請.
每一位的弟兄姊妹.
我們需要.
為教會.
有意象.
為教會.
有使命.
定了的.
就是.
錦宣家.
齊躍動.
上年就開始動了.
不過.
盼望在今年.
新的年度開始.

$^{361}$我們錦宣家.
無論我們這裡的堂會.
紅印堂.
嘉欣堂.
我們能夠有更大.
更大的躍動.
能夠真是.
更多的弟兄姊妹.
參與不同的事奉.
我先說說.
在這個祈禱區.
過往就只是在教會裡面.
一齊有祈禱.
但是.
在新的一年度裡面.
我增加了.
施段的祈禱會.
在葉醫生家裡.
他們照舊在星期三晚.
K段.
就在朱太家裡.
他們在星期四晚.
我會繼續聯絡.
可能在錦秀.
會再開一區.
兩區.
其實星期六早上.
有個老人家的.
祈禱會.
星期六就是九點半到十點.
還有一個錦秀區.
以外都有祈禱會.
我們盼望.
有更多的人.
能夠去加入事奉.
特別是加入.
祈禱的事奉.
以致到在我們.
錦秀堂,嘉欣堂.
紅印堂裡面.

$^{401}$我們能夠.
齊躍動,齊看到.
上帝要我們.
三間不同的堂會.
我們可以有.
甚麼使命.
看看下面幾段經文.
第一段.
在路加福音第三章.
耶穌洗完禮之後.
他正禱告的時候.
天開了.
你是我的愛子.
我喜愛你.
當時耶穌基督.
洗了禮.
然後他就準備走上.
事奉的路.
就在那個時間.
祂正在祈禱.
上帝就去.
堅定耶穌的心.
你就是.
我的兒子.
是我所喜愛.
是我所差派你.
我高承就救恩.
去建立.
耶穌基督的心.
然後在路加福音第九章.
就在這個登山變像.
經文是這樣說的.
耶穌帶著彼得.
約翰雅各.
上山禱告.
正禱告的時候.
他的面貌.
改變了.
然後就是.
彼得約翰雅各.

$^{441}$就看到.
有摩西.
和以利亞.
和耶穌一起在聊天.
知不知道下了山之後.
他怎樣.
然後他就準備.
上十字架的路.
準備去受苦.
這條路.
那又是.
當耶穌祈禱的時候.
上山祈禱的時候.
上帝就有.
特別的事發生了.
然後.
在《使徒行傳》第十章.
這次是彼得.
見這個異象.
在中午的時候.
彼得就上了.
這個樓房去祈禱.
在祈禱中.
他就看到了.
看到天開了.
有一塊大布.
裡面就有不同的.
動物在.
過往彼得都覺得.
這些動物是骯髒的.
但是在異象裡面.
使者告訴他你要吃.
但彼得說.
不可以的.
這些骯髒的東西.
是不可以吃的.
然後在《使徒行傳》.
十三章三節.
「安然禁食禱告後.
給巴拿巴和.

$^{481}$素羅按手.
然後就.
差派他們走了」.
就是巴拿巴和素羅.
他們要起行了.
踏上這個宣教的旅程.
是保羅他們.
第一次的宣教旅程.
也是第一次.
就將福音.
從耶路撒冷.
然後去到安提阿.
然後就去到亞西亞.
再去到歐洲的地方.
剛才我只是選了.
這三段的經文.
我們都看到.
在禱告的時候.
有特別的東西發生.
有特別的使命出現.
有特別的.
力量.
就臨到.
我們看到.
原來當教會.
在劃時代的一些情況.
要出現的時候.
都是由禱告.
開始.
都是由禱告開始.
所以.
我們就說.
大家想不想在禁宣家裡面.
前面不單止2024年.
2025年.
並且以後.
我們能夠有一個.
大躍動.
如果頂子妹不動的時候.
請多多同工.

$^{521}$我們都沒有辦法.
推動教會.
所以我們禁宣家.
要為上帝.
有一個劃時代的發展.
就要由.
你和我.
的禱告.
來開始.
一起禱告.
求上帝.
讓我們禁宣家.
怎樣行動.
也讓我們看到.
我們自己有什麼使命.
有一個人就去.
投考教育學院.
做老師.
是蒲公英金奇說的.
面試官就說.
為什麼你想.
做老師.
那個人回答.
因為.
我做不到很偉大的人.
我想.
但我做不到.
所以我想做.
未來偉大的人的老師.
我覺得這個態度.
很值得我們學習.
大家.
上帝沒有呼召你們去做.
傳道牧者.
不過你們.
可以成為傳道牧者.
背後禱告.
的支持.
上帝沒有呼召你們.
回來教會事奉.

$^{561}$但是你們.
可以為教會裡面.
每年的使命.
目標.
來去禱告.
成為教會的.
禱告勇士.
我相信.
藉著這個禱告.
聖靈會.
引導.
教會的發展.
引導著教會編排.
這個計劃.
藉著禱告.
上帝會指引我們每一個人的前路.
我記錄.
藉著某一次.
禁食祈禱裡面.
上帝很清楚很清楚.
讓我知道我要踏上.
全時間事奉.
我要去讀神學這條路.
所以靖姐妹.
有時當你感到.
前路模糊的時候.
鼓勵大家.
多去禱告.
我認識一些人.
很迷糊的時候.
每天早上五點鐘起床.
站一個小時.
站一個小時禱告.
兩個星期後.
上帝告訴他們.
所以靖姐妹不要想著.
禱告是.
看空氣.
在這裡.
我就正式和大家.

$^{601}$發出這個邀請.
在這個金星期裡面.
如果你可以的話.
最好在這個月裡面.
甚至在今年的裡面.
每天至少.
用.
十五分鐘.
為教會祈禱.
除了自己.
個人的祈禱之外.
為教會祈禱.
為教會的發展.
為教會各方工作.
為教會的.
童工去禱告.
不單止為了錦宣家.
為了洪恩家.
為了嘉恩家.
主要我們這個.
事公國的發展.
是這個.
上帝去引導.
當上帝一路.
引導著我們每一段路程的時候.
我們會一起.
去經歷到.
聖靈的風一路帶領著的時候.
不用很出力.
因為聖靈.
就一路.
這樣去帶領著.
我相信.
在這裡所有的五堂的頂子們.
願意每一天.
用至少十五分鐘.
為教會的.
發展.
為教會的童工.
為教會各方工作來祈禱的時候.

$^{641}$我相信.
在這一年裡面.
教會有一個很大的.
躍動.
我記得在.
幾年前教會在播放.
那套.
福音電影.
就叫做.
很難讀的.
叫做.
光鵲嶺.
很難讀的.
光鵲嶺.
雖然那套電影有一點.
小小的血腥.
但是很感動.
是一個真人.
真事的改編.
講到在第二次大戰的時候.
有一個美國的.
軍醫.
他是基督徒.
他當時就去到.
日本的.
沖繩島.
的地方叫做.
光鵲嶺那裡.
去隨軍打仗.
雖然他是軍醫.
他都可以開槍.
可以殺敵軍.
不過他自己.
就立了一個志向.
他就覺得十戒教了.
不可以殺人.
不可殺人.
所以立定志向.
就算面對敵人也好.
他不會開槍.

$^{681}$在這個打仗期間.
他又向神.
立了一個.
志向.
他就說要救得.
一個救一個.
因為這個戰場.
上高很危險的.
所以有時有些人.
受了重傷.
他自己動不了.
很多軍醫都不救了.
因為他自己都要保住自己的命.
但是他不理.
他從上帝那裡.
立了一個志願.
救得一個就救一個.
結果他將很多重傷.
和很多行動的.
一些他的同胞.
都救回.
這個安全的地方.
大戰結束了之後.
美國總統頒了一個勳章.
給他.
表揚去拯救.
這個同胞這個偉大的精神.
弟兄姊妹.
在這個搶救的.
靈魂是上高.
我們都需要有.
這個精神.
救得一個.
救一個.
大家都很怕.
跟人講福音.
大家更加怕.
叫人問人.
你信不信耶穌.
我告訴你講完都很怕.

$^{721}$呼召的.
沒有人舉手.
很尷尬的.
是不是.
又覺得被人拒絕.
但是我們需要.
有這一份的精神.
救得一個就救一個.
就算.
不是這樣個別的傳福音.
而是去看一下.
上帝我們每一個人的生命裡.
給我們.
有這個不同的恩慈.
有不同的才幹.
也都有不同的時間.
有金錢有生命.
上帝已經給了很多很多很多.
東西給我們.
但我們怎樣運用上帝.
給我們這些每個人.
不同的特質.
我們在工作裡面.
我們去見證上帝.
我們為上帝.
來去工作.
我們藉著上帝給我們.
不同的恩慈.
我們都可能有不同的方法.
可以來去接觸到.
一些未信的群體.
所以弟兄姊妹.
在新的一年裡面.
給我們有不同的才幹恩慈.
金錢時間生命.
但上帝.
給我們這麼多東西.
目的是怎樣.
不是自己享受.
而是為他做一個好管家.

$^{761}$用在教會.
用在.
人群的當中.
所以弟兄姊妹.
多去禱告.
去問下神.
問下神.
在我生命裡面.
無論我們年紀多大.
都好 我們問下神.
在我生命的當中.
在我剩下來的日子裡面.
你給我有.
什麼的使命呢.
就好像剛才.
我說我們的姐妹.
讀經計劃都好.
有些她說我不行.
我只可以讀完一卷的福音書.
但結果就讀完四卷.
上帝.
給我們多一點 我們就做多一點.
給我們做少一點.
我們就做少一點.
上帝不是要.
我們個個都要做這麼多.
上帝要求.
這個管家.
要忠心.
所以就算我說.
原來今年六月退休了.
但上帝說.
我給你仍然有精神.
仍然能走.
你就要繼續做下去.
做到多久.
上帝說你夠了.
你不要做了 你就走去.
去屢下行.
去下短孫 又加上屢下行.

$^{801}$上帝.
就要我們找回我們的使命.
當我們.
行在上帝使命裡面的時候.
我們會發覺到.
我們的生命.
會更加豐盛.
更加精彩.
更加燦爛.
我們的生命就一生能夠.
冠耀上帝的名字.
在這個時候.
我就想大家.
有一個個人安靜的時間.
去和天父問一下.
不用問.
我一輩子了.
問我在這一年的裡面.
我可以.
做些什麼.
我可以怎樣為教會祈禱.
怎樣在教會裡面.
我有沒有一些事.
我可以做一下呢.
我們一起有一個安靜禱告.
個別禱告.
有一個安靜禱告.
個別禱告.
(安靜禱告).
在哲學聖經裡面.
還有.
沒有意象.
民就放肆.
放肆這個意思.
就是.
很輕率.
自把自為.
沒有了目標.
喜歡做什麼就做什麼.
我剛才告訴你.

$^{841}$如果我不定了.
要去看聖經.
去完成舊約的時候.
我就不會這樣.
但是.
當我們沒有了意象.
沒有了使命的時候.
我們就可能會游游行.
也都可能.
當我們有時.
少了去教會.
少了去聽上帝的說話的時候.
有時我們會得罪神.
好像在出埃及的時候.
當摩西上了.
西乃山裡面.
他就想去拿.
法板.
十戒.
但當山下的.
以色列人他們沒有了.
摩西做這個屬靈的領袖.
他們自己就游游行.
結果就不知道怎樣搞.
就搞到自己.
去做了這個金牛毒出來.
代表了上帝.
結果上帝的懲罰.
就臨到.
但是.
在整節的聖經裡面.
就說到.
有意象.
民就放時.
下一句是甚麼呢.
為遵守律法的.
便為有福.
那怎樣呢.
遵守上帝的說話.
遵守上帝的誡明.

$^{881}$我們可以有福呢.
我們看看.
這個.
在《箴言》裡面.
請聽我讀出第八章.
孩子們.
要聽從我.
謹守我道的.
有福了.
要聽訓誨.
得智慧.
不可棄絕.
要聽上帝的說話.
要學習上帝的說話.
聽從我.
天天在我門口守望.
在我門框旁等候的.
那人便為有福.
即是我們時時都要有上帝的說話.
因為尋得我的.
就尋得生命.
他必蒙耶和華的恩惠.
所以常駐在.
耶和華的.
上帝的說話裡面的人.
是有福的.
他的生命.
是有力的.
我相信.
一個社會.
一個教會.
一個家庭.
只要在裡面所有的人民.
所有弟兄姊妹.
多聽上帝的說話.
多讀他的說話.
多祈禱的時候.
我們每一個人的生活.
都會是幸福.
美滿快樂的.

$^{921}$我們看看.
我們今年的口號.
都是「錦宣家齊躍動」.
我們今年的口號.
都是「錦宣家齊躍動」.
這個一會兒還有兩點.
有四點.
一本書叫.
《做理形象的建築師》.
很多點.
人生要定一些目標.
是李順祥博士寫的.
李順祥博士寫的.
在當中抄一些.
改一些.
我們看一看第一第二點.
然後看第三第四點.
然後看第四點.
我盼望.
這四點都成為了.
我們今年的立園.
但我不強迫的.
一會兒我們會一起去讀.
若果這四點都成為了.
大家今年裡面的目標的時候.
我們就有力地讀出來.
若果我不想.
我不想.
那你可以不讀.
我們聖靈會知道.
也盼望聖靈鼓勵我們.
每一位.
我們都要努力的.
在使命上.
我們要努力.
好 若果你看到了.
想過了.
我們一起 如果你願意的.
我們一起有力.
將這四點讀出來.

$^{961}$我生命最高目標.
是為了榮耀上帝.
並頌讚他的美事.
我願意多操練上帝給我的恩賜.
並且運用他來榮耀.
父上帝.
我會努力建立與弟兄姊妹間的友情.
彼此相愛.
我願意勤讀聖經.
祈禱.
讓上帝改變我的生命.
好 大家都忠心的讀完了.
我們一起唱「要忠心赴歌」.
要忠心 忠心.
此後何處可釋放.
忠心 要忠心.
終於你會了我.
要忠心 忠心.
總別人多不尊重.
無論何境況.
別站在主旁.
要尊心 連你也忠心.
要忠心 忠心.
此後何處可釋放.
忠心 要忠心.
終於你會了我.
要忠心 忠心.
總別人多不尊重.
無論何境況.
別站在主旁.
要尊心 連你也忠心.
我們重新禱告.
因為你給我們每一個人有不同的才幹.
恩賜 時間 金錢.
你也給我們有不同的特色.
願意你叫我們錦宣家每一位弟兄姊妹.
都安住你給我們所有的.
我們為你做一個忠心的管家.
在過往的日子.
或者我們會有冷淡了 懶惰了.

$^{1001}$求聖靈不斷的催逼.
以致我們都要成為主裡面的基督正兵.
為天國來去努力.
求主振興.
求主復興整個錦宣家.
求主工作.
多謝你垂聽禱告.
是奉靠主基督聖名.
阿們.
\newpage



\section{馬太福音 2:1-12}
\label{sec:Y_5AXHxN_yY}
\textbf{2024-01-07 早堂《星光引導》馬太福音 2\_1-12 講員:曾浩賢院長}
\newline
\newline
連結: \href{https://youtube.com/watch?v=Y-5AXHxN_yY}{\texttt{https://youtube.com/watch?v=Y-5AXHxN\_yY}} ~~~~ 語音日期: 2024-01-09
\newline
\newline
\hyperref[sec:SWoA7YUZj_4]{\small{< < < PREV SERMON < < <}}
~
\hyperref[sec:index_chronic]{\small{[返順時目]}}
~
\hyperref[sec:index_scriptual]{\small{[返順卷目]}}
~
\hyperref[sec:2um86Ip2zrs]{\small{> > > NEXT SERMON > > >}}
\newline
\newline
馬太福音 2:1-12
\newline
\begin{longtable}{cl}
\hline
\hline
章節 & 經文 (和合本修訂版)\\
\hline
2:1 & \begin{tabularx}{0.7\textwidth}{X} 在希律作王的時候，耶穌生在猶太的伯利恆。有幾個博學之士從東方來到耶路撒冷，說： \end{tabularx} \\ \\ \relax
2:2 & \begin{tabularx}{0.7\textwidth}{X} 「那生下來作猶太人之王的在哪裡？我們在東方看見他的星，特來拜他。」 \end{tabularx} \\ \\ \relax
2:3 & \begin{tabularx}{0.7\textwidth}{X} 希律王聽見了，就心裡不安；耶路撒冷全城的人也都不安。 \end{tabularx} \\ \\ \relax
2:4 & \begin{tabularx}{0.7\textwidth}{X} 他就召集了祭司長和民間的文士，問他們：「基督該生在哪裡？」 \end{tabularx} \\ \\ \relax
2:5 & \begin{tabularx}{0.7\textwidth}{X} 他們說：「在猶太的伯利恆。因為有先知記著： \end{tabularx} \\ \\ \relax
2:6 & \begin{tabularx}{0.7\textwidth}{X} 『猶大地的伯利恆啊，你在猶大諸城中並不是最小的；因為將來有一位統治者要從你那裡出來，牧養我以色列民。』」 \end{tabularx} \\ \\ \relax
2:7 & \begin{tabularx}{0.7\textwidth}{X} 於是，希律暗地裡召了博學之士來，查問那星是甚麼時候出現的， \end{tabularx} \\ \\ \relax
2:8 & \begin{tabularx}{0.7\textwidth}{X} 就派他們往伯利恆去，說：「你們去仔細尋訪那小孩子，找到了就來報信，我也好去拜他。」 \end{tabularx} \\ \\ \relax
2:9 & \begin{tabularx}{0.7\textwidth}{X} 他們聽了王的話就去了。忽然，在東方所看到的那顆星在前面引領他們，一直行到小孩子所在地方的上方就停住了。 \end{tabularx} \\ \\ \relax
2:10 & \begin{tabularx}{0.7\textwidth}{X} 他們看見那星，就非常歡喜； \end{tabularx} \\ \\ \relax
2:11 & \begin{tabularx}{0.7\textwidth}{X} 進了房子，看見小孩子和他母親馬利亞，就俯伏拜那小孩子，揭開寶盒，拿出黃金、乳香、沒藥，作為禮物獻給他。 \end{tabularx} \\ \\ \relax
2:12 & \begin{tabularx}{0.7\textwidth}{X} 因為在夢中得到主的指示，不要回去見希律，他們就從別的路回自己的家鄉去了。 \end{tabularx} \\ \\
[1ex]
\hline
\hline
\end{longtable}
$^{1}$各位弟兄姊妹早安 新年快樂.
今天是2024年第一個首個的主日.
也跟大家說聲新年快樂 願大家身體安康.
不如也和旁邊的弟兄姊妹祝福一下.
互相新年快樂 身體安康 新年蒙福.
讓我今早再一次有機會來到大家當中.
一同崇拜並且去做港道的服事.
同樣地也容許我美麗港道醫師.
先代表信義中神學院一眾師生同工.
向宣導會錦繡堂的眾弟兄姊妹.
參與崇拜的每一位文安.
願因為平安從我們的父上帝.
並主耶穌基督歸於在座的每一位.
今天選了段經文是.
《馬太福音》二章一到十二節.
講題是星光引導.
段經文相信很多弟兄姊妹都非常熟悉.
小朋友可能都會很熟悉.
因為可能在朱一學裡面.
或者在一些話劇裡面.
都會看到博士朝拜聖嬰這一幕.
今天為何會選擇這段經文和大家分享呢.
是因為在西方教會的傳統.
其實昨天一月六日.
我們稱它為顯然日.
十二月二十五日我們慶祝了聖誕.
過了十二日之後.
就會去到一月六日.
在教會傳統的年曆.
就會定這一天為顯然日.
顯然日是教會定出來.
是想提醒信徒.
是上帝將這個救恩向世界顯明.
而每一年的顯然日.
一月六日.
我們都會讀到《馬太福音》二章一到十二節.
就是博士從東方來朝拜聖嬰.
為何他們會來呢.
因為他們看到一顆星.
那顆星我們常常都放在聖誕樹的頂.

$^{41}$是上帝透過這顆星.
向這幾個東方的博士顯明.
救恩已經臨到這個世界.
所以當我們每一年幾月六日去讀經文的時候.
其實一方面是讓我們去提醒上帝顯明自己.
另一方面亦都透過這個經文.
讓我們去再想一想.
當上帝向我們顯明的時候.
人的回應又是怎樣呢.
所以就著這兩重的意義.
今日都和大家分享三點.
第一 顯明.
博士朝拜聖嬰這個故事告訴我們.
剛才都強調過了.
是上帝主動地向所有人顯明他自己.
是上帝主動.
我重複 是上帝主動.
有時有些信徒.
慢慢地就會覺得是我們信耶穌.
我們去尋求上帝.
我們尋找教會.
我們就去返教會.
好像是我們主動.
我們能夠得到救恩.
是我們的主動.
但是一月六日再一次提醒我們.
如果不是有那顆星.
博士不能夠去到伯尼克.
如果不是上帝顯明他自己.
我們根本不能夠去選擇去信不信耶穌.
同不同意嗎.
不是我們主動 是上帝主動.
講到這裡想和大家分享一首詩歌.
那首詩歌很靚.
那首詩歌叫做《沉著你的那天》.
大家聽過沒有.
如果沒有歌詞都可以搜尋一下.
《沉著你的那天》.
容許我唱幾句給大家聽.
傳說人生.

$^{81}$高了一點 不好意思.
傳說人生只不過是個夢.
無方向 無意義 沒內容.
貧窮或富足 結局都一樣相同.
如炊煙飄散 埋在糞土中.
靚不靚.
很感受到那種迷惘之中.
不知如何去向.
在迷霧之中不知如何去向.
意境很靚.
然後副歌就說.
直到一天 我再沉細沉著你.
徬徨空虛 悄然消失遠飛.
很厲害 你們馬上彈到歌詞出來.
專業就是專業.
不用我自己一個唱.
一起唱好嗎.
我們唱第二節 好嗎.
預備起.
茫茫人海 找一塊落美夢.
無傷痛 無怨恨 沒內容.
流浪在世間 試圖找圈難採紅.
但千山走遍 仍未見影蹤.
直到一天 我再沉細沉著你.
徬徨空虛 悄然消失遠飛.
與你偶遇 原是我畢生的福氣.
如黑夜盡頭 重現神氣.
在那一天 我再沉細沉著你.
從前孤單 我從此不記起.
你似盼望承諾借愛永沒盡期.
灰色人生變成絕美.
多謝大家 我們唱到這裡.
其實是一首很美的詩歌.
就是想告訴我們.
在我們人生裡面.
在一個重重迷霧.
甚至乎很困惑的時間.
我們能夠和上帝結連.
其實是給我們一個很大的福氣.
這首詩歌是一首很美的詩歌.

$^{121}$不過這首詩歌唯一一個美中不足的地方.
是在符國.
不是我們沉著耶穌.
是耶穌沉著我們.
不是我們偶遇耶穌.
不是偶然.
是耶穌上帝親自向我們顯明他自己.
今天你能夠回教會.
是什麼原因.
是不是你自己走來這個教堂呢.
還是有人傳福音給你呢.
如果上帝不顯明他自己.
又怎會有傳福音的人呢.
所以今天在1月6日.
我們讀星光引導.
馬太福音2章1到12節的經文的時候.
我們第一件事再一次提醒.
這個是恩典.
是上帝顯明.
是祂主動.
不是我們.
我們只是回應.
而且這個顯明還要告訴我們.
顯明不只是給猶太人.
是給世界上所有的人.
而博士來.
從東方來的博士.
是代表著外邦人.
上帝這個救恩.
不單止向猶太人顯明.
也都向外邦人顯明.
今天我們在香港.
如果按血統.
我們都是屬於外邦人.
我不知道網上有沒有.
或者在座當中有沒有猶太人.
大部分我們都是外邦人.
我們都明白原來上帝的救恩.
都臨到外邦人.
對我們來說不是一件驚訝的事.

$^{161}$不過在兩千多年前的時間.
其實是一個很震撼的訊息.
上帝的救恩.
不單止給猶太人.
也都給外邦人.
即是說今天我們都能夠.
得到上帝給我們啟示的救恩.
不單止給猶太人.
不單止給外邦人.
更加的在這個經文裡面.
我們看到還有兩類人.
上帝同樣向他們顯明.
這個寶貴的福音.
寶貴的救恩.
第一類是心懷惡念.
叛逆上帝的人.
我這裡所說的就是希律王.
希律王是一個怎樣的王.
他是一個不敬畏上帝的人.
我相信大家都會認同我的說法.
這位為保權位.
他會不惜一切去行惡的一位王.
即使他不懂得升將.
但是上帝仍然透過這幾位東方尼的博士.
告訴他.
星光之下引導著他們.
來尋找生尼作猶太人的王.
這個彌賽亞.
告訴希律聽.
彌賽亞已經來到.
其實都有點奇怪的.
是嗎?.
即是星光引導這班的博士.
去尋找希律王.
因為其實希律王是惡人.
為何上帝要引導這班博士.
去尋找這個惡人.
告訴他聖嬰已經誕生.
因為上帝不知道這位惡人.
會對這位聖嬰有惡意嗎?.

$^{201}$會想殺害他嗎?.
不過仍然.
聖經告訴我們.
上帝就引導這班東方博士.
先去希律王那裡.
告訴他.
彌賽亞已經來到.
為何上帝要向這位大惡人.
顯明基督的來臨.
為何一位殘害小聖嬰的人.
要讓他知道這個消息.
都是那一句.
上帝要向所有人.
顯明他的救恩.
包括犯罪作惡.
背叛上帝 敵黨他的人.
在今天的經文除了猶太人.
外邦人 懷惡念.
背叛上帝的人.
上帝都向不冷不熱的人.
去顯明他的救恩.
在經文裡邊.
哪些是不冷不熱的人呢?.
我指的是那班祭司長和文士.
祭司長和文士熟讀聖經.
所以當希律王問他們的時候.
基督該生在何處.
他們二話不說就引經據典地回答.
在猶太的伯利恆.
很熟.
不過對於這些宗教領袖而言.
他們日常生活都很宗教.
讀聖經 祈禱 講解聖經 獻祭.
我們可以相信.
在他們的言談之間.
常常都會談到彌賽亞來臨的應許.
就好像我們每逢聖誕節.
我們會說耶穌降生.
不過當他們聽到博士說彌賽亞.
有關他的聖.

$^{241}$希律問他們 他們回答了.
如果是一個愛慕上帝的人.
他應該是怎樣的呢.
是否應該會立即出發.
去尋找這位彌賽亞.
不過他們沒有.
他們沒有甚麼感動.
沒有甚麼雀躍.
又沒有立即隨從這些博士.
一起去找這個小聖嬰.
在伯利恆究竟在那裡.
去朝拜他.
找這位猶太人.
應該是引頸學待.
等待他來的基督.
對於這種反應.
就算我是用最正面的想法去想.
都只能用不冷不熱來形容他們.
同樣地.
為何上帝要向這些不冷不熱的宗教領袖.
顯明基督已經來臨呢.
答案除了同樣是.
因著上帝要向世人.
所有的人.
顯明他自己的救恩.
我也聯想到啟示錄當中.
上帝差派約翰.
寫信給那些不冷不熱的老弟家教會.
希望這些不冷不熱的人.
都有著一個改變.
回轉的機會.
頂子梅上帝的愛就是這樣.
這麼長闊高深.
他愛世間所有的人.
即使人處於.
不論你是甚麼種族.
甚麼膚色.
不論你的狀態是處於邪惡.
活逆不冷不熱.
上帝仍然會向我們顯明他的愛.

$^{281}$再說.
這幾位博學之士.
對於這幾位博學之士來說.
當他們知道生下來要作猶太人之王.
不是屬於希律王的後代的時候.
我心中其實也有一種疑問.
為何他們意識不到問題所在.
聞不到血腥的味道.
他們還好像相信希律的說話.
希律跟他們說.
你仔細去尋找那個小孩子.
找到就來報信.
我去拜他.
似乎這班博士是相信希律的說話.
拜訪耶穌之後.
真的打算回去告訴希律.
或者可以這樣說.
即或人的智慧不足.
即或人有時會好心做壞事.
但是上帝仍然願意向所有的人.
顯明他自己.
親愛的弟兄姊妹.
無論現在我們有著甚麼身份.
保障昔日的猶太人,愛邦人.
處於甚麼狀態.
我們遠離了上帝.
做一些他不喜歡的事.
不冷不熱,智慧不足.
有時好心做了壞事.
今日經文同樣告訴我們.
我們都在主耶穌這一顆星的引導之下.
這個星光仍然照亮我們.
上帝的愛,上帝的救恩和旨意.
都要向我們顯明.
除了顯明.
在讀經文.
讓我們反思的另一個重點.
是講到人的回應.
對於這個星光引導的回應.
要講故事裡面的人.

$^{321}$對上帝顯明這個救恩的回應.
同樣都可以從希律王講起.
希律王聽到幾位博士的說話.
經文描述他第一個反應是甚麼.
心裡不安.
對擁有王位的希律來說.
彌賽亞的彌林對他來說不是福音.
不是覺得高興,不是覺得感恩.
他是覺得不安.
為甚麼.
因為明顯他是憂慮到自己的王位不保.
受到威脅.
而且要消除這個威脅.
要保住他自己的既得利益.
他就對上帝的顯明.
他的回應就是剷除.
即使要剷除的對象是上帝而來的彌賽亞.
似乎希律王都在所不惜.
希律王腦裡不是沒有上帝.
他是那位大興土木修建聖殿的人.
到現在要起義剷除從上帝而來的彌賽亞.
他也懂得要向博士假冒為善一番.
說自己都要去拜他.
即是他知道是上帝.
不過他就是這樣去硬心.
他頭腦知道上帝.
但他心中沒有上帝.
他心中的上帝就是他自己.
因為他認為自己有能力.
可以隨己意去殺害這位彌賽亞.
至於腦裡那位猶太人所敬拜的上帝.
他可能只是當他是一個政治上的工具.
他的回應就是硬心 剷除 殺害.
繼續維護自己的既得利益.
說完希律 我們又說一下祭司長和文士.
剛才說過他們是一班不冷不熱的人.
這些宗教領袖對彌賽亞可能已經離林的訊息不冷不熱.
而他們的回應除了準確答了希律王的問題之外.
他們什麼都沒有做.
如果說博士是蒙福的話.

$^{361}$因上帝透過這顆星去引導這班博士去找到基督.
其實這班文士和法利賽人.
他們有機會去熟讀律法和先知書.
其實他們更有福氣.
因為上帝早早已經透過聖經向他們啟示和顯明.
不過很可惜.
他們就好像一班頭腦上有很多神學知識的一班神學家和牧者.
可以回答很多神學的問題.
可以在宗教活動上做領導 做得很好.
但是那些知識沒有驅使他們和上帝之間建立一種緊密的關係.
各位弟兄姊妹.
剛才提到希律他有些既得利益.
所以他硬心拒絕 顯明 甚至乎要剷除 殺害.
上帝有沒有呼召我們回轉歸向他的時候.
可能是因為我們一些不想放下的名利 既得利益.
我們好像希律那樣去硬心不回應.
文士和法理才人 祭司長這一班腦裡有很多宗教知識.
在教會裡面很熱心事奉.
都可以和人說很多關於耶穌的事情.
但會不會我們其實心裡面和上帝距離很遠.
而至到當上帝的訊息 他們的啟示 他們的呼召回轉來到的時候.
我們都好像這一班祭司長 文士等等 都是無動於衷.
還是我們好像這幾位博士.
博士對於上帝顯明這個回應.
他們是千里迢迢來到耶路撒冷.
再到伯利恆去找這位生下來作猶太人的王的英孩.
他們帶著寶貴的禮物 黃金 玉香 活藥.
不單止這樣 聖經描述說他們找到耶穌的時候.
他們是非常之歡喜.
而且這幾位又有錢 又有學識的成年人.
而話不說的就俯伏 俯伏去拜這個小孩子.
千里而來非常歡喜 俯伏敬拜 獻上禮物.
這些都和我們剛剛所說的希律王 祭司長 文士的回應有著很大的對比.
那些本來應該是最近上帝 最認識上帝的人.
那一刻和上帝很遠.
反而本來和上帝很遠的人.
這一刻和上帝很近.
就在他腳前俯伏敬拜.
各位弟兄姊妹 我們能夠成為基督徒.
就好像博士一樣 原本是遠的.

$^{401}$去到主面前很近.
但當我們成為基督徒之後.
會不會慢慢慢慢變成了希律.
祭司長和文士好像很近.
但其實很遠呢.
我們對上帝對我們的顯明.
是不是一種感恩 一種渴慕.
上帝的道來到我們當中向我們顯明的時候.
我們是不是用心聆聽.
然後在我們的生活裡面走出來呢.
親愛的弟兄姊妹.
願我們都好像博士一樣.
在星光引導下.
我們能夠可以歡喜 福福 獻上.
最後 說到回應.
剛才說的是人的回應.
我想第三點和大家分享的是.
想問一問大家.
用一個以人法的方式去問一問大家.
我們願不願意我們的回應.
不單止是去歡喜 福福 獻上.
我們願不願意都成為那顆引導人的星.
用一種以人法的方式去問一問大家.
我們的回應會不會再進一步成為那顆引導人的星.
作為基督徒或者你會覺得我是多此一問.
因為所有基督徒都不會反對.
因為所有基督徒都會知道.
我們傳福音就是我們的使命.
所以我們都應該好像一顆星一樣.
引導人去找耶穌.
對的.
我相信在座當中很多人都會願意成為那顆星.
不過我想讓大家有一個心理準備.
就是要當這顆星並不是人見人愛 車見車載.
剛才說到希律王討厭那顆星.
祭司長文士沒有因為那顆星有任何的行動.
猶太人只是為那顆星的出現感到不安.
沒有歡慶 沒有行動.
似乎那顆星在整個世界上 在這個故事裡.
只是成功引領了幾位從東方來的博士.

$^{441}$就連自己人 猶太人 即是上帝的子民.
都拒絕了那顆星的引導.
即是說原來我們去傳福音.
我們做這顆星 你不要想著很受歡迎.
有時是會受很多人的拒絕.
很多人的不理會.
甚至乎希律王懷著惡意去逼迫.
和這顆星有關的事情.
我們有沒有做好這個心理準備.
有沒有知道這個情況下.
我們仍然願意成為這顆星.
我不知道大家有沒有聽過.
有一些信徒很喜歡參加短宣.
去柬埔寨 泰國 非洲那些去做短宣.
為何他們很喜歡做短宣.
但在自己的地方 即是香港.
又不見他們很努力傳福音.
但他們就很喜歡去短宣.
有些原因是因為上帝的感動.
特別是關心某個地方的人士.
這個亦都有.
但我亦都聽過一些答案.
他告訴我 在那裡傳福音的滿足感很大.
因為你去那邊搞報道會.
呼召人信主 很多人信主.
很多人流淚 很多人舉手決志信主.
但在香港 很多人拒絕.
邀請人來報道會都很困難.
然後報道會裡面有幾個人舉手決志.
很開心 所以沒有甚麼滿足感.
所以喜歡去短宣.
因為短宣很弱滿足感.
各位靈性姐妹.
沒有滿足感 願不願意都成為這顆星.
有困難 受迫迫.
願不願意都成為這顆星.
上帝就是這樣用這顆星.
去啟示希律 祭司長 文氏 東方利的博士.
不是很成功 幾位而已.
不過祂願意.

$^{481}$我們願不願意呢.
我們同心祈禱.
天父上帝我們多謝你.
多謝你透過東方博士朝拜聖嬰這個經文再一次提醒我們.
提示我們.
我們是蒙福的一群.
不是我們主動去尋找你.
是你一早就去尋找我們.
向我們啟示你的救恩.
而且這個啟示告訴我們.
無論我們是什麼人.
現在是什麼狀態.
你都是愛我們.
邀請我們去到小聖嬰的面前.
去到耶穌的面前.
主我們多謝你.
多謝你這一份無條件這麼大的愛.
主求你幫助.
七獻聖靈臨到每一個弟兄姊妹當中.
去幫助我們.
讓我們對這個星光引導.
我們有回應.
不是希律.
不是祭司文士.
是好像博士.
願意去到你的面前.
俯伏獻上.
心裡也歡喜.
更求聖靈幫助我們.
不單止願意去到你的面前.
也願意被你使用成為一顆星.
去引導其他人來到你的面前.
縱然這顆星不是很受人歡喜.
但是我們都願意.
因為我們都是這樣被引導過來.
主我們再一次為到.
今天你所給我們的所有訊息獻上感恩.
也繼續求你讓我們在新的一年.
2024年.
我們就成為一顆星.

$^{521}$我們這樣禱告奉耶穌基督的聖名.
Amen.
\newpage



\section{希伯來書 11:8-9}
\label{sec:2um86Ip2zrs}
\textbf{2024-01-14 午堂《亞伯拉罕信心的成長》希伯來書 11\_8-9 講員:楊明麗姑娘}
\newline
\newline
連結: \href{https://youtube.com/watch?v=2um86Ip2zrs}{\texttt{https://youtube.com/watch?v=2um86Ip2zrs}} ~~~~ 語音日期: 2024-01-16
\newline
\newline
\hyperref[sec:Y_5AXHxN_yY]{\small{< < < PREV SERMON < < <}}
~
\hyperref[sec:index_chronic]{\small{[返順時目]}}
~
\hyperref[sec:index_scriptual]{\small{[返順卷目]}}
~
\hyperref[sec:Rlc0zB_br4Y]{\small{> > > NEXT SERMON > > >}}
\newline
\newline
希伯來書 11:8-9
\newline
\begin{longtable}{cl}
\hline
\hline
章節 & 經文 (和合本修訂版)\\
\hline
11:8 & \begin{tabularx}{0.7\textwidth}{X} 因著信，亞伯拉罕蒙召的時候就遵命出去，往將來要承受為基業的地方去；他出去的時候還不知往哪裡去。 \end{tabularx} \\ \\ \relax
11:9 & \begin{tabularx}{0.7\textwidth}{X} 因著信，他就在所應許之地作客，好像在異鄉，居住在帳棚裡，與蒙同一個應許的以撒和雅各一樣。 \end{tabularx} \\ \\
[1ex]
\hline
\hline
\end{longtable}
$^{1}$各位早晨.
今天是一個很特別的日子.
相信大家都會很忙碌.
其中一樣就是我們下午有一個特別聚會.
記不記得.
有差遣禮.
我們稍後都鼓勵大家.
稍後就盡量抽時間一起參加.
在差遣禮之前我們又提一個人.
這個人叫做阿巴拉罕.
阿巴拉罕他又得到神一個特別的差遣.
他很蒙恩地在他人生的前半段時間.
神就差遣了他去一個特別的地方.
做一件特別的差事.
他因為信.
在夢召的時候就遵命出去.
其實夢召的時候要能夠遵命出去.
真的要下定很大的決心.
去到那一處將來要承受為基業的地方.
他就向著那裡出發.
因為深信一定是可以得到的.
所以就算去到那裡.
他沒有一個固定的地方住.
他有一個好像作客一樣.
去到一個人家的地方.
一個異鄉.
住在一個像棚的地方.
但是他一個像棚的地方.
他好像是看到那個應許地.
阿巴拉罕是誰呢.
阿巴拉罕本名阿巴蘭.
在他99歲的時候.
神就為他改了一個名字叫做阿巴拉罕.
為什麼要改這個名字呢.
阿巴拉罕的名字的意思是多國之父.
本來阿巴蘭就是一國之父.
都已經很偉大了.
但是神幫他改了另外一個名字.
他就是多國之父.
他本來生於兩岸流域.

$^{41}$伊拉克米索不達米亞的吾爾的地方.
去到75歲的時候就開始遷居去加拿大.
最後他真的終身留在加拿大.
他死在加拿大.
應該就是死在加拿大的希伯倫這個地方.
也都葬在希伯倫的麥比拉洞.
在1075年告壽.
一生夢神福音.
聖經記載他的呼召是這樣的.
耶和華對阿巴蘭說.
你要離開本地 本族 父家.
往我所要指示你的地方去.
這句其實是很容易讀的.
我想小學生都懂得讀.
現在兩文三語他們什麼語言都讀得到.
不過要做是很不容易的.
很不容易的.
搬屋都不容易.
搬屋都要想過先.
要想得很清楚.
搬去什麼地方呢.
搬得更漂亮.
更漂亮都要我適應的 是不是.
搬屋不容易.
搬區呢.
我要找回一些朋友.
又不容易.
找一間熟悉的餐廳喝一頓茶.
都不容易.
現在是要離開這個地.
離開本族 離開父家.
所有親戚朋友都不會見到面.
我都不是 我都住在香港.
我只不過在元朗區.
受這裡服侍.
我住在天水圍.
我的夫家的親戚.
就在港島區.
港島東區.
港島東區到南灣半島.

$^{81}$與元代是一個斜角的關係.
我發覺很好.
他們現在個個都開始面臨退休.
時時約就閒日飲茶.
之後就說如果9點半.
去太古城哪裡飲茶.
我心想.
我都想去.
那麼遠.
你就算是我放假那天禮拜一.
我都不想那麼早起床.
走過去港島東區太古城.
雖然西鐵很方便.
是的.
離開父家.
有多少孤單.
有多少無助.
有什麼事要找人.
你為我祈禱.
因為大家都不方便.
還要是離開本族.
我想我和我家人就像離開本族.
過往四年未見過面.
電話都有望.
我都想起.
我最年輕的侄子有多高.
其實我都不知道他們有多高.
我只問他大學畢業了嗎.
還沒有啊姑姑.
有一段距離.
明明是姑媽.
不過都不是說廣東話.
我是姑媽.
但是大家那種認識.
好像這四年只能夠在電話裡.
我先生家人就很疼我.
經常問我.
你想吃什麼.
我就心想.
如果我空色想吃的東西.

$^{121}$你找不到給我的.
除非你家的傭人姐姐幫我煮.
所以我是不會說的.
是的 當你離開本族.
去到異鄉的時候.
很多東西都沒有.
不過感受的弟兄姊妹很好.
前幾天有弟兄姊妹請我喝椰汁.
吃椰肉.
是不是.
其實原來很多東西.
離開本族你要想過.
但是亞伯拉罕.
亞伯拉罕是離開本族.
離開父家.
離開他的本地.
真的離開本地的.
在你們的右手邊.
中下.
最後那個字「吾爾」.
那個就是他的出生地.
他的本地.
去到你的左手邊.
海之前的一些.
中下缺的地方.
那個地方有個事件.
那個就是他移居之後落腳地.
他去到就捉了一隻壇.
敬拜神.
因為那時候他們的敬拜的模式.
就是捉壇敬拜.
是不是很遠啊.
很遠啊.
中間全部沙漠地走不過的.
於是乎他們要繞一個長路.
北上一點.
向西北上.
接著向南下.
因為要沿著河流地域.
去那麼遠 去不去啊.

$^{161}$又有時差 天氣又不適應.
真的要想過才行.
還要是吾爾是一個繁榮的地方.
現在考古發掘出來.
他們四千年前的建築物.
已經是很有系統的流語了.
嘩 去到事件的地方.
他們做什麼呢.
先搭座棚住一下.
這個都是阿巴拉罕大半世的生活.
為什麼要去這個地方.
他不是一個人去的.
通常移民現在有兩種.
一種是拔起根.
全家一起去 屋都買了.
表明是去硬的了.
不會無情白事回流的.
另外一種不是的.
另外一種就是.
其中一個去住幾個月.
看看怎樣.
覺得OK嗎.
回來再看看.
帶個兒子去讀書.
兒子都適應嗎.
大家都適應了.
那爸爸才辭職.
或者退休了.
接著就想想房子.
還是留著先.
再過去那邊看看怎樣.
是不是.
要用幾年去考慮移民.
阿巴拉罕他是怎樣移民的.
阿巴拉罕的移民是拔起根的.
他帶著他所有的人.
所有的財物.
進入迦蘭地.
連他的侄子羅德.
他都帶過去.

$^{201}$因為這個侄子的爸爸去世了.
他就看著他.
阿巴拉罕還沒有兒子.
可能當他是兒子.
就一直帶著他照顧他.
過了這麼大的山嶺.
去到這個地方.
進去的感覺是怎樣呢.
不見得那個地方的人會很歡迎他.
所以他去到其實是搭帳棚住.
搭帳棚容易遷居.
他要養他的羊.
他的牛.
這個是阿巴拉罕的遷居.
神對他說.
你要離開你的本地本族和父家的時候.
他的反應是.
起行.
接著就去了.
接著神又對他說.
我一定使你成為大國.
我一定賜福給你.
使你的名為大.
你要使別人都得福.
這個呼召有後文好像更吸引.
不過進入那個地方.
完全不覺得有這回事出現.
而且還出現了一件事.
他帶著侄子去.
侄子一路成長的時候.
大家的財物都多.
接著就變成要分家.
侄子出去獨立了 自立了.
侄子自立 獨立了之後.
其實就出現一個問題.
其實阿巴拉罕你這麼努力做什麼.
你有這麼多人.
葬叮也有318個.
你有很多瀑被.
你有很多奴隸.

$^{241}$有很多牛羊.
正所謂打跛雙腳都不用休.
你勞碌來做什麼.
將來留給誰呢.
侄子都獨立了.
是喎 應該是.
阿巴拉罕為什麼還要努力不退休.
為什麼要留在那裡.
就是這一句話.
我定誓你成為大國.
我必定賜福給你.
誓你的名為大.
你要使到別人得福.
他還沒有一個兒子.
他連一個兒子都沒有.
他到八十歲都沒有.
他開始進入這個地方的時候.
是75歲到80歲都沒有.
於是乎我就在想.
其實他是為了什麼事呢.
但是他就是很認真去.
很認真去信這個神.
他因為信這個神.
接受了這個呼召.
於是乎他離開本地本家 父家.
於是乎就進入到這個迦南地.
進入到迦南地之後.
侄子離開他.
但是他仍然篤信這個神.
繼續在這個地上來努力.
他繼續在這個地上努力的時候.
不是為了下一代.
因為都沒有下一代.
他無謂為了那些瀑皮來努力.
照顧他們一輩子都夠了.
那麼辛苦做什麼.
但是他要信.
他仍然信神要叫他的名為大.
要叫他成為別人的祝福.
要叫他的國為大.

$^{281}$他的名有沒有為大呢.
有.
他就被人抓了.
因為打仗.
原來中東那邊經常都打仗.
他的侄子在打仗當中的時候.
被人擄了.
被人擄了之後.
亞伯拉罕就帶著他318位的壯丁.
去救他的侄子回來.
他的命一定宣揚得很大.
他也獻了一個很大的祭給神.
那就救吧.
留在這裡做什麼.
應該害怕的了.
打完仗.
是不是.
他繼續在那裡.
去事奉神.
但是神就跟他說.
應許他.
他要有個兒子的.
神應許他有很多財物的時候.
他跟神說.
你答應我這麼多事情做什麼呢.
不用了 天父.
神啊 你答應我的事情.
我將來留給誰呢.
不過我已經訓練了一個人.
來承繼我的了.
那個也是我的僕人而已.
將來承繼我的就是我的僕人.
你不用賜那麼多福氣給我.
哇 真的很少人這樣祈禱的.
但是呢.
神就跟他說.
不是啊.
我要賜個兒子給你.
你這個兒子才是你承繼人.
他太太知道之後.

$^{321}$都覺得不要阻礙著丈夫了.
於是就將那個婢女給了亞伯拉罕.
那個婢女真的為他生了一個兒子.
叫做以實瑪利.
亞伯拉罕你沒有信心啊.
不等神.
從撒拉幫你生了一個兒子.
人家等了十幾年.
75歲收到這個呼朝.
85歲納妾.
我們祈禱了十年.
未得應許的時候.
我們才說亞伯拉罕沒有信心.
他已經祈禱了十年了.
不過那個時候神沒有告訴他.
是撒拉生的才是.
於是他就生了一個兒子.
叫做以實瑪利.
後來到他99歲.
神就幫他改了個名字叫做亞伯拉罕.
他不再叫亞伯蘭.
就叫亞伯拉罕.
還要在這個時候.
神又跟他說.
你要從撒拉生一個兒子.
他才是繼承人.
不過你們就要在這個時候開始做一件事.
記不記得是甚麼事.
他們要行國禮.
行國禮的時候.
以實瑪利是有份的.
以實瑪利是跟他爸爸亞伯拉罕.
同一天行國禮的.
神沒有說歧視以實瑪利.
神沒有說這個不是我應許的.
我不理你.
你自己搞出來.
自己搞定它.
沒有.
神依然賜福給他看顧以實瑪利.

$^{361}$以實瑪利也是行國禮的.
跟他爸爸亞伯拉罕一起行國禮的.
不過到了以實出世之後.
撒拉開始接受不到.
以實瑪利跟他的兒子一起成長的時候.
想他們走.
於是亞伯拉罕不是馬上答應.
他說好啊.
我明天趕他走.
或者現在來.
不是.
這個也是他的兒子.
他疼愛的.
他真的陪著他成長了十三年.
但是.
十四年.
但是神就告訴他.
你不要跟撒拉有矛盾.
她是你的太太.
她的兒子才是你的繼承人.
以實瑪利你讓她出去吧.
我會看著她.
第十三節.
第十二節他說.
你不必為這個孩子.
和你的侍女憂愁.
撒拉對你說的話.
你都要聽從.
因為從以撒生的.
才要稱為你的後裔.
然後.
亞伯拉罕就把下甲和兒子.
讓他出去.
不過神有另一個應許.
第十三節他說.
至於侍女的兒子.
我也必使他成為一國.
因為他是你的後裔.
亞伯拉罕真是多國之父.
以實瑪利都成為大國.

$^{401}$他有十二個兒子.
每個兒子都有一族.
都有很多子女.
神有沒有歧視以實瑪利呢.
沒有.
他都是賜福給他的.
不過要做一樣東西.
就是讓以撒利順利地繼承了.
亞伯拉罕的那個應許.
亞伯拉罕他是憑著信心.
開始進入這個事件.
這個迦蘭地.
這個他不知道是什麼地方的城.
他的信心很堅定.
無論是他的侄子說要分家.
或者他的侄子被人擄走.
他經過打仗救回侄子.
或者一直都沒有子女.
他的信心都很堅定.
他很持續地堅定相信.
就算過了十年.
他在迦蘭地的角色身份都沒有改變.
他都繼續地相信.
他仍然沒有子女嗎.
他一樣是依靠這個神.
相信這個神.
他也相信神照顧他的.
也照顧以實瑪利的.
所以可以將兒子以實瑪利交託給神.
我們不可以只是用「無良心」這個字來形容他.
將兒子送出去.
可能將兒子送出去.
他才有好的發展.
不會在家族當中.
被視為「樹子」而受到影響.
還有一件事.
神是沒有偏心的.
對以實瑪利說.
亞伯拉罕 你讓他出去吧.
我看著他.

$^{441}$嘩 這個真是一個考驗.
對爸爸又是一個考驗.
對媽媽又是一個考驗.
對兒子又是一個考驗.
接著就跟亞伯拉罕說.
你兒子不是以實瑪利.
是你很疼愛那個獨子.
我說你的獨子 以實.
你帶他去.
帶他去摩利亞地.
在那裡我告訴你.
上哪座山.
上那座山就將他獻祭.
獻祭成甚麼祭.
獻祭成凡祭.
即是整個燒給我.
整個燒給我.
那亞伯拉罕怎麼做.
求求他吧.
不要吧 天父.
你講了那麼多.
這幾十年我都沒有撕過一句筆.
這次不要了 先想想吧.
想想可能會改變主意的.
不是 他帶兒子去.
他真的要奉獻兒子.
當兒子發現不妥.
全條路工人都被他放在山腳.
爸爸要跟我上山.
我就背著要來獻祭的東西.
那些柴木等等.
我已經背著tong 羊的刀.
那隻羊在哪裡 爸爸.
他問爸爸.
爸爸 那隻羊在哪裡.
我們找甚麼獻祭.
其實他不是問這個字.
他就說 爸爸.
你想做甚麼.
我和你 你想怎樣.

$^{481}$是不是.
年輕人很聰明 是不是.
嗯 爸爸說甚麼.
爸爸說 乖兒子.
神一定自己預備了他的羔羊.
這個兒子真的很乖.
他真的和他上山.
真的被他綁住.
那時候亞伯拉罕.
畢竟都一百十多歲了.
就算還沒有一百二十.
兒子十多歲的時候.
他綁住亞伯拉罕.
比亞伯拉罕快綁住他.
這個兒子很聽話.
真的被他綁住要獻祭.
但是神就真的為他預備了羊羔.
他真的獻了一隻羊羔.
亞伯拉罕.
他的信心轉變了.
由他擔心以實馬利將來怎樣.
要和撒拉在那裡.
大家有些疑惑.
不要讓他走.
神出聲了就讓他走.
還有一個階段就是.
他擔心自己生命.
但是當他擔心自己生命.
去到某一個地方.
他就要認妻作妹.
如果不是.
人家為了太太你漂亮.
他就會殺了我.
所以太太.
有人去到甚麼地方.
人家問我你是我誰.
我說你是我妹.
啊.
但是到現在說他的兒子要奉獻的時候.
他的信心一直成長.

$^{521}$他不是不看重這個兒子.
去到一百歲.
只有一個兒子你說多珍貴.
老來得子很疼愛的.
而且只有一顆在他身邊.
是最疼愛的時候.
神說獻祭.
他是相信神可以給他一個兒子.
他未必知道過程是怎樣.
但是他相信神會給他一個兒子.
這個是他的信心.
他的信心是否基礎於空白呢.
是他一直以來生命的轉變.
他信了幾十年.
好像信這件事畫不出來.
寫好像很簡單.
但是要用信去生活.
原來是很不容易的.
不過他很堅定相信.
神在他心目中的地位和能力.
他一直在轉變.
他對神的認識令他覺得.
不會影響到的.
獻祭不會影響到神的應許.
將以殺留給他做他的繼承人.
這個是神在他生命當中的應許.
神一定完成的.
他的信心一直轉變.
有一個很好的開始.
讓他好像甚麼都不知道.
進入一個不知道的四方.
但是很堅定地去面對很多環境的轉變.
很持續地去相信.
神的應許仍然不會改變.
甚至將他的兒子交託給神去照顧.
他現在信心成長到一個地步.
他相信神能夠使人復活.
在新約時期.
見過耶穌基督行神蹟就說信而已.
這個時候.

$^{561}$閻爾利亞行神蹟叫人復活都沒見過.
但是亞伯拉罕就是信.
你能不能夠信一些看不到.
不知道的東西.
見證好像由這樣而來.
人們說見證就好像是一些.
我們不知道 猜不到的事.
他祈禱就有.
都不知道是不是.
真的不知道是不是.
我媽媽癌症的時候.
第一個醫生一說.
我媽媽以後就不看醫生.
醫生就這樣.
給了一份報告我們.
接著他繼續跟人聊天.
不僅很不尊重.
我們怎樣都要等候問他.
要怎樣治療 有沒有藥等等.
他才說.
癌症嘛.
不用了.
回去他喜歡甚麼你就給他甚麼.
等待日子吧.
哇 那個病人只是在你旁邊.
我們母女完全嘔吐了.
不知道怎樣去面對.
第一次收到一個消息.
媽媽肚子痛原來是癌症.
接著他哄了很久.
我們兩個都覺得好像是絕望了.
末日了.
但哄了幾天.
哄了幾天的原因是.
我的同學打給我.
知道這件事就介紹醫生.
叫他一定要醫.
其實過程中有一次.
他肯進醫院了.
醫生就有一種藥.

$^{601}$他說這藥不是給這個症的.
這藥還沒有賣.
這藥是試藥.
你們肯不肯試一下.
試一下有兩個結果.
一個結果就是沒有了.
他現在都是保養沒有了.
另一個結果就是.
可以了.
可以的意思就是說.
可以醫.
可以有時間.
可以再等藥物的發明.
可以再調理一下.
可以令到他好一點.
你還有一段日子.
我們都祈禱過之後才回覆他.
最後我媽媽就經過治療.
第一次要打第一口針.
她要連打一個星期的針.
六日五夜.
第一晚要打針.
晚上八點鐘要打針的時候.
教會一知道.
其實就一直為我們禱告.
星期六的八點鐘.
好像是星期六還是星期日.
八點鐘要打針.
傳道人.
那位傳道人都是建道的校友.
建道的畢業生.
傳道人.
教會的團友.
個個都來醫院.
六點陪到他八點.
和他一起祈禱.
直到開始打針.
他們就離開.
然後就是.
個個鐘頭護士都進來看.

$^{641}$因為怕他馬上當晚有事.
醫生說第一針的反應.
就看他接不接受.
他一接受不到就出事了.
然後他平安過了第一天.
第二天 第三天.
都沒有什麼事發生.
第四天 第五天.
到第六天的早上.
就預計那天的十二點打完針.
就會拔掉那個針管.
那個靜脈輸入的針管.
在這個過程裡面.
第六天的早上的時候.
醫生護士來了一批.
因為學生都進來學習.
這個是試驗性質.
於是乎進來的時候.
十幾個人看完之後.
那個表情令我很害怕.
大家都很害怕.
又叫得很大聲.
我又不知道他們說什麼.
因為沒有什麼內容.
只說找醫生 找醫生.
找了醫生來.
找了醫生來了之後.
就再找另一批人來.
最後就過了一大輪時候.
我都想著要打輸數.
然後那些醫生護士最後.
就全部都在那裡 沒有走.
全部都在那裡就說了一句.
看著我說了一句.
Miracles.
神蹟啊.
真是神蹟啊.
Praise the Lord.
讚美主.
他說我們很奇怪.

$^{681}$就是為什麼.
我們摸不到那顆腫瘤.
他就說打了第一次針.
我們找不到那顆腫瘤.
找到最後最後.
是主診醫生來找.
才找到一個很小的.
他說已經是.
那藥對他很有幫助.
有沒有可能的事.
我們祈禱都不敢祈禱.
有什麼事.
當時對這個病又陌生.
這藥還要是.
剛剛開始.
試驗而已.
找白老鼠.
沒有發展的時候.
就要找媽媽做白老鼠.
打了輸數.
不過希望有一線生機.
有沒有祈禱.
有祈禱.
很多人為他祈禱.
但是.
就是在他原來.
有機會得的過程裡面.
我們是打了輸數.
這麼多人進來看.
這次不妨好事.
但原來就是這樣.
神的恩典就到.
你祈禱的事.
可能你不知道過程是怎樣.
你祈禱的事.
可能你都覺得沒有可能.
亞伯拉罕都覺得.
兒子可能沒有可能生存.
不過是會生存的.
是的.

$^{721}$信心是這樣成長.
當你去到覺得最絕望的時候.
最能夠令你成長.
你會不會願意.
好像他這樣.
有這個信心.
他不單止信.
他的兒子會存在.
他繼續在這個地方生活.
亞伯拉罕期望著.
這個地方.
將來會是一個很不同的城.
是他不知道的城.
你有沒有你盼望之城.
你盼望之城的根基是怎樣.
你能不能夠以你的信心.
來等候這一座城.
這一座城的根基.
是主耶穌基督為我們立了的.
整條道路會有很多.
我們不明白的.
不一定覺得這麼輝煌美麗的.
不一定處處覺得是花香的.
但是.
到後來你回頭看.
真是花香處處.
恩典滿滿.
我們學一件事.
學一下亞伯拉罕的信心.
真的把時間放在看聖經和祈禱.
亞伯拉罕的信心是由神的呼召那裡來.
唯有你把時間.
放在看聖經和祈禱的時候.
你能夠更加得到.
明白到.
神對你有些什麼應許.
有些什麼是真的幫助到你.
也是你需要特別去守望的.
特別去期待的.
特別去仰望神的.

$^{761}$我們為了我們的生命.
為了我們的未來.
我們都要仰望著基督.
和等候這個我們盼望之城.
你的盼望之城還想有什麼.
除了自己能夠進入這個盼望之城.
你的家人呢.
你的子女呢.
亞伯拉罕真的為他的兩個兒子.
其實不是 還有六個.
他太太過世之後.
她續言還有六個兒子.
她為她每一個兒子都有預備.
你為兒女都辛苦了很多.
最重要的是要更加辛苦.
帶領他們進入到這座城.
這個是最重要的.
讓他們得著永生.
在永恆當中.
可以一切 可以分享.
我媽媽都進入了這座城.
她走了很久.
她臨走之前.
我知道她很掛心子女的事.
我問了她一句話.
你放心走嗎.
她說我放心.
我說你騙我的嗎.
她說不是.
天父會看著你們.
我知道她真的預備好了.
她會走.
也都神建立了她.
讓她的信心真的很堅定.
她知道神會幫她看著子女.
你有什麼你要交給神看著的.
一定有的.
無論是父母或者子女.
都可能有一些是你要.
你要交給神去看著的人.

$^{801}$多點為他們禱告.
有沒有些事你要交給神看著.
你的事奉.
你的事奉不是你自己的工作.
是為神的工作.
我們都為著我們的事奉.
多去禱告 多去求吻.
求主幫助我們.
我們一起祈禱.
參見天父上帝.
我們來到你的面前.
我們求你與我們同在.
紀念我們.
主人我們真的很盼望這一座.
這一座是你預備給我們的城.
我們能夠在當中.
能夠與我們親愛的人一起在當中.
我們能夠在當中.
得到你呼喚我們.
得到你賞賜我們.
我們更加亦都很求你.
在我們等待這座城之前.
我們希望我們的生命被你建立.
我們手上的事奉被你建立.
我們又能夠以生命影響生命.
讓我們的親人.
一個一個能夠得到進入這座城的恩典.
求你看顧紀念.
最後我們都求你幫助我們.
求你幫助我們每一位坐在這裡的.
叫我們能夠.
為著你要賜給我們的這一座城.
我們去努力學習聖經.
努力多禱告.
叫我們能夠認識你多一點.
更加明白你的愛.
更加能夠將你的愛.
表現出來.
分享給人.
讓其他人都能夠得到.

$^{841}$我們祈禱仰望你.
奉主耶穌基督的聖名.
\newpage



\section{使徒行傳 26:1-23}
\label{sec:Rlc0zB_br4Y}
\textbf{2024-01-14 早堂《沒有違背那從天上來的異象》使徒行傳 26\_1-23 講員:李富成牧師}
\newline
\newline
連結: \href{https://youtube.com/watch?v=Rlc0zB-br4Y}{\texttt{https://youtube.com/watch?v=Rlc0zB-br4Y}} ~~~~ 語音日期: 2024-01-16
\newline
\newline
\hyperref[sec:2um86Ip2zrs]{\small{< < < PREV SERMON < < <}}
~
\hyperref[sec:index_chronic]{\small{[返順時目]}}
~
\hyperref[sec:index_scriptual]{\small{[返順卷目]}}
~
\hyperref[sec:dR8L9GrF4hQ]{\small{> > > NEXT SERMON > > >}}
\newline
\newline
使徒行傳 26:1-23
\newline
\begin{longtable}{cl}
\hline
\hline
章節 & 經文 (和合本修訂版)\\
\hline
26:1 & \begin{tabularx}{0.7\textwidth}{X} 亞基帕對保羅說：「准你為自己申訴。」於是保羅伸手辯護說： \end{tabularx} \\ \\ \relax
26:2 & \begin{tabularx}{0.7\textwidth}{X} 「亞基帕王啊，猶太人所控告我的一切事，今日得以在你面前辯護，實為萬幸。 \end{tabularx} \\ \\ \relax
26:3 & \begin{tabularx}{0.7\textwidth}{X} 更慶幸的是你熟悉猶太人的規矩和他們的爭論；所以，求你耐心聽我。 \end{tabularx} \\ \\ \relax
26:4 & \begin{tabularx}{0.7\textwidth}{X} 「我自幼為人如何，從起初在本國的同胞中，以及在耶路撒冷，所有的猶太人都知道。 \end{tabularx} \\ \\ \relax
26:5 & \begin{tabularx}{0.7\textwidth}{X} 他們若肯作見證，就知道我從起初是按著我們教中最嚴緊的教門作了法利賽人。 \end{tabularx} \\ \\ \relax
26:6 & \begin{tabularx}{0.7\textwidth}{X} 現在我站在這裡受審，是為了對神向我們祖宗的應許存著盼望。 \end{tabularx} \\ \\ \relax
26:7 & \begin{tabularx}{0.7\textwidth}{X} 這應許，我們十二個支派，晝夜切切地事奉神，都指望得著。王啊，我正是因這指望被猶太人控告。 \end{tabularx} \\ \\ \relax
26:8 & \begin{tabularx}{0.7\textwidth}{X} 神使死人復活，你們為甚麼判斷為不可信呢？ \end{tabularx} \\ \\ \relax
26:9 & \begin{tabularx}{0.7\textwidth}{X} 「從前我自己認為必須竭力反對拿撒勒人耶穌的名， \end{tabularx} \\ \\ \relax
26:10 & \begin{tabularx}{0.7\textwidth}{X} 我在耶路撒冷也曾這樣做過；我不但從祭司長得了權柄，把許多聖徒收在監裡，而且他們被殺，我也表示贊成。 \end{tabularx} \\ \\ \relax
26:11 & \begin{tabularx}{0.7\textwidth}{X} 在各會堂，我屢次用刑強迫他們說褻瀆的話，我非常厭惡他們，甚至追逼他們，直到外邦的城鎮。」 \end{tabularx} \\ \\ \relax
26:12 & \begin{tabularx}{0.7\textwidth}{X} 「那時，我帶著祭司長的權柄和命令往大馬士革去。 \end{tabularx} \\ \\ \relax
26:13 & \begin{tabularx}{0.7\textwidth}{X} 王啊！我在路上，中午的時候，看見從天上有一道光，比太陽還亮，四面照射著我和跟我同行的人。 \end{tabularx} \\ \\ \relax
26:14 & \begin{tabularx}{0.7\textwidth}{X} 我們都仆倒在地，我就聽見有聲音用希伯來話對我說：『掃羅！掃羅！你為甚麼迫害我？你用腳踢刺棒是自找苦吃的！』 \end{tabularx} \\ \\ \relax
26:15 & \begin{tabularx}{0.7\textwidth}{X} 我說：『主啊，你是誰？』主說：『我就是你所迫害的耶穌。 \end{tabularx} \\ \\ \relax
26:16 & \begin{tabularx}{0.7\textwidth}{X} 起來，站著，我向你顯現的目的是要派你作僕役，為你所看見我的事，和我將要指示你的事作見證人。 \end{tabularx} \\ \\ \relax
26:17 & \begin{tabularx}{0.7\textwidth}{X} 我也要救你脫離百姓和外邦人的手。我差你到他們那裡去， \end{tabularx} \\ \\ \relax
26:18 & \begin{tabularx}{0.7\textwidth}{X} 要開他們的眼睛，使他們從黑暗中轉向光明，從撒但權下歸向神；使他們因信我而得蒙赦罪，和一切成聖的人同得基業。』」 \end{tabularx} \\ \\ \relax
26:19 & \begin{tabularx}{0.7\textwidth}{X} 「因此，亞基帕王啊！我沒有違背那從天上來的異象； \end{tabularx} \\ \\ \relax
26:20 & \begin{tabularx}{0.7\textwidth}{X} 我先在大馬士革，後在耶路撒冷和猶太全地，以及外邦，勸勉他們應當悔改歸向神，行事與悔改的心相稱。 \end{tabularx} \\ \\ \relax
26:21 & \begin{tabularx}{0.7\textwidth}{X} 為這緣故，猶太人在聖殿裡拿住我，想要殺我。 \end{tabularx} \\ \\ \relax
26:22 & \begin{tabularx}{0.7\textwidth}{X} 然而，我蒙神的幫助，直到今日還站立得穩，向尊貴的和卑微的作見證。我所講的，並不外乎眾先知和摩西所說將來必成的事， \end{tabularx} \\ \\ \relax
26:23 & \begin{tabularx}{0.7\textwidth}{X} 就是基督必須受害，並且首先從死人中復活，把亮光傳給猶太人和外邦人。」 \end{tabularx} \\ \\
[1ex]
\hline
\hline
\end{longtable}
$^{1}$各位靈姐妹平安.
讓我首次在大家當中.
來分享上帝的話.
使徒保羅在面對殉道的關頭.
他在回顧為主奔跑一生的時候.
他寫信給提摩太.
他講了一段豪情壯語.
可能我們都很熟悉.
我們一起讀.
那美好的帳 我已經打過了.
當跑的路 我已經跑盡了.
所信的道 我已經守住了.
從此以後 有公義的冠冕為我存留.
就是按著公義審判的主.
到了那日要自給.
使徒保羅面對殉道的關頭.
他竟然講了這段豪情壯語.
為何保羅如此肯定.
甚至我們聽到.
他有沒有大言不慚.
有沒有說大了.
他說主一定會賜他公義的冠冕.
在《使徒行傳》第26章1至23節.
為我們提供了一個很清楚的答案.
原來保羅如此肯定.
他能獲得所與所賜的冠冕.
是因為他沒有違背.
那個從天上來的意象.
並且他終其一生.
他是以信服這個意象為他的職志.
可以說是赴湯蹈火 萬死不辭.
正如保羅在與以弗所的長老離別的時候.
他有一句這樣的自白.
我卻不以性命為念.
也不看為寶貴.
只要行完我的路程.
成就我從主耶穌所領受的職事.
證明神恩惠的福音.
今天早上讓我們透過.
保羅向阿基泊王的這段自辯當中.

$^{41}$去了解保羅沒有違背.
那個從天上來的意象的原因.
他帶來什麼結果.
以及他要為此付上什麼代價.
在為頂子妹我們今天.
去信服上帝的旨意.
去實踐得人使命的借鑒和激勵.
第一方面我們要思想的是.
保羅信服意象的原因是什麼呢.
在經文4至16節讓我們看到.
是由於他經歷了基督的救贖恩情.
他生命帶來了一個轉化.
大家有沒有留意剛才這個主題.
保羅向阿基泊王說.
固此沒有違背那從天上來的意象.
大家不要小看固此這兩個字.
這兩個字正正是說明.
保羅信服意象是有原因的.
他不是傻傻的.
不是一時衝動.
不是心血來潮.
走上了這條被人看起來很傻瓜的道路.
這個字表明保羅信服意象的原因.
是記載在第4至16節.
保羅向阿基泊王的自辯提到他原名.
我們一起溫習一下.
保羅原名叫什麼.
Solon.
他未信主之前.
他是師承經學大師伽瑪列.
可以說是出名門.
用今天的說話說.
即是名校出身.
他說他是希伯來人所生的希伯來人.
可以說是根正苗紅.
他是嚴守律法.
是法利賽教門的一個青年才俊.
可以說是前途似錦.
甚至他當時為了要消滅.
還是崛起的耶穌拿撒教派.

$^{81}$這些從他的角度看來是旁門左道.
異端邪說的.
他是大發熱心地不擇手段去迫害基督徒.
這個和耶穌是完全敵對的猶太狂熱分子.
他在前往大馬士革去捉拿基督徒的途中遇見主.
復活的主親自向他的顯現.
這個突如其來的經歷.
令到保羅永遠都不會磨滅.
在《士徒行傳》第22章.
當保羅向那群群情洶湧的猶太人.
致辯的時候.
亦都曾經提過這一段被主光照的經歷.
在第22章第10節.
保羅益述被主光照的回應是什麼呢?.
他說:主啊!我當作什麼?.
這個保羅對本來他要打壓的耶穌.
稱他為主.
對一個猶太人來說是不得了的事.
這個稱呼有包含耶穌就是猶太人所引頸以待的彌賽.
救主的意思.
保羅遇見主.
蒙恩得救.
生命轉化.
他本來人生的任務.
就是要打壓耶穌.
現在卻承認耶穌是救主.
還為了傳揚耶穌而賣命.
正如使徒保羅在《肥勒比書》第3章7至8節那裡這樣說.
他說:只是我先前以為與我有益的.
現在我因基督的事而懸顧我體為有損的.
不但如此.
我也將萬事當作有損.
因我已認識我主救主耶穌為至寶.
我為他已經掉起萬事當作糞土.
為要得著基督.
保羅亦在《哥林多前書》第15章9至10節.
很感恩地這樣說.
他說:我原本是使徒裡最小的.
我不配稱為使徒.
因為我從前是迫不得上帝的教會.

$^{121}$然而我今日成了何等人.
是蒙神的恩才成的.
為何保羅能夠信服二將?.
是因為他深深體會經歷耶穌基督的救贖恩情.
所以他的生命起了180度的變化.
可以這樣說.
耶穌的愛改編了保羅人生的樂章.
我首次和大家認識.
你猜我以前是從事哪一行的?.
我的姊妹說教書.
有樣子看的嗎?教書.
我在2013年選擇提早退休.
進入全時間的事奉行列.
在我提早退休之前.
我是從事了29年的小學教育工作.
我頭8年在宣道會陳元喜小學擔任訓導主任.
像不像訓導主任的樣子?.
在大家心目中訓導主任的樣子是凶神惡煞.
其實那些已經舊了.
新的訓導主任是不用凶神惡煞.
只是融化他.
在8年之後.
我去了仁濟醫院蔡英桃小學.
開始一間新的小學.
我做了21年的小學校長.
在這29年的小學教育工作裡面.
很多同學,老師,家長,校友.
都對我的工作有些詫異.
為何訓導主任這個校長.
特別疼愛學習能力薄弱的學生呢?.
大家聽不聽得懂甚麼是學習能力薄弱?.
這個名字是不是很漂亮?.
即是我讀小學的時候已經被老師趕出校.
沒有用,站在門口.
即是一班SEN的學生.
或者是一些專導彈的學生.
為何我會有這種反應呢?.
原因很簡單.
因為我是過來人.
我在小學階段的成績和操行.

$^{161}$比起我過去那29年見過的所謂很頑皮的學生.
他們和我還差很遠.
所以我對著那些頑皮的學生.
我一面要責備他們.
其實我心裡在笑.
你和我的東西還差很遠.
所以我很有容量去接納和體恤他們.
我成長在一條很偉大的屋村.
知道這條甚麼村嗎?.
這條村黑白相的時候是這樣的.
這條村是最近才宣布將要重建.
這條叫彩虹村.
這就是我小時候唯一找到的一張照片.
這張照片是我小六升中的時候.
還在考升中試.
在座有誰考過升中試的?.
舉手.
全部暴露了年紀.
因為學能測驗是1979年才開始的.
以前的升中試是考中英數的.
考科的成績就是你的.
沒有甚麼拉曲線的.
所以當時我們放榜要看甚麼.
買星島日報.
看自己如何在未上學時排位正前.
我出生在一個草根階層.
爸爸媽媽都要上班.
我上下午班的.
所以最少都是拉球底.
大家知不知道甚麼是球底呢?.
球底即是現在叫球場.
以前叫球底.
球底和球場有很大分別.
球底就是英雄地.
甚麼是英雄呢?.
你看過甚麼古惑仔電影講的橋段就是這樣.
因為整個屋村包括我們附近的大磡村.
牛池灣一帶的小朋友都會走來.
我們草形村唯一一個大的平地叫球底那裡爭場踢步.
所以如何分配場地呢?.

$^{201}$現在就是上網抽籤.
以前只靠兩樣東西.
就是拳頭和髒話.
當你要用暴力去分配場地.
自然有人來幫你.
幫你的人就要跟隨他.
就形成了同黨.
或者再進一步成為了社團的接班人.
其實我最少就在這個球底留年.
所以在球底裡面.
我發現的各種的壞行為.
例如講髒話,店舖盜竊,看色情書.
都成為了我的生活的一部分.
大家可以想一想.
我這個不喜歡讀書又留年街頭的小混混.
我繼續長大的時候.
我的下一步會是怎樣,是嗎?.
不過神很大的恩典.
在我小五的時候遇到這位女士.
這位女士就改變了我的生命.
你猜一下她是我哪位?.
問十次都答十一次是我太太.
她是我小五的代課老師.
我今年62歲.
你猜她幾多歲?.
她72歲.
這張相是我最近她回來香港探望媽媽.
她媽媽90高齡.
難得我終於在幾十年後.
請她吃一餐謝師宴的時候拍下.
她對我生命的影響是什麼呢?.
當時她來代課.
我原任的班主任就要進廠修理.
當時我們是不敢在我們原任老師裡面.
在課堂有什麼導彈.
因為當時可以看佛.
逢是可以看佛我們都會害怕.
但是代課老師是不敢看佛的.
所以代課老師來代課.
經常成為我們這些頑皮的出氣袋.

$^{241}$專門作弄她.
因為她不敢看佛.
但是這個老師你想像在50年前.
72歲都這樣的面貌.
你想像50年前是什麼樣子.
簡直是美少女般走進來.
我們全部同學都受驚一驚.
這個代課老師這麼漂亮.
我當時是小頭目.
我就說大家不要作弄她.
今天看她怎樣.
她教得很好.
很有心機教.
上課會說李富城你回答吧.
我就暈了.
到她要離開了.
她就跟我們說.
她說我是基督徒.
我們第一次聽.
她說你那一個小朋友想繼續見到我.
就跟我回教會吧.
誰想就舉手.
全部人都舉手.
隔壁班聽到都舉手.
我們說很難.
我們當時住在彩虹村.
我母會當時還沒搬去何文田.
還在深水Po .
怎樣去.
1972年.
她說不要緊.
如果你們夠人多肯去.
我可以叫教會安排一部旅遊車.
每個星期日來彩虹村接你們回教會.
然後再跟你媽媽說.
有老師帶.
即是她.
因為原來她也是住在彩虹村.
有老師帶 有旅遊車坐.
相信你家長可能會給.

$^{281}$於是我回去就跟我媽媽說.
我媽媽第一個反應就是.
臭小子 你懶惰不如跟你老師去教會吧.
就是這樣.
我的生命開始有變化.
她是我代課老師 叫劉月琴.
她現在在美國的.
一間美男浸訓會的中文部.
事奉了主35年.
她繼續事奉中.
我就開始搭足我的母會.
黑白相就是那間叫舊堂.
我回到兩個月之後.
哥哥姐姐說.
富仔我們搬吧.
我說搬去哪裡呢.
搬去何文田.
搬來這間在何文田的新堂.
不過新了50年.
就是因為這個緣故.
我能夠有機會接觸教會.
認識信仰.
我還記得那年.
當我升中那年有個叫夏令會.
大家有沒有去夏令會.
現在叫三日兩夜夏令會.
都很少人去.
因為覺得很長.
以前的夏令會五日四夜.
我們是無名要去.
因為可以不用在家裡睡覺.
去夏令會.
去夏令會的地方.
現在暫時封鎖了.
叫宣道園.
宣道園旁邊有一塊大草地.
我們踢球的人很喜歡.
聚會長到不得了.
大人小孩一起聽.
我們小五小六的小朋友就睡著了.

$^{321}$但我們懂得看紙.
我們知道時間表有自由時間.
我們熬完那個講道.
就出去踢球.
誰知我拿起球準備要走.
大哥哥又叫住我.
富仔 跟你聊天.
聊什麼呢.
大家想不想到跟我聊什麼.
那些叫個人報道.
第一句就跟我說.
現在我們去報道.
千萬不要第一句這樣跟別人說.
因為那個人就不會再跟你說.
他第一句跟我說.
你是不是罪人.
因為他知道我們的背景.
因為我們在教會還在說髒話.
在夏令會還被大哥哥罰站.
在籃球署旁邊.
所以他不怕一刀直入.
就說你是不是罪人.
我說是.
跟住很單純的信心去接受.
當然整個屬靈生命的成長.
和教會的栽培.
是物不可分.
就是這樣緣故.
這個街頭小混混.
就好像保羅在歌羅西書三章九至十節.
這樣說.
我們一起讀.
因你們已經脫去舊人和舊人的行為.
穿上了新人.
這新人在知識上漸漸更新.
正如做他主的形象.
我們由地位的更新.
去到屬靈生命的成長.
是一個過程.
我這個街童因為得著耶穌的愛.

$^{361}$改編了我人生上半場的樂章.
我學習去.
作為英才.
春風化雨.
同時間我也在我舞會有十分投入的事奉參與.
建立生命.
亦都是由於耶穌基督的深深的愛.
他繼續來改編我人生下半場的樂章.
他呼召我去全時間去服侍他.
用我人生仍有的日子.
去全時間去傳揚真道.
頂主妹.
你和我的生命故事都不同.
我們每一次浸禮.
可能都有機會聽一些頂主妹的見證.
她是怎樣.
切入這個信仰.
我們信仰的切入點都不同.
經歷都不同.
有些人因為子女信得很好.
就帶爸爸媽媽.
爸爸媽媽給點面子.
聽聽聽.
上帝的靈感動他們.
有些人在病患裡.
祈禱得醫治啟動他們的信心.
有些人是老師帶.
有些同學帶.
有些同事帶.
有些自己摸上門.
各式各樣都有.
不過我相信你和我.
都會認同保羅剛才的宣認.
然而我今天成了何等人.
是蒙主的恩才成的.
頂主妹讓我們去審查.
基督的救贖恩情.
是否繼續激勵我們.
願意信服上帝旨意.
實踐得人使命的原動力.

$^{401}$和唯一的原因.
我想強調的是今天.
我在教會生活了五十年.
我們見過很多弟兄姊妹.
曾經很熱心.
她的信仰停留在某一個階段.
這個持續的經歷神的身因.
就是我們每次崇拜的重點.
如果不是我們每個星期崇拜.
做什麼呢.
我們每一次的認信.
每一次的詩歌.
每一次講信.
都提醒我們耶穌基督.
這份深深的愛.
今天仍然要去改編你人生的樂章.
第二方面我們來思想.
在這段經文中.
我們很清楚看到的.
是信服異象帶來什麼結果呢.
對保羅來說.
就是他終生傳揚福音真道.
經文讓我們明白.
保羅所經歷的異象.
除了有蒙主拯救這個部分.
還有被主差派的部分.
在剛才我們讀的第十六節.
主說我要派保羅作福音的執事.
將生命轉化的經歷.
和主所指示的訊息證明出來.
在第十八節.
主又說我要差保羅.
去到猶太人和外邦人當中.
傳揚恩訊得救的福音.
請留意保羅所領受的異象.
不是單單他個人恩訊.
稱義得永生的那個確決.
還包含被主指示.
指派和差遣成為福音使者.
藉著保羅的見證宣講.

$^{441}$引領很多人從撒旦的權勢下歸向神.
以至能夠和眾聖徒得蒙救恩.
保羅領受被主差派傳福音的異象.
是很清楚的.
所以我們讀耳窟所書的時候.
讀到第三章.
保羅就描述自己做了福音的執事.
事實上就是由於保羅為了傳揚基督.
他一生勞苦盡心竭力.
用盡各種智慧.
勸誡人,教導人.
將各人在基督裡面.
完完全全的引導神的面前.
因此我們見到.
福音就隨著保羅的宣教團隊的足跡而廣傳.
不同城市的教會就一個一個地建立起來.
為了這二千年來教會的歷史.
揭開了一個輝煌的序幕.
可以說波瀾壯闊的宣教運動.
也都隨之而展開.
不知道大家有沒有看過.
或者聽過這位華人救藥聖經學者.
前中國神學研究院的院長.
李思敬牧師.
李思敬博士的著作.
或者聽過他的講道.
大家知不知道李牧師的爸爸.
就是李飛吾牧師.
李飛吾牧師是長期委身牧羊宣道會.
希伯倫堂的一位前輩.
原來李氏父子兩人蒙恩得救.
獻身侍主廣傳福音.
造就了無數的信徒.
其實他們是受一個人影響的.
哪個人這麼厲害.
這個人的名字叫做何二師教士.
他在華人當中有一個匿稱的.
他叫做何二姑.
如果現在叫你做阿姑.
你就想描他.

$^{481}$因為當他在華人地方服侍的時候.
姑只是代表小姐,女士的意思.
話說在20世紀初.
美國有一個基督徒的家庭.
他們很喜歡接待宣教士.
宣教士回來述職.
他們就請他們來家裡吃飯,聊天.
有一位在中國述職.
回來美國述職的何牧師.
來到一個家庭作客.
當時一個年僅十歲的小女孩.
一起在這個作客當中.
何牧師分享了在宣教工場.
在中國的一個小故事.
而這個故事令到這個小女孩不生難亡.
何牧師分享說有一天.
有一個村民推著一輛木頭車.
來到何牧師所處的報道所.
原來這個村民聽到說.
有些宣教士,有些鬼佬.
來講什麼耶穌.
就派他來.
邀請一些傳教士到村裡跟他們說.
究竟是什麼故事.
如果你是宣教士.
你覺得這是不是一個大好的機會.
他們主動開門想聽.
不過很奇怪.
何牧師說那次不知道為什麼.
報道所個個都去不到.
於是就說這個失望的村民.
就推著一輛木頭車.
慢慢在這個黑夜裡消失了.
聽到這裡我們都可能會說.
哦 是不是.
但這個小女孩當時聽著這個故事.
來到這裡突然站起來.
就說如果是我.
我一定去.
大家大人看一看她.

$^{521}$小孩坐下來.
這個木頭車的異象.
當時放在了這個小女孩的心裡.
不過隨著她成長.
她都忘記了.
到她入讀大學.
在史丹福大學讀商學.
她準備要做成功的商人.
在她升大二的時候.
不知道為什麼媽媽寫封信給她的時候.
裡面提到中國.
這個木頭車的異象.
又在這個小女孩.
或者叫小女子的心裡浮現出來.
經過一番的掙扎.
這個小女孩終於信服上帝給她的異象和差派.
在1913年11月.
這個21歲的小女子.
遠攝重陽來到中國.
輾轉間去到廣東省南海縣西朝的官山傳道.
建立禮拜堂孤兒院.
而原名叫做李珍濤的李飛吾牧師.
就是因為父母傷亡.
在9歲的時候.
就被人送去這個小女子所建立的孤兒院.
以致獲得照顧.
有機會聽民福音.
也都像這個小女子一樣.
將生命奉獻給主.
成為華人教會裡面的一代牧師.
這個就是希伯倫會的創辦人.
何以思教士.
何以思教士21歲踏足中國.
直至到1971年.
他78歲正式退休.
他在我們華人當中服侍了有60年.
因為他信服木頭車的意象.
他拋下了加州的陽光.
他捨棄了他高校的學位.
他埋葬了他從商的壯志.

$^{561}$他離別了他溫馨的家庭.
他千里迢迢的跑到一個陌生的國度.
他曾經經歷過中日戰爭.
太平洋戰爭.
國共內戰.
種種的困境和危險.
他始至終生傳揚福音真道.
因為他的信服.
造就了李氏父子.
李思敬牧師在他的自傳裡的序言.
他說如果沒有教士婆婆.
不會有我.
弟兄姊妹.
有一首青年聖歌的詩歌叫做萬世戰爭.
在這首詩歌的副歌.
一開始問了我們四個問題.
他說是否你見此異象.
是否你心激動.
是否你認住呼召.
是否你願停從.
真的唱得到.
會不會唱?.
試試跟我唱.
是否你見此異象.
是否你心激動.
是否你認住呼召.
是否你願停從.
要停從就要付代價.
信服異象帶來的代價.
對保羅我們看到的是他.
獻奉所有.
擺上他的一生.
保羅要付的代價是什麼呢?.
原來當他由一個猶太的陣營.
轉了去一個傳揚耶穌的陣營之後.
他當時所面對的迫害.
並不是來自當時的羅馬政權.
而是來自猶太的狂熱份子.
他觸怒了那班猶太人.
為此保羅要隨時付上被捉拿.

$^{601}$隨時會殉道的代價.
我們今天讀的這段經文的前文.
由第21章到第27節開始.
一路到26章.
作者路加就寄述到猶太的狂熱份子.
去慫恿那些耶路撒冷的猶太人.
他們誣衊保羅.
說保羅是踐踏猶太人.
他藐視舊約律法.
他污衊聖殿.
所以引致合城的猶太人都震動.
要一齊出來捉住保羅.
要拉他出聖殿.
想要殺他.
不過由於這個動亂.
引起了駐守耶路撒冷的羅馬千夫長的神經.
不能在他治下突然之間亂起來.
用我們今天說話說.
就是要維穩.
所以他唯一的方法.
就是將保羅從這群瘋狂的群眾裡扣押.
然後就看看如何處理.
這個扣押變相保護了保羅.
但是這群猶太狂熱份子.
並沒有因此而停止他們的攻擊.
按照《士大恆傳》第23章12-23節.
記載有四十幾個殺手.
雖然聖經沒有用殺手這個字.
但是裡面記載他們起誓.
如果我們不殺保羅.
我們就不吃不喝.
那不是保羅死就我死了.
還不是殺手.
這群人要除掉保羅.
是要立誓一定要將他剷除.
保羅的生命危在旦夕.
幸而這個消息不夠保密外洩.
因此這個駐防的羅馬軍官.
經過風險評估之後.
就派了一隊龐大的保安隊.

$^{641}$將保羅押送去一個安全的地方.
就是剛才我們讀的經文.
保羅受審的巡撫.
即是總督的官邸.
我們從那個龐大的保安隊.
我們就可以評估保羅的風險有多大.
我們看看保安隊.
千夫長便叫了兩個百夫長來說.
預備步兵二百馬兵七十.
長槍手二百.
即是多少人.
深算.
派了四百七十個武裝人員.
荷槍實彈.
要去押送保羅去一個更安全.
遠離被殺的地方.
因為保羅有羅馬的公民的身份.
所以如果他的命在千夫長手下無端端死了.
那個千夫長要問責.
所以他不能夠掉以輕心.
我們從這個保安隊的龐大的程度.
我們可以推想.
這個千夫長是一個職業軍官.
他不是傻的.
他一定評估了情報.
那些要來殺他的人.
有多大的動力.
他才出多少的押送隊.
保羅的命仔.
命是一線.
除了這些外部的攻擊之外.
保羅自己也曾經在哥林多後書第十一章.
去開列了一張.
他因為服事主傳陽福音的一張受苦的清單.
今天我們不會跟你讀.
你有興趣回去看一下這張清單.
有教會外的事.
有教會內的事.
外憂內患.
保羅要在這裡開列這張清單.

$^{681}$並不是用來演戲.
去告訴別人.
你猜我為主擺上多少那種.
其實他是很老實.
想跟哥林多教會的頂子妹說.
原來我們此至要服事主的人.
你要有一個準備.
你是要付代價.
剛才我分享.
我在2013年提早退休.
就全時間去服事主.
這個轉變對一個中年人來說.
其實也是一個頗大的轉變.
如果我要付一個全時間服事主的航道的轉變.
我要付什麼代價.
我第一個要付的代價很簡單.
就是辭職.
當時其實我領受了這個意象.
我本來打算在我開辦的這間學校.
做到20週年.
就是一個很好的階段.
交給大團體請另外一個校長接任.
20週年了 夠了.
當時校董會就說.
不要 怎樣也要做多一年.
因為當時我們是剛剛經歷了.
這個縮班殺校的漩渦.
而剛剛站穩 開齊班.
校董會就說怎樣也要做多一年.
所以就特別過多一年.
那我要放下校長的工作.
我要付什麼代價.
是一個什麼代價.
四個字就說完了.
叫高薪厚職.
校長的工作.
在一般華人社會都是備受敬重的工作.
而香港的津貼學校的校長工資是不低的.
這裡在在有沒有行家 教育界的.
如果你是常額的職位.

$^{721}$現在主任也有六七萬.
校長是過十萬.
就是年薪過百萬.
不過你又不要笑我.
李牧師你在這裡浪費命嗎.
說你幾千元工資 不是的.
我想告訴你.
我當時最大的困擾.
蒙召不是這份工作.
也不是高薪厚職.
因為我覺得金錢夠就行了.
我這套西裝都穿了五年.
我女兒結婚的時候穿到現在.
對男性來說.
一件深藍色的西裝.
什麼場合都可以去.
婚禮也可以去.
喪禮也可以去.
來港島也可以.
最大的困擾對我來說.
我漏了說.
當時激勵我最後要決定的.
就是戴紹增牧師.
不知道大家認不認識我們這位前輩.
戴紹增牧師是戴德生牧師.
第四代的孫.
戴德生牧師一家幾代.
在我們華人當中服侍.
戴紹增牧師在2009年的時候過身.
他年輕的時候.
十幾歲的時候.
因為隨著宣教的團隊.
被日本軍兵扣押在集中營.
在中日戰爭的時候.
看他的自傳他說.
當時他立志.
我不會像我的祖父和爸爸那樣傻.
去傳道.
現在做街下囚.
他說就算我要傳.

$^{761}$都不會來中國傳.
不過很奇怪.
當他被釋放之後.
他接受完訓練之後.
神又差派他去到華人的地區.
主要在台灣.
後來在東南亞.
大家如果聽過.
一些宣教的團隊.
他是一個很重要的負責人.
他晚年去到哪裡呢.
他晚年在四川.
透過服務.
進入他原生出生的地方.
他的普通話是流離過任何一個土生土長的人.
因為他在中國出生.
所以中國政府對他是.
不當他是外人.
因為他根本就是中國出生.
叫他做大博士爺爺.
當我看他生命的時候.
就激勵我還在想什麼.
踏上.
剛才我說.
高薪厚職不是最重要要考慮的代價.
是什麼呢.
原來是一個心理的因素.
是我要離開一個.
安書區.
對一個中年人來說.
要離開一個已經打拼了一段時間的.
安書的環境.
我在大埔工作了很多年.
在社區裡面有很多網絡.
我會叫書記做powerpoint.
我是不會做.
不過我會叫人做.
因為有足夠的文職人員支援我.
我有獨立的辦公室.
我已經在業界有一點點的影響力.

$^{801}$在一個這樣的階段.
要我放下一個賴以安全的環境.
特別我們在大埔喝茶一定有位置的.
因為門口的家長大底上都是你家長.
那個資客.
明明排隊說.
沒有位置了,我轉頭走.
校長,喝茶.
沒有位置,排隊.
校長當然有,這邊.
雖然這些是華人的陋習.
不過我都很喜歡.
有些排隊的說.
你那個人.
我是校長.
我真的想反映那個.
安書區的威力是很重要的.
不過,丙姐妹.
你想一想.
我剛才所說我所負的代價.
相對昔日.
無論在耶穌時代.
在使徒時代.
或者我們近代體的宣教士.
他們所擺上的是生命的代價的時候.
我就真的覺得九牛一毛.
丙姐妹.
如果我們今天都要去信服.
這個事奉的職志.
信服這個無為意象的教導.
你有沒有想過.
你都要負代價呢?.
我剛才八點半來到.
我已經看到.
聖學的人員在練師立.
他最低限度要負上.
早你們一個多小時回來練習的時間.
在我們目會的過程裡面.
我們有時候會過分地.
邀請丙姐妹事奉.

$^{841}$其實有參與事奉是我們的榮幸.
何須三催四請呢?.
當然我們做牧者要針對大家的恩賜.
要配搭.
但是那個負代價事奉的動力.
是應該我們每一個經歷過主大愛的人.
一種由衷的回應.
聽到這裡.
保羅說到這裡.
當時培同雅基柏王聽審的.
一個巡撫即是總督.
他叫做菲斯都.
他是一個外邦人.
他根本聽不明白保羅說什麼.
所以在剛才我讀的經文的第二下.
接著那節是這麼說的.
保羅這樣分數.
菲斯都大聲說.
即是插嘴.
保羅你癲狂了吧.
你的學問太大.
反叫你癲狂了.
因為對一個外邦的軍官來說.
他根本沒有辦法明白.
保羅好端端在猶太的社群裡面是明日之星.
現在要搞什麼.
耶穌搞到現在做了階下囚.
還口口聲聲說.
你如何如何信服他的旨意.
根本對這個外邦的軍官來說.
他一點都不明白.
丙姐妹.
我們有時候都癲癲的.
在微信的親友同事朋友當中.
他們會有.
有時候會覺得.
你的基督徒癲癲的.
禮拜日不睡晚一點回來祖堂.
這麼好天氣不是行山回來崇拜.
什麼十一奉獻.

$^{881}$自己還不夠用.
你多餘了.
你癲的.
你校長.
你好好的做校長都60歲就算了.
你搞什麼.
去做傳道人.
你知不知道你的公積金怎麼算的.
從不同微信的朋友.
他們會用他的價值觀來看我們.
但我也很坦白告訴大家.
如果我們沒有癲癲的行為.
跟他們來問我們問題.
那你都很有問題.
因為他都不覺得你有什麼不同.
就是因為我們的不同.
他們開頭會來笑我們.
接著他會來羨慕我們.
他羨慕你和諧的家庭.
他羨慕你在逆境裡的得力.
他羨慕你打不還口那份厚度.
他羨慕你對金錢那份不執著.
他羨慕你看透人生.
因為你有主.
他開始羨慕你豐盛的人生.
各位姐妹.
願我們都像保羅一樣的宣告.
我們就來宣告一次.
開始.
我固此沒有違背那從天上來的意象.
有一首詩歌將保羅整個人生都描繪得很清楚.
這首詩歌不容易唱.
但很好聽.
我們一起來唱完這三節的詩歌.
作為我們的回應.
大麻色 圖上寫幾多歉意.
榮耀盛光 忽照在人面前.
沒法解釋 其已刻骨的敬遠.
但卻深知 一生將要改變.
從前聽 並別把世上信仰.

$^{921}$還為自己當天所作示好.
然而人生 顯然往小路逃.
十格寫生的愛 今參上當.
自那天起 全力奔跑 崎嶇站路.
無懼關壁 不怕浪濤滔滔.
夜冷風高 仍未想過 失速卻步.
立志至上 眾主所學中路.
迢迢馬力 即使翻過暴雨路途.
茫茫夜晚 心守至處全步.
然而人生 便亦有破半路途.
讓我共進一山 將全愛回報.
為了幾多 無論經過多少事.
名利付出 一概拋卻身邊.
病痛較千 亦未改心底的信念.
夕陽打天 必得不留歡喜.
請你和喜立唱.
全力進一張福音 繼續必傳.
無問代價 至上天國萬年.
全拍十一 亂上 伴這數十年.
就算有天天出生 佢也無怨.
同心圖共.
多謝主你透過保羅向阿基帕王的自辯.
讓我們再一次看到.
他沒有違背意象的原因.
是經歷主深深的大愛.
又讓我們看到他的意象.
除了包含他得救.
確認他的信德之外.
他領受了繼續將福音傳揚.
更加讓我們看到.
他為了福音的緣故.
一生勞苦付代價.
求主讓保羅的生命.
今天再一次激勵我們.
讓我們沉睡的再一次興起.
讓我們冷淡的再一次火熱.
讓我們已經走在主的心意當中.
我們繼續力上加力.
祈禱奉耶穌基督的聖名.
阿們 各位姐妹請坐.

\newpage



\section{瑪拉基書}
\label{sec:dR8L9GrF4hQ}
\textbf{2024-01-18 周四晚查經班《瑪拉基書(2)》 講員:曾裔貴牧師}
\newline
\newline
連結: \href{https://youtube.com/watch?v=dR8L9GrF4hQ}{\texttt{https://youtube.com/watch?v=dR8L9GrF4hQ}} ~~~~ 語音日期: 2024-01-19
\newline
\newline
\hyperref[sec:Rlc0zB_br4Y]{\small{< < < PREV SERMON < < <}}
~
\hyperref[sec:index_chronic]{\small{[返順時目]}}
~
\hyperref[sec:index_scriptual]{\small{[返順卷目]}}
~
\hyperref[sec:_CjhJKqUcA8]{\small{> > > NEXT SERMON > > >}}
\newline
\newline
$^{1}$今晚我們在203這裡.
也有些桌子方便在現場的兄弟姐妹.
可以寫筆記等等.
我們一起有同心禱告.
多謝天父的恩典.
我們能夠可以行上.
逢星期四晚.
來去研讀聖經.
我們盼望我們一直累積.
對每一卷書的認識.
深化我們對神的真理啟示.
其實就是對神的屬性的認識.
以致我們在靈命上.
因為我們明白了很多的真理.
使我們的生命能夠豐富.
對屬靈的國度.
我們不感到陌生.
甚至我們極渴望.
在我們有限的人生完成之後.
將來去到永恆.
求主帶領我們.
我們今天晚上會再看《馬拉基書》.
這四章的經文.
其實很短的.
但是幫助我們能夠有系統地去研讀.
特別是一些可能跟神學有關聯的課題.
譬如我們提到的《簡選論》.
都讓我們可以透過聖經去掌握.
去明白.
我們仰望禱告.
奉耶穌基督的名.
阿們.
好.
上一個星期剛剛開始了.
大家在網上看到.
不好意思.
背景就是牆壁和電視.
大家忍耐一下.
或者你不用開著.
就這樣聽著一邊寫都可以.

$^{41}$不用看著我.
反正我也不會看著鏡頭的.
上一次.
即是前一個星期.
或者可能我要講得慢一點.
因為現在沒有麥克風.
可能聲音夠的.
錄音那裡是沒有問題的.
所以我可能就盡量減低聲浪.
因為怕喉嚨還是比較癢.
就是我們上一次提過.
因為以色列人.
他們來質疑神.
其實還愛不愛我們這個民族呢?.
這些是一些很大的課題.
居然聖經會觸碰到這些深層次的問題.
神到底還眷顧不眷顧.
這個祂揀選的民族呢?.
還有作者才知道.
祂挑起這個問題.
接著百姓當然會問一個很實際的問題.
就是何以見得呢?.
這是要解答的.
那你有沒有什麼充分的理據.
說神愛我們這個民族呢?.
這個愛是古代的時候.
還是現在還在愛呢?.
現在在愛是很重要的進行式.
那將來還愛不愛呢?.
這個問題.
其實如果我們從來都不敢問的時候.
原來聖經有這樣的問題.
甚至乎有解答.
不過那個解答.
比我們更加難以理解.
正正就是觸碰到一個絕對的揀選.
很有趣地.
我們上一個星期.
上一次就曾經開始看到羅馬書的經文.
其實在新約有四段聖經.

$^{81}$關係到這個耳掃的經文.
我們可以先看第一段.
其實也很清楚的分段.
第一個段落.
上個星期我們說過那個題材.
作者的寫作.
就不是說「耶和華如此說」.
接著就說一些審判的話.
或者好像這個.
撒迦利亞書那樣說一些將來的預言.
等等.
或者是透過一些.
神蹟或者是一些比喻來說.
這裡是一個好像是議論題.
用一些議論的文章的題材.
就是說我提出了一個東西.
你們認不認同呢?.
那些收信的讀者說.
我們當然不認同.
如果認同我們就不會有這樣的狀態.
其實作者透過一個這樣的設問的問題.
來點出他們現在的狀態.
不是神所喜悅的狀態.
當然是說屬靈的狀況.
和神的關係的狀況.
或者他們既然具備為以色列人.
神的選民.
那應該就是要在其他的民族當中.
表現出他們是蒙應許.
是很有promise的.
就是很有那個的應許的.
覺得自己是很亮麗的一個民族.
但是為什麼死著死著.
又沒有光彩.
沒有mood的呢?.
因為第六節開始.
就進入到很核心的.
那個宗教的敬拜.
那就罵的不單是那些敬拜者.
是那些帶領敬拜的人.

$^{121}$就是宗教人士.
變成了是多麼的靠近.
是碰到一些和我們今天教會.
日後的教會.
都息息相關的.
就是整個人生.
我們暫時看不到神.
祂的真相,真象.
但是我們是怎樣來到去.
以虔誠敬畏神和敬拜神呢?.
那就是說以色列人.
當時這本書的時代.
和我們今天.
其實他們都是憑信心這樣去做.
都是這樣來到敬拜神.
那個邏輯是完全一樣.
所以這些問題.
全部應該是我們內心要問的問題.
或者我們人生信了一段日子.
我們會不會有一種浮沉.
會不會有一種下降了.
都需要來到去重新檢視呢?.
那我們從哪裡開始.
是個人還是整個群體.
還是整個時代呢?.
這些正正就是.
經文提供了一個切入點.
讓我們可以來到去.
透過去研讀.
我們就有了一個方向了.
當然這裡應該是有一個結論.
神過去,今天.
到將來.
都愛祂所愛的以色列人.
就是說.
這個就真的可以變成我們每一個人.
無論你現在是什麼的光景.
什麼的狀況.
那個可能是短暫.
可能是比較長時間的.

$^{161}$譬如一些慢性病.
一直纏繞著你.
但是都可以進入我們這個結論裡面.
神一直都會愛每一個祂揀選的人.
這個是很重要的.
但是要怎樣醞釀去到這個結論呢?.
我們就要先看完這一段五節之後.
我們就要看新約四大段聖經.
接著我們就要透過一直看.
耳塑,牙角.
他們的人生.
變成我們抽絲剝繭.
就可以很有系統地.
喂!神為什麼你.
總是要愛牙角啊?.
為什麼你總是要憎惡耳塑啊?.
那就要讓我們掌握到這些經文了.
第一節開始.
其實第二節.
因為一開始就進入那個切文的問題了.
耶和華說.
當然是先知代表神說.
我曾經愛你們.
你們卻說.
就是說作者假設百姓現在的光景.
然後代表著百姓來回覆他們的內心.
感覺到神好像跟他們很遠.
不感到神愛他們.
所以祂說你們卻說.
你在何事上愛我們呢?.
在何事上.
就是How?你怎樣愛我們啊?.
在什麼事情上.
你表現到你是神聖的愛.
對我們的保佑啊.
干預啊.
譬如我們在哥山地.
他們進入迦南之前.
在出埃及的時候.
雅各的子孫.

$^{201}$就住在哥山地.
試過那個地方特別富庶.
特別蒙福.
試過哥山地.
甚至乎有日光.
其他地方都黑的.
很特別.
就像大光燈照著那個地方.
那個地方就是神的保護.
其他埃及人都沒有這樣的祝福.
你看《創世記》有記載的.
應該是五十章那裡說的.
這裡告訴我們.
你在何事上愛我們呢?.
如果我們透過話劇.
你會感受到.
說話的人.
好像去了一種.
心如止水的內心.
你曾經愛我們?.
是的.
曾經是出埃及的時候.
我們的老祖宗亞伯拉罕的時候.
或者大衛皇朝的時候.
但是現在我們是經歷了被擄.
整個聖殿沒有了.
現在重建那個聖殿.
重建完之後.
慢慢慢慢.
過了一百二十多年.
又有尼熙米重建那個城牆了.
沒錯.
是有自己的一門一戶.
有自己的城牆.
有自己的敬拜的中心.
但是.
金飛色彼.
或者很多的.
現在他們看到的狀態.
都是令到很多人的心都冷卻了.

$^{241}$或者去敬拜.
去奉獻.
去朝聖.
去獻祭.
都有一些的卻步.
或者.
可能有時某些大的原因.
小的原因.
都推卻去敬拜.
這些.
相信我們.
其實.
每間教會都出現的.
一定有很多人.
慢慢慢慢.
有時那種疏遠.
有親戚進來.
又去喝茶.
又回不了教會了.
或者有其他的原因又不行了.
可以有很多很多的因素.
但是.
原來.
歸根結底.
百姓將他們內心.
要說的.
內心的那個真實.
就透過這一句話.
來去.
否白出來.
你在什麼事情上面.
顯明你的愛和你的保守呢?.
就是說.
我們不感覺到.
你愛我們.
好像昔日.
那些的神跡性.
很玲瓏浮凸的.
那個干預.
那個保護.

$^{281}$所以呢.
神怎樣去拆解這個.
內心的問題呢?.
真是.
心情的否白.
百姓的內心.
不感受到.
神的愛.
接著就是我們要探討的這個問題了.
就是第二節的中間這裡.
耶和華就說了.
這個是解答他們.
你說.
我在哪些事情上面.
沒有那個Blessing.
有啊.
一直都有啊.
以素.
豈不是雅各的哥哥嗎?.
因為以素是出生的.
他渾身多毛.
出來之後.
接著雅各就是第二個兒子.
是媽兒子.
其實是不是算是媽兒子呢?.
一個出一個後.
就抓住以素.
抓住以素的手.
這樣來到接著出來.
接著你留意.
以素不是雅各的哥哥嗎?.
我卻愛雅各.
護以素.
使他的山靈荒涼.
把他的地業.
交給曠野的野狗.
這些是歷史.
我們已經知道的了.
無論民數記.
出埃及的時候.

$^{321}$又或者先知書.
尤其是這個俄巴底亞書.
整章就是講如何審判以東的.
如何講到.
他們無論多麼的.
在一些.
有些很險要的地方.
即是現在的皮特拉.
Petra那個地方.
有些很險要的地方.
神都會將他們拉下來.
他們就算有很多的聯盟.
最終都會失敗.
這些是有先知書的.
即是說古時.
神是一路對以東是有審判的.
這裡告訴我們了.
第四節.
以東人說.
我們現在需被毀壞.
卻要重建荒廢之處.
似乎這裡.
作者這樣描述的時候.
即是說這個是.
On the going的一個進行式.
即是說.
原來以東.
以東即是以素的後人.
創世記講過幾次了.
以素就是以東.
就來到去講說.
我們不會感覺到一兩次的失敗.
被敵人攻擊.
或者神要對我們.
有一些追討,審判.
我們就會氣餒了.
這裡告訴我們.
我們雖然被毀壞.
卻要重建荒廢之處.
但是萬君子耶和華如此說.

$^{361}$任他們去建造.
我必拆毀.
人必稱他們的地為罪惡之境.
稱這些裡面的人.
為耶和華永遠老老之民.
接著就是第五節.
你們必定會親眼看見.
即是回應第二節的問題.
神啊.
你在哪些事情上愛我們這個民族呢?.
神就用了幾節聖經.
透過先知的一個複述.
這樣來講述.
神一直就是這樣來到去.
在這個揀選的定論底下.
對雅各的後人.
就是以色列.
一直來到去保護.
這個保護不是代表.
溺愛的任他們無法無天的為非作歹.
胡作非為.
而是一路有先知的教導.
先知的警告.
先知的預言.
也有他們要被懲罰.
要有他們要面對國破家亡.
聖殿都被拆毀.
神的憲制的器皿.
都被擄到去巴比倫.
是有七十年.
他們要荒涼.
他們要荒廢.
但是呢.
揀選的定義就不是我們一般人所說的.
無論他多壞你都要保住他.
這樣才叫愛.
這個不是聖經所說的揀選.
和愛.
而是這個民族.
一直都會蒙主的管教.

$^{401}$蒙主的保守.
還有他們為什麼能夠可以.
好像否極泰來那樣.
能夠樹等.
就是只剩下樹等.
都可以再有枝條.
再有茂盛.
再有成長.
正正就是揀選的重要道理.
這個真理.
所以.
這裡就告訴我們.
你們第五節.
你們必定會親眼看見.
也必定會說.
願耶和華在以色列境界之外.
被尊為大.
這裡五節聖經.
已經解答了一個.
好像整個以色列民族.
很深遠的.
很長久的.
因為經過很多的年代.
很多很多的事件.
使得他們要面對神.
在他們的民族當中.
種種的.
無論是管教.
無論是一些困難等等.
神是沒有令到那些困難減少的.
有時是會失修.
有時是會饑荒.
有時是會有外敵的入侵.
因為這樣特別保護他們.
但是.
這樣的情況下.
反而讓我們看見.
神藉著這麼多事情的演變.
從來沒有離棄這個民族.
一直讓他們在種種的變化.

$^{441}$人生的不同階段.
他們仍然有一個.
他們必須要抓緊的事實.
就是神永遠不會撇棄.
他所揀選的子民.
這個民族.
這就是這段聖經.
作者要來回答.
以色列人很深層的內心的怨言.
那種現在整個民族.
好像落在一種.
可有可無的宗教敬拜生活.
這是很深層的.
當然說的是整個民族.
我們要明白之後.
就要化做我們自己個人的信仰.
當然猶太人就不會有這個所謂個人的.
他們是整個民族的一個共同信仰.
所以也很難怪他們從小.
我一出生我受了國禮.
我去憲制.
我就是亞伯拉罕的子孫.
我們就是聖約的子民.
每個人都是這樣想的.
所以你說他們不行.
其實是很難想像的.
因為是整個民族.
所以他們不會有個人的信仰.
當然到現代是有不同的.
回來我們要看.
正正就是關於耳塑.
關於雅各.
他們的人生的開始.
那個走向.
上一次我們引述的羅馬書.
普羅就很可圈可點.
他說兩個兒子都還沒有生下來.
他們的罪惡都還沒有形成.
就是剛出生犯什麼罪呢?.
不會的.

$^{481}$但是那個揀選的旨意是已經定了.
是已經很早跟他的媽媽.
這個李伯嘉說了.
所以這裡就讓我們想一想.
我們先看回羅馬書.
接著另外希伯來書也有的.
我們用一點時間看羅馬書第九章這一段.
明顯是神揀選這個課題.
羅馬書九章.
這裡告訴我們.
我在基督裡面說真話.
並不謊言.
有我良心被聖靈感動.
給我作見證.
為什麼這樣說呢?.
正正因為他一直說到.
外邦人在主耶穌的救恩裡面.
大大的蒙福.
一直從第三章的二十一節開始.
一直說到八章的結尾.
都是說外邦人可以因信而得到清義.
得到神的神兒子的名分.
而且是被神選召的.
將來是可以去承受天上的產業的.
對一個猶太人來說.
這一切其實應該跟我們猶太人說.
而不是跟外邦人說.
為什麼保羅會跟羅馬的信徒說.
這些福祿.
這些屬天的榮耀和身份呢?.
所以保羅在這裡.
他現在轉向了.
他講完羅馬書八章之後.
沒有東西可以使我們跟基督的愛隔絕.
講得很盡的,徹底的.
誰人在耶穌裡面.
沒有東西可以使你跟神的愛分開.
沒有東西可以使我們能夠離開神的照管和護衛.
因為我們就在基督耶穌裡面.
但是沒有信主耶穌的猶太同胞又是怎樣呢?.

$^{521}$這是一個在保羅裡面必須要解釋的問題.
因為他是猶太人.
他自己在猶太教裡面浸淫這麼多年.
他自己現在突然間叛教.
歸信了耶穌.
現在傳耶穌給外邦人.
他要交代給猶太的同胞知道.
如果解釋得圓滿.
是透徹又合乎舊約的一貫的真理.
能夠解得通.
至少是很理性地告訴猶太人.
不要再在舊有的宗教體系裡面.
停在那個安書區裡面了.
而是要有一個冒險精神.
就是排除可能你的父母.
可能你的同族的兄弟來怪責你.
你是叛教者.
但這才是新約啟示出來的真理.
就是在主耶穌裡面.
不再分猶太人和外邦人.
這是當時來說很新的道理.
需要有使徒很有系統地.
將舊約和神的心意抽絲剝繭.
來解釋給當時代的人知道.
令兩下猶太人和外邦人合一.
進入真理的救恩裡面.
對於我們後世的人當然沒問題.
因為聖經已經傳了二千年.
但當時代本來就需要有系統地解釋.
所以第九章到十一章.
這三章聖經就是要解釋.
為什麼猶太人現在是在救恩的門外.
而得不到神的救恩.
但不代表神撇棄這個民族.
因為正正就是那個揀選的問題.
將來這個民族是會全家歸主得救的.
所以現在本來就有系統地去講述.
所以一開始就好像先申報利益.
我在主裡面說真的話.
而且有聖靈感動他去說.

$^{561}$為什麼我一路說外邦人可以白白信主得救.
而猶太人為什麼在耶穌的救恩計劃裡面.
似乎現在是沒有份呢?.
保羅就解釋了.
第二節.
我是大有憂愁.
心裡面時常的傷痛.
為我的弟兄.
當然是猶太人.
我的骨肉之親.
就是自己被咒詛.
明明上文第八章說.
沒有人可以在基督裡面.
我們離開他的保護的.
咒詛都不行的.
但他現在說.
為了這樣的原因.
我就算自己被咒詛.
與基督分離.
我都是願意.
所以不是說一些很隨便的話.
他是內心很真誠地愛他自己的同胞.
很急他們還沒看到真理.
他們還停留在舊有的觀念.
舊有的價值觀.
舊有的宗教的信仰.
所以宗教.
如果我們對神學.
對真理不得全面.
我們停留在一套道理.
有時這套道理.
是會框著我們.
框著你們.
我們變成一種無形的枷鎖.
我們自己完全不察覺.
就會變成.
我們很怕.
望出這個框以外.
因為離開這個框.
我們是沒有安全感.

$^{601}$所以在宗教上.
是頗奇妙的.
是值得我們要小心的.
所以為何邪教的人.
你只要這樣跟他說.
他都不會理睬你的.
因為他離開那個框.
你說我們的教主.
你說我們的教義是邪教.
對他來說是生與死的問題.
所以我們必須要信主之後.
盡快對整傳的救恩.
全面的真理.
必須要掌握和融會貫通.
否則很容易被人誤導.
所以保羅說.
他的內心是很迫切.
希望自己的同胞.
跟他一樣.
或跟外邦人一樣.
來去恩訊基督得救.
即是說基督如果在猶太人的內心.
有那個地位.
有應該認順的身份地位.
猶太人就會改變了.
正正現在就不是了.
他是木匠的兒子.
是周圍去說自己是神兒子.
是該受咒詛的那個人.
最終被羅馬.
被祭司長定為死罪的一個死囚.
在猶太人的內心.
這個耶穌就是招搖撞騙的一個騙徒.
還敢說自己是彌賽亞.
所以保羅要怎樣令到猶太人.
內心裡面.
完全打開一個空隙.
接納主耶穌進入他的內心.
要有這一種完全降伏的認順.
難上加難.

$^{641}$你跟一些有見識彌賽亞的人說.
耶穌是彌賽亞.
很難.
所以保羅的人生.
雖然他的人生很短的傳道.
但是他做到一件很重要的事.
就是至少在他有生之年.
盡量在書信裡面.
告訴猶太人.
耶穌就是彌賽亞.
很重要.
所以我們今天.
我們很慶幸我們沒有這些包袱.
我們認順主耶穌就是彌賽亞.
對於我們來說都很抽象的.
救世主彌賽亞.
和舊約的各方面的Blessing.
但是我們新約也有我們的好處就是.
進入永生.
和眾物主真神的結連.
相交.
我們隨時可以去到致勝所.
天上的不是地上的.
隨時可以奉耶穌的名.
去到得到神的保守.
聽我們的禱告等等.
這些我們要有生之年.
要盡力地呼吸.
盡力地祈禱.
你要擴張你的祈禱的容量.
你不要.
祈兩句就已經不知道說什麼了.
應該一路擴大你的容量.
這個正正是.
耶穌是彌賽亞.
耶穌是救世主.
是神兒子給我們的應許.
今天有多少人.
我們真的奉耶穌的名.
去為一些很重大的事情去祈禱呢.

$^{681}$剛才我們有談論.
正生現在還是調查階段.
我們不便去提出有什麼質疑.
但是去年.
長洲的教會已經知道.
變賣那些資產.
賣了給一間教會.
非常貴的.
當年他們是用多少錢買回來.
去做這些青少年工作呢.
是地舖來的.
長洲的教會很早已經知道了.
我很早已經知道了.
但是當然沒有這樣聯想.
說儒家發言.
虧空等等.
這些是我們需要祈禱的.
有很多本來對福音未必抗拒的人.
是很反感的對福音.
對你們所說的耶穌是很反感的.
我們要幫他們.
你不要看那個人.
你看耶穌.
很難不看的.
去到英國.
即是看他們的Facebook.
新聞賣出來的那些.
去到英國還說.
我們有這個帶領.
如果真的有其事.
不是被人插贓嫁禍的.
是真的這樣犯罪的.
那又要講上帝.
上帝可以被人這樣用.
所以我們其實要祈禱.
當然我講遠了.
普羅這裡的意思就是想告訴我們.
怎樣令到猶太人.
心裡面有基督.
基督在他們的心裡面有地位.

$^{721}$很難.
真的很難.
所以第四節.
普羅首先來到去肯定一件事.
基於舊約神的揀選.
亞伯拉罕的揀選.
他的後裔是神的肉身的子民.
即是神在地上的子民.
不是神有肉身.
不要搞錯.
第四節.
他們是以色列人.
兒子的名分.
即是以色列這個民族.
榮耀.
即是可以得到神的光年.
敬拜神.
即是他們常常相傳.
致聖所裡面有榮光.
其實致聖所是沒有光的.
你是不會有光在致聖所的.
聖所就有燈台.
但有一個幕簾是完全遮住的.
致聖所是不會有光的.
但猶太人他們的相傳就是.
致聖所是有神的榮光的.
神的榮光是會臨在致聖所裡面的.
所以只是這些都有光.
再加上神在歷史裡面.
怎樣帶領這個民族.
將他們拯救.
將他們成立在迦南的地方.
接著是諸約.
即是各種的約.
無論羅亞的彩虹之約.
摩西的西奈山的十戒.
大衛.
神有應許他和他的子孫.
咦,是不是斷了線呢?.
請問是不是斷了線?.

$^{761}$真的,他現在重新連結.
可能那個上網那裡不是那麼穩定.
可能用203.
我試試用203那裡.
所以這個是一個問題.
有時真的不是那麼穩定.
是不是剛才不穩定斷了線?.
(錄影中).
現在是不是好一點?.
聽不聽到你們?.
可能還在連線.
我剛才剪了音.
聽不聽到你們?.
(聽到現在).
現在聽到,是不是剛才斷了線?.
可能上網那裡自己斷了.
現在我試試連接到另一個.
現在又沒有那個網絡.
203,204.
我試試去黃振輝社的中心.
不好意思,有這些限制.
不要緊.
OK,我們繼續下去.
我這個錄音筆是在錄音的.
可能你們遲些聽回錄音的訊息.
各種的約和律法.
很明顯就是十誡和衍生的.
拉比的解釋的那些條例.
據聞拉比說是有613條.
十誡衍生出來的那些條例.
另外就是一些禮儀.
他們會有很多不同的節期.
他們會去敬拜.
有很多其實是預表著耶穌的救恩.
還有應許.
神對以色列人的應許.
全部都是他們的.
記住,保羅說這裡的時候.
是在說主耶穌還沒來這個世界之前.
神在舊約和以色列人所訂立的契約.

$^{801}$和這些的應許.
在萬民當中就揀選了這個民族成為神的選民.
所以這裡告訴我們第五節.
列祖當然就是亞伯拉罕,以色列,雅各,約瑟.
你想到的,比如約書,一路下去的有時司.
然後又去到大衛.
一路是在堅定那個約的.
是沒有廢棄的.
一路都是這樣延續神對亞伯拉罕的應許.
那個Blessing.
所以這裡說列祖就是他們的祖宗.
按照肉身來說.
基督也是從他們出來.
基督就是大衛的子孫.
伯利恆降生.
他不是利未之派.
但他是永遠的大祭司.
因為他是按照默基西德的等次.
作為永遠的祭司.
他自己不是亞倫的後人.
更加不是利未之派.
耶穌是大衛,即猶大之派.
這裡告訴我們.
基督也是從猶太人那裡出來的.
那就是說耶穌基督的肉身.
肯定就是降生成為以色列人.
他是拿撒勒這個小城的地方的人.
不過我們說過.
他是在萬有之上永遠可稱頌的神.
阿們.
神主耶穌真正的身份.
所以有他的人性.
因為他要來這個世界成為人.
又斷線.
哎呀,真是.
好的就是我可以透透氣.
這裡的WIFI有很大的問題.
很大的限制.
我沒有帶水,隨便她.
抱歉抱歉.

$^{841}$不好意思.
又是斷線斷得很厲害.
他自己開回我不理他了.
繼續吧.
主耶穌他那個神人異性.
在這些地方.
在這些不是那麼喜眼的.
保羅的.
完全沒有想過要講神學的.
他只不過是想告訴猶太人.
你們本來.
你們這個身份.
肉身的身份來說.
你們是神的選民.
甚至乎基督都是從你們出來的.
就是說本來你們應該要很榮幸.
你們可以見證著彌賽亞.
在這個民族裡面出現.
不是空降的.
是在你們這個民族裡面.
不過基督不是普通人.
他在這個世界降生.
但是.
真是這個問題要解決.
(咳嗽).
不是不是.
我現在先解決那個問題.
那怎麼辦呢.
完全斷線.
WIFI問題.
是完全斷線.
都不知道是什麼原因.
我稍後再和Yoyo跟進.
因為今晚的WIFI我不知道為什麼.
會斷得那麼厲害.
可能要轉去別的地方.
也不一定.
將來.
我們回來再看這裡.
很重要的.

$^{881}$大家不要介意.
如果一直斷線.
我可能繼續說.
就不停了.
基督有地上的人性.
但更重要.
我們今天能夠有救恩.
是這位造物主.
化身成為人.
來到這個世界.
所以基督保羅的意思.
是他原原本本的身份.
他的origin.
是萬有之上.
永遠可稱頌的神.
這些完全沒有刻意講神學.
而這樣講述.
就更加證明.
這些神人異性.
在聖經裡已經充斥著.
所以我們今天要認順的主耶穌.
必須是這樣的屬性.
所以我們要留意那些異端.
如果不是這樣理解.
一定會有問題.
所以這裡告訴我們.
正正耶穌有這麼獨特的身份.
他就能夠救外邦人.
猶太人.
因為他是超越一切的神.
是宇宙萬物的主宰.
是他親自來這個世界.
將我們人類拯救.
到第六節.
這不是說.
神的話落空.
因為從以色列生的.
不都是以色列人.
也不因為是亞伯拉罕的後裔.
就都作他的兒女.

$^{921}$唯獨從以撒生的.
因為以撒的哥哥.
是什麼名字?.
就是以瑟瑪利.
就是夏甲的兒子.
將來.
這個以瑟瑪利都要被趕出去.
不可以承受亞伯拉罕的產業.
就是說.
舊約的時候.
都有這樣的安排.
現在保羅就將這個舊約發生的事.
來看為一個很重要的真理.
的一個啟示.
就是說.
就算是以色列人.
原來不是每個以色列人.
都是蒙那個簡選的.
唯有是神的簡選.
臨到的.
才是真正是他的子民.
這個當然就.
當然就這樣驟眼去聽.
各位不好意思.
當然就這樣驟眼去聽.
對猶太人來說.
是非常之噁心.
和是很冒犯的.
因為不是所有以色列人.
都是神的子民.
對猶太人來說.
是很難去理解和接受.
接著說.
唯獨從以撒生的.
才要稱為你的後裔.
這個是《創世記》的21章.
那裡這樣說.
接著第八節.
他說.
這就是說.

$^{961}$欲生所生的兒女.
不是神的兒女.
唯獨那應許的兒女.
才算是後裔.
為什麼呢.
因為所應許的話是這樣說的.
到明年這時候.
我要來撒拉必生一個兒子.
就是說.
以撒還沒有出生之前.
神就藉著兩位天使.
來向亞伯拉罕顯現.
亞伯拉罕就和兩位天使.
這樣來到去.
他甚至乎在《創世記》的18章.
做東西給這兩個天使吃.
接著他們走之前.
就來和亞伯拉罕說.
到明年.
你的太太就會生一個兒子.
這個就是應許的兒子.
因為其實當時亞伯拉罕.
已經是100歲了.
他已經生不了.
應該是99歲了.
但是.
撒拉.
他那時候.
都已經是差不多90歲了.
但是.
神就說.
他要做的事情.
應許的事情.
是不會有改變.
還有來到去限制的.
第十節.
不單止是這樣.
還有李伯嘉.
既從一個人.
就是從我們的祖宗.

$^{1001}$以撒懷了人.
雙子還沒有生下來.
善惡還沒有作出來.
只因要顯明.
神揀選人的旨意.
不在乎人的行為.
乃在乎召人的主.
這句話就很重要了.
就凸顯了揀選的問題.
這個真理的概念.
我們就要問.
當然我們人.
在我們人的水平.
在我們人的所謂判斷.
我們的分析.
由聖經給我們的描述.
內容.
我們看到的就是.
如果持平的說.
其實雅各是更加陰險.
更加的.
做事.
詭詐.
還有.
很部署.
很有機心.
而素呢.
就不是想像中.
表面的那麼.
想像中那麼差.
為什麼會這樣呢.
這是我們人.
透過有限的記錄.
來知道的事情.
所以.
如果我們.
比較持平的說.
就是.
兩兄弟.
在神的.

$^{1041}$那個.
救恩揀選裡面.
都是不合格的.
兩兄弟都.
不合格.
但他們.
在神的.
都是不合格的.
兩兄弟都.
不是哪個好一點.
稱先一點.
都沒有的.
其實兩個都是壞的.
甚至乎雅各是更加壞的.
甚至連神.
他都想欺控.
但是.
這裡告訴我們.
這句話是很.
很預言性的.
是神自己的宣告.
他都要揀選雅各.
他都是要來到去.
祝福雅各.
使他能夠蒙福.
這只能夠說.
神.
他既然有他自己的主權.
他要將什麼給人.
我們是不能說不行的.
當然對一個不信的人.
就會覺得.
那就是說我們全說了.
你這樣說.
但事實的確是這樣的.
那件事如果是屬於我的.
我如果有全權去擁有那件事.
其他人是不能說的.
我要怎樣處理那件事.
是真的應該由得我.

$^{1081}$其他人甚至乎連.
給意見都不准的.
這樣才叫完全的.
自主的主權.
所以這裡告訴我們.
這個就要顯明神揀選人的旨意.
不是在乎人的行為.
在乎這位呼召人的主.
神就對利伯加說.
他還沒出生的.
他兩個兒子.
還在媽媽的肚子裡.
神就對利伯加說.
將來大的要服侍小的.
其實是知道的.
將來呢.
其實是不知道的.
誰先爬出來比賽呢.
真的沒有超聲波也照不到.
但不知道為什麼.
真的是二數搶集先出來.
然後雅各出了.
但出現的聖經很簡潔.
一兩節聖經就說他們已經長大了.
出生.
接著兩句就已經說他們長大了.
一個就是喜歡打獵.
一個就是在家裡.
很安靜的為人.
就是這麼多了.
接著突然間的一幕就是.
二數打獵回來.
很渴.
問有什麼吃.
有什麼喝.
雅各就煲了紅豆湯.
二數就說.
給我喝吧.
雅各其實不是突然間想起的.
是一直有一個部署.

$^{1121}$你將你長子的名分賣給我.
我就給你喝.
二數是很輕佻地.
長子的名分有什麼用呢.
一點都不稀罕.
因為他知道他是長子.
他更加知道.
我們看子女行間.
他真的得到爸爸的歡心.
所以我說說而已.
你要就要.
現在我要你的湯.
實質一點.
說說就不給.
說說就可以給你.
你就想.
可能他真的是這樣的心態.
相信古代的人和我們今天有點不同.
當然如果你是在清朝.
在我們古代.
如果你是長子.
長子是繼承相分.
這些一定是.
現在現代就很少這樣.
古代一定是.
如果是大戶人家.
一定是會有這些想法.
代代相傳.
長子一定是繼承會多一些.
不停重覆.
不知道怎麼算.
這告訴我們.
第十三節.
天上所記.
雅各是我所愛.
以素是我所護.
這就是馬拉基書一章第二節.
和第三節的經文.
就是我們剛才說的經文.
但是保羅.

$^{1161}$他引這一節聖經.
來作出他的結論.
首先他一直說.
猶太人.
你們是神選民.
這是與生俱來.
肯定的.
是神與阿巴拉罕.
我們的列祖納了約.
而且給了應許之地.
迦南地我們.
一直都有我們的.
祖先承受那個約.
即是神不斷.
update與列祖的約.
即是一直都繼承.
但是一直說的時候.
原來將來的彌賽亞.
才是那個救贖的.
那個代理人.
和真正的忠心.
所以你們一直等候的彌賽亞.
而你們居然最終拒絕.
這個彌賽亞.
即是你們一直.
等待的那個應許.
那個blessing完全不能夠.
兌現.
因為你們最終要.
接受的那個彌賽亞.
你們不肯.
納入你們的心靈裡.
就是這位主耶穌.
所以保羅自己.
也這樣說.
你雖然是以色列人.
但是神的揀選.
是一直在那個.
應許裡面.
即是說.

$^{1201}$要有一些人有信心.
去承接.
神的應許.
才可以得到.
就不是有一個.
好像必然的.
你身為猶太人.
你就無論.
怎樣人生拒絕神.
你最終都可以.
得到神的恩秘.
這個保羅裡面.
都已經否定了.
所以他就告訴我們.
他引述兩代.
其實就是說兩代.
只不過是一個for example.
即是說其實這個兩代.
一直在延續一個.
方程式.
就是神不斷.
去揀選.
他自己.
的人.
有信心去接上這個揀選.
的以色列人.
首先就是說.
亞伯拉罕生的以實.
和這個以實瑪利.
唯有以實.
才是蒙應許的子孫.
得到神的.
福樂和延續的應許.
以實又有兩個兒子.
即是第二代.
神繼續延續.
這個.
運作的方程式.
就是兩個兒子.
但神要揀小兒子.

$^{1241}$小兒子才是.
蒙神揀選.
即是說是亞國.
和他的後人.
才可以蒙福.
一直下去的時候.
亞國有十二個兒子.
誰才是蒙福呢.
當然有好幾個.
藥失.
但最終是猶大之派跑出來.
就一路一路延續.
所以你會見到.
一路一路神去堅定.
他和後代的約.
最終去到大衛.
你會見到有這些字眼.
神怎樣和大衛再重新立約.
跟著和他的兒子.
所羅門.
一路延續.
揀選是一路沒有停止.
即是說.
這樣的理解.
馬拉基書的作者.
是很能夠明白.
就是神一直.
對這個民族是不離不棄.
所以.
不只是古時.
神曾經愛我們.
我們只不過是歷史.
我今天感受不到.
其實就不是.
是一直都沒有停過.
對這個民族身上的工作.
當然.
當時馬拉基.
作者的時代.
大概可能都是主前的.

$^{1281}$四百多年耶穌基督未來之前.
的那個時代.
已經有自己的聖殿.
重建好了.
也都經過了.
以色列尼希米的.
重建.
一路都是這樣.
繼續.
運作.
但是慢慢去到一個階段.
信仰成為了一個.
未同著立的.
例行公事.
沒有了那種感受到.
神的愛的熱切.
那種熾熱.
和那種.
在整個的身心靈.
都感受到.
自己是蒙福.
自己是蒙愛.
這樣的情況.
也都正正就是.
當時的以色列人.
面對的那個光景.
那這裡.
就是羅馬書九章.
這裡告訴我們.
那個的經文的.
簡選.
我們先看希伯來書.
有兩段的希伯來書.
第十一章.
希伯來書十一章.
第二十節這一段.
十一章.
如果我們看多一點.
可以看第十七節這一段.
亞伯拉罕恩爵信.

$^{1321}$被試驗的時候.
就把以薩獻上.
創世記二十一章那段經文.
這便是那歡喜.
領受應許.
將自己獨生的兒子獻上.
論到這個兒子.
曾經有話說.
從以薩生的.
才要稱為你的後裔.
第十九節.
他以為神.
還能叫人.
從死裡復活.
他也彷彿從死中.
得回他的兒子來.
就是說.
他真的想斬下去.
他真的獻給神.
但是神的天使.
就禁止他.
就說那個製身已經預備了.
在那些荊棘那裡.
就在那裡拿一隻羊.
出來獻祭.
就是神要試驗.
看阿拉罕是否真的聽從神的話.
人生去到這個階段.
他真的相信.
神的應許.
甚至他如果殺了他的兒子.
他也會使他復活.
就是說他真的認定.
這是神應許給我的.
現在神要收取.
他也是甘心樂意.
這個甘心樂意.
就真是人聖靈性的最高峰.
去到這個階段.
他能夠這樣做.

$^{1361}$不是神要試探他.
而是真的看出.
他對神是完全毫無保留的.
這樣去奉獻.
就是連一個.
他自己所愛的兒子.
都願意.
呈獻給神.
他明知道神是有預備的.
所以就叫耶和華爾納.
接著我們看第二十節.
二十節.
以撒恩著信.
就指著將來的時.
給雅各.
以掃祝福.
你看到.
一節這麼簡單的聖經.
希伯來書的作者.
將次序調轉.
首先得到福氣的.
是雅各.
其實雅各離開了之後.
以掃帶著獵物回來.
弄好了野味之後.
其實當時.
以撒就是老眼昏花.
他已經看不到東西.
以掃回來的時候.
吃了野味.
又摸著以掃.
終於有一個晴天霹靂.
剛才那個是誰來的?.
雅各.
他已經.
取得了我的祝福.
以掃就大驚.
剛才那個就是我的弟弟.
雅各他欺騙我.
他騙取了我長子的名份.

$^{1401}$現在還要來取.
我長子的祝福.
以撒就.
真的驚了.
而且以掃問他.
爸爸現在你再.
祝福我.
以撒說了一句話在創世記.
兒啊.
我的福已經給了雅各.
我不能夠再祝福你.
其實有人這樣理解.
就是說.
以撒他一心都是想.
祝福以掃.
但最終神攔了.
所謂陰差陽錯.
媽媽.
大家從中作梗.
做好了夜尾.
然後給雅各.
用羊皮包住.
就好像聖山是毛.
又令到氣味.
令到以撒就以為.
是以掃盡他的人生.
他有的福氣就給了雅各.
他以為幫了以掃.
誰知.
當以掃這樣來求爸爸的時候.
我們站回創世記.
應該是第28章.
對不起.
應該是第27章.
第30節.
以撒為雅各祝福以筆.
雅各從他父親那裡才出來.
他哥哥以掃.
正打獵回來.
也在了美味.

$^{1441}$那來給他父親說.
請父親起來.
吃你兒子的夜尾.
好給我祝福.
因為真的非同小可.
他們一定會相信.
那個老父親.
是一直接續著亞伯拉罕的福氣.
是直接從神那裡.
得到一些屬天的福氣.
所以是至關緊要.
臨死之前一定要祝福.
我想猶太人當時已經有這樣的傳統.
所以這裡告訴我們.
請父親起來.
好給我祝福.
第32節.
他父親以撒對他說.
你是誰.
他說我就是你的長子.
是他的戰經.
戰經不是因為他被人騙了那麼簡單.
是他的福氣已經給了那個小兒子.
你未來之先.
誰將夜尾拿來給我.
我已經吃了.
為他祝福.
他將來也必蒙福.
以撒.
在這裡.
似乎他.
人算不如天算.
他一心都想.
給他的長子一些祝福.
但是神一早就跟李伯嘉說.
兩個兒子還沒出世.
說明神愛的是雅各.
恨惡的是爾素.
還有會使到大的將來要服侍小的.
小的真正成為那個大靈.

$^{1481}$或者接續亞伯拉罕.
以撒的福氣.
的繼承人.
所以.
第34節.
爾素聽了父親的話之後.
就放聲痛哭說.
我父啊.
求你也為我祝福.
以撒說.
你兄弟已經用詭計.
來將你的福分奪去了.
爾素說.
他明雅各豈不是正對摸.
因為他欺騙了我兩次.
他從前奪了我長子的名分.
你看他現在又奪了我的福分.
爾素又說.
你沒有留下為我可祝的福嗎?.
以撒回答說.
我以納他為你的主.
使他弟兄都給他作僕人.
並賜他五穀生酒.
可以養生.
我現在還能為你作甚麼呢?.
很嚴重.
也可以這樣說.
以撒有些人這樣理解.
我覺得也可能是對的.
以撒在這裡.
他不敢再超越神的管理和揀選.
他終於明白.
神所謂陰差陽錯.
揀選了雅各.
他不敢超越神.
他知道這個長子真的不蒙福.
真的不蒙神的揀選.
他不敢再祝福他.
只給他一些普通的東西就算了.
不敢給他最大的.

$^{1521}$承接亞伯拉罕的名物的福氣.
來給雅各.
所以希伯來斯的作者很有趣地.
以撒給雅各祝福.
就是說長子是雅各.
得到真正的祝福是雅各.
也就是說以撒臨終.
做一件好事.
就是他不敢去祝福以素.
其實你想想.
如果我真的愛這個兒子.
我畢竟是神.
我就照樣去祝福.
現在就反過來就祝福雅各.
因為他還沒死.
他是爸爸.
他是家主.
那些東西是他的.
他不敢.
這個不敢就是他真正敬畏神.
他不能再超越神這個計劃.
真的沒辦法.
本來真的給他的.
但是神禁止.
是難了.
當然是人的陰謀.
因為李伯嘉和他的兒子.
合謀要來搶爸爸的祝福.
因為知道是好的東西.
是接受亞伯拉罕的福氣.
所以你可以這樣說.
雅各他很緊張屬靈的福氣.
不過他的態度和動機是不純全的.
他是想具備己有的享用.
不是想著自己以身作則.
其實他有沒有知道.
媽媽有沒有跟他說.
兒子你不是出生了嗎.
神已經向我顯現.
告訴我.

$^{1561}$你將來要成為大兒子.
蒙祝福.
他有沒有聽過呢.
我猜有.
但是他不想等那個時間到.
明明是沒可能的.
因為阿伯一直看他是愛自己的大哥.
多於愛自己.
很多的閃失.
很容易爸爸私底下捉了他.
就沒有他了.
如果家裡有很多兒女.
有些特別孝順的.
可能阿伯私底下分了一些給他.
都沒有人知道.
所以這個正正是雅各害怕.
他也信不過神的預言.
也很想用人自己的辦法.
先拿了就是先.
拿了就是自己.
所以他這個拿了.
也成為他自己.
接近二十年.
他要顛沛流離.
回到他的舅父拉班的家裡.
來受那個折磨.
惡人自有惡人磨.
他用詭計.
他又被人用詭計來對待他.
不要看他這個人生.
他將來真真正正成為以色列.
神的王子.
經過很多這些磨難.
他的生命就成熟了.
就不是一種二世祖的繼承.
而是真是很名正言順.
他的生命符合成為那個祝福的.
蒙福的人.
所以雅各臨終都祝福.
祝福十二個兒子.

$^{1601}$是很重要的.
創世記四十九章.
就是祝福十二個兒子.
就是說一定是他有福.
才可以祝福.
他一定要去到這樣的經歷.
人生.
他才能夠將這些福氣.
來給他的兒子.
種兒子.
這就證明了雅各他人生.
雖然是可能前大半段.
都是為自己去爭取.
但最終他當他成熟之後.
尤其是三十二章的那個比魯伊勒事件.
他求退之後.
他整個人生就真是以神為信靠.
真是憑信心去活.
他不再懼怕耳塑.
追殺他.
他也不再懼怕回到迦南地.
這也告訴我們.
他後半生真真正正是倚靠神.
他就不去抓住很多東西.
他最終是在找神.
我們看到希伯來書是這樣說.
接著十二章.
希伯來書的十二章.
連電腦也有問題.
(錄音有問題).
電腦也有問題.
現在電腦也開不了.
不知道是不是有種病.
(record in progress).
現在也不行了.
不要理了.
現在我錄音就行了.
我等下再處理.
現在我連電腦也開不了.
我先關掉.

$^{1641}$我錄音的沒問題.
有音錄的.
(聽回).
聽回就行了.
(record in progress).
不斷的.
(record in progress).
我不理了.
現在我們先回來這裡.
十二章的十四節.
希伯來書.
看完這一段關於耳掃.
你們要追求與眾人和睦.
並要追求聖潔.
非聖潔沒有人能見主.
又要謹慎.
恐怕有人失了神的恩.
恐怕有毒更新出來擾亂你們.
因此就叫眾人沾染了污穢.
恐怕有淫亂又貪戀世俗耳掃.
他因一點食物.
把自己長子的名分賣了.
後來他想要承受父所祝的福.
竟被棄絕.
雖然毫哭徹求.
卻得不著門路.
使他父親的心意迴轉.
這是你們知道的.
正正剛才我所說的就是這裡.
爸爸沒有回心轉意.
沒有回心轉意.
就是說耳掃的靈性.
到了那一刻.
他終於肯屈服神的預言.
神選擇了雅各.
他一生可能就是不服這件事.
他一直要幫耳掃平反.
所以他現在兒子怎樣哭哭徹求.
都不能使爸爸的心意迴轉.
爸爸終於明白神的計劃.

$^{1681}$他終於降服信服神的計劃.
這樣的人生雖然是一個負面.
但其實也是很高超的人生.
你終於肯承認.
神在你的生命裡是好的.
這是一個很高靈性的境況.
你很願意接受.
兒子怎樣哀求你.
這樣哭不行了.
兒子你是沒有福的.
我不可以祝福你.
因為我已經祝福了雅各.
那個真命天子.
你不是真命天子.
一個爸爸要這樣接受這件事.
是用慘痛的教訓來經驗這件事.
所以這就是耳掃的情況.
更加突顯.
希伯來書告訴我們.
不要效法耳掃.
他是貪戀世俗.
原來他一生一點都不稀罕屬靈的祝福.
各位我們來到教會.
這麼美麗的禮堂.
這麼動聽的聖樂.
這麼好的講員.
這麼多外來的講者.
有些小教會.
有一個老牧師和一個姑娘.
有時連一個姑娘都沒有.
只有一個老牧師.
堂堂就是那個老牧師講的.
未講完三句.
你都知道他講什麼.
大家一起準備.
十分鐘就開始睡了.
至少我們感受到有不同的牧者.
有不同的外來的講員.
很多不同的訊息.
又有不同的外來的.

$^{1721}$譬如阿克.
十四號來過.
土耳其的陣災.
或者那些事情等等.
有這麼多這麼好的牧者.
但是.
回到這個課題.
我們有多少人.
是一生.
只是貪戀地上的.
而我們對.
屬天永恆的祝福.
我們根本不感興趣.
我們根本無動於衷.
根本我們都不稀罕.
那就是.
價值觀就是.
耳掃的價值觀.
是一個很大的問題.
他最終.
去到.
他很緊張的時候.
其實是去到最後.
那個一百米衝刺.
但是沒他的份.
他很努力.
那時候真的拼命打獵.
打到回來.
打到很晚了.
一心想著可以得到那個福氣.
他的名份無所謂.
他本身就是.
但獎之的福氣就要拿.
但是.
殊不知就算.
爸爸未離開世界都不會給你.
因為爸爸自己.
靈性改變.
是不會再聽你講.
那這個.

$^{1761}$真是神的管理.
在整件事.
其實間接就告訴我們.
神的.
簡選又不是無跡可尋.
有些人.
真是一生.
他一改變信主.
整個生命又是愛主.
愛人的靈魂.
愛教會去傳福音.
你會見到這些人在神的.
簡選裡面是很明顯.
有些人未必明顯.
但都見到他很一心.
在教會裡面成長.
但是這裡就告訴我們.
耳掃.
他一直在那個所謂.
長子他以為的名份裡面.
他去生活.
他去逍遙自在.
這樣走他的人生.
去到最後.
得不到那個祝福.
雖然後來都會和弟弟和好.
但是.
在他的人生.
尤其是他後代.
是被神咒詛.
招致神的刑罰.
因為他的後代.
是那些外邦.
其實是迦南女子.
他知道爸爸.
剛才我們讀了27章.
知道爸爸將福氣.
給了他的弟弟之後.
他就知道爸爸.
是不喜歡迦南的女子.

$^{1801}$他就特意在迦南.
娶了兩個太太.
這樣去激自己的父親.
其實更加印證.
他.
反轉諸途.
神你不給我福氣.
爸爸不給我福氣.
我就更加肆無忌憚.
讓你知道我可以自己闖一條人生的路.
我可以自己.
不要你的福氣.
我可以自己打遍江山.
所以他就去到以東地.
以東人.
就是與神為敵.
所以他的後代.
都因為耳掃他自己的老亡的舉動.
招致後代.
得不到福氣.
之餘還是神說要永遠.
就坐的一個民族.
所以.
一個祖宗.
吳敬虔後代都受到影響.
雅各晚年.
是好事.
他也有子女.
都是好事.
你就可以想像到.
這個都是.
看父親.
在家族裡面的家長.
看他的靈性.
都是好重要.
所以耳掃一路一路.
他的長大,他的成長.
他的漂泊,他的離開.
他很快就在以東.
建立了他自己的山寨.

$^{1841}$慢慢變成一個王國.
但以色列人.
還是在漂流.
還未入迦南.
你可以想像到.
兩兄弟原來.
是有這麼大的人生.
有些是飛黃騰達.
但沒有神.
但是雅各繼續.
是很游牧民族.
沒錯,兒子約瑟.
做了宰相.
但其實都是很多痛苦.
被自己的兄弟害.
但最終.
雅各臨終之前.
他還是去埃及.
和他的兒子團聚.
最終都是這樣.
都是一個小小的.
部族.
但以東已經是.
體系的國度.
你可以想像到.
從地上.
表面.
耳掃沒有什麼咒詛.
沒有什麼損失.
沒有什麼得不到.
神所謂的祝福.
但長遠看.
在永恆看.
在神的眼光看.
其實是天淵之別.
我們不斷去鼓勵.
弟兄姊妹愛主.
追求,我們的信仰.
和我們實際生活的行徑.
要來到去.

$^{1881}$要有一個一致,幫人.
將福音傳開.
但我們只能夠這樣去勉勵.
今天我們難道有個.
在203.
弄一間宗教裁判所嗎.
你今天有沒有什麼.
沒有就要罰,沒有這件事.
不是中世紀.
沒有這樣的權力.
都是在說一個自由的選擇.
既然自由選擇.
就是各人自己面對神.
其實這件事.
也是一件很嚴肅的事.
每個人自己去面對神.
在你的內心的教導.
和帶領.
所以.
可能來到這裡.
先停一停.
我們今天不說這麼長.
因為錄音.
都有困難問題.
我也不知道什麼原因.
是電腦的問題還是上網的問題.
我們試試將來轉另一個房間.
看看有沒有這個情況出現.
我們一個簡單祈禱.
錄音筆是沒有問題的.
天父我們謝謝你.
讓我們看完五節聖經.
盼望我們能夠去明白.
神的愛.
在以色列這個民族裡面.
是一直沒有離棄.
是一直這樣.
對他們的關愛.
雖然今天.
我們都不知道猶太人將來會怎麼樣.

$^{1921}$他們現在那種.
面對這一場的戰爭.
他們的報復.
他們對哈瑪斯.
對巴勒斯坦的攻擊.
令到很多平民百姓.
這樣去傷亡.
有時候我們很難理解.
我們根本不能夠.
很簡單的對錯地分.
但是我們都盼望.
既然神在這一段聖經.
講到明.
你是對你的民族.
你揀選的.
是不會有改變的愛.
那就是說我們每一個信主.
在主耶穌裡面.
我們肯定是蒙恩的.
但是我們知道.
實際面對我們的生活.
有時候.
那些理論.
真的有時候很遙遠.
因為當我們.
面對我們一些處境.
或者一些困難的時候.
又會令到我們.
思前想後.
想到的可能是神.
未必好像介入.
或者未必會知道我們的困難.
但是我們都求主.
讓我們能夠慢慢學習.
透過很多的經歷.
尤其是一些回望.
過去和在我們.
面對的事情.
有時候那種很出人意表的化解.
就讓我們看到神.

$^{1961}$你很喜歡在暗中.
來幫助.
在暗中來施恩.
和聽我們禱告.
以致我們能夠.
更加在生活裡面.
認定.
這個真理.
就是主的揀選.
是堅定不移.
是不會有後悔的.
願主來施恩.
我們的技術.
令到網上的弟兄姊妹.
是有困難.
很抱歉的.
盼望神都保守.
這段的錄音.
都可以慢慢讓弟兄姊妹.
來聆聽.
我們多謝你聽禱告.
奉耶穌基督的名.
我們下次再見.
\newpage



\section{創世記 16:1-13}
\label{sec:_CjhJKqUcA8}
\textbf{2024-01-20 晚堂《沒有被遺忘的女兒》創世記 16\_1-13 講員:曾裔貴牧師}
\newline
\newline
連結: \href{https://youtube.com/watch?v=-CjhJKqUcA8}{\texttt{https://youtube.com/watch?v=-CjhJKqUcA8}} ~~~~ 語音日期: 2024-01-23
\newline
\newline
\hyperref[sec:dR8L9GrF4hQ]{\small{< < < PREV SERMON < < <}}
~
\hyperref[sec:index_chronic]{\small{[返順時目]}}
~
\hyperref[sec:index_scriptual]{\small{[返順卷目]}}
~
\hyperref[sec:DTr4fwkGAak]{\small{> > > NEXT SERMON > > >}}
\newline
\newline
創世記 16:1-13
\newline
\begin{longtable}{cl}
\hline
\hline
章節 & 經文 (和合本修訂版)\\
\hline
16:1 & \begin{tabularx}{0.7\textwidth}{X} 亞伯蘭的妻子撒萊沒有為他生孩子。撒萊有一個婢女，是埃及人，名叫夏甲。 \end{tabularx} \\ \\ \relax
16:2 & \begin{tabularx}{0.7\textwidth}{X} 撒萊對亞伯蘭說：「看哪，耶和華使我不能生育。你來和我的婢女同房，也許我可以從她得孩子 。」亞伯蘭聽從了撒萊的話。 \end{tabularx} \\ \\ \relax
16:3 & \begin{tabularx}{0.7\textwidth}{X} 於是亞伯蘭的妻子撒萊把她的婢女，埃及人夏甲，給了丈夫為妾；那時亞伯蘭在迦南已經住了十年。 \end{tabularx} \\ \\ \relax
16:4 & \begin{tabularx}{0.7\textwidth}{X} 亞伯蘭與夏甲同房，夏甲就懷了孕。她看見自己有孕，就輕視她的女主人。 \end{tabularx} \\ \\ \relax
16:5 & \begin{tabularx}{0.7\textwidth}{X} 撒萊對亞伯蘭說：「我因你受了委屈。我把我的婢女放在你懷中，她見自己懷了孕，就輕視我。願耶和華在你我之間判斷。」 \end{tabularx} \\ \\ \relax
16:6 & \begin{tabularx}{0.7\textwidth}{X} 亞伯蘭對撒萊說：「看哪，婢女在你手裡，你可以照你看為好的對待她。」於是，撒萊虐待她，她就從撒萊面前逃走了。 \end{tabularx} \\ \\ \relax
16:7 & \begin{tabularx}{0.7\textwidth}{X} 耶和華的使者在曠野的水泉旁，在書珥路上的水泉旁遇見夏甲， \end{tabularx} \\ \\ \relax
16:8 & \begin{tabularx}{0.7\textwidth}{X} 對她說：「撒萊的婢女夏甲，你從哪裡來？要到哪裡去？」她說：「我從我的女主人撒萊面前逃出來。」 \end{tabularx} \\ \\ \relax
16:9 & \begin{tabularx}{0.7\textwidth}{X} 耶和華的使者對她說：「你要回到你的女主人那裡，屈服在她手下。」 \end{tabularx} \\ \\ \relax
16:10 & \begin{tabularx}{0.7\textwidth}{X} 耶和華的使者對她說：「我必使你的後裔極其繁多，多到不可勝數。」 \end{tabularx} \\ \\ \relax
16:11 & \begin{tabularx}{0.7\textwidth}{X} 耶和華的使者又對她說：「看哪，你已懷孕，要生一個兒子。你要給他起名叫以實瑪利，因為耶和華聽見了你的苦楚。 \end{tabularx} \\ \\ \relax
16:12 & \begin{tabularx}{0.7\textwidth}{X} 他為人必像野驢。他的手要攻打人，人的手也要攻打他；他必常與他的眾弟兄作對。」 \end{tabularx} \\ \\ \relax
16:13 & \begin{tabularx}{0.7\textwidth}{X} 夏甲就稱那向她說話的耶和華為「你是看見的神」，因為她說：「他看見了我之後，我還能在這裡看見他嗎？」 \end{tabularx} \\ \\
[1ex]
\hline
\hline
\end{longtable}
$^{1}$我們一起同心祈禱.
我們可以來到去聆聽聖經的真理.
今天晚上我們會用聖經人物.
夏格這一位埃及人的侍女.
來作為講道的題目.
我們真的很盼望.
這些名不見經傳的人.
居然在聖經裡面是這麼突出.
我們很盼望.
我們的人生在禮拜堂裡面.
能夠與主相遇.
能夠認識你.
是那一位看見我們苦情的神.
今天晚上我們等候在你的面前.
奉主耶穌基督的名.
阿們.
都是真的送上祝福.
在紅恩堂差不多三個星期了.
又開始很多跟一些領袖.
領子妹來接觸.
去了解他們整體的運作 文化等等.
開始都會有一些的想法.
要開始準備來實施.
明天的堂委會.
就會有第一次.
來給每一位領袖了解一下.
在面對面前有什麼方向.
與同工一起為神的家去做夢.
如果要我講聖經人物.
我講了很多.
但很少講女性.
很感恩.
黎晟穆斯今季編了講聖經人物.
本來今晚不是我講的.
本來是上個星期的.
因為有些聚會要調了.
變了編了我今晚講下甲.
明天鄧全也是講下甲.
經文可能有些不同.
由一個男性的角度去講女性.

$^{41}$我都是很新的嘗試.
所以在座的女性的姐妹.
如果我說得不夠深度.
不夠肉麻.
道肉.
你們請原諒.
因為男性很難體會一個女性.
受苦的女性是怎樣的感受.
這個妹仔在埃及.
都不是用重價買回來的.
是亞伯蘭.
他姬芳.
整家人下了埃及.
下了埃及之後.
就跟太太講.
現在沒有人知道你是我的太太.
你要認做我的妹.
因為我不知道埃及那些人.
是不是不好的人.
如果不好的時候.
他就搶了你.
那就大件事了.
撒萊就說好啊.
應該沒事的.
如是者被法老見到.
撒萊是極其美麗的.
就要娶她為妻.
就是說這個大舅.
是妹妹.
娶她就是大舅.
送很多財物給她.
到洞房花燭夜那一晚.
神擊打法老.
十二章有記載.
第十二章的十七節.
耶和華因亞伯蘭妻子撒萊的緣故.
降大災於法老和他的全家.
名門正娶.
皇室的婚禮.
娶你的妹妹.

$^{81}$但是當晚就降大災.
在法老和他的全家.
什麼災呢.
他又怎麼知道會是神的刑罰呢.
他就召了亞伯蘭來.
你到底做了什麼事.
為什麼不告訴我她是你的太太.
為什麼你說她是你的妹妹.
以至我幾乎因為她的緣故.
遭致不測.
這樣的情況.
所以第二十節.
於是法老吩咐人.
將亞伯蘭和他妻子.
並他所有的都送走.
原來亞伯蘭是瘟神.
在法老的眼中.
我娶你的妹妹.
我遭致大的災殃.
然後就送走她.
十三章就說.
亞伯蘭帶著很多的財物.
亞伯蘭帶著他的妻子和羅德.
並一切所有的.
都從埃及上南地去.
回到以色列.
帶著很多很多的財物.
他說金銀生畜極多.
相信一直說下去的時間.
這個埃及人下甲.
就是亞伯蘭這一次的戰利品.
這一次法老送給他的一個物件.
大家可以想一想.
你作為一個女性.
小時候你已經在法老王的王宮.
不是做公主.
是做下人.
後來輾轉因為亞伯蘭這個人的緣故.
又被送去給亞伯蘭.
成為亞伯蘭的太太的妹仔.

$^{121}$因為一而再聖經的描述.
往後就是說.
埃及人使女 撒萊的使女.
這樣稱呼她.
接著就是我們今天看的經文.
其實我們還沒看之前.
再看第十五章.
十五章有一個很大很嚴肅的立約.
神同亞伯蘭立永約.
稱她為信心之父.
十五章第六節是我們很熟悉的經文.
告訴我們亞伯蘭信耶和華.
耶和華就以此為他的兒.
華箏說 你想想你是下格.
離開了法老.
不知道有沒有奴役他.
小時候就不在自己爸爸媽媽的家.
現在就輾轉由埃及.
出埃及 進入迦南.
跟隨的不是一個虐待人的家庭.
是信心之父的家庭.
是神向亞伯蘭顯現的家庭.
亞伯蘭和撒萊是很有在神面前的信心.
你以為你是下格.
天天有陪靈會.
亞伯蘭可以見到神.
神跟他說話.
告訴他 你要生誰的兒子.
等等這樣說.
十五章這裡已經告訴我們.
亞伯蘭已經跟神說.
我都沒有兒子.
怎麼能夠繼承你給我的產業呢?.
神就說 當然我們不說太長.
那個以便 以謝.
亞伯蘭的僕人.
想繼承他的產業.
神說不行.
一定要你所生的才可以繼承.
所以來到這裡的時候.

$^{161}$我們就見到十六章.
作者是很清楚告訴我們.
亞伯蘭那個時候已經是八十五歲.
這是第十六章的第三節.
於是亞伯蘭的妻子撒萊.
將侍女埃及人下格.
給了丈夫為妾.
那時亞伯蘭在迦南已經住了十年.
上文的兩節是很清清楚楚告訴我們.
十年有沒有家計會教你怎樣生育.
有沒有什麼偏方.
十年一定要由你生才可以繼承.
神給你的產業.
十年的光陰.
有多少個十年呢?.
亞伯蘭用了多少個辦法.
和太太去生孩子.
生不到.
又不能做DNA.
有什麼問題.
有沒有染血.
是男的有問題還是女的有問題.
都不知道.
都是靠猜.
等了十年之後.
撒萊終於緊張了.
就和丈夫說.
不如就要下格.
似乎下格在撒萊的身邊.
都是時時跟隨.
時時和撒萊一起.
撒萊能夠看到.
似乎這個埃及人.
可能是一線生機.
你想想會不會又是世紀婚禮呢.
信心之父.
信心的家庭.
娶一個妾侍.
而且得到大婆的允許.
會不會是一個迦南地.

$^{201}$一件盛事呢.
聖經完全一句都沒有記載.
到底這一句話.
一晚的行動.
還是有部署.
什麼時候復活節就結婚了.
就行禮了.
似乎沒有.
因為她是妾侍.
因為她是埃及人.
因為她從來都是一個物件.
從來都是一個沒有人知道.
沒有人會記住的一個妹仔.
從來沒有身份.
從來就沒有自己的社會地位.
更加不要說家庭的地位.
進入那個女性的世界.
進入一個被壓迫的世界.
的一個人.
一個沒有名字的人.
一個沒有自己的身份的人.
居然間輾轉輾轉.
有神的帶領.
去了信心之父的家庭.
甚至乎現在懷了他的骨肉.
一天都光了.
這個就是將來的龍種.
將來繼承亞伯蘭的產業.
就是以實瑪利.
神聽見你的苦情.
這樣的意思.
以實瑪利.
神幫他未出生的兒子.
改一個名字.
叫做以實瑪利.
我們再進入經文的時候.
下甲一夜之間.
有了新紀.
開始學人驕傲.
開始嘲笑自己的主母.

$^{241}$非同小可的大事.
從來都不敢的.
做丫鬟.
你想想上一代的人.
在一個壓迫的家庭.
從來都不敢多說一句.
不是現在我們這個現代.
是上一兩代的人.
是從來不敢的.
一個女性多說一句.
都會被丈夫罵的.
何況你是丫鬟.
我媽媽就是這樣.
多說一句都不敢的.
不會駁嘴的.
但是居然對主母.
這樣的嘲笑.
罪大惡極.
非同小可.
果然我們見到的後果就是.
信心之父的家庭.
可以是這麼殘忍的虐待一個孕婦.
信心之母殺來.
信心之父就說.
她在你的手中.
你想怎樣就怎樣.
都不理那個骨肉.
怎樣虐待呢.
只是一句.
我們有很多空間可以去思想.
信心之父的家庭.
照計不應該會發脾氣的.
照計不應該.
來虐待一個有生己的人.
那個有生己.
不是走出去.
去和那些不三不四的人.
是和你的丈夫.
是你允許的丈夫.
說明是為他留後的.

$^{281}$但是火遮眼.
我又不懂.
因為我們是男性.
我們不懂一個女性.
那種火遮眼是怎樣的.
所以男性說女性.
是很有限制的.
火遮眼到一個地步.
信心之母都發脾氣.
可能那些話比粗口還難聽.
虐待到這個下甲.
要馬上走.
走到一個水泉.
在那裡不知道去哪裡.
即是寧願回埃及.
我都不要留在信心之父的家.
這樣.
曾經看著這對夫婦.
見到神.
神和他們立約.
是多麼敬虔的家庭.
但是居然自己受害者.
被信心之父的這對夫婦.
這樣去糟蹋.
虐待.
到底有多大程度的虐待.
是言語的虐待.
還是身體的虐待.
聖經一句.
沒有再說.
你想想你是下甲.
感受到一個這樣的家庭.
你可以想像得到.
如果他不是基督徒.
他用的不是基督徒的方法.
沒有所謂.
但是基督徒對待那些.
下甲這一類的人.
本身就是埃及人.
是沒有信仰的.

$^{321}$不過跟隨了信心之父.
不到他可以選擇.
跟隨了來到迦南地.
見到自己的主人.
這麼敬虔的遇見神.
神和他立約.
這麼好的環境.
十年裡是多麼開心.
開始對神有些印象.
你的神真的有些東西.
我也想認識一下你的神.
但是突然之間.
十年後的一天.
你跟丈夫說.
我沒得生.
是你沒得生還是我沒得生.
現在要試一試.
就找第二個女人.
就是下甲.
一試就有得生.
那就是赫薩萊的問題.
不關亞伯蘭事.
這樣的時候.
對下甲來說.
應該是更加開心.
更加歡喜.
因為神的應許.
臨在我的肚腹身上.
我現在成為代孕的一個婦人.
多麼開心.
但是她自己還沒相信.
她自己沒有神.
是她主人的神.
所以她的脾氣.
她的驕傲.
一個女孩有多驕傲.
從來不敢驕傲.
原來心底裡是很驕傲的.
心底裡可以對主母.
可以說些什麼話.

$^{361}$單打十年都沒得生.
是不是有問題.
有沒有驗身.
這些話對一個女性來說.
很難聽的.
男性也是不懂的.
完全不知道怎樣感受.
很痛的.
說到你最痛的地方.
那些話.
一句就已經傷了你的心.
主母就是.
什麼屬靈.
什麼敬虔.
什麼夜日培靈會.
那些都暫時放下.
那些火就上了.
就爆了.
就跟丈夫說.
現在你准許我虐待她.
我現在因為她而受氣.
似乎從來都不需要受氣的.
一個主母撒內.
從來都是開開心心.
環境是很順暢.
丈夫是信心之父.
丈夫跟神立了永約.
其實是神跟她立了永約.
你作為她的太太.
多麼有面子.
全個迦南地最屬靈的就是我們的家庭.
就是我的丈夫.
但是一靠近.
有一些事情.
跟自己的切膚之痛.
一句話.
針埋你的內心.
一切的敬虔.
一下子一筆勾銷.
當面虐待自己的妹仔.

$^{401}$下格.
虐待一個有身己的婦人.
什麼光環都沒有了.
衷心希望.
靜身書院.
全個香港現在都知道的.
衷心希望是一個誤會.
不要是一個罪案.
不要是一個醜聞.
基督教 基督徒.
做在未信的人當中.
那種羞恥感是很大.
我們需要祈禱.
信心之父的家庭.
虐待妹仔.
應該可以上頭條了.
最敬虔的迦南的一個家庭.
居然間關上房門.
虐待妹仔.
妹仔要逃走.
你想想妹仔的苦情是多麼深.
都還沒算.
她本身都是很苦的人生.
本身都是很苦的一個命.
但是現在因為自己.
只不過是一個小小的過失.
不知道是不是少.
可能已經積埋了很多年的冤屈.
我的命這麼苦.
現在神給我一生.
我還不立刻將我自己的身份抬高.
有可能是這樣的心態.
作為一個女性.
可能是這樣的.
裡面的苦毒.
未有時間爆發.
現在雙方面.
兩府相爭.
一個是本來很位高權重的主母.
一個是沒有任何權的妹仔.

$^{441}$現在妹仔一句話針對主母.
主母就要立刻還以顏色.
本來有權的.
但是如果一直都是敬虔的.
她不需要去建議丈夫.
來納妾.
不需要幫丈夫完成神的旨意.
十年.
其實神都知道那個時間.
那個考驗可以再等一等.
果然去到十八章.
神就說明年就可以生了.
亞伯蘭.
而且是由你的太太.
你的元配撒萊所生.
你可以想像得到.
原來神是有時間.
那個時間不是人可以計算.
只需要還等片時.
你就可以真的很榮耀地.
既可以生.
你又可以不會一個這麼大的失腳.
在整個聖經裡面.
這個撒萊和亞伯蘭.
成為千古.
人們真是有一個唾棄的污點.
因為這麼敬虔.
這麼經歷神.
這麼親近神的人.
可以對一個下人.
說這麼重的侮辱說話.
和虐待的行動.
如果我們不去留意看下格.
隱藏在這麼大的聖經裡面.
創世記說很多神向人顯現.
首先就是神向阿當 夏娃顯現.
神向羅亞顯現.
神向以諾顯現.
很明顯記載的.
神向亞伯蘭顯現.

$^{481}$其他的兄弟沒有的.
你可以想像得到.
亞伯蘭真是繼承著古代的.
那些敬虔的人的善信心.
這樣走進去迦南地.
迦南人是罪大惡極的.
十五章那裡的記載就說.
你要確實的知道.
你們要寄居在埃及.
第四代你們就可以回來了.
因為迦南人.
阿摩利人的罪孽還沒有滿盈.
要等到他們罪惡滿盈.
你們就可以回來了.
你可以想像得到.
亞伯蘭是和其他的迦南人不同的.
不過真正的寶貴和重點.
就是在這裡出現.
就是去到第七節這裡.
耶和華的使者在曠野.
樞爾路上的水泉旁.
遇見他.
對他說 撒萊的使女夏甲.
你從哪裡來 要往哪裡去.
從來沒有人這樣來.
去稱呼她的名字.
直呼其名夏甲.
你看看撒萊每一次和丈夫說.
這個埃及人的使女.
沒有叫夏甲.
有名有姓的.
因為你是主母.
為什麼會叫你的名字.
你就可以想像得到.
有階級在他的眼中.
不過是丈夫在埃及獲得的戰利品.
是一個物件 是一個工具.
但是耶和華的使者.
對一個傷心透折的人.
走到水泉旁 去哪裡呢.

$^{521}$一個自己懷孕的 去哪裡呢.
他可以去的當然是回埃及.
在他的知識 在他的記憶裡面.
但是整個一個大的沙漠.
一個女性 怎麼自己一個人走回去呢.
果然發現當時是沒有人.
你可以想像得到那種荒涼.
那種淒涼.
那種描述不是特意襯托出來.
是真的走投無路.
有的就是一個水泉 還有水.
在一個乾旱的沙漠裡面.
但是在這個時候.
寶貴的地方就在這裡出現了.
上主的使者.
向這個失落的婦人.
來顯現.
首先問一問她.
你從哪裡來的.
你要去哪裡.
似乎真的要撩起那個創傷.
讓她感受到自己的事情.
到底有沒有人會關心呢.
和我內心的痛有共鳴呢.
這個都真的很重要.
現在中小學生的自殺.
不單只是質量 數量都越來越多.
是根本不知道怎麼解決.
他的問題未必是我們那代人的問題.
甚至我們不可能覺得這是問題.
而他們會覺得是問題.
這個就是一個很大的社會.
可能已經浮現出來.
是解決不了的問題.
這個急速性都不是我們任何一個人.
可以解決的.
你可以想像到.
居然現在耶和的使者.
似乎不認識她.
其實叫她的名字.

$^{561}$很久沒有人叫過你的名字.
如果真的這樣的時候.
那就非同小可了.
有大光向你顯現.
光到一個地步你知道.
而她叫你的名字.
我是一個小妹妹.
從來沒有人對我有一個尊稱.
現在不是了.
這樣殺來的使女下格.
你在哪裡來的.
你要去哪裡.
這樣的切問問題.
其實是引導她內心說話.
說出她內心的痛.
所以她能夠說的就是.
我從我的主母殺來的面前逃跑出來.
是逃跑.
不是出來散步.
不是出來默想.
是逃跑出來.
就是說很清楚.
她已經知道自己.
一方面也知道自己曾經的過失.
做了一些舉動.
令到她的主母.
很大的憤怒.
而對她做出這些事情.
她逃走出來.
但現在可以去哪裡呢.
誰可以解救我現在的困局呢.
誰能夠理解我內心.
我過去的苦.
導致我今天要面對的處境呢.
誰可以幫到呢.
親愛的弟兄姊妹.
真正的信仰就在這裡出現了.
真正的信仰就是.
她碰到你的心靈的最深處.
而你自己真是很有意識地.

$^{601}$知道她 知道你的內心.
這個才是真正的信仰.
下甲知道有了一位很好的神.
因為祂就在亞伯蘭的家.
看著自己的主人和主母的敬虔.
感受到他們真是信神的.
都感受到一些靈氣.
不過似乎對她自己還是很遙遠.
不過當來到這裡的時候.
居然這位使者問我問題.
她有什麼出路給他呢.
我逃跑出來又怎樣呢.
第九節就很有趣.
你回去伏在你的主母的腳下.
還以為是真知灼見.
原來是回去伏在她的腳下.
可能不是藤條拔豬肉.
可能是關進柴房三天不給你吃飯.
可能更大的虐待.
打掉你的胎兒.
都有可能.
信心之母都可能會發脾氣.
有很大的苦毒.
對一個女性來說.
搶佔自己最重要的丈夫.
和自己的身份的時候.
是一個很大的屈辱.
可能會殺人都會.
但是耶和的使者就告訴她.
那個解救的辦法就是.
你在哪裡跌倒.
就在哪裡起來.
回去.
當然附帶的條件.
就是灌注一個應許的盼望.
在下甲的心靈裡面.
最快的就是明年她生兒子.
幫她的兒子.
說明名字叫以實瑪利.
是神聽見你的苦情.

$^{641}$接著下甲.
全世界第一個人幫神改名.
替神改名.
一個小妹替神改名.
你看見我.
所以就叫比爾拉海萊.
就是神看見我.
是這個水泉的意思.
所以那個地方.
本來那個泉不是叫這個名字.
以後就叫比爾拉海萊.
是下甲改的.
下甲為神命名.
因為她有盼望.
因為她開始經歷到.
這位個人生命的神.
她要面對.
可能是如下.
回家之後.
繼續會被苦待.
但是內心深處.
神的光 神的盼望.
神的同在.
在她的內心.
她有力面對那個環境.
不再是.
半天之前的下甲.
原來突然間生命的扭轉.
是因為神遇見她.
親愛的弟兄姊妹.
剛剛年頭.
第三個星期.
第三個星期六晚.
第三次崇拜.
我們有沒有一個決心.
虛心.
神啊 我真的希望.
每一晚來到.
每一個星期.
我能夠遇見你.

$^{681}$你遇見我.
你知道我現在的內心的苦情.
你進入我的心靈裡面.
讓我能夠有力量.
面對我如下.
仍然一樣的困苦.
不同了的就是.
有神跟你一起去面對.
她幫神改名.
然後回去服在主母的手下.
第十五節.
這個後來.
就是說十月懷胎了.
幫亞伯蘭生了一個兒子.
亞伯蘭也幫他的兒子改名.
叫做以實瑪利.
當時就是八十六歲.
上文說的是八十五歲.
十年.
現在終於能夠有一個盼望出現了.
下鑒仍然是一個小妹.
在廿一章再有跑一次.
因為撒萊又要趕走他.
這次連他的兒子都要趕走.
不過仍然告訴我們.
下鑒的經歷.
是每一個受困苦的人.
我相信不僅是女性.
女性如果在一個男性的社會.
那個壓迫.
那個被隱藏.
是更加嚴重的.
今天當然是一個男女平等.
相當平等的時代.
始終我們不是亞伯世界的地方.
你要蒙頭 你要戴黑鯊.
但是或多或少.
你的成長.
你的家庭.
你的人生 你的事業.

$^{721}$未必真是如你所願的.
或者你真是內心有很多困苦的.
我們的信仰告訴我們.
做物主是聽見你的苦情.
他能夠看到你的需要.
他能夠在你仍然面對的事情上.
與你同在.
這個很需要我們個人自己.
來有經歷對上主的經歷.
以致我們的信仰.
不是知道這個世界有神.
而是你知道這個世界的神.
與你有生命的關係和接觸.
這個就是最重要的.
今晚藉著下甲這位主角.
一個女孩成為主角.
告訴我們原來她都可以遇見神.
很盼望神給我們每一位弟兄姊妹.
今天晚上我們能夠在神的恩典裡.
來有經歷.
我們一起祈禱.
今天晚上將我們每一位的心靈.
仰望在你的手中.
因為你是向人顯現的神.
你曾經向下甲.
顯現為那位看見人的神.
你聽到他的苦情.
你也曾經.
主耶穌在罪加的井旁.
與撒瑪利亞的婦人.
是一個真實被當時的人所唾棄的婦人.
主耶穌就近她.
進入她的心靈.
改變她的生命.
很盼望我們的生命都能夠得到.
主自己的更新和改變.
奉主耶穌基督的名.
阿們.
\newpage



\section{創世記 11:1-9}
\label{sec:DTr4fwkGAak}
\textbf{2024-01-21 午堂《從神的創造心意再思教會的大使命》創世記 11\_1-9 講員:楊穎兒姑娘}
\newline
\newline
連結: \href{https://youtube.com/watch?v=DTr4fwkGAak}{\texttt{https://youtube.com/watch?v=DTr4fwkGAak}} ~~~~ 語音日期: 2024-01-23
\newline
\newline
\hyperref[sec:_CjhJKqUcA8]{\small{< < < PREV SERMON < < <}}
~
\hyperref[sec:index_chronic]{\small{[返順時目]}}
~
\hyperref[sec:index_scriptual]{\small{[返順卷目]}}
~
\hyperref[sec:pNQz6hqXFi8]{\small{> > > NEXT SERMON > > >}}
\newline
\newline
創世記 11:1-9
\newline
\begin{longtable}{cl}
\hline
\hline
章節 & 經文 (和合本修訂版)\\
\hline
11:1 & \begin{tabularx}{0.7\textwidth}{X} 那時，全地只有一種語言，都說一樣的話。 \end{tabularx} \\ \\ \relax
11:2 & \begin{tabularx}{0.7\textwidth}{X} 他們向東遷移的時候，在示拿地找到一片平原，就住在那裡。 \end{tabularx} \\ \\ \relax
11:3 & \begin{tabularx}{0.7\textwidth}{X} 他們彼此商量說：「來，讓我們來做磚，把磚燒透了。」他們就拿磚當石頭，又拿柏油當泥漿。 \end{tabularx} \\ \\ \relax
11:4 & \begin{tabularx}{0.7\textwidth}{X} 他們說：「來，讓我們建造一座城和一座塔，塔頂通天。我們要為自己立名，免得我們分散在全地面上。」 \end{tabularx} \\ \\ \relax
11:5 & \begin{tabularx}{0.7\textwidth}{X} 耶和華降臨，要看世人所建造的城和塔。 \end{tabularx} \\ \\ \relax
11:6 & \begin{tabularx}{0.7\textwidth}{X} 耶和華說：「看哪，他們成了同一個民族，都有一樣的語言。這只是他們開始做的事，現在他們想要做的任何事，就沒有甚麼可攔阻他們了。 \end{tabularx} \\ \\ \relax
11:7 & \begin{tabularx}{0.7\textwidth}{X} 來，我們下去，在那裡變亂他們的語言，使他們彼此語言不通。」 \end{tabularx} \\ \\ \relax
11:8 & \begin{tabularx}{0.7\textwidth}{X} 於是耶和華使他們從那裡分散在全地面上；他們就停止建造那城了。 \end{tabularx} \\ \\ \relax
11:9 & \begin{tabularx}{0.7\textwidth}{X} 因為耶和華在那裡變亂了全地的語言，把人從那裡分散在全地面上，所以那城名叫巴別。 \end{tabularx} \\ \\
[1ex]
\hline
\hline
\end{longtable}
$^{1}$弟兄姊妹平安.
讓我可以在今天的早上.
在母會錦繡堂.
跟大家分享主的華語.
很高興可以參與教會現在的補位.
以及謹守在我做宣教崗位上.
繼續和弟兄姊妹.
特別是去分享關注大使命方面的訊息.
上個月我有機會在嘉欣堂.
在一個三語共融的崇拜裡事奉.
完結之後有一位弟兄走過來.
姊妹你的中文很流利.
你是不是非語的肢體.
於是我便說.
我也是回錦村家.
但我不是回非語的崇拜.
我是在本地做非洲人的宣教事工.
弟兄說.
很配合你的膚色.
很適合服侍非洲人.
讓我今天可以用這樣的方式來介紹自己.
也會在今天的主題.
讓我們先來禱告.
親愛的神.
讓我們今天聚集在一起的時候.
求你在我們當中同在和引導.
也藉著我們去探索巴別塔這個故事.
裡面充滿著你創造的智慧.
求你賜福孩子的口.
讓你的智慧可以比我們各人都清晰.
藉著你的智慧繼續去生活.
我們這樣奉主耶穌基督的名而求.
阿們.
大家都知道.
錦村家近來有一些事工的轉變.
所以在我預備經文的時候.
我也想和大家一起去深化一下.
教會如何去肩負神給我們的大使命.
如果要說大使命的話.
為什麼今天的經文不是《馬太福音》第28章19節.

$^{41}$這段經文大家都很熟悉的.
所以你們要去洗萬民作我們的門徒.
奉父子聖靈的名.
給他們施洗.
為什麼偏偏要選擇巴別塔的經文去說呢?.
如果像我們的題目這樣說.
是從創造的心意去看大使命的時候.
創世記也有很多經文.
有亞伯拉罕的故事.
有約瑟的故事.
為什麼偏偏要選擇一個巴別塔的片段.
就好像是神很憤怒去懲罰人的片段.
去再去說大使命呢?.
巴別塔和大使命有什麼關係呢?.
在我們過往的認識裡面.
巴別塔的故事.
好像是要讓我們後世的人引以為鑒.
人的功勞自大.
很想建築一座城一座塔.
而且塔是塔頂通天.
人以為自己的能力.
是可以和耶和華神相比.
藐視了神在大地的創造的權能.
而引起耶和華極大的不滿.
於是懲罰人類.
打散他們.
混亂他們的語言.
使他們無法溝通.
不能夠再建好這座塔.
我們就要問.
巴別塔在創世歷史裡面.
那個重點是不是懲罰人類呢?.
在今天的早上.
很想和大家進入創世的世界裡.
在神的創造心意裡.
我們再去思想教會的大使命.
是如何一起去承擔.
我們會從不同的角度.
去看巴別塔的重點.
是不是懲罰呢?.

$^{81}$第一.
我們從故事的脈絡去看.
要這樣形容我們的神.
為一個審判的神.
沒錯的.
因為從第一章到第十一章.
記錄的是人不斷墮落.
以至神要用洪水來消滅整個人類.
到了十二章開始.
就是亞伯拉罕信心之父的故事.
如果神作為審判的神.
這個角色開始模糊了.
因為下一次.
神再對神子文的審判.
可能要數到真正的被擄時期.
第二.
有人說巴別塔的審判.
是對著人類的驕傲.
我們就可以從原文的字義.
去看人類的驕傲是怎樣呢?.
第二節提到.
場景發生在一個叫士拿的位置.
當時是羅亞的後裔向東遷移.
其實編排這個佈景.
其實都知道是不好的事.
因為向東行是不好的.
該人殺了阿伯之後.
該人是向東行的.
要向東遷移.
羅德和亞伯拉罕分地的時候.
羅德也是向東行.
選擇了索多瑪.
阿摩拉.
這個都是麻煩的.
而且在士拿這個地方.
昔日的巴比倫.
今日的伊拉克.
這個地方是兩河交匯.
兩條河交匯.
其實河水不是一致的.

$^{121}$其實已經是襯托著.
當時人類和上帝的心意.
是不偏離的.
有一些切斷的感覺.
所以我們如果用電影去拍這一幕的話.
一定是陰陰暗暗的.
是不合神心意的.
這個鋪排已經是展示著.
人類產生的一些問題.
已知道神要降臨.
第五節.
耶和華降臨.
要看世人所建造的成和塔.
一看.
塔頂通天.
而且當時他們造磚的技術.
是非常成熟的.
而且他們是非常同心.
第三節.
他們說我們彼此商量.
來吧我們一起造磚.
還要將磚燒透.
他們就將磚拿起來造石頭.
然後又要將石屎黏在一起.
很有計劃.
不單止有計劃.
而且是人類是同心合意.
不要小看這個同心合意.
同心合意是可以很厲害的.
因為當神見到人的同心.
那時候他有一個這樣的說法.
第六節.
現在他們做什麼都可以.
沒有什麼可以阻止他們.
當時全地就是一種語言.
說一樣的話.
大家同心.
一起造磚.
一起建築.
建築那座城.

$^{161}$那座塔.
成為一個民族.
說一個語言.
這是代表著人類文化建築科技.
發展的興旺.
很多人馬上就聯繫到.
這麼厲害.
人類驕傲自大.
那是不是那個塔太高.
就代表人驕傲呢.
如果沒有那麼高.
是不是沒有那麼驕傲呢.
如果沒有那麼高.
是不是神沒有那麼生氣呢.
我們放在今天來看.
如果我YouTube有很多Followers.
或者我們的生活是充滿著AI的幫助.
那是不是代表人類驕傲自大呢.
或者是不是一件壞事呢.
更奇怪的是.
其實經文沒有一個字是講驕傲的.
第三.
可能你說.
不是啊.
第四節.
我們來看第四節.
他們為自己立名.
目標不正確.
應該傳揚神的名.
不是自己的名.
但是請大家留意.
我們現在在一個創世的世界裡面.
在頭十一章的創世故事裡面.
神是沒有叫人要傳揚神的名.
在創世以來.
神的尾線就是在大地的那種Shalom.
那種平安.
那種完整.
完美.
萬物都按其善.

$^{201}$按神的道.
按其規律去成長的時候.
已經在無言語.
無聲音的之下.
去傳揚神.
為自己立名.
是不是也是那麼邪惡呢.
如果我們來看看.
為自己立名.
原文或者英文.
我們會看到Make a name.
意思就是帶著有做.
有工作.
要製造.
去做的意思.
這個做字其實是為自己立.
即是Make.
其實是和神創天造地的做.
是同字根的.
然後我們又可以問一下.
既然人能夠去創造.
去工作.
這個不是神給我們的形象嗎.
人按著神的心意去做.
去工作.
又有什麼不對呢.
我們再看頭十一章的聖經裡面.
第四章 第五章 第十章.
甚至是說完巴別塔之後的經文.
都是有眾多的名字.
是家譜來的.
如果名字不重要.
我們還要頭十一章.
不需要花那麼多篇幅去說名字.
再說.
如果我們去到下一章.
第十二章.
第二節的時候.
亦都是看到耶和華神親自主動的.
向亞伯拉罕賜下祝福.

$^{241}$而其中一個祝福就是使祂的名為大.
所以有些學者就說.
其實Make a name.
為自己立名.
不是什麼彌天大罪.
我們就發現.
原來巴別塔在聖經的創世歷史重點.
未必是真的說懲罰人類.
又未必是說人的驕傲.
亦都不是指責人沒有傳神的名.
我們就要問.
聖經說什麼.
我們看回經文.
重要的事是會說三次.
第四節.
他們說.
讓我們建造一座城和一座塔.
塔頂通天.
我們要為自己立名.
免得我們分散在地上.
於是第七節.
耶和華降臨.
變亂了他們的言語.
使他們言語不通.
第八節.
耶和華使他們從那裡分散到全地.
他們就停止建造那座城.
因為耶和華變亂了傳遞的語言.
把人從那裡分散到全地.
所以打成巴別.
重要的事是要說三次.
剛才的經文.
耶和華神使他們.
要從他們在事拿的地方.
要分散到全地.
那時候的人.
知不知道神的心意.
知 當然知道.
不知不懂說.
免得我們分散到地上.

$^{281}$是這群羅亞的後裔他們說的.
知的 不做而已.
知而不做.
為何知而不做.
學者發現.
當時的傳統.
是有一個紀念死人的傳統.
而紀念死人我們都知道.
是要後世的人去紀念.
當人聚居的時候.
可以對先人.
有一些紀念.
甚至是一些祭祀的儀式.
但當人一分散在地面的時候.
問題就來了.
誰去紀念死人.
在我們的生活裡.
我們也有一些異教的信仰.
在我們的世界裡.
他們也有類似的祭祀方式.
人所關心的問題.
都是死後的生活.
有沒有人理.
如果我說到這裡.
大家都很明白.
自從創天造地的時候.
神就是要他們不要聚焦在死後的世界.
而是要活在當下.
要分散到全地.
所以問題不是那座城.
或那座塔有多高.
或是人的文明.
如何令人驕傲自大.
而是建築城和塔背後的目的.
那時候的人是免得我們分散在地面上.
我們就問.
經文有幾時說過.
分散到地面是神的心意.
如果我們回到創世記.
一章二十八節.

$^{321}$大家會記得.
神賜福給亞當和夏娃對他們說.
要生養眾多遍滿這地.
而在第九章.
神也賜福給盧亞和他的兒子.
說你們要生養眾多遍滿這地.
可能有人馬上想到.
生養眾多有點困難.
或不是當下的事.
但其實這是中文的譯本.
如果我們看英文的話.
英文是Be fruitful and multiply.
這是生養眾多.
Fill the earth 遍滿全地.
這個概念和分散到全地是有關的.
Fruitful 大家看到.
Fruit 即是果子.
和果子有關.
我們住在錦繡花園這個社區.
可能也有種植的習慣.
我們會種不同的果子.
會有大有小的果子.
果子要種得好.
要有光澤.
裡面要有營養.
又沒有蟲蛀.
我們會按時按季節去結果子.
而且一棵植物不會只出一粒果子.
總會結出很多果子.
如果我們要數聖靈的果子的時候.
都有很多特質.
人愛喜樂和平.
忍耐恩慈良善信實.
溫柔節制.
而當我們去看眾多的概念的時候.
這個概念是說由一而出.
由一變成很多.
這是神的心意.
人類從亞當一人而出.
神賜福給他們.

$^{361}$洪水過後.
人從羅亞一家而出.
神也賜福給他們.
要他們生養眾多.
耶和華呼召亞伯拉罕的時候.
也是由一族變成大國.
而且萬族也要因他得福.
當我們看到整個人類的救贖歷史的時候.
是神差派他的愛子成為人.
從一人的死而傳遞得到救贖.
大家看到嗎?.
這些是通的.
救恩已經完成了.
然而其實在世界上仍然有很多人不認識耶穌.
所以教會信徒是有責任繼續去秉承大使命.
向外伸展.
在萬民作我的門徒.
奉父子聖靈的名.
給他們施洗.
是散到全地的.
遍滿這地.
我們宣導會是一間帶著宣導會的DNA的教會.
繼續發展下去的.
是有一個去承擔大使命的心意.
教會過往的日子也是透過禱告.
參與和支持.
在這方面進行進行.
宣教士會去到前方.
將神國的福音.
或是看到的事.
或是遠方面對的屬靈需要.
帶回後方教會的群體.
讓教會的群體一起去守望.
紀念.
奉獻.
支持.
宣教就是要我們遍滿這地.
人人都可以參與宣教.
神的心意就是我們要散到全地.
這個概念就是要.

$^{401}$生養眾多.
遍滿全地.
Be fruitful.
然後multiply 變多.
我們可以從個人的靈明生活開始.
靈修很重要.
但考驗的是我們每天的群體生活.
由我們起床到睡覺的那一刻.
我們是否也有這些果子的特質呢.
再細分一點.
由我們起床見到家人.
我們見到我們的家庭傭工.
我們走出去見到我們的鄰居.
我們自己的個人習慣.
我們在職場在不同的崗位.
我們是怎樣的呢.
不要小看我剛才說的每一個位置.
如果在座有300人的話.
一年有365天的話.
每人用10小時.
每小時一次去顯示這些果子的話.
這裡已經有100萬次去展示神的那種美麗.
在職場上我們也想想.
其實結果子是怎樣呢.
結自己的果子.
結公司的果子.
我們看到甚麼果子.
大家可以互相影響呢.
或者在一個時間的空間裡.
怎樣去傳承交接.
這些果子怎樣延續下去呢.
在教會層面的時候.
當我們看到分散去傳遞的時候.
就是在說我們有沒有做社會工作呢.
我們有.
我們也會關注到本地弱勢社群的需要.
我們會做本地跨文化的工作.
也不只是顧本地.
我們會顧大地.
這個這地.

$^{441}$另外一個地方的.
去接觸不同的群體.
去猜派宣教士去海外宣教.
神的心意就是要我們去分散到全地.
雖然明明知道要分散到全地的.
但偏偏就去做那座城那座塔.
而不被神去分散去傳遞.
在這一點惹了耶和華的不滿.
耶和華叫人向橫向多元發展.
但如果一味向高.
成為一個很內在的.
大家都住在那個城市.
向高其實都不是問題.
但問題是他們不向外.
神的整個創造.
是有大地有其他地方的.
如果不向外的話.
由誰去關心管理遠方的事呢.
當神是交了管理的權力.
給人去做的時候.
但人又沒有按神的心意去做的時候.
這個是不馴服.
當不馴服的時候.
無論多同心都是枉然.
當神是要藉著一個.
要人散滿全地.
要這個福賜予大地和地上生活的人的時候.
人的不馴服會怎樣.
會為神帶來危機.
因為神說了的祝福.
就有阻礙.
不能夠將這個福賜予大地.
如果有這樣的危機.
神要不要出手.
要啦 需要啊.
第七節.
我們去變亂他們的語言.
讓他們彼此不通.
然後就要他們分散.
不再做那個城.

$^{481}$變亂了很多語言.
分散在全地上.
耶和華創天造地的時候.
造的植物不是一棵.
造的動物也不是一款.
而是多種類的.
我們在座的弟兄姊妹.
你看看你旁邊的弟兄姊妹.
大家五官都不同的.
如果神的創造心意.
是隱藏著這麼多元的元素.
我們就問巴別塔的故事.
最終變亂了全地的語言.
由一個語言變成很多語言.
我再問一次.
巴別塔是在說懲罰人類.
還是要延續神賜福給大地的心意呢.
也有學者從巴別這個字的字根.
因為都是變亂的意思.
去看這個字根可以有三個解釋.
可以是變亂 可以是混合 可以是調和.
這個字根在聖經出現了43次.
兩次是用在變亂.
也是用在這段經文裡面.
有32次其實是用調和的意思.
就好像當時做無糠餅.
要加些油去調和整個餅.
人的不信服.
使到神 也就是說神要出手.
去調和這個語言.
使人能夠分散到全地上繼續的.
生養眾多 遍滿這地.
上星期日我在教會.
參與了一個錦宣家.
為杜牧養士宮有轉變的差遣禮.
看到三個堂點.
其實都帶著自己的特色.
而是各自帶著不同的特色.
在社區服侍.
這樣的變動.

$^{521}$我們會看到神在做什麼工作呢.
我想大家心中都很期待.
錦宣家繼續在社區做一些新事.
教會的大使命.
會不會就是讓我們一起去學習.
去擁抱那種不同.
擁抱那種多元.
其實我們可以有一些練習.
去幫助我們開啟這個感官.
或者認識更多這種智慧.
我們可以首先從自我檢查開始.
可以看我們從我們的成長裡.
會不會有一些對某些種族.
言語文化有些愛恨.
或者是會有一些固有的偏見.
或者對某些群體的定型.
或者覺得很難接納某一類人.
而在我們日常生活的社交裡.
我們願不願意去開放.
去認識不同種族的文化.
不單止是接納.
不是一方面的接納.
但我們在這個和而不同的裡.
我們可不可以堅持緊守.
維護我們自身的信仰價值或者道德.
而我們能不能夠由家庭去做起.
去接納不同種族言語文化的人.
成為家人.
或者會不會將家庭的傭工.
也當作我們的家人.
而在教會的使命層面裡.
教會的活動可不可以有不同的種族.
不同的言語.
不同文化的人一起參加事奉.
我們可不可以在言語上調和一下呢.
這個很值得我們繼續去深思.
當教會同心.
而且目標較正.
分散在地上.
生養眾多偏門者地.

$^{561}$神必然在當中去賜福教會.
而且是賜福更多在地上的人.
早前教會一直都有一些.
例如唐慶.
或者三語的共融的崇拜.
在當中我們也會見到.
不同國籍的弟兄姊妹.
在特別唱歌跳舞的時候.
會跳舞.
會擺動.
會揮手.
會歡呼.
很不同的.
整個氣氛和我們平時崇拜不同.
可以調和了.
我們也會想的.
我們會想動一下.
最後我想用一個故事作結.
在九龍有一間教會.
四十年前.
一班會眾領受了.
要發展服侍智障人士的事工.
至今已經四十年了.
每年他們都有一個共融的崇拜.
在這個崇拜裡面.
有教會恆常的會友.
也有服侍的一班智力有問題的群體.
從他們的院舍來崇拜.
一起參與.
崇拜很特別.
崇拜不是很靜.
會很嘈吵.
因為當中會有一些新的朋友.
或是未見過面的朋友.
他們控制不了自己.
會發出一些聲音.
有時會有人說.
不知道崇拜說什麼.
很嘈吵.
不能集中.

$^{601}$但是我們會不會可以.
有一個眼光去看看.
他們都是神的創造.
在言語裡面.
其實是跟他們溝通不了的.
打亂了言語.
但是我們看著他們.
就看到他們是神的創造.
有一天.
我帶了一個弟兄去這個教會崇拜.
我不知道那天.
原來是這個共融崇拜.
就是跟一班智障人士共融崇拜.
我不知道.
我帶了那位弟兄去.
那位弟兄也帶了他的小朋友去.
其中一位弟兄的小朋友.
也是有智力有問題的.
這個弟兄一家.
就是第一次這樣去上教會.
他從佛教轉到基督教.
他第一次上教會.
他感受到教會裡面那種包容.
那種愛.
那種接納.
後來這位弟兄也繼續上教會.
也很感恩.
他其實是在感受堂那裡接受洗禮的.
當我們有這個能夠去接納.
能夠去看遠一點.
去看看原來神的創造.
不單是多元.
而且在缺陷裡面.
也是可以很美的時候.
我們今天.
我們會不會都一起轉念.
去用神的眼光去看這個世界.
我們不要做一個聽了.
但不懂得做的人.
讓我們一起同心同意.

$^{641}$親愛的父神.
我們真是來到你的面前.
你的話語歷久常新.
滿有智慧.
你帶著我們去一個不同的角度.
去看看教會.
社區是可以怎樣多元的.
我們相信你是走在我們.
引領我們的神.
你是從一位耶穌基督裡面.
將救恩賜下.
求你真是教導我們.
如果我們真是有知道而不做的.
我們求你去寬恕.
我們求你聖靈.
那種微聲.
是滿有你的憐憫.
你的恩典.
是繼續帶領我們教會.
去作新事的.
而在我們生活的每一日.
每一刻.
我們求你去轉化我們.
讓我們活出滿有智慧.
滿有力量.
就讓我們一起去擁抱.
神理次下這個大使命.
就是要我們去認識多元.
不是在我們的頭腦.
但是給我們能力去更新.
去改變.
去打破固有.
不合理心意的思想.
在我們將要完結這個崇拜的時候.
也都求你將這個福份.
賜給教會信徒.
我們這樣奉主祈求.
阿們.
\newpage



\section{創世記 16:8-14}
\label{sec:pNQz6hqXFi8}
\textbf{2024-01-21 早堂《憐憫人,讓人看見的神》創世記 16\_8-14 講員:鄧泰康傳道}
\newline
\newline
連結: \href{https://youtube.com/watch?v=pNQz6hqXFi8}{\texttt{https://youtube.com/watch?v=pNQz6hqXFi8}} ~~~~ 語音日期: 2024-01-23
\newline
\newline
\hyperref[sec:DTr4fwkGAak]{\small{< < < PREV SERMON < < <}}
~
\hyperref[sec:index_chronic]{\small{[返順時目]}}
~
\hyperref[sec:index_scriptual]{\small{[返順卷目]}}
~
\hyperref[sec:jLSZtQ5_sq4]{\small{> > > NEXT SERMON > > >}}
\newline
\newline
創世記 16:8-14
\newline
\begin{longtable}{cl}
\hline
\hline
章節 & 經文 (和合本修訂版)\\
\hline
16:8 & \begin{tabularx}{0.7\textwidth}{X} 對她說：「撒萊的婢女夏甲，你從哪裡來？要到哪裡去？」她說：「我從我的女主人撒萊面前逃出來。」 \end{tabularx} \\ \\ \relax
16:9 & \begin{tabularx}{0.7\textwidth}{X} 耶和華的使者對她說：「你要回到你的女主人那裡，屈服在她手下。」 \end{tabularx} \\ \\ \relax
16:10 & \begin{tabularx}{0.7\textwidth}{X} 耶和華的使者對她說：「我必使你的後裔極其繁多，多到不可勝數。」 \end{tabularx} \\ \\ \relax
16:11 & \begin{tabularx}{0.7\textwidth}{X} 耶和華的使者又對她說：「看哪，你已懷孕，要生一個兒子。你要給他起名叫以實瑪利，因為耶和華聽見了你的苦楚。 \end{tabularx} \\ \\ \relax
16:12 & \begin{tabularx}{0.7\textwidth}{X} 他為人必像野驢。他的手要攻打人，人的手也要攻打他；他必常與他的眾弟兄作對。」 \end{tabularx} \\ \\ \relax
16:13 & \begin{tabularx}{0.7\textwidth}{X} 夏甲就稱那向她說話的耶和華為「你是看見的神」，因為她說：「他看見了我之後，我還能在這裡看見他嗎？」 \end{tabularx} \\ \\ \relax
16:14 & \begin{tabularx}{0.7\textwidth}{X} 所以這井名叫庇耳‧拉海‧萊，看哪，它位於加低斯和巴列的中間。 \end{tabularx} \\ \\
[1ex]
\hline
\hline
\end{longtable}
$^{1}$各位姐妹早安.
今天和大家分享一個.
這段時間我們都會分享創世記裡面的人物.
今天分享一個比較特別的人物.
這個人物其實可以說在創世記當中是一個很小的角色.
一個小角色.
夏鴿是一個怎樣的人呢?.
如果我什麼都沒有說之前.
我的結論是他是一個困苦的人.
他是一個悲情的人.
他是一個完全被動的人.
你會發覺這個角色特別的地方是什麼呢?.
上帝是一個其實沒有主動權的當中的人.
他怎樣去彰顯他自己的榮耀和能力.
夏鴿是亞伯拉罕的一個埃及的奴隸.
即是埃及人.
所以其實聖經裡面主要去記載的.
我形容是在亞伯拉罕的故事當中一個插曲.
而他也都經歷了兩次和上帝他的使者相遇.
所以他的故事其實很值得我們參考.
為什麼呢?我會這樣說.
如果大家看聖經的時候你會發覺.
有些人都會問這個問題.
聖經好像一路說.
一路說這個世界很多創造很多.
後來一直去到亞伯拉罕之後.
就一直集中一條線繼續發展.
大家有沒有發覺?.
就只是去到以色列人.
由他們這裡一直去發展.
他基本上沒有怎麼說其他的.
多數都只是說外邦的軍隊來打仗.
所以你會發覺其實在聖經裡面去看的時候.
有主線和非主線.
其實大家要這樣看.
我們這班中國人其實是非主線.
不是在聖經當中的主流.
我想形容.
所以其實夏革在她這個故事當中.
也都不是主流.

$^{41}$但是我就是想告訴大家.
我們所信上帝奇妙的地方就是.
他是主流的上帝.
也都是非主流的上帝.
這個故事是怎樣的呢?.
我要先當有很多頂子妹都沒有聽過這個故事.
大概詳細和大家去分享.
究竟發生了什麼事.
上帝去應許了阿巴拉罕.
說他將會有子女.
然後時間呢.
聖經裡面第十六章就說.
撒萊套就有一個僕人.
一個奴僕就叫做夏革.
撒萊就走來和阿巴拉罕說.
上帝令到我沒得生.
我想大家留意這句話.
上帝令到我沒得生.
所以我現在有一個奴僕.
我就想他和你做妾侍.
好讓他在他身上我得到兒子.
所以你就會發覺.
跟著阿巴拉罕就和他同房了.
大概是這個事.
你會發覺撒萊的言詞當中.
其實是責怪耶和華的.
你又告訴我老公.
他將來有一個大族.
但是你想像一下我這麼多年.
當時他年紀都不小.
這麼久都沒得生.
現在的問題就是人的心情很急.
我不相信上帝的奇妙.
我不相信上帝你自己說過的話.
包括其實之後阿巴拉罕都不相信.
所以他們是試圖用他們的方法.
走去完成上帝的應許.
他們很想要上帝的應許.
很急於要上帝的應許.
所以才會發生這件事.

$^{81}$跟著他就和他說.
你就和他生個兒子吧.
古代的世界或當時的世界.
為什麼這麼重視生兒子呢.
因為生兒子和我們今天.
其實在香港的圍村差不多.
都是兒子才可以繼承一些東西.
所以生兒子對他們來說是很重要的.
所以對於他們來說.
也是一個女人價值的指標.
在香港現在慢慢.
這個社會我們不太感覺到.
有時我們說生女兒是好過.
特別是做爸爸的.
生女兒是開心一點的.
因為女兒會黏身一點.
會開心一點.
但是在一些農村裡面.
我不知道大家知不知道.
就算今天的中國國內.
都是一樣.
生不到兒子是很大件事.
我媽媽是客家人.
所以我客家人同樣都是.
我媽媽以前做女兒的時候.
是很珍貴的.
大家有沒有聽過一個字叫同養息.
同養息就是說.
因為我家窮.
然後很小的時候.
就將她許配了.
就是送給別人養.
將來就做高家人的老婆.
這就叫同養息.
你幫我養女兒.
將來就讓她嫁給你兒子.
所以我媽媽和我爸爸結婚之後.
其中有一樣東西就是一定要生兒子.
我形容是很幸運.
連生了四個.

$^{121}$所以我上面四個兄弟都是兒子.
之後應該很開心.
生到第五個是女兒.
第一個女兒.
我媽媽現在經常和我說.
她當時很怕女兒哭.
很怕我妹妹哭.
因為我爸爸做的工作.
要早上睡覺.
怕吵著她.
她會怕我爸爸不喜歡這個女兒.
然後我爸爸聽到之後.
說你都傻的.
我怎麼會不疼這個女兒.
但我所謂的我家很反轉.
我妹妹是最兇的.
我爸爸是最重的.
所以我經常形容我妹妹.
是可以站在我爸爸的頭上.
我們個個在爸爸面前是縮起來的.
你就可以想像到.
其實在女性的心目中.
其實兒子的重要性.
另外一樣東西.
她會覺得是羞愧.
如果我生不到的話.
我是很羞愧.
基於這些原因.
所以薩奈提出了這個建議.
原來當時處理母子生.
是有幾個方法的.
一是一夫多妻.
一是納妾.
一是離婚.
所以薩奈採取了納妾.
納妾有什麼不同呢.
一夫多妻.
大家的地位是平起平坐.
但是納妾就好像剛才的圖畫.
當妾侍生了兒子.

$^{161}$是歸妻子的名下.
有這個做法的.
其實你會看到一件事.
其實薩奈是非常聰明的.
我只是要結果.
但我不要有人和我爭.
由始至終薩奈在自己心裡.
是一個很重要的標準.
所以你會發覺.
為什麼我說夏革是一個很悲情的人呢.
第一他是奴隸.
他不能死.
就是說我的主母.
她叫我和她同房.
我就要和她同房.
而我也很清楚.
生了兒子之後.
那個兒子不是我的.
其實我不知道大家.
如果我用一個不太好的字眼.
就是薩奈真的當了夏革是一個貨物.
或者是一個工具人.
我這樣形容.
你幫他完成一些任務.
所以你就可以想像到.
薩奈對夏革不會好得去哪.
所以你就會明白.
當他們同房.
上帝好玩的地方就是這樣.
當他們同房.
然後就一現即中.
然後就懷孕.
你就會發覺聖經裡面.
夏革見到自己新有人.
就小看他的主母.
其實我想告訴大家.
是很正常的.
你明白我的意思嗎.
我被人這樣去對待.
接著的時間.

$^{201}$我現在可以終於吐氣揚眉.
在我的主母面前的時候.
我還也問也沒有.
明白嗎.
其實這些是很正常的.
如果大家有看那些.
我不知道是真是假.
早一段時間.
那些後宮的爭鬥戲.
就是全部這些.
生完之後就很威風.
類似.
你就會發覺.
其實夏革去到這裡的時候.
你就會發覺.
他有這種反應.
好了.
然後薩奈就去和阿巴拉罕投訴.
他說.
他這樣對我.
我要怎樣做.
阿巴拉罕的反應是.
你想怎樣就怎樣.
其實去到這裡.
我有一點點火.
我想告訴你.
女人就是.
神也是你.
鬼也是你.
你叫你老公和她生孩子.
生了之後.
接著.
他對你有一些面色.
接著你很不高興.
就和老公投訴.
我又形容老公是沒有用的.
你想怎樣就怎樣.
我不知道大家明不明白夏革的心情.
其實很慘的.
你會發覺.

$^{241}$為什麼去到後來的時候.
阿巴拉罕和他老婆去說.
你想怎樣就怎樣的時候.
聖經形容是苦待他.
聖經沒有說怎樣苦待.
但是可以令到一個人逃走.
你可以想像到有多苦待.
我形容他去到這裡的時候.
那種委屈.
接著他去走.
令到他後來的時候.
他感覺沒有出路.
但是最特別的地方.
原來夏革他的名字的原文.
是逃跑.
我想強調原來逃跑.
聖經的文字很有趣.
逃跑就是原本夏革他的命運.
他只可以逃跑.
但特別的地方就是.
德則哲堯說.
使者在曠野的樞爾路上.
水泉旁遇見他.
如果和合本的遇見.
很容易讓大家有個錯覺.
就覺得是剛剛碰到.
其實不是.
大家應該理解為.
就是上帝的使者周圍去找他.
因為圖畫應該是這樣.
所以你會發覺現在這個圖畫就是.
一個爸爸去見到一個女兒.
她受了很多苦.
接著她現在走了.
他馬上要走去找她出來.
接著就出現了剛才大家讀的那段經文.
這個使者就問她.
我想大家留意其中一個字眼.
撒來的使女夏革.
你是從那裡來的.

$^{281}$接著你現在要去哪裡.
夏革的回應就是.
我是從我的主母撒來面前出來的.
然後就跟她說.
你回到你主母那裡吧.
伏在她的手下.
其實去到這裡的時候.
可能我們大家都不太明白.
他逃跑 上帝苦待他.
上帝為什麼你會這樣.
原來當時的世界.
第一樣大家一定要理解.
一個奴隸.
一個懷有新人的女人.
離開了這個家庭之後.
其實我想告訴大家.
是死路一條.
因為她不會再有出路.
所以你就會發覺.
原來這個體為上帝.
很不近人情的原因.
原來是為了要去保護夏革.
讓她和她的兒子.
能夠在阿巴拉罕.
他們這個家族的庇蔭之下.
一路成長.
這個才是當時對她最好的方案.
但是上帝不斷強調的是什麼.
你是撒來的奴僕.
強調她這個身份.
但是上帝反過來.
上帝對阿巴拉罕有一個的應許.
你將來成為大族.
大家記得嗎.
但是去到這裡.
上帝竟然對一個這麼卑微.
外族奴隸的女人.
做同樣的應許.
其實這個是非常破格的事.
其實從這些裡面.

$^{321}$我想大家要學看聖經.
就是看這些地方.
我們的上帝不是像我們今天那樣.
幫我們去設定很多條條框框.
原來上帝是不斷無論在舊約還是新約.
他不斷在做很多破格的事.
在神的眼中.
其實下格的生命同樣都是這麼重要.
和阿巴拉罕和撒來是沒有分別的.
上帝給這個應許給阿巴拉罕.
同樣他都可以給這個同樣的應許.
給一個女人一個奴僕.
這個就是上帝.
不是你自己認為.
你在神的眼中是如何.
是我們每一個人在神的眼中都一樣.
但是反過來你的角色可以不同.
所以你就會發覺.
其實去到這裡的事.
上帝讓他們剛才那段經文.
一段這麼短的經文出現了四次.
看見看見.
我看見上帝.
所以他連井都看見.
我看見.
在這裡很強調的是什麼呢.
上帝讓人能夠.
第一 下格 看見出路.
而他也看見上帝.
人在他的困苦當中.
你能夠看見上帝.
所以以實瑪利是什麼意思.
以實瑪利就是.
神是一位聽見苦情的神.
我想大家要記得以實瑪利這個字.
是上帝幫他的兒子改的名字.
因為之後會再出現.
所以你會發覺.
到後來這個小朋友就順利出世.
其實我想在這裡和大家分享.

$^{361}$上帝創造這個世界原本是公平的.
當上帝創造這個世界的時候.
其實萬事萬物當中.
每一樣東西都是公平的.
動物是公平的.
人是公平的.
所有東西其實是沒有分界的.
沒有分任何東西.
你會看見是什麼呢.
每一個人上帝所給他的.
都是有他不同的崗位.
植物是要繁衍.
然後人要有一個管理動物的工作.
他有他自己本身的工作.
其實不公平是人自己創造出來的產物.
是人自己走去.
你會覺得不公平.
為什麼呢 從哪裡看得到呢.
還記不記得在創世記的時候.
兩兄弟 戈隱和阿伯.
當時他們獻祭.
戈隱獻他的禮物.
阿伯獻他的禮物.
上帝揀選了阿伯的.
我想大家留意.
但是戈隱就覺得不公平.
然後這種的草擬.
就令到他殺了阿伯.
上帝的揀選是他的主權.
弟兄姊妹.
其實上帝讓每一個人.
生在每一個地方是公平的.
但是上帝有權去揀選.
上帝可以揀選以實.
在他的計劃當中.
他沒有揀選以實瑪利.
所以阿夏格他不可以問為什麼.
他不可以說不公平.
因為上帝所賜的福.
其實不是理所當然.

$^{401}$所有的東西.
都是由上帝自己去決定給我們的.
很多時候上帝我們不斷強調.
我們有管理權.
但是我們沒有擁有權.
上帝賜給我們每一個人的東西.
在傳道書下.
我會在監牢開傳道書.
因為很強調告訴大家.
我們有份.
上帝給你每一個人都有不同的份.
你管理好你的份就可以了.
你不要和別人比較.
所以上帝揀選了阿巴拉罕.
以撒雅各成為了他救恩的主線.
但是同樣.
這樣是不是代表上帝不公平呢.
不是的 弟兄姊妹.
夏格原本他是一個奴隸.
但是因為上帝看見苦情.
所以上帝的拯救和上帝的恩典.
就留到他身上.
上帝是由他去給我們每一個人.
我們每一個人其實當要守住.
當刻的我們自己的身份.
所以剛才我說.
上帝和夏格一說的時候.
問他是誰.
撒萊的僕人夏格.
你是從哪裡來的.
他要強調告訴你.
你是僕人.
所以你是理所當然的.
你要回到家庭裡面.
繼續去履行你自己作為僕人的職份.
所以今天我們每一個人.
在神的面前的時候.
我很盼望我們能夠去管好.
我們自己上帝所交給大家的崗位.
我們今天其實你可以是別人的爸爸媽媽.

$^{441}$你可以是別人的子女.
你是別人的朋友.
還有很多是你公司的上司.
你是公司的下屬.
每一個人都應該在你裡面去各盡其職.
你有你自己的職責.
你做好你自己的事.
但是你不可以去侵佔你主人.
所給你的位份.
所以我們今天沒有辦法可以去.
控制我們自己的人生.
但是你有辦法可以好好去管理你的人生.
每一次當我有機會去參與安息禮拜.
特別是去到火葬禮的時候.
我去做禱告的時候.
或者我去做一些免禮的時候.
我都會說一句話.
就是今天我們每一個人仍然生存.
我們仍然有我們的生命.
其實無論你過去是怎樣都好.
你都可以去選擇去重新建立你的關係.
無論是夫婦關係.
無論是朋友之間的關係.
無論是子女和爸爸媽媽之間的關係.
你是可以重新去選擇怎樣去過好它.
而令到有一天當我們完成了離開這個世界的時候.
變得我們不會有遺憾.
在這裡的時候雖然聖經沒有再去說.
究竟阿霞鴿回去之後有沒有繼續被苦待.
沒有再說.
但是一跳到去的時候.
你就會發覺當阿霞鴿去生了.
當以撒萊去生了兒子的時候.
聖經有一個片段很有趣.
就是當以撒去斷奶的時候.
阿霞鴿的兒子以撒瑪尼.
那時候大概是十六,十七歲.
他在這裡形容是什麼.
見到他玩弄弟弟.
在聖經裡面和他嬉笑.

$^{481}$其實原文裡面有些叫和他玩.
但是我想問一下大家.
一個十六,十七歲的小朋友和剛剛斷奶的玩什麼呢.
同時也會令到撒拉很不高興.
意思就是令到他有一點點的氣援.
有一點點的玩他.
而接著的時間.
撒拉又走去和阿巴拉罕再投訴.
他說你叫他走.
我不想他和我得到產業.
我想告訴大家.
其實撒拉是一個藉口.
第一.
大家還記不記得.
剛剛他講了那個情況.
其實只可以由以撒去承擔.
承繼阿巴拉罕的位份.
但是問題就是.
現在撒萊用了這個理由.
其實根本是不會取代了.
他得到產業的機會.
但是他用這個理由叫他走.
但是又出現了一個很特別.
我想大家都看到很不舒服的地方.
上帝就和阿巴拉罕說.
你讓他走.
這裡有一個很重要的分水嶺.
當時那個小朋友是剛出生.
現在那個小朋友是十六,十七歲.
叫他走意味著什麼呢?.
如果大家看利未記裡面的關於奴隸.
主人是可以主動釋放奴隸的.
意味著這個走.
是給予自由,下格和伊薩瑪尼.
這是一個很重要的扭轉.
為什麼呢?.
因為當他叫他走的時候.
第一件事就是阿巴拉罕很憂愁.
神就叫他走.
從此之後他也會成為一個大國.

$^{521}$你記住.
你的產業是會留給伊薩的.
但是他也會成為里拉大國.
這樣令到阿巴拉罕很安心.
因為上帝看著他.
所以阿巴拉罕就給了.
有一個很容易令我們誤會的.
一皮袋的水給他.
有沒有搞錯?.
一皮袋水.
但其實這裡是告訴大家.
就是說.
下格在當時去這個地方.
要回到自己的鄉下.
不是很遠.
但你記住阿巴拉罕不是傻的.
他給他的糧食,給他的水其實是足夠的.
但是反過來.
為什麼後面又會出現不足夠呢?.
反而大家在這裡要看深層一點.
下格在當時.
其實他在忍辱負重了十幾年在裡面.
他一直都會覺得.
我會繼續去阿巴拉罕的家庭當中.
繼續下去.
我可以得到什麼.
但是突然間他這一刻的時候.
你叫我走.
對於一個女人來說.
他會突然間萬無目的.
他出去了.
他不知道怎麼辦.
他裡面會有很多的不忿.
很多的埋怨.
很多的不明白.
其實這裡你可以形容為下格.
拿到出去的時候.
他不是直接拿著東西就直接回鄉下.
因為他回到鄉下.
其實他也不知道怎麼辦.

$^{561}$但他就是在那裡徘徊.
時間一路過.
他的糧食,他的飲料.
他的水就慢慢地沒有了.
其實他到了那裡的時間.
他看到這些.
所謂叫做他走到那個絕景.
他發覺再沒有希望的時候.
然後他就在那裡痛哭.
其實他到了那裡.
聖經仍然叫做迷了路.
如果我看的時候.
我覺得他反而是心裡面迷了路.
他不知道有什麼出路.
但上帝再一次去捉捧他的時候.
上帝這次跟他說下格.
沒有再叫他殺來的盧埔下格.
是下格.
他說我聽見你的兒子的聲音.
聖經形容是下格在哭.
但上帝說我聽見你的兒子的聲音.
為什麼會這樣說呢.
為什麼不是我聽見下格你的哭聲呢.
大家剛才還記不記得.
黑阿二十萬的名字的意思是什麼.
我漢建嘛.
是上帝幫他改的.
上帝很想他在這裡.
記不記得我當日給你的應許.
你記不記得十幾年前的那個上帝.
你記不記得.
弟兄姊妹當你去面對困難的時候.
你的出路就是.
你要記得.
昔日昔在今在以後永在的那個上帝.
這個上帝以前是怎樣對待你的.
他不會變的.
儘管我們在這一刻見到的是困難.
在這一刻我們見到的是絕路.
但上帝不會變的.

$^{601}$上帝很想他們再一次去看回.
我就是你十幾年前的那個上帝.
而這次有什麼轉變呢.
這次就是.
他讓他見到那口井.
你記住這些全部是有關係的.
昔日他改口井叫做.
我漢建上帝.
這次又讓他看回那口井.
你看不看見我.
弟兄姊妹為什麼今天我去講這個主題.
叫做憐憫人和讓人漢建的神.
今天我們的信仰裡面最重要.
你見不見到上帝.
你漢建嗎.
上帝昔日所給你的東西.
你漢建嗎.
而接下來最重要的是什麼.
剛才說他給了一個自由釋放給他.
自由釋放代表什麼.
他有一個全新的身份.
下鴿和他的兒子終於可以做回自己.
神要你去做你自己.
你就會發覺如果他的兒子.
二十八里一直都在那個家庭的話.
我想告訴你.
雖然他是阿巴拉罕的後裔.
但他的身份永遠都是老婆.
但是他現在出來.
聖經形容他後來成為一個射箭的人.
大家不要看這些這麼簡單的形容.
其實當時很多國家的先祖.
全部都是這些人.
其實這裡就告訴你.
他終於成為一個民族的祖先.
他終於做回自己.
上帝要他活現回他自己的生命.
從一個奴僕原本不能改變的生命.
上帝的保守帶著他一路走.
走到最後一個適合的時間.

$^{641}$上帝帶他離開那個位置.
進入另一個人生的階段.
他能夠成為真真正正一個民族的祖先.
大家知不知道我們的生命.
是否可以這樣轉變.
我想問大家知不知道.
這兩段經文.
剛才那段時間.
最特別的地方就是.
有很多樣樣莫名其妙的東西.
上帝的手法是樣樣莫名其妙.
他受苦了 你叫他回去.
然後他叫他走.
上帝就叫他走.
在那一刻連下格自己都不明白.
他自己為何會這樣.
但上帝就在那時候出現了.
上帝找他.
上帝每一次當人去經歷這些時間.
上帝都會出現.
上帝要你看見他.
看見他最初的英雄.
所以我又回到最原初的那一點.
在座的都是弟兄姊妹.
都是基督徒.
為何你信耶穌.
昔日為何你信耶穌.
很多時候基督徒.
做著做著基督徒之後.
就忘記了最初為何我信耶穌.
最初是否你看見上帝.
在你生命當中給你的拯救嗎.
你見到這位是奇妙的上帝.
他愛你的上帝.
他拯救你的上帝.
所以你願意去決志.
放下自己.
將自己交給他 是不是.
這個上帝從來沒有變過.
上帝很想你再看見.

$^{681}$昔日所給你的英雄.
我答應了你.
我是你的主 我會帶你去走.
未來的每一個部分.
我都會繼續帶你走.
儘管你看不通 你看不透.
但你放心 我是在這裡.
弟兄姊妹.
這個就是我們作為基督徒.
最大最大的福份.
是不會改變.
因為我們的上帝不會改變.
所以我很盼望.
我們的上帝.
是很值得大家花更多的心思.
更多的時間去跟隨他.
我做了傳論這麼多年.
其實說真的.
有時候我自己最心痛的地方就是.
很多弟兄姊妹.
慢慢在他們的人生當中.
經歷了很多的事.
慢慢慢慢.
他的眼光.
他的心思意念.
已經遮蓋了上帝的恩典和大能.
他會看很多事都淡如.
深如止水.
回來後.
都是這樣.
我們的生活都是這樣.
教會都是這樣.
慢慢回來後.
教會的崇拜之後.
只不過變成一個形式.
信仰也不會對我有什麼很大的幫助.
給我一個很大的.
其實弟兄姊妹.
很可惜的.
你要記住上帝的救恩.

$^{721}$最初是救了你的.
是讓你得到新的生命.
上帝給你的應許.
你是可以得到生命.
變成更豐盛的弟兄姊妹.
不要這麼快就放棄.
你就求上帝.
告訴你.
我怎樣能夠過得更豐盛.
所以我長了一點.
我想用一個最後的例子.
跟大家分享.
剛好在星期五.
我們有一個敬拜訓練.
其中有一位.
其實他唱歌非常好.
在敬拜隊的一位姊妹.
他就和我們分享他的經歷.
他是讀音樂的.
他在美國畢業.
然後回到香港之後.
他是一個很乖的教徒.
然後教會就叫他.
我們教會開樂團.
我們找你幫忙.
他很願意事奉.
但我一直做.
很用心地做.
但很大壓力.
在裡面的時候.
我只知道做.
做 做 做.
但他覺得他心裡是空蕩蕩.
事奉並沒有帶給他.
任何生命裡面的激勵.
他只知道我就是做.
但後來他的妹妹.
她妹妹也是讀音樂的.
他妹妹就去了另外一間教會.
他覺得他妹妹很不同.

$^{761}$他形容一個字就是.
我們都是從小到大.
返回教會的人.
但他說他形容他妹妹.
在那一刻他找到耶穌.
他說他妹妹以前和他.
是勢不兩立的.
但反過來.
他妹妹竟然主動走來和他和好.
之後他發覺很不同.
他說為什麼會這樣.
後來他就隔了一段日子.
和教會的人說.
我暫時放棄了事奉.
之後他就好好地.
重新去尋找神.
現在這一刻的這個姐妹.
她在生命當中.
她找到神.
一個姐妹那天晚上和我們分享.
分享之前她不斷祈禱.
但她說她的組長提醒她.
不如用聖經和別人分享.
用神的話語和我們分享.
大概的內容是什麼呢.
你生命有那樣東西.
你才能夠去做到那樣東西.
所以弟兄姊妹.
我很盼望.
我們生命裡面.
都是有耶穌的人.
是真真正正.
我們生命有屬靈能力的人.
我們能夠從裡面而外.
去好好地過我們的生活.
和去事奉神.
讓我們的生命充滿神的喜樂.
好嗎.
我們求你幫助.
求你讓我們在你面前.

$^{801}$我們雖然未必是一個主流的人.
但是在神的面前的時候.
上帝的恩典.
是滿滿地會淋到我們身上.
我們很盼望.
我們藉著上帝所給我們的祝福和恩典.
我們可以去成為上帝.
我們所想過的人.
能夠充滿喜樂 充滿能力.
也都按著你的應許.
我們能夠得到一個豐盛的生命.
你我們靠主耶穌.
得勝名至其金.
好嗎.
\newpage



\section{瑪拉基書}
\label{sec:jLSZtQ5_sq4}
\textbf{2024-01-25 周四晚查經班《瑪拉基書(3)》 講員:曾裔貴牧師}
\newline
\newline
連結: \href{https://youtube.com/watch?v=jLSZtQ5_sq4}{\texttt{https://youtube.com/watch?v=jLSZtQ5\_sq4}} ~~~~ 語音日期: 2024-01-27
\newline
\newline
\hyperref[sec:pNQz6hqXFi8]{\small{< < < PREV SERMON < < <}}
~
\hyperref[sec:index_chronic]{\small{[返順時目]}}
~
\hyperref[sec:index_scriptual]{\small{[返順卷目]}}
~
\hyperref[sec:c0kCL7cehvg]{\small{> > > NEXT SERMON > > >}}
\newline
\newline
$^{1}$感恩我們簡單有一個祈禱.
多謝天父再一次.
雖然寒冷.
都讓我們能夠在家中.
或者在一個很溫暖的地方.
我們可以在教會來到.
有一起查考聖經的時間.
願神來到帶領我們.
能夠藉著《馬拉基書》.
都去明白神你的心意.
也都讓我們可以在其中.
特別當我們在我們自己的.
現在的敬拜生活.
都能夠符合神的心意.
願神帶領我們.
奉耶穌基督的名.
阿們.
不好意思.
不要緊 沒問題.
那.
要不要開一點風扇冷氣.
或者是.
開一點窗戶 看你們自己安排.
謝謝天父.
不要開太大.
可以 先開個窗戶.
要不要開冷氣.
不用 我不用.
好 那我們開始了.
剛剛祈完禱.
《馬拉基書》四章.
我們正式進入經文的討論裡面.
一直都提過.
就是神藉著先知馬拉基.
來到去.
透過一些反問的問題.
好像一些切問的.
先知來提出.
對以色列人.
整個民族的質詢.

$^{41}$這個好像很婉轉的問題.
但是當然附上了一個答案.
而很奇怪地.
讀者 或者先知構想讀者.
那個回應.
就是我們沒有.
你所描述的那個狀況.
那何以見得.
我們好像你所描述的情況呢.
那是連續.
一直延伸去到第四章都有.
當然去到三章尾.
四章是有另外的預言.
就是以利亞的降臨.
以利亞的出現.
就是這一位的先駛約翰.
我們其實用了一堂時間.
上個禮拜就說了.
這個耳掃不被神揀選.
那瓦閣被神去揀選.
似乎羅馬書講的是無條件.
因為兩個兒子都未出世.
還是在媽媽李伯家的肚腹的時候.
他們那些善惡都未作出來.
就是都未出世 未成長.
未受周圍環境的分途污染.
那怎麼會有犯罪的可能.
而遭致神的拒絕.
或者不被揀選呢.
新約的作者就往往告訴我們.
這個正正就是那個預定論的揀選.
當然我們昨晚有機會去聽系統神學.
預定論遲些都會講的.
會有一堂是講預定論.
就是嘉爾文的預定論.
當然是從聖經的歸納.
來探討這個課題.
很多人是不能夠接受.
為什麼會被他揀選.
而不揀選其他人呢.

$^{81}$這個都是一般我們人會問的問題.
或者會提出的反對.
但是聖經往往就告訴我們.
正正我們今天所看的神.
就是說我定要揀選雅各.
我要愛雅各.
和雅各後來衍生的.
以色列這個民族.
十二支派等等.
這個正正就是那個揀選的原初的意思.
還有一段經文.
上個星期就不夠時間看的.
如果看完就會再比較完整一點.
就是這個《若書雅記》24章.
就是若書雅年老了.
她自己110歲差不多要死了.
她就講了一個很重要的歷史的一些背景.
就是我們這個民族是怎麼來的呢.
在若書雅差不多死之前.
其實迦南地還有九個半支派.
都還沒有得到他們自己的土地.
就是神應許了.
好像已經規劃了.
他們需要自己的民族.
當然是那些南丁.
就要去打那些迦南人.
才可以得到的.
六半支派已經得到了.
就是劉變和迦德支派.
已經得到約旦河東的土地了.
所以劉變和迦德支派.
他們就會跟隨大軍去打完.
才回到約旦河東的.
所以後來也衍生了一個竹壇事件.
引發幾乎約旦河東和約旦河西.
有一些的紛爭.
以為他們都沒有吸取歷史教訓.
還在那裡好像迦南人那樣.
到處去拜那些偶像.
他說不是的.

$^{121}$我們怕我們的後裔和你們的後裔.
將來因為隔了一條約旦河.
大家彼此不認識.
所以就竹一個壇.
就是說我們是以色列人.
我們的祖宗一起打迦南人.
就有一個這樣的故事.
我們看《約書亞記》24章.
比較完整的約書亞.
講述一個很重要的民族起源.
其中很重要的都碰到.
雅各被揀選這裡.
24章的《約書亞記》.
由第一節開始.
約書亞將以色列的眾支派聚集在士堅.
照了以色列的長老,族長,審判官.
並觀長來.
他們就站在神面前.
約書亞對眾民說.
耶和華說以色列的神如此說.
古時你們的列祖.
就是亞伯拉罕和拿學的父親.
他拉即是亞伯拉罕的爸爸.
住在大河.
即是右法拉底河.
住在大河那邊事奉別臣.
《約書亞記》的24章.
《約書亞記》24章.
創出李文生書第六卷.
《摩西五經》.
接著就到《約書亞記》.
要背這些口訣.
你們沒有翻譯和主義學就不會背了.
所以就要補課了.
第三節.
我將你們的祖宗.
亞伯拉罕從大河那邊帶來.
令他走遍迦南全地.
大河就是右法拉底河.
和底格里斯河.

$^{161}$即是現在的中東一帶.
古代的近東的發源地.
即是文化的搖籃.
亞伯拉罕就由那裡開始.
或者當時叫做亞伯蘭.
就由那時開始被神選召.
呼召他離開本地.
本族 附加.
去到神應許的地.
就是迦南地.
來到他得到這個全地.
留意一下.
「由使他的子孫眾多」.
「把以撒賜給他」.
「又把雅各和以素賜給以撒」.
其實如果大家來不及.
用手機 手機就快些.
因為現在揭聖經.
都變成神功.
變成一些特技人才會這樣.
因為太多人已經不揭聖經了.
或者都不帶聖經出街了.
Apps是一個進步.
是一個歷劃時代的.
恩典的物件.
即是大家不用背.
都可以按到哪一個書卷.
第四節 留意這裡.
「又把雅各和以素賜給以撒」.
有沒有留意那個次序?.
原來約書亞記都已經有次序分了.
其實兩兄弟排名的.
一般都是由大到小數的.
譬如我家就是.
我哥叫做曾睿根.
根就是根源的根.
然後就是榮華富貴.
這樣排的.
然後又昌又雄.
很偉大.

$^{201}$如果你是女的.
譬如梅蘭菊莊.
或者什麼等等.
一定要這樣排的.
但是如果聖經刻意這樣調轉次序.
你就知道原來聖經的啟示.
是忠於創世記的事件.
是忠於神對以撒的應許.
和對利物加的應許.
就是說將來大的以素.
要服侍那個弟弟.
小的就是雅各.
雅各成為那個長子.
無論他用什麼手段都好.
這個是神命定了.
所以這裡你就會看到.
原來不只是希伯來書的作者.
有這種洞見.
原來約書亞都已經知道.
要這樣排列.
其實很不敬.
如果分開大兒子小兒子.
沒理由大兒子排第二.
所以絕對不是守文之誤.
不是排版排錯.
原來這個是忠於神的揀選.
那個預定.
有把雅各和以素賜給以撒.
將西爾山賜給以素為業.
後來雅各.
接著就不再講以素了.
一直講的就是雅各.
這個民族他們的繁衍.
一直數下去.
等於應許之地的希伯來人.
就是神的選民的意思.
要留意到後來雅各和他的子孫.
去到埃及去.
我差遣摩西亞倫.
又是倒轉的.

$^{241}$亞倫是大兒子.
又是倒轉的.
當然還有一個大姐姐叫米尼暗.
你可以看到摩西是被揀選的領袖.
其實阿倫是輔助的.
你會看到我差遣摩西亞倫.
其實是約書亞那一代.
約書亞看著他們回來.
帶領他們的族民.
離開埃及.
約書亞當時也是四十歲.
所以你會看到.
約書亞也懂得這樣分次序.
接著「照我在埃及中所行的降災於埃及.
然後把你們領出來」.
第六節.
「我領你們列祖出埃及.
他們就到了紅海.
埃及人帶領車輛馬兵追趕你們列祖到紅海.
一路如是者」.
就講這個民族的一路的發展.
他們在曠野.
他們怎樣來到去要漂流.
後來進入迦南.
約書亞就講到那些情況那些事件.
另外其實還有一個地方要留意的.
就是去到這個準備進入迦南之前.
有一個叫做以東.
留意一下下面.
第七節.
「你們列祖哀求耶和華.
因為追兵在紅海邊緣追到來.
是一個大的平原.
所以前無去路.
因為是紅海邊.
後面就是埃及的重兵追趕來到.
所以必須是一個神蹟.
就是紅海要分開.
當他們哀求的時候.
他就使你們的里門和埃及人中間的黑暗了.

$^{281}$又使到海水淹沒埃及人.
是淹沒埃及人的.
這是另一個角度去講.
那個紅海分開.
所以新派的.
可能大家很少會聽到.
新派的教會神學家.
他們就說其實那個不是神蹟.
那個叫做蘆葦海.
即是說其實是一些很淺的水.
淺的水怎樣去淹沒埃及的兵呢?.
所以如果你要幫聖經解圍.
讓那些神蹟不要令到人那麼難去接受.
你就盡量當時的地理各方面.
其實約書亞只是講得大了.
其實那些是蘆葦海.
不是紅海分開.
即是以色列人都可以走的.
就不會淹死的.
但這裡很明顯已經告訴我們.
是淹死了那些埃及的士兵.
不是那些婦孺.
不是那些竹蝦的太太.
是那些埃及的士兵.
接著說又使到海水淹沒埃及人.
我在埃及所行的事.
你們親眼見過.
你們在曠野也住了許多年日.
當然是說38年.
我領你們到約旦河東阿摩利人所住之地.
他們與你們爭戰.
我將他們交在你們手中.
你們便得了他們的地為業.
我也在你們面前將他們滅絕.
那時摩訶王西伯的兒子巴勒起來.
攻擊以色列人.
打發人召了比爾的兒子巴勒.
即是騎著驢子那個.
騎著驢子會說話的那個.
巴勒來救助你們.

$^{321}$我不肯聽巴勒的話.
所以他倒為你們連連祝福.
這樣我便救你們脫離巴勒的手.
輾轉.
接著一直說他們怎樣將迦南的七族.
最大的迦南國家的人.
特別是迦南人來滅絕.
一直說到約書亞爾加他們年紀老邁.
其中有涉及到以東的.
有涉及到以素的後人.
不過我們不詳細去看.
回到馬拉基書這裡的時候.
我們要很具體掌握到.
就是神來講.
我再讀一到第五節.
這個馬拉基書.
其實馬太福音之前.
舊約在我們中文的翻譯本.
最後的一卷書就是馬拉基書.
第一章耶和華.
馬拉基就是傳訊息的意思.
傳給以色列的默示.
默示又可以是一個burden這個字.
即是一個重擔.
不只是一個message這麼簡單.
或者是一些先知.
第二節耶和華說.
我曾愛你們.
你們卻說.
在何時尚愛我們呢?.
明顯是一個過去式的.
以色列人的領會就是.
你過去愛我們.
現在如何?.
將來會不會愛?.
過去愛我們會知道的.
有一些事件的.
例如剛才講的過紅海.
明顯就是一個過去.
曾經幫我們這個民族.

$^{361}$相信這樣的一個回答.
好像完全提不起勁的一種回覆.
顯然就是不滿意現狀的心情.
這樣的一個態度.
來到未必是跟神說.
就是直接跟先知反映這種心態.
就是是嗎?.
過去愛我們.
現在還愛不愛我們這個民族呢?.
或者古代的人.
跟我們不相同的地方就是.
很少會個人自己的感受.
因為古代都是一些幫國.
一些民族.
不會有你個人的自己的獨斷獨行.
例如等於圍村人.
你就可能會明白一點.
很多事都是阿公說了算.
祠堂說了算.
不到你個人有一些.
很強的自己主張.
這些就是一些強烈的族長的文化.
例如現在阿拉伯世界.
有些地方都會是這樣的.
林牧師我們去印尼宣教探訪.
他說那條村如果村長信主.
整條村就會跟著信主的.
就不能給你個人自己的喜好.
這些就很強烈的反映到那種集體性.
古代就更加是.
所以這裡的反映.
很明顯就是想問一問.
古代愛.
古代的人得到的幫助.
但是現在我們這一代人.
神怎樣愛我們呢?.
我們好像在現在他們的狀態.
他們的敬拜.
我們說過那個背景.
不是有很大的迫不.

$^{401}$未去到主前的第二世紀.
因為敘利亞王和埃及.
他們有一些競爭.
以色列夾在中間的一個和心.
所以成為了兵家必爭之地.
現在他們寫這本書的時候.
就未去到一個這麼嚴峻的政治.
外力的威脅.
所以就似乎現在.
他們未有這種太大的.
需要為他們猶太教的信仰付代價之前.
他們可能就要面對的就是.
那個生活開始慢慢趨於平淡.
趨於形式化的敬拜生活.
所以對於神他們似乎.
因為我們下面第六節開始.
就已經進入那個敬拜的中心.
探討什麼的敬拜才合神心意.
怎樣的獻祭才可以滿足神.
就是說又沒有外敵.
又可以恆常的獻祭敬拜.
就是說他們的生活一定是慢慢.
慢慢趨於一種形式的平淡.
這一份的平淡其實可以是一個危機.
所以現在先知馬拉基.
就來提一個很大的課題.
就是你要改變你自己內心的平淡.
就必須要再感受到神的愛.
那哪裡去感受呢.
從哪一個角度去得到神的愛呢.
這裡就告訴我們.
那個解釋就是我們已經用了一堂講耳掃.
和牙角不被神揀選的人.
和被神揀選的分別.
其實如果按人來看.
兩兄弟都是沒有分別的.
都是犯罪的.
都是不應該得到神的揀選的.
但是奇妙之處就是.
神的定義要揀選哪一個.

$^{441}$而不要他的哥哥.
接著這個就會引發.
百姓或者當時的以色列人.
會有很強烈的反應和抗議.
還有最終我們想問的就是.
怎樣讓我們可以經驗到神的愛呢.
這個我們從第三到第五節.
都有些少答案給我們可以去看到的.
當然要成為我們今天現代21世紀.
我們香港.
我們現在有我們自己的.
每個教會都有它獨特的文化.
會有慢慢營造的傳統的.
例如我們錦宣的傳統.
有些都有的.
譬如初一是有崇拜的.
不是每間教會都有的.
譬如我知道過渡會恩傳堂有的.
初一有的.
有些教會就算很大都未必有的.
譬如元朕初一我見他沒有開的.
因為他們沒有這個傳統.
每個教會都有它的傳統的.
所以我們要問.
透過這五節聖經.
給我們一個原則.
一個方向.
但是最終我們怎樣能夠.
落實在我們自己的生命裡面呢.
第三節就告訴我們.
以蘇系雅各的哥哥.
是人都知道的了.
沒有什麼困難的了.
但是我們剛才揭開的約書亞記.
希伯來書.
包括羅馬書.
都無論舊約新約都告訴我們了.
新約的作者都認同.
創世記的歷史背景.
第27章的創世記.

$^{481}$告訴我們.
以蘇興翰他長子的名分.
以蘇失去他長子的祝福.
因為都是兩次被弟弟雅各欺騙的.
折竹先登的.
但是原來心底的深處.
這個以蘇是一點都不稀罕神的blessing的.
看為是有礙於他自己的物質的私慾的.
但是雅各就識貨了.
緊張神.
是亞伯拉罕的神.
是列祖的神.
即是我的父親以撒.
我的爺爺亞伯拉罕.
還有其他的更加遠的先祖.
他們是經驗過這位神的.
相信他們家族的口傳歷史.
太公的太公 太公的太公.
就是阿當了.
即是他們有這個渴慕.
你應該要和這些先祖.
有信心的遠古的人.
要拉上真正的關係.
你就可以得到神的blessing.
所以現在雅各.
他知道自己不是那個大兒子.
即是說與生俱來.
按身體狀況.
他是不應該得到那個祝福的.
他不認命.
他用他的手段.
而且我們知道是不好的手段.
但是他就是要得到神給他的祝福.
聖經就沒有交代到.
兒子 媽媽有沒有告訴他.
天使已經不是報夢了.
是向你家顯現了.
媽媽有沒有這麼重要的訊息.
告訴兒子呢.
兒子 你將來是要長大的.

$^{521}$有沒有說過呢.
長大了 要不要自己做點功夫呢.
有 就是用手段.
媽媽也幫他用手段.
就是幫弟弟.
即是他自己的小兒子.
來做點功夫.
還有教他欺騙自己的父親.
即是丈夫.
有禍就我擔.
那個福氣就是你要.
媽媽愛兒子就愛到這樣.
你可以這麼說.
那些劇情都很符合我們那些八點鐘的電視劇.
就是這樣的.
即是什麼風暴那些.
類似那些.
第四節.
第三節的中間.
將他的地業.
交給曠野的野狗.
恨戶已掃到一個地步.
使他山寧荒涼.
他住在以東地.
其實都是很險要的.
能夠在這麼荒蕪的地方.
能夠打遍江山.
也可以印證到.
阿爾素和他的後人.
都是很強悍的.
都是很驍勇善戰的.
即是說在這麼缺乏的地方.
都能夠成就一番事業.
但是我們要留意就是.
第四第五節.
我們都很難從心裡面.
能夠對神的作為的佩服.
但是變成我們用時間去印證.
神的作為.
所以我們說完之後.

$^{561}$我們就要回來看回我們自己的信仰人生.
以東人就說.
我們現在雖然被毀壞.
卻要重建荒廢之處.
萬君知耶和華如此說.
任他們建造我必拆毀.
人必稱他們的地為罪惡之境.
稱他們的民為耶和華永遠老老之民.
你們必親眼看見.
也必說願耶和華在以色列境界之外.
被尊為大.
被尊為大又是在這本書出現很多次.
我們要認識我們的信仰.
由入門開始.
牙牙鶴語.
人家叫你認罪祈禱.
你跟他認罪祈禱 信主.
但是這個被尊為大.
是要我們一生要去追求.
要去讓神被尊大.
是很重要的一生的功課.
人見到我們 是否神得到榮耀.
我們的人生是否讓神得到榮耀.
神被高舉.
這告訴我們.
有一件很恐怖的事發生.
就是神說他要定義懲罰這個耳掃的後人.
他們甚至知道自己被毀壞.
他們的地圖荒涼.
但是那種不甘心.
看不到或者不肯接受.
是做物主的咒詛.
對耳掃的一個一直追討的刑罰.
使他們好像一生注定就是走壞運.
你很難想像.
為何神要這麼重的刑罰耳掃.
整個問題我們要抓得很實.
就是和以色列人說.
我罰這個人不是要告訴你我有多惡.
而是我要告訴你們.

$^{601}$我是多麼愛你們這個民族.
整個問題在這裡.
不是罰那個問題.
罰這個人 這個民族.
就是我愛你們以色列的明顯證據.
作出這樣來和我們討論.
將神的愛這樣告訴我們.
因為我們已經是信主的人.
在救恩的範圍保護下.
沒有東西可以使我們和耶穌的愛隔絕.
我們原來是一個很受保護的氛圍裡面.
只要我們長久在這個愛的氛圍.
我們就是神的愛的對象.
我們就要從這個.
記住這個不是安恕區.
這個就是在耶穌的愛的包圍底下.
的一個人生.
去看那些沒有耶穌的人.
他們的遭遇 他們的困境.
他們的人生的厄運.
這樣理解.
雖然那些不信的人會覺得.
你這樣說 曾牧師你很涼薄.
你們在愛的 在耶穌的氛圍底下.
一生蒙眷顧.
將來蒙神的祝福上天堂.
我們就注定好像耳掃的後人那樣.
抵擋神 我只不過是不信祂.
我只是覺得祂不存在.
我沒有意思要抵擋祂.
為什麼你們基督徒說得這麼重呢.
這個就是古代.
先知這樣和以色列人.
來辯論關於神的愛.
這樣的大方向原則.
能不能套用在我們基督教.
基督徒的身上呢.
能不能套用在我們新約的人身上呢.
沒錯 我們要記得神愛世人.
是不能再更改的了.

$^{641}$任何的民族都是蒙神的愛的.
所以在舊約有少許的限制.
就是因為尼塞俄未降臨.
以色列人仍然覺得自己的民族.
是唯一的蒙簡選.
所以難怪保羅都有這種扭曲了.
猶太人唯有自己是優越民族.
這個心態來抗拒.
外邦人可以說自己信耶穌.
就得到那個救贖.
這樣保羅未通.
因為他原來.
我慢一點說.
可以的.
就是他因為未通到彌賽亞.
就在耶穌的身上.
和他被舊約的人.
因為保羅就是舊約的人.
徘徊在憲制.
實物教材.
拖羊,牛去聖殿找祭司.
去憲制的人生.
即是說保羅小時候是憲制的.
他曾經是拿西爾人.
他去宣教都會上聖殿還願.
這是保羅的舊約,新約的過渡性人物.
所以他們就會有這種觀感.
就是很難接受外邦的基督.
外邦人可以信而得到.
和我們一樣的愛的保護的氛圍.
所以這一班人他們未明白.
所以我們在這一段經文.
有那個局限性.
就是說這個愛.
去到新約.
去到耶穌成就舊恩.
我們就要知道.
不僅僅是為我們信耶穌的人.
可以在這個圈裡面.
我們不可以因為在圈裡面.

$^{681}$就抗拒其他的人.
其他的民族進入這個舊恩裡面.
我們應該超越了馬拉基書的時代.
舊約時代.
他們一直的局限.
只是他們猶太人這個民族.
得到上主的任弊的祝福.
而導致他們就有這種錯誤的觀感.
所以我們記得.
這段經文的局限性.
耶穌的降生,耶穌的成就舊恩.
就已經超越了.
我們不可以來看其他未信的人.
為絕對的奏作.
因為舊恩就是在這個奇妙的地方.
或者預定論其中另外一種.
很重要的前設就是.
我們地上的人.
包括你和我已經信主了.
我們算是知道自己被揀選.
但是我們的限制就是.
我們不知道未信的人.
他是不是被神揀選.
我們不是神.
我們要帶他信主.
他最終都不肯信.
我們就可以肯定.
他是沒有被神揀選.
要搞清楚這件事.
所以我們只能夠做的就是.
盡量帶人信主.
我們就知道他就是被神揀選了.
這個限制就是.
神不會給我們有這種預知.
神就是預知.
我們沒有預知的.
我們只能夠帶他信主.
感動他信主.
他就是.
原來當他回望.

$^{721}$原來我是被神選召的.
得到神的預定的.
但是他信耶穌那一刻.
他不知道自己是在那個預定裡面的.
其實我們都不知道的.
這個就是預定論.
其中的奧妙之處.
所以如果我們記住.
在這一段經文.
有一個限制.
但是當我們回到經文的世界.
舊約的時候.
記住所謂的限制就是.
耶穌基督未成為彌賽亞降世.
未成就舊恩.
猶太人仍然活在舊約的氛圍.
所以也不能怪猶太人.
舊約.
也不能怪普羅.
因為他們只是停留在.
舊約的眼光的世界生活.
神是給予那個階段的啟示.
所以耶穌出來.
是整個的.
嶄新的啟示.
所以祂就是那個新約的.
很重要的使者.
這個希伯來書一開始已經說了.
所以這裡就告訴我們.
很重要.
就是神要給以色列人知道.
我對你們的愛.
是藉著兩兄弟和兩個民族.
好像又是一個示範的單位.
我定義愛你們.
這個愛就不是溺愛了.
是有管教的.
是有訓練的.
是有要遵守律法的.
是要聽從神命令的.

$^{761}$這樣才可以叫做在無愛的範圍裡面.
或者說繼續是被神揀選的.
那些有信心的以色列人.
所以普羅後來寫《羅馬書》.
我們上個星期看的第九章.
一到第六節.
正正就是說這件事.
普羅說了一句話.
以色列人不都是被揀選.
也不是因為亞伯拉罕的子孫.
都是全部被蒙揀選的.
因為以實瑪利就不是了.
只有以撒是.
所以原來就算是肉身的亞伯拉罕的後裔.
都不是每一個可以這樣去經歷這份愛的.
這個就很大部分.
就是看那個民族.
或者那些的支派.
或者那些的人.
他們有沒有跟隨神的摩西的律法.
來去生活.
來到迦南之後.
有沒有跟隨他們的族長.
來去敬拜神.
還有來去敬畏神.
過一個神規範的生活.
這個就是約書亞的之後.
當約書亞死之後.
就立即有時思的時代.
還有開始有撒母耳的王國的歷史的開始.
就是真的告訴我們.
到底摩西五經有沒有落實在每一個時代.
每一個家庭裡面的問題.
所以慢慢演變.
就成為了原來不是每個人都是被揀選的以色列人.
原來需要你們自己能夠去跟隨神的心意而活的.
才可以算得上是有信心的以色列人.
這裡就很重要告訴我們.
原來我們要去掌握神的愛.
我們不是那麼虛無縹緲的.

$^{801}$因為首先就是神已經在歷史裡面顯現.
神已經是藉著向阿當夏娃顯現.
向以洛顯現.
向羅亞顯現.
向亞伯蘭顯現.
向以撒.
向雅各顯現.
向約瑟.
向約書亞.
向跟著有其他的時思來顯現.
神從來沒有停止過向人類顯現.
祂是那一位全能的神.
而且還留下文字的記錄給我們.
讓我們可以透過研讀這些文字.
我們可以掌握到神的.
昨晚的系統神學就是講這件事.
讀系統神學就是要對神的attributes.
祂的屬性.
祂的nature.
神.
動物主義是怎樣的呢.
另外兩個課題.
就是在地上歷史祂做過什麼呢.
祂怎樣做呢.
藉著什麼去做呢.
去彰顯祂的救贖呢.
這個就是系統神學一定要掌握的兩個課題.
所以有救恩論.
也有將來的末世論.
我們遲些看完這本書.
我們就會看帖撒羅尼迦加前後書和末世論.
我們都會用William Goulden的那本巨著.
來參考.
因為整個topic都幾十頁了.
整個topic就是講末世論.
讀神學就一定要懂的課題.
所以都有機會和大家掌握到系統神學的課題.
所以如果你有興趣.
可以只聽郭院長一課.
他會講末世論的.

$^{841}$可能當然沒有講得那麼詳細.
因為他要探討其他東西.
關乎到靈因運動.
特別第五堂和第六堂的異端.
我不是賣告白.
我們回來看這裡的時候.
來講一下.
所以這一段聖經很重要.
以色列人就是感受不到神的愛.
那就是神怎樣去解答呢.
我的愛其實你們是可以見得到的.
神是一直在他的子民當中工作.
神一直不是在大自然那裡.
讓你自己去搜尋.
你看到地球多麼漂亮.
有些人就多願意去看北極光.
看完之後回來不會令到你敬虔了.
是不會的.
你去讚嘆完那些相片很漂亮.
你開影展等等.
你的相片.
你成為很出色的攝影家.
我沒有貶義.
當然我們是順著這個課題.
我都喜歡這些大自然的.
但是大自然是不能夠將神去傳息給你知.
必須是要在聖經.
必須是要在他的子民當中.
你才會見到神的工作.
所以神就是說.
你一直留意耳掃這個民族.
是你們的對頭.
我一直就是要就坐他們.
使這個簡選.
你們這個民族一直都有一個延續.
因為他們最埋怨的就是.
我們這個民族.
你有70年將我們遺棄.
被擄到巴比倫重建.
其實只有5萬8千人回來重建.

$^{881}$其實在巴比倫.
在波斯帝國.
是有很多的猶太人.
但是都沒有回去.
就是說其實是很零星落索的.
輾轉我們建立我們自己的敬拜.
我們有我們的聖殿.
也有我們的城牆.
有我們的一些領袖.
這個索羅巴伯.
這個大祭司約書亞.
也有以色拉.
接著有尼熙米等等.
但是這些都過去了.
過去了我們回去重建.
但是你能不能夠具體告訴我們.
你特別愛我們.
給我們有很多很多的那個祝福.
不需要我們面對.
可能是生活的艱難.
可能是氣候的艱難.
我們又要面對.
可能有很多其他的因素.
但是作者似乎沒有將那些民生困難.
來描述得很詳細.
一定是有很多這些因素.
來去編織很多百姓內心的埋怨.
來去將他們的困難.
來去看為神好像都沒有.
每天給我們有很多的祝福.
以致我們累積的這些困難.
都沒有辦法去疏解.
那何來講得上你是愛我們這個民族呢.
那似乎這個時候.
就差不多是萬能的關鍵了.
任何一個時代.
任何一個人都適合用了.
如果你面對一些很大的困難.
而那些困難三朝兩日都改變不了的.
那你就會有很多的苦不堪言的.

$^{921}$慢慢慢慢就變成沒有了那種動力.
來到去對於神的愛.
你會覺得是很感受的.
你會慢慢停留在一種.
每天起來都是那個困難.
有時那些困難就再踏上另一個困難.
每個人都不同.
就不是每人一個小確幸.
是每人一個小困難.
有些甚至是大的困難.
很難講得上.
我每天都很感受到神的愛.
大部分人都不會這樣.
那時候正正就慢慢慢慢就落在一種.
信仰的枯竭.
是內心的枯竭.
真是由內心感受不到神的愛.
那個問題而開始.
所以很好的就是第一個問題.
神就跟以色列人講.
我曾經愛你們.
而且是會繼續的來保守你們這個民族.
就是一直的在一個現在式.
這樣來到去好像要洗到耳掃.
和他們後人一直成為神的敵人.
神好像要留意.
不是祝福 是降禍.
洗到他們好像沒一口好吃.
好像黑社會那樣.
我見一次打一次.
這樣的已經沒那麼白.
很白就再難聽一點.
這樣的感覺.
事實上又真是這樣.
這樣的描述.
這個是.
喝口水.
因為碰到那個位置就會癢.
如果不說話就會好一點.
就是小聲一點.

$^{961}$這裡就再一次告訴我們.
為什麼我們要將這一段再透徹一點去演繹.
就是神不只是古代.
就是說神很想以色列人.
當然這個原則去到新約一樣.
很想我們留意的就是神的作為.
所以我們要留意神的作為.
大自然是其中一個辦法.
但是大自然只不過是普通的啟示.
這裡又涉及到系統神學的一個重要的分野.
就是普通的啟示和特殊的啟示.
普通的啟示就是說.
現在你見得到的東西.
我們有的定律.
例如萬有引力.
有其他的.
但是這些是不會令一個不信的人.
他會尋找到神的足跡的.
必須透過普通的啟示.
加上特別的啟示.
特別的啟示就是聖經的真理.
就是告訴我們.
我們可以透過聖經.
我們就可以尋找到神.
或者說正正就是神透過聖經.
透過歷史的救贖.
祂的啟示讓我們去認識祂.
當然最重要的啟示就是耶穌.
就是來這個世界成為神的愛子.
這個就是告訴我們.
透過這一段經文.
很想以色列人留意.
神的救恩在歷史裡面.
那個從來沒有缺席.
一直愛這個民族.
一直在這個民族的生活.
和他們的歷史裡面來發展.
所以在這個時候.
雖然都是主前的四百多年.
但是神從來沒有停止.

$^{1001}$祂在這個民族的介入.
那就是說我們要多肯定地相信.
今天的以色列.
神同樣有祂的計劃.
所以不是由誰的領袖.
政治的領袖.
現在的總理就是內塔尼亞胡.
由他或者他的幕僚來決定.
的確是他們發號施令的.
他們安排的軍隊.
進城等等諸如此類.
但是仍然整個局面是神掌管的.
我們就要從這個這樣的全面.
我們透過神一直的.
沒有改變的工作.
我們就知道.
祂從來沒有對這個民族.
以色列缺席.
那就是說新約更加.
透過主耶穌的神愛世人.
祂從來沒有對地球的人.
有任何的缺席.
所以救恩藉著不同的教會.
是繼續這樣去傳的.
所以教會是唯一不會.
變成夕陽公業的一個地方.
我們要記住.
我們可以過去.
我們可以回天家.
但是教會是繼續存在的.
神的救恩是繼續要傳的.
有一天就是末世論了.
耶穌回來就不需要救恩了.
那個時候就是不同的派別.
想的到底是七年大災難.
還是災前對題.
那些不同的派別.
就有不同的理論了.
我們去到末世論.
我們就會講這些派別.

$^{1041}$看到時大家是信哪一個.
接受哪一個方面了.
所以這裡得到一個結論就是.
我們必須透過這五節.
我們要認定神在歷史裡面工作的.
神有祂的計劃的.
所以我們必須要透過聖經.
特別是新約.
對於神的救贖的計劃.
我們要上心.
否則的話.
你就徒然說神不愛你.
那就.
的確.
你從你自己裡面的心靈去找.
你是找不到神的愛的.
我們要透過神的計劃.
神的一個在地上的工作.
我們就會看到神從來沒有對人缺席.
更加我們是那個其中一份子.
神從來都是愛他自己的選民的.
我們要記得.
新約再不只是以色列是神的選民.
是每一個被神揀選.
被主耶穌拯救.
有聖靈的人.
就是神的選民.
就是.
神的.
屬天的亞伯拉罕的兒女.
這個我們要肯定的.
首先有了這個肯定.
我們就要將我們的眼光.
放在.
那個神的.
救贖裡面的參與.
所以如果我們.
不去參與.
也不去留意.
神在地上.

$^{1081}$尤其是在人的工作上的留意.
那我們就很容易.
會停留在我們自己.
那個生活的順利.
順暢.
作為神愛不愛我的指標.
那就很危險了.
因為我們人.
受著很多環境影響.
甚至神都不會給我們.
特別有一些.
強烈的一個.
特殊的恩賜.
因為耶穌都說了.
那個定律就是.
神是照好人.
用太陽照好的人.
也照那些不好的人.
是普遍的.
是沒有例外的.
所以我們要記得.
既然耶穌都說了.
這樣的定律之後.
那就是說.
如果我們信了主耶穌.
被人灌輸.
這個第五堂.
穆斯會說的靈恩運動.
就是我們如果有一個.
深信不疑被人這樣改變了你.
你是神的子女.
你特別蒙福的.
你能夠可以很多東西.
化險為夷的能夠行神蹟的.
那你就會.
違反了主耶穌說的.
神愛好人.
又愛歹人.
他降於給義人.
又降於給不義的人.

$^{1121}$那個是正常的定律.
所以.
我們要記得.
我也提過了將來我會說.
我們信主.
我們事奉.
如果神要我去今晚.
要我去QE.
有一個教友.
他的家人被信的.
如果神要我.
有一個很特殊的原因.
要醫好他.
透過我按手祈禱.
他馬上康復了.
可以的 神可以的.
連那個醫生都覺得奇怪.
他已經是絕症了.
為什麼你祈禱他可以起來.
是不可能的.
醫學上是不可能的.
但是神如果要藉著我去做.
是可以的.
那我仍然都信這件事.
今天我都沒有經歷過.
神還沒有安排我這樣做.
但我不是不信.
我可以行神蹟.
但如果你變成一件事.
總之你是傳道人.
你就可以行神蹟.
這就不是普遍的聖經.
包括哥林多前書十二章都說.
這人蒙聖靈.
賜他有醫病的恩賜.
那人有神給他的.
一個安慰的恩賜.
這樣分開的.
所以如果我們將那件事.
變成本末倒置了.

$^{1161}$原來法號師裡面就是那個人.
是那個主張有神蹟的人.
那就不對了.
所以這五字聖經.
是一個很重要的結論.
要找到神的愛.
必須.
不要只是普遍的啟示.
即大自然.
不要只是從.
大自然.
不要只是從.
表面的人事物.
去找神的愛.
是很漂亮.
但很漂亮你.
不可以將它解釋為神的愛.
大部分人都不會覺得.
神的愛.
他看到一個很漂亮的雪山.
他不會覺得神的愛.
但是.
如果我們回去聖經.
在神給我們的.
救恩的耶穌的愛.
十字架的愛.
和聖靈在我們.
內心的感動.
我們跟隨神.
願意走在神的.
旨意下的基督徒.
我們就可以.
去詮釋.
神的愛.
甚至我們感受到神的愛.
在不同人身上的.
作為.
這個就變成很重要.
為什麼我們必須要在.
神的工作裡面.

$^{1201}$我們有參與.
有份.
你才會對神的愛.
有更大的.
經驗.
和你自己個人的經歷.
就是這個原因.
所以這裡有一個過去.
因為是.
神揀選了.
那個現在就是神繼續.
在以色列這個民族.
當中.
有祂的訓練.
有祂的工作.
使到以色列將來.
人肉彌賽亞的降臨.
所以你想想.
這個還不是救恩的歷史.
神在人類當中.
工作.
神就是在以色列這個民族.
當中工作.
你就會看到.
這個很重要的一個神.
從來沒有缺席.
將來.
神就說.
祂要在祂的子民當中.
就是五節.
你們必親眼看見.
也必定會說.
願耶和華在以色列.
境界之外.
被尊為大人.
外邦人.
因為一個猶太人.
耶穌.
就成為神的子女.
以色列人.

$^{1241}$如果懂得看.
看到這麼多外邦人.
信主.
這麼多外邦人不拜偶像.
敬拜真神.
敬拜獨一的真神.
就是猶太教最重要的信念.
而來歸向神.
又稱呼自己是神的兒女.
又得到將來.
天上的基業.
這個豈不是.
願.
耶和華在以色列的境界之外.
地域的境界.
生命的境界之外.
被尊為大呢.
這個不就是.
要在神的救贖.
持續的救贖的裡面.
有參與.
早期就停留在猶太教.
那個觀念裡面.
他就沒有參與了.
他就是這樣來到.
抵擋外邦人.
抵擋耶穌.
但當他通了.
他知道原來.
神已經藉著耶穌.
在以色列的境界之外.
就是外邦人.
得到.
被尊為大了.
所以保羅當明了這個真理之後.
他整個人.
就將他餘下的人生.
當然是神揀選他.
來奉獻.
給外邦人信主.

$^{1281}$這件事.
他就會更深的經歷.
其實他的生活.
因為傳道是很苦的.
猶太人又不認識他.
拉比會逼他.
有些.
控訴他.
那些人.
來亞利亞.
來告保羅.
後來甚至.
保羅要被.
押解到.
羅馬.
去上訴.
你可以想像.
保羅.
不是一帆風順.
的人生的生活.
遭遇有時甚至.
盜賊的危險.
江河的危險.
船壞的危險.
接著又要面對.
有些異端的影響.
但是.
對於整體來說.
他通了神的救贖.
在人類.
藉著耶穌.
他明之後.
他就深深的.
不單他自己感受到神的愛.
對世人的.
持續的關心.
和拯救.
更加.
他自己參與在.
這個救贖裡面.

$^{1321}$他寫.
他成長的.
你可以說.
他的.
心路歷程.
就是神的愛.
怎樣藉著他的宣講.
或者他的書寫.
來祝福.
萬世後代的人.
我們今天能夠讀到.
保羅的書信.
明白一些真理.
我們就明白神的愛.
就是因為藉著保羅的書信.
所以保羅.
正正第五節.
就是一個如果一路跟隨神的足跡.
來看神的作為的人.
他就會看到.
原來神的愛.
是古代.
現在將來.
是沒有窮盡的.
持續的愛我們這個民族.
所以他們會.
真是親眼看見.
必定會說.
願耶和華在以色列的境界.
之外.
被尊為大.
就是這樣的意思.
所以我們回到今天.
看回我們.
自己的人生.
我們當然不會那麼偉大.
我們成為使徒.
我們可以寫聖經.
也不需要.
因為聖經已經.

$^{1361}$封閉了.
66卷.
這個就是下.
第三堂就會說聖經無誤.
應該是下一堂.
聖經無誤.
就會說聖經的醞釀.
形成.
所以我們是不可以.
再自己作聖經.
因為66卷是已經.
封閉了的canon.
就是那個正典.
所以我們就要在.
神的救恩的場景.
去發現神的愛.
大自然.
只能夠很殘缺的.
找到神的愛.
一定要是.
夾雜.
大自然加上神的.
救贖的場景.
就是教會.
就是耶穌的十字架.
我們就可以找到神的愛.
所以這裡.
透過這五節聖經.
讓我們來看到.
神從來沒有缺席.
所以我們要認定這件事.
但是這裡的.
涵蓋的範圍.
其實我們看.
以色列的歷史一樣.
從來.
沒有保證這是一個.
的溺愛.
這是一個偏心的愛.
沒有保證.

$^{1401}$所以我們要記住.
我們要走在神的.
那個恩典裡面.
也都帶著一個對神的.
敬畏來生活.
進入第二個階段了.
第二個重要的問題.
作者很快就.
差不多入獄了.
為什麼你感受不到神的愛.
其中更大的部分.
就是敬拜的生活出了問題.
敬拜.
哪一個.
情況出了問題呢?.
第六節.
藐視我明的濟斯.
萬軍之耶和華對.
你們說.
咳咳.
兒子.
尊敬父親.
獨人敬畏主人.
我既為父親.
尊敬我的在那裡呢?.
我既為主人.
敬畏我的在那裡呢?.
尊敬和敬畏.
其實只不過是.
可能是程度上.
但是應該似乎那個.
措辭都是一樣的.
神要獲得那個被尊崇.
就是第五節.
那裡所說.
就是說你如果看到神的作為.
你必定會有一種.
內心 發自內心.
要記住發自內心.
那個對神的尊敬.

$^{1441}$和敬畏.
就很自然的了.
那你有這個這樣的.
追求.
和心態的時候.
你會.
經驗神的愛.
給其他的信徒.
給其他不信的人.
是會更加快.
和更加真實.
就是你會經驗到神的愛.
那是很重要的.
你在這個世界你能夠經驗到神的愛.
是很重要的.
你想想現在的社會.
經驗不到父母的愛.
經驗不到愛.
就已經死了.
因為覺得絕望.
但是如果我們經驗了神的愛.
所以.
多困難的人生.
當他經驗到神的愛.
他不會覺得自己的人生.
是苦的.
譬如上個星期六.
Happy姐的安息禮拜.
謝牧師講了一個例子.
在台灣.
一個牧師的女兒.
天生是殘缺的.
她的四肢.
都是.
可能是因為某些的病.
她是不能夠.
正常地長大.
但是她有很厲害的.
藝術天分.
成為了一個將來的.

$^{1481}$一個博士.
但是她讀小學的時候.
有些人笑她.
你說神這麼愛你.
為什麼你會搞成這樣.
她怎樣去回答.
她說我不會.
數我沒有的東西.
我數的是我有的東西.
我的父母是牧者.
他們愛我.
我的朋友.
我的老師愛我.
我覺得自己是神愛的.
我雖然不是像你們這樣正常的.
可以站在那裡.
但是我可以用我的嘴巴.
來畫畫.
後來她自己很勤奮地成長.
在美國成為了一個.
藝術系的博士.
到處去講見證.
是啊.
如果一般人看.
你真是慘了.
你真是悲哀了.
看到你就替你可憐了.
是的.
人一般的眼光就是四肢健全.
頭腦發達.
那就是好了.
家境富裕就好了.
入的是名校就好了.
有好的工作就好了.
有好的父蔭就好了.
沒錯.
沒有人說這些不好.
但是人生不容易.
就這樣.
更多大部分的人.

$^{1521}$絕大部分的人都是一般.
但是.
很不同的地方就是.
那個人他.
體驗到神的愛.
所以真是必須要在.
神的那個.
救贖的歷史裡面.
才能夠.
見到的.
然後就告訴我們.
神要獲得那個尊敬.
那個敬畏.
然後第六節.
你們卻說.
我們在哪些事情上.
污衊你的名呢.
這次就再有一個對話.
就好像一個.
尋找那個答案.
第七節.
你們將那些污衊的食物.
獻在我的壇上.
且說.
我們在何事上.
污衊你呢.
因你們說.
耶和華的卓子.
就是獻祭的祭壇.
就是似乎是那個同祭壇.
是可藐視的.
你們將那些.
核硬.
獻為祭物.
核硬就是有病的.
這不為惡嗎.
將瘸腿的.
有病的獻上.
這不為惡嗎.
你.

$^{1561}$獻給你的省長.
就是那些當時的官長.
他怎麼會喜悅呢.
怎麼會喜悅呢.
豈能夠看你的情面呢.
這是萬君之耶和華說的.
整段聖經.
就涉及到我們敬拜的態度.
首先就是這個是舊約的場景.
不是今天我們回崇拜.
舊約.
都有一個過渡期.
耶穌未來之前.
所有的猶太人.
我們說的就是在猶太人這個群體.
神的選民.
外邦人.
拜偶像的我們不理會.
我們說的就是.
神規定的獻祭.
這就告訴我們.
猶太人.
他們世世代代.
一直都是這樣教導的.
他們要獻祭.
首先就要有代理人.
必須要去到聖殿.
找到那個代理人.
就是祭司.
祭司也有他的手下.
就是利未人.
那麼大衛.
已經有聖殿建好了.
不需要那個流動的聖所了.
就是會幕了.
已經固定在耶路撒冷.
那個聖殿了.
所以不需要再搬搬抬抬了.
那就是說意味著.
祭司呢.

$^{1601}$總之你是亞倫的後人.
你就可以具備成為祭司.
等於你住在圍村.
你是藍兵.
你18歲就可以建房子.
就是這樣的意思.
總之你是亞倫的後人.
你就是.
利美派的亞倫的後人.
你就可以來到成為祭司了.
但不是那麼多事情要做的.
所以大衛真的很厲害.
他幫那些祭司和利未人.
就沒有什麼事情要做的.
獻祭獻不了那麼多的.
那怎麼辦呢.
他安排一些.
因為聖殿裡面不單只是獻祭.
而是有聚會的.
因為固定了那個聖殿不再搬了.
那麼大衛就設立了一些師班長.
訓練這些利未人.
來做師班員.
這樣來到聖殿來敬拜.
分了24班.
來到聖殿裡面來唱頌.
這是歷代至上的經文的教導.
所以大衛其實他本身是一個王.
他是猶大派的一個領袖.
他可以蒙神的喜悅和使用.
他可以任命那些師班長是利未人.
就是說你們如果有這方面的唱歌的恩賜.
你們就不用做那些其他的事情了.
因為真的沒有那麼多事情要做.
不用再搬不用再抬了.
那些器皿擺在聖殿裡面就不會動了.
不會經常動的.
所以就不是說今天我們要戶外崇拜.
沒有的 就是在聖殿裡面.
就設立了不同的工作.

$^{1641}$在聖殿裡面.
那你可以說大衛他的人生.
這些是他自己對神的敬畏.
而他想出來的那種巧妙.
使得他能夠流芳百世.
所以那個代理人.
其中就是大祭司.
和他的手下祭司和利未人.
就會在聖殿裡面來服侍了.
所以如果你是一個普通的以色列人.
你就去到聖殿.
你要獻祭的.
那個時候就不准有代用券的.
就是說你住在藍地.
例如別士巴.
很遠的沙漠.
你要上耶路撒冷.
你同樣都是.
所以是你很不值的了.
你拖著羊.
你要保牠大.
千萬不要走路扭傷.
扭傷呢.
那就上到聖殿.
祭司要驗那個祭物的.
他要驗證這個祭物是不是無殘疾.
如果有殘疾.
祭司是要毀滅的.
要將那個祭物毀滅.
這樣就不行了.
所以為什麼到了耶穌的年代.
他們的聖殿敬拜.
是可以買賣那些牛羊甲子的.
因為有更多的猶太人是住在.
不再是以色列了.
是住在歐洲的地方.
最遠去到羅馬.
但是他們都是敬虔的猶太人.
或者一些歸入猶太教的.
他們都會帶祭身來獻祭.

$^{1681}$但是就長途路遠.
又沒有飛機.
那怎麼能夠這樣讓那個祭身不會受傷呢.
所以路途這麼遙遠也不方便.
所以就有這種去到聖殿.
可以買賣這些獻祭的物件.
等於你去到許願樹.
你就不會自己拿東西去扔.
你家裡包著的金錢便宜一點.
十塊三個這樣.
自己一條繩紮著它.
到時去扔的.
一般人不會這樣做的.
是不敢這樣做的.
一定要去到現場買的.
人家可能幫你祈福的.
你扔上去就扔到就會有福了.
一般人會這樣想的.
所以他們去到聖殿.
就會買一些現成的.
那耶穌為什麼會這麼憎恨.
這些在聖殿附近買賣牛羊的攤檔呢.
原因就是貪污的問題.
那個攤檔是被人高價收購了之後.
再去賣給那些檔主.
所以就導致出現了很多的訛詐.
還有很多背後的黑幕等等.
所以耶穌不是反對人們獻祭.
去買一些現成的.
而是他要將那些.
在聖殿那些人.
只會帶來很多很多的埋怨.
你想想你去到.
我不是教大家去許願樹.
說笑的 當作是一個例子.
不要聽錯了.
你去到.
有沒有搞錯.
一百塊一個橙加一條繩子.
扔上去.

$^{1721}$接著那個人說.
是呀 你不買就不行了.
你自己帶來是不靈的.
那個人都害怕.
那些善信就害怕.
那當然要在現場買了.
但原來現場那些攤檔.
原來不是政府管的.
是當地的徒弟管的.
當然是亂說的.
來到那裡.
誰當越近哪棵樹.
就特別貴的.
就是這樣的.
這個時候.
如果你真的很虔誠.
要去扔那個補牒的人.
你就會有很多的埋怨.
因為原來裡面有這麼多.
衍生了的罪惡.
所以耶穌一樣.
所以我現在回來這裡的時候.
其實任何的以色列人.
都知道一個規矩.
必須.
你要去向神獻祭.
只有兩種祭.
兩個類別的.
但有五個祭.
有哪五個祭呢?.
凡祭,贖罪祭,贖獻祭,平安祭.
接著還有一些叫就職禮的祭.
有很多不同方式的.
你是哪一個類別.
祂就會幫你怎樣獻祭.
這個你不用擔心的.
因為利未記已經教了.
祭司有祂的手冊.
有祂的menu.
記住了.

$^{1761}$譬如我們牧師宣導的規定.
婚禮就是紅色的那個根.
白色就是安息禮拜.
我們不可以混淆的.
兩條東西容易一點.
慘了白色那條剛剛洗了.
用紅色的.
你帶紅色的去安息禮拜.
被人殺了你的份.
有規矩的.
是有安排的.
所以獻祭利未記是很清楚.
講得很清楚的.
其實由利未記的十六章已經很清楚.
講了獻祭那裡.
其中百姓只需要記住一樣東西就行了.
首先你獻什麼祭.
你是感恩.
還是贖罪.
是這兩個類別的.
其他不用你擔心的了.
那你無論是贖罪祭.
感恩祭都好.
需要要求的是什麼呢.
那隻祭生.
必須要.
無殘疾.
就是這麼多了.
誰決定無殘疾呢.
就不是醫生.
不是天父.
是祭司.
祭司決定.
當然是目測.
當然是按他的知識.
當然就不會說.
我先跟他驗一驗腸.
看看他的腸有沒有息肉.
那就不需要做得這麼仔細的了.
只需要從一個目測.

$^{1801}$他有沒有暗病.
有沒有其他.
表面傷痕.
有的.
原來在路上.
那個家主.
被打.
很生氣.
打他.
然後有瘀傷.
那祭司驗到.
屁股那位置.
有邊打過的傷.
瘀了的.
那祭司就說.
這些不是潔淨.
是有殘疾.
那就不可以獻了.
不可以獻就要焚花了.
就當你沒有來獻祭.
那就很麻煩了.
你可以預計到那個麻煩.
獻祭那個人.
我很遠來的.
我已經很有心了.
你現在這樣刁難我.
我現在這樣.
你不准我獻.
如果那個人是.
哪個地方的族長.
是哪個名門的望族.
說我現在大隻祭身.
還是公牛來的.
來到獻感恩祭.
為什麼.
祭司你這麼霸道.
說我的祭身.
有表面傷痕我找不到.
接著你說不准獻.
這些的亂章.

$^{1841}$就會一路滋長.
那個磨心.
似乎是那個祭司.
如果那個祭司只需要.
一隻眼開一隻眼閉.
我現在這麼有心來獻祭.
你收了我就行了.
為什麼你要這麼.
在這裡.
好像用放大鏡那樣.
來看那隻祭身呢.
整段經文你就會看到.
背後的角度.
(咳嗽).
祭司唯一的把關位.
就是那個獻祭者.
最終他離開.
他做完這個儀式.
他整個心態的動機.
他是愛神.
敬畏神.
還是他只是想做一些.
例行公事.
滿足他自己的.
裡面的內心的一些.
錯誤的想法.
例如我要表現.
我自己的富有.
我每次來獻祭都帶很多隻公牛上來.
告訴別人.
我多有錢獻給神呢.
或者.
再有其他的東西.
(咳嗽).
是不一而足的.
是大家人盡皆見的.
因為每個人上來獻祭.
都需要.
讓祭司查問的.
每個人都要的.

$^{1881}$所以.
你可以說那個隊伍.
都很長的.
去到一些特別的節日.
所以在這裡的時候.
大家都會明白.
有一個很大的張力.
在其中.
就是祭司.
是要為神把關.
那個獻祭.
如果.
那些祭物是有層質.
祭司都要.
幫他獻.
祭司都勉為其難幫他獻.
祭司礙於那個.
勢力.
礙於那些人的恐嚇.
都要幫他獻.
那就出現了一種.
此消彼長的.
那個落差.
而導致那個最終.
整件事.
原來不是為神的榮耀.
原來是為人的.
面子來去做一個.
獻祭.
所以我們回看這段經文的時候.
第一個指控.
就是第六節.
最終原來.
那個核心的.
問題的根源.
是為什麼最終令到.
獻祭的以色列人.
來對神有一個.
藐視呢.
根本的原因就是.

$^{1921}$那個祭司本身藐視神.
而導致.
到這一種的後果出現.
所以第一個罪名.
先知這次要.
譴責的.
首先不是以色列人.
是祭司.
來去.
藐視神的.
名次.
神的名.
被藐視.
這個就是第六節.
藐視我名的.
祭司.
神這樣說.
兒子尊敬父親.
僕人敬畏主人.
我既為父親.
尊敬我的在那裡.
是在問祭司.
我既為主人.
你們.
卻說我們在.
哪些事情上藐視你的名呢.
果然第七節.
就是說.
一直是說祭司.
你們將那些污穢的食物.
獻在我的壇上.
這個污穢.
去到什麼程度呢.
少少.
還是很多.
誰去釐定.
那個污穢呢.
咳咳.
這個就變成.
祭司似乎.

$^{1961}$礙於.
人情世故.
就來到去.
將神本來.
很清楚的要求.
來到去降低.
使得讓獻祭者.
不感覺被冒犯.
甚至乎.
還好像感覺.
很大的榮耀.
那就是祭司.
有沒有真正把關的原因.
和問題了.
所以神就藉著先知.
告訴他們.
你們將這些食物.
在神眼中是污穢的.
來到去獻在.
我的壇上.
他們當然就說.
我們獻祭.
是真的有做過檢查.
對那些祭物.
是有做一個詳細檢查的.
我們什麼時候.
在你們獻祭的過程.
對你有藐視呢.
第八節.
你們將.
咳咳.
核硬的獻為祭物.
這個豈不是惡的事嗎.
將那些.
瘸腿的有病的.
來到去獻上.
這豈不是惡嗎.
你們獻給.
你們的省長.
大家怎會喜悅.

$^{2001}$這種有病的祭物呢.
怎會看你們的情面呢.
這是萬君之耶和華說的.
我們就說到這裡.
原來已經一個鐘了.
過半鐘了.
已經九點鐘了.
下一個禮拜我們再繼續這個課題.
還有這個也是一個核心.
關於敬拜的態度.
我們下一個禮拜再繼續這個課題.
還有這個也是一個核心.
關於敬拜的態度.
還有這個也是一個核心.
關於敬拜的態度的一個問題.
這個心都是呼之欲出.
任何一個.
時代的敬拜.
都和這個核心有關連.
我們做一切事.
最終是否想.
神得到喜悅和榮耀呢.
獻祭就做了.
其實.
如果你.
燒過東西吃就知道.
那些煤的肉.
和很漂亮的豬扒.
一燒完.
都是一樣那麼香.
炒了蜜糖之後.
大家一樣那麼香.
但是.
可能吃下去會有些不同.
但是如果你燒盡了.
在祭壇上.
根本就不是人吃的.
是神聞的.
都是一樣.
燒了都是香的.

$^{2041}$但是它真的有殘疾.
你都要燒.
和你獻上上好的.
和獻上好的人.
內心的心態.
都是上好的.
那就.
正正德正.
但是.
心態又不對.
接著還打橫來.
你不准理我.
獻的祭物.
是否真的百分之一百的.
有殘疾.
那你父父就不是德正了.
父父最終原來是.
神得不到榮耀了.
那就是說獻祭者的態度.
和.
所做的那個行動.
是完全息息相關.
所以.
神得不到榮耀.
最終的受害.
是聖殿.
更加會.
荒涼.
人更加會.
對獻祭.
看為是禮行公事.
甚至乎可有可無.
甚至乎就是.
更加會對.
你的神.
會掉以輕心.
和看為.
是我給你面子.
我才來獻祭.
我不給你面子.

$^{2081}$我連這個獻祭我都不來.
那就是一種惡性的循環.
所以祭司只需要.
做回在.
神.
許可認可的.
喜悅的範圍.
其他我們就不需要.
去擔心.
那個獻祭者.
他有多少的獻祭.
神自然會感動.
一個人百姓.
會來找到.
那個.
祂獻祭的祭司.
好我們一起禱告.
多謝天父.
給我們透過.
頭十節聖經.
我們.
掌握到一個很重要的真理.
就是需要在.
神的.
宮牆裡面.
發現你的愛.
尋找你的愛.
和你的愛的.
足跡和真實.
幫助我們能夠.
真的有這一份的參與.
成為了基督徒之後.
不是.
交換去天堂.
而是我們人生.
成為神的兒女.
我們要在你的救恩裡面.
我們去參與.
這個參與過程.
是看到神的.

$^{2121}$愛的一個很重要的證據.
更加讓我們.
開始進入.
舊約.
對於敬拜的.
那個譴責.
對於錯誤敬拜的態度.
那個帶來的.
後患無窮的.
那種的惡果.
的一個.
警告.
我們需要來留心的.
下個星期我們詳細再講.
第六節的時候.
讓我們大家一起有一個.
認真的態度.
對於.
什麼才是合神心意的敬拜呢?.
我們來留心.
我們這樣禱告.
奉耶穌基督的名.
阿們.
\newpage



\section{約翰福音 12:26}
\label{sec:c0kCL7cehvg}
\textbf{2024-01-28 午堂《蒙福的事奉》約翰福音 12\_26 講員:陳一華牧師}
\newline
\newline
連結: \href{https://youtube.com/watch?v=c0kCL7cehvg}{\texttt{https://youtube.com/watch?v=c0kCL7cehvg}} ~~~~ 語音日期: 2024-01-31
\newline
\newline
\hyperref[sec:jLSZtQ5_sq4]{\small{< < < PREV SERMON < < <}}
~
\hyperref[sec:index_chronic]{\small{[返順時目]}}
~
\hyperref[sec:index_scriptual]{\small{[返順卷目]}}
~
\hyperref[sec:FhL_HPKlT9E]{\small{> > > NEXT SERMON > > >}}
\newline
\newline
約翰福音 12:26
\newline
\begin{longtable}{cl}
\hline
\hline
章節 & 經文 (和合本修訂版)\\
\hline
12:26 & \begin{tabularx}{0.7\textwidth}{X} 若有人服事我，就當跟從我；我在哪裡，服事我的人也要在哪裡；若有人服事我，我父必尊重他。」 \end{tabularx} \\ \\
[1ex]
\hline
\hline
\end{longtable}
$^{1}$弟兄姊妹 早晨.
我見到你們個個穿了重甲裝.
就知道今天的天氣都很冷.
包括我在內.
雖然這麼冷的天氣.
但都沒有難阻我們去敬拜主的心.
這個非常之感恩.
今早你看到牧師只說一次經文.
可能會有兩個反應.
一個反應是 怎麼說這麼短?.
真的不夠厚 可不可以說長一點呢?.
第二個反應是 真是好了 牧師很快就說完.
無論哪個反應都好.
我們都一起來看看這次的經文.
這次的經文我們不如一起讀.
因為只有一次 預備起.
若有人服侍我 就當跟從我.
我在哪裡 服侍我的人也要在哪裡.
若有人服侍我 我必尊重他.
這句話是耶穌親口說的.
也都很確實的應許.
若有人事奉上帝的時候.
上帝是會尊重他的.
大家有沒有留意到?.
也都告訴我們 若有人跟隨耶穌.
服侍耶穌 我們是蒙福的.
所以我今天的題目是蒙福的事奉.
大家知不知道成為上帝的子女.
真是有很多福氣 很多恩典.
為什麼呢?除了我們的生命重新得救之外.
原來當我們做神的子女.
聖靈內住在我們的心中.
聖靈做我們的屬靈導師.
聖靈又引領我們明白上帝的心意.
聖靈又給我們有恩賜.
叫我們可以事奉上帝.
可能你覺得牧師我舊底唱歌不行.
一唱起來就走音.
不是走音 是五音不全.
但信了耶穌之後.

$^{41}$上崇拜 聚會 團契 經常唱詩.
正是自己愛上了音樂 愛上了唱詩.
可能不久之後你會加入師班也不一定.
也都很想透過音樂 透過唱詩來事奉上帝.
這就是上帝給我們的恩賜.
各位弟兄姊妹妹.
上帝給我們恩賜.
我們可以事奉祂.
這是獨一無二的.
為什麼我這樣強調呢?.
是因為未認識主的朋友.
他們不明白什麼叫事奉神.
大家認同嗎?.
但我們重新得救基督徒.
我們能夠事奉上帝.
我們就明白其中的意義.
也都告訴我們.
是屬主的子女.
我們才有資格事奉天上榮耀的上帝.
所以我們今天事奉神.
真是很有福氣.
很簡單和大家今早看看.
原來事奉上帝給我們什麼福氣呢?.
第一個福氣就是上帝.
透過給我們的事奉.
我們能夠得到神的尊重.
因為耶穌說:若有人服侍我.
我忽必尊重他.
換句話來說.
當我們今天事奉神.
是神很重視的.
神是欣賞的.
如果我們按上帝的心意.
去事奉上帝的時候.
甚至我們服侍所做的工作.
是上帝建立的.
上帝會建立我們所做的工作.
所以各位各位頂智慕們.
你可以放心.
如果上帝讓你在教會有服侍.

$^{81}$讓你在家庭有服侍.
讓你在社區有服侍.
讓你在職場有事奉.
又讓你可能在一些社交圈子.
又或者在醫院,院舍裡.
有服侍的機會.
頂智慕們記得.
我們為上帝做的一切工作.
只要是按上帝的心意去做的時候.
上帝都會建立我們所做的工作.
所以無分貴賤.
也無分多少.
頂智慕們一說起事奉.
或者我再盡心一點跟大家分享.
事奉這個字在原文.
這個字的字根是什麼意思呢?.
就是一個侍立的意思.
侍立.
大家玩過猜皇帝的遊戲.
或者你看過皇帝的片集.
你會發覺在皇帝的左右邊.
眼下有些人是侍立的.
大家有沒有留意到?.
這個侍立就是一個事奉的意思.
換句話來說.
我想跟大家說事奉是有動態的一面.
也有靜態的一面.
靜態的一面就是告訴我們.
原來我們去敬拜上帝.
我們向上帝崇拜.
這個也是事奉.
又或者我們在家裡為人祈禱.
甚至我們關上門.
在房間很專心的禱告.
很安靜的去祈禱.
這個也是事奉.
所以大家的印象當中.
不要想著事奉是一定動態的.
有人覺得事奉一定很忙碌.
一事奉整桌子都是檔案.

$^{121}$打開時間表.
日程表.
編得很麻麻.
當然這是一個動態的事奉.
但事奉也有靜態的事奉.
所以各位弟兄姊妹們.
無論我們是動態還是靜態的事奉.
只要我們按上帝的心意去做的時候.
聖經告訴我們.
就會得到神的重視.
得到神的欣賞.
得到神的建立.
你說這是不是很有福氣呢?.
這是第一個.
事奉上帝也令我們靈明.
得到成長.
所以為何聖經裡面.
主耶穌鼓勵我們.
要去服侍上帝.
要去事奉上帝.
因為唯有事奉服侍的時候.
你才發現到自己的不足.
自己的缺乏.
所以當你去事奉的時候.
你會發覺過程當中.
原來發覺自己是兩桶水.
加起來也只有半桶水.
你明白我的意思嗎?.
有很多東西我是很不足夠.
如果不事奉的話.
我發現不到.
但我一事奉就察覺到.
原來自己不夠忍耐.
又或者自己對經文的了解.
認識也不是很透徹.
所以當你一邊事奉一邊服侍的時候.
你就越來越察覺到自己.
更加需要追求成長當中.
大家知道.
靈明的成長.

$^{161}$有時候不是一些理論.
而是實戰.
當你真實去做.
實際去經歷的時候.
你會發覺是很不同的.
舉個例子來說.
在教會裡面的事奉.
當你參與教會服侍的時候.
弟兄姊妹們.
你記得永遠都不會是One Man Band.
你認不認同?.
當你在教會服侍的時候.
或者其他機構服侍的時候.
你不會是單獨一個人去事奉.
幾個星期之前.
我參加一個參會.
場面很盛大.
很多位數在當中.
各位弟兄姊妹們.
我見到兩個MC.
一男一女.
我坐在這裡的時候.
覺得這對MC主持這個活動.
真的做得很好.
那些流程.
一個接一個.
井井有條.
間長說什麼話.
說得非常恰當.
各位弟兄姊妹們.
我坐在這裡的時候.
我真的很欣賞這對弟兄姊妹做MC.
如果我是認識他們的話.
就應該封他們做金牌司儀.
我相信都很搶手.
很多人會找他們做MC.
但你記得.
弟兄姊妹們.
這對弟兄姊妹做MC的時候.
他站出來.

$^{201}$台上一分鐘.
下面那句你懂得.
你知不知道他站出來之前.
他要跟很多人開會.
是不是.
他站出來之前.
他要跟很多人商量.
他站出來之前.
他要跟很多人有默契.
怎樣來做.
是否需要.
當你一同默契.
一同合作.
一同去配搭的時候.
那一刻你發覺.
我跟那些人服侍的時候.
這麼多巨星.
我說巨星是性格.
為什麼這麼多人做事這麼有作風.
為什麼這麼多人見解這麼獨到.
為什麼有些人很堅持他的看法.
弟兄姊妹們.
當你去事奉的時候.
那一刻你改變不了別人.
不過你可以改變自己.
可能透過事奉.
上帝就讓你學習.
什麼叫做謙卑.
什麼叫做忍耐.
什麼叫做寬大的胸襟.
去接納別人的意見.
在過程中.
自己就得到什麼.
成長.
自己就發覺到.
是喔.
如果我不去事奉.
沒有這些磨練.
我永遠都不會長進.
但當我跟別人去合作.

$^{241}$一起事奉的時候.
我發覺原來很多地方.
我都要努力改善.
無論你是一個專業人士.
無論你的職級.
怎樣高級行政人員.
但當你在上帝面前事奉的時候.
我們要學習謙卑.
學習忍耐.
學習接納別人的看法.
會不會比自己更加好的呢?.
所以各位弟兄姊妹們.
聖經真的非常豐富的教訓.
教導提醒我們.
基督徒當他事奉上帝的時候.
不單是他人際關係有很多學習.
原來他靈命上有很多成長和追求.
還記不記得彼得後書.
第一章說什麼呢?.
彼得這樣說.
我講幾句.
你們應該記起來.
彼得說什麼?.
有了信心.
就要加上德行.
有了德行.
就要加上知識.
有了知識.
就要加上忍耐.
有了忍耐.
就要加上節制.
一路一路加上去的時候.
弟兄姊妹們.
我們學習的內容豐不豐富?.
很豐富.
所以是透過服侍的時候.
透過事奉上帝的時候.
我們的靈命才會成長.
我們再看第三.
原來當我們服侍的時候.

$^{281}$事奉的時候.
我們能夠多結果子.
大家相信嗎?.
一個人怎麼能夠.
信了耶穌之後.
為神多結果子呢?.
你記得.
不是繞著手.
不是繞著腳.
什麼都不動.
什麼都不做.
你明不明白我的意思?.
這樣是叫做大食懶.
講得現代化一點.
這樣叫做躺平.
叫做吃.
如果我們吃.
我們躺平的時候.
可不可以結出果子?.
不可以.
弟姐妹們.
唯有當我們願意.
踏出事奉的步伐的時候.
你就發覺到.
原來我事奉上帝.
能夠榮耀上帝的.
能夠結出果子的.
有一個弟兄這樣講見證.
他說他初初回到教會的時候.
他不懂.
什麼都不太知道.
但後來教會有對夫婦.
很關心他.
告訴他.
怎樣讀聖經.
告訴他.
怎樣祈禱.
很有耐性地.
一步步地去教導他.
所以他明白得到.

$^{321}$什麼叫做靈修的生活.
他自己就有很大的得著.
這對廉長的夫婦.
不單是教導他.
是示範給他看.
告訴他.
當你有時候收到.
人家要求你.
請求你.
祈禱的時候.
這不時就有了.
你的手機一打開.
你的WhatsApp一打開.
你就經常收到什麼?.
代禱呼籲.
如果你有時候真的很忙的時候.
你記得.
什麼都不要做.
馬上為這件事祈禱.
為什麼?.
怕你會忘記.
說真的.
有時候我們因為太多事做太忙.
一收到人家的代禱呼籲.
代禱的時候.
很快就放進雪櫃了.
為什麼?.
忘記了.
但你馬上去祈禱.
馬上去禱告的時候.
你就能夠為有需要的人去代求.
弟兄說.
因為我這樣的緣故.
我發覺到自己在教會裡.
是很大的得著.
所以我願意.
現在開始.
我都留意一下.
教會裡有沒有初信的朋友.
如果有的話.

$^{361}$我都很樂意跟他們分享.
跟他們一起學習.
以致他都能夠幫到其他人.
弟兄姊妹們.
原來事奉上帝.
可不可以結出果子?.
是可以結出果子的.
因為當你事奉的時候.
你要留意一下.
總是有些人.
什麼?.
盯著你.
總是有些人留意著你.
如果當你的事奉.
能夠樹立一個美好的榜樣的時候.
其實你的服事就激發起.
或者激勵到其他.
初初十五,十六都不知道.
可不可以事奉的人.
因為你的榜樣.
又能夠得到鼓勵.
得到激勵.
而他都樂意去事奉上帝.
所以各位弟兄姊妹們.
很多時候.
這種就是生命的果子.
你影響了其他人.
好像一個生命影響生命一樣.
能夠多結果子.
這個就是上帝讓我們.
透過事奉的時候.
能夠為他結出美好的果子.
好.
我們再看第四.
原來透過事奉.
亦都能夠經歷恩典.
各位弟兄姊妹們.
你相信嗎?.
很多時候.
你發覺到.

$^{401}$初初你去事奉上帝的人.
有時候不夠信心.
亦都不夠把握.
不知道能否做到.
心情好像患得患失一樣.
但是你樂意.
很踏實地去服侍上帝的時候.
你發覺做著做著.
事奉的時候.
你發覺越事奉.
就越有恩典.
你發覺越事奉.
就越有喜樂.
你發覺越事奉.
上帝又加添你智慧.
讓你能夠知道.
如何做一些抉擇.
你又發覺當你去事奉的時候.
上帝又加添我們的心力.
寶不寶貴?.
弟兄姊妹們.
原來越事奉是越有力量的.
越事奉是越有恩典的.
所以我曾經聽過一個姐妹.
這樣來講見證.
讓她信了耶穌之後.
她覺得自己的生活.
是非常充實.
為什麼?.
因為她說我上班就上兼職.
但我下班就回家照顧家人.
但是我也可以抽些時間.
上教會參與事奉.
參加祈禱會.
甚至她也可以抽到時間.
跟弟兄姊妹們出去街頭報道.
做街報.
原來越事奉她就越有衝勁.
越事奉她就越起勁.
越有力量.

$^{441}$越有恩典.
當我聽完這位姐妹分享見證的時候.
為什麼呢?.
因為她真的從服侍神.
關懷人的過程裡面.
她體會到上帝的恩典,能力,智慧.
是源源不絕.
淋到她的身上.
所以各位弟兄姊妹們.
你說我們今天能夠事奉上帝.
是不是很蒙福呢?.
是不是很大的恩典呢?.
講到這裡.
可能你想起.
我也很想去事奉美丙.
弟兄姊妹們.
是的.
我們真的要起來.
要服侍上帝.
無論上帝讓你在任何崗位上.
我們都要去把握.
去事奉上帝.
千萬不要幻想以下這種病症.
在我結束的時候.
我想告訴大家.
我們基督徒.
有時候也很容易幻想這種症狀.
一起講一次.
什麼症?.
頭重腳輕症.
希望我們不要幻想這種症.
什麼叫頭重腳輕症呢?.
我舉個例子你就很明白.
我認識有一個弟兄.
他初初信耶穌之後.
很熱心,很有憧憬.
所以他信耶穌沒多久.
就去圖書館.
找很多書去看.
然後自己上網.

$^{481}$又找很多資料.
然後他又打開YouTube.
去聽很多名目的講道.
專題講座.
他搜集很多資料.
什麼叫宗教比較.
什麼資訊.
他都拿到很多資料回來.
我想問大家.
這位弟兄這麼做.
OK嗎?.
OK.
你千萬不要說他不OK.
OK的.
可以的.
不過後來我觀察這位弟兄.
他沒有實踐他所領受的.
他沒有去行道.
但他就什麼?.
他就只是拿很多資訊.
只是拿很多資料.
只是拿很多的知識.
他沒有將它融化成為他行道.
服侍上帝的一個很重要的步驟.
所以整整下.
我發覺這位弟兄連教會都沒有了.
為什麼呢?.
因為他沒有經歷上帝的恩典.
很可惜.
他學到很多東西.
但最後他連教會.
連崇拜都沒有參加.
這種是什麼病症呢?.
這種叫頭重腳輕症.
我在醫院工作.
有時候我都觀察一個現象.
都很特別.
我發覺那些臨出院的病人.
很多時候醫生巡完房.
就在排版上撩幾個字.

$^{521}$寫什麼字呢?.
就是寫這個病人去接受物理治療.
大家明白這個意思嗎?.
為什麼這個病人明明都快出院了.
還要去做物理治療呢?.
答案是因為這個病人.
在床上睡了一個星期.
睡了兩個星期.
他的腳還沒有力.
如果他要出院.
一腳起來.
接著就整個掉在地上.
你相不相信呢?.
因為他的腳兩個星期.
一個星期都沒有走過.
他的肌肉萎縮.
他沒有力.
頂指沒有門.
請你記得.
我們不要做這類的肌肉圖.
你問一問自己.
我上一次跟別人傳福音的時候.
是什麼時候呢?.
如果你告訴我.
牧師真的好了.
我很開心告訴你.
我上一次跟別人傳福音.
剛剛是上個月.
為什麼呢?.
因為上個月是聖誕節.
我出去報佳音.
我出去派單張.
我出去跟別人傳福音.
但如果你說.
牧師.
我上一次跟別人傳福音.
已經是一年多前的事.
我們就要什麼呢?.
要好好反省.
去反思.

$^{561}$我們會不會患上這種症狀.
可能我們要問問自己.
上一次.
我是什麼時候事奉上帝呢?.
上一次.
我是什麼時候服侍神呢?.
如果你說.
牧師.
我每個月都有參加服侍.
我每個月都事奉.
那就非常好了.
但如果你說.
沒有啊 牧師.
我上一次在錦繡堂事奉.
是一年多前的事.
或是幾年前的事.
之後我就沒有參與事奉.
如果是這樣的時候.
頂指沒有門.
要留意.
要小心.
不要患上這種頭重腳輕症.
因為當我們不去事奉上帝.
不去服侍神的時候.
你會發覺到.
我們屬靈的生命.
就會不斷不斷萎縮.
所以唯有當我們事奉上帝的時候.
我們的福氣.
我們的恩典是更加多.
更加大.
求主祝福我們眾弟兄姊妹.
我們一起低頭祈禱.
親愛天父.
多謝贊美你.
因為在聖見當中.
你不單止給我們有很豐盛的生命.
也讓我們的生命能夠活出來.
讓我們的生命成為一團火.
這團火叫我們熱心愛神愛人.

$^{601}$叫我們努力事奉上帝.
也樂意關心人的需要.
叫我們能夠將聖經所學到的.
能夠學以致用.
求天父上上提醒我們.
因為你應許過凡是事奉上帝的人.
你是必尊重他.
求主讓我們的生命.
更加成為多人的祝福.
奉耶穌名求.
\newpage



\section{路加福音 23:40}
\label{sec:FhL_HPKlT9E}
\textbf{2024-01-28 早堂《談敬畏》路加福音 23\_40 講員:陳韋安牧師}
\newline
\newline
連結: \href{https://youtube.com/watch?v=FhL-HPKlT9E}{\texttt{https://youtube.com/watch?v=FhL-HPKlT9E}} ~~~~ 語音日期: 2024-01-31
\newline
\newline
\hyperref[sec:c0kCL7cehvg]{\small{< < < PREV SERMON < < <}}
~
\hyperref[sec:index_chronic]{\small{[返順時目]}}
~
\hyperref[sec:index_scriptual]{\small{[返順卷目]}}
~
\hyperref[sec:Srg_1OhzOG8]{\small{> > > NEXT SERMON > > >}}
\newline
\newline
路加福音 23:40
\newline
\begin{longtable}{cl}
\hline
\hline
章節 & 經文 (和合本修訂版)\\
\hline
23:40 & \begin{tabularx}{0.7\textwidth}{X} 另一個就應聲責備他，說：「你是一樣受刑的，還不怕神嗎？ \end{tabularx} \\ \\
[1ex]
\hline
\hline
\end{longtable}
$^{1}$鄧小平早安.
前陣子有一段時間.
和一位教會的前輩吃飯.
吃飯的時候.
驚訝的解釋突然間問我.
你覺得崇拜什麼.
突然間這樣問.
有點奇怪.
你覺得崇拜什麼.
後來他就提起.
幾個月前留堂.
在星期六晚的一間比較年輕的教會.
我們有一個特別的崇拜主題.
叫做開學崇拜.
這個崇拜也有一段時間了.
我們邀請鄧小平一起穿上中學校服回來.
一起去崇拜.
全場有很多學生穿上校服回來一起去崇拜.
不說禮儀問題.
也是一種禮儀.
反而令我很難忘的.
就是這位教會前輩.
一句這樣的宿令忠告.
他跟我說.
陳偉安你要學習敬畏.
只要你是敬畏上帝的話.
你怎樣走都不會錯.
這句我覺得是一句很受用的說話.
對我很大的提醒.
所以後來我就預備了一些經文的研究.
有關敬畏這個課題.
今天也跟大家分享敬畏這個題目.
不知道你對於敬畏有什麼認識.
信了主這麼久.
當你提起敬畏的時候.
可能有一些對於敬畏這個題目既熟悉又陌生.
Fear of God.
英文叫做Fear of God.
一直也知道有這種態度.
不過似乎教會又不是經常提及.

$^{41}$如果你在網上Google搜索.
Fear of God這三個字的時候.
大家知不知道是怎樣的呢.
反而是一個時裝品牌.
其實也不便宜.
這個品牌裡面也有幾千元一件衣服.
原來創辦人也是想用Fear of God.
這種比較出於信仰的角度來開展品牌.
不過後果就剛好相反.
成為了一個比較奢華的意味的態度.
敬畏是一個我們今天的年代比較容易忽略的課題.
一個我們越來越少提的課題.
今時今日我們在教會裡面.
越來越傾向成為一種安慰宗教.
我們也強調憐憫,恩典,拯救.
恩典太美麗.
耶穌很愛你.
當我們很深深知道.
其實也是正確的.
耶穌上帝很愛我們的時候.
敬畏是一個我們不容易消化的題目.
事實上如果你去找一些書籍的時候.
你會發覺十車其九裡面的書.
都是叫你不要害怕.
不要懼怕.
愛裡面沒有懼怕.
要剛強壯膽不要害怕.
因此面對海量裡面的憐憫,恩典,拯救的訊息.
你實在不容易明白.
面對我們的神.
有些什麼是恐懼.
有些什麼是要畏.
有些什麼是要敬畏.
你都很難去拿捏敬畏的意思.
因此剛才讀的經文是一段很簡單的說話.
你還不怕神嗎?.
一段來自於路加福音裡面.
十字架上的一段說話.
另一個強度來問.
另一個強度.

$^{81}$你還不怕神嗎?.
這是臨死前的強度.
一個非常嚴肅的問題.
正是今天我們一起來思考的課題.
所以今天我們不只是講這段經文.
而是一個主題式的講道.
我們來認識聖經裡面敬畏的意思.
在舊約來說.
敬畏這個字眼.
其實叫做Yarit Elohim.
即是敬畏神.
敬畏其實是以色列人裡面.
對於信仰裡面最基本的課題.
《申命記》第四章第十節.
耶和華對摩西說.
我要叫他們聽見我的話.
使他們存活在世的日子.
我們可以學習敬畏我.
又可以教訓兒女這樣行.
學習敬畏耶和華上帝.
在以色列的信仰舊約裡面.
一個最基本的傳承.
不過對以色列人來說.
怎樣才叫做敬畏呢.
所謂敬畏對猶太人來說.
就是一個人不任意妄為.
是無形中一個非常崇高的上帝.
在他的心目中裡面.
所以這種敬畏去到兩約之間的時候.
Yarit Elohim就變成一種.
Yarit Shalayim.
對於無形中的上天.
有一種敬畏的心態.
我們未必能夠見到一個.
我們親眼見到的神.
不過對於他的存在.
無形中我們知道感受到他的存在.
從而去敬畏.
從而去影響我們的行為.
這種的看法.

$^{121}$中國人來說大概叫做什麼.
就叫做為天.
我們都有這樣的論語說話.
君子有三畏.
為天命 為大人 為聖人之言.
所以這種敬畏的態度.
我們都見到有幾段的.
舊約的小小的片段反映出來.
大眾經文就是初夏及記第一章.
講到兩位修身婆.
在河邊裡面.
去拯救以色列男孩的例子.
聖經說什麼.
修身婆敬畏上帝.
不按照埃及王的吩咐行.
竟存留男孩的性命.
聖經強調這兩位修身婆.
因為是敬畏神的緣故.
她就冒著危險公然抗命.
不去執行皇帝的命令.
反而是違反了掌權者的要求.
實不知你們為什麼要做這件事.
存留男孩的性命呢.
很明顯他們存留男孩性命的原因.
正正因為他們是敬畏神.
所以敬畏是成為了一個很大的約束.
給我們有勇氣去決定不做一些事情.
這是一個很大的例子.
這個例子剛剛相反.
就是阿巴拉罕去獻以殺的故事.
阿巴拉罕去獻以殺的故事.
我們經常都記得有關信心.
不過如果你去看經文的時候.
創世記念裡面其實記載著.
阿巴拉罕去完成這個試驗的時候.
耶和華是這樣跟他說的.
你不可在這童子身上下手.
一定不可害他.
現在我知道你是敬畏神了.
原來阿巴拉罕去獻以殺.

$^{161}$是一個敬畏的故事.
阿巴拉罕因著敬畏上帝.
堅決去做出一些.
可能是旁人覺得很瘋狂的事情.
因此修身婆的故事.
是一個堅決叫我們不要做的事情.
阿巴拉罕的故事剛好相反.
叫我們無論如何都堅決地去做一些事情.
兩件事都是很超乎我們的想像.
甚至乎是一些很瘋狂的事情.
不過我們還沒有回答這個問題.
對於猶太人來說.
究竟敬畏的畏是什麼呢?.
我們究竟要害怕什麼呢?.
我們要去fear什麼呢?.
我的意思是究竟那兩位修身婆.
或者是阿巴拉罕.
他們在那個過程當中.
他們究竟是害怕什麼呢?.
難道說那兩位修身婆.
因為很害怕上帝會擊殺他們.
所以相比起法老政權.
他們寧願更加害怕上帝.
所以就寧願去救小朋友.
難道就是說阿巴拉罕出於恐懼.
恐懼上帝又會罰他.
或者是會擊殺他.
就寧願殺自己的兒子.
當然不是.
猶太人的敬畏觀告訴我們.
fear of God這種態度.
不是一種叫做fear of punishment.
不是因為害怕被罰.
不是因為害怕受刑.
或者是害怕落地獄.
我想大家也是.
大家去敬畏神的時候.
或者大家願意回來教會做基督徒.
不是出於將來自己不去崇拜.
就會落地獄.

$^{201}$不會因為害怕自己被神罰.
或者做不到基督徒.
就會害怕被神擊殺.
這是我們應該做的事情.
我想大家的夢想或實際見證.
也不是這樣說的.
我們都是因為一種出於上帝的救恩.
出於上帝的愛.
往往是一種pull factor.
多於一種push factor.
不是因為害怕做不到.
所以會被罰而去做.
而是單單出於上帝的那份吸引.
可能你不知道有沒有試過.
我初信主義時也有這樣想過.
有時候偶爾肚子痛.
在廁所蹲的時候.
也會突然間想自己是不是做錯事.
不知道大家有沒有試過.
是不是我最近有什麼得罪了神.
神懲罰我去肚子痛或病.
可能你也會相信一種所謂報應神學.
因為自己做錯了.
或者自己不合神的心意.
所以就肚子痛或大病一場.
但是你知道這樣的敬畏.
絕對不是一種長遠的做法.
純粹基於一種懲罰的恐懼.
或者純粹一種功利的心態.
害怕上帝.
所以才選擇去做一些聽從神的決定.
這個不是真正的敬畏.
這是一種消極的恐懼.
或者是因為懼怕.
羅馬書3:3說.
作官的原不是叫行善的懼怕.
乃是叫作惡的懼怕.
上帝對於政權正正是這樣.
以配見就叫作惡的懼怕.
但是上帝自己並不是這樣.

$^{241}$這個敬畏我們就看回.
根據古代的教父很清楚地說.
這種是叫做奴僕的恐懼.
他認為我們面對上帝有兩種不同層次的敬畏.
最起初的時候.
第一種敬畏是奴僕的敬畏.
純粹是害怕上帝的緣故.
因為律法 因為上帝會懲罰而產生敬畏.
但他認為當我們越去成為基督徒的時候.
我們越來越去經驗上帝的愛的時候.
第一種的敬畏就會慢慢地消失.
取而代之的就是對於上帝的愛.
那種敬畏的態度.
正如張先生所說 愛裡沒有懼怕.
愛的完全就把懼怕除去.
所以這種敬畏 第二種更高層次的敬畏.
就叫做Trust the fear.
就叫做一種對於愛裡的敬畏的感覺.
這種昇華的層次.
這種敬畏正正是基於上帝的愛.
而不是一種純粹的恐懼.
因此我們對於上帝的愛和敬畏其實是分不開的.
Love of God和Fear of God是分不開的事情.
是用於同一件事情.
實際上有一個拉丁文.
就叫做Mysterium Tremendum Adversitens.
意思是一種令人既可畏又吸引的奧秘.
叫做令人既可畏又吸引的奧秘.
這句話正正是在描述上帝那種的奧秘.
正正因為上帝太過超越.
上帝的恩典 上帝的偉大 上帝的慈愛太過厲害.
永遠遠遠離到他凌駕於我們.
所以當我們面對的時候.
他同時間產生一種震驚而又羨慕的那種感受.
真正的敬畏是那種驚喜和所謂的恐懼.
是同時間出現的一種複雜的感受.
正如我們唐伯虎聽秋香首歌.
我的心又喜 我的心又慌.
這種快樂和敬畏的高不可攀的感受.
正正是我們面對上帝那種感受.

$^{281}$面對上帝本身的高度.
他的尊崇 他的榮耀.
上帝不需要做任何事情.
他的權威 莊嚴和崇高已經震懾著天.
叫我們帶來一種敬意.
即是近之的那種敬畏.
再一次來認識他.
所以這種恐懼又驚喜的坐鐘交集.
正正是詩篇第42篇所說.
你的瀑布發聲 深淵就與深淵響應.
波浪洪濤漫過我身.
八周夜我必向我師前外.
黑夜我要用歌頌禱告至我生命的神.
上帝的瀑布強而有力.
告訴你他那種龐大無比的力量.
上帝的深淵淹沒了你在十萬八千呎的水底之中.
上帝的波浪好像一個巨大的海嘯一樣.
遠遠覆蓋私人的生命.
上帝的氣場過盛.
他的慈愛淹沒我們.
我們同時間被上帝的慈愛同在所包圍.
這種大概就是叫人既震驚又羨慕的奧悲.
今時無論如何我們都不能夠失去這份震驚和戰慄的態度.
作為一個基督徒或一個新教徒.
我們不會害怕fear of punishment.
我們知道上帝恩典一次恩順稱義.
拯救我們.
我們都不會害怕落地獄.
不會害怕無端端做錯什麼被上帝擊殺.
但不代表我們不會有這份稱為神聖的恐懼.
單單親近上帝而又保持一種這樣的高度.
正正就是一種真正的敬畏.
正所謂耳熟耳知所講.
你要以敬畏而護華為之寶.
這是我們信仰裡面非常重要的一步.
它是我們信仰的第一步.
亦是關鍵的第一步.
你永遠都不能夠失去這個最關鍵的第一步.
可能大家都聽過奇恩典的詩歌.
這首詩歌的寫作其實是一個這樣的背景.

$^{321}$1748年的時候.
作者John Newton就面對著一個非常之可怕的暴風雨.
當時他在船上.
幾百年前去大西洋的船.
突然間一個暴風雨的來臨.
那艘船真的要快要沉在海底裡面.
面對著死亡的害怕和那種恐懼.
John Newton就大大地去呼求上帝.
後來就寫了奇恩典這首詩歌.
很有趣的.
如果你去留意詩歌英文版本的歌詞的時候.
第二節歌詞這樣寫.
The trace grace that taught my heart to fear.
And grace my fear relief.
How precious did that grace appear.
The hour I first believed.
當中講到上帝的grace.
上帝的恩典正正是去教我的心去恐懼敬畏.
同時間這個恩典又將這份恐懼排除.
當然奇恩典裡面所講的恩典.
永遠都是我們信仰的主線.
不過這首歌告訴我們.
沒有了那份恐懼和敬畏.
我們單單有所謂的恩典.
這份恩典就不是真正上帝而來的恩典.
這份恩典是出於我們那份敬畏和害怕的背景之下.
才是有意義的.
沒有了這份敬畏的高度.
任何純粹講親密.
耶穌疼愛你 哭泣你.
不要生氣 不要不開心的恩典.
就沒有了真正的意義.
所以這是猶太人最基本的信仰價值.
一切耶和華的信仰都起始於這一點.
所以有一句拉比的名言.
簡單來說.
任何人沒有了敬畏的時候.
就好像得到一份寶藏.
而欠缺了一條鎖匙一樣.
簡單來說就等同於.

$^{361}$你有一個不知大家是否熟悉的.
有一個比特幣的戶口.
很多錢 但你沒有密碼一樣.
等同於你有一個很大的房子.
但它沒有了一條鎖匙一樣.
所以我們今天重新思考.
我們自己是否一個敬畏神的基督徒.
面對我們今天的世代.
正是一個很大的提醒.
路加福音這樣提過.
那殺身體以後不能再做什麼的.
不要怕他們.
我要指示你們當怕的是誰.
當怕那殺以後又有權並丟進地獄裡的.
我實在告訴你們 正要怕他.
五個麻雀不是賣二分錢嗎.
但在上帝面前一個也不能忘記.
就算你們的頭髮也被數過.
不要懼怕 你們比麻雀還貴重.
敬畏的意思是什麼呢.
我們要好好地去搞清楚.
我們應該懼怕的對象.
我們應該懼怕真正你要懼怕的東西.
搞清楚你要敬畏的對象.
這樣我們就沒有什麼好害怕.
面對著一個很不容易的時代.
很多事情都令我們很害怕.
很多都擔心 怕踩界.
怕做錯了 或者怕認罪.
但是我們知道我們真正要去恐懼的.
真正要怕的.
當我們找到真正要敬畏的上帝的時候.
其他一切 聖經說.
我們就好像那句安慰說話.
連麻雀上帝都看顧的時候.
我們就沒有什麼可以懼怕.
事實上當我們再回到社會的時候.
更加重要的可能就是我們自己.
信仰裡面的初心.
正如我所說 敬畏是第一步.

$^{401}$也是最關鍵的第一步.
因為一切上帝的知識.
我們信仰的每一個之後的內容.
都是建基於這一步.
當你不能爬上去之後.
就扔下它.
當我們信仰的歷程.
回到教會 做了那麼多基督徒之後.
我們反而往往丟棄了最基本的那一步.
不知道你是否最敬畏上帝的時候.
正正就是你初信署的時候.
剛剛認識神.
對於聖經有一種默然的敬畏.
好好地來去珍重去看祂.
對於教會 對於每一件信仰的事情.
都是非常之鄭重認真的去面對.
其實我自己在受院經驗都是這樣.
通常那份敬畏的態度.
往往是隨著神學生的級數越高就會下跌.
不知道是否這樣.
一年班是最敬畏的.
大家剛剛進來神學.
他們就很敬畏地來參加學院每一個活動.
早上有六點半靈修.
慢慢二年班 三年班熟了.
慢慢就開始試功課.
有時候又會有很多夜夜發校規.
所以我要說我們的信仰.
往往都是這樣.
保持我稱之為一份菜鳥的敬畏.
純粹一個新手入門.
剛剛信耶穌.
剛剛認識神.
那種敬畏的態度.
不要因為自己的經驗.
閱歷豐富了.
人大了.
就忘卻這份敬畏.
特別是我們做人.
當我們的人生豐富經驗的時候.

$^{441}$基本上世上沒有什麼難倒你.
起碼你知道怎樣處理.
大概這些事情怎樣去處理.
你就開始不是太害怕任何東西.
不是太害怕任何的東西.
但是對於上帝.
一個永遠都是超越我們.
掌管我們生命.
永遠都是高過我們很多倍的上帝.
我們仍然保留這份敬畏的態度.
我想錦繡堂這裡.
我們不是天主教.
或者是那些什麼地方.
這裡沒有一塊神聖的碑.
要你去拜的.
也沒有一些聖水.
或者是一些聖物.
要你去摸摸的.
錦繡堂這邊沒有一個神聖的石頭.
要你特別去敬畏它.
這塊神聖石頭.
在我們每一個基督徒的心中.
你要好好地去保護著.
保存著這塊神聖石頭.
放在你的心裡.
將它帶入你自己工作的地方.
將它帶入你自己的家庭裡.
一群人的時候.
一個人的時候.
仍然好好地去保持著這塊神聖的石頭.
好好地來保護著它.
帶著它.
正如那位淑玲前輩跟我說.
都是跟大家共勉.
只要你保持著敬畏.
你怎麼走都不會走錯.
一起去的祈禱.
因為你很愛我們.
同時間你卻是一個值得我們去敬畏的上帝.
因為你的能力.

$^{481}$你的偉大窮滄.
遠遠地去壓榨我們.
讓我們在今天.
我們重新去思考.
我們對你的那份敬畏.
我們愛你.
同時間敬畏你.
讓我們好好地去.
讓這份敬畏影響我們的生命.
警戒我們不去做一些不應該做的事.
更加去回應.
要做一些你命令我們去做的.
更加讓我們知道.
我們面對著這個世界的時候.
我們當敬畏的是什麼.
讓我們有勇氣.
面對著我們生活裡面.
不同的艱難.
不同的對象.
讓我們用這份敬畏去面對任何的事情.
求你幫助我們.
逢主命求.
阿們.
\newpage



\section{瑪拉基書}
\label{sec:Srg_1OhzOG8}
\textbf{2024-02-01 周四晚查經班《瑪拉基書(4)》 講員:曾裔貴牧師}
\newline
\newline
連結: \href{https://youtube.com/watch?v=Srg_1OhzOG8}{\texttt{https://youtube.com/watch?v=Srg\_1OhzOG8}} ~~~~ 語音日期: 2024-02-02
\newline
\newline
\hyperref[sec:FhL_HPKlT9E]{\small{< < < PREV SERMON < < <}}
~
\hyperref[sec:index_chronic]{\small{[返順時目]}}
~
\hyperref[sec:index_scriptual]{\small{[返順卷目]}}
~
\hyperref[sec:vJ5BK_GiIWM]{\small{> > > NEXT SERMON > > >}}
\newline
\newline
$^{1}$我們一起祈禱.
再一次感恩天父.
讓我們可以學習聖經的真理.
盼望我們可以透過《馬拉基書》.
這本書已經二千四百多年了.
讓我們能夠從舊約.
特別是一些關乎到宗教敬拜的系統.
當時他們在崇拜裡面.
和神的關係.
和在這個民族一直發展之後.
神對於猶太人有什麼心意.
都讓我們能夠掌握,明白.
也讓我們自己在今天.
我們是一個二十一世紀了.
科技可以帶給我們.
大家可以在同步.
在不同的地方都能夠同時間學習.
在古代是不可能的事.
但也幫助我們大家的心靈.
正正就是我們能夠在任何地方.
只要有網路.
我們就可以大家的相交.
更加肯定就是.
做物主是無所不在.
是能夠隨時和我們有相交的.
幫助我們對於你是一位這麼榮耀的真神.
我們有認識.
我們這樣祈禱.
奉耶穌基督的名.
阿們.
我們回來看《馬拉基書》.
最後這三章聖經.
上一個禮拜我們都看了第一章.
可以再有些補充.
大家一齊再去溫習.
(補充時間).
希望我盡量講得溫柔些.
細聲些慢些.
這樣喉嚨不會太過繃緊.
就會咳嗽.

$^{41}$這個禮拜是好了很多.
(補充時間).
連續幾次講了.
這本書的其中一個特別之處.
就是先知代表神提出五至六個問題.
其實是預備有答案的.
是有model answer.
其實主要是譴責整個以色列民族.
很重要的一個問題.
就是沒有尊神為聖.
變成在很多事情上.
似乎有種金屬疲勞.
慢慢有些事情已經劃了瓦.
所以在不同的經卷裡.
都浮現出神對他們的責備.
如果這本書.
基本上大部分的解經家都認為.
《馬拉基書》是被擄歸回的最後一卷書.
所謂最後一卷就是神的默示.
給猶太人的經卷,經典.
所以之後.
要去到馬太福音的年代.
才有新約的啟示.
所以這個叫做兩約中間的.
大概四百年的歷史.
就是沒有默示的了.
所以這本書也在在告訴我們.
神對於猶太人.
記住比較重要的信息就是.
在主耶穌還沒降生之前.
還沒來到這個世界之前.
神仍然是和猶太人.
信守著一個約.
這個約就是.
例如是亞伯拉罕的約.
摩西的約.
西乃山的十誡的約.
所以神仍然是要求他的子民.
即是猶太人.
要來信守摩西的律法.

$^{81}$所以現在在那個標準.
怎樣才算對,錯.
怎樣才算神的心意.
這個就一定是在摩西的律法書.
那個精神.
那個要求下.
來作為準繩的標準.
所以往往就是針對他們.
那個已經是好像已經劃了牙的敬拜生活.
敬拜生活就是源於對神.
他的榮耀慢慢遞減了.
就會帶來在那個人應該負的責任的.
那個敬拜的可能叫儀式.
可能就是他們去到神面前的.
應該有的.
當然是先從裡面的態度.
到外面的行為.
來到去是神所關注.
是先知書.
是馬拉基先知.
來到去奉神的名.
當然是用一種責備的口吻.
來叫以色列人有一個回轉.
當然到底有多大的轉變.
聖經就沒有交代了.
進入了一個不再有神的默示的四百年了.
猶太人都認定了.
這本書就是那個天安門的一個結束.
不會再有.
當然兩約是有很多次經的.
但是那些次經是不能夠成為正典的.
所以我們要有這方面.
可能同步的對於聖經的默示.
都有少少的概念就可以了.
但是我們回到經文的世界.
在這個時候.
神仍然透過先知.
有一些的教導.
是有一些對民族的要求.
有要求就是仍然對他們有愛.

$^{121}$仍然對這個民族有他的計劃.
當然人就是始終是人.
他不會知道原來神的計劃.
不僅僅是這麼小的民族.
只是在耶路撒冷的一個.
其實可能那時候也只是幾十萬人.
這麼小的一個民族.
就是神全部的計劃.
這個是猶太人想多了.
以為自己就是得天獨厚了.
但不是的.
神的計劃仍然保留這個民族.
但是其實輾轉他們都要面對很多的.
皇后興起的那些.
其他的周邊國家的戰爭.
帶來的他們的困難的.
不過就不是現在馬拉基.
寫先知書的時候那個時期.
因為你要再看其他的兩約的歷史.
你就會知道有西留基的皇朝興起.
有埃及的皇朝興起.
有多利邁的皇朝興起.
那些正正就是令到當時近東.
特別是巴勒斯坦地.
是他剛剛夾在一個中心點.
幾個帝國的中心點.
所以是會帶來他們要面對很多的困難的.
所以這些輾轉.
譬如有瑪嘉比一二三書.
就是《刺經》來的.
很出名的《刺經》.
是主前的一百六十幾年的時間的作品.
但是我記得《刺經》不是《聖經》的正典.
是那個時代的歷史的書卷.
是真的有人寫的.
有一個叫做馬提亞家族所寫的.
來描述他們怎樣來抵抗.
當時的埃及亞古四世的事跡.
後來衍生了一個修奠節.
設立了一個新的節期.

$^{161}$就叫做修奠節.
這個Halucca.
這個是往後.
先知再過多二三百年之後.
就出現的歷史.
所以我們要記得.
我們現在看舊約.
去到一個這麼後的時期.
其實神的計劃從來都沒有停止過的.
繼續在地上.
當然既然耶穌還沒有來這個世界.
神仍然集中的默示,顯現,工作.
都是會透過猶太人的.
因為他已經選擇了這個民族.
所以你會看到.
耶穌基督都是由猶太人那裡而出的.
新約的作者都肯定的.
正正告訴我們.
原來就算人怎樣看自己的民族.
去到可能很巔峰.
很盛極一時.
或者去到很黯淡.
整個民族落到一個很困難的境況.
甚至乎好像覺得神都不再理會他們.
但是我們寶貴的就是我們後人.
我們是二千四百多年後的人.
我們就有聖經新舊約傳書.
給我們看到神的作為.
是從來都沒有停止的.
所以如果我們有這個眼界.
這樣去看我們還沒有發生的歷史.
就是明天,後天,將來的明天.
我們就要有這個信心和盼望.
神是繼續作供的.
不會停止的.
真正的停止就是末世論的說法.
我們完了《馬拉基書》之後.
我們之前已經投票過了.
我們就會講《帖撒羅尼迦前後書》.
這兩本書其中特別的地方.

$^{201}$是保羅很早期的傳道.
第二次的傳道寫的書卷.
很早期的.
但是整兩本書裡面.
每一章都講主在來.
很奇妙的.
所以研究末世論.
就一定要看這兩本書.
另外就是.
既然有機會講到關於末世主在來.
所以就會用一種系統神學的角度.
和大家去探討一下末世論.
有什麼主張.
有什麼派別.
給大家有個概念就可以了.
不然將來你去讀一些神學課程.
就會比較詳細一些.
我們不要講這麼遠了.
回到《馬拉基書》的時候.
我們就在這個.
所以大家要回到歷史當下.
就是馬拉基的時代.
我們要構想自己是一個猶太人.
是神立約的子民.
我們要有這個心情.
去聽先知跟我們講的話.
你要想法就是.
你現在是一個同神立約.
你去獻祭呢.
你敬拜神呢.
你是帶著牛羊去到聖殿的.
你要有這樣的幻想.
這樣進入經文的世界.
因為先知就是這樣.
跟那些猶太人這樣講話的.
所以我們如果有這種聯想.
當然我們這個是一個.
先進入經文的世界.
但我們不是那些人.
我們是現代香港人.

$^{241}$所以我們要回到現代.
我們要想想.
我們每個禮拜的敬拜.
是不是蒙神的月納呢?.
這是最重要的這本書帶給我們的問題.
還有什麼問題令猶太人得不到神喜悅的敬拜呢?.
我們就要避免出現這樣的犯錯.
所以看《馬拉基書》.
其實收穫是很豐富的.
因為至少對一個敬拜者.
你行常每個星期.
甚至你和我都是事奉者.
不單是敬拜.
我們有參與事奉.
接下來第一個禮拜.
在座的有些會獻唱.
有詩詩.
有其他的工作.
我是台上報告.
祝福.
我們所做的.
是不是蒙神的月納和清許呢?.
這本書第一章就罵得很厲害.
簡直是一無是處.
那個敬拜.
我們有責任要擔心我們的敬拜.
是不是一無是處.
是不是在神眼中的一文不值.
所以我們看這些書卷.
不要有太過抽離感.
真的要墮入他們的世界裡面.
因為做物主只有一位.
他要求的敬拜是怎樣的呢?.
我們就要去問.
我們現在是否在他.
符合他的心意的敬拜呢?.
所以當我們讀經文的時候.
其實很嚴肅.
第一個段落.
我們講了兩堂.

$^{281}$就是原來猶太人心底裡.
產生了一個共同的疑問.
就是神好像已經忽略了他們的民族.
不再愛他們.
所以一開始.
先知第一個問題就是說.
猶太人覺得自己被神愛.
先知就替神去講說話.
在第二節就已經講了.
我曾經愛你們.
他們立即很有共鳴感.
神愛我們是古代舊時.
不是現在.
所以先知就很清楚地代表著神去講.
解釋神愛你的發言.
是在哪裡找到.
然後你也要這樣代入.
我們講過了.
就是雅各,爾素.
就是以撒的兩個兒子.
他們兩個兒子的命運.
他們兩個兒子在神的約底下.
誰才真正得到繼承.
繼承亞伯拉罕的福氣.
一直延續的立約的人.
去承繼那個約.
你可以想像到.
這件事不是個人得失那麼簡單.
是涉及到神的應許.
一直透過他所揀選的人.
要將他的應許的福氣.
一直延續下去.
所以到底是爾素還是雅各呢?.
相當清楚是雅各.
因為是媽媽.
就是李伯嘉.
懷孕的時候.
神已經跟他說.
當然沒有告訴他有兩個兒子.
媽兒子.

$^{321}$更加不能照照超聲波.
照到是兩個兒子.
或者龍鳳胎.
那時候沒有那麼多科技.
什麼都要出來才知道.
但是已經很清楚告訴他.
誰最快出來.
就不是得到福氣的那個.
小的那個即是遲出.
因為只有一個門口.
遲出的那個.
就是真正得到福氣的那個.
有一年的時間.
給媽媽和以實去思想.
為什麼神這樣的說法呢?.
為什麼一定要小的得到福氣呢?.
然後就出來了.
輾轉我們見到他們.
其實一節聖經.
他們就馬上已經長大了.
好像《羽裳篇》那樣.
一晚就十八年後了.
出生了.
長大了.
馬上就到了一個場景.
就是耳塑.
其實他不是一天之間.
就喜歡吃紅豆湯.
而是他的為人.
說的是他的為人.
一直都是這樣的個性.
一直都是對.
似乎是暗示對神.
其實一點心都沒有.
但是雅各就很清楚.
祝福的來源是神.
是藉著爸爸的祝福.
我既然生在以實這個家庭.
我就已經肯定是蒙福的了.
但是我是小孩子.

$^{361}$我要是得到爸爸的祝福.
除非大兒子不要.
或者大兒子已經夭折.
死於非命.
他就可以繼承.
但是他知道.
哥哥比他更加生硬.
打獵過.
你作為小孩子.
如果你又懂得思考.
你就會很擔心.
你就會有很多計謀.
看看怎樣把他打中.
走雷.
諸如此類.
會有很多事情要思考.
因為聰明的人.
又沉靜.
心裡想很多計謀.
所以不是一個偶然的事件.
而是一個處心積慮的事件.
就要謀哥哥的福氣.
那個福分.
那個名分.
接著沒多久.
又到他們爸爸年紀老邁了.
眼睛都昏花了.
就跟大兒子說.
不是兩兄弟過來.
一直一心想著.
他的大兒子才是大兒子.
很合理的.
古代近東的文化都是這樣.
長子可以得到雙份.
再加上現在直接得到.
是亞伯拉罕和神的約的延續.
所以是家族的榮譽.
所以肯定是想給二嫂的.
也是陰差陽錯.
他們的對話.

$^{401}$被媽媽李伯家知道了.
李伯家就跟兒子.
小兒子就是雅各.
來合謀.
用了一些計謀.
做了所有的裝扮.
就欺騙了真真正正的.
那個屬天的祝福.
所以希伯來書和約書亞記都是.
都不斷的暗示.
以撒他是老眼昏花.
但是臨危就勒馬.
他不敢再超越神的命定.
不敢超越神的旨意.
神說了.
人原來用很多計謀想改變神的計劃.
包括以撒,包括李伯家,包括雅各.
唯一不想改變的就是二嫂.
從來都不在意自己的身份.
因為她覺得自己一定是這個身份.
整個家裡三個人想改變神的計劃.
但是最終以撒.
她真的知道她剛才祝福了她的小兒子.
摸著他按著他的頭來祝福.
她不敢再改變收回成命.
那個祝福就隨著雅各獲得.
她不敢再祝福給她的大兒子二嫂.
這個就是以撒.
你可以說臨死眼都昏花.
最好靈性的那個時間.
就是她不敢祝福給二嫂.
就是她最好靈性的時候.
就是她不敢得罪神.
不敢違反神的旨意.
大家我們盼望.
我們的人生不要這樣.
發現神的旨意.
我們很開心地遵行.
就是說以撒和利伯加.
應該在他們的孫子出生的時候.

$^{441}$就應該來擺滿月酒.
擺拜日宴的時候.
宣佈雅各.
James就是將來我的繼承人.
一早就立了遺囑.
大家不要爭了.
從來都沒有.
從來都不想聽這個吩咐.
利伯加也苦於丈夫是頭.
是族長.
都不敢來跟他挑戰.
大撒.
神說了.
雙子都未生下來.
善惡都未發出來.
將來大的要服侍小的.
小的要得到祝福.
要延續.
利伯加又不敢提.
只是不斷跟小的想想什麼手段.
你可以想像到.
他們全家人都想違背神的旨意.
不過奇妙的地方就是這個揀選論.
在這裡就見到神的揀選.
是堅定不移的.
賭羅說.
透過雅各和耳掃.
為什麼要再講一次這段經文呢?.
因為很重要就是.
神告訴猶太人.
我的愛是超越了人的好和不好.
神的揀選.
不會在乎猶太人.
他們在神的面前.
所做的事.
所以雅各和耳掃.
這些歷史人物.
已經很肯定地告訴.
現在修信的讀者知道.
神的愛不會隨隨便便廢棄.

$^{481}$神的揀選.
是不會後悔的.
所以他們生下來就是猶太人.
他們唯一應該要考慮的就是.
怎樣能夠讓他們自己這個民族的生命.
能夠反映到神的榮耀.
而不是思前想後.
我們這個民族神已經將我們廢棄了.
這是我們必須要在我們自己的心靈裡.
來學習的.
所以主耶穌的救恩.
在我們每一個弟兄姊妹的心靈裡.
是不會改變的.
所以我們一定要有這個很重要的學習.
就是神的愛是不會後悔的.
我們是永遠屬主.
沒有人可以從主耶穌的手裡.
將我們奪去的.
所以不是因為我們的際遇.
不是因為我們的疾病.
不是因為我們其他的事情.
就會使我們與神的愛隔絕.
這個愛是在主基督耶穌裡面的.
這是一個肯定的真理.
所以猶太人他們要問的問題.
就是他們為什麼不能夠經驗到神的愛呢?.
從第六節開始.
其實呼之欲出了.
就是那個敬拜出了問題.
猶太人很少會在舊約的時代.
今天沒空了.
很忙啊 過年了.
不上聖殿了.
兒子要懲婚事了.
都不去了.
因為他們是一個民族性的集體敬拜.
就不同我們現在新約的時代.
現代人我們都很個人的主義.
其實已經普遍了.
每個家庭有自己的繁忙.

$^{521}$有自己的生活節奏.
很少整個集體.
就很少的.
雖然還有的.
譬如我上個星期也提過.
我跟過林牧師我們去印尼傳福音.
那條村的村長信主.
整條村都要信耶穌.
沒得讓你選的.
如果那個村長是信伊斯蘭.
整條村都是跟伊斯蘭教的.
沒得選的.
所以這些集體性.
是我們香港人未必能夠明白的.
是現代的 譬如歐洲人.
我們未必會明白的.
我們不是那些圍村性.
不是那些集體性.
圍村都還有一點點這些概念.
就是阿公現在派豬肉.
大家一起去吃盤子.
那些集體性很強的.
但是現在因為已經工業化之後.
變成我們是家庭的單位.
很難就有整條村.
整個地方.
很難的.
特別香港這麼小更加.
如果地方很大的.
一個城市跟另一個城市的競爭.
譬如足球.
就很有榮耀.
譬如廣東散廣杯.
廣東贏了.
就很有榮譽感.
我們就要回頭看這裡的時候.
你要從這個集體性去看這件事.
他們不覺得神的愛.
其實慢慢會令到心靈有一種糾結.
好像蒙上了一種灰暗的薄紗.

$^{561}$就會變成很多時候.
沒有了那種力量.
在神的愛裡面.
去支取到一種.
令到心靈有一份平安.
而帶來的那種對前途的盼望.
第六節就告訴我們.
其實第五節就是那個總結.
神說了.
祂的名是要被高舉.
被尊為大.
要超越猶太人.
所以我們上一個星期說過第五節.
就應驗在保羅書信的羅馬書.
十一章.
保羅就是這樣說的.
很快揭一揭.
接著我們就進入第六節的經文.
因為其實真的回到.
最重要的就是我們敬拜神的態度.
我們看看羅馬書的十一章.
你可以肯定地相信.
保羅很明白.
救恩帶來的就是全個宇宙.
對做物主的偉大的敬拜.
十一章的三十三節.
其實如果再讀多一點.
你就會看到保羅是將猶太人的信主.
和神的榮耀拉上關係.
豈不是馬拉經書一至五節的邏輯呢.
你看看第25節.
羅馬書的十一章25節.
穆斯林和這裡都有關連的.
25節.
弟兄們,我不願意你們不知道者傲備.
恐怕你們自以為聰明.
寫的收信的讀者就是羅馬教會.
已經信主了.
這本書是告訴基督徒.
救恩的佈局.

$^{601}$神為什麼會起義有救恩.
救恩涵蓋哪些人.
什麼人可以得到神的救恩.
怎樣得到.
得到之後有什麼改變.
很有系統地將整個救恩講解清楚明白.
所以來到這裡是一個總結句.
你們不要以為你們現在蒙恩.
即外邦的羅馬信徒.
在這裡自高自大.
在這裡以為猶太人是被神撇棄.
他們是永不翻身.
其實不是.
這是一個傲備.
留意了.
就是以色列人有幾分是硬心.
等到外邦人的數目天滿.
於是以色列全家都要得救.
如經上所記.
必有一位救主從石安出來.
要消除那國家的一切罪惡.
這是《異象》書的59章.
有可能按字面解釋.
主耶穌再來就在石安.
腳踏橄欖山.
拯救他的子民.
雖然《異象》書已經早期應驗.
現在這裡再講.
第27節又說.
我除去他們罪的時候.
這就是我與他們所立的約.
再講神童列祖.
亞伯拉罕,以撒,雅各,約瑟,摩西.
一直所講的約.
就是約真真正正最終應驗.
即是暫時有一個.
好像是寬限期.
猶太人暫時以為靠自己的行為.
去換取神的恩典.
他們拒絕主耶穌.

$^{641}$因為人類唯有靠主耶穌.
憑信心才可以得到見神和神的救恩.
如果他們在耶穌以外.
立自己的救主.
那不是救恩.
所以我們今天新約我們信徒.
基督徒無論你是什麼民族.
只需要你歸向耶穌.
你就立即得到這個救恩.
猶太人一直拒絕這位彌賽亞.
我們一信就立即得到與神的和好.
原因就是這樣.
所以現在他們暫時因為他們的罪.
那個約未能夠應驗.
但是外邦人千萬不要有一份驕傲.
或者自大.
因為只有我們外邦人才可以得救.
因為第28節.
就著福音說.
就著福音來說.
就是說神在舊約.
透過憲制.
透過一些禮儀.
猶太人可以去到神的面前.
但是要有信心與神拉上這個關係.
才是真正的猶太人.
那個舊約過去了.
真正的救贖就是耶穌來了這個世界.
他宣稱自己就是神的兒子.
他就是那個贖罪制.
他就是那位大祭司.
他就是那位神的兒子.
開一條新的路.
使得透過他的人.
就可以去到天堂.
去到神的面前.
所以就著福音來說.
他們為你們的緣故是仇敵.
他們就是猶太人.
猶太人因為一直拒絕真正的尼西亞救世主.

$^{681}$所以他們就成為了福音的仇敵.
不是我們的仇敵.
是福音的仇敵.
記住.
就著簡選來說.
簡選.
他們又是說猶太人.
為列祖的緣故是蒙愛的.
剛才馬拉基書的邏輯就在這裡出現.
猶太人有一些生活的不如意.
或者慢慢慢慢.
去到一個習慣性的恩秦之後.
就感受到神好像和他們有疏遠.
有一段距離.
但是現在保羅再一次告訴我們.
就著簡選的角度來說.
神從來沒有廢棄和列祖所立的約.
是不斷地來確定約.
繼續堅定神不改變的心意.
就是愛這個民族猶太人.
再加上主耶穌來到這個世界.
就更加凸顯和真正的展現.
原來神不只是愛猶太人.
是愛世人.
愛每一個民族.
所以為了列祖的緣故.
猶太人是蒙愛的.
就算現在是犯罪.
被神好像放在一旁.
其實真正是蒙愛的.
所以要燃點猶太人有這份盼望.
他們將自己放在和耶穌十價救恩為敵.
明明只有一個救恩.
但他們卻不要.
就是說耶穌不是他們的彌賽亞.
他們不會相信這位耶穌.
所以這是一個很大的驕傲.
第29節.
因為神的恩賜和選照是沒有後悔.
原來選擇是夾雜著恩典的.

$^{721}$是恩賜.
就是說不是因為誰好就選擇.
原來這個選擇是無條件的.
所以我們不能夠明白.
明明選擇就是有條件.
我提過我小學的時候很開心.
六年級.
其實我那時候還很矮.
現在也很矮了.
我居然間入選小學.
在巴富街小學入選校隊的田徑.
開心到不得了.
因為一來你不用上課去練跑.
二來就可以參賽.
和其他小學的聯校去跑.
誰知道原來最後四成一百.
我是第五個.
所以從來沒有跑過.
但是你不要管.
選了就是有份.
為什麼會選呢.
田徑當然跑得快.
會計時.
其實大家都可以聯想到.
為什麼你去街市買蘋果.
你會選那些好的呢.
為什麼你要買一些東西.
都是你要千挑萬選呢.
人家是准你選的.
你一定會選一些好的.
你又不會選一些不好的.
所以在我們去想神的揀選.
當然是那些齊齊齊.
事事正正.
家族的榮譽.
幾代都是信主.
所以就選他.
其實不是的.
這裡告訴我們.
神的恩賜和揀選是堅定不移的.

$^{761}$是好像孖公仔一樣.
原來我們以為.
得天獨厚被神揀選.
當然要將來去到天堂才知道.
原來其實我們每一個人被揀選.
其實你第一個印象.
其實任何一個頂子妹都是.
是恩典行先的.
你回頭回望.
哇! 真是神在我的生命裡有揀選.
沒有人會說.
哎呀! 神你選我吧.
你沒有被揀選.
選我?.
那我們就搞錯了.
任何一個弟兄姊妹信主.
是恩典.
所以這個很奇妙.
所以保羅的神學其實很精妙.
神的恩賜和揀選.
是堅定不移的.
我們永遠不知道我們自己是不是被揀選.
這是清教徒很重要的信念.
這是改革中一個很重要的信念.
怎樣知道自己一直蒙揀選呢?.
說過了.
就是一直在主耶穌的救恩裡.
你就是蒙揀選了.
要保守自己在主耶穌的救恩裡.
所以一次相信.
一直相信.
就得救了.
記住是一次相信.
要一生的相信.
你才得救的.
這樣你就知道自己在揀選的行列裡.
所以不是我們做好蒙揀選.
而是我們在永恆裡.
原來神揀選了我們.
所以我們就會得到這個救恩.

$^{801}$我們在恩典,恩賜裡.
我們不斷發現原來自己是被揀選.
很奧妙的.
所以其實這個就是謙卑.
這個就是無條件的救贖.
所以保羅告訴我們.
神的恩賜和選照是沒有後悔的.
所以我們如果能夠對於自己.
還是那麼肯定自己是被揀選.
其實那個人的內心.
對自己的人生的前途有盼望.
對於永恆.
你是充滿了那種的福樂.
其實絕對不是想說什麼.
我在洪水橋幾個星期.
那些耕田的阿姐和一位教友.
那位教友很多年了.
因為十幾年在那裡直堂.
在處理感受的.
每個星期開車下去.
和那些耕田很熟.
我剛剛去幾個星期.
那個教友.
我不說是誰了.
那個教友.
那個耕田的阿姐就說.
你們來了新牧師嗎?.
這麼好笑容的?.
這麼開心的?.
諸如此類的.
和那個教友說.
那個教友和我說.
我的心想就是.
為什麼你要對著別人不開心呢?.
你應該是開心的.
我接著就會和他說.
有時候有些東西就會給他.
剛剛最近我們派了年糕給管理處.
都是一些教友送的.
他們很開心.

$^{841}$一來自己身份低微.
坐在耕田裡坐了12個小時.
你想想坐了12個小時.
當然有時候他們要巡視.
但是12個小時工作.
其實很悶的.
有些街坊.
有些教會.
關心他們.
走過說一聲問安.
這些其實是很少的事情.
但是對被問安的人.
就很感動.
這不就是因為我們被揀選.
因為我們看到這是神恩典的揀選.
所以是恩典和選照.
是神的堅定得宜.
就是祂的願望.
你想想我們如果有一個願望.
很強烈的要做一件事.
都常常能夠成就神的願望.
就是祂要在萬世以前.
揀選一些祂所揀選的人.
然後你和我出現在這個時空.
我們回望信主到今天.
肯定自己是被揀選.
你想想這是不是很大的福氣.
所以保羅在這裡告訴我們.
他講完之後.
其實我還沒講到最重要的那一步.
就是主的名要被尊為大.
在列國當中.
你留意一下保羅接著說什麼.
他的恩賜和選照是沒有後悔.
第30節.
你們就是白幫人.
就是集體的一個稱呼.
所有不是猶太人的.
就是沒有救恩的人.
你們從前不信服神.

$^{881}$有些是大偶像.
有些是信哲學.
保羅在《始祖行傳》中.
有些是未識之神.
在阿勒巴古中.
很多人拜不同的東西.
或者希臘人討論哲學.
不一而足.
這個就是外幫人.
你們從前不信服神.
如今因他們的不信服.
他們就是猶太人.
拒絕彌賽亞救恩.
他們不信服.
所以叫他們恩著.
思及你們的年殉.
現在也就蒙年殉了.
你就可以想像到.
賜下給我們的那個年敏.
可以臨到每一個人身上.
只需要每個人都能夠認定.
主耶穌在我們生命裡的地位改變.
接著就告訴我們第32節.
因為神將眾人都圈在不信服之中.
特意要連戌眾人.
你可以說原來每一個人.
一來到這個世界.
是一個墮落的世界.
每一個人出世.
到離開世界.
我們都是在一個不信服神的氛圍裡.
不同程度.
不同的抗拒.
但是我們在這個不信服裡.
在神的揀選.
我們就離開那個不信服的行列.
透過救恩.
是每一個人都可以信的.
不會設定了.
因為你是哪個民族.

$^{921}$你就不能信了.
就不會的.
所以這個就造成了.
我們在.
其實全人類都是被圈在不信服裡面.
沒有一個民族可以自己.
撇開了救恩.
聖經.
他可以認識真神的.
沒有的.
所以當我們今天能夠透過.
有一本聖經.
有人講解救恩.
我們信主.
我們在教會生活.
我們就應該明白了.
這個揀選的道理.
原來是從無條件的救恩開始的.
所以我們回到去.
我們知道原來我真的被揀選.
但是我們一生知道的.
其實就是恩典的美麗.
是主耶穌的救恩.
將我們拯救.
是一個無價的主.
為我們釘十字架.
所以這個就叫做.
如果用另一個角度.
這個就是揀選了.
如果真的要講條件.
這個就是條件了.
就是你接受一份無價的禮物.
這個就是揀選了.
當然一般叫做信耶穌.
但是其實回望.
這個就是神的揀選了.
所以你看看是多麼的公平.
因為全世界的人.
包括有信仰的猶太人.
對真神認識的猶太人.

$^{961}$但是拒絕耶穌.
我們是不認識神.
又拜錯了其他的偶像.
或者信自己.
但是我們都認識了主耶穌之後.
我們就生命改變.
我們就得到救恩.
所以神就特意要連述眾人.
33節到36節就是保羅最後的一個讚頌.
他明白救恩了.
所有人類原來都是罪人.
所有的人類都是在不信服的行列.
所有的人類原來都是生而核眼.
不知道有永恆.
不知道有美好的永生.
不知道有一位造物主是這麼偉大.
更加不會知道造物主甘心樂意將他的愛子犧牲.
所以第33節.
心齋,即是很偉大.
神豐富的智慧和知識.
所以救恩不是建立在感覺上.
是建立在真理,在教義,在系統神學上.
救恩不只是我們感動有一個好的故事.
信耶穌.
你那一刻感動.
如果你不再透過聖經的真理去建構你整個神學.
你是很容易失去救恩.
有些人洗禮不再回教會.
各樣原因可能情有可原.
可能不能夠接受.
無論怎樣都好.
牧者不是一生要跟著我們.
不是這個意思.
而是你盡快無論信主洗禮.
你自己要建立跟神之間一個很具體的.
很有基礎的信仰.
信仰是建立在神豐富的智慧和知識裡面.
救恩是有豐富的智慧和知識.
所以我們要記得盡快對真理有一個明白.
所以上初信班的弟兄姊妹.

$^{1001}$浸泥班的入門是很重要的.
整個救恩的邏輯原委要明白.
這樣我們就能夠對神有認識.
就不會變成有些道聽途說.
憑感覺.
你的信仰就不能夠堅固.
接著他告訴我們.
講到神的判斷是何其的難測.
判斷是他的judgment.
神做甚麼我們是無法知道的.
表面我們覺得很有底觸的.
其實我們不需要這樣看.
也不應該用這麼狹隘的.
我們自己有限的知識去看神的策劃.
祂的奇妙智慧和判斷.
祂的終積何其難尋.
祂的終積是甚麼呢?.
一定是救恩的佈置.
祂如何處理世界任何一個人.
這是祂的終積.
我們要不斷講的是.
你不會尋找到的.
你不會測度到的.
誰知道主的心.
誰做過祂的謀事.
即是說我們千萬不要幫神想一些.
我們以為必定對的事.
神為何會興起一些人.
他能夠生命有這麼大的改變.
譬如保羅自己就是過來人.
十多年前他還在大馬士革.
去抓基督徒坐牢.
有些人去到宗教法庭去定罪.
他在那裡做附和簽名.
來贊成支持.
有些人要下監.
甚至有些人可能要死.
保羅就是這樣的人.
現在他告訴我們.
他發現不是這麼簡單.

$^{1041}$原來神的計謀.
神的旨意.
至少在我保羅身上.
就體現到我測度不到.
我沒想過十五年之後.
我成為到處去到歐洲的傳教士.
他自己真的測度不到.
他只是想著自己一心在猶太教裡面.
做到最好的法利賽人.
就是嚴格遵守律法.
不斷地研讀.
不斷地將經卷背誦.
他是很熟悉的.
但原來當耶穌來了這個世界之後.
所有的舊約.
他所認識的舊約.
是完全一文不值.
不是聖經一文不值.
他解錯了聖經.
一個拉比一生自小的成長.
就是要去學習聖經的.
但當耶穌還沒出來這個世界.
他就非常之卓越的一個法利賽人.
但當耶穌來了這個世界.
當耶穌告訴他.
我就是舊約所預言的彌賽亞.
我就是你所迫伯的耶穌.
他整個神學立刻倒塌了.
他一直所讀的所背的那些舊約.
他一直都是對著神的.
對著耶和華神這樣來學習的.
但是當耶穌出現.
他整個神學完全崩潰.
你可以想像得到那種失落感.
那種挫敗.
他現在當然是正面.
是救恩使他整個人改變.
他不是屬靈破產.
是神重新將他打通他的任督二脈.
使他能夠完全通了.

$^{1081}$所以一個仇恨外邦人的保羅.
成為一生信主之後.
愛每一個外邦人.
他帶著的那些門徒.
除了提摩太.
大部分是外邦人.
歐洲人.
一個拉比,法利賽人.
不會接納外邦人的.
他可以帶著這些年青人.
這些門徒.
做他的手下.
信任他們.
他們去不同的地方教會去服侍.
你可以想像一下保羅.
哪裡會想到自己.
他收的學生全部是猶太人.
大家一起在猶太教裡面.
這樣鑽研那些經典.
他怎會到處去到外邦.
和不同的歐洲人.
來到去清卿道弟.
一起去傳福音.
告訴別人.
我們祖宗的神所差遣的彌賽亞.
就是釘十字架的那位.
拿撒勒人耶穌.
這個是保羅他一生的震撼.
所以他說.
誰做過他的謀事.
這些全部保羅自己有答案的.
他現在這樣說的時候.
就是告訴我們.
這是多麼偉大的造物主.
所以神的名是要在外邦中被尊為大.
就是馬拉基書一章五節.
這裡告訴我們.
使他後來常還.
接著第五十六節.
因為萬有都是本於他.

$^{1121}$依靠他 歸於他.
願榮耀歸給他.
直到永遠 阿們.
這個的頌讚.
保羅在這裡 寫到這裡.
一直在說救恩.
一直在說神堅定不移的計劃.
原來猶太人 外邦人.
本來就是不信服的.
但是主耶穌的救恩.
神的愛子.
可以令到一部分被揀選的人.
生命有改變.
至於哪些是被揀選呢.
我們不知道.
但是哪些人回應救恩呢.
我們肯定知道.
所以我們要宣講主耶穌.
是人類的救主.
我們就在幫助那個人.
知道自己被揀選了.
所以其實是一件很好玩的事.
當你這樣看的時候.
我們是將神的永恆救贖計劃.
每一次宣講.
就是在那個人身上.
他聽了可能十次後才信的.
第十次他就信了.
然後他一生的改變.
然後他再回望.
原來我在錦繡堂信主.
然後我不斷去發現.
原來我是被神揀選的.
我們就參與在神的救贖計劃裡面.
讓一些地上的人被神揀選.
一定是少之又少的.
人類能夠信主是很少的.
不是救恩不夠好.
不是有條件.
說明是白白的主耶穌的寶血.

$^{1161}$十價的救恩.
每個人都可以信就行了.
所以這個揀選.
不是神限制了.
是人自己拒絕.
所以你可以看到.
三十六字這個頌歌.
對眾物主偉大的頌揚.
萬有都是本於祂的.
救恩是祂想出來.
設立將人拯救.
犧牲的是祂的愛子.
不是天使.
是祂自己的兒子.
萬有都是本於祂.
我們是做物主所創造的.
祂是那個本.
我們根本是神.
所以信耶穌第一件事.
一生最重要的就是.
這是威斯敏斯特的教理問答.
問那些小朋友.
中國改革之後.
信耶穌第一件事要做什麼?.
就是榮耀神.
我們的人生.
因為萬有是本於祂.
是倚靠祂.
倚靠祂.
我們學習人生對神有一個信靠.
保羅所說的倚靠.
當然不是我們什麼都不理會.
坐著就會突然變.
這個倚靠是說我們都會有我們的勤奮.
讀經,靈修,祈禱.
在我們的生命裡盡量流露基督徒的品格.
所以保羅沒有衝突.
十二章之後就開始說品格的修為.
我們應該在我們的行動上.
怎樣能夠像基督徒,像基督.

$^{1201}$所以這個倚靠祂.
是我們整個生命的存在.
所以我們真的越來越應該.
到了一個年紀.
差不多是人生的退場.
我們應該開始有退場的心態.
不是悲哀的,不是負面的.
是正面的.
應該預備一個退場.
就是那個人應該越來越有一種感恩.
對你所經歷過的人生.
我們去回望這個月曆.
是帶著感恩的.
因為我們是倚靠這位神才能夠存活的.
所以記住.
本於祂,倚靠祂.
真是有秩序的.
我們認定神是我們的根本.
我們的人生是倚靠祂的.
很多人今天不信耶穌.
如果他一生都很平順.
又很富貴.
其實他可以自高自大.
以為不需要靠神.
是很有限的.
所以有那些物質的豐富.
是讓我們更加好的去事奉的意思.
我們應該更加將我們的人生.
聚焦在怎樣能夠榮耀神.
所以將來要歸於祂.
要交賬.
你看本於,倚靠,歸於祂.
你說你們真是信耶穌當然了.
你們基督徒.
不是的.
任何一個人離開這個世界.
都要回到神那裡.
所以是歸於祂的.
願榮耀歸給祂.
直到永遠.

$^{1241}$對神的重陽.
在救恩底下.
保羅在猶太人.
神安排的救贖計劃次序.
外邦人現在先蒙恩.
保羅是在說整個救恩的過程裡.
他來去重陽主.
所以差不多是代表著.
一個絕對的真理.
保羅看到救恩改變人類命運.
所以他能夠這樣來去.
都很崇高的了.
對於地球的人.
能夠這樣看到這個世情.
不是那麼多人能夠這樣看到.
人類離開了主耶穌.
就沒有救恩了.
猶太人他自己就是痛心的地方.
就是猶太人現在沒有救恩.
被撇棄.
所以他就很想.
他的骨肉之親.
能夠信主.
甚至他被咒詛.
他都願意.
所以來到十一章結尾這裡.
其實就是剛才我要和大家看的.
一章的第五節.
就是馬拉基書所說的.
描述的將來神.
要和外邦人的救贖.
要猶太人自己親口承認.
耶和華在以色列的境界之外.
被尊為大.
由保羅的口.
向造物主去重陽.
寫羅馬書的時候.
大部分人都說是在哥林多的時候.
他當時都處理著很多的困難.
但是那些困難.

$^{1281}$沒有影響他對救恩帶來的.
造物主的榮耀的重陽.
因為他真的知道.
本於他.
倚靠他.
歸於他.
所以這幾樣東西.
我們能夠一生這樣發現.
就是一個很寶貴的人生.
你的人生就很豐富.
你怎樣來的?.
我是本於神的.
你的人生靠誰?.
我倚靠祂的.
將來我要歸於祂的.
你整個人生就有一個敬畏.
對造物主的敬畏.
對造物主的敬愛.
對造物主在你生命裡面.
給你的恩典.
給你的恩賜.
你很想和人分享.
這就是我們人生的一個.
很大的屬靈的豐富.
所以來到這裡的時候.
第五節.
猶太人不要自我貶低.
表面上神好像沒有介入.
但其實從第六節開始.
神原來留意每一個人的敬拜生活.
那個心態是怎樣敬拜.
然後告訴我們.
神可以得著這麼大的榮耀.
要在外邦人當中.
祂的名字被尊為大.
所以我們錦宣家的弟兄姊妹.
現在當然包括了.
嘉欣堂.
孫德.
洪欣.

$^{1321}$當然再擴大就是宣道會.
再擴大的就有其他的宗派.
播道會.
浸信會.
這些比較大的宗派.
我們應該要給其他的.
未信的香港人看到.
神的名字被尊為大.
所以其實真的端扶我們.
信了主之後.
不信的人沒有要求.
信了主之後.
我們是否本於祂.
倚靠祂.
歸於祂.
我們的人生是這樣的走向.
信了主了.
在教會了.
但從來都沒有想過要歸於神.
從來都沒有想過要靠神.
我又不是想諷刺甚麼.
回祈禱會的人真是少之又少.
哀呀哀呀都沒有人敢回.
好像不想回.
不需要回.
當然每個人真的有自己的時間.
作息.
工作.
忙碌.
這個是完全明白的.
在一個工業化的社會.
每個人都有自己的忙碌.
但仍然在說.
我們這個本於祂.
倚靠祂.
歸於祂的生命成長.
還有我們有沒有用我們的人生.
奉耶穌的名去.
為一些國際大事祈禱.
捆綁.

$^{1361}$不是捆綁拜登.
不是捆綁特朗普.
不是那個人.
是捆綁那個.
敵黨神的.
那些背後的事情.
誰令到戰爭的.
我們能不能一起祈禱.
希望受害的人.
能夠得到神的恩典呢.
這些是我們在自由世界的人.
好像看完電視就跟我們沒有關係.
喂快點完.
劇集來了.
我們就會是這樣.
你試想想.
只是叫你坐飛機去巴勒斯坦.
去那個災區.
你看那些平民百姓.
突然間有炸彈炸下來.
你就可以想像到那些痛苦.
那些困難.
其實應該感動大家.
真的回來祈禱會.
崇拜.
崇拜就很少人說.
我不回去崇拜的.
很少人這樣說的.
祈禱會就很多人都不是很想回去.
從來都不是在他的意識裡面.
要回去.
有時甚至連真的很多時間都不回去.
這個我自己覺得.
會失去了很多福氣.
這個當然是給錢免禮.
進入第六節.
真是問題所在.
為什麼他們不能夠感受到神的名.
在外邦當中被尊為大.
和他們不感受到神的愛呢.

$^{1401}$其實諷刺就是.
全部的猶太人.
是有敬拜生活的.
是有向神奉獻.
獻祭的生活的.
「藐視我名的祭司」.
「萬君子耶和華對你們說」.
「兒子尊敬父親」.
「獨人敬畏主人」.
「我既為父親」.
「尊敬我的在那裡呢?」.
「我既為主人」.
「敬畏我的在那裡呢?」.
「你們卻說」.
「我們在何事上藐視你的名呢?」.
答案就在第七節.
他們就是這樣對神有一種藐視.
是哪些人啊?.
是事奉神的祭司.
「你們將污穢的食物」.
「獻在我的壇上」.
而且這樣說.
我們在哪些事情上污穢你.
根本好像已經慣性了.
做這些事情.
都完全不覺得神喜悅.
「因你們說」.
「耶和華的卓子是可藐視的」.
「你們將黑眼獻為祭物」.
「這不為惡嗎?」.
「將瘸腿有病的獻上」.
「這不為惡嗎?」.
「你獻給你的省長」.
即當時民間的地方官.
治理那些行省的官員.
因為在耶路撒冷.
當時波斯帝國統治的.
耶路撒冷只是一個行省.
不關事.
說得太多了.

$^{1441}$(你細聲一點).
我細聲一點說.
你獻給你的省長.
他怎會喜悅你呢?.
怎會來看你的情面呢?.
這是萬軍之耶和華說的.
接著第九節.
有什麼解救的方法呢?.
「現在我勸你們懇求神」.
「祂可以施恩給我們」.
這個忘獻的事.
原文就沒有忘獻這兩個字.
這些敬拜的場景.
或者那些敬拜需要有基本的手續.
都是由你們去經手的.
就是祭司.
他怎能夠看你們的情面呢?.
這是萬軍之耶和華說的.
唯一的辦法就是.
「甚願你們中間有人」.
有一個人有一種眼光.
「今日不要獻祭了」.
「因為我們要獻祭你的神」.
「祂不喜歡」.
「甚願你們中間有一個人」.
「一個人而已」.
「關上電門」.
「免得你們徒然在我的壇上燒火」.
燒也是徒然.
會不會有個祖堂的師士.
跟某人說.
「今日我們不要崇拜了」.
「甚麼事?為甚麼不可以崇拜?」.
「今日神不喜歡我們崇拜」.
「關上門」.
「Yoyo要開音響」.
「開冷氣,開燈」.
他到處走阻止他.
不讓他這樣做.
我們都一定以為.

$^{1481}$那個人是有精神問題.
我們就叫勵神牧師來跟他聊天.
說禁止人們崇拜.
還要關電門.
你可以想像到.
當然剛才這麼說是開玩笑.
但是當時的解決辦法.
如果不改變.
就寧願關了電門.
下面就說.
「我不喜悅你們」.
「也不從你們手中收納那些貢物」.
即是說就算你燒.
那些變成焚化爐.
不是輕輕的按鍵.
去到神的面前.
變成一般的燒垃圾.
燒了那些祭物.
即是一個好像.
好像等於人家做動物的火葬.
燒了那個動物.
可以這麼極端.
到底是馬拉基才知道.
他這個人太激進,太極端.
還是真是神期望.
一個情願關門的聖殿.
也不情願你們繼續在那裡.
不斷獻祭.
其實是香的.
其實燒任何東西都香的.
再加上蜜糖一定香的.
你燒烤過你都知道.
但是神聞到的不是香.
第十一節.
萬軍之耶和華說.
從日出之地到日落之處.
再第二次出現了.
我的名在外邦中必尊為大.
在各處人必奉我的名燒香.
獻潔淨的供物.

$^{1521}$即是他們的獻祭是不潔淨.
不蒙神月立.
因為我的名在外邦中必尊為大.
第三次說了.
第十二節.
你們卻竊讀我的名.
說耶和華的卓子是污穢的.
其上的食物是可藐視的.
這個藐視和污穢.
相信不是表面.
是他們已經降低了標準.
而符合了基本的民情.
所以才導致.
神說我要的是這個水平的獻祭.
你們給我的就是一般性的.
即是可有可無的.
神就不要這種敬拜的獻祭.
諸如此類.
其實下面都有關連的.
你們又說.
這些事是何等的繁瑣.
並且嗤之以鼻.
即是那個獻祭.
是一個對事奉者的心靈.
是一個很繁重的責任.
是一個去到例行公事.
快些完結後.
我就如釋重負的輕省了.
不是看到人來敬拜神.
感覺很大的滿足.
很大的快樂.
以及將百姓獻的祭.
燒到壇前.
讓神得到滿足.
為祂最高的崇高的即是的服侍.
看為是嗤之以鼻.
一個祭司.
一個幫百姓獻祭.
在最神聖的致勝所裡.
看這一切的活動.

$^{1561}$是嗤之以鼻.
類似這樣.
中文的嗤之以鼻應該是用什麼表情去做出來呢.
大家回去想想吧.
那個祭司嗤之以鼻.
覺得是一個繁瑣.
很麻煩.
即是越多人來獻祭就越麻煩.
最好你不要來.
但不來不行.
因為他們的生活.
和那個祭物掛鉤了.
所以大家如果不懂利未記.
就不知道那麼多的精彩地方.
因為原來越多人來獻祭.
那些祭司有下欄.
那些祭肉.
祭司可以吃的.
所以你們要明白.
就是原本那件事.
是很榮耀的.
你代表神去獻祭.
其實神要的祭物.
是很少的.
如果是獻一個贖罪祭.
或者平安祭.
祂要的是那個內臟的少少精華.
象徵性地燒在祭壇.
其他整個的祭身.
是歸祭司的.
平安祭的祭身.
就是獻祭者和獻祭的人.
就是祭司和百姓.
都可以吃的.
平安祭甚至祭司的家屬.
都可以吃的.
整個平安祭的祭身.
如果是一隻牛或者一隻羊.
或者有些窮的就獻班溝.
獻一些甲子.

$^{1601}$甚至獻一些麵粉.
祭司只是拿一撮.
這樣來灑在壇上.
其他整盤的麵粉是歸祭司的.
所以你們想想.
原來獻的時候.
你代表著神.
閱立這些祭物.
你已經有一份榮耀的了.
然後獻完祭之後.
你想想有多少人.
陸續將那些祭身獻祭.
不是完全整隻燒了的.
是給祭司的.
你可以想像到.
這個是神和人共享的福氣.
除非是全民族的贖罪祭.
整隻的祭身就一定要整隻燒了.
但也有手續的.
還有繁祭就一定是了.
其餘的大部分是要歸祭司的.
在我還沒看完這個之前.
回頭看看什麼叫潔淨.
22章的利未記.
簡單分有五個祭.
繁祭.
整隻祭身獻給神.
獻祭首先要整隻洗乾淨.
切好 放好.
然後將內臟拿出來.
然後就完成了.
整隻燒了.
無一可以保留的.
除了獻祭有一樣東西可以保留.
而且是歸祭司的.
就是繁祭身的皮.
如果你是獻祭.
可以獻老虎.
獻牛皮.
可以做衣服.

$^{1641}$就歸祭司了.
羊皮也是.
所以繁祭就是這樣的獻法.
整隻的.
所有要燒了.
不過燒完還燒完.
都要洗乾淨.
還有切開.
將內臟拿出來.
這樣來獻的.
第二個祭就是贖罪祭.
或者再加一個贖獻祭.
大的罪就是贖罪祭.
小的罪就是贖獻祭.
接著還有就是平安祭.
加起來.
一二三四.
還有一個叫素祭.
素祭就是非殺生的素祭.
素嘛 吃素嘛.
素祭.
所以譬如麵粉.
或者其他.
或者那些修成.
那些菊.
諸如此類.
總之是有殺生的.
非殺生的.
繁祭,素祭,平安祭,贖罪祭,贖獻祭.
五個祭.
當然在這五個祭之外.
有就職禮的獻祭.
有其他的祭.
譬如拿西爾人.
他們還俗又要獻祭.
有很多不同的東西.
第二章這裡就相當重要.
暗示一個大方向的邏輯.
你還沒記第22章.
耶和華對摩西說.

$^{1681}$你告訴亞倫子孫作祭司的說.
祭司不可為民中的死人.
瞻掩自己.
除非為他自己骨肉之親的父母.
兒女弟兄.
和未曾出嫁,作處女的姐妹.
才可以瞻污自己.
什麼意思呢?.
即是你去送殯.
誕喪時.
祭司就瞻污了自己.
除非很親的.
什麼意思呢?.
即是你有職在身.
你就不能參與這些事.
就要交由其他人去參與.
除非很至親.
你就可以這樣去參與.
第四節.
祭司既在民中為首.
就不可從俗瞻掩自己.
似乎有些繁文縟節.
可能哭喊.
可能要守一些日子.
祭司就不可以了.
不可以洗到頭光禿.
不可以剃除鬍鬚的周圍.
也不可以用刀去劃身.
要歸神為聖.
不可以褻瀆神的名.
因為耶和華的火祭.
就是神的食物.
是他們獻的.
所以他們都要成為聖.
就不可以娶妓女.
或者被污的女人為妻.
也不可以娶被休的婦人為妻.
因為祭司是歸神為聖.
所以你要洗到祂成聖.
祂就是摩西.

$^{1721}$洗到祭司成聖.
因為祂奉獻了你神的食物.
你要以祂為聖.
因為我洗到你們成聖的耶和華.
是聖潔的.
祭司的女兒.
如果行淫辱沒自己.
就辱沒了父親.
就要用火將她焚燒.
在弟兄當中.
作大祭司.
頭上倒了膏油.
又承接聖跡.
穿了聖衣.
就不可以逢頭散髮.
那就是哭喪.
也不可以撕裂衣服.
也不可以靠近死屍.
正正就是所說的.
不可以沾污自己的衣屍.
也不可以為父母沾染自己.
就是說.
當你在就職期間.
頭上已經有膏油.
有多大的事.
都是以神為先.
以聖跡為先.
有時候好像很不近人情.
但是邏輯的概念就是這樣.
就是說你不可以隨隨便便.
看你的職份是很輕.
每個星期都說自己不舒服.
星期六晚就說沒空.
其實就去喝酒.
諸如此類.
變成很多時候.
真的要當值的時間.
你就不做了.
這個責任感的問題.
其實就不對了.

$^{1761}$因為你已經是民宗的大祭司.
其實我們真的有要求的.
我們同工或者全職的.
全年只能夠放兩個星期的星期日.
每個星期日放假.
那你就是不用做了.
你的工作性質的工種.
是星期要做的.
很合理的.
其實沒有人會覺得不合理的.
如果你打算要放星期.
那不如不要做全度了.
除非將來不是星期.
是星期一崇拜就沒問題了.
其實這些是不講自明的了.
如果有些人要挑戰.
為什麼我放這麼少星期.
那是不行的.
因為你蒙召那一刻做全度.
就已經知道暫時地球上.
星期是星期.
就算有些教會是星期二晚.
星期六.
但是星期始終都是一個.
所有弟兄姊妹一起聚集的時間.
暫時吧.
可能將來改變.
但是這裡告訴我們.
重點就是.
作大祭司頭上倒了膏油.
又承接了聖積.
即是事奉的時候.
你剛好遇到有奔喪的事.
以及已經穿上了聖衣.
你就不可以逢頭散髮.
更加不可以撕裂衣服.
即是別人悲哀.
你就暫時不要悲哀了.
大祭在身.
其實這樣是很崇高的.

$^{1801}$這樣會令到一個事奉神的人.
感覺到自己身份的重要.
這些其實我們都要學習的.
所以星期見到人穿西裝.
一般都會傳遞給別人.
現在的社會越來越普通.
和簡單.
所以有時見到這些.
穿西裝的星期日的人.
一方面他未必是銷售員.
你就應該和他點頭.
打招呼.
默默為這些穿西裝的人祈禱.
他現在回教會了.
雖然好像開玩笑.
但其實不是的.
因為你看看.
舊約都這麼尊重這些大祭司.
他都要來到去.
對於自己戴著的積分.
他要看重.
其他的東西.
不是說要不近人情.
而是他的積分很重要.
所以他要暫時.
和其他人有點不同.
不可以為父母沾染自己.
不可以出聖所.
果然真是.
真是那個時候戴著積分.
也不可以褻瀆神的聖所.
因為神的高柔的冠冕在他的頭上.
不斷說榮耀.
是見得到舊約的時代.
戴著那個冠冕.
穿著那些禮服.
進到聖所.
似乎是在說.
致聖所的七月初十的贖罪日.
死了的父親都不可以參與.

$^{1841}$這真是神的一個要求.
我是耶和華.
他要娶處女為妻.
所以我不明白天主教.
為什麼說不可以結婚.
大祭司這麼神聖.
這麼崇高.
包括摩西自己也有.
寡婦或者被休的婦人.
或者被污為妓的女人.
都不可以娶了.
只可以娶本民宗的處女為妻.
不可在民宗辱沒她的兒女.
因為我是叫她成聖的耶和華.
這跟他的家庭生活有關連.
耶和華對摩西說.
你告訴亞倫說.
這裡說的是潔淨.
對於獻祭的要求.
你告訴亞倫.
世世代代的後裔.
凡是有殘疾.
都不可以近前來獻他.
獻他神的食物.
因為凡是有殘疾的.
無論是黑眼,瘸腿,塌鼻子.
肢體有餘,折腳折手.
駝背,矮錯.
特別矮的.
類似的.
眼睛有毛病.
長癬,長戒.
或者有腎病.
或者損壞腎子.
都不可以近前來.
他是祭司.
但這些似乎不是有歧視成分.
而是身份.
職份.
需要有這樣的想法.

$^{1881}$因為他將火祭獻給耶和華.
他如果有殘疾.
就不可以近前來獻神的食物.
神的食物.
無論是聖的.
至聖的.
他都可以吃.
因為他是祭司的家人.
祭司的家族.
雖然他有這些缺陷.
但他可以吃聖物.
所以剛才說的這些下欄.
獻祭者是可以吃那些祭物的.
所以這些概念.
到了耶穌基督保羅的年代.
保羅都不斷說.
牛在牆上踩谷.
都不可以.
綁著牠的口.
不讓牠吃東西.
一邊吃一邊耕田.
意味著.
傳道人是靠福音養生.
是應該有.
會眾對牧者的敬重.
至少現在香港.
還是弟兄姊妹的奉獻.
去支持教會的聖公.
不是由政府.
由國家出糧.
我們不是公務員.
至少我們知道.
舊約已經是這樣.
連那些有殘缺的家屬.
他是祭司.
但就不可以做祭司.
他是生下來有天生缺陷.
但是神道說明了.
那些致勝的食物.
聖物他都可以吃.

$^{1921}$有些是致勝所獻的.
有些是普通的民間獻的.
都可以吃.
但不可以進到萬子前.
所以他們都要有.
聖俗之分.
連你本身是祭司.
生下來是祭司.
你都不可以事奉.
這個殘缺.
似乎很不近人情.
但聽起來也很符合邏輯.
也不是對他們特別的歧視.
因為他有殘疾.
免得褻瀆了我的聖所.
我是叫他成聖的耶和華.
在神的眼中.
其實沒有人是殘缺的.
不過在人.
在百姓十二支派當中.
有沒有搞錯.
你們祭司找一個這麼殘缺的.
可能他每次見到都不開胃.
就批評論斷.
那個矮子走去獻祭.
說他鼻子又塌的.
有沒有見過.
我見過.
在元朗的時候.
那個男士做家長.
真的又矮.
鼻子又塌.
樣子真的很醜.
但我相信他很善良.
後來好像不見他了.
退休了.
真的很矮.
他的鼻子是塌的.
不是說鼻子沒有了.
是低了.

$^{1961}$但在聖經描述真的有.
你可以想像得到.
這些不應該叫歧視.
不應該叫歧視.
不過在這些聖職裡.
神就說不如.
有些人.
不要像被人標籤.
和被人歧視.
神就不會歧視人.
但人就會歧視.
人就會看 哇 搞錯了.
就變成有些人有殘廢.
例如.
他有瘸腿.
生下來.
有殘疾.
然後肢體有餘.
多隻手指.
你想想.
六指及摩.
被人取笑.
尤其是古代.
更加是.
所以其實舊約的利未記.
規範那個聖潔.
獻祭者要聖潔.
23,22章.
對不起 原來我看漏了.
我看了21章先.
其實我講的22章.
就是所有的聖物都要聖潔.
我們今天不夠時間了.
我們下星期再講.
我們禱告了.
親愛的天父 多謝你.
讓我們都能夠順利完成這過半小時.
我們真的很盼望敬拜的生活.
今晚尤其是強調的是.
大祭司的聖潔.

$^{2001}$帶動全民族.
對神的敬拜.
不要看為一個繁瑣.
和癡之以備的一個活動.
是看為每一次的聚集.
是歡慶.
是能夠去到神面前.
是能夠得到神的收納.
那些祭物 那些祭身.
以致神感到滿足.
而得到心靈的祝福.
盼望神都幫助我們.
盼望我們下星期.
未到過年之前.
我們都帶著一份虔誠的心靈.
去看我們的崇拜生活.
一起去檢視 一起去反思.
我們這樣感恩禱告.
奉耶穌基督的名.
阿們.
\newpage



\section{創世記 26:12-33}
\label{sec:vJ5BK_GiIWM}
\textbf{2024-02-04 早堂《以撒的神》創世記 26\_12-33 講員:黃頌言姑娘}
\newline
\newline
連結: \href{https://youtube.com/watch?v=vJ5BK_GiIWM}{\texttt{https://youtube.com/watch?v=vJ5BK\_GiIWM}} ~~~~ 語音日期: 2024-02-06
\newline
\newline
\hyperref[sec:Srg_1OhzOG8]{\small{< < < PREV SERMON < < <}}
~
\hyperref[sec:index_chronic]{\small{[返順時目]}}
~
\hyperref[sec:index_scriptual]{\small{[返順卷目]}}
~
\hyperref[sec:HWk9pevFs74]{\small{> > > NEXT SERMON > > >}}
\newline
\newline
創世記 26:12-33
\newline
\begin{longtable}{cl}
\hline
\hline
章節 & 經文 (和合本修訂版)\\
\hline
26:12 & \begin{tabularx}{0.7\textwidth}{X} 以撒在那地耕種，那一年有百倍的收成。耶和華賜福給他， \end{tabularx} \\ \\ \relax
26:13 & \begin{tabularx}{0.7\textwidth}{X} 他就發達，日漸昌盛，成了大富翁。 \end{tabularx} \\ \\ \relax
26:14 & \begin{tabularx}{0.7\textwidth}{X} 他有羊群牛群，又有許多僕人，非利士人就嫉妒他。 \end{tabularx} \\ \\ \relax
26:15 & \begin{tabularx}{0.7\textwidth}{X} 他父親亞伯拉罕在世的時候，他父親的僕人所挖的井，非利士人全都塞住，填滿了土。 \end{tabularx} \\ \\ \relax
26:16 & \begin{tabularx}{0.7\textwidth}{X} 亞比米勒對以撒說：「你離開我們去吧，因為你比我們強盛得多。」 \end{tabularx} \\ \\ \relax
26:17 & \begin{tabularx}{0.7\textwidth}{X} 以撒就離開那裡，在基拉耳谷支搭帳棚，住在那裡。 \end{tabularx} \\ \\ \relax
26:18 & \begin{tabularx}{0.7\textwidth}{X} 他父親亞伯拉罕在世的時候所挖的水井，在亞伯拉罕死後，都被非利士人塞住了，以撒就重新把井挖出來，仍照他父親所取的名為它們命名。 \end{tabularx} \\ \\ \relax
26:19 & \begin{tabularx}{0.7\textwidth}{X} 以撒的僕人在谷中挖井，就在那裡得了一口活水井。 \end{tabularx} \\ \\ \relax
26:20 & \begin{tabularx}{0.7\textwidth}{X} 基拉耳的牧人與以撒的牧人相爭，說：「這水是我們的。」以撒就給那井起名叫埃色，因為他們和他相爭。 \end{tabularx} \\ \\ \relax
26:21 & \begin{tabularx}{0.7\textwidth}{X} 以撒的僕人又挖了一口井，他們又為這井相爭，以撒就給這井起名叫西提拿。 \end{tabularx} \\ \\ \relax
26:22 & \begin{tabularx}{0.7\textwidth}{X} 以撒離開那裡，又挖了一口井，他們不再為這井相爭了，他就給那井起名叫利河伯。他說：「耶和華現在給我們寬闊之地，我們必在這地興旺。」 \end{tabularx} \\ \\ \relax
26:23 & \begin{tabularx}{0.7\textwidth}{X} 以撒從那裡上別是巴去。 \end{tabularx} \\ \\ \relax
26:24 & \begin{tabularx}{0.7\textwidth}{X} 當夜耶和華向他顯現，說：「我是你父親亞伯拉罕的神。不要懼怕，因為我與你同在，要賜福給你，也要為我僕人亞伯拉罕的緣故，使你的後裔增多。」 \end{tabularx} \\ \\ \relax
26:25 & \begin{tabularx}{0.7\textwidth}{X} 以撒就在那裡築了一座壇，求告耶和華的名，並且在那裡支搭帳棚；他的僕人就在那裡挖了一口井。 \end{tabularx} \\ \\ \relax
26:26 & \begin{tabularx}{0.7\textwidth}{X} 亞比米勒同他的顧問亞戶撒和他軍隊的元帥非各，從基拉耳來到以撒那裡。 \end{tabularx} \\ \\ \relax
26:27 & \begin{tabularx}{0.7\textwidth}{X} 以撒對他們說：「你們既然恨我，趕我離開你們，為甚麼又到我這裡來呢？」 \end{tabularx} \\ \\ \relax
26:28 & \begin{tabularx}{0.7\textwidth}{X} 他們說：「我們明明看見耶和華與你同在；因此就說，讓我們雙方彼此起誓，我們跟你立約， \end{tabularx} \\ \\ \relax
26:29 & \begin{tabularx}{0.7\textwidth}{X} 使你不加害我們，正如我們未曾侵犯你，素來善待你，並且送你平平安安地走。你是蒙耶和華賜福的！」 \end{tabularx} \\ \\ \relax
26:30 & \begin{tabularx}{0.7\textwidth}{X} 以撒為他們擺設宴席，他們就一起吃喝。 \end{tabularx} \\ \\ \relax
26:31 & \begin{tabularx}{0.7\textwidth}{X} 他們清早起來，彼此起誓。以撒送他們走，他們就平平安安地離開他去了。 \end{tabularx} \\ \\ \relax
26:32 & \begin{tabularx}{0.7\textwidth}{X} 那一天，以撒的僕人來，把挖井的消息告訴他，說：「我們得到水了。」 \end{tabularx} \\ \\ \relax
26:33 & \begin{tabularx}{0.7\textwidth}{X} 他就給那井起名叫示巴，因此那城名叫別是巴，直到今日。 \end{tabularx} \\ \\
[1ex]
\hline
\hline
\end{longtable}
$^{1}$各位弟兄姊妹平安.
我們一起有個祈禱.
親愛的主 多謝你.
你很愛我們.
我們很渴望你的話語.
願你的話語親自.
更新我們.
滿足我們 多謝主 奉耶穌名求.
阿們.
想問問大家當中.
有沒有父母.
是信耶穌的呢.
舉手來看看.
都有很多.
縱然大家.
年紀不小.
但原來大家的父母.
都信耶穌.
將大家帶到當中.
我想問問.
你的父母影響你所信.
有沒有影響到你所信呢.
你的神又是一位怎樣的神呢.
經歷了這麼多年之後.
究竟這一位神.
在你心裡面.
有沒有什麼改變呢.
還是其實是.
父母要求的.
從小時候你已經回教會.
去到今天.
你是否真的.
去相信這一位神呢.
當我們去看舊約聖經的時候.
很多時候我們都見到.
耶和華這樣介紹自己.
我是亞伯拉罕的神.
以撒的神.
雅各的神.
但其實聖經提到.

$^{41}$以撒的生命.
只用了五章.
基述亞伯拉罕.
是以撒的三倍.
即是十四,十五章.
都是講亞伯拉罕的生命.
但基述雅各的孫子.
竟然是四倍.
即是二十章.
這麼多.
坊間的裡面.
有很多書都是講.
亞伯拉罕和雅各的.
但偏偏.
真的很難找到關於以撒的書.
一個這麼特別.
乖巧的生命.
好像在人的眼中.
沒有什麼特別之處.
他整個生命都好像在複製.
重複他爸爸的生命.
好的壞的.
他都學.
我們一起來看看.
他們都是亞伯拉罕和以撒.
都曾經做過同樣的事情.
就是不承認自己的妻子.
因為他們的妻子都太漂亮.
他們很怕你來殺身之禍.
於是.
只說對方.
是我的妹妹.
但怎知.
都是惹來大禍.
令到.
當地的人很生氣.
第二件事.
就是和亞比米勒結盟.
亞比米勒和.
亞伯拉罕結盟的時候.

$^{81}$恩爵覺得.
你是被神祝福的.
後來他也和他的兒子結盟.
以撒結盟.
都是因為一樣.
讓大家都.
不要互不相害.
第三.
就是掘井.
以撒將他爸爸亞伯拉罕.
掘的井重新掘出來.
再.
以爸爸的.
所做的去改名字.
雖然以撒和他爸爸.
有著差不多的經歷.
但他們的生命.
截然不同.
因為亞伯拉罕是多國之父.
而以撒是兒子.
他擁有兒子的生命.
雖然.
篇章記錄的不多.
但以撒的生命.
其實都很值得我們去學習.
亞伯拉罕的神.
究竟怎樣成為以撒的神呢.
今天.
我們都會去看.
這五章的聖經.
看看以撒的生命.
讓我們可以從中學習一下.
第一點想和大家.
分享的是.
以撒有兒子.
完全接受.
恩典的心.
讓我們去重溫一下.
他爸爸.
和他一起.

$^{121}$經歷的事情.
我們先看這段經文.
是《創世記》26章.
3至5節.
這是耶和華.
祝福以撒.
祝福的.
就好像祝福亞伯拉罕一樣.
他重申.
要堅固.
他那裡說.
我要堅固我向你父親.
亞伯拉罕所起的誓.
我要使你的後裔增多.
好像天上的星.
又要將這一切的地.
賜給你的後裔.
並且地上的萬國.
都必因你的後裔得福.
因為亞伯拉罕聽從我的話.
聽從我的吩咐.
誡命.
律例和教導.
所以原來.
以撒得的福.
是因著他的爸爸亞伯拉罕.
他和爸爸的關係.
其實是盛載.
繼承爸爸的遺產產業.
和神的關係.
是得到.
耶和華的祝福.
因著亞伯拉罕的緣故.
他得到這些祝福.
所以無論原來和爸爸的關係.
還是和神的關係.
以撒的身份.
都是一個富二代.
信二代.
以及萬千寵愛在一身.

$^{161}$但是以撒.
其實不一定要接受.
爸爸的產業.
他未必一定要.
承繼爸爸的產業.
他要有足夠的容量.
才可以承擔.
那麼多的祝福.
以撒自己.
其實是一個承繼.
和接受恩典的人.
恩典不是每個人.
都因而接受恩典.
曾經我有一位朋友.
這樣和我說.
他說他的媽媽.
教導他不可以受到.
別人的恩惠.
不可以欠了別人.
所以他只可以做付出的那一位.
因此他遇到問題.
事情的時候.
都不找人去幫忙.
他自己一個很辛苦.
當我聽到的時候.
我就覺得.
很奇怪.
於是我就告訴他.
但是我相信的神.
是調動萬有的神.
祂會供給我一切的需要.
那些需要.
可能是透過我的朋友的幫助.
所以我是一個.
非常樂意接受人幫助的人.
原來.
真的不是每一個人.
都是那麼樂意.
接受人的幫助.
和接受恩典的.

$^{201}$而且有一些老人家.
他可能覺得不想麻煩別人.
不想別人幫忙.
不想別人去探望他.
但其實原來當他有需要的時候.
是必須學習.
接受別人的幫助.
因為接受幫助.
不代表你是軟弱.
不代表你是內弱.
他是因為你有.
更廣闊的勝在能力.
認知到這是神的恩典.
接受了.
讓神的恩典.
可以流動.
所以我們要學習.
以撒兒子的這一份.
接受勝在恩典的心.
接受地上.
爸爸的愛.
父母的愛 朋友的愛.
也接受天上.
爸爸的恩典.
以撒的神.
第二點要和大家分享.
就是耶和華爾了.
是一位.
必預備和供應的神.
我們在.
以撒的人生裡.
我們看到很多的故事.
都是神為他預備的.
第一個很出名的.
大家一定會記得.
他和爸爸一起在山上.
他小時候.
爸爸跟從神的吩咐.
將他帶上山去獻祭.
我們重溫一下.

$^{241}$這裡這樣說.
以撒對他父親亞伯拉罕說.
「我父啊.
亞伯拉罕說.
我兒 我在這裡.
以撒說.
看啊 火與柴都有了.
但繁祭的羔羊在哪裡呢.
亞伯拉罕說.
我兒 上帝必自己預備繁祭的羔羊.
於是二人同行.
當他想拔刀.
向以撒下手的時候.
神叫停他.
而且相信他是敬畏神的.
看到一隻羔羊.
兩角纏在松木裡.
亞伯拉罕牽著一隻羔羊.
去獻繁祭.
代替他的兒子.
亞伯拉罕稱那個地方為耶和華爾勒.
直到今天也有人說.
耶和華的山上必有預備.
以撒和他爸爸一起經歷耶和華爾勒.
他看見爸爸如何被神考驗.
他經歷了神如何拯救他.
神預備獻繁祭的羊羔.
和爸爸一起經歷的.
是耶和華爾勒.
那位神為他預備.
供應一切的神.
除了這件事之外.
神也為他預備妻子.
爸爸吩咐僕人在.
以撒媽媽的家鄉.
為他找一個妻子.
當時聖經形容以撒.
在田間的中間默想.
而僕人如何帶一個妻子回來呢?.
過程中.

$^{281}$聖經形容得很仔細.
但以撒卻一點也不知道.
他是安靜地在田間默想.
聖經沒有形容他在默想甚麼.
但或許他在掛念剛去世的媽媽.
神在他不知道的情況下.
已經為他預備了他需要的.
就是為他預備一位.
陪伴他一生的伴侶.
而且他一見到這位女子.
就非常喜歡.
神為我們預備的是最好的.
縱使在我們不知道的時候.
聖經繼續這樣說.
以撒領李伯嘉進入母親撒拉的帳篷.
娶她為妻並且愛她.
以撒自從母親離世之後.
才得到安慰.
神為我們安排的.
在我們不知道的時候.
他已經為我們安排.
人生大事來說.
其實找工場也是一件很大的事.
在我人生中我有兩次畢業.
兩次都很緊張因為要找工場.
然而兩次都是朋友介紹我.
進入他的公司.
在我不知道的時候.
原來神已經為我預備了工作的地方.
而且配套是我沒想過這麼適合我的.
原來無論愛情,工作,選學校.
所有的大事.
天父都會為我們預備合適的.
可能天父為我們預備一些人.
在我們身邊隨時幫助我們.
在我們還不知道的時候.
我們要做的就是全心的相信.
除了為他預備妻子.
神也為他預備飲食的.
聖經說那時候有饑荒.

$^{321}$但是我們看下去.
以撒在那個地方耕種.
在那一年有100倍的收成.
耶和華賜福給他.
讓他發達.
日漸昌盛成為大富翁.
以撒除了跟爸爸一起經歷神蹟之外.
他也經歷神很多的供應.
在他生命裡面.
他也經歷神不停的供應.
為他預備不同的事情.
最近家長團契當中.
家長也分享了很多神為他們預備的東西.
大家不停互相分享的時候.
大家都很開心.
例如有一位說他很需要炆怕水.
當他還沒說出來的時候.
就有人送給他了.
有一位媽媽說.
其實她一直都很需要一些嬰兒用品.
她還沒開口的時候.
就有很多人送給她了.
而有幾次.
我很快就要遲到的時候.
很需要的士.
我自己心裡面都求天父.
可不可以為我預備一輛的士呢.
下一秒就有的士的出現.
這些事情很多次都是這樣發生.
而我自己.
接下來會請大家去一個地方旅遊.
當地是會負30度的.
但是我還沒買大外套.
因為我覺得好像用一次.
又不太捨得.
要買雪外套.
但是一定要很保暖才行.
應該是很貴的.
我也不捨得去買.
然而有一位姐妹.

$^{361}$她說最近收拾房子.
找到一件外套.
有一件雪外套.
她很想借給我穿.
我覺得神真是為我們預備一切的神.
在小事上是可經可歷的神.
在饑荒的時候.
當我們覺得缺乏的時候.
原來神的恩典是一無所缺的.
第四是為她預備和好的筵席.
以撒覺得自己被懷恨.
因此被亞比米勒趕出去.
剛才我們主席也有讀到的經文.
就是講到以撒在那裡耕種.
有八倍收成.
他發達日漸昌盛成為了大富翁之後.
非利士人去妒忌.
於是把他父親挖的井塞住了.
而且亞比米勒要他離開這裡.
因為他很強盛.
原來以撒雖然他是富二代.
但他不是一帆風順的.
他覺得自己被懷恨.
甚至被欺負.
他的父親亞伯拉罕在世時挖的井.
也被非利士人塞住了.
以撒就重新挖出來.
而且按照他父親所取的名字.
去幫那些井改名字.
原來當被欺負的時候.
以撒所做的就是退讓.
不增敬.
被人搶他的水井.
他就另外再挖.
而且按照他父親所挖的名字.
再去重新挖.
相信他挖井的時候的心情是很委屈的.
然而他心裡也有那位神.
就是為他預備一切的耶和華以勒.
後來耶和華再次向他顯現.

$^{401}$其實是第二次向他表示.
會和他同在.
當夜耶和華向他顯現就說.
我是你父親亞伯拉罕的上帝.
不用怕.
因為我和你同在.
要賜福給你.
也要為我的僕人亞伯拉罕的緣故.
使你後裔增多.
然後聖經說.
以撒就在那裡築了一座壇.
求告耶和華的名.
並且在那裡支搭帳篷.
他的僕人就在那裡挖了一口井.
築了一座壇的意思就是在那裡獻祭.
以及求告耶和華.
向天父祈禱.
並且在那裡支搭帳篷的意思.
就是住在那裡.
停留在那裡.
並且做他要做的事.
就是挖井.
原來當以撒生命很無助的時候.
他照他父親所做的去做.
但是當父親不在旁邊的時候.
神向他顯現說.
會和他同在.
他讓神為他預備一切.
甚至幫助他去面對他的仇敵.
後來亞伯拉罕就和他的元帥去找以撒.
以撒對他們說.
既然你這麼憎恨我.
要趕我出來.
為何又來這裡呢?.
他們就說.
我們明明見到耶和華和你一起.
因此我們不如雙方起誓.
我們和你立約.
使你不要加害我們.
正如我們都沒有侵犯你.

$^{441}$我們素來都善待你的.
送你平平安安地走.
你是蒙耶和華賜福的.
以撒後來為他們擺設筵席.
和他們一起吃喝.
原來擺設筵席的意思.
即是和好的意思.
一起吃飯.
清早起來他們彼此起誓.
以撒送他們走.
平平安安地離開.
那一天以撒的僕人.
就將掘井的消息告訴他.
他們說我們得到水了.
他就將井起名叫「士巴」.
因此喇命叫「別士巴」.
直到今天.
這一次.
以撒親自的經歷.
神將他的敵人變為朋友.
亞伯拉罕的神成為了.
以撒生命中自己親自經歷的神.
他不是從父親那裡聽回來的.
而是他親自經歷神的幫助.
而且事情完結之後.
就得到了水井.
這個井叫「別士巴」.
即是盟約的意思.
他在這個井和他的敵人和好.
耶和華要和以撒立約.
兼顧他和亞伯拉罕的約.
而阿比米勒要和以撒立約.
不單止是因為他父親的約.
而是現在和他要立約.
而這個約.
無論是來自神和人的.
都是關於和平和賜福的約.
當我們和人的關係不和睦的時候.
耶和華以勒為我們預備一切的神.
會怎樣幫助我們呢.

$^{481}$我們可以學習以撒.
讓耶和華為你出頭.
不要和人爭敬.
你要做的事就好像以撒一樣.
捉壇 叩告耶和華.
然後住在那裡 停留在那裡.
做自己應該繼續要做的事情.
在我大專時期.
有一個好朋友.
不知道為甚麼無緣無故性情大變.
甚至惡言攻擊我.
我都不知道自己做錯了甚麼.
我只可以祈禱.
然後繼續完成我的畢業功課.
後來過了一段時間.
這位朋友主動找我.
原來他說他生命裡面.
他面對了一些事情.
所以他這樣.
因此原來是他的問題.
他告訴我不是我的問題.
後來我就明白到.
原來生命當中.
真的會有一些事情.
是會無緣無故的發生.
以致你不明白.
你因有而無辜地受害.
然而神是可以使人和好的神.
在合適的時間.
祂會為我們擺設筵席.
讓人和好.
我也有少許總結.
這個生命 這個約.
神向以撒所立的約.
向亞伯拉罕所立的約.
是由以色列民族到今天的約.
由亞伯拉罕的神.
以撒的神.
雅各的神.
整個的以色列民族.

$^{521}$去到基督.
然後到今天.
我們相信耶穌的人.
都可以成為神的兒女.
這一份的恩典.
其實是傳流了很多代.
去到我們今天.
天父也很想將這份產業給你.
將耶穌基督的愛給你.
將豐盛的生命給你.
今天以撒的生命.
去教導我們.
要以兒子的心.
去相信天父的供應.
你是不是也有以撒的心.
去接受這一切的恩典呢.
這一位以撒的神.
今天是不是成為了你的神呢.
讓我們現在也有一分鐘的時間.
我們去安靜在神的面前.
我們現在也有神兒女的身份.
當我們相信神的時候.
求主讓我們也有足夠的容量.
可以承載這個豐富的恩典.
或者我們也在等待神的預備.
求主讓我們有充足的信心.
去相信祂正在工作中.
就好像祂為以撒.
去預備利伯加一樣.
或許我們在難處的當中.
求主賜下信心.
相信祂會為我們擺設筵席.
讓我們一起祈禱.
親愛的天父.
你說萬代都要因亞伯拉罕得福.
以撒有兒子的身份.
如今我們相信你的人.
都可以有神兒女的身份.
得著你豐豐富富的恩典.
求你將我們有更大的承載能力.

$^{561}$接受恩典的心.
有兒子信服天父的心.
求主先讓我們真真實實的.
經歷你是耶和華已了.
那位必為我們預備.
供應一切的神.
以致我們可以和身邊的人.
分享這份愛.
多謝主奉耶穌名求.
阿們.
\newpage



\section{創世記 26:1-23}
\label{sec:HWk9pevFs74}
\textbf{2024-02-04 午堂《豐盛生命:依靠順服》創世記 26\_1-23 講員:陳麗蟬牧師}
\newline
\newline
連結: \href{https://youtube.com/watch?v=HWk9pevFs74}{\texttt{https://youtube.com/watch?v=HWk9pevFs74}} ~~~~ 語音日期: 2024-02-07
\newline
\newline
\hyperref[sec:vJ5BK_GiIWM]{\small{< < < PREV SERMON < < <}}
~
\hyperref[sec:index_chronic]{\small{[返順時目]}}
~
\hyperref[sec:index_scriptual]{\small{[返順卷目]}}
~
\hyperref[sec:ny6z0RWTL8g]{\small{> > > NEXT SERMON > > >}}
\newline
\newline
創世記 26:1-23
\newline
\begin{longtable}{cl}
\hline
\hline
章節 & 經文 (和合本修訂版)\\
\hline
26:1 & \begin{tabularx}{0.7\textwidth}{X} 那地有了饑荒，不是亞伯拉罕的時候曾有過的那次饑荒，以撒就到基拉耳，非利士人的王亞比米勒那裡去。 \end{tabularx} \\ \\ \relax
26:2 & \begin{tabularx}{0.7\textwidth}{X} 耶和華向以撒顯現，說：「你不要下埃及去，要住在我所指示你的地。 \end{tabularx} \\ \\ \relax
26:3 & \begin{tabularx}{0.7\textwidth}{X} 你要寄居在這地，我必與你同在，賜福給你，因為我要將這一切的地都賜給你和你的後裔。我必堅定我向你父親亞伯拉罕所起的誓。 \end{tabularx} \\ \\ \relax
26:4 & \begin{tabularx}{0.7\textwidth}{X} 我要使你的後裔增多，好像天上的星，又要將這一切的地賜給你的後裔，並且地上的萬國都必因你的後裔得福， \end{tabularx} \\ \\ \relax
26:5 & \begin{tabularx}{0.7\textwidth}{X} 因為亞伯拉罕聽從我的話，遵守我的吩咐、誡令、律例和教導。」 \end{tabularx} \\ \\ \relax
26:6 & \begin{tabularx}{0.7\textwidth}{X} 於是，以撒住在基拉耳。 \end{tabularx} \\ \\ \relax
26:7 & \begin{tabularx}{0.7\textwidth}{X} 那地方的人問起他的妻子，他就說：「她是我的妹妹。」原來他害怕說「我的妻子」。他想：「或許這地方的人會因利百加殺我，因為她容貌美麗。」 \end{tabularx} \\ \\ \relax
26:8 & \begin{tabularx}{0.7\textwidth}{X} 他在那裡住了一段很長的日子。有一天，非利士人的王亞比米勒從窗戶往外觀看，看哪，以撒在撫愛他的妻子利百加。 \end{tabularx} \\ \\ \relax
26:9 & \begin{tabularx}{0.7\textwidth}{X} 亞比米勒召以撒來，說：「看哪，她實在是你的妻子，你怎麼說『她是我的妹妹』呢？」以撒對他說：「因為我想，恐怕我會因她而死。」 \end{tabularx} \\ \\ \relax
26:10 & \begin{tabularx}{0.7\textwidth}{X} 亞比米勒說：「你向我們做的是甚麼事呢？百姓中有一個人幾乎要和你的妻子同寢，你就把我們陷在罪中了。」 \end{tabularx} \\ \\ \relax
26:11 & \begin{tabularx}{0.7\textwidth}{X} 於是亞比米勒命令眾百姓說：「凡侵犯這個人，或他妻子的，必要把他處死。」 \end{tabularx} \\ \\ \relax
26:12 & \begin{tabularx}{0.7\textwidth}{X} 以撒在那地耕種，那一年有百倍的收成。耶和華賜福給他， \end{tabularx} \\ \\ \relax
26:13 & \begin{tabularx}{0.7\textwidth}{X} 他就發達，日漸昌盛，成了大富翁。 \end{tabularx} \\ \\ \relax
26:14 & \begin{tabularx}{0.7\textwidth}{X} 他有羊群牛群，又有許多僕人，非利士人就嫉妒他。 \end{tabularx} \\ \\ \relax
26:15 & \begin{tabularx}{0.7\textwidth}{X} 他父親亞伯拉罕在世的時候，他父親的僕人所挖的井，非利士人全都塞住，填滿了土。 \end{tabularx} \\ \\ \relax
26:16 & \begin{tabularx}{0.7\textwidth}{X} 亞比米勒對以撒說：「你離開我們去吧，因為你比我們強盛得多。」 \end{tabularx} \\ \\ \relax
26:17 & \begin{tabularx}{0.7\textwidth}{X} 以撒就離開那裡，在基拉耳谷支搭帳棚，住在那裡。 \end{tabularx} \\ \\ \relax
26:18 & \begin{tabularx}{0.7\textwidth}{X} 他父親亞伯拉罕在世的時候所挖的水井，在亞伯拉罕死後，都被非利士人塞住了，以撒就重新把井挖出來，仍照他父親所取的名為它們命名。 \end{tabularx} \\ \\ \relax
26:19 & \begin{tabularx}{0.7\textwidth}{X} 以撒的僕人在谷中挖井，就在那裡得了一口活水井。 \end{tabularx} \\ \\ \relax
26:20 & \begin{tabularx}{0.7\textwidth}{X} 基拉耳的牧人與以撒的牧人相爭，說：「這水是我們的。」以撒就給那井起名叫埃色，因為他們和他相爭。 \end{tabularx} \\ \\ \relax
26:21 & \begin{tabularx}{0.7\textwidth}{X} 以撒的僕人又挖了一口井，他們又為這井相爭，以撒就給這井起名叫西提拿。 \end{tabularx} \\ \\ \relax
26:22 & \begin{tabularx}{0.7\textwidth}{X} 以撒離開那裡，又挖了一口井，他們不再為這井相爭了，他就給那井起名叫利河伯。他說：「耶和華現在給我們寬闊之地，我們必在這地興旺。」 \end{tabularx} \\ \\ \relax
26:23 & \begin{tabularx}{0.7\textwidth}{X} 以撒從那裡上別是巴去。 \end{tabularx} \\ \\
[1ex]
\hline
\hline
\end{longtable}
$^{1}$弟兄姊妹早晨.
知不知道下個禮拜日是什麼日子?.
是啦,是年初二.
在準備新年期間送大家一句聖經.
等候耶和華的必從新得力.
他們必如鷹展翅上騰.
他們奔跑卻不困倦.
行走卻不疲乏.
這樣怎麼能過一個有力的生命呢?.
有人說,有兩個很重要的.
第一個就是對上帝的完全的倚靠.
第二就是對上帝完全的信服.
今天我們說的這個聖經人物.
以撒在生命當中就有了這兩個特點.
先對以撒有一點認識.
他的名字的意思是「笑」.
很開心地笑.
為什麼他的名字叫「笑」呢?.
這麼開心地笑呢?.
原來他的出生可以說是一個神蹟.
是上帝特別賜給他父母的.
因為當他出生的時候.
他父親是亞伯拉罕.
已經一百歲了.
他母親就九十歲了.
一般來說已經過了生育的年紀.
所以是不可能的.
不過亞伯拉罕所投靠的上帝.
是凡事都能的.
他也可以使得「無」變為「有」.
以致他父母雖然年紀這麼大.
但是他仍然可以出生.
仍然可以生育.
當他一出生的時候.
聖經形容他母親「笑都不吭聲」.
就改了這個兒子.
以撒就是「喜笑」.
很開心地笑.
當我們問.
在這個創世記.

$^{41}$或者猶太列祖裡面.
我們認識哪幾個呢?.
通常第一個就是說哪個?.
亞伯拉罕.
然後呢?.
雅各.
以撒通常都被人忽略了.
但是在這個聖經裡面.
我們很多時候都會說.
就是這個亞伯拉罕的神.
以撒的神.
和雅各的神.
所以都是列祖之一.
在這三個.
很多時候聖經都說.
亞伯拉罕的神.
以撒的神.
雅各的神.
在裡面最長命的那個.
就是以撒.
最長命的是以撒.
但是整個創世記裡面記載.
記載得最少的那個.
又是以撒.
亞伯拉罕他175歲.
但是創世記記載他.
由第12章一直到25章.
足足有13章是他的.
他的兒子更厲害.
他的兒子還短命.
147歲而已.
但是由這個創世記第27章.
一直到49章.
主要都是由他做主角.
所以就說了22章的聖經.
這個以撒.
做主角就只有26章.
剛才啟仁弟兄幫我們讀了大部份.
其他都有的.
不過是一些配角.

$^{81}$不是他做主角.
這個以撒可以說是一個信義代.
因為他爸爸被稱為信心之父.
在對耶和華的信心.
是給了以撒和後代一個很好的榜樣.
我們看看這個以撒.
因著他爸爸影響的緣故.
他對上帝都很信服.
我們看看在當中聖經記載了三件事情.
第一件很出名的.
翻得久一點的弟兄姊妹都知道的.
就是當以撒大約二十多三十歲的時候.
有一天他爸爸就說要上山獻祭.
以撒就揹著柴和火.
但是一路走去很奇怪.
為什麼爸爸沒有帶製品羊羔呢.
問爸爸.
爸爸就告訴他.
上帝自己一定會預備的.
去到這個獻祭的地方.
他就跟著將以撒綁在祭壇上.
但是當時以撒已經二十多三十歲了.
我想他一定會奇怪的.
不奇怪的那個就傻了.
一定會奇怪.
為什麼無端綁在這裡呢.
也都可能心知不妙了.
但是完全沒有反抗.
為什麼他這麼乖這麼聽話呢.
第一 他很尊重他很愛的父親.
他很尊重爸爸.
他覺得爸爸要怎麼做就怎麼做.
第二 他很信得過爸爸所信的上帝.
我看過爸爸之前這麼多的經歷.
他會帶領爸爸一步一步經歷的.
所以他覺得.
耶和華吩咐了爸爸要這麼做的時候.
他就一定要去信服.
他也都知道上帝一定會負責任的.
就算到最後他成為了製品被獻上的時候.

$^{121}$他都覺得是應該的.
其實在那次的獻祭裡面.
上帝不是真的要以撒的性命.
而是要去看伯拉罕的心.
他要將這個兒子.
是上帝給過他的.
但也都變成了那段時間.
成為了伯拉罕和他的太太撒拉的心肝寶貝.
什麼都給了他.
他就要看上帝要看伯拉罕的心.
在以撒那裡.
還是繼續在耶和華那裡呢.
有時候我們都是這樣的.
上帝從我們的生命當中.
他會拿走一些東西.
我們覺得很寶貝.
當他拿走了的時候.
我們就會覺得很慘.
要生要死的.
但上帝其實就讓我們看到.
要看到我們.
對這件心肝寶貝的東西重要一點.
還是上帝重要一點呢.
我在教育學院畢業的時候.
拍了拖兩年了.
在我畢業之後.
我就發覺到.
我那個男朋友疏遠我.
又不再打電話給我.
又不再約我出街.
後來就知道.
原來在教會裡面找到另外一個人.
當時很難過.
很難過.
很多的埋怨.
很多的不忿.
但也在那次.
上帝開我的眼睛.
讓我看到人是會變的.
最初是他先找我.

$^{161}$但竟然兩年之後.
他就去找了第二個.
他找到第二個比我更加好的.
但是我一直在祈禱.
我一直在祈禱上帝.
幫助我.
因為那時候開始出來教書.
那怎麼辦呢.
我求上帝給我有平安.
我每一天都用一個多小時的時間去祈禱.
暑假嘛.
用一個多小時的時間去祈禱.
上帝很快就給我有平安.
並且都給我想得很清楚了.
在世界上唯一可以完全依靠的是上帝.
之前人的山盟海誓都是一個「港」字.
那次我跟上帝說要一生一世的事奉他.
那次就是我跟上帝第一次的立約.
伯拉罕就是想把兒子獻給上帝.
但以撒很乖很乖.
不要掙扎.
結果上帝就為他們預備了這個羊羔.
還有在婚事上.
以撒更厲害.
伯拉罕就離開了他的家鄉吳爾.
他又不想以撒嫁過迦南地的人.
外邦人.
所以他吩咐過魯埔人.
是老眼昏花的魯埔人.
回鄉幫他找老婆.
唯一的條件就是一定要信耶穌.
不是信耶穌.
是信耶和華的.
然後又肯來這裡跟他們一起住.
那這個以撒完全沒有反抗.
他覺得耶和華為他預備的東西.
為他安排的東西一定是最好的.
但現在我們怎麼會呢.
一定反抗到底.
然後就到了剛才我們所讀的經文.

$^{201}$就說到他當時在基拉爾這個地方住.
不是 對不起.
應該是在別士巴這個地方.
綠色箭嘴那個地方住.
有這個饑荒.
他就要去到這個基拉爾那裡.
紅色那裡.
想著再在基拉爾那裡.
就會下去埃及.
在我們左手邊那裡 埃及那裡.
就因為埃及被稱為世界的糧倉.
周圍的地方都是這個饑荒.
無論在這個別士巴或者基拉爾.
其實都面對著這個饑荒.
沒有雨下 在這個乾旱的情況裡面.
但是上帝就向他顯現.
他說你不能夠下埃及.
你要繼續留在基拉爾.
我會賜福給你.
我會賜福給你的.
各方面的祝福會給他的.
那你想一下.
其實這個基拉爾.
都是這個饑荒的裡面.
如果上帝賜福.
早在別士巴就賜福.
為什麼去到基拉爾.
準備下去埃及的時候.
才有這樣的吩咐呢?.
在不同的時間裡面.
我們都會有不同的難處.
不同的低谷.
但是當我們看著這些低谷的時候.
我們就會有很多的擔心.
如果以撒他在那裡基拉爾.
覺得這裡都是在饑荒.
你不讓我下去埃及那裡.
因為埃及那裡有條尼羅號.
所以他就有足夠的水來灌溉.
和給人飲用.

$^{241}$所以我們看到在聖經裡面.
每逢有饑荒.
人們就走去埃及.
但是在埃及裡面.
很多很多很多的偶像.
上帝就不讓他們去.
這次.
保羅也告訴我們.
他說我們作為信徒的人.
作事要憑信心.
不憑硬件.
當我們一看見一些艱難的事情的時候.
我們覺得慘了慘了.
我們就會覺得.
那個信心就會動搖.
當我決定了.
我要去讀神學的時候.
我家人個個都覺得我很慘.
甚至我媽媽爸爸我姐姐.
罵得我飛起來.
我覺得教書不是很好嗎.
鐵飯碗來的.
收入又穩定.
是不是一下子.
讀神學啊.
讀神學那四年沒有收入的.
還要做傳道.
不知道有多少錢.
經常問我多少錢一個月.
家裡的人個個都覺得.
我很淒涼的.
他們個個都這樣.
但是到了沙士那個年.
我們看到.
各行各業都很低落.
我家裡的人.
哥哥弟弟那些.
本來做工那些很穩定.
但是在沙士之後.
個個都很怕.

$^{281}$很怕會不會被解僱.
就算我們自己家裡的家族生魚蛋粉.
當年.
到後來一年之後.
我們都沒有辦法.
我們都要撿起來.
當時我媽媽就自己去說.
她說妹妹.
想不到最穩定的那個是你.
人就是憑硬件.
當我要去讀神學的時候.
她就覺得我很淒涼.
是不是不穩定.
但是當真真正正的艱難來到的時候.
他們看到.
在當中有上帝保護的時候.
最穩定的是屬於耶穌基督的人.
所以過程裡面.
有時候我們都會擔心.
或者過程裡面.
我們都會覺得.
怎麼辦呢.
但是堅持到底.
信靠 信服.
好了 以撒.
很乖.
是不是很乖.
但是生命當中.
始終都會有缺陷的.
當他面對這個危險的時候.
他仍然會有軟弱的.
他就說謊.
當他住在這個基拉爾.
是菲律賓人的地方.
他的太太很漂亮.
當他面對菲律賓人的王亞比米勒的時候.
他就告訴亞比米勒.
他說這個是我妹妹.
但是.
後來這個秘密就被亞比米勒王知道了.

$^{321}$發現是說謊.
明明是他的老婆.
所以菲律賓王就責備以撒.
你知不知道你差點害死我們.
因為有人真的想和他同枕.
有人想娶他做老婆.
以撒就承認了.
他說因為我怕你們娶了我太太.
又恐怕你因為我太太這麼漂亮.
你會殺死我.
其實很自私的表現.
因為怕自己會被殺死.
所以說了這個謊話.
就不理他太太的感受.
大家有沒有說過謊話.
應該倒過來問.
有沒有人沒說過謊話.
昨晚晚堂的人有沒有人敢舉手.
今天還是沒有人敢舉手.
說謊話好像一個天性.
一定會做的.
但是有時候在說謊話的過程裡.
會傷害了人.
又想想.
在你自己有這個思想的裡面.
你最小的時候說了什麼謊話呢.
我自己一個謊話.
在我自己有記憶的裡面是第一個謊話.
不過不是的.
之前很多的.
很多的.
之前是三年級的時候.
那時候因為自己很想去玩.
我那時候在灣仔讀書.
在聖後閣福屏會.
舉辦了一個叫做中秋賞月晚會.
晚上七點半到九點半就上去灣仔寶雲道.
我知道媽媽一定不讓我去.
但是我又很想去.
我就騙她.

$^{361}$我就說我在聖後閣福屏會那裡舉辦的.
九點半就完了.
還求我媽媽讓我帶我弟弟.
我弟弟那時候讀一年級.
去參加.
那就沒事了.
神不知鬼不覺.
誰知道那次上了寶雲道.
不知道為什麼.
他們十點鐘才下山.
我回到山上的時候.
已經是十點半有多了.
當解散的時候.
我就見到我媽媽在那條斜路那裡.
一見到我們的時候.
她真的一路哭.
她說嚇死她了.
我記得當時我媽媽在那裡哭.
那份緊張.
她說九點半在這裡解散.
我就九點半來這裡.
想著帶你們回去.
來到見到這裡完全沒有燈光.
沒有東西.
我多怕我兒子和我女兒不見了.
其實到今天說起來也有一點內疚.
就是因為這個謊言.
嚇得我媽媽不見了兩個兒女.
那份擔心.
所以求主幫助我們.
有時候會很順口開河就說這個謊言.
但求主幫助我們.
盡量保守我們的口舌.
免至你去闖成這個大禍.
最後我又繼續說以撒.
她又繼續遇到不同的艱難.
繼續在這個基拉爾這個地方.
生活非常的艱苦.
人們沒有水.
不夠水喝.

$^{401}$但正是以撒.
她所種的東西.
有百倍收成.
很厲害.
你種一粒米.
就有100粒出.
豐收期有足夠的水.
但不知道有沒有這麼多.
所以那些非利士人就開始妒忌她.
嘩 怎麼她這麼厲害.
通常妒忌人是怎樣.
想辦法踢人走.
於是她的王就走去和以撒說.
趕她走.
我們不要你在這裡.
不單止趕了以撒走.
還有當以撒離開之後.
她每次去刮井.
她刮的井就跑去爆.
其實這些井是之前她爸爸亞伯拉罕刮的.
開發了這些井.
又是被這些非利士人.
拿些泥土塞住.
到現在當以撒又跑去.
拿出來 開出來的時候.
她就不是愛她的井.
她就說這個井裡面的水是餓的.
又是非利士地.
很沒有功利.
是啊 你可以用井.
但你不可以拿裡面的水.
那就是沒用.
井對於當時的游牧人的民族來說.
是很重要的.
有井就有水.
有水就有生命.
對我們現在來說.
水也是很重要的.
很感恩一開水就有水.
不過久不久我們家裡就會有製水.

$^{441}$這個很麻煩的.
但既然水這麼重要.
但是很奇怪.
這個以撒沒有去跟非利士人爭水.
也沒有去跟他們打仗吵架.
她只是忍耐.
你說不讓我用這個井的水.
我就去開第二個 又去開第二個.
以撒並不是內弱.
也不是沒有能力跟非利士人去相爭.
不過只是因為在耶和華裡面.
她屬於上帝的.
她有一個廣闊的胸襟.
她寧願自己吃虧.
她也覺得自己走遠一點.
因為吵架打架都無謂的.
大家都有損傷.
但是她無論一路走一路走.
一路走開的時候.
上帝很大很大地賜福給她.
在聖經裡面說.
她所種的東西是百倍收成.
有很多的羊群,牛群.
刮井游順利.
一個,兩個,三個.
被稱為當時的大富翁.
她能夠有這麼好的收成的時候.
那個非利士人.
雖然趕走她 想著趕絕她.
誰知道她仍然這麼順利.
她又落在一個很困難的當中.
這個非利士人就開始走去跟以撒講和.
他就跟以撒說.
不如我們簽訂和約吧.
這個以撒.
給我們這些人.
之前又趕我走.
之前又一路欺負我.
不讓我拿井的水.
現在你就來求和.

$^{481}$我就未必願意了.
奉上這個金銀珠寶.
但是這個以撒.
她就有上帝那份謙和,忍讓的心.
她化敵為友.
我相信也是的.
在地上高多一個朋友.
少一個敵人.
不單止這樣.
其實她自己對上帝的信服.
對人的忍讓裡面.
她寫到上帝的名字得榮耀.
因為這個阿比米勒就說.
我們明明看見耶和華與你同在.
從這個以撒她所做的事.
她一直那個信服.
對他們自己這個非利士人那些的哈巴.
他們忍讓一路退一路退.
但上帝繼續的祝福她.
是有這個豐收.
她就說看到耶和華就與這個以撒同在.
弟兄姊妹.
我們信了耶穌那麼久.
有沒有人跟你說過.
你信耶穌是與周圍的人是很不同的.
我盼望.
盼望就算初信信了很久也好.
有人從我們生命.
從我們所作的.
就是這個.
真是看到我們在裡面有上帝的.
看到上帝在當中祝福我們.
以致到將榮耀歸給上帝的名字.
在《肥臘比書》第四章第五節.
保羅教導我們這樣說.
當叫眾人知道你們謙讓的心.
就是要我們在與人一起相處的時候.
走去這個征戰.
打仗.
沒有用的.

$^{521}$但在一起相處的裡面.
以謙讓來相處的時候.
這個可以叫更多來去榮耀上帝的名字.
在《羅馬書》裡面保羅更加講得清楚.
我們就去將保羅教我們這個榮耀.
忍讓的行動我們一起來讀.
1,2,3.
不要以惡報惡.
眾人以為美的事.
要留心去作.
若是能行.
總要盡力與眾人和睦.
親愛的弟兄.
不要自己伸冤.
寧可讓步.
傾憑主路.
因為經上記著主說.
伸冤在我.
我必報應.
所以你的仇敵若餓了.
就給他吃.
若渴了就給他喝.
因為你這樣行.
就是把炭火堆在他的頭上.
你不可為惡所成.
總要以善成惡.
先去上高的章.
我們看到耶和華.
雖然他有軟弱去講謊話.
但上帝也一路祝福他.
他的信服.
他對人的忍讓.
我們看到這個以殺.
是一個蒙福的人生.
雖然只是《詹聖經》裡面.
做主角只有一章.
但我們看到上帝臨到當中的時候.
他經歷那份豐盛.
其實想清楚一點.
在風調雨順的時候.

$^{561}$大家都有豐收.
不奇怪吧.
你有很多.
我又有很多.
但就出了出奇.
在乾旱的時候.
人人都沒有.
但唯獨以殺.
他就有很多的收成.
他有很多很多的羊群和牛群.
就讓我們看到.
原來上帝的能力.
是多麼大的.
上帝就是這麼奇妙的.
這樣祝福著以殺.
當以殺有信服.
明明要去世界糧倉.
埃及的時候.
就在當中才有發展.
但是呢.
但是上帝就說.
你要留在這個基拉爾.
所以唯有他留在.
當時人人都覺得不容易.
並且是艱難的地方.
但是呢.
上帝同在上帝祝福的時候.
上帝就給他.
看到他那個榮耀.
他那個能力.
在我讀完中學的時候.
剛好在那段時間.
就做這個劇集.
就講這個新聞界的女強人.
汪阿姐做主角的.
所以我就立志了.
立志我都要在新聞界裡面去發展.
但是其實我一點都不喜歡作文.
但是我看得很清楚.
上帝各方面.

$^{601}$來帶我進教育學院讀書.
所以在教育學院.
去面試的時候.
我特意教到最後最後的秒鐘.
我才出門.
還要穿一條又沒毛又很短的裙子.
我去面試.
我以為一定不行的.
誰知道呢.
上帝他的大能.
就不知道怎麼掩飾那些導師.
總之要我進就要進.
但是我覺得呢.
當時我真的很憤怒.
很憤怒上帝.
很不滿上帝為什麼要這樣.
我對我家人也都很不喜歡他們.
我生家人氣很久了.
但是當我後來回想.
我回想的時候.
我讚嘆上帝的奇妙.
我覺得他對我陳麗善這個人的認識.
比我自己的認識更加清楚.
就藉著在神學院裡面的訓練.
上帝就讓我去學了.
怎麼樣去教書.
怎麼樣去教導.
以致人可以能夠去吸收得更加的清楚.
原來上帝在每一段東西.
要安排我們去做的時候.
是為了對我們的生命更加有益.
在當時的那一刻.
可能我們不是很清楚.
但是當我們回望的時候.
我們能夠去發現得到.
所以弟兄姊妹.
凡願意好像以撒那樣.
你去完全的投靠上帝.
信服上帝.
在我們生命當中的安排的時候.

$^{641}$我們會經歷到一些出人意外的福氣和平安.
所以由上帝去掌管著我們的生命.
是會燦爛,是豐富的.
有一個全世界很出名的小提琴家.
有一次他舉行了一個音樂會.
他拉完之後.
他就說.
今天我來到你這個城市.
我下午就買了一個十幾萬的小提琴.
明天你們要再來聽.
要儘早入場八個好位.
我就會用這個十幾萬的小提琴.
去拉這些音樂.
到第二天.
這個音樂真的很好聽.
完結的時候.
所有的觀眾就立即起來.
就拍爛手掌.
他說好琴,好琴,好琴.
這個小提琴家就起來了.
突然間就把這個琴扔在地上.
琴就爛了.
然後就把這個爛琴拿起來.
就跟他們說了.
他說這個小提琴不是十幾萬.
只是幾百元從雜貨攤買回來的.
不過這樣做我要告訴你們知道.
剛才能夠有這麼美妙的音樂出來.
不是因為這個琴是十幾萬.
而是因為拉琴的人的技術,技巧.
所以各位姐妹.
無論我們的生命是屬於一個.
很貴很貴的琴.
還是一個很便宜幾百元的琴.
但只要我們在上帝的手裡面.
順服祂的安排.
對祂完全完全的依靠的時候.
我們生命當中.
都能夠奏出一篇美妙的樂章.
生命要完全的順服,依靠.

$^{681}$我們唱一首剛才的詩歌.
我們先想一想.
將我們生命裡面.
現在我們面對的任何一些艱難.
一些我們不明白的東西.
我們交給上帝.
然後我們才去唱這首歌.
請司琴給我們前奏.
我們有力地去唱.
心中盼望 全歸寄託我.
夜話全能記住.
你不畏懼 心裡不驚惶.
帶領我衝勁逃海浪.
隨時助我 面對迷路的.
心不孤單也不畏懼.
全能上帝同伴這世間所有.
世界各國四重圓圓.
贈你獨一生 心只有你.
我可一切心中寡慮.
我要盡全力 尊榮留中識.
尊重你 寄託救主.
我們一起禱告.
天父上帝我們謝謝你.
因為人生裡面.
必定有高有低的時候.
主祝福我們每一位頂子妹.
我們知道你們就是我們生命的主.
你們就是那位有大能大力的子.
主凡事都能的主.
是可以使動無邊為有的主.
主就讓我們在生命當中.
或高或低.
我們都要完全的依靠 信服.
求主幫助我們.
帶領我們走人生的路程.
在新的一年開始的時候.
我們都繼續將我們的生命交給你.
讓我們時時刻刻的尊主你為大.
多謝你瑞聽禱告.
是奉靠主基督的聖名.

$^{721}$Amen.
\newpage



\section{瑪拉基書}
\label{sec:ny6z0RWTL8g}
\textbf{2024-02-08 周四晚查經班《瑪拉基書(5)》 講員:曾裔貴牧師}
\newline
\newline
連結: \href{https://youtube.com/watch?v=ny6z0RWTL8g}{\texttt{https://youtube.com/watch?v=ny6z0RWTL8g}} ~~~~ 語音日期: 2024-02-09
\newline
\newline
\hyperref[sec:HWk9pevFs74]{\small{< < < PREV SERMON < < <}}
~
\hyperref[sec:index_chronic]{\small{[返順時目]}}
~
\hyperref[sec:index_scriptual]{\small{[返順卷目]}}
~
\hyperref[sec:OWpYTLtxpsA]{\small{> > > NEXT SERMON > > >}}
\newline
\newline
$^{1}$多謝天父的恩典讓我們今天晚上.
雖然都是近年晚.
大家都會忙碌的.
多謝你讓我們一起再有機會去研讀聖經.
讓我們能夠對《馬拉基書》有掌握,有明白.
更加能夠在我們自己的生活敬拜.
能夠符合神的心意.
願你帶領我們.
我們一起在你面前等候禱告.
奉耶穌基督的名.
廈門.
都是緊接上一個星期.
有些是讀得不夠的.
上個星期本來想跟大家讀二十二章的.
但是就記錯了.
記錯了看了二十一章.
二十一章是講祭司的守則.
祭司是怎樣才合資格事奉.
但是二十二章其實就跟《馬拉基書》.
第一章近尾和第二章是有直接關連的.
就是為什麼要責備當時的祭司.
就是他們獻的祭不符合規定.
什麼祭才是符合神的規定呢?.
二十二章就有這樣的教導.
二十二章你還沒記嗎?.
創出利民生.
即是摩西五經的第三卷書.
摩西五經的第三卷《利未記》.
《利未記》就是在抗耶的時候寫.
大概是一個月的時間.
就是他們出了埃及之後.
在西奈抗耶那裡頒佈律法.
第二十二章.
所以《利未記》就有二十七章.
所以是關乎到祭司的手冊.
祭司的手冊.
來做祭司.
二十二章是這樣說的.
耶和華對摩西說.
你吩咐亞倫和他子宣說.

$^{41}$要遠離以色列人所分別為性.
歸給我的聖物.
免得褻瀆我的姓名.
我是耶和華.
你要對他們說.
你們世世代代的後裔.
凡身上有污穢.
親近以色列人所分別為性.
歸耶和華聖物的.
那人必在我面前剪除.
我是耶和華.
亞倫的後裔.
凡長大麻瘋的.
或是有流症的.
不可吃聖物.
直等他潔淨了.
無論誰摸那因死屍不潔淨的物.
或是為精的人.
或是摸什麼使他不潔淨的爬物.
即是死老鼠,死蟑螂.
或是摸那使他不潔淨的人.
不驅那人有什麼不潔淨.
即是碰到那些人.
摸了這些人或那些物件.
必不潔淨到晚上.
如果不用水洗身.
就不可吃聖物.
即是那些建制元的.
祭肉歸祭司.
但如果那天你下班後.
不小心摸到那些不潔之物.
或是你碰到那些大麻瘋的人.
你又跟他一起.
那你就不潔.
當晚就不可以吃聖物.
一定要等到那個潔淨.
還有第二天就可以了.
所以第七節.
日落的時候他就潔淨了.
即是第二天了.

$^{81}$他們是黃昏六點鐘就是當一天.
接著由黃昏六點之後就當第二天了.
他就潔淨了然後可以吃聖物.
因為這是他的食物.
記住了.
祭司是有份於神的獻祭的.
獻祭給神其實某程度是獻祭給祭司的.
其實是給祭司的食物生存,維生.
所以提摩太前書第五章.
那些善於管理教會的長老.
當配受加配的敬奉.
那些勞苦全道的就更加如此了.
接著阿保羅就引述聖經.
就是你沒記經文.
那些牛在牆上踩穀.
就不要摟著牠的嘴.
即是牠一路在爬地的時候.
可能一路在吃東西的.
你讓牠吃.
即是牠是生畜來的.
但是都讓牠一路吃飽才有力去做.
那就工人得工價是應當的.
那是大概意思.
所以這裡祭物變成食物.
獻祭給神的聖物.
變成祭司的食物.
不過牠吃的時候就要有一個警告.
這個是獻了給神的祭物.
吃的人如果很隨隨便便.
你想想一個窮全道.
大江斯丹頓駕駛普什哲.
不是不行的.
不過好像有些.
在外國真的是這樣的.
有些教會的大教會.
特別是靈恩那些.
因為他提出的信念就是說.
神是豐富的.
所以神的僕人是不可以這麼寒酸的.
好像挺合理的.

$^{121}$但是這裡告訴我們知道.
那個獻祭的人.
他享用這些是神的聖物.
但是成為了他的食物.
就是他的衛生.
就是他的賴以衛生的工作.
但是他吃的時候要有講究的.
他不可以沾染污穢.
就是說他自己很警覺.
我是服侍神的.
那些百姓獻給神的祭物.
成為了我的食物.
所以這裡是大家都要留意的.
我們去敬拜神.
去尊榮神.
但是如果我們對牧者沒有了尊榮.
其實就不是這麼健康的.
所以這裡告訴我們知道.
這些的聖物.
因為這是他的食物.
是神已經給了祭司和他的後人的.
第八節.
致死的.
就是暗命死了的那些.
或者被野獸撕裂的.
他就不可以吃了.
更加不可以獻祭.
因此污穢自己.
就是去打獵,打野味.
那祭司就不行的.
如果他在街上看到.
哇 還是很漂亮的.
沒事的.
但是死了.
不知道什麼原因死了.
那都不可以吃的.
祭司都不可以拿回去煮來吃的.
第九節.
所以他們要守我所吩咐的.
免得輕忽了.

$^{161}$因此擔罪而死.
我食叫他們成聖的耶和華.
其實有利未記教導的.
待會我們一揭去《馬拉基書》第二章.
你就會知道.
其實完全犯禁.
就是照足這些禁忌去犯.
第十節.
凡是愛人不可以吃聖物.
愛人就是不是亞倫之派的.
他不是利美家族的.
就不可以吃聖物了.
記不記得大衛王逃走.
逃避他的兒子阿撒隆的時候.
對不起.
逃避蘇羅的追殺的時候.
在這個列王記下的十六章.
對不起.
《沉悟二記》上的十六章.
逃避蘇羅的時候.
他就去到聖殿會幕.
祭司就問.
很害怕.
見到大衛不是獻祭的時候.
大衛就說有沒有東西吃.
我們很久沒有吃東西了.
和他的手下.
祭司就說.
有的.
不過是一層切餅.
就是剛剛更換的層切餅.
其實是祭司才可以吃的.
大衛就說.
我真的很大需要.
祭司就勉為其難.
給了他吃層切餅.
本來供奉神.
是給祭司的下蘭.
祭司更換之後就可以吃的.
就給了大衛吃.

$^{201}$主耶穌在《馬太福音》.
都有引述這個事件的.
但這裡告訴我們.
原則上.
外人是不可以吃聖物的.
寄居在祭司家的.
或是僱工人.
都不可以吃聖物的.
所以你不要自己去領聖餐.
自己看不到.
牧師剛剛領開一個麵.
來時牧師領開一個麵.
你就拿.
馬上拿一個聖餅.
放在袋子裡.
然後拿回去給飛鴻吃.
這樣是不行的.
這樣他真的會吃了聖物.
是會被民眾剪除的.
所以你不要害他.
倘若祭司買人.
是他的錢買的.
那人就可以吃聖物了.
生在他家的人都可以吃.
祭司的女兒如果嫁了外人.
就是外嫁女兒.
就不可以吃了.
就不可以吃舉祭的聖物.
但祭司的女兒.
因為已經成為了寡婦了.
或者被休了.
又沒有孩子了.
又回娘家了.
就跟她做女兒的時候一樣了.
就可以吃她父親的聖物了.
因為她孤苦無依了.
已經沒有丈夫了.
沒有家庭了.
更加沒有子女養.
所以你看到原來那些律法.

$^{241}$都是很有人道的.
就是她回到爸爸.
爸爸是祭司.
回到爸爸那裡.
她就可以吃聖物了.
不是說現在我初二.
我回到爸爸家.
爸爸剛剛獻祭完.
他就可以吃.
這樣不行的.
你是路過.
你是放假回來.
你就不准吃聖物.
弄得很清楚的.
你是很規定的.
你已經是回娘家了.
就不是跟丈夫吵架.
是你真的孤苦無依了.
你要回去那裡終老了.
陪爸爸陪媽媽了.
這樣你就可以吃聖物了.
因為那裡已經變成了.
記住.
聖物和食物是交替的.
祭司的食物.
即是說是獻祭完的那些下蘭.
神已經享用了.
是那些精華而已.
是那些肝.
是那些小小的燒在壇上.
其實整隻的祭身是歸祭司的.
所以祭司其實嚴格來說是很豐富的.
但是當然掛鉤的就是.
百姓是不是有心對神.
他對神沒有心的.
就不獻祭了.
不獻祭自然就有一個惡性循環.
大家可以想到那個後果.
所以這裡告訴我們.
一定是有規定那個女兒.

$^{281}$即是說在當時的.
其實當時沒有綜援的.
沒有社會救濟的.
女性更加不會出去打工的.
不會是女強人.
你是沒有了丈夫.
沒有了家庭.
又沒有了兒女.
其實就要回到父家.
回到自己父母的家裡.
來得到那個基本的維生的供應.
所以記住是已經成為了祭司的食物.
才可以吃.
就不是隨隨便便去吃聖物.
兩樣東西是不同的.
這個第十三節.
就可以吃他父親的食物.
只是愛人不可吃.
如果有人不吃了聖物.
要照聖物的元素.
加上五分之一交給祭司.
即是要有20\%的懲罰金.
祭司不可以褻瀆以色列人.
說到五分之一.
贖獻祭就是小的罪.
就要賠償20\%.
因為我們那時我自己.
那時剛剛信主聽到.
傳到人就教導.
你如果是偷窺拐騙.
做了一些不是很嚴重的.
即是入屋打劫爆竊.
而是有時有些事.
譬如你搭車.
沒有買足車票.
你就騙了地鐵公司.
那怎麼辦呢?.
有些弟兄姊妹就寫封信給地鐵.
接著就買一張車票.
就不搭.

$^{321}$就賠償自己計算.
我都有做.
我都有騙過.
就買了車票.
就銷毀它.
就當賠償給地鐵公司.
要加上五分之一.
20\%.
這些就是我們要了結對別人的虧欠.
這裡告訴我們.
你不吃了聖物.
聖物未變成食物.
又講究了.
你可以想像到.
連祭司自己的家庭的食物.
都是有規矩的.
何況是憲制呢.
所以你看到利未記.
其實是很清楚教導祭司.
他要對自己事奉所得到的供養系統.
是神設立的系統.
他要有這種警覺性.
他的兒女可以吃.
他的僕人是他親手買的.
他可以吃.
因為真是屬於這個家族的.
和這個家族有關聯的.
就可以吃了.
所以這個觀念是很牢固在祭司的家族裡面.
亦都是慢慢滲入以色列人當中的.
第十五節.
祭司不可以褻瀆以色列人所獻給耶和華的聖物.
大家記得十五義記上的第一章.
這個俄弗尼菲尼哈.
就是叉那些祭肉.
未獻的他都叉走它.
人家講明這些獻給神的.
獻完你才要.
你是祭司來的.
不行.

$^{361}$我現在就要.
他可能拿來燒.
拿來燒烤都不一定.
人家獻了給神.
可能是平安祭.
就是整個祭司是你的.
他都不行.
他都不能夠容忍.
要立即強搶.
其實就犯了第十五節.
褻瀆了耶和華.
獻給耶和華的祭物.
所以這個第十六節就說.
免得他們在吃聖物上自取罪孽.
因為我是叫他們成聖的耶和華.
然後留意下面這裡.
就關乎獻祭的祭僧.
十七節.
其實看利尼面記都是很好看的.
很多跟我們的生活都有關係.
其實是一個價值觀.
要調教我們對神的心態.
其實很重要的.
耶和華對摩西說.
李曉遇亞倫和他的子孫.
並以色列眾人說.
以色列家中的人.
或在以色列中寄居的.
凡是獻供物.
無論是所許的願.
許願即是還願.
是甘心獻的.
或者是獻給耶和華作凡祭的.
凡祭就沒有任何的原因.
我要奉獻給神.
凡祭.
那隻祭僧.
整隻是要燒的.
任何人都不准吃的.
祭司和普通人都不准吃的.

$^{401}$獻祭者都不准吃的.
凡祭唯一是歸祭司的.
就是凡祭生的皮.
歸祭司.
這是利未記第六章.
還是第五章有說的.
這些告訴我們.
這是很重要的.
第十九節.
要將沒有殘疾的公牛.
或者綿羊.
或者山羊.
這是價值的問題.
公牛貴一點.
綿羊做次要.
山羊就再粗糙一點.
就會便宜一點.
如此荒蒙月立.
強調的是沒有殘疾.
第二十節.
凡是有殘疾.
你們不可以獻上.
因為這不蒙月立.
絕對不是說你那張銀紙.
全抄了.
沒有換新鈔.
獻祭奉獻.
我想今天不是說這些的.
因為現在都電子化了.
支票.
Payme.
轉數快等等.
而是那個心態.
對神的敬重這個問題.
所以這裡就說.
要獻得蒙月立.
第廿一節.
凡是從牛群或者羊群當中.
將平安祭獻給耶和華.
為要還特許的願.

$^{441}$有特許的願.
有些人有特別的祈禱.
特別的還願.
是他和神之間的一個協議.
他很想達成這件事.
他就要還祭.
還祂許願的祭.
耶復他的女兒.
耶復他做事師.
他就回家.
誰第一個出來.
我就獻祭給他.
獻給神.
他的女兒出來迎接他.
他就慘了.
哇.
怎麼辦呢.
他就好像女兒所說過.
不要說獻祭.
後來他的女兒就說.
不行.
答應了就要做了.
聖經就告訴我們.
他的女兒終身不嫁.
終身歸神.
就好像索法為尼.
或者是終身事奉神.
這是耶復他的女兒.
因為耶復他自己很輕輕許了一個願.
這裡說有特許的願.
所以巴尼亞保羅傳道.
他也說原來他有一個拿西爾人的願.
這是《使徒行傳》的23章.
他上去耶路撒冷還願.
所以就算他已經成為使徒.
已經不再守舊約.
他是一個天生的猶太人.
他們也有這樣的心跡.
所以我們基督徒.
當然我們不應該去拜偶像.

$^{481}$去黃大仙.
去許願.
去扔那些寶牒.
我們當然不應該.
但是對神的一個許願.
自己和神之間.
有一個協議.
其實很多聖經是根據的.
所以我忘記了.
我是不是很早年的時候.
初信又什麼都不懂.
就有一個願.
好像說想做傳道.
那時候很羨慕傳道人.
很羨慕.
那時候真的不要笑我.
我初信其實二十多歲.
初信主.
開始那些衣著像傳道人.
白恤衫藍扯褲.
二十多歲.
有些比我年長的弟兄就說.
你是不是想做傳道.
那時候已經羨慕.
羨慕.
還有很簡樸.
你知道我那時候.
未信主.
我用LV的.
手袋.
大的袋是LV.
錢包是LV.
鎖匙包也是LV.
戴一些粗的手鏈.
金鏈.
因為那時候是喜歡到處去玩.
去玩弄.
後來突然信了耶穌.
真的穿白恤衫藍扯褲.
還有大陸公司的藍扯褲.

$^{521}$真的好像一個同傳道.
我隱隱約約記得.
可能是我初信的時候.
不知道.
就許了這樣的願.
所以大家要留心.
這件事.
新約也教我們的.
不是不准許願.
但不要隨便許願.
耶穌說的.
什麼意思.
就是說原來同神許願.
是神所寶貴的.
是寶物.
覺得這件事.
你很認真去繼續去做.
你不要說完就一言不發.
那就不好了.
耶穌的意思是.
不是不准許願.
但不要許願之後還願.
許願要還.
耶穌這樣說.
這裡就告訴我們.
原來有些是特許的願.
或者是甘心獻.
這個就真是寶貴了.
每個人都信耶穌.
每個人都跟隨神.
但有些人是特別甘心樂意為神做.
那些就沒得說了.
他甘心樂意為神做.
你就可以想像到.
神都是這樣來到.
非常之欣賞這樣的人.
即是甘心樂意.
羅馬書十二章就是這樣.
我們要做得甘心.
我們是一個活祭.

$^{561}$我們獻給神.
所獻的必定要純全無殘疾.
才能夠蒙月.
即是說就算你有這麼偉大的願望,願景.
但當你要還願的時候.
你都不可以隨便.
你自己來到想怎樣就怎樣.
都要根據祭司的規定,要求.
所以這個要記住.
那個主次的問題.
我們永遠是對神要有一個恭敬.
做事我們要為神.
不要跟神計較.
跟神計較就很大件事了.
這個是慢慢我們由信主.
要建立這一份態度,這一份價值觀.
第22節.
這裡再很清楚了.
核硬,折傷.
即是扭傷,殘廢.
有生扭的,生鱔的,長戒瘡的.
都不可以獻給神的.
也不可以在壇上作為火祭獻給耶和華.
燒了就沒有人知道.
不可以的.
我提過養豬.
一到冬天就很害怕.
因為那些豬瘟.
有時整個場地幾百隻,過千隻豬.
幾天之內就完全奄奄一息.
有些收這些豬的人.
他收去用來做叉燒.
自己私人屠宰.
因為那些他死與未死之間.
還是可以吃的.
他就拿了那些.
你知不知道他怎樣.
正常是多少錢一啖.
即是有市價的.
即是一隻豬大概二千元左右.

$^{601}$那時候.
現在我想起碼五,六千元.
其實很貴的.
但是你現在有求於他.
他開車來.
嘩!一整批.
二千元.
譬如有十隻.
是奄奄一息的.
很瘦的.
很大隻的.
大豬來的.
每隻收你五十元.
這樣的.
即是說幫你扔的意思.
但其實他拿回去燒.
來做那些下架的豬肉.
叉燒.
諸如此類的.
那些其實是不准的.
因為去到屠房.
去到上水屠房.
是不准毀滅的.
因為不可以出市場的.
因為你有病.
其實這裡也是這樣.
不是說神那麼高要求.
那麼挑剔.
其實都是燒而已.
但在神的眼中.
那個祭身是代表著那個人.
那個贖罪.
或者呈現他的心意.
所以有這些的暗病.
有這些殘疾.
就不可以作為火祭獻給耶和華.
第四十三節.
無論是公牛.
綿羊羔.
或者是肢體有餘.

$^{641}$或者是缺少.
即是多隻腳.
少隻腳趾.
都不可以的.
只可作今生祭獻上.
用以還願.
卻不蒙月立.
臣子損傷.
即是有臣傷的.
或者是壓碎的.
或者是破裂的.
或者是滲尿的.
我不知道滲尿是什麼意思.
不可獻給耶和華.
在你們的地上也不可這樣行.
這類的物你們從外人的手.
一樣都不可以接受.
作你們神的食物獻上.
因為這些都有損壞.
有殘疾.
不蒙月立.
所以整章是說.
哪些的祭身可以獻給神.
所以這很反映到.
那個民族那個時代.
對神的心意是怎樣.
如果我們整個時代.
整個教會.
譬如香港的教會.
我們大家都追求對神.
有一個很忠誠的態度.
整個香港都會復興.
因為每個基督徒都很認真對神.
可想而知.
即是說不是神不想有神蹟在香港發生.
是不想有大量的人信主.
而是根本連基督徒都很輕視那個憲制.
很輕視對神的那個心意.
所以就會出現.
待會我們看的《馬拉基書》第二章.

$^{681}$神對百姓的譴責.
那就是二十六節.
耶和華曉如摩西說.
財山的公牛.
或者是綿羊山羊.
很人道.
七日就要跟著媽媽.
其實很快長大.
一出生就會走路.
那些牛仔你看過那些紀錄片.
那些羊也是.
人就當然不可能.
滿月都還未會走路.
叫都還未會叫.
七天就要跟著媽媽.
第八天之後.
就可以當為共脈.
蒙月了.
即是第八天才來奉獻給神.
這些就是最初生的那個.
作為耶和華的火祭.
無論是母牛還是母羊.
不可以同日載母和子.
即是不要一次過.
殺他的媽媽.
殺他的孩子.
給耶和華.
要獻得可蒙月立.
所以做事是做在神的心意裡面.
可蒙神月立.
在這一章就出現很多次.
可蒙神月立.
蒙神月立呢?.
是不是神月立呢?.
不要和.
很多時候會和傳道人.
跟著計較.
做多做少.
回到教會.
這樣不滿意那樣不滿意.

$^{721}$沒問題的.
同日計較.
但不要和神計較.
是不是可蒙神月立.
這個就要大家自己.
去默想.
我們要可蒙月立.
跟著就告訴我們.
第30節.
要當天吃.
一點不可留到早晨.
我是耶和華.
這個是在說平安祭.
跟著你們要緊守.
遵行我的誡命.
我是耶和華.
你們不可以褻瀆我的聖名.
我在以色列人中.
卻要被尊為聖.
待會馬拉基斯.
又出現這些字眼.
我是叫你們成聖的耶和華.
把你們從埃及地.
領出來.
作你們的神.
我是耶和華.
我們本來就蒙拯救.
從埃及領出來.
奴隸來的.
本來是死的.
但是被拯救之後.
成為神的選民.
在這個.
西奈抗也.
成為一個民族.
成為屬神的兒女.
準備入迦南.
就可以承受地土.
所以現在這個時間還沒有入迦南.
在出埃及一個月.

$^{761}$頒布這些命令.
先設立了.
祭司的制度.
設立了.
祭司怎樣去獻祭的制度.
所以他們已經有一個.
很清楚的.
神的心意.
神是怎樣的呢.
祂期望我們要獻祭.
透過祭司的袋贖.
我們可以透過.
祭司的獻上.
我們的罪就蒙神的.
穴立了 赦免了.
我們的許願.
感恩 酬恩.
就可以透過獻祭.
神也都穴立了.
甚至乎可以神和人.
都共享這種快樂.
就是那些的祭司.
平安祭的祭司.
就可以.
百姓和.
祭司可以一起在神的.
面前吃的.
去到第三天就不可以再吃了.
或者可能當時沒有雪櫃.
防腐.
但是那個酬恩.
都不可以.
每天都帶著感恩.
這樣就不是那麼好了.
好像還計較著.
很多年前的感恩.
這樣就不是那麼好了.
就是要不斷的更新.
有新鮮的.
我們現在看了這兩個背景之後.

$^{801}$我們就回來.
看這個《馬拉基書》.
一看你就會明白.
神為什麼會這麼激氣.
透過先知.
來去責備當時.
被老歸回的這些猶太人.
為什麼他們.
闖下這麼大的禍呢.
神說.
要在萬國當中.
他的名要被尊為聖的.
第二章.
眾祭司.
就是凡子.
所有當時事奉.
公職的神職人員.
這個誡命.
是傳給你們的.
什麼誡命呢?萬君之耶和華說.
我們看《馬拉基書》的第二章.
就是.
馬太福音之前的那一章.
也是舊約.
第39卷的最後一卷.
不過記住猶太人的聖經.
就是希伯來文的聖經.
是三分法的.
只有24卷.
沒有39卷的.
因為有些是兩卷.
合成一卷的.
《馬拉基書》其實不是最後的一卷書.
最後的那一卷書.
其實是歷代志下.
要留意這些.
是希伯來文的聖經.
和我們今天的.
舊約是不同的.
因為我們今天的舊約的翻譯.

$^{841}$是.
馬索拉版本的翻譯.
也是根據.
這個七十事易本.
那個的.
那個的排版.
七十事易本.
是主前的大概兩世紀.
主耶穌都是讀七十事易本的.
希臘文就叫.
Septuagint.
就是七十的意思.
所以我們今天是用.
Septuagint的那個排版.
就是39卷.
但是希伯來文的聖經是24卷的.
大家留意一下.
當然你不知道也不要緊的.
因為沒有影響大家的.
不過你如果看一些.
福音書你就會覺得.
憤外傳神.
因為歷代志下.
還有一個王.
對不起還有一個.
耶和耶帶.
他的兒子.
就在聖殿被殺.
就是這個.
我忘記了是約爾.
還是約什麼.
我突然想不起那個王.
就是殺這個.
耶和耶帶.
因為他是他的舅父.
他的兒子就在.
聖殿被殺害.
這個就是流先知血.
第二章.
我們看回這裡.

$^{881}$萬君之耶和華說.
你們如果不聽從.
也不放在心上.
將榮耀歸於我的名.
我就使到就坐臨到你們.
使到你們的福分.
變為就坐.
事奉神的祭司.
是蒙福的.
聖經明顯這樣說的.
其實事奉本身已經蒙福了.
能夠事奉神.
是蒙福的.
就是祈禱你的兒子.
做傳道.
現在應該沒有人這樣祈禱了.
應該是做醫生還是做律師.
因你們不把.
誡命放在心上.
我已經就坐你們了.
誡命.
和祭司.
和福分.
其實三者是物不可分.
你必須有誡命.
有神的召命.
你就有.
專心的.
根據神的心意.
去服侍.
你那個時代的百姓.
這個就是福分的.
來源.
做不到就是就坐了.
第三節.
我必斥責你們的種子.
因為祭司.
都有自己分到的田.
他們可以耕種.
就斥責你們的種子.

$^{921}$有把你們犧牲的糞.
沒在你們的臉上.
這是很誇張的警告.
是一個羞恥.
你憲制.
憲制不是隨便.
把牛放上焚化爐.
那些叫焚化.
就不是憲制.
真正的憲制就是.
有手續的.
其實應該首先就是.
那個憲制者.
假如我們回到舊約的時代.
你和我一齊.
來到憲制.
大家帶著羊羔,牛.
可能家族比較大.
帶一隻公牛.
那.
死了.
人兵.
比較單薄.
帶一隻羊.
或者家境沒那麼好.
甚至帶一隻白鴿.
都可以的.
首先第一件事其實是.
祭司要測驗.
是不是沒有殘疾.
如果是通過了.
目測.
就可以準備憲制.
那個憲制者.
要按手在那個祭物上.
表明.
如果他的平安祭.
就是說這一份的祭物.
我是獻給神.
可能你.

$^{961}$每個月的奉獻.
你都按手在那個奉獻封上.
是我甘心樂意.
向神奉獻的.
你當然不用告訴別人.
你不要拿著說我是甘心樂意的.
那你就是沽名釣譽.
你應該在心靈裡面.
做了都不需要告訴別人.
獻祭給神.
但是如果是贖罪祭.
你按手.
那個祭身就代你去贖罪.
神月立.
那個手續就完結了.
等於寄信到郵局.
你放了信進去之後.
就沒有你的份了.
由郵局處理.
你想拿回也拿不到.
法例規定.
然後祭司領了那個祭身之後.
祭司就做他的手續.
你獻什麼祭呢.
好像下單.
原來是鳳爪排骨飯.
那個廚師.
就去拿那些鳳爪排骨.
放在鍋裡.
做煲仔飯.
那個祭司就知道原來阮妮.
是獻平安祭的.
他就做平安祭的手續.
首先切開之後.
他就將那些乾.
少許的精華.
和一些油脂.
燒在祭壇上.
然後整個平安祭的祭身.
就歸祭司和獻祭者一起享用.

$^{1001}$那個分類有多少.
分到多少.
似乎沒有說到多少.
可能很視乎.
那個獻祭者.
分多少給那個祭司.
這裡就告訴我們.
這裡就說.
這個誡命我是傳給你們的.
但現在我剛才還沒說完.
獻祭.
你要處理那些髒東西.
它也有些內臟.
骯髒的.
有些排泄物.
凡是都是這樣.
你要處理掉它.
所以有個地方.
是倒灰的.
就在回望.
我忘記了是南面還是哪裡.
在一個倒灰的地方.
所以新約.
那個吸輪溪.
有一個叫做.
一個好像.
倒垃圾的地方.
都是一個像污穢的地方.
所以.
他們處理.
其實挺厭惡的.
聖經沒有說得那麼白.
假設你知道.
其實你應該有時間.
去逃房.
如果讓你進去參觀就最好.
你就知道每天那些.
tong 豬,tong 牛.
其實很厭惡的.
很多工作要處理.

$^{1041}$所以舊約的憲制就是這樣.
所以我想會有專人.
去處理的.
處理這些排泄物.
倒到一些髒地方.
他們收集的地方.
然後再去憲制.
但現在神說.
要將你們犧牲了的.
那些糞.
抹在你們的臉上.
聖經說得那麼白.
那麼噁心.
即是當時的祭司.
視這些憲制.
已經成為了.
他們可以去操縱的工具.
完全沒有一個意識去榮耀神.
很偶然有時聽到一些長輩說.
以前的時代.
有時有些教會.
我不是太知道那一間.
也不應該說.
有時偶然.
已經沒有傳道人.
或者教會已經很亂.
很多的艾煞.
慢慢連教會的奉獻.
都不是那麼認真去處理.
很容易有時會帶來很多的誤會.
甚至真有其事.
自己分了它就算了.
你更加令到整件事.
越來越嚴重的.
那種惡性循環.
等於更加對教會沒有信任.
對牧者沒有信任.
對你們的領袖沒有信任.
就變得更加的一個.
差不多是教會都會瓦解.

$^{1081}$這些其實是最不好的.
所以你們看看這裡就是說.
當他們落到一個這樣的情況下.
神就說我要洗你們.
本來是榮耀的.
是福分來的.
變成咒詛 變成羞恥.
你想想你的臉上.
有這些濟生的糞.
其實是很羞恥的.
你們要與那些糞.
一同的除掉.
第四節.
你們就知道.
我傳這個誡命給你們.
使得我和利微所立的約.
可以常傳.
這是永約來的.
是什麼永約呢?.
利未人為什麼會這麼蒙福.
可以在十二支派裡面.
被神揀選呢?.
是有兩次的重大事件.
一次就是.
摩西上了.
西乃山領受約版的時候.
山下.
他的哥哥阿倫.
就被那些百姓逼迫.
就說你要做一個偶像.
我們來到去.
翻埃及.
因為在山下受不了了.
太困苦了.
阿倫就被強迫.
做了金牛獨.
他們就拜金牛獨.
摩西在山上.
當時神就跟摩西說.
你快點下山吧.

$^{1121}$百姓現在放肆.
犯罪離開我了.
摩西沒辦法了.
下山就真的看到那些百姓.
在那裡很開心的載歌載舞.
在那裡敬拜金牛獨.
摩西就很生氣了.
就把兩塊藥板打碎了.
就叫百姓喝了.
把那些水.
就變成了.
好像要懲罰他們一樣.
接著摩西.
這是出埃及的三十二章.
摩西三十二,三十三章.
摩西就吩咐.
誰想站在神這邊的.
就過來我這邊.
當時利美的子孫.
就走過來摩西那裡.
接著摩西吩咐.
在日落之前.
你見人就殺.
這是聖經記載的.
利美的子孫就.
顧惜弟兄.
站在神這邊.
來把那些拜偶像的百姓.
殺害.
當時神的憤怒就止息.
摩西就再上山.
再領受那個藥板.
另外一次.
就去到民數記的第二十三章.
後來去到二十五章.
因為那裡有巴蘭的事件.
當時有一個利美之派的人.
他叫做菲尼哈.
當時已經有瘟疫.
瘟疫是很嚴重地蔓延在百姓當中.

$^{1161}$百姓和領袖都在祈禱.
認罪.
求神寬恕.
不要再講瘟疫.
但當時的支派的一個族長.
他叫做森利.
和另外一個.
其實是摩雅的女子.
叫做哥斯比.
在營裡出入.
帶一個外邦女子進去.
菲尼哈就看到.
拿一支長的矛.
一刺就刺死他們兩個.
瘟疫就止息.
神就和阿倫和摩西說.
因為菲尼哈的事件.
菲尼哈這個妓邪的心.
神就和利美立約.
所以利美支派就有這個約.
他們可以散居在百姓當中.
成為祭司利未人.
可以教導百姓.
當然後來還有六座屠城.
約旦河西有三座.
約旦河東又有三座.
總共六座屠城.
祭司在屠城裡面成為審判官.
這就是神和利美納的約.
這就是第四節所說的.
所以是一個永存的約.
本來利未人是被咒詛的.
是他父親雅各咒詛利美.
因為他們當時殺了事件的家族.
本來是咒詛來的.
但神就透過這兩次事件.
使咒詛變成了祝福.
所以這裡有一個約.
所以利未人 亞倫的後人.
永遠要作祭司.

$^{1201}$馬拉基書第二章第五節.
我曾經與他立生命和平安的約.
就是這裡.
我將這兩樣賜給他.
使他全敬畏的心.
他就敬畏我.
懼怕我的名.
真實的律法在他口中.
他嘴裡沒有不義的話.
他以平安和正直與我同行.
使多人回頭離開罪孽.
就是有神的平安.
有神的真理 正直.
其實就是負起教導的責任.
教導真理就可以使人回轉.
離開罪孽 歸向神.
所以祭司利未人.
平時他不是每天都有很多祭獻.
也不是經常有大的節日.
在平時的日子.
其實利未人就負責.
在不同的地方去教導百姓.
所以這裡告訴我們.
真實的律法在他口中.
他嘴裡沒有不義的話.
他就能夠使多人回頭離開罪孽.
祭司的嘴裡當傳知識.
人也當由他口中尋求律法.
因為他是萬軍之耶和華的使者.
這個使者被差遣.
就是將神的道深入淺出.
講解 明白.
使百姓知道神的心意.
所以就能夠在祭司的口中.
可以尋求律法.
祭司就是要自己去研讀摩西的五經.
在不同的處境.
因為有些誤殺就要由祭司去判.
譬如在打獵.
大家一起去到野外.

$^{1241}$然後用斧頭去斬柴.
一扔起斧頭就打死了隔壁的人.
到處都沒有人見到.
那應該怎麼辦呢.
那個人就要走去屠城.
由屠城的祭司去做判斷.
諸如此類.
也是將仇恨減到最低.
進入屠城誤殺人.
其實那件事不是有預謀.
是誤會.
但因為沒有人見到 也沒有閉路電視.
來到去 祭司就判斷.
祭司就說你進入屠城就不能離開.
因為你一離開.
那個報仇的 譬如他的兒子.
見到你就會殺你.
所以那個人就被軟禁.
離開自己家.
就等於是監獄的前身.
所以屠城其實是監獄.
不過是一個開放式的監獄.
你如果隨便你自己走.
別人的仇恨還沒消.
找到你報仇.
你就不要抓誰.
是你自己的問題.
但什麼時候可以走.
那我就一輩子都不可以見到家人.
等到那年的祭司死了.
你就可以走了.
那不知道會不會祈禱.
等那個祭司快點走他就死了.
有這樣的可能.
應該不會的.
因為那個祭司是保住你的命.
任何人都不可以進去屠城.
追殺他的那些人.
所以是一個很公平的.
叫做很健全的一個制度.

$^{1281}$大家都講一個信字.
你在屠城就繼續你自己一個人的生活.
就離鄉別井.
雖然以色列很小.
但都是叫做離開你的家.
那就是說報仇的就大家一人走一步.
他都不想的.
不殺了你爸爸.
他都有應得的懲罰.
那你報仇的.
那個年日都很遠了.
那個仇恨都慢慢消去了.
在網上偶然間都會見到這些真實的個案.
有些真的.
年輕的時候大家都衝動的時候.
喝酒殺了自己的兒子.
那個媽媽呢.
一直要追究.
那個媽媽自己又抑鬱.
痛苦.
但是見到這個判決.
又不公平.
又不是置她於死地.
一直都不肯饒恕.
殺害自己的兒子.
其實是他的朋友.
大家喝醉酒.
可能爭一些東西.
在內地的.
就殺了.
最終慢慢慢慢的媽媽呢.
不知道什麼原因.
過了一段時間.
她後來想通了.
我這麼想她死.
其實我自己現在這麼痛苦.
就是我失去了兒子.
我沒理由見到對方的媽媽.
讓她兒子受到懲罰死了.
變成她又失去了兒子.

$^{1321}$張心比己.
她就明白.
終於她肯寬恕那個兒子.
主動去探望他.
接納和解.
就帶來她自己的內心.
又釋除了那些仇恨.
而對方的兒子和媽媽.
說千多謝萬多謝.
你可以想像到.
有時都需要有一種時間.
讓那個仇恨減除.
所以屠城原來有這麼重要的醫治作用.
進了屠城就不能走.
除非等大祭司死了.
你就可以走.
你就可以回家.
原來這樣是真的很好.
所以你會看到這裡.
祭司有一個功能.
就是他的嘴裡有知識.
他要教導真理.
他要做一個民間的訴訟.
他要做判斷.
即是他自己要很公正.
要站在中間.
不可以偏幫任何一方面.
社會事件的時候.
我都已經盡量不要說自己是什麼.
很多次.
利源崇拜.
已經收到WhatsApp.
會都還沒散 都還沒祝福.
剛才那個堂導說了一句.
關我什麼事.
不關我的事.
有時就帶著這些無….
其實可能聽的人.
是很賞心那句.
但可能未必是全件事.

$^{1361}$不要看得這麼嚴重.
有時每件事.
當然最好沒有說.
這樣是最好.
但有時很難避免.
所以我們都要做一下.
好像和事佬的訴訟.
民間始終有時.
不知不覺一定會有一些.
大家意見不合.
或者不同的角度看事物.
總會有一些偏見或者偏差.
需要祭司來做一個擺平.
所以祭司的嘴裡當傳之色.
所以其實真的要對聖經學.
多研讀 多認識.
多投入 多淺出.
能夠融會貫通.
多用多聽.
所謂終身學習.
祭司的嘴裡當傳之色.
人也當由他口中尋求律法.
這是專明來的.
很多人問你的意見.
問你一些重要的人生方向.
如果你能夠解通給他聽.
給他一些方向 一些原則.
當然絕對不應該指指點點.
不應該說你要這樣做.
你要嫁給誰.
當然不可以這樣.
給他一些原則.
讓他知道你這樣去祈禱.
方向性 原則性.
千萬不要變成你欽點.
你是操控他的心靈.
這些是不行的.
也不應該的.
你應該有那個知識.
人在你那裡尋求對神的心意.

$^{1401}$你給他一些方向.
你為他祈禱.
這是一個很大的專明.
這是很寶貴的.
今早我去屯原天區的傳道人的祈禱會.
有一位年輕的牧者.
在教會.
因為他的主任牧師.
突然間身體要退休.
其實也不提早.
他年紀大了 要退休.
他真的很年輕.
三十多歲 我想.
也不要緊.
那個管理層很看不起他.
為什麼找一個這麼年輕的做我們的堂主任.
很大壓力 很難做.
問我 吃飯的時候.
應該怎樣做一個堂主任呢.
一來我的頭髮白就沒那麼白.
他就問我.
我就講一些方向 原則.
我說你衣著盡量要沉實.
講道不要花巧.
雖然你做青少年的.
但是講道要四平八穩.
讓他們感覺到你是很忠厚等等.
一些方向而已.
不是教會做什麼.
不過人家會問.
都是一些經驗的分享.
讓他們知道.
原來這裡是一個很寶貴的.
對祭司的一個應許的肯定.
但是沒有了.
應該本來.
人家在你的心靈裡面.
尋找到神的真理.
因為你們是萬軍之耶和華的使者.
這個身份很重要的.

$^{1441}$我們要記得我們是奉差遣.
神的僕人.
真是教會最重要的是神的僕人.
要讓這些神的僕人.
在他的崗位盡忠.
讓他被釋放盡忠.
剛好今天又聽到一位傳教人分享.
其實他在那個舞會裡成長.
已經離開十幾二十年了.
但是那個教會幾年都找不到傳道人.
無端端找他回去.
他的教會就很奇怪.
因為幾年都沒有傳道人.
是那些信徒已經做了這些.
找童工,找傳道,港道各方面.
他們已經有一種制度.
就是那個傳道人每天.
每個小時做什麼都要寫清楚.
只是做文件都要搞一天.
每個小時你做什麼都要告訴領袖.
還要放在伺服器上.
他們隨時要看.
這樣對一個牧者的操控.
還是因為你比牧者還強呢?.
我沒有意思做一個批判.
我說你要慢慢建立他們對你的信任.
你剛剛去兩個月.
其實是一個自小在那裡成長的會友.
但是你很難想像那些領袖會這樣想事.
對一個牧者會這樣一種的.
是關心呢?.
還是你很怕他每個小時不忠心呢?.
他做每一樣事都要讓你知道呢?.
為什麼會這麼奇怪呢?.
我第一次聽這樣的想法.
你可以想像得到.
這裡告訴我們.
如果一個正常的祭司.
他是萬君之而和的使者.
使者就是差遣到百姓當中.

$^{1481}$教導律法.
讓百姓能夠明白神的律法書.
所以不過反過來.
來到這裡的時候第八節.
你們卻偏離了正道.
使到許多的人在律法上跌倒.
你們廢棄了我與你未所立的約.
這是萬君之而和說的.
說了這麼多.
本來這麼好.
這麼有尊榮的身份.
還有工作的性質.
天天有人給你錢.
研讀聖經.
研讀完之後.
你將他融會貫通.
去教導弟兄姊妹,百姓.
將神的心意告訴他們.
每個人都有不同的處境.
不同的困難.
每個人都不同.
你怎樣能夠將他們內心的困擾.
來幫他們解除.
還有將他們帶到神的面前.
他們自己去倚靠神.
這是很大的尊榮.
是對人的心靈的一個牧養.
和醫治的.
和同行,陪伴的.
但是你們偏離正道.
使到許多的人在律法上跌倒.
所以你可以想像到.
馬拉基書的年代.
神透過先知來對祭司的責備.
即是顧名思義.
當時不會有人研讀舊約.
百姓對於獻祭已經是可有可無.
掉以輕心.
你可以想像到一個教會.
一去到復活節,長假.

$^{1521}$一百人的就只剩下二十人.
你可以想像到.
那個教會一定是.
那個教導不強.
傳統人都不敢去.
長假你們就很累很累.
一窩蜂就去旅行.
去完回來又很累很累.
因為玩完了.
然後又很累很累地上班了.
如是者就是這樣了.
平常的星期就變成.
平常的星期就很多人了.
一百人了.
不過個個都好像不是太打緊.
因為你已經等待下一次.
復活節之後的長假了.
可能是暑假.
就跟子女又去哪裡哪裡了.
那就慢慢.
為什麼會掉以輕心呢?.
原來如果看回這張聖經.
是祭司先在律法上跌倒.
偏離了正道.
沒有將神的正確的心意.
去教導百姓.
使得百姓能夠.
知道神的.
哪一方面是神喜悅.
第九節.
我使得你們被眾人藐視.
看為下賤.
因你們不守我的道.
竟在律法上.
沾深情面.
就是說.
祭司其實他知道神的標準.
但他自己因應.
對百姓的.
那個顏面.

$^{1561}$他就放下標準.
來討好百姓.
讓我不要得失那個百姓.
這樣就一線之隔.
有時就是.
差之毫釐.
你可以想像得到.
你越是討好百姓.
百姓只會對你越藐視.
甚至這麼嚴重的.
看你為下賤.
如果我們.
會有看到.
你的一生.
你的事奉.
都是在他們的眼中.
是可有可無.
這個就已經是.
構成了藐視.
甚至構成一個下賤.
是很痛苦的事.
所以都盼望神給.
教會.
當時的時代.
才這樣提出.
很想祭司能夠有一個反思.
來有悔改.
在哪裡.
惡化了人的情面.
就要在那裡有一個悔改.
不要將一些.
本來是聖經的真理.
遞減.
自己單方面的遞減.
來惡化人的情面.
惡化人的情面.
是一些.
就是一些.
梁院長.
我們以前見到的院長.

$^{1601}$他說你們.
通常叫我們畢業班的同學.
不要去到一些教會.
成為一些苗族.
苗族就是我不管你那麼多了.
你來獻祭我就開心了.
有香油了.
你做什麼不關我的事.
只是打掃一下.
這些叫苗族.
不應該這樣做的.
忠於聖經你要教導.
你要責備就要責備.
這樣你才能夠換取到.
百姓對這個積分的尊重.
現在這個問題就是.
他們可能各種原因.
作者沒有講到具體的處境.
為什麼?是不是當時的執事會很強?.
強到一個地步.
祭司勉為其難.
要面對這些可能惡勢力.
某些家族.
不知什麼家族特別惡的.
對當時的祭司有一個很嚴重的操控.
這些你不准講.
這些你不准什麼獻祭.
要就要.
不要那麼執著.
可能慢慢慢慢就有這樣的情況.
導致那個祭司或者那個祭司的家族.
就出現了這種此消彼長的.
一種惡性循環.
所以最終.
祭司都得不到神的月立.
這個是最重要最嚴重.
百姓可能一時之間.
會有一些情況.
剛才講那個年輕的唐柱蔭.
三十多歲.

$^{1641}$他的執事會.
去到一個地步就是.
這位同工 即這位弟兄.
遞信辭職.
後來要驚動到去到區聯會.
區聯會介入.
那些執事都不肯罷休.
後來會眾挽留這個傳道人.
即這位今天分享的那位牧師.
大家不要猜哪一間了.
免得麻煩.
接著帶來的就是最終.
那幾個惡的執事.
如果他不走我們就不做了.
自己不再做執事.
自己連事奉都不做.
去到一個這樣的處境.
其實很可惜的.
即神的家 神的工作.
我們看為好像是自己的事業.
自己的勢力.
是很可惜的.
完全沒有了對神的敬畏.
對神的同在.
那種的尊崇.
落到這樣的人和人之間的競爭的田地.
那個同工說.
他當時那種情緒.
差不多如果再不走我就會瘋了.
即一個年輕的牧者.
要忍受這樣的情況.
是因為什麼原因呢.
因為他年輕 做全隊人.
你沒有只有他.
你自然他做acting就去承擔.
未成熟給他空間.
讓他錯.
在他的可控範圍容許他錯.
不是一種處心積慮的錯.
是做事未成熟的錯.

$^{1681}$那又為什麼不可能呢.
每個人都會有這樣的成長.
需要這樣成長.
怎麼可能一出來就是一個頂天立地的大牧者呢.
不可能的.
每個人都要有他的成長階段.
你可以想像得到.
往往就是充斥著這些.
但是回到《馬拉基書》的情況.
神透過先知.
其實有很重對祭司.
他們可能真的不知道什麼處境.
慢慢慢慢放下了自己那個打關的位置.
沒有了為神的聖殿把關.
沒有了為憲制把關.
剛才我們讀過的利未記是二十一章.
二十二章.
什麼的祭是神的月立.
什麼是可蒙月立.
這些是很重要的價值觀.
我們來到神的家.
你什麼的敬拜是神所喜悅的敬拜.
是有教導的.
新約也有教導的.
所以如果我們的態度都不好.
其實你不要說你影響其他人.
你自己是慢慢疏遠神的.
如果嚴重來說.
其實是神疏遠你.
這樣是很嚴重的.
慢慢你不要以為很過癮.
是你自己被神疏遠.
所以我們要記得.
我們來到神面前.
我們的憲制是不是可蒙月立.
還有我們有沒有我們自己的.
甘心獻的特許的願呢.
神是喜悅的.
所以祭司這裡沒有說具體的原因.
為什麼他們越來越失去了這一份.

$^{1721}$本來是神使者的尊榮.
什麼原因我們不知道.
我們只見到那個狀況.
就是他們減低了對百姓的要求.
百姓就更加對你藐視.
原來是一個惡性的循環的情況.
接著我們看第十節.
二章的十節.
開始講另一個課題.
其實全部有關連的.
第一章近尾是講那些.
就是那些又跛又腎的你都獻.
你獻給那些省長.
反而那些地方官.
你都不敢去.
我們很少人會去拿那些.
荷蘭水蓋MBE.
有的 有些有頒的.
譬如我們都有弟兄.
會有那些在社會上.
在一些有那些的.
就是頒那些獎這樣.
你會不會到時又和那個單位.
斤斤計較.
為什麼不行.
穿牛仔褲 球鞋為什麼不行.
為什麼不可以.
現在還講這些權威的嗎.
我想不會啦 個個都做套西裝很漂亮.
整個頭設好.
盡量穿最漂亮的去領獎.
第一章就講這些.
你獻給你的那些官長.
你都不會這樣做.
你居然來獻給神.
所以就是和神計較.
慢慢他自己都會藐視自己.
那個對神的奉獻.
這個是第.
一章.

$^{1761}$一章第八節.
你們將核硬獻為祭物.
這不為惡魔.
將瘸腿有病.
剛才裡面的二十二章已經很清楚禁止.
這不為惡魔.
你獻給你的省長.
當時可能都要收稅.
他豈會喜悅你呢.
豈能夠看你的情面呢.
這個是萬軍之耶和華講的.
實際上他們不是不知道那些東西.
是不可以獻給省長.
但現在我們獻給神.
所以這些是心態的問題.
如果我們一初信就已經建立好一個心態.
就是對準神.
對神有一個敬畏.
對神有一個愛和敬畏.
我們做很多事.
將心比己.
其實最好就是由自己的內心出發.
不用別人逼你.
不用別人教你.
自己的內心出發就不同了.
一切從心開始.
來到第十節.
我們豈不是都是一位父嗎.
豈不是一位神所做的嗎.
我們各人怎可以詭詐待弟兄.
背棄了神與我們列祖所立的約呢.
其實想說什麼呢.
猶太人行事詭詐.
並且在以色列和耶路撒冷中.
行了一件可憎的事.
因為猶太人褻瀆耶和華所喜愛的聖潔.
原來就是下面這裡.
取事奉愛邦神的女子為妻.
凡是行這些事情的.
無論何人.

$^{1801}$就是獻供物給萬軍之耶和華.
耶和華也必從雅各的帳篷中剪除她.
這就是背棄了與神的約.
聖潔的族類.
歸神為聖的選民.
神的選民.
神將以色列分別為聖.
進入迦南之後.
絕對不可以與迦南的女子通婚.
原因就是他們是拜偶像.
他們沒有離去對偶像的敬拜.
所以如果娶迦南女子.
不知不覺就會跟隨他們的風俗.
去拜偶像.
這是價值觀的問題.
當然有例外.
奴德就是摩訶女子.
還有就是.
約書亞記有一個拉俠的妓女.
是耶利哥的妓女.
嫁給了誰呢?.
那兩個探子是其中一個.
應該就是.
猶太之派的.
這個的.
沒有人說約書亞不會娶了耶利哥的妓女.
不會是他們的.
是另外那兩個年輕人的探子.
突然想起.
我先想想.
我今晚告訴大家.
應該是那個名字我突然想不起.
我知道是誰的.
告訴我們.
各人以鬼詐待弟兄.
背棄了神和我們列祖所立的約.
說的是人和人之間.
已經去到一種沒有了那種互信.
沒有了那種建立的信任.
這就是當時.

$^{1841}$因為拜偶像.
大家會傾向貪婪.
因為就是想透過偶像.
得到自己想要的東西.
有時不擇手段.
有時在民間會用一些手段.
現在很多網上騙案.
最近這兩天.
神樂園的同學.
見到一個短訊.
就說有些積分要換取一些東西.
叫他一按進去.
就說給他一條連結.
要付款.
一萬元.
馬上知道原來中計了.
入那些信用卡資料.
知道中計了.
馬上停止了.
不再入他給的那個碼.
接著安全起見.
馬上打去信用卡公司.
馬上將那張卡取消.
不要了.
換那張新卡.
因為你一按進去.
其實你已經去到那些騙案中心的地方.
所以現在整個社會已經沒有了那種信任.
聽電話你都不敢聽.
電話如果沒有來電顯示你都不敢聽.
三字頭的電話一般我都不聽.
已經沒有了這種互信.
但當時就更加嚴重了.
各人以鬼詐待弟兄.
我們都是神所做.
大家都應該是在神的裡面.
大家是神選民.
大家是歸神為聖.
有神的律法.
應該就是建立了這種.

$^{1881}$大家在神裡面的那種相愛.
彼此建立.
好像新約的.
無論是羅馬書.
或者是臨前的十三章的教導.
這樣的一個態度.
但這裡告訴我們.
當人沒有了對神的敬畏.
也沒有了真理的教導之後.
人與人之間就以鬼詐待弟兄.
就必然會出現這種.
人性沒有辦法被駕馭.
所以大家在教會.
應該大家都是愛神敬畏神.
對主耶穌的救恩.
不斷地被傳頌.
大家來學習.
所以每一個月的聖餐.
其實是很重要的一個教育的機會.
對一代一代的.
無論年青人.
年青人未必會那麼賞心.
因為他未有那麼多經歷.
對成年人來說.
他有很多人生經驗.
憤懷體會到救恩的美麗和偉大.
但這個教育是很重要的.
小時候就要帶他回教會.
小時候就要告訴他.
我們的主.
我們的神.
是多麼的可愛.
就算他未需要領聖餐.
其實都讓他坐在那裡.
看看領聖餐的哥哥姐姐.
叔伯是怎樣領聖餐.
這是很重要的教育.
讓他都看到.
你的人生將來就是這樣.
我們就是這樣.

$^{1921}$對神有一個恭敬.
這樣來學習.
這是很重要的一個.
心靈的一個分途.
在愛的環境分途.
在神的同在下的分途.
那個小朋友的成長.
他至少會敬畏神.
很重要.
很感恩我的兒子說.
他自小在教會長大.
從出世開始.
他現在知道自己與神的關係.
他是自己直接與神.
去做一個對話.
父母再不再存在.
在他的重大事件.
不重要了.
我可以是一個諮詢人.
可以是一個禱告.
但他自己與神的關係.
所以我們要盡快將小朋友.
帶到神面前.
但問題就是.
成年人自己都以鬼咋相待.
不是以神以靈相待.
這是很可惜的.
大家失去了一個互信.
變成背棄了神.
本來設立我們這個民族.
大家是神的選民.
這種的福樂平安.
所以猶太人的行事是鬼咋.
並且在以色列當中.
還發生了一件可憎的事.
人與人之間已經失去那種誠信.
是用鬼咋對待對方.
接著有些人再扭曲婚姻觀.
沒有那麼多猶太女子.
就娶外邦的女子為妻.

$^{1961}$玷污神的聖地.
所以去到第十二節.
凡是行者事的無論何人.
無論何人即是意味著.
不只是普通的平民百姓.
尤其是官長領袖.
都做這些事情.
即是失去了以身作則.
提摩太后書這樣說.
在大戶人家有金器銀器.
有木器有瓦器.
接著那裡說了事奉.
人如果自潔.
脫離了卑賤的事.
就可以作貴重的器皿.
成為聖潔合乎主用.
所以一個條件.
事奉就是要乾淨.
這裡告訴我們.
他們這些獻的貢物.
都不合神的心意.
就是獻貢物.
神都不會要.
還要在雅閣的帳篷中剪除.
即是說整個民族.
又要面對再一次的被擄.
即是說民間已經一層層的.
由沒有了對祭司的敬重.
祭司沒有了對神.
對身份這個使者的積分的尊重和捍衛.
他不去傳一個合乎神心意的律法.
也都不可以在他的口中.
問到律法書.
即是你問一個傳到人.
應該怎樣怎樣.
我不知道的.
你自己甚麼吧.
你去抽支籤吧.
這樣.
即是口中沒有了律法.

$^{2001}$神都不可以在他那裡尋求神的律法.
去到這樣的階段.
百姓自然就會用自己的心去做事.
用自己的心做事.
就是大家用詭詐對待對方.
即是說這個詭詐就變成商業的行為.
大家 我講好了.
我過兩天就交易了.
過兩天走了不給錢.
這樣毀約.
給了你一些貨物.
到時沒有回款.
很多這些民間的瓜葛的商業糾紛.
就是各人用詭詐待弟兄.
背棄了那個約.
所以要令到教會美麗.
就是要讓大家都在神裡面有那個幸福感.
大家都感受到是神而來的身份.
所以如果我們每個人都健健康康.
很豐富 很富足.
有時未必會珍惜神給我們的身份.
當然不可以說一竹篙打一船人.
有些人都很好 很敬虔.
但是大體來說.
如果我們不珍惜神給我們這個屬天的身份.
我們只會很著意地上的身份.
如果著意.
你就會用人自己的手法,手段.
所以很肯定就是你沒有了神的道.
你就一定會用詭詐待弟兄.
很容易就會用人間的.
世人都是這樣的 有什麼問題呢?.
自己的利益就是最大的優先.
所以這裡真的告訴我們.
其實互為關聯的.
首先是祭司沒有了道.
沒有了對神的道的捍衛.
其實這樣我們星期三晚的.
不是賣告白.
過百人這樣聽這麼悶的三一論.

$^{2041}$個個都是成年人.
其實教會正正就是需要有這些真理.
你越講真理.
頂尖會明白.
他自然就會敬畏神.
敬畏神他自己那個詭詐的心.
都會慢慢被調整.
個個人都詭詐的.
我年輕的時候見到誰都不順眼的.
那時候又大隻.
經常好勇鬥狠的.
看完李小龍的電影.
更加想打人的 出街.
個個人都是這樣的.
如果我們社會看電視劇多.
看多了就會潛移默化.
又主妻又謀殺.
慢慢會在我們的腦袋裡萌芽的.
所以其實不要看那麼多那些東西.
只會煽動我們裡面的那種自我出來.
但是教會現在就要回頭了.
才是說要回到那個尊榮.
祭司本來是有神立約的尊榮的.
你是神的使者來的.
你好好的捍衛那個律法書的真理.
百姓就會敬畏神.
他就自然民間就會大家用愛去相待.
講個信字.
真是信的了.
那就互信了.
整個社會就會和諧.
現在就剛剛反過來.
是十級以下的那個靈性的破產.
當時的以色列人.
所以盼望我們都能夠在這個社會.
我們都發揮我們的見證.
還有我們的作用.
今晚我們就講到這裡.
我們過了年了.
下個星期.

$^{2081}$我不知道下個星期三還是初幾.
下個星期四.
對不起 下個星期四.
那現在是初六了.
應該就沒什麼特別的東西.
到時我們再看看.
如果有特別的東西.
我就再在群組裡講.
一起禱告.
感恩天父讓我們能夠.
透過第一章的結尾.
第二章的馬拉基書.
讓我們看到一個結構性的問題.
很多東西都不是個別那麼簡單.
原來是有牽連.
就是那個對神的敬畏.
出了一個問題的時候.
慢慢就會受到影響.
令到整個民族都會傾向於一種世俗.
傾向可能和外邦人.
沒有分別的那種價值觀.
和民間的那個人和人的交往.
就只會用很多的手段.
或者用一個猜疑.
帶來的一種甚至乎鬼詐.
互相的對待.
我們很盼望主讓我們能夠.
在教會裡.
我們不是一種無知的對別人的接納.
而是有意識地.
用神的真理去彼此的相待.
彼此在神的裡面.
來建立這份相交和互信.
和大家都是傾向於對神的真理的學習.
願主幫助帶領我們.
我們過幾天就會有過年了.
保守我們可以見到很多家人的時候.
都表達出我們這一份.
信仰帶給我們的生命改變.
願神施恩.

$^{2121}$這樣感恩禱告.
奉耶穌基督的名.
阿們.
\newpage



\section{創世記 26:1-23}
\label{sec:OWpYTLtxpsA}
\textbf{2024-02-03 晚堂《豐盛生命:依靠順服》創世記 26\_1-23 講員:陳麗蟬牧師}
\newline
\newline
連結: \href{https://youtube.com/watch?v=OWpYTLtxpsA}{\texttt{https://youtube.com/watch?v=OWpYTLtxpsA}} ~~~~ 語音日期: 2024-02-14
\newline
\newline
\hyperref[sec:ny6z0RWTL8g]{\small{< < < PREV SERMON < < <}}
~
\hyperref[sec:index_chronic]{\small{[返順時目]}}
~
\hyperref[sec:index_scriptual]{\small{[返順卷目]}}
~
\hyperref[sec:45mO3Udlc8E]{\small{> > > NEXT SERMON > > >}}
\newline
\newline
創世記 26:1-23
\newline
\begin{longtable}{cl}
\hline
\hline
章節 & 經文 (和合本修訂版)\\
\hline
26:1 & \begin{tabularx}{0.7\textwidth}{X} 那地有了饑荒，不是亞伯拉罕的時候曾有過的那次饑荒，以撒就到基拉耳，非利士人的王亞比米勒那裡去。 \end{tabularx} \\ \\ \relax
26:2 & \begin{tabularx}{0.7\textwidth}{X} 耶和華向以撒顯現，說：「你不要下埃及去，要住在我所指示你的地。 \end{tabularx} \\ \\ \relax
26:3 & \begin{tabularx}{0.7\textwidth}{X} 你要寄居在這地，我必與你同在，賜福給你，因為我要將這一切的地都賜給你和你的後裔。我必堅定我向你父親亞伯拉罕所起的誓。 \end{tabularx} \\ \\ \relax
26:4 & \begin{tabularx}{0.7\textwidth}{X} 我要使你的後裔增多，好像天上的星，又要將這一切的地賜給你的後裔，並且地上的萬國都必因你的後裔得福， \end{tabularx} \\ \\ \relax
26:5 & \begin{tabularx}{0.7\textwidth}{X} 因為亞伯拉罕聽從我的話，遵守我的吩咐、誡令、律例和教導。」 \end{tabularx} \\ \\ \relax
26:6 & \begin{tabularx}{0.7\textwidth}{X} 於是，以撒住在基拉耳。 \end{tabularx} \\ \\ \relax
26:7 & \begin{tabularx}{0.7\textwidth}{X} 那地方的人問起他的妻子，他就說：「她是我的妹妹。」原來他害怕說「我的妻子」。他想：「或許這地方的人會因利百加殺我，因為她容貌美麗。」 \end{tabularx} \\ \\ \relax
26:8 & \begin{tabularx}{0.7\textwidth}{X} 他在那裡住了一段很長的日子。有一天，非利士人的王亞比米勒從窗戶往外觀看，看哪，以撒在撫愛他的妻子利百加。 \end{tabularx} \\ \\ \relax
26:9 & \begin{tabularx}{0.7\textwidth}{X} 亞比米勒召以撒來，說：「看哪，她實在是你的妻子，你怎麼說『她是我的妹妹』呢？」以撒對他說：「因為我想，恐怕我會因她而死。」 \end{tabularx} \\ \\ \relax
26:10 & \begin{tabularx}{0.7\textwidth}{X} 亞比米勒說：「你向我們做的是甚麼事呢？百姓中有一個人幾乎要和你的妻子同寢，你就把我們陷在罪中了。」 \end{tabularx} \\ \\ \relax
26:11 & \begin{tabularx}{0.7\textwidth}{X} 於是亞比米勒命令眾百姓說：「凡侵犯這個人，或他妻子的，必要把他處死。」 \end{tabularx} \\ \\ \relax
26:12 & \begin{tabularx}{0.7\textwidth}{X} 以撒在那地耕種，那一年有百倍的收成。耶和華賜福給他， \end{tabularx} \\ \\ \relax
26:13 & \begin{tabularx}{0.7\textwidth}{X} 他就發達，日漸昌盛，成了大富翁。 \end{tabularx} \\ \\ \relax
26:14 & \begin{tabularx}{0.7\textwidth}{X} 他有羊群牛群，又有許多僕人，非利士人就嫉妒他。 \end{tabularx} \\ \\ \relax
26:15 & \begin{tabularx}{0.7\textwidth}{X} 他父親亞伯拉罕在世的時候，他父親的僕人所挖的井，非利士人全都塞住，填滿了土。 \end{tabularx} \\ \\ \relax
26:16 & \begin{tabularx}{0.7\textwidth}{X} 亞比米勒對以撒說：「你離開我們去吧，因為你比我們強盛得多。」 \end{tabularx} \\ \\ \relax
26:17 & \begin{tabularx}{0.7\textwidth}{X} 以撒就離開那裡，在基拉耳谷支搭帳棚，住在那裡。 \end{tabularx} \\ \\ \relax
26:18 & \begin{tabularx}{0.7\textwidth}{X} 他父親亞伯拉罕在世的時候所挖的水井，在亞伯拉罕死後，都被非利士人塞住了，以撒就重新把井挖出來，仍照他父親所取的名為它們命名。 \end{tabularx} \\ \\ \relax
26:19 & \begin{tabularx}{0.7\textwidth}{X} 以撒的僕人在谷中挖井，就在那裡得了一口活水井。 \end{tabularx} \\ \\ \relax
26:20 & \begin{tabularx}{0.7\textwidth}{X} 基拉耳的牧人與以撒的牧人相爭，說：「這水是我們的。」以撒就給那井起名叫埃色，因為他們和他相爭。 \end{tabularx} \\ \\ \relax
26:21 & \begin{tabularx}{0.7\textwidth}{X} 以撒的僕人又挖了一口井，他們又為這井相爭，以撒就給這井起名叫西提拿。 \end{tabularx} \\ \\ \relax
26:22 & \begin{tabularx}{0.7\textwidth}{X} 以撒離開那裡，又挖了一口井，他們不再為這井相爭了，他就給那井起名叫利河伯。他說：「耶和華現在給我們寬闊之地，我們必在這地興旺。」 \end{tabularx} \\ \\ \relax
26:23 & \begin{tabularx}{0.7\textwidth}{X} 以撒從那裡上別是巴去。 \end{tabularx} \\ \\
[1ex]
\hline
\hline
\end{longtable}
$^{1}$頂智慧晚安.
現在要大家去想一想.
在《創世記》裡面.
哪些人物你是最熟悉的.
你會想到誰呢?.
很快就想到.
亞伯拉罕.
雅各.
如果不是剛才讀以撒的時候.
我想都沒有什麼人記得他.
今天我們就要說.
以撒.
以撒他在.
在《創世記》裡面.
他都是叫做.
猶太人裡面的一個祖先.
很多時候聖經都會說.
耶和華是.
以色列的神.
是以撒的神.
是雅各的神.
以撒就剛剛在正中間.
亞伯拉罕就是.
以撒的爸爸.
然後呢.
雅各就是以撒的兒子.
三代同堂.
在正中間.
但我們往往就是.
只會認識.
亞伯拉罕.
亞伯拉罕裡面.
聖經裡面記載很多很多的事跡.
還有在《創世記》裡面.
原來由十二章.
一直到二十五章.
就開始記載亞伯拉罕的事情.
總共十章.
而以撒的兒子.
就是雅各.

$^{41}$就更厲害.
他由二十七章.
一直到四十九章.
都是關於雅各的事跡.
所以就是說.
就是以撒在全國的《創世記》裡面.
他只是記載了一章的事情.
是由他們去做主角.
但今天我們就用這段時間.
我們一起看完這章裡面所說的事情.
以撒這個名字就叫做.
笑.
他的意思就是笑.
很開心地笑.
那為什麼呢.
他的名字這麼開心呢.
就因為他能夠出世.
也是一個神蹟.
那他亞伯拉罕.
多少歲生他的呢.
是一百歲的時候.
那他媽媽已經是九十歲了.
所以當他們求了很久.
耶和華先賜了這個兒子給他的時候.
他出世的時候.
笑得合不上嘴.
所以他的名字就叫做以撒.
是開心到不得了的意思.
那在他出世的時候.
就被他們整個民族.
或者是之前或者之後看到的時候.
上帝真是很厲害.
在這個人已經不能夠生育的年齡的時候.
但是竟然都能夠生出這個以撒.
他們就看到了.
在上帝來講的時候.
凡事都能.
也都令到以後那些民族.
藉著這個以撒的出世.
他們看到了.

$^{81}$只要去堅信上帝的應許.
以及倚靠著他的能力的時候.
上帝可以使很困難很困難的東西.
或者使到不可能會有的東西.
他都可以變為有的.
所以有人就說.
這個以撒的生命.
他的特徵就是依靠信佛.
依靠和信佛.
我們來看看.
在這個以撒的裡面.
他整個人生給我們看到.
他很信靠上帝.
也都很信他自己的爸爸.
很聽話.
在《創世記》22章的時候.
那時候他還沒有做主角.
他的爸爸亞伯拉罕.
就帶他一起上山去獻祭.
在那裡走的時候就揹著柴和火.
在那裡走的時候他就覺得很奇怪.
為什麼爸爸沒有帶獻祭羊羔呢?.
他就走去問爸爸.
爸爸就只回答了他一句.
他說:耶和華神必有預備.
去到這個祭壇要獻祭的時候.
他爸爸就竟然把他綁在祭壇上.
當時那些學者計算過.
以撒應該都二十多三十歲了.
他都很厲害.
他都可以揹著柴揹著火.
去上山的時候.
力氣都很大.
但是為什麼他仍然很信服.
又不反抗呢?.
總之爸爸做什麼呢?.
他就聽了他的話.
在這個古代的社會裡面.
他很尊重他爸爸.
以至很尊重他很疼愛他的爸爸.

$^{121}$他不反抗.
他信得過.
他爸爸所信的耶和華.
一定有他的方法.
就算他爸爸所信的上帝耶和華.
要他爸爸來將他作為祭品.
他都要來信服的.
我想這是一個需要很大的信靠的心.
就算我們到今天.
有時候我們一有了什麼很大的病.
突然間知道自己有癌症的時候.
我們通常第一個的反應是什麼.
為什麼會是我呢?.
通常弟兄姊妹這樣跟我說的時候.
我不會立時跟她說的.
我心裡面跟她說.
為什麼不可以是你呢?.
在每一樣事情發生的時候.
就是上帝看我們.
看我們究竟那個心是怎樣的.
所以其實上帝那次要亞伯拉罕這樣獻祭.
他目的真的不是想要以殺生命.
他那次的獻祭.
其實他是要亞伯拉罕.
學會將他的兒子交給上帝.
由上帝來帶領一切.
有時候我們在生命裡面.
上帝會拿走我們一些很寶貴的東西.
有時候拿走我們的工作.
拿走我們的金錢.
拿走我們很愛的人.
但是上帝就是要看我們.
在我們沒有了這些很珍貴很寶貴的東西的時候.
我們愛這些東西多一點.
還是愛上帝多一點呢?.
除了這個被當作祭品獻祭.
他很聽話之外.
在婚姻的事情他都很聽話.
現在我們不是的.
現在我們作主的.

$^{161}$怎麼可以由父母 由上帝作主呢?.
所以在婚姻裡面.
他知道他爸爸就差派了那個老僕人.
回到他的鄉下梧意那裡.
說是命準帶外邦女子回來的.
所以那個老僕人就回到梧意那裡.
幫他找了一個女孩 李百嘉回來的時候.
當他在田間裡面默想.
默想著上帝會怎樣呢?.
他預備甚麼的時候呢?.
他就見到這個老僕人帶著駱駝.
帶著他的太太李百嘉回來.
很乖 很乖.
完全沒有抗議的.
所以一回來的時候.
他就帶了進去他媽媽的帳幕裡面.
他就跟她成為了夫婦了.
他也是很乖.
他信得過上帝為他安排的.
是最好的.
所以就不需要問那麼多東西.
也不需要反抗的.
接著呢.
剛才盧醫生幫我們讀了一段經文.
就是這個去做主角.
整章都是他做主角.
當他住在那個地方的時候.
就發生了這個饑荒.
他就要逃難了.
他就去逃難.
他原本住在別士巴這個地方的.
但是因為發生了饑荒.
他就要去到基拉爾.
上一點點的地方那裡.
接著去完基拉爾之後.
他就順勢就會去到埃及地.
這個基拉爾其實都是面對著這個饑荒的.
還有那裡住了很多非利士.
很多外族人的.
那他為什麼要去到埃及呢?.

$^{201}$因為在埃及裡面有尼羅河.
所以就算天氣下雨究竟多還是少.
那個影響都不會太大.
那個尼羅河就給他們在埃及裡面有足夠的水源.
所以當時這個埃及可以說是一個世界的糧倉.
凡是哪裡有饑荒就會走去那裡.
所以這個以撒都計劃去那裡逃避這個饑荒的.
但是就在這個基拉爾那裡.
上帝就向他顯現了.
他說你不可以去埃及.
你要繼續住在這個非利士人的地方.
繼續住在這個基拉爾.
他說他會賜福給他的.
他會賜福給他的.
那麼明明這個地方這麼艱難.
為什麼仍然要吩咐我留在這裡的地方呢?.
還有這個地方又是非利士人的地方.
很危險的.
他不明白.
但恩爵信.
也都很清楚知道這是上帝給他的聲音的時候.
他信服.
我想一個真正信服上帝旨意的人.
他是憑著這個信心.
不憑著這個眼見.
然後才能夠去經歷到.
真是那份上帝的安排.
我忘記了晚堂有沒有講過.
在我們人生不同的時間裡面.
你化導有時會遇到很多的難處出現.
我想現在就是經濟低迷.
很多工作做生意的人都很辛苦.
那時候我在讀完教育學院出來之後.
就很難找工作.
那些人就寫了很多很多信出去找學校.
但當時我堅信一件事.
我祈禱過.
然後也都看過那些學校.
是不是我自己合我這個科目的時候.
我才寫信的.

$^{241}$也都因著我自己信耶穌的時候.
我是讀中三的時候.
中三,中四的時候信耶穌.
所以我就跟上帝講.
第一我要教中學.
第二我還要教中三.
那些科目都是適合我的科目.
其實那時候真的很艱難.
那些同學發了二百多封信出去.
我的好同學住灣仔.
去了大澳教書.
他都胖胖的 幾塊石頭.
他就要走去教土風舞.
總之學校叫他做些什麼他都要做.
又有另外一個傳奇裡面很老的人.
他根本不會跳舞 又要去教中國舞.
暑假就拉他去學中國舞.
他學了一個暑假就去教學生.
當時我堅持.
所以我那次寄出了十九封信.
當然沒有人理睬我.
因為這是我跟上帝的堅持.
你亂發信出去是沒有用的.
他都求上帝給我.
每個人都九月一號開學的時候.
我都要九月一號開學.
上帝是你叫我進教育學院讀書的.
我最初都不是很想的.
上帝真是很好的.
他在八月二十九號有學校打電話給我.
他需要請一個代課老師.
因為他生了嬰兒.
所以前四後六總共多久?.
前四後六兩個半月.
他在一間學校裡面代課了兩個半月.
我覺得上帝就是這樣.
你只要信服.
你跟他領受什麼.
而他敢求又敢走的時候.
上帝就出人意表的平安.

$^{281}$和出人意表的情況出現.
代完兩個半月之後又怎樣呢?.
上帝又很奇妙.
就有一間中學請我.
他說有個人有個弟兄.
他原本正在申請政府工.
但是一直都沒有消息.
沒辦法之後他就不喜歡教中學.
於是都要去教了.
但是收到政府收他之後.
他就立即走人.
然後我剛剛代完那間學校之後.
我就進了那間學校教書.
是教中三.
我教七班的中國歷史.
我自己很喜歡教的科目.
也是我自己跟神說.
我在那段時間信耶穌的.
我都很希望教年紀的學生去信耶穌.
我們就看到.
當時我自己都覺得很麻煩.
我覺得上帝是你叫我進教育學院的.
到現在我又去學習去徵求.
就不亂發信.
又覺得這樣去求的時候.
為什麼人人慢慢陸續有人請.
但我卻沒有人請呢?.
其實這個就是一個信服和信心的功課.
保羅在哥林多後書裡面說.
我們行事是憑著信心.
不憑著硬件.
在最初的時候.
這些同學都覺得很慘的.
你找不到工作.
但上帝給我9月1號開學.
但我是最好的一個同學.
兩個半月之後.
我就在那間中學教書.
是我自己喜歡的科目.
也是我自己很想教的年紀.

$^{321}$有兩個小朋友在那裡聊天.
他就說.
我爸爸媽媽真的不了解我.
很不明白我那些心情.
另一個小朋友問.
他為什麼不了解你不明白你.
他說每一天晚上.
我想玩久一點.
不想睡覺的時候.
他總是叫我快點去睡覺.
到第二天早上的時候.
我想睡久一點的時候.
他總是來叫醒我.
叫我快點起來.
他都不明白我的心情.
他問我喜歡做什麼.
我喜歡晚上玩.
喜歡晚上不睡覺.
我喜歡早上睡久一點.
為什麼不行呢.
有時候我們就像這個小朋友一樣.
我們不明白上帝的安排.
在過程當中.
我們會生上帝的氣.
我們會埋怨上帝.
但是上帝看得最清楚.
祂知道哪一條路.
哪一個方法.
對我們每一個人的新生命.
對我們的生命是最好.
最有幫助的.
所以弟兄姊妹.
上帝知道一切.
我們只要的就是這個信服.
雖然這個以撒是一個很聽話的人.
很聽上帝的話.
很聽他爸爸的話的人.
但當他面對危機的時候.
一樣有軟弱.
我們剛才就看到了.

$^{361}$當他上帝要他留在.
這個非利士人的地方.
這個基拉爾.
他就害怕了.
我死了.
我的太太這麼漂亮的時候.
用什麼呢.
這幫非利士外邦人.
他爭著要搶我的太太的時候.
他會殺死我的.
所以他就告訴那些非利士人.
他只是我的妹妹.
但後來呢.
有一次他不知道為什麼.
他和他的太太.
大家在玩.
在做夫妻之禮的時候.
不知道是關窗還是怎麼樣.
就被這個非利士的王.
阿比米勒看到.
他就罵了.
這個以撒.
你知不知道你差點害死了我們.
你因著大衛.
不是 因著以撒說這個謊話.
差點害死了這個非利士.
不知道是誰.
原來這個阿比米勒.
都很正直的.
他一知道.
這個尼伯加是以撒的太太的時候.
他就下令.
不准傷害以撒.
也不准傷害他的太太.
這個謊話.
有時候會造成.
給人很大的傷害.
但呢.
你又覺不覺得.
說謊話好像是天性來的 是不是.

$^{401}$我想問一下大家.
有沒有人沒說過謊話.
誰夠膽說自己沒說過謊話.
又想一下.
你在你自己的思想印象當中.
你第一次說謊話的時候是甚麼時候呢.
我記得我自己.
我讀三年級的時候.
我講過一次謊話.
令到我媽媽很擔心.
因為我們那時候做大排檔.
就日日早上.
一早就出去.
剝蝦 挖魚.
就沒得出街玩.
所以我們學校.
那時候我讀聖雅閣的.
隔壁那個就是聖雅閣福群會.
教很多小朋友的東西.
他就在中秋節之前.
就搞了一個中秋賞月晚會.
是上寶雲道那裡的.
我就很想去.
還帶著我弟弟去.
但我知道我媽一定不讓我去.
我就告訴我媽.
我們在福群會那裡的.
有故事聽.
又有東西吃.
又有燈籠玩.
去吧去吧.
七點半去到九點半就炸了.
那一哀了很久.
那我媽就讓我們去.
其實是上寶雲道.
那就想著回家沒事了.
誰知道那天上了山之後.
遲了下山.
所以下到回來的時候.
去到聖雅閣福群會解散的時候.

$^{441}$已經是十點了.
一出了門口.
我們就見到我媽.
是那裡的斜坡.
在那裡很害怕地在那裡哭.
她說嚇死我了.
你們又說九點半下來.
我就想著更晚走來接你們回家.
誰知道來到什麼燈都沒有.
又不知道你們去了哪裡.
我害怕到哭.
我當時很內疚.
很內疚.
就是因為我的謊話.
令到媽媽更擔心.
有時候我們就是這樣的.
有時候這些謊話.
原來我們從小就懂得去說.
在說這些謊話的時候.
我們時時都會傷害到人.
所以求神幫助我們.
雖然真的很容易的.
是不是?不用想.
有時候都會說出來.
但求神憐憫我們.
保守我們的嘴唇.
免得你傷害其他人.
然後就有以撒了.
除了他在這麼聽話的人生裡面.
他說謊話差點害死人.
還有他那個很隱約的性格.
很隱約.
是生命的特質.
當時基拉爾生活得這麼艱苦.
周圍的人就沒有什麼修成.
但也就是說.
因著他的信服.
因著他的聽話.
被他所種的農作物.
有百倍的修成.

$^{481}$又有很多的羊群.
牛群.
當人這麼豐收.
這麼好的時候.
通常就會帶來人妒忌.
所以那些非利士人就趕他走了.
他們說你走吧.
因為你們太好了.
然後他們就去遠了一點點的地方.
他們都要生存.
他們就要挖一些井.
其實不是挖新的井.
這些井都是他爸爸亞伯拉罕.
之前挖了的.
不過又是被非利士的人堵住了.
但他挖一個.
那些非利士人又來搶了.
他就不是說這個井是非利士人.
不是說他們自己.
他說裡面的水是你的.
那你就死了.
你的井你可以挖.
但你不可以拿裡面的水.
逼得這個以撒他們就說.
又一路一路這樣走了.
開一個井又是這樣.
第二個井又是這樣.
第三個井又是這樣.
這個井對當時的遊牧民族來說.
是非常重要的.
有井就有水.
沒有水你就種植不了.
自己的命都沒有.
在現在的裡面.
如果你家裡有時製水.
或者洗水缸.
說沒有水的時候.
你就非常麻煩了.
但是看到以撒的那個信服.
和謙讓同人爭拗.

$^{521}$所以他寧願吃虧.
開了這個井.
你說裡面的水是你的.
不給.
於是他走去.
再開第二個.
到第三個的時候.
那些人都沒有再來了.
在這個大饑荒的時候.
周圍的人都種不了東西.
但他卻種每一樣東西.
是百倍收成.
很厲害的.
還有那些牛羊.
很多很多.
開每一個井都很順利地開始.
聖經剛才記載告訴我.
他成為了大富翁.
他因著這些經歷.
那些菲律賓人.
他們覺得這個人很厲害.
所以之前他們趕他走.
趕這個以撒他們走.
說要離開他們.
不要在這裡.
但是他後來.
發覺自己不妥當的時候.
他就走去找回以撒.
他就去求和.
簽和約.
給我們著.
說你之前趕我們走.
現在就來和我求和.
這個和好.
我當然不願意.
但是以撒呢.
以撒就跟他簽了和約.
以致原本是敵人的.
化敵為友.
也因著他這個忍耐謙和.

$^{561}$使得上帝的名得榮耀.
因為在26章第28節.
菲律賓人他們說.
我們從你的身上.
明明看見耶和華與你同在.
我們從你的經歷生命裡面.
我們看到耶和華神與你同在.
弟兄姊妹.
有沒有經歷過.
當我們信了耶穌.
或者到現在有人跟你說.
你信了耶穌之後.
你真是很不同.
或者你信了耶穌之後.
真是與其他人有點不同.
但這個以撒也是.
在之前.
你看見他每一件事情.
聖經裡面26章裡面.
沒有說他有任何的說話.
但到後來.
當這班菲律賓人來和他求和的時候.
他竟然開始說話.
他說什麼.
他說你們之前又走來.
說很憎恨我們.
又要趕我們走.
現在做什麼.
又來找我們.
要和我們一起來合作.
我相信他說話都大聲了.
有力了.
他真是經歷到上帝去祝福他.
所以我們看到.
他那個謙和.
他那個忍讓.
不是他武力走去打菲律賓人.
也不是因著他的內弱.
而是在耶和華裡面.
他有一個廣闊的胸襟.

$^{601}$你拿了我的井不讓我用水.
跟你吵架沒有用的.
這樣來打架.
於是他逐個逐個這樣.
到第三個.
那些菲律賓人沒有再去搞他們了.
腓立比書本來都教導我們.
當叫眾人知道你們謙讓的心.
在這個羅馬書裡面.
本來都教我們.
這個謙讓的行動是要怎樣呢.
我們一起讀了這些黑色的字.
好 一二.
有時黑色有時紫色.
好 那我們一起讀.
一二三.
不要以惡報惡.
眾人以為美的事.
要留心去作.
若是能行.
總要盡力與眾人和睦.
親愛的弟兄.
不要自己伸冤.
寧可讓步.
傾憑主路.
因為經上記著.
主說.
伸冤在我.
我必報應.
所以.
你的仇敵若餓了.
就給他吃.
若餓了就給他喝.
因為你這樣行.
就是把炭火堆在他的頭上.
為惡所成.
反要以善成惡.
所以.
凡願意你去信服上帝.
並且在與人相處的過程中.

$^{641}$以謙和來與他們對待的時候.
上帝的恩福就會大大的臨到.
願意我們都學校.
以殺他的生命.
我們一起來禱告.
天父上帝我們多謝你.
多謝你.
因為看到以殺的生平的時候.
我們看到他有軟弱的地方.
但也都看到.
他真的處處的為主.
你來去順從.
他對你的信心.
都是我們今天需要的學校.
主幫助我們.
在這個世代裡面.
不同的時間.
會有不同的艱難的出現.
好像避無可避.
主但在這些艱難的時候.
我們就看到.
上帝你在當中.
叫我們對你的信心的操練.
幫助我們.
何時我們堅持去信服你.
行你的心意的時候.
讓我們看到主你的開出路.
主給我們有這份信心.
給我們有這份經歷.
我們多謝你.
垂聽禱告.
奉靠主基督的聖名.
阿們.
謝謝大家收看 再見.
\newpage



\section{但以理書 6:1-3}
\label{sec:45mO3Udlc8E}
\textbf{2024-02-11 聯合崇拜《龍年生肖道》但以理書 6\_1-3 講員:陳韋安牧師}
\newline
\newline
連結: \href{https://youtube.com/watch?v=45mO3Udlc8E}{\texttt{https://youtube.com/watch?v=45mO3Udlc8E}} ~~~~ 語音日期: 2024-02-14
\newline
\newline
\hyperref[sec:OWpYTLtxpsA]{\small{< < < PREV SERMON < < <}}
~
\hyperref[sec:index_chronic]{\small{[返順時目]}}
~
\hyperref[sec:index_scriptual]{\small{[返順卷目]}}
~
\hyperref[sec:QYUt9NIhVQs]{\small{> > > NEXT SERMON > > >}}
\newline
\newline
但以理書 6:1-3
\newline
\begin{longtable}{cl}
\hline
\hline
章節 & 經文 (和合本修訂版)\\
\hline
6:1 & \begin{tabularx}{0.7\textwidth}{X} 大流士隨心所願，立了一百二十個總督，治理全國， \end{tabularx} \\ \\ \relax
6:2 & \begin{tabularx}{0.7\textwidth}{X} 又在他們以上立總長三人，但以理也在其中；使總督在他們三人面前呈報，免得王受虧損。 \end{tabularx} \\ \\ \relax
6:3 & \begin{tabularx}{0.7\textwidth}{X} 這但以理因有卓越的靈性，超乎其餘的總長和總督，王想立他治理全國。 \end{tabularx} \\ \\
[1ex]
\hline
\hline
\end{longtable}
$^{1}$好 各位姐妹們早安.
大年初二 祝大家在農年裡面.
大家能夠經歷耶穌基督的恩典.
能夠充滿基督的力量和身體健康.
今天我會講一篇新潮道.
其實我從羊年開始.
就在感受堂開始我自己的一個習慣.
雖然不是每年都在感受堂講.
但是今年到農年 我都保持這個習慣.
講一篇農年的新潮道.
不知道大家在十二新潮裡面.
大家最喜歡哪一個新少年.
有沒有哪一個特別喜歡.
怎樣問呢 不知道大家有何感覺.
在十二隻動物裡面.
其實有些動物會比較容易愛上.
例如虎年 羊年 兔年 馬年.
這些形象比較正面.
比較容易大家覺得愛上.
有些動物都很好的.
但要給予一些想法.
才會覺得這是一個年份.
小肥肥豬 豬年.
或者是蛇年.
都會給予不同的形式.
不過農年是很受歡迎的.
農年是一個在中國文化裡面.
非常非常愛好的動物.
在華人文化裡面.
農是象徵權威 力量 幸運.
農作為一個中國傳統裡面的神獸.
是君王的象徵.
都是代表著皇帝的力量和威嚴.
如果可以的話.
我們以前香港每逢農年的時候.
都特別多小朋友出世.
醫院特別多BB出世.
都追著要生小龍女 小龍仔.
所以都特別喜歡自己的小朋友.
在農年出世.

$^{41}$中國人認定自己是農的傳人.
所以我們香港以前都會被稱之為.
亞洲四小龍之類.
當然農的理解.
近來新聞所說.
和西方的文化很不同.
所以近來都開始嘗試.
不要叫龍 叫龍.
我就不說了.
因為會說不出龍這個動物.
因為龍在西方人裡面.
是代表著一個凶猛 邪惡的動物.
這個是在遠古的歐洲裡面.
已經是這樣的傳統.
龍常被描繪為一個很大型.
有翅膀的蛇.
鱗片是閃爍著很多寶石的光彩.
吐出來的都是猛烈的火焰.
很大的破壞力和傷害力.
特別在中世紀裡面.
隨著基督教的傳播.
龍開始被視為一個邪惡的化身.
星裡面的巨龍.
他們都純粹被認為是一個殺蛋的化身.
無論在啟示錄十十二章.
或者在不同的星裡面.
都是一個負面的代表.
所以很多時候就發覺.
開始有勇者救惡龍.
通常是一個勇士為了救公主.
成為了敵人 惡龍的對象.
所以你發覺後來.
有很多不同的理解.
西方裡面的龍都開始比較正面.
例如哈利波特.
是比較有智慧.
守護著一些寶物的形象.
甚至英國龍都成為了皇室的象徵.
成為了一個皇權和力量的記號.
所以發覺一大堆.

$^{81}$我們對於龍的歷史.
都是大家不陌生的.
說這麼久其實大家都知道.
其實龍是不存在的.
大家知道龍不是一個真實存在的動物.
是我們想像出來的.
不是上帝所創造出來的.
恐龍其實就不是龍.
雖然一個叫Dinosaur.
一個叫Dragon.
其實兩個完全不同的生物.
所以我們一直都是用一個.
一個是想像的.
一個是我們想像出來的動物.
作為我們一年的動物.
不過其實大家都不陌生.
對於一些不存在的動物.
其實大家都是很熟悉的.
無論是Pokemon 還是Doraemon.
還是Scarab 還是Godzilla.
Hello Kitty.
其實這些都不是神所創造出來的.
不過我們又有很多很多.
這些產品在我們的世界裡面.
大家想一下床頭裡面.
很多公仔.
全部都不是真的.
全部都是一些人想像出來的動物.
所以我們都習慣了.
活在一個不存在的生物裡面的世界.
一隻無口的Hello Kitty.
一隻會發電的比卡超.
一條沒有存在的龍.
都是我們在歷史上.
都已經很知道.
習慣它的存在.
所以發現從小到大.
你對龍都不陌生.
你知道雖然它不存在.
但是你對它非常熟悉.

$^{121}$甚至乎有感覺.
那裡面是很愛好.
甚至乎我們世上有成千上萬的事情.
和人都和龍是一種關係.
龍蝦 龍眼 龍井 龍貓 龍躉 龍豬 龍五.
龍舟 龍頭棍 龍鬚膛 龍虎門 龍捲風 龍鳳胎 龍應台.
龍金威 李小龍 小火龍 泡泡龍 法骨龍 九龍 成龍.
全部都是和龍有關係.
再不計算整個龍在江湖 龍在邊緣 龍騰四海.
還有陀鶴威龍 古惑仔之猛龍過江 古惑仔之駁龍爭虎鬥.
很多這些和龍有關係的東西.
其實一直都存在.
但會是連上一個不存在的動物.
所以驅龍年竟然是一個不存在的生物.
竟然能夠成為我們中國人的一年.
是一件很匪夷所思的事情.
說了這麼久 今年想說些什麼呢.
發覺既然是這樣.
要選一篇有關龍做主題的聖經.
其實就不容易.
一來就不多 二來就不太合適.
如果這樣找一個龍字的話.
最多出現的就是阿沙龍和沙龍這兩個人名.
沙龍是代表Shalom的意思.
一個普遍猶太人的人名.
阿沙龍就是大衛的兒子.
最多的就是剛才說的.
啟示錄第十二章經文.
在預言裡面所說將來的撒旦.
作為一個惡龍來象徵住.
不過既然龍既存在又不存在.
感覺上又好像沒有的生物.
我就選了一段很特別的經文.
一個叫做隱藏經文.
如果你看大約書十十四章的話.
是有龍的存在.
大約書十四章.
你會發覺你找不到十四章.
起碼這本聖經沒有的.
因為這十四章是一個隱藏經文.

$^{161}$一個在新教的聖經是沒有的.
基督教的聖經是沒有經文的.
不過如果找回天主教聖經的話.
是有的.
這段是稱之為次經經文.
即是說僅僅在天主教的傳統裡面.
作為一個次經一個符碌的存在.
它不在一個正典裡面.
它屬於一個附加的經文.
說回正經就好像今天大家看電影的彩蛋一樣.
並不在一個正史裡面.
這個當作一個平行宇宙裡面的彈射理.
一個是一個What if.
是一個另一個時空裡面彈射出現的經文.
同一個故事說的是一個不同的版本.
或者不同的原因.
同樣都是在說為什麼戴尼會被人捉到獅子坑.
但卻是出於不同的原因.
事實上這兩件經文.
十三章的經文是叫做次經.
它是由希臘文所寫成的.
不是希伯來文.
一個是後來有人用希臘文的版本.
一個是次經版本.
所以發現這些經文其實正正就是一個.
是一個人後來附加上去.
另一個What if的故事.
第一個故事稱之為Belle and the Dragon.
就叫做貝爾與龍的故事.
所以今天我們簡單來說.
我們會說故事.
大概就會說一下.
究竟這個隱藏經文裡面發生什麼事.
大概有兩個故事是講出來的.
一個叫做但爾尼金田一.
一個是偵探故事.
第二個就是一個叫做但爾尼屠龍記.
就是有關但爾尼屠龍的故事.
那我就說回這個故事.
在很久以前的波斯.

$^{201}$有一個叫做居魯士的國王.
居魯士在位的時候.
身邊有一個人叫做但爾尼.
他在王裡面是一個非常高的地位.
在居魯士在位的時候.
身邊有一個但爾尼.
他很得居魯士王的敬愛和信任.
同時間巴比倫人就有一個敬拜對象.
叫做貝爾.
貝爾是一個偶像.
每天都吃很多東西.
吃十二代的麵粉.
吃十隻羊和六桶酒.
那個王都很敬拜貝爾的偶像.
每天都去寺廟來敬拜他.
但爾尼並沒有這樣做.
他只是敬拜自己的耶和華上帝.
不是貝爾.
那個王就問他.
為什麼你不拜貝爾.
但爾尼就說.
我不敬拜人所做的偶像.
他是沒有生命的.
不會走的.
我只敬拜創天造地.
永活的耶和華上帝.
那個王就問.
你不覺得貝爾有生命嗎.
你不見得貝爾都吃很多東西和喝很多酒嗎.
但爾尼就笑著說.
大王 貝爾是一個由泥造成的偶像.
他怎會懂得吃東西和喝酒呢.
那個王就很生氣.
就立即叫那些祭司來見他.
就問那些祭司.
究竟貝爾懂不懂得吃東西.
祭司就說.
雖然貝爾吃東西是吃很多東西的.
大王 但爾尼忽然就想誣告我們.
我們就把食物放在寺廟裡.

$^{241}$就關上門.
你看看第二天食物會不會吃光.
吃光的話 大尼就說謊.
他就要去處死.
於是他就這樣做.
就把食物放在寺廟裡.
就鎖上門.
第二天清晨的時候.
一打開門.
那個王就回來寺廟裡.
就發覺那些食物竟然是吃光了.
貝爾是真神.
他吃光了東西.
比我女兒還乖.
突然間 但爾尼就站到最高就說.
一切的謎底已經解開了.
我奉我耶和華上帝之名.
解開這個真相.
真相只有一局.
大王就看看.
地上就看到有些腳印.
昨天我故意把廟裡的糞便都撒在地上.
我們就沿著腳印去看看.
就發現有個暗門在那裡.
你看看這裡的廟裡.
這個暗門是通往一個暗道裡的.
這麼多年來.
一家七十人.
其實是吃光了你所有的寺廟裡的食物.
其實吃光了的不是貝爾.
而是這班祭司和他的妻子和子女.
貝爾是假的神.
沒有生命和能力吃東西.
吃的就是這班人.
於是高王就覺得自己受騙.
然後就很生氣.
就殺了所有的祭司.
這個就是第一故事.
叫做貝爾與貝爾的故事.
第二個就是關於龍的故事.

$^{281}$叫做但爾尼屠龍記.
很久以前巴比倫王的城市裡住了一隻巨龍.
一隻很大的巨龍.
巴比倫人就把這隻巨龍視為神明.
視為一個有生命的上帝.
這個就不是沒有生命.
因為是有生命的神.
高王就命令大尉尼快點去敬拜他.
因為這條是一個非常之威嚴.
值得我們敬拜的龍.
大尉尼就說我只是敬拜我的上帝.
因為苗華是永活的神.
其他都是假神.
有一天會死的.
大尉尼就說這麼大條龍這麼容易死.
大尉尼就說.
大尉尼甚至可以不用刀不用劍.
都能夠殺死這條龍.
於是大尉尼就答應了他的要求.
就讓他這樣教了龍.
於是大尉尼就用了一些瀝青.
油和頭髮.
混合成一顆丸子.
我們香港吃的頭髮奶茶那些.
變成一顆很大的丸子.
餵給巨龍吃.
巨龍一吃.
胃就馬上爆裂.
馬上就死了.
大尉尼就說.
你看看這條龍已經死了.
你為什麼要敬拜耶和華上帝.
當巴布倫人聽見這個事情之後.
就大衛憤怒.
就說這個大尉尼既推翻我們的偶像貝爾.
又殺了這條龍.
我們要把他抓去獅子坑那裡.
於是大尉尼就被人抓去獅子坑那裡.
這樣就六天的時間.
但是上帝保守著他.

$^{321}$獅子都沒有吃到大尉尼.
最後和大家一樣這樣完結.
就被人救回來了.
這兩個就是隱藏的故事.
大尉尼裡面不存在的.
他不在聖經裡面.
一個是在《刺經》裡面的經文.
這兩個故事到底有什麼意義呢.
正所謂這兩個故事並沒有納入聖經正典裡面.
你找不到的.
我想這兩個不在聖經裡面有它的原因.
第一最少是什麼.
因為大尉尼沒有殺過這條龍.
因為龍是不存在的.
所以這個應該是想出來的故事.
不是真實的.
大尉尼都沒有那麼有型.
沒有一個像偵探那樣.
來揭穿一切人的詭計.
好像金田一一樣.
他沒有霸氣以葉和華的名去發誓.
去搗破敵人的詭計.
其實大尉尼都沒有辦法去證明.
其他他們宿敬拜的偶像是假的.
自己的神是真的.
但是你沒有這樣做.
但你更加沒有這樣的霸氣.
來公然在巴比倫城裡面.
去殺了人家那條巨龍.
親手去屠龍.
我們可以想一下.
其實金田一旦而來.
大尉尼屠龍記其實都不是我們所認識的大尉尼.
這個不是我們基督徒感覺到的大尉尼.
如果我們去打開正典.
打開大尉尼出的第六章的時候.
你會發覺整個故事裡面.
真實的大尉尼其實是怎樣的.
和我們是一模一樣.
都是一個很普普通通的人.

$^{361}$在這個世界裡面.
都是一個很微小.
甚至是一個很被動狀態的人.
在正經裡面大尉尼書是怎樣記載的.
大尉尼是什麼原因被人抓去獅子坑裡面.
大家知不知道.
為什麼.
清理文獸打開正經就可以了.
打開第十章經文.
打開正經就可以了.
根本講第六章第六節.
大尉尼就知道禁令蓋了玉璽.
就到自己家裡一天三次.
相室跪在他神面前.
為什麼人們會抓他去獅子坑裡面.
因為那時候全國頒布了道歉禁令.
禁止人去祈禱.
大尉尼他是堅持自己的信仰.
繼續去祈禱.
安於素常繼續祈禱.
所以就被人抓去獅子坑裡面.
這才是真正正經裡面的大尉尼.
這個其實和我們一樣.
在祈禱.
每天都在祈禱.
做基督徒的大尉尼.
大尉尼並沒有像這些知識所講.
成為一個名偵探金田一.
或者一個親手屠龍的勇者.
他沒有機會去拆穿世界的騙局.
都沒有能力去反擊.
攻擊他的人.
更加沒有權力去摧毀一些惡人所做的事情.
大尉尼其實只是在做一些.
信仰裡面應該要做的事情.
所以強調大尉尼屈膝在神面前.
正如小朋友所說.
大尉尼只是和平時一樣.
做一些很普通,很普通,很平凡的事情.
不過是一些很重要的事情.

$^{401}$沒有什麼花巧.
沒有任何的霸氣.
也沒有什麼攻擊力.
甚至這個世界似乎沒有什麼特別正面.
直接的影響.
不過這個單單來到去堅持去祈禱.
堅持去做一些很應該要做.
基本要做的事情.
做好它.
繼續做.
仍然去做.
其實已經是一件很了不起的事情.
親愛的兄弟姐妹.
真人.
聖經裡面的大尉尼.
他在祈禱.
一個很普通,平凡的祈禱.
他是充滿著巨大的震撼力.
我不是說祈禱本身有多厲害.
雖然很厲害.
也不是說有什麼神蹟會發生在祈禱裡面.
一定不會這樣做.
你說單單來到去堅持去祈禱.
安於素常.
按照我們平時所做的祈禱.
這件事本身已經是充滿著很震撼的事情.
如果你讀大尉尼出的六章的時候.
你會發現.
其實大尉尼本身這個人.
所流露出來的生命.
本來就已經是令到皇帝很讚嘆生命.
他不需要做很多的事情.
當人去誣衊大尉尼的時候.
皇帝怎麼說呢.
皇帝說他一心就受煩.
一心就要救大尉尼.
策劃去解救他.
所以大尉尼的power.
不是在於他有什麼獨特的神蹟奇事智慧.
或者有什麼勇氣大力.

$^{441}$而是他本身這個人.
安於素常地做回他自己應該要做的事情.
做好它.
keep住去做.
這個事情是更加震撼人的事情.
聽說在這個世界裡面.
你堅持你的信仰.
已經是一件很震撼的事情.
在一個這樣的世界和社會裡面.
你能夠可以很簡單地去堅持你自己所信的.
做你所信的.
每天都做這些事情.
已經是一件很震撼的事情.
你不需要殺一條龍.
更加不需要成為一個什麼英雄.
簡單微小的去做.
這個比殺一條龍更加厲害.
或者這個成為我們今年龍年的方向.
龍年提醒我們.
其實我們一直都和一些不存在,不真實.
或者不是那麼虛構的東西共存.
明明是沒有存在過.
世上的人都會視它為最尊貴的動物.
明明不是真相.
竟然沒有人將它一語道破.
明明是一個謊言.
卻每天很多人大肆宣傳.
我們活在這樣的世界裡面.
一直都是.
好幾十年,好幾百年,好幾千年.
先不說幾年前在同水大肆宣傳的JPEGS廣告.
或者不是說幾十年前.
金手指,家庭集團那些東西.
假雞蛋,假奶粉那些東西.
那些虛假,不存在的東西.
一直都肆無忌憚的存在.
不過鄧小妹.
世上最龐大,最虛構,最不真實的是什麼?.
最軟歷史的是什麼?.
不是一個集團.

$^{481}$甚至乎不是某些政權.
而是世上的黑暗力量.
世上來說這個叫做Nothingness.
虛無.
為什麼叫做虛無?.
因為這些東西是明明不存在的.
明明上帝是沒有創造出來的.
明明不是神所創造的.
它竟然可以暫時存在在這個世界上.
所以稱之為Nothingness,虛無.
不過這些不存在的東西.
它是很猖狂地以為它是真實的存在.
並且好像成為這個世界的主流.
不過不存在仍然是不存在.
無論它多麼的猖狂.
它最終都在天賦上帝永恆的角度裡.
它依然都沒有存在過.
這個就是Nothingness.
所以我們相信黑暗終必過去.
活在這個Nothingness的世界裡面.
你要去作為基督徒去堅持去做一些真實的東西.
一些合乎真理的東西.
一些很具體存在的東西.
一些美善的東西.
它顯得是無比的價值.
2024年我由年初基本上是病到現在.
年初第一個星期就發燒.
第二個星期就咳嗽.
然後又胃炎.
經歷過很多不同的小病.
病到現在喉嚨都有點聲音.
在這個病裡面我發覺.
人基本上是病到一個地步.
什麼都做不到.
很簡單地去走出客廳.
倒杯水給自己喝.
或者好好地做好一件工作上的一個小的任務.
已經是很不簡單.
這個病裡面我發覺原來我們人.
好好地去做好一件事.

$^{521}$完成了它.
很踏實地去做好一件事.
這個正正是一件很不簡單的事.
所以我明白我以前在德國裡面讀書.
德國人的做法.
可以在一分鐘裡面慢慢地寫好一個字.
到你去報名的時候.
就會貼上你的名字.
很簡單,很慢地寫好一個字.
原來很真誠,很用心去做好一件很小的事情.
已經是一件非常了不起的事.
特別是在信仰裡面.
當我們作為基督徒.
好好地去堅持.
做好每一件很微小的事情.
祈好一個土.
專心好好地愛一個人.
好好地讀好今天要看的聖經.
你做好它.
每天都繼續去做.
這些這麼微小的事情.
做了就做了.
我想這間教會也是一樣.
這間教會最吸引人的地方.
可能現在說的很多就是那些流量密碼.
什麼公關,什麼宣傳.
但我想作為一間教會.
最重要的.
就是真的做一間好的教會.
教會應該要有的東西.
簡單,用心地做好它.
起碼大家回來這裡的時候.
真實地彼此相愛.
好好地衷心問候.
我想這個教會裡面的人.
每一天都做一些基本上要做的事情的時候.
這個正是成為了一個很真實,很堅定的事情.
讓龍蓮,讓錦繡堂在錦繡花園裡面.
成為一個很真實的教會.
很多很多在這個世界裡面.

$^{561}$一間很真正發揮到教會功用的教會.
是一件很不離不棄的事情.
讓我們一樣在我們的生命上.
堅持做好每一件要做的事情.
這個我想正是龍蓮最蒙福的途徑.
正是我們龍蓮要去追求的事情.
我們一起祈禱.
因為你在我們龍蓮裡面有你的恩典.
我們懂得去回應你的恩典.
在龍蓮裡面我們學求的.
是很好的去追求我們的信仰.
做好每一件我們要做的事情.
愛好每一個人.
讓我們在同理的關係裡面.
堅持,堅守每一件小小的步驟.
讓錦繡堂能夠成為一個你所喜悅的教會.
一個真實,一個發揮到真正功用的教會.
敬拜群體.
一個彼此相愛的群體.
願意去傳揚你聖名的群體.
我們簡單的,微小的,謙卑的去做.
求你保守帶領我們.
逢春命寇.
阿們.
謝謝大家.
\newpage



\section{馬太福音 5:3-10}
\label{sec:QYUt9NIhVQs}
\textbf{2024-02-10 新春崇拜《福從天降》馬太福音 5\_3-10 講員:陳一華牧師}
\newline
\newline
連結: \href{https://youtube.com/watch?v=QYUt9NIhVQs}{\texttt{https://youtube.com/watch?v=QYUt9NIhVQs}} ~~~~ 語音日期: 2024-02-15
\newline
\newline
\hyperref[sec:45mO3Udlc8E]{\small{< < < PREV SERMON < < <}}
~
\hyperref[sec:index_chronic]{\small{[返順時目]}}
~
\hyperref[sec:index_scriptual]{\small{[返順卷目]}}
~
\hyperref[sec:oCSQzasymjk]{\small{> > > NEXT SERMON > > >}}
\newline
\newline
馬太福音 5:3-10
\newline
\begin{longtable}{cl}
\hline
\hline
章節 & 經文 (和合本修訂版)\\
\hline
5:3 & \begin{tabularx}{0.7\textwidth}{X} 「心靈貧窮的人有福了！因為天國是他們的。 \end{tabularx} \\ \\ \relax
5:4 & \begin{tabularx}{0.7\textwidth}{X} 哀慟的人有福了！因為他們必得安慰。 \end{tabularx} \\ \\ \relax
5:5 & \begin{tabularx}{0.7\textwidth}{X} 謙和的人有福了！因為他們必承受土地。 \end{tabularx} \\ \\ \relax
5:6 & \begin{tabularx}{0.7\textwidth}{X} 飢渴慕義的人有福了！因為他們必得飽足。 \end{tabularx} \\ \\ \relax
5:7 & \begin{tabularx}{0.7\textwidth}{X} 憐憫人的人有福了！因為他們必蒙憐憫。 \end{tabularx} \\ \\ \relax
5:8 & \begin{tabularx}{0.7\textwidth}{X} 清心的人有福了！因為他們必得見神。 \end{tabularx} \\ \\ \relax
5:9 & \begin{tabularx}{0.7\textwidth}{X} 締造和平的人有福了！因為他們必稱為神的兒子。 \end{tabularx} \\ \\ \relax
5:10 & \begin{tabularx}{0.7\textwidth}{X} 為義受迫害的人有福了！因為天國是他們的。 \end{tabularx} \\ \\
[1ex]
\hline
\hline
\end{longtable}
$^{1}$各位姐妹早安.
早安.
首先在這裡恭祝大家新年蒙福 身體健康.
我們也在旁邊祝大家八個年 互相恭賀.
平時大家很少見牧師穿得這麼鮮艷 是吧.
是一年穿一次.
看到牧師穿著這樣的裝束 大家都知道新年來了.
我們今天真的是大年初一.
來到新年一開始 大家的心就在想著上帝.
我們要什麼呢 去敬拜上帝.
所以我們錦宣家很多年來都有這個優良的傳統.
我們逢是年初一的大清晨.
我們第一時間就要去敬拜天上的父親.
敬拜天父 因為他是創天造地.
他大有能力.
他時刻眷顧保守 和引領我們的上帝.
所以我們要在新年開始 向他感恩讚美.
這是非常好的.
今天牧師講道的時候 順便提醒大家.
如果中途聽到前面兩個麥克風 前面兩個喇叭.
砰 砰 你不用怕.
這是娛慶節目 是炮仗聲.
我們沒有真炮仗 但有假炮仗聲.
所以大家可以專心聽牧師講道.
鄭華娟樂曲隊唱詩的時候.
她說到中國人的輝春很喜歡有一隻字.
我想大家都知道這是什麼字.
這是一個福字.
有些人唸上面 有些人倒轉來唸.
這個福字是中國人男女老幼 每個人都很喜歡這個字.
我們家裡很多時候都會將輝春.
將一些很好的 精簡的說話唸出來.
有些人說 福從天降 五福臨門 福升拱照.
我不知道你家裡有沒有唸過福字出來.
有唸的 舉手一下 你家裡有唸的福字.
哇 大部分都有唸.
每個人都貼了一個福字出來.
贊助是怎樣擺法 贊助是什麼福.
這個很好.
但基督徒通常福貼出來的時候.

$^{41}$就有這三種.
一 看你中不中.
有沒有人是福杯滿日的 有.
祝福滿滿 百福臨門 都有 非常之好.
證明大家新年前拿了很多輝春回家.
然後就選擇這些輝春 有個福周圍貼.
這個很開心的.
這個是我們一個很感恩的.
說起福 信主和未信的人.
他們的追求都有點不同的.
大家有沒有發現.
未信的朋友他們的追求.
起初是五福臨門的 是不是.
我們基督徒就多幾個福 是不是.
百福臨門.
為什麼會有一點點的不同呢.
因為是我們屬靈的生命.
我們得到上帝的拯救.
得到上帝的救恩.
重新得救之後.
我們是神的子女.
所以我們的看法 我們的認識.
我們的認知.
告訴我們原來我們追求的福.
是要百福臨門的.
一般朋友他們追求五個福.
在我們中國人的社會當中.
主要是長壽的福 是不是.
富貴的福 是不是.
或者康寧的福.
浩德的福 是不是.
善忠 追求這五個福.
後來那些人演變了.
改了一個浩德.
就變了這個福 是什麼福.
兒孫滿堂.
就很想能夠什麼.
兒孫滿堂 是不是.
各位頂次母們.
這五個福.

$^{81}$好是非常之好的.
我們每個人都很想得到.
但是我們又看回一個.
很真實現實的世界.
我們發覺原來在人世間.
很多我們很想追求的福.
都未必是一個永恆的東西.
大家認不認同.
譬如好像長壽.
最多都是長命百二歲.
你明白我的意思嗎.
沒有永恆的.
沒有長生不老的.
也都是這樣的情況.
又或者你說我想追求富貴.
但是富貴也不讓我們掌控.
所以你發覺原來很多東西.
都不讓我們掌握.
你說浩德.
有時候未必真的自己.
想做一些好的事情.
為什麼.
因為我們人始終都很有.
怎麼說.
很有私慾的問題.
很有情慾的問題.
所以保羅說.
我們心裡面有兩個律.
你明白嗎.
這兩個律在我們心裡面.
常常在增進中.
保羅說.
立志行善是甚麼.
由得我.
只是走出來是甚麼.
由不得我.
我所願意的善.
我賭不去作.
我所不願意的惡.
我賭甚麼.

$^{121}$去作.
有時候發覺人很想是浩德.
很想是浩斯,恆德.
但是很多時候我們都掌握不到.
因為我們人是很脆弱的.
我們的本性.
有時候是有情慾來控制著我們.
所以有時候你發覺到.
都不到我們去掌控.
你說康寧.
康寧我們每個人都期盼.
國家或者我們的社會.
安康,安寧.
大家過一個平安的日子.
大家都很想的.
但是你看回我們現在國際間的一些新聞.
你就發覺到.
都不到我們控制的.
忽然之間.
某一些的爭論.
某一些的利益的問題.
你就發覺.
大家就翻臉了.
大家就訴諸於武力.
所以我們在香港這個地方.
上帝都給我們能夠有什麼?.
有安居樂業的生活.
這個很感恩的.
但是你記得.
也都不到我們來到什麼?.
永遠的去掌控的.
好.
剛才有一個什麼兒孫滿堂.
也都不是的.
你知道現在的年輕人.
很遲才結婚的.
有些甚至不結婚.
有些結了婚.
又養隻貓養隻狗.
想不想要小朋友?.

$^{161}$又不想.
你明白嗎?.
又不到我們大人說了算.
所以有時候很多事情都不到我們掌握.
不到我們掌控.
也都不到我們掌握的.
所以還有一場情況.
原來我們追求那五個福.
五福臨門.
有時候你發覺.
這些福.
也都是什麼?.
是信息萬變的.
大家認不認同?.
就好像我們2008年出現了.
這個98年有金融風暴.
2008年有什麼?.
金融海嘯.
很多的疫情.
很多的一些因素.
令到我們發覺.
這個資產.
今天是屬於你的.
第二天是什麼?.
是屬於別人的.
大家相不相信?.
所以你發覺到.
原來我們追求五福是好的.
但是我們知道不到我們掌控.
也都不是永恆的.
所以今天牧師和大家分享的.
我們要追求的是什麼?.
就是提升我們的眼光.
我們的視線.
不是單單注視在地上的.
我們更加要追求.
渴求提升我們的目光.
是要追求更高層次的.
是一個永恆的.
是一個千真萬確的.

$^{201}$就是聖經給我們的應許.
也是主耶穌親口給我們的祝福.
這些福是什麼福呢?.
就是大家剛才讀的經文.
那個八福在當中.
今早牧師當然不會將八個福.
詳細和你們分享.
因為真的不夠時間.
我怕阻礙你們出去拜年.
平時的人這個八福.
可以講八堂道.
你相不相信?.
你有沒有聽過有些牧師講到.
將八福講八堂道.
你聽過沒有?.
因為內容很豐富.
內容是非常令人.
很扣人心緣的.
所以我們看到.
這八個福當中.
稍為和大家過目看看.
聖經說心靈貧窮的人有福.
大家都知道它是什麼.
它是很渴求,很追求屬靈的成長.
希望吸收到更多屬靈的話語.
餵養自己.
有個靈命的操練.
經常覺得不夠.
這是什麼?.
天國是他們的.
然後聖經耶穌又說.
哀動的人有福.
為什麼?.
這個人懂得為自己.
又懂得為別人難過的人.
你不要看他難過是一個弱者.
不是的.
為什麼?.
因為他真的有同理心.
他真的有敏銳的觸覺.

$^{241}$他發現原來自己有些地方不足.
有瑕疵.
所以他懂得為自己.
為別人哀動.
這些人.
主耶穌說什麼?.
他們必得安慰.
是不是?.
耶穌又說.
溫柔的人有福.
溫柔是說我們有方寸,有分寸.
懂得去管理自己的才能.
這樣的人有福.
為什麼?.
因為他知道什麼適當的時候.
就做適當的事情.
這些有分寸,懂得去管理自己的人.
有福.
為什麼?.
因為將來他得到的福是什麼?.
可以承受地土.
即是他能夠掌管更多的資產.
然後主耶穌又說.
飢渴無意的人有福.
為什麼?.
因為他在生活當中.
他是以神的國,神的義作為他的目標.
所以他追求什麼?.
去完成上帝喜悅的事情.
這些人有福.
為天國,為福音去打拼,拼搏.
這些人有福.
他們追求實行上帝喜歡的事情.
所以這些人就能夠得到飽足.
即是能夠得償所願的意思.
然後耶穌又說.
連戌人的人有福.
當他能夠關心別人.
當他能夠關愛別人.
當他能夠用行動去愛護別人的時候.

$^{281}$體諒別人的時候.
這些人有福.
為什麼?.
因為他們必得連戌.
然後耶穌又說.
清心的人有福了.
我見大家都很想回應一下.
這個福是什麼福?.
因為他們必得見神.
原來一個人心靈內心很純正,很單純.
很純潔的人.
他有福了.
為什麼?.
他的單純,他的純潔.
就好像小孩,小朋友一樣.
上帝讓他看到更多的神蹟奇事.
上帝讓他看到更多上帝的作為.
為什麼?.
因為他有一個單純的心.
清心的人有福了.
然後耶穌又說.
使人和睦的人有福了.
因為他們必得什麼?.
必稱為什麼?.
神的兒子.
即是說我們在生活當中.
我們經常能夠成為一個和平之子.
不是製造麻煩.
不是製造分端.
而是製造一個和平.
令人與人之間有一個和諧.
這個人有福了.
為什麼?.
因為他做上帝喜悅做的事情.
結果他就稱為神的兒子.
最後一個福.
大家都知道耶穌說什麼.
為而受逼不的人有福了.
因為什麼?.
天國是他們的.

$^{321}$即是說我們為了上帝.
為了福音.
為了天國.
為了造神喜悅的事情.
而遭遇到一些迫不及待.
遭遇到一些反對,抗拒.
別人不認同.
別人恥笑.
我們不要為這些事情寄掛在心裡.
為什麼?.
因為這些事情是必然有的.
因為這個世界是魔鬼來的.
控制著的.
所以地上的世界出現這些事情.
是不足為奇的.
所以我們不要因為這樣而放棄原則.
我們的立場.
我們要堅守聖經給我們的教導.
實踐聖經的說話.
主耶穌的英訊是什麼?.
天國是他們的.
好,牧仙很快就將這八個福.
跟大家略略的來溫習一下.
其實我想大家知道.
這八個福都是非常寶貴的.
這八個福不單是地上得福.
甚至將來在天上都得到賞賜的.
都得著冠冕的.
今天我想告訴大家.
原來領受八福的人.
這個人是一個什麼樣的人呢?.
原來這個人真是一個世界上被人很羨慕的人.
因為這個人是真正蒙福的人.
為什麼呢?.
因為領受八福的人有一個特點.
就是當他追求行神喜悅的事情.
而得到神的賜福的時候.
他不是停在那個階段.
他是很樂意將上帝給他的福.
能夠帶領人,鼓勵人,激勵其他人.

$^{361}$都同樣來到追求得到神的福.
所以他很願意將得到神的福.
這個恩惠與人分享.
我認識一個機構.
它都很特別.
這個機構的名字叫做奉獻力量.
我不知道大家有沒有聽過這個機構.
這個奉獻力量是怎樣呢?.
是由一班基督徒商人組成的.
他們組織不是很繁複,很嚴密.
他們是一班很愛主的基督徒.
做生意的人走在一起.
他們有一個目標.
就是每一年他們想賺到的錢.
將它累積了之後.
他們就全年奉獻出來.
他們的名字叫奉獻力量.
來推動福音的工作.
推動一些很需要關懷的人,家庭.
去關心他們.
我覺得這個機構是非常之好.
原來他們知道.
好像我們奉獻師所說的.
萬物都是從主而來.
他們將從主而來的.
所得到的上帝的恩典,祝福.
不是據為私有.
而是大家都拿出來奉獻.
然後去關心,去支持其他福音的工作.
我為他們很感恩.
因為他們真的明白什麼呢?.
他們明白到原來追求百福的人.
是以神為中心的.
樂意別人都得到什麼呢?.
神的恩典.
怎麼別人得到神的祝福呢?.
就是我們一班教徒.
願意將上帝的祝福.
將上帝的恩典.
來與人分享.

$^{401}$令到其他人都得福的時候.
這個福是一個真正.
是一個很有意義的福氣.
我們發覺到.
五福和百福.
是有明顯的分別在當中.
我們追求五福的.
你發覺到我們都是以個人為中心多一點.
長壽,富貴,康寧,子孫滿堂.
都是比較個人性質多一點.
但如果我們追求百福的就不同了.
是以神為中心.
這是一個很明顯的分別.
五福的通常都是想著.
自己得到這個福氣.
擁有這個福氣.
但是百福就是很希望.
更多人都得到福氣.
更多人都認識上帝.
都能夠成為上帝的兒女.
都得到從上而來的真福.
這也是第二個分別.
五福大家都知道.
這是一個地上的福.
而百福是一個屬天的福氣.
是福從天降.
當然大家都知道.
五福是沒有辦法永恆的.
但是百福是可以延續到永恆的.
所以我們明白到.
上帝所賜的福的時候.
我們就知道聖經說得非常之好.
它說你要認識神.
就得平安.
然後說什麼呢?.
福氣也必臨到你.
但當我們追求百福的時候.
這個百福也都帶我們去到一個更高的境界.
是什麼呢?.
就是《使徒行傳》二十章三十五節.

$^{441}$告訴我們.
是一句什麼聖經呢?.
是保羅說的.
保羅益述主耶穌曾經說過一句話.
是什麼呢?.
這句就是「施比受,更為有福」.
這是不是一個真正的挑戰?.
原來我們生活在人世間.
我們不是因為去找很多東西滿足自己.
而是我們學習.
不斷地來到什麼?.
施予.
不斷地來到去分享.
不斷地能夠將神的恩典去見證上帝的作為.
當我們這樣做的時候.
就好像《使徒行傳》提醒我們.
耶穌說過一句話.
最高的境界是什麼呢?.
施比受,更為有福.
正如有人經常問你一條問題.
你猜被人請吃飯的那個人有福嗎?.
還是請人吃飯的那個人有福呢?.
有沒有人問過這條問題?.
你懂不懂回答?.
各位弟兄姊妹們.
今天我們知道.
原來真正的福是什麼呢?.
是福從天降.
Amen?.
我們低頭祈禱.
主耶穌多謝你.
因為我們得到你的話語.
給我們一個非常寶貴的提醒.
叫我們在世間上追求的不是地上的福.
而是天上的福.
也讓我們知道你的應許永不改變.
你給我們的挑戰就是施比受,更為有福.
求主祝福帶領我們.
接下來這一年,2024年當中.
我們來實踐主你給我們的教訓.

$^{481}$讓我們的生活.
讓我們的生命活得精彩.
活得豐盛.
奉耶穌名求.
阿們.
多謝您收睇時局新聞,再會!.
\newpage



\section{瑪拉基書}
\label{sec:oCSQzasymjk}
\textbf{2024-02-15 周四晚查經班《瑪拉基書(6)》 講員:曾裔貴牧師}
\newline
\newline
連結: \href{https://youtube.com/watch?v=oCSQzasymjk}{\texttt{https://youtube.com/watch?v=oCSQzasymjk}} ~~~~ 語音日期: 2024-02-16
\newline
\newline
\hyperref[sec:QYUt9NIhVQs]{\small{< < < PREV SERMON < < <}}
~
\hyperref[sec:index_chronic]{\small{[返順時目]}}
~
\hyperref[sec:index_scriptual]{\small{[返順卷目]}}
~
\hyperref[sec:C6eEbdGd3Ao]{\small{> > > NEXT SERMON > > >}}
\newline
\newline
$^{1}$我們都可以有家庭的相聚.
幫助我們能夠回來在晚上的茶聽班.
盼望神的話語成為我們大家很多的提醒和祝福.
也盼望已經完成了的這些錄音.
在網上發佈的時候.
可以讓不同的地方的弟兄姊妹.
都得到透過聆聽 得到彼此的建立.
大家都是在神的話語裡面認識神.
還有一份愛神 渴慕服侍的心.
我們今天晚上等候在你的面前.
禱告奉耶穌基督的名 阿們.
都是再一次看第二章的《馬拉基書》.
過去也都再講那個背景.
都是被歸類在被擄歸回的歷史的那段期間.
大概是主前的460多年的時間.
很難確實一個確實的時間.
大部分猶太的學者都認為是460多年.
主前 也就是說還有大概400年.
主耶穌降生 去到這本書的第三章.
很快預言由這位彌賽亞的先鋒.
就是施洗約翰會來到這個世界.
成為主耶穌的準備.
是主耶穌自己確認了.
《馬拉基書》第三章講的那個使者.
就是講施洗約翰的.
我們可以去看第三章.
今天應該有機會去看第三章.
我們就會看這一位使者.
還有更重要的就是那一位彌賽亞降臨.
所以原來《十二小先知》.
或者歸回的歷史書 特別是撒迦利亞書.
講到將來的彌賽亞和將來的天國降臨.
就是第二次的降臨 是講得相當多的.
如果大家已經淡忘了.
有機會你就再去聽一下撒迦利亞書.
是有很多的預言的.
第二章我們就回頭看.
是當時才知道馬拉基.
馬拉基的原文的意思就是.
使者 Messenger這樣的意思.

$^{41}$剛好是個名字.
可能未必是特別要改這個名字.
而是很想告訴我們.
他是有一個信息.
就是猶太人他們被擄歸回了.
有聖殿了 有憲制了.
也都有尼希米的重建城牆了.
也都有他們的 回到一個規模.
就是大家都以自己是神的選民.
在耶路撒冷定居了.
也都有一個曾經的革新.
類似是一個復興.
但是慢慢的 當領袖一直淡去之後.
又再出現了百姓好像去到一個.
不知道有什麼可以去尋求的.
這樣的心情 或者這樣的狀態.
所以在那個敬拜裡面.
就開始出現了一種的停滯.
在敬拜事奉的生活裡面.
又會出現了一種.
好像慢慢又變成了一種例行公事.
自然就會 大家都習以為常的時候.
就不會感覺到一種敬拜的神秘感.
對敬拜帶來的那個無論是雀躍.
期待 或者是滿足.
這樣來到回家.
所以就很自然地.
敬拜帶不起那個生活的.
遇見神帶來的那個自己的更新.
所以又會出現了.
似乎這裡有 這一章都有四次.
是一個很重要的詞.
就是鬼咋 民間 百姓 人和人.
再加上就是對自己的.
很小的時候結婚的妻子虐待.
導致好像妻離子散那樣.
剛過來的時候.
就無意中看到新聞 國際新聞.
就是加拿大的總理杜魯多.
和他青梅竹馬的太太.

$^{81}$結婚的太太就正式的離婚 分居.
很長篇的 那些照片.
從很小的時候 青梅竹馬的時候.
一直拍拖結婚.
是所有人眼中的美滿婚姻.
最終現在都要離婚收場.
他主要想說的就是那些狗仔隊記者.
追蹤到他的分居太太.
去找另外的男朋友.
都是很有興趣報這些東西.
所以我們看第二章.
我們這次是從第十節開始看.
上一次就是講到祭司.
對於神的那個效忠出了亂子.
祭司扮演不了這個角色.
就不能夠有效地將神的真理.
傳遞給自己的百姓.
很自然就導致到第十節.
去到這一章的結尾.
第十七節就好像要.
不經意的按錯了按鈕.
那些警報響.
就要觸發了一級警報,二級警報.
好像不經意的按錯了按鈕.
這裡就是在講神透過祭司.
來將真理,憲制,關於神性的事.
透過他們去傳遞.
每一代的百姓都可以來到根據.
你和神的關係 透過憲制.
來得到神給你的那個教導.
但是當這個機制失效之後.
失效的原因上個星期我們講過.
就是他們違約.
祭司違反了神和尼微所立的約.
所以他們本來是會在百姓當中得到尊榮.
神就說將你們憲制的那些.
那些祭神的那些墳.
就是那些大便.
要塗在你們的臉上.
你想一下是多麼羞恥.

$^{121}$多麼被人藐視的一種 看為下賤.
這個是第九節.
所以我使得你們被眾人藐視 看為下賤.
就是說原來越討好那些百姓.
憲制不要那麼講究.
病的,瘸腿的都可以憲制.
你都收貨.
百姓不喜歡聽的某些的教導.
你又不去講.
你以為討好.
原來你去七除八扣.
這樣來有意識地.
來減低神的道的要求.
比如昨晚我們上課神學的講座.
國民節牧師都很坦白.
他說他神學院請教授.
請了五個都不對.
當然感恩他沒有說哪個名字.
不對的原因就不是那些學識,學歷不標稱.
是原來不再堅持聖經無誤.
那牧師就說在他任內.
他是不能夠讓這樣的事情發生的.
因為聖經無誤是一個神學院必須捍衛的.
你開放了這條線.
你就收不回了.
那要去到多大程度.
你才說他是錯呢?.
其實你打開了這道門.
其實是關不了的.
當然好像講遠了.
但其實祭司他們出現的.
這樣降低自己的教導.
就使百姓更加一直向下視那個底線.
更加對你不尊重.
更加藐視可有可無.
很簡單就是.
如果我們有聽過關於聖經無誤.
新派就已經沒有了這個底線.
什麼都是超自然的.
是會令到未信的人很反感的.

$^{161}$你說有神蹟,有天堂,有地獄,有審判.
有主耶穌的救贖.
是違反人性,違反一些理性的.
違反了人與人之間的博愛的.
那你說到這麼嚴肅.
人家是不會信的.
所以新派是沒有這些門檻的.
只是告訴你神愛你.
你做好就可以了.
是不會有超自然神蹟的.
因為超自然神蹟是不會有人信的.
這樣的宗教.
這樣的基督教.
是沒有殺傷力的基督教.
最終就是沒有人信.
因為既然你對我沒有要求.
我還怎麼去信呢?.
既然你都不相信有天堂地獄.
那我為什麼要信你呢?.
所以新派是不會再有市場的.
是不會改變到人心的.
所以什麼最導致那個導火線.
新派是沒落呢?.
是二次世界大戰.
是兩次的世界大戰.
原來人心是不可以透過教育.
可以改變的.
原來人與人之間都是會有仇恨.
國與國之間都會為了利益.
可以打生打死的.
所以如果沒有了一位真神.
他定立的標準和天堂地獄.
和有永恆的生命的存在.
其實是不會令到地上的人.
透過教育是會改變的.
那又好像說遠了.
關乎到聖經無誤的那個問題.
所以聖經無誤是不可以.
我聽梁院長講得很清楚很白的.
他說如果任何一個教授在他的任內.

$^{201}$一碰到聖經無誤.
就不需要談了.
就馬上要解僱了.
沒得談這個問題的.
因為這個是基要派 福音派.
賴以為福音派的基石來的.
你如果要和.
開始我們要討論一下.
聖經是不是真的無誤.
那已經是沒得談的了.
這個都是很重要的.
其實回頭看回這裡的時候.
你就看到上個星期我們講了.
祭司他嘴裡應該常傳知識.
第七節.
人應該由他的口去尋求神的律法.
因為他是萬軍之耶和華的使者.
就是差遣來奉神的名.
去傳遞真理的媒介.
所以他就是那個讓人去認識實的那個宗教.
當然絕對不是天主教所講的那些神父.
那個好像教皇那個宗保.
那個中間人.
你不透過他是不可以去到神那裡.
這個當然是後來天主教的謬誤.
那個等一下我們看到使者那裡.
再給大家一些類似的教導.
一些大家可能透過一些這樣的宗教的研究.
系統神學裡面的一些關乎到.
辨別一些的錯謬.
可以有一點點的涉獵.
大家可能從來都很少會聽到這些.
但是這裡就告訴我們.
原本祭司是應該要扮演著這個角色的.
第八節就是你們卻偏離正道.
使許多的人在律法上跌倒.
你們廢棄了我和你未所立的這個永約.
所以就是那個此消彼長.
就是當祭司自己的那個職份放低了.
所以你試一下想一下有一天.

$^{241}$你回來教會.
那個講台今天我們不講了.
講到太悶了 一講就每個人都睡覺了.
今天我們不講到.
我們講這套上次我看的電影.
大家評論一下.
還要不是我一個人講完.
是大家可以在位置上拿著麥克風.
大家講一下東西的.
接著我們唱一些流行曲.
你可以想一下一個地步就是.
那個教會越來越開心了.
不過記住開心完之後.
下個禮拜再下個禮拜再下個禮拜.
我用不著聽你說.
我聽其他人說更開心更好玩.
我就不再回教會了.
教會的存在就是要將真道.
來講明神的心意.
所以福音派很重要很強的.
或者再重一些就是宗教改革很強的.
就是《釋經》講道.
按上下文的正確的解經是很重要的.
就不是拿一節聖經.
將它無限的.
我今天覺得這節聖經對我有什麼幫助.
我和大家分享我的感受.
接著讀完那節聖經之後連解都不解了.
接著就講一些生活.
講一下電影,講一下時事.
這就叫做Pitching.
嚴格來說就不是Pitching了.
這個可能是你個人的在這裡分享.
嚴格來說這個就不是《釋經》講道.
所以大家要記得.
真是需要有《釋經》講道.
因為唯有神的道才可以帶領我們去認識神.
最終回來教會不是開心.
開心你可以打一張牌.
開心你可以去旅遊.

$^{281}$開心你可以.
當然我不是叫大家打牌.
開心是可以看一場電影都很開心的.
但是來到神的家,神的面前.
我們不是尋求開心.
我們是尋求認識,敬拜這位真神.
所以我們透過一個很開心的環境去敬拜神.
不是尋求終極的目標是開心.
所以我們要搞清楚那個目的.
如果不是其實就很可憐了.
所以我們要透過這件事再一次回到思想.
接著第十節開始出現了四至五次.
其實在這本書總共有六次.
一,二,三,四,五,六次講鬼詐.
在第二章出現了四次.
第十節開始了.
我們豈不是都是一位父嗎?.
就是這位真神.
猶太人豈不是一位神所做的嗎?.
我們各人怎麼以鬼詐代弟兄.
背棄了神與我們列祖所立的約呢?.
接著就講那個民間的日常生活.
就是各行各業.
不同的人.
猶太人行事鬼詐.
並且在以色列和耶路撒冷中.
行一件可憎的事.
因為猶太人褻瀆耶和華所喜愛的聖潔.
或者聖地都可以.
娶那些事奉愛邦神的女子為妻.
凡是行者事的無論何人.
就是獻供物給萬君子耶和華.
耶和華也必從雅各的帳篷中剪除他.
這是一個普遍現象.
猶太人行事鬼詐.
第十節就在講.
各人怎麼以鬼詐代弟兄.
我想印一兩段經文.
我今天在鬼詐這個字.
搜尋一下聖經出現的次數.

$^{321}$用法都頗多和頗明顯的.
真言就很多.
特別提到關於鬼詐.
和正直的對比.
和大家揭一揭真言.
比較多的是真言的十二章.
又出現好幾次關於正直,鬼詐.
會走的人生的路是很不同很詭異的.
十二章第一節.
他說喜愛管教就是喜愛知識.
恨惡責備卻是畜類.
畜生才不喜歡人責備他.
善人必蒙耶和華.
因為善人很有善心.
很kind heart.
聖靈所寄的果子.
人愛就是第一個.
設詭計的人.
耶和華必定他的罪.
我們再看第五節.
其實整篇都很好看.
第五節 義人的思念是公平.
惡人的計謀是鬼詐.
善惡就在這裡分野.
善不只是人好這麼簡單.
而是你思考的方式.
即是你的內心的價值觀.
內心的做人態度.
內心的思想.
暫時都找不到有什麼儀器可以照到那個思想.
心啊.
連那些芯片圖.
心的血管.
其他的都已經可以發明了.
但是思想就發明不了.
也都不需要了.
因為聖經就是真正的.
可以探究到人心靈.
這裡就特別提到.
好像將一樣東西極端化了.

$^{361}$就是那個人已經是這樣的人為人.
就是一個善人.
或者一個義人.
他是怎樣在神眼中的義呢?.
原來他的思想是很公平的.
即是他人很正直的.
他的人正直到一個地步就是.
聖經說的正直到一個地步.
就是守樂賢.
是一個在神眼中的義人.
但是惡人他的計謀就是詭詐.
即是如果經常想辦法的.
每樣東西都計算著人的.
不知不覺自己變成一個詭詐的惡人.
所以我們要很小心.
不要說一個惡念.
對別人應該要常常待人以善.
那你自己會變成善人的.
你自己會變稱義人的.
不知不覺你自己變成一個公平的人.
正直的人.
守樂賢的人.
這些不單是重要.
而是你變成怎樣的人.
這個是真的很重要.
很快又看看第17節.
說出真話的顯明公義.
作假見證的顯出詭詐.
又做了一個明顯的對比.
一個分野.
就是常常都講真話的人.
這個都簡單.
我不記得哪個明星講過.
講大話就很難.
因為你時時都要想一下自己有沒有講大話.
講真話就相對容易.
因為你不需要每次修飾.
不需要每次去想一下.
所以我們能夠不斷講真話.
似乎聖經不單是舊約.

$^{401}$新約都是這樣教導我們.
說謊言的嘴都是神所恨惡的.
就顯明我們是公義.
裡面是公義和真理.
主耶穌是誠信真實的主.
這裡說作假見證.
顯出詭詐.
我們就要留心.
我們去到甚麼情況.
有時要這樣作假見證.
這個真的要小心.
神好幾段的經文對我們的教導.
另外第20節.
說謊言的嘴就為耶和華所憎惡.
相信不是偶然間的誣塞.
是習慣性的講大話.
行事誠實的就為他所喜悅.
你可以想像得到.
背後內心那種心靈的動機.
和心靈的取向.
是會直接導致我們變成甚麼樣的人.
這個我們真的要留心.
我稍為很快揭一揭.
《箴言》剛才講的那段.
好像是22章.
22篇 我忘記了是22還是24.
給我一點時間.
因為突然想起這次聖經都重要.
相當重要的.
那些不是那麼穩定.
看看是不是24 可能是24篇.
我又不太記得了.
不然我今晚放給大家聽.
因為它這個排版又有些不同了.
但是它的意思就是我想表達.
剛才《箴言》12章幾句很有哲理的句子.
《箴言》其實就是說.
在神國度底下的期望的.
它的選民應該有的那個表達.
當然不是絕對.

$^{441}$但是這裡就告訴我們.
我暫時找不到.
找到我就再告訴大家.
今晚我可以回去找給大家知道.
但是它的意思就是.
你不要跟那些貪好美食的人一起.
跟他來往.
但是它的意思就是.
我久沒有再看這一段.
所以我暫時找不到.
不要緊的.
我今晚會找到的.
我告訴大家.
因為那次聖經是很可圈可點的.
告訴我們.
你要擔心你自己成為什麼樣的人.
這個很重要的.
你成為什麼樣的人呢?.
就是你平日是怎樣的.
《箴言》就教我們要小心這樣的人.
不要跟這些人來往.
你不要貪他的美食.
就是那些惡人.
因為他心怎樣思量.
我記得了.
因為他心怎樣思量.
他的為人就是這樣.
就是這個意思.
就是說原來他經常想些惡人惡事.
他自然會變成惡人的了.
好像剛才《箴言》所說.
所以我們回頭看這個鬼詐.
那就真的值得我們去想了.
整個社會已經變成一個鬼詐之都.
我們不斷在說我們是魅力之都.
我們是不夜城的大都會.
我們比如紐約,倫敦這些地方.
是國際的大都會.
是怎樣來到一些金融的地方.
比如杜拜是怎樣的美好的地方.

$^{481}$我們比較少去看聖經.
原來耶路撒冷在馬拉基的時代.
他們的宗教信仰已經普遍是一種很低沉.
但是找什麼原因找不到原因.
神說我告訴你真正的原因.
不過是不好聽的.
這個耶路撒冷已經變成鬼詐之都.
充滿了彼此的你如我詐.
充滿了大家是一種仇恨.
是一種大家沒有了互信帶來的後果.
做生意不是一落千金.
一握手成交就真的成交了.
是不會再有這件事發生.
你可以想像得到簽多少個約都沒有保證.
講明45分鐘出場都不出場.
握手都不握.
繞過特首去拿獎.
連獎都不拿.
是不是沒有了那個誠信呢?.
沒有了大家的諾言呢?.
是不是一些基本的人和人之間應該有的互信.
又沒有了呢?.
這裡就告訴我們.
剛才第一個段就已經在講.
整個猶大人行事是鬼詐.
整個以色列大家不再認定我們都是神的選民.
出埃及記的19章.
神和他們立約去到西乃山.
神就和摩西說.
我好像鷹一樣.
將你們帶到這一個地方成立這個國度.
所以以色列人就要歸我成為聖潔的族類.
是屬神的子民.
是有君尊的祭司.
這幾個身份是神賦予給以色列.
其實他們還是赤腳.
無衣的難民.
逃難離開埃及.
所以怎麼會有這麼尊貴的神的選民.
又是聖潔的族類.

$^{521}$又是君尊的祭司呢?.
你可以想像得到.
由那一刻開始.
他們就要學習做一個好像王子這樣的身份.
在神的面前.
他們要有對神,對人的真理,那個相待.
一直透過很多的律法去教導他們.
但是現在來到這個時候.
問題的根源不單止是他們憲制孖夫.
原來更甚的是.
他們為什麼會憲制孖夫的原因.
原來就是已經大家用一種詭詐.
去對待不同的人事物.
包括連對神都是詭詐.
所以帶來只會那個國家.
你想想六天的生活都帶著恐懼.
帶著彼此的不信任.
來到生活.
你很難期望在安息日.
他們去敬拜,去憲制的時候.
就會對神,對人都會有一份很大的雀躍,期待.
和彼此的相交.
來到有一個在信任,信心裡面.
來建立一份的相愛.
所以這裡就很顯示到.
宗教生活原來是有一個息息相關.
和我們的生活.
所以這裡告訴我們.
忘記了我們都是有一位父.
就是這一位天上的父.
這裡告訴我們.
對不起.
可能這樣說未必最對.
因為猶太人舊約的時代.
不敢稱神為父.
所以似乎是在說亞伯拉罕.
這一位眾人的父.
然後再加上不單止是亞伯拉罕.
還有就是我們都是同一位神所造.
就是阿當的後裔.

$^{561}$就是一路也是羅亞的後人.
也是亞伯拉罕的子孫.
一路承繼著神的約.
就是說他們是立約的子民.
在約翰福音.
猶太人經常這樣說耶穌.
我們是亞伯拉罕的後裔.
我們有摩西的約.
你是誰.
這樣罵耶穌的.
應該是第八章的約翰福音.
昨晚我們有機會看過的.
耶穌是神的兒子的那段經文.
這裡就告訴我們.
忘記了自己的本源.
忘記了自己的出身.
忘記了自己本來是一個遊牧民族.
是被埃及奴役的奴隸.
忘記了能夠入迦南美地.
是神的恩賜來的.
是神賜給我們的祖先.
我們在那裡寄居.
所以逾越節不要忘記你是奴隸.
是這樣的一個節日.
所以你們要趕緊.
那些家長守逾越節.
如果有機會去以色列.
看看那些猶太人守逾越節.
他們是要穿上那些衣服.
拿起一支棍子.
吃那個無教的羔羊.
要吃苦菜.
準備隨時離開埃及.
不要忘記自己是寄居.
駐彭寺又是另外一個.
不要忘記你自己住過帳篷.
當年在那個曠野和在埃及寄居.
其實這樣提醒猶太人.
這些都有它那個屬靈的含義的.
當然當沒有了那個敬拜的焦點.

$^{601}$就是對神的敬拜的實質.
自然那些就變成一種儀式了.
變成一個傳統了.
是啊 那個節日就為了做而做.
為了守那個節日而守節日.
已經慢慢淡忘了最重要的那個中心點.
就是自己和神之間的那個關係.
所以這個又是回過頭來看.
我們有沒有忘記我們那個本源.
我們大家都是由信主那裡入門.
大家能夠在教會.
雖然很多都是街坊.
因為你是一個社區教會.
不是在旺角那些.
大家都來自不同的區.
很難的 但是無論如何.
應該有一份親上加親.
既然就算未上主之前都已經認識了.
現在在主裡面.
在救恩裡面又再有這一份深度的認識.
每個星期都能夠走在一起聚會.
甚至和一個團契.
豈不是應該更加有那個默契.
更加有那份的相愛呢?.
不是從一個完全陌生的開始.
豈不是應該更加的相親相愛呢?.
大家都是來自那個本源.
如果是一個陌生人走在一起.
尚且都可以建立一份在主裡面的友誼.
更何況已經本身是認識的.
那為什麼有時候會去到一個.
互相沒有了那個互信呢?.
這個就是似乎作者透過這樣的描述.
告訴我們知道.
原來大家都忘記了.
自己都是神的兒女這個身份.
大家都是亞伯拉罕的子孫.
如果不是神的揀選.
其實沒有一個人可以得到神的各樣的Blessing.
Blessing可以包括你的與生俱來的恩賜.

$^{641}$你的成長的背景.
你的父母對你的孕育.
你在很多的事情上面.
神的保守.
如果我們都忘記了.
如果我們都覺得這些是必然的.
是你應份得的.
那有時候就會去到很盡了.
什麼都是你自己作為那個單位.
就忘記了一切都是領受.
一切都是恩典.
這個是都挺重要的.
我又講一下神學院.
因為通常畢業班.
院長都會特別和我們畢業班的同學講話.
當然是以院長和一位.
大家都很熟的老師這樣講.
很深刻有一堂在組會講的.
不要忘記你能夠讀神學.
是有很多人的恩典保送你進來的.
至少有神學院.
至少有教會的牧者簽那個同意書.
有長執會.
有可能你的團體弟兄姊妹的支持.
關心,代討.
我都提過.
我進神學院的時候有15封個人推薦信.
是不會有教會推薦我.
因為很特殊的原因.
我是有15封個人推薦信.
還要在建道4年的神學.
弟兄姊妹幫我設了一個基金.
只是可以用來交學費的.
不可以用的了.
4年的學費十幾萬.
是在那個pool裡面.
有弟兄姊妹來幫我交.
從來都沒有試過.
會有一個學生有個基金是為他.
是為我來有這樣的安排.

$^{681}$而且我是不知道哪個人.
奉獻在那個基金裡面的.
去到最後那一年第四年.
因為和那些建道的行政同工同事都很熟.
因為來不及就要進去.
那個叫鄧小姐已經退休了.
她說你畢業了.
我就將過去的名單給你看.
不要忘記別人的恩情.
很多當然是我認識的.
我舊教會的弟兄姊妹.
院長講了一堂道.
我很感動.
就是叫你們不要忘記.
你今天能夠讀神學.
是有很多人對你的那個blessing.
是因為信任你是神揀選.
所以來讀神學.
這個是很重要的.
我們不要忘本.
這裡就已經提了鬼詐的原來的原因.
慢慢忘記了自己.
猶太人神選民這個這麼寶貴的身份.
所以就很自然地看自己是最大.
要得到最多最好的利益.
而忘記了一切都是神所賜給我們的.
我們就忘記了就是.
沒有了那份對神的感恩.
自然就會對人有很高的要求.
對自己就沒有了一個太強烈的問責.
所以我們各人為什麼要鬼詐地待弟兄.
背棄了神和我們列祖所立的約呢?.
立什麼約呢?.
到底這位先知馬拉基是在講什麼約呢?.
是神在加勒底的吾爾呼召亞伯蘭.
離開他的本族 傅迦.
去到那個應許之地.
是不是那次的立約呢?.
還是亞伯拉罕奉獻以撒的時候.
他晚年的時候.

$^{721}$某程度上他的成熟信心高峰的時候.
神試驗他都是合格的.
那一次而來到立約.
論福我就會給大福你.
因為你沒有保留的.
將你的兒子以撒獻在壇上.
到底是指哪個的約呢?.
列祖是指摩西還是羅亞呢?.
本來羅亞入方舟.
沒有人可以生存的.
但是唯有羅亞一家八口.
可以來到他在方舟裡面.
逃避那個洪水.
這些都是恩典的約.
彩虹之約.
還有其他的.
所以到底是在說諸約.
這麼多的歷代歷史裡面的約.
還是特指哪一個祖宗呢?.
這就不是很詳細.
但是很奈何蓋炭的就是.
嘆息的就是.
猶大人他們已經在民間的日常生活.
已經是以自己的利益為最大.
並且在以色列和耶路撒冷中.
行一件可憎的事.
因為猶大人褻瀆耶和華所喜愛的聖地.
就是取事奉愛邦神的女子為妻.
只是求方便,快捷.
愛邦女子他們有的就是.
可以透過祂去求問他們國家的神.
而得到耶和華未必給到我們的福氣.
這個一定是猶大人那個民間裡面.
他們這樣來做的事情.
就是說我的心不會離開神的.
但是我的行動是會走去.
跟隨這些太太團去拜他們的神的.
所以你可以想像到猶大人他們的行事詭詐.
他們不敢全然否定耶和華的.
但是在他們實際的行動.

$^{761}$就行一件神所不喜歡,可憎的事.
就是他們來褻瀆這個聖地.
所以褻瀆這個聖地都可能適合那個聖潔.
神期望的就是百姓.
仍然生活在一個很原始的狀態底下.
要仰望神的公認,那個恩典.
還有在困難裡面.
例如失修,乾旱,饑荒,有蝗蟲瘟疫.
他們要回到聖殿.
回到所羅門的聖殿去祈禱.
那裡特別強的WiFi 6G.
他們可以在那裡去痛悔,悔改.
神就會賜下那個瑤書.
是歷代志下七章那段經文的精髓.
只要那些百姓仰望神.
就會得到那個重新的建立.
但是現在就有一個快捷的途徑.
就是表面我仍然敬拜神.
但是實際的行動.
我用腳去投票就是我透過其它的民族.
他們拜的假神.
其他人拜了不關我的事.
我跟它去而已 我其實不贊成的.
表面就和那些可能祭司.
可能就是一些領袖這樣講.
但是實際上神就說.
這些是和外邦的女子.
來尋求他們的神.
所以第十二節.
他說凡是行者事的人.
無論何人.
就是獻供物給萬君之耶和華.
就是兩手準備.
就是既來守這個節期.
來到聖殿.
但是實際上他們的行徑.
更多就是在那些偶像的敬拜裡面.
心不能專一對神.
所以這個詭詐.
可以把它來到去展示出來.

$^{801}$攤開的心靈.
你就會看到.
原來在實際的日常的生活裡面.
對神根本不是那麼專一.
那麼專注尋求神的喜悅.
這個特別是所講到的喜愛.
神所喜愛的事.
就是說敬拜不會令到六天的生活.
帶來有比較大的顛覆性的改變.
在生活裡面所做的事情.
可能就是做和信仰.
剛剛有抵觸的事情.
有時都不放在心裡面.
人都是容易善忘.
因為去到星期日準備穿衣回教會.
就忘記了六天做過的甚麼事.
都不會去刻意先做一個反思.
有一個沉澱.
和有一個心靈的悔改.
然後去到禮拜堂.
似乎猶大人已經是這樣習以為常的行事了.
所以你問他哪件事行事詭詐呢?.
他都一下子未必回答到你的.
因為他已經成為了他的日常.
他的生活的那種方法,態度.
以致做事一來都是以自己為中心.
很容易就會帶來這種情況.
所以可能開車開得很快.
過馬路紅燈照衝.
習慣性的遲到早退.
很多的生活裡面 去到他自己不開心.
遲到 去到餐廳見到那個夥計不喜歡.
又掌摑著人.
因為你是消費者 你覺得你應份的.
不知不覺就在那個日常生活的每一個態度上.
都是行事詭詐.
行事詭詐未必是惡形惡相.
可能是很斯文的 不過是敗類.
來對人又給一些難聽的話面色.
諸如此類都會有的.

$^{841}$所以這裡告訴我們.
其實一切都是回到宗教的根源.
那個使者不存在那個民族當中.
達至不到 其實帶領那些回眾去經歷神.
去對神有認識 是一種心靈的認識.
那個認識是會帶來生活有改變.
做不到的時候 百姓就只會靠自己.
因為要生存的 要生活的.
那就要靠自己去生活.
在生活裡 自然就是人與人之間.
我當然要為自己作保護.
我就要尋求最大的個人的利益.
所以猶大人就行事詭詐.
還有再做了一件可憎的事.
就是表面是敬拜神的 是不會不要的.
不敢的 但是實際的生活.
與神完全是不感覺需要向祂交代的.
這個是我們在《新約》裡面.
為什麼書信裡面的教義.
這麼多教我們在生活裡面.
都揭一揭 大家可能會好一些的了解.
其實這些是不斷的周而復始.
有教的講道是幾熱門的題目.
比如《而不所書》的四章.
你可以說保羅真的當那些.
愛邦信主的弟兄姊妹.
好像當他們不懂得做人一樣.
就教他們怎麼做人.
四章的二十三節.
他說又要將你們的心智要改換一新.
心智就是你們要心意更新而變化的意思.
並且要穿上新人.
不是新年穿新衣服這麼簡單.
是你當自己是一個新做的人.
果然真的下面這樣講的.
就是新人是照著神的形象做的.
這句話豈不是就是阿當呢?.
照著神的形象,樣式做的就是阿當.
現在我們就是新的阿當.
是新的人.

$^{881}$有真理的人意和聖潔.
就是說這個新人的穿著.
就是這個真理,人意,聖潔.
接著講了那個大原則之後.
就教你怎麼實際生活.
所以你們是新人了.
你們是神的創造了.
你們就要棄絕方言.
就馬上講鬼詐的問題.
「各人與鄰舍要說實話」.
又是人與人之間.
「由大人行事鬼詐」.
好像剛剛相反就是.
「各人與鄰舍要說實話」.
鄰舍就不是你隔壁屋.
B號門牌那樣.
是在講每一個都是我們的鄰舍.
因為我們是互相為肢體.
互相為肢體.
我們是神的揀選.
我們是在基督的身體裡面.
我們是互相為不同的Body.
其中一個一員.
接著就提到一些生氣.
你不可以犯罪.
不可以含怒到日落.
就是說永遠的搖樹的意思.
不是叫你看著那個標.
看著太陽下山.
你不可以讓魔鬼留地步.
到底是上文講的仇恨.
有可能是魔鬼是雪荒之人的父.
魔鬼是製造紛爭.
製造人與人之間的敵視.
《聖經新約》有一句挺重的.
提摩塔前書.
說末世有危險的日子.
就是不解怨.
什麼叫不解怨.
不解除冤仇.

$^{921}$就是說我恨你 我恨死你.
我這一生都不會原諒你.
原來是真的.
《聖經》說不解怨.
末世的人是不解怨的.
其實是很恐怖的.
這件事是令到很多人一生活得不開心.
捆綁著自己.
其實是很蠢的.
那個心靈浪費時間去生氣.
最慘是被你生氣的人.
根本從來都不會再介意你說你生什麼氣.
那你自己坐在家裡不就不斷生氣.
其實就真的很慘.
原來是出於魔鬼.
所以生氣不要犯罪.
有時候真的會發脾氣的.
很激氣的 很生氣的.
但是要馬上原.
原就要原.
大家都已經去到這個階段年紀.
說話要非常小心.
因為你說錯一句話.
其實你要收回是沒有辦法的.
收不到了.
甚至乎那些關係會破裂的.
所以要小心說每一句話.
很重要的.
這些說不可以含怒到日落.
很肯定就是在說魔鬼會製造這些糾紛.
會製造這些疑團.
會製造人與人之間的敵視.
大家是陌路人.
大家都不用同心似鳳.
他最喜歡這樣見到的了.
接著真的好像抓住手當他們是小朋友.
好像小學那樣.
不准喧嘩.
不准小息 亂跑.
真的一句句的.

$^{961}$好像那些班房裡面有那些日表的.
寫著什麼什麼的.
類似這樣.
從前偷東西的.
不用你們舉手了.
但是不要再偷了.
真的以前有賊仔的.
接著說總要勞力.
親手做正經事.
那你就有餘了.
可以分給缺少的人.
不單止去貪念.
去偷.
想盡辦法.
現在是想辦法去幫人.
用你的金錢.
所以金錢是很中性的.
用你的金錢.
去幫有需要的人.
所以我終於明白.
為什麼那些明星.
這麼有錢都要去回自己.
最應該得到的身價.
就是他可以有錢.
他可以做善事.
當然不是每個人.
我想這個都是重要的.
我們用盡辦法去賺錢.
但是原來是想辦法.
怎樣去幫人.
就可以有餘.
分給缺少的人.
就是說他以前是一個.
哇 夜媽媽到處去.
看誰可以偷東西.
現在不需要了.
他可以日光日白.
很理直氣壯.
抬起頭去做人.
甚至乎他可以洗底.

$^{1001}$可以幫助有需要的人.
你可以想像得到.
他多麼的開心.
那個人因為信了主耶穌.
帶來他的身份是多麼的被抬高.
他的家人是多麼的開心.
哇 我的先生現在不再偷東西了.
如果你認識那些吸毒的人.
賭錢那些.
你就知道.
可想而知.
信仰帶給那個人的生命改變.
是抬舉他的榮耀.
他有了那份尊嚴.
跟著那些污穢的言語.
果然真是一句都不可以出口.
只要隨時能夠講一些造就人的好話.
叫聽到的人得到益處.
造就人不是阿諛奉承.
不是講一些反話.
或者是講一些巴結人的話.
要小心的.
因為這樣的人.
由腔滑調.
有時內心是很不舒服的.
我們敬虔的人.
應該是講一些正直的話.
而且真是從對方去講.
所以這是很高的藝術.
看你用什麼角度去看.
那件事你可以從一個很正面的角度去看.
所以其實真是隨時.
你都可以講到一些造就人的話.
說話的藝術都是很重要.
不要叫神的聖靈擔憂.
你們原是受了祂的印記.
第一章已經出現過了.
二不所思.
我們是有聖靈作為我們的訂金.
我們是有將來天堂的進入.

$^{1041}$聖靈就是我們的訂金.
神下了大訂.
就說我們將來一定會得到那個基業.
應該是一章二十節.
二不所思.
對不起.
應該是十四節.
聖靈就是我們得到基業的訂金.
神已經下了大訂.
在我們將來得到基業的指紋.
所以現在這裡再提一次.
我們是曾經有聖靈作為印記的指紋.
所以我們是等候得贖的日子來到.
所以不要叫神的聖靈為我們擔憂.
我們要敏銳聖靈的工作在我們的內心.
就是上文講的那些東西.
怎樣做人.
信了主之後怎樣做人.
其實剛剛和剛才講的.
就是馬拉基書所講的.
就是他們沒有了信仰之後.
在人和人之間的生活裡面.
就用鬼詐去互相對待.
所以我想有生之年.
教會都會有這樣的情況出現.
都會有弟兄姊妹.
因為他們已經離開了基本的靈修生活.
祈禱聚會生活.
在聚會裡面得不到教導.
真是不好意思.
我們作為牧者在台上.
有時都看到人生百態.
來聚會這麼寶貴的一個小時.
坐下來都是在玩手機.
你回家玩吧.
不用來了.
你玩手機來做甚麼呢?.
聽到看著手機在玩.
看其他東西.
何苦回來教會呢?.

$^{1081}$你回來教會是想著.
令自己安心一點.
我來了打了卡.
我又出席了.
出席已經寶貴了.
不是不寶貴的.
真正的重點是.
在神的裡面.
令我們內心得到聖潔.
得到照明.
得到我們行事的有力.
做一個正直人.
這裡就是猶大人.
因為已經沒有了那個信仰.
帶給他們生命的更新.
他們就要下起.
和外邦人沒有分別.
即是其他的民族.
迦南人沒有分別.
就是行事詭詐.
雖然這裡沒有很具體.
講甚麼的事情詭詐.
但是新約.
保羅反而就是.
你們信了主耶穌.
不可以再做一個詭詐的人了.
你們的生活應該要怎樣呢?.
很仔細的教導.
即是保羅寫信的時候.
可能想起哪個弟兄.
哪個阿貴呢?.
以前偷東西的.
現在不要再偷了.
這樣教導.
你覺得挺好笑.
原來很實際.
剛才講了原則.
我們都是新人.
我們是神的形象所造的新人.
即是我們是新的阿當.

$^{1121}$我們就應該穿著這個新人.
就是真理,仁義,聖潔.
穿著完之後.
那不就是個個都曉飛了.
原來不是.
你們個個要由頭開始.
所以要上初信班.
所以要上成長班.
所以要不斷聽聖經的真理.
因為你們信了主耶穌之後.
就要怎樣了.
原來是很實際的教導.
所以來到這裡的時候.
神說要在雅各的帳篷當中.
當中的教導.
教導你們.
在雅各的帳篷當中剪除.
表面是獻供物.
但實際上.
生活裡面.
都是以鬼詐去彼此相待.
這樣的.
吾敬虔的信仰生活.
都是需要求神.
讓我們能夠不斷的去更新.
這個是很重要的.
最近有機會看.
清教徒的書.
其中相當重要的一個清教徒.
Richard Baster.
是很出名的.
非常重要的一個清教徒.
代表人物.
他講他自己的生命改變.
原來他是非常體弱多病的.
這位牧者.
很多病痛的.
病到一個地步.
經常都起不了床.
很弱.

$^{1161}$他很不開心.
但由他開始.
突然間有一個醒悟.
他在聖經裡面.
關於天堂的經文.
他用來做靈修.
他就說.
他自己好像置身在神的面前.
令到他自己.
越來越有那種心裡的力量.
他這樣的一個舉動.
後來就成為了.
他的牧區.
他的時代.
的一個大復興.
很多人因為他的講道.
而得到生命的改變.
就是這個Richard Baster.
BASTER.
我不知道中文應該怎樣譯.
但很出名的.
他的書在神學院一定有.
香港暫時難買到.
因為很年代久遠.
幾百年前的人.
清教徒的歷史.
5月花號.
由英國去美國的.
就是清教徒.
就是這一班人.
帶來美國的大復興.
就是他們這班人.
你可以想像得到.
原來這樣去思想.
聖經的真理.
自己的生命都會有改變.
所以盼望我們都能夠.
真是有一個認真的.
對神的追求.
我今天就給了一個PowerPoint.

$^{1201}$有些弟兄姊妹.
或者有些群體.
就是這個星期日.
下午在紅恩堂.
的一個意象分享.
和今年的事公的方向.
我神學院的同學.
都很熟的.
我給他們參考有什麼意見.
他給了一個很寶貴的意見.
就是說.
一切的事公.
最終是要帶領回眾.
去認識神.
敬拜神.
不是為了玩而玩.
不是為了開心而開心.
因為地下的地舖.
將會更新.
可能會招聚很多人來聚集.
最終是要帶人去認識.
敬拜神.
教會是一個敬拜神的地方.
這是一個很寶貴的提醒.
所以我給了一個.
很寶貴的提醒.
所以這裡告訴我們.
神怎樣看敬拜.
原來民間的生活.
他們這麼多的混亂.
居然沒有將他們的混亂.
帶回神那裡.
由神去幫他們.
種財.
由神去改變他們的心靈.
改變不了別人 先改變你自己.
去改變自己的心靈.
去做 啊 不行啊.
就帶來.
都欺控了.

$^{1241}$對神都詭詐.
在敬拜的生活裡.
都詭詐.
來打卡就算了.
當作是敬拜.
不是真正的朝著一個.
朝聖的心靈.
去敬拜神.
很難得能夠.
安息日去到獻祭.
獻的就是.
祭司一次兩次.
都不說我.
下次.
每個人都獻那些.
瘸腿的 暗病的.
阮明就說.
我拖一隻無殘疾的來.
又貴.
又要小心保養.
誰知道原來去到.
生死的 每個人都是獻那些.
瘸腿的都沒有人理會.
祭司說一樣可以的 沒問題的.
都是這樣燒的.
下次就不是這樣了.
慢慢慢慢就大家都放下了.
對神就.
不知不覺就落在一種.
起初就是無知的一種詭詐.
後來就是慢慢都已經.
神都不追究.
神都不會突然間說我們這個民族.
我們這樣做是不對的.
因為祭司不再說是對的.
和不對的.
祭司已經不判斷你這樣做不行的.
你這隻是不准獻的 你要快點拿走.
神是不要你的.
不稀罕你的 不敢說.

$^{1281}$祭司怕別人不獻祭.
所以不敢說.
慢慢就導致.
就導致大家都不獻祭.
大家的獻祭都是很馬虎.
很隨便.
這樣對神詭詐.
對人就一定會.
連神我都不怕了.
我還怕人.
自然人和人之間.
已經忘記了自己的本源.
所以我們看第13節.
連家庭生活都出了問題.
又再說另外一次的詭詐.
這裡告訴我們.
14,15節就有了.
你們又行了一件這樣的事.
使得前妻嘆息哭泣的眼淚.
遮蓋耶和華的壇.
以致耶和華不再看顧那貢物.
也不樂意從你們的手中收納.
那些不肖的財產.
你們又行了一件這樣的事.
使得前妻嘆息哭泣的眼淚.
遮蓋耶和華的壇.
以致耶和華不再看顧那貢物.
也不樂意從你們的手中收納.
這裡說的是.
前妻.
即是說娶了外邦的女子為妻.
使得前妻.
哭泣的眼淚.
遮蓋神的壇.
當然是誇張的.
女士.
根本不可以.
光天化日的走去聖殿.
來獻祭.
大家都知道.

$^{1321}$聖殿分開幾個地方.
有外邦人員.
即是最遠的.
好像我們在大球場看球賽.
坐在山邊的頂層.
那些一百多元一張票就很高.
最近的.
握到手的那些.
和那些明球星.
就九百元一張.
你可以想像到.
女子有一個叫女院.
是女孩子才可以去.
我們讀神學的時候.
去做猶太教的訪問.
去過猶太的.
在中半山的.
猶太教的廟.
即是那個聖殿.
去參觀.
要預約的.
去參觀.
女同學是坐在樓上的.
女孩子.
樓上要蒙頭的.
那些男士.
拉比那些猶太人.
在正中.
好像擂台一樣.
擂台.
是一個讀經的地方.
是站在最高最當眼的正中.
你可以幻想禮堂.
中間有一個高了.
幾級的.
裡面有一個四方的.
地方.
好像擂台一樣.
讀經員在上面讀經.
男孩子才可以讀經.

$^{1361}$女孩子不能讀經.
會將那些卷軸拿出來.
有一個致勝鎖.
將那些卷軸拿出來.
真的是手抄的卷軸.
在那裡讀經.
所以你可以想像.
怎可能會有眼淚流到.
耶和華的壇.
和同祭壇有多大.
我忘記了.
那個尺寸有定了.
那個同祭壇.
是所有的.
聖殿裡面的器皿.
神的器皿.
最大的.
那個volume.
最大的一個物件.
就是同祭壇.
因為我們獻祭的.
燒那些祭身.
所以怎可能.
猶大的女子的哭泣.
的眼淚.
可以遮蓋整個祭壇.
你可以想像到.
在一個男系的社會.
男性主導的社會.
尤其是猶太人.
有時哭泣.
可以去到祭壇.
家醜都不出外傳.
怎可能會這樣呢.
很誇張地.
神知道每個猶大的男子.
原來表面去獻祭.
實際上家裡就家暴.
虐待自己的太太.
太太在家裡.

$^{1401}$躲起來哭泣.
那些哭聲.
原來就去到神的祭壇.
獻的不是那些乾淨的祭身的血.
原來是.
一路獻的時候.
原來太太的眼淚.
已經去到神的面前.
就不是輕香的祭了.
所以你看到.
多麼的全神.
以至耶和華不再看顧那些貢物.
也都不會樂意從你們的手中收納.
不要 不收.
神根本很討厭這些祭物.
所以你看到.
多麼全神的講到家庭的關係.
如果以男士作為一家之主.
來獻祭.
譬如你看過《撒慕爾記》.
阿爾尼加拿和太太哈娜.
還有比利娜一起獻祭.
你會感受到哈娜那種哭泣.
那種悲哀.
那個苦情.
你可以想像到.
這個就是去到馬拉基的年代.
神好像想告訴百姓.
不是再是外敵入侵.
不是再是你被擄回來.
重建遇到的困難.
那些困難都容易克服.
更大的困難就是.
心靈已經慢慢對神沒有了敬虔.
對神已經沒有了那種雀躍.
沒有了那種信仰帶給你的生命力.
這個就是最大的困難.
所以復興必須是由祈禱屈膝開始.
不是人多就等於復興.
當然要小心說這句話.

$^{1441}$小心聽.
不然會被人罵死.
那怎樣才是復興呢?.
幾個人就復興嗎?.
不要將兩樣東西極端化.
人多是一個復興.
但是如果那班人不是心對準神.
那班人原來都不是真正愛神敬虔的.
其實那個多只不過是一種虛火.
不是真正的能夠長存.
更加重要的是在神的眼中.
怎樣看待這些人多呢?.
這是重要的.
所以我們不是否定人多人少.
而是告訴我們到底.
來到祭壇最神聖的獻祭.
神聞到的是輕香的氣.
還是聞到那些哭泣的慘叫聲呢?.
還是那些姐妹們在家中的冤屈聲音呢?.
這個其實真的很傳神.
真的反映到我們的生活.
是直接和我們的敬拜有關連的.
你想想你叫完大交再回崇拜.
如果你還要做雙禮.
那就更加搞笑了.
你剛剛才吵完架.
自己走得這麼快.
也不理太太也好 丈夫也好.
要趕回去做雙禮.
其實自己還沒有平復.
要擠笑容出來和別人雙禮.
這些長此下去其實是不健康的心靈.
不過你看看這件事.
神不收納是最大的問責.
神不願意看顧這些的共物.
第十四節還有那些托詞.
你們還說這是為什麼呢?.
作者這樣的一句.
好像萬不驚心的一句話.
將百姓那種對神的詭詐.

$^{1481}$對神的那種的膚險.
有獻祭.
但實際上的生活.
也是和自己的家人.
最親的人都以詭詐相待.
對外人也是詭詐相待.
以致到神都已經不收納了.
接著就說我已經獻了.
還是什麼事啊?.
什麼事啊?為什麼啊?.
我不是沒有獻啊.
這樣的情況.
接著神就說因耶和華在你.
和你幼年所娶的妻子中間作見證.
她雖然是你的配偶.
又是你盟約的妻子.
你卻以詭詐待她.
就已經說了那個根源了.
就是想盡辦法.
由於你內心對她已經沒有那種喜悅.
那種的羨慕.
來慢慢的看她是一個障礙.
你就用詭詐去對她.
逼令她離開這段盟約的婚姻.
使你可以理直氣壯的.
用一個堂而皇之的理由.
去娶外邦的女子為妻.
所以摩西是准猶太人休妻.
站一站19章.
講到婚姻這個問題.
馬太福音19章.
那就肯定不會錯了.
因為這段經文太重要了.
19章的馬太福音.
耶和華說.
耶穌說完了.
姐姐說就離開加利利.
去到猶太的境界約旦河外.
有許多人跟著她.
她就在那裡把他們的病人治好.

$^{1521}$有花里賽人來試探耶穌.
就說人無論什麼緣故都可以休妻嗎?.
耶穌回答說.
那起初做人的是做男做女.
《創世記》一章27節.
並且說因此人要離開父母.
與妻子聯合.
二人成為一體.
《創世記》第二章有記載.
這經.
《創世記》是經典的經.
這經你們沒有念過嗎?.
不是普通書那麼簡單.
因為耶穌看摩西五經是神的默示.
這經你們沒有念過嗎?.
不是特意要像德這些花里賽人.
而是你們很熟悉這些經典.
為什麼你問我這個問題呢?.
當然他們是振振有詞.
因為摩西准人分開和太太.
第六節.
既然如此.
夫妻不再是兩個人乃是一體.
所以神配合的人不可分開.
第七節.
花里賽人提出質詢.
為什麼摩西可以吩咐人休妻?.
寫休書就可以.
寫休書即是有黑白.
要具體列一些罪狀.
讓拉比去幫你做一個訴訟.
她煮東西不好吃 太鹹.
早上起床不叫我.
又不幫我倒水洗腳.
任何一個原因都是休書.
給一個理由.
不是理由也可以給.
總之給一個理由.
你認為她是值得你休的.
你給一個理由就可以了.

$^{1561}$所以告訴我們.
摩西為什麼吩咐給妻子休書.
就可以休她呢?.
是出自哪一段經文呢?.
《申命記》24章.
不如我們先看《申命記》.
我們先看《耳經解經》.
因為是花里賽人引的.
真是有聖經的.
24章.
看看所謂的alternatives.
在什麼狀態下可以休書呢?.
24章第一節.
人如果娶妻以後.
見她有什麼不合理的事.
不再喜悅她.
可以寫休書交在她手中.
打發她離開夫家.
婦人離開夫家之後.
可以去嫁別人.
後夫若恨惡她.
寫休書交在她手中.
打發她離開夫家.
或是娶她為妻的後夫死了.
打發她去的前夫.
就不可以在婦人.
招污之後再娶她為妻.
再回到自己的前夫那裡.
因為這是耶和華所憎惡的.
不可使耶和華理神.
所賜為業之地被玷污了.
開了一個很大的空間.
人如果娶了太太之後.
見她有什麼不合理.
這個不合理怎麼去詮釋呢?.
是誰的不合理呢?.
是對丈夫的不合理呢?.
還是對社會的不合理呢?.
是有不同的.
那你應該怎樣呢?.

$^{1601}$誰認為那個標準在哪裡?.
那個尺在哪裡怎麼量度?.
沒有說的.
其實拉比就有.
摩西這一節聖經有解釋的.
正正就是.
剛才我說的那些好像很搞笑的東西.
拉比就有詮釋了.
就是說原來.
你給了一個標準.
你認為不合理的事.
你就可以寫休書休她.
拉比就解釋了.
休書的內容可以是哪些.
可以發展到你覺得那些東西不好吃.
你就可以休她.
可以去到這樣.
所以到底摩西.
為什麼她是神人.
為什麼會制定這樣.
這麼不合理的那些姐妹.
太太的一個法例呢?.
那個時候.
女性是沒有生存的技能的.
她是要依靠.
還沒出嫁就是爸爸.
家庭.
嫁了之後就要跟隨丈夫.
其實是沒有自己的謀生能力的.
為什麼摩西會這樣呢?.
這個就值得我們看看主耶穌怎樣去解釋這段經文.
剛才我們已經講了.
主耶穌是講得很清楚.
那個真正的意思不是摩西的意思.
真正的意思耶穌已經講了.
神所配合的人是不可以分開的.
這個才是真正的神的心意.
然後他們就真的根據生命的24章.
摩西真的這樣講得很清楚.
講了的.

$^{1641}$為什麼呢?.
怎樣去理解呢?.
接著耶穌就解釋了第八節.
所以19章的第八節就很重要了.
因為耶穌幫我們解釋摩西她真正的意思.
摩西是因為你們的心硬.
所以許你們休妻.
而回過頭來就是起初神的心意並不是這樣.
這段聖經就應該是這樣理解的.
摩西因為你們的心硬.
你們是那些男人.
猶大人.
如果好像神的標準.
無論什麼原因一概不准休妻.
原來人性就很邪惡.
他就會虐待自己的太太.
虐待到她死.
你不准我休嘛.
我就虐待她.
我當她是阿四.
我天天這樣虐待她.
不給她好飯吃.
關她進柴房.
要她做一些粗重的事.
打打罵罵.
就令到那個被虐待的.
就都真是很慘了.
所以區聯會都更改了一些守則.
就是如果一個被嚴重虐待的家庭的其中一員.
當然經過有牧者的勸戒的輔導.
最終都發現這件事是改變不了.
都准許男女事大方提出離婚.
區聯會都更改了這個.
正正就是摩西因為你們的心硬.
如果不是呢.
因為當時是男性主導的.
不要用今天這個現代社會去想.
女性都可以有一個主導權.
古代是沒有的.
男性主導如果是鐵板一塊.

$^{1681}$他就會情願虐待到她死.
死了就解除婚姻.
神就不希望看到這樣.
所以摩西也都看到這個問題.
就說了.
至少你都要有合理的理由.
大家去評斷一下你這樣合不合理.
如果你認為合理的.
你都不再喜悅你的太太的.
那是你的心問題.
你自己將來面對神.
那你就休他.
所以摩西這樣說.
主耶穌就說.
起初不是這樣的.
回頭將這個問題放回男性身上.
其實主耶穌真的說得很好.
所以又是回到整件事的初心.
那個最高的設計.
是神所配合.
不可以再分開.
當然是有很多後來的因素.
會令到這件事會分開.
所以有時候又未必需要一個堅持.
第九節.
我告訴你們.
凡優妻另娶.
如果不是為了淫亂的緣故.
就是犯奸淫.
有人娶那個被休的婦人.
也是犯奸淫.
門徒就和耶穌說.
人和妻子.
原來這麼多手續.
這麼多事情要跟的.
既然是這樣.
倒不如不要娶了.
耶穌就說.
這句話不是每個人都可以領受的.
唯獨次給誰.

$^{1721}$即是上天給的恩賜.
有些人是獨生的.
誰才能夠領受的.
因為有生下來.
就有殘缺.
那些可能是不能夠生育.
或者其他.
有些人是有另外的.
譬如被閹.
有為天國的緣故自閹.
說誰能夠領受.
就可以領受.
有些人是為了神的國度福音.
自己手獨生等等.
這個問題.
當我們回到馬拉基書這裡的時候.
其實挺圓滿的.
耶穌說到.
特別是在新約的時代.
但是在舊約的時代.
沒有主耶穌的傳息.
沒有主耶穌的教導和警告.
他們是活在摩西的律法之下.
自然就想盡辦法.
寫休書.
休自己的太太.
其實真正的問題.
是十四節.
二章的十四節.
她是你的配偶.
是你的盟約的太太.
但是你用鬼詐對待她.
這個就是根源了.
再一次.
連自己的婚姻都用鬼詐去相待.
這個就是猶大人.
在這裡出現了四節.
五詞關於鬼詐.
整個民族都已經在鬼詐裡面生活.
鬼詐之刀.

$^{1761}$本來不是很好聽.
但是在神的眼中.
面對這個問題.
願意正視改變.
不是隨便獻祭就可以復興.
而是人心需要回到神的面前.
這裡再一次告訴我們.
十五節這裡很特別.
雖然神有靈的餘力.
似乎是說.
阿當和夏娃.
在阿當的肋骨做夏娃.
神有靈的餘力.
可以做多人.
他不是單做一人嗎?.
為何只做一人呢?.
乃是他願人得虔誠的後裔.
所以當謹守你們的心.
誰也不可以鬼詐.
代幼年所娶的妻.
耶護說:以色列的神說.
幼妻的是和以強暴代妻的人.
都是我所恨惡的.
你會留意到.
剛才我這樣解釋耶穌的教導.
摩西的准許幼妻.
原來與大相關.
你不准他幼妻.
他就強暴對待太太.
家暴.
以前沒有這個term.
但關上門的暴力更加嚴重.
是朝見口晚見面的時候.
太太處於一個很弱勢的家庭中的角色.
是很慘的.
所以我們小時候看的語場片.
每一套都是講家暴的.
娶了後代.
後代李香琴回來.
虐待 白煙.

$^{1801}$往往都是這樣的.
你可以想像到.
十六節原來是孿生兄弟.
幼妻和家暴原來是源於鬼詐.
你可以想像到.
神已經揭示了人與人之間.
是不再是耶路撒冷.
不再是神的聖城.
而是鬼詐之都.
在這裡大家的生活.
可能民間做生意的交易.
已經不再講信用.
已經不再講彼此的信字.
而是大家都可能你娛我詐.
實際去到家庭.
都是這樣想盡辦法.
准我寫休書.
我就寫休書.
你不准我就用強暴.
所以休妻和強暴.
原來是真的會伴隨的.
所以現在香港是一年.
雙方同意就可以解除婚約.
如果不行就多給一年.
我聽過官員說.
他們修訂這個法例.
背後的原因都是這樣.
與其這麼慘又有暴力.
不如再早一點.
所以兩年算是最短的.
以前是不行的.
我忘記那時候是多久了.
你可以想像到.
多離婚率.
某程度是少了強暴,家暴的.
但是神當然不期望.
要用這樣的手段去離婚.
神期望的當然是要對那個文約.
有那個忠誠.
所以第二章如果用一個焦點.

$^{1841}$是沒有了那個守約的諾言.
對神的守約.
就是宣佈民間的百姓.
不再看自己是神的選民.
所以今天的教會.
應該不只是華人教會.
很小的教會.
有沒有看重與神立了這個生命的約.
十字架的救恩帶給我們這個生命改變.
我們要感受到這份幸福.
這份榮耀.
以致在我們的生活裡.
是一個重諾言的人.
對神的諾言.
和人的諾言.
這個都很重要.
揭一揭到最後.
因為九點了.
詩篇24篇.
可能大家很熟的23篇.
24篇大家都應該不熟的了.
但是太重要了.
詩篇24篇.
地和其中所充滿的世界和住在其間的.
都是屬於耶和華的.
哈巴地建立在海上.
安定在大水之上.
想講的當然是講人類.
誰能夠登耶和華的山呢?.
誰能夠站在他的聖所呢?.
就是守潔心清.
不向虛妄.
起誓不懷詭詐的人.
他必蒙耶和華賜福.
又蒙救他的神.
使他誠意.
這是尋求耶和華的族類.
是尋求裡面的雅各.
你看到嗎?.
很清楚.

$^{1881}$誰才能站在神的聖山呢?.
就是守潔心清.
就是我做事光明磊落.
不向虛妄.
不會有時暴露是非.
暴弄是非.
製造那些緊張.
說別人閒話.
去做自己的資本.
然後散播謠言.
然後這句很重要.
起了誓不懷詭詐.
就是說一落千金.
做人要有這樣的誠信.
要不就不答應.
答應了就一定要做.
我們小時候就習慣被人騙.
哥哥姐姐說我們去哪裡玩.
最終原來不去.
說完當沒事.
你跟她追究她還打你.
現在不行了不去.
對小朋友的心靈傷害是很大的.
其中一個後果就是.
小朋友自小就知道原來.
說話不負責任是無所謂的.
這樣就很大件事了.
你就會在一個沒有誠信的家庭裡長大.
你長大了自己也不會有誠信.
真的 你就會慢慢.
不太明白什麼叫做聖潔.
什麼叫做這麼好的品格.
所以有時候你看到一些在屋村.
當然沒有貶義.
那些小朋友或者青少年.
慢慢慢慢去學壞.
其實你就看到.
他只不過是一個可憐的受害者.
因為他的家庭,他的父母.
根本就已經是這樣了.

$^{1921}$所以我們能夠信主.
慢慢透過福音,透過教會的改變.
聖經的教導.
慢慢自己改變.
其實絕對不是什麼家教等等.
什麼都沒有.
是神的恩典.
是很難的.
因為你自小就受了傷害.
所以如果你的父母是這樣的敬虔的家庭.
真的很寶貴.
我承諾了.
起誓我是不懷鬼詐的.
我說了就是了.
對那個小朋友的心靈.
是很大的建立.
我的爸爸媽媽原來是這樣的為人.
對他的長大.
他很自然就有這種正直.
很自然就有這種.
這樣的對神,對人的愛.
其實很感恩我的兒子.
他認為自己和神.
是有這種這樣的建立.
就是在家庭裡面這樣成長.
所以我們.
如果我們自己沒有這種.
自小的土壤.
不要氣餒.
聖經可以幫助我們的.
我們可以在教會裡面彼此建立.
我們就能夠讓我們的下一代.
或者讓我們教的其他的小朋友.
他們自小就是一種好的環境去成長.
好的環境不單止.
不單止就是供他最好的讀書補習.
那些是好的.
可以給他有好的裝備.
但是如果品格不好.
其實就很大件事.

$^{1961}$希望我們都能夠透過第二章.
其實很負面的.
是說鬼詐,鬼詐,鬼詐.
神對人的鬼詐,對太太的鬼詐.
但是最終告訴我們.
原來神很想先知.
來勸導他們.
第三章有好東西看.
第三章他見到一個很正面的圖畫.
我們下個星期再說.
我們一起禱告.
感恩天父給我們在過年.
都可以一起有相聚.
透過網上,透過現場,透過錄音.
我們可以來到一起學習聖經的真理.
願神繼續幫助建立我們.
讓我們能夠認定.
我們是有真理的仁義和聖潔.
的一個新的創造.
我們要穿上新人.
在我們的生活裡.
奉主耶穌基督的名.
阿們.
\newpage



\section{列王記上 19:1-18}
\label{sec:C6eEbdGd3Ao}
\textbf{2024-02-18 午堂《當你以為走到絕境時》列王記上 19\_1-18 講員:鄧泰康傳道}
\newline
\newline
連結: \href{https://youtube.com/watch?v=C6eEbdGd3Ao}{\texttt{https://youtube.com/watch?v=C6eEbdGd3Ao}} ~~~~ 語音日期: 2024-02-19
\newline
\newline
\hyperref[sec:oCSQzasymjk]{\small{< < < PREV SERMON < < <}}
~
\hyperref[sec:index_chronic]{\small{[返順時目]}}
~
\hyperref[sec:index_scriptual]{\small{[返順卷目]}}
~
\hyperref[sec:myBJVEoimI0]{\small{> > > NEXT SERMON > > >}}
\newline
\newline
列王記上 19:1-18
\newline
\begin{longtable}{cl}
\hline
\hline
章節 & 經文 (和合本修訂版)\\
\hline
19:1 & \begin{tabularx}{0.7\textwidth}{X} 亞哈把以利亞一切所做的和他用刀殺眾先知的事都告訴耶洗別。 \end{tabularx} \\ \\ \relax
19:2 & \begin{tabularx}{0.7\textwidth}{X} 耶洗別就派使者到以利亞那裡，說：「明日約這時候，我若不使你的性命像那些人的性命一樣，願神明重重懲罰我。」 \end{tabularx} \\ \\ \relax
19:3 & \begin{tabularx}{0.7\textwidth}{X} 以利亞害怕，就起來逃命，到了猶大的別是巴，把僕人留在那裡。 \end{tabularx} \\ \\ \relax
19:4 & \begin{tabularx}{0.7\textwidth}{X} 他自己在曠野走了一日的路程，來到一棵羅騰樹下，就坐在那裡求死，說：「耶和華啊，現在夠了！求你取我的性命吧，因為我不比我的祖先好。」 \end{tabularx} \\ \\ \relax
19:5 & \begin{tabularx}{0.7\textwidth}{X} 他躺在羅騰樹下睡著了。看哪，有一個天使拍他，對他說：「起來吃吧！」 \end{tabularx} \\ \\ \relax
19:6 & \begin{tabularx}{0.7\textwidth}{X} 他觀看，看哪，頭旁有燒熱的石頭烤的餅和一壺水，他就吃了喝了，又再躺下。 \end{tabularx} \\ \\ \relax
19:7 & \begin{tabularx}{0.7\textwidth}{X} 耶和華的使者回來，第二次拍他，說：「起來吃吧！因為你要走的路很遠。」 \end{tabularx} \\ \\ \relax
19:8 & \begin{tabularx}{0.7\textwidth}{X} 他就起來吃了喝了，仗著這飲食的力走了四十晝夜，到了神的山，就是何烈山。 \end{tabularx} \\ \\ \relax
19:9 & \begin{tabularx}{0.7\textwidth}{X} 他在那裡進了一個洞，在洞中過夜。看哪，耶和華的話臨到他，說：「以利亞，你在這裡做甚麼？」 \end{tabularx} \\ \\ \relax
19:10 & \begin{tabularx}{0.7\textwidth}{X} 他說：「我為耶和華－萬軍之神大發熱心，因為以色列人背棄了你的約，毀壞了你的壇，用刀殺了你的先知，只剩下我一人，他們還要追殺我。」 \end{tabularx} \\ \\ \relax
19:11 & \begin{tabularx}{0.7\textwidth}{X} 耶和華說：「你出來站在山上，在耶和華面前。」看哪，耶和華從那裡經過。在耶和華面前有烈風大作，山崩石裂，耶和華卻不在風中；風後有地震，耶和華也不在其中； \end{tabularx} \\ \\ \relax
19:12 & \begin{tabularx}{0.7\textwidth}{X} 地震後有火，耶和華也不在火中；火以後，有輕微細小的聲音。 \end{tabularx} \\ \\ \relax
19:13 & \begin{tabularx}{0.7\textwidth}{X} 以利亞聽見，就用外衣蒙臉，出來站在洞口。聽啊，有聲音向他說：「以利亞，你在這裡做甚麼？」 \end{tabularx} \\ \\ \relax
19:14 & \begin{tabularx}{0.7\textwidth}{X} 他說：「我為耶和華－萬軍之神大發熱心，因為以色列人背棄了你的約，毀壞了你的壇，用刀殺了你的先知，只剩下我一人，他們還要追殺我。」 \end{tabularx} \\ \\ \relax
19:15 & \begin{tabularx}{0.7\textwidth}{X} 耶和華對他說：「去吧，從原路回去，往大馬士革的曠野去。到了那裡，你要膏哈薛作亞蘭王， \end{tabularx} \\ \\ \relax
19:16 & \begin{tabularx}{0.7\textwidth}{X} 又膏寧示的孫子耶戶作以色列王，並膏亞伯‧米何拉人沙法的兒子以利沙作先知接續你。 \end{tabularx} \\ \\ \relax
19:17 & \begin{tabularx}{0.7\textwidth}{X} 將來逃過哈薛之刀的，必被耶戶所殺；逃過耶戶之刀的，必被以利沙所殺。 \end{tabularx} \\ \\ \relax
19:18 & \begin{tabularx}{0.7\textwidth}{X} 但我在以色列中留下七千人，是未曾向巴力屈膝，未曾親吻巴力的。」 \end{tabularx} \\ \\
[1ex]
\hline
\hline
\end{longtable}
$^{1}$第五姐妹,祝神.
在此時先講一講,請第五姐妹紀念郭牧師及師母.
因為我們昨日在六時多收到她的消息.
師母和郭牧師一起飛回香港的途中在飛機上暈倒.
後來時間,我們都很感恩,她後來都醒了.
但她最後送了去醫院,現在最後的情況,她發覺有些心三失的情況.
即是血管三失,所以她現在已經去了馬來西亞醫院.
所以請第五姐妹紀念她們的情況.
我想在講到之前為她們禱告,大家請提心吐氣.
七位天父,我們求你幫助,保守.
剛才在你的詩歌當中,你告訴我們你是我們的盾牌,你是我們的山寨,你是我們的伴石.
主在過去日子裡面,主在你如何帶領著郭牧師,你如何成為他的山寨盾牌和力量.
在這一刻裡面,我知道他很需要你,師母也很需要你.
最後我求主你親自保守,你去帶領著,保守他們兩位.
在神的面前的時候,是有恩典,是滿有眾神熱烈的平安.
最後我求主你親自特別的,將你醫治的手指頭能夠進到師母她自己的身上.
求你派最好的醫療團隊成為她的幫助.
我也求主你將平靜安穩的心賜給牧師,賜給他的兒女.
最後我們求主你親自帶領和看顧.
願意我們能夠在過程當中,我們完全看見上帝你如何帶領和看顧他們.
將一切所有的我們教託,其中一方教主耶穌,得勝名至其球,阿們.
我們今天讀這段經文,其實也是大家很熟悉的一段經文.
從小到大,如果小朋友回到教會,一定會聽過這段經文.
因為在以利亞,他在人生當中是一個很重要的高峰.
當我們去看這段經文的時候,我們必須要去比較十八章.
當時很義氣風發的以利亞.
當時的以利亞,他在一個,你叫他在一個事奉裡面是人生高峰.
他聽到上帝告訴他,將不下雨,他就會向著阿哈宣佈不下雨.
三年之後,上帝叫他再宣佈,下雨就下雨.
所以他有這一份底氣,去到當他面對阿哈的時候.
面對著他們,當時如果要記載,有八百五十個拜賈神的先知.
和他在山上一起去試,如果我想說的義氣風發就是什麼.
就是當時他們一起捉了兩個壇,然後我們一起用同樣的方法.
將祭神放在壇上,然後就叫他們,好啊,你們叫你們的神在天上降火下來.
聖經是形容這八百五十個先知用盡一切方法,跳舞都沒有辦法令這些火降在當中.
然後,以利亞他做了一個舉動,將水淋在祭壇上,淋濕他.
然後他就祈禱,上帝就將火降在上面,燒了整個祭壇.
所以當時,我不知道大家看不看到那個場面.
一個人對著八百多人,但是他一點都不害怕.
但是,一張而已,鏡頭一轉而已.

$^{41}$當他聽到耶駛別,差派使者來到他面前講這番話的時候.
突然間整個情況急轉直下.
其實我們很難去想像究竟兩者之間的差距,你明白我的意思嗎.
以利亞,我明明是跟隨著上帝,我看到你那些作為,所以我有這麼大的信心.
但是,竟然一個女人,他跟他說這番話,你令他害怕到一個地步,他馬上走.
聖經有趣的地方就是,他沒有跟你說究竟耶駛別有多可怕,沒有.
或許你可以在裡面有很多想像的空間,他是如何的凶殘.
他是如何的可以即時發動任何的情況,可以殺了他,都不一定,我不知道.
但是,在這裡我們只知道一件事.
以利亞去畏懼耶駛別過於上帝.
今天我想跟大家說的時候,在我們人生當中都有很多我們的耶駛別.
我們同樣都經歷了上帝很多很多過去的恩典,他與我們同行.
但很多時候,當一個危機,當一個困難難關臨到的時候.
接著你就會發覺,我們就馬上逃命,我們就馬上走.
所以,其實這段經文最有趣的地方就是,我們要重新繼續去認識.
我們在平時,我們以為自己,我們如何相信上帝.
甚至乎,我們在信仰當中有一份力量.
但其實,當這些困難危機來到的時候,就是一個很好的試金石.
你會發覺,其實以利亞他給予什麼情況,令到他發生這件事呢?.
當你看經文當中,我會這樣形容,就是他對上帝的認識毀了他.
他對上帝的過去認識毀了他.
意思是什麼?他心裡面有一套神學.
神學的意思就是你去認識上帝.
所以,他心裡面有一套神學,就是上帝是這樣的,是這樣的,是這樣的.
即是說,他心裡面的上帝是一個怎樣的上帝?.
就是你叫他不下雨,他就不下雨.
是滿有大能的.
還有一件事,只會贏,只會贏,是不會輸的.
但是,當他一聽到的時候,聖經裡面沒有很細微的告訴你,究竟發生了什麼事.
但是,或許,在裡面我們有很多想像空間.
當以利亞聽到耶洗別所說的話的時候.
他跟著他的反應,可能他就說,為什麼上帝不即時擊殺耶洗別?.
為什麼上帝不做一些事,讓我們看見?.
不是的,上帝好像突然之間消失了.
現在所發生的事,和我過去裡面,我對上帝的認識完全是兩回事.
他接受不了.
所以,你會發覺,當他逃命的時候.
當他去到那裡,無能為力的時候.
他就說,去到很出名的羅騰樹,向上帝求死.
然後說,算了吧,求你取了我的性命.

$^{81}$因為我比我以前的祖先好.
這句話其實很有趣,他又不敢罵上帝.
但是,他在裡面的時候,他又用了一番話.
他聽了這番話,他真的很誠心地說,我不及我的祖先.
還是在遞給上帝.
但是,你會發覺,之後的經文.
你就會發覺,解釋了,其實他是向上帝發悔氣.
他一路去裡面,這個就是他心裡面的神學.
我為耶華萬軍的神大發熱心.
因為是以色列人,不是我.
是那幫以色列人背了你的約,毀了你的壇.
用刀殺了你的才知道,現在只剩下我一個.
他們還要來取我的命.
他在裡面有很多的不忿.
很多就是上帝為你撲心撲命.
上帝為你做了很多的事.
為什麼現在你看到的時候,你不是保護我的?.
如果大家留意剛才的經文的時候.
基本上第一段的經文和第二段的經文是差不多一模一樣.
他講了兩次,這個是他很根深蒂固的裡面.
我已經盡了我的能力,我盡了力,難道你看不見嗎?.
我不知道大家在這裡有沒有一些痛感.
有時候我們面對困難,我們面對一個處境.
我們最難去渡過的是什麼呢?.
就是這件事不斷在我們的心裡面,不斷旋轉.
在自己的心裡面不斷碎碎念.
不斷地在說,不斷地在說.
上帝你去了哪裡?你是什麼神來的?.
我已經很撲心撲命為你了.
你還想怎樣?.
為什麼你還想怎樣?你不是公義的嗎?.
你不是慈愛的嗎?你不是幫我的嗎?.
以利亞當時他就一直在腦海裡面不斷地打轉.
他沒有辦法走出來.
但是你會發覺經文裡面很有趣的地方.
出了一段不解,很難明白的經文.
耶和華說,你出來,你出來殺我的面前.
然後他說,那時耶和華從哪裡經過.
在他面前有烈風大作.
崩山碎石,耶和華卻不在風中.

$^{121}$風後有地震,耶和華卻不在其中.
地震後有火,耶和華也不在火中.
火後有微小的聲音.
這段經文你不解,是很難明白的.
其實你會發覺,你記住.
以利亞去到河列山.
河列山就是當日上帝在上面頒佈律法.
如果大家還記得當時的圖像是怎樣.
就是崩山碎石,就是大風大作,就是地震.
這些圖像是代表著上帝將會出現.
上帝臨在的象徵.
這是以色列人每一個人一看這圖像就會知道.
上帝來了.
但最有趣的地方,在這裡說什麼.
上帝是怎樣,今次的出現是什麼.
我在微星之中.
上帝做了一件事.
現在是你以為我.
你認為的上帝是這樣.
現在我告訴你,上帝是這樣的.
明愛斯萬,弟兄姊妹.
其實這裡很重要.
上帝透過這件事,重新將焦點從自己身上.
轉移到上帝身上.
這亦是我們每一個人,當你面對困難.
你面對處境危難的時候.
最好的方法.
就是你去靈修,你去崇拜,你安靜.
其實就是上帝要你將你的目光轉到上帝身上.
上帝用微星跟你說話.
微星是什麼.
當你在很雜亂的情況下,你聽不到的.
大家有沒有試過.
現在我們的耳機很厲害.
我們現在有一個降噪功能.
現在我們用的耳機有一個功能.
你可以令到它不要降噪.
你又戴著又聽到外面的東西.
但大家就會見到分別.
就是你在外面很雜亂的環境下.

$^{161}$你裡面是很難聽到音樂的.
唯一就是當你將噪音降噪完全拿走了之後.
你就會很聽到.
弟兄姊妹,今天我們在困難的時候.
你以為自己走到絕境.
剛才我都提過.
魔鬼最好的方法就是在心裡面.
不斷地旋轉,不斷地徘徊.
不斷地跟你說,不斷地跟你說.
上帝不好.
你的上帝是差的.
上帝不愛你.
你以為你上帝很厲害嗎.
你看,他現在什麼都不做.
這個就令到我們沒辦法從那個困境裡面走出來.
但現在這個時候,就是我們要走出來.
上帝要他走出來.
在微聲當中,請你細心地去聽.
我和你說的這番話.
他很想告訴你.
你以為我是這樣.
但是這個上帝是怎樣的呢.
在剛才那段話還沒發生之前.
當他在尋死.
其實我想形容整段經文是很有趣的.
有趣的地方是什麼呢.
上帝從來都沒有正面回答過他任何問題.
從來都沒有.
他很想知道的答案,上帝從來都沒有說過.
當他求死之後.
上帝用了兩次的方法.
用了兩次什麼.
你起來吧.
吃東西吧.
上帝很清楚知道以利亞當刻他最需要什麼.
一個人在很不開心很恐懼的裡面.
他真的吃不下東西.
但他卻是真的需要吃東西.
上帝為他預備.
你起來吧,吃東西吧.

$^{201}$弟兄姊妹,這是我們的上帝.
上帝不急於和你爭辯你自己認為人生當中很重要的事.
上帝不急於和你爭辯.
上帝關心的是你.
我們經常強調我們的信仰不是doing,而是being.
是你的人是怎樣.
上帝很關注的是你,而不是你做了多少事.
上帝其實對於以利亞之前所做的事.
當然他欣賞和給他能力.
但上帝重視最重要的不是你做了多少事.
他不會因為你戰勝了那麼多先知.
說你很厲害,不是的.
而是你現在怎樣.
上帝很想你起來吃東西,兩次.
當你吃完之後,他就去睡.
第二次再叫你起來吃東西.
因為你要走的路還很長.
像不像一個爸爸和一個剛生育失敗的兒子.
回到家的時候,垂頭喪氣的圖畫.
暫時甚麼都做不到.
但我可以為你供應的是甚麼.
關心和愛,這是我們的上帝.
第二個,上帝問了一個問題.
以利亞你在做甚麼.
問了兩次.
以利亞你在做甚麼.
為甚麼要問這個問題.
其實當你問這個問題的時候也是.
等於回到以前.
如果有一天你們問我們童宮.
回過頭來問,為甚麼你要做傳道人.
其實要的不是那個答案.
要的是你會喚醒我們.
為甚麼最初我們要做傳道人.
是一個初心.
以利亞最初被上帝呼召出來的時候.
其實是很想以利亞能夠去為他工作.
才知道就是上帝的代言人.
是我叫你去說甚麼,我叫你去做甚麼.
你跟著做.

$^{241}$為甚麼這次你會跌到這樣.
因為你自己認為自己很厲害.
慢慢你改變了.
你是利用了上帝.
表面上的就是一樣.
我們的上帝多厲害.
其實背後是甚麼.
你看我多厲害.
我的上帝多厲害.
你越利用上帝去為你做事.
你下不來.
甚麼叫上帝的主權.
上帝的主權就是.
他可以供給你,他也可以收回.
是不是,弟兄姊妹.
這個才是上帝.
所以上帝很想重新去了解這件事.
所以上帝最後跟他說.
上帝跟他說的那些是不回應他的.
他很想知道甚麼原因.
但上帝給他的回應是甚麼.
你按你的原路回去.
其實這件事很難的.
可以從那裡逃走出來.
你按你的原路回去.
人在哪裡跌倒.
你在那裡重新起來.
你回去.
然後做甚麼.
你回去,高了幾個王.
亞蘭王,以色列王.
他想告訴你甚麼.
我是掌管外邦的神.
我是掌管以色列的神.
選召和廢掉都是上帝決定的.
而最厲害的是甚麼.
你回去,按納以尼沙接續你.
這件事就剛剛好回應他.
他說甚麼.
現在只剩下我一個.

$^{281}$現在上帝說.
我可以隨時不要你的.
你以為你是甚麼.
我不知道弟兄姊妹.
有沒有發覺很特別的地方.
如果大家看這段經文.
讀下去之後.
以尼沙不是發晦氣.
不發脾氣.
如果正常的話.
按照剛才那個思想邏輯的話.
他自己那麼厲害.
他的邏輯應該是甚麼.
上帝現在只是耍我.
我為你仆生仆命.
做了那麼多的時候.
現在你突然間不要我.
你算怎樣.
應該是這樣的邏輯.
但不是.
你會發覺見到這件事之後.
你會發覺突然間.
這個人突然間醒了.
他醒了.
他會發覺自己並不是那麼厲害.
為甚麼呢.
因為你會發覺.
最後一段經文就是.
我在以色列留下了七千人.
是未曾向巴力屈膝的.
未曾和巴力親嘴的.
七千人.
你是七千零一個.
在那一刻他醒了.
他真真正正知道.
其實原來他不是那麼厲害.
他不是那麼偉大.
上帝可以做到更多的事.
最諷刺的地方是甚麼呢.
就是在剛才.

$^{321}$我說他和先知大戰之前.
第十八章出現了一個人物.
叫做俄巴底.
俄巴底是當時阿哈的一個神僕.
聖經講得很清楚.
因為他敬畏耶和華.
所以他存留了一百個耶和華的先知.
是因為敬畏.
所以以利亞他突然間如夢初醒.
他脫胎換骨.
突然間變了另一個人.
他用了一個方法就是完全馴服.
如果大家再看下去的時候.
你會發覺.
當他去找以利沙的時候.
以利沙不是和他說了理由嗎.
你讓我回去做甚麼先.
伊麗亞他很隨意的.
他說不要緊.
我做了甚麼.
你便回去.
你會發覺.
我只是宣告上帝要我做的事.
但你會發覺之後他的厲害是甚麼.
大家還記不記得.
他很隨意的拿自己的外衣.
打在河水上.
分開了走過去.
然後以利沙跟著他.
最後以利沙.
最後結果是被火車火馬接走了.
跟著時間以利沙又用她的外衣.
打水之後她又走過去.
代表接觸了他的身份.
但你會發覺他整個身份是提升了.
而你亦會發覺在第21章的時候.
以利沙當時宣告.
向阿蝦.
向他以前所恐懼的耶駒別.
宣告上帝的審判的時候.

$^{361}$他是變回很平靜的.
他的說話是甚麼.
論到耶駒別.
耶華也說.
苟在耶斯列的外國.
必吃的耶駒別的肉.
跟著這件事在列王記下第九章就發生了.
當耶護他來攻打他們的城樓的時候.
跟著耶駒別的僕人.
就將他在城樓上丟了下來.
就死了.
你會發覺其實剛才那段經文.
他說將來逃避哈舌之刀的必被耶護所殺.
逃避耶護的刀的必被以利沙所殺.
上帝很想告訴你.
上帝要成就的事.
是一定會成就的.
弟兄姊妹今天我們在神的裡面.
上帝在他的計劃當中.
從來沒有離席.
當我們跟一些朋友去講福音的時候.
我經常說我們有一個觀念.
不知道大家知不知道.
我們今天所身處的地球.
其實是要有很多客觀環境下.
才可以正常運作的.
如果大家知道.
我們的潮汐漲退.
是因為我們的地球.
邪邪地在轉的時候影響潮汐.
即是說如果它歪一點點.
你就出事了.
如果它不轉.
你就出事了.
為什麼我們今天可以站在這裡.
不會全部浮起來.
因為我們有地心吸力.
為什麼會有地心吸力.
是整個地球的運作產生的地心吸力.
所以令到你站在這裡.

$^{401}$是剛剛好的.
我們很清楚.
上帝在這個地球上所做的.
全部都是一種繼續不斷的看顧.
不斷繼續在運作.
上帝從來沒有離開過.
上帝要做的事一定會發生.
所以為什麼我經常跟弟兄姊妹說.
你在性境當中.
你要掌握著上帝給你的應許.
祂給你的應許是從來都不變的.
祂對你的愛從來都不改變的.
所以弟兄姊妹.
當你去面對這些困難.
你覺得自己走到無力.
你走到自己無路可走的時候.
請你將你的焦點再一次回到上帝那裡.
因為祂一直都有掌管.
因為祂可以將一切不可能變成可能.
每一樣東西都可以有出路.
其實你不需要看很遠.
你看回你自己人生的歷史.
你就知道.
最簡單.
大家以前有沒有拍過拖?.
有多少人這麼幸運.
拍第一個就成為老公的?.
應該都不會佔太多吧?.
如果有的話真的要祝福你.
但是記不記得.
以前曾經你失戀的時候.
哭到死去活來?.
記不記得你呼天搶地?.
覺得這個世界沒有那個就不行.
記不記得?.
弟兄姊妹.
從來我們的人生都不是這樣的.
我們的人生當中有困難.
是我們去看那件事沒有出路.
不代表上帝沒有出路.

$^{441}$所以剛才我們在外面.
特別感恩到今天.
弟兄姊妹每次見到我都會問.
你身體怎樣?.
我的官方回應很簡單.
就是我的身體狀況和手術前沒什麼分別.
走路都繼續這樣走.
都是這樣挨.
都是這樣過生活.
但是弟兄姊妹.
從來我們都不應該將我們的困難.
放在我們的困難上.
因為上帝在當中.
上帝才會知道最終的答案.
你知實祂.
你會在實祂.
掌握祂的應許.
之後好好走祂的路.
不要偏差.
你一定能夠經歷到.
聖經裡面所有最豐富的應許.
我要得人.
我要拯救人.
我要讓他能夠得著豐盛的生命.
弟兄姊妹.
所以我們盼望我們一起.
好像以利亞一樣.
將你的心思重新改換一生去改變.
以至讓你能夠更加踏實的.
去成為上帝的代表,代言人.
我一直是同心祈禱.
差點我們要多謝神你自己.
主在你面前的時候.
我們要知道一切都在你的當中.
主我求主你自己讓我們能夠生命的改變.
讓我們眼睛的焦點的改變.
讓我們能夠可以完全的跟隨上帝你自己.
求你幫助我們.
主我們如果有軟弱.
我們有做得不好的時候.

$^{481}$主但我們知道主你要的是我們來到你面前.
很實實在在的我們去承認我們的罪.
我們承認我們自己做得不好的地方.
但我們多謝你.
你願意赦免我們.
你去幫助我們.
特別我求主你去保守.
在在座當中.
有很多弟兄姊妹正面對一些很不容易的處境.
在他們憑自己的能力.
沒辦法可以去勝過那個處境.
令到他們很懼怕.
甚至想逃走.
主但我求主你幫助他們.
我知道你會對他們.
跟對以利亞一樣.
利用你的溫柔.
利用你的恩典.
去提醒我們.
主人你是當中.
我們是一群相信你的人.
我們有你就足夠.
我們獻上感恩.
多謝讚美你.
其利奉教主耶穌,得勝名至其今.
阿們.
\newpage



\section{以弗所書 4:20-23}
\label{sec:myBJVEoimI0}
\textbf{2024-02-18 早堂《如果...可以重來》以弗所書 4\_20-23 講員:黃仲安傳道}
\newline
\newline
連結: \href{https://youtube.com/watch?v=myBJVEoimI0}{\texttt{https://youtube.com/watch?v=myBJVEoimI0}} ~~~~ 語音日期: 2024-02-19
\newline
\newline
\hyperref[sec:C6eEbdGd3Ao]{\small{< < < PREV SERMON < < <}}
~
\hyperref[sec:index_chronic]{\small{[返順時目]}}
~
\hyperref[sec:index_scriptual]{\small{[返順卷目]}}
~
\hyperref[sec:4SEW84zBGCU]{\small{> > > NEXT SERMON > > >}}
\newline
\newline
以弗所書 4:20-23
\newline
\begin{longtable}{cl}
\hline
\hline
章節 & 經文 (和合本修訂版)\\
\hline
4:20 & \begin{tabularx}{0.7\textwidth}{X} 但你們從基督學的不是這樣。 \end{tabularx} \\ \\ \relax
4:21 & \begin{tabularx}{0.7\textwidth}{X} 如果你們聽過他的道，領了他的教，因為真理就在耶穌裡， \end{tabularx} \\ \\ \relax
4:22 & \begin{tabularx}{0.7\textwidth}{X} 你們要脫去從前的行為，脫去舊我；這舊我是因私慾的迷惑而漸漸敗壞的。 \end{tabularx} \\ \\ \relax
4:23 & \begin{tabularx}{0.7\textwidth}{X} 你們要把自己的心志更新， \end{tabularx} \\ \\
[1ex]
\hline
\hline
\end{longtable}
$^{1}$早晨 各位節目好 今天很高興來到錦繡花園這個地方.
跟大家分享一些神奇的訊息.
我自己住在美孚.
所以今天早上很快就來到錦繡花園.
比我想像中還早了十分鐘.
有一次我媽媽在美孚帶我兩歲大的兒子回長沙灣.
帶他的時候發生了一件趣事.
我想跟大家分享一下.
在搭車途中 有一個太太說.
小朋友 你嘴巴張開說這麼多話 你今年多大了.
兒子說 我兩歲.
太太就鬱子一算.
現在是農年 鬱子一算.
那你不就是牛仔嗎.
牛年出世.
你猜一下我兒子有什麼回應.
他說 我不是牛仔 我是美孚少爺.
為什麼他會這樣回答呢.
原來我媽媽教一個兩歲小朋友的方法.
教他自理.
吃飯要擦嘴 要洗手.
就要像少爺一樣.
所以我兒子就潛移默化.
回答我是美孚少爺.
千萬不要誤會 我家裡不是富貴家庭.
所以他這樣說.
但是有時候我們反思.
我們生活有時候有很多的忙碌.
剛才聖經所提的.
有很多的罪 有很多的纏繞.
其實有時候令我們生命很多的重擔.
就像一隻牛多於一個人.
所以弟兄姊妹 你問一下自己.
你是誰 你知道你是誰嗎.
我是神的兒女 是不是.
有時候我們做一次牛.
忘記了這麼尊貴的身份.
就是我和你都是神寶貴的兒女.
所以弟兄姊妹 剛才給我們一些提醒.
就是時間其實不多了.

$^{41}$耶穌再來的日子可能很快.
在這麼短的時間裡.
我們不要再做牛了 是不是.
我們做一個神的兒女.
我們今天一起去看看這段經文.
給我們的啟發.
讓我們去想一想.
我們生命究竟是撐了些什麼.
或者我們現在的狀態要怎樣改變.
如果再來過 再來過.
新的一天 新的一年.
我們要怎樣過.
才能獲得此聖經所應許.
神兒女的身份 好嗎.
我們一起來去禱告.
主音請你讓我們今天.
回來教會裡面一起去敬拜.
是很多無比的事情.
現在我們可以在生活.
一些的忙亂和憂慮裡面去停一停.
重新去思想我們是神兒女.
這個寶貴的身份.
求主幫助我們去掌握這個福份.
讓我們的生命活得像聖經所應許.
你兒女那種的蒙福.
求主幫助我們.
我們禱告奉耶穌基督的名字.
阿們.
因為我們當說到神兒女的身份.
是很無比的一件事.
如果大家看到羅馬書裡面.
去應許神兒女身份那種特質.
就說到世界上沒有東西是控訴神的兒女.
也說到當我們不知道怎樣去祈禱.
心裡有很多憂慮的時候.
聖靈又幫我們祈禱.
這些都是神兒女寶貴的身份.
我們千萬不要忘記.
今天《以弗所書》裡面.
正正就是重新提醒我們.

$^{81}$《以弗所書》四章第20節裡面說到.
我和你都是跟耶穌跟基督的.
不可以好像世人那樣.
如果我們聽過耶穌的教導.
學過他的真理.
聖經裡面這樣去形容.
是脫去我們從前的舊人.
脫去是一個很生活化的形容.
說得好像穿衣服一樣.
我和你其實每一天都會更換我們的衣服.
所以說到脫去是每一天我們都要做的.
這個舊人是怎樣來的.
我們要清楚.
因為當我們知道這個舊人.
這個罪人.
這種前路是怎樣來的.
提醒我們要小心.
這種前路是怎樣來的.
聖經裡面說到.
是因為私慾.
私慾是來自世界的.
聖經說到各樣的好處都是來自上帝.
但唯有這個邪情私慾.
身體的私慾是來自撒旦.
是來自世界的.
這個世界的私慾.
慢慢令到我們在世界裡面生活的人.
變得大壞.
所以我們要留心.
世界給了我們價值觀.
然後聖經怎樣去提醒我們.
在這種價值觀的圍繞裡面.
我們要改變我們的心智.
即是改變你的看法.
然後二十四節裡面再延伸去說.
用剛才穿衣服的比喻.
我們要穿上新人.
那個新人是有神的形象的.
我們這樣去看的時候.
發覺當我們要去改變和更新.

$^{121}$我們不用等到年廿八.
我們每一天都可以這樣去更新我們自己.
就好像穿衣服換衣服一樣.
可能有弟兄姊妹會疑惑.
周圍有句說話.
江山易改 品性難移.
一個人是不是真的可以有這麼大的改變呢.
有時候我們都會有這種疑惑.
但大家要知道.
你知不知道現在耶穌基督.
今天在這個世界上的工作是甚麼.
我們很清楚耶穌基督在二千年前的工作.
即是拯救我們.
剛才詩歌都有很多提醒.
拯救我們.
但知不知道今天耶穌基督在做甚麼.
我們有時候都要知道.
耶穌現在在做甚麼.
我們才可以和祂的工作對照.
我們才知道祂現在是一個永活的神.
我們要看一看.
啟示錄裡面給了我一個很清晰的描述.
講到耶穌在做甚麼.
在末世的日子.
啟示錄第21章裡面.
講到坐在寶座的.
這樣說.
即是我們說耶穌基督.
他說看啊.
我將一切都更新了.
在末世耶穌基督會宣告.
祂已經將一切都更新了.
所以在耶穌再來之前這段時間.
祂不斷都在做一件事.
就是祂在更新我們這個世界.
更新我和你的生命.
所以我們已經知道耶穌基督今天的工作.
我們才可以和祂現在對應.
去match到.
我們才可以看到.

$^{161}$信仰的真實.
每一天都在我們生命裡面發生.
七弟姐妹千萬不要忘記這件事.
或者不知道這件事.
你要知道耶穌現在不斷地在更新.
有點像你的手提電話.
不斷地在更新.
所以你忘記了更新.
你會發覺好像很多事情都不順暢.
所以耶穌是不斷地更新我們世界.
不斷地更新我和你的生命.
所以你和我都可以改變.
任何時候都可以改變.
只要我們意識到.
神是有這樣的工作.
然後我們生命和祂的工作.
在推照的時候.
我們就察覺到更新.
在詩篇裡面記述大衛的一個心情.
一個禱告.
都是在回應這種警報.
在詩篇第51篇裡面.
有一個很出名的禱告.
我們都要學習一下大衛.
怎樣去祈禱去更新.
第51篇裡面講到大衛這樣祈禱.
聖帝啊求你為我做清潔的心.
使我裡面重新有正直的靈.
所以見到大衛這個禱告.
正正在反映我們生命的狀態.
你心裡面是有正直的靈.
但是被世界的邪情私慾去破壞.
所以大衛的禱告就是重新有正直的靈.
所以我們都要學習大衛這種禱告.
不斷地去禱告.
神去幫助我們去做這個清潔心.
然後讓我們重新有正直的靈.
每一日我和你都可以這樣去禱告.
生命有很多阻礙令到我們忘記了.
或者沒有動力去做這個禱告.

$^{201}$那個感覺是什麼呢.
就是一種遺憾.
和大家分享另外一件事.
就是當兒子兩歲多的時候.
香港有很多父母.
都要面對這個選幼稚園.
有沒有孫子孫女都面對這些情況.
或是子女.
現在的幼稚園是很艱難的.
大家有沒有看過一些幼稚園面試片.
沒有看過嗎.
有興趣看一下.
比恐怖片還要恐怖.
為什麼這樣說呢.
就是兩歲多的小朋友.
那些面試片是怎樣呢.
不是叫他們吃飯.
或者洗手那些.
去到科學館.
然後指著一隻翼龍.
然後問小朋友.
知不知道這隻翼龍的英文是什麼.
小朋友竟然說是terisoria.
我自己都不懂.
小朋友教我.
但這條都不是最恐怖.
最恐怖是另外一條.
是媽媽拿了一張數學工作紙出來.
兩歲多的小朋友竟然會計算加數.
兩個位的加數.
看到這些恐怖片.
父母都心寒.
竟然和這些人競爭.
近來有一次.
我自己的面試.
都遇到很大的挫敗.
一進到面試的房間.
竟然被我看到小學的校長和幼稚園的校長.
然後幼稚園的校長就說.
小朋友你要快一點.

$^{241}$不要賣摸摸.
現在面試中.
父母在當時遇到這些情況.
有種愛你的子女.
但又幫不到他.
那種愛莫能助的感覺.
就會產生那種後悔.
人生其實很多這種情況.
你遇到一些困難.
你很努力去爭取.
但到最後和機會可能是失之交臂的時候.
產生了很多後悔.
我們經常都在這些後悔裡面.
去糾纏.
你看以色列人出埃及.
他們都有這種情況.
他們基本上已經看到迦南地.
但仍然害怕.
看到那個地方.
每個人都巨人.
我們好像蚱蜢一樣.
不要搞那麼多事了.
不如走吧.
他們走了幾步就開始後悔.
然後他們就在抗議裡面兜轉很多年.
靖姐妹阻礙我們更新.
最大的敵人就是我們心中的後悔.
所以靖姐妹為什麼有時候我們更新不到呢.
其實我們心裡面有很多的愧疚.
是阻礙我們.
這些愧疚是怎樣來的呢.
因為我們經常去想像.
如果當初在那個時刻.
我做另外一個決定.
可能有一個更好的後果.
我就是被這種幻想.
去阻礙了自己去前進.
我自己都親身經歷.
如果那一刻.
我是怎樣幫我的小朋友的話.

$^{281}$可能當時的面子.
就不會是那麼糟糕了.
我經常在裡面去責備我自己.
令到自己都很難更新.
所以靖姐妹我們要小心.
這就是世界的價值.
阻礙我們去得到.
聖經所說的一無掛慮.
神兒女的身份.
我經常去想像.
當初如果.
我有另外一個決定.
可能有另外一個更好的後果.
我自己的反省就是.
我們有那麼多的如果.
我倒不如相信.
當時神已經在眾多的如果裡面.
幫我們選擇了最好的那個.
就是我們今天.
所身處的世界.
與其我去想像.
一些根本沒有辦法去證實.
的假設.
我倒不如相信.
今天我已經活在神為我們.
去選擇最好的一個結果裡面.
補充一句.
當初雖然是這樣的面試情況.
但最後那間學校都收了我兒子.
因為上次我在教會.
跟大家分享.
那些人很擔心.
我跟他們說最後那間學校都收了我兒子.
因為那間學校有六千人去競爭.
所以我自己當時都在壓汗.
與其我們相信.
可能有其他更好的結果.
倒不如相信今天神已經幫我們選擇了最好的.
所以弟兄姊妹你可以改變.
改變就是我們的福音.

$^{321}$我們今天無論在.
什麼情況都好.
你可以改變.
這就是神給我們的福音.
亦都是耶穌基督在末世裡面的工作.
但可能我們不察覺.
其實神在《以賽亞書》裡面.
都給我們一個啟示.
讓我們知道.
祂的心意是怎樣.
原來《以賽亞書》一早已經講了.
耶和華我們的神這樣講.
你們不要記念從前的事.
都不要記念.
古時的事.
看啊我要做一件新事.
如今要發現.
你們豈不知道嗎.
我必在曠野開道路.
沙漠裡面開光學.
原來神更新的心意.
由創世到現在都一樣.
只不過.
有時候我們生命有很多的前類.
有很多的顧慮的時候.
我們根本沒有留意到.
在新的一年我們都不要記念.
過去的事.
我們要看神在今年.
在你和我的生命裡面.
去做一些新事.
當我們.
願意去擁抱.
或者期待新的事的時候.
我相信大家馬上已經想到.
我們生命會有很多.
轉變會出現.
我們新事和轉變.
是分不開的一件事.
哥林多前書.

$^{361}$哥林多後書裡面.
給了我們一個提醒.
它說你們總是要.
自己審查自己的信心.
也要試驗自己的信心.
你們不可以棄絕.
有基督在你們的心裡.
所以我們要審查.
我們自己內在的警惕.
然後求神的真理.
和我們生命的真實.
是有一個相遇.
然後幫助我們生命.
去產生改變.
有時候生命的改變.
是來得很快.
很突然.
甚至是令人產生一種抗拒.
和大家分享一件.
在疫情間的小事.
當時我兒子剛出生.
沒多久就半歲.
所以大家做父母就知道.
頭半年.
和兒子難離難舍.
上班想著他.
下班馬上回家看他.
有一個星期.
星期六和他有一個節目去玩.
誰知那個星期三.
自己身體有點不舒服.
就去看醫生.
醫生很快就說.
先生你胃痛.
胃痛來看醫生是不是想騙病假.
沒人知道我是有良善.
有忠心的僕人.
就去搖曬他.
星期四就發覺還是不妥.
就去看另一位醫生.

$^{401}$醫生就和我說.
你應該吃腸胃炎.
胃痛.
星期四還是不好.
星期五就再看.
第三位醫生.
醫生就說.
先生你的盲腸都穿了.
來了幾天都痛.
我說不是吧.
什麼叫盲腸穿了.
就是腹膜炎.
大家有些知識的話就知道.
腹膜炎就是你橫膈膜對下.
那些器官就浸在濃水.
你趕快去醫院吧.
我還記得那星期要港渡.
我還問醫生星期日能不能出設院.
去港渡.
他說你做不去就死了.
當時就遇到一些.
這樣的轉變.
當你有人想著星期六.
和兒子慶祝.
或者很開心的家庭日.
誰知道在醫院裡面就渡過.
有時候轉變來得很快.
你完全是預計不到.
我當時的感覺.
甚至乎是轉變來的時候.
你是想抗拒的.
誰會想去住醫院做手術.
沒有吧.
有時候轉變來得很急很快.
甚至乎覺得抗拒是一種痛苦.
弟兄姊妹我想問.
覺得盲腸對你來說.
是好事還是壞事.
是好事.
長遠來說是一件好事.

$^{441}$轉變就是有這種特質.
它剛剛來的時候.
很痛苦很想抗拒.
但是你想深一層.
長遠來看其實它是為你生命帶來好處.
出庭姐妹.
我們要求神在今年.
在我們生命裡面帶來一些改變.
你可能即時想不到.
有什麼可以求神幫助.
因為剛才說你去祈禱.
你也要有個焦點才能留意到.
神在什麼位置更新你的生命.
是不是.
生命裡面其實有很多東西.
我們可以和神求改變.
我舉一些例子.
大家可以自己想想.
有沒有一些需要求神在今年改變的方向.
第一.
可能你和神的關係.
上年覺得和神有很多疏離.
有時候回教會又不穩定.
和神的連繫很弱.
求神今年幫助你.
第二.
可能是你自己的健康.
你的身體.
或者是你自己生命一些.
價值觀的排位.
會不會是一些關係.
你想求神幫助你.
或者是你自己個人的精力水平.
可能現在經常感覺到很累.
求神幫助你.
好像如英輾刺上藤一樣.
也可能是你自己的工作.
你自己的事業.
你對生命的看法.
甚至是你的婚姻.

$^{481}$去求神幫助你.
或者是你自己每天的習慣.
你自己生活的習慣.
你自己的興趣.
你自己對小朋友的方法.
你自己的自信心.
你的財政.
甚至是你的夢想.
可能很多東西.
你都已經先放下了.
就求神在今年幫助你.
重新去更新.
在小路告訴我們.
耶穌就是在做這個工作.
所以你一定要將你自己的需要.
放在神的面前.
求祂幫助.
改變你的心意.
一直是沒有停過.
所以求神去幫助你.
羅馬書裡面說到.
我們的信心要看得合乎中度.
千萬不要過份去看重自己.
我們其實是需要神.
需要神的幫助.
就算最微小的事情.
明天我們上班遇到一些挑戰.
都不要太有自信說自己能搞定.
將事情放在神的面前.
求祂去幫助.
這是我們可以對神做最好的事情.
我們並不完美.
聖經裡面其實都是這樣形容.
我們需要上帝.
需要神去幫助我們.
當我們去尋求神的幫助.
去更新我們的生命的時候.
我要告訴你.
你將會面對什麼事情.
人生有很多人說.

$^{521}$我經常聽到一句話.
失敗乃成功之母.
但有時候我們做了很多年人.
發覺失敗是否成功之母.
有時候不是.
有時候失敗之後.
都繼續是失敗.
當我們不斷面對失敗的生命.
剛才唱歌也提及過.
如果不斷跌倒.
我們的生命應該是怎樣.
我看到神在我們.
教會的群體.
賜下一個很寶貴的禮物.
就是你身邊的弟兄姊妹.
有人說到.
人與人之間的關係.
起碼有三種.
第一種.
就是大家可能.
為著共同利益而存在.
有時候你見到工作上的夥伴.
合作的拍檔.
或者你公司的同事.
大家都是為著.
共同的目標而存在.
我之前進了很多.
幼稚園備戰群組.
大家進去後.
竟然一個群組有幾百人.
每天有兩千條訊息.
是一個這樣的狀態.
但是一過了考試季節.
群組就全部散了.
因為大家都是為了.
過一過面試.
然後短暫地存在.
第二種人可能是.
比較長的.
大家可能是為了共同興趣而存在.

$^{561}$可能大家喜歡玩音樂.
或者玩體育.
或者喜歡攝影,行山.
那些可能是為了共同的興趣而存在.
可能有時間上的限制.
譬如可能我已經生命進入另一個階段.
要帶孩子,沒那麼空.
去玩攝影或者是玩音樂的時候.
可能我又在那個群組裡面.
去withdraw(退出).
這是第二種我們常有的關係.
但是我要告訴你第三種.
你需要珍惜的關係.
那種關係才可以幫助你在改變的路上.
成為你的支援.
那種關係是怎樣的呢.
那些人是單純.
為了你的.
幸福而感覺到快樂.
我們有沒有這些朋友.
單純為了你的.
幸福而感覺到快樂.
你想想那群人裡面其實有很多.
是我們身邊的弟兄姊妹.
這就是神給我們.
教會很寶貴的禮物.
一個人.
孤身一人.
想去改變他自己.
很容易跌倒的.
但如果身邊的弟兄姊妹跟你同行.
有神跟你同在.
你改變的可能性.
大很多.
所以弟兄姊妹當我們想去改變的時候.
我們千萬不要自己一個人去做.
可以和我們身邊的人一起去走.
為甚麼呢.
因為剛才提到.
當我們生命想去改變的時候.

$^{601}$其實失敗之後可能.
都會繼續失敗.
那怎樣去面對呢.
失敗之後繼續失敗.
我們最常的結果就是放棄.
算了 我活在當下就算了.
信仰的事或將來的事.
到時才算了.
很多世人就是這樣的想法.
然後放棄了信仰.
這個世代很多人就是活在一種.
無力.
生命不要想那麼多.
今天活成怎樣就算了.
但我們在.
基督的恩典裡面.
就算失敗之後.
繼續去失敗.
但神看的.
不是我們生命的成功.
就算我們失敗之後.
繼續去失敗 只要你肯繼續試.
你就顯示出.
你對神的.
決心.
神看的.
不是要你的成功.
是要你對生命.
對他的決心.
和信心.
這個才是神在信徒上.
想要的東西.
這亦都是我們聖經裡面.
所講的那種誠信觀.
宣道會所講的那種誠信觀.
好像你跑馬拉松.
你不一定要是終點.
第一個衝線贏冠軍.
你只要肯繼續跑.
跑到最後.

$^{641}$神就給你獎賞.
那種決心.
是我們信仰裡面.
要顯現出來的信仰特質.
宿命特質.
所以我們需要身邊有弟兄姊妹.
的扶持.
在我們最不濟的時候.
輕輕一聲 親切慰問.
拍一拍我們的肩膀.
不要緊 再來一次.
在信仰的路上.
不斷地再嘗試的時候.
可以向神.
展現我們對他的決心.
這就是他所閱立的.
所以求主這樣去幫助我們.
失敗不要緊.
再嘗試 再來一次.
聖經裡面羅馬叔這樣講.
我們這麼多人.
教會這麼多人.
在基督裡面成為一個身體.
互相聯絡.
作為一個身體 彼此去扶持.
這就是神給我們最大的禮物.
所以神是很願意.
幫我們去改變.
他亦都不斷地去改變這個世界.
他更加賜給我和你.
身邊這麼好的弟兄姊妹.
單純是為了你.
感到幸福.
而他們是快樂.
珍惜我們身邊的禮物.
珍惜神給了我們賞賜.
這樣我們就可以重新再來.
我們一起去禱告.
重新思想聖經的經文.
原來更新並不是一件遙不可及的事.

$^{681}$而是耶穌基督.
你每一天在這個世界.
現在做的事情.
求主真的幫助我們去辨識.
生命究竟是甚麼狀態.
是需要你去更新.
就求你在新的一年.
幫助我們在那些位置裡面.
去改變.
雖然在這段改變的路上.
有很多的抗拒.
甚或乎一開始都不想.
但我們知道當那些特質.
那些雜質被改變.
會帶來一個更好的影響.
你也試試.
弟兄姊妹在我們身邊.
互相去扶持.
在路上或會有跌倒或會有失敗.
但我們要互相扶持.
互相去勸勉.
起身再來.
這就是神你所閱立的決心.
求著這樣去幫助我們.
每一位感受堂的弟兄姊妹.
去經歷神的封聖.
跟著經歷神兒女的福分.
我們這樣禱告.
奉耶穌基督的名字.
阿們.
\newpage



\section{瑪拉基書}
\label{sec:4SEW84zBGCU}
\textbf{2024-02-22 周四晚查經班《瑪拉基書(7)》 講員:曾裔貴牧師}
\newline
\newline
連結: \href{https://youtube.com/watch?v=4SEW84zBGCU}{\texttt{https://youtube.com/watch?v=4SEW84zBGCU}} ~~~~ 語音日期: 2024-02-23
\newline
\newline
\hyperref[sec:myBJVEoimI0]{\small{< < < PREV SERMON < < <}}
~
\hyperref[sec:index_chronic]{\small{[返順時目]}}
~
\hyperref[sec:index_scriptual]{\small{[返順卷目]}}
~
\hyperref[sec:dYh7F8XFvZo]{\small{> > > NEXT SERMON > > >}}
\newline
\newline
$^{1}$這次我們一起祈禱.
感恩天父再一次招聚我們能夠上《馬拉基書》.
我們很盼望透過經文的內容.
認識當時代的背景.
和神向自己的選民的心意.
特別是透過先知的教訓.
讓我們能夠掌握到.
也知道透過聖經.
我們可以明白神的心意.
雖然時代已經是2400多年前.
但是經文的盛載是人的心靈.
所以是不會過時的.
幫助我們能夠明白.
也讓我們真的知道先知.
想要關注的事情.
我們在你面前等候.
禱告奉耶穌基督的名.
阿們.
我們上一個星期.
也試探討第二章第十到第十六節.
這一段我也提過出現了.
關乎到那個背役.
出現了很多次.
或者叫做鬼詐.
另外就是違背聖約.
用了幾個不同的受眾.
祭司的背約.
民間的百姓人與人之間彼此用鬼詐.
代替了諾言.
就是那個承信.
跟著就是婚姻的那個文約的背棄.
就三個層面.
都差不多涵蓋普羅大眾和宗教的層面.
一定是祭司自己都是.
在憲制的事情上.
很隨便.
隨便真的是一個很好的名詞.
但是就是最害死人的一個名詞.
因為我們人生如果什麼都隨隨便便.
就不會有壓力.

$^{41}$更加不用見父道.
但是變成了不能專注.
不能認真.
信仰如果不能認真也是很麻煩的.
所以有時怎樣去釋可而止.
也是一個藝術.
但是這裡無論如何.
是講到第二章這一大段.
是講了很多的鬼詐.
對神的鬼詐對人也會有鬼詐的.
所以就是源於對神的那個約的背棄.
就是不守諾言.
不守諾言就是說.
很普遍在教會界會出現.
美國人他們說笑.
結婚之後原來有一半一世都不離婚.
他們就當是一個笑話來講.
有沒有搞錯.
一生人就是一段婚姻.
就是一半.
當然後來有人告訴我.
其實不是的.
香港現在也是超過一半了.
就是結婚之後的離婚率也是超過一半.
似乎基督徒也不少的.
也多的.
這個也是那個諾言的問題.
上一個星期我們去到第15,16節.
第三個的層面.
關乎到那個婚姻的那個盟約.
我們探討到關乎到婚姻.
聖經怎樣看呢?.
我們就引述過《馬太福音》第19章.
很重的一句話就是.
耶穌說.
因為他們要抹黑摩西.
是他們的權威.
摩西准人休妻的耶穌.
為什麼有些拉比或者一些老師.
是不准人休妻呢?.

$^{81}$耶穌的那句話很重.
摩西因為你們的心硬.
但是耶穌說什麼呢?.
原初不是這樣的.
就是說神原初的心意不是這樣的.
這個就很重.
因為由耶穌親口說.
就是說摩西只不過是一個權宜.
兩害取其輕.
就是盡量避免了.
因為原來那個婚姻的那個盟約.
如果你堅決不准離婚.
不准休自己的太太或者丈夫.
在一個男系的社會就沒有所謂休丈夫的了.
現在就可能會有.
古代是不會有這樣的東西出現的.
所以是用這樣的做單位.
當然偶然間也有何東斯.
也偶然間會有.
但是一般都不會的.
在一個社會都不容許.
但是原來休妻和強暴.
是好像媽生兄弟一樣的.
就是你不准我用這一樣東西.
如果這一樣東西是行人止暴的.
我就用另外一樣東西.
我每天虐待他.
虐待到他死.
虐待到他傻.
虐待到他癲.
看你還准不准我休妻.
是不休的.
不過已經是每天做一些.
完全和婚姻的美滿是背道而馳的.
所以基於這樣的情況.
摩西某程度上也有很大的屬靈的智慧.
也擔心如果有些猶太的婦女.
日後在家囊的時候.
就會有這麼多的建議私遷.
這麼多的婚姻問題.

$^{121}$倒不如准許可以休書休妻.
但是似乎都要有一些法律手續.
也就是說都可能需要給一些民間的長老.
來去仲裁.
這個都是某程度上的一個小小的關卡.
去截一截.
不是這麼隨便這麼容易.
如果這麼隨便容易.
就超過一半以上.
可能九成的了.
這一頭結那一頭就來了.
就變成了一個很複雜的問題.
但是回來看聖經的時間.
主耶穌.
我再講第十九章這一段.
因為是耶穌基督引述神起初的心意.
這個就很重對於婚姻的一個盟約的維繫.
十九章這裡這麼說的.
有《巴釐塞人》第三節.
馬太芬的十九章.
有《巴釐塞人》來試探耶穌.
人是不是無論什麼緣故都可以休妻呢?.
耶穌就說起初做人是做男做女.
並且說因此人要離開父母與妻子聯合.
二人成為一體.
主耶穌是引創世記.
摩西都沒有出現.
都沒有出埃及.
還是阿當夏娃的年代.
神直接和阿當講.
當時只有一對夫婦.
只有阿當和夏娃.
沒有第三者.
我想問一下.
當夏娃生了那個人.
那個人的太太是誰來的?.
誰生的?.
哪裡爆出來的?.
這些是聖經難題來的.
如果是夏娃生就是她的妹妹.

$^{161}$拆開來講一下.
我主要講創世記這一段.
你們沒有想過這些聖經難題嗎?.
我們先回頭看一看.
耶穌引的創世記第二章這一段.
並且說因此人要離開父母.
與妻子聯合.
二人成為一體.
這經是聖經的權威.
不只是摩西的一些建議.
這經你們沒有念過嗎?.
既然如此.
夫妻不再是兩個人.
乃是一體的了.
所以神配合的人不可分開.
法利昭人說.
這樣摩西為什麼分父給妻子休書.
就可以休她呢?.
《申命記》的二十四章.
法利昭人有一個反對耶穌的.
這麼權威的宣告.
都來提出一個質疑.
很強的理據.
他們是引《聖經》《申命記》.
為什麼後來摩西又准人這樣寫休書.
是有背後寫在聖經裡面的.
第八節就很重要了.
耶穌說.
摩西因為你們的心硬.
所以許你們休妻.
他只是說是一個權宜.
但起初是神創立世界的時候定的秩序.
但起初並不是這樣的.
我告訴你們.
凡是休妻令娶.
如果不是為了淫亂的緣故.
就是犯了姦淫.
有人娶那個被休的婦人.
也是犯姦淫的.
接著門徒都聽得很明白.

$^{201}$肯定了耶穌所說.
原來起初神的心意是最高的那個旨意.
那個婚姻是有那個.
在神眼中的那個永恆的效力的.
至少去到死亡才可以解除那個婚約的.
所以這裡就告訴我們知道.
我們要明白.
到底婚姻在神眼中那個最高的旨意.
但始終人們面對人生不同的處境.
有不同的情況.
所以摩西就去到人生的木樣的時候.
就准許以色列人可以寫休書.
也就是說都給他有少少的限制.
如果完全沒有限制.
就更加形同虛設了.
所以這裡就讓我們看到.
這個是在耶穌時代.
猶太人真的碰到這個這麼倫理的課題.
這麼敏感的課題.
是逼耶穌親自去回答.
到底神的心意是怎樣的.
耶穌就回答了.
也都在某種程度是帶來一個.
對基督徒日後他對婚姻都有這一份的尊重.
除非你不尊重耶穌所說.
不然的話都是一個很嚴肅的課題.
後來到底使徒們又怎樣理解婚姻呢.
我們又要知道.
在什麼情況下.
有什麼選擇可以有離婚呢.
哥林多前書的第七章.
直接解這個問題了.
就是Q and A問題解答了.
林前的第七章.
肯定是保羅建立教會之後.
有教會因為他離開了.
去到其他地方傳道.
有教會的人士要處理一些教會上的.
可能紛爭.
可能是一些問題.

$^{241}$但是不懂處理.
就要問保羅應該怎麼處理呢.
又或者在當時他們成為了教會了.
但是很多未信主之前做的事情.
現在應該怎麼解決呢.
這些其實我們不斷都會發生的.
比如我以前是倒馬的.
現在還能倒嗎?.
買六合彩的.
金多寶成多少億啊?.
千多萬啊?.
是不是數進去.
擺張電腦票.
中了我不要自己要.
我給教會的.
行不行呢?.
那可不可以呢?.
這些很多倫理的問題.
其實是要解答的.
很多人都未必知道應該怎麼處理.
到底那個原則應該怎麼應用呢?.
這些都是每個時代都要面對的.
這裡第七章就這麼說.
林前的第七章.
明顯是有人寫信給保羅.
保羅現在不是哥林多教會.
有人寫信給保羅.
保羅收到信之後.
當然一定是教會的領袖寫信給保羅.
第七章.
論到你們信上所提的事.
我說男不近女都好.
這就是說我們只是有答案.
就不知道那個問題是什麼.
因為保羅就沒有將問題說出來.
就在字裡行間.
將他們問的關心的事情.
就來去解答.
那就是說現在他們在問一件事.
到底婚姻.

$^{281}$那就是似乎暗示.
哥林多當時的社會.
對於男女婚姻.
都是有一定的混亂.
所以信了主耶穌.
就要面對這個處境.
教會的牧者可能未必懂得怎麼處理.
就需要保羅給一些指引來去解答.
還有到底神的心意是怎樣呢?.
第二節.
但為了要免淫亂的事.
男子當各有自己的妻子.
女子也當各有自己的丈夫.
丈夫當用合兒之分待妻子.
妻子待丈夫也要如此.
那就是透過字裡行間.
讓我們看到不是單方面的一種服侍.
一種男尊女卑.
是平等的.
因為這樣說.
對當時的社會處境.
那當時的倫理是不可能發生的.
男的永遠是在外面.
包括在家裡都是權威的.
女性在當時的社會地位是非常低的.
只會成為男性.
或者男系的社會的一個附屬品.
但是保羅在這裡告訴我們.
女子也當各有自己的丈夫.
丈夫要用合兒之分待妻子.
妻子待丈夫都要合兒的.
大家是平等的.
在婚姻這個約裡面是平等的.
這是很石破天驚的.
對當時代的人來說.
是很特別的說法.
就是說保羅變得很前衛了.
第四節.
妻子無權並主張自己的身子.
乃在於丈夫.

$^{321}$丈夫也沒有權並主張自己的身子.
乃在妻子.
似乎說得很暗晦的.
這裡說的是男女夫婦之間的性行為.
那個生活.
就是你不可以任何一方去操縱.
去讓對方得不到婚姻裡面的歡愉.
那個滿足.
而帶來其他的問題.
然後就告訴我們知道.
一直都說的是男女.
丈夫和妻子是一個平等的看待.
如果不平等就會出現什麼問題呢?.
第五節.
夫妻不可以彼此的虧負.
就是說他現在這裡解答的問題就是.
如果正常的.
男大當婚 女大當嫁.
都是保羅的意思.
那你就應該要有進入婚姻的考慮.
進入婚姻的等候.
然後找到合適的對象.
就要進入婚姻裡面的.
那個約裡面的夫婦生活.
共同的面對人生不同的處境.
特別是強調.
為什麼這麼說呢?.
第五節就說了.
夫妻不可以彼此的虧負.
虧負什麼呢?.
除非兩雙情人暫時的婚房.
為了專心禱告.
以後仍要同房.
原來信了主耶穌的.
都是為著一時之間的那個專注.
但實際上既然兩個人是夫婦.
就要有夫婦之間的生活.
作息和大家的共同的興趣.
共同的攜手走人生的路.
達至那個.

$^{361}$當然信了主耶穌.
保羅在這裡說的就是.
尋求的是那個生命的理想.
就是說兩個人.
不可以太過一個走得很前.
一個就躲得很後.
這樣的婚姻.
特別在信仰上面.
就是一個很大的困難.
所以這裡特別強調.
如果其中一方很強烈.
我要祈禱.
我要不停地祈禱.
我要不停地自己一個人躲起來.
我要去.
那時就沒有那些.
引修啊.
去到做一些修士啊.
那時沒有的.
是去到第三,第四世紀才開始有的.
保羅的時代.
會不會有些人會有些奇奇怪怪的.
神學說法.
而導致有些扭曲了的.
對婚姻.
尤其是對性的那個問題.
會不會有些人會這樣呢?.
保羅就很.
你可以說他很持平.
他說專心禱告是應該的.
暫時的不要想男女.
或者日常世上的生活.
暫時擺下那些東西.
但是日後呢.
以後仍要同房.
免得撒旦趁著你們情不自禁的時候.
就引誘你們.
那就是說.
保羅說得很暗晦.
這個情不自禁呢.

$^{401}$就變成了.
可能是去到另外一種.
不合法的關係了.
可能就是說.
如果其中一方.
認為自己的身體是神性的.
不可以被玷污的.
那如果這些扭曲的神學呢.
就會帶來.
那個正常的夫婦之間的關係.
就會有一種的.
就是出現了一種這樣的矛盾.
那就會導致.
有另外一種的情況就會出現的了.
這個就是保羅提醒.
哥林多教會的弟兄姊妹.
要留心這件事.
不過下面就告訴我們知道.
第六節.
我說者說.
只不過是准你們.
不是命令你們.
我願意眾人像我一樣.
似乎這裡暗示.
保羅是守獨生的.
似乎暗示而已.
這個就永遠都是一個原案.
因為保羅沒有直接說自己是否獨生.
就是從來都沒有結過婚.
但是有些人就說.
他在公會裡面.
他是去到一個.
做一個指證的人.
那似乎是一定要結了婚的.
這樣的位份才可以這麼做的.
所以兩派最終都沒有人知道的.
要知道的就是將來我們回到天家.
就見到保羅就問他.
其實是不知道的.
他說你們就要.

$^{441}$我是准你們.
不是要命令你們.
第七節.
我願意眾人都像我一樣.
只是個人的領袖.
這個就要平衡了.
有些人真是為神的國度.
他守獨生.
終身都不結婚的.
譬如我自己的一個.
我初信的時候.
信主的時候.
一位傳道人.
後來有機會接觸過.
在教會裡面的一位姐妹.
都是很好的.
後來接著輾轉.
我知道他們.
都雙方面很清楚.
大家都不可以再繼續了.
這個牧者.
就是我的傳道人.
他到現在都沒有結婚的.
就是終身為神傳道等等.
對他來說是一個.
我覺得是挺自然的.
譬如我們認識的陳琴開牧師.
她是獨生的.
沒有結婚的.
已經很大年紀了.
事奉了很多很多年了.
我們晚上上神學課程的那位老人家.
那位傳道牧者.
退休很多年了.
殷存堂的很資深的牧者.
是裴靈慧的一位牧者.
她是獨生的.
這裡告訴我們知道.
保羅說有些人是這樣.
有些人是另外一樣的.

$^{481}$所以這裡就不可以勉強.
每個人.
所以婚姻絕對不是一個負累.
如果你個人認為婚姻是一個負累.
那可能就是一些個別的問題.
有些人就覺得自己逍遙自在.
可以一個人去事奉.
是很開心,很舒服.
可以怎樣去怎樣都可以.
其實我們應該訪問.
黎晨牧師.
有其他的同工.
譬如曾姑娘.
有些他們都還沒有結婚.
譬如仲妍姑娘現在還年輕.
其實應該訪問一下他們.
他們說一下他們的心路歷程.
我們作為男性.
又已婚.
又很難表達到保羅的心境.
保羅似乎在這裡暗示.
你們不要學我.
每個人是不同的.
他很自豪自己還沒有結婚.
但是他就說到.
婚姻裡面的張力.
他說到婚姻美滿就是.
彼此要看對方不是附屬品.
看對方是要尊重對方.
在這個約裡面的.
大家的平等.
你說我.
不要介意.
我開玩笑的.
我家裡的壞人.
他還沒有信耶穌的.
他怎會是我們約的平等.
第八節就處理這個問題了.
第八節就處理家裡面.
還沒有結婚的對方了.

$^{521}$怎樣能夠平等呢.
這個就出現了一個問題了.
我對著那些沒有嫁娶的.
和寡婦.
寡婦就是其中一方已經離世了.
若他們嘗嘗我就好.
果然又繼續說他的毒身.
倘若自己禁止不住.
就可以嫁娶了.
與其慾火攻心.
倒不如嫁娶為妙.
所謂慾火攻心.
就不是說經常想一些污糟的東西.
不是這個意思.
而是真的有那個生理需要.
就變成了.
如果你整天都只是想著這件事.
覺得自己很孤單.
很其他的欠缺.
必須要有另外一半去填滿你的心靈.
這些就是慾火攻心的描述.
所以.
否則就會去到一種更加扭曲.
就不健康了.
所以保羅就說.
在這樣的狀態下.
你的心都不是想著神的國度.
你的身體就變成了一個毒身.
保羅就說這樣不健康了.
不如你去嫁娶更好.
所以.
大家都認識的楊榮誠姑娘.
楊榮誠院長.
就是郭文賢院長的前一任院長.
楊榮誠院長.
一定是很熟悉的.
和我很認識的.
她其實一直都是毒身的.
後來轉轉轉.
就想著不可能再結婚了.

$^{561}$她就轉轉轉.
移居到蘇格蘭之後.
就和一位神學家.
外國人來的.
很出名的.
舉世出名的舊約的學者.
結婚.
當然她的丈夫都離開世界了.
安息主懷了.
所以不是說你真的覺得不可能的.
原來去到某個年紀.
都可以再.
她不是再婚.
她都可以結婚的.
所以有時神的時間配合.
又不是我們人可以知道的.
因為我認識她很久了.
很早期認識她的.
是我以前沒有會的牧者.
是相當有才華的一位姐妹.
在這裡告訴我們.
我們現在討論的就是婚姻這個課題.
反正我們碰到馬拉基斯.
我就要將新約出現過.
關於婚姻的處境.
和大家一起去討論.
我們就會畫面更加闊.
我們更加知道神如何看婚姻這件事.
如果真的剛剛說中你的案件.
你自己千萬不要介意.
聖經是有解答的.
聖經是給我們出路的.
聖經不是要批評我們的.
聖經是想我們明白神的心意.
所以這裡告訴我們.
保羅就說.
倒不如嫁娶為妙.
第十節.
至於那已嫁娶的.
我吩咐他們.

$^{601}$其實不是我吩咐.
是主吩咐.
重不重?.
很重.
之前就說.
我講我的意見給你聽.
我未必有主的命令.
但現在不是我吩咐你.
是主吩咐你.
吩咐什麼呢?.
妻子不可以離開丈夫.
第十一節.
若是離開了.
不可以再嫁了.
或是仍和丈夫和好.
丈夫也不可離棄妻子.
似乎是在未信的時候的狀態.
不是在信了主之後的情況.
即是說.
未信主之前.
已經出現了這種狀態.
譬如已經離婚.
是這樣的狀態.
似乎這些告訴我們.
如果離開了.
就不可以再嫁了.
或是仍和丈夫和好.
丈夫也不可以離棄妻子.
我對其餘的人說.
不是主說.
倘若某弟兄.
開始有一個案例.
講一些例子出來.
倘若某弟兄有不信的妻子.
似乎真的了.
婚姻已經有了.
兩個都是未信主.
但後來丈夫信了主.
有不信的妻子.
妻子知道丈夫信了耶穌.

$^{641}$妻子也情願和信了主的丈夫同住.
他就是男的丈夫.
不可以離棄妻子.
你未信耶穌.
你將來落地獄.
我不和你一起了.
普羅話不可以.
因為那個盟約.
雖然在大家都未信之前建立了.
但在神眼中的那個約是一路生效的.
要留意這一點.
13節.
一樣了.
相反.
妻子信了耶穌.
似乎這個案例.
香港人會比較多.
通常是太太比較容易信耶穌.
丈夫就比較難.
所以一般教會都是姐妹會多.
弟兄會比較少.
很平衡.
說了丈夫之後就說信主的太太.
即是說太太去完國陪禮報道會.
或者有其他的地方信了耶穌.
跟著回到家就面對一個未信的丈夫.
那個衝突就很大了.
未信的丈夫就是那個生命未改變.
當然太太一樣的未信.
他的生命未改變.
而你的生命改變了.
因為信了耶穌.
開始衝突就在這裡發生.
但是感恩的地方就是.
保羅在這個約裡面.
他涵蓋的是信和未信.
都在這個約裡面一並的討論.
原來這個約.
即是婚姻這個約.
是可以超越信和不信的狀況.

$^{681}$所以你會留意到.
如果妻子信了.
她就不可以離棄太太.
第十三節.
妻子如果有不信的丈夫.
其實是她信了耶穌.
丈夫都情願和她同住.
即是我意識到.
我知道你信了耶穌.
但是我仍然尊重和你的婚姻.
或者這個夫婦這麼多年的老夫老妻的關係.
我仍然和你同住.
我不會阻礙你回教會敬拜神.
之如此類.
當然你不要阻礙我禱罵.
我禮拜去哪裡.
可能會是這樣的.
她就不可以離棄丈夫.
第十四節.
這一節聖經就很寶貴了.
因為不信的丈夫.
就因著妻子成了聖潔.
並且不信的妻子.
就因著丈夫成了聖潔.
要解一解這一節.
什麼意思呢?.
是不是一個人信主一家都可以救?.
即是說妻子信了耶穌.
丈夫信了耶穌.
妻子和丈夫還沒有信.
自然就上天堂了.
成為聖潔.
所以其實不是這個意思.
因為一直上文處理的就是.
那個約不可以隨便的.
因為其中一方信.
立即去毀掉婚姻的文約.
反過來.
本來這十四節所講的.
正正因為這個約在神眼中是有效的.

$^{721}$所以未信的丈夫.
未信的妻子.
怎樣去因為信了的丈夫和妻子.
成為聖潔呢?.
而正正就是.
那個其中一方的生命改變了.
他可以每天為自己未信的丈夫妻子去禱告.
你的生活行為.
或者你認識的基督徒主裡面的弟兄姊妹.
可以慢慢去感染.
將這一份聖潔的生命.
去感染未信的丈夫或者妻子.
他們就因為這樣的緣故.
慢慢慢慢對福音就不抗拒了.
就可以去達至他信主而獲得生命的聖潔.
就是這個意思.
就是這樣的處理.
所以就不要說立即.
你這個未信的人.
抵你下地獄了.
現在我要和你解約.
這樣就不是聖經的教導的辦法.
沒錯.
一了百了.
他不再什麼了.
大前提是未信的那一方.
是容許你繼續有你的新的信仰生活.
因為下面第十五節就說.
未信的丈夫或者妻子.
是頂不住你信了耶穌.
而他要主動和你提出離開.
是兩個案件.
十五節就是這裡告訴我們.
倘若那個不信的人要離去.
就由他離去吧.
無論是弟兄還是姐妹.
遇到這樣的事情都不必拘束.
神召我們.
原是要我們和睦.
你們有沒有留意到.

$^{761}$要我們不是聖潔.
抵他死就走.
快點走.
用鎖拍你走.
不是.
神要我們和睦.
如果那個主導權在對方不信的一方.
譬如他很堅決.
很反對信仰.
而你本來很愛你的.
但是你居然信了耶穌.
我就不能夠接受了.
我堅決要麼你要選耶穌.
要麼我就走.
或者你走.
那麼呢?.
保羅居然教人分妻.
十五節說.
倘若那個不信的人要離去.
就由他走吧.
不要拘束.
所以你要很小心.
看上下文.
如果你拿著這句十五節.
你就回家和丈夫說.
保羅告訴我的.
拜託你走了.
明天律師樓見.
那不是保羅的含義.
因為上下文有很多的討論.
大進這裡.
他在說不同的處境.
現在這個處境就是.
對方他不信.
他堅決對你的信仰.
他極致的反感.
反感到他完全頂不住你的存在.
所以我們又要小心.
不要因為我們一信了耶穌.
過分的熱情.

$^{801}$我們就不定.
好像在我們未信的家人面前.
變成一個神婆.
變成一個.
都不知道是發了瘋.
這樣是我們自己製造出來的.
那個敵視.
就不好了.
我都曾經說過.
我初信了主耶穌之後.
那時回到星期六的查經班.
星期六的青年團體.
查事司記.
查到哪裡呢?.
就是這個.
機電.
他們去拆.
家裡的偶像.
我完了之後.
心裡很熱切.
回到家裡看到媽媽裝香.
整個香爐扔掉.
真的扔掉.
罵得我快瘋了.
有時你過度初信.
過度熱心.
就變成有時.
我們不懂聖經.
十五節就說.
不必拘束.
本來怎麼說的.
神呼召我們.
原是叫我們和睦.
在婚姻這個狀態下.
更加要和睦.
所以這個和睦.
原來是用愛去接納對方.
是讓對方.
慢慢明白我們信仰的真諦.
所以就不是用那些很厲害的.

$^{841}$祈禱.
本來開飯的.
十幾人等你.
你就是這樣祈半個小時.
都不開飯.
那丈夫完全沒有面子.
不休也很難.
有時我們千萬不要變成.
那種近乎無知的熱心.
但是.
聖經是給我們有教導的.
所以十五節.
不要拘束.
他如果堅決要離開.
就由他離開罷了.
神呼召我們.
原是叫我們和睦.
即是說你做了你的部分沒有.
原來你一信了耶穌.
崇拜又回.
回完就回主一學.
回完都不夠.
下午又在教會.
接著.
五六點都不走.
帶上女兒和兒子.
接著丈夫.
又週末有空.
接著你都不見你.
人影都沒有.
這樣的時候經年累月.
對那個未信的丈夫或太太.
是一個.
的冒犯.
所以其實我們要平衡.
不是教大家不要回崇拜.
不要什麼.
不是我們要平衡.
我們要給到.
未信的家人.

$^{881}$感覺到這份和睦.
本來真是很有智慧.
是神的啟示.
神召我們.
在現在的衝突里.
在這個.
的.
矛盾里.
是要保持和諧的和睦.
就是說你要.
放下一點點身段.
遷就未信的家人.
但是遷就到一個地步.
他要你否認信仰.
這些都不能遷就.
所以保羅就說.
不必拘束.
他走就由他走.
因為去到那個底線.
要你否認耶穌.
那你沒有可能.
只能夠他離開.
保羅就告訴我們.
是可以離婚.
這個保羅是這樣說的.
等於.
又是經過了很嚴肅的.
處理之後.
某程度.
任由對方去寫休書.
由對方做主導權.
就不是由我們.
信了主的一方做主導權.
去寫休書.
大家明不明白這一點.
當然見不到大家的額頭.
現場的就知道.
這就告訴我們.
第十六節.
你作妻子的.

$^{921}$怎麼知道不能夠救你的丈夫呢.
你作丈夫的.
所以兩邊這樣說.
兩邊這樣說.
在第十六節.
居然反過來說.
你不要這麼快.
隨隨便便就說.
由他走.
你怎麼知道你的生命的改變.
不能影響你的未信的丈夫或者妻子.
他可以透過你.
得到生命的改變.
而信住呢.
所以不要這麼快.
三兩個月就說意見不合.
他也是拒絕福音的.
算了算了.
甚至你都沒有用盡.
你自己的能力.
去改變他.
揭一揭.
先給前書.
第三章.
聖經教我們.
怎麼去將那些環石點頭.
彼得自己本身.
希利文就是石的意思.
石仔.
Petra.
Peter這個字本身就是石的意思.
第三章的彼得前書.
教當時的弟兄姊妹.
在家庭里.
怎麼來到有好的榜樣.
你們作妻子的.
要說得對.
你們作妻子的.
要順服自己的丈夫.
這樣如果有不順從道理的丈夫.

$^{961}$他們雖然不聽道.
都可以因為妻子的品行.
被感化過來.
原來妻子是感化官.
是感化自己的丈夫.
原來妻子是一個.
人肉的廣章.
廣道.
所以你不在家裡不停播.
《曾沐斯的道》.
就可以改變你的丈夫.
你的生命自己改變.
這裡告訴我們.
因為妻子的品行.
可以感化他.
他不聽道.
他也不順從道理.
第二節.
這正是因看見.
你們有.
貞潔的品行和敬畏的心.
所以這裡特別提到.
太太.
如果心靈是很溫柔.
很有那種愛神的.
敬畏神的心.
這種是一種裝飾.
是裡面的裝飾.
裡面的裝飾.
正是因為你聽道.
正是因為妻子聽道.
聽道就改變心靈.
改變心靈.
就變成裡面的裝飾.
越來越有好的品行.
越來越有貞潔的心靈.
敬畏神.
這樣不知不覺可以令自己的舉止行為.
影響未順的丈夫.
這是彼得所說的.

$^{1001}$第三節.
剛剛相反.
你們不要用外面的.
邊頭髮,戴金色,穿美衣.
為裝飾.
只要裡面全著.
裡面要全著.
全著即是說.
不是一兩天的.
你聽完培靈會是有改變的.
樣子也好看很多.
連續十天的培靈會.
但是要全著.
即是說你遇到一些事情.
你遇到每天.
那些繁瑣的生活.
你都能夠慢慢表現出.
你的真正虔誠.
你的丈夫是有眼看到的.
回來.
對你呼呼喝喝.
你仍然沒有改變.
完全面不改容的.
很客氣.
很溫柔,很有愛心.
那個丈夫會慢慢發現.
咦?裡面的神可能是真的.
他會相信的.
所以這裡告訴我們.
只要全著長久溫柔.
安靜的心.
為裝飾.
在神面前是極寶貴的.
因為古時.
養乃神的.
即敬畏神的聖潔婦人.
都是用這些作為裝飾.
就是長久的溫柔.
安靜的心靈.
來裝備自己的內心.

$^{1041}$即是說遇到任何的處境.
你裡面都會有平安.
你裡面都會有一份喜樂.
你都會有一份對神的盼望.
你都有對神的信心.
這樣就是長久的溫柔裝飾.
這樣就會帶來.
對其他的人.
譬如第六節.
就如撒拉聽從阿巴拉罕.
稱他為主.
你們如果行善.
不因恐嚇而害怕.
就是撒拉.
的女兒.
所以原來彼得都教我們.
如何做太太.
不知他太太是怎樣.
他教這些妻子.
即是說.
是肯定可以成功的方法.
即是說.
特別在古代.
如果你能夠這樣對自己的丈夫.
因為他下面講的就是.
那個恐嚇.
害怕.
即是說當時代的男性.
那個主導權.
是非常嚴重的.
但是那個妻子.
不是逆來順受.
而是從心裡面.
信服自己的丈夫.
從心裡面因為敬畏神.
愛神而信服自己的丈夫.
這個就是一般的社會.
都不會具備.
所以其實教會就是充斥著.
太多敬虔的姐妹.

$^{1081}$撐起著很多的家庭.
很多的倫理.
其實很多的丈夫.
信主都是因為太太的生命.
那種的敬虔.
所以真正的敬虔.
是在講內在.
當然內在重要.
不要外在就真的什麼都沒有.
你每天都.
逢頭散髮.
這樣就不行.
即是你走出去.
丈夫都沒有面子.
就真的覺得很受不了.
所以愛妻.
是要平衡.
所以千萬不要極端.
兩者要平衡.
聖經有教很多的裝飾.
不過又軟了.
不夠時間看這麼多.
我們回來第七章的哥林多前書.
還有一點點都講了.
我們先進入馬拉基書那一段.
其實這樣這些事情.
已經有很多的討論.
其實這樣這些事情.
已經有很多的討論.
其實這樣這些事情.
十七節.
七章的十七節.
只要照主的吩咐.
按各人所得到的.
來做.
和神所召各人去行.
我吩咐各教會.
都要這樣.
有人已經受國禮.
蒙召.

$^{1121}$就不要廢棄國禮.
即是他本身是生下來猶太人的.
那你說我信耶穌.
我不行國禮了.
原來在禮貌上.
他都要行國禮.
不過就不是靠國禮得救.
也不是靠國禮成為優越的等級.
其實任何一個人在神的眼中都是平等的.
所以不因為你的國禮.
因為以前就是.
猶太教.
你不受國禮.
你就不可以敬拜神.
所以外邦人要受國禮.
但是現在保羅說.
沒有啊.
真正的國禮就是在主耶穌裡面.
得到的那個永恆生命.
那個是真國禮.
所以保羅就告訴我們.
雖然那個意義已經不大了.
但是如果有人蒙召的時候.
是受國禮的.
就不可以來廢棄國禮.
有人未受國禮蒙召.
就不要受國禮.
就是說我們外邦人不需要.
除非是衛生的條例.
如果不是就不需要行國禮.
我們是外邦人.
根本不需要這一件事.
所以受國禮.
算不得什麼.
不受國禮都不算不得什麼.
都是沒有意思的.
這裡其實可以.
延伸到其他的地方.
就是我們在主耶穌裡面.
大家都是平等的.

$^{1161}$平等裡面.
不是每個人都要一樣.
有些人的學識是高一些的.
有些人的社會背景是高一些的.
有些人的家庭背景是高一些的.
那個是神安排的人生.
但是信主個個都是一樣的.
個個都是要向主交仗的.
個個都是要為主而事奉的.
你的事奉的足跡可能會多一些.
有些人可能是少一些.
所以每個人都不同的.
我都提過.
其實那位姐妹是沒有什麼特別的.
她們是沒有什麼特別的.
但是其實那位姐妹是會聽我們的講道的.
她會聽的.
都提過.
每個人都有不同的處境.
她原來是2018年的.
很嚴重的車禍.
我勉勵她.
現在你的處境.
沒錯你是四肢.
已經不能夠有任何活動.
而且那種痛是基本上是沒有停止過.
但是現在你.
既然開始尋求神.
神一定有祂的心意.
在你的生命上面.
我們每個人.
我們可以四肢活動.
我們可以有很多不同的事奉.
神將來會按照我們不同的跑道.
人生的跑道.
來賞賜.
和來審判的.
所以每個人都不同的.
如果你的足跡.
你涵蓋的範圍.

$^{1201}$是能夠影響到國家元首的.
那恭喜你.
你就可以有更大的空間.
但是神對你的要求.
又會再大一些.
我們作牧者去到不同的地方.
今天才和區內的NGO開會.
有什麼合作的空間.
每個人都有不同的事奉的領域.
總之.
每個人都一定要忠於.
自己神的托付.
保羅在這裡一直在講婚姻.
不要勉強.
你不要羨慕我獨身.
保羅.
我也不會羨慕你有婚姻.
每個人都有他的處境.
其實神給的份.
每個人應該要滿足於.
神給你的份.
這裡告訴我們.
信了耶穌之後.
一定要問自己.
神你給了我什麼份.
我要如何去服侍你.
在我們很多的限制下.
我們如何能夠在有限里.
力求的進步呢.
這些都是我們應該要去追求的.
不要滿足於.
自己現在的.
那種現狀.
就變成不斷地原地踏步.
這就是不思上進.
這就不好的.
基督徒應該是不斷前行的.
不斷學習.
有姐妹說想學英文.
我很開心.

$^{1241}$每個人都有不斷去探求的.
所以不要覺得自己很不足夠.
每個人都有不同的處境.
總之是不斷地成長.
是神喜歡的.
這裡告訴我們.
留意一下.
他其實是在說.
蒙召的時候的身份.
第20節.
各人蒙召的時候是什麼身份.
就要守著這個身份.
你是做奴隸蒙召.
就不要因為這樣來憂慮.
如果能夠自由.
那當然好.
就求自由.
因為作奴隸蒙召於主的.
就是主所釋放的人.
你說多麼超脫.
當時是奴隸制度.
你是奴隸信了耶穌.
你要記住.
社會不容許你離開你的主人.
但是你在神眼中.
你是自由的.
你想想原來我們為什麼沒有.
這份信了主之後的幸福感.
信了主之後的那種.
榮耀.
就是我們太少聖經的知識.
原來在保羅的眼中.
你想想他是一個.
很極端的白尼賽人.
是看不起任何外邦人.
他是極之自己.
有嚴格的對自己的要求.
看不起其他猶太人.
不守律法.
但是他現在因為耶穌改變他的生命.

$^{1281}$他看其他人.
每個人都有自己的位份.
每個人有自己神安排.
他的人生.
不要羨慕.
但是自己要鞭策自己.
保羅就告訴我們.
原來那些奴隸是主所釋放的人.
作自由之人蒙照.
就是基督的奴僕.
又是很重了.
你是一個自由人.
當時有這些階層的.
Free man.
有些甚至是.
奴隸主.
養了很多奴隸在家.
是社會階層很高的.
羅馬的貴族.
但是保羅說.
你信了耶穌.
你在神眼中.
你是要服事基督的奴隸.
你是基督的奴僕.
你不是自由的.
所以我們記住.
我們六天可以做我們自己喜歡做的事.
但是第七天.
我們每個人都要做.
但是第七天我們每個人都是神的兒女.
去敬拜神.
這個是神規定了的.
所以這裡告訴我們.
第23節.
你們是重價買來的.
每個人在神眼中.
所以信仰.
真是帶來一個很大的盼望.
所以如果你在地上.
你的人生.

$^{1321}$是感受到很大的捆綁.
教會生活就比較好.
因為教會生活.
教會生活就給你很大的釋放.
因為你六天.
都是在一種壓迫下.
所以為什麼.
雖然我們只是一個觀察.
譬如那些非裔的姐姐.
印尼語.
或者泰國語.
又或者是美國的黑人.
他們信主.
那些教會是特別活潑的.
因為他六天.
那種身份是很低下.
甚至受盡剝削.
我絕對不說香港.
不是說非庸受剝削.
不是這個意思.
我說非洲那些在美國的黑人.
所以他們的教會生活.
他們真的能夠完全擺下.
他們暫時的身份.
他們回到一個很重要的身份.
就是被重價買來的基督徒.
在神眼中是重價.
在你的那些.
社會階層.
雇主的眼中.
你是很低下的.
是剝削的.
現在好些了香港應該就沒那麼嚴重.
以前可能還是差一些的.
那你可以想象得到.
保羅說的就是重價買回來的.
就不要做人的奴僕.
弟兄們.
你們各人蒙照的時候.
是什麼身份.

$^{1361}$仍要在神的面前.
守住這個身份.
其實它意味著.
是想說婚姻的那個弱.
那個問題.
就是說你不要覺得.
早知道就不結婚了.
早知道就不嫁給這個壞人了.
早知道就什麼什麼.
保羅說不要早知道.
反而就是你信了主耶穌.
現在神眼中看你是重價的.
是主耶穌為你附上了.
那個屬價.
來得到神的公義.
所以在神的眼中.
他都會.
好像一個套件.
來結婚的.
因為你的緣故.
你的丈夫你的妻子都會成為聖潔.
所以千萬不要用那種.
很對立的眼光.
應該回頭.
我不知道完了歲月之後.
你回家或者關了電腦之後.
你望一望你未信的丈夫.
未信的太太.
你會不會突然間那些眼光不同了呢?.
聖經是這樣告訴我們的.
我們要記住我們是重價買回來的.
我們應該要來到去.
看我們的照明.
神都會知道的.
我都提過了.
我在牧區長大.
我那時候出世.
已經是十三兄弟姐妹.
我排第九.
去到我第七歲八歲的時候.

$^{1401}$我是不知道的.
我是長大之後很多年後才知道的.
早知道我就不在香港了.
有一個很有錢的.
我爸爸的.
不知怎樣認識的朋友.
說明瞭.
要收養我.
帶我去加拿大.
我當然不知道了.
知道了我就走了.
你想想幾十年前.
坐飛機都沒有人能坐.
一定要收養我.
轉轉我媽當然不肯.
雖然十三個沒有了一個.
始終是自己的骨肉.
整件事就這樣掩蓋著.
沒有人知道.
到我長大才告訴我.
我姐姐就知道.
那個梁太就住在我家附近.
我姐姐在加拿大多人多.
我去探望她的時候.
梁太當年收養你的那個.
她不就是住在附近.
驚到我快要死了.
如果我八歲知道我可以.
去加拿大被人收養.
不在家裡還沒有屋區.
很低下階層.
很窮.
那時候我是排第九.
現在也是排第九.
是受壓迫的.
是被哥哥.
姐姐就很少的.
是被哥哥欺負的.
有什麼事就要你做.
我要欺負弟妹就會被人罵.

$^{1441}$被人打.
我是剛剛夾縫.
所以那時候的成長是很粉碎.
執俗的.
塑造了我讀中學就被踢出校.
讀不到書.
是有家庭的影響的.
所以.
如果我信了耶穌回望.
我的父母.
都是沒有信耶穌的.
但是這樣訓練.
我這個成長.
我回望就很精彩了.
所以我不知道.
假如我成為了.
加拿大人在那邊讀書.
讀大學 博士.
我又會是怎樣的人呢.
我不知道因為我回不了頭.
不是拍戲的.
有兩個故事的結局.
現在只能夠說.
這段經文就醫治了我的內心.
神看我是.
重價的.
神知道我的原生家庭是這樣的.
因為他是永恆地.
看事物的.
我們只能夠看到面前的那一秒鐘.
和回望.
我們的成長.
但是現在你這樣回望.
我就感受到神的.
保守.
眷顧.
我在神的眼中是重價的.
所以.
弟兄姊妹如果你對自己的.
原生家庭有很多的.

$^{1481}$屈節.
有很多的不開心.
甚至很憤怒.
不要了.
因為你的人生.
很有限的.
你應該要感恩.
懷著感恩看你自己的從前.
一定是背後有神的.
安排.
肯定是有神安排的.
人家開口只要找我.
都沒有了這條路.
那就是說.
一定是神管理著.
就不是我的父母.
盟瑟.
所以我們要相信這個揀選.
神有祂的帶領.
所以這段聖經其實很寶貴.
其實再說下面.
就是說未婚和寡居.
不過.
夠不夠時間我們再看下去呢.
只是看這裡已經看了半個小時了.
就是逐節這樣去解釋.
那麼.
大家自己回去有機會再揣摩.
如果遇到一些.
你自己的一些困難.
或者朋友.
你不需要說自己的.
你說有個朋友這樣.
那應該怎麼處理呢.
聖經有教導的.
就是這樣大家去討論.
這些倫理有時都會困擾我們的.
其實我們認識了之後.
未必是我們自己.
而是我們幫助一些年輕人.

$^{1521}$進入婚姻的等候.
如果不是.
有時他們會有些扭曲.
有時他們會更加傾向一些.
世界的價值觀.
那就不是最好.
大家想不想看下面那段呢.
如果不是.
我們可以直接回到.
《馬拉基書》那段經文.
這裡就說未結婚.
怎麼等候.
寡居又怎麼處理.
那普羅都.
什麼話.
哦哦哦.
第七章二十五節.
其實如果將來我們看.
哥林多前書呢.
就更加我們整個.
整卷書去處理那些衝突.
因為哥林多教會.
正正就是他們的生命.
最大的問題就是.
沒有真正的分別為性.
歸給神.
就衍生了他們.
處在一個大都會的教會的環境.
二十五節開始.
論到同心.
就是等候婚姻.
都是臨前第七章.
其實整章都是處理那些倫理的關係.
論到同心的人.
我就沒有主的名字.
同心的人我就沒有主的命令.
但我既然是蒙主的連玉.
能夠做忠心的人.
普羅這樣看他自己.
就把自己的意見.

$^{1561}$告訴你們.
因為現今的艱難.
據我看來.
人不如守素安詳才好.
你有妻子前著.
就是上文所說的.
神安排了的每個人的份.
就是說你已經有妻子了.
就不要過分的.
就是說.
似乎有一些價值觀.
覺得有婚姻的約束.
就不能夠很專心事奉神.
好像是一個前類.
似乎有一些這樣的念頭.
在一些弟兄姊妹的內心.
所以保羅就說得這麼重.
你有妻子前著.
就不要求脫離了.
你沒有妻子前著.
就不要求妻子了.
所以保羅就說得這麼重.
你有妻子前著.
就不要求妻子了.
你如果娶妻.
就不是犯罪的.
處女若出嫁都不是犯罪的.
然而.
這等人肉身必受苦難.
我卻願意.
你們免受這些苦難.
所謂的肉身的苦難.
就是說.
始終你絕婚姻一定會有家庭.
有雙方面的性格的磨合.
人生有很多東西.
要拿來給別人.
要有一種.
大家生活的藝術.
擠牙膏有些人就由頭擠.

$^{1601}$有些人就由尾擠.
現在好些了.
現在是進化了.
那些牙膏是膠的.
就不會擠爛的.
以前是銅片還是鐵.
石子是石的.
你在前面擠後面就擠不出來.
所以妻子這麼激氣.
原因就是這樣.
你擠起來的.
就可以用得更多.
但是現在膠就沒有這個問題了.
你頭和尾擠都沒有問題了.
這些只是其中一點的生活藝術.
所以這裡告訴我們.
阿保羅不是誇張的.
有這些肉身的苦難.
就是說.
真的會有一些影響.
不夠專心的.
但是阿保羅的意思就是.
重點大前提.
是要看神安排你的人生的份.
這個要尊重.
這個才是最重要.
要尊重這個身份.
第29節.
弟兄們.
我對你們說.
事後減少了.
從此以後哪有妻子.
要像沒有妻子.
就是左手不要讓右手知道.
這意思.
主耶穌教我們奉獻.
左手不要讓右手知道.
怎麼可能不知道呢.
左手走去擠進奉獻箱.
右手擠在後面.

$^{1641}$他的衣服後面.
當然不是這個意思.
因為我們是整個人.
就是說那樣東西.
根本就是越低調就越好.
這就告訴我們.
好像沒有的意思其實就是有.
其實就是你很愛.
那個關係非常好.
以至你專心做那件事.
他也是跟你完全一樣認同.
是這個意思.
從此以後這裡.
第30節.
哀哭的要像不哀哭.
快樂的要像不快樂.
自買的要像無有所得.
用細物的要像不用細物.
因為這世界的樣子將要過去.
重點是我們終於明白.
人生原來短暫.
保羅的意思.
終於明白.
有些事其實真的不需要那麼介懷.
就是.
如果每個人有不同的處境.
比如什麼人.
原來真的不懂得這些智慧呢.
這些是智慧來的.
原來他那餐飯.
沒有牛扒.
是他當沒吃過東西的.
那這個人就變成.
為牛扒而生.
保羅的意思就是.
有牛扒開心.
沒有所謂.
是沒有所謂.
那些金訪大笑.
那些銀紙為什麼沒有燙過呢.

$^{1681}$沒有燙過呢.
誇張的那些廣告.
如果人生要你到這樣的水平.
才是人生的.
如果你不是這樣就不行了.
就完全發瘋了.
那個就完全沒有智慧.
不是說他不可以.
這樣的生活條件.
但人生怎麼可能是這樣呢.
保羅的意思就是.
你要好像.
有一個很大前提.
我們的出現.
已經是進入步向.
那個.
離世的死亡.
在這個階段.
你居然這麼看重.
短暫的那幾十年.
而忘記了永恆的.
那個參與那個服侍.
保羅就說這個就沒有智慧了.
所以他的意思.
就正正就是說.
整個人已經是完全.
受著那個真理的驅使.
去做一些.
合乎真理的事情.
對永恆有益的投資.
就是這個意思.
所以偶然間看下賀世波.
是沒有問題的.
就不是他們.
哇 每個禮拜我崇拜都不回去了.
我去看足球了.
看香港球 看什麼的了.
哇 這樣就.
我以前的弟兄就是這樣.
星期六晚.

$^{1721}$一定看通宵.
看英超.
後來開始回教會.
開始熱心.
接著發現原來.
原來不看還更加好.
專心第二天就教會.
回教會事奉.
他的靈明自己都有很大的改變.
就是原來有些事情.
做是沒有問題.
但你不停地做.
被那些事情捆綁.
你就忘記了原來永恆才是更加重要的.
保羅這裡的意思就是這樣.
就是說.
世界的樣子是要過去.
包括我們自己都會過去.
所以不要執著於.
短暫的得與失.
所以我們偶然間.
上山武是很開心的.
去看人家買東西.
是真的很開心的.
如果每個禮拜聚會又不回來.
那就是一個問題.
其實那個意思是這樣.
第32節.
我願你們無所掛累.
無娶妻的.
是為主的事掛累.
想怎樣叫主喜悅.
娶了妻的.
是為世上的事掛累.
想怎樣叫妻子喜悅.
他不是搞對立.
因為要看明相下文.
婦人和處女也有分別.
沒有出嫁.
出嫁是為主的事掛累.

$^{1761}$要身體,靈魂都聖潔.
已經出嫁的是為世界上的事掛累.
想怎樣叫丈夫喜悅.
我說者說是為你們的益處.
不是要牢籠你們.
不是他要強加在某個人身上.
所以我們每個人的份都是神安排的.
如果我們去到一個階段的靈性.
就是你滿足於你自己的份.
神給你的.
你在這個份裡面忠心,熱心去服侍.
洋日的一份平安,喜樂.
這樣就是你的成功成長.
譬如我經常都失眠.
晚上就醒了.
但是我很開心的.
我又不會說死了.
我又不會.
每個人都有自己的作息.
不同的事情要擔心,處理.
所以我想是整個人是不同的人.
這裡告訴我們.
乃是要叫你們行合宜的事.
得以殷勤服侍主.
沒有分心的事.
所以如果有人以為自己待他的女兒不合宜.
女兒也過了年歲.
事又當行.
他就可以隨意辦理.
不算有罪.
叫二人成親就是了.
倘若人心裡面堅定.
沒有不得已的事.
他自己可以做主的.
自己也很強的心.
要事奉神的.
並且由得自己作主.
心裡又決定了留下女兒不出嫁.
如此行也好.
本來你可以看見.

$^{1801}$他是多麼的通情達理.
一位的牧者.
他能夠看事物.
他看到每個人都有分別的.
每個人的性格不同的.
家庭環境又不同的.
所以不是一概而論的.
這些甚至給了一個原則之後.
都是有不同的處境去套用原則的.
所以這就告訴我們.
第38節.
這樣看來.
叫自己的女兒出嫁是好.
不叫她出嫁更是好.
丈夫活著的時候.
妻子是被約束.
丈夫如果死了.
妻子就可以自由隨意再嫁了.
只是要嫁.
這在主裡面的人.
然而按我的意見.
如果常守節更有福氣.
守節就是守貞節.
貞節牌坊那些.
就是不再嫁了.
這樣.
我又想自己是被神的靈感動.
就是按他的意思.
這就是這樣的情況.
就變成說了.
無論是獨身.
未婚.
和寡居.
應該怎麼處理.
反而最寶貴的就是.
在第七章說到關於婚姻.
保羅更加.
他的教導.
超脫了婚姻.
這麼有限的領域.

$^{1841}$就是每個人.
一定有你日常的身份.
無論你工作的身份.
家庭的身份.
你是男你是女的身份.
這個與生俱來.
給了我們.
無法改變.
但是你怎麼在這個身份裡面.
去經營你的人生.
在信仰裡面.
這個很精彩.
所以每個人都不同.
有些人追求得特別好.
就是滿足於自己的身份.
就在神裡面.
安於神給的那份.
去追求.
他就可以不斷地成長.
就像一般人這麼說.
給你一手好的牌.
你都打得很亂.
那就不行了.
給你很多資源.
你的人生很美滿的家庭.
很好的環境.
你讀很多的書.
但你根本不是用來榮耀神.
不是用來為神去事奉.
你的人生.
其實那副好牌.
被你打得壞了.
但是有些人出身很寒微.
成長我都說.
我是整個中學站在主教室門口的人.
我還好在家裡.
管教得很嚴.
我還沒至於入黑社會.
我的弟弟.
最小的弟弟.

$^{1881}$整個背後都是民生.
你可以想象得到.
家裡是沒有人來的.
我當然沒有民生.
我本來想的.
但是家裡禁止.
不想的.
我們那個時代民生一定是黑社會的.
所以.
每個人的成長都是不同的.
但是很奇妙地.
又走到一條全道的路.
到今天.
你就可以想象到神一定有祂的心意.
所以這一段聖經是最寶貴的.
就是七章的十七至廿四節.
保羅討論的就是.
任何的人.
你在神的眼中.
都是有用的.
都是有祂的安排.
和祂的旨意.
這一段是很寶貴的.
所以.
如果我們能夠在教會里.
認識自己在神眼中.
這麼寶貴的身份.
你就一定有在地上.
其他六天的影響力.
是很大的影響力.
譬如我來了一個多月.
在洪恩堂.
現在柏宇花園的職工.
有時看到我.
擺下的車.
抽著袋子.
他在一個奸廳里.
走出來開門給我.
看到我抽著東西.
兩只手都開不了門.

$^{1921}$你想想.
一個女士.
由奸廳里坐著.
離開奸廳.
走去大閘.
開門給我.
她當我是老闆.
其實不是.
這是一種彼此的尊重.
建立了彼此的關係.
你會看到.
原來我們對自己的身份.
有這份榮耀感.
滿足感.
幸福感.
是會影響到其他未信的人.
如果你自己都不覺得.
自己的身份寶貴.
回到教會.
你都會覺得自己像老諷一樣.
你又覺得自己是英勳.
你看樣東西都不順眼.
然後.
將六天.
你自己的眼界.
你自己的價值觀.
放在星期日的身份.
平日你是很尊貴的.
每個人都要巴結你.
你回到教會.
我都是很尊貴的.
每個人都要巴結我.
這就出事了.
你混亂了自己在神的眼中的身份.
其實這樣的人.
在神的眼中根本一文不值.
那個六天的尊貴.
是給你去事奉.
不是給你去炫耀.
你那個六天.

$^{1961}$是你去見證神的.
所以第七天.
每個弟兄姊妹都是神的兒女.
都是.
大家是完全一樣的.
就是公明天牧師所說的.
打回原形.
我們每個人回到教會.
我們沒錯.
你繼續可以有六天的身份.
但是大家是神兒女.
這是很重要的學習.
所以如果我們不學習到這件事.
你一定不可以在神的裡面有成長.
因為神會阻擋驕傲的人.
一切就是源於你內心的驕傲.
你就放不下你六天的身份.
如果你明白了.
通了神照國人.
你是什麼身份蒙照.
你是什麼身份.
保祿說你安於你的身份.
那你就會有六天的身份.
安於你的身份.
這樣你終於明白之後.
你就能夠.
那個運籌為厄.
在你自己的生命裡面.
我出身這麼寒微.
我家裡養豬.
小時候就要掃豬屎.
等等.
什麼都要做的時候.
你看那些基層.
我現在長大了.
看那些基層對人家很恭敬.
很尊重.
我未必需要再做那些事.
你想想我十八歲.
七歲十八歲.

$^{2001}$那個階段我要在彌敦推掃水.
拿掃水.
現在很流行叫廚余.
以前叫掃水.
你想想一個年輕人.
打大赤魯.
推著木頭車在彌敦道.
彌敦酒店後面的茂林街.
跟著你想想.
很多人看著.
你想想.
那個社會形象.
那個心態.
多麼扭曲,多麼粉碎疾足.
但是神慢慢會改變.
原來這些成長.
是給你將來對人的.
認識.
對人生的認識.
其實這是很寶貴的.
是否千金難買少年窮.
原來是的.
真的.
我都很滿意.
自己成長.
在這樣的出生.
雖然那個時候.
我們的成長都是很窮.
但是你就可以看到.
神居然在這樣的環境.
去呼召我.
使我能夠信主.
使我能夠成為祂的牧者.
這是很寶貴的一件事.
所以保羅講婚姻.
可以將它延伸到任何一個人.
在你的不同的處境.
你都要滿足於神給你的生命.
《馬拉基書》就在這裡.
其實講了這麼多.

$^{2041}$那個處境.
那個位份.
其實就是他們開始.
不再滿足於自己被老歸回.
在神面前的平淡.
不再滿足於自己.
千辛萬苦.
很多的前人.
所羅巴伯.
大祭司約書亞.
有哈該.
有撒加尼亞.
有尼欺米.
有以斯拉.
更加有以斯帖.
有慕底改.
這些敬虔的人.
他們所做過的.
帶來整個國民回歸.
耶路撒冷重建.
是為了將來的尼賽亞出現.
是為了尼賽亞出現.
周妻又和諧共處.
一定是大家都有共同的信仰在神里.
但一方面都不行.
因為始終是一個弱是雙方面.
都是神所恨惡的.
所以要緊守你們的心.
不可以行鬼詐.
所以鬼詐就是對神對人都不再忠誠.
是因為神的諾言違背了.
就是這個問題出現.
所以追本族人要探求的就是.
神的心意是這樣.
要回頭.
就是為什麼沒有那種氣氛.
沒有那種力量去維繫.
去建立,去經營.
這個就是神在這段經文提到.
這裡告訴我們.

$^{2081}$神似乎突然間加入了.
神只是做阿當和夏娃.
沒有做其他人.
所以剛才那個不是聖經難題.
那個人的太太.
就是他的妹妹.
真的是他的妹妹.
所以不是爆出來的.
當時世界只有四個人.
是可以通婚.
某程度上.
雖然我們後來叫亂倫.
埃及記就明文禁止了.
但是在那個時候.
因為神既然在.
他的創造裡面.
他認為這樣.
就是可以.
來到他凡顯.
所以頭那一陣一定是.
他的.
阿當夏娃的兒女.
就一定是這樣凡顯.
所以到第二代,第三代.
就已經是那些.
沒有那麼緊密的.
親戚關係.
所以一定是這樣發生來成長的.
所以為什麼說.
戈隱.
他的弟弟就死了.
為什麼他會怕.
有人會殺他呢.
那些人是哪裡來的呢.
又是聖經難題.
他的弟弟阿伯都死了.
難道阿爸阿媽追殺他嗎.
不會的嘛.
那戈隱怕什麼呢.
我要有一個記號.

$^{2121}$讓人看到我就不會殺我.
這是《創世記》第四章的經文.
那些人會殺他呢.
那就是阿當夏娃繼續生的.
繼續生的.
那不就有很多表兄弟.
然後就一直發展下去.
或者陳氏忠親會.
就一直發展下去.
就是這樣的.
這裡就是聖經說的.
有聖經根據的.
雖然神有零的餘力.
能夠造多人.
但是他不是單造一人.
為何只造一人呢.
乃是他願人不貼成的後裔.
大家都是來自一脈相承.
神創造的阿當夏娃.
大家都是在這個境前裡面.
去成長認識神的.
所以你會看到.
《創世記》第五章.
就有一個家譜.
阿當就在阿伯以外.
就生了一個家.
就是在阿伯的家.
阿伯以外就生了一個叫做失德.
就是另外代替了阿伯的.
就沒有再耕耘該人.
你留意到這個家譜.
然後失德就再生下去.
就是以羅氏.
你會留意到.
那個是境前的後裔.
就是這裡的意思.
所以就是說.
原來神體的婚姻文約.
都是很重視.
其實就是那個諾言.

$^{2161}$所以這一段經文.
其實由第十節.
到第十六節.
是三個諾言違背.
就令到社會充斥著.
詭詐,謊言.
和大家的強暴.
大家都不信守那個諾言.
所以我們信了主耶穌.
對於諾言.
我們都應該學習.
要去信守.
所以就不要隨便起誓.
新約因是這樣教的.
但就不代表我們不可以起誓.
意思就是說.
起誓要守.
神會追究的.
所以我們都是要學習這一方面.
我們繼續下去嗎?.
因為第十七節到三章五節.
是另外一個段落.
是講到到底公義的神在哪裡.
是那個問題來的.
到底神有沒有真的存在.
在一些公義的事情上.
應該說在一些不義的事情上.
神是怎樣來理解呢?.
第十七節我們先讀了.
你們用言語去繁瑣.
繁瑣就是.
我們想應該是對的.
日訛夜訛的.
英文是very.
很疲乏地去疲勞轟炸.
你們用言語去繁瑣.
耶和華.
你們還說我們在何事上繁瑣他呢?.
因為你們說.
凡是行惡的.

$^{2201}$耶和華眼看為善.
並且他喜悅他們.
即是扭曲了那一件事.
即是他們見怪不怪.
神都理不了.
他都沒有理這些事情.
講得多就變成.
好像金科玉律那樣.
就是說你看到社會.
那麼多的不公不義.
神都沒有立刻處理.
沒有立刻審判.
慢慢地就是這樣.
沒有理會.
看為善.
看為他喜悅.
最終就是.
先知就說.
公義的神在哪裡.
這件事.
公義的神在哪裡.
那件事因為我們自己有了.
一個所謂的公義的判斷.
而我們看不到神.
好像我們這樣要求的公義的判斷.
所以.
猶太人就這樣來.
繁瑣耶和華.
第三章.
一到第五節.
就立即去解答.
公義在哪裡.
或者神怎樣去彰顯.
他的公義.
不如我們下個星期再講.
我們就停在這裡.
我們一起禱告.
雖然我們用了很多的時間.
去解釋.
保羅在哥林多前書第七章.

$^{2241}$對於當時哥林多教會.
要面對一些.
倫理的處境.
婚姻的問題.
是很寬闊地讓我們有一個.
很新的念頭.
更寶貴的就是.
我們原來在神的眼中.
無論我們.
多麼的不堪不懼.
都是原來.
是主耶穌.
附上了重架.
將我們買熟.
我們成為神的兒女.
我們每個人都有我們成長的處境.
成長的家庭.
或者我們現在.
都要面對.
我們自己可能有一些.
根本自己都解決不了的一些困難.
無論是身體.
或者是一些可能人和人之間.
的倫理的處境.
或者是一些.
我們要面對的.
不同的社會事情.
在禰的眼中我們是重架.
我們是被買熟.
我們是有神為.
安排的人生的位份.
和人生的跑道.
多謝主讓我們.
能夠在.
不同的人生階段去認識禰.
盼望我們現在.
能夠在神的愛裡面.
我們追求.
尋找禰的.
心意以致.

$^{2281}$我們能夠用盡.
我們的精力.
我們的熱誠去為主.
為福音去效力.
去幫助.
多謝您聽我們感恩禱告.
奉耶穌基督的名.
阿們.
\newpage



\section{馬可福音 8:31-38}
\label{sec:dYh7F8XFvZo}
\textbf{2024-02-24 晚堂《作門徒的代價》馬可福音 8\_31-38 講員:褚永華牧師}
\newline
\newline
連結: \href{https://youtube.com/watch?v=dYh7F8XFvZo}{\texttt{https://youtube.com/watch?v=dYh7F8XFvZo}} ~~~~ 語音日期: 2024-02-27
\newline
\newline
\hyperref[sec:4SEW84zBGCU]{\small{< < < PREV SERMON < < <}}
~
\hyperref[sec:index_chronic]{\small{[返順時目]}}
~
\hyperref[sec:index_scriptual]{\small{[返順卷目]}}
~
\hyperref[sec:A8lXH1cPVBk]{\small{> > > NEXT SERMON > > >}}
\newline
\newline
馬可福音 8:31-38
\newline
\begin{longtable}{cl}
\hline
\hline
章節 & 經文 (和合本修訂版)\\
\hline
8:31 & \begin{tabularx}{0.7\textwidth}{X} 從此，他教導他們說：「人子必須受許多的苦，被長老、祭司長和文士棄絕，並且被殺，三天後復活。」 \end{tabularx} \\ \\ \relax
8:32 & \begin{tabularx}{0.7\textwidth}{X} 耶穌明白地說了這話，彼得就拉著他，責備他。 \end{tabularx} \\ \\ \relax
8:33 & \begin{tabularx}{0.7\textwidth}{X} 耶穌轉過來看著門徒，斥責彼得說：「撒但，退到我後邊去！因為你不體會神的心意，而是體會人的意思。」 \end{tabularx} \\ \\ \relax
8:34 & \begin{tabularx}{0.7\textwidth}{X} 於是他叫眾人和門徒來，對他們說：「若有人要跟從我，就當捨己，背起自己的十字架來跟從我。 \end{tabularx} \\ \\ \relax
8:35 & \begin{tabularx}{0.7\textwidth}{X} 因為凡要救自己生命的，必喪失生命；凡為我和福音喪失生命的，必救自己的生命。 \end{tabularx} \\ \\ \relax
8:36 & \begin{tabularx}{0.7\textwidth}{X} 人就是賺得全世界，賠上自己的生命，有甚麼益處呢？ \end{tabularx} \\ \\ \relax
8:37 & \begin{tabularx}{0.7\textwidth}{X} 人還能拿甚麼換生命呢？ \end{tabularx} \\ \\ \relax
8:38 & \begin{tabularx}{0.7\textwidth}{X} 凡在這淫亂罪惡的世代，把我和我的道當作可恥的，人子在他父的榮耀裡與聖天使一同來臨的時候，也要把那人當作可恥的。」 \end{tabularx} \\ \\
[1ex]
\hline
\hline
\end{longtable}
$^{1}$各位弟兄姊妹平安.
因為都是新年期間.
我想到明天.
各位講完都會祝大家在新的一年裡面.
有新的恩典.
每天都能夠經歷上自己與我們同在的恩典.
教會的年曆 我們很多時候.
就未必會按照教會的年曆來宣講訊息的.
教會的年曆的編排.
是將全本的聖經每三年或五年.
就會將全本聖經的重要經文讀一次.
即是如果你每個星期五十二日.
每個星期回來勤力教會的話.
三年到五年.
你就會聽到整本聖經.
重要的經文被重讀一次.
每一次重讀它有些經文.
它在摩西五經裡面有一段.
在先知書裡面有一段.
在智慧書,詩篇,傳道書,針耳瘤裡面又有一段.
然後去到新約裡面的.
保羅的書信又有一段.
然後福音書又有一段.
每個不同的段落.
有五六段的經文放在一起重讀.
它開始的點在哪裡呢.
是耶穌基督降生的時間開始.
所以在那幾個星期裡面.
它會讀很多關於主耶穌基督降生的事.
預言主耶穌基督降生的事.
然後主耶穌基督降生之後.
它有六到九個星期.
它們叫做主的顯現期間.
epiphany.
epiphany是主耶穌基督上帝顯現.
基督既然降生.
我們每年都似乎到聖誕節的時候.
就要將耶穌變成一個BB.
每年都來一波.
但在教會的年曆的經文裡面.

$^{41}$它就不容許我們這樣想.
耶穌基督降世的時候.
道成肉身住在我們中間.
世人怎樣認識祂呢.
所以主耶穌基督在世上生活的時候.
藉著不同的講道.
藉著不同的生活的表現.
藉著他和人間的往來.
藉著和詩篇門徒的往來.
將顯出這位道成肉身的上帝.
這位耶穌是怎樣的主宰.
過了6到9個星期之後.
就進入到一個期間叫做大齋期.
大齋期就叫做御苦期.
預備主耶穌基督降十字架的期間.
有40天那麼久.
今天是第三個星期.
御苦期的第三個星期.
在這個第三個星期.
我們今晚所讀的《馬禾福音》第八章.
是其中一段要讀的福音書.
它要表示什麼呢.
就是主耶穌基督臨在我們中間.
成為一個BB降世為人.
經歷人間的各種苦難.
但他同時愛上帝自己的兒子.
尊貴到不得了.
人間沒有一個字眼可以形容他多尊貴.
但同樣人間沒有一個字眼可以形容他多苦.
在又高又尊貴的上帝之子.
成為又卑賤和人共苦的一位仁子.
一個普普通通的人.
在這個期間的落差很大.
我們今天信主耶穌基督的話.
有人說我願意跟隨主耶穌基督.
是一位上帝之子.
無限的威榮 無限的權柄.
但這位主耶穌就說.
你想清楚.
我同時是一位受苦的耶穌.

$^{81}$你能不能夠像我這樣.
造成肉身由最尊貴 最有能力.
放下身段 放下一切榮耀.
甚至放下自己一切能力.
來到這個世界裡面.
成為人 接受人的各種挑戰.
接受人的各種苦難.
接受人的鄙視.
你能不能夠跟隨我去到這樣的地步.
今晚的經文主要是說一些這樣的意思.
教會的年曆裡面給我們有40天的時間.
在想這些經文.
在經文裡面 總是叫我們說.
你生活要簡單些.
所以很多教會在大齋期這段日子裡.
不是那麼鼓勵人去旅行.
不是那麼鼓勵人去吃自助大餐.
你可以吃 你可以吃.
昨天我去吃了自助大餐.
它不鼓勵你這樣做.
但主要就是說.
你要想 主耶穌基督怎樣為我們受苦.
我們與祂一起受苦的期間裡面.
節約地生活 簡單地生活.
再次操練自己的心靈.
讓自己同樣可以過一個簡單的生活.
生活一個無罪的生活.
來到這裡我們說做基督的門徒要付出代價.
我相信這些大家都明白的.
我們雖然頭腦裡明白.
但在生活裡面我們很難去接納擁抱這樣的道理.
信了耶穌 跟隨主耶穌基督.
做主耶穌基督的門徒會受苦嗎?.
會受到不公平的對待嗎?.
主耶穌基督的生命告訴我們.
是一位義者無罪的人.
但將各種的罪惡 編出來的罪惡.
加在他身上.
甚至他釘在十字架上的時候.
下面的人還在笑他 還在譏笑他.

$^{121}$你能夠醫治別人 你能不能夠醫治自己?.
你能夠救別人 你能不能夠救自己?.
你為什麼救別人救得這麼淒涼.
要死在十字架上呢?.
你竟然救人能夠脫離苦海?.
你能不能夠救自己脫離十字架的苦難呢?.
你下來吧 你下來吧.
但我們這位主.
他就讓我們看到他是一位怎樣忍耐的主.
在十字架上 他沒有說出惡言.
沒有說出兇惡的言語.
反而他在禱告 求上帝赦免他們.
他們所做的 他們不知道.
我們能不能夠在這四十天裡面想一想.
檢視一下我們自己的生活.
在我們的生活思維動機上面.
清理一下.
好像我們年廿八那樣掃邋遢 洗邋遢那些.
將自己的心靈在上帝的話語的對照之下.
讓我們可以明白.
我們在這段日子裡面.
能夠成為上帝的見證.
在這個世界裡面.
讓上帝可以公然的對這個整個的世界.
對這個罪惡的世界 不公義的世界.
主耶穌基督 道成肉身.
在十字架上面 上帝坐在天上的寶座上面.
向著所有的世人說 看一看我的教會.
我們能不能夠被上帝這樣尊崇.
在世人面前 被上帝說 看一看我的教會.
在我們今天所讀的經文 第三十一節那一句.
在上文那一句 祂就說.
耶穌基督在上文那一句.
就是去到哥澤尼亞 菲臘比的地方.
問門徒 兩個很撩心撩肺.
實際上我們都不是那麼撩心撩肺的.
人說我是誰 其他的人沒有跟隨我的.
旁邊大衛的那些人.
問門徒 人們說我是誰.
問門徒很容易回答 我聽到有人這樣說.

$^{161}$所以經文那一句就說.
人子 他們說我是誰呢.
有人說 是施洗約翰.
又有人說是以利亞.
又有人說是先子裡面的一位.
但是耶穌不知可否.
就跟著問門徒 你們說我是誰.
我相信在我們港台上 教會裡面.
已經講這篇道理.
講哥澤尼亞 菲臘比這篇道理.
講了無數次那麼多.
但是實際上 它裡面的意思是甚麼呢.
它叫我們真的對主耶穌基督有一個更深的認識.
不是單單投入知識裡面的認識.
不是單單停留在我們洗禮之前.
我們受洗禮班的訓練.
對主耶穌基督有很淺顯的認識.
我們每天在生命之中慢慢成長的時候.
要將主耶穌基督.
我們在聖經裡面得到的那種知識.
能夠在經歷上去認識它.
我不知道各位大眾節目有沒有這樣的感覺.
就是我們做見證的時候.
很多時候講不到我們對聖經的經歷.
我信上帝創造天地宇宙萬物的主.
我們怎樣在我們的生活裡面.
經歷這位全能的 全智慧的.
無所不知 無所不在的上帝呢.
我們怎樣將這樣的信念.
在我們的生活裡面有所經歷呢.
很顯顯的.
如果這位創造天地宇宙萬物的主.
什麼都知道.
馬達夫音裡面說.
你不用為世上各樣的事來擔心的.
要先求他的 告他他的義.
這句話.
他是創造天地宇宙萬物的主.
什麼他都知道.
那為什麼還要為我自己生活裡面.

$^{201}$一點點的 細小到無論的東西.
來為自己擔心呢.
我們能否經歷到這位.
創造天地宇宙萬物的主.
在我們生命裡面最微小的事情上面.
經歷到他的同在.
他的安慰.
與他的同行呢.
這些是可以的.
只要你在這四十天裡面.
靜心下來.
將我們所認識的.
認知的聖經的道理.
和我們生活裡面對上帝.
對這些知識的一種經歷.
你將它慢慢想 慢慢禱告.
如果可能的話.
慢慢用筆記下它.
在什麼事情上面.
我經歷到上帝自己的同在.
上帝釋除我一切的疑慮.
在我最擔憂的地方.
上帝給解答.
上帝給答案.
未必我能夠完全明白.
全部的事實.
但是在這一刻.
上帝讓我知道這麼多.
我生命得著釋放.
我們可以這樣經歷上帝.
我們生命裡面的奇妙的工作.
所以在這裡他就說.
這位這樣的耶穌.
你們恕我詞綴.
我們每一個人都要這樣回答.
用我們自己生命裡面的經歷.
來回答說.
我知道你是誰.
I know.
彼得就說.

$^{241}$你是基督.
是永生上帝的兒子.
跟著我們今天所讀的經文.
一開首.
第31節那一句.
他就說.
從此.
我們讀經的時候.
這些字眼.
所以我們要留心.
從此這兩個字代表什麼東西呢.
在從此之前.
這段話.
或者這段裡面的經歷.
或者文字裡面所表達的那種思維.
在31節之前.
是未講過的.
是不是.
現在已經七點半了.
最適合吃飯睡覺的時間.
從此他表示.
在他之前沒有講過這些東西.
31節打開的東西.
對門徒來說.
是很新鮮的東西.
如果我是當日在門徒之間.
當耶穌基督問.
人恕我詞綴.
沒有什麼殺傷力.
你們恕我詞綴.
其他人沒有回答.
只能說.
我知道你是基督.
永生上帝的兒子.
耶穌不知可否.
跟著耶穌就說.
從此他就教訓他們說.
如果你在場在那裡面.
你聽到耶穌講的這段話.
你要去死嗎.

$^{281}$我們跟你的命運是要這樣嗎.
我要寫記跟隨你嗎.
背在十字架上跟隨你嗎.
當初在海皮那裡.
來跟從我.
你只是叫我什麼.
得人如得魚.
有講這些嗎.
你講清楚一點.
當日在加利利的海邊.
我們都在那裡打魚.
生活日子非常舒適.
雖然有時打到魚.
有時打不到魚.
但不是每天都打不到.
有幾天魚獲非常好.
我們的生活無憂.
你突然殺人加利利的海邊.
突然殺人我們的生命.
對著我們說.
來跟從我.
我要叫你得人如得魚.
我們就跟隨你.
你當時沒有說.
從此之後我要去死.
你們又跟著我去死.
你們的生命有完全的改變.
沒有 完全沒有.
如果你在場的話.
你對耶穌基督這樣的講法.
你會做出什麼反應.
拜拜.
貨不對辦.
你都沒有講完.
如果你講完的話.
我們有得考慮.
但今日的聽眾.
今日道道經文的人.
你會怎樣想.
你現在知道了.

$^{321}$Too late.
可不可以跟牧師說.
牧師 我雖然洗禮.
不過現在信耶穌這麼艱難.
跟耶穌這麼艱難.
你有幾個出路.
你跟牧師說.
要不我就退縮.
把洗禮證還給他們.
拜拜了 我搞不定.
原來要付這麼大的代價.
這樣又不行 那樣又不行.
要為上帝 為信仰 為福音.
死去自己的生命.
這個要求過大一點.
太大了.
要為你的緣故.
放下所有的東西.
代價太大了.
可不可以退縮呢.
可不可以跟耶穌說拜拜呢.
可不可以在生命之中.
我可以用另外一個拜拜的方式.
不用這麼弱.
什麼都發力 又不發力.
用陰力來跟耶穌.
你不要這樣跟.
硬橋硬馬跟 死定了.
要用陰力.
做一點不做一點.
做表面不做實際.
這樣來跟都可以吧.
然後去到天堂的時候.
懂得念那個口訣.
最重要那個口訣.
我信耶穌 祂就讓你進去.
你有幾條路可以選擇.
你會怎麼選呢.
從此耶穌就教訓他們說.
彼得 你說我是基督.

$^{361}$當初耶穌沒有說他.
彼得 你都不知道你在說什麼.
你能不能夠付出代價.
耶穌一聲都沒有出.
到耶穌基督說完這段話.
從此.
人子必須受許多的苦.
被長老祭祀的民事氣絕.
並且被殺.
過三天後復活.
如果彼得聽到這些說話的時候.
聽到復活的話.
他應該 哦 還有點盼望.
我們不明白他為什麼聽一段.
只聽到前面那段.
你要死 不行.
我們跟你那時候不是這樣的.
他就即時和耶穌對撼.
他就責備耶穌.
責備耶穌.
你這樣說不合理.
我們畢竟所有的東西都跟隨著你.
你怎麼可以棄我們而去呢.
說跟隨你一定要寫起自己的生命.
耶穌就責備彼得.
退我後面去吧.
撒旦 退我後面去吧.
你只是體貼人的意思.
不體貼上你的意思.
兩使徒在對撼.
兩個觀念的對撼.
耶穌是很清楚說明.
你說我是基督.
但這個基督是要受苦的.
你行不行呢.
這個基督要為天國的緣故.
寫起一切的所有.
完成上帝的心意.
他有他自己的意思.
但他說按照你的意思而行.

$^{401}$我們能不能夠在生活裡面學習.
用生命的經歷.
和上帝的對照.
上帝說我願意跟隨你.
我願意寫氣.
我願意寫記.
我願意背起十字架來跟隨你.
寫記背起十字架跟隨主耶穌.
是需要我們每天在生活裡面.
慢慢去學習 慢慢去操練的.
實際上聖經裡面有很多的指引給我們.
告訴我們 教導我們.
要我們學習怎樣去寫記.
可以為信仰付出一些代價.
甚至極致付出生命的代價.
我想在我們現實的香港.
我們很難想到 很難經歷到.
要為信仰付出生命的代價.
但經文裡面從來都教導我們.
我們有一個心智.
或者有一個盼望.
有一個信仰裡面的一個很基本的要求.
雖然不是現在.
但始終有一天.
你可能經歷到那些苦難的時候.
經歷到那些困難的時候.
信仰受到極端的挑戰.
在極嚴寒的冬天裡面.
嚴寒之中 信仰受到挑戰.
生命裡面極大的困難壓迫著你.
作出你越來越大的挑戰.
你那一刻 臨到我說.
主啊 我仍然跟隨你.
我們聽到很多的故事.
都在我們的生命裡面做出一些的抉擇.
或者受到一些的苦難.
但那些苦難未必是很重的.
我們人都有一個傾向.
總是覺得我們自己受的苦是最苦的.
你那些不夠我苦.

$^{441}$我那些比你苦.
我們經常要鬥病重.
我的病比你厲害.
我中招而已.
我的中招比你厲害.
我們總是有一個這樣的心.
所以我說 我很苦.
我的苦難很苦.
但實際上我們的生命.
就好像啟示錄裡面所說的.
我們抵擋罪惡.
抵擋魔鬼撒旦.
他各樣的攻擊.
窒疫的攻擊.
貧困的攻擊.
失業 失戀等等很多的攻擊.
我們都未曾到流血的地步.
他說 你還不夠啊.
慢慢來吧.
慢慢粗俗一點吧.
當主耶穌基督和彼得這樣對話的時候.
耶穌就說.
我的一生為此而來.
我知道我自己來到世界上的目的是什麼東西.
要將要完成上帝給我的使命.
好像馬太福音那裡說.
要將自己的百姓從罪惡裡面救出來.
這是上帝給過我的責任.
給過我的一個召命.
我來到這個世界裡.
無論什麼苦都好.
我必須要完成上帝交代給我的使命.
上十字架為世人贖罪.
彼得 你呢.
你怎樣啊.
耶穌說.
現在我多給你一個使命.
你能不能為我受苦啊.
你能不能為基督的緣故.
為福音的緣故.

$^{481}$來寫幾本十字架跟隨我啊.
如果我們讀福音書一直讀下去的時候.
我們就發覺到.
彼得在生命裡有很多的陰暗面.
去到大祭司的花園裡.
三次不認主耶穌.
我不是彼得.
但是如果在那三次的挑戰裡.
使律對他說.
我知道你跟隨過耶穌.
你的口音就給我們露出來.
彼得三次說.
不是 NO.
發誓說 NO NO.
當他在說 NO 的時候.
口聲大得不得了.
NO 我不是跟隨他的.
我跟他毫無關係的.
我跟他必經關係的.
如果在他內心的深處.
他想起主耶穌基督對他說的話.
對他的挑戰.
若有人要跟從我.
就當捨己.
彼得啊.
你為你自己的安危.
你沒有背負十字架.
你為你自己生命的安全.
你沒有為我捨己啊.
為福音的緣故.
捨棄生命的.
你沒有啊.
為上帝自己更大的使命.
甚至將自己的性命放下.
你沒有啊.
當我們這樣讀經文的時候.
前後呼應.
我們說上帝啊.
如果昔日的彼得.
親眼見到你.

$^{521}$跟隨你.
生命裡受過那麼多訓練的彼得.
去到生命某一點的時候.
仍然有這樣的軟弱的話.
上帝啊.
求你幫助我.
我不敢自誇.
我自願上帝給我們足夠的力量.
在生命中.
在細微的地方.
慢慢操練我們自己.
讓我們能夠承擔苦難.
我們能承擔多少.
上帝啊.
求你在我們能夠承擔的苦難.
與挑戰裡.
慢慢磨練我們.
讓我們慢慢學習怎樣捨己.
讓我們學習怎樣背負十字架.
可以放下一點.
未必放下所有.
先放下一點.
慢慢進步.
慢慢挑戰自己.
在信仰上慢慢成熟.
如果我們能夠這樣跟隨主耶穌.
考慮生命裡付的各樣代價.
我們的生命是能夠成長的.
或者說我們的生命是能夠慢慢成熟的.
是被上帝自己所使用的.
如果我們能夠這樣成長的話.
我們就可以站在世人的面前.
成就上帝極大的挑戰.
說:這是我的教會.
實際上這句話.
可以在生命裡不同的層面來問.
上帝可不可以將我們.
擺在世人的面前.
擺在教會眾人的面前.
朱英華.

$^{561}$你是我的愛子.
是很沉重的承擔.
對所有的人說.
你是我的愛子.
我們要求主給我們力量.
能夠在這些環境裡堅持起來.
上帝不單止這樣對我挑戰.
趙小芬.
不應該叫她的名字.
我們每一個人.
上帝可以在眾人的面前說.
你是我的愛的女兒.
我們能不能夠在我們的生活裡.
在極少的地方慢慢學習.
上帝可以說.
趙小芬.
是我的女兒.
你看她多麼捨己.
你看她多麼肯背十字架來跟隨我.
牧師你看看她.
莫弟兄你看看她.
每個人都可以在上帝的跟前.
被上帝對整個不公義的世界.
黑暗的世界.
罪惡的世界.
在眾人的面前說.
你看看.
我的教會.
敢算價.
你看看.
我的僕人.
我的侍女.
我們可以在上帝的跟前.
承受上帝對我們的挑戰.
對我們向世界的推介.
完全源起於.
從此.
主耶穌基督說.
我要去受苦了.
我們能不能夠計算代價.

$^{601}$與祂一同受苦.
由細數學者.
我們真的要學習.
我可以舉一些例子.
幫助我們如何在生命中慢慢慢慢地成長.
未來的日子.
或者在我們生命中的挑戰.
是會越來越多的.
困難是會慢慢地越來越多的.
成長受到的挑戰是會越來越沉重的.
或者我們在教會裡久了之後.
我們對一些思想上的.
經歷裡的.
信仰的思考上.
理性的挑戰.
我們慢慢地集於圍牆.
我們不是很細心地去想.
但實際上我們每天的生活中.
都遇到這樣的困難.
而要求我們.
挑戰我們.
你會怎樣.
但很多時候我們就不假思索地.
將很多的困難掃到一邊就算了.
在第一世紀到第三世紀.
主耶穌基督.
大體上在讀新約的時候.
我們都知道了.
主耶穌基督三十歲左近出來.
呼召門徒.
傳揚天國的福音.
三年左近.
三十三歲.
做了三年功夫.
然後釘十字架.
那個是第一世紀的三十年代.
我們說從第一世紀到第二,第三世紀.
實際上第二和第三世紀.
是信仰極被挑戰的年代.
很多的書.

$^{641}$教會歷史的書.
或者教會神學.
基督教神學.
教義神學等等.
特別針對第二和第三世紀.
為什麼呢.
因為第一世紀.
第一世紀.
是那些人剛剛初信.
初信的時候.
慢慢信仰成長.
遇到很多挑戰.
他們不太懂得怎樣回答.
不懂得怎樣回應.
但是生命之中.
我們見到保羅.
保羅就說耶穌基督.
特別是在提瑟羅尼加前後書那裡說.
很多人挑戰他.
你說耶穌在第一世紀.
那些門徒都敢傳.
我信我主耶穌基督.
上帝獨生之子.
為什麼一位這樣全能上帝的兒子.
會死得這麼淒涼.
死得這麼不公義.
死得這麼罪惡呢.
釘在十字架上被咒了.
按照舊約來說.
被掛在木頭上.
都是被咒了.
掛在十字架上.
是被咒了的人才會這樣.
你說叫我信一個被咒了的人.
被上帝完全這樣棄絕的人.
相信他是上帝的兒子.
你回答我吧.
很多希臘人.
熟悉希臘哲學那些.
就會挑戰.

$^{681}$你怎樣叫一個神智清醒的人.
信一個這樣的人做上帝的兒子.
是媚菜 是救贖世人的.
信不下去的.
信的人是傻的.
或者用我們現在的語言來說.
因為神樂園是九龍塘.
遇到一些新來香港的朋友.
在我們來香港探望我們的時候.
我還要在金巴倫道上.
駕車走過的時候.
指幾間出名的大屋給他聽.
我們神樂園金巴倫.
不是金馬倫道.
尖沙咀是金馬倫道.
九龍塘是金巴倫道.
我希望寄奉獻的.
不要寄到金馬倫道.
我們那裡街上有幾間.
是出名的屋.
第一.
就是李小龍的屋.
就是我們街中間.
所以走過的時候告訴他.
這間就是李小龍的屋.
隔離沒多久就是周潤發.
周潤發住在這裡.
但是有一個不得不告訴他.
在右邊有間屋.
快要到街尾.
在什麼居士隔離那裡的屋.
是一個很出名的人物.
很夠膽.
直接走去香港的壽福屋.
問他拿十億的那個人.
他拿了十億裡面.
不知道用了多少億.
不知道用了多少千萬.
就買了九龍塘那間屋.
過了給他老婆.

$^{721}$所以警察抓人的時候.
就進那間屋抓人.
什麼陳什麼張什麼強.
不是小強.
一講那些人物.
我們就知道那些是名人.
出名的人物.
但是在我們的信仰裡面.
我們就說.
我們能不能夠跟隨主耶穌基督.
在我們的生命裡面.
我們未必好像那些人.
那麼出名.
特別是那位去到壽福屋.
他問他拿十億的那個人.
上帝啊.
在我們的生命裡面.
遇到一個這樣的人.
為什麼你不對付他呢.
終於他就被抓了.
在內地被槍斃了.
很多人都說.
這樣才順氣.
受到報應.
在我們自己的生命裡面.
我們會眼見到很多不公平的事.
不公義的事.
耶穌基督是一個義人.
如果罪人受到制裁.
很應該.
但是一個完全沒有罪的人.
被人冤枉他.
說他有罪的話.
很難受的.
公理何在呢.
公義在哪裡呢.
上帝啊.
如果你是公義的.
怎會有這樣的事情發生呢.
很多問題鋪天蓋地的來.

$^{761}$我們怎樣去面對它呢.
我們窮畢生之力.
有很多問題我們想都想不通.
想通這一點.
其他一點就跟著來.
不行啊.
不是這樣啊.
我們怎樣能夠面對生命裡面.
諸多沒有理由.
不知道怎樣解釋的.
我們能不能夠真的放下.
我真的不知道啊.
但我仍然能夠相信.
上帝知.
我不知道上帝何時執行他的公義.
肯定在我現在短暫的生命裡面.
未必能夠按照我的意思.
來執行我的公義.
但我絕對相信.
還有一段頗長的時間.
甚至做永恆的長遠的.
做永恆的短暫裡面.
長遠裡面.
上帝必定會執行上帝的公義.
我們生命太短了.
我們接納不了.
最好按照我們的有生之年.
這麼短暫裡面.
讓我們見到我們認為的公義.
從此上帝主耶穌就教訓他們說.
我要受苦.
你們跟隨我的人.
彼得啊.
你都要學習.
怎樣寫記.
背起自己的十字架來跟隨我.
若有人要跟隨我.
你就這樣吧.
做一個選擇.
做一個抉擇.

$^{801}$因此有很多人就說.
做門徒可不可以選的.
我只信耶穌不做門徒可不可以呢.
鋪天蓋地的書告訴我的是不行的.
你一天做一天不做.
你可不可以說.
我信耶穌不跟隨耶穌可不可以呢.
是可以的.
我只做基督徒不做門徒可不可以呢.
是不可以的.
你說基督徒的基督徒那個字.
是什麼意思.
那個字的意思是小基督.
小型基督的意思.
你做了基督徒.
是做一個小型的基督.
你可以不做門徒嗎.
No 當然不行.
你既然身份是一個小型的基督.
跟在你基督後面.
做一個門徒.
做一個基督徒.
你就盡上帝給你的恩典與力量.
慢慢學習.
慢慢經歷.
跟著主耶穌基督.
一直這樣走.
我願意在我們的生命之中.
能夠經歷上帝給我們的力量.
我們沒有辦法.
好像古代裡面的偉人.
那些偉人.
浮士德.
可能如果你有空.
在這四十天的大齋期裡面.
找些東西讀一讀.
你讀一下浮士德.
浮士德是一個很有意思的故事.
浮士德是一個德國的.
一個很出名的作家.

$^{841}$叫做戈德.
是一個很出名的哲學家.
出名的詩人.
用六十年的時間.
花了六十年的時間.
構思寫出這本小型的.
一本著作叫做浮士德.
這本巨著.
浮士德在生命之中.
他要找生命的意義.
好像我們今天要找生命裡面.
跟隨主耶穌基督.
做主耶穌基督門徒的意義.
他同樣的追究.
他在生命裡面追究到.
生命到底為什麼.
因為他總是覺得生命裡面.
得不到滿足.
他是一個很聰明的人.
但是每一樣事情做的時候.
都不能超越那個已知的界限.
做事情做到某一個階段.
他再也超越不了過去.
他在生命裡面.
那種不能夠超越.
他有很多的不滿足.
最好就讓我全能地.
我的生命裡面.
所遇到的各種困難.
我可以一躍而過.
即時可以解決它.
我願意有愛情.
我喜歡那個.
就可以得到那個.
手到拿來.
一點也不費功夫.
我願意這樣的.
我畫的畫.
過的難關.
只要我有心去做的話.

$^{881}$我一定超越到那些界線.
但是可惜生命裡面.
是超越不了過去.
最終於.
哥德就.
做了一個.
這樣的交易.
佛士德最終和魔鬼做了一個交易.
魔鬼就說.
行 你想每樣東西都得心應手.
每樣東西都可以超越.
行 我和你做個交易.
在24年裡面.
我給你24年的時間.
你心想事成.
我們工人不是經常心想事成.
心想事成是很可怕的東西.
我恨那個人心想事成.
多大劑.
我恨他.
心想事成多大劑.
我有一個你不愛的人.
心想事成多大劑.
佛士德就說.
我願意用24年來交換心想事成.
我想做什麼就即時能夠做得到.
24年.
但是24年的結果是什麼呢.
你將你的靈魂來交換.
24年之後.
你將你的靈魂交給我.
24年裡面.
你心想事成.
想什麼就得什麼.
佛士德說.
OK GOOD.
因為我生命裡面.
衝不過很多的難關.
起碼24年裡面.
我能夠嘗試到.

$^{921}$能夠經歷到.
生活裡面各種的困難.
迎刃而解.
爽呀爽呀爽呀.
衝呀衝呀衝呀.
無止境地.
所有的慾望.
可以得到應驗.
愛情.
歡樂.
神遊各種地獄的階段.
天堂的各種階段.
他樣樣都心想事成.
在生命最後一刻.
他要交出他的靈魂的時候.
他才領悟到一個道理.
生命領悟到一個道理.
人生的目的.
是要接受限制.
和不圓滿.
人生裡面有一個目的.
對於佛士德來說.
他領悟到一個極大的道理.
就是接受限制.
和生命裡面的不圓滿.
當然了.
佛士德在研究法律.
研究醫學.
研究神學的種種裡面.
是很有成就.
在他二十四年裡面.
在這些各方面都有成就.
但我用這一段經文來看.
佛士德所領悟到的.
生命裡面種種的限制.
和不圓滿.
我們會怎樣想呢.
上帝總是用他自己的話語挑戰我們.
盡你所能靠乃恩典.
造成上帝交代給你的使命.

$^{961}$那個未必是最圓滿的.
按照我們的力量.
按照上帝給我們的恩賜.
個人恩賜多或少.
按照上帝給我們的.
我們盡我們自己能力.
靠乃上帝給我們的恩典.
我們盡量去做.
我們有很多人做得比我們更加好.
但是我已經盡力去到這處了.
只要你誠心對上帝說.
上帝啊 你給我這樣的恩賜.
這樣的力量.
我做到這裡就行了.
我盡力去做到這樣.
就行了.
這是我最好的.
最圓滿的.
就這樣了.
但是離開那個超圓滿.
完美的階段.
仍然相差甚遠.
上帝啊 求你給我節納.
我這樣的限制.
但是在這樣的情況下.
是我最好的了.
我們願意靠乃上帝給我們的恩典.
盡上帝給我們各種的恩賜.
學習到怎樣寫記.
背起自己的十字架.
有人可以在寫記.
背起十字架裡面.
承擔很大的苦難.
給我們承擔更大的苦難.
但是在我承受的力量之下.
我已經做到這麼多了.
我已經做到這麼好了.
我用兩個故事來做結束.
千多年來大陸很多教會被拆十字架.
十字架被拆.

$^{1001}$我想大家都知道.
很多人十字架被拆的時候.
真的很痛苦.
代表上帝臨在的十字架被拆.
真的很苦.
信仰被受挑戰.
上帝你在哪裡.
有很多我們讀到的報告.
有些人因為信主.
在父母的墳頭上都有個十字架.
連墳頭的碑也被拆.
很痛苦.
生命為何這麼痛苦.
信仰為何這麼痛苦.
但最近我讀到一篇講章.
是德國的一位神學家.
叫做Helmut Tillich.
叫做底納基.
他的書只有一本翻譯成中文.
很小的.
講神學生怎樣做一個神學生.
很小的.
但他其他的著作.
都沒有翻譯成中文.
但在他的講道集裡.
他提出了一些事實.
因為他當時反對希特勒的政策.
所以他是大學教授.
但做不到大學教授.
他是教會牧師.
不可以讓他講道.
即是等於他甚麼都做不到.
但身邊很多人不斷為他.
不斷為他爭取.
爭取到在一段很短的時間裡.
只讓他講一次道.
即是講道的次數大大的縮短.
縮短到很少的時間.
但他在柏林的那間教堂.
每逢到他講道的時候.

$^{1041}$就坐無虛席.
但當時第四次大戰.
德國的飛機不斷在上面盤旋.
做著做著.
下星期剛剛唱完詩.
警報一響.
他就要逃難.
就要逃走了.
在他的記載裡.
他很多時候都要逃走.
有時候回來教會.
教會已經被炸了.
就去另一個教會.
我們認知有些教會被拆了十字架.
但原來另一個地方的人.
一個傳道一個牧師.
他連講台都要隨時變換.
我知道有些地方都是這樣.
有一次他講完道.
回家途中.
拿著鑰匙.
回家入門口.
發現自己的房子都不存在.
被炸爛了.
拿著鑰匙.
站在一堆灰燼裡面.
欲哭無淚.
有一次在講道途中.
飛機來了.
大家都要相爭要走了.
走進附近的防空洞裡面.
教堂被炸了.
詩琴被炸死了.
在灰燼裡面.
有個姐妹走上來跟他說.
她說博士教授.
上個星期你在這裡講道的時候.
我的先生.
跟我一起在這裡聽.
但是今天他被炸死了.

$^{1081}$但我多謝你.
你上個星期的講道.
為他的死亡.
做了一個最好的準備.
讓他在聽你講道的時候.
知道生命的短暫.
一切都有上帝的恩典.
講完之後.
那個姐妹握手就再見了.
當我讀到這些記載的時候.
我便說上帝啊.
你比他的訓練大很多.
知道我們有限制.
靠賴恩典與上帝自己給我們各種的恩賜.
做好 盡力的做好.
在能力範圍裡面做到最好.
就夠了.
完全憐憫我們.
在結束的時候.
我讀一篇土文給大家聽.
偶爾我都會用.
這本是《司土日新》.
是一位蘇格蘭的實學家.
叫做John Bailey.
所寫的一本很短的.
他用31天.
寫了62篇早禱晚禱.
每一天裡面有兩篇早禱晚禱.
今天是25日.
我選了25日沒有選24日.
25日這篇晚禱裡面.
一段經文.
土稿來做今天晚上講道的結束.
希望大家用祈禱的心.
聽我重讀這篇土文.
聖善的主.
我決志已久.
奉獻我的生命與心靈為你服務.
我在你面前憂傷痛悔.
因為我還是傾向罪惡.

$^{1121}$很少服從你的旨意.
我太注重肉體的歡樂.
而不注重靈性的事.
太貪道逸樂.
而不重視永恆的福分.
喜愛懶散而不起愛勞力.
快快由器慢慢禱告.
服侍自己太起勁.
服侍別人太疏懶.
太急於取而淪於語.
眼高手低能說不能行.
有志向善而不能功行實踐.
責人太嚴而對己太寬.
喜歡挑出別人的過失.
而憎恨別人揭露自己的過失.
對大事既無力應對.
對小事又不捨履行.
對逆境缺乏毅力.
處順境太易滿足.
離開你就沒有辦法.
然而又不願意依靠你.
仁慈的上帝啊.
求你再次饒恕我.
垂聽我憂傷的傾訴.
並用你的慈愛.
從你的賊女圖滿我的過失.
求你賜我信心.
使我能夠把握你的聖潔.
而得以在救主耶穌基督的公義裡喜樂.
靠主的功勞而不靠自己的功勞.
漸漸學上基督.
我的意志要像祂一樣.
順服你的旨意.
奉耶穌基督的聖名.
阿們.
\newpage



\section{馬太福音 3:1-12}
\label{sec:A8lXH1cPVBk}
\textbf{2024-02-25 午堂《國度,權柄,榮耀》馬太福音 3\_1-12 講員:李思敬牧師}
\newline
\newline
連結: \href{https://youtube.com/watch?v=A8lXH1cPVBk}{\texttt{https://youtube.com/watch?v=A8lXH1cPVBk}} ~~~~ 語音日期: 2024-02-27
\newline
\newline
\hyperref[sec:dYh7F8XFvZo]{\small{< < < PREV SERMON < < <}}
~
\hyperref[sec:index_chronic]{\small{[返順時目]}}
~
\hyperref[sec:index_scriptual]{\small{[返順卷目]}}
~
\hyperref[sec:__QTchIrjDM]{\small{> > > NEXT SERMON > > >}}
\newline
\newline
馬太福音 3:1-12
\newline
\begin{longtable}{cl}
\hline
\hline
章節 & 經文 (和合本修訂版)\\
\hline
3:1 & \begin{tabularx}{0.7\textwidth}{X} 在那些日子，施洗的約翰出來，在猶太的曠野宣講： \end{tabularx} \\ \\ \relax
3:2 & \begin{tabularx}{0.7\textwidth}{X} 「你們要悔改！因為天國近了。」 \end{tabularx} \\ \\ \relax
3:3 & \begin{tabularx}{0.7\textwidth}{X} 這人就是以賽亞先知所說的：「在曠野有聲音呼喊著：預備主的道，修直他的路。」 \end{tabularx} \\ \\ \relax
3:4 & \begin{tabularx}{0.7\textwidth}{X} 這約翰身穿駱駝毛的衣服，腰束皮帶，吃的是蝗蟲和野蜜。 \end{tabularx} \\ \\ \relax
3:5 & \begin{tabularx}{0.7\textwidth}{X} 那時，耶路撒冷、全猶太和全約旦河地區的人，都到約翰那裡去， \end{tabularx} \\ \\ \relax
3:6 & \begin{tabularx}{0.7\textwidth}{X} 承認他們的罪，在約旦河裡受他的洗。 \end{tabularx} \\ \\ \relax
3:7 & \begin{tabularx}{0.7\textwidth}{X} 約翰看見許多法利賽人和撒都該人也來受洗，就對他們說：「毒蛇的孽種啊，誰指示你們逃避那將要來的憤怒呢？ \end{tabularx} \\ \\ \relax
3:8 & \begin{tabularx}{0.7\textwidth}{X} 你們要結出果子來，和悔改的心相稱。 \end{tabularx} \\ \\ \relax
3:9 & \begin{tabularx}{0.7\textwidth}{X} 不要自己心裡說：『我們有亞伯拉罕為祖宗。』我告訴你們，神能從這些石頭中給亞伯拉罕興起子孫來。 \end{tabularx} \\ \\ \relax
3:10 & \begin{tabularx}{0.7\textwidth}{X} 現在斧子已經放在樹根上，凡不結好果子的樹就砍下來，丟在火裡。 \end{tabularx} \\ \\ \relax
3:11 & \begin{tabularx}{0.7\textwidth}{X} 我是用水給你們施洗，叫你們悔改；但那在我以後來的，能力比我更大，我就是給他提鞋子也不配，他要用聖靈與火給你們施洗。 \end{tabularx} \\ \\ \relax
3:12 & \begin{tabularx}{0.7\textwidth}{X} 他手裡拿著簸箕，要揚淨他的穀物，把麥子收在倉裡，把糠用不滅的火燒盡。」 \end{tabularx} \\ \\
[1ex]
\hline
\hline
\end{longtable}
$^{1}$請主內平安.
(主內平安).
進入農曆新的一年.
感恩.
我們不用好像道聽途說.
說甚麼農年.
聖經裡面說農.
說魔鬼的.
那就慘了 明年是甚麼年.
蛇年.
農蛇馬羊.
你背下十二生肖.
鼠牛虎兔龍蛇馬羊.
猴雞犬豬.
沒有甚麼好.
唯一一個是聖經沒有提的.
你知道是哪一個嗎.
全本聖經都沒有的動物.
老虎.
巴勒斯坦 迦南地.
是沒有老虎的.
所以唯一好的年度就是虎年.
這樣是迷信.
我們比普通人還要迷信得多.
聖經用這些動物.
其實是比喻.
好像羊走迷.
各人偏行幾路.
不是教我們去討厭那隻羊.
比德三次不認主.
雞就叫了.
以後不吃雞.
做基督徒其實不是這樣做法的.
所以小心.
有很多時候我們聽到一些似是而非的說法.
我們真的需要神的說話.
來幫助我們去分辨清楚.
所以今天我們跟大家讀的經文.
是施洗約翰.
傳道.

$^{41}$在約旦河.
施洗.
他所說天國近了.
按照馬太福音的記載.
天國的福音.
不是耶穌第一個來傳揚的.
大家知道吧.
是施洗約翰先傳講天國的福音.
然後耶穌接著出來傳道.
他重複同一句說話.
天國近了.
然後馬太福音第十章告訴我們.
當耶穌差遣他的門徒.
無論是十二個也好.
七十個.
兩個兩個出去也好.
你留心聖經的記載.
耶穌都是教導他們.
傳天國近了.
這個福音的訊息沒有改變.
天國的福音我們很少聽到.
我們比較熟悉的是.
保羅書信裡面所講解.
耶穌釘十字架.
耶穌為罪人死.
神的愛向我們顯明.
所以福音書的記載.
我們很少提.
於是很容易誤會.
我剛剛最近看到.
《經濟學人》這本周刊.
報導美國有教會.
靈恩的教會.
保守信仰的教會.
但他們說天國的福音.
就是要去投票選特朗普.
嘩.
是呀.
很認真的查經的結論.
你領頭也好.

$^{81}$他們覺得今年十一月.
就是天國和魔鬼的決戰.
我們做基督徒的.
就要去實踐天國的福音.
這篇文章逼我想一想.
弟兄姊妹明不明白.
甚麼叫做天國的福音呢.
有人說.
其實福音書亦都告訴我們.
天國和神的國是相同的.
不過猶太人被毀.
馬太是猶太人.
所以他用天.
他不用神.
這個被毀是正確的.
有這樣的傳統.
到今天.
猶太人如果在美國.
他們寫神.
這個英文字.
你認識的.
God.
他們不是整個寫出來的.
他們是寫G-I-V-E-N-D.
沒有了個VOW.
不能讀的.
是呀,這個是他們的傳統.
不過今天我們不是要講這件事.
馬太福音其實沒有這樣的被毀.
何以見得.
我們現在手上沒有聖經.
你可能不能馬上翻查.
我提醒大家.
馬太福音十二章二十八節.
同一卷福音書.
馬太福音.
沒錯,他說.
施洗約翰傳天國的福音.
沒錯,他說耶穌來傳天國的福音.
門徒出去傳天國的福音.

$^{121}$但是當他記載到.
耶穌和法利賽人辯論.
趕鬼的權柄.
法利賽人批評耶穌.
他靠鬼王別西卜趕鬼.
你聽過吧.
耶穌說不是的.
如果我是憑著神的靈趕鬼.
耶穌接著說甚麼.
那神的國就在你們中間.
就來到了.
所以馬太沒有被毀.
他照記載神的國.
其實說到神的國.
說到國度.
最多的地方是主禱文.
不過我們很少有這樣的留意.
剛才我們起立.
懂得背《使徒信教》.
你們不用背的.
看著就行了.
背的.
不過是哪個版本的意思.
我背的是宣導會.
生命聖詩的版本.
去到禮賢會閉上眼背就撞版了.
原來中文翻譯不同的.
翻譯可以不同.
不過主禱文應該大家都懂得背.
因為是閉上眼來祈禱的.
主禱文都有不同的版本.
有文言文的.
有白話文的.
白話文都有起碼三四個不同的翻譯.
你記得的.
主禱文開始的時候.
「與你的國降臨」.
為什麼呢.
因為國度權柄榮耀.
有貓子.

$^{161}$所以不如關掉PowerPoint.
不要通水.
「都是你的全是父的」.
主禱文首尾呼應.
「用神的國開始」.
「用神的國度結束」.
我們很少留意.
我們最記得的就是我們要用飲食.
父賜給我們.
待會兒爆位喝茶.
最重要留到好的位置.
不是不是不是.
聖經不是這樣記載的.
聖經怎麼說.
國度權柄榮耀.
三件事情一氣呵成.
幫助我們明白.
耶穌所說的神的國的福音.
到底是一件什麼事.
我們先從榮耀開始.
這句話.
不知道你們有沒有印象.
約翰福音十七章第一節.
耶穌知道他要上十字架.
他帶領門徒去喀西瑪尼.
三本《馬太馬可》路加.
記載喀西瑪尼耶穌和門徒祈禱的經過.
約翰福音特別在十七章開始.
記載了耶穌一段很長的祈禱.
似乎和其他福音書的記載不同.
是互相的補充.
約翰福音十七章第一節.
記載耶穌怎麼祈禱.
他說「父啊,時候到了」.
什麼時候.
我們大家知道了.
福音書記載到耶穌要上十字架.
但跟著耶穌祈禱的兩句說話.
你明不明白他在說什麼.
「願你榮耀你的兒子」.

$^{201}$也都讓兒子榮耀你.
明明是上十字架.
痛苦,羞辱的記號.
但卻成為聖父,聖子得榮耀的時刻.
不需要等到復活.
不需要等到升天.
人子被高舉掛在木頭上.
就吸引萬人來歸向基督.
這是福音書很清楚的焦點.
什麼是神的榮耀.
不是我們心裡想的那些.
八個警察騎著摩托車開路.
有八個在後面威風瀾瀾.
不是,那些是人間的榮耀.
對,你和我從來都沒有機會這樣做.
後面有一輛跟著你就害怕了.
但耶穌祈禱.
他明白什麼是榮耀.
什麼是父神的榮耀.
什麼是耶穌基督的榮耀.
原來按照神拯救的計劃.
為罪人死在十字架上.
這個時刻不再是人間視為痛苦羞辱.
我們唱詩唱到今天都是這樣唱.
「各擱他十字架,痛苦羞辱」記號.
原來不是.
是一個榮耀的記號.
所以為什麼宣導會要將十字架.
放在我們回歸的第一個位置.
因為耶穌基督是拯救之主.
你們還將祂升上屋頂.
你知道「敢宣」的標誌嗎.
這是很重要對神的榮耀的理解.
神的榮耀就是天上出現長雲.
就是有異象顯現.
現在來到就跪拜就降伏.
完全不是這回事.
好像我們說十二生肖.
不要搞錯.
那是民間宗教.

$^{241}$使我們洗腦了.
覺得神的出現就一定是這樣.
原來聖經不是這樣教的.
十字架底下有沒有人跪拜耶穌.
有沒有人相信耶穌.
十字架上面的強度.
相信耶穌而已.
奇怪了.
上帝的榮耀按照聖經的教導.
原來是一個我們看不到的事實.
今天我們明白.
的確十字架是神的榮耀.
因為從主耶穌到今天.
多少人就是因為這個救恩福音.
上帝的慈愛生命扭轉過來.
但是你看著這個十字架.
你看著主耶穌在上面斷氣.
你看到神的榮耀嗎.
在場的人.
敬虔的婦女門徒.
他們隨著空很傷心的離去.
看不到.
但原來這個是榮耀.
所以彼得教導我們.
彼得前書第二章二十一節.
你不要現在找聖經了.
除非你有手機.
點它就出來.
否則你找不到.
我就背完了.
彼得提醒基督徒.
「患蒙召」.
甚麼?.
「原是為此」.
為甚麼蒙召?.
上帝呼召我.
叫我出人頭地.
叫我成功.
叫我別人逆境.
我們順境.

$^{281}$不是的.
彼得接著說的呼召.
是基督為我們受苦.
留下的腳蹤.
原來耶穌基督呼召我們去受苦.
其他的宗教不會這樣說的.
是主耶穌先至夠膽這樣說的.
「你們蒙召原是為此」.
因為基督為我們受苦.
留下這個受苦的腳蹤.
你願意跟隨的.
原來是一位受苦的耶穌基督.
不是你和我心目中的榮耀.
看清楚.
但他的確是這樣叫我們.
跟從我.
並且耶穌為誰受苦?.
耶穌不是為那些信祂的人受苦.
曾經有一段時間.
歐洲的神學認為.
耶穌基督的受苦,受死.
只是為那些信祂的人受苦,受死.
這是異端.
雖然是很保守的神學教導.
聖經不是這樣說.
聖經記載的是甚麼呢?.
在我們還作罪人的時候.
耶穌基督為我們死.
甚麼是罪人?.
希伯來書說得好.
他忍受罪人頂撞.
耶穌基督為那些頂撞他的人受苦.
我們很懂得頂撞他.
輕輕動一點點.
我們就叫了.
上帝去了哪裡?.
你以為那是你的非庸,孕庸嗎?.
他回到屯契.
你也不能這樣叫.
他不是隨你叫.

$^{321}$就要出來Yes Sir.
聽你擺佈.
但我們就是這樣.
在我們頂撞上帝的時候.
他為我們死在十字架上.
這是一個榮耀的時刻.
我們想都想不到.
是沒得想的.
是聖經教的.
聖經幫助我們去明白.
原來耶穌基督看十字架.
是他的榮耀.
因為這個是父神的榮耀.
為甚麼是他的榮耀?.
原來神不需要像我們所想的那種榮耀.
我們沒有.
我們不是上帝.
所以我們要整色整水.
來榮耀自己.
上帝不用的.
所以他的榮耀.
就是他的拯救.
就是他的恩典.
就是他的愛.
馬丁路德稱這個為十字架的神學.
有沒有聽過?.
表面聽我們以為.
這個是教導我們受苦.
原來不是.
馬丁路德說.
這個是認識上帝的第一步.
因為我們太多迷信的觀念.
太多拜神,拜菩薩,拜偶像.
錯的出發點.
以為神的榮耀.
就是這樣的.
原來不是.
剛剛相反.
所以從榮耀開始.
我們要很小心.

$^{361}$讀聖經.
按照聖經的教導.
去明白.
你跟從耶穌.
就不會跟錯.
行不行?.
那甚麼是權柄?.
剛才我們提過.
《馬太福音》第十二章二十八節.
耶穌討論.
趕鬼是要權柄的.
趕鬼.
你沒有權柄的.
鬼上你身的.
聖經都有這樣的記載.
不過不是主要的記載.
所以你不用怕.
不要去碰這些東西.
不要去玩這些東西.
更加重要的另一段經文.
我相信我提出來.
大家會有印象.
《路加福音》第十章十七到二十節.
就是剛才我說.
耶穌差遣他的門徒出去傳道.
《路加福音》第十章就記載.
那班門徒回來了.
Debriefing.
你們有沒有?.
出去派完單章之後回來.
不是肚餓了.
大家去喝茶.
先來談一談.
剛才有什麼見證可以分享.
門徒分享什麼?.
你記得吧.
門徒不是分享人信耶穌.
門徒最興奮的是說什麼?.
鬼也服了我們.
我相信你也一樣.

$^{401}$如果你出去派單章.
你成功趕鬼.
你回來見證.
口沫橫飛.
說沒有試過.
連鬼也服了我們.
是不是值得高興呢?.
是不是值得高興?.
我們每個人都不敢說不是.
但耶穌說.
耶穌怎樣回答門徒?.
我想跟你讀.
但你沒有聖經,來不及.
耶穌說.
不要因為鬼服了你們而歡喜快樂.
不是不值得歡喜快樂.
但這不是最快樂.
更快樂的是什麼?.
這句你要記得.
要因為你們的名字寫在天上而快樂.
更重要的.
你要問.
你的名字是否記載在天上.
我出去傳道.
我出去趕鬼.
鬼也服了我.
我的名字還不在天上?.
又回到《馬太福音》.
登山寶訓.
臨結束之前.
一提.
你一定記得.
凡稱呼我主也的人.
什麼?.
不能到盡天國.
唯獨遵行我父子義的.
才能進去.
到那一天.
耶穌說.
很多人來跟我說.

$^{441}$我奉你的名趕鬼.
我奉你的名行神蹟.
耶穌說我不認識他們.
有沒有讀這段聖經?.
有沒有?.
耶穌說.
很多人來.
我曾經在錦繡派過單張.
耶穌會說我不認識你?.
開玩笑.
不可能.
你們.
我們今天要說普通話.
耶穌才能聽.
不是.
聖經是這樣教的.
神蹟騎士.
趕鬼醫病.
都進不了天國.
為什麼?.
懂不懂得讀聖經?.
為什麼?.
很簡單.
剛才耶穌說什麼?.
什麼人能夠進天國?.
遵行天父旨意的人.
沒有說遲到的人不能進天國.
如果不是我們就慘了.
但是遵行天父旨意.
所說的.
不是你準時或不準時出席崇拜.
也不是你有沒有打瞌睡.
打鼻鼾.
這麼表面.
聽得到嗎?.
你如何遵行天父的旨意.
第一.
原來這裡說的權柄.
不是說污鬼.
鬼一定服在神的權柄之下.

$^{481}$這是很矛盾.
很弔詭的事實.
是誰不服從上帝的權柄?.
你和我.
是人.
特別那些懂得趕鬼的.
特別那些給神用的.
你聽得懂聖經的記載嗎?.
是呀.
在教會裡面挺胸昂首.
走進來個個鞠躬.
那些要小心.
耶穌是這樣教的.
不是我說的.
耶穌說什麼?.
再聽一次.
凡稱呼我主.
凡是趕鬼醫病.
為什麼你不認識我?.
你有沒有遵行天父的旨意?.
這叫做神的權柄.
原來不是外表的恩賜.
不是事奉的果孝.
不是人的尊敬.
抑或人對你有什麼評價.
是你在你的生命當中.
甚至在人看不到的角落.
你有沒有遵行天父的旨意?.
所以主禱文是這樣祈禱的.
願你的國降臨.
然後呢?.
願你的旨意行在地上.
聽到那個關係了.
前面第一句.
願人都尊你的名為聖.
我們覺得不關我們事.
人,人地,其實是譯錯了.
其實中文不應該這樣譯.
文言文好一點.
願以名聖.

$^{521}$他沒有說哪個尊神的名為聖.
英文你會背的.
會不會?.
中學的時候有沒有要背英文主禱文?.
當然要的.
聖公會的學校.
我的學校從頭到尾.
每個人都背.
Our Father who art in heaven.
Hallowed be your name.
Passive voice.
哪個尊上帝的名為聖.
不是說人地的.
當然包括撒旦.
撒旦怕上帝的.
你看看約伯記第一第二章.
是哪個嘴滑滑不怕上帝.
就有你和我.
是呀.
我們是否尊神的名為聖.
我們是否懂得敬畏上帝.
遵行祂的吩咐.
所以耶穌在赫西瑪尼的祈禱求什麼?.
不是求那個杯拿去.
而是求你的旨意成就.
馬太夫人記載.
耶穌第二次回去祈禱.
說的話和先前的一樣.
不是那個是馬可和路加的記載.
馬太夫人記載了一句說話.
願你的旨意成就.
希臘文和主禱文這句說話完全一致.
耶穌不單止教我們這樣祈禱.
耶穌自己是這樣祈禱.
祂要成就的.
祂要進行的.
是天父的旨意.
不是祂的意思.
所以彼得弟弟子勸耶穌.
千萬不可以上耶路撒冷.

$^{561}$尤其是你知道有人在埋伏你.
耶穌說你不體貼神的意思.
看到嗎?.
我們只記得耶穌罵他撒旦.
其實你要明白原因在哪裡.
神的權柄是我們是否進行祂的旨意.
我小時候我們經常都記得一句說話.
人民眼睛是雪亮的.
今天沒有人說了.
周圍的人看得到.
你作為一個基督徒.
你是聽誰的?.
自己?.
還是上帝?.
看得到的.
不過不容易.
不認識你的人看不到.
最認識你身邊的人最看得到.
求主幫助我們.
甚麼是榮耀?.
甚麼是權柄?.
原來不是我們心目中那一套.
原來耶穌所教導的是受苦.
就是榮耀.
就是聽神的旨意.
就是見證神的榮耀.
在我的身上.
不是在鬼.
不是在他.
在我的身上.
這個才是神的權柄.
說甚麼是國度?.
我差不多過了時間.
不過給我多兩分鐘說完結論.
不然就很遺憾.
其實一氣呵成.
國度我們不明白.
因為我們不是當時的猶太人.
香港人應該明白多一點.
因為我們不是一個國家.

$^{601}$得了的.
香港是國家.
23條的.
香港從來都不是一個國家.
所以說國家其實是很遙遠的事.
殖民地的時候.
我們的大老闆在銀紙上印著.
在倫敦.
在白金漢宮.
在英女皇.
現在死了.
不過我們不會再印張King Charles的銀紙出來.
今天我們的國家在北京.
在中央.
他派人下來.
關心我們.
很遙遠.
對猶太人來說.
國家不是在耶路撒冷.
你知道嗎?.
他們亡國很久了.
經歷了馬代波斯的帝國.
希臘的帝國.
然後現在來到羅馬的帝國.
都不是在耶路撒冷.
聽得明白嗎?.
馬代波斯的帝國.
在書山城.
是首都.
今天的伊朗.
希臘的帝國的首都.
在今天的希臘.
愛琴海.
羅馬的首都.
在羅馬.
意大利.
所以對於猶太人來說.
什麼是國度.
其實是很遙遠的事.
很看不到的事.

$^{641}$改了嗎?.
不是馬代波斯嗎?.
輸了給亞歷山大了.
不是在眼前發生的.
沒有七一移交儀式.
原來已經是羅馬了.
這個帝國.
它的權柄.
它的榮耀.
會來到耶路撒冷.
但是.
不是發生在我們眼前.
猶太人摸摸頭.
OK.
那就等吧.
等它來吧.
所以你明白什麼叫天國要來嗎?.
願你的國降臨.
但是Kingdom come.
是它來的.
不是你去拓展天國,神國.
王明道先生寫過一本很薄的書.
你大概沒有看過.
我們讀大學的時候.
每個人都要看的.
人能建設天國嗎?.
這樣問答案很明顯.
不能.
猶太人在耶路撒冷.
能夠建設羅馬帝國.
建設希臘帝國.
建設馬代波斯帝國.
開玩笑.
不可能.
不是說這個.
說的是看不到.
但卻是一個事實.
耶穌如何進入耶路撒冷城.
祂騎一隻驢.
你記得嗎?.

$^{681}$這是一個暗號.
你知道嗎?.
不知道.
那你和彼拉多一樣.
羅馬人不知道.
因為他們不讀聖經.
讀聖經的猶太人.
讀到列王記上第一章.
大衛病重在床上.
然後知道他的第四個王子亞多利亞摧國.
得到軍隊的擁戴.
要繼位.
大衛很聰明.
雖然很軟弱.
他說吩咐人將他騎的那隻驢拉出來.
記得嗎?.
有沒有印象?.
讓所羅門騎在上面.
拉出去.
大家就知道.
誰才是用中國人的說法.
真命天子.
你說騎一隻驢出去.
你就不認得.
今天香港人就要.
開我的奔馳出去.
你去停車場.
光是看到後尾的燈.
你都認得自己的車.
不用看到車牌才認得.
所以城市人是認不到動物的.
但耶路撒冷一樣不認識.
你是電動車還是輕行車.
還是柴油車.
那個是文化背景.
但耶穌這樣進入耶路撒冷.
猶太人的反應是.
Hosanna.
記得嗎?.
逢主名來的.

$^{721}$怎樣?.
是應當稱頌的.
彼拉多還要審他.
你是不是猶太人的王.
耶穌回答得很清楚.
我為此而生.
我是猶太人.
我是王.
不過我的國不屬於這個世界.
所以神的國.
不是粉瑞黨的政綱.
要推翻羅馬政權.
要趕走羅馬駐軍.
神的國也不是瑞利.
祭司,文士.
跟帝國合作.
忍辱負重.
幫助百姓有好日子過.
是的.
你讀歷史.
去到公元70年.
當粉瑞黨.
好像哈馬斯.
真的去攻擊羅馬軍隊的時候.
結果耶路撒冷被毀.
第二聖殿再次被焚毀.
這是歷史.
耶穌基督的國度看不到.
但事實他已經是王.
天上地下所有的權柄.
你背下去嗎?.
都次級餓了.
用不用等他第二次回來.
才掌握天上地下的權柄?.
不是的.
不要搞錯.
他今天已經坐著為王.
但看不到.
為什麼看不到?.
上帝讓我看看.

$^{761}$你問錯問題了.
你當自己是老闆.
上帝是菲傭.
買了什麼送回來?.
讓我看看收據.
不是的.
上帝不是讓我們這樣問的.
作為神的僕人.
我們只能夠問.
不是問為什麼.
是問在這樣的情形下.
你想我們怎樣?.
有什麼吩咐?.
懂得問問題嗎?.
神的國是一個事實.
耶穌基督已經作王.
所以他吩咐我們出去.
我就尚與你們同在.
不只是耶路撒冷.
是直到地極.
所以今天我們在香港此時此地.
我們跟從耶穌.
我們遵行天父的旨意.
我們明白福音的真理.
別人看不到.
不是每個人都看得到.
他看不到的.
我們出去遊行.
別人懂得抬什麼神廟出去.
我們又抬十字架出去.
是沒用的.
看不到就看不到.
你明白聖經.
但是你的見證.
你的生命.
甚至你的受苦.
你願意為彼得的教導.
為義受苦.
耶穌的說法.
為義受逼迫.

$^{801}$為我的名受逼迫.
我經常說.
教會在外面的人看來.
是不會找他們笨的.
反過來說那句.
是我們可以去找教會笨的.
有沒有聽過.
在海外很多這樣的說法.
找教會幫忙.
很肯幫你的.
送你去買所有的東西.
那你就逃跑了.
不用留下來.
這樣被人找笨.
我們以後不讓人找笨.
其實被人找笨.
都是為義受苦.
如果你是為了耶穌的名.
你願意.
願意吃虧.
神的國.
神的權柄.
神的榮耀.
你再想清楚一點.
這個不單只是聖經的教導.
更加是主耶穌留下.
給我們跟從祂的人的腳踵.
我們同心一起祈禱.
多謝主.
其實你不是要我們.
做一些我們做不到的神蹟奇事.
醫病趕鬼.
你教導我們.
願意進行你的旨意.
願意信服你的大靈.
願意甘心.
歡喜.
感恩.
來經歷你給我們所經歷的.
不是偶然.

$^{841}$不是碰巧.
是要榮耀你.
是要彰顯你的權柄.
亦都是讓人認識真神.
可以進入神的國.
多謝主.
求你幫助我們.
聖靈教導我們.
將你的話藏在我們心中.
免得我們得罪頂撞你.
誰聽我們的祈禱.
我們也祈禱奉耶穌基督的聖名.
阿們.
\newpage



\section{創世記 25:27-34}
\label{sec:__QTchIrjDM}
\textbf{2024-02-25 早堂《當你以為是理所當然》創世記 25\_27-34 講員:鄧泰康傳道}
\newline
\newline
連結: \href{https://youtube.com/watch?v=_-QTchIrjDM}{\texttt{https://youtube.com/watch?v=\_-QTchIrjDM}} ~~~~ 語音日期: 2024-02-27
\newline
\newline
\hyperref[sec:A8lXH1cPVBk]{\small{< < < PREV SERMON < < <}}
~
\hyperref[sec:index_chronic]{\small{[返順時目]}}
~
\hyperref[sec:index_scriptual]{\small{[返順卷目]}}
~
\hyperref[sec:exEZqq1nt80]{\small{> > > NEXT SERMON > > >}}
\newline
\newline
創世記 25:27-34
\newline
\begin{longtable}{cl}
\hline
\hline
章節 & 經文 (和合本修訂版)\\
\hline
25:27 & \begin{tabularx}{0.7\textwidth}{X} 兩個孩子漸漸長大，以掃善於打獵，常在田野；雅各為人安靜，常住在帳棚裡。 \end{tabularx} \\ \\ \relax
25:28 & \begin{tabularx}{0.7\textwidth}{X} 以撒愛以掃，因為常吃他的野味；利百加卻愛雅各。 \end{tabularx} \\ \\ \relax
25:29 & \begin{tabularx}{0.7\textwidth}{X} 有一天，雅各熬了湯，以掃從田野回來，疲憊不堪。 \end{tabularx} \\ \\ \relax
25:30 & \begin{tabularx}{0.7\textwidth}{X} 以掃對雅各說：「我累死了，請你讓我吃這紅的，這紅的湯吧！」因此以掃又叫以東。 \end{tabularx} \\ \\ \relax
25:31 & \begin{tabularx}{0.7\textwidth}{X} 雅各說：「你今日把長子的名分賣給我吧。」 \end{tabularx} \\ \\ \relax
25:32 & \begin{tabularx}{0.7\textwidth}{X} 以掃說：「看哪，我快要死了，這長子的名分對我有甚麼用呢？」 \end{tabularx} \\ \\ \relax
25:33 & \begin{tabularx}{0.7\textwidth}{X} 雅各說：「你今日對我起誓吧。」以掃就向他起誓，把長子的名分賣給了雅各。 \end{tabularx} \\ \\ \relax
25:34 & \begin{tabularx}{0.7\textwidth}{X} 於是雅各把餅和豆湯給了以掃，以掃吃喝以後，起來走了。這樣，以掃輕看他長子的名分。 \end{tabularx} \\ \\
[1ex]
\hline
\hline
\end{longtable}
$^{1}$頂尖的早晨.
今天跟大家講一篇耳熟能詳.
整個月大家有留意我們的同工.
都會一起分享關於創世記裡面的人物.
今天會跟大家分享一個耳熟.
耳熟的故事不是很多.
裡面的時候其實.
但這個絕對是經典.
每一個小朋友從教會一直成長.
一定會聽過這個故事.
當這個故事一出現的時候.
奇怪的地方是什麼呢?.
先當作大家完全不認識他們兩兄弟.
他們兩兄弟大家很清楚.
媽生的.
即是媽生仔.
他出生的時候有一個很特別的圖畫.
就是哥出來.
耳熟出來.
之後就有啞哥.
他做了一個動作.
很經典的.
就是抓緊哥的腳.
所以啞哥的意思就是抓緊.
所以你會發覺.
在他們出來的時候.
最特別的地方就是.
媽生仔誰做哥.
誰做弟弟.
就差幾十秒.
就決定了.
整體上好像沒什麼特別.
但對於當時他們這個猶太世界裡面.
其實是很大分別的.
一會兒我會詳細的講.
究竟做長子有什麼福利.
有什麼好處.
但很特別的地方就是.
當他們說長大一講這段經文的時候.
介紹了他們一些很基本的性格之後.

$^{41}$突然間加插了一句很特別的.
就是以殺愛以掃.
因為時時食他的野味.
李伯家愛啞哥.
突然間插了這句下去.
如果大家一路去看完.
整個他們兩兄弟的故事.
和這個家庭的時候.
我幾肯定的有一個.
大家看完的感覺和印象.
這個家庭是一點都不和諧的.
你會發覺第一就是.
父母各自愛一個.
你偏重他我偏重他.
這裡加了一句.
這裡當然是為了之後交代下一個故事.
但為什麼爸爸愛大兒子.
因為時時食他的野味.
不知道是不是大兒子煮完之後.
不給媽媽吃.
這個我不清楚.
但總之這是一個很大的原因.
但你會發覺整個故事的發生.
我想大家留意.
從一些很小的事.
一路發展下去的時間.
去到最後的時候.
大哥想殺了小孩.
而在整個故事當中.
整件事的起始點.
是爸爸媽媽各自偏愛一個.
就是說其實這個不和諧的家庭裡面.
其實暗藏了很多非常不開心的事.
從下一個故事你就知道.
就是當以實想去祝福大兒子的時候.
其實祝福大兒子這件事是很理所當然的.
按照當時的習俗.
但他是暗地裡進行的.
他是暗地裡和大兒子說.
我的命不長.

$^{81}$你去打獵.
做一隻美味回來給我吃.
然後我就為你祝福.
他沒有通知老婆.
他沒有通知下一個.
聖經特別的地方就是.
又被李百嘉偷聽到.
然後媽媽就教弟弟騙老爸.
我想給大家看到的就是.
其實兩兄弟這次出現這件事.
不是突然間發生的.
是整個家庭的氛圍和教育出來的.
爸爸媽媽在家裡的角色究竟是什麼呢?.
很多時候我去認識一些爸爸媽媽.
特別是一些年輕的.
剛生小孩的爸爸媽媽.
很多時候我們都會有一個錯覺.
很感恩的.
在近幾十年其實好了很多.
以前的時候是更加的.
我爸爸媽媽那個年代是更加的.
我很努力賺錢回來.
養給你們食物.
供應你們.
供書教學.
然後時間幫你們洗衣服煮飯.
我就任務完成了.
就會覺得小朋友自動就會長大.
然後有一個假設.
只需要我做好基本條件.
子女就會好了.
會覺得小朋友的成長在父母眼中.
變得好是理所當然的.
但是是不是這樣呢?.
剛才我說過.
父母在家裡的角色是什麼呢?.
其實在聖經裡.
上帝所教導的.
是否我們帶動著整個家庭.
去能夠成為彼此相愛呢?.

$^{121}$但是現在出現的情況就是彼此傷害.
你會發覺其實今天做父母的時候.
是不容易的.
有些年輕人.
特別我做婚前輔導的時候.
他們問我一個問題.
就是我什麼時候才適合生小朋友.
我就反過來.
通常我就不回答.
我就先問了.
那你認為呢?.
我就說他心裡的憂慮.
他說什麼呢?.
他說我現在一個月.
才得萬多二萬元工資.
我什麼時候才儲夠錢生小朋友呢?.
我就問他.
我說生小朋友要很貴嗎?.
他說坊間外面告訴我.
當時生小朋友.
是要五百萬的.
我們去關注的反而是經濟上.
所以很自然的.
我們就會在經濟上.
我們就會在各樣生活上的條件.
在所謂的學術栽培上.
我們放了很多心機去幫這個小朋友.
我好記得.
在我剛剛讀完神學畢業之前.
我預備當時來實習.
就是來感受實習的時候.
我有一個很短的時間的放假.
我就說遲些就很忙了.
不如這樣吧.
遲些都沒什麼時間陪兒子了.
我就不如省下點時間.
跟他去泰國玩吧.
玩完之後就很開心.
當時我都很記得.
兒子當然是很開心.

$^{161}$在酒店跳跳跳.
在床上跳跳跳.
我就充滿著很多回憶.
之後到我兒子長大之後.
我問他.
我說你記不記得那時候跟他去泰國玩?.
他說不記得.
他說如果不是有照片.
他一點印象都沒有.
但他最記得什麼?.
他最記得的反而是.
很多陪他去沙灘對哨.
他好記得爸爸抱著他.
我想強調的.
小朋友今天我們.
當你以為你所提供給他的東西.
是最好的時候.
並不是.
他要的是聖經所說的那樣東西.
是愛.
那種關係.
我會形容的就是.
我們今天憑什麼去教導小朋友呢?.
作為基督徒.
我很肯定的.
我們是憑上帝自己的話語去教導他.
但我每次都不斷強調.
我們今天最重要的.
不是從你的口去告訴他聖經說什麼.
而是從你的生命.
從你的行為去行出來.
去告訴他上帝是如何.
你和上帝的關係是怎樣.
我可以告訴弟兄姊妹.
小朋友是看得一清二楚.
在這裡的時候.
我不妨分享.
其實我今天看到那對夫婦在這裡.
其實我很開心的.
之前我們一直祈禱會.

$^{201}$特別是這個年代.
我再說一句.
我們大人都尚且已經不回家.
我們都沒有什麼動力想回來.
但每一次祈禱會.
我們都見到一家四口.
爸爸媽媽帶著兩個小朋友.
兩個兒子回來.
大一個我估計應該是高中了.
應該要準備考DSE的時間.
小朋友可能還是小學都不一定.
我每次見到.
我心裡都很欣慰.
小朋友為什麼這麼好呢.
和爸爸媽媽每個星期從不間斷.
跟著他們一起.
在這裡一起祈禱.
特別是之前我們的祈禱形式.
我們不是太多互動.
我們真的同工上利.
在這裡一起公討.
帶大家領土.
就在下面的時候一起參與.
為什麼會這樣呢.
這個是不是今天我們很多做父母.
我們很希望小朋友身上見到呢.
為什麼會發生呢.
有一次當我們同工一起去觀塘宣道會.
去交流.
我們和他們談.
特別是青少年的模樣.
他們的青少年做得很好.
後來同工都和我分享.
這件事是很深刻的.
其實不是我們好.
是我們很感恩.
我們的弟兄姊妹的爸爸媽媽好.
因為在他們的經驗裡.
如果家庭的信仰生活做得好.
基本上小朋友回來教會.

$^{241}$和成長的機率是高很多.
是高很多.
因為我們不用和他們拆那麼多蛋.
在我的模樣下.
我也是.
我由年輕的時候模樣.
年輕人開始到現在.
我都很真實.
特別當時我記得.
裡面有些信義代.
很小的時候.
他走來和我說.
家裡有很多問題.
我和他談.
我說不是的.
和他分享鼓勵他.
後來他回來教會不好.
每次回來教會很好.
回到家之後.
一個星期再回來.
他就和我說.
鄧傳道你騙人的.
鄧導師當時叫鄧導師.
你騙人的.
我問他為什麼騙你.
你說上帝在這裡.
我說上帝是在這裡.
他說都不是的.
明明那時候我都已經原諒了爸爸媽媽.
但我回到家的時候.
他又是這樣.
其實我想說.
我們做年輕人工作.
要拆很多蛋.
他所見到的信仰.
裡面不是這樣.
當你口裡告訴他.
愛愛愛的時候.
他告訴我.
我們完全感受不到愛.

$^{281}$我爸爸媽媽同樣都是這樣.
我爸爸媽媽一樣.
我想強調不是不理解.
我想大家都明白.
我們成年人很明白.
這個世代是忙了很多.
包括學生都忙了很多.
但問題是.
我們有沒有留下時間.
特意去安排時間.
去和我們的小朋友陪伴他們.
所以有時候我們輕鬆一點說.
當今天有些照顧很小的小朋友.
由幼稚園到小學.
大家有時候做父母.
我們都會順一下.
很累啊.
很辛苦.
我說你享受一下辛苦吧.
不辛苦很久了.
因為年輕人.
像我兒子一樣.
一上中學.
他不黏著你了.
他就會跟同學一起.
弟兄姊妹.
其實今天我們有一樣東西.
我們不要以為.
幸福的家庭生活.
是理所當然.
其實我們是要.
放時間下去去學的.
每一個年齡階段裡面.
年輕人所面對的問題.
都不同.
我們究竟是陪伴.
還是我們只是口頭裡面的管教呢.
《政經》裡面只教我們很簡單.
教養孩童.
使他走當行的道.

$^{321}$就是道老也不偏離.
弟兄姊妹.
這是你的目標.
我們今天就算做.
一間這麼大的教會.
我們都是要說.
究竟我們的核心價值是什麼.
我們做的任何事.
都要達成我們的核心價值.
教小朋友.
就是教養孩童.
使他走當行的道.
就是道老也不偏離.
記住.
不是口頭上的教導.
而是很多很多.
我們要陪他一起走的時間.
有很多我們不理解.
是正常的.
因為我們的年代.
跟他們的年代不同.
而他們每一個年紀的階段.
所發生的事.
是不斷變化.
你以前覺得很work的方法.
不代表你今天會work.
所以其實這個很重要.
所以我以前去做一間教會.
比較上是基層一點.
爸爸媽媽真的完全很忙.
沒有時間可以陪.
但他們經常都跟我說.
我們的子女.
個個都走出去.
特別當時他們有些變壞了.
可能是黑社會或其他.
我每次都問父母.
我第一次走是何時?.
他第一次走是何時?.
然後他說.

$^{361}$跟我吵架的時候.
跟爸爸媽媽吵架的時候.
他就走出去.
然後我就有點兒無情地.
跟這些父母說.
其實是我們趕走他的.
是我們趕他離開這個家.
他沒有地方可以躲.
他就找他的朋友.
當他覺得外面的地方的愛.
比教會的愛好.
比家裡的愛好.
他就不再回來.
我記得那時我們教會裡.
有一些朋友是黑社會的.
後來有機會接觸順理耶穌.
我也問他.
為什麼你會參加黑社會?.
他說因為黑社會裡的朋友.
是絕對比我家裡的人好.
我問他為什麼?.
但原來很簡單.
一召集,即動.
然後是請喝請吃.
晚上不睡覺.
陪你不睡覺.
一起聊天聊到天亮.
很有趣的,弟兄姊妹.
原來這個原則從來都沒有變過.
就算放在黑社會也一樣.
適合用.
究竟我們有沒有這些時間.
有沒有這些心機.
去陪這些年輕人.
他們一起成長.
求上帝幫助我們.
我們是需要花心機的.
因為家庭是需要栽培.
第二件事.
耳掃一切都以為是上天注定.

$^{401}$剛才我提過.
他是上天注定了.
他是做長子的.
所以他就是比弟弟早一點出來.
就變成了哥哥.
我也很肯定的.
雅各最心生不忿的也是這樣.
是吧?.
差一點點而已.
但是長子的特權是什麼呢?.
原來不看也不知道這麼厲害.
一方面第一件事.
他可以獲得家庭的領導權.
第一件事就是如果爸爸不在.
長子就是一家之總結之長.
家裡一切的事都由長子作主.
第二件事.
這個最實際.
他可以得到雙份的產業.
就是說.
其實最公平.
我們心裡面.
我有100元.
我有兩個兒子.
最公平是什麼?.
分50/50.
每人分50.
Sorry.
在他們的世界裡不是這樣.
是要分三份的.
每一份都是33.333.
然後大兒子拿兩份.
小兒子拿一份.
差很遠.
第三件事.
當摩西的律法還沒有頒布之前.
原來長子是家族的祭司.
意思就是他代表了家庭和上帝的關係.
是一個很崇高的身份.
就是說他已經代替了祭司的身份.

$^{441}$第四件事.
他由出生就已經分別為聖歸給神.
就是說長子一出生的時候.
名利和權都拿了.
就是說其實差很遠.
而以蘇他去見到這個情況之下.
他就會覺得我是長子.
這件事不會改變.
中國人的說話叫做坐定粒碌.
我以後就是這樣.
但是弟兄姊妹你不要搞錯.
原來在上帝的國度.
由舊約至新約都一樣.
原來上帝的祝福和恩典.
雖然有一些很基本的框架.
但是最重要的是上帝的恩典和人的回應.
是成正比的.
不是的.
我想說的是上帝的恩典從來都在這裡.
我們都很清楚.
上帝的救恩是神愛世人.
他每一個都愛的.
但是你要得到的救恩條件是什麼.
是人的回應.
所以我們要決之.
我們要決定信耶穌.
要將祂成為我們生命裡面的主.
要人的回應.
但是現在出現的情況.
就是以蘇他覺得.
我就有了.
我本身就有了.
我不需要花什麼努力.
你從他剛才的回應就知道.
其實他是很輕佻的.
你聽就知道那些是誇張了很多的說話.
當他來到的時候問弟弟.
弟弟我快要死了.
給我喝點湯吧.
我快要死了.

$^{481}$然後弟弟就在那一刻.
趁你病拿你命.
那你賣個長子名分給我吧.
好像在開玩笑.
但是那個哥哥他怎麼樣.
我都快要死了.
那長子名分對我什麼用.
我只想大家留意.
究竟這番話.
甚至乎後來亞覺叫他去祈禱的時候.
你看到他那個反應.
他馬上就祈禱.
他馬上就發誓.
馬上發誓我要賣掉他.
他完全覺得沒事發生.
對他來說.
因為為什麼.
我們從之後所發生的故事.
當爸爸說要祝福他的時候.
他是很順理成章地接受的.
他完全覺得之前的事是沒有發生過的.
那個發誓是沒有效的.
但是弟兄姊妹.
我想你弄清楚.
不是這樣的.
今天其實你發覺.
亞覺在這裡很明顯.
他用了一個不應該的方法.
去搶奪這個長子名分.
第一件事.
當哥哥他在這裡很辛苦的時候.
他不是用一種憐憫.
他不是用一種憐惜去幫他的哥哥.
不是愛他的哥哥.
他在這裡的時候.
反而是趁機會.
看有沒有可以拿這個身分回來.
甚至乎要他發誓.
之後的時候.
他用盡一切騙的方法.

$^{521}$去騙這個回來.
但是你發覺整個故事.
到最後的時候.
沒錯.
亞覺是要承擔他自己的責任的.
他當時騙人.
所以你發覺後來的故事.
風回路轉.
他是被他的舅父騙回來.
他後來的時間.
是被他老婆騙回來.
整輩子被人騙.
他要承擔的.
要承擔後果.
但是上帝的確是揀選了他.
而很多學者都告訴你.
為什麼上帝會揀選他.
因為他渴慕上帝的祝福.
他和他的哥哥成為了.
一個很大的反差.
其實如果大家有細心去看.
之後他走去騙爸爸的故事.
其實是不容易會發生的.
你發覺他爸爸是一而再再而的.
去懷疑究竟這個是不是耳掃.
但是都過了骨.
其實後面那件事.
可以為他祝福到.
那個是完全是.
沒有可能發生的.
我想說.
因為最後他那句話.
明明你的聲音是啞角.
但是我身上又聞到他的味道.
又摸到他的毛.
是耳掃.
但是他都照樣祝福.
所以我想說.
不是上帝他自己參與在當中.
是不會發生了這件事.

$^{561}$所以弟兄姊妹.
耳掃他以為一些事情.
以為理所當然.
他以為一些事情是不會改變.
你搞錯了.
弟兄姊妹.
我們今天所信的這個信仰.
我想搞清楚.
上帝是不變的.
上帝的恩典是不變的.
但是你所得到的祝福.
是可以隨著你對上帝的態度.
而去改變.
你今天可以一個人.
繼續一路回來教會.
每個星期很恆常回來崇拜.
但是不等於你生命當中有平安.
不等於你已經得到了.
上帝祂所賜給你的福分.
你可以一樣的.
去為著這個世界裡面.
所爭取的東西.
而去勞累.
你亦都繼續可以為著.
人與人之間的關係.
你會不斷去受苦.
因為我們對信仰的態度.
並不是上帝想我們.
去對信仰的態度.
這麼認真和這麼好.
你會發覺.
以素去興翰上帝的主權.
他很輕佻.
剛才那句說話.
我快要死了.
大家知不知道.
其實有一個很諷刺的位置.
生死是誰決定的.
是上帝決定的.
他現在很輕佻.

$^{601}$我快要死了.
其實我想弟兄姊妹.
今天你也學一下.
不要這麼快就跟人說.
我快要死了.
不是你想死就死了.
是上帝容許你才死了.
但是弟兄姊妹.
他將.
我快要死了.
上帝這個獎之名分.
就是上帝的祝福.
還有什麼用.
你見到那種比重嗎.
所以他很輕佻.
發誓就發誓.
所以我們中國人有句.
很好說話叫.
發誓當食生菜.
你會發覺.
將上帝所給你的東西.
你當作兒戲.
但是你知不知道.
你信仰的上帝.
他是一個非常之.
看重細約的上帝.
你答應了他.
的東西.
他是一定會跟你.
追回來的.
鄧全道.
我今天五十多歲.
知不知道我是什麼時候.
決定奉獻給上帝.
我是在初中的時候.
當時我回荃灣宣道會.
我們宣道會.
每年都有一個夏令會.
夏令會.
我們那些是五日四夜的.

$^{641}$即是一個星期.
對於一個初中生.
進去夏令會.
每一段時間都是聽道.
還是聽成人的道.
一定是悶到抽筋.
然後我也不知道.
睡了多少次.
我都忘記了.
但我又不明白.
為什麼當時我去.
因為當時我的靈命不行.
其實當時我一個.
都很差的信徒.
雖然我回教會.
但是其實說真的一句.
表面上很乖.
導師叫我做什麼.
我就做什麼.
通常乖的.
是被人用的.
然後有時間.
我們叫職員會.
找我做職員.
我有一個很好的名字.
叫做靈修.
大家有沒有聽過.
有我差不多年紀的.
都聽過.
叫做靈修.
靈修做什麼的.
鼓勵人可以讀聖經.
但我是一個.
從來不讀聖經的人.
然後你知道.
最簡單的方法.
就是怎樣.
就是找些書.
每次截一句金句出來.
就拿出來.

$^{681}$我從來不看聖經.
從來不是一個好的狀態.
但是導師叫我.
進去下令會.
我就去.
為什麼呢.
下令會就是.
最棒的一件事.
大家知道是什麼嗎.
就是不用回家睡.
可以和同學.
和弟兄姊妹.
和小朋友一起玩.
下令會裡面.
有一個環節很重要.
就是最後一晚.
有一個獻心會.
獻心會有兩個心.
一個是簡單一點的.
就是弟兄姊妹.
我獻心.
我願意承諾上帝.
我要做一個.
好的教徒.
第二個心是什麼.
我願意全時間.
事奉你.
我忘記了那一晚.
究竟我說了什麼.
我完全忘記了.
你這樣問我.
我都不會記得.
但我只記得.
那一晚我哭得.
非常非常痛.
然後主席叫我.
獻心的時候.
我就十幾歲.
走了出去.
獻了一個.

$^{721}$我願意一世事奉你的.
的心出去.
弟兄姊妹.
我後來還離開了教會.
我後來還.
經歷了一段日子.
我到後來.
是二十幾歲.
二十幾歲中.
我才再回教會.
二十七,八歲.
我才回教會.
但弟兄姊妹.
上帝是不會忘記的.
你答應了祂.
所以弟兄姊妹.
你知不知道你決志.
當日.
你和上帝那個.
是一個納藥.
我們叫新約和舊約.
你願意信耶穌的時候.
是一個納藥.
你答應了什麼.
你答應了主.
我是罪人.
求你赦免我.
我願意將我整個生命.
我交給你.
求你成為我的救主.
我一生跟從你.
你答應了.
上帝不會忘記的.
上帝會用盡一切方法.
去令你達至這個教國.
所以弟兄姊妹.
大家記得.
我們都參加很多婚禮.
有婚約的.
今天在座的老婆.

$^{761}$在旁邊的.
還記得自己說過什麼嗎.
生死與共.
不會偏離.
富貴貧窮.
我們都一起走.
記得嗎.
你忘記了.
上帝記得嗎.
上帝很重視婚約.
很重視承諾.
特別有一個地方.
我們很清楚.
商業裡面.
我們做一些交易.
是要簽約的.
是嗎.
簽了約.
你違約的話.
是有後果的.
你可能要罰款.
你可能要坐牢.
所以弟兄姊妹.
我們今天其實.
我們不要將上帝.
跟他的關係.
跟他的承諾.
當做兒戲.
上帝跟我們很不同.
我們今天究竟.
用一種什麼態度.
去看我們自己的信仰.
是代表著.
你今後你自己.
整個生命是如何.
最近我多了.
因為隨著童工的調配.
我多了一些時間.
要回來幫忙.
我會幫忙在傳道外展部.

$^{801}$和栽培部.
跟著時間.
我跟弟兄姊妹.
跟這些執事.
跟他們一起坐下來.
聊聊天.
不斷地想這件事.
其實弟兄姊妹.
我們今天的狀態是怎樣.
我們去.
想弟兄姊妹成長.
究竟我們.
用什麼方法.
用什麼.
應該給她什麼.
讓她的生命.
能夠有效地成長.
我們今天去.
得到上帝的牧養之後.
我們有沒有.
將上帝給我們的恩典.
能夠走出去.
成為其他人的祝福呢.
這些都是很值得.
我們重新再思考的.
所以在這裡預告了.
如果所有的東西.
準備好之後.
到時候我會告訴大家.
接下來將會有一個.
叫信徒復興運動.
弟兄姊妹.
我們今天需要.
跟上帝.
履行上帝的承諾.
以素.
她因為對上帝的態度輕浮.
導致了.
她失落了原本.
應該擁有上帝的祝福.

$^{841}$我們今天對上帝的態度.
是怎樣呢.
我們有沒有對她.
回到最應該有的敬畏呢.
有沒有對她.
應該有的那份認真.
嚴謹呢.
而我們能夠.
真的在她面前的時候.
好好地.
在神面前的時候.
上帝問她.
我當如何去事奉你.
我當如何去過我的生活.
我當如何.
她能夠履行上帝.
我所對你的承諾呢.
求上帝讓我們能夠成為一個.
她真正的精兵.
願意從家庭的建立.
直至到我們整個生命的建立.
我們都能夠給建立起來.
然後就好像昨晚.
阿柱永華牧師.
他在這裡所分享的一樣.
我們真的很盼望有一天.
很盼望在上帝的眼中.
他會和其他人去介紹.
這是我的感受堂.
這是我的愛子.
這是我的兒女.
他能夠帶著榮耀地去介紹我們.
因為我們這間教會.
能夠成為他心裡面.
他所愛.
行出生命的教會.
好嗎.
求神幫助帶領我們.
很願意我們能夠.
也不要失落了上帝的祝福.

$^{881}$同心祈禱.
親愛的天父.
我們求你幫助帶領我們.
主恩我求主你去引領我們.
主恩我們在跟隨你的過程當中.
或許我們都經歷了很多的事.
或許我們的心裡面.
或許有時都封閉了.
我們都不敢再走多一步.
或者不想再走多一步.
主恩我求主你親自幫助我們.
主恩我們生命當中.
需要被復興起來.
求你叫我們在生命當中.
不要去看得理所當然.
因為唯有當我們和你一起同工.
和你一起去走的時候.
我們才知道真正的幸福.
真正的祝福.
能夠臨到我們身邊.
我們才能夠每一天.
我們能夠去經歷到主你自己的同在.
和你的恩典和你的愛.
所以我求主你去保守我們這裡每個弟兄姊妹.
願意我們的心向你呼求.
向你禱告.
去向你求主.
我們當如何的去生活.
如何的與你同行.
我們獻上感恩.
多謝讚美你.
其樂逢教主耶穌.
得勝名至其球.
好.
謝謝.
\newpage



\section{瑪拉基書}
\label{sec:exEZqq1nt80}
\textbf{2024-02-29 周四晚查經班《瑪拉基書(8)》 講員:曾裔貴牧師}
\newline
\newline
連結: \href{https://youtube.com/watch?v=exEZqq1nt80}{\texttt{https://youtube.com/watch?v=exEZqq1nt80}} ~~~~ 語音日期: 2024-03-02
\newline
\newline
\hyperref[sec:__QTchIrjDM]{\small{< < < PREV SERMON < < <}}
~
\hyperref[sec:index_chronic]{\small{[返順時目]}}
~
\hyperref[sec:index_scriptual]{\small{[返順卷目]}}
~
\hyperref[sec:oyj7tRcaGj0]{\small{> > > NEXT SERMON > > >}}
\newline
\newline
$^{1}$多謝天父再一次讓我們能夠在晚上.
雖然開始有點寒冷了.
但是我們都可以安坐家中.
還有在禮拜堂這裡.
來去查考《馬拉基書》.
很盼望一個比較短 還有相當長.
已經很多年的一些舊約的經卷.
我們可以掌握裡面的教導和那個精髓.
使得我們能夠在我們自己新約的時代.
耶穌基督還沒有回來之前.
我們能夠在教會的生活.
特別是那個事奉的態度.
敬拜的心靈 能夠蒙主的喜悅.
願你帶領我們大家的心.
更加讓我們能夠看到.
其實有很多的地方是沒有福音的.
很盼望神自己來到去工作.
興起弟兄姊妹在不同的教會.
不同的地方一起為神的福音去效力.
我們這樣仰望禱告.
奉耶穌基督的名 阿們.
可能我們用這次和下一次.
就完了這個《馬拉基書》.
就可以回去看新約.
還有因為關乎到末世.
都會仔細地看看末世論.
有不同的神學主張.
那個就不是絕對的.
但是在這個《帖撒萊尼加前後書》.
就是看聖保羅在早期傳道.
很明顯是第二次他的傳道旅遊報導的時候.
那個事蹟和那個對教會.
一個受苦難的教會.
怎樣來勸勉他們.
很特別就是用主在來去勸勉他們.
這個我們下幾次開始講新約的時候.
再詳細講.
我們第三章的《馬拉基書》.
上一個星期就給了一個問題.
就是到底公義的神是不是真的存在呢?.

$^{41}$是第三章的收尾.
如果公義的神不彰顯他的公義.
在他的子民 當然舊約就一定是耶路撒冷.
猶太選民和以色列國.
如果不彰顯的時候.
你敬拜是徒然的.
你忠心 嚴謹去跟隨.
都是沒有什麼好結果的.
不專心 很敷衍.
都沒有什麼事情發生的.
可能從那個經文的背景就告訴我們.
普遍在百姓的心裡面.
就有這個內心的鬱結.
所以他們問一個問題.
你講到自己這麼偉大.
都不見有什麼審判.
不見有什麼特別的東西.
只是說哪些不行 哪些不行.
特別是那個用鬼詐第二章出現很多次.
猶太人自己彼此用鬼詐.
領袖尤其是祭司向神鬼詐.
接著民間的百姓彼此鬼詐.
接著那些丈夫對自己的妻子鬼詐.
就變成了充斥著在社會裡面.
大家都不敢去講真心話.
大家都不能夠用真心.
用一個坦誠的心去彼此對待.
就出現了似乎就是內心有一個鬱結的問題.
就是公義的神在哪裡呢?.
這個是二章的16節的問題.
所以第17節 對不起 應該是16,17兩節.
或者公義的神到底在哪裡呢?.
他們不是不知道在耶路撒冷.
他們不是不知道在藥櫃裡面.
或者在致聖所.
但是政治就不在他們的心.
或者他們所見到的實際的生活處境.
不是很感受到神的同在.
神施行的那個公義.
不同當初他們所羅巴伯上來.

$^{81}$帶著百姓回去重建聖殿的時候.
怎樣能夠彰顯那個榮耀.
將那些本地人 即迦南的那些本地人.
透過很多的教訓.
將他們去自服了.
使得他們終於很戲劇性的可以重建聖殿.
接著有以色列的復興.
有後來的再過了100年.
120多年的尼希米的那個再重建城牆.
所以一切都安頓好之後.
也都似乎是馬拉基在一個最後的階段.
可能比尼希米稍微遲一點點.
不是遲很多的.
到底他們兩位有沒有曾經.
同個時代一起這樣服侍呢?.
聖經沒有交代 不過後人.
一些聖經學者就有這樣的認為.
他們兩個人可能是同那個時代背景有關的.
所以很多東西安頓了之後.
慢慢沒有什麼大的衝擊.
就使得猶太人回復在那個日常生活裡面.
敬拜 事奉 獻祭 或者日常的買賣.
人與人之間的交往等等.
慢慢就回到了一種很平淡的狀態.
所以挑起他們內心有很多的鬱結問題.
這個是從第二章給我們看到.
慢慢就開始可能對自己的保護機制要啟動了.
就有很多的鬼咋彼此相待.
這個是很難去拆解的.
因為安全感不是你說兩句就可以有的.
人越大就越要保護自己的.
那個安全感就越難去.
好像一個保護罩那樣 甚至乎是無形的.
不知不覺就會有一種的自我保護.
所以就一定是在民間就會慢慢又受到傷害.
所以第二章的鬼咋出現很多次.
所以最終他們將這件事化成一個問題.
公義的神到底是否存在呢?.
存在的 他們知道的.
不然他們都不會獻祭.

$^{121}$他們都不會真的這樣透過祭司去獻祭.
不過很多東西都慢慢習以為常了.
這個就去到第三章.
就有一個第二章是負面的.
講他們的那個狀態.
第三章就是講那個理想.
那個應該是這樣的敬拜 神才會喜悅的.
尤其三章十六節有一本紀念冊.
在神那裡給神紀念的.
不知你小學有沒有六年級升中士之後.
有沒有寫紀念冊呢?.
不知現在的小朋友有沒有呢?.
現在可能根本電子筆就不用了 用電話.
小時候我們真的什麼都要寫的.
那本東西很珍貴的 但是很快就不見了.
所以都有些可惜的.
但是十六節有些真的紀念冊.
在神的面前紀念冊.
所以你可以想像到.
第二章的那個不尋常的敬拜.
第三章的就是神喜悅的敬拜.
是怎麼樣的呢?.
那就很寶貴.
首先就是講那個審判.
審判完之後 神就期望.
要回到那種沒有自我保護.
要在神的裡面去學習.
真是開放自己的敬拜.
所以第三章很寶貴的.
值得我們一起去看.
如果第二章是你已經因為這麼多的詭詐.
而導致你有一個自我防護機制.
那試一下將這個制按一按.
取消它 第三章就帶你進入.
神的那個管治的世界裡面.
其實就是指到主在內.
有一位公義的聖者會出現.
立約的使者忽然會進入他的殿.
第一章就已經在講彌賽亞的降臨.
甚至馬可福音第五章就引述了.

$^{161}$那個立約的使者的那個僕人.
就是施洗約翰 是耶穌講的.
所以耶穌就解釋這段經文.
第三章開始了.
第一節.
「萬君子耶和說.
我要差遣我的使者.
在我面前預備道路.
你們所尋求的主必忽然進入他的殿.
立約的使者」.
這裡就是講施洗約翰了.
「就是你們所仰慕的快要來到.
他來的日子誰能當得起呢?.
他顯現的時候誰能立得住呢?.
因為他如煉金之人的火」.
煉金就是那些冶金.
要用很高溫的火.
將那些雜質燒紅之後.
那些雜質自然會浮起.
就會走的.
那些金再倒下去就變成四條九的金條.
這些是注金.
這些古代的人一定會懂的.
「他如煉金之人的火.
如漂泡之人的皂」.
就是洗作的那些布.
我做過染布的.
暑期工.
很辛苦的.
將那些布染色.
很重很臭.
化學的手指甲都長出了什麼.
因為那些化學劑很厲害的.
真的要用皂來洗的.
它要漂白的.
所以這句話「如漂泡之人的皂」.
我們小時候.
我不是說過我住在九龍塘的.
很多染布廠.
因為那裡那區.

$^{201}$很多地方.
因為那些布.
小時候.
他們不是用那些乾衣機去搖的.
有時候要在戶外曬的.
因為天然的.
太陽能的.
成本低很多.
所以很多的染布廠.
在我們小西湖村的.
也有豬油廠.
牙肘廠.
也有些家庭式的養豬.
譬如我家就是了.
有些種菜等等.
好像很多戶的.
幾千戶的.
很長的那個模區.
由北架山一直延伸到豬鼓仔.
差不多是大亨東的.
所以一個很大的地方.
所以那些染布.
就要曬那些布.
真的要染.
真的要洗白的.
所以那個漂布很重要.
這裡告訴我們.
要麼說煉金.
就是煉靜.
要麼就說洗到雪白.
有沒有想起44章的《二創亞書》.
有歌手唱的.
《廿二節》.
你當歸向我.
接著可以洗到你.
來到變成雪白.
有這些兒歌這樣唱的.
不過這裡告訴我們.
這個人的出現.
會洗到我們好像煉金之人的火.

$^{241}$如漂布之人的繭.
他必坐下如煉靜銀紙的.
必結證利未人.
利未人我們第二章不是說過.
他們帶頭去違背律法的.
祭司的嘴裡面應該傳知識.
人應該從他的口裡面.
來得到耶和華的教訓的.
現在第三章神開始做事了.
你說公義的神不在.
神現在首先要煉證這些利未人.
有份事奉的僕人.
不同崗位.
神會有那個煉證.
所以我們事奉.
一方面我們都是很謹慎.
事奉其中一個態度就是.
神是很想事奉的人.
能夠代表他.
不是隨隨便便的.
不是隨便來的.
而是真的能夠有一個結證的心靈.
所以這裡說結證利未人.
而他們就憐憫他們.
好像匠金銀一樣.
他們就憑公義.
獻供物給耶和華.
剛剛上文就說.
他們獻的那些.
百姓獻什麼他們就收什麼.
就是無所謂.
無要求.
百姓說這輪不行.
你不要這麼要求高.
那隻羊有一點瑕疵.
你不要這麼挑剔.
你收了它吧.
沒有人知道的了.
可能就塞錢給祭司.
你接受吧.

$^{281}$祭司說很久沒有人來獻祭.
有人來獻祭你當然.
當然什麼呢.
就好像你想想.
年三十晚王大仙這樣.
今年沒有人來獻.
沒有人來裝香這樣.
很嚴重的.
接著有人來.
你覺得.
你當他上賓.
你可以想像得到.
這是一個此消彼長.
如果百姓都不上來獻祭.
那祭司還這麼高要求.
有時候就像一個惡性循環.
所以越是我們港台不尊重.
或者求求其其.
百姓就更加求求其其了.
反過來如果我們很嚴謹.
嚴謹不代表沒有愛.
嚴謹每一樣東西都很有秩序.
人家就會看到你們很認真對神.
人家也會很認真對待敬拜,崇拜.
和那個心都會預備好.
所以這個是很重要的.
你可以想像得到.
如果你看到台上的講完.
那個祝福的主席.
在那裡晚堂不停地打陷路.
不停地打黑眼睛.
有沒有搞錯 一直講到.
那個祝福的不停地打黑眼睛.
在那裡釣魚.
台下的人怎麼想.
你就會看到其實是不行的.
所以我經常都要捏自己.
不可以的.
有時候做完事情很累.
但是你都不停地捏自己.

$^{321}$不停地搖頭.
有時候眨眼很累.
但是很快就要在那個狀態.
所以你可以看到.
神有工作.
神說祂要憐靜這些利未人.
他們就獻公義的祭.
就是蒙神月立.
上文第二章就說蒙神月立.
但是現在第三章就是.
這個使者出現.
當然這個是預言.
四百年之後耶穌才來.
彌賽亞和他的使者.
就是施舍爾汗才來.
但是這裡告訴我們.
祂的出現就會有這種憐靜.
所以我們信耶穌那一刻.
其實我們已經得到生命的拯救.
就是我們的靈魂救恩.
就是我們聚得到赦免.
我又沒有說過.
我初信的時候.
一信耶穌.
突然間整個人.
我晚上.
我不是說過在沙田.
拿完那個崇禎小學.
其實不關我的事.
我又不是讀小學.
是一個家長晚會.
是12月24日晚.
1987年我很清楚那個日子.
是我平生第一次.
都不是去教堂.
去那間學校的天台的禮堂.
來聽一些家長的講座.
其實是福音講座.
整個聚會是報道會.
信了主之後回到家裡.

$^{361}$有奇事發生.
突然間回到家裡.
很喜歡讀聖經.
但是接著第二天.
其實我之前是完全沒有興趣.
完全不知道說什麼.
但是我回到家裡.
其實一個多小時後就完了.
回到家裡.
就很喜歡讀聖經.
一直就開始自己讀聖經.
接著整個人很輕散.
接著試過有一次.
開始.
接著我回到元朗的一個教會聚會.
有一次聽完崇拜之後.
第二天禮拜一.
在家裡.
那時候我已經買了一本詩集.
在家裡唱那些詩歌.
開始靈修.
信了主都沒有兩個禮拜.
開始自己讀經靈修.
唱第一首詩歌.
是昨天的崇拜.
聽到覺得很好.
在家裡自己唱靈修的時候.
早上在自己的房裡靈修.
那首聖詩叫做.
「眾主聖歌」.
178首.
我還記得這麼深刻.
因為這些太久太久了.
何等稀奇一位救主.
他的題目叫做「奇妙救主」.
他說「何等稀奇」.
我經常走音.
「奇妙救主」.
我所沉著的就是主耶穌.
接著一直唱.

$^{401}$唱著唱著唱到第二節.
不停地流眼淚.
流得很厲害.
接著一直在裡面.
好像放映機放映我的罪.
逐件自己在神面前向神認罪.
接著很害怕.
很快到星期六就回團契了.
回去馬上跟傳道人說.
為什麼我會這樣的呢.
傳道人說.
這就是你的重生經歷.
你在神面前.
你自己神透過一首詩歌.
一節的詩歌.
的說話.
就來給我們看到.
給你看到你的罪.
你就向神認罪.
很輕省.
由那時候開始.
很喜歡快點到星期六.
回團契.
聽講道.
和自己平時讀聖經.
一路如是者.
就是這樣.
我自己的初信階段.
那時候我開始回夜校.
我都講過了.
開始在夜校裡面.
見到人就作見證.
有一個女孩子.
有一個人走過來.
她也是讀夜校的.
做不知道什麼的.
很有興趣問我信什麼.
為什麼你會這麼疑心.
原來她小時候也有回教會.
但是沒有再回了.

$^{441}$輾轉跟我辯論.
在下課的時候.
小息跟我辯論.
分享.
她都沒有信.
過了很多年之後.
透過我的教會.
一位弟兄做郵差的.
因為經常去他那個.
屯門的辦事處.
認識這個女孩子.
就是我那個夜校同學.
就說原來她信了耶穌.
等等.
所以很奇妙的事情.
這個就是我初信的階段.
所以立約的使者.
首先會在我們每一個人的心靈.
所以回教會.
你必須首先有你自己個人的.
信主經歷.
因為這個出現.
很明顯就是念靜.
很明顯是救恩帶來生命的改變.
不同的方式.
不同的說話.
有人說三章叫做重生.
人如果不是從水和聖靈而生.
就不能進入神的國.
這個是耶穌自己跟尼哥底姆說的.
也有另外的經文.
譬如你口裡認耶穌為主.
心裡信神.
叫他耶穌從死裡復活.
就是你口認心信.
他的降生和復活.
這個是一個很大的能力.
就必定得救了.
你只要這樣認信.
你必定會得救的.

$^{481}$就是說整個屬靈生命是為你而開.
是為你的生命改變.
最近有時候晚上我睡得不好.
就會聽新聞.
他YouTube自己不停的播.
聽著聽著聽到一個.
我不說是誰了.
免得你聽不懂.
很出名的一個律師.
女律師.
很出名的.
可能你猜到了.
她說她經歷見鬼.
她說真到一個地步.
是她沒法想像的.
我想她未信耶穌.
她信佛那些.
但是如果鬼尚且有這樣的奇異經歷.
說得很真.
似乎覺得是真的.
但是我們信神的人.
就更加要看到我們自己的生命改變.
所以每個人都不同.
不是每個人都要像我這樣.
哭得瘋了.
不是每個人都一樣.
有些人可能很開心很平安的信主.
但是我們需要看到這一段經文.
主耶穌的出現.
是一定要將他自己的子民.
帶進聖潔榮耀的國度.
而且這個不是將來才帶你去.
是現在由你信那一刻.
你已經開始有一種感應.
對聖靈的敏銳的感應.
這個是很肯定的經文.
所以舊約都已經有這種重生.
帶來生命改變.
這裡告訴我們.
那時猶大第四節.

$^{521}$那時猶大和耶路撒冷所獻的貢物.
必蒙耶和華月立.
剛剛上文說.
我不月立你們.
你們所獻的我看為可恥.
可憎.
現在這裡主耶穌的出現.
救恩彌賽亞降臨.
就帶來整個人類的改變.
所以是一個劃時代.
所以難怪保羅真的不明白.
他們做好了.
但是為什麼不能夠蒙神月立呢?.
因為其實如果沒有主耶穌的降生.
其實帶不了人的生命改變.
我們信一個宗教.
我們信一個禮儀.
信一個傳統.
是救不了我們的.
是耶穌才救了我們.
是祂的降生.
是祂的生命.
是祂的流血.
而且是要祂復活.
才可以拯救我們.
所以這裡就很寶貴.
神再次月立.
那就是說原來神的憤怒.
只不過是一息之久.
祂的恩典是直到永恆.
所以猶太人雖然這麼敵黨神.
這樣鬼咋對神.
這就好像我們第十節.
二章十節所說.
為什麼我們要用鬼咋去對待神呢?.
這裡現在神就解答了.
誰說我沒有公義.
我的愛子降生.
他的犧牲.
他的死.

$^{561}$就是神的公義的彰顯.
所以羅馬書很強調.
其實不是神的愛.
神的愛是去到第八章才說的.
其實一開始.
就是在說我們虧缺了神的榮耀.
我們現在因信基督耶穌.
就是神的義.
加給一切不義的人.
普羅一開始就說.
很嚴肅的.
就是神的義.
加給不義的人.
所以我們獲得的不是自己的義.
不是我們自己做好我們自己.
我們去修補我們自己.
不是.
是神的義給我們.
所以我們就能夠將來去到神的國.
因為神的義都已經算了我們.
我們就根本足夠資格去天堂.
所以不是修補.
不是靠自己的功德.
是靠神的愛子.
神的義的運作就是他的愛子的犧牲.
是這位無罪的替我們成為罪.
所以我們就能夠在神的面前.
成為神的義.
哥林多後書五章應該是二十一節的經文.
是很重要的教導.
所以這裡告訴我們.
回到一個狀態就是.
往不古時之日.
尚古之年.
不如我讀完才解釋.
不然我怕大家又.
第五節.
萬君之耶和華說.
我必臨近你們.
施行審判.

$^{601}$我必速速作見證.
警戒行邪術的.
辦奸淫的.
起假誓的.
欺負人之公家的.
欺壓寡婦孤兒的.
屈往寄居的.
和不敬畏我的.
你可以想像到.
一定是內心的道德水平很低.
神要提升水平.
所以這裡是說道德性的東西.
原來那把口說的話.
作假的見證.
屈往正直.
屈往那些寄居的.
虧負人之公家.
有時說了不算數.
各樣的形式.
虧負.
對別人有虧欠.
都是一個不好的.
做錯事.
我們人類最難說的就是對不起.
最難說的.
那個面子死要面子.
很難說的.
如果一個人去到一個人生階段.
能夠說對不起.
認真說.
不是隨便說.
亂說都沒有用.
那些叫游康或調.
那些人更加討厭.
真誠的.
看到自己的問題.
說對不起.
尋求別人的寬恕.
這個就不是虧欠別人了.
欺壓寡婦孤兒.

$^{641}$欺壓即是說.
你是有權的.
你可能用到盡.
明仔他的上司.
我不說太詳細了.
他有關教會的.
但是他一出街.
就好像當自己是一個公義的審判者.
看到有人吸煙.
人家在外面吸煙.
煙進來之後.
他都很認真地罵人.
投訴給商場等等.
他自己又衝紅燈.
他自己有時又做事.
人家的侍應送東西來.
送得慢一些又罵人.
自己就以為很公義.
公正.
但實際上兩把尺.
這些就是欺壓孤兒寡婦.
所以一個敬虔的人.
經過神的聖靈工作.
某程度上是一個靈魂的手術.
我們來到他被神改變.
應該對神對人.
都有一份聖靈所結的果子.
人愛喜樂和平.
忍耐恩慈良善.
信實溫柔.
溫柔不是陰聲細氣那些.
那些什麼那些什麼.
那型話 我沒有貶義的.
信實溫柔節制.
第一樣就是人愛.
我想對別人的一個包容.
一份的愛.
所以你看到這裡是很相輔相成的.
誰說神的公義不出現.
誰說公義的神去了哪裡.

$^{681}$神就告訴你.
我立即.
不過不是現在.
是四百年後的彌賽亞降臨.
他的降臨就是煉金人的火.
漂泊人的鹼.
神也有一些像醬和肝.
來帶領羊群.
所以尤其是煉淨利未人.
原來我們越是事奉.
越是蒙福.
因為越是事奉.
就越能夠近神.
你不事奉.
你就不會對自己有要求.
對自己有要求.
其實是一件美麗的事.
神就會有很多的祝福.
你會有很多經歷.
所以這裡是告訴我們.
我們要這樣學習.
第六節.
因我耶和華是不改變的.
所以你們雅各之子沒有滅亡.
果然就是這件事.
他的憤怒是一息之間.
一叔又哭泣.
到早晨就歡呼了.
我忘記了是二十七篇詩篇.
要回去找一找.
第七節.
萬君子耶和華說.
從你們列祖的日子以來.
你們常常偏離我的典章而不遵守.
現在你們要轉向我.
我就轉向你們.
你們卻問說.
我們如何才是轉向呢?.
你們可以看到那個心路歷程.
神已經說了.

$^{721}$因為祂的不變.
不是因為人的善變.
不是因為百姓的那種.
習慣性的那種憲制.
敬拜的隨便.
來到去.
神就好像跟他們一樣的見識.
所以一般人就是一般見識.
我都說過了.
屬靈人就不是一般見識了.
屬靈人在事物.
不要只看表面.
要看更加高的.
所以你看神是怎樣.
神更加不單止是屬靈人.
祂本身就是靈魂的神.
祂本身就是造物主.
所以是因為我是不變.
所以你們雅各的子孫.
才能夠留到今天.
祂才能夠容忍百姓的那種背影.
那種詭詐.
那種對事奉.
對敬拜的那種心淡.
你看見一個人完全敷衍來崇拜.
如果你有摩西的憤怒.
你真的想下台.
用法板打他的頭.
打醒他.
他好像完全不在乎自己崇拜什麼態度.
只在乎整個人躺下來.
去躺下來.
只在乎這樣.
不敢失禮別人.
但實際上他可能在看手機.
看看什麼新聞.
你來到教會還看新聞.
浪費自己的時間.
不如在家裡看吧.
家裡的螢幕更加大.

$^{761}$為什麼要拿手機來看呢.
證明他很敷衍對神.
其實不是對人.
你搞錯了.
其實不關我們傳道人的事.
是你敷衍對神.
你來敬拜是對神.
怎麼會對那個人呢.
所以這些是一般人一般見識.
肅靈人就不是這樣.
肅靈人就看到.
其實神就是有寬容.
神就是有萬般的忍耐.
所以你們雅各之子.
未有滅亡.
所以這是一個很大的祝福.
所以無論你裝模作樣也好.
你真心也好.
我們來到這個階段.
能夠有一點點的事奉.
是無盡的感恩.
感恩.
你隨時可以連講話都講不到的.
你隨時下一秒鐘.
可能就真的完全走都走不到的.
甚至你可能突然間就離開世界.
沒有人可以預測的.
但是有一些的事奉.
你帶著感恩的心去做.
做在神裡面.
這個就是事奉態度.
做在神的裡面.
耶和華是看著的.
你就可以想像得到.
你就可以抗衡那種.
在教會裡面的那種冷淡.
那個冷淡比冷氣還要冷.
你知不知道.
你有沒有感受過一個很冷淡的崇拜.
那種氣氛.

$^{801}$那種氣氛是很恐怖的.
大家都快點走.
很睏的那樣.
但是一個屬靈的聚會.
你感受到神同在.
你是不想離開的.
所以我們求主讓我們能夠看到.
原來這一段聖經告訴我們.
原來是神的憐憫猶太人.
神不是看不到他們的敬拜.
是這麼敷衍的事奉.
但是神就仍然給他們.
有一個恩典.
所以神就說了.
你們要轉向我.
我就轉向你們.
其實是神先轉向的.
所以百姓開始有一個這樣的學習.
他們就問了.
怎樣轉向呢.
接著就說錢了.
第八節.
所以有時候解這些經文.
上下文你解讀.
你就會明白.
不是這樣抽一節出來.
抽一節出來就沒有用了.
你抽了一節金句出來.
你不在上下文的處境.
那個張力底下.
那節金句就變成不是金句了.
第八節.
一直都在說.
百姓如果能夠明白神的心.
神對這個民族的那種不離不棄的愛.
和不斷的恩典挽回.
和神對百姓的轉向.
百姓現在就好像在問.
怎樣轉向呢.
第八節.

$^{841}$人豈可奪取神之物呢.
你們竟奪取我的供物.
你們卻說.
我們在何事上奪取你的供物呢.
就是你們在當立的十分之一.
和當現的供物上.
因你們通國的人.
都奪取我的供物.
就坐就臨到你們身上.
萬君之耶和華說.
你們要將當立的十分之一.
全然送入倉庫.
使我家有糧.
以此試試我.
是否為你們闖開天上的窗戶.
傾福與你們.
甚至無處可容.
萬君之耶和華說.
我必為你們斥責王松.
不容他毀壞你們的土產.
你們田間的葡萄樹.
在未熟之先也不掉果子.
未熟就不會自己爛.
不會有蟲咬.
或者生其他傳染的病菌.
令那些樹出現毛病.
即是說熟才會掉下來.
即是說有好的收成.
第十二節.
萬君之耶和華說.
萬國必稱你們為有福的.
因你們的地必成為喜樂之地.
喜樂這個字不是happy.
叫做Hepzibah是神所喜悅的意思.
Hepzibah是希伯倫文的字.
這個字為何我會認識呢?.
Hepzibah就是我的母會喜樂福音堂.
不是一個開心的意思.
Hepzibah是神所喜悅的意思.
這是二創亞書.

$^{881}$第49章的神所喜悅.
這裡很寶貴.
有一個argument.
首先就是神期望百姓知道.
他們未滅絕.
是因為神的珠般的恩典.
一直長久的對他們的忍耐.
對他們不離不棄的永生的愛.
所以神期望百姓明白了.
要回轉歸向神.
所以這個return其實就是心意更新.
所以羅馬書十二章.
這個心意更新是return.
所以我們要調頭.
我們的內心.
所以我們心靈不斷要由遠回轉歸向神.
所以這個回轉是神對百姓的呼籲.
所以這個呼籲之後.
他們就說怎樣回轉.
怎樣轉向.
我們說錢.
很具體就是.
本來金錢是很中性的.
本來金錢沒有什麼特別的東西.
但是金錢有它的魔力.
就是它會帶來一定的安全感.
錢會帶來安全感.
錢會帶來你可以駕馭一些事情.
所以萬萬不知覺.
錢就成為了神和百姓中間的一個媒介.
所以有沒有記得耶穌說.
我們不能夠又事奉神.
又事奉黃大仙.
不是的 又事奉馬門.
財利.
將神和財利掛勾.
等同.
那你可想而知.
原來古代,現代.
錢無論用什麼形式出現.

$^{921}$都是會帶來安全感.
我們不敬拜神.
就是慢慢失去了安全感.
所以我們要找馬門做我們的安全感.
做我們的神.
所以本來錢是中性的.
記不記得提摩太前書都說過.
六章.
錢是萬惡之根.
是不是.
六章十一節.
當然不是.
貪財才是萬惡之根.
不是錢是萬惡之根.
所以保羅是很清楚告訴我們.
錢是很中性的.
你怎樣看這件事.
貪錢的人才是萬惡之根.
在我們內心裡.
就是那個安全感要透過.
不斷要多些錢.
進入貪婪裡.
所以原來我們很多人看錯經.
以為錢是萬惡的.
其實不是.
完全不是.
清教徒看錢是看得很中性.
是很明顯的.
所以我們不要看錯了.
以為錢是錯誤的東西.
提摩太前書六章.
大家不熟悉嗎.
這些是很重要的金句.
提摩太前書的六章.
其實教提摩太.
去做傳道都要很小心.
六章這樣說.
六章十和十一節.
其實整段聖經.
就是說一個好的牧者.

$^{961}$和一個不好的牧者.
是一個對比.
第六章第三節.
你要以此教訓人.
第三節.
如果有人傳異教.
不服從我們主耶穌基督純正的話語.
和合乎敬虔的道理.
他就是自高自大.
一無所知.
專去問難.
奇奇怪怪的東西.
在那裡辯駁.
問難就不是.
問錢.
問難就不是.
問題解答.
是特意挑些東西來問.
爭辯言詞.
從此生出嫉妒.
紛爭,毀謗.
妄疑.
那些壞了心術.
失喪真理之人的憎勁.
他們以敬虔.
為得利的門路.
開始說得利.
第六節.
然而一個敬虔的牧者.
敬虔加上知足的心.
敬虔加知足的心.
下面說.
剛剛對比.
那些不好的牧者.
不好的心態的人.
貪念.
貪和敬虔知足是一個對比.
所以這裡告訴我們.
敬虔加上知足的心.
就是大利.

$^{1001}$大吉大利.
我們都會念些輝春.
保羅都說大利.
哇.
多少多少的利益都不夠這東西好.
就是你已經有一個.
對神的敬虔.
對自己的知足.
你家裡有多少雙鞋.
回家今晚數一數.
下星期搭一搭.
知不知足.
正常啦.
有一雙球鞋踢球.
有一雙單車鞋踩單車.
接著有一雙皮鞋.
禮拜崇拜.
沒理由你傳統人穿球鞋.
是吧.
可以的.
仍然傳統教會當然穿皮鞋好些.
是吧.
行山給你多一雙行山鞋.
是吧.
如果你要跑步多一雙跑步鞋.
因為那些功能性的鞋.
皮鞋當你有一雙黑色一雙啡色.
接著.
一雙夏天一雙冬天.
冬天又爭的.
數一數.
每樣功能一雙就夠了.
都五六雙了.
如果你在家超過十雙.
首先問一問.
自己是否知足的心.
你想想.
傳統人穿恤衫.
三四件沒問題.
夏天熱替換.

$^{1041}$三十件.
在你的櫃子裡.
又粉紅色又藍色又白色.
白色都超過十件.
是不是知足的心呢.
自己去衡量.
我們昨晚上的靈恩運動.
靈恩的傳道人.
他們推廣的.
就是一個叫.
Wealth Health Ministry.
Wealth 即是財富.
Health 即是健康.
神是這麼豐富.
神的僕人怎麼會這麼寒酸呢.
所以神的僕人是要開好車的.
住豪宅的.
帶勞力士的.
美國那些.
香港就沒有了.
所以.
敬虔加上知足的心.
保羅說.
提摩太.
可能你現在有缺乏.
有比較.
但是要守住這個心.
守住這個敬虔和知足.
有不同的人.
在不同的背景.
神安排不同的人生.
祂的知足和你的知足.
可能不同的.
但是.
你如果提摩太是傳道人.
你要記住敬虔加知足.
兩樣東西沒有抵觸.
而且是平衡的.
敬虔又不知足.
那就很辛苦了.

$^{1081}$你又很想學別人這樣.
現在電動車Tesla又說可以再延遲兩年.
我的車都沒有Tesla.
Tesla這麼漂亮.
最好就是不用油.
但可能你不是.
你是沒有需要的.
你是沒有了知足的心.
其實不要用這樣的理由去賄賂自己.
這就告訴我們.
第七節.
留意一下那個原因.
那個原因就是.
普世都適合用的.
因為我們.
不只是行經傳道人.
我們沒有帶什麼來這個世界.
也都不能夠帶走什麼.
也都不能夠帶什麼去那個地方.
所以這個就是敬虔加知足.
就是今生今世.
那個最重要的心靈的安全感.
那個座右銘.
那就很重要了.
只要有衣有食.
就當知足了.
所以真是有衣有食.
知足就是說常常感恩.
所以剛才不是說.
我們能夠有一點點的事奉.
你應該是.
那你就會帶來無盡的平安.
不過留意.
第九節又再講.
那個對比了.
那些假的.
不忠心的牧者.
但那些想要發財的人.
就陷在迷惑.
落在網羅.

$^{1121}$和很多無知有害的私慾裡面.
因為他內心的慾望.
被錢.
被想擁有的東西.
來遮蔽了.
那個心不能敬虔.
不能知足.
是不是一個對比.
保羅原來有一個.
可以照到內心的內視鏡.
比那些MRI還要厲害.
一兩句話.
已經將人心的心靈.
剖析出來.
所以人生在世.
不是貧窮富有.
每個人都有不同.
貧窮的人都可以很貪婪.
富有的人都可以很知足.
兩者是不會有抵觸的.
抵觸的就是我們那個心靈.
必須要有敬虔加上知足的心.
留意這裡告訴我們.
第十節.
這節真的是金句.
所以不要讀錯聖經.
金錢不是萬惡之根.
貪財是萬惡之根.
有人貪戀錢財.
就被引誘.
離開了真道.
而且用很多的仇苦.
把自己來刺透.
就是說.
貪戀錢財.
自然就會帶來.
那個心就不能滿足.
就不能平衡.
就會帶來很多.
無知有害的私慾.

$^{1161}$在心靈裡面.
來徘徊.
所以我們回來看這裡的.
《馬拉基書》.
很寶貴的地方.
就是神都要.
我們留心.
原來奉獻.
一個有敬虔的.
心靈的人.
神要我們回轉.
我們看到.
原來我們可以用我們的金錢.
神給我們所擁有的東西.
我們去奉獻在神的國度.
聖功.
將來是會蒙紀念.
是會有賞賜.
和今生.
原來這個就是最好對付.
我們的私慾的行動.
就是讓我們.
能夠保持敬虔.
加上知足.
原來就是奉獻.
所以有時我們講這段經文.
害怕醜醜的.
不敢告訴別人.
要奉獻十分之一.
其實不需要的.
整段去解明它.
好像沒有了那種張力.
帶來的份量.
整段就是要我們回轉.
怎樣回轉呢.
百姓這樣問.
然後神就說.
你們搶奪了我的國度.
應該得到的.
難道神真的稀罕那些牛羊嗎.

$^{1201}$難道神做物主.
祂要來不可以.
創造萬有嗎.
你想想現在我們發現那些稀有的金屬.
要用很多的辦法.
轉地底.
轉到很深才能挖出來.
如果你試試挖到一顆.
很天然的藍寶石.
是很多卡的.
你都發現原來價值連城.
這些不過是自然界的壓縮.
幾千萬年的結晶.
不過是玻璃石頭.
都是做物主.
讓我們看到.
祂的奇妙.
你想想鯨魚在海洋裡面.
多大隻.
可以想像到.
你見到的自然界各種的奇妙.
都是做物主給我們的.
所以祂既然賦予我們人類.
有這個媒介.
是用錢.
這個交易的平台.
的一個媒介.
來運作一個商業的模式.
祂就希望我們能夠.
用金錢.
來回饋神的國度.
再加上我們新約.
就是福音的.
新的聖公的需要.
所以一個教會.
如果每個月都要.
每個月都要籌款.
每個月都要求.
要問 要借.
不好意思.

$^{1241}$我們下個月又要借銀行的.
高息的貸款.
為什麼呢?.
因為沒有人奉獻.
那即是怎樣呢?.
即是我們完全不明白.
我們沒有.
將一份敬虔.
加上知足的心去教導弟兄姊妹.
讓弟兄姊妹看到.
金錢是中性的.
是一個媒介.
我們怎樣能夠用得其所呢?.
在這些聖公上.
所以清教徒的.
金錢觀是很奇妙的.
他們因為知道.
怎樣能夠保證.
自己是蒙神簡選呢?.
就是我一直的事工.
我的工作.
蒙神的賜福.
所以他們就很勤奮工作.
很勤奮的儲蓄.
然後很盡力的奉獻.
那就證明.
他們自己是蒙神簡選.
所以清教徒.
就變成新教的.
那個叫做資本家.
就是成為了.
一個時代.
的一個很重要的發言.
為什麼歐洲的.
資本家會是新教徒.
大部分是新教徒呢?.
就是離開羅馬天主教的.
新教徒呢?.
就是後來有社會學家研究.
就是因為簡選.

$^{1281}$他們很怕自己是不被神簡選.
所以他們要.
活在神所喜悅的.
那個教道裡面.
當然就是嘉爾文.
慈雲尼這些.
的牧者呢 中國改革的時候.
來去教導.
尤其是他們說的那個神的簡選.
所以.
很多很多的新教徒.
跟隨脫離羅馬天主教之後.
他們就要活在.
這種蒙簡選.
就是蒙福的那個生活裡面.
所以就.
不知不覺間.
就變成一個新的階層 就是資本家.
因為他們有錢嘛 又不亂花.
就帶來一個.
時代的一個產物.
主要不是因為他們想.
拼命存錢 是因為他們擔心.
自己是不是被神簡選.
所以你可以想像到.
如果一個人對神有一個.
充分的掌握呢 那他的人生呢.
都會朝向在神.
那個心意裡面生活.
所以這裡就是神告訴我們.
你要來到去.
不要搶奪神 應該得到的.
所以無論是我們.
我們要將人.
帶進神的榮耀裡面.
我們要來到去講論福音.
用不同的方法 不同的方式.
來到去承載這個福音.
那就更加讓我們.
能夠在這裡.

$^{1321}$告訴我們 神就好像在這裡.
這樣來挑戰百姓.
其實都是一種.
對他們的一個責備.
何必試呢.
神就是這樣的.
所以如果我們用現在.
現代呢.
我作為牧者 作為堂主任.
每個禮拜是要看.
每一個弟兄姊妹的奉獻的情況.
一來都是堂委會.
或者我有個責任要了解.
有些事我們.
快到四月了 我要簽.
每一張的奉獻收據.
這個是政府規定的.
堂主任加上那個師父.
要簽那個免稅收據.
每一個弟兄姊妹的奉獻.
某程度上我都知道.
但是我知道之後.
我轉過頭來 我忘記了.
我忽然忘記了.
我不會太在意.
每一個弟兄姊妹.
你那個奉獻的雜性.
每個人自己向神交代.
每個人自己面對神.
不是面對我.
我也不會因為你奉獻多.
我就特別好好你 不需要的.
我們作為牧者要有這種.
叫什麼.
叫Guts.
Guts就是什麼.
中文是骨氣.
有骨氣.
不要見高拜見低踩.
不要因為有錢就特別.

$^{1361}$特別探多些.
沒錢就快點走開.
千萬不要 這樣做全道都不適合.
但是這裡已經告訴我們.
這段聖經是說.
那個轉向是具體的.
那個轉向.
你自己得到神的恩典.
你還這麼淪陷.
對神的國度的工作.
其實你根本還沒有明白神的救恩.
一個明白神的救恩的人.
他不會介意自己.
能夠他手上所做的.
能夠找到的.
因為一切都是從神那裡來的.
所以這個整段的教導.
神就說.
這一幫人如果他們這樣去做.
最終告訴我們.
他們是必稱為有福的.
而且他們的地.
必稱為喜樂之地.
耶路撒冷充滿了喜樂.
充滿了平安.
充滿了恩典.
充滿了神同在帶來的.
那種的歡欣快樂.
所以這裡上文說的.
說錢不會失感情.
說錢是真的帶來.
神的恩典流到.
每一個家庭.
每一個子民當中.
所以這裡就是.
不只是第十節.
這樣說 試一試等等.
而是整個的.
教會.
對金錢的觀念.

$^{1401}$金錢原來.
背後的價值觀就是.
敬錢加上知足的心.
他就自然會.
對神的事情.
他會放在心裡.
因為我們沒有.
帶什麼來這個世界.
也不能夠帶走什麼.
有的是陪葬.
不過其實你都不知道的.
放在棺木裡幫你陪葬.
其實都不是你的.
接著.
那個聚會結束了.
有些不長進的回家就搶家產.
就一起.
去吵架.
可能在吃解慰飯都已經吵架了.
其實就很可惜.
所以.
我們要認定.
不是叫我們良薄.
不是叫我們做人做到.
這麼虛空的.
我不背了 也不需要的.
我們要記得正正就是因為.
我們沒有帶什麼來這個世界.
也不能夠帶走任何東西.
所以在世界裡有的.
世界裡有的就是.
敬錢加上知足的心.
就是大利.
所以第三章的.
馬拉基書.
在這一段裡.
神就說我會練淨的.
我會使你們.
能夠在我的面前.
重新獻神.

$^{1441}$月立的.
潔淨的祭物.
接著你們就要轉向我了.
因為我已經轉向你們了.
接著這個轉向.
就是實際上.
百姓.
轉向神之後.
他有沒有那個敬錢對神的心呢.
有沒有對.
自己的一個知足的心靈呢.
真的.
明白人生在世.
絕對不虛空.
你知道自己要離開這個世界.
可能會很虛空.
但是有神就不會虛空了.
有神你就會知道.
在地上還有一口氣.
就是還有那個機會和責任.
所以你有那個敬錢.
加知足.
你就能夠運用神給你的恩典.
其中一樣就是.
很明顯.
就是錢財.
來為神的.
福音的聖功.
所以神就說.
這樣的人.
表面上就好像.
去測試.
神的那個.
應許.
但實際上他內心.
如果一個不是有敬錢有知足的心的人.
其實是沒有可能.
經常持久的.
也不是誇耀什麼.
我信主就有很大改變.

$^{1481}$那時候我還沒有.
認識這些經文.
我已經定期.
那時候不是很多錢.
就十一奉獻.
後來更加有感動.
有一段時間我忘記了維持一年還是兩年.
我是十二奉獻.
將自己找到的.
那個月薪.
抽兩成奉獻.
給教會.
後來當然慢慢回到十一.
由我信耶穌到今天.
都是十一奉獻這樣做.
是你甘心樂意.
總之到一號.
自然就這樣做了.
明仔也是這樣.
他自己由這樣開始到現在.
他覺得我和神計.
神和我計我就不行了.
所以他自己已經做了.
那樣東西.
他覺得那樣東西不是他的了.
是神的了.
那時候內心反而充滿了喜樂.
充滿了平安.
這告訴我們.
神說你這樣試試我的時候.
第十一節.
我必為你們斥責蝗蟲.
即是你們耕種.
你們的田地特別好收成.
不會有害蟲.
不容許他們毀壞你們的土產.
怎樣做其實不知道.
神怎樣令到那樣東西發生.
你們的田產.
葡萄樹未熟.

$^{1521}$都不會掉到地上.
接著萬國都稱你們有福.
原來一個越肯奉獻給神的人.
越是蒙福.
在神的眼中蒙福.
在人的眼中.
你這麼蠢的.
奉獻錢給神.
很多微信的同事都是這樣笑我.
信耶穌十一奉獻的.
當你已經知道自己是蒙福.
心靈的福氣.
你根本不會介意.
所以感恩早期.
一路不是這麼有錢的時候.
你做了十一奉獻.
你就不會這麼肉痛.
但做著做著就習慣了.
每逢是一號就一定要抽十分一出來給神.
那你就可以想像到.
一生到現在都蒙福.
這個是經文的應許.
所以作為牧者我們不是經常說錢.
但是要教導什麼叫正確的.
神眼中的金錢觀.
其實是敬虔.
其實是敬虔加上知足的心.
就是大利.
沒有了敬虔.
知足來護衛我們的心.
貪念就會進入我們的心.
錢財本來中性.
變成惡性.
變成是萬惡之根.
所以保羅真的很好.
他說不是金錢是萬惡之根.
是貪財.
是萬惡之根.
所以今晚大家回去背熟.
提摩太前書六章十節.

$^{1561}$然後就數一數自己有多少雙鞋.
有多少件衣服.
就去問一問自己.
其實我是不是真的敬虔.
加上知足呢.
當然好像說笑那樣.
但是某程度上.
我們要回轉.
要轉向神的一個具體.
因為神說得這麼白.
正正就是當時的百姓.
一直都在.
還沒有肯行這個行動.
怎樣去回轉.
他講得很振振有詞.
第七節.
你們卻問說.
我們如何才是轉向.
這些好像很得意的.
其實他如果明白神憐憫他們.
挽回他們.
就不應該這樣問這個問題.
所以大家回來《查經班》.
你越聽你越恐怖.
《查經班》你越聽你越恐怖.
你不聽反而沒有東西.
越聽你越不行.
家裡真的很多衣服.
其實這是一個回轉.
正正就是我們有生之年.
我們真的一定要操練.
就是敬虔加知足.
真是根據.
六章五節.
《提案太前書》.
真的要讓這件事.
如果不這樣操練.
貪財就會進入我們的心.
就會進入我們的心.
剛剛恥笑比掌.

$^{1601}$就會帶來一個很大的落差.
所以慢慢很多人.
曾經愛主曾經事奉.
慢慢不知不覺他離開.
其實有跡可尋.
其實他已經沒有了敬虔加知足.
他就會被貪財.
霸佔了他的心.
帶來.
去日本有什麼問題.
櫻花這麼漂亮就快開了.
準備你一個月去.
那個就不是一次半次了.
每天都這樣做的時候.
其實慢慢不知不覺.
自己在騙自己.
其實真正的敬虔加知足.
就是你要做在神的角度.
想想神的角度.
想你怎樣參與幫手.
就不是偶然間一個調劑.
偶然間我們可以有一點點的放鬆.
是完全不同的.
這個就是敬虔加知足.
保羅說大利.
我們不斷追求大利.
但他想不到.
原來聖經說的大利是這麼簡單.
我們都比較情願.
拒絕這種.
心靈的滿足.
第十三節.
這裡就更加寶貴了.
原來有些人聽完馬拉基說.
就行動.
就做.
所以我們上茶經班其實都很悶的.
現在這個時段是甚麼劇集的.
因為我家裡沒有電視的.
我不知道TVB在做甚麼劇集的.

$^{1641}$是不是殺人放火那些.
還是11點鐘那些時段才有的.
諸如此類.
或者可能好笑的笑片.
那你一個星期沒有了一晚.
其實都很慘的.
但是你要留意一下.
馬拉基說.
有些人是真的會聽從.
第十三節.
你們用話頂撞我.
接著你們就說.
我們用甚麼話頂撞你.
你們說.
事奉神是徒然的.
原來這樣是頂撞神的.
遵守神吩咐.
在萬君之耶和華面前.
苦苦的齋戒.
有甚麼益處.
原來這裡就真的反映到.
原來人.
慢慢叫做.
世故.
慢慢有很多世故的基督徒.
慢慢看東西.
見怪不怪.
不要激氣.
盡量不要說話.
原罪快點走.
不要涉及太多東西.
找你事奉.
近來很忙.
其實根本從來都沒有想過.
想再事奉神.
因為受過傷害.
最好不要坐那麼前.
坐後一點.
唱完詩歌.
開始放人進來.

$^{1681}$你就坐最後.
有這樣的模式.
坐最後.
散會.
牧師一祝福完就立刻走了.
不要等報告了.
諸如此類.
原來是這種心態.
原來是這樣.
頂撞神.
就是他們用這種.
覺得.
神都不會再在意他們的敬拜.
一定是由一個很熱心齋戒.
很樂意去遵守神的命令.
到慢慢發現原來身邊的人.
都不再緊張.
為什麼自己還要緊張.
為什麼自己不斷去祈禱.
鼓勵人去祈禱會.
鼓勵人去崇拜.
但是慢慢.
你會發現原來自己是最.
在很多人的眼中是最傻的人.
有什麼益處呢.
開始覺得.
做這一切都是.
例行公事.
心都如止水.
這個就是.
很具體的去剖析.
一個已經內心.
好像已經麻木的基督徒.
那個心靈狀態.
神說這個是頂撞.
是對神的同在的一個頂撞.
所以要逼自己.
每一次你是參加者.
所以朝聖茶室.
真是很好.

$^{1721}$是朝聖.
我們去到神的面前.
好像從湖邊走回來.
我們好像一個上行的敬拜.
雖然是平路.
我們是朝見神.
所以記得你有崗位也好.
沒有崗位也好.
你都是朝見神.
你的對象.
你今天要見的主角就是神.
所以我們去朝見神.
所以如果不是就變成了頂撞.
你就變成了.
你的例行公事.
你又不敢不來.
慢慢不來你就什麼都沒有了.
所以你要維持逢禮拜日.
或者禮拜六晚.
那一個小時我一定要出現.
我一定要坐回我自己的位置.
我一定要看看你們有什麼錯.
真的回來收留.
收留你的錯.
不是他的錯.
這樣去明察秋毫.
其實這個頂撞神.
好像一個屬靈的判官.
這樣不對.
這樣不行.
那堂聚會就真的.
堂堂都興奮.
因為回來收留.
然後望程序表誰奉獻.
誰沒有奉獻.
其實你有沒有奉獻呢.
當然我知道.
那你有沒有見到.
這個原來是頂撞神.
感恩就是第十五節.

$^{1761}$如今我們稱狂傲的人為有福.
並且狹惡的人得建立.
他們雖然試探神.
卻得脫離災難.
那個心境去到這麼負面.
負面到一個地步.
狂傲是有福.
狹惡是得到建立.
試探神可以脫離災難.
不過感恩就是十六節.
不知為何三章十六節.
永遠都有很多金句.
新約的三章十六節.
大家很熟悉.
「神愛世人」那一節.
很多三章十六節都很好.
《提摩特後書》三章十六節.
《聖經》都是神所默示.
「與教訓督澤使人歸正教導人學義」.
都是有福的.
「叫屬神的人得以完全.
避免各樣善事」.
還有很多三章十六節.
不過這裡的三章十六節很多人不懂.
是真的很好.
「那時敬畏耶和華的彼此談論」.
「耶和華則以而聽」.
而且有紀念冊.
在神的面前.
記錄那些敬畏耶和華.
思念他明的人.
「萬君之耶和華說.
在我所定的日子.
他們必定屬我.
而且是特地歸我.
我必連戌.
這些敬畏他的人.
如同人連戌.
服侍自己的兒子.
那時你們要歸回.

$^{1801}$將善人和惡人.
事奉神和不事奉神的.
就分別出來」.
原來將來有一個真正的分別.
不過現在.
神就說.
祂有一個紀念冊.
紀念那些思念他明的人.
談論神的人.
所以你看看多寶貴.
上文就說.
有什麼用呢.
現在稱讚那些狂傲的有福.
行惡的得到保守.
得到建立.
試探神的.
可以遠離罪惡.
可以遠離災難.
不是罪惡.
是遠離災難.
這樣的負面.
充斥著在憲制的日子.
每個人都上耶路撒冷.
原來一讓祭.
就養百樣人.
很多人帶著這樣的心態.
去有什麼用呢.
去即管看看哪些人有錯.
我給神面子.
去而已.
所以我可以不去.
但是有些敬虔的人.
所以神自己幫我們做分別.
所以真的不要看人.
要看自己.
每一次的敬拜.
是否真心.
是否單純的敬錢.
加知足去敬拜神.
不要看人.

$^{1841}$真的不要看人.
甚至誇張到.
那個老牧者.
在台上不斷地釣魚.
你都不要介意.
可能他真的很累.
他真的很辛苦.
他年紀都很大.
但仍然很衷心地去祝福.
去守望教會.
從這個角度去看.
這些長輩.
反過來你自己.
十六節.
如果你可以跟我一起讀這節聖經.
我很少在茶經班叫大家一起讀經.
因為都想多說一些.
那時.
同樣都是那個時候.
原來神看到的.
不是那些不敬畏他.
和不尊重敬拜的人.
他看到的就是.
那些敬畏他.
那些彼此談論.
很想談論關於神的事.
你多說一些給我聽.
我很想知道神是怎樣的.
譬如你進到聖殿裡面.
有甚麼經歷.
我很想知道一下.
這種渴慕神的事.
是神閱立的.
人若善無善功.
是可閱立的.
這是提文派前書的三章.
是這樣說的.
他說.
人若渴慕善功.
是神所閱立的.

$^{1881}$所以.
往往一些.
如果他在教會不斷祈禱.
不斷傳福音.
往往這些人就是神呼召傳道的前傳.
他自己就去做.
我又說過.
我初信的時候.
就派去屯門醫院.
那時剛開院.
那些護士.
每天一個月下來.
那麼多護士.
下來拿支票.
因為未出到糧 未入到戶口.
下來拿支票.
我們叫阿打.
那些護士長.
有時他們又要下來.
那些單要付錢.
那時很熟悉.
有時我上去探病.
我下班五時準時收工.
沒有甚麼事做.
上去就向那些病人.
逐張床派單張.
當然.
兒科去不了.
婦科去不了 不讓你去.
通常去那些.
癌症的病房.
那時沒有那麼多人.
任你喜歡怎樣做都可以.
沒有人說你 再加上我是職員.
那些護士又認識.
護士長又認識.
向那些病人傳福音.
真的帶給我自己.
那一年多兩年.
在政府.

$^{1921}$真的有很大的生命改變.
看到人的靈魂的需要.
決心成為一個傳道人.
這個就是.
我們要.
談論神.
所以.
真的 上小一晚的.
看電視.
來查經班 不是我們要.
特意誇耀甚麼.
而是大家一起認真讀下聖經.
未必我們識原文.
但是讀下中文.
讀通它 我們已經感受到.
神的話語的甘甜.
那種的模樣.
所以你自己已經學到.
至少你學到人生最重要的是甚麼呢.
有錢 當然是重要的.
但是要有.
敬虔 對神的敬虔.
對自己的知足.
對自己對人的知足.
那你就會真是無往而不利的.
大利.
然後這裡告訴我們.
原來神留意著.
那時敬畏耶和華的彼此談論.
耶和華則以而聽.
你想想多傳神.
原來我們在談論神的事.
神則以而聽.
有紀念冊.
即是說.
不是一聽完就忘記了.
是紀念冊.
紀念著這些事情.
在祂的面前.
記錄了敬畏耶和華.

$^{1961}$思念祂明的人.
所以這個還不是金句.
肯定是舊約裡面.
其中一節重要的金句.
所以我們要背熟這些金句.
所以又是要擺上下文.
上文是很負面的.
頂撞神的.
不知道甚麼叫轉向神的.
神就說.
有這樣的時代.
即是說原來馬拉基.
表面上好像很多人.
是宗教混亂.
但原來也有一些事.
是小部分的人.
馬拉基的教訓.
他們聽.
他們願意去遵行.
所以有時一堂崇拜講道.
不是每個人都聽得明白.
聽完願意去走.
認真面對自己的心靈.
那些人就會蒙福.
聽完當作.
金堂堂很好.
那個例子不錯.
很傳神.
他們聽到的.
但有時不是全面的訊息.
其實要學習聽到.
那個訊息.
帶出的重要.
對我的心靈有甚麼的看法.
有甚麼的影響.
這是那個訊息.
例子是幫助我們明白.
那個訊息的一些材料.
所以不要因為沒有例子.
就睡著了.

$^{2001}$不一定需要這樣.
要學習.
聽到訊息.
就算是一個神學生.
我都在他內心.
在他講道裡面.
得到一些東西.
有時他講得很害怕.
很藤雞.
我就在台上.
為他祈禱.
為牧者祈禱.
因為神總有透過他的僕人.
來傳遞一些訊息.
即便那個訊息是很淺.
你都明白的.
一個護面.
所以.
這裡告訴我們.
這一節金句.
是在那個時代.
頂撞的時代.
不知甚麼是迴轉.
其實知道的.
來到思探神的時代.
但這裡告訴我們.
有很少聽從馬拉基先知的.
敬虔的猶太人.
他們會這樣來彼此的談論.
所以不是在聚會談話.
而是他們.
談論神.
用不同的方式.
WhatsApp 或者不同的渠道.
感恩的.
真是感恩.
在講美國那位女士.
陳女士.
她說遲些電郵給我.
因為不想WhatsApp問得太長.

$^{2041}$她說有很多神學的問題.
很想認真對這個信仰.
是很感恩的.
因為她都聽這些茶經的.
慢慢改變她的心靈.
希望她的身體的痛.
可以慢慢的減少.
所以盼望我們能夠.
還是有思想.
還是有教會生活.
我們一方面不是不要理人.
而是我們自己的內心.
但我們應該要和弟兄姊妹.
結連.
來幫助其他人.
最好多些去談論神.
多些在神裡面去討論.
使我們能夠.
得到神的眷念.
所以要記住.
有紀念冊在她的面前.
那本紀念冊不會不見了.
我們小學的紀念冊.
現在不見了.
真是價值連城.
雖然那些人已經沒有聯繫了.
但是如果有就很珍貴了.
五六十年前的紀念冊.
但是在神面前的紀念冊.
這個不是生命冊.
是紀念冊.
敬虔的人.
他們的一生.
在神的面前.
所以如果我們每一個禮拜的出現.
都是成為別人的祝福.
那就是一個很蒙福的人.
我們帶著的是.
有時很不明白.
為何那些人.

$^{2081}$他又不是真的那麼苦.
帶著很多苦來教會.
你說他真的很困難.
帶著苦來當然有人幫他.
他不是的.
再嚴重一點帶著苦毒來教會.
其實你想想多麼可惜.
是吧?.
很多事情其實算了.
到了這個年紀.
這些事情算了還看得那麼緊.
反過來.
敬虔之足他又不去追求.
所以這裡告訴我們.
那時敬畏耶和華彼此談論.
耶和華則以而聽.
有紀念冊.
然後記錄那些敬畏耶和華.
和思念他明的人.
所以思念.
即是你已經有了知識.
你要再去思想.
反覆地去想明白.
這叫思念.
所以這也是.
我們讀經.
大家都要認真.
因為你不讀經.
你沒有那麼多材料在你的思想裡.
你不夠那個思想.
去想神的道.
所以我們自己要多讀經.
或者聽用你的方法.
聽我就不行我會亂.
我一定要讀的.
要看的 我要揭的.
所以手機的apps.
是幫我背誦聖經.
我找的.
我不是靠apps讀經.

$^{2121}$是靠大本的聖經.
去研經真的要那麼大本才行.
所以.
如果你想自己.
聖經的根底慢慢再打好.
一定要有一本大本的聖經.
要投資的.
還有一個book stand.
45度角.
這樣看才舒服.
我們現在去到這個年紀.
會有老花.
有軟事.
有散光其實不行的.
所以這裡.
你羨不羨慕.
有紀念冊.
神紀錄.
敬畏祂的人.
思念祂名的人.
你羨不羨慕就是.
神在你的內心裡.
給你這種滿足.
不過將來肯定有.
一個這樣的分別.
神就說明了 這些人是屬於我的.
所以這個就是.
清教徒很重要的信念.
他們怕.
很簡單的邏輯.
宗教改革.
是帶來什麼.
是脫離羅馬的天主教.
天主教說什麼.
教皇說什麼.
教會以外沒有救恩.
是真的這樣說.
恐嚇基督徒.
嘩.
加爾文.

$^{2161}$馬丁路德.
我現在跟你們走.
脫離天主教 我們有什麼信心保證.
加爾文他們就說.
就是你們蒙簡選.
怎樣為之蒙簡選.
就是你們真的來到.
在你們的生活.
在你們的工作裡面.
看到神的祝福.
他們就需要勤力的工作.
勤力的.
來儲蓄.
在教會裡勤力的奉獻.
這樣就證明自己是蒙簡選.
要用這些.
具體的呈現的方法.
所以這裡告訴我們.
他們必屬我.
特特的歸我.
是說神.
我必定會連戌他們.
如同人連戌.
服侍自己的兒子.
你們就會歸回.
這是說將來.
將善人和惡人.
事奉神和不事奉神的.
分別出來.
所以神已經說得很清楚.
如果我們今生都不去事奉.
不用真心去事奉.
不用那份恩典去事奉.
保羅看自己的事奉.
是恩典.
真的 他看他自己的事奉.
羅馬書一章.
基督耶穌的僕人保羅.
我寫得.
一下子突然間講.

$^{2201}$就不記得背了.
一章.
奉召為使徒.
特派全神的福音.
特別被委派的.
這福音是從前.
眾先知在聖經上所應許.
論到他兒子.
我主耶穌基督.
按照肉身來說.
他是大衛的後裔生的.
按照聖善的靈.
即是良善聖潔的靈.
他從死裡復活.
以大能.
即是復活的大能.
顯明他是神的兒子.
即是既是人大衛的子孫.
在這個世界.
但復活的大能.
證明顯明印證他是神的兒子.
就是那位彌賽亞.
留意一下這裡.
第五節.
我們從他受了恩惠.
並使徒的積分.
在萬國之中.
叫人為他的名信服真道.
所以保羅看自己.
是一個.
蒙恩典的事奉.
很多卷書的開頭.
都是說.
他是一個蒙恩典事奉的一個工人.
所以是很甘心樂意的.
所以如果我們的事奉.
是不甘心不樂意.
證明我們對恩典不是太明白.
如果我們明白恩典.
我們便會很甘心樂意去服侍.

$^{2241}$任勞任怨.
沒有所謂.
因為神是會留意.
觀察服侍他的人.
他會則以而聽.
有生命.
有紀念冊在他的面前.
所以我們記得.
我們事奉是向準神便可.
當然我們一定要盡我們的力.
因為人看到的是你的盡力.
不是看到神的紀念冊.
如果他看到.
他便會有邪靈.
看到鬼.
他看不到紀念冊.
他看到你是否忠心.
對方.
所以我們要對準神.
你自然便會.
可以在人的面前顯現.
這段是第三章.
所以這樣有一個強烈對比.
那些.
甚麼叫.
轉向神.
頂撞神.
等等的人.
神都有解答.
甚麼是公義.
公義的神在哪裡.
很多這些問題.
都是源於他們在那個時代.
面對很多.
不是那些大的災難.
而是日夕的人際相處.
就要很多.
自我保護.
帶來的鬼壓.
慢慢令到.

$^{2281}$心靈蒙了一個很大的陰影.
所以事奉.
都變成一種.
可有可無的.
例行公事.
甚至乎.
索性根本不再這樣去做.
所以有時都很慨嘆.
你信主很多年.
你的兒童主一學的那些同學.
或者團契.
有些已經不見了.
有些已經不再信.
不是你特別好.
不是我特別好.
是我們蒙恩在神的恩典裡.
所以我們要記得.
我們透過這一大段經文.
《馬拉基書》.
就將普遍的.
那個社會狀況.
透過他的說話.
神的啟示.
來寫出來.
但是神自己有那個解答.
神自己有那個.
給的出路.
就是.
三章十六節.
所以.
彌賽亞的降臨.
就是要帶來那個廉政.
所以我們盼望信主.
不是徒勞的知識.
那個是必須的入門.
如果你沒有這個真理的知識.
你都不知道什麼叫信耶穌.
但是真理的知識.
是應該帶來我們的生命的一個改變.
今晚我們就講到這裡.

$^{2321}$下次我們講完第四章.
我們就可以完結了.
整個歸回的歷史.
和先知書.
如果你有時間你可以回聽.
沒有時間就讓他們放棄.
因為繼續還有很多書卷.
我們都還沒有講.
感恩給我機會.
其實都是蒙的.
今天我由早上開會開到五點.
在會裡面都睡著了.
人家開會講話我睡著了.
我還以為那些是外人.
不是自己的弟兄姊妹.
因為腦袋累了.
但是不要緊.
我反而講聖經就很精神的.
感恩.
我們一起簡單祈禱.
多謝主讓我們一起學習聖經.
多謝主.
一直我們看.
哈該書,撒迦利亞書.
和這個馬拉基書.
還有我們看到的.
歷史書.
有以色列記.
尼希米記和以斯帖記.
差不多六卷書我們都看完了.
還有最後一章第四章.
盼望我們能夠對主耶穌的降臨.
對主耶穌在地上的工作.
我們有認識.
幫助我們能夠對三章十六節.
這一節這麼寶貴的金句.
我們能夠明白.
以致我們在有生之年.
我們是一個很活潑的敬拜神.
我們是真的對談論神.

$^{2361}$是我們最大的滿足.
願神幫助.
奉主耶穌基督的名.
阿們.
\newpage



\section{創世記 32:1-12-22-32}
\label{sec:oyj7tRcaGj0}
\textbf{2024-03-03 早堂《雅各的祈禱》創世記 32\_1-12,22-32,33\_1-4 講員:謝靄薇牧師}
\newline
\newline
連結: \href{https://youtube.com/watch?v=oyj7tRcaGj0}{\texttt{https://youtube.com/watch?v=oyj7tRcaGj0}} ~~~~ 語音日期: 2024-03-04
\newline
\newline
\hyperref[sec:exEZqq1nt80]{\small{< < < PREV SERMON < < <}}
~
\hyperref[sec:index_chronic]{\small{[返順時目]}}
~
\hyperref[sec:index_scriptual]{\small{[返順卷目]}}
~
\hyperref[sec:a1w0x_rlHHg]{\small{> > > NEXT SERMON > > >}}
\newline
\newline
創世記 32:1-12-22-32
\newline
\begin{longtable}{cl}
\hline
\hline
章節 & 經文 (和合本修訂版)\\
\hline
32:1 & \begin{tabularx}{0.7\textwidth}{X} 雅各繼續行路，神的使者遇見他。 \end{tabularx} \\ \\ \relax
32:2 & \begin{tabularx}{0.7\textwidth}{X} 雅各看見他們就說：「這是神的軍營。」於是給那地方起名叫瑪哈念 。 \end{tabularx} \\ \\ \relax
32:3 & \begin{tabularx}{0.7\textwidth}{X} 雅各派使者在他前面到西珥地，就是以東地他哥哥以掃那裡。 \end{tabularx} \\ \\ \relax
32:4 & \begin{tabularx}{0.7\textwidth}{X} 他吩咐他們說：「你們要對我主以掃說：『你的僕人雅各這樣說：我在拉班那裡寄居，延遲到如今。 \end{tabularx} \\ \\ \relax
32:5 & \begin{tabularx}{0.7\textwidth}{X} 我有牛、驢、羊群、奴僕、婢女，現在派人來報告我主，為了要在你眼前蒙恩。』」 \end{tabularx} \\ \\ \relax
32:6 & \begin{tabularx}{0.7\textwidth}{X} 使者回到雅各那裡，說：「我們到了你哥哥以掃那裡。他正迎著你來，並且有四百人和他一起。」 \end{tabularx} \\ \\ \relax
32:7 & \begin{tabularx}{0.7\textwidth}{X} 雅各就很懼怕，而且愁煩。他把跟他同行的人和羊群、牛群、駱駝分成兩隊， \end{tabularx} \\ \\ \relax
32:8 & \begin{tabularx}{0.7\textwidth}{X} 說：「以掃若來擊殺其中一隊，剩下的另一隊還可以逃脫。」 \end{tabularx} \\ \\ \relax
32:9 & \begin{tabularx}{0.7\textwidth}{X} 雅各說：「耶和華—我祖父亞伯拉罕的神，我父親以撒的神啊，你曾對我說：『回你本地本族去，我要厚待你。』 \end{tabularx} \\ \\ \relax
32:10 & \begin{tabularx}{0.7\textwidth}{X} 你向僕人所施的一切慈愛和信實，我一點也不配得。我先前只用我的一根杖過這約旦河，如今我卻成了兩隊。 \end{tabularx} \\ \\ \relax
32:11 & \begin{tabularx}{0.7\textwidth}{X} 求你救我脫離我哥哥的手，脫離以掃的手，因為我怕他來殺我，連母親和兒女都不放過。 \end{tabularx} \\ \\ \relax
32:12 & \begin{tabularx}{0.7\textwidth}{X} 你曾說：『我必定厚待你，使你的後裔如同海邊的沙，多得不可勝數。』」 \end{tabularx} \\ \\ \relax
32:13 & \begin{tabularx}{0.7\textwidth}{X} 當夜，雅各在那裡住宿，就從他手中所擁有的拿禮物要送給他哥哥以掃， \end{tabularx} \\ \\ \relax
32:14 & \begin{tabularx}{0.7\textwidth}{X} 就是二百隻母山羊、二十隻公山羊、二百隻母綿羊、二十隻公綿羊、 \end{tabularx} \\ \\ \relax
32:15 & \begin{tabularx}{0.7\textwidth}{X} 三十匹哺乳的母駱駝和牠們的小駱駝、四十頭母牛、十頭公牛、二十匹母驢和十匹公驢。 \end{tabularx} \\ \\ \relax
32:16 & \begin{tabularx}{0.7\textwidth}{X} 他把每種牲畜各分一群，交在僕人手中，對僕人說：「你們要在我的前頭過去，使群和群之間保持一段距離」。 \end{tabularx} \\ \\ \relax
32:17 & \begin{tabularx}{0.7\textwidth}{X} 他又吩咐領頭的人說：「我哥哥以掃遇見你的時候，問你說：『你是誰的人？要往哪裡去？你前面這些是誰的？』 \end{tabularx} \\ \\ \relax
32:18 & \begin{tabularx}{0.7\textwidth}{X} 你就說：『是你僕人雅各的，是送給我主以掃的禮物。看哪，他自己也在我們後面。』」 \end{tabularx} \\ \\ \relax
32:19 & \begin{tabularx}{0.7\textwidth}{X} 他又吩咐第二、第三和所有趕畜群的人說：「你們遇見以掃的時候要照這樣的話對他說， \end{tabularx} \\ \\ \relax
32:20 & \begin{tabularx}{0.7\textwidth}{X} 你們還要說：『看哪，你僕人雅各在我們後面。』」因雅各說：「我藉著在我前面送去的禮物給他面子，然後再見他的面，或許他會寬容我。」 \end{tabularx} \\ \\ \relax
32:21 & \begin{tabularx}{0.7\textwidth}{X} 於是禮物在他前面過去了；那夜，雅各在營中住宿。 \end{tabularx} \\ \\ \relax
32:22 & \begin{tabularx}{0.7\textwidth}{X} 他夜間起來，帶著兩個妻子，兩個婢女和十一個孩子，過了雅博渡口。 \end{tabularx} \\ \\ \relax
32:23 & \begin{tabularx}{0.7\textwidth}{X} 他帶著他們，送他們過河，他所有的一切也都過去， \end{tabularx} \\ \\ \relax
32:24 & \begin{tabularx}{0.7\textwidth}{X} 只剩下雅各一人。有一個人來和他摔跤，直到黎明。 \end{tabularx} \\ \\ \relax
32:25 & \begin{tabularx}{0.7\textwidth}{X} 那人見自己勝不過他，就摸了他的大腿窩一下。雅各的大腿窩就在和那人摔跤的時候扭了。 \end{tabularx} \\ \\ \relax
32:26 & \begin{tabularx}{0.7\textwidth}{X} 那人說：「天快亮了，讓我走吧！」雅各說：「你不給我祝福，我就不讓你走。」 \end{tabularx} \\ \\ \relax
32:27 & \begin{tabularx}{0.7\textwidth}{X} 那人說：「你叫甚麼名字？」他說：「雅各。」 \end{tabularx} \\ \\ \relax
32:28 & \begin{tabularx}{0.7\textwidth}{X} 那人說：「你的名字不要再叫雅各，要叫以色列，因為你與神和人較力，都得勝了。」 \end{tabularx} \\ \\ \relax
32:29 & \begin{tabularx}{0.7\textwidth}{X} 雅各問他說：「請告訴我你的名字。」那人說：「何必問我的名字呢？」於是他在那裡為雅各祝福。 \end{tabularx} \\ \\ \relax
32:30 & \begin{tabularx}{0.7\textwidth}{X} 雅各就給那地方起名叫毗努伊勒 ，說：「我面對面見了神，我的性命仍得保全。」 \end{tabularx} \\ \\ \relax
32:31 & \begin{tabularx}{0.7\textwidth}{X} 太陽剛出來的時候，雅各經過毗努伊勒，他的大腿就瘸了。 \end{tabularx} \\ \\ \relax
32:32 & \begin{tabularx}{0.7\textwidth}{X} 因此，以色列人不吃大腿窩的筋，直到今日，因為那人摸了雅各大腿窩的筋。 \end{tabularx} \\ \\
[1ex]
\hline
\hline
\end{longtable}
$^{1}$各位弟兄姊妹平安.
在我們看神的話語之前.
我們一起仰望神.
每一天 每一次.
我們來到您面前的時候.
我們真的渴望您.
在當中幫助我們.
能夠思想您這份愛.
能夠常存在我們心裡.
若然我們這個時候.
有什麼是不滿意的.
求聖靈在當中剔除.
以致我們能夠專心一意的.
看到您自己的心意.
我們就願意去靠著您.
我們可以走出來.
求主帶領我們的心.
讓孩子心中的意念.
口中的言語.
蒙您如立.
奉主耶穌基督的名求.
阿們.
相信我們今天看這個雅閣.
大家已經很熟悉.
特別是.
他很喜歡跟人爭奪.
從小他一出生就要爭奪.
他要抓住哥哥的腳.
爭奪的腳要出來.
但事實上他輸了.
所以他就變成弟弟.
在人生裡.
大家會看到.
這一幕我們有時也很遺憾.
但原來就是一碗紅湯.
不知道是不是紅豆湯.
紅湯可以哄到他哥哥.
把長子的名分給了他.
然後.
在他長大一點的時候.

$^{41}$又騙到他爸爸的祝福.
好像跟他媽媽聯手.
他在母蔭之下.
很輕易的去騙到祝福.
但你要知道.
他帶來後患無窮.
仇恨的惡果就開始了.
你要知道.
當他爸爸的祝福.
原本是給哥哥的時候.
哥哥很開心地回來.
但完全什麼都得不到.
你可以想像他的心情.
於是他就殺了你.
但他媽媽很疼他的.
於是他媽媽就想著.
幫他想個方法.
可以逃避這個災難.
就使他去求福.
一去就20年了.
如果你最愛的兒子.
一去就離開了20年.
你媽媽心境如何呢.
她媽媽以為幫助他.
可以逃避哥哥的殺害.
但是.
你知道.
他應該是家裡媽媽好好照顧.
所以在路上.
他找個石頭去休息.
在那邊.
他在夢裡看到天梯.
看到天使上去下來.
原來他發現.
在他人生裡.
似乎.
在這個時候他很低落.
因為他要離開家.
離開他熟悉的地方.
在這個地方他發現.

$^{81}$原來.
雖然媽媽叫他去逃避這個災難.
但他看到.
他在逃避的裡面.
他很希望.
神可以保守他.
於是他發現.
原來.
這位神真的在當中.
神是與他同在的.
神可以在他最需要的時候.
達到神的面前.
而神也可以.
將祝福臨到他.
原來神是這樣的神.
今天我們就來看看.
祂的祈禱是怎樣的呢.
各位, 你想想.
你在哪個時候.
你會最迫切的去祈禱呢.
在什麼環境之下.
是不是走投無路的時候.
是不是沒有人幫你的時候.
不成功的時候.
你那個時候.
勉為其難.
去祈禱試試.
是不是.
回想他這20年.
真是詐騙.
很多詭計.
騙了爸爸的祝福.
但哥哥就發誓要殺他.
現在20年後.
他要回家.
哥哥就在那裡.
怎麼辦呢.
他很焦慮.
每次想到他哥哥就焦慮.
在回家的路.

$^{121}$越來越近.
他非常憂愁.
所以你會看到.
他要準備去迎見.
他哥哥耳熟的時候.
在經文裡看到.
就是雅各繼續走.
但是.
遇見神的使者.
神的使者遇見他的時候.
雅各就說.
這是神的軍營.
好像.
神在鼓勵他.
看到神的使者.
於是.
他就說什麼呢.
我把那個地方叫馬哈連.
是什麼意思呢.
就是有兩營的軍兵.
怎麼會有兩營的軍兵呢.
原來一營就是.
他自己.
在這個時候他很需要.
面前這個使者的營.
就很厲害.
是來鼓勵他的.
幫助他.
好像賽跑的時候有人鼓勵你.
努力啊努力啊.
加油的時候心境很好.
你想想.
這個時候.
雅各就派一些人.
去告訴他哥哥.
他回來了.
誰知道.
收到消息.
哥哥和四百人在一起.
他又害怕.

$^{161}$很受煩.
你想想.
四百人做什麼呢.
我見你而已.
有四百人.
東亞哥一聽到就很受煩.
很害怕.
你看他怎麼部署.
他的計劃.
想起.
在第八節說.
首先他將同行的人.
分兩隊.
你會看到.
如果這隊不行.
還有另一隊可以走.
他想得很.
他是一個很厲害的人.
所以他想到很多方法.
然後呢.
他會祈禱.
他祈禱是怎麼祈禱的呢.
你會發現.
我們要學他.
雖然他這個時候.
真的不行.
但他想到他所擁有的一切.
他去感激神.
竟然施恩給他.
因為他本來.
是一支弓出去的.
什麼都沒有.
但現在好像滿載如歸.
有兒女成群.
有兩個老婆.
兩個婢女.
有很多的.
生畜.
他很富有 也是財富.
他承認.

$^{201}$他自己不配得.
這些都是神供應給他的.
神賜福給他的.
所以他向神感激.
看到神怎麼眷顧他.
然後呢.
他去求神救他.
因為現在.
哥哥和四百個人一起.
他真的很害怕.
他一想就煩.
所以他求神.
去救他.
在這個光景之下.
他說.
不單是脫離他耳掃的手.
他說因為我怕他來殺我.
連母親和兒女都不放過.
因為他現在有.
家室.
有老婆.
有兒女 這麼多的兒女.
是用什麼.
殺光他們.
他以前是一個人.
一個人無所謂 喜歡去哪裡都可以.
但現在有很多人.
不再是一個人.
他要去.
為他們著想.
他的現實的.
計劃是怎麼樣的呢.
你會看到.
神說他好像他爸爸.
好像亞伯拉罕.
天上星海邊的沙這麼多.
神祝福他 厚待他.
他現實的裡面.
他想著送禮物.
送禮物就行.

$^{241}$可以送這麼多東西.
這些是我財富.
我送去.
送禮物代表什麼呢.
阿姐以為送了禮物.
就可以為他開路.
送了禮物就可以滿和.
他和哥哥的怨恨.
你想想.
二十年的仇恨.
是不是送禮物.
就可以一筆勾銷呢.
有沒有可能呢.
送禮物.
當他送禮物的時候.
然後你會看到.
他就躲在最後.
你會說他很害怕.
是不是.
男子漢在前面.
他躲在最後.
同時.
他的角色.
改變了.
他以前.
什麼都要爭第一.
現在他竟然稱自己.
是僕人.
他在這裡說.
他自己的稱呼.
是和他哥哥.
是他的僕人.
你看到他.
開始有些轉變.
忽然間很謙卑.
他為什麼叫自己是僕人呢.
其實有個意思.
就是我很尊敬.
我的哥哥.
表示在態度上對他很尊敬.

$^{281}$和有個謙卑的意思.
有個態度.
去見他哥哥.
你會看到.
他雖然害怕到.
走在最後.
老婆仔都在前面.
你會看到很奇怪.
他自己是躲在最後.
雖然大家都很認識.
溫偉耀博士.
他有一個小女兒.
有很嚴重的智障.
記得有年.
他要移民到.
多倫多.
由於各方面.
條件很好 很快就被批准.
但是他的小女兒.
因為有一個病.
就沒有批准.
那怎麼辦呢.
解決的方法就是.
需要持一張移民報.
發出的部長特許證.
需要是每年續期的.
因為加拿大的政府.
不肯定他女兒.
有什麼需要呢.
醫療開支怎麼樣呢.
健康照顧怎麼樣呢.
不清楚 所以要看察她五年.
才會給她一個永久居留.
這五年.
每一次.
他就要拿這個特許證.
記住 特許證的文件是.
86.
到94年的時候.
就是五年.

$^{321}$那五年的時候.
因為這五年裡 其實他所有的開支.
都是自己去承擔的.
就算他女兒有多麼的緊要.
都不會用太多政府的資源.
想一想.
一切應該很順利.
可以永久居留.
當他在等待.
政府通知的時候.
忽然收到一個紅色打印的.
公函.
內容是什麼呢.
要他女兒三個月內就要離開.
大家收到這個消息怎麼樣呢.
原來呢.
三個月要離開就不可以再回來.
他們很沮喪.
事後發現.
有過綜的原因.
政府的內部.
搞錯了.
不知道哪個官員.
將文件86抄成80.
80就剛剛相反.
就要立刻離開.
不可以住在那裡.
他在英國出生的.
在香港.
又沒有住滿七年.
他去哪裡呢.
加拿大又不要他的話.
他去哪裡呢.
而他的病情又這麼複雜.
他十幾歲了.
很難搞的.
這個時候應該要申請他去特殊的宿舍.
但是他又沒有身份.
那怎麼辦呢.
等下再交代.

$^{361}$我們又看看.
雅各他的祈禱.
他在這裡看到他對他哥哥的稱呼.
就是要告訴別人.
他是僕人.
他這個時候.
就將自己交託給神.
他要打點一切.
於是他就帶著.
兩個妻子兩個婢女.
和子女就要離開.
去到.
這個圖裡面可能你看不清楚.
就是去.
雅各道口.
他帶著他們送他們過河.
所有一切就過去了.
雅各是激烈的一條小河.
現在的名字叫做.
阿卡河.
就在死海.
死海的樣子.
死海北面約三十公里的地方.
向西流就會流到.
約旦河.
先打發他們過河.
河谷的南邊就是高地.
雅各似乎.
已經使得他家人.
很平安的過了河.
他自己留下一個人.
可能想著休息一下.
或者安靜一下.
誰知道這個時候.
在聖經裡面.
當然會看回經文.
有一個人和他摔跤.
直到黎明.
我看回.
荷西亞書十二章四字那裡.

$^{401}$就說到.
他和天使角力.
並且得勝.
他曾哀哭.
懇求私人.
在不得利遇見耶和華.
耶和華在哪裡吩咐我們.
你看到他哭泣.
他一個這麼驕傲的人.
在那個時候.
他肯哭泣.
一定是很淒涼.
表示他.
極之軟弱.
聖經裡面說到.
那個.
使者摸了他一巴.
他就跛了腳.
瘸了腿.
這個人和他摔跤.
到黎明.
你也知道他年紀不小.
你不要想著他是一個少年.
或者是青年.
其實他年紀不小.
他生的兒子.
倒數來看.
他生的兒子已經91歲.
你想像一下.
他去到埃及見他的時候.
是103歲.
他去到那裡的時候.
其實他年紀不小.
他怎麼摔跤呢.
我們現在要和人摔跤.
都沒有那個能力.
你會看到有個人.
這個人就是耶華使者.
要和他角力.
但是你會看到.

$^{441}$還要去到天亮.
叫你整晚不睡.
你行不行.
那個人天亮的時候.
他死纏著.
不讓他走.
他雖然沒有力氣.
但也要死纏著他.
為什麼.
我現在要見我哥.
我要你給我祝福.
你不給我祝福我怎麼辦.
你不給我祝福我不讓你走.
雖然摸了他一巴的時候.
已經沒有力氣.
但盡他一生的能力.
他也要被我纏著.
這個祝福.
我就不讓你走.
所以可以顯見.
他感受到這個使者是從天上來的.
他很需要這個能力.
所以神祝福他.
然後他才可以.
面對哥哥和四百個人.
那個人就問他你叫什麼名字.
他說他叫雅各.
那個人就說你的名字不要再叫雅各.
要叫以色列.
以色列是什麼呢.
就是你與神.
在教力的時候.
是有得勝的.
所以以色列的意思就是.
他與神在教力.
那比洛伊勒.
就是神的面的意思.
雅各很需要神的祝福.
雅各過去.
他每樣都是靠自己的.

$^{481}$自己都可以搞定.
但誰知道.
他去到舅父那裡.
他去到阿山高.
他舅父比他更奸.
變了他要再痛苦的.
再做多七年.
才可以儲到他最喜歡的拉結.
人生就是這樣.
你以為自己很厲害.
有人比你更厲害.
他在摔角裡發現.
他無能為力.
但他在這個.
這麼要命的裡面.
天使跟他說.
天快亮了 你讓我走.
他竟然都承認.
他有個厲害的.
所以你會看到.
這個這麼奇妙的發生.
這麼奇異的能力.
其實神在那裡.
透過這件事.
去管教他.
你過去.
是靠你自己.
神要去管教.
他的驕傲.
管教他的固執.
他靠自己.
正正就是要他面對這個挑戰.
讓他去反省.
過去是怎樣的生活.
所以現在被人.
摸了一把的時候.
他跛了.
跛了還有什麼作為呢.
他發現他不可以再靠自己了.
其實神正正是要讓他看到.

$^{521}$你肯信服我.
你肯倚靠我 倚靠神.
神讓他看到.
如果你肯這樣做的時候.
你會很順利的.
甚至在崇靈的層面裡.
你肯聽神的話.
肯信服神的話.
你會很成功的.
你會得勝的.
所以在摔跤的裡面.
他死捉住.
這個使者不放.
因為他知道沒有希望了.
哥哥和四百個人.
就在面前了.
沒有希望了.
我不拿到祝福怎麼辦.
你會看到他的人生裡.
要學習服從.
怎樣服從神.
在祈禱裡.
他會看到.
我唯有靠主.
我不可以再靠自己.
人生裡過去想了很多方法.
很多計劃.
但這些計劃都行不通的.
所以只要.
人生裡他肯聽神的話.
就好像我們讀聖經.
我們抓住神給你的應許.
神很多應許給你的.
如果你肯去讀的時候.
這個應許就是你的了.
所以他抓住應許.
我們祈求.
就得獎了.
接著.
他又問這個使者.

$^{561}$你叫什麼名字呢.
但是那個使者說.
你何必問我的名字呢.
神沒有透露他的名字.
但好像當日.
神都沒有向馬羅亞.
透露他的名字一樣.
其實是在警告我們.
不要干涉神.
神有做事方法的.
他真的被他.
掌握了.
這個祝福.
掌握了這個祝福的時候.
他忽然間看到.
我面對面看到神.
還讓我可以生存.
他看到人生裡.
他開始反省.
他不像以前.
跑得那麼快了.
他現在是跛了.
跛了的時候.
當我們去祈禱的時候.
我們去求神幫我們的時候.
我們必須要同神和好.
很多時候.
我們是求一樣東西.
但我們有沒有反省.
我們去祈禱的時候.
我們裡面還有沒有很多自己呢.
有沒有很多罪在裡面呢.
就好像聖經說.
如果想弟兄.
對你遺憾.
那你怎樣.
將這個祭物留在壇前.
先去和你的弟兄和好.
然後你再去獻你的禮物.
同樣的.

$^{601}$他現在最大的困難是什麼.
見他哥哥.
他必須做一樣東西.
是什麼.
他騙了祝福.
人家追殺你.
你必須和他和好.
這個仇人在面前.
雖然過了二十年.
二十年你還記得的.
難道二十年你不還錢.
他不記得嗎.
所以你看到這個仇恨.
二十年.
他見他.
這個時候.
我們知道.
當我們願意和好的時候.
神就會承傳你的祈求.
這個時候.
你兒女.
他有兩個老婆.
兩個瀑被都有子女.
打發他們.
你們先走到一邊.
還有.
他跪在那裡.
連續.
趴著.
瀑在那裡.
一直走到他哥哥面前.
七次.
走得苦服的.
跪地.
不要的態度.
你看到.
人生的一面.
很難看到.
他竟然可以.
兩個抱頭痛哭.

$^{641}$二十年沒見.
一個仇人.
他很害怕的人.
他終日焦慮的人.
終於見到了.
這一幕.
他們兩個可以抱頭痛哭.
是一個很美麗的畫面.
這個畫面.
不是送禮物那麼簡單.
而是有神的祝福在裡面.
是神的恩典.
臨到.
才有這個畫面出現.
剛才說回.
溫博士.
當他走投無路的時候.
教會的信徒為他祈禱.
有人建議他.
既然你女兒是在英國出生.
可以長期居留的.
就是英國.
於是.
人們很好.
幫他打聽.
幫他留意.
很快就查到.
有一間社會福利署.
屬下的特殊照顧宿舍.
是青少年到成年人.
還有一個位置.
機會很好.
他就快些將所有的資料.
交齊了.
去辦理.
他覺得神開了一條.
不太可能的路.
請教會為他祈禱.
大家都覺得.
這是很積極嘗試的好方法.

$^{681}$他心裡也有些平安.
同時.
他那時暑假又要去北歐.
去帶一個聯合的.
退休會.
剛好可以過去.
看看環境.
一齊安頓好.
前一天.
那個宿舍的主管打電話過來.
說不用擔心.
有了電話更加放心.
於是他第二天就搭機.
搭了九個小時的機.
你知道他的女兒.
不是正常的.
很反常的.
她十幾歲的力量很大.
她不安靜的.
在飛機裡.
坐過飛機就知道.
九個小時都在叫.
兩夫婦.
左右夾著她.
按著她.
九個小時.
終於到達目的地.
下機第一件事.
快點去找宿舍.
去到宿舍招呼她的人.
就是跟她打電話的主管.
但在這個時候.
這個主管.
對他們不友善.
第一句說什麼.
你們立刻離開.
你女兒不可以住在這裡.
有沒有聽錯.
他們兩個很愕然.
這麼辛苦來到.

$^{721}$他竟然叫她走.
怎麼會這樣呢.
很奇怪.
為什麼叫我走呢.
然後.
再說半天.
她在坐飛機的時候.
社會福利署的一個官員.
來巡視這個宿舍.
又不知道為什麼.
這麼巧合.
檔案在板台裡打開.
官員就坐下.
隨便揭開.
揭到.
她是來自加拿大.
不是英國.
她沒有資格.
可以住在這裡.
你立刻.
不要讓她留在這裡.
立刻下命令.
於是他們沒辦法.
官員已經通令.
這樣了.
那件事怎麼辦呢.
他們捱了幾天.
都是被人趕走的.
她很痛苦.
她在祈禱的時候.
她快要去.
退休會了.
怎麼辦.
人家不收留她.
叫她走.
不接受她女兒怎麼辦呢.
她當晚.
很痛苦流淚的.
去跟神祈禱.
她也承認她的祈禱.

$^{761}$是很不客氣的.
她跟神說.
她說什麼呢.
她說我的老婆得了癌症.
我都沒有怪神.
我願意接受.
她現在這麼嚴重的病.
我都咬著牙奸去過.
但是為什麼.
你明明開了一條路給我.
現在又趕我走.
她很痛苦.
她說了一段時間.
她說神你好像.
左手打我.
右手又叫我為你做事.
為你去講訊息.
她跟著祈禱之後.
好像沒有什麼事發生.
她就跟著要去廣道.
長話短說.
雖然經過很多波折.
她回到加拿大的時候.
有調查員就來找她.
聽她說了這麼多的時候.
就很同情她的遭遇.
檢查員就說.
好我決定幫助你.
根據國際保護兒童會.
就是加拿大唯一一個.
不會計較你的國籍的.
一個收容的地方.
一個機構.
因為他們基於人道的立場.
就收留了她.
給了她一個宿位.
她一住.
就住到26歲.
就離開世界.
她是最後一個住這個宿舍.

$^{801}$原來她進去住的.
第二個月.
安德烈省.
這個政府就改組改選.
新的政府下令.
一切的宿舍都不可以再收留.
不是加拿大籍的居民.
無論怎樣.
都不可以接受.
當她回想起來的時候.
其實她女兒不夠資格.
她收容的要很極端的.
有災難性的.
有危急的就會收的.
但她回想起來.
重整的時候.
原來.
神給她又被人趕.
又被人追捕.
搞了這麼多事.
原來神的計劃真的很奇妙.
竟然讓她女兒.
可以安心住這麼多年.
她得出的答案是什麼呢.
她提醒我們.
問題就是我們信不信神.
是不會搞錯的.
是不會誤事的神.
所以萬事都互相效力.
叫愛神的人得益處.
大英姐妹.
無論你今天在什麼情況下.
你相不相信.
神仍然與你同在.
神有祂的方法.
我們不要問這麼多.
我們就去信服.
我們聽祂的話.
行不行呢.
你今天甘不甘願.

$^{841}$縱然你的環境很不順利.
你仍然.
會不會是忠於.
耶和華的神呢.
我們一起祈禱.
愛我們主 愛我們的神.
當我們.
走投無路的時候.
你讓我們在經文裡看到.
神你仍然是.
沒有拋棄我們.
你仍然愛我們.
祈求你幫助我們.
我們前景是怎樣.
我們不知道這個疫情.
好像有時候會發生 有時候是沒有的.
我們身邊很多事情發生.
經濟又不是那麼好.
又看到打仗又不停.
人的心裡有很多矛盾.
但是主啊.
求你幫助我們仍然是.
堅定的仰望你.
堅心的去依賴你.
求你幫助我們.
正正的仰望你.
讓我們一生裡.
我們不後悔.
我們所走的路.
我們知道神.
你愛我們 愛到為我們流血的地步.
我們都還沒到這個地步.
但是主啊 求你幫助我們.
不要忘恩負義.
幫助我們始終能夠.
跟主跟到底.
求主憐憫 求主祝福.
我們這樣祈禱交托.
奉主耶穌基督的名求 阿們.
\newpage



\section{創世記 32:22-32}
\label{sec:a1w0x_rlHHg}
\textbf{2024-03-03 午堂《得著與失去》創世記 32\_22-32 講員:陳麗蟬牧師}
\newline
\newline
連結: \href{https://youtube.com/watch?v=a1w0x_rlHHg}{\texttt{https://youtube.com/watch?v=a1w0x\_rlHHg}} ~~~~ 語音日期: 2024-03-06
\newline
\newline
\hyperref[sec:oyj7tRcaGj0]{\small{< < < PREV SERMON < < <}}
~
\hyperref[sec:index_chronic]{\small{[返順時目]}}
~
\hyperref[sec:index_scriptual]{\small{[返順卷目]}}
~
\hyperref[sec:2A_RiSlG6Bg]{\small{> > > NEXT SERMON > > >}}
\newline
\newline
創世記 32:22-32
\newline
\begin{longtable}{cl}
\hline
\hline
章節 & 經文 (和合本修訂版)\\
\hline
32:22 & \begin{tabularx}{0.7\textwidth}{X} 他夜間起來，帶著兩個妻子，兩個婢女和十一個孩子，過了雅博渡口。 \end{tabularx} \\ \\ \relax
32:23 & \begin{tabularx}{0.7\textwidth}{X} 他帶著他們，送他們過河，他所有的一切也都過去， \end{tabularx} \\ \\ \relax
32:24 & \begin{tabularx}{0.7\textwidth}{X} 只剩下雅各一人。有一個人來和他摔跤，直到黎明。 \end{tabularx} \\ \\ \relax
32:25 & \begin{tabularx}{0.7\textwidth}{X} 那人見自己勝不過他，就摸了他的大腿窩一下。雅各的大腿窩就在和那人摔跤的時候扭了。 \end{tabularx} \\ \\ \relax
32:26 & \begin{tabularx}{0.7\textwidth}{X} 那人說：「天快亮了，讓我走吧！」雅各說：「你不給我祝福，我就不讓你走。」 \end{tabularx} \\ \\ \relax
32:27 & \begin{tabularx}{0.7\textwidth}{X} 那人說：「你叫甚麼名字？」他說：「雅各。」 \end{tabularx} \\ \\ \relax
32:28 & \begin{tabularx}{0.7\textwidth}{X} 那人說：「你的名字不要再叫雅各，要叫以色列，因為你與神和人較力，都得勝了。」 \end{tabularx} \\ \\ \relax
32:29 & \begin{tabularx}{0.7\textwidth}{X} 雅各問他說：「請告訴我你的名字。」那人說：「何必問我的名字呢？」於是他在那裡為雅各祝福。 \end{tabularx} \\ \\ \relax
32:30 & \begin{tabularx}{0.7\textwidth}{X} 雅各就給那地方起名叫毗努伊勒 ，說：「我面對面見了神，我的性命仍得保全。」 \end{tabularx} \\ \\ \relax
32:31 & \begin{tabularx}{0.7\textwidth}{X} 太陽剛出來的時候，雅各經過毗努伊勒，他的大腿就瘸了。 \end{tabularx} \\ \\ \relax
32:32 & \begin{tabularx}{0.7\textwidth}{X} 因此，以色列人不吃大腿窩的筋，直到今日，因為那人摸了雅各大腿窩的筋。 \end{tabularx} \\ \\
[1ex]
\hline
\hline
\end{longtable}
$^{1}$請姐妹平安.
不知道大家有沒有聽過一個故事.
講一個小朋友.
他的手伸進了玻璃瓶裡.
抓著一堆糖果.
當他想拿出來的時候 拿不到.
他媽媽就跟他說 你要拿出手來的話.
一定要放開手才能拿出來.
但他不肯 因為他要抓著這些糖果.
他很想吃這些糖果.
但抓著糖果的時候 拳頭就變大了.
瓶子就只有這麼多.
他怎麼也拿不出來.
這個故事告訴我們.
有時候我們抓著一些東西.
沒錯 我們很喜歡 很想得到.
但當我們不肯放手的時候.
我們的手就沒有自由了.
永遠就在瓶子裡.
在人生裡有時候也是這樣的.
我們抓著一些東西.
覺得我們自己很喜歡 很重要的.
但抓著的時候.
卻成為了我們生命裡很多的捆綁.
不放手的時候.
我們沒有自由.
我們生命得不到釋放.
今天我們講這個雅各的事情.
講他跟他的太太和11個子女.
一起回鄉的時候.
我們彼此有一個勉勵.
雅各這個人.
先講他的名字.
他的名字是抓著的意思.
人如其名.
他很喜歡抓著一些事情.
我們看到他是一個很厲害的人.
很有智慧的人.
但他就經常抓著這些智慧.
他就用他自己的一些力量.

$^{41}$方法來抓著他認為很重要.
又是他很喜歡的東西.
他經常以為自己很厲害.
因為他真的做事都很厲害.
處處都要走來爭第一.
我們看到他一出生的時候.
明明跟哥哥是雙胞胎的.
哥哥已經出生了.
但他不值得 為什麼我要第二.
他就嘗試抓住哥哥的腳.
把他推進去 是不是.
當然不成功了.
到他年少的時候.
他又接著一碗紅湯.
剛剛哥哥打完獵回來.
很口渴 很累.
他說好 你想喝我這碗紅湯都可以.
你給我長子的名分吧.
哥哥可能覺得.
哎 名字而已 名分而已.
就給了他.
有沒有用都好 總之他就抓住.
再大一點的時候.
他甚至搞到.
就把長子的祝福.
應該給哥哥的祝福.
他就用計謀來跟媽媽合作.
就搶了過來.
這個長子的祝福.
原本給二嫂的祝福.
就有土地肥沃.
在舊約的時代很厲害.
耕種土地的肥沃.
還有每一次種植都會五穀豐收.
哇 真是發達了.
種了東西都有豐收.
然後還有他說.
萬民萬族都會事奉你的後代.
萬民萬族都要拜他的後代.
就是說在這個長子的祝福裡面.

$^{81}$他的後代是非常有尊榮.
後代是非常有權威的.
但是他跟他媽媽就用這個計謀.
騙了這個祝福過來.
然後他就逃走了.
去到他的舅父的家裡面.
他雖然被舅父騙了.
這個舅父被他更加的證明.
騙了他做了14年的免費勞工.
他當然不放過你了.
所以他就用他的計謀騙了舅父.
那些純白色的羊群屬於舅父的.
那些黑羊或者有斑點的羊就屬於舅父自己的.
他就用這個方法.
使舅父那些純白的綿羊瘦了一點點.
而他自己那些斑點,黑羊就每一隻都肥脹脹的.
讓舅父知道的時候.
當然又不肯放過他.
但是我們看到這個雅各.
由他出生.
除了這個出世出不了第一之外.
其他的東西他去搶的時候.
他都可以成功搶得到的.
搶了一個長子名分.
搶了這個原本給耳掃的祝福.
搶了舅父那個封口的家財.
但是他得到這麼多東西.
他同時都失去了很多的東西.
他失去了跟爸爸跟哥哥的那個關係.
所以這個哥哥很生氣.
覺得你搶了我長子的祝福 是吧?.
所以他就跟他爸爸說.
他要追殺雅各.
他也失去了他的誠信.
在爸爸裡面來說.
覺得為什麼你把羊毛弄在頸上.
來騙取他的祝福呢?.
其實他也失去了跟他媽媽的時事在一起.
他媽媽很疼他.
因為他怕被哥哥追殺.

$^{121}$所以他媽媽教他北上去舅父那裡.
直到在這十幾二十年的裡面.
他不再見到媽媽.
媽媽死了他都不能在他身邊.
我們就看到了 雅各在他前半生的裡面.
他揹著很多他自己內心的慾望.
他要爭第一 爭好的物質.
爭好過那個的名分.
但結果來他卻失去了親情.
失去了誠信.
那究竟什麼是「得」?.
什麼是「失」?.
是怎麼樣來定這個案呢?.
我們來看看 在這個雅各爭的事情裡面的時候.
他騙取了這個長子的名分和長子的祝福.
當他知道哥哥要殺他的時候.
他就逃走了.
就去到北面的哈蘭的時候去舅父的家裡.
但當他逃走 去到一個地方.
在那個地方 他就已經很晚了.
不能夠再走了.
他就睡在地上.
拿一塊石頭當作他的枕頭.
想一下雅各當時的心情.
睡在這個荒山野嶺.
想一下 雅各會不會追上來呢?.
在那個時間 也會覺得很孤單的.
因為他將要去到一個地方.
雖然是去到他舅父的家裡.
但始終離開了他熟悉的家鄉.
離開了他的親人 爸爸媽媽和哥哥.
去到舅父的家裡 他已經沒有什麼親人了.
可能他也會感到徬徨.
感到前路茫茫 不知道該怎麼辦.
我想可能我們曾經出外國讀書.
自己一個人去外國讀書的時候.
可能我們也會經歷這一份心很混亂.
前面的日子是怎麼樣的呢?.
但在這個時候他很落寞 很孤單.
就是人在這個時候 感到無助.

$^{161}$感到沒有什麼打壓的時候.
上帝就會向人顯現.
所以在那晚的做夢的裡面.
上帝就讓他看到一條梯子很長很長的.
長到直達天上高的.
除了看到這條這麼長的樓梯之外.
他也看到有一個天使在這條樓梯裡面.
上上落落 上上落落的.
上帝就給這個眼角要知道.
這一條梯子是象徵著上帝和人的溝通.
在天上高 和人在這裡溝通.
而那個天使就象徵著.
這個天使落落落到地上高.
他看到地上高人的需要.
人的境況的時候.
他就從這條天梯一直上去.
就報告給上帝知道.
然後上帝就會按著每一個人的境況.
施給他們恩典.
將這個祝福臨到這個地上高.
上帝就在眼角.
在這個人很落寞 很孤單.
很前路茫茫的時候.
讓他知道上帝會與他同在.
上帝會帶領他的前路.
上帝會保護他的.
所以阿葛在這裡經歷了那件事之後.
他起來了 他心裡感到很安慰.
他就改了一個地方的名字叫做「伯特利」.
「伯特利」的意思就是上帝的居所.
上帝的殿.
在村口那裡有一座伯特利的中學.
所以上帝不只是在那裡.
而是這個意思就是說.
上帝在我們每一個人生命當中.
他時刻的看著我們.
在這裡上帝藉著伯特利的經歷.
他要慢慢讓這個眼角.
在他的人生觀 在他價值觀.
在他信仰的裡面.

$^{201}$都要去改變.
他要使眼角重新有一個新的方向.
不過在那裡他只學到一點點.
第二次就是我們剛才讀經文的裡面.
進去第二次很艱難的時間.
這次的裡面.
就因為他和他的舅父不相伯仲的.
同樣是一個很狡猾的一個人.
在那裡他就走了出來.
這次他就帶了他的太太.
他有兩個太太.
就是因為他的舅父騙他.
他很喜歡舅父的小女兒.
但是他舅父說.
行 你幫我做七年事.
我就會將她嫁給你.
為什麼做完七年之後.
嫁給他的太太.
原來是大女兒.
舅父就說.
我大女兒都還沒出嫁.
怎麼可以小女兒出嫁.
所以他又再做多七年事.
才可以將他自己很喜歡的女兒.
娶回來.
所以那次就帶著兩個太太.
還有11個的子女.
還有很多的牛啊.
很多的羊啊.
很多的駱駝啊 驢啊.
那些婢女啊.
一路這樣養回來.
就因為他當時.
他這個就是因為他騙了他的舅父.
舅父的羊全部都瘦了.
他自己的羊就肥嘟嘟的.
舅父就覺得.
哇 你繼續在這裡的時候.
不知道會再被他騙到什麼.
所以就說你走吧 你走吧.

$^{241}$就叫他走了.
他唯有就回到他的家鄉.
當他一路開始回去.
去到的時候呢.
就去到他哥哥的地方.
他的家鄉.
他都很聰明了.
哇 在十幾二十年之前.
他哥哥曾經說過.
一定要殺他的.
所以他就走去見一下哥哥.
找侍者去見一下哥哥.
打聽一下形勢.
就吩咐了那些僕人.
你看到哥哥的時候.
記住要謙卑一點.
說呢 稱呼哥哥做主人.
你說我們只是個僕人.
接著呢 派去的那些僕人就回來了.
哥哥就很心急.
怎麼樣 怎麼樣.
哥哥那方面有什麼行動.
那些僕人就告訴他.
你哥哥帶著四百人走過來.
哇 他就很害怕了.
他就很害怕了.
在第32章第7節裡面.
聖經說到.
雅各很懼怕.
並且受罰.
死了.
因為呢.
他哥哥是打獵人來的.
很大隻的.
很有力的.
但他自己是一種用路人.
他最厲害的就是有他計劃策劃的東西.
不能打的.
還有呢.
他對方帶著四百人來.

$^{281}$一定是尋仇的了.
還有這次呢.
他帶著兩個老婆.
還有所有的十一個子女在這裡.
然後接著呢.
他所有的家財都在這裡了.
所以呢.
他就先將他們打發過河.
這次怎麼辦呢?.
他就很辛苦.
很辛苦在這裡.
他又是自己一個人.
他擔心哥哥這次將他一鋪盡了.
殺光他的家人.
拿光他的財產.
甚至他自己性命都不保.
過往曾經騙過人.
搶過人.
心裡面自然會有想著的想法.
對方一定走來去報仇的了.
陳慎之先生.
他曾經得過香港十大青年傑出獎的.
他之前是吸毒了十幾二十年的了.
但是呢.
他信了耶穌之後.
改變了.
不僅沒有再做以前那些偷騙拐騙的事.
甚至在社會裡面有很多的貢獻.
所以政府才頒了一個十大傑出青年獎給他.
當他完全改變了之後.
一次逛街的時候.
看到後面有幾個人在那裡.
噗啦啦的向著他衝過來.
他說 糟了.
以前得罪過什麼人.
你現在是不是走來尋仇呢?.
於是他不管三七二十一.
他又走了.
走走走.
走著走著.

$^{321}$咦?.
沒有人追上來的.
他看著看著.
原來那幾個人在追的士.
在上的士.
原來以前你做過一些事情.
你在心靈裡面.
你都會認住.
哎呀 尋仇啊 尋仇啊.
這次亞國也是.
覺得哥哥一定會來追殺他的.
所以在那晚的裡面.
他想辦法 想不到.
他很害怕 很受煩 很孤單 很徬徨.
在那次裡面.
上帝又出現了.
上帝這次就藉著一個人的樣式.
來出現.
就不是再讓他看到天使上上落落了.
這次要跟他打架.
跟他摔角.
在摔角裡面 從晚上一路到天亮.
亞國一直都不肯放手.
一直在角力.
到天亮的時候.
他才發覺到.
這個不是人來的.
有兩隻大翼的.
是天使來的.
在摔角當中.
因為亞國一直都不肯放手.
這個使者就在大腿窩那裡.
摸一摸他.
就這樣摸一摸的時候.
他就失去了力量.
從此之後.
他走路就拐下拐下.
這個大腿窩是哪裡的地方呢.
雖然在上網看到很多圖片.
但我總是搞不清到底是哪裡來的.

$^{361}$唯一一個我自己都明白的.
就是說這個大腿和腰那裡.
和腰是在希伯來文是同一個字.
就是在腰那裡附近.
這個窩就去到腳底那裡.
總之就是整隻腳那裡.
就指到一個人支撐的最重要的步伐.
大家都差不多跟我年紀差不多的.
你會發覺有時候在天氣轉變.
或者下雨的時候.
你會發覺膝蓋會軟一軟的.
當你走路走著走著.
或者有時候想穿幫的時候.
當你的膝蓋軟一軟的時候.
你會怎樣呢.
沒有力氣了.
很容易就會慣了.
現在整隻腳從腰那裡一直到腳筋那裡的時候.
是沒有力氣的.
沒有力氣的.
這個天使只是摸一摸.
他的腰或者是屁股那裡的一部分的時候.
他的力量就沒有了.
我們先看一下.
先看一下上帝這個摸一摸.
拍一拍的時候有什麼意思.
在聖經裡面其實都出現了好幾次.
就是當他要人有一些改變的時候.
他會摸一摸他.
拍一拍他.
我們看耶穌在醫病的時候.
某一次當他見到某一個麻風病人的時候.
他問那個麻風病人.
他問你想我為你做些什麼呢.
這個麻風病人說.
我想得醫治.
於是耶穌就摸一摸這個麻風病人.
他的皮膚就立即好了.
其實當耶穌去醫這個麻風病人的時候.
他可以不用摸他.

$^{401}$說麻風病人又會傳染的.
那些皮膚又這麼肉酸.
他只要說一句.
你得醫治了.
但耶穌卻是要摸一摸他.
那一次耶穌醫人的時候.
不單止要醫好人的身體健康.
他更加要去醫治人心靈的創傷.
因為這些麻風病.
他們不可以在城市居住的.
他只可以去到郊外的地方.
還有當他見到有人在郊外走的時候.
他要大聲的說.
我是麻風病人的.
是不乾淨的.
你們千萬不要走近我身邊.
所以那些麻風病人很慘的.
一方面身體裡面很不舒服.
而另一方面不斷的告訴別人.
我是不乾淨的.
我是不好的.
心靈裡面其實真的很難過.
也很孤單.
所以耶穌這樣摸一摸他的時候.
讓病人知道.
他是接受他的.
他是憐憫他的.
所以這次使者和雅各角力的時候.
去拍一拍他的大腿窩.
叫他去軟一軟.
摸力.
他的目的是要叫他.
這個力量沒有了.
但他不是想弄死他.
而是很想這個雅各的生命.
可以能夠得到改變.
得到更新.
雖然身體的力量是少了.
但是他的生命才能夠得到真正的釋放.
因為過往雅各的生命是怎樣的呢?.

$^{441}$他說他處處都要爭第一的.
出世也要爭第一.
排名排長又要爭第一.
他是很驕傲的.
他是一個很自恃的人.
因為他做事又真的挺醒目的.
他凡事都是用自己的努力去爭取.
爭取他自己覺得應該有.
或者他自己喜歡有的東西.
但在爭取的過程裡.
往往就不理其他人的感受.
不理其他人的死活.
他爭取了這個長子名分.
爭取了這個祝福的時候.
他有沒有理會到他哥哥的感受呢?.
當他騙了他爸爸.
用羊毛劑在頸上.
他爭取了爸爸的祝福.
爸爸知道被騙的時候那種憤怒.
自己的兒子是這樣的時候.
是多麼的心痛呢?.
他有沒有想過呢?.
但是呢.
在他當日被這樣摸一摸的時候.
他整個人軟了 沒了.
並且從此之後.
他的腳就要走著走著.
拐著拐著.
但是就被他看到.
他自己那份軟弱.
他自己的限制.
他從此以後.
他都需要僅靠著上帝.
僅靠著上帝的那個祝福.
所以我們真是.
很多時候我們都會遇見.
甚至可能我們自己都是這樣的人.
我們就會靠著自己的力量.
靠著自己的經驗.
意志.

$^{481}$去覺得我要那個.
第一的生命 第一的生活.
但當我們到處去這樣挖的時候.
這樣去拿的時候.
就會很容易就一些的詭詐.
一些的自我.
一些的批評 一些的挑剔.
就會出現的時候.
我們就.
接著我們就會容易跟別人一些的不和.
我們要去爭.
要用我們自己的手段去爭.
但是當上帝要我們的生命有改變.
要釋放出來的時候.
他們要先經歷.
就是我們自己的軟弱.
我們自己的無能.
我記得有一次.
在這個.
那時候我都還沒有讀神學的.
就在團契裡面.
我們就說不如看一下這個.
《使徒行傳》.
保羅三次的宣教旅程.
當時就找了一組的弟兄姊妹.
她就負責就去研究那三條的旅程.
怎麼走.
每個旅程裡面做了些什麼.
但是給了兩個月的時間.
到最後做不出.
那後來我們開支援會的時候就說.
那組人做不到了.
不如找別人做吧.
我心靈當時覺得.
哇 差不差一點.
做了兩個月都做不到.
不知道為什麼那些職員會.
就說不如你做吧.
就想一下.
我讀歷史的嘛.

$^{521}$有些東西這麼簡單.
easy job 是不是.
那我心靈就說 好.
那我就做了.
人的驕傲.
誰知道結果.
我兩個月都做不到.
都做不到.
在裡面很驕傲.
覺得我一定可以的.
但是上帝就讓我看到.
你覺得你自己可以.
你看小人.
雖然我沒有出口傷人家.
你做什麼 你這麼差.
但是我心裡面已經說了.
你做什麼 你這麼差.
在那次裡面給我的經歷就是.
凡事不要憑著自己的經驗.
自己覺得那個學歷我可以.
就以為可以.
一定要依靠上帝.
在《中信月刊》裡面.
第742期裡面.
有一位姐妹.
她說見證.
她說她是一個.
重男輕女的家庭裡面長大的.
她一直以來.
她是一個很固執的爸爸.
一個很權威的媽媽.
一直壓著自己.
所以她一直成長.
她就覺得.
為什麼這個家庭裡面.
做女的就沒有地位.
為什麼做女的.
很多事情都不能夠做到.
心裡面藏著很多的苦澀.
以至於她的性情.

$^{561}$就變得很多的批評.
很多的挑剔.
對人很多都看不上眼.
她也對自己.
這個原生的家庭裡面.
她覺得.
她要脫離這個原生的家庭.
所以很快就結婚生孩子.
她人生的目標.
就是要能夠打理家裡很好.
照顧孩子很好.
就是人生的目標.
她覺得她以前的爸爸媽媽做不到.
害得她很辛苦.
所以她成家立室之後.
她就整理家庭很好.
雖然是這樣.
她心裡發現了.
她的丈夫有第三者.
她的丈夫說.
因為你只顧著兒子.
你都不疼我.
我出去找另一個吧.
她就覺得很生氣.
為什麼我為了這個家庭.
擺上這麼多.
但竟然我丈夫去抱怨我.
批評我.
並且在外面找第二個呢.
她最初就出去找打小人.
打小人打了她的丈夫周圍的蝴蝶.
就打走他們.
但越打的時候她越不安.
心裡的心靈很沉重.
也覺得自己這麼失敗.
她覺得她打理家庭這麼好.
為什麼最後會這樣呢?.
有一次有同事帶她去教會.
牧師就說到耶穌的一句話.
凡是勞苦擔重擔的人.

$^{601}$到他們裡面去就得到安息.
她覺得生命的裡面.
在她原生的家庭裡面.
她的心靈很沉重.
覺得父母不疼她.
女兒父母一直壓制著她.
到她家庭裡面.
本以為自己出力了.
自己做得很好.
但竟然丈夫是這樣.
她覺得很辛苦.
但因著耶穌的說話.
凡是勞苦擔重擔的人.
到他們面前就可以得到安息.
她就將這些擔子.
被欺壓的擔子.
被丈夫抱怨的擔子.
交給了耶穌.
她心裡說.
當她信了耶穌之後.
她過往的傷痛得到醫治.
她接納了自己的父母.
雖然她不可以再和丈夫在一起.
但她原諒她的丈夫.
她也和一些.
大家之前不是很好的同事.
他們再一次相交.
更加重要的是.
她也接納了自己.
所以在我們的信仰裡面.
上帝不是要我們.
真的拿很多很多的第一.
當然他要我們都要來努力.
但他不是要求我們.
什麼都要拿第一.
在我們的信仰當中.
是要我們.
接納我們自己的警方.
看我們自己合乎忠道.
而不是全部都要抓住.

$^{641}$更加重要的是.
去依靠上帝 信服上帝.
在這個當雅各.
他自己在這個河裡面.
那條河叫做瓦博河.
在那裡被上帝去摸一摸.
他整個生命由以前很強的.
從此以後變成武力軟弱的.
但那次之後.
他不再搶著地上高的地位財富.
他緊緊抓住的是上帝的祝福.
上帝給他的祝福.
賜給他一個新的名字.
在聖經裡面.
我們看到上帝時時都幫人改名.
原來是一個祝福.
我們看看雅各的爺爺.
亞伯拉罕.
他原本就叫做亞伯蘭.
但和上帝相遇之後.
就改了叫做亞伯拉罕.
亞伯蘭這個名字都很好的.
是一個尊貴的人.
但跟著他和上帝相遇之後.
他的名字就叫做亞伯拉罕.
這個名字更加厲害.
就是這個多國之父.
上帝祝福亞伯拉罕.
他將來的後裔子子孫孫.
成為了很多很多民族的源頭.
我們看看新約的時候.
這個耶穌基督的門徒.
他的大弟子.
本來叫做西門.
西門就是聽.
聽上帝的話.
也不錯的.
跟著他見到耶穌的時候.
耶穌就幫他改名.
叫做基伐.

$^{681}$在希臘話裡面就叫做彼得.
是我們很熟悉的一個耶穌的門徒.
這個基伐或者彼得的意思.
就是石頭.
石頭沒有用的.
不是呀.
這個石頭很珍貴的.
就是成為了教會的柱石.
所以彼得在《仕途行傳》我們就知道.
他就是成為了教會的領袖.
原來上帝給我們新的名字的時候.
就是給我們生命的改變.
他就叫這個雅各.
不是再是一個周圍地上高物質.
地上高的名份的東西.
以後他就要改為以色列.
成為了一個同神一起來治理.
以色列這個名字.
是為什麼呢?.
因為以色列就說.
他是一個王子.
是和上帝一起來治理一些事情.
名字是很好的.
這個名字其實就是一個祝福.
一個新的保證.
是一個記號.
是和他過往不好的東西一筆勾銷.
我記得我曾經說過.
說過好幾次了.
我以前小時候很討厭我的名字.
陳麗禪 陳麗禪.
那個禪字為什麼我這麼討厭它呢?.
就是因為是松字邊.
很醜怪的.
還有夏天的時候最討厭的.
在樹上吵得很厲害.
小時候不懂事的時候我很討厭它.
我曾經很想改自己的名字叫陳麗芝.
但我又不懂那個麗芝的麗字.
跟這個麗芝不一樣.

$^{721}$因為麗芝是紅色的.
還有那時候小時候很喜歡吃它.
但是長大後我相信了主.
又知道這個禪的生命的時候.
我又很喜歡這個陳麗禪這個名字.
所以我看到松字邊的禪的時候.
我就會兩眼發光.
原來這個禪我講過了.
它出生在地上的泥濕漉漉的時候.
要留多久呢?.
16年 17年.
然後它才能夠飛上樹枝上.
飛上樹枝裡面.
夏天的時候你會聽到很吵的聲音.
很難聽的.
但是它的生命只有60天到70天.
一過了兩個多月之後.
它就死了.
它出生在樹上吵得很厲害.
除了吵之外還做什麼呢?.
原來它的責任就是繁衍下一代.
在那裡有產卵,生孩子.
那些卵就下在樹上.
又16,17年才什麼.
所以沒什麼用.
不過我才知道從中醫裡面.
那個禪的退學很有用.
有什麼用我也不知道.
然後我相信了之後.
原來在生命裡面雖然只有很短暫.
但是在短暫的時間裡面.
我不斷地去宣揚耶穌基督的救恩.
入學講見證的時候.
之後有一個師兄.
給了我兩句話.
「厲日盡見神恩典.
禪明也道主恩心」.
就是鼓勵我.
在我日後的生命裡面.
我要不斷地去宣揚耶穌基督的福音.

$^{761}$當我按目的時候.
有一位弟兄又畫了一幅畫給我.
我真的很寶貝他.
那隻禪你看到嗎.
那隻禪逆透明的畫得很好.
然後下面就有一連串的果子.
祝福我在牧羊事奉的生命裡面.
可以不斷地為主耶穌基督去結出果子.
原來新的名字.
給人真的為之一振.
所以以色列.
到亞國.
他就從此知道.
他不是再為自己畫一些東西.
他是上帝的兒子.
他和上帝一起的童工.
所以亞國的生命裡面.
他在過往的裡面.
他真的得到很多東西.
但這些東西對他來說.
都是沒有益處的.
但在他第二次很孤單的時候.
他和使者裡面來角力.
在經文講.
他說他贏了這個使者.
但他不是真的這麼厲害.
贏了這個使者.
使者小小力拍一拍他.
他就那個大腿骨扭一扭.
他就沒力了.
而是講到.
他贏了他自己.
因為過往他抓緊一些東西的時候.
他的生命是不開心的.
很孤單的.
和人的關係都不好.
但當他認清了.
自己的軟弱.
自己不行的時候.
他投靠.

$^{801}$他抓住的是上帝的祝福.
是抓住上帝的祝福.
其實這段經文也讓我們知道.
我們整個生命.
是要先求他的國和他的義.
我們先以耶穌基督成為了我們生命的中心.
我們凡是要去決定任何的事情.
計劃任何的事情的時候.
我們都要求問上帝.
不要只顧著為自己最好的安排.
在地上高的無論事情.
愛情.
錢財.
只要我們尊耶穌基督為主的時候.
他會為我們去安排妥當的.
所以我們看到那次.
他的哥哥不是真的去追殺他.
相反他當國華送了很多禮物.
帶了很多禮物送給你.
他跟弟弟說.
不用了.
我自己都有.
你自己留著用吧.
並且他們兄弟兩個一起抱頭.
彼此的接納.
一幅很漂亮很漂亮的圖畫.
原來我們尊主偉大的時候.
我們放下我們自己要畫的東西的時候.
很美麗的圖畫.
和諧的圖畫就會出現.
我們一起土谷.
天父上帝我們多謝你.
多謝你.
原來生命裡面有你成為了我們的中心.
有你成為了我們給我們的依靠的時候.
我們的生命可以真正得到自由.
真正得到很多很多的福樂.
祝福我們.
如果我們仍然有些東西抓緊我們不肯放手的.
箍得我們自己很辛苦的.

$^{841}$求主聖靈幫助我們提醒我們.
以致我們很清楚知道.
在主的裡面.
我們有自由有力量.
幫助我們.
幫助我們能夠全部我們所有的人.
生命都以你為中心.
我們多謝你.
順天禱告奉主基督的聖名.
阿們.
\newpage



\section{瑪拉基書}
\label{sec:2A_RiSlG6Bg}
\textbf{2024-03-07 周四晚查經班《瑪拉基書(9)》 講員:曾裔貴牧師}
\newline
\newline
連結: \href{https://youtube.com/watch?v=2A-RiSlG6Bg}{\texttt{https://youtube.com/watch?v=2A-RiSlG6Bg}} ~~~~ 語音日期: 2024-03-08
\newline
\newline
\hyperref[sec:a1w0x_rlHHg]{\small{< < < PREV SERMON < < <}}
~
\hyperref[sec:index_chronic]{\small{[返順時目]}}
~
\hyperref[sec:index_scriptual]{\small{[返順卷目]}}
~
\hyperref[sec:AndPB0XPtoQ]{\small{> > > NEXT SERMON > > >}}
\newline
\newline
$^{1}$我們能夠用了很多的時間.
將歸回的歷史書和先知書.
今晚都是最後一章的馬拉基書的第四章.
雖然只有六節 但都給我們有很多的啟發.
尤其是去到那個日子.
耶和華大而可畏的日子.
都給我們能夠定精在神的國度.
一定是會成就的.
雖然未必在我們這一代去成就.
但是都給我們能夠有一個敬虔的心靈.
這樣來等候.
今天晚上我們仰望主.
給我們能夠有充足的精神.
有心靈的啟發.
也盼望信息能夠幫助很多的人.
我們仰望禱告奉耶穌基督的名.
阿們.
第四章和之前的三章的馬拉基書.
明顯有一些不同.
其中你發現到.
不再是一個提出一個問題.
接著自己解答.
來講百姓有什麼問題.
第四章這裡明顯主題就很清楚了.
耶和華的日子.
有兩個大重點.
就是提了很多次那天那天那天.
到那天提了好幾次了.
第二第三第四節都有提.
另外就是在這個日子來之前.
是會有一個徵兆的.
就是有一個人.
就是一個預備那個大日而來的使者.
就是第五節和第六節所講的.
這一位的施洗約翰.
因為新約當引述這段經文的時候.
是主耶穌去為我們解釋.
就是指到施洗約翰了.
所以這段聖經最後這一章.
就是作者一個很終極的宣告.

$^{41}$就是第三章我們上個禮拜講了.
有正反其實應該先講了反面.
就是那些不敬畏神的人.
最終都是要面對很大的審判.
但是三章的十六節開始.
就是有一些聽到馬拉基先知的預言.
和教訓而來進行的人.
他們會彼此談論耶和華的事.
就是說對神是很有興趣.
對認識神是很有抱負.
上個禮拜叫大家背.
那時敬畏耶和華的彼此談論.
在神的面前有一個紀念冊.
神是會紀念.
當然說笑了 神是不需要儲紀念冊.
對我們每一個人的生命.
祂是完全的掌握認識和明白.
也是掌管我們每一個人的出生.
那個成長.
和我們一生的人生都是在神的手中.
而且是一個美好的安排.
這個就告訴我們.
其實就是凸顯那個時代.
不是每一個猶太人.
都敬畏神進行摩西的律法書.
去到四章四節.
又再一次提摩西律法.
所以都是不斷的婦口婆心.
提到猶太人真正的認識神.
需要透過讀聖經.
不是靠傳統.
更加不是靠一些拉比的解釋.
而是需要摩西的律法書.
在自己的心靈裡面.
不斷的有聖靈的啟發.
又不斷的遵行而帶來的生命改變.
記住 舊約我們一而再都講過.
即使保羅這些這麼厲害的法律專家.
對摩西的律法 對舊約是這麼精通.
但是可以是讀死了那本舊約的.

$^{81}$當然我們要原諒的.
就是耶穌還沒有出現.
所以舊約都是很多東西好像蒙了一層沙一樣.
你只是守了字面.
你都以為自己是做了的.
但是的確有一些整個的舊約時代.
神是揀選一些有信心的先祖 先知.
或者一些跟隨神的人.
所以你可以想像到.
如果沒有新約的出現.
舊約就好像一個很大的一個牢籠.
令到人以為只需要按時按後的守律法.
守節旗 憲制 或者民間不要有太大的犯罪.
就已經是蒙神越立了.
但是原來我們發現.
去到新約主耶穌的降臨.
就是另外一回事了.
是真的去到一個很大的心靈上面的.
和神之間的接觸.
所以保羅的發言.
你可以說是一個很大的晴天霹靂.
他從來沒有想過自己是離開神這麼遠的.
他是一個教法師.
他是一個法利賽人.
他是一個拉比.
他是一個能夠精通舊約.
尤其是他引述的《二次亞書》.
詩篇 保羅引述很多的.
但是發現原來他在主耶穌的眼中.
在主耶穌的救恩裡面.
原來他一點都不認識舊約.
所以他當見到主耶穌.
他的生命就有一個很大的改變.
他整個人謙卑下來.
他要來到去再一次反思.
到底新約的這位耶穌.
是在舊約裡面怎樣來去預言呢?.
怎樣來去講述這位彌賽亞呢?.
所以保羅就變成一個對舊約重新去審閱.
重新去解釋的一個專家.

$^{121}$所以他的書信變得分外重要.
就是有這樣的原因.
你可以想像到.
神用一個這麼出色的人.
來解釋新約和寫新約.
使得我們後世的人.
對舊約都會多了一份的了解.
因為保羅是站在我們新約外邦人的角度.
幫我們寫這些書信.
所以我們就能夠對舊約.
我們一竅不通的外邦人.
我們能夠多了一份的了解.
所以我們讀通了保羅的書信.
是不會使得我們要死守律法.
因為讀通保羅的書信.
會使得我們知道是在乎零.
不是在乎疑問.
這是哥林多後書的四章的經文.
所以馬拉基書一而再講.
那個標準.
你們離開那個標準是什麼呢?.
其實馬拉基書如果我們濃縮四章來說.
其實就是他們離開了摩西的律法書.
過去很大一部分的先知警告.
就是他們要回到摩西的律法書來進行.
這樣就會得到地土的祝福.
五穀的豐收.
生產的無論牛羊.
都會有豐收.
或者不會有天災.
不會有一些收成的問題.
但是當他們不肯去進行.
神就會透過先知的警告.
也會興起一些外邦的列國.
來審判以色列.
到了中後期.
尤其是到了最後的耶利米的年代.
就是講一些不再是悔改.
就可以得到拯救的教導和預言.
耶利米大部分的預言.

$^{161}$就是說不用悔改了.
神都不閱立了.
你們現在應該怎麼做呢?.
就是開開心心的被擄去巴比倫.
去到巴比倫在那裡自買田地.
在那裡安居樂業.
你沒理由講這些這樣的信息的.
這些差不多是國安抓的.
叫猶太人不需要和巴比倫對抗.
要投降.
對當時的西底加王.
這個是叛國的.
所以耶利米就被那些領袖.
扔了進去一個枯井裡面.
又不讓他死.
要養著他.
來不及把餅掉下去.
你可以想像到.
其實我想講的是先知到了末期.
就是差不多到了被擄的時候.
講的信息全部不再是責備了.
是投降了.
就是說無可救藥.
就是說神都已經不再理這個民族了.
是必須要被擄.
去到外邦的70年.
所以現在這個時間.
馬拉基又不同了.
讀先知書就要留意這件事了.
不同的時代的那個預言.
那個大方向都不同的.
所以現在這個被擄歸回了.
也有了自己的聖殿了.
有了自己的城牆了.
也開始有自己的憲制了.
不過現在不是外族對猶太人的.
那個武力的威嚇了.
現在是過了一段長時間.
沒有了那種心靈的那個敏銳.
對神的真理.

$^{201}$還有那個慣性的敬拜.
已經成為一種例行公事了.
沒有了那個真誠的心去敬拜神了.
很多時候就會在民生上.
就會有很多的疲乏.
厭倦.
很多對自己的生活的不滿.
就投射在那個信仰都變得很習慣性了.
所以那些憲制有殘疾.
不潔淨.
都帶到祭司那裡.
祭司也慢慢沒有了那種的謹守.
沒有了那種的標準.
摩西的規定的標準.
那就慢慢的大家得過且過.
獻的人隨便.
奉獻的百姓.
又是諸多的理由.
叫祭司不要那麼高要求.
有總會沒有的.
所以帶來就是.
出現了一個好像大家都很.
心靈很疲倦.
所以這種心靈疲倦.
哀莫大於心死.
那個心已經去到一種很靜止的狀態.
所以馬拉基你可以說.
就是透過一些很有趣的.
寫作的手法.
來提出那個問題.
你們是這樣的情況.
那些百姓就好像被撩起他們.
有一種憤怒.
我們什麼時候是這樣.
提出好幾個問題.
好像引起他們的思想.
引起他們的關注.
來和先知馬拉基有一個對答.
其實是不是真是百姓的對答呢?.
還是先知說完那個問題.

$^{241}$就點出那個癥結在哪裡呢?.
可能是這樣的一個寫作的手法.
所以一至三章就重複.
在這樣的一個模式裡面.
來去出現的.
接著我們上一次說過.
三章十六十七十八節.
是說到當神要再回來的時候.
那個世界是會有一個劇烈的改變.
不過他還沒回來之前.
他會有一本紀念冊.
紀念哪些人是敬畏他.
馬拉基可能不知道的.
馬拉基只是去傳神的信息.
可能會有些人喜歡聽他說.
接著會自己私下跟隨去做.
但是馬拉基都不是很能夠立體.
知道每一個家庭.
每一個猶太人是怎樣的心靈光景.
但是很多時候先知書難讀的地方就是.
突然間有時候會將天上的那個景象.
告訴我們.
用一些寫作的手法告訴我們.
原來在神的裡面.
神的天上是有一本紀念冊的.
要寫多少頁呢?.
所以那個一定是心靈了.
一定不是紙張了.
有一本紀念冊.
紀念冊就是很珍貴.
紀念一些你很值得紀念的東西.
所以小學的畢業.
以前的畢業紀念冊是很珍貴的.
因為你一個人會長大的.
那本紀念冊就伴隨著你.
是你最幼年的時間.
又不是幼稚園.
但又不是去中學.
是你剛剛懂事.
開始大家分道揚鑣了.

$^{281}$那本紀念冊就非常珍貴.
那份友情.
那些學習階段.
但是現在在永恆裡面.
三章十六節就說.
我們有這一本紀念冊.
那你和我在不在神的紀念冊呢?.
這個就要問了.
所以新約有生命冊.
在這個白色大寶座.
這個是啟示錄的二十章.
有生命冊的.
所以我們很多時候.
沒有從這個永恆的角度.
去看我們自己.
是否符合神在我們生命裡面的.
即使用劇本.
有沒有符合這個劇本呢?.
我們是否在演活這個劇本.
那個角色呢?.
這個就是很重要的.
如果我們沒有永恆呢.
個個都是這樣生活.
個個都是這樣起身.
個個都是那個時間擦牙洗面.
擠巴士.
擠地鐵.
星期六就稍為放假.
可以行下山.
休息下.
星期就回教會.
如是者可能幾十年就是這樣了.
沒有什麼大變化了.
但是我們當然是這樣生活.
但是讀聖經寶貴的地方就是.
在神眼中.
我們都是活著的.
但是在神的眼中.
我們是否在祂的喜悅範圍呢?.
所以這個就是.

$^{321}$告訴我們三章的十六,十七節.
我們是否屬於神呢?.
是否特特的歸神呢?.
我們是否得到神的連率.
好像祂連率自己的兒子那樣的連率呢?.
我們是否真的一路事奉神.
被祂分別出來呢?.
這個就是將來才會真正的顯明.
但是現在呢.
我們就是要在神的紀念冊底下生活了.
所以我上個星期就說過了.
清教徒對金錢官是很重要的.
因為他們很怕自己不被揀選.
因為宗教改革就是離開天主教.
離開大公教會.
離開教皇.
你沒有了那個Blessing.
教皇說你離開不只是教堂.
你不只是不去天主堂聚會.
而是你跟隨那些宗教改革者.
加爾文,馬丁路德,慈雲利等等這些人.
離開教廷.
那你就沒有了教皇的Blessing.
他會警告你離開了在教會以外.
是沒有救恩.
沒有救恩就是蒙茫揀選.
所以他們就被教導.
在整個人生的.
他們一生的.
他們的Profession.
他們的專業.
他們人生的事業.
他們怎樣來看自己的事業呢?.
就是他們要拼命殷勤的工作.
就是怕自己不被神揀選.
所以這個就是清教徒的內心的守則.
那就是那一代的牧者.
馬丁路德.
加爾文比較嚴重.
特別是在瑞士一段長時間.

$^{361}$令到當時整個瑞士都很基督教化.
就真的帶來那個揀選論是很強烈的.
所以馬拉基書正正就是講這件事.
在地上大家都是猶太人.
但是有沒有在神眼中的紀念冊被紀念呢?.
有沒有是神連率的兒女呢?.
是他所得的連率揀選的呢?.
這個就是那個精髓就是這些地方了.
所以第二章第三章頭.
是有一個強烈的對比的.
很多人就說我什麼時候不敬畏你.
我什麼時候用鬼扎代弟兄.
我什麼時候有什麼.
但是第三章近尾.
就有一個很正面的給我們看到.
原來聽馬拉基先知的教導.
去思想神.
談論神的猶太人.
原來是在神的申命記念冊裡面被神紀念.
最後第四章我們六節快點講完.
就完我們這個系列了.
剛才提過了.
萬君之耶和華的日子.
當然了其實這是一個很大的課題.
不只是這六節有講的.
其實其他的先知書.
因為有些先知書我們沒有提的.
譬如亞摩斯書.
哈巴谷書.
是很多提到那個日子的.
分散在不同的先知書的寓言裡面.
所以這只是其中一小段.
這個其實牽涉到舊約怎樣看末世的.
所以是有些重疊的.
啟示錄講的末世和舊約講的末世.
是有少少的重疊的.
當然有些可能未是即時的.
在地上的應驗.
但是啟示錄就是一個比較終極的.
人類的結束的日子.

$^{401}$所以我們就從舊約去看一看.
那個日子會有什麼特徵呢.
其實你讀完這些.
你回去看福音書.
稍微揭曉一些.
譬如馬太福音的十三章.
帖撒羅尼迦的後書.
不過我們前面就會講.
講很多末日的相呼應.
和這一段聖經的相呼應.
要留意了.
到底那個日子會怎樣呢.
這裡又是一個對比.
惡人面對那個日子.
和一個義人面對那個日子.
怎樣是很慘.
和怎樣是很好.
很榮耀.
我好像沒有記錯.
道家福音說.
當那個日子到.
你們要歡喜跳躍.
這裡又講跳躍.
第二節也有講跳躍.
要好像那些肥牛讀.
離開那個牛圈.
跳躍的肥毒.
我讀讀第一到第六節.
第四章的《馬拉基書》.
標題都很清楚.
耶和華降臨的日子.
萬君之耶和華說.
那日臨近.
勢如燒著的火爐.
凡是狂傲行惡.
必如碎階.
在那日必被燒盡.
根本之條一無存留.
第二節留意了.
但向你們敬畏我名的人.

$^{441}$必有公義的日頭出現.
其光線有醫治之能.
你們必出來跳躍.
如圈裡的肥毒.
你們必踐踏惡人.
在我所定的日子.
再提到這個日子.
他們必如灰塵.
在你們腳掌之下.
他們就是那些被神厭惡.
和敵黨神狂傲的.
行惡的人.
被火爐燒盡.
然後連我們都可以踐踏這些惡人.
啟示錄有的.
啟示錄十九章.
那個白馬王子.
其實是說主耶穌.
十九章那個騎著白馬的那一位.
就會踐踏.
我們跟隨他的.
都是會這樣去踐踏敵人.
所以踐踏.
在啟示錄.
我想古代.
你打贏仗.
一定是踐踏敵人.
所以踐踏都是一個.
很形象化的勝利.
代表那個勝利.
Total Win.
你可以想像到.
連我們都可以.
因著這個公義的日頭出現.
我們都可以踐踏那些惡人.
所以不是今天我們就說.
我信耶穌的 抵你到地獄.
就不是這些踐踏.
這些是侮辱 不是踐踏.
這些是我們還沒有真正學得好的.

$^{481}$肅靈生命.
才會對人有一種懷恨.
千萬不要.
基督徒是不應該向人懷恨.
這裡說.
這是萬軍之耶和華說的.
留意第五節.
我特意不讀第四節.
第五節和第六節.
看啊 耶和華大而可畏之日.
未到以前.
我必猜顯先知.
以利亞到你們那裡去.
好像有以利亞的.
這樣的心智.
能力.
在哪裡說呢?.
路加福音第一章.
明顯解釋以利亞給我們知道.
因為這裡預言那個以利亞.
以利亞我們知道.
就是列王記相.
十七章說的那個神人.
突然間就出現在.
阿哈王朝.
而且當時是最黑暗的時代.
是這個.
阿哈的太太耶洗別.
這個皇后.
壟斷了朝綱.
令到整個以色列國.
是拜偶像.
拜巴力.
拜亞瑟拉等等.
所以其實是非常黑暗的時代.
那些巴力的廟.
在以色列供奉.
你可以想像得到.
是多麼困難.
突然間出現一個叫基列.

$^{521}$寄居的.
提斯比人.
以利亞.
他出現之後.
突然間去找阿哈王.
如果這幾年我不祈禱.
就不下雨了.
因為那裡是乾旱和饑荒.
因為越來越沒有水.
那你可以想像得到.
他不祈禱就不下雨了.
那就是要賣他.
怕他什麼時候會祈禱.
以利亞.
那他就去到一個溪.
在那裡等候乾旱.
就是那個叫基納溪.
在那裡有些烏鴉.
會給他擔那些食物.
烏鴉本來是吃腐肉的.
腐屍.
死屍.
烏鴉是污穢的.
但是烏鴉可以擔那些肉給仙子吃.
那些烏鴉又不會不飛來.
那些烏鴉又不會自己吃了.
是饑荒來的.
烏鴉都要面對沒有水喝.
沒有肉吃.
幾年了.
但是烏鴉是定期的.
擔肉給以利亞.
那去到那個基納溪乾旱.
神就叫以利亞.
哪裡都不准去.
去哪裡呢?.
去西頓.
西頓就是耶駒別的家鄉.
就是去他的鄉下.
那不就更加嚴重了.

$^{561}$那些巴力.
在以色列都拜了.
何況回到西頓就更加嚴重.
等於你去一個印度教的地方就更加厲害了.
以利亞就去到西頓.
撒勒法.
住在一個地方.
一個寡婦的家.
在那裡等待神的時間到.
到了.
他就要向這個阿哈王顯現了.
你可以想像到.
這個就是以利亞的那個能力.
因為告訴我們.
以利亞是一個.
以利亞的那個能力.
要告訴我們.
將來他會猜顯.
先知以利亞到你們那裡去.
所以很多人都會.
一路透過聖經的預言.
很多的猶太人.
都會等候.
誰是以利亞呢?.
有以利亞的能力.
可以使得整個國家.
行神跡.
來去改變.
而且是大家一心一意的.
向神.
這是聶王紀上的十八章.
所告訴我們的經文的世界.
所以.
大家都期望.
這位的以利亞是出現的.
所以這裡告訴我們.
原來他更加重要的.
地方的工作.
就是第六節.
他必使父親的心轉向兒女.

$^{601}$兒女的心.
轉向父親.
免得我來就坐遍地.
轉向的回報.
心靈的回報.
就是以利亞.
會帶來的.
出現.
就會帶來的.
生命的改變.
就是一個回報.
父親轉向兒女.
兒女轉向父親.
就是說大家是平等的.
不會是單向的.
就是說.
大家猶太人.
家庭的.
倫理的.
國民之間.
大家都有一個轉向的.
誰可以做一個這麼大的復興.
誰可以帶來.
一個人真正的願意.
在神的面前的改變.
這些.
給了那個預言在這裡.
那顆種子在這裡.
什麼時候這顆種子能夠發芽呢?.
這就是一個.
很大的.
他們的等待.
什麼時候來呢?.
誰是以利亞呢?.
這一定是一個問題.
一定是他們會問.
這個預言的以利亞什麼時候來呢?.
神說明的.
這個大日未到之前.
就會有以利亞的出現.

$^{641}$所以就意味著.
耶穌的降生.
施洗約翰的傳到.
就是末日的開始.
這是新約.
我們看末世論.
的一個很重要的訊息.
其實現在是末世的了.
不過是不是末世的末時呢?.
我們就不知道了.
也不應該猜了.
末世就是肯定的了.
聖經告訴我們.
主耶穌的出現.
和施洗約翰的出現.
就是末世的了.
所以這是大而可畏的日子.
未到以前.
我們說的大而可畏.
就是針對那些沒有主耶穌救恩的罪人.
要面對將來的公義的日頭的審判.
剛才我漏了第四節.
沒有說.
這個就很寶貴了.
第四節.
你們當紀念我僕人摩西的律法.
就是我在荷列山.
為以色列眾人.
所吩咐他的律例.
典章.
表面上好像正常的.
勸勉和教導.
但是.
這裡就是分辨到.
剛才第四章的第一節.
那些公義日頭降臨的時候.
有哪些人.
是要面對公義日頭的焚燒.
第二節就是公義日頭的光線.
可以有醫治的作用.

$^{681}$是兩個分水嶺.
有些人面對那個光線被燒盡.
但是有異人面對這些光線.
當然是主耶穌的.
那個.
牧羊 教導.
不是真的曬太陽.
就醫治你.
是的 增加你的維他命D.
令到你的骨.
就不會疏鬆等等.
你會有某些維他命.
維生素.
不是說這個物質.
而是說這位主的出現.
他會帶來我們生命的醫治.
所以用一個描述.
古代對太陽他們都是敬畏的.
因為不是太多研究不懂.
但是他們都知道.
我們今天都知道.
太陽直射.
如果你在當午.
如果你不戴太陽眼鏡.
戴都不行.
如果長時間直望太陽都會盲.
何況.
古代的人都沒有那麼多保護.
你可以想像到.
這裡有一個很大的分水嶺.
所以第四節.
就是將這兩種人.
帶來的.
那個.
就是怎樣能夠.
在.
這個人生.
你能夠有.
那個準繩.
就是.

$^{721}$所以不是讀一下聖經.
乖啦.
不是這麼簡單.
而是真是神的道.
是會帶領我們的人生.
去到成為.
在這個大而可畏的日子.
來臨之前.
我們就是.
在神的眼中.
得到祂的.
醫治,牧養.
但是相反.
如果我們將第四節.
看為又叫我讀經.
很煩啊摩西五經.
這些道理我認識很多.
不用你講了.
我小時候就已經回教會了.
這裡繼續.
告訴我們.
神的道是會帶給我們.
生命的一個改變.
其中一個回轉的改變.
是一個回轉.
帶來的改變.
就會使我們面對.
那個大而可畏的日子.
主耶穌這位公義的太陽.
公義的這位救主.
如果在耶利米蘇.
就叫做.
公義的描銳.
就是那個Seed.
這個描銳的出現.
會帶來我們的牧養.
所以你可以想像到.
我們基督徒.
只會越來越羨慕.
那個日子.

$^{761}$一個惡人.
只會越來越懼怕.
那個日子的降臨.
因為沒有.
內心的牧養.
沒有那個公義.
太陽的光線.
那個醫治.
你就可以想像到.
第四節有多重要.
不是再一次很煩.
讀一下聖經.
而是再一次說.
在人生裡面.
怎樣可以被神紀念.
寫在紀念冊的人.
就是第四節.
真是最重要.
如果不是.
我們就會跌落第一節.
那個經文裡面.
說到那些人面對那個審判.
將來是無法逃避.
的幾個痛苦.
所以我們.
回來看這裡的時候.
我可能先揭開一些經文.
這樣我們看回就容易些.
因為似乎.
新約的作者.
都是會根據這個.
的寓言.
來講到那個末日.
我看看.
譬如詩篇的21篇.
我引了好幾段經文.
這些我們不算太熟.
因為我們熟的23篇.
我引了好幾段經文.
這些我們不算太熟.

$^{801}$因為我們熟的23篇.
有時太熟23篇.
前後我們都不知道要講什麼.
因為真的不知道.
大衛的《思交與靈長》.
第21篇.
耶和說.
王必因你的能力搬起.
因你的救恩.
他的快樂何其大.
他心裡所願你已經賜給他.
他嘴唇所求.
你未嘗不應允.
你以美福迎接他.
把精金的冠冕戴在他頭上.
他向你求壽.
你便賜給他.
就是日子長久.
直到永遠.
他因你的救恩大有榮耀.
你又將尊榮威嚴.
加在他身上.
你使他有雄福.
直到永遠.
又使他在你面前.
興起快樂.
王依靠耶和說.
因至高者的慈愛.
必不搖動.
留意了.
惡人.
你的手要手出你的一切仇敵.
你的右手.
要手出那些恨你的人.
你發怒的時候.
留意一下.
和《馬拉基書》的四章.
火爐的烈焰.
完全一樣描述.
要使他們如在炎熱的火爐中.

$^{841}$耶和說要在他的震怒中.
噴滅他們.
所以火爐.
那個公義的日頭.
一定是對惡人來說.
是一個審判.
是一個神的烈爐.
其實更恐怖.
有一個苦杯.
是唇一不擇的.
唇到百分之一百.
小時候.
你吃過喝過苦茶.
很小的時候.
我們不懂得吞藥丸.
媽媽就很好.
就把它搗碎.
吃了之後覺得很苦.
就馬上送去喝水.
小時候就有很多辦法.
不吃.
媽媽就會買一些陳皮莓.
因為甜.
為了吃那些陳皮莓.
所以就要吃那些藥.
但又不懂得吞.
小時候在廣華醫院.
我們住在九龍塘.
經常去廣華醫院.
所以晚上小學的時候.
經常在廣華醫院排隊.
因為晚上很多人.
那時候很窮.
家家戶戶.
小朋友就在廣華醫院.
走來走去.
樓上樓下走.
當然現在裝修改變了.
你們不會想到那個樣子是怎樣的.
這個就是我們童年的成長.

$^{881}$是想要那個陳皮莓.
不是要那顆藥.
因為小時候它夠甜.
這裡告訴我們.
神的烈怒.
是會向惡人.
去發怒.
吞滅.
不過這裡強調的是.
那個炎熱的火路.
你記得.
有個歷史記載.
但以你.
他有三個朋友.
一起被擄.
第二章.
尼布格尼撒王就做了一個火搖.
其實應該說是做了一個金像.
就叫那些人一定要拜金像.
當有那些巴比倫的那些惡氣.
行奏的時候.
所有的人都要.
俯伏拜那個金像.
真的要拜.
但是不知道為什麼.
他們發現有三個朋友.
一起被擄.
第二章.
不知道為什麼.
他們發現有三個人沒有拜.
就是撒德拉,米薩,阿伯尼哥.
沒有拜.
就是大耳里的朋友.
大耳里就不在其中.
他可能有公務等等.
有人就報告給王知道.
王就因為這樣的緣故.
就很生氣.
很生氣.
居然不拜我的金像.

$^{921}$就是說這幫人.
怎樣教化.
改變他們的文化.
給他們高薪厚職.
給他們最好的飲食.
他們都不懂得抬舉.
還懷念著自己的故土.
就是說.
一直都養不熟.
但是尼布格尼撒王.
真的很有管理才華.
過來過來.
剛才可能你們誤會了.
聽不到.
給你多一次機會.
當你們再聽到.
那些割,滴,號的吹響.
你們拜.
之前就不算數了.
隨便他們.
接著就真的吹奏了.
但是仍然這三個人都不拜.
那麼尼布格尼撒王就沒辦法了.
就吩咐人.
燒熱那個火窯.
七倍.
就將這三個人.
要帶去台.
帶進去那個火窯.
所以.
因為皇命的緊急.
台的人.
都被那些烈焰燒死.
那三個人.
真的被扔進去那個窯.
一個很大的火爐.
來燒.
尼布格尼撒王.
以為自己.
下了那道火.

$^{961}$但是為什麼突然間.
看到那三個人.
在那個火爐裡面.
在那裡走來走去.
沒有燒死.
還不單止這樣.
為什麼我會看到四個人.
明明三個人.
難道彈劾跳下去.
為什麼那個人好像人質一樣.
你可以想像.
人世間.
都有這些酷刑.
這些火爐.
不過將來是神的火爐.
將來是神的.
憤怒的烈焰.
不是開玩笑的.
那個永恆的火.
在神的面前.
要面對.
那個審判.
如果不是有主耶穌.
洗乾淨我們的罪.
我們在神的面前.
有紀念冊.
被紀念有我們的名字.
我們就要面對這個.
公義的火爐.
所以公義的日頭.
對神來講是一個火爐.
對義人來講.
對屬神的人來講.
那個光線是有醫治的作用.
可以想像到.
一個此消彼長.
所以原來我們人生.
我們做基督徒.
要不然真的要認真去.
面對我們自己的信仰.

$^{1001}$我們的心和我們.
在人生的行動.
能夠符合神的心意.
還有其他的.
例如可以揭開.
13章的馬太福音.
這些一定是將來.
末世論要講的經文.
13章馬太福音.
第40.
36節這一段.
13章的36節.
13章有7.
對半的.
比喻.
很工整的.
馬太作者.
告訴我們.
末世天國的比喻.
用比喻去講天國.
這是馬太福音.
是一張很出名的聖經.
其實我們都應該讀一下.
會好的.
36節.
耶穌就解釋那個.
敗旨的比喻.
因為門徒就很感興趣.
為什麼會有這個敗旨的比喻呢.
首先.
我們趁這個機會.
有一點時間.
我們解釋這個敗旨的原先的講論.
然後耶穌就解釋.
那個敗旨.
24節.
13章的24節.
就是這個讀的敗.
敗壞的敗.
耶穌又設個比喻.

$^{1041}$對他們說.
天國好像人.
殺好種在田裡.
及至人睡覺.
有仇敵來.
將敗旨.
當然是不好的.
殺在墨紙裡.
就走了.
到長苗吐垂的時候.
敗旨都顯出來.
田主的僕人.
就來告訴他.
主啊.
你不是殺好種在田裡嗎?.
是啊 當然是殺好種.
從哪裡來的敗旨呢?.
因為有仇敵.
第28節主人就說.
這是仇敵作的.
僕人說.
你要我們去.
這個讀的好字.
去好出來摸.
拔走它.
29節主人就說.
不必 恐怕好敗旨.
連墨紙也拔出來.
容者兩樣一起長.
就是說.
未真正的飽滿.
結實 修成的時候.
就是說.
未修成的時候.
就容量存在.
接著.
等著修割.
當修割的時候.
我要對修割的人說.
先將敗旨好出來.

$^{1081}$菌成菌 留著銷.
唯有墨紙.
要收在倉裡.
比喻就很清楚.
那個主人的田.
是殺好種.
殺了一些敗旨.
對當時的農業.
的門徒.
他們是很明白.
很知道 但是不應該的.
為什麼不守衛著那個田地.
容釀仇敵.
來呢 說明是比喻.
要帶出那個contrast.
帶出.
那個人世間.
是有.
好壞.
就是比喻的精髓.
所以.
就不是說有沒有人守衛.
隨便讓人進來.
不是這個意思.
是要想帶出.
未到修割的時候.
這個主題.
修割的時候就有東西看了.
好.
任你們生長.
任你們去.
結果.
你結的是什麼果.
他結的是什麼果.
結的是敗旨.
是不能夠.
去吃的.
麥子.
結實的.
就是修成的麥子.

$^{1121}$是這個意思.
所以地球上.
公義的人還是罪惡的人.
就是在這個田地裡.
留意.
第36字.
門徒不問猶自可.
在門徒問耶穌.
你講了這個比喻.
你解釋給我們知道.
當下36節.
耶穌離開眾人.
進了房子.
他的門徒進前來.
就是說在房子裡的門徒.
才聽到這個比喻.
所以馬太就寫了下來.
說請你把田間.
就是你剛才講的那個.
敗旨比喻.
講給我們聽 門徒很有興趣.
為什麼會容許有敗旨呢.
留意耶穌的解釋.
第37節.
哈 就是耶穌.
回答了.
那個殺好種的就是人子.
就是耶穌.
田地就是世界.
留意.
這個今生今世.
地球.
總之是人類.
由出世到死.
的那個場景.
不是香港.
不是美國 是全個地球.
總之你是人.
你就在這個世界.
就是這個田地.

$^{1161}$記住 田地就是整個世界.
我們就是活在這個世界的人.
好種.
是天國之子.
你就是剛才所講的.
公義的日頭.
他的光線.
有醫治的作用.
就是對我們這些熟悉他的人.
在那個紀念冊裡面.
被紀念的人.
是談論耶和華.
思想神的人.
你可以想像到.
是一個.
原來舊約和新約.
都是這樣講這個真理.
問題就是.
誰使我們能夠成聖.
就是耶穌.
舊約就未懂這東西.
新約就懂了.
所以原來我們大家都是在這個世界.
但是我們沒有人知道.
原來那個使我們得生命的.
就是這位人子.
是他殺好種.
所以留意一下.
敗子就是那個惡者之子.
殺敗子的仇敵.
就是魔鬼.
原來在這個世界.
是有人.
因為未得到神的救恩.
是有罪性.
是罪人.
原來.
自己都不知道.
原來是魔鬼之子.
所以這個是.

$^{1201}$幾聖經的一個.
斬釘截鐵的.
好像很易言.
你不屬神你就屬魔鬼.
你告訴我們.
你看見就是.
收割的時候.
就是世界的末日.
收割.
上文就是講正常的.
那個田地.
大家都在那塊田裡面.
一路生長.
原來有些人告訴我們.
敗子和墨子.
剛剛殺的時候.
那些水子是差不多的樣子.
所以為什麼主人說.
不要拔出來.
因為你拔的時候.
很容易拔錯墨子.
大家都未生長成熟.
你拔了它.
連那些墨子都浪費了.
原來人生有些東西.
是要容許它存在.
我們很難想像.
為什麼這個世界會有監獄.
為什麼不將那些在監獄裡面的人.
全部殺了.
那就沒有監獄.
天下太平.
又不會再有惡人.
但問題就是沒有監獄.
我們在外面的人.
慢慢又是一個犯罪.
因為還在這個世界.
記住.
什麼時候才真正沒有罪惡呢.
世界的末了.

$^{1241}$就是這樣.
就是會將敗子.
拔走了.
所以收割的時候就是世界的末了.
收割的人就是天使.
和天使有關.
《小道行傳》有說.
帖撒羅尼加前書四章.
我們前兩個禮拜會說.
就是說天使.
會跟著回來.
來迎接我們這些.
的聖徒.
所以收割的人就是天使.
第四十節.
真正這個問題就在這裡.
將敗子好出來.
用火焚燒.
世界的末了也要如此.
也就是說.
世界的末日.
神的時間不會再寬容了.
不會再給人有悔改機會了.
那就需要.
進入永恆的審判.
就是永火了.
這個就是很嚴重的後果.
所以.
這就告訴我們.
第四十一節.
人只要差遣使者.
把一切叫人跌倒.
作惡從他國裡挑出來.
留意了.
再一次是雕在火爐裡.
在那裡.
必要哀哭切齒.
這個就奇怪了.
去到火爐就立刻燒盡.
何來還要哀哭切齒呢.

$^{1281}$就是說永恆的受苦.
不是一秒鐘.
跳到火爐裡.
死了就死了.
最多掙扎幾下.
燒死而已.
都是很痛的.
都是幾秒鐘就死了.
但不是的.
是繼續的哀哭切齒.
在神的國度以外的.
一個永恆的審判.
在那裡.
必要哀哭切齒.
留意第四十三節.
那時義人在他們.
乎的國裡.
要發出光來.
好像太陽一樣.
有耳可聽就應當聽.
主耶穌說.
很清楚解了.
世界就是田地.
田地就是世界.
到底我們現在.
可能有些人明顯.
見到祂是敗旨.
但有時不是那麼明顯.
真的不是那麼明顯.
就是說.
四章四節.
馬拉基書的重要就在那裡.
哪個人對神的話語.
對摩西的律法書.
能夠有一份這樣的敬畏.
和願意.
進行.
來修煉.
來修養.
自己的心靈.

$^{1321}$讓神的道在自己的生命裡.
這個就是.
分水嶺.
到底我們是物質還是敗旨.
到底好像清教徒所說的.
我們是蒙簡算.
還是不被簡算.
這個就是一個.
很大的分水嶺.
但仍然有火爐.
其他我未必再引.
因為前面我們.
下一次我們要說的是.
帖撒羅尼加前後書.
和講摩西論不同的理論.
有幾個派別的.
給大家有些少輪廓.
就可以了.
當然你如果能夠繼續自己再研究.
你是什麼派就好了.
有你自己那個.
對經文的持守.
譬如有.
栽前被提.
有栽中被提.
有栽後被提.
有後遷欺凌.
有無遷欺凌.
到底你是哪些呢?.
我都提過了.
近代的學者多數是無遷欺凌的.
我自己就是前遷欺凌的.
這個就.
我們要來去理解.
所以我們回頭看.
這個火爐.
你就可以想像到.
這裡有兩個對比.
所以你看完剛才那個敗旨的比喻.
耶穌自己親自去解.

$^{1361}$一定不會錯的.
真的對號入座的.
我們是脈子.
還是敗子.
大家都在這個田地裡.
大家都只有一生一世.
人都一生一世.
所以希伯來書.
九章二十七節.
按照著定命.
人人都有一死.
死後.
且有審判.
有沒有留意?.
二十八節大家不熟悉了.
照這樣.
其實不是想說我們人人都有一死.
其實不用說過都知道.
他想說的是二十八節.
耶穌照這樣.
基督都來到這個世界.
所以祂只需要死一次.
祂只能夠死一次.
就夠熟悉我們每一個人了.
所以耶穌不需要.
次次來釘十字架.
這個是不對的神學.
所以.
天主教的彌撒.
就是每次每個禮拜.
來釘十字架.
所以其實不對的.
嚴格來說是真理.
基督我揭一揭.
大家又不熟悉了.
大家就懂背.
不過就不懂二十八還是二十九節.
希伯來書九章.
是的.
應該由二十三節讀起.

$^{1401}$順便教大家.
在那個言經的過程.
都能夠有一個這樣的認識.
二十三節.
九章二十三節.
希伯來書.
按照著天上的樣式做的物件.
按照著天上.
天堂上的樣式.
在地上的那些.
藥櫃.
香壇.
陳設餅.
還有就是.
外面的洗浴盆.
祭壇等等.
按照著天上的樣式.
做的物件.
必須用這些祭物去潔淨.
但是天上的本物.
自然當用.
更美的祭物去潔淨.
二十四節.
因為基督並不是盡了人手.
所做的聖所.
這是真聖所.
天上真聖所的影像.
祂是進了天堂.
原來真有天堂.
是神的永恆居所.
就是我們將來要去的地方.
如今為我們顯在神的面前.
主耶穌不是沒有了.
祂現在是回到天上.
在神的面前.
為我們而顯在神的面前.
這是主耶穌現在的存在.
二十五節.
在天上.
祂在天上.

$^{1441}$是的.
二十五節.
也不是多次將自己獻上.
像那個大祭司.
每年都要帶著牛羊的血.
進入聖所.
是說那個每年的贖罪日.
祭司是不斷要為百姓獻祭的.
但是永恆的祭.
永恆的救主.
永恆的救恩.
那個祭物.
就是耶穌.
所以祂不需要多次.
來獻上自己.
祂只需要一次獻上就夠了.
所以第二十六節.
如果這樣.
祂從創世以來.
就必定要多次受苦了.
但如今.
在這末世顯現一次.
耶穌的出現.
就是末世的開始了.
這裡也有說.
如今在這個末世顯現一次.
就是他是拿撒勒人耶穌.
就是他是.
從貞女瑪利亞降生.
成為有血有肉的人.
他成長了.
所以他的時間.
當時沒有標.
都是每秒和我們今天一樣.
每秒地成長.
直到他要去到.
三十多歲的時候.
就在十字架上斷氣.
這個就是他顯現一次.
將自己獻為祭.

$^{1481}$很除掉罪.
留意二十七節.
作者做了一個總結.
所以其實不是說我們.
按照著定命.
按照著一個.
神定的定律.
人在這個地球的定律.
人人只能夠死一次.
不會死兩次.
所以不會投胎.
下世做牛做馬.
就不是.
這裡很清楚了.
人人只能夠死一次.
死後且有審判.
對的,不過不要停在這裡.
因為第二十八節才說.
那個死的問題.
他說像者樣.
像者樣.
什麼像者樣呢.
既然基督都是做人.
活在這個地球.
就受著神命定的.
生死的約束.
所以基督.
既然一次被獻.
就是死一次.
擔當了.
多人的罪,普世的人的罪.
當然是信的人.
將來.
就要向那些等候他的人.
第二次的顯現.
所以耶穌只需要顯現一次.
成為了贖罪祭.
他就已經回到上天堂了.
他第二次再顯現.
就是要將我們這些.

$^{1521}$等候他的人來拯救.
所以並與罪無關.
乃是為拯救他們.
所以基督.
我們不要想他不斷地死.
不斷地獻祭.
他只需要獻一次.
他就永遠救贖了我們.
因為他就是神的兒子.
神自己的本體.
他來到這個世界.
就是按照著人的命定.
他既然做人.
他要受地上的那些定律.
約束的.
所以他有生.
也有死.
但是.
唯一不同的就是.
他與其他人類不同.
就是他死而永恆復活.
這個就是信仰.
就是我們所信的耶穌.
我們會生.
我們會死.
但是我們不會復活.
除非神的命定我們改變.
生命改變.
就是我們信主的人就可以改變.
不信的人.
有生有死.
但是不會復活.
所以我們將來.
因為我們現在已經有主耶穌的生命.
有聖靈在我們裡面.
所以我們將來.
我們就會聽到天使的呼聲.
我們就會復活.
是身體的復活.
我們都是有一個永恆的生命.

$^{1561}$不過就不是.
回到變回永恆的神.
唯一的不同.
但是耶穌就是.
他從天上來到這個世界.
做人.
死然後復活.
升天然後回到.
他自己永恆的神聖裡面.
這個就是.
他與所有的人不同的地方.
所以這裡.
很想強調的就是.
那個一次性的贖罪.
耶穌基督已經完備的救恩.
所以我們回來看這裡的時候.
就是我們要記住.
作者告訴我們.
那個燒著的火爐.
對惡人就是狂怒行惡.
將來會變成.
碎街.
會被燒盡.
毫無保留的燒盡.
但是對於我們敬畏神的名的人.
必定又有另外一個.
的依治.
就是公義的日頭.
這裡其他的.
聖經的文章.
都講到.
我們要記住.
其他的.
先知書又有很多教導.
譬如我們先揭一揭.
耶利米書.
以創下書又有的.
耶利米書.
算是最難讀的一卷.
先知書.

$^{1601}$耶利米書.
又要背下那些.
聖經的目錄歌.
講出李文新書士得薩薩.
皇皇代代拉尼斯.
白詩.
曾傳歌.
蔡爺哀結.
但.
何以摸.
餓衲尼雄哈番.
該亞瑪.
39卷很容易背.
這些口訣.
耶利米書的23章.
33章都有的.
耶利米書是.
全舊約最長的一卷書.
不是詩篇.
按那個節數來講.
或者按原文的.
那個節數.
耶利米書很長的.
有些解經書我看過.
就是說耶利米書.
是全本舊約裡面.
全部舊約裡面出現最多.
耶和華如此說.
出現最多有700次.
耶和華如此說.
所以耶利米書沒有.
自己的.
什麼都是耶和華如此說.
那你就要聽了.
因為不是我講的.
不是我耶利米說的.
是耶和華要我這樣講.
耶和華如此說.
是權威中的權威.
但是一卷叫他們避老.

$^{1641}$叫他們投降的書卷.
有700次耶和華如此說.
所以.
耶利米書的欄在這裡.
23章.
對不起.
可能.
沒什麼可能有些乾.
對將來的希望.
耶和華說.
那些殘害趕散我.
草長之羊的牧人.
有禍了.
耶和華以斯利的神.
斥責那些牧羊他百姓的牧人.
如此說.
你們趕散我的羊群.
並沒有看顧他們.
我必討你們這下惡的罪.
這是耶和華說的.
我要將我羊群中.
所餘剩的從我趕.
他們到的各國內.
招聚出來.
令他們歸回本圈.
即神親自的牧羊.
他們也必生養眾多.
我必設立照管他們的牧羊.
他們不再懼怕.
不再驚惶.
也不缺少一個.
這是耶和華說的.
講的是一個公義.
一個他所立的牧羊.
這些意境.
在約翰福音的.
我是好牧人.
如果是故宮.
不會顧惜那些羊群.
當有狼來.

$^{1681}$故宮就會逃跑.
我只是打工.
你要我搶救那些羊.
當然不願意.
故宮逃跑.
因為他是故宮.
耶穌說我是好牧人.
好牧人為羊捨命.
所以這些意境.
其實在舊約出現很多.
所以主耶穌.
將舊約的事物.
應驗在自己身上.
祂就是好牧人.
第五節.
日子將到.
我要給大衛興起.
一個公義的苗裔.
呼之欲出.
肯定是說彌賽亞.
主耶穌是大衛的後裔.
所以馬太福音一開始.
開宗明義.
亞伯拉罕的子孫.
大衛的後裔.
就是耶穌基督的家譜.
公義的苗裔.
他掌王權.
在地上施行公平.
公義.
公平和公義.
有時好像孿生兄弟.
在舊約常常出現.
公平公義.
好像包含著國家裡面.
有牧羊.
有一個智慧的君王.
現在.
不說政治.
美國快有兩個老人家在爭.

$^{1721}$不知誰.
一選到.
就要入醫院了.
搞笑.
整個國家這麼多億人.
只有兩個老人家做總統.
你可以想像到.
有時.
那個制度未必.
真是能夠去到選賢.
去選一些.
真是優質的.
領袖.
你可以想像到.
將來.
有這個公義的苗裔.
在舊約已經有這樣的預言.
我們一看就知道.
肯定我們新約的人可以對號入座.
就是我們的主耶穌.
他有公義在他裡面.
他是大衛的苗裔.
即是說.
應許大衛的子孫.
可以坐王位.
那個苗.
即是那個屑.
那個根.
很多時候都有說這個根源本身.
將來的後裔.
所以這個裔字.
很多時候都出現在歷史.
我有一個裔字.
都是那個意思.
有那個.
這樣的含義.
所以這個苗裔.
都是很精進.
那個中文的翻譯.
告訴我們.

$^{1761}$將來他會掌王權.
所以現在耶穌被人嘲笑.
是現在.
因為.
田地是屬於.
神的.
不是屬於魔鬼的.
從來殺好種的人質.
魔鬼是偷進來.
仇敵.
偷偷地.
殺敗子.
所以我們記住.
我們成為了基督徒.
應該是很榮耀的.
活在地上.
應該展現的.
最近又有些人.
來談一些合作.
等等.
不知道他們是在開玩笑還是什麼.
牧師你為什麼不笑.
你為什麼不笑.
不笑的樣子都笑.
不知道他是誇大還是什麼.
但基督徒應該是這樣展現.
一種這樣的.
心境.
或者那個.
有主耶穌的生命那種.
活潑.
告訴我們.
我們屬神的人.
其實我們才是得勝.
行事有智慧.
所以我們相信了主耶穌.
應該求主給我們有智慧.
如果你每個禮拜都被人騙錢.
電話騙案.
那就真的不應該.

$^{1801}$是不是.
在地上就.
施行公平公義.
也就是說我們都應該要.
我們或多或少.
在我們自己的範圍.
譬如我也要管理一些.
下屬.
我是不是公平公義呢.
這些都是應該反映的.
在他的日子.
猶大必得救.
以色列也安然居住.
這些說話.
真是很好的.
就是說那個國家很太平.
大家的心.
晚上出街都不怕.
到了很晚.
你都不會怕.
你想想你在印度.
人們好像.
很瘋狂的.
你坐火車都可以被人強姦.
印度.
經常都出現.
你就可以想像得到.
那何來安然居住呢.
他們拜那麼多假神.
為什麼不會令到他們的民.
國民的素質.
道德會有改善呢.
那就是證明.
他們拜的不是真神.
是安然居住.
他的名必稱為耶和華.
我們的義.
這個就是神.
是我們的公義.
主耶穌的名字.

$^{1841}$所以是耶和華我們的義.
到了33章又再出現了.
下面我不再讀下去了.
33章.
和耶利敏書難讀的地方.
他的時序.
不是順序.
有時說一下.
在西底加最後的.
那幾年.
突然間會跳去說.
約西亞王.
西底加的爸爸.
你怎麼能夠將那個.
年期.
來做一個拼湊呢.
所以如果用那個.
那個巴刷.
這裡說約西亞突然間說西底加.
那就是說一個.
讀的經卷.
不是一個只是循序漸進.
所以很難讀的耶利敏書.
33章.
後面那裡.
夠不夠時間讀多一點.
對不起.
好了讀多一點.
反正大家都是.
不太熟悉的.
33章有些這樣的時序.
耶利敏還囚在護衛兵的院內.
耶和華的話.
第二次臨到他說.
成就的是耶和華.
造作為要建立的.
都是耶和華.
耶和華是他的名.
他如此說.
你求過我.

$^{1881}$我就應運你.
並將你所不知道又大又難的事指示你.
這個就是耶利敏的祈禱.
他的熱線電話.
就是333.
不是999.
耶利敏書33章.
三節是金句.
叫耶利敏你求過我.
其實他現在很不明白自己為什麼會被囚.
他不是很明白.
為什麼我講這些.
你叫我說的話.
我會面對這麼多人的攻擊呢.
神就說.
你求過我.
我就會應運你.
第四節.
開始解釋了.
就是撒韋.
為黨敵人.
高磊和刀劍的.
耶和華以色列的神如此說.
人要與迦勒底人爭戰.
正是那死屍充滿者房屋.
就是我在怒氣.
和憤怒中所殺的人.
因他們的一切惡.
我就掩面不顧者城.
看啊.
我要使到這城得以存如安書.
使到城中的人得醫治.
又將風聲.
的平安和誠實.
顯明於他們.
我又要使到猶大被擄的.
和以色列被擄的歸回.
並建立他們和喜初一樣.
講到預言.
神說你求過我.

$^{1921}$我就告訴你結局.
現在你去和巴比倫競爭.
你好像變成死屍一樣.
在他們的手中.
但是當神的時間到.
他會使到猶大.
以色列和以色列.
猶大和以色列.
可以歸回.
已經預言了.
可以被擄歸回.
和他們會有好日子過.
會有平安等等.
第八節.
我要除盡他們的一切罪.
就是向我所犯的罪.
又要赦免他們的一切罪.
就是干犯我.
違背我的罪.
這城要在地上萬國人面前.
使到我得眾讚.
得榮耀.
名為可喜可樂之城.
萬國人因聽見.
我向這城所賜的福樂.
所施的恩惠平安.
就懼怕賤驚.
但是耶利米在一個枯井下.
神告訴他們.
將來萬國都會清耶路撒冷.
是榮耀之城.
是要被眾讚.
但是現在耶利米在那個護衛院的枯井.
是要被眾讚.
很搞笑的.
現在衣衫藍柳.
告訴你你是一個大富翁.
現在不是看相.
是神的說話.
就是叫耶利米.

$^{1961}$你那些又大又爛的事.
你求過我吧.
就告訴他們.
將來這個民族會有這麼大的榮耀.
你可以想像得到.
那個落差對耶利米來說.
是很難接受.
但是相信.
他慢慢就明白.
神要他看到這一種.
將來永恆的榮耀.
但是目前.
沒辦法了.
因為他處身在一個被擄的時代.
他是要作為被擄時代的先知.
就是叫百姓.
你們沒辦法了.
神的審判來了.
唯一可以逃避的就是投降.
就帶來更大的苦楚.
到下面.
第十節 耶和華如此說.
你們論這地方說.
是荒廢無人民無新畜之地.
但在這個荒涼無人民無新畜的猶大城邑.
和耶路撒冷的街上.
必在聽見有.
歡喜和快樂的聲音.
新郎和新婦的聲音.
並聽見有人說.
要稱謝萬君之耶和華.
因耶和華本為善.
他的慈愛永遠長存.
這是詩篇的主題.
猶太人.
他們by default.
神.
就是這樣的一位他們的神.
他的本性是善的.
他的本質.

$^{2001}$雖然這樣說也有點冒犯.
神應該不是.
這樣稱呼的本質.
他的屬性是善的.
所以我們敬拜.
這位善的真神.
所以.
你今天.
事事都不如意.
你就不要罵神.
耶和華本為善.
他的慈愛永遠長存.
所以其實.
這個是我們一生都要去.
經歷去經驗.
耶和華怎樣本為善呢.
在你和我的生命裡面.
他怎樣本為善呢.
他怎樣永遠長存呢.
所以猶太人如果那些敬虔的.
就真的會看他們的人生.
就是這樣.
殿中的人的聲音.
因為我必駛到這地.
被擄的人歸回了.
預言將來的歸回了.
和起初.
起初就是那個亞伯拉罕的年代.
摩西的年代.
大衛的年代.
這是耶和華說的.
所以我們應該期望.
我們信了主將來.
在永恆的那個福樂.
地上是一個.
實驗場.
雖然這樣說有點誇張.
因為田地就是世界.
所以我們還沒離開世界之前.
我們記住.

$^{2041}$還沒去到修葛的日子.
我們都要保守自己的心靈.
在神的揀選裡面.
這個是很重要的.
接著.
十二節.
萬君之耶和華如此說.
在這個荒廢無人民無新築之地.
並其中所有的成一.
必在有牧人的住處.
他們要駛到羊群.
躺臥在那裡.
在山地的成一.
高原的成一.
南地的成一.
遍雅敏地 耶路撒冷四處的各處.
和猶大的成一.
這麼強調.
整個以色列.
沒有一處沒有牧羊人.
沒有.
那個豐收.
必在有羊群.
從數點的人手下.
經過.
就是說.
大家都是有最美麗的.
圖畫和光景.
就不會是.
那些羊群沒有了牧人.
上文就說了.
沒有了那些.
任避 任遊的.
豺狼來捕捉.
現在就是各家各戶.
都可以數點自己的群羊.
你可以想像到.
是一個多麼快樂.
多麼豐富的.
一個國度.

$^{2081}$接著十四節才是.
日子將到 我應許以色列家.
和猶大家的恩賢.
必然成就 當那日子.
那時候 我必使到.
大衛公義的苗裔長起來.
他必在地上.
施行公平和公義.
連續說兩次.
在那日子猶大得救.
耶路撒冷必安然居住.
他的名稱為.
耶和華我們的兒.
是說日子將到.
似乎就是主耶穌降生開始了.
就是帶來.
那個生命的改變.
所以保羅.
對這些經文憤外上心.
羅馬書的十一章.
他就說這個日子.
所以這裡告訴我們.
公義的日頭.
就看你是.
屬於哪一類別.
你是拜子還是麥子.
拜子就是.
一起生長在那個田地.
意味著.
都有那個的.
供應 即是說返教會.
輕道 有培靈會.
有查經班.
有獨經日營.
有暑假八月的.
培靈研經會.
有這些的供應.
但是本質你是拜子.
你不肯面對自己的本質.
所以你到修改的時候.

$^{2121}$都是拜子.
麥子就不同了.
麥子的偶然間的跌倒.
因為他內心是麥子.
他就會有一個改變.
因為還是可以有.
悔改的時間.
有神的道.
所以四章四節再一次告訴我們.
可以令到生命改變的.
唯有是神的道.
就不是一些神跡奇事.
所以第一第二節.
是一個對比來的.
是一個對立來的.
是生命的對立.
所以不要經常責怪.
牧者說你信了耶穌沒有.
其實有時會問一下自己.
是不是真的在神的.
救恩裡面.
很重要的.
如果一問你就已經火滾.
就說什麼都不談論.
對屬靈的事.
完全沒有興趣.
其實值得懷疑你自己.
是不是被揀選.
或者你自己要問一問.
你是不是被揀選.
是這個問題.
第三節就告訴我們.
這個公義的日頭出現.
就令到我們.
剛才我們說了.
安然居住是一個形容.
就是那種國泰民安的地步.
安然的.
內心是很坦然的.
比如說.

$^{2161}$剛才有一個緊急探訪.
在天蘇圍醫院.
我第一次去那個醫院.
那個老人家其實不大.
可能六十多歲.
因為末期癌症.
完全沒有心.
流眼淚.
他說我很想回家.
我很怕死.
我邀請他去信耶穌.
他說我不信.
那個人很upset.
很困難.
這些就不是.
安然居住.
安然居住就是.
就算我重病.
我是痛得很厲害.
甚至要用麻醉.
但一個基督徒他知道自己去永恆.
他會安然.
安然的心靈浮現出來.
連可能周圍的人.
醫生護士都感受到.
這個人有生命有永生.
真的.
這些是騙不了的.
因為真的痛.
但是安然居住是不簡單的.
所以舊約.
這樣的描述.
是在說神裡面的人.
然後這裡就再誇張一些.
好像一些公牛毒.
跳藥.
牛毒就是那些牛仔.
很生猛很有力.
我養豬.
你想一下.

$^{2201}$豬很蠢.
很懶惰.
吃完就睡.
是人類造成的.
因為給牠這麼小的圈.
不想牠在這裡跑.
我在養豬這麼多年裡面.
很有趣.
久不久有一些奇特的豬.
牠很有力.
因為牠經常跳.
通常我們做這麼高的鐵欄.
有東西攔住牠.
牠就不會跳.
牠就不會衝過去.
有一隻.
特別生猛.
吃飽之後.
牠就出來玩.
牠會跳出來.
真的很高.
牠可以一跳就跳出來.
走到處去.
有東西吃牠就回來.
牠自己會跳回圈裡.
跳藥.
這些是很麻煩的.
因為怕走出去.
但牠不會走出去.
總之每次吃完牠就跳出來.
牠就出來玩.
這樣很奇特的豬.
因為牠跳得多.
很健康很健壯.
因為其他那些通常養到越來越肥.
牠又不想動.
所以牠就睡著.
但這隻.
牠就特別經常跳藥.
所以這裡.

$^{2241}$告訴我們.
牠的光線有醫治之能.
我們在主耶穌的救恩.
在神的道的模樣.
滋潤之下.
我們應該有一種跳藥生命.
跳藥如圈裡的肥毒.
不是那些鴛鴦式的瘦猛猛.
是一個肥毒.
但是是跳藥的.
大家有沒有看過熱辣滾燙.
真人真事.
賈寧.
由二百多磅.
變回減了一百斤.
一年九個月.
為了拍這套電影.
不是教減肥.
是教你愛惜自己的生命.
所以很多.
內地的人一共鳴.
整個戲院的人都在哭.
那個描述.
很值得看.
那套電影.
我也哭了兩次.
就是.
你就看到.
那個跳藥生命.
這裡告訴我們.
在主耶穌裡面.
這位公義的日頭.
出現我們有這種活力.
生命力.
所以人子來.
要尋找拯救失喪.
人子來.
我來了.
要讓得到生命.
得到豐盛.

$^{2281}$生命和豐盛.
其實是兩樣東西.
這裡告訴我們.
如果我們在主耶穌裡面.
雖然我們.
都在一個大環境.
我們至少.
應該要有那個眼光.
那個價值觀.
應該要有一個比較全球性的視野.
和看市民.
不要看一個很灰暗的一面.
應該要有一種.
很積極的.
這時天賦世界.
這樣的心境.
這樣來生活.
因為最近.
很多不同的年齡層.
輕生.
不同年紀.
一定有他們的故事.
有他們的困難.
教會接觸不到.
但是.
我們屬神的兒女.
應該要有這種跳躍.
醫治了你.
你得到神的醫治.
你就應該有這種.
出來跳躍.
如圈裡的肥獨.
我們展現在微信的.
這個田地有.
墨紙.
有拜紙.
但我們展現的應該是.
一種生命力.
一種愛.
特別是主耶穌.

$^{2321}$公義.
耶和華的公義.
所以我們應該反映.
這件事在地上.
我們會踐踏惡人.
在所定的日子.
我們就好像.
踩著腳掌之下.
無論羅馬書.
十五章.
和啟示錄.
這些經文.
所以最重要的是.
第四節.
這裡真的要紀念摩西的律法.
所以不是重重複複.
而是.
這是令我們的生命.
得到永恆救恩的唯一辦法.
所以神留下的是.
聖經給我們.
祂不是要我們追求的神蹟.
要我們有.
對神的真理的.
認識.
和信服.
這樣就能帶來.
我們的生命可以改變.
最後就是講以利亞.
可能用多幾分鐘.
要解解.
盧家呼音這一段.
因為是盧家直接.
告訴我們.
他出現帶來的.
他直接引用這段經文.
盧家醫生.
一張.
我們看完這一段就完了.
施洗約翰這裡.

$^{2361}$就是他向.
第一章的第五節.
當猶太皇希律的時候.
盧家呼音.
亞比亞班裡面.
有一個祭司.
他的名字叫撒迦利亞.
他的太太就是.
亞倫的後人.
名叫伊利沙伯.
都是亞倫的後人.
即是大祭司.
他們二人在神的面前.
都是二人.
遵行主的一切誡命禮儀.
沒有可指責.
指示沒有孩子.
因為伊利沙伯不能夠生育.
二人年紀老邁.
撒迦利亞按班次.
在神的面前供祭司的積分.
照祭司的規矩.
祭簽.
得盡主殿燒香.
似乎進入至聖所.
燒香的時候眾百姓在外面禱告.
似乎那是一個.
重要的.
公眾的聚會.
有可能是主彭節.
有可能是贖罪日.
不知道沒有提到.
有主的使者.
站在香壇的右邊.
向祂顯現.
一個人在聖殿裡獻祭.
抽籤輪到他.
到這一班人進去.
接著.
亞比亞班.

$^{2401}$他就一個人在聖殿裡.
見到神跡.
天使向他顯現.
第十一節.
第十二節.
撒迦利亞漢賢驚慌害怕.
見到天使.
天使對他說.
撒迦利亞不要害怕.
因為你的祈禱已經被聽見.
你的妻子以利沙白.
要給你生一個兒子.
你要給他起名叫約翰.
你必歡喜快樂.
有許多人因他出世.
也必喜樂.
他在主面前將要為大.
探酒,龍酒都不渴.
從無福裡就被聖靈充滿了.
似乎他成長.
也是一個拿世以仁.
奉獻給神.
分別自己.
第十六節.
他要使許多以色列人迴轉.
就是《馬拉基書》四章六節說.
使許多以色列人迴轉.
歸於主他們的神.
他必有以利亞的心智能力.
他就未必有行神蹟.
施生約翰沒有行過一件神蹟.
但是他有那種心智能力.
他必有那種心智能力.
就是叫人迴轉的能力.
行在主的面前.
叫為父的心轉向兒女.
叫背逆的人轉從二人的智慧.
背逆的人也可以有屬靈的智慧.
二人的智慧.
不再做一個愚昧的人.

$^{2441}$背逆的人.
為主預備合容的百姓.
他又去修門徒.
那些人預備他們將來會跟隨耶穌.
有沒有記得拿彈葉和路加.
跟隨施洗約翰.
後來轉去跟隨耶穌.
就是彼得的哥哥安德烈.
他是跟隨施洗約翰的.
這是告訴我們.
路加醫生很清楚告訴我們.
馬拉基斯所預言的將來末日.
來到的以利亞.
就是施洗約翰.
天使現在向撒迦利亞說.
天使去幫助解釋.
說你的兒子將來有這種心智能力.
誰的心智能力?.
就是以利亞的心智能力.
所以我們可以看看舊約將來有機會.
我們好像說過以利亞這個人物.
以利亞的心智能力.
我們未必需要學習以利亞的行神蹟.
祈禱有雨落.
不祈禱就沒有雨.
我們不要那些神蹟.
甚至有些更大的神蹟.
曾經降火燒那些白夫長.
可能你記得了.
這是聶王紀夏的第一章.
很多很多那些神蹟.
但是我們這個時代缺少的就是.
心智.
那個心智.
我們基督徒.
很少這些心智.
年輕一代.
被整個世界的工作繁忙.
政治.
搞到我們沒有了以利亞的心智.

$^{2481}$它怎樣令到整個以色列.
回轉呢.
離開巴黎呢.
這個是很強的能力才做到的.
這個就是施洗約翰.
他就有這個時代的訊息.
和他背負的責任.
但都是為百姓.
為耶穌預備百姓.
他不是來塞瓦.
所以你可以想像到.
對他的描述原來是這麼厲害的.
原來他只不過是.
先鋒.
預備耶穌來的.
所以我們應該在人生中.
學習施洗約翰的心智.
以利亞的心智.
我們有沒有慰錦宣家.
我們有那個心智能力呢.
我們是不是因為.
是這樣的了.
不要那麼複雜了.
總之回完教會.
趕快祝福完就走.
就不是這樣了.
應該是我們要有心智.
我們一起祈禱.
多謝天父.
讓我們能夠完成.
這個先知的系列.
歸回的歷史和先知書.
盼望這六卷書.
成為我們大家.
繼續可以去認識.
雖然這些時代已經過去了.
猶太人已經復國.
現在.
都能夠有他們自己的耶路撒冷.
但是現在我們都不明白.

$^{2521}$以巴的戰爭.
帶來的.
各種的問題.
和世界對猶太人的針對.
我們不知道.
什麼時候.
真正可以見到神的將來的日子.
整個以色列.
被醫治和他們得救.
和也有神的.
道在他們的內心.
願主憐憫猶太人.
不同的.
向猶太人傳福音的組織.
求主都興起.
我們自己都有印一些.
希伯來文的新約.
希望猶太人.
可以看 盼望去到.
真正有心的.
猶太人的手中.
他們能夠去看.
去明白彌賽亞就是主耶穌.
他們去悔改.
多謝你們給我們這一段時間.
下一次我們會看.
帖撒羅尼迦的書信.
幫助我們都能夠掌握.
保羅講的末世論.
我們這樣仰望禱告.
奉耶穌基督的名.
阿們.
阿不!.
\newpage



\section{加拉太書 5:13}
\label{sec:AndPB0XPtoQ}
\textbf{2024-03-10 午堂《不被捆綁的自由》加拉太書 5\_13 講員:陳韋安牧師}
\newline
\newline
連結: \href{https://youtube.com/watch?v=AndPB0XPtoQ}{\texttt{https://youtube.com/watch?v=AndPB0XPtoQ}} ~~~~ 語音日期: 2024-03-12
\newline
\newline
\hyperref[sec:2A_RiSlG6Bg]{\small{< < < PREV SERMON < < <}}
~
\hyperref[sec:index_chronic]{\small{[返順時目]}}
~
\hyperref[sec:index_scriptual]{\small{[返順卷目]}}
~
\hyperref[sec:U06UnMRzdJ4]{\small{> > > NEXT SERMON > > >}}
\newline
\newline
加拉太書 5:13
\newline
\begin{longtable}{cl}
\hline
\hline
章節 & 經文 (和合本修訂版)\\
\hline
5:13 & \begin{tabularx}{0.7\textwidth}{X} 弟兄們，你們蒙召是要得自由；只是不可把這自由當作放縱情慾的機會，總要用愛心互相服侍。 \end{tabularx} \\ \\
[1ex]
\hline
\hline
\end{longtable}
$^{1}$丁子敏早安.
今天跟大家一齊思考聖經裡面一個很重要的題目.
但平時我們很少思考的題目.
就是自由的題目.
自由是甚麼呢.
我想要回答這個問題是不容易的.
因為自由其實是一個很廣闊的題目.
很多不同的範疇都會討論.
它可以牽涉哲學,政治,社會,心理,個人成長,神學的題目.
所以今天我們就看看我們的信仰裡面.
要去回答這些問題.
所以今天我們就簡單從一個神學的角度來講一下自由的題目.
所謂自由就是上帝創造我們.
我們擁有上帝的形象.
我們每個人都擁有自己的想法和意願.
這種的自由我們稱之為自由意志.
所謂自由意志其實都不是一些很高深難度的事情.
每個人都追求自由.
因為我們每個人都想去實踐自己的想法和意願.
我有一個想法對自己的想法.
我能夠實踐這個想法在自己身上.
這個就是自由.
是一個我們生命裡面與生俱來的一種能力.
所以嬰兒生下來就嚮往自由.
你知道嬰兒是最萌失的.
基本上是要什麼就要什麼.
得不到就會哭.
所以我們可以說從我們的嬰孩開始.
我們對於自由有一種很直接很強烈的控訴在當中.
當然小朋友的自由是有限的.
因為小朋友是受到父母的管教.
當你成為大人後.
相比起小朋友來說.
你的自由就隨之而增加.
基本上作為成年人.
我想買什麼就買什麼.
想去旅行就去旅行.
今晚想吃雞煲就吃雞煲.
真人雞煲是很好吃的.
所以你很想吃就吃.

$^{41}$當然貧窮都限制自由.
但基本上我們作為成年人的時候.
可以大概按照你的個人想法去過你的生活.
嘗試去實踐你想過的生活.
所以我們都是追求一種自由的生活.
存在主義大師沙特有句很重要的話.
他說存在先於本質.
意思就是說每個人生下來是一個很獨特的個體.
自己以自己的方式去生活.
其他人不能去設定你的生活方式.
所以你的存在.
僅僅就要為你自己的存在去做決定.
不知道大家覺得最自由是什麼.
我覺得最自由.
最自由之在就是我開電單車的時候.
她很不喜歡.
但都給我考了電單車牌.
近兩年我就開始開電單車.
繼陳一鳴牧師之後.
我是另一個開電單車去港島的牧師.
都很方便的.
基本上每個星期就開著電單車.
穿著西裝就去國會港島.
一停在教會裡面就停下來就走.
基本上開電單車是很自由的.
因為可以從塞車裡面得到自由.
見到Tesla還在停泊.
我就穿過Tesla就直接去.
所以是一個非常好的工具.
能夠在中環碼頭去旺角.
都是八分鐘時間就可以去到那裡.
非常之方便.
所以每次當我開電單車的時候.
都感覺到那種自由和奔放.
特別是當天氣好的時候.
在青嶼大橋那裡走.
是一個非常之.
你體會不到的感覺.
那陣風那陣路面的感覺.
是開車是感覺不到的.

$^{81}$只能夠是電單車才能感覺到.
不知道你覺得你最自由的一刻是什麼時候.
感覺到人生裡面.
不過當然剛才我們所說的自由.
只不過是我們所說的其中一部份.
在我們現實的情況之下.
我們所說的自由都會遇見一個問題.
就是自由的限制.
我想除了天堂之外.
天堂是一個真的令到你無拘無束的地方.
但我想當你存活在這個世界上的時候.
所謂自由最主要是一個.
你是面對著很多限制之下那種的自由.
剛才我所說的貧窮是一種限制.
資源短缺是一種限制.
其他人的家所.
其他人的規則都是一些限制.
你的老闆是你的限制.
一條法律都是你的限制.
因此當我們認真去思考這個題目的時候.
所謂自由最主要不單單是一種自由自在的感覺.
當然你能夠從某些事物的束縛中得到解放.
正正就是我們所面對的自由更加重要的部份.
所以free from something.
正正就是我今天所說的一個很重要的課題.
什麼叫做free from something.
我們以前讀中學的時候.
英文課程都經常學的.
那些proposition.
我們講的afraid of.
interest in.
我們都背了一些.
on the bookshelf.
go to school by bus.
一定是既定的某一個詞語.
既定每一個階詞的會是.
所以你說free是答哪一個詞.
正正就是free from.
free from一些東西.
所以我們發覺原來自由.

$^{121}$往往是從某些東西當中.
是能夠得到自由.
最基本的那種例子就是free from the prison.
從監獄裡面那種逃脫.
一個監獄叫你失去自由.
你監獄裡面你被囚禁.
你不能夠去你想去的地方.
你大概都不能夠做你想做的事情.
局限了你能夠做的事情.
所以你就要去free from.
這個枷鎖當中.
這個正正就是free from.
something那個意思.
面對今天這個世界.
自由的更加重要的就是.
你能夠可以去free from something.
那種可能.
就是從某些的限制.
某些的約束裡面得到自由.
從深淵黑暗罪惡的捆綁當中得到自由.
從絕望迷茫當中得到自由.
有很多不同所謂的自由.
所以就是言論自由.
集會自由 經濟自由 宗教自由.
私隱自由 思想自由 心靈的自由.
我們去嘗試去爭取的.
正正是從某些的限制當中.
不被那些枷鎖所捆綁.
所以自由最重要是一種free from的觀念.
free from something是我們基督教教授裡面.
很重要的觀念.
如果你去翻身經的時候.
你會發現聖經所講的自由.
絕大部分都是free from的概念.
都是我們從某些當中得到釋放.
這個正正就是聖經所講的自由.
譬如保羅在《羅馬書》裡面講到.
你們既從罪裡得了釋放.
就做了二的奴僕.
這個free from是甚麼.

$^{161}$就是free from sin.
從罪裡釋放.
保羅說 作基督耶穌裡的生命之律釋放了我們.
使我們脫離罪和死的律法.
就是free from death.
free from the law.
因此聖經所講的自由觀念.
所謂free from something.
無論是從罪中得到釋放.
或者不被律法所綑綁.
或者脫離死亡都好.
都是叫我們能夠從這些的核制當中.
得到那種真正的自由.
耶穌基督叫我們得到自由.
這是我們今天最重要要去明白的事情.
基督耶穌叫我們得到自由.
正如保羅在《加泰書第五章》所講.
基督釋放了我們 叫我們得以自由.
所以要站立得穩.
不要被奴僕的鴉所嚇制.
或者像保羅所講.
天父的兒子若叫你們得自由.
你們就真自由了.
耶穌基督是我們得到自由.
一個最重要和唯一的原因.
不過 主耶穌怎樣叫我們能夠free from那些東西呢.
這只是我們要去思考的問題.
真正的耶穌叫我們得到自由.
是什麼樣的自由呢.
以前我們信主都聽過.
基本上我們都熟悉的.
從罪中得釋放.
從死亡和律法當中得到自由.
這是一種free from death的情況.
不過我們都不是太過賞心.
我們都不以為然.
因為大概都是一些得救的觀念.
free from the sin.
不被罪群綁.
都是一些救恩的事情.

$^{201}$不過其實我們所講.
耶穌叫我們那份的自由.
不單單是叫你得救.
或者上天託那麼簡單.
而是能夠幫我們面對今天的世界.
面對今天的社會或者時代.
耶穌其實是呼召我們.
能夠真正的得著自由.
耶穌叫我們得著什麼樣的自由呢.
free from什麼呢.
其實是anything.
耶穌叫我們從任何事情當中得著釋放.
去享受 賜給我們這種自由.
所以我所講的anything.
真是anything.
耶穌叫我們從一些強權壓迫.
操控當中得著自由.
正如耶穌曾經在加勃隆所講.
主靈在我身上.
祂就叫我傳福音給貧窮的人.
差遣我報告被人得釋放.
乞顏得漢見.
受壓制得自由.
所以耶穌叫我們能夠從一些.
強權的壓制當中得著自由.
耶穌叫我們從深淵黑暗當中得著自由.
不被罪的綑綁.
不被苦毒操控.
要能夠從罪當中釋放.
這是讓我們得著救恩的自由.
不過在眾多的事情裡.
耶穌叫我們更加從一件事裡得著自由和釋放.
今天就說這件事情.
耶穌叫我們能夠從你自己當中得著釋放.
得著釋放和自由.
You can be free from yourself.
從你自己的裡面得著釋放.
今天經文第五章十三節裡.
我們一起讀.
這裡是經文.

$^{241}$預備一二三.
( 經文 ).
菠蘿在經文裡.
再次說明這份在基督裡面的自由是一回事.
菠蘿說我們這份基督裡面的自由.
不是一個放縱情慾的機會.
這裡要再一點點澄清.
所謂放縱情慾的機會.
不是純粹情慾問題.
而是一個occasion of flesh.
就是說是你整個人的血肉的思想和意志.
菠蘿強調的是.
當你能夠從你自己的思想意志意願的綑綁當中釋放.
是一種真正的自由.
真正弟兄姊妹會.
根據罪的緣故.
我們的自由其實並不是真正的自由.
相反我們的自由成為了我們的綑綁.
這是一個很矛盾的範圍.
我們的自由反倒成為我們自己的綑綁.
今天我們很多好像很自由的事情.
很多自由的選擇.
很多好像是我們自己決定的一些選擇.
其實都不一定是自由的選擇.
我們都見過很多有錢的名人.
之前我們都記得Johnny Depp和Amber Heard的官司.
很多有錢的人基本上可以做什麼都可以.
他們在財務上做什麼都可以.
但他們並沒有欣賞自己這份自由.
都是真正的自由.
他們被他們自己的意志和想法去操控.
他們的肉裡面就好像一個模型手一樣.
很簡單地控制他們的情緒和行為.
為什麼呢?.
我們所說的自由意志其實是一個敗壞了的自由意志.
明明是我決定的.
但我沒有辦法去拒絕.
所以今天我們有很多沒有意識的行為.
一些很不理性的行為.
一些甚至是很病態的行為.

$^{281}$表面上是我們自己決定的.
實際上我們不能夠真正去做選擇.
我們通常開車一鎖車.
通常我們啪完之後都會病態地開一開再啪.
你控制不了的.
總是覺得要生一次.
那些潔癖和很多不同的行為.
其實都不是我們能夠控制的.
你在床上不斷地滾動手機.
滾動了很久.
你很後悔因為很浪費時間.
但你發覺都不是自己能控制到的事情.
表面上這些是我們自由的選擇.
但其實你是沒有選擇.
你不能夠不去做.
你知道很浪費時間.
但你沒有辦法.
這個是你沒有這樣的自由.
不去這樣做.
情況更加嚴重的.
這些不自由行為成為了一些很沉溺的行為.
沉溺於購物.
不斷在淘寶買東西.
去到網購.
你發覺你是停不了.
沉溺於酒精.
沉溺於性.
沉溺於很多很多不同的事情.
這些是我們做不到走出來.
事實上這些沉溺的行為.
表面上是我們的選擇.
我們自己好像喜歡做的事情.
但這些是一些沒有自由的行為.
我們以為這些是我們很快樂的決定.
但其實我們被這些事情束縛住.
成為了老僕.
我們以為我們隨時都可以不做.
但我們不能自拔.
我們以為能夠找到出路.
但我們被這些行為捆綁住自己.

$^{321}$然後不知道是不是你的情況.
或者是被你的脾氣所操控.
你發脾氣.
戰戰鬥鬥我不能夠去控制自己.
一下失控之後.
就失去了自由去不做.
這個發脾氣的舉動.
或者你是被很多的恐懼.
或者焦慮所束縛住.
這些焦慮我無法不去想.
無法不去擔心.
這些其實都是我們不能夠有自由.
去自由自在這些事情.
甚至乎在一段關係裡面.
我們被其他人的想法和期望.
其他人的眼光束縛住自己.
我們無法去真開離開別人的眼光.
反過來.
或者很多所謂的怪獸家長.
他們無法離開自己對子女的捆綁住.
覺得這樣才好.
這樣才是這樣.
他們無法有一個自由.
能夠離開自己對別人的期望.
聽姐妹.
主耶穌基督.
祂要賜予一個最大的自由給我們.
你可以選擇做.
你更加有自由.
可以不去這樣做.
真人不做一些事情.
比做一些更加大的自由.
耶穌基督給我們的自由.
是一份更有能力的自由.
祂叫我們去不做某些事情.
即是保羅說.
住在我裡面的.
就是肉體中沒有善.
因為立志為善而不得我.
善而不得我.

$^{361}$只是行出來而不得我.
基督耶穌所賜給我們的自由.
正是從一個不能夠不做當中.
得到這個釋放.
我們可以去選擇不做一些事情.
我們不被那些任何事情束縛住自己.
不被事情去牽動.
不被任何事情去引誘.
不被任何事情去影響.
即是有能力去買一個極為名貴的手袋.
是一種自由.
但有能力去不買.
是更加厲害的自由.
這種能夠可以.
free from這種想法.
這種慾望的自由.
是更加厲害的自由.
真的不知道你面對著什麼困擾.
往往這些困擾都是源於你自己.
可能都是.
真的等於最捆綁著你的人.
往往都是你自己.
唯有你自己才有最能力去捉住你自己.
主耶穌正正是要來.
是要去釋放我們.
去醫治我們.
去要捉住你的手去放下.
不要再去克制你自己.
放過自己.
正正是耶穌基督來.
首先要我們去做到的事情.
當我們能夠得到這份自由的時候.
我們就能夠在這份自由之上.
去自由地去面對這個世界.
當你能夠在基督裡面去free from yourself.
你就能夠去free from anything.
面對著今天這個世界.
面對著這個烏煙瘴氣.
經濟衰退.
很多擔心的年代.

$^{401}$當我們能夠真正得到自由的時候.
我們就發覺.
其他外在的環境都不會影響到你.
這個社會這個世代.
所有人.
你的同事.
現在的股市.
現在的樓價.
現在的情況.
你的敵人.
你的朋友.
你的父母.
都沒有東西能夠真真正正去束縛著你.
因為你是自由的.
當你再一次去.
我們有一次一個很大的button.
你隨時都可以再一次去按這個button.
就可以reset.
重新再來過的時候.
主耶穌就給我們一個這麼大的button.
讓我們隨時都可以再按一下.
重新再開始.
不被任何的事情束縛著.
可能我們是會擔心的.
面對著我們現在的經濟環境.
很多很多令我們很擔心的地方.
但是主耶穌給我們一個很大的一個這樣的按鈕.
因為祂叫得到自由.
當我們隨時有按鈕.
祂隨時去按的時候.
我們就能夠可以重新的.
不被這些事情所綑綁住.
我們很熟悉的經文.
很多的很安慰的經文.
洪水泛濫之時.
耶和華坐著為王.
我雖然行過四季幽谷.
也不怕遭害.
白日太陽不必不相理.
日間月不不相理.

$^{441}$這些經文可以說是耶和華上帝保護你的經文.
但是這些經文更加可以說是一個自由人.
唯有一個真正自由的人.
一個在上帝裡面不被任何事情所影響的人.
才能夠見到這個世界.
這個不關悲觀樂觀的事.
完全是關你是不是一個自由的人.
洪水泛濫之時.
耶和華坐著為王.
這個不是因為你樂觀.
是因為什麼呢.
因為你是自由的.
外面是打風打雨.
亂七八糟也好.
你都能夠可以不被影響.
這個正正就是耶穌基督給我們在世的自由.
自由不單單是叫你上天堂.
自由更加是能夠在你現世裡面.
靠著這份福音的力量.
可以去面對這個烏煙瘴氣的世界.
最後我再說一個更加高層次的free.
你發覺free字有另一個proposition.
除了free from之外還可以free什麼.
一個更加高層次的proposition.
就叫做free for.
free for是一個最高層次的自由.
當你學會了free from之後.
你就可以繼續去學習free for.
你發覺free from都很好.
但你發覺躺平都是free from.
基本上我躺平.
我就完全不理這個世界.
我就放棄.
躺平都是free from.
基本上你就自由.
因為我都不被任何世界的事情所影響.
躺平是自由.
但它是一種對於世界的冷漠.
或者一種放棄的態度.
一個冷漠的人是自由的.

$^{481}$但這份自由沒有了愛在當中.
但主耶穌賜給我們的自由.
祂是指向著愛.
回到今天的經文.
經文裡面提到.
弟兄們你們夢照是得著自由.
只是不可將你們的自由放棄.
或者機會總要用愛心互相服侍.
經文的後半段.
本來從free from轉變到free for.
你們夢照是要得著自由.
為的是什麼.
為的是用愛心去互相服侍.
你的自由.
叫你可以真正去愛.
當你學習了free from之後.
你就能夠用這份自由去free for.
心運.
老實說雞煲是很好吃的.
我很喜歡吃雞煲.
但我這十多年來.
我吃的超過五次雞煲.
因為我的太太是一個非常介意熱氣的人.
她在全球有八成食物都是熱氣的.
所以我沒什麼機會和她吃雞煲.
通常都是和那些學生和小組員去吃雞煲.
但你會發覺.
因為愛的緣故.
你可以自由地去選擇.
去放棄你的自由.
我看海寧當年也跟我去德國讀書.
兩夫妻在德國裡面捱到哭.
很辛苦.
他也放棄他的自由.
他也跟著我去德國讀書.
所以當你去愛一個人的時候.
你就會選擇放棄一點點的自由.
不過這也是一個很自由的選擇.
我想每個父母都是一樣.
但你也知道作為父母.

$^{521}$為了小朋友也好.
你也是會甘願放棄你的自由去愛.
基本上你知道生了孩子之後.
那十年你也沒有什麼自由.
但這些都是你自己甘願的選擇.
是你最自由的選擇.
這個正正就是free for.
愛和自由是一件很矛盾的事.
自由是一個能夠free from something.
但愛卻能夠願意去為對方去失去自由.
甘願地去付出.
但我們的自由.
主耶穌基督教的自由.
正正叫我們用這份自由去愛其他的人.
全部的愛都沒有人逼你的.
如果你因為害怕犯法.
所以去照顧小朋友.
那就是很大件事.
純粹因為這些規條.
純粹因為這些規則.
害怕罰款才去愛的話.
這個就不是愛.
但真正的愛是你自由的選擇.
甘願地去付出.
這個是耶穌基督教我們的自由.
因此讓我們今天去思想.
兩方面.
一個是free from.
有什麼是我們正正被綑綁住的.
可能是很多的壞習慣.
可能是很多的恐懼.
很多你不能不做的事情.
其實你可以不做的.
耶穌基督給我們一個很大的制.
你一按下去.
耶穌給我們得到這份隨時隨地的自由.
你就resume最自由的狀態.
這樣你就可以重新做人.
這是第一個.
好好想想我們有什麼事情.

$^{561}$可以靠著主耶穌去free from.
然後再去思想.
我這份自由.
究竟可以怎樣去愛其他的人.
好好的甘願地付出.
你知這份愛是你甘願的.
是你自由地白白付出.
求主幫助我們去享受這份自由.
謝謝基督.
因為你是賜予自由給我們的上帝.
你給我們有自由.
我們真正有自由.
為我們一群頂智妹去禱告.
求你讓我們得到這份釋放.
不單單是從罪當中得到綑綁釋放.
更加因為這個緣故.
我們能夠不被任何事情所管轄.
無論面對這個社會環境.
面對這個黑暗的世代.
我們仍然能夠有自由在當中去生活.
雖然外面的環境有很多的合制.
但讓我們能夠心裡有這份自由.
更加有這份自由去幫助我們身邊的人.
撥出你給我們的這份呼召.
就是這份擁有自由去彼此相愛.
去幫助弱勢社群的這份愛.
求主幫助我們.
逢主命求 阿們.
\newpage



\section{馬可福音 1:14-20}
\label{sec:U06UnMRzdJ4}
\textbf{2024-03-10 早堂《作主門徒》馬可福音 1\_14-20 講員:董家驊牧師}
\newline
\newline
連結: \href{https://youtube.com/watch?v=U06UnMRzdJ4}{\texttt{https://youtube.com/watch?v=U06UnMRzdJ4}} ~~~~ 語音日期: 2024-03-12
\newline
\newline
\hyperref[sec:AndPB0XPtoQ]{\small{< < < PREV SERMON < < <}}
~
\hyperref[sec:index_chronic]{\small{[返順時目]}}
~
\hyperref[sec:index_scriptual]{\small{[返順卷目]}}
~
\hyperref[sec:pdca3AXE1WI]{\small{> > > NEXT SERMON > > >}}
\newline
\newline
馬可福音 1:14-20
\newline
\begin{longtable}{cl}
\hline
\hline
章節 & 經文 (和合本修訂版)\\
\hline
1:14 & \begin{tabularx}{0.7\textwidth}{X} 約翰下監以後，耶穌來到加利利，宣講神的福音， \end{tabularx} \\ \\ \relax
1:15 & \begin{tabularx}{0.7\textwidth}{X} 說：「日期滿了，神的國近了。你們要悔改，信福音！」 \end{tabularx} \\ \\ \relax
1:16 & \begin{tabularx}{0.7\textwidth}{X} 耶穌沿著加利利的海邊走，看見西門和西門的弟弟安得烈在海上撒網；他們本是打魚的。 \end{tabularx} \\ \\ \relax
1:17 & \begin{tabularx}{0.7\textwidth}{X} 耶穌對他們說：「來跟從我，我要叫你們得人如得魚一樣。」 \end{tabularx} \\ \\ \relax
1:18 & \begin{tabularx}{0.7\textwidth}{X} 他們立刻捨了網，跟從他。 \end{tabularx} \\ \\ \relax
1:19 & \begin{tabularx}{0.7\textwidth}{X} 耶穌稍往前走，又見西庇太的兒子雅各和他弟弟約翰在船上補網。 \end{tabularx} \\ \\ \relax
1:20 & \begin{tabularx}{0.7\textwidth}{X} 耶穌隨即呼召他們，他們就把父親西庇太和雇工留在船上，跟從了耶穌。 \end{tabularx} \\ \\
[1ex]
\hline
\hline
\end{longtable}
$^{1}$弟兄姊妹 主日早安.
早安.
弟兄姊妹 主日早安.
這次來到香港跟許多的牧者聊天.
跟許多的牧者聊天.
有一次有一個牧者跟我說.
當他反省過去這十年教會的經歷的時候.
有一個牧者跟我說.
他反省這十年的工作的時候.
他說當教會人數最多的時候.
反過來看發現的是最黑暗的時候.
他發現人數最多的時候.
是最黑暗的時候.
弟兄姊妹不要誤會我的意思.
教會人多很好.
大家不要誤會我的意思.
教會人多是很好的.
只是會不會一不小心.
我們就在維持組織的運作當中而耗盡.
但是會不會是因為這樣.
我們只是維持一些組織.
而耗盡我們自己.
服侍上帝很好.
但若不謹慎的話.
我們會不會在忙碌的事工當中自我膨脹.
我們會不會在忙碌的生活當中自我膨脹.
以至於在服侍當中我們傷害了別人.
我們也覺得別人傷害了自己.
今天我們回到《馬可福音》第一章.
重新回到耶穌呼召門徒的現場.
好像校正我們的生命.
再一次來看上帝到底呼召我們要怎麼前行.
再次去看看上帝呼召我們的時候要怎麼前行.
這段經文其實也提醒我們.
這段經文也提醒我們.
我們所傳的福音 領受的福音.
其實比我們想像的更有高度.
更有廣度.
也更有深度.
怎麼說呢?.

$^{41}$第一個第十四節我們看到耶穌宣講上帝的福音.
首先在十四節看到耶穌宣講福音.
耶穌說了什麼?.
耶穌說了什麼呢?.
耶穌說上帝的國近了.
很多時候我們把福音等同於個人的得救.
我們有時候把福音等同於個人的得救.
我們傳福音的時候講的是我們要認罪悔改.
所以有一天我們可以上天堂.
我們傳福音的時候是認罪悔改一天我們可以上天堂.
可是耶穌好像不是在講一個上天堂的福音.
但耶穌好像不是在講上天堂的福音.
耶穌是講一個上帝的國臨到的福音.
耶穌是在講一個上帝的國臨到的福音.
其實面對今天人生很多時候的困難.
面對我們今天很多人生的困難.
我們往往有兩種反應.
第一個是把問題歸咎給別人.
都是別人的問題.
都是大環境的問題.
另外一個極端則是總是覺得自己做不好.
另外一個極端就是總是覺得自己做得不好.
但是其實不論是歸咎於別人還是自己.
但是其實不論你是歸咎於別人或者歸咎於自己.
這兩種面對困難的方法有一個共同點.
這兩個面對困難的地方有一個共同點.
都是人就以我為中心.
都是以自我為中心.
但其實每一個危機,衝突跟挑戰.
但其實每一個危機跟挑戰.
都是上帝介入我們生命的一個機會.
我們不單單是只在當中看到我怎麼樣可以得救.
我們不單單是看到我們自己怎麼樣可以得救.
而是看到上帝的心意是什麼.
而是看到上帝的心意是怎麼樣.
幾週前我在默想經文啟示錄九章二十節的時候.
幾個禮拜之前我在默想啟示錄.
九章二十節.
九章二十節的時候.
那段經文講到上帝的審判臨到的時候.

$^{81}$這裡講到上帝的審判臨到的時候.
可是許多的人仍舊不肯悔改.
但是很多人都不願意悔改.
仍舊去拜魔鬼 鬼魔.
仍舊去拜鬼魔.
去拜那些偶像.
那天我在默想的時候就在問上帝.
上帝啊 我在敬拜的偶像是什麼.
那天我就在默想的時候問上帝.
我在敬拜的偶像是什麼呢.
大概在安靜七八分鐘之後.
大概安靜了七八分鐘之後.
有三個字越來越清楚浮現在我的腦海當中.
是什麼呢.
那三個字是.
這三個字是.
董家華.
這三個字就是董家華.
真的.
我發現很多時候我的情緒反應跟焦慮不安.
都是因為我太在意珍惜董家華這三個字.
很多時候我焦慮不安.
就是太在意我自己一個董家華這三個字.
我把這三個字.
這個人跟所代表的當作我敬拜的對象.
我把這三個字.
我人的代表就是成為我的偶像.
敬拜的偶像.
其實我們這個時代很容易把自己當作偶像.
很多時候我們把自己當作偶像.
但是神的福音提醒我們.
我們不是世界的中心.
就算有一天我身敗名裂.
但是在上帝的國度裡人就是有盼望的.
但是在神的國度裡我仍然是有盼望的.
正是這樣的盼望.
正是有這個盼望.
讓我們在面對人生的風浪的時候.
不再聚焦在自己的身上.
而是能夠看到上帝國度的豐富.

$^{121}$這是第一點.
這好消息是關於上帝的國.
當然也關乎我們的得救.
但是當我們迴轉離棄偶像歸向上帝的時候.
要我們迴轉離棄偶像歸向上帝的時候.
我們經歷了國度的美好.
我們就會經歷國度的美好.
第二個我想講到福音的廣度.
第十六節 耶穌沿著加利利海邊走的時候.
看見西門和西門的弟弟安德烈在海上撒網.
他們本來是打魚的.
什麼意思呢.
耶穌走到這些人所在之處.
在他們所在之地.
呼召他們跟隨.
在他們所在之地.
呼召他們跟隨主.
今天我們常覺得跟隨耶穌是要在特定的場合.
很多時候我們覺得.
耶穌呼召我們的時候應該在特定的場合.
好像來到堂會裡要跟隨耶穌.
講道的時候要跟隨耶穌.
服事的時候要跟隨耶穌.
可是做生意的時候呢.
但是做生意的時候呢.
工作的時候呢.
管教兒女的時候呢.
跟爸爸媽媽不耐煩溝通的時候呢.
很多時候我們的生活跟信仰是撕裂的.
我們週日在教會的生活.
與我們週間在職場跟平日在家的生活.
很多時候是撕裂開來的.
但耶穌沒有要門徒離開他們的處境來跟隨他.
耶穌是首先走入門徒的世界.
並且呼召他們在那裡跟隨.
服事上帝跟隨耶穌當然需要在教會裡面.
服事上帝跟隨耶穌當然在教會.
但也在我們的家庭裡面.
也在我們的職場裡面.
也在我們的人際關係當中.

$^{161}$做主耶穌的門徒如果不包含這些層面.
做主耶穌的門徒如果不包括以上所說的.
我們大概不是做基督徒.
我們大概只是做一個形式的門徒而已.
我們只是做一個形式的門徒而已.
福音是廣到包含我們生命的全部部分.
而不是我們生活的某一個面向而已.
大概八年前.
那時候我還住在美國木匯的時候.
那時候我們領受了一個.
我跟我太太有一個感動.
那時候我們領受到跟我太太有一個感動.
就是要回到亞洲來住.
回到台灣去服侍.
其實我跟我太太是很樂意的.
我也很期待.
我唯一的一個掙扎是.
我的小孩子的教育.
我小朋友的教育.
因為我知道回到台灣.
我們是讀不起英語學校的.
我知道我們回到台灣.
我們是讀不起那些英語的學校的.
只能讀一般的公立學校.
那我是在台灣公立學校長大的.
我自己是在台灣的公立學校長大的.
我恨死了台灣的教育.
我到三十歲的時候都還會做噩夢.
夢到考試.
我到三十歲我都夢到要考試.
然後我看到我身邊的很多朋友.
都是努力的搬到美國.
讓孩子接受另外一種教育.
我看到我身邊的朋友都是很努力的.
把小孩搬到海外讀書.
唯一的掙扎是.
我難道要把我的孩子帶回去.
我以前不喜歡的教育系統嗎.
我唯一的掙扎是.
難道我要把我的孩子.

$^{201}$帶回我以前不喜歡的舊約裡面嗎.
有一天當我在跟上帝對話的時候.
有一天我跟上帝對話的時候.
我覺得勝靈突然有一個感動問我.
我覺得勝靈有一個感動去跟我說.
說你能夠給你孩子最好的是什麼.
你可以給你孩子最好的是什麼.
我現在想是一棟房子嗎.
一人一棟嗎.
我想是不是給他們一人一棟房子呢.
可是我在美國看過太多的敗家子.
我在美國看到太多的敗家子.
我覺得給房子也無法保證他的未來.
我覺得給房子他們.
都不可以保證他們的未來.
我說是一個好的教育環境嗎.
你說是一個好的教育環境嗎.
可是我也看到每個人都很不一樣.
但我都看到每一個人都不同.
有些人在這種環境當中會長得很好.
有些人在這種環境當中可以長得很好.
有些人適合另外一種環境.
其實我也不知道我的兒子適合哪一種環境.
最後我發現我不知道.
我可以給我孩子什麼是最好的.
最後我發現我不知道要給我孩子什麼是最好的.
有一個意念就出來.
有一個意念就出來了.
說你能夠給你孩子最好的.
你可以給你孩子最好的.
就是跟隨耶穌的榜樣.
因為我們不知道未來會發生什麼事.
而且你的兒子不會是完美的.
他將來也會遇到挫折.
也會犯錯.
甚至可能會離開教會一段時間.
甚至會離開教會一段時間.
但是當他們在父母身上看到跟隨耶穌的榜樣.
當他們在父母身上看到跟隨耶穌的榜樣.
不是父母完美的生命.

$^{241}$而是在父母的軟弱當中去經歷上帝的恩典.
而是在父母的軟弱當中去經歷上帝.
看到上帝的真實.
或許當他們遇到困難的時候.
他們更有機會回到上帝的面前.
我發現那才是能夠給下一代最好的.
所以後來就跟我太太決定.
我們就搬回到台灣.
後來就跟我太太決定搬到台灣.
跟隨耶穌不是只是在教會堂會裡面.
而是在我們的家庭裡.
也是在我們的職場工作.
也是在我們的職場工作裡面.
也是在我們的親戚關係當中.
複雜的人際朋友關係當中.
不斷地來問耶穌.
你的心意是什麼.
我如何能夠在這些關係當中.
來衷心地跟隨你.
這是福音的廣度.
我們看到福音是關乎上帝的國.
我們看到福音是關於上帝的國.
不是關乎自我的實現.
不是關乎我們自我的實現.
而正因為我們看到上帝國的寬廣.
我們裡面有一種自由跟釋放.
福音的廣也是關乎生活的每一個面向來跟隨.
福音的廣都是關於我們生活每一個面向裡面.
接下來我們要來看福音的深度.
第十七節 耶穌找到這兩個人之後.
對他們說什麼呢.
第十七節 耶穌找到這兩個人之後.
對他們說什麼呢.
耶穌說 來跟從我.
關鍵字 我要叫你們得人如得魚一樣.
關鍵的字就是 我要叫你們得人如得魚一樣.
我們常常把焦點放在 來跟從我.
是對的 我們被呼召要跟隨耶穌.
可是我們常忽略了第二句話.
耶穌說 我要叫你們.

$^{281}$什麼意思呢.
跟隨耶穌的動力跟能力.
其實也是來自上帝自己的行動.
跟從耶穌的動力 都是來自上帝的行動.
福音從一開始到結束.
都是關於上帝的作為.
福音由開始到現在結束.
都是上帝的作為.
華人有一句俗話.
說師父帶進門 修行看個人.
師父帶進門.
帶入門.
修行看個人.
修行看自己.
好像一開始是一個恩典.
能夠進到上帝的國裡面.
看來好像恩典可以進入上帝的國.
但是後面你在這個國裡面.
可以做到多高.
那是看你自己的修行了.
但是後面你的國.
你要多高呢.
那就是看你自己的修行.
可是這不是我們所領受的福音.
但這不是我們所領受的福音.
我們之所以能夠進上帝的國.
是因為恩典.
我們所以可以進入上帝的國.
是因為上帝的恩典.
我們能夠在國度裡面持續的成長.
也是因為耶穌說.
「我要叫你們」.
我們在這個恩典裡面成長.
都是耶穌說我要叫你們.
所以自始至終是上帝恩典的行動.
坐在我們身上.
所以自始至終.
都是上帝的恩典在我們身上.
因此一個所謂屬靈的人.
是不會去自誇自己屬靈的.

$^{321}$所以一個屬靈的人.
不會自誇自己是屬靈的.
因為他知道一切都是恩典.
通常會覺得自己很屬靈.
然後覺得很驕傲的人.
通常都會覺得自己很屬靈.
然後覺得自己很驕傲.
可能真的不明白恩典.
可能不是真正的明白恩典.
而我們怎麼樣持續讓上帝的行動.
在我們生命當中更新我們呢.
我們怎麼樣持續讓我們生命.
是被神去更新呢.
就是透過屬靈的操練.
你說董牧師你不是自相矛盾嗎.
那你會問董牧師.
那不是自相矛盾.
你剛剛不是說.
這一切都是上帝的作為.
你剛才不是說.
這些都是上帝的作為.
怎麼說又要屬靈操練.
為什麼又要說屬靈操練.
又要讀經 又要禱告 又要靈修.
讓我舉個例子.
我講個例子.
大家今天起床到現在.
應該都有用過水龍頭吧.
大家由起床到現在都有用過水龍頭.
我想問一個問題.
請問我們打開水龍頭.
或自動感應水之所以流出來.
我們去打開水龍頭.
或者去感應水流出來.
是因為你的功勞還是別人的功勞.
是你的功勞還是別人的功勞呢.
大概不是你的功勞.
當然我們當中如果有水利工程師.
這可能有一點點是你的功勞.
當然我們當中如果有水利工程師.

$^{361}$可能有一點點你的功勞.
但是大部分的我們.
這跟我們沒有任何關係.
但是大部分的人.
都是跟我們沒有什麼關係.
是別人過濾了水.
別人建立管線把水引到我們的家.
別人去建立了這個管道.
然後我們家可以用到水.
別人設計了整個系統.
別人去設計了整個系統.
然後別人裝了那個水龍頭.
別人去裝了那個水龍頭.
我們只不過是使用它.
我們只不過是使用.
這是恩典不是嗎.
這是不是恩典呢.
好 那當你要用水龍頭的水的時候.
你要做什麼.
當你要用這個水龍頭的時候.
你要做什麼呢.
你至少要去感應一下.
或者是傳統的就是要打開它.
或者傳統的就是打開它.
感應跟打開就是屬靈的操練.
因為感應跟打開就是屬靈的操練.
我們讀經跟禱告.
我們讀經跟祈禱.
其實就是把手放在那個位置去感應.
就是把手放在那裡去感應.
就只是打開那個水龍頭去領受.
只是打開那個水龍頭去領受.
讀經跟禱告就是這麼一回事.
讀經跟祈禱就是這樣一回事.
本身沒有什麼好誇口的.
本身是沒有什麼好誇口的.
而是放在一個領受的位置.
當我們真的明白這一點.
當我們真的明白這一點的時候.
我們會開始發現讀經跟禱告是一件很享受的事情.

$^{401}$我們會發現讀經跟祈禱是一件很享受的事情.
而且不會讓我們產生比較競爭的心態.
而且不會讓我們產生比較的.
比較跟競爭的心理.
就好像你不會去跟別人比較.
今天你用過多少次水龍頭.
好像你不會跟別人比較.
你用了多少次水龍頭.
你不會說我今天用了五次.
你用了幾次.
你們不會問.
我今天用了五次你用了幾次.
不太行啊你只用了一次.
你不行了你只用了一次.
我用了五次啊.
你看我現在生命多厲害.
我們不會這樣.
因為我們明白能夠用水龍頭是恩典.
我們明白能夠用水龍頭是恩典.
跟隨是一個回應耶穌的行動.
而操練則是幫助我們在每日生活當中去領受上帝的工作.
領受上帝在我們身上的工作.
以及在我們當中的工作.
這是福音的深度.
不單單是改變我們外在的行為.
也改變我們內在的情緒.
我們其實今天非常需要去認識我們的情緒.
為什麼呢.
因為我們的情緒反映出我們真正所愛的是什麼.
我們的情緒就反映出我們真正所愛的是什麼.
因此我們會有正面的情緒.
我們會有負面的情緒.
很多時候我們用想法.
想法是可以騙人的.
很多時候我們用想法.
想法是可以騙人的.
也可以騙自己.
也可以騙到自己.
多少時候我們裡面其實很枯乾.
很多時候我們裡面其實很枯乾.

$^{441}$一點都沒有平安.
可是見到其他人的時候.
但是見到其他人的時候.
還是會說什麼.
都會說什麼呢.
平安.
你好嗎.
都好.
我好.
生意怎麼樣.
不錯.
生活怎麼樣.
非常棒.
非常好.
可是其實我們裡面是枯乾的.
其實我們裡面是枯乾的.
但是當我們很願意誠實面對我們真正的情緒的時候.
情緒好像一個指南針.
幫助我們看到生命真正的狀況.
以及我們到底所愛的是什麼.
而透過屬靈的操練.
透過屬靈的操練.
我們不斷地被上帝改變塑造.
所以我自己在讀經禱告前有一個習慣.
我自己在讀經祈禱之前有一個習慣.
我會先花幾分鐘跟上帝講講我心裡面的感覺.
上帝啊 昨天我其實蠻難過的.
上帝啊 我昨天是蠻難過的.
某某人講了一句話讓我非常的覺得被冒犯.
上帝啊 我今天自我感覺蠻良好的.
上帝啊 我今天覺得自己是自我良好的.
因為在事工上發生了一些重大的突破.
因為在事工上發現了一個突破.
上帝啊 我覺得我老婆今天早上看我的眼神怪怪的.
上帝啊 今天老婆看我的眼神怪怪的.
我覺得她是不是心裡面有不舒服的地方.
我覺得她心裡是不是有些不舒服的地方.
當我們願意把自己向上帝敞開的時候.
這個時候的讀經禱告操練才會進到我們內心的深處.
這個時候的讀經祈禱才會進到我們內心的深處.

$^{481}$你有沒有發現很多時候我們讀經禱告.
好像讀完之後不知道讀什麼.
你有沒有發現我們很多時候讀經祈禱.
我們不知道讀了什麼.
是因為我們沒有先向上帝敞開自己.
是因為我們沒有先向上帝去敞開我們自己.
就好像拿了一個茶杯去裝水.
水是流下來了.
水是流出來了.
蓋子沒有打開.
但我們的蓋沒有打開.
所以沒有領受.
這是第三個福音的深度要進到我們的內心.
這是第三個福音的深度要進到我們內心.
而當耶穌呼召門徒之後.
對他們說你們要得人如得魚了.
你們要得人如得魚.
門徒領受了一個新的使命.
其實每一個被上帝觸碰改變的人.
其實每一個被上帝觸碰的人.
改變的人.
也都同時領受了一個新的使命.
都同樣的.
是同時去領受一個新的使命.
他們的工作生活的終極目的改變.
他們的工作的終極是改變了.
彼得和安德烈本來是打魚的.
安德烈他們本來是打魚的.
但是他們生命的終極目的改變了.
他們領受了一個新的使命.
很多時候我們以為領受.
被上帝呼召就是離開我的職場.
很多時候我們以為領受.
被上帝呼召就是離開我們的職場.
因為一章十八節說.
他們聽了耶穌的呼召就捨了往跟從他.
在一章十八節裡面.
說到他們領受之後.
他們就捨了往跟從主.
我們常把這段理解為.

$^{521}$是耶穌呼召門徒離開他們正常的職場.
我們聽到這句話的時候.
好像是離開他們的職場.
但我認為這一段更好的理解是.
不是離開他們的職場.
而是轉換他們工作的終極目的.
以及所倚靠的事物.
看耶穌是說你們要得人如得魚.
耶穌說你們要得人如得魚.
也就是說.
耶穌其實並沒有否定他們過去的工作.
耶穌沒有否定他們過去的工作.
耶穌沒有否定他們過去作為漁夫的生活.
但是耶穌賦予他們現在的工作.
一個不一樣的使命.
但是耶穌賦予他們一個不同的使命.
以至於他們改變了生命的使命跟目標.
以至於他們改變了他們的使命跟目標.
而當使命跟目標改變的時候.
當使命跟目標改變的時候.
也很自然我們用的工具也改變.
很自然我們用的工具也改變了.
今天耶穌基督的門徒.
我們也都領受了這樣的使命.
今天基督耶穌的門徒.
都同樣去領受同樣的使命.
我們領受的使命是什麼呢.
大家立刻想到大使命.
沒有錯.
但我要說.
不完全.
不完整.
我要說是不完全不完整.
大使命是什麼.
大使命是什麼呢.
敘使萬民做主門徒.
是重要的.
是不可少的.
可是如果一個人他只是一天到晚叫別人來教會.
跟隨耶穌做門徒.

$^{561}$但一個人只是叫別人來教會去做主門徒.
你會不會覺得好像少了一種溫度.
你會不會覺得少了一種溫度.
好像不太識人間煙火.
好像不識人間煙火.
你也不能說不對.
你也不可以說不對.
可是總覺得少了一些什麼.
但好像少了什麼.
是的 是少了一些什麼.
少了什麼呢.
少了上帝給人類的一個其他面向的使命.
少了一個上帝給人類的其他面向的使命.
少了一個上帝給人類的一個其他面向的使命.
我用三個字詞來形容這整全的使命包含哪三個.
第一個是大使命 帶人做門徒.
第二個是大誡命.
愛靈社如同愛自己.
去關愛我們的靈社.
當門徒問耶穌說就約的總結是什麼的時候.
耶穌總結是愛上帝 愛靈社.
因此大誡命簡單來說.
就是關愛我們周遭的靈社.
上帝你把誰帶到我面前.
今天上帝將誰帶到我面前.
就好好地去回應關愛.
過去我在幕會的時候.
每個禮拜天回到家裡的時候.
有一段時間我覺得我的頭都昏了.
為什麼呢.
因為那天見到幾百個人.
每個人都想去關心一兩句.
可是在關心這個人的時候.
但在關心這個人的時候.
後面看到另外一個一個月沒出現的朋友.
後面又看到一個一個月沒見的朋友.
就很想趕快結束這個對話.
去關心那個人.
好不容易.
結束了去關心那個人的時候.

$^{601}$又看到另外一個直視出現.
又看到另外一個執事出現.
想到等一下開會要先跟他會前溝通的事情.
想到等一下執事會議跟他會前溝通.
然後就匆匆地結束跟這個人的談話.
結束跟這個人的談話.
然後去找那個人.
好像花蝴蝶一樣.
回到家我都忘記今天跟誰講話.
後來我學習一個操練.
就是上帝.
每個禮拜天.
每一堂崇拜.
求你讓我看到你把哪一個人帶到我面前.
求你幫助我看到哪一個人你將他帶到我面前.
讓我用那一天那個時間.
讓我用一天那個時間.
好好地回應那一個人.
好好地回應那個人.
我發現那改變了我禮拜天.
事奉的心情跟態度.
我發現那改變了我禮拜天事奉的態度.
這是大皆命.
關愛靈舍.
最後一個是大任命.
就是上帝一開始創造人類的時候.
賦予人管理,看守世界的責任.
賦予人看守,管理這個世界.
並且要我們生養眾多.
我用白話文來說.
就是創造美善的使命.
透過這美善.
使人得認識看到上帝榮耀的一部分.
上帝榮美的一部分.
當我們持續在生活中創造美善.
關愛靈舍.
帶領人做門徒變成一個很自然的事情.
帶領人去做門徒變成一個很自然的事情.
我不會說它會是一件容易的事.
我不可以說它會是一件容易的事.

$^{641}$但是會是一個自然的事情.
是很自然的事.
這是我們所領受的整全的使命.
因此今天最後我們要問自己.
上帝在我的生活當中.
我的職場家庭.
我的教會生活.
我的人際關係.
我如何回應你.
創造美善.
關愛你帶到我面前來的人.
並把握機會帶他們跟隨認識你.
並把握機會帶他們認識跟隨你.
我們一起祈禱.
天父謝謝你.
是你把我們從這世上呼召出來.
並且在這個世界.
帶著一個全新的使命而活.
不論我們的年紀如何.
不論我們有的資源是多是少.
不論我們的資源有多或少.
都是你給的.
也都是恩典.
幫助我們持續使用這一切獻給你.
在人生的路上.
持續地創造美善.
關愛靈舍.
並帶人來認識跟隨你.
帶人來認識跟隨你.
奉耶穌基督的名禱告.
Amen.
\newpage



\section{詩篇 139:1-13-18-23-24}
\label{sec:pdca3AXE1WI}
\textbf{2024-03-09 晚堂《在靈修和生活中經歷神》詩篇 139\_1,13-18,23-24 講員:廖炳堂牧師}
\newline
\newline
連結: \href{https://youtube.com/watch?v=pdca3AXE1WI}{\texttt{https://youtube.com/watch?v=pdca3AXE1WI}} ~~~~ 語音日期: 2024-03-13
\newline
\newline
\hyperref[sec:U06UnMRzdJ4]{\small{< < < PREV SERMON < < <}}
~
\hyperref[sec:index_chronic]{\small{[返順時目]}}
~
\hyperref[sec:index_scriptual]{\small{[返順卷目]}}
~
\hyperref[sec:RZv3JKPRqaw]{\small{> > > NEXT SERMON > > >}}
\newline
\newline
詩篇 139:1-13-18-23-24
\newline
\begin{longtable}{cl}
\hline
\hline
章節 & 經文 (和合本修訂版)\\
\hline
139:1 & \begin{tabularx}{0.7\textwidth}{X} 耶和華啊，你已經鑒察我，認識我。 \end{tabularx} \\ \\ \relax
139:2 & \begin{tabularx}{0.7\textwidth}{X} 我坐下，我起來，你都曉得；你從遠處知道我的意念。 \end{tabularx} \\ \\ \relax
139:3 & \begin{tabularx}{0.7\textwidth}{X} 我行路，我躺臥，你都細察；你也深知我一切所行的。 \end{tabularx} \\ \\ \relax
139:4 & \begin{tabularx}{0.7\textwidth}{X} 耶和華啊，我舌頭上的話，你沒有一句不知道的。 \end{tabularx} \\ \\ \relax
139:5 & \begin{tabularx}{0.7\textwidth}{X} 你前後環繞我，按手在我身上。 \end{tabularx} \\ \\ \relax
139:6 & \begin{tabularx}{0.7\textwidth}{X} 這樣的知識奇妙，是我不能測的；至高，是我不能及的。 \end{tabularx} \\ \\ \relax
139:7 & \begin{tabularx}{0.7\textwidth}{X} 我往哪裡去，躲避你的靈？我往哪裡逃，躲避你的面？ \end{tabularx} \\ \\ \relax
139:8 & \begin{tabularx}{0.7\textwidth}{X} 我若升到天上，你在那裡；我若躺在陰間，你也在那裡。 \end{tabularx} \\ \\ \relax
139:9 & \begin{tabularx}{0.7\textwidth}{X} 我若展開清晨的翅膀，飛到海極居住， \end{tabularx} \\ \\ \relax
139:10 & \begin{tabularx}{0.7\textwidth}{X} 就是在那裡，你的手必引導我，你的右手也必扶持我。 \end{tabularx} \\ \\ \relax
139:11 & \begin{tabularx}{0.7\textwidth}{X} 我若說「黑暗必定壓碎我，我周圍的亮光必成為黑夜」， \end{tabularx} \\ \\ \relax
139:12 & \begin{tabularx}{0.7\textwidth}{X} 黑暗對你不再是黑暗，黑夜卻如白晝發亮。黑暗和光明，在你看來都是一樣。 \end{tabularx} \\ \\ \relax
139:13 & \begin{tabularx}{0.7\textwidth}{X} 我的肺腑是你所造的，我在母腹中，你已編織我。 \end{tabularx} \\ \\ \relax
139:14 & \begin{tabularx}{0.7\textwidth}{X} 我要稱謝你，因我受造奇妙可畏，你的作為奇妙，這是我心深知道的。 \end{tabularx} \\ \\ \relax
139:15 & \begin{tabularx}{0.7\textwidth}{X} 我在暗中受造，在地的深處被塑造；那時，我的形體並不向你隱藏。 \end{tabularx} \\ \\ \relax
139:16 & \begin{tabularx}{0.7\textwidth}{X} 我未成形的體質，你的眼早已看見了；你所定的日子，我尚未度一日，都在你的冊子寫上了。 \end{tabularx} \\ \\ \relax
139:17 & \begin{tabularx}{0.7\textwidth}{X} 神啊，你的意念向我何等寶貴！其數何等眾多！ \end{tabularx} \\ \\ \relax
139:18 & \begin{tabularx}{0.7\textwidth}{X} 我若數點，比海沙更多；我睡醒的時候，仍和你同在。 \end{tabularx} \\ \\ \relax
139:19 & \begin{tabularx}{0.7\textwidth}{X} 神啊，惟願你殺戮惡人；你們好流人血的，離開我去吧！ \end{tabularx} \\ \\ \relax
139:20 & \begin{tabularx}{0.7\textwidth}{X} 他們說惡言頂撞你，你的仇敵妄稱你的名。 \end{tabularx} \\ \\ \relax
139:21 & \begin{tabularx}{0.7\textwidth}{X} 耶和華啊，恨惡你的，我豈不恨惡他們嗎？攻擊你的，我豈不憎惡他們嗎？ \end{tabularx} \\ \\ \relax
139:22 & \begin{tabularx}{0.7\textwidth}{X} 我恨惡他們到極點，以他們為我的仇敵。 \end{tabularx} \\ \\ \relax
139:23 & \begin{tabularx}{0.7\textwidth}{X} 神啊，求你鑒察我，知道我的心思，試煉我，知道我的意念； \end{tabularx} \\ \\ \relax
139:24 & \begin{tabularx}{0.7\textwidth}{X} 看在我裡面有甚麼惡行沒有，引導我走永生的道路。 \end{tabularx} \\ \\
[1ex]
\hline
\hline
\end{longtable}
$^{1}$敢宣家的弟兄姊妹平安.
很開心有機會來到大家當中.
跟大家一起敬拜.
以及思想神的說話.
我看到你們的連題是.
生命更新 傳揚福音.
大家今年一起努力.
生命更新 個人和集體都一起更新.
然後我們傳揚神的愛.
生命更新其實最基本就是我們和神的關係.
我們傳的福音不是一套理論.
不是一套道理 而是一個活的救主.
在我們生活上 我們相信他之後經歷到的.
一位復活的主.
所以今天跟大家聊聊大家的靈修生活.
靈修生活這個課題都很敏感.
我不知道大家有沒有弟兄姊妹之間.
有時候會彼此關心的時候 關心這部分.
但是這是一個很敏感的話題.
你最近有沒有靈修.
你會不會跟弟兄姊妹這樣說.
但我知道這是一個很關鍵的問題.
所以如果很簡單 問問你.
你滿意自己的靈修生活嗎.
不用大家舉手.
或者再進一步說.
你滿意你現在和神的關係嗎.
如果問弟兄姊妹這樣的問題.
關心她的時候.
通常弟兄姊妹會覺得很抽象.
我和神的關係好不好.
不好呢 怎樣會知道不好呢.
我不知道你會怎樣回應.
你現在和神的關係好不好.
你心裡面會給一個甚麼答案呢.
我沒有犯大罪.
但又沒有甚麼進步.
是好還是不好呢.
這樣就不太清楚.
我每天都有靈修.

$^{41}$但又好像不覺得自己很有動力.
如果沒有靈修.
就覺得這樣都不太好.
不知道你自己怎樣看.
今天和大家思想這個課題.
怎樣能夠在靈修的生活裡面.
和在生活當中.
不單止靈修.
其實我要說的是.
靈修是為了我們生活當中可以經歷到神.
你靈修的時候.
怎樣靈修才能夠經歷到神呢.
今天和大家思想這個很基本的課題.
但又很重要.
復活的主.
所以我們基督教的教會裡面.
通常都有十字架.
十字架是空的.
天主教是很不同的.
空的不是沒有耶穌基督.
而是因為他已經復活了.
所以他現在是活在我們裡面.
不是在十字架裡面.
最怕的是我們心裡面都不知道.
耶穌是不是真的活在我們裡面.
我們經歷不到他的時候.
那十字架就真的空了.
希望不是.
我們是真的有耶穌基督在裡面.
不單止活在這個世界上.
掌管這個世界.
他也掌管和引領著我每天的生活.
希望和大家思想這個課題.
今天用的經文.
就是大家讀了一次.
詩篇139篇.
我選了一些經文出來.
有些回教會很久的弟兄姊妹.
都挺多人喜歡詩篇139篇的.
我不知道大家當中是不是.

$^{81}$他講到神和我們.
是一個很親密的關係.
你首先就留意到.
這個經文.
就是講我和你.
你就是耶和華.
然後私人自己就講.
是我.
然後講我們.
你知道詩篇有些是講我們.
是集體敬拜神.
我們一起去讚美耶和華.
那就是我們.
我們有集體的時候.
我們崇拜都一起集體.
但也有個人的時候.
原來這個詩篇裡面是強調神和我們個人的關係.
第一件事和大家思想.
就是為什麼要靈修.
其實我和大家講三件事.
為什麼要靈修.
什麼叫靈修.
為什麼要做一個好的靈修.
今天講三件事.
為什麼要靈修.
最基本先.
因為神和我們有一個很個人性的關係.
神不單止關心你旁邊的弟兄姊妹.
或者不單止關心.
葛培里或者戴德生.
或者是誰做得很大事工的人物.
每個教徒.
都在眼中很寶貴.
不可以被另一個教徒代替.
所以他不單止.
只想你旁邊的弟兄姊妹有靈修.
有祈禱.
他想聽你有沒有將你的生活.
拖去祈禱.
和上帝交代.

$^{121}$或者和他去依靠他.
他關心你.
詩篇裡面講到.
你已經鑒察我.
認識我.
大家讀完之後有沒有留意到.
第一節的經文和最後兩節.
是一樣的.
都是講神鑒察和認識我.
這個個人.
最後兩節.
你記得.
23,24節最後的祈禱是甚麼.
整篇詩篇講完之後.
上帝求你鑒察我.
知道我的心思.
試煉我 知道我的意念.
然後就在.
看看我裡面有甚麼惡行.
引導我走永生的路.
你見到兩個經文.
有甚麼不同呢.
第一節經文說耶和.
根本就已經鑒察我.
和認識了我.
他根本你不求.
他都認識我 鑒察我.
為何詩人最後又再求.
神啊 你真的要鑒察我 認識我.
其實主要原因就是.
神就.
你求不求他都看著你.
不過你想不想.
知道上帝眼中的你.
今天的光景.
你想不想去引導你.
你想不想知道他怎樣看你.
他想.
詩人說我想你.
真的讓我知道.

$^{161}$我眼中 我靈命上你有甚麼喜悅的地方.
我還有甚麼缺陷的地方.
你讓我知道.
然後你引導我.
不要讓我走歪路.
做錯事都不知道.
我想知道你怎樣看我.
其實這個詩篇就是想講這件事.
他講這件事的時候.
他就用了中間.
十八節經文.
每六節是一個單元.
如果你熟悉詩篇一百三十九篇.
就知道了.
頭六節 一至六節.
就講到耶和華的無所不知.
他說你前後環繞我.
我一句說話未講出口.
你已經知道我想講甚麼.
然後他說我坐下.
我起來你都曉得.
如果你記得頭六節經文.
就是說神.
當然他自己是無所不知.
但他知道我們每日的生活.
很留意看著我們每日的生活.
我不知道你喜不喜歡被神看著你的生活.
你知不知道留不留意.
神都在留意你.
當然如果你很愛主.
很熱心事奉 原來上帝很開心.
知道我這麼熱心.
但如果你說我的生活.
原來離開了教會.
就不是一隻羊 是一隻狼.
回到辦公室就變成另一個樣子.
沒有人知道我是基督徒.
你都幾害怕 原來神都看著我.
怎樣生活 我家和老婆吵架.
原來上帝知道.

$^{201}$看著我 怎樣罵.
他說他看著我們.
那你就想想.
神為何這麼留心我.
全世界這麼多人.
現在七十多億人.
古往今來這麼多人.
他怎樣留意著我 看著我.
我坐下 我起來 他看著我.
就是說.
我整天的生活.
大小事他都很留意你.
神為何這麼留意我們 因為神愛我們.
下面會告訴你為什麼.
其實這個就是我們靈修的基礎.
希望大家去找回最基本的東西.
然後七至十二節.
就是無所不在.
大家記不記得.
有些人點頭 但有些好像不太熟.
可能有些人未看過CV139.
七至十二節.
是講無所不在.
我無論去到.
去到哪裡都可以逃避.
躲避你的面.
我就算去到海極居住.
你都原來在那裡.
我去到天南地北.
原來上帝一早就在那裡等我.
我沒有地方可以逃避上帝.
其實逃避是一個跨式的方法.
根本沒有人想逃避上帝.
其實斯人很想和上帝親密.
但他反過來說.
我想逃避你 其實逃避不到.
我去到海極 天南地北.
沒有人去的地方 原來上帝都在那裡.
然後上帝就說.
其實你無論去到哪裡 我都會引導你.

$^{241}$我是無所不在的上帝.
和你同在的上帝.
接著六節就講了.
上帝不單止留意我們 上帝和我們同在.
正如新約聖經裡.
我們每一個重生的基督徒.
聖靈都住在我們裡.
耶穌基督.
復活的主 透過聖靈和我們同在.
不會離開我們.
然後接著就講到.
今天主要講的經文.
13節至18節.
這六節經文.
講完上帝的無所不知.
無所不在.
上帝很留意我 知道我的生活.
很關心我的細節.
和我同在.
為什麼他那麼關心我呢.
我是一個小蝦子.
全世界那麼多人.
為什麼要那麼看著我呢.
原來13節的開頭.
在原文那裡.
如果你看NIV聖經 或者NASB的聖經.
英文聖經就看到.
有兩節沒有翻譯出來.
因為整篇詩篇.
去解釋.
為什麼神那麼留意我.
知道我一切的事.
然後又和我同在.
因為他創造我.
而他創造我 是因為愛.
所以創造我的.
詩篇第八篇是講.
人算什麼 你儘管眷顧他.
然後將他管理萬物.
然後上帝做人.

$^{281}$原來上帝做人和其他.
所有動物 受造物都不同.
是為了人有上帝的形象.
而上帝就是愛.
愛就是形象裡.
一個主要的成分.
神想人.
分享到他在創造這個世界之前.
神自己三位一體之間.
那個愛的關係.
聖父 聖子 聖靈.
已經在一個愛的關係裡.
而他希望人.
受造物裡 其中一個受造物.
就是人類.
可以進入他的團體裡.
經歷到他的愛.
而代表上帝管理萬物.
這個是神創造人.
但可惜人是犯罪.
他講的時候.
神為什麼.
創造 在這裡講.
留心.
為什麼上帝關心我們.
為什麼上帝和我們同在.
因為我們的肺腑是上帝所做的.
我們知道我們的內臟.
是上帝所做的.
他跟著回憶在媽媽的肚子裡.
在母腹當中.
神已經偏織我.
我們其實好像一件毛衣.
織毛衣織出來.
十個月.
我們現代醫學都知道.
用這個.
看到出來的BB.
兩個細胞結合.
然後變成一個人仔.

$^{321}$慢慢一個器官一個器官成長.
變成一個人仔.
奇妙.
真是.
我們在母腹的時候.
神好像織毛衣那樣.
然後他說.
我要稱謝你因為我受造奇妙可畏.
你的作為真是很奇妙.
這是他心中知道的.
BB仔的成長很奇妙.
兩個細胞.
這樣結合.
然後他自己就會懂得.
變成肺 心臟.
器官 血管.
全部都變成肺.
整個故事很奇妙.
我們自己現在仍然是一個人.
我們成長 出世之後.
每個人都會成長.
吃飯就會大.
成長之後 到我們今天都覺得很奇妙.
我不知道你覺得自己奇不奇妙.
其實大自然很奇妙.
我們都很奇妙.
我不知道大家吃了晚飯沒有.
吃了晚飯.
你不用記住要消化.
吃完之後自動會搞.
胃就會出胃液.
消化完之後膽汁就會出來.
自動運作.
不用記住.
如果記住就很煩.
如果你吃完飯之後聽到.
要送食物到小腸.
我忘記了.
他就不會動.
你不用記住.

$^{361}$很奇妙.
每天都在服務我們.
呼吸都很奇妙.
不用記住.
記住就大件事.
不能睡覺.
一睡著就窒息.
很有趣.
不用人意識去做.
但你最需要的呼吸自動運作.
上帝的創造很奇妙.
所以詩人說很奇妙.
我受造奇妙可謂.
我心心知道.
然後十五節.
我在暗中受造.
在地的深處被塑造.
那個時候我形體並不向你隱藏.
他在說什麼呢.
其實他在說.
他用一個很象徵性的文學表達.
其實我.
從一個細胞.
我們現在的說法就是兩個細胞結合.
成為一個生命體.
其實我初初這個生命.
由無到有出現.
是上帝做的.
當我一出現受造的時候.
全世界都沒有人知道我已經存在.
你知道你生命第一刻.
連你媽媽都不知道.
因為通常要過一段時間才查出來.
媽媽不知道你存在.
你自己又不知道自己存在.
因為你是人.
我們現在大概一歲半.
才開始意識自己的存在.
我們所知道.
你媽媽的肚子裡你不知道自己存在.

$^{401}$你又不知道自己存在.
你媽媽又不知道你存在.
全世界都沒有人知道你存在.
他用文學的方式表達.
你雖然受造了存在了.
但你好像在這個世界裡都不知道哪個角落存在.
其實他想說十五節.
全世界都沒有人知道你存在.
但上帝一早已經看著你.
上帝知道你存在.
看著你的形體在.
你這個BB細胞這麼小.
上帝都不會看不到你的.
當然不會看不到.
因為他造你出來的.
怎會看不到你.
所以他現在十五節就告訴你.
全世界都沒有人知道你存在.
你這麼微小.
可能存在第一個小時.
這麼小的細胞.
你媽媽都不知道你存在.
上帝已經知道你存在.
他跟你這麼熟悉.
他說上帝熟悉我們.
多過我們自己熟悉自己.
他想說這個奇妙.
然後我未成形的體質.
你的眼已經看見了.
明白嗎?這是我們.
一個BB還未完整.
發展中已經認識我們.
然後他跟著說了一樣東西很重要.
我未到過世上的一天.
即是未出世的那一天.
我的生命的年日.
原來上帝有一個計劃.
寫在冊子上.
所以我不知道你.
跟人傳福音的時候.

$^{441}$我猜你大概都會跟人說.
我們每個人上帝都有一個計劃.
那你信主.
到今天你找到你的計劃了嗎?.
你知不知道上帝.
你的生命有什麼計劃?.
我恐怕我們基督徒跟人說完之後.
我們自己都不太相信這件事.
或者不是很捉得實.
這件事上帝.
給我們有不同的人生經驗.
不同的知識.
出生在香港.
有這樣的家庭,教育,專業訓練.
一生他有什麼計劃呢?.
我不知道你有沒有問過這個問題.
我希望你聽完這篇道之後.
你開始去問這個問題.
上帝會告訴我們.
他很想告訴我們.
但問題是你想不想知道.
你記得這篇詩篇想說什麼嗎?.
上帝一向都有很多計劃在我們生命裡面.
你想知道你繼續問.
無論你現在是年青人.
是年長的.
不會太遲的.
你繼續問.
上帝有他的計劃在我們生命裡面.
17節.
上帝啊你對我的意念何等寶貴.
這個質上.
上帝對我們生命的計劃是很寶貴.
奇數何等眾多.
是量.
很多想法.
想我們一生過得豐盛.
經歷到它.
但你要問,尋求,想知道.
剛才我講了整篇詩篇.

$^{481}$就是詩人告訴我們想知道.
上帝.
對我生命裡面有很多計劃.
反過來.
用我自己做例子.
就是說.
其實上帝原來想我做一個傳道人.
當然現在我.
去到這個年紀就知道了.
但以前不知道.
因為我本來是做社工的.
但想做社工.
進大學的時候就經歷到神.
總之很奇妙的進了大學.
應該進不了.
第一次考了就不行.
第二次再考就自修.
進了大學.
那時候就.
開始成長.
就想做社工都挺好.
服侍神.
畢業的時候就.
全部是基督徒的機構去服侍.
就去了其中一間.
想不到.
第一年就經歷了很多事.
我做中學社工.
做學生輔導.
第一年.
神就讓我經歷到兩方面.
一就是個案很難做.
我們那時候還去四間中學.
我84年畢業.
所以早期剛剛開始.
就去考社工.
就在自殺,割脈.
黑社會開片.
我做東區四間中學.
你知道二十多歲.

$^{521}$剛剛讀完出來.
一堆都是理論.
你真的見到那些人怎麼處理.
談都談不來.
很大壓力.
有些案就出事了.
有個告訴我要去開片.
就勸他不要去.
他答應了我不去.
我真的相信了他.
就忘記了要通知上司.
因為在輔導守則.
他答應你的事不要全部相信.
一定要通知上司.
因為有生命危險就不可以全部相信.
怎知道我太傻了.
做了太多案件忘記了.
就沒有說.
怎知道第二個星期.
過了幾天之後.
就打電話來中心.
就說學生從那天開始沒有上學.
不見了去哪裡.
就大件事了.
他爸爸來找兒子.
校長就問我.
好像聽說有來見你.
就說是.
他有沒有談過什麼.
我說開片.
他說開片你不通知我.
你不通知家長.
大件事出事了.
那時候以為他死了.
或者出了大件事.
幾天都不見人.
就說為我試問.
有什麼事.
好在過幾天之後.
就全身都黃了.

$^{561}$內傷.
總之就躲起來不敢回家.
後來總之沒有死.
他最後爸爸就要帶他回大陸.
黑社會那些.
那裡令我很大壓力.
但那段時間有導師陪我一起祈禱.
就經歷到神.
怎樣托住我.
然後很奇妙的.
起碼解決了這件事.
另一個案件.
有人信了.
有個學生.
上學期的時候已經和他談過.
但解決不了問題.
很頑皮.
因為上課的時候搞事.
因為他的程度太低.
26個字都沒有學完.
當然是搗亂.
因為聽不明白.
老師又不會理他.
就投訴見他.
但這些案件是沒得做的.
因為有時候是老師針對他的問題.
但無論如何.
做不到都見見他.
但到下學期突然改了.
整個變好了.
我上學期見到他有些用.
可能開老筍.
現在開始明白我說的話.
照他們見.
原來他放學被人罵了.
在北角有一間教會.
那些牧者.
罵了他上去打乒乓球.
然後回到教會.
他爸爸也回來了.

$^{601}$這樣就變了.
你知道對我來說很大刺激.
原來這樣.
信耶穌和有神.
在生命裡有教會群體.
一個屬靈的家.
人就真的可以長遠地改變.
很徹底的.
見輔導室.
沒用的.
他不來見你.
你也沒辦法.
或者他亂說一堆.
其實也不是真的.
作個故事給你聽.
你沒得查.
其實你知道輔導室很被動.
於是乎成為了我一個很大的感覺.
做輔導也不是真的幫到人.
我說不要做社會工作人員.
我做靈性工作人員.
但靈性工作人員.
身心靈都要幫到人.
是什麼呢.
就是做傳道人.
但我很怕做傳道人.
因為那時候我所見到的傳道人.
穿白色衣服和藍色褲子.
在教會裡很正經的.
年輕人接受不了.
結果另一個故事拖了三年.
我只是說原來神要想我做傳道人.
三年後就見到了.
八百年就見到了.
又想不到另一個故事.
原來奉獻做傳道人.
都很擔心自己不適合做幕會.
因為我都比較內向.
教會裡要見那麼多人.
是很困難.

$^{641}$但神見呼召就靠著神的恩典.
畢業就讓我留在這裡教書.
是另一個故事.
但原來神在我的生命裡真的有計劃.
對我來說就是.
我每一次走一步.
我為了上帝去找一些.
只是基督徒的中心.
希望信仰同位.
做社工結合.
然後神原來要做傳道人.
願意做.
想做幕會.
原來可以教書.
更加適合我的性格.
教書然後可以讀到我自己喜歡的靈修.
然後在靈修裡除了教書.
還可以寫作.
還有做靈感運動的分析.
好像武俠小說.
人生一路走下去很精彩.
你要找.
你要問上帝.
然後你肯放下自己的喜好.
放下自己的計劃.
神的計劃有些可能和你的計劃差不多.
但很多可能不太相同.
你等他決定.
但你要問.
上帝求你監察.
我知道我心思刺濂.
我知道我的意念.
私人很想知道.
其實上帝一向已經有很多想法.
對於你.
對於我.
然後他說.
上帝啊.
看我的惡行會有.
其實最大的惡行就是.

$^{681}$我們自己不理會上帝.
過我們的生活.
就不只是貪心.
說謊.
不好脾氣那些.
那些是個別罪行.
整個人生的方向不是向上帝.
這個是最大的罪.
和最基本的罪.
其他罪的根源就是.
我整個生命.
沒有很認真.
問上帝怎樣容我.
而是我自己過著偏行紀錄.
偏行紀錄.
自己做自己的上帝.
做自己生命的主宰.
是一切罪的來源.
我希望大家去靈修.
就是第一件事.
先做一件事.
將生命交回給神.
信得過他.
帶領的人生是最好的.
靈修就是.
相信神會帶領.
在新約聖經裡面很清楚.
我們每一個重新得救的基督徒.
都有聖靈住在裡面.
大家知道羅馬書講了.
沒有聖靈的就不是真正的基督徒.
就不會得救.
所以你真心悔改認罪.
接受耶穌基督做你的救主.
和人生的主宰.
聖靈就住在裡面.
聖靈住在裡面只有一個目的.
就是每日引導我們.
透過真理去引導我們.
所以你要明白.

$^{721}$透過新約聖經.
聖靈每一日都要透過聖經.
和你講話.
要帶領你的生活.
我跟住就和大家講下去.
靈修是什麼呢.
就是每日去相信.
耶穌基督透過聖靈.
會給我們有很清楚的帶領.
這個就是靈修的目的.
所以我們靈修就是.
按住聖經這個教導.
每日去回應上帝.
上帝你會和我.
你會個人帶領我的.
所以我讀聖經.
我祈禱.
將每日教託給你.
所以總結剛才講的東西.
一個很基本的教義.
或者聖經的教導.
上帝和我們的關係.
不單止是集體的.
我們一齊.
整個群體神都會帶領.
所以我們每年可以有全年計劃.
一齊事奉.
整體一齊學習去傳福音.
一齊學習靈命更新.
譬如今年這樣.
集體性的.
但你亦都要個人性.
去跟從神.
即是說上帝.
不單止是想整個群體.
所謂集體的意思.
就是由一個個人組成.
即是說上帝帶領整體.
亦都帶領你個人的生活.
你個人的生活裡面.

$^{761}$你要記住.
剛才講了.
原來在聖經裡面.
我們每一個人.
都是神獨特的創造.
即是我出世的時候.
我為何有廖炳棠這個人.
存在這個世界呢.
因為上帝想起要創造一個人.
那個人就是廖炳棠.
他是男生.
生於香港.
然後就是生於某年某月.
我不告訴你.
要他出世.
然後他就會.
中四的時候.
有人介紹他.
就去了教會.
信耶穌.
這個就是上帝在我的生命裡面.
一步一步鋪排.
你明不明.
即是說上帝是.
因為想做一個廖炳棠出來.
所以廖炳棠就存在.
每個人都是他獨特的創造.
不可以有另一個生命代替.
而上帝愛我們每一個個人.
我們知道.
而聖經裡面講請教亦都是個人性.
九十九隻羊.
不見了一隻.
耶穌都怎樣.
放下牠.
九十九隻.
在羊圈裡.
然後他自己出去.
怎樣都要找回那隻.
我不知道你有沒有想過.

$^{801}$無論你怎樣信這個過程.
神是特意找你回來的.
神特意感動了弟兄姊妹.
要帶你回教會.
神特意感動了一個人.
可能要跟你傳福音.
神特意在那次聚會裡面.
報道會.
你聽得明那個訊息.
甚至你謙卑自己下來開放.
肯去接受這個福音.
神是特意要請教我們主動.
而神請教我們的目的是.
希望我們認清楚.
原來祂就是創造我們的那位主.
祂在我們生命裡面有計劃.
同時我們可以由初信主開始的關係.
真是經驗到祂跟我們互動.
而我們每日靈修就是跟神保持互動.
我希望一路講下去.
告訴你.
你怎樣做到這件事.
或經歷到這件事.
所以聖經裡面講創造是個人性.
請教是個人性.
聖靈內在是個人性.
聖靈內在整個群體.
但聖靈是住在你個人裡面.
所以去到最後.
剛才我一開始講的那件事.
上帝很關心你旁邊的弟兄姊妹.
有沒有祈禱.
有沒有靈修.
但他靈修.
上帝開心.
但如果你不靈修.
上帝也有幾天不開心.
因為旁邊那一代看不到你.
上帝想聽你.
懂不懂得去欣賞祂的創造.

$^{841}$祂的拯救.
你生活你願意和祂.
將生命交回給祂.
每日都問祂.
其實上帝你今天.
我應該怎樣在你面前用好今天.
怎樣依靠你.
你明不明我的意思.
所以個人性的帶領.
是我們整個信仰的基礎.
亦都是我們靈修的原因.
每日靈修.
上帝會帶領我們.
祂透過聖經帶領我們.
所以靈修基本上.
不只是那幾節經文.
所以很個人性.
耶穌基督每一個人都是擁抱我們.
一隻迷失的羊.
擁抱我們.
好像浪子一樣.
爸爸其實是走出去村口主動.
等浪子會不會回去.
所以我們每日就像浪子一樣.
睡醒覺的時候.
主人我仍然明白你想愛我.
你好像每日靈修.
上帝每日都問你一次.
我今天繼續願意做你的上帝.
你願不願意繼續跟從我.
我每日用靈修的時候.
就告訴你.
我願意.
我繼續.
你仍然是我的神.
你幫助我今天.
能夠懂得跟從你.
愛你.
所以這是我們靈修最主要的基礎.
所以靈修不只是那幾節經文.

$^{881}$你用什麼方法都可以.
來來去去都是靈修.
讀經祈禱.
唱詩.
其實方法是不重要的.
因為基本上聖經裡面說的都是那些元素.
什麼方法都可以.
但最重要是有沒有信心.
就是剛才說的東西.
你知不知道其實上帝是很想個人性.
聽我們.
怎樣將我們的生活交託給他.
我們心目中生活有沒有他的角色.
有沒有他在.
你問不問他.
所以其實每日我們去學習.
就是你讀那幾節經文.
其實一個象徵就是告訴你.
上帝你可以透過這幾節經文引導我.
也可以透過.
由現在這幾節經文開始.
我就將我整日生活交託給你.
你帶領我.
是什麼意思呢.
第一件事就是.
你讀那幾節經文裡面.
你真的要想聽到上帝的聲音.
正如你現在聽到都是.
第一件事你要運用信心.
上帝啊我不是隔壁那個聽到.
聽得好不好.
跟我沒有關係.
不會因為隔壁那個.
或者整班電視節目好聽到.
我可以分心.
游魂就等於崇拜上帝.
我可以沒有的.
你知道集體代替不了你.
隔壁那個代替不了你.
上帝是想知道你有沒有敬拜他.

$^{921}$所以譬如你崇拜.
跟大家說一樣東西很基本.
就是每一次回來崇拜.
就跟神說.
神啊我回來是敬拜你的.
無論一陣子的詩歌好不好聽.
不要緊.
因為你知道詩歌.
有些人喜歡某首歌.
某些人喜歡某類詩歌.
有些人不喜歡.
不重要的.
因為詩歌是唱給神聽的.
不是唱給我聽的.
所以神啊.
就算那首詩歌我不喜歡聽.
不要緊.
因為我現在不是來音樂會.
我來敬拜你.
你藉著那首詩歌.
你去想想你自己這個禮拜的生活.
有什麼要多謝神.
有什麼做得不好的.
有什麼地方是軟弱的.
求神力量的.
崇拜詩歌或者各樣的程序.
都是幫我們個人去敬拜神的.
所以第一樣東西就是.
你意不意得到神在這裡.
復活的基督現在在這裡.
說一個例子.
我爸爸信主的時候六十多歲.
信了一段時間之後.
我就考考他.
看他是不是真的信得很認真.
有一次早上和他吃早餐.
說爸爸這次不要我祈禱了.
每次都是我祈禱的.
你知道老人家信主.
這次是你祈禱.

$^{961}$很爽快的.
好啊 那我祈禱了.
一低低提起就要怎麼祈禱呢.
主耶穌說早晨.
多謝你賜給我們吃早餐.
奉耶穌的名 阿們.
我第一次聽到有人和耶穌叫早晨.
我從此記住一件事.
挺好的.
早晨見人就叫早晨.
見神就叫早晨.
他在這裡.
我反而覺得我爸爸心目中有神在這裡.
原來他知道耶穌在這裡.
和耶穌叫早晨.
所以大家回去崇拜.
早上崇拜就叫早晨.
晚上就叫晚安.
總之耶穌都說我回來敬拜你.
你心裡有沒有耶穌基督在這裡.
和有沒有和祂打招呼.
知不知道你回來做什麼.
你知道很多人可以參加整個聚會.
就是很被動的.
然後有詩歌就唱.
合心意就很開心.
不合心意那首詩歌就咕咕聲.
聽到也是.
那個人很多笑話.
或者很流利.
我又很開心.
他講得比較掉漆.
或者他講得比較什麼.
我又不開心了.
甚至開始心裡咕咕聲.
我們不是相信上帝使用.
不同的講源.
沒有一個講源是完美的.
神在教會裡面設立牧師和教師.
什麼性格都有.

$^{1001}$但是他講的是真理.
你就要聽.
而上帝可以超越任何一個講源講的東西.
如果你真的講得很悶.
你不如跟神說.
神啊 為什麼你呼著那些講源可以這麼悶.
但是我也想聽到你的說話.
你可不可以幫幫我呢.
誠心跟神求恩典.
也為那個講源祈禱.
可能他昨晚睡眠不足.
可能家裡有事.
不知道什麼事.
但是神給我謙卑的心.
我不是負責評定他是不是一個好的傳道人.
總之我反而想聽你的道野.
你可能聽到東西.
我們曾經有獻身營的時候.
有個弟兄姊妹說趕得太緊.
進來的時候差不多沒心機聽.
因為工作很多壓力那些.
誰知道最後的時候.
他整篇都聽不到什麼.
最後講完呼召的時候說.
你們誰還在今天裡面逃避神的呼召呢.
一句突然間他驚醒了.
他很感覺神是用這句話跟他說.
因為他真的在逃避神的呼召.
突然間醒了.
這個就成為他整個獻身營裡面.
經歷到神.
然後他就醒了.
身體很累.
但是靈魂就很想聽到神.
所以其實你有沒有祈禱.
我很想這篇道聽到神.
我接受靈修的時候.
神啊你真的在這裡.
帶著主耶穌一起去靈修.
其實最主要是信心的問題.

$^{1041}$是愛的問題.
你愛不愛主.
你想不想聽他的聲音.
你信不信得過.
他真的會藉著聖經跟你說話.
你聽到的時候.
是不是真的想聽到.
神啊我想聽到的東西.
我有什麼崇靈功課要學你.
你試一下心裡面真的這樣祈禱.
然後你就看一下.
那次的崇拜跟你以前的崇拜有沒有不同.
你那次的靈修.
跟另一次有沒有不同.
你明白我的意思嗎.
真的運用信心.
告訴神想聽.
不過你知道.
神未必說的東西.
我們是合心意的.
你那天早上可能有很多壓力.
上班.
早上要去跟老闆交代.
或者要presentation.
誰知道打開聖經.
你們要去.
望佈天下傳福音.
大件事.
上帝啊.
我明明很想你安慰我.
加力量給我.
要我同在.
誰知道還要我望佈天下許.
我哪有空想這麼多世界大事.
我今早的project都還沒搞定.
那你怎麼靈修.
靈修是一個很有人性化的過程.
上帝知道我們的軟弱.
就是說.
你坦誠跟神說.

$^{1081}$神啊.
我本來很想看聖經.
不過我想聽你怎樣跟我同在.
加力量.
安慰那些聖經.
但你還叫我望佈天下許.
我哪有空.
上帝你可不可以等我今早present完先.
我下午吃飯的時候.
我吃飯十分鐘.
到時再跟你談大使命.
我先交托我今天的壓力給你.
你幫幫我來救命.
如果你這樣的時候很真誠.
但上帝是有憐憫有恩典的神.
就是說你祈禱是兩個主體.
你講你自己想講的東西給神聽.
交托給他求恩典.
但又給上帝空間講他想講的東西.
你搞定完之後.
那個project下午最緊要不要走數.
最緊要真的找十分鐘出來.
上帝說我搞定了.
多謝你.
今早始終過了骨.
現在大使命.
我想一想.
你可以主動去想.
大使命.
我可能現在都敢跟人傳福.
但我不懂得關心宣教士.
想想自己現在在大使命裡實踐到哪一步.
每天多走一步.
上帝我多走一步.
我有奉獻給宣教士.
但我沒有真正為他們代禱.
我現在有奉獻又代禱.
你明白嗎.
每次走一步.
我們不用一次過走完.

$^{1121}$上帝說我可以做些什麼.
好一點的.
你明白嗎.
這個是一個很活的靈修.
是一個很人性化的靈修.
但每次向前進.
其實無論在基督教的誠信觀念.
他的看法都是一樣.
就是每天進步過昨天.
這個是最重要的.
就不用一次過走完.
因為很多信徒看聖經很害怕.
一看大使命.
我總有很多虧欠感.
我怎麼做都做得不好.
上帝不要我們一天就最好.
但你每天進步過昨天.
你上帝說他想說的話題.
你又可以將你自己的掛念跟他說.
這個互動就是靈修.
明白嗎.
在這個基礎之下.
想多說一件事.
在靈修裡面.
每日的靈修.
最重要的是幫助我們靈修之後的時間.
可以繼續用這個態度去跟從神.
所以不只是那節經文有沒有幫助.
或者只是那幾節經文.
上帝可以任何的聖經都可以提醒你.
而上帝可以活在你生活裡面.
藉著你生活的遭遇跟你說話.
聖經裡面說心裡面有道.
生活裡面任何的遭遇.
上帝都可以跟我們說話.
是甚麼意思呢.
剛才說了的就是.
我們的聖經可以怎樣應用.
大使命.
我想想我個人全方位可不可以.

$^{1161}$我家人上了主媒.
我的同事上了主媒.
或者我有沒有支持宣教士.
你可以逐格逐格想.
有時你想不到的.
上帝都會提醒你的.
突然間你腦裡面有個意念.
你應該可以奉獻一下.
這麼久一毫子都沒有奉獻過給宣教士.
人家這麼辛苦.
一家大小都去了外地都可以.
你自己出一點錢都這麼吝嗇.
可能上帝會有些事被動.
你都想到的.
這個就是在經文裡面.
想到怎樣應用.
另一邊就是.
當你每天都這樣去跟從神的時候.
最重要就是在那十五分鐘靈修.
或者二十分鐘靈修之後.
經常都要保持不住的禱告.
繼續運用信心.
你記得嗎.
這個是有上帝的世界.
我的生命是屬於上帝的.
所以我出門口生活和工作的時候.
我是帶著耶穌基督去上班和生活.
這個是靈修最重要的目的.
就不是那幾節經文.
是不是一定有很深刻的感動.
因為橫教會說得太多.
但是沒有什麼聖經根據.
因為有沒有感動.
是和你當時的情緒,壓力.
很多東西都有關係.
不可以做一個指標.
是不是一個好的靈修.
沒有感動的.
上帝都和你說了話.
大善明.

$^{1201}$你不用等到有感動才去想那幾節經文.
你明白嗎.
但是當我們心裡面.
裝載耶穌基督在.
祂會藉著聖經帶領我們.
接著你生活裡面.
什麼叫做遭遇呢.
就是當你這樣去依靠著.
帶著耶穌基督去上班的時候.
你就私下帶著祂去辦公室.
有弟兄姊妹的見證.
他說他工作裡面很大壓力.
老闆經常針對他.
所以很辛苦.
就站了很久肚子.
但是都沒有什麼解決.
但是突然有一天.
就有一個新同事來了.
發覺原來是基督徒.
又很談得來.
大家談得來.
因為大家都是基督徒.
都是比較認真的基督徒.
結果在他那裡得到很多鼓勵.
和一起祈禱.
甚至在公司開了一個銷售.
他突然間想起.
上帝幫我的方法.
原來不是將那件事拿走.
老闆就沒有走.
但是他現在有人支持.
就可以應付到那些壓力.
而且還可以在公司裡面傳福音.
上帝用了他的方法.
來幫他.
但很多時候.
在生活上.
如果你不是帶著耶穌在一起.
來了的那些.
又是基督徒.

$^{1241}$又不是基督徒.
那麼巧而已.
你明白我的意思嗎.
生活上很多事情都是.
那麼巧而已.
那麼發生.
星之加新.
那麼巧而已.
上帝在生命裡面做很多事.
你如果不是和耶穌基督在一起.
你不會覺得.
正如我剛才說.
你如果不是帶著耶穌基督在一起.
來崇拜.
你不和他打招呼.
其實上帝如果直接編導和你說了話.
你很多時候都覺得.
這次的講員很厲害.
所以他講到真的很有得著.
令到我這樣.
你不會覺得上帝直接和你說了話.
你明白我的意思嗎.
你去操練不住的禱告.
在生活裡面帶著耶穌基督在一起.
你就留意生活裡面很多事情.
上帝在介入.
我舉個例子.
其實剛才說的那些.
詩歌和聽道是我自己經歷.
我有興趣靈修.
就覺得自己靈性不是很好.
所以讀完之後.
開始留意了很多事情.
因為有些詩歌.
你知道我去不同的教會講道.
有些詩歌對我來說真的很難受.
其實我想大家都是.
我是傳統詩歌大.
我76年信主.
我現代的短詩我都喜歡.

$^{1281}$但如果打鼓打得太吵的教會.
我是很難受的.
因為那些鼓是影響我.
對年輕人很憤青.
但對我真的很嘈雜.
因為我不是這樣長大.
我不是聽那些詩歌.
更煩.
有一次那些詩歌還是那些.
流行曲的曲調.
改編了一點.
我聽完之後感覺很夜總會.
那個聯想是不能甩掉.
真的很騷擾.
我們進來這麼神聖的時刻.
弄了一首夜總會的流行曲的音調給我.
我真的很難受.
就鼓鼓聲.
我鼓鼓聲的時候.
就發覺心裡的聲音.
是上帝主動提醒我的.
這些詩歌是唱給我聽的.
不是唱給你聽.
我突然醒了.
是喔.
我說上帝你真是很厲害.
你心胸真的很闊.
這麼差的詩歌你都聽得到.
我不行了.
不過上帝你真的讓我繼續.
不要受騷擾.
勁快來.
總之不要受騷擾.
你愛那些人真是好.
那些人.
他們所唱的.
他們三文不曉得.
聽到也是.
你知道我們上了教會這麼久.
或者我們讀了神學之後.

$^{1321}$聽到還很容易.
你知道嗎.
學了這麼多色經的東西之後.
當然有很多偏偏都不滿意.
如果要挑剔起來的時候.
但是上帝就告訴你.
哪個人說的道.
不是你的學識的問題.
是上帝的僕人.
上帝要用他.
上帝在教會裡面設立牧師和教師.
而所有牧師和教師都是有限.
還有他的性格.
有他的表達.
你最重要就是祈禱.
上帝說完之後.
真的說成這樣.
上帝你真是.
為什麼這麼悶的說完.
你都會呼叫.
你可以抱怨一下.
但是結果就是.
那你又不奉獻.
你不是說話很厲害.
等你這麼久都不奉獻.
我就選擇那些說話差的.
可能是這樣.
回答你.
甚至有些就是.
我那時候有些領受天.
就是上帝其實是特意選擇一些.
很悶的人說道.
他想看你.
喜歡別人entertain你.
還是真的聽神的道.
如果你真是很有心聽神的道.
很悶的那些你都會認真聽.
來考考你.
特意有些悶一點.
不知道是不是.

$^{1361}$總之其實自己有些領受.
你真的帶著耶穌基督去生活.
你就留意到很多生活的東西.
去一個喪禮.
想一下.
上帝有沒有藉這個喪禮.
要和你說話.
你看新聞.
其實最好的是代禱時間.
是不是.
其實集體教會祈禱會.
除了集體回來祈禱之外.
一起看新聞時間就去祈禱.
是不是.
為巴勒斯坦人祈禱.
帶著耶穌一起去看.
你又不怕太不開心.
香港人看了這幾年新聞.
看到散了.
太多負面消息在香港.
因為我們靠自己去看新聞.
你明不明白我的意思.
我們不是和耶穌一起生活.
如果你真的受不了.
你和耶穌說.
哇 又是這樣.
這麼不公義的事情都會發生.
想打爛電視機.
和耶穌說.
神啊你幫助我.
我承擔不起這麼多.
這麼不開心的人生的事.
你明不明白.
然後我甚至為你代禱.
發生這麼多這些事情.
真的.
看新聞是一個最好的代禱時間.
帶著耶穌一起去看新聞.
帶著耶穌去生活.
工作就是我們靈修的目的.

$^{1401}$比一切的目的都重要.
所以保持著.
要保持著不絕的禱告.
去依靠著他.
耶穌你和我一起生活.
然後看看神在生活裡面.
怎樣介入.
好嗎.
我相信這個就是我們靈修.
每人想做的事.
所以今天和大家說的.
很簡單的幾個訊息.
第一就是.
真的靈修其實是一件.
不是很難操作的事.
是因為耶穌基督很愛我們.
祂和我們是有個人性的.
祂想聽你的祈禱.
不是隔離那個.
又不是只隔離你.
那些很有貢獻.
對基督教很有貢獻的人.
才會關心的.
其實好像.
靈修我經常覺得.
每次母親節和媽媽慶祝的時候.
其實最重要是你的心.
就是說.
其實你去見媽媽的時候.
你說是不是一個責任.
是一個責任.
每次母親節.
我們結婚那些回去找家人.
但如果你說只是一個責任.
就氣死你媽媽.
因為有愛.
那些責任是很開心去做的.
所以不在於裡面.
有沒有特別的感動.
或者什麼.

$^{1441}$每天去見一下.
媽媽定期回去探望她.
知道她在生命裡面.
很多的恩典給你.
然後去依靠她.
接著在生活上就是.
媽媽就很想聽.
不只是聽隔離.
大兒子很有成就.
她當然喜歡.
但如果二女沒有成就.
她根本就不是關心你有沒有成就.
她只是關心你想不想.
繼續和她聊天.
你明白我的意思嗎.
很有成就.
你回去的時候.
人到心不到.
媽媽都不開心.
上帝都是看我們是不是有心.
你每天是不是認真.
想依靠我.
想我做你生命的主宰.
每天和你聊天.
當你生活上覺得.
不和你聊天.
都照樣過得到生活.
那些就很糟糕.
你整個生活都不需要有上帝.
我覺得很開心.
我們究竟是一個什麼樣的生活呢.
所以上帝只是想聽我們的祈禱.
他甚至不理我們有沒有成就.
根本不重要.
你知道他只想你.
生活中有沒有他.
第二.
靈修就是每天帶著耶穌基督.
去上班和生活.
所以每一朝的靈修.

$^{1481}$其實就是一個重新立約的意思.
約就是恩典的約.
就好像聖餐.
每一朝確認我們和神的關係.
祂是我們生命創造的主.
祂是拯救我的主.
付了很大的代價.
十字架.
而祂現在透過聖靈.
每天帶領我.
而將來回到天家.
我們會和祂.
除了一個集體的敬拜.
我們每個人都和祂有關係.
我相信有個人和祂聊天的時間.
你看復活基督是和彼得聊天.
登山變象的時候.
有以利亞和摩西.
是個人性的關係.
所以將來回到天家.
不是只是集體唱詩.
祂和我們有關係.
個人會聊天.
我們今天就學習這個生活.
所以在生活中的時候.
就是從經文中知道.
有信心祂會藉著任何.
很熟悉的經文.
都可以和祂說話.
因為我們的應用.
是永遠學不完的.
彼此相愛.
永遠學不完的.
王普天下全副丹.
學不完的.
每天問祂.
我今天學了什麼.
每天學一步.
然後帶著一個耶穌基督的同在.
閱讀生活中的遭遇.

$^{1521}$上帝會透過遭遇和我們說話.
這個是聖靈.
因為我們不一定是透過聖經.
才想到應用.
去生活中遭遇.
去想回聖經.
你明白我的意思嗎.
新聞的時候就提醒我們.
要愛倫謝.
而那些倫謝.
可能就在巴勒斯坦很慘.
現在被人.
那些人生活在水深火熱中.
我怎樣可以幫到他.
起碼為他代禱.
這個就是我們每日靈修的意思.
今天和大家分享到這裡.
我們一起祈禱.
希望有半分鐘的時間.
大家去自己向神祈禱.
有沒有藉著這篇道.
有某些提醒對你是重要的.
你去和祂祈禱.
然後我帶領大家一個結束.
謝謝.
真是天父我們同心來到你面前.
我們每個人心裡.
有聲無聲的祈禱.
你都聽到.
你很想聽我們的聲音.
因為你愛我們.
你創造我們.
從無到有.
這個生命就是你自己所珍惜的.
你拯救我們.
你特意去尋找我們.
以致我們懂得回轉.
今天我們和你恢復這個關係的時候.
只不過是一開始.
你很想我們在生活上.

$^{1561}$真的經歷到你自己.
怎樣掌管我們的生命.
引導我們的生命.
我們每天希望透過靈修的時間.
去親近你.
去宣認你仍然是我們生命的主宰.
是唯一的主宰.
我們知道你同在.
你掌管著我們生命.
也都掌管著這個世界.
我們帶著你去生活和工作.
讓我們在你的話語上有領受.
知道怎樣每天在靈性上有應用.
也都讓我們在生活的經歷.
遭遇上聽到你的聲音.
你藉著周圍的事物.
去提醒我們.
怎樣可以在各樣的事上.
去見證你自己的名.
去榮耀你自己.
也都提醒我們靈明和事奉.
我們將自己交託給主.
我們每一個弟兄姊妹心裡面的心聲.
你都聽到.
求你按我們個別的情況.
繼續牧養和帶領我們.
我們多謝主.
奉耶穌基督的名而祈求.
阿們.
願神賜福給大家.
\newpage



\section{創世記 48:15-16}
\label{sec:RZv3JKPRqaw}
\textbf{2024-03-16 晚堂《一生牧養我的神》創世記 48\_15-16 講員:盧思源姑娘}
\newline
\newline
連結: \href{https://youtube.com/watch?v=RZv3JKPRqaw}{\texttt{https://youtube.com/watch?v=RZv3JKPRqaw}} ~~~~ 語音日期: 2024-03-19
\newline
\newline
\hyperref[sec:pdca3AXE1WI]{\small{< < < PREV SERMON < < <}}
~
\hyperref[sec:index_chronic]{\small{[返順時目]}}
~
\hyperref[sec:index_scriptual]{\small{[返順卷目]}}
~
\hyperref[sec:82RgV2iRB78]{\small{> > > NEXT SERMON > > >}}
\newline
\newline
創世記 48:15-16
\newline
\begin{longtable}{cl}
\hline
\hline
章節 & 經文 (和合本修訂版)\\
\hline
48:15 & \begin{tabularx}{0.7\textwidth}{X} 他就為約瑟祝福說：「願我祖父亞伯拉罕和我父親以撒所事奉的神，就是一生牧養我直到今日的神， \end{tabularx} \\ \\ \relax
48:16 & \begin{tabularx}{0.7\textwidth}{X} 救贖我脫離一切患難的那位使者，賜福給這兩個孩子。願我的名，我祖父亞伯拉罕和我父親以撒的名藉著他們得以流傳。又願他們在全地上多多繁衍。」 \end{tabularx} \\ \\
[1ex]
\hline
\hline
\end{longtable}
$^{1}$大兄姐妹平安.
願意在我們開始之前.
我們都有一個祈禱.
猜在我們再次到你面前.
天父一世間.
我們都要去聽你自己的話語的時候.
求聖靈與我們同在.
讓我們進入你自己的真理裡面.
讓我們被主經的話語得著.
被大主經的話語激勵.
以致我們更加行在你心意當中.
孩子去分享的時候.
求聖靈去引領.
以上禱告.
奉主耶穌基督明治祈求.
阿們.
大兄姐妹.
一生一世.
大家覺得對你來說有什麼意思呢.
或者這個字一生一世的時候.
你在哪裡會聽到多.
可能一些情侶.
又或者我們在婚禮裡面.
多時都聽到.
就是當一對新人.
和對方去許下一個一生的承諾.
就是一生一世去愛對方.
意思就是說.
在你整個人生裡面.
無論遇到任何的事.
都會去愛護對方.
這是一個很長久的承諾.
今天想和大家分享經文.
當中記載.
亞國即將要走到生命的盡頭.
我們見到這個他是在說亞國.
他為約瑟的兩個兒子.
嗎哪西和爾法蓮.
去祝福.
就收納為自己的兒子.

$^{41}$去祝福他們.
在當中我們見到亞國.
他提到他自己和神的關係.
我都highlight了紅色那裡.
他提到.
這個就是一生目讓我直至今日的神.
和救贖我脫離一切患難的那位使者.
在這兩句說話裡面我們見到.
在亞國他整個人生裡面.
這是他對神的總結.
而在聖經裡面.
這亦是最早以目讓去描述神和人的關係.
今天很想和大家藉著亞國的一生.
我們看看這位神.
今日亦是目讓你和我這位神.
究竟在亞國一生裡面.
他怎樣去目讓贖祂的人.
在當中我們有什麼學習呢.
說回亞國.
我相信大家都不陌生.
因為在這段時間的港島.
我們都提及過亞國.
大家想起亞國的時候.
就會記得他騙過哥哥二嫂.
以一碗紅豆湯.
騙取了長子的名分.
之後就騙了他爸爸.
拿了長子的祝福.
之後去舅父那裡拉班打工.
最後被舅父騙.
十次被他改工假.
之後回到家鄉.
去哥哥那裡.
在瓦博渡口和使者摔跤.
改名以色列.
我相信大家去看這些的時候.
都不陌生.
都聽了很多.
而我們如果回顧亞國一生的背景.
那個成長.

$^{81}$其實大家想想.
他從小就是在一個信仰虔誠的家裡長大.
即是他爸爸.
剛才說的以撒.
他爺爺是亞伯拉罕.
即是都是信上帝的人.
即是亞國.
基本上他是一個信三代的人.
即是關於對上帝的爺爺.
他應該一路以來.
聽他爺爺 聽他爸爸.
都說很多.
甚至都跟著他爺爺 阿爸去獻祭 祈禱.
如果我們看聖經.
在裡面我們見到.
在亞國一生.
上帝曾經七次向他顯現.
所以今晚很想我們看看.
究竟這七次顯現的時候.
上帝怎樣帶領亞國一生.
當我們去看聖經.
今日都會翻開幾多聖經.
我有經文.
如果大家有聖經可能去揭一下.
或者一路去看看.
我們亞國一生究竟上帝怎樣帶領.
上帝第一次向亞國顯現的時候.
是在創世記的28章.
在第十節的時候.
當時我們知道.
剛才說他騙了他爸爸以撒.
扮了他哥哥.
跟著拿了他長子的肅福.
所以以素就要去追殺他.
而他媽媽知道這件事.
知道這件事.
尼伯加就叫他快點走.
去舅父那裡 去哈蘭.
去舅父拉班那裡.
這是上帝第一次向亞國顯現.

$^{121}$這是他當時逃亡的路上.
聖經形容他當時.
孤身一人在半路上.
半夜上帝向他顯現.
去夢見神的天梯.
在經文28章13節.
他說 漢啊 耶和華站在梯子上說.
我是耶和華你祖父亞伯拉罕的神.
以撒的神 你現在躺我之地.
我要將他賜給你和你的後裔.
你的後裔必將地上的塵沙.
必向東西南北開展.
地上的蠻族必因你的後裔得福.
漢啊 我必與你同在.
無論你往那裡去.
我必保佑你 領你歸回這地.
我總不撇棄你.
直到我實現了對你所說的話.
我們看到當上帝第一次向亞國顯現的時候.
上帝告訴他.
他昔日和亞伯拉罕以撒的英雄.
他今日和亞國有這個英雄.
當時經文記載是怎樣呢.
如果大家看的時候.
亞國聽到的時候 他睡醒了.
耶和華真的在這裡.
當他第一次向上帝顯現的時候.
他的反應是 神真的在這裡.
他竟然不知道.
所以聖經說他立即拿起枕頭下面的石頭.
立作柱子 再敲油在上面.
將這地方叫伯特利.
如果我們再看下去.
亞國的反應是怎樣.
當第一次向上帝顯現.
他許願 如果神和我同在.
在我所行的路上保佑我.
讓我食物吃 衣服穿.
使我平平安安回到我父親的家.
我就必以耶和華為我的神.

$^{161}$我所立為柱子的這塊石頭.
必成為神的殿.
凡你所賜給我的.
我必十分之一獻給你.
我們大家看看亞國的許願.
他所求的其實很實際.
他求什麼呢.
他求我希望有衣服穿 有飯吃.
平平安安.
可能大家去想.
會不會和我們初初認識神的時候都很相似.
大家曾經有沒有做過類似的祈禱.
神啊 特別初信的時候.
神啊 如果你是真的.
你就保佑我平平安安.
工作順利 病就得醫.
如果是這樣.
我就真的相信你是真的.
大家有沒有試過.
可能初信的時候.
不多不少都會這樣去祈禱.
分享一點.
這段時間我爸爸.
他因為肺炎進了醫院.
其實之前都很嚴重.
醫生都叫我們有定心去準備.
但很感恩上帝的恩典和弟兄姊妹的祈禱.
他之前就在床邊的缺字洗禮.
現在的情況好一點.
不過還沒有康復.
最近很特別.
很特別是什麼呢.
他每天都打電話給我.
剛才也在打來.
每天都打來.
打來做什麼呢.
他每天打來.
你可不可以說我祈禱.
是每天 現在這兩個星期.
是每一天都打來.

$^{201}$想分前 祈禱啊.
我祈禱可以好睡一點.
他說原來每次他打來我和他祈禱的時候.
祈完禱真的好睡一點.
他覺得祈禱真的幫到他.
對一個初信的長者來說.
他覺得這個神真的幫到他好睡一點.
這個神是真的.
好靈啊.
就是這樣.
不過如果我們在座的人.
信著一段時間的時候.
我們的信仰是不是停留在.
啊 真是神真是聽啊.
好靈啊.
是不是停留在這個階段呢.
回到聖經.
我們看到雅各.
他雖然是信三代.
他爺爺是亞伯拉罕.
爸爸是以撒.
經歷了很多.
甚至乎他奶奶.
撒拉生他爸爸的時候.
大家知道是高齡產婦.
甚至當他爸爸.
其實曾經經歷過.
被爺爺差點成為一隻羊獻祭.
然後差點獻祭.
但臨下手的時候.
出現才獻祭一隻羊.
我相信在雅各的生命.
在他成長裡面.
他聽了很多這些.
關於上帝的經歷.
他知道上帝所作的.
以前所作的.
不過我們看到他許願的時候.
他怎樣呢.
他好像對上帝真實的同在.

$^{241}$好像有一點半信半疑.
當神跟他顯現的時候.
他對神開了一些條件.
他說如果我平平安安.
有得吃有得穿.
那就相信了.
如果上帝保證到我所要的東西.
他就以耶和華為我的神.
他就以上帝為他的神.
大家知不知道雅各說這番話的時候.
他幾歲.
在聖經裡面我們知道雅各.
他是147歲去世.
如果計算回頭.
雅各離開家裡.
他第一次顯現的時候.
他要離開家裡.
去拉班去舅父家裡.
那時候雅各已經是77歲.
他肯定不是一個小朋友.
在他77年裡面.
我相信在他的成長.
他聽了很多很多關於上帝的東西.
不過當上帝跟他顯現的時候.
我們看到他的信仰.
他所提的就好像停留在.
神啊你會不會保佑我.
你保佑我的話.
我就相信你了.
弟兄姊妹.
可能在座我們都有一部分.
我們可能是回教會一段時間.
甚至不知道有沒有信二代.
信三代.
或者我們都聽到很多身邊.
可能弟兄姊妹.
甚至家人的見證.
但有時我們的信仰.
是不是很容易就像雅各那樣停留在.
上帝有沒有實現我想要的東西呢.

$^{281}$我們所信的神.
是不是只有一個有求必應的神.
如果我們接著下去再看雅各的時候.
當上帝第一次向他顯現之後.
剛才說他去了拉班的家裡.
在當中我們讀聖經.
我們知道雅各是被騙.
明明就想娶個妹妹.
拉結.
誰知道他就娶了個姐姐尼亞.
為了這個妹妹被人騙了14年.
要服侍拉班.
之後又為他的牛群羊群.
又額外服侍多6年.
足足在20年裡面.
聖經裡面提到.
拉班不單止騙了他.
還要10次改他的工資.
即是10次騙他的價.
但在當中我們看到上帝的工作很奇妙.
上帝引領雅各就是.
剛才說昔日雅各.
他是說謊.
他欺騙了爸爸哥哥的時候.
騙了個長子名分.
我們有時看聖經會覺得.
為何上帝說謊.
都好像容許雅各.
不過我們看清楚一點.
其實上帝沒有認同他說謊的行為.
相反.
上帝藉著拉班.
給雅各經歷.
什麼是被騙的感受.
在當中我們看到上帝的引領是.
在神的愛裡面.
他都有管教.
上帝不是溺愛.
不是當我們做錯的時候.
他都說沒問題.

$^{321}$是在愛裡面有管教.
引導我們走對路.
但上帝很特別.
他不是說拉班騙你.
任由他騙你.
你騙人就被人騙.
那就算了.
剛才說有7次向雅各顯現.
而第二第三次.
其實就是雅各在拉班裡面.
住了20年之後.
剛才說因為雅各受不了.
再被拉班這樣虎待.
他很想走.
在創世記31章的時候.
耶和華就和雅各說.
你要回到祖先知地.
去你的本族那裡.
我必與你同在.
上帝說叫他回到他先祖的地方.
並且允許他與他同在.
而第三次就是.
上帝又繼續和雅各說.
這次和他說什麼.
除了叫他回到本族之外.
他告訴他.
拉班向你所做的事.
我知道.
他看到.
所以他叫雅各你回去.
你回去本族.
在當中我們看到什麼.
神真的一步一步帶領著雅各.
而且神沒有離開他.
他知道雅各所面對的一切.
他知道他受的委屈.
而且當我們看聖經.
其實神不單單和雅各顯現.
他也和拉班說話.
在31章24節夜間.

$^{361}$神來到亞蘭人拉班那裡.
在夢中和他說.
你要小心.
不要對雅各說好說歹.
如果用我們這些說話.
要凶他.
你要小心.
你見到雅各不要亂說話.
原來雅各想走的時候.
拉班本身是想追求他.
不過因為上帝和他說.
你見到雅各的時候.
你不要和他亂說話.
我們在當中不單止.
看到神帶領著雅各.
也在雅各不知道的情況下.
後期拉班當然和他說.
你的神和我說.
叫我不要和你亂說話.
但當時雅各不知道的時候.
上帝為他開了一路.
記得在我未讀神學之前.
或者預備讀神學.
接受了上帝的召喚的時候.
已經接受了上帝的召喚.
預備讀神學.
其實我之前的職業是.
物業管理.
我的職位當時要三個月通知離職.
而神學院的入學日期是九月中.
大家計算一下.
九月中入學.
三個月通知離職.
最遲就是六月中通知遞信辭職信.
不過當時的公司.
六月尾就會出一份bonus.
出了一份bonus.
正常你都遞了信.
你怎會有份bonus呢.
不過這份bonus.

$^{401}$坦白說也不少.
給你又如何呢.
正常你會想.
你計算一下這些價.
計算到六月尾.
然後你扣一下價.
看看可不可以得.
當時我沒有去祈禱.
究竟怎樣呢.
人是會這樣想.
但當我祈禱的時候.
很覺得我真的最遲最遲.
六月中一定要遞信.
因為我不可以延遲.
我神學院入學的日期.
就一定是六月中.
你說價有沒有扣呢.
我沒有問HR.
也沒有特別問究竟應該怎樣.
所以祈完禱之後.
就直接和我當時的上司.
我要辭職了.
因為我要去讀神學.
當我上司收到這封信的時候.
他很正常.
他就說.
你知不知道六月尾.
已經出bonus了.
你整年都是bonus.
你為何不再等兩個星期.
等了兩個星期.
你就先出吧.
我就告訴他原因.
我說不行.
要時間等等.
也分享了自己的信仰.
他不是基督徒.
不過他妹妹是.
一家甚至去了.
其他地方宣教.

$^{441}$過程中可以和他.
分享一些信仰.
分享完後.
他和我說.
其實上星期.
我曾經問過HR.
如果有人出bonus之前.
遞了信怎麼辦.
我不知道他為什麼這樣問.
可能其實是他想辭職.
所以他自己先問.
我肯定是這樣想.
我問了.
HR回覆他.
如果上司收到下屬的辭職信.
你不想和他走的話.
有一段時間.
可以當作negotiation.
即是你和他談.
談到的時候.
當他直接遞信給HR的時候.
Effective day仍然是.
回到6月中.
所以上司就說.
我幫你收了信.
我等到6月尾.
你的bonus出了之後.
才去給HR.
Bonus是不會收回的.
因為他問清楚了.
我就覺得.
你已經知道了整件事.
他就這樣.
如是者.
我真的給他信6月中.
等到我出了Bonus之後.
他就真的幫我遞了一封使信.
而我就很順理成章.
9月中就入學.
去讀神學.

$^{481}$這份Bonus.
也幫補了一些生活費.
因為大家知道讀神學.
其實我們是沒有收入.
在當中.
我看到什麼.
看到上帝是一步一步.
這樣帶領著我.
也都原來在我沒有想過的情況下.
上帝是為我去開路.
弟兄姊妹.
不知道今天.
可能你正在遇到什麼處境.
即是說.
他可能都會覺得.
看不到上帝很多的工作.
在自己身上.
但我們仍然.
我們仍然相信.
神必定一步一步.
這樣帶領著你.
而且可能在你不知道的情況下.
已經為你去開路.
如果我們回到聖經.
剛才提到.
雅各離開了拉班.
他就回到他的本族家鄉.
這個就是第四次.
上帝向雅各顯現.
就是當時在途中.
我想這個大家都很熟悉的.
上兩個星期的經文.
講道也有講過.
就是雅各一個人.
去到亞博道口.
很忐忑.
因為第二天就要見哥哥耳掃.
很害怕.
糟了.
哥哥會不會還在生他的氣.

$^{521}$如何面對他以前的事.
又很擔心.
還會不會害他的家人.
聖經裡面講了很多.
而聖經記載.
在這裡的時候.
神的使者向他顯現.
雅各就和神摔跤.
雅各的大腿被神摸了一下.
就變成了他的腿.
在當中我們也見到.
雅各捉住上帝要祝福.
捉住上帝要祝福.
而神在當中給了他一個新的名字.
就是以色列.
我們看到當雅各他最需要安慰.
最需要鼓勵.
內心覺得很懼怕的時候.
神的使者就在這一刻主動向他顯現.
並且賜一個新的名字給他.
由昔日叫雅各.
意思是捉住.
到上帝給一個新名字.
叫以色列.
意思是上帝和他爭戰.
讓雅各知道.
上帝一直和他同在.
神會保護他.
他不需要像以往那樣.
依靠自己去捉住很多東西.
而是相反.
他可以依靠神.
如果我們繼續看聖經創世記33章的時候.
我們看到當雅各重遇他哥哥以素.
我們知道以素原諒了他.
也見到兩兄弟和好.
以素沒有因為雅各之前所做而生他的氣.
而雅各也表示一個歉意.
之後兩個人就分道揚鑣.
以素去了西爾.

$^{561}$雅各就回到迦南地.
疏各就支搭槍柏.
其實我們見到上帝真的應運了剛才28章所說的.
雅各許的願.
上帝真的平平安安帶著雅各去迦南地.
弟兄姊妹.
不知道大家有沒有試過人生中.
經歷一些很懼怕面對的事情.
我想每個人都會有一些.
可能是一些人.
可能是一些事.
又或者像雅各那樣.
很怕面對自己曾經做過的事情.
可能覺得自己曾經做的事情不好.
之前大家不知道有沒有時間去看看網上的一些見證.
早前剛好看到一個星火飛騰的見證.
不知道大家有興趣可以上網看看.
都有很多很好的見證.
其中一集就說這個女士.
她13歲就已經絕學了.
沒有讀書.
當時就去做一個邊緣少女.
感情生活就很暖濃.
17歲就未婚懷孕.
之後還怎樣呢.
又去吸毒又去販毒.
基本上進入警局很多次.
後來就轉介到福音戒毒.
慢慢在當中她都經歷了很多事情.
但慢慢她感受到上帝對她那份的愛.
讓她可以重新做人.
而且更加一次機會裡面.
她重遇當時已經分隔了很多年的兒子.
令她覺得一切可以重新開始.
其中一段很深刻.
因為她在當中哭著分享.
她說雖然過去很多很多的黑歷史.
做了很多很多不好的事.
但發現神一直都沒有放下她.
她感受到神一直都在她身邊和她同在.

$^{601}$沒錯.
我們今天有些事都很怕去面對.
剛才說可能人可能是.
甚至乎怕面對我們一些過去.
我們那些亦都不想面對.
不過上帝允許我們所信的神.
他會陪我們經過.
他允許我們可以依靠他的時候去跨過.
接著如果我們繼續看聖經的時候.
大家記不記得已經顯現了第幾次.
現在是第五次.
即是即將來到第五次.
上帝第五次向雅各顯現的時候.
就去到35章.
35章的時候神對雅各說.
起來上伯特利.
住在那裡 在那裡捉一座壇.
給神 就是你逃避你哥哥耳掃的時候.
向你顯現的神.
為什麼上帝會叫雅各上伯特利呢.
這個就要剛才說.
我們看聖經的時候.
34章的時候其實發生了一些事.
中間提到雅各和利亞所生的女兒底娜.
就被當地的族長的兒子事劍污辱了.
所以底娜的哥哥西勉和利美就覺得很不甘.
為什麼我的妹妹被人污辱.
所以設計就將事劍城所有的藍丁都殺了.
也搶了他們的財物.
當雅各知道之後就很害怕.
害怕遭到報復.
就在這個時間.
上帝再次向雅各顯現.
再一次帶領雅各去經過.
當我們繼續看的時候.
35章五節那裡記載.
當他們起行離開事劍城的時候.
神使周圍城鎮的人都驚恐.
就不追趕雅各的兒子們了.
在當中我們看到上帝對雅各的應許.

$^{641}$就是他說他一直會保護他的一家.
就正如之前昔日他離開拉班的時候.
上帝向拉班發夢跟他說.
你不要跟雅各說好說大.
今天當雅各要離開的時候.
神使到周圍城鎮的人驚恐.
神再一次救雅各脫離危險.
之後我們再去看上帝對雅各顯現的時候.
就已經去到伯特利的時間.
當去到伯特利的時候.
在35章那裡.
我們看到伯特利是神第一次向雅各顯現的地方.
上帝在當中再重申一次對雅各的應許.
就是他有一個新的名字.
也是昔日他對亞伯拉罕耳撒的約.
就是要新養眾多.
將來有一國許多國從他而來.
而且應許他們可以得著上帝應許的地.
神再一次去同一個地方.
再一次的去顯現.
用他的說話去兼顧雅各.
我們很快去看最後一次的時候.
究竟上帝之後怎樣再繼續跟他說呢.
到第七次上帝向雅各顯現的時候.
其實隔了很多 剛才是35章.
到最後一次顯現的時候已經是46章.
當中距離上一次的超過20年.
如果我們看創世記的時候.
我們知道雅各在中間這段時間發生了很多很多的事.
其中一樣我們大家都很熟悉的.
就是他所愛的兒子耀失.
就是因為他偏心他.
因為一件彩衣令到耀失的哥哥妒忌耀失.
所以將他賣到埃及.
而我們知道耀失雖然被賣到埃及.
上帝都保守他.
成為埃及的宰相.
而當埃及出現大饑荒的時候.
只有埃及有糧食.
所以那個情節就是當雅各叫兒子去埃及拿糧的時候.

$^{681}$他們兄弟就相遇了.
而最後一次的顯現是什麼呢.
就是當時雅各決定離開這個迦南地.
要搬家 由迦南去一個沒有去過的地方.
完全陌生的埃及.
要和耀失見面.
不過大家在這裡停一停去想.
其實雅各在過去20年裡面.
其實他以為自己的兒子已經死了.
現在我們知道劇透了.
因為我們順著看.
他都還沒死.
他就在那裡.
但你要明白雅各不是.
他那20年以為自己的兒子已經死了.
以為已經沒機會再見他的兒子.
所以當他說原來他的兒子還在生.
還做埃及的宰相.
我相信雅各當時的心情其實很複雜.
一方面他當然是很想見到自己的兒子.
另一方面他要離開自己的地方.
搬去埃及.
對他來說真的一點都不容易.
對他來說有很多未知之數.
就在這個時間.
上帝再一次的去雅各顯現.
46章2至這裡說.
我是神 你父親的神.
不要害怕 下埃及去.
因為我必使你在那裡成為大國.
我要和你下埃及去.
也必定帶你上來.
若失要親自的.
及上你的眼睛.
在當中我們見到.
神再應許.
他一定會和雅各同在.
他不受地域的限制.
就算在埃及也好.
神也會繼續帶領他.

$^{721}$而且甚至上帝應許.
他最愛的兒子若失.
是可以和他去送中.
弟兄姊妹.
我們很快的看了.
雅各一生上帝七次的.
向他顯現.
我們見到每一次上帝顯現.
每一次上帝和他說話的時候.
都是他在人生遭遇一些難處.
或一些變動.
當他心情很不安.
很無助的時候.
上帝就是主動的向他顯現.
而且每一次顯現的時候.
上帝都是不厭其煩地和他說.
他是一位怎樣的神.
我們剛才看到.
我是你祖亞伯拉罕的神.
是以撒的神.
我是帶領你逃避你哥哥.
以素向你顯現的位.
我是你全能的神.
我是你父親的神.
上帝不斷的重申.
他是一位怎樣的神.
而且每一次上帝都會提到.
他給雅各的應許.
應許是什麼呢.
就是上帝必定會與他同在.
他的目的是什麼呢.
上帝的目的是要幫助雅各.
去認識他.
讓雅各可以更堅定.
可以信靠上帝給他的應許.
讓雅各不需要擔心.
不需要害怕.
因為上帝與他同在.
所以當我們回顧雅各的一生的時候.
由一開始雅各稱神為亞伯拉罕的神.

$^{761}$以撒的神.
但當他年老.
快要離開臨終的時候.
他已經親身經歷了.
在他人生裡面.
他經歷了很多上帝的恩典.
他知道對神的認識.
所以他將這份信仰.
確確實實化為他自己.
所以當他剛才說.
他為約瑟的兩個兒子祝福的時候.
他稱上帝為什麼呢.
一生牧養我直至今日的神.
以及救贖我脫離一切患難的那使者.
因為在他的人生裡面.
他經歷了上帝怎樣去牧養他.
經歷了每一次上帝怎樣帶領他.
經過一些患難.
他與神的關係.
不再是他父親的神.
而是他經歷了一生牧養他的神.
弟兄姊妹.
昔日牧養雅各的神.
今日也是牧養你和我的神.
雖然上帝不會像雅各那樣.
今天顯現給我們看.
神啊 你怎樣帶領我.
不過神藉著聖經.
藉著上帝的話.
讓我們知道他的應許.
這個昔日牧養雅各的神.
今日仍然不變.
仍然是同樣地牧養我們.
同樣地帶領我們經過.
他應許他與我們同在.
他應許他不會撇下我們.
他很愛我們.
就算我們還沒有認識他之前.
他已經認識我們.
而且他應許當我們遇到任何事情的時候.

$^{801}$我們可以依靠他.
但願當我們去看雅各的一生的時候.
我們都去追求和神這份.
自己和神個人和獨特的認識.
親身去經歷.
這個屬於自己.
你自己的信仰.
當我們每一次去讀聖經的時候.
我們讀到亞伯拉罕的神.
以撒的神.
雅各的神的時候.
我們都可以說出.
這個也是今日牧養我們一生的神.
但願上帝的說話成為我們彼此的面禮.
我們一同的禱告.
透過昔日你帶領著雅各的一生.
你讓我們知道你是那位一生一世牧養我們的神.
你是那位帶領我們經過所有困難危難的神.
主我們知道我們遇到任何事情的時候.
我們都可以依靠你.
求主幫助我們.
幫助我們在我們的信仰裡面.
我們真的經歷到對你的認識.
我們真的在捉緊你給我們的應許.
而且我們的生活.
在我們的行事.
在我們的為人裡面.
我們更深的去認識你.
更深的去信靠你.
更深的去愛你.
求主幫助我們.
幫助我們在你的愛裡面繼續去學習.
繼續去成長.
奉主耶穌基督名字祈求.
\newpage



\section{創世記 47:27-31}
\label{sec:82RgV2iRB78}
\textbf{2024-03-17 早堂《雅各的敬拜人生》創世記 47\_27-31 講員:曾裔貴牧師}
\newline
\newline
連結: \href{https://youtube.com/watch?v=82RgV2iRB78}{\texttt{https://youtube.com/watch?v=82RgV2iRB78}} ~~~~ 語音日期: 2024-03-21
\newline
\newline
\hyperref[sec:RZv3JKPRqaw]{\small{< < < PREV SERMON < < <}}
~
\hyperref[sec:index_chronic]{\small{[返順時目]}}
~
\hyperref[sec:index_scriptual]{\small{[返順卷目]}}
~
\hyperref[sec:ih_tESQH0P8]{\small{> > > NEXT SERMON > > >}}
\newline
\newline
創世記 47:27-31
\newline
\begin{longtable}{cl}
\hline
\hline
章節 & 經文 (和合本修訂版)\\
\hline
47:27 & \begin{tabularx}{0.7\textwidth}{X} 以色列人住在埃及境內的歌珊地。他們在那裡得了產業，並且生養眾多。 \end{tabularx} \\ \\ \relax
47:28 & \begin{tabularx}{0.7\textwidth}{X} 雅各住在埃及地十七年。雅各一生的年日是一百四十七年。 \end{tabularx} \\ \\ \relax
47:29 & \begin{tabularx}{0.7\textwidth}{X} 以色列的死期快到了，就叫了他兒子約瑟來，對他說：「我若在你眼前蒙恩，把你的手放在我大腿底下，以慈愛和誠實向我承諾，必不將我葬在埃及。 \end{tabularx} \\ \\ \relax
47:30 & \begin{tabularx}{0.7\textwidth}{X} 我與我祖先同睡的時候，你要將我帶出埃及，把我葬在他們所葬的地方。」約瑟說：「我必遵照你的吩咐去做。」 \end{tabularx} \\ \\ \relax
47:31 & \begin{tabularx}{0.7\textwidth}{X} 雅各說：「你向我起誓吧！」約瑟就向他起了誓。於是以色列在床頭敬拜。 \end{tabularx} \\ \\
[1ex]
\hline
\hline
\end{longtable}
$^{1}$弟兄姊妹 主女便平安.
今天能夠回來舞堂講道.
待會也會在嘉欣堂講道.
一起祈禱吧.
多謝天父的恩典.
我們能夠在禮拜堂敬拜你 朝見你.
我們大家都有自己的人生.
今早很盼望能夠藉著.
聖經創世記記載的雅各.
祂的人生如何在祂的生活裡面.
祂的遭遇 祂的人生際遇.
來去敬拜神.
剛才的經文.
祂是在臨終扶著像頭.
來去敬拜這位真神.
一生牧養祂的真神.
幫助我們明白.
什麼叫做一生的牧養.
幫助我們能夠明白.
如何能夠在我們的生活.
尤其是困難時候的生活.
來去經歷牧養.
我們大家都有.
我們自己生活的困難.
但是心信上主.
是在我們每個弟兄姊妹的心靈裡面.
有計劃 而且是美好的計劃.
我們就等候仰望.
藉著經文給我們的啟發.
使我們能夠有一個.
畜類旁通.
來去在我們自己的場景.
去應用 這樣禱告.
奉耶穌基督的名 阿們.
教會在說聖經人物.
就編到我要說雅各.
稍後五堂.
詩源姑娘都會說雅各的人生.
不過就不同的焦點.
在座早堂.

$^{41}$都應該有足夠對人生的體會.
如果年青的 暫時聽一下.
你的人生還沒有這麼多困難.
可能今早的講道.
你要慢慢再思考一下.
但是在座都老人家了.
這篇道 包括我自己.
都已經兩元了.
那就應該非常之合聽.
人生充滿了困難.
剛才我們提到.
有一個很美麗的圖畫.
雅各扶著拐杖敬拜神.
其實他是跛的.
被天使弄跛了.
創世記的三十二章.
他在這個叫比魯伊勒的地方.
來遇見神的使者.
整晚的摔跤.
摔跤完之後 天使要走了.
我不准你走 你要賜福給我.
天使就一摸他的大腿.
大腿窩就蹺腿了.
從此就不良於行.
那時候沒有輪椅了.
有拐杖的.
他的人生就是帶著拐杖.
這樣過他的人生.
三十二章去到五十章.
有很多章經文介紹這個雅各.
還有說很多很多的年期.
他幾歲 發生什麼事.
說了很多.
所以整個創世記五十章.
雅各是記載最多的男主角.
由二十五章開始出世.
和哥哥一起出世 是媽仔.
一路去到五十章都說他.
說他很多很多的遭遇.
當然今早我們只能夠收窄了範圍.

$^{81}$他的後半段的人生.
後半段才是最精彩的人生.
由三十二章.
神就說我為你祝福.
你的名字不要再叫雅各了.
雅各就是抓住.
我現在幫你改名天使.
三十二章摸到他攝取後.
幫他改一個名字.
叫做以色列.
所以以色列人就是由這個雅各而出.
叫做以色列.
希伯來文就是神的兒子的意思.
天使 王子.
這樣的意思.
不過他開始要回歸.
回到迦南地.
輾轉要面對自己的哥哥.
這個耳掃.
他很害怕.
後來輾轉終於要面對哥哥.
這樣來相認.
因為二十八章的時候.
他搶了 用非法手段.
搶了哥哥的兒子的名分.
媽媽怕他遭遇不測.
叫他回外家.
回巴旦亞蘭 寇府.
即是他的媽媽利伯加的地方.
就是巴旦亞蘭寇府拉班那裡.
輾轉過了二十年.
雅各這個叫做以色列.
從此由三十二章開始.
大部分的經文不再叫雅各.
只剩下很少經文會再叫雅各.
由三十二章開始.
聖經記載.
作者叫做雅各 叫以色列.
你可以回去再看一看那些經文.
不過剛才我都講了.

$^{121}$他的人生是一個.
由比魯伊勒開始.
他就要去一個敬拜的人生.
但敬拜的人生.
是否就是一帆風順的暢順.
從此我是神的王子 神的兒子.
我是否就是一帆風順.
很多很多的祝福無咒詛呢.
不過我們揭聖經.
有很多很多他的生活際遇.
有些是關他事.
更多是不關他事.
但是他要一一的.
作為家主 他要承受.
你想想這些傷痛.
如果你每一樣都和他一樣.
你的人生是非常的痛.
所以原來痛苦的人生.
如果有神.
如果你的敬拜.
是和這位神同步來接軌.
多痛你的人生就精彩了.
但沒有神同在同步.
好像WhatsApp突然間沒有WiFi.
同步不到.
有了WiFi之後.
那些WhatsApp有沒有經歷過.
就是這樣.
我們的人生如果沒有神和你同步.
你和我面對的痛苦.
你一樣要面對.
但是神的力量 神的同在.
就和你沒有關.
因為你沒有同步.
似乎那個WhatsApp還沒有到.
其實是有的 給了你.
不過你沒有WiFi 收不到.
當你一有WiFi.
他就來了 噓噓聲的來.
特別我們做牧者就是.

$^{161}$一去完這個地方 出來.
噓噓聲的那些WiFi就來了.
那些WhatsApp的群組.
因為太多人找我們.
我們看看35章第16節.
他們從伯特利起行.
來以法他還有一段路程.
拉傑臨產 甚是艱難.
在曠野野外.
這一位他最愛的拉傑.
為他來服侍舅父拉班二十年.
現在要回歸了.
回到迦南地 神應許之地.
拉傑臨產第17節.
正在艱難的時候.
修生婆對他說.
不要怕 你又要得一個兒子.
因為他的第一個兒子就是弱失.
第18節 他將近於死.
真是難產了.
古代根本不會趕上救護車.
立即打三條狗.
立即進去趕緊的.
來到去有最好的醫療.
是沒有這回事.
生產就是一個危險.
你在曠野生產.
很快就去到伯利衡.
但是都去不到.
在這裡的時候.
他說靈魂要走的時候.
就是死了.
就幫他的兒子改名叫做.
變俄尼.
希伯來文就是兒子的意思.
變俄尼.
俄尼就是艱難痛苦的意思.
這個兒子是令我艱難痛苦.
而且是死了.
真是生這個兒子難產而死.

$^{201}$雅各就幫這個兒子改名.
就不要叫變俄尼.
叫變雅文.
變雅文就是右手之子的意思.
拉傑死了.
葬在爾法他的路旁.
爾法他就是伯利衡.
即是很快就進城了.
都等不到.
就難產死了.
是自己四個老婆最愛的太太拉傑.
神的王子受打擊了.
沒多久.
第22節.
以色列住在那地的時候.
劉弁去與他父親的妾皮拉同寢.
以色列也聽見了.
自己的大兒子劉弁.
12個兒子最大的兒子.
同一個妾侍.
即是比較年輕.
你不會娶一個很大.
自己的元配的女士做老婆.
通常男士都是找一些年輕的.
妾侍.
小三.
這個妾侍和劉弁.
可能青梅竹馬.
相同年歲都不一定.
劉弁這個大兒子.
去與妾侍同床.
聖經很有趣.
以色列.
即是雅國聽見.
家丑聽見.
接著就是爸爸死了.
35章.
就是以薩180歲.
即是當時雅國已經120歲.
因為是60歲生這對雙胞胎.

$^{241}$接著告訴我們.
去到37章.
很詳細記載雅國最愛的太太拉潔.
最愛的太太的兒子.
因為辯雅文很小.
這個約瑟是爸爸最愛的.
一來就是太太都死了.
所以約瑟就出現在聖經的舞台.
不過好景不常.
約瑟被自己的哥哥.
不喜歡.
仇恨.
嫉妒.
以致將他賣了.
聖經告訴我們.
在第37章.
下面的經文就是說.
第14節.
因為突然之間.
眾兄弟報一個惡訊.
本來爸爸叫約瑟帶一些食物.
去到曠野.
放羊的哥哥.
去探望他們.
可能希望關係緩和.
誰知遠遠見到約瑟來的時候.
這個弟弟就要來下手.
想陷害他.
其中兩兄弟.
劉辯和尤大.
千萬不要.
我們不可以殺自己兄弟.
這個罪很大.
輾轉就將他扔到坑裡.
接著有些以實瑪利人.
去埃及做買賣.
反正見他見過一次.
都想打他一次.
見一次打一次.
這樣的意思.

$^{281}$那些兄弟.
就索性把心一橫.
將這個弟弟.
請上來.
賣了給以實瑪利人.
去埃及.
接著回來說一個這樣的消息.
給爸爸聽.
天衣無縫.
因為都走了去埃及.
回不了頭.
那個時代哪有這麼先進.
你可以隨時去監控他的行蹤.
你去找尋人.
沒有的.
十四節.
告訴他們.
藥室被撕碎了.
你看看這些衣服.
是不是你的兒子的.
那件彩衣.
爸爸幫兒子做的彩衣.
彩衣染了血.
有那些衣服在.
都不用驗DNA.
兒子的衣服.
現在染了血.
其實是野獸的血.
不是他自己的血.
那時不能驗DNA.
所以爸爸不知道.
兒子這樣說.
眾兄弟這樣說.
當然這樣相信.
接著告訴我們.
十四節.
三十七章.
雅各便撕裂衣服.
腰間圍上麻布.
為他兒子悲哀了多日.

$^{321}$他的兒女都起來安慰他.
他卻不肯受安慰.
說我必悲哀著下陰間.
到我兒子那裡.
弱勢的父親.
就為他哀哭.
實際上作者告訴我們.
米顛人這個以實瑪利人.
帶了弱勢在埃及.
將他埋了給法老的內臣.
這個護衛長波提弗.
所以我們讀者.
完全知道整件事.
但是當事人雅各.
這位老人家.
是什麼都不知道.
兒子最愛的兒子.
最愛的太太的最愛的兒子.
這個弱勢死了.
只有血衣.
請問各位.
你遭遇接連的打擊.
你會怎樣呢.
那件血衣你是不是燒了它呢.
還是保留.
時時拿出來去懷緬呢.
男性的眼淚未必很多.
但是那個心痛.
好像刺入心那種痛.
女士可能未必明白.
男士那種痛.
如果痛入你的心肺.
那些未必有眼淚.
女士容易一些.
不開心就流眼淚.
那個痛容易抒發.
男士難的.
在座不知道有沒有共鳴的弟兄.
有些好痛的事.
刺入你的心裡面.

$^{361}$那種痛是疏解不了的.
雅各回歸返迦南地.
太太在路上難產而死.
自己的大兒子和妾侍.
來去搞這些不倫關係.
最愛的兒子.
一下子這樣離世.
接連的打擊.
到底雅各是不是應該問一個問題.
為什麼是我呢.
Why me.
相信每個人遭遇這樣的不測.
自己的家庭的切膚之痛.
都一定會問Why me.
不過我們剛才看到的47張.
似乎雅各不僅僅問Why me.
如果一個人有信仰有神.
而持續都是問Why me的時候.
他沒有出路.
他的生命不能有更好的突破.
今天的講道.
是很想和大家一起去探討.
怎樣可以從Why me.
去到哪裡Why not me呢.
Why not me是一個進程.
如果Why not.
為什麼不可以是我.
和問為什麼是我.
是兩個人生態度.
是兩種敬拜.
如果弟兄姊妹.
今天你在你的遭遇困難.
你的厄運.
人際關係的厄運.
又或者可能在你的工作事業.
家庭都仍然有困局.
而你問的是Why me.
繼續問Why me.
神就不能夠在你的生命裡面.
成為一個很大的力量.

$^{401}$但當如果你能夠從一個Why me.
跳脫到一個因為經歷神.
而你能夠問為什麼不可以是我呢.
Why not是我.
為什麼不可以呢.
這樣就完全不同了.
眼光已經有一種.
因為神的同在.
而帶來的一個生命改變.
你開始會再問.
我怎樣可以.
來到這個我認為不好的處境.
來經歷經驗神的同在.
和生命的轉化呢.
就不再問Why not這麼簡單.
可能就開始問.
How我怎樣可以.
透過這件事.
我學到一些生命功課呢.
第46章因為長話短說.
不可以講太多.
若失輾轉飛黃騰達.
到底他是17歲就已經被他的哥哥.
這樣陷害.
因為聖經記載.
他30歲被法老召見.
到底是不是13年的含辛茹苦.
在埃及輾轉遇到很多奇妙的事.
不是他計劃.
沒得計劃.
一個奴隸在波提佛的家裡.
沒得他計劃的.
輾轉遇到很多奇妙的機會.
使他現在成為全埃及的宰相.
30歲見法老.
30歲見法老的時候.
就告訴法老.
你的夢我解給你聽.
有7年的大風年.
全地都會豐收.

$^{441}$不過7年的風年之後.
就會有7年的荒年.
荒年到一個地步.
沒有人再記得那7個風年.
法老召見若失的時候30歲.
接著輾轉我們去到第46章.
和大家記起這段經文.
45章尾先.
第21節.
以色列的兒子們就如此行.
若失照著法老的吩咐.
給他們車輛和路上用的食物.
因為已經召見完若失.
已經把戒指權力給了若失.
因為你既然解得通我這個夢.
又知道前面真的會發生.
現在你就幫我們拯救我們全地.
使到我們埃及和周邊的國家.
都可以得到保守.
不會在饑荒裡面.
大家在大饑荒死亡.
若失二話不說接受這個挑戰.
因為是神要他做的.
飛黃騰達.
第一件事是怎樣.
當然就回去迦南接爸爸.
第21節就告訴我們.
若失大照著法老的吩咐.
給他們車輛和路上用的食物.
又給他們就是那些兄弟.
各人一套衣服.
唯獨給便銀三百銀子.
五套衣服是親弟弟.
第23節.
送給他父親功勞十匹.
陀著埃及的美物.
母勞十匹.
陀著糧食與餅和菜.
為他父親路上用.
於是若失打發他兄弟們回去.

$^{481}$又對他們說你們不要在路上相爭.
他們從埃及上去來到迦南地.
他們的父親雅各那裡告訴他說.
這次講好信息.
你想想到底過了多少年.
如果17歲被陷害.
13年再加7個封年.
7個封年都未開始正式完結.
即是說有13至15年.
雅各是死了兒子的一個老人家.
自己最愛的兒子已經死於非命.
沒有見過.
但現在突然間那些兄弟回來講.
第26節.
告訴他說.
若失還在.
並請作埃及傳遞的宰相.
好像語場片張英才那些.
失憶了突然間回來認自己的父母.
因為不認識.
因為失憶了.
不知道自己的父母.
就是他的父母.
輾轉後來知道.
但父母已經死了.
因為語場片.
哭著跪著進了靈堂.
這裡告訴我們.
雅各心裡冰涼.
因為不信他們.
你會不會信呢.
叫兄弟仔去埃及掟米.
因為全地開始饑荒.
經過了七個封年.
經過了20年.
原來不是13至15年.
是20年.
這裡告訴我們.
雅各冰涼的內心.
不信.

$^{521}$第27節.
他們便將若失對他們說的一切話.
都告訴了他.
他們父親雅各.
又看見若失打發來接他的車輛.
心就掃成.
不是Tesla.
就是舒服過Tesla.
坐這些車下去見自己的兒子.
你想一下那個心情.
終於蘇醒了.
不再冰涼了.
見一下這個兒子.
第28節.
以色列說.
罷了 罷了.
我的兒子若失還在.
趁我未死.
以先我要去見他一面.
第46章.
敬拜的人生.
留意一下.
那種已經不是YM.
而是他怎樣去面對神.
留意第46章.
很精彩.
以色列帶著一切所有的.
起身來到別氏巴.
就獻祭給他父親以撒的神.
當然是他的神.
夜間神在異象中對以色列說.
留意聖經的描述.
在異象裡面對以色列說.
神怎樣叫以色列呢.
雅各.
他說我在這裡.
雅各這個名字.
由第32章已經不用了.
神的兒子.
神的王子.

$^{561}$就是以色列這個名字.
所以第46章這裡.
由第32章到第50章.
很少次數.
神親自叫這個名字.
叫雅各.
你有沒有花名的.
我小時候在木屋區長大.
有些花名街坊都認識.
很多年前.
在千禧廣場.
恆生銀行.
有個中庭的地方.
有些坐的街邊檔.
有個老人家坐在那裡.
他的腳不好.
腳痛.
我經過很熟面口.
認得是街坊.
他叫我那個花名.
小朋友的花名.
那個親切感.
如果大家不介意很噁心.
沒所謂的.
他叫我做豬屎櫃.
因為我家是養豬.
小時候.
有時沒有水.
我們要用腳去掃豬屎.
小時候家家戶戶都窮.
所以不用皺眉.
個個都是這樣.
他叫我這個名.
那種親切度.
因為我已經長大.
又傳道.
怎會介意這些.
小時候當然介意.
但長大了.
大家都經歷了很多人生.

$^{601}$他的兒子.
我和他一起長大的.
但已經不養他.
我都知道.
我的家姐和我說.
他就很坎坷地自己一個人住.
他叫我這個花名.
很深刻.
很親切.
以色列.
神叫他.
呼喚他.
雅各.
你留意神怎樣說.
他就說.
我在這裡.
神說.
我是神.
就是你父親的神.
你下埃及.
不要害怕.
因為我必使你在那裡成為大族.
我要和你同下埃及去.
我要和你同下埃及.
雅各的人生.
有神的陪伴.
有神的同在.
相信一個這麼利好的消息.
心臟受不了.
一個老人家.
最愛的兒子復活.
你不要以為很簡單.
心臟會爆發.
知道自己的兒子還沒死.
很嚴重.
但這次不怕了.
因為有神的應許.
有神的同在.
有神的陪伴.
一起下埃及.

$^{641}$所以你看到這裡.
也必定帶你上來.
約瑟必定為你誦宗.
所以輾轉.
就是剛才我們讀的經文47章.
約瑟在法老面前.
告訴法老.
我一生的年日.
是又短又苦.
其實有17年.
和他的兒子約瑟一起生活.
在埃及生養眾多.
而且等待整個民族.
出埃及 入迦南.
雅各就是這樣.
在埃及死後.
約瑟就將他的骸骨.
帶回上去迦南地.
到約瑟死的時候.
約瑟就叫他的族民.
要嚴嚴的起誓.
一定將來你們出埃及的時候.
不要留下我.
要將我的骸骨掘起.
起墳之後.
要帶回迦南地.
這是出埃及的12章.
約瑟的骸骨.
就跟隨著大軍.
進入迦南地.
弟兄姊妹.
痛的人生.
是必經的.
沒有人的人生是不痛的.
痛會使我們更加成長.
更加成熟.
但是痛的人生.
有沒有神在你的內心.
成為你的力量.
成為你的陪伴呢.

$^{681}$在大概兩個多月前.
未離開錦繡.
在洪水橋那裡事奉.
收到一個姐妹的電話.
當然是WhatsApp.
曾牧師.
我能不能將你的WhatsApp電話.
給我的妹妹呢.
她就說這個故事.
她說我在錦繡堂聽到.
覺得很好.
她就回去三藩市.
她在香港住的.
住沙田的.
她的妹妹很多年前.
已經移民去三藩市.
她就過去三藩市探望妹妹.
就將錦繡的鈴.
給她妹妹聽.
妹妹未信耶穌.
而且對福音都拒絕.
當然不會祈禱.
為什麼她要探望妹妹呢.
後來我就知道了.
是她妹妹親口跟我說.
2018年的時候.
她妹妹開車上班.
還有兩個街口就到辦公室.
被一個女大學生的車.
可能控制不了.
那裡都不是超級公路.
已經是內街.
那個女學生將她的車.
撞到推前.
推出去公路邊.
本來不是很傷.
但一推出公路.
剛好有車高速撞過來.
她整個車total loss.
人固然就.

$^{721}$由身體整個的以下.
變成終身的傷殘.
你可以想像得到這個痛苦.
姐姐就過去探望她.
跟她疏解她.
侮辱.
姐姐不知怎樣幫她.
因為真的很困難.
將圖拿給她聽.
就是我星期四晚的查經班.
一個不信的人.
這麼嚴重的車禍.
查經班是那些很深的信徒.
言經來的.
不是福音.
每一堂是過半小時.
妹妹是不是不知甚麼原因.
可能實在是太多時間.
每天坐在那裡.
等人去幫她做清潔.
幫她做一些日常生活的事.
她真的聽.
聽了一年之後.
就是剛才所講的故事.
叫姐姐.
求姐姐.
你能不能夠拿周牧師的電話給我.
我說當然無所謂了.
分享而已.
姐姐就問准我之後.
將我的電話給這個妹妹.
這個妹妹就WhatsApp我.
她每寫一隻字都說很難.
她說她夜晚三四點那些痛就來了.
痛到根本無法想像.
但她人就是動不了.
她很搞笑.
她說周牧師.
其實大家不認識的.
我聽了你的查經班聽了一年.

$^{761}$我完全不知道講甚麼.
因為根本我都沒有這些背景.
沒有這些聖經.
但不知為甚麼.
我的心裡面有很大的平安.
我的心裡面好像很多的困惑都沒有了.
我邀請她祈禱.
她說我嘗試吧.
她開始問一些神學問題.
輾轉和她對話這幾個月.
每朝早知道我差不多醒的時候.
開始工作的時候.
都不敢打擾我.
大概六點多七點就找我了.
問一些神學問題.
問我怎樣怎樣.
神是甚麼來的.
三一神是甚麼來的.
為甚麼她會知道呢.
原來有時聽完我的查經班.
她對這些概念很陌生.
她很想問下這些是甚麼來的.
現在她已經信了耶穌.
她現在為自己的丈夫祈禱.
她現在繼續有她的痛.
三四點醒又會痛.
但她說她現在不再迷惑.
心靈裡面不再有任何的迷惘.
她現在內心裡面有平安.
她有一天問我一個問題.
她說她很害怕.
神喜不喜歡我不回教會.
我回不了教會.
我跟她說.
你不是不回教會.
你是回不了教會.
有一天你的丈夫信主.
他可以推你去教會.
她就很開心.
釋懷了.

$^{801}$一個完全沒有行動能力的人.
怕自己回不了教會.
神不喜歡她.
弟兄姊妹.
你是不是還在問Why me的問題呢.
還是你的人生因為有主有神.
為甚麼不可以是我.
而當我問這個問題.
已經可以去到一個轉化.
就是How.
怎樣來榮耀神呢.
怎樣在我的生活困難.
際遇去經歷神.
祂的同在帶給我們的力量.
盼望呢.
我們一起祈禱.
多謝主.
讓我們能夠透過雅各的敬拜人生.
我們看到原來人生的際遇.
不是我們能夠掌握,預言和操控的.
我們知道每個人都有自己的困難.
信仰正正告訴我們.
有主的同在.
有主的力量.
有主在我們的生命裡面.
我們就能夠.
真的可以重新去解釋.
甚麼是困難.
幫助我們能夠有屬靈人的眼光.
和那個研判的能力.
以致我們看到的.
不是困難來壓制著我們.
而是困難使我們更加的經歷神的同在.
奉主耶穌基督的名.
阿們.
\newpage



\section{帖撒羅尼迦前書}
\label{sec:ih_tESQH0P8}
\textbf{2024-03-21 周四晚查經班《帖撒羅尼迦前書(1)》 講員:曾裔貴牧師}
\newline
\newline
連結: \href{https://youtube.com/watch?v=ih_tESQH0P8}{\texttt{https://youtube.com/watch?v=ih\_tESQH0P8}} ~~~~ 語音日期: 2024-03-22
\newline
\newline
\hyperref[sec:82RgV2iRB78]{\small{< < < PREV SERMON < < <}}
~
\hyperref[sec:index_chronic]{\small{[返順時目]}}
~
\hyperref[sec:index_scriptual]{\small{[返順卷目]}}
~
\hyperref[sec:u5__eddKCBg]{\small{> > > NEXT SERMON > > >}}
\newline
\newline
$^{1}$感恩天父給我們有新的書卷開始.
我們會看帖撒羅尼迦前後書.
還有我們都會去思想一下.
系統神學裡面關於末世論.
不同的派別對聖經的理解.
都給我們有少少的概念.
到底我們要怎樣去持守.
等候主耶穌的再來.
願神大力幫助我們.
我們今天晚上求主大力.
我們大家的新心靈都能夠建壯.
一起有一個渴慕神真理的心.
我們這樣祈禱.
奉耶穌基督的名.
阿們.
原因是為甚麼要看這兩卷書.
其實就因為很想講末世論這個神學的思想.
所以就選取了這兩卷書作為藍本.
因為都提過.
帖撒羅尼迦前後書.
特別是前書有五章聖經.
每章的結尾.
保羅都提到主再來的.
教會,個人.
怎樣去迎接主耶穌的再回來.
這個是帖撒羅尼迦前書的特色.
去到後書就更加講那個敵基督者.
更加講現在有一個攔阻他的.
這是第二章的經文.
那裡又是很明顯和末世論有直接關係的.
所以就很想透過這兩卷書.
當然一講到末世論.
就一定不可以不講的.
譬如馬太福音的24章.
是耶穌在鑒欖山的寓言.
末世的徵兆.
但是我們以前就比較有系統.
講過了啟示錄.
所以有些相關性的經文之前都講過的.
所以這次就會集中在.

$^{41}$新約保羅這一個書卷.
來看保羅的模樣.
給大家有一點這方面的了解.
開始帖撒羅尼迦前書.
我們又要看回那些背景.
帖撒羅尼迦城.
是在現在一個很現代的地中海的一個城市.
是希臘的第二大城市.
我們上一次都有機會去過.
當然只是去過海邊.
在薩羅尼加城.
比較重要的那個堡壘.
我們上次去那個 我忘記了那個名字了.
那裡來參觀.
看到整個愛琴海周邊.
是很美的地方來的.
帖撒羅尼迦城.
是亞歷山大大帝的妹妹.
來為他紀念而建的.
就是說是亞歷山大大帝的爸爸.
來打贏了加山德王.
幫他的女兒來建這個城.
所以是紀念他的妹妹亞歷山大大帝.
所以在薩羅尼加的海邊.
就有亞歷山大大帝的雕像.
所以大家如果有機會去過 你都知道.
我們上一次那個團的電影姐妹都有看過.
都有拍照留影.
就是今天的薩羅尼加.
就是這個帖撒羅尼迦城.
更加重要的是什麼呢.
就是因為在羅馬帝國的時代.
保羅的時代.
帖撒羅尼迦的城邦.
因為羅馬當時第一任的皇帝.
就是蓋薩奧古斯都.
或者叫做屋大維.
他和安東尼打仗.
帖撒羅尼迦城就剛好壓中了.
就是買屋大維贏.

$^{81}$很自然的.
那個皇帝掃平了敵對勢力之後.
安東尼被殺.
輾轉奧古斯都就是這個屋大維做了皇帝.
所以就給了當時帖撒羅尼迦城.
就有一個最惠國待遇.
就是免軍費.
還有很多的可能是駐幣或者駐軍.
都是不需要收那個軍費.
所以帖撒羅尼迦都是一個相當重要富庶的地方.
在保羅的時代.
你可以知道.
如果你看那個地理你就知道.
那個地方過去土耳其是一個必經的地方.
所以是一個很重要.
在保羅時代.
當然到了今天都是.
都是一個很重要的港口.
所以有機會大家如果有機會去.
是一個很漂亮的地中海的城市.
是僅次於希臘的.
很出名的.
這個是帖撒羅尼迦.
一些當時的背景資料.
我們看看《士道行傳》.
17章這段經文.
我們要多些聖經.
《士道行傳》的17章.
保羅和西拉.
16章就是說他們離開菲臘比.
所以菲臘比和帖撒羅尼迦都很近.
我們有機會去菲臘比.
在保羅的囚室.
我們有拍過照.
在那個舊城.
菲臘比.
都是一個算是富庶的地方.
我們有去過女底亞的教堂.
女底亞河邊的祈禱的地方.
所以16章的經文.

$^{121}$我們全部都去過了.
17章就是帖撒羅尼迦這一段.
他說保羅和西拉經過了暗非波利.
阿波羅尼亞.
來到帖撒羅尼迦.
就是說似乎他是由菲臘比.
他要連夜離開.
因為都不受歡迎.
甚至乎要坐牢.
因為他令到那個有錢人的.
那個侍女.
被鬼附的那個侍女.
就好回了.
那個主人就沒有了.
得到財利.
那些還財.
所以那個主人就抓了保羅坐牢.
和西拉兩個.
輾轉因為保羅是羅馬公民.
所以那些地方官都不敢.
沒有一個具體的罪名.
來囚禁保羅.
再加上就是.
在當夜.
那個監獄就有神跡.
大地震.
震爛了門.
但是那些人是沒有事的.
裡面的人一個都沒有事.
所以就很害怕.
就邀請保羅離開.
就去他家.
就帶領了他全家歸住.
後來輾轉就要提他出監審問.
那些人都知道保羅都不是善男信女.
所以就禁止他.
不要再亂說話.
就釋放他.
所以保羅就回去.
一放監之後就去了.

$^{161}$女底亞這位姐妹的家.
勸勉他們之後.
保羅就和西拉離開.
輾轉似乎當時都是走路的.
因為那些地方都是有些古代的王道.
都是很多經商的人.
做生意的人都會經過的.
所以不是荒山野嶺.
是很富庶的.
做生意的貿易的地方.
所以就經過了這幾個地方.
就來到帖撒羅尼迦加這個地方.
在那裡有猶太人的會堂.
保羅照他素常的規矩進去.
就是先向猶太的同胞去講福音.
什麼時候要去猶太人的會堂呢?.
這些冷知識.
一定是安息日.
星期五的黃昏.
就是開始安息日.
就是星期六的安息日.
到星期六的黃昏.
就是他們的安息日.
所以一定是安息日才去猶太人的會堂.
因為那個是聚集的地方.
我都提過我想問一下猶太人的廟.
在半山那裡.
中半山.
但是就沒有回覆我.
我們想團體去申請.
去參觀一下猶太的廟.
因為它不開放的.
但是如果團體呢.
以往我們試過一次開.
就是神學院的時候申請.
那時候都是我申請的.
現在就沒有聽到他回覆.
我看看遲些還有沒有機會.
這裡他們進了猶太人的會堂.
第二節.

$^{201}$保羅照他素常的規矩進去.
一年三個安息日.
安息日在這裡聚會的時候.
本著聖經與他們辯論.
因為猶太人安息日一定會讀聖經.
而且只能夠是男士讀經.
可以邀請不同的男士.
來上一個中央的台來讀經.
沒有固定的讀經員.
是會有安排就可以上去讀經的.
所以你路過都可以讀經的.
記不記得主耶穌回到拿撒勒.
人們就給他一本經書.
他就揭開.
揭開就是以塞俠書的61章.
主的靈告我.
他就是那個被釋放的.
差遣被擄的得釋放.
核硬得自由.
報告神越立人的欺凌.
接著不提審判.
那一章的第二節是講審判的.
還有神報仇的日子.
沒有提到那一章.
所以那些鄉親父老聽到之後.
哇 很好啊.
主耶穌講了這麼好的恩言.
所以他們個個都拍手掌.
因為是即興揭聖經.
所以都考你對聖經的認識.
如果你是猶太人.
當然一定要猶太人才能讀.
保羅又是被邀請.
來本著聖經.
我相信他自己都未必有那些手抄本.
一定是去到會堂.
有些卷軸.
那些拉比.
或者叫管會堂那些人.
就會給他經卷讀.

$^{241}$他就打開來讀.
這就告訴我們.
他就好像在那裡擦檔.
本著聖經和那些猶太人辯論.
辯論什麼呢.
講解陳明.
我必須受害.
從死裡復活.
又說我所傳語你們的這位耶穌.
就是基督.
講得很斬釘截鐵.
那就是說一定是基於舊約.
當時沒有新約.
是舊約.
那一定可能是以賽亞書五十三章.
或者詩篇.
譬如詩篇是110篇.
就是這位彌賽亞.
這位受告者.
諸如此類還有很多.
那保羅就很純熟.
本著聖經.
本著舊約的經卷.
再講得仔細一點.
是打開那些舊約的卷軸.
真的兩卷.
好像聖旨那樣.
那些卷軸來的.
在半山那個猶太廟.
是有摩西五經的真跡卷軸的.
所以是價值連城的.
是很重要的.
所以這裡告訴我們.
保羅就引述舊約聖經.
當然要講解.
要陳明.
就是說要evaluate.
要將它輾述.
舊約所預言的到底是誰呢?.
那保羅就著舊約那些預言.

$^{281}$來對號入座.
說你們所等待的彌賽亞.
就是耶穌了.
這樣的對號入座是很嚴肅的.
因為耶穌只是拿撒勒的木匠.
但是現在你說他是我們舊約等候的彌賽亞.
那就很大件事了.
你可以想像得到.
保羅就是預了.
可能受人歡迎.
更大的就是猶太同胞.
是他們心目中的彌賽亞.
和保羅所說的彌賽亞.
是差天共地.
怎麼會一個木匠的兒子就是彌賽亞呢?.
所以你可以預期得到.
保羅為了這個信息.
他一定會受很多的困難.
很多的苦楚.
因為大部分人都不會認為.
木匠的兒子是彌賽亞.
是救世主.
是神應許的救主.
所以你可以想像得到.
保羅就要為這個信息.
付一個很大的代價.
第四節.
他們中間有些人聽了勸.
就服從保羅和西拉.
並有許多乾淨的希利尼人.
或者希臘裔的人.
尊貴的婦女都不少.
貼薩諾尼加.
不要忘記.
貼薩諾尼加都很近菲律賓.
女底亞是第一個歐洲的婦女信主.
她是一個做生意的.
賣紫色布匹的.
女底亞都是一個有錢人.
我們去過女底亞的故居.

$^{321}$不知道是不是真的.
那個河邊.
那裡祈禱的.
據聞是女底亞的家.
她有個小教堂.
就當是.
但是的確有些尊貴的婦女.
所以你要留意.
保羅原來都做到一些.
有知識.
有社會地位.
有一定的社會階層.
這些人信主.
幫助保羅的宣教.
都有很大的幫助.
這個就是當時的情況.
不過第五節就很明顯.
因為猶太人期望的彌賽亞.
不是這樣的.
所以其實是打破了.
很多猶太人的.
他們那個幻滅了的理想.
因為彌賽亞應該是.
很英明神武.
很有影響力等等.
怎麼會是一個木匠呢.
所以我們經常看耶穌傳.
那個耶穌這麼帥.
頭髮又這麼型.
那個一定不是的.
真正的猶太人.
真正的古代的猶太人.
還要是奴婦大眾的階層.
不是拉比.
不是在猶太的議會裡面那些拉比.
所以我們拍耶穌傳.
都不是按真故事去拍.
因為你找一個又粗魯又矮.
整個臉都是皺紋.
這樣就不好看.

$^{361}$找一個帥哥來做耶穌.
有時候先入為主都麻煩.
其實主耶穌是沒有階型美容的.
其實是很滄桑的.
飽歷滄桑的一個青年人.
三十多歲.
肯定是三十歲就出來傳道.
又講遠了.
利命人出來事奉是二十五歲.
這些利命人是有教的.
二十五歲就可以出來.
不過做小孩 做學徒 做神學生.
因為要三十歲才可以正式在聖殿裡面.
古代就是沒有聖殿之前.
就是回膜那裡正式做.
所以有五年的實習.
所以神學院四五年原來適合聖經.
利命人都要五年.
跟那些老師 跟那些師父五年.
湯 牛 祭祀等等.
原來都要學的.
所以這裡告訴我們.
因為網上看不到.
下面就告訴我們.
不信的猶太人心裡面疾到.
焦聚了市井匪類.
搭火成群 聳動合成的人.
闖進耶孫的家.
要將保羅西拉帶到百姓那裡.
找不著他們就把耶孫和幾個弟兄.
拉到地方官那裡.
喊叫說那攪亂天下的也道者女來了.
這些罪名一定是叛國.
因為當時是貼索羅尼加城.
是得到當時羅馬帝國的一個.
你可以說是Blessing.
有一個最惠國的待遇.
惠就是優惠的惠.
你可以想像得到.
如果你這樣煽動.

$^{401}$對羅馬的政權受到威脅.
那就是搞亂天下.
所以猶太人都懂得煽動群眾.
找一些市井匪類來搭火成群.
好像給錢的臨時演員.
一車人帶著色變的帽子來吶喊.
說保羅在這裡令我們受到一個威脅.
好像衝擊到羅馬的政權.
所以第七節.
耶孫收留他們.
原來在耶孫的家裡.
來講解.
但耶孫是管會堂也說不定.
暫時我們還不知道.
第七節說.
這些人違背了該殺的命令.
說另外有一個王耶穌.
沒錯了.
真是一個煽動.
因為主耶穌是那位君王.
所以很特別.
我們看的貼索羅尼加前後書.
真是說主耶穌是再來的.
五章的聖經.
每章的結尾都說主耶穌是再來的君王.
保羅好像明知而故犯.
人家現在說你一個這樣的罪名.
說你們傳的耶穌.
就是取締羅馬帝國的.
是很忌諱的.
當時是尼祿做王的.
尼祿是一個很差勁的君主.
所以你可以想像到.
他們用這個堂而皇之.
就可以入罪了.
就可以告保羅了.
說有另外一個王耶穌.
所以這些你們要記住.
因為我們看完這些背景之後.
再讀貼索羅尼加前書.

$^{441}$或者後書.
每次都說主再來.
你就可以想像到.
在保羅的心底裡.
他所等待的彌賽亞.
耶穌君王的再臨.
不是地上現在的.
即刻的挑戰羅馬帝國.
而是他將會第二次降臨.
不過什麼時候呢?.
聖經就沒有說了.
到我們今天讀貼索羅尼加前書.
又過了二千年了.
聖經的應許還沒有兌現.
也就是說我們都應該和保羅一樣.
他是第一批等候主回來的人.
我們去到現在21世紀了.
我們是21世紀的聖徒.
等候有一天的日子.
整個宇宙.
整個世界是要變天了.
是真正改變了.
我們要等候這一天.
所以基督徒有生之年.
必須是.
要不就是你和我安息主懷.
如果不是呢.
就是主再來.
我們不需要安息主懷.
身體改變去見耶穌.
所以都是盼望來的.
我們所以人生.
真的盡快要整理好我們自己.
我們應該怎樣在很有限的人生去事奉.
直到我們不能夠了.
神說我們的時間已經圓滿了.
我們就準備安息主懷.
所以某程度對主再來.
我們預備自己安息禮拜的.
詩歌啊.

$^{481}$經文啊.
找誰講道啊.
都差不多要準備了.
不是負面的.
不是悲哀的.
是正面的.
正正就是我們要去到榮耀的天家.
進入永恆的福樂.
當然經過福樂之前是有些痛苦的.
但是基督徒就因為明白.
我們只不過是閉上眼.
就進入永恆.
所以我們應該要多一些這樣的思想.
這個思想成為了我們.
推動我們每天都把握我們的時間去事奉神.
當然我們講到這麼多.
我再補充.
第八節.
眾人和地方官聽見者說驚慌.
因為一般人想的.
所謂耶穌做王.
耶穌再回來.
一個死了的木匠.
猶太人的囚犯.
他說自己是猶太人的王.
這個罪名.
在十字架上.
現在說他復生.
說他會做王.
復生都不要緊.
還要做王.
做什麼王.
羅馬帝國.
當時他們都不知道有中國.
羅馬帝國的時候.
應該是漢朝過些的時間.
不知道有一個遙遠的東方.
有一個這麼大的中國人.
所以一定是挑戰羅馬政權.
所以他們就很害怕.

$^{521}$因為帖索羅尼加.
真的和國家的關係是很好的.
所以當他們這麼驚慌的時候.
於是就取了耶穌和其餘之人的寶藏.
就釋放了他們.
也都是不可以隨便去無故的審問.
或者未經定罪就來判刑.
所以叫做恐嚇耶穌.
或者可能要他們交出保羅和西拉.
所以第十節.
當時就見不到保羅和西拉.
可能他去了其他地方.
那些群眾就來到去捉拿了耶穌和其他的弟兄.
第十節.
弟兄們隨即在夜間.
打發保羅和西拉往比利亞去.
二人到了那裡.
就進入猶太人的會堂.
你可以想像到.
保羅離開了帖索羅尼加.
去到比利亞一個小小的地方.
特別強調這個地方的人.
很喜歡言經.
不單止喜歡.
而且他們是真的有一個.
很渴求對聖經真理的認識.
所以比利亞人.
能夠信主.
能夠很明白神的救恩.
就是他們因為天天考察聖經.
要曉得那個道是真還是假.
所以不單止是滿足頭腦知識的興奮.
而是很想尋道.
所以多有信主.
這個就是帖索羅尼加的背景.
你可以想像到.
保羅全程去到會堂.
似乎有一定的受歡迎.
因為一些尊貴的婦女.
一些地方官都喜歡聽.

$^{561}$甚至乎可能都已經信主了.
但是猶太人有他們的背景.
有他們的舊約的包袱.
通不過.
怎麼可能.
如果你說耶穌就是我們的彌賽亞.
祂為什麼要死呢.
為什麼會死得這麼慘呢.
為什麼祂的出生這麼寒微呢.
過不了他們的那個filter.
因為很多舊約的猶太人.
都已經假設.
彌賽亞的光臨就是我們國家的復興.
彌賽亞就應該是一個很英明的救世主.
怎麼可能是一個我們摸得到見得到的人呢.
還要祂的出生這麼寒微呢.
所以這些就變成自己製造了自己的基督.
自己對號入座.
舊約的基督不應該是這樣的形態.
所以這個神學就是猶太人他們自己對彌賽亞的誤解.
所以見到耶穌都要將祂置於死地.
就是因為他們心目中的彌賽亞不是這樣的.
所以我們都要小心我們自己的神學.
或者我們所謂理想的教會.
我們都要小心.
有時候好像不行了.
家庭教會都不行了.
不去教會聚會了.
在家裡自己讀經就行了.
不用這麼麻煩了.
我和幾個弟兄姊妹自己坐在家裡讀聖經就行了.
那就沒有這些問題了.
這個是你的理想教會.
那就大件事了.
因為原來你幾個人的理想.
其實就將真真正正的聖經要教導我們的扭曲了.
聖經反而告訴我們.
真正的理想就是每個人去實踐.
每個人在教會裡面彼此去學習相愛.
那個才是達至理想.

$^{601}$但是永遠都不會有的.
是我們一閉上眼睛進入永恆.
那個就是理想的開始了.
就真是去到永遠沒有罪的一個地方了.
就是神的居所了.
所以我們記得不要將我們自己的神學.
就當成是聖經的神學.
所以我們要讓聖經去修剪我們.
令到我們那些先入為主.
我們的主觀.
我們個人的單方面的意願.
都慢慢可以更新.
可以更新.
所以研經就有這個好處.
透過一些讀開的經文.
有系統地讀.
慢慢會令到你有一些改觀.
就不能夠一節金句.
因為一節金句未必適合任何人.
但是整卷整卷.
按它當時的氣氛背景.
來去研讀的時候.
慢慢會疏離了很多.
我們有時很先入為主的觀念.
因為觀念是要慢慢.
透過真理去解釋.
才能夠得到那個澄清.
不是要確定哪個人.
而是大家都客觀地.
看回聖經的那個真理.
那我們轉入貼薩諾尼加前書.
或者我們可以揭一些.
《少道行傳》的經文.
貼薩諾尼加城的人.
有些是後來跟隨保羅傳道的.
比如19章的《少道行傳》.
19章的就是23節.
那時.
這個就在爾忽所了.
爾忽所有一個大的騷亂.

$^{641}$那時因為這道起的擾亂不少.
保羅去到不同的城市.
都帶來不同地方的人有兩極.
有些是歸順耶穌.
生命改變.
有些是越來越硬.
那爾忽所那個舞台.
那個廢墟.
我們去過了.
那個真的很震撼.
全歐洲最大的廣場.
爾忽所.
可以同一時間坐上萬多人.
很大的.
那些設計.
那些建築.
在任何一個地方.
一說話.
任何一個位置都聽到.
真的很神奇.
這些就是.
你去到那個置身現場.
你會發現.
保羅當時那個傳道.
帶來的那個生命的影響.
這裡說.
因為這個道擾亂不少.
有一個銀匠.
名叫底米雕.
是製造亞底米神鵪.
亞底米女神.
還是.
蓄立在爾忽所的守護神.
我們當日.
11月見到的.
都已經是很輝煌的了.
何況保羅的時代.
那個亞底米女神.
是更加的香火鼎盛.
很多周邊的城市.

$^{681}$不是爾忽所的地方的人.
都去成為了世界七大奇景.
古代的一個地方.
爾忽所是很出名的.
真的香火鼎盛.
就是拜這個亞底米女神.
有些做神鵪.
就是神壇的底米雕.
這些銀匠.
因為他使得這個生意的人.
沒什麼事做.
因為個個都不再拜偶像.
信耶穌.
第25節.
他聚集他們和同行的人.
就說.
眾位你們知道我們是依靠這個生意來發財的.
這個保羅.
不單只是在爾忽所.
幾乎在亞細亞全地.
引有迷惑許多的人.
就說.
人手所做的不是神.
這是你們所看見所聽見的.
這樣不單止我們這個事業被人藐視.
就是這個大女神.
很高的一個柱.
亞底米的神廟.
都要被人輕忽.
連亞細亞全地和普天下所敬拜的大女神.
的威榮也要消滅.
你可以想像引起那個哄動.
引起那個改變.
你可以想像到.
譬如我們香港的大嶼山的大雄寶殿.
黃大仙廟.
雄磡的觀音謝父.
所有的都因為福音.
而很多人不再拜了.
歸順主了.

$^{721}$你就可以想像到那個復興是多麼大.
當然迫不及待就自然增加.
所以你會看到.
28次這樣說.
眾人聽見就怒氣填空.
因為有一個人.
這個似乎是銀像的工會的公頭.
這個有名的.
底米雕.
引起那些工會的同仁.
來去暴動.
要捉拿保羅.
要置他於死地.
因為令到他們沒有生意做.
不再信.
不再覺得這些神恩是靈驗.
很大件事.
衝擊到一個屬靈的爭戰.
要告訴我們.
滿城都轟動.
眾人就拿著保羅.
同行的馬其頓人.
該由和阿里達古.
齊心就湧進戲院裡面去.
就是那個劇場.
就是異不所的舊城區的劇場.
馬其頓人和帖撒羅尼迦很相近.
馬其頓是那個地方的省.
省會.
我們現在廣州廣東省.
廣州就是城.
我們香港就是保安縣.
你可以想像到.
馬其頓人有兩個.
是跟隨著保羅的.
該由和阿里達古.
另外第二十章.
亂定之後.
我就不說那個暴亂.
保羅請門徒來.

$^{761}$勸勉他們就辭別起行.
就是在異不所.
往馬其頓去.
走遍那一帶地方.
用許多話勸勉門徒.
然後來到希利尼.
在那裡住了三個月.
就是回到希臘.
將要坐船往敘利亞.
就是回頭了.
回到耶路撒冷.
猶太人設計要害他.
他就定義從馬其頓回去.
和他一起去的有亞細亞人.
回去的有比利亞人.
不羅斯的兒子.
索巴特.
比利亞我們說過了.
他離開帖撒羅尼加之後.
就去了一個小小的城市.
就是比利亞.
在比利亞之後.
這裡有一個叫索巴特.
即是說他去到那個地方.
有些年青人跟隨他做傳道.
現在有這個索巴特.
留意一下這裡.
帖撒羅尼加人阿里達古.
和西宮都.
就是跟隨著保羅傳道.
跟著還可以看.
剛才我們上文說的第29節.
馬其頓人該由和阿里達古.
這裡告訴我們原來阿里達古.
馬其頓的帖撒羅尼加城的人.
廣東省廣州市的曾蕊貴.
大概是這個意思.
即是說原來阿里達古.
是原原本本帖撒羅尼加人.
即是說都是屬於馬其頓省.

$^{801}$是一個大的統稱.
這裡阿里達古和西宮都.
都是帖撒羅尼加人.
還有這個特比人該由.
還有這個提摩泰.
提摩泰又是路斯德和特比地方的人.
即是說都是亞西亞人.
亞利桑那裡的地方.
其實應該叫土耳其.
是土耳其的地方.
當時帖撒羅尼加過去就是土耳其.
所以帖撒羅尼加是一個很重要的港口.
他們的船運其實很近.
地中海走很快.
人走可能會比較久.
所以帖撒羅尼加.
有比利亞,路斯德,特比.
這一帶的地方其實是很近的地方.
你可以想像到保羅不單止是提摩泰和西拉.
西拉本身是耶路撒冷猶太的拉比領袖.
但是跟隨保羅.
原來保羅在外邦人當中.
記住我剛才提的這些全部是外邦人.
路加醫生記載很多年輕人.
是來自不同的省市.
不同的城市的外邦人.
提摩泰就特別一點.
他是希臘裔和猶太人的身份.
因為他媽媽是猶太人.
他爸爸是希臘人.
所以當時有些年輕人跟隨保羅去傳道.
你可以想像到.
保羅從來沒有想過他自己年輕的時候做拉比.
他從來沒有想過十幾年後.
他信了主耶穌.
他成為主耶穌的使徒.
居然有些學生跟隨他.
用生命跟隨他.
是外邦人.
所以你看看神的國度.

$^{841}$神的工作多麼奇妙.
你沒有想過將來.
你想到就不是神的奇妙了.
是神的安排.
所以我們不要幫神去設我們的限制.
保羅沒有想過他會有外邦人做他的學生.
跟隨他鞠躬盡瘁來為他福音.
因為看到這個老師一心一意為福音.
所以這班年輕人來自五湖四海.
跟隨保羅去傳道.
其實不容易.
留意一下這裡的描述.
還有這個特比人 戈尤和提摩泰.
又有亞西亞人 推基古和特羅非摩.
這些人先走.
在特羅亞等候我們.
我們就包括了路加醫生.
第六節 過了除教的日子.
除教的日子 正月十四.
正月十四逾越節.
正月十五就是除教節.
所以大家要對舊約熟.
才知道路加醫生描述的事情.
過了除教的日子.
我們就從菲律賓開船.
五天就去到特羅亞.
所以真的坐船.
不要想他們走得那麼辛苦.
很快的.
地中海那邊繞來繞去.
很快的.
雖然他們不是用電動.
但是就算是爬.
用那些槳.
都很快的.
或者是帆船.
都很快的.
所以其實當時是很密切的來往.
地中海和土耳其.
是沒有什麼關卡可言的.

$^{881}$不需要什麼簽證那些東西.
所以你可以想像到.
保羅就走遍馬其頓一帶.
走遍.
三個月.
他離開了以弗所.
即是說土耳其.
他就是去探望教會.
去到馬其頓一帶.
探望教會.
接著就提醒有哪些人跟他.
就是周邊的外邦的城市.
都是那些歸主的年輕人.
來跟隨保羅.
他們先來到.
去到特朗亞.
保羅可能有一點點延遲.
輾轉由菲臘比坐船去特朗亞.
即是回頭的意思.
因為他本身是由特朗亞過去菲臘比的.
有沒有留意第十六章的經文.
馬其頓異象.
就是由特朗亞那裡.
得到那個馬其頓異象.
有一個人晚上.
就求保羅.
在夢裡面.
你過來吧 我們需要福音.
可能那個是路加醫生也說不定.
就不知道是怎樣.
保羅知道.
現在神為我們開路了.
我們現在就隨即.
第二天.
他在等候應該去北面的.
高加索那邊.
即是再北上的莫斯科.
蘇聯.
南下.
向西去地中海.

$^{921}$就清楚了.
去地中海.
所以正正式式.
第二次宣教.
就是去菲臘比那邊.
這個就是小說的背景.
我們對帖撒羅尼加前書.
先有這個概念.
我們現在就可以回到.
帖撒羅尼加前書.
這些有機會在講道的時候.
也有機會講的.
當然因為講到崇拜只能夠半個小時.
一定是集中一個題目.
但現在我們是整卷去看.
我們就會找回.
讀新約一定要留意.
每一個書信.
或明的書信.
一定是有寫作的作者.
有寫作的作者.
前人已經幫我們搞定了.
認為了這些書卷.
是誰誰寫了.
當然現在現代的學者就會質疑.
譬如以下所書不是保羅寫的.
就有很多的推論.
不過全部是推論.
沒有辦法證明.
因為那些原稿都不在了.
要留意了.
有作者.
有收信的對象.
可能是教會.
或者是個人.
一定要問一個問題.
我打開這卷書.
作者寫給第一批收信的人.
他想有什麼重點呢.
他要帶出他寫作的原因呢.

$^{961}$即是說作者的動機是什麼.
什麼引發他要寫這封信.
這是正常的.
你在英國.
你的兒子現在在讀大學.
他寫一封信回來.
當然因為現在是電郵和WhatsApp.
就沒有了這麼遙遠的.
寄信收信的動作.
因為一秒就收到了.
就沒有了這個.
但仍然都是.
你打開WhatsApp.
兒子現在WhatsApp我.
那你一定要看看他說什麼.
說完之後.
你要怎樣回應.
首先你要知道對方.
他那個WhatsApp的內容.
接著你可能想了一晚之後.
第二天你就回覆.
用你的WhatsApp回覆.
跟他寫的書信一樣的.
你一定要找到寫作的原因.
寫作的動機.
和他最終要達至什麼的結果.
一定是這樣的.
即是說他期望教會.
收到這封信之後.
會有什麼的改變.
一定是試圖寫信的動機.
和當時不是每個人都有權.
有資源 有錢去書寫的.
是很高的成本.
因為不識字的人大把.
不識字.
那個是很普遍的.
不是每個人都識字的.
根本不會每個人都識寫作的.
所以那個是一個成本.

$^{1001}$另外就是書信.
當時的紙張是極之昂貴的.
根本就沒有我們今天的紙.
是用當時的蒲草等等.
來作為書寫和傳遞的.
所以是很昂貴的一件事.
再加上送信的人.
郵差 郵局.
當時還不是很普及的.
所以是要很官銜 很重要的.
才可以這樣來傳遞.
民間很少這樣去飛鴿傳書.
是不可能的.
所以你可以想像到2000年的考古.
聖經比任何古代的文獻多得很離譜.
我們現在找到古代的名著文獻.
考古出來的 是很少的.
但是聖經新約的.
我們出土的考古是非常之多.
多到一個地步是遠遠拋離第二位的.
比如希臘的文人.
比如現在想起的.
比如我們希臘的哲學家.
蘇格拉底 柏拉圖.
這些出名的已經是.
還有很多其他不同的希臘的書籍.
但是那個比例是相當懸殊的.
所以你可以想像到聖經的重要性.
還有那些人繼續抄寫去傳遞.
但是保羅的時候絕對是昂貴的.
更加加增了寫作一卷書.
那個很重要的目的.
一定是有原因而寫的.
我們後世的人.
我們讀聖經.
我們一定要做一個還原的動作.
因為我們不知道.
帖撒羅尼迦加人.
他們問保羅什麼問題.
或者帖撒羅尼迦加這個地方.

$^{1041}$當時面對什麼困難.
我們不準確知道的.
因為我們沒有看到.
保羅他所見到的事情.
可能有人去到保羅身邊.
告訴保羅.
譬如三章就有了.
提摩泰在帖撒羅尼迦加城回來.
保羅就知道.
實際他們當時的情況.
所以我們只能夠透過新約的作者.
他們所寫的書卷.
來去還原那個教會.
到底遇到什麼困難.
保羅要寫信.
可能是真理的澄清.
可能是教會內部的一些問題.
保羅給他一些規範.
有一些教導.
有一些指引.
有些可能是個人的.
譬如提摩泰在提摩泰前書.
他應該怎樣在爾弗所的大教會.
去做傳道 做牧師.
保羅就給他有教導.
譬如第三章的提摩泰前書.
就是教提摩泰.
要在各城設立長老.
揀選那些領袖.
然後就提出什麼條件.
才符合執事.
做長老的要求.
很明顯是教提摩泰.
怎樣去自會.
所以你就要透過作者.
他們的教導.
去還原那個教會.
到底有什麼問題.
這個就是讀新約的書信.
大家必須要注意的地方.

$^{1081}$我們來看看.
這個帖撒羅尼加前後書.
剛才的背景已經說了.
就是保羅在耶穌的家裡.
和阿西拉.
但是有猶太人聚集.
來去詆毀.
攻擊保羅.
說他是搞亂天下.
說他們這群傳道人.
傳另外一個叫王.
名字叫耶穌.
所以我們就要留心這件事.
後來當耶穌被保釋之後.
他們回到家裡.
保羅和阿西拉就不想牽連太大.
他還有很多事情要做.
所以趕急離開帖撒羅尼加城.
去到另外一個地方.
剛才我們看過的就是.
去到比利亞的地方.
輾轉後來就去到.
希臘的雅典.
其實在下面.
也是不難的.
坐船去到那個港口就可以進去.
不過我們坐車就辛苦一些.
因為開了公路.
坐車就要很長的時間.
但是古代其實坐船很快的.
因為那些海岸線就已經定了.
我們看帖撒羅尼加的前書.
我們先看第一章.
保羅西拉提摩泰.
寫信給帖撒羅尼加.
在庫臣和主耶穌基督女的教會.
願因為平安歸於你們.
一方面我們也要問.
如果整個亞西亞.
土耳其.

$^{1121}$那個亞底米女神.
是這麼香火鼎盛.
有沒有影響歐洲的人呢.
因為只是加一個海岸.
他們的交易 他們的貿易.
很常有的.
不是封閉的.
所以土耳其的內陸.
這個亞底米女神.
其實如果你去過你都會看到.
雖然我們很遠的拍攝.
是一個很多乳房的女神.
是一條柱子.
但是古代的時候就不是只有一條柱子.
就是生育和耕種的豐收.
兩樣是實種的.
農業社會最需要人口.
勞動力.
還有風調雨順.
所以對於內陸土耳其的地方.
這個亞底米女神.
真是有求必應.
所以是非常香火鼎盛的.
但是正正本來的出現.
就打破了這種的迷思.
就令到很多的人.
我們提過.
以弗所.
史上有20章.
有很多以弗所人.
信了耶穌.
扔了家裡那些走邪術的書.
迷信的書.
拜偶像的書.
可能占卜.
星座.
其他的運程.
扔了.
燒了五萬塊錢.
古代的以弗所的遺跡.

$^{1161}$你可以想像到.
一個大規模的復興.
你看為復興.
人家看為就是.
你在搞亂檔.
我們是靠這個做偶像.
為生的.
所以保羅就要面對.
這個傳道帶來的迫害.
所以.
要留意.
十三封的保羅書信.
每一個的文案.
其實都是一個研究.
你要研究保羅.
他的文案.
一般都是公式.
都是這樣的格式.
當時的書信的格式.
因為有些學者都研究.
也有其他的出土的.
其他的不是使徒的書信.
是當時的格式.
但是保羅的書信.
有需要研究.
就是他的文案.
有時直接反映教會的情況.
比如哪一本書是完全沒有文案呢.
不文案.
就是大件事.
卷卷都問.
是這個文案.
還是那個文案.
卷卷都問.
卷卷都不問.
其實是加拉太書.
大家回去看看.
是不文案的.
一開波就罵.
罵得很嚴重.

$^{1201}$你就知道他們一定是出現了一個異端.
近乎他們會走錯路.
我們今天不說那裡.
我們看回這裡.
留意一下這個文案.
和主耶穌基督裡面的教會.
保羅強調不是一個地方.
一個建築物.
是這個建築物裡面.
這個建築物裡面.
是屬於父神的.
是屬於主耶穌的.
雖然我們今天每個人都知道.
但是往往.
一來到教堂就忘記了.
最重要的主角就是神.
就是主耶穌.
所以我們很少.
來到教堂.
會有一個肅然起敬.
開始有一個準備的心靈.
所以朝聖茶室很重要.
真是朝聖.
去到神的面前.
所以你看到保羅每一卷書.
都提到教會是神的教會.
不要小看這些文案.
特別是古代.
他們很緊張那些廟宇.
廟宇就很清楚了.
因為保羅是偶像.
每個人都知道.
就算我媽媽叫他黃大仙.
他都肅然起敬.
上一代都是這樣.
進廟宇就要拜他.
但是我們很多時候.
來到教堂.
一來可能是二樓.
甚至一個教堂.

$^{1241}$一個建築物教堂.
我們慢慢都忘記了.
最重要的心態.
就是對做物主的敬愛敬畏.
記住.
在父神和主耶穌.
裡面的教會.
當然我們有時候.
透過一些讓你見到的圖像.
譬如有一些擺設.
讓你有一份敬虔.
這些是好的.
但是記得.
不要靠那個外表.
是我們的心靈建立.
和神的關係和接觸.
記住.
在父神和主耶穌.
裡面的教會.
第二節.
禱告的時候提到你們.
在神我們的父面前.
不住的紀念你們.
留意一下這個信望愛.
這個組合.
在保羅書信經常出現.
有沒有哪一本書.
你們很熟悉的?信望愛.
臨前十三章.
如今常傳的.
如今常傳的有什麼?.
有信有望有愛.
哪樣最大?.
愛.
沒錯.
這個就是我們.
很多時候會看到保羅的信望愛.
不過你又要回去做一個研究.
保羅提到的信望愛.
每一本書都不同.

$^{1281}$所以不是一個公式.
只是這次.
第一章已經不同了.
和歌羅西書.
和臨前十三章.
已經不同了.
就算是前後書的信望愛.
都不同了.
你就會覺得.
原來這麼多樣性.
原來信望愛不是我們想像的那麼簡單.
你可以想像到保羅對於.
這三樣東西.
似乎.
是我們由信主那一刻.
開始要建立的.
要累積的.
屬靈的.
的保護.
是我們屬於.
信望愛不是物質.
信心.
盼望 愛心.
是由我們信主那一刻開始的.
如果這三樣東西.
是直到永恆的話.
我們就要問一個問題.
我們累積.
就好像你儲分.
比如你入油儲分.
你的apps一路加.
入了多少錢就加多少分.
我們有沒有這個apps呢.
信望愛.
保羅書信裡面說了很多.
今晚給大家一個功課.
保羅的信望愛.
你把它列出來.
你做一個研究.
你就會有多樣性.

$^{1321}$對保羅提的信望愛.
這裡開始了.
因信心所做的功夫.
當然是傳道.
或者是服事聖徒.
或者是教會裡面的服事.
因為信而來的工作.
功夫.
就是那個.
因為愛心.
所受的勞苦.
因盼望.
我們主耶穌基督.
所傳的忍耐.
所以信望.
愛都在這三個.
的範疇裡面.
不過提到.
信帶來的.
在帖撒羅尼加這裡.
信心帶來的就是會做事.
為神.
為教會做事.
為使徒做事.
愛心就會.
因為有愛心而來.
忍受那些苦楚.
可能是做得辛苦.
可能是有時人誤會了你.
而帶來的可能言語的.
衝突.
因為盼望.
有對象.
盼望是有對象的.
所以永遠盼望.
是有一個對象的.
這裡就很清楚了盼望是誰.
主耶穌基督.
肯定是說祂的再來.
肯定是說.

$^{1361}$對祂的再來的盼望.
而你心存的.
是忍耐.
所以你會看到.
保羅就開始示範.
信望愛.
他的理解.
和他的展述.
為什麼呢?.
因為帖撒羅尼加的前書.
是最早的保羅的書信.
就是第二次的旅行.
報道宣教.
他寫的.
所以很肯定是很早期的書信.
那麼早期.
他已經對信望愛有這樣的.
理解.
我們就要問了.
我們有沒有因為信心而來做事呢?.
問自己.
我們有沒有因為.
愛而受苦呢?.
愛.
你就願意忍受一些苦難.
一些可能是.
物質的苦難.
心靈的苦難.
等等.
有沒有因為盼望主耶穌.
而有心存的忍耐呢?.
忍耐等候.
是保羅羅馬書裡面說的.
如果我們羨慕那個日子.
就忍耐.
等待.
因為所見的盼望.
就不是盼望.
誰會盼望見到的呢?.
但如果我們盼望那個未見到的.

$^{1401}$就必定會忍耐等候.
這個就是保羅第一次說.
信望愛.
第一次將信望愛.
放在聖經裡面.
勸導他的門徒.
就是帖撒羅尼加教會的門徒.
我們首先.
因為還沒有內文.
還沒有裡面的內容.
詳細說當時教會面對什麼困難.
不過.
他先給了門徒.
帖撒羅尼加教會的門徒.
有了他自己對.
信望愛的.
理解.
就是說希望他們都能夠.
有這樣的.
愛心.
其實我要再補充多一點.
其實保羅這裡告訴我們.
是紀念他們.
現存的信望愛.
帖撒羅尼加人.
這個教會裡面的.
信望愛.
因為你們.
因信心所做的功夫.
你們因為愛心.
而受的勞苦.
因為盼望主耶穌.
而生存忍耐.
是對他們的誇讚.
留意第四節.
被神所愛的弟兄.
我知道你們是蒙簡選.
這些又是很重要的經文.
我們早幾堂都說過.
那個金錢官.

$^{1441}$財富.
金錢不是萬惡之根.
貪財才是萬惡之根.
所以清教徒.
看金錢看得很正面.
就是只需要你自己.
清楚你是蒙簡選.
神就會祝福你的事業.
你的人生.
所以他們就會很拼命工作.
很認真地去儲蓄.
然後很認真地去服侍神.
就是清教徒.
他們的原則.
對金錢.
所以這裡就呼應.
你們是蒙神所愛.
而且要認定你們是被簡選的.
這個很重要.
是一個受苦難的教會.
這個就更加好像.
黑暗裡面的明燈.
一點的燈光.
你認定自己是被簡選.
一點的燈光.
你認定自己是被神所愛.
你認定自己是被神簡選.
你就面對多大的難處.
你都會.
比永恆.
短暫的生命威脅.
短暫的痛苦.
不足掛齒.
你會.
聖靈會告訴你.
主耶穌都是這樣度過.
他最後的人生.
他看到的就是.
前面的喜樂.
就是他的犧牲.

$^{1481}$他的贖罪可以拯救所有的.
罪人悔改.
所以他的犧牲是必須.
但必須裡面也有.
有一種喜樂承載著他.
面對困難的.
那個痛苦.
這是希伯來書的十二章的經文.
那忍受罪人這樣.
頂撞的你們要思想.
免得疲倦灰心.
因為他前面的喜樂.
他就忍受.
那個羞辱.
背起十字架.
就是主耶穌.
所以大家要留意這個.
這個揀選.
大家有沒有.
認定自己是被神揀選.
你怎樣認定呢.
怎樣肯定你自己.
是被神揀選呢.
所以你必須對神學.
對聖經的教導.
我們若認自己的罪.
神是信實是公義.
必要赦免我們的罪.
洗淨我們一切的不義.
若人一輸的一張九折.
一定的.
必要的.
記住.
是神的信實.
是神的公義.
祂的信實.
即是說祂應許了一定兌現.
祂的公義就是.
耶穌已經犧牲.
必須要赦免.

$^{1521}$我們的罪.
這樣才叫待贖.
這樣才叫我們被揀選.
保羅的揀選論.
是很多的.
譬如二分所思又有.
一張.
所以揀選也是在保羅的神學裡面.
一個大的主題.
如果我們有機會再讀.
肯定帖撒羅尼迦人.
是被揀選.
被神所愛.
是好像雙父雙成.
既被揀選.
又是被神所愛.
誰先走呢.
不知道.
第五節.
因為我們的福音.
揀選之後就解釋.
你們怎知道自己被揀選.
第五節就解釋.
因為.
我們的福音.
即是我們去傳福音.
去到你們那裡.
不單止在乎言語.
也是在乎權能和聖靈.
所以靈恩派.
有時就會捉錯用神.
權能,聖靈,暴道.
有Power.
這個權能.
就是Dynamic.
Dynamic就是炸藥.
原文那個字就是炸藥的意思.
權能.
所以保羅這裡說的.
不是超自然的神蹟.

$^{1561}$是說福音帶來的.
生命影響.
和聖靈將人的.
心靈的錯誤的神學觀.
改變.
要改變一個人的神學觀很難.
要改變一個人的立場.
很難.
要改變猶太人對舊約的.
彌賽亞觀念也難.
外邦人就拜慣了偶像.
滿天都是神佛.
所以這裡保羅告訴我們.
我來到你們的.
帖撒羅尼迦城.
到你們改變.
有些人能夠深信.
自己是被選擇,被神所愛.
為什麼呢?.
因為我們所傳的福音.
不單止是言語.
Verbal這樣說.
而是有在乎.
In power.
In spirit.
在聖靈裡面.
在權能裡面.
神的權能彰顯.
所以記住那個彰顯.
不要太局限.
一定是大的神蹟.
一定是病要治好.
往往神有時都會這樣做.
尤其是一些落後的.
宣教地區有時都會這樣做.
譬如有些阿拉伯人會做夢.
見到耶穌.
但如果是一個城市.
很多福音已經傳了.
有時要改變那個價值觀才最難.

$^{1601}$改變一些人.
他們看世俗的眼光.
才難.
他們有錢.
有好的家庭,有好的事業.
要他們信耶穌是很難的.
往往這些價值觀改變.
是In power.
In spirit.
在聖靈裡面.
令到人的生命改變.
是.
古羅說.
福音就可以帶來.
這樣的一個轉變.
這個就是.
選擇的原因.
原來我們聽明白了福音.
很深信.
和認定.
降伏在福音神的.
兒子的面前.
我們就是蒙揀選.
反過來.
我們是信耶穌.
得永生.
但你一回望.
原來我是被揀選的.
千萬不要以為.
神沒有揀我所以我信不了耶穌.
你們不用勸我.
神揀選你了.
神沒有揀我.
首先我和你都不知道.
他有沒有被揀選.
那個笑我們被揀選的人.
就說你們被揀選.
我選不了.
所以我信不了.
其實他都不知道自己.

$^{1641}$是不是真的被揀選.
他可能將來就會信.
所以記住永遠在我們人的層面.
是明白了福音.
信耶穌.
而且感受到一種權能.
改變我們的生命.
改變我們的價值觀.
這個就叫信耶穌.
我們知道的就叫信耶穌.
但我們回望.
原來這就是神的揀選.
行不行.
很清楚將揀選論.
比較簡易的解釋出來.
就是我們.
原來我們能夠信耶穌.
原來真是恩典.
原來不是每個人都容易信耶穌.
有些人搞來搞去.
他都不信.
原因可能就是他沒有被揀選.
這個我們不知道.
所以我們千萬不要用這個.
作為一個理由.
活該你都沒有被揀選.
所以你信不了耶穌.
如果我們這樣說.
等於我們是神.
我們真的不知道.
我們唯一知道的就是.
必須保守自己在神的救恩裡.
所以記得一次相信.
一直相信.
就可以得到永生.
所以我們要保守自己在揀選裡.
我們要去到將來.
才知道自己真是永恆夢揀選.
當然.
我們一直信就已經肯定.

$^{1681}$我們是被揀選.
被神所愛.
但這個不能作為我們去論斷.
其他未信的人的一個理由.
的一個說詞.
不然人家就更加反感.
你叫我信耶穌我想想沒問題.
我沒有被揀選我就沒辦法接受.
首先就是我們自己都不知道.
你怎樣能夠斷定對方.
是沒有被揀選呢.
就算那一刻是很敵黨福音.
保羅是敵黨福音.
最終他都信耶穌.
神的時間一到就信.
所以我們記得這件事.
所以其實第五節就解釋了.
那個揀選的最基本定義.
尤其是.
不要說笑.
在極大的苦難裡面.
信徒能夠.
持守的.
不再是那些.
憧憬那些開心的日子.
而是你被揀選.
你能夠堅持到底.
是因為你知道自己要在那個揀選裡面.
你要是被.
神揀選的門徒.
所以你要堅持到底.
你才有份於.
神永恆的救恩.
這一點是.
盛載著在很多苦難的.
信徒坐牢的.
殉道的.
的信徒.
他們那個唯一的內心的.
你可以說就是他們裡面.

$^{1721}$那個神學.
如果用我們.
都是看倚天屠龍記大那些人.
就是有神功護體.
我們就有這個.
像張無忌那樣.
有平坤大羅兒.
神功護體怎樣打都沒事.
自己自我修復.
就是這個.
深信自己是被揀選.
這一點可以渡過.
你最黑暗的困難.
當然這裡是一個很簡單.
在禱告裡面.
一位的使徒.
看出他們有.
信忘愛.
他們的.
內心他們的事奉.
是有信有忘有愛.
是分析出來的.
信心所做的功夫.
愛心所受的勞苦.
因為盼望主耶穌而心存忍耐.
然後就提到他們被神所愛.
所以保羅對這個教會.
真的讚譽有加.
到現在還是讚他們.
好第六節.
並且.
你們在大能之中.
蒙了聖靈所賜的喜樂.
領受真道.
就效法我們.
有效法了主.
所以你們看到幾部曲.
在困難.
那個困難似乎就是那些猶太人的迪巴.
耶孫的家都被抄家.

$^{1761}$被人抓去審問.
這些是一個.
基本的屈辱.
更加是一個.
有生命威脅.
對信仰.
帶來的衝擊.
在這個大能之中.
真信心就有.
聖靈的喜樂.
好像主耶穌教過.
殺身體不能殺靈魂.
你不用怕他.
到時聖靈會指教你怎樣說話.
所以大家要背熟聖經.
如果聖靈到時指教不到你.
你又賴聖靈就不好了.
當然.
那是最深層的.
裡面的平安.
大家有沒有經驗過.
聖靈賜的喜樂.
這裡保羅告訴我們.
當你在大能之中.
蒙了聖靈賜的喜樂.
這些個人的遭遇.
例如你自己.
有一些病患.
是你自己面對的.
其實你要記住.
在月黑暗裡面.
主耶穌的同在.
憤外顯得突出.
這個要大家經歷.
所以基督徒正常.
你平日.
你自己好好追求.
去經歷將來.
在那些方年的日子.
就成為你自己很大的.

$^{1801}$心靈力量的支柱.
如果你現在都不好好追求.
去聚會.
去建立和主耶穌的關係.
去到那些方年的時候.
商界就說.
現在.
經濟好的時候.
豬都會飛.
經濟不好的時候.
豬就會掉下來.
因為牠們沒有了翅膀.
環境不好.
或者大家現在游泳.
大家都沒有問題.
退潮的時候.
誰沒有穿泳褲.
大家都看到.
商界就會有這些說笑.
現在香港沒有大巴.
每個人都開心的回教會.
每個人都齊齊齊.
沒問題.
訂位就去喝茶.
沒問題.
每一次的聚集.
有沒有建立你自己的生命.
這個就沒有人知道.
這個就要在患難.
方年的時候才知道.
就是個個游泳.
原來有些人沒有穿泳褲.
你就知道.
去到困難的時候.
為什麼有些人會否認耶穌.
就是他平日根本沒有裝備.
他來到教會原來不是朝拜神.
他是朝拜那些人.
所以那些人一不好.
他就不好了.

$^{1841}$我不回這個教會了.
因為這班人都不好.
首先問一問.
就算那些人不好都不關你的事.
你是朝拜神.
所以很多事我們要問一個問題.
我們在大能之中.
有沒有蒙聖靈所賜的喜樂.
平日我們就要.
經驗聖靈在我們生命裡的幫助.
到大能的時候.
就可以經驗到聖靈的喜樂.
接著還有.
聖靈給我們不只是無盡的喜樂.
是領受真道.
就是說.
平日的日子.
我們是在一份喜悅裡.
去學習真理.
是領受真道.
領受.
不只是聽.
領受真道.
就效法我們和效法主.
就是說使徒的教訓.
他們都明白了.
和實踐.
在實踐的過程.
更加他們會.
使徒帶他們去認識耶穌.
不只是效法使徒.
是效法主.
你可以想像得到.
這個就是帖撒羅尼迦教會.
得到保羅的愛戴的原因.
所以剛才提到.
阿里達古.
西宮都.
這兩位帖撒羅尼迦的年輕人.
信主 跟隨保羅傳道.

$^{1881}$你可以想像得到.
保羅在那裡.
雖然不是一個很大的城市.
沒有爾忽所那麼大.
但是都有很大的福音的果效.
接著就說第七節.
甚至.
你們.
真是誇張到一個地步.
你們成為了示範單位.
你們作了馬其頓和阿該亞.
阿該亞就是整個哥林多.
那個地方.
阿該亞省.
希臘一個很大的省.
阿該亞當時.
哥林多就是阿該亞省.
最東南面的地方.
甚至你們作了.
馬其頓和阿該亞.
所有信主之人的榜樣.
因為主的道.
從你們那裡已經傳揚出來.
你們向神的信心.
不單止在馬其頓和阿該亞.
所以阿該亞和亞西亞是不同的.
亞西亞就是.
敘利亞和土耳其那邊.
東面那邊.
西面這邊就是阿該亞.
就是歐洲南面的地方.
就在各處都傳開了.
這裡說.
他們有什麼被傳開呢.
就是向神的信心.
所以不需要我們說什麼話.
不需要我們誇耀你們.
我們就說.
你們是神的信徒.
我們不需要誇耀你們.

$^{1921}$第九節.
因為他們自己已經報名.
我們是怎樣進到你們那裡.
他們就是.
馬其頓一帶地方.
和阿該亞省.
一帶地方的弟兄姊妹.
教會.
對帖撒羅尼迦教會的稱讚.
對神的信心.
向神的信心.
有什麼具體的行動呢.
原來你們是怎樣離棄了偶像.
歸向神.
要服侍這位又真又活的神.
第十節.
等候他而止.
所以每一章的結尾.
都是說主在來的.
帖撒羅尼迦前書.
等候他而止.
從天降臨.
就是他從死裡復活的.
那位救我們脫離將來憤怒的耶穌.
這裡有很多的形容詞.
主在來會是怎樣呢.
從天降臨.
就是呼應《使徒行傳》第一章.
耶穌升天之後的應許.
天使和那些使徒說.
你看到他怎樣升上去.
就怎樣下來.
這個其中一樣.
現在我們說第九節.
就是馬其頓和阿該亞省的人.
馬其頓和阿該亞省的人.
誇耀帖撒羅尼迦教會.
他們是向神的信心.
向神的信心原來在第九節的解釋.
就是用一個行動.

$^{1961}$去呼應.
認識神帶來的那個信心.
就是離開.
拜偶像.
不再拜.
命運.
訴諸於神明.
拜的是真神.
就是這位又真又活的神.
就是說以前都拜的.
不過現在發現了.
神是唯一.
就拋棄了所有其他.
這些是很多.
譬如印度人最難.
因為是眾神的.
日本人最難.
也是眾神的.
尼泊爾人也是.
不拒絕的.
因為印度教就是什麼神都要信的.
尼泊爾就是很多印度教的.
所以.
中國人就.
不算什麼都拜.
但是還是難去拜一位.
獨一無二.
又真又活的神.
所以這裡就決定了.
向神的信心.
原來是必須.
離開偶像歸向.
獨一真神.
唯一要敬拜.
所以主耶穌說我們不可以又拜神.
又拜馬門.
就是說原來財利.
是令到我們有安全感.
原來我們拜的是安全感.
所以就變成.

$^{2001}$我們有時就.
沒有那麼專注在神那裡.
這裡就告訴我們.
向神的信心.
必須是要.
專一.
《申命記》的六章.
四節.
任何猶太人都會背.
將來我們有機會學希伯來文.
背那個六章四節.
以色列你要聽.
耶和華我們的神.
是獨一的主.
接著什麼.
盡心盡性盡意盡力.
愛耶和華你的神.
這個就是《申命記》六章四節.
猶太人他們有一個盒子.
拉比.
帶著一個盒子去安息.
有一個盒子.
要留意一下.
那盒子裡面的經文.
就是六章四節.
《申命記》.
這些冷知識.
是他們很重要的金句.
以色列你要聽.
就是這個.
耶和華我們的神.
是獨一的主.
所以一脈相承.
向外邦人說的.
是獨一的主.
不過多了一樣東西.
猶太人就沒有這一樣東西.
就是第十節.
所以猶太教是不可以上天堂的.
第十節.

$^{2041}$猶太人就沒有了這一句.
等候他兒子耶穌.
從天降臨.
就是神已經設立了一個.
的救恩.
就是差遣他的愛子來到這個世界.
所以猶太人只有舊約.
他們不信新約.
他們沒有新約.
所以更加不會相信你們所說的耶穌.
我問過那個拉比.
上一次我們.
2008年的時候.
我們神學院.
就有神學生去參觀那個猶太廟.
我親自問那個拉比.
在希伯來大學過來的.
你們怎樣看新約.
他很客氣的.
英文很好的.
又容易聽.
這個是你們基督教的.
耶穌不是我們的.
就是這樣的.
當然他未必否定.
他是猶太人.
但是這是你們的.
你們基督教說的救主.
和我們猶太人沒有關係.
所以保羅這裡你可以想像得到.
在保羅的內心.
他講這一節聖經.
他其實內心是很痛.
就是自己的同胞.
相信獨一真神.
不敢拜偶像.
但是問題就是.
最重要的.
神的唯一的救恩.
就是已經在他的愛子身上.

$^{2081}$人看見了我.
就是看見神.
《約翰福音》17章.
認識你獨一的真神.
認識你所猜來的愛子.
這就是永生.
我們講過《約翰福音》.
還是好像很莫言.
所以第十節.
是在這個時候.
所以第十節.
是保羅.
他已經跳出了猶太教.
他已經超越了.
他自己的同胞猶太教.
那個錯誤神學.
他過去是捍衛.
這個舊約的神學.
他這個猶太教的神學.
但是現在因為主耶穌向他顯現.
他現在傳遞的.
是一個真正人類的救恩.
就是等候他.
兒子耶穌降臨.
神的獨一的兒子.
這位又真又活的神的兒子耶穌.
就是從死裡復活.
你留意一下.
保羅強調再強調.
主耶穌是死而復活的.
你可以想像得到.
其實很難相信的.
死了就怎麼能夠復活呢.
一個人怎麼可能會復活.
還要永活呢.
所以這個一定是很重要的神學觀.
是我們基督教才有的.
其他的宗教都沒有的.
所以我們相信的就是.
一位死而復活的耶穌.

$^{2121}$也就是說如果他都能夠.
衝破死亡的生命的改變.
我們在地球上.
人類都是那些.
不外乎生老病死的東西.
我們都可以交託給神.
我們學習交託.
不要因為自己擔一些事.
就變成抑鬱.
我不反對輔導的.
不過如果每次都說是輔導.
一來你很有錢又貴.
不如你找傳統人聊聊天.
祈禱請我喝茶算了.
我不否定輔導的不要誤會.
真的.
不過如果每次都說.
我不開心我很苦.
那你想想主耶穌是死而復活的.
我們等候這位復活的耶穌.
我們所信的是一位復活了的耶穌.
是一位死了.
三日三夜再復活的耶穌.
現在在天上的耶穌.
我們的信仰.
應該有一種銀聖.
應該有一種.
可以經得起一些考驗.
當然有很大的困難.
需要有彼此的同行.
禱告.
不過這裡.
給我們看到一樣很重要.
保羅的神學已經超越了猶太教.
不只是舊約.
不只是以色列.
你要聽我們神是獨一的主.
我們盡心盡性愛祂.
不過.
到了保羅的書信.

$^{2161}$加了耶穌進入.
這個三一論裡面.
等候的是.
這一位.
復活的耶穌.
接著還要留意.
最後這句.
這個復活再來.
是要有憤怒.
對誰的憤怒呢.
是救我們脫離將來的憤怒.
是救我們信祂的人脫離將來審判的憤怒.
所以你想想我們屬主的.
帖撒羅尼迦的弟兄姊妹.
是這樣的心態去等候.
主耶穌的再臨.
所以.
整張聖經是對整個帖撒羅尼迦城.
是一個讚賞.
他們的信心.
傳到馬其頓省.
阿哥亞省.
他們向神的信心.
他們是怎樣向神的信心呢.
就是離棄.
那些偶像歸向.
又真又活的神.
和等候主耶穌的再來.
我們要問我們自己感受.
感孫家的弟兄姊妹.
或者將來聽到這一段經文.
的其他教會的弟兄姊妹.
歐洲的三藩市的.
我們要問一下.
我們是不是.
這個信心.
傳到其他人都看到.
這個人真是信耶穌的.
他有行動改變的.
生命改變的.

$^{2201}$而且他現在在地上.
沒錯 每天都忙的.
但是有一個終極的盼望.
就是等候主耶穌的再來.
而且這個再來不只是興高采烈.
有些人是痛苦的.
因為我們是被救贖的.
我們脫離將來的憤怒.
神永恆的憤怒.
但是我們已經脫離了.
所以我們是被脫離的行列.
有些人是要面對.
那個永恆的沉淪審判.
所以這裡又要逼使我們要傳福音.
和我們自己要一直在一個蒙簡選.
這樣的情況下.
今天就說到這裡.
剛剛說完一章 一起祈禱.
多謝天父的恩典.
我們開始了貼上羅尼加的前書.
我們很盼望.
我們看到這個教會.
是這麼得到保羅的讚賞.
他們的信望愛.
是我們需要學習的.
他們因信心所作的功夫.
是真的有功夫的.
他們因愛心.
所受的那個勞苦.
所面對的勞苦.
和因為盼望主耶穌.
而來得到那個忍耐.
幫助我們能夠在今生.
我們認定自己是蒙簡選的.
幫助我們能夠在今生.
我們認定自己是蒙簡選的.
更加我們能夠真是有一個.
學無主的再回來.
不是使我們古古怪怪的.
去深山避難.

$^{2241}$而是我們積極的.
在世界.
來得到見證我們的信心.
我們對人的愛心.
我們對福音的堅持.
願神保守幫助我們.
奉主耶穌基督的名.
阿們.
\newpage



\section{約翰福音 12:12-19}
\label{sec:u5__eddKCBg}
\textbf{2024-03-23 晚堂《耶穌榮耀進城》約翰福音 12\_12-19 講員:黃頌言姑娘}
\newline
\newline
連結: \href{https://youtube.com/watch?v=u5--eddKCBg}{\texttt{https://youtube.com/watch?v=u5--eddKCBg}} ~~~~ 語音日期: 2024-03-26
\newline
\newline
\hyperref[sec:ih_tESQH0P8]{\small{< < < PREV SERMON < < <}}
~
\hyperref[sec:index_chronic]{\small{[返順時目]}}
~
\hyperref[sec:index_scriptual]{\small{[返順卷目]}}
~
\hyperref[sec:94pJmtJo_JE]{\small{> > > NEXT SERMON > > >}}
\newline
\newline
約翰福音 12:12-19
\newline
\begin{longtable}{cl}
\hline
\hline
章節 & 經文 (和合本修訂版)\\
\hline
12:12 & \begin{tabularx}{0.7\textwidth}{X} 第二天，有一大群上來過節的人聽見耶穌要來耶路撒冷， \end{tabularx} \\ \\ \relax
12:13 & \begin{tabularx}{0.7\textwidth}{X} 就拿著棕樹枝出去迎接他，喊著：「和散那，以色列的王！奉主名來的是應當稱頌的！」 \end{tabularx} \\ \\ \relax
12:14 & \begin{tabularx}{0.7\textwidth}{X} 耶穌找到了一匹驢駒，就騎上，如經上所記： \end{tabularx} \\ \\ \relax
12:15 & \begin{tabularx}{0.7\textwidth}{X} 「錫安的兒女啊，不要懼怕！看哪，你的王來了；他騎在驢駒上。」 \end{tabularx} \\ \\ \relax
12:16 & \begin{tabularx}{0.7\textwidth}{X} 門徒當初不明白這些事，等到耶穌得了榮耀後才想起這些話是指他寫的，並且人們果然對他做了這些事。 \end{tabularx} \\ \\ \relax
12:17 & \begin{tabularx}{0.7\textwidth}{X} 當耶穌呼喚拉撒路，使他從死人中復活出墳墓的時候，同耶穌在那裡的眾人就作見證。 \end{tabularx} \\ \\ \relax
12:18 & \begin{tabularx}{0.7\textwidth}{X} 眾人因聽見耶穌行了這神蹟，就去迎接他。 \end{tabularx} \\ \\ \relax
12:19 & \begin{tabularx}{0.7\textwidth}{X} 法利賽人彼此說：「你們看，你們一事無成，世人都隨著他去了。」 \end{tabularx} \\ \\
[1ex]
\hline
\hline
\end{longtable}
$^{1}$各位弟兄姊妹晚安.
我們一起祈禱吧.
親愛的主 多謝你.
求你今晚更新我們的心.
賜下你的聖道.
讓我們聽見可以聽得明白.
多謝主 奉耶穌名求.
阿們.
按照基督教的聖力.
復活節前一個星期日.
就是中之主日.
是紀念人群拿著中樹枝.
迎接主耶穌進入耶路撒冷的日子.
然而中之主日.
又叫做苦難主日.
因為歡慶的日子背後.
其實是主耶穌基督迎接苦難的開始.
接下來的一周.
其實都是主耶穌的受難周.
而節期的意義.
是讓我們可以超越時空.
去經歷在其中.
今天我們再一同思想.
耶穌基督光榮進入聖城的情景.
去感受當中的張力.
主耶穌和當時各種勢力正面的交鋒.
耶穌的神的世界.
和神族的世界激烈的交戰.
神的心意和世人的心意彼此去角力.
而最終將主耶穌釘在十字架上.
我們今天看約翰福音十二章.
前面少許的背景.
其實當時是逾越節.
而逾越節是紀念以色列人.
整個民族被救贖的大日大事.
而這件救贖的大事.
是一個民族在受苦.
在哀求.
而神聽到他們的哀求.
解救他們的大事.

$^{41}$昔日這個民族在埃及受苦.
他們受政治的奴役.
受苦工的奴役.
他們哀求.
好像代表著世上所有在困苦.
被奴役的人一樣.
神聽到他們的哀求.
神解救他們.
這個救贖的日子.
是一個宗教的節日.
亦都是民族解放的解放日.
慶祝這個宗教節日.
同時紀念著民族的解放.
另一個說法就是.
紀念著人性的解放.
紀念著人性尊嚴的重獲.
亦都是紀念著能夠再次體會生命的意義.
和福樂.
這個多重意義的節日.
代表著民族解放的日子.
逾越節之前.
原來耶穌做了一件事情.
就是讓拉撒路從死裡復活.
死人復活這件事.
轟動了整個耶路撒冷.
今天我們看的經文是約翰福音12章.
12至19節.
耶穌榮耀進城.
但是原來榮耀進城之前.
都發生了一些事情.
我們看看前幾節.
就是約翰福音12章9至11節.
這裡沒有的.
如果是的話可以翻開聖經.
約翰福音12章9至11節.
可以聽我讀出.
有許多猶太人知道耶穌在哪裡.
就來了.
不但是為耶穌的緣故.
也是要看祂從死裡復活的拉撒路.

$^{81}$但祭司長雙耳.
連拉撒路都要殺了.
因為有許多猶太人.
為拉撒路的緣故.
回去信了耶穌.
原來有許多猶太人知道.
耶穌在那裡.
但是他們來到不是為了耶穌的緣故.
而是為了看從死裡復活的拉撒路.
大家都很想湊熱鬧.
非常好奇.
一直密謀要殺耶穌的祭司長.
連拉撒路都很想殺了.
因為猶太人去了信耶穌.
相信耶穌就是他們期待已久的彌賽亞拯救者.
所以他們蜂擁而至.
擁戴耶穌為王.
12至13節那裡.
Yoyo可以幫忙.
是的 謝謝.
在第二天.
有許多人上來過節.
他們聽見耶穌將到耶路撒冷.
就拿著鐘樹枝去迎接祂.
喊著說.
「我撒拿.
奉主明來的意識列王.
是應當稱重的」.
原來12章一開始.
是提到拉撒路和主耶穌一起吃飯.
是逾月節的第六天 即是星期六.
第二天即是甚麼時候呢.
第二天即是星期日.
逾月節其實是提醒.
以色列的民族猶太人.
他們的神是如何拯救他們.
亦都是提醒他們.
在先知書裡面描述那位彌賽亞.
他們將會出現一位拯救者.
縱然猶太人有一些都相信耶穌是彌賽亞.

$^{121}$但原來他們心目中彌賽亞拯救者.
是讓他們從民族的壓迫裡面.
去解放出來的拯救者.
他們很想獲得政治的解放.
民眾拿著鐘樹枝是充滿著政治的意味的.
在主前164年.
即是猶太人領袖瑪加比.
在希臘安提亞哥的政權當中.
重獲耶路撒冷聖殿.
並將其潔淨獻給神的時候.
原來他們是手執著鐘樹枝.
列隊遊行.
去淨化聖殿.
猶太人重獲耶路撒冷聖殿之後.
23年.
當猶太人國民在西門的領導之下.
在主前的141年.
完全獲得政治獨立的時候.
他們同樣都是拿著鐘樹枝.
猶太民族反抗希臘政權英勇的歷史裡面.
原來鐘樹枝象徵著宗教政治意義.
是不言而喻的.
除了鐘樹枝.
他們口裡面亦都喊一些口號.
「和撒拿」.
「奉主名來的」.
「以色列王」.
「是應當稱重的」.
而「和撒拿」是.
「求你拯救」.
「求耶和華拯救」的意思.
現在「和撒拿」我們都會用作讚美的用詞.
而這一句其實是引用詩篇118篇26節的經文.
但是原來.
原文當中是沒有「以色列王」這個詞.
是群眾自己加上去的.
這個富有宗教和政治含義的名號.
加上拉薩路從死裡復活.
這一群群眾真的非常之興奮.
對彌賽亞有極度的期望.

$^{161}$他們相信彌賽亞是帶領他們打一場聖戰的.
而耶穌就會是他們這場聖戰的領袖.
當時宗教領袖和當時的政治領袖.
在這個緊張的日子.
總是千方百計很想維持和諧的局面.
當時宗教領袖和政治力量聯手.
希望一切都可以和諧.
他們繼續做他們的宗教生意.
繼續享受當中的利益.
就在聖殿當中賣牛羊夾子.
群眾的「王」和「國度」.
和耶穌自己的理念原來是很不同的.
我們可以看見十二章.
除了說到拉撒路復活.
原來更加有說到瑪利亞.
用一根極貴的珍拿大香.
羔撫耶穌的腳.
耶穌預言.
「這是為我安葬之人去存留的」.
當耶穌使拉撒路復活.
耶穌同時也卻預言自己的死.
所以第十四,十五節那裡.
我們看到強烈的對比.
是耶穌強烈的對比.
「耶穌得了一個盧驅就騎上.
如經上記著說.
「舍安的民,不要怕.
你的王騎著盧驅來了」.
按照主耶穌慣常的做法.
祂是會走路的.
按照距離.
其實平常人都是會走路的.
本來如果主耶穌走入耶路撒冷的話.
就會淹沒在朝聖的人群當中.
平靜的進去.
平靜的出來.
但主耶穌這次卻故意的坐牢.
進入耶路撒冷.
如此這件平常的事就變成不平常了.
當群眾期待讓人從死裡復活的耶穌.

$^{201}$可以帶領他們去打一場聖戰.
可以騎著戰車,戰馬.
耶穌卻預言自己會死.
祂騎的是代表和平的盧驅.
《約翰福音》沒有記載.
耶穌是如何得到盧驅的.
然而這個舉動的意義卻是顯明不過的.
祂是以色列的王.
祂不像瑪加比那樣.
坐著戰車,戰馬進入耶路撒冷.
也不像所羅門王那樣.
在榮耀的進入城.
反而祂是像撒迦利亞先知所預言的王一般.
在撒迦利亞書九章九節中這樣說.
「赤安的民應當大大的喜樂.
耶路撒冷的民應當歡呼.
看你的王來到你這裡.
是公義的並且施行拯救.
謙謙和和地騎著盧.
就是騎著盧的驅子」.
主耶穌明白祂進入耶路撒冷的目的.
祂是要對抗假宗教對人的壓迫.
也要對抗不正義的制度對人的壓迫.
祂不單止用暴力.
祂不是用暴力的方式.
而是向人顯示真愛和和平的方式.
不會一下子改變世界.
祂是讓所有人可以看見.
可以改變世界的力量是何在的.
第十七至十八節說.
說到一群看神蹟.
臭熱鬧的群眾.
當耶穌呼喚拉撒路.
叫她從死復活出墳墓的時候.
與耶穌在那裡的眾人就作見證.
眾人因慶見耶穌行了.
這神蹟就去迎接她.
原來看神蹟臭熱鬧的這群群眾.
呼喊「和撒拉」.
擁戴耶穌為王.

$^{241}$同時原來他們都是喊.
「釘死祂!釘死祂!」的那群人.
有人相信耶穌是為了看神蹟.
有人相信耶穌是為了有飯吃.
有人是為了信耶穌的時候.
就可以一飯豐盛.
可以橫財就手.
有人相信耶穌是為了保平安.
然而耶穌來祂是要受苦的.
如果信耶穌是為了有飯吃.
當沒有飯吃的時候.
很容易就會信心軟弱離開耶穌.
會變成「由呼喊和撒拉」.
變為「喊釘死祂!」.
最近有一位初信的家長很可愛.
他問我.
「咦?為什麼信耶穌的時候他也會生病?」.
「連牧師也會生病?」.
「連傳道人也會生病?」.
他的想法我很明白.
因為我們經常為生病的人祈禱.
而神會聽我們的祈禱.
我們也會繼續為這些人祈禱.
但如果靈命好.
就不會生病的時候.
這其實是成功神學.
信耶穌就一定會成功.
如果是這樣的話.
人會很難接受信耶穌是有苦難的.
也很難背起十字架.
不要誤會.
信耶穌不是一件痛苦的事.
只是我們不會免疫.
我們也會感染COVID.
然而我們不同的是.
耶穌是和我們一同受苦的.
祂很明白我們.
縱使我們在風暴裡.
在困難的裡面.
我們也不會失去盼望.

$^{281}$不會失去平安.
這是我們救主拯救我們的方式.
是以愛的方式.
是以和平的方式.
在《約翰福音》十二章十九節.
講到法利賽人彼此說.
「看!你們是徒勞無益.
世人都追隨他去了」.
之前密謀要除滅耶穌的法利賽人.
他們互相指責.
然而一個重點是.
「世人都隨從耶穌去了」.
「世人」的意思是各種各族的人.
大概去迎接耶穌的人.
是很多不同種族的人都去了那裡.
耶穌拯救的是世人.
是每一個種族的人.
群眾給予耶穌熱烈的擁戴歡迎.
宗教領袖都控制不了這個場面.
而我們看到門徒第十六節那裡.
原來這些事門徒起先是不明白的.
等到耶穌得到榮耀之後.
才想起這些是指著他寫的.
並且眾人果然向他這樣行了.
門徒不明白耶穌為何要騎著驢進城.
亦不明白先知書所說的預言是指著耶穌.
縱然門徒跟隨了耶穌已經一段時間.
但他們不明白.
耶穌得榮耀的時候.
即是他釘十字架的時候.
從死裡復活的時候.
耶穌讓聖靈住在我們的心裡.
我們的保衛師聖靈.
會教我們去明白一切的真理.
我們的神常常都會用我們意想不到的方式.
去完成祂的救恩.
跟我們想像中的是不同的.
耶穌容許我們不明白.
但後來祂會讓我們明白.
祂的聖靈會幫助我們明白祂的奧秘.

$^{321}$讓我們看見祂的拯救.
最近我都有一件不明白的事情.
最近我去了芬蘭.
這是我追著要看極光的極光之旅.
去之前我跟主禱告.
我說我相信主你是大能的主.
你留一點光給我看.
我就會滿足了.
我留在北極圈裡面.
同一間屋裡面居住了四天.
我以為望出窗就會看到極光.
但無論極光指數是高的日子.
還是陰天的日子.
都讓我們看不到極光.
當我和同行的同伴.
覺得這個旅程應該看不到極光的時候.
其實心裡都想.
不要緊.
其實看到雪.
覺得那個景象已經很夢幻.
我們回程了.
坐火車走.
離開北極圈往南走的時候.
我們都睡著了.
凌晨一點鐘.
我驚醒.
看到我有個訊息.
有個朋友提醒我.
他說當時的極光指數很高.
現在你快點去看.
你是不是離開了北極圈.
我心想.
對 我已經離開北極圈了.
為什麼我要離開北極圈.
它才極光指數高.
我同行的朋友聽到我這樣說的時候.
他又望出窗.
誰知道真的看到極光.
原來.
哇.

$^{361}$Yoyo幫我.
我分享了給Yoyo.
她就讓大家看.
對了.
真的讓我看到極光.
這個是在車上拍出去.
用電話拍的.
其實眼望是沒有那麼光.
不過電話拍出來.
真的很奇妙.
好像彩虹一樣.
有一些綠色的光.
其實這個光.
是火車轉彎的時候.
讓我們看回北面.
看回北極那一邊.
讓我們可以.
看到大概十分鐘左右.
然後就慢慢消失了.
當我回想這件事的時候.
我覺得.
天父真的很神奇.
當我們的信心小的時候.
求天父留一點光給我們看.
他就真的給我們一點光看.
所以當求的時候.
就應該求見很多很多.
求大一點.
求多一點.
原本我也不明白當中的意義.
難道神留一點光給我.
他也做不到嗎.
但原來這件事讓我感受到.
無論我多努力.
有些事.
都不是靠我的努力就做得到的.
看極光.
都是全靠神的恩典.
神的做法往往是我們都意想不到的.
神要賜給我們的往往是出人意外的.

$^{401}$是我們意想不到的.
當眾人以為.
耶穌要騎著戰車.
戰馬去進城的時候.
耶穌卻是.
騎著驢.
謙卑地進城.
當大家手拿著.
鐘樹枝去呼喊的時候.
看似榮耀的背後.
原來大家都不知道自己在做什麼.
門徒也不明白當中的意義.
然而.
耶穌抗衡大家歡呼聲背後.
期望暴力的爭戰.
耶穌帶給大家的是和平和愛的訊息.
這種方式不會一下子改變世界.
卻是會讓所看見的人.
可以改變.
明白到可以改變世界的力量是何在.
我們的教會也只在社區當中.
平日也會開放讓大家去禱告.
也會有社區中心.
用一個教平教商的價錢.
去搞一些活動服務社區有需要的人.
或者在租金這麼貴的日子.
在賺錢艱難的日子.
有一天地產商也願意留下一角地皮.
去給窮人在那裡的茅廂.
感恩政府在城市規劃當中.
最近也留下一些空間.
讓窮人可以呼吸一些空氣.
未來在附近也會有簡約公屋和新田部屋.
去讓低收入的家庭入住其中.
讓我們也要預備好自己起來去服務社區的人.
或許會有很多新移民.
或者低收入的家庭.
去入住當中.
求主賜給我們有智慧.
怎樣將愛和福音帶給我們的鄰舍.

$^{441}$讓人可以看見有人騎著路子去經過.
就好像耶穌將愛和和平的訊息帶到當中.
我在這裡也預告一下.
在4月21日將會有暴行籌款.
為著這幾個新的屋村去禱告.
盼望大家都能夠參與在其中.
讓神開我們的心.
開我們的眼.
讓我們看見.
讓我們好像滿途一樣可以明白過來.
以前我們要很辛苦地去內地去傳福音.
現在卻是相反.
有一群有需要的人.
就入住在我們的中間.
傳福音反而可以成為我們平日的日常.
求主幫助我們可以把握時機.
和基督去同行.
和基督一同去受苦.
和基督一同去得榮耀.
讓我們一起祈禱.
親愛的天父.
因為主耶穌不是騎著戰車戰馬來到世上.
不是以暴力.
而是以愛和和平.
是謙卑的騎著驢驅在我們中間.
或許我們先前不明白.
然而求你開我們的心.
開我們的眼.
讓我們預備好與你同行.
縱使我們難以避免苦難.
但是我們相信你.
必定會拯救.
必定同在.
你有你的時間.
你有你最好的安排.
願錦繡堂的弟兄姊妹.
緊緊與你同行.
奉主耶穌基督的名求.
阿們.
阿蒙.

\newpage



\section{約翰福音 12:1-26}
\label{sec:94pJmtJo_JE}
\textbf{2024-03-24 早堂《先生,我們想見耶穌》約翰福音 12\_1-26 講員:李炳光牧師}
\newline
\newline
連結: \href{https://youtube.com/watch?v=94pJmtJo-JE}{\texttt{https://youtube.com/watch?v=94pJmtJo-JE}} ~~~~ 語音日期: 2024-03-26
\newline
\newline
\hyperref[sec:u5__eddKCBg]{\small{< < < PREV SERMON < < <}}
~
\hyperref[sec:index_chronic]{\small{[返順時目]}}
~
\hyperref[sec:index_scriptual]{\small{[返順卷目]}}
~
\hyperref[sec:NKxHUXgOsyY]{\small{> > > NEXT SERMON > > >}}
\newline
\newline
約翰福音 12:1-26
\newline
\begin{longtable}{cl}
\hline
\hline
章節 & 經文 (和合本修訂版)\\
\hline
12:1 & \begin{tabularx}{0.7\textwidth}{X} 逾越節前六天，耶穌來到伯大尼，就是他使拉撒路從死人中復活的地方。 \end{tabularx} \\ \\ \relax
12:2 & \begin{tabularx}{0.7\textwidth}{X} 有人在那裡為耶穌預備宴席；馬大伺候，拉撒路也在同耶穌坐席的人中間。 \end{tabularx} \\ \\ \relax
12:3 & \begin{tabularx}{0.7\textwidth}{X} 馬利亞拿著一斤極貴的純哪噠香膏，抹耶穌的腳，又用自己頭髮去擦，屋裡充滿了膏的香氣。 \end{tabularx} \\ \\ \relax
12:4 & \begin{tabularx}{0.7\textwidth}{X} 有一個門徒，就是那將要出賣耶穌的加略人猶大，說： \end{tabularx} \\ \\ \relax
12:5 & \begin{tabularx}{0.7\textwidth}{X} 「為甚麼不把這香膏賣三百個銀幣去賙濟窮人呢？」 \end{tabularx} \\ \\ \relax
12:6 & \begin{tabularx}{0.7\textwidth}{X} 他說這話，並不是關心窮人，而是因為他是個賊，又管錢囊，常偷取錢囊中所存的。 \end{tabularx} \\ \\ \relax
12:7 & \begin{tabularx}{0.7\textwidth}{X} 耶穌說：「由她吧！她這香膏本是為我的安葬之日留著的。 \end{tabularx} \\ \\ \relax
12:8 & \begin{tabularx}{0.7\textwidth}{X} 因為常有窮人和你們在一起，但是你們不常有我。」 \end{tabularx} \\ \\ \relax
12:9 & \begin{tabularx}{0.7\textwidth}{X} 有一大群猶太人知道耶穌在那裡，就來了，不但是為耶穌的緣故，也是要看耶穌使他從死人中復活的拉撒路。 \end{tabularx} \\ \\ \relax
12:10 & \begin{tabularx}{0.7\textwidth}{X} 於是眾祭司長商議連拉撒路也要殺了， \end{tabularx} \\ \\ \relax
12:11 & \begin{tabularx}{0.7\textwidth}{X} 因為有許多猶太人為了拉撒路的緣故，開始背離他們，信了耶穌。 \end{tabularx} \\ \\ \relax
12:12 & \begin{tabularx}{0.7\textwidth}{X} 第二天，有一大群上來過節的人聽見耶穌要來耶路撒冷， \end{tabularx} \\ \\ \relax
12:13 & \begin{tabularx}{0.7\textwidth}{X} 就拿著棕樹枝出去迎接他，喊著：「和散那，以色列的王！奉主名來的是應當稱頌的！」 \end{tabularx} \\ \\ \relax
12:14 & \begin{tabularx}{0.7\textwidth}{X} 耶穌找到了一匹驢駒，就騎上，如經上所記： \end{tabularx} \\ \\ \relax
12:15 & \begin{tabularx}{0.7\textwidth}{X} 「錫安的兒女啊，不要懼怕！看哪，你的王來了；他騎在驢駒上。」 \end{tabularx} \\ \\ \relax
12:16 & \begin{tabularx}{0.7\textwidth}{X} 門徒當初不明白這些事，等到耶穌得了榮耀後才想起這些話是指他寫的，並且人們果然對他做了這些事。 \end{tabularx} \\ \\ \relax
12:17 & \begin{tabularx}{0.7\textwidth}{X} 當耶穌呼喚拉撒路，使他從死人中復活出墳墓的時候，同耶穌在那裡的眾人就作見證。 \end{tabularx} \\ \\ \relax
12:18 & \begin{tabularx}{0.7\textwidth}{X} 眾人因聽見耶穌行了這神蹟，就去迎接他。 \end{tabularx} \\ \\ \relax
12:19 & \begin{tabularx}{0.7\textwidth}{X} 法利賽人彼此說：「你們看，你們一事無成，世人都隨著他去了。」 \end{tabularx} \\ \\ \relax
12:20 & \begin{tabularx}{0.7\textwidth}{X} 那時，上來過節禮拜的人中，有幾個希臘人。 \end{tabularx} \\ \\ \relax
12:21 & \begin{tabularx}{0.7\textwidth}{X} 他們來見加利利的伯賽大人腓力，請求他說：「先生，我們想見耶穌。」 \end{tabularx} \\ \\ \relax
12:22 & \begin{tabularx}{0.7\textwidth}{X} 腓力去告訴安得烈，然後安得烈同腓力去告訴耶穌。 \end{tabularx} \\ \\ \relax
12:23 & \begin{tabularx}{0.7\textwidth}{X} 耶穌回答他們說：「人子得榮耀的時候到了。 \end{tabularx} \\ \\ \relax
12:24 & \begin{tabularx}{0.7\textwidth}{X} 我實實在在地告訴你們，一粒麥子不落在地裡死了，仍舊是一粒；若是死了，就結出許多子粒來。 \end{tabularx} \\ \\ \relax
12:25 & \begin{tabularx}{0.7\textwidth}{X} 愛惜自己性命的，就喪失性命；那恨惡自己在這世上的性命的，要保全性命到永生。 \end{tabularx} \\ \\ \relax
12:26 & \begin{tabularx}{0.7\textwidth}{X} 若有人服事我，就當跟從我；我在哪裡，服事我的人也要在哪裡；若有人服事我，我父必尊重他。」 \end{tabularx} \\ \\
[1ex]
\hline
\hline
\end{longtable}
$^{1}$各位早晨.
非常高興今天在中樹節的時候.
跟大家一起敬拜.
當我來到這裡.
我為你們感恩.
因為第一件事.
看到這麼美麗的陽光.
看到這麼美麗的環境.
還有我進來之前.
看到這麼多人來敬拜.
但除此之外.
又聽到鳥語.
又聞到花香.
各位在這麼優美的環境.
敬拜神.
是不是神的恩典?.
真的很難得.
我希望大家今天在崇拜當中.
能夠感受神的同在.
能夠在崇拜的時候.
再次得著主耶穌.
好不好?.
我希望大家都能夠.
體驗到當年耶穌基督.
在中樹節的時候的遭遇.
和感受.
各位我常常說到.
人的生命.
本來是很平凡的.
但是因為接觸了耶穌.
就是不平凡了.
各位.
這些都是不平凡的信徒.
不平凡的主耶穌基督的兒女.
很多人都問一件事.
為什麼一個生長在.
二千年前的人.
發生這麼大的影響力?.
那撒勒的耶穌.
到底和他有什麼關係呢?.

$^{41}$我們常常提及他.
究竟有什麼關係呢?.
還有 究竟主耶穌.
是不是真的這麼重要?.
究竟主耶穌.
在這麼先進的時代當中.
還有沒有價值?.
他是不是真的影響了我們?.
我們沒有他可不可以呢?.
我希望大家今天思想一下這個問題.
但是有一句話.
是聖經說的.
我希望大家再次深入的體會一下.
除他以外 別無拯救.
因為天下人間.
沒有次下別的名可以靠著得救.
這句話直到現在.
是不是正確呢?.
美國著名的社會學家.
他曾經這樣說.
這個星球將來.
有關宗教最大的難題.
是有關耶穌本身和他的教訓.
提到宗教.
我們不能夠擺脫耶穌.
因為我們提到天文學.
正如我們提到天文學.
要擺脫這個Capitalist.
一樣 因為耶穌.
營活生生在人的心中.
英文是這樣寫的.
我不讀出來了.
就是耶穌基督.
他為什麼這麼重要.
這位學者 社會學者.
他用一個客觀的態度去看.
我們想起耶穌.
我們不能夠脫離信仰的問題.
同樣我們想起天文學家.
所提的一樣.

$^{81}$沒有他.
就像天文學也有些缺陷.
當年耶穌基督榮耀入京的時候.
其實中樹節.
是我們教會一個很重要的節日.
但是我們覺得.
很多時候不是太重視.
有時候我們認為.
只不過是一個普通的節日.
我們今天準備來的時候.
我看到我們的教會.
買了很多中樹枝回來.
每年中樹節的時候.
那些師班都拿著中樹枝.
一路像當年那樣.
一路說 一路唱.
這樣進來我們的教會.
如果你詳細研究中樹節當年的情況.
你會發現到.
當年的中樹節所發生的事.
實際上有很多東西值得我們思想.
我曾經說過一篇講道.
我說中樹節的眾生相.
中樹節.
如果你看當年在周圍的群眾當中.
每一個人都可以看到.
當年他對耶穌的反應.
有人覺得不外如是.
有人搖旗吶喊.
人云亦云.
有人五分鐘的熱度.
最初說和煞.
奉祂命來的.
是應當稱頌的.
但一個星期之後就說.
釘他上十字架.
你看看當年的群眾.
有很多東西值得我們思想.
有些人覺得.
又是一些嘈圈巴閉的事情發生.

$^{121}$無所謂.
有人覺得事不關己.
己不勞心.
其實在當年的群眾當中.
有很多都是我們社會人士的景象.
但很多時候我們讀聖經忽略了一件事.
原來我將四本福音書詳細的研究.
中樹節所發生的事之後.
我發現到有一個很特殊的現象.
不知道大家有沒有注意到.
原來大家一直辯論.
究竟這個是誰來的.
對耶穌的人就說.
你們要知道.
你們反對的人現在湧入城.
很失望.
但是在各種不同的反應.
或者各為鬼胎的情況下.
有一群人.
這群人不是猶太人.
這群人是外邦人.
竟然他們靜靜的找耶穌的門徒.
你猜說了什麼話呢.
他說先生.
我們想見耶穌.
先生我們想見耶穌.
奇怪了.
為什麼這麼多人在那裡搖旗吶喊.
很多人高聲歡呼.
很多人人魂亦魂.
但竟然有一群外邦人.
竟然說.
我們想見耶穌.
這個代表了什麼呢.
當年的人心渴望真理的熱誠.
亦代表了今天人類心靈的呼喚.
亦是更加.
今日世人向教會發出的質問.
他們想在教會當中看見主耶穌.
到底教會能否將耶穌基督表明出來呢.

$^{161}$他們求問教會有什麼答案呢.
最近有一群教會的領袖和我聊天.
他說現在香港的教會人數減少.
很多人都走了.
很多人移民.
我不知道貴教會有沒有影響到.
我的教會就有很大影響.
有很多移民到其他地方.
而且移民的很多都是教會的中堅份子.
我心裡很沉重.
教會在這樣的環境當中.
究竟我們用什麼反應呢.
我們是否默默無聲.
是否我們就這樣接受呢.
教會應該怎樣面對這樣的處境呢.
我們在這樣的環境當中.
我們怎樣見證主耶穌呢.
先生 我們想見耶穌.
反映出以下幾件事.
我想今天在中樹節的崇拜當中.
提醒一下大家.
我們彼此的勉勵.
我相信大家都相信.
教會唯一的根基是什麼呢.
是主耶穌基督.
沒有了耶穌就沒有今天的教會.
唯一的使命.
教會的使命.
就是見證主耶穌.
宣揚救世的福音.
今天教會的問題和弱點是什麼呢.
就是基督徒未曾真正的表彰基督.
今天很多人真的向教會來看看.
我看看你教會裡面做些什麼.
究竟教會在這個社會當中的影響力有多大呢.
還有 大家都知道.
英文說.
One of the witness of the church today.
is a Christian not making enough of Jesus.
not making enough of Jesus.

$^{201}$即是說.
我們未曾真正的被人感受到.
我們在教會當中見到主耶穌.
外邦人竟然在中樹節的時候問門徒.
我們想見一下耶穌.
的的確有很多人想見到主耶穌.
究竟他的真相是怎樣.
大家都知道.
世界各大宗教的特性是什麼呢.
是在乎宗教的教祖.
他的教祖的個性.
第一 印度教.
是宗一個理想.
孔教是宗一個傳統.
聖道教.
Saintism.
這個是日本的.
他們是忠於國家.
會教是忠於一個法典.
但是基督教呢.
是忠於一個人.
這個人是誰呢.
大家都知道.
這個話不是我說的.
是一個著名的神學家所說的.
其實我們對主耶穌的認識的資料.
並不是太多.
我忘記了有一句話是施布真說的.
施布真是一個著名的報道家.
他說每一節聖經.
都希望能夠與耶穌基督.
連在一起.
其實我們對主耶穌基督的認識的資料.
並不是太多.
只是藉著四本福音書.
略略知道他的生平.
四福音所描寫的主耶穌.
是以四種的見證.
四種的反應.
和四個人對耶穌基督的反應.

$^{241}$因為他們知道.
他是神的兒子.
大家讀四福音的時候.
你看見幾個作者所表達的.
方式是不同的.
馬太是表達耶穌是一個教師.
馬可表達耶穌是一個統治者.
一個克服者.
路加表達耶穌是一個醫治能力.
一個醫治人心靈的醫生.
約翰很簡單.
是表示耶穌基督是世界的光.
不少的書籍.
詳細地描寫.
不少的藝術品.
以他為主題.
各位你看一看.
世界上著名的藝術品.
著名的文章.
著名的書籍.
都和耶穌基督有關.
大家看到下面那幅畫.
達文西所畫的.
耶穌基督的最後晚餐.
其實這個也是他心裡想出來的.
最後的晚餐.
不是這樣擺設的.
去到耶路撒冷你會知道.
其實他們的錯位.
不是好像排排錯.
耶穌在中間.
門徒在旁邊.
原來耶穌是坐在角落.
而且他們坐的方法.
是沒有椅子坐的.
是坐在地上靠著.
這樣打側的.
我試過這樣坐.
其實我也很辛苦.
坐了一會兒.

$^{281}$我的耳骨會刺痛.
或者手會麻痺.
但是他們是這樣坐的.
猶太人這樣.
那麼.
彼得坐在哪裡呢.
寶罕坐在哪裡呢.
猶太人坐在哪裡呢.
其實達文西只是用他的意思表達出來.
我很喜歡收集耶穌基督的圖像.
你看看.
這些都是聖像.
這裡是.
很多人說.
這些叫icon.
即是說聖像.
聖像.
你不要看得這麼簡單.
原來研究聖像學的人說.
他畫出來的時候.
每一點每一眼的眼神.
都注意到的.
而且有很多東西說的.
我給大家稍為看看.
收集的圖像有很多.
這個是西乃山修道院的.
一個耶穌的像.
人家說是最古的.
不是.
四世紀的英國的.
一個在馬賽克的地板的耶穌像.
還有是土耳其的.
這間大教堂.
馬賽克的耶穌像.
但是這間禮拜堂.
曾經被不同的統治者.
用他的方法來改裝.
這間禮拜堂.
曾經做了回教堂.
現在.

$^{321}$這個新的總理.
或者總統.
本來我們在去看的時候.
就將那些被人遮蓋了的耶穌像.
拆乾淨看到了.
但是現在這個新的總理.
又要做回教的地方.
我們很難過.
這個飛虎畫.
有人說是最古舊的.
因為就在這間禮拜堂裡面.
拆掉那些被人塗了油之後.
就看得出來這幅像.
這個土耳其的.
索菲亞大教堂.
蘇菲亞大教堂.
現在就不知道改了什麼名字.
我不知道了.
還有耶穌登山變像的.
耶穌升天的油像.
耶穌和迦南婦人的油畫.
都是很有名的.
當然大家看到.
巴西理賽爾內盧的耶穌像.
是世界聞名的.
但是這個國家.
究竟對耶穌的信仰有多少呢?.
大家都心裡知道了.
我不想再提了.
耶穌在山上.
而且是做得最好.
最好看.
一伸出手來.
大家都感受到.
波蘭的聖母大教堂.
耶穌復活節的.
復活景象的模樣.
歷史上.
被我們證明到.
當祂來到世界.

$^{361}$整個世界因為祂的緣故而改變.
當人與祂接觸的時候.
人的生命是有改變的.
我相信大家都承認.
有誰接觸過耶穌.
他的生命是改變了.
在你讀聖經的時候.
你會發現到.
耶穌的影響力是這麼大.
各位請注意.
當年的希臘人.
來找菲力.
菲力找安德烈.
聯合引導他們見耶穌.
他們說.
他們在信徒的團計當中.
怎樣找到耶穌呢?.
原來這班外邦人.
竟然在信徒的團計當中.
想找耶穌.
即是說.
你想見到耶穌.
就在信徒的團計當中.
或者說.
是今天的教會.
教會究竟是一個什麼地方呢?.
教會是不是四面牆.
一個瓦頂.
一個建築物呢?.
其實.
如果我們詳細的思想.
教會的原文的意思.
就是一個團體.
不是一間建築物.
是你和我.
所以有些人批評教會.
是批評自己.
因為我們是教會的成員之一.
沒有你,沒有我.
就沒有今天的教會.

$^{401}$教會有人說.
是人的崇拜和見證.
表明耶穌基督在歷史當中的傳統.
有人說.
教會像家庭一樣.
父親,母親,兒女.
屈膝的禱告.
有沒有這樣的景象呢?.
我們家庭當中.
有沒有禱告的生活呢?.
我們家庭兒女.
屈膝一起敬拜耶穌.
當年的教會.
是在家庭當中發生的.
信徒的友誼,講談.
和宣教.
還有聖禮在那裡舉行.
平信徒日常生活.
一個信心的見證.
教會在一個不同的時代當中.
作出信心的見證.
是非常重要.
你看回教會歷史的時候.
當時教會的影響力是非常大.
還有.
我記得我讀神學的時候.
有一個美國的神學家.
在上課的時候問我們.
你可不可以說出.
教會是甚麼呢?.
當然用神學的角度來說.
有很多東西可以說.
但是我就盲拳打死老師傅.
我就說.
教會就是耶穌基督.
道成肉身的延續.
他看著我.
他沒有回答我.
也沒有說對與錯.
但定睛地看著我一會兒.

$^{441}$大家都靜了.
其實.
很多神學家.
用別的方法來看.
但是教會應該是甚麼呢?.
耶穌基督.
他道成肉身.
住在人群當中.
見證神的大愛.
這是耶穌基督.
這個世界上的使命.
今天我們的教會.
有沒有將耶穌基督.
道成肉身.
住在人群當中的經驗.
和實際的情況.
表現出社會當中呢?.
各位.
今天香港.
我們教會只是一個小群.
為甚麼呢?.
你看看.
如果統計一下.
香港基督徒的人數.
和香港人口的比例.
是差得太遠了.
而差距現在.
很多人提到基督徒.
只不過是小群.
在一個團體當中.
問有多少人信耶穌.
舉手的人.
是不多.
有很多人反對.
有很多人批評.
有很多人論斷.
有很多人對基督徒.
有一個負面的看法.
各位.
我記得.

$^{481}$有一次有一個小朋友.
來到我們禮拜堂的門口.
我們有沒有幫助.
看到耶穌呢?.
這個小朋友問了一句話.
甚麼呢?.
先生.
耶穌在不在裡面?.
我真的不知道怎樣回答.
我說是呀.
在裡面.
為甚麼我看不到呢?.
看不到的不是耶穌的本身.
是祂的兒女.
祂的門徒.
祂生活的表現.
各位.
我記得有一次.
有一個弟兄來找我聊天.
因為他和別人過不去.
有些衝突.
有些矛盾.
我和他分析.
我和他各自遙書.
和他一起祈禱.
他就說.
牧師多謝你.
我也要學習忍耐.
說完之後.
走出門口.
剛走出門口.
誰知樓上倒水.
倒正他.
他還沒離開禮拜堂門口.
他在門口破聲大罵.
甚麼粗話都說出來.
我聽到之後.
我好難過.
一方面認錯.
一方面表現出來.

$^{521}$不是耶穌基督的圖像.
真的很難過.
有一次.
有個信徒回來告訴我.
牧師.
我們教會有個傳道.
在巴士上和巴士司機吵架.
他甚麼粗話都說出來.
我說.
很難說.
一個人衝動起來.
是可能這樣說.
但衝動也沒有這樣說法.
的確是.
各位.
今天這個社會.
說粗話.
成為一種習慣.
如果.
特別在茶餐廳或巴士上.
你發現有些人.
三句一定要說一個字.
三個字還要多說一個字.
你試想.
我們基督徒.
能夠在社會中.
有甚麼好見證呢.
我和很多不同的信徒.
傾談關於報道的問題.
分享信仰的問題.
有人告訴我.
牧師.
大部份的反應都是.
你叫我信耶穌嗎.
你看看有沒有人信耶穌.
信成這樣的時候.
我真的不會信.
因為.
他信耶穌所表現出來的.
不是耶穌的榜樣.

$^{561}$有個信徒告訴我.
牧師.
我信了耶穌洗禮.
我本來很開心.
誰知我有一天回到家.
回到工場.
他們說了很多話.
我沒有回答.
他們又要我去不三不四的地方.
我就不去了.
他們說了我的綽號.
叫我耶穌佬.
我是否很不開心.
對呀.
叫我耶穌佬.
因為他們的表現.
好像耶穌一樣.
各位.
我不知道大家有沒有被人.
叫你信耶穌.
你信耶穌.
和他有什麼不同的地方呢.
有些人在不同的場合當中認為.
我們是做回平常的人.
禮拜一到禮拜六.
做回自己.
禮拜日就Holy Holy Holy.
你說這樣的時候.
我們怎樣能夠知道呢.
這個小朋友問一句話很深入.
先生.
耶穌是不是在裡面.
今天在門口有人問.
先生.
耶穌是不是在裡面.
沒錯.
這個答案我們值得深思.
不是一個耶穌像.
不是一個耶穌的形象.
不是在這裡找到耶穌.

$^{601}$的確.
是在你的身上.
在我的身上.
在我們的行動當中.
我們的工作當中.
被人家看到.
這個地方是神聖.
這個地方是屬於耶穌的.
這個小朋友問一句話.
是值得我們三思的.
請大家注意.
今天剛才.
上回剛才那幅相.
我想提一提.
那幅相大家很少見到.
耶穌笑的.
耶穌張開嘴巴笑的.
你看見耶穌的相都很嚴肅.
或者說很憂愁.
但是我有一次在美國三藩市.
一間聖堂.
這間聖堂.
我相信不是一間基督教的教會.
是天主教的教堂.
我進去裡面.
看到一幅就是這幅.
耶穌笑的像的形象.
其實.
我還說.
其實這個才是真正的耶穌的形象.
我們很少見到這幅相.
我收集了這幅相.
我曾經在這裡.
提出.
你有沒有見過耶穌笑的相.
他說看不到.
我拿出來看.
很多人問我拿.
結果我在講道的時候.
第二次.

$^{641}$我就將這幅相影印了出來.
今天我沒有影印.
對不起.
送給大家.
其實.
我們應該要常常有笑容.
讓人看見.
我們生命有活力.
給我們有平安有喜樂.
請大家注意.
我們想見耶穌.
但是主耶穌也想見我們.
因為他曾經說過.
當我從地上被舉上來的時候.
就吸引萬人來歸向我.
今天各位主內弟兄姊妹.
當我們高舉耶穌的時候.
就應該吸引萬人來歸向你.
但是我們有沒有高舉主耶穌呢?.
我們用耶穌基督向人介紹.
不是口述的介紹.
而是生命的介紹.
讓人感受到耶穌的真實呢?.
各位.
我們怎樣可以見證主耶穌呢?.
我們怎樣能夠領人.
來見到主耶穌呢?.
我覺得有以下幾個提議.
其實.
先生我們想見耶穌.
其實他想見到.
信徒生命的見證.
我們的生活.
和其他有什麼不同的地方沒有呢?.
我聽到一件消息.
我相信大家都曾經聽過.
就某教會.
我們就籌款.
打麻將比賽.
我覺得有很多方法.

$^{681}$為什麼要用打麻將呢?.
被人誤會了.
就是一個牧師說.
我們要與世人同行.
他們最喜歡打麻將.
就一起打麻將.
究竟這是否一個見證呢?.
我不敢說.
我不敢批評.
但是大家認為.
這是否一個見證呢?.
生命的見證呢?.
告訴大家一個很感人的見證.
這個見證.
我剛剛上星期提到過.
我退休了有二十年了.
去年開始已經有二十年了.
退休已經有二十年了.
沒錯.
去年我剛剛統計過.
我事奉主六十週年了.
就是說.
我今年幾歲.
大家可以想到.
我無謂再提了.
仍然使我有機會.
在不同的地方.
不同的教會.
不同的環境當中見證主.
我退休那一年.
弟兄姊妹說.
李牧師退休了.
你們洗禮快一點吧.
想找我和他洗禮.
結果過著聖誕節.
竟然間有八十多人洗禮.
我沒理由在崇拜當中.
和八十多人洗禮.
只過了一個多小時.
其他人就坐在那裡看.

$^{721}$太過辛苦了.
不如我在那一週下午.
舉行一個洗禮的崇拜.
就是說.
我特別為這八十多人洗禮.
並且請他的親友來.
想不到.
我那一週堂.
樓上樓下都坐滿了人.
原來那次的崇拜.
有很多有特殊的生命見證.
來參加洗禮.
我覺得.
無謂浪費機會.
就請幾個有特別象徵性的朋友.
出來作見證.
其中有一位朋友出來作見證.
可能大家都聽過.
這個賭徒的一個見證.
這個賭徒.
大家知道賭徒怎樣呢?.
什麼都忘記了.
大家記不記得.
最近有一個爸爸.
因為難道走投無路.
太太離開了他.
做一個賭徒的太太不容易.
因為他賭徒.
他欠人錢.
這個太太的家庭都受騷擾.
很慘.
這個做爸爸的竟然.
拖著她的兒子上天台.
第一件行動.
就把兒子扔在街上.
第二個行動.
就自己跳下去.
他就死了.
兒子就重傷.
我看這個故事.

$^{761}$很感動.
很傷心.
為什麼人這麼軟弱呢?.
為什麼自己被自己的軟弱.
被自己的嗜好影響了家庭.
影響了下一代.
你說小朋友長大之後.
心理是怎樣呢?.
知不知道?.
我希望小朋友好起來.
但是好起來.
他的心靈損傷.
是大到不得了.
我不是說這件事.
另外一件事.
有一個弟兄.
也是走投無路.
他的家庭受騷擾到不得了.
結果有一天.
他太太說要和他離婚.
沒辦法.
太太都不要了.
兒女都不要了.
他回到家的時候.
他太太把垃圾袋.
裝在他幾件內衣褲.
不讓他進門口.
推他出去.
就把垃圾袋推給他.
你走吧.
我們分開了.
這位弟兄.
拿著垃圾袋.
裡面滿了他的內衣褲.
去到哪裡呢?.
去到西櫳球場.
因為我禮拜堂.
是靠近西櫳球場.
坐了幾天.
想了很多事.

$^{801}$哭了很多場.
在那裡睡了很多晚.
想不通.
最後只有找到耶穌.
幸好.
有一個機構.
專門做戒毒的.
結果幫了他.
但是他怎樣找到耶穌.
就不是因為這個機構.
因為有一天他走過禮拜堂.
禮拜堂門口.
有一個電子板.
我們拍了一些東西出來.
歡迎大家來教會.
除了之外.
還有很多不同的經文.
這些經文.
都是法人心聖的.
誰知道這個弟兄走過.
我們禮拜堂門口.
看到犯勞苦擔重擔的人.
來到我們這裡.
可以得到安息.
他看了之後.
很奇怪.
我這間勞苦重擔.
不得了.
還要得到安息.
進去.
因為我重建禮拜堂的時候.
我特別在樓下.
建了一個祈禱室.
這個祈禱室是給人.
到禮拜堂可以祈禱的.
這個還可以反映出.
耶穌是不是在裡面見證你.
他進去裡面.
看到那個祈禱室.
有個十字架.

$^{841}$有很多聖經.
有很輕微的音樂.
很舒服.
他坐下.
他看到那本聖經.
插在座位前面.
他拿開來.
一揭開聖經.
剛剛揭開了.
《勞教福音》十五章.
一看就看到.
狼子回家的故事.
奇怪.
當我讀這個故事的時候.
忍不住眼淚.
哭了起來.
我同工.
我孫教師走過.
這樣就跟他傾談.
跟他分享.
結果幫助他脫離這個困擾.
當然剛才.
我幫他找了一個機構.
叫做「戒賭」.
「戒賭」不是簡單的事.
要有專人的幫助才行.
我不說「戒賭」的方法.
因為沒有時間了.
竟然他認識了耶穌.
他接受耶穌.
並且決定在過世的聖誕節洗禮.
我當然找他作見證.
我星期三有兩個講壇.
一個是講道的.
一個是做主席的.
我通常安排評論員.
在那邊作見證.
讀經等等.
他就出來作見證.
就說他的故事.

$^{881}$很感動人.
但說到這麼長時間.
突然間.
他說本來我離婚了.
但因為聖誕節.
我還沒有簽紙.
說到那時他突然停下來.
眼睛紅了.
看著那邊.
我聰明了.
我坐在那裡.
我說他太太來了.
看他洗禮.
我就走過去.
跟他說.
鄭大兄.
你太太是不是來了.
他點頭.
我就大聲說.
鄭太你今天是不是來了.
沒有人回答.
如果你在那裡.
請你舉手.
我叫了幾次.
然後在那邊.
我不是指你們.
不要誤會.
在那邊.
舉手.
我就說請你出來.
竟然他肯出來.
他站起來的時候.
全場拍掌.
站在那裡.
我就趁這個機會.
拉他們兩個一起.
我說鄭太.
今天你的先生在人家面前.
承認自己的罪.
願意跟隨耶穌.

$^{921}$願意改變.
你肯不肯接納他.
結果.
他太太沒有回答我.
看著她的丈夫.
又看著我.
一路流眼淚.
我就趁這個機會.
把他們兩個拉起來.
拉起來的時候.
他們兩個突然擁抱在一起.
這個景象多感人.
兩個都站起來.
我告訴大家.
我從未見過的場合.
差不多所有人都哭了.
見到兩夫婦.
由破碎而復合.
並且大家流出悔改的眼淚.
在這個時候.
你是不是感恩.
這個場合我們見到什麼.
我們見到什麼.
我們見到主耶穌.
是不是.
原來主耶穌能夠將一個罪人.
得救 悔改 更新.
生命改變.
所以.
如果有人說.
先生 我們想見耶穌.
各位我們教會有沒有這些見證呢.
有沒有這些見證.
信徒生命的見證.
就讓我們見到.
原來這些見證.
能夠影響別人.
各位.
有人說.
先生 我想見耶穌.

$^{961}$見到生命的改變.
那些人你說.
叫我不信多難啊 是不是.
還有.
信徒和教會的合一的見證.
大家有沒有見到教會的分裂.
有沒有見到教會的紛爭.
你現在怎樣見到耶穌呢.
有一次我去外國.
有位弟兄載我去一間禮拜堂.
誰知經過另外一間.
我見到.
前面有一塊紙牌.
後面一塊紙牌.
我說 咦 為什麼今天禮拜天.
是有一個罷工的行動呢.
原來紙牌寫著什麼呢.
我走過去看的時候.
寫著.
請不要入謝禮拜堂崇拜.
我說 為什麼這樣.
我們歡迎人來禮拜堂.
為什麼 原來呢.
這間禮拜堂.
大家彼此.
兩對人彼此鬥爭.
不想人進這間禮拜堂去別的.
嘩 怎麼會這樣呢.
當然有很多人的紛爭.
很多人的矛盾.
教會的紛爭.
是阻止人家見證耶穌.
信徒的悔改更新.
我做牧師這麼多年.
我呢.
是見到很多悔改的人.
生命的改變.
有人問我.
組織一個報道團.
喂 你說那麼多不厭嗎.

$^{1001}$我說 不是說道的問題.
而是我看見悔改更新的問題.
所以促使我不能夠不說.
老實說.
當我見到有人舉手決志.
願意跟隨耶穌的時候.
我內心的感動.
難以言喻.
各位弟兄姊妹.
我們幫助到人悔改更新.
信耶穌.
是不是一個教會的見證呢.
還有.
我見到教會的信徒生活的見證.
5月1號.
這個是以諾基金請我去說.
一些關於祈禱的見證的時候.
今年5月1號.
原來他們舉行.
百萬小時的祈禱.
為世界和平 人類的幸福.
為香港的前途.
嘩 這些多麼大的力量.
各位.
如果你知道教會的祈禱力量.
是多麼的重要.
我要說一些個人的見證.
我去年這段時間.
我突然間發現.
我的肺部有些事情.
三個醫生都說.
我將會.
要面對生命的終結的時候.
患上不治之症.
但是.
我太太.
在這段時間.
我要告訴教友.
她幫我作了一段文字.
大家說要祈禱.

$^{1041}$我只能將生命交給上帝.
三個醫生的決定.
就在這時候.
發出之後.
全教會為我祈禱.
有些我不認識的教會.
特別是國內的教會.
全教會都這麼認識 為我祈禱.
還有.
香港有很多.
認識你認識的.
或者不認識的.
一些著名的人物.
都打電話來問候我.
我覺得是很感動.
最難得的.
就是.
有一間國內的教會.
竟然.
我完全沒有去過那個地方.
全教會這麼誓祈禱.
結果.
我去見.
去醫院再一次檢查的時候.
去見醫生.
醫生.
竟然和我太太.
和另外一位姊妹.
陪在一起的時候.
醫生直接這樣說.
我在你的肺片上.
我看不到有癌細胞.
我睜開眼睛.
我還問.
那你看到什麼.
原來.
他們錯誤的判斷.
我看到不是醫生的錯誤.
而是神的大能.
是不是.

$^{1081}$各位.
我們教會在禱告當中的力量怎樣呢.
其實今天我們教會的祈禱會.
人數最少的.
是不是呢.
耶穌基督吩咐我們.
不要忽略了祈禱.
耶穌的生活表現.
聖經不斷提醒我們.
我們要不斷的禱告.
但我們有沒有這樣做呢.
我看到.
禱告的動力是何等的大.
各位.
怎樣可以令人相信耶穌.
我們教會是否禱告的團體呢.
聖經提醒我們.
有積極的行動.
耶穌說.
離開了我們.
不能做什麼.
當有人找耶穌.
想和他辯論關於.
耶穌的真實.
和是誰的手.
肥力不和他辯護.
不和他辯駁.
越辯駁越破壞.
他直接說.
你們看看.
請人來到教會當中.
讓人見到耶穌好不好呢.
今天現在.
我們有這麼多空位.
為什麼我們不請他來呢.
你得到耶穌這麼多恩典.
為什麼你不讓他來享受呢.
還有.
保羅.
他強調一件事.

$^{1121}$在人生這麼多的衝擊.
我知道所信的是誰.
這是見證.
當然.
在人生變相的時候.
他們見到很多不同的東西.
但是.
結果一句話.
觸動了我.
就是.
不見一人.
只見主耶穌.
各位尊敬的兄弟姐妹.
今天呢.
我們在這個.
中廚節的禮拜.
我想將這句話.
送給大家.
先生.
我想見耶穌.
你怎樣答覆呢.
這樣說.
你看看我吧.
你來教會看看吧.
你看看我們的生活吧.
可不可以這樣說呢.
今天.
我們有沒有預見主耶穌呢.
我們自己本身.
對耶穌認識了多少呢.
還有.
我們有沒有帶人來見主耶穌呢.
特別呢.
在這樣的環境當中.
香港真的需要大家.
努力祈禱.
讓人見到.
信仰的群體.
的表現.
被人家.

$^{1161}$不信是不行的.
好不好.
所以我希望呢.
各位弟兄姊妹.
今天在.
這個中廚主日.
中之主日.
裡面.
能夠感受到.
當年那些外邦人的問話.
提醒我們.
愛主更深.
我們同心祈禱吧.
親愛的主耶穌.
很多人想見你.
但是很多時候我們.
有意無意之間.
將你的形象.
遮蓋了.
很多時候我們失去見證.
我們愛主的心不夠.
人家在我們身上.
見不到主耶穌.
天父求你寬恕我們.
幫助我們的教會.
能夠活出主耶穌.
使更多人.
因為我們教會的生活.
來接受主耶穌.
天父啊.
我們不斷地報道.
不斷地工作.
為什麼這麼多人.
推卻不肯見耶穌呢?.
最大的原因.
就是可能他們.
見不到主耶穌.
主啊幫助我們.
讓我們今天在崇拜當中.
中之主日當中.

$^{1201}$這一班外邦人.
他們要見主耶穌.
讓我們關心到.
外面的群眾.
都能夠認識主耶穌.
多謝你.
今天提醒我們.
讓我們愛主更深.
更加願意活出主耶穌.
人的說話是講願.
願你的靈在我們生命當中動工.
奉基督聖名祈禱.
阿們.
\newpage



\section{士師記 16:4-31}
\label{sec:NKxHUXgOsyY}
\textbf{2024-03-24 午堂《勝過致命誘惑的秘訣》士師記 16\_4-31 講員:林誠信牧師}
\newline
\newline
連結: \href{https://youtube.com/watch?v=NKxHUXgOsyY}{\texttt{https://youtube.com/watch?v=NKxHUXgOsyY}} ~~~~ 語音日期: 2024-03-28
\newline
\newline
\hyperref[sec:94pJmtJo_JE]{\small{< < < PREV SERMON < < <}}
~
\hyperref[sec:index_chronic]{\small{[返順時目]}}
~
\hyperref[sec:index_scriptual]{\small{[返順卷目]}}
~
\hyperref[sec:xRsKG7u5A0E]{\small{> > > NEXT SERMON > > >}}
\newline
\newline
士師記 16:4-31
\newline
\begin{longtable}{cl}
\hline
\hline
章節 & 經文 (和合本修訂版)\\
\hline
16:4 & \begin{tabularx}{0.7\textwidth}{X} 這事以後，參孫在梭烈谷愛上了一個女子，名叫大利拉。 \end{tabularx} \\ \\ \relax
16:5 & \begin{tabularx}{0.7\textwidth}{X} 非利士人的領袖上去，到那女子那裡，對她說：「請你哄騙參孫，探出他為何有這麼大的力氣，以及我們要用甚麼方法才能勝他，將他捆綁制伏。我們就每人給你一千一百塊銀子。」 \end{tabularx} \\ \\ \relax
16:6 & \begin{tabularx}{0.7\textwidth}{X} 大利拉對參孫說：「請你告訴我，你為何有這麼大的力氣，要用甚麼方法才能捆綁制伏你。」 \end{tabularx} \\ \\ \relax
16:7 & \begin{tabularx}{0.7\textwidth}{X} 參孫對她說：「若用七條未乾的新繩子捆綁我，我就像平常人一樣軟弱。」 \end{tabularx} \\ \\ \relax
16:8 & \begin{tabularx}{0.7\textwidth}{X} 於是非利士人的領袖拿了七條未乾的新繩子來，交給她，她就用繩子捆綁參孫。 \end{tabularx} \\ \\ \relax
16:9 & \begin{tabularx}{0.7\textwidth}{X} 當時，埋伏的人正在她的內室等著。她對參孫說：「參孫，非利士人來捉你了！」參孫就掙斷繩子，繩子如遇到火的麻線斷裂一樣。這樣，人還是不知道他的力量從哪裡來。 \end{tabularx} \\ \\ \relax
16:10 & \begin{tabularx}{0.7\textwidth}{X} 大利拉對參孫說：「看哪，你欺騙我，對我說謊。現在請你告訴我，要用甚麼方法才能捆綁你。」 \end{tabularx} \\ \\ \relax
16:11 & \begin{tabularx}{0.7\textwidth}{X} 參孫對她說：「若用未曾用過的新繩子捆綁我，我就像平常人一樣軟弱。」 \end{tabularx} \\ \\ \relax
16:12 & \begin{tabularx}{0.7\textwidth}{X} 大利拉就用新繩子捆綁他，對他說：「參孫，非利士人來捉你了！」當時，埋伏的人在內室等著。參孫掙斷手臂上的繩子，如掙斷一條線一樣。 \end{tabularx} \\ \\ \relax
16:13 & \begin{tabularx}{0.7\textwidth}{X} 大利拉對參孫說：「你到現在還是欺騙我，對我說謊。請你告訴我，要用甚麼方法才能捆綁你。」參孫對她說：「只要用織布的線將我頭上的七條髮綹編織起來就可以了」。 \end{tabularx} \\ \\ \relax
16:14 & \begin{tabularx}{0.7\textwidth}{X} 於是大利拉用梭子將他的髮綹釘住，對他說：「參孫，非利士人來捉你了！」參孫從睡中醒來，將織布機上的梭子和織布的線一齊都拔出來了。 \end{tabularx} \\ \\ \relax
16:15 & \begin{tabularx}{0.7\textwidth}{X} 大利拉對參孫說：「你既不與我同心，怎麼能說『我愛你』呢？你這三次欺騙我，不告訴我，你為甚麼有這麼大的力氣。」 \end{tabularx} \\ \\ \relax
16:16 & \begin{tabularx}{0.7\textwidth}{X} 大利拉天天用話催逼他，糾纏他，他就心裡煩得要死， \end{tabularx} \\ \\ \relax
16:17 & \begin{tabularx}{0.7\textwidth}{X} 終於把心中的一切都告訴她。參孫對她說：「從來沒有人用剃刀剃我的頭，因為我一出母胎就歸給神作拿細耳人。若有人剃了我的頭髮，我的力氣就會離開我，我就像平常人一樣軟弱。」 \end{tabularx} \\ \\ \relax
16:18 & \begin{tabularx}{0.7\textwidth}{X} 大利拉見他說出了心中的一切，就派人去召非利士人的領袖，說：「請再上來一次，因為他已經說出了心中的一切。」於是非利士人的領袖手裡拿著銀子，上到她那裡。 \end{tabularx} \\ \\ \relax
16:19 & \begin{tabularx}{0.7\textwidth}{X} 大利拉哄參孫睡在她的膝上，叫一個人來剃掉參孫頭上的七條髮綹。於是大利拉開始制伏參孫，他的力氣就離開他了。 \end{tabularx} \\ \\ \relax
16:20 & \begin{tabularx}{0.7\textwidth}{X} 大利拉說：「參孫，非利士人來捉你了！」參孫從睡中醒來，說：「我要像前幾次一樣脫身而去。」他卻不知道耶和華已經離開他了。 \end{tabularx} \\ \\ \relax
16:21 & \begin{tabularx}{0.7\textwidth}{X} 非利士人逮住他，挖了他的眼睛，帶他下到迦薩，用銅鏈鎖住他，叫他在監獄裡推磨。 \end{tabularx} \\ \\ \relax
16:22 & \begin{tabularx}{0.7\textwidth}{X} 然而他的頭髮被剃以後，又開始長起來了。 \end{tabularx} \\ \\ \relax
16:23 & \begin{tabularx}{0.7\textwidth}{X} 非利士人的領袖聚集，要向他們的神明大袞獻大祭，並且慶祝，說：「我們的神明把我們的仇敵參孫交在我們手中了。」 \end{tabularx} \\ \\ \relax
16:24 & \begin{tabularx}{0.7\textwidth}{X} 眾人看見參孫，就讚美他們的神明說：「我們的神明把那毀壞我們的地、殺害我們許多人的仇敵交在我們手中了。」 \end{tabularx} \\ \\ \relax
16:25 & \begin{tabularx}{0.7\textwidth}{X} 他們心裡高興的時候，就說：「叫參孫來，逗我們歡樂。」於是他們把參孫從監獄裡提出來，在他們面前戲耍。他們叫他站在兩根柱子中間。 \end{tabularx} \\ \\ \relax
16:26 & \begin{tabularx}{0.7\textwidth}{X} 參孫對牽他手的僮僕說：「讓我摸摸支撐這廟宇的柱子，我要靠一靠。」 \end{tabularx} \\ \\ \relax
16:27 & \begin{tabularx}{0.7\textwidth}{X} 那時廟宇內充滿男女，非利士人的眾領袖也都在那裡，屋頂上約有三千男女觀看參孫逗他們歡樂。 \end{tabularx} \\ \\ \relax
16:28 & \begin{tabularx}{0.7\textwidth}{X} 參孫求告耶和華說：「主耶和華啊，求你眷念我。神啊，就這一次，求你賜給我力量，使我向非利士人報那挖我雙眼的仇。」 \end{tabularx} \\ \\ \relax
16:29 & \begin{tabularx}{0.7\textwidth}{X} 參孫抱住中間支撐廟宇的兩根柱子，左手抱一根，右手抱一根。 \end{tabularx} \\ \\ \relax
16:30 & \begin{tabularx}{0.7\textwidth}{X} 然後他說：「讓我與非利士人一起死吧！」他盡力彎腰，廟宇就倒塌了，壓住領袖和廟宇內的眾人。這樣，參孫死的時候所殺的人比活著所殺的還多。 \end{tabularx} \\ \\ \relax
16:31 & \begin{tabularx}{0.7\textwidth}{X} 他的兄弟和他父親的全家都下去收他的屍首，抬上去，葬在瑣拉和以實陶中間、他父親瑪挪亞的墳墓裡。參孫作以色列的士師二十年。 \end{tabularx} \\ \\
[1ex]
\hline
\hline
\end{longtable}
$^{1}$靜姐妹 早安.
今天我預備的這篇講章.
其實我都是最近在.
我想我們這個城市都有一些熱話.
有一些話題是提及關於.
就是一個大台的電視劇.
大家如果有.
我不知道有多少人有看過這套電視劇.
叫《婚後事》.
聽清楚 不是「身後事」.
是《婚後事》.
這個劇其實.
我看到那些編劇可能都用了一些心思.
用三個家庭的.
那些夫妻的問題.
甚至我們看到.
今天現代人的生活處境.
很多時候給了自己藉口.
就是來到在婚姻以外.
他們尋求一些所謂感情的滿足.
這是很危險的一件事.
這個危險的地方是什麼呢.
其中這套劇裡面.
我覺得它滲透了一個思想.
是和我們聖經的原則是很不同的.
據說婚姻出軌.
你不要真的去有肉體的關係.
這樣就OK了.
精神出軌沒有問題.
耶穌是不是這樣教導我們的.
耶穌說什麼.
你在街上見到一個婦女.
你動淫亂.
這就等於什麼.
就和她發生了淫亂的關係.
原來上帝看我們人.
祂知道我們有這種情感.
但上帝要我們去控制.
我們要好好去管理.
等於我們健康一樣.

$^{41}$上帝給我們有個身體.
是要我們好好去做管家管理.
上帝給我們有財富一樣.
我們要怎樣去理財.
但是如果我們來到是很隨便.
和亂七八糟的時候.
我們不用聖經神給我們的原則.
就一定會出很大很大的問題.
而這些問題在這套劇裡面.
上帝都帶出來.
不單止破壞夫妻的關係.
拆毀家庭甚至怎樣.
甚至連我們的兒女.
我們的下一代.
都會影響他們.
所以這件事帶來很重要.
今天我是用參宣他的人生.
給我們來看.
他這個典故.
我相信大家都會熟悉的.
在《士師記》.
而《士師記》也是讓我們看到.
神在起初《士師記》裡面.
這21章其實就是很多聖經學者.
都來和我們分析.
就是一個好像墮落的循環.
那三章裡面.
由第一章到第三章六節.
是說到第二代的領袖離世.
就是約書亞和一班長老.
他們是見證神怎樣帶領他們出埃及.
過紅海的.
但是他們要進入迦南地.
要得這個應許之地為業.
他們必須要聽神的話.
用信心他們才能夠去征戰.
以至能夠得著迦南地為業.
但是很可惜是因為以色列人離棄神.
而第三章七節到十六章三十一節.
我們就看到.

$^{81}$是有七個的事司.
他們是有不同的所謂的一個循環.
這個循環就是以色列人拜偶像離棄神.
神就透過仇敵去壓制他們.
懲罰他們.
後來以色列人他們痛苦不堪.
就來呼求神.
結果神也興起事司.
來拯救他們.
最後神容許他們的子民得享平安.
但是這個很可惜.
我們就看到人是不斷地重複墮落這些試探.
和忘記了神給他們的恩典.
而我們研究事司記.
他們這些循環式不斷地墮落.
其實有一樣很重要的.
就是他們為什麼會追隨迦南地的人去拜偶像呢?.
其實就是因為他們住進去之後.
就被迦南那七族.
他們來到是慢慢透過文化上.
讓他們來到是以如木掩.
日漸接受了他們那種價值觀.
我們人的軟弱墮落.
往往不是一夜之間你就墮落.
往往是一個過程.
你不要看輕.
但最微小的地方.
只要你稍一不慎.
你就慢慢擴闊試探的空間.
就越來越泥足深陷.
所以當我們來看回這個事司記.
特別我今天是專登將.
由十四章到十六章.
其實濃縮來跟大家說.
剛才多謝主席的弟兄.
幫我們讀了很長的一段經文.
這個就是說大利拉怎樣來引誘參孫.
最後是他真是致命的一個跌倒和失敗.
我們先再一次祈禱.
我們開始來看聖經.

$^{121}$神求你幫助.
叫我們今天打開你的聖言.
我們再一次來看見.
上帝的話語是發人深省的.
上帝亦都在這個末後的世代.
你要來警醒我們每一個.
讓我們知道我們今天.
我們不能夠自以為是.
甚至自高自大.
甚至我們以為我們自己被別人屬靈.
我們就站穩就一定不會跌倒.
主席其實我們今天沒有一個人可以自誇.
我們是一個屬靈的超人.
我們其實每一個人都陷在.
這個世界諸般的引誘罪惡當中的時候.
求主你幫助你的兒女們.
讓我們今天我們各人說的聽的.
我們都同德造就.
我們在神的面前.
我們願意謙卑自己.
領受你自己的憤悔.
再一次被提醒.
以致我們好好站在主.
你自己的話語和恩典裡面.
我們好好來跟隨你.
求主帶領我們下去的時間.
幫助說的聽的.
我們一同來得恩.
同心祈禱.
奉主名及.
阿們.
我們重溫一下參宣.
他做事司的一些片段.
在十四章一開始.
其實我們已經看見參宣那種墮落的.
是由他的眼目的情慾來開始的.
他在亭南這個地方.
他見到一個非利是的女子.
他就喜歡了人.
只是一個外表.

$^{161}$他都沒有跟人拍過拖.
只是在當中一見所謂鍾情.
然後就挑逗了他這個肉體的情慾.
我們都會看到.
是參宣堅持了.
他跟他父母說.
我一定要娶這個女子為妻.
他父母都說.
這個外邦的女子.
你不要與外邦人結婚.
我們需要找回我們自己的同族的人.
但是他說因為他喜悅她.
這個喜悅她.
在原文的聖經裡面.
其實都是有這種的.
是眼中看為對的一個這樣的意念.
新日本的翻譯就是說.
是因為參宣看中了她.
這個反映其實亦都是.
是士思記那個時期.
是整體神的子民的一個問題.
就是大家都是任意而行.
各人是行自己看為對的事.
這件事就像不像我們今天.
其實今天是那些異端邪說.
學術的講論.
今天我們講.
無論你講什麼理論.
今天是一個知識爆炸的時代.
甚至去到一個地步.
就是各說各話.
就說彼此尊重.
在基督教界裡面.
都有時會出現這樣的情況.
打開聖經.
竟然很明顯的神的話.
一些的原則.
竟然被人解釋成幾種不同的方式.
就好像你去茶餐廳.
你進去.

$^{201}$你看看午飯.
你說A餐 B餐 C餐.
悉隨尊便.
我們今天去跟隨神.
在這個彎曲薄茂的世道.
我們是不是這樣呢.
我想不是的.
我們今天需要去回聖經的真理.
我們查考神的道的時候.
我們要看究竟是我們自己的情慾.
是我們自己的驕傲.
然後我們找聖經去支持我們自己.
明明是罪的地方.
我們都可以解釋成不是罪的.
這些都是沒什麼的.
你不要太過緊張.
會不會呀.
有時甚至我們可能用文化去解釋.
但是我們看回神的話.
其實有幾樣東西是離不開的.
我們那種試探失敗.
是從眼目的情慾.
肉體的情慾.
還有什麼.
今生的驕傲.
這三樣東西都會不斷侵襲我們.
我們在參穿作時詩這個片段裡面.
我們看到在14到15章.
就是在說他從庭那裡墮落.
後來走到亞什基倫.
他怎樣去報復那些人欺控他.
搶奪了他所謂的未婚妻.
然後後來他又去到尼西的地方.
最後也去到加薩.
後來也去到希伯倫.
這片段我就不詳細在這裡說.
但是我只是和大家很快的回顧一下.
參穿那個復仇殺敵.
其實有一個特徵.
就是他以為他自己力大無窮.

$^{241}$他可以在庭那的地方.
赤手撕裂獅子如同山羊羔.
他以為他自己很本事.
他越來越發覺.
那個能力越來越是彰顯.
在亞什基倫擊殺30個非利士人.
到尼西這個地方.
他用盧衰骨殺了1000個非利士人.
如果大家看那個歷史的進程.
這些的時序都是重要.
讓我們看到就是.
原來我這麼厲害.
神的靈是與我同在.
其實神的靈大大感動參穿.
是為了去完成他作為事司的使命.
他作為事司的使命.
他應該要帶領神的子民以色列人.
同心合意去擊敗這些非利士人.
這些就是他們的敵人.
很可惜我們在當中就看到.
參穿不斷在這些事情裡面.
他看得很輕慢.
他自己作為一個拿西爾人.
他是有上帝作為一個選召的僕人.
所以當中我們看見.
這三個復仇的經過.
其實是神幫他的.
今天我們人很容易會驕傲.
你不要看教會.
我們看弟兄姊妹一樣.
你在商界在任何一個界別.
你是專業人士.
或者你是就算做一個廚師.
你很出色.
你能夠很有能力.
你很有成就.
我們記住.
永遠都要記住.
原來那些不是你厲害.
人家誇獎你的時候.

$^{281}$你真的要歸榮要給神.
你說這些都是神給我的恩典.
神給我的呼召.
是要透過神給我的恩賜.
能夠服侍到眾人.
是叫神的名得榮耀.
所以我們千萬不要搞錯了.
這個世界就由創世記.
始祖亞當夏娃.
食禁果的時候已經開始.
這三樣東西在約翰一書.
二章十六節都講得很清楚.
我們一起讀出來好嗎.
因為凡世界上的東西.
好比肉體的情慾.
眼目的情慾.
和今生的驕傲.
都不是從父來的.
而是從世界來的.
這三樣東西其實彼此是有關連的.
曾經有一個很出名的.
所謂靈修大師.
Richard Foster.
他多年前寫了一本書.
叫做Money, Sex, and Power.
他講金錢,性和權力.
這三樣東西自古至今.
原來在人類社會.
就是你說多厲害的人也好.
多有知識,才能的人也好.
甚至在教會界也好.
你久不久也會聽到.
為什麼一些很熟悉的領袖.
他都一敗塗地.
可能有些人就斂財.
可能有些人就是沒有來得謹慎.
就在男女的關係這些事情上.
出現了亂子.
亦都是墮落.
甚至是弄權,濫權.

$^{321}$這些權力往往就是吸引人的眼球.
當人坐在一個權力很大的位置裡面的時候.
如果我們不謹慎,不小心.
不清楚原來這一切.
都是神給我們的一個位份.
為的是要成就上帝託付給我們.
在這個位置裡面.
我們要完成祂的心意.
不是去完成我們自己的議程.
我們每一個人服事主.
都需要好好地提醒我們自己.
和我們彼此都要提醒.
這裡我特別用了一個.
在網上見到的.
我覺得都很幽默的.
他說The oldest computer was owned by Adam and Eve.
It was an apple with very limited memory.
Just one bite and everything crashed.
他暗示了.
他這個apple就是說那些apple的電腦.
其實是在說.
當然我們知道創世記的禁果.
不一定是蘋果.
不過在傳統上很多人覺得蘋果好像紅色.
很甜美 越人的眼目.
所以一般說的禁果就是蘋果.
但聖經沒有說是蘋果.
到現在我們都不知道是甚麼.
總之是知道它有那種特性.
所以這裡是借餘.
亞當 夏娃他們糟糕了.
明明神的話語是那麼清楚.
但他們就聽了那條蛇.
就是魔鬼的引誘.
是嗎 你吃了一定會死嗎.
不一定死.
神是怕你們好像上帝那麼有智慧.
結果夏娃吃完之後給亞當吃.
亞當就很驚訝.
你明明知道那些.

$^{361}$你就算不清楚也要問清楚.
他就是越人的眼目.
又很甜美.
他明明是違反神的心意.
我們來看參宣的致命誘惑裡面.
其實有三樣很關鍵性的東西.
這三樣的事.
第一樣 他是毫不謹慎.
你明明是拿世以仁.
你是分別為聖.
你要去專心去服侍神.
而且你是成為一個事司.
是帶著上帝一個很重要的歷史任務.
你要來幫助神的選民.
一起去擊敗那些非理士人.
但參宣就在這個過程裡面.
就是眼目的情慾.
他在加薩的地方見到妓女.
他就喜歡上了人家.
結果他就是這樣.
一步一步踏入這個圈套裡面.
我們在這個地圖裡面可以看到.
他慢慢由加薩這個地方.
我剛才那個地圖也給大家知道.
他由加薩走到塞列谷.
就是在第十六章我們會出現的.
我們會看到在那裡是一個很大的失敗.
在這個第一個墮落裡面.
首先我們要留意一件事.
為什麼說要謹慎呢.
就是意思我們要更加留心.
在希伯來書第二章說.
我們需要謹慎聽從神的基督所說的話.
越發鄭重.
這個謹慎留心是越發鄭重的意思.
我們不要輕忽輕看神給我們的真理.
如果我們不認真去對待神的話的時候.
其實我們就很容易越飄就越迷失.
其實今天我們在教會界裡面.
時不時我們做牧者會遇到一些人.

$^{401}$原來遇到一個新朋友來.
他不是誤導者.
他原來在這個世界兜兜轉轉幾十年.
他說他中學的時候他信過耶穌.
不過出來社會之後.
他就隨波逐流.
接受了世俗化的價值觀.
結果他就一直迷失離開教會.
但為什麼他現在會回來教會呢?.
他經過了很多年的漂流之後.
他不單覺得自己很迷失.
他覺得他自己人生沒有那種力量.
他很想回到神的家裡面.
重新去找到那種永恆的真正的愛和價值觀.
這件事對我們今天來說都是十分重要.
在參宣的墮落裡面.
參宣是做事司做多少年?.
剛才經文大一讀的時候.
你都看到20年了.
不是一個短的時間.
而他的墮落失敗就是一個漫長的過程.
但這個漫長的過程我們可以有跡可尋.
就是看到他是雜飛成仕.
就是他一直興奮神對他選照.
這種身份.
他沒有去謹守作為拿西爾人當守的誡命.
他一見到在庭蘭非利士的女子.
他就一見鍾情.
而他來到就是看見人家的美貌各方面.
然後他就想去跟人親近.
想做就去做.
這種是很危險的.
衝動沒有約束.
我們看到這是第一方面.
第二方面我們來看.
就是參宣他毫不警醒.
一連四次大利拉來跟他說.
要哄他要說出那個秘密.
剛才主席帶我們來看.
大利拉大家知不知道希伯來文的原意是什麼意思.

$^{441}$Dalila 原來舞台劇也有.
這個名字其實本身的意思.
意思是虛弱.
是軟弱無力.
原來一個的大利拉.
所謂他軟弱無力.
就是他道德上無能為力.
他自己是很隨便.
而是他這種女性.
當他一步一步來給非利士人利用他.
就是要找出沒有能力的參宣.
為什麼你這麼有能力這麼有力量.
竟然參宣一次兩次三次都頂得住.
到第四次他就真的投降了.
為什麼呢.
因為這個大利拉很厲害.
他收了別人的錢.
5500 赤霞.
其實就是相比當時.
一年的工資只是10個赤霞.
在當時.
5500 赤霞也有多少工資.
550年 天文數字.
如果我們今天來說.
我們的MPF都肯定沒有.
你想想.
如果你說他拿了這筆錢.
他可以提早收工.
他退休.
他都安枕無憂.
你就是看到大利拉.
他是有這種金錢上對他的誘惑.
他也都不介意.
就是要來引誘這個參宣.
說出他力氣的源頭.
結果我們看到.
他真是一步一步.
可以說是得寸進尺.
當時參宣他沒有警醒祈禱.
就算前文十五章都說他口渴的時候.

$^{481}$向神呼求 求神給他水喝.
結果神都是有供應他.
不過他當時不是真正的悔改.
其實參宣都是吃喝自己的罪.
我們看大利拉引誘參宣的過程.
有幾個字眼.
十四章第七節.
逼著他 天天是纏著他 求他.
十六章十六節怎樣.
天天用話催逼他.
這件事男人是最怕的.
不要這麼囉嗦.
不要一直纏著.
男人就爽快一點.
一日一日.
我都說了不告訴你.
他第二天又再來纏繞他.
一天是這樣.
為什麼他脫不了呢.
因為他太喜歡大利拉.
之前那些經文說到他在庭蘭.
喜歡妓女.
甚至去到塞列谷.
他又去找另外一個妓女.
這個參宣風流快活.
其實到處流程.
這個就是參宣.
但聖經只有一個女人.
是留名的.
就是大利拉.
而聖經在十六章甚至說.
是參宣愛這個女子.
他喜愛這個女子.
大利拉在十六章四節.
所以我們有合理的所謂推斷.
參宣為什麼最後屈服呢.
給一個他心裡.
他以為他心愛人家.
但人家是不是心愛他呢.
人家是愛著那五千五百的赤霞.

$^{521}$這個是風流場所的女子.
被非利士人利用.
結果陷這個參宣.
在一個很大致命的誘惑裡面.
大家知不知道北極熊是怎樣捕獵的.
原來北極熊如果和一般獵人.
要來搏鬥獵人一定輸的.
因為北極熊起身的時候.
一般來說都高過你和我.
然後他可以有幾百公斤.
那些是成年的北極熊.
所以很危險.
所以獵人很聰明.
他們就找一些魚.
在那些冰捉了上來之後.
然後就宰了牠.
不是吃.
宰了牠就讓那條魚很多血.
血腥那些魚血.
插在一把刀.
鋒利的刀插了下去.
但在刀柄就塗了這些魚的血.
然後他就走開了.
就等北極熊出現.
北極熊一聞到那些.
有美食.
有魚生吃.
一來到.
原來在刀柄上面.
好味.
馬上去舔.
不像吃那些魚生.
又吃那些血.
一邊吃的時候.
其實北極熊他不知道.
他越吃越好味.
那些血其實是戒傷了他自己的舌頭.
北極熊的舌頭.
所以這個最後.
他為什麼可以.

$^{561}$暴殺這些北極熊.
不費吹灰之力.
到北極熊就慢慢享受.
他這一種所謂魚生.
血腥的用膳之後.
他就昏倒在地上.
然後被這些獵人.
就很容易捉到.
完全好像大利拉引誘參孫一樣.
參孫最後.
他自以為自己力大無窮.
大利拉跟他說.
參孫.
非利是人.
要來捉你了.
這是第四次.
就是說參孫一從睡中醒來的時候.
好像前幾次一樣脫身而去.
他卻不知道.
耶和華已經離開他了.
我們作臣的兒女.
做到一個這樣的地步.
參孫應該是一個屬靈的領袖.
但竟然這麼昏迷.
這麼糊塗.
因為他驕傲.
驕傲在敗壞已先.
他內心那種高傲.
讓他跌倒.
甚至一敗塗地.
這個令人覺得很心傷.
最後我們很快的提一提.
參孫的失敗.
他毫不謹慎.
他是拿西爾人.
雖然你說拿西爾人可以有一個期限.
但他是與神立約的.
而更加重要的.
參孫不只是拿西爾人.
參孫在母胎的時候.

$^{601}$神已經揀選他.
神揀選他要分別為聖.
神要參孫是終身的.
來好好跟隨和事奉神.
很可惜他不謹慎他的身份.
他也不警醒.
他放縱他的眼目和肉體的情慾.
甚至他自己的驕傲.
以為自己力大無窮.
最後他成為那些非利士人.
將他帶到大菌廟中戲耍的對象.
這是在《時詩記》16章最後.
你看到這件事的情況.
真的好像那些溫水煮蛙一樣.
參孫跌倒的過程.
是到最後他自己死期臨近.
他也不知道.
就好像那些青蛙找盆溫水.
慢慢煮,慢慢加熱.
青蛙不會跑出來.
他以為多好.
慢慢以為在這裡捉菇.
誰知到最後被人家戮熟了也不知道.
我們今天人生有沒有都好像這樣.
我們是不是也會跌倒呢?.
在參孫死之前.
他的祈禱是這樣的.
他來求神幫助他.
我們剛才看那些非利士人捉了他的時候.
很慘的.
刮了他的眼睛出來.
然後鎖住他.
然後要他在監獄裡推磨.
然而他的頭髮被剃之後.
又慢慢長起來了.
其實神用這個拿西爾人留長髮.
而不去剪髮.
是做一個記號.
要參孫重新去醒悟.
他是一個分別為聖.

$^{641}$是要來歸神.
帶領神的子民去事奉.
很可惜.
參孫沒有謹慎他自己這個位份.
而他死的時候.
他的祈禱.
你發覺這個祈禱有什麼特色呢?.
這個祈禱不是主禱文.
這個祈禱是一個悲哀的人生的結局和呼求.
主耶穌教導我們祈禱是以神為中心的.
作為一個門徒應該怎麼祈禱.
而這個禱告參孫是怎麼求呢?.
他說求神給他有力量.
叫他可以向非利士人捕拿刮他雙眼的仇.
看不看到?.
不是以神為中心.
臨死的時候還是以自己.
我真的不值得.
我一生這麼有能力.
為什麼現在壞到被人鎖住.
還被人刮了雙眼.
我一定要一步所謂恥恨.
你看見參孫抱著中間支持廟宇的兩根柱子.
左右一抱的時候.
然後說讓我與非利士人一起死吧.
他一盡力彎腰.
整個廟宇就塌下來了.
然後同歸於盡.
整個廟宇有多少人呢?.
原來這些非利士人的領袖和信眾.
按照聖經的記載.
是有三千人.
然後聖經就說.
原來當時死的時候.
被活著所殺的還多.
不過這個很悲哀.
用這樣的方式來終結.
我們看見參孫.
他讓我們今天真的要引以為鑒.
司馬剛有一部很著名的著作.

$^{681}$叫做《資治通鑒》.
他曾經說到.
我們用人要用才德.
我們才能善用人.
這四類人是聖人,愚人,君子和小人.
才德全盡的叫做聖人.
才德甚麼都沒有的叫做愚蠢的人.
德勝過才就是君子.
而才勝過德.
即是有才剛但沒有德行.
這些叫做小人.
我想問問大家.
你覺得參孫是甚麼人?.
他是否聖人呢?.
當然不是了.
你說他是小人.
他是君子?.
也不是了.
他是愚人.
很肯定了.
真的很愚昧.
但他真的在當中.
他不懂得來跟隨神.
我們覺得這個是很可惜的.
他的失敗到最後的終結.
實在令人覺得很可惜.
而我們今天所謂用人為才.
我們要好好善用.
是講聰明.
而且不單止有聰慧.
要明察,要堅強,要有果毅.
這些都是人才的一些本質.
還有就是要正直,要公道.
這個忠就是公道.
和就是待人能夠有這種真正的忠和.
即是有親和力.
如果你今天在教會跟人服侍.
你很有才能.
但誰跟你一起合作都死定了.
我們就說這個真的很糟糕.

$^{721}$我們求主幫助我們.
最後我們再總結一下.
致命的誘惑.
你必須有一個謹慎的心.
這三個錦囊你一定要記住.
聖經真言書說我們要怎樣.
我們要保守我們的心.
勝過保守一切.
因為一生的果效是從心而發出的.
第一個就是謹慎的心.
因為魔鬼都好像爬開了獅子.
到處尋找可吞吃的人.
我們今天這個末世裡面.
其實真是魔鬼四處要我們跌倒.
而我們中國人都說這個色字頭上有什麼.
有把刀.
所以你不要來玩.
不要以為自己很剛強.
第二就是警醒.
警醒這件事.
我覺得今天我們都要好好學習一件事.
就是有時候你自己警醒不夠.
你要找一個小組.
找我們團契.
甚至如果傳道人可以跟你一起.
做一個守望者.
一定是有一個accountability group.
我們教會都很著重團契要小組.
所以為什麼要那些人不要只是來崇拜.
你加入團契要有小組.
小組就是弟兄姊妹彼此關心.
每個星期大家在團契.
無論什麼內容之後.
大家會有一些分組.
一起來彼此代討.
彼此代討不是叫我們家常便飯那些.
所謂飲飲食食那些.
是講你心底裡面一些的需要.
是要講出你自己最需要有人紀念你的東西.
所以我們今天說師友同行.

$^{761}$同輩同行都是重要的.
正如保羅提醒他自己的屬靈的兒子提摩泰.
就是要來逃避年輕人的私慾.
同那以純潔的心求告主的人追求公義信實.
仁愛和和平.
最後不單止有謹慎的心.
警醒的心.
還有一個是敬畏的心.
這個就是我們要專注為大.
今天我們能否是很好的來專注為大呢.
我們發覺參孫他自己真是一個獨行俠.
你在《士師記》裡面.
你看參孫他真是由頭到尾.
好像他自己one man band那樣.
一個人好像.
就是因為他自己自負.
他以為自己很叻很大力.
他不需要人幫他.
這個就是他的死證.
如果我們今天要別人跟你一起同工.
要好好跟人一起有好的合作.
你必須要是一個謙卑的人.
你也都需要是專神為大的人.
在這裡我給大家看一看.
他面對在被擄的當時生活的處境都是艱難的.
但是他和他的三個朋友就成為這個所謂accountability group.
他們一起可以去祈禱.
遇到任何危機的時候.
他們彼此待討.
所以他看見神是真正的掌權者.
神能夠超越我們人一切的軟弱.
包括所有錯誤愚昧的一些選擇.
今天大家來看看.
藥室他的遭遇都是.
舊約的藥室給他自己的親兄弟出賣他.
但是他竟然見到他們相認之後.
他說你不要害怕.
我豈能代替神呢.
雖然你們對我不好.
但是原來上帝要我來到去.

$^{801}$被賣到去埃及.
後來能夠升為宰相.
就是為了保存雅各我們的父親一家.
他們70人在大饑荒的時候可以去到埃及.
他們得以保存性命.
延續他們的後代.
這個就是神好奇妙的一個恩典.
在人諸般的軟弱失敗裡面.
最後最終極的仍然是神在掌權掌管.
我最後用一個見證結束今天這個講道.
我為什麼會用這一句說話來總結呢.
你不能夠選擇你在哪個人的家出生.
但你可以選擇你怎樣活出你的使命人生.
不知道大家同不同意.
其實我們今天沒有一個人可以選擇在哪裡.
有些人你不需要羨慕人家.
他含著金鎖匙出生.
一出生都不用憂了.
不做事他都可以過一世.
你很羨慕嗎這些.
我不會很羨慕.
其實最精彩的人生就是.
你如果能夠好好的照著神給過你的託付.
善用你的地上的每一天的光陰.
和神給你的恩賜.
你能夠貢獻社會,國家.
甚至在我們的社區裡面.
能夠去愛神愛人.
我覺得我們能夠不枉此生.
這個是最快樂的一件事.
能夠去恩著人在我們身上.
看見主耶穌基督而他受吸引.
而他願意去跟隨主.
這個是最幸福的一件事.
這個是關乎到永恆的.
我將我自己一個小小的經歷和大家分享.
我自己和師母有兩個乾兒子乾女兒.
乾女兒就是我們以前在加拿大牧會的時候.
她三歲就來到我們教會.
我們見到她.

$^{841}$師母因為教會言同主學.
一見鍾情.
師母覺得和她很有緣份.
後來我們去探訪.
亦了解原來他們是一個單親的家庭.
這個女孩子很可愛.
她是中美的一個混血兒.
不過很可惜她一出生的時候.
爸爸就離開了她.
直接不要她這個家庭.
所以是媽媽一手把她帶大.
當時我和師母因為我們沒有孩子.
沒有自己的小孩子.
我們就考慮認她做我們的乾女兒.
我們也可以填補一下.
這個小孩子成長的過程.
做父親的一種.
給一種家裡的一種愛.
我們也和我們的岳父岳母一起住.
所以這個女孩子.
你看見她很趣緻.
你見到她有一點像外國的樣貌.
但其實她媽媽是中國人.
所以眼睛是棕色的.
你就會發覺她其實很聰明.
她讀書各樣東西.
不過她成長的過程裡面.
她就開始去問.
為什麼別人有爸爸我沒有呢.
她到了她在少年人.
她真的經歷少年風暴的時候.
她很掙扎.
她自小在教會大.
她認識神的話.
不過她覺得為什麼.
好像是不是上帝對她不公平呢.
我和師母不斷用除了神的愛之外.
我們用我們的愛.
都讓她知道.
其實神對我們沒有任何虧欠.

$^{881}$是因為我們人的軟弱.
人的錯誤.
但是神今天不會因為我們人的失敗.
我說你今天有天上的父親.
有地上的契爺契公.
我們都同樣很疼你.
而你人生的價值.
就是你是無愛的.
你那個無愛.
是因為你是神的女兒.
你是有一個很尊貴榮耀的身份.
所以在她少年的風暴期.
她曾經有很反叛.
她和她媽媽頂到癢.
我和師母那時候已經回到香港事奉.
我們有時候用電話.
就跟他們調停一下.
跟他們談.
這張照片是她中學的畢業照.
她現在不單止是大學畢業.
那幾年的大學全部是獎學金.
一毫子都不用給 加學費.
她讀書讀得很好 很優秀.
但是這麼優秀.
她媽媽又很害怕.
她說怕她學壞.
我說你什麼都怕一頓的話.
你不行 我們還是交託給神.
把女兒交給神.
她大學畢業已經有三年.
她也做得很成功 做得很好.
她最初準備要選擇.
進大學讀什麼的時候.
我每次回加拿大.
我會和她吃一個長時間的早餐.
那次我還記得.
我和她談.
我問她喜歡讀什麼.
她說其實她媽媽想她讀醫生.
其實她有這樣的資質.

$^{921}$我們都相信她一定可以讀到.
不過她不是很喜歡讀醫.
我說你又不讀醫.
那你想怎樣呢.
你喜歡讀其他的專業嗎.
她媽媽就在律師樓工作.
她媽媽也很想她.
或者考慮.
她又不會.
她說律師.
她雖然有能力去讀.
但結果她選擇讀商科.
工商管理.
她進了加拿大最出名的大學.
叫Western University.
離開多倫多.
她走到去倫敦.
在安大略省另外一個城市.
就在那裡讀.
她讀完之後.
其實還未畢業之前.
已經有很多公司聘請她.
到後來她畢業的時候.
她其實答應了一間銀行.
但結果美國有另外一間.
更出名的大行.
要來聘請她.
她心動.
她說人家願意給她罰款.
要她毀約.
毀了合約.
你去那裡吧.
我跟她說.
我說 乾女兒.
其實你已經簽了約.
你也答應了人家.
雖然看來好像很吃虧.
那邊人工高一點.
那間大行說出名一點.
我說 你信得過神嗎.

$^{961}$如果神要你去一間.
比較沒有那麼出名的銀行去做.
其實也是一個訓練.
我說你不需要心急.
你照上帝原本給你的.
你就好好去享受神所供應.
結果她真的聽我說.
在人看來覺得很吃虧.
有什麼所謂呢.
人家都說賠賣那些罰款.
結果她真的在那兩年裡.
她很好的 有很多好的學習和鍛鍊.
然後她今天去到三藩市.
去一間更大的行裡.
她做得更成功.
所以我有時覺得.
不要看眼前任何的東西.
其實我就是用這些.
都是聖經的價值觀.
去感染我這個乾女兒.
我另一個乾兒子.
在香港也是三歲大起.
在教會認識的.
不知為何乾兒子乾女兒都是三歲.
就和我們結緣.
但是神有很多恩典.
這個乾兒子也很懂事.
他今年在英國一間大學讀完書.
他的心智就是做那些醫學研究.
他也讀得很好 很優秀.
他的成績全部都是超越的.
所以我們除了愛他們 關心他們.
一直其實有時我們說不要比較.
如果我們的兒女和人家的兒女.
你不要說你不是像隔壁B仔那麼厲害.
千萬不要這樣.
我們年輕人對每一代.
我們下一代的年輕人.
我們要鼓勵他們.
你去敬畏神.

$^{1001}$你去專心的謹慎和一個警醒的心.
我們人生的路.
其實是照著上帝的帶領.
和給過我們各樣的恩賜 才剛.
我們去服侍神 服侍人.
能夠榮神益人.
而且能夠帶領別人去認識我們的主耶穌.
其實我們的人生就不枉此生.
而且我們能夠真正的幸福快樂.
所以今天很多事情你不要那麼白閉.
我知道很多做家長.
很多時候是很為下一代.
會很擔心這樣那樣.
其實我們不需要.
我們都是和神一起同行.
不過有很多時候人家就會說.
其實這件事都是上帝.
牧師你在台上是這樣說.
或者上帝特別保守你.
祝福你.
你又不要這樣比較.
不要迷信.
我們今天就是.
我們每一個人的份.
其實在神裡面.
神已經量度給我們了.
願神幫助我們.
願神祝福 感受和每一個弟兄姊妹.
帶著神給我們的託付和使命.
我們活出祂的期望.
願一切榮耀頌讚歸於神.
我們一起祈禱.
天父上帝多謝你.
求你繼續施恩加力.
保守我們眾弟兄姊妹.
在這個彎曲勃謬的世代.
我們要保守我們的心.
不要真的被世俗的潮流.
一直帶到我們迷失.
甚至失敗的境況.

$^{1041}$主啊 我們真的要引以為鑒.
不要像參孫一樣.
願上帝在當中.
順領著我們每一個人.
求主 堅固我們的信心.
也讓我們持著永恆的盼望.
我們能夠帶著基督耶穌的愛.
去過地上每一天的生活.
能夠榮神益人.
同心合意祈禱.
奉主耶穌基督 聖靈求.
阿們.
\newpage



\section{帖撒羅尼迦前書}
\label{sec:xRsKG7u5A0E}
\textbf{2024-03-28 周四晚查經班《帖撒羅尼迦前書(2)》 講員:曾裔貴牧師}
\newline
\newline
連結: \href{https://youtube.com/watch?v=xRsKG7u5A0E}{\texttt{https://youtube.com/watch?v=xRsKG7u5A0E}} ~~~~ 語音日期: 2024-03-30
\newline
\newline
\hyperref[sec:NKxHUXgOsyY]{\small{< < < PREV SERMON < < <}}
~
\hyperref[sec:index_chronic]{\small{[返順時目]}}
~
\hyperref[sec:index_scriptual]{\small{[返順卷目]}}
~
\hyperref[sec:Su045rsIjTY]{\small{> > > NEXT SERMON > > >}}
\newline
\newline
$^{1}$我們一起祈禱.
多謝天父.
有一次我們可以來打開聖經.
去有一個系統的研讀.
一些新約的書卷.
貼上羅尼加前後書.
都讓我們看到.
保羅他作為一位牧者.
對教會.
一個遭患難的教會.
那個牧羊的掛心.
更加提到.
主的再來.
為什麼會不斷提到.
都是我們在這個課程裡面.
需要去研究的.
願神都幫助我們.
亦都給我們對於末世論.
可能我們遲些要講的時候.
對幾個理論我們都能夠掌握.
願神帶領我們.
在長的假期.
弟兄姊妹會有不同的考慮.
可能是休息.
可能是回鄉.
可能是數目.
或者有其他的安排.
盼望神自己.
都在我們大家的心靈裡面.
叫我們不忘亂.
在不同的生活.
都有主的同在.
我們這樣祈禱.
奉耶穌基督的名.
阿們.
上一個星期我們就開始講過.
這個教會.
保羅在貼上羅尼加.
其實就應該說在菲律賓.
他被害.

$^{41}$他坐牢.
被獄中的禁卒來打.
接著保羅和西拉.
在《少行傳》的十六章.
很詳細記載.
他們在深夜的時候.
普告唱詩歌.
在監獄裡面.
普告唱詩歌.
應該就不會有本詩集.
又不會有YouTube.
都是可能背誦一些詩篇.
那時候還沒有新歌等等.
他們都是可能用希伯來文.
因為西拉都是猶太人.
用希伯來文來唱頌一些舊約的詩篇.
唱詩讚美神.
保羅都教人的.
在聖靈的感動裡面.
第一樣東西就不是滾地.
又不是大叫大跳.
是心被恩感歌頌神.
但是他唱完詩歌之後.
地帶震動.
神要在監獄裡面顯一個神跡.
使監獄的根基都被打開了.
都被動搖了.
人們都沒事.
你想想如果在一個密室有地震.
日本就會經常有.
香港都很少.
你要立刻躲在桌底.
怕被壓死.
因為外面很大震動.
坐牢都不能走.
但那些獄卒和囚犯全部沒事.
門都打開了.
囚犯都不走.
很難的.
那些獄卒就很害怕.

$^{81}$就挑燈.
看什麼問題.
看到監門都開了.
囚犯都不在.
不知道去哪裡.
其實反正都死了.
因為你失職.
所以就想自殺.
保羅就知道這個情況.
就大聲喝止他們.
不需要我們所有人都在.
原來在聽保羅講道.
所以這個就是.
菲律賓的那個被害受辱.
後來就獄卒全家信主.
因為看到保羅和西拉.
這麼有信心.
又給他們很大的安慰.
就邀請保羅和西拉.
去他們家裡吃飯.
就帶了他們全家歸主.
即是獄卒的全家歸主.
那些囚犯有沒有信就不知道了.
到第二天.
那個官就要提保羅出監審問.
保羅就說我們是羅馬的公民.
你這樣來無緣無故將我們抓了審問是不對的.
不合法的.
非法拘留.
所以那個官都害怕.
所以就輾轉恐嚇.
就釋放了保羅和西拉.
他們回到呂底亞的家裡.
勸勉弟兄姊妹之後.
就離開.
就來了帖撒羅尼迦加.
在帖撒羅尼迦加這裡.
就有一個監察員.
在帖撒羅尼迦加這裡.
就在安息日進入會堂.

$^{121}$菲律賓就沒有會堂.
沒有猶太人的聚集點.
但是在帖撒羅尼迦加就有.
在帖撒羅尼迦加就向這些猶太人傳道.
但是猶太人就攻擊,詆毀.
結集一些市井的匪類.
即是黑社會.
來逼迫保羅.
就把他抓了坐牢.
去見官.
後來都輾轉.
那些人就叫保羅離開.
所以保羅就離開了帖撒羅尼迦加.
就一路輾轉去到經過比利亞.
跟著暗飛玻利之後.
就落到去瓦點.
在瓦點那裡逗留.
所以這本書.
帖撒羅尼迦加前書.
就是在瓦點那裡寫的.
就是保羅.
打發提摩泰.
可能沒有那麼張揚.
找個年青人上去教會.
避開那個風頭.
因為他們針對的是保羅和西拉.
所以保羅就打發提摩泰.
去帖撒羅尼迦加.
了解一下.
現在那個官方的鎮壓.
是不是越來越嚴峻.
還有最主要是教會的弟兄姊妹.
那個信心是不是動搖.
是不是因為一些患難.
就有後退.
甚至乎會改信.
不再信了.
那就是保羅所擔心的.
所以這本書嚴格就不是講主在來.
這本書很重要的寫作動機.

$^{161}$其實很單純.
就是保羅急速離開帖撒羅尼迦加.
有兩個人.
會不會那些初信者誤會了保羅.
一有難處就先離開.
就是.
放下教會不理.
自己走了.
有好處就來.
要沾那個光.
有困難就自己第一個走.
這些就不是好牧人.
好牧人是為羊捨命.
所以保羅就擔心這件事.
不僅僅是他自己個人的名義.
那個聲望.
而是擔心教會會不會在患難裡面.
受到動搖.
這個是比較重要的寫作動機.
剛才都提過了.
既然現在每一章的結尾.
記住他們的原文.
新約是不會有分段的.
是不會有分章的.
更加不會有格行等等.
如果有格行的話.
就不會有結尾.
因為他們的原文.
是沒有分章的.
等等的.
如果真正的原文.
是全部大寫希臘文字母的.
是根本不會有格行的.
是我們後人幫他們做分段的.
所以.
這個我們要留心.
為什麼每一個的轉接位.
都講到主在來.
講了五次.
後書還更加詳細.

$^{201}$這個是一個比較重要的主題.
但是整本書的寫作動機.
就很簡單.
一位的牧者.
他突然離開.
很擔心教會 初信的教會.
帖撒羅尼迦.
他們有沒有困難.
而令到他們失去信心.
失去對主耶穌的效忠.
跟隨.
記得帖撒羅尼迦.
本來就是一個很得到.
皇帝的任避.
一個Blessing的地方.
因為當年他們.
押注在奧古士都.
做王.
打贏了安東尼.
所以.
這個奧大維就成為了第一任.
羅馬的皇帝.
就是蓋薩.
所以你押中了.
有派彩.
你買六合彩.
你中了就有派彩.
所以帖撒羅尼迦城.
是有一個最惠國的待遇.
是有免的.
軍費.
還有那些駐軍.
在這裡證明.
這個地方一定有保護.
帝王的保護.
帖撒羅尼迦城有沒有印象.
是哪個皇帝的妹妹.
紀念他的妹妹.
而蓋的.
就是亞歷山大大帝.

$^{241}$你去到帖撒羅尼迦的海邊.
你會看到.
有一個紀念的銅像.
就是亞歷山大大帝.
騎著馬.
那個英姿.
是非常出名的.
所以是紀念他的妹妹.
亞歷山大大帝的妹妹.
來建立這座城.
但後來.
因為他們一直.
一直都是一個很獨到的.
港口.
地中海的沿岸.
是一個非常方便通商.
做貿易等等的地方.
所以.
亦都可以是四季宜人.
因為近地中海.
所以是一個很重要的地方.
但是我們有這些背景之後.
我們就要留意一下.
在第二章.
保羅都提到.
對君王那方面的.
不崇拜.
那個.
都是你可以說.
頗敏感的.
因為這個城本身.
就是有這麼大的.
官方.
和帝王有的關係.
關係是好的.
好的關係.
如果你一直傳說有一個耶穌是.
另一個王.
其實都頗大的事.
所以你看《四道行傳》的十七章.

$^{281}$他們的罪名就是說.
這個保羅叫人信.
另外一個耶穌就是王.
我們可以看回《四道行傳》十七章.
你對那些背景.
有些了解.
你就會大概想一下.
如果你是當地人.
當時你會不會有一個擔心.
如果我們真的堅信耶穌.
好像保羅所說的.
這樣信.
會不會受到官方的迫害.
或者官方的誤解.
以為我們在煽動一個.
拜領一個君王.
十七章.
《四道行傳》.
保羅和西拉.
離開了菲律賓.
經過了暗非波利.
阿波羅尼亞.
去到這個帖撒羅尼迦加.
在那裡有猶太人的會堂.
菲律賓明顯是沒有的.
所以保羅在安息日.
去到河邊.
見到一個婦女.
女底亞帶著一些婦女團.
在那裡祈禱.
在河邊祈禱.
因為根本沒有禮拜堂.
沒有猶太人的會堂.
他們是未信主的.
但是他們敬畏神.
這個女底亞.
所以保羅就去到.
他們的祈禱團裡面.
帶領他們信主.
和他們講解耶穌的救恩.

$^{321}$所以在帖撒羅尼迦加.
是很不同的.
所以他在這裡有猶太人的會堂.
第二節.
保羅就照他素常的規矩.
即是說他每到一個地方.
首先他要尊重有沒有猶太人.
因為猶太人應該.
先有機會聽福音.
然後信主.
他很希望讓猶太人.
得到神的救恩.
因為主耶穌自己.
也是由猶太那裡出來.
或者他自己.
也是一個猶太的拉比.
他很想讓他自己的同胞.
信主.
雖然神已經告訴他.
你是要奉差獻.
去到外邦.
成為外邦人的使徒.
但是保羅照常的規矩.
就是進入猶太人的會堂.
一年三個安息日.
本著聖經.
與他們辯論.
本著其實不是有一本大字聖經.
更加不會有一個.
IPAD 更加不會有電話.
本著聖經.
是保羅我不知道他有沒有.
手抄的卷軸.
舊約.
因為價值連城.
很昂貴的.
也不方便的.
如果有的話.
也不稀奇的.
因為在《小道行傳》第八章.

$^{361}$那個太監.
在車上.
敬拜完神之後.
就在車上.
回去埃塞俄比亞的時候.
就在馬車上.
來讀.
以賽亞書.
所以可能有些小小的流動.
可能是一些.
手抄本的.
迷你聖經.
一定是舊約.
沒有新約.
所以保羅本著聖經.
所謂本著.
要不你會背.
要不就很大的卷軸.
是很麻煩的.
本著聖經.
當然我應該補充的.
在會堂一定有手抄本的卷軸.
保羅自己沒有.
但進到會堂就有.
就有些卷軸.
女士對不起.
你要蒙頭.
要在正中間讀經的地方.
不能出入.
後面一層全部是男士.
即是中環半山的.
女士是在樓上的.
其實那個景觀更加好.
從頂部看下來.
正中.
好像打博士的擂台.
正中.
就是一個地方.
讓弟兄上去讀經.
所以保羅有身份.

$^{401}$總之是男士.
你露經也可以.
你讀經他也會邀請你讀經.
所以保羅是本著聖經.
但不是這樣讀.
是本著聖經和這些猶太人辯論.
即是他很熟悉舊約.
他隨便.
可以揭開哪一卷書卷.
來和他們討論.
關於彌賽亞.
所以你可以想像.
保羅經過聖靈的工作.
他自己的生命更新.
重新去研讀舊約的時候.
他有很大的亮光.
他重新.
透過主耶穌.
復活的耶穌再讀入舊約.
完全不同了.
那個理解也不同了.
等於我們信了耶穌.
有了聖靈.
我們再讀聖經其實已經不同了.
如果你是一個未信耶穌的.
中學生.
考Bible.
他真的當是一個故事來聽.
對他來說沒有甚麼感動.
更加不會有甚麼共鳴.
對神蹟等等.
他可能是依書直講.
為了考.
一來Bible都沒有甚麼人讀.
拿A是很容易的.
人們是想拿A才讀Bible.
變相.
他在考Bible.
他在考.
實際上.

$^{441}$連繫不到.
那個生命.
但本來就不同了.
他在未信耶穌之前.
已經很熟悉舊約.
有他自己一套的解釋.
有他們拉比一套的解說.
但是.
很牢固的改變不了.
但是耶穌的顯現.
聖靈進入他的內心.
和看到.
有很多的新生命.
例如在大馬士革的阿拿尼亞.
這個小弟兄.
或者其他在大馬士革.
認識的.
其他外邦人的信主.
令到他有很大的眼界開了.
原來這樣的人都可以信神.
得到天堂.
我們猶太人以為只有我們才有.
原來.
很多的禮儀是不需要的.
令到阿保羅.
由那一刻開始.
整個天堂為他重新打開.
或者真真正正.
開啟.
他越來越看到.
神在外邦的工作.
所以.
現在他已經是第二次的.
宣教旅程.
他更加認識.
更多的外邦的信徒.
尤其是年輕的信徒.
令到阿保羅.
越來越對.
他自己的民族的主義.

$^{481}$猶太人才可以信耶穌的.
對不起.
猶太人才可以進入神的國的.
那個謬誤.
徹底的打破了.
任何人只要在神的救恩裡面.
透過耶穌.
就可以得著這個永恆生命.
所以.
阿保羅在這裡看到.
他是本著聖經.
和他們辯論.
《小藤傳》的十七章.
講解陳明.
基督必須受害.
從死裡復活.
接著.
我所傳給你們的.
這一位耶穌.
就是基督.
這一位耶穌.
就是這一位在當年.
拿撒勒出生.
成長.
是木匠約瑟的兒子的這一位.
耶穌.
是不能推翻的.
這個人.
這個耶穌.
就是我們.
猶太人.
等待的救世主.
彌賽亞.
所以耶穌是基督.
是很重要的.
對猶太人來說.
耶穌只是耶穌.
耶穌只是木匠的兒子.
是出現過.
但是耶穌是基督是沒有辦法接受.

$^{521}$你講這句話的人.
禮貌地.
叫你快點離開這個會堂.
再嚴重一些.
就要把你逐出會堂.
甚至再極端一些.
就要像現在.
跳下梨家的人就是這樣.
要逼迫要抓保羅坐牢.
就是因為.
你說這位耶穌是基督.
這個就是天大的冒犯.
對於猶太人來說.
但是.
保羅沒有辦法.
因為他已經.
他的心靈的眼睛已經開了.
他已經看到真理.
看到神的.
真實的救恩.
所以.
他必須要這樣說.
耶穌是死而復活的.
彌賽亞.
基督救世主.
接著第四節.
黑門中間有些人.
聽了勸.
服從保羅和西拉.
所以保羅的陳詞.
不是一種很嚴厲的.
是勸說的.
讓聽的人有機會.
去反思去思考的.
這個耶穌.
會帶來這樣的.
一個生命力.
如果保羅這樣說的時候.
這裡是沒有用一些神跡去.
印證保羅所說的道.

$^{561}$所以不是每一次.
都需要神跡的.
每一次的講道.
都要加一些.
好像是附加的條件.
或者福利來叫人.
信要因為神跡.
是不需要的.
保羅的勸說已經有人聽勸.
所以就服從.
保羅和西拉.
並有許多的乾淨的.
希利尼人.
尊貴的婦女都不少的.
當然有人信.
也有.
有些人是有某些.
某些特性.
有些人是.
有某些特性.
有些人是.
有某些特性.
有些人是.
有某些特性.
有些人是.
有某些特性.
有些人是.
有某些特性.
有些人是.
有某些特性.
有些人是.
有某些特性.
有些人是.
有某些特性.
有些人是.
有某些特性.
有些人是.
有某些特性.
有些人是.
有某些特性.

$^{601}$有些人是.
有某些特性.
有些人是.
有某些特性.
有些人是.
有某些特性.
有些人是.
有某些特性.
有些人是.
有某些特性.
有些人是.
有某些特性.
有些人是.
有某些特性.
有些人是.
有某些特性.
有些人是.
有某些特性.
有些人是.
有某些特性.
有些人是.
有某些特性.
有些人是.
有某些特性.
有些人是.
有某些特性.
有些人是.
有某些特性.
有些人是.
有某些特性.
有些人是.
有某些特性.
也到這裡來了.
耶孫收留他們.
這些人都違背.
該殺的命令.
該殺甚麼命令呢?.
當然是全世界只有一個該殺.
當時如果沒有搞錯.
要不就是提比留.

$^{641}$或者是已經去到.
這個尼祿了.
可能是在兩個之間.
到底是該殺的提比留.
還是該殺的.
這個.
這個尼祿.
所以.
你就要看看.
哪個王.
如果是尼祿王.
就是極之針對猶太人.
但其實提比留都不是好人.
對不起.
提比留之後還有另外一個王.
就是這個.
格魯雕.
第四任的羅馬皇帝.
十八章就有了.
這個Goddess 格魯雕.
提比留之後.
其實有一個叫加里居拉的.
就是凱撒大帝.
接著就去到這個格魯雕.
接著就是尼祿.
所以到底這個時間是不是.
似乎十八章這裡說的是格魯雕.
那就是說現在這個時候.
那些人說的該殺的命令.
該殺是那個王的稱呼.
實際該殺就是.
羅馬皇帝的尊稱.
就叫做該殺.
所以現在這個就是.
格魯雕.
這個皇帝.
在保羅的任內.
接著就到尼祿王.
尼祿王就是.
根據羅馬古事.

$^{681}$尼祿王就是燒了舊的羅馬城.
他要建一個新的羅馬城.
就嫁禍給基督徒.
然後就把這個羅馬城.
嫁禍給基督徒.
就大規模的迫害.
捉拿基督徒.
坐牢.
就是這個尼祿.
所以較為明顯.
較為大規模的.
對基督徒的迫害.
就由尼祿開始.
但其他之前的王一點都不友善.
因為他們不斷.
聲稱自己就是.
可以和神去親近.
所以這裡這樣的一個.
政治的.
控訴其實是.
頗嚴峻的.
這些人違背該實的命令.
就說有另一個王耶穌.
所以你想想.
對帖撒羅尼迦城的人.
本來享有.
這麼高的一個.
最惠國待遇.
又有那些軍費的免差費.
差役費.
所以一定是和國家.
雖然帖撒羅尼迦.
離開羅馬很遠.
但其實坐船.
或者他們行軍都很快的.
所以如果那些消息.
傳到羅馬.
其實是很嚴重的.
所以他們很厲害.
一想一說.

$^{721}$就是最高.
最高的顛覆政權罪.
你可以想像得到.
對於當時的基督徒.
是很懼怕的.
所以保羅這樣.
突然間的離開.
某程度是減輕.
那個控訴.
不要令到教會.
受很嚴重的虧損.
但是又會帶來.
另一個角度就是.
有些人就會扭轉.
這件事來說.
他不是使徒.
或者就算他是使徒.
他見到困難就扔下教會.
不理了.
好像撇下羊逃走了.
故宮逃走.
因為他是故宮.
不愛惜羊群.
好牧人.
是為羊捨命的.
這是主耶穌給予的標準.
一個近乎是.
一個好的牧者的標準.
我們了解了.
這些背景.
我們再讀.
帖塞羅尼加前書.
有些深度.
有些對當時.
當地人.
他們的想法.
他們的擔心.
又會明白多一點.
真的.
四個星期.

$^{761}$一個月突然間就不見了.
而且周圍的.
逼迫可能他們.
有些通緝令通緝保羅西拉.
說他煽動.
帖塞羅尼加人.
要拜耶穌不要拜蓋撒.
那就百詞莫辯.
你剛剛信了主.
你就突然間.
要面對兩難的抉擇.
那應該怎樣呢.
牧者又不在.
牧者已經走了.
已經離開了帖塞羅尼加城.
其實是一個很大的懸案.
而且.
都是會對.
當時的信徒.
有一定的困惑.
所以這本書.
應運而生.
保羅當他去到雅典之後.
他不是要補妃.
他是很擔心.
教會現在的情況.
怎樣呢.
相信擔心的不是教產.
而是那班信徒.
的信心.
是否堅固.
如果堅固他就不擔心了.
因為神自己會保守.
這一班的.
帖塞羅尼加教會的弟兄姊妹.
是能夠一路面對困難.
都成長了.
所以這本書很簡單的寫作動機.
就是一位好牧人.
愛惜他的羊群.

$^{801}$看他們是神的兒女.
很擔心他們的情況.
所以就.
虛寒問暖.
不過他不可以自己親自回去.
就透過提摩太回去.
其實當時保羅.
在雅典已經看到.
通街都是那些偶像.
他就已經很焦急.
記不記得在雅典.
在巴特隆神廟.
下來之後.
有一個叫阿略巴古的地方.
有很多人.
拜那些偶像.
其中有一個偶像叫.
微色之神.
那個地方一個小的山丘.
就是那個阿略巴古.
我們去過了.
去年11月的時候.
這告訴我們.
貼前這裡我們要有.
一個心理準備.
保羅其實要幫忙.
解說給弟兄姊妹知道.
患難.
和那個.
信主.
其實是不能夠分得開的.
尤其是.
在一個患難的時期.
時代.
那個信主.
就要預謀這個.
計算.
就不是必然.
好像沒有困難.
所以在這本書裡面.

$^{841}$你會看到有幾樣東西.
包括就是.
牧者對群羊的擔心.
再加上就是.
這本書講很多苦難.
怎樣面對苦難.
再加上就是.
很特別就是每一章的結尾.
都有.
主在來的.
安慰.
應許.
初信的教會.
是否保羅當時.
這個提出了其中一個.
的見解.
保羅是否在這個時代.
他那個時候.
他這個傳道.
第二次.
宣教的旅程.
他很強烈.
感受到.
主要再回來了.
很快回來了.
所以他就在他的.
作品裡面.
不斷去講.
主的回來.
尤其是一個患難的教會.
就更加需要那個盼望.
去講主的再來.
所以我們讀到.
這些經文的時候.
我們又可以大家一齊去思想一下.
所以.
我們回來第一章這裡的.
帖索羅尼加前書.
開始已經講了.
一方面就是對這個教會的.

$^{881}$感恩.
還有就是那個的信忘愛.
大家有沒有回去.
組織一下新約.
出現過關於保羅講的.
信忘愛.
是有不同的pattern.
不是一律的.
是有不同的描述.
即是說原來是很多樣化的.
即是說保羅講的.
信忘愛.
不是只是好像千篇一律的.
一個格言.
又或者是.
好像是不能改變的.
原來不是的.
譬如在這個歌羅西書.
即是貼後一點.
貼回前一卷.
歌羅西書.
著重的.
反而是講那個盼望.
第三節.
第一章的第三節.
這裡講的信忘愛.
和貼上來講的完全不同.
我們主耶穌基督的父.
上上為你們禱告.
接著.
一聽見你們在基督耶穌裡面的信心.
並向眾聖徒的愛心.
為甚麼呢.
在基督耶穌裡面的信心.
對眾聖徒的愛心.
原來呢.
原來就是因為.
是為了給你們存在.
天上的盼望.
這個盼望就是你們從前.

$^{921}$在福音真理上.
真理的道上面.
所聽見的.
原來就是福音帶來盼望.
盼望帶來對基督耶穌有信心.
對眾聖徒有愛心.
這樣講信忘愛的.
和貼上來講是不同的.
有沒有留意到.
所以這個歌羅西書.
特別強調盼望.
這個盼望是怎樣來的.
就是聽見.
這個福音.
福音的真理.
的道上面所聽見.
這個福音傳到你們那裡.
亦都傳到普天之下.
結果增長.
好像在你們中間一樣.
所以.
原來信心.
愛心在保羅的理解.
是因為聽福音.
聽明而信福音.
而帶來對.
屬天有盼望.
當你有了這個屬天盼望.
這個盼望原來會生出.
信心對神的信心.
當然是對耶穌基督.
對別人.
對眾人的愛心.
是由盼望而來的.
所以.
大家如果真的有時間.
你回去做一做.
關於保羅的信忘愛.
例如.
哥林多前書十三章.

$^{961}$大家很熟悉的.
如今常傳的有信有望有愛.
其中最大的是愛.
那你又看看保羅怎樣在臨前.
去理解那個信忘愛.
現在貼上來.
我們有機會去看看.
保羅怎樣用信忘愛.
第三節.
貼上來的前書.
在神我們的父面前.
不住的紀念.
你們因信心.
所作的功夫.
這次提到不是對耶穌的信心.
是因為信了耶穌.
而有的工作.
即是那些工作.
因愛心所受的勞苦.
愛心.
這個對神還是對人呢?.
似乎是對神.
愛心所受的勞苦.
即是說.
在一個愛裡面.
不會覺得很難.
很苦.
這個是要有愛.
所以原來.
事奉就不是.
洗禮.
你幫我派程序表就行了.
因為派程序表.
甚麼人都有.
可能那些來聚會的人會令到你難堪.
那你就開始.
因愛心.
被踐踏.
那你就會.
感覺那個.

$^{1001}$辛苦.
那你慢慢就會甚至乎.
可能會最容易的.
因為你是義工.
義工就說明是義工.
我一定是隨時.
第一時間就可以不做.
那你一不做.
其實就因為你的愛心.
受到少少的踐踏.
所以很多時候我們不是太明白.
背後的原因.
以為義工其實誰都可以.
其實就變成了.
這樣的一個情況.
但這裡保羅告訴我們.
貼上羅尼加的教會.
一信耶穌.
他們因為信而開始.
有一個工作.
就是他們有為神而做事的工作.
因愛心.
他們願意忍受.
這裡是可以忍受那些勞苦.
跟著.
盼望主耶穌.
而全忍耐.
在地上.
他們已經知道.
在地上我們要學習忍耐.
就是說.
那些環境.
有時我們可以去.
駕馭.
但是更多.
未必是我們可以.
操控駕馭的.
那就盼望.
忍耐了那個力量.
可以繼續忍耐.

$^{1041}$很多人都會.
跟你所期望的.
不對現.
你就會失望.
跟你所等候的出現.
不是那回事.
你就會絕望.
很容易就有這個落差.
但是因為盼望.
主耶穌.
而心存忍耐.
這樣理解信望外.
這裡就很實在了.
一個信徒在信了主之後.
應該有功夫.
應該有.
愛而來的.
那個勞苦.
應該有盼望帶來的.
內心的忍耐.
就是說原來愛和忍耐.
是息息相關的.
你沒有忍耐.
原來你對主耶穌的愛不夠深.
明天晚上先告白.
又受苦之崇拜.
哼哼哼.
最好讓你流眼淚.
沒有人看到.
領聖餐.
你就要去思想一下.
耶穌為什麼這麼愛你.
為你釘十字架.
就是說你要加油.
你要加.
耶穌的愛的燃料.
在你的心裡.
你才會有忍耐.
所以.
對的義工不做.

$^{1081}$太容易了.
太容易放棄.
其實你背後不知道.
其實是.
耶穌的愛這個燃料.
你的內心已經越來越少.
當你有耶穌的愛的燃料.
你就能夠忍耐.
所以保羅說的這裡.
是對這個受苦的教會.
一個很貼切的描述和稱讚.
他們很快就已經因為信主.
而來到有功夫.
愛心.
而來到有勞苦.
然後盼望主耶穌.
而心存忍耐.
這裡就是.
一個內心很實淨.
準備.
去服務.
去事奉.
不然.
教你怎樣站位.
教你施事一和施事二.
是做什麼的.
功能性.
重要的是心靈上.
我們要有這種裝備.
接著.
第四節就開始說.
這一班弟兄姊妹.
我們上個星期已經說了.
所以不再詳細說揀選.
一定要很清楚.
在苦難裡面.
什麼令我們.
能夠過得了.
我們的困難.
那個困境.

$^{1121}$我們是被揀選.
為什麼會被揀選呢.
其實是沒得解的.
因為是神在萬世以前揀選我們.
還沒有這個地球就已經揀選你.
我們怎樣能夠明白呢.
還沒有這個地球.
就是說我在1963年出生的.
1950年代.
根本從來都不知道有這個人會出現.
是不是.
怎麼會有人知道.
我將來會變成今天的曾睿桂呢.
是沒有人知道的.
包括我爸爸都不知道的.
我爺爺都不會知道的.
但是萬世以前神已經揀選你和我.
只能夠是.
被神所愛的弟兄.
只能夠是.
原來萬世以前.
被揀選的人.
是將來.
很多世代之後才信主.
才出現在這個世界的.
我們.
能夠明白福音.
信主.
但原來唯一的原因就是.
被神在萬世以前已經愛了.
已經點名.
我們將來就會出現在這個世界.
這些是很玄妙的.
其實嚴格來說比佛教更加玄妙.
不用說什麼輪迴.
我們在萬世以前.
神已經揀選我們.
就是因為愛.
保羅現在就告訴我們.
好像代表神告訴我們.

$^{1161}$被神所愛的弟兄.
我知道你們是蒙揀選的.
是蒙神的揀選的.
蒙揀選的人.
會有一個表現的.
就是他們會聽明白福音.
信主.
所以因為我們的福音傳到你們那裡.
不再讀乎言語.
也都在乎權能.
in power.
和in spirit.
在神的聖靈的感動下.
就是說那些說話.
會打動你的心靈.
令你明白真理.
明白福音和你有關.
所以原來.
每一次的報道會.
有時候都是每一次的撒種.
總有些人.
你不知道為什麼他會信.
他不知道什麼原因.
可能他來到會場.
可能是那首詩歌.
可能是看到一些情景他會信.
所以每一個機會.
都是一個撒種.
或者很多年前.
他聽聞現在他真的有一個改變.
所以不知道.
不一定是那個講道的人.
很多的因素.
聖靈都可以用的.
就是in power和in spirit.
所以這個重要了.
就是為什麼我們要祈禱.
就是盼望神的聖靈.
有力.
進入那個人的心靈.

$^{1201}$無論他怎麼樣混亂也好.
聖靈會給他一個.
突然之間.
豁然開朗的明的.
有些人可能很大反應的流淚.
有些人很開心的信主.
每個人都不同.
但是保羅在這裡告訴我們.
記住有幾個步驟.
被神所愛.
是神所揀選.
他會聽明白福音.
講福音的人.
原來不知不覺.
是in power.
和in spirit的去講.
使聽的人會相信.
明白 接受.
原來我們在製造.
一些神在萬世以前.
揀選的人.
我們就透過我們.
此時此地.
這個星期日.
的下午報道會.
利弊光牧師.
會在紅恩堂講.
我們製造這個機會.
有人真的舉手決戰.
如果他真心信.
原來萬世以前神的揀選.
就在星期日發生.
我們就在做這件事.
記得我們是.
不知道任何一個人.
是誰被揀選.
誰沒有被揀選.
額頭是沒有作字的.
是不會有的.
有些人是很貼近基督教.

$^{1241}$甚至從小在教會長大.
他可以是不得救的.
不信神的.
但有些人可能是.
一直都很拒絕的.
但不知為何他後來就信了.
所以我們不能概觀定論.
不要這麼快看死一個人.
也不要這麼快看好一個人.
有些在教會成長的.
很多都後來不再信了.
沒有什麼譴責的.
完全沒有的.
揀選的.
多少年青人.
現在都不在教會了.
我自己沒有跟他一樣.
跟以前的一樣.
很多曾經很熱心的在青年團.
一個都不回的.
都沒有了.
只能怪導師做得不夠.
或者我們令他們不明白福音.
但是其實更嚴重的是.
他根本沒有被揀選.
這個才是最心軟.
最恐怖.
但是我們記得.
我們不是神.
我們不知道誰沒有被揀選.
我們知道的就是.
我們盡力將福音.
勸說.
讓人明白.
去得救.
這是很重要的.
所以保羅就告訴我們.
很重要的神學.
保羅的神學就是這樣理解.
揀選.

$^{1281}$為什麼會被揀選.
不能解釋神的愛.
愛他 定義愛他.
不能解釋.
我和你都是神定義愛的.
這個愛就是揀選.
我們就在.
出現在這個世界.
我們就會成長順著.
就是因為我們聽明白.
這個福音.
所以對帖撒羅尼迦人來說.
他們如果堅固了保羅.
收到這封信.
知道保羅的講解.
保羅是愛我們這一班人的.
他不是見到困難就馬上走.
是他只不過是不想.
那個情勢去到失控.
去到被人.
可能官方有機可乘.
要鎮壓整個教會.
就不想這樣可能他要離開.
或者當時他有感動.
神要他離開.
都不出奇的.
但是更重要的就是.
他對這個教會的不離不棄.
肯定是一個.
保羅要跟進.
因為他帶領他們信主.
一定要跟進這個教會的成長.
和到底情況是怎樣.
所以.
好像應該叫做.
受後服務.
終身保養.
作為一個牧者就應該對信徒.
有這一份這樣的愛.
有時在街上偶然間見到一些沒有回教會的.

$^{1321}$弟兄姊妹.
有時都會默默為祈禱.
有時可能打個招呼走過.
有些都不想看你.
真的在元朗有位姊妹.
小時候已經認識她.
她父母都在教會很多年.
成家都回教會.
後來就沒有回教會.
中學之後就沒有回教會.
後來就知道她已經嫁了.
好像前晚在街上見到她.
她看一看我.
就轉了另一面.
沒辦法.
因為已經不回教會不信.
其實教會很大.
我的母會的地方.
想像得到.
一定要留心這件事.
所以我們要告訴別人.
傳福音不是說.
你揀選了.
神揀選你了 你快點信耶穌.
這樣會令到他困惑.
只能夠說明白救恩的內容就算了.
他要信了主.
他要回望.
原來.
我是在萬世以前被神揀選的.
他就覺得很榮幸.
還有那種榮幸.
他就覺得很榮幸.
還有那種榮譽感.
信耶穌.
當你去到有榮譽感.
你不介意說自己是基督徒.
你很開心.
告訴別人你是基督徒.
因為你已經明白揀選的真理.

$^{1361}$這個神學.
很重要的.
這是我們的內功心法.
令我們能夠支撐.
在任何的環境.
是因為神在萬世以前.
愛你.
揀選你.
所以你能夠成為.
祂的兒女.
不是偶然.
我們要記得.
人生有太多的事情.
很多事情我們根本不知道.
所以.
真的不知道.
如果我們能夠有肯定.
我們能夠握住的.
無論我們去到.
什麼環境.
甚至那個時候還有沒有基督徒.
沒有教會.
那些都不重要.
因為我握緊了神愛我.
我是被神揀選的.
就已經足夠我們面對不同的人生的遭遇.
譬如有些人可能快要移民.
他要知道.
神和我一起.
我有神的揀選.
就會很不同.
這裡就是告訴我們.
保羅就是對他們的肯定.
對他們一個很大的安慰.
就是將神的揀選.
灌注在他們的心靈裡.
在信徒的內心.
有神的揀選.
我經常說笑.
好像張無忌.

$^{1401}$他有九陽神功.
護體.
所以其他的武功.
打進去都不影響他.
因為他真的有一個功.
在他裡面.
我們就需要有這些神學.
所以如果我們很輕易地說信.
你也很容易面對困難.
你就不信.
因為信仰來得太容易.
你也未明白背後的真理.
隨著你的感動.
一定會慢慢地消退.
消退之後.
那個就不再是感動.
就是對真理的堅持.
這個就是我們為什麼一定要在真道上.
站立得穩.
輾轉保羅再提到.
他們在這個信裡.
又會得到另外一樣東西.
就是有聖靈在他們裡面.
聖靈一定會將神的愛.
囂觀.
聖靈會使他們經驗到一份.
屬天的喜樂.
所以信主的人.
必然會面對.
有聖靈的同在.
帶來的喜樂.
那個不是開心.
是心靈的.
平安而來的.
那種滿足.
心靈的喜樂.
所以喜樂和開心.
是兩回事.
是不同的.
這個喜樂.

$^{1441}$是會包括對真道的領受.
明白和實踐.
所以.
聖靈是不會令到你.
越來越不讀經.
只追求超自然神蹟.
那就不健康.
偶然一次可能不出奇.
但長此下去.
經常是這樣.
你就離開聖靈.
你不要說你自己被聖靈感動.
你離開聖靈.
因為聖靈一定會令到你有.
心靈裡面的平安.
平安日於言表.
就是喜樂.
喜樂是會令到你更加對真理.
更加明白.
更加認識真道.
所以這裡是保羅告訴我們.
第六節.
他們的經驗.
在困難裡面的經驗.
接著不單止這樣.
這個教會他們好到一個地步.
當時證明.
當時商業貿易.
那些航運.
都是很發達的.
那些書信來往.
都是很有規模的.
如果不是.
那些風聲怎會這麼快傳到馬其頓.
那時沒有WhatsApp.
也不會有Zoom.
怎會有即時.
怎會有機會傳到外面.
而且是好的明星.
你要記住.

$^{1481}$人類.
發明了互相傳遞訊息.
好的訊息.
通常是很慢的.
通常是七除八扣的.
有些訊息.
不真確的訊息.
就立即傳到全個地球.
壞話.
是非.
所以這裡告訴我們.
他們真是很好.
第八節.
因為主的道從你們那裡已經傳揚出來.
你們向神的信心.
不單止在馬其頓和阿該亞.
阿該亞即是最南面的哥林多.
我們那次探索.
土耳其的探索之旅.
就是由哥林多.
走回頭.
由雅典去哥林多.
接著上去帖撒羅尼迦加.
再過去對面的菲律賓.
剛剛反過來.
現在.
很厲害.
連歐洲南面的哥林多.
這一帶地方.
都知道帖撒羅尼迦加教會.
他們在患難裡面.
繼續有一個好的見證.
即是說.
見證主耶穌.
真是他們心靈的君王.
真是他們盼望的.
那一位主.
尤其是.
連續五次說主再來.
其實都是一個.

$^{1521}$頗尷尬的位置.
因為.
那個再來.
其實君臨天下.
有聽過.
帝王的出巡.
就是.
就是這個意思.
所以.
該撒去到一個.
微服出巡.
或者官式的出巡.
就是一個這樣的臨在.
所以現在不斷說.
耶穌的回來.
好像就是一個.
令到.
當時的官方.
是會懷疑.
你說的耶穌是否真的活生生的一個人.
真是一個.
頗顛覆羅馬政權.
所以.
你可以說保羅.
真是在說.
不過他在說的是世界末日的再來.
就不是.
一個.
另一個政權的興起.
所以.
保羅都是很清楚.
自己所做的.
他不是要顛覆.
地上的羅馬政權.
甚至乎.
我們說過.
他是在那些.
菲律賓的官員面前.
來爭取自己.
把羅馬的身份證拿出來.

$^{1561}$我是羅馬公民.
你無緣無故監禁我們.
是不可以的.
即是說羅馬公民.
原來是有一定的.
優越和.
有保護.
所以這個也是.
保羅都會用的.
所以他又不是真的和羅馬政權為敵.
他都是會用.
那些身份的.
不過有一樣東西他不能夠不說.
就是主的回來.
這個第二次.
他就真是很清楚在說.
世界末日的主再來.
所以其實幫助了我們.
因為第四章.
是說那個次序.
主再來的次序.
世界末日的主再來的次序.
所以.
這五段的說主再來.
肯定是世界末日的主再來.
所以.
在保羅的神學裡面.
對於世界末日的主再來.
是很清楚不含糊的.
所以.
我們作為基督徒.
我們都要有一個心理準備.
世界末日.
我們都要等待.
雖然我們應該等待的.
就是我們自己的安息.
我們可以有些可能計算到的.
譬如你突然間患癌症.
醫生說你只剩下半年.
你的肝有事.

$^{1601}$可能三個月.
這些我們可以計算到.
是沒有辦法的.
因為就算現在最高最先進的醫學科技.
都改變不了.
那沒有辦法.
那我們就安息去立刻建築.
但是其實另外一個心態就是.
既然.
末日的主再來.
是這麼具體.
我們就應該要.
在我們的心靈裡面.
要有這個位置.
去渴慕耶穌的再回來.
就不是我們的離世.
是我們活著見耶穌.
就不是死了.
安息靈魂去到天上.
是兩樣東西都對的.
同時間都要這樣祈禱.
我們要這樣去祈禱.
都要這樣祈禱.
這樣等待的.
所以如果教會.
少了說主再來.
末世的主再來.
其實是未必最完全.
應該越是有困難的教會.
或者那些時代.
越應該說.
末世的主再來.
所以這裡就告訴我們.
亞該亞和馬其頓.
一帶的地方.
都誇耀.
帖撒羅尼迦城.
這個教會.
他們的信心.
愛心和對主耶穌的盼望.

$^{1641}$保羅說.
都不用我們怎樣為你們宣傳了.
第九節 因為他們.
就是這些亞該亞省.
和馬其頓一帶.
其實是很大的地方.
包括可能保加利亞.
即是現在的.
再北一點的地方.
東歐一帶.
那些地方.
很多地方 例如現在的希臘.
已經包含了亞該亞省.
因為他們已經報名.
我們是怎樣去到你們那裡.
你們是怎樣離棄偶像.
歸向神.
要服事又真又活的神.
這裡就是.
帖撒羅尼迦的弟兄姊妹.
的生命改變.
就是知道不應該再拜偶像.
到底在帖撒羅尼迦.
有甚麼偶像呢.
保羅沒有詳細地說.
因為似乎外邦人.
拜的都是來來去去的.
就師.
他們在羅馬有個眾神廟.
希臘有個巴特隆神廟.
都是這些.
當時比較.
偶像的.
都是這些.
當時比較.
多人拜的.
如果在亞洲 例如二匹鎖.
就是拜那個女神.
亞底米女神.
或者.

$^{1681}$英文就應該翻譯做.
戴安娜.
戴安娜原來是女神.
那你就可以想像到.
都是那些.
但羅馬帝國越來越強大.
所以蓋薩.
都當自己是神.
後來就應該還沒有尼祿的時間.
可能尼祿之後.
他們規定.
國民都要拜.
蓋薩是神.
所以.
現在他們這樣離棄偶像.
雖然可能是一般性.
真的不再拜他們自己拜開的.
但是.
就歸向了這一位真神.
就是猶太人的.
真神耶和華.
猶太人本身自己拜.
這一位獨一的神.
但是現在.
外邦的貼索羅尼加.
歐洲的人.
都開始歸信這一位真神.
和猶太人不同.
上星期我們都講過了.
凸顯了這個最重要的問題.
猶太人為什麼不信.
信耶和華.
而不可以上天堂.
而不能夠得救.
是因為他們只是信神.
在整個救恩裡面.
神已經猜顯了他的愛子.
耶穌成為人類唯一的救主.
沒有第二個中保.
沒有第二個救恩的路.

$^{1721}$所以猶太人.
被自己的絆腳石.
絆倒了.
因為他們拒絕了這一位神的兒子.
我們就不同了.
我們因為信耶穌.
我們連神都信了.
我們就敬拜獨一的神.
雖然我們今天不需要獻祭.
但是我們向耶穌祈禱.
其實已經是獻祭了.
我們已經去到神那裡.
透過耶穌.
我們就可以去到真神那裡.
當然.
應該說.
我們每個人.
隨時可以去到神那裡.
敬拜禱告.
但是有時候.
在我們那個信心和.
我們那個祈禱的建立裡面.
有些地方如果特別神聖.
其實是對的.
猶太人.
我都說過了.
猶太人看聖殿是WIFI特別強的.
他們在那裡祈禱.
直上天庭.
那些的祈禱.
這是猶太人的信念.
我不敢說.
一個地方我們分別為聖.
那個地方是說明.
只能夠我們敬拜神的.
我不敢說.
特別強WIFI.
直上天庭的那些祈禱.
但是至少你可以在那個禮拜堂.
安靜.

$^{1761}$天主教是很接受這件事.
他們無論如何.
中午他們都要打開.
那個教堂.
讓人去坐.
他們不會怕人弄壞那些東西.
從來沒有試過有人這樣弄壞那些東西.
其實那些風琴,鋼琴都很貴的.
他們不會覺得.
你浪費我的電.
他們不會的.
他們會覺得你安靜.
有時候.
那個傳道的心態不同.
其實我們是應該.
讓人去到神面前.
所以有時候.
我們要令到.
早期我們禮堂.
變成一個.
靜思靈裡.
打開禮拜堂.
給那些街坊來.
開頭那些人以為機關都不敢進去.
因為長期都封閉.
好像怕人進來幹什麼.
走開.
很多時候有時候會有這樣的感覺.
不停地叫都不敢.
後來慢慢慢慢.
不斷地發現.
真的開門真的歡迎我.
真的沒有人騷擾.
真的慢慢.
當然不會多 因為始終是感受區.
有些人.
其實不懂的.
真的來安靜祈禱.
所以其實中環就更加.
發揮到這個作用.

$^{1801}$因為真的很繁忙.
很多人其實工作很亂.
那個心靈.
你不要管他在那裡睡覺也好.
他安靜下.
那個心靈的一種洗滌.
其實就很大幫助.
所以這裡告訴我們.
我們歸向神.
很重要.
第十節.
我們和猶太人不同的地方.
就是我們有耶穌.
有神.
但是猶太人.
他們世世代代有神.
但是沒有耶穌.
那就是不得救.
這兩樣東西是天淵之別.
所以我們很少人會信耶和華.
有的.
有些跟人嫁了猶太人.
他們就信耶和華.
香港人.
你上網YouTube有些人.
他信入了猶太教.
諸如此類.
但是他不得救.
因為他信神.
不信耶穌是不得救的.
我們剛好相反.
我們是透過認識耶穌基督.
我們就得著神.
耶穌自己說的.
人有了我.
就有生命.
信了我就是信這位差我來的神.
所以.
這裡很重要.
第十節.

$^{1841}$我們有這個盼望.
猶太人等待彌賽亞.
是不是很諷刺呢.
猶太人等待救約.
神應許的彌賽亞.
但是一直都等待不到.
我們.
是等待一位已經回到天堂的耶穌.
他是末世再來的救世主.
所以是大家都不同的.
所以猶太教.
如果他自己是信神是不得救的.
就算他在哭場不斷祈禱.
也是不得救的.
但是我們等候的是神的兒子.
將來從天降臨.
這個肯定是第二次.
最後的末世的再來.
是世界末日的再來.
所以我們一定要定性.
這一次是世界末日的再來.
就不會再有人類的歷史了.
這就是.
等候他兒子從天降臨.
就是他從死裡回到天堂.
等候他兒子從天降臨.
就是他從死裡復活的.
那位救我們脫離將來憤怒的耶穌.
所以我們是免去將來的憤怒.
我們已經在主耶穌的救恩裡面.
我們是免去將來的憤怒.
有多憤怒呢.
那個憤怒.
神的憤怒有多大呢.
在《禱告》的15章裡面說的那個神的烈怒.
純一不雜.
你可以想像得到將來的恐怖,可怕.
所以我們是在主耶穌的救恩裡面.
已經是有憤.
我們是被揀選.

$^{1881}$我們是蒙神所愛.
我們是會在人生對真理明白.
我們會有聖明給我們的喜樂.
這一切.
更加重要的是.
我們在這個世界有盼望.
這個盼望就是主會再回來迎接我們.
很快揭曉第四章.
就是那個次序.
四章是十三節.
很詳細說主再來.
說那個次序.
關於遂了的人.
即是在地上已經死去的人.
我們不願意弟兄們不知道.
恐怕你們憂傷.
有指望的人.
似乎是有一個.
衍生了一個.
不知道什麼想法.
就是死了的人.
耶穌回來.
跟他沒有關係.
不知道他會去哪裡.
有很大的擔心和憂傷.
就好像那些沒有指望的人.
即是那些可能不信的人.
會有一個這樣的謠傳.
或者一個誤解.
所以第十四節.
保羅就在這個牧羊的過程.
聽到這樣的問題.
他就要去澄清.
他就要去解釋.
到底死了的基督徒.
將來是怎樣的呢.
當然他應該說現在是怎樣的.
第十四節.
我們若信耶穌死而復活.
那已經在耶穌裡遂了的人.

$^{1921}$很清楚了.
已經即是說.
在主耶穌還沒有回來之前.
就已經安息主懷的.
那他現在怎樣呢.
神也.
必定將他.
與耶穌一同帶來.
這裡就告訴我們.
我們不需要擔心.
那些物質的變化.
一同帶來.
即是說.
他不是將帶的灰回來.
耶穌不是帶的骨灰回來.
一定是帶那些人回來.
所以.
我們就不需要擔心.
已經一百年.
二百年.
根本在歷史上都沒有人記得.
沒有人認識.
三百年前的.
就算是頗為出名的人.
最多也只是十多個.
有記載在歷史上.
大部分都沒有記載.
就已經離開了.
沒有人知道.
但耶穌會與他們帶回來.
神使這些人.
與耶穌一同回來.
甚至再詳細解釋次序.
就是在說復活的次序.
我們現在照主的話告訴你們一件事.
照主的話.
即是有真理的.
我們這些活著.
還存留到主降臨的人.
他即是說堅信.

$^{1961}$在他那一代.
耶穌會再回來展開那個活著.
那些人會再回來.
耶穌會再回來展開那個末日.
但是.
即是說.
我們都應該要學習保羅.
這樣的一個想法.
不只是安息主懷.
我們是活著見耶穌.
這裡告訴我們.
斷斷不能夠.
在那已經遂了.
的人之先.
原來這裡有個次序.
主回來的時候.
是那些死了的聖徒.
復活一定是復活.
不是拿些骨灰回來.
第十六節.
因為主必親自.
從天降臨.
還要強調Himself.
他自己.
親自.
沒有代理人.
即是說復活的耶穌.
在天堂的永活的耶穌.
又要離開天堂.
要來到這個地球.
因為主必親自.
從天降臨.
有呼叫的聲音.
是主耶穌呼叫.
還是天使長的呼叫.
和天使長的聲音.
似乎這個呼叫是神聖的呼叫.
因為第二個聲音.
就是天使長的聲音.
天使長.

$^{2001}$即是說很高階的.
不是普通的天使.
最高級的天使.
來叫 叫咪.
你可以想像得到.
有神的號.
吹響 一定是天使吹.
那在基督裡死了的人.
這裡是說.
再回來和耶穌一起回來.
就是這些在基督裡.
死了的人.
必定會首先復活.
以後我們.
這活著還存留.
說了兩次.
活著還存留.
即是說還未肉身離世的聖徒.
總有一個時代是這樣的.
因為耶穌總有一個時代.
一個時間回來.
在那個時間還未離開世界的人.
就會活著見到耶穌.
但是他們活著見到耶穌之前.
因為他們不需要死.
即是說還活著.
但是原來.
第一批是那些死了.
即是說古往今來.
在耶穌裡面死了的聖徒.
摩西 阿巴拉罕.
阿當.
他們會復活.
和耶穌一起回來.
到底是否夠奇怪呢?.
天地滿滿的人.
到底是靈體還是怎樣呢?.
這就不知道了.
那個復活的靈體.
似乎不需要佔據空間.

$^{2041}$因為不然會很擠迫.
你想想突然有五萬人.
湧進來感受.
你都不行了.
你都想像得到.
那個壓迫感.
即是在大球場四萬人坐滿.
突然間.
要走了有炸彈.
你都不行了.
你可想像古往今來死了的.
全部要復活.
全部跟耶穌一起回來.
你可以想像.
那個情景.
是多麼美麗.
接著.
我們這些活著還存留的人.
必定和他們.
這些復活的人.
一同被提到雲裡.
所以.
似乎真的是靈體.
不佔據空間.
不然天空就很壓迫.
很多人.
排隊.
你想想從香港排到.
溫哥華.
那些.
你可以想像到.
天空就不用有紅綠燈.
因為大家高度可以.
一些是飛一萬公尺.
一些是飛三萬公尺.
但是將來.
那個靈體是怎樣.
飛的 怎樣出現.
誰在三萬公尺.
誰在一萬公尺.

$^{2081}$不是這樣量度.
所以似乎都是靈體.
但是這個靈體.
應該大家都可以見到.
不受到各方面的影響.
即是說.
這些靈體就在天上.
我們就被提到雲裡.
在空中與主.
和那些聖徒相遇.
這樣.
我們就要和主永遠同在.
這個就是.
世界末日的主在來.
因為是和主永遠的同在.
不是見了見面就回來再做人.
而是.
我們都會身體改變.
突然間你會變了你的身體.
變成一個靈體.
就去到空中.
就和主相遇.
接著就進入那個永恆.
所以.
如果保羅不是在開玩笑.
這個就是一個絕對的真理.
還沒有兌現的基督教的應許.
所以你們當用這些說彼此的勸諱.
為什麼基督徒有時候會這麼軟弱.
我們都不說主在來的那種興奮.
我們身體改變的那個生命.
新生命.
不需要再有任何的限制.
進入永恆.
我們都很少去提到.
我們身體改變的那個生命.
新生命.
進入永恆.
我們都很少去談論.
根本不知道有.

$^{2121}$那裡才是最麻煩.
保羅說我們要談論.
初信的教會.
四個星期.
不是信了主五六十年的老的聖徒.
是剛剛初信.
還是壁壩的教會.
這個就是第四章說的主在來.
這個次序很清楚.
還要是.
主親自降臨.
有呼叫聲音.
有神的號吹響.
接著那些.
死了的聖徒首先復活.
和主耶穌在天上.
一定是我們變聖靈體.
不受時空的限制.
不受物質的限制.
接著就是我們.
這些活著的.
身體改變.
又變聖靈體.
又去到可以永遠活在天堂的那個靈體.
來去永活.
一直都在說永活.
是再.
他們從死裡復活.
接著永活.
就好像耶穌.
死的時候.
接著復活的那個形體.
這個就是.
主在來的問題.
所以第一章的第十節.
是開始說.
一個教會.
是必須要教導信徒.
我們固然是需要.
在地上勤奮的工作.

$^{2161}$有的.
貼手來的前書第四章就有教導.
不要.
因為主在來很快.
人們就懶惰不做工.
四章有教導.
就要責備那些人.
但是.
這裡實際告訴我們.
我們需要有一個永恆的盼望.
信了主耶穌.
原來不只是地上.
我們開開心心的.
可以入除就林.
有份.
去黃塵.
接著就去.
上中學去哪間呢?.
基督教好一點的.
陳瑞芝.
我們都期望是好的.
不是壞事.
一間好的學校的分途訓練.
都是好事.
將來可以更加被神用.
但是地上有的.
福利.
還加上.
可以去基督教墳場.
落葬.
如果不是受浸.
沒有位置.
給多少錢都不可以進去.
不浸就沒有份.
這些全部.
是地上的福利.
沒有錯的.
不過我們要追求的是.
將來憤怒的.
耶穌.

$^{2201}$現在用祂的愛.
用祂的釘痕的手.
去拯救我們的耶穌.
現在祂願意.
用一份無條件犧牲的愛.
去招納我們.
進入祂的救恩裡面.
所以.
我們在主耶穌裡面.
真是一次得救.
就一定會得到.
這個永恆生命.
記住是一直相信.
就得到這個永生.
將來.
這位來救我們的.
耶穌.
就會迎接我們的.
耶穌.
就是四章剛才所說.
就不是用祂的憤怒.
對待的那些.
敵對祂的人.
所以這裡第十節.
就是告訴我們.
一方面我們是基督徒.
因為耶穌.
我們有了這位真神.
又真又活的神.
猶太人他們只是信一位.
又真又活的神.
但他們不信神已經設立的.
救主救恩.
他們是不會得到神的拯救.
所以這是很.
重要的一個分水嶺.
所以我們今天就是要.
透過耶穌的救恩.
逃避將來的憤怒.
這個是.

$^{2241}$耶穌基督祂所給我們的.
這個拯救.
第一章很快地.
幫大家複述了一點.
希望第十節.
作為我們來看.
而來到更加有把握.
對我們自己的信仰.
和救恩.
第二章我們可以.
稍微看一點.
弟兄們.
你們自己原曉得.
我們進到你們那裡.
並不是徒然的.
我們從前在.
被害受辱.
這是你們知道的.
然而還是靠.
我們的神放開膽量.
在大.
爭戰中.
把神的福音傳給你們.
我們的勸勉.
不是出於錯誤.
不是出於污穢.
即是提到動機的問題.
也不是用鬼詐.
但神既然.
嚴重了我們.
把福音託付了我們.
我們就照樣說.
不是要討人喜歡.
是討這位.
察驗我們心的.
神喜歡.
因為我們從來都沒有.
用過淺味的說話.
這是你們知道的.
所以這裡真是將他內心.

$^{2281}$作為好目人.
來表白.
即是說.
透過保羅現在的言說.
即是說.
帖撒羅尼迦的教會.
極有可能會被.
這些暴徒.
來威逼利誘.
要不就相信我們的該殺.
要不就不要再相信這位耶穌了.
這個保羅.
說得這麼好.
早就走了.
一定是他見到困難.
自己先走了.
你還相信他.
可能一定是這樣.
稍為解釋.
他那個傳道的.
內心和動機.
保羅原來.
一方面就是我們在肥納比.
我們都是.
遭遇那個.
辱罵.
即是無妄之災.
一個有錢人.
那個小妹.
是懂得法術的.
可以預知.
占卜未來的.
使得主人.
可以發財的.
不知當時有甚麼賭馬.
還是買甚麼.
彩票.
後來.
保羅見到這個.
小妹.

$^{2321}$那個小妹就知道保羅是神的僕人.
因為他裡面有鬼.
就.
不斷跟著保羅.
不斷去和人說.
到底他在.
被鬼附之.
的情況下.
說還是他在.
即是.
特意詆毀保羅.
嘲笑保羅.
保羅都不理他.
由得他.
但後來覺得很厭煩.
可能引起一些輿論.
引起一些人的討論.
所以保羅.
就將他裡面的鬼趕走.
不讓他再騷擾.
其實是趕走裡面的鬼.
即是意味著.
可能是這個使女裡面的鬼.
在說保羅.
他們和西拉兩個.
所以因為這樣.
主人發現.
使女不可以再.
幫他做一些.
投資.
買一些彩票.
就沒有那些法力.
就遷怒於保羅.
保羅和他無仇無怨.
現在令到我.
沒有了這個賺錢的門路.
就似乎有一定的勢力.
就抓了保羅和西拉.
去見官.
坐在菲臘比的監裡.

$^{2361}$這個就是.
保羅說他在菲臘比.
被害受辱.
就是這件事.
這些都是你們知道的.
即是意味著.
原來使徒.
和他們的傳道.
都叫眾所周知.
他們不是.
為財為利.
不是為了出名.
不是為了自己的利益.
來傳道.
這是一個很重要的推薦信.
很重要的一個.
在一些不同的地方.
讓人看到.
你一如既往.
的傳道的.
風格.
那個心態.
會對教會.
有一定的.
美好的一個榜樣.
所以保羅在這裡.
就真的去.
講他內心的處境.
和.
希望讓貼索羅尼加的教會.
不要誤會他.
為甚麼會見到難處就馬上走.
就不留下.
陪伴弟兄姊妹.
面對那個誤解.
那個.
對耶穌基督的.
君王.
這個錯誤的指控.
有一個澄清.

$^{2401}$但是保羅.
是要突然離開.
都不夠時間去解釋.
所以就令到.
教會可能會受到.
一些的影響.
有人會在批評保羅.
但是唯一.
可以解釋的就是.
將自己的內心的動機.
都去.
讓人知道.
這裡就在.
第三節就說.
在大真賤裡面.
將神的福音去傳.
我們的勸勉.
不是出於錯誤.
不是出於污穢.
不是用詭詐.
是神託付了我們.
這個福音.
我們就去傳遞這個福音.
保羅就是帶著福音的訊息.
有神的託付.
去傳遞.
他不知道誰是被神揀選.
但是他的福音去到哪裡.
誰人的生命得到改變.
那個就是神的揀選.
所以.
我們是將.
神的揀選.
來去具體地.
實現在地上.
所以.
我們要有這個心態.
就是我們去傳福音.
我們去報導.
都是將人帶進.

$^{2441}$他自己去揀選神的.
那個決志裡面.
所以又是一樣.
我們要信了主之後.
都要成長.
我們要看到.
神有沒有嚴重我們.
嚴重即是說.
看得起你才會嚴重你.
經過那個測試.
我們是.
頗堪被神用.
跟著他才會託付.
都是一個過程來的.
由你甚麼都不懂.
開始慢慢讀聖經.
越來越有興趣.
開始對教會對聚會.
你有那個感動.
熱心其實都是一步步.
累積上去的.
開始你自己對神.
很想告訴別人.
慢慢慢慢.
驗證.
驗完之後就託付.
所以.
一個有託付的基督徒.
才是一個社會.
最重要的棟樑.
今日很多人都有.
那個託付.
不過可能是他自己的.
家族.
可能是為錢.
可能是其他的東西在社會裡面有貢獻.
但我們基督徒.
要有神的託付.
所以如果我們.
很多基督徒.

$^{2481}$能夠.
感受到自己有託付.
那個教會就很重要.
因為.
每個人的心都有這一份負擔.
有時我們.
我現在就不是在.
錦宣 是在紅欣.
有時跟公會.
有些資料.
現在我見到.
四月份有幾個同時間.
的主一學.
看誰多些誰少些.
現在還未到.
下星期才是.
加起來.
都沒有30個.
可能還未報名.
三班加起來都沒有30個.
我在紅欣.
他們每一班上主一學.
最低限度都有15個以上.
只有80人的教會.
有15個人上主一學.
有70-80人的教會.
30個人.
你計算一下比例.
其實我想說的不是要做什麼比較.
完全沒有這個意思.
而我想說的是.
那個教會要訓練到弟兄姊妹.
是深深感受到.
自己被託付.
是被神驗證.
你看到那些教會.
是極之興旺.
就一定是這個問題.
連他自己都不覺得被神驗證.
你不要驗我.

$^{2521}$你不要干涉我.
這樣很難的.
你不驗證還何來託付你呢?.
託付即是說.
我交給你這個信物.
是很重要的.
我託付你了.
你好像古代那些運標的標局.
你去做保鏢.
託運.
我們今天很少基督徒.
或者教會.
說這麼嚴肅的真理.
我們每一個都要裝備自己.
成為有神的託付.
在這個時代.
是要有一個信息要去傳.
就是我們要捍衛著神的真理.
很感恩.
上次展秋那個.
馬丁路德的.
那個中學改革的神學.
有50多人回主教學.
應該要延續.
有些可能.
很有趣.
有很多歷史.
另外一些可能是.
當時的背景等等.
這些是很令人興奮.
興趣的.
但是聖經科.
或者其他的學科.
都應該要有這樣的裝備.
其實是很重要的.
託付和驗證.
福音託付.
如果我們的教會.
沒有福音的託付.
根本大家都不緊張福音.

$^{2561}$那教會是否存在都沒有用.
都沒有需要.
因為你們唯一的教會可以講福音的.
不講福音.
那你想哪個教會去講呢?.
哪個社團去講呢?.
關愛隊去講福音?不會的.
我們一定要記住.
教會存在在這個地上.
就是神驗證我們.
託付福音給我們.
我們去傳.
所以.
如果教會已經不傳福音.
大家都不熱心去傳福音.
其實就得不到神的很大的Blessing.
我們繼續要製造那些歡樂.
就是人為的歡樂.
我們以為很開心.
但其實可能已經離開了.
神最重要的心意.
所以我們要記得.
這個託付的.
又是一個很重要的真理.
有了託付.
我們就會裝備我們自己.
有些東西不懂.
自己找書.
自己再去看.
去想深一點.
早期我讀經.
很多不明白.
真的要逼自己去想想上下文.
去想通它.
再去問人.
你就有這種追求.
如果不是.
聽完就算了.
也不影響我的工作.
對我的工作完全沒有關係.

$^{2601}$當你這樣想的時候.
你連那種興趣.
好奇都沒有.
這才是基督教的悲哀.
譬如年輕人.
對很多事都有興趣.
唯獨對聖經沒有興趣.
這真是悲哀.
是一個時代的悲哀.
盼望神讓我們能夠.
看到福音的寶貴和重要.
所以第四節.
今天和大家看到這裡.
神既然嚴重了我們.
把福音託付我們.
我們就照樣說.
照樣不要理那個時代.
是順著基督教.
還是反對基督教.
因為.
誰該殺.
是善待基督教呢.
是沒有的.
由奧古時都開始已經沒有了.
所以無論提比留也好.
對格魯雕也好.
對尼祿也好.
都是不善待基督教.
所以既然神託付了我們.
我們就照樣說.
不理那麼多了.
不討人的喜歡.
包括當時羅馬的皇帝.
乃是要討.
測驗我們心的神喜歡.
就說到這裡.
第四節.
我們一起禱告.
感恩天父給我們又過半小時.
很詳細的將.

$^{2641}$《貼實羅尼加前書》的一章和第二章.
去講解.
求主給我們堅定我們的.
揀選在神的裡面.
因為我們認識了.
主耶穌是我們的救主.
是神你.
差遣來這個世界.
將我們拯救的救主.
多謝你揀選我們.
多謝你在萬世以前.
用你的永恆的愛去愛我們.
使得我們今天.
得到主的選照.
和我們可以在地上.
有這個對永恆.
對主再來的盼望.
和將來我們真的.
可以進入永恆的真光.
天上的天家.
去到主的懷中.
幫助我們在今生.
教會是受託付.
我們不是被主.
所驗證.
而不能夠合格的.
相反我們是被主驗證.
我們是能夠.
真的有主的託付.
將福音我們是.
照著神的喜悅去傳遞.
而不是照著人的喜歡.
去刪減.
或者去扭曲.
令到福音變成為.
好像沒有任何的要求.
等等 願神幫助我們.
奉主耶穌基督的名.
阿們.
\newpage



\section{羅馬書 6:1-14}
\label{sec:Su045rsIjTY}
\textbf{2024-03-29 受苦節崇拜《耶穌為你死了,你呢?!》羅馬書 6\_1-14 講員:鄧泰康傳道}
\newline
\newline
連結: \href{https://youtube.com/watch?v=Su045rsIjTY}{\texttt{https://youtube.com/watch?v=Su045rsIjTY}} ~~~~ 語音日期: 2024-04-02
\newline
\newline
\hyperref[sec:xRsKG7u5A0E]{\small{< < < PREV SERMON < < <}}
~
\hyperref[sec:index_chronic]{\small{[返順時目]}}
~
\hyperref[sec:index_scriptual]{\small{[返順卷目]}}
~
\hyperref[sec:ez2giRGwvx8]{\small{> > > NEXT SERMON > > >}}
\newline
\newline
羅馬書 6:1-14
\newline
\begin{longtable}{cl}
\hline
\hline
章節 & 經文 (和合本修訂版)\\
\hline
6:1 & \begin{tabularx}{0.7\textwidth}{X} 這樣，我們要怎麼說呢？我們可以仍在罪中使恩典增多嗎？ \end{tabularx} \\ \\ \relax
6:2 & \begin{tabularx}{0.7\textwidth}{X} 絕對不可！我們向罪死了的人，豈可仍在罪中活著呢？ \end{tabularx} \\ \\ \relax
6:3 & \begin{tabularx}{0.7\textwidth}{X} 難道你們不知道，我們這受洗歸入基督耶穌的人，就是受洗歸入他的死嗎？ \end{tabularx} \\ \\ \relax
6:4 & \begin{tabularx}{0.7\textwidth}{X} 所以，我們藉著洗禮歸入死，和他一同埋葬，是要我們行事為人都有新生的樣子，像基督藉著父的榮耀從死人中復活一樣。 \end{tabularx} \\ \\ \relax
6:5 & \begin{tabularx}{0.7\textwidth}{X} 我們若與他合一，經歷與他一樣的死，也將經歷與他一樣的復活。 \end{tabularx} \\ \\ \relax
6:6 & \begin{tabularx}{0.7\textwidth}{X} 我們知道，我們的舊人和他同釘十字架，使罪身滅絕，叫我們不再作罪的奴隸， \end{tabularx} \\ \\ \relax
6:7 & \begin{tabularx}{0.7\textwidth}{X} 因為已死的人是脫離了罪。 \end{tabularx} \\ \\ \relax
6:8 & \begin{tabularx}{0.7\textwidth}{X} 我們若與基督同死，我們信也必與他同活， \end{tabularx} \\ \\ \relax
6:9 & \begin{tabularx}{0.7\textwidth}{X} 因為知道基督既從死人中復活，就不再死，死也不再作他的主了。 \end{tabularx} \\ \\ \relax
6:10 & \begin{tabularx}{0.7\textwidth}{X} 他死了，是對罪死，只這一次；他活，是對神活著。 \end{tabularx} \\ \\ \relax
6:11 & \begin{tabularx}{0.7\textwidth}{X} 這樣，你們也要看自己對罪是死的，在基督耶穌裡對神卻是活的。 \end{tabularx} \\ \\ \relax
6:12 & \begin{tabularx}{0.7\textwidth}{X} 所以，不要讓罪在你們必死的身上掌權，使你們順從身體的私慾。 \end{tabularx} \\ \\ \relax
6:13 & \begin{tabularx}{0.7\textwidth}{X} 也不要把你們的肢體獻給罪作不義的工具，倒要像從死人中活著的人，把自己獻給神，並把你們的肢體獻給神作義的工具。 \end{tabularx} \\ \\ \relax
6:14 & \begin{tabularx}{0.7\textwidth}{X} 罪必不能作你們的主，因你們不在律法之下，而是在恩典之下。 \end{tabularx} \\ \\
[1ex]
\hline
\hline
\end{longtable}
$^{1}$透過詩歌剛才帶我們再次回到.
主耶穌基督釘十字架的場景.
如果剛才大家很認真很留心.
投入在裡面一起去聽的時候.
你應該在你的心裡面.
在你的面前都會有一幅很大的圖畫.
是主耶穌釘十字架的時候那種張力.
他那種掙扎.
主耶穌基督釘十字架這個訊息.
我很相信大家一點都不陌生.
特別感受到傳統每一年我們在壽苦一節.
我們都會不斷和大家重覆又重覆.
再重覆這個事情.
但是作為傳道人我就會去問.
我們聽了這麼多年.
耶穌基督為我們死了.
你呢?.
接受了救恩的信徒.
這裡講得很清楚.
我們是和基督一同歸入他的死裡面.
信主的一刻.
從前的你應該是已經和基督同釘十字架上面.
這個你今天認知嗎?.
旁邊有一個很好的十字架.
請大家看著這個十字架.
大家看著它.
去問.
你的救我釘在十字架上面了嗎?.
沒有死是不會有復活的.
死是為了復活.
復活是為了讓你的一舉一動有新生的樣式.
耶穌的復活.
他做了一個榜樣.
是一個完全的改變.
是超越一切.
你會看到當耶穌復活的時候.
他是可以隨時出現.
他隨時可以消失.
是和他仍然未復活之前完全兩個形象.
而我問大家.

$^{41}$如果你復活的話.
你改變了嗎?.
你和以前有改變嗎?.
還是我們仍然和以前一模一樣.
和未接受救恩的時候是一模一樣.
什麼是改變?.
剛才在經文裡說是和基督聯合.
不要再做罪的勞瀑.
但是弟兄姊妹.
如果我們不願意死的話.
我們根本沒有辦法脫離罪.
很多時候當我們做傳道人都告訴大家.
我們都有軟弱.
其實我們都有罪性.
我們到死的一天.
我們根本沒有辦法離開罪這個綑綁.
我們都承認我們有這個軟弱.
但在宿命的層面上.
耶穌基督.
他就挑戰你.
你在靈明的時候.
你在你以前的舊惡裡面.
你死了嗎?.
你不願意去死.
你沒有辦法脫離罪對你的綑綁.
然後他提供了一個方法給大家.
你要向神活.
你就能夠向罪死.
向神活.
就是說.
你只為神而活.
弟兄姊妹 這個信息大家一點都不陌生.
我們透過很多不同的門訓課程.
透過很多不斷告訴大家.
我們要為神而活.
你作為神的門徒.
這個是最基本的要求.
耶穌基督挑戰你.
你要背我的十字架.
不然你不配作我的門徒.

$^{81}$這些都是一些很核心很重要的信息.
弟兄姊妹.
我們聽了很多.
弟兄姊妹 請你再看十字架.
你再問自己.
我今天為了神而活了沒有?.
我今天是不是為神而活?.
我做的每一個舉動.
有沒有為神而活?.
我每天所思所想.
究竟是神的事.
還是你的事?.
是你所關注.
還是上帝所關注?.
保羅在這裡挑戰大家.
問了一個很深層的問題.
他說.
你將自己獻給了誰?.
他用一個例子.
一個比喻.
你是一個器皿.
你這個器皿成為誰的器皿?.
器皿是用來裝著東西的.
曾經保羅透過一個經文告訴大家.
你不管是金器,木器,瓦器.
但你要是一個純潔的器皿.
我就會問大家.
究竟今天你這個器皿.
裝著些什麼?.
我們的器皿裡面.
我們的心.
我們的思想.
裝著些什麼?.
保羅說.
你要作義的器皿獻給神.
上帝一直告訴我們.
我們要做義的器皿.
上帝一直告訴我們.
我們做基督徒.
請大家聽見.

$^{121}$基督徒這三個字.
是滿有榮耀.
滿有作用.
滿有功能.
因為人們一聽到你是基督徒.
就應該見到主耶穌.
當今天人們見到你的時候.
見到的是主耶穌.
還是.
是一個醜惡的你.
還是.
是完全看不見基督徒應該有的生命.
所以.
他提醒你.
作義的器皿.
去獻給神.
獻給神.
就是要代表耶穌.
去成為別人更多的祝福.
你今天的生命.
是不是在祝福其他人.
所以我們求聖靈幫助我們.
求神讓我們徹底改變.
徹底改變.
求神讓我們成為祂的器皿.
承載著祂的愛.
承載著祂對人的祝福.
去繼續祝福我們身邊每一個遇上的人.
包括你的伴侶.
你的爸爸媽媽.
你的子女.
你的朋友.
你的兄弟姐妹.
你的親人.
所有一切碰上你的人.
都希望得到基督的祝福.
因為我們只為神而活.
因為我們所有人.
都在上帝的恩典之下.
我們每一個人都是蒙福的人.

$^{161}$我們每一個人都不配跨口.
所以我們只需要做一件事.
成為上帝真正基督的門徒.
每天行在祂的道路上.
去祝福祂.
求上帝幫助我們.
死得透徹.
讓我們在三天之後復活.
成為一個處身的人.
\newpage



\section{馬太福音 28:1-10}
\label{sec:ez2giRGwvx8}
\textbf{2024-03-31 復活節崇拜《復活:衝破生命的枷鎖》馬太福音 28\_1-10 講員:陳麗蟬牧師}
\newline
\newline
連結: \href{https://youtube.com/watch?v=ez2giRGwvx8}{\texttt{https://youtube.com/watch?v=ez2giRGwvx8}} ~~~~ 語音日期: 2024-04-02
\newline
\newline
\hyperref[sec:Su045rsIjTY]{\small{< < < PREV SERMON < < <}}
~
\hyperref[sec:index_chronic]{\small{[返順時目]}}
~
\hyperref[sec:index_scriptual]{\small{[返順卷目]}}
~
\hyperref[sec:_PqGg3XbEA0]{\small{> > > NEXT SERMON > > >}}
\newline
\newline
馬太福音 28:1-10
\newline
\begin{longtable}{cl}
\hline
\hline
章節 & 經文 (和合本修訂版)\\
\hline
28:1 & \begin{tabularx}{0.7\textwidth}{X} 安息日過後，七日的第一日，天快亮的時候，抹大拉的馬利亞和另一個馬利亞來看墳墓。 \end{tabularx} \\ \\ \relax
28:2 & \begin{tabularx}{0.7\textwidth}{X} 忽然，地大震動；因為有主的一個使者從天上下來，把石頭滾開，坐在上面。 \end{tabularx} \\ \\ \relax
28:3 & \begin{tabularx}{0.7\textwidth}{X} 他的相貌如同閃電，衣服潔白如雪。 \end{tabularx} \\ \\ \relax
28:4 & \begin{tabularx}{0.7\textwidth}{X} 看守的人嚇得渾身顫抖，甚至和死人一樣。 \end{tabularx} \\ \\ \relax
28:5 & \begin{tabularx}{0.7\textwidth}{X} 天使回應婦女說：「不要害怕！我知道你們是尋找那釘十字架的耶穌。 \end{tabularx} \\ \\ \relax
28:6 & \begin{tabularx}{0.7\textwidth}{X} 他不在這裡，照他所說的，他已經復活了。你們來！看看安放他的地方。 \end{tabularx} \\ \\ \relax
28:7 & \begin{tabularx}{0.7\textwidth}{X} 快去告訴他的門徒，說他已從死人中復活了，並且要比你們先到加利利去，在那裡你們會看見他。看哪！我已經告訴你們了。」 \end{tabularx} \\ \\ \relax
28:8 & \begin{tabularx}{0.7\textwidth}{X} 婦女們急忙離開墳墓，又害怕，又大為歡喜，跑去告訴他的門徒。 \end{tabularx} \\ \\ \relax
28:9 & \begin{tabularx}{0.7\textwidth}{X} 忽然，耶穌迎上她們，說：「平安！」她們就上前抱住他的腳拜他。 \end{tabularx} \\ \\ \relax
28:10 & \begin{tabularx}{0.7\textwidth}{X} 耶穌對她們說：「不要害怕！你們去告訴我的弟兄，叫他們往加利利去，在那裡會見到我。」 \end{tabularx} \\ \\
[1ex]
\hline
\hline
\end{longtable}
$^{1}$弟兄姊妹早晨.
Hallelujah 主耶穌基督已經復活了.
藉著主耶穌基督的復活.
我們可以能夠衝破了死亡的權勢.
藉著主耶穌基督的復活.
我們都可以衝破我們生命裡面的枷鎖.
這裡有一隻陶瓷碗.
在這個《皓國月刊》第260期裡面.
來介紹的.
這隻陶瓷碗是裂開了的.
但經過修補之後.
那些藝術家就在裂痕裡面.
就塗上了金色的油在裡面.
這個是日本一個新興的一種修補的方式.
以致這隻修補了的陶瓷碗更加的名貴.
可以再一次被人去欣賞.
有人就將這些破碎了的陶瓷碗.
用金油來油上.
與我們這個信仰拉上了一些關係.
就說到其實我們每一個罪人.
我們都是帶著一個破碎的生命.
也是因著這個破碎.
主耶穌基督為我們釘在十字架上.
藉著祂的寶血.
我們可以能夠與上帝恢復關係.
人與人之間只要我們願意的時候.
都可以能夠有復和的關係.
因著主耶穌基督的修補.
我們進入永恆的天國的時候.
我們就像金色油油在破碎的陶瓷碗裡面.
我們的生命可以能夠更加的燦爛.
活得更加的釋放.
這個陶瓷碗就成為了一個新的救贖的故事.
我們就看一下在這個二千多年之前.
耶穌基督復活的日子裡面所發生的事情.
我們彼此的勉勵.
我們在當中先看一下.
在文章裡面的一些枷鎖.
在第一節裡面這樣說.
「安息日過後」.

$^{41}$七日的第一日.
天快亮的時候.
「茉大拉的瑪利亞和另一個瑪利亞」.
離去到墳墓裡面.
當耶穌基督死了之後.
一群門徒和婦女.
他們生命當中突然間失去了.
這三年多裡面的一個忠心人物.
一直有耶穌帶著.
但現在他死了.
好像就沒有盼望.
沒有力量.
沒有朝氣.
我想他們的心情是很痛的.
很沮喪的.
一種感到很迷失.
很疑惑.
甚至可能會感到很憤怒.
為甚麼一個這麼好的人.
他醫了這麼多人.
又幫了這麼多人.
他講道又講得這麼好.
走過這麼多神蹟.
幫了這麼多的人.
為甚麼他要被釘死.
為甚麼那些羅馬兵兵這麼殘忍.
當時可能他們就整個世界.
好像就倒塌了.
可能我們都會經歷過.
當我們的親人離開我們的時候.
我們都會覺得整個人生.
是天昏地暗.
曾幾何時.
或者就在這個現在裡面.
我們都會處於一個這樣的情況.
當我們失去了生命裡面.
我們以為一些很重要的東西的時候.
我們的人.
或者一些物品.
或者一些的崗位.

$^{81}$我們就會覺得.
糟了 好像失去了那份安全感.
我們感到很擔心.
失去了那份倚靠.
當我記得.
在我出來事奉了.
大約半年的時候.
主任牧師就告訴我.
他說一年之後.
他就會移民了.
到時你自己來.
他說的時候.
嚇得我半死.
我自己一個人來.
我怎麼會做.
我剛剛出來工作.
我怎麼會做.
我記得當晚我就發了一個噩夢.
彩色的噩夢.
我就看到了.
早上星期日崇拜.
好像我們感受教會外面有個水牌.
寫著當日的講員.
就是陳麗蟬姑娘.
回到來的時候.
一看到.
今天陳麗蟬講道.
她就馬上走.
當時我們的教會在大衛.
他們就因為等不到.
182巴士.
我很清楚.
等不到182巴士.
於是就寧願走路.
出獅子山隧道.
都要快點離開這裡.
我記得當時.
全部的會眾都走了.
我就自己一個人.
坐在一個空地裡面.

$^{121}$我們教會外面都有一個空地.
我自己就在那裡大哭.
我很擔心.
很難過.
我自己怎麼頂得住這個教會.
怎麼會做.
我覺得主任牧師走了之後.
我就好像沒有了安全感.
沒有了力量.
我又埋怨了.
為什麼這個牧師.
在我入職之前.
不跟我講清楚.
他早點告訴我的時候.
我就不理會了.
結果.
我就屈下屈下.
屈了一個輕微的抑鬱症.
在當時的日子.
很難捱.
昔日這班婦女.
這班門徒.
他們沒有了主耶穌基督.
成為了他們隨時的幫助.
不能夠時至和他們一起.
他們很憂心.
很難過.
所以就在那一天.
七天的第一天.
他們就拿著一些香膏.
去到耶穌基督的墳墓.
心裡面是很想.
為耶穌基督去做.
最後最後的一件事情.
但去到墳墓的時候.
有一件很特別的事情發生.
第二節.
忽然地帶震動.
因為有主的一個使者.
從天上下來.

$^{161}$把石頭滾開.
坐在上面.
那個情景很特別.
當那班婦女很憂愁.
沒有盼望.
沒有朝氣的時候.
突然地帶震動.
然後石頭慢慢地滾開.
然後那個使者.
還坐在上方.
你可以想到一些卡通的場景.
那個使者有兩隻翅膀.
站在那塊大石的上方.
兩隻腳就一起.
搖搖晃晃.
很開心地嘗試.
但是.
地上的婦女.
看到地帶震動.
驚上加驚.
什麼事發生.
不過地帶震動這個場景.
其實在幾天之前.
他們都經歷過.
當耶穌基督死的時候.
這個聖殿的萬子.
就從上到下.
裂為兩半.
聖經就這樣形容.
地帶震動.
盤石崩裂.
在聖經裡面.
有這個地震發生的時候.
即是說.
上帝的大能將會來去顯現.
在舊約的時候.
當摩西帶領.
以色列人出埃及的時候.
他們去到這個西乃山.
準備去.

$^{201}$接受上帝和立約.
和摩西上西乃山.
領這個十誡.
當時聖經形容.
全山無煙.
因為耶和華在火中降臨山上.
整座山劇烈震動.
那次之後.
以色列就和上帝立約.
耶和華就要成為.
以色列民的主.
要保護他們 帶領他們.
以色列民就要.
順服上帝一切的安排.
到耶穌基督死的時候.
有地帶震動.
而這次呢.
在七天的第一天.
這次的地震.
就象徵著耶穌基督.
已經勝過了死亡權勢.
這個死亡權勢.
已經完全的被打敗.
死亡不再是人類的終結.
因為反信耶穌基督的人.
都可以有這個永生.
當耶穌基督在十字架上.
死亡的時候.
這個地帶震動.
就象徵著.
他已經完成了救贖的工作.
打破了罪的權勢.
對人類的捆綁.
以致我們人類.
可以能夠和上帝的關係.
可以復和.
人和人之間的關係.
都可以靠上帝的恩典.
只要我們願意的時候.
都可以能夠彼此的恢復.

$^{241}$大家有沒有試過.
和周圍的人.
有些過節.
大家會不抽不彩的.
我相信一定有的.
不要告訴我沒有.
騙人的.
其實在這樣的情況的時候.
一定是雙方都不好過的.
雙方的心裡面.
都會覺得辛苦的.
我記得在幾年之前.
有一套劇集叫做.
「法正先鋒」.
有沒有看過.
每一個故事.
都會講到.
有些事情為什麼會殺人.
每一次的破案.
都發現到每一個人的犯罪動機.
就是因為不肯放過.
之前一些的仇恨.
每一個被抓到的時候.
他們都會很生氣地這樣說.
這些過往的人.
怎樣對不起他們.
怎樣令他們很痛苦.
就是因為人.
將很多的仇恨.
擠在心靈裡面.
不肯饒恕.
不肯釋放.
不肯放過.
當我們不肯饒恕其他人的時候.
其實都是在幫自己.
在捆綁著自己.
耶穌說.
不要只是看見你弟兄眼中有刺.
卻看不到自己眼中的良目.
古代也都教導我們.

$^{281}$愛是不計算人的惡.
耶穌的愛徒約翰也都這樣說.
有愛能夠遮掩許多的罪.
在主耶穌基督的愛裡面.
我們可以能夠去原諒一些得罪我們的人.
在主耶穌基督的愛裡面.
我們也可以學習到.
我們不再只是看對方很多好不好的地方.
去怎樣刻印真的地方.
主耶穌基督在十字架上的無私的大愛.
就是我們最偉大最高超的榜樣.
耶穌基督復活大能的愛.
不單止使地大震動.
墳墓的石都要滾開.
更加的廢掉了冤仇.
使人和上帝之間可以能夠復和.
人和人之間都可以能夠去重新恢復那個關係.
我們去想一想.
在我們生命當中有沒有一些人.
我們一直都不肯饒恕.
一直都不肯去寬恕 接納他呢.
靠著自己的方法.
靠著自己的力量.
一定做不到的.
但我們告訴這位復活的耶穌知道.
祂會幫我們的.
只要我們肯去遷就.
祂會給我們機會.
幫助我們去經歷那份復和.
那份釋放 那份生命.
接著這班婦女就看到.
在這個墳墓裡面有使者在.
天使就跟一班婦女說.
不要害怕.
我知道你們是來尋找那位釘十字架的耶穌.
祂不在這裡.
祂已經復活了.
祂說你們來看看安放祂的地方.
你們快點去告訴其他的門徒.
告訴他們.

$^{321}$耶穌基督已經從死人中復活了.
並且耶穌會讓他們先去到加利利的地方去.
在那裡你們會再一次跟耶穌基督會合.
你一定會再一次見到耶穌基督.
所以你們要快點去.
我已經告訴你們這個消息了.
天使說耶穌基督不再在這個墳墓裡面.
這個墳墓代表了死亡.
死亡代表了斷絕.
代表了關係的破裂.
代表了人生的絕望.
使者就叫這些婦女.
你們要去告訴你們的弟兄聽.
叫他們去到加利利.
在那裡就會見到耶穌.
加利利是耶穌基督那裡呼召那十二位的門徒.
也都在加利利裡面.
就行了很多的神蹟奇事.
第一個神蹟水變酒.
就在這個加利利裡面發生了.
當時這個辦喜事的家庭.
他們不交酒飲.
因著這個缺乏.
但當耶穌基督知道的時候.
祂就將水變酒.
以做整個的婚宴.
這個本來缺乏的.
現在到你去變得滿足.
在加利利裡面.
也都行了很多的醫病的事情.
叫盲的看見.
跛的走走.
也都去趕出很多的邪力.
在這個地方很多人都經歷到.
人在這個絕望.
沒有依靠的時候.
他就經歷到耶穌的幫助.
他們的生命就從一個很絕望的裡面.
就能夠去變得滿足.
變得豐盛.

$^{361}$所以對門徒和婦女來說.
加利利這個地方是充滿了希望.
充滿了這個恩典.
所以有復活的耶穌.
在我們生命的時候.
我們原本是一個很倒楣的.
很倒楣.
很絕望的.
很不足夠的生命.
可以能夠變成一個充滿了希望.
充滿了幸福的生命.
我們看看這個經文.
它裡面其實講出我們的信仰.
很重要的幾個字.
來 看 去.
它告訴人聽.
就是當我們要來.
我們來接觸一下耶穌.
我們跟別人說耶穌的時候也是.
你不去試一下.
不去跟祂站著的時候.
你經歷不到耶穌.
所以就來.
你來去看一看.
當婦女去到這個墳墓.
裝個頭進去看的時候.
裡面沒有了耶穌的屍體的時候.
她們就知道耶穌基督真的復活了.
當她們去看到 經歷到是真的時候.
她們的責任就是怎樣.
要去 並且要快去.
告訴其他人知道.
婦女發現了.
耶穌基督真的復活了.
在這個墳墓裡面.
再沒有這個死人的屍體的時候.
她們很開心.
她們的責任就要去告訴這個叛徒知道.
將這件事宣佈來告訴所有人知道.
耶穌基督已經復活了.

$^{401}$這個去.
這個去告訴.
就是一個分享.
這個分享就是這個人和人之間.
一個最美最美的關係.
當我們吃那些東西.
很好吃的時候.
我們會怎樣?.
不是藏在袋子裡自己繼續吃.
而是當我們和人的關係好的時候.
我們就會和人分享.
剛才回來有個姐妹問我.
你吃了早餐沒有?.
有多比味.
一個愛的分享.
又或者當我們知道.
哪裡的東西又便宜又好.
很好吃的時候.
我們就說.
你去吧.
我和你訂位一起去吃.
這個就是一個愛的分享.
所以當我們經歷過.
在我們生命裡面.
我們曾經迷茫.
曾經無力.
曾經孤單.
但復活的耶穌就幫過我們.
以致我們生命不再是污濁濁.
不再是污味合遂.
我們就告訴別人.
我生命可以今天變成這樣.
是因著復活耶穌基督的大能.
我們知道.
在我們這個錦繡附近.
將會有兩座臨時房屋.
將會落成.
有些什麼感受?.
有些人就和我說.
我們在這裡都是中產.

$^{441}$但是這兩座臨時房屋.
他們好像較為基層.
會不會搞亂我們的社區.
帶給我們很多麻煩.
沒錯.
這些都是一些.
等候上公屋的家庭.
是較為基層的.
但他們不等於麻煩.
不等於令人害怕.
當我們有這種害怕.
這樣想的時候.
我們要記得一件事.
耶穌基督是復活的主.
祂來到這個世界上的時候.
就要將人從軟弱.
從不足夠.
從很多懼怕的黑暗地裡.
將我們帶到一個.
人愛豐盛滿足的永恆的國度.
耶穌基督來到地上的使命.
亦都是攀爬給我們.
所有門徒.
所有信徒的使命.
耶穌基督有這個能力.
有這個權柄.
能夠將人枯乾的生命環境.
扭轉過來.
能夠使我們變得勇敢,剛強.
充滿這個能力.
我們信不信呢?.
有一位姐妹.
她已經回到天父家了.
她有一位親戚進了醫院.
然後又進了老人院的時候.
她很疼愛這位親戚.
不久就探望她.
不過她說.
每次探望她都很辛苦.
因為大家都知道.

$^{481}$醫院或者老人院.
那一股味道.
有時候真的會很辛苦.
她有時候真的試過一兩次.
她真的吐了.
要衝進洗手間吐.
她心裡跟上帝祈禱.
她說求上帝幫助.
因為她真的很疼愛她這個親戚.
她想多點去看她.
但這樣真的不行.
她求神幫助她.
她已經穿過了.
一聞到這種味道.
她很深的那種警方.
只要我們有信心去感求.
帶著使命去感求的時候.
上帝就是這麼奇妙的.
祂的權能就是這麼厲害的.
從此以後這位姐妹.
她不再覺得有什麼難聞.
不再覺得有什麼難頂.
她的姐妹到最後.
她跟我拍檔去醫院.
去老人院那裡探訪.
有時候我都覺得很辛苦.
她就說有味道嗎.
我剛才都不覺得.
我想上帝就是這麼奇妙.
只要我們信.
只要我們求的時候.
上帝就可以幫助我們.
衝破我們的擔心.
各位姐妹.
有時候我們真的忘記了.
我們的上帝.
我們的耶穌那份大能.
祂衝破了死亡的權勢.
祂那份復活的大能有多厲害.
所以有時候我們只看著.

$^{521}$我們自己的困難.
我們自己的限制.
我們對其他人的不信任.
以致我們的信仰生活.
我們不敢向前.
總覺得死了.
我怕什麼.
哎呀 我不行的.
我不敢.
哎呀 我不行.
我的性格不是這樣的.
以致我經常看著那些怕.
看著那些擔憂的時候.
我們經歷不到.
上帝復活的大能.
久而久之.
我們的屬靈生命.
就好像一趟的死水.
不再蠕動.
在這個傳福音.
拓展天國的事情上.
我們知道.
頭腦上知道.
是上帝給我們的大使命.
耶穌基督升天的時候.
吩咐我們的.
我們知道.
很明白.
也都很識背那兩節的聖經.
但是我們有時.
我們不再有這個行動.
我們求主憐憫.
我們那時候小學,中學.
寫這個紀念冊.
我們就說.
萬如逆水行舟.
不進則退.
我們的信仰都是.
當我們不進步.
那些屬靈的生命.

$^{561}$不進步.
不讀經.
不祈禱.
不傳福音.
不以信心來向前走一步的時候.
是不進則退.
求主幫助我們.
再看看.
這兩座臨時的房屋.
我們用上帝的角度.
來看一看.
哇.
這麼多人.
突然間在我們附近.
到處出現的時候.
這個很難得的.
傳福音的機會.
耶穌說過.
讓小孩子到我這裡來.
不要禁止他們.
因為天國是他們的.
是屬於他們的.
我們就將他們.
配合我們現在的境況.
讓所有人.
都要有機會.
來去聽耶穌的福音.
不要禁止他們.
因為上帝的福音.
要傳過給世上高.
所有人.
天父上帝的心願.
是不願一人沉淪.
乃願人人悔改得救.
你和我.
都已經上了主耶穌基督.
救贖的之名.
所以我們時時都要去.
並且要快去.
告訴周圍其他的人.

$^{601}$告訴他們.
耶穌基督.
就是那位復活的主.
祂幫助我們.
跨過很多的恐懼.
很多的自卑.
很多的自憐.
很多心裡面的枷鎖.
仇恨.
暗的罪惡的綑綁.
祂絕對有能力.
來釋放我們每一個人.
將我們帶到去.
永恆榮耀的國度.
耶穌基督升天的時候.
祂再一次吩咐門徒.
在今天也都吩咐你和我.
我們一起讀這兩節的經文.
準備起.
所以你們要去.
使萬民作我的門徒.
奉父子聖靈的名.
給他們施洗.
凡我所吩咐你們的.
都教導他們遵守.
看那我天天與你們同在.
直到世代的終結.
使萬民作我的門徒.
是主耶穌基督.
給所有信徒的使命.
我們不能夠說.
「安感動」來去行.
在今天.
是主耶穌基督復活的日子.
我們有沒有一些實際的行動.
去實踐主的大使命呢?.
復活主的能力.
是奇妙莫測的.
靠著主復活的大能的時候.
我們每一位弟兄姊妹.

$^{641}$都要走出我們的按書區.
一起去活出.
在主耶穌基督的裡面.
我們的生命是自由.
是釋放.
是豐盛的.
所以弟兄姊妹.
求主復活的大能.
充滿著我們每一位.
我們現在就請.
常青和常恩的師班.
去唱出「因他活著」.
去宣告復活的主耶穌.
會帶給人剛強喜樂.
有意義的人生.
唱到第三節的時候.
我會出來邀請大家.
一起喜樂.
加入他們的行列.
來宣告主耶穌基督復活的大能.
我們在騰起南寮一起唱第二生.
我們同心禱告.
主耶穌基督我們要Allelujah.
稱讚主你復活的大能.
你不單止衝破了死亡的權勢.
我們多謝你幫我們去恢復了.
跟上帝的關係.
恢復了跟人的關係.
所以我們多謝你.
也恢復了我們跟自己的關係.
我們多謝你.
祝福我們的生命.
是一個時時有主.
你活水在當中.
一路的湧流.
讓我們所說的.
所想的 所恨的.
我們都能夠活出主你的美善.
活出主你復活的大能.
祝福 祝福我們看到.

$^{681}$我們真的周圍人的需要.
我們經歷過主恩的滋味.
主就願意.
我們願意快的 去的.
去講給周圍其他的人聽.
就是我們心靈裡面有什麼的懼怕.
求主幫助我們釋放出來.
主就願意我們每一個站在你面前的人.
我們的生命是滿有耶穌基督復活的大能充滿.
以致我們的生命是有朝氣的.
是有活力的.
是要見證主耶穌基督你的榮美和大能的.
我們就求主你使用我們一生.
多謝你垂聽禱告.
是奉靠主基督的聖名.
阿們.
\newpage



\section{創世記 30:14-16}
\label{sec:_PqGg3XbEA0}
\textbf{2024-04-06 晚堂《缺乏了愛的家庭–雅各,利亞,拉結》創世記 30\_14-16 講員:鄧泰康傳道}
\newline
\newline
連結: \href{https://youtube.com/watch?v=-PqGg3XbEA0}{\texttt{https://youtube.com/watch?v=-PqGg3XbEA0}} ~~~~ 語音日期: 2024-04-09
\newline
\newline
\hyperref[sec:ez2giRGwvx8]{\small{< < < PREV SERMON < < <}}
~
\hyperref[sec:index_chronic]{\small{[返順時目]}}
~
\hyperref[sec:index_scriptual]{\small{[返順卷目]}}
~
\hyperref[sec:TfOT_yj_ii8]{\small{> > > NEXT SERMON > > >}}
\newline
\newline
創世記 30:14-16
\newline
\begin{longtable}{cl}
\hline
\hline
章節 & 經文 (和合本修訂版)\\
\hline
30:14 & \begin{tabularx}{0.7\textwidth}{X} 收割麥子的時候，呂便到田裡去，找到曼陀羅草，就拿給他的母親利亞。拉結對利亞說：「請你給我一些你兒子的曼陀羅草吧。」 \end{tabularx} \\ \\ \relax
30:15 & \begin{tabularx}{0.7\textwidth}{X} 利亞對她說：「你奪走了我的丈夫還是小事嗎？你還要奪取我兒子的曼陀羅草嗎？」拉結說：「今夜他可以與你同寢，來交換你兒子的曼陀羅草。」 \end{tabularx} \\ \\ \relax
30:16 & \begin{tabularx}{0.7\textwidth}{X} 到了晚上，雅各從田裡回來，利亞出來迎接他，說：「你要與我同寢，因為我真的用我兒子的曼陀羅草把你雇下了。」那一夜，雅各就與她同寢。 \end{tabularx} \\ \\
[1ex]
\hline
\hline
\end{longtable}
$^{1}$今晚同大家講一個《異屬能詳》聖經裡面大家一定會聽過的故事.
這個故事裡面特別的地方是什麼呢.
其實我想給大家看到的是一路延伸下去的一個家庭問題.
之前和大家分享過雅各.
雅各他是因為爸爸媽媽或者各自的偏心.
而後來出現了一個情況就是和他的兄弟不和.
所以後來才會發展到今天這個故事.
在一個家庭裡面如果你問什麼是最重要.
絕對不是錢,經濟,其他.
絕對是對於特別的基督徒來說很清楚.
是愛.
一個沒有愛的家庭是多麼可怕呢.
今天透過這段經文就和大家講.
這個家庭裡面有三個很主要的角色.
一個是利雅,一個是雅各,一個是拉傑.
你就會發覺我每一個人都做了一個主角.
雅各是一個…利雅是一個怎樣的人呢.
我稱之為他叫做失去愛的妻子.
大概很快地講一講他們的故事.
當雅各被追殺.
去到他舅父家的時候.
如果大家還記得那段經文其實挺有趣的.
就是他去到井旁.
然後他就去問那些人.
你們認不認識我舅父拉班.
然後他說認識.
他的女兒還沒來.
拉傑就帶著羊來.
那裡有一個很有趣的圖畫.
就是當時因為在這些曠野的地方.
水源是非常重要的.
所以每一個水源裡面.
他一定會有一個很大的大石.
去蓋住井口.
這個大石原則上一定要有幾個牧羊人.
一起去搬才能搬走.
但是在那一刻當他見到拉傑的時候.
其實已經有點眉目了.
他做了一個很誇張的舉動.
就是當那些牧羊人說現在還沒夠時間.

$^{41}$等一等吧.
然後雅各就走去自己一個.
搬開了井的那塊石頭.
就叫拉傑去餵羊喝水.
我不知道大家有沒有看到.
作為一個少女.
見到前面一個這麼大力威猛的男士.
一個人舉起這塊石頭.
然後還要餵我.
你明白那種的面子嗎.
多厲害.
然後又出現了另一個圖畫.
雅各見到拉傑之後.
一見到他就大聲痛哭.
就跟他說他的情況.
拉傑就立刻帶他去他的舅父拉班.
故事是這樣開始的.
就是說其實你發覺.
他們第一眼見到的時候.
有些像當時的以撒和李百嘉.
你會發覺他們一眼就認得.
有些惡意在裡面.
所以當雅各幫他拉班工作的時候.
拉班就非常客氣.
就跟他的兒子說.
你都幫我做了一個月.
這樣的情況下.
我沒理由要你白做的.
那你要什麼呢.
大家看到的時候都不錯.
你看看後面的拉班.
絕對沒有可能是這樣的.
又騙他的兒子的東西.
其實他是一個很奸詐的人.
就是說這個月裡面.
他絕對見到雅各對他的女兒是有心的.
他預了雅各怎麼回答的.
所以你就會發覺雅各也很不客氣.
他說我什麼都不要.
你只要把女兒嫁給我就可以了.

$^{81}$他說好呀.
然後雅各自己提出.
因為娶他的緣故.
我願意跟你做七年.
聖經有一個主角很有趣.
因為他愛拉基的緣故.
所以他覺得七年是幾日.
大家有沒有曾經愛過一個人.
愛到這樣的地步.
你可以見到雅各對這個女子的一心一意.
但是整件事好玩的地方就是.
這件事在香港是不會發生的.
因為香港是有光污染的.
就是說你走到街上.
就算沒有任何燈光也好.
你都會認得那個人.
即是很遠也好.
我一晚看過去.
就算沒有燈.
我都認得那個是Alan.
你會認得那個人.
但是你曾經去過很黑的地方.
是身受不濟五指.
大家有沒有試過去過這些地方睡覺.
一關燈之後你什麼都看不到.
七年之後.
當他那晚去洞房的時候.
他的舅父拉班就將姐姐推給雅各.
然後洞完房.
他才說你為什麼騙我.
然後舅父就說了.
不行的.
我們的規矩就是.
姐姐不嫁妹妹不嫁.
然後因為這個緣故.
他為了要再娶這個妹妹.
又跟他做了七年.
在裡面的時候.
其實聖經裡面做了很多很有趣的註腳.
當時他介紹這兩姐妹出場的時候.

$^{121}$是怎樣的呢.
29章第17字.
利亞的眼睛沒有神器.
拉潔卻生得美貌盡秀.
這個是和諧本.
另外一個版本是怎樣的呢.
利亞雙目無神.
拉潔卻體態嬌美.
容貌可人.
你見到那個比較嗎.
我想告訴大家.
利亞是怎樣的呢.
她是先天書舌.
所以雅各一眼就喜歡美女.
他完全對這個姐姐一點興趣都沒有.
所以這樣的情況下.
其實聖經形容她是什麼呢.
她是失寵.
甚至乎有些譯本告訴你.
她是避恨.
一個女人嫁了給她的老公.
她對我一點愛意都沒有.
甚至乎我感覺到她恨我.
所以後來就出現在經文裡面.
你會發覺這個女人.
因為她一直都渴望丈夫的愛.
所以令到讓子女生子.
成為了爭寵的工具.
你從她的子女的改名就見到.
其實很可怕的地方是怎樣的呢.
你要記住他們是有上帝的人.
利亞她生第一個子叫做女辯.
女辯的意思就是我有兒子.
西冕第二個兒子的意思就是聽見.
看不看得到啊.
上帝聽見我的祈求.
利美的意思就是聯合.
但是其實有另外一個譯本譯為.
依戀或者戀慕.
就是說我現在跟你生了三個了.

$^{161}$你還不愛我.
第四個尤大讚美.
表面上是說很好.
我要讚美上帝.
其實不是的.
你看看妹.
我生了第四個了.
你看看上帝多麼疼我.
我要讚美他.
第五個以撒迦獎賞.
是因為她做了我的注腳.
因為你看看我把我的僕女.
我的婢女都送給我老公了.
所以上帝獎賞我.
你見不見到這個扭曲.
第六個西部倫叫尊榮.
我幫她生了六個了.
現在我的老公一定會尊重我.
她生第一個的時候.
她說因為我有兒子了.
所以我丈夫一定會愛我.
生第六個的時候.
我現在丈夫一定會尊重我.
我不知道特別在座.
作為姐妹的你.
你覺不覺得這個女人很可憐.
她一世裡面.
她就是圍著丈夫在轉.
她人生裡面的目標.
只有一個.
我就是要我丈夫愛我.
所以你會發覺.
她對她的兒子.
是一點都不是愛的.
她的兒子只是成為了.
她達到一些目標的工具.
你可以想像.
你很難見到一個這樣組合的情況.
這個媽媽和這些子女.
會得到幸福是不可能的.

$^{201}$其實今天有很多作父母的.
我們一不小心.
其實你和利雅是一樣的.
大家有沒有經歷過.
如果我們都有子女的話.
有時候特別.
到我們子女一直讀書的時候.
讀哪間學校.
然後人們說你的子女考哪間.
你考到一間很心儀的.
大家都很喜歡的時間.
很開心.
她考到哪間了.
然後時間如果她考不到.
心儀的那間又如何.
你不要問.
再大一點.
小學如是 中學如是.
去到大學.
你知道大學的學位都不容易.
你的子女能夠讀到大學的時候.
你知道我們人多虛偽.
然後人們問你的時候.
你的子女怎樣.
我們這些是蠢的兒子.
不過剛好被她考到大學.
其實你是開心的飛起.
因為她能夠幫你拿到面子.
大學生.
但當你的子女考不到的時候.
可能倒過來的.
是成為你丟架.
人們問你的時候.
讀成怎樣.
不要問了.
都不是讀書.
但我們要很小心.
今天我們這些的心態.
其實是很容易出現.
我們喜歡和別人比較.

$^{241}$慢慢不知不覺之間.
我們子女所做的事.
成為我們今天.
在別人眼中的時候.
我是一個有榮譽的人.
還是我是一個丟架的人.
所以不知不覺之間.
你亦都會將這些事.
影響你的小朋友.
美其名在口中跟他說.
你用心讀書.
讀好書的時候.
對你將來會有好的發展.
但其實實際裡面.
我們的心是否真的這樣想.
還是我們更在乎.
究竟他帶著什麼光環給我們.
所以弟兄姊妹.
你不要小看這件事.
其實孩子這些成長.
是非常受影響的.
我以前去做輔導的時候.
我接觸過很多成長的年青人.
一些可能出來做初職的人.
其實在心裡面.
他一生都被這些環繞著.
究竟我在爸爸媽媽的眼中.
我是什麼.
他是不是真的愛我.
其實如果這個家庭裡面.
沒有這份真正爸爸媽媽.
對子女的愛的話.
這個小朋友是不會在一個.
很健康的環境當中長大的.
我曾經去接觸過一些家庭.
三兄弟姐妹.
三個都這麼厲害.
一個是醫生.
一個是做大生意.
另外一個.

$^{281}$我忘記了.
總之有兩個都是做大生意的.
真的很厲害.
其實妹妹已經不假了.
她已經是香港裡面.
最出名連鎖集團的合夥人.
但這個妹妹永遠都在家裡面.
面對著哥哥和姐姐.
她都抬不起頭.
因為她最小.
讀書不是最厲害.
從小到大的時候.
爸爸媽媽都覺得你是最蠢的.
所以每次我有機會接觸到.
這個姐妹的時候.
我都跟她說.
說真的.
這是比較出來的.
我給你一個凳子.
你記得多聰明.
我想強調的是.
但她今天都揮之不去.
她已經是幾個小朋友的媽媽.
但這件事.
她每一次回到家裡.
每一次吃飯的時候.
她仍然永遠都是那一個.
幾歲的小朋友.
在媽媽眼中.
她永遠都是家裡面最蠢的那個.
弟兄姊妹我要問的就是.
在這裡.
其實我們看見上帝.
祂是一個看顧人.
扶持人.
幫助困苦人的神.
你會發覺.
其實祂是看見利雅的苦情.
所以祂出手幫助她.
但是利雅看不看到.

$^{321}$利雅她完全看不到.
上帝的出手.
只是成為了.
祂對其他人耀武揚威的工具.
弟兄姊妹.
今天其實我們一個屬於上帝的人.
我們究竟為誰而活呢?.
我們想討好.
讓祂喜悅的對象.
究竟是誰呢?.
是你的配偶?.
是你的子女?.
還是其他人?.
其實我們今天不知不覺之間.
我們就被這件事困住了我們的心靈.
很多時候我們爭奪不到.
我們再努力都好.
我們只想追求到那件事的關注.
就好像我一樣.
我經常說我小時候.
我不是一個讀書的人.
因為我記性很差.
我到今天的時候.
大家覺得我是藉口.
我都是這樣說的.
我真的記得大家的名字.
頭都痛了.
所以我很佩服陳牧師.
因為牧師真的.
每個人都記得大家的名字.
站在門口的時候.
一叫就叫得出.
但是我差到.
現在年輕人在唱完歌之後.
每次問他之後.
下次再見到他.
現在他們懂得玩我了.
得知到我叫什麼名字.
我讀書很差.
我記憶力很差.

$^{361}$所以我從小到大的時候.
我就不放在讀書上.
我就放在其他那裡.
就是什麼呢.
很多表演.
就是故事演講比賽.
文歌比賽.
幫學校做這個做那個.
目的我只有一個.
我家裡有七兄弟姐妹.
你會明白一個七兄弟姐妹的家.
媽媽哪有那麼多時間理你.
你明白我的意思嗎.
接著的情況.
我只是做那麼多事.
我讀書不厲害.
讀書不厲害的時候.
我就做其他事.
想拿獎.
想拿獎做什麼.
吸引媽媽的注意.
我好記得.
我中學的時候.
我在學校不知道拿了什麼獎.
我好開心.
我回家跟媽媽說.
媽媽我拿獎了.
你會不會來看我頒獎.
我媽媽怎麼跟我說.
你花多點時間去讀書好過了.
搞一些古靈精怪的東西.
所以我從小到大的時候.
我媽媽是沒有來看過我拿獎的.
我到今天都記得這件事.
是證明了這件事.
在我心目中有多麼重要.
今天我不再被困住了.
因為我好知道.
一直看著我的那個是上帝.
上帝一直由我成長.

$^{401}$一直看著我.
我無論在任何環境當中.
祂都照顧我.
我以前很生我爸爸的氣.
因為我爸爸只會上班,下班和打人.
我完全不懂得欣賞我爸爸對家庭的盡力.
讓我重新去改變.
上帝告訴我.
讓我看見爸爸對家庭的忠心.
我們如果能夠在這裡脫離困住我們的東西.
我們能夠很清楚今天為上帝而活.
上帝祂才是我們追求的目標.
因為祂從來都不會變.
我們就有出路.
第二個角色是一個偏心的丈夫 雅各.
偏心是雅各的特色.
你會發覺這個人的偏心是一直維持到最後的.
的確我們要很公平地說.
他是被迫和無奈地接受了利雅做他的老婆.
所以你可以說他不喜歡是不明白的.
但是我想強調.
你記住屬於上帝的女兒就是這樣.
你除了個人的喜惡之外.
你也需要尊重盟約.
結婚是一個婚約.
婚約是有責任的.
而且當時雅各是答應的.
大家記得嗎.
雖然是他舅父冤枉他的.
但反過來是他答應的.
好啊 我娶了他.
你既然答應了他的時候.
你就要遵守這個婚約的裡面.
利雅作為一個太太.
其實她應該得到一個最基本的尊重和照顧.
但是很顯然是沒有.
其實你從剛才那段時間.
大家去看的時候.
如果不是真是神的恩典.
她不可能生到六個子女.

$^{441}$我可以告訴大家.
絕對是.
你看到後來的時間.
她就要用封茄.
在和合本叫封茄.
去換取和她老公可以同房.
即是她老公從來都不找她.
你明白嗎.
就好像皇帝一樣.
後妃其實要看皇帝何時來.
我才有機會生一個王子.
但問題是你從來都不來找我.
我怎麼生王子.
但上帝好玩的地方就是這樣.
我找你這麼少你都生六個.
所以你會發覺一件事.
其實這件事是丈夫和太太之間.
很基本應該有的尊重和責任.
但你會見到.
景然利雅要用封茄.
才可以獲得和丈夫同房的機會.
你會發覺這個人.
其實才是整個家庭裡面.
令整個家庭破碎的始作俑者.
由這裡她的偏心.
大家記不記得.
弱失彩衣的事件.
哪有一個老爸.
我形容這麼瘋狂.
你有11個子女.
你在11個子女面前.
你只給一件彩衣給最小的兒子.
還擺明因為他是拉結的兒子.
所以我特別疼愛他.
所以是你令到最初出現的弱失.
這麼沒腦.
還要發夢和哥哥去演戲.
我發了個夢.
我見到你將來要拜我.
你給我做.

$^{481}$我都想殺你.
你明白我的意思嗎.
你不單止搶走.
奪走爸爸對你的偏愛.
你還在我面前演戲.
兄弟之間是不會有和睦的關係.
所以到最後就發生了.
弱失差點被殺.
但上帝的保守才被賣的結果.
一個丈夫應該是一個家庭的頭.
應該是帶領著這個家庭.
走向上帝的道路上.
他應該是妻子的頭.
是兒女的榜樣.
其實作為一個丈夫.
他在一個家庭裡面.
一個家庭到底幸不幸福.
丈夫的身份,角色是非常重要的.
所以我真的很盼望.
今天每一位我們作為自己的丈夫的弟兄.
我們好好在神的面前.
我們真的要回到神的裡面.
我們去問上帝.
我怎樣才能將這個家帶到你的面前.
讓我整個家都在上帝的恩典當中.
我們會不會有缺乏呢.
有,我們有缺失,我們有缺乏.
但我想大家記得.
我們每一個人都有上帝.
無論最後的境況.
只需要你不放手的話.
你一定能夠做到這個頭的角色.
只需要你堅持.
第三個.
一個不辱的太太拉潔.
其實她一樣.
你會發覺在整個的創世故事裡面.
幾代當中.
同樣都出現了差不多的事情.
但你會發覺.

$^{521}$她面對不辱的時候.
她並沒有去尋找上帝的幫助.
而是她反而選擇了他們祖先撒拉的路.
去尋找代母.
如果大家回看聖經裡面的時候.
聖經講得很清楚.
她去和雅各投訴的時候.
原因是什麼呢.
只因為我嫉妒姐姐.
我沒有兒子生.
如果雅各你不和我生一個.
我死給你看.
整個經文是這樣寫的.
但雅各已經提醒她.
聖經說她發怒.
雅各和她說.
她說你怪我嗎.
生孩子是誰決定的.
上帝決定的.
其實她提醒了她.
生孩子這件事是上帝的事.
但她並沒有去改變她的心意.
她沒有去尋求上帝的心意.
她用自己的方法.
為什麼她用這個方法呢.
因為她只是將她的焦點.
放在自己的榮辱上.
她不想輸給姐姐.
她一點都不愛雅各.
一點都不愛家人.
其實你看回拉傑.
是一個很自私的人.
雖然她丈夫很愛她.
為了她可以14年.
但在這裡我們完全看不到.
她有其他的東西.
她為丈夫去回饋.
她沒有.
她只是去愛自己.
她完全想生孩子.

$^{561}$她不會考慮上帝的心意.
我只有一心一意.
我要贏我姐姐.
所以她可以用風檐出賣她老公.
換成今天聽得明白的話.
即是說有一個人給你幾千萬.
你就可以賣了她老公.
我想買一層樓.
我可以達到我的目的.
老公不重要.
換一個.
所以她就能夠出賣她丈夫.
你亦都會發覺.
當他們離開舅父家的時候.
他做了一件事.
他偷了他爸爸的偶像.
他的神像.
如果看回解經家的時候.
應該是當時因為這個神像.
是代表著平安.
是代表著風調雨順.
他想他未來都能夠風調雨順.
所以他就偷了這個神像.
但是他差一點.
就令到整族人都有殺身之禍.
他的世界裡面.
只有他自己.
頂至每一個家庭當中.
一個人我只是愛自己.
我將自己放到很大.
我完全不用顧及其他人的感受.
我喜歡怎樣就怎樣.
我只是追求我自己的東西.
這個家庭裡面.
怎會有幸福.
其實你會發覺.
當拉傑去生弱勢的時候.
作者做了一個主角.
創世記三十章二十二節.
神顧念拉傑.

$^{601}$應允了他.
所以他就生了弱勢.
其實很簡單.
原來上帝就是要你回轉.
你找我吧.
你不要用你自己的方法.
上帝就是要我們每一個人.
都能夠去找回他.
其實為何我們作為基督徒.
在港台上我們經常都很強調.
弟兄姊妹.
我們每一個人都需要去走回.
去找上帝的路.
去問回上帝的意思.
我們每一個人所做的事.
究竟合不合乎你的意思呢.
因為這個將會影響在未來當中.
無論你自己個人.
你的家庭.
和你整個的群眾.
我們都需要一起去尋求上帝的意思.
整件事由拉傑.
利雅至拉傑.
雅各的幾段.
三章裡面大概的故事.
其實只是環繞一個主題.
就是生孩子.
裡面總共生了十一個人.
但是最諷刺的一點.
生孩子永遠都只是上帝的主權.
現在這裡就是人用盡自己的方法.
去搶奪上帝的主權.
到最後換來的是甚麼呢.
家庭的破碎.
所以在這個家裡.
三個都輸.
孩子都輸.
所以弟兄姊妹.
我好盼望今晚透過這個.
好簡單的.

$^{641}$好直接的故事.
和大家一起重新再思考.
我們自己的家庭.
沒錯.
我們在哪個家庭裡面出生.
我們無得選擇.
哪個是你的爸爸媽媽.
你無得選擇.
哪個是你的子女.
你無得選擇.
但是我昨晚.
我去教傳道書的時候.
我和弟兄姊妹說了一樣東西.
我覺得是非常重要.
就是我們今天所有的事.
都在處境當中.
但是我們第一.
在處境裡面.
你可以決定你怎樣面對這個處境.
這個處境再差.
你都可以用好的方法去面對.
而另外更加重要的是甚麼呢.
請你不要自己親手去創造一個.
壞的處境.
亦都請你好好盡量的.
去創造一個好的處境.
你明白我的意思嗎.
今日我們的小朋友.
就好像十隻手指一樣.
有長有短.
有些聰明些.
有些沒那麼聰明.
有些可愛些.
有些沒那麼可愛.
有些頂撞些.
有些怪些.
甚麼人都有.
但其實我們做爸爸媽媽的.
我們唯一可以做的是做甚麼.
你創造一個幸福家庭的環境給他們.

$^{681}$這個就是我們能夠做得到.
你不要自己在家庭當中製造麻煩.
你不要自己親手在家庭當中.
親手去做很多破壞的事.
為了達到我們的目的.
為了我們自己很個人私人的一些理由.
而我們不顧其他人.
而去做一些事情.
最簡單的例子.
最直接.
今日有很多人.
未信耶穌.
都很喜歡賭錢.
最初的時候的理由都是甚麼.
做得那麼辛苦.
玩玩嘛.
然後讓自己可以去放鬆一下.
最初打麻將.
買下六合彩.
然後玩玩玩玩.
後來發現不是的.
澳門很過癮.
澳門的時間一鋪就五百.
嘩 夠刺激.
是吧.
然後我為了滿足我自己.
這些東西的時候.
這種所謂的得到賺到額外的橫財.
又或者是反過來的.
這種薄煞的興奮的感覺.
然後你把你全家人.
就陷入一個危機裡面.
是你自己創造出來的.
而今日最簡單.
最直接的就是.
不要讓你的家人.
讓自己陷入危機當中.
在哪一方面都是.
做的每一樣事情的時候.
去考慮上帝是怎麼想的.

$^{721}$你做這件事的時候.
你的家人會受甚麼影響.
我們能夠這樣做的時候.
我相信我們家庭.
能夠得到幸福的機會會更大.
當然最直接簡單的方法.
求上帝讓我們所有的家庭成員.
都在上帝的裡面.
我們每一個人都在上帝.
他自己尋求旨意的道路上.
我們這個家就一定會幸福.
求神幫助我們.
我們要同心企圖.
我們求上帝你親自帶領我們.
讓我們去看見.
有很多不理想的事情.
讓我們今天千萬不要重蹈覆轍.
因為我們能夠去看見.
在神的恩典裡面.
我們怎樣能夠找到幸福.
怎樣會讓我們陷入泥沼當中.
我求上帝你親自保守.
帶領著每一個弟兄姊妹.
很願意我們的家庭.
因為我們認識你.
去得到恩典和幸福.
我們也在這裡.
如果家裡有家人未信耶穌.
我們求你幫助他們.
讓我們尋找到你.
讓我們能夠一起進入幸福當中.
將一切的我們交託.
祈禱奉靠主耶穌.
得勝名自祈求.
阿們.
好嗎?.
\newpage



\section{路加福音 24:36-43}
\label{sec:TfOT_yj_ii8}
\textbf{2024-04-07 早堂《轉化我們情緒困擾的耶穌》路加福音 24\_36-43 講員:區祥江博士}
\newline
\newline
連結: \href{https://youtube.com/watch?v=TfOT_yj-ii8}{\texttt{https://youtube.com/watch?v=TfOT\_yj-ii8}} ~~~~ 語音日期: 2024-04-10
\newline
\newline
\hyperref[sec:_PqGg3XbEA0]{\small{< < < PREV SERMON < < <}}
~
\hyperref[sec:index_chronic]{\small{[返順時目]}}
~
\hyperref[sec:index_scriptual]{\small{[返順卷目]}}
~
\hyperref[sec:CWxLFOO62pc]{\small{> > > NEXT SERMON > > >}}
\newline
\newline
路加福音 24:36-43
\newline
\begin{longtable}{cl}
\hline
\hline
章節 & 經文 (和合本修訂版)\\
\hline
24:36 & \begin{tabularx}{0.7\textwidth}{X} 正說這些話的時候，耶穌親自站在他們當中，說：「願你們平安！」 \end{tabularx} \\ \\ \relax
24:37 & \begin{tabularx}{0.7\textwidth}{X} 他們卻驚慌害怕，以為所看見的是魂。 \end{tabularx} \\ \\ \relax
24:38 & \begin{tabularx}{0.7\textwidth}{X} 耶穌對他們說：「你們為甚麼驚恐不安？為甚麼心裡起疑惑呢？ \end{tabularx} \\ \\ \relax
24:39 & \begin{tabularx}{0.7\textwidth}{X} 你們看我的手和我的腳，就知道實在是我了。摸摸我，看，因為魂無骨無肉，你們看，我是有的。」 \end{tabularx} \\ \\ \relax
24:40 & \begin{tabularx}{0.7\textwidth}{X} 說了這話，他就把手和腳給他們看。 \end{tabularx} \\ \\ \relax
24:41 & \begin{tabularx}{0.7\textwidth}{X} 他們還在又驚又喜、不敢相信的時候，耶穌對他們說：「你們這裡有甚麼吃的沒有？」 \end{tabularx} \\ \\ \relax
24:42 & \begin{tabularx}{0.7\textwidth}{X} 他們給了他一片烤魚， \end{tabularx} \\ \\ \relax
24:43 & \begin{tabularx}{0.7\textwidth}{X} 他接過來，在他們面前吃了。 \end{tabularx} \\ \\
[1ex]
\hline
\hline
\end{longtable}
$^{1}$第一天早上.
剛剛過了復活節.
我今天選了耶穌一段故事.
是他復活之後發生的.
我們都可以一起去看看.
其實主耶穌不單止復活.
祂更加活在我們的生命裡面.
在我們的心裡面.
甚至在我們的情緒困擾裡面.
我不知道大家怎樣看待情緒.
情緒其實是我們人.
神創造我們的時候給我們一份禮物.
雖然很多時候我們說我們有情緒困擾.
但情緒其實是反映我們當時的狀況.
是一個很好的資訊給我們了解我們自己.
很多時候我們都會傾向將情緒壓抑.
導致有時候都會有一些情緒困擾.
我不知道大家有沒有想過.
其實我們的天父.
我們的主耶穌基督.
其實祂是很著緊我們的情緒.
祂會幫我們去經歷.
渡過 甚或轉化我們的情緒.
今天我們一起來看那段經文.
我們順著經文去看一看.
這裡說 正說這些話的時候.
要說一說上文下理.
正說這些話 當時耶穌是復活了.
那些人在說話呢.
其實有兩個門徒.
如果大家記得有一段很出名的故事.
就叫做「以孟思路上」.
有兩個門徒.
因為知道主耶穌基督死了.
復活 有人傳 但未見到.
所以他們一路走 在路上大家在談那件事.
耶穌臨近他們 跟他們走了一段路.
還跟他們一起吃了一頓手聖餐.
當兩個門徒突然間.
主耶穌就突然間不見了.

$^{41}$他們經過了一段跟主耶穌同行的時間.
他們心裡很火熱.
所以趕著回去跟門徒談那件事.
所以背景是這樣.
正說這些話 即是說那兩個門徒.
好像經歷了一件很特別的事.
很想我們叫做Debriefing.
要說回那件事.
剛剛在說的時候.
耶穌就親自站在他們當中.
然後就說願你們平安.
其實願你們平安.
也是猶太人慣常的問候語.
Shalom.
其實在這個世界這麼多戰爭 地震.
這麼多的挑戰.
其實問平安 或者祝福遇到的人平安.
是很重要的事.
我不知道今天你和旁邊的弟兄姊妹.
有沒有說平安.
其實主耶穌說平安這件事.
原來對於門徒來說也有一個特別的意思.
因為當時的門徒不是太平安.
為什麼這樣說呢.
我們看37節.
他們卻驚慌害怕.
以為所看見的是雲.
這裡就是說他們即時的反應.
但其實當時的門徒.
我們看福音書.
有時候會看婦女福音.
約翰福音有另外一個記載.
約翰福音20章19節特別提到.
當時的門徒其實有一點風聲鶴唳.
耶穌被人抓了然後釘了十架.
連彼得也有三次不認耶穌.
所以他們其實是一個惶恐的狀況下.
所以約翰福音20章19節特別提門就關了.
即是說耶穌出現之前.
門徒是圍爐取暖.

$^{81}$或者在一個不是公開的地方.
關上門的地方去談那件事情.
為什麼要關上門呢.
其實就害怕有恐懼.
所以主耶穌一開始.
願你們平安這句話.
其實是說我們主耶穌基督.
除了是一個慣常的問候語.
我覺得他是很知道門徒的狀態.
所以說願你們平安.
因為門徒真的不是太平安.
不知道大家有沒有想過.
其實我們的天父.
我們的主耶穌基督.
或者聖靈三位一體的神.
其實是很清楚很明白我們內心的感受.
即是說在我們未需要表白之前.
其實他已經知道我們的情緒.
所以他主動地向我們問安.
這裡門徒的反應是有一點點令我們覺得奇怪的.
如果你看到你老師.
他跟你打招呼然後問你安.
你應該笑著歡迎他和他握手.
正常的反應是這樣.
但門徒的反應是驚慌害怕.
這裡有兩個情緒字眼.
驚慌和害怕是不同的.
驚慌就等於.
感受花園不會.
在外面如果那些雷槍在上面掉下來.
啪一聲.
驚慌的字眼其實是驚嚇.
彈起.
好像地震一樣.
令我們嚇一跳.
除了嚇一跳之外.
其實還有一種恐懼.
害怕就是一種恐懼的感受.
主耶穌基督其實見到門徒有這樣的反應.
其實是反映了實際他們正在經歷的一件事.

$^{121}$雖然看上去好像不是很明白.
但我剛才也給了大家一些線索.
就是說.
當時他們.
不知道會不會被人抓之外.
其實當時耶穌基督復活了.
有傳到他們面前.
但他們還沒有證實.
所以他們還沒有100\%肯定耶穌復活了.
所以看到耶穌的出現.
和問他們安.
這裡說了.
以為他所看見的是雲.
為什麼呢.
因為我剛才說了.
在約翰福音說門都關上.
所以耶穌的出現是突然間出現的.
穿牆過壁.
什麼是可以穿牆過壁的呢.
我們中國人的殭屍鬼魂.
穿牆過壁.
那耶穌是不是呢.
當然不是.
為什麼他可以呢.
因為他是復活的主.
是復活的神.
所以門徒的反應是可以理解的.
我不知道大家在你的生活裡.
有沒有遇到什麼情緒狀況.
你有沒有將這個情緒狀況.
告訴別人.
和別人分享.
讓別人了解你的狀況呢.
其實情緒呢.
我本身是教輔導和做輔導工作的.
很多時候都要聽別人說他的情緒.
但原來情緒我們說有不同的層次.
有些是比較外面一點.
有些是裡面一點.
這個驚慌害怕是一個表層一點的情緒.

$^{161}$真正的情緒是什麼呢.
其實主耶穌基督就是很清楚我們的狀況.
他幫我們處理我們內心的情緒.
我們一起看看.
38節.
耶穌對他們說.
你們為什麼驚恐不安.
為什麼心裡喜而華呢.
這裡有個驚字.
如果大家有留意.
我選了和修本.
去描述片段的情緒字眼.
有三對情緒字眼.
剛才是驚恐不安.
剛才是驚慌害怕.
這裡是驚恐不安和疑惑.
這兩個情緒.
其實就是說.
他們驚慌害怕背後.
有一些內心的不安.
當時的和修本是用受煩來形容的.
就是有些困擾的受煩.
還有心裡有疑惑.
就是有疑惑.
疑惑其實也是一種情緒的描述.
就是你很困惑,很煩惱,很疑惑.
其實耶穌為什麼要問這個問題呢.
其實是在幫助門徒處理他們底層內心的感受.
我舉一個例子.
可能大家會容易明白.
這裡也有些媽媽.
如果你的兒子擒槍.
即是快要掉下來.
你看到的時候.
你會怎樣反應呢.
會不會大聲喝他.
臭小子 快要掉下來.
當你呼喝大聲叫他的時候.
你的兒子感覺到你有什麼情緒呢.
生氣.

$^{201}$很生氣.
媽媽為什麼這麼生氣.
但生氣是不是你最底層的情緒呢.
不是.
是你的secondary emotion.
是一個表層的情緒.
你的primary.
你的內心真正的情緒是恐懼.
即是你怕你的兒子跌.
所以你就大聲喝止他.
給他感覺是生氣.
但其實你裡面是恐懼.
所以門徒他們驚慌害怕的背後.
是很受煩.
不知道耶穌是否真的復活了.
還有疑惑.
因為未親眼見到.
行不行.
其實在我們的人生裡面.
我們都容易流露一些外面的情緒.
但我們內心深處的情緒.
我們有沒有去接觸呢.
譬如我想.
如果我講男女之別的話.
其實女性是比較善於跟別人分享自己的情緒.
男人就比較少.
所以有些心理學的研究就說.
因為這件事所以女人長命一點.
因為她不會壓抑情緒.
而是很容易跟別人分享流露.
而男士其實特別.
我們感覺到脆弱的情感情緒.
我們是比較少表達的.
譬如很難承認自己害怕.
又或者很難承認自己有些事情不行.
所以我們說容易死撐.
或者覺得羞恥羞愧.
男人最容易表達的情緒是什麼.
憤怒和煩躁.
我們很容易被人感覺到這件事.

$^{241}$但譬如有個人他被同事在公司被人誣蔑.
其實他可以發火.
但他不可以說我受傷.
我很傷.
說自己很傷是很難的.
但你真正的情感就是受傷的情緒.
所以我們的主耶穌基督是很明白了解我們的實況.
所以他面對的門徒就是問這些問題.
有什麼受煩有什麼疑惑.
幫他們接觸自己裡面真正的感受是什麼.
你今天願意開放你自己.
打開你的心.
讓主耶穌基督進入你內心的深處.
不是停留在一個表面的流露.
而是將你最核心的感受都跟他說.
我們的主耶穌基督其實是很明白了解我們.
和他也很關心我們的情緒.
所以他問了這個問題之後.
其實就不是停留在那裡.
他還會用說話來安慰我們.
我們一起來看第39節.
39節請.
不如我們一起讀.
開始.
你們看我的手和我的腳.
就知道實在是我了.
摸摸我.
因為魂無骨無肉.
你們看我是有的.
這裡說了一番說話.
進來這間教會的時候都看到釘痕的手.
主耶穌基督就是將他的手和腳.
為什麼要給他們看.
就因為有釘痕.
是不是.
就是我有釘痕.
就是我經歷過死亡.
然後復活.
還有他知道他們怕什麼.
一開始那一節就說.

$^{281}$因為看見以為他是魂.
耶穌就告訴他們.
因為魂是無骨無肉的.
你們看我是有的.
對不對.
就是在對應他們的疑惑.
在解決他們的疑惑.
你疑惑我是鬼魂是殭屍.
我就給你看.
還有有骨的.
接著40節.
說了者說.
他就把手和腳給他們看.
不單止用說話安慰.
還走過去給他們看.
行不行.
有行動的.
我想給大家看.
其實我們日常生活裡面.
遇到很多困難,困擾.
有些時候我會聽到基督徒說.
他可能避過一些困難.
就說真是好彩.
其實我就說沒有好彩這回事.
不是好彩.
你說我沒有門徒那麼幸福.
因為主耶穌基督在他們面前說話.
又把手提給他們.
我們2000多年後.
我們的生活裡面.
我們怎樣知道耶穌很實在呢.
他是否在幫我們處理人生裡面.
有的情緒,困擾呢.
其實大家如果豎起我們的天線.
去看我們日常生活裡面所做的事.
你會發覺你以為或者說給人聽是一些偶然的事.
其實你細心想是神的安排.
為何你有一天.
一日向左走一日向右走.
平時你向左走.

$^{321}$那天你向右走呢.
你向右走之後.
你會遇到某一個人.
那個人可能在解決你面對的一個問題.
又或者你很困擾的時候.
你突然間內心浮起一些詩歌,一些經文.
是否偶然的呢.
為何會出現的呢.
那些都是我們的神的供應.
祂是在幫我們處理我們內心的困擾.
為何回來教會又碰到那位姐妹.
那位姐妹又拉我一邊和我分享呢.
又為我祈禱呢.
不是偶然的.
是我們的主.
我們的神.
是很細心的每一步.
每一步這樣.
體察到我們的困擾.
祂的介入.
這裡就是說主耶穌的說話的安慰.
和親身.
所以我們在透過讀聖經.
透過神的話語是會得到安慰的.
然後觀察你周圍發生的事情.
為何讓我遇到這些事情.
有沒有什麼訊息告訴我呢.
透過我們肯開放自己底層內心的情緒.
然後讓神去安慰.
有時候我們不讓神安慰.
我們要靠自己做這個做那個.
原來神已經安慰了你.
你有沒有發覺呢.
然後我們的情緒會得到轉化.
我們看41節.
黃修本很有趣.
三對情緒都有個驚字.
第一個驚慌.
第二驚恐不安.
這裡是什麼.

$^{361}$又驚又喜.
都是個驚字.
其實是不同的.
第一個驚慌就是嚇一跳那種.
第二個驚恐不安就是煩惱的那種.
這裡又驚又喜.
英文的翻譯是讚嘆.
驚嘆.
很驚訝.
可以嗎.
很驚訝.
即是面前的主耶穌基督是真的.
我聽婦女說.
或者聽這兩個門徒開始和我們分享.
我都不相信.
現在在我面前和我說話.
又把手遞給我看.
又摸到她有骨有肉.
很驚嘆.
和開心.
又驚又喜.
所以等於當你像過山車一樣.
經過了一些情緒的起跌.
心情的轉化的時候.
不要忘記要去多謝神.
要去感恩.
是祂將你的內心的情感轉化過來.
有沒有留意到自己的起承轉合呢.
你的情緒的起承轉合呢.
即是走過死任的幽谷.
然後內心有個變化.
我們看看.
他們還在驚又喜的時候.
不敢相信的時候.
是不是真的.
真的嗎.
耶穌對他們說.
你們這類人有什麼吃的沒有.
整段經文.
這個故事令我最感動的是這一段.

$^{401}$四月一節.
你這裡有什麼吃的.
那些門徒的反應是什麼.
我們再看.
他們給了一片烤魚給他.
他就接過來.
在他面前吃了.
耶穌是不是很餓.
他是不是很餓.
當然不是餓了.
他為什麼要問這個問題.
門徒又驚又喜的時候.
不敢相信.
還差一點點.
好幾個月之前.
無線有一套戲叫做靈氣逼人.
現在應該沒有人看無線劇集.
那套戲是說一個歌女.
被人殺了之後.
後來復活.
那個歌女很奇怪.
除了穿牆過壁之外.
她什麼都吃.
但正式我們的理解.
是殭屍鬼魂是不吃東西的.
我們中國人都是燒一些衣紙給先人.
所以其實耶穌再做多一步.
為我們走多一步.
我吃完東西給你看.
你還信不信.
我吃完東西給你看.
真的吃.
沒得不信.
有沒有想過我們的主耶穌.
是不是對我們好到加零一.
為我們走多一步.
所以你有的情緒困擾.
你不用怕.
我們的主耶穌基督是知道的.
祂很想和我們同行.

$^{441}$祂很想轉化我們.
想將我們從幽谷帶出來.
你願不願意被祂捉摸呢.
你信不信祂是很真實地.
我們復活的主真的臨近我們呢.
祂對我們有多好.
有沒有什麼你應該值得感恩.
你回望你的人生走過的路.
很多起起跌跌.
好像竹印那首歌.
是誰.
為什麼只有一雙竹印.
原來那雙竹印是我們的神背起我們的.
祂為我們所做的事.
是值得我們感恩回顧.
巴不得我們勇敢些.
將我們的內心世界.
向我們的天父主耶穌聖靈開放.
讓祂捉摸.
讓祂轉化.
讓我們過一個驚歎的驚.
又喜的生活.
我們一起低頭祈禱.
你是復活的主.
你更加活在我們的生活.
在我們的心裡面.
情緒是我們最個人.
最獨特的一份禮物.
你都很關心我們情緒的起跌.
巴不得我們.
懂得怎樣去捉摸自己的內心.
又有勇氣的向你傾訴.
我們相信.
當我們傾訴.
當我們帶到你面前的時候.
你會安慰我們.
你會幫助我們.
你會轉化我們的情緒.
讓我們過一個.
由憂愁變回喜樂的生活.

$^{481}$阿們.
\newpage



\section{帖撒羅尼迦前書}
\label{sec:CWxLFOO62pc}
\textbf{2024-04-11 周四晚查經班《帖撒羅尼迦前書(3)》 講員:曾裔貴牧師}
\newline
\newline
連結: \href{https://youtube.com/watch?v=CWxLFOO62pc}{\texttt{https://youtube.com/watch?v=CWxLFOO62pc}} ~~~~ 語音日期: 2024-04-13
\newline
\newline
\hyperref[sec:TfOT_yj_ii8]{\small{< < < PREV SERMON < < <}}
~
\hyperref[sec:index_chronic]{\small{[返順時目]}}
~
\hyperref[sec:index_scriptual]{\small{[返順卷目]}}
~
\hyperref[sec:_2QFVuN84qA]{\small{> > > NEXT SERMON > > >}}
\newline
\newline
$^{1}$好 我們一起同心祈禱.
感恩天父讓我們在晚上.
我們可以再一次去研讀.
貼撒羅尼加前書.
都幫助我們能夠對於.
仕途保羅和提摩泰.
或者是.
西拉等等.
都能夠對於當時.
二千年前他們在.
這個希臘.
在雅典.
來寫這封信的時候.
他們的內心是對貼撒羅尼加城的弟兄姊妹.
是掛心的.
都讓我們能夠從哲理行間.
感受到.
牧者和弟兄姊妹之間.
那份牧養的恩情.
我們可以來掌握.
也在我們現代.
一個繁忙的社會.
大都會裡面.
我們彼此的關顧牧養.
應該怎樣能夠.
可以透過一些大的方向原則.
來找到我們自己的方法.
來彼此的關顧.
願神帶領我們.
我們都會碰到關於主再來的問題.
無論我們之後會看的末世論.
但是在整卷書都在在的充斥著.
主耶穌要回來.
我們要準備這樣的概念.
這個教導.
到底保羅他的用意.
和當時的他的處境.
是為什麼會這樣想呢?.
願主都帶領我們大家明白.
我們仰望等候禱告.

$^{41}$奉耶穌基督的名.
阿們.
好,我們來看第二章.
因為停了一次.
上一次是清明.
我們來看第二章.
我很快地讀一讀.
因為這次我們應該由第十三節開始看起.
但是十三節之前有幾節未講得太詳盡.
我可以再補充.
因為有時候陸陸續續.
都會有新朋友加入我們的茶經班.
當然他們是在網上聽.
所以有時候都需要多講一些背景.
或者是有些重要的經文.
我們稍為停留久一點.
第二章很快地講了.
為什麼保羅要寫這封信.
剛才國祈禱的時候都講了.
他很快地輾轉去到雅典.
在雅典那裡.
他有一些事工.
但是當他完聖所有事情之後.
他就很掛心.
這個貼索羅來加城的弟兄姊妹.
尤其是提摩泰.
很早期就跟隨保羅到雅典.
保羅就不讓他留在雅典.
而是就打發他回去.
因為他可能是年青.
不是主要被針對的對象.
所以貼索羅來加城的暴徒.
或者那些反對的領袖.
就不會對提摩泰.
有太多的為難或者逼迫.
所以他也方便了解當時教會的情況.
這是記載在《士徒行傳》的18章經文.
大家可以有時間去看那些歷史背景.
對於我們掌握貼索羅來加教會的近況.
和這本書的寫作原因.

$^{81}$也是相對比較容易掌握的.
當然我們稍後才講關於主在內.
有時說得太多,太強調.
又會衍生一些副作用.
就是四章開始就說.
有些貼索羅來加的弟兄姊妹.
主都回來了,不要做事了.
翹起雙手等主耶穌回來了.
但是仍然要生活.
要生活的時候.
他們就在教會裡面帶來很多的怯.
就是在那裡不做事.
又令到保羅要從另一個角度去講.
主在內和你恩勤工作.
是不單止沒有抵觸.
而且是相輔相成的.
這是一個新概念.
初期教會特別是剛剛信主的教會.
其實就很需要明白這一點.
因為很多時候以為.
他回來我們就甚麼都不用管.
這個概念其實一向都是錯的.
否則主的在內就變成養了懶人.
甚至連一些未信的人.
都會覺得教會很奇怪.
所以我們一定要對真理有一個平衡.
恩勤工作不等於不等待主在內.
又是一個很平衡的.
所以原來我們對真理都要有一個.
比較掌握得全面一些.
譬如這個星期我在洪恩講道.
特別強調保羅叫我們不要論斷.
不要評斷那些牧者.
但是兩者是不是應該要有一個緩衝.
和怎樣能夠調和呢.
是不是甚麼都不要判斷呢.
那又不是的.
所以如果我們偏向某一個道理.
而在那裡不斷放大.
其實又不健康.

$^{121}$所以真理的平衡.
其實是信徒不斷需要去做一個調教.
否則偏向某一樣東西.
又會帶來一個很大的.
如果是小的可能是一個笑話.
如果大的可能是一個異端就不好了.
第二章開始.
他說弟兄們你們自己原曉得.
我們進到你們那裡並不是徒然的.
我們從前在菲勒比被害受辱.
這是你們知道的.
然而還是靠我們的神放開膽量.
在大爭戰中把神的福音傳給你們.
保羅看他的傳道.
很多的難阻.
菲勒比就是那個侍女.
那個主人因為有地方勢力.
又因為那個發財的.
即是交軌的那些靈媒.
令到他發達致富.
現在沒有那個靈媒去幫.
所以他就抓他們去見官.
保羅就要坐在菲勒比的監獄裡面.
都是一個屈辱來的.
是一個保羅看的戰爭.
一個真戰來的.
但是他仍然是靠著神的恩典.
他第二天就離開菲勒比了.
勸勉完這個女底亞的教會之後.
他就輾轉離開.
之後坐船就去到暗非波利.
去到貼索羅尼加.
在貼索羅尼加那裡有猶太人的會堂.
暗示菲勒比是沒有猶太人居住.
所以我們去年的11月去.
又發現菲勒比真的沒有猶太人的社群.
挺吻合聖經的教導.
貼索羅尼加就不是了.
保羅就看到有會堂.
他就進去會堂.

$^{161}$照著他素常的規矩.
在會堂裡面去讀經.
男士猶太人是可以.
路過都可以在安息日讀經的.
這裡告訴我們第三節.
我們的勸勉不是出於錯誤.
不是出於污穢.
也不是用詭詐.
似乎講到他的傳道在貼索羅尼加城.
有一個純傳的動機.
他特別強調這件事.
似乎和他急速的離開.
因為有極大的困難.
很難一下子解釋給每一個人知道.
他這麼急速的走.
所以一定會被人攻擊他的動機不純傳.
所以在這裡他就要提醒他們這班使徒.
進到貼索羅尼加城.
我們甚麼原因去你當中傳道.
或者甚麼原因我們急速四個星期就要走.
所以這是要解釋的.
他要讓弟兄姊妹知道.
不是見到難處.
沒甚麼好處.
沒甚麼油水等等.
就走了.
就不是了.
這裡告訴我們.
他被神嚴重了.
第四節.
但神既然嚴重了我們.
把福音託付我們.
我們就照樣講.
不是要討人喜歡.
乃是要討.
察驗我們心的神喜歡.
所以這個簡選.
保羅在這裡就用察驗.
即是神託付保羅.
這班使徒.

$^{201}$他們有那個責任.
神好像驗過他們.
即是試驗過.
因為是驗證.
試驗過保羅.
覺得他是忠心的.
可以承擔一個這麼大的使命.
就是神國度的福音.
藉著保羅.
可以傳到外邦人當中.
即是你不可以索託非人.
神嚴重了我們.
你可以想像到.
保羅在過去.
由他在大馬士革.
在路上被大光照耀之後.
那個開始.
他就接觸福音了.
接觸耶穌了.
原來從另一個角度.
是神嚴重他.
是神已經考驗過保羅.
他經過了很多訓練.
磨練.
困難.
現在他正正式式.
第二次宣教的旅程.
就來到歐洲.
大陸傳道.
遇到的是一個比較急的攻擊.
實在不知道什麼原因.
保羅不是每個地方都要走的.
有些地方他不走的.
比如十八章說他去到雅典之後.
最終他去到哥林多.
他是不走的.
他住了一年半的.
所以當然那裡又有另外一個神的顯現.
就是耶穌在晚上.
在他家裡.

$^{241}$向他顯現.
告訴他.
你不需要懼怕.
你要只管說.
因為在這個城裡面.
哥林多城.
是有很多很多我的百姓的.
保羅就在那裡住了一年半.
所以原來走與不走.
都是有神的時間.
神的心意.
作為牧者就應該要有這個.
最重要的想法.
當然客觀因素都可能會是的.
但是更重要的是內心裡面有帶領.
有神的印證.
引導.
安排.
這些我們大家都要學習的.
所以有些事情我們學一下祈禱.
或者很簡單就是.
這件事.
神啊我應該怎麼做呢.
問一下也好.
問一下之後你再用你的腦去想.
不是錯的.
都要用腦的.
但是用腦之前都要祈禱.
神啊.
我這樣做.
是不是對呢.
去或者留.
都是要這樣學習祈禱.
尤其是那些婚姻對象更加.
就是終身大事.
是.
我們說遠了.
嚴重了保羅.
和他們這一幫使徒.
嚴重之後就託付了.

$^{281}$你看看神.
都不是傻的.
神都是嚴重.
接著託付.
我們在教會.
我們不要.
下盤輕輕.
別人叫你做你就做.
你要自己都要承擔到.
所以神是嚴重.
接著託付.
接著呢.
他們就照樣說.
不懼怕人情.
不是說不看人情世故.
而是.
不是說你喜歡我說我就說.
你不喜歡我說我就不說.
那就不是.
反過來.
你不喜歡我說的真理.
我都要說明.
是.
我不是要抱人的喜歡.
否則就太容易做了.
全世界最容易做的工作就是傳道人了.
因為你喜歡我聽什麼我就說什麼.
其他什麼都不說.
如果嚴格來說.
一個星期最忙就是星期的上午.
其他時間都是做一些行政.
那很空閒的.
很沒事的.
但是.
如果你真的要說神要你說的道.
是說到人心的.
那這個就不是時間那麼簡單了.
是那個張力的問題.
是那個你要令到人.
聽到神的道的問題了.

$^{321}$所以普羅說.
既然嚴重了我們.
就是說.
不能走.
走不掉的了.
在我的人生.
他已經選擇了我做傳道.
做這個使徒.
嚴重之後.
接著神就不只是嚴重那麼簡單.
是有託付的.
所以.
其實.
我們每一個弟兄姊妹由信主那一刻.
其實你自己不知道而已.
你的心靈裡面的學務.
成長.
裝備.
其實在.
我也經常提到.
我們有一個平衡的宇宙.
在永恆的天家.
神是一直在揀選.
在地上.
在感受.
在洪恩,在嘉欣.
在浸信會.
在其他的教會.
被嚴重的弟兄姊妹.
然後他就會託付.
他驗過你都不合格.
他怎麼託付你呢?.
你可以明白那個道理.
你不追求.
你不上主一學.
其實是你自己的損失.
神都不嚴重你了.
你很渴慕神要嚴重你的.
那麼.
是相輔相成的.

$^{361}$被嚴重的人往往是自己渴慕事奉的.
譬如剛才阮薇.
去到這個階段的年紀了.
都還是很渴慕做小朋友工作.
這是心靈裡面的渴慕.
是不可能假的.
你就知道.
原來兩者是沒有衝突的.
所以普羅就說.
既然神嚴重了我們.
將福音託付給我們.
我們是受託人.
受託的是福音.
是價值連城的寶貴.
將人的生命帶進永恆的好消息.
所以普羅就不會再照人自己的喜歡.
他喜歡的就是他不喜歡外邦人.
他看外邦人不起.
他是一個拉比.
他是猶太人.
他是教法師.
他是不會和外邦人接觸的.
去完街市回家要洗手的.
這個是猶太的拉比.
外邦人是豬來的.
猶太人.
一天的祈禱.
一個男士.
三個祈禱.
一起床就是.
因為我們讀律法書.
我們識字認識神.
因為我不是一個女人.
猶太人就是三個祈禱.
所以全部是歧視的.
全部都是一個.
對其他人看不起的.
這個是猶太人的藍丁.
阿保羅就是喝這樣的奶.
但現在神厭重他.

$^{401}$要去將福音傳給外邦人.
所以阿保羅現在能夠對外邦人.
他的生命可以改變到一個地步.
他現在喜歡外邦人.
他要照直說.
關於猶太人.
如果你們再不信耶穌.
你不在耶穌的救恩裡面.
你不看耶穌就是我們的彌賽亞.
那你就不能得到神的救恩.
這是很尖銳的.
得罪猶太人的.
所以保羅就說.
我的照樣說.
就是不討人的喜歡.
所以所謂不討人的喜歡.
就不是叫我們骨雷鎚.
就是這樣不理.
就不是這個意思.
而是我們真的要.
向不同的人.
講明白那個救恩的原委.
因為不同的人都有不同的背景.
社會的背景.
他的學問.
他的職業.
所以是需要我們很有智慧地.
將那個救恩解明的.
所以這就告訴我們.
他跟這些撒羅尼加人說.
我們不討人的喜歡.
原因是要討.
測驗我們心的神喜歡.
測驗心靈的神.
所以永遠講道者.
永遠牧者.
最難治理的就是.
自己的心靈.
因為是要神測驗.
所以你講道的動機.

$^{441}$比你講的內容重要.
你的動機.
為什麼要這樣講這條道.
那個動機才是最難去準備.
就不是那些例子難準備.
最難講的是那個動機.
所以很多時候.
是先要為自己祈禱很久.
你有神的感動.
要講什麼訊息.
你就要去醞釀那篇道.
所以是等候那個訊息.
比你講的過程久.
因為要討的是誰喜歡.
討誰喜歡.
在這個時代我想大家也要學習.
既然我們有這個信仰.
唯一的信仰是真實的信仰.
我們只能夠建立在一個.
不斷學習去討神喜歡.
討神喜悅的最重要的大前提.
第五節.
因為我們從來沒有用過淺微的話.
這是你們知道的.
相當不容易的.
阿普羅很坦誠.
從來都沒有用過淺微.
淺微即是巴結.
油腔滑調.
你喜歡聽什麼我就講什麼.
不是這樣的.
阿普羅是很直接.
將神的真理.
當然照著不同的人.
你對一個小朋友.
和對一個老人家.
和對一個有學問.
和對一個可能學識不是很多的人.
你的對答都有不同的.
如果你硬邦邦對著小朋友.

$^{481}$小朋友見到你都害怕的.
好像覺得你是一個生人物近的.
你連基本的都不行.
因為主耶穌是喜愛小孩子的.
小孩子是很喜歡去到主耶穌的面前.
我從來都沒有做過幼兒工作.
我傳道信主到現在.
從來都沒有試過做.
但是我真的很喜歡小朋友.
喜歡教那些小朋友.
教那些小朋友見到我.
現在在雄恩也是.
很開心的.
那些家長告訴我.
(好誠意).
曾牧師很好笑.
他們經常都很喜歡和我聊天.
小朋友走過來.
這些是我們牧者要做的.
但是你對著一些人.
有時你要很嚴肅.
這個你就可以想像得到.
其實真是我們每一個人.
都需要學習這個淺微.
是不是我們自己已經慣性.
現代人我們有很多的.
可能包袱.
可能面具.
淺微.
傳道和淺微.
這個很文言.
即是托大腳.
即是你講一些人家喜歡聽的東西.
令到對方不斷捧到他上天.
傳道本身就不應該是這樣的.
有些是沒有辦法的.
你去做服務員.
譬如剛才在會所吃下午茶.
Lily職員就說.
哇 曾牧師我見到你就很害怕.

$^{521}$我不敢說笑話.
其實那句話本身就是說笑話.
你做服務員可能就需要說這些話.
這些是工作上大家都無雙大雅.
但是傳道和淺微.
保羅經常掛鉤的.
保羅經常.
他自己都是這樣講.
其實主耶穌都有提的.
大家有沒有讀過《馬太方》23章.
我們站一站.
很實在的將保羅這裡的心態講出來.
23章的《馬太福音》.
24章就是《鑒欖山》的寓言.
末世.
23章就是耶穌已經去到耶路撒冷的時間.
23章.
那時耶穌對眾人和門徒講論.
就說民事和范理賽人.
坐在摩西的位上.
凡他們所吩咐你們的.
你們都要謹守遵行.
但不要效法他們的行為.
為什麼呢.
因為他們是能夠講.
不能夠做的.
他們就將蘭丹的重擔.
就捆綁起來.
來去擱在其他人的肩頭上.
但是自己一個指頭都不肯動的.
很全神地.
閉上嘴講.
講完之後.
自己做不到之餘.
從來沒有想過做.
只是壓制其他人做.
那意味著耶穌教門徒將來做傳道.
他升天之後.
你們就不要學這些范理賽人了.
你們要學的.

$^{561}$因為你們要教聖經.
教真理的.
譬如彼得.
從來都不是一個傳道人.
他是漁夫.
將來就要學人寫聖經了.
要學希臘文的.
他要講道的.
但是他們要避免的不是不講.
但是要講完之後.
實踐你自己講的真理.
還有千萬不要.
好像講完之後.
根本就沒有自己的實行.
可能是用一種危言聳聽.
令人懼怕神的角度.
而來去遵行.
但實際上自己都不是真正那麼懼怕的.
阿保羅就是這樣來實踐.
所以耶穌在臨他升天之前.
就是這樣警戒門徒.
你們看著法理賽人民事的.
即是他們看著耶穌.
怎樣來跟這些法理賽人.
在受難州的時候.
或者受難州之前.
都有很多類似的衝突.
但是你看到耶穌很尊重他們.
因為他們是代表著神.
坐在摩西的位置來作教訓.
耶穌不會對這些人.
來去很輕視.
很侮辱等等.
千萬不要.
因為重點是他坐在摩西的位置.
他是代表著神.
某程度是講聖經.
救藥聖經.
但耶穌就說.
不是對那個人.

$^{601}$而是他們做不到的態度.
千萬不要學.
就是這個意思.
所以日後你學了那個樣子.
民事是怎樣裝扮自己的.
在聚會裡面.
怎樣來講聖經.
你不要學外學的外表.
而是內在的心態.
耶穌就教內在的心態.
既然講摩西五經.
講神的道.
你自己就要以身作則.
去成為門徒.
或者百姓的仿效對象.
所以保羅在這裡.
他就來跟這些帖撒羅尼迦城的人講.
他自己傳道的心態.
一個傳道者.
最重要的就是.
給信徒看到.
那個內心的態度.
不是淺味.
不是講一些淺味的話.
應該要講的就要講.
接著就告訴我們.
不是.
「這事你們知道.
亦都無藏著貪心」.
貪心是要藏著的.
就是你如果.
人家講甚麼.
你都說好.
那你就是說.
你就是.
讓人看到你的內心的動機.
那你就.
不能夠騙到更加多人.
一般人都會藏著貪心的.
所以你看到保羅.

$^{641}$對於他自己的內心.
來到他否白出來.
真是非常之.
重要的一個傳道的守則.
或者我們去找一些.
事奉的人.
他們的內心的動機.
是要經過這一類的試驗.
其實保羅講了這麼多.
他自己的內心.
其實就是證明.
神已經驗證他.
試驗過他.
而他能夠通過這些.
所以處理金錢又要很小心.
要很透明.
不可以隨隨便便.
否則又會令到人一個誤會.
就會變成一些風波.
或者謠言等等.
所以保羅說.
不單止不是淺微.
淺微就是言語.
言語背後就是內心.
原來你不是敬畏神的.
原來你怕某些人的.
原來你很想得到那個人的心.
去歸你.
你不是想將他帶到神那裡.
這些有時很微妙的.
很微妙的一個人和人之間的關係.
尤其是牧者.
是一個牧羊.
就會有一定的人對牧者的信任.
對這個積分的信任.
所以我們更加不應該將人.
帶到自己的心去歸你.
應該要帶他去歸神.
所以這個就涉及到.
那個淺微的問題了.

$^{681}$貪心固然就是內心.
一切都是自己裡面對神的供應.
對神的榮譽看得不夠大.
我們才會緊張人的榮譽.
如果我們緊張神的榮譽.
我們就不會看得那麼重.
人怎樣看我們.
這是神可以作見證.
不單止是自說自話.
有神作見證.
你想一下.
只是這幾節聖經.
已經將一個傳道者.
一個在教會裡面的牧者.
那個積分,那個心態,那個品格.
已經表露無遺地.
來到淋漓盡致地講出來.
就是說讓帖撒羅尼迦的教會.
不需要來猜度保羅,提摩泰.
他們為甚麼急速離開的原因.
因為自然一定會有人講保羅的壞話.
大家保羅就跟他們講他現在內心裡面.
他自己的那個事奉的態度和動機.
很寶貴的.
一個牧者能夠將他內心的動機.
很坦誠地去分享.
其實這個真是一個很大的福氣.
神可以為我們作見證.
第六節.
我們作基督的使徒.
雖然可以叫人尊重.
卻沒有向你們或向別人求榮耀.
你想想這些智慧行間.
保羅有的是信心.
對自己的積分的自信心.
對自己在神面前的猜險有信心.
即便我得不到你們的物質供應.
甚至心靈的安慰.
我也不介意.
我不求你們的榮耀.

$^{721}$所以有時真是此消彼長的.
我們越是阿諛奉承.
那你就可以不用問阿貴.
那些會眾就不尊重你了.
因為你自己.
downgrade你自己的積分.
那個身份.
反過來.
當然我們不是在扮抬高自己.
又不是另一種的.
就不對的.
而是本於我們自己的傳道.
既然答應了神.
走上這條人生的事奉路.
我們就要在神的裡面得到.
那個供應,感動和榮耀.
所以這些是我們事奉者一定要學的.
一定要尋求的是神的榮耀.
所以保羅說.
我不是要向你們求榮耀的.
所以你會看到.
我們看過一些經文.
《四道行傳》.
帖撒羅尼迦都有幾個年輕人是跟隨保羅的.
就是說.
他講明了我不是要求你們的榮耀.
但是反過來.
反而就會有更多的帖撒羅尼迦城的人.
會對保羅事奉的動機明白之後.
就會對他讚賞有加.
所以這個也是某程度的相輔相成.
就是說.
傳道者一定要搞清楚自己的事奉的動機.
不斷在這個動機裡面.
是純傳.
自然一定會得到神的榮耀.
和眾人的認同.
最終的肯定等等.
這些全部跟事奉有關的.
所以頂子們你們參加或大或小的事奉.

$^{761}$都必須要尋求神的榮耀.
第七節.
只在你們中間全心溫柔.
如同母親預養自己的孩子.
預養就是說還沒過戒奶期的.
這樣的嬰兒.
帖撒羅尼迦城是信主四個星期.
保羅就突然離開.
很掛心初信的教會.
他們面對Big Dark.
面對一些謠言.
面對一些是非.
他們能不能夠站得住呢.
保羅是掛心信徒.
他們的靈命.
他們對真理的認識.
或者在小小的Big Dark的面前.
是否能夠站穩呢.
保羅掛心這些的.
不是掛心他自己.
那些人怎麼看我.
其實不用太過顧影自憐的.
那些人怎麼看他.
他自己都控制不了.
但他可以控制的就是.
他自己的內心的.
崇傳的動機.
所以他說他全心溫柔.
這位使徒.
他可以在哥林多教會.
他說如果你們不肯改變.
我就帶著權杖.
去到你們中間的了.
但是對帖撒羅尼迦城的弟兄姊妹.
他說我好像一個母親那樣.
一個溫柔的心.
全心溫柔.
所以傳道者.
是真的隨時要變臉的.
對一些人你要惡.

$^{801}$對大部分的人你要溫柔.
將神的道講得.
人家聽得很舒服.
很多安慰.
尤其是有困擾的人.
得到安慰.
其實是我們的責任來的.
有些人.
不覺得需要安慰的.
那些道是要針到他的心靈.
他怕神.
那又是要的.
那你怎樣拿捏呢.
那就是對牧者的挑戰.
對牧者對群羊的了解.
如果你從來只懂得安慰.
講一些道.
不痛不癢的.
講完你都不知道你在講什麼.
大家都很開心地繼續去喝茶的.
那可能你從來都沒有做過一個.
嚴肅的傳道者.
應該講的那些訊息和真理.
當然.
如果一個人硬邦邦.
都不理人的內心.
那也是不合格的.
所以其實最需要的就是.
牧者的心靈的柔軟.
是能夠感受到信徒內心的情況.
所以普羅在這裡說.
全心溫柔.
如同母親與養自己的孩子.
如同.
就是說他又很明白.
一個小朋友需要媽媽的那個語釀.
需要媽媽的那個溫柔.
比如小朋友.
小時候我想大家都在座的父女.
你們知道.

$^{841}$BB有時候他晚上醒.
最需要的就是抱.
一抱就不哭了.
你抱著他放下他就馬上哭了.
有時候就是這一種.
有時候他所謂的那些叫做害怕.
以前小時候那些母區.
經常傾門大雨等等.
小朋友很害怕的.
你抱著他就會有安全感.
原來真的是.
所以不同人生階段.
都需要有這種語釀.
所以普羅不只是一個硬邦邦的使徒.
他也很明白.
初順的好像一個嬰孩一樣.
他需要母親的語釀.
語釀就是用媽媽的奶去餵養.
語釀.
所以中文都翻譯得很好.
但是如果全個教會所有人都是這樣語釀.
就不行的.
你應該慢慢教會長大.
有些人已經成為了精英.
他可以在任何的環境.
都經受得起挑戰.
那教會才成長的.
每個人都要語釀的.
回來後說重一點都不行.
我走了.
我不信耶穌了.
那教會就慘了.
每個人都要語釀.
信了二十年還要語釀.
這樣就不行了.
那個教會就沒有前途了.
這裡是說了四個星期.
和面對這麼大的困難.
和使徒他們要立即離開.
不夠時間去解釋.

$^{881}$但是普羅他都很明白.
初信的教會需要在神的真理上.
慢慢去紮根成長.
但是如果你說已經二十年三十年.
你還慢慢去紮根.
還慢慢說你不要說這麼深.
你每一節用人夫三章十六節就行了.
其他都不用教我們了.
那教會怎麼會有希望.
怎麼會有成長呢.
應該是隨著信主的年日.
那個歷年.
應該是越來越自己可以謀生.
自己可以來到獨當一面.
面對不同的困難.
都可以來到依靠神的去事奉.
那教會就很不一樣了.
這裡告訴我們.
第八節.
我們既是這樣愛你們.
不但願意將神的福音給你們.
連自己的性命也願意給你們.
因你們是我們所疼愛.
愛就帶來這種犧牲.
這個犧牲有的.
最背後的價值觀就是.
因為更加有神的託付的福音.
所以既然福音是這麼偉大的.
主耶穌的降生.
來背十字架.
所以傳道者最終.
都是要這樣背十字架的.
去養大那些信徒.
讓他們能夠.
將來他們自己能夠獨當一面.
所以在這個過程.
一個生養的過程.
怎麼可能沒有犧牲呢.
半夜起床餵奶.
突然間三更半夜發高燒.

$^{921}$要趕去急症室.
這些是在在.
在一個貧窮的年代就會出現.
多少的地方.
怎麼可能小時候我們有飛蛹.
有印蛹哪有的呢.
就是那個媽媽.
什麼都是那個媽媽背著去的.
我十二歲.
不是十二歲.
八歲還是九歲.
背部燒到.
不知道是一級燒傷還是怎麼樣.
三點多半夜.
家裡這麼多兄弟姐妹.
最終誰跟我去.
就是我媽媽背著我去.
那時候九龍塘地鐵站還沒有建造.
背著我出去.
三點多.
出去到外面那個兵房.
就是現在解放軍的那個地方.
就是窩底的路道.
三更半夜哪有車.
突然間有的士停下來.
有乘客的.
問什麼事.
媽媽背著小朋友.
媽媽當然說燒傷了背部.
馬上送去廣環醫院.
去到醫院.
對面就是窩底路道.
就是那個.
那個什麼學校.
信義學校下面.
那個叫什麼堂.
幼齡堂還是什麼.
我忘記了什麼堂.
信義會那裡.
停在斑馬線.

$^{961}$就叫我們下車快點跑去對面.
對面就是急症室.
那又不收我們錢.
走了去.
一個好心的乘客.
還有那個司機.
很歷歷在目的.
那就馬上進去急症.
馬上做.
然後馬上要留醫.
起碼一個星期.
那個背部好像.
大家吃過那些乳豬脆皮.
我是癢的.
那時候廣環醫院沒有冷氣的.
小時候.
要怎樣呢.
因為有時候被衣服壓著.
要不准穿衣服.
打噴嚏.
然後上面開那個大的牛角線.
風扇.
吹著降溫.
背部.
然後有時候小孩子.
還沒到十二歲.
撕到那些蕉爛了流血.
那些護士那時候.
很兇的那些護士.
很大年紀的.
綁著兩隻手在床邊.
即是小朋友的.
嬰兒床.
綁著兩隻手不准動.
有人來才准.
整個人就躺下.
背部向上.
要降溫.
小時候是這樣的.
所以是媽媽的餘養.

$^{1001}$我又試過那時候剛剛十二歲.
中一在新發書院.
貪玩.
在那區那些街坊.
那些工廠.
大家扔石.
不是扔石是扔豆腐渣.
豆腐渣是不會扔到的.
在那裡追逐.
沒有穿拖鞋.
不停跑的時候.
整隻腳踏在玻璃瓶上.
玻璃瓶是尖的.
大家看到那個.
一拆下去.
差不多一個半月.
那隻左腳是不可以走路.
要用木頭車推出去.
然後上課.
又是我媽媽.
背我上三樓.
去到新發書院.
太久不上學不行.
背我上去.
這樣的日子.
這些就是餘養.
地上也尚且有.
我特別多.
有試過被狗咬.
又是我媽媽背我去廣南醫院.
所以腳有.
腳板也有.
背部也有.
現在打開衣服.
看不到紋身.
但看到有些疤痕.
腳也有不同的疤痕.
這些就是餘養.
我想不同的弟兄姊妹.
你也有你自己人生.

$^{1041}$一些別人對你的餘養.
或者你自己.
現在對其他人的關愛.
所以.
如果我們都能夠將心比己.
如果教會裡面.
我們稍為有成熟的眼光.
我們都要扮演.
這個餘養.
其實沒什麼.
你回來坐完.
祝福我也不稀罕.
我走了.
OK的.
但你的人生就只能夠維持這樣.
餘養一定是痛苦.
餘養就用很多心機.
很多時間.
就能夠取得這個積分.
你感受到.
對其他人的一種.
不知什麼原因.
我總是對他有些興趣.
有些留意.
你就開始在禱告裡面.
為對方祈禱.
為這位朋友祈禱.
慢慢認識了.
有機會再這樣接觸.
這些餘養.
其實就是你開始扮演著.
神救贖的.
角色的其中一個點.
這是很大的.
對你的生命的成長.
所以如果你沒有積分.
沒有事奉.
是輕鬆的.
是很開心的.
回到教會.

$^{1081}$禮堂的音樂是最好的.
師班是最好的.
你享受這麼好的音樂.
那邊的廣道又這麼好.
已經訂了位在會所.
又這麼舒服.
其實你失去了很多.
經歷神的恩典.
和痛苦帶來的滿足.
那個生命的建立.
那份喜悅.
譬如.
我們在座的Alex.
對小朋友.
小朋友有時很頑皮.
走來走去.
但你不要理會.
總是建立了一種關係.
看著他們成長.
由無心到有心.
你好像能淋一點水.
在他們的生命裡.
將來他們長大能夠記得.
這些其實就是.
價值連城.
其實就是愛.
所以普羅話既然我們這樣愛你們.
不單止願意將.
神的福音傳給你們.
連自己的性命都願意.
給你們.
因為那裡關乎到疼愛.
所以一定是愛的工作.
教會為何能夠.
立到今天.
教會是有愛的.
彼此的傳遞.
才能夠建立生命.
我們不單止建立一個.
頭腦的知識.

$^{1121}$基督教.
特別是.
宗教改革很強調.
真理,理性.
那個pitching.
所以.
福音派教會的講台是放在中間.
禮儀教會.
是放在旁邊的靠邊站的.
不要放在中間那麼礙眼的.
你留意一下.
你去禮儀教會就知道.
那個講台在旁邊.
但是沒錯.
我們強調真理,強調理性.
強調對真理的堅持.
因為正正是.
要避免天主教的謬誤.
但是絕對不可以.
沒有了愛和犧牲.
不然就變得.
硬邦邦的太過理性.
太過只是一道道理.
所以兩者.
不可以分割的.
傳道者.
為什麼有些人要去到.
遠方的去傳道.
正正就是因為愛.
這個民族有一個託付在這裡.
所以保羅這裡也都講了.
他雖然是這樣突然間離開.
但是.
歸根結底他是對.
這個教會.
這個剛剛成立的教會.
是有一份很強烈的愛.
是從神那裡而來的愛.
其實不是很認識的.
其實按保羅來說.

$^{1161}$愛邦人他根本不會.
放在他的心裡面.
但是現在愛邦人因為福音.
放在他的心裡面.
成為他所愛.
所疼愛.
你就知道其實是神改變了保羅.
不是那些愛邦人特別可愛.
最重要其實.
保羅在神眼中.
其實不可愛.
這個才是最大的.
危險.
大家聽過周柱.
除三害.
有個南方的蛟龍.
在一個深山.
有一隻很兇猛的老虎.
村民怕得要死.
有三害.
這個周柱原來.
很厲害的很powerful的.
很好打的.
就說好啊我幫.
大家除害.
除了老虎打死了.
那隻蛟龍跟周柱.
搏鬥.
搏了三日三夜.
還沒回來.
村民已經不停慶祝.
慶祝到很開心.
周柱最終.
幾經辛苦打死.
蛟龍回來了.
那些人就不開心.
就怕.
周柱知道之後.
就死了.
原來第三害是他.

$^{1201}$他是第三害.
他本來跟蛟龍一起.
自盡.
村民就拍手.
燒炮仗很開心.
他一回來.
兩害都除了.
誰知原來他才是.
真正的害.
他後來就.
我不知道是真實還是假的.
他就切心悔改.
真的除三害.
所以那些村民.
很喜歡他.
他做了官.
讀書做官.
這個就是普羅.
普羅就是周柱.
他就是這樣.
在神眼中.
原來最不可愛的是他.
所以他現在對.
帖撒羅尼迦教會.
其實四個禮拜的交情.
怎會有很熟呢.
普羅說因為你們是我所疼愛的.
你可以想像到.
普羅因為福音的緣故.
因為主耶穌.
他看外邦人.
已經是另一種眼光.
他已經好像自己做了.
一個白內障手術.
看東西現在看得很清楚.
原來神愛是世人.
不是神愛猶太人.
神愛每一個民族.
這些都是我們需要消化.
信主來了.

$^{1241}$在教會內我們還有沒有.
一些偏見呢對一些不同的民族.
有沒有一些我們覺得.
我們很優越呢.
那真的要大家.
要小心.
不然神會審判的.
因為神愛.
每一個民族.
正正就是我們既然領受.
這麼多恩典就要分享.
這個就是.
關乎普羅對於.
這個教會.
他們的誤解.
或者一定是有心人會.
控告普羅.
普羅就要好像.
其實不是保妃.
而是很想講他那個牧羊的內心.
給喬薩羅尼加教會.
知道.
所以這些經文都是對傳道者.
最好的提醒.
對事奉神的人.
最好的一個教導.
接著第九節.
弟兄們你們紀念我們的辛苦.
勞碌.
晝夜作供全神的福音.
給你們免得你們.
一個人受累.
我們向你們信主的人.
是何等聖潔公義.
無可指責有你們作見證.
也有.
神為我們作見證.
又再一次講他的內心.
那個動機和態度.
這次又是另一個角度.

$^{1281}$你們都要曉得.
我們怎樣勸勉.
我們提到炎父.
和那個真理.
之前講母親這個比喻.
就是講了一種禹養.
所以你會看到.
第十一節.
所以是很平衡的.
你們要曉得我們怎樣勸勉你們.
安慰你們.
祝福你們各人.
獨習.
你們各人.
好像父親.
待自己的兒女一樣.
不是幼兒.
即是說你已經是成年人.
我要跟你講的是道理.
講的是真理.
是父親應該有的.
例如家教.
一些規範.
一些做人的.
成長的一些.
很重要的價值觀.
品格的建立.
這些就是父親的形象.
所以如果我們是父親的成長.
我們又要保阿非.
我們要尋找一些智慧人.
所以如果我們教會.
只是母親.
沒有父親形象.
又是一個嚴重的缺失.
如果只是一個.
只是講道理.
沒有了愛沒有了那種的.
勸勉關懷.
又是不健康的一個失衡.

$^{1321}$失去了平衡.
每個人的成長.
順處的階段成長.
有些角色的弟兄姊妹.
一些主裡面的先賢.
哥哥姐姐.
或者叔叔阿姨.
但也都需要有父親.
父親的形象.
好像父親待自己的兒女一樣.
要叫你們行事.
對得起.
召你們進他國.
得他榮耀的神.
所以父親是帶來的傳承.
父親是要帶來.
回到神的國度.
那個考慮.
所以不只是安慰.
都是有平衡的真理的傳遞.
所以兩者.
這樣就平衡了.
信徒的靈命就穩固了.
有媽媽的愛.
一種的餘養.
含深如虎的.
但也都有爸爸的威嚴.
有那個價值觀.
有那個品格.
更加重要有那個永恆盼望.
所以.
要對得起.
召我們進.
他的國度.
得他榮耀的神.
這個.
品格相稱的人生.
所以.
基督徒其實是很有榮譽感的.
因為我們就是要得到.

$^{1361}$神的國度.
神的榮耀.
所以我們就要對得起.
你想想.
如果.
現在.
整個社會都進入.
第二,第三代.
甚至第四代.
很多的接班.
因為就快要跟.
業主的新的接班人.
來見面.
其實到處都是.
譬如周大福.
新宏基.
都有接班人在做.
你就會看到.
新一代他們的做法.
看事物的視野.
他們是飽讀詩書.
在外國回來.
就不是創業的了.
又怎樣能夠在現有的家族的事業.
再擴大呢.
譬如現在需要的就是.
可能社會的經濟的轉變.
怎樣能夠多的.
所謂的ESG.
那個社會的責任.
我想第三代他們會有這種.
眼光.
因為有這麼多的.
職類的實力.
其實是很好的一個.
推展.
這些福利事業的一個機會.
這裡一樣.
所以我們要.
怎樣來對得起.

$^{1401}$我們的行為.
我們的品格.
是怎樣能夠符合.
即是你走出來.
那個暴發戶.
一看你就會看到.
這個人是很有教養成長的.
第三代了.
第四代了.
有眼睛看的.
看得到的.
人家從小就在這樣的環境的.
分途底下.
一出來留學回來.
就大家都知道這個是哪個家族的.
小朋友.
是有一種氣質的.
跟你在.
心身來說的基層家庭.
完全是兩碼子的事.
我們不是要做比較.
不是要做歧視.
而是這裡聖經告訴我們.
保羅強調.
要叫你們.
行事.
即是為人.
對得起.
召我們的神.
這位神就是召我們.
去他的國度.
得到他的榮耀.
他的國度.
就是永恆.
國度的榮耀.
所以我們有多少弟兄姊妹.
所以我們說.
讀遍了羅尼加前後書.
必須我們讀完之後.
你要時時在心裡.

$^{1441}$可能你從來都不會這樣想的.
但你開始要這樣想.
今晚耶穌回來.
我就會進入他的國度.
那我是不是.
他看我是有這個榮耀.
而稱讚呢.
我們一定要問這個問題.
問這個問題其實很.
你可以說很嚴肅.
但應該很興奮.
因為你已經開始有一個警示.
開始你會有一個警示.
我今晚見他.
他今晚回來.
那會怎樣呢.
這裡就已經說了.
要叫第十二節.
要叫你們行事.
對得起那召你們.
進他國.
得他榮耀的神.
對得起.
他是召我們去他的國度.
就是說如果你.
很快下個月.
你要長住在白金漢宮.
那你沒理由.
那個行事還是.
一定要有一個.
王子公主.
那種修養.
學你都要的.
學一學就像了.
這個就是基本地上.
都尚且是這樣.
何況我們是神召我們.
進入他的國度.
這裡已經是第二次再說到.
他會回來.

$^{1481}$和我們要去.
所以一直都不是說.
我們人只是死了就算了.
是說進入永恆.
那個國度.
是神召我們去的國度裡面.
所以強調的.
反而不是將來那個世界.
是怎樣.
強調的是現在我們有了.
這個盼望之後.
應該怎樣行事.
怎樣每天去裝飾自己.
所以是一個.
很激動的.
很動態的.
不是坐在那裡等.
就是說我們要參與事奉.
我們要學習禹養.
我們要學習好像父親的.
那種教誨的真理.
因為好像父親.
勸眠.
安慰 祝福.
這樣的意思.
就是說我們需要.
在這個人生.
善用神.
給我們的恩典.
恩典帶來的就是.
你的恩賜的發揮.
盡用我們的恩賜.
我們的人生.
這樣來最終都是幫助有需要的人.
進入神的國度裡面.
十三節.
為此.
這就是我們今天講的經文.
因你們聽見.
我們所傳神的道.

$^{1521}$就領受了.
不以為是人的道.
乃以為是神的道.
這道實在是神的.
並且運行在你們.
信主的人心中.
抵門你們曾效法猶太中.
在基督耶穌裡.
神的國教會.
你們也受了本地人的苦害.
貼上羅尼加城的當地人.
像他們.
受了猶太人的苦害.
一樣.
這猶太人殺了主耶穌.
和先知.
又把我們趕出去.
保羅就講了他們離開的原因.
他們不得神的喜悅.
在講猶太人.
和那些匪徒.
即當地的土豪.
不得神的喜悅.
且以眾人為敵.
即以福音為敵.
不許我們傳道給外邦人.
明顯是講猶太人.
那些拉比等等.
或者會堂裡的人.
使得外邦人可以得救.
上上充滿了自己的罪惡.
神的憤怒.
臨在他們身上.
已經到了極處.
其實他們感受不到的.
他們感受到的.
是神的氣焰.
迫打保羅.
保羅要走.
猶太人在當地.

$^{1561}$可能結合政府的勢力.
攔阻.
使徒他們.
講耶穌是猶太人的.
彌賽亞.
迫黨福音.
恐嚇他們.
不可以講福音.
但這是人.
看到的.
現實上的.
形勢.
但保羅在神的眼中.
看到的是.
猶太人.
你不接受耶穌.
你攔阻人接受耶穌.
你就在神的震怒下.
已經去到極處.
可能他們大部分人.
都是受宗正寢.
沒有迫打.
但是永恆的沉淪.
因為神的審判.
非常嚴重.
所以我們不要只是從.
臂上的眼界.
去看得和失.
成功和失敗.
要從保羅.
他們在宿天的眼中去看.
福音.
如果我們站在福音這一邊.
永遠我們都是.
在神這一邊.
保羅就是站在福音這一邊去傳.
講了這麼多節.
由第十節開始到第十六節.
都是講猶太人.
這個這麼難堪的誤解.

$^{1601}$對主耶穌的身份的不認同.
堅決地在他們有生之年.
要拒絕自己信耶穌是彌賽亞.
不接納耶穌就是他們的救救世主.
還要現在加上一個.
他們的一個熱心.
就是攔阻猶太的傳道人.
就是保羅.
或者提摩泰是半個.
西拉猶太人.
向外幫人.
已經不向猶太人了.
不打擾你們了.
你們不喜歡我們在會堂了.
我們就向外幫人傳道.
都不行都要趕.
都要報官趕走他們.
正因為你們是猶太人.
正因為更加重要的是.
你們講的是那個訊息本身.
不單止你們猶太人這麼簡單.
你們講的那個福音本身.
就是耶穌是猶太人.
是神的代理人.
是神設立的救主彌賽亞.
這個就是令到猶太人.
堅決要抵制和攻擊.
但是我們要記住.
保羅這樣講.
第十五節.
「這猶太人殺了主耶穌和先知」.
這個就是在耶路撒冷的.
逾越節那幫猶太人.
但是這個錯誤的真理.
繼續傳遞到其他地方.
就是說很純正的猶太教的人.
都是這樣抵擋耶穌的.
所以又把我們趕出去.
就是說一直延續到.
耶穌已經升天了.

$^{1641}$但是這個錯誤.
繼續發展下去.
所以保羅既然你傳遞耶穌.
是救主.
保羅就和耶穌一樣的下場.
面對迫德.
面對那個人的攻擊.
是猶太的同胞對他的攻擊.
所以他們又把我們趕出去.
「他們不得神的喜悅.
且以眾人為敵」.
不准許我們傳道給外邦人.
使得外邦人可以得救.
傳道聽明白會得救.
所以我們記得要講.
無論如何都要講.
講明白.
那個人一聽.
聖靈會開他的心.
他就會相信.
所以這裡告訴我們.
很重要真的要傳.
叫外邦人可以得救的是道.
是真道被解明.
所以很重要.
所以現在他們就常常.
充滿了自己的罪惡.
在他們來說.
我們是替神來懲治這些異端.
猶太人看保羅就是異端.
但原來他們是充斥著自己的罪惡.
在神的眼中.
神的憤怒.
淋在他們的身上.
已經去到極處.
最恐怖的是.
完全感受不到.
好像已經挖了瓦.
沒有那種感受.
不覺得我們有生之年的猶太人.

$^{1681}$是神所不起悅.
甚至乎要面對毀滅.
很快就去到主一號七十年.
整個耶路撒冷.
整個以色列.
被羅馬的提多將軍攻陷.
聖殿被拆毀.
那些金融了.
因為燒了聖殿.
那些金箔融在石縫裡.
所以羅馬人就挖開那些石.
拿那些金.
所以整個耶路撒冷被拆毀.
整個聖殿都被夷為平地.
最後在馬薩大.
那些猶太人最終要自殺而死.
不可以讓外邦人來殺.
所以羅馬人就做土鼎.
一直繞著上到馬薩大那個地方.
最終有差不多一百人.
最後自殺而死.
因為如果被外邦人殺.
他們就會覺得不能去到神的國.
這是一些傳說.
但事實的確當時.
羅馬人真的這樣做.
繞道上到山頂上.
殺那些猶太人.
最後的那些抵抗軍.
整個以色列在主要70年.
猶太人就血流成河.
全部被擄.
分散到天下.
所以這裡告訴我們.
神的憤怒臨在他們身上.
已經到了極處.
保羅一方面他不是危言聳聽.
因為他也知道.
主耶穌也說過.
猶太人逼迫耶穌.

$^{1721}$要面對他們自己的審判.
所以很快十多年後.
提多將軍.
當然保羅已經死了.
他已經回天家了.
提多將軍就在.
我不知是否多米團的任內.
就來攻擊耶路撒冷.
就真的將猶太人徹底地移為平地.
這就是歷史的真確發生.
所以這裡說.
神的憤怒臨在猶太人身上.
已經到了極處.
他們這樣迫害教會.
迫得福音.
迫得這些使徒.
所以也是一個挺大的問題.
教會也遭受過很多迫害.
不同的年代.
近代我們中國人的文革時代.
比較大的.
其實早期也有的.
譬如有飛機的時代.
一九一幾至一九二幾年的時候.
非基督教.
不是飛機.
不是Alpine.
非基督教運動.
就是中國人要剷走基督教.
輾轉就去到五四運動等等.
接著就開始有國共內戰.
接著有日本侵華.
輾轉都有這些對基督教的不友善.
但是神的工作.
往往在一些困難的時候.
更加蓬勃的發展.
信仰的時候就是要經受一種考驗.
才能夠更加茁壯成長.
所以這裡猶太人.
他們完全看不到自己面臨的危機.

$^{1761}$但是保羅表面就是受到迫害帶來的不方便.
還有那個影響.
但是保羅看到的就是.
神站在他那邊.
神就是那個審判者.
他們要面對神的憤怒.
所以貼薩諾尼加前書也說很多神的憤怒.
所以我們不要以為.
信耶穌信不信也不要緊.
其實不是的.
某程度我們信耶穌.
也不要覺得這個宗教很好.
有很多好處.
生養死葬.
是的真的有的.
你進了學校也有人有婚.
你能夠進去基督教的墳場.
也是因為你有洗禮證.
你沒有洗禮證就不行了.
你快點搞定它.
不然你將來只能去和合石,臨灣那些.
但是仍然是說一件事.
信耶穌是要逃避神的憤怒.
這是其中一個重要的點.
如果我們很多人說.
信耶穌是希望逃避神的憤怒.
而因為耶穌而我們得到這個救恩.
這是健康的.
你不要以為不健康.
反過來你只以為這個宗教是很開心.
其實某程度上不是很健康.
這是一個入門是沒有問題的.
任何人入門都沒有問題的.
但是入了門之後.
你要知道原來神是有憤怒的.
我們要逃避這個憤怒.
這就是透過耶穌.
真是主耶穌擔待我們世人的罪.
就是這個憤怒就解除了.
所以都是要全面.

$^{1801}$所以任何人怎樣信耶穌是沒有問題的.
甚至動機不良都可以的.
但是入了門之後.
必須要更正很多錯誤的想法.
否則你永遠都停留在一個.
只是很好啊很好啊.
其實基督教往往很多東西.
是要我們付代價才能換取的.
所以這個就是我們要留意.
經文其實對一個初信教會.
普羅都很坦誠地說了.
你不站在神那邊.
你就面對神的憤怒了.
猶太人就不是站在神那邊了.
又是猶太人.
又是舊約的神的選民.
但是居然拒絕耶穌.
耶穌就說了.
有了我就有神.
你離開我你就什麼都沒有了.
所以這個是一個挺嚴肅的問題.
所以神已經給了耶穌全權的.
將人類的福祉和救恩.
在十字架上已經成就了.
所以這裡就告訴我們.
我們就必須要來學習這樣的功課.
我們多讀一點.
十七節.
弟兄們,我們暫時與你們離別.
是面目離別.
心裡卻不離別.
我們極力的想法是.
很願意見你們的面.
所以我們有意到你們那裡.
我保羅有一兩次要去.
只是撒旦阻擋了我們.
我們的盼望和喜樂.
並所誇的冠冕是什麼呢?.
豈不是我們在主耶穌來的時候.
你們在他面前站立得住嗎?.

$^{1841}$因為你們就是我們的榮耀.
我們的喜樂.
如果我們一直從這個.
自律行間去理解保羅.
我們可以很深的感受到.
他不是要來講一些好聽的話.
去欺控這班弟兄姊妹.
而是他真的講述他的內心.
那種困難.
好像神沒有為他開門.
而且也有一個.
可能是他認為客觀的原因.
他必須離開.
可能就是令到教會.
不需要在這麼大的風波底下.
來去被當權者.
當時可能是羅馬的官員.
或者是猶太一些有勢力的人.
來去要將整個教會.
可能都要來去關閉.
就是說如果保羅走.
教會還可以繼續存在.
在這樣的情況下.
保羅就認為他離開.
都不是太大的一個問題.
所以他自己這樣離開.
原來他的內心說.
我是極力的想法子.
很想見你們的面.
所以他是有意去到他們那裡.
回到他們那裡.
但是暫時目前.
不是一個好時機的時候.
他就不去和那個形勢.
太大的對抗.
反而他自己現在.
暫時讓那個風波稍微緩和.
他就可以將來日後.
再去部署.
所以他自己就離開.

$^{1881}$只要落到雅典.
其實他經過比利亞.
再落到雅典.
在雅典那裡開始穩定了.
雅典人因為很多是外邦人.
還有很多是講道理.
講哲理的人.
所以就不是那些猶太人的死硬派.
所以當地的氣氛就相對比較緩和.
保羅就有機會寫信.
叫提摩泰回去.
還有看一下教會現在的狀況.
所以就是這樣的原因.
不過第十九節.
就是將保羅他想講的東西.
講出來.
他盼望的不僅僅是他們相安無事.
我們都期望教會風平浪靜.
沒有問題.
但是原來問題是會令到我們更加成長.
所以保羅的祈禱就是.
我們的盼望和我們的喜樂.
和他的祈禱就是.
他所誇耀的他們使徒.
最想要的冠冕是什麼呢?.
大家有沒有想過呢?.
有沒有想過只是個人的冠冕.
還是在事奉帶來的喜樂.
那個將來的獎賞呢?.
我們所誇的冠冕是什麼呢?.
是什麼呢?.
又將他這班使徒.
他們的內心原來渴望的東西.
原來那個層次是這麼高的.
我們經常都想天堂是怎樣的.
是不是真的不用做.
怎樣飛的.
我們就很好奇這些.
但是聖經從來都不講這些的.
從來都不講你怎樣飛.

$^{1921}$你飛一千公里是不是比飛機快.
還是戰鬥機.
都不講這些的.
怎樣閃來閃去是怎樣的呢?.
我們很好奇.
那個黃金街是怎樣黃金的.
我們很好奇.
但是從來我們不會羨慕.
阿保羅說我們所誇耀的冠冕.
是什麼呢?.
就是說他們這班使徒.
是很肯定知道將來.
有這樣的榮耀獎賞.
是神國度的冠冕.
我們很少人渴慕去爭冠冕.
教會我們爭的是自己的義氣.
我永遠是對的.
你們永遠是錯的.
每間教會都是這樣.
沒有說哪一間的.
我們很少會提.
冠冕是什麼呢?.
沒有人會問這些問題的.
但是我們讀這裡已經有了.
是什麼呢?.
豈不是主耶穌來的時候.
仍然都是你們.
你們就是說是我們的福音果子.
你們不僅不是不得救.
而且是能夠在主耶穌的面前站得住.
所以因為你們就是我們的榮耀.
就是我們的喜樂.
就是說在耶穌來的時候.
耶穌會稱讚這班帖撒羅尼加人.
他們的成長.
其實稱讚這班人的成長.
就是稱讚保羅這班使徒的奴婦.
那個服侍的盡忠.
來去養他們長大.
原來他們的成長.

$^{1961}$就是我們的榮耀.
是我們的喜樂.
所以冠冕是什麼呢?.
保羅說.
就是你們成長.
我們就有榮耀.
我們就有喜樂.
我們真的要追求.
一種這樣的他者的成長.
是我們的榮耀.
大家走過一些街.
沿路這麼多教會.
你看到一些教會.
你就要為他們祈禱.
因為你有榮耀的.
說希望他們坐滿人就好了.
真心祈禱.
希望教會更加多人.
希望傳到人們講到很振奮人心.
你這樣祈禱.
就是為人祝福.
人們得到榮耀.
其實最終都是那個祈禱的人.
得到滿足.
保羅說.
我們受苦有什麼所謂.
最重要是你們在主耶穌回來的時候.
你們不是全部都進不了神的國.
而且是進得了.
而且還是很榮耀地進去.
被主耶穌稱讚.
稱讚你不就是稱讚我.
保羅的意思.
你們得到榮耀不就是我得到榮耀.
所以保羅重要的人生的觀念.
其實就是追求他所做的.
去到主的面前.
又是那句.
又良善又忠心.
又良善.

$^{2001}$良善就是說.
你是一個很善良.
還有你是一個良好的人.
就是說你的人生事奉.
你的動機.
是人盡皆知的美好美麗.
就是人民的眼睛都是雪亮.
看到你最終這麼多十年.
原來你是鞠躬盡瘁為神的福音.
這些叫良善.
忠心就更加明.
就是始終如一.
不會懈的.
這樣的一個擺上.
所以最高的榮耀.
就是良善忠心.
是耶穌說的.
所以保羅這些就是.
他絕對不是去抄福音書.
而是他自己的座右銘.
就是他自己面對逆境.
為什麼能夠屹立不倒的牧者.
就是這樣.
就是因為你們將來.
在主耶穌回來的時候.
能夠站立得穩.
我們就能夠喜樂.
我們就能夠榮耀.
所以這裡再一次第二章說主在來.
主在來就是要在耶穌的面前.
站立得住.
還有一個事奉的人.
是要經過主耶穌的那個測驗.
整個教會沒有人.
只有你一個進天堂.
做的全部沒有一個.
可能真的不關你的事.
是那些人自己不信.
但是保羅說.
如果你們全部都能夠站得住.

$^{2041}$在耶穌的面前.
我們就是那個榮耀.
和我們一生所做的喜樂.
所以官面是什麼呢.
今晚大家都要回去.
坐那個枕頭自己想想.
你的官面是什麼.
為主想得到什麼.
保羅說得到那些人.
能夠在主耶穌回來的時候.
站立得住.
站立得住就是說.
這幫人真是屬神的.
真是在人生無論如何都跟耶穌的.
而且能夠在耶穌的眼裡站立得住.
就是說他們都是為主打拼的.
你能不能將你的信念.
和另外的人.
另外一個的下一代.
他都繼續為主打拼呢.
這個就是傳承了.
福音的傳承了.
就是那個真理.
和那個事奉態度的傳承了.
所以我們期望的下一代.
是不是能夠真的有那個力量.
來一同的為福音.
無論什麼時代.
我們其實不需要太擔心什麼時代.
每個時代都是敵黨神.
敵黨真理,敵黨教會的.
但是我們需要的是有一幫人.
在主耶穌來之前.
我們是站得住的.
我們是為真理站住的.
還有主耶穌真的回來的時候.
我們是在祂的面前的.
還有曾經牧養我們的人.
是會因為我們的站得住.
得到榮耀的.

$^{2081}$所以這裡就大家要一同去追求學習了.
所以這裡第二次講主在來了.
主在來站得住.
主在來你的事奉.
被稱為良善中心.
這些不複雜了.
完全不會覺得是虛無縹緲了.
就是我們人生怎樣為主去事奉了.
不同的人生.
不同的崗位.
不同的起跑點.
不同的人生際遇.
譬如我沒有我小時候那麼多的我區經歷.
也不會講那麼多實際的東西給大家笑一下.
如果我根本是住九龍塘.
不是隔壁的小西湖那區.
是真的住在九龍塘的.
我的故事可能未必適合你聽.
因為我根本從來都養尊處優.
你想想70年代就已經住在九龍塘的.
就不得了了.
每個人都真的有不同的故事的.
所以神用不同的人.
就會幫到很多不同的人.
所以其實現在反而反過來.
很回味小時候的那些經歷.
就成為今時今日.
牧羊上面對別人的一種鼓勵.
因為我都是由一個赤貧那樣成長.
家裡有16個人.
16個人吃飯.
那道菜一轉身就沒有了.
16個人跟你爭.
所以我吃飯是快到不得了.
你不吃飯是沒有了.
就是這樣的.
你可以想像得到.
現在就很好笑.
但當時就很激氣的.
神就是安排我在這樣的環境成長.

$^{2121}$所以現在我們要追求的是.
你牧羊的人.
將來是不是站在神的面前呢.
祂是不是你的榮耀.
祂是不是你的喜樂呢.
這些是很高級的.
你一定要用心去做.
別人才會得到.
所以這個就是.
其實就很考驗我們在人生.
怎樣開始為神一起去做功了.
我們祈禱.
多謝主給我們今晚.
很快又過了一個半小時了.
多謝你給我們透過使徒保羅.
對帖撒羅尼迦教會的血面.
其實是肺腑之言.
講祂的內心的事奉.
講祂面對的困難.
其實祂講的很少.
很少的困難.
但是祂描述的是祂內心的.
事奉的喜樂.
事奉的態度.
和對主的福音的動機的傳承.
但願我們能夠對真理有認識.
我們就加入一個傳福音的行列.
一起在神的國度裡面.
神一定會容我們在不同的恩賜.
在不同的崗位.
願主來幫助.
多謝你給我們今天晚上.
一同學習真理.
奉主耶穌基督的名.
阿們.
\newpage



\section{路加福音 10:25-37}
\label{sec:_2QFVuN84qA}
\textbf{2024-04-13 晚堂《關愛鄰舍再思》路加福音 10\_25-37 講員:陳一華牧師}
\newline
\newline
連結: \href{https://youtube.com/watch?v=-2QFVuN84qA}{\texttt{https://youtube.com/watch?v=-2QFVuN84qA}} ~~~~ 語音日期: 2024-04-16
\newline
\newline
\hyperref[sec:CWxLFOO62pc]{\small{< < < PREV SERMON < < <}}
~
\hyperref[sec:index_chronic]{\small{[返順時目]}}
~
\hyperref[sec:index_scriptual]{\small{[返順卷目]}}
~
\hyperref[sec:iqMY9wa3DY0]{\small{> > > NEXT SERMON > > >}}
\newline
\newline
路加福音 10:25-37
\newline
\begin{longtable}{cl}
\hline
\hline
章節 & 經文 (和合本修訂版)\\
\hline
10:25 & \begin{tabularx}{0.7\textwidth}{X} 有一個律法師起來試探耶穌，說：「老師！我該做甚麼才可以承受永生？」 \end{tabularx} \\ \\ \relax
10:26 & \begin{tabularx}{0.7\textwidth}{X} 耶穌對他說：「律法上寫的是甚麼？你是怎樣念的呢？」 \end{tabularx} \\ \\ \relax
10:27 & \begin{tabularx}{0.7\textwidth}{X} 他回答說：「你要盡心、盡性、盡力、盡意愛主—你的神，又要愛鄰如己。」 \end{tabularx} \\ \\ \relax
10:28 & \begin{tabularx}{0.7\textwidth}{X} 耶穌對他說：「你回答得正確，你這樣做就會得永生。」 \end{tabularx} \\ \\ \relax
10:29 & \begin{tabularx}{0.7\textwidth}{X} 那人要證明自己有理，就對耶穌說：「誰是我的鄰舍呢？」 \end{tabularx} \\ \\ \relax
10:30 & \begin{tabularx}{0.7\textwidth}{X} 耶穌回答：「有一個人從耶路撒冷下耶利哥去，落在強盜手中。他們剝去他的衣裳，把他打個半死，丟下他走了。 \end{tabularx} \\ \\ \relax
10:31 & \begin{tabularx}{0.7\textwidth}{X} 偶然有一個祭司從那條路下來，看見他就從另一邊過去了。 \end{tabularx} \\ \\ \relax
10:32 & \begin{tabularx}{0.7\textwidth}{X} 又有一個利未人來到那裡，看見他，也照樣從另一邊過去了。 \end{tabularx} \\ \\ \relax
10:33 & \begin{tabularx}{0.7\textwidth}{X} 可是，有一個撒瑪利亞人路過那裡，看見他就動了慈心， \end{tabularx} \\ \\ \relax
10:34 & \begin{tabularx}{0.7\textwidth}{X} 上前用油和酒倒在他的傷處，包裹好了，扶他騎上自己的牲口，帶他到旅店裡去，照應他。 \end{tabularx} \\ \\ \relax
10:35 & \begin{tabularx}{0.7\textwidth}{X} 第二天，他拿出兩個銀幣來，交給店主，說：『請你照應他，額外的費用，我回來時會還你。』 \end{tabularx} \\ \\ \relax
10:36 & \begin{tabularx}{0.7\textwidth}{X} 你想，這三個人哪一個是落在強盜手中那人的鄰舍呢？」 \end{tabularx} \\ \\ \relax
10:37 & \begin{tabularx}{0.7\textwidth}{X} 他說：「是憐憫他的。」耶穌對他說：「你去，照樣做吧！」 \end{tabularx} \\ \\
[1ex]
\hline
\hline
\end{longtable}
$^{1}$弟兄姊妹晚安.
多謝主席剛才讀了一段頗長的經文.
這個故事我相信信耶穌的.
未信耶穌的.
對這個很失馬的事情.
有人所做的事情都是耳熟能詳的.
特別我們基督徒對耶穌講這個比喻.
是更加非常熟悉的.
所以我今晚特別來跟大家分享.
一個你們熟悉到背回頭.
都這麼熟悉的故事.
看看天父在當中.
有什麼東西是讓我們能夠學習的.
我們開始之前.
一起我們有個祈禱.
親愛的天父.
願你獨人口中的說話.
心裡的意念.
在你的面前蒙月立.
猶言天父上帝聖靈的光照.
叫講的聽的同感一靈.
同德造就.
奉耶穌的名求.
阿們.
這個律法師雖然不太知道.
他叫什麼名字.
不過這個不要緊.
但他來到耶穌面前問的問題.
是一條明知故問的問題.
為什麼說是一條明知故問的問題呢?.
因為這條問題的答案.
他已經十拿九穩了.
這條問題的答案.
也都在他的專業範圍裡.
這條問題的重心.
他也都不用再翻查答案.
隨口都背得出來.
因為他問耶穌的一個問題.
就是一個好像是心靈上.
或者信仰上的問題.

$^{41}$他說夫子.
又或者翻譯做老師.
他說我應該怎樣做.
才可以承受什麼?.
永生.
這條問題問得好.
問得好有深度.
怎樣做才能夠承受永生.
就是說耶穌說.
我這條問題跟你是有密切關係的.
因為我這條問題是靈性的問題.
不是科學的問題.
不是一些其他的法律的問題.
這條問題跟你有密切關係的.
是你才能夠給我答案的.
各位電視節目.
實際上他都不需要耶穌給答案.
因為他自己老祖就已經擁有答案.
所以當耶穌要回應他的時候.
耶穌就問他.
律法司上記的是什麼.
你讀的又是什麼呢?.
就是說耶穌信他的意思是什麼呢?.
他就純粹推奏給他有機會發揮.
這個律法司聽完耶穌的回應.
二話不說就將整個舊約的律法的精義.
總括為兩點.
大家都很熟悉的.
第一點.
律法的精義不外乎什麼?.
盡心,盡性,盡義,盡力愛主你的神.
第二點.
愛,鄰舍,如同,自己.
所以這兩條問題在他的心目當中已經掌控了.
不過大家留意.
他如果停在這裡都沒有什麼問題.
為什麼呢?.
因為他回應耶穌這條問題的答案是100分的.
拿了100分的.
拿distinction的.

$^{81}$不過最笨拙的是他多嘴.
他多問一條問題.
他多問第二條問題.
就露出他的馬腳.
大家知不知道他第二條問題問什麼?.
第二條問題他問什麼呢?.
這些經文我們鄭華都讀過了.
對了.
他多問一條問題.
這條問題是什麼呢?.
誰是我的鄰舍?.
就是說耶穌說我的答案這麼標準.
這麼高分.
他為了要顯出自己更加有理由.
他推前多問一條問題.
誰是我的鄰舍呢?.
很多時候人的失敗就在這裡.
已經是去到自己很顛峰的狀態的時候.
天下無敵的時候.
他還沒有滿足.
他更加要顯出自己那種霸道.
或者那種王道.
他多問一條問題.
就是問什麼呢?.
誰是我的鄰舍?.
第一次每次問這個問題.
不問又是可惡.
一問就看出他的弱點在哪裡.
因為你稍微精準一點.
分析他這個問題是什麼意思呢?.
誰是我的鄰舍.
你要記得.
問題出現了什麼呢?.
我這個字.
大家有沒有留意到?.
誰是我的鄰舍?.
換句話來說.
我律法師.
作為一個律法師.
誰是我的鄰舍?.

$^{121}$就是說我是主角.
我是主體.
我是主導.
大家有沒有留意到?.
因為有誰可以做我的鄰舍呢?.
所以那個的重心.
中心點就是什麼?.
就是我這個字.
他想著自己成功的地方.
是我成功了.
但是敗也敗在什麼呢?.
就是我這個字.
各位弟子姊妹們.
當他這樣去問的時候.
問題就出現了.
就是說當我要去關懷鄰舍的時候.
當我要來做符合聖經的教訓.
愛鄰舍如同自己的時候.
現在因為人性的軟弱.
就暴露了我的問題.
露出我的馬腳.
原來人本性的驕傲是什麼呢?.
就是表露無遺.
大家有沒有留意到?.
所以你看到.
誰是我的鄰舍.
接著演變下來是什麼呢?.
哪一個有資格做我的鄰舍.
大家有沒有留意到?.
做得我鄰舍的人.
這個人對我的肺腑有什麼呢?.
則貴樂.
或者和我的身份都相稱.
或者他的社會地位.
都有他被人尊重的地位.
或者他有某些專業.
就算不是律法書.
可能他也是一個精算師.
可能他是一個建築師.
可能他是一個工程師.

$^{161}$大家明白這個意思嗎?.
誰是我的鄰舍呢?.
換句話說.
有資格做我鄰舍的人.
應該是半斤八兩.
然後他才有資格做我的鄰舍.
好,我們再看第三.
當我說到.
誰是我的鄰舍的時候.
換句話說.
我才有最後的決定權.
選擇誰是我的鄰舍.
所以從律法師的角度來看.
簡約公屋的人.
不應該放在我們感受的村口.
應該放在藍生圍那邊.
就會比較適合他們.
為什麼呢?.
因為簡約公屋裡的基層的人.
那批草根的階層.
那一班新移民.
應該去遠一點的地方.
因為他們沒有資格.
放在我的門口.
放在我的村口.
做我的鄰舍.
大家明白這個意思嗎?.
這個就是我們看到.
人性的軟弱就顯露出來了.
原來鄰舍不是別人決定的.
是我決定的.
是我選擇的.
所以各位電視節目們.
在律法師多口多問的時候.
結果就出現了問題.
耶穌也忍住了一口氣.
氣定神閒.
跟他說了一個故事.
沒有直接回答這個問題.
耶穌說的故事.

$^{201}$接下來大家就耳熟能詳.
滾瓜難熟.
我不需要在這裡再重複了.
是一個很撒瑪利亞人的故事.
大家很清楚知道.
那個焦點就是.
這個撒瑪利亞人.
他怎樣去關懷扶持.
去救援這個重傷的人.
然後把他扶在他的身束上.
帶他去酒店.
不是酒店 說錯了.
那時候還沒有酒店.
應該是客棧.
來到那裡休息休養一下.
讓這個受傷的人.
好好休養休息之後.
然後他告訴客棧的老闆.
他住多少天.
不要緊.
這裡有兩千元.
如果不夠的話.
我回頭會還給你.
這是什麼.
大家很熟悉的.
整個撒瑪利亞人的故事.
就在這裡.
接著耶穌說完故事之後.
又是什麼.
順便又問一條問題.
這個時候才正式回答.
這個律法師的問題.
接著耶穌問一條什麼問題呢.
耶穌的反問.
就是哪一個是受傷者.
落在強盜手中的鱗石.
大家有沒有留意到.
第一條問題和第二條耶穌的反問.
是不是有很大分別.
看不看得出分別在哪裡.

$^{241}$原來第二條的問題.
是180度改變過來.
完全不同了.
因為耶穌問的第二條問題.
是問受傷者的鱗石是哪一位.
如果這樣問的時候.
弟子們.
你發覺.
角色就改變了.
角色改變了.
就是說我不再是一個照顧人.
關心人 關懷人的這樣一個人.
我再不是一個重要人物.
不再是一個主體.
不是一個主導 不是一個主角.
如果按這條問題去問.
我關懷人的角色.
不過是一個配角.
主角是誰.
主角是受傷者這個人才是主角.
如果是這樣的時候.
我是一個客體的話.
我發覺我的表現.
我要表達出來的.
和受助者.
就要有一個對稱的比例在當中.
所以在這個時候.
我的態度.
我的表現.
和剛才第一條問題就很不同了.
為什麼呢?.
因為現在的處境是改變了.
所以我的角色都改變了.
接下來.
就是態度的改變了.
這個態度的改變.
不同之前.
誰有資格做我的鱗石.
誰是我看得上眼.
我就選擇他做我的鱗石.

$^{281}$現在不是了.
現在是受傷者.
他是我要關懷的對象.
所以他才是最重要的對象.
因為這樣的緣故.
他是一個中心人物.
而我呢.
不過是一個輔助的角色.
開始我要問一個問題.
到底這個受傷的朋友.
他的需要是什麼呢?.
各位節目們.
可能我們遲些會問.
這一批住在簡約公屋裡的朋友.
到底他的需要是什麼呢?.
就不是說我將我那套東西.
強加於在他的身上.
不是說我認為怎樣去.
傳遞訊息就由我決定.
或者我做什麼行動呢.
由我決定.
不是這樣的.
我們應該去關心之前.
多些了解一下.
到底這一班我的鱗石.
他在我的附近.
他住在我的旁邊.
是我的鱗石.
他操到跪頸.
人家伸梳.
他有什麼需要呢?.
這些家庭有什麼問題.
是一個最常見的問題.
或者這些鱗石.
他們最急切需要面對.
是什麼問題.
什麼需要呢?.
可能這些我們需要什麼呢?.
多些了解.
或者可能我們切身處地.

$^{321}$透過一些探訪.
透過一些行動關懷.
了解他們多些.
認識他們多些.
然後當我們真的要去關心.
關懷他們的時候.
我們提供的服務.
或者我們提供的行動.
對他們來說.
才是一個真是落地的.
有地氣的一種關懷的行動.
所以各位電影節目.
我們發覺到.
我們的態度.
再不是高高在上的一個態度.
而是一個平等的.
而是一個謙卑的.
而是一個願意忍耐去聽聽.
對方有什麼要傾訴.
或者有什麼想和你分享.
我們忍耐去聽聽.
以致我們更加知道.
明白原來他們的處境是這樣的.
原來他們發覺.
來到這個地方真的很不習慣.
每天坐車的交通時間很長.
沒什麼機會.
沒什麼時間照顧子女.
如果你們教會有人幫忙補習.
讓小朋友有機會.
可以在他們課餘時間裡.
都有人和他們玩玩.
或者關心他們的功課.
關心他們的成長.
這是他們的需要都未敵.
各位電影節目.
我舉一個例子.
意思就是.
我們的態度是謙卑的態度.
好像什麼都不熟悉.

$^{361}$我要重新去學習.
重新去了解.
知道他們的處境是什麼.
然後我就可以真正進入到他們的世界裡.
第三個改變.
這個二應該是三.
第三個改變是什麼.
是身份的改變.
這個身份的改變就是說.
當我們無論錦宣家也好.
無論家欣堂的弟兄姊妹也好.
我們去到簡約公屋.
去到部屋的地方的時候.
原來這一批的人.
雖然我們不認識他.
雖然他們的學歷背景.
沒有我們那麼高.
雖然他們的專業程度.
沒有我們那麼明顯.
雖然他們的經濟能力.
都很大的缺乏.
雖然很多不同的地方.
但是電影節目你記得.
原來當我們用一個.
調轉了的角色.
以關懷 以受助者為中心.
那個角度去做的時候.
你發覺到我們要去關心他們的時候.
我們的身份就什麼.
就不再是我們一個高高在上的身份.
而是一個跟他們很平等.
跟他們大家都好像這麼高這麼大.
大家都是一個居民的角色.
來跟他們做朋友.
雖然我們的背景是一個中產的階級.
在社會上別人會看我們錦繡花園的居民.
是不同的角度.
但是當我們關心他們的時候.
我們的身份就要改變了.
我們怎樣以一個朋友的角色.

$^{401}$朋友的身份.
等他們接觸我們.
我們接觸他們的時候.
他不會感覺到自卑.
他不會感覺到好像有高低之分.
他會感覺到我們是一個.
關心他們的友善的朋友.
來關懷他們.
所以我們無論有什麼身份.
在那一刻我們都學習什麼.
學習放下這一切.
能夠跟他們一起.
好像同舟共濟.
所以這就是第三個改變.
第四個改變就是責任上的改變.
責任上的改變就是說.
當我以一個受助者.
作為我一個主要的關懷對象的時候.
這個時候我發覺到.
我關心他們是責無旁貸的.
所以我要學習主動一點去關心他們.
不是被動.
你有什麼需要.
你過來吧.
你有什麼需要我們幫忙.
你過來我們這裡談談吧.
不是的.
我們應該去到他們當中.
去跟不同的途徑去關懷他們.
因為這種關心是一個主動的.
是一種自發的.
是我的責任.
如果我不關心他們.
是我失職.
大家明白這個意思吧.
因為上帝已經將這一班的鄰舍.
放在我們的左鄰右鄰.
所以我們有一個責任是什麼.
要採取主動的行動.
因為我們住在這個地方.

$^{441}$感受花園的日子比他們久.
我們熟悉這個環境比他們多.
我們了解如何能夠適應這裡的生活.
也比他們多.
所以從任何角度來看.
我們關懷他們是一份責無旁貸的職責.
所以我們能夠用這種改變的態度去看的時候.
我們就發覺到我們做這個撒瑪利亞人的角色.
就不再用自我,自己做中心了.
而是以一個受傷的人,受助者作為中心.
當我們了解了主次的角色之後.
接著可能我們要問.
耶穌說的這個故事.
到底有沒有一些什麼.
在關懷人的行動上.
給我們一些啟發呢.
我發現到是很寶貴的.
原來耶穌在關懷人這個撒瑪利亞人的身上.
告訴我們原來關心人的途徑.
有六個道是很重要的.
我們今晚試試看有哪六個道.
第一個道是什麼.
是眼道.
上帝讓我們看到.
現在你坐車也好.
開車出入感受也好.
你一去到油站.
還沒過閘口.
你已經發現旁邊已經是什麼.
已經是大興土木了.
是不是.
遲些那裡一建好的簡約公屋.
有多少個單位.
如果你上個星期留心PowerPoint.
原來有2100個單位.
在我們村口的地方.
這2100個單位.
他們住多少人.
起碼有5000,6000或者7000個人.
在那個地方.

$^{481}$等於節目無人.
你沒有眼看到.
為什麼上帝會無端端.
也不應該無端端.
不是無端端的.
上帝一定有祂的目的.
上帝一定有祂的心意.
也不是早不是遲.
是剛剛這個時候.
讓我們感宣家可以很意外地.
發現到原來有一班新的鄰居朋友.
將要出現在我們旁邊.
在我們身邊.
你有沒有看到這是一個廣大的禍場.
你有沒有看到這六,七千人的靈魂得救的問題.
你有沒有看到這六,七千人家庭所需要的東西.
你有沒有看到他們在情緒上.
在心靈上有一個很大的需要.
除了物質之外.
等於節目無門.
眼到是很重要的.
所以你現在明白為什麼耶穌行神蹟之前.
祂為什麼會出手去行神蹟呢?.
除了祂動了慈心之外.
耶穌才是什麼?.
才看到眼前受助者的需要.
無論祂看到的是麻風的病人.
無論祂看到的是盲的人.
無論祂看到的是患了血流的婦人.
當祂看到人的光景的時候.
等於節目無門.
祂發覺這是祂義不容辭.
要去關心的對象.
同樣今天不得不上帝給我們加恩堂.
給我們感宣堂.
看到這麼遼闊的和場.
上帝已經安排在我們身邊.
讓我們再一次來迎接挑戰.
迎接這一批新的朋友.
讓他們的心靈在這裡過渡期間.

$^{521}$都能夠得到福音的好處.
得到教會 認識教會.
是一個在他們人生的旅程中.
一個不能夠忘記的心路歷程.
所以眼督是很重要的.
我想感宣家弟兄姊妹.
每一個人 大家都看到這種需要.
看到這種挑戰.
所以我們下星期日就要報行籌款.
大家記不記得.
有多少人已經拿了表格的.
舉手給我看看.
有一半人已經拿了.
還有一半人還沒有拿.
你要祈禱了.
你要想一想.
為什麼我們要去關心.
去參與這個淪陷的關懷的行動呢?.
弟兄姊妹們.
這個是需要我們眼睛看到的.
不單止看到.
我們的心都要悼.
所以這個撒瑪利亞人.
看到受傷的人.
躺在地上的時候.
很辛苦.
心仁痛苦的時候.
聽而不聞.
有些人是視而不睹.
所以有些人就這樣走過.
但是這個撒瑪利亞人沒有就這樣走過.
他是停下腳步.
他的心動了.
他看到受傷的朋友這麼淒涼.
這麼可憐.
遍體鱗傷.
差不多半條命都沒有了.
那一刻他的心有什麼呢?.
就動了持心.
各位弟兄姊妹們.

$^{561}$原來上帝又讓我們看到一個福音禍場.
在這麼大的時候.
我們的心不單止開始心動了.
我們都要開始什麼呢?.
心要感恩.
要興奮.
因為上帝在我們有生之年.
都有一次能夠開方的機會.
雖然我順著耶穌到現在還沒有開方.
沒有試過怎樣去做傳福音的工作.
但是現在已經有一個這麼好的禍場.
在我的附近.
讓我有很好的學習.
很好的參與.
所以弟兄姊妹們.
孫家的每一位.
我們都能夠有一個心.
能夠同樣.
同這班人大家能夠可以心連心在一起.
接著第三個關懷的途徑是什麼呢?.
是一個守道.
因為你看到這位撒瑪利亞人.
就將酒和油倒在受傷者的傷口那裡.
然後就將它清潔.
將它消毒.
將它包紮好.
然後扶起它在牢區上面.
弟兄姊妹們.
我們看到整個行動當中.
守道是非常之重要.
給我們一個很好的學習.
就是說如果我們關懷我們的淪陷的時候.
我們不是勢必要講很多話的.
又或者我們不是勢必要會講很多話.
我們口才要頂級.
完全都不需要.
為什麼呢?.
因為現在的人是很聰明的.
他知道你是不是真的關心他.
是不是知道你是不是真的有誠意呢?.

$^{601}$他不是聽你說多少話.
而是看你什麼.
有多少行動.
有多少實際你付出的行動.
讓他體會到.
原來基督徒那種愛心.
原來就是學校你們的耶穌那種無私的愛.
各位弟兄姊妹們.
這個是我們透過我們的實際對關懷者所做的.
讓他們能夠更加經歷.
原來上帝的愛是很真實的.
第四個到就是腳要倒立.
大家看到這幅圖畫就知道.
好撒馬利亞人自己走路.
拖著爐區自己走路.
然後就將傷者放在生畜上面.
一路走走走走到入村的地方.
去到有客棧的地方.
其實在過程當中.
他都不是說什麼話.
你有沒有留意到.
他只是默默在做.
默默陪伴著受傷的人.
一路盡他的力量.
盡他的能力來關心他.
一路去到客棧的地方.
所以各位弟兄姊妹們.
將來我們關懷的鄰舍.
不是去一次就完.
可能我們去兩三次.
四五次甚至去十次.
十幾次都需要的.
為什麼呢.
因為我們要多一點主動的.
跟他們有接觸.
有交往.
有交流在當中.
這個行動是一個美好的行動.
這個很現實的了.
弟兄姊妹們.

$^{641}$原來還要錢都要倒了.
當我們關心人的時候.
不是每天都說.
怎樣關心你.
怎樣愛你.
怎樣關懷你.
但是一說到錢就退了.
就縮了.
不是的.
原來我們連錢都要倒.
所以這位撒馬尼亞人.
他拿了兩千銀子出來.
就是當時打工仔一天的工資.
他拿了出來.
交給老闆.
黑店的老闆.
告訴他.
這個就是房租.
如果他住得久.
不要緊.
我會再給你錢.
再還錢給你.
所以弟兄姊妹們.
現在大家明白.
為什麼我們有步行籌款了.
是不是.
為什麼呢.
因為我們知道.
當我們有計劃.
要去關懷這些鄰舍的時候.
我們需要籌募一些經費.
讓我們這個行動計劃.
有好的部署.
有好的安排.
有人手.
有物力.
有資源.
將來能夠去到這些鄰舍的手中.
所以我們每一個人.
在這個時候.

$^{681}$我們做好準備.
這次我們知道.
簡約公屋也好.
部屋也好.
都不是五月份就開始了.
而是什麼.
而是可能是九個月.
或者半年之後就開始.
所以仍然都給我們時間.
讓我們做好多些的安排.
計劃準備.
所以這個步行籌款.
又讓我們教會.
能夠在見到弟兄姊妹的愛心關懷當中.
或者你將這個意象.
將這個挑戰.
告訴你身邊相熟的朋友認識.
都讓他們知道我們教會.
準備這個行動計劃.
要關愛我們的鄰舍.
所以弟兄姊妹們.
這個就是第五個榜樣.
第六個榜樣是什麼呢.
就是賢道.
賢道就是說.
這位好撒馬利亞人.
他會不會回頭.
他會的.
他言出必行.
我雖然現在要走了.
趕著要做其他事.
不過我一定會回來.
而且我回來一定會還錢給你.
所以這個賢道.
就是說是有承諾的.
說得出就是什麼.
做得到的.
各位弟兄姊妹們.
對我們來說.
關懷我們的鄰舍.

$^{721}$關愛我們的鄰舍.
其實我們是需要.
做多一些跟進的工作.
我們要有多一些後續的工作.
我們不是只去一兩次探訪.
就完了我們的關心.
可能是一段比較長時間.
我們是繼續的.
來關懷這些有需要的家庭.
或者關懷那個小社區.
這個是我們教會.
不單只是我們一個挑戰.
是我們的責任.
也是我們義不容辭.
要做的跟進栽培的工作.
讓我們今天晚上.
看到這個故事的時候.
好撒馬利亞人的故事.
讓我們大家都能夠有一個.
很好的機會再想一想.
反思怎樣去關愛我們的鄰舍.
老祖聖經已經有一個.
很好的榜樣讓我們去學習.
不得不得我們弟兄姊妹們.
從今晚經文裡面所提及的.
就讓我們能夠.
好好地去裝備自己.
我們一起同心低頭祈禱.
我們讚美你.
高過我們的道路.
你的旨意也都高過我們的旨意.
你在我們最合適的時候.
你就預備了我們有一個.
非常之遼闊的一個屬靈的和場.
就是有一批的鄰舍.
很快就搬進來.
我們感受花園.
做我們的好朋友.
或者我們部屋的朋友.
也都會很接近我們的嘉欣堂.

$^{761}$將這兩批不同的對象.
在這個時刻.
來到出現在我們教會身邊的時候.
我們相信不會是偶然.
不是這麼湊巧.
上帝你對我們每一個屬你的兒女.
一個計劃.
一個屬靈的挑戰.
讓我們能夠看到.
你就是莊家的主.
是你打發工人出去收你的莊家.
多謝主.
不是我們揀選你.
而是你揀選了我們.
給我們一個上好的福分.
無比的恩典.
能夠去學習成為鄰舍的關懷者.
奉耶穌的名頭.
阿們.
拜拜!.
\newpage



\section{申命記 1:1-18}
\label{sec:iqMY9wa3DY0}
\textbf{2024-04-14 午堂《同行共拓主國度》申命記 1\_1-18 講員:戴浩輝牧師}
\newline
\newline
連結: \href{https://youtube.com/watch?v=iqMY9wa3DY0}{\texttt{https://youtube.com/watch?v=iqMY9wa3DY0}} ~~~~ 語音日期: 2024-04-16
\newline
\newline
\hyperref[sec:_2QFVuN84qA]{\small{< < < PREV SERMON < < <}}
~
\hyperref[sec:index_chronic]{\small{[返順時目]}}
~
\hyperref[sec:index_scriptual]{\small{[返順卷目]}}
~
\hyperref[sec:y77Sfudracc]{\small{> > > NEXT SERMON > > >}}
\newline
\newline
申命記 1:1-18
\newline
\begin{longtable}{cl}
\hline
\hline
章節 & 經文 (和合本修訂版)\\
\hline
1:1 & \begin{tabularx}{0.7\textwidth}{X} 以下是摩西在約旦河東的曠野，疏弗對面的亞拉巴，就是在巴蘭、陀弗、拉班、哈洗錄、底撒哈之間，向以色列眾人所說的話。 \end{tabularx} \\ \\ \relax
1:2 & \begin{tabularx}{0.7\textwidth}{X} 從何烈山經過西珥山到加低斯‧巴尼亞要十一天的路程。 \end{tabularx} \\ \\ \relax
1:3 & \begin{tabularx}{0.7\textwidth}{X} 第四十年十一月初一，摩西照耶和華所吩咐他一切有關以色列人的話，都告訴他們。 \end{tabularx} \\ \\ \relax
1:4 & \begin{tabularx}{0.7\textwidth}{X} 那時，他已經擊敗了住希實本的亞摩利王西宏和住亞斯她錄、以得來的巴珊王噩。 \end{tabularx} \\ \\ \relax
1:5 & \begin{tabularx}{0.7\textwidth}{X} 摩西在約旦河東的摩押地講解這律法，說： \end{tabularx} \\ \\ \relax
1:6 & \begin{tabularx}{0.7\textwidth}{X} 「耶和華－我們的神在何烈山吩咐我們說：你們住在這山上已經夠久了。 \end{tabularx} \\ \\ \relax
1:7 & \begin{tabularx}{0.7\textwidth}{X} 要起行，轉到亞摩利人的山區和附近的地區，就是亞拉巴、山區、謝非拉、尼革夫、沿海一帶，迦南人的地和黎巴嫩，直到大河，就是幼發拉底河。 \end{tabularx} \\ \\ \relax
1:8 & \begin{tabularx}{0.7\textwidth}{X} 看，我將這地擺在你們面前。你們要進去得這地，就是耶和華向你們列祖亞伯拉罕、以撒、雅各起誓要賜給他們和他們後裔為業之地。」 \end{tabularx} \\ \\ \relax
1:9 & \begin{tabularx}{0.7\textwidth}{X} 「那時，我對你們說：『我獨自一人無法承擔你們的事。 \end{tabularx} \\ \\ \relax
1:10 & \begin{tabularx}{0.7\textwidth}{X} 耶和華－你們的神使你們增多。看哪，你們今日好像天上的星那樣多。 \end{tabularx} \\ \\ \relax
1:11 & \begin{tabularx}{0.7\textwidth}{X} 惟願耶和華－你們列祖的神使你們更增加千倍，照他所應許你們的賜福給你們。 \end{tabularx} \\ \\ \relax
1:12 & \begin{tabularx}{0.7\textwidth}{X} 但你們的擔子，你們的重任，以及你們的爭訟，我獨自一人怎能承擔呢？ \end{tabularx} \\ \\ \relax
1:13 & \begin{tabularx}{0.7\textwidth}{X} 你們要按著各支派選出有智慧、明辨是非、為人所知的人來，我就立他們為你們的領袖。』 \end{tabularx} \\ \\ \relax
1:14 & \begin{tabularx}{0.7\textwidth}{X} 你們回答我說：『你說要做的事很好！』 \end{tabularx} \\ \\ \relax
1:15 & \begin{tabularx}{0.7\textwidth}{X} 我就將你們各支派的領袖，就是有智慧、為人所知的人，立他們為領袖，作你們各支派的千夫長、百夫長、五十夫長、十夫長等官長，來管理你們。 \end{tabularx} \\ \\ \relax
1:16 & \begin{tabularx}{0.7\textwidth}{X} 「當時，我吩咐你們的審判官說：『你們聽訟，無論是弟兄之間的訴訟，或與寄居者之間的訴訟，都要秉公判斷。 \end{tabularx} \\ \\ \relax
1:17 & \begin{tabularx}{0.7\textwidth}{X} 審判的時候不可看人的情面；無論大小，你們都要聽訟。不可因人而懼怕，因為審判是神的事。你們當中若有難斷的案件，可以呈到我這裡，讓我來聽訟。』 \end{tabularx} \\ \\ \relax
1:18 & \begin{tabularx}{0.7\textwidth}{X} 那時，我已經把你們所當做的事都吩咐你們了。」 \end{tabularx} \\ \\
[1ex]
\hline
\hline
\end{longtable}
$^{1}$願因平安 從我們天父上帝並救主耶穌.
歸於在座的每一位弟兄姊妹.
阿們.
你陌生的面孔.
我是第一次在大家當中分享上帝的話語.
戴皓輝是我的名字.
我是兩個成年的女兒的爸爸.
我有兩個孫.
我人生下一年就到一生的年日了.
大家知道什麼叫一生的年日嗎.
七十歲.
很感恩.
如果我是強壯的話.
是不是可以活到八十 我不知道.
但是上帝的恩典夠我用.
什麼時候我說不到就說不到.
今天很想和大家分享一個題目.
叫做同行共創主角度.
是分享申命記的第一章1到18節.
在過去幾年.
我們好像摩西走去曠野帶著以色列人.
為什麼呢.
因為我們面對疫情 社會事件.
很多的政策改變了.
加上我們現在開始面對很多經濟上的困難.
有時候我們自己覺得我們不知所措.
按我自己教會.
至少我自己是屬於信義會的.
我們有很多人其實都離開了.
那些事奉人員就大減了.
經濟當然都開始出問題.
當我們在今天去思考同行共創主角度的時候.
可能我們都會看到.
是不是過去這幾年才令到香港是這樣的呢.
我們的教會是不是因為過去幾年的原故意戰.
它開始是這樣呢.
我告訴大家其實不是的.
或者我們可以看一下這個統計表.
如果你看到這個統計表.
從2014年開始.

$^{41}$應該就是我們說香港教會崇拜人數最高的時候.
有30萬多一點.
如果你看看其實在2009年.
它還是29萬.
但是2014年就30萬.
不過可以說那個趨勢一直一直這樣跌下來.
包括了2019年已經跌到26萬.
2021年5月的統計已經跌到只有18萬.
我們這裡其實都還沒有計算那些在過去幾年的當中.
移民的弟兄姊妹.
其實很多都是年青的青壯年的.
有家庭的多於移民的.
我再看一下一個年齡的分佈圖.
你就會體會在這個分佈圖當中告訴我們.
香港教會的人已經踏入老年的時刻.
你可以看一下1999年.
那個15歲至44歲.
加起來的時候差不多是60\%.
但是在2019年的時候.
加起來的時候只有3\%.
即是退了25\%.
相反的45歲至65歲以上.
你就會看到本來只有25\%.
但是在2019年已經進入超過一半.
從這個統計當中告訴我們一件事.
香港教會一直都是衰敗.
數字不講大話.
我們香港教會你可以說正在衰落.
我們很多青壯年的人.
他們就離開了.
連同他們的家庭換句話來說.
下一代甚至連零到14歲的人都可能減少.
在這樣的情況當中.
教會的支柱特別是事奉的人都沒有了.
經濟上很大的可能.
我們開始面對慢慢的困難.
這個是香港教會的現象.
我不知道桂堂是不是這樣.
但是我想至少香港的教會現象是很普遍的.
在這個困惑的時代.

$^{81}$有些人去比喻香港現在好像是一個被擄的時候.
教會被擄的時候.
不過我很想不是那麼消極去看這件事.
我很想看香港的教會.
其實就好像摩西進入迦南之前的那個情況.
換言之我要看到的是.
究竟上帝給我們的未來是怎樣的呢.
換言之就是我們好像摩西帶領以色列民.
經過了曠野的掙扎.
到了摩西這個摩亞的平原.
準備進入迦南美地.
當以色列人到了摩亞地.
準備要進入上帝所應許的地方的時候.
他們需要些甚麼呢.
如果我們說今天我們是一個那麼衰敗.
好像在曠野的時候.
我們需要些甚麼可以進入到.
上帝應許給我們美好的地方.
換言之我們怎樣可以真是.
同行共託主的國度呢.
今天我們讀申命記的第一章.
1到18節可以說是申命記的開始.
其實在這裡提醒了我們有三個方向.
我們可以看到我們怎樣可以.
同行共託主的國度.
第一件事 以史為鑑.
申命記是五經的最後一卷.
五經有些人翻譯為律法書.
其實它在希伯來的意思就是.
「陀拉」意思就是指導的意思.
換句話來說.
五經或者律法書.
它是直接律法.
直接勸告,祝福,咒詛.
來指導將來要進入應許之地.
得為業的上帝的子民.
在這個申命記當中.
你可以說也是摩西臨終的時候的叮嚀.
好像申命記你讀到尾的話.
你會讀到摩西死的.

$^{121}$由頭到尾都是摩西的講話.
在這一切的講話當中.
好像最後的叮嚀一樣.
但是同樣的.
這是一個歷史的回顧.
英文有一句說話.
.
中文就是那些沒有從歷史當中學習的人.
會注定重蹈覆轍.
當我們看聖經第一節至第五節的時候.
我再讀一次.
.
我們讀舊約聖經.
有一個不是很中意的例子.
在第五節裡面.
那些名字太多.
在第一節,第二節.
詳細記錄地名.
不單單給信徒會覺得.
都不知道它在哪裡.
都不知道在講甚麼.
甚至很多解經家.
其實都面對過困難.
因為這些名字.
如果你說現在的考古學.
根本你對不到在哪個位置.
在教學當中.
那究竟這個申命記的開始.
它是不是只是說了一些地名呢.
還是這些地名有其他的意思在裡面呢.
當我們細心去想的時候.
我們就不難發現到.
其實這些地名.
都是歷史的背景.
是甚麼歷史.
是成功和失敗的歷史.
交織而成的一個背景.
你可以說是上帝指紋歷史的回顧.
就以巴蘭哈西路為例.
其實代表了甚麼呢.

$^{161}$如果我們詳細去讀民數記的.
十一章到十四章的話.
我們就知道這兩個地方.
其實是以色列失敗的歷史場景.
在哈西路.
因為米利安和亞倫心裡面驕傲.
不滿意摩西的大寧.
於是米利安就長了大麻蜂.
所以這個地名哈西路.
其實你可以說有另一個意思.
就是要謙卑.
另外巴蘭在以色列人探子回來的時候.
他們報告.
他們做甚麼.
沒有信心.
沒有信心去承繼這個應許之地.
他們發怨言.
所以巴蘭某程度上.
亦都提醒以色列人信心的功課.
最終只有迦勒和約書亞能夠進去.
還有底薩哈這個詞.
底薩哈在希伯來語意思就是金.
在猶太人解經的書叫做Midrash.
它就說底薩哈是指以色列人拜金牛讀的事件.
第二節它說.
由荷列山西爾山的路.
到迦底斯巴尼亞.
有十一日的路程.
不少解經家都認為這個是正確的.
所以是二百五十六公里左右.
不過問題就是.
從西乃山.
從它的路線來說.
二百五十六公里為甚麼可以走了這麼久呢.
多久呢.
四十年.
十一日的路程走了四十年.
換句話來說.
很多冤枉路.
以色列人走了很多冤枉路.

$^{201}$我今天去思考這個問題的時候.
我會和大家去看看.
香港教會的狀況.
其實都很大可能走了很多冤枉路.
最少香港的教會.
今天的警方是退回三十年前.
為甚麼呢.
因為一九九四年有教會崇拜的人數的統計.
是已經有十九萬九千人.
差不多二十萬.
比對二一年的統計.
五月都只剩下十八萬人.
是否已經是退回三十年都不止呢.
這幾年的疫情.
社會事件.
經濟的下滑.
走了幾十萬.
內中其實有不少的教會的支柱.
我們知道.
我都去不同的地方去講道.
很多人現在找事奉人員.
是不容易的.
經濟接下來.
是要擔心的一件事.
回顧這些過去.
面對我們現在的現狀.
其實我們怎樣去看神給我們未來的應許呢.
正如以色列文.
他得應許之地.
他就在這個曠野.
在摩亞地面對這個應許之地.
如果我們從哈西路.
從巴蘭可以學到功課的話.
其實提醒我們.
首先帶領的人.
我們教會的教睦.
我們教會的執事長老們.
其實他們帶領教會是不容易的.
我們的導師們.
他們不是容易的.

$^{241}$我們需要甚麼.
我們需要謙卑的.
有信心的.
跟他們一起去建立.
另一方面.
在第四節當中提了兩個得勝的例子.
提到摩西戰勝了巴山王岳.
和阿摩利王這個西弘.
其實是一個得勝的印證.
這個印證是告訴我們.
指明你是可以得勝的.
好 使道以色列人.
他可以進入神所應許的地方.
可以使我們今天進入.
神給我們未來的地方.
換句話來說.
勝敗得失.
其實都可以成為我們的借鏡.
上帝要我們得他所謂我們預備的地方.
記住這個不是一個地方的問題.
這個預備的地方.
這裡已經很好了.
不是 而是得未得之物.
這個才是最重要.
我們怎樣看到.
神在我們的過去和我們的未來.
帶領著我們得著未得之物呢.
另外當我們繼續去讀.
第一節 第三節 第五節.
我再重複一次.
以下的是摩西在約旦河東曠野.
一直講到的就是.
他所說的話.
跟著第三節.
出埃及記第四十年十一月初一.
摩西照耶和華.
借著他所吩咐的.
以色列人的話.
曉喻他們.
第五節說.

$^{281}$摩西在約旦河東摩雅的地方講解律法.
這三節其實都講到.
摩西是主角.
他的主要的工作是甚麼.
向他們講神的話.
他主要的工作就是.
第一節和第三節.
就是向他們講神的話.
而第五節他就會解釋那個律法.
因此在這三節當中.
猶太人的解釋家就說.
其實這裡是在說猶太人.
每個禮拜都需要讀兩次希伯來文聖經.
然後讀一次亞蘭文.
即是他們的他爾根的.
意譯的聖經.
解釋聖經.
所以他們在他們的猶太人的會堂裡面.
是這樣做的.
其實對我們來說.
對我們來說可以很簡單去體會的就是.
研讀上帝的話.
其實就是我們探索前路的關鍵.
因為神的話提醒我們很多.
我們自己走過的失敗的道路.
和我們體會神給我們得勝的方法.
研讀神的話.
是我們每一個基督徒.
我們每一個人.
承擔神給我們的託付的一個重要的方法.
除了歷史為鑑之外.
這段經文提醒我們第二件事.
就是承擔使命.
對於以色列子民來說.
進入應許之地.
是他們的目標和使命.
為要完成神對列祖的應許.
在第六節裡面.
這裡寫了第五節是錯的.
在第六節裡面說.

$^{321}$你們在這個山的日子已經夠了.
我讀的是一個老的和學本.
這裡翻譯是一個新的修訂版.
從領受上帝話語.
到行動承擔使命.
使到神對列祖的應許可以實現的話.
摩西提醒他們.
你們在這裡已經夠了.
我下一年就要離開杜豐山.
因為我住杜豐山的日子都已經夠了.
意思就是說我要下山.
我住了三十多年.
在我們的日子已經夠了.
這可以說一方面提醒我們.
我們在現在這樣的情況已經夠了.
我們需要做些事.
如果我們說神要賜給我們這一代.
祂的應許的話.
神所立的約和列祖所立的約.
就是這個國度的開始.
不過要去承受這個國才能夠實現的話.
這樣我們可以怎樣做呢.
神不希望祂的子民滿足在這個曠野的生活.
曠野的現況.
祂要我們繼續.
直到我們看到承受應許的地方.
記住不是place.
現在是說另一個理念.
今天上帝也很希望我們滿足.
不要滿足我們的現況.
祂要我們踏出信心的一步.
去承受祂要我們得著的地方.
主耶穌其實一開始祂的事公的時候.
祂第一句話說了些甚麼呢.
我們可以看一下.
馬可福音的第一章第十五節.
祂說日期滿了.
神的國近了.
你們要悔改.
順福音.

$^{361}$其實就好像摩西在第六節裡面所說.
你們在這個山上住的日子夠了.
對於我們現代的教會來說.
我們承擔使命.
作洗萬民造門徒.
其實是一個我們都覺得很自然的事.
不過大家可能不知道.
其實對現今的基督徒來說.
不是所有人都知道.
他們需要承擔使命.
至少我看回2018年.
Bana Research有一個調查的結果告訴我們.
在美國至少有51\%的基督徒.
他們根本不知道甚麼是大使命.
有25\%的基督徒.
他們不確定知不知道甚麼是大使命.
如果加起來的時候.
至少有四分三的基督徒.
很大可能他們沒有傳過福音.
在第八節.
上帝的賜予應許之地.
擺在了以色列人的面前.
神已經給了.
不過問題是有另外兩個動詞.
這兩個動詞就是.
進去得.
進去得為業.
其實是要行動的.
換句話來說.
是需要有信心去承受的.
我相信擺在我們今天面前的.
有好幾樣重要的事.
第一就是.
我們有很多的信徒可能已經移民了.
在他們不同的地方當中.
沒有辦法參與教會的話.
其實我們作為教會的.
在這個互聯網的時代.
我們很應當繼續藉著互聯網去牧養他們.
直到他們可以參與地方的教會為止.

$^{401}$第二就是.
面對不少.
我覺得說從內地移港的同胞.
我們同樣的承擔著這個使命.
或者我們都可能會體會的.
就是香港的轉變.
是在於香港有很多青壯的人離開.
但是我可以肯定告訴大家.
也有很多青壯的人.
會從內地來到這裡.
除了150人.
還要加上我們叫做高材的人員.
那些人仍然都是青壯的人.
他們來到這裡.
人生路不熟.
仍然需要很多的關顧.
不過我覺得他們最重要的一件事.
是需要福音.
他們最重要的是福音.
不過我們看看.
我們這個開始衰落的教會的時候.
我們面對的是一個什麼樣的困難.
我們可以看一下統計.
這個困難是什麼.
你看看這個統計你就知道.
這個困難就是.
神學院全時間的神學生.
在十年裡面一直在下跌.
如果我們看到2009年.
就有1203個全時間的神學生.
2014年已經跌到1039個.
2019年連千位數字都沒有.
929個.
教務的職位需要是逐年增加的.
到現在2024.
我自己看回我自己教會信義部.
缺乏的人的同工至少是23個.
在50間的堂會當中缺少23個.
即是你可以比對一下.
10個茶壺只有6個蓋.

$^{441}$10個6個.
換言之當我們體會在今天的日子當中.
如果大家會覺得.
你們在現在這個情況已經夠的話.
各位年輕人.
想一下神會不會呼召你們參與其中.
第三 當我們體會.
除了我們看到以使為鑒承擔使命之外.
我們看到這段經文告訴我們.
共託未來.
整本申命記.
你可以說是一本有關立約子民的策略書.
從第9節至第18節.
特別去強化了這一點.
摩西語重心長地提醒以色列民.
向神要盡本分.
這段經文的記載.
其實是重複了在出埃及記的18章.
和民數記的11章.
我們可以將這三段經文.
同一個講述.
把它們放在一起的時候.
我們可以發覺有一些很特別的事可以出現.
民數記的11章記載.
這個揀選人去服善.
是神的旨意.
而在出埃及記.
即是摩西的岳父葉蒂羅提議的.
即是兩個記載不同.
究竟有甚麼意思呢.
如果兩個記載不同.
我會體會就是.
上帝會藉著不同的人.
去實行上帝的旨意.
無論是出於某人的旨意.
都是神的靈的感動.
同樣我們今天.
我們有時回應上帝的話.
有些人說你很適合回應.
但亦有些人心裡面.

$^{481}$神的感動說我適合回應.
都可以的.
第二在出埃及記和民數記.
都說摩西揀選人去做.
但在申命記.
他告訴我們是人民選出來.
摩西最後同意的.
在第十四節裡面.
換句話來說.
無論是很民主的申命記.
又或者是之前的兩個記述.
是神的揀選都好.
其實都出於神的旨意.
不過出於神的旨意.
是等於他無憂的.
譬如說當神揀選素羅的時候.
聖經告訴我們.
最終素羅都可以失落的.
第三出埃及記.
是以信心和行為作為標準的.
裡面說要敬畏上帝.
誠實無妄.
恨不義之才的人.
而在申命記的理念.
是很理性的.
因為這裡翻譯了.
和合本的翻譯是.
他有知識的.
而不是為人所熟悉的.
在民數記.
真的告訴我們.
這是上帝的默示.
換言之.
無論是他的行為準則.
又或者是從上帝的默示.
又或者是他自己擁有的知識.
這三樣東西.
都構成了人的知識和神的默示.
很大可能都彼此相輔相成.
其實為何在不同的技術當中.

$^{521}$有不同的解決方法呢.
其實我就會覺得.
好像耶穌基督的福音一樣.
你有馬太,馬可,路加.
三個不同的角度去講同一件事.
因為上帝的國.
不是一個人能夠成就的.
上帝的國是需要我們很多人參與的.
各位弟兄姊妹.
在我們今天面對未來.
有更多需要的時刻.
教會繼續衰落的時刻的時候.
我們可以怎樣做呢.
怎樣共拓未來.
共拓神的國度呢.
神的國度不是只有一個領袖.
不是只有教務的.
我們無論喜不喜歡.
管理,治理其實都需要人.
管理,治理做得好.
其實不等於成功.
相反的.
管理,治理做得不好.
一定失敗.
以色列人要進入所應許的地方.
他需要管理.
我們要突破今天教會的衰敗的限制.
其實我們很需要弟兄姊妹.
整個肢體都同心合意.
在未來的日子當中.
在這裡我有一個插曲.
北區叫做都會區.
在未來的日子當中.
信義會在新田有一個大的屋苑計劃.
叫做新田部屋.
在我們附近將會有不同的新居民進來.
換言之.
錦繡堂有很多未來要得的應許之地.
所以各位弟兄姊妹.
共拓未來.

$^{561}$不只是靠我們的執事,我們的墓者.
而是我們弟兄姊妹一起.
在這裡特別聖經說到.
如果有難斷的案子可以呈到我這裡來.
我就判斷.
換句話來說.
這些也是作為領袖承擔的一個重要條件.
我們要負上責任.
今天我們就好像摩西.
在摩亞地和以色列民去分享.
最後可以怎樣得到應許的地.
我們看到有些應許.
不過似乎我們是選兵結帳的軍旅.
面對很多屬靈的征戰.
我們怎樣可以操練我們的信心.
使上帝的應許在我們中間成就.
未來香港有很多人.
未必懂得聽廣東話.
所以我們也要學習.
要常講普通話.
使我們真的可以跟他們有溝通.
幫助他們成為香港人.
幫助他們成為主的子民.
未來的日子.
是我們需要共託未來的.
沒有了我們.
未來好像現在這樣繼續的衰敗.
各位弟兄姊妹.
香港的教會需要有一個新的思維.
是國度的思維.
國度的意思不是信義會的國度.
不是浸信會的國度.
不是宣道會的國度.
信義會,浸信會,宣道會.
都是神國度的一員.
當我們能夠彼此協作的時候.
福音的廣傳.
教會就能夠在香港當中.
有一個茁壯的成長.
我們上司最終的.

$^{601}$就是讓我們依靠聖靈在今天提醒我們.
以使為鑑.
不要走冤枉路.
承擔使命.
香港每一個人.
其實都應該得福音的好處.
共託未來.
讓我們依靠聖靈的參與在當中.
成就主的旨意.
成為主合用的器皿.
榮歸聖父,聖子,聖靈.
昔在,今在,永遠常在的上帝.
我們低頭祈禱.
慈悲人愛天父上帝.
因為你知道我們的不足.
你又看到我們在這個時刻.
在香港的困難.
教會的不濟.
求主真的幫助我們興起.
我們中間的弟兄姊妹.
讓我們參與在其中.
讓我們共託主你自己的國度.
使到神你自己的國度.
在香港當中能夠茁壯成長.
因主幫助我們.
讓我們有眼光去看未來.
不是看我們過去輝煌的成就.
而是看現在我們的衰敗.
而是看主你對我們未來的應許.
求因主加添我們力量.
奉耶穌基督的聖名.
阿們.
\newpage



\section{創世記 37:1-36}
\label{sec:y77Sfudracc}
\textbf{2024-04-20 晚堂《主必追討》創世記 37\_1-36 講員:陳麗蟬牧師}
\newline
\newline
連結: \href{https://youtube.com/watch?v=y77Sfudracc}{\texttt{https://youtube.com/watch?v=y77Sfudracc}} ~~~~ 語音日期: 2024-04-23
\newline
\newline
\hyperref[sec:iqMY9wa3DY0]{\small{< < < PREV SERMON < < <}}
~
\hyperref[sec:index_chronic]{\small{[返順時目]}}
~
\hyperref[sec:index_scriptual]{\small{[返順卷目]}}
~
\hyperref[sec:kK4Sp5yFq94]{\small{> > > NEXT SERMON > > >}}
\newline
\newline
創世記 37:1-36
\newline
\begin{longtable}{cl}
\hline
\hline
章節 & 經文 (和合本修訂版)\\
\hline
37:1 & \begin{tabularx}{0.7\textwidth}{X} 雅各住在迦南地，就是他父親寄居的地。 \end{tabularx} \\ \\ \relax
37:2 & \begin{tabularx}{0.7\textwidth}{X} 這是雅各的事蹟。約瑟十七歲與他哥哥們一同牧羊。他是個少年，與他父親的妾辟拉和悉帕的兒子們常在一起。約瑟把他們的惡行報給父親。 \end{tabularx} \\ \\ \relax
37:3 & \begin{tabularx}{0.7\textwidth}{X} 以色列愛約瑟過於其他的兒子，因為約瑟是他年老生的；他給約瑟做了一件長袍。 \end{tabularx} \\ \\ \relax
37:4 & \begin{tabularx}{0.7\textwidth}{X} 哥哥們見父親愛約瑟過於他們，就恨約瑟，不與他說友善的話。 \end{tabularx} \\ \\ \relax
37:5 & \begin{tabularx}{0.7\textwidth}{X} 約瑟做了一個夢，告訴他哥哥們，他們就更加恨他。 \end{tabularx} \\ \\ \relax
37:6 & \begin{tabularx}{0.7\textwidth}{X} 約瑟對他們說：「請聽我做的這個夢： \end{tabularx} \\ \\ \relax
37:7 & \begin{tabularx}{0.7\textwidth}{X} 看哪，我們在田裡捆禾稼；看哪，我的捆起來站著；看哪，你們的捆圍著我的捆下拜。」 \end{tabularx} \\ \\ \relax
37:8 & \begin{tabularx}{0.7\textwidth}{X} 他的哥哥們對他說：「難道你真的要作我們的王嗎？難道你真的要統治我們嗎？」他們就因他的夢和他的話更加恨他。 \end{tabularx} \\ \\ \relax
37:9 & \begin{tabularx}{0.7\textwidth}{X} 後來他又做了另一個夢，告訴他哥哥們說：「看哪，我又做了一個夢；看哪，太陽、月亮和十一顆星都向我下拜。」 \end{tabularx} \\ \\ \relax
37:10 & \begin{tabularx}{0.7\textwidth}{X} 約瑟告訴他父親和哥哥們，他父親就責備他說：「你做的這是甚麼夢！難道我和你母親、你的兄弟真的要俯伏在地，來向你下拜嗎？」 \end{tabularx} \\ \\ \relax
37:11 & \begin{tabularx}{0.7\textwidth}{X} 他的哥哥們都嫉妒他，他父親卻把這事存在心裡。 \end{tabularx} \\ \\ \relax
37:12 & \begin{tabularx}{0.7\textwidth}{X} 約瑟的哥哥們到示劍去放他們父親的羊。 \end{tabularx} \\ \\ \relax
37:13 & \begin{tabularx}{0.7\textwidth}{X} 以色列對約瑟說：「你哥哥們不是在示劍放羊嗎？來，我派你到他們那裡去。」約瑟對他說：「我在這裡。」 \end{tabularx} \\ \\ \relax
37:14 & \begin{tabularx}{0.7\textwidth}{X} 以色列對他說：「你去看看你哥哥們是否平安，羊群是否平安，再回來告訴我。」於是他派約瑟出希伯崙谷，約瑟就往示劍去了。 \end{tabularx} \\ \\ \relax
37:15 & \begin{tabularx}{0.7\textwidth}{X} 有人遇見他，看哪，他在田野走迷了路。那人問他說：「你找甚麼？」 \end{tabularx} \\ \\ \relax
37:16 & \begin{tabularx}{0.7\textwidth}{X} 他說：「我找我的哥哥們，請告訴我，他們在哪裡放羊。」 \end{tabularx} \\ \\ \relax
37:17 & \begin{tabularx}{0.7\textwidth}{X} 那人說：「他們已經離開這裡走了，我聽見他們說：『我們往多坍去。』」約瑟就去追哥哥們，在多坍找到了他們。 \end{tabularx} \\ \\ \relax
37:18 & \begin{tabularx}{0.7\textwidth}{X} 他們遠遠看見他，趁他還沒有走近他們，就圖謀要殺死他。 \end{tabularx} \\ \\ \relax
37:19 & \begin{tabularx}{0.7\textwidth}{X} 他們彼此說：「看哪！那做夢的來了。 \end{tabularx} \\ \\ \relax
37:20 & \begin{tabularx}{0.7\textwidth}{X} 現在，來吧！我們把他殺了，丟在一個坑裡，就說有惡獸把他吃了。我們且看他的夢將來怎麼樣。」 \end{tabularx} \\ \\ \relax
37:21 & \begin{tabularx}{0.7\textwidth}{X} 呂便聽見了，要救約瑟脫離他們的手，說：「我們不可害他的性命。」 \end{tabularx} \\ \\ \relax
37:22 & \begin{tabularx}{0.7\textwidth}{X} 呂便又對他們說：「不可流他的血，可以把他丟在這曠野的坑裡，不可下手害他。」呂便要救他脫離他們的手，把他還給他父親。 \end{tabularx} \\ \\ \relax
37:23 & \begin{tabularx}{0.7\textwidth}{X} 約瑟到了他哥哥們那裡，他們就剝去他的外衣，就是他身上那件長袍。 \end{tabularx} \\ \\ \relax
37:24 & \begin{tabularx}{0.7\textwidth}{X} 他們抓住他，把他丟在坑裡。那坑是空的，裡頭沒有水。 \end{tabularx} \\ \\ \relax
37:25 & \begin{tabularx}{0.7\textwidth}{X} 他們坐下吃飯，舉目觀看，看哪，有一群以實瑪利人從基列來，用駱駝馱著香料、乳香、沒藥，要帶下埃及去。 \end{tabularx} \\ \\ \relax
37:26 & \begin{tabularx}{0.7\textwidth}{X} 猶大對他的兄弟們說：「我們殺我們的弟弟，遮掩他的血有甚麼好處呢？ \end{tabularx} \\ \\ \relax
37:27 & \begin{tabularx}{0.7\textwidth}{X} 來，我們把他賣給以實瑪利人，不要下手害他，因為他是我們的弟弟，我們的骨肉。」他的兄弟們就聽從了他。 \end{tabularx} \\ \\ \relax
37:28 & \begin{tabularx}{0.7\textwidth}{X} 那時，有些米甸的商人從那裡經過，就把約瑟從坑裡拉上來。他們以二十塊銀子把約瑟賣給以實瑪利人，他們就把約瑟帶到埃及去了。 \end{tabularx} \\ \\ \relax
37:29 & \begin{tabularx}{0.7\textwidth}{X} 呂便回到坑旁，看哪，約瑟不在坑裡，就撕裂自己的衣服， \end{tabularx} \\ \\ \relax
37:30 & \begin{tabularx}{0.7\textwidth}{X} 回到他兄弟們那裡，說：「孩子不在了。我往哪裡去才好呢？」 \end{tabularx} \\ \\ \relax
37:31 & \begin{tabularx}{0.7\textwidth}{X} 於是，他們宰了一隻公山羊，拿了約瑟的那件外衣染上了血， \end{tabularx} \\ \\ \relax
37:32 & \begin{tabularx}{0.7\textwidth}{X} 派人把長袍送到他們的父親那裡，說：「我們發現這個， 請認一認，是不是你兒子的外衣？」 \end{tabularx} \\ \\ \relax
37:33 & \begin{tabularx}{0.7\textwidth}{X} 他認出來，就說：「這是我兒子的外衣，惡獸把他吃了，約瑟一定被撕碎了！」 \end{tabularx} \\ \\ \relax
37:34 & \begin{tabularx}{0.7\textwidth}{X} 雅各就撕裂衣服，腰間圍上麻布，為他兒子哀傷了多日。 \end{tabularx} \\ \\ \relax
37:35 & \begin{tabularx}{0.7\textwidth}{X} 他的兒女都起來安慰他，他卻不肯受安慰，說：「我必哀傷著下陰間，到我兒子那裡。」約瑟的父親就為他哀哭。 \end{tabularx} \\ \\ \relax
37:36 & \begin{tabularx}{0.7\textwidth}{X} 米甸人把約瑟賣到埃及，給法老的官員，就是護衛長波提乏。 \end{tabularx} \\ \\
[1ex]
\hline
\hline
\end{longtable}
$^{1}$弟兄姊妹平安.
我們在教會的上半年.
主要是講舊約的人物 特別是創世記.
今天我們就講到這個約瑟.
約瑟其實佔創世記的篇幅都很長的.
從37章 即是今天我們講的那裡.
一直到50章.
大部分都是在講約瑟的事跡.
不過很奇怪.
在37章的第一節 他是這樣說的.
「亞國住在迦南地.
就是他父親寄居的地.
這是亞國的事跡」.
在第一節他就在講亞國.
他就告訴我們 亞國住在迦南地.
而迦南地是誰之前住的呢.
就是他的爸爸爾薩.
以前住在那裡的.
由第二節開始他就說.
「下面那些就是亞國的事跡」.
不過接著他就開始講約瑟17歲的時候.
一直發生的事情.
為什麼好像有點混亂呢.
原來這個作者這樣鋪排.
就很想講每一個人.
他自己生命裡面的事情.
一定會和上一代有關.
也都會影響著下一代的.
他主要開始講約瑟的.
約瑟的爸爸就是亞國.
第二節他說「這是亞國的事跡」.
「事跡」這個字也可以指到是歷史.
就是說他在講約瑟的事.
其實是在講亞國這個家族的歷史.
而亞國我們知道他有12個兒子.
在這裡就一直在講他們12個兒子發生的事情.
不過主要都是約瑟.
而這12個兒子.
也都牽涉到他們的祖先亞伯拉罕.
因為上帝曾經應許過.

$^{41}$這個亞伯拉罕的後裔.
是好像天上高的星星.
或者海面裡面的那些沙那麼多的.
然後他就開始講這個約瑟的事情.
我們來看一下這個.
當時是一個少年青年的約瑟.
17歲的時候.
這個約瑟是怎麼樣的.
這個約瑟他的性情.
都有一些驕傲 有一些自負.
從經文第二節裡面講到.
他就說這個約瑟.
以色列愛約瑟過於其他的兒子.
就講到這個約瑟的爸爸亞國.
他又叫做以色列.
他就很疼愛這個約瑟.
帶來的結果是怎麼樣呢.
就是這個哥哥見到爸爸.
是那麼偏心 那麼疼愛這個約瑟.
他們心裡面就很不喜歡.
為什麼這個以色列或者這個亞國.
那麼疼愛這個約瑟呢.
我們知道亞國有12個兒子.
有兩個太太.
在很不情願的情況下.
他就先娶了一個姐姐叫做莉亞.
然後就生了10個.
之後他就娶了他自己很喜歡的那個拉傑.
他很喜歡很疼愛這個女孩子.
到他娶了拉傑之後.
一段時間他才生了這個約瑟.
他就特別疼愛這個兒子.
也都因為他生這個約瑟的時候.
他已經是年紀大了.
那麼亞國疼愛這個約瑟.
他疼愛到一個什麼地步呢.
他就為他這個心肝寶貝.
就做了一件的衣服.
叫做長袍.
聖經講他是一件的長袍.

$^{81}$這件長袍是他特別為這個約瑟而做的.
我們就略略看看這幅圖.
你看到那個穿著長袍.
又多漂亮的衣服就是約瑟來的.
其他那些都是哥哥.
全部都穿著一些沒有袖子的衣服.
也都看上去就知道是一些工作服來的.
就是說他爸爸偏心到一個情況就是怎樣呢.
就是這個約瑟是不用工作的.
穿著這件長袍.
又一件長袍.
又一件這麼漂亮的衣服的時候.
他是不用工作的.
他就享受世界.
不過也不是享受世界的.
他要管理家裡所有的一些大小的事情.
他差不多是一個管家來的.
所以這件長袍就代表了約瑟在家裡的地位.
記住他在12個兒子裡面.
他是排第11的.
但是他卻是一個領袖的地位.
也都凸顯了爸爸非常非常的疼他的.
也都相信約瑟有了這件這麼漂亮的衣服.
這件長袍的時候.
他天天在家裡 在外面周圍走來走去的時候.
他的哥哥一路看的時候.
可能就看到眼火爆了.
我們做得這麼辛苦.
你又不用工作.
穿著衣服周圍走走.
第四節告訴我們了.
約瑟的哥哥.
看到爸爸這麼疼他.
開始就很憎恨約瑟.
就不跟他說一些和睦的話.
每次說這些話.
大家都是擇來擇去.
你激我 我激你 你逮我.
大家是很和諧的.
爸爸對父母 對子女的偏愛.

$^{121}$一定會帶來整個家庭的不和諧.
在我們這一輩.
很多都是很重男輕女.
我在港台上已經講了很多次了.
我們是潮州人.
這個重男輕女.
帶給我們家庭其實很多的事情.
在那個年代的時候.
整個香港經濟都很貧困的.
所以我們一家八口.
就住在一個房間裡面.
分三層床的.
上層是我和我兩個姐姐住的.
然後就有一個碌架床.
中間那層就是我哥哥和我大弟弟住的.
最低那層就是我爸爸,媽媽.
和我最小的弟弟在那裡住.
當大家一直長大的時候.
我兩個姐姐和我哥哥.
都開始讀中學讀書了.
那時候媽媽就去收拾一下房間.
就收拾空了這個地方.
她就想 現在姐姐和哥哥都開始讀中學了.
就想一下 買一個書架.
讓他們多放一些書.
還是買一張小桌子給哥哥呢.
那你猜最後會怎樣呢.
哥哥說 我都讀中學了.
我想專心一點讀書 媽媽這樣.
結果就買了一張小書桌子.
就給哥哥讀書了.
這張小書桌子都成為了我們家裡.
這六兄弟姐妹.
一些的「硬淵」的地方.
當我兩個姐姐.
看著她要專心讀書.
我們兩個就夠要了.
當我姐姐有時候想一下.
那些中學的書多了 厚了.
都不知道放在哪裡的時候.

$^{161}$她就覺得 為什麼他們連一個小書架.
每人有一層 可以放書都不行呢.
我想在十年八年之前.
有時候我們一想起那個書架.
我們的兄弟姐妹都有一些怒氣.
特別是我哥哥 因為都叫長子的孫.
雖然他不是全家最大 他排第三而已.
但他都以一個長子的孫的地位.
有時候他那些說話 或者所做的事情.
他就要全對的了.
但是給我們家裡 都很多不和諧.
就算到現在我們個個都長大了.
想通了很多了.
但有些在當中 有些的家人跟他一起.
都不是很合得來 都仍然存在的.
這個亞國 他那麼偏愛這個弱勢.
就都帶來了家庭很多的不和諧.
也都造成了這幫兄弟想殺這個弱勢.
其實除了這個爸爸的偏愛之外.
還有這個弱勢自己的性格的自負.
他都要負一些責任的.
我相信弱勢他是比較直率的人.
當他跟那些哥哥一起去放羊的時候.
他最喜歡什麼呢? 就是打小報告.
總之他那些哥哥在放羊的時候.
有些什麼事情做得不對的.
他一定會去告訴爸爸的.
我相信是不是全部都是他哥哥不對.
而他又是對的呢? 我又相信不是的.
不過總之這個哥哥有什麼不對.
他就拼命一定會告訴爸爸的.
他就好像做一個監督人那樣.
另外 他又發了兩個夢.
在夢裡面 他就看到.
好像有12個和菌 他自己其中一個.
其他的11個都要向他下拜.
然後他又發了另外一個夢.
就是在這12個和菌裡面.
再加上這個太陽和月亮.
這11個和菌和太陽和月亮.

$^{201}$就一起向著這個弱勢下拜.
這是兩個很特別的夢.
我們看這段經文的特色.
我們發覺到 裡面形容到 描述到.
弱勢是在說話的.
一方面向爸爸打小報告.
告哥哥做錯了什麼.
然後他講這個夢 他就不斷在講話.
在第六節他說.
喂 你們要聽我說話.
在中文裡面說 和合本的修訂本說.
請聽我說 我所做的夢.
但原來其實正式在希伯來文裡面.
它是一個很命令的說話.
喂 你們現在要聽我說話.
聽我說話 我發了一個很特別的夢.
我們看看這段的.
在第五節到第八,第九節裡面.
他說 就是這個.
弱勢他跟他說.
你現在要聽我說這個夢.
他說 看啊 看啊.
我看到在這個田裡面有這個和菌.
你看啊 你們看啊.
就是你聽啊.
在我的和菌裡面.
在我其中的和菌.
還有你看啊 你們其他的十一個和菌.
都向我下拜.
到了第九節 是另外一個夢了.
有了這個太陽和月亮了.
他說 你們看啊 你們聽啊.
我又發了一個夢.
你看啊 太陽,月亮.
還有你們這十一個和菌.
都向我下拜.
你看看這個看啊這個字.
出現了多少次啊?.
一,二,三,四,五 五次.
當有人跟你說話的時候.

$^{241}$他突然間這樣跟你說.
喂 喂 喂 看看外面吧.
那你會怎麼樣?.
你會覺得外面真的有東西看的.
是不是啊?.
那你的眼睛就很容易集中去那裡看了.
當你說 看啊 看啊.
會帶著一個命令式.
就是說 若是他要去哥哥.
甚至之後去爸爸.
聽他說那個夢的時候.
他很興奮.
他有些自負 有些驕傲.
說我發了一個這樣的夢.
我多聰明啊 你們每個人都要向我下拜.
哇 帶著一個沙塵的味道在當中.
有人說 這些沙塵.
這些驕傲 就好像口臭一樣.
自己說的人 自己沒有感覺.
自己不知道的.
不過在周圍的人聞到的氣味.
或者聽到這麼沙塵的說話.
是很難受的.
這個弱勢就是.
不知道自己這些沙塵.
得罪了哥哥 得罪了爸爸.
這些說話 當我們說話的時候.
有時候我們真的要去小心 要去留意一下.
在一個農莊裡面.
有一隻公雞.
牠每天早上起來的時候.
就叫醒人的.
有一天 當在農莊裡面走來走去的時候.
牠突然間看到.
原來有隻狐狸躲在那裡.
牠就很害怕了 就立即走走走走.
那隻狐狸就很溫柔的跟那隻公雞說.
親愛的朋友 你為什麼這麼害怕呢?.
我今天特地這樣偷進來.
我是很想你聽一下.

$^{281}$你在那裡咕咕咕的 那些很清脆的聲音.
我真的很想聽.
那隻公雞就說 有人這麼欣賞牠.
於是牠就站在那裡.
深呼吸 開始想咕咕咕的.
當牠還沒有咕到出口的時候.
這隻狐狸一撲向前就抱緊牠.
就走啊走啊走啊 就咬緊牠.
那隻公雞當然是大聲叫 救命啊 救命啊.
於是全個農莊的豬啊狗啊鴨啊.
都是大聲的叫.
吵得農莊的主人都知道發生什麼事.
所以這個農莊的主人就出來追啊追啊追啊.
就想追回這隻公雞回來 救回這隻公雞.
在那隻狐狸一直咬緊這隻公雞的時候.
這隻公雞靈機一觸 咦?.
怎麼可以走呢?.
牠又跟那隻狐狸說了.
那隻狐狸大哥 哇 你跑得好快啊.
你看後面那些動物和農莊主人跑得那麼慢.
你不用怕 牠們一定追不上來的.
你跑得好漂亮啊.
不如你就停下來.
看看那些豬狗鴨啊那些.
還有那個農莊主人跑得多笨的樣子.
不如你停下來看吧.
那隻狐狸又以為自己真的跑得很快.
又以為自己跑得很漂亮.
於是牠又停下來.
看看後面 笑笑牠們.
當牠停下來 一笑牠們的時候.
牠的口一鬆了 那隻雞就鬆了出來.
就飛飛了很遠.
農莊主人就追上來.
就拿著棍子一直打那隻狐狸.
結果就抓了牠進去農了.
這個故事教訓我們甚麼呢?.
教訓我們 我們說話的時候要小心.
要想清想楚.
要沉默的時候.

$^{321}$就千萬不要張開張嘴去說話.
否則就會大禍臨頭.
可能連生命都沒有的.
這個弱勢.
他就沒有錯 哥哥做錯了事.
你就要想想怎樣去說.
來告訴爸爸.
而不是一句一句地去尖著他 釘著他.
他發了兩個夢.
他很開心 很興奮.
但怎樣去說的時候呢.
會令人的感覺會良好一些.
而不會令到人感覺到討厭呢?.
在《箴言》裡面說.
「多言多語 難免有過失.
節制嘴唇是有智慧」.
所以我們說話的時候.
求主給我們有智慧.
一些事情我們真的要說的.
但怎樣說 能夠令人更加舒服.
怎樣去說的時候呢.
是被人吸收到這個內容.
在《箴言》十八章 它也有說.
「生死在舌頭的掌握之下.
豈愛弄舌的必吃他所結的果實」.
生死在舌頭的掌管之下.
有時真是當我們說錯話的時候.
這個大禍就會臨頭的了.
另外 我們也知道.
人與人之間 我們在一起相處的裡面.
我們一定會大家有彼此不和.
彼此有得罪的時候.
絕對不會每次都是自己對.
人家錯了.
所以這個弱勢 每次都去打小報告.
他以為他自己什麼都對.
這樣就令到哥哥對他很討厭.
所以我們很清楚知道.
有時我們自己都會有得罪人的時候.
有時我們也都會有被得罪的時候.

$^{361}$但當我們被人去得罪的時候.
我們一定的 心裡面一定會生生不忿.
但保羅教導我們是.
生氣卻不要犯罪.
「含怒不可道日樂」.
因為當人家得罪了我們.
就算人家說了一些話.
令到我們感到很討厭.
感到很生氣.
但是我們千萬不要被這些生氣.
這些嫉妒.
一路一路這樣去拖延下去.
因為當我們一路拖延下去的時候.
這些怒氣 這些嫉妒.
會令人生出罪惡來的.
所以中國人都有句話.
「得饒人處且饒人」.
這個是與人一起相處的道理.
少一些的仇恨 多一些的恩情.
人生會活得更加的靈.
我們又繼續看法律律師.
繼續看一下.
剛才我們說了爸爸的偏心.
又說了弱勢的自己性情裡面的怒氣.
那哥哥又怎樣呢?.
哥哥的嫉妒是很重要的.
在第四節說了.
因為爸爸的偏愛.
哥哥很生氣 弱勢.
甚至跟他說的話都是惡言相向.
你責我 我責你.
到了第五節.
哥哥又說了.
他更加的恨他.
又加深了一重.
到了第八節的時候.
當弱勢說完他那兩個夢的時候.
他又再加多一重下去.
就很生氣 很生氣弱勢.
然後到第十一節.

$^{401}$他就說他的哥哥責到他.
那個恨 那個憤怒.
是一直積聚下去.
他覺得大家都是這個亞國的兒子.
為什麼爸爸這麼偏心.
為什麼這個弟弟可以有一件這麼漂亮的衣服.
為什麼他可以不用做那些粗重勞動的事.
為什麼他在家庭裡面要做大亞國.
心裡面有很多的不平.
很多的不忿.
但是當這些不平.
這些不忿一直存在心裡面的時候.
就日而久之.
那些苦毒就會出來.
那些苦毒一直出來的時候.
就會帶來罪惡.
在詩篇37篇第一第二節這樣說.
他不要為作惡的心懷不平.
很多時候我們看到那些人.
我們不看得的.
覺得他這麼壞.
為什麼他的日子可以這麼順.
也不要責到那些行不義的人.
有時候我們就覺得.
為什麼他這麼壞.
他可以這麼順利.
但是上帝說.
他們必如草快被割下.
猶如綠色的亂草快要枯幹.
所以不用激氣.
不用對這些惡人.
這些壞人感到心懷不平.
我們相信.
相信上帝一定會來對付.
一定會來審判的.
我們又看回亞國整個的家族.
好像覺得整個家族.
整天都在說謊話.
你騙我 我騙你的.
看回這個家族的亞國.

$^{441}$他的假話.
在他少年的時候.
他為了要得到這個長子的名分.
他媽媽就教他.
用一隻小山羊的毛.
包著他的頸 包著他的手.
因為哥哥有很多毛.
就騙了這個長子的名分.
因為當時亞國的爸爸叫以撒.
他的眼睛很濛 差不多瞎了.
所以看不清楚 靠一隻手摸.
摸到亞國的時候 毛很滑.
以為他是哥哥.
就把長子很多的祝福.
都給了亞國的.
後來亞國因為怕哥哥尋仇.
所以就逃難去了舅父的家.
到現在 亞國自己.
他有12個兒子.
但那10個兒子就出了去放羊的時候.
亞國就叫約瑟去探望.
看看那10個哥哥安不安全.
因為去到那個事件的地方.
是那裡 第二兒子西面.
曾經殺了很多人.
就看他們安不安全.
到現在 亞國自己去嘗一下滋味.
就是這10個兒子因為很恨約瑟.
很恨這個弟弟.
所以就想把他殺了.
但後來就沒有殺.
就把他賣了給一些路京的商人.
他們當時怎麼騙亞國呢.
就用山羊的血.
《聖經》很有趣的.
之前亞國就用山羊的毛來騙他爸爸.
到現在他10個兒子就用山羊的血.
就染了在亞國自己做的那件漂亮的衣服.
給約瑟的漂亮的衣服.
就弄到他血淋淋.

$^{481}$就騙他爸爸.
哎 弟弟 第一十一弟弟.
就被這些惡獸咬死了.
從這個事蹟給我們看到.
每一件事情 上帝都會追討的.
你昔日 你走去騙你爸爸.
拿了長子的名分.
然後你又逃難去了舅父的家裡.
你自己已經不能再見你的爸爸媽媽.
去了一個陌生的地方了.
但到現在呢.
他10個兒子對他更加淒慘.
他以為兒子死了 約瑟死了.
以後都不能再見.
在《聖經》裡形容到說.
這個亞國他那個終日的難過.
無法得到安慰.
就算其他的女兒都來安慰他.
他都得不到這個安慰.
我整個家族就是這樣.
我們就看到.
上帝是一個審判的主.
所有的事情 他都會來追討的.
他對這個亞國 你昔日做錯了事.
他會懲罰.
他也會寵回他的兒子來懲罰他.
我們看到上帝是那位公平公義的主.
人犯罪 人欺凌.
一定有後果 一定會有報應的.
上帝是一位無所不知.
無所不能的上帝.
他絕對不會以有罪的為無罪.
所以當我們作為上帝的子民的時候.
我們要傳一個敬畏上帝的生命.
來去存活 來去榮耀上帝.
最後我們就來看看.
我們的上帝是一位.
他的摯愛永遠長存的.
我們繼續看亞國這個家族.
亞國的這個家族有12個兒子.

$^{521}$我們今天就只看他第三個兒子.
他第四個兒子叫做尤大.
我就是沒有了那一邊.
好 不要緊.
他第四個兒子叫做尤大.
他第一個兒子叫做呂便.
但是因為他很淫亂.
他竟然和爸爸的妾侍上床.
所以亞國很生氣.
就不給他這個長子的名分.
這個長子名分在聖經裡面說了.
就給了他生愛的太太所生的藥室.
就是這個藥室.
然後接著亞國的第二個兒子.
就是西冕.
他因為他的妹妹在事件這個地方.
就被人污辱了.
所以他就很生氣.
就想計劃又是騙人.
就殺了很多人.
也都因為這樣.
亞國都很生氣他.
唯有是利美.
利美是上帝特別揀選的.
成為了祭司的一路的延伸下去.
到這個尤大了.
尤大是上帝特別祝福的一個人.
和特別祝福的一個支派.
不過尤大都不是一個好人來的.
講一下尤大.
尤大他因為爸爸這麼偏心.
也都知道爸爸不喜歡他們的子女.
娶外邦的女子.
所以他就特意走了去迦南地.
就娶了一個迦南的太太.
和這個迦南人的太太.
又生了三個兒子.
因為在這個舊約的時期.
這個外邦人是上帝所不約立的.
當尤大生了三個兒子.

$^{561}$他第一個兒子.
他長大之後就幫他娶太太.
又幫他娶了一個外邦人.
這個外邦人就叫做他媽.
娶了一個太太叫做他媽.
但這個大兒子又是很不聽話.
聖經形容他.
他行耶和華眼中看為惡的事.
所以上帝就擊殺了這個大兒子.
在舊約的規條.
大兒子死了之後.
二兒子就要來娶大嫂的.
和照顧他.
二兒子表面很聽話.
就與嫂子同房.
但是他覺得.
這個是哥哥的太太.
幫他生了兒子的時候.
這樣不是吃虧嗎.
他卻和她同房.
卻不肯讓她生兒子.
所以他就.
還拿了他哥哥的產業.
上帝又覺得這樣做是不對的.
於是就將他擊殺了.
眼角看到.
死了 自己的大兒子.
二兒子都被上帝擊殺.
現在只剩下第三個兒子.
如果又是這樣的時候.
那是不是變成沒有後代了.
所以就和這個媳婦Tama說.
她說Tama.
這個小孩很小.
還沒娶到你.
不如你先回娘家吧.
等他長大了之後.
我才叫他娶你吧.
但其實這個猶太.
是想趕走他.

$^{601}$免得他去駝衰街.
第三個兒子都死了.
但是這個Tama.
雖然她愛幫女子.
但是她也覺得.
她嫁進了猶太家裡.
她就是猶大的媳婦.
所以她回娘家之後.
一直都穿著寡婦的衣服.
她很想能夠.
猶太的第三個兒子.
娶了她.
為了這個猶太家族.
來有後代.
但一直等.
這個猶太都沒有找她.
沒有叫她回頭.
所以她就知道.
原來這個猶太又騙她.
她就想辦法.
又騙回她了.
她竟然扮演一個妓女.
就和她發生了關係.
然後就幫她生了這個兒子.
生了這個兒子出來.
後來人們就說.
糟了.
這個Tama是寡婦.
竟然有了孩子.
就告訴這個猶太.
你的媳婦是寡婦.
不知道為什麼現在竟然有了孩子.
猶太當然義正嚴詞說.
任婦來的.
要將她燒死.
但誰知道.
這個奸婦是誰.
是猶太自己.
整個故事很精彩.
不過今天說不出.

$^{641}$整件事有證有據出來.
說這個人就是這個猶太.
所以當這件事揭露出來的時候.
猶太沒有話好說.
不過上帝欣賞她.
勇敢地說了一句話.
她說這個Tama.
這個媳婦.
說得有道理.
做得有道理.
然後她做了什麼呢.
從此猶太不再與她同寢了.
這是猶太真心的悔改.
她的悔改有兩方面.
第一方面.
她沒有守她的諾言.
第三就是她長大了.
她應該要娶回這個媳婦.
讓她延續她的後代.
她沒有守她的諾言.
第二就是猶太自己的太太死了之後.
她就走去找妓女.
在當時的古代社會裡面.
男孩子有責任去照顧這些寡婦.
自己家族裡面的寡婦.
所以猶太做錯了.
Tama也沒有錯.
她想猶太的家族有後一路延續下去.
但她做的手法是不好.
不過上帝欣賞的是猶太.
她肯承認自己的錯.
並且她不再做這個奸淫.
以致竟然猶太和Tama.
都在耶穌的家譜裡面出現了.
在《馬太福音》第一章第二第三節.
一直在說耶穌基督的家譜裡面.
是這樣說的.
亞伯拉罕生以撒.
以撒生雅各.
雅各生猶太和他的兄弟.

$^{681}$這個猶太就是雅各的第四個兒子.
然後猶太生從Tama時.
即是他們兩個的關係.
從Tama裡面生了法勒斯.
跟著跟第三個兒子再生謝拉.
從法勒斯裡面再生希斯倫.
一直下去生大衛.
大衛一直去到耶穌基督.
我們看到在整個耶穌基督的家譜裡面.
其實很多人的生命都很不光彩.
很骯髒很不行的.
但是上帝他所體重的是人做錯了事.
無論你多大的事都好.
你犯了這個奸淫.
你不守誠信.
但當他被揭發.
他自己肯認錯.
肯去改變的時候.
上帝會接答.
上帝會洗肉.
我想其實我們在這裡每一個人.
我們都好像雅各的十二個兒女.
亦都好像猶太這樣.
生命裡面有很多軟弱.
很多的污穢.
很多的不見得光的東西.
我們其實都不配成為上帝的兒女.
不配成為信徒.
不配被耶穌來揀選的.
但是就是耶穌基督的寬宏.
對世人的恃憐.
祂給我們的恩典是多麼的大.
以致祂竟然可以能夠去接受我們.
成為了他聖潔國度裡面的一份子.
所以 靈姐妹.
我們的生命會有軟弱的.
會有不光彩的.
會有失敗.
不要緊.
上帝要求我們的.

$^{721}$就是能夠有這個肯去承認.
肯去改變.
肯去悔改的時候.
上帝會接納.
上帝會給我們有一個精彩的人生.
繼續使用我們的.
求主親自到來的.
鼓勵我們每一位.
我們同心的禱告.
親愛的天父上帝.
我們多謝你.
因為你就是那位滿有慈憐.
滿有恩典的上帝.
我們多謝你.
縱然我們每一個人都是不潔的.
我們都有很多的軟弱.
很多的問題.
很多不見得光的事情.
但多謝主耶穌基督.
你仍然要恕我們.
仍然等我們每一個人的迴轉悔改.
多謝你.
不要嫌棄我們.
求主幫助我們.
求聖靈時刻的光照我們.
以致我們知道做錯了的.
我們立即認錯悔改.
以致我們的生命.
不斷被耶穌基督使用.
我們這一生能夠榮耀上帝的名字.
我們多謝你瑞聽禱告.
是奉靠主基督的聖名.
阿們.
\newpage



\section{創世記 37:2-4}
\label{sec:kK4Sp5yFq94}
\textbf{2024-04-21 午堂《約瑟的華服》創世記 37\_2-4;39\_1-6,19-23;41\_14,38-43 講員:楊明麗姑娘}
\newline
\newline
連結: \href{https://youtube.com/watch?v=kK4Sp5yFq94}{\texttt{https://youtube.com/watch?v=kK4Sp5yFq94}} ~~~~ 語音日期: 2024-04-23
\newline
\newline
\hyperref[sec:y77Sfudracc]{\small{< < < PREV SERMON < < <}}
~
\hyperref[sec:index_chronic]{\small{[返順時目]}}
~
\hyperref[sec:index_scriptual]{\small{[返順卷目]}}
~
\hyperref[sec:ZVwFpbgglfY]{\small{> > > NEXT SERMON > > >}}
\newline
\newline
創世記 37:2-4
\newline
\begin{longtable}{cl}
\hline
\hline
章節 & 經文 (和合本修訂版)\\
\hline
37:2 & \begin{tabularx}{0.7\textwidth}{X} 這是雅各的事蹟。約瑟十七歲與他哥哥們一同牧羊。他是個少年，與他父親的妾辟拉和悉帕的兒子們常在一起。約瑟把他們的惡行報給父親。 \end{tabularx} \\ \\ \relax
37:3 & \begin{tabularx}{0.7\textwidth}{X} 以色列愛約瑟過於其他的兒子，因為約瑟是他年老生的；他給約瑟做了一件長袍。 \end{tabularx} \\ \\ \relax
37:4 & \begin{tabularx}{0.7\textwidth}{X} 哥哥們見父親愛約瑟過於他們，就恨約瑟，不與他說友善的話。 \end{tabularx} \\ \\
[1ex]
\hline
\hline
\end{longtable}
$^{1}$我們一起低頭祈禱.
參天父 楊帝 我們來到你面前.
我們要多謝你.
多謝你在風雨當中帶領我們每一位.
來到這些樹 我們一起敬拜.
一起來聽你的話語 求你賜福.
讓我們無論聽的 說的.
都蒙你的恩典 可以得著.
你的話語和生命的造就.
祈禱仰望你 奉主耶穌.
基督的聖名 阿們.
今天我們一起.
看藥室的華服.
大家呢.
你喜歡的衣服是怎樣的.
你喜歡的呢.
是豪艷的晚裝.
還是.
運動服.
有一位街坊呢.
我跟他聊天.
他就說.
我家有幾雙的運動服.
嘩.
原來呢.
你喜歡的.
就會有很多.
我聽過姊妹說.
其實呢.
我都未穿過 牌子還在.
不過呢 都要早點.
處理了.
什麼衣服我們都可能是.
要來看多過要來穿的.
對於.
藥室的衣服 我們.
又了解多少呢.
可能我們最熟的呢.
他有他爸爸幫他做的.
一套是彩色的.

$^{41}$衣服 一向呢.
我們信主多幾年的.
原生的和合本呢.
就知道那件叫做彩衣.
現在叫做.
長服 長袍 長衣.
原來呢 它是一件長的.
起碼呢.
都過膝的了 衣服.
是有顏色的.
彩色的 也都長袖的.
這個是.
一件什麼的服裝呢.
這一件的彩衣呢.
是在表達它的地位.
當然啦.
無論是去宴會也好.
上班也好.
或者是運動裝也好.
這件衣服呢.
都在表示你的地位.
這件衣服呢 對他來說.
是很特別的 因為呢.
是爸爸 親愛的爸爸.
做給他的 而且呢.
是一件在當時來說.
是一件行政的衣服.
他十.
個哥沒有的.
他們十二兄弟.
有個大媽.
一個小媽過世了.
就是他的親生媽媽.
還有兩位是妾侍.
所生的兒女呢.
個個都沒有的.
連他弟弟都沒有的 因為太小了.
還沒有開始工作.
他開始工作的時候呢.
他十七歲的時候.

$^{81}$他爸爸就幫他做了一件彩衣.
這件彩衣是很特別的.
很有意義的.
這件彩衣在表達一些東西.
在表達他的地位.
在這麼多哥哥當中.
他爸爸是.
認同他的做法的.
在他的工作能力上面.
他爸爸是嘉許他的.
給他一個主管的職位.
而十個哥哥是沒有.
這件彩衣的.
如果十個哥哥有不同顏色的彩衣.
不過是他那件特別漂亮.
那都容易一點.
去明白和解釋.
但是似乎呢.
十個哥哥都沒有.
這個主管職位的袍.
即是十個哥哥上班呢.
就是穿著.
勞動裝.
而他穿著西裝.
開會的時候.
就是他說了算.
這些職位呢.
我們叫他做什麼呢.
他是主管.
他是副總經理.
他是未來的代表.
在上市公司.
在上市公司來說.
他是將來的主席.
那他主席是誰呢.
就是他爸爸.
十個哥哥看到的.
十個哥哥看到的時候.
感覺是怎樣.
就是說你寫了遺產.

$^{121}$是嗎.
你寫了遺產會怎麼做呢.
雖然你沒有像爺爺那樣.
祝福你.
不過.
你將你的祝福寫了出來.
將衣服.
做一件.
就給他.
而這個人是怎樣呢.
這個人是真的.
好像來看著我們工作.
有什麼時候我們有什麼.
做得不太合你心意的時候.
他一定會告訴你.
我們做得好的時候.
他好像沒有回報.
以主管來說他一樣做得失職.
在哥哥的眼中.
他會是一個.
無情百思.
空降下來的主管.
而且這個.
空降下來的主管.
還要在家裡面.
差不多年紀是他最小.
除了他最小的弟弟.
父親對他.
一向溺愛.
從小到大.
他都不用出來做看養.
照顧養.
和獅子.
熊搏鬥的.
他現在來.
就是做主管.
他來到的時候.
我們就要很小心.
如果不是呢.
就會有人.

$^{161}$戳我們背脊.
他是未來的CEO.
不過是現在的二五仔.
是這樣的時候.
爸爸這件的彩衣.
戴在藥室.
是地位.
但是是一個更加不好的.
兄弟的關係.
其實.
他之前的十個.
比他大的十個哥哥.
他爸爸都有教過他.
他爸爸.
在幫外父做事的時候.
都教他們.
怎樣.
騙外父的牛羊.
那些羊.
是.
要到要.
培育的時候.
肥的就讓牠們生順色的.
瘦的呢.
就讓牠們生雜色的.
那些兒子知道的.
我想最熟的最大的四個.
因為最幫得.
啞角多年.
知道了怎樣.
在行裡面.
出古惑.
你現在爸爸.
聽弟弟說.
說我們出古惑.
他們的心裡面是怎樣想的呢.
這件彩衣.
戴在藥室的.
是一個我們想像不到的.
家庭關係.

$^{201}$或者我們.
對一些新聞.
上面對於某一些商業.
機構的.
多妻的主管.
的家務事之多些的時候.
可以想像到.
藥室的家庭是怎樣的.
但是藥室很喜歡這件衣服.
那是的.
誰不喜歡高積呢.
還要是父親的肯定.
那麼多兄弟.
父親只肯定自己.
當然很開心.
這件是他很開心很喜歡的衣服.
可惜他不是穿得很久.
他還有很多件衣服.
他第二件衣服.
就是中間那件.
他第二件衣服.
應該沒有那麼乾淨.
被人出賣的時候.
扔進坑裡面再拉上來的時候.
應該是髒髒的.
甚至衣服有些穿爛.
後來他就去到.
波提佛的家裡面.
做奴隸.
不過做奴隸也好.
他也是晉升得快的奴隸.
比其他原生的奴隸.
他晉升得快.
做了奴隸的主管.
不過.
身份也是一樣 他是奴隸.
可以被出賣.
這個身份是不會變的.
不過待遇不同.
甚至乎他應該會有些奴隸來服侍他.

$^{241}$他也不一定.
不過他仍然是一個奴隸.
後來就因為.
主母出賣他.
主母是很想勾引他.
挑逗他.
但是他就.
很正義.
也不懂得脆弱.
於是他就沒有犯罪.
沒有得罪神.
不過得罪了主母.
主母就誣衊他.
說是他要來逼主母.
當波提佛聽到的時候.
他應該怎樣呢.
佛的太太.
自己很沒面子.
太太想在外面做些事.
想做外面那棵杏樹的花.
波提佛護衛長的面子放在哪裡呢.
無論波提佛信不信.
一定要罰藥室.
於是乎.
是沒有經過任何審訊.
沒有調查.
藥室就收了監.
為什麼呢.
收監是安全的.
收監太太就不會飛了出去.
於是乎藥室就.
入了監.
由奴隸的身份改變了.
他有另一件衣服.
就是監等.
奴隸的衣服.
已經沒有什麼衣服.
奴隸的衣服.
應該只有下半身.
他做了奴隸的主管的時候.

$^{281}$應該還有一件全新的衣服.
但是進到監獄的時候.
未必有衣服.
不過圖畫通常.
要給大家看.
是畫得多了.
這是他另一件的.
畫服.
他在監獄當中.
坐監休息.
好像不是.
他不像在父親家.
他家什麼時候.
都有事情做.
坐監的時候.
都有事情給他做.
坐監有事情做也正常.
不過坐監還有什麼事情做呢.
坐監也給他做了主管.
其實我們也可能想跟他說.
弱勢 你坐監了.
乖乖地坐監就算了.
有事情做就做.
你為什麼那麼多事呢.
做到被人欣賞你.
然後給你做主管.
做監獄的主管.
不過你不是懲教官.
仍然是監等.
兒子他就是弱勢.
他去到什麼情況下.
他都很積極努力地工作.
其實去到監獄.
撇下.
埋怨下.
向神申訴下.
很正常的.
是不是.
耶和華神啊.
我做錯了什麼.

$^{321}$我有一件彩袍.
那件彩袍沒有得罪誰.
難道得罪了你嗎.
耶和華.
有時候我們會這樣跟神說.
做什麼事會發生.
這些事在我身上.
我都沒有做過什麼事.
我又敬虔.
又守你的律例.
發生什麼事.
這些事發生在我們身上.
甚至乎是.
不好的事發生一次就算了.
為什麼一而再再而三.
這樣來呢.
一而再再而三不止.
一而再.
再而三之後.
好像無了期.
無了期的.
主母.
誣衊他.
之後他被.
關在監獄裡.
沒有刑期的.
沒有任何審訊.
沒有任何審訊的時候.
不會有刑期的.
若生一生要在監獄裡面過.
倘若我們身上有一根刺.
好像保羅說的有一根刺.
這根刺要與他共存.
但是呢.
若生這根刺是在監獄裡面.
我們會怎樣.
如果有一個病.
是不好受的病.
我們要與他同存.
我們都會覺得很辛苦.

$^{361}$如果我們怎樣做.
都不會晉升.
我們都會覺得失望.
若果無論做什麼事.
做得多好.
多用心.
得來都是負面的批評.
我們都會難過.
他是這樣的.
他做得多好.
都是在監獄裡面的.
你是主管.
奴隸都不是.
不過你出了監獄就是奴隸.
這對弱失是一回事.
他怎樣感受呢.
他在這裡應該很長日子.
由他收到這件彩袍.
那件彩袍都未到霉爛.
應該不是穿得很久.
接著就被出賣.
然後去了波提佛家.
後來去了波提佛那裡的監獄.
那裡是屬於天奴.
專門來找高官的.
即是在那裡.
你還要做惡事服侍人.
你得罪誰.
他有刑期的.
刑滿之後你就知道未到.
這是弱失的盼望.
每放月下.
這就是他的未來.
他應該像是要這樣.
在這種情況下.
他怎樣過呢.
應該埋怨.
應該甚至會想得更加負面.
我自己了結自己.
或者是.

$^{401}$我在監獄裡裡製造一堆麻煩.
但他沒有.
他沒有.
他對著這一堆的衣服.
帶來給他一個又一個的問題的時候.
他仍然很努力.
很努力在這裡工作.
小時候.
他的認真.
有父親的認同.
不過帶來兄弟對他的嫉妒.
他在監獄裡面.
做得與眾不同.
老實說監獄裡面的人.
不會因此而給他一個獎賞.
監獄裡面的人.
如果他得罪了.
他又不是很好.
但他仍然很積極.
很主動地工作.
甚至乎他見到.
這個施酒長和施膳長.
他們不開心的時候.
他會問候他.
你做關顧.
用不用做到這個地步.
你不是去監獄報道的教務.
你坐在監獄裡面.
他關心他們.
為甚麼你們兩個不開心.
我們會不會這樣.
為甚麼他會這樣.
他仍然在監獄裡面.
做一件事.
就是發揮他的恩賜.
領導.
在最不開心的時候.
他仍然在監獄裡面.
做一件事.
就是發揮他的恩賜.

$^{441}$領導.
在最不開心的時候.
最沒有辦法去解決問題的時候.
似乎問題要永遠存在的時候.
他繼續發揮他的恩賜.
去領導.
可能他很年輕.
十多歲的時候.
他已經明白了夢境.
於是乎他時時都幫助到人去解夢.
亦都出現了一些情況.
他很看到神在他身邊的工作.
他相信神給他的生命.
他要認真去生活.
即使與監獄一起共存.
他的生命.
就算是監獄.
都榮耀神.
神又真的在監獄裡面.
都賜給他榮耀.
不過他始終是監獄裡面.
他的工作很順利.
亦都他一定是很有品德.
所以能夠在一個天牢裡面.
那麼多不同的高官.
第一級的監獄裡面.
那麼多不同的高官.
他仍然工作順利.
其實他十七歲之後被埋.
去到他最後被釋放出來的時候.
三十歲.
當中這十幾年.
是怎樣過的呢.
我相信在十幾年的負面環境裡面.
是很難生活的.
但是他沒有自暴自棄.
他仍然是一個很可被信任的人.
亦都發揮他的恩賜.
他的智慧.
繼續過他每一日的日子.

$^{481}$他不能盼望.
有得釋放出來.
可能波提佛都相信他無辜.
不過為了家裡安全.
不可以讓他回家.
如果不是呢.
其實哪一個坐牢的人.
能夠得到看顧監獄的人的信任.
派給他做主管呢.
一會兒他自己不知道走到哪裡去.
他是很被信任.
其實是誰栽培了他呢.
是誰栽培了他呢.
那時候沒有兒童主義學.
沒有天崇.
17歲而已.
18歲他可能來到和我們一起崇拜.
他是怎樣.
是誰培養他的信仰呢.
其實第二節.
是被我刻意扣起了一句.
沒有給于唐讀的.
因為我想現在才講這一句.
這一句就是.
「這是雅各的事蹟」.
其實整個36章就講.
雅各的哥哥耳掃的事.
他整個的家族有多少代.
有多少人.
有甚麼的支派出來了.
然後一開37章就講.
雅各的後裔.
雅各的事.
一講雅各的事就講弱勢.
其實一邊講弱勢的時候.
其實是在講雅各的事.
第38章整章是講猶大的事.
39章再開始講弱勢.
講到40多章.
主要都是講弱勢.

$^{521}$那雅各的事是甚麼呢.
可能雅各12個兒子當中.
最能夠寫照到雅各的信心的是弱勢.
當時是誰栽培弱勢的呢.
其實以撒在弱勢被賣的時候.
或者當他死了的時候.
以撒還在生存的.
不過很老了.
但是弱勢應該不是由以撒得到教導.
因為他很長時間.
他出生的時候.
都是在哈蘭的地方.
是後來跟爸爸回來的.
所謂應祖歸宗.
那他的信仰是從哪裡來的呢.
不是他媽媽.
他媽媽是一個多神的信仰者.
他的信仰應該由雅各那裡教導.
弟兄姊妹.
有時候教兒女.
或者教兒童主義學.
我們不知道他吸收了多少.
好像看不到.
但是他吸收了.
你的努力是沒有白費的.
你跟他說的.
帶給他的.
是沒有白費的.
他吸收到.
甚至乎他再記得這些話的時候.
我們都不會知道.
小朋友到他長大.
他上教會.
他受洗的時候.
我們不一定都是在那裡.
因為他可能去了外國讀書.
在外國受洗也不一定.
我們有一位弟兄.
應該是以前約書亞團的團友.
我們在座有沒有約書亞團的在這裡.

$^{561}$前約書亞團的團友.
他去了外國讀書.
應該說他是做交換生.
去了外國讀書一段時間.
就在那裡.
他在香港有回來感受.
他去了外國.
他也有回教會.
原因是在外國.
他覺得很得到牧養.
於是就在外國受洗.
回來之後他做什麼.
他不斷求父母要回教會.
不斷求他爸爸媽媽要回教會.
弟兄姊妹.
你教主教學是很有用的.
你帶青年團少年團是很有用的.
是教會的未來的盼望.
終於他爸爸媽媽上初訓班.
他也結婚了.
你想想是一段很長的路.
不過你不只是帶著青年人.
他會信主.
原來影響他父母.
這個青年弱失.
是他父親的教導影響了他.
很積極很認真去信神.
他信得很真.
可能弱失是從他爸爸口中.
聽到很多他爸爸的經歷.
是聽得比他的哥哥多.
他爸爸沒叫他出去工作.
對他的教導是正面一點的.
對十個哥哥是教他怎樣古惑的.
所以他得到多一點.
他得到多一點.
得到正面的教導.
除了很有恩賜去做領導之外.
他在信仰上真的很認真.
無論什麼環境都認真地.

$^{601}$去跟隨這位神.
依靠這位神.
弟兄姊妹.
我們的工作.
無論哪個年紀是我們的對象.
我們都要很積極認真去做.
神在他們生命當中所做的.
是我們不知道的.
而我們自己本身.
有時候我們都會面對苦難.
好像弱失他是一個未來CEO的地位.
原來可以去到這樣的境界.
坐在監倉裡邊.
開頭是等人派飯.
尾巴是幫人派飯.
但無論怎樣.
苦難總要過去.
恩典會長存.
苦難會佔人生多少的年日呢.
如果以弱失來計.
我當他從17歲開始被埋.
因為聖經沒有寫是幾時.
都不知道是幾時.
他被釋放的時候是30歲.
他應該做了13年的奴隸兼監等.
然後他的生命有110年.
他苦難的日子是一成多一點.
他大部分的日子是甚麼呢.
是真的穿著華服.
是他爸爸做給他的華服.
和天父做給他的華服.
雖然他是做奴隸.
雖然他在監獄當中.
人一直尊重他.
這是神給他的華服.
比任何一件美麗的衣服更加重要.
但不單止是這樣.
他為著人家所做的.
亦都帶來給他恩典.
有時我們會覺得.

$^{641}$為何苦難總不過去呢.
其實可以受少兩年苦.
那位是法老的施酒長.
若果幫法老做施酒的時候.
一早記得藥室.
就可以救到藥室.
不用再坐兩年監.
不過藥室不坐監之後.
出來做甚麼呢.
可能都是做奴隸.
因為身份不能跳.
是奴隸就是奴隸.
不會成為自由人.
到法老做夢的時候.
找誰解夢呢.
不知道去哪裡找藥室.
做奴隸的時候.
名字可以被改.
不是自己想不想改.
被出賣去到哪裡.
就去到哪裡.
甚至乎不會在埃及.
可以再將他反賣去不同地方.
於是會發生的一些事.
就是埃及有七年的豐收.
享受得很好.
到七年的旱災的時候.
甚麼都沒有剩下.
埃及人 法老.
都要捱這個饑荒.
其他的國家.
雅各的後裔 以色列.
都要捱這個饑荒.
那兩年.
在施酒長忘記了藥室.
或者施酒長不知道藥室是甚麼水.
我告訴法老.
你出來之後做甚麼事.
你負責任的時候.
我又要再進去坐牢.

$^{681}$又或者他真的忘記.
我們不知道.
但是神仍然沒有忘記藥室.
祂知道的.
這兩年游藥室在困境的時候.
是為了兩年之後.
他可以做埃及一人之下.
萬人之上 宰相的職位.
不只不是他一個人的成就.
我們常說萬事要甚麼.
互相效力.
愛神的人的益處.
過了這兩年.
他為法老解夢的時候.
埃及人有救了.
法老有救了.
他的父家有救了.
不只這樣.
隔離很多不同的國家.
中東的國家有救了.
很多民族有救了.
神為他安排了一件.
法老給他的華服.
做宰相.
讓他能夠成為一個很高的位置.
也能夠以後.
饑荒之後.
他仍然是很高的位置.
也讓以色列國能夠在那裡.
一直孕育成為一個大國.
後來才離開埃及.
進入迦南地.
他父親為他預備了一件彩衣.
但我覺得神為他預備了一件彩衣.
更加榮耀.
無論是在奴隸的身份.
或者監獄的身份.
他都被尊重.
他都被視為主管.
他仍然繼續操練他的恩賜.

$^{721}$他要操練的.
因為一出了監獄.
他就是做宰相的了.
那十幾年是沒有冤枉過的.
是他拿了他的PhD.
於是乎他就穿著神給他的華服.
那件彩衣.
過他的餘生.
直到如今.
我們都記得他.
弟兄姊妹.
我們要將神交託給我們的下一代.
好好去教導.
另一方面.
當我們在困境的時候.
我們記得神.
祂沒有改變.
祂沒有離開.
祂一直與我們同在.
祂照顧著我們.
在困境中.
祂仍然讓我們受尊重.
仍然會帶我們脫離困境.
幫助我們.
\newpage



\section{帖撒羅尼迦前書}
\label{sec:ZVwFpbgglfY}
\textbf{2024-04-25 周四晚查經班《帖撒羅尼迦前書(4)》 講員:曾裔貴牧師}
\newline
\newline
連結: \href{https://youtube.com/watch?v=ZVwFpbgglfY}{\texttt{https://youtube.com/watch?v=ZVwFpbgglfY}} ~~~~ 語音日期: 2024-04-28
\newline
\newline
\hyperref[sec:kK4Sp5yFq94]{\small{< < < PREV SERMON < < <}}
~
\hyperref[sec:index_chronic]{\small{[返順時目]}}
~
\hyperref[sec:index_scriptual]{\small{[返順卷目]}}
~
\hyperref[sec:a_rRhGfuxUw]{\small{> > > NEXT SERMON > > >}}
\newline
\newline
$^{1}$一起同心祈禱.
感恩天父給我們再一次.
我們可以看貼實羅尼加前後書.
表面是一卷比較簡單的書卷.
和每一章都涉及到主耶穌的再回來.
這個所指的再回來.
是明顯說的末世第二次的再來.
也是世界的末日.
也是信徒要進到永恆榮耀的地方.
所以我們都很盼望神給我們重視.
末世論這個課題.
相信和我們今天處於社會.
我們可能都會很少去想.
世界會有一個終結.
但我們很盼望神幫助我們去明白.
去掌握.
我們將貼前的經卷,經文.
都仰望在你的手中.
奉主耶穌基督的名.
阿們.
大家先拿聖經.
好.
最熟那段就是哥林多前書十三章.
大家應該都相當熟的了.
我們可以當作讀一下這些經文裡面的意思.
然後我們就會擴大.
就是保羅書信裡面.
他講到信忘愛的.
我已經列出了那些經文.
有部分的經卷只是講兩樣的.
例如信心,愛心.
很少會只是講一樣.
譬如信心加盼望.
都這樣不多的.
但是如果只是講一樣的當然很多.
差不多每一卷書都有講.
但是因器苛成講幾方面的關係.
就是這幾卷書卷有講.
我們先看哥林多前書十三章.
那個叫做不是因此.

$^{41}$原來是愛心是最妙的道.
因為這裡是由第十二到第十四章的經文.
十二和十四章的確是講恩賜的來源.
哪些人可以有恩賜呢?.
恩賜怎樣去運用才叫做恩賜呢?.
即是有它的局限性.
如果你離開了教會這個場景.
那個就不可以叫做恩賜了.
或者神給你有那些才幹.
或者你自己先天後天有的才幹.
但你如果不是用在神的教會去祝福.
去幫助教會去成長.
那個可能是你個人的一個才華.
就不能夠去對別人有一個貢獻.
這個似乎在《古羅書信》第十二,十四章裡.
它很想帶出這個局限性.
十二章其中更加重要就是.
要釐清恩賜的來源一定是出於神的.
因為有些可能哥林多教會的人.
他未信耶穌之前.
他們拜很多偶像.
他們可能會有很多靈界經驗.
當然是邪靈.
所以保羅就要撇清這一切.
都要搞清楚.
所有的恩賜都是由神.
三一神那裡而來的.
這個是大前提.
所以如果我們在未信主之前.
曾經有過一些靈界經驗,經歷.
那些我們要放在神那裡.
求神來光照.
甚至要悔改.
有些可能是神婆,教鬼,靈魂出竅.
這些要求主去掃清它的.
不是說基督教沒有.
保羅都說有辨別諸靈的恩賜.
有這樣的恩賜.
在初期教會的時候.
甚至在近代都有些特別.

$^{81}$在靈界上有服事的傳道牧者.
但不是每一個.
所以一定要搞清楚.
恩賜的來源一定是從三一神那裡而來.
而且是為著教會.
每一個信徒都一定有神賜給你的恩賜.
一定有的.
所以是多還是少?.
是甚麼類別?.
但更重要的是.
那個目的不是用來炫耀.
是用來服事的.
這是幾大前提.
就是十二章的哥林多前書.
去釐清這些大原則.
接著去到第十二章的結尾.
要留意一下.
它一直都在說那個比較.
即是說哥林多教會的弟兄姊妹.
就是忘記了恩賜是從神那裡賞賜而來的.
當然這個賞賜也都可能包含後天我們的努力.
我們的去耕耘,去發展.
但仍然那個大前提都是神賞賜給不同的人.
是要造就那個教會.
所以如果我們要去比較的時候.
就很可惜了.
第十二章的二十八節.
神在教會所設立.
第一就是使徒.
第二是先知.
第三是教師.
其次是那些行義林的.
再次是那些得恩賜.
可以醫病的.
幫助人的.
治理事的.
說方言的.
豈都是使徒嗎?.
豈都是先知嗎?.
豈都是教師嗎?.

$^{121}$豈都是行義林的嗎?.
豈都是得恩賜醫病的嗎?.
豈都是說方言的嗎?.
豈都是翻方言的嗎?.
有沒有留意到?.
首先就是說那個次序.
在教會裡面.
例如有不同的恩賜.
有些是顯著的.
譬如使徒,先知,教師,傳福音.
治理事,教導人,安慰人.
諸如此類的.
一路去數.
當然不是很具體的.
每一個恩賜都數.
只是例如這樣說.
顯著的,重要的.
尤其是哥林多教會.
爭的那些恩賜.
特別這樣有按次序去數.
然後他就強調.
你看到那個修辭.
你看中文都看到的.
豈都是.
即是說.
怎麼可能只有一種恩賜.
我們要爭呢?.
怎麼可能只有.
整個教會裡面.
只有一種的恩賜.
一類的人.
一類的職份.
我們要去學武.
我們要爭兵頭呢?.
大概的意思就是這樣.
那樣就.
本來就說.
其實不是的.
是有不同的恩賜.
每個人的領受都不同.

$^{161}$聖靈會按照著他的意思.
來分給不同的弟兄姊妹的.
然後呢.
就是第31節.
突然間改變了那個.
那個寫作的手法.
你們要切切的求更大的恩賜.
切切要求更大的恩賜.
什麼是更大的恩賜呢?.
14章第一節就說.
就是先知講到.
然後下面那裡就說.
我現今把最妙的道指示你們.
所以這個愛不是一個恩賜.
是最妙的道.
所以這裡就告訴我們.
第13章.
我如果能夠講萬人的方言.
即是說當時教會的人.
爭講不同國家的方言.
並天使的話語.
或者有些人可能見過天使.
可能他以為自己是在講天使的話.
有些三重天的經歷.
異象.
卻沒有愛.
就成了明的羅.
響的白一般.
我如果有先知講道之能.
即是能夠講預言.
或者去將舊約的聖經.
深入淺出的去講解.
明白.
不過當時似乎都是預言多.
也明白各樣的奧秘.
各樣的知識.
而且有傳遞的信心.
叫我能夠移山.
耶穌講過.
路加坤第十章.

$^{201}$你吩咐那個山移開.
你只要信心不軟弱.
就可以成就.
保羅在這裡好像講一個移山的信心.
不過留意.
卻沒有愛.
即是十三章保羅講的最妙的道.
就是講愛在教會行使恩賜的.
好像一種黏合劑.
一種帶來合一的.
很重要的那個way.
那個道.
所以一直都是整個第三章都是講愛.
然後再有.
如果我將所有的周濟窮人.
在教會裡面關心人.
又捨己身叫人焚燒.
即是燃燒自己.
叫其他人效法.
即是一種捨身.
一種很拼命去事奉去幫人.
幫到一個地步.
都好像能夠燃點起別人的愛心.
不過卻沒有愛.
這樣很難的.
兩者很微妙.
拼命拼命這樣做.
但原來心裡面是沒有愛的.
這樣就很高難度.
那個沽名釣譽都沒有人看得出.
連自己都騙了.
保羅說有這樣的人.
他如果內心沒有了那份愛.
只有個「做」字.
可能何時會顯露這個底蘊呢?.
可能他燃燒的時候.
當他做很多事.
最終原來想別人得到.
別人對他的稱讚.
但得不到.

$^{241}$甚至有人覺得.
你做就應該焚.
可能這句說話就勒慶他.
那個焚燒就是他內心火滾.
原來背後.
保羅說你做這麼多事.
原來你的動機都不是為了有愛.
是為了愛.
其實就不能夠有益處的.
所以仍然與我無益.
接著說愛的定義.
其實可以這麼說.
一個事奉者.
他就應該有這份.
傳遞心靈的態度.
價值觀.
一投入事奉裡.
他就應該有這樣的內心.
是不斷相輔相成的.
隨著事奉.
因為你的心底有一份愛神愛人.
你就會相輔相成.
不知不覺間.
你其實不需要追求謙卑.
忍耐.
其實就已經在你的事奉裡.
就會成長了.
愛是恆久的忍耐.
愛是又有恩錢.
愛是一個不疾度.
即是神高舉你的好朋友.
不高舉你.
你拍手掌的.
你很開心見到別人.
得到神的恩典.
這個就叫不疾度.
即是今季有三班住一行.
每一班 我不知道現在實際是多少人.
一班有一百人上.
一班只有十個人上.

$^{281}$如果我是開十個人那班的.
我有沒有為到隔壁那班坐滿人的.
過來我們那一行借椅子.
借了椅子去.
然後那些人坐著坐著不對勁.
連那十個都過到隔壁廳.
你會不會不疾度呢?.
雖然好像開玩笑.
但是我們就很懂得保.
可能心底疾度不敢告訴別人.
這些其實是事奉裡面.
其實要預了的.
要真是有這種胸襟的.
所以愛就會令到我們.
可以在那個成長過程.
由可能很疾度.
到慢慢發現一切工作都是神的.
神有恩典在不同的人身上.
豈不是更好嗎?.
就帶來你就會不疾度了.
這個就是一個成長的人.
隨著他的愛心越來越增長.
越來越純全的事奉的動機.
就會出現了這些的恩慈.
很kind的 對人很kind的.
那個心靈是很kind的.
有時對事可以很嚴肅.
對人是很kind的.
有時我都不想這樣說.
我這幾個月.
紅恩那些小朋友.
小學幾個.
很喜歡我.
一來見到我就笑.
衝過來.
其實我很想退休就要申請.
和Oti申請教兒童主義學.
但是我有生之年到現在.
我都未真的做過兒童工作.
其實是全世界最開心的工作.

$^{321}$長者也是.
你可以想像得到.
愛是不自誇,不張狂,不作害羞的事.
不求自己的益處.
愛本身的定義就是這樣.
不輕而發怒.
所以愛不是一種的溺愛.
不是一種的包庇.
包容和包庇.
大家都知道已經是完全不同的態度.
包容是對別人這個人仍然對他有信心.
仍然覺得神的恩典可以改變他.
包庇是自己趕緊的.
將他的罪抱住.
使他不會看到自己的問題.
也不會悔改.
這叫包庇.
所以包容是讓他有空間慢慢成長.
不要計算別人的惡.
不喜歡不義.
只喜歡真理.
凡事都是包容,凡事都是相信.
凡事盼望,凡事忍耐.
後面才是所謂的三角.
保羅的信忘愛就在這裡出現了.
很簡單.
整張這裡都是在說愛.
就是那個最妙的道.
使所有的恩賜能夠黏合.
使所有有恩賜的弟兄姊妹.
能夠在教會裡面一起向一個目標.
就是因為有愛.
所以大家的恩賜發揮的時候.
就會祝福整個教會的弟兄姊妹成長.
就是說在愛裡面的合一.
來彼此服侍.
接著說.
愛是永不止息.
先知講道之能終必歸於無有.
說謊言之能終必停止.

$^{361}$先知講道是現在在教會裡面的恩賜.
在地上需要的恩賜.
所以一個傳道者.
他需要有先知講道的能力.
否則都是那段經文.
睡著了.
因為那些人都沒有答案.
其實不只是答案這麼簡單.
而是他有這種感染人心的能力.
讓聽道的人很想有一個改變.
所以每一次的道都是一個呼召.
都是一個changer.
心靈的changer.
譬如悔改的報道會.
培靈會.
現在進行的恩福的培靈會.
或者是其他的不同的崇拜聚會.
都是希望聽的人有一個changer.
心靈的一個改變.
或多或少的改變.
所以先知講道之能終必歸於無有.
就是說.
他現在說的是現在的短暫和永恆的對比.
因為下面是真的在說這件事.
說方言也是終必會停止.
所以這句話不要小看.
就是方言的終止論的最重要的一節經文.
所以大家如果沒有留意那些靈因運動.
你就不會留意到原來在哥羅德前書13章.
講到關於方言.
到底今天還有沒有方言呢?.
終止論.
終止論就是不會再有.
就是根據這裡.
說方言之能終必停止.
知識即地上現在我們需要有的對神.
對這個世界.
對人的認識.
也必歸於無有.
不要留意.

$^{401}$歸於無有出現了兩次.
停止出現了一次.
然後第九節.
我們現在所知道的有限.
先知所講的也有限.
等那完全的來到.
這有限的必歸於無有了.
一直都在講短暫的.
有限的.
無論我們多厲害的人.
都是只能夠盡我們的能力.
去無論是服侍那代人.
或者是將我們僅有的知識.
神給我們就這麼多.
我們就盡力去做好.
去讓弟兄姊妹能夠掌握真理.
掌握福音的內容.
就這麼多了.
我們不可以有無限的智慧.
知識和恩賜.
而且還有就是.
人的時間的限制.
就是說在現在有限的地上的時間.
我們就只能夠這樣去服侍.
不過有一天會完結了.
就是那個完全的來到.
所以到底這個完全是代表什麼呢?.
有兩個不同的看法.
這個完全可能就是將來.
整本聖經的正典的完成.
或者更加直接就是主耶穌的再來.
就一切都不需要了.
不需要再有講道.
不需要再有講方言.
不需要再有傳遞知識等等.
所以這個就是這個意思了.
不過我們暫時再看下面.
第十一節.
我作孩子的時候話語像孩子.
心思像孩子.

$^{441}$本來這裡是講一個自己的類比的比喻.
就是講他自己心思像孩子.
意念像孩子.
既成了人.
就把孩子的事丟棄了.
模糊不清.
就是說用一個類比.
小孩子的時候.
心智都還沒有成熟.
做很多事情可能都有很多的幼稚.
但是既成了人.
就把孩子的事丟棄了.
就是不再這麼幼稚了.
不再這麼自我中心了.
不再去看人事物看得這麼簡單了.
這個是保羅大概的意思.
當然他更加想講的就是第十二節.
我們如今.
就是所有信徒.
在這個有限的地上.
神給我們的了解和知識.
模糊不清.
到那時就是那個完全來到.
或者到那時就是主耶穌的降臨.
或者我們去見主耶穌的那個時刻.
到那時就要面對面了.
我如今所知道的有限.
到那時就全知道.
如同主知道我一樣.
就是將那個有限.
將那個地上和永恆的.
全被的.
不再有任何的需要去第三者.
就可以見到耶穌.
見到神的那個永恆的天家.
全知道的那個將來的那個警方.
那個的場景.
做一個強烈的對比.
那就是說有限的和永恆的一個對比.
所以現在有限裡面.

$^{481}$我們要在教會裡面去追求彼此的服事.
其中怎樣能夠融合呢?.
就是那個教會必須所有人都有愛.
看愛為最妙的道.
看愛為大家都要擁有.
要明白的.
要掌握的.
因為很重要.
我們想解釋的就是十三節.
十三章十三節要背熟,保熟.
因為保羅就是講信忘愛.
他說如今常傳的.
就是有信有忘有愛.
他沒有講到三者的關係.
沒有講到怎樣可以得到信.
怎樣可以得到忘.
怎樣得到愛.
他只是說如今.
就是說現在在這個地球上.
每一個信主耶穌基督的人.
你都需要運用信忘愛.
他沒有講到怎樣運用.
怎樣等等.
不過他強調一樣東西.
這三樣東西.
信忘愛這三樣東西.
哪樣最大?.
他就講了.
其中最大的就是愛.
就是這樣.
為什麼?.
他沒有解釋.
沒有告訴我們.
愛因為比較重的.
不是的.
信忘愛都是如今常傳的.
就是說如果回應上文.
將來去到神的永恆.
其實就不需要再信.
不需要再盼望.

$^{521}$但是愛是繼續要存在的.
還有留意新約聖經.
神從來沒有提過.
祂是神就是信.
沒有這件事.
神就是盼望.
沒有這件事.
但是神就是愛.
就講了很多次.
尤其是《約翰一書》.
講了很多次.
神就是愛.
《約翰一書》四章十九節.
你可以想像得到.
原來不同的使徒.
都是一脈相承.
來看愛是神的屬性裡面.
甚至是永恆裡面.
都是會存在的.
是我們必須要去擁有和學習的.
第七節.
親愛的弟兄們.
四章七節.
我們應當彼此相愛.
因為愛是從神來的.
凡是有愛心的.
都是由神而生的.
並且認識神.
會有愛心的就不認識神了.
接著他就很清楚了.
因為神就是愛.
很清楚的.
當然我們不是敬拜愛.
愛不是神.
神是愛.
可以反過來的.
愛就不是神.
我們是敬拜神.
我們不是敬拜愛.
敬拜這位獨一真神.

$^{561}$但是神的一個屬性.
很重要的一個屬性.
我想是全部的.
是愛.
但是神就不需要信心的.
神不需要對我們有信心的.
我們就反過來.
我們需要對神有信心.
因為這是地上的.
在這個物質的地球.
還是一個墮落的世界.
是神正在經營的墮落的世界.
所以我們今天.
由我們信耶穌開始.
我們就要一直發展這個信心.
所以是一次相信.
一直相信.
其實應該說一生相信.
我們就必然的得救.
盼望是伴隨著信心的.
所以是相輔相成的.
神就不是盼望.
但是我們人需要盼望.
其實就算不信耶穌的人.
他們都需要用信心和盼望的.
他們都知道愛心的.
其實信望愛.
不只是存在於信徒.
其實存在於世界.
是神賦予人類.
有的那種的.
內心的一種特質.
才是為人.
所以我們就需要.
所以動物是不會盼望的.
沒有這些東西的.
他都不會怎樣運用他的信心.
他只靠本能.
但是我們人類就不同了.
所以信望愛.

$^{601}$如今常存的.
在地上.
其實運行了一樣東西.
很重要的.
就是由我們信主耶穌的一刻.
我們就必須盡快發展信望愛.
如果我們的信心.
我們的信仰.
根本就沒有信望愛.
是很大的一件事.
逐一的就很大件事.
因為去到永恆.
我們都會繼續有這些屬性.
但是如果我們根本都不會發展.
在有生之年.
我們都看輕我們的信心.
我們根本都不介意.
我們自己有沒有那個信心.
盼望.
我們根本完全都不介意.
這樣就很大件事了.
再加上愛心就更加.
因為神就是愛.
我們在神的生命裡面.
如果我們是沒有愛.
所以聖經說的愛是很重的.
你所見到的弟兄.
你就等於沒有愛.
你不是神所生的.
在《約翰一書》這樣說的.
所以原來信望愛.
是由我們信耶穌的一刻.
就要發展的了.
就要不斷去檢視一下.
我們有沒有這樣的特質.
最近衛生署不是做了一個新的表.
就說過去那些太舊了.
現在有一個新的量度.
就是現在的青少年的指標.
原來真的高了.

$^{641}$比起以前的那些年輕人.
你真的要有一個量度的.
當然有平均.
也有平均之上平均之下.
證明我們的信望愛.
是要與日俱增的.
那個成長.
這個正正是保羅所說的.
所以第十三節是這樣的意思.
如今常傳.
是有信有望有愛.
其中最大的是愛.
所以我們就不是要反過來去問.
為什麼?為什麼愛是最大的?.
其實不用問那件事.
保羅的意思就是說.
我們現在在地球上.
還沒有離開這個世界之前.
還沒有進入永恆之前.
我們必須信徒要發展信望愛.
要令我們的生命.
能夠很成長和成熟.
但是愛是永恆.
也不會再改變的.
我們去到永恆.
我們一樣是會用愛.
我們不需要再因為見不到而相信了.
我們不需要因為還沒有出現.
還沒有兌現我們要盼望了.
不需要了.
去到永恆就應該沒有這種信望了.
但是愛是永恆的存在.
由做物主做人之後.
因為他自己本身就是愛.
所以他所創造的宇宙.
所創造的人類.
都是賦予了這一份的愛.
這是《異不所思》一章所說的.
有因愛我們.
就來去借助他的愛字.

$^{681}$將我們拯救.
這裡就是第一個.
保羅所說的信望愛.
他沒有解釋其中的互為關係.
這裡就是他用信望愛的.
為何我要大家一起去做這個分析呢.
原來不是機械式的.
保羅的信望愛.
原來是可以保立體的.
所以我們要去看一看歌羅西書.
歌羅西書特別強調盼望.
盼望令信心和愛心去運行.
你就可以想像到.
保羅用的信望愛.
就是一個很靈活的.
或者很豐富的詞彙.
所以我們其實都是要學習.
歌羅西書第一章.
第三節開始.
首先就是哥羅西教會不是保羅建立的.
他聽到哥羅西的信徒非常好的靈命.
也是透過爾巴弗.
和他一起坐牢的傳道人.
就知道哥羅西教會的情況.
也都可能有些異端.
開始在教會裡面出現.
所以保羅知道之後.
就寫這封歌羅西書.
讓頂子妹認識.
尤其是對基督的救贖.
他的榮耀的本質.
那個認識.
第二章就提到.
可能他們會將基督的元秘降格.
一章的第三節開始.
我們看看保羅如何在歌羅西書裡面用信望愛.
我們主耶穌基督的父.
上上為你們禱告.
因聽見你們在基督耶穌裡面的信心.
並向眾聖徒的愛心」.

$^{721}$留意第五節.
「是因為那給你們存在天上的盼望」.
然後解釋這個盼望是甚麼呢?.
「這盼望就是你們從前在福音真理的道上所聽見的」.
即是說這個盼望是源於福音的真道被聽明白.
被傳遞.
所以就產生了這個盼望.
而這個盼望就會帶來信徒對誰有信心呢?.
因信心來到去在基督耶穌裡面.
並向眾聖徒的愛心.
因為這個真理的道被傳遞.
被聽明白.
信徒就擁有了這個盼望.
而這個盼望就帶來對主耶穌有信心.
對眾聖徒有愛心.
這樣的意思.
是源於這個盼望.
我們很少這樣想.
那個真理的道會使我們有盼望.
而這個盼望大到一個地步.
使我們不知不覺就直接和主耶穌有一個生命的關係.
不再輕易隨便嘲笑耶穌.
反而尊重耶穌是我們的救主.
看耶穌是一件很嚴肅的事.
不再當一個笑話來講.
是對主耶穌有信心.
然後不知不覺就會對眾聖徒有愛心.
因為盼望.
是因為這個存在天上的盼望.
原來這個盼望不是那麼短暫.
不是那麼局限性.
是天上的盼望.
這個盼望是怎樣來的呢?.
為何會擁有呢?.
是因為傳給我們的真道.
所以你看到一個層次.
一個層次這樣講.
保羅這樣豐富的運用信望愛.
或者他理解的信望愛.
可以是這麼豐富.

$^{761}$所以我們如果信耶穌.
信耶穌是甚麼來的?.
就是信望愛.
可能你只是念口黃而已.
你根本就不明白.
真正的信望愛那種互為關係.
互相的那種來去共融帶來的福氣.
首先就是保羅強調.
是因為這個真理的道.
就是這個福音被傳講.
而帶來這個盼望.
這個盼望原來一直就蘊存在天上.
現在被聽明白的人.
就可以來擁有了.
即是說今天我們每一個弟兄姊妹.
我們現在有一個很永恆的盼望.
共同的.
就是對天堂的那個憧憬.
不是一個海市巡流的幻覺.
是真真實實的一個地方.
就是我們要進入的那個天家.
神的居所.
在那個永恆裡面.
是有你和我的位置的.
這個就叫存在天上的盼望.
是因為這個福音被我們聽明白了.
然後我們就會對主耶穌有信心.
然後就會對眾聖徒有愛心.
所以就很闊了.
就是我們大家都需要追求了.
所以信望愛我們就要來去運用了.
接著去到《貼撒羅尼加前書》.
或者我們先看《羅馬書》.
又比較少去看《羅馬書》第五章.
又是講信望愛一併出現.
《羅馬書》第五章.
我們既因信清義.
就直著我們的主耶穌基督.
得與神相和.
直著.

$^{801}$就是因為耶穌的意思.
在這一段十一節裡面.
出現直著很多次.
甚麼都是直賴耶穌的.
一切你能夠蒙的福.
都是直賴主耶穌的.
即是直著.
第一個直著就是可以與神相和.
to find in peace.
在神的裡面.
與神之間的和好.
即是意味著一章十八節.
神的憤怒臨在天下.
我們就是那個憤怒之子.
我們就是與神為敵.
現在在主耶穌裡面.
所以由一信那一刻.
我們就已經立即.
to find peace.
在國裡面.
所以這個就是與神相和.
第一個就是直著耶穌.
第二.
我們又直著他.
就是耶穌.
因信得進入現在所站的恩典.
這個是一個恩典的地位.
蒙恩寵.
特別被神愛的一個身份.
我們現在因為信進入.
所以是一個進程.
所以一個教徒.
應該越來越進到神的致勝所.
應該進到那個恩寵裡面.
我們是否真的感受到自己是蒙恩寵呢?.
這個就要學習了.
要追求了.
要成長了.
你如果出生在一些富三代.
富四代.

$^{841}$那沒有什麼問題.
你自小都不知道為什麼.
你就要讀一些名校.
中學你可能已經去英國留學.
然後出入的都是一些上流社會.
那你不難理解你自己是蒙恩寵.
家裡一出生就有很多工人.
出入一定是坐車.
去的通常是一些很漂亮的地方.
那你不難覺得自己是蒙恩寵.
如果你是生在一些很貧窮的地方.
住tong 房的.
出生你就以為在tong 房裡長大.
你家裡轉身都轉不到的.
那你很難一走出街.
你說我是蒙恩寵.
很難的.
接著你長大之後.
你發現你是單親家庭的.
你更加難覺得自己是蒙恩寵的.
在學校可能你開始會被人欺凌你.
你更加難覺得自己是蒙恩寵.
所以信徒要進入.
保羅說.
我們要因著他.
因信得以進入現在所站的恩寵.
得以進入.
人生我們有多少的進入呢?.
每一個星期來禮拜堂.
我們有沒有一個進入聖殿.
希望能夠感覺到自己是蒙恩寵呢?.
到你自己真的發現.
我的人生是多麼的豐富.
多麼的蒙福呢?.
我們每一個都想要福.
每一個都想躋身在很美的地方.
人們尊敬我們.
一見到我們.
牧師你來真開心.
就像會所那些.

$^{881}$經常開玩笑.
和西餐的職工.
經常開玩笑.
牧師你是不是調了回來呀?.
經常都在的.
我說不是開玩笑.
又不是的.
那些是工作關係的打招呼.
但是我們信徒.
我們是要進入這種恩寵.
是要恩信.
這裡就說信這個部分.
接著不單止.
並且歡歡喜喜的盼望神的榮耀.
既在恩寵的那個.
不是感覺.
是那個身份地位.
所以你自己要.
行為自己是一個蒙恩寵的人.
當然不是要穿很美很貴的衣服.
裝飾.
是我們的生命的.
信望外的成長.
這樣叫做裝飾.
這樣來感覺自己是蒙恩寵.
接著是歡歡喜喜的盼望.
不是感覺絕望的一種灰心.
是歡歡喜喜的盼望.
神的榮耀臨到.
神的榮耀怎樣臨到呢?.
這個本來就已經用這麼豐富的詞彙.
去說有信心的人.
感受到自己是蒙恩寵.
是進入那個恩寵的身份.
那個地位.
那個感覺裡面.
你想想那個人生有多豐富呢?.
接著還要是歡歡喜喜的盼望.
神的榮耀.
已經將信和盼望拉上那個關係了.

$^{921}$就是基督徒直著耶穌.
就是這個大前提了.
如果不直著耶穌.
我們根本說不上甚麼恩寵.
我們是敵對.
我們是與神為敵.
我們是在神的震怒底下.
但現在直著耶穌.
我們一定要感受到我們是蒙恩寵.
我們要找到這個恩賜.
這個是第二節的意思.
接著是要歡歡喜喜的盼望.
歡歡喜喜的盼望.
好像亞伯拉罕他望到將來的國度.
他就興翰他現在做寄居者.
希伯來書11章那裡說.
11章好像是11節.
這個就是讓我們看到.
保羅用信心和盼望的那些拿捏和豐富.
所以不僅是裝飾一些詞彙的修辭.
而是他真的將永恆的國度.
救恩的豐富榮耀盡他的能力.
用文字去講解給我們知道.
即是說保羅自己寫羅馬書的時候.
他自己是深深感受到.
作為使徒 愛邦人的使徒.
寫信給羅馬教會.
鼓勵他們一起參加宣教的運動.
那種內心的雀躍.
那種內心的榮耀.
和他得到這個使徒的身份.
那個滿足和那份雀躍.
這個就是保羅的意思.
接著第三節不僅是這樣.
所以人生怎會只是歡歡喜喜呢.
魔鬼,世界是會令我們不歡喜的.
所以保羅說你令我不歡喜.
我就不歡喜?不是的.
第三節不但如此.
就是在患難中也是歡歡喜喜.

$^{961}$你可以想像到原來人生的患難.
是一個藉著耶穌的信徒的內心.
因為他已經掌握到.
我在神的面前.
我是蒙恩的.
我是蒙恩寵的.
我是等待榮耀的顯現的.
所以地上的苦難.
原來是神訓練我們的工場.
這樣看人生就完全不同了.
再立體一點.
原來有一個平衡的世界.
不只是地上.
不信的人只有地上.
只有今生今世.
信徒是有地上的.
的確的我現在站在二樓.
如果現在不走樓梯.
我跳下去是會痛的.
因為有地生級力的.
事實是我們真是在這個世界的.
但是信徒不同的地方.
是有一個宇宙永恆.
是我們已經現在正在發生的.
神已經這樣看我們.
是祂的兒女.
所以我們為何能夠.
不是阿Q精神.
而是我們可以在患難當中.
都能夠歡歡喜喜.
所以也是牧者的責任.
要告訴信徒.
人生一定不會沒有苦難的.
但是苦難只會鍛鍊我們更加成熟.
更加看我們這個恩寵的地位.
更加真確.
這裡說原來患難可以使我們生出忍耐.
生出老練.
即是經驗.
我們能夠有了寶貴的經驗.

$^{1001}$可以成為別人的鼓勵和一個例子.
一個人蛋.
然後再生出盼望.
而且這個盼望不是一個羞恥.
或者好像絕望.
不真實的.
虛構的.
只是自己的幻想.
不是的.
是真的.
為甚麼呢?.
因為這裡說愛了.
冠著愛了.
信和盼望.
能夠一路越來越成長的原因.
第五節就說了.
因為所賜給我們的聖靈.
將神的愛囂冠在我們的心裡.
因為我們還是軟弱的時候.
基督就按所定的日期為罪人死了.
所以就將信和盼望.
又來到同一時間這樣說出來.
令到我們可以繼續有信心.
在自己的身份.
在自己的地上.
現在擁有的地位.
本來不是屬於我們的.
本來我們根本沒有資格享有一個恩寵的.
但是現在因為我們藉著耶穌.
我們就在神眼中蒙恩寵.
而且我們是歡歡喜喜的盼望.
神將來的榮耀的顯現.
然後我們實際是在面對人生的.
保羅就說.
對的.
人生就是一個訓練場.
訓練我們的信心.
訓練我們的盼望.
使我們的盼望和信心.
更加的成長.

$^{1041}$如果好像小朋友的一個成長比喻.
我們的橫顱越多.
那個成長就會越更加堅固.
原因原來是我們的心靈.
是有燃料.
使到我們的生命.
不會好像變成了滑了.
好像沒有了那個動力.
原因原來神的愛.
會藉著聖靈.
囂觀在我們的心裡.
怎樣囂觀呢?.
又是很抽象的.
我們都不懂.
聖靈怎樣將神的愛囂觀呢?.
我們都不太知道.
每天是不是幫你打點滴呢?.
這些我們知道.
你住在醫院.
你吊著一條繩子.
那個喉.
在那個手上.
那個叫什麼?.
黃豆是吧?.
這樣插下去打點滴.
我們知道了.
一滴一滴這樣滴下去.
我們知道這個叫點滴.
囂觀.
但是抽象的.
是靈界的.
聖靈怎樣囂觀神的愛呢?.
最終是使到我們的信心.
我們對自己的身份.
對自己的地位.
那個蒙恩寵.
和歡歡喜喜.
心靈歡歡喜喜的.
有盼望對榮耀降臨的主.
我們有一個熱切的期待.

$^{1081}$就是這個聖靈的工作.
在我們的生命裡面.
所以聖靈的工作.
不是要這麼突然的.
不是一定要神蹟的.
不是一定要講方言的.
不是一定要按手醫病的.
聖靈在每一個信徒的心裡.
來將神的愛囂觀.
像針水一樣.
淋花.
所以是生生不息.
將神的愛.
源源不絕地灌注在我們的心靈裡面.
所以一個基督徒.
如果他真是清楚自己的身份和地位.
是不會有一個灰心和絕望的.
這就是羅馬書運用信望愛.
我們又豐富了這一點.
我們回來看《貼撒羅尼加前書》.
因為涉及到和貼錢有關.
所以也很想大家能夠看看這裡.
一章可能是第五節.
或者是第三節開始.
《貼撒羅尼加前書》一章第三節.
在神我們的父面前不住的紀念.
你們因信心所作的功夫.
這次這裡就很實在了.
有信心的基督徒.
信了主就會為神作功.
為神去作甚麼功呢?.
相信一定是和其他人的生命.
信主,靈明得到建立的功夫.
一定是正面的做在別人身上的奴婦.
所以無論甚麼崗位也好.
是幫助到那個人能夠成長的.
那個教會能夠健全的.
這就是因信心所作的功夫.
一定是正面的了.
所以如果我們去拆毀教會.

$^{1121}$保羅的意思很重.
如果我們去拆毀教會.
神都會拆毀我們.
所以我們真的不要做一個傳遞錯誤消息.
煽動別人或者是擴散紛爭.
千萬不要.
我們應該成為一個和平之子.
所以因信心所作的功夫.
然後因愛心所受的奴婦.
有愛就不覺得是奴婦了.
在座的媽媽,婆婆,奶奶就很明白了.
其實爸爸也是一樣的.
在外面多辛苦都是為了自己的兒女,家庭.
所以有些研究很恐怖的.
一個如果沒有負擔.
沒有家庭壓力的人.
自殺率是高一些的.
因為他沒有什麼需要牽掛.
沒有什麼需要他去拼命.
被人罵兩句都要撐住.
因為家裡的化骨龍很重要.
你寧願忍都可以面對.
原來這樣會令男兒氣概不知不覺就會建立.
反過來一個完全沒有負擔的人.
只為自己就已經夠了.
他的自殺率是高的.
這是有研究的.
某程度上是很順理成章的.
我都沒有什麼需要為人去幫我自己.
我喜歡就做.
我覺得這個世界不好玩我就走.
就沒有任何所謂他要負的責任.
所以責任帶來的一定是和那個奴婦拉上關係.
有愛心就會勞苦.
所以我們其實要在教會裡面.
建立一份對神對人的愛.
不要太快去做事.
先建立了這一份對神的愛.
為什麼我們會成為基督徒.
神為什麼這麼愛我.

$^{1161}$愛大家.
我們明白了之後.
我們就看到其他人都是很蒙愛的時候.
我們大家就明白.
那個功夫那個做是為神去做.
是為神的福音為教會去做.
這樣你就不會再有一些比較.
和會有很多的埋怨.
過程都會有的.
但是慢慢你就會知道.
因為那個更大的價值觀是為神去做.
所以因愛心所受的奴婦.
因盼望我們主耶穌基督所傳的忍耐.
同時間發生了.
有信有愛有盼望.
是正正常常的.
因為基督徒.
這裡似乎是說一個信了主耶穌的人.
他就會具備這個特質.
就是因信所作的功夫.
愛所受的奴婦.
盼望主耶穌基督生存的忍耐.
就是說盼望對方見不到.
你見到你相信.
人家就會嘲笑你.
所以你要忍耐的.
你要告訴人家有天堂.
他說你是傻的.
怎麼會有天堂呢這個世界.
在哪個宇宙的地方呢.
物質的宇宙在哪裡呢.
我們都不能夠證明給他看的.
但是有盼望的人就會忍耐.
所以我們是需要發展信望愛.
是一個基督徒.
一信耶穌就與生俱來有的了.
就是說救恩帶給我們.
立即其實是令到我們有信望愛.
或者這裡叫信.
這裡就是信望愛.

$^{1201}$所以你可以想像到.
剛才的歌羅西書.
是從那個盼望而來的.
而導致到我們在信心.
愛心上的增長.
但是貼在羅尼加前書.
阿保羅又不是機械式的.
他又會將這樣的一個基督徒.
貼在羅尼加教會的弟兄姊妹.
這樣來描繪.
自從福音進入他們的生命裡面.
因為聽見他們因信所作的功夫.
因愛所受的勞苦.
盼望主耶穌來的時候.
那個的忍耐.
生存的忍耐.
這樣就是說在一個健全的基督徒.
其實信主只有四個星期.
一個月時間.
貼在羅尼加.
保羅寫這封信的時候.
可能是半年.
你當半年吧.
可能沒有半年.
這樣的基督徒.
就已經因為福音而帶來的.
那個信望愛了.
所以基督徒我們要追求信望愛.
我們要對救恩明白.
而帶給我們對信心,愛心和盼望的.
很實質的我們自己都要盤點.
我們真的需要這樣來成長.
但是我們透過這幾段經文.
讓我們看到.
保羅運用的信望愛.
是很豐富的.
在不同的處境.
他能夠講到.
那個的.
不同的豐富.

$^{1241}$和那個特質.
使得我們可以在追求的過程.
更加的去體驗.
神永恆的那些恩賜.
是給我們每一個人去追求.
就是說.
信主之後.
原來是有一些成長的階梯的.
雖然不是物質去量度.
不是說頒個證書給你.
但是剛才幾段聖經都講了.
是有不同層次的信望愛.
這個就真是我們一生要追求.
記住.
十三章十三節.
我們去分禮都會背的.
但是那裡不是說分禮.
如今常傳的是有信,有望,有愛.
其中最大的是愛.
那裡是想告訴我們.
地上現在包括去到永恆.
我們都需要有這個信望愛.
所以如果基督徒在很有限的幾十年.
不去追求信望愛.
不在自己的生命裡面去發展信望愛.
不單止是可惜.
甚至將來受到很嚴重的虧損.
因為你可以想像到.
如果我們沒有愛心.
我們帶著一種憤怒離開這個世界.
你面見的是一位愛的神.
你怎麼辦?.
你還帶著一個遺憾.
我怎麼都不會饒恕.
遠離的了.
一輩子我都不會饒恕他.
我死都不會饒恕.
你講完之後就真的死了.
你去到神那裡.
神就是愛.

$^{1281}$你情何以堪?.
所以我們很多人.
根本我們基督徒.
整個世界不能夠看到有神的原因.
是很多人看到基督徒沒有愛.
看不到神的真實.
看到很多基督徒爭來爭去.
為了一些無謂的東西去浪費人生.
很可惜的.
做牧者我們要教導.
如果我們沒有愛.
而我們有憤怒有仇恨.
你都不用想了.
你的靈命一定會下跌.
如果有排名榜.
當然一定沒有.
你會跌出十大.
甚至每年有很多大學排名.
或者什麼的排名.
如果基督徒可以排名.
你的樣子可能很高.
但是你的生命可以很低.
是因為你裡面有一絲仇恨.
因為神就是愛.
這是告訴我們.
如今常傳的是有信有望有愛.
其中最大的就是愛.
所以一定要追求成長.
信望愛的成長.
怎樣成長呢?.
保羅在其他書卷就說了.
羅馬書說了.
歌羅西書說了.
傑瑟羅尼加前書又說了.
其他的書卷都有說.
所以你就可以想像到.
這麼豐富的詞彙.
不只是念口黃.
不是機械式.
不是基督教就是信望愛.

$^{1321}$那些不信的人都是這樣分析的.
對的不是錯的.
不過他們所謂的信望愛.
只不過是一種術語.
是沒有內涵的.
盼望基督徒我們弟兄姊妹.
真的要在信望外面成長.
盤點又去到五月了.
真的要盤點.
我們的信心是否越來越增長呢?.
我們的盼望.
盼望甚麼呢?.
有沒有想過希望教會.
能夠在我們這個時代有大復興呢?.
有好像馬丁路德的人出現呢?.
有嘉爾文這些改革家的出現呢?.
帶來整個社會有一個改變呢?.
我們要有這個盼望.
我們要有那個愛.
對人對神的愛.
要純全.
沒有雜質.
不會有一些錯誤的動機.
這些全部都需要追求.
所以盼望我們都能夠在禱告裡面.
去追求這個信望愛.
還有在我們的實踐裡面.
我們的下一代他們應該都敏銳.
我的爺爺,我的媽媽,我的爸爸媽媽.
有沒有信望愛呢?.
他們應該會看得清楚的.
所以我們是否帶了這些.
很重要的屬靈的遺產給他們呢?.
我初信的時候.
初信就是我講的二十多三十年前.
我1988年洗禮.
1988年的6月3日洗禮.
就是信主之後的第二年我就洗禮了.
1987年的12月24日信主.
平生第一次去教會.

$^{1361}$1988年的6月3日就洗禮.
那個時候就聽了一個見證.
很深刻的.
因為他們整個家族.
幾兄弟都在我們的母堂教會裡面.
很熱心事奉的.
姓王的.
這個不知道是不是大哥還是甚麼呢?.
叫王欣華.
1988年的時候已經離世了.
南正離世了.
就是死了很多年了.
他臨死之前告誡他的弟弟妹妹.
那些兄弟一定要愛主,事奉主.
這樣離世.
沒有遺憾的.
突發性的.
很忠心事奉的.
在教會裡面很多人都很景仰他們.
但是很年輕就離世了.
他幾個兄弟都很熱心.
在教會有時幫忙講道.
很深刻的.
你想想二十多年前的一個小小的見證.
我都還這麼深刻.
你可以想像得到.
就是這些的榜樣.
就是這些的信脈外的人.
人生.
如果我們能夠看到有些人有信脈外.
我們一定要與他做朋友.
所以與智慧人同行.
必得智慧.
和愚昧人作伴.
必受虧損.
是這麼說的.
所以我們已經很有限了.
我們要認識那些要智慧的人.
要和真正愛神,國度的人交往.
如果你的好朋友都是那些講是非.

$^{1401}$愛世界的人.
你都不用問了.
你的心靈.
你自己都喜歡那些人.
你才會和他們在一起.
你不喜歡就不想和他們在一起.
所以盼望我們都能夠追求.
用了少許時間講信脈外.
我們回到第三章.
《跳瘦羅尼加》前書的第三章.
我們既不能再忍.
就願意獨自等在雅典.
打發我們的兄弟.
在基督福音上.
作神執事的提摩太前去.
就是回到跳瘦羅尼加.
堅固你們.
並在你們所信的道上.
勸慰你們.
免得有人被諸般患難搖動.
因為你們自己知道.
我們受患難.
原是命定的.
這樣的教導真的挺健康的.
初信已經告訴你.
你要跟隨耶穌背實字架.
可能真的會有苦難的.
是你自己想清楚才相信.
所以是預定的.
不可能好像我們要告訴別人.
信耶穌就什麼都改變.
什麼都開心.
如果不開心的時候.
他不相信.
你就說不要這樣.
是你告訴別人開心.
所以保羅也很正面的.
當時根本是一個逼迫的時代.
你也不需要說.
因為羅馬帝國也是逼迫基督徒.

$^{1441}$第四節.
我們在你們那裡的時候.
預先告訴你們.
沒錯 預先告訴了.
我們必受患難.
以後果然應驗了.
不是要說他們懂得怎樣預言.
而是要給你們一個心理準備.
背實字架永遠是信耶穌不能逃避的.
是耶穌也說的.
誰要來作門徒.
必須要否定自己.
背實字架跟隨耶穌的.
第五節.
因此我既不能再忍.
就打發人去要曉得你們的信心如何.
恐怕那誘惑人的.
到底誘惑了你們.
叫我們的奴婦歸於途人.
接著就說提摩太回來.
但提摩太剛才從你們那裡回來.
寫這封信的一個比較主要的原因.
就是保羅在雅典.
很擔心他離開.
喀薩羅尼加教會的弟兄姊妹.
面對沒有了牧者.
又有一些攻擊.
也有些人傳出一些錯誤的消息.
說保羅扔下你們.
自己走了去.
因為有患難.
所以就見異思遷.
所以保羅就很擔心.
打發提摩太他回去.
即是他自己現在在雅典這個地方.
那裡有一個很大的神殿.
巴特隆神殿.
當時是一個很熱鬧的地方.
人們敬拜那些假神.
保羅就在那裡等候.

$^{1481}$第六節說.
但提摩太剛才從你們那裡回來.
將你們信心和愛心的好消息報給我們.
又說你們常常紀念我們.
切切想見我們.
如同我們想見你們一樣.
所以弟兄們.
我們在一切困苦患難之中.
因著你們的信心就得到安慰了.
你們如果靠主賺得穩.
我們就活了.
我們在神面前.
因著你們甚是喜樂.
為這一切的喜樂.
為你們去報答神呢?.
我們晝夜切切祈求.
要見你們的面.
補滿你們信心的不足.
最後這一段關於主再來.
願神我們的父和我們的主耶穌.
一直引領我們到你們那裡去.
又願主叫你們彼此相愛的心.
並愛眾人的心都能夠增長充足.
如同我們愛你們一樣.
好使你們當我們主耶穌.
和祂眾聖徒來的時候.
相信是所有的眾聖徒.
不只是一些特選的精英的.
是所有在主耶穌裡面.
死了的聖徒.
都會跟隨耶穌一起回來.
這是四章裡說的.
《喬治羅尼加前書》四章.
主耶穌和祂眾聖徒.
再回來的時候.
你們在父神的面前.
將來每一個信徒.
特別提到貼心羅尼加.
在父神的面前.
心裡面堅固成為聖結.

$^{1521}$無可襲地.
即是說貼誠的弟兄姊妹.
需要在現在的人生.
就算沒有目的也好.
他們都不可以荒廢自己的信仰.
原來有一個日子.
要看你是不是有忠心.
是不是有可指責.
即是說你有沒有成長.
信主正正是因為我們.
沒有任何好的行為.
在神面前可以得救的.
就是每個人都是沒有好行為.
因為耶穌而得救的.
但是信了耶穌之後.
保羅說我們有一個日子.
將來神要驗證我們.
是不是無可襲避.
是不是心靈純全.
堅固.
堅固即是說你的心對信仰.
是很穩固的.
這個信念改變你的生活態度.
改變你的待人接物.
對事物的研判.
這叫心裡堅固.
成為聖潔.
在神的眼中.
不只是乾淨這麼簡單.
成為聖潔.
而且是無可襲避.
保羅這樣說主在來.
不是我們開心.
馬上生對翼.
飛到天堂.
他不是說那些.
那些可能不在話下也說不定.
反而是今生我們要怎樣裝備好自己.
在主耶穌回來的時候.
我們具備這個升天的那個人.

$^{1561}$不是升天的那個方法.
升天就一定是.
但是這個人的生命.
怎樣去到神的面前.
去到永恆呢?.
尤其是強調無可襲避.
保羅是不是一個很務實的人呢?.
這麼著重你的成長.
為何不說一些我們很興奮的事呢?.
怎樣突然間變呢?.
怎樣可以完全不需要再有空間的限制呢?.
這些是重要的科學的發現.
你想想你同時間在203.
你又不可以跳過202.
有牆阻住你要開門出去.
怎樣可以這樣穿牆過壁呢?.
我們好奇這個世界是沒有的.
可以走在一個不用走的路上.
一去就去到對面的大交樓了.
我們都好想知道.
但是保羅都不感興趣這些事.
反而是我們每一個基督徒.
在主耶穌和眾聖徒來的時候.
我們具不具備這種的聖潔.
具不具備這種的心靈的.
很堅定的那個信念.
應該叫信仰.
還有我們是否真的無可責備在神的面前呢?.
這個就是我們有幸的人生要追求.
所以原來主在來.
不是叫我們好奇等待那個日子.
計算一下通勝什麼時候來呢?.
不是這樣的.
不斷都不是計算.
是我們裝備自己.
盡我們的能力去多認識主.
多去愛人.
多去傳福音.
多去準備自己進入那個永恆的榮耀.
所以第三章這裡.

$^{1601}$其實這裡又發生了一個很大的.
值得我們思想的課題.
今天每個時代都有不同的方法門徒訓練.
小組訓練或者門徒的訓練.
或者什麼什麼的查經班,言經班訓練.
這些全部不是錯的.
但是保羅怎樣跟他的弟兄姊妹的關係.
能夠很密切.
很能夠有一種牧養的生命的灌輸成長呢?.
原來就是對弟兄姊妹的敏銳的心靈.
關心和同行,陪伴.
這些特質.
所以任何的課程都是好的.
但是如果課程只是為授課而授課.
而沒有了那種生命關係呢?.
所以就是一個問題了.
雖然好像很敏感.
我們很熱衷參加不同的課程.
但是很少大家有彼此相交的生命成長.
這些課程變成只是一種知識的傳授.
其實不足夠.
不可以說錯是不足夠.
必須是有生命的大家的接觸.
大家的磨合.
大家的成長.
必須是要在教會裡面發生的.
所以很多信徒很熱衷去參加晚上的課程等等.
不同的神學院開的課程.
全部是對的.
完全是幫到那個人的.
但是回到教會呢.
冷冰冰.
跟人一點都不接觸.
這樣就出現了一個問題.
你根本就沒有一個場景去訓練你.
讓你的知識可以成為一個實際的操作.
可以在彼此相愛的時空裡面.
真正的生命成長.
所以這一大段聖經.
是一位暫時肉身跟這個教會分開的牧者.

$^{1641}$他怎樣在他的智慮行間.
講到他內心對牧羊的那個情懷.
對這一班弟兄姊妹的那份愛.
那份擔心而來的愛.
因為怕他們會跌倒.
保羅現在暫時在礙於可能很多因素.
甚至有神的帶領.
他要離開貼薩諾尼加.
導致到他.
他都不是介意人家怎樣看他.
是不是真的背棄了這班人.
而是他現在在雅典.
他很著緊要來安慰弟兄姊妹.
所以都是要人.
都是要年青的傳道者提摩泰.
帶著保羅的那個安慰去.
等於代表著保羅.
回去探望貼薩諾尼加的弟兄姊妹.
都是需要有牧者的牧羊.
教會必須是要有牧者的牧羊.
如果沒有牧者的教會.
其實可以的 人工都不需要了.
每個人都義務的 省很多的.
如果那層樓又不用交租.
更加沒有負擔.
不過那個教會就不能夠成長.
因為沒有了這份牧羊的愛在其中.
因為每個人都是義工.
很多衍生的問題.
這裡告訴我們.
保羅就是這樣示範.
其實他不是要特意.
好像來做一些教書教人.
而根本就是他生命上的一個情懷.
他成為了牧者.
就和這個教會帶著他們信主.
他就有這個義務 有這個責任.
有這份的愛.
就是這樣來衍生了.
所以我們如果帶著人信主.

$^{1681}$我們自己都起碼要有一段時間.
陪伴他們.
讓他們能夠在那個教會成長.
能夠穩固.
才離開就會好一些.
當然要陪伴他們一起成長.
信主的時候初期會有很多的困難,問題.
如果能夠有和他們一起同行的人.
就會帶來一個很大的幫助.
所以這一段聖經.
其實是一個很重要的牧羊的教導.
保羅就是這樣.
來趕緊的打發提摩泰回去.
首先就是了解他們的信心.
還要保滿他們信心的不足.
就是說牧者總是有些東西可以教弟兄姊妹.
總是有一些安慰的說話.
提醒的說話.
可以來勸勉弟兄姊妹的.
總是有些東西.
是很想讓他們能夠得到那個教導.
所以這個就是第二節.
提摩泰前去.
這個作神執事的提摩泰.
在基督福音上面.
作神執事.
這樣來不是誇耀提摩泰.
提摩泰是這樣的一個牧者.
在保羅的眼中.
這個年輕人.
都還沒有在這個二福所教會.
來去住宅.
第三次宣教旅程.
保羅就留提摩泰在二福所的大教會.
現在是跟隨保羅在第二次宣教旅程.
剛剛在路斯德見到這個年輕人.
為弟兄姊妹所愛戴.
所欣賞的這個年輕人.
這個提摩泰就帶他去傳道.
在保羅的眼中.

$^{1721}$保羅認定這個提摩泰.
這個年輕人.
有一個這樣的身份.
在神的基督的福音工作上.
他是神的執事.
作神執事.
你可以想像得到.
這個身份.
是何等的尊貴.
神的執事.
即是神的僕人.
為神工作的.
重點是為神的基督的福音的執事.
是神的執事.
所以好像都是想教會的弟兄姊妹.
看傳道人.
年輕的提摩泰還有很多成長的空間.
很多的甩漏.
甚至他的性格比保羅都沒有那麼剛強.
但是他有的是那個身份.
和會成長的一個傳道者.
保羅要求弟兄姊妹這樣看提摩泰.
他是在神面前的執事.
而且是基督福音的執事.
這樣專注在福音上.
去傳遞,去牧養.
所以感恩的.
宣道會對牧者的身份的敬重.
都是很清楚的.
對牧者有的領袖的牧養的權柄.
都尊重和欣賞的.
如果能夠讓這件事.
可以很專職地在教會裡面發生.
這個就是教會的福了.
所以教會一定是必須要有牧者.
不可以全部是義工.
是義工就很大件事了.
因為他不會覺得自己在任何的環境.
他都要具備這個神的執事.
有困難我就走.

$^{1761}$有人不喜歡我就不做.
我義務的.
一句義務就大了.
但是如果我們是很清楚神.
既然選擇,訓練,安排.
我就不輕易說我喜歡怎樣就怎樣.
是神想我怎樣我才能怎樣.
這個叫做神的執事.
是為基督的福音工作的執事.
不同的 兩件事是很不同的.
有了這個認定的身份.
他就會鞠躬盡瘁.
所以會耐心很多的.
某程度你讚我也好.
你彈我也好.
你笑我也好.
你侮辱我也好.
我會自省當然.
正常是要自省.
但是我不會太過介意.
因為我清楚我在神所領受的職份.
這個和義工不同的地方就是這樣.
所以保羅這樣去不是要誇耀提安泰.
而是要讓教會.
要有一種這樣對傳道人的身份.
這一份由神賦予的身份的尊重.
所以要對那個牧羊的職份的尊重.
記不記得《馬太方言》23章.
耶穌來到和大家讀一讀.
《馬太方言》23章.
耶穌教的門徒.
那時耶穌對眾人和門徒講論.
就說.
「問事和番理猜人.
他們坐在摩西的位置上講解舊約聖經.
講解律法.
凡他們所吩咐你們的.
你們都要謹守遵行」.
因為講的是聖經.
神的道.

$^{1801}$但是不要效法他們的行為.
為什麼呢?因為他們能夠講.
不能夠行.
這是生命的影響力問題.
是他們不能夠以身作則.
不是他們講的東西是錯的.
所以要分很清楚.
所以耶穌基督都教導門徒.
要尊重這樣的身份.
他們坐在摩西的位置上.
事奉的時候.
你就要尊重他.
即是說他的為人.
如果真的做得不好.
你也不需要藐視他.
你帶著謙卑的心.
警惕的心去看神.
在他生命裡面的可能是改變.
可能是光照.
但是我們自己就要有一個.
看這些人而帶來的警惕.
所以我們需要有這種融會貫通.
但是在職份上.
正正就是提摩泰.
他就是具備保羅.
我跟大家熟讀一下.
我怕部分不是太熟悉聖經的人.
十六章的《士徒行傳》.
保羅來到特比.
又到勞斯德.
這是在土耳其一帶的.
在那裡有一個門徒.
名叫提摩泰.
是信主的猶太婦人.
就是這個有尼基的兒子.
他的父親是一個希臘人.
即是說他是一個混血兒.
父親是希臘裔.
母親是猶太的婦女.
即是嫁了一個希臘人.

$^{1841}$勞斯德和以哥廉的弟兄.
以哥廉是一個比較大的省.
勞斯德是其中這麼多的省.
廣東省的我家是博羅縣.
香港是寶安縣.
那個大的省就是以哥廉.
其中一個城就是勞斯德.
特比等等.
這裡說以哥廉的弟兄都稱讚他.
留意第三節.
保羅要帶他同去.
只因那些地方的猶太人.
都知道他的父親是希利尼人.
就給他行了國禮.
接著他們經過郭城.
在耶路撒冷使徒和長老所定的條規.
就是那個十五章裡說的那個大會.
詳細我不說了.
來到他交給門徒遵守.
於是宗教會信心越發堅固.
人數天天加增.
成為保羅的同工.
勞斯德以哥廉的弟兄稱讚他.
即教會裡稱讚他.
保羅亦都見到提摩泰.
有些可以造就的特質.
就帶他一齊去傳道.
未帶之前.
首先幫他行國禮.
即是提摩泰這麼大的兒子.
都未行國禮.
即是猶太的那個律法的影響力.
在這個家庭不算很大.
或者可能原因就是爸爸不準.
因為爸爸是希臘人.
不是兩個都是猶太人.
所以保羅就意識到這件事.
為甚麼要幫提摩泰行國禮呢?.
就方便他去到猶太的會堂.
就不要被那些猶太那些比較硬的人來挑剔.

$^{1881}$說你是猶太兒子.
為甚麼你不行國禮?.
就令到提摩泰可能會有點尷尬.
所以就幫他行國禮.
就不是因為他要行國禮才得救.
不是這件事.
所以這是一個規矩上.
讓他更加可以在一些猶太的人面前.
猶太裔的基督徒面前.
可以來表白他是一個愛神愛人的人.
提摩泰.
你可以說保羅初次認識提摩泰.
透過不同的口碑.
透過不同的接觸.
可能保羅就在傳道的過程.
就發現這個年青人.
可以也可能是經過祈禱.
就開始保羅效法主耶穌.
去到傳道一段時間.
他知道自己的大限將到.
盡快要物色接班人.
一個真正的有眼光的領袖.
一上任就要找接班人.
就要找到可以將他的生命灌輸給這個年青人.
所以看教會是否成長.
也可以看到有沒有第二第三代的年青人.
接上福音的棒來做呢?.
或者願不願意跟隨一些長輩.
來出入效法的榜樣呢?.
這個是指標來的.
嚴格的指標.
因為耶穌訓練十二個使徒.
保羅訓練提摩泰提多的.
還有其他的特羅非摩.
路加等等.
真是一脈相承.
如果沒有門徒.
教會真是很大有限.
是生命的傳遞.
所以保羅現在揀選提摩泰.

$^{1921}$我想他經過很多的祈禱.
很多的觀察.
還有很多在神面前的領袖.
看到這個年青人.
裡面是有一個可堪重用的特質.
將來就會被神大用.
所以去到帖撒羅尼迦前書.
其實很快的.
沒多久就去到帖撒羅尼迦.
在土耳其一坐船就過去對面.
接著就開始歐洲的宣教.
在菲律賓.
接著就去到帖撒羅尼迦.
接著沒多久就要走去雅典.
大概是半年多一點.
但現在保羅怎樣稱呼提摩泰呢?.
這個年青人是在神面前的執事.
是甚麼執事呢?.
是神福音 基督的福音.
工作上的執事.
這樣稱呼提摩泰.
弟兄姊妹我們要追求.
我們的人生都是被主用.
被主去揀選 被主去用.
我們為下一代祈禱.
希望能夠找到像提摩泰一樣的年青人.
真是對於自己的人生願意奉獻給神.
為神的福音願意做執事.
在基督的福音工作上.
這樣擺上的.
跟隨保羅去傳道.
你想像得到那個困難.
不要說那個築跡很大.
比主耶穌還要大.
而且所受的困難還更加多.
但是提摩泰就是這樣.
可以被神揀選.
也不可以欠缺的是.
保羅有這個眼光.
栽培年青的牧者.

$^{1961}$教會其中一個很大的工作.
就是必須栽培下一代的領袖.
去接續福音的工作.
今晚就說到這裡.
我們一起禱告.
感恩天父讓我們能夠梳理了.
在保羅書信裡面的順忘案.
我們很盼望我們能夠在信心.
我們的盼望和愛心.
對準我們的神.
我們的生命是一路的成長.
一路的成熟.
我們能夠對主有認識.
對福音可以改變人生命的這個能力.
我們是雀躍和參與.
更加我們能夠等候這位主再回來的時候.
不單止我們是有份.
到時是會點名這樣叫我們.
但是更重要的.
聖經一而再說的.
不是叫我們好奇的去計算日子.
或者是怎樣進入那個過程.
而是我們這個生命的成長.
在神的眼中.
我們是否成為聖見.
無可責備和心裡堅固.
願神幫助帶領我們.
多謝你給我們有機會去有系統地去研讀.
貼上羅利加前後書.
和我們前面會說的末世論.
願神繼續保守.
叫我們雀躍學無.
主你的再回來.
我們這樣感恩禱告.
奉耶穌基督的名.
阿們.
\newpage



\section{羅馬書 8:26-30}
\label{sec:a_rRhGfuxUw}
\textbf{2024-04-27 晚堂《再思天路歷程》羅馬書 8\_26-30  講員:曾裔貴牧師}
\newline
\newline
連結: \href{https://youtube.com/watch?v=a-rRhGfuxUw}{\texttt{https://youtube.com/watch?v=a-rRhGfuxUw}} ~~~~ 語音日期: 2024-04-30
\newline
\newline
\hyperref[sec:ZVwFpbgglfY]{\small{< < < PREV SERMON < < <}}
~
\hyperref[sec:index_chronic]{\small{[返順時目]}}
~
\hyperref[sec:index_scriptual]{\small{[返順卷目]}}
~
\hyperref[sec:w0vSwuU8gCk]{\small{> > > NEXT SERMON > > >}}
\newline
\newline
羅馬書 8:26-30
\newline
\begin{longtable}{cl}
\hline
\hline
章節 & 經文 (和合本修訂版)\\
\hline
8:26 & \begin{tabularx}{0.7\textwidth}{X} 同樣，我們的軟弱有聖靈幫助。我們本不知道當怎樣禱告，但是聖靈親自用無可言喻的嘆息替我們祈求。 \end{tabularx} \\ \\ \relax
8:27 & \begin{tabularx}{0.7\textwidth}{X} 那鑒察人心的知道聖靈所體貼的，因為聖靈照著神的旨意替聖徒祈求。 \end{tabularx} \\ \\ \relax
8:28 & \begin{tabularx}{0.7\textwidth}{X} 我們知道，萬事都互相效力，叫愛神的人得益處，就是按他旨意被召的人。 \end{tabularx} \\ \\ \relax
8:29 & \begin{tabularx}{0.7\textwidth}{X} 因為他所預知的人，他也預定他們效法他兒子的榜樣，使他兒子在許多弟兄中作長子。 \end{tabularx} \\ \\ \relax
8:30 & \begin{tabularx}{0.7\textwidth}{X} 他所預定的人，他又召他們來；所召來的人，他又稱他們為義；所稱為義的人，他又叫他們得榮耀。 \end{tabularx} \\ \\
[1ex]
\hline
\hline
\end{longtable}
$^{1}$我們同心祈禱.
多謝天父我們再思想這個題目.
盼望一段相當重要和熟悉的經文.
26至30節.
聖靈的幫助.
聖靈讓我們認識.
我們的生命從神而來的新生命.
是如何運作.
幫助我們對聖靈的敏銳.
我們可以在生活中.
經驗到這一份.
超過了我們能夠用言語.
所來表達的.
這一份內心的嘆息.
是神所知道.
來去成就的禱告.
願主幫助.
也真的讓我們去思想天路歷程.
我們去看看.
檢視一下我們的人生.
在走在什麼軌跡上.
這樣禱告奉耶穌基督的名.
阿們.
相信第26節.
已經說了一個因為這個字眼.
在這段聖經裡面.
有強調的是聖靈的幫助.
在22,23節.
就說到原來凡是眾生.
每一個在地上受造之物.
都是嘆息勞苦的.
所以第22節就說.
我們知道一切的受造之物.
都是一同嘆息勞苦.
一直去到如今.
但是有一個盼望.
就是將來有一個永恆.
是會有一個榮耀.
23節保羅就特別強調.
就是我們信了主耶穌的基督徒.

$^{41}$不但如此.
我們這些有聖靈初見果子.
即屬於神的人.
也是自己心裡嘆息.
等候得著兒子的名分.
乃是我們的身體得屬.
所以已經規範了一個.
很重要的範圍.
就是第26到30節.
特別很清楚強調.
哪些人會感受到.
經驗到萬事的互相效力.
其實講得很清楚.
還要用兩節聖經29,30節.
來描繪這一類的人.
到底是什麼人呢.
就是第28節.
什麼人會經驗到互相效力呢.
28節.
我們曉得萬事都互相效力.
清楚了.
只有愛神的人.
這個範圍的人.
才可以感受到萬事萬物的互相效力.
得到當然是整個生命成長的益處.
這裡告訴我們.
然後再解釋是什麼人.
強調神的很重要的救恩作為.
就是按照著神的旨意.
來被呼召的人.
然後第29,30節.
神的愛展現在這個如知如定裡面.
29.
因為他如先所知道的人.
就是我們這些今天.
雖然這麼軟弱.
不懂祈禱.
或者不能用言語祈禱的.
可能是困難.
內心的困苦.

$^{81}$有逼迫而不能夠去面對.
這樣的人.
在愛神底下的困苦.
有神的聖靈.
在最深處的心靈裡面幫助.
為你向神祈禱.
所以監察萬物的神.
就能夠知道這樣的祈禱.
幫助這樣的人.
這樣很戲劇性.
怎樣發生呢.
本來就不太詳細講運作過程.
不過講了什麼是愛神的人.
什麼是在這段範圍裡面.
聖靈幫助在那個人的心靈裡面.
他就講了那個很重要的身份.
這個是神做的.
不是我們自己可以去.
用很重的價錢買到回來的.
因為他如先所知道.
即是在創世以前.
他知道他預定.
就預先定下.
定下了一個規模.
就是原來我們每一個弟兄姊妹.
不是那些特別精英的神學院.
很高分畢業的才可以.
是每一個基督徒都可以去得到.
就是每個人都定下了.
效法他的兒子的模樣.
我們一生從信耶穌開始.
我們要像耶穌.
主耶穌的.
當然不是那些耶穌轉頭髮那些.
那些是戲劇明星.
真正的主耶穌是很滄桑的.
是無階型美容的.
他是常經憂患的.
是多受痛苦的.
那我們學什麼呢?.

$^{121}$學他的橫欄?.
相信都不是了.
是學他那個屬靈的品德.
我們的屬靈生命.
原來是有個規模.
是有個工程.
神定下每一個基督徒.
要效法他兒子的模樣.
以致他的兒子就可以在許多弟兄中.
作長子.
即是原來我們每一個基督徒.
都是神所收納.
如果再嚴格些.
adoption這個字應該是收養.
我們是被收養.
我們本來不屬於神國度的人.
我們本來不屬於神的兒女.
但因為他的愛子的犧牲流血.
我們就被收養.
我們就能夠有一個人生.
要學效他兒子的模樣.
這就是人生的每一個頂尖的工程.
都還沒講完.
保羅說 使得他的兒子在許多弟兄中作長子之後.
三十節又再次講.
講了三次的預定.
預先所定下的人.
就是我們基督徒.
又召他們來.
原來我們信耶穌.
是被神呼召的.
接著就是.
所召來的人.
首先就是經驗到清義.
所召來的人又稱他們為義.
即是神不再追究了.
因為透過他愛子的寶血.
已經一筆勾銷了.
已經可以在主耶穌的救恩裡.
得著神的義.

$^{161}$這就是被呼召的人.
經驗到的第一個初信班.
就是問你得救了沒有.
得救了就可以上初信班了.
如果還沒清楚決志.
就要去探索班.
因為真的要這樣分.
這就告訴我們.
基督徒原來被呼召.
被清義.
接著呢.
所稱為義的人.
又召他們.
不是精英分子.
是每一個基督徒.
是每一個有聖靈初結果子的弟兄姊妹.
得到什麼呢.
榮耀.
將來永恆的.
不再有罪惡的地方.
我們每一個在人生裡面.
學校耶穌的榜樣的一個基督徒.
還不是天路歷程.
還不是我們人生從信耶穌開始.
就有一場屬靈爭戰.
要面對各種的難處試探.
務求的我們要去到天城.
這就是本人約翰.
那本暢銷的書籍.
天路歷程.
一初信就要看這本書.
不過越看越怕.
很多危險很多困難.
他很多很多的患難.
但是令到他不灰心.
最終他望到.
和進入天城.
當然是他在坐監的時候.
寫的這本書.
我今天就用了這個題目.

$^{201}$來和大家一起給錢免禮.
26到30節.
很重要說的是.
聖靈幫助那些人.
有神生命.
聖靈初結果子的人.
是被神修養.
收納.
和得到清義.
和來盼望將來榮耀的基督徒.
這個是保羅.
在這裡來勸勉.
教會的弟兄姊妹羅馬.
要來這樣理解自己的.
得救的工程.
生命的工程.
你對自己有什麼期望呢.
這個是很清楚.
由被清義.
去到將來一個永恆的榮耀.
就是我們每一個基督徒.
必須要一起學習.
聖靈的幫助.
本身耶穌說的時候.
《約翰福音》14章說了很多次.
這個聖靈.
是叫做保衛師.
或者安慰者.
在你身邊隨時幫助你的那一位.
這樣的意思.
我試過一次經歷.
是三秒鐘內.
立即要打999.
很恐怖.
傾盆大雨.
接近黑雨.
駕車.
進入隧道的時候.
沒有雨.
我就稍微關了雨擋.

$^{241}$那個水撥.
差不多出隧道了.
開回那個水撥.
失效.
黑雨底下的失效.
那種懼怕.
你是無法想像.
你的車不敢行.
因為你完全在你面前.
這裡都看不到.
我媽咪和飛龍坐在後面.
在公路的快速公路.
長征隧道出口.
怎麼辦呢.
停在快線不敢行.
看不到任何東西.
停了兩秒.
三秒.
一 二 三.
突然間叮一聲.
為什麼還不報警.
立即拿起電話打999.
都不用三十秒.
就立即有工程車收到order.
立即來拯救我.
帶我去那個中心.
隨時的幫助.
如果你面對火災.
你想一下.
消防員架起雲梯.
見到你.
將你拯救.
這個是隨時的幫助.
是那個時候你最需要的人幫助你.
本來這裡說基督徒.
我們需不需要聖靈幫助呢.
保羅說是很迫切要.
我們敢不敢受到.
我們自己面對人生困難.
我們要有聖靈的幫助呢.

$^{281}$甚至有些狀態.
我們連祈禱都不行.
那個幫助隨時就來.
保羅說這樣說.
最近因為家裡經常煮東西吃.
忘記了.
煮了一隻蛋.
燒焦了.
忘記拿.
記性.
上網淘寶買了個電子鐘.
按了十五分鐘.
真的很開心.
一小時就響了.
但那個鐘不是平安鐘.
你不會按了鐘有人找你.
平安鐘不是一個普通的鐘.
其實是一線通的平安鐘.
平安鐘是什麼用處呢.
就是有可能獨居的老人家.
當他不舒服的時候.
或者不開心的時候.
他一按那裡.
其實是一條熱線.
那個熱線就立刻有人專人來找他.
打給他.
婆婆你有什麼事.
有什麼不開心.
或者你有什麼難處.
我立刻找人來幫你.
這樣叫平安鐘.
所以我家裡的電子鐘.
不是平安鐘.
是告訴你現在時間到了.
我們今天信耶穌的弟兄姊妹.
保羅告訴我們.
我們是有聖靈初結果子.
我們是神修養的人.
其實第八章已經說了.
十五節這個概念.

$^{321}$就是我們被修養.
你們所受的不再是奴隸的心靈.
其實原文是那個靈的字.
不是心.
不是普通的心臟那個心.
是說那個靈.
我們不再是奴隸的心.
仍舊的害怕.
我們所得到所接受的是兒子的心.
因此我們呼叫阿爸父神.
接著第十六節.
聖靈與我們的心同正.
我們就是神的子女.
既然是子女.
就是後製就是修養.
就是神的後製.
神所修養.
因為當時的羅馬帝國.
是有這些當時的社會制度.
是有些有錢的人.
是會修養很多的人.
做他的食客.
這些叫做Kinds.
有個回報.
那些Kinds就是食客.
就要在主人有些重大的需要.
譬如我要做議員.
要投票.
你就要站在我這邊.
所以在羅馬帝國的制度.
有這些修養人制度.
但是保羅所說的.
借用社會的制度來說.
我們信主的人.
我們現在得到的.
不是奴隸的心懼怕.
得到的是兒子的心靈.
不單止是這樣.
聖靈在我們裡面.
與我們共同的印證.

$^{361}$你真是屬於神.
好像吸了個印.
我們是屬於神.
這樣的一種修養.
就是說那份文書.
很確定.
是真的吸了個印.
你是屬於神.
我們剛才已經說了.
修養的概念.
我們是神的後製.
而且是與基督同作後製.
所以如果我們與他一起受苦.
必定會與主耶穌.
一同得到榮耀.
就是由這裡開始說.
那個修養的概念.
接著.
剛才我們有說到的就是萬物.
其實都是勞苦嘆息.
因為我們生活的.
是一個墮落了的世界.
是一個犯罪.
還未真正完全的.
可以在一個完整狀態下的世界.
是進入了當墮落的世界的我們每一個人.
所以我們每一個人都是罪人.
救恩使我們得到神的拯救.
所以第26節到30節.
只是與一類的人說.
就是有聖靈初結果子的基督徒.
我們的人生朝向一個榮耀.
以及在我們的生命裡.
憤外經驗到聖靈工作的人.
只能夠是愛神的人.
本來似乎在這裡規範了一個類別.
就是說.
有聖靈初結果子.
可能我們會有不同的際遇.
會使我們對神的愛.

$^{401}$未必這麼敏銳和熱切.
但這個都可以有一個改變.
因為我們是被神選召.
所以拍望弟兄姊妹.
我們能夠在生活裡看到.
我們是神用祂永恆的愛.
將我們呼召.
將我們拯救.
如果用一些會員概念.
可能大家就會比較容易明白.
今年的4月5日.
在北歐有一間航空公司.
破天荒.
邀請了6000個中心會員.
即是大客的中心會員.
他說我們有一個行程.
這個行程是不會告訴你去哪裡的.
你想想吸引程度就難以想像了.
一定不會送我去KK園.
我是中心會員.
還要是三萬分以上的才能報名參加.
就有6000人報名了.
其實只有150個位置.
可想而知那種雀躍.
有很多技術性.
出機票的時候都知道去哪裡.
你是去日本,韓國還是去歐洲呢?.
連機票都做到手腳都沒有人知道.
接著問一下航空公司的職員,地勤,空姐.
對不起,我都不知道.
誰知道呢?原來只有一個人知道.
就是那個開飛機的機師.
開飛機之後在北歐.
兩個小時後就廣播了.
各位乘客,這個盲游呢.
其實現在我們去的地方就是雅典了.
雀躍到整個飛機都瘋了.
當然那個旅程是非常優質.
非常預先計劃好那些行程.
是真的要超級會員才能享受到的待遇.

$^{441}$穿什麼衣服,住哪間酒店.
他都幫你安排好.
這樣的一個盲游.
親愛的兄弟姐妹.
如果你是其中一個被選上的.
我相信你都很雀躍.
你一定不會擔心死了.
抓住老公會不會送到KK園呢?.
接著把我的器官剝了呢?.
一定不會的.
因為你知道那個航空公司是很忠實的.
還有這次的行程不是開玩笑的.
是真的有一個很高的優質的待遇.
各位兄弟姐妹.
很少我們會想到我們屬於神.
看到神這麼偉大的愛.
將祂的愛子這樣犧牲.
我們會感覺很大的雀躍.
我們能夠被輕易.
不再追究你任何的罪.
不再看你是有罪的人.
有時候回教會.
都很不容易的.
因為如果你跟一些團體有些少的摩擦.
見到面大家有時候都好像有條刺一樣.
但是神看每一個罪人的悔改.
不再追究任何的罪.
因為祂已經稱我們為義.
其實後面第八章.
保羅繼續說.
有誰可以來控訴神所揀選的人呢?.
由基督耶穌已經死了.
而且在天上為我們祈禱.
沒有人可以再控告神所揀選的人.
包括神自己.
所以這裡告訴我們.
我們是特選的.
在地球上被神揀選的人.
只是基於愛這個大的原則.
是神的愛將我們拯救.

$^{481}$我們成為了真是神的收納.
收養的兒女.
我們應該要明白一下.
這位造物主為什麼要這樣愛我們.
剛剛最近看到一個挺感人的故事.
一個報導.
有一位他的英文就是Jenna Cook.
她就在耶魯大學讀三年班的時候.
她做一個project.
她聯絡了有四個和她一樣背景的女孩.
回去中國湖南.
來做一個計劃.
就是去孤兒院做義工.
當然就是要做一些研究.
為什麼要做這樣的project呢?.
有經費的.
後來拍了一部電影.
拿獎了.
原來這個Jenna.
她很小的時候.
沒有意識地.
就被一個美國的家庭收養了.
她開始懂事了.
她的媽媽.
收養的媽媽.
美國人.
盡力幫助她對中國的文化.
對中國的事情了解認識.
但是在她成長的那個地方.
那個小鎮.
是沒有中國人.
她的成長.
很大的一個conflict.
因為她不明白.
為什麼.
她是被遺棄.
這個是她內心很深的刺.
所以她要做到最好.
讀書要最好.
來彌補她.

$^{521}$可能是不是小時候.
她哭得很厲害.
很頑皮.
被自己的父母.
這麼殘忍的遺棄.
一路長大.
能夠在那個小鎮.
有的很少的中國人.
一路耿耿於懷.
很想回到中國.
到底我的根在哪裡呢?.
後來去到耶魯大學.
第三年班的時候.
終於因為她讀的是一個人類的研究.
那個人類學.
和女性主義那些課題.
她的主修.
她真的拿了獎學金去做這個project.
她回去尋根.
登報紙.
去回那間孤兒院.
因為她的養母知道.
回去那個孤兒院裡面去尋根.
哇 登報紙之後.
當時的武漢的報紙.
幫她登這個故事.
最終篩選了有50個家庭.
仍是她的父母.
來到去相認.
輾轉輾轉這樣篩選.
最後最好的辦法.
最科學的辦法.
就是做DNA的鑑驗.
很可惜50個家庭.
沒有一個是她的親生父母.
但原來在這個過程裡面.
她感受到原來一孩政策.
表面看文字就很殘忍.
但實際上原來很多的處境.
她真的去做一個研究.

$^{561}$去跟這些家庭做訪問.
和看到他們過去這麼多年.
捨而不捨的要尋找自己.
失散的兒女的時候.
那種其實不是兒女.
全部是女孩子.
大家知道一孩政策.
所衍生的問題.
就是你生女兒.
快點扔掉再生一個.
因為要報人口.
報那個戶籍.
所以在80年代就出現過.
這樣的情況.
當然現在越生越好.
因為人口減少.
但政策有時帶來很多家庭的悲劇.
珍娜就一直這樣訪問的時候.
不知不覺雖然她找不到自己的父母.
但她對自己的身份認同.
又強了很多.
她終於敢承認自己既不是美國人.
雖然她也不是中國人.
但她能夠在她自己現在這樣的發言裡面.
這種尋根之旅裡面.
她能夠認識自己原來有個位份.
就是幫助很多還未能夠跳出.
這樣的身份危機的中國女孩.
被美國人收養.
有總共八萬多人.
她自己的身份不能夠認同.
說的字是滿口的英文.
但自己原本就是一個中國人.
但不知自己的親生父母家庭在哪裡.
什麼原因要遺棄自己.
她就在這個project裡面幫助很多人.
找回自己的身份.
親愛的弟兄姊妹.
我們的根當然就是在我們的出生.
我們的地上.

$^{601}$但這段聖經告訴我們.
我們是屬於神揀選.
神愛我們.
預定我們.
得到祂的榮耀.
所以聖靈的幫助.
就是讓我們不斷在困難.
在我們的逆境.
不放棄.
我們認同自己的身份.
是屬天的身份.
很重要就是聖靈的介入幫助.
所以這裡告訴我們.
一切的困難.
是不會使我們和神的愛隔絕.
這個是由下面的31節開始.
一直講到39節.
是保羅來排除任何的事情.
沒有東西可以阻礙我們.
屬神屬於主耶穌的人.
在人生裡面的成長.
困難都不會打倒我們.
我提過有位在三藩市.
遇到嚴重車禍的姐妹.
她現在差不多每天.
有時怕我忙都不敢找我.
問我一些神學問題.
有一天她突然問我一個問題.
我很害怕.
你贊不贊成安樂死.
一個嚴重傷殘.
車禍令她殘缺的一個女孩.
一個姐妹.
她問我這個問題.
她說她家人不贊成.
問我的意見.
我當然不贊成.
我當然覺得她不應該這樣想.
原來她只是想問.
知道我怕她會做傻事.

$^{641}$她說曾牧師你不用擔心.
耶穌在我的心裡.
我不會.
我知道自己去天上.
祂不會做傻事.
信主不是很久.
是聽勵善牧師.
我的查經班.
他現在每個星期.
聽錦繡堂 洪恩堂的講道.
他說如果我回不了教會.
神宗不喜歡的.
我說你不是回不了教會.
你暫時要人陪.
要人推你才可以回教會.
你自己是不能回教會.
他問這些問題.
和神的關係.
我應該怎樣可以在我現在的狀態.
去維持.
他說你不用擔心.
我不會安樂死.
他知道自己現在存在的價值.
弟兄姊妹.
這段聖經告訴我們.
我們由清二到得榮耀.
都是神選擇我們帶來的後果.
所以我們要開展我們的沉根之旅.
我們有沒有感受到聖靈.
在我們軟弱的內心的幫助呢.
我們一起禱告.
去研讀這段經文.
很盼望對於我們自己的.
或者景福時代.
我們的天路之旅.
能夠不斷對主的救恩和大愛有認識.
也使我們能夠在人生裡面.
我們朝向一個.
每一次都是一個朝聖之旅.
回到禮拜堂之後.

$^{681}$我們所得到的教導,牧養.
是給我們有力去面對這個世界.
和將福音有效.
有力地去告訴人.
邀請人進入救恩的筵席.
這樣感恩禱告.
奉耶穌基督的名.
阿們.
\newpage



\section{腓立比書 2:17-30}
\label{sec:w0vSwuU8gCk}
\textbf{2024-04-28 早堂《同心事主的知己》腓立比書 2\_17-30  講員:馮耀榮牧師}
\newline
\newline
連結: \href{https://youtube.com/watch?v=w0vSwuU8gCk}{\texttt{https://youtube.com/watch?v=w0vSwuU8gCk}} ~~~~ 語音日期: 2024-04-30
\newline
\newline
\hyperref[sec:a_rRhGfuxUw]{\small{< < < PREV SERMON < < <}}
~
\hyperref[sec:index_chronic]{\small{[返順時目]}}
~
\hyperref[sec:index_scriptual]{\small{[返順卷目]}}
~
\hyperref[sec:jRV45ettAd0]{\small{> > > NEXT SERMON > > >}}
\newline
\newline
腓立比書 2:17-30
\newline
\begin{longtable}{cl}
\hline
\hline
章節 & 經文 (和合本修訂版)\\
\hline
2:17 & \begin{tabularx}{0.7\textwidth}{X} 我以你們的信心為供獻的祭物，我若被澆獻在其上也是喜樂，並且與你們眾人一同喜樂。 \end{tabularx} \\ \\ \relax
2:18 & \begin{tabularx}{0.7\textwidth}{X} 你們也要照樣喜樂，並且與我一同喜樂。 \end{tabularx} \\ \\ \relax
2:19 & \begin{tabularx}{0.7\textwidth}{X} 我靠主耶穌希望很快能差提摩太去見你們，好讓我知道你們的事而心裡得著安慰。 \end{tabularx} \\ \\ \relax
2:20 & \begin{tabularx}{0.7\textwidth}{X} 因為我沒有別人與我同心，真正關懷你們的事。 \end{tabularx} \\ \\ \relax
2:21 & \begin{tabularx}{0.7\textwidth}{X} 其他的人都求自己的事，並不求耶穌基督的事。 \end{tabularx} \\ \\ \relax
2:22 & \begin{tabularx}{0.7\textwidth}{X} 但你們知道提摩太是經得起考驗的，他與我為了福音一同服侍，待我像兒子待父親一樣。 \end{tabularx} \\ \\ \relax
2:23 & \begin{tabularx}{0.7\textwidth}{X} 所以，我一看出我的事怎樣了結，我希望立刻差他去， \end{tabularx} \\ \\ \relax
2:24 & \begin{tabularx}{0.7\textwidth}{X} 但我靠著主自信我不久也會去。 \end{tabularx} \\ \\ \relax
2:25 & \begin{tabularx}{0.7\textwidth}{X} 然而，我想必須差以巴弗提到你們那裡去。他是我的弟兄、同工和戰友，是你們差遣來供應我需要的。 \end{tabularx} \\ \\ \relax
2:26 & \begin{tabularx}{0.7\textwidth}{X} 他很想念你們眾人，並且極其難過，因為你們聽見他病了。 \end{tabularx} \\ \\ \relax
2:27 & \begin{tabularx}{0.7\textwidth}{X} 他真的生病了，幾乎要死。然而神憐憫他，不但憐憫他，也憐憫我，免得我憂上加憂。 \end{tabularx} \\ \\ \relax
2:28 & \begin{tabularx}{0.7\textwidth}{X} 所以，我更要盡快送他回去，好讓你們再見到他而喜樂，我也可以減少憂愁。 \end{tabularx} \\ \\ \relax
2:29 & \begin{tabularx}{0.7\textwidth}{X} 故此，你們要在主裡歡歡喜喜地接待他，而且要尊重這樣的人， \end{tabularx} \\ \\ \relax
2:30 & \begin{tabularx}{0.7\textwidth}{X} 因他為做基督的工作不顧性命，幾乎至死，為要補足你們供應我不夠的地方。 \end{tabularx} \\ \\
[1ex]
\hline
\hline
\end{longtable}
$^{1}$各位早晨.
今日同大家分享一段經文.
《肥立比書》第二章.
裡面講到保羅有兩個同伴.
一個是提摩泰 一個是伊巴忽提.
保羅和這兩個人的關係.
有一次經文表達.
在我個人來說有一個反思.
所以在這段時間和大家分享一下.
這段經文其實都很清楚記載保羅晚年的時候.
他的心情和提摩泰和伊巴忽提之間的關係.
對不起 把經文記得小一點.
在第20節裡面提到提摩泰也好.
或者伊巴忽提這兩個同工.
他表達沒有哪一個人與我同心.
好像這兩個人一樣.
保羅在人生裡面有很多同工.
或者有很多弟兄姊妹在他周圍.
他年老的時候曾經說過這句話.
其實在我人生裡面 在這個階段裡面.
只有兩個人與我同心 是我的知己.
這個就是今天和大家分享一下心境.
我們回來教會一段長時間.
究竟再反省一下.
我們接觸這麼多弟兄姊妹或者朋友的時候.
有哪些人是我們的知己或者我們同心的人.
所以今天和大家分享心路歷程.
其實這段經文提到同心.
在呂鎮忠的譯本裡面不算太清楚.
他把同心這個詞翻譯成知己.
我沒有知己 沒有其他人.
人好像這兩個人 好像是知己一樣.
這個翻譯本身在希臘文裡面.
在新約只有出現過一次這個字.
最好的翻譯應該都是和本或者和修本的同心.
或者是同一個掛念的事情.
呂鎮忠他翻譯了我們中文裡面.
比普通朋友更加好的一個用詞 叫知己.
我太太最近都問我.
她見我比較進入了退休年齡.

$^{41}$我兩年之前退了休.
但有一段時間我都做兼職在神話院幫忙.
我們開始進入另一個人生階段.
沒有公事上和人家的接觸.
突然間多了一些時間.
我不需要在公事上和同事或者同學相處.
以前我是需要有關心小組.
所以這個同學或者這班同學.
我一個月有兩三次要一起小組.
或者單獨面對他.
現在退了休之後就可以不用做.
究竟我平時有什麼時間找朋友或者找人一起去相處.
我退休之後我太太就問我一個問題.
究竟我有沒有朋友?.
通常太太都會問丈夫這個問題.
特別是年紀大了的丈夫.
多數比較孤獨.
我是一個很典型的內向的人.
如果有工作的話.
社交能力都可以的.
都會約人聊天.
但一退了休.
其實就沒有那麼大興趣.
不是沒有時間 時間是有的.
但要我約人見面聊天.
不如我在家裡坐一會.
休息一下 看一看書 聽聽音樂.
所以可能漸漸就會更加內斂 退縮.
對我來說其實都不是太大問題.
我一個人 特別我們是一種教書型的老師.
我不是一個牧養型的老師.
所以平時坐在這裡上課都是很正常的事情.
所以退了休是多一點時間上課.
少一點應酬.
所以當這段經文說到保羅.
他沒有什麼知己的時候.
我有點共鳴.
其實我朋友都不多.
所以更加難說有很多的知己.
特別到了這個年齡.

$^{81}$是不需要公事上或者職責上去接觸太多的人.
人生可能要找到朋友聊聊天.
其實都不是很多.
所以當我看到保羅都說他沒有什麼知己.
他只有兩個知己.
一個就是提摩泰 一個是爾伯伏提.
我覺得原來人是比較孤獨.
這個都不是個性的問題.
一個都是正常人.
特別到了一個年齡.
沒有了職責上要和外面的世界相處.
是有這樣的現象.
當然這個在我們想像當中不是很正常.
我們回到教會應該多一點朋友.
要多一點相處.
但既然保羅其實都不是很多知己.
如果我們像保羅 我們學校保羅.
其實應該都不是很大問題.
當然我再想一想.
我有沒有知己 我有沒有朋友.
我最後想到一個.
其實我有一個同班同學.
由畢業到現在.
我記得我們一起在神樂園讀書.
那時候二十多歲.
他是一個韓國的華僑.
我們讀書的時候.
他帶我去請過他去韓國的家.
韓國那時候還是很窮.
我記得很有印象.
就是一去到那裡天氣又比較冷.
所以難不倒.
頭暈暈 出不了街.
直到他帶我去一些洗澡的地方.
韓國好像北面都有些公共浴場.
洗完澡之後.
熱了你 然後冷了你.
身體會好一點.
最深刻印象就是.
沒有去過那裡.

$^{121}$大概一個星期.
當時他請我去那裡參觀教會.
也有一些茶經我負責.
這個是幾十年前的印象.
然後我們畢業後他回到韓國事奉.
最近這十多年回來香港事奉.
我們去到近五年.
他退休了 我又退休.
我們一個月 或者一年有兩次約出來吃飯聊天.
他有個小朋友也是健道的校友.
也是我的學生.
所以我再想一下真的只有一個.
如果要找出來聊天吃飯.
還是他打電話找我.
我都沒有特別有空打電話找他.
所以我不是很多朋友.
所以還好這段經文.
呂振中譯了這個童心是知己.
在上文下理.
他其實只有提摩泰和爾巴忽提.
這兩個人是他的知己.
我覺得少一點朋友少一點知己.
似乎好像是保羅的現象.
我們在座當中有些可能進入這個年齡.
當然如果我們說我們認識的人是相當多的.
可以聊天的人其實都不是說聊不到的.
不過年紀大了會對那個所謂朋友和知己的要求高了很多.
他明不明白你 你明不明白他.
人生久了 要求高了.
朋友就少了 知己就差不多沒有.
這個是這個現象.
我今天說的主要原因就是.
因為我們其實好像我們這個年齡退了下來.
其實還有十年二十年或者三十年的生命.
我們怎樣面對沒有職責的人.
他怎樣跟外面的世界聯繫.
或者跟我們教育電影節目聯繫.
我有很多校友很多同學.
有很多教育電影節目.
我有幾間教會做顧問.

$^{161}$所以要認識朋友不是沒有機會的.
也都不是沒有時間.
問題就是怎樣去把握跟自己跟別人的關係.
所以今天我想在這段經文裡面.
這個呂鎮忠的譯本去說知己的字.
我覺得這個譯本是給我們信徒.
我們雖然可能在作當中有些不是那麼老.
不過你有父母比較老一點.
所以你開始要了解一下他的想法.
我們現在這個老化了的世界.
教會也有很多老化了的弟兄姊妹.
或者一些弟兄姊妹的父母.
所以我們可以給他們一些想法和思想.
說到知己這個呂鎮忠.
他大概十幾年之前去世.
他去世的時候.
他透過一些親人來問過見度.
把他的參考資料想捐給見度.
但當時見度的圖書館已經不夠位.
我們要把舊書或者重複的書拿走.
然後再問.
我們新書上不到架.
所以最後沒有接收到他剩下來的書.
我其實很有興趣.
他究竟藏了些什麼書.
因為呂鎮忠這個人我沒有見過.
不過在我教書的過程.
無論是新約還是舊約.
我都很經常去參考他的翻譯.
所以嚴格來說.
我會覺得如果說到知己.
他是其中一個我會有興趣約他出來聊聊天.
究竟他是怎樣去思考聖經.
這一節是經文.
他這個翻譯.
我有一個很想問他一個釋經的原則.
因為這個字只出現過一次.
而大部分的中文或者英文的翻譯本.
都沒有他這樣翻譯成知己.
大概都是同一個掛念的心.

$^{201}$所以他怎會翻譯成一個.
沒有什麼參考作用.
或者沒有其他地方的作用.
所以今天講到其實是一個.
牽涉到解經的問題.
我比較有興趣的東西.
因為聖經的翻譯是我其中一個.
比較有興趣的課題.
所以所謂和自己有一個同一個興趣.
或者想做的事情.
想了解 想進一步去推敲的.
其實能夠有的人.
你能夠去和他聊天就不多.
曾實華牧師都是在建道一段很長時間.
張美眉牧師都是我們同事一段時間.
在行政上 在教學上.
不過他們是負責猜拳的.
所以我大概都不會問他們.
究竟這個翻譯是怎樣.
因為我是聖經科老師.
他都不會問我他教的跨文化的事情.
所以我們可能辦公室就在樓上樓下.
或者在隔離.
每個下午吃飯都有機會坐在那裡吃飯.
不過大概都不會講到這個翻譯的問題.
這個呂振中這裡的翻譯.
是不是最準確 最合適的.
所以其實不是很多人問的.
因為每一個人的興趣是很不同的.
所以我有的知己都不是在我旁邊裡面活著的人.
他可能已經過世了.
所以我都很欣賞他.
因為他是一個人翻譯了所有的舊約和新約.
這個其實是很高難度的.
能夠將整個聖經翻譯的一個人做.
不要說中國人 就算是外國人.
我們的King James 本 和學本背後的譯本.
全部都是一群學者去翻譯的.
很少一個人可以將由創世記到啟示錄去翻譯.
所以呂振中是一個.

$^{241}$我想如果我說我孤獨 他更加孤獨.
我想他都不會答應出去講道.
因為他應該都是喜歡坐在那裡.
在書房裡面用了他整個人的人生.
去將聖經裡面每一段的經文去翻譯.
這個要很有耐性 很有孤獨感.
他大概都不需要朋友 亦都不需要知己.
他才有時間去專心去看聖經 然後再去翻譯.
我都未去到這樣 我有很多職責上.
要做很多人與人之間的接觸.
所以我能夠比較翻譯的聖經.
是一兩卷書的.
不可能是這麼長的時間.
所以呂振中他是一個很天才.
在翻譯方面.
我們當然中國人無論是在著色書.
各方面都是剛剛開始起步.
但我們有一個在五十多年之前.
已經為中國人教會.
貢獻了一個相當之好的中文的翻譯聖經.
我間中都有幫忙一些新譯本的翻譯.
去校對其中一兩卷書.
都要花很多時間.
所以呂振中如果講知己.
所謂知道自己.
就是自己的興趣與對方的恩賜興趣是夾在一起.
你才覺得他可以與你溝通.
你想再見他.
但保羅他兩個朋友兩個知己.
與我的知己的判斷是有少許不同的.
首先他第一個知己.
與他的年齡差了一代.
提摩泰是他的兒子.
我們有多少人會將我們的下一代.
成為朋友成為知己呢.
多數都不會的.
我們可以聊天.
不過要說服對方理解自己是很困難的事情.
雖然我們都是做傳道人.
對不同的年齡都需要去溝通.

$^{281}$但到了自己的子女.
其實可以說的話.
那個障礙都頗大.
他都知道你想說什麼.
不過從你的口說.
他多數不聽.
我們對他的要求.
是比普通大姐妹的要求高.
所以特別是跨了代的人.
可不可以做知己做朋友呢.
其實有一定的難處.
當然我們理論上.
都應該將我們的家人.
變成知己朋友.
但實際上是不容易的.
所以你問我和我太太的關係.
她是不是我的知己呢.
理論上應該是的.
不過年紀大了之後.
她和你說的話是相差很遠的.
這是很典型的.
首先我太太應該是很典型的一個姐妹.
她一說話是從頭開始說到最後.
她好像很簡單.
她是做阮木的.
她前兩天和支持阮木的教會長執開會.
談了一談.
她回來後會將她談的事情.
由頭開始.
即是說約她在什麼時候見面.
到談了什麼.
結果是什麼.
她一直很詳細地告訴我.
其實她之前告訴我.
她有見面的.
她告訴我究竟她是否支持這個士工.
她不需要說約了在哪裡吃飯.
跟著是誰遲了來.
我因為要有和諧的婚姻家庭.
所以我要忍一忍.

$^{321}$慢慢聽她.
如果不聽她的說話.
她應該會不開心.
我說話其實沒什麼可以說.
學院所有的事物.
我都不可以說給她聽.
因為所有重要的事物都是人事和學的問題.
所以沒什麼可以說.
我和太太的關係相當好.
不過我又不能說話給她聽.
她說了話給我聽.
我要忍一忍.
忍到她說完為止.
這是很典型的男女不同的地方.
究竟她是不是我的知己.
她和我受同一個神學訓練.
同一個地方.
我們都是教會事奉者.
當然她是院務.
我是神學院的老師.
但是我們能夠理解對方的機會應該多了很多.
但是如果說到朋友或者知己.
倒不如說我們只是夫婦的關係.
或者同事奉的關係.
所以興趣很不同.
我們最近有一個問題要處理.
去旅行應該去哪裡.
我去旅行當然不是每個人都喜歡去的地方.
我喜歡去的地方不是很多.
不是很多在中國的地方,城市.
或者去日本.
我大概去哪個地方.
過一兩年我都會再去一次.
所以我太太去過的地方就不想再去.
所以我們下一次旅行就是我一個人旅行.
再等她去旅行的時候.
要去一個地方是她沒有去過的.
我們以前上班.
大家都上班.
所以大概幾年才有一次可以合得來的時間.

$^{361}$可以去一個比較長的旅行.
我們兩個現在開始進入.
我退休了.
她過兩年就退休了.
我也有些煩惱.
究竟要去哪裡.
因為我想去的地方和她的地方又很不同.
所以大家成長受到同一個訓練.
相處了幾十年.
但是距離是越來越遠的.
所以不敢說叫做知己.
通常朋友你容易找到去旅行.
因為你喜歡去的地方的人.
你就喜歡去.
我們不是很合得來的.
因為我們是夫婦的關係.
我們是比較親密的關係.
如果距離這麼大.
好像我們需要再重新思索一下.
不過為什麼我說我這次講到.
其實是對我有幫助的.
就是我想到保羅.
其實他的知己都不是我想像中的知己.
因為他和提摩太的距離相差很遠.
他講到提摩太的時候.
其實提摩太沒有寫過書信.
我們知道聖經有提摩太前書和後書.
是保羅寫的.
所以提摩太保羅寫了十幾卷書卷.
其中兩卷提摩太前書和後書.
即是提摩太不可以繼承保羅.
有一種聖經著作的能力.
提摩太我們教會裡面.
通常年輕的團契都叫做提摩太.
如果是年長一點的團契叫做保羅團契.
所以是隔了兩代的.
如果你是有年長的年齡.
你都不會去提摩太因為相差很遠.
或者提摩太的人不會去保羅團契.
所以提摩太和保羅在年齡上.

$^{401}$其實是差了一代.
他怎樣可以成為朋友.
怎樣可以成為知己.
這個就是保羅他有兩個知己.
但這兩個知己的選擇.
他為知己.
不是因為他和差不多的經歷,年齡.
或者興趣.
所以保羅當提摩太是知己.
都是一個很大的想像.
我在建道教書的時候.
由年輕到年長.
我是最近這十年.
開始和我的學生的距離相差很遠.
我可以即時和他們相處.
但不可以當他們是朋友.
因為大家的興趣很不同.
幸好我就退休了.
所以我感覺上.
我最後一組的學生.
長大了四年.
他好像我的小朋友一樣.
比我的小朋友還小.
所以我們開組.
大概都是他們聊天.
或者在旁邊聽.
因為如果我一出聲.
我主導的內容.
他們就不聊了.
他們聊些什麼呢.
他們聊些打機的東西.
當然上課一定是講聖經的.
教會聚會.
但如果是私下聊聊天.
他們很喜歡打機.
特別是女生都很喜歡打機.
這是一個很獨特的文化.
所以到最後我感覺.
如果是同樣的興趣.
知己其實就不太好.

$^{441}$所以保羅和提摩太.
其實是隔了一代.
他都可以將和自己不同年齡的人.
都當為知己.
又不用他的恩賜興趣.
作為我的選擇.
或者認知他是我的好朋友.
好知己.
這是我今天訊息裡面.
對我自己有一個反思.
如果我們再看另一個人.
應該大家不算很熟.
大家都沒有什麼團契的名字.
叫做爾巴弗提.
沒有.
這是保羅另一個知己.
而在這段經文裡面.
講得最多的是爾巴弗提.
反而提摩太很簡單地講一講.
在經文裡面.
保羅和爾巴弗提是一個知己.
有些獨特的.
因為這個爾巴弗提.
比較猜他去關心保羅.
當時保羅坐牢.
或者有身體上各方面的需要.
他猜了一個比較年輕的人去幫他.
不過這個人的身體有病.
所以他很快要回家.
保羅就找人送他回家.
所以保羅對爾巴弗提.
有一個很好的評價.
他說他是為了基督的工作.
不顧惜性命去幫助保羅.
不過他本身自己病毒快要死.
這個就是爾巴弗提.
爾巴弗提最後一節經文裡面說.
他不求自己的事情.
只有求耶穌基督的事情.
這一點和提摩泰有點相似.

$^{481}$和保羅和爾巴弗提.
這三個人都有些相似.
他們是為了基督的事情.
他們將生命放上.
他們不求自己的事情.
這個就有些獨特.
不求自己的事情.
剛才我所說的知己.
都是找一些朋友.
和自己的興趣一樣.
個性一樣.
是為了自己的世界觀.
去凝聚周圍我可以接觸的朋友.
或者知己.
但是保羅說提摩泰也好.
爾巴弗提也好.
他們兩個在年齡上和保羅有距離.
在能力上也有很大的距離.
保羅寫了很多聖經的經卷.
而提摩泰呢.
提摩泰的前後書不是提摩泰寫的.
所以他繼承不了他的職責.
今年是我最後一年在建道可以.
參加畢業禮.
正式來說.
今年有一個同學.
他完成了他的博士論文.
我和他相處了十年的時間.
因為他有一部分時間去寫作論文.
他寫的是撒慕爾記的內容.
我本身二十多年之前進入博士論文.
本來也想寫撒慕爾記的.
最後是完成不了.
我就轉了另一個經卷.
寫了創世記中的《約瑟的故事》.
到這個同學寫的時候.
我做不了的東西.
希望他做得出來.
他今年畢業了.
他通過了論文士.

$^{521}$所以這次的畢業禮是很有意義的.
因為我做不了的東西.
他做得出來.
他是我的知己嗎.
他是一個女同學.
不過起碼如果我想說聖經的內容.
和他聊應該比較有共鳴.
因為大家都上了經卷相當長的日子.
提摩太太其實她沒有像我的學生一樣.
繼承到保羅 超越保羅.
她在能力上沒有保羅那麼有才幹.
通常我們找朋友和知己.
都是和自己的能力或成就差不多.
起碼你會覺得你可以和他聊到.
或者可以讀到.
但保羅都不是找一個和他差不多的威望的人.
起碼他和彼得讀書不多.
不過彼得在十二個使徒當中是很出色的.
但保羅和彼得都有爭吵過.
大家都有個性.
所以才幹兩個人在歷史上.
彼得雖然書卷少一點.
但天主教都用教皇.
都不用保羅 用彼得.
因為彼得他的個性.
不是他的訓練是否如敷的問題.
彼得的個性在當時的門徒當中.
都是頗有才幹 頗有領導能力.
頗有動見.
和耶穌的關係.
他是最踴躍和耶穌聊天的.
所以保羅 耶穌都當他和另外兩個門徒是他的知己.
有三個的知己.
耶穌有什麼事情都找他們聊天.
所以保羅有兩個知己.
耶穌有三個.
所以從某個角度.
我們不需要很多知己.
不過一定要有一點.
否則我們又會有一定的孤單.

$^{561}$所以我印象這裡的保羅.
他選擇的知己.
和我選擇的知己的想像是不同的.
他不是說一個人.
我自己有這個才幹 有這個成就.
或者我有這個興趣.
我要找一些人和我有差不多的興趣.
差不多的才幹 或者差不多的成就.
他是找提摩太 找爾巴弗提.
爾巴弗提我們是不清楚是怎樣.
但是聖經在這段經文有提過.
這兩個人是不求自己的事情.
他是求耶穌基督的事情.
這一點就是我們所有回來教會的弟兄姊妹.
回來崇拜之後.
你會不會和一些人一起去吃飯.
你會和哪些人去吃飯.
通常都很明顯的.
師班的人多數找師班的人一起吃飯.
因為大家一起唱詩很久.
如果有這個團體就找這個團體.
這個是正常的.
我們的教會都有少少聯誼的性質.
你的工作 如果你是做醫護的.
你和相關的服事行業的人.
可能會約出來吃飯 聊聊天.
你是老師的 你是做生意的.
你可以接觸所有人.
你會多些和你的行業附近的人.
可以一起聊天 或者回教會.
或者去茶經 甚至是團契.
你都可能和你的背景很相似的人.
這個不是壞事情 也不是不正常的事情.
不過到了保羅去選擇朋友和知己的時候.
他的選擇就不是用他的背景.
而是用他的興趣.
他的選擇是說這些人好像我一樣.
是為了基督的事情將人生放上.
所以不求自己的事情.
這些就是知道我的知己.

$^{601}$不知道大家明白嗎.
保羅的知己就不是我想像中.
當我年紀老的時候.
我再想想哪些人和我一起.
一樣有興趣去茶經.
一樣有興趣去寫書.
其實不是這樣的.
回來教會的時候.
如果你覺得孤單.
你覺得沒有人和你的興趣一樣.
是真實的.
年紀大了之後.
就算同一個行業的人.
你可能都不是那麼賞心的.
但是保羅給我們一個很重要的榜樣.
他選擇知己.
不是因為他和他的年齡差不多.
或者他的才幹結果是怎樣.
其實這個巴菲提是相當之弱的.
本來想幫人.
後來要被人幫.
保羅還擔心他.
所以我有些學生我要擔心的.
有些學生不用擔心.
有些學生擔心他寫不到功課.
畢不了業.
有些學生你不用擔心.
因為他比你厲害.
我記得有些學生畢業後回來見教書.
有印象他入神學院的時候.
我已經覺得這個人不用特別關心他.
因為他都有能力去處理自己的人生.
所以要付出的.
要和他相處的.
其實我們作為一個牧者.
或者作為一個信了主先的人.
其實下一代的人.
他們的才幹也都有比自己強的.
我有三個小朋友.
有一個是我比較擔心要幫他.

$^{641}$有一個他就擔心要幫幫我.
他經常想著何時我退休的時候.
他要考慮一下買樓的時候.
要和他一起住.
其實我就不是很有興趣.
做朋友都不要經常黏著.
黏著其實很難做到朋友.
因為他會認真理解你.
你又認真理解他.
相處好和處難.
我喜歡一個人.
我最多和我太太在一起.
我太太都不是特別想帶孩子.
所以她的優點就是我和她一起.
各有各的興趣.
所以保羅他選擇朋友知己.
他選擇一個我們比較容易有共通的目標或者內心.
菲拉比書其實由開始到結束.
就好像他對伊巴忽提和保羅說.
這些人在人生裡面不求自己的事情.
他是求耶穌基督的事情.
這個就成為我們的知己.
所以嚴格來說.
所有神學生進來見道.
他們都是為了基督耶穌的原因去服上.
剛剛我們有一個培靈的音樂會.
都有好幾十個人年輕人.
去考慮全職事奉.
其實他們肯放下工作進入全職事奉.
為了基督的原因做事情.
這個其實可以成為我的知己.
因為這個也是保羅的榜樣.
放下自己以基督耶穌的心為心.
以基督耶穌的福音為念的人.
這些人就是我的知己.
所以我就不需要再找哪個朋友和我一樣有興趣.
做一些事情.
或者會去旅行的時候去哪裡.
去吃些什麼.
一起去哪一間去吃好一點.

$^{681}$其實我又沒有什麼興趣吃東西.
所以如果說我們屬靈的知己.
其實每一個都可以成為.
只要你願意回來教會.
將基督的福音為念.
就是我們的知己.
就是自己的知己.
你就不會覺得太孤單.
有時候我覺得很多弟兄姊妹回教會幾十年.
回來之後他們都是用世上的習慣.
原來回來都沒有什麼人和我聊天.
沒有什麼人和我一樣的興趣.
所以回到家裡變得只有我和神的關係.
我和其他人只是一種疏離的關係.
這種疏離的關係.
因為大家的尋找知己都用世上的準則.
就是和我的年齡差不多.
差不多興趣.
差不多成就.
或者差不多的社會經驗.
我們要超越這種選擇.
我們好像保羅.
他是一個很屬靈的人.
在《肥勒比書》裡面.
其實最屬靈的就是世上沒有什麼需要知己的.
他說我以基督耶穌的心為心.
我有耶穌基督的時候.
就是我的一切.
其實這個也是我們的問題.
究竟我們怎樣面對人的中腦.
對我來說.
其實最重要就是我能夠接受自己.
還有耶穌基督在我心裡面.
我印象人走到一個階段.
特別是面對離開這個世界的時候.
你只有神是你唯一的朋友.
所以有時候回教會有些人越來越孤獨.
就因為某個角度.
他的屬靈是越來越高超的.
只有基督是我的滿足.

$^{721}$所以我不需要朋友.
也不需要告訴其他人很多東西.
我坐在這裡都可以跟神有相處.
如果這種準則是對是錯.
保羅就是這樣的人.
保羅他對人生以基督耶穌的心為心.
他就是他的珍寶.
他可以將地上所有東西都變成糞土.
所以一個屬靈的人.
其實是可以不需要自己.
甚至不需要朋友.
配偶是不錯的.
不過其實也是一個負擔.
所以保羅選擇不結婚.
這個也是有一個問題.
保羅他是最屬靈的人.
所以我在世上沒有人可以跟我一樣.
跟神這麼親密.
我以基督耶穌的心為心.
當這段經文說.
我好像一個獻祭.
我的生命就是一個囂電的祭.
我將我整個人生燃燒在神的祭壇面前.
所以我的生命只是為了神.
這是保羅的生命.
所以其實他不需要太多人知道他.
理解他.
因為只有神理解他就行了.
這個其實也是我的人生.
我覺得我自己孤獨的其中一個原因.
其實我覺得要討好人不如討好神.
要遷就人不如遷就神.
遷就人大部分時間我都沒有什麼問題.
因為我是一個比較遷就人的人.
但我都有個性.
所以我要遷就我太太的時候.
比我遷就其他人要很多倍的力量.
因為你要遷就一個跟你同住的人.
但我們作為神的人都要做一個見證.
所以我都要有耐性對我的親人.

$^{761}$自己的兒女或自己的伴侶.
所以有很熟的人在一起.
其實也是一個壓力.
是一個重擔.
所以需要很多朋友很多知己嗎.
其實不需要.
但保羅是兩個世界的極端.
他可以跟神在一起.
但同時他當提摩泰,巴伐提這些人是他的知己.
所以我的心要怎麼遊走.
就是我需不需要有朋友.
還是只有耶穌是我的生命最關的最重要.
但當我們有朋友的時候.
我們會不會用我的興趣成為我的焦點.
知己本身就是一個很自私的表達.
就是要自己的配偶.
要自己的弟兄姊妹去知道自己.
我的生活是不需要別人知道我應該想做什麼.
因為這樣我會很霸道.
所以知己這個字其實是很霸道的.
要知道自己.
所以我會贊成呂鎮忠的翻譯有一點錯.
在上文下理裡面應該都是同心多於知己.
同心的意思是我們是不是有同一個心思.
為了耶穌基督的事情.
我跟神同心.
提摩泰跟神同心.
提摩泰知不知道我.
我知不知道他.
我需不需要他明白我呢.
其實是不需要的.
沒有這麼霸道的.
你不可以叫你丈夫或太太明白你.
就算他明白我也好.
他也忍不了我.
或者我也忍不了他.
所以明白和接受是很相反的.
所以我太太明知道我不是很喜歡吃東西.
但她每一次煮東西都是煮相配給我.
我明知道她是為我好的.

$^{801}$所以這種的相處就是在知道自己.
或者以自己為中心的生命裡面.
所產生一些跟你最親的人產生的衝突.
所以我不會贊成呂鎮忠的逆叛知己.
寧願用和本的好一點.
或者和修本的同心.
我們同有基督的心.
才是我們能夠掌握我跟神的關係.
我跟我的弟兄姊妹.
我的配偶.
我的親人的關係.
如果你想找知己的時候.
你會不滿足.
因為很多人甚至你的親人.
不是很了解你.
甚至不接受你.
所以知己是難求的.
我想保羅講到的這段經文.
呂鎮忠如果有機會.
我應該將來會再見到他的時候.
問他我不贊成你的解讀.
因為雖然幫到我這次準備這個講章.
但是這個翻譯.
知己是想世界的人去理解自己.
這是相當容易會走歪路的現象.
我們年長的弟兄姊妹.
或者我們見我們的父母.
他的世界會進入了別人不理解他的現象.
因為他孤單 他沒有朋友.
所以我們不應該以自己為中心.
我們應該跳出自己.
以基督耶穌的心為我們的心.
今天借呂鎮忠的翻譯.
跟大家去將這種心路歷程.
這個今天是我的一個叫做感受多一點.
解經上和修本的翻譯會好一點.
我們一起祈禱.
請上天賦予我們多謝.
讓我們有機會聚集在這裡.
一起去敬拜你.

$^{841}$我們在你旁邊.
或者我們周圍的弟兄姊妹.
我們都是有同一個心思 一個意念.
將生命擺上.
以神你為我們的救主.
也將福音成為我們人生一個很重要的目的.
我們多謝你.
因為你讓我們能夠在屬靈上同心.
以致我們超越了在世上找朋友.
找知己這種的需要.
我們多謝你給我們屬靈上的追求.
我們在每一個弟兄姊妹.
都能夠以福音這個心為心.
我們可以成為保羅的知己.
保羅一切的同工.
我們懇求 交託.
奉耶穌聖名祈求.
阿們.
\newpage



\section{撒迦利亞書 2:6-8}
\label{sec:jRV45ettAd0}
\textbf{2024-04-28 午堂《上帝眼中的瞳人》撒迦利亞書 2\_6-8  講員:陳韋安牧師}
\newline
\newline
連結: \href{https://youtube.com/watch?v=jRV45ettAd0}{\texttt{https://youtube.com/watch?v=jRV45ettAd0}} ~~~~ 語音日期: 2024-05-01
\newline
\newline
\hyperref[sec:w0vSwuU8gCk]{\small{< < < PREV SERMON < < <}}
~
\hyperref[sec:index_chronic]{\small{[返順時目]}}
~
\hyperref[sec:index_scriptual]{\small{[返順卷目]}}
~
\hyperref[sec:sCNFnlmgxC0]{\small{> > > NEXT SERMON > > >}}
\newline
\newline
撒迦利亞書 2:6-8
\newline
\begin{longtable}{cl}
\hline
\hline
章節 & 經文 (和合本修訂版)\\
\hline
2:6 & \begin{tabularx}{0.7\textwidth}{X} 耶和華說：「來，來！你們要從北方之地逃回；因我曾把你們分散到天的四方。這是耶和華說的。」 \end{tabularx} \\ \\ \relax
2:7 & \begin{tabularx}{0.7\textwidth}{X} 來！住巴比倫的錫安百姓啊，逃吧！ \end{tabularx} \\ \\ \relax
2:8 & \begin{tabularx}{0.7\textwidth}{X} 萬軍之耶和華在顯出榮耀之後，差遣我到擄掠你們的列國那裡，他如此說：「碰你們的就是碰他自己眼中的瞳人。 \end{tabularx} \\ \\
[1ex]
\hline
\hline
\end{longtable}
$^{1}$我們一起去祈禱.
弟兄姊妹我們一齊來聚集.
求主你去你的靈降下.
讓我們每個人在生命裡面.
來聽見,看見你自己對我們的恩典.
你藉著你昔日在撒迦利亞書.
先知的說話裡面提醒我們.
讓我們能夠看得見.
幫助孩子都能夠清楚地說出自己的話語.
2024年就過到差不多中間.
不知道大家這一年半年裡面過得怎樣.
我們都經歷了幾年的不太好.
股市,樓市都跌.
股市好像現在好一點.
不知道大家怎樣去看自己的將來.
看香港的將來.
在網上有一句說話都很流行.
都說人家悲觀.
就說我看得到幼稚及經的好處.
當中其實幾令我有印象.
反而是看見這個字.
看得到一些好處.
我們都和我們一齊來思考.
聖經裡面有關看見.
vision這個字眼.
如果你有vision的話.
其實都不難的.
基本上你花兩萬多元.
就買到一個vision pro.
apple的vision pro.
不單止有vision.
更加有pro.
你就可以看到一些看不到的東西.
當然你的樣子是有點奇怪.
但是你會看得到更加很多很多的資訊.
所以我們今天去思考一下這個課題.
有句說話這樣說.
我們的視野是決定我們的高度.
不知道你明不明白這句說話的意思.
視野決定高度.

$^{41}$當然這句說話你都聽過相反的.
高度又決定視野.
你有多高呢.
這個物理學問題.
你有多高就能夠看得多遠.
欲窮千里目.
更上一層樓.
視野決定高度是一個人生的哲理.
在我們人生裡面.
特別是我們基督徒.
當我們的眼光能夠去決定.
我們去看事情的高度.
和你所見到的遠像.
你看到有多遠.
你的生命就可以有多高.
你看到有多闊.
你的人生裡面就可以有多廣闊.
似乎看見是整件任何事情裡面的第一步.
你見不見到.
所以我們一起來思考這個課題.
講題叫做上帝眼中的同仁.
不知道大家對於這句經文.
有多了解.
其實你發覺這裡面都出現過幾節經文.
是關乎於上帝眼中的同仁.
這句說話.
我先講講經文的典故.
這句說話其實英文叫做甚麼.
叫做Apple of His Eyes.
很有趣的.
Apple of His Eyes.
中文翻譯為眼中的同仁.
如果你翻譯了很久.
你都大概聽過.
上帝看顧我們.
如同保護眼中的同仁.
這句說話.
經文裡面出過五次.
第一段是申命記的三十二章第十節.
然後說.

$^{81}$遇見他在曠野荒涼野秀洪澆之地.
就圍繞他看顧他保護他.
如同保護眼中的同仁.
第二段是詩篇十六篇第八節.
一個禱告.
詩人這樣祈禱.
求你保護我如同保護眼中的同仁.
將我隱藏在你翅膀的陰下.
第三段是箴言.
箴言這樣說.
遵守我的命令就得存活.
保護我的法則.
好像保護眼中的同仁.
第四段是耶利米安哥一段的哀求.
耶利米安哥這樣說.
是安的城牆.
願你流淚如何.
晝夜不息.
願你眼中的同仁流淚不止.
最後一段就是今天的經文.
《撒基拉書》第二章第八節.
然後說.
萬安之和說.
在顯出榮耀之後.
差遣我去懲罰那.
擄掠你們的列國.
摸你們的就是摸他眼中的同仁.
經文裡面.
其實英文亦都發現.
譯作Apple of his eyes.
這樣的一個翻譯.
不知道大家有沒有想過.
為何英文會叫Apple.
為何不叫Orange.
不叫Banana呢.
其實Apple of the eyes.
是一個英文的俚語.
意思是一個極度喜愛.
極度所愛慕的人.
眼中比任何人看來更加寶貴的一位.

$^{121}$這個俚語出現在公元第九世紀.
是一個非常古老的英文俚語.
後來去到莎士比亞的時候.
這個仲夏夜之夢.
就用了這個很出名的face.
Apple of the eyes.
後來這個英文的King James Version.
就翻譯到聖經裡面.
所以就見到Apple of his eyes這個英文.
說了這麼久.
其實我們都不知道.
為何會叫Apple.
不過當我們看回原文的時候.
希伯來文的聖經裡面.
原文就不是Apple.
原文解作甚麼.
就是解作瞳孔的意思.
就是我們眼球裡面.
即是黑色的部分.
不過無論是希伯來文.
或者希臘文的七十四年版.
都可以解作甚麼呢.
很有趣的.
解作女兒.
或者解作小人兒的意思.
所以不知道有沒有關係.
大家記得有關係嗎.
即是瞳孔英文是甚麼.
大家記得嗎.
瞳孔英文是甚麼.
是叫做pupil.
其實是學生的意思.
所以發覺都是一樣.
即是我們英文裡面的瞳孔.
都可以解作一個學生.
一個人兒的意思.
所以這裡有一個很特別的部分.
即是我們眼中裡面.
黑色的部分.
可以解作蘋果.

$^{161}$解作瞳孔.
解作學生.
解作女兒.
都可以的.
所以和本這樣翻譯.
翻譯成瞳人這個字.
我覺得很有意思的字眼.
即是中文裡面沒有將apple翻譯出來.
眼中的蘋果.
而是特意將它翻譯成眼中的瞳人的意思.
文言文的聖經就翻譯成母子.
中文的本就叫眼中之女.
一個很有趣的東西.
叫眼中之女.
就將daughter這個字翻譯出來.
所以和本這個翻譯.
瞳人都是一個很有味道的字.
原來古代的中國.
一早就有瞳人這個字眼.
所謂瞳人就是瞳孔.
都用一些古文.
唐朝裡面的道家唐兒歌裡面這樣寫.
骨寵神寒天妙氣.
一箱瞳人展秋水.
或者宋代裡面的一個詩歌這樣寫.
瞳人展水腰如涌.
一虎烏紗裡寒玉.
所以發覺原來不單止是十八瞳人.
瞳人都解作瞳孔這樣的意思.
所以無論是唐詩還是聖經.
古人一直都將我們眼球裡面的瞳孔.
都和一個人作出一些比較.
或者apple.
總之發覺是一個很多串連裡面.
眼中的那個人.
就是那個瞳人 瞳孔.
明白了這個apple of the eye的典故之後.
我們就要開始明白聖經裡面.
撒迦利亞書裡面這句話.
apple of the eye的意思.

$^{201}$先簡單講一下撒迦利亞書的背景.
撒迦利亞書是一本後期的小說.
一本記載了以色列人被擄流亡之後.
將要歸回的小占之書.
撒迦利亞對流亡海外的以色列人講話.
同時對著一班耶路撒冷的漁民講話.
你知道公元前五百六年.
耶路撒冷淪陷.
以色列人被擄.
其實分了幾批人.
有些是流亡海外的.
有些是移民到國外的.
前者是被迫被擄走的.
有些是自己移居的.
但也有些仍然在耶路撒冷.
一個頹垣敗瓦的廢墟裡面.
繼續生活的一些以色列人.
所以撒迦利亞書和黑海書.
有一個同一個的小說目的.
就是代表著上帝.
去吩咐一班身處異地的以色列人.
一班殘留在以色列裡面的以色列民.
吩咐他們要一起去重建耶路撒冷.
恢復聖殿敬拜.
所以發覺整本書.
其實都是關乎一個很重要的概念.
就是vision這個概念.
撒迦利亞書強調耶和華上帝的漢見.
藉著撒迦利亞先知的漢見.
講述上帝漢見耶路撒冷的願景.
所以首八章經文.
都是藉著八個不同的意象.
來表達上帝望著頹垣敗瓦所望到的願景.
正正就是他望到的另一番景象.
所謂漢見或者所謂vision或者意象.
是一些不容易消化的東西.
明明是望著一些頹垣敗瓦.
但是他要望到並且宣告.
這是一個極度興旺的事情.
等同於當股市去到萬五萬六的時候.

$^{241}$有一個大師就說望到是四萬點.
不知道有沒有聽過.
今年會去到四萬點.
這些vision是不容易說的.
因為好像是很不可能的.
所以同樣的道理.
在廣東被人笑.
所以同樣的道理.
當日撒迦利亞先知都是這樣.
望著一個頹垣敗瓦的廢墟.
撒迦利亞是要說他望到的.
說這是你悲觀.
我望到不是這些.
我望到是一個完全不同的一番景象.
我看到樓市股市升得很厲害.
我看到香港是十年前的景象.
店鋪不斷加租.
晚上很多人.
這些景象是沒有人相信的.
甚至乎還叫一班移民的人回來.
他們去到英國已經開始找客人買樓.
回來香港重建香港.
這裡很好.
他們會笑.
他們不會相信.
因為他們看不到那個景象.
在這個現實景象和意象看見.
那個相差甚遠的前提之下.
就補充這句話.
希望這班海外的市民能夠安心.
不要害怕.
回來耶路撒冷.
就是這句話.
第二章一到六節.
我們聽聽我說出來.
我從前分散你們在天的四方.
現在你們要從北方之地逃回.
這是耶和華說的.
與不被鄰人同住的錫安文娜.
應當逃脫.

$^{281}$萬君子耶和華說.
在顯出榮耀之後.
差遣我去懲罰那老弱你們的之列國.
摸你們的就是摸他眼中的同人.
這裡耶和華說.
上帝要求這班流亡海外的市民歸回.
叫他們不用怕.
因為上帝會保護他們.
因為他們就是上帝眼中的同人.
The apple of his eye.
甚至乎上帝更加多說一句.
一句很生動的描述.
他說摸你們的就是摸我眼中的同人.
簡單來說.
一個極度繪形繪聲的描述.
誰敢摸鬥這班色情文.
等同於你插我眼中的同人.
這麼模範一樣.
誰敢搞你們.
就等同於搞我眼中的同人一樣.
所以這是一句非常厲害的話.
耶和華說誰敢搞你們.
就等同於搞我的眼球.
我的同人.
因為你們就是上帝眼中的將來.
任何人傷害你們.
就是傷害我的眼球.
我的vision.
色尼亞書正正要表達你們不要害怕將來.
雖然你們還未見到這個將來.
但是你們就是我眼中的同人.
這是一個切膚之痛的愛護.
說起這件事.
我先說一下我上年發生的單件事.
上年我帶我女兒去了屯門安定友愛村.
不知道大家有沒有去過.
友愛村有一個商場.
天台有一個玩Xgame的滑板場.
也很有名.
整個屋外有一個很大型的Xgame場.

$^{321}$可以讓你踩滑板,踩滑板,踩BMX.
整個場沒有人.
我們兩父女就整個早上在那裡玩.
她踩滑板,我踩滑板.
差不多都累了.
我都想走.
我女兒就說玩多十分鐘.
誰知道十分鐘之內就出事了.
差不多走的時候.
我的滑板突然間失腳.
滑板就飛到我女兒那邊.
我女兒踩得很懶.
她怕怕就突然間就跌倒了.
失平衡了.
跌倒的時候就神思凝重.
我就心知不妙.
她沒有哭.
那時候她就戴口罩.
她脫下口罩就發現.
她口罩裡面有一個三分二的門牙.
在口罩那裡.
我就嚇死了.
整個門牙就沒有了三分之二.
就立刻流血.
滿口都是血.
那時候是非常的後悔.
作為一個爸爸.
因為她是一個行使.
有我在村立刻就買了一瓶冰水給她清洗.
做很多事情.
打電話看看醫生.
後來就立刻找了我以前的大學同學.
立刻幫她補牙.
不過最令我發現是很心痛的.
當然是有期的.
補得很漂亮.
但是發覺那時候我就很明白.
什麼叫做Apple of the eyes.
She is the apple of my eye.
因為她是我的眼中之女.

$^{361}$你明白到眼中的同人.
正正就是這個意思.
你寧願自己撲崩牙.
都不想自己的女兒撲崩牙一樣.
她是我的心肝寶貝.
如珠如寶.
萬千寵愛的女兒.
她是我眼中的同人.
你寧願自己受傷害.
都不想眼中的同人有任何的傷害.
所以領子妹.
如果我們說天父上帝.
是看我們為眼中的同人的時候.
這是一個非常非常貼切的形容.
誰搞她的心痛.
那種愛護.
正正是說到這一點.
不過今天這篇道.
不單單是要停下來.
今天其實不單單跟大家說.
你們是上帝眼中的同人.
就完結.
這些說話是真的.
但我想都不是單單.
是一句安慰的訊息.
我們要認定.
我們是上帝眼中的同人.
不過我們更加要知道.
在同一雙上帝的眼睛那裡.
撒迦利亞所說的是另一番的視野.
一個更加高層次的視野.
同一雙的上帝眼睛.
他見到我們是上帝眼中的同人.
他更加見到更加多的事情.
撒迦利亞書第八章第六節.
我們一起去看看.
.
第八章第六節.
先知撒迦利亞描述了一個.
另一個上帝眼中看見的異象.

$^{401}$小路子記載的是.
同一雙上帝的眼睛.
上帝所見到的.
以至於文章要重回.
這個重建的城市.
這個城市當中的景象.
上帝見到什麼.
上帝見到年老的男女.
坐在耶路撒冷的街上.
見到街裡面很多的小朋友在那裡玩耍.
這個是耶路撒冷和平的景象.
這個是上帝從輝煌裡面所見到的景象.
所謂的vision 所謂的異象.
所謂的看見.
正是一種視野.
一種尚未實現.
尚未發生.
是以行動的之先.
首先必要見到的因素.
不過請留意第六節.
對於視野.
關於眼睛的問題.
而我上帝提出一個值得我們深思的問題.
上帝反問這班以色列的愚民.
這些事在聖餘的眼中看為稀奇.
在我眼中也看為稀奇嗎.
即是這班被上帝眼中看為同人的以色列民.
竟然對這個異象感到很驚奇.
很出乎意料.
覺得是很不可能.
上帝就反問他們.
難道在我眼中也看為稀奇嗎.
實在可圈可點.
這句說話.
翻出來的意思是什麼.
上帝說這些事情你們沒見過嗎.
你覺得很陌生嗎.
你覺得很出奇嗎.
你以為你和我是同一般見識嗎.
你們從未見過覺得很陌生很出奇.

$^{441}$就等同於沒有嗎.
這句話正正是導出了.
以色列民和耶和華上帝的眼光之間的差距.
這班原民見到的只是一個被摧毀的廢墟.
然而上帝在這個頹垣敗瓦的地方裡面.
他見到一個和平的國度.
一件很諷刺的事.
這班人是被上帝眼中看為同人的人.
自己卻是看不到上帝要看的東西.
其實這班人不是看不見而是沒見到.
這班人不是瞎的.
他們是有眼的.
他們見到很多很多東西.
不過他們因為見到很多很多東西.
令他們見到很多很多很多見不到的東西.
即是我們今天見到很多東西.
令我們反而見不到一些重要的東西.
經文告訴我們這班人不是看不到.
而是看得太多.
令他們忽略了見不到.
當我預備這篇經文的識經的時候.
發現一個古代教父.
對於這段經文的識經很特別.
這個教父叫做狄狄莫斯.
一個盲的基督徒.
一個教父扮的.
這個叫盲人.
對於這個的《十一日書》是這樣說.
他說當我們心靈的眼睛被啟迪的時候.
我們就能夠見到上帝給我們的意象.
重點就是說.
我們其實重點不是一個肉眼裡面的看見.
而是我們有沒有能夠可以.
一個所謂心靈裡面的看見.
相反我們不知道大家.
眼睛重要.
可能到四十多歲開始有老花.
有近視有遠視.
有很多不同的東西.
這個世代裡面我們要看的時候看得太多.

$^{481}$大家沒有發現.
基本上從Vision Pro到手機.
到很多不同的東西.
其實太多太多東西給你看.
很多人都是晚上拿著手機.
可以看了很多時間.
你看的東西其實不難.
今天特別是這個世界.
我稱之為TikTok化.
以前我們都說已經沒有人看文字.
現在連片都沒有人看.
片要十五秒就看完.
十五秒都太長了.
以上都太長了.
因為太多太多東西給你可以看得到.
令到我們反而看不到很多應該要看得到的東西.
當你的眼球裡面有太多資訊.
真正的看不見不是你瞎了.
而是給很多很多的資訊.
令到你看不到你應該要看得到的東西.
因此看見的相反詞不是看不見.
而是見到太多之後忽略.
明明在你眼前.
明明就擺在你眼前.
你卻是有眼無珠.
眼睜睜看著龍就看不到.
所謂視而不見.
正正就是這個意思.
面對著我們今天的香港.
這個社會.
你可以看到很多很多眼前的事情.
你會發覺元朗大馬路的店舖.
以前都沒有這個情況.
很多都是吉舖.
你會看到很多很多的資訊.
很多很多的消息新聞.
但是這些看見都不是最重要的.
上帝要我們看見的是看見上帝所看見的東西.
擺在我們眼前的.
但是我們很多時候因為我們太過多.

$^{521}$令我們看不到.
我們看見一些應該要看見的東西.
不要給一些不應該看見的東西.
來阻礙我們的心靈.
我知道這個不容易.
我們很難突然間會醒覺.
或者突然間看清很多的事實真相.
很多的事情就出現了很多的所謂.
全部都是一個見字.
偏見 成見 短見.
這些看見其實都不是看不到.
而是看錯了.
或者看得太多.
我爸爸上星期離開了世界.
回到天家.
在過程裡面.
其實都有機會分享.
但是其中一個令我很深刻的是什麼呢.
其實人都是很有限.
我爸爸在去年暑假就發現有肚血.
胃就以為胃潰瘍.
然後發現原來是胃癌.
暑假就切了胃.
好端端的.
在今年新年的時候.
他爸爸煮的豬肉是很厲害的.
就為一班親戚在年初二煮了大餐飯出來.
年初十就開始痛.
我發現原來我們很多東西都看不到.
他們說肚子很痛.
不知道是食道問題.
還是因為沒有胃.
還是食道問題.
痛從哪裡來呢.
才知道原來是肚子裡面已經滿是擴散了的癌細胞.
我們很多時候都看不到.
我們COVID都是一樣.
COVID的最大問題是什麼呢.
就是你看不到病菌在哪裡.
你不知道哪裡有病菌.

$^{561}$所以全世界都戴口罩.
為了避免看不到病菌.
我們很多東西都看不到.
在腦裡面發生什麼事都看不到.
小小的很小的病毒你看不到.
將來都看不到.
但是在爸爸在病床裡面的時候.
很痛很痛很痛的時候.
醫生打多少針都幫不了他的時候.
他一開始都問為什麼天父要這樣養.
他都學會了原來他見到耶穌的恩手.
他就說我現在就願意去馴服你.
你什麼時候接我走就走.
我想我們在我們很多的人生裡面.
見到一些很重要的東西.
是很重要的.
很多東西我們都見不到.
但唯有我們能夠見到一些最重要的東西就足夠了.
正如昔日猶太人一樣.
他們眼前的全部都是一些很荒廢的景象.
但是上帝要求他們見到一件事情.
你見到自己是上帝眼中的同人.
見到上帝的城市裡面.
上帝所見到的.
弟姐妹我們有沒有什麼東西值得我們去看見呢.
在我們生命裡面.
有什麼東西值得我們去看見呢.
這些東西絕對不是一個高科技.
或者任何的事情.
更加多的資訊能夠幫助我們.
唯有我們憑著我們信心的眼睛.
能夠看見上帝恩典在當中.
我一起祈禱.
因為你一直都在.
讓我們能夠看見你自己的恩手.
看見你在我們生命裡面所過的大靈.
縱然有一百萬的事情.
一樣事情我們可以知道.
我們可以看見.
我們可以掌握.

$^{601}$但求求你讓我們能夠撥開這些事情.
單單可以看見你自己的恩典在當中.
讓我們能夠去存活下去.
讓我們能夠看見你的恩手.
我們樂觀不是因為看見其他的東西.
而是單單因為看見你自己的應許和景象.
願意讓我們看見你的看顧和愛護.
讓你能夠開啟我們的雙眼.
奉主並求.
阿們.
\newpage



\section{帖撒羅尼迦前書}
\label{sec:sCNFnlmgxC0}
\textbf{2024-05-02 周四晚查經班《帖撒羅尼迦前書(5)》 講員:曾裔貴牧師}
\newline
\newline
連結: \href{https://youtube.com/watch?v=sCNFnlmgxC0}{\texttt{https://youtube.com/watch?v=sCNFnlmgxC0}} ~~~~ 語音日期: 2024-05-03
\newline
\newline
\hyperref[sec:jRV45ettAd0]{\small{< < < PREV SERMON < < <}}
~
\hyperref[sec:index_chronic]{\small{[返順時目]}}
~
\hyperref[sec:index_scriptual]{\small{[返順卷目]}}
~
\hyperref[sec:U5CngR8W_hQ]{\small{> > > NEXT SERMON > > >}}
\newline
\newline
$^{1}$我們一起有個簡單的祈禱.
謝天父再一次給我們在晚上.
過了幾天都挺天盤大雨.
甚至在元朗都有下冰雹的現象.
也都看見在梅州那裡有特大的山泥傾瀉.
導致現在都有四十多人來到罹難.
我們想到有時天災人禍.
好像在以色列,在俄羅斯.
都一樣繼續有戰爭.
都是我們不想見到的事情.
很怕望處都被我們屬處的兒女都祈禱.
求神來止息戰爭.
讓平安臨到有需要的人的心靈.
也都想到現在在美國有很多間大學抗議.
我們都不希望帶來一種暴亂的一個反政府.
我們希望在一個很理性,很平和的情況底下.
讓事情可以得到最好的解決.
回來我們今天晚上我們重新查看.
我們再看貼薩諾尼加前書.
幫助我們對於主在來和當時代的信徒.
他們怎樣面對在羅馬帝國.
在猶太教那個雙重的影響底下.
他們仍然持守的那個信心,愛心.
都讓我們能夠在其中有所得著.
幫助我們能夠在我們自己現在香港的21世紀.
我們能夠可以了解怎樣來去處我們的生活.
和我們為主去作言證.
求神帶領我們奉耶穌基督的名.
(主) 阿們.
上個星期都比較詳細將保羅的信王外.
不同的結構,不同的組成.
比較深入揭露,講解了.
所以盼望大家如果上一個星期沒有聽.
可以利用一些時間聽.
因為我都差不多用了一個小時講關於信王外.
保羅在幾段聖經.
羅馬書,歌羅西書,貼薩諾尼加前書.
特別是哥林多前書大家比較熟悉的.
我們這四段經文都是重要的關於信王外.
是連續出現的那個案件.

$^{41}$有些只是講信心,愛心.
有些可能是講關於另外一個配對.
但一定有它那個經文的處境的.
所以大家都可以揣摩一下.
其中就是不是死板板的.
保羅是因應那個群體的需要來去運用.
即是說不是一些詞彙,修辭那麼簡單.
而是實實在在信王外這三樣東西.
是信徒必須由信主之後就要去追求.
要去建立的.
而且傳到永恆的.
去到永恆都一樣有信心,有盼望.
但是唯有愛是最重要的.
這個是我們要留意的.
尤其是保羅或者其他的使徒.
清晰一件事.
雖然我們都不是太知道做物主到底是怎樣的.
但是我們知道一樣東西就是.
愛是從神而來的.
神就是愛.
《肉眼一書》四章十九節.
因為神就是愛.
從來不會有一次聖經.
因為神就是信心.
從來沒有的.
因為神就是盼望.
都是沒有的.
我們有一個叫盼望的神.
但是不會有信心之神.
但是神就是愛.
這個是很重要的一個永恆的屬性.
所以其實信心,愛心,盼望.
是基督徒一生信主之後.
必須盡力去追求的.
就是說原來是有不同程度的成長的.
我們的信心.
對神,對人的信心.
我們的盼望.
對永恆的盼望.
是不是與人俱增呢?.

$^{81}$我們對於愛.
今天就會講比較多愛了.
第三章這裡.
保羅自己講述的.
對當時的貼薩諾尼加教會.
他們處於的那個時代背景.
貼薩諾尼加其實與羅馬帝國的關係是相當好的.
他們得到最惠國的待遇.
就是免那些軍費的.
所以是很得到君王的那個歡心的.
因為正正就是他們壓著在.
那個羅馬的皇帝奧古斯都.
就是這個安東尼.
來到去打敗了.
接著屋大維來到去成為了羅馬皇帝.
正正貼薩諾尼加的地方的官.
就壓著在這個戰爭上面.
所以當奧古斯都做了皇之後.
很自然就會favor這個貼薩諾尼加城.
所以就免他們的差餉,軍費.
還有會有很多的待遇等等.
還有君王會去到出巡.
都會去貼薩諾尼加城這個地方.
是一個地中海今天非常之美麗的地方.
叫做薩諾尼加.
一個很出名.
地中海很美麗的一個地方.
就是保羅當年在這裡.
不過是四個星期 他就要離開.
所以第三章其實就是這一本書.
貼薩諾尼加前書的主題.
就是保羅西拉去到瓦典.
第二次的宣教旅程.
輾轉去到瓦典之後.
就很擔心到底貼薩諾尼加的弟兄姊妹.
會不會受到那些逼迫而來到去退卻呢?.
會後退那個信心.
不再回教會跟隨著呢?.
所以就打發了提摩太.
回去那個地方去探望教會.

$^{121}$了解他們的信心 他們的情況.
所以這一本書寫作的原因.
就是提摩太在貼薩諾尼加.
回到瓦典會合保羅.
保羅聽到提摩太所提到.
貼城現在的狀況.
就是他們的信心和愛心.
都是能夠很穩固.
保羅就很開心.
那個開心就是一個喜樂.
就是這一本書寫作的原因.
但是我們又提過了.
每一章 包括後書都有的.
每一章的結尾都提到主耶穌回來了.
主耶穌因為已經死了升天了.
當時代的人是不會相信.
有一個死了復活的人的.
但是《新約聖經》試圖他們.
再再告訴我們有一個人死了.
但是已經復活了.
而且現在在天堂.
甚至他會再回來.
所以這裡說的主在來.
就是說主耶穌第二次.
再回來的榮耀降臨.
所以這個是末世論一個很重要的主題.
差不多是唯一聖經還沒有應驗的啟示.
所以你想想如果信徒.
一生從來沒有聽過主在來.
從來沒有想過主耶穌會再回來.
那麼你的信仰其實是假的.
不是真信仰來的.
主耶穌是會再回來的.
是充滿了在聖經的.
有些人數過在新約.
三百多次說主耶穌再回來.
還要是還沒有應驗的唯一.
神應許的那個預言.
那你就可以想像得到.
這是幾種重要的聖經的新約的主題.

$^{161}$差不多是整個的脈絡.
無論我們怎樣悲觀對世界.
對政治,對戰爭,對人類的災難,瘟疫.
無論怎樣悲觀都好.
主耶穌再回來.
是一定會帶來一個新天新地.
一定會帶來一個很大的整個人類的改變.
或者說正是那個結局的.
進入那個光明的永恆.
而且這個就是我們將來在天堂的歸宿.
所以如果我們從來都不去想.
和都不去準備.
準備甚麼呢?.
就是準備我們的順忘愕.
就是準備我們在人生.
我們有沒有盡忠為主.
來去站穩在真理的福音.
所以你可以說雖然第三章很淺白的經文.
但是裡面的含義都很深的.
上一個星期我們只能夠講到第二節.
講到我們這一位的提摩太.
他是基督福音的執事.
他是為主來到牢府的一個執事.
在基督的福音上作神的執事.
這一位就是提摩太.
這位年輕人.
其實保羅認識他.
可能真的不是很久.
一年都未必夠的.
因為第二次的宣教旅程.
保羅就去到這個盧斯德.
就見到這個年輕人.
那個地方的教會的人就稱讚他.
保羅如果真的看到他內心有一份的單純.
一份可造之材的材料.
心靈.
就帶他去傳道.
接著又要幫他行國禮.
原來這個年輕人未行國禮的.
猶太人第八天就行的了.

$^{201}$保羅自己都是第八天行國禮的.
就是很根正苗紅的猶太人.
因為他是猶太教的.
但是提摩太就不是在猶太教的氛圍長大的.
因為他爸爸是外邦希臘人.
所以他是半個猶太人.
都有這些影響.
但是未必是像保羅這些這麼嚴重的訓練.
所以保羅就很想他在猶太的一些圈子裡面.
不會讓人去講是非.
你說你自己是猶太人.
又不受國禮.
一句而已已經沒有聲音.
所以保羅就幫他行國禮.
邀請他行國禮.
不是關乎得救.
是關乎因為你是猶太人.
所以我們應該要活在猶太的文化底下.
讓那些其他的猶太教的人.
不能去指指點點.
這個年輕人說自己是猶太人.
其實都不敬虔的.
一句你已經令到很多的人受影響.
所以提摩太就讓保羅幫他行國禮.
所以記住提多同樣是外邦人.
保羅是沒有幫他行國禮.
甚至堅持不幫他行國禮.
因為他是外邦人.
他不需要受猶太教的文化的影響.
所以保羅是分得很清楚.
很有眼光.
很有識別.
很有接受的文化.
但不會被這個文化蓋過了福音.
福音的真理是蓋過文化.
所以不要反過來.
我們由第三節開始看.
三章第三節.
首先是提摩太去兼顧他們.
保羅在雅典打發提摩太回去歇薩羅尼加城.

$^{241}$去探望兼顧他們.
並且在你們所信的道上勸慰你們.
免得有人被朱般的患難搖動.
朱般的患難.
你可以想像到有多少種類.
猶太教其中一個.
因為保羅是在提薩羅尼加城的猶太人會堂.
有會堂.
保羅先去會堂.
向那些猶太人講解神的道.
然後真的有猶太的婦女.
和外邦的希臘婦女信主.
也有一些其他的外邦人等等.
這裡可以想像到.
猶太教如果還堅持在摩西的律法.
還是未能夠.
那個心靈是有一種通達.
對那個道不是一下子的.
不是每個人都能夠像保羅那樣.
保羅都經過很大的神的工作.
有耶穌的顯現.
他才會改變的.
不可以就這樣討論一下.
寫一篇文章就可以改變保羅的.
他那種的他自己對猶太教的眼光.
需要真的活的主是會顯現的.
這裡就讓我們看到.
朱本的患難來自猶太教.
也都來自地方官.
所以他們就向當時.
在貼壽羅尼加城的那些官告狀.
就說這個保羅是搞亂天下.
你想一下.
一個最為國待遇的貼壽羅尼加城.
和羅馬的皇帝是很友好的.
關係是很好的.
說這個人來搞亂天下.
所以你想一下那個罪名是很大的.
等於是煽動分裂國家的.
這個是一個很嚴重的罪名.

$^{281}$《少林寺傳》.
我很快讀給大家聽.
17章.
大家不用看了.
17章的第五節.
但那些不信的猶太人.
心裡面窒道.
從心裡面衍生出的那種因為不滿意.
自然就會想一些很惡毒的計謀.
或者說辭.
來去給保羅.
就說招聚了一些使者匪類.
就是找一些人來去你一句我一句.
搭伙成群.
聳動合城的人.
闖進耶穌的家.
要將保羅西拉帶到白城那裡.
找不著他們.
就把耶穌和幾個弟兄.
拉到地方官那裡.
就是當時的羅馬的駐城的官.
接著就喊叫.
留意這裡.
那攪亂天下的也道者女來了.
來到我們帖撒羅尼迦城.
耶穌收留他們.
就是在那裡業造.
控告耶穌.
這些人都違背了該殺的命令.
接著留意一下.
說另外有一個王.
叫做耶穌出現.
所以眾人和地方官聽見這話.
就驚慌.
於是取了耶穌和其餘之人的寶藏.
就釋放了他們.
驚慌就是因為擔心.
這個這樣的傳福音的行動.
變成一個政治的事件.
羅馬的君主.

$^{321}$首先就是將那個罪為國待遇收回.
和派羅馬的士兵直接來干預.
就會帶來一個很大的影響.
所以地方官都不敢擔保誰.
或者包庇保羅等等.
所以你會看到保羅在傳講的主耶穌.
其實是帶來.
你可以說是真是幾令當時代的人.
有一個很大的考驗.
因為某程度保羅真是在說另一個王.
你看看《貼索羅尼加前書》.
五次講主耶穌再回來.
是回來的再來君王.
的確再來君王是在說耶穌.
不過不是回來.
是好像要來和羅馬去吞併他.
而是全世界都要臣服在主耶穌的威權底下.
所以你可以說對當時那些不知道.
以為只是政治上來干預.
是一個很敏感的題目.
是帶來跟隨者一定會受到諸般的患難.
來去搖動.
搖動有沒有影響.
有沒有動搖到信徒的內心呢?.
這個就是牧者要關心的地方.
就是關心他們不是將那些患難.
去患難法減輕基督教福音.
那個需要受苦.
而是在受苦裡面.
怎樣能夠牧養,陪伴和一起渡過.
一起渡過可以在監獄裡面.
保羅自己在菲律賓都被人打到半死.
那些傷痕累累都還是身傷舊患.
都還是在他的身體裡.
因為沒多久就去了貼索羅尼加.
在貼索羅尼加是一個月.
那就是說是一兩個月的事.
保羅還沒有完全康復.
那些傷都還在.
他在這個提摩泰後書都說.

$^{361}$我們從前在安提摩,爾戈廉,路斯德.
所忍受的逼迫,苦難.
和在菲律賓被害,受辱.
是你們知道的.
你可以想像得到.
不是的 應該是二章.
二章就已經說了.
沒錯 果然在二章已經說了.
二章就是說.
「第一門你們玄曉得.
我們進到你們那裡並不是豪言.
我們從前在菲律賓被害,受辱,收辱.
是你們知道的.
然而是靠我們的神放開膽量.
在大的征戰裡面.
將神的福音傳給你們」.
所以原來福音不需要幫助他們減輕受苦的部分.
而是要將福音明明白白講明白.
就有那種很大的威力.
那些人受苦都願意跟隨主.
這才是福音.
不是說「信耶穌就沒事了.
不時都沒事了 你有病嗎?.
不怕的 我為你祈禱就會好.
你很窮嗎?.
我為你祈禱你就很有錢了」.
這就變成了我們扭曲了的福音.
福音可以令到我們的生命改變.
生命改變帶來生活改善是肯定的.
我已經不再吸煙了.
我吸兩包煙的.
信了耶穌.
但真的由那時候開始.
幾十年沒有再吸煙了.
你想想你幫我算一算這條數.
兩包煙現在多少錢一包?.
然後閉癌那裡怎麼算呢?.
這條數是不是?.
那時候我吸煙真的很多.
經常咳嗽 經常要看醫生.

$^{401}$你可以想像得到.
生命改變帶來生活改善.
但是從來沒有保障.
從來沒有允許說.
每個信耶穌的人都要免去受苦.
所以其實這本書.
是在苦難裡的初信教會.
他們怎樣來建立他們的信望愛.
牧者和信徒之間.
怎樣建立這份愛心信心的互動.
使信徒仍然得到生命的建立呢?.
雖然牧者已經不在他們身邊.
但他們仍然因為對牧者的信靠.
牧者所傳遞的真理.
他們可以透過寫信.
不是今天WhatsApp半秒鐘就收到.
而是好幾個星期甚至乎好幾個月.
才收到那封信.
你可以想像得到.
那個緩慢有很多空白位.
需要他們的信心.
需要他們自己去堅忍.
需要聖靈的工作在人的心裡面.
這些都是我們要留意的地方.
因為今天太多的按個鍵.
就可以見到成效的世界.
我們失去了那種等候神的空間.
我們需要學習的.
有時包括我自己也是.
發信息給別人.
希望他快些回覆我.
但變成慢慢一個惡性循環.
好像變成你又快 別人又快.
大家就變成建立在一種很快的文化底下.
不知不覺間 其實都不健康的.
有時你一個小時不看電話.
自己都不覺得不舒服 覺得很慘.
你就可以想像得到.
初期教會其實他們在這些環瀾裡面.
牧者不在場.

$^{441}$牧者擔心他們會不會失腳呢?.
會不會後退呢?.
這個就是在這本書裡面.
很重要的精髓.
即是間接探討了一個問題.
就是牧者怎樣去牧養受苦的弟兄姊妹呢?.
所以第三章裡面.
保羅原來沒有甚麼所謂法術.
沒有甚麼所謂他讀過甚麼MBA.
只是一個牧者.
是我帶你們信主.
你們得到神的福音.
現在因為各種原因.
我要立即離開.
你們是否繼續在這個福音上去成長呢?.
保羅很想盡快去探望他們.
但是在現在的處境是不行的.
所以無論如何.
保羅都打發他自己信任的提摩泰.
去牧養貼守羅尼加城的弟兄姊妹.
所以當回來的時候.
第四節.
「我們在你門那裡的時候」.
預先都說了給你們知道.
「我們必受患難 以後果然應驗了」.
保羅不是說要誇耀他自己.
有甚麼神經妙算.
有甚麼微卜先知.
而是眾所周知的.
跟隨耶穌 背十字架的耶穌.
被人這樣去凌辱的主.
就是我們那位真正的跟隨者.
我們要跟隨他.
所以背十架跟隨主.
是必然的事.
基督教我覺得最寶貴的地方.
基督教不是要鼓勵人.
隨便找些苦來受.
好像一個虐待自己的宗教.
就不是的.

$^{481}$而是在神安排命定下的.
命定即是你和我都不知道.
而我們可能進入了那個時代.
而那個時代是一個苦難的時代.
我們是你想逃避都不能逃避的.
要麼就信 要麼就不信.
只有這樣對錯的情況.
我們就在我們還沒進入真正的.
大的苦難之前.
其實就不是要我們信還是不信.
是要我們裝備我們的信心.
裝備我們對主的盼望.
裝備我們的愛心.
我們當進入了那個神命定的.
那個時代的苦難.
所以有些人就站穩了.
有些人就被淘汰了.
是因為在平安的日子.
我們根本沒有裝備.
甚至在我們的腦海裡.
以為患難是不應該出現的.
我們搞錯了基督教.
基督教其中精彩的地方就是.
就算在患難.
那個信仰都是這麼真實.
和會更加令到那個在苦難裡的人.
那個基督徒更加成長.
所以患難只會令到基督徒.
那個信望愛更加建立.
更加堅固.
第四節就這樣說.
果然應驗了.
你們都知道.
為此.
為了這件事.
我們就不能夠再忍了.
因為既然患難是命定.
既然患難保留一早已經說了.
就是說你們要信耶穌.
是伴隨著困難的.

$^{521}$伴隨著可能是帝國帶給你們的迫害的.
要你們跪拜偶像.
要你們敬拜羅馬君王.
要你們做一些違反真理的事.
如果你是猶太人.
再加上自己猶太的同胞.
對舊約的誤解.
以為信耶穌的猶太人.
是離開自己的真神.
離開自己的信仰.
舊約的信仰.
所以其實是要面對幾個困難的.
所以初信者.
其實我們都真的要問自己.
為何要信耶穌.
回到這裡.
為此我既不能夠再忍.
就打發人去要了解你們的信心如何.
恐怕那誘惑人的.
到底誘惑了你們.
誘惑人的似乎是眾數.
叫我們的奴婦歸於徒然.
本來很隱晦.
哪些人是一些明顯的逼得.
因為《少林寺傳》提到明顯的猶太人.
因為他們以為自己很熟悉十誡.
很熟悉摩西的律法書.
所以以為保羅變了.
以為保羅離經貫道.
既然離經貫道.
自然就教導其他人都離經貫道.
所以他們要打擊保羅.
所以聽聞保羅在不同地方.
已經傳這個假道.
所以他們就在心裡面嫉妒保羅.
所以就想出可能很多的錯謀.
來對保羅加以攻擊.
但無論如何在這裡讓我們看到.
保羅擔心的是.
貼成弟兄姊妹的初信的內心.

$^{561}$如果真是四個星期.
保羅就離開.
這封信輾轉.
在雅典上去的時候.
可能是幾個月的時間.
所以保羅就要來打發人.
要了解貼成的弟兄姊妹.
他們的信心到底是怎樣.
信心真是要來保持.
「恐怕那誘惑人的到底誘惑了你嗎?.
叫我們的奴婦歸於徒兒了」.
「今朝早郎一條兩條」.
(英文翻譯).
好 再來我們看第六節.
「但提安是剛才從你們那裡回來.
將你們信心和愛心的好消息.
報告給我們.
又說你們常常紀念我們.
切切的想見我們.
如同我們想見你們一樣」.
這裡提到提安太的回報.
一方面不只是人數這麼簡單.
而是這班弟兄姊妹他們的信心,愛心.
這裡提到兩樣.
不是提信忘愛.
提到他們的信心和愛心.
當然他們切切想見保羅.
都是一份盼望.
不過特別強調.
貼撒羅尼加城的弟兄姊妹.
在信心上,在愛心上.
是一直增長的.
所以是一個好消息.
告訴保羅他們現在的狀況.
這真是令人很振奮的一件事.
牧者到底我們喜歡.
或者我們內心的歡喜.
是在哪個狀態呢?.
看到弟兄姊妹甚麼表現.
我們感覺雀躍,開心.

$^{601}$相信一定是信心,愛心的成長.
教會的弟兄姊妹整體的信心很堅強.
愛心是能夠有木階建的彼此相愛.
無論對初信的或者內的.
我們都有這份的愛心.
這些一點都不容易.
一方面每個人的性格.
每個人的背景都不同.
來自不同階層.
所以如果越是不同.
那個困難度又會再提高.
所以如果信心,愛心是一個好消息.
保羅很擔心教會的問題.
可以解決了.
令到他可以很開心地.
很歡喜地去為著這班弟兄姊妹的靈命.
所以一方面保羅很想去見他們.
但礙於現在不能夠.
所以提摩太太帶著好消息回來.
他們當時應該都可以行路.
不過船會更加快.
因為坐船在地中海其實很快.
應該坐船都可以更加快見到保羅.
因為保羅在雅典.
如果是陸路就會比較遠.
但在這個時候.
第七節留意一下.
所以弟兄們.
我們在一切困苦,患難之中.
去到雅典去傳道的時候.
又遇到那些欺凌極哲學家.
對保羅那種嘲笑,批評.
保羅向他們傳遞一個叫「微色之神」.
在這個叫阿勒巴古的地方.
雅典本身就已經有一個眾神廟.
一個很宏偉的建築.
應該當時是很多人聚集的.
所以保羅面對的是不同的困苦,患難.
但是保羅說.
因著你們的信心.

$^{641}$我們就得了安慰.
就是說弟兄姊妹那個靈命的表現.
或者生活上反映基督的信心.
如果是堅固.
就帶來穆者的心靈.
得到很大的安慰和建立.
這也是一個很實在的一件事.
聽弟兄姊妹很熱切的傳福音.
作見證.
這是一個令穆者感到欣慰的地方.
第八節.
你們如果靠主站立得穩.
我們就活了.
我們在神的面前.
因著你們的甚是喜樂.
為這一切的喜樂.
來為你們報答神呢?.
好像一個反問句.
其實就是不需要.
因為保羅自己已經是極其的喜樂.
這個喜樂是因為見到弟兄姊妹.
站穩在那個真理.
在那個初信的階段.
信心和愛心都有成長.
有堅固.
所以他就非常之感到欣慰.
然後為了這一份的喜樂.
他所…就是一個好消息.
帶來就是應該.
心靈會得到很大的快樂.
保羅說不知道應該用什麼話.
為他們去報答神.
不是不知道的.
怎麼會不知道呢?.
不過好像令到貼手來加城的弟兄姊妹.
更加感受到保羅這一份的愛.
和對他們的欣賞.
和內心放下那個心頭大石.
我想牧者是最開心.
真是見到教會是信望內成長.

$^{681}$信心,愛心都成長.
這裡提到信心,愛心.
最大的一個重要的資產.
就是人的心對準真理的主.
在主耶穌的感召下.
知道自己的身份.
知道自己現在所處的地位.
和知道自己將來要去哪裡.
而你努力為這個國度來效忠.
如果牧者能夠傳遞到這樣的真理.
這是最大的一個喜樂.
和應該說是一個回報.
因為為他們去報答神.
報答神的恩典.
就是神給保羅有生之年.
看到自己所做的有什麼成果.
就是貼手來加城的弟兄姊妹.
有信心和愛心.
自己都很多年了.
說一點點吧.
又不是什麼花樣.
在武堂很多年了.
離開了很多年.
有一位弟兄.
我讀神學的時候.
他可以無條件將他整間屋.
讓我在元朗住了四年讀神學.
然後一毛錢都不用給租.
讀神學我收入都沒有.
而且是一個過千呎的屋.
在元朗的.
當然什麼都不用簽.
什麼都不用說的「信」字.
對牧者的信任.
連太太也是很熟的.
都很恩賢的夫婦.
是同心的讓屋給我們住.
住了四年 租了四年.
完了之後就給他.
這麼多年了.

$^{721}$過去這麼多年已經二十多年了.
到今天.
偶然間他又知道在雄恩.
他又帶一些人來聚會.
不是追隨的問題.
而是他現在在現在的教會.
元朗的教會.
是很熱心的事奉.
很熱心的教主一行.
很認真的教主一行.
他現在也有一些年輕的夫婦.
跟他去查經.
這樣的一種回報.
是一個很大的喜樂.
是很開心的一件事.
有理所不欲的人.
看著他信主.
看著他回教會.
到現在他已經退休了.
他可以提早退休.
他現在將他的時間.
更好地去幫這些年輕的夫婦.
帶他們去查經.
很認真地.
有些讀書會.
大家一起去看一些神學的書等等.
所以是很感受到.
保羅這裡所說的.
那種回報的喜樂.
那種報答神.
偶然間他.
有時候我講到他又會下去.
帶一些人下去.
他的新朋友.
去看看不同的教會.
你就會想像到.
那個福音可以一直傳遞開去.
有這種延伸.
其實對牧者來說.
是最大的一個喜樂和滿足.

$^{761}$我絕對不是那麼明星.
不會有人這樣跟隨.
跟隨到當神一樣拜.
這些一定要拒絕.
我不是那麼有魅力.
有這種魅力.
但總是有一兩位.
他們真的看到牧者.
然後他真的有這種跟隨.
跟隨之餘.
是他自己現在做他神給他的.
那個崗位的工作.
這個是我自己.
最感到欣慰和滿足的地方.
看到有弟兄姊妹自己成長.
甚至跳出你對他的影響.
而他自己再度當一面.
這個是牧者必須要這樣做.
因為每個人都要過去.
但過去之餘.
而他能夠繼續他自己去有那個福事.
這個就是最寶貴.
是最重要的.
所以這裡說報答神.
盼望我們每一位弟兄姊妹.
在我們自己得到神的話語的教導之後.
我們自己又能夠獨當一面.
我們自己能夠融會貫通.
我們自己能夠將神的道.
去用不同我們自己的人生去承載.
去見證.
去使更多人得到神的話語.
這個就是報答神.
所以我們應該要找到一些.
我們的一些貼實羅尼斯加城的教會.
我們去用心盡心去牧養.
他們將來就成為了我們的喜樂.
還有因為他們的站穩.
為你們的成長我們去報答神呢?.
所以我們都應該要有.

$^{801}$如果我們一生都沒有一兩個人.
真的對我們有景仰,有效法.
甚至超過了我們的能力.
他們更加獨當一面.
某程度牧者真是一個很大的失職.
一兩個都沒有.
沒有理由的.
你說一百個當然不行.
除非是公民簽那些.
一開港道過千過萬人.
但是總是每個牧者一定有.
自己的牧羊的責任.
所以這裡是保羅很坦誠地分享.
不是要奪取他們的心.
而是要牧養他們的靈命.
使他們能夠知道牧者和他們的關係.
是這種牧羊人的付出.
捨命的關係.
所以他們對牧者的信任.
就帶給他們一個很強的信心.
愛心的成長.
所以他們說切切要想見保羅.
等於保羅也很想見他們.
這就是第十節.
所以我們晝夜的切切祈求.
要見你們的面.
補滿你們信心的不足.
已經有信心了.
保羅認為還可以再補滿.
可以再給他們有另外的信心.
即是說保羅其實對永恆.
對靈界.
他真的有很多神給他的啟示.
因為我還可以給一些你們從來不懂的東西.
令你們的信心再加增.
保羅沒有說得那麼詳細.
即是說保羅原來他那個啟示的領袖.
是隨著他傳道繼續增加.
所以去到提摩太前後書.
或者在哥林多後書.

$^{841}$保羅他得到的三層天的經歷.
那個見到異象.
去到一個地方.
是現在還沒有發生的.
因為寫這本書的時候.
他都還沒有在哥林多.
在雅典寫完書才去哥林多.
在哥林多一年半.
過了一年半才寫哥林多前後書.
你可以想像到原來保羅所得到的啟示.
是繼續增加的.
他去到羅馬坐牢的時候.
寫的這筆所書.
對天上的永恆的屬靈福氣.
他能夠如數加增地數出來.
和他的禱告等等.
超過了現在神學的難界.
即是說原來保羅這個使徒.
他們在神裡面繼續會有新的領受.
即是說我們不要小看基督教所謂信忘愛.
其實我們對永恆主的認識.
是應該無窮無盡地增加的.
是應該繼續在我們人生越成長.
就要對神又多了認識.
因為神根本是無限的.
我們是有限的.
所以我們應該讀經.
剛才聽到有些弟兄姊妹說讀經.
傳道書每次讀都有不同.
每個階段有不同.
事實是這樣的.
因為關乎那些哲理的書.
隨著我們人生的閱歷.
隨著我們有些東西看通了.
或者我們自己通了我們自己的人生.
不再執著了.
有些東西我們就會突然豁然開朗.
這些是事實的.
聖經都是那本的.
沒有改版的.

$^{881}$只不過是翻譯改了.
但是問題是我們的人生.
是應該不斷地有成長.
所以保羅說我們切切地見你們的面.
是要補滿你們信心的不足.
哪方面不足呢?.
當然他沒有詳細說.
哪方面的信心是不足呢?.
要補滿呢?.
但是保羅似乎就是.
總是有一些屬靈的可能是經驗.
可能是他的領袖.
都是很想來分享給弟兄姊妹知道.
又岔開說.
聽到真的不要選同工選牧者.
說得好與不好.
那個演繹不同是對的.
真的會有的.
的確的.
那個是牧者盡自己能力的責任.
不斷去完善.
不斷想方設法.
來講明白.
是那個責任來的.
是應該的.
應份的.
但是不同的牧者有不同的恩賜.
領袖,人生,際遇.
所以如果我們只是停留在一些.
好的出色的港元.
我們就憤外精神.
其實是那個憤外精神救了你的.
不是那個港元救了你的.
因為是你自己已經有準備.
我去見神了.
聽這位僕人教導我.
其實是你自己已經準備好.
到你知道原來曾蔭貴這次來港島.
你已經無精打采了.
你已經都是那些東西.

$^{921}$沒什麼興趣聽了.
儘管去敬拜神了.
你自己已經沒有進過禮堂.
你已經關了門了.
心門了.
當然好像開玩笑那樣講.
其實在在有無數的人.
可能那個聚會觀和聚會態度就是這樣.
所以這裡告訴我們.
保門那個信心的不足.
我都說我自己不斷去勉勵自己.
多麼的差.
所謂差是那個演繹.
除非他講錯了那個真理.
那個就真的差了.
就算很完美的演講的口才.
他如果講錯那個真理.
那個道本身就已經不要了.
根本就不可能被神用.
但是不好的.
講得一一咳嗽的.
但是那個道本身.
我們要學習.
一來學習忍耐.
他講得這麼難聽.
哇 睡著了.
那學習忍耐.
學習包容.
學習為牧者祈禱.
但是其實每一篇信息.
我自己勉勵自己.
告誡不是勉勵.
真的要從牧者的內心.
他這樣講這一篇信息.
我可以得到什麼樣的教導呢?.
要變成我們.
其實不是挑剔.
我可以從他整個主題的演繹裡面.
我得到什麼樣的教訓.
或者什麼呢?.

$^{961}$這個就是我們聽道裡面.
要用心和用神的地方.
就是這樣.
所以雖然大家忙碌裡面.
回教會就是想放鬆.
就是想聽一些又令你很好笑.
又發人心醒.
又令你很有提醒.
然後走出去就很開心的道.
是的 每個人都想的.
沒有人不想的.
但是如果不是切切.
都要學習包容.
總之就是我們自己都要.
重點反而是我們自己.
就是我可以從那個信息本身.
我的信心有沒有一個提升呢?.
或者完全沒有的.
完全沒有都是一種學習.
所以這個告訴我們.
保羅很想補滿他們信心的不足.
保羅還是有.
他不是特意藏起一些東西不說.
而是他自己都有新的經歷.
他自己對神都有新的認識.
他自己在經過了很多處境.
傳道.
對著真的不同的人.
對著雅典那些有知識的哲學家.
他是怎樣的方法去傳道.
對著猶太人.
對著一些尊貴的希臘婦女.
真是有頭有臉的.
他又怎樣傳道呢?.
你可以想像得到.
但是這些應變已經是很高的藝術了.
保羅又可以慢慢完全放下.
自己猶太人那些固有的偏見.
那些驕傲.
來將自己放在一個外邦人的社群來傳道.

$^{1001}$已經很不簡單了.
已經沒有那些所謂猶太人的優越了.
但是他自己在猶太人面前.
他又是完全的尊重猶太人的文化.
所以他為甚麼要幫提摩太行國禮.
就是他有這種靈敏.
心靈的敏銳.
怕猶太人被那些叔伯兄弟排擠.
所以就預先幫他.
未出去傳道已經幫他做了行國禮.
就是告訴人家.
你不用問他甚麼時候行國禮.
不用給他證明.
他就行了國禮就行了.
做了這個手續就行了.
你是不是第八天啊?.
當然不是了.
我二十多歲才行國禮.
但是都行了.
還有自己心甘情願.
傳道之前行國禮.
很不同的.
你看到阿保羅.
你看到他整個人.
他這個神學觀.
耶穌改變他.
他可以.
即是說他可以初原禁貶.
能夠向外邦人.
他就做外邦人.
向猶太人.
他的行為就變回猶太人.
所以他上到耶路撒冷.
第三次宣教旅程.
回到耶路撒冷之後.
他去到聖殿還原.
大家有沒有留意《聖行傳》.
22章那裡.
他去還原.
他有拿世人的緣.

$^{1041}$這個就是阿保羅.
所以他每不斷的人生都有成長.
我們很難勉強.
姐妹.
你有不斷的新的經歷.
因為大家的工作.
關注的事情.
都可能是家庭.
是生活.
是工作.
是事業.
很難要大家有這些.
這麼強的屬靈的領悟等等.
因為要大家看一些很深的神學書.
你根本都不會有興趣.
可能那些小品.
雞湯.
就很受歡迎.
書室是最好賣的書.
又簡單又薄.
字又大.
但是很艱難的神學書.
你一揭開就想睡覺.
還要探討哪個神學家.
他的神學觀.
這些是很複雜的東西.
但是很多時候.
隨著我們的人生.
神都會給我們在不同的人生階段.
都應該有所的領悟.
應該有所信心的提升.
所以真的不是一成不變的.
信望愛是真的會成長的.
如果我們已經自問在教會30年40年.
而我們的信望愛.
還停留在一個很嬰孩階段.
這個是個人的困難或者悲哀.
其實更加是教會的悲哀.
如果我們的長則已經停留在一種.
幼稚園階段的信望愛.

$^{1081}$那就很嚴重.
因為整個教會的眼界.
是受領袖的影響.
這些不是開玩笑的事.
你想想如果我們的特首.
好像一個嬰孩一樣.
很賴皮的.
那就很大件事了.
因為政策的制定.
是需要很多的智慧.
很多的商討.
要很通達.
要很識時局的分析和研判.
這些全部是需要智慧的.
這個不是說我們要特意歧視一些.
比較低智商的人.
不是這個意思.
而是因此上大家不同的崗位.
就需要有不同的對崗位的意圖.
所以聖經往往都告訴我們.
初入教的人不要做監督.
恐怕他會自高自大.
就會帶領教會走向一個極端.
往往是這樣的.
因為必須要有很多的經歷.
有很多人生的苦難.
更加磨練那個人.
純全的就是為主為福恩去捍衛.
去到第十一節.
「再講主再來」第三次講了.
這次是從哪個角度講主再來呢?.
「願神我們的父和我們的主耶穌.
一直引領我們到你們那裡去.
又願主叫你們彼此相愛的心.
並愛眾人的心都能增長充足」.
所以教會如果沒有愛.
很難繼續成長.
是一個亂七八糟的教會.
甚麼都來去爭風相對.
甚麼都來去爭一頓.

$^{1121}$很難教會能夠成長進步.
愛心能夠增長充足.
你可以想像是雙輔雙成的.
似乎保羅看管教會是在一個愛心裡面建立著.
所以苦難是外面的.
苦難原來對他們來說不是太難.
最好的地方就是裡面自己那班人有信心.
大家都是共同信主耶穌的那班人.
有那個的愛心.
所以「願你們彼此相愛的心.
愛眾人的心都能夠增長充足」.
愛眾人即是愛教愛.
教愛的人 教愛即是適合到Big Dad.
都能夠有這份顯出來的愛.
如同我們愛你們一樣.
第十三節就是講主在內.
所以原來教會在一個很順暢的環境.
就是那班人 一班共同有信念的那班人.
信仰 不是信念.
信仰那班人能夠彼此相愛.
彼此相愛的適合度就是可以包容.
可以寬恕 可以接納.
可以坦誠 這樣去成長.
所以往往教會帶來的.
不是Big Dad令到教會分化.
而是沒有了愛心的包容.
就會帶來教會的分化.
但是教會那個不能夠成長.
所以這裡也提到既然那個愛心.
當時羅馬帝國 這個很重要的.
我們讀神學一定要讀那個時代背景.
就是那些背景就是.
當時是有一個叫做私恩者的制度.
私恩者就是那些叫做Patreon.
Patreon現在那些YouTuber.
你自己開的專頁.
就要收會費 會員費.
就會有一些完整版的.
或者有些事先可以給會員可以聽的.
Patreon就是那個意思.

$^{1161}$就是私恩者 就是說我私恩給你.
你就是那個Client.
你就是那個 不是叫顧客.
叫受惠者 恩惠的惠.
這個制度就是保證了那個帝國的平穩.
避免那個過度的貧富懸殊.
其實是貧富懸殊的.
因為他通常收養的那些受惠者那些Client.
就是社會上比較窮苦的.
那個私恩者就會將一些恩惠.
例如錢 例如一些工作.
就會給那個Client.
那個Client就要有一個回饋.
就是假如那個私恩者.
他要競選做議員.
或者在某些社會上.
有些重要的可能是發言或者其他的東西.
那些Client就要出來支持.
你可以想像到這個私恩的制度.
其實是一種換取利益的制度.
某程度是帶來一個社會相對穩定.
就是派餅仔 吃齋餅麵.
就是類似這樣的社會的穩定.
是的 都對的.
政府找一些在鄉郊有一定影響的叔伯父老.
他們本身就已經有那個穩定性.
我聽二叔公說 比聽那個官說好.
所以那個官 那個李文虎或者那個政務官.
只需要跟那個鄉親父老談妥.
那就解決了很多問題.
那個叔伯兄弟 二叔公.
就會在祠堂裡面跟他的鄉民談那些政策.
所以這個是帶來那個穩定的.
所以帝國都容許這個私恩的制度的存在.
所以普羅根本就是在那個時代出現的人.
他自己本身.
我和大家看看羅馬書的16章.
以前我們講羅馬書的時候已經講過了.
如果大家手頭上有英文的羅馬書.
你就更加清楚普羅他講的意思了.

$^{1201}$就是這個菲比.
我對你們舉贊我們的姐妹菲比.
她是堅格里地中海一個很大的地方.
堅格里是一個島嶼.
堅格里教會的女執事.
請你們為主去接待她.
合乎聖徒的體統.
她在何事上要你們幫助.
你們就要幫助她.
留意一下這句話.
因她數來幫助許多人.
也幫助了我.
原來這個菲比是普羅的私恩主.
她是有恩於普羅的.
原來這個菲比是做了許多人的私恩主.
所以我寫英文的羅馬書16章.
稍等一會.
我看看NASB有沒有這個解釋.
羅馬書的16章第二第三節.
你要接受她.
在法律上.
在聖徒的體制上.
你要幫助她.
在任何方法上.
這裡沒有提到.
她所說的Hupper.
其實就是Patreon的字.
很多人的.
我自己也有.
我先去King James看看有沒有.
稍等一會.
King James.
這裡有很多告白.
大家先等一下.
有沒有?.
英文.
英文有.
King James.
不知道是不是King James.
稍等一會.

$^{1241}$羅馬書.
我先寫16章第二第三節.
這裡有.
這個古老字我也不懂得解釋.
Subcurra of many.
算了.
這個我們先不看了.
今晚我回去再翻譯給大家看.
卑彼就是在做這個私人主角色.
他是保羅的私人主.
保羅是從他那裡得到一些幫助.
所以那個幫助不是普通.
在主裡面接待他.
就不是說請他在會所.
喝一頓茶好一點.
吃得好一點.
記得叫一隻燒鵝.
不是這一種的接待.
是真的要看到他有恩於很多人.
這個就是卑彼.
所以這裡保羅說.
在教會裡面.
不應該只是禮尚往來的.
愛心的互相付出.
就是你有恩於我.
我就要回報你.
這是社會基本的羅馬帝國.
帶來社會穩定的一個制度.
保羅說教會要超過這個制度.
是真正基於神的救恩.
帶來的彼此相愛.
所以大家越是不同背景.
都要學習彼此相愛的心.
和愛眾人的心.
都要增長充足.
增長.
可能是數量的增長.
可能是質量的增長.
和要充足.
就是要充滿的.

$^{1281}$所以我們對人應該不要若即若離.
不要豐富的恩獻.
而是應該帶著一份的熱誠.
當然我們要有智慧.
不可能我們每個人.
三百人我們都能夠發揮到這一種的愛.
是有可能少少的不同的程度.
但是仍然在我們的內心.
我們應該對弟兄姊妹.
大家都是有共同信仰.
大家都是追求神的國度.
和大家一起傳福音.
我們不應該有一些門派.
你們是振信會.
我是宣誓會.
就不應該有這些所謂門派之見.
除非他們對真理有模糊.
如果不是的話.
應該超過社會基本的那個制度.
當時就是施恩主的制度.
就是說那個感情的建立.
是因為利益的輸送.
本來就不希望教會是利益輸送.
而是一種出於真誠的.
大家都是神兒女的建立的相交,相愛.
而且那個要求是要增長.
是要充足.
所以其實一生我們就是學習在信望愛上去成長.
是很激動的.
是一個不斷成長的生命.
我回望我三十多年前信主.
現在的我.
是真的那個程度不同了.
所以我們每個弟兄姊妹都應該.
在主的救恩裡面.
有我們自己那個成長的跑道.
這些經文是很簡單.
但是要進入裡面就很不容易.
要我們能夠彼此相愛和愛眾人.
都要來到去有那個成長.

$^{1321}$接著就是那個結果.
好使你們.
就是那個日子來的時候.
你們能夠無可責備.
一定是在說我們今生的信望愛.
好使甚麼?.
好使就是當主耶穌和他的眾聖徒來的時候.
就是說已經死了的那些聖徒.
是會復活的.
是會永活的.
他們不是神周十七日回來地球.
他們不過是住了幾個月在外太空.
現在說的是古往今來的聖徒.
你數的阿當,赫拉娃.
還有他的兒子.
瑟特,阿伯,羅亞.
有些我們根本不知道是甚麼人.
他們都會復活.
都會跟隨耶穌回來.
你可以想像到那個日子的榮耀.
是接我們.
整班聖徒是來接我們.
如果我們還沒有死.
如果我們死了.
就是我們一齊回來接將來的人.
所以你想一下是多麼雀躍的一個.
真是賞賜的日子.
今天的基督徒.
我不知道全世界.
我太有限了.
香港有多少的教會講主在內.
講到那些信徒是寶義設去傳福音.
來希望耶穌快點回來.
這就告訴我們第三次講了.
好使你們當我們的主耶穌和眾聖徒.
回來的時候.
在我們的父神面前.
在神的面前.
心裡面堅固.
心裡面即是信忘我愛.

$^{1361}$我們裡面的人堅固.
堅固不是頑固那些堅固.
不是的 怎樣都.
我沒有任何的分析或者數據.
是不是人越大了就越固執呢.
越是有些東西.
明明都未必是真理的.
就堅持到人家不是這樣.
你就受不了了.
會不會是這樣呢.
這些不是真是心裡堅固.
是心裡固執.
我們要留意就是.
這裡講的是真理的心靈的堅固.
即是對信心,愛心,盼望的堅固.
心裡面是堅固即是完全不動搖的.
這個人是跟耶穌的.
他死了都要跟隨的.
這樣的堅固成為聖潔.
在神面前無可責備.
在神的眼中.
我自問我是可以來到去.
得到神的光照.
我都是不會得到祂的責備的.
即是連我自己的良心都不會責備我.
這個意思.
你想一生我們有多少東西要追求呢.
有時我們後悔講了一些說話.
回不了頭的.
求神去赦免,原諒.
和自己想起都覺得羞恥的.
這些回不了頭的說話.
你就真是無可責備.
我們就要追求.
所以要尋求對方的饒恕.
我說的那句說話是很重的.
請你饒恕我吧.
去解除這種的內疚.
這些我們都要做的.
往往我們為面子就不講了.

$^{1401}$死都不講了.
忘記了它就算了.
其實是忘記而已.
都不是饒恕的.
洗乾淨.
洗乾淨和忘記是不同的.
因為忘記只是你的內心仍然有內疚感.
用寶器洗乾淨.
就沒有內疚感了.
神都已經饒恕了.
就不會再有這個追求.
所以這裡說.
心裡面堅固成為聖潔無可責備.
所以是說主耶穌來的時候.
我們的內心的信和愛是堅固的.
對於提撒羅尼加城的人.
保羅有這樣的對他們的要求.
是不是當時的保羅是很確信耶穌.
很快就會回來呢?.
有可能是的.
有可能是.
所以他就這麼熱切地每一章都講.
其實當時是沒有分章的.
他們希臘文的抄寫是不會有隔行的.
是不會有分段的.
是後人幫他分段的.
所以保羅講到某一個重點講完.
就好像加了最後一個祝福語.
就是主耶穌的回來.
接著就進入第二個話題.
第四章.
你可以想像到保羅對於主在內.
可以改變一切是看得很重的.
譬如我們看看《譚安第後書》四章.
就是保羅最後的.
你可以說是他人生最後階段的.
對於主在內的那個期盼和渴求.
四章的第七節這一段.
第六,第七,第八節.
「我現在被囂張.

$^{1441}$我離世的時候到了」.
剛才其實羅尼加城是他第二次宣教旅程.
還是很有精力地到處去的.
但現在他已經被所拘捕在羅馬的監獄裡面.
應該都是一個頗高度設防的羅馬監獄.
「我離世的時候到了.
那美好的將我已經打過了.
當跑的路我已經跑盡了」.
全部是一種過去完成.
即是說已經過去了.
我已經做完了.
「所信的道我已經守住了.
從此以後有公義的官臉為我存留.
就是按著公義審判的主.
到了那日要賜給我的」.
這裡又多說了一點.
「公義的官臉賜給保羅」.
不單止賜給我的.
第八節這裡.
「也賜給凡是愛慕他顯現的人」.
所以這次又多了一個心情.
原來主耶穌是回來.
不是硬生生的審判.
而是我們等待一位愛我們.
為我們犧牲的耶穌.
愛慕他顯現.
愛慕.
因為我都完全對於追明星.
我都完全過時代了.
無論舊的新的我都不懂.
也都沒有那個mood.
如果你還是很熱衷看演唱會的.
可能隨著時代.
比如我們的院長梁家倫院長.
很喜歡林子祥.
他就會和太太去看林子祥的演唱會.
因為是那個時代.
新一些的 比如五月天那些他不懂的.
你有這個愛慕.
如果有人給你一張很貴價的.

$^{1481}$可以握手的門票.
你去看.
你就是很喜歡的那個明星.
你想想他現在開演唱會.
那張可以握手的門票.
是人家送給你的.
你一定要去看.
甚至你轉場的時候.
你可以去後台去探班.
你想想那個愛慕那天出現.
你是多麼的日於言表.
這些我們尚且都能夠去了解.
是我們可以去投入去同情和去等待.
所以愛慕的意思.
保羅說到的他自己的離世.
他是說到自己愛慕主耶穌的翻臉.
他將來要得到那個公義的冠冕.
所以你看到他其實是很苦的.
但是苦裡面他看到榮耀的盼望.
那些苦就完全失色.
他看到的是榮耀的賞賜頒獎的日子.
叫保羅來領取那個公義冠冕.
那種榮耀不是要裝模作樣.
不是要驕傲.
而是他要得到那個.
在所有的聖徒面前.
他走去主耶穌那裡得到那個公義冠冕.
所以這些是我們一生頂尖有沒有追求呢?.
最近星期日報經常說娛樂界的明星.
說最近有一個負三代的名媛.
說她去旅行.
那些行李搭得那麼高.
好像整個手推車.
不是那些.
是說那些好像酒店的地下大堂那些.
大的車裝滿行李.
是她一個人.
應該是這兩天的星島.
你回去看看.
詳細的介紹.

$^{1521}$我們不是要追求這些榮耀.
我們要追求主耶穌回來的時候.
我們愛無他顯然的人.
愛無他顯然.
我們講不講得上我們愛無呢?.
還是逼於無奈他來就來.
我都沒辦法了.
我都不是很稀罕了.
我現在這麼開心在這個地上.
當然我們有一天我們就不開心了.
因為身體又一路一路走下坡了.
走都走不了了.
你想上華山想都不用想了.
你想走高一點.
還要爬那些天梯根本不用想了.
因為已經過了我們那個人生的時間了.
不要說那些平地.
那些都沒有精力去走.
我們都一定會去到這些日子的.
我們會走下坡的.
但是愛無他顯然.
愛無.
應該是不是如日俱增的那種心情的愛無呢?.
這個是值得我們基督教要去想.
今天我們講的基督教到底是一個怎樣的宗教呢?.
是不是一個很開心很開心.
唱的歌又是很完美的.
很興奮的.
那個情感上.
每次都有一種自如的.
這個不是錯的.
不過長期都只是停留在這些階段.
沒有跳出一個好像舒適圈那樣.
一提到苦難呢.
不要提了.
基督徒不應該有苦難的.
那不是聖經不得平衡.
平衡的就是既有我們情感的滿足.
對的絕對沒有錯的.
但是如果每次都是情感的滿足呢.

$^{1561}$那我們怎樣能夠經受起那些外在的苦難呢?.
還有教會是不是只是靠這些好聽的音樂去維繫呢?.
而是不是真的會鼎鼎相愛的心靈.
是真誠的這樣大家成長呢?.
這樣去建立那個教會呢?.
這個就是帖索羅尼加的前書.
給我們看到的一個框架.
不過人家是信主四個星期.
一個月多一點.
本來就要離開.
那些剛剛初信班都還沒有上完.
就已經沒有了牧者了.
就已經是外面的Big Brother進入來教會了.
但是人家都已經有這種信望愛的成長了.
所以這個是保羅對頂子妹的那個勉勵要求.
從來沒有減低做基督徒要受苦的這個真理.
今天我們不是要去一個受苦難的宗教.
或者一個虐待人的宗教.
但是在地上我們相信一位獨一真神.
這個本身就是一個矛盾的了.
世界是不容許我們講獨一真神.
由主耶穌回來審判世人.
這樣的一個真理.
所以只會受到排擠和攻擊.
盼望我們不要減低了福音對我們的要求.
所以如果我們有頂子妹.
他們要去到極權的國家傳道.
我們應該很清楚她真是有神的帶領.
我們就要送她一程.
真是要祝福她.
真是要一路支持她.
讓她能夠將福音傳到那個地方.
我有同班同學,夫婦.
他們很厲害的.
我不可以說太多.
因為始終公開在極權的國家周邊.
做著福音工作.
有一定的生命危險.
不可以跟他們WhatsApp.
不可以說太多教會的事.

$^{1601}$因為向著極權的地方傳道.
所以真是我們需要支持和需要祈禱.
但是說的是全基督教苦難是必然的.
是命定的.
是一定會有的.
所以我們在教會裡面.
我們由成長.
如果你是從小朋友開始的.
你一路長大.
如果教會都沒有告訴你.
信耶穌是可能要受苦的.
要為主耶穌作見證的.
如果都沒有的.
可能他一直在聽的一個不是很純正的宗教信仰.
我很深刻有一個姐妹.
在很多年前.
小時候我們都一輩子的.
她回到港大讀書.
整班老師.
她其實讀數學而已.
不是讀哲學那些.
她分享說那些老師.
很不喜歡基督教的.
趁她起床.
在眾人面前質問她.
很難受.
很難堪.
詆毀基督教.
她那一刻沒有甚麼可說的.
因為她不能辯駁.
但是她感受到福音那種的受苦.
其實是一個榮耀.
在眾人幾百人面前.
笑你們的基督教.
現在的現代社會應該很包容.
也不是.
都是這樣嘩虎嘲笑你們的基督教.
故意挑些難題.
但是這個小姐妹.
那時候她還沒有結婚.

$^{1641}$讀大學的時候.
就感受到福音帶給她.
原來有這種要受苦的.
其實不是打她.
只是用一些說話.
令她難堪.
要為福音來受這些屈辱.
保羅說我們在肥牛比被害受辱.
保羅說我們告訴你們.
受患難是命定的.
命定是命中注定.
我們必然會有的.
後來果然應驗.
保羅的意思是.
真的我們傳的福音.
是伴隨著貝什自駕.
跟隨耶穌.
就是說它是沒有減少的.
今天我們盡量希望.
令基督教去苦難化.
得出的結果就是.
基督教我信也行.
不信也行.
因為原來苦難是不存在的.
信就不用受苦難.
那我用不著信你.
我信自己就行.
因為我都可以解決我自己的問題.
所以往往.
基督教最令人感到.
欣慰的地方就是.
無論我現在怎樣都好.
這個信仰.
這位主耶穌從來沒有怠惰過我.
我怎樣都要跟隨祂.
這個就是很含金量高的信仰.
希望我們能夠在我們有生之年.
都能夠為主耶穌去責問.
我們的信心可以增長和充足.
第四章我們可不可以再看多一點.

$^{1681}$現在都是49分.
我們稍為補讀.
不如我們先看第13節.
頭一段我們先不看.
因為有些當時的處境.
我們下一個星期再詳細看.
第13節就很詳細講主在來的次序.
我們現在在安息聚會.
去到火化禮我們都會讀這段經文.
「論到垂了的人.
我們不願意大興門不知道.
恐怕你們憂傷.
好像那些沒有指望的人一樣」.
這裡提到指望.
「我們如果相信耶穌死而復活.
那些已經在耶穌裡面垂了的人.
神也必將他與耶穌大同.
為何要這麼卑鄙不說死呢?.
要說垂呢?.
似乎主耶穌都是這樣說.
或者猶太人他們也有這樣的傳統.
大家都知道垂就是已經死了.
我們中國人這麼卑鄙都是一樣.
但是後面15節.
「我們現在照主的話告訴你們一件事.
我們者活著還存留到主降臨的人.
活生生眼見到耶穌回來的人.
將來總有一代人會見到的.
斷斷不能夠在那些已經垂了的人之先」.
即是甚麼意思呢?.
第16節就解釋了.
「因為主必親自從天降臨」.
特別強調他親自.
即是不需要經代理人.
他親自來.
有呼叫的聲音到底是甚麼聲音呢?.
是主耶穌呼叫?.
還是整個諸天來呼叫?.
因為後面就說第二個層次.
就是天使掌的聲音.

$^{1721}$天使掌即是加巴列.
如果打仗的就是米加勒.
掌管軍隊的.
傳遞訊息的就是加巴列.
還有就是四門密.
還有十二位長老在聖路基督.
哪個天使掌呢?.
是否眾天使掌的聲音呢?.
我們要回去看原文才知道.
又有神的號吹響.
保羅說的天使掌是誰呢?.
接著有神的號吹響.
即是眾天使吹號.
「那在基督裡面死了的人.
必定首先服他」.
開始說那個次序.
在基督裡面死了的人.
即是已經死的聖徒.
就必定首先復活.
即是說有一個復活的身體.
「以後我們這活著還傳流的人」.
即是只有兩批.
一批就是在主裡面已經死了的.
一批就是還活著.
活生生在地球的.
有生命氣息的在地球的人.
信耶穌的.
那就不會復活了.
那些死了的人先復活.
接著呢.
「以後我們這活著還傳流的人.
必和他們」.
就是那些復活了的聖徒.
「一同被提到雲裡.
在空中與主相遇.
這樣我們」.
我們就是所有無論死了.
仍然活著的.
所有都復活了.
所有都去到天上了.

$^{1761}$所有都和主永遠同在.
我們就要和主永遠同在.
「所以你們當用這些話彼此勸諱」.
勸諱就是因為當時他們有憂傷.
有對人死不能復生的一個懼怕.
就是他們所愛所寶貴的家人.
親人.
他們受迫害.
可能是患病而死了.
這些人怎樣呢?.
在主耶穌回來的時候.
這些人怎樣呢?.
會不會叫他們的名字復活呢?.
保羅說會.
而且是首先復活.
那就安慰了歷世歷代的聖徒.
原來我們在主耶穌裡面.
安睡,安然的離世.
是不需要害怕的.
你沒有了你的名字流傳後世.
沒有了你的墳墓.
沒有了你的骨灰.
Nothing了 沒了.
真的在歷史裡面沒有你的名字.
但有人在.
1517年我們知道的就是馬丁路德.
但是馬丁路德身邊的人.
我們一個都不知道的.
只知道這些重要的人物.
但是五百年裡面.
有很多的聖徒.
沒有人認識的了.
但是所有的人都要復活.
保羅想說的這個點.
也是當時可能有些異端.
就來到危言聳聽.
不要跟著耶穌.
如果你死了你就什麼的了.
就帶來很大的一個巨砲.
就是說初期教會他們很擔心.

$^{1801}$那個逼白之後的去處是哪裡呢?.
那時候肯定就沒有天主教說的煉獄.
沒有那個發明的.
是後來很多年後才有的.
就是那些主張.
所以主耶穌的教導.
《路加分第十六章》說的.
那個拉撒路和財主的那個比喻.
一個是樂園一個是陰間.
是猶太人所相信.
基督教亦都相信這樣的教導.
但問題就是說.
在主耶穌回來的時候.
那些已經死了的聖徒.
是不是會復活呢?.
這個是那個困惑.
這個困惑到底是怎樣形成的呢?.
保羅就沒有詳細說.
這個是牧養這一班弟兄姊妹在生的.
他的親人已經離世了.
他們現在在哪裡呢?.
將來他們會怎樣呢?.
這個就是保羅很想來這裡說的.
所以真是一個很大的安慰.
在生無論是什麼光景.
因為有些人很慘的.
我們做牧者真是有一個很慘.
見到很多人很慘的.
我自己很心虛的突然間閃到腦海裡面.
我也說過的 在港島.
有一位在博愛醫院的一位.
我不記得名字了.
今晚我給他一個名字.
他是無依無靠的.
他有一天透過醫院打來教會.
找了負責人找我 讓我聽.
那個護士說我們有一個病人想洗禮.
就馬上找我去他醫院.
要我穿上全套的保護衣服.
那個時候不是SARS 不是新冠.

$^{1841}$是好幾年前 新冠之前.
原來要我穿上是因為怕我傳染給他.
不是他傳染給我.
他虛弱到這樣.
虛弱到這樣的程度.
但是很簡單說一兩句.
他真的很清楚 順序和他洗禮.
接著就離世了.
過了一個禮拜多一點就離世了.
老人院就再找我幫他辦安息禮拜.
是沒有親人 一個都沒有.
是孤苦無依的.
這位女士.
幫他辦了安息禮拜.
最後都拿了骨灰等等.
幫他處理好所有的東西.
很多這一類的人有時.
但是你要記住.
他將來復活.
他是完全正常的人.
有一個榮耀的身體.
所以這是我們基督徒的盼望.
是很寶貴的一件事.
所以將來所有的聖徒復活.
原來又不需要擔心人滿知患.
因為靈體是沒有這東西的.
不只是靈魂.
不是像我們經常怕的鬼魂.
不是那些.
是靈體是怎樣的呢?.
古往今來這麼多.
都沒有不認識的人.
不會說你穿古裝.
你穿唐裝都不會.
個個都是榮耀的身體.
個個都是將來在耶穌帶回來的人.
全部復活.
所以這裡就是給那些.
仍然對於真理有擔心的弟兄姊妹.
一個的安慰.

$^{1881}$所以我們去做安息禮拜.
幫很多很多人辦理安息禮拜.
其實每一次都感受到一份榮耀.
因為將來他就要在主裡面復活.
這是聖經的啟示.
和一個可以彼此勸慰.
沒有人想人生是坎坷的.
沒有人想人生是孤苦無依.
孤苦伶仃的.
甚至去到人生階段.
中國人說生榮死哀.
生前古榮耀.
很多人圍繞在身邊.
死的時候是很哀.
哀即是很多人為你哭.
是很大的榮耀.
但是其實更多人死了都沒有人知道.
但是你不要管.
他在主耶穌裡面.
他有那個福音.
他有主的那個恩典.
將來就會榮耀的復活.
還要是比我們那些活著.
在地上的人首先復活.
所以我們當然要了解一下.
為甚麼保羅突然間會講這一段.
講得這麼詳細的主在來.
和那個次序.
今次講的主在來是講那些信徒的復活的次序.
所以這個也是值得我們去想的.
譬如你想想在梅州五一.
開車回福建.
沒有想過就掉進了山坑.
然後火災.
死不了.
48個人突然離世.
你都沒有想像.
之前又沒有預兆.
那些三泥傾瀉.
整條路倒塌了.

$^{1921}$這些都不是我們可以控制的人生的命運.
長本在上帝的手中.
但是無論如何就是再一次提到.
主在來帶給我們人類的安慰.
好 今晚我們就到這裡.
多謝天父給我們有機會去看.
貼撒羅尼加前書.
和我們將來講的關於末世論.
不同的解釋.
幫助我們能夠對於主在來.
對於我們愛慕主顯然的基督徒.
我們帶著一份渴望.
幫助我們每一位在現場聽的.
或者遲些可以透過聽錄音聽到的訊息.
其實都是聖經很清楚的教導.
總是幫助我們大家建立一份對主在來的渴慕.
對於我們自己現在裝備我們的信心.
愛心,盼望的熱切.
我們真是來到有一個這樣的渴慕的心靈.
願神幫助我們.
多謝你聽我們感恩禱告.
奉耶穌基督的名 阿們.
阿們.
\newpage



\section{創世記 39:1-6-19-23}
\label{sec:U5CngR8W_hQ}
\textbf{2024-05-04 晚堂《變幻人生主同行》創世記 39\_1-6,19-23  講員:陳麗蟬牧師}
\newline
\newline
連結: \href{https://youtube.com/watch?v=U5CngR8W-hQ}{\texttt{https://youtube.com/watch?v=U5CngR8W-hQ}} ~~~~ 語音日期: 2024-05-07
\newline
\newline
\hyperref[sec:sCNFnlmgxC0]{\small{< < < PREV SERMON < < <}}
~
\hyperref[sec:index_chronic]{\small{[返順時目]}}
~
\hyperref[sec:index_scriptual]{\small{[返順卷目]}}
~
\hyperref[sec:ltEd3gYtiPs]{\small{> > > NEXT SERMON > > >}}
\newline
\newline
創世記 39:1-6-19-23
\newline
\begin{longtable}{cl}
\hline
\hline
章節 & 經文 (和合本修訂版)\\
\hline
39:1 & \begin{tabularx}{0.7\textwidth}{X} 約瑟被帶下埃及去。有一個埃及人波提乏，是法老的官員，是護衛長，他從那些帶約瑟下來的以實瑪利人手中把約瑟買了去。 \end{tabularx} \\ \\ \relax
39:2 & \begin{tabularx}{0.7\textwidth}{X} 約瑟在他埃及主人的家中，耶和華與他同在，他是一個通達的人。 \end{tabularx} \\ \\ \relax
39:3 & \begin{tabularx}{0.7\textwidth}{X} 他主人見耶和華與他同在，又見耶和華使他手裡所辦的事都順利， \end{tabularx} \\ \\ \relax
39:4 & \begin{tabularx}{0.7\textwidth}{X} 約瑟就在主人眼前蒙恩，伺候他主人，主人派他管理家務，把一切所有的都交在他手裡。 \end{tabularx} \\ \\ \relax
39:5 & \begin{tabularx}{0.7\textwidth}{X} 自從主人派約瑟管理家務和他一切所有的，耶和華就因約瑟的緣故賜福給那埃及人的家；凡家裡和田間一切所有的，都蒙耶和華賜福。 \end{tabularx} \\ \\ \relax
39:6 & \begin{tabularx}{0.7\textwidth}{X} 波提乏把他一切所有的都交在約瑟手中，除了自己所吃的食物，其他的事一概不知。約瑟英俊健美。 \end{tabularx} \\ \\ \relax
39:7 & \begin{tabularx}{0.7\textwidth}{X} 這些事以後，約瑟主人的妻子以目送情給約瑟，說：「你與我同寢吧！」 \end{tabularx} \\ \\ \relax
39:8 & \begin{tabularx}{0.7\textwidth}{X} 約瑟拒絕，對他主人的妻子說：「看哪，一切家務我主人一概不知，他把所有的都交在我手裡。 \end{tabularx} \\ \\ \relax
39:9 & \begin{tabularx}{0.7\textwidth}{X} 在這家裡沒有人比我更大，除你以外，他也沒有留下一樣不交給我，因為你是他的妻子。我怎能行這麼大的惡，得罪神呢？」 \end{tabularx} \\ \\ \relax
39:10 & \begin{tabularx}{0.7\textwidth}{X} 她天天這樣對約瑟說，約瑟卻不聽從她，不與她同寢，也不和她在一起。 \end{tabularx} \\ \\ \relax
39:11 & \begin{tabularx}{0.7\textwidth}{X} 有一天，約瑟進屋裡去辦事，家裡沒有一個人在那屋子裡， \end{tabularx} \\ \\ \relax
39:12 & \begin{tabularx}{0.7\textwidth}{X} 婦人就拉住他的衣服，說：「你與我同寢吧！」約瑟把衣服留在她手裡，逃出外面去了。 \end{tabularx} \\ \\ \relax
39:13 & \begin{tabularx}{0.7\textwidth}{X} 婦人看見約瑟把衣服留在她手裡逃到外面， \end{tabularx} \\ \\ \relax
39:14 & \begin{tabularx}{0.7\textwidth}{X} 就叫了家裡的人來，對他們說：「看，他帶了一個希伯來人到我們這裡戲弄我們。他到我這裡來，要與我同寢，我就大聲喊叫。 \end{tabularx} \\ \\ \relax
39:15 & \begin{tabularx}{0.7\textwidth}{X} 他聽見我放聲大喊，就把他的衣服留在我這裡，逃出外面去了。」 \end{tabularx} \\ \\ \relax
39:16 & \begin{tabularx}{0.7\textwidth}{X} 婦人把約瑟的衣服放在身邊，直到他主人回家， \end{tabularx} \\ \\ \relax
39:17 & \begin{tabularx}{0.7\textwidth}{X} 就用這樣的話對他說：「你帶到我們這裡來的那希伯來僕人進來要調戲我， \end{tabularx} \\ \\ \relax
39:18 & \begin{tabularx}{0.7\textwidth}{X} 我放聲大喊，他就把衣服留在我身邊，逃到外面。」 \end{tabularx} \\ \\ \relax
39:19 & \begin{tabularx}{0.7\textwidth}{X} 主人聽見他妻子對他說的話，說：「你的僕人就是這樣對待我」，就非常生氣。 \end{tabularx} \\ \\ \relax
39:20 & \begin{tabularx}{0.7\textwidth}{X} 約瑟的主人把他抓起來，關在監獄裡，就是王的囚犯被關的地方。於是約瑟在那裡坐牢。 \end{tabularx} \\ \\ \relax
39:21 & \begin{tabularx}{0.7\textwidth}{X} 但耶和華與約瑟同在，向他施恩，使他在監獄長的眼前蒙恩。 \end{tabularx} \\ \\ \relax
39:22 & \begin{tabularx}{0.7\textwidth}{X} 監獄長就把監獄裡所有的囚犯都交在約瑟手下；在那裡的一切事都由他處理。 \end{tabularx} \\ \\ \relax
39:23 & \begin{tabularx}{0.7\textwidth}{X} 任何交在約瑟手中的事，監獄長一概不察，因為耶和華與約瑟同在，耶和華使他所做的都順利。 \end{tabularx} \\ \\
[1ex]
\hline
\hline
\end{longtable}
$^{1}$各位兄弟姐妹平安.
開始訊息的時候.
我想大家聽一聽首歌.
可能大家都很回味的.
也都很有同感的.
記不記得什麼歌呢.
《知否世事常變》.
(歌詞).
當王瞻寫這首歌的時候.
我想他很經歷到人生.
甚至你和我都會經歷到.
人生是會不斷地變的.
所以中國人在那些很古老的學說裡面說.
世上唯一一變的東西是什麼呢.
就是世事常變.
聽清楚了沒有.
世界上的事是不斷地變的.
世界上唯一不會變的.
就是世事常變.
有時候想起來都很驚訝.
不知道會變成怎樣的整個世界.
今天是5月4號.
原本在這個時間.
我已經在我的辦公室裡.
把我的書和行李收拾好.
回家了.
準備開始我的人生的另一個階段.
就是退休了.
很輕鬆.
但是因著天父的恩典.
又讓我在這裡繼續承擔一個大使命.
一個很大的變化.
我做夢都沒有想過.
還有大家很熟悉的孫偉雄弟兄.
他一向做事很穩陣的.
有什麼危險動作不小心的.
他一定會指著我們大罵的.
小心一點的.
誰知道想都想不到.
今年年初.

$^{41}$他在工作上爬高做事.
竟然不知道是扇腳還是怎樣.
就一下子從兩米高跌下來.
結果住了幾個月醫院.
可能還有一段時間住的.
這些事情是做夢都想不到的.
有些病人的家人甚至說.
醫生明明說家人正在好.
回家吧.
但是當家人回家的時候.
洗完澡.
接著醫院就會打電話來.
你的家人停止呼吸了.
快點回來.
有時候我們真是.
這一刻鐘.
我們都不知道下一刻鐘.
下一刻鐘會發生什麼事.
約瑟的一生就是這樣了.
他就是一個.
讓我們看到.
由這個37章.
一直到創世記37章.
一直到50章的時候.
就展現了這個變幻的人生.
這個波浪有起有跌.
有幸有不幸.
但從約瑟的生平裡面.
我們看到.
在他很多不幸的遭遇裡面.
又有上帝很多的恩典.
上兩個星期我都講到了.
他原本在家裡是天之驕子.
他是奴隸.
但因為哥哥的妒忌.
就把他賣給了一些商人.
就做了奴隸.
但在做奴隸的裡面.
他竟然做得很有成就.
管理了所有的事情.

$^{81}$接著又因為他被主人的太太.
去引誘他.
他不接受這個引誘.
就去坐牢.
在坐牢的裡面.
蒙神的恩典做得很好.
監獄長就把事情都給他去打理.
也都因著這個成功的解夢.
他又做了埃及的宰相.
看到約瑟的一生.
又起又跌.
我們來看一下.
逐樣東西我們來看一下.
他在家裡是天之驕子.
爸爸就送了一件很漂亮的衣服給他穿.
是長袖的袍.
我們看到原來他什麼都不用做的.
在家裡.
哥哥就做很多粗重的事情.
看養.
但是他就是十指不沾陽春水.
什麼都不做.
但是當他被賣到埃及的時候.
就去到一個大官叫波提弗的家裡.
做奴隸.
做奴隸很慘的.
不止今天.
當天什麼粗重的功夫.
都要去做.
所以說很辛苦.
對他來說.
那個適應都要很大.
天天做苦工.
但是他竟然在做奴隸的期間.
恩著耶和華的同在.
使他所做的事情.
都樣樣順利.
他竟然被奴隸的主人.
一個大官波提弗.
他從約瑟的身上.

$^{121}$他看到耶和華與他同在.
他看到約瑟是一個通達的人.
所辦的事情樣樣都順利.
什麼叫做通達的人呢?.
這個約瑟是一個怎樣通達的呢?.
在《箴言》裡面說到.
他心中有智慧.
必稱為通達人.
這個通達人所說的話是怎樣的呢?.
嘴津的甜言.
加憎人的學問.
這個通達人是說到.
他與人的關係.
他心中有智慧.
這個智慧不單單只是讀那些學識裡面.
讀那些現在你沒有中文科,英文科那些學識.
而是說到他.
是懂得怎樣與人相處.
他有一個辨別的能力.
就不會與人相處的時候.
時時板起臉來指責人,罵人.
而是懂得在與人相處的裡面.
無論是鼓勵人或提醒人的時候.
就會怎樣呢?.
嘴中的甜言.
聖經就不是說到他嘴甜舌滑.
是不是這樣?.
嘴甜舌滑.
不是那些東西.
而是說的話.
他懂得積極的鼓勵人,建立人的.
以致人的生命大有進步.
因為聖經裡面說到這個甜言.
可以說是人的一些潤滑劑.
以致他與人說話的時候.
就算人家做得不好也好.
他怎樣可以用一些合宜的語氣.
合宜的言詞來提醒人.
給人鼓勵去建立人.
這個就是這個約瑟.

$^{161}$在做奴隸的裡面.
被波提佛看到.
他是一個通達人.
所以為什麼在這個大觀.
這個波提佛的家裡面.
他可以被這個波提佛.
將所有東西交給他來管理.
就懂得處理人的關係.
我們中國人裡面.
我們最尊敬的孔子.
他的其中一個弟子.
就叫做顏回.
有一次顏回在街市裡面與人吵架.
與一個賣布的商人吵架.
這個賣布的商人說.
8X3=23.
想不到那時候就已經有乘數表.
所以我不知道那時候.
還是人家說故事.
8X3=23.
但是顏回就很清楚.
8X3=24.
所以兩個人在街市吵架.
最後就找了最尊敬的老師.
孔子去理論.
回到那裡就告訴他們的老師.
孔子聽了.
他說我們在這裡輸了.
顏回就說.
8X3=24.
如果我輸了的時候.
我就將我頭頂帽子送給你.
送給賣布的商人.
賣布的商人通常吵架是串鬥的.
我說是8X3=23.
如果我輸了的時候.
我就將我的頭頂帽子送給你.
孔子聽了之後就哈哈大笑了.
然後他說.
顏回你輸了.

$^{201}$顏回心裡咕咕叫.
8X3=23.
老師已經腦矇洞了.
但是因為那時候那些人很尊敬老師.
老師說一不敢他們不敢說二.
所以沒辦法.
就將他頭頂帽子送給賣布的商人.
但是這個顏回就很不服氣.
就覺得這個老師已經腦矇洞了.
以後不再跟他學東西.
就去收拾包袱準備走人.
走人的時候就去跟老師說再見.
孔子就跟他說.
顏回你是不是很生氣呢.
他說是啊.
8X3明明是24的.
怎麼會是23呢.
孔子就告訴他.
如果你輸了.
你只是將頂帽子送給他.
但如果他輸了的時候就大件事了.
他就要將他的頭割下來.
給了你.
那你說顏回.
你說頂帽子重要一點.
還是人的頭重要一點呢.
所以孔子的學生就明白了那個大道理.
就是說原來老師看重的.
就是那個事情的輕重.
輸了頂帽子而已.
但是人命的重要.
所以覺得原來老師看重的.
是那個事情的輕重.
而人永遠就是爭一口氣.
所以顏回就明白了.
不過雖然我看完這個故事.
我都有點不太順氣.
以後賣布的商人.
永遠都算錯了那又怎樣呢.
不過還是命重要一點.

$^{241}$不過其實無論這個通達人.
裡面說的人際的關係.
怎麼懂得去建立人.
鼓勵人.
或者孔子顏回這個故事.
給我們看到一樣東西.
這個人際.
這個人情.
會比一些業績.
會比一些經濟.
會比這個憲制更加的重要.
耶穌基督在馬太方都這樣的教導.
當你在祭壇上去獻這個祭物的時候.
如果你想起這個有弟兄對你有懷恨的.
就要把祭物留在祭壇前.
先去跟弟兄和好.
然後才去獻上這個祭物.
所以耶穌基督說.
要去看重的不是我們獻的祭物.
而是看我們跟人的關係.
有沒有一個和好的關係.
我們有沒有來得罪人.
看回這個藥室.
它在埃及這個地方.
一個拜偶像的地方.
它的主人.
這個波提佛.
是埃及裡面的一個大官.
他也是拜偶像的.
但當這個藥室.
被賣到這個大官波提佛裡面.
做這個奴隸.
他沒有自怨自艾.
相反在這個奴隸群裡面.
他怎樣做一個通達人.
我相信在這班奴隸裡面.
不要每個人都在叫你窮啊 頂硬啊.
而是怎樣在做奴隸的當中.
怎樣去鼓勵他們積極.
以致這個波提佛.

$^{281}$就將這班奴隸的管理權交給他.
甚至他所有的事務.
交給這個藥室管理.
波提佛就是見到這個藥室.
和那些人.
和他所辦的事情.
都非常的順利.
波提佛對這個藥室.
信任去到一個什麼程度呢.
這個波提佛就將他一切所有的.
都交在藥室的手裡.
這個波提佛嘆到是怎樣.
除了自己所吃的.
知道自己那天吃雞還是吃魚還是吃蝦之外.
其他的事可以一概不知.
他完全放心藥室管理的.
因為藥室是一個通達人.
有耶和華的同在.
耶和華大大賜福給他.
我們繼續看藥室.
這個耶和華賜給藥室.
不單是他做事很成功.
在第六節裡面.
他更加講到.
又賜給他英俊健美.
也就是說這個藥室很英俊.
很強壯.
不過這個英俊.
這個強壯.
就帶給藥室一個另外一個很大的考驗.
就是這個波提佛的太太.
這個老闆娘.
聖經說她日日以目全情.
向他斬眉弄眼.
然後又叫藥室跟她一起睡覺.
甚至有一天.
這個老闆娘趁著家裡沒有人.
不過其實學者說.
其實不是趁著家裡沒有人.
而是把家裡的人都調走.

$^{321}$因為她是大官.
家裡除了很多奴隸之外.
還有很多婢女很多人.
所以家裡裡面不可能沒有人.
所以學者相信.
是這個老闆娘把家裡的家僕婢女都調走.
然後趁著藥室進入家裡做事的時候.
她就把衣服脫下來.
跟我一起睡覺.
想一想.
藥室當時只有二十多歲.
在一個血氣方剛的時候.
她怎樣可以受得了這麼大的引誘呢?.
還有只要她肯就範的時候.
可能她就會連升幾級了.
有人說男人最大的成功是什麼呢?.
勝過美色的引誘.
還要說她一定要堅持住.
在《箴言》裡面所說的.
就是時時都記得.
要喜悅幼年的時候所娶的妻.
一定要很堅持喜歡自己.
第一次去娶的太太.
為什麼藥室可以能夠勝過這麼大的誘惑.
這麼大的引誘呢?.
她對老闆娘這樣說.
第九節.
她說你是老闆的妻子.
我怎樣可以行者大惡得罪上帝呢?.
她堅持一件事.
她知道這件事是上帝很憎恨的.
這個奸淫在十誡裡面是其中的一誡.
不可奸淫.
她知道上帝的生意.
其實在今天我們都應該很清楚知道.
上帝要我們遠避淫行.
倘若什麼時候有任何誘惑出現的時候.
這個遠避就要我們有這麼遠.
走這麼遠.
千萬不要被魔鬼留地步.

$^{361}$所以老闆娘拉著她的衣服.
她寧願不知道怎樣脫下了衣服.
就扔下了衣服給老闆娘.
除了她也好.
她也要快快走離這間屋.
結果老闆娘就在波提佛裡面.
去誣告藥室.
藥室就要入獄了.
但是她去了哪裡呢?.
藥室的主人就將她抓住.
就被關在監獄裡面.
關在哪裡呢?.
就在黃的囚犯被關的地方.
我們明白到.
這個波提佛是一個很高級很高級的官.
在他的家裡面有這個監獄.
這個監獄就給那些人住呢?.
不是普通的囚犯.
是黃的囚犯.
一些重犯.
給他們去看著.
看著的.
不過從這本經文裡面看到.
在當中沒有經過審判.
要抓就抓.
是沒有公理的.
是沒有公義的.
那些大官喜歡怎樣就怎樣.
還有呢.
他不就範.
他不去向太太就範.
但是呢.
這個不是大罪來的.
但都要被囚禁在黃的囚犯監獄裡面.
上帝就是一個這麼奇妙的上帝.
他就行歧視上帝.
就在這個囚禁黃的囚犯的地方.
在這裡.
這個藥室就關了.
兩個被關在監獄裡面的人.

$^{401}$第一個就是施酒長.
第二個就是施膳長.
這兩個囚犯呢.
都是在黃帝吃東西之前.
他就要先吃.
看有沒有毒.
黃帝喝的酒呢.
他就先喝.
又看一下有沒有毒.
所以這兩個都是黃帝的心腹來的.
就算出去黃帝出征遠行.
都會帶著這兩個人.
不過不知道為什麼他們有什麼得罪黃帝.
所以就被監禁在黃的囚犯的監獄裡面.
在這裡呢.
上帝就給藥室認識了這兩個人.
也都透過解夢.
自己回去看這個創世記.
第39,40,41章那裡.
藥室可以去到法老王的面前.
因為有一天法老王又發了奇特的夢.
於是藥室就被推薦去幫法老王解夢.
在那次解夢的裡面.
就因為他成功解了法老王的夢.
然後就實質在埃及地發生.
就是在埃及地.
藥室告訴法老王.
他有七年是豐收的年.
你要儲多點糧食.
然後就有七年是饑荒的年.
在當中你就可以用之前那七年所儲的糧食來用了.
那真是當法老王在那七年的荒年裡面.
很多人都不夠食物.
這些食物不單止可以應付埃及地居民的生活需要.
還有周圍的人不夠.
都要來到埃及那裡向法老王買糧食.
所以法老就將藥室提升成為了.
一人之下 萬人之上的宰相.
我們看到相對就這麼奇妙.
在這個很不幸的坐牢.

$^{441}$我想這張照片那個牢都好一點.
我想在埃及的時候那個牢更加難受.
但是這個逆境裡面.
上帝有很多的恩典給藥室.
在上一年4月18號.
一個美國很著名的電視報道家Charles Stanley.
他就回到天父家.
他的生命很坎坷.
當他九個月大的時候.
他的爸爸就有腎病回到天父家.
到他12歲的時候.
媽媽就再婚.
不過就嫁給了一個醉酒鬼.
還有他有暴力傾向.
所以那時候整個家就打打殺殺.
連Charles Stanley都曾經拿刀和這個爸爸在打.
到他長大的時候.
他在教會的幕會.
在教會的幕會裡面.
因為就有兩派的鬥爭.
就是保守派和溫和派.
在神學上高處理教會的事情上高.
很多事情來鬥爭.
開董事會的時候.
Charles可以被那些領袖一拳打過去.
開執事會.
又或者當有一次.
當他成功就去留任做這個主任牧師的時候.
有大半的執事就離職.
所以在整個幕會裡面.
很多很多的艱難.
神卻很使用他.
在牧羊的教會.
他在1969年的時候.
他有五千人聚會.
到2020年的時候.
他88歲了.
他的教會是一萬五千人.
在幕會的過程裡面.
他同樣製作很多的廣播.

$^{481}$是一個很出色的電視報導家.
他在電視裡面的報導.
他就去在這個180多個國家裡面.
當然是極權或者進不去的地方.
也都有去到台灣或者周圍的地方.
人們都有看的.
他的那些電視報導.
也都翻譯了120多種的語言.
上帝很使用他.
不過整個生命.
無論小時候或者在幕會裡面.
都很多的艱難.
上帝就是這樣.
在艱難的當中.
只要他堅持的話.
上帝會在艱難的當中來祝福他的.
他就說過一句話.
他說上帝容許我們經歷這個苦難.
為的是要我們經歷他的信實.
加增他的信心.
上帝要我們在他的苦難當中.
經歷他的信實.
加增我們的信心.
所以我們在《約瑟》裡面我們看到.
他有這個信境.
天之驕子.
爸爸很疼他.
他做奴隸 做囚犯 很慘.
但當中他堅持不去做一些得罪上帝的事.
他堅持行上帝的話.
但上帝就祝福他.
在逆境的裡面.
他滿有主的恩典.
詩篇25篇第三節.
凡等候耶和華的.
必不羞愧.
惟有那無辜行奸詐的人.
必要蒙羞.
所以弟兄姊妹.
無論我們現在處在一個什麼樣的光景裡面.

$^{521}$我們等候上帝.
行上帝喜悅的事情的時候.
我們必不抱愧蒙羞.
他剛才開始播那首歌.
知否世事常變.
變幻原是永恆.
聽起來很有懷舊.
但我覺得無論歌詞或者音樂.
都會帶給我們心靈裡面有些悲願.
好像人生就是這樣.
好像沒有無奈.
好像沒有什麼出路.
在耶穌基督的裡面.
我們就知道.
當人生有這個起跌.
有這個痛苦的時候.
只要有主的同在.
我們可以體會祂的平安.
經歷到這個萬事互相效力的事實.
如此認識我這一首的詩歌.
就講到在人生裡面.
在人生裡面.
都是的.
會有很多的風浪.
有日子很快地過.
有春光.
有春天的時間.
有夏天很炎熱的時間.
有感到很落寞的秋天.
有這個令我們很冰冷的冬天.
但上帝的右手.
總會扶持著我們.
我們一起唱吧.
《春光》.
年月有如閃電飛過.
但你的手從未離開過.
看春光 看艷陽.
看落葉 看冰河.
你右手總就亦扶持著我.
我在地上年日你已經數輸過.

$^{561}$你心底的熱如海沙那麼多.
我在地的雨天也願夠看顧著我.
這孤身的愛可擦得透我.
年年將來雖比不過.
但我有主同在能衷過.
有北風 有巨浪.
有和煩 有艱難.
我從新主任下望而無望.
大家唱得不錯了.
去除了剛才那首《世事常變》的那種憂鬱.
可以再唱好一點的.
我們就停在這一頁裡面.
我記得那時候在午堂.
會在晚堂說過.
這些北風 這些巨浪 這些和煩和艱難.
就代表著我們生命裡面的逆境.
所以一定要有力.
有北風 有巨浪.
還要打死他們.
所以還要強力地打死他們.
有北風 有巨浪.
因為有主耶穌基督和我們同行的時候.
無論在什麼逆境裡面.
主會和我們同在.
讓我們在逆境當中有恩典.
我們要快一點.
那種味道要興奮一點.
因為無論逆境也好.
我們勝得過.
唱兩次好不好.
1 2 Go.
迎面將來雖被不過.
但我有主同在能衝破.
有北風 有巨浪 有和煩 有艱難.
我創新主任下望見到國.
再來迎面將來.
迎面將來雖被不過.
但我有主同在能衝破.
有北風 有巨浪 有和煩 有艱難.
我創新主任下望見到國.

$^{601}$我想再來多一次.
我創新主任下昂然報國.
我經常想到.
那時候中國大陸的樣板戲.
那些人就嗯.
那些人是吧.
在這個詩篇27篇就想.
在這個禍患當中.
我仍然可以昂身走過.
去到那裡就要這樣.
好 迎面將來最後一次.
迎面將來雖被不過.
但我有主同在能衝破.
有北風 有巨浪 有和煩 有艱難.
我創新主任下望見到國.
我們一起禱告.
再來多謝你.
人生當中有你和我們一起的度過.
以致這個世界上無論怎樣變.
我們的經歷怎樣變.
怎樣在高位 怎樣在低位.
怎樣在逆境當中.
我們仍然可以昂然這樣去度過.
我們多謝我們的主耶穌基督.
多謝主的恩典.
祝福我們時時刻刻有主的力量和勇氣.
感恩奉主基督聖名.
阿們.
\newpage



\section{加拉太書 2:20}
\label{sec:ltEd3gYtiPs}
\textbf{2024-05-05 早堂《基督在那裡?》加拉太書 2\_20  講員:朱浩廉博士}
\newline
\newline
連結: \href{https://youtube.com/watch?v=ltEd3gYtiPs}{\texttt{https://youtube.com/watch?v=ltEd3gYtiPs}} ~~~~ 語音日期: 2024-05-07
\newline
\newline
\hyperref[sec:U5CngR8W_hQ]{\small{< < < PREV SERMON < < <}}
~
\hyperref[sec:index_chronic]{\small{[返順時目]}}
~
\hyperref[sec:index_scriptual]{\small{[返順卷目]}}
~
\hyperref[sec:1QRoYKXVPvQ]{\small{> > > NEXT SERMON > > >}}
\newline
\newline
加拉太書 2:20
\newline
\begin{longtable}{cl}
\hline
\hline
章節 & 經文 (和合本修訂版)\\
\hline
2:20 & \begin{tabularx}{0.7\textwidth}{X} 現在活著的不再是我，乃是基督在我裡面活著；並且我如今在肉身活著，是因信神的兒子而活；他是愛我，為我捨己。 \end{tabularx} \\ \\
[1ex]
\hline
\hline
\end{longtable}
$^{1}$鄧小慧 早安.
今天很開心有機會來到大家當中.
去分享上帝的說話.
多謝教會的邀請.
多謝麗善牧師的邀請安排.
讓我來到這裡當中.
很少來到錦繡花園.
來到當中大家的湖.
真的很舒服可以走一圈.
特別昨天下這麼大雨.
今天太陽出來.
有沒有覺得心情很不同了.
不同了 是不是.
想像一下如果昨天不下雨.
今天這樣的太陽.
你的心情又如何.
可能有些不同 是不是.
正正就是可能下完大雨.
生活很多各樣的不方便.
今天見到太陽.
我們覺得 哦 神的恩典 是不是.
真的有這個感覺 是不是.
如果天天都是這麼曬.
夏天了 是不是.
三十多四十度 溫室效應.
我們就很多這些.
我們很多時候心情.
都會被環境所影響.
同樣的太陽 同樣的溫度.
但是當你經歷過不同的事情.
你有不同的解讀.
你和神的關係.
可能都是因為這樣受影響.
小弟在杜豐山工作.
很長的名字 主席幫我介紹.
容許我簡單說一下.
大家認識杜豐山嗎.
多少人去過杜豐山 舉舉手.
有沒有人去過 大部分去過.
杜豐山原來其實年紀都不小了.

$^{41}$杜豐山1930年開始開山.
都幾乎一百年了.
杜豐山是一個.
其實最初都是一個宣教地方.
向一些和尚 道士來宣教.
希望藉著一個中國式的建築.
讓他們在這裡比較舒服.
比較有親切的聽福音.
所以宣教其實很多東西配合.
所以流傳下來.
也因為過往有和尚 道士.
或者一些尋找耶穌的人.
在這裡信了主.
需要有些栽培 門徒訓練.
建立了現在有一些的靈修課程.
小弟現在在杜豐山.
主要有幾個工作.
剛才說過.
會做一些一對一的靈修指導.
關於你祈禱 靈修.
或者我們和神的關係的各種東西.
你可以上來找我聊聊天.
另外都會做一些.
可能大家剛才舉手.
都去過杜豐山退修.
我會帶退修.
無論在山上的退修.
或者如果教會覺得.
邀請我來這裡都可以的.
我會去不同教會的地方.
如果有合適的.
都會去帶一些退修.
一起親近主的聚會.
另外都有一些教學的工作.
因為我們都繼續培訓.
一些不同的靈修導師.
所以大家記得的時候.
紀念一下因為這個時代.
正正很需要我們多些.
從神那裡能夠得著能力的一個時代.

$^{81}$說回我們今天的題目.
基督在哪裡.
很多人上來杜豐山或者接觸我.
都會找耶穌基督.
可能有未信的人.
好像以前杜豐山那些和尚道士.
他們可能很多年修道.
學習他們的道理.
但他找不到耶穌基督.
有些信了耶穌基督.
但是缺了志的.
上教會的都來找耶穌基督.
為什麼.
因為我們頭腦上有耶穌.
在我們生命裡不是很經歷到.
不知道你是否這樣的.
不知道你是否有這個需要.
去找耶穌基督.
在我們經常說靈修的時候.
我們就會想.
我們每天看聖經祈禱.
然後就再想起什麼.
就再想起我們很多時候.
好像不能夠持續.
然後我們又再想起什麼.
每年年尾我們就在神面前立志.
我明年要做很多靈修.
多點親近主.
然後我們不斷循環.
一年又一年.
我不是說所有基督徒都是這樣.
有些基督徒是這樣.
我一想起靈修生活.
就是在說看聖經祈禱.
其實我們學習靈修.
就是學習我們知道耶穌基督在哪裡.
剛才我們讀過一段聖經.
傳統和合本其實有上面那句.
我們再讀多一次好不好.
這麼短不怕.

$^{121}$我們一起讀多一次.
1,2,3.
我已經與基督同釘十字架.
現在活著的不再是我.
乃是基督在我裡面活著.
並且我如今在肉身活著.
是因信神的兒子而活.
他是愛我.
為我捨忌.
我們再一次低頭祈禱.
七主多謝你.
多謝你愛我們.
你為我捨.
為著我們捨棄了你的生命.
叫我們這些原本是等待著刑罰.
等待著死亡的人.
我們今天得以養起念.
仰望你的十字架.
我們得歸入你的名下.
一起領受你的餅.
一起領受你的杯.
讓我們能夠.
可以承受你所為我們應許的.
這一份的平安喜樂.
所以主我們在你的愛裡面.
願你更新我們的靈.
讓我們常知道你的愛.
以致我們知道你的愛常與我們同在.
我們一起去思想.
我們在《太書》這句聖經.
求主的聖靈給我們有啟示.
給我們有感動.
讓我們明白你說話裡的奧秘.
我們祈禱奉耶穌基督的名.
阿們.
是不是很熟悉這句聖經.
我已經與基督同釘十字架.
現在活著的不再是我.
我會來找耶穌基督的弟兄姊妹說.
耶穌基督在哪裡.

$^{161}$在哪裡.
在這裡.
因為有這句金句.
我們就記得.
基督在我裡面活著.
所以叫大家多讀一次.
我們小時候上主學就要背這句聖經.
一定要懂得背.
我們與基督同釘十字架.
現在活著的不再是我.
我背了幾十年這句聖經.
所以不用看都應該OK.
但當我們不用看.
有時會不記得.
有時會找不到.
我們知道基督在我裡面活著.
找不到的.
會不會找不到.
不會的.
當你細心找.
你就會找到.
耶穌基督在我裡面活著.
怎樣在我裡面活著呢.
保羅在《哥林多後書》四章八至十節這樣說.
他說我們身上常帶著耶穌的死.
使到耶穌的生.
也都顯明在我的身上.
所以我專登和大家讀前面那句.
我已經與基督同釘十字架.
怎樣在自己的生命裡面.
找到活著的耶穌基督.
其中一個我們需要可以去學習操練.
或者提醒自己.
就是我們身上常帶著耶穌基督的死.
什麼叫做帶著耶穌基督的死呢.
保羅在前面有說.
我們四面受敵卻不被困住.
我們帶著耶穌基督受苦的生命.
我們面對著我們人生各樣的敵人困難.
每個人都會有.

$^{201}$信了耶穌都會有.
教會的牧者都會有.
執事都會有.
領袖都會有.
我們都會有敵人.
就算多好人的人都會有敵人.
你信不信.
信了是吧.
小弟自問.
我們學輔導覺得自己親和力比較高.
親和力高的人都有敵人.
這個世界總有人不喜歡你.
必然會有.
保羅都有.
耶穌有沒有.
耶穌都有.
是吧.
耶穌都有敵人.
不過保羅說我們四面受敵卻不被困住.
當我們常帶著耶穌基督的死和耶穌的復活.
我們都會面對敵人.
但我們不會困在那裡.
有些基督徒或者有些朋友.
因為有這些困擾.
為什麼會這樣.
為什麼我信了耶穌都是這樣.
為什麼很多情緒.
為什麼好像這麼多麻煩.
但我們找到我們耶穌.
耶穌的死和生在我們身上發生功效.
我們不被困住.
心裡作難卻不至失望.
我們都有作難的時候.
我們不是24小時充滿信心喜樂平安.
你是不是呢.
我就不是24小時都會這樣.
你說靈修導師都會這樣.
是的.
所以我們才要去靈修.
就是知道我們會有軟弱.

$^{241}$我們都會有時不開心.
我們有情緒.
我們有遇上不同的困難.
我們不是必然活在平安喜樂的環境裡面.
當我們經歷現在這種的世代.
是不是.
很多不同的轉變.
世局的變化.
經濟的起跌.
我們身體的健康.
社會的衛生上上都有改變的時候.
我們都會心裡作難.
但是不至什麼.
我們不至失望.
為什麼.
因為我們盼望在哪裡.
在上帝那裡.
人生最多就是什麼.
最多是什麼.
就是死.
是不是.
我不怕說.
最多就是死.
人生的困難.
無論經濟也好.
我們工作上也好.
家庭也好.
情緒也好.
身體健康也好.
最多最多就是什麼.
死.
但是我們死後有什麼.
我們死後有永生.
是不是.
我們的盼望在於主.
不是說我們不會遇到困難.
不是說我們一定會勝過.
因為第九節說.
我們遭逼迫.
我們卻不被掉棄.

$^{281}$我們遇到有逼迫.
但是神不會掉棄我們.
我們有時會被打倒.
會不會.
我們卻不至死亡.
當我們身上常帶著耶穌的死.
耶穌的生.
也都會顯明在我們身上.
所以各位弟兄姊妹.
我們面對現在很多變.
不穩定的一個世界.
我們更加要在我們心裡找到什麼.
有耶穌.
當我們順境的時候.
我們知道耶穌和我們在一起.
這個當然知道.
當我們遇到逆境的時候.
耶穌有沒有和我們在一起.
有啊.
希伯來書說.
祂為我們凡事受過苦難.
當我們經歷人生每一個苦難.
我們帶著耶穌基督的死.
我們知道在十字架上的耶穌.
和我們一起經過每一個人生的片段.
你有沒有信心.
耶穌基督在你裡面活著.
祂不是閉上眼睛看不到你.
祂又不是躲起來.
專門來訓練你.
其實在耶穌基督的死和生裡面.
我們是和耶穌基督在一起.
來經過我們的人生.
耶穌基督在哪裡.
耶穌基督不單止住在我們心裡面.
好像很渺小.
反而我們是住在耶穌基督的裡面.
耶穌住在我們裡面.
我們又住在耶穌的裡面.
怎樣解釋呢.

$^{321}$聖經說我們在肉身活著.
是因順神的兒子而活.
我們肉身活著.
心卻在主的裡面.
這句說話我們怎樣解釋呢.
我們活在肉身裡面.
意思是我們都有生老病死.
我們都會有肚餓.
我們會有不舒服.
我們肚餓要吃飯.
我們冷就要穿衣服.
我們都需要下大雨的時候.
我們需要有一個遮蓋.
我們在肉身活著.
但我們卻不是為肉身而活.
保羅說我們不是為肉身而活.
乃是為誰而活.
為耶穌而活.
活在主裡面.
我們肉身有各樣東西.
我們每個人都有任務.
即使我們做牧者,做傳道人.
我們都有我們的任務.
我們都需要上班.
我們需要發工資.
我們都需要坐車.
每樣東西都很實際.
但我們不是為這些而活.
我們在肉身活著.
乃是為著主而活.
聖經裡面說.
我們乃是信神的兒子而活.
這個信字在原文裡面.
它並不是一個動詞.
它是一個名詞.
說的是我們心裡面帶著一種信仰.
帶著一種信心.
帶著一種我們在人生裡面.
作為我們做人的原則.
作為我們的價值觀.

$^{361}$作為我們的目標.
這個就是信字的解釋.
這個信字在原文裡面.
亦都可以解著一種的誓約或者忠誠.
當我們肉身活著.
沒錯我們要吃飯.
我們要賺錢.
我們要養家.
我們要面對我們人生九萬多樣的問題.
但這些是我們其中一部分的責任.
並不是為這些而生存.
今日我們為這些生存.
是為著因信神的兒子而活.
這個信就是我們在主裡面的一個立約.
就是當你當日舉手信耶穌.
說我們接受耶穌基督.
做我們的主和個人的主的時候.
我們納了的一個約.
你還記得那個約嗎.
弟兄姊妹.
就是讓主成為你真正的主.
我們小時候學傳福音.
有一本書叫做什麼.
你有沒有看過四個屬靈的定律.
有沒有學過.
這本書我以前就是這樣去學傳福音.
當年就是跟我們教會的牧者和柳振平牧師很熟.
邀請柳振平牧師.
我是親自和柳振平牧師的身上學.
怎樣用這個四個屬靈定律.
原來很簡單.
你照著讀.
由第一頁讀到最後一頁就行了.
你用過了沒有.
當時我是一個中學生.
我就拿著這本小書.
中學生走出去找什麼人傳福音.
小學生.
難道找大人嗎.
很害怕.

$^{401}$小學生拿著這本書.
你聽過屬靈的四個定律沒有.
然後就由第一頁讀到最後一頁.
原來不用讀神學.
到最後一頁問他.
你願不願意接受耶穌基督做你的主和個人的主.
我也是照著讀.
原來他說我願意.
原來真的可以.
最後一頁有一個圖畫.
圖畫有一個圓圈.
中間畫了一張椅子.
在椅子旁邊有一個十字架.
有一個我字.
第一幅圖原本的我字坐在椅子上.
耶穌在旁邊.
沒有放在椅子底下那麼壞.
旁邊有個十字架.
然後他就解釋什麼叫接受耶穌基督做他的主.
就是將我字放在上面.
然後放一個什麼上去.
放一個十字架上去.
由我們第一天看著這本又明不明的小書.
然後聖靈感動我們.
我們接受耶穌做我們的主.
信就是立了這個約誓.
對主忠誠.
今天我們生活有很多的問題.
我們相信耶穌基督會給我們有能力帶我們經過.
但我們最重要的並不是為了解決這些困難.
以致我們進入一個更美好更舒適的生活.
而活我們這個人生.
而是因為我們立志將耶穌基督放在我們生命的寶座上.
所以怎樣叫活在基督裡面.
好像這個第二點裡面說.
就是活在有上帝兒子的信仰裡面.
在你的生活裡面.
有沒有上帝的兒子在你的信仰裡面.
當你做人生抉擇的時候.
有沒有上帝的兒子在你做人的原則和信仰上面.

$^{441}$很多時候我們以為我們去靈修.
我們去找耶穌.
好像是一些很宗教的東西.
好像是很神秘的.
我們一定要進入某種狀態.
然後我們就在我們的意象裡面.
我們經歷些什麼.
沒錯那個都是靈修的一部份.
但那種靈修最終是要我們將我們和神的關係.
放在我們的生活裡面.
那個靈修才是有用的靈修.
最後聖經裡面說什麼呢.
他是愛我為我捨己.
我們在愛和捨己裡面找到耶穌基督.
耶穌基督在我們裡面.
我們也活在耶穌基督的裡面.
而這一種主又活在我們裡面.
我們又活在他裡面.
最終有個結果就是我們能夠可以讓我們的生命.
成為一個愛和捨己的出口.
聖經說主就是愛.
所以找到愛你就找到主.
你說有沒有主基督在你的生命裡面呢.
那你可以問一下你自己.
當我們靈修當我們安靜.
當我們默想.
你數一下你生命裡面有多少愛.
對神的愛 對人的愛.
甚至對自己的愛 你有多少.
找到愛你就找到主.
因為約翰一書說 神就是愛.
如果你和人去相處都是一樣的.
你在人的身上找到愛.
你會找到主.
我未必和你很熟.
但在我和你相處的裡面.
和你傾談的裡面.
和你一起祈禱和你一起合作事奉.
我找到愛.
我知道有主在我們當中.

$^{481}$在我們群體的裡面.
所以找主不一定要躲在深山裡面.
當然我歡迎你來杜豐山.
杜豐山並不是一個深山.
現在人多到不得了.
每個人都上來打卡.
甚至那些小紅書.
有些網紅介紹.
怎樣上來杜豐山 哪個景點.
然後下去山下排頭村.
吃哪間茶餐廳都好吃.
多人到不得了.
我們不是再去那些標誌性的地方找主.
是在我們生命裡面找主.
這是我們學靈修一個重要的方向.
寫記 寫字 原文其實有託管的意思.
甚麼叫託管.
你們幼稚園有沒有託管的.
還是只有幼稚園.
託管的意思.
即是他不是放棄 是甚麼.
託管 把小朋友託管給你.
不要小朋友 是甚麼.
先放在你身上.
將我們生命裡面最寶貴的.
我相信父母應該都是.
我們放在一些老師.
老師是否在一些專業人士的手裡面.
雖然理想上是這樣.
實際上個個心態有點不同.
託管 我們將我們的生命交在神的手裡面.
我們很怕說寫記.
因為我們很怕只有一張嘴.
好像說了做不到.
寫記很難 是不是.
我牧養了幾十年.
鄧應先輩都覺得很難.
我們一說寫記就想起少年的官.
一說寫記就以為耶穌叫我們甚麼.
你變賣一切所有分給窮人.

$^{521}$然後你再來跟從我.
寫記 我盡量都不要說.
或者我們輕輕帶過.
但原來 皆太書裡保羅說寫記.
就是將我們的生命託管在神的手裡面.
神不需要你拿了你的命去.
為甚麼.
耶穌基督都為我們寫命.
祂做甚麼要拿了我們的命呢.
耶穌基督花這麼大的代價拯救我們.
祂說祂來了 要叫人得甚麼.
得生命 並且得更豐盛.
祂怎會要去毀壞 或者破壞我們.
要我們過一個很艱難的生活.
寫記的意思就是我們不要執著自己的生命.
跟我剛才說的那一點有很連貫的關係.
就是我們活在肉身裡面.
但卻是不為肉身而活.
我們將我們的生命交在神的手裡面.
我們相信上帝向我們所懷的意念.
都是持平安的意念.
祂不會害我們 祂不會作弄我們.
更加不會去奴役我們.
寫記就是我們不再為我們自己而活.
今天我們已經將十字架放在椅子上.
我們就是要先求祂的國和祂的義.
我們相信我們整個生命放在神的手裡面.
上帝給我們的恩典超過我們所放在的更多.
你相信嗎.
耶穌基督復活了.
當我們去找耶穌基督的時候.
祂不是再在十字架上.
當然不再在墳墓裡面.
卻是活在與基督同釘十字架的人的生命裡面.
你有沒有與基督同釘十字架.
耶穌基督復活了.
祂活在那些對主忠誠的人裡面.
你是否對主忠誠的人.
耶穌基督活在那些愛主和渴望主的愛的人當中.
你是否愛主和渴望主愛的人.

$^{561}$耶穌基督活在那些將生命傳言交託在主的人的生命裡面.
你是否將生命傳言交託在主的手裡面呢.
讓我們低頭祈禱.
我們有一分鐘安靜的時間.
當我們低頭安靜的時候.
我邀請你回想一下今早.
你聽到這個陌生的講員.
跟你講一節你很熟悉的聖經.
今日在這個宣講裡面.
有沒有哪一句的經文.
哪一句的說話.
正在與你的心去說話呢.
我想邀請你在這個位置裡面.
就為著這個感動我們有個回應的禱告.
求主幫助.
過往我們生命裡面.
我們有忽略了的.
我們有困難裡面有迷失的.
求主幫助我們.
輕輕我們的心.
重新的來到去仰望主的十字架.
因為這個世界.
所有的東西.
最終都會過去.
不是我們能夠捉得實.
不是我們能夠雙手可以保存得到.
唯有我們在神裡面的信心.
和關係和愛心.
永遠都不會改變.
所以求主堅固我們的心.
在你的話語裡面.
常常的紮根.
面對人生各樣的轉變.
即使有更壞的天氣.
我們仍然相信主你活在我們裡面.
我們又活在你的裡面.
當我們每天在日常的生活和人相處.
無論是工作無論是事奉.
我們找到愛.
我們找到有犧牲的精神.

$^{601}$我們又知道主在我們中間.
我們同心祈禱.
奉耶穌基督的名.
阿們.
\newpage



\section{馬太福音 11:2-19}
\label{sec:1QRoYKXVPvQ}
\textbf{2024-05-05 午堂《沒有最大,只有更大》馬太福音 11\_2-19  講員:張雲開老師}
\newline
\newline
連結: \href{https://youtube.com/watch?v=1QRoYKXVPvQ}{\texttt{https://youtube.com/watch?v=1QRoYKXVPvQ}} ~~~~ 語音日期: 2024-05-08
\newline
\newline
\hyperref[sec:ltEd3gYtiPs]{\small{< < < PREV SERMON < < <}}
~
\hyperref[sec:index_chronic]{\small{[返順時目]}}
~
\hyperref[sec:index_scriptual]{\small{[返順卷目]}}
~
\hyperref[sec:VvxSAXf_Lww]{\small{> > > NEXT SERMON > > >}}
\newline
\newline
馬太福音 11:2-19
\newline
\begin{longtable}{cl}
\hline
\hline
章節 & 經文 (和合本修訂版)\\
\hline
11:2 & \begin{tabularx}{0.7\textwidth}{X} 約翰在監獄裡聽見基督所做的事，就派他的門徒去， \end{tabularx} \\ \\ \relax
11:3 & \begin{tabularx}{0.7\textwidth}{X} 問耶穌：「將要來的那位就是你嗎？還是我們要等候另一位呢？」 \end{tabularx} \\ \\ \relax
11:4 & \begin{tabularx}{0.7\textwidth}{X} 耶穌回答他們：「你們去，把所聽見、所看見的告訴約翰： \end{tabularx} \\ \\ \relax
11:5 & \begin{tabularx}{0.7\textwidth}{X} 就是盲人看見，瘸子行走，痲瘋病人得潔淨，聾子聽見，死人復活，窮人聽到福音。 \end{tabularx} \\ \\ \relax
11:6 & \begin{tabularx}{0.7\textwidth}{X} 凡不因我跌倒的有福了！」 \end{tabularx} \\ \\ \relax
11:7 & \begin{tabularx}{0.7\textwidth}{X} 他們一走，耶穌就對眾人談到約翰，說：「你們從前到曠野去，是要看甚麼呢？看風吹動的蘆葦嗎？ \end{tabularx} \\ \\ \relax
11:8 & \begin{tabularx}{0.7\textwidth}{X} 你們出去到底是要看甚麼？看穿細軟衣服的人嗎？那穿細軟衣服的人是在王宮裡。 \end{tabularx} \\ \\ \relax
11:9 & \begin{tabularx}{0.7\textwidth}{X} 你們出去究竟是要看甚麼？是先知嗎？是的，我告訴你們，他比先知大多了。 \end{tabularx} \\ \\ \relax
11:10 & \begin{tabularx}{0.7\textwidth}{X} 這個人就是經上所說的：『看哪，我要差遣我的使者在你面前，他要在你前面為你預備道路。』 \end{tabularx} \\ \\ \relax
11:11 & \begin{tabularx}{0.7\textwidth}{X} 我實在告訴你們，凡女子所生的，沒有一個比施洗約翰大；但在天國裡，最小的比他還大。 \end{tabularx} \\ \\ \relax
11:12 & \begin{tabularx}{0.7\textwidth}{X} 從施洗約翰的日子到今天，天國受到強烈的攻擊，強者奪取它。 \end{tabularx} \\ \\ \relax
11:13 & \begin{tabularx}{0.7\textwidth}{X} 眾先知和律法，直到約翰為止，都說了預言。 \end{tabularx} \\ \\ \relax
11:14 & \begin{tabularx}{0.7\textwidth}{X} 如果你們願意接受，這人就是那要來的以利亞。 \end{tabularx} \\ \\ \relax
11:15 & \begin{tabularx}{0.7\textwidth}{X} 有耳的，就應當聽！ \end{tabularx} \\ \\ \relax
11:16 & \begin{tabularx}{0.7\textwidth}{X} 「我該用甚麼來比這世代呢？這正像孩童坐在街市上向同伴呼喊： \end{tabularx} \\ \\ \relax
11:17 & \begin{tabularx}{0.7\textwidth}{X} 『我們為你們吹笛，你們不跳舞；我們唱哀歌，你們不捶胸。』 \end{tabularx} \\ \\ \relax
11:18 & \begin{tabularx}{0.7\textwidth}{X} 約翰來了，既不吃也不喝，人們就說他是被鬼附的； \end{tabularx} \\ \\ \relax
11:19 & \begin{tabularx}{0.7\textwidth}{X} 人子來了，也吃也喝，他們又說這人貪食好酒，是稅吏和罪人的朋友。而智慧是由它的果子來證實的。」 \end{tabularx} \\ \\
[1ex]
\hline
\hline
\end{longtable}
$^{1}$多謝主席帶領我們讀了這段今天要思考的經文.
這段經文表面上好像是在說施洗約翰.
由於他坐牢所以引起一個話題.
耶穌也說有關他的一些事情.
耶穌對約翰有一些看法.
耶穌對所有人都有看法.
包括你和我當然.
但是一會兒我希望能夠說服大家.
其實這段經文不是說施洗約翰.
施洗約翰只是一個引子.
其實是說主耶穌自己.
這段經文和我們有沒有關係呢.
表面上又好像是沒有什麼關係.
說施洗約翰當是說主耶穌基督.
和我們有什麼關係呢.
一會兒我也希望能夠讓大家看到.
這段經文耶穌所講論的.
提到施洗約翰也好.
提到他當時當世代的人也好.
其實今天和我們息息相關.
焦點我一會兒和大家分享的焦點.
是會放在Unpack第十一節那裡.
就是主耶穌所說.
說到施洗約翰的時候.
說了一句比較奇特的話.
表面上就好像有一點點自相矛盾.
主耶穌對當時的門徒說.
凡女子所生的沒有一個比施洗約翰大.
當時周圍的人可能也會聽到.
但是在天國裡面.
興起的沒有一個.
最小的也比他還大.
我們主要到最後也希望能夠明白.
到底主耶穌在這裡說的是什麼.
為什麼焦點不是約翰.
其實是他.
好像和我們沒有關係.
又和我們其實是直接息息相關.
我就將自己在這裡的思念.
對這段經文的一些思念.

$^{41}$呈現在大家面前.
也在崇拜當中呈現在主的面前.
向他擺上再一次低頭祈禱.
天父我們在這裡敬拜的時候.
願意你自己心滿意足.
不單止我們知道不單止只是我們.
而是在這個城市這個時區.
這個世界裡面不同時區.
我們這群屬於你自己的子民聚集.
向你敬拜的時候.
願意你也向我們施恩.
你自己的話如啟迪我們.
開啟我們的心你的靈.
也感動我們開我們的眼睛.
也給力量我們.
讓我們能夠在平常的日子.
能夠在身體力行.
將我們的信仰.
我們對你的信能夠實踐出來.
誰聽我們禱告.
奉主名求.
阿們.
這件事整段經文引發事情的原因.
當然就是第二節.
約翰坐牢的時候.
以前和現在一樣.
都可以有人探.
聽見基督所做的事情.
這裡很特別的.
馬太的記載.
不是說約翰聽見耶穌所做的事情.
他已經說明是基督.
但如果你看這節經文.
接著約翰聽見基督所做的事情.
就派兩個門徒去.
第三節就說要問.
這裡就轉了做耶穌.
要問耶穌.
將來要來的位置是不是就是你.
還是我們其實都要等另外一個.

$^{81}$你都不是.
所以馬太起講的時候.
開頭承認耶穌已經是基督.
倒轉引發這件事的本身.
反而是施洗約翰.
耶穌的表親.
對耶穌出現懷疑.
所以第二個問耶穌都合理的.
因為施洗約翰門徒.
就不會稱呼耶穌為基督.
他們的師父正正懷疑.
是不是那位基督.
那位彌賽亞.
那位先知應許要來.
所以整件事引發就是因為.
約翰坐牢.
他對耶穌起了疑問.
坐牢的人可能都會對很多事情產生疑問.
包括對自己的人生都會產生疑問.
為什麼會坐牢.
我是不是真的做錯事.
是不是蒙不白自怨.
抑或是死有餘辜.
到底是一件什麼事情.
我們要明白約翰坐牢.
到底是一回什麼事.
在馬太福音.
這裡沒有記載.
但三章之後.
我們就會看到約翰坐牢最終的下場.
在盧伽福音就記載得仔細一點.
他最終會被希律所殺.
但是他為什麼坐牢.
就是因為他的嘴巴.
如果我們中國人說的話.
是有些意思的話.
就是禍從口出.
他坐牢的災禍.
就是由他的嘴巴招致.
因為他是責備當時的昏風王希律.

$^{121}$不是耶穌出世的時候.
那個大希律.
就是把伯利恆附近的城鎮.
一歲以下的小朋友通通殺光.
那個是希律的爸爸.
這個是大希律的兒子.
希律做了什麼事.
使約翰的嘴巴針對他.
就是他娶了他的嫂子.
他的哥哥菲利普的老婆希羅底為妻子.
問題就是.
菲利普還沒有死.
我們知道猶太人的規矩.
對面為了中國人所說的.
某個程度上都是繼後鄉燈.
所謂繼後鄉燈.
其實大家姓氏都一樣.
是承受遺產.
猶太人的內婚.
就會兄弟的配偶.
由於兄弟過身.
是由哥哥或弟弟娶過來.
讓他的財產不會流出去.
其他外人的地方.
我看魯德紀都可能聽說過.
接觸過這類事情.
但問題是哥哥還在生的時候.
就變成亂論.
再加上希羅底和菲利普和希律.
都是有其他關係的.
事源大希律他們的爸爸.
很喜歡將他生出來的子女.
互相分配.
很奇怪的.
希律的子女都是不同的太太.
不同的媽媽生出來的.
所以其實希律和菲利普.
是同父大希律.
但二母的兄弟.
而希羅底也是希律的大哥.

$^{161}$阿利多布的女兒.
所以你想一下.
這個情況就很複雜.
這是他們的父親大希律.
安排了他的兒子的女兒希羅底.
嫁給他另一個兒子叫做菲利普.
然後他第三個兒子希律.
去到羅馬探望菲利普的時候.
有一次因為菲利普住在羅馬.
見到他的嫂子生了一隻貓入眼.
兩個都互相感情.
結果希羅底跟了希律回巴勒斯坦.
跟他的丈夫菲利普離婚.
所以希律就娶了他的嫂子.
但希羅底是希律的大哥.
也是同父二母的阿利多布的女兒.
所以希律的老婆是他的姪女.
也是他的嫂子.
現在變成他老婆.
很亂啊.
皇族帝王之家的貴族.
我們知道在埃及的法老傳統宮廷裡面.
他們都會有亂倫這種情況出現.
但是聖經裡面我們可以看見.
約翰針對他主要就是他娶了他的嫂子.
當他哥哥還在世的時候.
他娶了他的嫂子.
結果希律就把他扔進監獄.
因為不想聽.
為什麼是施洗約翰的嘴巴.
又很難呢.
這個當然是和施洗約翰自己的工作.
他的身份有關係.
這個我們就要看馬太福音第三章.
沒有powerpoint的.
我讀給大家聽.
其實對施洗約翰我們都會有些認識.
我讀給大家聽.
第三章是馬太這樣記載.
那時我讀的是和合本.

$^{201}$耶路撒冷和猶太傳遞並約旦河一帶的地方.
都出去約翰指著施洗約翰那裡.
承認他們的罪.
施洗約翰宣講天國近了.
你們都要悔改.
後來主耶穌剛出來的時候.
都是宣講這個訊息.
那些人就去約翰那裡承認他們的罪.
在約旦河裡面受他的洗.
約翰見到那麼多的法利賽人和殺都該人都來受洗.
就對他們說.
篤蛇的種類.
誰指示你們逃避將來的憤怒呢.
你們要結出果子泥和悔改的心相稱.
不要自己心裡說.
亞伯拉罕是我們的祖宗.
我告訴你.
上帝能夠在這些石頭裡面替亞伯拉罕興起他的子孫來.
現在把斧頭已經放在樹根上.
凡不結好果子的樹都會撼下來.
吊起火裡面.
我是用水.
接著下來他就說.
耶穌說.
用水和你們施洗.
叫你們悔改.
但是那位在我以後要來的能力比我更加大.
我就是和他提鞋.
都不配.
他要用聖靈和火為你們施洗.
他手裡面拿著波姬.
要揚正他的牆.
將墨紙收起倉裡面.
將腔用不滅的火燒盡了.
這段是經文.
尤其是約翰講到後來.
他的表弟老表主耶穌和他相差六個月.
因為瑪利亞探訪於伊利沙伯的時候.
他已經有六個月的身孕.
他對他的認識.

$^{241}$就是說那位要來的比我更加大.
我和他提鞋都不配.
我是用水和你們洗.
下藥蛋河.
叫你們悔改悔罪.
天國近了.
人都應該悔改.
悔改是預備迎接天國的基本的人類的態度.
如果不是就沒有人能夠進入天國.
所以洗禮悔改.
接受他的洗.
但是後來那位.
他是用聖靈和火.
和你私洗.
約翰一開始出來.
就把口.
我們說口沒遮攔.
指著個個都有罪.
你試一下.
今天出街.
在地鐵站門口巴士站隔離.
沒有阻礙人上巴士.
總之在隔離拿個大聲公.
天國近了你們都要悔改.
每個人都有罪.
世人都有罪.
規矩上你試一下說.
我們傳福音.
我都沒見過人夠膽這樣說.
說話是對還是錯.
是對的.
但是能夠說嗎.
你夠膽這樣說嗎.
約翰就是這樣說.
他夠膽說.
也就是這樣說.
甚至有人來到.
要聽他說話.
又或者都有點感動.
要讓他接受他的.

$^{281}$在約旦河洗禮.
他要罵他們.
你們洗什麼.
你講一套做一套.
除非你的行為結出.
和你悔改的心相稱的果子.
不然白洗沒用.
上帝的斧頭已經放在樹根上.
凡結好果子的就斬下來.
扔在火上.
所以約翰服侍的特徵.
就是他很有火.
他的火是從口中噴出來的.
所以大家記得啟示錄裡面.
第十一章的兩個見證人.
突然第十一章開頭.
量度一下那個聖城.
然後突然不知道從哪裡來的.
就出現兩個見證人.
又像以利亞的穿著.
令人想起以利亞.
令人想起施洗約翰.
他們兩個見證人的口是噴火的.
凡敵對他們 敵黨他們.
這兩個見證人用口中的火.
是消滅了這些人.
夠了火.
而主要不是透過暴力.
是透過口.
事實上如果我們看.
不單止他施洗約翰有把火.
他對主耶穌的理解.
對他這位老表的理解.
都是由火來理解.
約翰指出耶穌的三味真火.
我們中國人道家所說的三味真火.
沒有什麼是燒不爛的真火.
剛才我讀給大家聽.
馬太福音第三章的第十節.
是當約翰向那些要來接受洗禮的人警告他們的時候.

$^{321}$第十節他說.
現在把斧頭已經放在樹根上.
凡不結好果子的樹就砍下來.
雕起火裡面.
這個就不是他自己會做的.
這個可以說是上帝做的.
第十一節我就說到耶穌.
我是用水替你施洗 叫你悔改.
但那個後面來能力給我大.
我就算和他提鞋都不配.
他用聖靈和火.
主耶穌都是和火有關的.
和你們施洗.
再加上進一步.
他說耶穌手裡面拿著波基.
要揚正他的牆.
把墨紙收在倉裡面.
剩下的來看批.
就用不滅的火燒了.
所以他將耶穌將來會來的那位的服侍.
兩次用火 聖靈和火.
這個火就沒理由理解成.
是士團傳第二章聖靈降臨.
好像聖靈降臨 好像火焰一樣.
是臨到各人的頭上.
就不是那些火的樣式的東西.
而是真是火.
而且他說要把人篩.
可能意思就是那些不願意悔改的人.
就會用火 不滅的火燒盡.
而且耶穌將來的火的服侍.
聖靈與火和不滅的火的服侍.
都是和上帝的斧頭.
放在樹根上的斧頭.
對著不願意悔改的人.
把那棵樹不結出果子的人.
把那棵樹砍了 扔在火裡面.
同出一轍 是同一條道路.
所以我們看到約翰是這樣看待耶穌的.
他自己有火.

$^{361}$但他預期要來的那位耶穌的火更加厲害.
我都是用水 其實是水去陪火.
水是和他施洗 叫你悔改.
他就不同了.
這是在馬太福音第三章.
約翰對耶穌的看法.
我們剛才讀第十一章.
你就會發覺 耶穌對約翰都有看法.
老表互相對對方都有些看法.
耶穌對他的老表 表哥的看法又如何呢.
我們看剛才讀過的十一章第十節.
我們就這樣抽第十節和第十四節來讀.
第十節之前就說 你們先出去看甚麼.
看蘆葦 看被風吹的蘆葦.
不好意思.
手機的問題就是手動一動 又不見了.
還是穿細軟衣服的人 還是要先看先知.
接著主耶穌在第十節說.
前面說主耶穌說 沒錯他是先知.
而且比先知更大.
第十節說 經上面記著說.
我要差遣我的使者在你面前.
為你預備道路.
所說的就是這個人.
這裡是引用馬拉基書第三章第一節.
接著第十四節 馬太福音第十一章第十四節.
主耶穌再說多一次.
說到是婦人所生的 無一個大得過約翰.
但是在天國裡面最小的都比約翰大.
這番說話之後 在第十四節主耶穌就說.
你前面說是先知和預言到約翰為止.
你們如果願意領受的話.
這個人 即是指著約翰.
就是應當來的 以利亞凡有耳的就應當聽.
這是引用馬拉基書的四章五節.
馬太福音前面那一卷就是了 很近的.
舊日最後的一卷書.
而事實上如果我們看.
路加福音第一章第十七節.
當施洗約翰的爸爸.

$^{401}$施洗約翰出世.
起名的時候 終於他爸爸撒迦利亞.
祭司來的 開口了.
上帝鬆了他的口.
因為他太太懷著嬰兒期間.
上帝不讓他開口.
因為他要笑 老婆這麼老了.
怎麼可能還生孩子呢.
不讓你出聲.
都不知道是不是一種懲罰.
不過都很辛苦的 應該是.
到了出世的時候.
要改名的時候 開了他的口.
其中上帝感動他.
他就說我這個兒子.
就要讓父母的心歸向兒女.
兒女的心歸向父母.
亦都是引用馬拉基書第四章第六節.
所以可以說施洗約翰和以利亞的關係.
其實由出世的時候.
他的爸爸撒迦利亞.
已經是受上帝所感動.
是連在一起的.
以致他的老表.
是耶穌自己都承認.
在這裡剛才我們看的十一章.
第十節和第十四節.
未必每個人都承認.
但最重要是你能夠領受的話.
這個人就是應當來的以利亞.
施洗約翰就是馬拉基書裡面所說.
那位的以利亞要為將來.
那位尼徹亞預備道路.
預備人心.
亦都會將父母的心歸向兒子.
兒子的心歸向父母的那一位.
以利亞和施洗約翰有什麼關係呢.
就是撒迦利亞說歸說.
但我們理解這樣的關係來說.
以利亞到底為誰預備地方.

$^{441}$以利亞為誰預備.
為他的徒弟預備.
為以利沙預備.
沒有理由.
所以如果約翰為基督預備道路.
那以利亞又沒有為人預備道路.
就著預備這一談.
以利亞和施洗約翰就沒有什麼直接有比較.
到底稱施洗約翰為以利亞.
那個比較點在哪裡來呢.
我都認為基本上都是一樣東西.
就是火.
大家知道以利亞的嘴巴.
以利亞不是宮廷裡面的先知.
他是地方人士.
但是他的勇到一個地步.
施洗約翰都離開希律的宮廷.
向他講話.
說他犯罪應該悔改.
以利亞直接在阿哈王面前篤口篤鼻.
對他的皇后出言不遂.
皇后夜洗別就最憎恨他.
一方面在宣講上叫以色列人回歸.
挽回以色列人的心.
這是以利亞工作.
但另一方面以利亞的服事當然和火有關係.
在加密山上和巴力先知鬥法的時候.
巴力先知做什麼都沒事發生.
他們獻的祭.
但是以利亞看夠了.
夠了嘛.
可能巴力睡了覺.
可能他走開了.
可能去了廁所.
遞他們一番之後.
以利亞祈禱.
天上有火來燒盡了.
以利亞擺在壇上面的祭物.
所以以利亞無論是口.
或者實際上最令人記得的大場面的服事.

$^{481}$都是和火有關係.
我覺得其實以利亞和施洗約翰.
連在一起.
都是他們的宣講叫人悔改.
叫皇帝悔改.
還有象徵著這種精神.
這種態度.
內心裡面的火.
記住.
一直以來都是和火有關係.
施洗約翰認為.
他的表弟耶穌的服事.
那把火比他還大.
開始明白.
為什麼在約翰坐牢的時候.
他內心有一點動搖.
我們回到第二第三節.
第十一章.
將來他叫他的徒弟.
去問耶穌.
將來要來的那一位.
我做這麼多事.
是不是為你預備.
還是我其實是白做.
或者為一個不知道是誰.
在預備他的道路呢.
為什麼會這樣說呢.
為什麼會導致這樣的懷疑呢.
我們由耶穌基督的回覆裡面.
可以看到一些端倪.
既然約翰的門徒來問耶穌.
耶穌當然要答覆.
這就是第四節.
他就和約翰的徒弟這樣說.
耶穌回答說.
你們回去將所聽見.
所看見的事情.
告訴約翰.
就是盲的看見.
跛的行走.

$^{521}$生大麻風的得到潔淨.
聾的聽見.
死人復活.
窮人有福音傳給他們聽.
這就是我所做的.
但是其實這些東西.
有一點點redundant.
有一點點多餘的.
耶穌說這番話.
因為第二節那裡說.
約翰在監裡聽見基督所作的事.
他正正就是聽見耶穌在做這些事.
他就知道.
耶穌又叫他的門徒回去說.
說同樣的事給他聽.
所以問題就是最後那句.
是第六節.
最後那句就是.
凡不因我跌倒的就有福.
這句就可圈可點.
因為耶穌所說的其實約翰已經知道.
而約翰開始懷疑.
正正就是因為聽見耶穌就是做這些事.
做這些事有什麼問題.
耶穌治病趕鬼.
有什麼問題才行.
後面我們剛才讀的那裡.
還說到他和瑞利聚人一起吃飯.
所以被人批評他貪食好飲.
約翰從來不參加宴會.
你自己飲好了.
我就吃蝗蟲野蜜.
自己一個人在那裡.
他就不會去別人那裡.
所以就被鬼附他.
這個人不群.
有問題.
所以主耶穌所做的事.
是約翰有少少看不慣.
為什麼和我差這麼遠.

$^{561}$而且我們說治病趕鬼.
是不是好事.
問問大家是不是好事.
上帝的恩典能道.
是不是好事.
被擄的得釋放.
是不是好事.
但是對於約翰來說.
OK.
不是不行.
但是缺少一樣東西.
缺少什麼東西.
缺少什麼.
火.
火從何來.
我現在是為了天國.
為了宣講上帝.
為了叫人悔改.
得罪了昏風王希律大帝.
不是希律大帝.
是他的兒子.
而坐牢.
有沒有命都不知道.
他有點擔心.
可能會做完.
做完了.
但是要為預備的人.
還在那裡摸摸.
摸摸這個套口水.
擦一下眼睛.
在那裡驅鬼.
什麼事.
他看不到那把火.
看不到好像他那樣.
叫人悔改.
責備人.
甚至三味真火.
包括他會履行上帝.
他就是上帝的使者.
把斧頭放在樹上.

$^{601}$把他斬下來.
用火燒那棵樹.
沒用的.
斬了他.
因為如果枯木慢慢就.
大家都知道大風一吹.
倒下了反而會壓到人.
壓到屋子就不好.
所以要斬了他.
此其一.
用聖靈同火先洗.
上帝的靈.
大能的靈臨到.
以利亞正正就是.
以利沙所求的.
就是感動你的靈.
加倍的感動我.
所以以利亞的靈.
讓他有能力.
但是火呢.
聖靈和火.
用衛利利斯洗.
接著就用波基 楊俊 張康.
用不滅的火燒盡.
喂 大哥.
這些就不是自病趕鬼.
不是飲飲食食.
不是坐在酒吧那裡.
跟人聊天.
這段時間在哪裡發財.
好像聽說.
你的兒子身體不太好.
要不要我看看.
幫到你.
就不是做這些事.
當然耶穌也有宣講.
天國近了 要人悔改.
但是耶穌的人跟耶穌.
自病趕鬼之後.
還有得吃.

$^{641}$起約翰福音第六章那裡.
那些人跟他.
跟跟耶穌說.
他跟門徒說.
你知不知道那些人跟我做什麼.
質詢那些人.
你們跟我.
就是因為跟著我.
有得吃的意思.
耶穌是分東西給人吃.
行神蹟.
分東西給人吃.
門徒也是沒有福氣的時候.
耶穌居然有福氣.
是讓跟他 聽他到處.
聽他宣講的人得飽.
這些跟火沒有什麼關係.
跟審判沒有什麼關係.
跟上帝的憤怒.
是沒有什麼關係.
這個令到約翰極其困擾.
那我為什麼要坐牢.
你不要想我為上帝大發熱心.
好像以利亞到了一個地步.
他depress.
我為你大發熱心.
不如我的祖宗為你大發熱心.
結果他們還要殺我.
為上帝大發熱心.
那些人就會殺你.
就是追斬你.
以利亞是這樣的.
約翰也是這樣的.
但是他與其.
耶穌也是這樣.
但是我們要花多一點精神.
看耶穌向約翰的回覆.
尤其是第五節.
將你所見到的事回覆約翰.
就是 瞎子盲的看見.

$^{681}$跛的行走.
長大謀反德潔正.
聾的聽見.
死人復活了窮人.
傳福音給他們.
他們兩個怎麼溝通.
我不知道你們夫妻.
尤其是夫妻兩家都信主.
我不知道你們家怎麼溝通.
這對老表他們怎麼溝通.
其實是用聖經溝通.
我引用這段.
你們回來就引用那段.
互相用聖經溝通.
在這裡 沒錯.
主耶穌的確是將他所說的說出來.
不過叫死人復活.
其實耶穌真的即使在當時.
都沒有做多少宗.
我們看約翰福音就看到.
拉薩路比較明顯.
那姻成寡婦.
女兒可能是.
沒有多少宗.
出名的都不是在這些位置.
但是耶穌特別提到.
叫死人復活.
窮人有福音傳給他們聽.
原來都是引用.
又是引用另外一個先知.
這個我們會熟一點.
就是以賽亞書.
引用以賽亞先知.
兩個地方是混在一起.
來引用事實上耶穌所做的事.
但亦都是在說.
先知預告尼賽亞將來會做的事.
這個我讀給大家聽.
以賽亞書第35章3至6節.
以賽亞這樣說.

$^{721}$引用耶華的說話.
你們要使軟弱的手堅壯.
武力的膝穩固.
對膽怯的人說.
你們要剛強不要懼怕.
看你們的神必來報仇.
必來施行極大的報應.
祂必來拯救你們.
第5節35章第5節.
以賽亞書.
那時核子的眼必睜開.
龍子的耳必開通.
那時瓊子必跳躍.
象鹿啞巴的舌頭必能過牆.
這是第一部分.
第二部分是耶穌所說的.
另外一段也加在一起.
這就是後面的.
以賽亞書61章的第一節.
主耶華的靈.
更出名.
因為主耶穌在.
《路加福音》記載.
在那撒勒會堂裡面.
一打開卷軸一讀.
就是讀這段.
61章的第一節.
以賽亞書.
主耶華的靈在我身上.
因為耶華用膏告我.
叫我傳福音給貧窮的人.
傳福音給貧窮的人.
後面還有猜險.
我醫好傷心的人.
報告被擄得釋放被囚的出監牢.
報告耶華的恩賢.
和我們上帝報仇的日子安慰.
一起悲哀的人.
耶穌是用了.
以賽亞書35章和61章兩段.

$^{761}$叫門徒去福施洗若翰.
施洗若翰聽到.
你猜他知不知道.
主耶穌是用了.
以賽亞先知所說的話.
因為若翰一聽.
這些事我早就聽過.
知道你在做什麼.
你重複這些事.
做什麼才行.
如果耶穌所說的.
若翰不知道.
他未必一定會聯想到以賽亞.
但問題是.
傳回去的說話.
若翰已經知道.
所以他就不會.
要麼就說.
大哥說這些做什麼.
我已經知道.
如果他聰明一點.
我相信若翰很聰明.
就會想多一點.
他在用以賽亞和我溝通.
用以賽亞溝通.
就很大問題.
我們說.
Not so much what is said.
but what is not being said.
很多時候重要的東西.
不是說了什麼.
而是沒有說什麼.
這裡剛才引用的.
以賽亞書裡面.
是有些很重要的東西.
是沒有說到的.
其實是同一個段落裡面.
若翰是對聖經滾瓜爛熟.
他是祭司之家成長.
所以他知道那段經文.

$^{801}$說什麼.
當然包括沒有說.
報告被擄.
得釋放被囚.
出監獄.
主耶穌說.
你不用怎麼想.
幫不到你.
你坐牢幫不到你.
應該沒有這個意思.
但起碼沒有了這部分.
但更加重要的.
兩段經文.
以賽亞書35章和61章.
裡面都提到同一件事.
主耶穌回覆若翰的時候.
都沒有提的.
就是35章第4節.
對膽怯人說.
你們要剛強不要懼怕.
但後面.
看啊你們的神必來報仇.
必來施行極大的報應.
他必來拯救你們.
拯救是報應.
審判.
復仇.
憤怒當中進行的.
61章後面那段第二節.
報告耶華的恩賢.
好啊好啊.
和我們上帝報仇的日子.
但報仇日子.
又和安慰一切悲哀的人.
是關連在一起的.
上帝的釋放和上帝的憤怒.
是連在一起的.
沒有說這兩段東西.
所以耶穌是否正正印證了.
施洗若翰的懷疑.

$^{841}$就是上帝的復仇.
上帝的憤怒.
就是上帝的什麼.
火.
就是他的火.
就是那個斧頭.
就是燒盡不滅的火.
燒盡康毗的火.
偏偏主耶穌回覆若翰的時候.
就沒有說.
說是瓊子 核子 龍子.
有福音傳給窮人.
都說了.
沒有說上帝的復仇.
上帝的憤怒.
上帝的審判.
你說就是囉.
所以耶穌最後才加一句.
凡不因我跌倒的.
就有福了.
第十一章的第六節.
這句話一向都令我很折騰.
為什麼要加多一句.
說給若翰做什麼.
現在若翰因你跌倒嗎.
若翰似乎是.
他懷疑主.
因為我認為你應該做那些事.
我看不到你做那些事.
即是不是我認為你是那個.
我應該跟第二個.
我跟兩個坐牢.
怎麼辦.
有少少不是那麼穩陣.
但主耶穌在這裡加一句警告說話.
你千萬不要因我做這些事.
你看不到某些事而跌倒.
耶穌所說的盡是耶穌私恩的事情.
上帝藉著他的兒子在地上.
施展都要能力的.

$^{881}$都要大能.
私恩的事情就沒有火.
沒有審判.
到底是不是耶穌真的欠火.
五行獻打.
耶穌贖水欠火.
抑或是.
你都沒有機會看.
在若翰無緣看到.
有沒有可能.
若翰其實無緣看到耶穌的火.
即是耶穌在這個階段.
他做事是沒有什麼火.
因為還未到時候.
火可能是在一個階段才來.
但不好意思.
你已經回到天家.
我的說法是.
若翰其實真的還未看到.
因為耶穌的火其實都是上帝的火.
是上帝的憤怒.
上帝的審判.
藉著耶穌在地上私行.
就等同上帝的恩典.
藉著耶穌在地上私行.
所以剛才讀以賽亞書61章.
第一節35章第四節.
兩度都講明上帝報仇的日子.
然後就講安慰一切悲哀的人.
上帝的恩典和上帝的憤怒.
都是由上帝的憤怒和恩典.
都在上帝那裡來的.
看你們的神必來報仇.
必來私行極大的報應.
審判報應.
他必來拯救你們.
拯救的恩典和審判報應.
都是來自上帝的.
所以我們說如果耶穌是私行上帝的恩典的話.
我們基本上必須要肯定.

$^{921}$上帝的報應.
對罪人的報應.
上帝的憤怒.
都是耶穌負責去私行的.
不過約翰就看不到.
所以他就懷疑.
你說約翰看不到.
那誰看得到.
你和我有沒有看到.
很久了2000年之後.
看不看到.
看不看到耶穌的那把火.
看到耶穌的那把火就舉手.
死了.
大家都好像約翰一樣.
是不是這一位還是等另一個.
還在等.
我們看看羅馬書.
這裡都沒有做powerpoint.
我懶了一點.
讀給大家聽.
因為那段經文其實都很熟的.
羅馬書.
我們看看保羅怎樣理解耶穌的服事.
三章二十三節其實很熟的.
前提就是耶穌的火其實就是主耶和華的火.
三章二十三節很熟.
因為世人都犯了罪.
怎樣.
虧缺了上帝的榮耀.
接著二十四節可能沒那麼熟.
如今卻蒙上帝的恩典.
因基督耶穌的救贖就白白地稱義.
恩典先.
白白地稱義蒙救贖.
第25節.
上帝設立耶穌.
作為耶穌這個人.
作為挽回祭.
是憑著耶穌的血.

$^{961}$藉著人的信.
要顯明他的義.
接著下面那句.
下面那句才是可圈可點.
因為他用忍耐的心.
寬容人先前所犯的罪.
重點呢.
三章二十四節最後這一句.
就是因為這兩個字.
保羅說上帝設立耶穌做挽回祭.
是憑著耶穌的血.
即是他的死.
藉著人的信.
要顯明上帝自己的義.
因為.
因為什麼呢.
有個問題.
有個說法.
然後將原因給出來.
因為上帝用忍耐的心.
寬容人先前所犯的罪.
意思就是說耶穌的血呢.
因為保羅說.
是上帝用忍耐的心.
是寬容人先前所犯的罪.
所以就以耶穌的血.
來顯明他的義.
意思就是說耶穌的血.
是不是他容忍的一部分.
就不是.
因為他先前.
他先前就是容忍的心.
但是來到耶穌那裡.
就不再容忍了.
這個是我們信仰裡面.
很重要的一部分.
我們看十字架.
這個十字架就.
不是最正規的十字架.
我們是透光.

$^{1001}$但都認得出是十字架.
我們見到十字架的時候.
你見到是什麼.
大家想像.
耶穌掛在十字架上面.
見到上帝恩典的舉手.
舉手.
見到上帝恩典.
耶穌掛在十字架上.
你見到上帝恩典的舉手.
見到耶穌憤怒的舉手.
一般都是不會舉手的.
但保羅在這裡說.
耶穌的死.
正正就是上帝憤怒的彰顯.
如果不是主在十字架上.
為什麼要說我的神啊.
我的神啊.
你為什麼離棄我.
他為什麼要這麼叫.
上帝離棄我.
掩面不看他的兒子.
他兒子的死.
是上帝對罪人的罪.
由於他兒子背上.
所有人的罪.
憤怒的表現.
所以他必須死.
一定要死.
必須流血.
所以十字架是上帝恩典.
對相信的人成為.
上帝救我們的恩典.
但對於不相信的人.
上帝對所有的人.
在我們相信之前.
仍然是代表著上帝的憤怒.
如果不是耶穌.
為什麼要死.
他的基督為什麼要死.

$^{1041}$上帝的憤怒全部倒在.
耶穌基督的身上.
甚至他的獨生子說.
你為什麼離棄我.
所以保羅說得很對.
他先前就用忍耐的心.
他用忍耐的心.
寬容人先前所犯的罪.
意思就是到了耶穌死的時候.
他就不再寬容了.
上帝的憤怒.
在耶穌身上發盡.
所以耶穌其實他的死.
就是上帝憤怒表現.
上帝的火.
在耶穌釘十字架的時候.
全部表現出來.
而這個火也是耶穌的火.
上帝對罪人的審判.
對罪的審判.
也都等同主耶穌基督.
對罪人對罪的審判.
而這件事約翰見不見到.
見不到.
在馬太福音第十四章.
約翰已經被殺了.
希羅底和肥力山的女兒薩羅米.
她辜滑了她現在的後父.
希羅底和肥力山的女兒薩羅米.
趁她生日哄到她開心.
結果就問你要什麼.
什麼都給你.
我要約翰的頭.
這樣就慘了.
希羅底和肥力山不想殺她.
但還是要殺了她.
所以約翰看不到這一點.
更加看不到後來上帝的審判.
對世界的審判.
意思就是說.

$^{1081}$現在你看楊佩音.
相信的人就有永生.
不信的人.
約翰怎麼說.
第三章 罪已經定了.
在耶穌身上.
耶穌同時有恩典和上帝的憤怒.
是同時集於一身的表現.
不過在約翰的時候.
耶穌還在地上宣講的時候.
是看不到憤怒的部分.
看不到憤怒的部分.
完全看不到.
耶穌罵幾句.
個個都不是什麼憤怒.
看不到.
第一次看到.
是耶穌的死在十字架上.
上帝的憤怒.
是去到那個位置.
我們開始就明十一章第十一節.
終於回來unpack.
我們說沒有最大只有更大那一節.
十一章十一節.
主耶穌說我實在告訴你.
凡女人所生的.
沒有一個比施洗約翰大.
每個人都是女人所生.
施洗約翰不是最大.
後面說天國裡面.
最小的都比他大.
那是什麼意思.
天國裡面全部都是女人.
不是 凡女子所生的男人女人都包括在裡面.
天國裡面就沒有人.
如果有人的話就矛盾.
除非天國裡面的人都不是女子所生.
只是天使.
沒有地球人.
那又不行.

$^{1121}$什麼意思.
女子所生的沒有一個比施洗約翰大.
原因是什麼.
即是說他比以前的大.
後來的都是.
包括我們都比他大.
in a sense.
in what sense是比他大.
先看他比之前的人大.
阿當厚比老祖宗大.
阿布拉罕 以撒雅各 以塞 耶利米 以西結 但以理.
全部比他大.
in what sense.
我要知道約翰的基本工作是什麼.
就是為彌賽亞開路.
他將門徒過給耶穌.
看上帝羔羊.
是要來除去世人的罪.
不過除去世人的罪是透過火.
他就不知道是透過死.
所以他是為基督開路.
將人指向基督.
我問你 在約翰之前.
有沒有人能夠好像約翰那樣.
指向具體的那位耶穌基督這個人.
有沒有人能夠這樣做.
沒有 因為耶穌還沒出世.
他們先知阿布拉罕都是從遠處.
希伯來書說他們從遠處看見.
天上的星 海邊的沙.
他的後裔.
不單止是以色列國的子民這麼簡單.
他們完全看不到他們說預言.
都是遠處看見.
以賽亞61章35章.
拉基的預言全部都是遠處看見.
他為基督做見證.
向以色列人為基督做見證.
但指不到向任何一個人.
他看不到的.

$^{1161}$唯一一個看到.
而又指向基督的.
連耶穌的門徒都做不到.
你會發覺是烏下烏下的.
指向基督的就是施洗約翰.
所以由這個角度.
施洗約翰的確大過所有女人生的人.
連我們都指不到耶穌出來.
耶穌在哪裡 你指一下.
在尼敦島 奧斯丁島.
在皇后大道中.
指不到給你耶穌 哪裡看見.
連我們都沒有辦法.
好像施洗約翰動篤企耶穌.
將人指向耶穌.
in that sense.
沒有一個人比施洗約翰大.
他能夠為耶穌做見證.
是沒有人能夠做到的.
事實上是沒有.
所以大勢是憑甚麼分.
就不是身家財產 學問地位.
起率大一點.
就不是了.
大勢之分在耶穌這個說法裡面.
單純是憑你能夠.
怎樣為耶穌做見證.
你的見證是多具體 多實在 多清晰.
來分大勢的.
但後面那句.
天國裡面最小的.
都比他還大.
最小的都大過他.
我們一定要知道.
上面那句和下面那句的邏輯.
一定要一樣.
大小的邏輯應該是一樣的.
同樣的原則來做比較.
否則真的不合理.
沒有人會明白耶穌說甚麼.

$^{1201}$如果上面那句是因著施洗約翰所做的見證.
是沒有人能夠做到的.
下面那句亦都應該.
是天國裡面的人.
個個都能夠做到見證.
是約翰做不到的.
所以約翰做到的事.
我們做不到的.
指向具體的耶穌的位置.
但我們做的一些事.
亦都是約翰做不到的.
甚麼事.
就是我們看到耶穌的火.
是約翰看不到的.
我們能夠為耶穌所做的見證.
亦都是約翰所做不到的.
所以兩家都做到一些事.
是對方做不到的.
在天國裡面所有人都知道耶穌死.
而復活.
都知道上帝在耶穌身上.
所展示的審判的火是怎樣一回事.
若不流血.
罪就不得赦免.
希伯來書作者這樣說.
所以我們全部看到.
聖靈與火.
不是五旬節那把火.
是上帝審判的火.
這是約翰看不到的.
以至他質疑.
但約翰能夠這樣指法.
亦都是我們做不到的.
所以他仍然大過我們.
但我們在另一個角度.
亦都大過約翰.
所以由頭到尾.
回來跟我們有甚麼關係.
就是我們必須要意識到我們的身份.
意識到我們全部都是基督的見證人.

$^{1241}$我們在天國裡面的大小.
就在乎我們怎樣為基督做見證.
當然既然是在乎我們怎樣為基督做見證.
最終可能這件事都跟我們無關.
因為所有的事都跟約翰無關.
跟我們無關.
直接就是因為基督.
是事情關乎主耶和華.
和他的彌賽亞.
而我們身為上帝的兒女.
天國子民的我們.
就是上帝和他的基督的見證人.
今日你和我都需要知道.
我們天國裡面的大小.
天國裡面的大小是一回事.
是否多幾間大屋,小屋.
賞賜之類.
我們不知道大小還有沒有其他意思.
但起碼就著做見證這個位.
我們知道在上帝眼中是有大小的.
我們能夠在地上.
為主做見證.
在上帝,基督的眼中.
我們就視為大.
如果不是就視為小.
這個不是一個恐嚇.
這是一個激勵.
是上帝講到施洗約翰的時候.
正因為他所能夠做的見證.
是沒有人能夠做到.
所以他視為大.
而我們今日必須要意識到.
我們所經歷的.
即使不能夠施洗約翰做見證.
但我們主的生命.
在我們自己生命當中所產生的.
恩典,變化,經歷.
是我們自己才能夠講出來.
你不能夠指人向皇后大道中站某個人.
那個是基督.

$^{1281}$你不能夠這樣做.
但你卻能夠藉著你的生命.
將人帶到上帝的面前.
見證基督在你身上的恩典.
這個是我們看到.
我們上帝赦罪的平安.
我們悔改的見證.
人要靠什麼而生存.
靠上帝口中一切所講的話.
靠信,恩信而得生.
這些全部都是我們能夠做到的見證.
而這個也是我們在天國中身份之大的最基本的原因.
所以保羅在《羅馬書第十章》十六節講得很有意思.
我不以福音為恥.
為恥不是在地鐵派傳單.
很醜啊.
我很怕被拒絕.
現代人很怕被拒絕.
不是在講這件事.
保羅面對的是坐牢的羞辱.
其實是.
我不以福音為恥.
就是指著我願意為基督做見證.
為基督的福音做見證.
因為這個本來就是上帝的大能.
要救一切相信的.
然後他就講原來上帝的憤怒.
第十八節就立刻又講上帝的憤怒.
講完要救一切相信的人.
第十八節就立刻就講上帝的憤怒.
向地上所有不權不義.
阻擋真理的人.
發出來.
所以這個是我們要做的見證.
而這個是直接連繫著我們在天國裡面的身份.
我們都是見證人.
彼得前書二章九節說.
我們被簡書族裡有君尊濟私.
聖潔的國度無礙子民.
說要宣揚那位召我們出黑暗.

$^{1321}$入奇妙光明者的美德.
所以就讓我們在這一年無論剩下多少時候.
都能夠膽壯的.
剛強的.
有把火.
能夠為基督做見證.
因為我們見過上帝把火.
見過基督把火.
如何為他做見證.
一齊低頭祈禱.
天父我們在你面前謙恭屈膝.
將自己再一次擺上保羅畫.
將身體獻上當作活祭.
我們其實你所要求的不是要整死我們.
基督已經在十字架上.
為我們留下寶血.
為我們成為贖罪祭.
你要的是我們的生命.
你要的是我們對你的忠誠.
在地上腳踏實地過你要我們過的生活.
天父加力量給我們.
叫我們明白我們的身份.
叫我們明白我們身份的尊貴.
在你眼中我們是捍衛寶貴.
因為我們是你兒子的忠實見證人.
加力量給我們.
叫我們手所作的功能夠結出果子.
奉主耶穌基督命求.
\newpage



\section{帖撒羅尼迦前書}
\label{sec:VvxSAXf_Lww}
\textbf{2024-05-09 周四晚查經班《帖撒羅尼迦前書(6)》 講員:曾裔貴牧師}
\newline
\newline
連結: \href{https://youtube.com/watch?v=VvxSAXf-Lww}{\texttt{https://youtube.com/watch?v=VvxSAXf-Lww}} ~~~~ 語音日期: 2024-05-10
\newline
\newline
\hyperref[sec:1QRoYKXVPvQ]{\small{< < < PREV SERMON < < <}}
~
\hyperref[sec:index_chronic]{\small{[返順時目]}}
~
\hyperref[sec:index_scriptual]{\small{[返順卷目]}}
~
\hyperref[sec:BHhp4UIkEH0]{\small{> > > NEXT SERMON > > >}}
\newline
\newline
$^{1}$好 我們一起祈禱.
感恩我們再一次可以聚集.
無論在網上,線上的.
或者是在現場.
還有當我們成為了錄音之後.
在不同地方的弟兄姊妹.
甚至乎可能是不同宗派的教會的弟兄姊妹.
可以聽到主的道的時候.
可以彼此的勉勵.
我們真的很盼望主的聖靈.
在我們大家的心靈裡面去工作.
使我們明白.
少徒保羅在雅典.
所要向傑瑟羅尼加城的弟兄姊妹.
傳遞的那個真理.
尤其是提到主的再來.
是講了很多次.
前書,後書都講.
為甚麼在第二次的旅行報道.
就已經這麼強調主再來的真理.
這個都值得我們很認真去思想的.
求主給我們對於主的再來.
不是一個天方夜譚.
不是一個都不在我們的心思意念裡面.
一個不會出現的真理.
我們能夠常常在這樣的教導下去生活.
使我們的生活更加具體.
更加有建設性.
更加合乎主的心意.
願神帶領我們每一位.
我們這樣禱告奉耶穌基督的名.
阿們.
我們上一個星期第三章.
近尾最後一節就提到三章的十三節.
我們又用了一點時間.
去跳一跳.
就去到第九節到十二節.
講到關於主的再來的情況.
應該是十三到十八節那一段.
我們待會兒會再補充的.

$^{41}$但是我們就涉及了一段很重要的經文.
就是一到第九節.
就探討了愛和合一.
或者說愛和聖結.
這兩個重要的生命的基督徒的特質.
是保羅在第三章結尾.
和第四章下半段講到主再來.
然後夾雜的一個這麼重要的教導.
是一個性教育.
都很露骨的.
似乎又不是責備性的.
去罵跳上羅尼加城的人.
而是他們要明白神的旨意.
就是說他們第一章的時候.
保羅說他們是離開偶像敬拜真神.
所以離開偶像.
就是說他們以前都是拜偶像.
是有機會去涉及那些供奉偶像.
和那些妙技的.
哥林多就一定有的了.
跳上羅尼加都很近哥林多的.
我想那個時候他們應該都很繁盛.
那些是大都會.
所以就不是說很遠很遠的.
都是在地中海附近的地方.
如果坐船就更加快過馬車.
或者是其它的陸路.
似乎當時他們那些.
因為都已經是羅馬帝國了.
所以那些的道路網其實都是很發達的.
所以哥林多所面對的問題.
可能跳上羅尼加城的人一樣會面對.
所以在整個的教導裡面.
第四章才是正式對教會的一個教導.
怎樣做基督徒.
就好像長話短說.
特別聚焦是聖的教育.
所以又很值得我們去想一想.
保羅為甚麼會這麼煞有介事呢?.
明明他自己很掛心頂 姐妹.

$^{81}$信了主四個星期.
保羅自己就要跟同宮.
跟阿西拉就要離開.
去到雅典之後稍為平靜了.
就打發提摩泰回去.
去了解教會實際現在的那個屬靈的光景.
有沒有被人來到去危言聳聽.
逼迫抓你坐牢.
令到你不信耶穌.
回來的好消息就是提摩泰說.
他們仍然保持這份信無礙的信心.
所以保羅就因為他們這樣的情況.
保羅就說我們活了.
我們那個內心很歡喜.
很有平安 很喜樂.
因為弟兄姊妹能夠在患難之中.
都能夠站立得穩.
跟著再勸勉他們要來到去繼續成長.
跟著保羅說如果我能夠見到你們的面.
我就會來到去加添你們信心的功夫.
就是十二節.
「又願主」這個三章是十二節.
「又願主叫你們彼此相愛的心.
並愛眾人的心.
都能夠增長充足.
如同我們愛你們一樣」.
跟著就是「好使你們當我們主耶穌.
和祂眾聖徒來的時候.
在我們父神面前.
心裡堅固成為聖潔無可責備」.
這個「心裡堅固成為聖潔無可責備」.
正是第四章的第三到第九節的經文的重點.
怎樣能夠在神的眼中將來不受責備呢?.
我們上一次特別提到.
十三節說到關於睡了的人.
即是死了的聖徒.
和到底那個復活次序.
第三章這裡我們可以再補充一點.
就是主耶穌來的時候.
已經很多次說是和眾聖徒一起來.

$^{121}$眾聖徒一定是在現在的墳墓.
靈魂現在在主耶穌的懷裡.
雖然我們不能夠具體.
到底那個現在的存在是怎樣呢?.
因為沒有人去過回來告訴我們.
通常那些死了又死不去的那些.
又說見到人搶救他.
回來就繼續很有盼望去做基督徒.
最多是這樣而已.
沒有一個很實實在在.
真的去過天堂.
然後回來這樣的人.
所以我們都只能夠透過聖經.
給我們有限的啟示.
來讓我們知道.
有可能未必是不想讓我們知道.
而是我們現在地球的人.
根本不會明白.
到底那個多維度的空間.
是怎樣的一回事.
所以我們怎樣說我們都不會明白.
所以只能夠說那個概念.
是有一個很實在的地方.
是永存的 是永恆的.
是再會有一個復活的身體的.
不會再有疾病痛苦的了.
而且是與神永遠同在的.
就是這麼多了.
不可以說得太詳盡.
因為再遙距天堂.
我們都未必能夠掌握到那個具體.
我稍微關掉靜音.
這是給我們有一個很重要的啟示.
首先就是眾聖徒.
以瑞的眾聖徒.
是一定會與耶穌一同回來.
四章十三節就說那個次序了.
就是這些死去的聖徒.
將來一個一個就會復活.
就會與主耶穌一同回來.

$^{161}$所以這裡就告訴我們.
原來我們在有限的地上的人生.
信了主之後.
我們要有一個交承職表的時間.
就是使得我們在主耶穌.
和祂眾聖徒回來的時候.
我們好像拿了獎一樣.
如果我們能夠感受到那個榮耀.
而且那個榮耀不是一種很虛假的.
是你經過了很認真的努力.
而你能夠得到的.
那個榮耀就很大.
所以在我們的人生.
將來在父神的面前領獎.
我們能夠心裡堅固.
我們的信仰很堅固.
我們的信心很穩固.
不會被一些異端邪說來動搖.
貼上羅尼加城的頂指會.
成為聖傑.
似乎是當時初期教會.
頗需要面對的一個真理.
頗需要直接確認那些偶像敬拜.
所以我們將來有機會看哥林多前書.
很明顯針對那些拜偶像.
而拜完偶像之後.
就會有一個集體的淫亂.
就在那個廟裡面發生的.
和那些妓女.
那些是廟妓.
據聞看一些書.
有五百個在那個廟裡面.
如果你看回羅馬帝國古代的西片.
它都會有這樣的描述.
所以當時希羅的文化.
都會有這種情況.
成為了他們一個.
好像拜偶像和淫亂拉上關係.
成為聖傑無可責備.
即是說將來到底是要得到讚賞.

$^{201}$拿到那個冠冕.
還是會有虧損.
會有那個責備.
所以不是離開世界.
就好像虛無縹緲.
沒有存在.
如果是這樣的時候.
真的大家都不需要努力.
甚至回不回教會都可以了.
你說自己是基督徒都可以了.
個個都去到天堂了.
似乎不是這樣.
保羅告訴我們救恩是無條件.
任何人憑信心接受主耶穌.
就可以得到永恆的救贖.
就可以得到永恆的天家.
但是還有一個就是.
神的面前會有一個獎賞.
和會有一個佛的.
不知道那個到時的狀態.
那種虧損是怎樣呢?.
那種是否還有一種永遠的羞恥感呢?.
我們不可以太掌握到.
當然肯定就不是那些去到永恆火湖裡面.
那種哀哭切齒.
那種是不會再有回頭的機會的.
永遠沉淪的那些.
就不是那個狀態.
但似乎這裡是說基督徒.
明顯是說基督徒.
所以到第四章.
差不多這章聖經的最重要教導.
就在第四章這一段.
當然第五章都有的.
因為你可以想像到.
說了這麼多主在來.
又會帶來一個反效果.
就是會有一些好食懶飛的基督徒.
不工作等主耶穌回來.
等運到.

$^{241}$就會扭曲了主在來.
其實應該要我們更加積極面對我們的生活.
把握著每一天的生活.
來去工作.
來去實踐.
神要我們人生的照明.
當然不是說不可以娛樂.
不是說不可以運動.
不是說不可以休息.
我們身心靈要調整.
要平衡.
但是說的是.
他們真的有一個扭曲了的神學.
就帶來一種好像義正言辭那樣.
就是不用做了.
等主耶穌回來就行了.
那就是第五章的警告.
那就真的會在教會裡面有這樣的人.
本來他已經不敢這麼大聲了.
你說這麼多主在來.
還用做這麼多事嗎.
不如就一起等候耶穌回來.
那就可以改變一切.
但是正正就是我們不知道什麼時候回來.
終於二千年後了.
都還沒有回來.
所以這些經文都是永恆地.
給我們每一代的基督徒要警醒.
第四章我們開始說這個段落.
一到第十二節是一個段落.
接著十三到十八節.
就是再一次補充主在來的次序.
這個是我們上個星期有機會說了一點點的.
所以你可以想像到.
第三章的十三節.
說主的在來.
突然之間又不再說了.
又插入了一個教導.
甚至說的是神的旨意.
不是保羅的個人見解.

$^{281}$不是他個人的道德觀.
而是神面前的性教育.
以及怎樣待人.
要有那種的愛心.
我們讀第四章.
這個教訓都相當重.
還有保羅好像很強調那個新般.
四章第一節.
弟兄們 我還有話說.
我們靠著主耶穌求你們.
勸你們.
你們既然受了我們的教訓.
知道該怎樣走.
可以討神的喜悅.
就要照你們現在所行的.
更加勉勵.
到底這個現在所行的.
是他們的信望外繼續增長.
知識繼續增長.
還是有其他的一些.
在實際的生活裡面.
他們的操守.
要怎樣去做呢?.
但是保羅其實可以有很多東西教導.
但是他突然間在第三節.
就收緊了那個.
他講的要更加的勉勵.
更加的來到要討神喜悅.
就將第三節.
將神的旨意.
要關乎聖潔的生活.
聖潔的生命.
來到去講出來.
這個是一個命令.
告訴我們知道.
到第二節了.
你願曉得我們憑主耶穌.
傳給你們什麼命令.
他又不詳細講了.
耶穌傳給我們什麼命令呢?.

$^{321}$愛人如己.
彼此相愛.
彼此洗腳.
還是講關於聖潔呢?.
耶穌好像很少講.
關於這麼強烈的聖潔的要求.
你看科幻書.
都不是很強調的.
但是你接著就看到第三節.
講完耶穌的命令.
接著第三節.
神的旨意.
這些差不多是寫得很清楚.
很白的.
每一個基督徒.
每一個時代.
每一個宗派.
任何一個基督徒.
不論你信什麼的.
什麼教會.
你都要去遵守.
因為保羅說這是神的旨意.
是最高的權威.
等於十誡.
等於十誡的教導.
神的旨意就是要你們成為聖潔.
遠避人行.
要你們各人了解.
怎樣用聖潔尊貴.
守著自己的身體.
不放縱私欲的邪情.
像那不認識神的外邦人.
似乎這裡不是一個責備性的口吻.
是一個正面的性教育.
一個基督徒要這樣理解聖潔.
這樣的意思.
就不像哥林多前書.
整個口吻是責備性的.
就是說有人已經乾犯了.
這裡就是說你們信了主耶穌.

$^{361}$你們現在要有一個新的.
好像新的老闆.
你就要怎樣看神.
在你的生命裡面的那個膚照.
這個膚照.
怎樣化做你的生活裡面具體的.
你真的能夠碰到的.
你的身體的這個人呢.
所以特別強調那個身體.
特別強調要守住自己的身體.
守住.
似乎這個是.
比較.
就是負面一點的.
就是遠避.
積極的.
不消極的.
剛才那個可以叫消極.
不是負面.
很消極的就是不斷避開避開.
但是更加基督徒的人生的積極就是.
你要在自己的身體.
就是說原來是一個整體來的.
我們有靈魂.
我們有肉身.
原來是一個整體.
英文叫holistic.
是一個很整全的生命.
在神的眼中的.
就是說我們的.
應該這麼說.
一個人叫being.
就是有知識.
有情感.
有意志.
知情義.
這三樣東西構成一個人.
動物就不是了.
所以我們人類才是being.
神的造的人類.

$^{401}$可以與神溝通的.
所以我們那個will.
我們那個emotion.
我們那個情感.
我們那個內心.
當然更加是我們裡面的靈.
是整個人的.
是有相關和有影響的.
所以這裡就.
保羅一直這樣提.
我先讀完.
特別強調與外邦人.
與外邦人.
這個與外邦人就不是一般.
你是希臘人.
那你就是外邦人.
你是猶太人.
你就是神的選民.
不是這樣分.
是強調不認識神的外邦人.
就是說.
是信耶穌和不信耶穌的對比.
所以不是那個民族.
不是那個教會的宗派.
是信有聖靈.
或沒有聖靈的基督徒.
沒有聖靈的人.
就是不認識神的外邦人.
要告訴我們.
第六節.
不要一個人在這件事情上.
越分.
越分就是過了界.
過了界就是說原來有個標準.
欺負他的弟兄.
因為這一類的事主必報應.
剛才說命令.
主耶穌的命令.
神的旨意.
這裡說報應.

$^{441}$我們都很怕報應.
中國人都很怕報應.
小心你收尾那兩年.
小心你什麼什麼.
我們有時候會詛咒人.
撲在街上.
很多時候我們會這樣.
我們會怕報應.
就算不信神的人都怕報應.
所以鬼片永遠有些人不喜歡看.
又很想看.
這個就是.
但是這裡聖經新約告訴我們.
將來不單止是有一個.
在神面前的賞賜.
而且是有一個報應.
這裡說正如我預先對你們說過.
又次次地祝福你們.
祝福這個字就是一個命令的字.
charge的字.
接著第七節.
神照我們.
再一次說那個初心.
一個基督徒.
為什麼會成為基督徒呢?.
成為基督徒之後.
其實應該怎樣行動呢?.
怎樣在神的面前.
來有那個表現呢?.
神照我們.
本不是要我們尖而污穢.
乃是要我們成為聖傑.
所以乃氣絕的.
不是氣絕人.
乃是氣絕那次聖靈給你們的神.
所以會有一個危機.
就是說有其他的學說.
如果不是根據保羅這樣的教導.
可能會滲入弟兄姊妹的心靈.
因為保羅不在.

$^{481}$因為教會現在沒有了那個領袖.
可以繼續教導他們.
很容易就會有那些.
似是而非的道理.
就入了貼壽羅尼加.
所以就會有這種.
不知不覺間慢慢放鬆了.
那個對聖傑的看法.
就容易會帶來一種.
好像溫水煮蛙那樣.
所以保羅告訴我們.
有些人是故意的挑剔.
故意的說保羅.
是不是你太過度強烈的禁慾主義呢?.
對這件事你認為是徹底的污穢.
誰都不可以做這些關乎性的事情呢?.
是不是因為你自己不結婚.
所以你就這樣呢?.
所以保羅都已經說了.
如果那些氣絕的.
其實不是氣絕我.
氣絕人的意思就是說.
不是你氣絕我的教導.
其實是氣絕了主耶穌的命令.
神的旨意.
和你難免將來要面對主的報應.
所以這個都是我們必須要來看.
由第三到第八節.
說了六節聖經.
來到說關於性的教育.
負面.
消極的就要禁止.
積極的就要看自己的身體.
是一個乾淨的器皿.
心思意念都要在神的聖靈的同在下.
活得坦蕩.
看這件事是一件在神的眼中.
是神賜給我們人類的一個恩賜.
不單是繁衍後代.
傳宗接代.

$^{521}$所以說到性方面.
其實保羅就在臨前的第七章.
就很詳盡的說了.
到底男女應該怎樣看性的功用.
很詳細的講解說.
好像我們小時候.
讀小學的時候.
突然間有一天.
其實我們不知道.
老師又不說.
女孩都走了.
不見了.
然後怎樣呢.
我小時候才知道.
原來分開男女.
教你男孩的器官是怎樣的.
女孩是怎樣的.
好像這些教導.
給一個正面一點的價值觀.
那時候有這些健康教育其實都好的.
其實這些如果越來越能夠正面的去探討.
就會好.
基督徒其實聖經既然都有這樣的教導.
尤其是當時的社會.
是充斥著一個強弱的.
很剝削的一些貴族.
奴隸.
男的好像操控著女性.
還有一些強大的帝國.
去南征北討的羅馬帝國.
那些欺凌的文化.
又有很多假神.
那就好自然.
那個神學會賦予了一些價值觀.
所以聖經的教導.
其實伴隨著對抗背後的價值觀.
背後的神學.
所以如果我們中國人.
我們不說太多.
我們就變成.

$^{561}$令到這件事越來越神秘.
基督徒很少.
我從來沒有在港台聽過人講聖教育.
即是崇拜.
講這些經文.
正面的教導.
怎樣看待自己的身體.
很少.
如果對年輕一代.
可能會有更大的一個神秘.
那就會有個空隙.
被一些似是而非的學說.
就會進入了那個心靈.
那就會帶來可能會有一種扭曲性的.
對神的旨意的誤解.
講了那個purity聖潔.
或者Holiness.
再重一些的字眼.
接著到第九節.
就是去到關乎相愛.
論到弟兄們相愛.
當然包含了是全個教會的弟兄姊妹.
「不用人寫信給你們.
因為你們自己忘了神的教訓.
叫你們彼此相愛」.
所以一個信了主的人.
其實差不多是與生俱來.
一信耶穌.
信得好.
所謂信得好其實就是給耶穌的愛感動.
悔改.
不是交換利益.
不是為了.
所以永遠如果有些人本來是想著.
哪些利益才信耶穌.
神都會給他有個空間的.
將來就要他交賬.
他一定要面對一些的臨界點.
原來去到信耶穌是沒有利益的.
那一刻他都相信.

$^{601}$那一刻的信是真的.
因為很多時候.
有時候會有些錯誤的觀念進了門.
我也說過.
我中二.
不是 升中四的時候.
我的學校已經很差了.
有一間在深圳某區的基督教學校.
有三個同學.
現在很大了.
在溫哥華.
突然轉校過來.
那就.
你想想.
來我的學校當然是很差的了.
成績等等.
他們原來有上團契的.
教會團契.
就問他為什麼.
他說認識女孩子最容易的.
那時候我的觀念就是很鄙視基督徒.
因為我根本沒有信耶穌.
中學的時候.
原來是這樣的.
上教會原來是這樣的.
這幾個同學都還沒有信耶穌.
到現在都還沒有信耶穌.
一來都很少機會再見.
那就.
當然我有些責任.
要為他們祈禱.
但是.
這就告訴我們.
我們的入門.
所以其實一個信耶穌的弟兄姊妹.
本來就告訴我們.
其實你們自己忘了神的教訓.
這裡似乎在說的.
不是一個道理這麼簡單.
是有了聖靈.

$^{641}$有了聖靈.
聖靈會在你的心裡.
將神的愛挑冠.
所以.
一個信了主耶穌的人.
必定會從神而生.
愛生他之神.
和愛神所生的兒女.
這是《約翰一書》三章的經文.
主為我們寫明.
三章十六節.
《約翰一書》三章.
我們從此就知道何為愛.
我們也當為弟兄寫明.
這就告訴我們.
是很實在的.
所以本來這裡說完.
關乎到聖潔.
聖潔原來就是.
如何看待正面的.
神學的.
看待自己的身體.
自己的存在.
在這個世界.
要很平衡.
不要偏激.
要恰如其分地.
看待自己的生命.
然後.
好像拆開了.
其實不是拆開的.
Galto姐妹可不可以龍胸呢?.
為何要這樣做呢?.
會不會變成了.
一種很扭曲的.
可能現在的技術很成熟.
沒有什麼風險.
甚至可以很漂亮.
但是是否需要這樣呢?.
是否要這樣取悅對方呢?.

$^{681}$諸如此類.
這些其實是倫理的問題.
其實歸根結底.
就是我們如何看待聖潔.
如何看待我們這個.
用正確的神學.
去看自己的身體.
身體是靈魂身體.
是整個人.
是你的全人.
所以這裡是告訴我們.
我們需要.
這樣去理解.
神起初對我們.
這個人的結構.
例如我是.
1.69 也想高一點.
1.72就好了.
但是高不了 只有這麼多.
8歲就已經抬珠了 抬得越來越矮了.
那個發育時期.
你根本沒有機會讓你發育.
又不能打籃球.
又不能去玩 在學校中學的階段.
放學後就一定要回家.
根本不准你有任何的.
社交 任何的.
運動 參加學校的運動.
沒有 沒機會.
很自然就在這樣的.
貧乏下 你就自然.
那個長高沒那麼高.
會是一個介意的.
因為同學很高.
你矮呢.
其實又有些影響.
我想.
我想真的每個弟兄姊妹.
或多或少.
會介意自己的身體.

$^{721}$這段聖經.
如果再正面一點去看.
就是如果我們.
真的看自己的身體.
是神的創造.
和神的.
預定的預知.
我們應該.
欣然的接受.
這個不完美.
譬如我們有時成長.
聽過.
有些女孩子.
為了愛那個男孩子.
甘願去做賣人.
出賣自己的身體.
明知那個男孩子.
是姑爺仔.
騙她錢的.
她都為了這樣去做.
可想而知.
她看待自己的身體.
當然她未信耶穌.
更加.
一定是她有很多的坎坷.
家庭的沒有溫暖.
以為.
在男朋友身上.
會找到.
任避.
她明知是出賣她.
騙她.
她都情願在她的.
騙下.
生活.
可想而知.
她根本沒有一個很正確.
看待自己人生的價值觀.
所以這段聖經.
非常寶貴.

$^{761}$告訴我們.
我們要這樣積極.
去看待自己的存在.
在神眼中的你.
或者我.
如果我們在.
這個整全的神學裡面.
這樣去看待自己的時候.
有些突然間的患難.
或者有些人.
是長期的皮膚病.
不說太多.
是哪個名字.
其實看著都辛苦.
那種辛苦是.
比我想像的.
不斷的拗.
不能停止.
都帶著她的辛苦.
一定是那個價值觀.
有很多東西.
自然就會不斷的拗.
拗來就會硬.
其實是很困擾的.
我想.
為數都不少.
現在有這麼多化學的.
清潔劑等等.
社會有污染.
可想而知.
我們或多或少.
未信耶穌之前.
怎樣看待自己.
這個人.
這個自我.
一定受著很多的社會觀念.
社會的思潮.
價值觀來影響.
一定會有一定的扭曲.
不會有一個人.

$^{801}$十全十美的看自己是一個完美的人.
可能有的.
但是.
基督徒就.
應該不是這樣看.
我們到下面這一段.
第十節.
你們向馬其頓傳遞的眾弟兄.
固然是這樣走.
但我勸弟兄們.
要更加的勉勵.
當時整個馬其頓有多少人呢?.
整個馬其頓省.
即是說由帖撒羅尼迦.
一直到雅典.
都是一個馬其頓大範圍的地方.
當然他們當時.
那個行政區分了多少呢?.
哥林多就是.
阿該亞省.
但可能都是統稱是馬其頓一帶.
其實很多.
很大的地方.
現在的保加利亞.
再進一點的.
就是克羅地亞.
即是東面的東歐的地方.
如果去到帖撒羅尼迦.
一直打下去.
去到哥林多.
去到雅典.
這些地方欺負馬其頓.
一帶地方其實很大.
保羅說你們在.
馬其頓人的面前.
已經是固然這樣走.
但保羅說.
要更加的勉勵.
更加勉勵的.
要來彼此相愛.

$^{841}$即是表現出.
我們是沒有門戶之見.
我們不會有.
我們自己的成見.
我們不會有我們自己.
一個民族很.
求滅的.
那種保守的主義.
排他的.
一定有的.
每個民族.
中國人.
都有我們的排他性.
你就可以想像到.
這裡讓我們知道.
保羅這次在這一段.
9到12節.
是說到關乎到彼此的相愛.
愛這個課題.
說的是實的愛.
在我們的生命裡面.
怎樣去愛其他人.
然後第11節.
又要立志的作安靜人.
辦自己的事.
親手作工.
正如我們從前.
所吩咐你們的.
叫你們可以向外人.
行事端正.
自己也就沒有什麼缺乏了.
原來在這個.
彼此相愛的裡面.
有一個這樣的.
假設就是.
不要以為藉著.
愛去剝削對方.
去令到.
對方受到剝削.
愛背後.

$^{881}$就好像是互相分享.
所以原來.
每個人保羅就.
很平衡地說了要.
這麼博愛要這麼彼此相愛.
甚至乎誇了地域.
誇了大家不同的.
可能當時的民族.
希臘人和馬其頓人.
都有些不同.
可能有自己的城邦的特色.
因為當時也有很多的城邦.
統稱是羅馬帝國.
但是有不同的人種.
你可以想像得到.
大家一定有一定的文化的差異.
小小的不同.
但是現在在一個大同裡面.
就是神的兒女.
以神為主的.
神的兒女.
以神的愛彼此相待.
但是又不可以.
沒有自己的責任.
自己應該要盡的義務.
這就告訴我們.
要立志地.
作安靜人.
不多是非.
不多一些.
將一些.
可能不盡不實的東西.
來傳遞.
要辦自己的事.
親手的作供.
即是說原來.
彼此相愛的.
能夠達至.
美滿.
在神眼中的彼此相愛.

$^{921}$原來是每一個個體.
基督徒都要自己.
是殷勤的.
來守住.
不是身體.
這麼簡單.
是言語.
和勞動.
勞動.
基督徒積極的勞動.
我們再一次說.
就是在宗教改革.
帶來的.
清教徒的運動.
其中最重要.
就是工作的倫理.
就是要親手作供.
要每個人.
勤奮的工作.
在神給你的職業的照明.
你來到去.
殷勤的發展.
去發揮.
使到.
那個的照明.
所謂的Profession的照明.
就能夠.
可以榮耀神.
所以就是這一節聖經.
就是告訴我們.
彼此相愛的.
能夠達至.
就不是彼此在這裡.
來到去.
你要愛我.
我做不到事.
你要幫我.
這樣就是一個剝削.
就不是彼此相愛的.
真正背後的精神.

$^{961}$沒錯.
當我們發現.
那個人不是到處去問人.
要愛他,要施捨給他.
要周濟給他.
而是發現他已經盡了他的力.
而他也有困難.
缺乏.
那個時候我們就要去.
愛,幫助這樣的人.
所以剛才我們看的.
《約翰一書》的三章十六節.
大家記不記得?.
主為我們寫明.
我們從此就知道.
何為愛.
然後我們.
當為弟兄寫明.
接著十七節.
有世上財物的.
看見弟兄窮乏.
卻塞住那個憐恤的心.
愛神的心.
怎能夠在他裡面呢?.
這是約翰人這樣說的.
所以原來這裡的彼此相愛是很平衡的.
正常的情況.
我們要在彼此相愛裡面.
但是不可以去剝削別人.
我們要發展.
神在我們生命裡面的那個照明.
神給我們甚麼工作.
我們要在我們的崗位上.
盡忠.
自然保留告訴我們.
當我們向外人行事端正的時候.
自己也就沒有甚麼缺乏.
很平衡地.
彼此相愛又大家沒有缺乏.
原來建立在一個這樣的互信.

$^{1001}$一種的愛裡面.
一種彼此的.
來盡自己的責任.
這樣的一種彼此相愛.
所以你會看到是一個很平衡的一件事.
不如可能大家有些困惑.
不如大家揭一揭剛才看的.
《若人一書》那段經文.
真的有使徒這樣教導的.
即是說原來他們理解的彼此相愛.
是有這樣的含義在其中的.
《若人一書》三章.
十三節開始.
「弟兄們!世人如果恨你們,不要以為稀奇..
我們因為愛弟兄,就曉得是已經出死入生了,.
沒有愛心的仍住在死中..
凡是恨他弟兄的,就是殺人..
你們曉得凡是殺人的,沒有永生在他裡面的.」.
第十六節「主為我們寫命,.
我們從此就知道何為愛,.
我們也當為弟兄寫命.」.
十七節「凡有世上財物,.
看見弟兄窮乏,卻塞住憐恤的心,.
愛神的心怎能存在他裡面?.
小子們!我們相愛,不要只在言語和舌頭上.」.
肯定是回應十七節.
只有說話 只有說話.
所以是要具體的行為和誠實.
要這樣一個具體的行動.
所以彼之相愛原來有這樣不同的層次,技巧.
還有你要審時度勢 每個人都不同.
所以回到保羅這裡.
他的教導是在正常的情況下.
我們的彼此相愛是端賴於每個基督徒.
都是行事在神眼中.
行事端正.
即是我每天有一個很強烈的照明.
去為主而生活.
那個生活就一定是.
在工作上就是你的崗位.

$^{1041}$所以問自己.
神給你甚麼崗位?.
你要盡責,盡忠.
這樣你便能夠在神的供應下.
你便能夠開始有這種富裕.
然後你便可以開始在教會裡面.
在一個最好的場景.
有很多意外我們都可以去發展.
所以近年連香港.
似乎都是的 先進的國家都是的.
會有很嚴重的貧富懸殊.
所以那些先進的國家.
反而那些扶貧的工作的NGO機構很多.
譬如我有一次去三藩市.
我都見到有個救世軍.
就在Downtown.
真的有些人去吃免費午餐.
似乎都是這樣.
所以教會就更加是一個重要的場景.
我們去幫一些.
他真的不能夠去行事端正.
去得到可能是他的溫飽.
所以是要很有智慧去運用.
神給我們的金錢.
去彼此相愛的.
所以如果我們幫助一些懶惰的.
不了解自己的人生照明的.
這樣是一個對他來說.
他的剝削就會助長他繼續的.
有這種走精扁的心態.
和更加會令到彼此相愛.
是會扭到不能夠真正的落實的原因.
但是剛才約翰所講的.
對每一個基督徒.
你已經是一個行事端正的人.
你已經在神的祝福下有那個賦予.
而你眼見到那個人在你面前.
神已經性命感動你了.
你都塞住這個連戌的心.
你就真的不能夠說得上.

$^{1081}$你是在實行叫你彼此相愛的人生.
所以約翰都是很重的去教導.
當時異不所教會裡面的弟兄姊妹.
不要因為可能試過失敗過.
被人騙過.
就塞住了那個連戌的心.
可能一次被蛇咬.
終生都怕繩.
會有這種錯誤的念頭.
但是要保持那個赤子的心.
就要兩方面的平衡.
所以這一段就講到我們彼此相愛.
先是提到個人我們自己要有這樣的裝備.
剛才講的就是我們活在這個地球.
神生我們出來.
哪年哪月哪日你出生.
一直成長.
你的成長都是在神的旨意.
譬如我真是1.69.
那個高度不能再高了.
你拉長都拉不到.
要接受.
可能幾十年都不接受.
慢慢終於你明白了真理.
原來神愛我這個人.
不是因為我1.78 1.幾.
是因為我這個人神就愛我了.
你慢慢接受.
滿足於自己的存在.
在自己的存在裡面.
繼續發展.
那個聖潔的觀念和成長.
愛一樣.
彼此相愛是一個很大的課題.
對每一個人我們都要愛.
但原來愛裡面.
都有這樣的一些小心的界線.
有些人是剝削彼此相愛這個觀念.
他要來到特意拿便宜.
我們就要有智慧地.

$^{1121}$我們就要小心.
但保羅這裡就強調.
我們既然有這麼強的大家庭.
馬其頓整個省的人.
都見到弟兄妹姑姑這樣走.
要更加勉勵.
即是怎樣呢.
我們更加要因為這樣的緣故.
我們自己平日就要好好地裝備自己.
就是做安靜人.
親手做工.
和就能夠在神的面前行事端正.
在愛人的面前.
愛人即是可能是不信的人的面前.
我們這班基督徒.
既是有主在來的盼望.
那個渴望.
但是是一班非常勤奮工作.
勤奮的來致富的人.
但是又是好樂於隨時去幫人的人.
這樣那個時代就不會有.
因為那個時代是說.
我們上個星期說過.
那個叫做施恩者.
和那個受恩的人的一個制度.
就是令到那個社會穩定.
就是那些貴族.
那些有勢力的人.
就會促養一些人.
來給一些好處他.
可能是食物.
可能是錢.
可能是工作.
但是是要有回報的.
當那個貴族有什麼事情.
那些人就要來去.
可能投票.
可能其他的東西.
這個是當時的制度的交換.
這個叫做類似受養人的一個制度.

$^{1161}$但是保羅說.
神的愛就不是這樣.
就不是我們中國人說的施恩忘報.
是大概的意思.
但都不是一種好像很命運式的.
很空泛的.
而是我們平常都已經是這樣的.
尤其是當時初期教會.
他們每個星期的聚會.
是一個很緊密的群體.
後來都要地下.
因為尼祿都要攻擊教會.
不准有教堂等等.
不准有聚會的地方.
那時候沒有教堂的.
但是這個群體本身就已經賦予了.
要凝聚.
要彼此大家的扶持.
特別是在逼迫的時期.
所以當試圖離開.
有時一定會有一些未必是這麼理解.
甚至已經去到.
五章就真的說那些人.
就是那些不作功的人.
保羅說警告過他一兩次.
就要來趕走他.
就是說有些人看到基督徒.
有這麼大的理想.
就會剝削.
就會濫用.
就會借機會說.
我都是神的兒女.
你要幫我.
但原來他自己背後.
沒有親手作功.
沒有行事端正.
沒有去發展他的照明.
所以原來2000年前.
已經有人生規劃了.
職涯規劃.

$^{1201}$那就很好.
我們讀書的時候畢業.
哪裡有這麼多社工.
哪裡有這麼多講座.
根本就沒有.
你畢業的時候.
你考不到大學.
是你的事了.
沒有人會理你的了.
根本就沒有這麼多學位.
是沒有人告訴你.
現在你想走哪條路.
這樣吧 你幾分啊.
幫你吧.
沒有的 那個時候完全沒有的.
但原來保羅的年代.
已經教導弟兄姊妹.
有這個職涯.
就是職業是一個神聖的照明.
就算你是一個農夫.
你為神耕種.
都是一個神聖的照明.
每一個工作是沒有分貴賤的.
不會貴賤的.
所以原來我們這樣看自己.
又是整個人的存在.
剛才上文講的就是聖潔.
下文講的就是愛.
原來都是和我們這個人.
總之我們還存在這個地球.
我們都是要發展在神眼中的無可責備.
靈魂的救恩的聖潔.
我們就不可以換取.
因為主耶穌為我們十字架的寶血洗乾淨.
我們就立即在神的眼中.
因為聖靈的內著.
我們的生命就已經是成聖清兒.
哥林多後書的七章第二節.
我們已經在生命上成聖了.
但實際的生活.

$^{1241}$實際我們的人生.
神怎樣給我們一生.
我們怎樣去追求聖潔和愛呢?.
這就是保羅在這一章裡.
想教會學習的地方.
我們可以再看多一點.
我們稍為讀了這一段之後.
我們先讀13到18節.
「論到誰了的人.
我們不願意弟兄們不知道.
恐怕你們憂傷.
著那些沒有指望的人」.
又是提到關乎到盼望.
亦都是保羅似乎從提摩太的口中.
或者他過去的模樣裡面.
又或者當時代充斥著一些這樣的思潮.
令到可能教會的弟兄姊妹.
信主的年日比較前.
根基未必很穩.
所以對於人死不能夠復生.
這些觀念他們有多少的影響呢?.
所以保羅在這個可能他內心的等候.
一直去關乎到主再來的時候.
就會想到可能有一些在耶穌還沒回來.
已經死了的親戚朋友.
又是信了主的.
到底他們現在的狀況是怎樣呢?.
以及將來他們又是怎樣呢?.
這個現在和將來.
現在即是已經離開世界.
將來他們會不會再復活呢?.
這些是當時初期教會.
可能受著一些的哲學觀念.
因為當時真的很盛行的哲學.
如果在地中海一帶.
特別是這些這麼富庶的貼索羅尼加地方.
是一個很繁盛的城市.
不是一些內陸很窮乏的地方.
所以希臘的文化和思潮哲學.
一定是很繁盛的.

$^{1281}$很多講學的人到處去講.
所以一定是會影響著很多人的思想.
所以保羅就很詳細地講那個次序.
主的復活,再來和聖徒.
現在還沒死的,已經死了的.
是怎樣的次序的復活.
就講得很詳細.
好像你都沒有什麼空間.
需要再懷疑和需要再擔心.
是一個很大的盼望.
保羅就說你們要用這個盼望.
給予彼此安慰.
即是說可能初期教會是有這些思潮的影響.
保羅就似乎這幾樣東西.
現在你怎樣看待你這個身體.
身體不只是講這個肉體.
是講你的思想,你的靈魂,你的內心.
所以你的意志都是全部受那個影響的.
比如現在我們見到社會上.
每個年代都有不同的明星,歌星.
MIRROR我一個都不認識.
叫得出的可能一個半個.
即是有意識地去包裝.
有意識地令到他們.
每個人見到都會說都瘋狂的.
這樣的追星的效應.
是填補了可能很多人.
內心裡面的渴望,渴求不能夠達到.
就透過看它,追它就可以有一定的滿足感.
但這些其實是一個很虛無縹緲的.
真的不是會實踐到的.
不是每個人都能實踐到的.
但你就會用了很多的時間.
很多的精力,很多的金錢.
但是保羅告訴我們.
原來我們每一個在主耶穌裡面的基督徒.
每一個在神的創造裡面的人.
都有你自己的身體.
都有你自己的人生的存在.
你的生活,你的存留.

$^{1321}$我們有多少歲是實定了的.
我們如果這樣理解我們這個存在.
而我們在聖經裡面這樣守著自己的身體.
這樣就很健康了.
就是說原來曾銳貴你還有三十年命的.
即是你九十歲才死的.
一來我不需要去找那些體相.
二來就是我就應該盡力保持我的身體.
身心靈健康.
盡力我怎樣來繼續與時並進的.
學習一些新的科技與人溝通.
諸如此類等等.
變成你不會停留在一個階段.
你自己進步被淘汰.
你不要賴神.
是怪你自己沒有那個成長.
是你自己窒礙了你自己那個.
繼續用你的身體去服侍神.
所以這裡其實很正面去看.
我們每一個頂子妹.
存在在這個社會在這個地球.
特別提到這個是沒有指望的人的憂傷.
14節我們若信耶穌死而復活了.
那已經在耶穌裡面睡了的人.
神也必將他與耶穌同帶來.
一切的動員的根源都是從神做物主那裡.
所以是神賦予那個日子.
到了耶穌基督再回來了.
帶著天君帶著天使.
帶著已經死去的聖徒等等.
然後就要這樣來到他所有的眾生都要復活.
15節我們現在照主的話告訴你們一件事.
我們者活著還存留到主降臨的人.
可能本來都覺得自己是會等候到主再來.
即是他未死之前.
「斷不能夠在那已經睡了的人之先.
因為主必親自從天降臨.
有呼叫的聲音和天使掌的聲音.
又有神的號吹響.
那在基督裡死了的人必先復活」.

$^{1361}$就是講上文「睡了的人」.
以後就是第二個次序.
以後我們即是頭為尾為那些.
首先復活就是那些死了的聖徒.
「以後我們者活著還存留的人.
必和他們」就是那些死了的聖徒.
復活回來的和耶穌回來一起.
在天上的這些聖徒.
「必和他們一同被提到雲裡.
在空中與主相遇.
這樣我們就要和主永遠同在」.
即是說目前真是這裡有一個階段.
死了的現在在樂園 未在天堂.
所以肯定未在天堂.
因為還有我們這些總有一代人是不會死.
就見主的.
所以現在是一個樂園的狀態.
即是說一個離開世界的基督徒.
他的靈魂現在未有身體.
未有那個復活身體.
他的靈魂應該就不是孤魂野鬼飄來飄去.
走去佛墓 走去嚇人 就不是這樣的.
這些是民間的信仰.
應該聖經的教導就是說.
就現在在主的懷中.
即是亞伯拉罕的懷裡.
一個安睡的地方.
所以是不是靈魂靜止呢?.
因為很熱切地在天上看著我們.
我們亂說的.
因為根本聖經沒有這樣教導的.
就是告訴我們那個狀態.
是等到有一天神的時間到.
主耶穌再回來.
首先就是已睡的聖徒就要復活.
一定要有一個身體.
跟著和主耶穌一起榮耀的降臨.
去到天上.
這裡就守望朋友在做文章.
就說他們是那些十四萬四千.

$^{1401}$來跟主耶穌回來.
他們就是那個Chosen People.
Chosen 即是被選照的唯一的十四萬四千.
在第六章的經文.
其實就不是的.
所以我們就要留心.
是每一個基督徒將來死了的.
現在在主的懷裡是安息的了.
我的理解.
遲些再揭曉不同的經文.
因為我們遲些會說末世論的神學.
那個系統神學.
真是有一天.
突然間主的聲音有號筒吹響.
有天使長的好像召集靈.
來叫醒了那些聖徒.
古往今來有多少人.
不同的民族一起復活.
有新的身體.
復活的永恆的身體.
來跟主耶穌回來.
幾個浩兼浩蕩的行列.
這樣來到再來.
就變成我們不是一種幻想.
或者是一種幻覺.
而是我們要有這個.
是保羅自己說的.
當用這些話彼此勸慰.
保羅說你對號入座就行了.
你真是這樣彼此勸慰.
我們今天基督徒.
似乎很少人會說.
主耶穌回來.
那個興奮.
那種滿足.
所有死了的人.
無論他怎樣死.
有些人病死.
是夭折.
有些人是被逼白死.

$^{1441}$各種都有的.
所有都要復活回來.
有些人是很榮耀.
很有錢.
有些人是很坎坷.
所有都復活回來.
大家是完全平等的.
這樣來到我們會見到他.
所以保羅說要彼此的勸慰.
真是少.
基督徒今天茶餘飯後.
在聚會.
離開聚會.
談論主的在來.
自己在甚麼位置.
自己怎樣將來.
享受到那種稱讚.
神在天使眾人面前.
對基督徒的稱讚.
這些是使徒他們的追求.
是保羅他們自己渴慕的.
將來保羅說.
東邑的官面為他存留.
還得到主面前的讚賞.
所以為何他們會這樣勞苦一生.
正正是一個很榮耀的盼望.
所以保羅告訴我們.
將來我們一定會有一個.
所謂的空中的相遇.
就是到底在哪一個天空.
在多高呢.
二萬呎還是三萬呎呢.
你想想文學機都有規定.
去的和返的都有距離.
剛開始我們小時候就想.
天上沒有街燈怎麼辦呢.
原來不是.
長大之後就發現原來是定了的.
怎樣飛原來是有規定的.
相隔多少分鐘原來是有秩序的.

$^{1481}$就不會撞機了.
將來我們就不需要擔心這些細節的問題.
將來會有一個大團聚.
所以姜濤終極的盼望就是主的再來.
就是有這個團聚.
所以這張聖經談論到那個聖結.
和彼此相愛.
和夾雜的都是主在來.
甚至掃除那個擔憂.
原來是很重要的.
姜濤如果還有擔憂.
還要對死了的人擔憂.
似乎又真是過份了.
因為一來已經沒有辦法再改變了.
但是原來這個沒有辦法改變.
我們就要盡力在有生之年為的祈禱.
帶他信主.
原來那個不需要改變.
就是我們不需要擔心他將來不會再出現.
是不會出現.
這樣的不需要擔憂.
所以這張聖經表面上很淺白.
但是背後補充了一個很正確的神學.
就是神拯救我們每一個人.
原來都已經有一個計劃.
這個計劃包括我們的身體是怎樣呢?.
都是在神的眼中有它的設計.
它是最奇妙的智慧的創造主.
所以我們每一個人.
我們的身高肥瘦.
都是在它的定子先見裡面.
當然絕對是有空間.
我們要盡量保持我們今天能夠掌握的.
那個資訊的BMI的標準.
過分的,過度的.
一定不是神的意思.
如果是因為暴飲暴食而帶來的不適.
肯定是需要在有生之年去求神幫助.
否則就很難去說得通.
神創造我 為甚麼我會這樣.

$^{1521}$會不會是你每晚吃的都是垃圾餐.
麥當勞 喝汽水.
每餐都要喝汽水.
你可以想像得到.
脂肪的糖尿一定會增加.
你不可以怪神.
這是我們真的要守住自己的身體.
不過特別強調.
當時社會思潮扭曲的偶像敬拜.
一定會帶來拜偶像和肉體得到的滿足.
似乎就是性的滿足.
那個身體的所謂高潮的滿足.
保羅說不應該這樣去看.
我們是一種為了慾望的動物.
應該是神給我們這個身體.
我們應該怎樣去負面的消極的就是逃避.
但更好的就是要守住.
守住即是我們應該發揮到極致.
神怎樣在我們的人生去用我們這一生.
這就是守住自己的身體.
不是自戀狂.
越來越對自己有感恩.
神多謝你讓我生在這個家庭.
我的成長,我的經歷是獨一無二的.
第二個人不可以翻版的.
我的存在總是可以為神的國度盡一分力.
基督徒一定要這樣去想.
你這樣想.
你就會開始想你的人生,你的照明.
你就會積極面對不同的處境.
你就會將你可以幫到的人帶到他去順主.
所以保羅在這裡特別強調.
真的不要被那些拜偶像的人.
或者那些欺凌文化的思潮影響.
似乎他們一聽就明白.
他們真的不是做這些事情.
貼守羅尼加城的人.
似乎保羅這裡是一個客觀的告誡.
不是已經有人從事這樣去廟裡拜偶像.
跟著和那些妓女來行刑.

$^{1561}$即是一些公開性的.
似乎他們不是這樣.
哥林多就是這樣.
但是這裡告訴我們.
強調的就是守住自己的身體.
成為聖潔.
在神的創造裡面.
去用自己整個的生命.
所以似乎這裡告訴我們.
身體的節制也是很重要.
節制包括你的飲食.
我們怎樣去用我們的時間.
節制可以牽涉到整個人.
怎樣能夠哪樣東西是最優先.
我就不斷做最優先的事.
這樣的人生.
不要永遠都去到限期.
重要到不得了的事.
你都要去到限期才做到的.
其實可能那個時間分配就不行了.
即是說你不是守住自己那個人.
令到整個人是服你.
很多東西是可以學習的.
很多東西是可以彼此勉勵教導的.
譬如我們有時牧者禮拜.
有時忙到都沒空吃飯.
你又會不覺得餓.
你又會感覺到.
你看到頂智媒回來.
大家很開心地分享聚會.
你會有一種滿足.
你好像忘記了你自己暫時那個需要.
有時我們都需要這樣的.
所以原來我們在平日.
我們多去留意我們自己的身體.
留意我們內心的情緒.
留意我們內心為主而活的動機.
那個意志.
這樣來到去.
有沒有東西已經拉我們離開了.

$^{1601}$那個熱心事奉的心情.
那個情緒等等.
這些就是守住自己的身體.
可以讓我們能夠一直在神的旨意裡面.
我們有一個熱切的參與.
因為有時慢慢地.
順轉來那個動力沒有了.
那種越來越慢下來的那種.
連那個mood都沒有了.
變成一個好像聚會都是例行公事.
甚至去到一天.
連這個例行公事都不需要人家怎樣看你了.
那個人就更加不回教會了.
所以其實不知不覺冷卻的過程.
其實就是我們一直守住自己的身體.
守住我們自己對於永恆的價值觀.
所以這裡很明顯就是說.
他們要逃避那些拜偶像.
而引申的淫亂.
必須是要很消極地一定要逃避.
但積極地就要在神的恩典裡面.
來留意自己整個的存在.
所以我當然經常都說笑.
姜濤有時自己做一些違反常理的事.
我們牧者就一定要去.
譬如我們經常去殯儀館.
我們去都不只是做事.
有時我們都會留意下那些人.
我自己要下葬 去墳場.
去骨灰位 上位.
我有時會早去.
即不是甚麼秘密.
我早去一方面就自己心靈預備.
其實我望下周圍的那些骨灰位.
你會看到有些其實有些故事的.
有些很年輕的 有些很老的.
多數都是老的.
你會看到原來人生是很無常的.
我都講過一次在鑽石山的骨灰.
是我的表哥.

$^{1641}$他的人生都是很坎坷的.
我們都幫他撿骨灰 灑灰.
我幫他灑灰 在花園.
安上位.
跟著我和我弟弟去幫他上位.
我弟弟就差一點點.
叫我過去望一望隔壁附近等等.
嘩 不看由自可.
一個家庭 整個家庭.
爸爸媽媽 還有幾個子女.
看到那些姓氏.
為甚麼幾歲又有 十歲又有.
我弟弟說是不是滅門慘案.
可能是的 沒理由突然同日同日死的.
是全個家庭.
你會看到很多故事.
所以有時了解人生.
你就能夠守住自己基督徒的身體.
那個人生觀.
所以不只是很負面.
甚麼都不准做.
而是我們大家都活一生.
有些人只是為了很少的事.
他叫生叫死.
但基督徒我們應該再進化一點.
我們應該看到神在我們生命裡面有一個召命.
有一個旨意.
我們應該為主而活.
我們應該用聖潔去看待我們自己的身體.
所以表面上我們不要放縱.
但更重要的是我怎樣準備我自己.
我那時好像有神的感動.
我一信耶穌第二年就已經感覺到有那個呼召.
在教會不知不覺就學傳徒人穿衣服.
自己一個年輕人.
我那時在百貨公司.
大陸的百貨公司.
買那些大地牌.
穿白恤衫 藍扯褲.
希望那時開始學一些簡樸生活.

$^{1681}$穿到像大陸仔那樣.
那時是很古怪的.
但在教會裡面有些弟兄很資深的.
突然間問我.
你是否想做傳道?.
我說你怎知道的?.
我內心有一個感動和呼召.
我看這個樣子就覺得你是了.
類似這樣.
那時是沒有人教我的.
我自己有這種好像模仿.
但原來那時是很粗糙的.
原來不知不覺是裝備自己.
想走這條人生的路.
一個這樣的以聖潔.
當然是傳道的照明.
去裝備自己.
去守住自己的身體.
至少我有這樣的信念.
在當時我想這樣為主而活.
當然可能有些偏激.
但慢慢慢慢你去看.
對聖經越來越認識多了.
你會知道原來很多事情.
原來人生一來已經過了這麼多年.
原來很多事情是命中有神的安排帶領的.
你在神的引導下.
去活你自己神給你的光位.
因此我想是一個最大的滿足.
所以不要和其他人比較.
一定要和自己過去比較.
有沒有更加愛主.
有沒有更加很想得到神的那個獎.
即是那個讚賞.
不單是問心無愧.
是更加很想得到神的那個讚賞.
接著就是彼此相愛要有很平衡.
保羅的含義是每個人都應該愛不同的人.
但這個愛裡面不要被人剝削.
你自己要先做好自己的那個本份.

$^{1721}$就是你要立志的做一個安靜的人.
立志的親手的作功.
要在外人的面前要行事端正.
就是說講求你每天八個小時的時間.
一般性的說.
當時更加是現在就不是了.
有倫班現在根本就不會有夜晚的觀念.
家家戶戶都有燈光.
古代是沒有的.
夜晚就是夜晚.
是甚麼工作都要停頓的.
所以那個時候他們看的工作觀念.
其實200年前都還是.
根本不會有街燈.
不會有燈 不會有電燈.
你就可以想像到他們那個照明.
一定是你每天怎樣殷勤的發展神給你的恩賜.
無論是工作,教會等等.
所以彼此相愛原來是這樣的一個層次的累積上去.
這樣來彼此相愛.
所以偶然間出現一些剝削的人.
五章就說了.
這個人 保羅就說不可以去容許他.
第十四節 五章的十四節.
告訴我們.
「我們又勸弟兄們要警戒不守規矩的人.
勉勵那些灰心的人.
輔助那些軟弱的人.
也要向眾人忍耐」等等.
後書就告訴我們.
我再站一站在後面.
後書的三章的第六節.
「弟兄們 我們奉主耶穌基督的名.
吩咐你們.
凡有弟兄不按規矩而行」.
一直說那個規矩 就是那個order.
就是對自己的人生沒有去守住.
不遵守從我們所受的教訓.
就當遠離他.
你們自己原知道應當怎樣效法我們.

$^{1761}$因為我們在你們中間.
未嘗不按規矩而行.
不斷說規矩.
也未嘗白吃人的法.
賭時辛苦勞碌.
晝夜作工.
免得叫你們一個人受累.
這並不是因為我們沒有權柄.
說的是使徒的權柄.
乃是要給你們作榜樣.
叫你們效法我們.
所以這裡說.
「我們在你們那裡的時候.
曾經吩咐你們說.
如果有人不肯作工.
就不可吃飯.
因我們聽說.
在你們中間.
有人不按規矩而行.
什麼工都不作.
反倒專管閒事.
我們靠主耶穌基督吩咐.
勸誡這樣的人.
要安靜地作工.
吃自己的飯」.
就是這個意思.
就是從剛才我們看的四章的意思.
不斷說的是親手作工.
按規矩而行的意思.
所以原來彼此相愛.
基督教的彼此相愛.
不是一種剝削.
是一種我們在神給我們的祝福裡面.
去分享給其他人.
使其他人透過我們的幫助.
而得到神的愛.
更加積極.
更加成長.
我們需要有這種智慧.
所以首先我們自己.

$^{1801}$我們自己就是一個親手作工的人.
來服侍其他人.
我們能夠有神給我們的豐富.
即有餘.
不是缺乏.
我們成為別人要幫助的對象.
這是一個打破了當時的社會制度.
一種很牢固.
其實就是強弱.
其實是一種社會的貧富.
稍微有一個叫做.
不要令到問題惡化的一個暴亂的.
一種的社會制度.
這些受助人制度.
但是是一種交易.
是一種沒有愛心底下的交易.
所以保羅建立這個初期教會.
大家成為神兒女.
這個概念是很新的.
大家都是活在一位神.
一個救恩底下.
大家都是只有耶穌.
和將來大家都去到永恆.
但是在地上的短短的幾十年.
又是怎樣來彼此相待.
彼此來愛呢?.
開始不同的使徒.
就會有原則性的教導.
所以這裡特別強調.
我們要看待我們自己的人生.
是一個神的召命的人生.
所以要幫助弟兄姊妹信了主之後.
要發展他們自己的成長.
和神給我們的恩賜和我們的召命.
和要看每一個職位的工作.
都是神的祝福和恩賜.
所以當然有些倫理的問題.
譬如馬會投注站可不可以做呢?.
是不是神的召命恩賜呢?.
這些倫理問題其實不容易.

$^{1841}$一眼看三言兩語就解釋得到.
譬如我們都聽過那位牧師.
說澳門的弟兄姊妹.
很大試探.
你做賭場.
你的工資是低幾倍的.
那怎麼辦呢?.
賭場就最容易賺錢了.
是一個很大的倫理的困難.
等於保羅說在哥林多教會.
可不可以吃「濟澳藏之物」呢?.
保羅說凡是市場上買的肉.
都是比偶像更大.
那怎麼辦呢?.
保羅說除非我們離開世界.
原來保羅也會看.
整個社會現象的問題.
原來是看你到底.
涉入多深的問題.
這些又是很大的倫理課題原則.
成為我們的參考和看法.
保羅說很快的.
給大家稍為有一點這方面的知識.
哥林多前書的八章.
八章這一段.
「論到濟澳藏之物.
我們曉得我們都有知識.
但知識是叫人自高自大.
唯有愛心能夠造就人.
如果有人以為自己知道什麼.
按他所當知道的.
他仍是不知道」.
先講了原則.
原來當時有些人自吹自擂.
覺得自己甚麼都不怕.
百毒不侵.
很夠信心.
去踩他們的牆都不怕.
他就說了.
其實那個課題不只是你個人那麼簡單.

$^{1881}$是你的見證和影響性的問題.
第三節.
如果有人愛神.
乃是神所知道的.
接著說.
「論到吃濟偶像之物.
我們知道偶像在世上算什麼.
也知道神只有一位.
就是我們的眾物主.
再沒有別的神.
須有稱為神的」.
即是人造出來的那些.
「在天或在地.
就如那許多神許多主.
然而我們只有一位神.
就是乎萬物都本於他」.
就是神的主.
「我們也歸於他.
並有一位主.
就是耶穌基督.
萬物都是藉著他而有的.
我們也是藉著他而有的」.
我們被創造.
所以用聖結去守住自己的身體.
都是一個很重要的神學.
接著第七節.
「但人不都有這等知識.
有人到如今因拜慣了偶像.
就以為所吃的是濟偶像之物.
他們的良心既然軟弱.
也就污穢了.
其實食物不能叫神看中我們.
因為我們不吃都沒有損害.
吃了也不會有益.
只是你們要謹慎.
恐怕你們這個自由.
這個自由就是說.
我什麼都不怕.
我是沒事的.
我的信心很堅固的.

$^{1921}$我去那些偶像廟裡吃.
是完全不怕的.
他們說明拜鬼我也不怕的.
是不會影響到我.
我睡得很好.
阿普羅說不是你.
是你這個見證影響了那些軟弱的弟兄.
這裡說.
「只是你們要謹慎.
恐怕你們這個自由.
竟成了那軟弱人的絆腳石.
如果有人見你這個有知識的.
可能是那些長老.
或者那些覺得自己很厲害的人.
我吃給你看.
我都沒事的.
這樣似乎就輕忽了一件事.
就是這人的良心如果軟弱.
豈不放膽吃那濟偶像之蜜蜂.
它又跟你一起.
然後它吃完之後.
它自己又有良心不安.
所以第十一節.
因此基督為他死的那個軟弱弟兄.
也就因你的知識沉淪了.
因為你這個行徑.
你這個所謂全心去挑戰別人偶像.
這個信念.
就令到那個軟弱弟兄沉淪.
沉淪即是跌倒.
不是真的落火湖的那個沉淪.
你們這樣得罪弟兄.
傷了他們軟弱的良心.
就是得罪了基督.
所以食物如果叫我弟兄跌倒.
我就永遠不吃肉.
好像很誇張那樣.
即是說保羅是很自由的.
他但是會很遷就到軟弱弟兄的良心.
所以他一定不會去那些廟裡來吃.

$^{1961}$免得叫我弟兄跌倒了.
你可以想像到.
保羅他這樣去運用他自己的權柄.
我們很快就讀第十一章.
有少許教導的.
讓我看看.
應該是第十一.
等一下.
我有沒有看漏了少許呢?.
我暫時找不回那段經文.
它提到凡在市場上賣的肉.
就只管吃.
是良心的緣故令他不吃.
但如果有人.
應該是聖餐那段經文.
在聖餐那段經文有這樣提過.
我暫時這裡我看不見.
我稍後補回.
我今晚回去補回給大家.
大概意思就是這樣提到.
我們要涉及到自由的問題.
不要濫用自己的自由.
否則就會帶來讓其他人受到虧損.
所以今天總結了第四章.
其實是很寶貴的一件事.
就是基督徒我們應該要在我們的人生.
越早信主.
我們應該越早教導.
我們要有正確的觀念在自己的身體.
不是負面這麼簡單.
反而是積極的.
在自己整個人生.
包括我們的原生家庭.
對我們的影響.
和我們的裝備.
如果我們的弟兄姊妹.
我們的下一代.
是大家都裝備.
還是要為主而活呢?.
其實根本不需要擔心.

$^{2001}$無論香港也好.
英國也好.
無論什麼政治.
什麼前途.
什麼都好.
不需要擔心.
因為每一個基督徒.
他都知道自己要為主而活.
他要發展他的照明.
他要在他的人生去榮耀神.
不會再為自己的身體不漂亮.
去唉聲嘆氣.
為自己和別人比較.
我的工作又低下一點.
收入沒那麼高.
那些全部是社會問題.
是很現實的問題.
但是我相信保羅的意思.
就是告訴我們.
我們信了主之後.
那個真正的生命改變.
就是我們的靈魂救恩.
是應該影響著我們的生活上.
雖然我們不會突然有一個樣子.
讓人看到我們是屬於基督的人.
但是我們慢慢的那個價值觀.
是從裡面的生命改變.
而我們有一套正確的神學.
無論是身體的神學.
職場的倫理神學.
和我們的人生.
譬如我們那個什麼為之標準的健康.
在神的面前的.
我想我們會少了很多很多的煩惱.
就會用很多的時間.
可以怎樣在神的國度裡面.
有份參與.
因為如果不是我們只是為健康而健康.
譬如美國人.
他們現在很多都不回教會.

$^{2041}$有些人家裡一天吃多少種維他命丸.
多到一個地步好像都很難想像.
其實他已經是很健康的.
但是還要不斷地吃.
那就變成有時候有一種過渡了.
可能就是出於一種很大的恐懼.
覺得自己的身體要和其他人做比較等等.
有些譬如美國南部那些黑人區.
那些癡肥的嚴重性是你難以想像的.
為什麼這麼不平均呢?.
為什麼白人有錢的就會不斷地吃那麼多.
那些保健的丸.
甚至很多又要看輔導各方面的東西.
後來慢慢就要尋獲毒品等等.
為什麼呢?.
是不是源於根本沒有一個正確的神學觀念.
去看自己的存活.
整個用聖結看待自己的身體.
所以這個告訴我們.
我們需要去留心.
彼此相愛也是.
我們需要在神的教導下.
我們要建立這份由心而來的愛神愛人.
但是個人的成長.
個人的用行事端正.
和親手作供.
彼此相愛也是不可或缺的根基.
晚了一點我們一起祈禱.
再一次多謝主.
讓我們今天可以來看第四章的《帖撒羅尼迦前書》.
我們懇求主讓我們不要受社會思潮的嚴重侵蝕.
反而我們能夠像保羅的教導.
用聖結,用神的觀念,神學去看待自己.
去看自己的存在.
我們能夠有一個很正面的價值觀.
這樣去塑造我們的人生.
在神的家裡我們彼此的服侍.
所以聖結和彼此相愛都是息息相關.
我們個人正確地看待自己.
必然會很殷勤地親手作供.

$^{2081}$在愛人面前我們會是行事端正的人.
幫助我們能夠常常活在主的聖靈的引導帶領下.
過一個基督徒健康的人生.
願主幫助保守這篇信息.
成為很多在困難裡面弟兄姊妹的一個祝福.
提醒和安慰我們禱告奉耶穌基督的名.
阿們.
\newpage



\section{創世記 3:20}
\label{sec:BHhp4UIkEH0}
\textbf{2024-05-11 晚堂《一個蒙恩的母親》創世記 3\_20  講員:楊穎兒姑娘}
\newline
\newline
連結: \href{https://youtube.com/watch?v=BHhp4UIkEH0}{\texttt{https://youtube.com/watch?v=BHhp4UIkEH0}} ~~~~ 語音日期: 2024-05-14
\newline
\newline
\hyperref[sec:VvxSAXf_Lww]{\small{< < < PREV SERMON < < <}}
~
\hyperref[sec:index_chronic]{\small{[返順時目]}}
~
\hyperref[sec:index_scriptual]{\small{[返順卷目]}}
~
\hyperref[sec:rgd8tUJeioY]{\small{> > > NEXT SERMON > > >}}
\newline
\newline
創世記 3:20
\newline
\begin{longtable}{cl}
\hline
\hline
章節 & 經文 (和合本修訂版)\\
\hline
3:20 & \begin{tabularx}{0.7\textwidth}{X} 那人給他妻子起名叫夏娃，因為她是眾生之母。 \end{tabularx} \\ \\
[1ex]
\hline
\hline
\end{longtable}
$^{1}$大兄姐妹平安.
很開心.
這次可能是我第一次.
在晚堂和大家分享主的話語.
當我們去準備母親節之際.
我們都透過夏娃.
夏娃的生命故事.
去明白上帝創造的偉大.
可能你會在想.
為什麼說夏娃.
母親節.
夏娃是被蛇用詭計去欺騙.
以致犯罪.
以致人被逐出伊甸園.
為什麼我們還有什麼可以向夏娃學習.
所以可能這是我們對夏娃的印象.
就是她犯下的錯.
今天都想和她平反.
正如我們對著我們的母親.
我們的婆婆 岳母.
我們的家人.
可能都會墮入同一個陷阱.
可能只記得某一方面的事件.
所以我們看看.
我們可以在夏娃身上學到什麼.
而我作為一個.
我的小朋友都步入青年.
我都開始明白到.
其實我在他們的生命中.
已經沒有那麼重要.
我都在習慣.
所以可能你們都會有這樣的經歷.
可能你們都去到一個階段.
子女在遠方.
或者母親都在遠方.
或者離開了.
但是我們都會因著母親對我們的愛.
在當中的成長.
都有些影響.
所以我今天想在經文中.

$^{41}$和大家一起去發掘一些亮光.
所以第一.
我會為夏娃去平反.
因為我們有時很快就作出一個判斷.
覺得她沒有一個角色.
第二.
我們嘗試從阿當身上去學習.
怎樣去看待夏娃.
第三.
我們亦會回到上帝的恩典裡.
去看待夏娃所蒙的恩典.
今天我們讀的經文不只是一節.
我們會看很多節的經文.
讓我們一起祈禱.
請許天父.
今天晚上我們一起進入夏娃的生命故事時.
我們會發現當中很多屬天的智慧.
求你打開我們的心扉.
教導我們當中.
你要我們明白的道理.
亦都更加去體會到你給人類的恩典.
以致我們更加可以看重母親這個角色.
多一份愛.
多一份尊重.
我們這樣的祈禱.
奉主明求.
阿們.
第一點.
我們說說.
大家對夏娃都很熟悉.
所以我會慢慢解釋一些經文.
第一點.
夏娃是蒙恩成為第一位的女性.
當我們回到創世的.
即是墮落之前的世界.
我們會發現.
夏娃是那時候唯一一個女性.
在一個無罪的世界裡.
她真是只有她一個.
那時候上帝給了她甚麼心意.

$^{81}$要她怎樣去生活.
如果我們細心去讀一下創世記.
可能我們已經譯成中文.
未有覺得她有這樣的思義.
但其實她在原文寫成的時候.
是一篇很優雅的經文.
但翻譯的時候可能失去了一些味道.
但我們仍然會看到.
神是逐一步一步將生命注入這個世界.
使混亂的世界去恢復上帝想要的平衡和和諧.
祂首先創造了一個伊甸園.
即是世界.
然後又創造了伊甸園.
然後將祂心愛的人類.
放在這個很美麗的花園裡.
我們看看經文.
我們看看第一章26節.
我們看到上帝是用話語去講出這個世界.
去創造這個世界.
是由講開始的.
神說我們要帶著我們的形象.
按照我們做人的樣式.
使他們管理海裡的魚.
天空的鳥.
地上的生畜和全地.
以至地上爬的所有一切爬行的動物.
亞當當時是代表著人類.
是按照著上帝的樣式去做.
我們去講上帝的樣式的時候.
就不是說祂這樣的樣子.
而是說祂的設計裡面.
是穩含著上帝的良善.
神的樣式.
祂的形象.
在所有萬物之中.
都是有著上帝的形象.
人類是賦予可以管理大地的權柄.
我們繼續看下去.
第27節.
上帝是會用行動.

$^{121}$祂不是只會說.
會用行動去創造這個世界.
神就按照著祂的形象創造人.
照著神的形象創造他們.
祂創造了他們.
有男有女.
其實這一節的經文.
其實在我們翻譯到希伯來文的時候.
其實是以一支詩詞的手法去寫.
所以你會發現.
其實三句所說的東西.
都是差不多的.
都是有神.
有祂做人.
這三個組合是不情有的.
所以我們這樣去看.
這個寫作手法的時候.
其實它是象徵著.
上帝的設計裡面.
是跟人很親密的.
那個聯合是很親切的.
這個是上帝的心意來的.
所以在經文裡面反映了.
他們每句話都有祂們.
大家會看到那種緊密性.
所以我們看到.
既然神是這麼想跟人緊密的時候.
祂在創造人的時候.
也有這個心意.
祂其實是希望男和女.
都能夠擁有這麼親密的禮物.
這個親密的關係.
其實是夫婦可以經驗到的.
而不是某一方是特別經驗多一點.
或者少一點.
但是在上帝的設計裡面.
其實兩者是不羞愧的.
而且在當中是享受上帝的同在.
甚至我們的聖經也有一本書.
《雅歌》其實也是在說.

$^{161}$夫妻之間的親密.
而它不是一種禁書.
而是在說那個高峰.
就是上帝和人的愛的高峰.
我知道在我們的文化裡面.
我們可能很避忌.
我們未必會這樣跟大家說.
尤其是在講壇上.
但是我想說的就是.
這種親密的關係.
其實是存在我們的人裡面.
在夫妻關係好和壞.
其實都是有一個正比的.
儘管我們在我們的生活裡面.
都有各種限制.
或者還未能夠做到這種親密的關係.
但是也不能否定.
神創造的心意.
到了28節的時候.
神也會賜福.
祂會說 然後祂會做.
然後祂也會賜福.
神賜福給他們.
神就跟他們說.
要生養眾多 遍滿全地 治理.
還要管理海裡的魚 天空的鳥.
地上的各種活動和生物.
我們看到.
其實這種親密 神聖的禮物.
其實是神要我們去承受.
也會賜福使萬物繁殖後代.
而唯獨是神挑選人.
跟神有一個這麼親密的關係.
所以我很鼓勵大家.
在夫妻關係裡面.
去珍惜這個關係.
真的要抽出時間去約會.
去說一個時間.
何時有一個親密的時間.
何時一起跟上帝有一個親密的時間.

$^{201}$而當中可能大家沒有配偶.
可能繼續跟神有一個美好的相交.
仍然可以在神的華語裡面.
得到這種親密的關係.
而夏娃也不是第一個.
夏娃也是唯一一個.
跟上帝和阿東.
那時候在無罪的世界裡面.
是這麼親密.
而且她也是唯一一個.
不是由女性而生的.
她是由男性的身體而出的.
這是一個大圖畫.
我們說了上帝的創造人的畫面.
這是第一張.
當我們回到第二張的時候.
我們就有一個更加近距離的圖畫.
去看看上帝為何要用這種方式去創造萬物.
如果我們看看《創世記》第二章十八節.
提到上帝為何要在阿東的肋骨裡面.
去找夏娃.
因為上帝有原因的.
他見到那個人單獨一個不好.
就要做一個配偶給他.
其實這件事之後.
下一章就提到神給他工作.
他要做他的工作.
還要為那些動物起名.
阿東在做的時候.
發現他沒有找到一個適合他的配偶.
這個經文我們留意到甚麼呢.
其實在阿東還沒有意識到.
他需要一個伴侶去協助他.
去治理這個大地的時候.
其實上帝已經有這個想法.
他已經想創造一個女人給阿東.
所以還未結婚的弟兄姊妹.
我們可以在這段經文學習.
就是神其實知道我們的配偶是誰.
但很多時候他需要我們去明白.

$^{241}$我們為何需要有配偶呢.
我們人生的目標是甚麼呢.
如何為神工作呢.
然後這個配偶才能大到在你的身邊.
當夏娃被安置在這個原址裡.
神給她有一個目標去做.
他們就是要去工作.
去看守這個原址.
所以其實我在教導小朋友的時候.
我都會這樣教的.
你不要追求美貌.
不要追求財富.
即是你們找配偶的時候.
你要找一個能夠跟你一起去完成上帝工作的配偶.
因為這個就是上帝設計裡面一個很重要的部份.
我們會見到阿東和夏娃當時是一個很平等的夥伴.
關係在地上生活.
阿東可能是第一個出現.
然後是夏娃.
但我們看到他們都是要一起去運作這個世界.
所以我們看到男人和女人其實是平等地去分擔一個工作.
就好像我們在家裡.
父母他們都有一個很平等的角色.
很責任和權柄.
其實是成為一個伴侶一起去追隨神的話語.
當我們懂得這樣去看的時候.
其實是很美好的.
我們會看到在一個親密的關係裡面.
沒有人是為了另外一個人的快樂而被做的.
但我們是共同參與在上帝神聖的計劃裡面.
這個就是夏娃在墮落之前的生活.
其實她的故事也是很值得我們慢慢去學習的.
蒙恩成為眾生之母.
我想我們最清楚夏娃就是她的罪.
我講過了.
但壞事是否定斷了她就是這麼壞呢.
我們去深入經文的時候.
我們會發現在女人被做之前.
在《創世記》第二章提到.
就是這樣的.

$^{281}$耶和華在東方的伊甸園栽了一個原子.
放了人安放在那裡.
耶和華使各樣的樹從土裡面長出來.
可以越人的眼目.
好作食物.
原子當中也有生命樹和知善惡的樹.
原子裡面是有兩棵樹的.
到第十五章.
耶和華神就把人放在伊甸園裡面去監魂看管.
耶和華吩咐那個人.
原中的各樣所出的.
你都可以隨意吃.
只是知善惡樹所出的.
你不可以吃.
因為你吃了的日子必定死.
耶和華神很清楚.
講了哪一棵樹不可以吃.
大家留意.
在當上帝對男人講這句話的時候.
女人是不在的.
到第三章的時候.
試探來了.
蛇去誘惑這個女人.
這個女人我相信她所認識的.
就是從亞當爾來的.
或者是她自己再領受.
但是我們不知道.
我們只知道她怎樣回答.
原中的果子我們都可以吃.
只是原中那棵果子.
哪一棵.
剛才她講了兩棵.
原中那棵果子.
神說不可以吃.
又不可以碰.
免得你們死.
很明顯我們看到.
我們看到什麼.
她是講不出哪一棵樹.
還有她加了一些東西.

$^{321}$不可以摸.
你有沒有玩過一個遊戲.
就是你排隊.
第一個人就得到一個訊息.
就一直傳傳傳到最後那個.
通常最後那個人講出來的東西.
都變成另一回事.
大家有沒有玩過這些遊戲.
這個就好像這樣.
究竟為什麼當中有些誤會.
上帝和亞當說有兩棵樹.
一棵叫生命樹.
一棵叫分辨善惡樹.
但是考一考的時候.
答不中.
而且加了一些東西.
當我們在生活裡面.
我們去分辨一些事情的時候.
我們去求問神的心意的時候.
我們會祈禱.
我們會等待.
但是通常我們等到.
想到辨別出來的.
又好像不太肯定.
這句經文可能給了我們一個答案.
夏娃究竟在哪裡得到神的訊息.
會不會是亞當告訴她.
亞當告訴她.
為什麼中間會不一樣.
會不會是亞當接收得不清楚.
我們在分辨的時候.
我們會不會加了很多.
人自己添加的規則.
所以上帝在做男女夫妻.
是有原因的.
因為要大家好好地互補.
亞當從上帝裡面.
知道有什麼可以吃.
有什麼不可以吃.
但是女人在面對誘惑的時候.

$^{361}$亞當是站在旁邊.
因為一起吃.
站在旁邊又不出聲.
但是上帝責怪的時候.
就說你給這個女人我的.
是她給我吃的.
當中我們不知道發生了什麼事.
但是我知道的就是.
神要夫妻一起.
是要大家互相支持.
互相幫助.
不要陷入一個陷阱裡.
我們究竟有沒有很認真去看待這個角色呢.
在我們的關係裡面.
我們可能都有陷入一個想責備的誘惑.
這個都是出於人的天性.
我們很容易就.
因為說你當然容易.
責怪你嘛.
我們責怪你就不關我的事.
但是我們有沒有.
我們上次去責怪伴侶.
或者我們的家人是什麼時候呢.
可能家裡有些小事.
小小的可能是一個垃圾.
或者是那道菜少放了一些鹽.
我們都會陷入一些責怪.
你經常都是這樣.
這些話語.
但是我們會看到.
很有趣的.
神.
當亞當 夏娃犯罪之後.
神叫第一個出來是叫誰.
叫夏娃.
不是啊.
是叫亞當出來的.
亞當出來.
但是亞當當時是沒有去承擔這個責任.
而當上帝去做這個裁決的時候.

$^{401}$他有一個判決.
他說女人.
你要生育是痛苦的.
你要丈夫是去管理你的.
而亞當也是這樣說.
你就要.
因為你吃了你妻子給你的果實.
你就要受咒詛.
終身勞苦.
你要歸於塵土.
亞當當他看到自己.
原來那個結局是要流汗.
去工作到死.
變成塵土的時候.
但是他看著上帝給女人的判決.
上帝縱然給了她痛苦.
但是他也承擔著賦予生命的責任.
如果當亞當沒有女人的話.
亞當根本就不會完成到.
我們剛才說的一章二十六節.
上帝是要繁衍後代這個目標.
當時這個女人是沒有名字的.
但是亞當現在幫她起名字了.
知道怎樣起了.
他命她為夏娃.
意思是生命.
亞當甚至去解釋.
她是眾生之母.
這個時候我們見到.
亞當變得更加像神.
他看到神的奇妙.
他看到神的心意.
儘管是有咒詛.
但是背後仍然是有祝福.
如果大家要去想想.
上一次和你愛的人.
或者你的家人去吵架的時候.
你有沒有看到上帝在他身上的作為.
上帝擺他在你生命中的美意.
還是看到他就憎恨他死呢.

$^{441}$我們能不能夠接受.
我們家人 我們伴侶的好和壞.
然後在背後看到神的心意呢.
當亞當真的看到夏娃.
神在她身上的心意的善的時候.
她就是一個蒙福的女人.
要明白神對一個家庭的使命的肯定.
其實是需要夫妻一起去領受.
一起去溝通.
在過程中會有磨合 會有不對.
但是仍然是向著同一個目標去做.
我們有沒有做到呢.
在我的肯定來說.
我自己的呼召是做宣教的.
現在也是很感恩.
也是錦繡堂 差遣.
在本地做非洲宣教的事工.
非洲人宣教的事工.
我的mission 其實我的calling.
其實也是和我先生一起去做這個mission.
而這麼多年來也是在邁向這個目標.
很多時候可能是小時候的一些訓練.
都是衝得很厲害的.
有時候在服侍的裡面.
我就衝得忘記其他東西.
在生活中其實是有一個很大的不平衡.
第一可能少了時間去陪先生和子女.
自己的治理 健康也是忽略了.
所以這幾年來.
這麼多年來 不是幾年.
先生其實也一直幫我去恢復這個平衡.
很感恩現在我們可以一起去做運動.
因為之前我曾經有一個背脊痛.
痛了十年.
但經過做運動 做了半年又康復了.
所以我們現在也會一起做運動.
我們其實在加州跑的.
可能你們如果住加州也會看到.
還有自從我加入了堂去做傳道.
一個part time的傳道的時候.

$^{481}$我也需要當值.
我先生是很聰明的.
他五點半回家.
他就會帶著我小兒子 兩歲的.
就來到監堂.
在那裡玩玩.
其實意思就是叫我要回家.
因為如果他不來.
可能我就一直做下去.
我兒子來到教會.
他兩歲當然會搗亂.
所以我也是逼著要走.
然後回家繼續作為母親的工作.
因為家裡還有其他小朋友.
我們就是這樣慢慢找到一個節奏.
以致我們也看到將來可以怎樣去服侍.
當丈夫和妻子一起能夠找到同一個目標.
在神裡找到同一個目標.
真的是一個很美好的關係.
我也很鼓勵大家.
積極去為你的婚姻去努力.
因為你已經進入了這個關係.
沒有人比你有一個更加合適的能力.
或者崗位去做這個工作.
還有是去教導你的子女去培養這種觀念.
而怎樣教導他們呢.
就是要活出來 要展示出來.
他們就可以在當中學習.
第三點.
夏娃亦都是蒙恩去承傳神的賜福.
夏娃在上帝的計劃裡面.
她是生命 要生養眾多.
夏娃被阿當稱為眾生之母.
當然她有一個生育的能力.
而阿當是沒有的.
但她亦都是因為自己兩個兒子的殺害.
即是個人殺了阿伯.
而她是失去了自己的兒子.
我們很難想像到她的痛苦.
我們看看戈忍和阿伯.

$^{521}$其實他們看起來不是那麼差 邪惡.
因為他們都懂得向上帝獻祭.
可能他們都跟隨父母去獻祭.
這是我猜測的.
但他們都不是那麼差.
只不過在一次的妒忌裡.
上帝是鑽沽阿伯和他的祭物.
而不鑽沽戈忍的祭物.
戈忍是生氣.
而戈忍是殺了他的兄弟.
我不知道夏娃作為母親.
去經歷這個片段是怎樣的痛苦.
因為我想在座都不會經歷到.
但我們去想想.
當這件事發生的時候.
她的痛苦是怎樣.
阿當的痛苦是怎樣.
夏娃作為父母的痛苦是怎樣.
那個煎熬 那種內疚.
那時候沒有人.
但如果在今天來說.
這個社會是怎樣去看待母親呢.
我們的孩子都不會是殺人犯.
但我們的孩子都可能做過一些事.
去傷害他自己或傷害其他人.
有時候我們都會陷入一些情況.
其實我們都做不到甚麼.
他們已經長大成人.
他們可以為他們自己的行為去負責任.
他們可以選擇去做他們的事.
我們再做下去.
其實只會令事情越來越壞.
身為父母其實就是有這些無助的痛苦.
就是這樣在我們心裡面慢慢積聚.
但又出不來.
如果我問你.
你在教養孩子的過程裡面.
經歷過一些這樣的創傷.
你會不會再想有其他孩子呢.
對於阿當夏娃來說.

$^{561}$她目睹她的兒子這樣的悲劇之後.
還能不能夠去承擔那個眾生之母.
那種生育的使命呢.
可不可以有更多的孩子呢.
你去想一想.
很難吧.
我想和大家分享一個.
都是一個生命的點滴.
在八年前.
那時候我其中一個兒子.
他三歲.
他發燒了.
發燒了幾天.
都不退燒.
後來就不太妥了.
發現他只看著燈.
頸部又歪了.
帶他去診所.
診所說不行了.
你不是來這裡的.
馬上入院.
馬上送他去醫院.
上了去ICU的隔離病房.
經過醫生的診斷.
他患了腦膜炎.
而且他還處於半昏迷的狀態.
醫生和我們說了很多話.
我們看了MRI CT Scan的片段.
腦部就像合桃一樣.
我們看了Scan之後.
完全沒有合桃的形狀.
壓在一起.
醫生也和我們說.
你的孩子應該是不能動.
腳沒有力量.
眼睛聽東西有問題.
所以我們那時候是很沉悶的.
我們那時候剛來到感受堂.
但因為都不和弟兄姊妹熟.
所以沒有一個熟齡的支援.

$^{601}$當時心情很差.
我記得有一天.
我和我先生在吃早餐.
我們想不要太低落.
我們過著平靜一點.
過著我們會做的事.
我們就去吃早餐.
我問他現在怎樣.
我先生說如果他癱的話.
我們也有打算的.
他說最好就不行.
如果不行的話.
反而照顧就更加困難.
既然他這樣說.
我也開始預備自己.
不要緊.
這個孩子沒有了不要緊.
慢慢去預備.
我們都已經準備好放棄.
誰知手術又很成功.
他就繼續在醫院住了兩個月.
但隨後兩年多的時間.
每個星期都要出入醫院.
因為他要做很多物理治療.
他的腳沒有力量.
他的頭會比其他人大.
要檢查他的腦有沒有問題.
耳朵有沒有問題.
所以他的童年都沒有不同的.
只有家庭和醫院.
我們有很多事情不能跟他玩.
例如跟小朋友拋棄.
玩鞦韆其實是不能玩的.
兩年多就這樣過了.
有一天我再問我先生.
會不會再考慮有第二個.
這是第三個兒子.
我問他會不會考慮有第四個.
我先生就說當然不會.
如果又是這樣.

$^{641}$如果再這樣多年怎麼辦.
就是這樣.
但我們最後都有第四個.
我想說的是.
阿當和夏娃的小朋友是這樣的.
他們的小朋友是殺人犯.
再生的是殺人犯.
另外一個是死了.
是什麼感覺呢.
我們繼續看下去.
阿當與妻子同房.
他就生了一個兒子.
他叫做失突.
神給我納了另一個子子.
代替阿伯.
因為那個人殺了他.
在希伯來文的同房.
意思其實是認識.
其實在關係中認識對方.
這個也挺好的.
如果我們回到原文這樣看.
因為你同房的人.
你當然不是陌生人.
是你的配偶認識你的.
而且當時來說.
阿當知道神給夏娃的身份.
阿當是按著神的心意.
為夏娃改名.
他沒有將她改為殺人犯的母親.
沒有將她改為死人的母親.
而是憑著信心去改這個名.
說她是眾生之母.
而另外就是納了.
這個納了的意思.
其實是有一個.
好像appointed.
給了一個工作.
指派她做這件事.
上帝指派了夏娃去完成.
賦予生命的使命.

$^{681}$所以夏娃其實是.
她有承認的.
她承認自己經歷過的過去.
但她也接受了自己的使命.
去跨過那種痛苦.
跟阿當再生了一個兒子.
叫做失突.
我們參考耶穌基督的家譜.
這個失突是在家譜裡面的.
所以不是阿伯的.
是失突.
所以耶穌的祖先一直追溯到.
眾生之母夏娃.
而夏娃就是在她的後代裡面.
有一位神應許的救世主.
成為了我們今天的生命之源.
我們去想.
上帝去創造一個很完美的世界.
但是罪的出現.
然而神是沒有輸給這個殺蛋的.
他讓阿當去意識到夏娃這個使命.
她自己的使命.
當阿當的意義是死亡.
但是夏娃是一個生命.
上帝就是這樣跟人聯手.
繼續繁殖這個後代.
阿當最終是能夠用上帝的角度.
去看待他的太太.
而最終這個應許的救世主.
也是跟著這個家譜繼續出現的.
上帝沒有一開始就讓我們耶穌基督為救主.
但是他是讓阿當知道.
他們是因罪的限制.
他們是不能夠自己去幫助自己.
然後在兩千年以後.
這麼多世代的人類.
一起去期待著彌賽亞的出現.
以致他真的成為肉身.
來到我們這個世界.
我們在我們的生活裡面.

$^{721}$我們會面對很多的困苦.
我們會有困難.
我們去看著我們的配偶的時候.
我們去看著我們的父母的時候.
我們能不能夠超越.
我們在他們身上看到不好的東西.
我們能不能夠看到.
上帝在你的家庭裡面所做的工作.
放在你的家庭裡面的心意.
當你的家人做了一些你不喜歡的時候.
你還能不能夠讚美神呢.
是因著神擺他在你的生命中.
而你繼續去敬虔的去愛他們.
明天是母親節.
但我們不要只懂得母親節.
才去愛我們的母親.
我們的妻子.
我們很希望大家能夠.
藉著今天的講道.
去想想我們自己的環境.
我們怎樣可以做好.
我們在家庭裡面的角色.
怎樣可以一起去看到.
擁抱上帝那種設計.
祂在你生命中的旨意.
在你家庭中的旨意.
你能夠做到這一點的時候.
你家裡的女人或者任何人.
她們都成為這個蒙福的家庭成員.
你所體驗到的不是一個咒詛.
而是上帝對你們各人的祝福.
然後每一天都是母親節.
每一天都是父親節.
每一天都是兒童節.
我們一起祈禱.
親愛的天父.
你藉著這個母親節.
讓我們再次去回到.
你創造裡面的大能.
你的設計裡面.

$^{761}$賦予了這種特殊的心意.
就是在女性身上.
你去要她們.
去賜予她們生命.
亦都能夠繁殖後代.
今日我們能夠再去.
反思我們自己的生命的時候.
你讓我們每一個人.
生命都是母親.
無論她是近.
無論她是遠.
我相信你放在我們的心裡.
我們的生命裡面.
都是有你救贖的旨意和計劃.
願我們都在當中繼續努力.
在當中找到神.
你給我們各人的意義.
意義 責任.
每天都能夠感受到你的帶領.
我們這樣的奉主名求.
阿們.
\newpage



\section{撒母耳記上 1:10-18-20}
\label{sec:rgd8tUJeioY}
\textbf{2024-05-12 午堂《母親的禱告》撒母耳記上 1\_10-18,20  講員:黃頌言姑娘}
\newline
\newline
連結: \href{https://youtube.com/watch?v=rgd8tUJeioY}{\texttt{https://youtube.com/watch?v=rgd8tUJeioY}} ~~~~ 語音日期: 2024-05-14
\newline
\newline
\hyperref[sec:BHhp4UIkEH0]{\small{< < < PREV SERMON < < <}}
~
\hyperref[sec:index_chronic]{\small{[返順時目]}}
~
\hyperref[sec:index_scriptual]{\small{[返順卷目]}}
~
\hyperref[sec:D8TmPX_lJJQ]{\small{> > > NEXT SERMON > > >}}
\newline
\newline
撒母耳記上 1:10-18-20
\newline
\begin{longtable}{cl}
\hline
\hline
章節 & 經文 (和合本修訂版)\\
\hline
1:10 & \begin{tabularx}{0.7\textwidth}{X} 哈拿心裡愁苦，就痛痛哭泣，向耶和華祈禱。 \end{tabularx} \\ \\ \relax
1:11 & \begin{tabularx}{0.7\textwidth}{X} 她許願說：「萬軍之耶和華啊，你若垂顧你使女的苦情，眷念不忘你的使女，賜你的使女一個子嗣，我必使他終生歸給耶和華，不用剃刀剃他的頭。」 \end{tabularx} \\ \\ \relax
1:12 & \begin{tabularx}{0.7\textwidth}{X} 哈拿在耶和華面前不住地祈禱，以利注意她的嘴。 \end{tabularx} \\ \\ \relax
1:13 & \begin{tabularx}{0.7\textwidth}{X} 哈拿心中默禱，只動嘴唇，聽不到她的聲音，因此以利以為她喝醉了。 \end{tabularx} \\ \\ \relax
1:14 & \begin{tabularx}{0.7\textwidth}{X} 以利對她說：「你要醉到幾時呢？不要再喝酒了！」 \end{tabularx} \\ \\ \relax
1:15 & \begin{tabularx}{0.7\textwidth}{X} 哈拿回答說：「我主啊，不是這樣。我是心裡愁苦的婦人，清酒烈酒都沒有喝，只在耶和華面前傾心吐意。 \end{tabularx} \\ \\ \relax
1:16 & \begin{tabularx}{0.7\textwidth}{X} 不要將你的使女看作不正經的女子。我因極其難過和生氣，所以一直禱告到如今。」 \end{tabularx} \\ \\ \relax
1:17 & \begin{tabularx}{0.7\textwidth}{X} 以利回答說：「平平安安地回去吧。願以色列的神允准你向他所求的！」 \end{tabularx} \\ \\ \relax
1:18 & \begin{tabularx}{0.7\textwidth}{X} 哈拿說：「願你的婢女在你眼前蒙恩。」於是婦人上路，去吃飯，臉上不再帶愁容了。 \end{tabularx} \\ \\ \relax
1:19 & \begin{tabularx}{0.7\textwidth}{X} 他們清早起來，在耶和華面前敬拜，就回去，往拉瑪自己的家裡。以利加拿和妻子哈拿同房，耶和華顧念哈拿。 \end{tabularx} \\ \\ \relax
1:20 & \begin{tabularx}{0.7\textwidth}{X} 時候到了，哈拿懷孕生了一個兒子，給他起名叫撒母耳，說：「這是我從耶和華那裡求來的。」 \end{tabularx} \\ \\
[1ex]
\hline
\hline
\end{longtable}
$^{1}$各位弟兄姊妹平安.
(平安).
大家比較少在午課見到我.
因為平時午課這段時間.
11點到12點半.
我們在幼稚園會有兒童聖經學校.
我們會培育我們的小朋友.
所以我多數會在那邊服侍.
今天很榮幸在這裡和大家一起分享信息.
我們一起祈禱.
親愛的天父,多謝你.
我們很渴望你的話語.
願你開啟我們的心.
開啟我們的耳朵.
讓我們聽見,我們聽得明白.
多謝主,奉耶穌名求.
阿們.
今天和大家分享一個主題.
叫母親的禱告.
不知道大家有沒有試過.
為一些事情祈禱很久.
但好像天父都不聽.
我小時候有一個經驗.
小學的時候.
我經常都祈禱.
我和天父說.
天父,我可不可以所有課都是美奴課.
因為我覺得很辛苦,讀書中英數.
我很喜歡畫畫.
可不可以所有課都是美奴課.
我和小朋友.
我模樣的小朋友.
分享我這個祈禱的時候.
他們說,天父一定不聽你這個祈禱.
因為是天方夜譚.
其實我長大的時候.
我開始忘記了我自己的這個祈禱.
後來我發現原來.
我忘記了,但天父沒有忘記我這個祈禱.
我一會兒會說多一點給大家聽.

$^{41}$或者大家可能曾經都為你的家人祈禱.
希望他可以快點信主.
又或者可能你所愛的人.
身體有些毛病.
你為他禱告,但他久久都未得到醫治.
其實當祈禱很久很久.
但天父好像沒有回應的時候.
的確是會令人灰心喪志.
但我們今天來看一看.
撒慕爾記上的一章十至十八和二十節.
我們看看這位婦人.
她去禱告蒙英韻的經歷.
剛才我們都讀過了.
那個故事的背景其實是.
當時以色列人進迦南地之後.
神揀選的選民.
其實他們進入了一個非常混亂的士司時代.
當時的人聖經形容他們任意而行.
即是他們喜歡做什麼就做什麼.
很多的罪惡發生.
這個又形容為黑暗的士司時代.
原來這個時代完結.
去到另一個時代了.
進入君王時代之間.
原來神興起了一個偉大的先知.
叫做撒慕爾.
去引領他的子民跪向他.
這位領袖的誕生.
源自於以利亞和哈拿的家庭.
而哈拿就是一個不孕的婦人.
她不能生育.
這個神跡是怎樣誕生的呢?.
以利加拿的家庭.
當時來說是非常敬虔的.
因為士司時代的人任意而行.
但她的家庭竟然一年三次去敬拜耶和華.
以前的人要走很遠路程才能去到聖殿.
不像我們這麼方便.
每個星期都可以來教會崇拜.
所以一年三次已經對他們來說是很多了.

$^{81}$哈拿因為不孕.
她的丈夫以利加拿.
娶了另一個妻子叫做皮利娜.
我們看看圖片.
左邊和丈夫一起就是哈拿.
後面有很多小朋友的就是皮利娜.
其實神的心意是一夫一妻的.
但當時社會裡的人口勞動力是來自人的.
所以他們是容許家庭裡.
如果你的妻子不孕.
就可以另娶一個妻子去生養子女繼承家業.
當時結婚有小朋友其實是代表著神的祝福.
但皮利娜看到神不使她生育.
神不使哈拿生育.
當時覺得不潔的人就是神不喜悅.
神就不會賜子女給她.
所以學者認為皮利娜應該是捉住這一點去刺激哈拿.
令哈拿大大的激動.
大概她可能會這樣說.
一定是你不潔.
你做了神喜歡的事.
連神都不喜歡你.
所以你沒有子女生育.
令哈拿真的很不開心很委屈.
甚至不想吃飯.
而哈拿的丈夫很愛她.
每一次去獻祭的時候.
祭肉都會加倍給哈拿.
他也安慰哈拿說.
有我不是比有十個兒子更好嗎.
大家幻想一下如果你是哈拿的話.
聽到你的丈夫這樣安慰你.
究竟你覺得安慰嗎.
覺得安慰嗎.
雖然這位丈夫真的很有心安慰.
但他好像不太明白哈拿的心.
哈拿依然很傷心.
當時長子才可以得到相配的東西.
譬如你看到.
我幫你生育了這麼多小朋友.

$^{121}$你都不當我其中一個是你的長子.
你竟然相分的都是給哈拿.
於是他更加要去激哈拿.
這是一個惡性的循環.
所以我們可想而知.
這個家庭每年去敬拜耶和華的同時.
其實他們充滿了紛爭.
這個家庭是家嘈屋閉的.
我們可以欣賞哈拿一些東西.
第一點想和大家分享的是.
她沒有埋怨.
向神傾心討義她的苦情.
哈拿的不孕令她很無助.
亦被毗尼拿取笑她不潔不孕.
而她丈夫亦安慰不了她.
甚至在祈禱的時候.
她在聖殿裡很用心祈禱.
祭司竟然誤會了她喝醉酒.
好像沒有人明白她.
縱然面對這麼大的痛苦.
但她沒有埋怨神.
她亦沒有怪責她的丈夫.
亦沒有發洩在毗尼拿的身上.
她選擇向神傾心討義.
我們一起讀這段經文第十一節.
1,2,3.
萬君知耶和華.
你若睡姑婢女的苦情.
眷念不忘婢女.
賜我一個兒子.
我必使他終身歸耶和華.
不用剃刀剃他的頭.
好,可以暫停.
不用剃刀剃他的頭.
我們待會再說.
是甚麼意思.
但我們可以欣賞一下.
哈拿有行切的祈禱.
縱使被人取笑.
她不能生育.

$^{161}$但她甚至是不潔的人.
連製肉都不能吃.
所以哈拿連食物都不能吃.
很冤屈.
但她沒有埋怨神.
她繼續祈禱.
她沒有放棄不住的禱告.
但在祈禱中.
我們看到.
她覺得自己被神忘記了.
每一年她都這樣祈禱.
其實她祈禱了19年.
或者研究了.
大概她祈禱的年數是19年.
每年她去聖殿都這樣禱告.
每天都這樣向神祈禱.
不住的祈禱.
但神沒有回應她.
這次她有些轉變.
她改變了祈禱的角度.
她想得到這個孩子.
不是為了取悅她的丈夫.
也不是為了對抗彼尼娜.
也不是為了在人面前得到面子.
是為著神的國度的緣故.
不用剃刀剃她的頭.
意思是不幫她剪頭髮.
其實這是把孩子分別為聖.
分別為聖獻給神的意思.
這些人稱為拿西爾人.
就是一些向神許願.
分別為聖歸主用的人.
拿西爾人很多時候都是分時段.
例如我獻給神幾個月.
我願意做拿西爾人幾年.
但哈娜在這裡祈禱.
是許願她的兒子終身歸於耶和華.
拿西爾人和平時的人有什麼不同呢.
第一遠離清酒濃酒.
不會喝酒.

$^{201}$第二不會吃不潔之物.
神分了很多聖潔和不聖潔的食物.
她是不會吃的.
第三不能靠近不潔的物質.
例如屍體.
第四她們不會剃頭髮.
可能是很長的頭髮.
這個黑暗的時代.
其實就是需要一個終身歸耶和華的人.
哈娜這個祈禱得著了神的心.
所以有時候.
我們的生命裡面.
未必有人明白我們的冤屈.
甚至可能我們身邊最親的人.
我們的丈夫.
我們的太太.
我們的愛人.
都未必明白我們的冤屈.
但是我們可以親自來到神的面前.
向神傾心吐意.
原來神會將我們的苦變成福.
關鍵是我們是否將我們的祈禱.
連繫於神的國度.
很多時候.
這些福會埋藏在苦的裡面.
有很多時候我們只是停留在埋怨.
但原來在神的國度裡面.
可以扭轉一切.
只要我們相信.
關鍵在於我們是否將祈禱.
連繫於神的國度.
第二點我們可以欣賞哈娜的是.
縱使缺欠.
她都樂意奉獻.
不知道大家有沒有發現.
其實哈娜還沒有小朋友的時候.
她已經樂意將自己將會有的小朋友.
獻給耶和華.
不知道大家會不會覺得自己有些缺欠的東西.
而不敢去奉獻.

$^{241}$覺得自己都不夠.
我怎麼奉獻.
小時候我也有這些想法.
例如.
我想去宣教.
我想去傳福音.
因為我已經知道.
我要傳福音是基督徒要做的事.
但我很膽小.
不知道該跟別人說什麼.
很害怕.
會不會很想獻詩.
做詩班.
做領唱.
做敬拜隊.
但我又不懂音樂.
又不懂唱歌.
怎樣去敬拜呢.
或者會不會是.
很想關心別人.
但我有些害怕.
有些蛇恐.
不敢去走近人群.
或者很想奉獻.
但其實我最近經濟上有些困難.
奉獻不了.
原來我們可以試試這樣祈禱.
我們可以這樣祈禱.
主啊若果你讓我在這件困難的事上.
經歷你的時候.
你可不可以讓我.
可以有能力.
樂意去為你傳福音.
作見證呢.
我們也可以這樣祈禱.
主啊若果你讓我有機會在台上帶詩的話.
勝過恐懼的話.
我願意全心全人一生去敬拜你.
或者我們可以這樣祈禱.
主啊如果.

$^{281}$你讓我去看見人的需要.
你讓我有足夠的能力的時候.
我樂意去主動關心別人.
又或者可以這樣祈禱.
你遇到經濟困難.
你可以這樣說.
主啊若果你讓我在經濟上.
我可以在財政上經歷你的給予.
經歷你的供應的時候.
我樂意將這個金額的錢.
捐到去教會當中.
我們可以試試這樣禱告.
看看我們怎樣經歷神的供應.
小時候.
因為我一出生就有棉毯.
我一直都不敢上台.
甚至我小時候已經很羨慕別人.
可以上台表演真是太好了.
但我很害怕在台上.
一站上來就會被人笑.
於是我一直都不敢.
後來神給了我一個機會.
就是我中學的時候.
祂給我機會做敬拜隊.
在敬拜隊裡面帶領敬拜.
我很害怕.
很恐懼.
當時我向神祈禱.
我說主啊.
若果你要我在台上帶領敬拜的話.
求你幫助我.
不要顧別人的目光.
幫助我可以專一敬拜你.
也可以將會眾帶到你的面前.
我很多很多的祈禱.
因為很害怕.
當日我經歷了一個非常喜樂的經驗.
神改變了我的心.
我將我自己的缺陷的地方.
帶到神的面前.

$^{321}$神將我這個缺陷的地方.
這個苦變為祝福.
後來神讓我成為一個敬拜者.
讓我非常享受敬拜讚美.
你會不會都很想經歷.
在神裡面有這麼豐盛的生命.
讓我們都為神勇敢的夢想.
勇敢的去禱告.
我們第三點可以欣賞哈拿的是.
信守承諾.
我們一起讀.
《以利說》.
「你可以平平安安地回去.
願以聖餐的神允許你向他所求的哈拿傳.
願彼女在你眼前無一人.
於是苦人走去吃飯.
面上再不戴壽桶」.
哈拿這個名字其實是有恩典的意思.
他真的很相信神透過祭司以利.
允許他的祈禱.
所以他聽完祭司以利跟他說.
「你平平安安地回去吧」.
「神允許你的祈求」.
然後他就去吃飯.
很開心.
再不戴壽桶.
神將他的苦變為祝福.
也將他的受煩變為喜樂.
很多時候我們祈禱的時候很受煩.
祈完禱的時候還是沒有甚麼喜樂.
繼續去受苦.
我們可以檢視一下我們自己.
是否真正有交託給主呢.
這裡我們見證了哈拿是真正很有信心的.
我們一起讀.
哈拿1 2 3.
(漢語).
《奧義記》上一章二十節.
這裡說哈拿其實是預養小朋友.
神賜給他的兒子.

$^{361}$叫做撒母耳.
Samuel這個名字就是.
跟神母求來的意思.
或者神的名字的意思.
哈拿預養這個小朋友.
直到他斷奶大概三至六歲.
他就將孩子再帶到耶和華的殿中.
這三至六年間.
大概是我們的小朋友讀幼稚園的時間.
哈拿親自去教導他.
要預備他進入聖殿.
撒母耳記上記載了有兩個祈禱的人.
一個是哈拿.
另一個就是撒母耳.
當我翻閱的時候.
我才知道這個消息.
原來撒母耳記上只有兩個祈禱的人.
撒母耳長大之後.
神賜給他大大的重用.
更加將他成為一個偉大的先知.
他也為著他的民族祈禱.
為著立王的事去禱告.
原來一個祈禱的母親.
會有一個祈禱的孩子.
最近很感動.
有幾位家長各自來分享給我聽.
他們晚上會和家人手牽手一起去祈禱.
為著那一天的事情去禱告.
後來他的孩子.
當有一天還未捉住他去祈禱的時候.
他會自己走過來說.
我們要祈禱.
他會主動要求我們一起祈禱.
這個其實是很美好的畫面.
因為孩子平時說不出的心聲.
他可以透過祈禱去說出來.
而家長的日常.
可以在禱告中不知不覺之間.
告訴孩子,分享.
當你為著孩子的事去感恩的時候.

$^{401}$他會知道原來你察覺到他的優點.
這些都是你去肯定他的時候的時間.
其實我們有孩童奉獻禮.
其實孩童奉獻禮的意義也是一樣.
就是家長願意一生將孩子帶到神面前.
用神的話語教導他成長.
而後有一句金句.
「教養孩童使他走當行的道.
就是道路他也不偏離」.
是《箴言二十二章六節》.
哈拿在這三至六年的當中.
他很用心栽培撒母耳.
我們大概會明白.
其實孩子和家長要分離的時候.
是很困難的.
幼稚園第一天上學的時候.
家長哭得比小朋友還要厲害.
一定要預告給孩子.
「你要上學了」.
自己也要預備好.
「你可不可以不上學」.
心也要預備好自己.
才可以和孩子分離.
我有一個朋友.
她和她的兒子說明.
五歲生日的時候就要分開睡.
她兒子生日前哭得很慘.
每天都哭得很慘才睡覺.
但到了五歲生日當天.
他一滴眼淚也沒有哭過.
原來孩子比我們想像中更加堅強.
所以我們要學習放手.
不要捉緊所愛的人.
要將他交託在神的手中.
讓神去建立幫助他成長.
我們不可以帶孩子成長太久.
很多時候.
可能到了孩子的小學.
很多家長已經很擔心.
「我已經管不了他」.

$^{441}$「他手機在看甚麼」.
所以我們只能將孩子交託在神的手中.
將愛的人交託在神的手中.
讓神可以無時無刻提醒和帶領他.
哈拿每年都會見撒母耳一次.
或者三次他們去聖殿敬拜的時候.
可能會去探望他.
但聖經記載.
他每年都會為撒母耳做一件小小的袍.
其實母親多少歲都會為孩子禱告.
即使小朋友多大都會是你的孩子.
所以無論如何都會掛念你的孩子.
相信哈拿也是.
一直都會為著撒母耳禱告.
我們看到撒母耳的信心非常堅定.
哈拿自己是一個信守承諾的人.
神很眷顧他.
讓他真的有一個孩子.
其實他可以不還願.
如果他不還願的時候.
他的孩子可能只是成為家中其中一個.
母親在家中鬥爭的工具.
但撒母耳的人生如果只是這樣的話.
就會不知所措.
但當哈拿信守承諾.
他真的將撒母耳分別為聖獻給神的時候.
他的一生就在神的眷顧下成長.
我前幾天訪問我的母親.
因為我是一個信三代.
我母親信主.
我的婆婆也信主.
所以我問我的母親.
你有沒有為我祈禱.
剛才我們早課的時候.
我問小朋友這個問題.
你猜我母親有沒有為我祈禱.
小朋友說沒有.
可能沒有禱告習慣的小朋友.
會覺得難道母親沒有為我祈禱.
他可能覺得我母親也沒有為我祈禱.

$^{481}$但我母親就說.
我有 但我不是為你祈禱.
我是為你們祈禱.
因為我有一個哥哥和一個妹妹.
我繼續問.
母親你為我們祈禱什麼呢.
她說我為著求神為你們預備合適的配偶.
這個我很明白.
因為我哥哥的女朋友經常都惹怒我母親.
希望她找一個合適她的配偶.
第二件事我母親繼續說.
為著你們做的每一件事都合神心意去禱告.
我很感恩我有一個祈禱的母親.
她很放手.
其實很相信神帶領我們的生命.
從小我習慣跟母親聊天.
以致我會主動跟母親分享我的心事.
今天母親節.
相信母親都很渴望聽到小朋友跟她聊天.
將她的事告訴母親.
奉勸各位做兒女的.
多點跟母親分享你的心事.
祝福大家母親可以認識我們的天父.
以及你的生命當中.
都可以有一位為你禱告的屬靈母親.
最後神其實更加祝福家天.
阿哈拿有五個小朋友.
我們一起讀這一句.
一二三.
(聖詩).
原來神很祝福哈拿.
讓他最後加上姆爾一共有六個小朋友.
當我們樂意奉獻的時候.
原來我們不是失去.
而是神要祝福我們得著更加多.
哈拿是一個信守承諾的人.
在黑暗的時期時代.
需要一位信守承諾的人.
在我們現在的時代當中.
神正在尋找像哈拿一樣的正直的人.

$^{521}$願意一生奉獻給神使用的人.
今天我們從哈拿的禱告中.
可以明白到我們可以將所有冤屈.
在神面前去洗掃.
神可以將苦變為祝福.
也可以將受煩變為喜樂.
當我們缺失的時候.
我們也可以有奉獻的心.
當中更加可以經歷神的豐盛.
當我們樂意奉獻的時候.
我們不是失去.
而是得著更加多.
當我們經歷了神的大能的時候.
聖經告訴我們.
我們必須要還願.
要做一個信守承諾的人.
不單只是對神.
我們對任何人.
都需要去遵守承諾.
剛剛我也提到.
我小時候經常有一個祈禱.
就是求神讓我所有課堂.
都變成美奴課.
其實我將當時覺得很苦的事.
跟天父說.
其實我不想上中英數課.
後來其實神讓我讀到大專的時候.
我可以讀設計.
所以真的所有課都是美奴課.
但這個從小學到大專的時候.
起碼有十年的時間.
當我忘記了自己這個祈禱的時候.
發現原來天父竟然沒有忘記我這個祈禱.
後來我做了玩具設計師.
然後我發現.
其實跟小朋友一起玩遊戲.
建立他的生命.
好像比設計一件玩具更加重要.
因此我就站在這裡.
成為了錦繡堂的Odi姑娘.

$^{561}$我們每一個的生命.
其實都很需要這位神.
讓神使用我們的生命.
勇敢地為主夢想.
讓我們勇敢地向主許願.
讓我們將心愛的人都交託在主的手裡.
讓我們一起祈禱.
親愛的天父.
我們都為著我們的母親去禱告.
讓我們每一個都有愛我們的母親.
為我們犧牲了很多.
縱然地上的母親未必每一個都已經認識你.
我們渴望你祝福他們.
祝福每一個的家庭都有你的外遺當中.
求你教導我們禱告.
不單止只是看見自己的苦情和缺乏.
讓我們都相信你.
讓我們相信你是可以將苦變為祝福.
將受煩變為喜樂.
所有的奉獻在你面前的國度裡.
我們都不是失去的.
而是得著更加多.
求主幫助我們成為樂意奉獻的人.
奉主耶穌基督名求.
阿們.
\newpage



\section{約翰福音 16:5-8}
\label{sec:D8TmPX_lJJQ}
\textbf{2024-05-18 晚堂《看啊!聖靈駕到!》約翰福音16\_5-8 講員:楊明麗姑娘}
\newline
\newline
連結: \href{https://youtube.com/watch?v=D8TmPX-lJJQ}{\texttt{https://youtube.com/watch?v=D8TmPX-lJJQ}} ~~~~ 語音日期: 2024-05-21
\newline
\newline
\hyperref[sec:rgd8tUJeioY]{\small{< < < PREV SERMON < < <}}
~
\hyperref[sec:index_chronic]{\small{[返順時目]}}
~
\hyperref[sec:index_scriptual]{\small{[返順卷目]}}
~
\hyperref[sec:AlakgRywpkk]{\small{> > > NEXT SERMON > > >}}
\newline
\newline
約翰福音 16:5-8
\newline
\begin{longtable}{cl}
\hline
\hline
章節 & 經文 (和合本修訂版)\\
\hline
16:5 & \begin{tabularx}{0.7\textwidth}{X} 現在我要到差我來的父那裡去，你們中間卻沒有人問我『你去哪裡？』 \end{tabularx} \\ \\ \relax
16:6 & \begin{tabularx}{0.7\textwidth}{X} 只因我對你們說了這些事，你們就滿心憂愁。 \end{tabularx} \\ \\ \relax
16:7 & \begin{tabularx}{0.7\textwidth}{X} 然而，我把真情告訴你們，我去對你們是有益的。我若不去，保惠師就不會到你們這裡來；我若去，就差他到你們這裡來。 \end{tabularx} \\ \\ \relax
16:8 & \begin{tabularx}{0.7\textwidth}{X} 他來的時候，要為罪、為義，為審判，指證世人； \end{tabularx} \\ \\
[1ex]
\hline
\hline
\end{longtable}
$^{1}$弟兄姊妹 吃飯了嗎?.
真是太好了 為何這麼早吃飯呢?.
因為知道有晚堂.
原來我們知道我們會有預備.
聖靈駕到 我們要有甚麼預備呢?.
可能有人會問 聖靈?.
是誰來的? 好像不太熟.
跟聖靈熟不熟呢?.
好像不是很多人跟聖靈會很熟.
在約翰福音裡面有十六章.
十六章的意思是.
耶穌基督快要被捉拿.
快要受難的時候了.
耶穌基督才跟門徒特別說關於聖靈.
其實不是這裡才開始說.
約翰福音第三章就有提過.
教會裡我們時時認信都有讀.
我信聖靈.
但好像.
我信聖靈是怎樣呢?.
好像沒怎麼說下去.
說得多一點 我們又怕別人說.
你為何經常問關於聖靈? 你是不是靈因?.
於是我們對於聖靈的認識.
就變成去了一個樽頸位.
不是很清楚.
又好像知道.
又不敢告訴別人不知道.
於是就變成了基督徒.
為何今天會特別說到關於聖靈呢?.
你猜是甚麼原因?.
還說聖靈駕到.
其實明天是一個節日.
明天是聖靈降臨節.
有這樣的事嗎?.
教會明天沒有甚麼慶祝會了嗎?.
沒有.
聖靈降臨節好像不是很多教會一定會慶祝.
是哪些堂會會慶祝呢?.
通常是禮儀教會.

$^{41}$聖公會屬於禮儀教會.
宣道會不屬於禮儀教會.
不過我們偶爾又會提一下.
我們又沒有提得像靈因派那麼多.
我們又不是每年.
好像禮儀教會.
聖公會.
秦道為理.
每年都有慶祝.
每年在這個節期的時候.
都會有一些特別的講道.
因為明天就是聖靈降臨節.
所以今晚我想.
我們也說一些關於聖靈的事.
大家認識一下.
明天是聖靈降臨節.
我們又知道我們要怎樣去.
面對這個甚麼的日子.
耶穌基督快要受難了.
然後跟門徒說.
我快要離開你們了.
不過我去對你們好.
我去了之後.
我就會猜聖靈來.
聖靈就是保衛師.
就會保護和安慰你們.
門徒其實聽完之後是甚麼反應呢?.
有人就猜.
門徒會不開心的.
因為知道耶穌基督.
快要離開他們.
有人就會覺得.
其實門徒不是很清楚是甚麼.
不知道耶穌在說甚麼.
他們還在想.
耶穌不會被人捉拿的.
耶穌會對抗的.
所以大家對於.
耶穌說.
會猜顯保衛師來.

$^{81}$那個意識不是很強.
也都不是很記得.
最後是耶穌基督復活之後.
要跟那些婦女說.
你要去提醒我的門徒.
要去哪裡見我.
然後在升天前跟門徒說.
你要在哪裡等.
等聖靈降臨.
然後就開始做事了.
似乎門徒對這件事.
聖靈降臨都不是.
捉得很緊.
不是看得很緊.
於是要被提醒才記得.
我們是否記得勝靈降臨是何時?.
這裡應該有一個人知道.
他讀完神學.
讀完神學要知道的嗎?.
我們都應該知道的.
大家有沒有聽過五重節?.
是的.
主耶穌基督復活.
然後就升天.
復活之後計.
五重就是五十.
第五十日就應該是聖靈降臨節.
即是七個七星期.
所以是放在星期日.
然後就是聖靈降臨節.
然後我們就應該進入另一個世紀.
門徒就是這樣進入另一個世紀.
耶穌基督跟他們說.
我快要差遣聖靈.
聖靈要來做甚麼呢?.
我們聽完這句都會背.
可能會背舊的日本多一點.
為義為罪為義為審判.
然後下一句是甚麼?.
大家都在背自己責備自己.

$^{121}$不過這段不是這樣說.
那一刻是指正世人.
不過不是我們.
我們不是去做指正世人.
是聖靈來.
祂會為罪為義為審判.
指正世人.
這個是在聶文的譯本.
亦都是我們和合本現在的譯本.
前一個和合本的譯本.
後面就是.
為罪為義為審判.
自己責備自己.
都不是說我們自己責備自己.
是說聖靈來感動我們.
知道錯而我們會自責.
原來不是出自聖靈的自責.
我們可以放下.
不需要時時有一個罪疚感.
但是當聖靈責備我們的時候.
叫我們自責的時候.
我們就要謹慎.
我們又看看一些關於.
到底聖靈是誰呢?.
我剛才說過.
約翰福音第三章.
耶穌已經有提到聖靈.
是誰呢?.
他是和一個拉比.
一個叫做阿尼德姆.
跟他說將來你就會知道.
其實聖靈會來在我們當中工作.
聖靈來.
那聖靈是怎樣的呢?.
因為人要重生.
好像在水,血和聖靈底下.
來到重生.
那聖靈是怎樣的?.
耶穌基督和阿尼德姆說完之後.
我想尼哥德姆也不知道是甚麼.

$^{161}$為甚麼呢?.
他說聖靈就好像進風一樣.
那風你是看不到的.
那有沒有風呢?.
好像也有一點.
甚麼顏色呢?.
今天的風穿了甚麼顏色呢?.
我們真的不知道.
是嗎?.
如果日本正在下櫻花的時候.
我們就說風穿了粉紅色.
是嗎?.
但櫻花也有白色的.
那到底風是甚麼顏色呢?.
那個沒有的.
風你是看不到的.
不過吹過了你會知道.
風吹過了我們知道.
微風我們知道嗎?.
也知道.
大風我們知道嗎?.
知道.
強風我們知道嗎?.
不知道.
電視台也說你知道.
天文台也說你知道.
是嗎?.
其實風吹過一定知道.
無論那時候是溫文的風.
或是強烈的風.
它也讓我們知道.
那聖靈是溫文還是強烈呢?.
咦?.
待會我們來看看.
然後你想想.
到底聖靈是溫文還是強烈呢?.
那它駕到的時候.
你要怎樣的心情.
怎樣的心情去預備迎接它呢?.
那聖靈它又在舊約聖經說過.

$^{201}$我們不要想著.
它是約翰福音才開始出現的.
約翰福音第四章也有說.
我們敬拜神要用心靈誠實.
原來不是在那裡.
是很早的.
如果你讀《創世記》.
只是讀第一頁.
你也讀了的了.
《創世記》第一章的第二節就說.
「地是空虛混沌.
深淵上一片黑暗.
神的靈運行在水面上」.
它在水面上做什麼呢?.
它不是坐著一艘遊艇去旅遊.
那時候這個世界還在創造的過程當中.
第三節開始就有光了.
那時候聖靈在做什麼呢?.
它在水面上醞釀著.
接著要做的事情就是創造.
就是說它的能力那時候是很強的.
那時候風應該是很強的.
空虛混沌.
你怎樣把它整理成可以重建.
或者可以建造呢?.
大家可以去一些建築地盤看看.
怎樣將那片空虛混沌的建築地盤.
變成高樓大廈.
萬多人一呎呢?.
原來聖靈是有這樣的大能.
創造的時候.
天父主耶穌基督聖靈都一起在做.
去創造這個世界.
由那個時候開始已經有記載聖靈的工作.
不過我們少留意.
我們有時候讀過神的靈.
耶和華的靈.
我們都當作不是很清楚.
繼續讀下去算了.
pass.

$^{241}$但其實舊約很多地方都有記載.
新約有沒有?.
你猜一下新約有沒有記載?.
原來只是說新約傳書.
都記載聖靈超過200節的經文.
很多啊.
為什麼我們不知道?.
我們都沒有讀到那260多節的經文.
求主幫助我們.
我們讀到的時候能夠了解多一點.
這兩節我特別將它拿出來.
震風是很強的.
讓地盤煽成大廈.
是很強的風.
讓一個冤冤一片黑暗.
變成一個被創造出來.
有樹木花草.
有動物有人的地方.
是一股很強烈的力量.
以西結書的37章第14節.
是講到神和以西結才說.
你去那裡那裡的山谷那裡看看.
你看看看到什麼.
去到看到什麼?.
一堆的骷骨.
一堆分散了的骷骨.
然後神就令這些骷骨再組合起來.
好像一個人骨.
然後加添給它有筋有皮.
然後還是一個死人.
一陣風吹過.
神的靈來到他們那裡.
這班的骷骨變成有生命的人.
走回到城裡.
他們還生活了一段日子.
神就直著這樣告訴以色列人.
你們覺得現在你們被打仗.
被欺負嗎?.
被欺凌得很慘嗎?.
當我在你們當中的時候.

$^{281}$你們就要再得著恩典.
福氣和生命.
所以這一陣風是很強的.
神在這裡.
聖靈所顯現出來的力量是很強的.
當我們再讀到這些經文的時候.
我們要留意一下.
這一節講關於聖靈的.
其實它是說很強還是很弱呢?.
耶穌基督介紹了聖靈的工作.
聖經裡也有一些地方說到.
主耶穌基督說.
祂是為了門徒好.
也是為我們好的.
就是要差派保衛師來.
保衛師是什麼來的?.
原來它是很強.
你去感受一下.
但是它是很溫馨的.
保衛師是很溫馨的.
它是令到你傷感的時候得到安慰.
絕望的時候得到盼望.
它是令到你覺得危險的時候得到保護.
這個是保衛師.
那真理的靈是怎樣的呢?.
真理的靈是博學多才的.
聖靈是什麼都知道的.
關於耶穌基督的教導.
舊約的教導.
我們不明白背景的教導.
聖靈可以讓我們明白.
我們怎樣去迎接呢?.
如果是這樣.
我們要求聖靈幫助我們.
幫助我們在讀聖經的時候.
我們能夠得到祂為我們解釋.
能夠讀得明白一點.
讀得感興趣一點.
能夠記得多一點.
這個是不是我們的需要?.

$^{321}$我相信是的.
由我們未信到信了很久也好.
這個也是我們的需要.
祂是三位一體的神當中的第三位.
第一位是天父.
第二位是主耶穌基督.
第三位是聖靈.
祂作為一個位格.
祂有祂的性格.
祂的性格是怎樣呢?.
祂的性格是.
祂又主動又溫馨.
又跟人有很多接觸.
天父跟我們的接觸沒有那麼直接.
但是聖靈跟我們很直接接觸.
當你信主耶穌的時候.
你決志那一刻.
帶你祈禱的人就祈求.
幫你祈求.
帶你祈求.
請你一句句跟著祈求.
聖靈就被你請進心中.
求祂帶領你.
祂一路與你同在.
幫助你感動.
讓你明白神的心意.
明白聖經的話語.
而另一方面.
祂跟你接觸很直接.
所以祂感受到你的事.
也讓你接觸到祂.
你容易傷害到祂.
或者我們容易.
令到祂覺得是替我們擔心.
什麼事替我們擔心?.
原來我們不是很長進.
不是很勤力回教會.
不是很想祈禱.
見到聖經還想蓋起來的時候.
聖靈會替我們擔心.

$^{361}$這是我們需要學習.
當心裡有把聲音催促我們.
是時候回教會了.
很久沒有打開聖經了.
你有沒有為自己祈禱?.
有沒有為家人祈禱?.
當聖靈在我們心靈裡.
催迫我們的時候.
我們不要令祂擔憂.
我們要迎接祂.
接受祂的呼喚.
由心裡開始我們要向祂回應.
原來祂擔心我們.
就是怕我們.
心軟離了神.
祂不是擔心什麼呢?.
祂擔憂不是擔憂我們有沒有奉獻.
祂不需要我們的金錢.
祂需要什麼?.
祂需要我們毀身.
祂想我們將心.
完整對信靠神.
祂想我們怎樣呢?.
祂想我們能夠從聖經的話語裡.
我們能夠得到心靈裡面的力量.
得到屬靈的飲食.
這個是聖靈為我們擔心.
會有這個感覺.
祂也會看我們有沒有強烈的意志去追求.
祂有時候會感動我們做一些事情.
你有沒有試過心裡有兩把聲音.
好像自己和自己在爭論.
我應該做一下這個.
回去做一下事奉.
不過又不好.
我這麼忙我哪有空.
其實我們心裡就和聖靈一直在爭論.
你感動我事奉嗎?.
不好啦.
你休息一下吧.

$^{401}$你休息一下我也可以休息一下.
我們很想和聖靈這樣說.
這樣祂就會替我們憂心.
但我們經常推祂的時候.
我們會涉滅了祂的感動.
祂也會被我們抗拒了之後.
就少出聲.
或者是我們對祂的感覺越來越弱.
我們很想將祂.
不要啦.
我給你一個很漂亮的房間.
你休息在我心房裡面.
不要出來.
不要騷擾我的心房.
聖靈就被我們抗拒了.
但是聖靈祂是很想我們能夠直著祂得到自由.
大家記不記得這幾個星期一直在背的金句.
神是一個靈.
然後怎樣.
一會兒記得背完它.
背吧.
背完這兩個星期.
很短的.
神是一個靈.
是一個聖靈.
聖靈祂令我們能夠對罪惡說.
我有自由.
不受魔鬼的試探誘惑.
我們是否能夠得受聖靈的感動.
跟著祂來做.
能夠不受誘惑.
能夠得到真正生命的自由呢.
門徒試過.
經歷到的.
經歷到一些他猜不到的大風.
就是耶穌基督說.
叫你們在一個地方.
你要等聖靈降臨.
那一天聖靈降臨的時候.
就是耶穌基督復活之後的第五十天.

$^{441}$他們聚在一起在一間屋裡面.
祈願土.
那一天.
聖靈來到他們當中.
來到他們當中的時候.
他們得到一個恩賜.
就是他們能夠很有口才.
說了一些別人聽起來是自己國家的語言.
然後就信了主.
原來聖靈可以讓我們說話.
可以學習多一些不同的語文.
不但如此 讓我們能夠得到口才.
那一天.
只是那一天.
他們十二個門徒.
或者我們叫做十二使徒.
就在彼得的帶領下.
就帶領了三千人信主.
或者應該說三千個家庭.
嘩 一天的報道會.
我們回去都要好好為紅恩堂.
明天的報道會祈禱.
我們都很想明天有很多人信主.
原來我們需要.
我們要讓很多人信主的時候.
我們需要聖靈在我們當中動功.
聖靈讓這一班門徒.
可以說到別人很信服的話.
於是那天有三千人信主受洗.
這三千人信主受洗之後.
接著做什麼.
回家休息.
他們有的.
他們有回家休息.
但他們每天都做一件事.
他們每天都互相鼓勵.
大家一起學習神的話語.
大家一起彼此代禱.
時時都守聖餐.
這個是我們都要學習的.

$^{481}$原來信仰不是一個人的事.
這十二個使徒.
他被聖靈充滿了之後.
他傳福音之後.
這三千個家庭信了主.
就自己回家了.
整家人都信了主.
很蒙恩的 是不是.
不夠的.
這三千個家庭.
是一起學習和成長.
信仰我們不是單單一個人的事.
是大家都要一起去參與.
和互相幫助.
互相激勵.
聖靈就是這樣帶領著初期的教會去成長.
這次是三千.
接著後來又繼續有更多的.
五千.
聖靈來到我們當中的時候.
讓人信主進入人的心靈.
你決志那一刻.
他就進入你心靈做些什麼呢.
原來是做一個印記.
這個印記很重要的.
這個印記.
讓你買一層樓.
或者買一塊地.
接著之後你的地契 屋契.
那個印更加重要.
因為這個印記.
是你將來升天堂.
去天國的時候要的入場券.
我們信了神.
主耶穌基督讓我們能夠得到這個恩典.
我們被稱為義.
然後聖靈來到我們的生命當中.
成為我們的印記.
我們得了這個印記.
我們就可以一直在聖靈的帶領下.

$^{521}$在神的恩典當中去學習和明白.
祂的教導.
然後我們又能夠讓榮耀歸於神.
我們可以見證神.
見證神.
怎樣去見證神呢.
見證神除了你是講見證之外.
明天就應該是陳友會講見證.
除了你用說話去見證神之外.
更重要的是你用生命去見證神.
在《加拉太說五章》22-23節.
我相信大家都會背的.
那九個其實只有一個果子.
那個是聖靈的果子.
有九個詞語來形容聖靈的果子.
人來喜樂和平.
人來恩慈良善.
信實溫柔節制.
當你有這九個果子.
你很見證到神.
這九個原來叫做聖靈的果子.
聖靈來到我們的生命當中.
祂要在我們的生命裡面結出果子.
果子是何時結出來的呢.
現在就好像是荔枝季節.
如果你家裡有荔枝樹.
可能開始有荔枝果子.
不過真的不知道它何時結出來.
大致上我們知道是月份天氣.
但你不會知道明天會有一粒荔枝出來.
可能是一鼓作響.
它是由生命一直醞釀出來.
聖靈的果子就是聖靈在我們心靈裡面.
在我們的生命當中.
由信主那一刻開始.
祂進入我們的生命 祂塗造幫助我們.
讓我們能夠明白甚麼是義 甚麼是罪.
讓我們能夠堅持對魔鬼撒旦說.
我已經得自由.
我不需要受你的試探誘惑.

$^{561}$而我們能夠生命一直進步 結出果子.
這些果子不是一次兩次.
而是生命一直彰顯出來的時候.
越結越豐富 越結越甜美.
所以我們的生命不是要結九粒果子.
我們要繼續不斷地結果子.
讓人愛喜樂和平.
忍耐恩慈良善.
信實溫柔節制.
能夠一直在我們生命中彰顯出來.
可不可以呢.
聖靈有沒有這些果子.
原來聖靈也有.
在我們生命中忍耐很重要.
所以祂為我們擔憂.
祂有忍耐的.
祂很謙卑.
祂來的時候從來沒有打鑼打鼓.
好像風吹過 你就見到有事情發生了.
人有轉變 是不同了.
聖靈在我們的生命當中.
也要這樣幫助我們 建立我們.
聖靈來讓我們得罪恩典.
能夠明白罪 明白義.
而選擇義 放棄罪的時候.
祂做了一件事 祂指正世人.
這件事我想是很多基督徒.
或者不只是基督徒.
很多人都很喜歡做的.
指正別人有錯 是不是.
聖經有沒有教我們繼續指正別人有錯呢.
原來聖經教了我們一件事.
不要彼此論斷.
原來我們發現什麼是罪之後.
我們又能夠避開不犯罪.
不過又容易做一件事.
容易彼此論斷.
弟兄姊妹如果我們不是在一個.
需要去指正別人的崗位.
那個不是我們責任的時候.

$^{601}$我們就學習一件事.
欣賞別人的果子.
要學習自己不要受引誘.
而做一些要被指正的事.
反而聖靈穿著保羅.
都提醒我們一件事.
我們要用和平彼此聯絡.
竭力保守聖靈所賜.
合而為一的心.
原來聖靈想我們合而為一.
負責指正世人的是聖靈.
不過他還是做一件事.
他是指正世人.
他不是想指正基督徒.
他指正世人都不是想指正你有罪.
而是讓他能夠得到感動.
能夠認識神.
我們作為在主裡面的弟兄姊妹.
我們能夠有聖靈在我們生命當中的時候.
我們可以有機會結果子.
既然有這麼多個糾結果子的模樣的時候.
求主都幫助我們.
讓我們能夠合一.
我想合一是很不容易的.
合一是只要有兩兄弟或兩姐妹.
或兩個兄妹也好.
姐弟也好.
都很容易會發生的爭執和批評.
但是聖靈就是說.
祂想我們合而為一.
祂有沒有做到合而為一這件事呢.
原來祂有的.
天父主耶穌基督和聖靈.
他們不爭執的.
大家的工作不同.
角色不同.
在我們生命當中.
他們出現的時候.
做的事情也都不同.
但是他們不爭執的.

$^{641}$他們就是合一地.
來建立我們.
幫助我們.
求主也幫助我們合一.
合一可能是.
我想比九個果子更加不容易.
所以我們更加要盡心去建立.
在創世記當中.
剛才我們說到神.
聖靈用了很大的能力.
將一個冤冤之地.
都是黑暗的地.
變成一個可以重建的地方.
而以西結書說到.
一排排的骨.
沒有生命的骨.
能夠因聖靈而得著生命.
合一是不容易的.
但是合一不只是放在一起.
合一是有生命的.
求聖靈在我們當中幫助我們合一.
無論是在我們的教會.
在我們的群組.
團契.
或者是家庭.
在公司.
合一都需要放在禱告當中.
求主去幫助我們.
讓我們能夠恩著神的同在.
聖靈在我們當中的工作.
我們能夠彼此合一.
彼此合一能夠恩著我們.
認識神多一點.
明白聖經多一點.
就能夠做得更好.
求主幫助我們.
在民數記當中的時候.
有一件事.
就是民數記十一章.
不知道前幾天.

$^{681}$大家有沒有聽到.
民數記十一章發生什麼事.
李思敬牧師他都提到我們.
以色列民離開埃及.
出到曠野的時候.
有一件事經常都做的.
就不是合一.
那件事叫做什麼呢.
埋怨.
那件事叫做埋怨.
彼此埋怨.
埋怨摩西.
為什麼救我們出來呢.
我們在埃及都不捨得.
你救了我們出來曠野.
我們就死定了.
我們不需要自由.
我們只要不捨得.
這是他們的埋怨.
在埋怨當中的時候.
神就跟摩西說.
你覺得很辛苦嗎.
你都做到想有埋怨.
我找人幫你.
你不要一個人做.
然後他就在七十個長老當中.
楊聖靈降臨在七十個長老當中.
於是這七十個長老.
就一起陪伴著摩西.
帶領以色列人.
預備要進入迦南地.
對著六十多萬人.
一個摩西當然很吃力.
對著六十多萬.
不是六十多萬人.
是六十多萬個家庭.
七十個長老都做到想有埋怨.
但就是在神的恩典下.
這七十位長老.
就一直跟摩西一起.

$^{721}$帶領著以色列人.
在曠野裡幾十年.
然後就預備要進入迦南地.
我們需要有聖靈的同在.
去學習合一.
我們需要有聖靈的同在.
去參與我們事奉.
祂幫助我們.
帶領我們進入我們要去事奉的崗位.
是祂預備要給我們的.
我們一定要在祂的帶領下.
去接觸不同的人.
這七十位長老.
都要對著六十多萬個家庭.
行不行啊?很辛苦的.
是啊.
沒有人說事奉很享受.
不辛苦的.
不過辛苦得來可能都很享受.
當聖靈感動你的時候.
不要讓祂覺得擔憂.
教會真的很多東西需要有人事奉.
在國度更加需要很多東西.
有人事奉.
有報道的工作.
有人去報道宣講.
有人去跟進栽培.
有人是很強力組織的.
可以組織得教會很好.
有人很強於做關懷的.
建立的.
讓信徒能夠由初信開始.
一路成長.
我看到今天.
我們的執事.
我們的堂董.
很多原來一直在感受當中.
感受堂當中.
由未信到信主.
受洗.

$^{761}$參與事奉.
後來成為了執事長老.
也有做了團長.
做了栽培部的教師.
有些是幫忙.
我們栽培年輕人.
少年人.
和兒童.
是的.
大家都是由未信開始.
去到信主.
去到受主的差遣.
聖靈在我們當中幫助我們.
感動我們.
讓我們在祂的召喚下.
在不同的地步.
不同的崗位當中.
一起服侍.
見證神的奇妙.
讓榮耀歸於神.
藉著聖靈在我們當中的工作.
今天我們可以做些什麼呢.
今天我們可以.
馴服聖靈的引導.
不要令祂擔憂.
馴服聖靈的引導.
不要去抗拒祂.
接受祂的引導.
我們一起參與事奉.
教會需要有你有我.
一起參與事奉.
沒有什麼.
沒有什麼可以沒有了你.
我們有崇拜的時候.
需要你一起來.
不是我們需要有一個人坐在那裡.
不然就少了一個人.
而是我們需要一起歌頌.
我們天上的神.
我們需要一起去學習.

$^{801}$所以當我們有獨經日營的時候.
我們需要大家一起在那裡.
我們一起去學習.
讓我們整個群體.
能夠得著神的話語.
得著聖靈的感動.
我們有不同的崗位.
無論什麼崗位.
我們都需要大家一起參與.
大家一起合作.
我們一起低頭禱告.
親愛的天父上帝.
我們來到這一面前.
我們要多謝你.
多謝你知道我們的需要.
你知道我們的需要.
讓我們不只是能夠.
藉著耶穌基督.
我們能夠得著救恩.
天父我們能夠藉著聖靈.
能夠時時得到祂的感動.
祂的提點幫助我們.
我們求你紀念我們.
紀念我們實在是軟弱的.
我們實在是需要幫助的.
我們實在是很需要聖靈的引導.
但是我們時時都不是很收到這個引導.
求主憐憫幫助我們.
讓我們能夠得著祂的召喚.
能夠同心合意一起事奉.
互助互勉.
主啊求你幫助我們.
主啊求你給我們一股力量.
讓我們永遠都不改變對你的愛慕.
永遠都不改變要接受聖靈的感動.
求主幫助奉主耶穌基督的聖名.
阿們.
\newpage



\section{撒迦利亞書 4:1-10}
\label{sec:AlakgRywpkk}
\textbf{2024-05-19 午堂《靠主的靈方成事》撒迦利亞書4\_1-10 講員:陳麗蟬牧師}
\newline
\newline
連結: \href{https://youtube.com/watch?v=AlakgRywpkk}{\texttt{https://youtube.com/watch?v=AlakgRywpkk}} ~~~~ 語音日期: 2024-05-22
\newline
\newline
\hyperref[sec:D8TmPX_lJJQ]{\small{< < < PREV SERMON < < <}}
~
\hyperref[sec:index_chronic]{\small{[返順時目]}}
~
\hyperref[sec:index_scriptual]{\small{[返順卷目]}}
~
\hyperref[sec:0Pem3SkLOJI]{\small{> > > NEXT SERMON > > >}}
\newline
\newline
撒迦利亞書 4:1-10
\newline
\begin{longtable}{cl}
\hline
\hline
章節 & 經文 (和合本修訂版)\\
\hline
4:1 & \begin{tabularx}{0.7\textwidth}{X} 那與我說話的天使又來叫醒我，好像人睡覺時被喚醒一樣。 \end{tabularx} \\ \\ \relax
4:2 & \begin{tabularx}{0.7\textwidth}{X} 他問我：「你看見甚麼？」我說：「我看見了，看哪，有一個純金的燈臺，頂上有燈座，其上有七盞燈，每盞燈的上頭有七根管子； \end{tabularx} \\ \\ \relax
4:3 & \begin{tabularx}{0.7\textwidth}{X} 旁邊有兩棵橄欖樹，一棵在燈座的右邊，一棵在燈座的左邊。」 \end{tabularx} \\ \\ \relax
4:4 & \begin{tabularx}{0.7\textwidth}{X} 我問與我說話的天使說：「主啊，這是甚麼意思？」 \end{tabularx} \\ \\ \relax
4:5 & \begin{tabularx}{0.7\textwidth}{X} 與我說話的天使回答，對我說：「你不知道這是甚麼意思嗎？」我說：「主啊，我不知道。」 \end{tabularx} \\ \\ \relax
4:6 & \begin{tabularx}{0.7\textwidth}{X} 他回答我說：「這是耶和華指示所羅巴伯的話。萬軍之耶和華說：不是倚靠勢力，不是倚靠才能，乃是倚靠我的靈方能成事。 \end{tabularx} \\ \\ \relax
4:7 & \begin{tabularx}{0.7\textwidth}{X} 大山哪，你算甚麼呢？在所羅巴伯面前，你必夷為平地。他安放頂上的那塊石頭，人就歡呼：『願恩惠、恩惠歸與這殿！』」 \end{tabularx} \\ \\ \relax
4:8 & \begin{tabularx}{0.7\textwidth}{X} 耶和華的話臨到我，說： \end{tabularx} \\ \\ \relax
4:9 & \begin{tabularx}{0.7\textwidth}{X} 「所羅巴伯的手立了這殿的根基，他的手也必完成這工，你就知道萬軍之耶和華差遣我到你們這裡。 \end{tabularx} \\ \\ \relax
4:10 & \begin{tabularx}{0.7\textwidth}{X} 誰藐視這日的事為小呢？他們見所羅巴伯手拿石垂線就歡喜。這七盞燈是耶和華的眼睛，遍察全地。」 \end{tabularx} \\ \\
[1ex]
\hline
\hline
\end{longtable}
$^{1}$各位姐妹平安.
在二千多年前.
當主耶穌基督復活了之後的五十天.
在一個樓房裡面.
就有一百二十多個人.
都是跟隨耶穌基督的信徒.
聚集在一起去祈禱.
就在那個時間.
主的聖靈就按照耶穌基督先前的應許.
降臨到這個地上宮.
充滿在當中.
每一個跟隨主耶穌基督的人.
在這個聖靈降臨的時候.
我們從這個使徒行傳裡面.
我們就看到了.
一班很軟弱的門徒.
他又重新得到這個勇氣.
一班他自己說自己是口舌笨拙的人.
他竟然有口才來去見證上帝.
一些本來不是很願意去奉獻的.
捐獻的人.
一班的論識者.
他們也變了生命.
他們願意為著上帝的國度.
來去解開他們的前路.
一些憂傷的人.
他們在聖靈充滿的時候.
他們得著這個喜樂.
聖靈更生人的生命.
不單止在這個新約聖靈降臨之後.
其實在舊約的時候.
主的靈已經在地上宮來去工作了.
去增強軟弱人的心靈.
撒迦利亞書在當中有八個的意象.
耶和華就藉著這些的意象.
叫當時那班很軟弱.
很沮喪的猶太人.
能夠重新的有力量去重建這個聖殿.
我們就藉著在當中第五個的意象.
我們彼此的勉勵.

$^{41}$彼此的同行.
在意象一開始的時候.
就說到先知就去被主的靈來叫醒他.
當時他就好像一個剛剛睡了覺的人.
就被這個人叫醒一樣.
其實撒迦利亞先知當時見到的意象.
就好像迷迷糊糊.
其實就表達在當時以色列人那種的警況.
是夢察察的.
都好像不知道在哪裡.
不知道在做什麼.
大家有沒有試過.
很累很累的時候睡著.
一定試過.
無論你睡在床上還是躺在桌子上.
你突然間會驚醒.
或者突然間有人會叫醒你.
喂!這樣.
你那時候會怎樣.
就會覺得.
糟了,我去哪裡了.
或者.
現在幾點了.
糟了,天這麼亮.
是不是我不記得起床了.
我不記得上班了.
糟了,遲到了.
有時就是這樣.
當你在睡夢當中.
突然間驚醒的時候.
你會很慌張的.
很混亂的.
不知道自己去了哪裡.
也不知道那時候是幾點鐘.
其實就是代表著.
在當時的社會.
猶太人的社會就是這樣.
當猶太人回歸耶路撒冷.
要重建聖殿的時候.
按照以色列記裡面.

$^{81}$就說到大約有五萬人左右.
他們原本就充滿了熱誠.
就回到耶路撒冷.
重建聖殿.
但回到沒有多久.
就遇到他們去種植一些農作物的時候.
又失修.
然後在當他們重建的過程裡面.
一些外族人.
當然了.
他們霸佔了那個地方已經十幾年了.
現在你們這群人又說要回來這裡的地方.
當然不行了.
所以這群外族人.
時時就來搞破壞.
又笑他們.
你們人這麼少.
在這裡做什麼.
有時當他們在這裡做事.
在這裡挖石頭 做磚的時候.
他們又來來去去搞事.
又打人 推人這樣.
最麻煩的.
就是他們又寫信上去給波斯王.
說這群人回來這裡想造反.
千萬不要讓他們重建聖殿.
這群猶太人.
回到耶路撒冷的猶太人.
看著這些困難.
他們的士氣就越來越低落.
沒有力量再去做了.
於是一停 停工.
停了17年.
當時就是這群猶太人的境況.
其實當我在今年的年初.
我要負責去做這個錦繡堂的負責人的時候.
當我真的看到.
哇 這個教會的財務.
以前我都不會看的.
但是當我做負責人的時候.

$^{121}$我就要來看著.
在今年2月中.
我一看到程序表最後那頁.
它就說我們可以用的現金.
只有18萬的時候.
我又知道了.
原來我們每一個月的支出.
有童工的薪金.
又有租金.
家嬰 各方面.
差遣 交通津貼.
加起來每個月的支出.
差不多有100萬的.
我想一下.
2月中只有18萬.
那兩個星期就要出糧.
又要交租.
又要很多很多的東西支出的時候.
哇 看到.
又想到現在香港的境況.
香港的境況是怎樣的.
經濟低迷.
每個人都北上消費.
然後又有這個移民潮.
教會怎麼辦呢.
我說看著看著.
這些艱難的時候.
心裡面冷一冷.
有些後悔.
為什麼會答應做這個位置呢.
但是上帝就藉著這個.
撒迦利亞的先知.
他和那班回歸的百姓去說的時候.
他就說了.
有一個金燈台.
這個金燈台上有一個燈座.
在這個金燈台裡面.
又有七盞燈.
每盞燈裡面又有七條的管子.
旁邊又有兩棵的橄欖樹.

$^{161}$然後兩棵橄欖.
一棵的橄欖樹在燈台的右邊.
一棵在燈台的右邊.
這個就是上帝給撒迦利亞先知.
一個很沮喪的先知.
他看到的異象.
說完之後知不知道說的是什麼.
不知道.
我們看看這個異象是怎麼樣.
就是在這個燈台兩邊.
有這個兩邊都有這個橄欖樹.
這個金燈台對以色列人來說.
一點都不陌生.
因為在會幕或者在聖殿裡面的時候.
就一定會有這個金燈台.
這個金燈台是必須要長年累月都要點著的.
祭司要按著時候.
按著這個班次.
就來將這個橄欖油.
就倒入這個燈台裡面.
這個燈台就象徵著猶太的百姓.
也象徵著今天的教會.
要在這個黑暗的環境當中.
為主作這個明光.
照耀這個大地.
路易約翰在他年老的時候.
他被擄掠到這個拔摩島上高.
在這個拔摩島裡面.
他就見到這個人子耶穌.
右手就拿著七火星.
和七個金燈台.
在啟示裡面就解釋了.
七火星就代表著七間教會的使者.
就是七間教會的領袖.
七個燈台就代表了七個的教會.
這個教會不是一個組織.
不是一個建築物.
也都不是一個機構.
這個教會是象徵著你和我.
在這個地上高要為主來去發光.

$^{201}$在哥林多前書第六章19到20節.
保羅也都這樣說.
你們豈不知道你們的身體.
就是聖靈的殿.
這聖靈是從上帝而來的.
住在你們的當中.
所以今天你和我都有聖靈在我們裡面.
而且你們不再屬於自己了.
所以你們這條命不是再屬於自己了.
是屬於那位將我們寵價和罪惡的權勢.
所買回來的主耶穌基督.
所以我們要在我們的身體上高.
要來去榮耀上帝.
教會就要為主來去發光.
但是要想一想.
當教會現在面對的是甚麼.
經濟的困難.
人們又老化又移民.
所以實在是不夠人手去工作.
不夠資源.
再加上在教會一起事奉的時候.
人手非常複雜.
性格不同 背景不同.
做事手法不同.
再加上近來再加上政治立場不同的時候.
裡面很複雜.
可能也會有人事的鬥爭.
權力的鬥爭.
又搞小圈子.
怎麼為主發光呢?.
但是如果那個金燈台.
這個燈台不發光.
是因為甚麼呢?.
沒有油.
是因為沒有這個鑑鑫油在這個燈台當中.
所以這個燈台.
它沒有油在裡面.
它也不發光的時候.
無論這個燈台是多麼漂亮.
用的那些金是多麼純正.

$^{241}$多麼重都好.
是沒有用的.
只是一個燈台.
就放在這裡.
我們又看看.
看看上帝給撒加尼亞先知.
見到的金燈台是怎樣的情況呢?.
兩旁都有這個橄欖樹.
我們知道橄欖樹是油的來源.
在古時的時候.
他們要將那些橄欖的果子壓榨了.
然後就成為油.
然後就倒入這個燈台裡面.
這兩棵橄欖樹在燈台的兩邊.
即是說.
這個燈台時時都會有油.
源源不絕不絕地流入燈台.
所以這個燈台.
是可以時時都會有足夠的油.
來發光.
照亮黑暗的地方.
在夜中.
當我看著那十八萬.
也有一些肢體來找我聊天.
想想兩個星期之後.
怎麼辦呢?.
但很快上帝就讓我想到.
我們的上帝是耶和華以納的上帝.
祂一定為我們預備足夠.
耶和華也是我們的牧者.
我們必不自缺乏.
很感恩.
在之後的每一個星期的獻金.
就不斷的增多.
再加上我們在為了教會的方事公.
為了這個新田部屋.
我們來去有這個步行籌款的時候.
很感恩.
在今天我們看到.
我們的程序表裡面說到.

$^{281}$我們這個可動用的獻金有二百多萬.
但為什麼千萬不要停留在這裡.
哦 夠了 這麼多.
我們要繼續去仰望上帝.
繼續為主的國度我們去派上.
耶和華就應許那班子民.
他說 雖然你們人力不夠.
資源又不夠.
但是上帝一定會來幫助他們.
完成這個重建聖殿的.
所以在第八節.
上帝就這樣應許.
這個重建聖殿的所羅巴伯的領袖.
他說 所羅巴伯會安放.
這個聖殿的頂那塊石頭.
就是說這個聖殿一定會建好的.
然後最後那塊石頭.
都會由這個所羅巴伯.
來去安放好.
到時當這個聖殿一建好的時候.
全部人就會歡呼.
願因位因位歸於這殿.
就是說願榮耀榮耀歸給這個.
去眷顧他們 加能力給他們的上帝.
我們知道靠著主的力量去工作的時候.
只要那件工作是屬於上帝的.
只要我們這班作上帝子民的.
我們肯參與在其中.
肯盡力的時候.
上帝一定會幫助我們去成就的.
並且是超過我們所想所求.
上帝的使者又繼續解釋這個.
撒迦利亞先知所見到的金燈台的異象.
在這個第六節說.
他不是依靠勢力 不是依靠才能.
乃是依靠我們的寧方能成事.
在原本的文字裡面.
是沒有這個依靠這個字的.
他的經文就是這樣.
不是藉著勢力 不是藉著才能.

$^{321}$乃是藉著主的靈.
才能夠完成上帝的事情.
在一次的打仗裡面.
一個神槍手去射槍射得很準確的.
但與他一班隊伍在走來走去的時候.
不知為何失散了.
就在這個森林裡面.
他找來找去 找不回他的同伴.
但也在那個時候.
他發覺到很多的敵軍.
就在這個森林裡面去搜索他們.
他就很害怕.
於是就躲進了其中一個的山洞.
他就想.
雖然我射槍很準確.
但這麼多的敵軍.
開槍射得了多少個.
還有我這個子彈也是有限的.
所以就不停地在祈禱.
求上帝來幫助.
千萬不要讓任何的敵軍.
進入這個山洞找到他.
他一定會死的 一定會死.
當他不斷地祈禱的時候.
他看到在這個山洞裡面.
有一隻蜘蛛在這裡織網.
他就跟上帝說.
上帝你真是搞笑了.
我現在需要有一個很堅硬的門.
在這裡鎖緊這個洞.
你弄一隻這樣的蜘蛛來織網.
有什麼用呢.
上帝你真是搞笑了.
但是他很害怕.
他繼續地祈禱.
突然間他聽到.
敵軍有人走近這個洞口.
他就說.
上帝這次我死定了.
我死定了.

$^{361}$但是呢.
他接著就聽到.
敵軍的軍官在大聲喊.
喂 你傻的嗎.
他罵那個士兵.
你傻的嗎.
你看不到在這個洞口裡面.
有一個大的蜘蛛網在這裡嗎.
這個大的蜘蛛網這麼完整.
就是說沒有人進來.
就進去吧.
那你竟然還進去找什麼呢.
你傻的嗎.
不要浪費時間.
去隔離那個山洞守.
在那次經歷之後.
他找到其他的士兵.
隊伍.
他就說.
他現在很明白.
什麼叫做不是勢力.
不是才能.
而是靠主的令.
方成事.
在打仗的時候.
需要有這個勢力.
有那個軍事勢力.
這個勢力就是指軍事的勢力.
但是古時候是要靠什麼.
是靠人多.
靠人多勢眾.
現在就靠坦克車.
靠這個空軍.
什麼那些.
在當中也不是靠勢力.
這個勢力就是指人的智慧.
人的知識.
人的才幹.
人的各方面.
但是在當日.

$^{401}$這個所羅巴巴.
去重建聖殿的時候.
他真的只有五萬人回歸.
十七年十八年之前.
但當中當然有生養多了.
不過不要忘記.
有很多已經老了.
甚至很多死了.
所以你說.
真的靠人是沒有的.
是不夠的.
雖然波斯王在十八九年之前.
拿他們的錢給了他們很多.
聖殿的器具回來.
但實在是不足夠的.
在那段時間裡面.
他們知道.
人不夠.
錢不夠.
資源不夠.
財力也都不夠.
更加比那些周圍的外族.
在那裡一路來去笑他們.
來去嚇他們.
他們的心力也都不夠.
他們很軟弱.
在兩個多星期之前.
宣道會就區聯會開了周年大會.
在當中周曉輝牧師他分享.
現在香港教會的那個情況.
他說每一間教會就是.
童工不夠.
領袖不夠.
金錢不夠.
個個都移民.
也都很多的弟兄姊妹都退休.
他問一下是不是所有教會都是這樣.
哪一間教會不是這樣的.
警方請舉手.
沒有一間教會舉手.

$^{441}$但是我們大家都發出這個會心的微笑.
但周牧師說.
只要我們靠著上帝的靈的時候.
上帝一定為我們去開出路的.
靠著主的靈方能成事.
當我經過了祈禱.
答應了做教會的負責人的時候.
我最初心想.
我心裡都向上帝認錯的了.
過渡期而已.
上帝告訴我是過渡期的.
所以不對我有任何的要求.
不需要有大發展.
總之你呢.
他沒有說這句.
對不起 我想的.
他說不需要一定要有很大的發展.
我心裡想.
總之我做這個過渡期間裡面.
教會不會倒塌.
不會倒閉.
這樣就行了.
我對我自己的要求是這樣.
但聖靈提醒我.
聖靈提醒我.
「焉知你得了皇后的位分」.
「不是為了今天的機會嗎?」.
我的英文名是Esther 以斯帖.
就是當這個猶太人.
他在面對這個猶太人將會被滅絕的時候.
以斯帖皇后的叔叔.
就走去找他.
給他一個挑戰.
他就要走去圍禮.
進去見這個皇帝.
因為當時皇帝沒有叫的時候.
沒有人准許見他.
於是他就吩咐了周圍的人.
整個猶太的民族.
三日三夜不吃也不喝.

$^{481}$然後就去見這個皇帝了.
我跟上帝說.
我年紀漸漸老了.
我又沒有行政的恩賜.
又沒有領導的恩賜.
我問上帝.
在這個期間裡面.
我有甚麼可以做呢?.
我又不想白佔土地.
阻止地球轉.
但這個萬軍的耶和華.
他就藉著這段經文.
不是藉著勢力.
也不是藉著才能.
乃是靠上帝的令.
方能成就自己的事情.
高明謙博士來解釋經文.
解釋這兩次的經文.
他說其實當時所羅巴伯.
去重建這個聖殿.
比起他先前的聖殿.
規模小很多.
因為當時的人力不足.
財力也不足.
還有人們說很多負面的話.
說你現在建的那個殿規模這麼小.
比起巴比倫他們之前建的殿宇.
多麼大.
你又說你自己的上帝.
是萬軍的耶和華.
你這樣建這個殿有沒有用.
行不行啊.
我又想到我自己.
心靈裡面的掙扎.
小女子一名又將要退休了.
有沒有任何的恩賜.
行不行啊.
做不做得到啊.
但聖靈又告訴我.
昔日這個以斯帖皇后.

$^{521}$她面對這個民族將要被滅絕的時候.
她不只是叫那些猶太人.
去禁食三日三夜.
她自己和宮女.
在這個宮廷裡面.
都要去禁食.
都要去多禱告.
然後接著她就說了一句.
死就死吧.
藉著這個聖靈.
藉著這個聖靈的工作的時候.
上帝一定會幫助我們.
完成祂所要做的事情.
在第七節.
撒迦利亞先知就說了.
大山啊 你算什麼.
在所羅伯伯的面前.
你必而為平地.
就是當他過往.
他們在那十七年之前.
他們為什麼建不成這個聖殿.
停了工.
就是你看實很多很多的艱難.
在第七節裡面說.
大山啊 你算什麼.
就是說你說.
你有錢不夠.
人力不夠.
人又老的時候.
這些是艱難.
在上帝的裡面.
算得是些什麼呢.
他說這些艱難.
都將要被移為平地.
當他在耶路撒冷.
這個山上高建這個聖殿的時候.
是艱難的.
比在地上高.
平地上高建是艱難的.
但藉著主聖靈的應許的時候.

$^{561}$這個所羅伯伯.
他建這個殿的時候.
就好像在這個平地上高建殿的時候.
這麼順利.
自從一月開始.
我接任了這個負責人的位置.
到今天五個多月了.
見到我的時候就問我.
這個責任大了.
工作大了.
人的意見又多了.
你是不是很辛苦啊.
我通常都問他們.
你看到我的樣子的時候.
是不是辛苦了.
有沒有殘了.
沒有化妝的.
但他們說.
有很多人都告訴我.
你變了.
開心了.
釋放了我.
上帝告訴我.
靠他的靈才能夠成事.
以斯帖皇后的.
祂就要你去禱告.
我就學習.
每一天五點多.
我就起床.
來去禱告.
每一天.
我先去靈修.
藉著主的靈.
要隨時多方的禱告.
要警醒不倦.
為眾聖徒來去祈禱.
我深信主耶穌基督.
不斷在這位父神的寶座右邊.
不斷的為我來去禱告.
有很多的頂子妹.

$^{601}$我自己也去學習.
要多為教會.
為教會的領袖.
教會的同工來去禱告.
唯有這樣的時候.
教會才能夠得到這個福興.
上一個月.
我們一班的傳道同工.
和Leo.
就一起去杜豐山.
有一天的退休.
去杜豐山的聖殿.
他們的教會.
原來在他們的聖殿的下方.
就有一個祈禱洞.
就叫做蓮花洞.
這一座有很多的石頭.
原來上方就是剛才說的聖殿.
在這個聖殿的下方.
就是這個祈禱洞.
的入口.
這個就是他們在當日.
在這個聖殿的時候.
他們在打根基.
然後建了一個祈禱洞.
然後就建了一個聖殿.
你看看.
看不到那些字.
你們看不看到那些字?.
Yoyo可以幫到我的.
在1932年.
在建這個聖殿的時候.
他們就仿效.
這個西方.
東西方他們的寺院.
或者很多的教會裡面.
原來他們的聖殿的下方.
就是一個祈禱洞.
他們就象徵著.
這個聖殿.

$^{641}$或者這個教會.
他們的祈禱.
就是這個靈性的心臟.
所以弟兄姊妹.
多禱告 多為教會.
為教會的同工 為教會的領袖.
為教會的弟兄姊妹.
為教會的事工來去禱告.
我們幾位的同工.
是進入這個祈禱洞.
其實進入這個祈禱洞.
前面就有四個字.
「放下重擔」.
因為太黑 拍不到我們.
然後我們在當中祈禱了一個小時.
然後我們走了.
那幾個字就是.
「背起十甲」.
我們盼望.
不是單單是我們同工.
一起來去為主背十甲.
為主發光.
而是我們每一個弟兄姊妹.
我們都是上帝的教會.
所以當我們說教會不好.
不行的時候.
其實都是在說我們自己不好.
不行.
所以我們一起努力.
為主來去發光.
我們一起禱告.
今天是紀念聖靈降臨的日子.
所以如果我們願意.
我們以信心.
將雙手伸出來.
就讓聖靈充滿在我們裡面.
軟弱.
貧乏.
枯乾的生命.
願聖靈充滿在我們的生命裡面.

$^{681}$使我們每一個信徒.
都能夠重新的得力.
能夠愛轉愛人的心.
再一次被燃點起來.
燈台是順金做的.
象徵著你和我的生命.
要是純潔的.
要為主發光的.
要踏上那個成聖之路的.
就願我們裡面的生命.
被聖靈來去充滿.
求聖靈更生我們.
在這個成聖的路上.
會有很多的艱難.
有很多的難阻.
但是如果聖靈在我們生命裡面.
我們可以經歷.
經歷到活水在我們裡面來去湧流.
我們經歷到生命的釋放.
經歷到如英展翅上騰那份力量.
今天就願你和我的生命.
被聖靈來去更生.
被聖靈來去釋放.
以致我們有能力在錦宣家裡面.
和所有弟兄姊妹同工.
一起更生主的教會.
求聖靈興起我們每一個.
燕點我們的生命.
使我們每一位肢體.
都能夠成為基督的精兵.
一起參與在聖靈重整教會的工程裡面.
以致我們能夠經歷.
我們的生命是蒙福的.
有份去參與重整教會的殿.
我們感到蒙福.
感到喜樂.
感到力量.
求主聖靈充滿.
求主聖靈工作.
多謝主耶穌基督.

$^{721}$您瑞聽禱告.
是奉靠主基督的聖名.
阿們.
謝謝大家!.
\newpage



\section{約翰福音 11:35}
\label{sec:0Pem3SkLOJI}
\textbf{2024-05-19 早堂《耶穌的眼淚》約翰福音11\_35 講員:陳韋安牧師}
\newline
\newline
連結: \href{https://youtube.com/watch?v=0Pem3SkLOJI}{\texttt{https://youtube.com/watch?v=0Pem3SkLOJI}} ~~~~ 語音日期: 2024-05-23
\newline
\newline
\hyperref[sec:AlakgRywpkk]{\small{< < < PREV SERMON < < <}}
~
\hyperref[sec:index_chronic]{\small{[返順時目]}}
~
\hyperref[sec:index_scriptual]{\small{[返順卷目]}}
~
\hyperref[sec:Z6ZZRqsI0K0]{\small{> > > NEXT SERMON > > >}}
\newline
\newline
約翰福音 11:35
\newline
\begin{longtable}{cl}
\hline
\hline
章節 & 經文 (和合本修訂版)\\
\hline
11:35 & \begin{tabularx}{0.7\textwidth}{X} 耶穌哭了。 \end{tabularx} \\ \\
[1ex]
\hline
\hline
\end{longtable}
$^{1}$今晚我們一起思想基本課題.
就是眼淚課題.
所以今天選了一段經文.
全本聖經裡面最短的經文.
只有四個字 耶穌哭了.
可能大家未必很熟悉.
所以我讀一次整段給大家聽.
大家可以自己打開聖經.
也可以聽我讀.
用第十一章來讀經文.
用第十一章.
有一個患病的人名叫拉撒路.
住在白大利.
就是瑪利亞和他子子.
馬大的村莊.
瑪利亞就是用香膏抹主.
又用頭髮擦他腳的.
患病的拉撒路是他的兄弟.
他姐妹兩個就打發人去見耶穌說.
主啊 你所愛的人病了.
耶穌聽見就說.
這病不至於死.
乃是為神的榮耀.
叫神的兒子因此得榮耀.
耶穌數來愛馬大和他妹子並拉撒路.
聽見拉撒路病了.
就在所居的地營住了兩天.
然後對門徒說.
我們再往猶太去吧.
門徒說.
拉比 猶太人近來要拿石頭打你.
你還往那裡去嗎.
耶穌回答說.
白日不是有十二小時嗎.
人在白日走路就不至跌倒.
因為看見這世上的光.
若在黑暗走路就必跌倒.
因為他們有光.
耶穌說了話.
隨後就對他們說.

$^{41}$我們的朋友拉撒路碎了.
我去叫醒他.
門徒說.
主啊 他若碎了就必好了.
耶穌說是指著他死亡的.
他們卻以為是說照常碎了.
耶穌就明明地告訴他們說.
拉撒路死了.
我們有在那裡就歡喜.
這是你們的緣故.
我叫他們相信.
如今我們可以往那裡去.
當馬又稱為大土馬.
就對那合作的門徒說.
我們也去和他同死吧.
耶穌到了.
就知道拉撒路在墳墓裡已經四天了.
不過來耶路撒冷不遠.
若有六里路.
有好些猶太人來見馬大和瑪利亞.
要看他們的是弟兄.
安慰他們.
猶太人聽見耶穌就來了.
就說出去迎接他.
瑪利亞卻仍然在在的家裡.
馬大就對耶穌說.
主啊 你若早在這裡.
我兄弟必不死.
就算現在.
我也知道你無論向神求什麼.
神也必賜給你.
耶穌說 你兄弟必然復活.
馬大說 我知道在末日復活的時候.
他必復活.
耶穌對他說.
復活在我 生命也在我.
信徒的人雖然死了也必復活.
凡活著 信徒的人必永遠不死.
你信這話嗎.
馬大說 主啊 是的.

$^{81}$我信你是基督賜神的兒子.
就是那要臨到世界的.
馬大說了這話.
就回去安穩地叫他妹子瑪利亞說.
夫子來了 叫你.
瑪利亞聽見了 就急忙起來.
到那裡去.
那時耶穌還沒進到村子.
仍在馬大迎接他的地方.
那些和瑪利亞在家裡安慰他的猶太人.
見他急忙起來出去.
就跟著他.
以為他要往墳墓那裡去.
瑪利亞到了耶穌那裡.
看見他.
就伏在腳前說.
主啊 你若早在這裡.
我兄弟必不死.
耶穌看見他哭.
並看見與他同來的猶太人也哭.
就心裡悲嘆 又甚憂愁.
便說 你們把他安放在那裡.
他們回答說 請主來看.
耶穌哭了.
我們一起去看這段經文.
有關耶穌在這段故事裡的經文.
本身打算去選一段.
有關聖經裡提及到哭的經文.
其實發覺很多.
發覺聖經裡有關哭的字.
出現過209次.
眼淚出現過30次.
流淚出現過24次.
相反 笑這個字.
在聖經裡只不過出現了98次.
所以發現聖經裡.
哭比笑壓倒性的多.
所以要說一篇有關眼淚的信息.
其實很容易.
我想特別很多經文都提及到.

$^{121}$哭泣的時候.
不過今天也特意去選這段經文.
有關耶穌哭了.
經文跟大家分享.
這段經文其實我自己在神學院畢業的時候.
應該是我自己寫第一篇的講章.
很多年前 二十多年前.
我神學生的時候寫第一篇講章.
今天來到這裡再一次寫的時候.
發現其實有很多地方都很值得我們去思考.
基本上約翁福音第十一章.
是一個有關喪禮的故事.
來自於伯大利的家庭.
馬大 瑪利亞和拉薩路.
耶穌參加了一場慈禮的喪禮.
應該說因為耶穌慈禮.
就變成了一場喪禮.
本來不是一個喪禮.
本來是一個很普通的聖經裡面.
耶穌醫病的故事.
一個很典型的耶穌醫病的故事.
拉薩路病了.
姐姐就打發人去找耶穌.
懇求耶穌來醫治他.
親自來叫拉薩路能夠得醫治.
然後耶穌聽見拉薩路病了.
就故意停留多兩天時間.
這是聖經所說的.
過了兩天之後.
耶穌就悉心言的從丹河走到猶大的伯大利.
花了十二個小時.
所以拖延了總共四天的時間.
耶穌辭了.
原本是一個耶穌醫病的故事.
變成了一個喪禮.
所以當耶穌到了.
拉薩路就在墳墓已經四天.
這就是聖經裡面的描述.
於是整個經文就開始了一個過程.
首先是耶穌見到瑪大.

$^{161}$然後就是瑪利亞.
最後就是拉薩路復活.
耶穌見到瑪大.
瑪大並沒有哭.
耶穌和他聊天的時候.
瑪大並沒有哭.
瑪大就討論拉薩路的死這件事情.
他說主啊 你若早在這裡.
兄弟必不死.
耶穌就跟他說.
兄弟必然復活.
復活在我 生命也在我.
所有的人雖然死了也必復活.
你信者說是嗎.
這就是第一個的場景.
瑪大比較事務.
大家姐的角色.
沒有哭.
然後就和他好理性的聊天.
如果你來的話就沒事了.
然後耶穌就去到喪禮的第二個人.
就是二姐那裡.
瑪利亞.
瑪利亞相反是哭得很厲害.
聖經記載.
那些同在的猶太人見到瑪利亞出去.
就以為他要往墳墓那裡哭.
這句話說什麼呢.
就是說瑪利亞哭的程度.
哭到身邊的人都覺得擔心.
知道出去都擔心他會做什麼傻事.
就陪他一起去.
怕他不知道去哪裡哭.
哭得太厲害.
所以就是一個這樣的情景.
當瑪利亞見到耶穌的時候.
第一個的反應就是哭.
瑪利亞哭的時候.
想不到身邊一群猶太人都一起哭.
一個淚崩的情況.

$^{201}$在這個時候.
聖經就這樣記載.
我們留意一下.
耶穌看見他哭.
並看見與他同來的猶太人也哭.
就心裡悲嘆.
憂心憂愁.
全世界都哭.
耶穌就去到拉薩路墳前.
聖經就這樣記載.
今天我們來看的經文.
就是耶穌哭了.
雖然是一段很短的經文.
但其實是一段非常難解的經文.
如果你去看很多有關讓福音的識見書的話.
這段經文其實有很多不同人.
嘗試去解釋的角度和方法.
總括來說.
問題就是為什麼耶穌會哭呢.
表面上來看.
耶穌哭是很正常的.
反正都是喪禮.
反正都是一個很正常的死了人.
自然就會哭.
耶穌是一個真實的人.
又有同理心同情心憐憫的心.
怎會不哭呢.
不過大家不要忘記.
耶穌來到這裡是做什麼的.
其實耶穌知道拉薩路是不會死的.
甚至耶穌來的目的.
本來就是要叫大家不需要再哭.
能夠叫死人復活.
所以耶穌是懷著這種動機和鋪排.
叫眾人能夠見證拉薩路復活.
試想像如果有人患了絕症.
在醫院裡一群親人都聽到這個情況.
大家就哭得很厲害.
在病房裡受人慘磨齊哭.
而你正正就是唯一一個知道事情真相的人.

$^{241}$所謂的絕症原來只是一場誤會.
原來只是醫生搞錯了一份報告.
很老套的.
總之是醫生搞錯了一份報告.
你的任務就是要告訴這群人.
原來報告錯了.
原來是一場誤會.
你就帶著這種任務目的.
走到病房裡叫大家不用再哭.
原來是搞錯了.
當你進到病房的時候.
看到大家一起哭.
請問你會不會哭一份呢.
你應該不會哭一份.
因為你來的目的就是叫大家不需要哭.
你知道你來到這間病房要做什麼.
就是告訴他們其實他們沒有死.
不用死.
沒有病.
所以耶穌來的目的.
正正就是要用行動去告訴這群人.
拉什路將要復活.
耶穌來你們其實不需要哭.
因為祂將會復活.
這是一個神蹟.
更加是一個讓人可以振奮的消息.
當然這二千年來一直有不同的人.
解經的學者來嘗試去解釋.
耶穌哭的原因.
不知道你一直怎麼看這段經文.
有人認為耶穌哭了.
基本上就是要證明耶穌是人.
證明祂是一個真實的人.
有人認為耶穌哭不是為了拉什路哭.
而是為了全人類哭.
為了生命的短暫.
人的生命的脆弱而哭.
有人認為耶穌哭不是為了拉什路.
也不是為了瑪利亞哭.
而是為了當時不信的人.

$^{281}$情緒很激動.
正所謂氣到哭.
所以是一個很憤怒的表現.
所以這些真的是.
找到聖經的學者嘗試去解釋.
但其實也不是太完滿.
也不是一個完美無瑕的解釋方法.
也不是一個完美的答案.
所以究竟耶穌為什麼會哭呢?.
耶穌為了什麼原因而哭呢?.
當我再次來到寫這篇講道經文的時候.
我自己近日有一個新的看法.
我發現原來哭或者是眼淚.
其實未必一定需要很多合理的原因.
未必一定需要很充分的理據.
甚至乎一個很合理的解釋.
哭是什麼事情呢?.
我分享一下我自己的事情.
我爸爸近日去世.
簡單來說.
爸爸去年患上胃癌.
暑假的時候開刀切了胃.
直接由吃到駁小腸來生活.
基本上不容易.
要靠小腸來消化.
體重突然由170多磅.
驟降到120多磅.
適應了大半年的時間.
今年年初的時候.
大概新年年初十的時間.
他就開始說痛.
痛得很厲害.
以為他不適應.
小腸可能吃東西覺得很辛苦.
頂住這樣.
他發現原來情況很不理想.
醫生就說癌細胞擴散.
擴散到整個肚腹都是.
所以就令他痛.
基本上動到一點點就覺得痛.

$^{321}$然後就出院.
從內科到外科.
從外科到腫瘤科.
腫瘤科到舒緩科.
大概在一個月之前.
醫生就打電話來.
星期五的時候.
說爸爸的情況似乎不太理想.
然後就開始那天叫很多親戚.
來探望他.
因為星期六午晚的時候.
他就開始睡覺.
之前他就睡不到.
因為痛得很厲害.
星期五晚上就開始睡.
因為就影響到肝的情況.
就開始睡得很多.
星期六就陪了他一天.
在床邊那裡.
基本上全日都是睡了.
就是都是睡.
醒了一次半次.
晚上的時候護士就安排給我們.
有一間私人的房間.
就是一間可以關燈的房.
最後就和媽媽一起.
和爸爸在身邊祈禱.
就離開.
關燈就讓他睡覺.
星期一我就出去港島.
其實都不想說.
因為都知道爸爸可能隨時都會離開.
但是都接了.
就去港島.
港島之後就聽到電話.
爸爸就離開了世界.
我想說.
我發現在這一年來.
由去年暑假到這段時間裡面.
其實我自己都沒有怎麼哭過.

$^{361}$認識的人都知道.
其實我都很容易哭.
以前我在神樂院裡面教書.
教到某一課.
我都會很感動地哭.
以前在教會裡面.
就是令聖餐都會很容易哭.
就算我上台看驅魔人.
我都會哭.
總之都算是很容易哭的人.
我發現.
我自己都令我覺得自己很驚訝.
這次我自己都沒有認真地哭過.
由去年爸爸知道有癌症開始.
到陪他一起去做切換.
首先前一晚大家一起吃最後一餐晚飯.
做完之後.
跟護士醫生去談爸爸的康復.
適應沒有胃的時候怎樣去生活.
扔完的奶.
很多很多的東西.
到年初的時候.
爸爸開始就說不舒服.
陪他照PET scan.
發現擴散.
一起想想怎樣照顧爸爸.
如何跟醫生和護士聯絡.
找醫生聯絡.
Google爸爸的情況.
很多很多的事情要去處理.
在這大半年裡面.
我發現自己都沒有怎樣哭過.
然後在一個月前.
當我說完之後.
就跟女兒和太太一起趕去醫院.
去見爸爸最後一面的時候.
在車上就開始哭.
去到醫院的時候.
進到病房.
見到爸爸的身體的時候.

$^{401}$我就不斷地眼淚流.
在這大半年裡面.
第一次來到他.
很認認真真地來到他.
不斷地哭哭哭.
哭了幾分鐘.
嗚嗚嗚嗚.
哭得很厲害.
但當我哭的時候.
我發現有一個很特別的感覺.
我發現自己流出來的眼淚.
好像不是我這樣.
好像不是我控制的.
甚至乎我感覺上面的眼淚.
好像是不知道什麼地方流出來.
然後大概哭了很久.
哭了幾分鐘之後.
妹妹就問我問題.
不知道問什麼問題.
然後我的眼淚就突然間cut out.
突然間就像山水喉一樣.
突然間說停就停.
然後就回答問題.
然後就開始回到事無無.
得目思離.
然後就處理爸爸之後的事情.
眼淚真的好像一個水喉水.
一流出來就很大聲地流.
然後說關掉.
突然間就cut out.
不知道大家有沒有試過.
有沒有試過這樣的情況.
所以這個就是我第一次.
爸爸這大半年裡面.
一個哭的經驗.
然後爸爸離開世界之後.
第二次哭.
就是上次我回來錦宣港島的時候.
回來錦宣港島.
弟兄姊妹傳道人就問候我.

$^{441}$當大家感覺到大家關心的時候.
我發現我自己不是一個.
來到港島的目思傳道人.
而是一個很普通在教會裡面.
被愛的孩子.
然後我的眼淚又再一次流出來.
我想說什麼呢.
我發現原來我再一次來到.
去認識自己的眼淚.
究竟眼淚是什麼呢.
弟兄姊妹.
我記得我幼稚園的時候.
有個同學叫做鄧成肯.
很小的時候.
突然間流鼻涕.
跟我說你打我吧.
然後我就用積木打他.
他就哭哭哭.
然後我就被老師罰.
罰抄他的名字.
我記得名字叫鄧成肯.
就是這個意思.
留堂就要罰抄他的名字叫鄧成肯.
一直抄一直哭.
等媽媽回來接我.
我中四的時候失戀.
哭得很厲害.
三更半夜在家裡哭.
然後我媽媽就說你不要再吵了.
這是我中學的時候哭.
幾年前自己面對一些委屈.
事奉上的一些不開心.
都在哭.
但我發現這一次的哭.
有些不同.
我再一次來到去認識.
什麼是自己的眼淚.
我發現眼淚的意義.
其實不一定是悲傷.
眼淚真正的意義其實是愛.

$^{481}$當然悲傷是會流眼淚.
但眼淚真真正正的來源.
不是因為傷痛.
眼淚的來源是因為心裡有愛.
眼淚分兩種.
第一是你自己為自己流的眼淚.
第二是你為別人而流的眼淚.
我們每個人生下來.
就立即哭.
生出來的剎那.
就大聲地哭.
目中無人地哭.
我記得我女兒出生沒多久就哭.
哭了一百天.
我們說這條淚叫做八日紅.
是真的.
一百天之後就開始停.
因為一個嬰兒哭是你的生存伎倆.
你想喝奶你哭.
你想睡覺你哭.
你想爸爸陪你哭.
當你慢慢長大的時候.
你的眼淚一直都是為你自己而流.
考試不合格你就哭.
考不到想讀的大學就哭.
受到委屈.
工作不開心.
家裡情況不好.
你就會哭.
這些眼淚都是我們人生裡.
為我們自己而流的眼淚.
這都是我們第一種眼淚.
然後生命裡面有一剎那.
你發現你的哭泣.
你的眼淚並不是為你自己.
這是一個很奇妙的一剎那.
當你的眼淚不是為你自己而流.
你的眼淚是為別人而流.
為你所愛的人.
甚至乎是為你一些價值.

$^{521}$為你一些理想.
為你所愛的地方而流眼淚的時候.
這種眼淚其實是一種.
我們人生裡面很特別的體驗.
這些眼淚其實是溫暖的.
這些眼淚不代表悲傷.
原來眼淚是代表你心裡面有愛.
所以你問為何耶穌會哭.
主耶穌的眼淚正是告訴你.
耶穌哭了其實不是純粹因為悲傷崩潰取弱.
而是因為他心裡面有愛.
為什麼耶穌哭了.
下一句說話是什麼.
用第十一章三六節.
「你看他愛者日,是何等很遲?」.
耶穌的眼淚告訴你.
眼淚為別人而流.
因為你心裡面有愛.
耶穌的督是上帝愛的道成肉身.
如果眼淚只是人世間受苦和憂傷的記號的時候.
上帝本來就不需要哭.
因為上帝是沒有眼淚的.
但如果眼淚是愛的記號的時候.
耶穌的眼淚就不需要特別什麼特別什麼合理的原因.
因為一個人心裡面有愛.
眼淚就在這裡.
眼淚是愛的記號.
因為愛是分擔別人的痛苦.
你願意或多或少去承擔別人的痛苦的時候.
這個就是愛.
所以我們英文說fall in love.
是墮進去的.
是一種結種去受傷的體驗.
這個意味著什麼.
意味著愛本來就是願意去開放自己.
隨時會受到或多或少的傷痛.
所以愛很難.
因為愛就是要去承擔一些傷痛.
你願意這樣做.
一個拒絕去愛的人.

$^{561}$他不會受到任何傷害.
因為他就和這個世界隔絕.
變得冷漠無動於衷.
所以看電影很少看到反派會哭.
我們什麼時候見到反派會哭.
反派通常會怎樣.
都是笑的.
極少見到反派會哭.
就算哭都是假裝.
宮中計那些.
宮廷劇那些太后哭都是假裝出來的.
反派最真正的成績就是笑.
因為他不需要去背負別人的痛苦和傷痛.
由蝙蝠俠的小丑.
到宮女,太后.
全部都不會哭.
相反電影裡面一些好人.
電影裡面一些正直的人.
當他願意去為別人想的時候.
願意和他所愛的人去分擔痛苦.
眼淚自然就會在那裡.
所以你明白為什麼耶穌會哭.
耶穌哭是一件很不得了的事情.
因為他的眼淚不是為他自己的.
上帝本身是不需要流眼淚.
上帝的眼淚正正是為世上一切受苦的人流.
十字架也是一樣.
耶穌丁十字架.
耶穌丁十字架是為世人的緣故.
耶穌的眼淚也是一樣.
叫世上每個人每一滴的眼淚都不是白流.
如果這個世界上沒有上帝的時候.
如果上帝不在的時候.
我們的眼淚只不過是一些生理反應.
一些失敗和脆弱的表現.
但如果上帝願意為我們流眼淚的時候.
我們的眼淚就變得不再一樣.
教徒的眼淚正正因為這樣而變得不得了.
前一陣子聽到天主教神父的一番說話.
叫恩保德神父.

$^{601}$一個上年離開世界的天主教神父.
一個滿口流離廣東話的意大利人.
說了這樣的見證.
在一套利瑪竇音樂劇上台分享眼淚的題目.
他說他小時候是一個很頑皮的小朋友.
媽媽很慘.
每天都要抓著他上學.
一放手他就立刻逃學.
離開學校很晚才回家.
每天都是這樣.
幾個月之後他媽媽就說.
你不喜歡上學就算了.
他沒有生氣沒有打他.
他說了一句話令他很記得.
他說你記住你越頑皮我就會越愛你.
然後他說了一句這樣的說話.
他說媽媽的眼淚是子女的肥料.
每個媽媽都會為自己的子女流眼淚.
流很多很多的眼淚.
這裡很多的媽媽都知道.
基本上子女五歲前都是一些體力勞動.
但長大後的辛苦再不是體力問題.
而是媽媽抽心地感覺到悲傷.
看到兒子不長進.
怕事懶惰.
顧著打機不聽話.
頂著自己.
很多這樣的眼淚不斷流出來.
你發現你越愛你的小朋友.
就越同意這些為小朋友流的眼淚.
其實是會很多很多.
但這個神父就說.
媽媽或父母的眼淚正正就是子女的肥料.
所以當我們眼淚有愛的時候.
是不得了的事情.
所以今天這篇講道.
不是要恭賀大家繼續哭下去.
不過是我們知道我們哭什麼.
知道我們哭的目的.
當我們發覺眼淚不再為自己流.

$^{641}$開始為我們所愛的人流.
為我們著緊的理想.
甚至為我們這個地方流的時候.
這些眼淚是不得了的.
這些眼淚將會成為上帝使用的肥料.
讓這個地方變得更好.
讓我們所愛的人可以繼續建造.
聖經解釋錄最後這樣說.
上帝要拆去一切的眼淚.
不再有死亡.
不再有哀號.
哭泣.
不代表眼淚是不好的東西.
不代表眼淚是和上帝國度相違背的東西.
眼淚正正是去築成我們天國的肥料.
我們的眼淚慢慢流出來的時候.
不是徒然.
這些眼淚因為耶穌的哭泣的緣故.
這些眼淚會成為上帝使用的力量和恩典.
流淚殺腫的必歡呼收割.
基督徒的眼淚將會成為上帝使用的東西.
當我們心裡有愛的時候.
所以想想我們今天.
不知道上一次你哭是何時.
哭得很淒涼的事.
想想我們是為了什麼而哭.
我們為了誰而哭.
當我們為我們所愛的人.
當我們為我們一些執著的事情.
去感覺到傷痛的時候.
如果你的眼淚是有愛的時候.
這件事我們作為基督徒.
是一個不得了的光榮.
你在聖經裡特別的記載.
你哭泣因為你就是愛.
你叫我們每個人生下來有愛.
為別人為我們所著緊的事情.
一起願意去承擔痛苦的時候.
求主你安慰我們.
更加讓我們心靈裡面.

$^{681}$仍然去擁留這份的愛.
為我們身邊所愛的人.
為一些我們覺得重要的事情.
我們願意去承擔當中的後果.
一起為他們著緊的時候.
求主你繼續讓我們心靈裡面有愛.
讓我們的愛去流露.
讓我們的愛能夠流露到.
我們這些人的關人裡面.
願主你這樣來陪伴我們.
奉主名求 阿們.
\newpage



\section{申命記 34:10-12}
\label{sec:Z6ZZRqsI0K0}
\textbf{2024-05-25 晚堂《面對面的認識》申命記 34\_10-12  講員:盧思源姑娘}
\newline
\newline
連結: \href{https://youtube.com/watch?v=Z6ZZRqsI0K0}{\texttt{https://youtube.com/watch?v=Z6ZZRqsI0K0}} ~~~~ 語音日期: 2024-05-28
\newline
\newline
\hyperref[sec:0Pem3SkLOJI]{\small{< < < PREV SERMON < < <}}
~
\hyperref[sec:index_chronic]{\small{[返順時目]}}
~
\hyperref[sec:index_scriptual]{\small{[返順卷目]}}
~
\hyperref[sec:luI8_7RFevk]{\small{> > > NEXT SERMON > > >}}
\newline
\newline
申命記 34:10-12
\newline
\begin{longtable}{cl}
\hline
\hline
章節 & 經文 (和合本修訂版)\\
\hline
34:10 & \begin{tabularx}{0.7\textwidth}{X} 以後，以色列中再沒有興起一位先知像摩西的，他是耶和華面對面所認識的。 \end{tabularx} \\ \\ \relax
34:11 & \begin{tabularx}{0.7\textwidth}{X} 耶和華差派他在埃及地，向法老和他的一切臣僕，以及他的全地，行了各樣神蹟奇事， \end{tabularx} \\ \\ \relax
34:12 & \begin{tabularx}{0.7\textwidth}{X} 又在以色列眾人眼前顯出大能的手，行了一切大而可畏的事。 \end{tabularx} \\ \\
[1ex]
\hline
\hline
\end{longtable}
$^{1}$大英姐妹平安.
願意我們開始之前.
我們都有一個祈禱.
一同去敬拜.
一同聽你自己的話語.
一時間孩子要去分享你自己信息的時候.
求真理的聖靈.
都進到我們的心裡面.
打開我們的心.
讓我們進入真理的當中.
我們被主你自己的話語去得著.
我們被主的話語去提醒.
被主的話語去幫助.
我們恭敬的.
將以下時間交上.
奉主耶穌基督明智祈求.
阿們.
大英姐妹有沒有聽過這句話.
看得明白.
人生得一知己足矣.
就是說如果用口語化.
就是人生得一知己死而無憾.
這句話原來反映.
朋友原來對我們生命.
都是很重要的.
有些研究資料顯示.
一個人在十歲至二十四歲的時候.
他們的社交圈是最大的.
平均大約有九個朋友.
互動的朋友.
經常會聯絡.
不過到了三十歲的時候.
就開始.
降低了.
只有七個.
之後可能逐漸減少.
可能因為工作.
家庭.
或者照顧朋友.
都沒有很多時間去維持這些關係.

$^{41}$而在研究中.
人擁有朋友最大的數量.
調查.
因為有受過時間限制.
大腦可以粗略地處理.
大約一百五十段友誼的關係.
不過這些友誼.
都有分親密的程度.
友誼的深度.
都分了四個層次.
第一個就是.
熟人.
什麼是熟人.
熟人就是一些你經常見到的人.
可以一起聊天.
譬如問問吃飯了嗎.
最近怎麼樣.
可能是一些閒聊.
不過這些可能只是礙於點頭之交.
可能是打個招呼.
雙方沒有一些很想加深聯繫的慾望.
或者想著再建立.
這些稱為熟人.
第二個.
普通朋友.
都會跟你一起活動.
定期你都會見到他.
你對他都有基本的了解.
甚至乎在活動中.
你可以跟他又說又笑.
可以聊天.
又或者除了活動之外.
你們偶爾都會約出來見面聯繫一下.
不過.
如果不見面.
你又不會覺得特別牽掛.
這些就稱為普通朋友.
第三種.
親密朋友.
Close friend.

$^{81}$親密朋友多數是跟熟人開始的.
兩個人會互相欣賞.
可能彼此的信任.
慢慢可能會分享多點自己的事.
自己的私事.
可能都會去分享.
亦都很享受大家一起相處的時間.
簡單來說.
當然如果什麼是親密朋友.
就是你睡不著.
兩點鐘.
可能突然間打電話給他.
哭著打給他.
他仍然會接你的電話.
那個都應該是你的親密朋友.
都是陪伴你的人.
第四種.
摯友.
摯友比親密朋友更加緊密.
就好像你的家人一樣.
他知道你內心很多深層的事.
甚至乎可能一些秘密.
他都會知道.
而且這些摯友.
他不會令你失望.
或者不會陪伴你.
你知道他永遠都是陪伴你.
都是會支持你的.
而在美國有在大學研究過.
分析了255位成年人.
搬到一個新環境.
亦都分析了112位大學生.
即是入到大學的時候.
他們了解他們怎樣認識陌生人.
即是看他們怎樣相處.
原來發現一個陌生人.
成為朋友.
普通朋友.
陌生人成為普通朋友.
原來他們需要.

$^{121}$大約40至60個小時的相處.
如果要發展成親密朋友.
就需要大約90個小時.
如果摯友.
大家覺得多久呢.
原來他們的調查.
即是在大學研究裡面.
調查的時候.
要成為摯友.
大約要200個小時.
大家立刻想想.
原來亦都在當中說.
如果兩個人相處.
超過80至20個小時.
基本上.
他由普通朋友.
成為親密朋友的機會率.
是超過50\%.
簡單來說.
你有時間的時候.
即是跟他有80至100個小時.
基本上你都可以由普通朋友.
成為親密朋友.
不過大家要記住.
這個不包括工作時間.
即是可能我上班.
我對著我的老闆.
對著我的同事.
我由早上九晚五.
我都對著他.
但這個調查是不計工作的.
即是如果一個關係.
只是留在工作的情況下.
不容易成為朋友.
現在是說工作以外的相處.
而且在整個調查裡面.
原來一個人.
人生如果有三至五個親密朋友.
其實你都有足夠的情感支援.
又假若如果你有一個朋友.

$^{161}$當你是好朋友.
你呢.
你的生活的滿意度.
是會提高的.
即是簡單來說.
即是原來朋友.
對我們生命都是很重要.
絕對可以提升我們的幸福感.
回到今天.
想跟大家分享的聖經.
當中提到摩西.
我相信大家都不陌生.
摩西是一個怎樣的人.
即是摩西.
他昔日是被神重用的人.
他帶領以色列人.
離開了埃及遺老之家.
在西乃山上接受神.
那個.
神十誡的頒佈律例.
接受訓誨.
而且我們知道.
摩西和以色列人.
他經歷了四十年.
在曠野的生活.
今天和大家看申命記.
大家有聖經.
都鼓勵大家去翻閱一下.
申命記.
提到.
以色列人在曠野.
滿了四十年之後.
準備入迦南之際.
摩西要向新的一代.
以色列人.
去重新展述那個律法.
和他們立約.
教他們進入迦南地的時候.
應該怎樣.
而我們今天看到.

$^{201}$剛才說申命記的三十四章.
就是申命記的最後一章.
也是摩西五經裡面最後的一章.
當中.
對摩西有一個總結.
和一個評價.
經文.
三十四章十字這樣說.
以後.
以色列中再沒有興起一位先知像摩西.
他是耶和華面對面所認識的.
如果我們看.
呂鎮忠譯本.
這個譯本.
這句經文.
這句聖經是這樣譯的.
以後在以色列中.
再沒有興起神賢人像摩西.
他乃是永恆主.
面對面的知教.
我們看到聖經除了在這段.
提到摩西和神面對面之外.
其實還有另外兩處的經文.
都提到摩西面對面.
在《出埃及記》三十三章.
十一節那裡.
他說.
耶和華與摩西面對面說話.
好像人與朋友說話一般.
也在《民數記》裡面.
十二章七至八節.
這裡說.
我的摩西.
我的僕人摩西不是這樣.
他是在我全家盡忠.
我要與他面對面說話.
但是我們記得.
如果讀聖經的時候.
我們記得神曾經說過.
沒有人可以見到神的面.

$^{241}$是的.
在《出埃及記》.
這句話都是神和摩西說的.
祂說.
只是你不能看見我的面.
因為沒有人看見我的面.
還可以存活.
那麼.
我們想.
摩西.
如何與神面對面說話呢.
所以我們知道.
這裡說的面對面.
不是說摩西真正見到神的面.
而是在形容.
摩西與神的關係.
如果回到剛才我們說.
《民數記》十二章那裡說.
這裡的面對面.
如果在原文希伯來文.
是意思口對口.
口對口不是代表.
嘴對嘴.
口對口的意思.
其實就是在說.
神與摩西的對話.
是直接無間的.
摩西與神的關係.
是非常之密切.
而且神與摩西說的話.
是明講的.
很清晰的.
不需要解釋.
不需要糾纏.
是很清楚的.
所以我們見到摩西.
與神的關係.
就好像剛才我們提到.
他們好像是自有的關係.
他們很清楚.

$^{281}$對方.
大家.
想不想.
有與神的關係呢.
不過.
如果我們.
回看摩西的一生.
其實摩西一開始的時候.
他與神不是有這麼親密的關係.
我們知道他經歷了很多.
今天的經文.
如果我們.
翻開前一點去看.
第七節看到.
其實摩西死的時候.
是一百二十歲.
曾經有些學者.
將摩西的人生.
分為三段.
三個的四十歲.
他們說摩西第一個四十年.
摩西的態度.
就是我能救.
因為.
頭四十年.
摩西是在皇宮長大的.
他出生的時候.
聖經記載當時法老要下令.
將所有以色列所生的男嬰都要殺.
但神保守了摩西的生命.
最後.
被法老的女兒收為兒子.
如果我們在《士徒行傳》七章二十至二十二節說.
就在那時.
摩西生了下來.
神看為盡美.
在父親家裡被撫養了三個月.
他被雕器的時候.
法老的女兒十了他.
當自己的兒子撫養.

$^{321}$摩西學了埃及人一切的學問.
說話辦事都有才能.
在這裡我們看到.
當時摩西身為一個埃及王子.
他有很高的學識.
他.
正常有這樣的背景.
應該就是可以為神做大事.
也可以改善他自己以色列人同胞的生活.
不過如果我們繼續讀聖經的時候.
我們看到.
在《出埃及記》裡面提到.
有一天摩西因為他知道自己是以色列人.
是希伯來人.
所以他看到一個埃及人.
打他自己的希伯來人兄弟的時候.
他看到旁邊沒有人.
他就將埃及人殺了.
就藏了屍.
最後這件事我們知道他被法老發現了.
所以法老要殺摩西.
摩西要逃避法老的時候.
就去到米殿.
曠野居住.
我們看到在這段時間.
摩西原本.
他關心他自己的同胞.
見義勇為.
為他們抱不平.
理應是對的.
不過.
很可惜.
他這樣做.
他不是光明正大地做.
他用了一些不正確的方法.
去做了一些.
他覺得自己認為合理的事.
最後我們看到的結果就是.
不單止他的同胞拒絕他.
甚至乎他被人追殺.

$^{361}$最後.
他去到曠野米殿.
結束了他第一個階段的40年.
如果他看第二個40年的時候.
那些學者會說.
第二個40年摩西的態度是.
我不能 我不行.
因為我們看到在這40年的時候.
摩西在曠野做什麼呢.
牧羊嘛.
摩西長時間在曠野學習.
不過.
摩西在曠野生活裡面.
是一個很重要的階段.
因為原來.
神就在這段時間去鍛鍊他.
摩西以為他自己.
很失敗.
覺得現在都被他自己的同胞拒絕.
甚至乎被人追殺.
覺得自己不行了.
誰知道在這個時間我們也知道.
正正好像一個.
很普通的過程的時候.
神就在這段時間預備他.
我們知道.
摩西在曠野的時候.
他遇見了一件.
很奇妙的事.
他看到在曠野的荊棘.
被火燒著.
不過沒有被燒毀.
神就在荊棘裡面向他顯現.
也在當時神.
吩咐.
呼召摩西.
去領以色列人出埃及.
在《出埃及的三章》那裡說.
上帝和摩西說.
我要打發你去見法老.

$^{401}$使你可以將我的百姓.
以色列人從埃及領出來.
我們看到.
在這個時間上帝就.
呼召摩西.
要去.
帶以色列人.
出埃及.
不過如果我們看整章.
《出埃及的三章》.
我們看看摩西有什麼反應.
我簡單截錄了.
摩西的反應是什麼呢.
他的心態是.
我不行的.
他向神推辭.
摩西和神說.
我是誰.
怎麼可以見法老.
怎麼可以將以色列人.
從埃及地領出來.
如果我們繼續看.
《出埃及的三章》.
我們看看神怎樣回應他.
神說.
我必與你同在.
不過.
你猜摩西又怎樣反應呢.
當時.
摩西繼續說.
那些百姓不會信我的.
也不會聽我的.
他們一定會說.
耶和華沒有向我顯現.
於是.
我們也知道.
上帝怎樣去肯定他.
就是叫摩西將他的像.
丟到地上.
變成一條蛇.

$^{441}$然後拿回來.
那條蛇又變回一隻像.
神叫摩西可以行到一個神蹟.
憑這個權柄去見法老.
帶領以色列人出埃及.
但我們看到摩西怎樣回應呢.
摩西繼續說.
主啊.
我素日不是能言的人.
就是從你瀑人說話以後.
我也是這樣.
我是絕口不洩的.
摩西是這樣應對神.
但神怎樣回應他呢.
神說.
誰做人的口呢.
是我.
我必賜你口才.
指教你當說的話.
即是上帝是這樣回應他.
但你又看看.
如果看整個《出埃及的三章》.
摩西又怎樣回應呢.
摩西繼續說.
主啊.
你願意打發誰就打發誰.
即是摩西說你不要叫我了.
你找別人吧.
這裡聖經裡面形容.
耶和華法老.
他就跟他說.
我現在叫你哥哥一起跟你去.
即是阿倫跟你一起去.
我必賜你和你哥哥口才.
阿倫要代替你去說話.
等一會.
如果我們單單看.
摩西頭兩個四十年.
覺得跟我們都好像有點像.
有時可能我們都好像摩西一樣.

$^{481}$有時可能我們覺得.
我們可以自己靠自己去做一些事.
特別那些事.
如果我已經有想法的時候.
最好.
上帝你不要干預.
即是我們可以的.
讓我們自己搞定.
我們可以做.
讓我們做就好了.
又或者很容易去到另一個極端.
相反.
就是.
跟神說.
我不行的.
即是特別那些可能你是不想做的事.
不知是不是可能是使訪或者其他事.
就是不行的.
上帝你找另一個.
另一個適合些.
我不是很適合的.
找其他人.
這些都是很人性的狀況.
我們人很多時候就是這樣.
沒錯.
如果我們看摩西.
其實他起初都是這樣.
他起初都是.
覺得自己可以.
或者覺得自己不行.
不過.
如果我們仔細再看聖經的時候.
我們見到摩西的轉變.
而這個亦都是摩西之後.
最後那40年的階段.
我們見到摩西.
他不再著重他行不行.
他行.
和不行.
都不重要.

$^{521}$最重要他知道.
神能救.
當我去看出埃及的神呼召摩西的時候.
帶領以色列人出埃及.
剛才說的.
摩西怎樣回應.
我不會說話的.
即是你不要找我.
所以神就差阿倫.
去代摩西去跟百姓說話.
剛才是出埃及的第三章.
但看到第六章的時候.
摩西又是這樣說.
摩西要去見法老.
摩西就跟神說.
看看以色列人上邪都不聽我.
法老怎會聽我呢.
摩西又說自己不會說話.
然後去到三十節.
摩西又再說一次.
即是神知道他不會說話.
他就再說.
看看我真的不會說話.
法老不會聽我的.
又再說.
然後神就說.
我要使你好像在法老面前.
好像神一樣.
你的哥哥就做你的代言人.
又再這樣跟他說.
之後如果我們一直看聖經.
我們讀聖經的時候.
我們知道.
神藉著十災去擊打法老.
而每一次.
神耶和華對摩西說話.
吩咐他對法老說什麼.
的時候.
摩西都這樣去說.
但如果看十災.

$^{561}$當中有一個發言.
原來在頭三災裡面.
都是耶和華.
曉喻摩西你對亞倫說.
然後亞倫就跟耶和華.
摩西跟他說的話去做.
就照耶和華所說去做.
不過當我們看的時候.
到了第四災開始的時候.
發現耶和華是對摩西說.
是直接.
叫摩西跟法老說話.
他沒有再中間經亞倫.
而摩西也沒有再說.
你找第二個吧.
你找我不懂說話的.
摩西就照著耶和華的吩咐去走.
我相信摩西在過程裡面.
他看到神的同在.
他切切地經歷到.
神帶給他的能力.
以致他有勇氣的.
有膽量的.
從沒有信心.
到他有信心有膽量.
去對著法老說話.
如果我們接著下去看.
剛才說我們今天不是說十災.
如果大家回去看聖經.
看那十災之後的七災.
很精彩的.
摩西對法老的對話.
在當中你會看到.
摩西絕對不是一個絕口笨舌的人.
相反.
他是一個非常能言善辯的人.
弟兄姊妹.
有時候我們的信仰.
都可能是這樣.
大家可能回到教會一段時間.

$^{601}$不知道大家會不會曾經.
可能對神抱有些懷疑.
會覺得.
是不是真的.
舉個例子.
我們經常都說.
神很愛我們.
但是我們是不是真的.
很經歷到神的愛呢.
當我們遇到一些.
不好的事的時候.
會不會我們都會覺得.
神你是不是真的很愛我們.
為什麼會這樣經歷.
又或者.
聖經裡面應許.
上帝是公認的神.
祂是耶和華耳.
但是當我們打開戶口.
打開銀河.
上帝是公認的.
但是我好像很缺乏.
弟兄姊妹.
有時候.
會不會心裡面都會有些懷疑呢.
因為.
原來信仰.
是需要經歷的.
如果我們只是有.
頭腦上的知識.
是不夠的.
有時候和一些弟兄姊妹分享的時候.
可能分享到.
大家.
上一次經歷神的時候是何時呢.
大家都會好好去分享.
有些可能在說.
我上一次經歷神.
是我幾年前.
遇到怎樣怎樣的事.

$^{641}$又或者有些.
我經歷到神.
就是我信主那一刻.
上帝怎樣怎樣.
很好的.
這些都是一個很真實的經歷.
但是.
如果我們已經信了主一段時間.
可能我們要想一下.
為甚麼我們信主一段時間.
為甚麼我們的經歷.
好像那麼少.
可能我們的經歷是說.
很久以前.
甚至我們信主的經歷.
是不是.
上帝不讓我們經歷呢.
還是我們經歷了神.
我們不知道呢.
經歷是需要時間.
但是更需要建立那個關係.
剛才聖經說.
摩西和神有面對面的認識.
不是單單因為神揀選了摩西.
而在過程裡面.
我們見到摩西和神有很多的經歷.
有很多共同的經歷.
以致摩西可以在這些經歷裡面.
去認識神是怎樣的神.
回到今天.
可能大家會問.
我們怎樣可以經歷神呢.
其實沒有甚麼秘訣.
我相信大家上了教會很久.
都聽到很多傳道人或是講道人.
秘訣就是.
透過我們去讀聖經.
透過我們去讀神的話語.
去明白神的心意.
透過我們祈禱.

$^{681}$我們將我們的事情交給神.
我們抽時間去和神見面.
去認識神.
去和神建立一份關係.
以致當上帝去作供的時候.
或者去回應我們禱告的時候.
我們知道.
我們知道這是神的工作.
我們經歷到.
這是神的作為.
如果回到聖經的時候.
我們見到摩西和神有這麼親密的關係.
其他人好像沒有的時候.
究竟是不是神.
不想和百姓有這麼親密的關係呢.
不過我們看聖經原來有時候.
不是神不想和我們這麼親密.
是人可能不是想和神太親密.
我們看《昭希及記》第20章18-21節的時候.
當時上帝在西乃山頒佈十誡.
給摩西.
當時大家在經文裡形容.
以色列人走到很遠.
百姓見到雷軍閃電.
各星戰鬥.
遠遠站著.
他們和摩西說.
你和我們說話就行.
不要神和我們說話.
我怕我會死.
摩西和百姓說.
你不用怕.
神降臨要考驗我們.
要敬畏他.
不要犯罪.
但是百姓繼續遠遠站著.
不過摩西靠近神所在的幽暗當中.
我們見到.
這裡很有趣.
百姓說.

$^{721}$我遠遠站著.
你和我們說話就行.
神不用和我們說話.
如果我們再看另一段經文.
也是《昭希及記》第33章7-11節.
當中提到摩西在營外支搭帳棚.
這個教會幕.
不過大家不要誤會.
這個會幕不是tabernacle.
不是有聖所至聖所的會幕.
是說英文的the tent of meeting.
這是一個地方.
當凡求問耶和華的時候.
摩西就去到這個會幕當中.
當摩西進去的時候.
我們見到聖經形容.
雲柱降下來.
耶和華就和摩西對話.
不過其他人就站在自己的帳篷門口.
聖經形容.
當神和摩西對話的時候.
是好像朋友一樣.
大家和朋友聊天是怎樣的.
有來有往.
不是匯報式.
我今天做了什麼.
之後就完.
是大家你一句我一句.
是有來有往.
我們見到摩西和神就好像朋友一樣對話.
相反那些以色列人就站在自己的門口下拜.
弟兄姊妹.
剛才我們說.
當中我們見到摩西可以和神有面對面的認識.
不單止是因為神認識摩西.
揀選摩西.
而是摩西他都很想去認識神.
去親近神.
有時我們會想.
會不會有時我們好像認識的以色列人.

$^{761}$可能我們都不是很想和神太親近.
保持少少距離都可以.
因為太親近.
好像什麼都讓對方知道.
有些保留好些.
我曾經聽過有人分享.
真的這樣說.
和神我都不要太親近.
因為我覺得自己有時做的事不是那麼好.
和神太親近.
好像讓神知道好像不是太好.
所以很多時候可能都會回來崇拜.
回來聽道.
甚至可能都會參加聚會事奉.
不過在他生命裡面.
有些位置.
他從來都不會讓神觸碰.
他對神有些保留.
原來有時我們和神的距離.
不是因為上帝不想和我們很親近.
而是可能我們有些位置對神有保留.
我們不是很想和神很親近.
回到聖經的時候.
摩西可以和神有面對面的認識.
剛才說除了他和神有經歷.
他很想去親近神之外.
我們在聖經裡面看到.
摩西的為人.
都很值得我們去學習.
剛才文數第十二章都讀過.
這件事前面的背景是說.
摩西的姐姐米利安和哥哥亞倫.
批評摩西娶了一個古實女子為妻.
他們就說摩西.
難道神只和摩西你講話嗎.
為何不和我們講話.
當時聖經形容.
神在雲柱降臨.
停在會幕中.
神就向摩西出頭.

$^{801}$神就對著米利安和亞倫說.
你們聽我說.
你們中間若有先知.
我耶和華必在.
異象中向他顯現.
在夢中和他說話.
不過摩西不是這樣的.
他在我全家是盡忠.
我和他是面對面說話.
我是清清楚楚不用謎語.
甚至看到他的形象.
你們為何批評摩西而不懼怕呢.
當中我們看到.
上帝給摩西的評價是.
他是全家盡忠.
如果我們看回前一點.
這節的經文.
我想大家都有讀過.
在民數記裡面的十二章.
聖經有一個經文.
這樣評價摩西.
摩西為人極其謙和.
勝過地上的任何人.
我們看到.
上帝對摩西的評價是.
他盡忠 他很謙和.
如果我們看回摩西五經.
其實當中有很多句.
提及耶和華吩咐摩西.
曾經有人數過.
起碼有六十五次.
摩西就照耶和華吩咐的去行.
如果大家去看利未記.
大家知道利未記裡面有說很多.
頒布一些例典章.
真的很複雜 很長.
另外會幕點起.
尺寸是怎樣.
當中都寫得很清楚.
我們看到摩西.

$^{841}$如果看回經文.
摩西照耶和華吩咐的去行.
不知道大家有沒有曾經玩過.
一些遊戲 以訛傳訛.
你講一句說話給對方聽.
對方又要講給對方聽.
很多時候第一個講到第五個.
可能已經整句說話.
已經走樣了.
都試過 是不是.
但我們看到摩西不是.
耶和華怎樣吩咐他.
他就是怎樣去行.
我們記在米利巴水的時候.
因為當時的以色列人埋怨沒有水.
所以摩西去擊打了盤石兩次.
最後因為這件事.
進不了加拿大.
不過雖然我們看到.
摩西進不了加拿大.
但他沒有埋怨神.
他最後都是信服神.
大家可以想像一下.
摩西當時.
如果我們現在去想.
他是一人之下 他是上帝的代言人.
所有人都要聽他講.
他理應可以去驕傲.
他可以覺得他自己很厲害.
不過聖經不是這樣形容他.
聖經形容他是.
他極其謙和.
他很謙卑的去聽從神.
但我們今天看聖經.
可能我們覺得.
他是神人摩西.
他當然可以和神.
有面對面的認識.
不過剛才我們很粗略.
我們去看回摩西的一生的時候.

$^{881}$其實摩西和我們.
都是有一樣的性情.
不過.
摩西他願意去認識神.
他願意去親近神.
他願意去信服神.
以致他和神有一個很親密的關係.
弟兄姊妹.
今天你願不願意和神建立這種關係呢.
一開始的時候.
我講人和人成為摯友.
是要多少個鐘.
是需要200個鐘.
去建立那份關係.
今天.
我們和神建立關係的時候.
我們又用了多少時間呢.
雅各書四章八節.
提到.
你們親近神.
神就必親近你們.
聖經裡面應許當我們每一次親近神的時候.
神必然會親近我們.
就好像你約個朋友.
你不需要擔心.
我約他 他會不會不應約呢.
神不是這樣.
神不會說沒空.
上帝是在等.
等我們去親近他.
帕望.
藉著摩西的經歷.
弟兄姊妹都成為我們的勉勵.
我們抽時間.
我們真的同時.
建立那份親密.
個人的關係.
讓我們真的更深的去認識.
我們一同的祈禱.
賜下你自己寶貴的話語.

$^{921}$昔日你藉著摩西.
你讓我們知道.
上帝你是很願意.
很願意讓我們親近.
求你幫助我們.
幫助我們在我們信仰的路裡.
我們更深的去.
經歷你.
我們更深的去認識你.
我們更深的去愛你.
你也幫助我們.
在我們生命裡.
我們真的將你放在我們生命當中.
讓我們真的知道.
我們.
你可以成為我們隨時的.
依靠 因為我們知道.
我們來了你 我們什麼都不能做.
求主讓你的.
話語成為我們的激勵.
讓我們真的.
在我們的生活裡.
我們常常記起上帝.
你是那位與我們同在的神.
當我們親近你的時候.
你必然去親近我們.
願我們各人都在你裡面繼續去學習.
繼續去成長.
讓我們的生命更深的去認識你.
聽我們祈禱 奉主耶穌基督的名字祈求.
阿們.
\newpage



\section{}
\label{sec:luI8_7RFevk}
\textbf{2024-05-26 早堂 樓恩德牧師《耶穌復活的真諦》}
\newline
\newline
連結: \href{https://youtube.com/watch?v=luI8_7RFevk}{\texttt{https://youtube.com/watch?v=luI8\_7RFevk}} ~~~~ 語音日期: 2024-05-28
\newline
\newline
\hyperref[sec:Z6ZZRqsI0K0]{\small{< < < PREV SERMON < < <}}
~
\hyperref[sec:index_chronic]{\small{[返順時目]}}
~
\hyperref[sec:index_scriptual]{\small{[返順卷目]}}
~
\hyperref[sec:QGag1Zr1Cdk]{\small{> > > NEXT SERMON > > >}}
\newline
\newline
$^{1}$各位弟兄姊妹平安.
我們一起祈禱.
仍然有生命氣息.
能夠來到你的面前.
求主賜福.
讓每一個尋求你的人.
能夠尋見.
等候你的人能夠重新得力.
求主恩代使我們演下.
現在每一位.
無論我們口裡的言語.
心裡的意念都能夠夢你的月亮.
奉耶穌名次祈求.
阿們.
神愛世人.
但神是否閱立所有人呢?.
希伯來書說.
人非有信.
是不能得神的喜悅.
因為到神面前的人.
必須信有神.
並且相信.
祂會賞賜尋求祂的人.
等於知會今天早上.
我相信大家到這個地方來一同聚會.
多多少少.
不敢說很多.
都應該有些信心.
如果你沒有信心.
坐在這裡.
剛才主席和大家讀了聖經.
你不知道讀了什麼.
待會我可能會和大家再讀一次.
為什麼呢?.
因為你到神的面前來.
除了神的話語之外.
我請問你還可以信什麼?.
聖經告訴我們.
這個世界有些東西是看得見的.
有些東西是看不見的.

$^{41}$看得見的.
看得見的是暫時的.
看不見的是永遠的.
看得見的是暫時的.
看不見的是永遠的.
問題是你信耶穌.
是要得看到的東西.
還是得看不到的東西.
這是關鍵.
我們早不久之前.
吃過粽子.
我突然將師兄拉到你去吃粽子那裡.
那個叫什麼節?.
端午節.
接著我們要吃什麼?.
月餅.
那個叫什麼節?.
中秋節.
香港人端午節又放假.
中秋節又放假.
中間夾了個節你都放的.
復活節.
好東西.
一疊菜.
我想問一個問題.
這個復活節你怎麼過的?.
過節而已.
你過得恰不恰當?.
以你今時今日住在錦繡花園.
我每次來錦繡花園.
回來的時候.
有個特別感受.
因為我五十年前住在這裡.
五十年前.
我看上去很多.
當時跟我差不多年紀.
現在都不在這裡.
有些病了.
真的.
所以我見到大家好笑.

$^{81}$五十年前在家裡有祈禱會.
五十年前.
上帝讓我在錦繡花園學習.
怎樣接待孫教士.
蒙福.
上帝讓我在這個地方.
早上祈禱的時候.
見到一些我從來不發覺的東西.
我從來不發覺的東西.
所以我每次回來.
都好有些感動.
最令我感動.
就是我進來的時候.
我是在房客的車上.
去哪裡?.
我去教堂.
很少現在你去香港.
去任何一個屋苑.
說我去教堂.
是不是?.
你去看看全香港.
那麼多豪華極都好.
他沒有教堂.
我們有教堂.
只有這樣.
你不知道故事.
就看回你的歷史.
是神的恩典.
坐落教堂的地方.
在湖邊.
懂不懂得善用這個地方?.
不懂得善用.
像我這樣後來搬出去之後.
那時候真的挺好.
那時候挺好.
我那時候沒有經常.
像現在這樣我的腿已經有了.
神運.
運動一下.
不動就不動.

$^{121}$很貴的.
那時候我買單車都沒有機會.
很忙的時候.
今天和大家看回復活節.
復活的時候.
我今天的題目叫.
耶穌復活的真體.
不是你不知道.
而是你.
不要說你.
那幫門徒.
那幫愛耶穌的婦女.
那些仕徒.
都沒有一個知道耶穌的復活.
今天你知道耶穌復活.
是知道耶穌的復活節.
不過.
要搞清楚.
到底那是甚麼事情.
你搞不清楚.
你把耶穌復活當成拉薩路復活.
你就大錯特錯了.
應該吃粽子.
吃月餅.
吃粽子.
行不行.
行 怎麼不行.
我過中秋節的時候.
我吃粽子.
我紀念屈原.
我吃粽子的時候.
粽子裡面.
有沒有餡料.
有條紙.
不同的東西.
耶穌復活之前.
已經跟門徒說了一清二楚.
一而再再而三.
我要受死.
去耶路撒冷.

$^{161}$被由.
祭司長文士他們出賣.
不要錯了.
陷害了.
然後三天之後我要復活.
說了一次兩次三次.
說到有一次耶穌說的時候.
你猜猜彼得怎麼回應.
不行耶穌.
你不要搞錯.
不行耶穌怎麼罵彼得.
你記不記得耶穌怎麼罵彼得.
你這麼關心我.
撒旦退我後面去吧.
說得清清楚楚的.
他要死.
三天之後復活.
剛才經有讀到.
天使再提馬利阿們.
門有兩三個馬利亞.
其中一個馬利亞.
就是莫達拉馬利亞.
這裡出過問題.
所以看完聖經.
就發覺不要說我們今天.
不知道復活是什麼.
我們只知道.
復活節假期前坐地鐵.
聽女兒兩個人聊天.
我坐在旁邊聊.
我都聽到.
這個復活節.
你怎麼過.
你又說什麼.
他不知道是北上還是日本.
我忘記了.
你怎麼過.
復活節那天應該是慶祝.
還是應該是.
哭得七彩.

$^{201}$跟你有什麼關係.
今早希望我不會過時.
就將.
很多經文給你讀.
讀完之後你自己回去想再看.
希望.
這個復活跟你有關係.
如果這個復活與你沒有關係.
弟兄姊妹.
你信耶穌得永生.
是不到的.
只是羅馬書第十章第九節.
你背兩遍你就知道了.
我背給你聽好不好.
羅馬書十章九節.
你若考慮以耶穌為主.
就必得救.
再說一次.
你若考慮以耶穌為主.
就必得救.
對.
打開看.
不敢碰嗎.
不要拿手機出來.
手機一定有叮叮噹噹.
其他東西出來.
羅馬書十章九節.
很多人以為.
信耶穌就得永生.
太簡單了.
是很簡單的.
不過這個恩典裡面帶一個條件.
有沒有找到.
找到就大聲讀一讀.
給大家聽好不好.
不要只是我說了.
一二三.
你若考慮以耶穌為主.
心裡相信.
他從死裡復活.

$^{241}$第二句去哪裡.
心裡相信.
什麼.
他從死裡復活.
不是死裡翻生.
弟兄姊妹.
拉撒路就是死裡翻生.
耶穌是死裡復活.
這個活.
唱一首詩歌.
活什麼.
你當耶穌死了.
跟他過節.
耶穌死了.
你可以很受苦.
很不開心.
很不快樂.
很不平安.
很沒有盼望.
但是弟兄姊妹.
耶穌活了.
活到今天.
活到將來.
聖經中的讀.
我和你.
如果在基督耶穌裡.
都一定會復活.
然後耶穌基督的榮耀身體.
這樣.
所以今天早上.
我希望用這句話.
所以今天早上.
我希望用這段時間讀聖經.
剛才.
《盧家輝文》24章.
我想可不可以很快.
再讀一次.
因為我給弟兄姊妹看幾個字.
七日第一日.
黎明的時候.

$^{281}$婦女帶什麼去到.
墓穴那裡.
我告訴你.
一會兒我就.
舉例說.
我會入院.
我會痊癒.
我出院的時候你帶什麼來.
帶香膏來.
帶壽衣來.
我說這些話.
說得粗俗一點.
不過你明白.
婦女們帶著.
耶穌已經被我師父包緊了.
他們想做好一點.
帶些香膏.
去墓穴再去高一次.
他們信不信耶穌復活.
我說我會出院.
家人帶什麼來.
不是帶壽衣給我.
帶回我應該要穿衣服.
是不是帶錯東西.
他們愛不愛耶穌.
愛不愛.
你愛不愛耶穌.
他們根本不信耶穌復活.
好在下去.
石頭滾開再下去.
他們不信耶穌復活.
石頭滾開再下去.
看不到耶穌的身體.
有些聖經新翻譯.
看不到耶穌的遺體.
身體還好.
就是說.
耶穌被釘死的那一幕.
他們親眼看到.
這是事實.

$^{321}$不是魔術幻術表演.
看不到.
然後掉下去.
好困惑.
好多人對耶穌的復活.
覺得好困惑.
困就是不容易解釋.
或受到迷惑.
有兩個人站在旁邊.
這個不看了.
好害怕.
耶穌已經告訴我復活了.
我沒病.
我出院了.
你怕什麼.
兩個人.
對他們俯伏在地.
為什麼在死人中找活人.
說到活人.
死和活的分別是什麼.
死就死了.
活的不單會活.
耶穌的活不是像小孩一樣.
跳出來.
活到今天都活著.
再下去.
他不在這裡已經復活了.
要記得.
還在加利利的時候.
怎樣告訴你們.
耶穌說過給你們聽.
你今天在聽就聽過了.
耶穌復活已經是歷史故事.
還聽過耶穌說還要再來.
他已經說了.
不過這班婦女.
記不記得.
不記得.
今天很多基督徒.
歷史還記得.

$^{361}$但我認為將來還要再來.
不記得了.
你記得不會這樣生活的.
你記得去見耶穌的時候.
不會帶著壽衣去見他.
再下去.
他說.
人子必須被交載在人手裡.
釘在十字架上.
第三天復活.
這個聖經除了這裡記載之外.
馬太也有記載的.
這班婦女們.
門徒們 使徒們.
是聽過的.
不過.
沒有裝載在心裡.
不相信.
當信主耶穌.
你和人家一家都必得救.
你信不信呢.
不是你做不做到.
你做不到的.
他才做到的.
只要你信他 他就跟你行事.
在錦瑟花園幾個家庭.
我親自聽過的見證.
一直都不肯去教堂.
到最後流離的時候.
信了.
幾個頂智會做到這樣的見證.
問題就是.
今天我們的信仰不是叫你做什麼得到.
是要你信什麼得到.
說了一句話.
他們不信.
他們見到耶穌釘死那些慘劇.
見到他躺在墓穴裡.
他見到的.
不見的東西他們就信.

$^{401}$看不見 但是耶穌說的話.
他們不信.
這是不是我們的光景.
我們見到世界上的東西.
哇.
這裡亂亂 那裡亂亂.
這是天災 這是以巴 以哈戰爭.
擔心那麼多幹什麼.
以哈戰爭還有一段時間.
你說不是 你想咒詛他.
我不是咒詛他.
阿不拉.
以撒的子孫和以撒瑪利的子孫.
都是神應許給他們的.
以撒瑪利是.
上帝叫下鴿.
你回到你老婆那裡.
讓她繼續糟蹋你.
不過明年今天你會生個以撒瑪利.
以後說要跟他改名.
以後說改名的人.
你會滅絕他 行不行.
我不是替政治.
說話.
我只是就聖經立場看.
神愛世人.
但你沒有見到異象.
是尼波格.
尼撒王.
見到異象.
不過不懂解.
但你見了解了.
神都一對恩對很多人.
愛世人.
他們就想起耶穌的話.
這樣才想起.
你什麼時候才想起神的話.
你沒有災難.
沒有痛苦.
你不會想起神的話.

$^{441}$我沒有懷疑.
我有很多同學.
以前做校長的時候.
學生.
他不遇到難題.
有一個同學聊天的時候.
他要開大刀做手術.
我不敢跟他說醫學的事.
因為那些事不是我專長.
我只問他一件事.
你手術前.
你明天要做手術.
你這一陣子能不能睡覺.
他沒有回答我.
我說那我和你祈禱好不好.
他不出聲.
多剛硬.
那些學校學生.
很厲害.
那我和他再祈禱.
他又沒有出聲.
到祈禱完了之後.
我臨修煉之前說了一句話.
如果這一刻你還睡不著覺的時候.
你自己試試祈禱吧.
第一次聽到這句話.
第二天.
我不敢打電話.
又不敢WhatsApp他.
大手術.
誰知道他到黃昏的時候.
發一張照片給我看.
他和他主修醫師一起拍照.
他睡在床上.
主修醫師在旁邊拍照給我.
我心想這麼快就.
不是翻生.
那一副很瘋狂.
這麼快就醒了.
我知道他醒了.

$^{481}$因為我和他看的醫生.
也是我的同學.
他信得過他.
那我就.
馬上WhatsApp給他.
當我WhatsApp給他的時候.
他有回應給我.
這次我好像經過一場.
大考驗.
我發覺原來很多事情.
都沒有意思.
開始醒覺.
那我接著就WhatsApp給他.
我不敢和他對接.
我說.
昨晚的睡眠如何.
誰知道他怎麼回覆我.
他沒有回覆我.
我心想他不回覆我算了.
我又不是.
昨天一定很靈.
第二天.
我去探望他.
是我以前的同事.
他將他患病前的.
一些掃瑜的經歷.
全部告訴他老師.
他老師聽了之後很開心.
馬上打電話給我.
告訴我.
我心想豈有此理.
原來他昨晚有祈禱.
又睡了覺.
但他不告訴我.
但我不會生氣.
我知道這個工作要做很久.
讓他親自經歷.
你祈禱吧.
神會聽的.
弟妹朋友.

$^{521}$你聽著.
你有困難嗎.
你向他祈禱吧.
他活著的.
他不是死了的.
死了他聽不到的.
他活著他聽到.
他會替你解決問題.
所以.
想起耶穌的話.
從墳墓回去.
將這一切的事.
告訴11個使徒.
和其他人.
開心嗎.
很開心.
然後.
這班名字不提了.
這些名字大家以後再看.
他們這些說.
使徒以為是.
什麼叫胡言.
亂說.
我估計.
其中一個是莫迪拉瑪尼亞.
如果門徒見過他.
以前的景況.
都以為他舊病復發.
知道他說什麼嗎.
知道就好 不知道就回去看聖經.
是不是.
你見到幻影.
你思覺失調.
再下面.
彼得聽了之後.
這個.
心靈中人.
非親眼見過都不相信.
他立刻衝過去.
俯身一看.

$^{561}$只見世麻布.
什麼人 世麻布.
不重要 只見.
不是瑞衣.
只見壽衣不見人.
這樣你明白了.
是不是.
因為發生的事.
而心裡.
什麼.
信 驚.
什麼叫驚.
這樣都行 我舉個實際例子.
你會不會驚.
你看過那些幻術表演嗎.
一塊布一蓋.
不見了 我弄給你看.
我都會弄.
這塊布.
這支筆在這裡.
你只是笑.
你信不信.
如果你要拍掌 你拍不拍掌.
拍掌就不拍掌.
你見到魔術表演你不拍掌.
拍掌.
哇 厲害.
你信不信.
很多人信耶穌就是這樣.
拍掌 厲害.
不信.
我希望 你不要弄錯.
我真的不懂這些.
我手法.
都不見了.
是有拍.
真的在這裡.
我拍掌就不拍了.
耶穌不是快手.
弄東西給你看 不是.

$^{601}$心裡驚訝 仍然不信.
再下去.
接著有兩個門徒.
同一天.
二千多年前.
耶穌復活那一天.
上午就發生了事件.
人們談得七彩.
同一天 有兩個門徒.
向村子走.
那個地方叫二萬五十.
離開耶路撒冷約二十五里.
二十五里大概等於.
十一二公里.
如果要人的出道.
距離大概要走三個小時.
到四個小時.
一個人步行.
大概走得快一點.
可以一個小時走五公里.
或者走四公里.
四公里只是十五分鐘.
走一公里.
你試一下 從哪裡走.
可以計算得到.
所以三四小時的距離.
當一走的時候.
你猜那兩個人做什麼.
他們不是說去哪裡消費.
不是說去哪裡遊覽.
是回鄉下.
因為那個是非地在耶路撒冷.
哇 整個耶路撒冷.
大家都知道了 知道什麼事了.
有個人叫拿撒拉人耶穌.
他有行神蹟.
有騎士 講得很恰當.
是好人 我們盼望他.
來拯救我們尼徹亞.
誰不知.

$^{641}$大祭司那些將他釘死了.
判給羅百人釘死了.
唉 正經點而已 看下去.
他們看似談論.
所發生這一切的事.
都是在講耶穌的事.
不過心裡不信.
已經離開耶路撒冷是非之地.
要走了.
再下去快點.
正義與談論是.
耶穌親自走近他們 和他們同行.
耶穌和他們同行.
他知不知道 下一章.
可知他們的眼睛.
模糊了 沒認出他.
可能他們兩個年紀大.
白內障.
等一會.
人會模糊的.
看得見的當作永遠的.
看不見的永遠的.
當作是暫時的.
你說是不是將他是非顛倒了.
害怕了嘛.
就離開耶路撒冷.
結果後來怎樣.
我想很誇張 耶穌說.
你們談論什麼事.
彼此談論什麼事情.
他們就站著 走著走著.
老兄你都不知道.
面上戴著什麼.
你現在是喜樂的樣子.
你照照鏡子.
相由心生.
你的心.
開心就開心.
你的心不開心.
你要扮開心.

$^{681}$那個叫什麼笑.
你怎麼知道.
心裡很苦.
面帶愁容.
沒有盼望了.
彌賽亞被人釘死了.
他們盼望的那位救世主.
被釘死了.
不是彌賽亞被釘死了.
第十八節.
媽媽 沒了.
對不起 媽媽.
回去不要看了 很好看的.
於是他們走著走著.
越聽越心裡火熱.
不是火滾 心裡火熱.
聽到後來日落平息.
正午走.
走到黃昏.
這兩個人回到家裡.
老兄你先不要走.
留下來 叫他留宿.
耶穌就留宿.
一坐下吃飯.
開始祈禱的時候.
餅一打開 第二個去.
聖經記載.
三個人 怎麼只剩下兩個.
不見了.
復活的耶穌是這樣的.
可以出現 可以不見了.
於是這兩個人就.
聖經記載 馬上.
回到家裡.
飯都不吃.
馬上回到耶路撒冷.
再走兩三個小時.
我們急不及待.
走回去.
回到家之後 告訴門徒.

$^{721}$聖經有些門徒聽了之後.
哦.
大家對口真的復活了.
你猜這幫人信不信.
看到了.
見到了.
仍然不信.
我今天想和大家.
透過這些經文.
接著下去 路加文24章.
我將會.
我將會分享幾個大綱.
不是這裡.
耶穌復活那一天.
門徒的光景是怎樣.
沒有一個信.
不記得整個笑話.
他們憑眼見 看到耶穌被人釘死.
看到的東西.
觸目驚心.
影響了他們的信心.
他們不信.
耶穌復活之後 前後曾經有一個門徒.
講過什麼話 待會和你讀一次.
讓你知道.
當時耶穌復活 未死之前.
所講的話 他們不信.
耶穌講過之後的說話 他們信不信.
這個讓我看 第三個.
耶穌復活對你和我信仰有什麼影響.
好 接著下去 經文.
24章.
6到9節.
剛才讀過了.
24章11到12節.
剛才都讀過了.
24章的17.
20 23節.
這裡是講到.
當時門徒的警方.

$^{761}$惶惶驚.
甚至乎.
是憂愁.
沒有一個是表示.
雀躍歡喜 哎呀.
他真的康復出院了.
不是.
還是在見到他.
入院之前.
插了喉管.
接著下去 聖經.
害怕.
以為這是當.
耶穌再出現中間的時候.
他們見到耶穌之後.
不是開心 是害怕.
以為他暈了.
於是耶穌示範.
講話給他們聽 願你們平安.
見到他們還這麼驚暈.
好.
拿一條魚 有沒有東西吃.
問他拿東西吃.
吃一條烤魚 吃給他看.
他能吃到身體嗎?.
能吃到魚嗎?.
能吃到.
頂智滿 這是37節所講.
好 接著下一章.
下一章.
耶穌基督在他復活之後.
曾經對馬尼阿們.
和他們說過一句話.
我還沒有.
升上去見我的父.
當時馬尼亞想.
抱他.
他說 先別這樣.
我還沒有見我的父.
耶穌還要升天.

$^{801}$還要回去父神那裡.
因為他從上頭而來的 回到上去.
下一章 拜託.
耶穌對他們說.
願你們平安.
父怎樣差遣了我.
我也照樣差遣了你們.
這是耶穌復活之後說的.
耶穌復活之前說的那些話.
門徒們.
聽不聽過? 聽過.
信不信? 不信.
到耶穌復活之後.
他說 願你們平安.
父怎樣差遣了我 我也照樣差遣了你們.
這是耶穌說的話.
結果這班門徒們信不信?.
到聖靈降臨之後.
他們接受了差遣.
今天詩歌有唱.
在這個地方你居住.
神保你這麼好的環境.
有一個有教堂的屋苑.
裡面居住.
是你的機會.
將人帶回來這個地方.
不是單單你的聚會地點.
帶回來讓他們認識神.
弟兄姊妹.
我們看不到我們.
聽到的我們不信 我們不會做.
你聽到的信.
我下一章再給你看.
28章.
耶穌進前來對他們說.
天上地下.
所有的權柄都似給我了.
下一章.
所以你們要去奉父子之名.
給他們施洗.

$^{841}$下一章.
凡我所吩咐你們的.
都教訓來遵守.
我尚願你們同在.
直到世界無聊.
你以下的時間 連日.
你希望誰跟你同行.
你整天看著看著這些新聞.
看見張 看得見.
這裡打仗 那裡瓦蘭.
那裡加沙 天天看.
你生氣都生氣了.
有沒有什麼幫助.
為什麼不想回.
就像基督所說 到末世的時候.
各要攻打各民 要攻打民.
時候就到了.
你不準備.
你和世界門徒.
復活節那天早上一模一樣.
和二萬五十路上兩個門徒.
一樣.
所以.
跟耶穌三年.
緩分都是這樣的結果.
問題在於.
你信不信聖經的說話 下一章再下去.
24章.
你回去再看.
不說這章.
好了 我看看.
哥林多傳書15章.
普羅話 我當日傳給你們的.
是重要的.
就是照聖經所說.
耶穌為我們的罪死了.
再下去.
而且埋葬了 又照聖經所說.
一起讀.
第三天復活了.

$^{881}$福音就說這些.
耶穌為我們的罪死.
被埋葬.
使徒信經都這樣說.
然後說第三天復活了.
弟子們.
很多人.
包括.
信徒做見證.
不敢說耶穌復活.
通不過這裡.
什麼叫復活.
他現在還活著.
你唱就懂得唱.
活著活著 活什麼.
你有難處.
他知不知道.
你病了 他知不知道.
我們將.
理性 將我們的感受放在.
看得見的事情上.
看不見上帝的英雄.
你漠視.
怪不得真不想到你.
再下去.
顯及希望.
再下去.
我不說你翻譯沒問題.
五百多人看.
都有.
再下去 麻煩.
再下去.
到最後都獻給我看.
他說他自己是一個.
產其未到出世的英雄.
保羅說自己.
甚至保羅說自己是罪魁.
因為他曾經被北戈教會.
所以以前做什麼.
弟兄姊妹 過去的歷史.

$^{921}$耶穌去了北戈教會.
今天迴轉歸向他.
仍然蒙恩 再下去.
而且神的恩典.
使我成今日我.
並且所賜給我的恩典不是徒然的.
我比眾信徒更老.
更愛勞苦.
不是我 是神的恩典.
與我同在.
是我們認識耶穌的恩典.
再下去.
下一章.
無論是我.
或其他使徒 我們如此傳.
你們也都如此信了.
我們今天相信.
耶穌復活 因為傳到你的福音.
聖經告訴我們我們相信.
還有你自己可以.
主觀經歷.
好像我跟學生說 你睡不著覺.
你祈禱吧.
真的睡著了 我聽過.
好幾個見證都是這樣.
信耶穌 再下一章.
若基督我有復活.
我們所信的.
所傳的是謊言.
白費力氣了.
胡說八道 只是神話.
你們所信的.
都是謊言 根本你信的.
謊言是無效的.
得不到永生的.
考慮認耶穌為主.
叫主也 主也的人不能都進天國.
要相信他從死裡復活.
下一章.
基督若沒有復活.

$^{961}$我們一起讀 一二三.
基督沒有復活.
你們的信就是土言.
你們應活在罪中.
罪人得不得救.
下一章.
頂門.
我要告訴你們的是.
血肉之軀不能承受神的國.
必扭壞的也不能承受不扭壞的.
身體會不會扭壞.
一定會扭壞.
但是裡面的靈更新之後.
靈魂甦醒之後.
我們裡面的身造的人.
不會扭壞.
下一章.
我現在把奧秘的事告訴你們.
我們不是都要睡覺.
這個睡覺的事.
很多時候是指死.
我們都不是要死 不是亡 是死.
是都要改變.
如果你不是很明白.
你看哥林多前書第十五章.
整章裡面.
說到最後.
什麼叫改變.
種一粒種子下去.
那粒種子生命出來的時候.
就是一棵大樹.
無無從爬下爬下.
結了繭出來.
就是一隻飛得的昆蟲.
自然間都有這樣的變化.
神要他大能不可以改變.
下一章.
就在一剎那之間.
浩同末次吹響的時候.
一定吹響.

$^{1001}$死人要復活.
成為不扭壞的.
我們也要.
什麼?.
改變.
人如果不是從水和聖靈生的.
是不能進神的.
重生是要靈裡面.
復甦的.
所以詩人說.
靈魂蘇醒.
靈魂沒有蘇醒.
你不是活在基督耶穌裡面.
詩篇23篇第四召.
你經過死亡的幽谷.
經過死亡的通道.
你一定遭害.
詩人說.
我不怕也不怕遭害.
因為什麼?.
你與我同在.
復活的主義你同在.
但你經過死亡的.
沒有復活的主義你同行.
你經過不了死亡的.
你死了就亡了.
復活的意思在這裡.
不是每個復活節.
一年吃一次粽子.
一年吃一次月餅.
一年吃一次.
是不是叫復活蛋?.
是不是?混帳.
吃復活蛋.
下一章請.
羅馬書一起讀.
1,2,3.
使耶穌從死人中復活的神的人.
若住在你們裡面.
使耶穌從死人中復活的.

$^{1041}$有意接著住在你們裡面的聖人.
使你們必死生還救我.
這個身體必死的.
我沒有指那一段東西.
我只是說我自己.
這個身體必死的.
根據聖經.
我這個身體.
會死的身體.
活過來.
怎麼活我不知道.
耶穌怎麼從同類瑪利亞生出來我不知道.
神怎麼創造天地我不知道.
但聖父工作我見到了.
聖子的降生我見過了.
我這個身體的復活.
耶穌的復活都顯示了一次給我看.
你信不信?.
再讀一次.
羅馬書八章十一節.
回去慢慢自己看.
時間不到我已經交時間了.
我還要去.
什麼堂?嘉恩堂.
我還要去嘉恩堂.
下一章請.
一起讀.
一二三.
基督使我們的心.
他已滅的時候.
你們也要與他一同做榮耀.
你顯現過沒有?.
基督有一次榮耀顯現.
當他顯現的時候.
你與他一起在榮耀裡面.
聖經是應許來的.
你信不信?.
你不信就相信他復活.
他復活你不信你自己怎麼會復活?.
和理性上分析.

$^{1081}$在信仰上真的七折八扣.
七折八扣即是多少折?.
五六折都沒有.
是不是?.
下一章請.
你最後一起讀.
一二三.
你等著什麼?.
第一次節目你回答我.
你等著什麼?.
人生是.
有一條路人以為正.
結果成為死亡之路.
心言講的十四章.
人人等.
是等那個字.
但我說不.
那個字我會經過.
但我等候主耶穌基督從天降.
祂來的時候我要復活.
見祂.
哎呀!你在講神話.
等於那些門徒.
耶穌三番四次說了.
我要死要埋葬.
三天之後復活.
發個瘋.
幸好.
耶穌知道這群人的弱點在哪裡.
等呢.
等到五十天之後.
聖靈降臨.
你們就必得能力.
翻生.
我三月三十一日.
我去一間教會講道.
那天是復活節.
我想不到.
不知為何教會安排了我.
五月十九日.

$^{1121}$再去講道.
我那天一上台.
那邊我還.
我上去給你看.
我問那些頂尖.
我問他們今天幾號.
他說五月十九.
我說.
五十日之前.
是什麼日子.
他們沒有人計得到.
我說.
十九日.
要加多少日.
才等於五十日.
計不計得到.
加三十一日.
四月份剛剛三十日.
三月三十一日.
最後那天.
又剛好復活節在那裡講.
到聖靈降臨節又在那裡講.
當我告訴他們的時候.
哇.
頂尖會.
聖靈降臨.
詩歌有唱過.
沒有聖靈你沒有真正的自由.
沒有聖靈你不會.
認耶穌為主.
沒有聖靈你沒有力量.
在這個世代裡經得起考驗和試探.
最後一章.
他要按著那使萬有.
歸服自己的大能.
將我們這個.
悲戰的身體.
改變形狀.
不用死了.
和他自己榮譽的身體相似.

$^{1161}$先敬敬他們.
下一章.
這章一起讀.
讀完之後我們結束.
(念).
不要打這章.
他們夠不夠膽講.
包括傳道同工.
包括我自己在內.
我讀完之後我死了.
哇.
美好的仗打不打得過.
我屢敗屢戰.
靠著這一年.
開始戰勝我也有.
當跑路跑不跑得順.
所順的道.
什麼道啊.
耶穌基督的降生我接受.
耶穌基督的受死我也接受.
耶穌基督的升天.
我半信半疑.
耶穌基督的復活.
耶穌基督的升天.
我沒見過不知道是不是.
一塊雲上去都不用穿太空衣嗎.
耶穌基督再來.
唉 有一段時間了.
這個世界多美啊.
弟子們 下面一句.
從此以後.
有公義的官員為我存留.
就是按照公義審判的主.
到了哪一日.
要賜給我的.
這裡之後是一個分號.
如果這裡寫一個句號.
保羅對提摩太太說的話.
已經說完了.
接著說 不但.

$^{1201}$一起讀 一二三.
不但賜給我.
也賜給凡愛母他天地的.
記住這句.
不然你沒機會.
不但賜給我.
保羅說 也賜給.
什麼啊.
愛母他顯現啊.
不是賜給那些.
擔心要死的人啊.
死亡通道個個都要過的.
如果沒有耶穌的同行 你走不過的.
但你愛母他顯現的時候.
弟子們 這是一個激勵.
讓你能夠在今生.
不要重蹈昔日.
那班門徒們 婦女們.
很愛耶穌那班人.
那些行為.
視而不睹 看而不見.
看見的.
作作永恆的.
今天我們需要神的面前.
重新調整我們的信仰.
希望.
明年國復活節.
不是國復活節.
耶穌復活那天 你重新思考.
你的信仰問題.
今天開始的時候 一起去祈禱.
讓我們今天早上.
透過昔日門徒們.
在耶女復活那天早上.
發生的事件.
給我們一個反省.
看到將來我們的光景如何.
求聖靈繼續不斷地工作.
好讓我們能夠在你的面前.
由心意更新變化.

$^{1241}$測驗何為神的善良.
來承傳你的旨意.
祈禱教託.
奉求耶穌明子祈禱.
謝謝 結構.
\newpage



\section{約珥書 2:28-32}
\label{sec:QGag1Zr1Cdk}
\textbf{2024-05-26 午堂《預備耶和華大而可畏的日子》約珥書 2\_28-32  講員:郭鴻標牧師}
\newline
\newline
連結: \href{https://youtube.com/watch?v=QGag1Zr1Cdk}{\texttt{https://youtube.com/watch?v=QGag1Zr1Cdk}} ~~~~ 語音日期: 2024-05-29
\newline
\newline
\hyperref[sec:luI8_7RFevk]{\small{< < < PREV SERMON < < <}}
~
\hyperref[sec:index_chronic]{\small{[返順時目]}}
~
\hyperref[sec:index_scriptual]{\small{[返順卷目]}}
~
\hyperref[sec:eejVMLNyY8I]{\small{> > > NEXT SERMON > > >}}
\newline
\newline
約珥書 2:28-32
\newline
\begin{longtable}{cl}
\hline
\hline
章節 & 經文 (和合本修訂版)\\
\hline
2:28 & \begin{tabularx}{0.7\textwidth}{X} 「以後，我要將我的靈澆灌凡有血肉之軀的。你們的兒女要說預言，你們的老人要做異夢，你們的少年要見異象。 \end{tabularx} \\ \\ \relax
2:29 & \begin{tabularx}{0.7\textwidth}{X} 在那些日子，我要將我的靈澆灌我的僕人和婢女。 \end{tabularx} \\ \\ \relax
2:30 & \begin{tabularx}{0.7\textwidth}{X} 「我要在天上地下顯出奇事，有血，有火，有煙柱。 \end{tabularx} \\ \\ \relax
2:31 & \begin{tabularx}{0.7\textwidth}{X} 太陽要變為黑暗，月亮要變為血，這都在耶和華大而可畏的日子未到以前。 \end{tabularx} \\ \\ \relax
2:32 & \begin{tabularx}{0.7\textwidth}{X} 那時，凡求告耶和華名的就必得救；因為照耶和華所說的，在錫安山，在耶路撒冷將有逃脫的人。凡耶和華所召的，都在餘民之列。」 \end{tabularx} \\ \\
[1ex]
\hline
\hline
\end{longtable}
$^{1}$我們一起同心祈禱.
慈悲人類的天父上帝.
因為你就是那位.
全能創造天地的上帝.
你對我們充滿恩典和憐憫.
而且你同時是一位公義的上帝.
就求你讓我們的內心.
被你的靈去觸動.
讓我們能夠看到.
我們自己現時的狀態.
我們和神的關係.
以及我們如何走我們前面人生的路.
無論我們現在是甚麼年歲.
甚麼狀態 甚麼處境.
求你都親自向我們說話.
讓我們看見.
我們在你的計劃之中.
是有你自己給我們的角色和作用.
願意我們真的能夠好好地.
活出上帝對我們人生的計劃和心意.
我們的祈禱奉耶穌基督的名.
阿們.
很開心和大家見面.
不過首先要說一說.
我做的PowerPoint是很素色的.
是沒有顏色的.
所以你們繼續為我們祈禱.
可以讓我再提升一下.
這個題目是.
你聽完之後好像有一點恐嚇性.
其實是為甚麼我會這樣說呢.
就是在我們.
麻煩我們的弟兄.
姐妹.
在今年2024年4月8日.
在北美洲出現了全日蝕.
不知道大家有沒有印象.
它是經過了墨西哥 美國和加拿大.
當然了.
最近兩個禮拜前.

$^{41}$應該是.
兩三個禮拜左右.
如果有些在英國的朋友.
弟兄姊妹 親戚.
都會做一些極光出來.
我又不是很清楚那些是甚麼東西.
不過我是比較注意日全蝕.
其實我不是天文學家.
首先我有興趣就追一下資料.
原來日蝕是分有三種.
第一就是日全蝕.
另一個叫日環蝕.
第三種叫日偏蝕.
日全蝕是甚麼意思呢.
月球遮住了太陽.
所以白天的時候太陽變成黑色.
日環蝕是怎樣呢.
你想一下.
如果我用圓形的東西.
遮住下面發光的東西.
旁邊的圈還有一個光環.
那個就叫日環蝕.
即是太陽的位置距離地球遠一點.
遮不住太陽.
日偏蝕是怎樣呢.
半月彎.
月亮遮住太陽一截.
原來在香港以前曾經出現過一個日環蝕.
有位弟兄對天文很有興趣.
所以寫了一篇文給我聽.
有沒有人在1958年的時候.
看過香港的日.
是 大家馬上有反應.
我還未出生.
好表示我們是年輕一代.
在1958年4月19日.
香港的上空出現過日環蝕.
而且有報紙報道.
這位弟兄說他的姐姐親眼見過.
第二次是2012年5月21日.

$^{81}$那天是天陰的.
應該不是很看得到.
下一次是甚麼時候.
是2320年.
到時讓我們下一代去看.
講回來.
為何講這件事呢.
原來在2023年.
大家聽過土耳其和敘利亞發生地震.
大家有沒有印象.
在2024年4月3日台灣都發生地震.
4月5日美國紐約發生地震.
大家有沒有留心日期的關係.
所以在英語的傳媒裡面.
他們提出4月5日美國紐約發生的地震.
和4月8日日環蝕.
是否有關係的問題.
炒作的東西可以有很多炒法.
我們當它是一個好的問題.
不過我是不懂得回答.
因為我不是研究這個方面.
其實目前沒有明確的答案.
不過科學界的人認為.
這個月頭會影響潮水漲退.
這個大家都知道.
有可能會影響到地震的出現.
都說不定.
這件事真的沒有人知道.
我真的不懂這些事.
我沒有跟進美國紐約發生地震.
和日環蝕是否有關係.
我在說什麼好.
其實我首先告訴你.
這個題目我沒有答案.
但我想的是反而是信仰的問題.
我關注的就是我們作為基督徒.
我怎樣從一個基督徒的眼光去想這些事.
包括地震和日環蝕這些大自然的現象.
我反而關心這個問題.
經過新冠病毒之後.

$^{121}$聯合國世界衛生組織宣告了.
人類是面對目前仍然未知的病毒和細菌.
將會構成一個廣泛擴散.
叫做X疾病.
這些可以查到.
下面小小的就是我的footnotes.
在哪裡找出來的.
我聽完這個真的很嚇人.
將會有種病.
不過現在我不告訴你是什麼病.
那種病是會廣泛流傳的.
現在經常有說.
哪種病毒會流傳等等.
《使徒行傳》第二章二十節.
記載了這段經文.
一頭要變為黑暗 月亮要變為圈.
這都在主大而可 大而明顯的日子未到以前.
這段經文是《使徒行傳》新約.
我們今天講到經文是舊約.
但你留心一下.
這兩段經文其實內容是非常相似的.
一頭要變為黑暗.
可以是什麼呢.
你想一下是什麼.
就是日全食.
這個現象就是.
其實聖經以前都有提過這件事.
這個現象提醒我們.
要預備迎接主耶穌再來的日子.
我自己在祈禱和思考之後.
我講到的題目就是說.
預備耶和華大而可畏的日子.
我用的經文是舊約.
《約珥書》二章二十八到三十二節.
講到有三部份.
第一部份是《使徒行傳》第二章二十節.
和《約珥書》第二章三十一節的關係.
第一件事看看.
新約《使徒行傳》和舊約《約珥書》的經文的關係.
第二部份就是說.

$^{161}$我們面對世界這麼奇怪的事.
我們要怎樣呢.
經文告訴我們我們要有異象.
大家可以去問這些問題.
什麼是異象.
我應該有什麼異象.
第三個就是說.
我們要信服聖靈的帶領.
我也嘗試跟經文的主題.
跟大家分享.
內容方面就是《使徒行傳》第二章二十節.
和《約珥書》第二章三十一節.
我用的是和合本的舊版本.
你們用的是新的和合本修訂本.
其實意思是一樣.
《使徒行傳》第二章二十節.
日頭要變為黑暗.
月亮要變為血.
這都在主大而明顯的日子未到以前.
到《約珥書》第二章三十一節.
日頭要變為黑暗.
月亮要變為血.
你聽到這兩句是一樣的.
這都在耶和華大而可畏的日子未到以前.
這兩節經文除了是《使徒行傳》用主.
《約珥書》用耶和華之外.
和《使徒行傳》用大而明顯.
《約珥書》用大而可畏之外.
內容完全一樣.
只是想給大家聽.
原來聖經是有這樣的特點.
彼得在五秦節的講道.
他引用舊日聖經的約珥書來講.
記載在《使徒行傳》第二章十六到二十節.
大家要記住.
彼得講原道很多人信主.
就是這段經文.
我們一起讀一讀.
讓大家可以投入講道的佈審.
一二三.

$^{201}$這正是先知約爾所說的.
神說在末後的日子.
我要將我的靈挑冠凡有血氣的.
你們的兒女要說預言.
你們的少年人要見異象.
老年人要做夷夢.
在那些日子.
我要將我的靈挑冠我的僕人和使女.
他們就要說預言.
在天上我要顯出奇事.
在地下我要顯出神蹟.
有血有火有煙霧.
日後要變為黑暗.
月亮要變為血.
這都在主大而明顯的日子未到以前.
到那時候.
凡求告主明的就必得救.
這段經文很有意思.
彼得向三千人講道的時候.
很多人都有反應.
究竟誰可以在主來到的日子.
我們可以承受得住.
究竟誰可以在這樣的情況下.
我可以站立得住.
史特蘭傳第二章38節.
彼得說你們各人就要悔改.
奉耶穌基督的名去受洗.
叫你們的罪得赦.
就會承受所賜的聖靈.
彼得講得很清楚.
有這樣的日子會出現.
那些事是會這樣發生.
你想想你可不可以應付呢.
如果不可以應付.
你今天可以怎樣呢.
彼得就說你要悔改.
奉耶穌的名去受洗.
還有你就可以罪得赦免.
你可以領受神應許的聖靈.
面對主大而明顯的日子.

$^{241}$未到以前.
凡是求告主的就必得救.
這句話你看看上文下.
是甚麼意思.
凡是那麼大可畏.
很可怕的日子.
未到之前.
你去求告上帝的時候.
你就得救.
這才是甚麼含義.
這才是不是到那時候.
真的發生的時候.
我才去求他.
是不是這個意思.
這個是提醒和預告.
你想想彼得當時是向一班人.
去提醒和預告.
而當時有一班人是有反應.
就真的一起去受洗.
我們今天怎樣呢.
我們會不會已經見到很多這樣的現象.
在我們傳媒新聞裡經常看到.
不過我是從來沒有想過.
這件事的現象和我的信仰有甚麼關係.
或者我從來沒有想過.
耶穌會有甚麼行動.
他大而可畏的日子會不會來.
或者會不會有甚麼其他的事發生.
或者我更加不會想我今天應該怎樣.
可能我們未去到這個地步.
所以我希望大家一起看這段經文.
不過我真的要告訴大家.
其實我是不喜歡港獨去嚇人的.
那些說耶穌快要回來.
你快點去信主.
不然你日後上不了天堂.
這些不是我港獨的風格.
我也不覺得要這樣去嚇人.
不過我也要這樣說.
聖經真的有這樣記載的時候.

$^{281}$我們又不能夠不去說.
我希望我講出來的時候.
不是用一種像煽情或者嚇你的方法.
而是用上帝的說話去告訴我們.
有些事是正在發生.
有些事是將會發生的.
而關鍵是甚麼.
就是我今天怎樣去準備.
就是這個位置.
凡是求告主明的.
首先我們一定要全心全意.
你想想怎樣去求告主明.
是很認真.
我會覺得這件事我搞不定的.
我不可以用我自己的智慧.
經驗和能力去處理.
面對一個將來我也不肯定的時間.
一個事件.
我只有謙卑的去求告上帝.
只有這樣我才能夠解決問題.
所以這一刻我能夠做的就是悔改.
我這一刻能夠做的就是奉耶穌的名去受洗.
可能在座已經是基督徒了.
已經受洗.
這件事對大家沒有甚麼作用.
但是經歷罪的赦免.
可能不是每個人都經歷的.
我不是讀輔導.
不是讀心理學.
但是人的罪疚感.
或者如果你拉闊一點.
我們童年成長的歷史.
創傷 陰影等等.
那些記憶.
其實可能在我們的心裡.
我們仍然未能跳出來.
也會有這些情況.
如果我們是經歷罪的赦免的時候.
其實你是應該明白保羅所說的.
我今天活著的不是過去的我.

$^{321}$我現在在基督裡面.
我是新造的人.
意思是甚麼呢.
他明明知道以前是怎樣的.
怎會不知道呢.
不過他現在不用被以前的東西綁住.
這個你和我有沒有經歷呢.
彼得就勸我們去想這件事.
還有你要領受神給你的聖靈.
這是禮物來的.
你有沒有領受這件事呢.
當然我不想說得太虛幻.
其實領受聖靈.
你可以讀路加和史滕轉.
你可以讀保羅的書信.
所以有兩個不同的層面.
一個就是你行神蹟 治病.
很適逢的能力.
很事工上.
但是保羅所說的聖靈.
你的聖靈的臨在和充滿.
是說你內在的生命.
你的生命得到更新.
你有很多承托的能力.
這些就有機會你們在主義學裡面上.
或者是其他的部份去學.
但是每個基督徒都需要經歷聖靈的能力.
聖靈的更新.
以致我們能夠有新的眼光.
看我們現在的世界.
約珥書第二章28到32節.
你會發覺我讀很多經文.
沒辦法了.
因為講道不是講我自己的東西.
講聖經的東西.
你試一下剛才讀了史滕轉.
你現在看史滕轉.
其實他引用約珥書二章28-32節.
一齊嘟一嘟那些味道 好嗎 我們又一齊讀了 好嗎.
好 123.

$^{361}$以後我要將我的靈曉冠 凡有血氣的.
你們的兒女 要說御言 你們的老年人 要做義蒙.
少年人 要有意象.
在那些日子 我要將我的靈曉冠.
我的僕人和使女.
在天上地下 我要顯出奇事.
有血有火有煙柱.
日頭要變為黑暗 月亮要變為血.
這都在耶和華大而可畏的日子 未到以前.
到那時候 凡求告耶和華明的 就必得救.
因為照耶和華所說的.
在石安山耶路撒冷必有逃脫的人.
在聖下的人中 必有耶和華所召的.
這樣你就看到 彼得他在《少林寺傳》說的時候.
他如何將《若意書》第二章28-32節的經文重點抽取.
這裡我不說那麼多 我怕你們會太過辛苦.
《若意書》第二章28-32節 講到上帝的靈會曉冠他的兒女.
少年人又有會說義言.
老年人有意象 有義蒙 中年人都有.
其實說這些是什麼含義呢.
我們現在是在說我們有沒有什麼願象.
有沒有vision.
這是人間的語言.
我們一定有願景才行.
沒有願景就很可憐.
但聖經裡面說的不只是人的願景.
而是神給我們的意象 神給我們的義蒙.
神給我們從祂裡面領受的寓言.
在耶和華大而可畏的日子 未到以前.
自然界有很多奇怪的事.
血 火 煙柱代表什麼呢.
在舊日聖經《若意書》裡面是說一個城市被燒毀.
日頭要變成黑暗.
剛才說了全日食的現象.
月亮變為血是代表著大自然的災害.
一說到大自然災害 大家又想到很多東西.
這一兩年都很多.
我希望通過《若意書》第二章28到32節.
和《史特恆主義》二章16到20節兩段經文.
來幫助我們的大英資媒去想一想.

$^{401}$就是說2024年4月3號台灣發生地震.
4月5號美國紐約發生地震.
和2024年4月8號全日食.
這些現象我們有沒有把它們放在一起去想.
究竟對我們的信仰有什麼提醒和作用.
這個是我很希望大家今天可以思考的.
可能你沒有想過這些事情天天都發生.
這裡說水浸那裡說旱災.
其實你可能沒有什麼感覺.
但是當你將這些現象放在一起.
和在經文裡面一起看的時候.
或許你會看到其他東西出來.
第二部分我想說.
面對世界這些奇怪現象.
其實經文說的一件事不是叫我們不要想.
是要我們有異象.
是不是很有趣呢.
這個很有趣的.
我們活在耶穌降臨和耶穌再來這個時間的中間.
耶穌再來是一個大而可畏的日子.
他來到是要審判世人的.
這個不是開玩笑的.
吸引我注意力的地方就是.
耶穌再來之前大自然是有些不尋常的現象之餘.
人是會有異象和異夢的.
很有趣的同時發生的.
那裡有很多問題.
但也有些人他看東西是不同了.
你覺得是不是很有趣.
上帝做事的時候是很奇妙.
氣候的變化令到2024年有成百分之二十二的地方.
可以說我們成為最亂的一年這麼多年以來.
同時有百分之七十九的.
厄爾尼諾的現象.
回復中性.
一會兒說的名字你可能覺得比較複雜.
但我真的找到2023年.
聯合國世界氣象組織的報告.
你可以上網去找.
有中文的.

$^{441}$北半球的鈣雪量就接近平均.
但南極的冰量一月份減少了五分.
你一聽到冰融化了.
水位就會上升.
其實影響很大.
地球會越來越熱.
其實老實說現在也是五月.
什麼時候都八月了.
因為拉尼娜的氣候現象將要結束.
這兩個名字可能大家都很陌生.
厄爾尼諾現象期間.
全球氣溫會升攝氏0.2度.
到拉尼娜現象的時候.
全球氣溫下降0.2度.
世界就是這樣.
厄爾尼諾的現象期間.
亞洲很多地方都經歷乾旱.
非洲南非都有.
你想想近年這幾年很多地方乾旱.
拉尼娜現象期間.
南部發生乾旱.
但加拿大和亞洲有大量雨水.
你去歐洲一樣.
水浸得嚇死人.
2022年10月澳洲破紀錄的洪水.
正是拉尼娜現象的結果.
2024年4月4日.
津巴布韋總統宣告.
國家陷入災難狀態.
乾旱造成糧食短缺.
這就是厄爾尼諾現象的結果.
贊比亞馬拉威共和國都準備.
宣告他們進入災難狀態.
拉尼娜現象會造成風暴.
洪水氾濫等等.
你有沒有發覺.
現在的水浸很厲害.
你會發覺.
是為什麼.
最近阿聯酋的杜拜.

$^{481}$經歷了非常嚴重的水浸.
這些都是氣候反常變化的結果.
2024年4月23日.
世界氣象組織表示.
2023年在亞洲地區出現了.
極端的氣候現象.
你不要只想著是北半球.
亞洲都中.
伊朗連續三年乾旱.
5到10月阿富汗糧食短缺.
中國西南部都乾旱.
其實我們應該很感恩.
我們香港要吃什麼就什麼.
又不會缺水.
但真的世界各地有這樣的事.
2023年.
圖蘭低地的大部分地區.
烏茲別克斯坦.
喀薩克斯坦等等地方.
興都富什.
我自己找了出來.
都不知道在哪裡.
喜馬拉雅山恆河附近.
這裡我不讀這麼多.
總之那裡都發生很多缺水的情況.
中國西南部都有乾旱.
2023年幾乎每個月的降水量都降低.
而且印度夏季風相關的降雨量都低於平均.
所以其實在亞洲一樣有事.
其實大家心裡面想.
其實我東南西北走哪裡.
東面有事我就去西面.
西面有事去南面.
南面有事就去北面.
其實你發覺到處都有事.
那怎麼搞呢.
亞洲高山地區.
在青藏高原那裡.
都是極地之外的含冰量大地方.
有覆蓋的面積是十萬平方公里.

$^{521}$但過去幾十年裡面.
這些冰川大部分都在退縮.
速度很快.
你發覺有些地方是不可以讓你滑雪的.
很有趣.
在黑潮洋系.
洋流系統.
阿拉伯海.
巴倫茲海.
南克拉海.
以及拉普提夫海東南部的海面溫度升溫.
是全球平均的海面溫度的三倍多.
你想一下這些東西.
就影響氣候海洋的變化.
這些我都是簡而言之.
總之你一看.
周圍都有很多大自然的問題.
這些都跳過去了.
讓大家不要那麼煩惱.
總之.
剛才那些大自然的東西.
加上經濟方面.
全球的經濟仍然放緩.
伊朗革命衛隊扣留以色列有關的貨櫃船.
伊朗也發動對以色列進行無人機和火箭炮的襲擊.
這些是報復.
以色列人在敘利亞和伊朗那裡.
領空被轟炸的地方.
人類現在是活在一個不穩定的狀態.
我們是感受著一件事.
世界是逐漸走下坡.
大家同不同意嗎.
這是一個事實的情況.
我相信你再想多一點的時候.
你會發覺我們正在調教.
我們平時晚上靜下來.
都是在想這些事情.
史特蘭傳第二章16到20節和約爾舒二章28到32節.
它是說面對大自然奇怪的現象的時候.
我們是需要有意象.

$^{561}$我有什麼意象呢.
我們需要從上帝那裡來的意象.
面對一個非常態的世界.
我知道有很多人是用另一個方式去面對的.
我想想自己心情好一些.
或者做一些可以做的事情.
或者去旅行等等.
絕對沒有錯的.
這些絕對沒有錯.
但是大家都知道有些事情.
那些事情不持久的.
雖然是好的事情.
但是不一定持久.
我們需要從上帝而來的意象.
面對一個非常態的世界.
所以約爾舒提出耶和華日子的意象.
史特蘭傳提出主在內的意象.
耶和華日子就是審判的日子.
主在內的日子都是審判的日子.
這就是這兩段經文一個很重要的提醒.
上帝不是要來毀滅人類.
他是要審判世界.
啟示錄最後那兩章的經文.
大家不要以為啟示錄是講世界末日.
啟示錄最後那兩章的經文是講新天新地.
你明不明白.
有很多人對啟示錄就以為.
什麼文要攻打文.
國要攻打國.
總之就是那時候打到七彩.
其實關鍵都不在那裡.
雖然篇幅很少.
最後那兩章其實是講新天新地.
新耶路撒冷.
今年三月復活節期間.
我自己一個人去了美國三環市灣區講道.
我就和建道神學院的校友見面.
有校友就分享.
他為什麼決定去灣區那裡事奉呢.
他有掙扎的.

$^{601}$有些傳道人是真的能撲到去別處的.
但是這位傳道人其實是.
他不是很想去的.
他在香港的教會想安立他.
總之就傳一番後不說那麼多了.
最後決定的原因.
他發覺他對AI元宇宙的興趣很大.
他就很希望探討這些新科技.
如何影響人的生活.
究竟基督徒如何回應.
然後他就說.
他說因為Google Facebook.
這些科技公司是發展AI.
GPT等等這些東西.
這些東西都在改變人類的生活.
現在搞AI都很厲害.
搞到我都不學.
其實我都沒興趣學了.
我打算退休.
但是似乎是不可以不去學.
現在講的老師人人都有AI.
很厲害的.
我現在的powerpoint那麼差.
你也知道肯定我沒有AI.
是不是可以學呢.
我們當然要有耐性的人去教我們才行.
我等下神學院有沒有這些安排.
究竟我們基督徒在信仰上.
如何回應這些世界的變化呢.
我查過美國灣區有沒有大學.
神學院研究這些課題.
答案是有的.
講座也不少.
我未必會做第一手的研究.
但是我很喜歡看別人的東西.
讓我面對這個急變的世界.
各種對人生生活的困擾.
我覺得神讓我能夠有智慧去應付.
和去適應.
並且我覺得神讓我有意象.

$^{641}$這個我不是要講自己.
而是藉著這個經驗和大家講.
你如何得到神的意象.
關鍵重點在這裡.
當社會急劇變化.
舊的基督信仰解釋的方式.
是和現實的生活.
我們叫disconnect.
是失聯.
就是不對位不對口.
所以牧者是需要首先要.
重建一種基督信仰.
在現在現實生活.
適合的切入點.
這件事真的高難度不得了.
真的很難.
如果我們還是講60年代.
70年代那些講道的內容或者方法.
我想走了一半人都可以.
你明白嗎.
現在是不行的.
但問題是我們要有一套新的.
比較好的解說.
又沒有那麼快可以.
那怎麼辦呢.
其實是鬱悶在這個位置.
如果我們不正視.
舊的基督信仰的解說.
和現在的現實生活.
有失聯的現象.
我們就會發覺什麼現象呢.
現在香港就會出現離堂會的潮流.
很多人不上教會.
特別是年青人不想上教會.
逼到你不去面對這些問題.
我真的要講一點點.
我怕我過了時間.
但我以前的時代.
在教會的文化裡面.
多數人是不問問題的.

$^{681}$是不問問題的.
問問題就是說你沒有信心.
所以那時年青人問問題.
通常都會被熟齡長輩.
和我們聊天.
其實你應該多點祈禱.
多點依靠主.
意思是什麼.
你不要問那麼多.
你先馴服了.
所以我們是在這個氛圍.
孕育出來的.
但我們的本性是喜歡問問題的.
所以你明白我們的掙扎.
在這幾十年裡面就是這樣過活.
但我們見到的就是.
香港教會的文化是變了.
其實大家都開始問很多問題.
不會你問問題就立刻給你一頂帽子.
這位弟兄真的沒有信心.
他靈命比較差一點.
你明白嗎.
真的很可憐.
無端端問問題就被人說我靈命差.
到現在有個好處.
在現時的港獨還沒有即時發明的時間.
如果有都很難搞.
如果你在網上就立刻在chat room有回應.
但我想說的是.
其實現在我們是要優先處理的.
基督信仰的解說色淺性的問題.
就不只是活動的方式吸不吸引人的問題.
或者我會不會對一些我想要的目標人物吸引到他.
其實你想想再這樣下去的時候.
你是不長久的.
因為那些人來教會是要找的是上帝.
你給到他的東西.
你給到他入門的東西.
你給不到他深層的東西.
這是我自己的體會.

$^{721}$如果是這樣我們一定是在打防守波.
宗教會都是一樣.
我們在做的事全部都是人家走一步.
我又退一步.
我又試試這樣.
但關鍵的問題是我有什麼信息是必須要去傳揚的.
我們今天有什麼信息是人人都需要聽的.
而這個信息我是很深信來自上帝的.
以至到我去到哪裡我都會這樣去做.
我就會覺得我們基督徒的群體教會.
就能夠有個位置站得好.
我站在這個位置.
我就知道我有使命去傳揚.
我自己個人就是性格內向.
我不是很喜歡和人交際.
其實我來港獨.
我都覺得和我本性是完全相違反的.
但我有一樣東西是專注去讀聖經和神學.
其實樣樣東西都不專注.
就專注這些性格.
第二我看到的異象就是.
教會要更新復興.
現在疫情之後宗教會是大洗牌.
大家都知道.
我們看看哪間教會少了人.
其實我們已經收了很多風.
教會人數如何增長呢.
其實一個是重要的課題.
但另一個問題就是.
屬靈的質素你如何去提升.
這個才是根本問題.
我自己就會用實事求是的角度去看教會的制度.
90年代,千禧年代.
適用的教會模式.
在2024年都不一定適用.
我這樣說是不是很顛覆.
但「顛覆」這個字也不是我先用的.
我的同事高明軒牧師剛剛出了一本書.
顛覆上帝的形象.
你這樣看人.

$^{761}$上帝不是這樣看的.
這段時間說顛覆.
顛覆的比較多.
我做神學研究.
我是可以做港獨的牧師.
你叫我做駐堂的那個.
我想那間教會都應該不會怎樣發展的.
做主任牧師更糟糕.
我都不敢說有什麼後果.
在神學院裡面我們有童工團隊.
而這班童工是有經驗的.
他幫我打開眼界.
所以你覺得我懂很多東西.
其實原因是什麼.
我後面有一班厲害的人.
面對急速的變化世界的時候.
我們首先要看清楚我們站在哪裡.
我覺得數據的分析是很重要的.
我不是數字派.
因為你要看到實況.
然後你量力而為.
但你不要退縮.
因為你知道現在是在打仗.
你要把機會去拓展上帝的國.
這是必須有的態度.
因為你不動的時候.
你是留什麼給後面的人.
教會是有會友有穩定的奉獻.
但你想想神學院和教會機構是沒有會友的.
我們做少一點就沒有人支持.
所以我們變成和大家的心理不一樣.
我們要快速回應信徒的需要.
希望得到大家的認同和支持.
所以在疫情期間.
我們用了Zoom做了很多服事.
發覺我們突破了地域的限制.
所以我們現在想事情和別人很不同.
由於英國和加拿大有很多香港的新移民.
和很多屬靈的需要.
所以我們可以考慮支持他們來建立教會.

$^{801}$我上次去加拿大很開心.
我看到很多大英姐妹都還未見齊.
一call就十幾人出來.
如果是的話你會在我的Facebook也看到那些照片.
我說這樣真的厲害.
你會發覺原來有些牧者告訴我.
那些去了其他地方的大英姐妹.
繼續參加本地教會網上的團契.
又崇拜.
我就開玩笑地說.
那他們有沒有奉獻.
這樣開玩笑.
原來你發覺我們的服事可以跨地域.
你不要只想香港.
如果你加上普通話的form需要的時候.
你真的發覺.
你絕對是超越香港走向世界.
我們真真正正是實踐普世宣教的使命.
現在真的可以.
最近有牧者告訴我.
原來在北區火車沿線.
有很多教會去直堂.
這是公開的秘密.
有一間教會準備在上水開一間教會.
應該是以普通話為主.
為什麼呢.
我不知道感受有沒有這些大英姐妹.
有些在國內的大英姐妹.
星期天過來這裡崇拜.
有人還好笑說.
以前在那裡開不了書局.
現在可以開了.
因為所謂人數增加了.
所以你可以說.
你不是說香港的人去消費.
我們打擊香港的零售.
人來人往.
所以有很多人已經看到這個現象的時候.
已經部署著怎樣把握機會去傳福音.
建立教會.

$^{841}$去直堂的是我們以前的校友.
不過也不要說這麼多.
你明白嗎.
人家已經做好了.
已經整個計劃出了.
人去怎樣動員去.
遲一點你又說不好聽.
其實也不用怎麼做.
如果我們這樣的時候.
你怎樣去想將來的事.
我也要承認.
我和大家都是那個年代的人.
在在有些比我年輕.
但我也是想著退休.
大家怎樣去想未來的日子.
我就很想告訴大家.
我又有龐大的寫作計劃.
我用這個字很嚴重.
龐大.
真的不是開玩笑.
給我真的去征服神.
我今年也會有兩本書出版.
書展也應該有兩次的新書發佈會.
我們神學院有另一本.
去年一本書叫我做新書發佈.
為什麼要做這件事.
我剛才也和一些弟兄說.
我們香港很多教會的同道.
你叫他搞活動不用擔心.
他們一定比我厲害.
搞活動 搞宣傳 搞很多東西.
但我覺得有什麼可以留給華人教會.
當然現在是後文字時代.
但你不只是搞活動.
你要有些思想的東西.
所以書本和思維性的東西是重要.
你怎樣重新解釋基督信仰.
對我們現實生活的意義.
為什麼歷代有很多人.
重新講出為什麼人要信耶穌.

$^{881}$因為時代變了.
你要面對這件事的時候.
你就要解釋.
所以出版可能是一個服事的資源.
現在還不是實體書.
還有e-book 網上的書.
無縣佛界.
而且可以復興很多教會.
所以我就去想.
看來我應該要做這些才對.
我很想和大家說.
神都會用你過去的經驗和你現在的處境.
第三樣.
我們要馴服聖靈的帶領.
當你要去思想你心裡面的願望.
但你的願望是不是聖靈給你的意象.
意象不是我們而來的.
意象是從聖靈而來的.
我們需要進行一個屬靈辨別.
英文叫spiritual discernment.
即是你要找個屬靈的前輩和你聊天.
來分析我們裡面的心理狀態.
你的感受.
你是否聽得對上帝的聲音.
還有你要去找上帝的印證.
對我來說.
我自己對德國的宣教很有心.
但其實我還不是很能做事.
我成立了中歐基督使團.
歐洲宣教團隊.
一加入WhatsApp group.
其實沒有事工.
你搞這些東西.
根本我沒有時間做.
其實已經超過我所想像.
我的禱告就是說.
不要去建立山頭.
大家都明白我說的話.
這些到處都是山頭.
如果這些東西不合神的心意的話.

$^{921}$神就要蓋起來.
要摺起來.
你明不明白.
當你在教會圈子內的時候.
你就知道東一個山頭西一個山頭.
個個都是大哥.
我想我們現在要想的就是.
無論你做什麼事.
都是以榮耀上帝來做的.
如果去到某個年紀.
你還是想做大哥的時候.
你都想著你之後都不會很好.
因為上帝不喜悅這些東西.
我不可以憑著感覺去做事.
不可以憑著我覺得有負擔去做.
而是說上帝想做什麼.
我們要跟著做.
上帝開什麼門.
我們就跟著去做什麼.
這個很重要.
如果我們不過到這一關.
我們就仍然是用自己的領袖.
去事奉我們上帝.
小聰明和大智慧是兩回事.
你想想就知道了.
我們要去問上帝.
你在你的國度的拓展之中.
有些現在什麼東西.
有價值的東西我可以做.
就是上帝正在做的那件事.
你跟著去做.
就不是我想做的那件事.
上帝你去跟我的想法去帶我走.
不是這樣.
面對急變的世界.
我們不要浪費時間.
特別是快要退休的人.
他更加會想著.
我剩下來的時間是有限的.
我就不會浪費在一些沒作用的東西上.

$^{961}$我們就是要集中精力.
去為上帝去服侍.
大家同不同意嗎.
你應該這樣計算的.
那件事你發覺做了也只有一個空隙.
你不要去做.
不過這樣說.
如果上帝是要你去做那件事.
那你就要做.
但是正路來說.
當然要做一些對上帝的角度產生影響力.
作用的事.
對我來說.
每一次的講道.
我都是盡力而為的.
你看我寫完的.
雖然寫的字不吸引.
沒有圖畫.
但你都應該明白.
我都沒有騙.
沒有找品.
這個很重要.
我們要從上帝而來的聖經的信息.
你嘗試去回應時代.
同時你要引發一班弟兄姊妹.
你去為耶穌基督去呼召門徒.
我自己覺得講道就是要做這些.
就是要讓大家不要浪費時間來聽.
而是你聽了真的有衝擊的.
我們的事奉.
目的就是要將人帶去上帝的面前.
讓上帝親自和大家講說話.
所以其實你不記得我說什麼.
我一點都不介意的.
你明不明白.
但當這個是一個渠道.
而上帝和你說話.
這樣就我成功了.
你明不明白.
我的服侍而上帝正在使用.

$^{1001}$這樣是不是很興奮呢.
所以我很希望.
各位你自己聽到上帝和你說話的聲音.
面對這麼龐大的屬靈需要.
真的不是.
全世界都是.
我們不能夠只靠自己的能力.
我們一定有並肩作戰的同工.
我現在去溫哥華和三環市灣區服侍.
其實我背後有非常龐大的禱告團隊.
不是他們禱告.
我也很不濟.
如果沒有他們的支持.
我就根本沒有辦法專心去應付講道.
天天都這麼多事情.
我都不明白為什麼.
大家可能都知道.
進來的時候有弟兄姊妹問我太太怎麼樣.
那次真的很驚險.
在溫哥華回來.
你知道吃完早餐就要一個小時才下飛機.
那時候我太太去洗手間.
她現在還在.
她很精神.
在去洗手間的時候.
出來的時候有點暈.
然後整個人在飛機上掉到地上.
嚇死了.
然後服務員叫醫生來.
後來發覺小腿的血管又硬化.
變成血凝固了.
然後上到去.
但去了是右肺.
你知道這些東西可大可小.
可以上去心.
可以上肺.
肺可以是一邊.
又可以塞在中間.
左右肺都去不到.
上腦就中風.

$^{1041}$你明白我的意思嗎.
是非常危險.
所以我有一次非常特別的經驗.
但我要這樣說.
我們香港的醫療系統相當高水準.
我是不需要入境.
入境處的職員.
是直接登機.
我和太太的身份證.
然後我們上救護車.
就去北大嶼山醫院.
通常大家坐飛機回來有事就先去那裡.
然後做完檢查之後.
就轉去瑪嘉烈.
系統是這樣運作的.
最後去了瑪嘉烈心臟科的心切治療部.
最後檢查了這麼多東西.
就發覺山室的問題去到肺部.
當然是嚇死我們.
你明白嗎.
風險是大到不得了.
所以我感恩有禱告的團隊撐住.
所以結果去灣區那次只有我自己去.
所以是這樣的情況.
所以你問我.
其實我不想自己一個人去.
因為我很想兩夫婦一起.
今年七月都會有一次.
和牧者和實學生去德國瑞士法國交流.
我們不只是看風景打卡.
好像Gordon太太都應該知道.
領教過我.
兩位都去過.
我上一次很勇敢.
我們整天晚上去短孫.
去很多教會.
有個弟兄余景灣的弟兄說.
國牧師我們年紀比較大.
當時2017年.
你晚上行程比較密.

$^{1081}$我又擔心我受不了.
我說很簡單.
你那晚不去留在酒店就可以了.
結果那位弟兄每次都去.
最後.
因為我們去很多地方.
由萊比錫到Marktburg.
很多教會不是旅行團是短孫隊.
當時我們去.
加上他們唱德文歌.
在德國教堂那裡一起獻唱交流.
這裡是七月.
我自己覺得.
現在是要做一些拓展上帝國的時候.
我們做這件事.
我亦都知道.
不要忘記堂會的需要.
地方堂會有很多肅靈的需要.
所以在地方堂會的服侍.
我們有做農夫的心態.
很有耐性去殺種子.
你等他成長.
真是很不同.
有些時候是欲速則不達.
我們需要對上帝的心意要有馴服.
檢查自己的內心.
是不是在服從上帝的意思.
其實我都掙扎了很久.
這些事情.
我告訴你.
香港的情況很獨特.
元朗區都很獨特.
其實我們是需要了解.
全香港的教會發展的趨勢.
還有元朗區的特色.
人口的資料全部拿到.
統計處有給你都可以.
我們是有實事求是的心.
為教會去把脈.
你都要身不檢查.

$^{1121}$你去尋求對教會發展.
一個最適切的模式.
這個很重要.
每個牧者.
每位信徒領袖.
都有自己對教會的理念.
這是很正常的.
我又有 人人都有.
但最重要的是.
在我們嘗試和實踐之中.
你一定在過程裡面.
你會看到很多東西.
我們一定要保持客觀的態度.
願意去查找不足.
願意去面對問題.
願意去堵塞破口.
你今天不塞那個破口.
十年後就不得了.
在神學院的事奉.
我們是收到很多.
不同教會的情況的消息.
有些情濟不前.
有些做得好.
我們可能讓大家知道多一點.
我有一個提醒自己.
可能我不是一個先知先覺的人.
但我都不能夠後知後覺.
更加不能夠不知不覺.
如果是這樣的時候.
我只能用幾個字形容.
就是不負責任.
因為你發覺不動.
你不前進.
你就一定會後退.
這個是不能說的.
這是鐵一般的事實.
我的禱告就是教會的領袖.
不單止是有先知先覺的心.
他最低限度都有後知後覺.
但真的不要不知不覺.

$^{1161}$教會領袖包括牧者和信徒的領袖.
是不斷地進步.
在事工技巧上進步.
在團隊溝通合作上不斷進步.
在屬靈的品格上不斷進步.
在和上帝的關係上都需要長進.
你想想這些是基本的東西.
連這些基本的東西都沒有.
就不用說做其他東西.
很多人著重在事工上.
但最後發覺自己都burn out.
做到死.
我們需要面對的是什麼呢.
是我內在生命的真實狀態.
發覺我們壓力很大的時候.
我真的鼓勵大家.
你要去找人幫.
或者找人去尋求輔導.
不要逼死自己.
逼死自己是不行的.
每個熱心的基督徒.
都曾經有孤單的感覺.
那個無助.
不被了解.
不被肯定.
我們是需要有上帝的心靈醫治.
你試想想.
你的心未被神醫治的時候.
你去想些事 做些事.
其實都不是真的很有愛.
當我們回歸上帝的時候.
我們要懂得在上帝懷裡安息.
我們才可以重新得力.
我們在上帝面前.
我自己很坦白的說.
祈禱你不需要假裝.
你用一些很粗俗的話.
或者不加修飾的話.
上帝是很喜歡.
因為你拿你的真我和祂聊天.

$^{1201}$因為上帝知道你.
其實祂比你更加知道自己.
我所以覺得.
在上帝面前就不用打扮.
即是不用的.
你沒有需要這樣做.
我只要在上帝面前承認.
我自己本來就是這樣.
我是這麼不配的.
我的內心就會感到被釋放.
我不用做到某個款式.
你才接納我.
當我們經歷上帝的醫治.
我們就會成為受傷的治療者.
盧雲神父說的.
你才可以幫其他人走出.
他那個心靈的幽谷.
當我們的內心清潔.
我們沒有私心的時候.
我們就會看到上帝.
是會用你去做奇妙的事.
但你自己都沒有整好自己.
你又想做大事.
通常就會出事.
如果我們希望能夠.
獲得其他人的肯定.
我們首先要放下身段.
放下自己的尊嚴.
大家明白我的意思嗎.
你和我都是人.
你用什麼心去對待人家.
人家用什麼心去對待你.
教會就是這樣.
你用真誠就可以留下人.
這句話是在元朗某一間教會的牧者說的.
這間教會在元朗開辦了不是很多年.
大家都應該知道我在說什麼.
我不能說名字.
但這間教會的增長是很快的.
這位牧師在他還沒來元朗之前.

$^{1241}$已經告訴我他會來元朗.
他還告訴我.
另一位牧師也會退休來.
我說這樣就行了.
元朗的教會會大亂.
果然現在這間教會已經有500人以上.
我不說什麼了.
所以我只想說.
很多人在元朗區對信仰是有期望的.
有期望的人就會去有期望的教會.
很簡單這個道理.
多遠都不要緊.
那位牧師告訴我.
有人很遠坐車.
那地方地鐵下車也要下一輪.
這樣都行?.
因為他說那間教會有愛.
這句話真是打得正.
好了.
下一道.
過去20多年.
事奉上帝教導我.
和人有合一的事奉.
所以我夠膽說.
在錦繡堂和人合作.
我真的沒和一個傳道童宮叫過交.
我都沒和各位堂董執事叫過一次交.
因為我都沒開會.
但不開會平時也會碰面.
應該沒有.
沒試過和人吵架.
這些人通常選了做友誼先生.
友誼小姐那些人.
選不到香港小姐香港先生那些人.
是什麼呢.
你是最沒殺傷力的.
其實都OK的.
我覺得上帝是幫助我們學會和不同的人溝通合作.
其實很簡單.
你只要不要有野心去事奉.

$^{1281}$你面對人事的問題.
多複雜.
你都要去學習用慈愛公平公義的態度.
你就會處理得到.
就不會有那麼多人和你過不去.
所以我們說可以再見也是朋友.
所以有些童宮.
很好笑.
不要說錦宣.
孫德都有些.
離職的傳道人.
他找我去寫推薦信.
錦宣都有.
所以我自己會這樣想.
幫得到就幫.
每個基督徒都可以是不同層次的領袖.
我們是需要時間去時時刻刻的警醒.
一定要保持單純的心去事奉.
一定要用單純的心.
如果你是一個管理層的人.
你必須要向持份者負責.
每一個弟兄姊妹的一分一毫都是血汗.
我們要這樣想就行.
這就是我們的誠信和我們的公信力.
你只有這樣才可以讓弟兄姊妹知道.
我可以投身在這個教會裡面.
作為一個領袖的時候.
他一定要謹慎自己.
管理自己.
先謙卑自己.
讓上帝去帶領.
然後他去領受從上帝而來的恩賜去帶領人.
意象是上帝給的.
作為一個領袖必須要公平.
公開公正.
你才能夠得到人家的尊重.
究竟你怎樣保持體諒個別的人.
和整體發展兩樣的平衡呢.
這是很困難的事.
非常困難的事.

$^{1321}$我們唯有求上帝賜給我們從上而來的智慧.
按著上帝的心意去思想和行事.
聖靈的充滿不只是推動事工.
而是更新我們的生命.
事工是doing.
生命是being.
是我們內在的生命去承托事工.
不是事工去承托我們的生命.
不知道大家明不明白這句說話.
很重要.
上帝不偏待人.
只要你願意倒空你自己.
上帝就會願意藉著聖靈充滿你.
在現在這個時代.
這麼多複雜的事.
你猜上帝想不想做事.
你想想.
你覺得上帝想不想做事.
上帝要做事.
那會怎樣呢.
他期待些什麼呢.
他期待一班願意真心跟隨他的人.
願意倒空自己去事奉他的人.
願意聽他的說話.
在現在應該用什麼新的方法去做事的時候.
我想這件事是重要的.
如果我們有這個心願.
我真的很希望大家多走一步.
你再聽聽上帝有沒有什麼具體的說話跟你說.
我說這些只是一些一般的訊息.
我希望這個訊息能夠幫到你.
我現在都不理那麼多.
講到都像陪靈會那樣說.
總結.
《史唐尊》第二章二十節.
他在引用《若意書》的二章三十一節.
是說什麼呢.
主的日子大而可畏.
面對世界那麼多奇怪的現象.
我們不要被它嚇怕了.

$^{1361}$反而你要知道聖經已經講了.
這些事是會發生的.
那我們要怎樣呢.
我們要信服聖靈的帶領.
去經歷上帝在我們生命裡頭的工作.
然後經歷上帝的大能.
以致我們不要白佔土地.
以致我們雖然年齡生理是自然會退的.
你會衰老的.
但是我們是仍然有剩餘的價值.
大家明白我說什麼嗎.
我們是仍然有剩餘的價值.
就是說年輕的人有青春有魅力.
但是我們這些拿兩塊錢的.
我們仍有剩餘的價值.
在上帝的國度裡邊.
還有些工種我們可以去做的.
上帝仍然願意用我們.
我們一起同心祈禱.
我們親愛的天父上帝.
因為你就是那位愛我們的上帝.
你自己來帶領我們看到世界的變化.
還有很多的事情超出我們可以控制.
但同樣的你也是在我們心裡激動我們.
給我們一種負擔.
就是說神要做事.
我們有時候都很矛盾.
我們都不知道怎麼去回應.
想做又不想做之間.
我們覺得那些事情根本是沒得做.
或者很艱難.
但是神也在我們的內心常常都在搗亂著我們.
給我們有這種掙扎.
我們就很希望這些經文能夠給我們有幫助.
是你給我們有意象.
不是我們的想像.
是你去感動我們.
你猜謙我們.
你是要去復興我們.
你給我們有意象.

$^{1401}$給我們有能力.
給我們有專注.
求主你幫助我們自己生命的更生.
以致我們看到的不是世界的混亂.
我們看到的神的心意.
正是想將這些混亂的東西變成正確.
讓那些被這些混亂的現象搞到心煩意亂的人.
可以認識上帝.
求你讓我們看到神的心意.
神在向我們說話.
祈禱奉耶穌基督的名.
阿們.
\newpage



\section{帖撒羅尼迦前書}
\label{sec:eejVMLNyY8I}
\textbf{2024-05-30 周四晚查經班《帖撒羅尼迦前書(7)》 講員:曾裔貴牧師}
\newline
\newline
連結: \href{https://youtube.com/watch?v=eejVMLNyY8I}{\texttt{https://youtube.com/watch?v=eejVMLNyY8I}} ~~~~ 語音日期: 2024-05-31
\newline
\newline
\hyperref[sec:QGag1Zr1Cdk]{\small{< < < PREV SERMON < < <}}
~
\hyperref[sec:index_chronic]{\small{[返順時目]}}
~
\hyperref[sec:index_scriptual]{\small{[返順卷目]}}
~
\hyperref[sec:7PHC2DLBPCo]{\small{> > > NEXT SERMON > > >}}
\newline
\newline
$^{1}$今晚我們一齊來禱告.
多謝天父.
再一次我們可以在晚上.
大家來讀聖經.
貼撒羅尼加的前後書.
以及關乎到末世主再來的一些神學的討論.
無論是一些千禧年.
或者是不同的前遷或者是後遷的討論.
一些系統神學.
都盼望我們能夠慢慢掌握一些概念.
歷史歷代都有不同的神學家.
對於主耶穌的再回來.
都有不同的討論.
或者有不同的結論.
都幫助我們認識.
無論哪一個派別.
大家都肯定的就是.
主耶穌一定會肉身回來.
這是我們每一個基督徒很重要的盼望.
願神都讓我們透過看貼撒羅尼加的前書.
都能夠掌握這個重要的真理.
也都真是透過這樣的形式.
去在我們的生活裡面帶來一個改變.
就是我們有意識地知道主耶穌會回來迎接教會.
使我們能夠在地上的生活.
都能夠貼近主的心意.
我們一起在你面前等候.
禱告奉耶穌基督的名.
上一次第四章.
貼前的第四章.
我們就大概講了.
聖傑的生活和彼此相愛兩方面.
是使徒保羅來對教會的勸勉.
我們可以由第九節.
四章的九節再讀下去.
大家可以看看那些經文.
他說:論到弟兄們相愛.
不用人寫信給你們.
因為你們自己蒙了神的教訓.
叫你們彼此相愛.

$^{41}$弟兄們向馬其頓傳遞的傳遞卿.
固然是這樣行.
但我勸弟兄們要更加勉勵.
又要立志作光正人.
辦自己的事.
親手作功.
正如我們從前所吩咐你們的.
叫你們可以向外人行事端正.
自己也就沒有什麼缺乏.
沒有缺乏就是他們自己勞力作功.
自己來到去.
信了主之後.
有一些務實的生活的改變.
這裡是正面的.
如果已不所書第五章.
那裡又叫人勤力工作.
那些是比較負面的.
他說你從前偷東西的那些小偷.
不要再偷了.
信了主耶穌之後.
那就要勤力工作.
你就會有餘了.
就可以知道怎樣去分配給其他人.
就是說不單單是不可以再偷東西.
就是說要改變你的行為你的習慣.
有些人是習慣偷東西的.
習慣的反而是.
我有一段時間在中學的時候.
中四中五可能中四的時候.
有其中一個同學很喜歡去七級市場偷東西.
後來我也跟了他一段時間.
真的會有一種滿足感.
你能夠偷了放在袋子裡.
當然是那些小的.
不是拿尿片那些.
是拿小的朱古力,糖果.
那些放在袋子裡或者放在書包裡.
那時候又不會有防盜.
基本上你出去了就可以了.
所以也維持了一段時間.

$^{81}$我不知道最後那個同學有沒有被抓到.
會有這樣的慣性.
偷偷有一個人.
那個《以弗所書》.
保羅就直接教導信徒主耶穌的.
《以弗所》裡面的弟兄姊妹.
你從前是偷東西的.
現在不可以再偷了.
這樣也很醜.
《以弗所書》的弟兄姊妹.
誰在上教會之前是偷東西的呢?.
他又沒有說是誰.
就是說你自己也去入座了.
但是這裡就強調在彼此相愛之下.
弟兄姊妹就要.
尤其是第十一節.
這裡說立志的作安靜人.
為什麼安靜人要立志呢?.
要有一個崇高的志向.
來安靜的在教會裡.
不要多是非.
不要有很多的胡亂的評論和猜測.
或者好像是「廉恥鬥非」那樣.
要做安靜人.
辦自己的事.
親手的作工.
這個似乎不只是教會.
而是在平日的生活裡面.
跳下來的這個地方.
很早就說過是亞歷山大大帝的妹妹.
來紀念他.
為他而建的一座城.
這個地方其實是相當富庶.
因為它真是在一個很重要的港口.
也是亞歷山大大帝的爸爸.
當時在加沙德王朝.
來建造.
所以一直到羅馬帝國.
他們都是有那個最惠國的待遇的.
即是說羅馬的駐軍.

$^{121}$是不用給那些軍費的.
正因為他們擁戴當時的奧古斯都.
所以是一個相當富庶的地方.
不是一些內陸貧困貧瘠的地方.
所以你去到切索羅尼加.
你會看到亞歷山大大帝的銅像.
相當出名的.
這裡特別提到.
貼城的弟兄姊妹.
他們信了主耶穌.
仍然要保持他們自己的勤奮工作的態度.
要親手作工.
如果這裡隔閘了主在內.
有一些錯誤的扭曲的心態.
主耶穌都回來了.
還這麼勤奮幹甚麼呢?.
你還說甚麼十年計劃呢?.
還說甚麼大規模的一個研究呢?.
不需要讀書.
不需要做一些很長期的投資.
可能會有些扭曲的想法都不頂的.
所以主在內說得這麼詳細.
這個是13到18節.
就一定會有些餘波的.
所以第五章就有了.
到後書直接說第二章.
貼城後書第二章.
直接說叫那些人不要這樣想.
主在內跟你做事不可以掛勾的.
主在內跟你勤奮工作.
繼續在你的崗位盡忠.
是兩者沒有衝突的.
這個才是真的不容易.
所以很多時候就是會將那個真理.
變成一種片面的去理解.
就會帶來出現了一些好像要補鑊一樣.
就是要勤奮.
但是這裡似乎又沒有說太多那個背景.
為甚麼要立志的做安靜人.
為甚麼要辦自己的事.

$^{161}$要親手的做工.
好像教人要勤奮工作.
不知當時那個社會的氣氛是怎樣.
是不是崇尚一個五天工作.
古代不會有這樣的計算的.
特別是一個農業的社會.
不是這麼高度的商業化.
更加不會說24小時營業.
晚上沒有油燈.
街道已經漆黑一片.
是不會有人出夜街的.
所以一定是根據那個日出而作的.
那個作息的.
不同現代現在的社會的.
所以這裡提到的親手作工.
似乎保羅是不是有些的預見.
教會需要在一班人等候主耶穌再來.
我們都說過了.
每一章的結尾都提到主再來的.
在《潛受留念》的前書.
每一章都說的.
都說主耶穌再來的信徒要怎樣裝備的.
所以這裡這樣的時候.
保羅是不是意識到.
說主再來和我們在我們的生活.
我們的那個召命.
我們的崗位.
不應該有衝突.
不應該有一個抵觸.
反而就是更加要給那些不信的人.
有一個很清晰的圖畫.
或者是一個見證.
就是等候主再來的基督徒.
就是這班「潛受留尼」加成的人.
他們在社會上不是缺席的.
不是避世的.
不是在生活裡面很奇怪的.
所以這裡很需要我們能夠有一個平衡的信仰.
保羅好像在這本書.
我們感覺就是很強調主的再回來.

$^{201}$和「潛受留尼」加成的弟兄姊妹的關係.
叫他們不斷不斷地等候.
要警醒 五宗說了很多次警醒.
要好像一個白晝之子那樣.
不要沉睡.
等一下我們會說.
雖然一兩節的勸勉.
特別很強調就是.
你自己沒有缺乏.
這個是12節的下半節.
行事端正 向外人行事端正.
似乎這個行事端正就是說.
我們不會讓其他被信的人.
找到一些我們不好的地方.
如果整班基督徒都不做事的.
整班基督徒只是天天聚會 查經.
那就不健康了.
你應該有你自己的生活.
有你自己的家庭.
可能有你自己的田地.
你要耕種.
可能有些基本的.
簡單的商業的貿易.
你們都要從事的.
或者有些是當兵的.
他要打仗的.
要保家衛國.
在當時的「潛受留尼」加成的.
一定會有這樣不同的崗位的.
雖然聖經很難說.
每一個基督徒.
當時他們是什麼背景.
現在不是一個會有調查.
而是保羅在一個這麼緊急.
要離開 趕著去雅典.
在雅典現在就稍微平靜了.
他又要寫一封信去勸勉.
「潛受留尼」加成的弟兄姊妹.
所以只能夠是很一般性的勸勉.
一般性的勸勉裡面.

$^{241}$都要求他們在生活裡面.
要勤奮工作.
所以基督徒真的要調查.
我們是不是勤力工作呢?.
我們有沒有一些「呃眾援」不做事.
來去罷工的.
有沒有這樣的人呢?.
可能我們感受就不會有.
但如果有些基層教會.
可能實在要教導.
因為這裡很明顯是說.
要納住安靜的人.
似乎就是自己要有這樣的修養.
自己要在彼此相愛的裡面.
不要變成一種濫用.
不要變成我關心你.
其實可能關心他.
想聽聽別人的是非.
或者你想知道他的秘密.
這些有時就超越了彼此相愛.
然後說辦自己的事.
親手的做工.
似乎這裡就是說日常沒有聚會的時間.
那六天怎樣工作.
如果大家有看過《出埃及記》.
十屆的第四屆.
六天要勤力工作.
但第七天就要安息.
所以那個大前提是做六天.
然後第七天要安息去敬拜神.
所以這個都是一個很重要的.
我們人生周而復始的一個循環.
真的需要勤奮工作.
所以很多人就是不斷做不斷做.
慢慢發現自己不知道做來為了甚麼.
因為他沒有了第七天的安息.
去到神裡面的敬拜.
不只是回教會.
而是意識到為甚麼你要存活在這個世界.
你的工作的原因.

$^{281}$各方面是為神而活.
所以那個是一個循環.
其實原來對身體都是一個很大的復原.
如果一些社會是不停工作.
不能休息.
其實最終那個公司或者僱員都受虧損.
一定會有影響.
甚至會慢慢遺失.
原來休息是更好的工作的準備.
所以這裡表面上好像很簡單.
其實是工作的倫理.
當然我們保羅當時沒有講到這麼深奧的神學.
即職場的倫理.
職場的神學.
沒有發展到這麼複雜.
但這裡他已經告誡弟兄姊妹.
必須在等候主在來.
在彼此相愛的同時.
在生活裡面.
在自己的家庭.
在自己的工作崗位.
相信是佔據最多時間.
除了家庭就是工作崗位.
最多時間佔據的.
你怎樣去投放心力去達到精一求精.
達到可以繼續在崗位裡面盡忠.
所以這也是對基督徒一個勉勵和提醒.
有沒有在崗位裡面做好自己的本份.
所以這裡告訴我們.
叫你們可以向外人.
外人肯定是微信的人.
行事端正.
即是說你看到我們.
都是有一個恩懇.
勤奮,盡忠.
所以這個「盡」字.
可以在我們的生命裡面.
來成為我們一個重要的座右銘.
一個「盡」字.
很多年前有一位老牧者.

$^{321}$是我們的長輩.
很出名的 在香港的教會界.
當然是二,三十年代的牧者.
他很大年紀的.
那時候我剛剛學習傳道.
他進來港道 進來元朗.
他本身在深水Po .
進來元朗港道.
他們夫婦進來港元道.
送他去搭車.
他給我一個字.
就是一個「盡」字.
即是一個事奉主的人.
要記住這個字.
「盡」字.
盡心,盡性,盡義,盡力.
愛主你的神.
即是十誡裡面主耶穌說的那個字.
那個「盡」字.
所以原來一個人.
你有多少去為主去工作.
所以職場的神學.
職場的倫理.
其中最重要的就是.
為甚麼要這樣做.
為誰而做.
這個就是職場最重要的神學.
否則你會迷失的.
就算你去到很高的位置.
其實你這樣做.
更加會帶來那個影響.
所以我們都要有這個學習.
同時端正自己.
就沒有甚麼缺乏.
一定是在金錢上.
一定是勤奮而帶來的自負.
所以也都是在說.
那個宗教改革之後帶來的.
就是清教徒的自負.
就是殷勤的工作.

$^{361}$殷勤的去儲蓄.
殷勤的在教會裡面去奉獻.
這個就是那個.
清教徒的新教神學觀.
或者說他們的工作倫理.
帶來的那個改變.
所以很早已經說了保羅.
要自己沒有缺乏.
沒有缺乏之後.
所以我們都說以不所書.
其實挺有趣的.
貪心的人才會偷東西.
貪心的人就是自我中心.
貪心的人就是甚麼都是自己要的.
但是奇妙的地方就是.
剛才我們還沒讀過.
以不所書五章.
我翻到前面來看看.
保羅這樣教導.
有時都挺搞笑的.
在一些書信裡面.
講的東西這麼白的.
你應該說教會現在個個都是聖人.
不是甚麼.
他說不是的.
以前有些人是這樣.
我讀給大家聽.
第五章.
對不起 第四章.
四章是尾章.
四章的28節.
他說從前偷東西的人.
就不要再偷了.
總要勞力勞力.
這個字真是.
用你的汗水.
用你自己的手腳.
古代現在可能用手機都不行.
做事都不用用電腦都可以了.
因為有網絡的世界.

$^{401}$總要勞力.
真是出汗.
真是用你的時間.
親手的做正經事.
正經的事.
以前不正經的.
以前貪心.
以前就貪便宜.
越偷就越多.
我是八歲還是九歲.
不是 應該是中一的時候.
讀新法的時候.
中一的時候.
就開始.
十一歲還是十二歲.
十二歲左右.
就開始學會偷東西.
晚晚有些人睡著了.
是爬在地上滾.
滾去公事桌那些櫃子.
那些手現在當然很粗.
很小的時候.
你想一下現在讀六年級的小朋友.
那隻手可以放進辦公桌.
就是那些櫃子.
主要那張桌子.
有鎖的桌子.
手在裡面遊走.
就去找那些紙幣.
每次拿一張.
星期天上學去午餐吃飯.
又不在家拿.
但是又要偷.
慢慢發現.
為什麼哥哥打開櫃子.
為什麼這些文件都亂了.
因為我的手一直在找.
偷端上癮.
最終有一次是被發現的.
打到飛起來.

$^{441}$以後就不敢了.
你想一下.
一直這樣建築上去.
就越來越有貪婪.
背後就是心底的貪婪.
所以這裡告訴我們.
親手作證經史.
接著很奇妙的.
保羅的教導.
他不是干涉你做人這麼簡單.
就是你在教會裡面.
你已經信了主耶穌.
你說就開有餘.
就開有餘之後.
是分給那些缺少的人.
不單不再自我中心.
不單不再貪婪.
不再偷.
還是有餘.
接著就去留意.
以前偷東西就是周圍留意哪些東西可以偷.
現在就留意哪些人有缺乏.
你就去幫那些人.
這樣去令自己的心.
不單不再自我中心.
而是有一種主耶穌的那種愛.
所以那個眼目去留意人.
你看看聖經給我們那種生命的改變.
是多麼奇妙.
不是說自己顧自己就算了.
我現在不再問人了.
我再不再貪心了.
我不會理人了.
不是的.
就開有餘之後.
就不是你可以買LV.
就是你可以去留意幫那些有缺乏的人.
你看看.
這個就是金錢觀.
財富.

$^{481}$還說了一個偷東西的賊.
你可以想像到.
聖經原來這樣教導我們做人.
所以我們回到這裡.
這個工作的倫理.
在這裡又很細緻地教導了.
是建立在主在尼.
和在彼此相愛之下.
我們要自己沒有缺乏.
這裡就說向外人可以行事端正.
教會裡面每一個弟兄姊妹.
都是在社會有好的明星.
不是被人感覺到我們很怪誕.
不斷說主在來.
但自己其實就不做事.
也不能夠去和人分享生活的一些工作的成果.
這些不是正常基督徒.
正常基督徒應該在生活裡面.
能夠反映我們的信仰.
帶給我們的生命的改變.
所以過度的工作.
現在近代.
特別是青年一代.
我們叫享世代.
享受那個享.
這個本身不是負面的.
是正面的.
開始他們要探索.
不可以用盡時間精力去工作.
而失去了自己的平衡.
例如自己的一些興趣.
一些的夢想等等.
或者一些參與的一些可能社會的公義等等.
所以現在這個世代.
你叫他不斷加班.
他不願意的.
他寧願辭工.
都有他一定的合理性.
所以我們都需要有這種的平衡.
我們上一代.

$^{521}$包括我自己這一代.
我們都是沒有那個.
所謂要很平衡.
我們的教導就是要盡量做.
你不做可能就沒得做.
但是有時過度了.
其實有些東西又不夠太平衡.
所以這裡就是告訴我們.
這個親手作工.
肯定是和六天工作有關連.
就要我們既有教會生活.
是我們很重要的.
我們那個安身立命的基督徒的生命.
我們的基督徒的照明.
但是離開教會之後的那個工作倫理.
就是我們要為了神而工作的那個.
重要的那個學習.
所以有工作.
有合適的作息時間.
有重要的一些的興趣.
或者是運動去調整那個新心靈.
都是需要.
如果過度了就不健康.
所以如果全都人不運動.
都是一個問題.
不出汗不運動.
只顧著坐牢.
長期都不健康.
會神經衰弱.
有時會想歪了.
所以那個作工都是要平衡.
這就是保羅對當時弟兄姊妹的一些勸勉.
其實認識她們都不夠兩個月.
頂多是一個半月的弟兄姊妹.
保羅就要急速的和提摩太離開.
一路由貼壽羅尼加.
輾轉一路走.
就去雅典.
在雅典才叫穩定.
就打發提摩太和寫信去安慰她們.

$^{561}$使到貼壽羅尼加的弟兄姊妹.
可以有這個意識到.
保羅不是因為逼迫而離開.
而是有他個人的一些神的帶領.
穩定之後.
保羅其實是很關心她們的靈命.
很想她們在一個整全的神學觀上.
來有一個成長.
所以這些都是牧者和信徒之間的關係.
那個屬靈的餵養的關係.
去到第十三節.
我們其實在其他不同的時間都有提過.
特別強調就是要問一個問題.
為什麼保羅在這一段裡面.
第四章第四次講主在來.
這次是講得非常詳盡.
唯一不詳盡的就是.
那個日子那個時間.
哪一年哪一月哪一日.
主耶穌回來.
唯一沒有講.
其餘的就是講到突然間.
有聲音之後.
耶穌回來的那個程序.
那個過程是怎樣.
就很詳盡的講.
而且特別是因為針對一個課題.
就是有一些頂子妹的家人離開世界.
他們現在在哪裡.
將來他們會怎樣.
因為這樣一個牧養的原因.
保羅就來講述主在來的過程.
還有哪些人會得到主在來的那份恩典.
所以朝這個方向去講去探討.
還有就著這個再來.
聖經絕對一點都不含糊的.
講了很多次的.
我引述了都有起碼.
1,2,3,4,5,6.
有八段聖經.

$^{601}$等下再和大家揭一揭.
都是講主回來的.
一定回來的.
還有用甚麼心態去迎接的.
這個就是肯定到不得了.
耶穌一定會回來.
是肉身回來.
是我們見得到的回來.
但是一樣東西我們要記住.
就是不要去計算主何時回來.
因為是一個奧秘.
是不會有人知道的.
一猜就錯了.
一定是會錯的.
是沒有人可以計算到的.
記得這個我們就會知道.
就是告訴我們.
保羅的時代講到這麼實在.
保羅在他一生都以為主會回來.
但是他們都會過去了.
他們都真的過去了.
那個訊息流傳在後世了.
每一個時代都有人計算的.
但是都是會失敗,錯誤的.
只會尷尬地修訂.
甚至變異的.
例如守望台就是這樣計算.
或者耶穌說見證人就是這樣計算.
主在來已經回來了.
然後就說原來還有一些時間.
他在天上.
這樣你怎麼扭曲都沒有用的.
因為當你一計算那個日子.
地上的日子就一定是錯的.
如果你堅持錯當然沒有問題.
但是錯你都不肯認錯.
那就是異端了.
因為沒有人可以計算到的.
但是這個就反而值得我們去問的問題.
那個這麼確鑿的一個情景.

$^{641}$主回來的時候.
那個復活的次序.
是很值得我們去反省.
這一段聖經對我們自己.
那個生命的影響.
我都在群組裡分享過.
上星期前的一個禮拜.
我的姐夫離開了.
一直再去思想的時候.
你想一下中午的一點鐘.
醫生在會議室跟我們講解.
講完之後.
那個醫生是很好的.
廣東話的那個女醫生.
當然不是很多.
我問她可以多久呢.
可以多久呢.
我的姐夫說按她的這麼多年.
因為她是善終病房的女醫生.
按她的經驗.
可以是幾天就會離開世界.
甚至可以去到兩個星期都有的.
所以她都告訴我們都講不了.
聽了這個訊息之後.
就走回病房.
其實隔幾個巷而已.
就是從203走到201.
這麼短的距離.
開始有感動.
就是如果我姐夫.
我心裡一直想.
如果我姐夫拖兩個星期.
糟了我下個星期三就要回香港了.
機票訂好了.
如果她不死怎麼辦呢.
心裡是有一個祈禱.
當然我不是詛咒她死.
我向神祈禱.
是感動的.
既然我姐夫都清楚.

$^{681}$已經洗禮了.
我下個星期三就要回香港了.
按你的時間讓她能夠解脫.
是我自己默討的.
我沒有在別人面前講的.
接著差不多要吃午餐了.
我就和所有的家人.
包括我姐夫的大家姐他們都來了.
就一起有個公禱.
在姐夫的耳邊.
我從來沒有試過這樣祈禱.
是開口.
我說姐夫.
如果你看到耶穌.
如果你看到大光.
你跟著祂去吧.
你不用擔心了.
祈禱完之後.
我就和我姐姐下去地下.
去坐下休息一會.
喝杯咖啡.
其實在二樓走下去坐電梯.
三分至五分鐘而已.
還沒坐下.
兒子就打電話下來.
叫我們馬上上去.
上到去.
姐夫是很安詳的慢慢斷氣.
前後我去祈禱.
到姐夫離開是五分鐘.
那是不是他真的看到大光.
他跟著那個光去呢?.
我看不到.
因為我不會看到突然間.
好像恐怖片那些.
有些東西閃進來.
我看不到.
我們所有人都看不到.
兩個護士進來量度血壓.
量度脈搏等等.

$^{721}$確認他真的安詳離開.
沒有任何掙扎.
他也沒有醒過.
一直是昏迷著的.
但是就是這麼奇妙地.
那個祈禱完了五分鐘.
他就安詳的離開了.
這個我們要回頭看.
和我們今天有關.
因為我們還是在生.
保羅是說對那些在貼壽羅尼加城.
帶著一些懷疑.
擔心的弟兄姊妹.
他們有什麼可以成為他們活著的盼望呢?.
我們再用這個角度去看回這大段經文.
論到那些死了的人.
我們不願意弟兄們不知道.
恐怕你們憂傷.
像那些沒有指望的人一樣.
沒有指望的人.
肯定就是那些未信的人.
但是如果基督徒我們不對主在來.
有這個熱切的渴慕.
可能我們都不知不覺.
我們不知不覺都會沒有那個指望.
沒有那個熱切的渴慕.
所以那個程度上不同.
我們可能在那個意識上.
知道是會有耶穌回來.
但是沒有那個迫切性.
我們好像必須要來到去.
有一個確認.
比如用一些可能大家都會明白的.
比如你真的明天要上飛機了.
其實你要確認的.
你要來到去訂機位.
或者要準備check in的東西.
這個就是你要double confirm.
這裡告訴我們.
第十四節本來就說到關於那個真理.

$^{761}$就是最好的安慰就是耶穌回來.
但是耶穌回來.
對於那些在生的.
特別是這裡針對的是那些已經死了的.
他們應該怎樣呢?.
他們是不是沒有盼望.
已經是灰飛煙滅呢?.
中國人的想法是怎樣呢?.
這段正經就很詳細的講論到.
每一個人.
就是所有的眾生.
都是一定在這個範疇裡面.
面對耶穌回來.
即是說有一個很大的神聖的能力.
掌控著所有古往今來的人.
死的,未死的.
所有人類都是在這個範疇裡面.
這段經文的重要性就在這裡出現.
首先就是安慰已經失去了家人的基督徒.
現在怎樣呢?.
接著就14節.
「我們若信耶穌死而復活了.
那已經在耶穌裡面睡了的人.
神也必將他與耶穌一同寨來.
一同回來」.
神會將這些已經死了的人.
與主耶穌一同回來.
即是說那些死了的人一定會復活.
這是很清楚的假設.
所以夏文尼也很清楚.
宣告所有的人都要復活.
所以第15節.
「我們現在照主的話告訴你們一件事」.
照主的話暗示不是我保羅的個人見解.
是保羅他得到的啟示.
即是那個嚴肅性是神的啟示.
不是我個人的羨慕.
是我自己的一些想法.
「我們現在照主的話告訴你們」.
就是這些哀傷的弟兄姊妹.

$^{801}$「一件事 我們這一活著.
還存留到主降臨的人」.
即是說總有一批人一定會肉身未死見到主.
可能是我們這一代.
可能是幾千萬後的一代也不奇怪.
我們不知道的.
「你告訴我們 到主降臨的日子.
斷斷不能夠在那已經垂了的人之先」.
又再一次說那些死了的人.
是要首先復活.
第16節.
「因為主必親自從天降臨」.
強調親自.
強調是肉身的親自降臨 耶穌.
「有呼叫的聲音和天使掌的聲音」.
這個呼叫似乎是靈界天上的呼叫.
然後再加上天使掌的呼叫.
然後又有神的號吹響.
然後哪在誰吹神的號呢?.
又似乎是天上的靈體.
「哪在基督裡面死了的人」.
不斷說垂了.
現在正式說死了.
即是通用的 共用的這個字.
在保羅的神學裡面.
垂了只不過是剛剛睡了.
他會起床.
但是死了就更加清清楚楚.
一定有一種很神聖的能力.
令到一個死了的人有了生命.
這個就是基督教的重要的神學.
沒錯 安人是不行的.
沒有人可以令到人復活的.
死了就沒有了 那個生命.
但是神聖的能力.
首先在耶穌身上死而復活.
所以我們信主耶穌的人.
將來就是首先會復活.
所以這裡說是那些死了的人.
在基督裡面死了的人首先復活.

$^{841}$然後以後這個首先復活.
是要多久呢?.
似乎是同一時間.
同一時間整個地球.
你想一下在和合石有一大批人.
在博弗林有一大批人.
都沒有計算全世界的那些墳墓.
還有一些在海裡會海葬的.
沉下去了永遠都找不到屍體的.
被魚吃了的.
你想一下有多少的人.
在基督裡面來復活呢?.
似乎是同一時間.
就不需要有時差.
先到香港 接著第二天天亮了.
到加拿大就沒有這件事了.
一定是即時的 同時間的.
你可以想像到神聖的能力.
你想一下這件事你都覺得很奇妙.
神聖的能力讓所有.
自從有人類開始.
即是阿當 喬娃這些要復活了.
同一時間.
你可以想像到.
不是有些穿古裝 有些穿青裝.
不是這樣的.
所有的基督徒.
是有一個新的肉身的生命.
當然這個是靈界的肉身生命.
是真的有身體的.
不是靈體 不是靈魂.
是有一個身體這樣復活.
是一個榮耀的身體的復活.
每個人都有.
到時我是不是六呎呢?.
還是怎樣呢?.
我羨慕的小時候.
現在就算了.
因為已經沒有再回頭了.
跳高 又說甚麼增高.

$^{881}$後來都放棄了.
因為你的腫是這麼矮的.
所以每個人到底怎樣的呢?.
那個身體將來是怎樣的呢?.
但是榮耀.
你可以想像到.
那個神聖的能力有多麼大.
接著以後活著還存.
即是未死的那一代人.
都要用身體改變來復活.
因為我們要被提到去雲裡.
不是坐無人機.
不是那些.
是被提到去雲裡.
是你自己會騰雲駕霧上到天上.
所以我們覺得是很匪夷所思的.
超自然的科幻片.
但其實這個是聖經二千年前已經的應許.
是保羅得到的啟示.
那一刻 耶穌回來那一刻.
我們就很有興趣.
什麼時候那一刻.
哪一年哪一月哪一日.
但是聖經就從來不說那個時間.
但是說那一刻.
耶穌一回來.
就會有什麼情景呢?.
就有那些聲音.
有神的號.
有天使掌.
諸如此類.
接著所有死了的人.
立即復活.
所有.
無論你怎樣死.
所有的都復活.
你想想全世界出名的基督徒很少.
不出名的就更加多.
所有都要復活.
衛斯理諾罕也復活.

$^{921}$馬丁路德也復活.
但是可能很普通的一個婆婆.
住老人院的.
在錦山洗禮的.
她什麼都不懂的.
都復活了.
都是有那個榮耀身體.
你可以想像得到.
這件事是多麼的.
你想都覺得奇妙.
我們很少想到這樣的規模.
我們很少想到.
整個宇宙.
到時會有這麼榮耀的一個國度的顯現.
因為這班人不是一些烏合之眾.
曇花一現.
而是永恆的復活.
大家有一個榮耀的身體.
去到聚集.
去到耶穌的面前聚集.
在天上見到耶穌.
你可以想像得到.
古往今來的每一個聖徒.
整個人生.
就是這樣去到聚集裡面.
在天上提到雲裡.
在空中與主耶穌相遇.
當然包括那些天使.
包括那些已經死了的復活的人.
你會見到所有的人.
昨天在英國的一個弟兄.
他移民到英國.
他說突然間原來他家附近.
就是維斯利約翰的墳墓.
紀念碑.
拍給我看.
等等.
他好像突然間尋獲至寶.
原來在他家附近就是.
他很興奮拍給我看.

$^{961}$原來也是.
因為維斯利約翰對英國的大復興.
帶來很大的幫助.
在牛津等等.
所以這些告訴我們.
我們要提到去雲裡.
你是有份的.
你不要覺得自己關我甚麼事.
一向我在教會沒有人認識我.
一定是那些名人才會站得很前.
不要這樣想.
所有在主裡面.
安睡的人全部要復活.
去到天上.
不會漏了你的名字.
叫錯了你.
你忘記了.
不知道要復活.
沒有這回事.
你可以想像那種能力是多麼大.
是不會叫錯.
不會叫漏名.
你可以想像到.
如果你去覆診.
護士一直都叫漏了你的名字.
你一直坐著等.
等了很久.
差不多下班都還沒有叫你.
你都會著急.
但不是的.
將來不會這樣.
神的號吹響.
有一種能力就是.
任何的基督徒.
他都會記得.
是會有一種能力進入那個生命裡面.
其實已經死光了.
沒有了.
你無論怎樣想.
都是沒有了.

$^{1001}$但是可以突然間有回你自己.
屬於你自己的永恆生命.
還有你的位置在天堂.
你就可以想像到.
這個榮耀這個盼望.
是多麼在這個時代需要去提倡.
需要去教導.
基督教如果我們沒有對永生的渴慕.
基督教有時候就沒有那種能力.
你不能夠告訴其他人.
我們基督教是很強大的.
是很有那個盼望的.
這個就是困難所在.
就是我們太少去羨慕主的再生.
去解決一切我們解決不了的問題.
戰爭,瘟疫,疾病.
因為這個地球是受咒詛的.
我們解決不了的.
你怎樣研發都不夠那些病毒的變種來得快的.
還有很多人為的災難.
我們改變不了的.
一個地震就已經可以帶來很多人的死亡.
但是這個告訴我們.
當主的再回來.
這個特別強調.
這個過程完結之後.
完結之後不是聚一聚集.
我們就回去做一杯.
是這樣我們包括死了復活和未死的被提.
所有的人.
當然那些未信的人.
這裡就沒有涵蓋.
所以就衍生了又災前被提.
災中被提的那個情況.
還有地上的未信的人.
還要在地上有那個生活的.
這個就遲些說到末世論那裡.
有幾個主張給大家知道一下.
這裡就說我們就要和主永遠同在.
和主永遠的同在.

$^{1041}$當然也不能夠說現在你不愛主你不追求.
將來就很甲甲不柔.
很生矮很甲膜.
似乎將來都不會是這樣的.
但是無論如何.
你要記得我們進入永恆.
其實就是和這位主永遠的同在.
在天堂裡面.
你還要吃點東西.
因為有點乾.
好好喝的.
稍微等等.
抗泉水.
好好喝的.
是甚麼水.
可能是這八個星期.
天然的.
天會活著.
抗泉水.
很舒服的.
很後悔.
我們的結局是永遠同在.
被提不是一個只是很超然的一個奇觀.
所有人復活去到天上.
去到天上只不過是一個聚集令.
然後要跟隨耶穌進入永恆的天家.
天堂.
我在星期一洗完禮.
星期二我姐姐問我在醫院姐夫的床邊.
阿貴可不可以這樣跟她的丈夫說.
我說說甚麼.
你如果離開世界.
你會見到你的女兒.
我說當然可以.
因為Lily她的大女兒.
是九年前死的.
我又要趕過去幫忙辦喪禮.
28歲就死了.
現在就葬在她的旁邊.
所以星期二姐姐這樣跟我說.

$^{1081}$在姐夫的耳中說.
我說你離開這個世界.
你就會見到Lily.
見到耶穌.
星期二那天她是回光返照.
是可以張開眼睛看著自己的兒子十分鐘.
張開眼睛.
星期一只是閃.
叫她閃一閃.
叫她又再閃一閃.
是有反應給你看到.
但星期二是很明顯的回光返照.
到星期三再去探望她.
因為要一個小時車程.
她已經是完全昏迷.
看不到東西.
也沒有再眨眼.
但是她可能真的明白到這個信念.
她離開就會見到自己的大女兒.
去接她.
迎接她.
有主耶穌迎接她.
所以各位這是一個很重要的盼望.
我們是與主和我們所愛的人永遠的同在.
基督教最寶貴的地方就是.
承載人的生命結束之後的歸宿.
雖然我們很少人會突然間.
每天都想著何時死.
沒有人這樣想的.
這樣想也不健康.
但是終極的關懷.
正是基督教唯一的基督教宗教.
夠膽這樣探討和很清晰的宣告.
這是神的啟示.
這才是基督教最重要的地方.
這才是真正終極的關懷.
當然終極之下有很多的.
譬如入小學有婚.
結婚走教堂.
死了就能撞上基督教墳場.

$^{1121}$這是終極關懷之下的模樣.
但是我們不要將最重要的事情.
我們反而不要認真去思考.
普羅說.
所以你們當用這些話.
彼此的勸諱.
這些話就是說.
我們要對這個過程.
看為是真理和事實.
真正理解我們如何離開世界.
不是最重要.
我們要在主裡面安睡是最重要.
不要離開主.
你就會有永生.
我們的人生為主而活.
就是一生的內容.
當然每個人的結局都不同.
有些人是疾病.
有些人可能是殉道.
諸如此類.
但是無論如何.
那是一個過程.
重要的是我們的靈魂.
會進入天間.
和在主耶穌的懷抱裡面.
但那還不是最後的結局.
這裡就是有一個榮耀的身體.
復活的榮耀身體.
然後跟隨神的號角的聲音.
總之是有一種能力.
你就會復活.
我就會復活.
就去到耶穌的天上見到主耶穌.
可能是天使,天君.
我們就成為了焦點的主角.
然後就一起去到永恆.
跟隨耶穌去到永恆.
所以我們要彼此用這些真理來勸惟.
我們要有這個學習.
知道有這樣的過程.

$^{1161}$不是那個時間.
不是那個邊年邊月.
公元三千年.
還是公元四千年.
不是那個.
而是最後的那個日子.
主回來的時候.
就是那個過程.
這個就是要給到那些對於自己的親人離世.
而擔心他們現在在哪裡.
那個的一個安慰,一個勸勉.
而保羅也都帶出了那個真理.
就是對於每一個的基督徒.
我們都有這個自己的復活的身份.
我看看是否夠時間.
有足夠時間和大家揭一揭.
關主在來了.
我們就逐段去揭一揭.
《馬太福音》24章.
至少有十段.
大家稍安毋躁慢慢來.
24章是《馬太福音》.
所以絕對不是一些旁門左道.
鑽孔鑽瀨才找到一節.
這樣就有些危險.
因為如果只有一節聖經.
就想確立成為一個重要的真理.
就怕不夠保險.
但是關於主在來.
基本上任何的使徒.
包括主耶穌.
都是很清楚的宣告.
一定會有這個日子.
44節.
這裡上文是主動他們說的.
或者我們看由36節開始.
否則就不夠完整.
「但那日子那時晨.
沒有日知道.
連天上的使者也不知道.

$^{1201}$只也不知道.
唯獨父知道.
羅亞的日子怎樣」.
不說哪年哪月哪日.
但是說那個狀態.
我們大概可以知道.
那個情況那個時間.
是會怎樣的.
社會的狀況.
人心的那個情況.
這裡說.
「羅亞的日子怎樣」.
羅亞的日子即是羅亞入方舟.
預備方舟.
他要入方舟的時候.
「人子降臨也要這樣」.
也要怎樣.
第38節就解釋了.
「當洪水以前的日子.
人是照常吃喝嫁娶.
即是生活如常」.
因為神安排了人.
是要嫁娶.
是要繁衍後代.
是要透過工作耕種.
才可以有食物.
這個很正常.
這個我們理解.
人照常吃喝嫁娶.
唯一不正常的就是.
有一個像傻子一樣的羅亞.
到處說我們要入方舟.
因為就會有洪水來.
似乎當時不是經常下雨的.
所以大家都習以為常.
到底羅亞是說了四天.
四十年還是怎樣呢?.
會不會是120年呢?.
因為神跟羅亞說.
人可以的日子.

$^{1241}$第六章《創世記》.
還可以有120年.
所以雖然我們不確實知道.
是不是120年.
但是似乎羅亞那個傳的異讀.
人是看為掉以輕心.
因為我可以照常的.
記住是照常.
即是說是很正常的每日.
我們根本不會怕有洪水.
有個天窗會打開.
因為《創世記》第九章.
是天上的窗戶打開.
是這樣形容的.
我們後代的人看不明白.
這些搞笑的神話.
天上怎會有個窗戶呢?.
是雲聚集.
然後降雨的.
你懂不懂的?.
這些基本的小學常識都不知道.
天窗打開.
聖經是這樣描述的《創世記》.
你可以想像到.
那個照常是一個出事的原因.
人根本不稀罕羅亞的警告.
更加看到那個方舟.
只不過是一個笑話.
方舟是一個笑話.
耶穌是可以被人嘲笑的.
現代人就是叫做好一點.
不會用武力.
阻礙傳道的人.
即是不同以前.
但是這個照常是.
人基本沒有神的狀況.
直到羅亞進方舟.
不知不覺洪水來了.
照常和不知不覺.
你可以想像到.

$^{1281}$主耶穌這樣去描述.
這樣去解釋《創世記》第九章的經文.
人照常.
不信神的人照常.
羅亞不正常.
羅亞不斷說.
老婆心叫人進方舟.
她就不正常了.
就帶來不知不覺.
根本那些人下雨.
下雨有什麼問題呢?.
就下了.
方舟可以浮起來了.
多好玩.
羅亞你真的沒有說錯.
真的下雨了.
不知不覺.
人的心硬到.
不怕有審判和滅亡.
所以主耶穌說.
那些人是照常的吃喝.
但是不知不覺.
自己沉淪.
自己被洪水淹沒.
已經沒有辦法.
因為方舟那道門終於關了.
人就進不去了.
水一直升高.
所以這個表面已經是一個福音信息.
耶穌說的這裡.
不過是和那十二門徒說.
因為這個時候.
在鑒嵐山.
鑒嵐山我們去過聖地.
我們知道.
就是剛剛看著對面.
就是聖城.
就是聖殿.
就是那個金碧輝煌的.
全世界第三的清真寺.

$^{1321}$鑒嵐山就是剛剛看著.
一個小的山丘.
看著聖殿.
耶穌就是和那些門徒.
那些門徒就問耶穌.
你的降臨和這個世界的末後.
是怎樣的?.
那耶穌就說了.
你就和他們說.
所以將他們全部都衝去.
人子降臨也要這樣.
說了兩次.
人子降臨也要這樣.
那是不是意味著.
今天的人都是照常的生活.
如常.
沒有什麼大的衝擊.
偶然間有的.
譬如地震就有.
但是慢慢有如常的生活.
但是最終又是不知不覺.
不知不覺.
就是說根本沒有意識.
只是圍著地上的可能是湖口.
找吃的.
找吃的之後.
一到長假期.
就蜂擁地過關.
休息.
放鬆.
按摩.
唱卡拉OK.
吃東西.
慢慢慢慢就是照常.
和不知不覺.
這個是教徒要留心.
要給未信的人留意.
四十戰.
那時兩個人在田裡.
娶他一個撇下一個.

$^{1361}$兩個女人推磨.
娶他一個撇下一個.
這麼齊齊整整.
一半一半的.
特別強調是在田裡.
這裡說.
推磨就是男性可能做戶外工作.
女性在室內做這些收割.
收成.
來將那些麥.
大麥來磨.
接著是四十二戰.
所以你們要警醒.
因為不知道你們的主.
是哪一天來到.
特別強調就是.
娶他一個撇下一個.
有信心的就被娶去.
不論你在田裡.
或者是推磨.
或者是日間或夜間.
因為第五章也會說日間夜間的概念.
這裡特別強調.
是真的會娶他一個撇下一個.
我們緊張的是.
我們是否在主裡面.
是否有份.
這是很重要的.
我們如果不是在主裡面的人.
是一個很大的警號.
接著是第四十三節.
家主如果知道幾更天有賊來.
就必定警醒.
不容人挖透房屋.
這是你們所知道的.
所以最後一個結論.
所以你們也要預備.
因為你們想不到的時候.
人子就來.
強調這裡.

$^{1401}$人子來是人想不到的.
但是要準備.
還有一些社會的狀況的信息.
比如剛才說的.
一般人只是地球的生活.
人一世民一世.
甚麼都去享受.
甚麼都去拼命去做.
無論為自己的身體.
為自己的健康.
為自己的家人.
為自己的下一代.
為甚麼都是.
但是從來沒有為自己的靈魂.
沒有為將來永恆的歸宿.
自己的存在.
所以耶穌就強調.
是會有復活.
有祂的再來.
接著就是說.
真的會娶祂一個騙下一個.
所以連我們這個.
其實是十二使徒.
所以你們要警醒.
因為你們不知道.
你們的主是何時會回來.
接著有一個比喻.
就是一般的家主.
一定是有一定社會地位.
家財的人才會警醒.
我甚麼都沒有.
做乞丐的.
你不就要偷.
你就拿.
是沒有問題的.
但是一定說家主是一個.
那個家庭裡面.
是一定有很多的財富.
他才會害怕.
那些賊.

$^{1441}$所以他們會有那個準備.
所以原來主在來.
是和我們警醒.
是會直接拉上關係的.
所以多黨主在來.
信徒就會警醒.
因為真是.
哪個人可以掌握到自己的生命.
新加坡的航空公司.
突然間會有亂流.
就會一個人死了.
引發他的心臟病死了.
就是沒有人會想到.
這些的意外.
是沒有人可以預測的.
有些事我們可以防備的.
你不要開快車.
不要來到去設線.
但是有些事就算我們很安全.
都不代表我們一定平安.
但是這裡主耶穌就說.
所以你們要預備.
要預備.
因為你們想不到的時候.
人子就來.
是想不到的.
是意料之外.
才叫意外.
所以這段就是說.
不需要計算那個時間.
但是要留意身邊的人事物.
羅亞這樣存異到.
沒有人感覺興趣.
只有他一家人.
八個人人.
方舟這麼大隻.
居然沒有一個的乘客.
是他自己的同胞,村民.
這個是很可惜.
但是責任不在羅亞身上.

$^{1481}$他已經盡力做好他的工作.
但是因為人.
掉以輕心的就是.
他只是有人的一生.
沒有神的永生般.
所以照常吃喝嫁娶.
所以不知不覺就被洪水沖去.
是怎會不知不覺呢?.
整隻方舟不是一個幻覺.
是真的每天有人搬.
應該不只是羅亞一個人做.
可能請了很多工人幫忙做.
但是那些幫忙的人.
可能一邊做一邊笑一邊發工資.
但是從來自己不入那隻方舟.
這個是令人感覺到坎坷.
都很可悲.
很多人曾經一生在教會.
去過教會.
但是最終都是和神沒有關係.
這個是很坎坷的.
第二段.
《約翰福音》第十四章.
最後晚餐的時候耶穌說.
當時門徒是非常憂愁.
很自然的.
因為都不知做甚麼.
突然間宇宙心腸.
好像自己近乎離流.
說這些話.
耶穌都摸不著頭腦.
然後就像打暗號一樣.
給你一些餅,油帶.
你吃就甚麼了.
所以彼得都叫約翰靠近耶穌.
誰出賣你呢?.
耶穌說我給那個人一點餅.
那個人就出賣我了.
其實十二門徒是很害怕的.
當時外面的風聲都說.

$^{1521}$那些人會抓拿耶穌.
置他於死地.
現在在內室裡的最後晚餐.
其實氣氛是很凝重的.
第十四章第一節.
你們心裡不要憂愁.
你們不信神.
先說上面.
因為第36節.
西門彼得就問耶穌.
十三章的36節.
主啊 你要去哪裡啊?.
耶穌就說.
我所去的地方.
你現在不能夠跟我去.
後來就要跟我去了.
越聽越不明白.
你欲言又止.
不如你說白一點.
為甚麼我現在不可以跟你去啊?.
我為你死都可以啊.
很真誠的.
已經按他的愛主的能力去說了.
就是說你要去哪裡.
我一定會跟你.
甚至如果你的性命受到威脅.
我願意為你的緣故.
犧牲我的性命.
第38節耶穌就說.
你願意為我死命啊?.
我實實在在的告訴你.
雞叫而先.
就是天光,雞啼.
你就要三次不認我.
你這裡這麼有那個客人.
這麼有心來跟你共赴一刻.
現在你就潑我冷水.
你到天光那些雞還沒叫完.
你就要否認我.
下毒誓.

$^{1561}$這個意思.
耶穌這樣來跟彼得這樣說.
好像他一切所做的.
都真是在主耶穌的眼裡是不合格.
disqualify了他的表忠.
對主耶穌表達的忠誠.
這樣是很黑和很尷尬的.
當面這樣硬邦邦的.
在12個門徒面前這樣說他.
當然有一個猶大已經出去了.
出去了找祭司長的.
11個門徒.
你想想多尷尬,無地自容.
你要三次否認我.
這樣跟彼得說.
所以你想想.
當時的氣氛是極之沉重.
在樓房上.
最需要安慰的其實就是耶穌.
你們心裡不要憂愁.
你們信實也當信我.
在我父的家裡.
有許多住處.
若是沒有.
我就早已告訴你們了.
所以記住是天家.
不是回到天上.
你喜歡你不要理我了.
我飛遠一點走來玩玩.
不是好像孫悟空那樣.
打幾個關斗就去到有保雷宮那麼遠.
不是的.
是去到天家.
是父的家.
耶穌特別強調.
這是父的家.
有許多住處.
是為我們預備的.
我去完事為你們預備地方去.
我若去為你們預備了地方.

$^{1601}$那多久呢?.
我們不知道了.
但耶穌在天上為聖徒禱告.
希伯來書9章說.
現在祂在天上為我們預備.
或者等候神最終末的時間.
神沒有啟示給我們知道的時間.
而是我們人類在地上.
不斷地猜測哪年哪月哪日的時間.
但那個時間到了.
真是按照地上的時間.
有一天是到了.
但耶穌就會回來.
我若去為你們預備了地方.
就必再來接你們到我那裡去.
我在那裡.
叫你們也在那裡.
我往那裡去.
你們知道.
哪條路你們也知道.
這就是耶穌基督在樓房的勸勉.
我在那裡.
叫你們也和我在那裡.
就是這個父的家.
所以是一個確確實實永恆的天家.
我們要進入的.
所以這是我們每一個基督徒.
共同一定會有的盼望.
接著到希伯來.
先到《小道行傳》.
《小道行傳》的第一章.
耶穌升天之前的景象.
當時某程度上是天使來責備門徒.
《小道行傳》第一章第六節開始.
他們聚集的時候.
他們就是十二使徒和另外的跟隨者.
他們就在橄欖山聚集.
他們聚集的時候.
問耶穌說.
主啊.

$^{1641}$你復興以色列國就在這時候嗎?.
當時耶穌已經向他們復活.
顯現四十天了.
耶穌現在升天了.
準備升天了.
他們關心的就是.
你復興以色列國是不是在這時候呢?.
耶穌對他們說.
「復行著自己的權柄.
鎖定的時候日期」.
不是你們可以知道的.
但是聖名降臨在你們的身上.
你們就必得著能力.
並要在耶路撒冷.
猶太傳地和撒瑪尼亞直到地極.
作我的見證.
他們很緊張以色列的復興.
但是耶穌緊張的就是.
你們不知道我的再回來是何時.
但是很快就是.
聖靈會降臨在你們每一個人的身上.
你們就會得著聖靈所賜的.
似乎是全港耶穌是復活主的能力.
不是行神蹟的能力.
行神蹟只不過是印證真理.
耶穌是復活的救主.
所以行神蹟可以有可以沒有.
如果神要我行我也可以行.
如果我的祈禱真是實聽.
而我的姐夫真是見到大光.
某程度神藉著我.
來令到姐夫見到神蹟.
但是重要的不是神蹟.
所以我們不是追求神蹟.
是追求無論用什麼方式.
去表達真理的能力.
是那個真理的能力改變人心.
所以你們必得著能力.
是說有勇氣夠膽量.
在耶穌都不在的時候.

$^{1681}$你們十二個使徒.
能夠繼續為耶穌作福音見證.
這個是聖靈的能力.
所以我們今天追求的是.
講福音的能力.
在任何的環境.
講明白給人知道.
福音的大能可以改變生命的能力.
所以就在不同的地方.
這裡說有了這個能力之後.
就要在耶路撒冷,猶太,全地,撒馬利亞.
直到地極.
很清楚了.
這個能力是為耶穌作見證.
記住不是追求某一種你自己羨慕的能力.
我很想唱歌.
我很想唱得很好的.
我要追求這種.
是那個唱歌的能力.
是為著耶穌作見證.
不是炫耀你的口才.
炫耀你的聲線.
炫耀你的歌聲的美麗.
是要榮耀主耶穌的救恩.
記住不要搞錯了.
任何一個人.
無論你什麼恩賜都好.
都是為著要作耶穌的見證.
否則你始終都會不成功.
因為這是你的心結.
也是一個錯誤的扭曲.
神給你的恩賜.
你為自己去炫耀.
我多厲害.
我講道理多出色.
那你就搞錯了.
是給那個能力.
你是為耶穌作見證.
所以說了者說.
他們正漢的時候.

$^{1721}$他就被取上升.
他就是耶穌.
有一朵雲採把他接去.
真是好像孫悟空騰雲掛霧.
有一朵雲採來將他接去.
將主耶穌接去.
如果我們看著那個情景.
一路升上去.
隨隨升上去.
不是輕易地升.
是有個雲載著耶穌走.
接著看不到.
因為越來越遠.
接著第十節.
當他往上去.
他們定睛望天的時候.
似乎那些未信的人就看不到.
是這班仕途看得到.
忽然有兩個人身穿白衣.
站在旁邊.
站著不是半空.
和他們一起.
有多恐怖突然間出現.
穿白衣的.
肯定是天使.
加里里人.
特別這樣強調.
他們是加里里人.
即不是耶路撒冷人.
即是你們這些是漁夫.
你們現在跟隨耶穌.
成為耶穌將來的見證人.
你們成為獨當一面的使徒.
但在神的眼中.
其實你們本身是加里里漁夫.
你們為甚麼站著望天呢?.
所以我們如果不為耶穌作見證.
最終榮耀自己.
其實只是加里里人.
我們在神的眼中.

$^{1761}$你多厲害都是微不足道.
所以我們如果能夠盡我們最好的恩賜.
為著的訊息就是為主耶穌作見證.
人跟隨了耶穌.
不稀罕你的講道.
有甚麼所謂呢?.
因為那個講道帶他看到耶穌.
耶穌得到榮耀.
得到那個人的心.
就是那個的重點.
所以人生應該是這樣.
加里里人你們為甚麼站著望天呢?.
還繼續想望到耶穌再回來.
「這離開你們被接升天的耶穌.
你們見他怎樣望天上去.
他還要怎樣呢?.
又是「駕著魂降臨」.
就是《切剎萊利加前書》四章.
剛才我們讀的經文那裡.
主耶穌又是這樣回來.
所以我們今天不是要站著在橄欖山.
等耶穌回來.
因為上文講的是甚麼呢?.
「精靈會降臨在你們身上」.
你們就會得著為耶穌作見證的能力.
在哪裡呢?.
在不同的地方.
所以我們神安排你在哪裡.
你可能有些人在地極.
有些人在耶路撒冷.
當時沒有飛機.
當時沒有航運,沒有高鐵.
他們是用腳,可能用雷仔去走的.
可以走多遠?他們去到多遠?.
譬如多瑪居民就去到印度.
他們每個人都有不同的.
譬如使徒彼得.
他們要留守在耶路撒冷.
所以他們將來會各自精彩.
按聖靈的帶來的.

$^{1801}$因為是聖靈的差遣.
所以我們的人生都是.
所以如果基督徒信了主耶穌.
只是望天等主耶穌回來.
就忘記了上文所講的.
「但聖靈降臨在你們身上」.
所以今天我們應該要.
在事奉裡等候主的再來.
在人生為主耶穌有能力的.
聖靈的能力的裝備,辯證.
去等候主再來.
這樣等候主再來.
就不會被天使責備了.
那兩個人就責備他們.
還望天?他怎樣升上去.
他就會怎樣回來.
即是你們現在開工要去傳福音.
大概是這個意思.
因為真是羨慕的.
你看著主耶穌升天.
和你一起40天的主耶穌.
是你們的主.
祂真是從死在墳墓復活.
40天同在不斷講關於神國的事.
現在終於有一天要分開了.
真是眾目睽睽看著祂升天了.
真是多看兩眼.
但不知道怎樣的.
但天使說開工了.
聖靈要降臨在你們身上了.
你們要得著能力.
要在你們的人生為耶穌作見證.
這個就是《小丑行傳》第一章開場白.
還有希伯來書.
還有7分鐘.
我們還有8段經文未讀.
希伯來書的9章.
希伯來書9章的.
似乎應該是27,28節.
那裡講審判的.

$^{1841}$9章.
希伯來書9章.
由23節開始.
23節開始.
照著天上樣式作的物件.
必須用這些製物去見證.
但那天上的本物.
自然當用更美的製物.
是在講主耶穌成為贖罪祭.
相對於地上的公牛的血.
那個獻祭.
譬如他們在贖罪日.
譬如他們在就職典禮.
祭司就要獻祭.
有一個制度就是.
那些祭司代表著那些獻祭的人.
贖罪.
這裡告訴我們.
照著天上樣式作的物件.
即是地上的聖殿.
是複製天上真正的聖殿.
必須用這些製物去見證.
但那天上的本物.
自然當用更美的製物去見證.
因為基督並不是進了人手.
所造的聖所.
這個聖所.
地上的聖殿.
不過是真聖所的影像.
即是一個複製的.
是一個真正的複製.
天上是怎樣呢?.
地上是神吩咐摩西造的.
所以是要反映天上的實體.
那個真實.
所以這些我們都很難理解.
將來去到天堂.
我們就會明白.
那個完全的本像是怎樣.
接著就說.

$^{1881}$因為基督並不是進了人手所造的聖所.
乃是進了天堂.
《聖經》只有這裡講天堂.
真正的.
真是在講天上的榮耀的聖殿.
如今為我們顯在神面前.
耶穌現在在天堂.
在神的面前.
永生神的面前.
而且是為了我們.
即是為我們守望祈禱.
等候神的時間.
祂就會回來了.
真是不是一個物質界的.
光年的宇宙.
是永恆屬靈的存在的天堂.
接著就是第25節.
也不是多次將自己獻上.
像那大祭司妹.
就是在說贖罪日那天.
帶著牛羊的血進入聖所.
其實這裡是想和地上的祭物.
和耶穌的犧牲做一個比較.
到底耶穌的犧牲帶來甚麼改變呢?.
一個徹底的改變呢?.
所以現在做一個比較.
第26節.
如果這樣.
他從創世以來.
就必定要多次受苦了.
因為大祭司是每年贖罪日.
去贖罪.
獻祭.
是要令到地上的猶太人.
每年都想起自己.
其實是虧欠神的.
全個民族都要贖罪的.
這樣的意思.
如果這樣的時候.
如果基督都要這樣做的時候.

$^{1921}$不是每年都要死嗎?.
這是不可能的.
因為他是永恆的主.
根本不需要這樣.
所以這裡作者就說.
如果這樣.
他從創世以來就必多次受苦.
但是如今.
就是主耶穌.
在這個末世顯現一次.
將自己獻為祭.
就是在十字架.
各國他山.
很除掉罪.
這個就很清楚.
說到主耶穌的流血犧牲.
就帶來一個永恆的救贖.
不需要不斷不斷地死.
一次就已經夠了.
27節.
按照著定命.
按照著神的命定.
某程度上是一個人類的律.
人人都有一死.
死後且有審判.
人人只有一次死.
不會有九條命.
不會死九次.
每個人只死一次.
這是物質世界.
人類一次.
所以我們要怎樣死都要求問神.
人人都有一死.
死後且有審判.
所以像這樣.
基督既然都是人.
他也只能夠或只需要死一次.
就是十字架各國他的犧牲.
所以像這樣.
基督都是人性.

$^{1961}$在這個世界.
像這樣基督既然一次被獻.
擔當了多人的罪.
現在又在榮耀的神的寶座右邊.
為我們預備.
所以當他再回來的時候.
將來要向那等候他的人第二次顯現.
並與罪無關.
乃是為拯救他們.
作者在這裡特別強調.
屬罪的罪已經獻過了.
主耶穌已經擔待眾人的罪.
而且因為他是永恆永生的主.
他只需要死一次就已經解決了.
一個叫Once and for all.
一次性就已經完成了.
不需要多次多次的受苦.
要這樣死的了.
所以十架的救恩是永遠有功效的.
所以這又是一個問題.
我們在人生.
很多機會會讓我們掉以輕心.
基督的十架救贖.
有時候有很多大型的挫折.
可以使我們與基督隔離.
我們慢慢就不珍惜.
或者覺得不真實.
但是這裡特別強調.
基督是一次已經完成了.
他現在在天上.
將來就要向那些等候的人.
等候他回來的人.
第二次顯現.
不再是追究你的罪.
而是拯救我們.
我們終於可以脫離地上物質.
還是受罪.
受咒詛的世界的驅控.
我們進入真正永恆的自由.
是這一種拯救.

$^{2001}$但是我們已經有靈魂救恩.
我們已經與耶穌一次性的犧牲.
是完全認同.
所以一次得救.
永遠得救是錯的.
我們是一次相信.
一直相信.
就必然得救.
是要一次相信.
一直相信.
才可以得救.
就是保守自己在神的救恩.
救恩只有一個.
一次性已經完成了.
就是主耶穌在十架的犧牲.
我們與祂有生命關係.
認同祂的死.
是我們的贖罪.
因為祂就是贖罪制.
我們認同這個贖罪制.
我們就永遠在我們的生命裡面.
有這個救贖的功效.
所以將來就是.
我們等候祂的再來.
與罪無關.
那些不信的人.
就與罪有關.
他們將來就要面對審判.
因為他們拒絕.
承認耶穌的十字架的贖罪制.
所以第九章是想強調.
主耶穌的這個贖罪.
是一次永遠有功效的.
所以對一些.
如果有懷疑自己的信心不夠的人.
這些經文都是很重要的.
因為看的是耶穌.
不是我們自己.
我們之後就講到這裡.
下星期再繼續.

$^{2041}$看一些新約有講到.
主回來要具備什麼心情.
那個狀態的基督徒.
禱告一下.
多謝主.
盼望主的再來.
成為我們大家共同的學務.
我們在人生.
應該真的要有聖靈的能力.
去為耶穌作見證.
偶然抬頭望天是必須的.
但是這個抬頭望天.
正是給我們有一個永恆的盼望.
去在地上因勤為主.
作福音的見證.
奉耶穌基督的名.
阿們.
感恩.
\newpage



\section{以弗所書 4:17-32}
\label{sec:7PHC2DLBPCo}
\textbf{2024-06-01 晚堂《黑暗與光明》以弗所書4\_17-32 講員:曾裔貴牧師}
\newline
\newline
連結: \href{https://youtube.com/watch?v=7PHC2DLBPCo}{\texttt{https://youtube.com/watch?v=7PHC2DLBPCo}} ~~~~ 語音日期: 2024-06-04
\newline
\newline
\hyperref[sec:eejVMLNyY8I]{\small{< < < PREV SERMON < < <}}
~
\hyperref[sec:index_chronic]{\small{[返順時目]}}
~
\hyperref[sec:index_scriptual]{\small{[返順卷目]}}
~
\hyperref[sec:Il1UURetEws]{\small{> > > NEXT SERMON > > >}}
\newline
\newline
以弗所書 4:17-32
\newline
\begin{longtable}{cl}
\hline
\hline
章節 & 經文 (和合本修訂版)\\
\hline
4:17 & \begin{tabularx}{0.7\textwidth}{X} 所以我這樣說，且在主裡鄭重地說，你們行事為人，不要再像外邦人存虛妄的心而活。 \end{tabularx} \\ \\ \relax
4:18 & \begin{tabularx}{0.7\textwidth}{X} 他們心地昏昧，因自己無知，心裡剛硬而與神所賜的生命隔絕了。 \end{tabularx} \\ \\ \relax
4:19 & \begin{tabularx}{0.7\textwidth}{X} 既然他們已經麻木，就放縱情慾，貪婪地行種種污穢的事。 \end{tabularx} \\ \\ \relax
4:20 & \begin{tabularx}{0.7\textwidth}{X} 但你們從基督學的不是這樣。 \end{tabularx} \\ \\ \relax
4:21 & \begin{tabularx}{0.7\textwidth}{X} 如果你們聽過他的道，領了他的教，因為真理就在耶穌裡， \end{tabularx} \\ \\ \relax
4:22 & \begin{tabularx}{0.7\textwidth}{X} 你們要脫去從前的行為，脫去舊我；這舊我是因私慾的迷惑而漸漸敗壞的。 \end{tabularx} \\ \\ \relax
4:23 & \begin{tabularx}{0.7\textwidth}{X} 你們要把自己的心志更新， \end{tabularx} \\ \\ \relax
4:24 & \begin{tabularx}{0.7\textwidth}{X} 並且穿上新我；這新我是照著神的形像造的，有從真理來的公義和聖潔。 \end{tabularx} \\ \\ \relax
4:25 & \begin{tabularx}{0.7\textwidth}{X} 所以，你們要棄絕謊言，每個人要與鄰舍說誠實話，因為我們是互為肢體。 \end{tabularx} \\ \\ \relax
4:26 & \begin{tabularx}{0.7\textwidth}{X} 即使生氣也不要犯罪；不可含怒到日落， \end{tabularx} \\ \\ \relax
4:27 & \begin{tabularx}{0.7\textwidth}{X} 不可給魔鬼留地步。 \end{tabularx} \\ \\ \relax
4:28 & \begin{tabularx}{0.7\textwidth}{X} 偷竊的，不要再偷；總要勤勞，親手做正當的事，這樣才可以把自己有的，分給有缺乏的人。 \end{tabularx} \\ \\ \relax
4:29 & \begin{tabularx}{0.7\textwidth}{X} 一句壞話也不可出口，只要隨著需要說造就人的好話，讓聽見的人得益處。 \end{tabularx} \\ \\ \relax
4:30 & \begin{tabularx}{0.7\textwidth}{X} 不要使神的聖靈擔憂，你們原是受了他的印記，等候得救贖的日子來到。 \end{tabularx} \\ \\ \relax
4:31 & \begin{tabularx}{0.7\textwidth}{X} 一切苦毒、憤怒、惱恨、嚷鬧、毀謗，和一切的惡毒都要從你們中間除掉。 \end{tabularx} \\ \\ \relax
4:32 & \begin{tabularx}{0.7\textwidth}{X} 要仁慈相待，存憐憫的心，彼此饒恕，正如神在基督裡饒恕了你們一樣。 \end{tabularx} \\ \\
[1ex]
\hline
\hline
\end{longtable}
$^{1}$我們同心祈禱.
我們能夠在神的家裡面.
剛才我們以聖餐紀念主十家的大愛.
我們是同樣有著永生生命的弟兄姊妹.
我們真的很盼望.
我們能夠在主的救恩裡面.
再一次思想我們的生活.
是怎樣能夠反映救恩帶給我們的改變.
不單止是我們自己要發現.
而是要讓未信的朋友.
看到我們的生命改變.
將榮耀歸給我們的神.
今天晚上我們看.
《使徒保羅》這一段經文.
願主讓我們有聖靈的引導.
我們能夠明白.
其實很多都是很淺白的教訓.
是我們需要在生活裡面去具體的執行.
不是甚麼高深的神學.
而是我們能夠讓聖靈不要擔憂.
在我們的內心裡面.
今天晚上我們等候在你的面前.
同心禱告奉耶穌基督的名.
阿們.
用這個題目相當明顯.
是有負面的沒有神生命的人.
與保羅由四章的十七節.
一直延展到五章.
都是教導基督徒.
既然你已經有神的生命.
你要怎樣行事呢.
所以黑暗是那些沒有神生命的人.
四章剛才我們沒有讀到的.
就是第十七節這一段這樣說.
他說「所以我說」.
「而且要強調不是他個人見解」.
「是在主裡面我確實這樣說」.
「你們的行事就不要再像那些外邦人」.
「傳一個虛妄的心來行事」.
他們是怎樣虛妄呢.

$^{41}$他們是心地昏迷的.
為甚麼會心地昏迷呢.
是在說那個「mine」.
那個知識對真理的否定.
為甚麼呢.
保羅就說了.
是與神所賜的永恆生命.
是隔絕的.
原來這些叫做外邦人.
這些叫做沒有神生命.
活在虛妄 活在黑暗裡面的這樣的人.
即是說充斥著在社會.
在每一天的生活.
而不所教會的弟兄姊妹都會遇見.
都會可能跟他們做買賣.
可能在生活裡面會有交往.
但是保羅提的就是.
你不要再學他們這樣.
即是以前是.
每個人以前都是.
現在不是.
因為有了主的生命.
然後再說他們是怎樣與神的生命隔絕.
是因為自己無知.
心裡面剛硬.
良心既然喪盡.
就放縱私慾.
貪行種種的污穢.
接著第20節.
保羅就很正面說.
一個光明生命的內涵.
那個的發展是怎樣的.
其中就是神的道在那些人的心裡.
是會生根的.
是會一路建造.
是會反映到神的生命.
特別是提到第23節.
要將你們的心智改換一生.
可以改的.
可以換心的 基督徒.

$^{81}$心智可以改變.
為甚麼呢.
因為24節.
我們穿上了這個新人.
相信是指到上文第4章所說的.
一主一信一洗一實.
這個重要的信仰.
落實是一個基督徒.
透過聽明白了真理.
他去接受他去得到這個新生命.
留意一下這裡保羅特別帶我們去上.
人類的創造的時候.
創世記第二章的場景.
保羅說這個是新創造來的.
第4章的24節.
穿上這個新人之後.
這個新人是怎樣的.
是照著神的形象做的.
我們不就是新的阿當.
阿當不就是照著神的形象做.
創世記一章第26節.
神就說這個人.
我們要做一個人出來.
要和我們的形象樣式一樣.
所以就創造了阿當.
現在這裡告訴我們.
有一個新創造.
就是我們每一個基督徒.
都是照著神的形象做的.
有真理的人而.
和聖潔.
似乎就是25節剛才主席所讀的經文.
就是反映真理人而聖潔.
在一個基督徒的生活裡面.
原來沉根很重要.
第4章保羅告訴我們.
每一個頂子妹知道.
我們是因為透過主耶穌的犧牲.
而入門成為基督徒.
去到24節再一次說.

$^{121}$這個新創造是照著神的形象.
就是說我們每一個頂子妹.
一信耶穌.
你就是神的創造.
其實2章10節已經說了.
現在這裡再具體地說.
這個新創造.
所以今晚我們要搞清楚.
你是屬神的.
還是不屬神的.
不屬神的就是與神的生命.
所賜的生命.
耶穌基督在十誡犧牲的生命.
隔絕的人.
他有很多他們自己的觀念.
他們有很多他們自己的想法.
是與神的生命隔絕.
所以尋根是一個很重要的開始.
你是為什麼會信耶穌.
你是怎樣可以得到神的救恩.
每一次的聖餐.
再再告訴我們.
是因為曾經有一個很恐怖的犧牲.
使我們成為他的兒女.
是一個很關鍵性.
人類的救主就是耶穌.
大家都知道我上兩個星期.
要趕去多倫多.
因為我的姐夫離樓.
他要差不多離世.
很奇妙的一個星期內.
由決志信主洗禮.
安排安息禮拜.
第二天我就可以回香港坐飛機.
其中一樣東西是很令我開心.
就是見到已經長大的第二代.
有四個在加拿大出生.
完全對中國文化.
對很多的東西完全不知道.
但他們已經去到二十多歲出來工作.

$^{161}$他們的朋友全部是外國人.
說話全部是流利的英文.
但他們不約而同.
在吃解慰酒的時候.
很有興趣想知道家族上一代是怎樣的.
說對了.
你們應該要知道你們自己的根.
其中一個原因就是.
大家都知道上一代的移民.
很多都是.
特別我的家庭.
我的哥哥姐姐.
全部是低學歷.
低學識.
小學都未畢業.
所以他們長大之後.
他們的生活和家庭有很大的落差.
所以這次有個牧師來.
為什麼有個叔叔舅父會做牧師呢.
為什麼又會有這麼多東西可以懂得呢.
他們很有興趣想知道.
這種尋根.
基督徒你有沒有想想你自己的.
一主一信一洗一神.
同不同的弟兄姊妹.
是共同擁有這個永恆生命.
是神的形象所做.
這個起點.
就斷定了我們不是屬於黑暗.
我們不是與神生命隔絕的人.
不過第二方面.
保羅用了很長很長的篇章.
去講到一個有神生命.
一個照著神形象做的人.
應該是怎樣呢.
所以第20節.
你們學了基督之後.
就不可以這樣了.
有沒有留意那些措辭.
好像一個爸爸.

$^{201}$其實保羅在這個時間已經在坐牢.
他都知道自己的人生.
是差不多要退場了.
他現在是苦口婆心的勸導.
已不鎖的弟兄姊妹.
要怎樣能夠在生活地上.
為主耶穌作見證呢.
第25節.
所以你們要棄絕謊言.
17節.
所以我說.
在主裡面確實的說.
有沒有留意到那個心情.
措辭是在在一個.
很想弟兄姊妹生活裡面.
表現出你自己是有屬天的生命.
這件事很重要.
在這個月的12號.
即是6月12號.
當然還未到父親節.
有一位都算是我很敬重的老師.
魏智健老師.
他安息主懷一周年.
為何突然間會提起他呢.
因為他很年輕.
比我還小.
就有突然間的癌症.
很快就離世.
但這個病沒有令他身體很差.
他有足夠的時間.
來退場.
叫了他的親人.
他說我不知道為什麼.
我明明很健康.
不知道為什麼突然間有這個病.
要令到我知道一定要離開世界.
回天家.
很多事情還沒有做完.
他說你不要問為什麼.
我也不知道為什麼.

$^{241}$還有要禁止問為什麼.
因為將來會知道.
你們要做的.
他勸每一個家人.
勸每一個來探望他的.
教會的弟兄姊妹.
不要問為什麼.
要問的就是.
怎樣更加認識基督.
怎樣能夠在有限的人生.
去服侍主.
為什麼我會提起這位危老師.
他不是我的關心小組的老師.
是我們另外一班的同學.
同班同學我們分了兩班.
有兩位關心小組的老師.
會和我們一起.
在生活裡面和我們相交.
帶領我們.
另外一班同學就看這個危老師.
是他們的屬靈的爸爸.
所以他有一個綽號.
在神學院裡面.
就叫做Timba.
就是真的我的同學的老師.
和好像父親一樣.
為什麼我會對他.
有這麼深刻的印象呢.
上幾個月和Jason傳道.
聊開WhatsApp.
間接了解危老師的生平.
他的為人.
Jason不是見到神學院的.
吳傳道現在在溫哥華.
他說原來當年他結婚.
就是魏健志老師.
做他們的婚姻輔導.
輾轉可能是間接的認識.
當Jason傳道.
去年移民去溫哥華讀書的時候.

$^{281}$魏老師知道他們移民.
主動來探望Jason傳道.
來幫助他們.
無論融入教會.
找一些生活的安排等等.
對Jason傳道來說.
好像一個這麼窩心的售後服務.
其實Jason傳道沒有和他說.
他移民去溫哥華.
而是魏老師知道之後.
主動去找Jason傳道.
我聽到Jason這樣和我說.
我就對魏老師的為人.
更加敬重.
因為基本上不是有很深的認識的相交.
都尚且這樣主動去探問.
這位傳道來到溫哥華之後.
你的生活是否能夠安頓呢.
可以這麼細微的關心.
弟兄姊妹.
保羅在這裡很重的語氣.
那些沒有神生命的人是隔絕的.
隔絕即是說.
他的有限的智慧.
以為可以過到自己一生.
可以是很富有.
可以是很有學問.
可以是很高的知識.
但是和誰隔絕呢.
與神所設的生命隔絕.
這是很大的一件事.
你沒有了地上的生命之後.
神就要追究你的生命.
我們在救恩裡面.
已經是新創造.
所以在我們的生活裡面.
保羅就有東西要教了.
而且是很直接的教.
所以你們不要像以前那樣.
這是20節.

$^{321}$你們學了基督卻不是這樣.
接著第22節說.
就要脫去你們從前行為上的舊人.
從前的生活的舊有觀念.
舊有觀念帶來的那些行為.
我去多倫多的Downtown.
我的外甥女告訴我.
那個地方.
你知道是什麼地方嗎.
原來是吸大麻的地方.
是很多人的.
是很普遍的.
你一出火車就聞到.
在Downtown就聞到大麻味.
是合法的.
你可以想像得到.
保羅在這裡告訴我們.
我們要小心.
我們已經是新創造的人.
但是有更多人是與神生命隔絕.
我們應該怎樣逃避呢.
這個舊人是因私欲的迷惑.
漸漸地變壞.
所以我們要將我們的心智.
要改換一生.
接著就25節開始.
差不多你這樣讀.
你會感覺到那種壓迫感.
因為你們要棄絕謊言.
各人與鄰舍說真實的話.
豈不是聖經24節所說.
神創造的新生命.
在我們裡面有什麼.
真理的仁義和聖潔.
立即就說我們的口.
要怎樣說真話.
無論在任何情況都不可以說謊.
讓我騎騎肚子.
能否有力量說謊.
保羅說不是.

$^{361}$要棄絕謊言.
接著呢.
不可以生氣.
直到日落.
真的有個鐘可以直到日落就不生氣.
相信跟下面說的.
不要讓魔鬼留地步.
就是我們內心.
積蓄別人的恨.
對別人懷有一種.
不寬恕的態度.
一直不肯解除那個冤心.
這是人類與生俱來都有的.
但原來信了主耶穌.
有了這個新生命.
保羅也告訴我們.
我們不可以這樣對待其他人.
我們自己也需要有這種.
在公義下.
但不會成為我們的憤怒.
一直在我們內心積蓄.
所以你會看到.
保羅去到後面31節.
也再一次逐樣拆解.
什麼叫憤怒呢.
什麼叫內心的憤怒會積蓄.
會爆發呢.
很多的形容.
對一個人的內心.
一切的苦毒.
一切的惱恨.
憤怒.
讓罵.
毀旁.
並一切的惡毒.
都要在你們中間除掉.
立即cut off.
不可以有這樣的為人.
這樣的言說.
我們的言說要這樣的.

$^{401}$要帶著一份溫柔的調和.
你想想一個基督徒多麼美麗.
今早有一位教友.
她和她未信的先生.
開車去停車場.
繞了幾個圈.
趕時間要交文件.
她先生就發火了.
男士發火.
通常要麼就拍桌子.
要麼就說髒話.
她在這個時候立即祈禱.
她分享給我知.
她說她以前都會很緊張.
都會很生氣.
為什麼會這麼不幸運.
一個車位都沒有.
祈禱之後.
出去停車場.
對面街剛剛有個位.
有車走.
原來一個基督徒.
可以這樣來到.
不斷的內心.
有這種愛帶來的人而聖潔.
自己內心聖潔.
你想想多難.
對人一定會有意見.
有不同的判斷.
不同的判斷.
但可以判斷之後.
不帶有仇恨.
這個就是基督徒.
這個就是神的新生命.
可以判斷.
對很多事物應該有研判.
研判之後.
仍然保有一份內心的純潔.
真理的人意和聖潔.
這是其中一件事.

$^{441}$還有就是.
原來已不鎖教會的人.
這麼喜歡偷東西.
從前偷東西的不要再偷了.
總要勞力親手作正經事.
你就可以有餘分給缺少的人.
這麼直接的教導.
不知道收到這封信.
在教會中中毒的時候.
那些曾經偷東西的人.
會怎樣看呢.
這裡意味著.
原來基督教是接納罪人.
實實在在的罪人.
坐過牢的.
都可以信主.
去年有機會和劉恩德牧師聊天.
真的聊天吃飯.
他說他在安儲中學.
這麼多十年裡面.
都收留過兩個坐過牢的朋友.
但是信了主.
做校工.
非常用心去做.
很多工程打點得很好.
很低調地去做.
沒有其他老師知道.
個別一兩位的主任.
會知道他過去的背景.
剛好又和富哥聊天.
同一張桌子.
有時他們有些授課員.
盡量很低調地幾個人知道.
是做娼妓.
當然是從良了.
你可以想像到.
基督教這一份這麼實在.
對人性的信任.
基於主救恩的愛.
但從前偷竊的.

$^{481}$現在不要再偷了.
所以教會是接納罪人.
但教會傳遞的十價救恩.
可以改變人的生命.
親愛的弟兄姊妹.
我們錦宣家的弟兄姊妹.
可能你的生活.
未需要你去偷騙拐騙.
貪小便宜.
走過白街.
帶包朱古力落袋.
我們比較少這樣的情況.
有些基層的.
基於他們內心的貪婪.
可能會這樣做.
但這裡告訴我們.
在這一段經文.
大部分不是要我們去思考.
可不可以這樣做.
或者我祈禱.
神給我有力量不要做.
是斬釘截鐵.
不准 不准 不可以.
因為你有新生命.
這樣的教導.
盼望弟兄姊妹我們一同.
在神的裡面去學習.
偶然間看到一個故事.
是明報的記載.
在英國一個小地方.
有一位老闆.
有一天看到一個年青人.
在他的店舖徘徊.
他開一些雜貨店.
很好奇走過去.
他很快就走了.
他走了之後.
留下了一封信.
他打開了一封信.
這封信是很工整的英文寫的.

$^{521}$講述他過去年青的時候.
在店舖裡面偷了一個ice cream.
然後沒有給錢 走了.
後來他良心發現.
所以他寫了一封信.
道歉.
附上一張十鎊的英鎊.
這個老闆就很開心.
感覺很滿足.
因為有一個人沒有被人捉到.
而甘願自己去承認自己的錯.
向這個老闆認錯.
和承諾以後不會再犯.
其實不認識他的.
沒有人知道的.
也沒有閉路電視的.
但是這個年青人.
願意這樣去和過去一刀了斷.
離開自己這個.
當年可能是小朋友讀書的階段.
來犯下的錯誤.
所以這個老闆看到這封信.
真的很滿足.
人的生命改變.
親愛的兄弟姐妹.
教會豈不是不斷見到.
不同的生命改變嗎.
剛剛我們在紅恩堂的初信班.
有一個家庭是年青的夫婦.
和小朋友一起上初信班.
是泰國語的.
很認真學習真理.
很感恩.
真的看到生命改變.
兄弟姐妹.
我們是陪伴著.
很多人的生命改變.
昨天晚上我們出席安息禮拜.
慰安牧師的爸爸.
陳國榮弟兄.

$^{561}$他的人生很簡單.
但是他所表達出的信仰.
在生活裡面.
三十年沒有試過一次的意外.
又拿了三十年.
全香港的九巴司機裡面.
只有兩個人.
他是其中一個.
是非常之一個很實的見證.
他將主耶穌的救恩.
可能未必用真理講出來.
但就是在生活裡面.
愛神 愛教會 愛家人.
和愛敬重他自己的職業.
親愛的兄弟姐妹.
這段聖經是在講.
黑暗的 無神生命.
與神生命隔絕的人.
但講很少.
只是很概括地講了.
但講信主的 有神生命的.
就講很多.
還有從前是這樣的.
現在不應該了.
即是說原來我們所有人.
其實本來就是屬於黑暗.
現在屬於光明.
是主救恩的榮耀.
我們一起禱告.
多謝主.
你就是這樣將生命改變.
多謝主的聖靈.
在我們裡面來帶領 感動.
使我們能夠得著永恆生命.
願我們錦宣家的弟兄姊妹.
每一次我們的聚集.
我們的見面.
是反映我們在主的救恩裡面.
那個學習.
那個生活的力量.

$^{601}$讓人看到.
這個信仰是真 是活的.
奉耶穌基督的名.
阿們.
祝福各位.
\newpage



\section{列王記下 5:1-19}
\label{sec:Il1UURetEws}
\textbf{2024-06-02 午堂《無名的小使女》列王記下5\_1-19 講員:楊明麗姑娘}
\newline
\newline
連結: \href{https://youtube.com/watch?v=Il1UURetEws}{\texttt{https://youtube.com/watch?v=Il1UURetEws}} ~~~~ 語音日期: 2024-06-04
\newline
\newline
\hyperref[sec:7PHC2DLBPCo]{\small{< < < PREV SERMON < < <}}
~
\hyperref[sec:index_chronic]{\small{[返順時目]}}
~
\hyperref[sec:index_scriptual]{\small{[返順卷目]}}
~
\hyperref[sec:ui5VIns3kj4]{\small{> > > NEXT SERMON > > >}}
\newline
\newline
列王記下 5:1-19
\newline
\begin{longtable}{cl}
\hline
\hline
章節 & 經文 (和合本修訂版)\\
\hline
5:1 & \begin{tabularx}{0.7\textwidth}{X} 亞蘭王的元帥乃縵在他主人面前是一個偉大的人，得王的喜悅，因為耶和華曾藉他使亞蘭人得勝。他雖然是大能的勇士，卻染上了痲瘋。 \end{tabularx} \\ \\ \relax
5:2 & \begin{tabularx}{0.7\textwidth}{X} 亞蘭人成群出征的時候，從以色列地擄了一個小女孩，她就服事乃縵的妻子。 \end{tabularx} \\ \\ \relax
5:3 & \begin{tabularx}{0.7\textwidth}{X} 她對女主人說：「我希望主人去見撒瑪利亞的先知，他必能治好主人的痲瘋。」 \end{tabularx} \\ \\ \relax
5:4 & \begin{tabularx}{0.7\textwidth}{X} 乃縵去告訴他主人說，從以色列地來的女孩如此如此說。 \end{tabularx} \\ \\ \relax
5:5 & \begin{tabularx}{0.7\textwidth}{X} 亞蘭王說：「你可以去，我也會送信給以色列王。」於是乃縵手裡帶十他連得銀子、六千舍客勒金子和十套衣裳去了。 \end{tabularx} \\ \\ \relax
5:6 & \begin{tabularx}{0.7\textwidth}{X} 他帶著這信給以色列王，說：「現在你接到這信，看哪，我派臣僕乃縵到你這裡來，你要治好他的痲瘋。」 \end{tabularx} \\ \\ \relax
5:7 & \begin{tabularx}{0.7\textwidth}{X} 以色列王讀了信就撕裂衣服，說：「我豈是神，能使人死使人活呢？這人竟派人來，叫我治好一個人的痲瘋。你們要知道，看，這人是找機會來跟我吵架的。」 \end{tabularx} \\ \\ \relax
5:8 & \begin{tabularx}{0.7\textwidth}{X} 神人以利沙聽見以色列王撕裂衣服，就派人到王那裡，說：「你為甚麼撕裂衣服呢？讓那人到我這裡來，他會知道以色列中有先知。」 \end{tabularx} \\ \\ \relax
5:9 & \begin{tabularx}{0.7\textwidth}{X} 於是乃縵帶著車馬到了以利沙的家，站在門前。 \end{tabularx} \\ \\ \relax
5:10 & \begin{tabularx}{0.7\textwidth}{X} 以利沙派一個使者，對乃縵說：「去，在約旦河中沐浴七次，你的肉就必復原，你會得潔淨。」 \end{tabularx} \\ \\ \relax
5:11 & \begin{tabularx}{0.7\textwidth}{X} 乃縵卻發怒走了。他說：「看哪，我以為他必定會出來，到我這裡，站著求告耶和華－他神的名，在患處上搖手，治好這痲瘋。 \end{tabularx} \\ \\ \relax
5:12 & \begin{tabularx}{0.7\textwidth}{X} 大馬士革的亞瑪拿河和法珥法河豈不比以色列的一切水更好嗎？我難道不可以在那裡沐浴而得潔淨嗎？」於是他生氣，轉身走了。 \end{tabularx} \\ \\ \relax
5:13 & \begin{tabularx}{0.7\textwidth}{X} 他的僕人近前來，對他說：「我父啊，先知若吩咐你做一件大事，你豈不做嗎？何況是吩咐你去沐浴，得潔淨呢？」 \end{tabularx} \\ \\ \relax
5:14 & \begin{tabularx}{0.7\textwidth}{X} 於是乃縵下去，照著神人的話，在約旦河裡浸了七次。他的肉復原，好像小孩的肉，他就潔淨了。 \end{tabularx} \\ \\ \relax
5:15 & \begin{tabularx}{0.7\textwidth}{X} 乃縵帶著所有跟隨他的人，回到神人那裡，站在他面前，說：「看哪，我知道，除了以色列，全地沒有神。現在請你收下僕人的禮物。」 \end{tabularx} \\ \\ \relax
5:16 & \begin{tabularx}{0.7\textwidth}{X} 以利沙說：「我指著所事奉永生的耶和華起誓，我必不接受。」乃縵再三請他收下，他卻不肯。 \end{tabularx} \\ \\ \relax
5:17 & \begin{tabularx}{0.7\textwidth}{X} 乃縵說：「你若不肯，請把兩匹騾子能馱的土賜給僕人，僕人必不再把燔祭或祭物獻給別神，只獻給耶和華。 \end{tabularx} \\ \\ \relax
5:18 & \begin{tabularx}{0.7\textwidth}{X} 惟有一件事，願耶和華饒恕你僕人：我主人進臨門廟在那裡叩拜的時候，他總是扶著我的手，所以我也在臨門廟叩拜。我在臨門廟叩拜的這事，願耶和華饒恕你僕人。」 \end{tabularx} \\ \\ \relax
5:19 & \begin{tabularx}{0.7\textwidth}{X} 以利沙對他說：「你平安地回去吧！」乃縵離開他去了。走了一小段路， \end{tabularx} \\ \\
[1ex]
\hline
\hline
\end{longtable}
$^{1}$各位,我們今天看一個無名氏.
我們認識了亞伯拉罕.
認識了雅各.
認識了以撒.
我們認識了很多人.
但今天我們認識一個叫做無名氏.
因為不知道他叫什麼名字.
聖經叫他做小侍女.
就無名的小侍女.
這個小侍女是什麼人呢.
這個小侍女是以色列人.
是被亞蘭的軍隊擄走的難民.
被擄走的難民.
他難民.
為什麼不說楊姑娘說的難民呢.
今天我們就看看這個難民怎麼樣.
這個難民後來去了元帥的家裡.
做了婢女.
是一個小侍女.
她的情況可能好一點.
她是做大元帥的太太的小侍女.
不用那麼多操活.
但都應該要看大元帥和他夫人的臉色.
這個小侍女不知道為什麼她很勇.
有一天她就和她的主母說.
其實我們的大元帥.
長了一身麻風.
不如去找以色列的先知.
等先知醫好她.
不知道為什麼.
她和這個主母的關係可能很好.
主母真的把這件事和乃萬說.
因為乃萬大元帥.
其實他是當時很出名的元帥.
在亞蘭國當中.
亞蘭國的君王都很棄重他.
有什麼事他就是他的左右手.
他打了很多場勝仗.
不過他有一個問題.
他贏不了這個問題.

$^{41}$就是他身上的病.
這個病是麻風病.
不是我們現在認為是大麻風.
是一個身上長期的皮膚病.
這個皮膚病會令他傷損.
會有機會傳染給人.
不過未至於到大麻風那種.
會掉了手指.
掉了手.
但是又有傷損又很痕癢.
都是他的問題.
而且一個大將軍一站出來.
而容就受到影響.
打起仗來.
那種痕癢對他都有很多影響.
又日曬雨淋的情況.
令到他很不舒服.
在家裡有人服侍.
軍隊出去的時候更加不方便.
他不是沒有求醫.
我相信他有求醫.
不過全個亞蘭國.
沒有人治好他.
亞蘭國是不是一個小國.
所以沒有什麼名人呢.
其實亞蘭國是一個強國.
亞蘭國經常欺負以色列.
不過在這段時間.
以利沙的時候.
和以色列的關係比較好.
因為以利沙用很多不同的方法.
神已經讓以利沙知道.
軍隊去到哪裡.
怎樣就可以打敗他.
軍隊又去到哪裡.
怎樣就可以打敗他.
於是亞蘭王和以色列王.
有一段時間是和好的.
但是亞蘭不單止對以色列.
他對當時由大馬士革.

$^{81}$一直去到以色列的地方.
他都是控制得很好.
歸行於他.
受他管理的.
有三百多個諸侯地.
你想想他國境多大.
就算每一個諸侯地是小小的.
三百多個都很多.
他也控制了當時的社會.
將近一千年.
他是很大的一個地方.
在那個地方當中.
他出現的問題.
是後來他被亞述對滅了.
但他是一個很強的國家.
在那個情況下.
大將軍很重要.
大將軍是有地位的.
現在我們打開電視.
一向是誰的國家.
去訪問誰的國家.
兩位君主怎樣去點閱.
軍隊是很重要的.
大元帥是在做這些工作的.
除了領兵之外.
他時時都在皇帝的旁邊.
但他一身麻風病.
令他很無助.
一種無助令他連一個奴隸小侍女.
說的話他都要聽.
他真的很認真地聽.
這個小侍女其實是老人家.
他去找誰呢.
他去找阿難王的君王.
跟君王說.
我為夫人的小侍女這樣建議.
她是以色列人.
叫我去以色列求助.
這個乃曼的地位高到.
君王說你要去以色列.

$^{121}$那好吧 我寫封信引誡你.
叫以色列王幫你安排醫治.
一點都沒有去疑惑他.
其實小侍女說的話是真的嗎.
我還要寫封信去叫以色列醫好你.
如果成為一個笑柄那怎麼辦.
他沒有去質疑什麼.
他相信乃曼說.
真的寫封信.
好像國事探訪那樣.
去了以色列.
你知道很大陣仗.
國事那樣出現.
他帶了很多禮物去.
這些禮物乃曼帶過去.
他要讓以色列的君王.
能夠幫他安排醫治.
他帶了禮物去的時候.
他請以色列的君王為他安排.
誰知道以色列君王知道這件事後.
他的反應不是說.
好 我叫以色列醫好你.
他說我們以色列哪有這樣的人.
你分明不是來鬧交.
你是來鬧打仗.
皇帝和皇帝說了一件事.
你幫我做這件事.
是做不到的時候.
就引起國戰.
於是他就很不開心.
甚至他極度哀傷 撕裂衣服.
這件事傳到給以利沙知道.
以利沙就派了一個人.
去和以色列的君王說.
他說你為什麼要撕裂衣服.
你很哀傷嗎.
你叫將軍來找我.
以利沙對君王的信心.
有什麼對比.
大家一看就知道.

$^{161}$以利沙就安排了我.
去見大將軍.
要去見大將軍的時候.
這個大將軍帶齊他的東西.
要去見以利沙.
但是他就帶齊了他的人.
帶齊了他的物品.
要去見以利沙.
但是以利沙就發生了一件事.
以利沙就叫他做一件事.
他就叫他去約旦河沐浴七次.
去約旦河沐浴七次.
你的皮膚病就好了.
這個情況就是將軍帶齊他的車馬.
帶齊他的禮物.
整隊儀仗隊去到.
可能是小小先知.
在以色列當中.
連皇帝都不認識他.
去到先知以利沙的門口.
排隊在那裡.
他期望什麼呢.
最低限度以利沙都自己出來.
或者帶齊他的工人出來.
歡迎一下他.
不是 以利沙就叫了一個使者.
這個使者未必是他的工人.
可能是他的學生.
出來見大將軍.
大將軍 我們的主領.
以利沙沒空.
叫我來見一下你.
如果是兩間公司.
兩間公司要談一個生意的話.
大家都會是相等的職位的人去見面.
去談合約.
如果兩國見面.
都是相等的軍官去見面.
這個大將軍來到的時候.
以利沙不是一個軍官.

$^{201}$但他都很尊重以利沙.
竟然以利沙自己都不出來.
就找一個使者.
或者找一個學生出來.
就跟他說.
我的老師這樣說.
你去吧 你去就會得到醫治.
這個乃曼就非常之火滾.
他的反應是非常之怒氣.
轉身就走了.
為什麼他這麼怒氣呢.
除了一來完全感覺到不被尊重之外.
以色列王也不太尊重他.
以利沙也不太尊重他.
然後他說.
你去約旦河沐浴七次.
他說 你對我說這些話.
去約旦河.
那條黃金的河.
難道我大馬士革沒有河嗎.
我大馬士革那兩條河不是更漂亮嗎.
如果我是去大馬士革那裡.
去沐浴的話.
不是會好起來嗎.
但他轉身走的時候.
那個僕人就跟他說.
他說 不是啊.
他說 大將軍.
如果他叫你.
上去以色列最高的那座山.
叫你一定要自己爬上去.
過一晚夜.
第二天早上的朝露.
令你康復.
這麼難的事你一定去做.
他叫你去禽足七次.
你都去做.
但叫你做一些難度都沒有的事.
做完之後.
都不知道你做了什麼.

$^{241}$或者做完之後.
人家會覺得他做了這些事.
你就不想去做.
但你試一下.
你去做一下.
可能你沐浴七次之後.
會不會好起來呢.
這個乃萬將軍的僕人.
未必是奴隸.
他的僕人可能是他的副官.
隨著他出行.
隨著他打仗.
他說的話.
很適合乃萬聽.
乃萬下了一口氣.
他就真的試一下去約旦河.
他去約旦河.
真的沐浴七次.
我不知道什麼叫沐浴七次.
是否要分開七天呢.
還是每次要起來.
抹乾淨再下水呢.
總之他就沐浴七次.
七次之後就好了.
每次要浸多久呢.
不知道.
不過應該他和以利沙溝通好了.
他浸完這七次之後.
有一個很特別的結果.
可能全部人都嘔吐了.
他不是皮膚好了.
他的皮膚變得像小朋友一樣.
滑溜溜的.
可能他中東人不是白雪雪.
但還是滑溜溜的.
嬰兒呢.
抱嬰兒的時候.
嘩!他的肌膚多漂亮.
做多少次Facial都不行.
沒有這樣的效果.

$^{281}$但他就是這麼漂亮.
乃萬沐浴完七次之後.
這個結果他反對.
可能他太太就說.
將軍,我也想去沐浴七次.
不用做Facial.
他就真的沐浴完七次.
那他是不是很開心.
馬上回亞蘭呢.
其實可以的.
以利沙什麼都沒有見過他.
他不是.
他真的回到以利沙面前.
帶了他的禮物去.
原來那些禮物不是送給以色列王.
帶了他的禮物去.
就要送給以利沙.
怎麼求以利沙呢.
原來以利沙怎麼求她都不收.
原來不收禮物有一個結果.
他怎麼求的時候.
以利沙都不收的時候.
乃萬就跟他說.
如果你不收.
你送份禮物給我.
你不收我的禮物.
那你送份禮物給我.
你送給我一些泥土就行了.
你不用送金,不用送銀.
你送泥土給我.
我兩隻女孩背得起的泥土.
你送給我.
因為我以後要敬拜耶和華.
不知道他拿來怎麼做.
他拿這些泥土.
他就會向他們獻祭.
敬拜耶和華.
他說因為我知道.
哪裡都沒有神.
只有以色列有神.

$^{321}$所以我要敬拜這個神.
不過有些事.
我一定要先跟你申報.
就是我一定要陪我的君王.
敬拜他的神.
而他敬拜叩拜的時候.
我又不可以高過他.
他叩拜我又要陪他叩拜.
這件事你就跟耶和華說一聲.
讓他先原諒我.
我什麼時候都要這樣做.
只要我的君王去到他的臨門廟.
我就要這樣做.
以利沙就很簡單地.
只說一句話.
她說你平安地回去吧.
她沒有刁難他.
沒有跟他說你信得不清不楚.
她跟他說你回去吧.
神知道.
這個人就回去.
是的.
這個人他信了主.
他說他信耶和華.
他只獻祭給耶和華.
其實小侍女沒有跟他傳福音.
小侍女沒有跟他傳福音.
不過經過小侍女.
引帶他去見一位先知.
其實小侍女認不認識以利沙呢.
小侍女沒有說他是以利沙.
不過就告訴他他去撒瑪利亞.
你去撒瑪利亞找那位先知吧.
撒瑪利亞都不小的.
去到找哪裡呢.
原來他找到以利沙.
由以利沙那裡他經歷了神.
他經歷了神之後.
他相信了.
因為他不是沒有拜祭過其他神.

$^{361}$亞蘭有很多神.
他都要陪他的皇帝去敬拜.
他不會沒有為自己的身體祈求過.
他求醫求神一直都沒有好結果.
竟然現在.
他可以得到醫治.
這個才是真實.
他相信了.
原來小侍女只是引解他這件事.
這位大將軍就相信了主.
這個小侍女是什麼人.
我們剛才說他是難民.
是奴隸.
其實他應該自卑的.
我跟大將軍說你去哪裡求醫.
大將軍醫不好回來.
那條命就沒有了.
他沒有擔心沒有自卑.
而且他是出於真心去關心這位將軍.
所以他跟他太太說的時候.
我真的很想他能夠去到那裡找那位先知.
他就可以得到醫治了.
我相信他不止說了一次.
說了一次主母都不出聲.
出了聲用什麼將軍.
他不開心的時候就出事了.
他很有恆心.
也很有方法慢慢哄哄主母.
有花心思的.
令到主母能夠相信他.
然後將這件事跟主人說.
因為他不自卑.
因為他的真誠.
他帶了是不是一個人去信主.
不是啊.
我相信乃萬回去影響他太太.
影響他的軍隊.
他是亞蘭的一位大元帥.
他做的事他沒有去哪裡拜祭.
很快大家都知道了.

$^{401}$上了頭條了.
是啊.
這位就是乃萬.
他家裡那位無名的小侍女.
到現在我們都不知道他的名字.
現在都只知道他是一個以色列人.
但是他就是這樣.
他的小小的努力.
他就幫助這位大元帥.
這個大元帥他得到幫助之後.
他身體一切的機能就好了.
似乎他告訴我們一些事情.
原來帶人信主.
向人傳福音.
不一定做很多事.
首先第一件事你要認清楚.
你要帶他去信主.
這位小侍女原來不是跟他的主母最近.
他最多跟他接觸的人跟他說.
你信我的耶和華吧.
而他是很看到乃萬有個需要.
他看到乃萬有需要的時候.
他就說乃萬要去以色列.
要去撒瑪利亞.
在那裡求醫.
這樣的一個行動.
他帶了這位大將軍信主.
他的真誠打動了主母.
也打動了乃萬.
他不怕死的.
不怕乃萬殺了他.
但是他就很真心關心他.
關心你行養是很慘的一回事.
以前都沒有那麼多飼養的藥物.
有一些皮膚病真的會損壞流血的.
影響兒童的.
他很關心他.
他的關心乃萬收到.
他也在抓這個時機.
就是兩國沒有戰事.

$^{441}$以色列和亞蘭國沒有戰事和平的時候.
他就可以去求醫.
打仗的時候誰來醫你.
他也很有堅定的信心.
一定可以的 你去吧.
他要在背後做很多鼓勵的工作.
而且他能夠不斷接再厲.
我想這是我們要做的.
他跟主母說的時候.
不一定第一次主母就說.
好吧 我跟乃萬說.
不是的 主母也要想一想.
考慮一下環境.
看看他心情好.
也要主母心情好.
主母相信他才會跟乃萬說.
所以小侍女要很努力去做.
我們一起做一件事.
我覺得這個小侍女其實是一個傳道者.
她雖然無名.
她雖然沒有直接跟乃萬傳福音.
沒有跟他說.
我認識三福又認識四律.
你坐下來聽我說吧 沒有.
但她就是引進了他去一個先知的家.
讓先知去幫他.
弟兄姊妹 你一定有一些對象.
是未信主的.
甚至你覺得是帶不到他信主的.
大家的關係不去到那裡.
小侍女是一個奴隸.
怎樣跟一個大將軍 大元帥叫他信主呢.
但原來真誠的時候.
其他的東西是神做的.
我們只不過是負責說幾句話.
這個小侍女就是這樣.
做了一件事.
做了一件事到現在是千八年之後.
我們記得她 雖然不認識她.
我們都知道她是小侍女.

$^{481}$而且知道她是無名無姓的小侍女.
乃萬的小侍女.
乃萬太太的小侍女.
這個就是我們對她的記憶.
但她的生命完結.
她的名聲繼續流傳下來.
我們都學做一個無名的傳道者.
或者我們不是很多機會走上來台上講見證.
我們不是很有機會在這一說講到呼召.
但我們可以引人回來.
帶他回來 引帶他去信主.
我們認定哪一個是我們要邀請的對象.
我們多方嘗試.
又為他祈禱.
要向他傳福音.
要他信主這件事交給神.
神一定會感動你什麼時候是最好的.
祂會製造時機給你.
為他禱告.
神會讓他心靈裡面覺得需要.
他會信主.
我們將他交給神.
神會幫助他.
在一些我們估不到的時候.
神會令他信主.
甚至乎他信了主之後.
你會覺得.
是真的嗎.
是不是啊.
我和我媽媽講福音講了十幾年.
她都不肯信.
但原來一個很突然的機會.
就是她病的時候.
有一個不懂得講鄉下話的傳徒.
因為我媽媽只懂得講鄉下話.
和少量的菲律賓話.
她不懂廣東話.
她找了一個懂得講鄉下話的人.
那個傳徒是廣東人.
她找了一個懂得講鄉下話的人.

$^{521}$去和我媽媽傳福音.
那晚我媽媽就信了主.
我心想.
她不是.
她只不過是想打發人走.
過了一兩天我就試了我媽媽.
我試了我媽媽.
我知道她真的信了.
多開心.
然後她就四處找團契回.
她自己四處找團契回.
找人帶她回教會.
我們只是做一點點的事.
一點點的事.
將事情交給神.
讓神帶他們回教會.
讓神帶他們信主.
我們交給天父.
22號報道會.
你有沒有想到要邀請誰.
你是不是想到一個很難邀請的對象.
你交給天父.
交給天父.
讓天父幫你安排.
我們一起祈禱.
將天父上帝我們來到你面前.
我們要多謝你.
你拯救了我們.
你給我們有機會聽福音.
有機會決志信主.
有機會.
現在我們有機會帶人信主.
主我們獻上感恩.
求主你幫助我們.
幫助我們將你給我們的恩典.
帶給我們身邊的家人朋友論社.
讓他們都有機會聽聞福音.
求主你感動他們信主.
主也你幫助我們.
我們做不到什麼.

$^{561}$但是求你幫助.
甚至連邀請.
主也我們都不能夠很有把握.
求主你紀念幫助我們.
求主也都預備.
讓那天有美好的天氣.
讓那天每一個人都能夠有好好的心情.
願意回來.
求主預備.
預備他們回來.
預備他們信主.
求主也都紀念那天報道會.
每一個事奉的人.
我們仰望你奉主耶穌基督的聖名.
阿們.
\newpage



\section{歷代志上 13:9-14}
\label{sec:ui5VIns3kj4}
\textbf{2024-06-02 早堂《前車可鑑》歷代志上13\_9-14 講員:陳一華牧師}
\newline
\newline
連結: \href{https://youtube.com/watch?v=ui5VIns3kj4}{\texttt{https://youtube.com/watch?v=ui5VIns3kj4}} ~~~~ 語音日期: 2024-06-05
\newline
\newline
\hyperref[sec:Il1UURetEws]{\small{< < < PREV SERMON < < <}}
~
\hyperref[sec:index_chronic]{\small{[返順時目]}}
~
\hyperref[sec:index_scriptual]{\small{[返順卷目]}}
~
\hyperref[sec:DjV4K1Dg8GU]{\small{> > > NEXT SERMON > > >}}
\newline
\newline
歷代志上 13:9-14
\newline
\begin{longtable}{cl}
\hline
\hline
章節 & 經文 (和合本修訂版)\\
\hline
13:9 & \begin{tabularx}{0.7\textwidth}{X} 到了基頓的禾場，因為牛失前蹄，烏撒就伸手扶住約櫃。 \end{tabularx} \\ \\ \relax
13:10 & \begin{tabularx}{0.7\textwidth}{X} 耶和華的怒氣向烏撒發作，因他伸手扶住約櫃而擊殺他，他就死在那裡，在神面前。 \end{tabularx} \\ \\ \relax
13:11 & \begin{tabularx}{0.7\textwidth}{X} 大衛因耶和華突然衝出撞死烏撒就生氣，稱那地方為毗列斯‧烏撒，直到今日。 \end{tabularx} \\ \\ \relax
13:12 & \begin{tabularx}{0.7\textwidth}{X} 那日，大衛懼怕神，說：「我怎能將神的約櫃接到我這裡來呢？」 \end{tabularx} \\ \\ \relax
13:13 & \begin{tabularx}{0.7\textwidth}{X} 於是大衛不將約櫃接進大衛城他自己的地方，卻轉送到迦特人俄別‧以東的家中。 \end{tabularx} \\ \\ \relax
13:14 & \begin{tabularx}{0.7\textwidth}{X} 神的約櫃停在俄別‧以東家中三個月，耶和華賜福給俄別‧以東的家和他一切所有的。 \end{tabularx} \\ \\
[1ex]
\hline
\hline
\end{longtable}
$^{1}$各位姐妹早晨.
今早我們回祖堂崇拜很有福氣.
大家知不知道.
因為我們今天可以聽到錦繡堂最高級數的師班.
來獻詩.
這個是講年紀.
我們真的很欽佩一班長者.
因為他們加起來都一千多歲了.
雖然年紀大.
但一樣都這麼積極的釋放.
也都很敬佩我們的長者.
牧師今早跟大家思想的訊息.
就叫做前車可鑑.
剛才主席秀慧都已經讀了出來.
大家都知道這段經文.
記載是一個失敗的例子.
也都是一個反面的教材.
為什麼這麼說呢.
因為記載的經文裡面提及到.
大衛想將藥櫃搬回去耶路撒冷.
即是大衛城.
為什麼無端端藥櫃要搬回去耶路撒冷呢.
背景就是因為當時.
曾經以色列人和非利人打仗.
結果非利人打勝仗.
以色列人節節敗退.
在那個時候有人想到一個方法.
不如就將他們最厲害的殺手簡.
將藥櫃搬到戰場的地方.
看來這次我們都贏定了.
為什麼呢.
因為藥櫃在那裡.
這是當時以色列人的想法.
不過他們想得太幼稚了.
為什麼呢.
因為他們單單外表有個藥櫃.
而沒有神的同在.
是形同虛設的.
大家認同嗎.
所以在這樣的情況下.

$^{41}$大家都知道那場仗.
以色列人敗得很慘烈.
而藥櫃都被人搬走.
為什麼呢.
因為以色列人離棄了上帝.
所以上帝就離開了以色列的地方.
藥櫃自此就被非利人搬去.
大捆神像的廟宇.
結果第二天大捆神像就撲倒在藥櫃的面前.
大家有印象嗎.
讀過這段經文就記得這個情景.
嚇死非利人了.
為什麼呢.
既然他們自己的神像都撲倒.
所以他們覺得非常之可怕.
這次很嚴重.
所以馬上就要將藥櫃搬離開.
送回以色列人的境界.
這就是當時的背景.
不過很可惜.
藥櫃沒有搬回耶路撒冷.
藥櫃只去到基列耶林的地方.
聖經記載.
藥櫃停在那地方有二十年這麼久的時間.
換句話來說.
那種光景可以說是樂得淒涼.
既然忘記了上帝的藥櫃.
有二十年的日子.
藥櫃就放在基列耶林的地方.
直至大衛造王.
大衛造王的時候.
大家都知道大衛王是很愛上帝.
當然上帝也很愛錫他.
所以大衛王所做的事.
都希望是土神的喜悅.
大家都認同的.
不過人始終都是人.
有他軟弱的時候.
有他閃失的一刻.
這次的例子是一個很明顯的例子.

$^{81}$說到大衛王做一件事.
竟然做得這麼淪盡.
做得喜事變成悲事.
好事變成壞事.
所以在我這次的題目當中.
就定了叫做前車可鑒.
頂止無門.
大衛竟然我不知道他是心急.
不知道他是一時疏忽.
或者他沒有時間在上帝面前等候.
或者他沒有時間和祭司利未人商討.
還是不知道什麼原因.
竟然他忘記了摩西五經記載.
上帝的藥櫃是怎樣搬運的.
大家都知道是要利未人才有資格.
是什麼? 抬神的藥櫃.
不過這次可能他覺得自己.
現在時代都一日千里了.
要有新的創作 新的創舉.
竟然他想起做一架車.
找一架車 然後前面找一隻牛來拉 去運送.
旁邊派兩個人來護衛.
運送上帝的藥櫃.
他想這個方法應該是效率最高.
他想到這個新的方法.
應該是一個很嶄新的一種行動.
他忘記了藥櫃是象徵神的同在.
那個的精義.
這次出現了好幾個的錯失.
這次跟大家思想一下.
我又是忘記了給大家看PowerPoint.
第一個錯失是什麼呢?.
是政策的錯失.
正如我剛才所說.
大衛他決定搬運上帝的藥櫃的時候.
他做了一個很錯誤的政策.
他是用車 用牛 來運送上帝的藥櫃.
結果離開了上帝的心意非常之遠.
在這樣的情況之中.
他就種下了這次失敗的原因.

$^{121}$其實政策也是很重要的.
各位電影節目們.
原來我們做每一件事情.
我們都需要能夠想清楚.
或者深思熟慮.
到底這個政策 這個方式 這個方法.
對不對呢?.
是不是上帝喜悅的呢?.
這樣的做法 是不是令人得到幫助的呢?.
我曾經在上個世紀的時候.
八九十年代裡面.
我曾經認識一些宗派.
當他們要提供一些社會的服務.
或者社區的服務的時候.
他們的政策很特別.
他們的政策是沒有包括到心靈的關顧.
或者沒有福音的元素.
純粹就是一個服務.
純粹就是一個社區的服務.
一個社會的服務.
但是他們沒有把一個重要的.
全人的關懷.
一個心心社靈的關懷的元素.
來結合在一起.
他們沒有這樣做.
所以是很可惜的.
為什麼呢?.
他們這樣做的時候.
令到受益的街坊朋友.
他們只是在物質上 身體上.
得到一些好處.
但他們卻不能夠在當中開始.
來滿足他們心靈上的需要.
所以各位電視節目們.
原來我們做一些的政策是很重要的.
我們再看大衛第二個錯失是什麼呢?.
第二個錯失.
大衛是一個什麼呢?.
情緒的錯失.
大家知道嗎?.

$^{161}$當大衛知道烏撒被上帝擊殺的時候.
他的情緒有什麼呢?.
有很大的反應.
和合本翻譯了大衛心裡面非常的受害.
和收本翻譯成大衛生氣.
大家有沒有留意到?.
換句話來說.
他很不明白上帝為什麼對他這麼好.
他這麼想念他.
他這麼想念他的約櫃.
想做一件好事.
但他竟然來到什麼呢?.
竟然擊殺了他的神僕.
各位電視節目們.
這是一個情緒的反應.
其實大衛又做錯了一件事情.
為什麼呢?.
他沒有首先去檢視一下自己.
他沒有去檢討一下自己.
我這樣的一種方式.
來搬運上帝的約櫃.
是不是合神的心意呢?.
他沒有做一件這樣的檢視行動.
相反他就用他的情緒來表達.
他的不滿.
所以你會發覺到.
大衛接著他的情緒又犯了第二個毛病.
很多時候當我們要和別人合作的時候.
我也發覺有一個很大的通病.
或者一個很大的致命傷.
就是當你和別人合作的時候.
誰不知合作的計劃失敗了.
通常你會怎麼想呢?.
通常你會想.
當然是我的合作夥伴做得不夠徹底.
當然是他有些東西漏了.
當然是他有些事情瞞著我.
不告訴我.
所以這件工程,計劃就失敗了.
很多時候我們發覺.

$^{201}$當有些不順利的事情.
有些失敗的現象出現的時候.
我們很多時候都很容易歸咎其他人的錯誤.
我們沒有去想.
到底我自己在這件事情上.
我有沒有做錯了的地方呢?.
或者有沒有一些程序我做得不對呢?.
各位電視節目的觀眾.
我自己在外間和很多機構合作.
或者和教會合作的時候.
我都是常常提醒自己.
當有些不太理想的結果出來的時候.
第一個行動我們應該先去想想自己.
在我的本份上.
我有沒有遺漏了呢?.
我有沒有做得不夠好的地方呢?.
然後再看看這次的失敗.
這次不理想的結果.
其他人是不是都有因素在當中呢?.
所以我們要常常謹記.
不要重蹈大衛的覆轍.
因為這個大衛是第二個錯失.
就是他什麼?.
情緒的錯失.
我們再看看大衛第三個錯失是什麼呢?.
大衛第三個錯失就是關係的錯失.
因為聖經記載.
自從這件事之後.
大衛就什麼呢?.
懼怕上帝.
為什麼上帝這麼兇惡呢?.
為什麼上帝不可以留一手.
不去擊殺我的神僕呢?.
他覺得整件事情上帝表現得很不近人情.
所以他就懼怕上帝.
懼怕上帝的引致是什麼呢?.
他就開始和上帝的關係疏遠了.
更加說不上要把藥櫃搬回大衛城.
結果他沒有把藥櫃搬回大衛城.
為什麼呢?.

$^{241}$因為他懼怕上帝.
所以他就漸漸疏離上帝.
在這樣的情況當中.
我們發覺真的非常可惜.
他不但政策錯失.
情緒錯失.
連關係都錯失了.
他沒有在這件事情上馬上悔改.
求上帝寬恕.
求上帝憐憫.
再一次地緊緊扣上和神的關係.
不是,他是怕了上帝.
不再接近上帝.
曾經我也見過一些弟兄姊妹.
他們向上帝祈禱.
舉個例子來說.
他們向上帝祈禱.
求神祝福他們做一盤小生意.
上帝說我的生意也不是很大.
也是十萬八萬投資下去.
祈禱求上帝讓他們的生意興隆.
讓他們的生意做得成功.
讓他們的生意做得有氣息.
希望上帝祝福我.
他這樣禱告是OK的,是對的.
但他想不到原來當他做了這盤生意之後.
突然之間原來香港人的口味變了.
個個都喜歡北上.
做的生意沒有人光顧.
越來越少.
好像上帝不祝福我.
上帝好像不再將恩典給我.
所以心裡面很不舒服.
心裡面很不忿氣.
我這樣求上帝.
為何上帝不聽我的祈禱呢?.
結果有些弟兄姊妹.
因為一些不順利,不如意的結果.
他們開始跟神的關係疏離.
我不再讀聖經了.

$^{281}$我不再祈禱了.
上帝都不祝福我.
他都不看緊我.
我不再回教會了.
我不再敬拜上帝了.
為甚麼呢?.
因為上帝都沒有聽我的禱告.
又不祝福我.
所以各位弟兄姊妹們.
在這樣的情況當中的時候.
我們要加倍地提醒自己.
遇到不順利的時候.
不如意的時候.
我們會否想到.
我們更加需要加倍地.
回到神的面前.
求上帝再一次憐憫自己.
幫助自己.
在困境,逆境當中.
能夠可以再突破這些困難呢?.
各位弟兄姊妹們.
原來你發覺.
真是一樣接一樣都發生在大衛的身上.
我們再看第四個的錯失是甚麼呢?.
就是連福氣都錯失了.
你說是不是真的很不值呢?.
弟兄姊妹們.
為甚麼這樣說呢?.
原來當大衛將上帝的約櫃.
不搬回去的時候.
他就這樣放在一個利未人的家裡.
臥別以東的家裡.
那就算了.
放在那個地方多久呢?.
放在那個地方有三個月的時間.
上帝的約櫃就留在臥別以東的家裡.
但是你知不知道.
臥別以東因為願意去服侍上帝.
他就在這樣的情況當中.
得到上帝很大的祝福.

$^{321}$上帝就賜福臥別以東和他的全家人.
由此我們見到.
原本這個福氣,這個祝福.
應該臨到大衛的身上.
是不是?.
如果他真的好好地做對程序.
做對規矩.
有將上帝的約櫃搬回去大衛城的時候.
我相信上帝會有更大的祝福臨到甚麼呢?.
大委身上.
但是現在大衛一點福氣都得不到.
為甚麼呢?.
因為他將約櫃放在臥別以東的家裡.
他沒有將這次失敗的問題.
整個原因.
能夠從一個根本的地方著手去處理.
所以他就一步一步地錯失了.
而最後連上帝的祝福都失去了.
所以你說是不是一個很大的可惜遺憾的事情呢?.
各位電視節目們.
今早這段經文所記載的.
是一段以色列人歷史當中所發生的事情.
這件事情其實.
如果我們應用在今天我們的生活當中.
其實有很多地方你都覺得是很相似的.
就是當我們做任何的工作.
如果我們真的沒有很好的等候.
沒有很好的禱告,思考的時候.
很急躁去做一個決定的時候.
很多時候就會引起我們一個甚麼呢?.
最基本的問題.
就是政策就錯失了.
但你要記得.
錯失了第一步就千萬不要錯失第二步.
要馬上止血.
不要再錯下去.
這個情況有點像在座當中.
弟兄姊妹如果你喜歡下棋的話.
你就很明白.
你下棋的時候你要走得很精準.

$^{361}$走一步要多想幾步.
為甚麼呢?.
正所謂一子錯,然後又怎樣呢?.
滿盤皆落索.
各位電視節目們.
今次大衛一子錯.
繼續錯下去再錯下去.
是非常可惜的.
所以我們求主幫助我們.
如果我們日後遇到一些事情.
錯了第一步的時候.
不要再錯下去.
馬上回到上帝的面前.
再安靜,再等候,再禱告.
檢討一下自己哪一個環節.
我做得不是那麼好呢?.
我是有些忽略了,大意了呢?.
如果我們能夠在那個時候.
大衛馬上悔改.
知道自己做錯了.
不是別人錯.
是自己錯的時候馬上悔改.
我相信上帝都會憐憫他的.
因為他的動機是好的.
他的出發點是好的.
不過他的方法是用錯了.
各位電視節目們.
今日我們祈求上帝幫助我們.
透過這段經文.
讓我們前車可駕.
我們同心祈禱.
多謝天父.
因為你是一位又慈愛又有公義的上帝.
你在大衛的身上.
你是很愛你的僕人.
但是我們知道我們人始終都是人.
有時候我們有軟弱,閃失的時候.
求主就提醒我們.
叫我們在一些錯誤當中.
我們及時地回轉,悔改,檢討自己.

$^{401}$以致我們不要讓這個錯誤再延伸下去.
求天父幫助我們.
叫我們做每一件事情的時候.
都繼續地對準你的心意去做.
以致我們所做的一切是蒙你的賜福.
也都有你自己的恩典,祝福在其中.
我們同心仰望,禱告.
奉耶穌的名及.
\newpage



\section{帖撒羅尼迦前書}
\label{sec:DjV4K1Dg8GU}
\textbf{2024-06-06 周四晚查經班《帖撒羅尼迦前書(8)》 講員:曾裔貴牧師}
\newline
\newline
連結: \href{https://youtube.com/watch?v=DjV4K1Dg8GU}{\texttt{https://youtube.com/watch?v=DjV4K1Dg8GU}} ~~~~ 語音日期: 2024-06-07
\newline
\newline
\hyperref[sec:ui5VIns3kj4]{\small{< < < PREV SERMON < < <}}
~
\hyperref[sec:index_chronic]{\small{[返順時目]}}
~
\hyperref[sec:index_scriptual]{\small{[返順卷目]}}
~
\hyperref[sec:CutYNgBop3c]{\small{> > > NEXT SERMON > > >}}
\newline
\newline
$^{1}$我們可以真是有系統地去研讀.
《貼實羅尼加》的前後書.
一卷是頗重的木樣的書信.
試圖保羅對於初生的教會的關心.
他們在迫不還難裡面.
到底那個信心有沒有動搖?.
保羅怎樣來到去在信望外.
勉勵當時的弟兄姊妹去追求.
去捍衛這個信仰.
都是值得我們今日去效法的.
求主都幫助我們.
我們都有一兩堂講到關於主的再臨.
第四章提到的復活的次序.
都是值得我們去深思.
這個世界的人如果不信的.
他們是不會相信有復活.
也都不會在他們的意識裡面.
有想到將來身體,靈魂.
會有一個永恆的歸宿.
所以這個是一個真是石破天驚.
是基督教才有的真理.
我們很盼望神給我們不要掉以輕心.
我們真是很認真去看這件事.
以致真是影響著我們地上的有限人生.
真是幾十年.
我們需要在主的眼中去過我們的生活.
求主幫助.
我們今天晚上再一次等候在你的面前.
我們還有好幾段的經文.
關於主再來要看的.
都幫助我們能夠透過不同的使徒.
我們去梳理經文的當時的內容.
和那個勸勉.
我們這樣仰望禱告.
奉耶穌基督的名.
阿們.
稍微補充第四章.
如果我們真是對於那個復活的盼望.
其中保羅不是為了講復活而講的.
最大的那個關注.

$^{41}$就是因為當時有一些弟兄姊妹.
很擔心自己已經死了的家人.
他們到底現在身在何處.
和將來是怎樣呢?.
所以保羅基於這個模樣.
來跟他們去講.
所以這個整個主再來.
是在這一段經文.
不是特別要探討的.
而是講到那些真是離開了世界的人.
他們到底現在是怎樣的情況呢?.
和將來主耶穌回來的時候.
又會是怎樣呢?.
就一個非常之.
你可以說那個涵蓋的內容.
那個關心和那個層面.
是非常的廣闊.
只要在主裡面.
不論你是甚麼背景,甚麼階層.
只要你在主裡面.
在救恩裡面.
你就能夠在那個復活的行列.
這個是真的很重要的.
因為死了的人.
證明地上他有的東西完全沒有了.
所有死的人就甚麼都沒有了.
甚麼都不關你的事.
但是復活的時候又關你的事.
無論是五千年前還是五百年前.
無論你是一個皇帝.
還是你是一個奴隸.
都關你的事.
你可以想像到主耶穌的救恩.
帶給我們屬於主的人的那個盼望.
是多麼的終極和永恆.
所以這個基督教.
真是和任何宗教不同的地方.
一方面我們不是佛教的律.
那些是人的構想出來的.
但這裡是真的在說.

$^{81}$我們是會有一個永恆的生命.
所以這個記住是.
保羅說的是免得我們憂傷.
好像那些沒有指望的人.
這樣來去憂傷.
所以足以安慰任何要面對.
喪葬的人的心靈.
所以這段經文.
某程度上都應該是在安息禮拜.
在火化場.
火化禮的時候來去做讀的.
所以我們都改了.
不知道大家有沒有留意的.
可能你們參加的少.
很早期張牧師年代.
就是用《哥羅門島前書》15章的.
後來我們都改了.
就是用《跳獸來的家前書》4章這段經文.
在火葬場我們會做讀的.
就是真是給到每一個喪家的弟兄姊妹.
家人有這個安慰.
將來必定會有一個復活.
所以保羅這段是一個牧羊的關心.
不是主要想講主在來的.
無論是那個次序.
但必然一定會碰到那個課題.
所以保羅真是照主的話.
來去講這個啟示.
因為不是人自己想出來的.
另外想和大家看看《雅各書》有一段.
第五章.
似乎這個就是耶穌的弟弟雅各.
肉身的弟弟雅各.
他提到.
第五章這個可以讀一點點.
讓你感受一下當時的氣氛.
寫這封信的時候.
很實際處理的課題.
教會人士應該怎樣避免一些不好的錯的事.
這裡第五章一到第六節.

$^{121}$和下面應該不相同的.
但是值得我們去想一想.
關於那個警醒,忍耐,等候.
和上文講的很實際的生活裡面.
《雅各書》五章.
大家不知道有沒有機會讀過這些經文.
五章一到第六節.
就是講那些為富不仁的人.
為富不仁即是很有錢.
但是仍然要掠奪那些窮人.
第五章.
你們這些富足人應當哭泣,號圖.
因為將有苦難臨到你們身上.
你們的財物壞了.
衣服也被蟲子咬了.
即是落葬了會.
你們的金銀都長了壽.
那壽要證明你們的不是.
即是已經可能是死了.
又要吃你們的肉如同火燒.
好像特別是強調死了的時候的狀態.
但是這些人是為富不仁的人.
這些是在地上富足的人.
他們應該要趁沒有審判之前.
要哭泣,要號圖.
號圖即是很大聲的哭.
即不是普通的忍泣.
接著說 為什麼呢?.
因為第三節,下半節.
你們在這活世 只知道積賺錢財.
工人給你們收割莊稼.
你們虧欠他們的工錢.
即是乞購工錢 即是為富不仁.
即是他們是地主來的.
他們是田主來的.
他們是僱工人去耕種的,去收割的.
但是不付錢.
即是用各種藉口去乞購工人的工錢.
「這工錢有聲音呼叫.
並且那收割之人的冤聲.

$^{161}$已經入了萬君之主的耳了」.
所以這些冤情是會上達.
不是去包青天 上達去天庭.
因為神是則以而聽困苦窮乏人的哀告.
舊約很多的 利面記有很多的.
所以你有沒有讀那些律法書呢?.
那些律法書其實挺好看的.
它就提到你借錢給人.
那個人就沒有錢還了.
接著你借錢給他 那叫抵押的.
聖經這樣教的.
你見要站在門口.
他給你什麼是抵押 你就要收什麼的.
你不可以進去自己拿的.
很有趣的 那個債主不可以進去裡面的.
接著就是如果那些是棉被.
就不可以過夜你不給他.
因為他如果冬天冷 沒有那個衛生的.
那罪就歸你的了.
不是的 我借錢給他.
這些利面記的律法書.
有教那些債主應該的.
你可以借你錢 甚至你可以向外邦人取利.
就是收利息 但是對自己的同胞.
就不要收利息 只是免息借貸這樣.
有教的 但這裡似乎作者看到.
一個社會上的一個結構性的剝削.
不知道什麼原因呢?.
他們剋扣那些工錢.
這裡就告訴我們.
那些工錢就有聲音呼叫.
當然那些工人呼叫 那些錢是不會叫的.
你可以想像到一個富戶繼續的享受.
那些可能是收成的成果.
大批延職 但是那些工人就很慘地.
都不能夠戶口了.
這些一定會追究的 第五節.
你們在這世上享美福好艷樂.
當宰殺的日子 竟驕讓你們的心.
宰殺的日子就是秋收.

$^{201}$冬天我們小時候就要納那些臘肉.
不知道你有沒有做過.
就是很富庶的.
將那些變成臘腸,臘肉這樣.
來預備過冬.
我想在舊約的年代.
那些農業的社會都會這樣.
比如外國那些草留給那些馬,牛去吃等等.
中西部的那些大的農莊.
這裡告訴我們 當宰殺的日子.
這幫有錢人竟然驕讓他們自己的心.
你們定了二人的罪.
把他殺害 他也不抵擋你們.
揭露了社會的不公.
要他們趕緊去哭泣,號圖.
要他們悔改 如果不是的話.
那些錢就會過去了.
他們那些人 那些漂亮的衣服.
那些錢都會變成生鏽.
甚至會吃你的肉.
就是說他們真的埋葬.
可能葬的時候都是風光大葬的.
不過就進入那個審判的領銜了.
那個神去追究那些冤星.
那些窮乏人受欺壓的冤星.
這一段都很繪形繪聲的說了.
主在來的審判.
接著帶來第七節 就是說主在來.
弟兄們 你們要忍耐直到主來.
看!農夫忍耐等候地裡寶貴的出產.
直到得了秋雨,春雨.
你們也當忍耐,堅固你們的心.
因為主來的日子勤了.
接著就說在教會裡面.
跟人的相處應該怎樣呢?.
那裡可能涵蓋的是一些很實際的實務.
我們不詳細說了.
這裡特別強調.
要好像一個農夫他下了那個種子.
在猶太人的社會.

$^{241}$就是以色列地中海氣候.
他們有晚雨,有早雨.
或者說有那個叫秋雨.
按那個季節.
他們的秋天就可能是我們夏天的時間.
秋雨,春雨.
就是說他們去到冬天.
應該是1,2,3月就會有最後的雨.
那些農作物就可以有豐富的收成了.
這個是地中海的氣候就會有的這樣的狀況.
這裡告訴我們.
農夫他的忍耐就不是白等的.
一定是他首先就是弄鬆土地.
讓那些石頭不會阻礙那些種子的生長.
還有那個土壤要肥沃.
做完這些了.
然後才播種的.
播種之後.
雖然他等待雨水的滋潤.
但是也有很多東西要等候和警醒的.
例如那些飛鳥.
例如可能會有一些蟲等等.
這些都是要考慮的.
所以這個形容主在來的時候.
我們要忍耐.
這個忍耐就不只是一個白等.
應該都是說我們都會繼續為主耶穌.
記住耶穌說的.
這個田地就是這個世界.
我們都是在這個田地裡面去播種的.
就是說我們都是代表著這一位真正的大地主.
不是香港政府 不是特區政府.
實際上是的.
香港所有地都是政府的.
但是真正的地其實是神的.
所以耶穌說的.
田地就是世界的.
《馬太福音》13章說的那個殺種的比喻.
田地應該說是那個敗子的比喻.
所以整個宇宙萬物都是神的.

$^{281}$所以現在我們就是大地的管家.
我們是代表著神在這個地上播種.
所以求主讓我們能夠看到.
保羅都是這樣教導我們的.
他就是那個殺種的.
保羅就是囂觀的.
唯有誰叫他生長?.
唯有神才可以叫種子生長.
所以保羅就很懂得自己是.
他刻如其分自己是做什麼恩賜的崗位.
神要他去外邦 他就去外邦.
他不會去猶太人的地方.
和彼得一樣的去傳道.
所以我們每個人都是領受神給我們.
某程度上是其中一部份的地圖.
在我們那個地圖上成為大地的代表.
成為全地的一個代表.
我們去傳遞神的真理.
神的平安的福音.
所以一個穩定的社會.
才能夠有事可以做.
如果你是民不聊生.
你想一下現在你在哈馬斯.
你在巴勒斯坦地 你可以做什麼呢?.
吃都不能吃了.
每天排隊就起來去倒水了.
那你還可以做什麼呢?.
每天都怕突然間會死.
那你在一個這麼不穩定的世界裡面.
就什麼都做不了.
你要顧那個生存.
你都不會再想其他事情.
所以保羅時時都教我們.
《彭斯前書》第二章.
為君王和一切在位的也該如此.
然後怎樣呢?.
使我們可以平安無事的度日.
這是好的在神.
我們救主面前可無月立.
所以教會是要為在上的禱告.

$^{321}$如果沒有一個穩定的政權.
我們就不能夠彰顯那個平安的福音.
當然神可以用戰亂.
神可以用一些甚至溫逆的時候.
神可以用不同的方法.
但是大體來說.
都是要一個很穩定的社會.
讓神的道可以在其中滋生.
所以這裡表面上是說一個農夫.
但其實是在說我們是一班.
真是有那個唯一有主在來.
盼望的一班人.
世界的人是知道這個盼望的.
不會羨慕.
不知道會不會羨慕.
他不會稀罕的.
他只會知道人死就是死.
最多是想想辦法.
不要想那樣東西.
麻醉自己.
就算是發達國家.
落後的國家.
都不會想死亡之後的.
但是基督教就不同了.
真是要認真看這件事是必然會發生.
所以這裡提到的是忍耐等候.
直到得到秋雨,春雨.
作為一個農夫.
盼望到那些雨水.
是他所做的農作物.
必需的那些雨水.
就是說真正的豐收就會來了.
所以我們都不要這麼狹隘去看.
我們個人的靈性.
當然這是很關鍵的.
你的靈性不好.
你根本沒有心去為實做.
但是其實我們要看的是.
田地是世界.
田地就是實的經營的地方.

$^{361}$這個地球真是神給我們有一個很清楚的真理.
我們的主耶穌是會第二次回來.
所以我們今天要等待的.
就是好像一個勤力的農夫.
是帶著盼望每天出去殺種的農夫.
所以你有沒有起來就感覺很悶,很累.
心靈又很疲倦,沒有心情.
偶然間沒有問題的.
就是你病或者很想去旅行.
那當然沒有問題的.
但是你經常性都是這樣.
可能你對於永恆的盼望.
要做一些功課去補習.
因為你應該起來.
就是你能夠改變這個世界.
其中一個人就是你.
你應該有這個渴望.
你應該看到你在這個世界裡面的管家.
每個人都是在主那裡領受了一個恩賜.
和一個的光偉.
所以我們都要等候,忍耐,兼顧我們的心靈.
因為主來的日子近了.
似乎這裡作者雅各都告訴我們.
這個主一定是在說耶穌的.
一定是在說那個第二次再來.
祂都是看著耶穌升天.
和彼得他們一起在橄欖山.
看著耶穌坐雲升天會再來的.
所以他都是這樣去勉勵當時的信徒.
就是說在當時的社會道德.
這個似乎都還在猶太人的圈子.
和似乎都還在耶路撒冷.
福音還沒有很廣傳的情況.
因為雅各似乎一生都在耶路撒冷.
他沒有去到很遠的外邦.
這一位耶穌的弟弟.
肉身的弟弟.
很難想像你自己的哥哥,哥大佬原來是神.
很難想像的.
我想他將來耶穌升天之後.

$^{401}$對他來說的衝擊是很大的.
自己小時候一起玩橄欖山的那個哥哥.
原來是神來的.
你很難想像.
那個哥哥從來都不發脾氣的.
從來都不犯罪的.
你很難想像.
你想一下都很美的一件事.
在這個時候他意識到.
他雖然在責備當時的社會.
指出那些不公義的問題.
但是他仍然對教會.
一定要懷抱著主在來的忍耐,等候.
所以這個忍耐,等候其實是經營這個世界.
至少在你自己.
好像現在的人很喜歡去那些星期六的耕種.
有自己一塊可能幾尺的地種一些菜.
這樣的概念.
因為好像伏耕那樣.
但是我們應該這樣看我們自己的人生.
就是一個和場.
我們就是在主安排的世界.
我們去為主作見證.
所以我們要好像一個勞力的農夫.
來去殷勤的等候主的在來.
接著我們再看彼得候書的這一大段新聞.
彼得候書還有兩三段.
彼得候書.
差不多每一個使徒都有講到主在來的.
彼得候書的第三章.
本來是在第十節的.
但是也要跟大家梳理一下上下文的.
三章第一節開始.
「親愛的弟兄姊妹.
我現在寫信給你們是第二封信來的.
這兩封信都是要提醒你們.
激發你們誠實的心.
叫你們紀念聖先知御先所說的話.
和救主的命令.
就是使徒所傳給你們的」.

$^{441}$這麼詳細的說.
一定是和那些異端有關.
才說主的命令 使徒的教訓等等.
聖先知 舊約的先知.
他們的預言.
這樣一而再再而三的訴諸於權威.
證明一定是和那些異端.
在當時傳遞的錯謀有關聯.
第三節 第一要緊.
「各位知道 在末世.
必定有好機肇的人.
隨從自己的私欲」.
就是自己的私心.
自己有閒的眼光.
熟地的那種淺見.
只是以為自己是很聰明的.
這些叫私欲.
「出了機肇」.
就是還要到處去跟別人說.
自己有這樣的想法就算了.
還要到處告訴別人.
這樣叫機肇.
又譏諷又不屑的來去論斷.
第四節 他說.
「主要降臨的應許在哪裡呢?.
因為從列祖歲了以來.
萬物與起初創造的時候形式一樣」.
就這樣驟眼看又好像是.
我們人 我們只能夠活幾十年.
我們看不到以前的.
秦始皇都不會看到我們今天 是不是?.
我們不會看回以前的.
我們只看到自己的.
我們都不會看到2500年後的事的.
他說「主要降臨的應許在哪裡呢?.
因為從列祖歲了以來.
一直到現在.
哪裡有萬物復興呢?.
哪裡會跟創造的時候一樣呢?.
又似乎這樣的見解.

$^{481}$這樣的邏輯.
這樣的命題又好像確立了.
真的沒有人見過復活了.
到底是彼得已經知道在民間.
有人這樣機肇 似乎是的.
還是他擔心他的離開.
就會有些假的使徒.
會來傳一些似是而非.
可能起初的心不是想這樣.
特意拆毀基督教的真理的.
但是久而久之.
人們喜歡聽的時候.
就越來越有恃無恐的.
這樣來傳遞.
所以這裡可能開始是機肇的說.
他說這樣說第五節.
留意一下 彼得就在拆解了.
「他們故意忘記.
從太古行神的命有了天」.
即是未有世界以前的那個世界.
叫太古.
「並從水而出.
藉水而成的地」.
就好像Out of the water那樣.
即是說整個地球.
好像是被重水那樣掩蓋的.
神要陸地又翻出來.
這些是挺科學的.
不過就很簡略的說了.
總之就是最終的結局就是.
一個可以適宜人類居住的地球.
創世記創造論就是這樣的意思.
一個有空氣,有土地.
有我們可以呼吸.
有更加不同的生命.
有那些無論是樹木,花草和食物.
可以讓我們生存在這個地球上.
我們可以繁衍下一代的一個形態.
所以我們一出生根本不需要想.
地心吸力,萬有引力是多少呢?.

$^{521}$Gravity是多少呢?.
月球是多少?.
我們地球是多少呢?.
你都不需要想.
你一出生你慢慢就會懂得走路.
你會腳踏實地.
你去到月球 如果你生在月球.
你就不是腳踏實地了.
因為月球的引力是比地球輕的.
所以你是會跳來跳去的.
所以這個就是一個從水而出.
即水而成的地.
第六節.
故此當時的世界被水淹沒就消滅了.
但現在的天地還是憑著那命存留.
神的命定.
神的時間要這個地球.
要這個物質的宇宙再存在在祂認為的時間.
所以按神的命存留.
不是我們人.
無論我們怎樣去維繫都改變不了的.
過年的時候有一套電影叫甚麼《拯救地球》.
忘記了.
就是用一些爆炸的東西.
引地球離開太陽.
這些是科幻片.
沒有用的.
因為神命定了.
要它存留到今天.
所以我們今天看著這個地球.
我們有責任去維護.
但是最終都是神的保守.
接著說直到不敬田之人受審判.
遭沉淪的日子用火焚燒.
似乎上文講的被水淹沒.
特別強調是羅亞的洪水.
是經過了一次毀滅性的審判.
現在告訴我們第七節.
但是現在的天地還是憑著那命存留.
憑著神的命定維護而存留.

$^{561}$直到神認為要徹底解除.
那個人的罪惡就要審判.
所以將來就是一個天火焚石.
就是所有的眾生.
不在主的救恩裡面的.
都要被進入永恆的永遠的火壺來焚燒.
所以第八節就是那個渲染.
講了這個大前提就是.
已經有人在彼得的年代.
已經出來恐嚇.
沒有擔心這個世界不會毀滅的.
是不會有一個極端終極的審判的.
仍然存在的.
彼得說不是的.
古代都已經有洪水審判.
將來就會有火的永恆審判.
接著就講那個安慰了.
第八節.
親愛的弟兄.
有一件事你們不要忘記.
就是主看一日如千年.
千年如一日.
主所應許的順便成就.
有人以為他是單言.
其實不是單言.
是寬容給你們.
不願有一個人沉淪.
乃願人人都悔改.
但主的日子要好像賊一樣.
那一天就是主再來的一天.
天必大有享聲廢去.
有形則的都要被烈火消化.
有形則.
即是物質的.
所有物質的.
都要被烈火融化.
地和其上的物.
即一切.
眾生.
都要消盡.

$^{601}$這一切既然都要如此消化.
你們為人.
該當怎樣聖潔.
怎樣敬虔.
切切仰望神的日子來到.
原來那個的審判.
都和我們地上每一個人有關聯.
即是我們意識到有這個審判.
我們就要在神的救恩.
真理.
那個愛的生命裡面.
來過一個敬虔的生活.
所以這裡特別強調的.
本來物質的世界的消亡.
就應該和我們道德沒有什麼關聯的.
但是彼得就是將基督教.
基督徒那個的敬虔生活.
和將來我們去見主.
又是拉上一個密切的關係.
所以這裡真是值得我們再一次來想.
有形則都要被烈火融化.
但是我們照他的應許.
我們有神應許的人.
即是我們在救恩裡面的人.
就盼望的是新天新地.
有義居在其中.
不是一個虛無縹緲的狀態.
是一個有主 有真理.
有公義的永存的地方.
就是神的天堂.
就是這個新天新地.
不再是舊有的了.
所以這裡又很詳盡的.
來藉著一些人的基俏沒有復活.
彼得又帶進了一個.
從一個審判的角度去講那個新天新地.
那個是一個真是Brand New的.
不再是修修補補的.
是神的日子來的.
一切都要在神的憤怒底下.

$^{641}$完全的消亡.
但是唯一不消亡的.
就是我們這班人.
就是我們這班屬神的兒女.
所以這裡又是很值得我們去看到.
這個多重要的責任.
一方面我們的為人就要敬虔.
還有聖潔.
切切去仰望神的日子來.
就是說地球一定是會毀滅的.
但是我們其實在我們成為神兒女之後.
我們就有這個盼望.
所以這個盼望會激發我們地上的生活.
是要貼近神的心意.
所以第十四節.
親愛的弟兄.
你們的盼望這些事情就當是恩勤.
使自己沒有玷污.
無可指責.
安然的見主.
並且要以我主長久忍耐為得救的恩由.
就是說我們在哪裡跌倒.
就要在哪裡再起床.
來繼續追求神要我們達至的標竿.
不要灰心.
不要喪志.
我們應該來到.
因為主是長久忍耐.
成為了我們的力量.
成為了我們繼續可以做好的那個恩由.
跟著就如我們所親愛的兄弟保羅.
照著所賜給他的智慧.
寫了信給你們.
他一切的信上都是講論這些事情的.
信中有些很難明白的.
那些沒有學問又不堅固的人去強解.
好像強解別的經書一樣.
就自取沉淪.
所謂強解似乎是.
bypass了保羅的危言聳聽.

$^{681}$真理的啟示和教導.
就是掉以輕心.
根本就是不會當是神的聖經的真理.
你不信就要沉淪.
就是你否定了耶穌是唯一的人類救主.
你不要你就沉淪了.
是自取沉淪.
所以這裡就透過彼得這一大段.
又給我們再闊一些思想.
主在來.
沒錯.
是應該成為了我們的安慰和盼望.
因為有復活的日子.
但是彼得這裡似乎是一個.
頗有一個機動性去追求.
在逃避將來審判.
特別是用熱火來審判的結局之前.
我們基督徒就要來到去.
在我們的人生貼近主的心意.
你們為人該當怎樣聖潔.
怎樣敬虔來到去切切盼望主的日子來.
所以大家都是切切盼望的.
我們的盼望就不是進入那個火焰的火湖.
而是進入那個有意在其中居住的新天新地.
所以又是兩個極端.
就是兩個不同的歸宿.
基督徒的歸宿就是新天新地.
就是講到那個意義.
所以這裡就沒有講那個復活次序.
彼得就在這裡強調的.
反而他就強調要審判那些機肖的人.
我們如果在這個時代都有人機肖.
我們就要很小心.
我們首先就是不要自己去自我的強解經書.
就是說慢慢你都接受了.
是的 我上了書幾十年.
真的沒有見過基督教有什麼改變.
真的沒有怎麼見過基督徒特別敬虔.
特別盼望天堂.
都是和我們一樣的.

$^{721}$和不信仰的人都一樣的.
有時候我們不知不覺.
我們自己都會矮化了我們自己的信仰.
那個盼望.
所以彼得這裡其實是挺重的.
告誡信徒都要看到.
原來要教我們看.
原來主看一日如千年.
千年如一日.
不是他沒有時間觀念這個意思.
而是那裡上下文是在講.
他等待 本來他應該來的.
但是他不來.
但是他要來的話隨時都來得到.
是這個意思.
所以我們就應該要以我們的主長久忍耐.
為我們得救的恩由.
所以我們都是繼續在哪裡跌倒.
就在哪裡起來.
什麼時候沒有再參與事奉.
就再起來事奉.
不要蹉跎我們的人生.
做一些在天堂裡面.
一些價值都沒有的工作.
或者那些事情.
我們不是要去只是好像賣告白.
勸人哄人事奉.
而是每一個弟兄姊妹你事奉.
其實是為你自己個人.
其實就不可以說是措積分的.
但是其實真是為我們個人.
我們自己不斷在那個所謂圍爐取暖.
其實好像我們那個炭.
不斷在那個火爐裡面燃點.
如果我們離開那個火爐.
那個炭無論多大塊燒得多旺盛.
一定會慢慢那些灰會令到它熄滅的.
我沒有見過一個人.
我不回教會的.
我不需要再跟你們回教會的.

$^{761}$但是我的靈性是超級好的.
好過你們所有傳道人.
好過你們所有弟兄姊妹.
我從來都不相信這件事.
根本從來都不會這樣發生.
肯定的不用想一定不會發生的.
一個離開了教會的人.
其實他不可能有信仰的.
所以這個要很小心.
一定要以我主長久忍耐.
為我們得救的恩由.
至少在教會裡面.
講道都是一個每個星期的提醒.
或者彼此的勸勉.
尤其是更加有意識地.
我要強調.
弟兄姊妹是要有意識地來崇拜.
你要看神 你要看人.
你看看不同的人的敬拜.
讓你自己回去沉澱.
去舉一反三的反思.
你不要混混樂樂的.
來到都沒有所謂的.
這樣你學不到東西.
是要有意識的.
看到每個人的心靈的.
這個人他經歷了一些事情.
那個內心真是踴躍了.
有些人現在開始面對一些困難.
開始有些內心有些壓抑了.
有些人還是子罪甘霖.
完全不怕自己不怕神的.
他都來聚會的.
這樣就變成了令人高視.
你就要為他擔心.
因為他不懂得為自己擔心.
他不懂得為自己的靈性擔心.
不是他現在不好 可能很好.
但是牧者就有這些眼光.
超越了他.

$^{801}$好像你看到他裡面的那個靈魂.
所以如果我們弟兄姊妹你成熟.
我希望大家成熟.
來到教會 不是批判.
是你去感受一下.
用你的靈去感受一下弟兄姊妹的那個靈.
這樣你就會引發你去感恩去禱告.
這些不是小事的.
這些我們不是在說驅魔的靈界的那些東西.
是很實在很自然的.
是你所熟的.
是你所愛的弟兄姊妹.
有時候你聽他講幾句話.
你感受到他的那個靈性.
這些其實是重要的.
所以這一段經文其實是挺重的.
對基督徒的警告.
我們要小心.
固然我們不要是那些機腸沒有復活的人.
但是我們不要慢慢下去.
就好像萬物和現在一樣.
一樣就是甚麼都不用擔心.
沒有事發生.
但如果不是.
如果真的有一個將來火燒的宇宙要毀滅.
我們真的那個叫做不知道甚麼地球.
流浪星 流浪地球.
記得了.
就是那個地球遲早要越來越近太陽.
是會毀滅的.
所以科學家就想辦法.
怎樣用那些核爆引走那個地球.
拖走那個地球.
是科幻片.
但是那個意思是很有意思的.
就是知道地球一定會毀滅.
但是就用科幻電影的題材做給我們看.
應該是吳孟達也有做的.
不過這裡說的是彼得說.
將來真的宇宙都要毀滅.

$^{841}$我們是.
其實我們不是在那個行業裡面.
被拖下去那個火壺.
我們是一班有新天新地的盼望的人.
所以我們現在只需要在這個地球上.
為主來做這個見證.
好 和大家再看《若人一書》這一段.
《若人一書》是三章.
所以其中很多都是不同的作者.
不同的使徒.
講到關於主的再來.
就著在他們那個牧羊的範疇.
去帶進主的再來.
我們要注意的地方.
第三章的《若人一書》.
就是這位老若翰寫.
《若人》中的作者.
就是這個漁夫若翰.
就是這個.
不是安德烈.
雅各的弟弟.
可能是哥哥.
第三章.
「你看父賜給我們是何等的慈愛.
使我們得稱為神的兒女.
我們也真是他的兒女」.
似乎這個是和一些信主.
都久的弟兄姊妹去講神的慈愛.
就是說不是一種理論.
是你經驗過神何等的慈愛的.
我們是被他牧養過的.
我們是在生活裡面經驗過神.
在我們生命裡面的那個建立.
那個幫助.
我們感受到祂的愛的.
這個基督徒要感受到神的愛.
是很重要的.
如果不是的話.
我們就硬邦邦.
我們跟祂好像很抽離.

$^{881}$很抽象那樣.
那就不健康了.
所以有些人你聽他祈禱.
他一講天父.
憤愛是有感情的.
憤愛是感受到他跟天父之間的密切.
這個就是第一節所講的.
「你看」.
「看」這個字不是就這樣翻譯的.
是細心的留意的意思.
是contemplate這個字.
「你看父賜給我們是何等的慈愛.
使我們得稱為神的兒女」.
慈愛都已經夠了.
夠大了.
「何等的慈愛」.
作者約翰用到很大的形容.
就是說似乎在他的內心很澎湃的.
感受到神對人類.
對神的兒女的愛.
現在這位魯約翰.
他想一下跟耶穌在幾年的傳道.
見過耶穌做過這麼多事.
現在他似乎在這個時候.
他是唯一死剩的仕途了.
他流亡出去突摩海島.
然後回來在二不所教會.
來到中路.
不知道哪間二不所教會.
是不是我們去的遺址那個地方.
那個廣場那個劇院.
可以坐三萬二千人.
很大的 那個時候.
那時候沒有那麼多人.
人口都不會那麼多.
但是那個劇場都可以坐三萬多人.
約翰在晚年的時候就在那裡墓會.
所以他要告訴我們.
這位神的慈愛是何等的慈愛.
我們能夠成為神的兒女.

$^{921}$我們也真是祂的兒女.
西人 這個當然是說.
似乎是未信的人.
或者否認信仰的人.
二不所教會附近就有這個.
那個Diana 即是那個什麼神子.
我忘記了名字了.
抵文是那個銅像.
我想起了.
最大的那個神像就是在二不所.
我們都有拍照的.
西人所以不認識我們.
是因為未曾認識祂 就是神.
親愛的弟兄 我們現在是神的兒女.
將來如何呢?.
還未顯明.
還未顯明就是說.
我們那個形態未轉變.
但是我們知道.
主如果顯現.
我們就必要像祂.
因為必定會見祂的真體.
凡是向祂有這個指望的.
就立即再提醒我們自己.
有意識地去打理我們自己的心靈.
潔淨自己.
好像祂潔淨一樣.
這個就是作者約翰.
一方面訴諸於就是.
我們要和神去比較.
和神的愛去比較我們自己的身份.
新的生命.
不要和西人去比較.
西人是不認識我們的.
他不會覺得我們有光環的.
他看不到的.
西人所以不認識我們.
原來因為不認識祂.
不認識做物主.
但是約翰就很想我們.

$^{961}$我們要去留心神的何等慈愛.
我們就會不知不覺.
對我們自己的新生命.
新的身份.
這個身份帶給我們內心的盼望.
就會有一個很不同的踴躍.
記住是西人不認識我們.
因為他們不認識神.
反過來我們認識神.
我們知道祂給我們的愛是何等的大.
所以我們第二節就是說.
親愛的弟兄.
現在是神的兒女.
將來如何還未顯明.
就不是說將來我們是不是神的兒女.
就不知道了.
沒得說了就不是這個意思.
將來是說身體的改變.
現在未顯明.
即是未是顯現的日子.
但是主顯現.
我們就會見到祂.
而且會見到祂的真體.
這個就真的很奇妙了.
我們將來身體改變之後.
就會見到永恆.
見到主耶穌.
現在就見不到了.
現在頂多只能在異象裡面.
異象裡面見到.
我真的不知道是真是假.
似乎是真的.
因為真的對她的生命有改變.
以前我的舞會.
有一位姐妹做護士.
她崇拜.
看到天使.
我不知道是真是假.
她說是真的.
特別很多來流的人.

$^{1001}$我聽牧者說過.
老前輩說過.
有些人跟她說.
幾十年前.
香港應該是深水Po .
來探望她.
她都來流.
她跟牧者說.
指著天上來流.
一些來流的人.
已經差不多可以.
還有短時間就進入永恆.
似乎神都給她有一個.
見到真體的狀態.
是開心平安的見到.
是很快樂的.
不是很驚恐的.
我聽過這兩件是真事.
如果我姐夫.
我為他祈禱.
我很少這樣祈禱.
你見到光 你見到耶穌.
還有一個.
你見到Lily 她的大女兒.
你跟著她.
五分鐘後她就離開世界.
安全地離開.
我都是很少這樣.
為來流的人這樣祈禱.
但其實某程度上都應該是對的.
因為這裡說.
主耶穌來的時候.
我們會見到祂的真體.
一定是生命改變才會見到.
地球上的人應該見不到.
而且這是一種很美麗.
很平安的那種福樂的見到.
所以我們就進入永恆.
帶著這樣的快樂進入永恆.
我們見到的是一位迎接我們的救主.

$^{1041}$西人見到的就是審判的公義的那一位主.
所以兩者是不同的.
所以凡是向家有指望.
地上我們有意識地要見證自己.
天堂是不可以容許罪惡.
所以我們有意識地要在地上見證.
我預備的聖經有說的.
就是這些經文.
我們可以先喝水.
我們有時間的.
我們第五章又有另外一個很重要的渲染.
如果我們能夠比較仔細地逐段看.
我們都會得到一個結論.
這本書保羅對一個初生的教會.
都是有一個很重的父親的愛的那種著緊.
又好像媽媽去養自己的嬰孩那種的體貼.
那種的關心.
要打發提摩太去留意他們會不會因為逼迫.
而後退冷卻了那個信心呢?.
保羅都知道見到他們在馬其頓的地方.
都成為一個好的榜樣.
2000年前的貼壽羅尼加都是相當富庶.
因為他本身的地區就是地中海.
那裡是地中海.
再出一些才是地中海.
那個地方是一個很富庶的一個港口.
所以一定是很繁盛的.
歐洲的頻繁的一個商業的運輸的地方.
所以他們都要留意一些的日常的生活.
上了主之後的生命.
怎樣能夠配合在生活裡面的轉變,轉化.
第五章:弟兄們 論到時候日期.
不用寫信給你們.
因為你們自己都明明曉得.
主的日子來到.
好像夜間的賊一樣.
人正說平安穩妥的時候.
災禍忽然臨到他們.
如同產難臨到懷胎的婦人一樣.
他們絕不難逃脫.

$^{1081}$這樣的比喻很特別的.
又說賊 又說快生孩子的一個婦人.
似乎是強調準備.
差不多你的懷胎是十個月.
古代都是十個月的.
你不會古代是三個月.
現在是十個月一直累積的.
神的結構人類是十月懷胎的.
最多是差兩天.
或者有些可能是不足月.
無論如何大體上都是十月懷胎.
為什麼會忽然呢?.
為什麼會沒有留意呢?.
災禍要臨到 突然間臨到呢?.
接著就說那個賊會夜間來.
似乎那個家主一定是.
有些是身家 有些是家財.
他才會擔心.
如果家逃四壁的.
那個賊來都不會偷不到東西的.
那不用擔心.
一定是他一定會有足夠的防護.
保安和警醒.
可能派一些僕人守住某些地方.
可能是那個金庫.
或者是那個時候沒有革萬.
怎樣來守住.
不要被賊進來你都不知道.
應該就是這樣的.
但是似乎這裡告訴我們.
用這樣的比喻去強調.
明明有這樣的危機可能會發生.
那個家主都掉以輕心.
完全不理.
那個懷胎的婦人就更加誇張.
就快要生了.
什麼都沒有收拾.
什麼都沒有.
我當然不知道古代的分娩.
醫學常識去到哪裡.

$^{1121}$但是也不是說每個人生出來都有事的.
可能順產的機會更大.
雖然那時候古代的醫院沒有那麼先進.
但是神都會讓那個生命力更加強悍.
就算是一些窮的國家.
譬如非洲那些生存率很大的.
可能後天戰亂瘟疫.
反而後天的死還多過那些夭折.
這裡就是正式告訴我們.
一個臨到產期的婦人.
懷胎的婦人居然沒有去做任何的準備.
似乎是很難想像.
但是保羅和不同的使徒都強調.
主的日子來好像是表面就是忽然間.
因為我們從來都不去想.
在我們有生之年都不去想主的再來.
我們當然不是叫你每秒每分鐘去想.
譬如每天可能都是奢侈的.
可能每一個月或者每一個季.
我們去想一下主的再來和我們的生命關係.
或者我們是不是走前到主的再來的那些行列.
因為已經越來越好像脫鉤那樣.
越來越遠遠的墮後.
這個是我們人生的軌跡.
我們一定是走去死亡的路.
走去那個地方.
但是我們又要想一下.
死和主再來可以是同時間發生的.
所以如果我們地球的人只是想到死.
想到我們那個晚年的退休保障.
而我們從來都沒有想到.
我們應該要準備主再來.
而我們不會忽然被發現.
這裡就是保羅在第五章提到.
是忽然的降臨.
我們是沒有準備地面對這個日子.
第四節.
「你們卻不在黑暗裡.
要那日只臨到你們像賊一樣」.
原來不是那個實際的物質的準備.

$^{1161}$原來是心靈的光的準備.
是心靈裡面要活在日光.
不是活在黑暗.
不是你白天就警醒.
晚上要睡覺的那個黑暗.
很清楚我們怎樣逃避那個忽然的降臨.
叫主的日子好像賊一樣.
原來就是我們在屬於黑暗.
我們是屬光明.
很清楚了第四,第五節.
「你們都是光明之子」.
「都是」每一個基督徒都是.
「都是白晝之子」.
特別強調光明和白晝.
是說心靈的光.
不是黑暗.
不是活在一個黑暗的裡面.
所以我們是不屬於黑夜的.
也都不是屬於幽暗的.
所以我們不要睡覺像別人一樣.
所以一定不是字面的意思.
總要警醒緊守.
因為睡了的人在夜間睡.
醉了的人在夜間醉.
但我們既然屬乎白晝.
就應當緊守.
所以不是說物質的白晝夜晚.
是說心靈的白晝夜晚.
我們的心靈.
要怎樣去讓我們的心靈.
常常在白晝是光明之子呢?.
接著就說信望愛.
和第一章說的信望愛又有不同.
所以保羅用的信望愛.
原來是很生動的.
原來是很遷就著那個處境.
這裡說的是怎樣逃避那個黑暗.
怎樣逃避我們的心靈進入黑暗.
怎樣讓自己是光明.
是白晝.

$^{1201}$是能夠常常有這個盛載主載來的盼望.
接著就說信望愛了.
所以我們又要學這些信望愛了.
所以不是念口枉的.
不是一成不變的.
這個我們以前也說過了.
第八節.
既然我們是屬乎白晝.
即是我們從救恩那裡.
我們得到拯救.
我們得到靈魂的救恩.
我們就是神的兒女.
我們已經是白晝之子.
所以就應當怎樣呢?.
怎樣緊守呢?.
這裡說用信和愛.
當作護心鏡遮空.
得到盼望信望愛了.
當作頭盔戴上.
護心鏡頭盔.
似乎說的是一個戰士.
一個羅馬的軍人.
一個保羅能夠意識到.
那些有裝備的士兵.
似乎當時的羅馬軍人是最精銳的.
和最有紀律.
所以他們有的長矛.
來打仗的時候.
或者刀.
如果你有一個護心鏡.
不知道是什麼物料.
不知道是幾厚的防彈.
或者是避彈衣.
或者是盔甲.
無論如何.
保羅很想強調心靈和思想.
那個重要.
對真理的捍衛.
這個真理不是頭腦的知識.
是進入心靈裡面.

$^{1241}$讓我們有信有愛.
在我們的人生.
似乎不只是自己的信和愛.
是對人的信心.
愛心.
哀莫大於心死.
都是幾真實的.
一個人悲哀到最極致.
那個人的心都死了.
沒有什麼可以再改變他.
是很悲慘的.
人類最悲哀的事情就是心死.
所以這個告訴我們.
第八節又是重要的金句.
我們既然一熟蝴蝶咒.
就應當謹守.
把信和愛當作護心鏡遮空.
所以真的要問一問.
我們信了多久.
我們的信和愛.
有沒有很活潑地增長呢?.
成為一個護心鏡.
我們的心房是最重要的.
心靈是和神溝通的地方.
心靈承載著對神對人的.
很重要的信念的地方.
信仰的地方.
所以千萬不要心硬.
法老實在令到他心硬.
一開始他也害怕.
他的神蹟很厲害.
水面血 他很害怕.
叫摩西快點改變.
我準備什麼什麼.
摩西一走完神蹟.
他《聖經》出埃及記第七章第八章.
經常這樣描述.
摩西走完神蹟之後.
法老就張身一硬.
他原來沒事的.

$^{1281}$原來是民間術士.
來到去笑摩西.
下次呢.
接著呢.
水面血 接著呢.
青蛙.
然後是蝨子.
蒼蠅.
粟就是那些.
可能是瘟疫.
粟生的瘟疫.
或者是豬瘟等等.
蛙 蝨 蠅 粟.
接著蒼 蝨.
蝨就下蝨.
黃就黃蟲.
黑暗 全個埃及都黑了.
實在令到法老越來越心硬.
後來他發現連他自己的術士都做不到.
他都不肯改變.
最後要用滅掌字.
他才肯放以色列人走.
所以其實摩西的十災.
是令到法老越來越心硬.
所以護心鏡原來是令到我們人生.
不要心死.
不要覺得基督教都是不外如是.
隨便就好.
不要這麼熱心.
信和愛要用護心鏡.
所以信心 愛心.
在這裡是主在來.
你不要變成你忽然間才發現.
我沒有準備.
本來說不是的.
我們要時時以信和愛.
當作一個護心鏡.
來遮住我們最重要的部位.
我們的心房.
似乎不是物質的心臟.

$^{1321}$擔心別人偷你的器官.
應該不是這樣.
而是我們的心靈.
心靈真的可以同時溝通.
然後就是將得到的盼望.
當作頭盔戴上.
似乎那個狀態.
時時都有這個盼望.
不是過去式的.
也不是將來的.
是現在的.
這是得到的盼望.
這樣用信和愛.
不是對抗的武器.
是護衛自己的頭和心是最重要的.
至少是一個重要的器官.
但這裡說的是心靈和思想.
和對真理的融會貫通的形式.
所以你看保羅這樣用經文.
信望愛.
不要看那個日子是忽然來到你.
時時都是你的渴望裡面.
令到你那個信望愛.
不斷地有一種增長.
這就是保羅的意思.
原因就是第九節.
因為神不是預定了我們受刑罰.
是預定了我們.
我們的主耶穌基督得救.
這個得救不是信耶穌的得救.
是再來的主.
我們被祂去迎接.
這個就是得救.
所以這裡告訴我們.
我們是預定了.
要進入那個海船的航列.
你想想.
你可以復活.
升去見主耶穌.
升天.

$^{1361}$不是鎖飛機等等.
是你自己轉變之後.
就可以去到天上.
所以我們是預定得救的.
第十節.
祂替我們死.
叫我們無論醒著睡著.
都與祂同活.
或者某程度就是將信望愛這一段.
去用這個生活的解釋.
我們無論醒著睡著.
無論我們是有意識的.
無意識了.
因為我們需要睡眠.
去調整我們的節奏等等.
我們都是與主同活.
我們都是屬主的兒女.
所以你們該彼此的勸慰.
互相的建立.
正如你們素常所行.
就是在教會的生活裡面.
彼此的扶持.
彼此的勸勉.
重點都是建立.
互相的建立.
建立可以包括擇地.
當然這個擇地背後是有關係.
是有愛.
是要很小心.
不是一種論斷性.
指麻式的所謂的建立.
真正的建立.
是父母心.
是出於一種愛.
是出於一種神秘的感動.
所以這裡又是再一次這樣說.
好像回應上人所說的.
復活的次序.
但是寶來說.
事後日期不是我們可以知道的.

$^{1401}$主的日子來只有一個形態.
就是對等候的人.
是一個歡喜的迎接.
對從來都沒有準備.
活在黑暗的人.
是慘 突然降臨.
是慘難的苦人.
突然要生.
那種束手無策.
雙不及防備的那種狼狽.
那種不能回頭的.
好像十個童女.
有五個沒有預備由.
她回來的時候門已經關了.
那五個童女就被趕在外面.
哀哭切齒.
那五個聰明的童女.
馬太芬25張.
就預備由和主人一起在內室.
來享受主人回來的喜悅和快樂.
又是主人所說的.
在你何時沒有為意的時候.
主就來了.
這段經文又是很清楚地說.
我們在人生要怎樣保護我們的頭腦.
就是思想.
對真理的認識和心靈.
讓我們能夠敏銳神和人的需要.
或者自己的需要.
所以需要有護心鏡.
需要有頭盔.
和以弗所書所說的軍裝又不同.
完全是一樣東西.
再多了很多描述.
好像一個軍人的全副武裝.
所以這裡不是說軍人的全副武裝.
是借用護心鏡和頭盔去說到「信望愛」.
說到「信望愛」在主回來的時候.
那個重要.
原來那個人會養育自己的心靈和思想.

$^{1441}$就是護心鏡和頭盔.
養育著很敏銳主的再來的情況.
敏銳就是很融會貫通.
他又是那個家主.
在主人還沒有回來之前.
他就會按時分糧給家庭.
按時的.
根據不同的季節.
按時的.
季節是不同的.
你早上和你吃火鍋.
吃麻辣火鍋.
是不是傻的.
早上起床吃麻辣火鍋.
就算四川人都不行.
所以按季節.
冬天吃夏天的東西都不適合.
夏天如果吃一些很適合夏天的東西.
你就覺得很舒服.
譬如很美味的西瓜.
你已經覺得很開心了.
所以原來那裡說的按時分糧.
是按季節.
因為家主那個僕人.
他要知道這個時候冬天了.
可能要吃羊腩煲.
你可以想像到.
按季節的分糧.
還有身體又是按季節的.
我今天胃又有點不舒服.
很不舒服.
所以吃一些東西都想反胃.
你還給我弄一塊炸雞.
不是那個炸雞有問題.
是我的腸胃今天有問題.
又是按季節.
所以按時分糧其實很講究.
如果一個很有經驗的牧者.
按季節去牧養.
不同的人他現在是什麼處境.

$^{1481}$有些人剛剛失戀.
有些人是剛剛喪流.
有些人可能剛剛升職.
不同的季節其實都不同.
所以這裡告訴我們.
我們要養育自己的心靈.
要有順忘愛.
如果心靈既然是敏銳.
我們都會敏銳到其他人的需要.
我們會將神給我們的恩典去分享.
按時分糧的忠心.
有見識的管家.
忠心和要見識.
好像見過世面那樣見識.
其實是對人性的一種了解.
所以真人一個有經驗.
有牧養的情懷的牧者.
世上難求的.
很難找到一個人.
拼命為整個教會.
為那一代的人去守望,去禱告.
這些是絕對不可以缺少的.
一個教會或者一個時代.
所以這裡告訴我們.
我們自己要追求.
我們的順忘愛.
所以不是公式化的.
基督教就是順忘愛.
其實不是那麼簡單的.
保羅這裡說的是.
為什麼有些人會在主的日子.
突然間淋到他呢?.
因為他從來都是活在黑暗裡面.
但是保羅說我們是光明.
我們是白晝之子.
我們不屬於黑夜的了.
所以我們要怎樣去護衛自己.
我們要怎樣謹守呢?.
就是以信和愛去作為護心鏡,遮空.
所以第五章第八節.

$^{1521}$又是送給大家的金句.
我們很少人會這樣留意自己的心靈.
有個護心鏡.
我們很少人會戴頭盔.
說錯話戴頭盔就有.
說笑的那些戴頭盔.
但是保羅教我們要戴頭盔.
但是是得救的盼望.
對將來永恆.
我們是有份於那個地方.
而且是永恆的福樂地方.
最佳兒臣的地方就是天堂.
我們要有得救的盼望.
作為頭盔戴上.
但願神都讓我們能夠.
從這些保羅的勸勉裡面.
找到原來主在來.
是和我們息息相關.
是真的很切身.
一直提到那個時候的日期.
就突然帶回我們.
再要留意貼壽羅尼加的弟兄姊妹.
你們已經不再屬於黑暗.
現在我們是屬白晝.
我們是屬於神.
所以我們就是要有信,有愛.
去有護心鏡.
要有得救的盼望.
作為我們的頭盔.
來捍衛自己的信仰.
第十二節.
弟兄們,我們勸你們敬重那在你們中間奴婦的人.
就是在主裡面治理你們.
勸戒你們的.
又因他們所作的功.
用愛心格外尊重他們.
你們也要彼此和睦.
我們又勸弟兄們.
要警戒不守規矩的人.
勉勵灰心的人.

$^{1561}$輔助軟弱的人.
也要向眾人忍耐.
你們要謹慎.
無論是誰.
都不可以以惡報惡.
或是彼此相待.
即是以牙還牙那些.
或是待眾人.
常要追求的是良善.
要常常的喜樂.
不住的禱告.
凡事的謝因.
因為這是神在基督耶穌裡.
向你們所定的旨意.
大體上都是想在教會裡面.
建立一個和睦的群體.
一個大家可以彼此去勸勉的群體.
所以仍然是教會的生活.
應該是每個星期的優先.
不是一個可以有可以沒有的.
一般就是六天的工作.
現在是五天.
五天的工作是最重要的.
佔據我們差不多每天八個小時.
某程度上那個時間的涵蓋已經很大.
家庭.
其實似乎這個保羅.
其實都不只是他.
舊約也是.
我們應該以神作為我們的生活中心.
沒錯.
我們在教會的時間可能是一個半小時.
但是重點不是那個時間.
而是我們的心靈.
是應該以神.
以神安排我們在哪一個家.
神的家作為那個中心.
然後就在我們的生活裡面.
落實我們在信仰裡面所得到的.
那個教導.

$^{1601}$然後去實踐.
在教內我們怎樣輔助弟兄.
教外我們怎樣不要以惡報惡.
你會遇到一些不好的事.
遇到一些令到你.
甚至乎官非的事.
我們都要.
總要追求良善.
那就差不多說完了.
對內對外都是要這樣的人生.
你肯定就是.
誰是那個中心點.
一定是以神.
我不是說二錦宣家,洪恩家.
那就太狹隘了.
你應該以神.
作為我們真正生活的最重要的中心點.
然後環繞著神的.
就是我們的家庭,工作.
我們的生活,我們的興趣.
我們一些的休息.
我們一些的運動.
或者我們一些喜好的某些事.
譬如我喜歡行山攝影.
這些都是以神的中心點去看.
我們整個人生.
如果不是這個焦點.
有時候就會.
很容易擺錯位置.
會亂.
所以你會看到.
神的旨意就是要我們.
在我們的人生.
大體上就是我們要常常的喜樂.
就是說.
要令到你喜樂的事情很多.
但是記住.
那些都是會發生的.
但是仍然我們都要喜樂.
就是喜樂可以蓋過了那個憂慮.

$^{1641}$喜樂可以令到我們有力量.
好像撥開雲霧見到神.
來到去.
不要將憂慮擺過了神那個位置.
喜樂就會.
不知不覺你就會.
在我這邊我還怕什麼呢.
那我靠神就可以了.
有些事情真的不是我可以操控的.
神為我解決了它.
所以要常常喜樂.
不住的禱告.
所以這些是連帶關係的.
一個不喜樂的人.
我也很難想像他會祈禱.
長期不喜樂的人.
其實憂慮霸滿了他的心.
其實很難禱告的.
一些憂心,苦痛,痛苦的祈禱.
其實真的萬中無一的.
很難的.
所以不住的禱告.
就不住的謝恩.
又是孿生兄弟.
祈禱就會帶出那個德性.
出路.
其實先在你的心裡面有出路.
先是讓神充滿了你.
請你有平安.
你就會感恩凡事謝恩.
所以喜樂,禱告,謝恩.
其實是互為關係的.
因為這個還要是神的旨意.
所以其實一個以神為中心的人.
他就會活在喜樂.
活在一個不住禱告裡面的心靈.
自然他會凡事都會謝恩.
小朋友回家煮飯就容易的.
每個星期去唱兒歌.
回家有人煮飯給你吃.

$^{1681}$很開心的回到幼稚園.
那個兒童主義學.
沒有什麼考驗的.
一路長大開始讀書多了.
開始面對有不同的競爭.
開始見到原來有很多不公義.
開始有些不公平.
其實在考驗我們喜樂的指數.
我們是否能夠真是凡事交托.
我們是否能夠真是可以去對神.
安排我們的人生來去謝恩.
一個誠實的基督徒.
可以有很多限制的.
很多限制不是不可以不謝恩的.
有些人沒有什麼限制.
都不等於他可以謝恩的.
什麼都是他可以自己得到的.
他就不會謝恩.
我多謝自己 多謝誰.
但是一個屬於神的人.
越是有經驗 有經歷.
他的感恩是會與人俱增.
不是僥倖.
真是神的保守和神的恩典.
不要消滅聖靈的感動.
不要藐視先知的講論.
即是聽道.
但要凡事察現.
所以是要帶意識去聽道.
帶著意識去崇拜.
知道自己現在是在朝見神.
即是去演唱會.
以前是螢光棒 現在不用了.
手機就開燈.
關掉燈 大家揮舞.
大家怎樣都不肯走.
有意識地去演唱會.
因為喜歡那個偶像.
沒有人不喜歡的.
我快點走 我很吵.

$^{1721}$不會這樣的.
這麼貴的擺張票就是想他吵.
就是想聽你想唱的歌.
你叫譚詠麟唱英文歌.
個個都走都不及 都走了.
不可能的 一定是愛在心秋.
人們就如癡如醉的 是不是?.
是這樣的.
但是沒有理由我們去崇拜.
我們很多人都沒有意識.
沒有意識去敬拜神.
所以你擺脫不了崇拜帶給你的感動.
我們要有意識地去朝見神.
所以穿得漂亮一點是對的.
是代表著你很認真去看這個聚會.
去敬拜神.
然後有意識地和頂尖去交往.
用心靈交往.
不是用你的眼睛去批判.
是用心靈和人交往.
所以這裡叫我們不要消滅聖靈的感動.
聖靈的感動可以是很主觀和客觀.
可以是很奇妙的.
感動可以是你會看到不同的事物.
令到你有感動.
尤其是一些有為的青年.
你去到一個崇拜.
你見到有很多教友,頂尖姐妹.
有些甚至和你相差很遠.
有時那個感動還來得強烈.
神就說我就是感動你.
為我來服侍.
你就會感受到那種呼召和帶領.
所以這裡說在教會的生活.
尤其是崇拜.
不要藐視先知的講論.
藐視即是你已經有先人為主.
你還未講你已經塞住的鹽.
現在就再容易些.
因為這麼多堂.

$^{1761}$紅因就好些只有一堂.
不過現在都不是.
有Zoom.
我不喜歡這個團隊.
我就在Zoom.
如果在錦宣就好些.
我不會晚堂.
我去早堂.
星期六晚就沒空.
其實原來根本不喜歡那個團隊.
不過不要藐視.
不要藐視即是不要先人為主.
因為下面就說.
但要凡事察驗.
善美的就要持守.
即是不要先人因人而廢言.
是他我就一定不聽.
那你就真是因人而廢言.
應該說每一個牧者.
都有神給他的說話去教導.
所以不要藐視的背後.
是要凡事去察驗.
察驗就是不戴著有色眼鏡的.
去聆聽,去默想.
去根據聖經去判斷.
這樣叫做一個凡事察驗.
接著你會明白就是善美的.
這個人說這件事.
我從來沒有想過.
他這個角度我從來沒有想過.
我要持守.
我要記下他.
我要真是嘗試去體驗一下.
這個叫做善美的要持守.
所以如果你崇拜都不捨不忌.
接著一出門口就已經忘記了.
那你就五十二堂都浪費了.
因為善美的要持守.
聖靈感動了你.
你都不肯去面對.

$^{1801}$過了惡事就要禁戒不作.
相信就是聽到神的道.
這裡說願賜平安的神.
親自使你們全面成聖.
又願你們的靈與魂與身子.
就是整個人全人得無保守.
在我主耶穌基督降臨的時候.
完全無可指責.
這是第五次講主在來.
完全無可指責的是我們的新心靈.
都是很有意識地等待祂的再來.
所以等候主在來不是一種很被動的.
是我們整個的新心靈.
全人靈魂體都是無可指責的.
就是說我們真的很有意識地等候主的再來.
這個不是一個僥倖.
不是基督教的一種僥倖主義.
避世主義.
所以有些異端就要撿好救生包.
就要在教會買的.
接著就去到哪個地方.
等候主的降臨.
那些一定是異端.
一來你計算了日子就錯了.
二來就不是這樣的一種逃生的僥倖.
而是靈魂體的裝備.
所以其實一個離開世界之前的人.
如果他是一個很敬虔的信徒.
其實他的思考.
他的心靈應該是很活潑的.
當然身體一直衰殘.
甚至呼吸都有困難.
但其實那個心靈應該是.
願你們的靈與魂與生子都得無保守.
在主的再來的時候完全是無可指責的.
今晚就說到這裡.
我們一起禱告.
感恩我們能夠處理了.
好幾段重要的主在來的經文.
都提到我們千萬不要活在黑暗.

$^{1841}$成為了黑暗之子.
因為我們已經因為救恩.
我們已經進入了光明的國度.
我們已經是白晝是光明之子.
求主幫助我們能夠真是以信,愛和望.
去建立我們的生命.
是一個活潑的信仰.
使我們能夠在神的豐厚的恩典裡面.
來處世為人.
和成為了在地上的管家.
去幫助有需要的人.
成為一個忠心有見識的管家.
能夠按時去分量.
感恩禱告奉耶穌基督的名.
阿們.
謝謝大家.
\newpage



\section{詩篇 84:1-12}
\label{sec:CutYNgBop3c}
\textbf{2024-06-09 午堂《永恆的居所》詩篇 84\_1-12  講員:馮耀榮牧師}
\newline
\newline
連結: \href{https://youtube.com/watch?v=CutYNgBop3c}{\texttt{https://youtube.com/watch?v=CutYNgBop3c}} ~~~~ 語音日期: 2024-06-11
\newline
\newline
\hyperref[sec:DjV4K1Dg8GU]{\small{< < < PREV SERMON < < <}}
~
\hyperref[sec:index_chronic]{\small{[返順時目]}}
~
\hyperref[sec:index_scriptual]{\small{[返順卷目]}}
~
\hyperref[sec:AT0lX1Ebpa4]{\small{> > > NEXT SERMON > > >}}
\newline
\newline
詩篇 84:1-12
\newline
\begin{longtable}{cl}
\hline
\hline
章節 & 經文 (和合本修訂版)\\
\hline
84:1 & \begin{tabularx}{0.7\textwidth}{X} 萬軍之耶和華啊，你的居所何等可愛！ \end{tabularx} \\ \\ \relax
84:2 & \begin{tabularx}{0.7\textwidth}{X} 我羨慕渴想耶和華的院宇，我的內心，我的肉體向永生神歡呼。 \end{tabularx} \\ \\ \relax
84:3 & \begin{tabularx}{0.7\textwidth}{X} 萬軍之耶和華－我的王，我的神啊，在你祭壇那裡，麻雀為自己找到了家，燕子為自己找著菢雛之窩。 \end{tabularx} \\ \\ \relax
84:4 & \begin{tabularx}{0.7\textwidth}{X} 如此住在你殿中的有福了！他們不斷地讚美你。（細拉） \end{tabularx} \\ \\ \relax
84:5 & \begin{tabularx}{0.7\textwidth}{X} 靠你有力量、心中嚮往錫安大道的，這人有福了！ \end{tabularx} \\ \\ \relax
84:6 & \begin{tabularx}{0.7\textwidth}{X} 他們經過「流淚谷」，叫這谷變為泉源之地；且有秋雨之福蓋滿了全谷。 \end{tabularx} \\ \\ \relax
84:7 & \begin{tabularx}{0.7\textwidth}{X} 他們行走，力上加力，各人到錫安朝見神。 \end{tabularx} \\ \\ \relax
84:8 & \begin{tabularx}{0.7\textwidth}{X} 萬軍之耶和華神啊，求你聽我的禱告！雅各的神啊，求你側耳而聽！（細拉） \end{tabularx} \\ \\ \relax
84:9 & \begin{tabularx}{0.7\textwidth}{X} 神啊，我們的盾牌，求你觀看，求你垂顧你受膏者的面！ \end{tabularx} \\ \\ \relax
84:10 & \begin{tabularx}{0.7\textwidth}{X} 在你的院宇一日，勝似千日；寧可在我神的殿中看門，不願住在惡人的帳棚裡。 \end{tabularx} \\ \\ \relax
84:11 & \begin{tabularx}{0.7\textwidth}{X} 因為耶和華神是太陽，是盾牌，耶和華要賜下恩惠和榮耀。他未嘗留下福氣不給那些行動正直的人。 \end{tabularx} \\ \\ \relax
84:12 & \begin{tabularx}{0.7\textwidth}{X} 萬軍之耶和華啊，倚靠你的人有福了！ \end{tabularx} \\ \\
[1ex]
\hline
\hline
\end{longtable}
$^{1}$各位早晨.
今天我出來講道的時候.
我太太就吩咐我多帶一對襪子.
因為在長洲.
我們在神樂縣的山頂.
要下山的時候.
多數是濕透了的.
因為她下來的時候.
她有些水會沖下去的.
所以很難避開流水.
今天我也多帶了一對襪子.
幸好出門之前開始沒什麼雨.
到我出來香港.
搭地鐵又開始多了一點雨.
來到這裡.
經過湖邊的時候.
就下了一點雨.
所以我最後在湖邊走來的時候.
是濕了一點的.
我也不需要換襪子.
今天我講的訊息.
和回來教堂有關係.
一個敬拜神的歷程.
由家裡或者起步的地方進去敬拜.
其實需要一點時間.
這種朝聖的歷程.
最高峰的時間可能是崇拜的時候.
或者一起祈禱的時候.
但是由家裡回到教堂.
到去敬拜神.
這是一個朝聖的旅程.
在舊約裡面就是這樣的.
其實他們也不是每個星期都回到聖堂.
以色列的聖殿.
可能在耶路撒冷.
一年有幾次節期.
以色列人在不同的地方.
就會有一個朝聖之旅.
如果不是平時.
他們也會在會堂裡面去祈禱.

$^{41}$這是耶穌的那個年代.
如果是舊約的年代.
這首詩篇就是大衛的詩.
所以他多數是住在耶路撒冷附近.
所以當他有這種心境.
去耶華的殿裡面去敬拜.
就和我們都會比較方便.
我們一個小的城市.
我們每個星期都可以去敬拜神.
所以今天和大家分享這首.
也不是很陌生的詩歌.
它分開有幾個段落.
幫助我們進入.
我們每個星期都有敬拜的心路歷程.
首先第一個段落.
他說耶華的居所和耶華的言語.
是很可愛的.
感受堂是一個很蒙福的地方.
因為有一個很好的教堂.
不單止是坐得舒適.
而且是在湖邊裡面.
其實在香港某邊境教會.
特別是在都市化.
教堂的旁邊.
好像恩福堂很大顆.
北宣都是很大的.
但是它旁邊不是很有詩意的.
它是在商業的樓宇.
工業的樓宇.
或者住宅區裡面.
是比較方便.
因為它近地鐵.
你可以不用帶雨傘.
如果你家住在地鐵上街.
你可以不用帶雨傘.
因為由你家坐地鐵.
去到北宣那裡.
北宣還要走一點路.
恩福堂可以衝過馬路.
你就不會下濕.

$^{81}$也可以去到教堂.
但是這裡除非你是開車來.
其實你走過一個階段.
都會經過雨水.
而這首詩篇裡面.
其中有一段是講.
經歷秋雨的陰.
一個叫做雨的季節.
所以這首詩歌很適合大家.
特別在雨季的時候.
回去敬拜神的時候.
可以襯托詩人.
當他處理耶華的殿.
是很可愛的時候.
整個詩歌.
如何將一種心路歷程表達出來.
所以今天我教聖經科是舊約.
我主要是教聖經的故事.
好像撒母耳記.
或者創世記的故事.
我不是很擅長教詩歌.
所以詩篇都是比較難一點.
因為它比較抽象.
它好像很簡單.
但其實要進入他心境.
就沒有那麼分析性.
他比較情感上去襯托.
因為他每一段經文都不是很複雜.
但如何去表達.
詩人去渴求神的時候.
他藉著聖殿.
或者藉著他描寫的.
看見的事物.
去接觸看不見的神.
所以詩篇就是有這樣的表達.
所以你當今天是一種上課.
上課.
而且這個課不是我專唱的詩歌.
但對我來說是重要的.
因為我們敬拜神的人.

$^{121}$最重要都要明白.
聖經的詩人.
他如何可以有一種敬拜神的心路歷程.
所以我都盡量去思考.
和大家分享這段經文.
首先我們的定位.
耶和華的居所是可愛.
這個詩人或者我們敬拜神的人.
是會羨慕耶和華的言語.
這個就是間接性.
我們最重要的就是要羨慕神.
學慕神 與神親近.
但其實中間裡面.
這段經文開始.
其實我們不直接去羨慕神.
我們是間接地.
透過神的居所去接觸神.
如果純粹是居所漂亮.
我們可以建一個很漂亮的教堂.
很宏大的教堂.
其實它和神還是有一個差距.
雖然我們感受堂很蒙福有這樣的教堂.
但如果我們去一個小一點的地方.
小一點的教堂.
我們是不是等於沒有那麼蒙福呢.
如果居所只是一個工具.
它不是真正的終極的學慕神.
我們就可以解決了.
所以我們不一定要很宏偉的教堂.
我們不一定要有一個教堂.
就在湖邊.
我們只要能夠透過教堂.
神的居所去接觸神.
其實都是可以做到的.
所以這個就是詩歌.
他一定是想我們能夠親近神.
但是他透過居所.
是一個工具.
一個間接的東西.
但是它是很重要的.

$^{161}$因為整首詩都是用這種形式表達.
雖然我們最終極的就是與神親近.
但是神的居所.
或者我們安排崇拜時候的那些情緒.
或者我們坐得舒不舒服.
其實這些東西都在影響我們.
很具體地去敬拜神.
如果太冷.
現在通常教堂或者任何地方都很冷.
所以如果沒有東西蓋著.
坐到差不多就開始冷.
就沒有精神去敬拜神.
所以教會負責讓人去親近神的時候.
最終極都是神自己.
但是整個的程序.
或者居所能夠讓人有一個途徑去敬拜神.
這個就是這首詩篇裡面去表達的.
所以我們不需要走回耶路撒冷的聖殿.
因為它已經拆毀了.
但是這首詩歌還是有一個很永恆的價值的.
所以我們在這首裡面.
看慕神的居所.
它是一個工具,一個途徑.
但是這是一個必然的途徑.
所以有些人習慣了網上崇拜.
我們現在有沒有網上.
多數有些人如果習慣了.
現在沒有了.
他就不會回教會.
因為純粹是敬拜神.
在不同的地方.
在家裡都是敬拜.
都可以聽到.
他們不太著重居所.
就是因為他們未了解舊約.
他們為何這麼著重神的殿.
因為他們是和神的關係.
通過神的殿去接觸神很有關係的.
當然在疫情之後.
有些人沒有見很多很久的弟兄姊妹.

$^{201}$回來會熱鬧一點.
崇拜完之後可以聊聊天.
這些都是重要的.
不過這個不是最重要的敬拜神學.
敬拜神學是網上是不可以取代的.
因為裡面是包括了神的殿.
這個居所是神在我們當中裡面.
讓我們能夠透過可愛的居所.
羨慕的耶和華的殿.
去羨慕神,去接觸神.
所以這個都不需要很大力氣.
都對一個敬拜神的人.
可以解釋到那個次序或者重要性.
在疫情或者現在我們神學院教書.
方便很多的.
所以我們有一部份的課程.
都一定會放上網.
因為純粹是知識的交流.
網上是有效的.
有時要來到課室一起.
其實是浪費時間的.
但是如果敬拜是另一件事情.
在神的殿裡面是無所謂浪費時間.
因為你要去神的殿.
那個過程本身就是一種朝見神的過程.
所以在其他的學習.
其他的人生裡面.
特別是比較有效的社會.
最重要是有效.
能夠達到目的.
至於那個過程,那個工具.
我們可以跳過去的.
所以現在來到這首詩篇.
我們開始或者掌握到.
親近神和神的殿.
是有一個很間接的關係.
我們想一想這首詩歌.
這個人他對神的殿.
他投入的時候.
他的心腸,他的肉體.

$^{241}$向永生神呼籲.
我們就要調教一下我們內心世界.
我們在神的殿裡面.
不單止是頭腦或者心靈.
還有肉體.
你的比較感官的部份.
不是精神上的部份.
你都是有敬拜神的心.
所以我們不會像以前那樣.
或者像一些大宗派.
在座位前面有一個跪的地方.
因為在舊約裡面是有的.
去到神面前要俯伏,要跪拜.
我們祈禱.
上一代的人會跪在這裡祈禱.
現在就方便些.
沒有那種肉體上的姿勢.
當然其實只有肉體.
只有具體的神的居所.
如果我們對神的關係不好的時候.
其實也沒有什麼用.
所以我也不認為現在我們一定要跪.
因為如果我們只有外在的肉體形式.
沒有內心世界和神有關係.
這個也不是有什麼大作用.
不過就是因為我們已經相信.
或者我們知道.
我們最終極的就是敬拜神自己.
但是那個過程我們也要反省一下.
聖經裡面的作者.
他用肉體去向神呼籲.
他是在表達行動上.
具體上他是在敬拜神.
這個就在詩篇裡面.
當然心腸,心境也要襯托.
不過我們現代人.
特別是在敬拜裡面.
是比較著重思想或者理念.
所以對我們的肉體的表達.
是比較弱一點.

$^{281}$所以這首詩篇一開始.
有三個層次.
第一個就是神自己.
萬軍的耶和華.
接著是神的居所.
第三個就是敬拜者的身體行動.
除了心靈方面之外.
所以這個也是值得我們去思想.
所以我今天也是來敬拜.
所以我的敬拜開始.
就是我太太提醒我.
要戴多一對襪子.
這個對我來說.
是一種叫做肉體.
實際一點的情況.
我要坐船.
接著坐地鐵.
我年紀大了.
所以我坐地鐵的時候.
是要去元朗的香港站.
因為我自己的舞會是在大角咀.
所以我習慣了.
回自己的教會崇拜的時候.
就在奧幻站出去了.
一出去之後.
我今天好像不是回教會.
不是回自己的教會.
接著我就走回去.
因為去南昌站才轉這個部分.
所以一個人的敬拜.
或者一個習慣.
到時就會有這樣的行為.
所以我其實還在看書.
或者在看聖經的時候.
但是知道大概那個時間.
就要出閘了.
所以我習慣了.
回自己的舞會教會崗道.
就在奧運站.
很自然地就走出去.

$^{321}$接著就想今天我去哪裡.
有些不妥再回來.
所以就沒有遲到.
這個就是.
由你起身去到真正去教堂.
其實中間的過程.
都是一種朝見神的一個階段.
我們需要有些心靈準備進入去.
我們多數香港人是比較有效率.
知道敬拜就是.
有時去到崗道的時候才敬拜.
之前那些都是叫做等人的時間.
等齊人的時間.
因為我們比較著重精神的層次.
不處理我們由起身開始準備.
調教我們的內心世界.
所以有些人回來教會.
去崇拜之前都不坐在其他地方.
因為他.
到唱詩的時候.
我們就開始進入那個狀態.
主席帶領我們進入那個狀態.
其實不應該的.
應該就是我們回來的時候.
可能跟其他人打招呼都可以的.
但最主要就是.
你要面見的是耶華自己.
所以當然主席或者負責領導的人.
是幫我們進入那個程序.
但我們自己去到神的殿裡面.
應該都坐下.
才在神面前有一個祈禱.
有一個安靜.
這個就是我們主動參與崇拜的歷程.
是我們負責的同工或者帶領的人.
他是比較輔助性.
他不是那個主動.
因為朝見神是每一個人的責任.
每一個人敬拜的重點.
所以甚至是講道的人都不是主動.

$^{361}$因為神本身是我們敬拜的對象.
他是主動.
而我們是最重要的一個成員.
不是可有可無的.
不是回來聽聽上堂的事情.
所以敬拜的程序應該由起身開始.
今天我可以去到神的殿裡面.
我準備我的內心世界.
你不需要穿特別的衣服.
或者不需要有什麼特別的.
你的心境去準備.
就好像這首詩歌.
我當爺爺說的句子是很可愛的.
很難得的.
當然我記得我喜歡一個教堂.
我覺得是在香港裡面.
最有印象的一個教堂.
他就是在我讀神學的時候.
我們做監獄工作.
在喜靈州的監獄裡面去傳福音.
我第一次去的時候.
在我分組的時候.
他是在教堂裡面.
喜靈州是一個麻風病院.
是一個基督教或者應該是一個基督教.
即是有宗教的.
所以他開始的時候是醫治一些麻風病人的.
後來香港這方面就沒有那麼多.
他就轉去變成監獄.
所以我人生裡面最有印象的就是那個教堂.
其實都不是可以經常去到的.
我大概有一年的時間.
在長洲坐那些街道.
那些街道都搖啊搖啊.
幸好我是水上人.
所以我不怕潦倒.
不過都挺刺激的.
我們去到喜靈州那裡.
然後有些監獄的車就送我們上山.
那裡的路很遠的.

$^{401}$比長洲遠很多.
長洲在碼頭上到神學院大概十多分鐘.
但我想如果那裡要走上去監獄.
大概都要一個小時.
比較長一點.
所以我們去到那裡.
因為那裡是在山上.
所以它在崖邊裡面有一個教堂.
它對著就是南中國海.
在那個教堂裡面.
它又建築得挺好的.
因為坐在教堂裡面.
有個窗口可以看到外面的海.
所以我覺得比在湖邊的教堂.
它是無敵海景旁邊的教堂.
它更加覺得有一個寧靜.
有一個跟大自然接觸.
其實跟大自然接觸.
你會覺得你是跟神接觸的.
如果躲在這個地方.
你的內心世界都可以想像和去敬拜神.
都可以的.
其實我們的敬拜神可以很心靈.
所以不一定要外在的幫助.
但這首詩歌就是這樣.
有外在的幫助.
亦都是比較正常.
我們一個平凡人.
不是那麼超然.
我們需要見到大自然.
才看到神的創造.
在湖邊裡面看到人的寧靜.
所以我們一個屬神的人.
他不應該靠感官.
他可以很精神性.
靈魂方面去敬拜神.
不過我們屬神的人都是一個平凡人.
平凡人有時都要有一個好環境.
或者一個可愛的殿去敬拜神.
這個我不太確定.

$^{441}$有時教書都會講得不清楚.
究竟我們是著重心靈誠實.
還是著重肉體的感受.
記住我們要找一個較好的教堂去敬拜神.
是肉體的感受.
如果單是聽這部份的訊息.
這個就有些問題.
不是問題 很大問題.
因為你沒有了心靈.
但是我們最沒有錯的地方.
凡事都去到最屬靈的部份.
不用回教會.
我跟神有關係就行了.
我是去敬拜神 不是去敬拜人.
這樣就很絕對.
當然聖經都很清楚.
如果你真的要去敬拜神.
你應該跟弟兄姊妹一起去敬拜神.
跟弟兄姊妹都不是神.
所以如果講得極端一點.
是很精神性的 很屬靈性的.
我們兩邊都要兼顧.
我們怎樣說呢.
如果我們太著重感官.
我們要反省一下.
因為感官是會無優質.
我們希望教堂很可愛.
就應該更加可愛.
重點就是在感官上.
去加深脫離了跟神的關係.
但是如果純粹跟神的關係.
我們就高估了.
我們人都是需要有感官的經歷.
其實都不是很複雜的道理.
但是我們需要有兩邊的互相表達.
我們看看下一個段落的表達.
神是我們最終敬拜的對象.
然後就是神的殿.
第二個段落裡.
這個詩人看到麻雀.

$^{481}$在耶華的殿裡找到房屋.
燕子找到抱初之窩.
這首詩歌就是一個層次.
進入一個比較感官上.
或者比較肉體上的演著.
它由一個耶華本身的永恆.
去到比較神聖的居所.
神的殿.
然後他的眼睛轉移到燕子抱初的窩.
你想一想為何聖經的作者.
在說一個敬拜神的歷程.
耶華的殿語的時候.
會提到一個燕子的窩.
或者鳥仔的巢.
這些都不是很永恆的.
只是一個暫時的地方.
所以我們香港有一個名詞叫做窩居.
而不是豪宅.
我們很多樓宇的名字都是豪宅.
他用他的價錢去定義.
它是好像深水Po 的tong 房 窩居.
所以這個窩居應該都是一些動物.
小的或者脆弱的動物住的地方.
一個很暫時性的地方.
這個詩人不是直接住在耶華的殿裡面.
他是想像一下.
我如果住不到在殿裡面.
但是剛好見到一些麻雀.
在耶華的殿裡面.
捉了自己一個抱初之窩的窩居.
他覺得這是很美好的事情.
我自己住不到在殿裡面.
我可以投射我自己.
好像一些麻雀一樣.
可以在耶華的殿裡面.
有一個暫時居住的地方.
一個窩居.
這個就是將永恆和屬世.
都有一個很詳細的描寫出來.
永恆就是耶華.

$^{521}$中間有耶華的殿.
但是在殿裡面.
有一隻麻雀.
有一個燕子.
有一個居所.
這個就是.
麻雀有多少年的壽命.
它可不可以過到冬天.
燕子有些熟悉的人可以找找.
但應該都不會太長的日子.
所以他們的居所和窩.
都不是像聖殿那樣.
用的材料比較長久.
你可以看到.
作者是由永恆裡面.
看到一個短暫.
一個細微的地方.
一個燕子.
它本質上是脆弱的.
是暫時性的.
但由於它是寄居於神的殿裡面.
所以他都羨慕這個燕子.
不知道大家是否明白.
作者是一個詩人.
他的主題是神自己.
然後是神的殿.
不過他很懂得.
再具體一點.
再實際一點.
因為我們不是每個人都可以在殿裡面服侍.
我們不是每個人都可以.
經常在殿裡面.
我們只可以想像一下.
我們好像動物世界的燕子.
捉了一個蝸居在殿裡面.
我們都很有滿足感.
不知道大家是否明白這個比喻.
所以聖經的作者.
他有永恆和心靈的一個層次.
但他亦都很懂得.

$^{561}$人在肉體上.
在感覺上.
在我們實際上.
有一個蝸居.
這個蝸居可以通往永恆.
可以在耶華的殿裡面.
這個是我們可以調教的心境.
我記得我在英國讀書的時候.
我教會的執事.
他的媽媽都八十多九十歲.
我教會執事是新界人士.
在英國做餐館.
大半生.
後來他媽媽年紀大了.
爸爸過世了.
他接他回英國去照顧他.
但他一來到就開始有點不舒服.
病了.
他們餐館人士不是很多時間可以照顧.
所以來了很短時間.
一起住之後.
他就送了他入老人院.
我就去探望這個老人院.
第一次去探望的時候.
我覺得他的處境.
有點像剛才那首詩歌.
在耶華的殿中找到一個蝸居.
因為他的老人院是一個教堂.
是改建式的.
因為英國由於他們下一代.
即年輕一代很少回教會.
回教會都是上一代的人.
所以很多地方的教會都沒有聚會.
因為他們一個很小的地方.
有好幾間教會.
起碼有聖公會.
然後有秦道慧理公會.
我在英國那幾年.
我們教會買了一個秦道慧理公會.
我們只有八十多人.

$^{601}$教堂可以坐三百多四百人.
很大.
差不多不用錢就給你.
不過你要維修主位的暖氣.
是很貴的.
跟你買的教堂差不多貴.
所以要維修教堂.
維持教堂.
沒有人用或者沒有用.
很容易荒廢.
所以他們就改建成一個老人宿舍.
它挺高的.
它應該有三層.
連電梯也做了.
所以我去探望它的時候.
我就記得這首詩篇.
這首詩篇說在耶華的殿裡面.
找到一個萵苣.
我跟他說.
我們看到這首詩篇.
你是很幸福的.
因為我也做不到.
如果我將來退休.
我要在耶華的殿裡面.
找到一個萵苣.
在教堂的角落.
在崇拜的地方.
可以住在那裡.
住到離開世界見到神.
這個也是.
起碼照著聖經的教導.
因為那個教堂要改建.
也要很多錢的.
能夠24小時住在耶華的殿裡面.
我做牧師在英國.
很可想的.
他們有牧師樓.
他們有一個崇拜的地方.
旁邊有一個小一點的牧師樓.
前面有些花園.

$^{641}$我也很想華西.
我喜歡在長洲住.
看到海.
有個花園.
當然我知道這個老人家.
千里遙遙和家人相聚.
最後沒有在一起住.
其實很可憐.
心境上很可憐.
他還未相信.
所以來到之後.
我們慢慢幫他去相信.
一個人的屬靈的醒覺.
其實是很難預測的.
我教這麼久的主學.
他是我第一個反過來告訴我.
他為我祈禱的.
當教導他信耶穌.
然後叫他讀聖經.
他的眼睛不是很好.
所以找一些弟兄姊妹幫他寫一首詩篇.
一些金句寫得大隻一點.
放在那裡.
每一次去探望他.
他就說我也有為你祈禱.
其實是他安慰我.
因為我安慰他.
我也有些困擾.
我也跟他說.
因為他既然問我有什麼代表事成.
我說我現在要寫論文.
博士論文.
我也沒有什麼頭緒.
不知道怎麼做.
他就為我祈禱.
他其實是相當快捷.
掌握得到敬畏神的心境.
起碼在他人生的遭遇裡面.
家人是對他好.
不過本來是想著重逢.

$^{681}$因為其實相隔了很久.
他這個兒子去了英國都三十年.
到年這麼大的時候.
想著人生最後可以再相聚.
原來又要分開.
當然人在這個處境.
在這個患難裡面.
他的心境可以沉下去.
跟著這個環境去消沉.
但同時人的生命.
是可以看得開一點.
看得永恆一點.
所以這個老人家.
對我來說是一個很大的鼓舞.
他可以是最有效地.
去彰顯這首詩篇講的訊息.
這個是耶華居所.
神是帶我來到神的面前.
能夠回到教會.
能夠信了耶穌.
能夠住在耶華殿.
他以前的失望.
剛剛來到第一次住老人院的時候.
失望很快能夠轉變.
這首詩篇裡面.
所要表達的.
這首詩篇我們在心裡面.
在耶華殿裡靠神幫助我們.
我們能夠在朝聖.
在這個食安大道裡面.
覺得自己是有福的人.
這首詩歌說了.
我們其實有福.
但並不是表示我們環境上是純利.
它告訴我們.
其實我們的環境.
是走過死刃的幽谷.
不過我們當這個不是死刃的幽谷.
是一個秋雨之谷.
它之前說我們是流淚之谷.

$^{721}$所以其實一定是有很多的難處.
身體上,經濟上,社會上,文化上.
所以本來是流淚.
你就可以變成秋雨.
所以你見下雨的時候.
你是有兩種心境的.
你覺得很麻煩.
你會昏迷.
你會不舒服.
會潮濕.
這個是很正常的.
因為始終都是不方便的.
你要很超越.
才能夠由一個不方便.
變成一種秋雨之福.
能夠可以栽種的.
世界上所有的事物.
都有這種兩重的果效.
如果你看得它悲觀.
它就很悲觀.
看得它好.
它就很好.
所以我現在宿舍.
我退休之後.
我本來想搬去一個小一點的宿舍.
小一點的宿舍.
就可以打理少一點.
現在有空了.
不知為何經常都見到.
地上有很多頭髮.
有很多皺紋.
以前不太察覺.
你一打理的時候.
為何這麼多.
其實就沒有多到.
以前沒有時間去照顧.
一有照顧.
經常都要彎腰去整理.
其實都幾麻煩.
所以我就希望.

$^{761}$找一個小一點的地方住.
但現在暫時未搬到.
如果你覺得麻煩.
就算很漂亮的.
很好的地方.
你會覺得很多東西麻煩.
但如果你覺得.
它是一個物質世界.
就算是一個蝸居.
你只要這個蝸居建築.
在耶和華的殿裡面.
在你的表達裡面.
它就是一個福氣.
一個詩歌上的福氣.
所以我們是有能力.
將這個流淚谷.
變成一個秋雨之聲.
它有能力的原因.
就是因為它在學想朝見神.
它在神的殿裡面.
在神的蝸居裡面.
一個懂得敬敬.
一個懂得在星期日崇拜的人.
它就會有能力.
面對你工作上.
你家裡面的.
所有世上的事物.
裡面的流淚.
你都可以變成一個秋雨之聲.
它是有實際果效的.
起碼看這首詩篇.
是帶我們去看出這個過程.
我們不是為了敬拜神而敬拜神.
這個有點錯誤的表達手法.
我們敬拜神.
是在影響我們世界每一日的生活.
我們能不能夠由一種流淚谷.
變成秋雨之谷.
這種轉變的能力.
我們一起在祈禱.

$^{801}$張天賦我們謝謝你.
你讓我們有一個詩人.
他都經歷過人生的.
現實上的難處.
或者流淚.
他直接敬拜你的時候.
他在永恆的居所.
在這個地上.
一個暫時的蝸居.
我們都能夠可以.
去親力你對我們的祝福.
保守.
我們將佐助梅弟兄姊妹.
他們能夠在星期日.
在平常的生命裡面.
去接近你.
去敬拜你.
求你親切的幫助我們.
挑撥我們的心.
去學上你.
我們這樣求.
奉耶穌聖名.
阿們.
\newpage



\section{路加福音 9:46-48}
\label{sec:AT0lX1Ebpa4}
\textbf{2024-06-09 早堂《接待小子的服侍》路加福音 9\_46-48  講員:蔣文忠傳道}
\newline
\newline
連結: \href{https://youtube.com/watch?v=AT0lX1Ebpa4}{\texttt{https://youtube.com/watch?v=AT0lX1Ebpa4}} ~~~~ 語音日期: 2024-06-11
\newline
\newline
\hyperref[sec:CutYNgBop3c]{\small{< < < PREV SERMON < < <}}
~
\hyperref[sec:index_chronic]{\small{[返順時目]}}
~
\hyperref[sec:index_scriptual]{\small{[返順卷目]}}
~
\hyperref[sec:w32Ot_b7FcE]{\small{> > > NEXT SERMON > > >}}
\newline
\newline
路加福音 9:46-48
\newline
\begin{longtable}{cl}
\hline
\hline
章節 & 經文 (和合本修訂版)\\
\hline
9:46 & \begin{tabularx}{0.7\textwidth}{X} 門徒互相議論，他們中間誰最大。 \end{tabularx} \\ \\ \relax
9:47 & \begin{tabularx}{0.7\textwidth}{X} 耶穌看出他們心中的議論，就領一個小孩子來，叫他站在自己旁邊， \end{tabularx} \\ \\ \relax
9:48 & \begin{tabularx}{0.7\textwidth}{X} 對他們說：「凡為我的名接納這小孩子的，就是接納我；凡接納我的，就是接納那差我來的。你們中間最小的，他就是最大的。」 \end{tabularx} \\ \\
[1ex]
\hline
\hline
\end{longtable}
$^{1}$弟子們 早安.
很感恩 雖然今天大家讀剛才的經文都感受到.
外面的雨水或者那種翻騰都是一個大自然很大的力量.
但也感恩我們神的保守 我們都可以聚集.
更加感受到上帝我們一份敬畏對祂.
或者一種對祂的話.
可能我們都預備心靈去聆聽.
跟大家分享這段經文其實都想了很久.
基本上在想怎樣跟大家分享這段經文的訊息.
而跟大家分享這段經文的訊息的時候.
可能大家讀路加福音剛才的經文都會.
如果你熟悉負累福音.
即我們所說的馬太馬可路加福音.
三卷類似的福音書的時候.
都有類近的記載.
待會我主要會用新漢日本來跟大家分享.
這個記載裡面其實就有幾卷福音書.
所以在我的前言當中.
其實都很小心.
不想跟大家處理太多一些很技術性的問題.
也不是很想跟大家說.
其實我們要傳統我們有時候讀經的時候.
就會說參考其他的經文.
就會找回一個真實的情況.
其實我們不需要這樣去讀經.
主要是集中去看我們要讀的那段經文.
其實有什麼訊息要跟我們說.
當我們要無可避免地參考其他的經文的時候.
其實我們的心態都是希望幫助我們.
去了解我們今天讀路加福音這段經文裡面.
其實去比較之下想突出.
作者路加想突出一些什麼訊息.
給我們去認識.
在馬可,馬太都有相近的記載.
一般我們會認為馬可是最早寫成的書.
所以馬可的福音是比較短.
而通常我都會說馬太和路加.
是有參考馬可的說法來表達.
拍幕我們待會所分享的裡面.
都能夠更深刻.

$^{41}$當我們要比較一點點馬可和馬太的時候.
其實更深刻讓我們了解路加福音裡面所說的訊息.
我們一起去禱告.
然後我們都預備心靈去接收上主今天的讀.
天父多謝你讓我們今天.
雖然外面風雨飄搖.
但你也保守中弟兄姊妹.
來到你的殿裡面.
我們感受到那種力量.
感受到你那種威嚴.
也感受到你的同在的寶貴.
盼望你保守中弟兄姊妹的心靈.
讓我們整個崇拜當中.
由詩歌到聽道.
到我們向你回應的時候.
都謙卑地開啟我們.
專注在你的面前.
願你保守我們以下聽道的時間.
同心禱告奉主名求.
阿們.
這段的記載裡面.
小孩子是一個重點.
我們說一個小子.
有時候聖經也有一個Little One.
一個不一定是年紀很小的小孩子也可以.
其實是一個很卑微的人物.
一個這樣的身份.
我先用《馬太》《馬可》裡面的背景.
和大家作一點點的引言.
就是他們也有說到關於.
好像一個小孩的教導.
大家應該也不會陌生.
在《馬可》裡面說.
如果我們不像小孩就不能夠接受天上的國.
將小孩子的特質.
和我們去接受.
去領受上帝的國.
去拉上關係.
大概是說一個孩子那種溫柔謙卑.
那種很單純.

$^{81}$接受好就馬上全心全意的領受.
上帝給他的恩典和國度.
在說這種的特質.
這種接受.
很欣然接受的特質.
去到《馬太》.
我剛才說了通常《馬太》就將《馬可》.
會有一點點的expand.
將它說多一點點給大家聽.
所以除了剛才說的那個很核心的像小孩.
這樣才能夠進天國之外.
好像一個條件或者一個特質這樣說的時候.
再強調多一點.
可能《馬太》的讀者.
猶太人多.
心裡面有很多的環境.
所以叫他們要回轉.
其實回轉就是悔改.
你還不回頭.
你還不去回頭是岸的時候.
就好像回到孩子那樣.
當然不是叫我們全部返老還童.
不是叫我們已經做了成人.
可以回到小朋友這樣的年紀.
沒有什麼可能的.
其實是在說那種心境或者信仰的態度.
要我們可能有一種很自我.
很自意的心.
特別猶太人以為自己認識了很多聖經.
特別是舊的聖經.
他們預測到彌措是哪一位等等.
要他們回轉.
要他們悔改.
這樣的時候.
好像孩子那麼單純.
剛才馬可夫這樣說的時候.
你才能夠接收到天國.
所以《馬太》說得很白.
然後再下面那句就說.
誰自我謙卑.

$^{121}$像這個孩子.
你才能夠進入天國.
將馬可的信息再引伸再解釋給我們聽.
所以當我們這樣去很概括地去想的時候.
這個小子.
這個小孩的角色.
剛才說了好像一個特質.
好像一個進天國的特質和條件.
馬可就說要單純倚靠.
好像孩子那樣.
馬太就說要回轉謙卑.
不可以以為是.
不可以自誇的轉變.
特別對一班當時猶太人.
或者有時候我們信了很久的人.
都有容易跌入馬太的對象的那種狀態.
所以就提醒我們要回轉悔改.
而馬太馬可除了說到孩子的特質之外.
他也有說到一種天國裡面.
這些小子的地位.
所以有一種接待天國.
他們天國裡面被看重的這班小孩.
這班小子.
好像馬可十四章.
讓小孩到我這裡來.
不要攔阻他們.
因為天國裡面只有這些人.
馬太再詳細又是.
馬太力度又加很多.
如果誰因我的命.
對孩子接待了我.
如果他.
令這些相信我的小人物.
去叛島.
倒不如將這個叛島的人.
好像你想像不到.
將大磨石放在頸上.
石沉大海.
大家聽我們中國傳統故事.
聽很多這樣的景象.

$^{161}$很嚴厲地對待他.
所以不要攔阻這些小人物.
進到上帝的國裡面.
不要叛島他們.
這是馬太和馬可.
除了說到孩子的特質之外.
也提醒讀者提醒我們.
神的國是屬於這些人的.
我們作為他旁邊的人.
他說到.
好像是你不要做錯事就幫到忙.
什麼意思呢?.
你不要叛島他們.
你不要攔阻他們.
讓他們進到上帝裡面就可以了.
很確固地將馬可和馬太的訊息.
告訴大家.
你記住這個馬可和馬太的訊息.
我們轉入路加.
今天我們讀的經文裡面.
路加可以說他講的主題.
其實是完全不同的.
螢幕上是新漢易本正經.
大家可以繼續看你手上教會的聖經.
翻譯上不會差很遠.
你見到紅色上都和大家突出了.
其實路加的記載.
是強調一個群體.
或者說是一個內部的關係.
一班門徒之間.
一班信了主.
在教會裡面一起生活的人.
那種之間的問題.
之間的矛盾.
所以你見到之間當中的事情.
其實在三節聖經裡面.
每一節都有講.
門徒之間在自言自語.
在吵鬧.
開始爭論誰做大.

$^{201}$誰做大佬.
一個這樣的問題.
然後耶穌就看見他們之間.
心中的意念是眾數.
每個人裡面都在想.
耶穌看穿了每個人的內心.
他看見每個門徒裡面.
都在爭論.
很想做大這個問題.
是否他自己呢.
就另一個孩子來.
讓這個孩子站在自己的身邊.
然後就交到門徒.
無論很熟悉.
像剛才馬可馬太都有類似的經文.
因我名接待這個孩子.
就是接待我.
但多了一句.
接待我的時候.
就是接待那差我來的那位父.
他再將這個孩子的重視程度.
再推進多一步.
其實是很寶貴的.
你接待了孩子.
直接就好像對現.
就是接待了上帝一樣.
這麼重要.
這麼矜貴的一件事.
然後最後一句.
因為在你們中間.
就是那句.
你們當中間裡面.
誰是為小的.
就是在上帝眼中為大.
路加的記載裡面.
那個要處理.
或者路加的記載.
要突出的問題.
是這班門徒.
開始在吵.

$^{241}$開始爭論.
當中誰是大.
大家可以細心想一想.
這種狀態.
就是一群人裡面.
他們在爭論誰是大的時候.
就是說有些人.
很想在別人之上.
這樣的看法.
很想有可能更大的體面.
更大的話事權.
更大的曝光率.
更大的權柄.
更大的權力.
去凌駕於人.
去壓倒於人.
去控制很多手上想要做的事.
或者要成立.
要成就自己.
很想達到的目的和事情.
很多種可能性.
他們想做大.
有很多心裡的意圖.
很想在這個門徒之間.
這班人裡面.
其實只有12個.
他們不是說整個教會.
他們說這個團隊.
這個事故的團隊.
一個小組裡面.
這班人.
他們之間.
想凌駕於其他人之上.
一種這樣的爭論.
一種這樣的意圖.
一種這樣的試探.
所以可以說.
聖經說這班門徒之間的問題.
其實是表示一個事奉的群體裡面.
都要面對一個.

$^{281}$人性上很真實的問題.
其實沒有一個聖人.
也沒有一班義人.
在我們地上出現.
通常一班人.
很有心一起去服侍也好.
走在一起.
例如我們崇拜裡面.
有不同的角色.
一起走在一起.
可能大家的心態.
大家的狀態都未必一樣.
很多人性的軟弱.
可能很累.
或者一早起床.
下大雨就不想回來.
很多這些很真實.
人性的卻步.
人性的軟弱.
是我們不需要避免的.
是很真實出現的.
反而我們真實去面對.
是更加重要.
更加需要.
所以盧家福音.
就算這班跟了耶穌三年的門徒.
去到吳忌位.
直接去趕這班門徒.
其實都有問題.
有人性的軟弱.
跟了耶穌多的時候.
可能見到耶穌行得多神蹟.
是吧.
對一多人說話.
然後他都很有地位.
很有力量.
好像很威武.
他們自己想到耶穌.
將來要做天國的主的時候.
他們都很想.

$^{321}$是吧.
其他的福音書就說.
我們是一個造友.
大家記得嗎.
都是一種想做大的.
一種試探.
一種人性很想得到滿足的地方.
是很真實的.
我想說.
就算一班門徒群體裡面.
都要面對人性.
這些很真實的挑戰.
而究竟我們要去思想的時候.
究竟作為一個事奉的群體.
我們這些高低的關係.
是吧.
都必然存在.
我們究竟怎樣去執行呢.
參考什麼模式呢.
怎樣去面對這些問題呢.
而又沒有離開到.
上帝要一班人一起事奉的情況.
福音書裡面都有其他記載.
就說.
好像馬可裡面.
應該是剛才第十章之後再說的.
是吧.
他說外邦人的社會裡面.
就真的要想一些領袖出來行使權柄.
去管治.
去用權力去壓倒其他人.
指示在你們中間不是這樣.
原來我們天國的群體.
不是參考外面世界.
或者社會.
或者甚至大家工作裡面.
那些上下之間.
那種行使權力的群體.
天國裡面只有一個參考.
其實回到主耶穌.

$^{361}$他對門徒的教導.
他對門徒的提醒.
其實那個參考點就在福音書裡面.
就在究竟耶穌要怎樣提醒門徒.
應該要怎樣由他們那種.
世俗化的問題.
回到一個很淑齡.
回到聖經裡面.
怎樣看一班人走在一起.
所以耶穌的時候必須介入.
如果不介入就大件事了.
是吧.
由得他們繼續吵呢.
分中發酵呢.
其實你再想像下去.
這個時候.
這個社會群體就必然出現紛爭.
是吧.
那十二個門徒.
可能就分成兩班人.
或者有兩個大佬出來.
然後就各自.
是吧.
有一個對質.
有一個針鋒相對.
一直鬥氣.
或者角力都不一定.
整個群體就會分裂.
就會出現很多的問題.
很多爭權.
內鬥的問題.
面對這個情況的時候.
耶穌就帶了一個孩子.
帶了一個小子在身邊.
然後就說了剛才我說的那番話.
就是耶穌說的那番四十八字的話.
無論誰因我的名接待這個孩子.
就是接待我.
無論誰接待我.
就是接待那差我來的那一位.

$^{401}$因為那在你們當中的最少.
其實不一定是最少.
當中為少的才是大的.
那為什麼又要帶一個孩子來呢.
這個小子出現在耶穌身邊.
和剛才馬可,馬太.
只是說像個孩子謙卑.
和要像個孩子才能進入天國.
有什麼分別呢.
在路加這個記載當中.
路加是不是好像一個實例.
帶一個孩子到他身邊.
然後就開始教導門徒.
當然了.
剛才說了路加都會參考馬可.
所以那個共同性都存在.
就是說這個小子在天國的尊貴身份.
其實路加都有說.
剛才我說無論誰接待我.
因為我接待孩子就是接待耶穌.
接待耶穌就是接待了差我來的那位父.
這個推到這麼遠.
是直接說更上一層.
接待天父這樣的表達.
其實是在表示在天國裡面.
這些小子應該得到重視.
這個可以說是馬太,馬可,路加.
三卷的福音書是一個共同信息.
不過路加就再推一層.
說那個重視程度去到天父.
天上一樣的情況.
或者地上怎樣對待他們.
在天上如同一樣.
在上帝國當中.
他都得到一樣的被尊重的對待.
但是我們要再想.
在路加裡面.
他有一個什麼獨特的信息呢.
這個就是我接下來要集中和大家分享的一個重點.
路加裡面有什麼獨特的信息呢.

$^{441}$這個小子在當時的社會地位裡面.
身份地位又如何呢.
耶穌要提他們.
天國他們這麼重視.
但在當時身處的社會又如何呢.
如果要凸顯天國是要被重視的時候.
都好像彷彿暗示.
其實他們在社會裡面不受重視.
所以天國裡面要額外去重視他們.
這些小子,小人物.
沒有特別的價值.
你記得馬可.
他講的時候.
那幫門徒是集被帶孩子來的人.
不要礙事.
為什麼帶一些不相干的人走過來.
是一種這樣排斥他們.
或者是拒絕他們的一種反應.
不是一個很大的貢獻.
也沒有什麼經濟價值.
所以當時這些小人物.
或者一些弱勢的人.
在當時猶太社會裡面是邊緣化.
不得到一些很重視的肯定.
也不會特別尊重去看顧他們.
有什麼額外的對待他們.
不受重視.
特別是成人的眼中.
他沒有很重要的地位.
可以說是微不足道.
不放在眼內的小人物.
有時我們思想的時候.
我們都不自覺.
將人和人之間.
會分等級.
經常都評分.
有什麼級別.
很多時候都有這些意識形態.
在我們身邊.
如果我們不清醒.

$^{481}$就會用那些思維.
用那些分級別的角度.
去看我們身邊的人.
看我們身邊的弟兄姊妹.
好像低級一點就沒有優惠待遇.
高級一點.
銀行都是這樣.
你坐飛機都會舒服一點.
類似這樣.
你覺得自己可能尊貴一點.
這種待遇.
這種按人在社會的地位價值.
這樣去評級.
其實是無孔不入的.
所以頂尖子母我們都要很清醒.
在天國裡面.
我們怎樣去看這些小人物呢.
怎樣去看這些小子呢.
去到耶穌裡面最核心的教導.
當他們是很重要.
你接待他們.
如同接待了主.
如同接待了主的時候.
接待了差距來的父.
是有解釋的耶穌.
就是整段經文最後一句.
因為那在你們當中最小的.
才是最大的.
有個因為.
好像給我們一個基礎.
我們為何要這樣做呢.
為何要這樣看重那些小人物呢.
路加在這裡.
按住門徒記住.
不要忘記背景.
他們之間爭奪著做大.
最終其實路加的記載.
是要回應門徒這個問題.
所以最終講了一段.
帶了孩子來的時候.

$^{521}$講出一個很重要.
在天國裡面.
怎樣分大小的標準.
怎樣跟外面的社會.
不同的一個標準.
我們這兩句大家去想想.
首先前半句.
在你們中間成為最小的.
Be in the list.
Among you.
在你們當中這班人.
記住是留心.
是說那班門徒.
依然都是跟住耶穌的那班人.
好像我們一樣.
我們這班人之間.
圍繞著我們這班人裡面.
要成為小的.
本身就未必是小的.
即是說我們不是小朋友.
但我們要甘願樂意成為小.
其實是一種這樣的表達.
那甚麼意思呢.
關乎當下這班人.
這十二個門徒.
我們聽到的頂子妹.
並不是說天國.
並不是說將來.
並不是說上帝在來臨的時候.
這班人要怎樣成就.
是說當下.
是說我們現在聽到的時候.
我們自己生命.
跟身邊的人那種關係的狀態.
要成為最小的時候.
其實有幾個意思.
一會兒我會再跟大家.
在應用上再深層去想想這個問題.
可以是為了一個群體.
而去完成.

$^{561}$去做一些可能是不顯眼.
而它是很重要的一些事.
譬如剛才經文裡面記載.
就是接待這些孩子.
個個都想接待主.
每人會將接待孩子.
放為一個很重要的事.
有時候我們.
有時候我記得在教會崇拜裡面.
現在到崇拜完了.
詩歌完畢了.
到孩子就要.
被導師帶去.
有些教會是這樣做的.
即刻一起敬拜.
到了聽到的那一部分.
他就將分開了.
小朋友就跟了一些導師出去.
是不是.
你可能會覺得.
怕他們會妨礙.
如果這樣就真的接待不了他們.
但不是的.
是令一些更加好的東西.
讓孩子聽得更加清楚.
這樣就真的很重視他們.
看看教會的做法是想怎樣.
所以接待孩子可能.
可能很多特別當時的社會來說.
是一些沒什麼用的東西.
豬頭骨的工作.
沒有人將射燈放在旁邊.
接待孩子和他玩.
這樣的工作.
個個都看著耶穌.
當時看看耶穌想說什麼.
旁邊有什麼人在.
跟他對話等等.
做一些不顯眼.
他是重要的東西.

$^{601}$今早大家回來崇拜.
清理了很多水跡.
或者很多東西.
讓我們都能夠進來.
都很順利.
都平安.
不會滑倒.
我們看不到的.
看到就在崇拜當中.
前後的事情.
大家都未必看到.
完了之後可能大家.
同工再打點一切.
我們留意一下.
其實有些不顯眼.
不重要的東西.
好像漢威.
其實是有人甘心樂意.
這樣去做了他.
這個就是所謂.
在《信經》裡面所說的.
成為小的一種表現.
然後再說.
當這樣的時候.
才是大.
才是最大.
是一種什麼模式呢.
這個剛才說了.
這個最大是在說.
在上帝眼中是最大.
不會說在人當中最大.
但是他講明.
聖經當中.
當我們成為小的時候.
是大.
這個是上帝判斷的.
在上帝眼中.
可能很多人看不到.
但是上帝自己在天上.
無所不在.

$^{641}$他每樣東西都能夠.
監察到的時候.
去看見我們每一個.
細微所做的一切.
所以我們能夠這樣去.
甘願去服侍的時候.
在上帝眼中是大.
不單只是在說謙卑.
馬太.
馬可的問題.
或者那個出路.
而是更加願意.
永遠去承擔一些.
沒有成效的職責.
我在這裡形容裡面.
想到兩個情況.
先說事.
再說人.
我怎樣能夠去實踐這件事呢.
哪些是一些.
不顯眼的東西.
是少的東西.
其實是需要去承擔的.
是甚麼來的呢.
沒有讚美.
剛才我說了.
沒有讚美.
沒有答謝的聲音.
不是站在台上去.
也不是要問公不公平.
為甚麼要我做呢.
為甚麼要我近呢.
是我們願意.
主動地去做了.
不用問這個問題.
甚至有時是.
有人不明白.
甚至不解釋你.
其實你都願意去承擔.
我分享一個經驗給大家聽.

$^{681}$有時很多時候都講道.
當然都有試探.
有機會有人說你講得很好之類的.
這些是傳道者心裡很清醒的一種操練.
但除了站在台上講道.
我記得我自己也做過一個.
類似這樣的工作.
是有一個群體發生了一些事.
是事奉人員.
要調查一些事情.
可能大家明白.
我不講得很白.
是甚麼事.
我就成為了其中一個成員.
去調查一些事.
你明白嗎.
平時可能大家握手.
都會分享訊息.
這些教導的角色.
但當我自己坐在那些會裡.
在那些關上門裡.
可能有三四個成員.
對著一兩個要問的人的時候.
其實就有很多人性.
對方就會有很多激動的語氣.
他不會尊重你.
但你知道為了整個群體.
你那個調查一定要進行.
而且因為當我有公信力.
即是他尊重這幾個成員.
坐在那個會裡.
是能夠判斷到那件事.
所以我們一定要承擔那個角色.
後來有一位中途受不了.
離開了.
但我們知道不可以完的.
如果你不完整地完成整件事.
那個群體就不會有好結果.
所以一定要很多人.
在等待這個調查.

$^{721}$出到甚麼判斷.
然後再繼續去作一些公佈.
所以是一些沒有人知道的.
關上門.
晚晚靜靜.
可能連續開幾個會.
沒有人會留意到.
但我自己的領袖就是.
如果我不去做.
亦都有人退出了.
為了群體緣故.
就沒有人能夠去承擔.
或者繼續去堅持下去.
我想起一首流行曲.
不知道大家有沒有聽.
有沒有聽過這首歌.
都幾好聽的.
大家可以聽聽現在現代的歌聲.
我不是聽很多.
但這首歌聽了.
就整個腦裡面loop.
大家自己聽歌的時候.
很有意思.
你可以搭訕到歌詞.
歌詞大家詳細慢慢回去聽.
我get到的意思就是.
有時我們常常想做.
只是在一些社會的規範之下做好人.
但其實最後一句.
要扮演壞人或者壞人.
是不敢.
但有時原來我們是需要.
我自己熟悉地get到這首歌的意思.
我們知道是這樣.
我們那種好和壞之間.
有時是一線之差.
有時很多社會的壓力.
很多群眾的壓力.
我們不喜歡或者不想.
將自己放在一個.

$^{761}$被人指罵的地方.
或者被人覺得你做得.
為什麼要這樣做.
為什麼要這麼堅持.
為什麼要這麼執著.
為什麼你要將不好的東西說出來.
這些角色.
但有時你發現.
是需要一些人去扮演角色.
所以說要扮演的時候.
好像我剛才說的經驗.
是頂不住的.
但你禱告.
在上帝裡面知道自己發生什麼事.
可能在那個情況.
你就要扮演一個壞人.
原來在上帝眼中.
他可能是會知道.
心裡面是發生什麼事.
這個就是和大家分享一個.
有些不顯眼.
雖然大家未必需要去到這樣的情況.
但有時是.
你也會問為什麼要我做呢.
為什麼要這麼辛苦呢.
就是那些情況.
經文拍網也給大家一些安慰和鼓勵.
原來大家所做的.
在天國是大的.
這個角度.
另外就是接待小子.
剛才說了.
在當時的社會裡面.
這些小子.
是沒有人重視.
沒有人去看重他們.
或者等到他們成人之後.
才將他們的意見.
將他們的想法.
放在眼內.

$^{801}$但耶穌卻是剛才說了.
在馬太.
馬赫路加.
往往都是帶一個孩子.
當時他們不會放在眼內的人物.
成為一個提醒這群門徒.
的一個教材.
或者一個教導的方向.
我們可以想.
對於我們來說.
我們身邊的這些小子是什麼人呢.
我們怎樣去應用.
這段經文呢.
可能是一些被忽略的人.
一些被遺忘的一群.
大家想想.
有沒有一些人.
我們有一群人很開心.
但往往被孤立了.
一些小組組員.
或者其他一些街坊.
或者其他的人.
同學.
我們有沒有額外的.
關注關懷他們呢.
有時是付出很多.
但你得不到回報的.
其實接待孩子就是這樣.
你全心全意接待他.
你不會想著.
我給你一杯茶.
你請我喝咖啡.
你不會有一個公平交易.
你不是帶著一種交易性的心態.
拿到好處的心態.
來去接待一個人.
何況接待一些小人物.
何況接待一些不顯眼.
或者是一些你覺得被忽略的人.
你是全心全意.

$^{841}$願意.
你決定了這樣去接待他們.
接待裡面最奧秘的.
最難的就是.
有時我們都會站在一個高低之分.
我接待你.
我就幫你.
我是服侍你.
然後我就做了一部份.
沒想過其實是.
接待的時候是雙方.
最神聖的時候.
你會從對方身上.
你得到好處.
得到好處是一種屬靈的好處.
得到上帝給你的提醒.
所以我最後就分享一個.
我自己的經驗.
就結束今天的講道.
我自己在自己的機構裡面.
都會常常服侍一些大學生.
這個是一個.
他們每次開學的時候.
就會在這個烏溪沙林.
這個營會.
他們都會找我.
去幫他們一起大大查經.
幫一班組長和職員.
預備新一年在大學團契裡面.
去服侍的大學生.
那天早上我就進了營地裡面.
和他們做了一個查經.
我都很熟悉這些事.
所以就拍了這張照片.
但是其實是很失敗.
想大家說.
那次的.
我的一個很大的衝擊.
所以基本上.
這個經驗常常都還在我心裡面.

$^{881}$我發覺沒有接待他們.
沒有和他們查.
就是查今天這段經文.
就星路加夫音這段經文.
我預備了我要講的東西給他們聽的時候.
我就即將結束了.
時間差不多了.
我就想著.
祈禱結束就結束了那個session.
突然之間.
有一位騎士在我身邊.
戴著口罩的弟兄.
就輕輕在我身邊說了一句話.
其實在你們中間.
剛才說我突出的那個字.
接著我聽完之後.
我沒有接納到他提這件事.
剛才說不要攔阻他們.
其實我就攔阻了他提出.
以致我就要自己完結這次的查經.
我就結束禱告.
就完結了.
好像一個教導.
好像一個總結就完結了.
其實表面上沒有什麼.
拍了照片.
大家都很開心.
聽完我說.
他們不會有什麼問題.
但後來我回家的時候.
就沒辦法了.
裡面要我自己回傳.
上帝也要.
其實是祂在提醒我.
這件經文裡面最重要的東西.
而我沒有在當中.
去深入了解.
而我為了要完結那個教導.
就停止了跟他們之間.
有一個生命交流.

$^{921}$我樹立了一個障礙.
在他們之間.
就沒有跟他們.
就算在某些時間上.
你都可以繼續下去.
但我那次就是一種這樣的心態.
完結了一位小子.
跟我的分享.
我沒有保持一個開放.
或者對他一種.
真的有的尊重.
剛才說的那種平等.
真正看待對方是一種平等.
你接待他的時候.
其實他也是一個.
在上帝眼中寶貴.
剛才說了.
接待他如同接待主.
如同接待天上的父王.
是不是.
不是要吃豬頭骨那麼簡單.
你是一個很尊貴的服侍.
當時我沒有帶這種心態.
後來很悔改.
當然我馬上那一晚.
所以我才有今天和大家分享的訊息.
我即時在他們的WhatsApp.
將一個重新.
再將經文演繹.
在那一晚就發給他們.
他們不知道發生什麼事.
但我自己給上帝一個.
很寶貴的提醒.
什麼是真正的接待呢.
當我們真正放下了自己.
無論對方是什麼樣的人.
你都不帶著一個條件.
你都不帶著一個框架.
你都不帶著一個標籤.
去看待對方的時候.

$^{961}$你全心全意.
做一個接待小子的服侍.
你在上帝眼中就捍衛最大.
好,我們一起在這裡禱告.
主耶穌多謝我們今天.
藉著我們一段很熟悉的經文.
但它是有一個很深層的意思.
你看到門徒中間.
出現了一些紛爭.
你都看重.
你都想幫助門徒.
去明白.
在天國裡面.
那種群體之間.
怎樣去彼此的接待.
彼此的承擔.
並不是按著我們在社會當中.
在外面的環境裡面.
學習了一套模式.
套入去上帝你的教會.
上帝你的國度當中.
多謝你.
藉著耶穌基督的教導.
不單止重新提醒這班門徒.
也提醒我們每一個.
怎樣在接待小子的服侍當中.
無論是面對這些小人物.
或者是做一些我們沒有人願意做.
不顯眼的事情都好.
我們因為你的緣故.
我們都會更多地學習謙卑悔改.
去承擔這些為著群體更好的緣故.
這些責任 這些功課.
求你繼續帶領我們.
同心祈禱 伴著我們.
好嗎?.
\newpage



\section{帖撒羅尼迦前書}
\label{sec:w32Ot_b7FcE}
\textbf{2024-06-13 周四晚查經班《帖撒羅尼迦前書(9)》 講員:曾裔貴牧師}
\newline
\newline
連結: \href{https://youtube.com/watch?v=w32Ot_b7FcE}{\texttt{https://youtube.com/watch?v=w32Ot\_b7FcE}} ~~~~ 語音日期: 2024-06-14
\newline
\newline
\hyperref[sec:AT0lX1Ebpa4]{\small{< < < PREV SERMON < < <}}
~
\hyperref[sec:index_chronic]{\small{[返順時目]}}
~
\hyperref[sec:index_scriptual]{\small{[返順卷目]}}
~
\hyperref[sec:1kzKFPo0unA]{\small{> > > NEXT SERMON > > >}}
\newline
\newline
$^{1}$好 我們一起重新普告.
感恩再一次我們星期四晚.
又可以有系統的來研讀.
帖撒羅尼迦的前後書.
還有關乎到末世論的一些神學.
幫助我們能夠透過「士徒保羅」.
對帖撒羅尼迦城的教會的牧養,關心.
來在書信裡面所流露的這份牧養的恩情.
讓我們能夠掌握,認識教會2000年前.
都已經是一個這麼富庶的地方.
來得到福音.
而且在馬其頓一帶和哥林多,希臘,雅典.
一直那個福音的廣傳.
都是「士徒」和一些當時的信徒的努力.
我們很盼望這些真理.
成為我們今天在香港的信徒.
來去依靠和盼望的.
特別是對永生都是在這本書不斷講主在內.
對信徒的盼望和影響.
願神幫助我們.
讓我們真是有一個平衡的時空.
我們既是活在實實在在的香港生活的.
每一個不同的處境.
但我們實際上是有永恆,國度的神的兒女.
主的幫助我們.
我們今天晚上再一次等候在主的面前.
禱告奉耶穌基督的名.
阿們.
第五章要再重新讀多一次.
相信第五章最後的經靈.
我們還沒有計算後書.
相當清楚是勉勵信徒要警醒.
等候主的再回來.
這個警醒裡面有一些要注意.
今天我特別想強調.
和大家去看一看保羅所看的預定信徒.
將來在主耶穌回來的時候.
我們會有甚麼得到的福氣呢?.
這個在保羅的書信裡面.
絕對不是一些例外的.

$^{41}$是差不多很多的書信都講到關於預定的.
第五章我們再重新 很短的.
是一個比較重要的勸勉.
弟兄們 輪到時候日期.
不用寫信給你們.
因為你們自己明明知道.
明明知道主的日子來到.
好像夜間的賊一樣.
所以不是計算哪一年哪一月哪一天.
是那個狀態.
就是如何準備.
你如何去準備.
就不是說哪一天 何時.
就不是那個叫主的日子.
所以一提到主再來的日子.
是信徒那個等候的那個來臨.
我們的心靈怎樣等候.
所以你在英國 你在內地.
你在極權的國家.
你在阿拉伯世界.
都是在講我們怎樣安排我們的人生.
來等候那個主的再來.
所以不是那個哪個地方.
不是說你要上去大帽山.
或者去橄欖山.
或者拆開講了.
有一段時間.
有一些的基督派.
就弄了很多的聖經.
當然很多年前就躲在橄欖山.
希望主再來.
那些猶太人會走到那裡.
就會讀到聖經.
就悔改.
但是那些聖經都已經荒廢了.
殘舊了 變亡了.
但是有這個心願.
很想猶太人得救.
這個好像不同的人.
都很想去為主再來.

$^{81}$做一些事.
但是保羅這裡就很清楚.
和耶穌所講的恩物相承.
不要計算那個日子.
不要在那裡預測.
有些什麼特別的方程式.
你知道哪一年哪一月哪一日.
他就會回來.
這樣計算.
我都說一定會變異端.
如果你又不聽人說.
那就更加變成大魔王了.
你自己死都沒有所謂.
你帶著一些信眾跟你就慘了.
你就令到人下地獄.
所以講到主的日子.
就絕對是那個信徒.
那個教會.
那幫人的內涵.
內心的那個狀態.
所以你看到明顯就是用幾個比喻.
就是有賊來.
如果你是一個家主.
你應該用什麼心態去保護你自己的家財呢?.
好像夜間的賊一樣.
人靜說平安穩妥的時候.
災禍就會忽然淋到他們.
如同慘難淋到懷胎的婦人一樣.
他們絕不能逃脫.
絕不能逃脫的背後.
正正就是沒有警醒.
正正就是沒有去準備.
完全不會在乎那個災難的人質.
還有不斷告訴自己 告訴其他人.
現在是平安穩妥.
就是不關我們這一代人的事.
不是我們這一代人需要擔心的.
我們要擔心的就是.
到底衣食住行哪方面有缺陷呢?.
各方面.

$^{121}$我們可以用任何的事物.
人事物取代了警醒主在內.
所以保羅在這裡接續耶穌在馬太宗24章.
所說的警醒.
第二個比喻就是慘難的婦人.
按道理是不可能的.
你十月懷胎.
你要準備東西的.
那時候是明姐出生.
我都準備了一大袋東西.
早上就入院.
因為似乎是順產的.
不知道的.
入了院之後.
當然拿了醫藥假.
回到家裡.
沒有事做.
百無聊賴.
醫院又不會讓你進去.
又去看不到.
就在家裡.
誰知道一睡就睡到五點鐘探病時間.
原來已經有很多事情發生了.
所以都不到你去準備了.
很多事情是突然間發生的.
按道理你是一個預產期.
你是會知道的.
連慘難.
雖然可能你不知道什麼時候會陣痛.
什麼時候會穿羊水.
諸如此類的.
姐妹可能會比較熟悉.
你就會有一個準備的.
但是這裡說.
去到這麼急的一個臨界點.
還是懵然不知.
主已經來審判.
所以你可以想像得到.
那種平安穩妥.
牢牢霸佔了信徒的心靈.

$^{161}$這是一個大問題.
教會如果沒有了這個警醒.
只是不斷發展.
不斷傳福音.
這完全不是錯.
但是我們的信徒.
原來從來都沒有想過.
會有主再來的末日.
真的有違反聖經.
很大路的啟示.
尤其是普羅書信.
很大部分都說到這一方面.
所以我們要留意.
第四節.
弟兄們.
你們卻不在黑暗裡.
叫那日子就是主的日子.
臨到你們像賊一樣.
其實本來就不應該是這樣的.
如果我們正正常常有真理的教導.
教會是一個警醒的教會.
教會是一個等候主耶穌回來的教會.
我們很少會有張椅子.
放在講台上.
這張椅子是主耶穌回來坐的.
大家要有這個意識.
回來後為什麼有張空椅子.
那是為了主再來的.
但我們很少會這樣.
這麼有趣的.
好像開玩笑的安排.
我們也不會.
甚至連那個心態.
來到禮堂.
我們覺察.
覺察就是心靈的.
Jesus是在這裡的.
我們也未必每個人都有這種心態.
所以我上星期跟大家說了.
你一定要有一個很強烈的意識.

$^{201}$你來到這個地方就是敬拜神.
你想想.
每個人都到齊.
三百人坐滿.
紅恩就一百三十至一百四十坐滿.
但是沒有人會想到.
神才是最重要的主角.
那個教會就很慘了.
因為最重要的是我們完全沒有意識.
所以我們不可以擔保每一個人.
但至少我們查經班的弟兄姊妹.
我們真的要每一次有強烈的意識.
主耶穌已經到了.
我們是來敬拜祂的.
這個就是我們不斷地經歷.
就會變成在生活裡.
我們慢慢就會有一種敏銳.
這個是要這樣學習的.
所以這裡告訴我們.
我們不是屬於黑暗的.
我們不在黑暗裡.
叫那個日子來到我們.
就好像賊一樣.
原來這裡有一個情況.
原來不是他不想避開那些賊的光顧.
或者是暗中的搶奪.
原來這裡又提了一樣東西.
是在說如果你在黑暗裡.
你就不會意識到有這種的防備.
所以這裡就告訴我們.
到底我們是屬於黑暗還是屬於光明呢?.
所以第五節就說了.
你們都是光明之子.
都是 每一個信徒.
不是你要洗禮才是光明之子.
現在你還沒有得住 你還沒有洗禮.
你先上初信班吧.
去浸禮班吧.
都不行 還要成長班.
不是 你們都是光明之子.

$^{241}$都是白晝之子.
所以你可以想像得到.
連續說了兩個狀態.
光明和白晝.
相對於剛才第四節的活在黑暗.
所以兩種狀態.
就是說一個如果有意識地.
知道自己在主救恩裡面的弟兄姊妹.
應該就一定是好像餃豬一樣.
將你放在主在來的餃豬的盤裡面.
你就已經屬於那個了.
會抽籤的.
我們其實就已經是在主的日子.
等候的一班人.
所以原來我們一信了主.
已經是屬光明的了.
不是因為你上了多少班主一學.
不是因為你上了多少卷查經班.
那個日子會令我們更加認識.
對真理的透徹 掌握和實踐.
如果我們沒有真理.
我們就變成甚麼都不知道.
只會按我們自己的喜好來生活.
甚至很容易就會人云亦云.
人家說甚麼你就信 你就聽.
因為你從來都沒有讀過聖經.
所以這裡告訴我們.
我們屬主的兒女.
本身一成為神的兒女.
就是光明之子.
就是白晝之子.
很強烈地 雙重地說了這個身份.
既是光明 但同時間又是活在白晝.
相對於第四節所說的黑暗.
生活的黑暗.
所以我們是不屬於黑夜的.
也不是屬幽暗.
黑夜幽暗 又重覆.
本來就是很強調白晝光明.
相對於沒有神生命的黑暗 昏暗等等.

$^{281}$似乎這裡要有少許竅門.
因為屬黑暗的可能都是基督徒.
如果堅持要活在那種自我滿足的心態裡.
他不知不覺又變成屬黑暗.
所以保羅的意思就是很想告訴我們.
我們是光明 是白晝之子.
但我們要有一個選擇.
我們的人生要揀選活在白晝.
這就是第六節所說的.
所以我們不要垂膠像別人一樣.
總要警醒緊守.
這是主耶穌對十二使徒的警告.
對後世門徒的勸勉.
藉著使徒們他們的傳道.
繼續去堅持這個事實.
教會是一個只能被動等待主回來的教會.
我們繼續接續前人.
好像傳承給我們的財富一樣.
我們接收了前人給我們的信仰.
我們要將這個信仰在我們有生之年廣傳.
我們去活出.
然後當我們過去了.
離開這個世界.
我們又將這個信仰原原本本.
沒有打折扣的.
沒有加添的.
交付下一代接行的人.
可能是青年的牧者.
可能是下一代的信徒領袖.
就要這樣一代一代交接.
這裡很重要告訴我們.
我們不可以睡覺.
好像別人.
別人是完全不行的.
似乎這裡不是說完全不行的人.
是說活在黑暗裡面的人.
所以保羅這裡說的別人.
沒理由突然間說教外那些.
貼手羅尼加城的一些扮偶像的人.
可能的 但是範圍似乎這一段的常夏文.

$^{321}$是說一些不知不覺沒有了盼望的人.
或者沒有了以主在來.
說得比較清楚一點.
沒有了以主在來.
為他人生的主軸的生活形態.
這樣就是別人.
所以這裡真的要求主讓我們能夠常常.
因為不是每個禮拜都要說主在來.
說多了你崇拜都不行.
但是你不說.
弟兄姊妹真的從來不會這樣想.
不會自己讀聖經.
就變成了一個挺大的掙扎.
所以保羅告訴我們.
總要警醒緊守.
似乎這裡總要警醒緊守.
不是說面對羅馬的逼迫王.
不是那些周圍現在在教會外面.
來抵毀使徒那些.
當然是阻礙他們信主.
但是這裡是說.
你自己怎樣去面對永恆.
在臨到的時候.
或者你要去永恆的時候.
那個心靈的付足.
那個才是警醒緊守的個人自己的裝備.
就不只是有逼迫.
那個現在實際在貼城.
貼壽羅尼加城的外面的官長.
那些猶太教的人.
這樣來抵擋信徒.
所以這個總要警醒緊守.
是每一個時代每一個信徒.
必須自己要有這個裝備.
真的每天都懷著感恩的心.
不是僥倖 是感恩的心.
蒙主的保守.
繼續我是開心活潑的.
去為神的福音盡一點綿力.
服侍神 求神顯明我有甚麼恩賜.

$^{361}$因為聖靈說明了.
哥林多前書的十二章說明了.
這位聖靈會隨己意的運行.
將各樣的恩賜分給人.
所以這個就是我們必須盡快.
順主或者去讀一些課程.
去發掘你的恩賜.
有些已經與生俱來.
有些是後天神賜給你的.
你去發掘之後.
一方面去發掘 發現.
然後去發展你的恩賜.
一定是那個場景.
就是在主的教會裡面去服侍.
當然主的教會可以再闊一些.
只是禮拜堂裡面.
可以是其他神給你的一些機遇.
或者你有些傳福音的恩賜.
可以去到不同的群體等等.
所以總要警醒謹守.
就是我們要很意識地.
我們是有閒的人.
是會離開這個世界的.
進入永恆的.
我們的警醒謹守.
就是以一種心靈.
我們準備好.
有這樣的福樂去到永恆.
但是還沒有去到永恆.
因為永恆沒有事要做的.
去到永恆之前.
其實有很多事要做.
我們要為神的國度去效力.
不要浪費我們的時間.
過多在地上物質的世界.
記住物質的世界是很實在的.
每分每秒在我們身邊.
在我們面前發生的.
沒有可能我們突然間.
像雲駕霧去到第三度空間.

$^{401}$是不會的.
我們實在是.
你想一下.
如果你晚上失眠.
三點起床.
逐秒到天亮.
看一看手錶還是四點.
看一看又是五點.
如果你有過這些經驗.
苦不堪言.
但是實實在在.
但是永恆就沒有時間了.
所以我們要把握.
我們一路年歲漸長.
經驗累積越來越多.
我們應該真的越來越看到.
能夠事奉神是最寶貴.
這些詩歌絕對不是騙人的.
是真的過來人.
尤其是我們越來越發現.
時光飛逝.
再不事奉其實就最可惜.
所以這裡總要警醒謹守.
不是一種很負面的.
一種捱打的狀態.
而是很正面的.
是很看穿世情.
很明白人生的.
那種真正純全的服侍.
不為名不為利.
真的這樣去服侍神.
所以一個剛剛初信二十多歲.
他又不明白這些事奉.
很多事情都是要慢慢神去雕琢.
如果我們都五六十歲.
還要雕琢呢.
其實從來沒有事奉過.
沒有失敗過.
所以我們不要怕失敗.
更加不要怕人對我們的批評.

$^{441}$我們盡力去做就行了.
最重要問心問神無愧.
這裡告訴我們第七節.
因為睡了的人.
就是在夜間睡.
醉了的人是在夜間醉.
有一個不斷講這個對比.
不斷講不警醒的後果.
或者甚至好像是一個循環.
既是沒有警醒.
自然就很容易會揀選一些.
隨手可得的.
及時可以行樂的.
來去享受.
你就會錯過了事奉.
有時可能需要的眼淚.
可能是禱告.
可能是擺相.
甚至可能是人對你的誤解.
但是你可以想像到.
保羅這裡的意思就是說.
睡了的人在夜間睡.
如果只講一句.
說了好像廢話一樣.
當然是在夜間睡.
睡了的人當然在夜間睡.
如果你看《上下文》.
其實是他的心靈的昏暗.
根本他活在一個.
實際是很健康的人.
譬如美國人.
特別是中產的白人.
吃維他命丸.
玩健身.
他們健身得不得了.
他們的人生根本就沒有神國度.
這樣他們可能是非常健康的.
但是不等於他有心靈的熒光.
相比一個屬神的兒女.
可能很多病痛.

$^{481}$人生都很坎坷.
但是他的內心有永生.
和他的盼望是很實在.
完全是一個有永恆.
為甚麼一個只是今生今世的人.
所以你可以想像到.
保羅這裡說的睡了的人在夜間睡.
醉了的人在夜間醉.
相對是八折.
就是我們在八咒之子.
我們應該怎樣呢?.
一定是在說我們心靈的護衛.
所以上一次上星期都說了.
信忘愛不是一種口頭禪.
保羅說的信忘愛是很機動性的.
是很立體的.
信和愛是要護衛心靈的.
讓我們的心靈對永恆有一個感動.
有一種感性的感覺,感知.
即是心靈沒有神的人就是沒有神.
但我們有神,有永生的人.
我們的心靈是不同的.
所以保羅這裡說的是多麼重要.
不是和我們研究古代的軍事裝備.
他只是隨手可得的就地取材.
將當時他知道的羅馬軍人的裝備.
就可以很快的將它轉化成.
對於我們的心靈的護衛的重要.
就是我們必須以信心,愛心.
來守衛自己的心靈.
對神,對人的信心.
很不容易的.
越認識人多.
你就越知道人有時去騙你一次.
你就慘不忍睹.
你就更加不敢去信人.
但對神,對人的信.
信之後就要有愛.
不只是一種信念.
是愛即是有一個具體的行動.

$^{521}$這些其實一點都不容易.
你要擺實際的.
因為這個講的信和愛.
一定是在講投射在身邊的人.
或者你不認識的人.
都要這樣來去實踐.
所以牧者就要面對很多這些事情.
今天我們有同工.
屯門,元朗,天水圍的宣道會的同工.
我們每兩個月就會開祈禱會.
今天剛好去到我們雄恩堂.
我又會做接待教會.
分享.
有一個同工吃飯.
就說有一個教友.
是男士來的.
不斷打給他.
問很多問題.
他都知道他有情緒病.
有時三更半夜.
很慘的.
你面對這樣的事情.
他說有時聽電話.
不知道第二天會不會在報紙上頭條.
誰誰誰.
有時你都好像沒辦法.
你都沒可能.
好像一個救世主.
任何時間去接聽這些電話.
明知他在這裡好像真的阻礙你的時間.
但是這個正正就是有時牧者要面對這些兩難.
但是無論如何.
這裡告訴我們.
信徒要避免活在一個黑暗和夜間的罪.
就是需要我們時常保守自己在白晝.
我們是屬乎白晝.
我們要保守我們的心靈.
和似乎那個頭就是在說那個思想.
就是對真理的認訊.
那個真理是能夠越來越融會貫通.

$^{561}$當然不是叫你現在阮明.
拜託你快點去學希臘文.
學希伯來文.
我想不是那些.
我們已經過了那個時代.
神都沒給我們這樣的裝備的需要.
每個人都不同的.
但是對基本真理的那個持守捍衛.
那些是必不可少.
如果不是.
我們如果沒有真理.
只是建立在一種感覺上面.
感覺很快就會失去的.
但是有真理護衛著.
我們就不會走去一個錯漠或者異端.
所以真的不要小看這裡.
信和愛是會令到我們心靈常常有一種敏銳.
哀莫大於心死.
中國人都很有智慧.
就是那個心一切的動力都是沿著那個心靈.
沒有了那個動力.
其實做什麼都是好像行屍走肉那樣.
那就多麻煩.
所以信和愛當作護心鏡遮空.
得救的盼望當作頭盔戴上.
就是很想告訴我們.
我們有這個對永生的盼望.
而且這個盼望是得救的.
就是說我們是得到救贖的.
我們是得以.
當然是開玩笑的.
真的有一張入場券.
就是我們在主耶穌的救恩裡面.
在耶穌裡面.
我們才可以進入天堂.
所以得救的盼望.
我們一生都要守住.
好像就是信者得救.
初信已經知道.
很多小朋友.

$^{601}$回幼稚園的主要學校都知道.
知道是.
問題是他20年之後大學畢業就不知道了.
理性將他完全的得救盼望.
完全蓋過.
好像一個大浪那樣.
有些人就突然間.
這個世界是不是真的有神.
又蓋過了那個大浪.
所以得救的盼望.
有時最後.
有時一些絕對沒有貶義.
有些比較基層的婦女.
她的信念更加強.
就是對神的那種單純.
那種堅持.
有時比起一些有學問的人.
不是學問不對.
絕對不是.
而是太多東西複雜了.
就沒有了得救的盼望.
就令到自己的思想混雜了.
很多人就是長大之後發現.
原來基督教我初頭信的.
信得這麼蠢.
現在我就不會再這樣信了.
但有人在的.
很多真的讀完書就以為自己識了很多很多的東西.
識東西是完全沒有問題的.
神創造知識.
整個宇宙是有結構.
是有秩序的.
證明裡面是有很多的理論.
那些科學上的定律等等.
一定是神賦予的.
但是我們如果崇拜了這些定律.
就忘記了這位造物主.
那他就沒有了得救的盼望.
所以那個心靈和那個思想.
是需要任何時刻的捍衛著.

$^{641}$所以保羅說.
屬伯奏的人就會心靈和思想.
要一致.
要能夠讓真理.
不要成為一個知識.
而是會影響我們的心靈的.
而是我們真的會活出有信有愛的.
這個就是我們的盼望了.
一章其實他都講了.
就是因信心所作的功夫.
因愛心所受的勞苦.
因盼望我們主耶穌基督.
所傳的忍耐等等.
但是這裡第五章.
再將信望愛又好像重新撿回來.
但是他完全這樣轉來轉去.
其實是在講一個很活的基督徒.
對永生的等候.
那個生命的改變.
所以是很不同的.
一個有永生盼望的基督徒.
他的生活一定是或多或少.
是越來越像耶穌.
一定要活在白晝.
你要記住這個我們要追求的.
所以第九節開始講預定了.
預定論了.
我都預備了很多段經文給大家看看.
都是保羅講的.
第九節.
因為神不是預定我們受刑.
是講那個得救盼望了.
將來我們不是受刑的.
那就是說將來是有些人會受刑的.
而且那個是神親自的刑罰.
永恆的審判.
所以不是預定我們受刑.
不是預定我們什麼呢?.
講了兩次預定.
第九節.

$^{681}$乃是預定我們藉著我們的主耶穌得救.
藉著我們的主耶穌.
這個預定.
仍然都是在主裡面.
享受那個永恆的天家福樂.
還有那個榮耀可以得到.
所以第十節.
就再一次講基督的得救了.
基督的救恩了.
祂是我們死.
叫我們無論醒著睡著都與祂同活.
很寶貴的涵蓋.
我們的生活的每一個層面.
就是所謂三衰六旺.
有時候你會突然間很精神.
很強壯.
很厲害.
但是或多或少有時候我們會病到死呀死呀.
我們會可能的際遇不好.
很upset.
甚至乎家庭有些變故.
人生一定會有一些高低起伏的.
所以無論醒著睡著都與祂同活.
這個是真的涵蓋了一個人生.
很寶貴的人生.
無論你是現在很風光.
但你不可能到永恆的.
如果你現在很困難.
但是你仍然都有那個盼望.
我們為一些時常都只能活在困難裡面的人去祈禱.
如果我們能夠時時保持自己有心靈.
身體的健康.
和各方面都穩穩陣陣.
我們是要帶著感恩.
但有更多人其實是活在困難裡面.
所以原來很寶貴.
祂替我們死.
有這位死而復活的耶穌.
我們與祂同在.
與祂同活.

$^{721}$所以叫我們.
即是我們已經屬於耶穌.
必須要認定這個這麼寶貴的皈依.
就是我們是屬於耶穌的.
所以叫我們無論醒著睡著都與祂同活.
即是說現在你在倫敦.
現在你在香港.
現在你在香港的35度熱到傻.
潮濕到不得了的那種悶熱.
或者你現在在一個很清新的地方.
很舒服的環境.
記住無論醒著睡著.
無論人生現在是好景.
或者是坎坷.
都與祂同活.
即是說耶穌知道我們現在的狀況.
這是很寶貴的一個牧養的恩典.
其他人未必知道.
但耶穌一定知道.
我們可以將我們的情況.
我們的焦慮.
我們的懼怕.
我們感到灰心.
都可以跟耶穌說.
因為祂說的.
無論醒著睡著都要與祂同活.
即是說同行人生路.
就是我們的主.
應許了我們一路走到上天堂.
所以不知何時.
人家會跟你說.
「啊 玉桂 你見到大光就跟著去」.
遲早都要的.
看誰跟你說.
但是現在有一個這麼精神的狀態.
有一個這麼強烈的意識.
知道自己是屬於耶穌.
所以我們就要記住.
無論醒著.
醒著即是說.

$^{761}$你現在可能是很能夠運用.
你自己的時間,金錢,你的自由.
各方面.
你做自己喜歡的事.
對的 醒著.
但是或多或少我們都會睡著.
我們人生一定會走下坡.
身體的狀況大不如前.
但是睡著的時候.
耶穌跟你同活.
即是說不是大包圍.
而是我們更加感受到這份很大的恩典.
即是他應許了照顧我們每一個的基督徒.
所以保羅也告訴我們.
等候主再來.
就要有這種的白晝觀.
有這種與主同活的人生.
無論去到什麼地方.
牧者都不會擔心信徒.
因為最重要的是耶穌跟你在一起.
或者說你自己有意識地.
耶穌跟你在一起.
這就是信徒力量的來源.
十一節.
彼此勸慰.
互相見諒.
正如你們素常所行.
就是要將剛才上文所說的.
要意識到主來的日子好像賊一樣.
不是要我們驚慌.
來去慌失失.
而是我們要知道我們是屬白晝的.
我們不是預定我們受刑的.
是預定我們在主耶穌裡面得救的.
所以這裡就是保羅對於信徒的勉勵.
所以提到預定.
既然第九十節出來兩次.
說到關於這個預定.
我和大家都揭曉一點.
關於預定這個課題.

$^{801}$可能先看《路加福音》這裡.
22章.
《古記》22章22節.
22章22節.
晚餐耶穌設立晚餐.
「時候到了,耶穌坐直,.
使徒也和他同坐,.
耶穌對他們說:.
「我恨願意在受害以先.
和你們吃這魚魚子的筵席..
我告訴你們,我不再吃這筵席,.
直到成就在神的國裡面.」.
耶穌接過杯來,.
祝謝了,說:.
「你們拿這個大家分著喝,.
我告訴你們,從今以後.
我不再喝這葡萄汁,.
直到神的國來到.」.
又拿起餅來,.
祝謝了,就張開.
遞給他們說:.
「這是我的身體為你們捨的,.
你們也應當如此行,.
為的是紀念我.」.
所以聖餐是紀念主十價的犧牲.
第20節.
「飯後也照樣拿起杯來說:.
這杯是用我血所立的新約,.
是為你們流出來的,.
看那賣我之人的手,.
與我同在桌子上.」.
因為猶大還未出去.
正在吃聖餐餅.
這就是出賣耶穌的猶大.
第22節就是預定.
「人子固然要照所預定的去世,.
但賣人子的人有禍了,.
貪活彼此,對問:.
是誰要做這件事?」.
留意這裡.

$^{841}$耶穌自己親自跟12個使徒說.
包括猶大.
原來他們出賣是神預定了.
被照著所預定的安排去世.
即是說猶大出賣耶穌.
都是在神的預知裡面.
不是神令猶大要這樣做.
是猶大自己這樣做.
但是原來主耶穌來這個世界的人生.
是被神預定要這樣犧牲.
應該說要再仔細一點.
是要被人出賣的這樣犧牲.
即是被自己最愛的門徒出賣.
置於死地.
你想想這樣的預定.
如果有個上士告訴你.
你就是這樣這樣.
你都想快點打他兩拳.
然後走出去,不要聽了.
通常我們會聽的.
說了你很慘,你就會很害怕.
所以聖經就禁止我們看相.
對自己的預測.
但是耶穌很早已經說了.
其實舊約教有很多經文.
預定耶穌要這樣被人出賣.
來去世.
所以這個就是預定的一個很重要的基礎.
我們的主是預定在這個世界.
是要以最痛苦的人生.
來去成長和擔待世人的罪.
這個就是我們要留意.
聖經是這樣理解耶穌.
不是一個偶然.
一個陰差陽錯.
突然間出現猶大.
不是,是神預定了.
耶穌要這樣來受苦.
我們開始看關於救恩的預定論.
有《少林傳》有好幾大段.

$^{881}$可能先看第二章.
第三章先.
《少林傳》第三章.
就是使徒的講道裡面.
不斷提這個預定.
有些人真的不夠堅決抵擋真理.
三章就是生下來就瘸腿的.
在美門門口的四十多歲男士.
生下來就已經不懂得走路.
因為彼得約翰經過.
見到這個男士乞淺.
突然間性命感動.
就覺得要治好他.
彼得和約翰牽他的手.
那個人還以為他會施捨.
可能攤大了手.
彼得和約翰說金銀我們都沒有.
但是我奉拿撒勒人耶穌的名.
叫你起來行走.
然後捉住他的手.
請他起來.
那個人突然間.
生下來就已經不懂得走路.
他的腳根本沒有力.
可能是扭曲了.
畸形.
突然間可以有肌肉.
有力.
四十多歲.
四十多歲從來沒有走過路.
又跳又開心地叫.
衝入聖殿裡面.
應該是會堂.
不是聖殿.
因為他不可以進聖殿.
在會堂裡面走來走去.
很開心.
所以個個都知道.
那個在美門口乞食的四十多歲.
是人抬來的.

$^{921}$所以是一個很大的神蹟.
當然也有人相信.
三章.
在第二十節之前.
看多一點.
可能是要關乎彼得說的道.
第十一節.
三章十一節.
《士徒行傳》.
那一人就是剛才說的四十多歲的僑隊.
那一人正在稱為所羅門郎.
所羅門郎.
耶穌在那裡說過道.
《約翰福音》第十章.
那篇經文.
那些人質問耶穌.
你叫我們猶豫不定到何時.
如果你是基督.
你就明明告訴我們吧.
就在所羅門郎冬天的時候.
拉著彼得入汗.
眾百姓一起跑到湯門那裡.
很覺稀奇.
彼得看見就對白勝算.
以色列人.
為甚麼把這事當作稀奇呢.
他覺得很自然.
耶穌的能力在我們現在兩兄弟身上.
我們叫這個人站起來.
因為神.
似乎給他真正有那個靈感.
有那個感動.
就真是他們立刻可以行神蹟.
是挺超自然的.
因為你有甚麼可能.
你不過是漁夫.
過了也是幾十天.
突然間變成一個.
可以行這麼大神蹟的使徒.
真是不簡單的.

$^{961}$但彼得就完全不是飄飄然.
很平常而已.
你怎會看這件事看得這麼誇張.
這麼稀奇.
為甚麼一定不徵汗我們.
以為我們憑自己的能力和血性.
能夠使得這個人走呢.
這個說話真是幾可圈可點.
彼得之前甚麼都要認厲害.
現在就是說.
你沒有看錯.
我們都是普通的門徒.
為甚麼你覺得我們兩個有甚麼神蹟呢.
以為我們的能力呢.
第十三節.
亞伯拉罕爾薩雅各的神.
因為這班人是猶太人.
就是我們祖宗的神.
已經榮耀了他的僕人耶穌.
即是他的兒子.
你們卻把他交付給彼拉多.
彼拉多是當時的巡撫.
仍然在位.
彼拉多定義要釋放他.
因為已經審問了.
覺得他沒有罪.
最主要彼拉多是想置身於事外.
即是不出聲沒有人說你啞巴的.
就叫他走.
走就不要再搞得很震.
現在暴亂怎麼辦呢.
但是走不了.
那些領袖說你不釋放.
你不置他於死地.
你不釋放巴拉巴.
你就是銅鑼合污.
彼拉多就說你們說的.
那我就將耶穌釘十字架.
所以彼拉多不是好人.
彼拉多定義要釋放他.

$^{1001}$但是你們竟在彼拉多面前.
戲絕了他.
所以是猶太人的領袖.
要置耶穌於死地.
彼拉多只不過是.
想自己置身於事外.
不會受到羅馬.
該省即是上級的審問.
即是你是一個地方的.
譬如元朗的政務專員.
李文虎.
現在村民有些東西反對.
你是那個政務官.
你又要擺平了.
如果你擺不平.
他就說我向特首告發你.
那你很煩的.
類似這樣.
彼拉多就是這樣的情況.
所以第十四節.
你們戲絕這位聖潔公義者.
反求著釋放一個兇手.
就是這個巴拉巴.
當時的強盜.
已經定了死罪.
給你們.
即是寧願放一個死囚.
也寧願放一個無辜的耶穌.
可以想像得到.
當時的領袖是多麼的活在黑暗裡.
然後告訴我們.
你們殺了這位生命的主.
神卻叫他從死裡復活了.
我們都是為這事作見證.
留意第十六節.
打後路.
開始說預定了.
我們因信他的名.
他的名便叫你們所看見.
所認識的這人見蹤了.

$^{1041}$正是他所賜給我們的信心.
所以這個信心就是彼得和約翰的信心.
記住 彼得現在在說神的思行原則.
不是一個方法.
不是你現在跟著一步一步.
你便可以繼續翻版.
是不可以的.
最大的重點是.
神在那一刻.
要透過彼得和約翰.
見到這個機會.
要令這個人的抱翻.
引發一個討論.
引發一個講道的契機.
彼得和約翰就著.
因為那些人已經湧向彼得和約翰求問.
還不是一個好的契機.
去講福音.
講耶穌.
榮耀耶穌的救恩.
所以其實是神製造這個神跡的機會.
給彼得可以即時的.
都沒有機會寫講章了.
對不起 今天我先回去寫個講章.
然後再來.
不需要了.
他已經現醒了.
就已經可以來到去即刻.
透過去榮耀耶穌的聖名.
救恩使你們見到這個情景.
所以那個情景只不過是一個.
穿針引線講福音的機會.
重點不是要大家.
大家現在快些追求神跡.
這樣就搞錯了.
靈恩最早期的蠻點就搞錯了.
追求神跡忘記了賜神跡的機會的那一位主.
那個救恩實的榮耀.
所以正是祂賜給我們的信心.
這個信心不是每次都可以的.

$^{1081}$這次要就行了.
所以我們要追求在主裡面有信心.
所以我們都要求神給我們有這個信心.
如果神要我們做一些.
超乎我們所做的事.
神會給我們有感動.
所以我們應該要追求與神同工的那份感動和信心.
就叫這個人在你們眾人的面前.
傳言的好料.
所以是很肯定的一個大神跡.
在你們的面前.
不用用布遮住.
然後就一二三的站在那塊布.
不是的.
在你們面前.
彼得和約翰就這樣拉著他的手.
就立刻站起來.
立刻又跳又跑.
這樣進入會堂裡面.
所以這個是一個引發講福音的契機.
所以盼望我們都能夠.
方聚會最重要就是.
要讓來聽的人產生一個對耶穌的感覺.
那個對耶穌的救恩的好奇.
甚至他已經預備好心靈要送.
這個就是我們需要大家一起去禱告去醞釀的熱誠.
然後就是說使這個人傳言好料.
彼得就在這裡說.
就是說他還他.
你們要面對你們自己的救恩.
弟們.
我曉得你們作者是出於無知.
出於不知.
你們的官長都是不知.
不知就是說不知道原來耶穌是真實的.
以為祂只是拿撒勒人耶穌.
木匠的兒子.
所以既然你們現在是這樣的時候.
現在就誕臣曾經直眾先知的口.
預言基督將要受害.

$^{1121}$就這樣應驗了.
是曾經預言的.
眾先知.
舊日的所有先知.
所以第十九節.
所以他還他.
你個個拍手掌.
彼得你真是厲害.
行神蹟.
可不可以教我.
那些人就是避開面對自己內心.
要悔改的問題.
所以通常你要陪談.
你要懂得巧妙地帶回去看.
你自己現在面對神.
你的心是怎樣呢.
那個見證好好.
陳友廣也好好.
嘴巴好動.
說完他回家就什麼都沒有了.
因為他根本不肯面對他自己.
他沒有將自己的心打開.
彼得就說他是好回.
弟兄們你們現在要面對.
現在我要說的是那個哭音.
不是說那個見證.
幾多秒內他的腳會慢慢復原.
不是說那些東西.
不說那些東西.
那些東西現在完全沒有用了.
引發了你的興趣之後.
引發了你的心靈有一個好奇之後.
現在就要踢你進去那個救恩裡面.
所以我們作陪談的.
都要很有那種敏銳.
心靈的敏銳.
和寶貴一個人.
如果那個心在那一刻.
他真的有一點猶豫不定.
我們幫他一把去決治.

$^{1161}$是很重要的.
所以現在彼得就說.
關你們事了.
不是那個人了已經.
因為第十九次.
所以你們當悔改歸正.
明明那個人神跡已好.
關我什麼事.
為什麼要搞我呢彼得.
我們敬拜神嘛.
這個是當時的新初禱告的時候.
在會堂.
當然他們是聖殿.
彼得就把握住這個機會.
既然現在說到關於耶穌.
藉著耶穌去醫好這個病者.
彼得就現在說.
你們沒有錯.
你們身體很健壯.
你們沒有瘸腿.
但是你們的心靈與神隔絕就是瘸腿.
你們的心靈根本就是黑暗.
你們就需要悔改歸正.
使你們的罪得而圖滿.
所以很冒犯啊彼得.
三句不合.
立即就說你們這幾千人需要福音.
你們現在是去聖殿.
守新初祈禱.
你說多麼的傲慢.
守新初祈禱的猶太人.
去敬拜神的猶太人.
彼得說你們要悔改歸正.
你們是出於無知.
不知道耶穌的榮耀身份.
所以你們置他們於死地.
但是現在要結帳.
你們現在要趕緊接受這位復活的耶穌.
所以是復活的耶穌.
使我們兩兄弟可以行這個神蹟.

$^{1201}$是神要我們在這一刻.
此時此刻行這個神蹟.
行這個神蹟原來.
最終我們終於明白了.
是為你們幾千人信耶穌來做的.
不是我們現在故意要演戲.
讓你們知道我們有多厲害.
所以彼得和約翰說完道之後.
回家什麼都沒有.
打回原形.
都是一個普通的門徒的使徒.
所以這裡就引發了一個.
很重要的製造了一個屬靈的干渴.
所以信徒或者未信的朋友.
如果有一個屬靈的干渴.
那個報道會就非常有功效了.
因為他已經看到自己與神的距離是隔絕.
如果不透過主耶穌.
他就永遠與神隔絕.
所以現在悔改歸正.
你們的罪就可以逃沒.
這樣永恆的安殊日子.
就是與神和好的日子.
就必定從主的面前來到.
主都必定會差遣.
所預定給你們的基督耶穌降臨.
主耶穌會降臨.
預定了的.
天必留他.
這個有些難解.
等到萬物復興的時候.
就是說原來主耶穌現在.
是在神安排的時間還沒有到.
所以祂不會來到這個世界.
我們只能憑著心靈去與主交往.
實際的面對面.
就是要在萬物復興的時候.
所以原來在信徒的角度.
這個是萬物的復興.
創世的時候墮落了的世界.

$^{1241}$萬物的復興就是主在來.
所以在未信的人.
在不信的人就是審判.
受刑的日子.
但是到萬物復興的時候.
這裡告訴我們.
就是神從創世以來.
藉著眾先知的口所說的.
然後有一節.
第22節.
摩西曾說.
主神要從你們弟兄中間.
給你們興起一位先知像我.
這是《申命記》十八章.
應該是十八節.
凡他向你們所說的.
你們都要聽從.
十五節.
原來我這本聖經是由經文傳出的.
凡是不聽從先知的.
我就要從民中傳滅絕.
這是十八章的《申命記》.
十九節.
第24節.
從撒母耳以來的眾先知.
撒母耳是一個比較指標性的先知.
似乎由他開始就有先知學校.
有很多先知門徒.
像使徒濟那樣.
開始先知的這個積佛.
所以撒母耳是一個頗劃時代的人.
因為那時候有濟師.
開始有君王.
但原來神興起一個濟師.
之後就是先知這個積佛.
所以從撒母耳之後.
陸續都會有先知的興起.
譬如我們想到的.
有拿丹,迦德.
這是大衛年間的先知.

$^{1281}$即是宮廷先知.
似乎拿丹先知.
都是和大衛一起成長的.
可能真是很熟的.
一起玩.
但拿丹就成為了神的僕人.
還有其他的.
譬如我們認識的就是.
以利亞,以利沙.
明顯以利沙就是跟隨以利亞.
很明顯的.
更加是先知門徒.
所以有那個傳承.
已經有這個訓練,裝備.
有先知門徒的學校.
所以告訴我們.
從撒母耳以來的眾先知.
凡是說預言的.
都說到這個日子.
即是原來一直都是在舊約裡面.
神已經很早就預言.
是會有基督耶穌預定來這個世界受刑.
祂會預定將來萬物復興.
主就會再回來.
所以這裡告訴我們.
25節其實對猶太人是一個.
你們一直都讀舊約.
小時候已經懂舊約了.
但居然對這個日子.
對這位救主完全陌生.
那你就慘了.
外邦人都可以了.
25節留意一下.
你們是先知的子孫.
承受了神和你們祖宗所納的約.
就是亞伯拉罕的約.
摩西的約.
大衛的約等等.
你們居然掉以輕心.
就是亞伯拉罕說的.

$^{1321}$地上萬族都要因你的後裔得福.
引創世記.
阿普羅.
所以對不起.
比特印.
第26節.
神既然興起他的僕人.
就先猜他到你們這裡來.
賜福給你們.
叫你們各人迴轉.
離開罪惡.
所以這就是一個很完整的訊息.
藉著這個瓊腿的神徵.
引發群眾的那個熟人的胃口.
或者引發群眾開始有一點光.
進入他的心靈裡面.
所以人始終都不會完全黑暗的.
總會透過神藉著不同的機會.
總會有一點光的.
這一點光我們怎樣來利用.
使他更加掌握而得救呢.
這個就是牧者傳道.
都有很重要的港道的責任.
所以比特的港道.
是有很多人信著的.
第四節.
四章四節.
但是聽道之人有許多信.
南丁的數目約到五千.
很奇怪的.
怎可能對著這麼多人的港道.
聽不聽得到呢.
等等呢.
在那個地方有多大的廣場可以這樣講呢.
臨時的.
沒有安排的.
你可以想像得到.
這些這麼臨時的.
都可以帶來這麼大的信主.
肯定是聖靈工作.

$^{1361}$所以不是說我們不可以有各方面的設計.
準備,宣傳.
不要變成一種太極端.
但是如果我們只是宣傳.
只是人為.
只是很多很多的佈置.
去到那個時候就要唱什麼歌.
彈奏什麼.
一切都不是錯.
但是沒有了聖靈最重要的工作.
在人心.
一切都是人為.
人為就即是說.
很容易就會掉落.
但是如果有聖靈工作.
那個人自己有這個呼求.
我們只是很純粹的推州.
推他一把進到救恩裡面.
就很不同了.
所以這個預定.
說的是基督耶穌被預定.
在末世來做這個世界受刑.
照著實的那個預定.
祂要來受那個刑罰.
但是當萬物復興的時候.
主就會再降臨.
這個就是預定論和我們的那個關係.
三章看完.
四章我們又揭一揭.
四章近尾了.
四章的28節這一段.
第23節開始.
那件事就帶來五千人信主.
但是彼得約翰就引發那個的哄動.
甚至乎好像一個官長.
這樣看為很難去調適的一個群眾運動.
所以就抓了彼得和約翰.
要他們不准再說關於耶穌的事.
23節二人既被釋放.
就到會有那裡去.

$^{1401}$法濟師長和長老所說的話.
都告訴他們警告他們不准再做.
第24節他們聽見了.
就同心合意的高聲向神說.
主啊你是造天地海和其中萬物的.
你曾藉著聖靈託你僕人.
我們祖宗大衛的口說.
很特別的.
這樣來講述這個詩篇.
這是詩篇第二篇.
彼得這樣理解詩篇.
是神藉著聖靈託付大衛在靈感裡面寫的.
即是詩篇的靈感不只是大衛.
而是聖靈感動他.
外邦為甚麼爭鬧.
萬民為甚麼謀算希望的事.
世上的君王一起起來.
神宰也聚集.
要敵黨主並主的受高者.
這是詩篇第二篇.
神立了他是那個受高者.
第27節引了詩篇之後.
講現實.
現在的現實.
就是彼得約翰被監禁的現實.
希律和本丟彼拉多.
外邦人和以色列人.
果然在者成侶耶路撒冷聚集.
要攻打你所告的聖母耶穌.
成就你手和你意旨.
留意下所預定必有的事.
又是那個預定.
又是主耶穌這位神所告的聖母.
她所要面對的.
在彼拉多和希律和的手.
受這些刑罰.
所以原來彼得說.
是預了的.
因為是實定了.
所以彼得就適逢在這個時候.

$^{1441}$跟隨耶穌.
又要面對耶穌的死.
他們又曾經背叛.
現在恢復了.
那個信心.
耶穌又升天了.
所以現在他們就將復活的耶穌.
不折不扣地.
傳給那個時代.
你想想那些人.
還是見過耶穌.
釘十字架.
見過耶穌死了.
那些人說祂復活.
但祭司一定是不相信耶穌復活了.
只會繼續用一個.
沒有說服力的說法.
那些人偷了那條屍體.
繼續只能夠用一種強硬去鎮壓彼得約翰.
他們不要再說了.
希望慢慢淡忘.
耶穌復活這件事.
所以他們只能夠這樣做.
但是居然他們沒有一個意念.
就是我們會不會有一點改變.
是不是耶穌真的復活了呢?.
是不是彼得所說的.
真的有一定的對呢?.
所以他們把心一橫.
無論多大的阻力.
都寧願阻礙彼得他們再說這個真理.
其實你想想是很懸殊的.
官方的壓力.
官方的資源.
和彼得約翰.
不過是沒有學問的小物.
只是漁夫.
一直都在加里尼的窮困的加伯隆的地方.
這樣去作業.
你怎能夠躋身耶路撒冷.

$^{1481}$這些宗教氣氛權力的中心呢?.
但是你看到他們兩個.
有一種裡面心靈的力量.
一方面清清楚楚自己是蒙召.
所以蒙召真的好像有臉色金牌一樣.
所以牧者如果我們不肯定自己的蒙召.
你事奉就要面對很多困難.
因為你自己都不信自己的身份.
你就很難有說服力說服其他人.
但如果你有很強的蒙召.
其實難有多少呢?.
耶穌也是這樣.
你就會面對那些困難.
你就不會覺得太難.
所以彼得約翰他們真的有很強烈的呼召.
所以我們怎可能不說呢?.
而且我們真的自己都改變了.
我們都是曾經灰心.
曾經否認耶穌復活的.
但是現在我們真的被改變.
所以他們就是要將這個這麼寶貴的信仰.
去傳給這些的官長.
所以他們是第一站見過耶穌死而復活的門徒.
他們就有這個很強烈的使命.
責任,託付.
要來將耶穌說明給這些官長知道.
總有些人無論多麼黑暗.
都會有一點光進去改變.
所以這個就是每個人的選擇.
但這裡告訴我們.
耶穌基督是被神預定.
祂要來接受這一切的安排.
就是祂的死是一個預定.
死在哪些人的手上又是一個預定.
祂要面對眾叛親離又是一個預定.
耶穌就是這樣來去成就.
所以想像得到主耶穌的救恩.
是一種有意識地進入死亡.
進入痛苦的一位救主.
所以這個救恩真的很偉大.

$^{1521}$我們又看看十三章.
還有幾段可以看.
十三章這一段.
又是後一點.
這次是保羅講道.
保羅在居比路.
即是塞浦路斯那個島.
講了整個島的.
很小的居比路.
因為是巴拿巴的故鄉.
所以第一站就去到居比路.
回到巴拿巴的故鄉.
向那些鄉親父老傳道.
十三章的四十八節.
我們看的是四十二節這一段.
他們出會堂的時候.
眾人請他們到下安息日.
再講這些說話給他們聽.
這個時候他們已經進入了土耳其.
因為已經在安提安.
比西底的安提安.
即是土耳其的一個中部的地方.
第四十二節.
他們出會堂的時候.
眾人請他們到下安息日.
再講這些說話給他們聽.
因為保羅,巴拿巴他們有一個想法.
就是我們要先傳給猶太人.
所以很自然就是去到哪個地方.
如果有猶太人的會堂.
他就會在安息日去會堂.
和那些猶太人辯論聖經真理.
然後如果他們拒絕.
才去那些當地的外邦人當中.
所以他們離開這個會堂之後.
第四十三節散會以後.
猶太人和敬虔進猶太教的人.
多有跟從保羅,巴拿巴.
二人對他們講道.
勸他們不要貧狗.

$^{1561}$在神的恩典當中.
然後去到下一個安息日.
四十四節.
到下安息日.
合城的人幾乎都來聚集.
要聽神的道.
但猶太人看見人者用當.
就滿心疾度.
顏薄保羅所說的話.
並且毀謗這是猶太人.
建立在他們舊有的猶太教的解釋下.
保羅自己以前也是.
所以保羅很諒解.
很明白他們的盲點.
現在猶太人.
其實這些不是耶路撒冷的猶太人.
是安提摩,比西底的猶太人.
更加不是很強的宗教中心裡面的猶太人.
但都是一知半解.
從小信仰到現在.
即是說其實不知道救恩.
更加不認識耶穌.
所以他們就硬駁保羅所說的話.
並且毀謗.
第四十六節.
保羅和巴拿巴放膽說.
留意這裡.
神的道先講及你們原是應當的.
只因你們棄絕這道.
決定自己不配得永生.
原來是這個意思.
不是負氣.
是保羅意識到.
這些人不是牛皮燈籠那麼簡單.
是他們根本不配得永生.
我們就轉向外邦人.
因為主曾經這樣說.
我已經立你作外邦人的光.
叫你施行救恩直到地極.
那段是《彌賽亞書》49章.

$^{1601}$講到彌賽亞會傳揚她的威望到外邦.
保羅直接引用.
是講主耶穌的福音.
傳給外邦人.
這個光.
我已經立你作外邦人的光.
叫你施行救恩直到地極.
外邦人聽見這話就歡喜了.
讚美神的道.
即是說如果沒有先日為主.
沒有成見的日信主.
就會得到生命的光.
得到改變.
會歡喜.
因為心靈會有一個很大的平安和喜樂.
然後就讚美神的道.
所以都是一個心靈的反照.
不是聽到誰見證就開心.
那個是一個開始.
是要使我們明白那個道.
所以一定要那個道.
是最實在在我們的心.
才能夠使我們有平安.
福樂和歡喜.
留意這裡48章.
凡是預定得永生的人都信.
於是主的道傳遍那一帶地方.
但猶太人挑唆乾淨尊貴的婦女.
和城內有名望的人.
迫得保羅巴拿巴將他們趕出境外.
二人對著眾人.
躲在腳上的塵土.
猶太人就知道是甚麼一回事.
就望耳歌念去.
門徒滿心喜樂.
又被聖靈充滿.
路加醫生就是描述.
信了主的人那個心靈狀態.
有些猶太人.
而且是在這個.

$^{1641}$比西米的安提柯會堂裡面的猶太人.
當然也有些入猶太教的外邦人.
所以原來不是猶太人.
與生俱來就可以得救.
猶太人.
沒錯是有救藥的神的.
應許的藥.
但原來主耶穌來了這個世界.
所有的人.
不論你是猶太人還是外邦人.
一定要在耶穌裡面.
你才可以得到與神建立這個生命關係.
一定要透過耶穌.
所以預定得永生的人.
一定是明白接受耶穌的人.
這個就是預定得永生.
所以你可以想像得到.
猶太人是多麼的堅決.
要否定保羅巴納所說的.
完全是基於聖經舊約.
這樣去勉勵他們.
來邀請他們去接受.
我們祖宗一直傳下來的.
但你們拒絕.
所以耶穌基督已經被設立成為外邦人的光.
叫他施行救恩一直到地極.
即是說很多很多不同國籍的人.
都可以順處.
不同的時代的外邦人.
所以如果你想想.
從耶路撒冷傳到中國.
都要2000年.
1900年開始.
就開始大量宣教士來中國.
19世紀初.
所以都要1900年.
去到地極.
即是來到中國.
不是因為飛機未發明.
是因為福音一直傳不到歐洲.

$^{1681}$傳到歐洲了.
保羅已經傳到歐洲了.
但是傳不到亞洲.
要有1900年才傳到亞洲.
因為開始進入天主教的時代.
封閉了.
真是去到15,17年的宗教改革.
真是福音要傳遍.
經過很多的醞釀.
都要200多年.
18世紀才是普世福音運動.
才發現意識到有幾個黑暗的大陸.
非洲,印度,中國沒有福音.
所以當時的歐洲才叫醒覺.
大量的去猜險宣教士.
去到這些地方去傳道.
所以我們都是承受了.
歐洲,美國的宣教士.
他們當時19世紀初期.
甚麼資源都沒有.
中國基本上是清末的時候.
我們可以想像到那種污濁邋遢.
未開化.
有知識的那些博博齋的.
很鄙視這些洋人.
去到內陸就更慘了.
因為迷信.
很多的風俗.
基本上對宣教士.
初期好奇.
後來你再講關於耶穌.
他們會逼迫宣教士.
所以很多教案.
慘死的宣教士.
都是這樣來到.
特別是義和團.
1900年.
有一個很嚴重的教案.
就是慈禧太后默許這些權匪.
其實那些玩十打的人.

$^{1721}$來到去.
即是說刀槍不入.
可以去殺宣教士.
到處去抓宣教士.
這個就是包括受牽連的.
中國的跟隨宣教士的人.
那些基督徒.
都要受迫迫.
所以在1900年.
一個比較大的教案.
宣教士被殺.
教堂被封.
所以這個就是.
預定得永生的人.
就會接受那個道.
原來有些人拒絕.
如果真的不知不覺.
不斷地拒絕.
原來從另一個角度.
他們是沒有被預定得永生.
所以這個都是很恐怖的事.
所以我們都不要說笑.
談論到神和福音.
我們要很認真.
活在神的真理和主的保守裡面.
就做這麼多時間.
還有一些時間.
我們可以看臨前.
這段都有幾可圈可點的一節經文.
臨前的第二章.
哥林夫前書第二章.
這些都是保羅的教導.
第二章的第七節.
第二章第一節開始.
上下文幫大家解釋一下.
大家稍微有些印象.
就會很重要.
因為如果只抽一節.
你不明白上下文就不是那麼好了.
弟兄們.

$^{1761}$從前我到你們那裡去.
地沒有用高言大字.
對你們宣傳神的奧秘.
當然不是說保羅亂說話.
沒有去準備.
不是這個意思.
要對那個時代背景有些認識.
他們說的高言大字.
就是那些修辭學.
當時很多那些.
東都港的遊行報道的演說家.
當時是很流行的.
到處講你的學說.
用你的修辭.
用你的口才的技巧.
有一套的方程式.
保羅說我不用這些.
我去講很簡單.
很淺白.
講明給你們知道那個福音.
因為我曾定了主義.
在你們中間.
知道別的.
只知道耶穌基督.
並他丁十字架.
相對當時那些人.
講一些虛無縹緲的哲學.
講了好像都不能夠兌現的東西.
一味就是修辭.
口才好.
保羅說我不是用那些東西.
我要講明.
講白主耶穌和丁十字架的故事.
這件事.
所以第三節.
我在你們那裡又軟弱又懼怕.
又甚震驚.
我說的話講的道.
不是用智慧委婉的言語.
所以那個委婉就是這個意思.

$^{1801}$好像令到你聽起來很興奮.
很開心.
但實際可能空洞的沒有內容.
乃是用聖靈和大能的明證.
所以相對又軟弱又懼怕.
保羅是靠自己的口才.
所以我們不要這樣看.
以為保羅是一個奇貌不揚.
不修蝙蝠.
完全沒有任何的科理的人.
其實不是.
他靠科理.
不靠當時人認為標準的那些口才學.
他靠的就是聖靈和大能的明證.
叫你們的信.
不是在乎人的智慧.
是在乎神的大能.
所以這裡就可圈可點.
你要講.
這個我好心想就是楊木公牧師.
你要令到一個小朋友都聽得懂你說什麼.
這個就是很高的智慧.
不是你要扮到變成你像個BB那樣.
而是你要令到那個受眾聽得懂.
所以就是說.
如果你向那些大學很有學問的人講.
你就應該要因應他們的那個程度.
可能講一些他們有共鳴的東西.
所以一些很出色的演說家.
有時才是引起回眾對你的那個共鳴.
你來引發他的興趣.
然後講論一些很深度的東西.
所以要分享一個想法.
一個訊息.
是最重要的中心.
但是引起他們想聽的那個意欲.
又是一個很高的技巧.
所以保羅告訴我們.
他有一個大原則.
他就說我知道現在坊間.

$^{1841}$現在的時代是流行什麼.
是崇尚什麼.
高舉什麼.
我就正正我有的.
但我不做這些東西.
我要讓你們知道.
福音才能改變你們的生命.
不是我的口才.
即是說易明哲保羅有口才.
易明哲保羅是有修養.
是有高的學問.
這個是他的背後.
但不是他用.
他用的就是令到你.
既然哥林多教會是一般人.
我就用你們能夠聽得懂的.
去將道講明.
使到聖靈在你們的心裡面.
產生大能的見證.
這樣就是這段經文的意思.
接著第六節.
然而在完全的人中.
我們也講智慧.
這個完全不是perfect.
是成熟的人的意思.
我們也講智慧.
但不是世上的智慧.
也不是世上有權有位.
將要敗亡之人的智慧.
我們講的乃是從前所隱藏.
神奧秘的智慧.
這裡好像分界了.
有實屬靈的智慧.
地上的智慧.
保羅其實要講解的就是.
這個十字架的奧秘.
主耶穌成為人類救主.
唯一救主的這個奧秘.
其實已經揭開了.
不神秘了.

$^{1881}$但保羅就說這個才是真正神的智慧.
是人會發明.
是想不到的.
但神為人預備的.
所以保羅現在既然領受了這個職份.
他要講解十字架的奧秘故事.
他就要用他自己的真材實學.
他的成長.
他由太教後來變成基督教.
他由敵黨迫北耶穌.
變成服侍耶穌.
這些心路歷程他自己都尋見了.
他現在就著哥林多教會.
你們那麼崇尚那些所謂智慧.
所謂哲學.
我就不用那些東西.
我就單純講耶穌的十字架救恩.
這個奧秘.
這世上有權有位的人.
沒有一個知道的.
如果知道就不會將榮耀的主釘在十字架上.
如經上所記.
神為愛他的人所預備.
是眼睛未曾看見.
耳朵未曾聽見.
人心都未曾想到.
又是引以暢亞書第52章.
講彌撒亞受苦那段.
第十節.
只有神藉著聖靈向我們顯明.
所以如果不是神向我們啟示.
如果不是神的聖靈感動我們.
其實沒有人信耶穌.
沒有人信到耶穌.
信到都是一種道理.
這裡告訴我們.
因為聖靈參透萬事.
就是神深奧的事都參透.
除了在人裡面的靈.
誰知道人的事呢?.

$^{1921}$像這樣做一個比喻.
做一個比較.
人自己內心.
甚麼冷暖自知.
又忘記了那個成語.
甚麼冬天喝水你又忘記了.
甚麼冷暖自知.
每個人自己內心有甚麼苦楚.
有甚麼其實你自己內心就知道.
大概意思就是這樣.
所以他的意思就是說.
像這樣除了神的靈.
沒有人能夠知道關於神的事.
我們所領受的並不是世上的靈.
乃是從神來的靈.
很肯定地說了.
我們信主的人.
是得到從神而來的聖靈住在我們的心.
所以我們知道神的事.
所以我們不怕面紅.
告訴別人有天堂的我們一定會進入.
因為這是神啟示給基督徒知道的.
我們是真的掌握著永恆.
永生的.
所以我們因為有聖靈.
所以叫我們知道神賜恩給我們的事.
接著第二章.
剛才我不是說了第二章七節這一段.
第二章告訴我們.
神在萬世以前預定使我們不榮耀.
保羅說的不是地上的智慧.
就是萬世以前所隱藏神奧秘的智慧.
剛才我沒有讀這一節.
就是神在萬世以前.
即是未有宇宙萬物之前.
這個萬世是這個意思.
即是未有任何物質世界之前.
那個開始已經是神的預定了.
預定了使我們得榮耀.
保羅這樣理解基督徒和十誡的故事.

$^{1961}$將人的生命的抬舉改變.
我們就是被神在萬世以前預定得榮耀.
保羅這樣理解.
所以對哥林多教會是一個挺諷刺.
因為上下文都是罵他們的.
他們整個心態就是追求那些哲學.
說得玄妙動聽.
但是空洞無物.
保羅說不是的.
我不要這些.
我講白給你知基督十誡犧牲的故事.
這個就是神萬世隱藏的奧秘.
而且是神在萬世以前.
就已經預定我們得到.
而且將來是得榮耀的.
這個就是在保羅底的這道關符裡預定.
還可不可以多講一兩段呢?.
因為都九點了.
而不所書就更加明顯了.
講了兩次.
不如我們下一次再補充一點.
我們看完第五章.
下個星期我們可能都可以準備一章的後書.
我們一起簡單的禱告.
今晚我們處理了關於生命的問題.
屬白晝和屬光明.
相對於活在黑暗.
在黑暗裡面的那個狀態.
願神幫助我們.
對於我們已經被預定得救.
因為我們活在主救恩裡面.
就肯定我們是預定神萬世之前的揀選.
以及我們將來要得榮耀.
以及萬物復興的時候.
我們就可以看到主耶穌的降臨.
我們就進到永恆裡面.
願主保守我們的心靈.
有信有愛.
我們的思想.
我們的意念是有盼望.

$^{2001}$這個是得救的盼望.
不是短暫的地上的.
是永恆的盼望.
幫助我們能夠在生活裡面.
無論醒著睡著.
都與主同伴.
願神幫助.
感恩的就是我們不同弟兄姊妹.
真的不是每個人都一樣的人生.
有些人真的是困難的.
而且那些困難都不是短暫.
有時是心靈的.
有時是人際的.
有時是要面對一些我們都很難去面對靠自己.
但是多謝主.
因為真的是一個很大的安慰.
醒著睡著是可以與主同伴.
願你保守.
藉著主耶穌是被預定受刑.
我們就充滿了感恩.
我們都是帶著一份很真誠的心.
去到台前領受聖餐的餅和杯.
因為是來之不易的.
是有血有肉的救贖主.
為我們這些罪人犧牲的.
願我們能夠充滿感恩.
真是服事是甘心樂意的.
是有意識地裝備好自己去服事的.
成為基督的精兵.
在不同的崗位的.
求主幫助.
奉主耶穌基督的名.
阿們.
\newpage



\section{民數記 14:24}
\label{sec:1kzKFPo0unA}
\textbf{2024-06-15 晚堂《專心跟從神》民數記 14\_24  講員:陳迦勒導師}
\newline
\newline
連結: \href{https://youtube.com/watch?v=1kzKFPo0unA}{\texttt{https://youtube.com/watch?v=1kzKFPo0unA}} ~~~~ 語音日期: 2024-06-18
\newline
\newline
\hyperref[sec:w32Ot_b7FcE]{\small{< < < PREV SERMON < < <}}
~
\hyperref[sec:index_chronic]{\small{[返順時目]}}
~
\hyperref[sec:index_scriptual]{\small{[返順卷目]}}
~
\hyperref[sec:pibg95moZiM]{\small{> > > NEXT SERMON > > >}}
\newline
\newline
民數記 14:24
\newline
\begin{longtable}{cl}
\hline
\hline
章節 & 經文 (和合本修訂版)\\
\hline
14:24 & \begin{tabularx}{0.7\textwidth}{X} 惟獨我的僕人迦勒，因他另有一個心志，專心跟從我，我要領他進入他所去過的那地；他的後裔必得那地為業。 \end{tabularx} \\ \\
[1ex]
\hline
\hline
\end{longtable}
$^{1}$各位姐妹 晚上好.
晚安.
今晚下了很大雨.
之前非常大雨.
我們為天體祈禱.
真的很感恩.
在我們聚會開始前半小時.
雨就停了.
風也停了.
我們在講道之前.
我們一起低頭禱告.
才藝天父上帝.
你此刻救恩給我們.
你減少我們當中.
每一個弟兄姊妹.
主啊 你亦帶領我們.
讓我們要作你的門徒.
主啊 在這個分裂的世代當中.
主啊 我們的心很容易被.
其他的事物所牽引.
主啊 求你去指教我們.
在這個幕後的世代.
主啊 求你以我們同在.
我們亦都懇求你.
求你的聖靈將你的話語.
讓我們能夠聽得清楚.
能夠明明白白.
亦都幫助我們能夠.
活出基督的樣式.
改變我們的生命.
我們願意將一切的榮耀都歸給你.
我們謹禱告.
是奉至耶穌基督的聖名祈求.
阿們.
弟兄姊妹.
我們在2020年的時候.
我們就經歷了幾年的新冠的疫情.
我們可以看到一波又一波的.
變種的病毒.
政府亦都有不同的限聚令.

$^{41}$我們就是這樣過了幾年.
但現在我們仍然都看到.
有很多弟兄姊妹都會戴口罩.
因為這個疫情雖然完結.
但亦都有流感.
和這個病毒亦都同時存在.
是夾雜.
讓我們亦都有很多的擔心.
亦都是我們現在香港.
現在正在經歷97年以後.
第二波的移民潮.
在過去這幾年.
亦都是香港最多人移民的時候.
以前97年是年輕人走.
現在是全家人一起走.
有些是幾家人一起走.
有些是教會弟兄姊妹和牧者一起走.
我們就在這樣一個環境當中.
當然要移民亦都不是那麼容易.
你要面對一些外面.
一些很遙遠的地方.
一些不同的生活環境.
要重新去找工作.
重新去找居住的地方.
所以亦都面臨著很大的挑戰.
大家都知道香港現在都是缺人才.
有一個高才通計劃.
早幾天報道.
有8萬到10萬人.
在一年當中申請高才通.
這些人全部是高學歷.
是名牌大學.
就來了香港這裡工作.
亦都是令到.
在以後的一些競爭當中.
是更加大的.
所以可以預想到.
未來我們香港在就業上面.
我們亦都會遇到很大的困難.
很多弟兄姊妹都在想.

$^{81}$究竟我們要離開香港.
還是留在香港呢.
CPA已經告訴我們.
在第46篇第5節怎樣說呢.
甚在其中.
成必不動搖.
道天一亮.
神必幫助者成.
聖經告訴我們.
有神同在的時候.
我們在哪裡都不怕.
只要甚不是我說的.
是聖經說的.
確實我們在這幾年當中.
我們很多人都是不想去回顧的.
真是不容易.
面對很多的艱難.
各種的環境正在急劇的.
在轉變的當中.
我們需要怎樣呢.
我們需要.
預備好我們自己.
我們需要調整我們自己.
我們又怎樣去調整我們自己呢.
我不是政治家.
我亦都不是經濟學家.
我不知道怎樣可以.
讓到大家.
討論一些思想.
立場和一些政治的方向.
不過我查考聖經.
我將聖經的人物嘉勒.
和大家一起分享.
到底我們身處在逆境當中.
聖經教導我們怎樣去面對這個困境呢.
我分享的主題就是.
專生庚從臣.
請看民數記十四章二十四節.
嘉勒可能大家都比較熟悉.
在以色列盟.

$^{121}$在七個埃及的歷史當中.
有四個最重要的人物.
有摩西.
有阿能.
有約斯亞和嘉勒.
嘉勒和約斯亞都是同一個時代的.
大家知道十二個炭治.
其中有兩個就是約斯亞和嘉勒.
在聖經當中.
在申命記的第一章三十六節.
在約斯亞記的十四章十一節.
都有類似的經文.
是講到嘉勒這個人物.
這個內容是差不多的.
亦都是臣對嘉勒這個僕人的一些評議.
唯獨我的僕人嘉勒.
因為他專生庚從.
有一個生字就是專生庚從臣.
三段的經文都同樣提到嘉勒.
和專生庚從這四個字.
所以可見嘉勒的獨有性.
亦都可以看到.
到底嘉勒又是怎樣的一個人物呢.
我們再回到聖經的一些背景.
在民數記的第十三章一到十六節.
就講到貓摩西打發十二個探子.
去窺探迦勒美地.
亦都提到.
第六節就講到.
這個迦勒是猶大支派的族長.
他是耶夫麗的兒子.
接著十七到二十節就講到.
摩西給他們一個窺探的任務.
簡單來說窺探迦勒的情況.
百姓的強弱人數.
還有他的成衣防守的情況.
還有當地的七產的情況.
好了他們到了二十一節和二十五節就講到.
這些探子就窺探了四十天.
然後他們就完成了任務.

$^{161}$就回來了.
到了二十六節到三十三節就講到.
他們回來之後就報信.
這個果然是流瀨耳蜜之地.
但是有十個探子卻報惡信.
令到那些民眾非常灰心.
我們再看回聖經三十一節.
但那些和他同氣的人.
卻是我們不能上氣攻打那百姓.
因為他們比我們強大.
三十二節.
於試探子中有人向以色列人.
散佈所窺探之地的謠言.
我們所走過的所窺探之地是吞嚇基民之地.
好像食人族一樣.
我們看到那些百姓身材高大.
他們看見巨人.
巨人就是那些雅納族人.
他們說我們看到巨人.
我們就像草蜢一樣.
而他們的眼中看到我們自己.
亦都是這樣.
很自卑.
雅納族人到底有多大呢.
我就想到大衛和哥利亞對戰的時候.
所看到的身形.
那個比例.
完全是很大的反差.
其實他們所遇到的高大身形.
他們的敵人強大.
其實都是必然的.
但我們想說的第一點是甚麼呢.
迦勒先生跟從神的聖訊.
以色列民陣入迦勒地遇到第一個問題.
就是我們的居民高大衛猛.
成一堅固的問題.
很困難.
於是他們就非常恐懼.
但迦勒沒有害怕.
他沒有退縮.

$^{201}$在文述第13章30節怎樣說呢.
迦勒在摩西面前安撫百姓說.
我們立刻上去得罪他們.
我們必能夠征服他們.
迦勒不是說我們想想辦法.
我們從長計議不是這樣.
迦勒安撫百姓我們立刻上去得罪他們.
不是想想.
立刻上去就是急不及待的.
我們必能征服那地.
其實是多大的聖身.
他知道那塊地是神賜的.
是英起之地.
他沒有想其他的.
他沒有想那些敵人多強大.
所以我們要學習的是.
生靠神福.
我們要有真正的聖身.
在面對困難的時候.
迦勒是專心跟從神.
他以聖身去面對神的帶領.
到底聖身又是怎樣呢.
在雅各司2章17節.
聖身若沒有行為是死的.
如果你的聖身沒有行為.
你只不過是一個知識的層面.
你的信仰和你的生命是沒有關係的.
沒有發生聯繫的.
所以你的聖身是死的.
你的聖身是需要有行動.
行動是甚麼呢.
就是神從神.
生靠和神從是有關係的.
你生你的頭.
你相信神.
你的外表一定會有活七個行為.
你對神,你見神.
你就一定會遵守神的懷疑.
這個就是真正的聖身的表現.
第二次我回想起.

$^{241}$我在2022年讀中華神學院的時候.
我需要辭去教會的前職工作.
而是前職去讀神學.
我的家庭就變成了生無人言.
為甚麼會生無人言呢.
第一是沒有積蓄.
第二是前職讀神學是沒有工作的.
第三是太太不是香港人.
沒有香港身份證是沒有工作的.
人家是帶著積蓄去讀神學.
我是甚麼都沒有的.
我計算過現在要讀四年的話.
生活費和學費起碼都要一百萬.
但錢在哪裡我不知道.
於是我就向神禱告.
我說神啊.
你確定你要讓我去讀神學嗎.
前職去讀.
然後神就透過馬太福音14章31節告訴我.
你這小省的人為甚麼要疑惑呢.
我真的不知道錢在哪裡.
疑惑得很正常.
你沒有錢.
但神提醒我不要看眼前的困難.
不要擔心經濟的問題.
神會必定供應你所需的一切.
我生不生.
生不生.
你生就要行.
你生靠就要盛從.
我就是憑著「生生」報道了中華神學.
在希伯來斯11章第一節說.
「生是所望之時的實體,未生之時的確具」.
前面的路是怎樣我不知道.
沒有人知道四年後會發生甚麼事情.
我相信能夠七醫神的話.
神一定會負責.
一定會負責到底.
所以當我決定讀神學的時候.
下一個月我就經歷了一個神蹟.

$^{281}$就是中華神學已經免學費兩年.
就是我要決定讀神學的時候.
神就是這樣讓我的學費解決了.
然後我看到神為我開路.
但生活費我不知道在哪裡.
生活費更大.
那時候剛好是教會財赤的時候.
知道移民教會的財政都差了.
這其實不是讀神學的好時機.
教會都缺錢.
每間教會都很多奉獻不足.
就在這個時候去讀神學.
但神就會感動弟兄姊妹去奉獻我的生活費.
我看到神是繼續為我開路.
我亦收到每個月的奉獻是不同的.
因為教會開了帳號.
弟兄姊妹就奉獻.
不是教會奉獻給你的.
所以不是每個月是固定的.
是這個月有多少給你一個數.
下個月多少給你一個數.
神是要我每個月都要仰望祂.
每個月都要倚靠祂.
兩年來我就是這樣來的.
不單是經濟上的困難.
身體上也有隱約.
我試過一個月去四次醫院.
我在醫院做什麼呢?.
靈修啊.
在醫院是我靈修的好時間.
是我讀經的好時間.
是我做功課的好時間.
我依然要這樣去努力去讀書.
神也讓我經歷了很多的恩典.
有聖僧是很重要的.
但我們所信的對象是更加重要的.
我們信的對象又是誰呢?.
弟兄姊妹你信的對象是誰呢?.
有一首聖詩.
我知道我所信的是誰.

$^{321}$大家有沒有聽過.
應該是.
大家應該會唱很多.
會唱幾首.
這麼少?應該不止.
為我心知所生的是誰.
並且也深信.
他必定能夠.
無收我所交付他的.
道前臂直到那一.
有沒有以前的氣聖靠的呢?.
我們也要留意.
我們生靠的不是自己的自生生.
那種生生是一定要從神而來的.
因為你的生生要有行為.
否則你的生生會滯死.
我們也要非常留意.
在教會當中我們不正確的生生.
在教會當中擴展.
我們在前福音的時候.
你要聖耶穌.
是不是我們在前福音?.
你要聖耶穌.
怎樣聖?.
總之你要聖.
你聖就行了.
你好好聖不要問太多.
再問多兩句.
聖耶穌有天堂的.
有天堂我就聖耶穌了.
是不是這樣?.
你是不是這樣?.
還有人家前福音.
前福音的時候我們要告訴人家.
其實聖是誰.
耶穌是誰?.
他有沒有真正的認識耶穌?.
我最擔心的是什麼?.
他只要天堂他不要耶穌.
他有一天來到天堂.

$^{361}$他見到耶穌我不要你.
我就是要天堂.
那種聖就是有問題的聖生.
那不是真正的聖生.
因為聖生是沒有行為就死的.
所以很多人他回到教會沒多久就走了.
他不再回.
他不是真正的聖生.
他不是真正的認識耶穌.
所以迦勒所聖的就是萬君之耶和華.
他知道神是統鼓萬有的.
雖然現在是很困難.
亞納族人是有很大的困難.
聖一是很堅固的.
但我們聖神是萬君之耶和華.
我們就不用怕.
祂們就是我們的食物.
有神同在我們不用怕.
我們不會經過任何的困難.
迦勒是以自己的行動去回應神的.
我們再看第二點.
迦勒天生跟從神的困難.
迦勒是一位對神充滿上心的人.
作為一個領袖.
他也會遇到生死有關的困難.
在文授記的十四章一到四節.
當探子說了這麼多負面消息的時候.
傳會中的百姓就開始肅殺.
開始對摩西和亞能賣弦.
要另外找一個領袖.
他們不再跟從摩西和亞能.
第五節就說到摩西和亞能要俯伏於地.
連摩西和亞能都受不了.
當時那些民眾的壓力.
我們再看第六節到第十節怎麼說.
窺探那地的人中年的兒子約斯亞和.
耶夫麗的兒子迦勒.
撕裂衣服對以色列前會重視.
我們所走過所窺探之地是極美之地.
第八節.

$^{401}$耶和華若喜愛我們就必領我們.
贈日那地把那流乃耳蜜之地賜給我們.
但你們不可背叛耶和華.
也不要怕那地的百姓.
因為他們是我們的植物.
第九節.
保護他們的已經離開他們.
耶和華卻與我們同在.
神一同在就不用怕.
當前會重正要拿石頭打死他們的時候.
耶和華的榮光在回眸當中.
向以色列眾人顯現.
我們看到有聖身的人.
也會遇到困難.
我們所看到迦勒,摩西和.
亞蘭他們要面對超過六十萬的回眾.
要打死他們這兩個唱反調的人.
弟兄姊妹.
有沒有想過.
面對死亡和威脅的時候.
那麼多人反對你.
你會不會再堅持自己的立場呢.
識時務者為政結.
看到風馬上就變戰了.
作為一個神的管家.
作為一個領袖真的不容易.
你處理的方法是少數服從多數.
還是你終生的去臣服.
去臣靠臣的呼召和臣的帶領呢.
這個就是迦勒.
戰生耕從臣.
至死都要忠心.
當我們面對困難的時候.
我們要懂得倚靠上帝.
在這裡我有張紙.
我想給大家看看.
大家有沒有看到這張紙有什麼.
有沒有.
你們近一點看.
有一座山.

$^{441}$對啦.
有一座山.
除了一座山還看到什麼.
空白.
沒錯.
這個黑點就是一座山.
迦勒山大.
對了.
這個白紙是什麼呢.
就是神的偉大.
這點就是你的困難.
這點就是神的偉大.
我們面對困難的時候.
往往有兩個導向.
第一個導向是什麼.
就是我們忘記困難.
你忘記了神的能力很大.
於是你的困難就越來越大.
這座山迦勒山大越來越大.
另外一種導向.
當你看到神很大的時候.
神的能力很大的時候.
中間這座山就只不過是一個點.
親愛的弟兄姊妹.
我們都會遇到困難.
我們在困難的時候.
我們有沒有依靠神勝過眼前的困難呢.
有沒有看到神的偉大呢.
但我們一眉間看這個環境.
很難.
然後就靠自己的能力.
靠自己的經驗.
靠自己的技巧去做.
你有沒有想過去靠神.
有沒有想過禱告.
有沒有想過神與你同在呢.
大兄姐妹.
我們再看第三點.
迦勒專心跟從神的結果.
回到我們的子題經文.

$^{481}$我們一起讀七好不好.
我們一起讀七.
一二三.
唯獨我的僕人迦勒因他另有一個心志.
專心跟從我.
我要令他贈入他所去過的那地.
他的後裔必得那地為宴.
神稱我的僕人迦勒.
迦勒的特點就是專心跟從.
他有一個很特別的心志.
他是與眾不同的.
他專心就是專一的生思而念.
集中他的心思.
他的目標就是很明確的.
就是向著這個標光直跑.
他是緊緊的跟從神.
他不是跟著跟著就掉了.
他的目標是非常明確的.
他始終而一的是跟從.
迦勒是神從神的呼召.
他專心的跟從神.
他也是蒙神所紀念的.
所以他才能得地為宴.
所以我們也要知道.
專心跟從神.
也是需要耐心氣等候的.
得地為宴.
我們知不知道.
約斯亞和迦勒等了多久.
原本他們氣到家冷地.
有些學者說是三天到十一天的時間.
結果因為會眾不相信神的英氣之地.
他們背叛了神.
然後他們就在曠野中漂流了多久.
是的 四十年了.
很厲害的.
這麼小的小朋友也知道.
凡是二十歲以上的人最後都死了.
發然而言.
下一代才能夠正式的進入迦勒地.

$^{521}$而上一代剩下誰和誰.
好像考大家一樣.
剩下約斯亞和迦勒.
用今天的話來說.
剩下兩個.
像李靖種那樣.
在曠野中兜兜轉轉.
兜轉了四十年.
多有耐性.
這四十年迦勒是不是什麼都沒有做.
是不是躺平.
什麼都沒有做.
算了 不做了.
不是的 迦勒是族長.
他依然帶領以色列盟.
兜就兜吧 也要照做.
神給他一些功課.
他也是一個資深的領袖.
他繼續的尊心.
他耕從神.
他到了八十五歲.
他還要去打亞納族人.
指動請英.
迦勒他是神.
給他的上身.
他就勇往向前.
我們很佩服教會有很多的牧者.
也都很強壯.
我們知道退休是六十歲的.
今次有很多牧者.
每人不夠牧者就六十五歲才退休.
也有些教會沒有牧者.
七十歲退休.
現在牧者荒.
我聽到有個教會.
有個牧者八十歲都七生.
我們要為牧者祈禱.
一定要身體好才行.
如果不是的話.
七八十歲港獨都沒有力.

$^{561}$求子氣幫助我們.
希望有更多的人起來.
去做牧者.
迦勒因為他尊心去耕從神.
他有這樣的生智.
他和他的後裔就得到.
那個希伯來山地.
在日西雅記的十四章.
十二到十四節.
我們到底又有什麼.
是我們的迦勒美地呢?.
有什麼又是我們的產業呢?.
在四月二十一日.
我們都參加了一個步行籌款.
在那個游博路的.
簡易公屋.
我們為那裡簡易公屋.
我們為那個生田屋.
步屋.
去籌集資金.
因為我們耿秀花已經附近.
一下子未來都會多了幾千戶的.
哇那些就是幾千戶的.
什麼?.
未得之地.
未得之盟.
我們會不會像迦勒那樣.
我們立刻上去.
要得著那地.
哎呀好累.
又要做多點事.
我們要立刻上去得著那地.
在詩篇一百二十七篇.
第三節.
康啊.
耶和華.
而你是耶和華所賜的產業.
所懷的胎.
是他所給的賞賜.
我們下一代.

$^{601}$同樣都是我們的產業.
在這裡有一本詩.
左邊那本詩.
是一個學者.
在1900年就做了一個研究.
他比較兩個家族.
就寫成一本詩.
詩名就叫做.
《尤克斯愛德華之》.
他追蹤了兩百年以來.
兩個家族的繁衍.
一個就是馬克尤克斯家族.
馬克尤克斯是一個無神能者.
他出生在1703年.
另一個家族就是愛德華之的家族.
愛德華之是一個虔誠的基督徒.
其他一個偉大的清教徒.
神學家.
他生於1703年.
無神能的馬克尤克斯.
對愛德華之曾經說過.
你信那個耶穌.
我永遠都不會相信的.
並且我的後代.
是永遠都不會相信的.
結果200年後.
那個作者對他的家族.
就作一個統計.
最左邊這個就是愛德華之的家族.
總人口是1394人.
其中13位是大學的校長.
65位大學的教授.
75位金官.
80位公務員.
60位作家.
60位醫生.
30位法官.
100位牧師.
100位牢書.
3位美國參議員.

$^{641}$1位美國副總統.
而馬克尤克斯的家族.
總人口903人.
其中310個是黑衣.
150位是罪犯.
7位是殺人犯.
100位是酒徒.
190位是巨女.
透過以上的一個案例.
我們看到.
相信神的家族.
信從神的家族.
和不相信神的家族.
是很大的不同.
很大的差別.
弟兄姊妹.
生人是需要虔誠的.
是需要代代相虔的.
神賜我們權柄是作神的兒女.
沒有作神的孫女孫子.
屬靈是沒有孫子的.
從愛德華茲的身上.
讓我們看到.
專心跟從神的人.
才能夠有敬虔的後代.
配有敬虔的後代.
不是教會的責任.
不是單單教會的責任.
而是每位基督徒.
做父母的責任.
特別在信仰上.
我們有沒有上上.
和我們的下一代.
一起讀經,祈禱.
去開展家庭崇拜呢?.
親愛的弟兄姊妹.
真正的聖身是需要有考驗.
經得起考驗.
亦是需要有行動的.
我不知道在這一刻.

$^{681}$你有沒有經歷困難.
無論是你身體上.
疾病的困難.
你經濟的困難.
又或者是你.
要不要移民的困難.
又或者是家人的困難.
你有沒有呢?.
可能你所有的方法.
你都嘗試了.
你都沒有辦法去解決.
這個問題.
弟兄姊妹.
在演繹當中.
你要專心跟從神.
僧言三章第六節.
在你所行一切事上.
都要認定祂.
祂必指引你的路.
和修館說.
祂必使你的道路平直.
而我們當中.
每位弟兄姊妹.
在這個黑暗背役的世代當中.
我們要孝乏加拿.
不看環境.
專心去依靠神.
以聖身去回應.
每一個困難.
每一天的困難.
都要眼前的.
交在神當中.
你要好好地去禱告.
專心跟從神.
成為眾人的榜樣.
我們一起去.
同受屬天的福分.
我們一起低頭禱告.
所有的天父.
讚美你.

$^{721}$你幫助我們.
讓我們能夠在大環境當中.
直至今天晚上.
你的聖息.
求你的話語.
藏在我們的心裡.
成為我們腳前的燈.
和路上的光.
照亮我們未來的方向.
亦賜給我們.
聖身有盼望.
有力量.
有能力.
勇氣.
去勝過這個世代的挑戰.
而我們當中.
每位弟兄姊妹.
都能夠專心地去跟從你.
不聽其他人的聲音.
而是聽你的聲音.
在眾人面前.
能夠活出美好的見證.
我們願將一切榮耀都歸給你.
我們今天禱告.
是奉子耶穌基督的名求.
阿們.
\newpage



\section{詩篇 23:1-6}
\label{sec:pibg95moZiM}
\textbf{2024-06-16 午堂《有爸爸的愛,你不孤單》詩篇 23\_1-6  講員:鄧泰康傳道}
\newline
\newline
連結: \href{https://youtube.com/watch?v=pibg95moZiM}{\texttt{https://youtube.com/watch?v=pibg95moZiM}} ~~~~ 語音日期: 2024-06-18
\newline
\newline
\hyperref[sec:1kzKFPo0unA]{\small{< < < PREV SERMON < < <}}
~
\hyperref[sec:index_chronic]{\small{[返順時目]}}
~
\hyperref[sec:index_scriptual]{\small{[返順卷目]}}
~
\hyperref[sec:Yzl3vxypdqk]{\small{> > > NEXT SERMON > > >}}
\newline
\newline
詩篇 23:1-6
\newline
\begin{longtable}{cl}
\hline
\hline
章節 & 經文 (和合本修訂版)\\
\hline
23:1 & \begin{tabularx}{0.7\textwidth}{X} 耶和華是我的牧者，我必不致缺乏。 \end{tabularx} \\ \\ \relax
23:2 & \begin{tabularx}{0.7\textwidth}{X} 他使我躺臥在青草地上，領我在可安歇的水邊。 \end{tabularx} \\ \\ \relax
23:3 & \begin{tabularx}{0.7\textwidth}{X} 他使我的靈魂甦醒，為自己的名引導我走義路。 \end{tabularx} \\ \\ \relax
23:4 & \begin{tabularx}{0.7\textwidth}{X} 我雖然行過死蔭的幽谷，也不怕遭害，因為你與我同在；你的杖、你的竿，都安慰我。 \end{tabularx} \\ \\ \relax
23:5 & \begin{tabularx}{0.7\textwidth}{X} 在我敵人面前，你為我擺設筵席；你用油膏了我的頭，使我的福杯滿溢。 \end{tabularx} \\ \\ \relax
23:6 & \begin{tabularx}{0.7\textwidth}{X} 我一生一世必有恩惠慈愛隨著我；我且要住在耶和華的殿中，直到永遠。 \end{tabularx} \\ \\
[1ex]
\hline
\hline
\end{longtable}
$^{1}$大型節目 早安.
我事奉的時候.
我最開心和最感恩的是.
能夠見到.
什麼叫做同感一靈.
我完全不知道斯班尼獻醜過的.
是的.
他們都不知道我說這篇.
所以上帝是很奇妙的.
我想說.
你們見到上帝會工作.
今天是父親節.
先說一次.
父親節快樂.
和你身邊的.
如果見到爸爸.
父親節快樂.
好.
各位父親問候完了.
想問.
你是不是真的快樂.
OK.
為什麼我會這樣說呢.
我想告訴大家.
或者很正式的.
和大家宣佈.
這個世界上.
最會說謊的物種.
就是爸爸.
為什麼呢.
因為你每一次.
我不知道大家有沒有留意.
當你問他的時候.
你有沒有事.
沒事.
需要嗎.
不需要.
痛不痛.
不痛.
辛苦嗎.

$^{41}$不辛苦.
為什麼我會這樣說呢.
這是典型的爸爸的例子.
我爸爸是一夫當關.
我家有七兄弟姐妹.
你可以想像.
要養七個法骨龍.
我爸爸是做沙廠的.
沙廠很熱的.
所以他每天的工作.
是要做十二個小時.
所以小時候.
我爸爸說了一句.
是真的不懂溝通.
一定不會的.
他只有一個任務.
就是賺錢.
回來養家.
大家要吃飯.
子女要讀書.
就是這麼簡單.
所以我爸爸小時候.
你問他什麼.
他就這樣回應.
所以我媽到今天.
經常都說.
爸爸那時候是怎樣的呢.
做到胃潰瘍.
做到吐血.
但反過來.
當媽媽問他.
痛不痛.
不痛.
你信不信.
所以其實你會發覺.
當我預備.
今天這邊到的時候.
現在多了一個功能.
大家都可以玩玩的.
最近很喜歡和AI聊天.

$^{81}$我就和AI聊聊天.
我就給了題目他.
我說 作為現代男性.
身為爸爸和丈夫.
可能面對的壓力是什麼.
之後不說就自可.
一說的時候.
就一堆出來.
我先和大家說一說.
這個人工智能和大家說.
它說什麼呢.
在身體層面方面.
工作壓力和長時間的工作.
導致身體的疲倦.
因應工作和家庭的需要.
而失去了個人的運動和健康時間.
生活作息不規律.
飲食習慣不好.
在心理層面.
需要負起家庭的經濟重擔.
維持家庭的生活.
在工作和家庭之間.
努力取得平衡.
時間分配困難.
面對孩子的教育.
和家庭管理的挑戰.
感到壓力.
照顧配偶和家人的需求.
自我照顧的時間不足.
靈性層面.
缺乏個人的靜修和探索生命意義的時間.
在道德和價值的選擇上.
和家人產生衝突.
無法滿足個人成長和自我實現的需求.
缺乏宗教和靈性生活上的寄託.
我就說給多一點資料給我.
我又給多一句.
你進一步告訴我是怎樣.
然後他又告訴我.
經濟壓力.

$^{121}$必須擔起家庭的主要經濟支柱的角色.
需要維持家庭生活足夠的費用.
包括住屋,教育和保險.
收入是否足夠.
是否能夠滿足家庭的需要.
工作和家庭的平衡.
工作繁忙常常需要加班和出差.
難以抽出足夠的時間陪伴家人.
如何在工作和家庭的責任之間取得平衡.
子女教育的壓力.
如何有效地去管理和引導子女的成長.
擔心無法充分了解孩子的需求和心理.
擔心教導方式是否恰當.
夫婦關係的維繫.
如何維繫夫婦之間的溝通和親密.
面對伴侶情緒的問題或矛盾時的處理.
協調家庭的決策和分工.
自我照顧的壓力.
缺乏個人的空間和獨處的時間.
難以保持運動鍛鍊和良好的作息.
情緒管理和壓力的釋放.
家庭角色的轉型.
傳統父權社會角色的轉變.
如何適應更多家庭照顧的責任.
自我認同和價值的重塑.
只是聽完都覺得崩潰.
所以你會發覺.
《團道書》難怪會說.
因為多有智慧就多有受範.
增加知識就增加憂傷.
你花了這麼多時間.
你的感覺就是崩潰.
所以再問多一句.
各位爸爸累不累?.
今早小朋友獻詩.
他獻的詩歌其中有一句.
是很觸動我的.
這句話說.
當爸爸有一天.
你的頭髮白了的時候.

$^{161}$我要反過來照顧你.
我就哭了出來.
為什麼呢?.
我爸爸是一個這樣的男人.
小時候不要說牽手.
不要說擁抱.
談話也很難.
所以我最印象的.
就是我爸爸在一次很生氣的時候.
我想他應該是很生氣.
打了我一巴掌之後.
我就流鼻血.
之後我媽媽就馬上拉著他.
她說你不可以這樣打他.
我爸爸還很生氣.
他是我的兒子.
為什麼我沒有權力打他?.
我媽媽就說你有權教他.
但是你沒有權力打他.
引起了一次小小的衝突.
但是我很記得.
當我爸爸年老了.
真的退休了.
我爸爸就很.
你可以說有一點不幸.
他在退休那年的時間中風.
所以後來的時間.
他的活動也是差.
但是我爸爸很堅強.
同樣道理.
你問他有沒有事.
沒有事.
之後他爸爸.
我爸爸真的很厲害.
所以後來他康復之後.
他霸佔了整個廚房.
他要霸佔著廚房煮飯.
不讓我媽媽進去.
之後就一直.
我很感恩.

$^{201}$照顧到自己.
我們家庭裡面.
又不需要太吃力.
但是到了我爸爸年紀大一點.
大概八十歲左右.
真的行動不方便.
大家都經歷過.
我爸爸開始走不了.
我跟他說.
用拐杖好不好.
大家知道了.
不用.
之後的時間.
肯用拐杖.
後來的時間.
你刺一下.
你也知道.
這樣的速度.
由家裡刺到診所.
要多久.
最初的時候.
我用輪椅推你好不好.
當然不好.
但是到了後來的時間.
也很感恩.
哄到他.
推輪椅.
我很記得.
特別有一次.
我記得.
我推他去.
除了推他去覆診.
我也推他去剪頭髮.
有一次我推他上去輕鐵道.
他對著窗口.
我在後面拍了一張照片.
在後面拍著他白頭髮.
在後面拍了一張照片.
這張照片我放在Facebook.
所以每一年都會走出來.

$^{241}$我記得那次的圖畫是怎樣的.
我看到頭髮.
我忍不住.
你記住.
以前我怎樣也不敢碰我爸爸.
然後我就在摸他的頭髮.
我在兩手摸他的頭髮.
然後我爸爸.
他就怎樣說呢.
喂.
你做什麼.
然後他就說.
你當我是小孩嗎.
我記得那次上帝給我智慧.
我真的彎腰在他耳邊跟他說.
我當然當你是小孩.
你現在在我手上.
然後我還記得我爸爸的笑容.
那種甜.
那種他在做很硬的男人的時候.
感受不到的甜.
我第一次可以碰爸爸.
我第一次可以抱爸爸.
在最後那幾年的時間.
都能夠可以去照顧他多一點.
幫助他多一點.
和他一起多一點.
這是上帝給我來引典.
這是上帝給予我很多的修補.
所以當爸爸累不累.
累.
這篇詩篇其實是大衛寫的.
有些學者估計應該是大衛未作.
大衛王的時候.
他還是一個牧羊人的時候的作品.
即是說他當時是一個大自然的裡面.
他很單純.
他很真實地去體驗和感受.
當時的氣氛.
所以也可以從這篇詩篇裡面.

$^{281}$你可以體會到.
其實大衛他一生裡面.
其實他和上帝的關係.
為什麼能夠令到他可以去跨越.
這麼多不容易的困難.
所以我今天特別將這篇詩篇.
送給每一個的爸爸.
坦白說我自己都是男人.
我們其中有一個.
我不知道是我們的優點還是盲點.
我們總是會覺得有很多東西.
我們不要連累人.
我們不要麻煩到人.
所以很多時候我們都是.
很多東西都是自己吞了就算了.
但如果你問我們作為一個男人.
我們自己有沒有需要.
我們有需要.
我們都需要一個地方.
我們可以去休息.
我們都需要一個地方.
我們可以傾心吐耳.
這篇詩篇他一開始的時候.
他說耶和華是我的牧者.
我必不自缺乏.
他使我躺在清草地上.
令我在可安潔的水邊.
大衛他很清楚.
上帝創造他.
上帝創造你是寶貴的.
上帝創造你不是要你受苦.
請大家記得.
今天雖然我們的生活裡面.
我們會覺得很辛苦.
但你要記得.
這個不是上帝的原意.
如果大家還記得.
這個是罪帶來的結果.
為什麼今天我們汗流浹背.
是當日當他們犯罪.

$^{321}$得罪上帝的時候.
上帝給男人的一個懲罰.
是罪的結果.
所以當今天如果你問我.
我們怎樣才可以不苦.
回到上帝面前.
在大自然的創造底下.
其實我們每一個人都很渺小.
我們的能力其實很有限.
來到上帝面前.
我們不需要再去裝作自己強大.
我們只需要來到上帝面前.
變回一個小朋友.
變回一個很微小.
被他創造的一個小孩就行了.
我們有天父,爸爸在.
這個是你的底氣.
所以如果當你認定.
他才是你的牧者.
在他的帶領之下的時候.
你是什麼都不會缺乏.
因為這段經文說.
他會帶領我們去到一個.
最適合需要的地方.
他會提供一切.
他會帶我們去到青草地.
去到水邊.
青草地和水邊.
是上帝才知道你的生命裡面.
最需要的是什麼.
有時候我們去到某些地方.
我們自己都不明白.
我們都不明白為何上帝.
會帶我們去到這個地方.
但你要放心.
上帝知道那個地方.
才是你最需要.
是你生命裡面最需要.
可能不只是物質.
可能還有其他東西.

$^{361}$所以他繼續說下去.
他使我的靈魂蘇醒.
為自己的名引導我走二路.
我雖行過死刑的憂骨.
但我受害因為你與我同在.
你的障你的肝都安慰我.
天父爸爸是你生命的引導者.
所以你不需要懼怕.
你不需要擔心.
靈魂蘇醒在另外一個版本.
新普及版裡面.
它譯為是重新得力.
上帝讓你重新得力.
當你累的時候.
你一定要去到神的面前.
你一定要定睛到神的面前.
你才能夠找回你的出路.
我記得早一段日子.
我自己都遭遇到一些很不開心的事情.
我自己都會覺得有些事情很遺憾.
然後的時間.
在這裡去想是沒有用的.
我就與勵士牧師請了幾天假.
我就在那個時間.
我見到兩個人.
一個是我兒子.
所以我兒子在台灣.
第二個的時候.
我很想見天父.
勵士牧師就很害怕.
他說想見天父.
我說不是我要安靜.
我就飛去台灣.
當然第一晚當然是找我兒子吃飯.
但我想告訴大家.
我與我兒子吃飯.
一句都沒有透露.
之後他吃完飯很安慰.
第二天的時候.
去了一個我夢寐以求的地方.

$^{401}$是我很多年都想去.
但我沒有機會去.
就是那個叫做.
沙溪還是清溪我不記得了.
十分左右的地方.
我很想去.
去到那裡之後.
其實我很感恩.
沒有什麼人.
很少遊客.
我見到每一個都很真誠.
每一個的當主都很真誠.
坐下來和他們聊天.
他們會很友善.
會和你講.
我又問問他們.
你們的生活如何.
很好.
好像朋友一樣.
他們問我很好奇.
你從哪裡來的.
我說香港來的.
香港啊.
現在很少香港遊客.
聊了一段時間.
然後另外一個環節.
我自己一個慢慢走.
走到那裡.
大家有看過電影節目.
知道很多鐵路.
舊的鐵路.
有很多舊的火車.
有很多一個很美的地方.
我在那裡坐了很久.
有些人去到那裡.
他覺得很悶.
我說我不說.
我在那裡坐了四五個小時.
我在那裡和上帝聊天.
當我去看回上帝.

$^{441}$你去創造這一切的時候.
我突然間會發覺.
我所面對的東西原來很小.
很少.
當我今天.
我自己覺得我自己面對的東西.
好像大到我自己沒辦法.
可以再去撐的時候.
原來我知道天父.
原來你什麼都知道.
原來在你整個創造當中的時候.
我們需要來到神的面前.
你會令我重新得力.
所以那次坦白說.
我臨走的時候.
其實我不捨得走.
然後我問天父一個問題.
可不可以讓我留在這裡.
當然上帝當然不會回答你.
當然不理你.
然後我就很黯然地離開.
很短的行程.
臨走就和兒子再吃一頓飯.
然後就回到香港.
因為你發覺大衛在這裡.
他肯定一件事.
他會引領你去走正路.
你跟著上帝是不會走錯路.
所以上帝是你腳前的燈.
是你腦上的光.
你是要跟著他.
其實一切所有解答的答案.
就是跟著他.
所以當我們今天面對著.
我們人生裡很多迷茫的時候.
其實你找上帝.
你問上帝.
有時候你不知道怎麼辦的時候.
可能答案很簡單.
只不過挑戰你.

$^{481}$你肯不肯跟上帝的意思去做.
有些隱友來到的時候.
你見不見到上帝不喜歡.
當你見到一些事情.
令你很猶豫的時候.
你有沒有見到原來上帝告訴你.
一早已經跟你說了.
你應該怎麼做.
其實他是知道的.
所以他就說.
就算你走到人生的谷底的時候.
你都不需要懼怕.
因為有他與你同在.
因為他會成為你的什麼.
安慰.
所以每一個做爸爸的.
其實你要記住.
你最大的動力是上帝.
是創造你的天賦.
我再強調多一次.
他創造你.
你是寶貴的.
他不是要你出世是辛苦的.
是受苦的.
其實這個也正好反映了.
我們每一個的爸爸.
我們自己作為人的爸爸.
雖然我們這麼軟弱.
雖然我們這麼差.
但說真的.
你自問吧.
你生了個兒子出來.
你是想他受苦嗎.
我們這麼多事情不跟他說.
我們這麼多事情吞下去的時候.
就是為了不想讓他受苦.
這個也是我們每一個人.
共同的原因.
所以今天有些人很有趣.
我兒子跟媽媽的關係比較好.

$^{521}$以前很多事情都會跟媽媽說.
但他不會跟我說.
通常我兒子找我.
一定是一些大事.
所以我接聽電話.
我們都會害怕.
當然現在不會有什麼.
通常會談什麼呢.
最近他要不要升職.
公司想升他.
但他很猶豫升不升.
然後他就找我談.
但我就說我很感恩.
跟他是一個好像朋友一樣.
可以跟他維持這種關係.
他找我談話.
但我就是想說.
其實我們今天雖然.
特別是今天做子女那些.
你一定要很明白你的爸爸.
如果他能夠跟你談得來.
其實你要Hallelujah.
但如果反過來有很多事情.
他不知道怎樣跟你表達的時候.
請你明白其實有很多事情.
他很疼你.
他不想帶你給你一些困擾.
所以他自己背了.
第三件事.
有時更加重要.
在我敵人面前.
你為我擺設筵席.
你用油膏撩我的頭.
使我的福杯滿人.
有時候好像我剛才所說.
我們面對一些敵意的攻擊.
和含冤受屈的時候.
其實他一切都知道.
而這裡的經文是說.
他一定會作出回應.

$^{561}$他會讓對方知道上帝的存在.
他要知道.
讓你受傷的人.
他要知道上帝的存在.
有時候反而危機.
是可以令到上帝.
得到榮耀的時候.
大家還記不記得.
當藥櫃被搶走的時候.
記不記得大菌廟發生什麼事.
那個大菌掉了幾次.
後來掉到連頭都掉了出來.
之後那些人就很害怕.
就知道那個才是真正的神.
急急忙忙.
賠罪賠禮.
送走了這個藥櫃.
大家還記不記得但爾尼.
當他受委屈.
被人迫害的時候.
被扔到獅子坑.
記不記得大理烏的反應.
大理烏的反應是.
你的上帝必救你.
必會拯救你.
所以當但爾尼沒事的時候.
大理烏是看見上帝的榮耀.
有時候我們未必會明白.
為什麼上帝讓我們在那個位置.
去面對事情.
但其實很多時候.
在上帝當中.
可以彰顯到上帝的大能.
和上帝他自己的榮耀.
所以當你在工作上.
在你的生意裡面.
在你的生活當中.
你遇到一些事情不明不白.
不知道為什麼會這樣的時候.
其實不用怕.

$^{601}$上帝在這裡.
上帝會為你出頭.
上帝會讓對方知道.
所以更加重要的是.
他說你會讓你看見神的榮耀.
讓你也可以得到榮耀.
因為祂會祝福你.
讓你福輝滿人.
所以當你看見.
但爾尼沒事的時候.
他就被大理烏王去重用.
他能夠成為上帝.
他自己一個很重要的見證.
所以大衛他自己去肯定一件事.
我只需要一生一世.
必有恩惠慈愛隨著我.
因為我永永遠遠.
都要在神的殿中.
在祂的恩典當中.
因為祂的恩典和慈愛.
會永遠伴隨著我.
所以其實詩篇23篇.
我們不需要很用力去解.
因為大衛.
他剛才我提過.
他是在一個大草場上.
他是用他真實的感受去經歷.
上帝在草場當中.
怎樣去照顧他和他的羊.
怎樣去帶領他.
怎樣去照顧他.
怎樣去幫助他.
怎樣去保護他.
其實上帝在這裡.
上帝從來沒有變過.
有時候反過來我們作為男人.
背得太多了.
所以我也想說.
我們作為爸爸.
我剛才也說.

$^{641}$我們最心底最深的那裡.
其實我們很愛我們的家庭.
但我們不懂得表達.
但我想告訴大家.
唯有當你的生命裡面.
你是一個很有健康.
很有活力.
充滿陽光的生命的時候.
你的家庭才能夠真正得到祝福.
所以今天這個主題.
有兩個意思.
一個就是跟爸爸說.
跟各位父親說.
有天父的愛.
你放心.
你不孤單.
第二個是跟每一個家人說.
有爸爸的愛在.
你從來都不會孤單.
因為每一個人的表達都不同.
每一個人用的方法都不同.
就像我爸爸一樣.
他用了一生一世默默付出.
當然我們還很小的時候.
什麼都不明白的時候.
我們經常都會投訴,埋怨.
爸爸都不喜歡跟人聊天.
都不疼我.
經常罵人.
到今天你自己再看的時候.
你知道你爸爸有多偉大.
他為你們的緣故.
他吃了多少苦.
但是也有些爸爸.
如果他能夠在神的裡面.
他能夠更有力量的時候.
其實你將會祝福到整個家庭.
讓他們感受到安全感.
讓他們感受到被保護.
讓他們感受到你的隨時同在.

$^{681}$所以求神幫助我們.
讓到這個父親節.
大家都能夠拿一份最好的禮物回家.
雖然教會也有禮物送給大家.
但那份不是最好的禮物.
最好的禮物就是.
你在神的面前拿到這個肯定.
去能夠很有力量,很有盼望.
去面對未來的生活.
求神幫助我們.
我們一起同心祈禱.
新加坡天父求主.
你親自帶領我們.
我們知道作為現代人.
裡面有很多父親.
特別是我們的知識去增長.
我們知道其實真的一點都不容易.
我們所面對的事情.
眾多煩惱.
你就告訴我們.
我們在你的牧羊之下.
我們在你的人靈之下.
所以當我們回到最基本.
最簡單的一個環境的時候.
我們知道我們一樣都不缺.
主人因為你每天都看著我們.
你最清楚我們.
我們所有的每一樣.
你都最清楚.
主人想我們回到神的面前.
我們求你讓到每一個父親.
能夠他們都有安靜.
可以來到神的面前.
去安殊找回你的時間.
讓他們的心能夠尋見你.
能夠找到你的愛.
找到你的安慰.
找到你的明白.
也都求主你因為這樣的緣故.
祝福每一個家庭.

$^{721}$因為每一個家庭都急需有一個.
是滿有能力.
滿有健康.
滿有陽光的父親.
在這個家庭當中.
我就求主你親自去報仇.
因為這樣的緣故.
讓每一個家庭都能夠得到祝福.
基督教主耶穌.
得勝明智其及.
好.
\newpage



\section{路加福音 15:11-32}
\label{sec:Yzl3vxypdqk}
\textbf{2024-06-16 早堂《父親的心》路加福音 15\_11-32  講員:許亞光牧師}
\newline
\newline
連結: \href{https://youtube.com/watch?v=Yzl3vxypdqk}{\texttt{https://youtube.com/watch?v=Yzl3vxypdqk}} ~~~~ 語音日期: 2024-06-19
\newline
\newline
\hyperref[sec:pibg95moZiM]{\small{< < < PREV SERMON < < <}}
~
\hyperref[sec:index_chronic]{\small{[返順時目]}}
~
\hyperref[sec:index_scriptual]{\small{[返順卷目]}}
~
\hyperref[sec:aAKtjwj8FbQ]{\small{> > > NEXT SERMON > > >}}
\newline
\newline
路加福音 15:11-32
\newline
\begin{longtable}{cl}
\hline
\hline
章節 & 經文 (和合本修訂版)\\
\hline
15:11 & \begin{tabularx}{0.7\textwidth}{X} 耶穌又說：「一個人有兩個兒子。 \end{tabularx} \\ \\ \relax
15:12 & \begin{tabularx}{0.7\textwidth}{X} 小兒子對父親說：『父親，請你把我應得的家業分給我。』他父親就把財產分給他們。 \end{tabularx} \\ \\ \relax
15:13 & \begin{tabularx}{0.7\textwidth}{X} 過了不多幾天，小兒子把他一切所有的都收拾起來，往遠方去了。在那裡，他任意放蕩，浪費錢財。 \end{tabularx} \\ \\ \relax
15:14 & \begin{tabularx}{0.7\textwidth}{X} 他耗盡了一切所有的，又恰逢那地方有大饑荒，就窮困起來。 \end{tabularx} \\ \\ \relax
15:15 & \begin{tabularx}{0.7\textwidth}{X} 於是他去投靠當地的一個居民，那人打發他到田裡去放豬。 \end{tabularx} \\ \\ \relax
15:16 & \begin{tabularx}{0.7\textwidth}{X} 他恨不得拿豬所吃的豆莢充飢，也沒有人給他甚麼吃的。 \end{tabularx} \\ \\ \relax
15:17 & \begin{tabularx}{0.7\textwidth}{X} 他醒悟過來，就說：『我父親有多少雇工，糧食有餘，我倒在這裡餓死嗎？ \end{tabularx} \\ \\ \relax
15:18 & \begin{tabularx}{0.7\textwidth}{X} 我要起來，到我父親那裡去，對他說：父親！我得罪了天，又得罪了你， \end{tabularx} \\ \\ \relax
15:19 & \begin{tabularx}{0.7\textwidth}{X} 從今以後，我不配稱為你的兒子，把我當作一個雇工吧。』 \end{tabularx} \\ \\ \relax
15:20 & \begin{tabularx}{0.7\textwidth}{X} 於是他起來，往他父親那裡去。相離還遠，他父親看見，就動了慈心，跑去擁抱著他，連連親他。 \end{tabularx} \\ \\ \relax
15:21 & \begin{tabularx}{0.7\textwidth}{X} 兒子對他說：『父親！我得罪了天，又得罪了你，從今以後，我不配稱為你的兒子。』 \end{tabularx} \\ \\ \relax
15:22 & \begin{tabularx}{0.7\textwidth}{X} 父親卻吩咐僕人：『快把那上好的袍子拿出來給他穿，把戒指戴在他指頭上，把鞋穿在他腳上， \end{tabularx} \\ \\ \relax
15:23 & \begin{tabularx}{0.7\textwidth}{X} 把那肥牛犢牽來宰了，我們來吃喝慶祝； \end{tabularx} \\ \\ \relax
15:24 & \begin{tabularx}{0.7\textwidth}{X} 因為我這個兒子是死而復活，失而復得的。』他們就開始慶祝。 \end{tabularx} \\ \\ \relax
15:25 & \begin{tabularx}{0.7\textwidth}{X} 「那時，大兒子正在田裡。他回來，離家不遠時，聽見奏樂跳舞的聲音， \end{tabularx} \\ \\ \relax
15:26 & \begin{tabularx}{0.7\textwidth}{X} 就叫一個僮僕來，問是甚麼事。 \end{tabularx} \\ \\ \relax
15:27 & \begin{tabularx}{0.7\textwidth}{X} 僮僕對他說：『你弟弟回來了，你父親因為他無災無病地回來，把肥牛犢宰了。』 \end{tabularx} \\ \\ \relax
15:28 & \begin{tabularx}{0.7\textwidth}{X} 大兒子就生氣，不肯進去，他父親出來勸他。 \end{tabularx} \\ \\ \relax
15:29 & \begin{tabularx}{0.7\textwidth}{X} 他對父親說：『你看，我服侍你這麼多年，從來沒有違背過你的命令，而你從來沒有給我一隻小山羊，叫我和朋友們一同快樂。 \end{tabularx} \\ \\ \relax
15:30 & \begin{tabularx}{0.7\textwidth}{X} 但你這個兒子和娼妓吃光了你的財產，他一回來，你倒為他宰了肥牛犢。』 \end{tabularx} \\ \\ \relax
15:31 & \begin{tabularx}{0.7\textwidth}{X} 父親對他說：『兒啊！你常和我同在，我所有的一切都是你的； \end{tabularx} \\ \\ \relax
15:32 & \begin{tabularx}{0.7\textwidth}{X} 可是你這個弟弟是死而復活，失而復得的，所以我們理當歡喜慶祝。』」 \end{tabularx} \\ \\
[1ex]
\hline
\hline
\end{longtable}
$^{1}$各位弟兄姊妹 早安.
早安.
今天非常感恩榮幸.
可以來到錦繡堂和大家一起分享.
我從小已經在宣道會成長.
所以對宣道會的教會也有一些認識.
當然我還是很小的時候.
錦繡堂也是我聽聞很慈福的一間教會.
今天有幸到此 為大家感恩.
我也跟教會說 今年我只是出去一次廣道.
我只是來這裡 我大病之後一段時間休息.
所以這次也是我大病之後第一次出去廣道.
我一來到的時候 我很感應得到.
心裡面體會到神那份恩典.
透過電樂進唱.
我們的敬拜 小朋友的歌聲.
還有那些 還有那個什麼教導.
我覺得很好.
待會你們父親節 子女做給爸爸.
這個就是令你心歡喜.
謝謝你賜我快樂 教導.
有時候我們中國人 尤其是男性.
母親節就大張旗鼓.
又賣花 又唱歌 又跳舞 跳火圈.
但是一到父親節 我們就變成了沉默的男人.
心領了 明白了 怎樣了.
行了 不用那麼噁心了 行了.
我記得我爸爸還在的時候.
我現在回想起來很後悔.
我還是覺得他說行了.
有時候做些事給他 說些事給他.
他都沒有什麼反應.
但是當我爸爸離開了.
或者我自己一年多月越大的時候.
我就明白 其實這些動作.
做多些是不會令人做噩夢的.
你真心去做 你真的想讓他歡喜.
你真的多謝他教導.
你真心去做 好像小孩那樣.
你是開心 你意會到的.

$^{41}$跟他說句多謝 好像很難開口.
但實際上 爸爸到某個年紀.
某個時刻某個處境 他記得起來.
回憶起來 可能對他是無價的.
所以我希望在座的弟兄姊妹們.
珍惜神給我們有的聲音.
神給我們有的體感.
守之 捉之 舞之 導之.
都是古然上分.
神給我們有的東西.
我們好好地去運用.
今天父親節 在座的父親.
你可不可以舉手.
多謝你 多謝你們.
今天這個訊息分兩部分.
第一部分 我們在經文裡面.
耶穌用了三個比喻.
失約 失錢 失去兒子 浪子回頭.
這兩個失而復得的比喻.
是要表明在《路加福音》第十五章.
二至三節 其實耶穌是因應一個提問.
或者一個處境.
祂帶出這兩個失而復得.
我們應該歡喜高興 高興歡喜.
主要的原因是因為.
「法利賽人和民事私下以論說.
這個人(耶穌)接待罪人.
又同他們吃飯」.
所以 耶穌說這三個比喻的原因.
最直接的原因.
是回應著有些人不明白.
甚至質疑耶穌.
為何你會和罪人去吃飯 接待他們呢?.
在漢語聖經的譯本是這樣翻譯.
「法利賽人和律法師.
咕咕嚕嚕發言說.
這個人接納罪人.
還跟他們一起吃飯 咁嘆嘍」.
所以 你會看到.
其實在今天我想大家去了解.

$^{81}$這個浪子回頭的比喻的時候.
耶穌是刻意地去回應當時一個提問.
就是耶穌 為何你會和罪人去接觸.
當時來說不會是這樣.
你還和他吃飯.
這個吃飯是代表了新層次的接納和認同.
不是現在我們去完餐廳 去牛奶 去大卡羅 去大快活.
吃飯 很多人一起吃飯.
當時不是這樣.
吃飯是代表著一種很深刻的認同 認受.
來到我家 可以和我一起.
這樣的共聚聚餐.
這是一個很難得的事情.
耶穌是要解釋 為甚麼.
天父的心是為甚麼.
這是第一樣我想和大家一起去體會.
因為天父怎樣看我們.
是反映著 或是指導著 影響著我們怎樣做一個父親.
或者我們怎樣去面對一位父親.
所以當我們明白這段經文的核心.
是天父怎樣看世人 怎樣看父親.
是幫助我們怎樣去做一個父親.
和怎樣去面對在世每一位的父親.
第二方面 我又想透過這個比喻.
看到人成長的三個關鍵的位置.
透過浪子回頭.
這個整個比喻的內容裡面.
我們找出人生在世.
我們都和浪子一樣 有不同的階段.
這三個階段是幫助我們認識.
或者我們作為父親 我們明白自己是這樣的.
或者我們面對著我們的子女兒孫.
我們明白是這樣.
所以 第一 我們要知道天父是怎樣.
第二 我們要知道自己和別人是怎樣.
今天講到的重點 就是這兩方面.
開始之前 我想請大家一起去祈禱.
可能今天你來崇拜 你有你的目的.
有你自己的期望.
今天我在這裡想邀請你.

$^{121}$好好地和天父有一點點整理的時間.
因為最大的目的是天父來.
和你一起.
讓他告訴你 你是寶貴.
他憐憫你 他愛你.
所以你是與天父相遇.
讓我們將我們的agenda 想法放下.
讓我們一起和天父講.
我們願意聆聽天上的父親.
對我們的愛 呼召和引導.
我們一起低頭.
我們有一段安靜禱告的時間.
告訴天父 我在這裡.
爸爸 讓你來引導我.
親愛的天父.
人看人是看外表.
你看人是人的內心.
這個心是對神的心.
求主你幫助我們對準你.
我們對天父你的心是信一 真摯 無雜質.
讓我們回應你對我們無盡的愛和恩典的同時.
我們內心裡面沒有任何牽掛.
其他的動機或期望.
我們只是單單焦點於你.
尊崇你 告據你.
願你在我們當中得顯偉大.
得著榮耀稱讚.
聽我們的祈禱 掌管我們的心.
奉耶穌基督得勝的名字 阿們.
第一個階段.
這位浪子他有他自己的期望 想法 生活的方式.
這個階段是一個要別人為自己負上責任的階段.
我做的事有爸爸 有媽媽 有其他人來承擔.
第一個階段 浪子出場的時候.
他要和爸爸分家產.
分家產後要過自己的生活.
錢不是他的 是爸爸的.
他要求過的新生活本身是要爸爸負上責任.
第一個階段是代表了生命裡面的渴望和欲求.
是人的本性 是人的需要.

$^{161}$但不是為自己爭取回來 累積回來.
而是要別人為自己去負上.
我記得我小時候有一個經驗.
那時候沒有那麼多超級市場 百貨公司.
我們逢星期日就會一起有些家庭聚會.
出到佐敦道 有間叫大大百貨.
我不知道大家應該不知道的.
你們看這個樣子 沒人救我老.
所以你們不知道的.
大大百貨 很厲害的.
有兩層 地牢有一層.
很誇張的.
有一次我們一班表弟 表哥去玩.
去到一個玻璃廚房用具那些.
大家玩到興高采烈 女追我族的時候.
整棟掉到地上 皮亂不爛.
我當然呆在那裡.
還有我想是誰做的 是不是我呢.
還是其他人 表哥 表姐.
但是每個人都指著是我 那就是我了.
要付錢的.
在那個場合 我見到爸爸要付錢的時候.
我就明白 原來要付錢的.
還有爸爸付錢付錢.
及時讀書成績不好 要請補習 見家長.
各方面諸如此類.
從小到大總有些階段.
我們是要人幫我們負責任.
因為我們沒有這個能力.
也沒有這種醒覺.
神為我們預備 願意為我們負責任的人.
帶領我們 幫助我們去成長.
今天我們回想看.
其實很多時候.
如果不是有很多人默默的.
甘心為我們付出這些代價.
其實我們未必成為今天這樣.
但很多時候我們只是集中於現在.
我們的需要.
就好像這個老子一樣.

$^{201}$他想這樣.
他的焦點是這樣.
他的期望是這樣.
他心裡面充滿了這份物慾.
他被很多需要 被很多擁有.
充斥著內心.
他忘記了其實是有別人為他負責任.
所以這個第一個階段.
由別人為我們負責任.
這個成長的階梯裡面.
是我們很重要.
是幫助我們認識原來我們是這樣走過來.
同時間今天我們會不會仍然有一部份.
我們仍然有這個未曾長大.
不斷地要別人為我們負責任的地方.
性格.
我們的生活模式.
這個是很值得我們思想.
因為當我們越了解自己.
或者了解父親 了解我們每一個人的時候.
我們就會越來越明白.
天上的父親對我們的憐憫.
寬恕有多深和有多真.
第二個階段.
當這個浪子去到窮途末路.
錢花光 人物用盡.
好朋友都離他而去.
珠瓶狗有.
到他要去到一個窮途末路的時候.
聖經它用了一個名詞.
就是「醒不過來」.
「醒不過來」在希臘文當中.
直譯為「回到自己的身邊」.
意思就是說.
以前的浪子.
那個他不是一個真正的自己.
他被物欲所充斥.
被很多生活上的.
很多不同的滿足.
引誘著 牽引著.

$^{241}$人家說買樓就買樓.
人家說去吃自助餐就去吃自助餐.
人家說要去回鄉下.
他就去回鄉下.
人家說這個好 他就去做.
我們似乎是自由.
但是無意之間.
有意無意地都被很多不同的.
物質的牽引著我們.
但是聖經在這裡說.
第二個階段.
他突然之間醒不過來.
回到自己的身邊.
暗示著浪子.
經歷了極度的痛苦和困難之後.
終於恢復了理智.
開始反思自己的處境.
每個人都會反思.
第二個階段是.
從一個別人為自己負責任.
幫我埋單的人生.
第二個階段提到.
自己為自己負責任.
做回自己.
承擔回自己應該要面對.
和解決的事情.
做回自己應該要做的事情.
他醒不過來.
他醒不過來原來.
我爸爸對我這麼好.
我醒不過來.
為什麼今天我會在這個處境.
在這個和豬爭食物的困境裡.
醒不過來.
這班朋友.
豬朋狗友 走遲欲臨.
那時候是好朋友.
現在是棄我於不顧.
醒不過來.
是透過自身一種痛苦的經歷和體驗.

$^{281}$去明白原來自己是這樣.
對過去所擁有的不珍惜.
對那些珍惜自己幫助自己的人.
視而不見 忽略.
而最終.
他願意去承擔棄自己.
願意面對自己生命的真相.
願意承認這個事實.
這個痛苦.
是自己拿來.
是自己的問題.
我想問問大家.
這樣好不好.
一個願意為自己負責的人.
好不好.
好 你是不是這樣的人.
為什麼你們不立即點頭.
是因為你覺得.
害怕點頭後被人發現.
剛好有人問你.
你是不是這樣的人.
想不想是這樣的人.
立即多點反應.
這個問題真的挺難回答.
如果我們每個人.
只是為自己負責任.
已經很好.
不過我覺得還不夠好.
如果這個世界.
每個人都只是為自己負責任.
我們都會容易.
是好的 不會成為別人的負累.
但都會容易成為一個自我的人.
第三階段 最後一個階段.
是這個悔改.
明白自己的處境.
承認今天我的處境.
我看到這樣的一個警方.
是我自己拿來的.
我願意去面對.

$^{321}$我願意去承擔.
而京雲告訴我們.
他做了一個很重要的動作.
就是回家.
我試過離家出走一年.
我還記得.
我真的願意回家.
但那個時刻是很難.
對不起 因為我電療化療之後.
喉嚨受傷了.
所以抱歉.
有時你要看著我喝水.
希望你們不要介意.
知道自己錯.
承認自己錯.
願意面對自己的錯.
說是很容易.
願意回頭.
回家.
我覺得這個腳步是難.
意味著不單止承認.
知道 承擔.
自己做錯的事.
要面對之餘.
別人的看法.
別人是否給你機會.
你是否真的做得到.
你願意改變.
是否真的變得到.
還是很快就做回自己.
去回第一個階段.
就是繼續要人幫你負責任.
所有的事.
我們看見.
其實帶著很多不定的因素.
這就是真實的人生.
各位弟兄姊妹.
聖經的道理.
不是讀出來 計出來 說出來.
說一下就算.

$^{361}$而是在生活裡.
我們都明白.
我們是有很多不定的因素.
而是在生活裡.
我們都明白.
我們都是曾經.
或者現在你都是一個.
要別人為你負責任的人.
可能我們形態上.
或者我們在做很多事.
都是為自己負責任.
我們是一個盡責的人.
稱心的人.
但是在我們內心深處.
可能我們都會體會到一種無奈.
一種張力.
一種距離.
是我們自己難以逾越的.
你的脾氣.
你的價值觀.
小時候沒有錢.
可能我們看錢很重.
對教會的奉獻.
對別人的時間.
可能我們都很計較.
因為我們受著很多的影響.
受著有很多.
在性價值成長上的一些關聯.
所以我們生命裡面.
都有些位置.
是需要更新和不斷地突破.
所以今天這個堂道.
仍然對你有價值.
不是因為我們已經知道.
這個故事的來龍去脈.
掌握了所有有關浪子和父親.
大兒子和小兒子的關係.
大兒子是反理賽人.
小兒子是罪人.
天上的父親.

$^{401}$在路加福音的架構裡面.
三重的比喻.
都叫我們高興歡喜.
都提到失而復得的重要.
信仰不是只是資料.
和所謂我們掌握的真確性.
信仰更加是這些資料和真確性.
和我們的關係.
和你們的關係.
你怎樣對你的父母?.
你怎樣做父親?.
你怎樣對你的父親?.
你怎樣對教會的父親?.
牧者?弟兄姊妹?.
這些才是天父的心最關注.
所以第三階段.
為別人負上責任.
願意承擔自己的需要責任.
而且他知道爸爸的需要.
他知道哥哥的需要.
他知道其他人的需要.
他更加知道得罪了天.
得罪了父.
他知道天父的需要.
所以他願意為別人負上責任.
走上悔改之路.
來到最後.
我要結束今天的講道.
如果在座當中.
你是爸爸.
或者你是新朋友.
我想鼓勵你.
今天我們都是一個可以回轉.
回頭.
慰駕.
回到天父的身邊.
我相信這裡有些長者.
和我一樣.
人生路高高低低.
經歷過很多事情.

$^{441}$可能我們已經決志信了耶穌.
就好像耶穌在這裡和罪人的相處.
是為了尋找拯救失望的人.
所以今天如果有這個機會.
你可以去到耶穌面前.
接受祂的這份愛和寬恕.
這樣非常好.
父親節.
去做這個接受主.
相信主.
跟隨主的決定.
如果在座當中.
很多作為父親.
或者你有你的父親.
你有你的家人.
我們可以這樣去看.
天上的父親.
透過這個訊息去告訴世人.
他對人的愛.
恩慈.
憐憫.
和寬恕.
是很確實的.
我結束的時候.
我集中講寬恕.
在這裡.
我們看到這個比喻裡面的父親.
他在家門口.
等待他的兒子.
存著一個盼望.
並且他信任這個兒子會回來.
而他回來之後.
對他所作所為.
沒有再提.
當他說.
兒子跟我說.
我不配做你的兒子.
我只可以做你的僕人.
他穿上好的衣服.
和牛羊.

$^{481}$為他去擺相.
開派對.
高興.
歡喜.
在這裡.
我想各位弟兄姊妹.
你明白.
耶穌講這三個比喻.
是鼓勵我們.
人世間總有失前.
失陽.
人的失去.
人心的失去.
生命的失去.
但是天父.
是帶著期盼.
帶著信心.
帶著愛.
和寬恕.
等待我們回轉.
無論今天你是哪個階段.
要人為你負責任.
自己為自己負責任.
你願意為別人負責任.
為社會.
為人群.
為地方.
去謀幸福.
你都要記住.
是神的愛.
恩典.
憐憫.
寬恕.
臨到我們.
呼喚我們.
所以我們可以學習.
去做一個地上的父親.
分享神的愛和恩典.
我記得我那時候沒有書讀.
我的家人和親戚幫助我.

$^{521}$讓我有機會讀書.
那副校長在那個面試室.
看我的成績.
又被舊的那間學校趕出來.
踢出來.
叫做什麼.
Excommunication.
就是那些叫做.
港卒校.
是那些校.
他就說到很真實.
你的成績那麼差.
科科都不合格.
對一科中文都合格.
還是57.
他說科科紅到.
評語惡劣還緊.
什麼爛順.
很難不給這個機會.
雖然你是誰介紹.
親戚又說什麼表弟.
但是我們做人很公正的.
我們是一間基督教學校.
不緊要.
不如你說一下.
為什麼我們要給你機會讀書.
不知道為什麼.
可能我那時候還沒有相信.
不過剛剛有試.
小時候很小的時候.
就回過教會.
所以可能我們要聖靈感動.
我不敢這樣說.
我作為牧師這方面.
要很小心.
不過大家都明白.
按照我們宣導會.
我們的解釋都是.
被聖靈充滿掌管.
我們願意過一個合神生命.

$^{561}$並且願意聖靈帶領我們的生活.
這個解釋是.
但是那時候我沒有這些.
我只是覺得.
是啊.
老師你說得對.
其實我真的不值得給我這個機會.
我知道其實我真的很差.
不過我也想告訴你.
如果你今天給我這個機會.
我可以用日後.
我的生活.
或者我有什麼改變.
去讓你知道.
究竟為什麼你要給我這個機會.
我心想這個應該.
我現在回想.
這個應該是很頂級的銷售員.
或者是一個叫做.
很機智.
很敏銳的人.
才可以說到這些.
好像黃子華的對白.
而且我還很記得.
那個老師是這樣.
然後說.
行了 你出來回去等消息吧.
我也不知道.
但後來這個老師就教我.
中國文學.
我就很用心地讀.
還拿了幾個A.
寬恕對每一個人都很重要.
因為你和我都錯過.
或大或小.
或明或惡.
或知或未知.
我們都得罪過人.
得罪過天父.
如果我們的能力為自己負責任.

$^{601}$或者為別人負責任.
而忘記了.
其實由始至終.
是天父為我們負責任.
為我們釘死在十字架上.
我們所謂的人生.
只不過是用自己的方法.
努力營運.
那個你覺得好.
或者世人覺得好的人生.
唯有我們甘心承認.
其實完全是天父的恩典和寬恕.
所以今天我有機會在這裡講道.
分享我少少的片段.
我覺得不值一親.
只是上帝的恩典.
上帝的憐憫和寬恕.
求主的憐憫,恩典和寬恕.
常在你們生命當中.
讓我們在人生不同的階段.
都能夠明白天父的心.
愛你.
等待你.
讓你成為別人的祝福.
我們一起低頭.
我們祈禱.
親愛的天父.
幫助我們回到你身邊.
如果我們和你的關係若即若離.
求主吸引我們.
讓我們快跑跟隨你.
如果在座當中.
有新來的朋友.
未曾認識主耶穌.
我鼓勵你.
你打開你的心.
跟隨這位天父.
接受這位天父.
祂對你的愛.
祂差遣獨失愛子耶穌基督.

$^{641}$為你生,為你死,為你復活.
擔當起你一切的罪過.
賜下憐憫,賜下寬恕.
讓我們得著這份愛和救恩.
讓我們一起在永恆裡.
得著這份盼望.
願主賜福我們眾人的心.
讓我們一起回到天父的身邊.
多謝主耶穌.
多謝天父.
多謝聖靈.
聽我們的祈禱.
奉耶穌基督得勝的名字.
\newpage



\section{使徒行傳 1:1-15}
\label{sec:aAKtjwj8FbQ}
\textbf{2024-06-23 午堂《復興的等候》使徒行傳 1\_1-15  講員:陳麗蟬牧師}
\newline
\newline
連結: \href{https://youtube.com/watch?v=aAKtjwj8FbQ}{\texttt{https://youtube.com/watch?v=aAKtjwj8FbQ}} ~~~~ 語音日期: 2024-06-25
\newline
\newline
\hyperref[sec:Yzl3vxypdqk]{\small{< < < PREV SERMON < < <}}
~
\hyperref[sec:index_chronic]{\small{[返順時目]}}
~
\hyperref[sec:index_scriptual]{\small{[返順卷目]}}
~
\hyperref[sec:vr_pYc_qVuc]{\small{> > > NEXT SERMON > > >}}
\newline
\newline
使徒行傳 1:1-15
\newline
\begin{longtable}{cl}
\hline
\hline
章節 & 經文 (和合本修訂版)\\
\hline
1:1 & \begin{tabularx}{0.7\textwidth}{X} 提阿非羅啊，我在第一本書中已論到耶穌從開頭所做和所教導的一切事， \end{tabularx} \\ \\ \relax
1:2 & \begin{tabularx}{0.7\textwidth}{X} 直到他藉著聖靈吩咐所揀選的使徒後，被接上升的日子為止。 \end{tabularx} \\ \\ \relax
1:3 & \begin{tabularx}{0.7\textwidth}{X} 他受害以後，用許多確據向使徒顯明自己是活著的，在四十天之中向他們顯現，並講說神國的事。 \end{tabularx} \\ \\ \relax
1:4 & \begin{tabularx}{0.7\textwidth}{X} 耶穌和他們聚集的時候，囑咐他們說：「不要離開耶路撒冷，但要等候父的應許，就是你們聽見我說過的。 \end{tabularx} \\ \\ \relax
1:5 & \begin{tabularx}{0.7\textwidth}{X} 約翰是用水施洗，但過了不多幾天，你們要在聖靈裡受洗。」 \end{tabularx} \\ \\ \relax
1:6 & \begin{tabularx}{0.7\textwidth}{X} 他們聚集的時候，問耶穌：「主啊，你就要在這時候復興以色列國嗎？」 \end{tabularx} \\ \\ \relax
1:7 & \begin{tabularx}{0.7\textwidth}{X} 耶穌對他們說：「父憑著自己的權柄所定的時候和日期，不是你們可以知道的。 \end{tabularx} \\ \\ \relax
1:8 & \begin{tabularx}{0.7\textwidth}{X} 但聖靈降臨在你們身上，你們就必得著能力，並要在耶路撒冷、猶太全地和撒瑪利亞，直到地極，作我的見證。」 \end{tabularx} \\ \\ \relax
1:9 & \begin{tabularx}{0.7\textwidth}{X} 說了這些話，他們正看的時候，他被接上升，有一朵雲彩從他們眼前把他接去。 \end{tabularx} \\ \\ \relax
1:10 & \begin{tabularx}{0.7\textwidth}{X} 他升上去的時候，他們定睛望天，看哪，有兩個人身穿白衣站在他們旁邊， \end{tabularx} \\ \\ \relax
1:11 & \begin{tabularx}{0.7\textwidth}{X} 說：「加利利人哪，你們為甚麼站著望天呢？這離開你們被接升天的耶穌，你們見他怎樣升上天去，他也要怎樣來臨。」 \end{tabularx} \\ \\ \relax
1:12 & \begin{tabularx}{0.7\textwidth}{X} 有一座山，名叫橄欖山，離耶路撒冷不遠，有安息日可行走的路程。那時，門徒從那裡回耶路撒冷去， \end{tabularx} \\ \\ \relax
1:13 & \begin{tabularx}{0.7\textwidth}{X} 他們一進城，就上了所住的樓房；在那裡有彼得、約翰、雅各、安得烈、腓力、多馬、巴多羅買、馬太、亞勒腓的兒子雅各、激進黨的西門，和雅各的兒子猶大。 \end{tabularx} \\ \\ \relax
1:14 & \begin{tabularx}{0.7\textwidth}{X} 這些人和幾個婦人，包括耶穌的母親馬利亞，和耶穌的兄弟，都同心合意地恆切禱告。 \end{tabularx} \\ \\ \relax
1:15 & \begin{tabularx}{0.7\textwidth}{X} 那時，有許多人聚會，約有一百二十名，彼得在弟兄中間站起來，說： \end{tabularx} \\ \\
[1ex]
\hline
\hline
\end{longtable}
$^{1}$各位姐妹早晨.
孫家在上半年的講道.
我們自己同工的講道.
主要都是講這個舊約裡面的人物.
特別是這個創世記裡面的人物.
藉著他們的生命.
就成為了我們今天提醒.
今天那個的尷尬.
到下半年了.
我們由今天開始.
同工就會主要是講這個使徒行傳.
新約的使徒行傳.
我們很想藉著初期教會他們的一些經歷.
成為我們今天教會的借鏡.
在剛才第一節裡面他講道.
他就說他在第一本書下.
作者下是盧家醫生.
在第一本書已經提到耶穌基督所有教導的事情.
第一本書就是盧家福音.
我們知道這個盧家福音和使徒行傳.
都是盧家醫生所寫的.
他的修信人就是一個人叫做提亞飛羅.
提亞飛羅這個真正是什麼人呢?.
沒有人知道的.
因為他只是提了個名字.
但我們在歷史,在聖經其他地方看不到他.
但就知道他一定是一些高望重的人.
因為他講到提亞飛羅大人.
是一些地位比較高的人.
他的意思就是愛神的人.
或者是神所愛的人.
所以盧家寫這封使徒行傳.
是去給一些神所愛的人.
而他也是愛神的人.
剛才讀了第一章第一至第十五節.
講到耶穌在升天之前.
就吩咐那些門徒.
他們要去將這個福音傳到地極.
命令不單止給昔日的門徒.
也在今天裡面給過你和我.

$^{41}$要將耶穌基督救贖的恩典傳到地極.
我有時候也喜歡看漫畫.
這位就是蔡頌輝牧師.
他是一位牧師.
但他很擅長畫這些漫畫.
這些漫畫在當中有時給我們很多反省.
講到信徒.
我們看到那間教會有很多錢.
有很多鑽石 有很多音樂.
很多東西.
信徒就說.
哇 這間教會很復興.
但是他另外一個聲音就說.
這不是復興.
只是這間教會是很熱鬧的教會.
其實從使徒行傳裡面.
我們一直看下去的時候.
我們發現到.
若果一間教會要興旺起來.
要復興的話.
他一定要從禱告開始.
藉著聖靈的充滿.
給我們有動力.
將這個救恩一直傳開它.
所以若果那間教會.
沒有真實的傳講上帝的話.
沒有實際的行動來傳福音的時候.
這裡所有辦的一切.
都不是一個真正的復興.
其實當耶穌基督升天之前.
他就叫那些門徒.
你們要專心的去等.
等什麼呢?.
等聖靈降臨.
但其實他們不知道聖靈是什麼來的.
只是在約翰福音的時候提過.
這個聖靈就是真理的聖靈.
是保衛師教我們怎麼去敬拜.
是去提醒我們知道我們有罪.
所以實質的聖靈是什麼呢?.

$^{81}$那些門徒都不知道的.
但是耶穌就說.
你要去專心的去等候.
到日期滿足的時候.
祂就會成就各樣的事情.
我們來看看.
在等候的期間.
這一幫門徒他們在做些什麼呢?.
在第一章的第三,第五節.
他這樣說.
他說耶穌基督受害之後.
用很多的確據,用很多的事情.
證明祂是復活了.
祂和那幫門徒生活了40天.
在那40天裡面.
祂有教導上帝的國度.
究竟是怎麼樣的國.
有吩咐他們不可以離開耶路撒冷.
要等候父的應許.
這個父的應許.
就在約翰福音裡面.
耶穌基督去請求天父.
將這個聖靈臨到地上.
給過那幫門徒.
他說這個就是我曾經和你們說過的.
還有他們說.
他們這幫門徒已經受了約翰的洗了.
但是他們遲些就要受聖靈的洗.
所以這幫門徒聽了一頭霧水.
都不知道他們在說些什麼.
整個使徒行傳.
我們看到路加醫生.
他很著重說神的國度.
和聖靈的降臨.
這兩樣東西很有關係.
在神的國度.
「神的國」這個字.
在路加福音裡面出現了32次.
也很多了.
在使徒行傳就出現了8次.

$^{121}$我們看到醫生很看重神的國度.
在使徒行傳外.
第一章就說到神的國將要降臨.
到了第28章使徒行傳最後那章.
他仍然是在說保羅雖然在坐牢.
但是他仍然在傳上帝的國.
為什麼路加醫生.
他這麼強調神國的降臨呢?.
要傳揚神的國呢?.
使徒行傳就寫在大約公元62,63年左右.
當時耶路撒冷一帶的地方.
羅馬政府就差派了希律阿基伯一世.
去統治猶太人.
希律阿基伯一世有猶太血統.
所以他希望拉攏更多猶太人來支持他.
於是那些猶太人時時做一些小動作.
去追逼基督徒.
因為基督教傳得很快.
發展得很快.
所以希律一直支持猶太人去打擊基督徒.
所以那些基督徒當時還沒有去到最嚴峻的迫迫.
但是因為希律王加上那些猶太人一起的時候.
就給他們很多誣告的控訴.
隨時可以抓人.
你看看保羅當他未信耶穌之前.
他可以拿著令旗到處去抓基督徒.
可以不經審訊就抓他們去坐牢.
所以在整個形勢裡面.
這個情況越來越嚴峻了.
路加不斷提神的國.
就在提醒當時的信徒.
其實也都提醒我們今天.
當我們面對一些艱難.
一些挑戰的時候.
我們要堅定的相信一件事.
神自己掌權的國度.
是和平的 是公義的.
祂一定會臨到地上的.
但是雖然耶穌講了很多次了.
但是這幫門徒都好像不太明白.

$^{161}$他們仍然問.
當耶穌升天那天.
仍然問.
主啊 你復興以色列國.
是不是就在現在了.
是不是這個以色列國.
就要脫離羅馬的統治了.
他們真的很想快點脫離羅馬人的統治.
因為羅馬他們的政府.
又有不同的壓迫.
其實很自然的.
當我們面對一些.
不是很好的環境的時候.
一定是想快點走的.
是不是.
譬如你在工作的當中.
你面對的上司很壞的.
你想辭職.
你面對的同事很難搞的.
又不合作的 是不是.
你一是就想辭職.
一是你就想自己走.
或者我們有時工作很大壓力.
我們也都想離開的地方.
我們也都看到現在香港的情況.
我們有些很特別的年輕人覺得.
哇 香港真的不行了.
好像經濟越來越衰退.
政治越來越不知道會怎樣的時候.
同樣很多人都會離開.
但是耶穌祂不斷說.
神的國和地上的國是不同的.
在地上的國裡面.
真的會有很多不公義.
很多壓迫.
很多不合理的事情.
但是上帝的國.
就是不同地上的國.
地上的國就有地域.
香港就有新界.

$^{201}$一直不知道去到哪裡的地方.
有國土的.
但是神的國就沒有國土的.
它不會受住區域的限制.
神的國也都在當中沒有政治.
沒有名利 沒有金錢的.
神的國是超越了這個地域.
超越了時間.
超越了世上任何的物質的.
合理來說.
神的國究竟是什麼呢?.
神的國就是一群願意信服上帝.
帶領和安排的人.
他們走在一起.
他們願意學校上帝給我們律例典章.
剛才我們讀了十課了.
這群神國的人.
就是他們願意去跟從上帝給我們的教導.
就是說.
今天你和我已經加入了上帝的國度裡面了.
所以你和我都要活出這個天國.
這個上帝的國裡面的價值觀.
按照上帝的人情和平憐憫公義.
去展現在上帝的國度裡的輝煌和和平.
大家認不認識Ken 弟兄?.
徐正寧弟兄.
他是我們錦宣家的弟兄.
不過他現在搬了出去筲箕灣.
和他工作就在柴灣.
所以他就沒有再回來錦宣家了.
他就在外面的教會來崇拜和事奉.
Ken 我昨天也打了電話問准他.
我今天港獨會提醒你.
他很樂意.
他三十多歲就中風.
他告訴我.
他之前的生活是很頹廢的.
喝酒 吸煙.
還很安靜的告訴我吸了毒.
他說這些東西樣樣都齊.

$^{241}$但是他中了風之後.
醫生告訴他.
你會成為永久傷殘的人士.
他在床上因為不能動 走不動.
他就開始反省他自己的人生.
他反省到他之前的人生.
真的很頹廢 很沒用的.
他除了喝酒 吸煙 吸毒之外.
還有很多的頑皮事不告訴你們.
他還經常爭第一.
樣樣東西都希望自己最威風的.
但是病床上高.
他自己覺得這樣的生活.
其實傷害了他爸爸媽媽.
他覺得爸爸媽媽供養他讀書.
又去外國讀書.
但是他讀完書.
他竟然在外國的生活這麼敗壞 這麼頹喪.
浪費了爸爸媽媽很多錢.
也有很多的行為.
令到他爸爸媽媽很擔憂.
特別是他媽媽真的非常擔心他.
他覺得原來之前自己的行為是這麼差的.
他都覺得他對不起爸爸媽媽之外.
他對不起天地間的神.
他說他那時候還沒信耶穌.
他只覺得天地裡面有一個主宰.
他覺得自己這樣做人法.
這樣得罪了上帝.
他心裡面又向天地間的上帝說.
真是對不起了 對不起.
當他信了耶穌之後.
剛才有些人都點頭認識他.
一看到他整個人很不同了.
他說他中風前的生命.
就在地上高嶺去打滾.
爭權 爭第一.
他又暗中做了一些不合真理的事情.
但是當他成為了永久傷殘人士.
那時候他就說朱子一又去探望他.

$^{281}$另外去了宣教的那位周小嫻姑娘.
又去探望他 回來又跟他一起經.
後來他信了耶穌.
我想大家都看到.
他信了耶穌之後.
你沒有辦法想到他之前是怎麼樣的事情.
他的生命有很大的轉變.
醫生就跟他說.
他說你現在才三十多歲.
你現在中了風了.
你不能動了 走得不好了.
他說但是你這麼年輕.
你一定要努力去過一個自立的生活.
所以無論多艱苦都好.
你都要去學習.
我記得他開始第一份工作.
去爸爸的朋友.
在葡萄園那裡有一家餐廳.
很多踩單車路過.
就去喝杯水.
教會有些是弟兄姊妹.
來不及過去支持他.
他就在那裡做侍應.
不過是一個正正式式的祈禱.
他真的只是站著.
因為他手無力.
寫不了單.
手無力 他捧不了餐 捧不了碗.
他只是站著.
有人來到的時候.
他就說歡迎請坐.
歡迎請坐.
他真的只是做祈禱.
等到告訴他媽媽.
他穿得很斯文 很整齊.
又穿著背心.
我說 哇 他就真的很有禮貌.
整個紳士感.
他真的整個生命完全的改變.
他不單止努力去學.

$^{321}$去過一個自立的生活.
他更加努力要去學習.
做一個在神國度裡面的人.
將之前的迂腐.
將以前的不好的東西.
全部都扔掉.
他要做一個聖潔.
和平 滿有主持愛的人.
他說 如果讓他再選擇.
他都會選擇中風這條路.
因為當他還沒中風之前.
他就覺得自己很厲害 很厲害.
但原來整個生活是這麼不好的.
但當他中了風.
去反省人生得到救恩的時候.
他整個生命是不再一樣.
這個神的國就是這樣.
在這個世界的國裡面.
可能我們追錢 追地位 追名 追利.
但是這個神的國.
是一個不再去談這些東西.
我們只是談.
怎樣走出上帝的人愛和平的慈愛.
我們要在一個.
雖然未必一定很有錢.
但又未必一定很沒有寵.
看上帝的心意.
不過在這個神國的掌權底下.
我們可以過一個又豐盛.
又有意義的生命.
又看一下這幅畫.
這次說到不知為何.
就想起這個蔡仲輝牧師.
看他的畫.
他就說到這個信徒.
他就見到一個很慘的人.
那個人就說.
哎呀 請你停下來幫幫我吧.
但那個信徒怎麼樣?.
很忙碌.

$^{361}$他說 不好意思.
我現在趕著去開一個宣教的會議.
我遲到了 我要先走了.
你看看那隻小老鼠說什麼.
宣教?.
有時就是這樣.
我們在聖經的真理裡面.
我們很多的上帝教了.
我們要去傳福音 要去宣教.
但是當實質要去.
真是這個傳福音宣教.
不是就這樣將一套的理論告訴別人.
而是要在當中去實踐.
上帝的人外憐憫出來.
當人這麼有需要的時候.
卻是為了開會.
不肯停下來去幫助人.
這個是不是一個真正的.
實踐到上帝給我們.
那個傳福音宣教的使命呢?.
這個就是我們實踐神國裡面.
一個生命的反省.
接著我們看看.
在實踐神國這個使命的時候.
我們在等候.
因為聖靈還沒有降臨.
那麼他們又在做些什麼呢?.
那麼耶穌吩咐.
你們不要離開耶路撒冷.
因為你們要等候父的應許.
直到聖靈降臨的時候.
你們就要將這個福音傳.
從猶大,撒瑪利亞.
一直傳到地極.
要等候.
當時那幫門徒其實不知道要等多久.
我們等耶穌死了之後.
三天復活了.
等三天.
短短的.

$^{401}$接著跟祂生活了40天.
現在又要等候.
等候聖靈.
很辛苦 因為你不知道要等到什麼時候.
現在好一點.
我們等村巴.
你剛剛去到這個村巴站.
村巴走了.
你就很生氣了.
不知道下一班車什麼時候來.
不過現在很好.
有些Apps可以看看.
下一班車.
我今天在新元宗那裡.
剛剛走了.
所以我排第一.
我最憎恨的就是排第一.
排第一那個.
最憎恨的.
那麼我看看.
下一班車什麼時候.
星期天很疏的.
五分鐘.
有沒有搞錯.
五分鐘.
其實在這個急速的社會裡面.
我們最不喜歡的就是等.
是不是.
這個社會現在是說什麼.
說速度的.
說要快的.
祈禱的時候.
我們希望很快就能夠蒙英運.
我們跟人相處的時候.
希望能夠很快就建立到這個關係.
我們住樓的那些.
等電梯.
是不是.
特別有些很心急的那些人.
一出去按電梯.

$^{441}$五秒鐘的時候.
電梯還沒到就說.
電梯壞了.
於是再按按按.
按多幾下.
以為按多幾下.
電梯就快點就會來到.
其實經常電梯壞.
就是因為人按得太多.
我們都不喜歡等.
這個慢的節奏.
慢吞吞的節奏.
令我們很辛苦.
有些電影姐妹提醒我.
有時星期天.
見我走路.
那些腳步很急速的.
很快的.
有時見我.
腳步走得不夠快.
不過頭和頸都要走快一點.
有弟兄提過我.
不過弟兄就禮貌一點.
不是禮貌一點的.
他不想傷我的心.
有個姐妹.
你知不知道.
你的頭和頸先走.
很難看.
我就上網看.
當我要出來講道.
由那個位置走出來的時候.
咦 有時真的是這樣.
是不是很難看.
我們不喜歡等.
我真的發覺到原來我真的這麼快.
有同工跟我一起坐穿巴.
在總站下車.
他經常問我.
你走這麼快趕著去哪裡.

$^{481}$我趕著回辦公室.
他還沒夠衝.
於是有時他就拖著我.
他因為要拉著我.
不准我走這麼快.
但其實我覺得他盲吞吞.
我很辛苦.
我們熱心地去為神工作.
希望快點有果效.
希望快點有果效的時候.
就讓人看到.
我有努力的.
我有成就的.
不過上帝.
上帝往往就要我們安靜地等候.
去等.
所以在這個升天前.
耶穌就吩咐這班門徒.
在耶路撒冷裡面去等.
等多久.
不知道.
等聖靈降臨.
祂是甚麼.
都是不知道.
總之就是等.
這班門徒跟隨了耶穌這麼久.
他們知道耶穌的心意了.
他們等在裡面.
要做些甚麼呢.
他們就回到這個牢房裡面.
他們一起的禱告.
施行傳這樣形容他們.
他們在這裡行切的禱告.
同心合意這樣來禱告.
這個行切的意思.
就是要不斷地.
不停地.
不斷地去求神.
求神.
你之前給我們的應許.

$^{521}$你快點到來吧.
這個同心合意.
很重要.
就是要有同一的思想.
同一的理念.
等候這個聖靈的降臨.
這個同心合意.
在這個.
《路加福音》和《使徒行傳》.
總共出現了十二次.
其中十一次.
就在《使徒行傳》裡面出現了.
路加醫生為甚麼這麼著重.
這個同心合意呢.
就因為初期教會的時候.
他們面對這個南亞政府的政治壓力.
面對猶太教很多無理的.
在這裡周圍的搗亂.
拉他們來去坐牢.
這群初入教的信徒.
他們就更加需要.
來同心合意這樣來禱告.
在第一次這個教會的祈禱.
就是說在這個初期教會.
在《使徒行傳》裡面形容.
他們第一次的禱告會是多少人呢.
120人.
我相信如果有120人.
一起來去同心合意.
行設禱告的時候.
這個力量是非常之大的.
這個路加醫生就列出了.
這些人其中的11個.
再加上耶穌的媽媽瑪利亞.
我們看看這11個門徒是什麼人.
記一下他們.
有彼得.
我們知道彼得是那個粗心大意.
又懶的又失敗的一個門徒.
還有那個經常信心不穩定.

$^{561}$不信耶穌是復活的多瑪.
還有一個叫做什麼.
憤慮黨的西門.
就是說這個是衝動派.
還有那兩兄弟爭第一的約翰和雅各.
我們看到這班門徒.
都不是一些很完美的人.
但他們竟然可以在上帝的面前.
同心合意的來祈禱.
求聖靈的降臨.
他們等到了.
就在十天之後.
在五場節他們等到了.
其實這個教會.
每一間的教會.
都是來自不同文化.
不同背景.
不同色彩的人.
不同弟兄姊妹.
他們每一個人都是不完美.
不過我們就要走在一起的事奉.
這個才是教會.
在這麼完美這麼不同的底下.
我們還要怎麼樣呢.
是要同心合意來祈禱.
來等候神.
我記得上個星期.
我和Patrick聊天.
他告訴我看過宣道會的年報.
你們有沒有看過.
他說.
錦宣家祈禱會的人數有多少人.
四十幾人.
我們錦宣家有多少人.
應該六百多人.
七百多人.
四十多人祈禱.
是不是少了一點呢.
昨晚又是我們這個音樂報道會.
在過往逢有聚會的時候.

$^{601}$我們就會設計一些橫額.
掛在禮堂外面.
讓街坊們都知道.
弟兄姊妹回來都知道.
昨晚的音樂的form聚會.
記不記得我們外面有沒有橫額.
有一兩位弟兄姊妹有表示.
其他不知道.
所以就是說.
教會有沒有宣傳的橫額.
你們全部都看不到.
你不回答 我懶得回答.
答案是沒有的.
因為今年我就會覺得.
放一個橫額在那裡.
哦 有啊.
或者甚至看不到的.
是不是.
但我又會覺得.
一個在這些form聚會裡面.
最好的宣傳就是什麼呢.
禱告.
最好的宣傳就是這個禱告.
所以在過幾個月之前.
我已經在這個五分鐘的代禱事項裡面.
其中教會的那一欄.
每一次我都在為這個報道會.
有時就為了事奉人員.
有時就為了同心圓的那些負責人.
弟兄姊妹帶微信的朋友.
或者微信朋友的心靈.
還有每一個的祈禱區都鼓勵.
甚至另外發起一些特別的小組.
來一起為這個報道會來祈禱.
還有我邀請了大家.
在過幾個月之前.
我們的同工邀請大家.
在這個月裡面.
就每一天都要為自己家人.
或者一些微信的朋友.

$^{641}$每天都為他未得福音來祈禱.
祈禱你快點相信.
不知道有多少會做了.
不過我知道很多弟兄姊妹都做.
雖然未必能夠每天.
我自己坦白都未必能夠每天.
但我起碼有行動.
在全方位事情上.
你和我都要有份.
其實在兩個星期之前.
就快到這個報道會了.
我心裡面有些害怕.
慘了 這次沒有做這個Banner.
我又好像沒聽弟兄姊妹說.
有沒有帶到人來.
我只聽到有些人說.
我正在邀請人.
好像沒有任何人告訴我.
他邀請到.
我心裡面慘了.
要不要做一條Banner呢.
幾百元.
我又看到同心圓.
出了一張Poster給我們的時候.
看到我們的講員.
這麼年輕.
我說.
青年人或者少年人的報道會.
差不多.
也有些弟兄姊妹來跟我說.
那個講員這麼年輕的樣子.
本來是想邀請他的姐姐.
和一些同事來的.
他說還是不要了.
我想都是青少年人來的.
所以害怕了.
從來沒有人正式告訴我.
邀請到人來的.
我說不要緊.
祈禱了.

$^{681}$發起了這麼多的禱告.
昨晚有多少人出席呢?.
232人出席.
當中一開始的時候.
問一下有多少第一次.
或者是未信耶穌的來到的.
有二十多隻手舉手.
當中有少年人.
有青年人.
有成年人.
也都有不少的老人家.
最後我拿到的缺字.
是有六位的缺字.
在裡面也都有老人家.
有三位慕道.
其實有一對女孩信了耶穌的.
不過她不肯出來.
又說要走.
就走了.
所以我覺得.
是上帝要我們去學習的.
就是要更多的禱告.
每一樣事情我們要去等候的時候.
不是就這樣坐在那裡呆等.
要去禱告.
並且要跟周圍的弟兄姊妹.
同心合意.
行切的來去禱告.
保羅在《哥林多前書》第九章裡面說到.
「我傳福音.
完事沒有可誇的」.
在這個報導裡面.
真的帶了很多人信耶穌的.
保羅.
他說他不能夠不傳.
因為上帝就給了這個使命他.
他接著下一句說什麼呢?.
「如果我不傳的時候.
我就有禍了」.
弟兄姊妹.

$^{721}$要當我們覺得信了耶穌很久了.
我們都覺得很怕跟別人講福音.
我就這樣坐在那裡的時候.
記住這句話.
「如果我不傳的時候.
我就有禍了」.
這個「傳」.
可能在言語上的「傳」.
也都可能在我們生命見證裡面的「傳」.
也都有位姐妹告訴我.
她說真的跟家人不傳福音很難.
她姐姐很難講.
但是她經過很多祈禱之後.
她終於可以跟她姐姐講.
她說以前跟她講耶穌兩個字都不敢.
就這樣是一個見證.
所以弟兄姊妹.
要教會復興.
我們需要有更多更多的人.
加入這個禱告的行列.
所以弟兄姊妹.
每一個星期我們的程序表裡面.
五分鐘的代禱事項.
還有每個星期.
我們有這個猜拳網站.
猜拳的一些代禱事項.
很需要弟兄姊妹.
用這個禱告來支持他們.
他們在前線工作.
我們在後方.
我們以禱告.
跟他們一起打這個屬靈的爭戰.
在最後裡面.
等候的結果是怎樣呢?.
就是這一班祈禱的人.
120個人.
他們都得到聖靈的充滿.
得到聖靈的充滿.
在當中可能遲了兩個星期.
童宮就會講了.

$^{761}$這一班本來很懦弱的.
是一個武膽匪類來的.
但是他很有勇氣.
被人抓都好.
他都要起來 站起來講道.
原本是一班絕口笨舌的人.
他們要出來見證耶穌.
原本很吝嗇的人.
他為主的國度賣這個田產.
原本是一班很優秀.
沒什麼力的人.
但是他主的裡面變得喜樂.
所以弟兄姊妹.
我們想一想.
在今天裡面我們的生命.
我自己屬靈的生命是怎樣的呢?.
我自己滿不滿意我這個生命呢?.
就想一想.
耶穌滿不滿意我跟祂的關係呢?.
弟兄姊妹.
你想不想在這個生命裡面得到更生呢?.
你渴不渴慕這個聖靈充滿呢?.
這個聖靈充滿.
不是像一些宗派那樣說.
說有守母足道.
又大笑又大哭.
或者要講謊言.
又要行神蹟.
可能有些會的.
但是這個聖靈充滿最重要的.
就是將我們裡面生命的雜質.
都倒出來.
我們過往的驕傲.
我們過往的腐敗.
過往的那些頹喪.
全都扔出來了.
然後接著又聖靈充滿在我們生命的裡面.
現在每個人信了耶穌的人.
生命裡面都有聖靈的了.
不過我們有多少空間.

$^{801}$又去掌管呢?.
當聖靈帶領著我們的生命的時候.
我們的生命變得有色彩.
變得有力量.
但是這個開始的時候.
我們最需要的.
就是這個和其他的信徒.
一定要行切的.
同心合意地來禱告.
只要我們肯這樣做的時候.
我們就可以經歷到.
上帝的能力.
淋到在我們當中.
更新我們的生命.
也都可以看到我們全部的教會.
我們盼望我們教會.
不單止120人是禱告的勇士.
我盼望我們教會有600多700人的時候.
我們都能夠成為了.
這個天國的禱告勇士.
自己得到更生.
教會得到更生.
上帝的國就能夠興旺.
將榮耀歸給上帝.
我們唱一首詩歌.
求主賜復興的火.
求主賜復興之火.
求主賜眩暈火焰.
燃燒教會能夠無窮力量.
在這黑暗時代.
臨在你光輝難解冷漠.
重點在你的心.
求主賜復興之火.
求主賜眩暈火焰.
燃燒教會能夠無窮力量.
在這黑暗時代.
臨在你光輝難解冷漠.
重點在你的心.
求主賜復興之火.
我們現在每一個人都安靜.

$^{841}$自己去想一想.
你願不願意為主的國度擺上禱告.
你願不願意成為在主的國度裡面.
其中一個祈禱的勇士呢.
你不要想很多的艱難.
就是跟上帝說你願不願意.
你要想一想.
你願不願意.
願不願意被聖靈充滿.
將生命裡面的雜質都等了它.
你要為主的國度成為不一樣的生命.
你要繼續被主耶穌基督你去使用呢.
如果你願意的時候.
我們再唱這一首詩歌.
你就有這麼大聲.
只有這麼大聲來唱這一首詩歌.
謝謝.
求主賜復興之火.
求主賜聖靈可憐.
燃燒也會能夠無窮力量.
在這黑暗世代能作你光輝.
容佳難忘.
重點愛你的心.
求主賜下國興之火.
你們剛才唱的聲量.
耶穌源源稱道.
你不用怕.
你也不用想怎麼做.
耶穌會引導你們怎麼做.
待會奉獻的時候.
我們繼續唱這一首詩歌.
謝謝.
好.
\newpage



\section{馬太福音 5:13-15}
\label{sec:vr_pYc_qVuc}
\textbf{2024-06-23 早堂《好ye 介紹返》馬太福音 5\_13-15  講員:陳韋安牧師}
\newline
\newline
連結: \href{https://youtube.com/watch?v=vr_pYc-qVuc}{\texttt{https://youtube.com/watch?v=vr\_pYc-qVuc}} ~~~~ 語音日期: 2024-06-25
\newline
\newline
\hyperref[sec:aAKtjwj8FbQ]{\small{< < < PREV SERMON < < <}}
~
\hyperref[sec:index_chronic]{\small{[返順時目]}}
~
\hyperref[sec:index_scriptual]{\small{[返順卷目]}}
~
\hyperref[sec:4L9GLJKp4e4]{\small{> > > NEXT SERMON > > >}}
\newline
\newline
馬太福音 5:13-15
\newline
\begin{longtable}{cl}
\hline
\hline
章節 & 經文 (和合本修訂版)\\
\hline
5:13 & \begin{tabularx}{0.7\textwidth}{X} 「你們是地上的鹽。鹽若失了味，怎能叫它再鹹呢？它不再有用，只好被丟在外面，任人踐踏。 \end{tabularx} \\ \\ \relax
5:14 & \begin{tabularx}{0.7\textwidth}{X} 你們是世上的光。城造在山上是不能隱藏的。 \end{tabularx} \\ \\ \relax
5:15 & \begin{tabularx}{0.7\textwidth}{X} 人點燈，不放在斗底下，而是放在燈臺上，就照亮一家的人。 \end{tabularx} \\ \\
[1ex]
\hline
\hline
\end{longtable}
$^{1}$丁子母早安.
我們一起祈禱吧.
天父我們求你的聖靈在我們當中.
正如我們剛才所認信的.
我們所祈禱的.
我們所唱的.
我們再一次懇求你的聖靈在當中.
去塗鑿我們的生命.
去轉化我們的生命.
更新我們的教會.
能夠將一些.
你古舊的十家好消息.
我們能夠可以再一次認真.
思考,看待.
將你的福音.
能夠再一次認真看待.
從而去宣揚和見證.
求神你幫助我們.
逢轉名詞.
阿們.
好,今天我們的講題就叫做.
好東西介紹.
我很少用廣東話講題.
但今天真的有些好東西.
跟大家介紹.
大家都知道.
福音其實就是.
好消息.
這個字其實就是.
無論在舊約.
或者在新約裡面都出現.
Evangelion這個字.
在余文裡面.
這個字其實不是一個很宗教的字眼.
大家都知道.
本身在舊約.
當素羅王死了的時候.
大衛就.
有兩個人就.
將素羅王這件事情.

$^{41}$笑嘻嘻地跟大衛說.
大衛就說.
你竟然將這個好消息.
笑著說.
我就要殺了你.
這個所謂的Good News.
在舊約裡面都出現了.
在羅馬政權的時候.
當時候這個好消息.
都不是解造耶穌基督的福音.
而是當時候.
羅馬帝國的皇帝.
來一個政治文宣的意思.
一個皇帝普遍天下.
來做一些公告.
都是叫做Evangelion.
好消息.
直到我們的聖經.
保羅.
我們的一些教會初期的信徒.
就將這個好消息.
來稱之為耶穌基督.
上帝的兒子的降生.
和受難的拯救.
所以當我們說到好消息的時候.
其實本身這個福音.
本身就是一件很好的事情.
今天我們一起來.
再一次來思考這個好消息.
這個好東西.
怎樣和我們的關係.
今日第一經文就是.
登山補佛一段經文.
你們事實上都炎.
炎藥失了味.
怎能叫它再咸呢.
以後無用不過.
雕在外邊.
別人踐踏了.
你們事實上的光.

$^{81}$成造在山上.
是不能隱藏的.
人點燈不放在斗底下.
是放在燈台上.
就照亮一家的人.
你們的光也當這樣照在人前.
叫他們看見你們的好行為.
便將榮耀歸給你們.
在天上的父.
一段大家很熟經文.
今日我們就不說.
大家都認識的東西.
那些你們是炎是光.
要見證的東西.
大家都知道的.
大家要做就行了.
今日和大家說什麼呢.
今日和大家說一段.
當我發現留意一個很重要的字.
就是當中的好字.
好字的希臘文是什麼呢.
今日我教大家.
一個很簡單的字.
好字就是Kalos.
多一次好不好.
Kalos.
這個就是希臘文裡的好字.
以前我有個同學叫做Kalos.
就是比較好.
如果你細心讀聖經的話.
當你比較三卷的婦女福音.
就是馬太,馬可和路加的時候.
你會發現在炎和光這段經文裡.
馬可和路加福音.
是特別有一個形容.
是馬太福音裡沒有的.
就是有關那個好字眼.
馬可福音.
九章五十節和路加福音十四章三十四節.
都特別提到一個經文.

$^{121}$就是炎本是好的.
若失了味.
可用甚麼叫他再喊呢.
路加和馬可都特別提醒和提到.
炎本是好的一句說話.
你會發覺這個很炎.
這個描述.
是燈神寶訓裡沒有寫出來的.
這個好字.
炎本是好的.
今天我們來到思考的課題.
雖然和本的聖經翻譯做炎本是好的.
但其實原文裡有一個很簡單的字.
就是Kaluntun Halas.
很簡單就是好的炎.
Salt is good.
這樣的意思就是一個.
馬可和路加不單單談及門徒系世界上的炎.
更加給這個炎本質上的評價.
就是炎是好的.
Salt is good.
好的炎.
甚麼叫好炎呢.
世上有沒有不好的炎呢.
有沒有炎比其他炎更加好呢.
說一下我女兒.
前段時間我女兒就跟學校去了日本的學校交流團.
現在才開始流行.
去交流團.
她不是去彭福公園.
我們小時候就去彭福公園.
她就去沖繩的交流團.
七千多元也不是便宜.
不過是一個很好的經驗.
讓她能夠自己照顧自己.
自己收拾東西.
去旅行.
住酒店.
自己和人相處.
去日本就給一點日本錢.

$^{161}$去花一下.
大概也給了一萬日圓.
也不是很多.
幾百港元.
去買零食.
買紀念品.
買東西.
當然一去到就馬上買小吃.
扭蛋.
買一些小玩具.
不過她也很懂事.
她也買了一些手信給我們.
她買了一包鹽回來.
因為在沖繩.
買一包日本鹽.
買一包鹽.
沒有什麼包裝.
只買一包鹽回來.
不過不倫不說.
作為父親.
收到女兒的手信.
也很感動.
買一包鹽回來送給我.
基本上我一段時間裡.
間中也會拿一包鹽出來.
自己用一下.
吃飯的時候.
就灑兩下鹽.
就變成日本食物.
整頓飯吃都不同.
變成日本定食.
加鹽就變成海鹽雪糕.
一包鹽.
一包很好的鹽.
日本沖繩鹽.
很珍貴的.
女兒送給我.
不過其實你知道.
其實一包鹽.
說真的.

$^{201}$你跟維康白雞買鹽.
應該沒什麼分別.
起碼科學構成是一樣的.
起碼成分是一樣的.
不會有什麼特別.
不會有什麼紫薯黑糖.
草莓味那些.
有的話都是加上去的.
都不是鹽本身有的東西.
所以如果你暗地裡.
把鹽換了.
換了白雞買.
我都分不清.
我很開心地以為.
日本那包鹽.
明明就是會是.
所以鹽就是鹽.
鹽本來就是好的.
所以和本的翻譯很好.
特別強調.
加了個「本」字.
鹽本是好的.
鹽本來是好的.
它是與生俱來.
就是好的.
鹽就是鹽.
但是鹽.
只要是鹽就是好的.
沒有壞的鹽.
這個就是鹽本是好的的意思.
所以鹽本是好的.
和今天天氣很好.
是兩個不同意思的說話.
是不是.
前者是一個.
本質上的評價.
後者是一個可以變的狀態.
今天天氣好.
上星期天氣就不好.
你知道是這樣.

$^{241}$鹽本是好的.
這是一個本質上的問題.
正如創世記所說.
神漢萬物是好的.
這是一個本質上的問題.
它是一個根本上好的問題.
所以就明白整個問題的核心.
明明是一些好的東西.
一些好的東西.
怎會無端端變得不好呢.
這正正就是耶穌在經文裡面.
向門徒要展示的問題.
有關鹽和光的最關鍵的問題.
你們是世上的鹽.
若失了味怎能叫它再含呢.
你們是世上的光.
神造在山上是不能隱藏的.
耶穌列出了一個假設性的狀況.
鹽如果失去了味道.
失去了味道的鹽.
這是一個如果.
放在豆底下的光.
耶穌說了一個假設性的可能性.
一個不可能的狀態.
一個不太可能的可能.
一些本來是很好的東西.
怎會變得不好呢.
這正正就是今天要問的.
這是一個甚麼狀況.
一個甚麼情況呢.
要回答這些問題.
首先要問.
好的東西是一回事.
好的本質在哪裡.
我嘗試用一個例子去講解.
我預備了一條英文題目給大家下做.
下一張看看.
在世界大都應該看到.
一條小學的英文題目.
I am a punctual person.

$^{281}$然後有一個空格.
I never hand in my homework late.
我是一個準時的人.
我從來都不會遲加功課.
中間應該加一個甚麼連接詞呢.
so 還是 because.
誰覺得是so的舉手.
其實兩種都很難過.
如果你說是so的話.
是甚麼意思呢.
我是一個準時的人.
所以我不會遲加功課.
如果是because是甚麼意思呢.
我是一個準時的人.
因為我不會遲加功課.
所以這個不是有答問題.
而是一個哲學問題.
一個準時的人.
是因為他不遲加功課.
所以是一個準時的人.
還是相反.
他是一個準時的人.
所以他不會遲加功課呢.
這是一個非常哲學的問題.
兩種都有道理的.
視乎你怎樣去定義.
你用一個行為去定義一個人.
還是用一個人的本質去定義一個人的行為.
當然這個不是講哲學.
如果你看聖經的話.
聖經似乎是贊同哪個當地人.
聖經似乎是贊同前者多一點.
有經文根據.
最基本的經文根據.
好樹就結好果子.
因為你是好樹.
所以你會結好果子.
壞樹是會結壞果子.
好樹不能結壞果子.
壞樹不能結好果子.

$^{321}$所以你想如果你是一個好樹的話.
你自然就會結好果子.
所以重點不是在乎於你怎樣結果子.
而是你是一個好樹還是壞樹.
所以這個似乎是聖經對於好的邏輯.
所以認為一些本來是好的東西.
它不會隨著一些其他的因素改變而改變.
一件好事,一個好言是會好到底的.
按照這個好的邏輯的話.
我們回到問題問.
明明一些好的東西.
怎會無端端不好呢.
或者問這麼好的東西.
為何沒有人覺得好.
為何沒有人覺得好去分享.
我想說什麼呢.
福音.
Evangelion.
Good News.
好消息,好東西.
這個是上帝在耶穌基督裡面最大的好事.
我們的信仰的根本和我們的根基.
我們的信仰因為這個好消息而起.
教會都為了這些好消息而傳.
我們都因為這些好消息而去.
大家來到聚會去認識.
身為基督徒,身為教會.
我們基督教信仰的開始.
就是由一班人去宣告世上的好消息.
不斷地去傳揚.
傳到今天.
就是這樣.
我們就是這樣而去存在.
不過不知道大家有沒有發現.
今天可能大家.
或者香港教會裡面.
我們對於這個好消息.
對於福音的傳講.
又似乎沒有這麼好.
不知道大家有沒有發現.

$^{361}$大家對於好消息.
這個傳福音是知道是好的.
但又沒有怎麼做.
不知道對你是不是這樣.
起碼今天很多人都比較.
不是太過覺得傳福音是一件很重要的事情.
特別一些年輕人.
很明白的.
因為這幾年裡面.
太多事情要關懷.
或者我們太多太多.
我們好像很多事情.
教會要顧.
但是你知道的.
可能我們未必會是.
很多事情要處理.
香港的問題.
教會裡面很多不同的事情.
我們似乎將這件好事.
有少少是丟下了.
我覺得教會傳福音.
都是成為了一件很適合.
但又有少少老土的事情.
好像做.
但又不是你會很熱心地去做.
我們將傳福音變成了一件.
很program的事情.
教會有傳道部.
教會有一個方法program.
這樣去做.
我們似乎純粹是加了差.
或者來到去做了.
fulfill了.
這樣就叫做了.
我們似乎沒有那個好事.
那種興奮的感覺.
回到經文裡面.
如果耶穌說.
「言本是好的」.
「言」就是言.

$^{401}$「光」就是光.
福音就是作為上帝.
給世人最好的消息的時候.
福音是天父上帝.
在基督律裡面.
為世人做過最好的事情.
祂叫世人發現這份美好.
都叫世人因著這份美好福音.
能夠叫生命變得美好的時候.
我們真的要好好地.
重新去思考.
我們所擁有的那份福音.
福音是教會一個卑微壓棄裡面.
一個最美好最美好的寶貝.
我們錦繡堂也是一樣.
錦繡堂最美麗的.
未必.
應該說不是.
肯定不是這個教堂.
也不是十年前.
或者二十年前.
那個美好的時代.
而是我們仍然一直保留的.
這份福音的信息.
祂和我們開始.
七零年.
八十年代開始到現在.
一直都擁有的.
最美好的好消息.
上帝去呼召我們.
在這裡建立教會.
正正是為了這個原因.
我相信一件好事.
好好地自自然然.
精神上.
就會變得美好的.
相信一件好事.
如果是好的話.
好到一個地步.
你會忍不住跟別人說.

$^{441}$正所謂好事不怕跟別人說.
這裡我們最見證耶穌.
最基本最基本的方式.
就好像聖律便井旁的婦人.
那個被格拉森.
被鬼附的人.
當我們認識主之後.
他們自自然然跟別人分享.
沒有人叫的.
耶穌都沒有命令他們.
他們自自然然會跟身邊的人說.
同樣道理.
基本的.
看過一套好看的戲.
九龍城寨.
好看就叫別人去看.
吃一間好的餐廳.
好吃.
你自然就會帶朋友一起去吃.
認識了耶穌.
回到一間好的教會.
認識了福音.
自自然然就會叫別人去認識.
除非你覺得這件事不是很好.
如果你沒有跟別人說的時候.
只有一個可能性.
就是你覺得福音都不是很好.
其實都是一些不好的東西.
讓我承認吧.
不要介紹給你.
你不會這樣.
福音其實是一些好的東西.
你應該會跟別人分享.
所以一件好事.
如果它本質上是好的話.
讓你自自然然.
不費力氣.
很開心的.
你就會跟別人分享出來.
這個正正就是福音本身的好處.

$^{481}$你不需要誇大.
不需要現過其實.
你只需要將你覺得信仰裡的好.
跟別人分享.
你只需要將裡面的好.
好好地將別人帶到你面前的時候.
這個就是做了他本身的事.
所以今天的題目叫做「好ye 介紹返」.
一方面是介紹好東西.
同時也是介紹他回來.
既然這裡是一個好的地方的話.
你每個星期都願意回來的時候.
你其實都可以很自然地帶他回來.
經文最後這樣說.
「你們的光也當這樣照在人前.
叫他們看見你們的好行為.
便將榮耀歸咎於你們上天之府」.
這個所謂的好行為.
不是一些很特別,很厲害的好行為.
而是一些很基本的.
你感覺上的好.
你自然地將這個好分享出來.
不需要強行來的.
不需要特意去做作的.
很自然地將這份好介紹給人.
嚴格來說我們可以說叫做.
「好ye 介紹返」.
有好東西便介紹給人.
無論是格拉森比鬼的婦人.
井旁的婦人,安德烈等等.
我們就說將這份好東西去介紹給人.
帶他回來這裡.
你可以說這個叫做「友誼報道」.
我們教會的策略叫做「友誼報道」.
其實我不是太喜歡這個名字.
所謂「友誼報道」其實不是一個報道策略.
這是我們基督徒最基本的行為.
這個不是一個報道方法.
而是一個最基本的生活方式.
和生命的行為和見證.

$^{521}$所謂「友誼報道」其實就是.
在你生命裡面,在你生活裡面.
很casual的.
你和福音之間一個最authentic的關係.
一個最真實的關係.
因為你和你的朋友和你的親人.
在一個最真實的關係的時候.
你的信仰都是真的話.
你將兩件事拉在一起.
這樣就叫做「友誼報道」.
既然你是很真心相信福音是好的話.
你是真心和你的朋友和你的家人和你的同事.
在一個真實關係的時候.
你只不過將你生活裡面的兩件事拉在一起.
這個就是你的「友誼報道」要做的事情.
不需要什麼報道技巧.
不需要去學會什麼複雜的理論,神學.
這個不是不重要.
而是最基本的.
你就是帶人回來這裡.
將耶穌來到他大眾面前.
那些五分鐘,十分鐘的報道技巧.
可以學到.
但最基本的就是你自己將一些好東西.
帶到人回來見到一些好東西.
多謝弟兄姊妹關心.
上個月「白白安息禮拜」.
我爸爸其實也有回教會.
我自己不是家裡信耶穌的家庭.
我是第一個家裡信耶穌.
我是第一個回來錦繡堂的人.
1999年的時間.
我自己是第一個回來教會的人.
當時開始就有好和.
這裡錦繡堂是一個很好的地方.
信耶穌很好.
開始就自然邀請我兩個妹妹回教會.
當時我大約是中六.
那時有一次我生日.
就跟妹妹說.

$^{561}$今年生日你不用送東西給我了.
你就跟我回來教會團契吧.
她們兩個就說這麼好.
不用送東西就跟我回來吧.
很自然她就回來我們禮拜五晚的.
當時的約留單團.
就這樣她就回來這裡.
就自己認識了一班朋友.
發現了耶穌的好處.
家裡就馬上頓時變成了三比二的信主比例.
三個基督徒.
兩個讚美三使美信.
當我兩個妹妹信了主之後.
就跟我妹妹說.
我們既然做了基督徒.
就要好好地去見證上帝.
我們在家裡就互相提醒要做好一點.
以前中學的時候.
回家就不幫忙摺衣服.
晾衣服.
家裡不小心就馬上做好一點.
不要那麼隨便.
就馬上把家裡收拾好一點.
真的爸爸媽媽會看到那個不同.
接著就到了我爸爸.
那時候我就帶我爸爸回來教會.
其實他那時候也拜了很多年觀音.
有一次他就跟我回來教會.
參加我們教會的舞蹈班.
那時候我們有舞蹈班.
長期都是進駐拜祭的.
因為人家是講福音的.
那時候有一次報道會.
他就信了耶穌.
那時候他很厲害.
回家就馬上把觀音的偶像.
吳室牧師幫忙摺好.
馬上就變成4比1.
當4比1就很容易.
事情就很簡單.

$^{601}$剩下一個的時候就會是.
那時候我媽媽因為她要上班.
星期天要上班.
就回不了教會.
有一次她就.
可能她的身體不好.
就沒有上班.
有一次她就做夢.
突然就哭醒了.
哭得很厲害.
就突然間做夢.
她自己去一個很黑的地方.
很孤單.
她說我們其他四個都在天堂裡.
她自己一個人在那裡.
她就跟我回教會.
回教會第一次就唱詩歌.
就很感動地在那裡哭.
記住我這個不是一個.
帶家人信主的見證.
我是一個帶家人回教會的見證.
只不過是.
只不過是一個好東西.
就自然帶她回來這裡.
我沒有做過任何報道的工作.
我只不過是將一些好東西.
介紹給人.
如果是好的話.
我相信她回來.
她聽,她參與.
自己來面對.
她自然會覺得是好事.
問題是我們有沒有.
將這份好東西介紹給人.
真人.
今天很多的教會.
特別是一些大教會.
或者是一些.
很好的,很漂亮的教會.
我們都未必是.

$^{641}$做得最好的這件事.
相反.
我認識一個學生.
在上水那裡.
開了一間新的教會.
幾個人的教會.
因為他每個禮拜每天.
正正做一些開荒工作.
跟一些上水的年輕人.
一些新移民去傳福音.
這樣來開設一間新的教會.
很艱難.
今天仍然維持不了.
但這件事是一個非常可恥可敬的事.
福音.
或者教會.
正正是為這件事而存在.
今天我們有沒有很好的來到.
將這件事.
放回我們的.
那個守衛裡面.
不需要做到很誇張.
不需要有任何的.
很硬來的事情.
只不過將這份好事.
放回它應該好的裡面.
以前我回錦首堂.
其實每個禮拜都會有新朋友回來.
因為當時.
每個人都很願意帶新朋友回來.
我想這也是我們如果是.
想教會有改變的時候.
或者不這樣說.
做回我們最基本做的事.
我們將這份福音的.
好事去介紹給人.
不要稱之為.
當然可以稱之為福音的責任.
福音的債.
但我覺得快樂一點開心一點.

$^{681}$將這份好事介紹給人.
是一件很輕鬆的事情.
想想有什麼人可以帶回教會.
不回不要緊.
回一次不回都不要緊.
因為你覺得是好.
你自薦而.
繼續會的一齊去享受這件好事情.
我們一起祈禱.
因為你將這份好消息.
去給我們.
我們能夠知道福音的時候.
我們確實是一生受委.
更加是永生受委.
求主你幫助我們.
讓我們能夠去享受這份好處.
願意和人分享.
求神你幫助我們.
讓我們去見到你自己的座位.
願意我們見到你自己的好消息.
福音能夠繼續傳揚.
在我們生命裡面.
在我們生活裡面.
和我們一起去生活.
讓我們能夠常常帶著這份好消息.
和人相處.
逢主明求.
阿們.
\newpage



\section{帖撒羅尼迦前書}
\label{sec:4L9GLJKp4e4}
\textbf{2024-06-27 周四晚查經班《帖撒羅尼迦前書(10)》 講員:曾裔貴牧師}
\newline
\newline
連結: \href{https://youtube.com/watch?v=4L9GLJKp4e4}{\texttt{https://youtube.com/watch?v=4L9GLJKp4e4}} ~~~~ 語音日期: 2024-06-29
\newline
\newline
\hyperref[sec:vr_pYc_qVuc]{\small{< < < PREV SERMON < < <}}
~
\hyperref[sec:index_chronic]{\small{[返順時目]}}
~
\hyperref[sec:index_scriptual]{\small{[返順卷目]}}
~
\hyperref[sec:uJ_qz3eVEf0]{\small{> > > NEXT SERMON > > >}}
\newline
\newline
$^{1}$我們一起先祈禱.
多謝天父再一次恢復查經班的時間.
盼望藉著《帖撒羅尼迦前後書》.
對於仕途保羅在早期第二次宣教旅程的報導.
以及對於教會的模樣.
都讓我們能夠藉著一些大的方向原則.
找到一些重要的線索.
尤其是這本書很多時候講主的再來.
而且主的再來要我們在有限的人生做準備.
都幫助我們能夠明白,掌握.
亦都讓我們在前面有機會去看.
末世論這個神學.
有不同的派別.
對於主再來的想法都讓我們從其中大概知道.
甚麼是災前,災中,災後的被提.
亦都讓我們看到七年大災難.
與教會的關係等等.
願神都帶領我們.
雖然我們那個時代已經不同了.
但是那個模樣的大方向原則.
關於耶穌基督的再來.
是每一個世代的基督徒都需要這樣來警醒等候的.
願神都幫助我們.
我們再一次在晚上來等候.
盼望不單止是延長.
而是成為了錄音.
可以祝福不同的教會弟兄姊妹.
不同的地區的弟兄姊妹.
都可以在家中慢慢靜心聆聽.
去得到神話語的教導.
幫助我們能夠有系統地按照正義去分解.
按照著上下文.
保羅的一路講解.
我們去尋索作者.
他講這些內容背後的意圖.
背後的原因等等.
我們仰望等候禱告.
奉耶穌基督的名.
阿們.
今晚我們有機會總結一下.

$^{41}$一到第五章前書給我們什麼教訓.
我們可以在這次停了一次.
再用多一點時間.
先看上半段的查經.
再看十二到第二十八節.
這相當明顯就是教會生活.
要學習的種種原則.
在一個群體的生活.
我們需要有群體的禮儀.
那個態度.
應該要注意要具備的事情.
十二節開始.
五章十二節.
我們講完這一段之後.
我們做一個總結.
對於前書有什麼重點.
我們甚至可以深入研究.
弟兄們.
我們勸你們敬重那在你們中間奴婦的人.
就是在主裡面治理你們.
勸戒你們的.
相當明顯就是傳道者.
使徒不會停留在一個地方.
比如耶路撒冷的使徒.
他們就會主力成為了.
耶路撒冷教會的長老.
因為他們必須都是猶太人.
神都安排他們.
所以比如保羅,巴拿巴那些.
他們都是在外面宣教.
他們就不是耶路撒冷的長老.
彼得,約翰,雅各等等.
這些就是耶路撒冷的長老.
說的雅各就不是彼得和約翰的哥哥.
而是耶穌的弟弟.
肉身的弟弟.
那個雅各.
或者是寫雅各書的作者雅各.
這些就是在耶路撒冷的長老.
就是這些使徒.

$^{81}$在外邦保羅建立的初期教會.
這本書本身的動機.
寫作的原因.
就是保羅很急速地.
四個星期就要.
因為一些很大的逼迫.
他要離開.
輾轉去到比利亞.
然後就去到希臘的雅典.
就在雅典有一些相對穩定的時間.
寫信給帖撒羅尼加教會.
原因就是.
保羅留下提摩太在帖撒羅尼加城.
他繼續這樣去牧羊.
後來提摩太就在帖撒羅尼加.
下去會合保羅.
在這個雅典.
保羅聽了提摩太的匯報.
弟兄姊妹仍然不會因為謠言.
逼迫退後.
信心失去.
反而很堅固.
所以保羅就因為這樣的緣故.
就非常喜樂.
所以就寫了這封牧羊的信.
透過提摩太.
再帶回去給帖城的弟兄姊妹.
讓他們的信心得到堅固.
也有一個這樣的大家的溝通.
一個大家的對話.
因為收到一個好消息.
收到一個回覆.
而且是需要人用腳去走.
或者坐船.
但是也是需要好幾個星期.
才收到書信.
所以和我們每秒鐘收到WhatsApp.
氣氛就完全不同了.
他們會憤外,感覺雀躍.
因為離開的保羅.

$^{121}$原來不是見異思遷.
見到困難就走了.
反過來就是.
弟兄姊妹的生命表現.
又令到保羅覺得原來.
他所做的不是白費.
不會因為很短暫就失去.
就是不再信.
有時候我們去都會經歷的.
報道會有十個人缺席.
然後逐張紙打電話給他.
他說沒空了.
不來初訓班了.
有時候就會慢慢的.
其實不是水分.
是那個落差.
他有時候還沒準備好.
他那一刻是很有心想信.
但是要他承諾一個.
就是十個星期的初訓班.
再上浸制班.
他可能會說我都要想一下.
所以這個也是一個.
我們要留意.
所以保羅看到.
貼守羅尼加城的弟兄姊妹.
最擔心就是他們因為苦難.
逼得而退後.
所以他們現在能夠表現出.
這一份的信心和堅固.
就令到保羅真的感覺到.
那個牧羊的果子.
還有那個信心更加提升.
一方面就是能夠令到外邦的教會.
能夠不在猶太人影響下.
可以慢慢的穩固.
這個是對保羅來說是很驚訝的.
因為他自己是一個.
原汁原味的一個猶太人.
一個拉比.

$^{161}$所以他見到外邦人信主.
憤外感到是神的差遣.
神的命令他成為外邦人的使徒.
就真的憤外的有那個果效.
所以保羅的感覺就會強烈很多.
因為猶太人不信.
他認為是很合理.
但是外邦人.
他們一直都沒有這些包袱.
而能夠離開偶像歸向真神.
已經是很難能可貴.
所以這個也是保羅.
他對於那個牧羊上.
有的那份奇妙的愛.
我們看看這裡.
第十三節說.
「又因他們所作的功.
用愛心格外尊重他們」.
對於當時的貼城的傳道者.
又或者一些信徒領袖.
其實大家都是四個星期信主.
何來有那個屬靈的位分.
或者是累積的時間.
其實可能都是比較有心的.
未必是每個人都信了很多年的弟兄姊妹.
但是他們都已經是這樣.
為著福音來擺上.
來服侍保羅的教會.
所以保羅勸弟兄姊妹.
一定要對這群牧羊的人.
他們那份牧羊的恩情.
來用愛心.
用弟兄姊妹的愛心.
格外尊重這些傳道者.
下面就說「你們也要彼此和睦」.
這個也是教會必須有的一個重要的形態.
教會的存在應該是帶來平安.
不是帶來分裂.
不是帶來為了一些很無謂的事情.
來彼此去爭執.

$^{201}$要和睦.
「我們又勸弟兄們.
要警戒不守規矩的人.
勉勵灰心的人.
輔助軟弱的人.
也要向眾人忍耐」.
教會本身的存在就是眾人.
就是不同類的眾人.
不同的民族的眾人.
不同的背景的眾人.
不同的大家性格的眾人.
就算是唐人.
就算是廣東人.
都已經有不同的.
你是順德 我是墮羅.
都已經有不同的飲食習慣.
如果我們再闊一點.
你去到一些國際一點的教會.
譬如我這次去多倫多.
有些是英文 有些是國語.
有些是其他的語言.
是很正常 是很普遍的.
所以教會是需要彼此學習接納.
所以如果有些是不守規矩的.
就要警戒.
如果是有些灰心的就要勉勵.
軟弱的就要輔助.
所以是要因人而異的模樣.
是要敏銳每個人的情況.
那個的光景.
等於你大病初愈.
馬上跟你去腳踏.
比較辛苦 不是不行.
早上就要去吃火鍋.
就很辛苦.
剛剛病得很辛苦.
你還要去吃燒烤.
那些咳嗽都還沒完.
更加令到他更加慘.
但是如果他很健康.

$^{241}$身體情況都很好.
偶然吃一餐真的很豐富的.
就會感覺很開心 很滿足.
所以其實每個人.
都真的有不同的境況.
不同的際遇.
不同的那一輪的心情.
所以傳道者就要敏銳.
所以這裡說警戒.
勸勉 輔助.
還有要向眾人忍耐.
無論是牧者.
或者無論是信徒領袖.
都要在一些可能.
政策上面你推行.
都要讓人慢慢去接納.
不要一下子.
好像一個威權式的管治.
或者好像一刀切.
什麼都不理.
就會比較大的一個問題.
譬如你說有些很強的女強人.
她來做領袖.
有時就很難.
文姐有時都提到.
那些上司如果是女性.
特別要求高.
本來是好的.
但是有時就是.
哇 七點都不下班.
原來沒有拖拍.
那些下屬就慘了.
哇 我趕著約了人放工.
我下班五點半.
接著你五點四給她工作.
有時有很多這樣的抱怨.
有時真的反映了.
我們忘記了.
有些少少的對人的體諒和勉勵.
所以很特別.

$^{281}$它是向眾人忍耐.
這真是反映了我們在教會裡面.
真的需要用包容和愛.
先不要用一種判斷.
判斷帶來的.
如果我們負面就變成了論斷.
判斷還是正面的.
判斷未加你自己很深的.
個人限制的評價.
那還是一個判斷.
但是如果加了個人的一種定性.
可能那件事不是這樣的.
但是你定性了.
那就變成了論斷.
這就變成很可惜.
這就告訴我們要謹慎.
無論是誰都不可以以惡報惡.
這又是另一個角度.
就是知道一些事情.
甚至你知道有人用一些惡言.
毀謗不存在的事實來說你.
都要彼此化解那個誤解.
那個誤會.
所以這裡很明顯是在說教會生活.
因為如果你出來教會.
其實你不需要學這些東西.
你一個人生活學什麼呢?.
是吧?.
星期你都大覺睡.
11點半的鬧鐘響.
你才去排隊喝茶.
沒有什麼需要學的.
學的就是爭位.
被人搶了你的號碼.
你生氣的意思而已.
沒有什麼需要你去輔助.
沒有什麼需要你去勉勵.
所以如果我們的心.
如果心是一個群體.
那就恭喜你了.

$^{321}$你開始有做傳道人的呼召.
那個感動.
因為其實任何一個傳道人.
就是心是神家裡面的羊.
對那個教會會上心.
這個就是牧羊人.
跟一個故宮不同的地方.
所以有生之年.
傳道人都是擔驚受怕的這樣牧養的.
都是在那個牧養上面.
總是會有很多東西掛心.
所以很佩服那些.
真的可以又掛心又大覺睡的傳道人.
就是這樣真的很交托.
他很平衡.
但是大多數其實.
很多牧者都是睡不夠.
還有一些甚至乎有情緒的問題.
因為太多東西.
甚至乎有很多的壓力等等.
所以要謹慎.
無論是誰.
關乎到一些的真側.
就無所謂誰威權一點.
誰有高地位一點.
而是大家都應該要平情而論.
不要訴諸於權威.
不要訴諸於一些的結黨.
而是將那件事用在主的救恩的愛裡面.
來到去.
就是怎樣呢?.
不要結惡.
就解決了那個問題.
就是主耶穌講的寬恕人七十個七次.
這個一定要一生去學.
因為哪個人可以做到呢?.
今早九點半.
因為另外一個同事他翻眼.
我就站在門口等了半個小時.
等那些運輸送貨來.

$^{361}$這次這個運輸就很嘈吵.
跟上一次搬上去很重的那些傢俬.
大件的傢俬搬上去禮堂.
你想想要走幾段的轉彎的樓梯.
都沒有任何怨言.
這個運輸就跟那個導航.
就停錯了地方.
停在那個所謂臨時街市那邊.
其實根本就在外面.
他就好像很生氣.
一直等都不見有人來接.
他自己就用那些摔車.
拖了六七件的那些傢俬包裹.
就一直走那條橫街就過來了.
我一直接電話一直等都沒有來.
後來通了電話打給他.
就以為問他去哪裡.
他為什麼還沒有到.
我等了他半個小時.
就在教會旁邊走廊的位置.
他就開始嘈了.
開始爆粗口了.
以為我責怪他.
為什麼這麼久都沒有來.
這個誤會了.
我就說真的很兇.
那你牧者啦 教會啦.
還有一來就是誤會而已.
無所謂我站在這裡.
我怎會怪責你呢.
我只是想有什麼可以幫忙.
他就不斷地嘈.
不停地罵我.
誤會就說不好意思.
千萬不要介意.
沒事 不是說你.
他就慢慢地把那些行李.
那些傢俬去到地鋪門口.
就放下.
然後給我一張單.

$^{401}$然後我幫他簽.
還有一件他又去拿.
不斷地跟他說對不起.
就是說其實彼此誤會而已.
其實我等了你半個小時.
絕對不是怪責你這麼遲都沒有來.
而是我想幫你去接貨.
他就氣消了.
那就沒事了.
一來他們這些勞動階層.
有時都會很容易會嘈.
所以如果你以嘈歸嘈.
那就又找交換了.
那就不好.
不是因為我是牧者.
而是根本很小事.
他都在拖那些貨來.
我一出來想他把貨拖近一點.
就不用他搬那麼遠.
他以為不可以泊車.
一來已經有私家車泊在那裡.
就變成有一個彼此的誤會.
如果我用粗口還給他的.
一來失見證.
二來只會製造更多不必要的矛盾.
所以無論是誰.
都不可以惡報惡.
這個都是很重要的.
或是彼此的相待.
我稍微揭開一下.
我看到他們不停停了.
我不知道是什麼原因.
團體的聊天.
我不知道是什麼原因.
是你們收得不好還是什麼原因.
我怕你們停了又出去.
沒事吧?聽到嗎?.
有沒有少少回應?.
淑英或者Jenny.
是不是聽到?沒問題.

$^{441}$給少少反應.
或者在群組.
我不管了,我繼續說.
是沒問題吧?.
聽到嗎?.
聽到吧?.
行,OK.
我們繼續.
或是彼此相待.
沒錯了.
你不要因為別人惡.
你就用惡相待.
要沉靜先.
要自己反思.
明顯是一個誤會.
大家以為去責怪對方.
他說你為什麼還不來接貨.
我就說你的車還沒到.
我一直以為他會開進來.
所以有這個彼此的誤會.
但如果我們在那些.
很容易就會用一些很激烈的語言.
所以那一刻就要真的依靠神.
沉靜下來.
不是我不懂說髒話.
而是不需要說.
有什麼需要說呢?.
很小事而已.
為什麼要搞到這麼大件事呢?.
所以我們不要讓這些誤會.
製造成為一些彼此相惡的地方.
譬如有時你看到報紙,新聞.
在地鐵打架.
只是為了坐的位置而已.
然後失去儀態.
還要去見警察.
其實最終都是彼此受害的.
所以我們要學習.
謹慎.
就是說不是這麼容易的.

$^{481}$有些事情.
每個人都有自己的死穴.
自己都要很小心.
無論是誰.
保羅說得很好.
都不可以惡報惡.
就算你自己都正確.
都要去彼此包容.
彼此接納.
或是彼此相待.
就是說你做初一.
我就做十五.
你要這樣對待我.
我就要這樣對待你.
這是一般人的水平.
我都說了.
我們是有屬天和屬地.
兩個身份的.
平衡的世界的.
我們在地上.
其實要學習活出天上的樣式.
要有主耶穌的品格.
或者代眾人也是.
所以代眾人.
代眾人.
就是說不只是教內.
所有人.
常要追求良善.
你要有那個kindness.
你要很kind.
不是一種陰聲細氣.
而是一種內心的仁慈.
追求良善.
就會帶來.
福音是一個見證.
不然的話.
原來你是教會來的.
你是誰.
你是誰負責人.
你又跟我說髒話.

$^{521}$那你就麻煩了.
你有時那一刻火來.
你就說髒話.
那一刻火來就是不夠謹慎.
那些事件其實很小事.
能否接到那十分鐘算得了甚麼.
又不是趕著甚麼.
就是收貨的問題.
還有他也有司機看車.
所以其實根本就很小事.
有時那些勞動階層.
他就因為習慣了這樣的語氣.
就很容易就會.
他會在一種激氣裡面.
他一來就分析到.
其實我不是這個意思.
我是等了你很久.
就是怕你不懂得找.
他以為沒有人去接.
他接著還說我導航錯.
其實不是.
所以就變成要追求良善.
相信不是我們自己本身就可以做到.
都需要真是有神的恩典.
在我們生命裡面.
我們才可以追求到良善.
你在一個冷氣房.
穿西裝.
你當然不會激氣.
你也不會燥.
人家這麼辛苦又熱.
三十多度.
一句話就可以立慶.
所以其實我們都要將心比己.
看到人家的情況困難.
十六節.
常常喜樂.
我們很多時候就抽了這段出來.
當然是金句.
當然是金句當中的金句.

$^{561}$要常常的喜樂.
不住的禱告.
凡事借因.
因為這是神在基督耶穌裡.
向你們所定的旨意.
仍然在教會生活裡面.
會使我們近神.
而且越來越像.
耶穌基督的洋式操作.
教會的生活.
真是不外乎我們要追求主的同在.
追求我們成為祂在地上的一個代言.
就是我們在禱告裡面.
可以見到神的工作.
因為我們為一些事情祈禱.
如果是對神的旨意.
神會透過我們的祈禱去成就.
所以這是最直接.
最近我們在教會生活裡面成長.
和明白神的旨意.
在我們的生命.
在教會裡面的神的安排等等.
所以其實是這幾樣東西.
在喜樂裡面.
不住的禱告.
去明白尋求神的旨意.
所以一個在平安裡面的基督徒.
內心是不住的禱告.
不住的禱告會帶來他內心有喜樂.
喜樂不是嘻嘻哈哈的大笑.
不是看完一些笑片的大笑.
而是內心裡面有一份平靜.
對永恆的認知而來的滿足.
這是聖經所說的喜樂.
一般人是不會有這種喜樂的.
如果他沒有聖靈沒有神.
沒有救恩.
更加談不上真正的喜樂.
一定會有笑.
神製造人類.

$^{601}$為何會不斷有人要拍笑片.
每年都有市場.
古今中外都有.
不會只是文藝片.
不會只是悲情片.
一定會有笑片.
原因是人需要喜樂.
但那只是令你歡笑.
不是令你內心的深處.
有一份因為永恆而帶來的平安.
和真正的喜樂.
所以教會生活離不開.
即是祈禱會.
尋求對神的認識.
譬如查經丹.
崇拜.
和參與福音的事奉.
這些都會幫助我們.
很貼近神的旨意.
接著又好像很零碎.
其實不是的.
都是一個套餐.
給信徒一個教會的治理.
來有規矩地.
有秩序地.
讓信徒能夠穩定地成長.
所以教會的設計.
不要消滅聖靈的感動.
不要藐視先知的講論.
但要凡事察驗.
善理的要持守.
各樣的惡事要禁戒不作.
願賜平安的神.
親自使你們傳言成聖.
又願你們的靈與魂.
與身子得無保守.
在我主耶穌基督降臨的時候.
完全無可指責.
即是在我們有限的今生.
每個人都有自己的大限.

$^{641}$神定了的.
不能讓你去.
希望哪個人幫你求多十年命.
沒有的 因為神定了.
不能去有一些法術.
符咒.
沒有的.
那些都是騙你的.
但是在我們每個人的今生.
我們要完全無可指責.
在主耶穌再回來之前.
再一次提到主再來之前.
原來那個等候.
真的等候.
不過不是用計算機去計算.
而是相反的計算.
在我們很有限要見主之前.
盡力去事奉.
阮妮她怎樣都在離開之前.
要參加我們洪恩堂的.
第一次聖經班.
我當然開心.
這就是告訴我們.
我們原來真的要為主再來準備.
所以這裡一併說的.
其實我們可以做到一些歸納總結.
那個教會.
跳索雷加教會.
保羅勉勵他們必須要有純正的道.
必須要有對港台的尊重的聆聽.
那個聆聽就不是一面倒.
追星.
跟哪個偶像.
哪個崇拜的講員.
而是這裡很清楚.
不要消滅聖靈的感動.
不要藐視先知的講論.
但是要凡事的去測驗.
就是說你要有一個很敏銳的想法.
來去聆聽.

$^{681}$就是說原來就算是名人.
他的講道都可能會錯的.
但是一個可能是神學生.
都可能講到很出色的道.
去感染你.
去有聖靈的感動.
令你明白了一些東西.
甚至是神透過他跟你講的.
所以這裡就不是哪個人.
而是那個道本身.
這個就是大方向原則.
教會的生活必須要有純正的道.
無論是哪個人講出來的.
我們都要善美的持守.
要凡事的測驗.
這個就是相信來查經班的弟兄姊妹.
都相當明白.
其實我們絕對不是.
跟隨哪個人的講章,講道.
而是我們尊重每一個牧者的教導.
他一定是有他的領袖去講.
他一定有他的限制.
無論他的口才.
無論他的演繹.
無論他的例子.
一定有他的不夠成熟.
未到爐火純青.
但是總有一些道.
你可以從他的生命裡面.
或者他的講道裡面去得到的教訓.
所以這裡就給我們一個很完整的教會生活.
首先就是必須要有純正的道.
你不要靜.
所以講道的人.
大家不用擔心.
我先靜一靜.
行了 靜一靜.
所以我們.
為什麼突然間又有聲音.
不是 首先是另外一個 不是你.

$^{721}$我們就要記得.
一個教會能夠建立.
必須要有穩定的神的道的供應.
每個星期都需要有神的話語的供應.
教會裡面因為有神的話語.
大家都知道我們是尊敬神的.
我們是有聖靈在我們中間感動的.
所以我們就應該在教會的人.
跟人之間的相處.
聚集的時間.
就要學習在十誡的救恩裡面的那份愛.
包容.
那個接納.
所以我們是建立一個和睦的.
要與眾人和睦的.
我們是要避免以惡報惡的.
我們是要盡量學習.
在主的愛裡面.
來成長.
因為你成長.
你就會開始看到原來人.
真的有不同的類別.
是有不同的那一輪的際遇.
所以每個人都有不同的.
可能這次是他需要你來安慰他.
扶持他.
可能下一次是輪到你.
需要被其他人來安慰.
所以睦者都要學習被安慰.
被鼓勵.
甚至乎是被警戒.
你要讓人知道.
你也是人.
你也是需要弟兄姊妹.
所以這個就是教會生活都很重要.
歸根結底.
就是教會.
因為大家能夠聚集的時間.
可能一個星期只有半天.
甚至乎有時是整天.

$^{761}$但是始終還有六天的時間.
都是自己的家庭生活,工作.
當時是一樣的.
所以我們能夠聚集的時間.
應該利用不同的配套.
譬如現在是高科技的時代.
我們已經有不同的群組.
不同的團契的群組.
不同的時工的群組等等.
都是一個彼此可以去.
牧羊安慰的一個平台.
但是人與人之間的見面.
就會真正可以消除很多的誤解.
因為文字沒有看到樣子.
有時都可能會誤會的.
面對面彼此的.
這樣去面談.
或者去分享.
肯定是可以帶來.
我們消除了大家的誤解.
消除了大家一些的隔膜.
就會帶來你可以認識到不同的人.
大家都是神兒女的身份.
而一起在主裡面學習.
一起在主裡面成長.
以及一起來事奉.
這是相當寶貴.
我們能夠不只是自己一個人.
在這個地上生活.
我們來到看第23節.
回到聖經最後一段.
保羅的一個竹園教會的弟兄姊妹.
在帖壽來家.
他們只不過是四個星期的信主.
就因為保羅要離開.
他們絕對不是無所適從.
而是他們慢慢都學習到.
在教會裡面的生活.
相信更重要的.
絕對不是提摩泰一個年輕傳道人.

$^{801}$因為當時提摩泰真的相當年輕.
在第二次宣教的旅程.
他都是在路斯德被發掘.
跟隨保羅去傳道.
接著輾轉就有馬其頓的異象.
輾轉就去了歐洲.
其實由土耳其坐船去歐洲.
輾轉在菲臘比.
接著在帖壽來家.
所以提摩泰自己.
其實都是出來傳道很短的日子.
所以似乎不是保羅.
或者不是提摩泰.
就可以令到一間如此初信.
又面對迫伯的教會.
來到去站穩.
所以保羅的擔心是有他的理由.
但是很重要的是.
有這位赤平安的神.
他不是默默的.
他一直看著教會.
所以願似平安的神.
親自使你們傳言的誠信.
又願你們的靈與魂與身子.
都無保守.
這裡為甚麼要分開靈,魂,身子,身體呢?.
是否真的有三元論呢?.
即是原來靈和魂是不同的.
是否因為這一句.
我們就要這樣去分開呢?.
靈就是靈的部份.
魂就是魂的部份.
身體就是身體的部份.
如果靈界就是我們內心的靈.
與聖靈的契合.
魂就好像是我們知識,感情的部份.
是否要這樣分割呢?.
當然透過這一節.
表面上就像是保羅自己分割了.
但是在其他的新約聖經.

$^{841}$很少會這樣分得這麼仔細.
這樣去分開三個部份的.
我們人類的結構.
其實靈魂加身體.
其實已經是一個整全的人.
靈魂牽涉的可能就是涉及到.
甚至你的知識,你的感情世界.
你的感性上.
理性和感性.
還有我們那個.
譬如我們的personality.
那個性格.
與生俱來那些.
譬如你是衝動型的.
文靜型的等等.
每個人都不同的.
有不同的性格.
也有我們身體的身體.
所以其實又不需要分割得這麼重要的.
靈魂加我們的身體.
保羅這裡到底有沒有特別的意思呢?.
看字面就不是很明顯.
不過重點就是.
這一位造物主.
這一位是會將平安賜給我們的神.
他本身就是平安之君.
平安的神.
所以我們接觸神.
第一件事令我們獲得的.
救恩帶給我們的.
就是內心與神相和.
和好.
會帶給我們內心有一份.
是你從來無論用錢.
或者用你自己認為的幸福.
都換取不到的平安.
一個信主的人立即會有的平安.
因為你與神相和了.
這個是差不多每個人都會經驗到的.
這份平安有些人就會用很大的喜樂.

$^{881}$有些人可能就是會流淚.
但是滿足喜樂的眼淚.
這樣來反映流露的.
就好像你在內心湧入出來的.
所以平安的神遇見一個.
沒有平安的罪人.
就會帶來救恩的生命改變.
所以原始平安的神.
親自使用你們傳言的誠信.
相信保羅將他下面那裡分開的.
靈魂,身體說的誠信.
有一點點的關聯.
因為神會透過救恩.
透過耶穌基督.
會使我們生命成長.
會使我們對價值觀.
對事物的看法.
會有一種嶄新的.
嶄新不是那樣東西變了.
是我們的心變了.
因為我們有永恆.
因為我們有永生.
我們會看那個人生的時限.
會拉長了.
很多人都說.
他真的會明白.
離開世界只是打開另一道門.
進入永恆.
沒錯.
肉身真是在一種很大的痛苦裡面離開.
但是如果真是嚴格.
我們如果看到靈魂.
如果有機器放在床邊.
看到那個靈魂.
那個靈魂真是打開另一道門.
去一個所謂的極樂世界.
當然不是那些佛教的極樂世界.
是說去到永恆.
所以如果有這個機器真是好了.
聖經就是這個機器.

$^{921}$告訴我們.
我們真是離開這個世界.
進入永恆的天家.
是不會再有眼淚.
不會再有痛苦.
最近真的帶了一個弟兄.
一個圍村很資深的父老.
很難信主的.
忽然間在床邊舉手指公來信主.
很清楚.
接著又是上個星期二.
很緊急去探望他.
應該是星期六.
大家在報道會.
我就有感動不參加報道會.
我就過去天宿圍醫院再探望世伯.
真正是已經.
星期二信主幫他洗禮.
星期六身體很虛弱.
我去探望的時候八點多.
又是有感動.
又是在床邊.
因為家人都知道不會再搶救.
不會再甚麼的.
大家都接受到.
要等他慢慢這樣離開.
但一直都持續很久.
好幾天.
星期二晚洗禮.
到星期六.
已經還有四天在天宿圍醫院.
在床邊又有感動.
我說天父他都是你的兒女.
這麼辛苦不如你接他.
又站好八點多就回家.
三點凌晨.
接著就很安詳地離開.
又是很平靜地離開.
聽著音樂離開.
一首聖詩.

$^{961}$聖詩就是你們晚堂唱的那首.
《平安》.
這麼奇妙.
你可以想像到.
平安之君.
平安的神會洗到我們全年的聖聖.
現在安排了七號在村公所出殯.
又願你們的靈與魂與生子都能保守.
所以平安的神預見我們這個凡人.
由罪人開始慢慢.
我們其實正在超越的.
你回想你的價值觀.
從信主到現在.
有些事情沒那麼執著了.
有些關係不要咬得這麼緊.
好像一定要自己正確.
有些事情可以慢慢放開.
重要的就是.
反而對福音越來越看重.
越來越有人跟你講福音.
去傳福音.
你就很想去幫忙.
這個原因就是.
因為我們看到靈魂才是真正不滅和永恆.
我們參與在神的聖功上面.
令到人可以得到救恩.
其實很有趣的.
為甚麼我叫他舉右手指的手指公呢?.
因為他的手指已經包著東西.
那些夾子.
可能是心電圖那些夾著.
然後他再用東西包著 怕它掉下來.
有時他睡醒的時候.
他會扯那些東西.
因為迷迷糊糊.
那個手指公很大隻.
好像腫了一樣.
其實是包著那些紗布繞著.
他在床上已經很虛弱了.
七點二十萬的時間.

$^{1001}$我說不如這次信主吧.
他說好啊 他點頭.
我就怕他點頭都不是很真.
不如這樣吧 你舉手指公吧.
他真的很大隻地舉起.
差不多是不可能你說他是錯的.
說他是假的.
你想想一個鄉紳.
一個為頭的父老.
是很迷信的.
是那一刻他真的清楚自己信主.
還有他可能是真的很辛苦.
他的癌細胞已經擴散了.
所以他真的決心信.
你可以想到那個救恩帶給他的生命改變.
這裡就是神.
平安的神進入我們的生命裡面.
保羅的意思就是靈與魂與身子.
即是你的女兒案外.
你的知識,你的感情.
所以我們有時應該要為主的福音而哭.
見到人信主你應該流淚.
你應該很開心,很雀躍.
見到人說我不再回你的教會了.
我不再信耶穌了.
你來耶穌是垃圾來的.
假如我當然不是這樣說.
如果有這樣的人.
其實你應該要為他難過.
沒錯 我一下子都不可以解釋給你知道.
希望你回心轉意.
但是我會記住你 我為你祈禱.
應該是這樣的心腸.
這就是靈與魂與身子都得無保守.
我個人這樣理解保羅的意思.
我想他這裡沒有意思說.
其實不是的.
人分了三個部分.
靈一定是與神溝通的.
魂就是我們的感情世界.

$^{1041}$我們的知識世界.
我想他未必分得這麼仔細.
未必需要這樣分.
但是重點仍然是平安的神.
進入一個人的生命裡面.
應該帶來我們要有一個標準要求.
就是我們首先.
原來是神會保守我們.
叫我們能夠得無保守.
接著在主耶穌再回來之前.
降臨的時候.
我們要盡量追求完全的無可指責.
我們要在主的裡面.
要有一個豐盛的生命.
譬如主耶穌說的十章十節.
「若人福音:我來了是叫人得生命」.
「得豐盛」.
所以我們應該追求一個豐盛的生命.
豐盛不只是物質上.
肯定是與那個靈魂.
與你看到別人的救恩.
它能夠被記錄在主的名下.
肯定有關連.
所以你看到別人能夠信主.
其實是一件很開心的事.
這就是我們一起去服事帶來的功效.
所以這位召我們的本事信實.
所以救恩是由他開始.
由他去賜下.
我們都是因為主的愛.
我們可以得救.
所以這個追求的成長.
追求成為一個成熟的基督徒.
絕對不是一種驕傲.
絕對不是一種我們靠自己的行為功德.
而是這位信實的主.
他會在我們生命裡面成就這件事.
肥牛秘書有一句說話很重要.
我們的立志行事.
都是神在我們的心裡運行.

$^{1081}$為了成就祂的美意.
這個肥牛秘書應該是二章.
To will and to do.
我們will就是我們的立志.
然後to do就是我們坐而起行.
我們決心做這件事.
原來是神在我們的心裡運行.
為了成就祂的美意.
原來我們這樣被感動.
背後已經是有神的工作.
所以在大體來說.
保羅說這位召我們的主.
他是信實的主.
他必定會成就我們生命的工程.
他必定會成就我們要放心的交託給神.
我們要做的就是我們盡力的.
好像上文所說.
教會必須的是我們要聆聽主的道.
我們要持守.
我們要在主的道裡面去經歷,去成長.
所以預備今晚講到.
今早就出現了一個運輸事件.
你就知道原來那些事是突發性的.
你有沒有想過.
原來你這句話是令到他誤會的.
半個小時都還沒有來.
原來他誤會了.
以為是他責怪他.
他就說你這麼久都不化妝,不接貨.
你都不見人.
你還說我拍得這麼遠.
其實這個是大家的誤解.
所以原來我們真的隨時都要學習神的道.
令到我們能夠處之泰然.
對任何不同的人都是.
甚至我們可以表達出.
正是因為我們裡面有神的平安.
所以可以駕馭我們本來的性格.
即是說原來我們的性格都可以慢慢改變.
如果我們被人感覺到.

$^{1121}$我們已經和以前的「我是」不相同.
這個是神的恩典.
是必然會改變的.
因為平安之主已經在我們的生命裡面.
開始祂的工程.
接著就是保羅的很簡單的文案.
最後的祝願.
請弟兄們為我們禱告.
與眾弟兄親嘴文案.
這個是他們當時的.
似乎是歐洲人的禮儀.
現在都是的.
即是基本的禮儀.
男女都是的.
興文這樣的.
我有時去加拿大.
我的外甥的朋友.
或者外國人.
他們都會輕輕撫一撫.
這個是很基本的禮儀.
但興文.
記住.
興文.
與眾弟兄親嘴文案務要成結.
我只著主祝福你們.
要把這信廉給眾弟兄聽.
即是說這是公開的信.
即是說要在聚會的時間.
不知道初期教會他們怎樣運作.
似乎他們未有很具體的形態.
有講台.
有講道.
因為未必有傳道者.
每個星期都在.
尤其是當時根本沒有聖經.
記住沒有新約聖經.
新約聖經根本沒有寫.
有的可能都很少的.
福音書.
前書,後書.

$^{1161}$後書根本沒有.
你可以想像到.
他們可能會重讀一些保羅的書信.
又或者是看舊約的抄斷等等.
所以其實那個階段有很多神蹟.
很多神特別的工作.
因為甚麼都沒有之前.
所以他們必須會重讀使徒的書信.
這個似乎是當時他們的慣例.
所以有路過的弟兄或是牧者.
他們就會邀請他們講道.
所以到了後期.
在約翰的年代.
就說不要隨便.
千萬不要隨便.
給那些來講道的人.
一定要分辨到底他們有沒有講異端.
這是約翰三書所講的.
所以這裡告訴我們.
保羅對牧羊,對貼城教會的弟兄姊妹.
其實未必有很多人.
那個牧羊都是基於他自己寶貴的使徒的身份.
看到救恩可以改變這些愛邦人.
貼塞羅尼加都是一個很富庶,很漂亮的地方.
似乎當年更加漂亮.
很天然.
亦都很富庶.
因為亞歷山大大帝的妹妹.
亞歷山大大帝為他的妹妹而建造的那些座城.
一定有很大的文化底蘊.
希臘的文化.
更加有羅馬的統治.
貼塞羅尼加城不是說了.
他們壓縱奧古斯都.
這個奧大維.
安東尼戰敗.
所以變成他們有那個.
既然你擁戴我做皇帝.
我又真的打贏了而做了皇帝.
是羅馬的第一位的國王.

$^{1201}$所以很自然地.
貼塞羅尼加城就享有很高的政治方面的地位.
以及那些待遇.
所以貼塞羅尼加地方都是一個相當繁榮富庶的地方.
今天都還是.
你去到塞羅尼加真是很漂亮的城市.
這本書其中幾件事.
保羅他所說的「信望愛」.
我們要留意.
我們很早期已經說了.
第一講已經說了.
「信望愛」原來不是只是一種基督教的術語.
而是在保羅的模樣裡.
他可以在這個信心.
愛心盼望上很豐富地去形容.
一個肅靈生命應該怎樣去追求.
所以原來「信望愛」是每一個基督徒.
必須每段時間要檢視.
是否一直在成長.
而且是到永恆.
「信望愛」都是這麼重要的.
所以當然其中最大的是愛.
因為你看見神,看見耶穌.
你的信心和盼望.
你都去到天堂了.
還要盼望甚麼呢?.
所以其中最大的是愛.
是很合理的.
但是每一個基督徒.
真的使我們成為基督徒.
我相信「信望愛」.
一定與其他人不相同.
一般人都一樣會有信念.
一樣會有盼望.
譬如下個月就去很美的地方.
一個月旅行.
而且是一個很開心.
和一群很多年沒見的好朋友去.
你想想你盼望這個月.
你天天都看著那些事情.

$^{1241}$很明顯你會明白甚麼叫盼望.
你會很羨慕那個盼望實現了.
就是下個月的旅行.
愛更加.
你又很愛你的家庭.
很愛你的家人.
又很愛你的國家.
又很愛你的生活.
很圓滿的.
不過沒有神的信望愛.
沒有救恩的信望愛.
就不行了.
你只能夠是今生今世.
你就不會有永恆了.
所以就算是很出色的人.
很有錢的人.
如果他沒有信仰.
最終他都是一樣面對.
他有一天.
他要甚麼都帶不走的離開一個世界.
還有就是進入永恆的黑暗.
所以不信的人都有信望愛的.
你就更加要明白.
原來基督徒的信望愛.
在保羅的眼中.
是非常激動的.
是不斷可以成長的.
是應該我們自己要苦心自問.
我是不是已經失落了.
那個信心.
對主,對人.
我的盼望.
是否已經暗淡了.
就是今天有酒.
今天醉了.
已經對永恆沒有甚麼期望了.
愛就更加不用說了.
你沒有信,沒有望.
自然你的愛就會窒息了.
你的窒息就是.

$^{1281}$你要看到周圍的人都不愛你.
你都不會愛周圍的人.
即是安慰大於心死.
所以基督徒的信望愛.
是應該和一般人的信望愛是不同的.
這是保羅在《貼壽羅利加前書》.
給我們看到的教導.
當然關乎到牧羊.
這本書有一個很好的示範.
保羅他要急速離開.
其實有很多要解釋的.
最想解釋你的動機.
為甚麼你要離開.
但是他反而不是那麼迫切.
要為自己辯護.
反而他很關心.
所以他在第三章就說得很明白.
他離開是有個別的原因.
但是當他去到雅典的時候.
他就很擔心弟兄姊妹的情況.
所以就打發提摩泰去了解.
所以這本書的寫作就是第三章的原因.
提摩泰在貼壽羅利加城回來.
去到雅典來會合保羅和西拉之後.
保羅聽到提摩泰的匯報.
他們仍然信主和站立得穩.
保羅就說我們活了.
我們的信心,我們的愛心活了.
他對這班弟兄姊妹只是一個月的形式.
其實只是路過而已 在那裡傳道.
也有猶太人對他的批評和逼迫.
保羅要離開.
但是這位牧者有那麼深厚的.
對弟兄姊妹的愛,靈魂的關心.
可想而知保羅絕對不是.
他們所眼中被人誤導的.
保羅見到利益就會連累他.
有困難就立刻走.
表面的形態就是這樣.
但實際上就不是.

$^{1321}$所以保羅後來會回到.
馬其頓這一帶去牧羊.
第二次完結之後.
他就回去第三次傳道之前.
主要在亞洲的爾弗所 在土耳其.
但是他去的時候.
都走遍馬其頓一帶地方.
這是二十章的《少林傳》所說的.
就是說在馬其頓一帶.
就再回去兼顧宗教會.
有機會再去接受羅尼加.
和弟兄姊妹的相交等等.
這個第二方面就是.
剛才第一方面就是.
我們要對信望愛.
基督徒的信望愛.
必須要與人與時的成長.
真的要自己檢視的.
因為真的去到永恆.
都還是信望愛的.
沒有其他東西的了.
你帶不走任何東西的.
唯有的就是我們的信心.
在主救恩裡面的信心.
對永恆的,對主的這個盼望.
和愛是永遠維繫著我們這個人.
成為人不是一個動物.
那個最重要的品格,品質.
所以基督徒的信望愛是要成長的.
是要越來越從你的面容.
都反映你有信望愛的.
基督徒的信望愛.
問你一點東西.
甚麼都不懂的.
甚麼都不知道的.
對很多的聖經蒙查查.
這樣就變成你的信心.
就建立在一種感覺.
不是建立在一個真理.
所以其實互為有關的.

$^{1361}$我們要讓我們的信仰.
不是只是一種口頭禪.
而是很實質我們經驗到的.
這位肅靈的主可以讓我們經驗到的.
所以盼望我們能夠在基督徒的信望愛.
在牧羊上面我們有份於參與.
看到保羅用這麼長的篇章.
講述弟兄姊妹在碧白患難裡面.
但仍然持守那個信心.
這個難能可貴的一種情操.
就是無論傳道者和被牧羊的弟兄姊妹之間.
的一份恩典的情誼.
這份情誼是需要不同的後來者都要學習.
就是說我們都要在教會裡面.
就算不是傳道人也好.
你都應該要成長.
你要關心那些後來者.
關心他們會不會跌倒.
關心他們會不會因為一些困難而退後.
這些都是我們要大家共同負起的責任.
就是說如果你都不知不覺有一種牧羊的愛.
恩典其實就代表著你都已經有一個主的心腸.
和一個成長的生命.
這是很不容易的.
因為要牧羊就是說有不同的人.
他們會落在不同階段的困難的時候.
我們就要很敏銳他的情況.
這是第二大方面.
第三方面就是最突出.
就是關乎到那個主在來.
每一章前書後書都有的.
每一章的結尾都不斷提到主的在來.
為何要這樣刻意.
是不是早期的保羅的神學.
他有一個很強烈的意識.
認定主耶穌很快在他有生之年就回來呢?.
有可能是的 這樣的命題.
但是我們都沒有什麼具體的答案.
因為後書他似乎主在來就說得少了.
反而就是他越來越知道自己就快要離世.

$^{1401}$他就自己特別提一提後書.
大家都相當熟的就是《納美好擊仗》我已經打過了.
當跑的路我已經跑盡了.
所信的道我已經守住了.
從今以後有公義的官面為我去傳流.
開始他說就是他自己等不到主回來了.
他要去聚集了 去到天堂聚集了.
他反而就是去盤點.
他自己現在他的信仰是怎樣呢?.
年紀大了要開始盤點了.
《美好的仗》我們有沒有已經打過了?.
或者我們開始打了沒有呢?.
所信的道我們是怎樣呢?.
我們當跑的路我們有沒有跑盡呢?.
就是開始要真的做一個盤點了.
所以保羅似乎去到晚年了.
所以晚年其實五十多歲而已.
他的晚年就開始等不到主回來了.
或者他已經知道不是他去計算.
而是神的安排時間.
他知道他不會灰心.
不會因為主不回來他就生氣.
而是他歡喜在他有限的人生.
他走完他人生的路了.
他要做的做完了.
或者神要他寫十三封書信寫好了.
成為後來.
後來說起來是過百年之後的傳道人開會.
二不所大會.
這本就是保羅書信了.
要放在聖經的正典.
才開始有新約.
在一個研討會裡面.
當然是一個很嚴肅的教父聚集.
因為當時已經有很多異端的假的書信.
所以要找一些真的出來做聖經.
新約聖經.
所以保羅都不知道自己十三封信會被人.
後世去讀的.
過了一百多年.

$^{1441}$再久一點.
他的書信才被人平頭品足.
這本就是保羅書信了.
這本就要放在聖經裡面了.
是這樣的意思.
所以其實有一本叫.
《哥林多前書》和《後書》中間的一本書信.
是遺失了的.
是不知在哪裡的.
就沒有放在正典裡面的.
所以我們就要明白.
你可以想像到保羅的人生.
他不斷講主在內.
但有一天他知道.
他自己要離開這個世界.
他就要盤點他自己的信仰.
和安排他特別的崗位.
他有沒有已經完成了.
所以又是回到那個問題.
盡快要發現自己在哪個崗位被神用.
你的恩賜在哪個崗位被神用.
因為臨前的十二章是講得很清楚的.
這一切的恩賜都是那一位的聖靈.
隨他的己義而分給人.
為著要使教會得到造就的.
所以每個人的存在.
都一定是有你的恩賜.
和有你的崗位去滿足你那個人生的.
所以我們真的要不斷追問.
有甚麼是神正在用你.
有甚麼是神特別一定要你用的.
你就只用那一樣東西.
你就不要去彈琴了.
你不要浪費自己的時間.
你彈琴的話那些牛都不會上樹的了.
你不會做那些不是你要做的事情.
要專責去傳道.
所以每個弟兄姊妹都真的要求主.
讓我們去學習.
又是困難的.

$^{1481}$因為來到我的曾經班的.
都是一些比較成熟的弟兄姊妹.
可能你已經錯過了你年青年幼的時候.
可能還未信耶穌.
還未認識主.
但是不要管了.
現在就是把握當下.
開始發掘你的恩賜.
發現你的恩賜.
更重要的是最終要發揮你的恩賜.
去服侍神.
恩賜不在教會給神用.
是你的才幹.
你會飛都沒有用的.
恩賜是一定要在教會裡面給神去用.
那就是叫恩賜.
去祝福其他人的.
這才是恩賜的定義.
所以這個第三方面就是主在來.
表面上保羅說了很多.
很快就會落實.
很快就會實現.
但是他知道他等不及主在來了.
他自己真的越來越知道自己身體會走下坡.
那個逼迫越來越嚴重.
當時是尼祿王.
已經把他拉進內監.
他在監裡面真的知道自己要離開世界了.
他真是最後的那個海哥.
盤點.
生命的盤點.
就是說你交付我的.
我已經盡力了.
雖然中國我還沒有去到.
也沒有船沒有飛機.
怎麼去呢?.
他一個人這樣走.
都還沒有發現大陸.
還沒有發現這些地方.
靠人的帆船怎麼去到中國呢?.

$^{1521}$所以都不關他的事了.
神沒有要他去做中國的人.
但是他寫下的書信.
2000年後就成為我們今天去研讀了.
所以都不是沒有傳遞給我們.
神的話語留到今天.
我們已經翻譯成中文了.
百多年前.
1918年整本的和合本正式面世.
有中文的聖經.
我們可以讀到保羅的書信.
所以他雖然離開.
但是他繼續說話.
這就是神奇妙的地方.
所以保羅他講主在來.
但是他知道主在來不是他去定的.
是神去定的.
所以保羅不斷講主在來.
他都不會灰心.
他都不會感覺絕望.
好像騙人的.
都沒有的.
後書就這樣說.
後書第二章就這樣說.
或者彼得後書也有這樣說.
那些不信的人就會嘲笑你們基督教.
又說主在來很快.
總之就是保羅他帶著一份不是遺憾.
帶著一份是滿足.
就是他反而在很短的十幾年的傳道.
他當跑的路他跑中了.
所信的道他守住了.
他要將基督教不是一個易拉的宗教.
不是猶太教的一個附屬.
如果保羅不講那些羅馬書的經文.
猶太教是會慢慢慢慢吃了基督教.
但是很清楚是由猶太教出來.
但是是獨一無二的新約.
所以保羅說自己是新約的執事.
是和舊約完全不關的.

$^{1561}$由舊約出的.
因為主就是一直這樣引用這個舊因.
但是歷史的轉變就是耶穌的降生.
一切都更新.
不再在舊約的陰影底下.
是完全新的.
所以外邦人就可以因為耶穌就得救了.
不需要進入猶太教了.
所以你可以想像到.
士多保羅他就有這樣的生命的改變.
所以我們去看《索羅尼加前書》.
有一位不斷講主在來的聖徒.
他帶著滿足地離開世界.
他帶著就是交完那份人生的卷.
就是那個考試交了卷.
然後很安然地很開心地離開.
分數是由實去評的.
不是每個人都要100分.
而是盡力就可以了.
所以我們每個人一生.
真是一個「盡」字.
盡心盡性盡義盡力愛主你的神.
其次就是要盡力愛人.
都是那個「盡」字.
所以我想去到任何一個時間.
一個基督徒如果他能夠把握每一天「盡」.
所以休息有時,運動有時.
去旅行休息放鬆有時.
都是講那個「盡」字.
就是你做這一切都不是為了只是那一樣東西.
而是為了盡力愛神愛人.
所以盼望我們看《索羅尼加前書》.
看到這位初期第二次宣教旅程的傳道者.
他已經有對主在來那麼強烈的教導.
還有他給到我們的觀察就是.
原來有次序的.
主的在來是有次序的復活的.
死了的人才復活.
未死的人見主面的就會身體改變.
是第二個階梯的復活.

$^{1601}$就讓我們看到這一日子是必定會降臨的.
就是說如果過了2000年.
可能我們這一代人都要等待.
但是這個只是等待.
沒有保證說這個時代一定會回來.
但是我們需要的就是盡我們的人生.
在我們的信望愛上面去成長.
在教會的場景.
神安排我們的場景去服事.
也都是帶出那份平安.
帶給其他人的那個詛咒.
然後我們記住.
主是否回來不是我們去決定.
但是我們要去到那裡聚集.
就是我們的責任.
要在臨離開之前要做一個盤點.
是不是我們已經在主的眼中.
我們已經成為一個盡力的門徒.
這個應該都是我們大家共同的.
很重要的那個責任.
所以這一本書不難讀的.
但是裡面的牧羊的情懷就很深.
看到保羅對弟兄姊妹.
其實真的是素未謀面的一個月的時間.
但是對他們很深的愛.
所以主耶穌可以令到我們.
對一些你不認識的人你都很有愛的.
如果我們稍為多一點時間.
每個每季都有人洗禮的.
那些名單都會登在程序表的.
我們看完之後就算了.
其實不是的 看完.
這八個人今年信主這一季信主是甚麼人呢?.
我又有時間 其實個個都不認識的.
其實全世界不斷歡迎.
歡迎大家回來洗禮聽.
我坐在一邊聽他們講他們的見證.
接著繼續為他們往後的成長去祈禱.
那份表千萬不要扔掉.
拿回家記下它.

$^{1641}$有空就拿出來.
他們現在的情況是怎樣呢?.
默默為這些新僕祈禱.
其實已經是在事奉一個很崇高的一個境界.
積少成多 不斷累積.
接著有一天見到這個姐妹 這個弟兄.
我為你祈禱的.
他很愕然的說.
我是初信其實甚麼人都不認識的.
這麼多人.
你居然認識我 為我祈禱.
為甚麼呢?因為我聽見你的見證.
很感動.
繼續默默為你祈禱.
其實這些的模樣比較少人做.
通常好像全球人一定要做.
固然一定是我們在做.
但問題是弟兄姊妹都有這些責任.
嚴格來說信徒的領袖.
你既然現在成為領袖.
你就更加要為教會弟兄姊妹祈禱.
某程度《神話大前書》列舉了那些.
誰可以做監督 誰可以做執事.
列舉了全部 很清楚的.
我們應該要對號入座.
對不起 哪個弟兄你其實不適合做.
因為你不符合.
大家回去讀一讀那個清單.
很詳細的.
要怎樣怎樣怎樣 很詳細的.
全部是在講品格上的.
以及對神的道的認識的.
一句就是善於教導.
監督,執事必須善於教導.
善 善就是你很擅長.
很熟悉.
對主的道是你很有那種愛慕追求.
從來都不上柴經班的.
從來都不去甚麼的.
突然間可以做領袖.

$^{1681}$你就覺得奇怪了.
就一個問題要大家反思的問題.
當然你不要對號入座.
這個是我們大家一起要學習的地方.
所以你看到保羅的恩情.
他對弟兄姊妹的愛.
日於言表.
這個是我們真的必須要學習的.
教會其實最大的那個存在.
就是很多弟兄姊妹都有愛心.
你就會令到很多初信.
或者面對困難的弟兄姊妹.
她就會感受到教會的愛和關心.
特別是移民教會.
加拿大現在無論英國都是.
很多新人移民.
很感受到那些舊的弟兄姊妹.
教他們怎樣融入社會.
買車,交通規則.
有需要就在沒有車之前就接送去教會.
很多這些機會令到他們感受到愛.
很多這些機會製造出來.
他們得到最快的資訊.
還有很多的幫助等等.
很多這些的見證.
所以你會令到那些團契.
很快就會很多人.
很快就會大家成長.
在那個群體裡面.
就是近這兩年.
香港當然不斷有人走了.
但是在那邊明顯見到那個成長.
就是這些細微很小的事.
就感受到愛.
感受到弟兄姊妹的模樣.
那個關心.
就帶來在這些這麼小的事情上.
就製造到彼此的愛.
香港我們不是這樣的形態.
我們不需要刻意去製造這些形態.

$^{1721}$但是總是我們要找機會,找辦法.
來實踐這些不只是在口裡說.
而是有具體的行動的愛.
所以求神幫助我們能夠透過前書.
後書就負面一點.
後書就開始罵一些人不守規矩.
前書就很正面的探討主在內.
不過再補充多一點.
還有14分鐘.
講主在內其實是基於第四章.
因為已經有信徒知道主會回來.
就反而擔心自己死去的家人去了哪裡.
還有狀況是怎樣.
所以原來保羅講的主在內.
是基於一個的安慰,一個的模樣.
對當時還沒有聖經.
傳道人都不是經常在一個地方.
駐守的模樣.
是很容易帶來那種誤解,懼怕和憂心.
不知道自己離開世界的家人會怎樣.
所以在主裡面就是真正的安穩.
所以我們要看這一段的主在內.
其實是基於模樣的.
基於一種的安慰的.
基於不要被一些異端影響的.
所以主在內在這裡.
貼壽羅尼加是一個模樣,是一個安慰.
就不試探討那個時間,計算那個日子.
有甚麼方法.
就不是這樣計算.
而是那個事實.
但怎樣將這個事實套用在很多人的心靈裡面.
其中一樣就是很擔心主在內.
自己死了的親人會怎樣呢?.
保羅就說他們會首先復活.
所以不需要擔心.
還有嚴格來說.
他們現在比我們更加開心.
因為靈魂再沒有痛苦,再沒有眼淚.
所以我們現在仍然在這個地上.

$^{1761}$是要滿足神要我們完成的責任.
所以保羅自己最終是被斬頭的.
被斬頭而死.
所以整個主在內在他的內心.
還有他知道死了的人是首先復活.
所以尼祿的斬頭只是會傷害到他的頸.
沒有傷害到他的靈魂和他的永恆的盼望.
他已經安睡在神的懷中.
使徒保羅.
他已經有公義的冠冕為他傳能.
所以我們要看主在內成為我們有限人生的.
大家的終極的盼望.
所以我們今天如果大家想得到一句.
袋落袋的就是我們要盡我們人生.
盡心,盡性,盡義,盡力.
愛主你的神.
其次就是要愛人如己.
相信也是盡力的愛人如己.
很盼望大家一起去服事.
接下來三壽我們在紅印堂又有福音主日.
因為教會盡量要多福音主日.
就可以帶不同的人回教會去信主.
會有傑霞姐妹講她的生命故事,她的見證.
很盼望有更多人透過不同的方式.
讓那個福音可以傳遞.
我們一起簡單禱告.
感恩天父給我們能夠在過去的一段長時間.
我們都可以有系統地去讀聖經.
我們很盼望教會的生活.
本來就應該是平安的主.
在罪人的內心帶來的是復和.
帶來的就是有永恆的福樂.
和內心的平安.
和感受到對永恆的憧憬帶來的喜悅.
和那個盼望.
很盼望我們能夠不要被地上的人事物影響.
這個肯定.
如果不是聖經都不會教我們.
千萬不要以惡報惡.
幫助我們能夠建立一個內心的和睦.

$^{1801}$和內心裡面的一份有主的聖靈的感動.
一份的愛.
以致我們在面對不同的事情.
都能夠有主的同在.
有主的恩典.
這樣來不是逢凶化吉.
而是使壞事都可以變成美善.
願主幫助我們.
有機會有系統的一起讀聖經.
願本來不是那麼完善的錄音.
都可以祝福不同的弟兄姊妹.
大家都有不同的需要.
不同的場景.
願神都幫助.
唯有你知道我們每一個人的需要和心靈.
願神來保守.
下一個星期我們會進入後書的時候.
又看看主在來.
為甚麼會帶來一些人的誤解.
而令到他懶散.
都讓我們從其中又發現.
原來真理不可以一面倒的.
我們要求主幫助.
奉靠主聖名.
阿們.
\newpage



\section{帖撒羅尼迦後書}
\label{sec:uJ_qz3eVEf0}
\textbf{2024-07-11 周四晚查經班《帖撒羅尼迦後書(1)》 講員:曾裔貴牧師}
\newline
\newline
連結: \href{https://youtube.com/watch?v=uJ_qz3eVEf0}{\texttt{https://youtube.com/watch?v=uJ\_qz3eVEf0}} ~~~~ 語音日期: 2024-07-12
\newline
\newline
\hyperref[sec:4L9GLJKp4e4]{\small{< < < PREV SERMON < < <}}
~
\hyperref[sec:index_chronic]{\small{[返順時目]}}
~
\hyperref[sec:index_scriptual]{\small{[返順卷目]}}
~
\hyperref[sec:ntrA1aTVgRk]{\small{> > > NEXT SERMON > > >}}
\newline
\newline
$^{1}$我們一齊恢復晚上的參景班.
一齊禱告.
我們能夠在天氣相當炎熱.
但也相當的晴朗.
我們紀念在湖南洞庭湖.
都有很嚴重的水災.
還有一些地方有很嚴重的乾旱.
甚至乎現在好像在美國.
也曾經有一些很嚴重的水患.
好像看到天氣氣候變化.
都在在的.
讓我們看到.
那個似乎修復修補的問題.
都未必能夠自我復原.
很盼望神都保守.
尤其是讓世界的領袖.
在那些氣候制定的減排.
碳的排放.
都能夠真是有一個果斷的心.
不要將一些問題.
留給下幾代的人.
來承擔,來承受.
天氣越來越炎熱.
都是反映世界需要真是有一個轉變.
求主幫助.
我們其實一直都在說主在來.
這個很重大的課題.
當我們處理完這兩本書之後.
我們就會探討.
在聖經裡面這個很重要的真理.
求主讓我們能夠憑著信心去明白.
今天晚上我們進入後書.
幫助我們能夠透過保羅他的模樣.
怎樣來到他有前書之後.
再有一些的安慰和提醒.
都是讓我們能夠從其中得到教訓.
特別今天晚上.
我們又再一次看保羅怎樣勸勉他們.
在苦難裡面.
去有的信心,信靠和堅忍.

$^{41}$願主都保守我們大家的心靈.
奉主耶穌基督的名.
阿們.
我們稍為先讀一讀聖經.
我們之前都說了.
貼索羅尼加前書全卷.
有五章.
如果大家要稍為記得.
貼索羅尼加前書.
每一章的結尾都提到主耶穌回來.
和教會的關係.
亦提到兩次提.
貼索羅尼加教會.
他們的信望愛是很活潑的.
還有是真的有具體的一個行動.
所以這個信望愛.
似乎在保羅的書信.
保羅自己的神學的理解裡面.
絕對不是一式的文字.
用來裝飾.
而是很激動的.
信望愛在保羅的運用裡面.
在不同的處境都有不同的.
所以我們都要學習.
我好像曾經說過.
我們無論如何.
去到我們人生終結.
給我們能夠存留的.
只是有信望愛.
所以如果我們在這麼有限的人生.
不去發展我們自己的信望愛.
或者當然群體都應該有信望愛.
但我們自己就肯定.
是我們永恆裡面的美德.
屬名的美德.
我們不是與生俱來就有的.
當然每個人都有基本的信心.
基本的盼望和愛心.
但是那些是未信的人.
天然就有.

$^{81}$但是信了主耶穌之後.
是要經過在舊恩.
在主耶穌的生命的改變裡面.
去發展的新的信望愛.
存到永遠的.
所以等於我們給了一個戶口.
我們要存這些信望愛.
如果那是錢 那是貨幣.
我們要存那個信望愛.
所以初信的信心.
應該和你三十年後.
信主之後的信心.
應該有些不同的.
因為你已經有了很多經歷.
盼望一樣.
愛心就更加一樣.
都是需要不斷的發展的.
所以我們不要只是有一個戶口.
有信望愛.
但從來我們都不儲錢下去.
所以將來就真的空手見耶穌.
即是攤開雙手 什麼都沒有.
其實是很大的損失.
所以今天第一章.
正正就是講這個問題.
正正就是那個信望愛帶來的生命.
不知道大家有沒有聽過.
加爾文神學.
他有一個叫做Tulip.
Tulip即是鬱金香.
荷蘭的 是不是薄花.
鬱金香.
鬱金香就是Tulip.
T U L I P.
那個P 當然這個是加爾文.
加爾盟約它弱化了.
加爾文神學有五個重點.
我今天想強調這個.
普羅書信裡面.
今日講的《跳傘羅尼加》後書.

$^{121}$有一個叫做聖徒的堅忍.
堅定的忍耐.
所以是忍耐到底必然得救的.
所以這個忍耐.
即是在苦難的世界裡面.
在面對迫迫裡面.
這樣叫做堅忍.
所以在加爾文神學裡面.
一個聖徒是要一直相信.
一生相信.
這樣才能夠得救.
所以加爾文神學是很強調.
那個聖徒是要堅忍等待主.
和今晚我們第一章要講的.
就有密切關係.
我稍為讀一讀.
很短的十二節.
你看到有大概三個段落.
都容易分的.
一間很好的教會.
一位很公義很好的神.
第三個就是.
一班很好的人.
為受苦的聖徒祈禱.
所以第一節開始.
保羅西拉提安泰.
寫信給《跳傘羅尼加》.
在神我們的父.
與主耶穌基督的教會.
願因為平安.
從父神.
和主耶穌基督.
歸與你們.
我拆開講.
如果你研讀保羅的十三封書信.
其實有一樣東西很好玩的.
就是他每一封信的文案.
都有一些不同的.
當然有專書寫了.
你可以去看看.

$^{161}$保羅的文案.
和其他的書信有什麼不同呢?.
你會發現.
有些端倪.
你會看到.
例如《迦拉太書》.
很明顯的完全沒有問安.
因為一開始第三節就罵.
罵到不得了.
罵那個教會.
接近傾向異端.
靠行為得救.
保羅就完全沒有任何問安.
所以你可以想像到.
不同的保羅書信.
可能早期.
後期都有不同的.
例如保羅寫給提摩太太.
提多.
特別多了一個促詞.
多了一個足源.
你回去做功課吧.
我說了就沒有意思了.
唯有寫給個人的提摩太太.
和提多的書信.
就會多了一種文案.
你回去看看.
其實很簡單.
一看就看得到.
一般寫給教會的書信.
就不會這樣寫.
寫給個別的傳道人.
他就會加這一句說話.
其實是兩個字.
你回去發現一下.
第三節.
弟兄們啊.
這本是合宜的.
因為你們的信心格外增長.
並且你們眾人彼此相愛的心也都充足.

$^{201}$甚至我們在實的教會裡為你們誇口.
都因為你們在所受的一切迫不還難中.
仍夠全忍耐和信心.
這個忍耐就是堅忍.
在迫不還難當中.
喬治羅尼加的聖徒.
仍然有一個堅忍的心.
而且信心也能夠成長.
所以保羅說.
因為合宜的.
是很合適的.
因為你們真的有這樣的表現.
所以更加印證.
印證加爾文神學裡.
其中一樣很重要的就是.
無論你處在什麼環境.
你信了主之後.
你要忍耐到底.
直至主耶穌回來.
大家有沒有聽過林美香.
很粗俗但是很傳神.
廣東話.
雖收尾.
就是信仰的破產.
其實大有人在的.
很多人初信信得很好.
但慢慢經過了很多很多的人生.
際遇,困難,否定主.
不再相信.
是決心不再相信.
也有這樣的人.
所以不要小看那些患難.
好像真的會淘盡我們的心靈.
令我們枯竭.
所以如果一個老聖徒.
仍然保持很單純.
記住是心靈的單純.
不是腦袋的單純.
很容易被人騙.
腦袋單純.

$^{241}$是心靈,心胸很單純.
對主很有信心.
對人也很有盼望.
這樣的聖徒.
他真的有這樣的堅韌.
因為他不是沒有絕緣於那些患難.
而是那些患難伴隨著.
但他深信主是短暫的.
而且是一個過渡的.
神有他的計劃.
有他的時間.
我脫離地上的帳篷.
我就會得到永存的居所.
永存的新生命.
所以這就是堅韌.
所以保羅在這裡.
如果作為一個現代的牧者.
你怎能夠說.
你這麼多難處.
這麼多人生不如意的事.
被人眼巴巴的逼得.
真是要投訴.
要甚麼甚麼.
因為你們在患難裡面.
仍然存有信心和堅韌的心.
這樣會的.
這樣會養弟兄姊妹的.
那個要求原來.
不是保羅自己特別提高了那種門檻.
本來信仰應該是這樣.
所以是不是我們今天的基督教.
說得太過便宜呢.
太容易呢.
很容易就說.
希望你不要有這麼多患難.
保羅說患難令你們更加反映.
你們那個堅韌的心.
更加你們那個信心是很增長.
所以這裡告訴我們.
真是有一間很好的教會.

$^{281}$到底前書和後書相距多久呢.
我們就沒有資料了.
保羅也不太需要給很多資料.
但是無論如何.
貼索羅尼加前書寫作的時間.
是保羅離開幾個星期.
去了雅典之後.
就寫這封信.
透過提摩太再帶回去.
然後了解了情況.
他就再寫後書.
來去兼顧安慰他們.
在這裡的時候.
仍然很讚賞他們在苦難裡面的表現.
保羅沒有去批評那個苦難.
更加沒有抓那些製造苦難的人出來鞭屍.
去控訴,去咒詛.
沒有.
保羅還很欣賞你們面對苦難.
仍然有那個真正的信心.
即是說這是真正的信仰.
在你們的生命裡面流露.
基督徒其實真的需要看你們的表現.
所以不是靠行為就能夠救人.
但是你的好行為.
是順了主之後的表現.
順了主之後.
你就應該整個人有所不同.
你對仇恨會有新的角度.
不再是以牙還牙.
以眼還眼.
而是以愛還愛.
而是你看到主耶穌都是這樣學習.
我們在人生.
有時面對一些誤會.
有些刻意不會化解的冤仇.
記住在提摩太前書.
應該是四章說.
有些叫不解怨.
有些罪叫不解怨.

$^{321}$聽過沒有.
你可能不熟悉.
本來列舉很多種罪.
其中一種罪叫不解怨.
甚麼叫不解怨呢.
就不是說你都沒有解的.
不是這個意思.
是不肯解除那個冤仇.
即是我氣你氣到死為止.
這樣的一個存心.
香港大把這樣的人.
不信的人大把.
成為脈絡人不再見面.
都有這樣的人.
但基督徒 基督教.
保羅說有的 小心.
那些人你要留意 要擋開.
末世就有這些.
對不起 這是第五章.
你就看到那些問題.
所以這裡告訴我們.
保羅首先稱讚這間教會.
第三到第五戰.
是一間非常好的教會.
沒有傳道者.
但碧白仍然存在.
按道理土壤很難生長.
但原來.
很多人都告訴我們.
千金難買少年窮.
在草根成長.
墳外的心靈.
或者胸襟.
或者見識.
可能會好一些.
因為經過很多歷練.
等等.
至少你的生命.
那條命都耐得住.
動不動就怨天尤人.

$^{361}$當然這是基本分析.
不是每個人都一樣.
很多階層都一樣.
很奮世疾逐.
很多有錢人都很好的品格.
雖然不可以一概而論.
但無論如何.
一間受苦沒有停過的教會.
一間在保羅要離開.
不可以和他們一起.
共渡患難的時間.
只是提摩泰在那裡幫忙.
輾轉提摩泰都要回雅典.
接著保羅再寫後書.
仍然是這樣去肯定.
他們在苦難的時代.
尼祿造王的時代.
那個苦難伴隨著.
所以原來基督教從來不需要懼怕苦難.
我們不需要去強求.
不要有苦難在我們這一代.
這些真是強求.
我們是強行來的.
每個時代都有它的苦難.
每個時代的苦難都有不同形式去包裝.
總之令我們離開實的.
那個就是苦難.
我們要記住這一點.
可以是身體的.
可以是心靈的.
但無論如何不要令我們離開實.
所以現在歐洲美國.
是說不是逼迫.
誰逼迫你呢?.
你信耶穌不信耶穌.
是很尊重你的.
每個人都很尊重你的信仰.
但是他用其他的世俗.
來令你對信仰已經不感興趣.
對返教會覺得是一件很過時的事.

$^{401}$教堂的維修是要很多錢的時候.
就將它變成放物館.
將它變成可以給人收費.
為了要保護的一個地方.
變成一個古蹟.
不是再是人的心靈.
去到朝見實的地方.
這個是歐洲現在.
任何一個地方都很普遍的.
所以我們記住苦難已經是.
不同的形式出現.
但是在帖撒羅尼加後書這裡.
仍然是忍受著初信的時候.
前書的時候那些光景.
他們是面對著很大的苦難.
但是有沒有傳道人.
有沒有好的牧者.
一直陪伴.
但是他們最重要就是.
這群弟妹有一個團結的心.
有神的靈在他們裡面.
有他們對信仰的堅持.
所以對信仰的堅持是最大的動力.
首先就是你找到一個真正的信仰.
你就要為信仰堅持.
信仰就真的是一個很重要的信念.
整個生命的依歸.
不是我喜歡就聽聽.
不喜歡就對這個宗教沒興趣.
這個就說不上是信仰.
真正的信仰是你要為真理.
來付代價的.
所以這個就是告訴我們.
保羅說我可以在各個教會裡面.
不同的教會裡面去稱讚你們.
甚至我們在神的各個教會.
為你們誇口.
保羅當然不是來吹捧.
吹噓.
其實我現在都說這是一個很富庶的.

$^{441}$歐洲的一個很漂亮的城市.
真是值得去的.
真是去看一下.
當時真的是亞歷山大大帝.
為他的妹妹所建造的城.
因為是他爸爸建造的.
真是一個很漂亮的地方.
風景如人.
但是這裡告訴我們.
這個教會的人.
他們信得真是很堅定.
保羅一說他們信.
他們就持守那個信仰.
所以保羅說.
現在我要為你們誇口.
所以保羅不是誇耀哪個人.
保羅誇耀他們的信心.
是他們那個堅忍的心.
所以我們都要真是我們自己.
平日不是很覺的.
當有一些難處來.
或者有一些事情來.
我們是不是仍然站在真理.
這個就是很重要的.
所以使徒們或者傳道者.
最開心的就是看到信徒的成長.
看到弟兄姊妹的靈明的增長.
面對一些事情的價值觀有改變.
不再是一般見識.
我們在《邊卷書》曾經說過.
一般人就一般見識是很合理的.
但是我們是屬靈人.
對不起 哥林多前書.
我們說過我們是屬靈人.
應該就有屬靈的見識.
我們其實同時間既是屬靈人.
同時間又是地球人.
所以我們是兩個世界遊走的.
我們什麼時候祈禱依靠神呢?.
我們就去到靈界了.

$^{481}$我們什麼時候用我們自己的血氣.
去決定一些事情呢?.
我們就是屬世的人了.
所以隨時這樣遊走的.
我們那個身份是很重要的認同.
自己是基督徒.
所以真是要讓我們那個靈明.
每天都要有經歷.
所以保羅為他們誇口.
正是因為他們在什麼呢?.
在《碧白環瀾》中.
這是第四節的下半節.
營救全堅忍的心.
這個Persuadance就是堅忍.
保羅書信裡面都是常出現的.
所以你不要以為保羅的福音.
信就行了.
其實信之後是要付代價.
為信仰去付代價.
保羅的恩信正義是說.
我們都要為信仰付代價.
當然不是我們釘十字架.
是我們要為主耶穌.
給了我們的救恩.
我們用我們的生命.
去見證耶穌的復活.
在我們的生命裡面去流露.
所以這裡就告訴我們.
這是一間很好的教會.
當然我們不是說.
不用傳道人都可以.
也不是這個意思.
否則保羅都不會趕緊叫提摩太.
帶這封信去.
看看他們的情況.
繼續勸勉他們.
可能有其他的牧者去到.
不過我們不是太清楚的意思.
所以一個教會如果沒有牧者.
嚴格來說.

$^{521}$基本上都很難成長.
如果我們去到有一天.
教會是不需要牧者的.
牧者你麻煩出去打工.
禮拜天大家都是傳道人.
當然好.
不要太重負擔.
沒有問題.
這樣就最好了.
到時你喜歡怎樣就怎樣.
他喜歡怎樣就怎樣.
根本就不可能有一個.
在穩定下的結構下的成長.
似乎由主耶穌升天之後.
已經驗證了.
神都安排了.
安插了.
呼召了.
牧者的存在.
在每一個時代.
就是要為著神的兒女去效力.
今早我在黃塵.
去聖經班.
跟Odi姑娘學習.
好深感受.
初一,二,三年班.
很乖.
叫甚麼都有反應.
舉手說鬥快.
答問題.
接著第二節.
四,五,六年班.
六年班都坐在那裡.
好像空神樂煞那樣.
坐在那裡攤開.
Odi姑娘說甚麼都不理.
不認.
不唱詩.
要站起來.
不站不行.

$^{561}$就這樣了.
我內心很大感受.
都有些牧者.
有些老師.
都明知是這樣.
都默默地繼續去守望著這班年輕人.
將來希望他們偶然長大.
有機會再聽.
那個生命的扭轉.
好像今天就有一個小小的種子.
落在他們身上.
可能聽完就水過鴨背.
根本不會再放在心裡.
當然偶然間有些小數.
是很有心不斷地有反應.
你知道原來這些小朋友.
未了解珍惜生命的道.
未了解看到.
其實和你的生命有關.
未看到就是說.
這位做物主一直愛你.
我們都做過小朋友.
我們一樣做過那些很反斗的人.
但是神的工作.
神的時間.
但是你看到有一些很用心的牧者.
和一些老師.
默默地去做.
你想想做那個part time.
用了多少時間.
用了多少心血.
去找資料.
我就坐在旁邊.
學習.
我就看到這樣的對比.
你就可以想像到.
這個工作.
神透過牧者的牧羊.
是由主耶穌升天之後.
到主耶穌第二次回來.

$^{601}$肯定神會在每個時代.
呼召不同的傳道人.
所以如果有些信徒領袖.
膽敢說這個世界不需要傳道人.
那麼他的教會很快就會倒下.
不再是教會.
沒有傳道人的教會就不是教會.
當然有傳道人.
要加上有神.
就不是傳道人.
所以其實是一個正面的.
有規劃的牧羊.
才能夠經歷聖靈的工作.
這個才是最重要.
不可以只喜歡這樣.
我又喜歡這樣.
這樣的打游擊.
一定不會成長的.
當然我說遠了.
不過保羅就是看到他們的需要.
雖然他們的人不能回去.
但是他也要寫信去牧羊.
首先頭五節.
肯定教會的美好.
美好是甚麼呢?.
從來你見到保羅不會誇耀建築物.
不會誇耀哪些有錢人.
誇耀的就是他們的信心.
這裡誇耀他們在迫不還難中.
營救全堅忍的信心.
這就是真正重要的資產.
基督能夠抗衡尼祿王的.
就是很單純的對主耶穌的信靠.
無論如何我不會否認我的耶穌.
這就是最大的武器.
就是內心的信念,信仰.
所以第一節,五節.
可以告訴我們是一個很好的教會的名正.
你可以想像到.
第四節說的是.

$^{641}$營救全人類和信心.
所以盼望我們的堅忍.
我們在家庭.
如果有很多未信的家人.
在工作地方有很多未信的.
或者在教會時.
又面對一些誤會,難處.
我們要學習堅忍.
學習信靠神的信心.
因為明顯是對付入座的.
保羅說你們在一切的迫不還難中.
營救全堅忍的信心.
所以是相輔相成的.
即是說那個苦難.
會更加凸顯你內心裡面的信仰.
所以是騙不了人的.
有知於內就迎知於外.
有時我們牧者就是這樣.
有時一些人突然發脾氣.
你都沒辦法想像.
好像原來他幾十年的功力就沒有了.
其實不關功不功力事.
是裡面仍然存藏著很多血氣的性格.
其實很需要神的恩典.
所以通常牧者都是在背後祈禱,流淚.
不是為自己流淚 是為這些人流淚.
這是很重要的.
所以記得一些長輩在教會裡面.
未必你的行動有很大的活力.
其實你們最寶貴的地方就是你們的眼淚.
不是多愁善感的眼淚.
是你看到教會的荒涼,流淚.
這些就真是很重要的.
這些屬靈人未必要到處打鼓,打鑼.
我們要祈禱了.
未必需要的.
你看到一些人的情況.
就慢慢為他祈禱守望.
或者你看到一些人很被神用的.
你更加為他祈禱求神保守他.

$^{681}$其實一個屬靈人就有這種眼光.
特別是年紀大了.
教會久了就應該有這種眼光.
有這種眼界.
我們應該有時間可以這樣祈禱了.
我們這個年紀.
開始第五節了.
神 這一位是怎樣的神呢?.
為何這麼多地上的不公,不義,患難,迫不.
而且是針對著衝著教會來的.
到底有沒有神?.
保羅會問這些問題嗎?.
從來都不問這個問題的保羅.
接著就開始回答.
到底神是怎樣做事的?.
第五到第十節.
就是答案了.
就是這位公義的神.
怎樣來等待那個時候到.
祂的兒子降臨.
就會履行終極的審判.
所以原來不是沒有答案的.
原來不是不來上善罰惡的.
原來不是沒有終極的審判.
第五節.
所以大家可能要對這一章熟悉一點.
因為解釋了很多我們一直纏繞著的問題.
又或者有些不信的人.
他只能看表面.
他當然只能看表面.
他都未接受聖靈.
他都未有主耶穌在他內心.
他一定是表面的.
但是我們怎樣去梳理他的內心呢?.
第五節開始.
「這正是神公義判斷的明證.
叫你們可算佩德神的國」.
出現了兩次這個佩德.
第五和第十一節再出現的.
保羅真的很誇張.

$^{721}$你們這樣的堅忍.
正正就是你們這班.
接受羅尼加教會的人.
佩德過神的國.
即是說原來神的國不是那麼兒戲的.
神的國不是沒有門檻的.
神的國度原來不是說.
你隨便就可以了.
都可以進去的.
不是的.
原來保羅說.
這就證明神公義的明證.
神用公義來判斷.
哪些人才可以佩入他的國.
就是你們這些貼羅尼加的人.
這樣的一個層次.
入到神的國的人.
是否要佩呢?.
佩是什麼原因呢?.
就是神公義的判斷的明證.
即是說神作了一個 witness.
作了一個裁決.
你們這班貼成的人.
就是真的佩入神的國.
所以嘉義文在這一章聖經.
拿到很大的靈感.
所以他認為.
這個簡選.
預定論.
真是因為背後有神的安排.
所以那個聖徒.
是會經歷很多的患難.
他都仍然有那個信心.
去信靠神.
當然太簡化了.
就令到大家可能未必會全面掌握到.
嘉義文神學的精髓.
因為當時也有另外一種.
叫做亞米利亞的.
另外一種學說.

$^{761}$彼此有一些衝突.
彼此有爭論.
也有他們的門徒.
兩派的門徒.
是有很大的爭辯.
但是我們在這裡.
我們看到.
本來一點合乎都沒有.
神公義的明證.
神作了一個裁決.
你們真的配得我的國度.
這個國度是屬於你們.
接受羅尼加城的人的.
就是因為你們.
全個堅忍的信心.
這是第五節.
接著你們就是為這個國度受苦.
所以我們是不是很少聽到這些爭道呢?.
原來信耶穌上天堂.
是要為這個國度受苦的.
原來我們要為福音來效力.
原來如果這個世界是敵黨主耶穌.
敵黨真理.
我們要站出來說.
不行啊.
我是和主耶穌一起的.
我們要為這個國度受苦.
我們要為這個國度受辱.
所以你會看到.
保羅是不會降低他們的要求.
叫他們利祿王.
避重就輕.
保羅不是這樣的.
當然絕對不是說.
你不可以去別的城市.
如果在《大白天》裡面.
主耶穌也有教的.
在24章的《馬太福音》.
但是重點就是.
無論你去哪個城市.

$^{801}$你都要傳一個.
為這個國受小的心切行動.
這樣就是重要的.
那個配得神的國度.
就是worthy這個字.
那個字根是worship.
所以中世紀所說的worship.
是說對這個造物主的價值.
形式.
而它是worthy.
是值得你去敬拜.
這樣叫做worship.
所以我們要搞清楚.
人家看這個worship.
是背後那個worth.
那個字價值.
worth.
W-o-r-t-h.
worth.
worthy就是很值得你去做的.
配得.
所以要配得是很對的.
中文這個配.
真的很對.
所以你會看到.
所以我們敬拜.
如果我們根本不認識.
我們的主.
我們的神.
說不上敬拜的.
所以新朋友來的.
福音主人是對的.
是福音性.
因為他真的不認識神.
他來到真的格格不入.
但當他信了主.
他有生命.
有了聖靈.
他就是敬拜了.
因為他已經認識這位造物主.

$^{841}$所以真的有不同.
所以這裡普羅告訴我們.
這一位公義的神.
他怎樣做事呢?.
就是會收納這些.
在為他的國受苦的基督徒.
接著.
第六節.
神既是公義的.
再一次強調神的公義.
在苦難裡面強調神的公義.
即是有兩個情況.
即是原來神不知.
在苦難裡面弟兄姊妹的困難.
痛苦,無助等等.
即是神的時間未到.
容許苦難.
所以不要以為神缺席.
也不要以為神是無能.
絕對不是.
我們在無助的時候.
不要以為神也是無助.
神是有他的判斷.
所以留意時間.
神既是公義.
就必將患難報應.
那家患難給你們的人.
即是上善罰惡的時間.
是在神的公義的判斷裡面.
最後再說主在來.
這裡說.
報應那家患難給你們的人.
第七節.
也必使你們受患難的人.
與我們同得平安.
似乎是在說主在來的時候.
有永恆的安息.
接著說主在來.
那時.
因為與平安平衡地說.

$^{881}$那時主耶穌與他有能力的天使.
即是第四章的時序.
降臨的次序.
肉身降臨的次序.
主耶穌.
那時主耶穌與他有能力的天使.
從天上在火焰中顯現.
要報應就是上文所說的.
要報應那不認識神.
和不聽從我主耶穌福音的人.
不認識.
不聽從福音.
雙重.
不認識不一定會迫使信耶穌的人.
但是不認識神.
再加上不聽從神福音.
耶穌的救恩的人.
你信耶穌.
笑你.
罵你.
再嚴重就迫使你.
這就是在地上現在出現的兩個情況.
所以我們必須要有心理準備.
基督徒.
其實應該由他信耶穌.
決志.
福音主人就要說.
你信耶穌.
你要想清楚.
跟你決志沒有問題.
但是你要為這個國度受苦.
我們很少人會這樣說.
保羅就在這裡.
他就說這位很好的主.
這位神.
原來從來都是.
用公義.
用他的聖見.
去做一個研判.
他一直看著世界的發展.

$^{921}$他一直都來到.
絕對不是監控.
他一直都是來到.
和他所創立的教會.
和弟兄姊妹同在.
就是說2000年從來沒有缺席.
這是給我們一個很大的安慰.
首先我們是在受苦的行列.
保羅已經再一次重新肯定.
你們是配得進入神的國.
因為你們是為這個國度受苦的.
所以我們就不是在將來.
來與神為敵的那個行列.
我們就是原來一直.
在與神同行的行列.
神現在是在判斷哪些是惡人.
哪些是善良的.
將來他就要來報應那些.
加患難給基督徒的人.
所以如果我們今天已經是基督徒.
今天已經決心在這個信仰底下.
為這個國度受苦.
如果我們有這個決心.
我們就已經在走在這條.
通往永生的道路上.
如果不是我們就要面對什麼呢?.
就是報應哪些人呢?.
報應那些不認識神.
和不聽從我主耶穌福音的人.
哪些人會傳耶穌的福音呢?.
一定是保羅.
一定是歷世未來的教會.
一定是持守真神信仰的基督徒.
所以人家不聽不認識.
就不代表他沒有行動.
他們是逼迫.
一代一代都會逼迫基督徒.
所以我們就要記住.
既然神將來都報應.
報應即是他們有行動.

$^{961}$不是在腦裡要求這麼壞.
不是這樣想完就算了.
是有行動的.
是有排擠的.
是有攻擊的.
所以我們要記得.
永恆的審判是在將來.
首先我們自己千萬不要.
對別人有一些超過神的控訴.
或者我們自己不肯解除一些冤仇.
記得我們已經一把年紀了.
不要介意正面的.
我們應該內心坦蕩蕩.
不會有冤仇.
應該是很開心.
任何人都是我們的朋友.
雖然不可以心交.
但是一定是不可以有仇人.
絕對不可以.
這是我們基督徒一定要學的.
所以你會看到這裡.
就是要報應那些不認識神的人.
要留意第九節.
他們就是這些不認識神的人.
原來是要有報應的.
原來是有審理根的.
他們要受刑罰.
永遠沉淪的刑罰.
很恐怖的.
離開神裡面的永遠刑罰.
是今生今世不聽從福音.
不認識神的人.
當然再加上是加患難給基督徒的人.
這是第十節.
「這正是主降臨.
要在他聖徒身上得榮耀」.
再一次肯定了在苦難裡面的基督徒.
在暫時在一般人的眼中.
何來有榮耀呢?.
這個宗教這麼弱雞.

$^{1001}$又手無寸鐵.
又任人魚肉.
任人宰割.
怎麼可能講榮耀呢?.
榮耀應該就是很強大的.
有軍隊等等.
不是的.
有的就是信念對神的國度的效忠.
所以就叫堅忍的信心.
或者加以文叫聖徒的堅忍.
Perseverance of the Saints.
就是「Tulip」最後的P字.
今晚給大家「Tulip」有五個向度.
如果你曾經讀過《圍讀聖經》的課程.
就有這個「Tulip」加以文的揀選論.
所以如果你是加以文主義的.
你就是無揀選.
你就是注定一定得到榮生.
不過有條件的.
就是聖徒的堅忍.
這裡告訴我們.
將來主耶穌要在我們身上得到榮耀.
在他的聖徒身上得到榮耀.
相信這是一個很大的榮耀.
你看我們這群錦繡堂的弟兄姊妹.
洪恩家的弟兄姊妹.
他們在患難裡仍然堅忍.
仍然沒有否定我的命.
而且誇耀他們是屬於主.
屬於基督的.
是跟隨基督的.
對主耶穌來說是很大的榮耀.
有一群人堅貞的跟隨你.
這樣的效忠.
所以今天我們其實不是說.
真的在乎太多人信主.
當然多人信主一定是好事.
但是香港有多少的基督徒.
願意為這個國度效忠.
什麼是這個國度?.

$^{1041}$怎樣效忠?.
我們都沒有興趣聽.
何況真的要去面對苦難.
而真正的效忠.
更加需要一些委身的基督徒.
不要灰心.
一定要盼望.
總是希望下一代.
興起一些工人.
興起一些年青人.
是對主的真道有心的.
這樣去為下一代的人祈禱.
我們這樣去帶著信心在神面前.
因為主耶穌將來.
總是找到一些人.
他得到榮耀的.
在他們的生命裡得到榮耀.
那你和我是不是呢?.
我們要自己問自己.
我們會不會有苦難.
我們就馬上走開呢?.
有榮耀.
你就要走過來呢?.
慢慢神都會改變我們的心態.
所以第五到第十節.
原來保羅告訴我們.
做物主原來不是不做事.
是暗中做事.
其實一直做事.
一直看著那些人是真正效忠.
所以他不容許苦難.
沒有了.
原來真是很有原因的.
原來人生沒有一些困難.
我們是珍惜的.
我們會覺得是很廉價的.
所以.
我這樣講例子.
也不是甚麼秘密.
北宣呢.

$^{1081}$時間到了就鎖門了.
你進不了是你的事.
不會稀罕.
當然他不是這個意思.
你敲門.
因為他還有下一堂的.
因為太多人了.
你不來就甚麼的了.
時間到了就鎖門.
施善就很乖的.
這樣.
問題當然是你不同的角度.
我們稍為站久一點.
別搞錯了.
每個人都不同.
當然好像開玩笑.
但在這段經文.
我自己是很深刻的.
我想我們為信仰.
是要建立一些受苦.
當然不是自己拿苦來受.
而是我們要建立一個.
為神國度效忠的基督教.
基督徒.
所以大家都要學習.
柴經班重要的地方就是.
可以有系統的教導.
一定會碰到這些題目的.
一定會碰到這些經文的.
你不能逃避的.
所以你要逃避就不要來柴經班了.
因為柴講到這些.
比較有深度一點的要求.
信仰對你對我的要求.
所以記住.
一直是神公義的判斷中.
哪些人是配得他的國度.
還有就是.
那天 記住那時.
那時即是說.

$^{1121}$原來今生沒有真正的報應.
或者實際的賞賜.
我們就不要灰心.
因為未必是今生就有即刻的判斷.
但是那時就肯定是最後的了.
肯定是神的公義的判斷.
就是會報應那些將患難.
加在基督徒身上的人.
即是說.
我們如果不是那些報應的人.
即是說.
我們在人生總是有為主.
我們要受的一些苦.
或者有些位置.
我們要為主守住的.
可能是信仰.
可能是見證等等.
記住了.
第十節就是說.
主耶穌回來的時候.
會在聖徒.
祂尋找的聖徒身上得榮焉.
即是說.
這班人聲稱自己是屬於我的基督徒.
我在他們的生命裡面.
我感受到很大的榮焉.
因為他們在苦難.
沒有否認我.
作為他們的救主.
作為生命的主宰.
所以我們很快.
11號就要洗禮了.
8月.
我們希望每一個基督徒.
信主的.
都有一個很堅穩的信心.
所以這裡告訴我們.
第十節的下半節.
又在一切信的人身上.
顯為稀奇的那個日子.

$^{1161}$我們對你們作的見證.
你們也信了.
信的人身上.
顯為稀奇的日子.
到底想說什麼意思呢?.
和上文說的榮耀.
再說上文的報應.
再說上文的.
配受苦難的國度.
所以你可以想像到.
保羅的意思.
就是說.
將來.
我們會覺得稀奇.
見到榮耀的主降臨.
我們會覺得很稀奇.
這個稀奇到底是正面還是負面呢?.
保羅沒有詳細說.
在一切信的人身上.
顯為稀奇的日子.
就是說原來我們在人生.
經過了一些困難.
為主的名我們去站穩.
我們也是很少的.
但是可能將來得到的賞賜.
稱讚榮耀是很大的.
大到傳到永恆.
我們真的會看為稀奇.
其實我們沒有做過什麼.
主耶穌.
我只不過做一點點事.
我為你的名.
我只不過是應份要做的.
為什麼要得到這麼大的榮耀呢?.
顯為稀奇的日子.
所以.
我們如果去一些重大的頒獎禮.
我們就會感受到那份榮耀.
那份榮譽.
任何人都想的.

$^{1201}$比如你喜歡看足球.
歐洲盃,國家盃.
現在明姐開心到跳起.
英格蘭也進入決賽了.
很開心.
然後就.
你看看你在什麼興趣裡面.
你認為什麼是最榮耀的.
比如我上個星期.
去錦西幼稚園的畢業禮.
整個大的劇院.
關了燈.
你站在上面其實很大的榮耀.
幾百人.
那裡已經是差不多一千人.
看著你站在正中間.
如果你是全場的焦點.
去榮耀你的時候.
是很榮耀的.
你可以想像得到.
將來主耶穌對這些堅忍的聖徒的讚賞.
主耶穌說.
我在你們身上得到榮耀.
你可以想像得到.
將來在永恆裡面的嘉許禮.
說得少了.
不是金銀的冠冕那麼簡單.
是做物主對我們的讚賞.
所以在一切信的人身上.
顯為稀奇的日子.
所以我們今生這麼短暫.
為主而活.
不要做物成為一個口號.
而是實際的行動.
我們真的要多走一步.
尤其是為那些未信.
不珍惜福音的人.
為他們流淚.
為他們祈禱.
這是信仰的種子.

$^{1241}$我們的眼淚就是種子.
可以令他們得到生命的改變.
尤其是更加和我們有關係的人.
所以保羅說.
我們對你們作的見證.
你們都相信了.
保羅自己就是一個見證人.
他自己就是這樣為神的國度.
周圍向不同的民族.
特別是因為歐洲那麼多種族.
都不是猶太人.
保羅都為這些外邦的人.
鞠躬盡瘁地傳耶穌的福音.
再沒有任何的種族.
帝國的偏見.
再沒有了.
你看保羅多麼豁達.
福音改變他的生命.
看自己是生育的一個執事.
你可以想像到保羅.
福音怎樣改變他.
不會看一個外邦人格格不入.
猶太人對外邦人是很拒絕的.
但保羅已經完全沒有這種味道.
完全對外邦人是很歡迎的.
內心擁抱這些外邦人.
所以很多年輕的外邦的傳道者.
跟隨保羅去傳道.
保羅的手下有很多年輕的傳道者.
來自不同的國家民族.
例如提波塔是安提阿路斯特的人.
即是土耳其人.
提多又是其他地方的希臘人.
你可以想像到還有不同的民族.
這就是保羅的見證.
你們也相信了.
保羅的見證.
接受了尼加城的弟兄姊妹.
相信和實際活出這個見證.
就是為主耶穌的福音受苦.

$^{1281}$所以如果我們建立了一班弟兄姊妹.
甚麼都不追求.
追求的是榮耀主和福音.
那個教會真是有希望.
榮耀主的福音.
就算自己有吃虧.
也要榮耀主的福音.
要休息一下 喝口水.
有沒有甚麼對於一到第十節.
一間很好的教會.
第五到第十節.
一位很好的神 賞善罰惡的神.
記住要報應.
保羅用的措辭是很重的.
所以不是沒有報應.
所以那些不信神的人.
不要以為自己很得意.
以為沒事 其實不是.
所以其實任何的講真理.
都是神一定會追究的.
我們最後講這兩節.
十一,十二節.
保羅再次講他們是很好的牧者.
很好的牧者一定是祈禱.
因此第十一節.
我們常為你們禱告.
願我們的神.
這個心願.
看你們配得過所蒙的恩兆.
配得過.
第二次頭尾呼應.
再一次你們本身已經.
在苦難裡面仍然堅忍.
已經證明你們是配有神的國度.
但保羅繼續為他們祈禱.
我繼續為你們禱告.
願我們的神.
看你們配得過所蒙的兆.
原來那個呼兆.
原來那個決志.

$^{1321}$是要讓人看到.
讓神看到你們是配得過.
這次就是那個地上.
有血有肉的一個人的生命.
在帝國.
在皇帝尼祿的迫不得已的社會.
政治的氣氛下.
仍然要站在神的國度這一邊.
這樣就叫保羅說.
希望你們在神的眼中.
你們是配得過所蒙的兆.
所以叫人信耶穌.
不是你信也好不信也好.
算了我都不要勉強了.
這個世界現在很講自由的.
其實不是的.
真的要讓他看到.
你信是一個很大的榮耀.
你是被選兆.
還有你還要用你的生命.
一生去印證這個呼兆.
是你配得過在神的眼中的一個基督徒.
這裡就是告訴我們.
然後又用大能成就你們一切所羨慕的良善.
和一切因信心所作的功夫.
用大能保羅說.
他的心願就是神的揀選.
在你們生命裡面被印證.
是配得過的.
另外就是神繼續用大能去成就.
你們所羨慕的.
這裡說到的良善.
一個基督徒的生命的內心價值.
很善良.
裡面的那種良好.
那種善良.
其實要反映神的屬性.
反映神的生命.
神的公義.
神的愛.

$^{1361}$神的全能.
在一個基督徒裡面.
聖靈在我們裡面的工作.
其實大部分都是讓你思想耶穌的愛.
耶穌的救恩.
令到你自己裡面就會有變化.
好像張無忌下去練那些《九音真經》.
突然間裡面打通了《淫篤二物》.
有個武功心法.
其實我們基督徒都是不斷思想耶穌的愛.
和十價的救恩.
聖靈就會在我們裡面工作.
令到我們不知不覺會有一種良善.
屬天的神的性情會在我們裡面去滋養.
我們會有一種良善.
所以原來江山易改.
品性是可以疑的.
信耶穌.
是真的.
有些姐妹都說.
我為甚麼這麼喜歡小朋友.
其實不是與生俱來.
信了主耶穌.
自然就會很自然.
有一個小朋友在學校畫了一張圖畫.
叫媽媽告訴我.
畫了一張圖畫.
我帶他幾十年.
又不是教他.
他喜歡將自己的圖畫和我分享.
小朋友就不能騙人.
他自己感受到牧師的親切.
聖文師喜歡靠近.
這是很真實的.
小朋友真的不會作利益等交換.
一定是內心裡面那種仁慈.
或者那份親切,愛.
諸如此類.
所以我們真的要為做兒童工作的牧者祈禱.
因為如果小朋友都感覺不到你的愛.

$^{1401}$其實與年紀無關.
與大年紀.
小朋友要喜歡.
如果老人家很親切.
基本上小朋友是很開心的.
我們在座的Alex.
做小朋友工作.
肯定小朋友喜歡他.
他才會繼續做.
肯定的.
沒得選擇.
你可以想像到.
原來要用大能的祈禱.
希望那個良善在我們裡面.
所羨慕的良善.
我們有沒有羨慕呢?.
羨慕主耶穌的良善.
我們去追求.
原來保羅說.
願神的大能.
用祂的大能.
去做你們在身上所成就.
你們所羨慕的一切良善.
以及一切因信心所作的功夫.
所以我們要用信心去做.
做在神的裡面.
記得一定要榮耀歸神.
所以做不到.
在哪裡跌倒.
然後在哪裡再跌倒.
都不要緊.
很多的工作.
很多的事工呼發不到.
不要緊.
繼續不要灰心.
繼續去做.
所以一切因信心所作的功夫.
叫我們主耶穌的名字.
在你們身上得到榮耀.
反過來.

$^{1441}$之前是主耶穌在你們身上得到榮耀.
現在是主耶穌的名字.
在你們身上得到榮耀.
不同了.
即是你們有良善的內心品格.
有用信心所作的功夫.
即是一直憑信心去做.
不要灰心.
不要一步十行.
默默地做.
不是要打鑼打鼓告訴別人.
而是你要做在神的角度裡面.
這樣主耶穌的名字.
在你們身上得到榮耀.
就是說.
原來真是印證你們這樣尊榮耶穌.
又得到榮耀.
「你們也在他身上得榮耀」.
今次再重複一次.
說「相輔相成」.
第十節已經說了一次.
現在又說兩次.
即是原來苦難只會使我們得榮耀.
原來阿保羅他真是很好.
你們有這麼堅忍的信心.
你們是配得神的角度.
但我仍然作為牧者.
阿保羅說我要為你們禱告.
願你們在神的眼中.
配得過所蒙的恩兆.
你們真是證明你們是被揀選的.
就在你們的生命裡面流露出來.
我用這樣的例子也應該是對的.
去奧運的參賽者.
每個都已經是最初的一歲.
在他的項目裡面.
例如跳水那個叫全紅禪.
有沒有聽過.
一個很窮很窮的女孩.
上一屆的東京奧運會.

$^{1481}$14歲.
她高台的下水.
是完全沒有水花的.
即是你把水放進湯裡.
也有水花.
她是高台跳下去沒有水花的.
而今年是17歲.
已經是世界冠軍.
是一個很窮很窮的鄉村成長.
很瘦弱的.
但是很神奇.
她能夠拿到這個跳水的金牌.
每一個應該是配得過那個呼召.
因為你一定是在那個項目裡面.
經過千除百練.
刻苦的訓練.
才有這樣的成績.
所以去奧運一定是很高的水平.
就配得著這個入場券去參賽.
所以基督徒.
我不知道將來有多少要被淘汰.
在錦書家我們都有過千的名冊.
有多少是配不得神的國呢?.
他得救上天堂.
但是不配稱為在神眼中的榮耀.
所以說的是榮耀這個部分.
在苦難裡面仍然堅忍.
所以我們是要求主幫助我們的.
以前有些古老詩歌說.
坐著巴轎上天堂.
甚麼都不用付出.
一信就行了.
判若完全沒有改變的那個人.
其實就不是了.
所以《普羅書信》裡面.
這一章其實是很重的.
說的是那個堅忍的品格要成長.
所以這裡告訴我們.
我們在主耶穌身上都會得到榮耀.
而主耶穌的名字.

$^{1521}$在我們的見證的生命裡面.
主耶穌又會得到榮耀.
所以都是照著神和主耶穌基督的恩.
原來都是恩典.
所以原來每個人都可以得到.
你說你走那條路當然不是了.
你走那條路是特別隊.
我們才是慘.
原來不是.
大家都是這樣得到神的無條件恩典.
福音.
每個人都是這樣聽福音.
每個人都是因為這位神並主耶穌的恩典.
每個人都是一樣.
為甚麼有些人會成長?.
有些人那麼愛主.
有些人對福音完全若即若離.
根本完全沒有興趣.
這個原因就是我們對恩典的理解.
是那麼深入.
所以這個都需要神給我們改變.
所以這兩節是講保羅的禱告.
保羅對貼撒羅尼加城的弟兄姊妹的祈禱.
願他們繼續蒙受那個良事.
和他們用信心所做的功夫.
可以在神的面前傳得住.
另外叫主耶穌的名字.
可以在我們的身上.
貼撒羅尼加的弟兄姊妹的身上得到榮譽.
所以你可以說這個祈禱.
是提升門徒對信仰的要求.
保羅原來在訓練精英的基督徒.
所謂精英不是甚麼階級觀念.
每個人都是這樣蒙受那個恩典.
但有些人對恩典珍惜.
不用他的人生為主而活.
只是掛名的基督徒就有很多.
但有些人認定.
我們跟隨這位愛我們的主.
祂要我怎樣我就怎樣.

$^{1561}$因為這個世界不是由地上的人主宰.
是由造物主來說領著.
記著這章是說那時主再來的時候.
祂就會賞善罰惡.
和真正的永恆的焦點.
那個中心其實原來就是主耶穌.
因為那時是主耶穌回來.
帶著有能力的天使.
那個時候就是賞善罰惡.
是最終的終極的公義的裁決.
所以原來從來焦點的中心都是耶穌.
不過現在的人不承認.
所以世界都是這樣亂.
沒有一個國家可以一錘定音.
叫停止不准再打仗.
根本沒有人理會.
繼續這樣來仇殺.
但是神看著的.
神管理著的.
但是從來那個中心是耶穌.
整個將來的世界都是主耶穌.
是焦點的所在.
所以如果我們今天整個人生.
就是為主就是我們的焦點.
我們就一直走那條向實的路.
一定會在我們的生命裡得到榮耀.
耶穌也會在我們的身上得到榮耀.
第一章說到的.
保羅對貼城教會的勸勉.
簡單來說一到第四節.
很好的貼城教會.
第五到第十節.
這一位非常公義的神.
他有他的判斷.
所以我們這麼有限的人.
不要只是看五年,十年.
就覺得沒有審判的.
就不要這樣看.
我們一生追求的就是神的榮耀.
福音被人認識.

$^{1601}$主耶穌被人高舉.
這就是我們要做的.
有沒有榮耀?有沒有賞賜?.
有沒有立刻要報應?.
那個就不是我.
我們不要扮演上帝.
我們記得我們是凡人.
我們去倚靠神.
第十一到十二節.
就是為什麼信徒這麼好呢?.
有好的牧者.
有好的眼光.
他仍然為他們祈禱.
願他們可以得到.
配得過所夢的照.
很重要.
你為這樣弟兄姊妹祈禱.
我願你.
願阮妮,願Della,願Alex.
願你配得過所夢的照.
這個祈禱一定是和他現在的狀況.
為他們更加的順靠神去禱告.
這些都是很重要的.
我們也要為一些.
事奉很久的弟兄姊妹禱告.
願他們能夠繼續.
配得過所夢的恩照.
又願他們那個品格的良善.
不斷的增長.
神的大能.
使這些良善增長.
他們那個信心所作的功夫.
沒有說哪些功夫.
但總是為神.
憑信心去做的功夫.
繼續的來到去.
就是這樣增長.
所以這裡就是.
有很好的牧者.
看到的是神的榮耀.

$^{1641}$在教會需要神的榮耀.
在教會的弟兄姊妹.
需要有這樣的良善的品格.
和充足的信心所作的功夫.
還有去羨慕.
那個目照帶來的.
那種配得的失明.
這就是保羅的禱告.
好好的牧者就為教會守望禱告.
所以真的.
弟兄姊妹用一些時間為教會祈禱.
真的可以讓一些事情.
在我們能夠有這種.
因為祈禱會使我們的眼界開啟.
我們的眼界越來越有神的角度.
我們其實可以再闊一些看事物.
大家繼續下去第二章嗎?.
可以做一些時間的.
我們用自己的時間.
今天有些乾了.
第二章立即說.
又再說主在內.
不過這次就一個比較負面的.
就是到底.
為什麼會有些人.
有些錯誤的神學呢?.
保羅也有一個幫他們矯正.
第二章.
頂門論道我們主耶穌基督降臨.
和我們到他那裡聚集.
又是再回應第四章.
降臨主在內的次序.
我勸你們.
無論有靈有言語.
有夢有明的書信.
說主的日子現在到了.
如果今天有些不知名的短信.
有些WhatsApp.
你就不要聽了.
似乎當時已經湧現了.

$^{1681}$這些人傳人的訊息.
以訛傳訛.
或者有些人招搖撞騙的.
假的使徒.
去冒充.
說很多這些危險的主在內.
所以正牌就一定會有A4.
開始有人猜測.
來估計.
和製造很多駭人聽聞.
關於主的在內等等的東西.
說主的日子現在到了.
沒錯了.
真是這樣計算.
用了不同的計算方法.
這樣計算.
你就不要輕易動心.
也不要驚慌.
對的.
真的會突然很雀躍.
如果那個人有份量的.
說得這麼言之鑿鑿的.
你就更加動心.
但是反過來又會很害怕.
因為他說得這麼實在.
什麼都未準備怎麼辦呢?.
第三節.
人不驅用什麼法子.
你們總不要被他誘惑.
因為那日子以前.
必有彌度反教的事.
並有那大罪人.
這裡又真的在說那個迪基督.
就是沉淪之子.
顯露出來.
明顯是在說一個人.
所以我這幾時說的那個迪基督的.
邪惡的三位一體.
那個叫海獸.
其實真的有些根據的.

$^{1721}$所以這裡是在說那個大罪人.
大家如果看那個空少.
我很小的時候看的.
那時未信耶穌.
說那個666.
那個鬼仔.
是很正常的人.
後來成為了世界的政治家.
他就是那個666.
他就是那個魔鬼化身.
當然是演戲.
根據聖經的某程度的橋段.
當然他在醜化了聖經.
你可以想像得到.
這個迪基督.
就會怎樣呢?.
有些行徑的.
他是抵擋主.
這個叫大罪人.
他叫沉淪之子.
耶穌基督就是那一位的光明之子.
跟隨他的就是光明之子.
耶穌基督就是那一位.
常經幽惑的幽惑之子.
Man of Sorrows.
二賽書53章第三節.
這些是沉淪之子.
Demnation.
Demnation就是注定滅亡的意思.
滅亡之子.
保羅用這些話都相當蠢.
到底這個人是甚麼人呢?.
他要甚麼化身呢?.
我們不知道.
但無論如何.
這裡保羅說的是一個人.
在這段的上下文.
很肯定是說一個人.
他的行徑就是他是抵擋主.
高抬自己.

$^{1761}$超過一切稱為神的.
和一切受人敬拜的.
甚至坐在神的殿裡自稱是神.
將來的基督就是這樣的行徑.
我還在你們那裡的時候.
曾把這些事告訴你們.
當然有些歷史觀的解經家.
認為是離陸.
不過說得暗晦一點.
但如果這樣對號入座.
那經文的世界就不得全面.
這樣理解就會有一種強解.
你強行將那個人加入去.
就當是那個「迪基督」那個大罪人.
那個將來的應驗就很難實現.
你可以說他是類這樣的行徑的「抵擋」.
但是離陸包括最後的希特勒都不是.
一定是有他的那個行徑「抵擋神」.
但是真正的「迪基督」是在將來.
這裡告訴我們他自稱是神.
我還在你們那裡的時候.
曾把這些事告訴你們.
你們不記得嗎?.
現在你們也知道那拿了他的是甚麼.
是叫他到了的時候才可以點路.
因為這段就開始難解了.
因為若有所指.
但是又說得比較隱晦.
可能保羅自己也不知道.
所以他只能夠隱晦.
這個不法的隱義已經發動了.
只是現在有一個拿了的.
等到拿了的被除去.
那時這個不法的人就必顯露出來.
主耶穌當然就是說主在來.
就要用口中的氣滅絕他.
用放立的榮光廢掉他.
這個不法的人來果然.
所以他是真是鬼仔.
是照撒旦的運動.

$^{1801}$行各樣的義能,神蹟.
記住這個不法的人是不斷地出現的.
是說那個「迪基督」.
是一個人來的.
不是一隻鬼或者靈界.
但是記住是有撒旦的運動.
並且在那沉淪的人身上.
行各樣出於不義的鬼咒.
他會行各樣的義能,神蹟.
和一切虛假的歧視.
也就是說讓那些沒有真理的人的心靈.
很容易會煽動和吸了去.
因為如果你沒有真理去維持.
你自己的信仰的辨別力.
你就會被這些虛假的歧視.
保羅看這些是各樣的義難,神蹟.
真的做到的.
但是原來那個心態,那個價值觀.
那個出發點完全不是為了榮耀耶穌.
是為了使人跟隨魔鬼.
是兩個很不同的對比.
所以我們不要看到神蹟就馬上開心.
因為原來有些可以是假冒.
所以這裡告訴我們.
告訴我們第十節說.
在這些沉淪的人身上.
行這些出於不義的鬼咒.
因他們不領受愛真理的心.
使他們可以得救.
原來我們不愛真理的心.
就令到你很難對信仰有一個心悅誠服的明白.
而且還去到一個階段.
原來是會有人搶你的心靈.
進入魔鬼撒旦的行列.
他會用一些神蹟奇事.
令到你受迷惑.
所以保羅告誡.
接受萊尼加城的弟兄姊妹.
要小心.
無論有不明的書信.

$^{1841}$或者有其他的言辭.
不要輕易動心.
要留意,要判斷.
不要輕易動心.
大家說明白了.
廣東話叫「吶反應」.
你都沒有想清楚就「吶反應」.
輕易動心.
很容易就有這種不設防對真理的捍衛.
在自己的心靈.
第十一節.
故此神給他們一個心發錯誤的心.
所以不愛真理.
只會傾向於錯誤的.
叫他們順從虛荒.
在一切不信真理.
賭喜愛不義的人.
都被定罪.
一直都在說要小心.
要留意,要提防.
一直都在說其實透過負面.
有些人被完全吸了去.
被撒旦的運動吸了去.
被這個不法的人帶了去.
但原來在另一個角度.
保羅說其實是他們不愛真理.
其實是他們愛不義.
其實是他們被定罪.
是因為他們內心已經有一個心發錯誤的心.
內心已經錯誤.
錯誤就很容易物以類聚.
那些事情聽起來都好像很吸引.
但想心裡好像有些事情不對.
譬如美國有人拜撒旦教.
拜撒旦.
以撒旦那種的敵黨作為他們的教義.
總之聖經說的良善,寬恕.
他們就反過來.
我們要凶惡,我們要報仇.
我們要以顏頑固.

$^{1881}$真是實行教的教義.
因為我讀神學的時候就做了這個功課.
老師真是很好笑.
叫我找撒旦教的人跟他們對話.
怎樣對話?.
去被人容了都有份.
而且每個人都不是凶神惡煞.
有些律師,有些專業人士.
美國真有其事,真有其人.
而且是合法的宗教.
拜撒旦.
所以你可以想像得到.
因為基督徒的基督徒.
要上車嗎?.
OK.
我們完了.
我們祈禱了.
多謝主讓我們能夠在神的恩典裡.
我們得到這個永恆的救恩.
盼望我們能夠心靈裡.
保羅第二章就很重的說.
主在來不是一件開玩笑去計算的事.
是真的會有敵黨.
會有撒旦的運動.
尤其是對基督教的攻擊.
會對基督徒帶來一個.
可能真是一個很嚴重的判斷.
所以我們求主讓我們要持守真理.
一個持守真理的心才是最重要的.
願主幫助奉耶穌基督的名.
阿們.
謝謝大家.
\newpage



\section{帖撒羅尼迦後書}
\label{sec:ntrA1aTVgRk}
\textbf{2024-07-18 周四晚查經班《帖撒羅尼迦後書(2)》 講員:曾裔貴牧師}
\newline
\newline
連結: \href{https://youtube.com/watch?v=ntrA1aTVgRk}{\texttt{https://youtube.com/watch?v=ntrA1aTVgRk}} ~~~~ 語音日期: 2024-07-19
\newline
\newline
\hyperref[sec:uJ_qz3eVEf0]{\small{< < < PREV SERMON < < <}}
~
\hyperref[sec:index_chronic]{\small{[返順時目]}}
~
\hyperref[sec:index_scriptual]{\small{[返順卷目]}}
~
\hyperref[sec:ewzMm4Y0R5I]{\small{> > > NEXT SERMON > > >}}
\newline
\newline
$^{1}$我們一起祈禱.
再一次感恩.
我們能夠有機會去看.
《貼壽羅利加前後書》.
我們都很想去了解.
使徒保羅在面對這麼多教會的事務.
和外面的一些困難.
明白阻礙福音工作.
但是他仍然能夠在神的豐厚恩典裡.
他有這一份心智.
還有很多身邊的年輕的牧者.
跟隨他,支持他.
和跟他同心同功.
我們真的很盼望主的聖靈.
讓我們能夠透過經文.
我們可以跨越一些時代.
跨越一個時空.
我們能夠再一次掌握.
上主在地上.
在不同的時代.
你仍然坐著為王掌權.
叫我們有這一份信心.
的而且確我們每天都有大小的新聞.
或者一些事情.
會在我們身邊發生.
或者離開我們比較遙遠.
我們未必能夠掌握.
或者能夠去改變.
但是深信我們這一位主.
是賞善罰惡的神.
好像我們上個星期看過.
主耶穌將來就是那一位.
執行審判的主.
在這本書都有很多關於這方面的教導.
原主都給我們有一個警覺的心.
對於你的再來.
我們真是不是一種.
只是一個信念.
而是影響著我們大家的生活.
真是以這個主回來的日子.

$^{41}$作為我們生活的動力.
和人生的一個很大的盼望.
的而且確我們生活上.
是有一些更加緊迫.
每天都發生的事情.
排滿了在現代人的行事曆.
但是我們真是很盼望主.
給我們整個的新心靈.
都完全仰望在主的再臨.
世界的局勢改變.
我們仰望等候主.
奉耶穌基督的名.
阿們.
好 我們再簡單的溫習一點.
這個第一章我們曾經講過.
就是.
我們講過就是.
《索羅尼加後書》.
就是仕途保羅.
他寫了前書之後.
似乎又是提摩太回來.
告訴他現在這封信帶來的影響.
那個改變.
所以保羅亦都在寫這封後書的時候.
再勉勵弟兄姊妹.
但是當然一定是在第二章.
有些我們上星期稍為讀一讀.
都讀過了.
有些已經是冒充保羅的書信.
來開始湧現在初期教會.
似乎《索羅尼加後書》.
都是保羅擔心.
甚至可能已經真是有那些書信.
冒名的書信.
來影響順途.
所以保羅就需要.
來寫這封後書.
來勉勵《索羅尼加後書》的弟兄姊妹.
要留心.
到底講甚麼呢?.

$^{81}$有甚麼危言聳聽呢?.
有甚麼錯誤的教義.
去影響順途呢?.
我們在第二,第三章.
就會再有很詳細的交代.
第一章就是保羅再一次.
來勸勉教會.
他們其實非常好的.
那個好就是在迫不患難當中.
營救全一個的信心.
和這份彼此相愛的心.
這是第四節.
三到四節.
到第五節就已經說了.
這一位的主.
我們這一位的神.
是真是那一位認真看待人類.
每一個的行動.
每一個的人生.
到底是否真是為公義.
為真理而活.
所以我們一直都盼望的.
那個所謂的公義的裁決.
絕對不是現在.
有一天.
所以後書第五節.
一章第五節一直到第九節.
就是說公義的神.
或者說這一位.
真是看著宇宙萬物發生的神.
將來他就會有一個賞善拔惡的時間.
所以你可以說對基督教來說.
是最大的一個信仰.
不信的人一定會嘲笑我們.
將來可以審了.
可以分了.
的確經文告訴我們.
將來就是神會賞善拔惡.
所以本來就給初期教會.
初信的貼索羅尼加教會.

$^{121}$有這個定心丸.
有這一份的安慰.
大家可能如果沒有再重溫.
上一份我說得比較急和快.
我怕大家吸收不到.
第五到第九節.
這裡說.
首先就是因為這間教會實在太好了.
在混亂當中仍然都持守著這個信心.
所以他們的信心是格外的增長.
愛眾人的心又是充足.
很值得保羅去稱讚教會.
他們面對這些困難.
不會雞飛狗走.
不會因為一些謠言.
來去分化.
而是很專一的愛主愛人.
當然這個信心.
包括他們為主耶穌作見證的福音.
所以到第五節.
將焦點放在神這位造物主.
這位真真正正有權.
去賞善拔惡的神.
第五節.
這正是神公義判斷的明證.
叫你們可算配得神的國.
即是說他們在患難中.
真金不怕紅爐火.
他們的信心.
他們的堅忍.
證明他們真是神的揀選.
所以不是因為你的行為.
是因為你的信心的表現.
正正就是神的揀選的明證.
是一個Witness.
這裡告訴我們.
他們可以配進入神的國度.
你們就是為這個國受苦.
所以要搞清楚.
我們為神的國度受苦.

$^{161}$相信人類最高的情操.
一般人都會有一種國家觀念.
愛國.
所以基督徒就是神的國度的指紋.
當然你不要極端化了.
你只是做神的國度.
不做自己的地球上其中一個國度.
譬如你是法國人.
你不認自己是法國人.
這樣是不行的.
最近有一個新聞.
就是阿根廷的球隊贏了冠軍.
美洲盃.
在旅遊巴那裡.
唱一些帶著種族歧視法國的歌詞.
現在就越來越大風波.
國際足協都介入調查.
當事人本身在英國踢球.
所以現在有很大的風波.
就是阿根廷的總統支持這個球隊.
不用道歉 不要道歉.
所以有時牽涉到地上的國度.
其實都複雜的.
但基督徒我們有所謂雙重國籍.
我們既是地上的公民.
所以我們要行公義.
但更加要記得.
這個就比較抽象.
很少基督徒會經常.
除非有問題.
否則一般我們很少在日常生活.
會經常意識自己是天國指紋.
我們是為神的國度受苦.
但《少途寶來》真的講了很多.
講基督徒當時的教會.
其實是屬神.
不是屬於哪一個人 哪一個群體.
所以你們就是為這個國度受苦.
再一次提升他們的層次.
看表面的受苦.

$^{201}$當時很明顯就是尼祿.
還未到最白熱化.
就是當地的地方的.
可能是匪徒.
類似現在的黑社會.
就結合猶太人.
要逼迫這些基督徒.
如果你公開說耶穌是這位彌賽亞.
明明祂只是普通猶太人.
對一些猶太教的人來說.
是很大的褻瀆.
所以如果這些外邦的.
貼索羅尼加教會的人.
他們都這樣宣稱.
就真的會帶來當地的猶太教.
那裡有一個比較大的猶太教的組織.
貼索羅尼加有的.
所以肥臘比反而都有.
就會帶來當地人.
可能是有錢有勢力的.
就會結合那些匪徒.
來逼迫保羅和跟隨保羅的人.
所以保羅就很想告訴他們.
第五節這一段是很重要的.
你們不單是算是在這個國裡面.
而且你們現在是在為神的國度受苦.
這樣提升其實是很高的尊榮.
因為原來你的受苦是為神.
是為神的國度.
沒有東西是再有高價值過這個國度.
神的國度.
所以教會如果能夠.
弟兄姊妹有這個意識.
教會就無論在甚麼處境都不需要害怕.
因為就算有沒有傳達人都不要緊.
因為所有弟兄姊妹都已經有這個意識.
我是神國度的子民.
我所做的事是要榮耀神的.
是要為這個國度去打拼的.
當然不是用刀槍.

$^{241}$是用愛和信心.
其實就是福音.
所以保羅這樣來提升.
對他們來說是一個很大的安慰.
即便有苦難.
但你為神就很不同了.
所以第六節.
神既是公義.
這裡就訴諸於神的屬性.
所以就是一個受苦的教會.
到底你是在面對.
那些逼迫在你面前.
還是逼迫都改變不了的.
就是你在做物主.
公義的神的面前.
去生活作見證呢?.
兩樣東西同時也存在.
的而且確.
真是他們要面對的是逼迫.
不是打的.
但是原來更加要看到的就是.
神是一直都看著地上的教會.
神既是公義.
就必定會將患難報應.
那加患難給你們的人.
所以是很實實在在的賞善罰惡.
保羅沒有什麼需要懼怕.
就是他真的將頂子妹.
帶到審判裡面.
帶到審判的角度的報應觀裡面.
所以你就想像得到.
那些頂子妹就有很堅強的心靈.
因為我現在為誰受苦.
是為保羅?不是的.
是為神的角度.
那誰逼迫你呢?.
逼迫你那個一般人.
是一個很橫瀾的尼祿.
或者有些匪徒.
但是誰真正去審判呢?.

$^{281}$就是這位公義的神.
接著就是七字.
必定會使你們這些受患難的人.
與我們同得平安.
你可以想像得到.
將來就會有一個安慰.
將來就會有一個真正賞賜,冠冕.
所以既是要知道我們為誰而活.
我們為誰而戰.
戰也是對的.
不過是心靈的用十字架的福音.
去得到人的心靈.
所以現代就不是用刀,槍.
絕對不可能是一種武裝革命.
譬如中美洲的解放神學.
就真是鼓勵教會要有刀有槍.
去推翻那些邪惡的政權.
也都是因為很貪污.
中美洲.
所以就會有這些革命神學.
但是聖經從頭到尾都在說不是這樣的.
所以一直都是訴諸於神的拯救和公義的裁決.
所以保羅說叫你們這些受患難的人.
當然是將來與得平安.
至少你知道你站在神的陣營.
你就能夠有這份平安.
但是這份平安將來就會轉變成賞賜.
八字.
就是說那時.
相信一定是說主在來.
這個那時我們上星期有提過.
那時第七節的中間.
主耶穌與他有能力的天使.
從天上在火焰中顯現.
所以那個將來永恆的代理人.
一定是主耶穌.
所以是最重要的.
所以現在弟兄姊妹.
如果你都不向耶穌祈禱.
現在你都不去敬拜聚會.

$^{321}$即是聖餐去思想主耶穌.
我暫時都比較少見到.
有人思想主耶穌思想到哭.
在聚會裡面.
其實是應該的.
十字架一哀這麼大.
有時你給你一些偶然間的感觸.
是很正常,很好的流淚.
如果我們能夠在聚會裡面有聖靈感動.
應該有很多弟兄姊妹會感動流淚.
但記住將來主耶穌就是唯一.
神安排的代理.
就是審判活人死人.
所以保羅再次提到.
主耶穌與他有能力的天使.
即是加伯尼,米加勒.
這些是我們在聖經舊約裡面.
有出過的名字.
但還有很多可能不知名的.
是甚麼階級的我們都不知道.
但無論如何那些是力體.
全部都是受造之物.
要伏在主耶穌的大靈底下.
主耶穌就是那位萬主之主.
萬王之王.
所以今天信耶穌.
應該我們初信.
很怕別人說你是耶穌.
信耶穌.
但如果信得久.
你還很怕連祈禱.
謝禮也不敢.
半秒就馬上閉眼.
很怕別人看不到你.
不知甚麼原因.
但也不好.
因為耶穌就是將來審判的代理人.
祂是很榮耀地降臨.
所以如果我們今生.
我都不覺得信耶穌是榮耀.

$^{361}$不覺得自己有榮耀的身份.
這是一個真正的問題.
你不要只是受浸那天.
你才覺得很榮耀.
然後你已經不再覺得身份是榮耀.
這是很大問題.
教會也是一個問題.
就是那個模樣為何會有這麼大的落差.
很多浸禮完之後.
真的不再回教會.
這是很大的問題.
所以你想想.
將來主耶穌就是那個代理人.
你跟祂完全沒有關係.
那怎麼辦呢?.
祂要不要審你呢?.
你能不能通過祂的審判呢?.
這些是真的要問的問題.
所以我們很認真地一起去研讀聖經.
我們就會有很多機會去探索到.
關於主耶穌的身份.
不只是恩信稱義這麼簡單.
主耶穌和祂有能力的天使.
在火焰當中顯現.
不再是彩虹.
不再是這麼浪漫.
是在火焰當中顯現.
真是要用祂的公義來審判世界.
應該這麼說.
祂的再來.
再沒有恩典了.
再沒有給罪人悔改的恩典.
對於我們信主的人.
就是一個所謂尋冤得雪.
真是我們苦盡甘來.
真是我們受的苦.
得到主的獎勵.
所以再一次提到報應.
已經是第二次提到.
第八次.

$^{401}$要報應那些不認識神.
和不聽從我主耶穌福音的人.
即是不認識.
拒絕認識.
詐假意義不認識.
明知而不認識.
迫不得已,嘲笑.
但也不接受福音的人.
福音去到他的生命.
他要以輕心.
對他來說是一個廢話.
是完全沒有關係.
我不做誰給我們吃.
很多人會這樣.
似乎很肯定地說.
殊不知他沒有想到.
人生誰給我們有生命氣息.
那個真正的本源.
我們都不去追求.
尋求明白.
很多在一個這樣的無知裡.
世界差不多有八十億人.
但世界一直在運作.
嚴格來說.
沒有神,沒有主耶穌.
世界都在運作.
很多人從來沒有聽過關於耶穌.
他一樣是依食住行.
一樣這樣生活,讀書.
很聰明.
但最大的缺失.
就是對自己的存在.
一無所知.
自己的本源.
將來要去哪裡.
完全都不知道.
所以基督徒有一個責任.
我們自己要活好自己的生活.
但我們要將我們所擁有的.
屬天的盼望.

$^{441}$很有意識地告訴別人.
讓人看到這個人真的有點不同.
和我們一樣這樣生活.
他不會特意地心吸力對他沒有影響.
他可以漂浮.
不是都是這樣生活.
我們一樣面對生老病死.
但基督徒應該有這個.
屬天永恆的.
因為呼召而來的生命.
和有永恆的盼望.
和對主耶穌應該有一份心悅誠服.
越來越愛他.
越來越因為認識他而愛他.
願意效忠於他.
所以很多弟兄姊妹其實很好的心態就是.
牧師你不要跟我說這些東西.
不用說多謝的.
我們為神做事是應該的.
這些就是正常正確的價值觀.
基督徒的價值觀.
牧者找你去事奉.
而你又真的可以做到的.
基本上嚴格來說.
當然人與人之間這些基本禮貌當然要說.
但嚴格來說弟兄姊妹說.
你不用說多謝的我們為神做的.
其實是對的.
這個態度本身就是應該的.
所以你也告訴我們要報應這些人.
神的報應.
主耶穌的報應.
主耶穌一直讓我們看到.
在地上被人嘲笑.
侮辱.
毫無能力.
完全是一個笑話.
一個笑談.
這個很難說.
但是原來將來他有這樣的能力去報應.

$^{481}$即是說他不是一種黑社會的報仇.
完全失去理智.
不是.
是公義的審判.
公義的主.
即是.
啟示所形容他是有七眼.
有七靈.
甚至遍體都是眼睛.
即是每一個方位.
每一個寸.
每一個角度.
是他看不到的.
做物主就是這樣.
就報應.
接著就是說那些受報應的人.
他們要受刑罰.
就是永遠的沉淪.
所以是很清楚將來的永恆結局.
不是說一些似是而非的東西.
是很具體.
是有這個日子要來.
所謂的永遠沉淪.
就不是你有你.
我有我.
我們在地獄那裡.
多開心.
沒有了那個做物主.
再在那裡搞事影響著我們.
我們以為是.
那個變成天堂.
我們在神的角度就變成地獄.
因為神每天都和我們一起.
一般的人就會這樣看.
說笑就是這樣說笑.
但這告訴我們.
真正的永遠沉淪.
是離開主的面.
和他權能的榮光.
即是有一個地方再沒有做物主的Blessing.

$^{521}$再沒有做物主的同座.
那個地方就稱為永恆的沉淪.
地獄.
火活的地方.
不是那些未信的人.
不信的人要消失在永恆.
而是在永恆裡面繼續存在的受苦.
所以這個是我們都不能夠用.
我們現在物質世界可以理解.
可以明白的.
有時卡通片就會反映到的.
但始終都不行.
這就是十節.
「這正是主降臨.
要在他聖徒身上得榮耀」.
所以報應和榮耀是兩個極端的.
我們一直等待的是審判帶來的榮耀.
我們經歷了那些苦難.
我們為神的角度.
所以我們真的要有我們自己.
不同的程度的對上主的角度的衛生.
你一定要找到你自己的崗位.
你自己一個的召命.
是神安排你的.
你要為這個角度的召命去衛生.
甚至死而不足惜的服侍.
真的要這樣為神去做.
因為將來.
你想想.
第九節就是那些不信拒絕福音的人.
第十節是我們這些為這個角度效忠的人.
我們就會在他的聖徒.
即是他揀選的人身上.
這個審判帶來的.
原來就不是我們要受虧損.
反而是得榮耀.
他要在我們身上得榮耀.
我們因為他的再來.
我們又會得榮耀.
這個就是你可以想像得到.

$^{561}$多麼的發酸.
我們還有十多二十年命.
希望不要那麼長.
他們九十歲就真的還有四十年.
三十多年都很長.
當然不是我控制.
但其實你想想.
三十年.
然後去到永恆.
那個時間真的很懸殊.
所以如果我們仍然還未找到自己的召命.
其實就有點可惜.
第二,第三章開始說工作倫理.
就真的探討到我們實際的生活.
因為主在來.
原來有些事故.
原來有些錯誤.
因為第二章有些假冒保羅的書信.
所以原來保羅都真的意識到.
似乎講主在來.
就會令到一些基督徒.
會受到一些錯誤影響.
所以保羅就要更正.
所以正確的工作倫理.
就在第三章就會再出現.
所以怎樣令到這件事平衡呢?.
需要透過保羅的教導我們去明白.
所以這裡告訴我們.
我們是一班將來要得榮耀的人.
這個將來得榮耀.
就不是你坐在那裡等見主.
那就是你會得榮耀.
一定是在說你怎樣能夠.
用你的人生為主的角度.
說了每個人都自己去面對.
每個人都自己去審視.
自己有多少參與.
當然有很多人參與.
就是在教會裡面搞一個陣.
影響了士氣.

$^{601}$大家不要追求不要事奉.
一回來就搞是非.
都是他們的參與.
不過是一個惡性的參與.
記住真正是要將來.
神才會審判.
所以洗人和睦的人.
我越來越深體會.
洗人和睦的人才是神的兒子.
這個登山寶訓的八福.
是很重要的.
因為這個地方永遠都是說.
大家彼此相愛.
大家都是同心去服侍.
大家都是義工.
既然是義工.
自然每個人都喜歡說甚麼就說甚麼.
那就變成很多人性在那裡發生.
就會令到真正想服侍的人.
有時會心灰意冷.
所以這個就要很小心.
記得我們都不是用我們人間的血氣方法.
仍然都是交託給主.
因為主仍然看著教會.
所以就會在一切信的人身上.
顯為稀奇的日子.
就是將來原來每一個弟兄姊妹.
都會得到自己神審查之下的榮耀.
記住我們是期望得榮耀.
當然我們香港很多運動員.
真的從來沒有想過.
可能個別有幾個.
他們想過一定我可能不是在一就忘二.
但我想大部份的參加奧運.
已經感覺是一個榮耀.
他都知道自己根本跟水平還差一截.
個別有些可能劍擊.
現在何詩蓓游泳.
可能有機會爭奪奧運金牌.
但其他的運動.

$^{641}$我雖然不是這方面的專家.
但大部份我想他們去都是想體驗.
但基督徒都要小心.
我們講明將來是得榮耀.
所以我們的人生應該期望.
我們在神的眼中得榮耀.
我們不要做給人看.
但我們要在神的眼中.
期望我們將來是得榮耀.
《哥羅不辭書九章》.
保羅應該有29次.
我們在場上都跑.
但得獎賞的只有一個人.
但我們都努力地跑.
不要叫你們都得獎賞.
所以正正就是保羅所講.
每個人都有自己的金牌.
就是神怎樣看我們.
下面我們為你們作的見證.
你們都相信了.
所以保羅一直緊守著這個信念.
所以他現在會到處去不同的國家.
在歐洲傳道.
一方面就是神很清楚帶領他.
由馬其頓異象.
要他去到歐洲.
由土耳其過對面的海.
去到歐洲.
在菲律賓上岸.
一直掃到雅典.
一直在整個歐洲的南部.
東南部.
去沿途的傳福音.
接著就進入.
希望他去到西班牙和羅馬.
這個就是保羅他自己很清楚.
神給他的那個召命.
那個差遣.
所以他就很清楚.
自己要向這些外邦人傳道.

$^{681}$其實很不簡單.
他能夠自服自己的心.
對於自己是猶太人的那些優越感.
已經完全沒有了.
你想想多麼厲害.
一個保羅.
學貫中西對於猶太教.
是非常有權威的一個解經.
一個律法師.
但是在他的人生.
自從成為主耶穌的僕人之後.
你聞不到他有任何猶太教的氣味.
完全沒有.
但是當他一回到耶路撒冷.
第三次旅行報道.
《少林寺》21章.
他一回到耶路撒冷.
就立即去做一個還願的祭.
因為他是猶太人.
他上到耶路撒冷.
就去到聖殿還願.
似乎是一個拿西人的願.
他自己原來也有這樣的謹守.
但是在外邦人的眼中.
完全聞不到猶太教的味道.
即是說不會令到外邦人感覺.
一種冒犯.
一種厭惡.
不知道大家有沒有試過.
我很深刻.
很早期.
我的妹妹和我.
就去了天主教教會聚會.
後來應該是她去了荃灣堂.
後來荃灣堂就差遣一些人去直堂.
就去了天水圍直堂.
她初期在那裡聚會.
後來她想走.
想轉其他地方.
就問她為什麼.

$^{721}$她說那些人在我面前撩鼻屎.
很噁心.
崇拜.
她先生就說.
我們有奉差遣來事奉.
那個基層的地方.
當然是這樣的.
人家是新移民.
那些婦女沒有什麼學識.
又不是什麼雙大雅.
當然是雙大雅.
是很噁心.
但是你做秘書.
你每天回中環穿衣服.
但是你應該有這個差遣.
有這個意識.
我就是想說.
我們每個人都一定有自己的感覺.
別人對你的厭惡.
或者我們對別人的厭惡.
但是保羅真的沒有他自己猶太教的.
我和你有些距離.
如果我們作為傳道.
別人感覺我們有距離.
其實這也是我們要學的通俗.
我很深刻楊蒙古牧師.
他說如果你能夠連一個小朋友都喜歡你.
聽你說話.
聽得很清楚很明白.
他都聽得明白.
你這個牧者就做得很成功.
因為小朋友都聽得明白你說話.
證明你真是有意識地說一些跟他的語調一樣的東西.
很深刻的.
這些牧者是.
你想想楊牧師是多麼有學問的一個牧者.
他的過來人.
所以為什麼耶穌喜歡小孩呢.
真的.
如果你是一個傳道人.

$^{761}$小朋友見到你就走這麼遠.
你要問自己為什麼小朋友不喜歡你.
小朋友是很憎恨的.
他不會講假的.
所以這裡告訴我們.
我們要學習.
我們對你們作出見證.
你們不需要信任.
這個就是保羅看到.
這位做牧主.
原來一直看著整個地上的事情.
沒錯了.
2000年過去了.
一代一代的梟雄.
碧北教會的人都已經過去了.
但是神的道仍然安在.
主的教會仍然繼續.
永遠到主再來之前.
都是不完全的.
都是有很多困難的.
都是有內外的交戰.
內部的分化.
外面的攻擊.
一定有的.
時代也有很多的不同.
但是當我們看回這些經文.
總是有一個盼望.
主耶穌一定會回來.
所以最終改變一切的.
都是耶穌的榮光.
和這個審判的日子.
我們可以做的就是.
盡我們自己能夠為神的家.
為教會的守望.
我們是洗人和睦的人.
叫人做一個和之子的人.
這肯定就是神期望教會要這樣做.
所以保羅就祈禱了.
所以第三段就是.
這位很好的牧者.

$^{801}$繼續為教會的情況祈禱.
所以說了之後.
他自己也祈禱.
我們常常為你們禱告.
願我們的神看你們.
配得過所謀的恩兆.
再一次肯定你們要在這個兆命裡.
來到有這份榮耀.
有這份值得的.
看你們配得過的.
即是說一定是和生活的事奉.
掛鉤的身份.
的生命.
即是你穿一個龍袍都不像太子.
你是富二代.
但你完全都不讓人感受到.
你有一種氣質.
或者叫暴發戶等等.
但我們教徒就應該要.
慢慢由信主.
那個救恩帶給我們的.
就是配得過所謀的兆.
又用大能成就你們一切所羨慕的良善.
和一切因信心所作的功夫.
再一次加人.
因為這個兆命.
因為你感覺到自己是配.
很榮耀的一個匹配.
所以你就更加的殷勤的.
去追求神的屬性的其中一種良善.
和繼續用信心所作的功夫.
記住用信心不是用手段.
用手段就初頭別人看不到的.
但慢慢慢慢發現.
日久見人心.
見到你的動機原來這麼純全.
就打了一劍折扣.
但你用信心在神裡面做.
表面別人看不到.
但慢慢久了.

$^{841}$別人發現這個人原來真是為神國度的.
所以保羅說的是很重要.
追求良善自然就很想追求服侍主.
不是我們自己本身有的.
是我們能夠在主裡面.
是救恩帶給我們的生命的改變.
接著我們就能夠在.
這裡告訴我們.
叫我們在主耶穌的名裡面得到榮耀.
主耶穌的名在我們身上得到榮耀.
主耶穌.
別人會看到我們的生活.
主耶穌的名被別人稱頌.
所以主耶穌的名得到榮耀.
亦都是我們在祂身上得到榮耀.
你們也在祂身上得到榮耀.
因為是祂改變我們.
是祂令我們現在成為這樣的人.
我們現在可以說髒話.
現在我們不說.
不是因為律法.
是因為我根本沒有這種壓力.
沒有這種憤怒要說.
不是不懂.
隨便都可以說.
不會忘記.
但是根本不需要.
現在可以在這個救恩裡面.
我們在祂身上得到榮耀.
我們改變了.
充滿了平安喜樂.
盼望對人有真實的愛.
這些不是我們自己有.
是我們在祂身上得到的榮耀.
是因為祂改變我們.
所以你可以想像到.
主耶穌的名得到榮耀.
我們在祂身上得到榮耀.
都是照著.
沒錯.

$^{881}$是神的恩典.
是照著我們的神.
並主耶穌基督的恩典.
這就是很清楚告訴我們.
原來一切都是救恩帶給我們的生命改變.
所以一個教會.
如果大家都很賣力去傳福音.
大家都很想去服事主.
其實很多不同的性格.
不同的看法.
其實那些很小事.
不要將那些放得那麼大.
其實正正就是不認識救恩.
自己是幾個蒙恩.
才會這樣放大事物.
而每個人都有不同的看法.
有些不同的意見.
其實這樣才是豐富了做事的豐富.
就不是分歧.
我們記得.
真的要用恩典去看自己的生活.
以及用恩典去對待其他人.
第二章開始進入那個氣慾.
或者說這本書後書.
寫作的原因.
第三章就是要解答那些很懶散的人.
有什麼可以幫到他們呢?.
第二章就是這個神學如果錯誤.
就會引發這些懶散的人的原因.
所以本來就要處理這個很重要的實學.
到底主在來是怎樣運作呢?.
這一段我們上星期很快地看完.
那個大罪人的出現.
我們很快地讀一讀.
第二章.
弟兄們論道我們主耶穌基督降臨.
和我們到他那裡聚集.
我勸你們.
無論有靈有言語.
有沒有我們必須信.

$^{921}$說主的日子現在到了.
不要輕易動心.
也不要驚慌.
人不虧用什麼法子.
你們總不要被他誘惑.
因為那日子以前.
必有彌道反教的事.
並有那大罪人就是臣之子顯露出來.
似乎這個肯定是單數.
他是抵擋主高抬自己.
超過一切稱為臣.
和一切受人敬拜.
甚至坐在臣之中自稱是臣.
我還在你們那裡的時候.
曾把這些事告訴你們.
你們不記得嗎?.
現在你們也知道.
那難阻他的是什麼.
是叫他到了的時候.
才可以顯露.
因為那不法的引義已經發動.
只是現在有一個難阻的.
等到那難阻的被除去.
那時這不法的人必顯露出來.
主耶穌要用口中的氣滅絕他.
用降臨的榮光斃掉他.
這不法的人來.
是照撒旦的運動.
行各樣的義林實跡.
和一切虛假的歧視.
並且在那沉淪的人身上.
行各樣出於不義的詭詐.
因他們.
他們就是眾數.
就是那些跟隨那個沉淪之子的人.
他們的生活或者他們的態度.
就是因為他們領受一個愛真理的心.
使他們可以得救.
就是說原來我們信耶穌.
是不知不覺我們的生命.

$^{961}$是揀選真理.
是尋求那個真實.
以致我們就能夠.
因為有這個渴慕的心靈.
我們就可以得到神的拯救.
所以救恩都不是所謂.
藥食亂投.
即是扔下來.
誰撿到就甚麼.
不是拋獸球.
而是揀選.
加爾文的揀選論.
其實就是神預知到.
哪一個人會信主.
然後就會揀選那些基督徒.
那些被揀選的.
這個很簡單的.
將加爾文的學說.
這個圖例這樣講一講.
那些被揀選的人.
是一生持守住這個信.
一直這樣相信.
堅忍到底的.
就是一直會被這個救恩所拯救.
保護.
所以就是聖徒無保護.
就是加爾文挺重的牧羊.
特別當時因為中華鬼革.
很多基督徒信得住.
跟了馬丁路德,慈雲尼,加爾文.
這一班的實學家.
離開了天主教.
天主教說你們離開了.
你們沒有救恩了.
教廷才是救恩.
你在教廷以外.
你就沒有救恩了.
那個聖徒真的害怕.
就是說你已經沒有了天主給你的Blessing.
你已經沒有了那個救恩.

$^{1001}$你離開了大公教會.
這個教廷.
所以對於當時日內瓦,瑞士,德國.
很多的基督徒.
其實他們真的要問.
我們跟了你之後怎樣呢?.
還有這麼長的人生.
人生都不要緊.
我將來怎樣呢?.
所以當時的那個簡選論.
帶給他們永恆的保守.
那個Persuasion是很重要的.
就是說你告訴我.
我是跟隨你們說的恩順清義.
信這位主耶穌.
我一直信是一直保守.
如果你們說是真的.
我們就跟隨這樣信.
因為離開了教廷.
真的是一件很嚴肅的事.
就不同今天.
我不喜歡這個傳聞.
我就轉教會.
沒問題.
這麼多教會.
但當時不是這樣的氣氛.
離開了教廷.
其實是很嚴肅的.
他們那個永恆的永生怎樣呢?.
教會就說沒有.
我們不保守了.
但神學家就說.
嘉義文就說不是.
我們一直在主的救恩裡面.
一直就沒有保守.
一定可以得到永生.
所以就是這個堅韌.
所以這裡告訴我們.
保羅就是這樣提到.
給我們有保守.

$^{1041}$但相反那些繼續拒絕.
而信從虛方的人.
其實就被自己的方言.
這樣去欺控了.
有些人就不想聽.
有些人就刻意不聽.
拒絕這個福音.
所以其實背後.
他們根本就是一個.
生發錯誤的心.
所以他們是信從虛方.
寧願喜歡一些假的東西.
不願意面對自己的心靈.
不願意面對永恆的真實.
所以就用一切.
無論享樂,工作,家庭.
地上的事物.
來去麻醉自己.
所以工作都要在實的眼中去看.
會不會工作成為我們的偶像.
很多人會這樣說.
很多人會這樣想.
但是真的.
有些人有空又不回教會.
很忙又不回教會.
原因就是.
已經將他的生活成為了偶像.
他完全沒有方向.
但他不會承認自己沒有方向.
只會說沒空又不回教會.
教會很悶.
沒有需要.
但其實就是.
他完全不懂得去梳理自己.
有些人根本就不懂.
因為他根本沒有這樣的學問.
但是有更多的人.
只是用生活,用他的工作.
整個工作的人生.
浪費了追求實.

$^{1081}$和追求內心的平安的永生.
所以這裡說.
是使一切不信真理.
喜愛不義的人.
被定罪.
就是在這個失彈的運動.
所以原來一個人能夠信主.
真是一個聖靈工作.
在這個大前提下.
當然是說將來.
運動越來越嚴重.
有很多人跟隨.
不喜歡真理.
接著就轉口問.
就是關乎到我們這些信徒.
主所愛的弟兄.
上文那些就是.
神給他們一個生發錯誤的心.
但對於我們.
是主所愛的門徒.
因為從起初.
揀選了你們.
叫你們因信真道.
又被聖靈感動.
成為聖傑.
能以得救.
這個就是聖人一生在我們內心的工作.
神值我們所傳的福音.
呼召你們去到這個地步.
就是現在這個生命的光景.
無論是事奉.
信心,愛心,盼望的增長.
好得著我們主耶穌基督的榮光.
所以弟兄們.
你們要站立得穩.
凡是所領受的教訓.
不拘是我們口傳.
是信上寫.
剛剛上文才說有些假信.
有些假的留言.

$^{1121}$等等.
現在就說.
我們所說的.
口傳的.
書信所寫的.
你們都要堅守.
所以原來兩批人.
一批人.
很容易被人哄了就下答了.
但一批人.
是因為裡面有真理的心.
被揀選.
有聖靈在我們內心.
能以得救.
成為聖傑.
就是原來一直都有神.
在我們生命裡面工作.
但有一批人.
就是撒旦在裡面工作.
所以保羅這裡.
透過他自己的假冒書信.
即是有些人假冒他的書信.
在這裡說出.
兩類人.
所以我們記得.
我們是第十三節.
到十七節的.
被稱讚.
一生不斷成長的基督徒.
是一直成長.
昨天就去了老人院.
探訪了很多老人家.
有一位在病床.
講普通話.
在內地的老姐妹.
很勁竭.
和身邊.
大家都是不同床位.
隔一格這樣的床位.
有些是九十多歲.

$^{1161}$但那個迷信觀念都很重.
有些又不肯信.
這位姐妹.
不是很大.
但都老了一點.
不太老.
我想六十多七十了.
原來在內地上班很久了.
很勁竭.
見她講話.
一起祈禱.
雖然她現在在老人院住.
但又感受到.
人生真的有不同階段.
所以我們要持續成長.
趁現在還有記憶.
背多些聖經.
將來幫助我們.
因為你有那個建立.
將來你都會有這種意識.
否則突然間會沒有.
當然我們不擔保.
突然間會老人癡呆.
甚麼都忘記了.
不奇怪的.
神會保守的.
但我們在有意識地.
本來就是這樣教信徒.
要這樣追求.
為甚麼呢?.
因為有假冒.
有冒他名的書信.
所以保羅說.
使你們真的能夠成長.
是因為聖靈.
聖靈在哪裡運作?.
聖靈要在真理裡面運作.
即是那個火車.
無論是高鐵也好.
或者是綠皮火車也好.

$^{1201}$都要在路軌運作.
沒有路軌.
那個火車走三百五十公里都沒有用.
它根本沒有路軌讓它行.
大家記住.
聖靈的工作就在我們的生命裡面.
透過我們實踐真理.
會使我們越來越成長.
所以這裡就告訴我們.
福音呼召你們去到現在這個地步.
很得著的就是主耶穌基督的榮光.
那個光榮.
那個榮耀.
不過用另一個說法這樣說.
所以教徒我們不只是做一些很低層次的事.
又要我讀聖經?有沒有搞錯?.
經常都有人這樣說.
每天都要靈修?有沒有搞錯?.
叫你回來祈禱吧.
不是全都人祈禱嗎?.
我們覺得好像信仰變成了一種專人去示範給我們看.
做給我們看.
我們回來欣賞他們的祈禱.
我們不是參與者.
保羅說我們每一個頂指都是參與.
我們自己去追求認識耶穌.
和愛這位耶穌.
和聖靈在我們裡面呼召我們.
繼續孕育我們的生命成長.
所以我們得著的是主耶穌基督的榮光.
每一個基督徒我們得著的是主耶穌的榮光.
所以就是保羅所說的.
要堅守他的教訓.
不單止是我們口傳的.
信上所寫的.
都要堅守.
所以基督教永遠建立在一個真理的宗教.
不是一種道聽途說的憑感覺的宗教.
所以教會不可以經常憑感覺.
不可以經常為自己感動就要做甚麼都可以.

$^{1241}$不健康的.
尤其是全都人更加不可以經常憑感覺.
就會影響頂指很深.
我們一定建立在教訓上.
要深入淺出.
要具體的去明白和去實踐去遵行.
這才是基督教的健康.
當然不只是不斷講道理講文字.
而是沒有了生活的經歷.
這又不是的.
兩者都要平衡.
所以這裡說要堅守.
接著就是一個祝願.
十六,十七節.
但願我們的主耶穌基督.
和這位愛我們.
開恩將永遠的安慰.
美好的盼望賜給我們的父神.
安慰你們的心.
所以神的安慰.
哪位的神呢?.
這位開恩將祂永遠的安慰.
神永遠的安慰.
即是祂真的很體恤我們.
很明白我們人性的那種卑賤.
軟弱,無助.
完全沒有身份.
你想想還是那些奴隸時代.
如果你只是出身一個奴隸家庭.
你完全沒有身份的.
但被一個人抬舉.
你可以想像到你那個生命的改變.
現在神開恩在祂的愛子的救恩裡面.
將祂永遠的安慰.
美好的盼望.
美好.
保羅經常說美好的.
希利文的《Galatians》.
我忘記那首歌是怎樣讀的.
講到很美麗的盼望.

$^{1281}$所以那個盼望.
我們經常做夢.
有沒有一些美麗的夢呢?.
美麗的盼望呢?.
美好的盼望呢?.
我們有沒有想想天堂.
我們很羨慕可以去呢?.
這些我們很少見.
中世紀中國改革的人.
是很羨慕那個盼望的.
譬如五月花號他們要離開英國去美國.
因為受到聖公會的迫迫.
他們要走.
要尋求新的大陸.
在那個新的地方的傳道.
最早期的五月花號那班聖徒.
那班清教徒.
就是有這種美好盼望.
所以我們教徒有沒有耐心.
時時羨慕這些美好.
美善的盼望呢?.
因為我們太現實.
生活太過壓迫.
很少會這樣做夢.
很少會對永恆有這種.
渴望的途徑.
因為我們從來都不稀罕.
這個都要記得我們思想一下.
就是這位賜給我們的神.
安慰你們的心靈.
並且在善行善言上兼顧你們.
誰的善行善言呢?.
是我們的善行善言.
我們實踐這種良善的行為.
美善的言語.
善行善言.
善言.
所以我們很少讀聖經.
我們就不知道原來我們的口.
都是神會保守和會掌管.

$^{1321}$尤其是我們要面對.
就是我們說的話.
千萬不要有些話說出去收不回.
其實很嚴重的.
你傷害了一個人.
一個可能好朋友.
或者是一個弟兄姊妹.
但是原來保羅叫我們追求的.
就是要善行善言來兼顧我們.
所以這裡第二章.
整本書的靈魂就是.
因為有假冒保羅的實學.
書信或者一些的言說.
在教會裡面出現了影響.
所以你可以想像到.
保羅是比較著重教導的.
主要在教會裡面的一個恩賜.
還有教會要成長.
真的需要有神的道.
真的需要對神的道有明白.
這裡一直在說主的再來.
一直在說教會面對困難.
還能不能成長呢?.
是否只是挨打逃避.
既然這麼大的困難.
不如解散教會.
不要再聚會了.
大家各自自己.
喜歡就騎騎肚.
喜歡就去河邊養生館讀讀聖經.
是否這樣呢?.
不是的 教會仍然存在.
無論多大的兵敗.
教會都要為主耶穌十字架作見證.
別人不准你聚會.
不要緊 你不准我聚會.
我不能否定耶穌.
你不准我在一個公開的地方聚會.
我就在你安排的地方去聚會.
即是坐牢.

$^{1361}$我們在牢裡為主耶穌作見證.
是這樣的.
基督教繼續會存在.
一直會存在.
何時不存在呢?.
其實不需要存在.
就是因為主耶穌回來了.
否則教會永遠會繼續存在.
所以我們要求神幫助.
所以大前提就是.
保羅提到有靈有言語.
有無他明的書信.
但剛才我們說了.
第十五節.
你們要站立得穩.
凡是所領受的教訓.
不單止是我們口傳的.
曾經告訴你們的.
面對面說的.
或是我信上所寫的.
或其他的使徒.
信上所寫的.
你們都要堅守.
即是原來兩班人.
一班人很容易受這些假的書信影響.
保羅說其實是他們內心.
根本不接受真理.
根本不會有一個喜歡真理的心靈.
即是會有一種分辨的能力.
沒有的.
所以很容易有實行的運動.
他們就會受影響.
更深的影響.
但原來基督徒有性命.
有一個渴慕真理的心.
就會自然懂得分辨.
這些是保羅的教導.
這些不是旁獨左道.
所以保羅說我的教導.
你們要堅守.

$^{1401}$因為我是奉主.
在主的默示.
感動下去寫的.
是那個真理.
所以真是兩班人.
是兩種結局.
所以我們求神幫助.
渴慕神的還原.
我經常都說.
有很多是自己可以分享的地方.
同弟姐妹彼此去面對.
我每一次我是認真聽到.
不因為那個人是高級.
或者是出名.
或者是不出名.
或者他講得很明顯.
可能很沉悶.
都有一些教訓可以學到.
真的.
所以如果我們作為牧者.
都這麼認真去聽每一篇道.
因為我們不只是聽他的口才.
而是他今天神透過他和我說什麼呢.
這個才是大前提.
所以頂朝更加應該這樣學習.
你想想你只有一個星期.
一個小時.
你能夠坐定聽到.
其他時間你很忙.
有你自己的工作.
有你自己的事業.
這是合理的.
完全沒有問題.
但是最不合理的就是.
居然有人這麼蠢.
那個小時最重要就是.
神要透過他和你說話.
透過那個牧者和你說話.
你就完全看不到這個大前提.
看到的就是.

$^{1441}$這麼悶.
有沒有搞錯.
我等下趕著喝茶.
悶得我死了.
你就被那麼小的東西影響你的心.
而不是你為那位牧者祈禱.
神今天透過他和我說什麼呢.
其實這是很美的一件事.
聖經就是神一定有些事和他說.
然後就要和你說.
你要有這樣的浪漫聽道才行.
每天餐餐你都覺得這個要說得好一點.
那個要再說得好一點.
你的心態已經有點不好.
很少人會在講台上這樣教導頂子妹.
所以這個平台其實多好.
我們可以講到一些很深入的話.
面對頂子妹.
其實頂子妹.
任何人一定有他的好.
有他的限制.
我們要接納別人的限制.
看到他已經盡了他的力.
神都收貨了.
居然是你不收貨.
這樣不行的.
神都收貨了.
他已經盡了力.
神說我整晚在祈禱.
我已經很認真預備.
但我資質那麼有限.
我盡力講了.
勵志牧師又要我講道.
我不講也不行.
我已經盡力了.
神已經收貨了.
他的限制就是他已經盡了.
但如果我們永遠都覺得.
不好呀不行呀.
我不是為任何一個牧者.

$^{1481}$包括我自己去辯護.
這個就是聽到最重要的.
那個背後的精神.
神今天有什麼要和你講.
這個就是最重要.
當然你說神會經常和你說的.
那為什麼你不做傳道去講呢.
這是一個問題.
如果經常有時候聽到.
神和你直接講的.
而你居然還看不起任何一個牧者.
這是一個很大的問題.
為什麼神會和你說話.
而你的心靈沒有改變呢.
稍為我們成熟一點.
我們知道待人要有恩怨.
對人要多接納.
還有要看外面人的限制.
譬如有些牧者你看到他真的盡了力.
你就欣賞他盡了力.
其實你的心靈已經會有所不同.
每個牧者背後都有很多故事.
很多困難.
記住不是為大家辯護.
這真的很實在.
因為我們會看到牧者的困難.
有他自己的掙扎等等.
所以兩種態度.
這班人很容易就會受到影響.
就是假冒輸信.
其中最吸引的就是主的日子.
大家去猜測.
大家去計算到底有什麼驚天動地.
大家就聚焦在主回來的計算日子.
不是去準備好自己.
怎樣去迎接這位審判的主.
所以你會看到.
這裡說牧他的名輸信第二章第二節.
就說主的日子現在到了.
又是強調現在.

$^{1521}$這個就是末世論的終末論.
到底是怎樣的呢?.
很多人就有很多不同的計算.
我完後書.
我就會將好幾個的實行學.
關於末世論.
給大家去做一個.
當是一個討論.
讓你知道在神學界.
有這幾個大的.
關於主在來的神學思想.
神學院都有需要讀的.
保羅說現在到了.
不要輕易的動心.
不要驚慌.
兩者深情.
一些人著弱到一個地步.
如獲至寶般的去攀天起地.
一些就過份的一個懼怕.
因為主在來.
變成一種很大的壓力.
壓迫.
就是一切都沒有了.
所以保羅說.
不要過份的開心.
又不要過度的驚慌.
人不理用甚麼法子.
你們千萬不要受迷惑.
因為乃弟子以前.
必定有離道反教的事.
要稍為停一停.
保羅都意識到.
不是他的時代主回來了.
即是說他自己都沒有領受到.
一些清楚的寓言.
那個的默示.
他要來叫弟兄姊妹.
來動心.
來到去驚慌.
因為主現在來了.

$^{1561}$這個現在.
似乎在保羅有意識地.
不是他的那個時代.
所以那個日子.
似乎保羅都沒有暗示到.
一定不是他那個時代.
但至少保羅在這裡後書的時候.
還是他比較早期.
第二次的旅行報道.
宣教的旅程的時間.
就是他出來傳道.
大概十多年的時間.
來到去講.
即是說他都意識到.
不是在他那個有生之年.
主耶穌回來了.
但是已經這麼快.
湧現了有人在亂說.
說主日子到了.
現在就是了.
所以就帶來.
就有些人心動.
有些人就恐慌.
所以保羅就必須要澄清他的教導.
特別是第四章的.
《貼索羅尼加前書》.
我們講了那個很具體的.
那個降臨的次序.
就是那些已經死了的.
首先復活.
那些活著的.
是會身體改變.
那似乎這裡就會帶來.
不知道當時.
是否已經有那封書信.
貼索羅尼加的教會.
已經手頭上有原裝版本.
因為當時的插款都不是這麼容易.
就馬上有了.
所以到底這些講道.

$^{1601}$保羅的書信.
有沒有實際去到教會領袖的手.
而繼續在保羅不在之後.
繼續教導去解釋.
保羅講的將來主祖來的次序.
會不會就帶來一個空隙.
讓這些似乎又以為自己知道.
但是實際上可能借這個名字.
借主的再來.
來影響著信徒.
所以就真的出現了現在這個局面.
就是有些人就驚慌.
有些人就輕易的動心.
整個人好像瘋了一樣.
主的回來到處跟人說.
就出現了一個很恐慌性的危機.
所以保羅就在這裡再一次面對他們.
有什麼可以帶來的澄清呢?.
由這裡開始.
第三節開始.
一直到第八節.
都是在說離道反教.
都是在說有一個沉淪之子.
會顯露出來.
一直都在說.
那個世界趨向的.
不是一個世界會偏向那個真理.
而是會偏向跟隨那個沉淪之子.
來偏向那個虛方和敵黨真理.
所以到底這裡這一段是難解的.
因為保羅似乎說得很明白.
但實際上是有很多問題我們要推敲的.
到底是最後那個階段的大罪人出現.
好像啟示所說的那個海獸的基督.
還是說每個時代都會有這種形態的一個敵黨呢?.
是當時實際一個敵黨呢?.
就是尼祿等等.
還是那個最末後呢?.
這裡是不清楚的.
很難你憑這一段就可以確定到.

$^{1641}$所以我們看到就是.
講主在來.
而不是在說那個等候的人的生命的一個情操的改變.
就著這個意識而活的生命,生活.
這樣就出現問題了.
就是說我們好奇那個日子來.
因為你知道自己上了岸.
你就好奇那個日子來.
你就到處去傳.
因為你知道自己上了岸.
你就是信耶穌.
所以那個來就很厲害的告訴人家.
這個就是容易出現問題了.
保羅一直上下文.
這一大段都沒有講到很具體是怎樣.
但是下文講到主在來.
信徒就是不是興奮的那個日子.
是我們仍然的屬靈生命.
每時每日的那個成長.
這個成長就緩慢了.
這個成長就不能誇耀了.
因為每個人都是為自己而活.
為自己的靈命在神面前活.
就是說你可能下個月突然間你會有些意外.
你就出現了人生變故.
又或者你可能明年有些大的轉變.
每個人都不想面對這些日日面對的人生.
希望一下子按個鍵就變了.
這個就是整張聖經的分水嶺.
所以後面那段才真正是主在來.
對於信徒最重要的等候.
就是每天好好生活.
每天好好地為自己的身份而活.
身份就是我們被神揀選的人.
這樣來活 大家都不喜歡.
因為太平淡了.
天天都是這樣做基督徒.
天天都是這樣讀經.
沒有什麼特別的.
所以其實根本就不需要羨慕什麼什麼時候回來.

$^{1681}$真的有人問過這個維斯利約翰.
問他如果主耶穌今晚回來你會怎樣?.
他就說我會繼續來到去港道.
港完道就回家吃飯睡覺.
第二天早上起來.
沒有什麼特別.
這個沒有什麼特別.
正正就是他很明白使徒所教導的.
所以這一段首先我們翻轉看第三到第八節.
很難解釋的.
因為真的涉及了很多的時代.
但肯定就是有一班如果不愛真理.
即愛慕真理的人.
就會受到這個撒旦運動的影響.
只會越來越對真理沒有興趣.
和越來越和真理有一個生命的距離.
所以我們真的要求主幫助我們.
我們信耶穌不能走自敬.
真是我們每天自己的成長.
在神的面前去檢視.
我們讀經祈禱,靈修聚會,作見證,傳福音.
做一個愛神愛人的人.
是很不容易的.
因為尤其是在教會裡面.
因為真的很容易人會有不同的誤會.
你要保持和睦,寬恕,彼此接納.
這些就是生命功課.
這些就是等候主在或頂止會.
自己要面對的.
不能假,不能逃避的.
但這裡說有些人因為這樣的訊息來受影響.
但保羅說在那個日子來之前.
他就說有離道反教.
所以我們看這些經文都會感覺到.
原來是有一場熟悉的戰爭.
從來都沒有停止過.
主耶穌未來這個世界.
就已經有魔鬼攔阻主的降臨.
拒絕祂的降臨.
甚至祂來到了.

$^{1721}$猶太人都要置祂於死地.
但是原來主再回來之前.
地上仍然是神容許的等待.
希望罪人可以回轉悔改.
但是原來有一個運動.
似乎是不受限制.
似乎神都容許.
但其實是讓那些不愛真理的人.
只會更加證明他們真是不愛真理.
因為救恩同時間都是苦口婆心地在地上傳遞.
總有些人是會慢慢相信的.
但是有些更加.
只會令那些不信或者是敵黨的人.
更加顯露他們內心的邪惡和拒絕.
是真實地處心積慮的拒絕.
所以不能怪說沒有恩典.
其實一直都有恩典.
所以這裡告訴我們.
有離道反教.
有這個大罪人.
有這個沉淪之子.
所以你可以想像到保羅這樣的形容.
到底他寫的時候是不是若有所指呢?.
這個就最難解釋了.
就是我們不知道他是不是因為政治原因.
他不敢說得這麼白.
說的是地上的君王.
其實我想這個就比較牽強一點.
似乎說的是將來的那個真是敵黨神的人.
他是怎樣呢?.
抵黨主高抬自己.
超過一切稱為神.
和一切受人敬拜.
然後自己坐在神的殿自稱是神.
早一輩的人都會看這個是敵基督.
是那個真是魔鬼的化身.
是那個帶領整個地上的人都去拒絕神.
你可以想像到地上我們不同的時代.
我們都會見到有這樣的人.
他是類敵基督.

$^{1761}$沉淪和高抬自己.
大罪人.
抵黨主.
現在自己稱自己是神.
超過一切神和受人敬拜.
他自己稱呼是無所謂的.
但是既然這樣稱呼.
一定是要帶起這個漩渦.
帶領很多人歸向他.
所以任何的邪教主教都有這樣的性格.
都有這樣的自視自大.
用宗教去包裝自己.
怎樣得道.
所以如果一個牧者盡力.
人們都不聽其實是一種祝福.
每個人都聽你 每個人都崇拜你.
可能是一種罪.
好像說得誇張了 其實是的.
因為那些邪教教主是洗腦式.
人們是完全懂事.
蜂窩地跟他.
韓國也有不少這些邪教.
一定有這些行徑.
抵黨主一定是高抬自己.
一定會自稱為神.
坐在神殿裡自稱是神.
一定覺得自己有某一種神力.
神諭.
可以講語言.
可以指出哪個人的問題.
好像能夠改變其他人.
這些要很小心.
牧者就有這種.
本來這個積焚是希望改變人歸向神.
要將人帶到神那裡.
但是似乎這個撒旦的沉淪.
這些大罪人.
要將人帶到魔鬼.
帶到自己那裡.
所以牧者的大忌就是.

$^{1801}$不要累積自己的粉絲.
應該有粉絲.
有人對你敬重是沒有問題的.
因為他知道你經常對人敬重.
但絕對不是將人帶到自己的心靈裡去操控.
這個我自己一生.
我很謹守這件事.
我很清楚自己不可能這樣.
這是很重要的.
當然我沒有權說其他人.
因為神會審判.
但自己一定要謹守這件事.
所以這些就是撒旦運動.
要很小心.
我們越是出名.
記得要將榮耀歸給神.
第五節.
「我還在你門那裡」.
曾經將這些事情告訴了你們.
你們忘記了嗎?.
因為有假冒的書信.
所以現在真的出現了危機.
不只是嚴重的失敗離開了真理.
但一定已經去到一個.
在教會裡面討論這個課題.
而這個討論不是帶來生命的益處成長.
是帶來對這件事的好奇和驚慌.
所以本來就勸勉他們.
不要接受這些假冒的教導.
第六節.
「現在你們也知道.
那難阻他的是甚麼.
是叫他到了的時候才可以顯露」.
很多的代名詞.
所以我們根本不知道他.
或者是誰難阻他.
這個他肯定是上文說的大罪人.
那個沉默之子.
誰是難阻他呢?.
現在未能夠顯露出來.

$^{1841}$大部分人會覺得是教會.
本來是說現在有一個神的屬地.
即是地上的教會.
或者教會的時代.
是難阻著這個不法運動.
是全面的.
佔據所有地球上的每一個地方.
一般會說是教會.
當然不是某一間教會.
不是全部貼在一起的教會.
是現在在地上還有一個福音去傳的地方的教會.
所以是這個難阻.
暫時是這樣解釋.
另外有一派就說是聖靈.
也說得通的.
但根本不知道保羅所指的是誰.
所以我們也不能夠很肯定地說.
一定是教會或者一定是聖靈.
你說神一定是對的.
沒有說錯的.
但這裡告訴我們.
不法的引義其實已經發動了.
只是現在有一個難阻.
又繼續說第二次了.
說到這個難阻.
說到這個不法的引義.
記住是隱藏的一個意識.
在運作中.
是不法的.
只是現在有一個難阻.
等到這個難阻被除去.
所以就說教會被提了.
信徒被主耶穌接去了.
就開始.
所以這個就是其中一派的大災難.
七年大災難.
這一派就是教會被提了.
地上再沒有福音.
再沒有教會.
這個不法的人就全面顯露出來了.

$^{1881}$就是基督的出現.
等到這個難阻被除去.
那時這個不法的人就必顯露出來.
再一次說那時.
就是最黑暗的時代.
主耶穌當然是回來才會.
用口中的氣滅絕他.
用降臨的榮光廢掉他.
普羅為什麼可以這麼仔細的描述.
主耶穌口中的氣.
是不是氣功.
當然我們不需要這樣聯想.
就是說用一個很顯而簡單的動作.
就可以將盤踞在地球.
在地上可以影響眾多的人.
歸向虛方拒絕真理的一個世代.
一個時代.
以掌控的大罪人為首.
主耶穌很輕易的.
就已經可以將他自伏.
口中的氣已經可以自伏他.
降臨的榮光去廢掉他.
你可以想像到.
不是神沒有那個能力.
是時間未到.
就是說現在仍然容許罪惡.
但其實容許罪惡的原因.
仍然是福音廣傳的時代.
我們要記得.
所以仍然是福音得勝.
不是罪惡得勝.
仍然是我們可以傳.
總有人會聽他說.
我們要找到那些總是會信的人.
就是原來有些是實揀選了的人.
所以加爾文的學說.
其中奧妙的地方就是.
我們千萬不要以為.
哪個人一定信不了耶穌.
因為我們知道誰會被揀選.

$^{1921}$但加爾文的神學.
我是加爾文主義.
我都接受的.
我不是亞敏尼亞思想的.
加爾文思想其中一個很重要的精髓就是.
你和我永遠不知道自己是不是被揀選.
別人是不是被揀選.
但是福音去到那個人.
如果他是被揀選.
神是會在他的生命裡工作.
他就會回應那個呼喚.
他就會被揀選.
他一生為這個呼喚而活.
就是肯定他被揀選.
這也是加爾文當時的模樣.
因為剛才不是說了.
天主教說教會以外沒有救恩.
你們離開了.
就是分散的羊.
你們沒有救恩.
我們不會有任何的祝福給你們.
就是那些神父.
彌撒等等.
所以他們問加爾文.
我們怎麼辦.
加爾文說.
你在主的救恩裡面一直堅守.
你們就是被揀選.
這是很重要的盼望.
所以這裡告訴我們.
主耶穌很輕而易舉地.
就可以將這個大罪人.
你想想他盤踞了這麼多世代.
甚至最後帶著這麼多人拒絕真實.
但是主耶穌只需要用口中的氣.
只需要用降臨的榮光.
就可以將他治服.
所以現在.
所以真的是半杯水和一杯水的道理.
現在是罪惡越來越嚴重.

$^{1961}$還是神仍然容許福音繼續在這個罪惡世代傳.
總是可以選到一些人.
至少可以選到你和我.
現在能夠這樣聽到.
至少我們的生命有改變.
至少我們現在喜歡幫人多於害人.
我這樣估計.
當然應該不是一定想幫人.
因為我們有神.
因為我們有了信仰.
我們知道神是賞善罰惡.
我們不會想著害人.
這已經是救恩帶給我們的一個挑選帶來的改變.
所以這個似乎好像很誇張地盤踞這麼多世代.
容許地球上的人受虛方影響.
其實表明就是那個人影響其他人.
其實是那些不領受真理的心的人.
自己拒絕福音.
拒絕真實.
這就告訴我們這個實行的運動.
用一些各樣的異能,神蹟和虛假的歧視.
在這個大罪人沉淪之子的身上.
行各樣出於不義的詭詐.
所以全部東西.
超自然我們都要問一問.
為什麼持續有超自然.
而這個出現超自然的人.
為什麼和真理和愛和和平和平安是背道而馳的呢?.
就要問這個問題.
那些實跡虛假的歧視就真的虛假了.
因為沒有改變那個實跡的人自己的生命.
所以邪教的教主就是這樣.
說了很多超自然的東西.
但是自己的生命沒有改變.
這很明顯就是一個不同的地方.
所以就會令那些愛真理的人.
不領受愛真理的心.
使得他們得救的心靈.
就會在這些行虛假歧視的神蹟.
這個大罪人身上.

$^{2001}$就會得到他們的心.
所以記住.
如果我們只喜歡那些興奮,很開心,很吸引的東西.
基督教永遠就是這樣.
多走一些神蹟給我看.
我就會很不同了.
那你要很小心了.
如果你從來都沒有一個愛真理的心.
使得我們能夠對真理,對主,對愛,對福音.
產生濃厚的興趣.
而且改變我們的生命.
那我們可能走錯了,選錯了領袖.
因為主耶穌叫我們.
我來了是要叫羊得生命.
然後是得豐盛.
豐盛在裡面的生命.
就是說就算周圍的環境沒有改變.
就算我這個人在這個環境.
我的情況都不是那麼樂觀.
沒有想像的.
每樣東西都順暢.
都不要緊.
因為我裡面有豐盛.
因為我有主耶穌.
我就有豐盛.
所以真的不同.
信主久了會看到你裡面有沒有豐盛.
初信未必感覺到.
有些人信得很好馬上就改變.
有些人慢慢沉著.
但是日子久了.
你是否越來越愛那個生命的豐盛呢?.
還是你只是越來越喜歡那些千奇百怪的東西呢?.
要很小心.
因為這是你的胃口.
內心的縮力胃口.
所以求主幫助我們.
這裡說的是那個撒旦的運動.
沉淪之子會做很多的實跡異能.
虛假的歧視.

$^{2041}$在這些沉淪的人身上.
各樣出於不義的詭詐.
這樣來牢籠人心.
所以主耶穌永遠是用愛.
用真理去致服我們.
所以任何一個聖徒.
你剖開他的心.
只會看到他愛耶穌.
當然不能剖開.
因為不是真的剖開心臟.
而是你問每一個聖徒.
為什麼你要這樣做?.
每一個愛主服事主的弟兄姊妹.
為什麼要這樣做?.
根本不需要問為什麼.
因為我愛我的耶穌.
所以我們建立的就是這樣的門徒訓練.
我們建立一群愛主耶穌.
愛真理的門徒.
否則撒旦的運動就在這些人身上.
來做一個影響.
繼續偏離真理.
繼續一個愛真理的心.
這個就是保羅很緊張地寫後書.
就已經有這些假冒書信出現.
而那個內容就是說.
主的日子現在到了.
所以那個大前提就是.
可以用不同的方式.
仍然問我們.
我們甘不甘心過個星期.
默默地在教會裡服事.
越來越沉靜地去幫助別人.
幫助也不需要打鑼鼓.
不需要別人說.
為什麼你不稱讚我?.
我做了這麼多事.
其實不需要的.
因為做在主的身上.
主是會知道的.

$^{2081}$還有我們真的越來越傾向的是那個真理.
不是要尋求自己的公理.
而是尋求神的真理.
所以就算被人誤會.
船啊很小事這麼簡單.
何必這麼介意呢?.
不如介意主的真理有沒有被傳開.
所以這裡有兩個很大的分水嶺.
都是保羅的書信.
第二段的人13到17節.
就是一聽到保羅的書信.
就願意遵循的人.
所以你看到兩個大分水嶺的分別.
我們就停在這裡吧.
今晚我們一起祈禱.
感恩我們再一次進入.
貼實羅尼加後書.
保羅思想的世界.
和這位牧者.
他是如何緊張教會.
不受到那些錯謬影響.
錯謬帶來的.
連生活的行徑都可能受影響.
所以我們懇求主幫助我們.
有一個愛慕真理的心.
愛慕遵循你的旨意的心靈.
和真是在我們的生活裡面.
實踐我們所信所明白的.
我們這樣望感恩.
禱告奉耶穌基督的名.
阿們.
\newpage



\section{帖撒羅尼迦後書}
\label{sec:ewzMm4Y0R5I}
\textbf{2024-08-01 周四晚查經班《帖撒羅尼迦後書(3)》 講員:曾裔貴牧師}
\newline
\newline
連結: \href{https://youtube.com/watch?v=ewzMm4Y0R5I}{\texttt{https://youtube.com/watch?v=ewzMm4Y0R5I}} ~~~~ 語音日期: 2024-08-06
\newline
\newline
\hyperref[sec:ntrA1aTVgRk]{\small{< < < PREV SERMON < < <}}
~
\hyperref[sec:index_chronic]{\small{[返順時目]}}
~
\hyperref[sec:index_scriptual]{\small{[返順卷目]}}
~
\hyperref[sec:bkamj1iGLMk]{\small{> > > NEXT SERMON > > >}}
\newline
\newline
$^{1}$我們今晚一起來禱告.
多謝主讓我們能夠在不同的地方.
在不同的時段都有不同的聚會.
現在在城鎮有現場的培靈年經大會.
都是香港不同的教會一個比較重要的聚會.
我們真的很盼望上主來祝福.
使到每一個的信息.
特別是一個有系統的研讀.
某些經卷的題目.
能夠祝福更多的弟兄姊妹.
一方面我們盼望一個已經是.
行了這麼多年的運動.
可能對新一代的信徒.
他們未必那麼習慣這樣的單向的聽道.
又或者可能現在年輕的弟兄姊妹.
他們面對的問題.
或者他們內心對於公義,對於民主.
對於一些其它的事情.
尤其是一個落實的信仰在生活裡面.
都是有可能很強烈的內心的想法的.
我們都盼望海外的牧者.
他們的十天的信息.
能夠適切本地的需要.
也都真的讓更多的青年人.
被主的聖靈感動.
可以研讀聖經.
找到我們自己新一代現在這個社會.
或者在21世紀關注的議題.
無論是氣候變暖.
或者更多的一些倫理道德的問題.
都讓我們能夠有那個出路.
有新一代的牧者可以被興起.
可以來傳遞一個與時代並存的信息.
今天晚上多謝主給我們.
我們不多人 但是我們可以繼續去研讀.
慢慢地將《貼十勞尼加》的書信去講解.
希望能夠明白到保羅.
他原初寫的時候寫作的原因.
以及找到那個出路,答案.
願主幫助我們.

$^{41}$也都盼望當往後成為了YouTube的記錄之後.
有更多海外或者弟兄姊妹.
可以有空的時間可以聽.
去得到建立 願主幫助我們.
我們恭敬地打開聖經的時候.
盼望主的聖靈都在我們大家的心裡面.
來工作 我們仰望禱告.
奉耶穌基督的名 阿們.
我們又再重返第二章.
第二章的《貼十勞尼加》後書.
來再去溫習.
因為上個星期是暑期遊樂園.
我們稍微停一停.
第二章都相當明顯.
是有些經文我們沒辦法解得明白的.
因為保羅似乎講得比較含蓄.
但是當然講了一個就是.
總是有一個日子 不是現在是正在發生的.
但是似乎不是在保羅那個年代.
就是馬上應驗的.
我們就一定要拿捏住這件事.
意味著主再來.
不是我們想像當中.
是真的好像有一個完大色那樣.
去到某年某月就會停頓.
這個就變成真是呼喚回耶穌講的.
是沒有人知道什麼時候主再回來.
我們必須要學習的就是.
不要扮演那個曉得去預知未來.
和去分析 去做一個研判.
而是在在福音書講主再來.
是講我們每個人警醒自己那個人生.
等候 就是這麼多.
在這個警醒和等候的過程.
就是要忠於自己那個照明.
在那個照明裡面做好那個買賣.
所以主耶穌不斷在用那個僕人.
就是一個懶惰的僕人.
還是一個有智慧有見識的僕人呢?.
所以稍後我們講完這三章之後.

$^{81}$就會開始有一些關於末世論的不同的學說.
給大家有一個了解.
也都不會有具體的答案.
當然我會講我自己的一些見解.
所以也都涵蓋關乎到末世論.
必須是由24章的《馬太福音》.
那裡的《鑒欖山》的那個儀仗語言那裡講起.
所以涉及到好幾個《馬太福音》裡面的比喻.
我們又要認識.
然後我們就進入那個末世論的神學的討論.
第二章就很明顯講到.
保羅警告帖索羅尼加的弟兄姊妹.
有一件事情是要很小心.
稍後再看的時候.
這裡原來又涉及到聖經的新約書信的形成.
那個過程.
這些真的是很深奧的東西.
要專門在神學院裡面讀.
一些聖經研究或者是那個叫做正典.
但是我們稍微了解一下都可以的.
其中有一樣東西要留意就是.
保羅說已經有一些假冒我的名的書信.
和一些想法,一些學說.
已經充斥著在哥林多.
或者是當時的歐洲的哥林多.
或者是帖索羅尼加.
就是北一點的歐洲.
又沒有去到很北的.
剛剛是這個馬其頓,比勒底,德羅亞.
對陣的地方.
所以帖索羅尼加就是在這個很重要的.
一個地中海貿易的地方.
那個航運都是相當發達.
這裡告訴我們我們就要很小心.
就是說保羅不單止是勸勉他們.
要小心當地那些影響他們的信仰.
令他們灰心後退的那些攻擊,逼迫.
甚至是一些似是而非的學說.
現在這裡甚至關註到主再來這個課題.
他們要很小心.

$^{121}$就是說在前書四章講得很清楚那個次序.
雖然我們不可以計算那個日期.
但是我們知道當主耶穌一回來.
那個預兆是清楚的可以看到之後.
就有次序地死了的人.
那些叫亡靈呢?.
是死了的人其實靈魂已經在天上了.
那怎樣在不同的地方.
有自己榮耀的屬天的身體.
來相繼的復活.
似乎又不只是墳墓打開這麼簡單.
不是那些土葬的突然間那些棺材飛開.
那些恐怖片就會了.
而是每個已經離開世界的信徒就會復活了.
有回自己的身體.
但是到了主耶穌回來的時候.
仍然在生的.
就必然會有一個身體的轉變.
就是不需要再死再復活.
而是在身體的轉變裡面復活.
這個也是一個很奇妙的一件事.
所以這麼清楚在第四章講了.
就是說那封書信帶到貼索羅尼加.
就帶來很大的一個震撼.
但是原來會有反效果.
原來有人就會冒充寶龍.
來講這一課的Blessing.
是已經開始發生了.
是帶來很大的.
到底是振奮還是困惑呢?.
很偶然間.
大家不知道有沒有聽過.
曾經有雙氧水的事件.
叫他們買那些教會的求生包.
然後就說約定要上大帽山頂.
在那裡等候主再來.
我們就覺得很搞笑.
但是如果你是經常被那個宗派洗腦.
你就一定會買.
你就一定會覺得.

$^{161}$我們才是那幫被救贖的人.
後來當然經過傳媒的輾轉就沒有這件事了.
也都沒有了.
仍然都是的.
譬如二端守望台.
仍然都講他們自己.
就是那14萬4千他們唯一.
Chosen people被選的.
譬如安思耀會.
他們現在不敢這麼說.
他們說他們的救主已經降臨了.
不過在天上預備.
所以現在我們安思耀會的人要準備了.
後來輾轉修訂.
說主耶穌還沒有回來.
就是不攻自破的.
當你一計算好那個日期.
自然就會成為二端的了.
因為沒有人能夠預測的.
當然絕對不可以說今晚就不是了.
今晚我起床看奧運就真的不是什麼決賽.
其實我們不要這樣想.
這樣想就是否認了主耶穌隨時會回來.
因為《聖經》說是隨時.
但這個隨時不是三千年後.
也都不是幾年幾月哪一天.
所以有時那種弔詭.
那種無伶兩可.
就會被外邦就是不信主的人嘲笑我們.
或者根本哪裡會有主的再來.
世界一來就是亂亂亂.
二來世界都不知道多麼美好.
還用靠你的耶穌嗎?.
兩樣都不對的.
完全美好又不對.
完全敗壞又不對.
世界還有很多很多的美善.
至少還有基督徒.
現在默默的為神作見證.
諸如此類.

$^{201}$所以這個原來是無名的書信.
說到主再來.
原來每一個時代都有市場.
原來保羅說.
真的必須要小心.
不要隨便相信.
所以他提到.
有他的親筆簽名.
這是第三章裡說的.
三章十七節.
我保羅親筆問你回安.
但是我的信都以此為記.
即是有他的筆跡.
就是這樣.
即是意味著原來我們留意新約.
原來透過這一節聖經.
我們知道原來很多新約的書信.
因為當時沒有這麼先進的文字記錄.
有電腦打字.
又有人什麼的.
原來有錄音.
沒有的.
是有口訊.
多數是口訊.
接著有一個類似文士.
幫作者立刻盡快.
好像默書一樣.
去撰寫.
所以這個過程一定會有人為疏忽.
錯謀.
需要後來的《釋經》.
或者做那些聖經原文研究的人.
比較不同的抄本.
來盡量找回原意.
這個過程.
其實你可以想像得到.
如果你是不信的人.
那就是聖經充滿了錯謀.
充滿了人為的限制.
包括連保羅口訊的時候.

$^{241}$他自己都會錯.
幫他寫的人.
比如羅馬書.
是誰幫保羅寫的呢?.
有一個叫做德丁.
十六章羅馬書.
那個叫做德丁.
你可以想像得到.
是有這個口訊的.
原來在古代是成立的.
是可以這樣成為作品的.
尤其是那些嚴謹的文學作品.
多數是會有人去做撰寫的.
所以保羅告訴我們.
他親筆寫的這封信.
但是第二章的《節錄行間》.
我們就要留意了.
大家先聽聽這一段聖經.
弟兄們.
論到我們主耶穌基督降臨.
和我們道他那裡聚集.
我勸你們.
無論有靈有言語.
有沒有我明的書信.
說主的日子現在到了.
不要輕而動心.
也不要驚慌.
兩者都是兩個心情.
兩個想法.
一個就是非常興高采烈.
因為你會動心.
因為萬事萬物是一下子的改變.
完全得到神無窮盡的任弊的開展.
你是其中的人.
你想想.
你是多麼的光彩榮耀.
例如香港奧運完結了.
拿了十面金牌.
你想想.
你的興高采烈是很大的榮耀.

$^{281}$所以也有另一種情況.
就是會驚慌.
因為很多人可能覺得還沒有準備好.
什麼都還沒有預備.
甚至他自己有很多地上的掛慮.
都還沒有解決.
哇.
什麼就來了.
他就驚慌.
可能驚慌是他自己的東西都沒有了.
不是驚慌他去不了耶穌那裡.
所以其實保羅就將心情說出來了.
大前提是什麼呢?.
其實就是前書第四章那段聖經.
說得很清楚主降臨的次序.
最終降臨的時間.
首先有天使掌的呼聲.
有神的號吹響.
接著有主耶穌的呼叫.
接著祂降臨.
接著似乎就是一個聚集.
就是那些死了的聖徒.
古往今來的.
不管多少千年了.
都全部同時間復活.
似乎又不佔據空間.
不會說數千萬,數億人.
怎麼可以擠在天上這麼大.
是不是夠不夠位置呢?.
沒有這些困難.
沒有這些問題.
似乎那個靈體.
又不是有物質的空間.
你可以想象到那種奇妙榮耀.
但是.
現在保羅好像要補妃了.
因為正正就是這件事.
在民間帶來發酵.
有很多的誤傳.
所以.

$^{321}$這一種誤會.
如果不盡快保羅意識到.
馬上解決.
就會成為一個謠言.
又會成為一個錯謬.
所以.
保羅一意識到.
他就立即去解說.
所以.
那個大前提就是.
輪到我們主耶穌的降臨.
以及我們.
我們就是說.
死了和還沒死的信徒.
去到祂那裡聚集.
所以是一個聚集.
是一個將來天上的團聚.
你這樣想天家.
就很美了.
我們回神的家裡面去團聚.
跟所有的聖徒團聚.
所以.
其實我們有生之年.
真的應該.
尤其是我們越來越近離世的日子.
我們應該更羨慕回天家.
更羨慕團聚.
但是記住.
這個羨慕帶來的就是.
我們很重要的就是.
新教的工作倫理.
因為第三章立即說那些懶惰人.
可能就是那些錯誤的異端.
影響帶來他們對工作倫理的錯誤理解.
所以.
羨慕回天家和工作倫理.
是完全等號.
新教的工作倫理.
新教工作倫理當然就是加爾文神學.
我今天臨回來之前.

$^{361}$在洪恩的辦公室.
特意拿了加爾文的那兩套鉅著.
已經由一本翻譯成兩本.
兩本講了一千頁.
當然它是分類的了.
開始講到工作的照明.
原來工作倫理.
最重要就是你的照明.
神在你的工作.
在你的人生.
安排你的崗位.
你在那個崗位裡面盡忠.
這個就是新教倫理.
工作的倫理.
你盡忠.
很自然你就會一路一路的完善你的工作.
你就會越來越.
因為這個工作而帶來一定的滿足.
滿足帶來的當然就是財富.
或者你的社會地位改善.
所以到一九零幾年的時候.
因為2002年.
就是一本很出名的鉅著.
其實那本書很薄的.
就是社會學的書.
我也提過一下的.
就是新教倫理與資本主義的精神.
就是這個馬克思,韋伯.
這個社會學家.
觀察到為什麼基督教的國家會這麼發達.
連天主教的北歐的一些國家.
比如法國等等.
他們的智庫不是這麼大.
原來他們對工作的態度是不同的.
原來天主教的弟兄姊妹.
他們對工作的倫理.
與加爾文的教導下的工作倫理是完全不同的.
所以這個馬克思,韋伯.
其實他本身不是信主的.
這個也是德國的猶太人.

$^{401}$這個馬克思,韋伯觀察到.
就將這個社會現象寫了一本鉅著.
就是後世社會學的一個很重要的文件.
就是新教倫理與資本主義的精神.
就是說為什麼會有資本主義,資本家興起呢?.
就是因為新教的倫理這個工作觀.
所以這個觀察是相當具震撼性的.
所以其實就是保羅第三章教我們.
不做工的不要給他吃飯.
其實就衍生第二章關於主宰內的問題.
原來主宰內的問題要恰如其分地說.
這樣去領會,領受.
然後我們要在我們自己有生之年.
一方面一定要找到自己的跑道.
一定要找到自己神在你生命中的照明.
所以那個Calling.
在加爾文神學裡面是一個很重要的詮釋.
就是對於加爾文神學的那個揀選論的詮釋.
就是說你怎麼知道你自己被揀選.
就是你知道你的工作.
就是你每天差不多是佔據你整個人生.
你在你的工作照明裡面.
來到恩勤的.
你就自然知道自己等於被神揀選了.
他們是這樣理解的.
所以很自然就變成了雨後春筍.
離開了天主教.
但是反而更加激發他們.
恩勤的去工作.
盡心的星期天回教會.
繼續在教會裡面有美好的服侍.
這個是加爾文主義的精神.
所以我們回到這裡看的時候.
那個大前提就是主的降臨.
和我們去到那裡聚集.
也沒有否定就是死了的復活.
未死的見主面的.
就會身體改變這樣的聚集.
所以將來是一個真是永恆的聚集.
第二節就是寫這卷第二章的原因.

$^{441}$就是我勸你們無論是有靈,有言語.
有沒有我明的書信.
似乎有幾個層次有靈界.
可能是那些人說自己有靈感.
有預言.
有一些神的託付要講.
當然我們可以聽.
但是聽完之後.
仍然都是審慎透過理智去思考.
和聖經有沒有違背.
那個是要神思靈辨.
才知道我們講到第五章.
不要消滅聖靈的感動.
接著不要藐視先知的講論.
但是要凡事測驗.
善美的就要持守.
不善美的.
亂說的.
異端的.
我們就要否定它.
所以很清楚的.
聖靈感動.
和你有神思靈辨.
嚴格來說是沒有衝突的.
基督教似乎兩件事都是重要的.
不可以只是聖靈感動.
你又要有對真理的思考.
測驗和分析.
所以幾件事是沒有抵觸的.
你可以既有感動.
但同時你可以問其他一些屬靈長輩的意見.
再自己去尋求.
聖經是不是這樣說.
幾件事加起來.
就更加客觀地對神的旨意的了解.
所以這裡告訴我們.
這裡有靈有言語.
什麼言語呢?.
可能是人的說話.
例如自稱是先知的.

$^{481}$跟你說.
哪年哪月哪日.
就會有特別的意象.
諸如此類.
又或者有些人對這些方面很有興趣.
現在萬軍攻打以色列了.
因為它發動戰爭.
第三次世界大戰.
似乎看起來越來越像.
它炸死了伊朗.
炸死了那些領袖.
可能伊朗就發動聖戰了.
那就變成了.
不只是以色列和巴基斯坦的那個地方.
而是可能整個阿拉伯世界.
已經來到戰爭了.
萬民攻打以色列.
可能就是預言了聖經的應驗的地方.
當然我們這樣看下去.
似乎局勢就是這樣.
但是我們不要太過用政治新聞.
或者時勢,時局.
去讀入聖經裡面.
這些我們都要留意的地方.
重點就是.
要小心留意.
有沒有我明的書信.
果然就是冒充保羅寫的信.
因為當時實在是.
那個資訊一定不會像今天這麼流通.
但是一封信真的要人手帶去.
最多可能是馬車.
或者保羅的時代已經有船了.
但是真的很簡陋的帆船.
船難也是容易有的.
所以似乎已經有郵局這個制度了.
但是沒有今天這麼方便.
這麼快遞等等.
第二就是收到一個包裹.
你可以想像到.

$^{521}$一定會有這種延遲帶來的影響.
所以他們要留心的就是.
這一封信.
當然是說教會的領袖.
那幫帶領的人.
他們要把關收到的不同的信息.
有些人可能是路過傳道.
你都會容讓他可以去講.
但是那個被邀請的.
或者邀請的人.
就要留意他講的是否合乎真理.
所以那個把關是每一個教會的領導人都要把關.
因為你不能說我看穿你不是使徒.
你不是傳道者 你就不可以講.
那是很難的.
因為你不能這樣判斷.
當時也沒有什麼神學院.
有證書,有畢業等等.
所以一定是會有很多不同的人.
都會按著那個感動.
譬如在《聖堂主義19章》.
有一個叫阿波羅.
他在哥林多.
他很能講的.
非常懂得主的道.
講解很詳盡.
但是阿居拉夫婦聽完阿波羅講道.
就邀請他晚上到他家裡跟他們聊天.
發現他原來對某些真理不是很全面的透徹.
所以就將神的道再幫他講解得更加詳細.
後來輾轉阿波羅就去了這個義肢所.
這個是《聖堂主義19章》.
阿波羅已經是一個年輕有為的傳道者.
不過阿居拉夫婦都聽到有些端倪.
是需要給他有更好的詮釋得清楚一點.
涉及到的似乎是《聖堂主義19章》說的.
那個施洗約翰的事件.
有些人只知道約翰的洗禮.
不知道關於耶穌的洗禮.
到底是不是阿波羅自己也是這樣呢?.

$^{561}$這個聖經就沒有太詳細交代.
所以這裡告訴我們.
今天我們就比較少.
偶然間都會有新聞說.
曾經有一個叫猶大書信.
又說主耶穌原來是什麼什麼的先生子.
猶大書信裡面這樣說.
所以當時的聖經學者.
就是早幾年前.
應該是疫情之前的了.
那聖經學者.
特別是新約的學者就有責任了.
就出來講解什麼是正典.
為什麼我們不需要擔心這一封書信呢?.
絕對我們不需要看它是一個.
對基督教的一個拆毀信心的動員.
但是那個學者就必須要.
透過他們一直這麼多年的.
神給他們的崗位的那個恩賜,資源.
他們就去講解什麼叫做猶大福音.
還有我們根本不值一提.
但是就要講解給人知道.
尤其是對於微信的人他們會有興趣的.
但是保羅的時代.
就已經有這種情況出現.
一定是你有正面的真理.
自然就會有人很想濫用.
無論是善意或者是惡意.
去濫用這些真理.
或者在這個真理之下他自己加上他自己的詮釋.
所以就要很小心.
無論就是今天我們就少.
再有這些新約的書信出土.
發掘出來.
我們很快就會知道一定是假的.
一定是錯的.
我們就很少會受影響的.
但是初期教會因為正典都還沒有形成.
甚至連使徒都還沒有離世.
就已經有這些流通.

$^{601}$所以他們就有責任.
必須盡快將這些假冒來澄清.
最重要就是讓信徒不要受迷惑.
現在仍然是那個課題.
主耶穌基督的降臨.
和我們去到祂那裡聚集.
這個課題就已經有人去誤傳了.
所以保羅就說.
他們說這封書信.
說那些無名書信裡面.
其中提到的就是.
可能繪影繪聲.
可能說到現在全面開始了.
可以帶來很多人就會很開心.
但是也有些人同時也很恐慌.
所以說主的日子現在到了.
重點是那個現在.
就是說在保羅的時代.
所以第四章.
有些解經家就說保羅.
其實他這樣理解.
主的回來了.
但實際上就沒有回來.
還沒有回來.
也不是在保羅的有生之年回來.
所以就覺得保羅有些錯了.
但其實保羅在第四章.
根本就沒有說到.
他的有生之年就是主在來.
所以我們其實就不應該這樣誤會保羅.
但是現在.
無論如何在他自己寫這封信的時候.
就已經聽聞.
甚至乎可能有證據.
就說主耶穌的日子現在到了.
這就是很大的影響.
不是開玩笑的.
所以保羅就說.
沒有輕易的動心.
動心就是發達了.

$^{641}$興高采烈.
你就可以想象到.
接受羅尼加成昨天的初信.
信主也是一個月多一點.
最多現在給他.
再計算遠一點.
可能也是半年以內.
保羅就已經發現有人誤會了他的信息.
接著又用了他的名字.
來說主的日子現在到了.
所以這個現在到了.
是會帶來很大的影響.
所以不要輕易動心.
當然更加不要驚慌.
但是我們又不動心又不驚慌.
可能就是我們這一代的人.
基督教.
根本我們都不知道什麼叫主在來.
完全不知道那個是很重要的.
唯一未兌現的神的應許.
神的預言.
這才是一個問題.
又不開心.
動心.
又不驚慌.
那你想怎樣呢?.
其實就是我們要恰如其分.
主在來是我們每一個信了主的頂子妹.
自然就是.
你已經好像如果用一個大抽獎.
你的名字就已經自動放在那個抽獎箱了.
每個人都可以等候主在來的.
除非我們最終沒有被神揀選.
最終我們根本沒有信到主.
那就一定不是在那個抽獎箱裡面.
所以這裡就告訴我們.
叫頂子妹千萬不要動心.
以及不要驚慌.
第三節.
然後就解釋不要動心和驚慌的出路.

$^{681}$或者是方法.
「人不驅用什麼法子.
你們總不要被他誘惑.
因為那日子以前.
必有離道反教的事.
並有那大罪人.
就是尋人之子顯露出來.
他是抵擋主.
高抬自己.
超過一切稱為神的.
和一切受人敬拜的.
甚至坐在神的殿裡.
自稱是神」.
所以如果我們真的按字面的意思去理解.
似乎保羅真的按字面去講解.
雖然不是在他這一代的時間.
就有這個大罪人.
居然斗膽地說自己是神的.
去到神的殿裡說自己是神的.
這樣的一個褻瀆的人出現.
但是保羅就在這裡.
好像講了一個主在來.
真名天子和假名之子的對比分別.
尤其是強調是誘惑.
誘惑就是說.
是令你心思思.
是令你很容易落踏的.
令你失去了防備的.
真理就不需要誘惑人的.
真理就是明碼實價.
主耶穌一定會回來.
你只需要在你有生之年.
記住加爾文.
我加爾文學說的.
我加爾文.
連我姐姐幫我改名都叫加爾文.
入神學院都很尷尬.
你叫自己做加爾文.
其實與我無關.
我姐姐經常在網絡上教.

$^{721}$那時候我們讀英文中學.
就要改名 改英文名.
回去問家人.
什麼都不懂.
問家人 姐姐就叫加爾文.
我不知道是什麼來的.
是入神學院才知道.
原來是16世紀的宗教改革家.
僅次於馬丁路德.
其實馬丁路德是同時期的.
加爾文是整個瑞士的一個改革家.
整個瑞士的大復興都是因為加爾文.
你可以想像到那個震撼性.
但是加爾文神學最重要的地方就是.
神的預知和神的揀選.
神的揀選在信徒身上.
他是一定會信的.
一生的持守他的信心.
所以他們一生的持守.
就在於他們的工作倫理.
因為離開了天主教.
失去了那個Blessing.
我們說過了.
教會以外沒有救恩.
這是天主教很重要的理念.
你離開教廷就沒有救恩了.
你們離開了天主教.
天主的那個Blessing.
基督教說不是的.
我們靠主耶穌就行了.
因信請義和為了信心.
所以怎樣能夠使得弟兄姊妹.
中世紀的時候他們能夠對自己的得救.
對自己的神揀選.
有清楚上心呢?.
就是你在你的工作裡面.
其實工作本身就是那個召命.
每個人神都安排他的工作.
所以.
你是何詩蓓.

$^{761}$你是游水.
我是做傳道.
我就講明白聖經.
大家在不同的崗位.
去接受這個Calling.
在這個Calling裡面.
去忠於上主.
你就知道自己是被揀選.
所以這裡告訴我們.
相對於真理的啟示.
更多就是假冒.
假冒就是從那個可能是神蹟啟示.
可能是會有一些超自然的現象.
其中我們要留心的就是.
不要動心,不要驚慌.
因為那個誘惑.
誘惑就是說令到你.
有一種心思思.
很想真的去想.
可能真的計算一下.
都遲了.
主再來.
還有那個人真的越來越行神蹟.
他講的話真的很有紋路.
那些異端的洗腦就是這樣的.
特別邪教就更加嚴重.
我們在過去早幾個月.
國民慈牧師的神學院.
就是播道神學院.
上過一個課程就是神學的初階.
關乎系統神學.
他就放了一些韓國的邪教.
給我們看在課堂上.
原來韓國都有很多邪教.
就是這樣來洗腦式的讓信徒.
而且那些信徒全部不是那些不識字.
沒有學問.
是很厲害的人都有的.
原來誘惑.
就會不知不覺轉入了我們的心.

$^{801}$就是我相信一定是內心裡面.
對傳統的真理.
好像吃慣了媽媽的餐餐都是梅菜.
煮豬肉就已經味道脹立了.
不好吃.
出外面吃就有不同的口味.
就會很新鮮.
相信這份新鮮感.
對於傳統的教會弟兄姊妹.
就會有這種衝擊.
所以為甚麼保羅是用誘惑這個字.
誘惑就是你很想試一些新的東西.
你從來沒有經驗過的.
就會帶來有一種不知不覺.
你下了踏你自己都不為意.
所以保羅就告訴我們.
因為那日子以前.
必有離道返教的事.
原來那些人帶我們思想主宰來.
原來那些是離道返教.
是帶我們去跟那個敵基督者.
是跟那個褻瀆神的人.
自稱為神.
自稱為自己是那一位獨一的主.
不是告訴你.
你不要這樣看我.
我不過是神的僕人.
任何一卷新約.
無論使徒甚至連天使.
都是這樣理解自己的.
就是這樣的異端.
這樣的將來的迷惑.
就是指向終極的.
他說自己是神.
這樣就可以想像到.
我們必須要小心留意.
就是這個大罪人.
離道返教大罪人沉淪之子.
就是不斷他自己沉淪.
他自己是離道.

$^{841}$他就會帶領那些他的信眾.
就是他可以洗腦的信眾.
就會離道.
就會返教.
所以我們真的要小心.
一方面我們要聽真理.
我們要對恆常的真道.
我們要有那種屬靈的胃口.
和那個渴慕的心.
然後我們才能夠分辨到.
就是說一定會顯露出來.
一定他們的行徑.
保羅已經講得很清楚.
和主耶穌在《福音書》講的.
是完全一樣的.
所以離道返教.
是對基督教的一個攻擊.
尤其是給一些可能已經沉睡了的信徒.
要小心的一個提醒.
當然這裡主要是在講那個時候.
保羅在講那個時候.
保羅還沒死.
已經有人開始傳.
說主的降臨已經開始了.
所以保羅就說千萬不要動心.
千萬不要受迷惑.
然後描繪這個人.
第四節.
他是抵擋主.
高抬自己.
超過一切稱為神的.
和一切受人敬拜的.
甚至坐在神的殿裡.
自稱是神.
所以第四節可以作為我們去敬拜神的一個.
正常的敬拜生活的一個態度.
因為一切稱為神的唯有是主.
受人敬拜的也是主.
我們去到聖殿.
神的殿.

$^{881}$我們對著這位造物主.
我們有沒有稱祂為神去敬拜祂呢?.
去重揚祂呢?.
我們有沒有因為來到.
你坐得不舒服.
你就開始很多很多胡思亂想的想法.
又怪責這個怪責那個.
將整個焦點都放在敬拜裡面呢?.
第四節就是我們當頭彈劾.
不過.
這裡就是在說將來那個大罪人.
就要人敬拜他.
他明明自己是大罪人.
居然說自己是神.
居然有這麼多人.
完全沒有那種雪亮的眼睛.
看他是神.
拜他.
你想想是多麼的罪大惡極呢?.
那個人本身已經是該死的了.
說自己是神.
在神的殿裡面.
但是居然會迷惑了一些信眾.
看他是神.
而拜他.
拜一個很現在的.
甚至連偶像都不是.
活生生的一個人.
你拜他是真神.
你想想是多麼墮落到這樣的田地.
但是如果我們反過來.
當然這樣抽開上下文.
上下文就不是在說敬拜生活.
是在說千萬不要拜這樣的人.
但是.
反過來就是.
我們稍微退後一點.
敬拜本來就是拜神.
我們本來去到神的殿.
所以我們的教堂真的是神的殿.

$^{921}$一個二樓的教會.
奉主名聚會.
每一次主席宣告.
你不要以為他在開玩笑.
現在我是奉聖父,聖子三一神的名.
來去聚會.
就有神在.
所以我們來是敬拜神.
所以一個二樓教會.
是敬拜神.
一個座堂的教會.
也是神的殿.
大家一定要這樣理解.
你不要去到梵蒂岡才是神的殿.
是很麻煩的.
你每個星期坐飛機過去都不行.
你搬去住都不行.
所以.
耶穌已經說了.
祂的出現.
祂的離世.
祂的升天.
你們敬拜也不是在耶路撒冷.
也不是在這裡基里森山.
因為敬拜神要用什麼?.
心靈和誠實.
是在靈裡面.
在真理裡面敬拜.
所以每一次你去到那個地方.
只要你是奉主的名.
那個地方就是神的殿.
就是敬拜神.
所以第四節.
也是作為一個敬拜的反思.
我們去到.
不要拜錯了人.
今天李斯敬來.
不好意思.
我們去聽.
對的沒問題.

$^{961}$聽到是好的.
但是.
我們去到神的殿.
敬拜神.
不要拜一個人.
也不是拜一個.
我們也有責任的.
經常找一些名人來講道.
好像突然間會開心一點.
其實我們也是想.
多一些牧者.
不同的心靈的教導.
去有不同的詮釋.
那個背後.
所以千萬不要是一個名人的效應.
重點仍然是大家成熟地.
我來到教堂.
至少我真的要立即吞一吞自己的心靈.
我來到朝見神.
敬拜神.
所以穿得漂亮一點其實是好的.
穿得漂亮一點是代表我們內心的敬虔莊重.
當然穿得漂亮一點是想炫耀.
那就另一個問題了.
就另計了.
那個就不對了.
所以我們要記住這件事.
原先第四節不是這個意思.
不是教我們敬拜神.
是在講將來必然會去到一個極端.
有一個人夠膽說自己是神.
他真的在神的殿裡自稱是神.
而且一定有人去拜他.
他才能夠這樣講.
如果不是.
尤其在青山裡面也有些人說自己是神的.
但是不會有人拜他.
你就可以想象到將來主在來之前的黑暗之夜.
就會出現.
一直到現在.

$^{1001}$本來也沒有說在他有生之年.
這個人就會出現.
不過他告訴我們.
似乎也是有靈感裡面.
聖靈的感動裡面.
他會知道那個狀況就會是這樣.
在主在來之前的黑暗.
所謂暴風雨.
就會是這樣的形態.
就是會有人.
令一些對真理已經不再感興趣.
真理打動不了他的心的人.
就會受迷惑.
受影響.
這個真的是一個很大的問題.
所以.
那個人說自己是神.
固然他一定要為自己負責任.
但是你聽他講.
跟隨他.
然後你以為自己在走真理.
那個問題背後.
下面的保羅就說了.
第五節.
我還在你們那裡的時候.
曾把這些事告訴你們.
你們不記得嗎?.
似乎保羅在他一直的教導裡面.
很重要.
因為使徒真的有很特殊的權威.
使徒真的有特殊的行神跡的能力.
你有沒有讀過《使徒行傳》呢?.
彼得的影子也很治好人的.
彼得的手帕也可以叫人得到醫治的.
保羅看到在《使徒行傳》的十六章.
菲勒比有一個侍女.
被鬼迷的.
被鬼附的.
可以講預言.
說這兩個人.

$^{1041}$保羅和西拉.
他們是神.
是至高神的捧人.
對你們講的是救世之道.
保羅覺得很煩.
就斥責他趕走那隻鬼.
那隻鬼是能夠幫主人發財的.
所以你可以想像到.
保羅很清楚弟兄姊妹.
如何理解他們這一班的使徒.
是好像有特殊的榮譽.
好像是另一個階層.
好像所謂聖品那樣.
其實保羅就很清楚告訴他們.
我們不過是神的捧人.
這個哥羅多後書就不斷講.
他自己是一個很謙卑的捧人.
甚至乎是一個不配的捧人.
這樣理解自己.
但同一時間.
神有時叫他行一些很超自然的神蹟.
是要令到當時的人.
不是受吸引.
是能夠見證那個福音的大能.
所以他們是徘徊在神國那種能力的彰顯.
和他們自己對自己的認識.
對自己的認識.
是因為對主耶穌的榮耀的身份的認識.
而這樣看自己的事奉的心態.
所以如果我們是這樣理解.
其實我們根本不知道什麼叫做事奉.
就是說不提你的名字.
不講你的名字.
你就發脾氣.
那怎麼可能你還會明白事奉.
我們能夠事奉已經是很大的一個榮耀.
基本上是不應該我們事奉的.
我們不配的.
這樣是一種深深的.
是真實內心的一個心態.

$^{1081}$不是作假.
是真的不配.
不過既然神驗證我們.
選擇了我們.
我們就盡力終身去做好.
就是這個意思.
所以保羅就是徘徊在這個心態.
所以保羅就說.
我在你們當中教導你們的時候.
我都這樣講.
就是說他不是基督.
他也不會扮演基督.
他不是救世主.
他更加知道自己的限制和軟弱.
所以在他們當中.
保羅真是恰如其分地將自己的僕人身份.
和他所領受的對外邦人的託付.
就是要將耶穌是救世主的這個福音.
告訴外邦人.
他要盡力將自己猶太人.
對外邦人的義務性格清除.
一個猶太人在外邦人是豬狗.
不配領受神國.
猶太人是神選民.
所以保羅從小就有這些優越觀念.
等於這個世界還有貴族.
我是貴族 我的出生.
一出生家裡有幾十個工人.
是貴族的階層來的.
有奴隸的.
如果你是一個奴隸.
在你的家族裡面.
你可以想像得到你會怎樣看.
但是現在保羅如果他是一個貴族.
但是耶穌改變了他.
他就已經沒有了猶太人的這種優越.
這種錯誤的價值觀.
他要盡力向外邦人展現他是弟兄的一個身份.
這樣來講解福音.
你可以想像到保羅就是這樣的一個心情.

$^{1121}$所以你們還不記得了.
第六節.
現在你們都知道這位難了的是誰.
是叫他到了的時候才可以領路.
因為那不法的隱義已經發動.
只是現在有一個難了.
也是剛才第六節所講的那個難了.
等到這個難了的被除去.
那時這不法的人必顯露出來.
就是上第四節所講的.
那個人自稱是神的那個不法的人.
主耶穌當然就是主再來的時候.
似乎是再來之前.
主耶穌要用口中的氣滅絕他.
用降臨的榮光廢掉他.
到底是不是同一時間呢?.
天使掌呼叫.
神的號吹響.
主耶穌又再呼喚.
跟著是不是同一時間.
那個不法的人.
那個坐在神的殿的自稱是神的人.
就被主的榮光來廢掉呢?.
主耶穌口中的氣滅絕.
你可以想像到.
差不多是兩個敵對陣營.
真理的陣營.
和不法的陣營.
但是是召集著很多人的.
是很大的集團的.
是很強烈的對神國度的.
那種褻瀆和攻擊的.
但是保羅說.
不用打十年八年.
到時候也不知道慘勝還是怎樣.
還是有多少人戰死沙場呢?.
是不需要的.
保羅說.
主耶穌的榮光就可以廢掉.
口中的氣就可以滅絕.

$^{1161}$只不過是一個時間.
只不過是容讓罪惡繼續在地上.
好像肆虐肆無忌憚的.
這樣來迷惑人.
原來口中的氣.
光明的榮光.
就已經足夠將整個罪惡瓦解殆盡.
所以我們很難理解.
那你做這麼多是幹什麼的呢?.
現在有榮光.
現在口中的氣就不會有罪惡了.
大罪人就不會出現了.
不出現就好了.
但是神的工作不是像我們這樣想的.
很多時候我們就說.
你不要擺出分別善惡樹.
人們吃不到.
阿當吃不到.
夏娃吃不到.
那就不會犯罪了.
那就不需要有救恩了.
每個人都可以什麼的.
我們很容易就會用我們自己現在表面的想法.
但是我們稍微限制一點點的東西.
叫他不可以選擇.
他就跟你拼了.
根本就沒有可能.
他以為你收回我們的自由意志.
讓我們不用選擇.
叫他不要選擇.
他已經不行了.
當然我們要尊重每個人的思想.
在認識神之前.
都有很多的思考.
這些正面的是好的.
但是有時候為反而反.
這種不拉肩的抗議.
其實根本你認真回答他.
是浪費力氣的.
浪費你的時間.

$^{1201}$你認真回答他.
原來他根本就是很輕佻的.
其實是不需要的.
當然小朋友他們很單純地去問這些問題.
我們就應該要很單純地讓他知道.
其實你這個問題問得很好的.
正正就是容許我們自由的選擇.
就是說我們可以選擇善.
可以選擇惡.
這樣才是正常的一個人.
神對我們的尊重.
就不是給我們沒有選擇的.
你只會選善的.
其實一選善就是選擇.
就是自由意志.
你沒有選擇.
連一個都沒有選擇.
那根本就不是人.
所以這些表面上是很哲學的東西.
其實是很簡單.
每個人必須一定會面對這件事的.
所以神就是全知.
正正知道我們人只會說話.
收回這個自由意志.
根本就不變成人了.
所以神是容讓阿當.
可以聽不聽神的命令.
吃不吃那個分別的蟬膜,樹果,罪.
就因為蛇的引誘就進入了人的裡面.
所以就會開展神的救恩計劃.
所以這些告訴我們.
現在我們就要去思想這個難阻.
這個難阻被除去.
第八節就說這個不法的人就會顯露了.
不法的人和第七節的不法的隱秘.
隱藏的一個意念已經發動.
所以是慢慢越來越明顯.
越來越那個人就成形.
越來越那個人是肆無忌憚的.
坐在神的殿自稱是神.

$^{1241}$保羅的年間.
當然他肯定知道.
曾經在猶太人差不多滅國的時候.
就是主前的165.
可能是166的時候.
聖殿曾經給一個君主.
敘利亞王.
他叫埃提亞古四世.
這個人.
其實他有另一個名字叫以比凡尼.
就是君臨天下的意思.
他打贏仗.
去到耶路撒冷.
去到聖殿.
接著他要在聖殿獻一隻豬.
來去褻瀆.
污穢聖殿.
所以猶太人當時有一個家族叫馬加比家族.
就看到這件事是極之嚴重.
所以他們誓死抵抗.
最終打敗了埃提亞古四世.
撤軍.
所以聖殿.
當時據聞說那個聖殿.
不知為何.
聖殿的金燈塘的光.
一直都有油沒有停過.
一直都點著.
所以他們去到聖殿的時候.
發現一件這麼奇妙的事情.
他們就潔淨聖殿.
就是重光聖殿.
所以就在聖殿裡面點燈.
所以就叫做重光節.
或者叫做聖殿的修殿節.
《約翰福音》十章.
耶穌就在所羅門的樓下.
那個時候就是修殿節.
哈魯克就是每年的12月.
猶太人的.

$^{1281}$就是現在的洋曆的12月.
就是哈魯克點燈.
重光節.
叫做重光節.
或者又叫做修殿節.
就是保羅都很熟的了.
因為保羅都是主後的.
大概30多至40多至50多年的人.
所以那段歷史不是很久的.
100年前左右.
就是那個重光節.
就是埃及雅古四世.
他又叫做大罪人.
就去到神的殿裡面.
又自稱是神.
又是獻珠.
所以每個時代都有那些大罪人.
所以保羅的年代.
他如果現實的.
就是尼祿.
尼祿差不多開始嫁禍基督徒.
就說舊的羅馬城.
是基督徒放火燒的.
其實就是他洗了手.
找一些可能是黑社會.
燒了舊的羅馬城.
就想建一個新的羅馬城.
就被元老院禁止他.
他就要嫁禍給基督徒.
所以就大規模逮捕基督徒了.
就是這個尼祿王的時候.
所以保羅就在那個年代.
被尼祿捉拿.
監禁.
最終都是殉道.
所以在這裡告訴我們.
我們要記住.
特別關乎到主耶穌降臨.
以及我們將來去到永恆的聚集.
一定會有一個似乎是一個時間.

$^{1321}$就是不法的隱形.
以及不法的人會出現.
保羅就不是說他有生之年就會有這個人.
可能這個不法的運動.
是透過不同的時代都會出現.
都會來攻擊.
比如希特勒殺600萬的猶太人.
某程度上是大罪人.
對那個時代來說.
每個時代都有他們這樣的情況.
所以我們就要有這種意識.
主耶穌回來之前.
是會有尼祿反教的.
這個似乎是任何時代都會出現的.
所以我們要留心.
保羅這裡就很含蓄地沒有說到具體說.
到底是教會.
未被除去.
仍然在阻礙著大罪人的全面出現.
還是福音.
還是我們基督徒.
也不是很詳細的.
第六節這裡說.
現在你們都知道.
那個攔阻他的是什麼.
是叫他到了時候才可以顯露.
他就是那個大罪人.
第七節.
因為不法的引義已經發動了.
只是現在有一個攔阻的.
等到那攔阻的被除去.
所以是不是教會被提呢?.
基督徒被提呢?.
答案似乎是.
但是始終要保險.
保羅沒有說得那麼清楚.
到底是地上的真理.
或者是主耶穌.
這裡不詳細說得清楚.
所以我們只能夠說.

$^{1361}$這個大罪人原來是有時間才出現的.
原來現在還沒有全面出現.
當他出現的時候.
也是他滅亡的時候.
真正很快就要被主的榮光去廢掉了.
被主耶穌口中的氣去滅絕了.
這個就是我們要留意的事情.
說到第八節的那時.
就是那個攔阻被除去.
所以在這裡有些神學的論述.
就是關乎末世論.
就是說教會被提了.
這個攔阻被除去了.
就是說現在教會存在在這個地上.
是攔阻住那個大罪人的全面出現.
有一派會是這樣理解的.
這個就是災前被提的理論.
七年大災難.
就是對應啟示錄所說的那個教導.
但是要記得.
那個大罪人一出現.
主耶穌的降臨就毀滅他了.
所以我們要看到.
原來罪惡無論多麼黑暗.
都在主的旨意.
在主的預定裡面.
和在主的.
不是打開了那個門.
收拾不了.
沒得搞了.
他已經是肆無忌憚.
不是的 主的降臨.
榮光和口中的氣.
已經可以收拾了.
就是說簡直是易如反掌.
口中的氣.
和降臨的榮光.
都不是主自己的.
不是他自己要一鎚打他要打多少回合.
不是的 榮光.

$^{1401}$口中的氣.
就是說意味著.
絕對不是我們用一些物質去想.
那些口氣很重.
降臨的榮光很殘暗.
我們不用這樣理解.
就是說在主的時間到.
這個大罪人.
這個沉淪之子.
這個攔阻地上的人.
敬拜真神.
敬拜他這個人.
是會完全滅絕.
完全去到永恆的火葫蘆.
就是那個獸和假先知的地方.
也是魔鬼要去的地方.
所以有時候我們想想魔鬼學.
《神學院》有一讀的.
有一個關乎到.
什麼叫做天使.
什麼叫墮落天使 就是魔鬼.
魔鬼現在的情況是怎樣.
他的形態是怎樣.
和他的結局是怎樣.
有一讀的 有這些神學的.
其實《聖經》很多的對於魔鬼.
其實很多教導的告訴我們.
首先我們要記住.
《聖經》肯定是有靈界的存在.
是我們看不到的那個世界.
不過不是中國人說的那些靈界.
《聖經》真的告訴我們.
有屬靈的世界.
有屬靈的世界 也有靈體.
所以天使是靈體.
所以我不知道為什麼.
人們會看到有翅膀.
我從來沒有見過.
我也很想見 但是沒有需要.
也不需要羨慕.

$^{1441}$可能就是舊約有說到.
關於天使的出現.
譬如基督帕.
有三對形.
一對遮眼.
一對遮腳.
一對是用來飛翔的.
你可以想像得到.
原來三對是有它的作用的.
無論是《倚創》的書也有說.
或者是《以色列》的書也有.
這裡告訴我們.
這個不法的人.
可言是有天使的身份.
不法的人.
先喝一口茶.
(主持人:是的 真是好).
因為沒有停下來.
(主持人:我也想提醒你).
剛才上文說的是.
不法的忍耳這個運動.
現在第九節就很清楚.
更加具體說.
是這個撒旦的運動.
一定是黑暗的靈界.
黑暗的靈界.
對神 做物主的攻擊.
所以原來魔鬼是不會再悔改的.
是他根本不會悔改的.
天使也不能悔改的.
因為天使是沒有救恩的.
只有人類才有救恩.
這些我們都需要知道.
天使只不過是服役的靈.
它有不同的階級.
但都是服役的靈.
是為信主的人效力的.
所以有些神學就說.
顯生我們每一個人.
都有自己的天使在天上.

$^{1481}$當然聖經沒有說得那麼詳細.
所以我們要記得.
我們要留意.
原來這一切.
就是撒旦的運動.
所以不怪那個大罪人.
似乎有一定的能力.
能夠行各樣的異能.
神蹟.
和一切虛假的騎士.
這些異能神蹟.
我們都要用在真理的尺度裡面.
去驗證.
所以不要有神蹟.
有異能.
就一概傳言.
照單傳修.
必須要驗證.
就算是神蹟也要驗證.
是不是出於神.
因為這裡已經說了.
原來神還容許.
容讓魔鬼.
有一些如魚靈.
能夠行神蹟.
其實它的動機.
不是將人帶到神那裡.
而是將人帶到錯謀.
帶到那個大罪人.
帶到魔鬼的撒旦的運動.
那一定是敵黨.
撒旦的另一個名字就是敵黨.
抵黨.
所以撒旦這個字.
本身就是敵黨的意思.
所以我們要記住.
如果不是將人帶到神裡面.
去行讀耶穌基督的福音.
和對主耶穌的敬拜.
那些神蹟絕對沒有任何的效力.

$^{1521}$沒有任何的值得羨慕.
因為它會令到你.
去跟隨了虛方和撒旦.
所以那些不同的民族.
都有一些超自然的現象.
我們要小心.
不要去亂說.
因為我們是在.
傳教的.
我們要小心.
不要被那些現象.
去迷惑了那個心靈.
第十節.
並且在這些沉淪的人身上.
行各樣出於不義的詭詐.
各樣.
但是動機是不義的.
是詭詐來的.
以致你信了它.
你就會受到影響的了.
是會令到你進入.
不是那個真理的道路.
越來越明白真理.
是會令到你走下去.
一個不義的那個結局裡面.
後面那裡我們都提過了.
原因就是.
你自己本身已經不再喜歡真理了.
使得我們得救了.
所以你就會很傾向於那些.
似是而非.
令到你很興奮的.
一些學說.
所以一些不再回教會的弟兄姊妹.
他們越來越不讀聖經.
他們就會越來越喜歡聽那些.
無無畏畏的東西.
講講佛偈他們也很有興趣.
講講迷信的那些風水.
他們也很有興趣.

$^{1561}$覺得很靈驗.
他們已經沒有了一個愛真理的心.
可以得救.
可以得救的意思不是說.
你終於可以洗禮這樣的意思.
越來越對真道.
有一種融會貫通.
有一種本身就好像住在真道裡面的.
那種心境.
就是你經常出入的那些地方.
你就會有那種氣質.
就是現在我們奧運.
整個世界.
現在就準備買劍.
學什麼.
然後就去學游泳.
就是在那個氛圍裡面.
其實就是那個意思.
我們不領受愛真理的心.
不過愛就很重要了.
我們原來對真理是完全不感興趣.
就算是誰說的.
我也不感興趣.
慢慢這種不感興趣.
是令到你連回教會都不想再回.
甚至到了一個地步.
有關乎教會的人事物.
你都感到討厭.
感到作嘔.
這樣的情況.
其實就很危險.
那個時候就因為心靈已經真空了.
就會傾向那些不義.
讓他感到很迷惑的事.
就會受到那個影響.
所以第十一節.
故此神給他們一個生發錯誤的心.
叫他們順從虛妄.
使一切不信真理.
賭喜愛不義的人都被定罪.

$^{1601}$故此就是上文告訴我們.
如果我們沒有去耕耘我們自己的內心.
讓神的真理在我們的內心.
來不斷的洗滌,提醒,勸告.
我們就會慢慢傾向於一個不喜歡真理.
賭喜歡不義的心.
譬如你看到保羅在提摩泰後書提到.
有些人.
保羅就叫提摩泰.
務要全道.
得時不得時都要專心.
因為很多人在末世厭煩純正的道理.
耳朵發癢.
就是說不要再聽真理了.
喜歡聽那些.
我讀給大家聽聽.
提後的三章.
教提摩泰.
他要這樣做.
對不起 第四章才是.
第四章.
我在神的面前.
並在將來審判活人死人的基督耶穌的面前.
憑著祂的顯現和祂的國度去祝福你.
這個祝福有很多形容詞.
將來審判活人死人的基督耶穌.
關乎到祂的顯現和祂的國度去祝福你.
就是說你是將人帶進神的國度.
他是活人還是死人.
他是活人將來就得到賞賜.
如果他是死人.
將來就要面對永恆的沉淪.
就是憑著你的國度.
所以傳道者其實一點都不好玩.
一點都很嚴肅的一個職份.
阿母要傳道.
就不是叫提摩泰.
那你不如不要做傳道了.
你這麼懦弱.
阿母要傳道.

$^{1641}$因為本來自己要走.
自己要回天家.
要宗道.
母要傳道.
無論得時不得時.
總要專心.
專心就是專注於你自己.
你要對自己的時代.
對自己那班的回眾有一個敏銳.
譬如你只懂得做乾炒牛河的.
你就不要在那裡扮人去做什麼.
做片皮鴨.
你就是自己拿來壞.
你應該做好你自己的乾炒牛河的責任.
母要傳道.
要專心.
要洞察你的回友的需要.
那個時代的需要.
那個時代.
那個季節.
現在是夏天.
你就不斷燉那些熱飲.
大家都不想吃了.
很不想.
你當然要做一些清涼的.
爽口的.
時令的飲品.
所以總要專心.
是在說你對那個時候的敏銳.
對弟兄姊妹的敏銳.
並用百般的忍耐.
所以就要這樣做.
一個傳道者.
最難的就是你敏銳弟兄姊妹的心靈.
所以要百般的忍耐.
以及各樣的教訓,責備,警戒,勸勉.
所以是很平衡的.
為什麼呢?.
因為時候要到了.
人必厭煩純正的道理.

$^{1681}$保羅現在很明顯地教提摩太.
不是那個迪基督來的.
而是時候.
其實保羅還沒死.
人是必定會厭煩純正的道理.
以躲法養.
很全神.
聽.
因為當時是一個說的.
傳遞真理的過程.
是傳道者去傳的.
是用口說的.
當時是沒有那麼多聖經.
他們每個人家裡自己會有一本聖經.
是靠教會聚集的地方.
傳道者去口述說出來的.
所以以躲法養.
根本是要離開那一面.
不想聽你說.
很快.
說來說去都是那些東西.
不想聽.
當然.
說來說去那些東西是傳道者的問題.
不是那個以躲法養的問題.
所以你不要看到有人嘮叨.
就好像不是聽你說話.
所以其實保羅.
因為你必須要上文.
你自己是專心的.
你有百般的忍耐.
以及你已經有各樣的教訓.
就是說你已經有充足的心靈和準備.
去說婦口婆心的勸勉.
使得更多弟兄姊妹.
能夠得到教導.
但是保羅也說得很清楚.
安泰你要記住.
你所做的.
是真的會有怨聲的.

$^{1721}$因為以躲法養.
就會隨從自己的情慾.
就是說它裡面有太多東西.
好像滾動.
很多東西.
它會有很多內心的想法.
會使它很雀躍.
每樣東西都覺得很新鮮.
很想去試.
唯有神的道.
不想去認真去聽.
去遵行.
所以它會增添好聲的師父.
並且掩耳不聽真道.
偏向的.
所以不聽就不是.
我很客氣的.
我回來很安靜的.
我很欣賞這個寧靜的星期初.
我坐在這裡.
我就睡得著.
我就能夠很安靜.
我在六天的日子.
我很煩亂我的心靈.
我需要的是詩歌給我洗滌我的心靈.
所以你說道.
我沒有意思的.
你繼續說.
不過我是不會聽的.
原來不是的.
你不聽真道.
其實是偏向荒謬的言語.
所以不是開玩笑的.
我們不會對神的道有一個渴慕和追求.
就會偏向荒謬的言語.
所以提摩太就說.
你卻要凡事謹慎.
忍受苦難.
作全道的功夫.
記住.

$^{1761}$又是那個Calling.
盡你的積分.
這個就是嘉義文神學的精髓.
或者說清教徒離開了天主教之後.
那個神學很重要的核心.
就是那個照明.
盡你的積分.
別人聽不聽.
你盡你的積分.
你令到它由荒謬變成喜歡真理.
這個就是我們一生的責任.
我們回到貼索羅尼加後書這裡的時候.
我們可以完第二章的經文.
我們很快說一說.
當然我們上一星期提過.
是一個相對的.
主所愛的弟兄.
因為從起初揀選了你們之後.
叫你們恩順真道.
又被聖靈感動.
成為聖潔能而得救.
神藉我們所傳的福音.
召你們到這地步.
好得著我們主耶穌基督的榮光.
所以主降臨的榮光.
在那些不信敵黨神的人.
跟隨大罪人的人.
就是滅絕的榮光.
消滅的榮光.
但是對於跟隨福音.
能而得救的人.
主的榮光已經在我們大家的心靈裡面.
已經是上加榮了.
我想跟大家讀一讀.
哥納多後書的第四章.
應該是第六節.
這節聖經很重要.
第六節.
四章的六節.
長話短說不可以太長.

$^{1801}$將第四章講解.
第六節.
「哪吩咐光從黑暗裡照出來的神」.
當然就像宋世紀那樣.
要有光就要有光.
那個大概是典故.
「哪吩咐光從黑暗裡照出來的神」.
已經照在我們心裡.
就是說原來我們信主的人.
是有神的榮光在我們內心.
然後.
「叫我們得知神榮耀的光.
顯在耶穌基督的面上.
雙重發光.
我們自己既有神的光.
因為我們已經離開了黑暗.
進入了光明的神的兒女.
就是這位造物主照在我們的心裡.
然後.
「叫我們得知神榮耀的光.
已經照在耶穌的面上」.
我舉一個例子.
大家現在很少看電影.
因為現在家家戶戶都上網了.
跟我們小時候一樣.
一定要去戲院.
不能讓你自己私人.
自己塗鑽.
在家裡就可以看電影,看電視.
又有YouTube.
我們要去.
小時候不會有手提電話.
手提電話.
很多偉大的發明就是演唱會最後.
大家開著電筒.
一起照.
這樣轉的.
這是一個很偉大的演唱會的發明.
因為每個人都有手提電話.
我們回到沒有手提電話的世代.

$^{1841}$我們進戲院.
你會跟誰坐下你的位置呢?.
那個大委員.
是他才最偉大的.
他有一支光管.
他有一支電筒.
還有戲院特別黑.
每個人都沒有燈.
靠他照著你的位置.
帶著你走.
第六節就是這個意思.
原來我們大概都有這種責任.
我們將我們自己裡面的光.
神已經照在我們裡面.
我們要告訴別人.
神榮耀的光已經照在耶穌的面上.
我們自己都要這樣看到.
我們就有那個電筒.
我們就照著耶穌的那個.
其實不是我們照.
其實是神的光已經照在耶穌的面上.
所以我們為什麼會越來越愛主的弟兄姊妹.
會越來越看到耶穌基督的偉大.
任何一個時代的基督徒.
那些牧者.
如果他看到耶穌的偉大.
他就會心悅誠服的.
至死忠心.
為這位主的福音去效力.
因為我們已經看到神榮耀的光.
照在耶穌的面上.
我們就要告訴那些微信嘲笑耶穌的人.
告訴他們.
其實你不明白.
耶穌現在是神榮耀的光.
照著的那一位唯一的偉大的救主.
所有人的見證都是為這個榮光做的.
所以我們一定要將我們的榮光.
不要照在我們自己的面上.
你看我多厲害.

$^{1881}$你看我多偉大.
你不找我事奉.
你不找我講道.
那你真是你的事了.
不是這樣的.
我們是要將我們的榮光帶給人看到耶穌.
所以第六節是很可圈可點的.
如果有人問我們.
為什麼你這麼開心.
你信耶穌.
就是因為第六節.
我們看到主的榮光.
神現在給我們這個恩典.
我們可以得到這份榮光.
我們要告訴你.
主才是真正的榮耀.
意思就是這樣.
所以保羅說.
真是兩極.
其實上一次我們也講過了.
不過這次我們再詳細一點總結.
第十三到十七節.
是一群很單純.
愛主.
在保羅的內心.
這一群是主所愛的弟兄姊妹.
因為你們從開始就已經被揀選.
而且這個被揀選.
是因信真道.
又被聖靈可以感動.
成為聖潔能言得救.
保羅這樣稱讚貼壽羅尼加的教會.
就是說一方面你不要動心.
不要受那些假的主在內影響.
你們自己本身其實已經有了這個榮光.
已經照在你們的內心.
現在你要帶人去看到耶穌的榮光.
所以一定要有這個持守.
神所傳的福音.
仍然是所傳的福音.

$^{1921}$所以不要小看福音.
很簡單.
恩信清義.
信耶穌得永生.
其實不是這麼簡單.
就是這位主的一生將來帶進救贖裡面.
所以福音照我們去到得著耶穌基督的榮光.
所以要站穩.
凡是所領受的教訓.
不區是口傳.
是信上寫.
都要堅守使徒的教訓.
就是說對真理的渴慕.
對真理的辨別.
因為上文是口傳的.
信上寫的.
二能的.
不要輕易動心.
但是保羅說.
這裡不是輕易動心.
是要測驗這些真理.
接著就要持守使徒的教訓.
要領受.
所以是領受.
領略之後.
真的要完全接受.
不要聽完就根本.
我沒有想過要聽你走的.
我自己有一套.
神未必能夠拆毀你.
所以我們要一個不斷謙虛的去領受神的道.
要持守.
所以要堅守.
第十六節.
但願我們主耶穌基督.
和這位愛我們開恩將永遠的安慰.
美好的盼望賜給我們的父神.
安慰你們的心.
由始至終都是說安慰.
開恩將永遠的安慰.

$^{1961}$安慰你們的心.
並且在一切的善行善言上.
兼顧你們.
也是對神的道.
那個渴慕.
有善行善言.
有美好的良善的行為.
不同的榜樣.
不同的見證.
我們來學習.
一起去堅守.
你可以想像到.
十三到十七節.
是非常正面的欣賞.
這群初信的.
貼十來黎加的弟兄姊妹.
他們就是這樣信忘愛.
他們就是這樣對保羅的教訓.
有一種這樣的學習和持守.
所以我們也求神幫助我們.
在不同的時代.
透過現在的網絡.
不同的地方都可以聽這些錄音.
一起學習.
甚至可以在現場.
但是重點就是我們對保羅的教導.
是要有這個學習和持守.
我們一起禱告.
再一次我們多謝主.
讓我們能夠再看回深入的第二章.
有兩個極端的對比.
就是一些不聽真理的人.
在那個不法的忍耳.
殺人的運動.
就會被迷惑.
被影響而離開真道.
就是十三到十七節.
保羅對貼十來黎加教會的欣賞.
和羽從心腸的勉勵.
他們繼續在神的道上去成長.

$^{2001}$盼望我們在《查經班》裡面.
我們不斷其實累積了很多.
對無論是新約.
或者是一些聖經人物的教導.
很盼望我們身體力行地去揣摩思想和持守.
使我們行在主的真道上.
求你讓我們能夠清楚.
我們自己生命裡面最重要的照明.
以致我們用心地在我們有限的人生去服侍主.
將我們裡面的光帶人去認識耶穌基督.
才是真光.
才是神的榮耀的光.
求主幫助奉耶穌基督的名.
阿們.
\newpage



\section{帖撒羅尼迦後書}
\label{sec:bkamj1iGLMk}
\textbf{2024-08-15 周四晚查經班《帖撒羅尼迦後書(4)》 講員:曾裔貴牧師}
\newline
\newline
連結: \href{https://youtube.com/watch?v=bkamj1iGLMk}{\texttt{https://youtube.com/watch?v=bkamj1iGLMk}} ~~~~ 語音日期: 2024-08-16
\newline
\newline
\hyperref[sec:ewzMm4Y0R5I]{\small{< < < PREV SERMON < < <}}
~
\hyperref[sec:index_chronic]{\small{[返順時目]}}
~
\hyperref[sec:index_scriptual]{\small{[返順卷目]}}
~
\hyperref[sec:IJNXTQniUCA]{\small{> > > NEXT SERMON > > >}}
\newline
\newline
$^{1}$都是一起開始祈禱.
我們能夠再一次可以在晚上安靜的.
或者在線上都來到去.
在家中一同的來到去等候.
我們很盼望一些可能我們都曾經聽過.
有機會去研讀的經文.
再一起來到去靜心思想的時候.
能夠更加深入.
更加了解明白當時的歷史的狀況.
和當時的信徒他們到底是抱持著什麼樣的心態.
需要保羅來去在那個教導.
或者是一個警告上面.
對信徒有一個指引.
願神幫助我們.
的而且確我們不是活在那個時代的人.
但是我們都知道聖經那種跨越時代的原因.
就是因為這些大的原則.
就是人性面對聖經的真道.
面對神的這個拯救.
和永恆的那個天家.
都是每個人一樣的.
而且不分什麼種族.
不分什麼貴賤.
或者是厲害一點還是不厲害.
是我們的心怎樣對待神.
和神的真理.
願主幫助我們.
我們今天晚上等候在你的面前.
我們更加盼望每一段錄了音的教導.
是成為不同的弟兄姊妹的幫助.
我們一起在你面前等候,祈禱.
奉耶穌基督的名.
反覆我們也將第二章講得比較詳盡.
第二章其實是一個很危急的警告.
就是保羅突然叫人不要隨便信那些書信.
一定是冒名的書信.
或者有些人招搖撞騙的.
一些口述的來講到.
可能一些預言等等.
我們也都不是太知道當時.

$^{41}$初世紀使徒未全部離世之前的世代.
他們有沒有很特殊的經驗.
會有一些的默示等等.
這些都是一些很神秘的.
但是無論如何.
至少由保羅自己的口中.
這樣來警告帖撒羅尼迦的弟兄姊妹.
千萬不要隨便什麼文字.
或者是一些可能是他說的一些異象.
都不要相信.
不要相信的原因就是有來到反教.
就是有大罪人.
就是有一個風潮是會起的.
就是撒旦的運動.
這樣一路延伸就去到第十二節.
撒旦的運動是會吸引到一些人的原因.
就是原來有一些人是注定是不會明白真道.
因為他們的心根本就不想改變.
不想讓神的道改變他們的心.
所以這樣來作出警告.
這些人將來要面對自己.
從來沒有在地上為神的道.
在他們的生命來改變.
做過任何的宋土.
做過任何的晦改.
所以這裡說第十二節.
一切不信真理.
喜愛不義的人都要被定罪.
而且他們是會有一個被生發錯誤的心.
就是說凡是聽什麼進入心裡面.
好像攪拌機一樣.
正常你放一個蘋果進去.
你現在是出一些蘋果汁 是不是?.
但不是的.
放一個蘋果進去出一些橙汁.
因為他們已經有一個生發錯誤的心.
聽任何神的道都會有一些反駁.
都會有一些歪曲.
還有來對立.
原因就是本身他們有一個不喜歡真理的心.

$^{81}$所以這些運動為什麼能夠得到這些人的認同.
和追隨呢?.
說的是撒旦的運動.
所以正正就是那幫人.
根本從來就沒有一個晦改的心.
所以真的不要希望在我們有生之年.
教會我們所認識的朋友.
團契裡面所認識的好朋友.
如果他從來只是喜歡團契的那個商交.
那個老友記的聚餐.
一起去旅行 去踩單車的開心.
而他從來沒有想過要被真道去改變.
這個真的不好的.
所以是真的有這樣的警告.
不過13節到17節.
我們就很清楚看到.
保羅馬上轉一轉那個心情.
他說 你們是主所愛的弟兄.
應該的.
因為你們是從起初被揀選.
你們是因信真道被揀選.
因信真道而帶來聖靈的感動.
生命裡面成為了那個潔淨.
所以你們是能以得救的.
得救是說將來永恆的救贖.
先是靈魂的救恩.
先是在你們的內心有聖靈.
有這個永生的確具.
所以問心事.
就是要問你內心信什麼.
所以這裡就很肯定的將上文的那些不愛真理.
一個生發錯誤的心的.
可能是基督徒來的.
或者當然是聲稱自己是基督徒的人.
在教會裡面.
所以就導致到.
會有兩個就是耶穌所說的.
墨子和拜子的那個比喻.
在天國未降臨之前.
神好像容許兩樣的種子都是生長.

$^{121}$一來就是未成熟之前很難分辨.
所以這個都是做物主的智慧.
或者說都給足夠的空間人去思想.
悔改.
始終都不肯悔改.
那就真真正正證明.
他們是一個生發錯誤的心.
不喜歡真理.
最終的結局就是.
使一切不信真理.
喜愛不義的人都被定罪.
被神將來的結局性定罪.
是說將來的沉淪.
因為上文整個第二章.
都是說那個撒旦的運動.
所以不是三朝兩日那麼簡單.
是將來永恆的定罪.
但是我們要留意.
13到17節是很肯定.
帖撒羅尼迦的信徒.
他們是能而得救的.
因為他們明白外無真道.
所以有聖靈的感動.
所以他們就能夠.
內心先得到生命的改變.
穿上了新人.
有一個聖潔的生命.
所以第14節.
神藉我們傳的福音.
呼召你們來到這個地步.
所以是生命上的地位.
好得著我們主耶穌基督的榮光.
所以基督徒初信主的開始要追求.
不只是說我信了耶穌.
很開心了.
原來還有更大的前面.
就是主耶穌的榮光.
你要得到主耶穌的榮耀.
這位主曾經是無名小子.
曾經被人嫌棄.

$^{161}$曾經是木匠的兒子.
曾經在十誡上.
被人嘲笑.
曾經被人誤會.
他是褻瀆永生神的.
來褻瀆我們猶太教的.
但是這位主.
已經是肉身升天.
去到永恆神寶座右邊.
我們的主耶穌.
所以有時候真的要學習.
在我們自己內心的自律行間.
不要隨便呼呼喝喝.
耶穌.
這樣是不好的.
我們自己.
別人說當然沒問題.
我們應該要親切一點.
就是主耶穌.
或者你叫基督也可以.
因為名字.
耶穌已經是頭銜了.
所以基督已經是什麼了.
所以這些從我們的內心.
很自然對主有一種敬畏.
有一種羨慕.
很想去見他.
認識他.
所以保羅說的是.
基督徒的人生.
原來不只是追求.
我每個星期日.
回教會已經很給你面子.
你不要這麼要求高.
其實保羅說不是.
每個星期要求高的原因.
是你要得著主耶穌的榮光.
那個才是最高要求.
就是你和主耶穌之間.
你要有他那份的Glory.

$^{201}$要有一份的榮耀.
所以基督徒其實是一個很亮麗的人生的追求.
絕對不是一個重擔.
大家背著基督徒這三個字.
你會不會覺得是一個擔子呢?.
會不會怕人笑你是基督徒呢?.
會不會人認得你是基督徒.
就會覺得有些詫異呢?.
或者人見到你的時候.
都真的會感覺到.
你是真的有一份的平安.
是我沒有的呢?.
你是有一份的光彩.
是我沒有的呢?.
這個就是你反映到主耶穌的榮光.
保羅一直教導.
貼著羅尼加的信徒.
其實他們不是資深的幾十年的會友.
其實信主半年多一點.
你可以想像到.
福音使你們來到這個地步.
你會看到一個基督徒的進程.
神使到都是大家聽的福音.
我們所傳的福音.
不會多一點少一點的.
大家都是相信.
主耶穌就是人類的救主.
那撒勒人耶穌.
就是舊約所應許的救世主彌賽亞.
所以外邦的信徒.
貼著羅尼加的弟兄姊妹.
這樣信了之後.
就離棄偶像.
歸向真神.
這是一張前書的一張糾結.
要服侍這位又真又活的神.
就是這個意思.
所以福音呼召你們來到這個地步.
這個地步就是.
我們在一個屬靈的地位上.

$^{241}$我們有一個地位.
就是追求榮光.
追求榮光就是說.
一些表面看不到的不是物質.
是一個心靈.
有些人要做公職.
他追求的是一個服侍.
帶來的一個被人肯定.
或者很多做這些公職帶來的.
很多人對他的稱讚.
這些是非物質的.
不可以用物質去衡量的.
但是那份滿足感.
就會由然而生.
所以真心做公職的人.
做公益事業的人.
你會看到他的面容是很平和的.
每個人都是這樣的.
但是如果你看到那些人又做公益.
但是面容又不好.
這些可能是沽名釣譽.
就是為了拿位.
拍張照.
當然我們不要這樣評論.
因為每個人都不同.
但是無論如何.
我們要記得.
基督徒是要追求榮光.
我們每一個人的生命.
要每天追求的是主基督的榮光.
我們要在我們的面容.
反映耶穌基督的榮耀.
十字架的救恩其實是很榮耀的.
我們就要去學習這個功課.
十五節.
接著又有教導.
弟兄們你們要站立得穩.
怎樣站立呢.
當然就是神的道.
去改變我們的方方面面.

$^{281}$改變我們一些.
我們在生活.
我們在習慣.
我們在一些做事的態度上.
從前偷過錢的就不要再偷了.
小的偷錢的就不要偷了.
每個人都有不同的習慣.
凡是所領受的真道.
不拘是我們口傳.
是信上寫.
都要堅守.
第二章的第一節.
就說無論有言語.
有冒我名的書信.
有什麼的異象.
不要隨便信.
但現在保羅就說.
凡是這些口傳的.
信上寫的.
當然就是使徒親筆的書信.
或者是其他的使徒流通在教會裡面.
即是說我們的心靈.
和我們的耳朵.
要教台.
要調教真台.
真道台.
第二章的第一節.
就是那些教錯了台的人.
弟兄們.
論到我們主耶穌基督的降臨.
和我們道他那裡聚集.
我勸你們.
無論有靈有言語.
有冒我名的書信.
就說主的日子現在來到.
不可以輕易動心.
所以原來那些人.
真的不是想聽道.
是想聽下什麼叫末世論.
想聽下什麼時候是末世.

$^{321}$什麼時候那些人.
是怎樣厲害到一個地步.
怎樣那些人被提.
他們很想去尋求.
那個時間.
他自己就一路去研究.
保羅就說.
千萬不要沉迷於.
這些這樣的假設.
和這些的.
越是神秘的東西.
我們就越感興趣.
千萬不要這樣去追求.
就說不要輕易動心.
也不可以驚慌.
和第十五節.
這裡所說的.
剛剛相反.
就是甚至是一個很平和的基督徒.
是一個已經有了神生命的改變.
有聖靈在心裡面.
還有.
去到一個所謂的地步的意思.
就是說.
自己現在.
已經有了神的改變.
意思就是說.
對自己的存在的認識.
就是說.
我是神的兒女.
我是基督徒.
我要追求的.
當然我每天都要生活.
但更重要的是追求.
得到主耶穌基督的榮光.
就是說在他的眼中.
我是他所喜歡的.
我所做的事情.
他是喜歡的.
這個才是最重要的.

$^{361}$所以.
有了這個決心的人.
很簡單.
神的道就會很敏銳.
想聽.
想學習.
就算我還沒走得到.
我也很想去有一個學務.
去成為這樣的基督徒.
所以這裡告訴我們.
所以站立了之後.
所領受的.
不論是我們口傳.
信上面寫.
我們要去堅守.
就是說.
不論是哪個傳道者.
不論是什麼的傳道者.
如果他說的是神的道.
我們就要學習和堅守.
因為重點是.
在你的心靈裡面.
有一個很想學務.
對主耶穌的認識的心.
有一個例子就是.
譬如那個郵差.
有些是你很熟悉的郵差.
他送一封信給你.
你很自然就走開去看.
或者有一個郵差.
剛剛代工的.
你不認識他的.
他剛好又給你一封信.
而你很看見他的樣子就覺得很生氣.
不喜歡他這個郵差.
衣衫藍柳.
又負氣.
但是他給你的那封信.
是你的親人給你的.
你不會因為那個郵差.

$^{401}$那麼噁心那麼難看.
你就不看那封信.
有時候我們回教會聽到.
這個郵差很好.
神給他這個信息很好.
這個郵差就不行了.
根本還沒講.
我已經不喜歡聽他講.
反正今天又很睏.
今天睡覺.
你就錯過了.
那個信息.
神透過這個郵差給你的信息.
所以任何一個傳道人.
站上港台.
其實嚴格來說他應該是一個郵差.
他在那一刻.
是將神給他的信息.
傳遞給回眾.
所以.
這裡告訴我們.
保羅說不要理會什麼人.
無論是口傳的.
信上寫的.
或者不同的.
牧者給你的教導.
如果那個是真理.
是聖經的教導.
那就要堅守.
要堅定的遵守.
接著就一個祝願.
就是.
但願我們主耶穌基督.
和這一位愛我們.
開恩將永遠的安慰.
並美好的盼望.
賜給我們的父神.
形容父神.
愛我們.
父神怎樣愛我們呢.

$^{441}$將他的愛子.
來到他伯伯的犧牲.
至於死地.
剛剛.
報紙報導.
好像是內地.
好像是廣西.
有一個媽媽牽著兒子過馬路.
然後.
兒子很小的.
就.
發脾氣.
不肯過馬路.
在馬路上.
媽媽激氣到一個地步.
就扔下兒子.
兒子就不理會.
在地上躺著睡.
誰知有輛車.
轉彎.
車是高身的.
看不到他.
前輪沒有事.
後輪就滾著頭.
媽媽就衝過來.
就馬上.
拉著兒子就過了.
那些網民就罵到.
為何會這麼無知.
等等.
所以我們要記得.
什麼是愛.
父母的愛.
媽媽發現危險.
那一刻.
忍不住發脾氣.
以為沒事.
就忘記了.
小朋友不懂保護自己.
應該拉他一邊.

$^{481}$打他屁股.
但他母愛.
一發現.
就衝出來拉兒子.
到馬路上.
你就可以想像.
這是一個不合格的母愛.
尚且這麼強烈.
這麼強烈的母愛.
何況永恆的主.
保羅好像說得很簡單.
但其實我們要.
在人生裡經驗.
父神怎樣愛我們.
父神的愛.
是怎樣的呢.
我們看不到這位造物主.
他怎樣愛我們.
那個愛有多大的程度呢.
如果這個世界有高智慧的外星人.
為何他不愛外星人.
而愛我們人類呢.
而且人類有這麼多問題.
為何他都要愛呢.
愛到犧牲他的愛智呢.
這是其中一種.
你想想創世記的時候.
為著人類可以出現在這個地球.
鋪排那些.
好像那張嬰兒床.
生個兒子.
你就知道了.
第一胎就特別緊張.
又吃奶瓶.
又去預備.
尿片等等.
買好多罐奶粉.
怕不夠.
各方面要溫度計量.
水溫夠不夠熱.

$^{521}$很多東西又要消毒.
又要什麼.
過了三個月你都受不了.
什麼都隨便.
沖一沖就算了.
還要煮嗎 沒可能.
你可以想像得到.
天父為人類.
預備的.
新天新地.
和伊甸園.
未墮落之前.
那種美麗.
那種可宜居的.
舒適.
每一個果子.
可越人眼目.
每一個樹.
每一種動物.
每一個海洋生物.
全部可以讓人.
享受.
人不是萬物.
其中一種.
是照著神形象做.
唯有人.
才照著神形象做.
說了三次解創世記.
你可以想像得到.
這是父神對人的愛.
但願.
這位.
我們的主耶穌.
和這位愛我們.
開恩他的恩典.
將永遠的安慰.
美好的盼望.
永遠的安慰.
我們不是一時間得到得救.
是我們在神的愛裡面.

$^{561}$我們得到永恆.
永遠的安慰.
所以安息主懷的人.
真是好得無比.
人世間始終有.
我們的 只是睡不著.
頭痛.
你已經覺得很辛苦.
何況.
如果我們面對很多.
例如工作 煩惱.
人事 假設香港有.
戰爭更加麻煩.
聖靈塗炭.
顛沛流離.
的日子.
但真是永恆的安慰.
是我們在神裡面.
得到的 所以保羅在這裡.
告訴我們 永恆的安慰.
和美好的盼望.
美好這個字就是美麗.
那個字 希臘文.
good那個字的.
形容.
那個盼望是很美的.
很美好的.
很美麗的那個盼望.
所以我們基督徒.
有那個永恆盼望.
是因為我們有一個永遠的安慰.
所以是孖公仔來的.
就是因為父神的愛.
令到我們得到.
永恆的安慰.
當然是因為主耶穌.
祂的犧牲.
賜給我們 去安慰我們的心.
並且在一切的.
善行善言上面.

$^{601}$兼顧你們.
這個善行善言.
似乎是我們生命的一個改變.
善行善言.
心靈裡面的善.
那個goodness.
那個.
一個.
內心裡面的一個改變.
聖經很.
很重的.
說 特別.
在《大約書4章》.
說有些人.
一定要改變.
不要再不斷由你的說話.
產生很多的.
問題.
不信主的人.
內心有些可以很.
苦毒的.
他自己苦 接著就播毒.
就會令到別人.
受很深的傷害.
所以言語的傷害是很難.
改變的 甚至乎.
你不信的人都知道 你說出的.
事情是回不了頭的.
所以我們要很小心的說.
話的 所以這裡告訴我們.
善行善言.
去兼顧我們.
所以你想想基督徒.
是多麼寶貴呢.
有永恆安慰.
有美好盼望.
就是說我們將來的歸宿.
只是永恆的光亮.
所以但是今生.
這麼有限的人生.

$^{641}$我們就要潔淨我們的心.
我們就要在這份愛裡.
我們就要帶出.
一份善行.
一份善言.
所以保羅真的說到.
原來行事為人.
是裡外要合一的.
你裡面是有惡毒.
你不可能.
你說的話是乾淨.
是美麗.
是造就人的.
你是假到自己也呃了.
這樣才會這樣.
不然一定是.
你先是裡面被神的愛.
改變 感動.
你就有這種力量.
在你的生活裡面.
能夠有善行善言.
使你能夠.
得到神的喜悅.
如果善行善言.
是神的.
那個屬性.
帶給我們的改變.
那神給我們.
什麼善行 什麼善言.
來兼顧我們.
我們可以想像得到.
整本聖經.
詩篇119篇告訴我們.
神的道好像蜜的甘甜.
好像比上.
堂的.
常過的.
食物的甘甜.
你可以想像得到.
神的道對我們的改變.

$^{681}$的偉大.
這就是第二章.
兩類人.
原因就是.
保羅就透過.
原來在當時的.
教會.
因為講主在來.
講得很.
有那個.
危機感.
所以貼實羅尼加的人.
一方面原先他們是.
懼怕 不知道.
自己死去的親人去了哪裡.
保羅就安慰他們.
但是會帶來.
一些可能是異端.
可能是一些假的使徒.
或者一些人.
可能一知半解.
就將這件事無限的上綱.
就告訴我們何時回來.
有什麼.
的預言.
保羅說不要相信.
無論那些人.
冒我的名的書信.
都不可以輕易的動心.
因為有一個.
離道反教的日子.
有這個大罪人.
會出現.
也有撒旦的運動.
會令迷惑那些.
本來的心.
不是在真理裡面.
或者沒有愛慕.
真理的心的人.
即是沒有得救的人.

$^{721}$就會被迷惑.
所以真正的試金石就是13.
到17節.
我們是主所愛的.
我們是蒙神的愛的.
我們是在聖靈的感動下.
已經成為聖潔.
能以得救的.
這個試金石.
其實每天起身.
你都要確認.
自己是神的兒女.
你的行事為人.
是反映著.
基督的榮光.
這樣的試金石.
或者說清教徒就是這樣.
怎樣知道自己得到永生呢.
一次相信.
一直相信就必定得救.
就不是一次相信.
就得救.
即不是一次得救.
永遠得救.
是一次得救一直相信就必定得救.
所以中間加上的是.
一定的相信.
是一直的相信.
這個是清教徒.
他們的信念.
也都上個禮拜.
前個禮拜講過.
清教徒的定義.
就是如果我們收窄那個範圍.
宗教改革之後.
他們是沒有了.
最根本的天主教的信仰.
是離開天主教.
教會的以外.
是不會有救恩.

$^{761}$你乖你就回來天主教.
你不要跟馬丁路德走.
你不要跟這些改革家走.
你回來有救恩.
因為天主教會任避你.
但是離開教會就沒有救恩了.
所以無論是嘉爾文.
馬丁路德.
慈雲尼.
或者其他的後來的.
清教徒的領袖.
一定要安慰他們.
我們是相信在主耶穌裡面就有救恩.
不是在教會裡面.
不是在教廷裡面.
在主耶穌裡面唯獨基督.
唯獨恩典.
唯獨信心.
唯獨聖經.
唯獨榮耀神.
這五個唯獨就是當時清教徒的根基.
所以就令他們有一個信心.
所以後來就沉澱成為一次得救.
一直相信就必定得救.
所以要一直持守這個信仰.
所以保羅就告訴我們.
我們要追求的是主基督的榮光.
所以第二章就有一個對比.
但第三章似乎是一個後補.
就是教會裡面有些懶惰人.
就幾玲瓏浮突地罵.
罵得都幾嚴重.
甚至乎如果不聽勸.
就要當他們是外邦人.
就要跟他們疏遠.
所以我們要留心.
第三章我們看看這一段經文.
弟兄們,我還有話說.
請你們為我們禱告.
好叫主的道理.

$^{801}$快快行開,得著榮耀.
正如在你們中間一樣.
神的道可以在人的心裡面產生改變.
這個時間應該就是第二次的傳道.
保羅應該是在雅典.
我不知道後書的時候.
他是不是還在雅典.
還是在哥林多.
甚至乎在哥林多住了一年半的時間.
在寫這個後書.
時間就不是很清楚.
前書我們就知道.
可能相差無幾.
可能都是三幾個月.
接著第二節.
也叫我們脫離無理之惡人的手.
因為人不都是有信心.
但主是信實的.
要堅固你們.
保護你們脫離那惡者.
我們靠主深信.
你們現在是遵行我們所吩咐的.
後來也必要遵行.
願主引導你們的心.
叫你們愛神.
並學習基督的忍耐.
在患難裡面.
有些事都不是你想改變就改變到的.
因為有惡者.
有些是特意衝擊基督教的人.
他們是敵視基督教的.
是反對福音的.
例如猶太教就是反對.
因為他們要維護自己猶太教的根本.
他們是拒絕耶穌是他們的彌賽亞.
這些是不能改變的.
所以只能夠零零星星的猶太人信主.
所以就會有這種逼迫.
或者是信主之後.
帶來的一種籠罩著的受苦的氣氛.

$^{841}$所以每一間教會都要面對.
每一個時代的教會都要面對.
這個保羅所勸勉他們.
所以要愛神.
和要學習基督的忍耐.
學習基督的忍耐.
基督是怎樣忍耐罪人.
怎樣忍耐十二使徒.
怎樣忍耐在自己的對宮裡面.
有一個叫猶太.
要忍耐.
第一天就把他抓出來.
當然沒有問題.
你是神.
你知道他的心.
其他門徒都不知道.
連他出去買東西.
人們都以為他是周濟窮人.
在最後晚餐.
你想想他的隱藏是多麼的出色.
但耶穌從.
《約翰福音》大家讀過.
耶穌從起初就知道誰要賣他.
《約翰福音》六章裡說的.
作者好像不斷提醒讀者.
我現在要告訴你.
往往聖經的書信都是馬後炮.
發生了一些事.
作者會記錄那些事.
不是一路寫著.
還沒發生.
現在不是寫劇本.
還沒拍的.
現在要馬上寫.
不是的.
我們看電視劇就是.
是還沒看.
但要寫劇本.
但聖經就是事後那些事發生了.
作者會寫出來.

$^{881}$所以作者是很清楚.
跟著耶穌的人生的時候.
是耶穌說過的.
記載下來.
很想讀者知道.
其實耶穌早就知道了.
知道而仍然讓他存在.
一方面是神的旨意.
要藉著他被出賣.
他的門徒.
他所愛的門徒.
用腳踢他.
69篇的詩篇.
不是要親托猶大的無奈.
好像命中被神操弄.
其實不是.
完全是猶大自己的自由意志.
但是這裡反映了主基督的忍耐.
對罪人的忍耐.
這些真的要學.
因為如果我們真的明知道有些事要發生.
而我們知道神都要讓它發生.
我們都要學習忍耐.
因為這樣才能看出救恩的偉大.
所以有時神容許我們明知道有些事要繼續去做.
他都容讓.
所以不要說神為什麼不提醒我.
其實不是不提醒.
是你好像一匹不受約束的雷子.
來做這件事.
所以這一段才是安慰.
對帖撒羅尼加的信徒.
最後的叮嚀.
保羅擔心的是.
他們必須要去時刻的警醒.
因為惡者會用不同的方式.
進入教會或帖撒羅尼加的信徒的心靈.
所以一方面就是無理之惡人的手.
是實實在在有形的手.
即是說實質的.

$^{921}$在地上的.
就伴隨著他們的全道.
他們的聚會的人生.
是會有一些逼迫的.
是會有一些反對的勢力的.
這個是沒辦法面對.
因為人不都是有信心.
但主是信實的.
就要兼顧你們.
保護你們.
所以在這個困難下.
我們就不是以牙還牙.
你做初一我們就做十五.
而是我們要用基督的愛.
去包容.
去寬恕.
去接納.
重要的就是去忍耐.
有些事是改變不了的.
我們就要學習去接受那個處境.
但是記住.
人不都有信心.
但主基督是信實的.
即是說祂會保守我們.
祂如果容讓我們面對這些難處.
祂一定會保守.
我們只要問心無愧的.
根據真道而行.
神就會有對祂子民的保守.
一定會兼顧我們.
而且最終是保護我們.
脫離惡者.
或者脫離罪惡.
這裡我們要記得.
又是這樣教導.
我們靠主心信.
你們現在是遵行我們所吩咐.
後來都必定會遵行.
即是說保羅對帖撒羅尼加教會的弟兄姊妹.
是充滿了一份信任.

$^{961}$一份對他們的信心.
他們原來真是生命改變.
很單純地被主的救恩改變.
而這個救恩是由保羅宣講出來.
他們就信.
信得很好.
他們來到這個階段.
就算保羅有很多限制.
不可以在他們中間繼續的供養.
教導他們的心靈.
但是他們已經有了那個生命.
他們就能夠一路的.
透過不斷的認識真道.
他們就會來到成長.
所以保羅就仍然的勉勵他們.
你們現在遵行.
後來必要遵行.
即是說有沒有牧者在你身邊.
只要你有這個渴望真道的心.
你們就會找到食的了.
你們就會找到哪些是神的道.
來到去成長的了.
所以我們不要偏食.
不要偏某一些的講源.
我們是要足類旁通.
我們要透過不同的牧養.
來到去豐富我們自己的生命.
重點是我們的屬靈生命.
要有這個饑渴.
要有這個對神的愛的渴慕.
自然我們就會好像一個永遠的海綿.
不是一些已經爛泥.
你怎麼曬它都不會柔軟的了.
那些泥已經混雜了很多東西.
但是一個海綿你怎麼整它呢?.
它都是繼續這麼柔軟的.
我們的心靈要時時保守自己.
是一個很柔軟的心.
能夠不斷可以讓主的道.
在我們裡面改變.

$^{1001}$所以第五節.
願主引導你們的心.
叫你們愛神.
和學習基督的忍耐.
全部有關連的.
我們那個心讓神去引導.
我們很自然就一定會有一個愛這位造物主.
愛神的心.
臨前第八章第三節.
講一講一些事情.
如果有人不愛神.
這個人就是可助可救.
就是這樣講的.
突然間講一些很普通的東西.
他突然間就說好像很有感動.
所以我們信主.
應該和這位造物主不會越來越隔膜.
而是越來越有一種很想愛他.
很想事奉他.
很想被他不斷引導.
這樣的一個內心.
愛神帶來的就是越來越像基督.
像基督的就是主在地上面對困難的時間.
他有那個信靠.
他知道一切的公義是在神那裡.
所以他不會怕被人奪去公義.
《二次元書記》五十三章十一節.
他的公義從地上被奪去.
主耶穌這位受苦的君王.
這位被牽到宰殺之地的羔羊.
從地上被奪去.
地上沒有你的公理.
沒有你的公義.
你可以想像得到多慘.
含冤.
完全沒有辦法得到解脫.
神幫你去伸冤.
包青天都不幫你伸冤.
這樣的一個耶穌.
所以耶穌他仍然有那個忍耐.

$^{1041}$他仍然知道他縱然被父神離棄.
被人釘十字架九個小時.
這樣來斷氣.
但是他知道他仍然可以忍耐這些痛苦.
因為他清楚是神的救恩.
藉著他的犧牲他的死.
才能夠成就在罪人的身上.
所以學習基督的忍耐.
都是我們一生的功課.
神引導我們愛神.
學基督的忍耐.
這些真是.
如果教會的領袖全部是這樣的心靈.
應該是這樣的.
因為你應該是身先士卒.
讓人看到你的愛神.
讓人看到你的人生是有主的引導.
讓人看到你有基督的忍耐.
這樣才是領袖的級別.
才是這樣的特質.
任何一個走去服事主的人.
應該就是要這樣.
所以這是一個很重要的基督徒.
或者教會裡面最寶貴的資產.
當然很多教產都是重要的.
沒有了這個建築物我們都不能聚會等等.
但是在苦難的日子.
例如尼祿攻擊教會的時候.
他們哪裡有聚會呢.
在羅馬在地底下.
因為很多洞穴就在裡面聚會.
哪裡有教會呢.
沒有硬件 但沒有硬件不要緊.
有一班被主引導的基督徒.
愛神的.
他們願意效法基督的忍耐.
這樣就是教會了.
這樣就能夠吸引很多.
當時在碧白時期.
很清楚的.

$^{1081}$你跟隨耶穌就要受苦了.
但都有這麼多人跟隨.
就證明了他們看到這些基督徒.
他們的心裡面有神的愛.
因為他們愛神.
同時時有主耶穌的引導.
在他們的心靈裡面.
而且他們願意學習忍耐.
很多故事流傳下來.
有些羅馬的兵丁.
在雪地捉了40個基督徒.
就要行刑.
其中一個就是先脫光他們的衣服.
冬天冰天雪地.
不准他們穿衣服.
然後就活活的凍死.
接著那些兵丁就聽到天使唱歌.
就說到40個殉道者.
有40個官免.
接著一個一個真的凍死了.
後來去到第39個.
聽不到天使在唱.
不知道什麼原因.
突然間那些兵丁.
有一個人已經凍得不停的在顫抖.
就跟那些兵丁說.
我不信耶穌了.
你把衣服給我穿.
我否認耶穌了.
那個兵丁就說.
我馬上給你穿這件衣服.
你穿.
接著那個兵丁就變成赤身露體.
在冰天雪地.
突然間有天使唱歌.
40個殉道者.
40個官免.
那個兵丁就拿了官免.
因為馬上信耶穌.
因為他真的看到異象.

$^{1121}$是當時流傳下來.
初期教會.
大碧白時期.
主後300年的碧白期.
君士旦丁未信耶穌之前.
東羅馬的國家.
碧白基督徒的十大碧白時期.
就有這些故事流傳.
那個兵丁就說.
你不信?我信了.
馬上給他一件衣服.
接著他聽到天使唱歌.
他就得到殉道之後的官免.
所以這裡告訴我們.
如果這些是真的異象.
是神真的在一個很特殊的碧白時期.
而有的異象.
所以這裡告訴我們.
有一群人.
至少現在貼在內.
現在教會的弟兄姊妹.
他們真的被神引導.
他們的心很單純.
因為福音令他們很單純的心靈.
去認識主.
接著不斷被神的愛充滿.
有神的愛在他們的內心.
他們一定會回應愛神.
愛神就是說.
我在地上為神作見證.
我在地上.
好的不好的境遇.
我都不會否認神.
會學習的.
很肯定就是學習這位神的兒子基督.
在地上的人生.
他怎樣默默的為那些罪人去投歡笑.
在十字架上.
七言的第一句.
父啊赦免他們.

$^{1161}$他們就是兩個囚犯和下面的群眾.
因為他們所作的.
包括那些兵丁.
釘他們.
鞭打他們.
笑他們.
因為他們所作的.
他們不曉得.
在十字架上被釘住.
這樣去說的第一句說話.
所以基督的忍耐.
有多深的深度呢.
你想想我們去探測一下.
地球下面有多深呢.
我不記得太平洋哪一個海溝.
全世界最深的地方.
那些蛟龍號去到.
因為越深水壓就越大.
人類的極限都去不到.
所以你可以想像到.
基督的忍耐在地上.
因為他要成就救恩.
是有多深呢.
在這個世界有多少基督徒.
有基督的忍耐呢.
要學一生只要學這件事都夠了.
通常就是.
神說你想學就對了.
就放一些不可愛的人在你面前.
激你的.
是真的要心要肺的.
你要學基督的忍耐.
就不是愚昧地打我左面右面.
而是你上心的人.
令到你很痛的.
你要學忍耐.
這些我想真是每一個基督徒的功課.
但本來這樣說的時候.
可想而知貼薩諾尼加的教會.
是多麼願意行真道呢.

$^{1201}$本來是多麼對他們的稱讚呢.
但是第6節到第15節.
就開始罵一些懶惰人.
我先喝杯茶.
我們看看第6節.
弟兄們.
我們奉主耶穌基督的名.
奉主的名.
很嚴肅的.
在主同在下.
去說什麼呢.
去吩咐你們.
凡是有弟兄不按規矩而行.
不遵守從我們所受的教訓.
就當遠離他.
你們自己原知道.
應當怎樣效法我們.
因為我們在你們中間.
未嘗不按規矩而行.
似乎在說一些實質的問題.
就是關乎到工作.
第8節.
也未嘗白吃人的飯.
倒是辛苦勞碌.
晝夜作工.
免得叫你們一個人受累.
這並不是因我們沒有權柄.
乃是要給你們作榜樣.
這裡有很大的原則開放出來.
牧者是有權柄.
這裡沒有詳細說.
哥林多前書第8章.
就詳細說供應的問題.
牧者.
或者第8第9章.
領受薪津的問題.
所以這裡只點出那些.
現在貼在教會的一些人.
肯定是教會裡面的人.
不是那些外面.

$^{1241}$他們現在出現了一個行為的落差.
就是不做事.
而且還要合理化自己不做事.
但我們實質上.
又不太知道為何會去到這麼盡.
甚至不做事.
他沒有了供應的時間.
不可以用愛去當飯吃.
那怎樣去維持呢.
可能會濫用了帖撒羅尼迦城.
弟兄姊妹的愛心.
會用一些他們認為的理由.
所以如果看回整本書.
有可能就是主在來.
不要做事了.
快點預備主見主.
可能就是這樣也說不定.
始終都是猜測.
因為保羅沒有詳細說.
為何會導致有些人不肯做事.
就要叫人受累.
第9節.
「這並不是因我們沒有權柄.
乃是要給你們作榜樣.
叫你們效法我們」.
保羅也是一個很有智慧.
和很落地的牧者.
他一開始就說我們是使徒.
我們的身份是超賢.
我們要做傳道.
我們要將神的奧秘講述出來.
所以教會還沒有成立.
你們現在就要出這樣的新津給我們.
保羅就不是這樣.
他反而要做一個榜樣給他們看.
他五天的時間.
他會織帳棚.
他會有其他的工作.
或者有其他成熟的教會.
供應保羅.

$^{1281}$比如肥立比教會.
馬其頓其他教會.
供應保羅的需要.
但一個初教會.
剛剛建立的.
他就完全不會立即啟動.
他所謂的權柄的福利.
那你就可以想像到.
保羅是一個很有智慧.
很願意為別人著想的一個牧者.
接著留意一下這裡.
第十節.
我們在你們那裡的時候.
曾經吩咐你們說.
如果有人不肯作供.
就不可吃飯.
因我們聽說.
在你們中間有人.
不按規矩而行.
什麼供都不作.
辦到專管閒事.
不作供那個人.
就不會只是坐在那裡.
看電視.
沒有電視.
也不會有互聯網.
也不會自己坐在家裡.
他不作供的時間.
就會專管閒事.
那管完之後.
聽完聽完的閒事.
就會變成很多的是非.
所以其實是一個惡性的.
螺旋的循環.
所以不單止他自己不按規矩.
而且會繼續破壞規矩.
破壞教會的愛心和合一.
所以那個問題.
不是吃不吃飯那麼簡單.
就算他不在教會吃飯.

$^{1321}$不問其他人的要錢.
他只要不作供.
不在他人生中有神給他的差事.
他不來看.
所以清教徒看.
那個工作.
不是工作.
不是一個做了賺錢就算了.
他們是一個叫做照命.
是一個叫做vocation.
所以vocation 勞.
就是職業介紹所.
職業先修.
訓練你去工作.
其實嚴格來說.
vocation 是在說你終身的職事.
是一個終身的事業.
叫做vocation.
所以vocation 這個字眼.
是清教徒那裡演變而來的.
所以今天的職業先修.
或者叫做IVY.
vocational center.
訓練不是叫你只找份工作.
是訓練你一個技能.
終身的.
就叫vocational.
所以這裡告訴我們.
這裡有些人不作供.
就會專管閒事.
所以真的是孖公仔.
我們要小心.
我們不作供.
我們就會管閒事.
其實很需要求主憐憫的.
當然.
很拼命做事的人.
一樣會有這樣的情況.
但這些肯定會更加嚴重.
第十二節.

$^{1361}$如果我們揭一揭.
保羅的教目書信.
有提的.
我們先讀完.
我們就揭一揭第五章的提摩太前書.
就有些年長的婦女.
第十二節.
我們靠主耶穌基督.
吩咐勸戒這樣的人.
我們要勤奮地去做.
勸戒這樣的人.
記住.
奉主耶穌的名.
第十二節.
靠主耶穌基督吩咐.
是一個祝福.
是一個charge.
勸戒這樣的人.
要安靜地作供.
吃自己的飯.
他能夠有重地而來的修成.
或者如果有一些勞動的工作.
當時不是那麼多的工種.
可能幫一些貴族打工.
有錢.
就有生活.
或者在他家裡做管家.
做一些服役.
又會在他家裡.
好像上一代的傭工.
那些阿媽.
就是.
沒有菲傭的之前.
其實現在菲傭都是.
都是在你家裡.
做一些工作.
親手作供.
吃自己的飯的意思.
第十三節.
弟兄們.

$^{1401}$你們行善不可喪志.
如果有人不聽從我們這些信上的說話.
就要記下他.
不和他交往.
叫他自覺羞愧.
但不要以他為仇人.
要勸他如弟兄.
這裡又說了教會的紀律.
慰爵不是要紀律.
慰爵是透過一些.
讓他知識自己的問題.
的一個方法.
使他最終.
認識自己.
使他最終.
回去和神.
和人的和好.
所以不是為紀律而紀律.
或者根本沒有紀律.
做什麼都可以.
那教會就一定倒了.
太嚴重的.
嚴厲的紀律.
紀律也是會令人.
受很深的傷害.
所以你們留意保羅的教導.
不要以他為仇人.
我們如果拿著雞毛就當作令箭.
你最不喜歡那個人.
現在你就可以動他.
你就真的當他是仇人.
來處理.
那是不對的.
所以保羅就告訴我們要留意.
我想和大家聊聊.
提摩亞太前書的五章.
提前的五章.
保羅就勸提摩亞太.
因為他是一個年輕的牧者.
就是剛剛出道.

$^{1441}$沒多久.
肯定沒有保羅的年資.
歲數.
不可以嚴責老年人.
因為他真的比你大很多.
是你的長輩.
他當然在那個真道上.
沒有你的那個恩賜.
和那個心的教導.
裝備.
但是他始終是一個老年人.
所以不可以嚴責老年人.
只要勸他如同父親.
要勸.
就是說他也有錯的.
不是老年人就不會錯.
他有錯但是你要勸導他.
用安慰的話勸導他.
好像勸父親那樣.
就是說這種尊敬的心態.
去作出那個模樣.
勸少年人如同弟兄.
在教會上.
勸他如同父親.
如同弟兄.
在主裡面的弟兄.
勸老年婦女如同母親.
勸少年婦女.
如同姐妹.
總要清清潔潔的.
要尊敬那真為寡婦.
若寡婦有兒女.
或有孫子孫女.
便叫他們先在自己家中.
學著行孝.
就是說讓那些子孫.
去幫助.
那些子孫.
幫助.
供應自己.

$^{1481}$那就是說意味著.
那些寡婦真寡婦.
是有.
教會的供應.
因為他真是.
舉目無親.
真是.
完全沒有人.
真是孤兒寡婦.
那教會當時.
因為可能社會沒有.
這些機制所以教會.
就不可以.
袖手旁觀.
要教會去行使.
牧養的責任.
對這些信主的寡婦.
來去供應.
但如果有孫子孫女.
有兒女.
保羅就說.
由自己的家庭盡那個.
孝道盡那個責任.
為什麼呢.
報答親恩.
因為這在神面前.
是可閱立的.
對父母的孝順.
真是神閱立.
當孝敬父母.
十誡裡面的第四誡.
不是第五誡.
對不起.
第五誡.
十誡如果兩塊法板.
頭那四誡是對神的.
但第五誡.
是第一條誡命對人的.
就放在第一塊石板裡面.
所以保羅說.

$^{1521}$當孝敬父母.
使得你得恩典.
在世長壽.
這是第一條.
大應許的誡命.
十誡裡面的第五條.
在約板.
就是神給摩西的法板.
如果左手邊.
是第一塊.
右邊第六到第十條.
才是對人的人倫關係.
所以保羅就說.
看第五條的.
當孝敬父母.
是第一條.
大應許的誡命.
就是對人.
你做了你會得到神的喜悅.
閱立的.
所以聖經都講孝道.
我很久沒有和媽媽喝茶了.
都是要悔改.
因為太多兄弟姐妹.
所以不用輪到我.
但是.
沒有了.
這個上心的孝道.
你可以想像在神面前.
是可以閱立的.
第五節.
留意.
第五節.
那獨居無靠.
真為寡婦養賴神.
晝夜不住的祈求禱告.
所以你會看到.
一直都在講.
真寡婦.
無依無靠.

$^{1561}$但是他們的心是依靠神的.
所以我們要敬重.
在教會裡面.
不是抬舉.
也不是要褒揚.
而是我們要留意.
我們要關心.
我們要尊重.
這些一心養賴神的.
敬虔的婦女.
所以真寡婦.
在保羅的教導裡.
但是有些假扮真寡婦.
第六節.
但是那些好燕落的寡婦.
真是寡婦.
不過好燕落.
正活著的時候.
都是死的.
原來他們的人生.
他們的專注.
沒有了老公.
沒有了家庭.
不用每天煮飯.
又被老公罵.
又有空.
就好燕落了.
所以正活著的時候.
他們的寡婦.
不會帶給他們傷痛.
反而帶給他們.
很多時間去燕落.
所以保羅就叫提摩太.
要很小心.
可能這些都是老寡婦.
沒有說到年輕.
老的都可以好燕落.
有時死.
所以保羅其實都很重.
他說.

$^{1601}$這些事你要祝福他們.
叫他們無可指責.
人如果不看顧親屬.
就是背了真道.
給不信的人都不好.
不看顧自己家裡的人.
都是更加如此.
接著立即就說.
真是有個制度.
寡婦既在冊子上.
必須年紀到60歲.
從來.
可能那時60歲很大年紀.
都未死.
又是寡婦.
就要公認.
尊敬他們.
尊敬裡面包括她的生活.
都要給她.
有尊嚴地過.
那個生活.
從來只作一個丈夫的妻子.
又有行善的明星.
即是說原來.
都有些客觀條件.
因為如果她真是.
沒有好明星.
大家都覺得她不好.
那你教會再這樣.
好像一個很僵化的機制.
其實就不是那麼好.
有些人就算奉獻.
給教會去行使.
這個公認.
都不鋸.
因為你.
安慰這些全部.
又沒有好明星.
你又給她那邊出去燕樂.
那邊又出去講是非.

$^{1641}$其實帶來教會傷害.
受虧損.
所以在行善.
有明星.
養育兒女.
接待院人 洗聖徒的腳.
救濟遭難的人.
竭力行各樣的善事.
這個明星都很多.
即是說有目共睹.
淋漓盡致.
其實一個女性本身就是這樣.
是很好客.
很低調地.
很殷勤.
很有心去服侍人.
等等.
其實已經差不多.
你看到這些對號入座.
就已經可以.
接著第十一節.
保羅教提摩太.
都要留意.
至於年輕的寡婦.
未夠六十歲.
就可以辭退她.
不要幫她登記.
因為她有手有腳.
有氣有力.
就給她做事.
就可以辭她.
情慾發動.
違背基督的時候.
就想要嫁人.
這裡登記了.
她三十歲.
那裡還是深夜夜.
又見到如意郎君.
接著就不用你登記了.
不用你救濟了.

$^{1681}$帶來一個好像搞出一個笑話.
保羅就說.
要小心.
她從其他地方.
去受供應.
繼續學習在寡居裡面.
寡居不等於不可以再結婚.
所以就不要隨便登記.
所以沒有說不准結婚.
而是她這樣的時候.
因為.
為什麼呢.
如果你有一個供應的制度.
我多麼說自己是寡婦.
但實際上心底.
還是情慾發動.
想要嫁人.
就變成一個很大的.
大家都好像.
被一個無謂的.
登記制度.
約束住.
那不如保羅就說.
辭退她.
她想要嫁人就嫁.
第十二節.
就是因為這裡.
他們被定罪.
是因為廢棄了.
當初所許的願.
所以這些.
為什麼在神的面前.
許願.
包括耶穌.
包括保羅都說不要隨便許願.
所以提摩泰.
你都不要立一個.
這樣的在教會裡面的.
規矩來影響.
那些年輕寡婦.

$^{1721}$隨隨便便.
我沒有下一任.
我現在一心.
這樣走這條路.
我和那些60歲的寡婦一樣.
你幫我登記吧.
我是很慘的.
這個願就害死她.
所謂害死她就是.
她在神面前.
好像違背了.
原來她很快.
就有那個心動動.
就要嫁人.
當初就不要這樣許願.
你這樣許願.
只是為了那個救濟就不好.
所以保羅就教提摩泰.
千萬要小心.
她廢棄了當初.
所許的願的罪.
保羅說.
他們被定罪.
當然眾人都會笑她.
這裡登記.
那裡就去嫁人.
教會就很尷尬.
這個是人對她的定罪.
但是許願.
我們就不要隨便.
真的不要隨便.
我們要在神面前.
很嚴肅看這件事.
神會追究.
第十三節我們要想說.
他們習慣了懶惰.
挨家行遊.
逐家逐戶去.
不但是懶惰.
又說長道短.

$^{1761}$好管閒事.
說些不當說的話.
所以.
為什麼保羅就說.
這個叫提摩泰要很小心.
不是他要殘忍.
不是教會那麼忍心.
明明是寡婦.
真寡婦來的.
但是那個真.
未能夠用時間去考驗.
六十歲.
就因為真的.
調查過.
真是無兒無女.
她也真是有好的明星.
這樣去照顧.
教會去照顧.
是很合情合理.
但是年青.
真是還有很多空間.
她自己的緣原來未是那麼實.
所以提摩泰作為.
學者也要勸告她.
當下這個.
許願我很欣賞.
但是.
人生其實很漫長.
很快就變.
你又慘.
教會又慘.
那就會不好.
所以第十四節.
所以保羅說.
我願意年輕的寡婦.
嫁人.
原來是可以的.
寡婦也不需要登記.
去拿救濟.
是可以嫁人.

$^{1801}$生養兒女.
留意一下治理家務.
不給敵人辱罵的把柄.
辱罵相信就是上文說的.
因為他們習慣懶惰.
因為有救濟.
就挨家行遊.
不單是懶惰.
懶惰是對自己.
也不是欺騙.
因為他已經登記了.
這個後遺症很厲害.
就是.
那個人就會很多的.
錯誤的想法.
層出不窮.
所以拿縱緩是不是真的會差一點呢.
社會福利署.
可能也要讀一下這些經文.
就是經常說.
縱緩會令人懶惰.
如果一個很穩定的機制.
是你什麼都不用做.
你就拿到這個收入的.
這個生活的.
可能真的會有.
沒有了那種動力.
他會有另一種動力.
就是講閒話的動力.
每次都是這樣說的.
聖經說一懶惰的人.
那個腦子就特別發達.
要講話.
真的不知道為什麼.
是一個很大的我們要小心的問題.
是保羅這樣講出來.
他作為一個牧者.
一個使徒.
明明教提摩太.
就是說要.

$^{1841}$觀察人要有這樣的.
淑寧的眼光.
和對人性有一種.
透徹的理解.
年輕寡婦.
可能真的有很多心事.
可能真的有很多故事.
可能真的有很多苦水.
可能真的很想和人吐.
吐得多就慢慢變成.
挨家這樣走油.
所以.
所以有時.
學下不同班組.
去釣釣魚也好.
在家裡閒懶.
也是一個問題.
又沒有子女去照顧.
因為說明寡婦.
又沒有其他的家人.
可能女性.
感性的情況會多一些.
所以保羅說.
要生養兒女.
要自理家務.
有把柄.
接著.
留意十五節.
因為已經有轉去.
隨從撒旦.
越講越重.
保羅.
我做寡婦.
我挨家行油.
我說長道短.
我好管閒事.
怎會隨從撒旦呢.
那個心.
慢慢就越來越.
偏向魔鬼.

$^{1881}$所以第十六節.
信主的婦女.
若家中有寡婦.
自己就當救濟她們.
不可累著教會.
好使到教會能夠救濟.
那真是.
無依靠的寡婦.
所以教會是真的要.
有這個.
收集弟兄姊妹的獻金.
來代表神.
去幫助.
真正有需要的人.
所以要有智慧地設立機制.
有智慧地勸退.
一些.
不合適的.
不是你的身份不合適.
是你現在這個心態光景.
不合適.
時間會塗鑿你的心.
你會慢慢就心野.
所以保羅說.
你再做寡婦.
可能你嫁人.
你就不要這麼快.
嫁人的寡婦就不需要了.
這裡.
只是說懶惰.
已經在這裡說關於寡婦.
關於救濟.
就已經有這些經文.
這裡已經是保羅在後期.
在二不所教會.
提摩泰在這個大教會裡面.
保羅在監獄裡面.
勸提摩泰.
所以貼在後書這一段.
就更加是一個.

$^{1921}$代表性.
就是如果.
主在來.
真是背後的問題.
而他們覺得很快.
主回來我們要被提.
有些人就走精面.
不做事了.
保羅就說要嚴肅地責備這些人.
是奉主的命.
靠主耶穌吩咐.
第六和第十二節.
這樣去作出警告.
人生.
是有照命的.
當然當時不是.
中世紀的中國改革.
但是在保羅的內心.
Vocation.
工作是神聖的.
每個人神一定給我們有工作的能力.
每個人.
都真是用工作去榮耀神.
其中就是.
神讓這個地球.
工作.
可以有修成.
可以有農作物.
接著有.
食物.
可以來分配.
可以衍生不同的.
商業.
一路一路的不同的.
工種一路一路的來發展.
去衍生出來.
所以每個人應該在.
自己的人生找到自己.
神給的崗位.
那個Vocation.

$^{1961}$所以記住Vocation是.
神眼中的.
那個照命.
工作我們要看為一個照命.
所以.
清教徒就已經是一個很明顯的.
一個example.
所謂富不過三代.
某程度是真的.
第一代的清教徒很殷勤.
覺得整個人生.
神給我們能夠有工作.
我們現在能夠坐五月花.
去到美國.
發現新大陸.
又獻一隻火雞.
去一個新的地方.
他們很殷勤的工作.
但是第二代第三代.
就慢慢以工作成為了神.
工作就是賺錢.
謀生享受的.
那個.
終極.
所有的Vocation.
都是為著榮耀神.
所以你會見到.
美國人很勤力工作.
很勤力去玩很勤力去健身.
很勤力吃那些保健的.
產品.
但是他們會很勤力的.
最後整個人生最重要的.
就是你禮拜日.
還會不會去敬拜神呢.
美國人就不會了.
你想想.
1776立國.
現在總共還不夠300年.
已經.

$^{2001}$你當十幾代.
都已經忘記了工作.
最終是榮耀神.
工作變成是賺錢.
要賺更多的錢.
要更多的時間去享受.
去享樂.
跟著.
最終工作變成.
就是他們的神了.
清教徒完全不是這樣的.
第一代完全不是這樣的.
我們能夠工作.
是證明神的恩賜.
我們去工作.
殷勤的工作.
忠誠的工作.
是為了榮耀神.
誠實.
一心為神.
所以他們一定會為著神.
他們會有奉獻等等.
但是第二代.
第三代已經忘記了.
工作的背後.
是要榮耀神.
還有工作本身就是.
神揀選你的名正.
清教徒很強這個信念.
所以.
你就可以想像到.
保羅都看工作.
在懶惰.
是很嚴重.
他要奉主的名.
要警戒那些不守規矩的人.
要遠離這些人.
如果他繼續在教會裡.
用這種行徑.
來影響其他人.

$^{2041}$去到一個.
他已經成為一個信念.
就要遠離這樣的人.
所以要小心.
叫他去.
去做一些.
要小心的.
叫你不工作的人.
當然香港就少.
我們就是.
叫人不要這麼忙.
叫人真的要多些休息.
很少叫人.
做事.
因為整個香港就是拼命做事.
很少人會.
懶惰不做事.
少的真的少的.
少到一個地步就是.
整個香港都會對你有一些歧視.
30歲有手有腳.
你還拿著碗.
通常如果他有殘廢.
你不會覺得他有問題.
有手有腳你還是這樣.
那就很奇怪.
他的心態很奇怪.
現在我都見過.
真的有些男士.
不做事的.
他有回教會的.
真是懶惰的.
你看到他的樣子.
都感到他懶惰.
未至於.
真是有重病.
人是很皮膚的.
所以這裡告訴我們.
似乎這裡.
如果是一些真是關乎神學.

$^{2081}$就更加嚴重.
就是說他們是很合理化.
我們現在等候主回來.
整段沒有講過.
不過我們推敲上文.
一般都會這樣理解.
他們就是以為.
被提很快到了.
所以我們現在就要準備救急包.
隨時見主.
但是問題就是.
那個時間還沒有到.
他們又不做事了.
又這樣等待.
等待得來自然就會不斷去傳.
你們到底在反映什麼哲學觀念呢.
什麼心態呢.
所以似乎.
包含了不只是不做這麼簡單.
不只是說.
一個現在的.
叫做.
躺平.
現在比較流行的字眼.
就是坐著什麼都不做.
就不是這樣.
而是.
告訴我們.
這個信念.
整段是在講.
因勤工作是一個照明.
是我們.
賴以.
讓神的.
恩典在我們的生命裡面.
原來我們能夠找到食.
是祝福 是恩典.
是照明.
但是記得找到食.
不是那個終極.

$^{2121}$美國人.
已經是一個寫照.
立國.
基督教很開心地.
感恩節獻火雞.
他們隨手捉到的.
當地的農作物.
就獻情給神.
今天的感恩節.
今天的美國人.
只會因勤的工作.
貧富懸殊.
但是.
工作之後.
因勤的去享受.
去享樂 發展.
發明不同的娛樂.
神已經.
不在工作裡面了.
已經不在他們的心靈裡面了.
工作只是.
令他們更加拼命賺錢.
可以更加的去享受.
所以.
現代人必須要牢記.
工作的相反.
懶惰的相反.
就不是一個悠閒.
而是.
要有一個悔改.
所以這一段.
都是一個很值得我們去思想.
我們未必是工作狂.
我們未必是.
只是為了賺錢.
而是.
真是.
禮拜日這個名字是很好的.
要你放下.
一切你的身份.

$^{2161}$你的.
六天的.
穿著的制服.
或者應該說.
面具 你要放下.
你禮拜日.
所有的人.
飛鴻.
其他的弟兄姊妹.
都要是一個虔誠的朝拜.
朝聖者.
這樣就健康了.
你每個禮拜一個循環.
這樣想自己就健康了.
很多人拼命去做工作.
去到禮拜日.
就慢慢不再去教會.
睡到十一點.
因為我已經太辛苦了六天.
然後就去拼命排隊.
去拿籌去喝茶.
喝完茶下午.
百無聊賴.
寧願走兩個小時商場走到腳軟.
不會去找教會.
這個就是工作已經變成了神.
所以如果我們能夠明白.
我們要工作.
是一個照命.
是我們整個人生.
我們有工作的時間.
那個能力.
譬如六十五.
去到身體走下坡.
都是很合理的.
但是可以是另外一種形式.
譬如用我們的智慧.
經驗.
去扶持後輩.
都是工作來的.

$^{2201}$都是讓那個照命去到終結為止.
所以.
這裡提到的工作.
是很嚴肅地.
告訴我們原來我們整個人生.
是要榮耀神.
我們整個人生.
不只是工作,睡覺.
享樂,保障.
而我們忽略了.
工作.
我們休息.
我們去有安全的保障.
但是.
更重要的是.
我們是為了去榮耀.
這位神.
所以求神幫助我們.
因為美國人已經沒有了.
最重要的.
工作.
為神而工作.
所以盼望我們能夠.
明白我們真的要有一個.
為神所做的心靈.
今天我們就說到這裡.
一起一個簡單的祈禱.
再一次多謝主.
我們能夠.
貼上萊加後書.
第二章的兩類的人.
愛慕真理.
和.
不喜歡真理的心的人.
來擺在.
我們的面前.
是保羅.
來看到.
而來警告.
多謝主讓我們看到第三章.

$^{2241}$關乎到工作.
我們是屬主的兒女.
我們要學習.
基督的忍耐.
我們要能夠常常的有主的.
引導.
以致我們愛神.
我們能夠在工作裡面.
看到不只是工作.
有沒有榮耀.
這麼簡單.
不是工作的待遇.
那個.
明星好與壞.
這麼簡單.
幫助我們能夠在我們的人生.
看到原來我們能夠.
工作是真是神的恩賜.
是神的恩典.
是神的揀選.
但是我們盼望.
真的我們要記住.
一切都是為了要去榮耀神.
我們能夠工作.
賺到錢.
我們可以致富.
甚至我們可以有有限.
這一切最終都是要.
願主幫助我們.
讓我們不要覺得一切.
都是我們自己一生的努力.
是的的確.
我們是流汗.
甚至我們用了很多的心力.
很多的勞苦.
但是如果不是神賜給我們.
的生命氣息.
如果不是神給我們看到這個永恆.
我們根本只不過是地球.
上面一個.

$^{2281}$存在很短時間的生物.
願主保守.
讓我們看到我們是有神形象的.
神的兒女.
以致我們重新去看.
我們的工作.
這個照明的時候.
是會有很不同的眼光.
我們仰望禱告.
奉耶穌基督的名.
阿們.
\newpage



\section{帖撒羅尼迦後書}
\label{sec:IJNXTQniUCA}
\textbf{2024-08-22 周四晚查經班《帖撒羅尼迦後書(5)》 講員:曾裔貴牧師}
\newline
\newline
連結: \href{https://youtube.com/watch?v=IJNXTQniUCA}{\texttt{https://youtube.com/watch?v=IJNXTQniUCA}} ~~~~ 語音日期: 2024-08-23
\newline
\newline
\hyperref[sec:bkamj1iGLMk]{\small{< < < PREV SERMON < < <}}
~
\hyperref[sec:index_chronic]{\small{[返順時目]}}
~
\hyperref[sec:index_scriptual]{\small{[返順卷目]}}
~
\hyperref[sec:I4c89OnZC6M]{\small{> > > NEXT SERMON > > >}}
\newline
\newline
$^{1}$我們一起祈禱.
再一次感恩天父給我們晚上.
我們一班弟兄姊妹.
無論在線上的.
或者在現場.
我們都可以一起來聚集.
去研讀這個《帖撒羅尼迦前後書》.
還有我們都很想在前面的日子.
對於末世論這個課題.
都有少許的輪廓的涉獵.
很盼望主都給我們能夠對於主的再來.
能夠有一個認識.
我們去看貼後的時候.
特別可能和他們當時.
對於主再來的一些誤解.
而引申了一些可能在倫理上.
教會的生活上.
需要來試圖提點,教訓,督責的.
都讓我們能夠從其中去找到一些原則.
在教會生活裡面的原則.
以致我們大家在不同的教會.
不同的地方來聚會的時候.
我們都能夠透過這些大的方向原則.
來親近主.
和在教會裡面敬拜.
可以來到有規有矩的去敬拜神.
今天晚上我們再一次仰望在主的手中.
求你帶領奉耶穌基督的名.
阿們.
都想在我們完結.
《帖撒羅尼迦前後書》之前.
來看第三章這一點的一段.
今晚可能下半段.
我們就可以開始去探索一下.
末世論這個大的神學課題.
雖然相對比較少會講的.
但是遲些都會深入淺出.
可能看看用三堂還是四堂.
和大家去講一講.
有不同的學派的.

$^{41}$有哉前,哉中,哉後.
有所謂前遷,有後遷.
也都有一個叫做無遷.
無遷希年.
大概是這幾個學派的神學論述.
各有它的經文支持.
所以不會有一個百分之一百肯定的.
唯一肯定的 大家都肯定的.
三派都肯定的.
唯一肯定的就是主耶穌一定會回來.
但是怎麼回來.
那就有不同的解經的學派的論述.
我們回頭看最後這一段聖經.
上個星期我們都稍微涉獵過的.
很記得第六和第十二節.
都是兩次提到.
靠著基督耶穌.
奉主耶穌基督去吩咐.
在教會裡面不守規矩的人.
所以你可以看到保羅.
他是有那個power 有那個權柄.
有那個神賦予他的那個位份.
來讓弟兄姊妹要有規矩.
尤其是引申到夏文.
就是讓那些不信的人.
不要錯誤的理解.
以為基督徒就是這樣的.
好食冷飛.
說主再來就不做事.
就帶來一個混亂.
如果主耶穌當晚回來.
其實都不等於你那天不需要做事.
或者你那個月不需要做事.
所以我們真的要記得.
不要再去計算那個日子.
一個確實的日期.
是計算不了的.
只能夠給我們帶著一個警醒的心.
耶穌一定會回來.
這個穆斯林必須要牢記在心裡面.

$^{81}$我們成長的年代.
80年代信主很普遍.
大家都去計算主耶穌何時回來.
或者見到當時所謂歐盟的十國.
十隻腳趾.
就將《但義理書》第九章的經文.
這樣來讀進去時事.
或者是新聞裡面.
很容易就會有這種錯誤的理解.
所以我們都要很小心.
千萬不要計算.
當然我們會有一些原則的.
去到穆斯林.
有幾段重要的經文.
成為我們很重要的規範.
和那個脈絡.
我們再看回這裡.
其實經文很簡單.
不過要實行起來.
其實就很需要有智慧.
有膽量.
有那個的慧份.
來去做.
所以保羅特別是用.
他不是用那個權.
他是用那個榜樣.
來去勉勵當時教會.
必須不要讓那些不循規蹈矩的人.
來到教會.
或者在教會以外.
來去帶來那個滋擾.
帶來那個讓真理.
被人誤解.
所以其實都需要有這種智慧.
和那個膽量.
來去特別是奉主耶穌的名.
靠主耶穌的吩咐.
來去勸誡這樣的人.
我們再讀最後這一次.
第6到第18節.

$^{121}$三章的《貼索羅尼加》後書.
二章很明顯就是我們上一次.
都一而再講了兩堂.
都想講到就是.
有兩類人.
有一類是很愛慕真理的人.
是跟隨聽從保羅的教導的.
他們所領受的教訓.
他們會遵守 會堅守的.
還有他們很明顯.
因為他們對神的話語的愛慕.
反證他們是被神揀選的.
他們是信真道.
而且是被聖靈感動.
成為聖潔能以將來的得救.
所以你會看到.
好像真是有一個所謂的天淵之別.
愛慕真理的人.
會越來越對神有渴慕.
有生命的流露.
但是沒有真理的人.
或者根本不信的人.
就會讓魔鬼可以在他們的內心.
來去搞擾.
記住一個叫做撒旦的運動.
二章的九節.
將來是會行各樣的異能神蹟.
和一切虛假的歧視.
而且是在那些注定沉淪的人身上.
會行出於不義的詭詐.
那些詭詐.
說明詭詐了.
還要是不義的詭詐.
就是說.
這些人不能夠對真理有一個透徹的明白.
融會的貫通.
所以原因就是因為他們不是不聽道.
是不領受愛真理的心.
使得他們得救.
但是跟隨主.

$^{161}$主所愛保羅教導的那幫人.
他們能已得救.
但是那些人就是不能得救.
你可以想像到.
第二章就明顯說.
要很小心.
保羅的年代都已經有了.
何況離開了世界的使徒.
之後的時代.
更多這些異端的出現.
所以去到主後的.
大概是三百多年.
開始必須盡快確立正典.
已不鎖的會議.
教父時代.
要盡快確立聖經的正典.
就因為已經充斥了很多異端.
很多人假冒那些書信.
還有開始人云亦云.
對於真理.
因為大家都沒見過耶穌.
使徒都死了.
那些書信還沒緝錄成為聖經.
所以我們其實很幸運.
我們已經有了66卷的聖經.
當時代他們是要怎樣醞釀.
這39卷就不用醞釀了.
因為舊約已經有了.
27卷的新約.
就需要有這些大公會議.
來緝錄.
其實應該是民間去搜尋.
哪些經卷已經在教會流通中.
已經有權威性.
例如羅馬書.
肯定是保羅所寫.
就不需要去爭論.
所以希伯來書,啟示錄那些.
去到很後期才放在正典.
你可以想像到.

$^{201}$保羅的年代.
已經有那些無拒名的書信.
而衍生的就是那個撒旦的運動.
後期,後代更加會再多一些.
所以必須有這些教父去捍衛那個真理.
確立.
所以善徒信經就由那時候開始.
大家要在聚會裡面.
最簡單的.
因為那時候識字的人不多.
最簡潔的就是對三一神.
是怎樣的呢.
那個真理的確認.
一個確信.
教會是甚麼來的呢?.
聖靈又是甚麼來的呢?.
到底我們信耶穌的人是不是有永生的呢?.
很多這些基本的教義.
就是這樣來的.
從那個時候開始就一直流通到今天.
當然.
使徒信經以後.
會不斷的衍生不同的信經.
有很長的.
尼西亞信經.
阿他拿修信經.
還有其他的.
加克敦信經.
很多這些的.
接著就會開始有很多的教理問答.
問你人生.
這個是.
教理問答.
這個是.
宗教改革的時代.
第一條.
人出現在這個世界.
唯一要做的事.
最重要的事.
不是生老病死.

$^{241}$是.
小朋友已經要認識了.
是很重要的.
第一條就是你的人生就是要榮耀神.
是很重要的.
所以那些小朋友.
教那些小朋友.
也是 離開天主教.
他們就開始有這些兒童主義學.
所以就是這樣.
循序漸進.
將真理去教導那些小朋友成長.
對真理 將很多的神學.
化做一些教理問答.
就是我問你答.
就變成一個.
他們就一定要認識.
所以這裡真的告訴我們.
在教會裡面.
必須是要有一個穩固的.
信仰根基.
和真理的.
教導的系統.
讓弟兄姊妹.
能夠.
不是憑感覺.
而是憑真道.
去成長.
所以第二章是一個很強烈的對比.
保羅的年代.
已經這樣來警告.
你不愛真理.
你不喜歡聽道.
你又不讀經 你要小心.
如果你裡面沒有一個愛真理的心.
那就是一個問題.
那你就不得救了.
但是反過來.
13到17節.
這些就是真的弟兄姊妹.

$^{281}$真是神的兒女.
他們有什麼行徑呢.
他們是怎樣的呢.
因信真道.
被聖靈感動.
成為聖潔.
即是內心不住要追求.
要像耶穌.
得到神的喜悅.
不敢得罪神.
能以得救.
不是靠行為得救.
是他們已經有這樣的成長.
因為他們有了新的生命.
他們開始很想行神的旨意.
所以神藉我們.
所傳的福音.
照你們去到這個地步.
得著的是主耶穌基督的榮光.
所以.
不再是追求.
自己的自我中心.
要怎樣能得到主耶穌基督的榮光呢.
所以.
仍然都是站立得穩.
用神的道.
給自己去勸勉.
你就可以想像得到.
一個基督徒.
有神生命.
和一個沒有神生命的基督徒.
他是.
根本不會愛慕.
神的真理.
沒有一個這樣的心.
即是一個洋娃娃.
穿得多漂亮.
扮得多漂亮.
穿得多漂亮的衣服.
她是永遠都這麼漂亮的.

$^{321}$因為她根本不會長大.
一個嬰兒.
她就算到處拿東西玩.
沒有人理會.
又搞這些東西.
髒東西.
你不要理會她.
慢慢她會長大.
她會慢慢起來.
她會越來越是一個成人.
一個.
回教會的人.
理應是有生命的人.
才會回教會.
所以沒有生命而回教會.
其實是一個很大的危險.
因為撒旦會.
抓這樣的人.
進入和他一樣的.
沉淪.
這是第二章.
我講多了.
本來應該講三章的六到第十八節.
再溫習.
第十節.
在這些沉淪的人身上.
行各樣出於不義的詭詐.
是撒旦的運動.
原因.
就是因為.
既然你不愛真理的心.
你就喜歡那些詭詐.
喜歡那些不義.
喜歡那些很奇妙的東西.
例如那些異能 神蹟.
虛假的歧視.
很喜歡這些.
很有那種口味.
去看這些神蹟.
但是心底裡.

$^{361}$沒有愛真理的心.
所以神會.
將一個生發錯誤的心.
叫他們順從虛方.
你說神這麼壞.
他不行就要給他一個.
真理的心.
其實應該說與神無關.
他根本就不愛.
根本就沒有生命.
所以他自己會一個.
不斷生發錯誤的心.
順從的只是虛方.
記不記得十災.
不是十界.
是十災.
摩西就經常和白鷺.
在那裡鬥智鬥力.
水變血.
白鷺說.
我的術士都行.
不用怕你.
跟著青蛙.
然後那些瑟拿 蒼蠅.
那些術士都做到.
很快就沒事了.
好像都不是很怕.
摩西那樣.
但是你聽聽.
《出埃及記》怎麼描述呢.
每一次當法老.
那些術士.
就是他的手下.
能夠解除了那些.
所謂魔咒.
神的審判的災禍.
解除了之後.
跟著加了一句.
每一次的神蹟之後.
都加了一句.

$^{401}$神使得法老的心.
僵硬.
法老的心是越來越硬.
就是我都不怕你.
摩西 後來有些他真的.
做不到 他就求摩西.
摩西就再去求神.
神就使得那些災患.
就舒緩了.
但是舒緩之後.
法老又把心.
一硬.
不撇了.
那你就可以想像到有一個人版.
經過這麼多的異能.
是真的異能.
是神蹟的十災.
但是都不要了 都站一站.
因為跟這裡第二章的十節.
這一段都有一點關聯.
希望教會裡面.
如果這堂課有人聽到.
我就是還沒信那個.
那我要快點信了.
不然.
一個心法錯誤的心就死了.
其實很多的.
第七章已經有了.
摩西出埃及記.
耶和華對摩西說.
我使得你在法老面前代替神.
你的哥哥亞倫.
是替你說話.
凡我所吩咐你的.
你都要說.
你的哥哥亞倫要對法老說.
容以色列人出他的地.
留意了.
我要使法老的心剛硬.
也要在埃及地多行神.

$^{441}$即其是.
神法老必不聽你們.
我要伸手重重的迎罰埃及.
這裡說神使法老的心剛硬.
跟著去到後面.
法老自己把心一硬.
不是把心一橫.
是把心一硬.
即我根本不怕你.
我都做到你的法術.
有哪幾段呢.
我們稍微看一看.
都可以看到的了.
留意第八章已經有了.
八章的十災.
記住要記熟背.
小朋友要記熟背.
你不要記背第一個和第十個.
因為是實識的.
水變血和殺長子.
之後的八個你分兩對就行了.
蛙失凝蓄.
如果小朋友.
兒童主義學就要跟我讀一次.
這裡不是兒童主義學.
不用了.
蒼薄黃黑.
記住十災就背好了.
失災.
八章的十六節.
耶和華吩咐摩西說.
你對阿倫說.
伸出你的杖擊打地上的塵土.
使塵土在埃及遍地變作蝨子.
即是跳蚤.
很癢的.
他們就這樣行.
阿倫伸杖擊打地上的塵土.
就在人身上.
和牲畜身上.

$^{481}$就有了蝨子.
埃及遍地的塵土都變成蝨子了.
留意第十八節.
行法術的也用邪術.
要生出蝨子來.
卻是不能.
於是在人身上和牲畜身上.
都有了蝨子.
行法術的就對法老說.
這是神的手段.
已經去到不做了.
術士都做不了了.
不可以.
因為這些災難可能太大了.
他們做不了了.
留意一下.
法老仍然心裡一硬.
不肯聽摩西阿倫.
正如耶和華所說的.
這只是其中一個描述.
所以.
法老是不斷把心一硬.
把心一硬.
去到一個.
已經是回不了頭的.
一個這樣的情況.
所以用回保羅的說話.
就是神會使他們有一個.
生發錯誤的心.
所以不要小看.
我們日積月累.
都完全對神的道.
沒有興趣.
這個要很小心.
這個可能是.
反映我們裡面.
有一個錯誤的心靈.
我們是一個生發錯誤了的心.
我們順從虛方.
一切不信真理.

$^{521}$喜愛不義的人.
都是這樣被定罪.
歷世歷代.
由法老的開始.
一直到保羅的年代.
都是每一個時代.
總是有些人接觸過基督教.
接觸過耶穌.
接觸過猶太教的真神.
但是都是離救恩的門.
遠遠去.
這是很可惜的.
我也曾經說過.
在之前的《查經班》.
我初信的時間.
有一個弟兄.
很資深的.
雖然他跟我差不多大.
有一次見他.
都少回教會.
剛好又見到他.
就問他.
跟他一起吃飯.
他就很認真跟我說.
我不會再信這位神.
可能他有很多的際遇,經歷.
自此就沒有機會見.
比我更資深.
那是幾十年前了.
你可以想像到.
人生往後.
不再在教會.
不再在真道上去成長.
所以這些是要.
引以為戒.
但聖經一早已經說了.
說到主再來.
不要輕易的動心.
這是一個很吸引的題目.
會動心的.

$^{561}$人家會告訴你什麼時候主回來.
所以我們要留意.
原來我們.
對真理的.
研讀,追求.
和在生活裡面.
去實踐神的道.
就是最大的.
等候主再來的.
根本.
否則就會.
很容易受到那些不法的.
運動,撒旦的.
運動的影響.
那裡有描述.
這裡到底.
是否真的所指將來.
那個大罪人.
因為這段其中都是末世論的一個經文.
二章我們之前說過.
那個沉淪之子.
到底是否真的.
啟示了講的十一章的.
那個迪基督.
那個叫海獸.
是否呢?.
我們很難將幾段經文.
將它們融合的.
但又很像那個氣氛.
就是說.
似乎末世論.
不是越來越傾向好.
是傾向不好的多.
就是沉淪之子.
的出現.
會帶來教會的.
那個逼迫,受虧損.
我完結了這裡.
當第六節開始.
我們再看回一個生活的倫理.

$^{601}$就是因為那個神學.
的一個扭曲.
而導致在教會裡面.
已經有一些人.
他們.
是振振有詞的.
以為.
主回來我們就可以.
專心去等待.
出現的後果.
就是有些人不做工.
來去.
好食冷飛.
第六節.
弟兄們我們奉主耶穌基督的名.
吩咐你們.
凡有弟兄不按規矩而行.
不遵守從我們所受的教訓.
就當遠離他.
基督教是有我們特定的真理的.
基督教是有一個.
的道統的.
是有一個規模的.
是由使徒一直傳下來.
所以三一真神.
已經不需要再討論的了.
我們要知道.
我們就不是去質疑.
因為質疑不到.
因為我們根本沒有這樣的智慧.
去明白什麼叫三一.
怎樣三位神又變成一.
又說是獨一的神.
已經是我們不討論的了.
必須這樣相信的了.
基督教我們自己有這樣的教義的.
其他的宗教是不會有的.
所以使徒已經給了我們一套教訓.
也都經過不同的年代.
將這些神學繼續.

$^{641}$傳適得更加清楚.
將這些神學繼續傳適得更加清楚.
所以這些告訴我們.
如果那些人不遵守.
由使徒所傳下來的教訓.
就要遠離他們.
你們自己原知道.
應當怎樣效法我們.
因為我們在你們中間.
未嘗不按規矩而行.
所以這是一個很好的榜樣.
使徒他們也未嘗白吃人的飯.
倒是辛苦勞碌.
晝夜作工.
免得叫你們一個人受累.
你可以想像到.
保羅在第二次宣教旅程.
他特別強調.
他在哥林多的時候.
當時有一對夫婦叫阿居拉夫婦.
阿居拉夫婦.
白基拉和阿居拉.
兩個由羅馬.
因為要被當時的皇帝.
應該是第四任的皇帝.
就是格魯丁.
應該是第四任的皇帝.
來到去要驅逐猶太人.
離開羅馬.
阿居拉夫婦離開羅馬之後.
輾轉去到哥林多.
見到保羅.
保羅就是那個織帳篷.
就是猶太人的手藝.
平日保羅有個地方.
去織帳篷.
可能是那些.
波斯地氈.
來做售賣.
阿居拉夫婦和他的同行.

$^{681}$跟他一起.
在那裡織帳篷.
在那裡記載.
所以保羅這個時候.
就是第二次旅行.
宣教旅程.
所以保羅在貼索羅尼加.
給他們一個榜樣.
他自己也是.
在那裡辛苦.
勞碌晝夜的作工.
當然他不只是去做.
他在安息日也會去到.
不同的會堂.
去傳道.
為何大人信主.
這是他更加重要的.
原本的工作身份.
那個神.
托付他的.
位份.
所以現在就這樣.
去講.
保羅也挺有智慧.
挺謹慎.
他去傳道.
對一個不太認識的教會.
他不敢運用.
沒有說不敢.
有一種敏銳的心靈.
很自律地.
不太認識.
不太有心交的一班.
貼索羅尼加教會的人.
他就在他們中間.
來展示.
他也不會運用.
神給他的.
差遣,權柄.
來勉強人.

$^{721}$要給他那些.
工錢.
反而他自己是親手.
做的.
貼索羅尼加這樣講.
但肥勒比書完全不是這樣講.
你可以想像到.
保羅在不同的群體.
他會展現.
這種彈性.
牧羊上的彈性.
你可以想像到.
保羅他真的對於錢財.
他不是看得那麼緊要.
那麼重.
再留意.
第九節.
這並不是因我們.
沒有權柄.
乃是要給你們.
作榜樣.
叫你們效法我們.
他這樣身體力行.
以身作則這樣做.
就是會讓.
那個真道他去講的時候.
是會不會受到.
人的責難.
誤會等等.
等一下我會站一兩段經文.
剛剛在另外的經文.
就真的講.
關於.
供給傳道人的.
遺訓上的教導.
我先讀完.
等一下再貼一兩段經文.
是保羅自己引述.
舊約申命記二十五章.
講到那些牛在踩谷的時候.

$^{761}$耕田的時候.
或者可能他轉圈.
在磨谷的時候.
就不要.
不要摟著牠的嘴.
要做的話要給牠吃.
他自己在田裡.
咬一些谷來吃.
一直在那裡消化.
就是這麼多.
牠吃幾粒就讓牠吃.
吃得你幾多.
我記得我那時候.
很小的時候 十八歲左右.
家裡是養豬的 我們要去酒樓.
收那些廚餘.
以前叫哨水.
沒那麼清高的.
好聽的名字 去收哨水.
在一間酒樓.
我們在後巷進去.
上到去.
是後巷.
上到去就是.
廚房.
大廚房.
很大的廚房.
那時候不會在深圳做 是在大廚房做.
點心.
跟那些伙計很熟.
上到去 從小時候.
十七八歲 看到一些人推東西.
準備去酒樓.
因為有幾層.
在後樓梯有電梯.
推那些餐車.
我一上去.
看到蝦餃就拿來吃.
其實我是倒哨水的.
但是那些伙計很熟.

$^{801}$看到蝦餃就拿了一粒.
他就要幫我補.
不要弄著牠的嘴.
當然是開玩笑.
但本來就印這節經文.
《新聞記》25章.
等一下我會和大家揭.
在《新聞記》這段經文.
在提摩泰前書五章.
在哥林多前書九章.
明顯地引述.
引述完之後還要解釋.
關乎到要供給傳道人的那個教導.
這樣的意思.
這樣就是說.
一定是衍生了一些問題.
才要這樣教導.
但是在這裡完全沒有教導.
就是說我就是你們的榜樣.
我是親手勞碌的 作工的.
我就是你們的榜樣.
我是親手勞碌的 作工的.
似乎因為這個教會我們說了.
四個星期保羅就要離開了.
因為有逼迫.
他就要走了.
根本就沒有在他那裡.
有什麼很重要.
停留一段很長時間的一個工作.
所以未必能夠.
和他們真是有很深入的交往.
和認識等等.
所以保羅就.
透過他自己在哥林多.
當然他還沒有到哥林多.
如果剛剛傳道的時候.
在第二次的旅程.
大家要記住.
《小徒行傳》的經文.
大家都要有一點點的輪廓.

$^{841}$當你看這些新約的經文的時候.
對於你能夠了解保羅.
他們的行程.
和他們所做的.
要有一點點的串連.
就會有幫助的.
所以他在.
這個.
被趕走之後.
輾轉.
他就去到.
比利亞.
接著就去到阿波羅尼亞.
帖撒羅尼迦.
接著就去到哥林多.
所以.
他是很趕急的.
所以他日間所做的事情.
就是真是.
做一些這樣的跡象棚的工作.
或者幫其他人.
當時的教會.
帖撒羅尼迦教會.
受到教會.
因為教會剛剛開始.
根本都沒有一個很有規模的.
所以就沒有一個.
恆常的穩定的一個機制.
來對當時的.
駐堂的一些牧者.
的那個的功級.
所以留意第十節.
我們在你們那裡的時候.
曾經吩咐你們說.
如果有人不肯作供.
就不可吃飯.
到底是一個恆常的.
有些人本身的.
習性是好食冷飛.
是.

$^{881}$習慣這樣.
還是那個神學導致他們.
主載來什麼都不要搞了.
要等候了.
會不會也是這樣呢?.
這裡就不詳細說.
到底背後什麼原因.
但是.
就出現了一個這樣的後果.
在帖撒羅尼迦加教會.
就有一些的.
似乎都是信徒.
有一些的迷思.
就是不作供.
就不是不吃飯.
不是禁食.
是不作供.
人家供應他.
所以保羅就說.
他們如果不作供.
就不要給他吃飯.
就是不要供應這一類這樣的人.
就不是不供應其他人.
有困難的人.
所以就是11節.
因我們聽說在你們中間.
有人不按規矩而行.
就是說知道規矩.
而不按規矩行.
那就.
明顯的有他自己一套的.
可能是神學.
可能是他自己一套的.
他自己認為的生活的.
他認為的標準.
這些都是不健康的.
就是香港暫時.
雖然這個.
又是有中性的.
又不可以說是一個絕對的案例.

$^{921}$香港如果.
丈夫不工作.
都有一點點社會壓力.
因為.
華人的觀念還是.
男的主外.
要勤奮工作.
沒問題的.
如果那個女的.
又主外又勤奮工作.
而能夠大家都有一個平衡的.
就沒問題了.
男的就主內做家務.
可能這樣是更加好.
每個家庭不同的.
但是如果兩個都這樣的情況.
不按規矩而行.
什麼工都不作.
又不可以說得太白.
太白就變得不太好.
因為這些私隱.
因為我真的見過有一個稱為弟兄.
真的不工作.
就.
又拿綜援.
又什麼的.
真的經過這裡的教導.
好像牧者也有.
警告過他.
當然就離開了.
沒有再回來了.
真的不是說他.
身體有毛病.
或者他不是不工作.
而是他真的沒有想過要工作.
還要年輕搖搖.
你可以想像得到.
根本就不要說神學.
不是神學問題.
根本就是一種.

$^{961}$懶惰的問題.
當然這些是比較極端的.
在一個華人社會是比較極端的.
基本上是有一個很大的.
壓力的 群眾壓力的.
這裡就告訴我們.
第十一節就說.
有人就已經有不按規矩而行的.
那個行徑.
一直都是留意.
在教會裡面發生.
所以保羅才這麼緊張.
記住.
凡是奉主耶穌的名 靠主耶穌的名.
一定是在說教會內部的人.
外人.
靠什麼耶穌的名.
一來都不信耶穌.
那裡就有社會的制度去干預.
或者他自己要去面對.
那個收入.
他自己要面對.
那個生活.
我們說的是教會.
因為教會是有一個慈善的機制.
是幫助那些有困難的.
比如無依無靠的寡婦.
我們上個星期就提過.
和柯特前書第五章.
原來上個星期我們看漏了.
跟著下面那段.
就教導怎樣去.
供給長老和傳道人.
所以你就可以想像到.
教會是有一個健全的.
完善的機制.
是真正去.
收集了弟兄姊妹的獻金.
是幫助有困難.
有真正困難需要的人.

$^{1001}$而且是彰顯了一份的愛心.
和教會以外.
沒有人會做的那種美德.
因為無依無靠的寡婦.
在當時的社會是無助的.
是絕望的.
是不會有社會保障.
是不會有社會福利處.
是不會有社會的.
是不會有關愛隊去派飯.
沒有的.
教會就要做這些.
所以時代不同了.
但是那個美德.
那個的留心.
哪些真正有需要的人去幫助.
是從來都沒有減少的.
所以.
如果現在的社會說.
一些獨居老人.
一些失救或者照顧者.
受到很大的壓力.
已經有這樣的社會問題.
教會應該要支援.
以及.
應該要分一些資源.
怎樣去做到這些關愛.
所以這裡就圍繞著.
一定是奉主耶穌的名.
在教會內部.
要令到那個制度要健全.
不要有一些.
類似是違規的人.
去想搞亂了那個制度.
讓這個制度.
受到挑戰.
人的心就是.
原來教會的錢是這樣用的.
用在這些人身上的.
那就不行了.

$^{1041}$那你帶來的就是.
原來拿綜援可以不用做.
還可以到處去旅遊的.
那我也拿綜援了.
意思就是這樣.
何想而知 本來就是關心的.
不是要針對那個人.
而是針對那個制度有沒有被人濫用.
去挑戰.
所以我們要留意.
基督徒.
我們都真的要有智慧.
不要因為我們說基督徒.
我們就要.
對一些人.
沒有了那個.
小心謹慎.
特別是已經是已經.
不按規矩而行的人.
似乎不是一次.
他已經有這個習性.
那就要留意.
我們要怎樣去處理.
因為他在教會裡面.
帶來的.
而且不是他個人的問題.
而是他已經開始.
一個不做工的人.
真是太有空.
就會專管閒事.
所以.
似乎又是孖公仔.
一不做工.
他就很有空就專管閒事.
12節.
我們靠主耶穌基督吩咐.
勸戒這樣的人.
所以在主裡面.
祝福 勸戒.
要安靜的做工.

$^{1081}$吃自己的飯.
自己的勞力 自己用神給的時間.
因此.
工作無分貴賤.
記住那個無分貴賤.
是你要知道.
這個召命的人生.
神安排你的人生崗位.
那個無分貴賤.
我們要記得.
否則.
你做多高級.
如果你的心根本不是為神.
其實在神眼中.
你會低下過一個.
可能一個勞工.
他很默默地工作.
一來就是他很理直氣壯.
賺錢很辛苦.
滿身大汗.
但是為家人去謀福祉.
而且他很熱切.
很想來服侍.
這樣的勞工.
比一個可能很出色的經理.
跨國公司的經理.
如果他根本沒有想過去為神.
高尚的.
那個工作.
人心靈.
對神已經忘記了.
那個工作本身的背後.
是那個問題所在.
所以在這裡 保羅就勸戒這樣的人.
要安靜的做工.
吃自己的飯.
弟兄們 你們行善不可喪志.
所以.
在教會裡.
要讓那個機制健全.

$^{1121}$我們不要因為一個.
壞分子.
變成我們一拼.
什麼都不要管.
教會原來這麼複雜的.
我們不背了.
我們沒理由被那個人.
破壞的一個行動.
那個行為就影響了整個.
本來是神喜悅的一件事.
就是教會.
要幫助真正有需要的人.
的這個美德.
所以保羅就.
真的勸勉那些有心的弟兄姊妹.
所以我們很多時候.
我們一見到一點點的東西.
我們就馬上不再回教會.
我們有時候都好像.
沒有了這種學習.
保羅就很明白.
人會受影響而灰心.
所以行善不要喪志.
如果有人不聽從我們這些信上面的說話.
就要記下他.
不和他交往.
叫他自覺羞愧.
所以要讓那個制度完善.
不是要將制度由你去瓦解它.
你帶頭.
灰心.
然後就不再行善.
就不是這個意思.
反而我們就應該更加有弟兄姊妹.
有這種公正.
有這種公義.
來處理這些事情.
所以叫他自覺羞愧.
讓他看到.
他自己的問題.

$^{1161}$但不要以他為仇人.
要勸他如弟兄.
又說不要交往.
又說要記下他.
叫他自覺羞愧.
自覺羞愧其實都是他自己被感動羞愧.
這個過程可能需要緩慢.
需要有人私底下勸說他.
和他講.
不要在講台上指名道姓.
誰誰誰.
你這麼懶惰.
這樣很難自覺羞愧.
因為這樣的時候.
真的是在眾人面前羞辱他.
這樣其實就不行.
要很小心.
不然又會帶來一個誤解.
所以這裡告訴我們.
要很有智慧地.
有幾樣東西.
教會.
行善.
建立健全的機制.
這個美德是神喜悅的.
神要求我們做的.
我們就要做了.
不要因為一個人.
或者一小撮的人.
或者一些錯誤的想法.
來挑戰那個機制.
可能他是有意的.
可能更多是無意的.
但是我們就要小心.
我們是那個管理者.
你不是這樣.
這裡說很慘很慘.
然後你轉頭就走去馬會.
那怎麼辦呢.
我們知道之後.

$^{1201}$首先不要懷疑我們自己的愛心.
懷疑我們自己的團隊.
我們的教會的機制.
不要將那個機制打垮了.
不要因為一個人賣馬就打垮機制.
而是處理那個人.
要怎麼處理呢.
調查過.
你這樣其實不對.
是什麼原因.
了解了原來他真的這樣做.
他肯悔改了.
自覺羞愧.
再勸他如弟兄.
不要看他是仇人.
往往我們就用群眾的壓力.
有時未審就先判.
有時審了都不肯饒恕.
多的事的教會的紛爭.
往往在台下.
來傳那些是非.
傳到一個地步.
整件事失實.
你叫人怎麼自覺羞愧呢.
教會的機制.
蕩然無存.
就會帶來將事情更加複雜.
更加麻煩.
我們就容易.
對一些真心想悔改的人.
都因為那個輿論.
被逼要離開.
可能他更加是.
把心一橫.
有恃無恐.
這個教會這樣對我.
真心想悔改.
我做錯了一次.
當我是一個仇人.
當我是永不超生.

$^{1241}$我怎麼能夠在這裡生存呢.
他就走.
可能帶著的就是仇恨.
帶著的是一種憤怒.
他就不能自覺羞愧.
他是自覺被人羞辱.
兩樣東西是不同的.
你是一個成年人.
你會很有智慧地.
記住那個機制.
是永遠都需要存在的.
那個美德.
我們怎麼來把持.
一代一代傳下去.
教會的確是講愛心的.
我們是要幫助.
有需要的人.
所以我們不要質疑那個機制.
反而我們就要.
怎麼樣讓那個機制更加透明.
更加完善.
使得更多神要我們去幫的人.
得到幫助.
這樣的一個運作.
所以這段經文.
其實很簡單.
但其實很寶貴.
講到教會內部.
我們真的需要奉主的名去看很多事情.
裡面所包含的就是.
神的教會是有規矩.
就是說有真理.
就是說有一個標準在.
我們就不看標準.
我們很多時候.
就會看.
我們很多時候就會覺得.
什麼這麼嚴.
什麼這麼鬆.
我們就好像都不去信任.

$^{1281}$管理團隊.
我們很低調地.
默默地想挽回一個人的.
那個功夫.
往往.
真的很容易就會傳開.
那些的是非.
因為.
是非永遠是最好聽的.
道理.
變成很容易就會對.
無論是機制的挑戰.
當事人.
就會感覺到無地自容.
而受到.
近乎的迫害.
他就一定要離開.
離開是不會帶著自覺羞愧的.
離開只會帶著一種憤怒.
他不能夠回到神那裡.
這些都是.
真的會.
在教會會發生.
所以記得.
這裡有一個精髓就是.
後面夾雜著第六第十.
二節.
在主耶穌基督的.
明吩咐下去做.
也就是說.
容讓管理層.
去做這些事情.
低調地帶著禱告.
帶著愛心.
去做.
信任信納.
管理層的調動.
安排.
使得大家.
不知道的.

$^{1321}$未能公開的.
信任管理層的.
守口如瓶.
因為有些事情.
暫時不能公開.
當那件事.
明朗化.
那個人又真的自覺羞愧.
他又回到神的裡面.
那件事就解決了.
所以這樣.
才是聖經的.
教導我們.
對一個人的挽回.
和對教會的.
健全機制的.
一個uphold的捍衛.
否則.
只會又帶來機制.
又受到挑戰.
挑戰就是很多人心灰意冷.
你們這班人都不知道怎麼搞的.
奉獻給你.
錢亂用的.
我不再奉獻了.
很快就會彈到極端.
機制就沒有了信任.
沒有了信心.
其實那個人不吃飯.
不做工.
就不要這麼快賴機制.
但是很多時候.
就不容許那個時間.
那個人又不會改.
因為被你們千夫所指.
那個機制.
就令很多人心灰意冷.
又受到挑戰.
兩者都不好.
但是保羅的教導就真的很好.

$^{1361}$既有權威.
保羅說要奉主的名.
有權威.
奉主的名其他的弟兄姊妹.
要不就參與.
這個管理團隊.
否則就盡量.
默默的等待.
那個事情的發生.
和那個處理.
那就是一個信賴和信任的問題.
所以記住.
奉主耶穌基督的名.
我們要背著.
這個擔子.
是神要我們去做.
其實將來.
都是要向神交代.
所以否則.
很多時候民間已經先審先判.
未審先判.
只會搞亂機制.
和那個人.
會改不了.
所以這個都是我們.
大家一起要去學習.
記住裡面就包含了.
不只是這個這麼小的事.
就是教會裡面.
太多其他複雜的問題.
我想我們在這個大原則下.
就是第六.
第十二節靠主耶穌.
奉主耶穌的名.
的底下去處理.
所以我們.
管理的.
是真的需要在主的真理下.
去看一些事情.
清楚.

$^{1401}$要去了解.
要去調查.
要找回.
雙方面的真相.
盡量尋求那個.
盡量的真相.
如果還沒有就不要.
這樣去未審先判.
或者有了我們自己的.
judgment(判斷).
先日為主 又懂得眼鏡.
那就很快很容易就會.
有一個所謂的結論.
往往那個結論.
都是我們對那個人的喜好.
和那個人的不喜好.
而作出的決定.
因為已經有一個印象在那裡.
很多時候就不得客觀.
所以這一段經文.
就是關乎到.
保羅的教導.
也都真的有些案例.
是有人不肯作供.
說開這裡.
就關乎到.
保羅自己第九節說.
這並不是因我們沒有權柄.
關乎到.
獲得.
教會的支持.
在金錢上 在他的生活上.
乃是要給你們作榜樣.
我們看兩段聖經.
當然有其他.
不過我就不看那麼多.
譬如哥林多後.
對不起 加拉太書五章.
不知道你有沒有留意到.
那裡都有教導.

$^{1441}$反正我都揭到了.
我讀一讀給大家聽.
聽一聽就行了.
對不起 是六章的第六節.
先是六章.
加拉太書六章.
弟兄們 如果有人偶然被過犯所聖.
你們屬靈的人.
就當用溫柔的心.
加拉太書的六章.
關乎到金錢的奉獻.
就當用溫柔的心.
把他挽回過來.
就是這個.
要很柔軟的心靈.
一方面聆聽.
一方面盡量了解對方有什麼困難.
所以屬靈的人.
屬靈的人.
這裡說的似乎是有智慧.
有內心.
有眼光.
因為聖文說的屬靈的果智.
是有眼光的.
屬靈的人就當用溫柔的心.
把他挽回過來.
也就是不偏不倚.
又當自己小心.
恐怕也被引誘.
你們各人 屬靈的人.
恐怕也被引誘.
你們各人的重擔要互相擔當.
聖文那裡就是互相扶持.
如此就完全了基督的律法.
其實加拉太書.
就是叫我們要有一個.
溫柔的心靈.
我們要有一個溫柔的心靈.
加拉太書就是.
叫我們不要再遵守律法.

$^{1481}$不要再守國禮.
這個是加拉太書三章.
一直給我們的教導.
整本書就是.
責備那些行國禮的人.
但是保羅說.
原來我們互相去擔當.
互相去扶持.
就完全了基督的律法.
就說基督的律法.
不是在文字上死守.
而是在我們的實踐裡面.
見證著.
我們已經超越了.
基督的律法.
所謂基督的律法.
不是說現在有一套新的.
就不是十誡.
就不是這個意思.
保羅的意思就是說.
基督本身就完成了.
神的要求.
已經滿足了神.
對於罪.
對於人要做的.
他已經完成了.
所以我們在主耶穌裡面.
那個就不是一個.
條文的誡命.
而是一個愛的法律.
所以.
我們互相擔當彼此的.
重擔.
這個就是完全了.
或者完全滿足了.
基督的律法.
第三節.
人如果無有,自己還以為有.
就是自欺.
各人應當測驗自己的行為.

$^{1521}$這讓他所誇的.
就專在自己,不在別人了.
因為各人.
必擔當自己的擔子.
又說互相擔當.
這裡又說要自己擔當.
兩者的微妙.
其實字裡行間都已經很清楚.
就是你自己.
根本就不要看自己.
以為自己是什麼.
你是沒有.
你根本就是沒有.
所以沒有就要追求.
就要擔當.
自己的擔子.
不要一種狐假虎威的基督徒.
來以為自己.
就已經達到.
第七節.
不要自欺.
對不起,第六節.
在道理上受教.
當把一切.
需用供給施教的人.
不要自欺.
神是輕慢不得的.
人重的是什麼.
信著情欲殺種.
必從情欲收敗壞.
信著聖靈殺種.
必從聖靈收永生.
信著聖靈.
殺種.
就是說.
就是第九節.
我們行善不可喪志.
若不灰心.
到時候就要修成.
所以有了機會.

$^{1561}$就當向眾人行善.
向信徒一家的人.
你可以想像.
保羅在這裡.
教導在教會裡.
如何超越.
律法的要求.
成為一個緊密的群體.
如何在教會裡.
用愛.
互相款待.
用愛互相擔當.
但是自己.
有行為的表現.
有實際上.
為什麼你會這樣.
關心其他人.
有這個愛.
就是因為你信著聖靈去殺種.
就因為你已經知道.
自己是按照著.
聖靈的指引.
教導和感動.
去幫助有需要的人.
所以這裡就告訴我們.
必定會從聖靈.
去收穫.
這個永生.
收穫這個永生.
不是去換取功德.
好像一個優惠券.
而是.
按照著你內心.
已經有的新生命.
去活在.
教會的生活裡.
所以就不要灰心.
到時候會收成.
在整段的.
夾雜裡面.

$^{1601}$又是說到有一個這樣的系統.
就是要互相擔當.
各人的需要.
和尊重.
在教會裡面.
那個教導的系統.
在道理上.
受教的.
當把一切的需用.
供給私教的人.
當時保羅已經開始.
建立了一個健全的.
在似乎加拿大.
土耳其一帶的地方.
的教會.
已經成熟了.
已經開始.
有他們的運作.
這個重點就是.
要讓真理的教導.
成為恆常化.
成為一個健全的.
穩固的一個體制.
令到很多的.
真道可以.
有神恩賜的人.
呼召的人.
可以來到去.
在教會或其他地方.
來讓更多的信徒.
得到那個.
教導,牧養.
使他們能夠.
願意按照著.
神的聖靈.
來生活在當時的世界.
所以這裡六章的六節.
似乎和上下文無關.
其實是很有關連的.
我們看到.

$^{1641}$保羅是很理直氣壯的.
去勉勵這個教會.
必須要建立一個.
供給私教的人.
的一個.
需要的一個機制.
使他們能夠.
安穩地.
有系統地.
將神的.
真理來教導.
給其他的人.
我們揭開哥林多前書九章.
揭開後.
和大家看看經文.
大家平時也算是.
未必會用心留意到這些細節.
第九章的.
第一節開始.
開始談論到.
使徒他們的自由.
我不是自由的嗎?.
我不是使徒嗎?.
我不是見過我們的主耶穌嗎?.
在大馬士革的路上.
你們不是.
我在主裡面所作之工嗎?.
就在哥林多的傳道.
一年半.
假若在別人.
我不是使徒.
在你們我總是使徒.
因為你們在主裡.
我作使徒的印證.
就是那個戰書.
你們就是我的戰信.
保羅這樣來.
和哥林多教會說.
關係也很好.
保羅帶了他們信主.

$^{1681}$在那裡也有一段時間很長.
第三節.
我對那盤問我的人就是這樣分數.
難道我們沒有權.
並靠福音吃漢嗎?.
難道我們沒有權.
並取信主的姊妹為妻.
帶著一同往來.
彷彿其餘的使徒.
和主的弟兄.
並姬法一樣嗎?.
其實引述彼得有娶老婆的.
但是為什麼教皇.
是不娶老婆呢?.
又說彼得是第一任教皇.
經文就是這樣.
很清楚的.
彷彿其餘的使徒.
可能安德烈 可能馬太.
是有太太的.
不過聖經沒有記載.
主的兄弟.
耶穌基督的弟弟.
雅各 約西 西門 猶大.
他們都有太太的 有家庭的.
最後就引述了姬法.
姬法也是有家庭 有家室.
保羅就說.
即是意味著.
保羅是有這個權柄.
但是他不結婚.
他為神的緣故 獨身.
難道我們不可以嗎?.
他真的是可以很自由的.
這樣做的 我不是說過.
馬丁路德晚年.
因為太多那些.
修女 即我們叫老姑婆.
修女.
真是脫離天主教.

$^{1721}$跟了你們宗教改革.
這些先賢.
即馬丁路德 加爾文.
慈雲尼 這一幫的人.
跟了你們離開天主教.
他們.
一直都在教廷裡面生活.
來做修女.
那怎麼辦呢?.
馬丁路德是帶頭.
晚年的時候.
帶頭娶了一個修女.
即是說神父是一個結婚的.
不是一個晚年.
很孤單的 想結婚.
其實他說他自己.
你看他的著作.
他說他自己.
他自小已經是一個神父.
要做一個神父.
所以他知道自己是要獨身.
他要娶一個修女.
即是告訴其他人.
要照顧這些.
他又不是寡婦 因為還沒結婚.
但是他們真是.
因為宗教改革而這樣的下場.
記住教會以外沒有救恩.
是很嚴重的一件事.
你們作為.
宗教改革 你就帶了人走了.
你就不去照顧.
這樣說不通.
所以其實你可以想像.
是一個很美麗的圖畫.
可惜我出世太遲了.
不然我參加他的婚禮.
華丁路德 你想想.
多麼神聖 一個神父娶一個修女.
是一件多麼美麗的事.

$^{1761}$記住 是因為恩信清義.
是因為這個大原則.
是因為誰說教廷.
你們就去捍衛.
你們去專利福音.
當然不是 每個人在耶穌裡面.
就是救恩 就是得救.
他們要傳這個真理.
就要顛沛流離在歐洲.
這樣去面對很多的逼迫.
就是那個信心.
就是這個唯獨信心和基督.
所以你可以想像到.
保羅都是因為福音的緣故.
他要用他的時間.
去不同的地方 歐洲去傳道.
我想當年的歐洲.
比今天更加漂亮.
因為2000年前.
更加天然 更加人少.
當然很多東西很落後.
還是耕田 還是要靠騎雷 騎馬.
但是始終沒有那麼多人.
應該是傷害破壞沒那麼大.
但是保羅他看到的不是巴黎鐵塔.
當然當時沒有.
他看到的是很多歐洲人.
是沒有福音.
外邦人 他說他欠外邦人的債.
羅馬書一章的17節.
你可以想像到.
他看到的就是很多福音的需要.
所以他放下這個權柄.
難道第九章的第五節.
難道我們沒有權.
去娶信主的姐妹為妻嗎?.
第六節.
獨有我與巴拿巴沒有權並不作供.
就是說其實他們有權並.
因為神很清楚在使徒當中.

$^{1801}$要為我分派巴拿巴和蘇羅.
記不記得13章的使徒行傳.
他們在安提阿禁食禱告的時候.
在那個祈禱團裡面.
神就開聲了 聖靈說話了.
要為我分派巴拿巴和保羅.
去作我所差遣的供.
所以是有權並神要他們專心.
全時間去做一個使徒.
所以巴拿巴和保羅.
聽到這個呼召之後.
隨即教會就差遣他們兩個.
第一時間想起巴拿巴人有感動.
不如回我鄉下去安提阿.
對不起 去塞浦路斯的居比路.
因為巴拿巴是居比路人 塞浦路斯.
第一次的傳道走遍整個塞浦路斯島.
一路由東去到西邊 去到柏弗.
你可以想像到.
所以我們對使徒行傳熟悉.
我們看回這些經文就會有印象.
所以大家應該要由第一次宣教旅程.
去到第三次宣教旅程.
第四次就上耶路撒冷.
你應該要很熟悉那些經文.
當然最好是有時間去那些地方.
所以這裡告訴我們.
第七節 有誰當兵自被良享.
有誰栽葡萄園不吃園裡的果子.
有誰在酒樓的大廚房.
他們在做事的時候不會偷吃點心呢.
一樣 有誰牧養牛羊不吃牛羊的奶.
我說者說 豈是照人的意見.
即是說不只是一般的見識.
律法就是申命記25章.
律法不也是這樣說嗎.
這樣引咎於我.
就如摩西的律法記著說.
牛在牆上踩谷的時候.
不可以擋住牠的嘴.

$^{1841}$接著他解釋 保羅.
其實25章的申命記是沒有解釋的.
很特別的.
那裡是說怎樣對待那些弟兄.
不過突然間加了一句.
第四節 牛在牆上踩谷.
就不要擋住牠的嘴.
那段聖經很特別的.
但是保羅現在幫我們解釋.
摩西這樣說.
難道摩西說.
神掛念的是牛嗎.
不全是為我們說的嗎.
明明是說那隻牛這麼辛苦.
你都不要綁住牠的嘴.
不准牠吃東西.
要牠做完事才可以回到牛圈裡.
才給牠 才餵牠.
原來在做事的時候.
咬一下就讓牠咬一下.
的意思.
保羅說不是為那隻牛.
不是叫我們有一個牛主.
有一個惻隱的心.
不是這樣的.
很明顯這一節聖經.
上下文是不關事的.
但是保羅現在說到關於.
他的權柄.
使徒的權柄.
其中說到壽身的問題.
他就說有法律根據.
有舊約的聖經根據.
第十節.
不全是為我們說的嗎.
就是我們這些蒙召的使徒.
分明是為我們說的.
因為耕種的.
當存著子望去耕種.
打場的 即是收割的時候.

$^{1881}$打谷.
把衣服.
拿走剩下的穀物.
打場的.
也當存得良的子望去打場.
我們若把屬靈的種子.
撒在你們中間.
就是從你們收割.
奉養肉身之物.
這還算大事嗎.
若別人在你們身上.
有這個權柄.
何況我們這些使徒呢.
不過留意 說了這麼多.
因為哥林多教會.
當時對他的誤解.
對他的身份的誤解.
他說 然而我們沒有用過這權柄.
倒凡事忍受.
免得基督的福音被阻隔.
就是說要很有骨氣地.
我又說一下.
我們中穆斯年代都教我們很清楚.
很好的 因為原來真是發現.
是會有些微妙的.
別人會對教會有些誤解.
中穆斯年代已經很清楚說了.
我們去和一些家屬.
特別是那些離世的家屬.
第一句一定要說清楚.
我們幫你們辦這些安息禮拜.
是一毫子都不收的.
所有的童公義務做的.
是本於愛心.
就算你們未信耶穌.
我們都會幫你做的.
幫你做基督教的儀式等等的.
我們帶你去那些殯儀館.
講明我們不是因為有一個收回佣等等.
而是你喜歡哪一間沒有問題的.

$^{1921}$不過如果你沒有什麼意見.
比如九龍殯儀館比較熟.
我們就去而已.
完全沒有任何的.
你一定要去九龍殯儀館.
你一定要這樣.
我們一定要講得清楚.
原來很多家屬都發現.
原來他們真是會很擔心.
會問 因為他們從來沒有試過的.
他們會問多少錢呢.
原來外面的有些牧者是收錢的.
講得很清楚要給多少錢.
我沒有權去判斷.
但是教會我們就很清楚的運作.
所以從我的時間.
我們就覺得張牧這樣的教導是很好.
我們都是這樣很清楚的告誡.
當然他會奉獻給教會.
我們是會接受的.
給教會.
但是如果奉獻給個別的牧者.
我無時無刻都會收到的.
我就會全數給教會.
但是教會就必須替我.
當然替教會開一張收據和Thank you card.
給對方.
就證明你們的愛心.
我們就奉獻給教會了.
因為教會都有規定的.
牧者800元 現在升了800元.
以前500元.
800元以下牧者是可以收了.
是沒有問題的.
但是800元以上就要向教會講一聲.
就由教會去決定.
這個機制是很好的.
避免了很多的問題.
所以我們就是很清楚.
第一時間要讓他知道.

$^{1961}$我們是愛心幫助.
不是去給你收錢.
避免了他們會有一種擔心.
又不知道是不是會收很貴.
給完殯儀館又要給教會.
我們講得很清楚.
讓對方是很心裡面平安.
去接受教會的幫助.
所以很多時候他們現在.
給到一些專業的協助.
這個原來是很重要的.
原來免得基督的福音被阻隔.
所以大前提是在這裡.
保羅說其實我有權.
我連娶老婆都有權.
何況這個小小的物質的供應.
我引咎約神的話語都是這樣講的.
難道關心一隻牛.
不關心牧者使徒.
保羅說我這樣講不是要逼你們.
而是我更加讓你看到.
我們是將屬靈的種子.
塞在你們的中間.
在你們當中收養一些肉身之物.
你們都會覺得是大事.
保羅說不過.
這些權柄我們沒有用.
你可以想像到保羅.
他的真理.
他的神學.
他的智慧.
和他願意放下自己.
兩樣東西.
你想一下那個的平衡.
是很偉大的.
往往如果你有一個很大的權.
你一用那個權用到盡.
那就不好的了.
而保羅能夠放下自己.
他看到你還不能夠明白.

$^{2001}$我收起這個權.
我為著福音不要受影響.
在你們這一幫人的內心.
所以這個就是那個Gus.
真是為著基督的福音不被攔阻.
求主幫助我們.
頂子們我們都真是要學習.
你和我都真是需要捍衛福音.
為著福音不要受攔阻.
如果不是.
你想一下.
周末思原來已經扣了15分.
要停牌.
然後要坐牢.
你想一下.
基督的福音為什麼.
當然你可能一次有些特別的事.
是不關我的事的.
那另算吧.
如果說你每個月都扣了分的.
那你經常開快車.
危險駕駛.
那會不會別人看到.
基督的福音受到阻隔呢?.
比如我們有時候都需要去法庭.
陪一些人上去.
有時候尷尬就是離遠都見到法官.
都和我點頭.
因為認識 因為是教友.
那你可以想像得到.
在他每天審案.
這次周末思陪他來.
他要面對這個官司.
作為這個法官的教友.
他會怎麼想呢?.
雖然是個人的事.
但是我想大前提仍然都要告訴我們.
保羅這些事奉神的心態.
是很值得我們去效法和去欣賞.
今天的社會教我們要據理力爭.

$^{2041}$千萬不要低聲下氣.
社會教我們就是要爭到盡.
但是正經不是這樣.
保羅他是很有分寸.
很成熟地看這些事情.
再加上提摩太前書的五章.
我們就說了.
因為剛才貼上後書就是說那些不按規矩的人.
現在我們看到保羅.
他非常克制自己的那個位份權柄.
提摩太前書的五章.
我們上個星期正是說了這個寡婦制度.
供應那些寡婦.
所以年輕的寡婦.
其實就不要這麼快在那個機制裡面.
受到那個供給.
因為他還有很多的potential.
很多的空間他可以找工作.
或者他可以再嫁人.
不要這麼快有一個很大的承諾.
要起一個獨世.
那就沒有需要了.
要在神的面前這樣起世.
說我不嫁了.
我是一心到底的.
但是很快就他的情慾發動.
就是說他慢慢慢慢發現.
原來還有很多事情還沒有處理得好.
其實就不需要這麼做了.
不過留意下面.
第十六節我們上一次說完了.
信主的婦女如果家中有寡婦.
自己就當救濟他們.
不可累著教會.
好使到教會就能夠救濟那些真無依靠的寡婦.
整段就是在說.
那個教會要有一個行善的機制.
其中就是這些真是獨居無靠的寡婦.
又是年老.
真是需要有教會.

$^{2081}$當然他是信主和很敬虔.
教會要照顧.
我自己親手都經歷過.
一位老姐妹將她的遺產.
要由我去執行.
信教會 信傳道人.
不過問題是不信兒子.
兒子就是因為沒信.
所以就要由我去執行.
搞到都是很困難.
因為寫了遺囑.
律師都寫好了.
誰知道博愛的靈養倒閉了.
律師說不行 名字都沒有了.
不可以給博愛.
要轉去屯門.
屯門說不是我們的名字.
我們不要的 到時告我們怎麼辦.
大家都不要 博愛也不要 屯門也不要.
搞了差不多10年.
那些錢就繼續在我戶口裡.
怎麼辦呢.
很煩啊.
因為我要託管人.
我有責任和法律效力託管.
因為遺囑寫著我的名字.
我做執行的.
真的很麻煩.
後來要找區議員搞定.
區議員再找大律師.
大律師都搞了很久都不行.
後來區議員就說.
你不要這麼麻煩.
你們屯門收了它.
現在根本沒有靈養.
只有屯門有靈養.
屯門已經說 好吧 好吧 收吧.
搞了差不多10年才肯收那些錢.
我當然不敢放在儲蓄戶口.
有利息的 我放在支票戶口.

$^{2121}$其實意思就是.
一定要有一個機制.
教會一定要有一個機制行善.
就是要這樣安排.
讓這些年老的寡婦.
得到救濟的供應.
有這個機制.
有人就會挑戰這個機制.
我是寡婦 只是年輕.
保羅就說 要很謹慎.
因為他這一刻的心願.
可能會有變.
其實不是教會受到虧損.
是他受虧損.
是他和神之間的協定.
許願受到挑戰.
因為來日方長.
很多事情.
他有時候沒想得那麼周全.
或者慢慢沉澱下來.
可能就會有很多變數.
所以保羅就說不要了.
保羅留意第17節.
一直去到近尾那裡.
善於管理教會的長老.
這裡是教題魔態.
怎樣在已不鎖教會.
來去牧養和帶領管理.
作為一個屬靈的管理者.
這裡說.
善於管理教會的長老.
有智慧和心.
很專一去服侍神.
當以為配受家配的敬奉.
那奴婦全道教導人的.
更當如此.
接著又引這裡.
哥林多前書都引過.
又是他引 保羅.
因為經常說牛在牆上踩谷的時候.

$^{2161}$不可擋住他的嘴.
再引多一節利未記.
又說工人得工假時應當的.
利未記十九章十三節.
所以這裡就是兩節聖經.
再一次凸顯了.
教會需要有完善的管理.
教會需要有commit的人.
一心要為神去奉獻他的智慧.
他的管理在教會裡面.
讓那些機制越來越完善.
這些是教會的福氣.
還有人一心為神的道.
放下自己的喜好.
可能是生活習慣.
來奴婦去傳道.
讓教會能夠在健康裡面去成長.
所以相反了.
十九節控告長老的情子.
就是那些black and white.
非有兩三個見證就不要收.
所以不要聽一個人說.
哪個哪個不好.
你應該是雙方面去尋求答案.
所以犯罪的人.
即是有證據確鑿了.
就當在眾人面前去責備.
叫其餘的人都要懼怕.
我在神和基督耶穌並無揀選的天使.
的面前祝福你.
哪些被無揀選的天使呢.
加伯列 米加勒.
不知道了.
要遵守這些說話.
不可以傳成見 管理者.
不可以戴著油色眼鏡.
不可以聚眾的搞小圈子.
不可以有自己的黨派.
因為你有黨派就會有成見的了.
行事都不可以有偏心.

$^{2201}$不要喜歡那個人就側側膊.
不喜歡那個人就置他於死地.
這樣就偏心了.
給人行按守的禮 不可急促.
不要在別人的罪上有份.
要保守自己清潔.
諸如此類的情況.
是一個完善的教會機制的各方面.
人是重要之中的重要.
還有那個機制如果能夠完善.
就會帶來教會是一個很美善的德行的見證.
所以教會真的需要做這些辭位和幫助.
但是要很小心運用教會信徒的獻金.
在這些美善的德行上.
在神的面前留有神的稱讚和美好的明星.
教會在很多未信仰的人的內心.
都是有這樣的美德的.
比如我們還是這一代.
50年代 60年代是最艱難的時候.
香港因為國共內戰之後的大遷徙.
100萬的難民.
那個時候政府根本沒有什麼能力.
民間基督教就帶頭去做.
很快就看到就做.
其實大家都艱難的.
大家都在那個最困難的時候.
西教士就派奶粉 派麵包之類的.
真的是一個美德.
西教士做了 中國人都沒有.
當時其實有些聖公會.
是很大的宗教會.
但是完全看不到這樣的.
只是西教士在外國來的.
特別是美南浸信會的西教士.
就是做這些的扶貧.
當然也都因為當時1949年的時候.
所有中國的宣教士.
在中國的宣教士 全世界的.
要立即撤走.
不可以再在裡面 因為開始那個逼迫.

$^{2241}$所以當所有宣教士來到香港.
就變成香港成為一個福音的基地.
但是他們看到的就是立刻辦醫院.
辦其它的教育等等.
所以成為了70年代.
教會的大復興的一個開始.
所以我們都需要透過歷史不斷去反思.
現在這個時代人需要什麼呢?.
什麼是教會的美德 美善.
可以來幫助呢?.
所以這個教會永遠是在一個不穩定裡面.
去幫助人是應該的.
就不應該是越來越穩定.
變成一個俱樂部的會所.
大家在裡面不需要有外人的.
只是我們一幫人 那個已經不是教會了.
整個理念都不對了.
是應該不斷有些很頑皮的小朋友.
怎麼教他都不聽的了.
每堂都要給你倒彈的.
教會就應該是這樣的.
就不是趕走這些.
剩下的全部很乖的.
自己的子女就行了.
教會永遠不應該是這樣的.
當然不是鼓勵他們去倒彈.
而是你在幫助他們的過程.
他們將來成長.
他們對教會又會有不同的看法.
所以這個都是一個.
透過今天講的.
不知不覺的 透過第三章.
又講了一些教會的運作的體制.
其實是很寶貴的.
我自己這麼多年的經驗.
跟很多不同的前輩.
是學了很多東西的.
我做不做傳道呢?.
我根本不會涉獵這些的範疇.
我的生命不會這麼豐富.

$^{2281}$因為對人其實是非常高的藝術.
是很難的.
但是會使自己更加的成長.
所以其實是很寶貴的.
所以記得不是機制一被動搖.
我們就灰心.
應該說我們怎麼更加修補那個機制.
不要讓那個機制失效之後.
大家就無所適從.
教會就出現了一個更嚴重的問題.
所以只要是需要用愛心去勉勵.
勸勉那個違反機制的人就行了.
就不要搞亂那個機制.
今天就差不多去到這裡.
我下個星期開始就和大家去看.
這幾個的神學.
關乎到末世論的.
我們一起禱告.
多謝主.
今天晚上我們又用了一個半小時.
去將劫後的第三章.
來和大家彼此去學習.
我們很懇求主給我們對第六,第十二節.
基督在教會裡面的權威,權柄.
教會應該怎麼去運用.
怎麼有智慧的,有愛心的來運用.
其實重點是為了那個機制完善.
是要有美德.
是要讓教會能夠將福音.
透過很具體的,實際的幫助.
可能是一個麵包.
可能是一個愛心的提攜,說話等等.
這個都需要神給我們.
無論是管理者或者弟兄姊妹.
都要有這一份的成熟.
求主保守.
不要讓教會充斥了很多的流言.
很多的是非.
而是讓大家都能夠同心.
大家放下自己以為有的權柄.

$^{2321}$我們來學習基督.
默默地來服侍.
真的將人帶進救恩裡面.
不要讓基督的福音受到攔阻.
是我們每一個弟兄姊妹都需要學習的.
我們仰望你.
下個星期我們會開始.
看一些關於穆斯林神學的論述.
幫助我們能夠明白裡面的意思.
稍微了解一下.
以致我們對於主的再來.
是有一個渴慕,等候的心.
多謝你聽我們的感恩禱告.
奉耶穌基督的名.
阿們.
\newpage



\section{}
\label{sec:I4c89OnZC6M}
\textbf{2024-08-29 周四晚查經班《淺談末世論(1)》 講員:曾裔貴牧師}
\newline
\newline
連結: \href{https://youtube.com/watch?v=I4c89OnZC6M}{\texttt{https://youtube.com/watch?v=I4c89OnZC6M}} ~~~~ 語音日期: 2024-08-30
\newline
\newline
\hyperref[sec:IJNXTQniUCA]{\small{< < < PREV SERMON < < <}}
~
\hyperref[sec:index_chronic]{\small{[返順時目]}}
~
\hyperref[sec:index_scriptual]{\small{[返順卷目]}}
~
\hyperref[sec:a8KnWbij7Uk]{\small{> > > NEXT SERMON > > >}}
\newline
\newline
$^{1}$跟大家一起祈禱.
感恩天父讓我們在晚上.
再一次能夠放下工作.
來到你的面前去學習.
去聆聽聖經的真理.
多謝主讓我們能夠透過《貼撒羅尼加前後書》.
我們看到一個很重要的題目.
就是仕途保羅在他第二次宣教旅程.
向一個初信的教會.
來強調主耶穌的再來.
相信這是一個在保羅的內心.
和仕途們的內心裡面.
一個相當重要的題目.
或者是他們信心的確據.
主耶穌回到天上之後.
必定會再回來.
所以我們又看到.
真的有不同的新約聖經.
其他的作者.
都不約而同地去講道關於主再來.
所以我們都能夠有幾堂的時間.
去淺談一下活世論這個課題.
一個都是和我們息息相關的題目.
讓我們能夠明白.
後來我們會提到一些神學的不同的派別.
怎樣去看這個千禧年.
但願主都讓我們能夠掌握明白.
和我們有一個心之確信.
等候主的再回來.
今天晚上我們在你面前等候禱告.
奉耶穌基督的名.
阿們.
是 阿力斯.
這樣又有新的朋友加入我們的茶經班.
我們大部份因為很多經文要揭.
所以每一次都反正都開了個聚會.
不如就長一點.
所以一個半小時.
我又不怕站這麼久.
一個半小時都是能夠比較可以鬆動一點.

$^{41}$有系統的將一些經文比較深入的講解.
所以大家如果想去茶經班.
預了一個半小時.
如果你是上聽的.
你就可以隨便按照你那個時間.
你方便的時間去聽.
也是一個不錯的一個溫習.
甚至乎可以整個系列你去重複聽都可以.
我們已經有系統的講了約翰福音.
又曾經講過一些聖經人物.
但後來停了一段時間就沒有再講.
遲些我們有機會我們會再講都可以的.
可以大家去看不同的書卷.
你的需要.
因為我們大概三堂就會完這個淺談末細論.
所以剛才我剛剛在群組裡都給了幾個提議.
看大家來選擇.
我想也是一個投票.
除非有些姐妹你很想某一卷書.
你很強烈想.
你可以私下告訴我.
讓我了解一下.
不然大家按那個投票.
看誰多一些的我們可以先去查考.
舊約我就給了兩卷.
新約就有希伯來書.
因為我都知道後來有很多人加入了羅馬書.
但因為我們最起初就是查羅馬書.
那個時候還在嘉欣樓.
都查了很久.
但是因為很多人沒有聽過.
還有就是那個時候我們沒有那麼先進錄音.
所以就有些浪費了.
所以我就想一下會不會有機會.
反正現在很多弟兄姊妹都沒有上過羅馬書.
記住是有系統地.
由第一章讀到第十六章.
你必須要知道整卷羅馬書.
到底寫作的動機.
寫作的原因.

$^{81}$以及保羅怎樣有系統地將那個救恩講解.
這些必須要拿捏得很準確.
很清楚的.
所以我又有感動想再講回羅馬書.
當然剛才打了.
我反而更加感動就是很想講一下神學.
講一下馬丁路德時期的宗教改革.
有很多東西講的.
包括那個來龍去脈.
為什麼會有改革.
因為星期二我帶了小朋友去長洲.
有一個小朋友有很多問題.
他在讀一年級而已.
但他已經在問很多神學的問題.
他問了一個問題在船上.
他說天主教偉大一些還是基督教偉大一些.
他已經問這個問題.
一年級.
小學一年級.
我就講解給他聽.
有很重要的日子,年份.
他聽完.
當然也不知道我說什麼.
因為一年級.
但是他解決了他的疑團.
他就知道應該來說基督教偉大一些.
給他一個初步的印象.
所以突然間剛才打字的時候.
又想起.
如果跟大家去看看.
這個宗教改革時期的神學.
和當時的思潮.
和為什麼要改革.
改革帶來什麼正面的影響.
改革帶來什麼負面影響.
這些可能都要知道.
因為我們今天是500年之後的宗教改革的製成品.
我們現在有分浸信會.
有分宣導會.
播道會.

$^{121}$或者其他的獨立堂會.
或者是高派的聖公會.
這些我們可能也要知道.
和宗教改革完全有關係的.
導致我們今天宗派的臨立.
不會再有一個統一的基督教.
到底是好還是不好呢?.
這些在宗教改革裡面.
其實我們都可以了解的.
不過我不是遊說.
來到大家一定要上這個神學課程.
不過可以給大家幾個選擇.
在這個末世論提過了.
因為《帖撒羅尼迦前後書》.
特別第四章全章.
是講到關於主再來的次序和原因.
是因為當時有帖撒羅尼迦教會的弟兄姊妹.
他們的親人離世.
所以保羅為著要牧養安慰他們.
以致他們不會有憂傷.
甚至失去那個救恩的盼望.
就告訴他們.
不是要滿足他們的好奇心.
而是告訴他們.
那個人死去之後.
他要去到哪裡呢?.
還有還沒死的人在主裡面的.
又是怎樣的復活呢?.
就講得相當詳細的.
所以勾起我們大家.
不如我們在神學院.
都有機會涉獵關於末世論的這個課題.
也有好幾個不同的主張的派別.
有前遷 有後遷.
有一個叫做無遷.
所以只有這三個.
不會有第四個的了.
如果你是第四個.
根本不會有主再來.
那你就不是基督教了.

$^{161}$你是異端來的.
你就是一定會沉淪的.
所以主再來是必須要明白的基本真理.
但是三個派別裡面.
都共同的 一致的.
是不會有任何的含糊的.
都相信主耶穌回來的.
不過問題就出在這裡了.
什麼時候回來.
到底是已經回來了.
還是災難之前回來.
還是災難過了之後回來.
這個我們可以仔細的去分析一下.
給大家有一個小小的概念.
大家純粹去聽一下都可以.
因為其實不涉及到你信耶穌得救的問題.
不過如果你不信耶穌回來.
這個就涉及到你的得救問題了.
這個大家要明白.
所以三派對於主再來.
他們都是共同相信.
確信耶穌一定會回來.
不過是什麼時候回來就有差別了.
這個差別也挺大的.
因為涉及了一千年.
千禧年.
千禧就是一千年.
Millennium.
就是一個拉丁文.
一千年的意思.
所以你可以說相差一千年.
如果沒有千禧年的話.
就不相信耶穌會首先回來.
接著就開始一個千禧年.
這個就叫做無千禧年.
在這樣的情況下.
我們就要先揭幾段聖經.
是相當具代表性.
關乎到主再來.
給大家有這個概念.

$^{201}$就是福音書三卷福音書.
都不約而同的說了.
馬太福音的24章.
這些就沒有懸念了.
三個派別都是相信主耶穌的應許.
唯一最大的差別就是.
有千禧年還是沒有千禧年.
是災難之前的千禧年.
還是災難之後的千禧年.
這個就是再細分的不同派別的意思.
暫時來說.
見到還沒有改他的神學立場.
如果我讀的時候.
仍然是相信前千禧年的.
即是有千禧年的.
這個是我讀神學的時候.
因為是頗重要的他們的信仰宣言.
24章的馬太福音.
你就可以知道這個叫做.
鑒欖山的寓言.
因為當時主耶穌還沒有釘十字架之前.
他就和門徒準備要入城.
很快就是逾越節.
是一個很大的節期.
即是說將燈結彩.
聖殿會有很多準備.
很多人準備來守這個很大的逾越節.
當時耶穌基督和門徒.
就在鑒欖山看著對面的聖殿.
當時門徒就很奇怪的問了一個問題.
因為耶穌這樣說.
第一節開始.
我給大家有少少的概念.
因為我怕有些人都不是很熟.
全章就是講這個聖殿.
耶穌就在鑒欖山和十二個門徒.
講關於將來的末世.
耶穌出了聖殿.
正走的時候.
門徒進前來.

$^{241}$把電雨指給他看.
耶穌對他們說.
把電雨即是說.
你看看我們猶太人這個電雨多麼的輝煌.
多麼的平安.
多麼的我們猶太人感到賴以自豪的地方.
神的殿 神的家.
百姓在那裡來敬拜.
不過耶穌就講了一句話.
第二節.
耶穌對他們說.
你們不是看見這電雨嗎.
我實在告訴你們.
將來 很快的將來就是主後七十年.
是應驗了.
就是耶穌死了之後.
其實到保羅他們都死了之後.
到主後七十年.
就是當時的羅馬皇帝提多將軍.
當然提多將軍就不是皇帝.
是羅馬這個.
我忘記了名字了.
這個其中的將軍提多將軍.
就帶兵去踐踏.
甚至把整個聖殿拆毀.
然後這裡說有一塊石頭留在石頭上.
不被拆毀.
因為他們將整個聖殿燒了.
那些金箔就溶在那些石縫裡.
所以羅馬人就挖爛了那些城牆.
和那些牆身.
將那些石頭撬開.
拿那些金箔這樣拿走.
所以整個聖殿是夷為平地的.
這個是主後七十年的時間.
所以第三節.
耶穌在橄欖山上坐著.
門徒暗暗的來說.
請告訴我們什麼時候會有這些事情呢?.
當然就是講到聖殿被拆毀.

$^{281}$接著門徒又問了一個很精準的問題.
其實你問一下關於聖殿就夠了.
但不是.
你降臨和世界的末了.
有什麼預兆呢?.
很特別.
耶穌明明坐在他們當中.
為什麼又突然間講到你的降臨呢?.
和世界的末後呢?.
為什麼突然間會問這樣的預兆?.
一個這樣特大的神蹟呢?.
為什麼會相信耶穌會再回來呢?.
很奇怪.
很矛盾的.
因為當耶穌就快過幾天.
他釘十字架之後.
那些門徒就不再相信了.
就害怕了.
不敢承認耶穌是能夠復活的那位救主.
那你可想而知.
他不復活.
不升天.
又何來會降臨呢?.
何來會有那個預兆呢?.
世界的末日的預兆呢?.
他們理解耶穌的降臨和世界的結束.
是一個拉上關係的緊扣的.
所以真的很矛盾.
為什麼過幾天你們又會失去完全的信心呢?.
耶穌在的時候又會問這麼偉大的問題呢?.
第四節就展開了解答的序幕.
你們要謹慎.
免得有人會迷惑你們.
因為將來有好些的人冒我的名來.
說我就是基督了.
所以這是一個很指標性的末世論.
我們要留心的.
就是似乎基督教不是一直偏向去到好的.
而是敵基督,假基督.
是會在那個末後的日子出現的.

$^{321}$有好幾段經文其實證實了這件事.
譬如我們上個星期提過的.
貼上來後書的第二章.
就提到那個大罪人.
那個沉淪之子.
那個照著撒旦的運動.
去走出於不義的真理的人.
是會迷惑那些不肯接受真理的人.
二章一至十節那班人.
到第二班人就是在主裡面.
有信心,堅固的信仰的人.
就不受迷惑的.
這是我們上兩個星期提過的.
這裡耶穌已經很早很早告訴這十二個門徒.
要小心.
因為將來會有奉主名而來的假基督.
來去迷惑許多人.
所以假的基督,假的信仰.
是會在心智上迷惑.
在神學上迷惑的.
是令你對一些真理.
慢慢就似是而非的.
慢慢越來越模糊的.
越來越沒有了一個界線的.
越來越不能夠分黑和白的.
這個就是假基督了.
我們要很小心這件事.
很快我們因為要看的都相當多的經文.
所以我不可以逐段和大家仔細去看.
我先點出了那個重點.
就是三派.
前遷,後遷,無遷.
記住這三個概念.
前遷就是災難前的千禧年.
應該說災難開始的千禧年.
後遷就是過了那些叫做Tribulation.
特大的災難之後.
才會有一個千禧年的.
過了千禧年耶穌回來的.
接著無千禧年就最簡單的.

$^{361}$就是說耶穌升天之後.
就是教會的時代.
耶穌第二次回來就是世界末日.
不會有千禧年.
這個叫做無千禧年.
我們揭一揭後面的四十多節.
或者我們看到第36節打後.
這一段和大家最後看看這裡.
24章的《馬太福音》.
提到主再來的確實的應許.
第36節.
但那日子那時辰.
沒有人知道.
連天上的使者也不知道.
子也不知道.
唯獨父知道就是說到主第二次回來的時候.
羅亞的日子怎樣.
人子降臨也要怎樣.
羅亞是什麼日子.
是怎樣的日子.
羅亞是預備方舟的日子.
是預備有一個大審判的日子.
是預備人子要順進入那個方舟.
就可以得到逃避洪水的毀滅.
這個就是羅亞的日子.
但是這裡告訴我們.
第38節果然提到準備.
逃避一個大災難的審判的信靠.
就是洪水的方舟.
成為一個唯一的拯救.
當洪水以前的日子.
人照常吃喝嫁娶.
洪水以前.
即是說馬土薩拉.
他是969歲.
是他2月18日.
我忘記了是2月18日還是27日.
洪水就來了.
應該是.
先幫大家揭曉.

$^{401}$馬大夫.
倉西記.
因為講得太快.
倉西記的第7章.
洪水的時間.
(記載中).
沒錯了.
7章的11節.
當羅亞600歲.
2月17日那一天.
其實就是馬土薩拉死的那天.
969歲.
你回去算一算那個家譜.
羅亞600歲.
真是剛好馬土薩拉969歲.
你就說他帶漏了他的太公.
沒有進方舟.
他可以是正月初一死的.
可以是2月死的.
是2月17日才有洪水來.
才有大淵的泉源打開.
所以有些人批評.
羅亞你這樣.
帶了自己三個兒子三個媳婦和太太.
忘記了帶馬土薩拉進方舟.
其實就不是.
有足夠的時間.
馬土薩拉是壽終正寢離世.
記住.
新約的作者告訴我們.
其實是耶穌告訴我們.
當時的人.
在洪水以前.
他們是正常的吃喝嫁娶.
正常.
生活是很悠哉遊哉.
不會擔心有任何的地球的改變.
氣候的改變.
不會擔心有大海嘯.
不會擔心有地震等等.

$^{441}$我們今天太多了.
所以我們已經很容易知道.
幾現代的科技都控制不了地球.
天然的災難.
我們經過了2000年.
我們終於知道.
多厲害的科學家.
都只能夠預測.
但不能夠防止.
也不能夠逃避.
所以你可以想像到.
洪水以前.
是應該沒有地球上任何的大規模的災難.
這才是一個恐怖.
因為從來沒有試過.
所以主耶穌說.
他的降臨.
和洪水未來之前的日子.
羅亞的日子.
是完全有吻合的.
其中最關鍵的就是.
不知不覺.
沒有人會預測到.
因為從來都沒有地震.
從來都沒有海嘯.
何來預測呢.
現在你想想海洋.
和很多地上的地震儀.
擺了很多不同的地方.
來偵測.
所以有少許風吹草動.
都偵測到.
甚至那個級數都可以計算到.
但是那個時代是沒有的.
用一個很簡單的生活的常態.
就是神父與人.
有的生老病死.
在這個地球上活的.
吃喝嫁娶.
正常的生活的狀態.

$^{481}$羅亞的時候.
已經有很大的秩序.
不是說.
好像魚無飲血.
住在山洞.
打獵.
吃生的東西.
應該不會的.
吃喝嫁娶.
即是開始有文化.
開始有秩序.
開始有很多的排場.
開始有很多的講究.
甚至當時的物質的豐裕.
那個地步.
根本不需要擔心.
夠不夠訂金.
夠不夠禮金.
夠不夠什麼.
我想不用那些東西了.
因為當時的人類很少.
和地球的豐富的資源.
也是相當多.
記住是未墮落之前.
你可以想像得到.
我所指的未墮落之前.
就不是當下那個墮落.
是未有大規模的.
毀滅地球的那個事件之前.
所以直到羅亞進方舟那天.
不知不覺洪水來了.
所以出現了兩次的不知不覺.
就把他們全都衝去.
他們是整個地球的人.
只有八個人在方舟裡面.
留種給後世.
所以羅亞.
當然他們都是帶著罪性.
所以羅亞出方舟之後.
聖經第一次.

$^{521}$第一件事記載他們開始有犯罪.
也開始仍然需要那個救恩.
來拯救這個家族.
接著下面第40節.
那時兩個人在田裡.
娶他一個撇下一個.
兩個女人推磨.
娶他一個撇下一個.
就是說在田裡.
在家裡面推磨.
在田裡面可能是男丁.
去耕種.
女士推磨.
就是打鞠.
也是娶他一個撇下一個.
誰被娶去.
誰被撇下呢.
相信主耶穌在這裡.
都很清楚讓門徒知道.
因為是跟門徒說.
因為福音在主耶穌的時候已經開始傳.
相信肯定就是相信有審判.
相信有永生的人.
就是會被娶去.
被撇下的.
大家都是這樣生活.
基督徒也是這樣生活.
非基督徒也是這樣生活.
但是一個被娶去一個被撇下.
而且是正常的勞動.
不是說懶惰的就被撇下.
勤奮的就被娶去.
也不是.
兩個人都是推磨.
兩個人都是在田裡耕種.
一個娶去一個撇下.
所以這樣去分.
是生命的問題.
不是說你勤力與否的問題.
也不是我們去換取救恩.

$^{561}$基督教從來不是一個換取.
你儲夠優惠券.
就給你印花.
每次吸一吸.
回來教會吸一個.
回來茶經班就吸兩個.
儲夠積分.
你就可以上天堂了.
基督教完全不是這樣.
基督教從開始告訴我們.
沒有人可以靠自己去得救.
我們得救是本乎恩.
恩著信.
不是出於自己.
這個很重要.
不是出於我們這個人的自己.
連我們想救恩的方法.
都不是我們自己.
所以這裡告訴我們.
因為一直說到下面.
就去到最重要的44節.
所以你們要警醒.
因為不知道你們的主.
是哪一天來到.
第一個問題就要問.
我到底是被撇下還是被取去呢.
你要自己計算.
什麼令你有一個信心.
你是被取去呢.
唯一的信心就是.
羅馬書所說的.
為義人恩信得生.
這個二章的經文.
很重要的.
所以你可以想像得到.
被撇下只有一個原因.
不是他不勤力.
不是他不是好人.
而是他在神眼中是一個罪人.
不可以靠自己得救.

$^{601}$所以就是被取去.
和撇下的唯一原因.
所以主耶穌說.
你們如果認定自己是跟隨我.
在這個世界是一個小群的一班人.
你們將來被取去.
所以要警醒保守自己.
在神的救恩裡面.
這樣就會被取去了.
所以我們一直都說過.
一次得救.
一次相信.
一直相信.
就必定得救.
記住.
不是一次得救.
永遠得救.
這樣就很容易被人混淆了.
哦 我信了耶穌.
我就不用回朕禮班了.
我就不用初信班了.
我一信了耶穌.
我就不用回教會了.
我就永遠得救了.
是.
你一次相信.
你堅定的相信.
你是可以永遠得救的.
但問題是什麼去支援.
你一直的相信呢?.
就是你必須活在神的揀選裡面.
所以.
回查經班.
回崇拜.
那個不是去緩救恩.
而是讓你更加知道你是屬神的兒女.
這樣的意思.
所以我們在真道上要同歸於一.
所以你要留意一下.
耶穌說你們要警醒.

$^{641}$跟這十二個人說.
因為不知道你們的主.
是哪一天來到.
如果知道哪一天又會有賊來.
他必定會警醒.
不容人去挖透房屋.
這是你們所知道的.
也是當時的.
大家一定會很容易理解.
當時那些賊是真的會挖透了那些灰牆.
很容易就會.
不同於今天的防盜系統.
那些警鐘就已經一按就立刻響得很厲害.
等等.
我正在看的是《馬太福音》的24章.
所以他們知道.
就是說重點耶穌說的這個例子.
家主的警醒.
就是等於門徒.
你既然確信耶穌基督會再回來.
在末日.
而且上文一直都在說.
是會有假的基督的時候.
就更加要警醒.
要來到對自己所信的.
是要有一個明白.
和持守.
所以.
唯一可以做的.
就不單純去益測.
去計算.
耶穌什麼時候回來.
而是我們心靈裡面有一個準備.
所以第44節就是那個結論了.
所以你們也要預備.
因為你們想不到的時候.
人子就來了.
這一節.
相信大家都可能要記一記它.
背一背它.

$^{681}$24章的44節.
是關鍵的一句.
因為耶穌基督很清楚說了.
沒有人可以想得到祂什麼時候來.
但是耶穌基督就必定要來.
我們能夠做.
門徒在地上能夠做的.
就是我們必須準備我們自己.
常在這個信仰的裡面.
使我們能夠成為被神揀選的一個人.
所以這裡就是24章的44節.
其他我不夠詳細時間.
和大家再仔細看每一段.
我們來看看約翰福音.
約翰福音是另外一段經文.
講到主的再回來.
往往這段經文.
我們通常在安息禮拜會讀.
很特別的.
有可能是對在憂傷裡面的家庭.
一個安慰.
這個沒有錯的.
不過整個上下文.
14章其實就是流防的講論.
最後晚餐在這個喀西馬利園.
完了.
耶穌基督回到流防.
對不起.
應該是在流防吃完雨節晚餐.
接著就去喀西馬利園.
所以在去喀西馬利園之前.
耶穌和11個門徒.
因為猶大已經出去了.
找祭司長他們來捉拿耶穌.
所以11個門徒.
陪伴著耶穌.
在最後晚餐.
其實是非常的落魄.
憂傷和疲倦的.
但是耶穌也知道.

$^{721}$這些門徒他們的內心.
軟弱.
招架不了前面來的一個極大的黑暗.
那個勢力.
那個逼迫.
那個對耶穌那種無情的指控和定罪等等.
所以門徒是不能夠招架的.
但是其實最應該憂傷的就是耶穌.
在這個時候.
來安慰門徒.
他們的內心.
第一節.
14章的第一節.
所以嚴格來說.
不是在安息禮拜讀的.
應該是在一個我們對於前途茫茫的時候.
的信徒的時候.
應該要讀的經文.
例如現在有大逼迫來了.
我們要趕緊離開了.
這些經文就真的很適合了.
你們心裡不要憂愁.
你們信神也當信我.
在我父的家裡有許多住處.
為什麼要說父的家呢?.
因為13章已經開始說.
你們是我的朋友.
天父就是我的神.
就是我的天父.
猶太人.
基本上任何一個猶太人.
絕對不敢稱神為父的.
現在耶穌在約翰福音.
一而再稱神為他的父神.
所以我們今天叫父神.
是非常正確的.
我們可以親力親為地稱呼我們的造物主.
做父神.
因為耶穌告訴我們.
祂和神的關係就是這樣的.

$^{761}$在我父的家裡就是永生的天堂.
有許多的住處.
如果沒有.
我就早已告訴你們了.
有許多的住處.
不過要再確認.
是真的有許多的住處.
如果沒有.
早就告訴你們了.
就是說其實真的有許多的住處.
其實當時只有11個門徒.
他們所理解的多.
有多大的概念呢?.
我去原是為你們預備地方去.
我若去為你們預備地方.
就必再來接你們到我那裡去.
我在那裡.
叫你們也在那裡.
門徒完全.
一來在憂愁裡.
完全不理解耶穌說的去.
去是否想逃走.
離開耶路撒冷呢?.
去到加利利?.
或者再去遠一點?.
還是什麼去呢?.
但是耶穌知道自己的離世.
13章已經開始說了.
13章的第一節.
逾越節以前.
耶穌知道自己離世歸附的時候到.
他既然愛世間屬自己的人.
就愛他們到底.
接著就幫他們洗腳.
所以.
作者約翰很清楚的告訴我們.
整個敘述.
一開始已經很清楚.
耶穌知道自己離世歸附的時候.
在作者的.

$^{801}$當然也是時候很明白.
因為耶穌已經死了.
他才寫這封信的.
他已經是很年紀大的門徒.
約翰.
他現在回顧回.
記住作者是有一個.
很清楚的神學的立場.
耶穌是離世歸附.
不是死於非命.
不是死於殘酷的人.
祭司長的手.
不是死於被出賣的猶大的手.
耶穌知道自己是離世歸附.
就是說.
他是要為世人拯救.
而上十字架.
如果不是的話.
沒有人可以來.
抓到他.
來將他置於死地.
接著呢.
耶穌在這裡再繼續這樣說.
他說我去.
是為你們預備地方.
耶穌基督的離世.
原來是去一個更美的地方.
當他完成了救恩之後.
他就可以帶著.
這些有他生命的門徒.
進入耶穌本來的住處.
就是天家.
就是天堂.
所以這是勝利的海船的宣告.
不是一個.
沒辦法.
我就快要被人出賣了.
我也不想的.
但是.
必須要這樣做.

$^{841}$其實就不是.
在耶穌這個時候.
他那種信心.
那種知道自己的身份.
和自己的使命.
在這個世界.
就是為了這一刻.
就是為了要上十字架.
但是他的內心是很清楚.
我去是為你們預備地方.
接著我必定會再來.
接你們到我那裡去.
是耶穌的住處.
是與耶穌永遠同在的地方.
所以當然安息禮拜絕對適用.
絕對很大的安慰.
就是離開世界的人的門徒.
我們的親人.
他就是與主同在.
但是這段聖經.
嚴格來說的上下文.
是流放的.
最後晚餐的.
離別的感言.
所以.
約翰福音也告訴我們.
主耶穌去的是預備地方.
我們要有這個眼界.
我們要知道比世人知道更多.
世人知道耶穌在十字架.
釘十字架.
可能他們也會知道.
耶穌是死在十字架.
不過他們不相信.
耶穌是他們的救主.
這就是一般人對耶穌.
對基督教的理解.
但是我們要再深層一些.
我們知道耶穌釘十字架.
與我們有關.

$^{881}$因為我們已經認順了耶穌.
但是我們要從另一個更深層的角度看.
就是原來耶穌.
回去預備地方.
原來他的死.
他的復活.
是回去預備地方.
預備好之後.
他就會再來接我們.
所以這是一個很寶貴的迎接.
團聚的一個迎接.
所以我們應該要帶著.
沒有一個基督徒是想快點死的.
除了已經患了重病絕症.
當然想快點死.
但是正正常常.
沒有一個基督徒想.
我明天就遜尿就好了.
很少人會這樣想的.
某程度上也要輔導一下.
什麼原因的.
但是這裡告訴我們.
原來我們要善無這個重聚.
原來耶穌說我預備好之後.
我必定會再來.
去接每一個門徒.
就是說約翰和其他的使徒.
日後所帶順主的門徒.
就是那些子子孫孫.
要渴無這個重聚.
就是耶穌預備好地方之後.
就會再來去接我們.
去到那裡.
叫你們也在那裡.
這個就是想像得到.
一至第三節是一個最重要的安慰.
就是對當時在極之哀傷的門徒.
又不知道耶穌到底想去哪裡.
有沒有說得很清楚的時候.
而耶穌在這裡就說了.

$^{921}$去一個地方.
就是父的家.
就是天家永生神的住處.
那裡預備地方.
這個就是約翰福音提到的主再來.
留意喔.
全部是耶穌親自說的.
好了 接著我們又再看.
我都是已經.
因為太多經文.
我只能夠選擇性的幫大家去看一看.
關於主再來這個課題的.
希伯來書的九章.
二十八節.
希伯來書.
希望大家對主再來.
對一個天家.
有多一層的憧憬.
就是耶穌基督的住處.
天父的住處.
是永恆的 是永生的.
是沒有任何的.
我不知道是那個.
擁擠.
因為根本是足夠地方.
給所有基督徒住的.
不會覺得你要住tong 房.
你會很擠迫.
應該沒有這個問題的.
也都不是我們只是一個靈體.
似乎聖經不是這樣理解.
但是.
剛才提到的約翰福音.
就是一個歡喜的重聚.
我們要有這個心態.
要培養這個心境.
渴慕這個重聚.
我想移民的人.
就會有這個心靈.
心情.

$^{961}$就是會再見到自己所愛的人.
那種心情.
九章的希伯來書.
本來不是說主再來.
但是不知道為什麼.
突然之間說了一件事.
就是說到主耶穌的.
國族的.
主教的.
就是說到主耶穌的.
那個贖罪祭.
是永恆有功效.
是會改變了我們的生命.
這段大家應該相當不熟悉的了.
比較少人會說的.
九章23到28節.
我們要用一點時間.
去解解這段聖經.
但是和主再來是相當息息相關的.
23節.
因為上文是在說.
這個摩西的憲制.
在那個贖罪日.
當然不是摩西.
是摩西頒佈了之後.
是由大祭司亞倫.
和他的子孫.
必須是要利美茲派裡面的.
亞倫的後人才可以做大祭司.
大祭司才可以進入聖殿.
每年的7月初十.
猶太人的節期.
7月初十.
就是那個贖罪日.
他才可以進入至聖所裡面贖罪.
平日的日子.
只能夠在那個.
外院的同祭壇.
去贖罪.
但是一年有一次.

$^{1001}$他們要進入至聖所.
當然每天他們都要進入聖所.
聖所和至聖所.
是有一個門簾遮隔著的.
他們不可以進入至聖所.
一年只可以進入一次.
所以在聖所.
他們每天都要進入.
就是要整理油燈.
更換陳設餅.
還有燒香.
就這麼多了.
其他時間.
他們就不准進入至聖所.
所以希伯來書.
就提到主耶穌.
就是新約的宗寶.
祂就是那個完美的祭僧.
所以祂的獻祭.
祂的贖罪.
就一次.
就永遠有功效的了.
其實是想說什麼呢?.
是想安慰.
牧養當時信主的猶太人.
因為當時希伯來書.
還是有聖殿的獻祭.
希伯來書的作者說.
不可以再去到聖殿獻祭.
那些祭僧是救不了你的了.
因為真正的祭僧就是耶穌.
祂已經成為了永恆的贖罪祭.
不需要再靠著牛羊的血了.
所以要有一個.
這樣的過渡的信心.
是很不容易的.
你可以想像得到那種情義結.
對一個猶太人來說.
是很難的.
再加上就是.

$^{1041}$那些叔伯兄弟會逼迫你的.
你在這麼大的節日.
你不去獻祭.
你說我在家裡憑信心信耶穌就行了.
那些叔伯兄弟怎麼看你呢.
你在祠堂.
現在要很重要的打召.
要去燒香.
你說我不燒的.
你已經被人罵了.
逼迫你了.
何況現在你不獻祭.
贖罪日你不獻祭.
你就要面對神的永恆審判.
你不靠大祭司靠誰呢.
那個小朋友說.
不用了.
約翰伯伯說.
我信耶穌就行了.
我洗乾淨了我的罪.
我就可以憑信心得到神的拯救了.
你說這番話自己都會臉紅.
都沒有確立到的信心.
所以希伯來書.
整個寫作動機就是要勉勵.
希伯來的猶太人.
他們必須要將舊有的體系.
舊有的觀念.
轉變過渡去到新的救贖的系統.
就是耶穌基督.
是人類唯一的救主.
對於我們外邦人是沒有這個矛盾.
要麼信.
要麼就不信.
一信就根本就是信耶穌.
我們不會還有一些情意結.
要怎樣靠其他的祭物.
我們沒有這些問題.
但是猶太人就有.
所以第23節.

$^{1081}$九章23節.
照著天上陽色作的物件.
必須用這些祭物去潔淨.
但那天上的本物.
這裡好像告訴我們.
原來摩西班布的聖殿和回望的擺設.
原來是照著真正的天上的物件來做的.
就是說地上的聖殿.
原來是一個Certified True Copy.
是一個和天上一模一樣的陽色.
給我們又有些新的想法.
原來不是摩西自己突然間.
很具體的創意.
自己去設計一個聖殿出來.
原來摩西所做的.
是根據天上的聖殿的陽色來做的.
所以這裡作者告訴我們.
如果地上的聖殿系統.
尚且要潔淨.
但每天要做很多手續.
首先要倒灰.
當然大祭司未必需要做的.
是由利未人幫手的.
因為這些下欄的東西應該是利未人幫手的.
大祭司就要進到聖所裡面.
就要更換陳設餅.
那些香油要加些油燈.
金燈台不斷的.
金燈台是不可以熄的.
是要點著燈的.
在香壇燒香的.
每天就要做這些事.
接著就要開工了.
人們就要來獻祭了.
大祭司就要來幫他按手.
其實應該說獻祭者自己要按手.
接著祭司就要測驗.
要檢查一下那些祭僧有沒有殘疾.
沒有了 合格了 就可以獻了.
你獻什麼祭啊?.

$^{1121}$獻梵祭.
我獻贖罪祭.
我獻感恩的酬恩祭.
手續都不同的.
原來是A餐的手續.
B餐是平安祭.
又不同的.
所以如果獻了贖罪祭.
而吃那些祭僧.
要從民中剪除.
所以祭司要很小心的.
不要搞錯了.
要聽清楚.
要下單.
贖罪祭.
手續又不同的.
獻平安祭.
那些祭肉可以帶回家吃的.
獻祭者和幫手獻祭的.
就是祭司和獻祭的人.
因為平安祭.
就可以吃那些祭肉.
可以回家吃.
去到第三天就不可以再吃.
就要燒了它.
家裡的人.
甚至乎家丁都可以吃.
就叫做神人共享祭.
但是贖罪祭就不是的.
獻祭者是不准吃的.
祭司就可以吃.
你可以想像得到.
但是如果是全民的贖罪祭.
祭司都不可以吃.
這個是利未記的第三到第五章.
就講得很清楚了.
所以祭司是有一個守則.
是有一個menu的.
是要知道要獻什麼祭的.
就是你乾炒牛河.

$^{1161}$或者是廈門炒米.
都有不同的.
那些鍋子都要洗乾淨.
或者怎麼樣的.
你可以想像得到.
留意一下第23節這裡.
照著天上樣式作的物件.
必須用這些祭物去結證.
但那天上的本物.
自然當用更美的祭物去結證.
開始講那個更美.
這個字在希伯來書是一個關鍵字.
更美這個字出現很多次.
這裡告訴我們.
因為基督並不是盡了人手所做的聖所.
只不過是真聖所的影像.
乃是盡了天堂.
如今為我們顯在神的面前.
原來主耶穌的獻祭.
祂的犧牲.
祂在神的面前的贖罪.
祂現在已經去到天堂了.
在那個真正的聖所.
為我們顯在神的面前.
祂不是一個死了的冤魂野鬼.
祂是復活.
回到天堂.
而且祂在神的真正聖所裡.
是永遠的贖罪祭.
永遠有功效.
為我們在神的面前顯現.
第25節.
也不是多次將自己獻上.
「將那大祭司每年帶著牛羊的血」.
就是在說那個贖罪祭.
對不起.
那個贖罪日的獻祭.
進入聖所.
「將那大祭司每年帶著牛羊的血」.
進入聖所.

$^{1201}$第26節.
如果這樣.
祂就是耶穌.
從創世以來.
必定要多次受苦了.
就是說原來每年的贖罪日.
只不過是一個預表的象徵.
象徵著人類的罪.
必須透過第三者的救贖.
才可以得到神的拯救.
所以是生生不息.
每年都要重複這樣做.
是要告訴人知道.
要你想起自己是有罪.
所以我們來到禮拜堂.
想起自己有罪.
然後立刻靠耶穌潔淨.
是重要的心靈的發現.
不信的人沒有這個動作.
沒有需要做這個動作.
但是一個基督徒.
每次來到.
安靜.
神讓你想起自己.
剛才真的激氣.
罵了兒子一頓.
罵到我又激氣.
他又哭得七彩.
都知道自己不對.
立刻想起.
立刻悔罪.
想起罪是很重要的.
立刻靠耶穌的寶血洗乾淨.
然後記住.
回家要跟兒子說.
對不起.
剛才媽媽大聲了.
你不要怪我.
因為我有些激氣.
這樣挽回大家的關係.

$^{1241}$這樣的贖罪是很健康的.
所以這裡告訴我們.
原來贖罪日.
每年都要設定.
是為人類的緣故的.
是要讓人類有一個.
敏銳和神關係的心.
所以現在我們更加好了.
不用一年一次.
每個星期來到神的家.
來到神的殿.
其實你沒有事奉.
你的迫切性不大.
其實如果你稍為事奉過.
你就知道.
每一個事奉者.
要上港台也好.
參與事奉也好.
其實第一件事就是安靜.
我有什麼令到神不起悅的.
你立即要在心靈裡面.
清除了才可以事奉.
否則你不要說笑.
你事奉一位永生神.
你不要下巴輕輕.
真的.
真的每一個牧者你會留意到.
是很謙卑地.
醒察自己.
然後才可以去服侍神.
這些動作.
如果我們能夠越敏銳就越好.
是對我們的新心靈絕對健康的.
不要積壓.
更不要掩飾逃避.
這告訴我們.
原來耶穌基督是一次就完成了.
所以是永遠有功效.
就是說原來這個贖罪制.
永遠在我們的心靈裡面.

$^{1281}$永遠我們只要一信.
就立即可以解除我們和神的隔絕.
就是那個罪.
所以你想想多麼重要.
所以我真的奉勸查經班弟兄姊妹.
每一次來到神的家裡.
最好就是在座堂教會.
又不用搭電梯上去五樓,十四樓.
那些多層的教會.
來到整個座堂.
一個很明顯.
好像真的走進了聖殿裡面的一個動作.
就不要輕飄飄的.
好像來看演唱會一樣.
是應該安靜的.
沉默的.
來預備自己的心靈.
如果你參與事奉就更加要.
如果不是.
你真的不要開玩笑.
你用你自己的驕傲去事奉.
神甚至會擊殺我們.
所以我們要求主給我們帶著一個敏銳的心.
你可以留意到.
作者在這裡就做了一個很強烈的比喻.
耶穌基督已經是一次性完成了這個永恆的救贖.
所以第26節.
如果這樣.
祂從創世以來就必定多次受苦了.
如今.
如今就是十字架那一次.
這個黑西瑪利元和那個各國他山那一次.
但如今在這末世.
祂只是顯現一次.
把自己獻為贖罪祭.
好除掉罪.
就是這麼簡單.
主耶穌將所有過去的舊約時代.
一切繁瑣的律例.
那些做法.

$^{1321}$全部好像清卡數那些一筆勾銷.
耶穌基督.
祂一次的犧牲.
就將所有舊約的東西.
完全更新.
靠賴主耶穌.
祂的流血.
立即得到永恆的救贖.
和神的那一刻.
恢復和好關係.
留意27節和28節.
突然講到主再來.
按照著定命.
定命就是說.
在地球上.
天然的神安排的秩序.
人生老病死.
只可以死一次.
人人都有一死.
死後且有審判.
第28節.
像這樣.
既然人人只能夠死一次.
因為上下文說耶穌基督.
祂是一次過獻祂自己成為贖罪祭.
所以像這樣.
基督既然一次被獻.
擔當了多人的罪.
將來祂要向那等候他的人第二次顯現.
耶穌基督來這個世界.
是一次的顯現成為贖罪祭.
成為一個待罪羔羊.
在十字架上犧牲.
所以既然基督也是人.
祂只可以死一次.
不需要再死.
所以這次聖經.
直接對天主教彌撒神學的批判.
因為彌撒真正的意思.
就是每一個星期.

$^{1361}$耶穌基督再死給你看.
祂要再釘十字架.
所以你吃的聖體.
千萬不要太大力.
爆的時候聖體會流血.
因為耶穌再一次犧牲.
但是我們不接受這一套神學.
因為耶穌基督只需要來這個世界死一次.
就成為了永遠的救恩.
所以我們只是用信心.
和主耶穌的救恩接軌.
每一次.
所以每一次其實是分分秒秒都可以.
你剛剛講完道.
覺得講得很好很驕傲.
立刻就要認罪.
在你下台立刻要認罪.
哎呀真是不對啊神.
赦免我吧.
或者你講完之後.
你很灰心很絕望.
覺得自己講得很差.
又祈禱.
不要這樣想.
神自己會負責.
所以其實分分秒秒都可以和神和主耶穌的赦罪.
貼上關係.
這個就叫做一次性的永恆救贖.
所以一個基督徒.
不謙卑成長成熟的基督徒.
他只有一個.
他真是.
誇口的.
必定要指著主誇口.
每一個人能夠.
事奉到神.
都是神的恩典.
所以這裡就告訴我們.
基督既然一次被獻.
擔當了多人.

$^{1401}$就是世人的罪.
當然是信的人的罪.
記住將來就講第二次.
就是我們講的這個課題.
將來.
當然我們不知道什麼時候.
將來要向那些等候他的人.
第二次顯現.
就是說.
這個就是主再來的時間.
雖然這裡沒有提到到底是不是有千禧年.
因為這一段聖經.
主要想討論的.
就是在講耶穌基督.
只需要死一次.
永遠有效.
就是說對猶太人有效.
對外邦沒有這些.
摩西律法信仰的人.
同樣有效.
所以每一個人只要你和主耶穌.
結合了這個生命關係.
認順他的十價犧牲.
你就是.
他的兒女.
你就是那個等候他回來的.
那一位的兒女.
他第二次顯現會怎樣.
按照著.
每個人都要死.
死後是要面對審判.
但是我們在再生之年.
已經不斷.
透過主耶穌的寶血.
洗乾淨了我們的罪.
所以將來就不是再是一個.
這樣沉淪的審判.
將來是什麼呢.
第二次的顯現.
等候他的人.

$^{1441}$很清楚了.
與罪無關.
為了拯救我們.
所以.
這個是對於我們信主.
有信心的弟兄姊妹來說.
這是一個.
渴慕.
他再來的一個.
拯救.
與罪無關.
不再追究你的罪.
因為你每一次靠主耶穌的寶血.
你已經是完全結證.
所以.
其實就是不斷將那個.
贖罪制永遠.
伴隨著我們.
所以什麼時候我們不再信耶穌.
你就要面對自己的罪.
神的審判.
所以為什麼一次相信.
一直相信.
必然得救.
就是這個意思.
要一直維持在神的救恩裡面.
這一段其實也是很重要的.
因為.
我們很多時候都會擔心.
到底主回來的時候.
我們是不是會再一次被他審判呢.
其實就不是.
希伯來書已經很清楚說了.
對一個有信心的信徒來說.
特別是針對當時受信的讀者.
就是希伯來人.
他們有很大的信仰包袱.
很難一下子放下.
不靠那些物質的牛羊.
不靠那個祭司.

$^{1481}$而靠一位已經釘十字架的耶穌.
所以他們要建立這個信心.
的恩信清義.
其實一點都不容易.
當然有先行者.
保羅自己就是這樣.
他就說.
先行者.
保羅自己就是這樣.
完全撇棄了自己舊有的宗教體系.
建立一套新的.
完全的憑信心的.
羅馬書的恩信清義.
這個新的體系.
就由保羅來確立出來的.
你可以想像到.
保羅自己對於主耶穌的這個贖罪.
是多麼的.
感動和受到影響.
而他願意.
佩服地.
一生的為主耶穌服侍.
應該我們越來越大年紀.
在信主內.
應該越來越對於.
信耶穌.
或者見到人信耶穌.
你的內心有莫名的興奮.
和感恩.
因為你真的看到生命的信主.
絕對不容易.
他能夠建立了一套的恩信清義.
其實是很多人.
很多人用過的功夫.
出過的力.
所以我們要記得.
我們在這個救恩.
在這個寶血裡面.
我們將來.
不是再審判.

$^{1521}$來面對那個痛苦.
那是與罪無關.
耶穌是要回來拯救我們.
就是說.
脫離這個人世間.
物質世界.
未被救贖的世界.
那種的難處.
那種的痛苦.
那種的無奈.
那種的無助.
無力感.
那這個.
耶穌一回來就什麼都沒有了.
就完全的釋放了.
所以我們應該都要.
渴慕主回來的這個日子.
這幾段聖經.
我很快地幫大家.
揭一揭之後.
再看多兩段.
我先喝一口水.
如果有問題.
這些解經有些問題.
太深.
你一是在WhatsApp群組.
或者你私下再問我都可以.
彼得候書第三章.
彼得候書的三章.
親愛的弟兄啊.
我現在寫給你們的是第二封信.
這兩封信.
都是提醒你們.
激發你們誠實的心.
叫你們紀念.
聖先知御先所說的話.
和主救主的命令.
就是使徒.
所傳給你們的.
當然他自己也是使徒.

$^{1561}$第一要緊.
知道在末世.
必有好機肖的人.
在那裡.
無論是亂說.
或者造謠生事.
隨從自己的私欲.
出來機肖.
主要降臨的應許.
在那裡呢.
因為從列祖誓了以來.
萬物與起初創造的時候一樣.
他們這樣理解.
就是說.
沒有什麼地動天移.
沒有什麼.
根本性的改變.
為什麼會說.
有一個特別的日子.
會有這樣的情況呢.
第五節.
作者就和讀者說.
不要不相信這些機肖.
機肖就是說.
隨口這麼說的.
興之所至的詆毀的.
但是彼得在這裡就告訴我們.
有些歷史.
根據.
他們故意忘記.
從太古.
憑神的命就有了天.
就是神說的動作.
那個神的說話.
就是一個能力.
憑神的命有了天.
並從水而出.
藉水而成的地.
就是創世記的第一章.
關於天地的來歷的問題.

$^{1601}$第六節.
故意忘記.
第六節.
故此當時的世界.
被水淹沒就消滅了.
被水淹沒.
就是羅亞的方舟洪水.
但現在的天地.
還是憑著那命存留.
那命就是神的命定.
就是說.
主還沒回來之前.
地球仍然可以生存.
人類的.
因為神仍然給人有機會.
去接受這個福音.
仍然憑著神的命定.
來存留.
直留到不敬虔之人.
受審判遭沉淪的日子.
用火焚燒.
將來地球要被毀滅.
第八節.
就是要留意.
親愛的弟兄.
有一件事你們不可忘記.
就是主看一日如千年.
千年如一日.
主所應許的尚未成就.
有人以為.
乃是擔言.
其實不是擔言.
乃是寬容你們.
不願有一個人沉淪.
乃願人人都悔改.
所以第十節.
但主的日子.
要像賊來到一樣.
那日.
就是主再來的那一日.

$^{1641}$天必大有響聲廢去.
有形則的.
都要被烈火消化.
上下文都提到.
烈火消滅.
有形則.
所有的物質.
都要被消化.
整個地球整個宇宙.
物質的宇宙.
都要過去.
地和其上的物.
都要消盡.
這一切既然都要如此消化.
你們為人.
該當怎樣.
聖潔怎樣敬虔.
切切仰望神的日子.
來到.
在那日天被火燒.
就消化了.
有形則的都要被烈火融化.
有第二次提.
應該第三次.
但我們照他的應許.
盼望的是新天.
新地就在這裡出現.
啟示錄會再出現新天.
新地的有異.
居在其中.
告訴我們.
會過去會更新.
我們會被完全的消毀.
再沒有任何的住處.
但是我們有的是.
新天新地.
有神在其中.
其中特別強調的就是異.
因為上文第二章.
是說那些假先知的不異.

$^{1681}$和第三章.
有些人在機肖主的在來.
這個世界.
為什麼需要信耶穌.
為什麼要迫切.
需要逃避將來神的審判.
機肖的說話.
彼得就要來強調.
天堂.
不是我們想到的.
好像是.
什麼什麼人的樂園.
而是那裡真的有.
神的異在其中.
當然不單只是神的異.
特別強調的.
也有神的愛.
神的同在.
永生神成為我們的牧者.
等等.
所以我們可以逃避.
這個地球上的毀滅.
但是我們應該要有.
這個盼望.
等候新天新地的到來.
我們是有份的.
所以第三章.
其實重點就是.
想告訴我們.
我們要在有限的地上的人生.
面對機肖的說話.
什麼去捍衛.
我們那個信念.
就是我們對神的.
這個應許.
原來是寬容人.
在地上.
趕緊過一個公義的生活.
靠主.
去過一個公義的生活.

$^{1721}$我們自然就不會怕那些機肖.
怕那些人的.
嘲笑詆毀.
甚至用我們的公義.
去印證我們.
真的有永生盼望的一班人.
因為將來.
天堂就是有異.
居在其中的.
所以天堂從來.
不是一個逃避的.
避難所.
天堂是一個在地上有信仰.
生命改變的.
基督徒.
一直在主的愛底下生活.
然後去到將來.
就是團聚.
在永恆的公義的世界裡.
團聚.
所以不是一個避風港.
不是一些可以逃稅的.
加勒比海.
什麼三不館的地方.
沒有司法的交換權.
那些地方.
這些我們聽得太多了.
天堂不是這樣讓你去逃避.
在今生.
是已經有神的異.
透過耶穌.
我們已經有神的異的一班人.
我們繼續活在這個機肖的世界裡.
但是我們.
繼續因為信仰.
因為耶穌的真實.
我們活在這個.
受壓的一個環境.
所以我們不是逃避的.
我們是越來越展現這個世界.

$^{1761}$有基督有公義.
去到一天.
可能就是我們的生命結束.
其實就是新天新地的開展.
就是我們進入那個永恆裡面.
所以對基督徒來說.
全部是正面的.
全部都是在說.
我們要羨慕.
這個新天新地.
羨慕主的再來.
再看多一段.
暫時我們就先看到這裡.
約翰一書.
第三章.
大家留意到.
每一個重要的使徒.
他們都不約而同.
對於主再來絕對不含糊的.
是非常肯定的.
你可以想像.
初期教會受到很大的逼迫.
他們要面對.
那個困難.
耶穌又走了.
回天家了.
猶太教又沒有改變.
很少部分的祭司信主.
大部分的猶太人都是拒絕的.
再加上.
外圍的勢力.
羅馬的皇帝又加重那些逼迫.
保羅都已經殉道了.
約翰呢.
都是年老.
到達摩海放軸.
你可以想像到.
那些艱難.
地中海那些地方.
是非常困難的.

$^{1801}$不過留意第三章的第一節開始.
又提到主的再來.
你看.
父賜給我們.
是何等的慈愛.
使我們得稱為神的兒女.
我們也真是.
他的兒女.
世人所以不認識我們.
是因未曾認識他.
所以這裡和認識.
有很大的關係.
這個漢字不是普通的.
你看一下.
這個漢字是要我們去思想.
就是說我們要很認真.
如果你有個小朋友.
小小一個.
七星仔.
幾乎養不大.
在醫院的保安商.
都養了幾個月.
一路長大.
很多的問題.
一路很難.
去吃東西.
很孱弱.
很多病痛.
他慢慢終於長大.
大學畢業.
你就可以想像到.
他對父母的含深如虎的養育之恩.
他會憤外感受.
那種強烈.
因為他自己本來.
基本上養不大的.
但現在能夠大學畢業.
你可以想像到.
父母那種.
從心裡面.

$^{1841}$歡欣的眼淚.
看著自己的小朋友.
上台拿畢業證書.
還拿獎學金.
那種歡欣喜悅.
就是這節聖經.
李漢父賜給我們.
是何等的慈愛.
就是將他的愛子.
白白的犧牲.
將我們拯救.
稱為神的兒女.
我們也珍視他的兒女.
為何要強調.
珍視他的兒女.
我們有沒有感受到.
主耶穌的含深如虎.
有沒有感受到十字架的痛.
主耶穌的白白犧牲.
無條件的愛.
所以.
當我們長大.
兒子你要相信.
父母.
很辛苦養大你.
你現在能夠出人頭地.
是很多人.
為你曾經.
做過很多事.
所以我們珍視他的兒女.
世人是不認識的.
所以有一個很強烈的比較.
所以基督徒們.
我們《查經班》的弟兄姊妹.
我們不要讓世人.
知道多一點.
我們是應該很深的感受.
十字架和我們的生命關係.
我們應該很深的感受.
我們應該很深的.

$^{1881}$被主耶穌的救恩.
這份的愛.
來感動.
每一天我們是有動力.
我們說自己是基督徒.
不要忘記這件事.
因為世人.
他不會對這個神學.
對這個救恩.
他根本就是不信的世人.
所以他是不認識的.
不認識這位.
不認識耶穌.
第二節.
親愛的弟兄啊.
我們現在是神的兒女.
將來如何還未顯明.
這個意思絕對不是說.
其實我也不知道.
不過是還未正式的.
向我們來彰顯.
但是我們知道.
主如果顯現.
我們必定要像他.
因為必得見他的真體.
凡向他有者指望的.
就潔淨自己.
像他潔淨一樣.
原來我們.
其實像保羅在十三章的.
哥林多前書說.
無湖不清.
到將來就面對面了.
約翰在這裡又這樣說.
原來他很想說.
帶我們進入那個.
與主耶穌相交的生活裡面.
你在那個生活裡面.
你越來越.
很想認識耶穌.

$^{1921}$很想像主耶穌.
你就會慢慢有這種.
所謂氣質.
所謂那種.
生命的聯繫.
的聯合.
這就是這段聖經的意思.
所以不是突然一變.
我們在地上.
我們是認識神的兒女.
我們是透過.
與主耶穌的相交.
我們越來越渴望.
見祂.
認識祂的真體.
到將來不再是.
無湖了.
是我們立即變成祂的身體.
因為必定要像祂.
因為到時我們會.
見到祂的真體.
主耶穌會向我們顯現.
接著第三節就是.
再一次告訴我們.
我們不是去修煉.
我們是兒女.
我們就是在神的家裡面.
去活的一班基督徒.
所以我們很自然地.
我們就很想.
像主耶穌一樣.
凡是有這個指望的.
就潔淨自己.
所以原來等候主再來.
不是等候.
是在我們日常的生活裡面.
如常地過著一個.
要來到學校.
耶穌基督的人生.
我們的良善.

$^{1961}$我們的喜樂.
我們的平安.
我們撥散的那份對人的愛.
這樣來讓.
其他人感受到.
從你的身上.
我感受到你裡面是有基督.
有基督的形象.
有基督的形體.
所以將來我們.
不是未知.
而是將來我們肯定知道.
我們是會.
和主耶穌.
有一樣的.
當然神聖的真體.
和我們將來的靈體.
是有根本性的分別.
不過.
我們將來都會有主的形象.
在我們的天家裡面.
約翰就講到這裡.
其實就很想.
勉勵.
真的基督徒.
和上下文講的假基督徒.
的根本的分別.
所以就想告訴我們.
我們學無主的.
再回來的弟兄姊妹.
就必須要有.
這樣的裝備自己.
其實不外乎.
在生活裡面.
過一個基督徒的人生.
讓主在我們身上.
被彰顯出來.
我們就要.
在我們身上.
彰顯出來.

$^{2001}$幾個段落.
就是對於主再來的.
不同的角度的.
教導.
都是在他們神學的模樣.
上面.
對當時代的教會.
的一個勉勵.
或者是一個的.
吩咐.
特別是當時已經有.
一些異端.
和一些不是.
適合聖經真理的.
那些生活形態.
所以.
希望像撥亂反正.
那樣.
給到真基督徒.
就要有這樣的.
考慮和想法.
這裡就.
幾段已經足夠涵蓋.
新約.
一個還沒有完全應驗的.
就是主耶穌.
第二次回到地球.
肉身回到地球.
我們有些.
新朋友加入了.
我們再看多一次.
帖撒羅尼迦前書的四章.
記住我們等候的.
不是一個靈體.
而是主耶穌肉身.
要再回來.
而且所謂再回來是親自的再回來.
這是聖經肯定的.
這樣說的.
帖撒羅尼迦前書的四章.

$^{2041}$四章這一段經文.
第十三節開始.
論到誰了的人.
我們不願意弟兄們不知道.
恐怕你們憂傷.
好像那些沒有指望的人一樣.
我們若果信耶穌死而復活了.
那已經在耶穌裡誰了的人.
神也必將他與耶穌一同帶來.
就是那些在主裡面.
已經安寧的.
已經在耶穌裡面.
已經安寧的.
那已經在耶穌裡面.
就是那些在主裡面.
已經安睡.
即是已經離世的信徒.
將來不需要擔心.
他不是消失了.
是一定會與主耶穌回來的.
第十五節.
我們現在照主的話.
告訴你們一件事.
我們這活著.
還存留到主降臨的人.
斷不能在那已經睡了的人之先.
因為主必親自從天降臨.
有呼叫的聲音和天使掌的聲音.
又有神的號吹響.
那在基督裡死了的人.
必先復活.
以後我們這活著還存留的人.
必和他們一同被提到雲裡.
在空中與主相遇.
這樣我們就要和主永遠同在.
所以你們當用這些話.
彼此勸諱.
保羅告訴我們.
這些說話是要時時彼此分享的.
是有一個確實的.

$^{2081}$身體復活的盼望.
而且是因為主耶穌.
他列天而降.
肉身再回來這個地球.
改變整個形態.
因為死了的人變灰.
無論你是一年還是一千年.
都是灰.
什麼都沒有.
完全找不回你.
但是原來在永生的神的全能下.
任何一個.
活在地球上信主的人.
任何一個.
會有回自己的身份.
有回自己的榮耀.
復活身體.
這是奇事.
但是將來會有.
這個日子.
所以這是一個絕對.
基督徒很大的盼望.
對於死亡的.
最好的解釋.
和福音.
接著.
未死的身體會改變.
在空中.
在哪裡的空中呢.
在耶路撒冷.
一秒鐘就去到.
在空中.
空中很大.
你沒有辦法量度.
香港的上空還是紐約的上空.
都很遠.
你飛去多倫多.
五月的時候去探望我姐夫.
探望我姐夫.
飛了十五個小時.

$^{2121}$人都瘋了.
累到不得了.
現在身體改變.
是不是又要飛十五個小時去耶路撒冷呢.
似乎不需要.
立即就可以在空中.
和主耶穌相遇.
不需要多少光年.
不需要多少功利.
不用儲積分.
換那些飛行里數.
我們就可以永遠.
而且這個相遇之後.
就是永在.
相遇完之後就不是再見.
相遇完之後是跟隨耶穌.
浩浩蕩蕩地回天家.
是和主永遠同在.
真正的.
德勝的海船的軍隊.
記住.
首先復活的是死了的.
那些聖徒.
接著在耶穌回來還沒死的.
那一代.
身體改變.
你可以想像得到.
在地上.
意味著.
在地上又沒死.
又不信耶穌的人.
其實某程度就是暗示.
進入千禧年的災難.
我們就進入永恆.
我們相信了.
這個可以想像得到.
是真的會有一個次序.
所以這段間接就好像.
告訴我們原來沒有千禧年.
有可能不對的.

$^{2161}$我不希望誤導你.
因為下個星期就開始講無千禧年.
接著就講千禧年和後千禧年.
可能我們分了次序.
就是先簡單說了.
沒有千禧年.
接著就講後千.
最後我們才講千禧年的前千.
變成了我們幾段聖經都告訴我們.
這是我們基督徒.
要很大的安慰.
所以和大家揭一揭.
本來今天應該不夠時間揭的.
就是路加福音第十二章.
耶穌叫我們要挺身昂首.
就是說要開心到飛起.
十二章應該是二十二節.
我們先看看.
路加福音的.
對不起.
是後面那段才是.
路加福音的.
對不起.
是後面那段才是.
路加福音的.
對不起.
是後面那段才是.
路加福音的.
對不起.
是後面那段才是.
路加福音的.
對不起.
是後面那段才是.
路加福音的.
對不起.
是後面那段才是.
路加福音的.
對不起.
是後面那段才是.
路加福音的.

$^{2201}$對不起.
是後面那段才是.
路加福音的.
對不起.
是後面那段才是.
路加福音的.
對不起.
是後面那段才是.
路加福音的.
對不起.
是後面那段才是.
路加福音的.
對不起.
是後面那段才是.
路加福音的.
對不起.
是後面那段才是.
路加福音的.
對不起.
是後面那段才是.
路加福音的.
對不起.
是後面那段才是.
路加福音的.
對不起.
是後面那段才是.
路加福音的.
對不起.
是後面那段才是.
路加福音的.
對不起.
是後面那段才是.
路加福音的.
對不起.
是後面那段才是.
路加福音的.
對不起.
是後面那段才是.
路加福音的.
對不起.

$^{2241}$是後面那段才是.
路加福音的.
對不起.
是後面那段才是.
路加福音的.
對不起.
是後面那段才是.
路加福音的.
對不起.
是後面那段才是.
路加福音的.
對不起.
是後面那段才是.
路加福音的.
對不起.
是後面那段才是.
路加福音的.
對不起.
是後面那段才是.
路加福音的.
對不起.
是後面那段才是.
路加福音的.
對不起.
是後面那段才是.
路加福音的.
對不起.
是後面那段才是.
路加福音的.
對不起.
是後面那段才是.
路加福音的.
對不起.
是後面那段才是.
路加福音的.
對不起.
是後面那段才是.
路加福音的.
對不起.
是後面那段才是.

$^{2281}$路加福音的.
對不起.
是後面那段才是.
路加福音的.
對不起.
是後面那段才是.
路加福音的.
對不起.
是後面那段才是.
路加福音的.
對不起.
是後面那段才是.
路加福音的.
對不起.
是後面那段才是.
路加福音的.
對不起.
是後面那段才是.
路加福音的.
對不起.
是後面那段才是.
路加福音的.
對不起.
是後面那段才是.
路加福音的.
對不起.
是後面那段才是.
路加福音的.
對不起.
是後面那段才是.
路加福音的.
對不起.
是後面那段才是.
路加福音的.
對不起.
是後面那段才是.
路加福音的.
對不起.
是後面那段才是.
路加福音的.

$^{2321}$對不起.
是後面那段才是.
路加福音的.
對不起.
是後面那段才是.
路加福音的.
對不起.
是後面那段才是.
路加福音的.
對不起.
是後面那段才是.
路加福音的.
對不起.
是後面那段才是.
路加福音的.
對不起.
是後面那段才是.
路加福音的.
對不起.
是後面那段才是.
路加福音的.
對不起.
是後面那段才是.
路加福音的.
對不起.
是後面那段才是.
路加福音的.
對不起.
是後面那段才是.
路加福音的.
對不起.
是後面那段才是.
路加福音的.
對不起.
是後面那段才是.
路加福音的.
對不起.
是後面那段才是.
路加福音的.
對不起.

$^{2361}$是後面那段才是.
路加福音的.
對不起.
是後面那段才是.
路加福音的.
對不起.
是後面那段才是.
路加福音的.
對不起.
是後面那段才是.
路加福音的.
對不起.
是後面那段才是.
路加福音的.
對不起.
是後面那段才是.
路加福音的.
對不起.
是後面那段才是.
路加福音的.
對不起.
是後面那段才是.
路加福音的.
對不起.
是後面那段才是.
路加福音的.
對不起.
是後面那段才是.
路加福音的.
對不起.
是後面那段才是.
路加福音的.
對不起.
是後面那段才是.
路加福音的.
對不起.
是後面那段才是.
路加福音的.
對不起.
是後面那段才是.

$^{2401}$如果你是有份代表香港參加奧運.
你是被選照到.
而你可以得到入場券.
其實已經是很大的榮耀.
你能夠這樣進入.
所以將來去到這些.
真的沒辦法再面對.
自己的力量.
去逃避這些審判的時候.
我們基督徒的身份.
就真的很超然地.
顯明出來.
所以這種顯明.
正正就是挺身昂首.
所以今天我們應該要活得有尊嚴地.
始主耶穌.
我們在我們的行事為人裡面.
應該要停了.
因為是9點03分.
我們下一個禮拜就正式講.
無千禧年.
多謝主給我們在過半小時.
第四五段關於主在來.
重要的經文.
都大家一齊來分享.
很盼望我們能夠.
對於挺身昂首.
對於渴慕主回來.
我們的身體改變.
我們能夠不再在.
罪裡面被審判.
反而是主.
拯救我們.
和罪無關的.
我們進入這個永恆的天家.
可以永遠的同在.
我們內心有一個很大的渴慕.
的而且確在生活裡面.
有時我們很難.
在每一天都想到.

$^{2441}$主的回來.
有時工作,家庭,生活.
或者有很多.
不同的社會議題.
都會很困住我們大家的心.
或者我們有兒女.
兒孫.
有時我們都有很多掛心等等.
但懇求主給我們.
對這幾段的真理.
有很認真的.
認識和持守.
使我們能夠.
在很有限的短暫人生.
為主而活.
為真理.
這樣來見證.
在我們的生命裡面.
多謝您聽我們的感恩禱告.
奉耶穌基督的名.
阿們.
\newpage



\section{}
\label{sec:a8KnWbij7Uk}
\textbf{2024-09-19 周四晚查經班《淺談末世論(2)》 講員:曾裔貴牧師}
\newline
\newline
連結: \href{https://youtube.com/watch?v=a8KnWbij7Uk}{\texttt{https://youtube.com/watch?v=a8KnWbij7Uk}} ~~~~ 語音日期: 2024-09-20
\newline
\newline
\hyperref[sec:I4c89OnZC6M]{\small{< < < PREV SERMON < < <}}
~
\hyperref[sec:index_chronic]{\small{[返順時目]}}
~
\hyperref[sec:index_scriptual]{\small{[返順卷目]}}
~
\hyperref[sec:tl2Gc71LinQ]{\small{> > > NEXT SERMON > > >}}
\newline
\newline
$^{1}$請致命我們一同的禱告.
讓恩天父給我們可以再恢復星期四晚上的茶經班.
求你讓我們透過言讀活世論.
都給我們對於主在內.
我們有多一些的資料,認識.
重要的都是刺激我們的心靈.
對於主的永恆的國度.
我們是有留心的.
雖然在地上我們每天看到的新聞.
世界的局勢.
有時都很難讓我們看到有一個永恆的.
光明的,和平的國度.
但是我們很盼望神給我們有從聖經的啟示而來的信心.
一同的去等候.
這個等候是讓我們殷勤,熱心.
去為神的國度作見證.
為我們要進入的將來的天家.
來有美好的見證.
使我們能夠將福音傳開.
求你保守我們.
我們會看一些關乎神學的教導.
願神給我們在分析過程.
都有神的保守和恩典.
我們一同等候,禱告.
奉耶穌基督的名 阿們.
可能要細聲一點.
我自己喉嚨還沒有清醒.
有些癢.
今天開會比較多.
下午也咳得比較頻密.
現在可能好一點.
剛剛打了幾行字.
很簡單的將.
到底有千禧年,無千禧年.
其實前千和後千的分別不是很大.
但是有千禧年和無千禧年.
分別就很大了.
所以我們就要去了解.
到底我們是支持哪一個學說.
到底最後那一堂.

$^{41}$如果無千禧年.
會帶來什麼影響呢?.
這些都是我們要問的問題.
我們先看無千禧年.
比較簡單.
因為千禧年的學派,學說.
經文相當多.
因為我也提過.
無千禧年是關於以色列復國.
關於舊約的一些.
關於千禧年的預言.
他們是完全不會有任何的指涉的.
只會說新約的經文.
但是千禧年的學說.
前派的學說.
就會有以色列的復國.
就會涵蓋舊約一些的預言.
尤其是以賽亞書,詩篇.
有一些預言.
那些經文表面上就不是說現在的耶路撒冷.
是說將來.
那個指向是將來的一個美好光景.
所以新約的反過來理解.
就正是千禧年的那個解說.
譬如以賽亞書第十一章.
以賽亞書第五章.
有些經文你是解不開的.
那些豺狼,獅子和小朋友一起玩.
那些毒蛇.
小朋友好像和牠一起玩一堆.
那些經文是一個寓意.
還是說千禧年是有一個突變的.
一個和平的,美好的世代呢?.
千禧年的學說.
就會這樣理解這些經文.
所以為什麼無端端又會有千禧年.
無千禧年就一定要回到啟示錄.
因為始作俑者.
啟示錄的第20章.
20章就說了兩次有千禧年.

$^{81}$就是有一千年的禧年.
在地球上.
是從來都沒有過的.
但是無千禧年的學說.
就否定這一章聖經.
要按字面去解釋.
所以很簡單.
就是不需要這樣解釋.
那是一個象徵.
所以不需要理解為.
為什麼主耶穌又要第二次回來.
又要開始千禧年.
然後才來毀滅撒旦.
那個問題就在這裡出現了.
所以其實一直以來.
都是千禧年主導著神學界.
近30多年.
近代的學者.
大部分支持無千禧年.
不相信有千禧年.
地球上有一個和平的.
禧年的世界.
似乎見到還沒有改變他的立場.
見到的立場就很清楚了.
是千禧年的前派.
有千禧年的.
宣導會也是有千禧年的.
所以大家.
加入宣導會已經是千禧年了.
當然你可以有另外的解釋.
你可以不認同.
大家有些很熟悉的牧者.
我不能夠改名字.
他是無千禧年的.
所以有些有趣的故事就是.
去到系統神學的第三.
就是會講末世論.
講其他關於教會論的那些.
不准教的課.
這是見道的傳統.

$^{121}$因為他是無千禧年.
彼此尊重.
因為這不是絕對的真理.
這是相對的.
你可以認為是有千禧年.
你可以認為是沒有千禧年.
所以這不是異端.
不過理解主在來.
是會有這個小小的.
一千年的問題.
所以大家要理解.
絕對不會說千禧年才能上天堂.
無千禧年就上不了天堂.
這不是的.
我們去看一些經文.
就是啟示錄的第20章.
由第一節開始.
差不多是這張聖經引發的爭論.
就是有和沒有的背後的理據.
千禧年就會再多一些的支持點.
再多一些的解釋.
去豐富整個千禧年.
第20章.
所以我們今晚的開始.
會先看無千禧年的學說.
他們的論點.
大家理解這樣的論點.
一般就會有大概六個方面.
他們反對有千禧年的解釋.
或者不需要這樣理解聖經.
我們看第20章開始.
我又看見一位天使.
從天降下.
手裡拿著無底坑的藥匙和一條大鏈.
他捉住那龍.
他就是那個天使.
就是那個龍是誰呢.
就是古蛇.
我們回想起.
創世記第三章.

$^{161}$蛇引誘夏娃的故事.
又叫魔鬼.
也叫撒旦.
把他捆綁一千年.
相當清楚.
是指到這位魔鬼.
這位敵黨神.
這位不是黑暗代表.
是一個敵黨的代表.
把他捆綁一千年.
首先出現了一千年.
是否算得上的數值呢.
是否按字面理解的一千年呢.
在這裡就開始了.
無千一年.
來去理解.
第三節.
泯在無底坑裡.
到底在什麼地方.
我們沒有人知道.
將無底坑關閉.
用印封上.
古代的人要封一些封條.
就是沒有人開過.
其實很多那些信.
古代都是這樣的.
現代一些重要的信.
都會用過油桌.
那些漆.
蓋上去.
一開過就知道.
一定是相當重要.
重要到一個地步.
是絕對機密才會蓋上這些章.
古代就更加清楚.
所以用印封條.
沒有人.
可以開.
這些無底坑.
因為裡面有什麼呢.

$^{201}$就是撒旦在裡面.
被關閉一千年.
使得他不能夠再迷惑列國.
然後聖經又再一次提到.
強調.
是那個一千年好像似乎是計算的數值.
等到那一千年完了.
好像沒有告訴我們.
這是象徵.
這是代表的寓意.
就是很多很多年.
好像沒有這樣理解.
要我們說明是一千年完了.
然後.
以後必須暫時.
釋放他.
好像.
出師表那個.
孔明.
他就捉那個萬末.
七禽七種.
這就告訴我們.
似乎要.
捉拿這個.
敵黨神的勢力.
由創世記開始.
一直到今時今日.
都是容讓他.
迷惑列國.
影響.
試探地上的人.
包括我們的主耶穌.
但是原來去到有一個時候.
似乎的暗示.
暗示是一個很準確的.
時候就是魔鬼.
敵黨的.
勢力的中心點.
是要被關閉.
關閉多久呢?.

$^{241}$是可以數的一千年的數值.
似乎經文的意思.
是頗自勉的理解.
好像不是很想我們用一種.
抽象的聯想.
其實很久了.
反而是想我們去計算.
那個時間.
接著是四節.
我又看見幾個寶座也有在上面的.
並有審判的權.
並賜給他們.
我又看見那些.
因為給耶穌作見證.
並為神之道.
被斬者的靈魂.
和那沒有拜過獸與獸匠.
這是十七章的啟示錄.
當時因為.
邪惡軸心.
邪惡的.
三位一體.
神的三位一體.
就是聖父,聖子,聖靈.
十七章的啟示錄.
就是邪惡的三位一體.
就是魔鬼,地獸和海獸.
所以這裡告訴我們.
有拜過獸和獸匠.
也沒有在額上.
反而是沒有拜過.
沒有在額上和手上.
受過他的印記.
人的靈魂.
因為已經死了.
也復活了.
這裡似乎好像.
遙指帖撒羅尼迦前書.
我們之前說過的.
復活次序.

$^{281}$所以這裡告訴我們.
是有些人首先復活.
他們復活.
與基督一同作王.
一千年.
即是說第五節才是重要.
這是頭一次的復活.
即是說.
人的第一次復活.
是沒有拜過獸匠和受印記.
你不受印記就不能買賣.
即是今天用一些簡單的.
你沒有八達通.
你沒有支付寶.
你沒有微信支付.
你連信用卡都沒有.
你只有現金.
對不起現金我們不收了.
你不能買賣了.
這我們能夠理解.
在古代他們不能夠理解.
你沒有跟隨.
君王敬拜.
你不可以吃那些製肉.
你甚至乎被排擠.
你記得但以理的時候.
他三個朋友.
哈德拉,米薩,阿拜尼哥.
尼泊爾沙皇.
做了一個像.
叫所有全國的人都要拜.
當你們聽到.
那些割,滴,塞.
這樣彈奏的時候.
所有人都要跪拜金像.
等於跪拜尼泊爾沙皇.
眾多人跪拜之後.
有三個人站在那裡不拜.
給你一次機會.
剛才你可能聽得不清楚.

$^{321}$現在我再講一次.
通常你不敢不聽.
因為你已經引起群眾的敵視.
那時候.
他們仍然堅持.
這些就叫不拜獸匠的下場.
就是怎樣.
就是不拜獸匠的下場.
就是怎樣.
抓他扔進火窯.
火窯燒熱七倍.
扔的人都走不及.
都燒死了.
扔了之後他們在火窯裡面.
想著變成灰燼.
誰知道在裡面跳舞.
這個就是.
很戲劇性的神跡保守.
但是如果不拜獸匠.
其實死得很慘的.
十七章說他們.
面對很大的逼迫.
就是說原來有一個時候.
會進入一個千禧年.
他們享受.
復活帶來的.
那個跟隨主的賞賜.
留意第五節.
是很重要的一節.
一個暗示.
應該是按字面解釋.
他說這是.
頭一次的復活.
為什麼呢.
其餘的死人還沒復活.
就是說.
原來帖撒羅尼迦前書.
第四章.
對不起 我好像一直在說.
千禧年前派的學說給大家聽.

$^{361}$所以一直好像.
解釋經文的時候已經在引導大家.
其實現在我們在看的是.
毛千禧年那個學說.
所以.
大家不要介意.
這裡告訴我們.
這是頭一次.
不用不用 這是頭一次的復活.
強調.
上文說的頭一次復活.
是哪些人呢.
你們看到的.
就是沒有拜過獸像.
和沒有在額上手上.
受過獸的印記的人.
就是這些人.
才可以第一次復活.
那就是意味著.
有第二次復活.
第二次復活就是後面說的.
白色大寶座的審判.
那就是說.
我們應該這樣理解.
經文強調.
又強調第一次第二次的復活.
那這個就變成了.
為什麼千禧年前派.
是會加插了.
有主再來第一次再來.
和第二次再來之間.
是有一千年的那個時間.
所以第五節是一個重要的鑰節.
這是頭一次的復活.
就是那些.
二人的第一次的復活.
說的是死了的人.
那就是沒有抵觸.
貼索羅尼加前書.
四章說的復活次序.

$^{401}$必定是死了的人.
首先復活的.
那裡就沒有提到.
到底是所有的人復活.
還是說是二人才復活呢.
但是啟示錄二十章.
就很清楚的告訴我們.
頭一次的復活.
就是這些.
的二人.
他們復活之後.
做什麼呢.
又很清楚了.
接著就說回第三次說的.
那個一千年.
他們的復活就是.
與基督一同.
作王一千年.
那千禧年前派就肯定是.
說的是地球作王一千年.
就是說主耶穌.
回到地上作王一千年.
就是真真正正的.
那個管治整個大地.
在這個時候.
當然.
一千年前派就會再引進.
一些經文.
譬如二創亞書有兩段重要的經文.
另外詩篇.
都有重要的經文.
在那裡引進這些.
舊約預言了.
一直都兌現不到.
你今天怎麼看都沒理由兌現.
不應該這樣去.
實踐.
因為你這樣實踐.
你走去那裡玩毒蛇.
你跳進去那裡玩熊貓.

$^{441}$我想馬上就變了.
一手掌就拍死了那個小朋友.
但是二創亞書告訴我們.
將來是這樣的.
如果你不介意.
很快給大家看一下.
其實我們今晚是想說.
無前幾年.
二創亞書第十一章.
這些大家是讀過的.
這些經文.
不過如果我們現在擺在.
這個千禧年.
這個context裡面.
那個意義就相當大了.
十一章.
還有一會兒我們看六十五章.
都有的.
二創亞書的.
十一章.
從耶西的本.
必發一條.
從他更新的.
之子必結果實.
我已經預言了的就是主耶穌.
耶西的後裔.
大衛的子孫.
耶和華的靈必住在他身上.
主耶穌有主的靈.
就是使他有智慧聰明的靈.
謀略和能力的靈.
似乎在說執事的資格.
真的很重要.
就是你有靈性.
那個靈性.
是令你更加有.
屬天屬靈.
有一個大愛的智慧.
和對教會的存局.
你會有一個.

$^{481}$很大的胸襟.
這些是很重要的.
如果你的胸襟很窄.
甚至乎你很偏私.
那個教會就被你影響了.
是不是.
知識和敬畏耶和華的靈.
有知識但是敬畏神.
愛神愛人.
他的知識是拿來敬畏神的.
引導人敬畏神的.
這樣的.
當然我們不是在說執事.
是在說主耶穌.
他有這樣的靈.
他必以敬畏耶和華為樂.
敬畏不是怕他.
不是敬而遠之.
是敬而樂.
是充滿了一種的愛.
而甘心樂意去馴服.
你看到小朋友.
叫他做事.
不准偷吃.
那顆糖.
他哭著不吃.
和他很聽話.
是兩樣東西.
很生性.
行審判不憑眼見.
可以想像得到.
多麼的大公無私.
有內心的那種正直.
斷是非.
不靠耳聞.
我們聽完就當是.
聽完是非.
三夫六婆就當是.
這樣不對.
是一種謊言.

$^{521}$卻要以公義審判貧窮人.
貧窮人都要審判.
不是永遠富有的人.
都是錯的.
所以千萬不要有一種仇富.
以正直判斷世上的謙卑人.
以口中的仗擊打世界.
以嘴裡的氣殺戮惡人.
以口中的仗擊打世界.
以嘴裡的氣殺戮惡人.
公義必當他的腰帶.
公義的腰帶.
信實必當他脅下的帶子.
所以保羅的靈感.
全副軍裝.
都不是無端的.
熟悉舊約的這位拉比保羅.
見到軍隊士兵捉拿他.
在羅馬監獄.
他想起的.
寫給二分所教會.
村大臣所賜的全副軍裝.
就是對福音全面的形式.
對外外頭頭腳腳的了解.
對福音有一種擁抱.
能夠有熟天的眼光.
第六節就是千禧年.
當然這是前派的千禧年.
我們要公平一點.
今天還未可以發生.
任何一個以色列國的時代都不能發生.
以賽亞的時代都未發生.
豺狼與綿羊同居.
還要是綿羊羔.
不是那些老的山羊.
是那些綿羊羔.
豹子與山羊羔同餓.
少壯獅子.
和那些牛獨.
並肥畜同群.

$^{561}$少壯獅子不會感覺有大餐吃.
餐餐有送家.
小孩子就牽引這些牲畜.
牛和熊同食.
牛與小熊同餓.
獅子吃草.
與牛一樣.
獅子是吃肉的.
怎會吃草呢.
如果是吃草的.
非洲就不用那麼擔心了.
吃奶的孩子.
在哪裡玩呢.
橫灑在.
蟋蟀的洞口.
蟋蟀就是當時的毒蛇.
等於今天的飯鏟頭.
金腳帶.
青竹蛇.
大人看到都害怕.
小朋友和牠玩.
斷奶的嬰兒.
毒蛇的穴上.
在我聖山的遍處.
這一切都不會再傷人.
不會再害物.
這些押韻.
好像我們香港人說的.
傷天害理.
傷人害物.
這些押韻的字.
因為認識耶和華的知識.
要充滿遍地.
好像水充滿洋海一般.
到那天.
耶西的根.
就立作萬民的大旗.
外邦人必尋求他.
他安息之所.
大有榮耀.

$^{601}$當那日主必二次伸手.
救回自己百姓中.
所餘城的.
就是亞述,埃及,巴提羅.
古實,以蘭.
斯拿,哈瑪.
並眾海島所剩下的.
你可以想像到.
整張聖經.
如果我們放在任何一個現在的世界.
都是匪夷所思.
都是天方夜譚.
甚至你應該用靈異解了它就算了.
你不要按字面解.
如果你按字面解.
我們就現在下個禮拜.
我們的兒童主義學就去玩毒蛇.
由曾牧師帶隊.
根據二次教書十一章.
那就死了 要抓警察都要抓.
因為你們變成異端.
邪教.
帶小朋友去玩毒蛇.
所以如果這些經文.
千禧年的學說.
就會這樣將它未應驗的部分.
就要放在千禧年裡面.
就會有一個劇變.
整個世界的劇變.
某程度上就是回到.
墮落之前的樂園.
所有的動物.
都可以被人當改名.
當時的動物.
沒有分說.
有些是.
純良的.
有些是兇猛的.
墮落之前是沒有.
這樣的分別.

$^{641}$墮落之後.
可能真的有一個劇變.
是我們不知道.
聖經都沒有交代.
似乎不是一開始出來.
毒蛇就是毒蛇.
獅子就是猛獸.
似乎不是這樣的.
如果不是 阿當都不會走隻獅子來改名.
根本抓不到他.
全部自己很乖很純.
走到阿當的面前.
是不會改名的.
所以很多這些經文.
千禧年前派有很多.
這些助證.
他都會理解為這些舊約未應驗的經文.
必須要在千禧年.
發生.
所以.
某個程度我們.
叫做復活離世.
去到主的懷裡.
但是主耶穌再回來的時候.
我們是跟著耶穌.
回來的.
我們信千禧年.
我們都會在地上過千禧年的生活.
所以又豐富了我們對帖撒羅尼迦前書的理解.
這是其中一段.
我們嘗試保留65張.
給大家看一下.
有些我們覺得很奇怪的經文.
為什麼二創會這樣預言.
他寫的時候.
他自己有沒有覺得.
很難想像.
那些最兇猛的.
動物界的.
毒蛇 走獸.

$^{681}$他拿來做一個這樣的比喻.
因為他理解.
他不知道什麼時候會發生.
但是將來會有這樣的日子.
會發生.
就是當那個世界改變.
當主耶穌作王的時候.
就會有這樣的.
光景.
這樣的美麗的一千年.
我們站一下.
65張.
有一兩段經文.
大家可以看一下.
第17節.
前面那些不讀了.
看啊.
我做新天新地.
從前的事不再被紀念.
也不再追想.
你們當因我所做的.
永遠歡喜快樂.
因我.
做耶路撒冷.
為人所喜.
做其中的居民為人所樂.
我必因耶路撒冷歡喜.
因我的百姓快樂.
其中.
必不再聽見哭泣的聲音.
哀號的聲音.
其中.
必沒有數日腰亡的英鞋.
也沒有.
愁愁不滿的老者.
因為百歲死.
刑算孩童.
有百歲死的罪人.
算被奏詐.
他們要建造房屋自己居住.

$^{721}$栽種葡萄園.
吃其中的果子.
他們建造.
別人不得住.
他們栽種別人不得吃.
因為我民的日子.
必像樹木的日子.
我選民親手.
勞碌得來的.
他們必不徒然勞碌.
所生產的也不遭災害.
因為都是蒙耶和華賜福的後裔.
他們的子孫也是如此.
他們尚未求告.
我就英文.
正說話的時候.
我就垂聽.
然後又再引述十一章.
豺狼必與羊羔同食.
獅子必吃草與牛一樣.
塵土必作蛇的食物.
豺狼必與羊羔同食.
獅子必吃草與牛一樣.
塵土必作蛇的食物.
在我聖山的遍處.
這一切都不傷人.
不害物.
這是耶和華說的.
然後更特別的是.
六十六章是說新天新地.
就很特別.
就好像第六十五章.
這一大段經文.
如果你不理解是千禧年的.
禧年的福樂.
就很難解釋得到那些的意思.
所以千禧年前派.
不是說為反對而反對.
而是千禧年前派才真正有那個理據.
就是啟示錄那裡.

$^{761}$是說了好幾次.
在五節的經文裡面.
說了三次.
有一千年的時間.
還有一些人是未復活的.
而未復活的那些人.
就是要將來面對.
火湖沉淪的那些人.
現在我們就再進入.
第二十章的啟示錄.
我們讀完後面那段經文.
接著我們就說.
毛千禧年.
希望今晚夠時間說完.
毛千禧年.
他們那個理據.
大概有六個點.
我先喝口水.
如果不是太快.
我就繼續.
如果我真的說得太快.
大家跟不上.
你們就在群組裡.
私下WhatsApp我.
我了解一下大家有哪些問題.
就在下一次我們去解答.
那就會好一些.
可能都慣了.
大家都是單向.
如果你問問題.
可能會打亂了我的節奏.
可能延後的.
大家聽完再自己消化一下.
我們從第五節的下半節讀起.
其餘的死人.
還沒有復活.
等到一千年完結之後.
他們就會復活.
那個時候.
就是一個大審判的開始.

$^{801}$所以這裡是很清楚.
第五節分了.
有一群人.
這群人已經界定了.
是沒有受印記.
沒有在額頭手上.
有任何的證據.
沒有在額頭手上.
受印記.
從來都沒有拜過獸像.
就是說單單拜.
真神的一群人.
當然不只是限於猶太人.
而是所有敬拜神.
的人類.
所以這裡是.
界定得很清楚.
一千年的復活.
只是給那些.
信主在救恩.
裡面的人.
接著就是.
一千年的開始.
第六節.
再頭一次復活的.
有福了.
聖潔了.
這樣來稱讚.
誇許.
好像我們香港人說的.
發達了.
的意思.
就是聖潔了.
因為他們曾經無論如何付代價.
都要跟隨主.
拜獸.
不受印記.
某程度上都是在地上.
受盡肉身.
或者一些誤會.

$^{841}$被人排擠.
但是.
他們最終含冤得雪.
就是.
復活有份.
這個很重要.
所以.
我們每個人都有得死.
死後要面對審判.
但是在主的救恩裡面.
我們都一樣會死.
但是將來.
天使掌吹號你有沒有起來.
如果你沒有起來.
就是你不能復活.
你就不是在第一次復活裡面.
你將來都會復活的.
不過就是面對白色大寶座的審判.
現在我們講第二次.
第六節要告訴我們.
已經強調.
第四次這樣說.
在頭一次復活的.
這一班人.
是有福了.
因為進入千禧年.
第二次的死.
在他們身上是沒有權定.
因為他們已經是永活了.
他們必作神.
和基督的祭司.
並且.
和基督一同.
作王一千年.
已經是第四次這樣說.
和基督.
一同的一千年.
而且是作王.
而且是在神面前的.
祭司.

$^{881}$就是說.
可以將人帶進去救恩奉獻的一班人.
這個就是.
第一次復活的人.
然後聖經就很快.
第七節.
那一千年完了.
七節聖經已經重重複複說了.
五次一千年.
撒旦必從監牢裡被釋放.
就是說原來有.
一百集的.
撒旦等了一千年.
終於可以被釋放了.
記得是一個.
被動式的.
撒旦不是逃獄.
撒旦是被釋放.
記住.
所以我們要有這種.
眼光見識.
爭戰.
我們是站在勝利的一邊.
我們已經是勝利.
我們去fight for the victory.
所以我們要很小心.
這件事.
對不起 我說錯了.
應該是fight for the victory.
我們是由勝利去爭戰.
我們不是fight for the victory.
不是我們去fight.
是耶穌已經為我們頂十字架.
已經完成了救恩.
我們是勝利的行列.
的一個精兵.
我們不是要打一場勝仗.
所以撒旦是.
必須.
來被釋放.

$^{921}$出來他也沒有.
任何的改變.
一千年的默想.
你給他多少聖經.
他背都沒有用.
因為他根本就是撒旦.
他根本就是敵黨.
他根本就是要用盡一切的辦法.
智慧去敵黨神.
所以他出來就繼續.
迷惑地上的四方的.
就是哥格和.
馬國.
叫他們聚集爭戰.
他們的人數多於.
海沙 似乎就是千禧年.
出世的那些人.
又是一張白紙.
又是要面對和我們一樣的.
試探等等.
當然我們應該要問的問題.
就是為什麼神要.
容讓 這是千禧年要回答的問題.
為什麼神要容讓.
一千年地上又重複.
魔鬼又在地上.
迷惑這些人呢.
為什麼不索性.
好像沒有千禧年就最好.
皆大歡喜 沒有的.
沒有千禧年的 耶穌來.
就是開始教會的年代.
只會越來越好 耶穌回來.
就解決了 新天新地.
不會有魔鬼的迷惑.
這就是最簡單的解釋.
但是聖經就多了.
這麼詳盡的千禧年的描述.
逼令我們.
沒辦法不去想.

$^{961}$到底是有千禧年還是沒有千禧年.
這就明顯地告訴我們.
過了一千年之後.
原來在第二次復活之前.
還是會有一個.
這樣的爭戰.
第九節 他們上來.
遍滿了全地.
圍著聖徒的營.
與蒙愛的城.
似乎就是將來的.
新的耶路撒冷.
或者地上的千禧年.
國度的耶路撒冷.
就會比較完整一點.
因為將來的新耶路撒冷.
從天降臨.
就是開展新天新地.
那應該是十九章的經文.
這裡告訴我們.
就有火從天降下.
消滅了他們.
就是這一幫跟隨魔鬼.
在千禧年裡面的.
我按照現在的字面.
當然是千禧年前的解釋.
但是按照字面的意思.
我先去解釋這段經文.
就是.
去到千禧年最後的日子.
撒旦就被釋放了.
他就去迷惑.
能夠.
短時間內.
繼續迷惑地上的四方列國.
而且人數是遍滿全地.
不是十個八個.
不是只有幾個.
而是這麼多的人.
你可以想像得到.

$^{1001}$主耶穌三年半在地球上.
在耶路撒冷.
祂有多少門徒呢?.
十二個.
減一.
因為猶大一直都是厭濟.
然後就死了.
吊頸死了.
有十一個門徒.
三年多才有十一個門徒.
撒旦一出來.
就遍滿全地.
迷惑那些人.
你可以想像得到.
那些邪惡.
對人心會有一個很大的.
好像一個吸的.
叫什麼?攝石.
那種的磁力.
那個磁吸磁力.
可以這樣去背叛神.
但是追求真理.
永遠是少數的.
真的.
教會永遠都是少數.
認真追求神的角度的人.
更加是少數當中的少數.
所以我們不要小看自己.
大家這麼經年累月.
我們差不多幾年了.
每個星期四.
不計那些假期.
病了.
一個半小時去聽聖經.
其實都不簡單.
聽完之後再去實踐.
更加反映我們對神的心的真實.
所以這些告訴我們.
撒旦很容易就可以召集.
遍滿全地的跟隨者.

$^{1041}$圍住蒙愛的城.
就有火來降下.
消滅了他們.
迷惑他們的魔鬼.
很清楚的魔鬼.
這次就泯在流亡的火湖.
就是獸.
和假先知.
所以地獸海獸.
地獸就是地上的政權.
海獸就是假先知.
就是宗教的背叛.
就是所在的地方.
他們必晝夜受痛苦.
直到永遠.
這個就真是撒旦.
邪惡的三位一體的最終結局.
不過記住.
他就召集了很多敵黨神的人.
那些真是人來的.
有血有靈魂的人.
將來就要去到火湖.
接著就去描述第二次的死.
就是說第二次復活就是第二次的死.
復活之後死.
你了解到嗎?.
第一次復活就是我們跟隨耶穌.
耶穌呼召我們復活.
我們就進入千禧年的享受.
接著有些人是沒有復活的.
死了但是沒有復活.
但是那些人在千禧年就沒有復活.
但是一復活之後就立刻進入審判.
是白色大寶座的審判.
第十一節.
我又看見一個白色的大寶座.
與坐在上面的.
從他面前天地都要逃避.
在無可見之處.
就是說沒有一個人可以避席.

$^{1081}$在這個審判台前.
我又看見死了的人.
無論大小都站在寶座前.
案卷展開.
並且另有一卷展開.
就是生命冊.
死了的人都憑著這些案卷所記載的.
照他們所行的受審判.
接著就說明有個生命冊.
說明有一個不能逃避的地方.
不會說我在那裡處邊.
沒有人看到我.
是沒有這件事.
你不會走到天涯海角.
是沒有了 因為已經沒有天涯海角了.
天地都要逃避了.
所有事情都要面對在白色大寶座的面前.
這樣的意思.
就是說無一倖免了.
唯一的倖免就是在救恩裡面的人.
就是這個生命冊.
他說死了的人都憑著這些案卷記載的.
照他們所行的受審判.
於是海交出其中的死人.
死亡陰間交出其中的死人.
他們都照個人所行的受審判.
死亡和陰間都被拎在火壺裡.
很奇妙.
為什麼那些是抽象的東西.
為什麼又會成為一個物質.
丟在火壺裡.
就是說不會再有死亡.
不會再有陰間.
就是說不會再有任何一個地方.
可以躲避.
造物主的面前.
那個presence是真的永恆的presence.
你可以想像得到.
我們將來在永遠.
我們見神基本上不需要朝思暮想.

$^{1121}$快點到星期日我們可以見神.
不需要了 將來在永恆裡面.
我們就會見到造物主.
因為連死亡 連陰間都要被拎進火壺裡.
那些死人就是.
哇 我是作惡多端的.
我在死亡裡面 我在陰間裡面.
我不想見到你 不用見到你 多開心.
沒有啊 沒有啊.
死亡和陰間都被拎進火壺裡.
不過留意十四節.
這裡告訴我們.
這火壺就是第二次的死.
就是說第二次復活那幫人.
是立刻進入第二次的死.
哇 那真的很恐怖.
你叫醒我 原來想我死.
想我下火壺.
經文就是這個意思.
所以其實我們要很珍惜.
很感恩.
我們在聖餐.
其實不應該只是聖餐.
我們很感恩我們能夠是第一次復活.
有份的人.
這裡就說了.
第十五節.
若有人明志活記在生命冊上.
他就被拎進火壺裡.
這裡就很清楚的告訴我們.
這個就是那個第一次復活的人有福.
第二次復活的人沒有福的原因.
就是他的名字沒有記在生命冊.
記在生命冊.
就不是順利交接.
不跟教會把你的洗禮名單.
過進去.
好像有個USB手指過進去生命冊裡.
是的 正常是的.
但是要記住.

$^{1161}$真心跟隨主.
一直不改變的.
在地上的教會名錄上.
才是一直得救.
你掛了個名字在那個名冊.
你不要說多少號.
因為大家真的有號碼的.
如果我亂說對號入座就死.
我又不記得自己多少號.
不知道是哪一號.
你信了之後.
你不再回教會.
不再信耶穌.
你的名字還是在感受堂的名冊.
但是你過不了生命冊.
因為你根本就不信.
所以清教徒很緊張這件事.
一次得救 一直相信.
就一定得救.
要是一直相信.
保守自己在神的救恩裡面.
這是我們自己要做的責任.
所以是一次相信 一直相信.
就一定得救.
不是一次得救 永遠得救.
這樣就太過扭曲了聖經的恩信清義.
所以這裡就讓我們看到.
毛遷希年在這裡.
他一概否認了.
20章應該按字面去解釋.
不應該按字面去解釋.
他們的理據就是告訴我們.
整個新約講毛遷希年.
只有20章在這裡講.
還有啟示錄.
很多的上下文.
是應該按靈意去解釋的.
不應該按字面去解釋的.
是解不通的.
為什麼這20章.

$^{1201}$突然要按字面去理解呢?.
這就是他們最大的argument.
他們的理據.
大家可以這樣看的.
他們有一段經文這樣引述.
《馬太福音》.
應該是第12章了.
記不記得那個比喻呢?.
耶穌說被趕走的那個人.
看回前面的經文.
第22節開始.
12章的22節.
當下.
其實他們沒有那麼長的.
不過我詳細的解釋這一段經文.
待會我們也有時間.
當下有人將一個被鬼附著.
又黑又啞的人.
帶到耶穌那裡.
耶穌就醫治他.
甚至那個啞巴.
又能夠說話.
其實他本身不是啞的.
是鬼令他又黑又啞.
就醫治了他.
而且又能夠看見.
又能夠說話又能夠看見.
眾人當然很驚訝.
就說這個不是大衛的子孫嗎?.
但法利賽人聽見.
就說這個人趕鬼.
無非是靠著鬼王被西北.
第25節.
耶穌知道他們的異念.
就對他們說.
凡是一國自相紛爭.
就成為方牆.
一城一家自相紛爭.
必站立不住.
如果撒旦趕逐撒旦.

$^{1241}$就是自相紛爭.
一國怎能站得住呢?.
我若靠著別西北趕鬼.
這些就是法利賽人.
無理取鬧的.
為反對而反對.
就說他是靠著鬼王別西北趕鬼.
就是說這些是掩眼法.
你訴諸於一個更大的勢力.
去趕這個人的鬼.
他們又認為那個人是被鬼附.
因為似乎是.
真的人都看得到.
鬼附明顯到.
令到那個人又乞又啞.
所以不是一些精神病.
是真的鬼附.
但是他們不接受耶穌有這樣的能力.
去驅逐這些靈界.
所以他們只能夠說.
耶穌是靠黑暗的勢力.
別西北鬼王去趕鬼.
那麼耶穌的解釋.
如果我靠別西北趕鬼.
那麼你們的子弟趕鬼又靠誰呢?.
法利賽人都是信有靈界.
殺刀的那個人就不信.
不信有鬼 有天使.
這些可能大家對當時的背景都要了解.
就是一些無神主義.
但是又看別人的.
他們信有神 但不信有鬼.
這些去到馬太福音26章那些.
有時作者會將那些背景告訴大家.
去到受難州.
他們去抵擋耶穌的時候.
法利賽人和那些民事.
民事有些是殺都該的.
本來大家不arm arm 的.
大家是兩個信仰系統.

$^{1281}$但是會走在一起.
為了一個共同利益就是要逼迫耶穌.
陷害耶穌.
所以變成好朋友去找一些很難駁斥的把柄.
他們認為無懈可擊的.
去攻擊耶穌.
但是現在這裡法利賽人是信有鬼魂的.
所以他們只是說.
耶穌不是沒有趕鬼.
不過是靠鬼王趕鬼.
接著看下去.
你們的子弟趕鬼是靠誰呢?.
這樣他們就要斷定你們的是非.
就是說你們是錯的.
因為屬靈的聖潔.
怎麼可能容讓魔鬼的介入呢?.
耶穌怎麼會訴諸於這些靈界.
來做這些事情呢?.
一定是從他自己的能力.
或者是從父神的旨意.
來知道在這一刻要醫好被鬼附的人.
去印證福音的大能.
沒有什麼是不可以做到的.
盲的,啞的.
可以立刻把鬼魂趕走.
立刻完全恢復意志.
福音書有很多這些例子.
所以大家不要怕鬼.
大家要有屬靈的眼光.
當然叫你不怕鬼就不要看鬼片.
你不要沾染那些自己嚇自己.
嚇完之後又信有鬼就不對了.
我們要記得這些.
我不知道基督徒是不是很喜歡看鬼片.
如果沒有市場的就不會經常都可以做.
經常都可以做的那些恐怖片.
來來去去都是那些.
陰陰森森,鬼屋.
諸如此類.
或者特別多小朋友的鬼仔.

$^{1321}$古今中外都是這樣的.
但是聖經說的不是這樣的.
靈界的介入干擾.
這些告訴我們.
接著耶穌就說了第29節.
人怎能夠進狀氏家裡搶奪他的家具.
除非先捆住那狀氏.
才可以搶奪他的家財.
這一節聖經就是毛遷熙年他們認為.
可以拿來解釋啟示錄第20章.
明顯耶穌是說了一個比喻.
這個比喻就說到那個拘禁.
說到那個攔阻.
所以20章整章的意思大概就是這樣.
不過是說黑暗勢力被阻礙.
直到主耶穌第二次回來就沒有這個問題了.
就進入那個新天新地了.
所以不需要理解為遷熙年按字面去理解.
現在是捆綁著撒旦的意思.
就是說現在福音已經開始了.
耶穌已經傳揚這個福音了.
所以現在就是教會的時代.
無所謂有另外一個遷熙年的時代.
教會的時代一完結.
就是新天新地了.
這個就是毛遷熙年唯一的賣點.
教會一直到耶穌再回來.
就是新天新地了.
所以現在教會只會一直傳福音.
因為壯士已經被捆綁了.
可以來到大地得到福音了.
所以朝向一個越來越美好的新天新地.
這個就是毛遷熙年的理想主義.
就不接受又會有一個實質的一千年.
然後魔鬼再被拎入火壺.
他就不能夠接受這個東西了.
主要就是他根據這一節.
問題就在這裡了.
這一節聖經只是給這一段的比喻.
給這個神蹟.

$^{1361}$和這個法利賽人對主耶穌的無理取鬧的挑戰.
不拉肩的敵黨.
所以耶穌才說這個比喻.
意思就是說.
如果撒旦沒有被捆綁.
怎麼能夠驅逐他呢.
人怎麼能夠進壯士家裡搶奪他的家具呢.
除非是先捆住那個壯士.
才能搶奪他的家財.
意思就是說.
撒旦已經失去他真真正正的權力.
耶穌基督已經對屬靈的黑暗有一個捆綁.
大家有看過路加福音第十章.
耶穌曾經說過這樣的說話.
很快的 大家讀一讀就行了.
十章的路加福音.
第21節.
對不起 應該是17節這段先.
他說那70個人歡喜的回來.
就是耶穌叫他們兩個兩個出去傳道.
這70個人就歡喜回來跟耶穌說.
主啊 因你的名啊.
就是鬼都服了我們.
耶穌就跟他們說了.
我曾經看見撒旦從天上墜落像閃電一樣.
我已經給你們權柄.
可以踐踏蛇和蠍子.
又勝過仇敵.
一切的能力.
斷沒有什麼能夠害你們.
然而不要因鬼服了你們就歡喜.
要因你們的名記錄在天上歡喜.
記錄在天上就是那個生命冊.
所以不要不斷去誇耀自己的事奉的能力.
要誇耀的是自己獲得無條件的救恩.
和可以參與服事的這一份福氣.
而來歡喜.
耶穌這樣說.
不過耶穌說了.
為什麼你們可以經驗到屬天的能力呢?.

$^{1401}$因為原來祂告訴我們.
撒旦從天上已經墜落了.
和已經賜給我們屬主的兒女.
有踐踏蛇和蠍子的能力.
如何踐踏呢?.
當然我們就不知道.
你不是說見到一個人不聽福音.
不肯收丹章的.
你一拔就打他.
踐踏吧 一定不是這樣.
所以這裡是一個讓我們看到.
好像打開了一個屬靈的眼界.
讓我們看到.
我們去傳福音.
原來是和屬靈的黑暗挑戰.
但是記得耶穌很想我們保持那個現實.
就是不要因為鬼服了你們而歡喜.
所以事奉有果後的下一步.
就是事奉的驕傲.
接著慢慢就變成一個墮落.
往往傳道人的失敗.
是由傳道很成功開始的.
所以要很小心.
要時時都保留著自己的赤子之心.
就是我能夠成為神的兒女.
我能夠參與在神的角度裡一點點的事奉.
來去感恩.
榮耀是要歸給神的.
這個真的要銘刻在自己的心靈裡.
那個事奉就會來.
因為你每一次參與.
你都會知道你全力以赴.
其實是應份的.
但是全力以赴的另一個角度.
原來是神的角度彰顯著.
所以你想想.
是因為撒旦墜落.
是因為神給他的時間已經沒有了.
我們才能夠做到這些拯救人靈魂的工作.
所以我們必須要一生保持這一份謙卑.

$^{1441}$有謙卑就會有熱誠.
有謙卑就會有一份.
你真正從內心裡看到福音的寶貴.
你真心參與.
不介意別人怎麼看.
不介意一些誤會.
用你的人生去澄清那些誤會.
耶穌就是這樣.
我們就應該要這樣一齊.
為神這個國度去傳揚.
所以耶穌說的是很重要的.
不要因為鬼服了你們而歡喜.
要因為你們的名字.
被記錄在生命冊上歡喜.
記錄生命冊其實就是信耶穌舉手而已.
這樣有什麼好歡喜的.
要我五十年都是這樣歡喜.
其實你不知道.
你不斷維持這一份初信的單純.
其實比你有能力難.
是很難的.
信了久了就不單純了.
那個才是問題的開始.
所以主耶穌的說話.
和這些使徒說.
其實不是潑冷水.
而是讓他們看到.
真正應該要這樣來服事.
所以這裡就告訴我們.
撒旦的罪惡.
這裡是耶穌親自說的.
剛才那個只是一個比喻.
是對那些法利賽人的無理指控的一個責備.
告訴他們.
如果撒旦沒有被捆綁.
怎麼能夠趕鬼呢.
怎麼能夠神的能力.
可以在這個被鬼附的人身上.
彰顯到那種能力呢.
所以不是邪惡的勢力附在耶穌身上.

$^{1481}$而是耶穌知道神的旨意.
和運用神給他的權柄.
去將那個啞和盲的人.
趕出他內心的污鬼.
所以這一節的比喻.
嚴格來說.
不應該這樣對號入座.
去二十章的啟示錄.
因為二十章的啟示錄.
真的很詳盡說那個千禧年的描述.
尤其是有第一次復活的那幫人是什麼人.
第二次復活又立刻死的是什麼人.
是那幫是什麼人.
是講得很詳細.
是分得很清楚的.
所以如果你硬說.
全本聖經只有二十章講千禧年.
你就說一段經文就當有千禧年.
無千禧年就不接受.
他給的解釋就是說.
馬太福音十二章也有這樣的理解.
捆綁了.
自然就是一個欺凌.
所以教會時代就是進入欺凌.
就無所謂真的有一千年的欺凌.
教會時代的開始就是千禧年.
不過他們不用這個term.
因為他們不接受二十章的字面解釋.
所以教會時期就是撒旦被捆綁的時候.
就是等到耶穌第二次回來.
就會有新天新地.
就是這個意思.
所以波羌時代就是撒旦被捆綁.
另外他們就在二十章做一點點功夫.
他們的解經.
二十章的第四節.
因為我要先看我的筆記.
第二十章的啟示錄.
第四節.
我又看見幾個寶座也有在座上面的.

$^{1521}$並有審判的權並賜給他們.
我又看見那些因為給耶穌作見證.
並為神知道被斬者的靈魂.
和那沒有拜過獸和獸像的那些靈魂.
他們復活.
與基督一同在皇一千年.
他說這個復活是在天上的.
不是千一年地上的.
因為與基督同在在天上.
所以他們理解這個復活是天上的.
所以是靈體.
所以他們不接受地上物質的一千年.
是在天上復活.
如果真的有這個復活.
就是將來天上的復活.
這是第一個點.
第五節.
與基督一同復活一千年.
第一次復活是在他們去到神的面前.
又是靈體去到神的面前.
他們活在神的面前.
所以不應該看為地球上的一千年的開始.
另外他們否認有一千年的千禧年.
他們有好幾段經文.
但這些經文全部都不是說關乎千禧年.
是關乎復活.
他們只能夠相信一次復活.
不相信有兩次復活.
我們有幾段經文.
看看約翰福音第五章.
第五章.
不事大辭的神蹟.
之後就有些經文.
耶穌與他們辯論.
所以神蹟才引發辯論.
為什麼你有這個能力.
大家可以看看第20節.
可能要長一點的經文.
你就會明白這個意思.
五章的十九節.

$^{1561}$到近尾的三十節.
這段是難解的.
約翰福音難解的經文.
其中的一段.
發生在不事大辭.
我們去過耶路撒冷.
真的去過不事大辭.
不過當然已經變成了死談.
但是真的有七個口.
是真的出土出來的.
不是大辭.
根據經文後來的出土.
那個描述就是耶穌說的約翰的描述.
是有七個詞.
我們看看下面.
第十九節這一段.
讀完之後再去分析.
耶穌對他們說.
就是那些敵黨耶穌的猶太人.
不能夠接受.
耶穌在安息日.
好像是很吹毛求疵.
為反對而反對.
Stick to oneseek日.
什麼都不准做.
這些宗教迷思不能夠與時並進.
不能夠有一種融會貫通.
耶穌來已經是一個新的世代.
他去重新去解釋律法.
或者不是他否定廢掉律法.
完成律法.
使律法圓滿.
但是他們就不理那麼多.
你在安息日.
叫一個病人拿玉子.
起身行走.
就是犯安息日.
工作嘛 拿個玉子回家就工作.
我們覺得很好笑的事.
但是猶太人就是這樣.

$^{1601}$當時的教派.
十九節.
他們說.
我實實在在的告訴你們.
只憑著自己不能做什麼.
唯有看見父所作.
子才能作.
父所作的事子也照樣作.
父愛子.
將自己所作的一切事.
只給他看.
還要將比這更大的事.
只給他看.
叫你們稀奇.
父怎樣叫死人起來.
使到他們活著.
子也照樣.
隨自己的意思.
使到人活著.
父不審判什麼人.
乃將審判的事.
傳交與子.
叫人都尊敬子.
如同尊敬父一樣.
不尊敬子的.
就是不尊敬差子來的父.
這個不是語帶相關.
是必然的一個真理.
原來神已經將.
全權所有的榮耀.
都交付給耶穌.
但你去尊敬耶穌.
其實等於尊敬父.
是要這樣理解的.
所以異端往往就在這些地方出事.
就是不會覺得.
耶穌與父神.
同等.
這樣的情況.
所以第23節.

$^{1641}$叫人都尊敬子.
如同尊敬父一樣.
不尊敬子的就是不尊敬差子來的父.
我實實在在的.
告訴你們.
哪聽我說.
由信差我來者的.
就有永生.
不至於定罪.
是已經出死入生了.
我實實在在.
再強調這個.
出死入生的更清楚的定義.
我實實在在的.
告訴你們.
時候將到.
現在就是了.
死人要聽見神兒子的聲音.
因為父怎樣在自己有生命.
就似給他兒子.
也照樣在自己有生命.
並且.
因為他是人子.
就似給他行審判的權柄.
你們不要把這事看作稀奇.
時候將到.
凡在墳墓裡的.
都要聽見他的聲音.
就出來.
行善的.
復活得生.
作惡的.
復活定罪.
這裡.
有些考究.
我憑著自己不能夠作什麼.
我怎麼聽見就怎麼審判.
我的審判也是公平的.
因為我不求自己的意思.
完全是想告訴猶太人.

$^{1681}$耶穌基督所做的.
尤其是在約翰福音.
一而再.
再而三強調.
只有這樣的權柄.
但是這個權柄是由.
父神那個時候到.
他就會執行的.
主耶穌是完全.
明白.
神的時候.
神的心意和神交付的權柄能力.
約翰福音是一而再這樣說的.
他不會僭越.
神的時間.
和那個位份.
他一而再說自己是被差遣來的.
他一直都說自己是子.
但是父將一切榮耀.
權柄.
全部交付給耶穌.
所以約翰福音說的.
基督論是相當重要的.
真理.
但是我們要留意這裡.
就是剛才看的.
生和死.
義人和罪人.
那個分別.
所以似乎這裡.
好像暗示.
告訴我們.
我們應該理解.
有兩次的復活.
因為.
這裡告訴我們.
凡是在墳墓裡.
都要聽見他的聲音.
就出來了.
行善的復活得生.

$^{1721}$作惡的復活定罪.
似乎那個定罪是永恆的.
所以.
這裡的聖經.
其實應該這麼說.
是千禧年前派.
復活只是在天上.
是一個.
你可以說是一個.
的理解.
毛千禧年這樣理解.
復活只是在天上.
是不對的.
約翰福音這裡很明顯說了.
行善的.
復活.
得到真正的永生.
作惡的.
是要定罪.
這個定罪會不會就是剛才說的.
第二次復活的那幫人.
其實已經死了.
但是復活之後.
要再經過白色大寶座.
再一個公正.
嚴明的審判.
然後再死.
死了復活再死.
這個就是那些罪人.
死了復活永生就是行善的人.
所以我們一生.
一定要寶貴我們所有的福音.
無論什麼困難.
我們都要.
否認耶穌.
因為如果不是復活.
就要再死.
就要被扔進火壺.
火壺就是永死.
不是死了就沒事.

$^{1761}$是永遠在火壺裡面受苦的死.
這樣叫永死.
這個就是.
聖經給我們這個理解.
是耶穌自己親自說.
是有兩種的審判.
所以如果我們在這裡說.
其實是有肉身的.
是有時間的.
分開的都說得通的.
這個就是將.
剛才說的第一點.
說不會有千禧年.
就算是說的第四節.
復活不過就是在天上.
沒有了.
地上的.
那個經文的.
那個理解.
還有一些經文的.
比如《少行傳》.
是普羅講的.
在伊柏王面前講的.
24章.
可能大家平時就沒有留意這些經文.
我們從第十節開始.
秦府點頭叫普羅說話.
他就說.
在這個千夫長面前說的.
那個申辯.
我知道你在這國裡.
斷事多年了.
所以我樂意為自己分數.
你查問就可以知道了.
從我上耶路撒冷禮拜.
到今天.
不過有十二天.
他們並沒有看見我在.
殿裡或者在會堂裡.
或者在城裡.

$^{1801}$跟人辯論.
聳動眾人.
他們現在所告我的事.
並不能對你證實.
但有一件事.
我向你承認.
留意這裡.
就是他們所稱為異端的道.
我正按照著那道.
事奉我祖宗的神.
是按照著那個道.
就是那個真理的啟示.
去事奉神.
所以我們記得.
是有道去守.
有道去信.
有道去傳.
這樣才是事奉.
我們這一位造物主.
當然本來是在和猶太人說.
事奉我們老祖宗的神.
就是耶和華神.
就是舊約所啟示的神.
但是我們今天新約.
我們是有道去信.
去傳.
去守著.
這樣去事奉我們的神.
所以我們記得我們的事奉.
是要有道.
是有真理.
是要知道自己信什麼.
這次是合乎律法的.
和先知書上的一切記載.
所以真是有理據.
保羅是知道自己信什麼.
和知道自己在傳什麼給人.
諸如此類等等.
接著.
當然要留意.

$^{1841}$並且靠著神.
盼望死人.
無論善惡.
都要復活.
就是他們自己也有這個盼望.
我因此自己勉勵.
對神,對人.
常存無虧的良心.
接著就說.
他自己的傳道的生涯.
保羅在這裡特別強調.
有一個日子.
就是他盼望.
無論善惡.
死人都會復活.
這就是猶太人堅信.
和真是一定會發生的.
所以.
保羅也這樣理解.
將來的復活.
是義人.
和惡人.
都要復活的時候.
就一定要問這個問題.
到底.
是不是會有一個次序呢.
啟示都就很清楚的告訴他.
那個次序就是一千年.
但是無千禧年就說.
不需要介意那個次序.
是復活.
但復活就在天上.
就不會在地上發生.
那就看大家怎樣去理解這個解釋.
覺得無千禧年是不是牽強呢.
還是說千禧年多此一舉呢.
真的多此一舉.
多了一千年在地上.
接著撒旦就綁在無底坑.
一千年後.

$^{1881}$無底坑就解禁.
出來.
但是記住.
那一千年是沒有.
死了的惡人的.
只有.
千禧年.
只有義人.
信主的復活.
然後他會繼續生養.
接著才會出現後代.
就是會被迷惑.
被試探.
那你說是不是多此一舉呢.
但無千禧年就沒有這件事.
根本就不存在.
現在就是教會的年代.
耶穌再回來就是新天新地.
當然是大歡喜.
所以無千禧年就某程度.
是挺樂觀的.
對世界不悲觀的.
因為耶穌時代.
就是撒旦被綑綁.
就是進入那個.
一直只會等待千禧年的彌臨.
對不起.
等待那個新天新地的彌臨.
是不會有千禧年的.
就是撒旦繼續迷惑地上.
可以再看多一段.
但願你輸.
十二章.
耶哀結但何.
但願你輸.
但願你有這樣的理解.
就是將來的末日.
那時保佑你本國之民的天使長.
米加納.
米加納通常是打仗的.

$^{1921}$十二章.
最後那章但願你輸.
加巴列就傳遞訊息.
但都是天使長來的.
原來蜀寧的領域.
都有階級的.
都有層級的.
當時羅馬都有.
八夫長 千夫長.
米加納必站起來.
並且有大艱難.
從有國以來直到此時.
沒有這樣.
你本國的民.
凡名錄在冊上的.
必得拯救.
誰在塵埃中.
必有多人復醒.
其中有得永生.
有受羞辱.
永遠被憎惡.
智慧人必發光.
如同天上的光.
而人將來.
將來就像太陽一樣.
發光在天上.
主耶穌的比喻.
都和.
但以理書的描述有關.
將來的智慧人.
發光如天上的光.
那使到多人歸宜的.
必發光如星.
直到永永遠遠.
但以理.
你要隱藏.
這些說話.
直到末時.
必有多人來往奔跑.
知識就必增長.

$^{1961}$留意第二節.
但以理都認定.
將來在塵埃裡.
即是塵土裡睡著的人.
必定會蘇醒.
然後其中.
有得永生.
有些是受羞辱.
永遠被憎惡.
即是不斷描述.
原來將來.
約翰福音又這樣描述.
保羅在.
《士徒行傳》的24章.
說猶太人是相信.
有審判 有將來的.
行善得生.
作惡定罪.
所以配合.
啟示錄20章說的.
似乎是說得通 解得通.
舊約這些的寓言.
這些的先知.
他們.
看到末後 當然他們絕對不知道.
什麼時候.
他們看到會有一個.
這樣的復活的.
概念.
善惡是分開的.
所以.
除非你說.
不理那麼多 一男子.
復活就一班人.
善惡 但問題是.
啟示錄幫我們分得那麼仔細.
有一千年在那裡.
來做了一個.
兩個次序.
現在叫優先醜.

$^{2001}$老人家優先排隊.
接著就再.
第二些醜.
這樣幫我們分了.
第一次復活.
是那些沒有拜.
受像 沒有.
受印記 第二次復活.
就是那些惡人.
其中 海.
陰間 死亡.
都交出自己的人.
很清楚 行善.
和行惡的那兩條隊.
所以.
如果我們這樣理解.
這個就是.
前派 千禧年派.
應該是有千禧年派.
因為千禧年再分前後.
災前和災後.
變成了.
有和沒有的分別.
是非常之大的.
看你接受.
哪一個派別.
其實我從頭到尾.
都說我是千禧年派.
有千禧年派.
那就.
你不需要受我影響.
因為近代的學者都是沒有千禧年派.
比我更有學問.
更有知識的學者.
都是沒有千禧年派.
有時聽人家講 怎樣解釋.
遲些我們會聽一個講座.
國民字幕師.
有講末世論.
我相信他都是千禧年派.

$^{2041}$有千禧年派.
到時我聽聽他講什麼.
可能他無千禧年派也不一定.
到時我聽完就告訴大家.
其實上一輩的牧者.
全部有千禧年派.
近代的學者.
多數是無千禧年派.
但不是絕對真理.
是不同的解經.
是不關上天堂的.
上天堂是走善還是走惡.
是有救恩還是沒有救恩.
我們可以看看第三個類別.
意思大概是這樣理解.
還有十分鐘.
他說千禧年那些聖徒和罪人.
仍活在地上.
是不能夠接受的.
為什麼復活再和那些罪人同流合污.
何來千禧年.
那些聖徒和罪人.
仍活在地上.
是不能夠接受的.
為什麼復活再和那些罪人.
同流合污.
何來千禧年.
有千禧年的解釋.
就用《少路行傳》.
第一章.
我不知道這樣是不是牽強.
但某程度上也不是錯的.
我們很少用這個角度去想.
但當我們集中去看這個對比的時候.
又似乎有些人.
有些人是有些人.
在對比的時候.
又似乎好像.
都講得通.
第一章的第一節.

$^{2081}$這些經文我們很熟悉的.
不過我們集中這樣看主耶穌的復活.
又好像有些不同的提示.
提亞菲羅.
我已經作了前書.
那就是路加醫生.
寫第二卷.
新約.
這些冷知識大家要懂.
一定是路加醫生.
因為路加福音24章.
《少路行傳》28章.
所以他寫的聖經.
比保羅還要多.
這個是路加醫生.
按章數.
和他運用的希臘文.
保羅加上的13封書信.
都沒有那麼多.
還有路加福音.
新約路加福音.
是全新約最長的經卷.
馬太28章都沒有那麼長.
當然是算原文的希臘文.
還有.
《少路行傳》和《路加福音》.
那個修辭.
希臘文的修辭.
是非常有學問的人寫的.
和沒有學問的人寫的.
就很不同.
什麼什麼的.
不懂中文的人.
很多口語病.
但是《路加福音》和《少路行傳》.
是非常優雅的文筆.
因為他真的是一個學者.
是的.
論到耶穌開頭一切所行所教訓.
直到他藉著聖靈.

$^{2121}$吩咐所揀選的使徒.
以後被接上升的日子為止.
他受害之後.
用許多的憑據.
將自己活活顯給使徒看.
40天之久.
一個復活了的靈體.
40天在哪裡.
在耶路撒冷.
向他們顯現.
講說神國的事.
又可以講話.
又可以隨時去到任何一個地方.
這是一個靈體.
所以這就是第三點.
對沒有天欺人的人.
他們認為.
你們這樣搞天欺人出來.
是自己騙自己.
自己在玩自己.
天欺人.
好人和罪人一起.
天欺人.
好人和罪人一起.
是怎麼可能的呢.
為什麼你們要製造一個天欺人出來.
自己令自己這麼混淆呢.
復活了.
怎麼還要和那些罪人一起呢.
但是.
主耶穌他自己.
都有40天已經復活.
是永活的靈體.
他都在地上.
如果神要製造一個真是欺人的.
一千年.
每個人都是復活的靈體.
在地上有自己的身體.
為什麼又不行呢.
這就是有千禧年的解釋.

$^{2161}$對於他們認為.
既然復活在天上是靈體.
不應該再在地上.
不應該再和那些罪人一起.
地上現在就不同了.
因為我們還沒有復活.
現在和地上的罪人一起.
你隔壁屋的.
或者你隔壁位的.
都是罪人.
還沒有去到一個真正不再犯罪的階段.
但是既然已經復活了.
為什麼還要有一千年.
和那些罪人一起呢.
那不就是復活沒有意義嗎.
但是有千禧年的.
其中一個解釋.
主耶穌都是以這個靈體.
在地上和他的門徒.
和有500個弟兄.
都見過耶穌.
在地上40天.
40天.
這樣去聽耶穌講道.
所以.
如果40天都尚且可以.
這樣的時候.
為什麼千禧年.
當神都改變了世界之後.
為什麼又不行呢.
這個就是有千禧年的解釋.
去解釋.
其實不需要這樣想.
復活之後就永遠在天上.
不應該再有.
和罪人為伍的時間.
但是我們有千禧年.
就是說有一千年.
都會有試探.
都會有將來的.

$^{2201}$撒旦的再出現.
再被釋放.
所以無千禧年就說.
不要這樣想.
但是耶穌40天都是一個靈體.
復活.
向門徒顯現.
就是說如果神要這樣做.
是可以這樣做.
是可以的.
第四個論點.
無千禧年.
也是這樣說.
神都在地上掌權.
為什麼還有罪人.
他們還要堅持犯罪呢.
這個正正就是.
撒旦了解.
人的自由意志.
的問題.
所以當神一直都沒有.
再收回.
而那個自由意志繼續運作.
這樣的時候撒旦再被釋放.
就會出現.
這個相同的問題.
就是墮落之後的人類.
就會受撒旦的迷惑影響.
所以千禧年完結.
再有這個試探出現.
為什麼又不可以呢.
既然有千禧年.
就必然會有地上的人.
仍然可以選擇的空間.
所以又是不需要這樣想.
為什麼神.
還要堅持.
這個自由意志.
為什麼神.
還要繼續容許罪人.

$^{2241}$選擇.
還要堅持犯罪.
而不堅持去揀選神呢.
這個正正就是.
神繼續容讓一千年之後.
都給地上的人.
他還有一個選擇.
要不要跟隨神的權.
但是魔鬼.
實在太厲害了.
又是迷惑了很多的人.
剛才的第三點.
我可以補充一節經文.
可能看完這一段就完結了.
馬太福音的27章.
對了.
是很特別的經文.
當耶穌復活之後.
有這樣的描述.
詳細的不看了.
去到第50節.
耶穌大聲喊叫.
27章的50節.
大聲喊叫.
當然就是.
以羅伊拉瑪撒拜過大利.
斷氣.
我的神為什麼要離開我呢.
然後就乘了.
斷氣.
忽然電裡的萬子.
就是至聖所和聖所.
那個莫萬.
記住是由上到下.
列為兩半.
不是人撕.
人撕是從下面撕上去.
由上面到下面撕.
斷開的.
而且猶太人就認為.

$^{2281}$沒有人可以撕爛的.
除非你用很大的力量才可以.
然後就下列為兩半.
就是說.
至聖所的門進入.
可以打開.
地又震動 盆石又崩裂.
墳墓都開了.
留意這裡.
以遂聖徒的身體多有起來.
到耶穌復活以後.
他們從墳墓裡出來.
進了聖城.
向許多人顯現.
不是那些喪屍.
是那些復活了的聖徒.
這是很特別的描述.
耶穌的復活.
似乎那些.
勁的光.
可以令到墳墓都打開.
可能就是附近的墳墓.
耶路撒冷.
在橄欖山的墳墓.
全部向東的.
向著聖殿.
大家去過就知道.
全部向著聖殿.
他們等待那個復活.
似乎他們真的首先.
在那裡復活.
到時應該同步.
古往今來的人.
任何地方.
無論你在海南島.
或是在哪裡都會一起復活.
不會有次序.
真正的次序.
就是善和惡的次序.
很特別的經文描述.

$^{2321}$在當時也有.
這樣的短暫的應驗.
耶穌的復活也帶來周邊的墳墓.
都打開了.
死了的聖徒.
多又起來.
到耶穌復活的時候.
他們就從墳墓裡出來.
也就說是印證著.
有一個肉身的復活.
的一個階段.
當然這是一個很短暫的.
他們怎樣的結局也沒有人知道.
就是這樣描述出來.
突然間有一班人復活了.
走到聖城.
向人顯現.
似乎後文下文所說的.
就是那些人都不相信.
繼續不相信.
有主的復活.
有些印證他們都不相信.
這裡就間接.
印證了.
如果這裡也有一個尚且的復活.
有一個生命.
何況將來.
神的大能.
使主耶穌和他的門徒.
在一個劃時代的復活.
開始的時候.
就是千禧年的一個開始.
就是所有.
聖徒所有行善的.
就復活了.
所有聽到耶穌和天使長.
聲音的就要復活了.
這是第一次的復活.
就是啟示錄的20章.
所描述的第一次復活.

$^{2361}$接著就是.
一千年過去之後.
就有第二次的復活.
其實就是第二次的死.
就是死了一次.
就是正常肉身的死.
第二次的死.
就是永恆火狐的死.
就是第二次.
沒有復活的人.
沒有第一次復活的人.
第二次進入火狐.
所以復活之後死了.
那些人就很慘了.
這個就是後遷的解釋.
其實解釋後遷.
就差不多等於將.
有千禧年.
作為一個平衡的對立.
這樣就對了.
下一個星期.
我們就說的.
可能是千禧年前.
再加一些的補充.
還有我們會看.
後遷的解釋.
就算是應該.
都夠我們了解的了.
當然最終你要怎樣理解呢?.
這個是神學的理解.
和真理是沒有直接關係的.
所以.
你一定要相信有復活.
這個才是最重要的.
有永恆的復活.
和你一定有份於神的復活.
那個永生.
和那個生命.
這個才是真正的.
重要的重中之重.

$^{2401}$我們一起簡單祈禱.
感恩天父.
讓我們透過體末細論.
又清楚一些經文.
我們又知道.
原來我們真的有一個屬天的國度.
不論我們是信有千禧年.
還是沒有千禧年.
但是我們知道.
我們會有一個新天新地.
是永恆的居所.
是真正存在在這個宇宙.
雖然我們不能夠了解.
是一個物質的宇宙.
還是將來的屬靈的宇宙.
但是我們知道.
做物主就是為我們預備了.
這個救恩.
保守我們常在你的恩典裡.
保守我們常存.
這個感恩的心.
常存一份服事主.
因為我們有份於.
神國度裡面的那個服事.
我們這樣感恩禱告.
奉耶穌基督的名.
阿們.
\newpage



\section{}
\label{sec:tl2Gc71LinQ}
\textbf{2024-09-26 周四晚查經班《淺談末世論(3)》 講員:曾裔貴牧師}
\newline
\newline
連結: \href{https://youtube.com/watch?v=tl2Gc71LinQ}{\texttt{https://youtube.com/watch?v=tl2Gc71LinQ}} ~~~~ 語音日期: 2024-09-27
\newline
\newline
\hyperref[sec:a8KnWbij7Uk]{\small{< < < PREV SERMON < < <}}
~
\hyperref[sec:index_chronic]{\small{[返順時目]}}
~
\hyperref[sec:index_scriptual]{\small{[返順卷目]}}
~
\hyperref[sec:H4W4IN0LZRY]{\small{> > > NEXT SERMON > > >}}
\newline
\newline
$^{1}$我們一起祈禱吧.
我們可以再一次在晚上安靜的環境.
也都是在家中,在線上.
我們可以學習聖經的真理.
我們很盼望我們看的這些神學的論述.
關乎到千禧年.
我們都能夠掌握相關的經文.
提出這些見解的神學家.
都是有他們的心意.
都是用心去研讀.
但都讓我們能夠對經文.
我們有一個全面和客觀的分析.
以致我們不會將自己的心意.
讀入聖經裡面的時候.
就會將聖經有少許的扭曲.
願神幫助我們.
今天晚上我們再一次等候.
求主給我們身體,心靈都建壯.
更加我們盼望在線上.
成為了錄音的檔案.
可以有更加多人有興趣去聽一下.
因為都是大量看很多聖經的時候.
對聖經就不知不覺都有一份熟悉.
以致都能夠有一個對神的渴慕.
願主幫助.
我們這樣祈禱.
奉耶穌基督的名.
阿們.
上星期我們都有第二次.
講到關於這個毛遷熙年.
他們的主張大概有六個方面.
都跟大家講了.
那就不重複了.
我們就看第二個.
就是侯遷熙年.
所以我們最後才會看遷熙年.
遷熙年和毛遷熙年.
其實是剛剛的一個對比.
那侯遷為什麼又會有遷熙年呢?.
那侯遷和前遷到底有什麼分別呢?.

$^{41}$那個根本的分別.
最簡單的定義.
就是遷熙年的前派.
就是說主耶穌降臨.
接著就開始遷熙年.
主耶穌降臨.
就是撒旦被捆綁在無底坑一千年.
復活的聖徒.
就是說未信不信的人.
是不會復活的.
這個是有遷熙年的主張.
就是不是所有的人都復活的.
這個就是最根本的分別.
但是侯遷呢.
我再補充多一點.
就是有遷熙年.
就是一個悲觀的世界觀.
但是侯遷熙年.
就是一個樂觀的世界觀.
這是一個根本性的分別.
接著一路去到一千年完結.
就是主耶穌真正的第二次降臨.
就是說第一次降臨.
就是已經是釘十字架那一次了.
第二次就是.
耶穌在天上回來.
接著就是新天新地了.
這個就是侯遷.
所以其實都容易去界分的.
但是無遷就根本不會再跟你算一千年.
當然困難的地方就是.
怎麼從什麼時候開始去算一千年呢?.
然後主耶穌回來呢?.
這個就是一個比較難去回答的問題.
或者他們都不嘗試去釐清這件事.
因為他們又是在啟示錄的20章那裡.
既然你說了一千年.
所以他們就不是用一個實際的數值.
去看那個遷熙年.
而是反過來他們是根據幾段聖經.

$^{81}$就是說侯遷的時候.
那個世界是趨向於越來越基督化.
越來越全世界不同的民族.
都是被基督來統治.
所以就叫做遷熙年.
所以不是一個物質的.
你可以計算的.
一個更新了的地球.
耶穌在地上跟他的聖徒.
復會的聖徒來享受一千年.
所以就遷熙年前半.
你可以說是最需要神職的.
因為你要記住那個次序.
當然我們下一課講的就是遷熙年.
詳細講他們的見解.
不過始終有對比.
遷熙年是最難的.
就是耶穌回來.
跟著那些所有死了的人復活.
信的,在生的都要身體改變.
就開始一個正式的遷熙年在地上.
撒旦就被捆綁.
是這樣一直遷熙一千年的地球的年份.
這樣去計算.
接著才是耶穌再跟他的聖徒第二次回來.
就是新天新地了.
或者說才是那個大審判.
白色大寶座的審判.
這個就是遷熙年前派的.
應該說叫遷熙年派.
有遷熙年派.
因為再講前遷.
就牽涉到那個大災難.
就跟那個大災難有關的那些經文.
但是後遷和無遷就不相信有這樣的大災難.
世界是趨向越來越好.
好到完全被基督統治.
最後基督就降臨了.
所以我們看後遷有幾段重要的經文.
來看一看.

$^{121}$你就會知道他們是這樣理解經文的.
我們先看第一段.
三段都是來自馬太福音的.
馬太福音的第28章.
28章的最後那段經文.
因為他們提到的就是基督的那個管治.
大家其實很熟悉的了.
是大使命的經文.
不過後遷的支持者就認為這段聖經.
是講到基督的國度是由他升天之後.
一直就是這樣來到他.
越趨美好.
越趨於基督化.
整個世界每一個民族.
第16節這一段.
是最後的叮嚀.
十一個門徒.
往加里尼去.
去到耶穌約定的山就是鑒欖山.
他們見了耶穌就拜他.
然而還有人疑惑.
耶穌進前來對他們說.
天上地下所有的權柄都似給我了.
所以你們要去.
使萬民作我的門徒.
奉父子聖靈的名.
給他們施浸.
凡我所吩咐你們的.
都教訓他們遵守.
我就常與你們同在.
直到世界的末了.
他們這樣理解的這個末了.
就是他第二次回來.
就是說整個的.
因為十八節說.
天上地下所有的權柄.
都已經似了給耶穌了.
所以門徒.
是去好像攻城掠地那樣.
福音是一直帶來.

$^{161}$只會得勝.
只會每一個民族.
因為是向萬民作見證.
來帶領人歸向神.
透過洗禮.
將他歸入神的名下.
還有他們會遵守神的誡命.
一直到世界的末了.
就是主的第二次回來.
所以他們就這樣理解.
這段經文.
其實就是講後天的那個理解.
但是他們就不以那個作為.
那一千年.
作為那個物質的計算的數值.
他們只是會說.
由主耶穌得到權柄之後.
那個世界是趨向越來越基督化.
越來越基督教化的.
世界是會越來越美好的.
這個是後天其中是較為樂觀的信念.
這個就是他們的主張.
而繼續有的.
另外就有兩段聖經.
就是在《馬太福音》第十三章.
稱為天國的比喻的其中兩個.
就是「界賽種比喻」.
和「面教比喻」.
是一對的比喻.
因為好種和敗子又是另外一對的.
第31節.
13章總共有7個比喻.
或者嚴格來說叫做三對半的比喻.
因為是一對一對的講的.
這裡告訴我們.
第31節.
他又設一個比喻對他們說.
對那些群眾.
因為當時他在船上.
天國好像一粒界賽種.

$^{201}$有人拿去種在田裡.
這原是百種裡面最少的.
真是一粒很小的.
我們在聖地都摸過.
那些種子是很小.
就是界賽種來的.
所以我們不要想我們自己種的.
吃的那些界賽.
不是那個意思.
完全是一棵樹來的.
等到它長起來的時候.
卻比各樣的菜都大.
卻成了樹.
天上的飛鳥來縮在他的枝上.
接著再講另一個比喻.
他又對他們講一個比喻.
說天國好像面教.
有婦人拿來.
藏在三斗面裡.
直等全團都發起來.
似乎兩個或者這一對的比喻.
都是提到天國.
是一路會膨脹.
一路會大.
甚至大到一個地步.
是可以跟他原先的完全.
你不會察覺到那個變化.
而是大到一個地步.
是完全出人意表的.
好像一個面教.
又是這樣發得很大.
這一對的比喻.
是特別強調天國的無限擴大.
所以正正就是後遷的信徒.
或者當然的神學家.
提出這一對比喻本身就是說.
天國是一路的擴展.
擴展到一個地步.
是一個大到一個這樣的狀態.
結合剛才說的28章的馬太福音.

$^{241}$就是說後遷的人.
他們就認為基督的得勝.
基督的得到天上地下所有的權柄之後.
那個的開始.
就是基督教化在地球上.
一路得著不同的民族.
不同的人.
這樣來歸向神.
所以是一路的會長大.
一路的這樣來成長的.
根據就是這三段聖經.
大家可以想得到.
這三段聖經能不能夠.
這麼肯定地.
可以承載一個這麼美好.
越趨完美的一個世界呢?.
基督教是不是真是我們想象的.
由基督得勝之後.
升天之後.
一路整個宇宙.
先不要說宇宙.
說地球的人類.
是一路向基督.
向福音來降服.
來被基督的權柄征服.
來得勝.
所有的國家都是基督教呢?.
相信我想小朋友都知道不是.
所以後遷.
他們經歷了兩次的世界大戰.
基本上都沒有什麼市場.
因為世界原來這麼醜陋.
這麼殘酷的.
人與人之間.
國家與國家之間.
可以去到世界性的戰爭.
聖靈塗炭的.
怎麼會世界越趨越美好呢?.
而且是基督教的國家.
與基督教的國家打仗.

$^{281}$尤其是第一次世界大戰.
你可以想像得到.
那個問題不是我們.
這樣的解釋聖經.
就等於我們可以套用下去.
仍然那個始作俑者.
就是他們要解釋.
啟示錄的20章.
那個經文.
但是又不能夠完全的接受.
為什麼基督要降臨在這個地球.
一個改變了的地球.
然後開展在地球上的一千年的希年呢?.
他們覺得那個希年.
其實現在是在天上的管治.
那個天上的管治.
藉著基督的德性.
他領受的權柄.
和交付給教會.
教會的歷世歷代.
是一路帶來地上的人.
越趨基督化.
越趨基督教.
如果我們用美國做一個例子.
美國繼續說自己是基督教的國家.
現在輸出的是福音.
宣教士.
我們還聽不聽到有很多很美麗的宣教士.
在21世紀.
有的 摩門教就有了.
摩門教是定期的出來的.
除了摩門教之外.
基督教有沒有宣教士大量的來.
被輸出呢?.
年輕人受感動去到世界.
去到有這麼窮得這麼窮的國家去傳道呢?.
19世紀大部分是美國的.
宣教士 傳教士.
大部分是美國.
甚至乎是美國南部的教會.

$^{321}$派遣宣教士.
我們宣道會.
都是宣信博士.
19世紀末期.
大復興的時候.
他就開那個宣教學院.
只能夠給你兩年的授訓.
年輕人.
接著你們自己按感動去.
某程度是沒有支援的.
你去到那些地方.
就是要面對當地的困難.
面對那些天氣.
疾病.
語言 文化.
所以很多宣教士的子女是死的.
如果你看回.
你有機會去到那些警察會.
去看那些記錄.
真是逐筆的記錄.
你會看到有很多.
記下那些宣教士和他的家人.
離世 有些是殉道.
有很多根本不是殉道的.
根本頂不住那些疾病等等.
已經是離開世界的.
其實很多的.
所以你可以想像得到.
宣教士去到這些窮苦的地方.
都有很多很多的掙扎.
去印度 非洲 中國.
這三個黑色大陸.
完全沒有福音.
你可以想像得到.
那個時候有這麼多的宣教士.
去改變幾個國家.
幾個大洲的文化.
那個信仰.
需要經過多少的年日.
但是現在越來越方便.

$^{361}$越來越先進.
各樣東西都好了.
但是還有沒有大量的宣教士.
美國輸出呢?.
有的是荷里活電影.
有的是香煙.
有的是毒品.
有的是美國文化的放淫放縱.
你可以想像得到.
還有沒有那些19世紀末期的.
那些宣教的可歌可泣的故事呢?.
可能有 但是零零星星.
因為連他自己本土.
那些靈命都已經很軟弱.
很軟弱何來輸出呢?.
何來有這麼多的宣教士呢?.
歐洲就更加沒有.
你可以想像得到.
我們要問的這個很直接的問題.
侯遷就是非常正面的樂觀主義者.
這個世界因為基督的德性和管治.
他得到的權柄.
就是會令到整個世界越趨美好.
這個就叫做千禧年了.
所以是一個概念性的千禧年.
不是一個物質的可以量度計算.
如果你是一個電子錶.
你可以一千年都還能走的.
他就由按鍵開始.
一千年基督就在地上了.
這樣的千禧年.
所以有千禧年的弟兄姊妹.
其實就是這麼悲觀.
要等到基督降臨.
我們才能享受這個千禧年.
所以基督還沒降臨之前.
其實是非常之苦難的地球.
是非常之敵黨福音的地球.
因為撒旦在地上.
撒旦什麼時候不在呢?.

$^{401}$唯有就是千禧年他才不在.
所以千禧年和侯遷的根本分別就是.
悲觀主義和樂觀主義的分別.
所以我們就要看這些經文.
我們怎樣去反駁這些經文呢?.
我們要看一看《馬太福音》.
也是《馬太福音》.
好幾段都是這一段經文.
給我們一些啟發.
第七章.
這些真是耶穌的教導.
祂來告訴我們.
13到14節.
你們要進窄門.
因為引渡滅亡.
這是第七章.
所謂登山寶訓的經文的教導.
那門是寬的.
去到滅亡是寬的.
路是大的.
進去的人也多.
引渡永生.
那門是窄的.
路是少的.
找著的人也少.
所以某程度是很悲觀的.
因為原來.
當然這裡沒有說到.
什麼時候會發生窄和闊的對比.
這個對立.
是在說人類對福音的.
如果這樣理解就對了.
人類對福音的興趣.
就是原來神已經設計了.
進永生的門這麼窄.
路又少.
找到的人又少.
很自然就變成了小眾的玩意.
永生.
當然我們看這段經文.

$^{441}$為什麼要製造這些麻煩給人呢.
但是我們要全面看其他經文.
其實就是人自己拒絕福音.
拒絕去選擇.
合神心意的生活.
和盡快認識神.
解決自己一生的罪.
這些就是指的.
窄和少的真正背後定義.
不是神刻意設計一個窄的.
甚至乎沒有門牌.
找不到的.
我想不是這樣的物質理解.
而是心靈上對永生的渴求和羨慕.
一直在說的就是.
對這位主耶穌.
祂來到這個世界.
祂告訴我們.
要怎樣的新倫理觀去看和神的關係.
如果是這樣的話.
連猶太人都不可以得救.
整個段落其實根本就是針對猶太人.
如果你還憑著.
以為靠你自己的肉身.
就可以進入永恆.
其實是不行的.
你以為你自己是猶太人.
已經是順理成章.
就好像哥哥在徐小的.
妹妹就肯定可以進入徐小.
根本不用填的.
你只是報了你自己.
當年的學籍就行了.
大家明白我的意思.
因為這是法例.
就是說.
讀書不讀書都可以進入徐小的.
你多厲害.
如果你哥哥姐姐都不在徐小.
那你就要競爭了.

$^{481}$那你就扣門.
你就搞豬.
好不好彩.
但這裡告訴我們.
如果是這樣的時候.
就很端乎我們怎樣去理解福音.
為什麼我們有些人得到.
有些人得不到.
很肯定就是對福音的認識.
對福音的認識的大前提.
就是對自己的心靈的認識.
自己心靈裡面的空虛.
我們在成長裡面.
不斷用不同的物質去填滿.
用愛情,用事業.
用世上的享樂去填滿.
誰知道到頭來都是填不滿.
唯有神才可以成為我們真正的滿足.
所以這就是澤的真正定義.
不是神故意令大部份人都不行.
而是這個救恩.
神願意萬人得救.
萬就是ALL.
所有人明白真道.
所以願意人得救.
但需要人明白真道.
選擇神.
這個選擇的信心.
基本上不是條件.
是每個人都可以得到.
所以如果得不到的.
只是一個結論.
就是人道滅亡的門是寬.
路是大.
進去的人也多.
這句說話其實是很可圈可點的.
告訴我們.
我們千萬不要灰心.
因為大部份的人.
其實都是喜歡拒絕神的恩典.

$^{521}$或者不會因為耶穌的救恩.
看為是他必須得到的恩典.
因為他有的.
是神給他的一生的可能是富足.
可能是成功.
可能是有很多他可以運用的資源.
以致他浪費了大部份的人生.
就在這些物質的享用上.
而忘記了次物質的這位生命的主.
他根本從來都沒有興趣和他交往.
所以我們基督徒.
我們是和這位次恩典的主交往的.
所以那些物質是恩典.
所以我們不要抱著那些恩典.
就忘記了次恩典的主.
這是一個很重要的情況.
所以這裡就告訴我們.
透過這段經文.
都讓我們看到.
似乎基督教不是一個大眾化.
是越趨於越來越統治整個世界的一個宗教.
只是告訴我們.
得到福音的人.
只會是真正生命改變的人.
真正願意放下他的想法.
他的哲理.
他的價值觀.
他願意放下符合神的標準的澤門.
符合神的標準的得救.
他才可以進去.
如果不是就進不去.
你們基督教這樣說.
不知道你們有沒有聽過這些經文.
這些人這樣說.
是你們這樣說.
你們認為這樣的永生.
我不是這樣認為.
沒有問題的.
因為人是有理智的.
是有理性的.

$^{561}$我們也不能強迫人改變人.
但是最終他自己要看到.
神的真實.
和他真是心靈打開.
他遇見神.
他會有生命的改變.
所以那個澤和那個闊.
是人的心什麼時候打開的.
一打開.
原來就很寬闊的了.
原來神的恩典.
原來一直都存在.
在那個人的人生裡面的.
所以這一段.
我們來到去.
也不能叫做駁斥.
就是說其實太樂觀了.
基督教自從耶穌得勝.
升天.
就一直像所向無敵一樣.
太樂觀了.
但是馬太福音自己本身.
也是耶穌親口說的.
信主的人.
其實是極之少數的.
接著我們又再揭一揭.
路加福音第十八章.
這些經文某程度大家可能讀過.
不過你沒有這樣想過.
其實應該這樣理解的.
十八章的.
我們看長一點.
是由第一節開始的寡婦和不義的官.
這個比喻來的.
十八章.
因為耶穌說完這個比喻.
解完之後.
耶穌有一個慨嘆的.
他說當人子來的時候.
其實地上.

$^{601}$不是真的每個人都有這個信心.
耶穌說.
設一個比喻.
是要人常常禱告.
不要灰心.
說某城裡有一個官.
不懼怕神.
也不尊重世人.
就是一個不義的官.
不懼怕神.
通常如果那些不尊重世人的.
他也會懼怕神.
他也會拜神求神問卜.
很誇張.
又不懼怕神.
又不尊重世人.
他不尊重的意思是秉公行義.
審判.
因為他是官.
第三節.
那城裡有一個寡婦.
常到他那裡說.
我有一個對頭.
求你給我伸冤.
正正常常的.
擊鼓鳴冤.
就是說你開工審訊.
做一個判斷.
這個很正常的.
尤其是寡婦.
她也不是在當時有什麼社會地位.
就希望有一個法律的途徑.
去得到一個合理的仲裁.
但是第四節.
他多日不準.
後來心裡說.
我雖不懼怕神.
他自己這樣說.
耶穌的比喻裡面這樣描述.
這個人知道自己不懼怕神.

$^{641}$又知道自己不尊重世人.
真的好像貪官污吏.
但是也不會說出口.
但是比喻裡面.
耶穌說這個官是意識自己的狀態.
你很難會說出口.
我不懼怕神.
我又不尊重世人.
不過只因這寡婦煩擾我.
我就給她伸冤吧.
免得她常來纏磨我.
就是磨爛隻這樣磨.
纏磨.
原來真的有這樣的成語.
又纏你.
死纏爛打.
磨就是一直這樣磨著你.
不讓你走.
為什麼會被一個婦人.
寡婦這樣纏她.
她就會心軟.
她就會有良心發言.
比喻強調的是.
第一節耶穌說.
叫人常常禱告.
不要灰心.
說了一個極端的例子.
極端到不敬畏神.
不懼怕神.
不尊重世人的官.
根本不會理會寡婦.
隨便找她的士兵.
或者那些衙役.
打發她走就算了.
但強調寡婦天天纏磨她.
要她伸冤.
她好像過意不去.
良心發言下.
不是趕走她.
而是幫她伸冤.

$^{681}$原來一直都知道誰對誰錯.
現在終於幫她審.
這樣的一個比喻.
其實重點不是伸張正義.
是免得她再來禁鐘.
又在激鼓.
很奇怪的一個比喻.
耶穌說了這麼極端的東西.
似乎當時不會有這樣的官.
猶太人一定是敬畏神.
一定會跟法律去做.
所以保羅說自己是猶太人.
本來猶太人要打多少下.
40下 鞭打40下.
因為你是猶太人 打39下.
有1\%的優惠.
真的 哥林多後書.
哥林多後書那裡這樣說.
鞭打.
可想而知 原來猶太人對保羅.
都會跟律法去做事.
你就可以想像到.
在這樣的比喻裡.
很難會出現這樣的極端.
一個婦人根本沒有社會地位.
一次兩次不行.
你一定會放棄.
你就會錦懷身世.
誰叫你是寡婦.
誰叫你沒有社會地位.
寡婦這樣守著一定會被人欺負.
在你的村裡.
這些我們都見過.
我們小時候成長的時候.
都見過這樣的家庭.
所以這樣的極端的比喻.
其實是想說那個信心.
持之以恆的信心.
這麼極端.
接著第六節.

$^{721}$主耶穌說.
聽不義的官所說的話.
神的選民晝夜呼籲他.
他縱然為他們忍了多時.
豈不終究給他們新冤嗎?.
是要不斷忍耐的.
原來地上一直沒有公義.
沒有絕對公義.
耶穌還沒回來之前.
地上原來是很多很多的冤.
還沒申的.
但是神的時間還沒到.
就是還沒申.
還沒申就要等候那個新冤.
就不要灰心.
經文的意思就是.
叫人常常禱告.
不可灰心.
是這個比喻的解釋.
所以耶穌就用這麼極端的比喻.
比喻裡面這麼極端的例子.
地上很少會發生.
有這樣的人 有這樣的官 有這樣的制度.
而繼續.
重點是繼續竊而不捨.
繼續死纏爛打.
繼續纏磨那個官.
所以這段聖經是想告訴我們第八節.
我告訴你們.
要快快的給他們新冤了.
然而人子來的時候.
遇得見世上有信得的人嗎?.
就是說原來耶穌基督清楚知道.
如果我們不持續的.
憑信心的繼續的懇求.
堅持去祈禱.
那些事情的成就.
這裡說的不是你想的.
你個人的那些物質.
這裡說的是對地上不義的新冤.

$^{761}$是對基督徒我們作為神兒女.
為著神國降臨的持續的祈禱和傳福音.
但是一直如果主未回來.
都是未真正的完善和圓滿.
所以這段經文.
又是可以去抵消後遷他們那些樂觀主義的主張.
因為耶穌已經說了.
他回來的時候.
世界上其實有信心持續去等候神國的.
是極之少數.
而且如果我們不繼續誠心的去禱告.
不灰心的去禱告.
沒有寡婦的決心去祈禱.
我們就更加沒有信心.
因為告訴我們.
那個世界不會讓我們看到那個美麗.
那個美麗其實在我們的信心等候裡面.
等候主真正的降臨.
而帶來的世界真正的改變.
所以這段經文其實是很悲觀的前纖.
由千禧年的悲觀的解釋.
當然整段聖經是說要我們信徒常常禱告.
不要灰心.
所以如果我們從開始.
其實我自己都真的有一個要審查.
我最近終於真的有一本祈禱本.
有些人叫我為他祈禱.
我上一次去講福堂.
講完道之後.
其實那個是福音聚會.
講完道之後有個弟兄在台上領唱.
原來是執事會的副主席.
他走過來跟我聊天.
完了之後就問我.
你可不可以為我爸爸媽媽祈禱.
因為我講的例子就是Terence的爸爸信主.
原居民信主和在醫院洗禮.
原來他的父母都是原居民.
就說很想我去探望他.
我就由八月講完道之後.

$^{801}$我就寫下了他的名字.
原來結下我就記得了.
我就為他祈禱.
真的會想一下.
什麼時候有些聚會可以邀請他.
因為住新界的.
叫他來聽聚會.
原來這本祈禱簿這麼重要.
是提醒自己.
你自己許下的承諾.
不只是祈禱手就打發人走.
原來真的有這樣的作用.
所以我發現現在我也想這樣去學習.
如果我們一直看到那個記錄.
而有改變.
其實就很驚訝.
其實這裡告訴我們.
原來上上禱告不可灰心.
原來就是告訴我們.
我們要有一本我們自己的祈禱簿.
要好像這個寡婦那樣.
去憑信心去到神的面前.
神肯定不是不義的官.
神肯定快快要伸冤.
但是快快的裡面.
原來還容許地上有很多災難.
有很多的敵黨神的勢力.
某程度上就是對後天樂觀主義的當頭棒喝.
一盤冷水.
就是說如果我們不堅持用我們的方法.
我們要持續的禱告.
還有信禱告.
很多基督徒已經不信祈禱了.
已經不信祈禱有效.
已經不會覺得.
回到祈禱會是一件很榮耀的事.
回到祈禱會的人才不是榮耀.
這樣想的.
我見到很多教會領袖是不回到祈禱會的.
寧願晚上排滿自己的節目都不回到祈禱會.

$^{841}$但是很熱心做領袖的.
不只是我們.
很多其他教會都是的.
所以沒有任何針對性的.
回到祈禱會好像是一些初信就對了.
那些初信要回的了.
我們資深的人不用回到祈禱會的了.
其實不要這樣想.
其實你應該想想.
我為神的緣故我分別一個小時出來.
和弟兄姊妹一起祈禱.
是一件榮耀的事.
很少人這樣想.
或者可能港台也很少這樣教.
但是這段聖經是耶穌親口用一個比喻告訴我們.
在我們人生這麼有限的人生.
你要面對的挫折.
面對那些潑冷水的禱告沒有效.
是經常出現的.
所以才令我們沒有效不祈禱.
如果我們一祈就馬上成就的.
每個人都回到祈禱會.
坐到滿了.
因為一進完星期三晚的祈禱會.
一出門就馬上成就.
那我不回來當然要回來.
帶著十件回來.
但是問題就是.
一年又是這樣 兩年又是這樣.
慢慢就 原來祈禱.
只不過變成一個例行公事.
既然例行公事.
不如星期三晚都不要有祈禱會了.
大家都很忙.
不如索性星期三晚沒有祈禱會了.
起碼知道原集是沒有的.
我因為看不到他們開燈.
因為他們現在都是用線上的.
是隔個禮拜的.
還是一個月一次的 我忘記了.

$^{881}$燈火通明的祈禱會已經沒有了.
很多教會都沒有了.
既然沒有人回來.
不如就不要有祈禱會了.
耶穌就說 正正就是沒有人.
就要更加鼓勵人.
設一個比喻.
是要人常常的禱告.
不要灰心.
就是說地上有太多.
令我們灰心.
令我們不需要感覺.
要祈禱的一個世界.
所以不是樂觀的世界.
是悲觀的世界.
是悲觀的阻礙.
禱告的那種光景.
不是你一帶回來就立刻成就.
就是祈禱很多年都未必成就.
那就更加要祈禱.
所以如果剛剛聽到的弟兄姊妹.
真的弄一本帥一點的簿.
自己寫下來.
一想起就拿出來.
就為這些人祈禱.
原來是對自己的一個承諾.
對自己的一個行動.
原來真的挺好的.
我又加了一個新的人名.
那麼盼望我們大家一起去學習.
第三段.
貼上來後書.
我們之前讀過的了.
好像越講越清.
好像那些痕的感覺好像少了.
好像暢通了.
貼上來後書.
我們讀過了這些經文.
現在我們集中在後遷這個階段.
這裡的經文是回應後遷.

$^{921}$在聖經裡面有更多經文.
比我們悲觀不是樂觀.
世界不是趨向基督教化.
是趨向更加的敵黨神.
後書的第二章.
我們講過有幾個段落.
頭這一段是講那個末後的日子的迷惑.
接著有些人會接受迷惑.
接著第十三節是另外一班人不受迷惑.
因為他們願意在聖靈裡面去生活.
所以是兩類的人.
第二章的第一節.
弟兄們,論到我們主耶穌基督降臨.
和我們道他那裡聚集.
相信肯定是在講主在來.
我勸你們.
無論有靈有言語.
有沒有我明的書信.
很肯定當時是由那裡開始.
由保羅的使徒時代開始.
就說主的日子現在到了.
但不要輕易動心.
也不要輕易驚慌.
不要有連結就按進去.
人不驅用什麼法子.
你們總不要被他誘惑.
因為那日子以前.
日子就是第一節講的基督降臨.
必定會有離道反教的事.
相當悲觀.
保羅.
並有那大罪人.
就是沉淪之子.
要顯露出來.
這個沉淪之子.
他是抵擋主.
高抬自己的.
超過一切稱為神.
甚至坐在神的殿裡.
自稱是神.

$^{961}$當然後遷的人就說.
就是主後七十年.
羅馬的皇帝.
他們這樣來抵擋神.
如果這樣局部應驗.
就抵消了經文的意思.
每個時代都有這一類.
抵擋神的人出現.
但似乎不是.
泛指一個過去式的.
一個抵擋神的人.
例如尼祿.
例如其他一些君王.
或者近代有人說希特拉.
諸如此類等等.
但是這些經文.
應該是指向真正的幕後大罪人.
那個迪基督.
所以.
第五節.
我還在你們那裡的時候.
曾經把這些事告訴你們.
你們不記得嗎?.
現在你們也知道.
那攔阻他的是甚麼.
是叫他到了的時候.
才可以顯露.
因為那不法的引義已經發動了.
只是現在有一個攔阻的.
等到那攔阻的被除去.
似乎有些人解釋的就是.
教會被提.
那時這不法的人就必顯露出來.
主耶穌要用口中的氣滅絕他.
用降臨的榮光廢掉他.
這個不法的人照撒旦的運動.
就是說原來他的內心.
他整個的神學.
整個的內心裡面.
都是敵黨.

$^{1001}$都是撒旦的運動.
還有一些超自然的.
行各樣的異能神蹟.
和一切虛假的歧視.
一直對照著.
啟示錄的十七章的.
那個基督的出現.
第十節.
並且在這個沉淪的人身上.
行各樣出於不義的詭詐.
就是令到他.
是充滿了可以令到人信服的.
但是虛謊來的.
是詭詐來的.
因為因他們.
就是那些跟隨這個大罪人的人.
因他們不領受愛真理的心.
使得這些人不救.
這些人就是沉淪的人.
所以沉淪之子.
滅亡之子.
那個大罪人.
他會帶領一大班人.
跟隨他而沉淪.
第十一節.
故此神就給他們.
一個生發錯誤的心.
就是說原來每個人都有思想的.
但是有些人的思想.
原來傾向於撒旦魔鬼.
和那個敵黨的.
會追求這些異能神蹟.
虛假歧視的.
因為為什麼呢.
原來他們不領受一個愛真理的心.
其實每個人都不愛真理的.
就算我們信主.
都是的.
但是很奇怪很奇妙的.
我們打開了心.

$^{1041}$接受耶穌之後.
聖靈就改變我們了.
越來越愛慕真理的心.
相信在座的.
都已經是一把年紀了.
信主內.
我們給初信的時候.
雖然有些事我們還未行得到.
但是越來越真之確信.
越來越沉澱.
我們是真的有永生.
進入永生的一班人.
所以大部分的基督徒.
如果在這些年紀.
整個心靈.
知道敬畏神.
是過於一切的.
是重於一切的.
如果沒有了神.
其實什麼都沒有.
這個肯定是我們心底的說話.
但是居然有一班人.
他們都是在追求真理.
但是原來是一個生發錯誤的心.
什麼都信唯有不信真理.
這樣叫做生發錯誤的心.
所以叫他們順從虛方.
結論就是十二節.
使一切不信真理.
喜愛不義的人.
都同樣被定罪.
所以這裡告訴我們.
大罪人會顯露出來.
會對一大部分的人.
似乎是走去闊路.
很闊門的人.
就會向這些人招手.
所以就得不到一個真正的.
真理在他們的內心.
和剛剛十三節另外一班人.

$^{1081}$就另一個世界了.
所以保羅說你們帖撒羅尼迦.
就是一班愛真理的心.
就是被神揀選.
叫你們因信真道.
所以揀選和你去聽明聖經.
是沒有抵觸的.
你不要賴我都沒有被揀選.
所以我就不會聽.
所以我去福音聚會.
我是不需要去的.
去我都是睡覺.
我去我的心都是不斷地去搏他.
因為他的理由就是說.
我都沒有被你的神揀選.
所以我就信不了耶穌.
所以我就這麼多的問題.
我就這麼叛逆.
其實不是的.
其實那個揀選裡面.
那個人明白真道.
好像似乎.
又不是50/50.
是神的絕對揀選.
和人的打開心靈的100\%的相信.
其實是沒有抵觸的.
就不是200\%.
很特別的.
我們每個人都感恩.
今天我拒手絕志.
因為我是被神揀選的.
從來沒有這樣的人.
因為他根本都不知道什麼叫揀選.
我們才明白.
我們才知道原來我們信主的人.
是被神揀選的.
是被神嚴重的.
在創立世界以前.
就已經揀選我們了.
所以.

$^{1121}$兩類人.
不過他現在告訴我們.
第二章的經文的脈絡.
是叫我們很小心.
我們要保持自己.
在這個常在真理的行列裡面的那班人.
但是因為末後.
是會有那個大罪人顯露出來.
現在只不過是有一個攔阻他.
等到這個攔阻被除去.
我自己都這樣理解.
就是教會被提了.
這樣才可以.
這個大罪人就會出現了.
所以這個就是.
我們要了解.
真正的背影.
對神聖的真理的抵擋.
是一直都存在的.
而且是存在的.
是很誇張的.
所以.
千禧年永遠是悲觀主義者.
不是樂觀主義.
還有一段.
提摩太后書.
其實我們一會兒再看24章.
就更加悲觀的悲觀了.
因為耶穌說到.
將來末世地動天搖的情況.
提摩太后書其實都是這樣的氣氛.
三章.
提摩太后書的三章.
提摩太說.
這個時候真是最後的遺言.
最後的叮嚀.
因為他已經說.
我離世的時候已經到了.
提摩太后書的三章.
提摩太你該知道.

$^{1161}$末世.
必定有危險的日子.
來到.
因為那時.
人要尊顧自己.
貪愛錢財.
自誇.
狂傲.
自誇和狂傲.
相信都是孿生兄弟.
對自己不斷認叻.
就會不斷的.
狂妄的驕傲.
驕傲就是我們.
什麼都看不起別人說什麼.
就是自己永遠都對.
就算你自己沒有那個知識.
都覺得自己都對.
那就叫做狂傲.
接著就蚌毒.
誹謗和.
這個字都很深很文言.
和言語有關.
就是說你的內心充滿了自己.
你很自然就.
目空一切.
很自然就會對其他人.
很多的批評論斷.
違背父母.
如果父母都違背.
自然就忘恩負義.
所以是很多的東西.
是全部有關連的.
是一個.
好像一個心靈.
這個世界.
只照到心臟.
你的血管有沒有塞.
照不到那個心靈.
但是原來2000年前.

$^{1201}$已經有一個可以讀心的.
一個的.
工具.
就是聖經的真理.
就是神的啟示.
他將人性的內心剖析出來.
就是危險就在這裡.
記住.
不是那些.
孫悟空那些.
八怪那些.
是正正常常的每一個時代的人.
而且可能是.
社會上出類拔萃的人.
是正正常常穿西裝打領帶的人.
要記得.
可能是有大學畢業.
可能哈佛的.
可能是高材生.
忘恩負義.
這期就最.
最熱的就是.
關欣娜的家族增產.
又放報紙.
所以不是什麼批評 報紙已經賣了.
就是她的.
不知道是婆婆還是什麼.
不給她們幾個女兒.
就要說這個為足無效.
打官司.
那就.
是不是會跟這些有關呢.
我就不知道了 違背父母.
忘恩負義.
心不勝結.
就是不敬虔.
自然就沒有親情.
不解怨 你都沒有解的.
是不肯去解除那些的冤仇.
就是說我生你氣到死.

$^{1241}$生你氣到見耶穌.
我都不會跟你解怨的.
這樣的意思.
就是不會set loose.
不是說沒有解的.
是不去解除.
你想想.
這個很嚴重的.
因為這個將來神會審判的.
不要帶著一個.
跟人的結怨去天堂.
是去不了天堂的.
這個問題很嚴肅的.
所以我們要盡力與人和好.
好說殘言.
不能夠自虐.
就是self-control.
已經失控了.
接著就會盛情空步.
因為你不自虐就空步.
不愛良善.
保羅是不是在說提摩泰.
在教會裡面.
一回到禮拜堂.
看看是不是這個這樣.
每個人都是這樣.
不愛良善賣主賣友.
任意妄為.
自高自大.
愛燕落不愛神.
有敬虔的外貌.
似乎開始說不只是教外的人.
教外的人沒有意思.
還沒信耶穌你說什麼.
還沒信耶穌每個人都是這樣.
信了耶穌之後.
仍然有很多態度還沒轉變.
慢慢就變成.
根本不覺得需要轉變.
這個也是基督徒要留心.

$^{1281}$有敬虔的外貌卻背了敬虔的實意.
這等人你要躲開.
牧者不是應該要接觸嗎.
要躲開這些人.
哪偷盡人家牢籠.
無知婦女的正是這等人.
這些婦女.
突然間說這些婦女.
擔負罪惡被各樣的私慾引誘.
常常學習.
終究都不能夠明白真道.
接著就說一個例子.
這些在.
創世記沒有出現過.
出埃及記也沒有出現過.
這些人.
從前雅尼和楊比.
怎樣敵黨摩西.
似乎是那些婦女.
其中的兩個表表者.
這等人也怎樣敵黨真道.
怎樣敵黨摩西.
這樣就很多了.
出埃及記.
和民數記就很多例子.
在馬拉的水有苦.
沒有嗎哪.
每餐都吃嗎哪.
接著去到一些地方.
最終的結局.
就是準備入迦南之前.
要去探路.
接著傳媒眾.
都說要回埃及.
所以其中有兩個敵黨摩西.
即是說.
這等人.
就是提摩泰時代的人.
都是這樣敵黨真道.
原本的原則就是.

$^{1321}$心地壞了.
在真道上是可廢棄.
好像我們見到一些船拋錨.
一些車拋錨.
即是壞車.
然而他們不能夠再這樣敵黨.
因為他們的愚昧.
必在眾人面前顯露出來.
像那二人一樣.
提到是.
摩西的光景.
就提到.
他們的愚昧.
必在眾人面前顯露出來.
像那二人一樣.
提到是.
他們的愚昧.
必在眾人面前顯露出來.
像那二人一樣.
提到是.
他們的愚昧.
必在眾人面前顯露出來.
像那二人一樣.
提到是.
他們的愚昧.
必在眾人面前顯露出來.
像那二人一樣.
提到是.
但你已經服從了我的教訓.
品行自向信心 寬容愛心 忍耐.
所以你會看到.
有神敬畏神.
走擇路的基督徒.
就是第十節開始.
所以我們不要欣賞羨慕那些.
他不回教會.
繼續喜歡怎樣就怎樣.
喜歡怎樣就怎樣.
將來他就要面對他自己的人生.
保羅現在就是.

$^{1361}$完全扭轉了360度.
對提摩太的警告.
勸勉.
和提醒他.
但你已經服從了我的教訓.
即是你不要跟那班人.
反過來你更加殷勤地去追求.
品格的成長.
品行自向信心 寬容愛心 忍耐.
原來要不斷的探索自己的生命.
探索自己的性情.
靈性.
以及不只是你自己做好.
原來有逼迫苦難.
令我們更加要學習忍耐.
我在安提摩.
我在以哥廉.
在路斯德.
這是第二次的傳道.
在路斯德被人用石頭打.
在路斯德呼召提摩太.
我所遭遇的逼迫苦難.
我所忍受的是.
何等的逼迫.
但是你要學習了.
這一切的苦難裡面.
主耶穌將我救出來.
即是說原來.
際遇是使我們.
遇見主 是使我們的.
信心成長 是使我們的.
生命成熟.
所以基督徒沒有辦法逃避.
苦難的 逃避不到的.
如果你一面對.
苦難就說我不再信你的耶穌了.
原來令到我這麼難.
這麼慘的 我不信你的耶穌了.
那你就不能成長了.
保羅說連我是使徒.

$^{1401}$我都經歷了.
在傳道裡面.
是你親眼目擊.
見到的.
我在路斯德 在以哥廉.
在安提摩.
我所受的苦難.
是何等的逼迫.
不過在這些逼迫裡面.
原來有主和我同在.
原來有主耶穌.
時候未到.
他將我解救.
不單止是這樣 十二節.
不但如此 凡是納至在.
基督耶穌裡面敬虔道日的.
也都要受逼迫.
即是說我們站在.
真理這一邊.
我們就要受逼迫.
只是作惡的和迷惑人的.
必越久越惡.
他欺控人也被人欺控.
因為又是一個.
生發錯誤的心.
不愛真理的心.
所以.
越聽聖經.
只會越更加的.
靈性的明白.
願意去學習.
面對神.
面對自己.
不聽聖經.
不再渴慕親近神.
其實就不覺得.
需要改變.
所以這裡告訴我們.
十四節 但你所學習.
所確信 要存在心裡.

$^{1441}$因為你知道.
是跟誰學的.
固然是跟保羅.
其實是跟從小.
跟他敬虔的.
外祖母和母親.
和敬虔的.
那些的聖徒學習的.
所以提摩泰.
現在你可能在一個十字路口.
心靈的十字路口.
你要很小心.
讓那些人走闊路.
走那個大的門.
你應該走窄路.
你應該走那個.
小的路.
窄的門.
並且你知道.
你是從小明白聖經.
這個聖經能夠使你.
因信基督耶穌.
有得救的智慧.
這個聖經是講舊約.
當時沒有新約.
這個聖經是使你.
信基督耶穌.
很有巧救.
你就這樣讀舊約.
如果沒有聖靈.
你只不過是一個守律法的主義者.
但是你讀聖經.
因為你讀到主耶穌.
你就有得救的智慧.
我們外邦人沒有舊約聖經.
小時候在小學.
怎麼會讀舊約聖經呢.
沒有的 猶太人就有了.
但是我們.
進入基督.

$^{1481}$所以我們就沒有了這個包袱.
我們不需要這個困難.
但是猶太人就是.
這聖經使你.
因信基督耶穌.
所以真的有巧救的.
原來不是每個猶太人.
都能讀到基督耶穌.
必須要有那個引導.
他透過講解.
他明白.
神就開到他的心.
他以往一直守的.
律法舊約.
突然之間豁然開朗.
原來舊約所講的.
那個彌賽亞.
就是我們的耶穌.
他打開了心信.
因信基督耶穌.
記住一個猶太人.
因信基督耶穌就有得救的智慧.
我們因為一開始就清楚.
耶穌是人類的救世主.
我們就不需要經過那個包袱.
我們就一信.
我們就立即因信.
就得永生.
不過我們這麼容易因信.
就容易失信.
因為我們沒有經過那些困難.
所以為什麼.
基督教必須是.
一次相信.
要一直相信.
記住這個一直相信.
因為阮妮快要走了.
見不到的.
只在鏡頭.
但是一次相信.

$^{1521}$就不是.
一次得救.
永遠得救.
是一次得救.
相信一直得救.
就可以永遠得救.
所以是強調.
一直的堅持.
所以這段經文是講我們要在.
信仰的路上.
忍耐有苦難.
上文那班人.
就是沒有苦難的人生.
或者說沒有苦難的基督教.
如果基督教沒有了苦難.
沒有了逼迫.
沒有了需要面對自己面對神.
那只是一個俱樂部.
不過名字叫基督教.
但是真正跟隨主的人.
是有那個苦難的.
所以這裡告訴我們.
得救的智慧就是來自聖經.
所以第十六節.
大家會背的經文.
聖經都是神所默示.
於教訓督澤使人歸正.
教導人學而.
都是有益的.
叫屬神的人得以完全.
預備行各樣的善事.
所以相對於上文講的那班人.
就是聖經從來沒有在他們的心靈裡面.
留下過一些感動.
聖經改變不了他們.
那就困難了.
每個人要面對自己.
將來面對神.
所以聖經只是會令到那些人.
愛慕真理的人.

$^{1561}$會有一個真正的改變.
所以教訓督澤使人歸正.
歸向神的正路.
還有教導人.
學什麼是易.
是有益的.
是profitable.
所以我們是屬神的人.
更加成長.
更加成熟.
所以大家的mindset.
來到.
這些階段人生.
沒理由停留在低層次的.
回來的位置被人坐了.
就已經整堂崇拜不開心了.
有沒有搞錯.
不是你的位置.
你就不開心.
你的層次低到這樣.
你應該回來.
看看周圍的人.
你應該為他們感恩.
為他們祈禱.
有什麼事情.
為今天的聚會祈禱.
守望.
爭戰.
為什麼還停留在這些層次.
甚至有些人回來.
我一見到那個人.
我就已經整堂很憤怒.
不解怨.
在面上告訴你.
我是不解怨的人.
我不是借題發揮.
我們人生.
所謂何事.
就是與神和好.
與人.

$^{1601}$我們在基督的恩典裡.
可以和好.
所以很重要.
所以我們看到.
我們在神的裡面.
就是要有這種改變.
現在我們說的是後遷.
我們某程度是.
駁斥後遷.
絕對不是傾向於.
樂觀去思想基督教.
其實是不樂觀的基督教.
真正走擇路.
要忍耐等候的.
是少數的人.
所以提摩泰.
我都要這樣去勉勵你.
我要走了.
我完成了我人生的路.
但是現在那個千斤重擔.
是交付給你.
不是只是教會.
是你自己.
你自己.
教會是你自己.
這個是很重要的.
是你自己.
要這樣來納至.
在基督裡面.
敬田道逸.
所以這裡就是.
保羅對提摩泰的.
勉勵.
所以至少保羅眼中的.
末後是悲觀.
不是樂觀.
基督教不是普遍的.
是人一聽就馬上接受.
馬上生命改變.
一生願意跟隨耶穌.

$^{1641}$不是那麼樂觀.
還有一些經文.
四章的第三第四節.
都是提摩泰後書.
因為時候要到了.
人必定厭煩純正的道理.
耳朵發癢.
聽到不喜歡聽.
真的會喜歡撩耳朵.
好像很全神.
保羅可能也見識到.
這樣的人.
搞錯還未完.
快點說完吧.
他又不敢舉手.
很少這樣的人.
也有修養.
他開始撩耳朵.
隨從自己的情慾.
增添好些師父.
並且掩耳不聽教.
並且掩耳不聽真道.
掩耳.
偏向荒謬的言語.
所以.
他的不從.
人一定會有東西.
想追求的.
所以他會追求那些荒謬.
追求那些.
越是怪誕的他越有興趣.
什麼五夜怪談.
特別多人喜歡聽.
總是嚇嚇你.
蒙古人.
你就會睡得好些.
這就是告訴我們.
厭煩.
他不會厭煩之後.
下星期再來.

$^{1681}$他厭煩之後.
他會隨從自己的情慾.
有些人就慢慢喜歡.
看相片也不是壞事.
星座也有些哲學.
慢慢就越趨向.
那些.
很想去經驗到.
那個靈界.
但是真正的.
那個教訓.
他就不太喜歡.
所以第五節保羅就說.
你卻要凡事謹慎.
忍受苦難.
作全道的功夫.
盡你的積分.
要盡職.
去忍受苦難.
你可以想像到.
不是一個樂觀的傳教事業.
絕對不是夕陽工業.
是一個苦難的工業.
是要孕育一個生命成長.
由信到成熟.
到能夠獨當一面的.
無論什麼環境.
我都要為耶穌作見證的人.
是很難去孕育.
十月懷胎的人.
十月懷胎的孕婦.
所以每一個生命的出生.
洗禮和成長.
成為一個獨當一面的領袖.
是經過很多人的孕育.
所以如果大量量產的製造.
那些出來的未必是真正的.
影響生命.
所以這裡都是告訴我們.
不要這麼快開心.

$^{1721}$很多人相信.
太多人說相信.
但其實那個信不是那麼有質地.
現在說的是新質.
新質的成長.
那些工業,商業都要有質的.
不只是量的.
當然量也是重要.
其實我們剛才讀的這麼多經文.
都是趨向基督教悲觀的成長.
悲觀裡面.
真正有心的人.
仍然在多困難都跟隨著.
你可以說這是悲觀.
但在神眼中.
根本就是他的揀選.
他真的要地上的人有信心.
堅持堅忍到底.
所以你可以說.
我看這些經文.
我絕對不是想.
令到大家受我影響.
而是我真的覺得.
千禧年.
那個經文是多過後遷.
多過無遷.
這些書信都告訴我們.
我相信不只是保羅的時代.
因為有尼祿.
的事業.
一直都在說耶穌回來之前.
傳教本身就是.
你要想清想楚.
是非常困難.
所以為什麼基督教不會悲觀呢.
而要想到這麼樂觀呢.
我就不能夠認同.
後遷的.
只是三段的聖經.
就說那個世界只會越來越大.

$^{1761}$基督的影響力.
和他的管治只會全面.
我就覺得.
這幾段經文.
不是說.
一個這麼實實在在的.
基督教的內涵.
我們再看多一段.
最重要的那段.
馬太福音的24章.
就是主耶穌說的.
這個鑒冷山的寓言.
胡錫進一定是.
從這段經文.
不讀不可能.
第十五節.
24章的十五節.
那個上文就是.
耶穌在鑒冷山.
上面坐著.
這個是.
叫做.
受難州.
就快是.
那個逾越節.
門徒暗暗的來說.
請告訴我們.
什麼時候有這些事情.
因為耶穌寓言.
聖殿會被拆毀.
那就是說.
很不老理.
就快逾越節.
全個猶太人最繁忙.
最燈火鼎盛的.
那個節期.
聖殿是最.
香火鼎盛的.
就是說.
年初三.

$^{1801}$沙田你不要去.
那裡有大災難.
年初三赤口.
年初三是車公最繁忙的日子.
車公廟是最忙的.
你這樣寓言.
等於門徒不問都很難.
你是門徒.
你都會問.
什麼時候有這些事情.
你降臨.
和世界的末了.
有什麼兆頭.
很對號入座.
他不是說下個星期.
逾越節.
不是.
說主耶穌的降臨.
世界的末了.
他問到這麼後的日期.
他其實知不知道.
他們不知道不要緊.
耶穌解釋給他們知道.
第四節.
這個就是背景.
第十五節就是我們要留心.
第四節.
耶穌回答說.
你們要謹慎.
免得有人迷惑你們.
因為將來有好些人冒我的名.
就是迪基督.
假基督.
就說我是基督.
他不是說我是耶穌.
我是拿撒尼的兒子.
那個就沒有power.
救世主.
能夠行神蹟奇事.
所以他們要的是超然的基督.

$^{1841}$不是地上的神人異性的耶穌.
所以耶穌又是人.
又是神.
這是基督教很重要的真理.
一般的異端就不會接受耶穌.
這個人性的部分.
真是一個木匠.
真是要在拿撒尼成長.
他的手錶真是每秒都在走.
他真是要被人白眼.
他是窮苦人家出身.
鬥木.
甚至他的爸爸.
藥室是早死的.
即是說他是大哥.
都要擔起成頭家.
所以你可以想像到.
主耶穌他成長.
都是受著他自己的.
原生家庭的成長的限制.
感懷身世應該會有的.
但他就只會感懷自己是神的兒子.
但就不會影響他對於自己的成長的粉飾疾逐.
因為那個落差越大.
越大對成長的身份的囚禁就很大不滿.
很多男生.
他由十一二歲到長大.
那個所謂暴風少年.
就是看他的家庭環境.
如果那個落差大.
周圍身邊的朋友都很幸運.
而他很不好.
家庭又破碎.
學壞的機會差不多100\%.
因為那個落差太大.
那個時候他在無數自己的成長性格.
沒有人在他身邊成為一個指路的人.
他就大部份機會都學壞的了.
因為他的家庭背景和社會.
和他天天見到的朋友是不相同的.

$^{1881}$所以就會有這種對世界的憤恨.
為什麼我不是李嘉誠的兒子呢?.
這樣的對話就會在他的心裡出現.
所以這裡告訴我們.
主耶穌說我們要小心.
因為迷惑的人會出現.
第六節再告訴我們.
你們也要聽見打仗和打仗的風聲.
總不要驚慌.
因為這些事是必須有的.
只是末期還沒有到.
當時已經有羅馬的軍隊不斷的去打仗.
成為了一個帝國.
很多的外族的俘虜.
都是被帶到意大利羅馬.
就會準備被處死.
因為是俘虜.
為了用他們的強權去鎮壓其他的民族.
所以猶太人見到這些打仗.
根本就是加利利也有些軍人.
納薩勒其實就是一個駐軍的地方.
所以一般的猶太人.
有教養猶太人是不會去納薩勒的.
很合理有羅馬的駐軍出入.
很正常的.
我們見到有解放軍.
都是太平盛世.
當時有打仗.
但是記住.
這些情況似乎不只是限於當時的耶穌的時代.
更加不是限於主後的70年.
因為猶太人背叛.
終於要反抗羅馬.
後來輾轉他們就有革命軍.
最終被提多將軍移為平地.
就有這個馬薩大的地方的殉道.
全部集體自殺.
因為他們不能夠被外邦人殺死.
要自己去死.
所以似乎不是這麼簡單的.

$^{1921}$當時代就應驗了.
一定是在說遙遠的將來.
這裡就繼續告訴我們.
哈馬斯第九節.
這個災難的起頭.
因為打仗自然就會多處有饑荒.
有地震.
現在由哈馬斯的地方再蔓延到黎巴嫩.
一天之內又有幾千人被炸死.
當場死.
你都沒得想的了.
就埋在瓦礫下面.
有男有女.
無差別的了.
可能那裡有一個是真主黨的人.
但大部分是平民百姓.
這樣的無差別.
不過這些是災難的起頭.
那時人要把你們陷在患難裡.
也要殺害你們.
你們又要為我的名被萬民恨惡.
那時必有許多人跌倒.
也要彼此陷害彼此恨惡.
且有好些假仙子起來.
已經有壁壩了.
還有宗教的迷惑.
假仙子的起來.
迷惑多人.
裡應外合.
外面有肉身的迫害.
又有思想神學的迫害.
這樣就令到你的信仰體無完膚.
如果你不堅持到底.
你說又沒有苦難.
信仰又很純正.
每個星期的道都很棒.
每個都受感動.
那信心怎麼會軟弱呢?.
這裡告訴我們.
有很多假仙子起來.

$^{1961}$迷惑多人.
只因不發的時真多.
很肯定.
許多人的愛心就會冷淡.
唯有忍耐到底的必然不夠.
這天國的福音要傳遍天下.
對萬民作見證.
然後末期才來到.
你可以想像到.
又是趨向於傳福音.
悲觀地進行.
極之悲觀地.
很少人可以信到.
接著15節是我們要讀的經文.
所以這張聖經.
是侯謙和毛謙.
是不能解釋的.
因為實實在在後來.
災難是非常之大.
非常之實在.
才是主耶穌的降臨.
這段是戲肉.
差不多可以直接抵消.
侯謙他們樂觀主義的主張.
你們看見先知但以你所說的.
那行毀壞可憎.
但以你所九章.
那個大罪人.
有人說當時那個.
埃及阿古四世.
但未必的.
因為如果將歷史讀進去.
有時會淡化了但以你所說的寓言.
當然近代大部分的學者都會這樣理解.
這裡說站在聖地.
讀這段經文的人要回憶.
沒有括弧可能更容易明白.
有了更難說什麼.
領會什麼呢.
大罪人站在聖地.

$^{2001}$在神的聖所.
那時在猶太的應當途道山上.
在房上的不要下來拿家裡的東西.
在田裡的工作緊.
也不要回去取衣裳.
當那些日子懷孕和奶孩子的有禍了.
大妻子和剛生的要餵奶的有禍了.
走不及了.
有很多限制.
第20節你們應當祈求.
叫你們逃走的時候.
不遇見冬天或安息日.
因為那時必有大災難.
從世界的起頭直到如今.
沒有這樣的災難.
後來都沒有.
若不減少那日子.
凡有血氣的總會有一個得救.
只是為了選民的緣故.
這個日子就減少了.
第23節.
那時若有人對你們說.
基督在這裡或說基督在那裡.
你們不要信.
因為假基督假先知將要起來.
顯大神即大騎士.
倘若能行連選民也就迷惑了.
看啊我預先告訴你們了.
若有人對你們說.
看啊基督在曠野.
你們不要出去.
或說看啊基督在內屋.
你們不要信.
這麼大段經文是說.
那個信仰的挑戰.
要有一個很真確的信心的信仰.
是對基督真基督的那個持守.
不斷說有人告訴你.
基督在哪裡在哪裡.
你不要信.

$^{2041}$接著第27節.
閃電從東邊發出.
直照到西邊.
人子降臨也要這樣.
東方閃電就是根據這一節.
很有趣的.
在東方出現了基督.
整段聖經就叫我們不要隨便信.
誰知道就有人信了閃電教.
就是東方閃電.
這個當然在內地.
據聞是在河南的一個姓陳的女士.
她說是女先知.
但是沒有人知道.
實際是不是這樣都不知道.
第28節.
我看過他們的一些文集.
因為我們教會曾經有一個婦女.
就是跟了東方閃電.
後來都沒有變.
給我看一些東西.
叫我讀.
就是這個女先知.
在香港.
不是在內地.
閃電從東邊發出.
直照到西邊.
人子降臨也要這樣.
屍首在哪裡.
鷹也必在哪裡.
當然那些麻鷹是等待那些屍體.
沒有任何掙扎.
它還沒死.
都不會走過去跟它搏鬥.
鷹是很聰明的.
看著它差不多死.
它就會吃它的腐肉.
它不介意它死了久一點.
因為它根本不需要跟它掙扎.
等它死了.

$^{2081}$就會低飛下來.
來刺它的肉.
這些就是那些麻鷹.
這個意思就是說.
你會感受到基督的降臨.
耶穌這樣說的意思.
所以第29節.
那些日子的災難一過去.
日頭就變黑了.
你可以想像到.
這裡就是整個地球的地動天搖.
月亮都不放光了.
眾星要從天上墜落了.
天勢都要震動.
那時人子的笑頭就要顯在天上.
地上的萬族都要哀哭了.
他們要看見人子.
有能力有大榮耀.
駕駛天上的雲降臨.
他要差遣使者.
用浩同的大聲.
將他的選民.
從四方.
從天這邊到天那邊.
都焦聚了來.
這裡會不會就是.
帖撒羅里加前書說的.
主在來.
信徒的被提呢.
就是災難過了之後.
那怎麼會是樂觀的基督教.
普遍的傳遞呢.
為什麼不是假基督.
繼續這樣蔓延.
直到真的基督降臨呢.
真的笑頭顯現呢.
為什麼不是有千禧年的悲觀的傳教呢.
然後才真正開始地球上的千禧年呢.
真正沒有罪惡的更新了的地球.
撒旦是在這個無底坑呢.

$^{2121}$所以我是這樣按字面去理解千禧年.
已經有五段經文去駁斥這些後遷的太樂觀.
由基督升天就是千禧年了.
就是世界就趨向美好了.
這個就是解釋不到的.
就是一個大災難之後的基督降臨.
明顯是經文來的.
是耶穌在橄欖山預言給他的門徒知道.
他們的問題就是.
你降臨世界的末了是什麼時候呢.
有什麼預兆呢.
這個就是人子的笑頭就要顯在天上了.
那時的意思就是上文的地動天搖.
假基督的傳播.
大災難和那些戰爭,瘟疫,饑荒地震都過去了.
那時第三十節.
就是由那些完了之後才是那時人子的笑頭.
所以很難不是按字面去理解這些經文.
如果真的按字面去理解.
只能夠說我們比較相信傾向於有千禧年.
是比較合理去看到這個世界的演進.
這個演進似乎是想拒絕基督.
多過是全世界都基督教化.
如果我的結論就是這樣.
9點01分我們要完了.
感恩給我們能夠在這一個半小時.
看了後千他們的理論.
一方面我們尊重每個神學家.
不同對聖經的理解.
至少我們知道主耶穌是一定回來.
三個的理論 無千,後千和千禧年.
都是相信主耶穌會再來.
不過是時間軸的問題.
但我們懇求主讓我們看到.
我們現在活在地上.
其實真是面對敵黨和福音的難阻.
幫助我們好像十八章的盧家福音.
要常常禱告不可灰心.
在地上是一班等待基督降臨的人.
是一班熱切去傳福音的基督徒.

$^{2161}$我們這樣感恩禱告奉耶穌基督的名.
阿們.
\newpage



\section{}
\label{sec:H4W4IN0LZRY}
\textbf{2024-10-03 周四晚查經班《淺談末世論(4)》 講員:曾裔貴牧師}
\newline
\newline
連結: \href{https://youtube.com/watch?v=H4W4IN0LZRY}{\texttt{https://youtube.com/watch?v=H4W4IN0LZRY}} ~~~~ 語音日期: 2024-10-04
\newline
\newline
\hyperref[sec:tl2Gc71LinQ]{\small{< < < PREV SERMON < < <}}
~
\hyperref[sec:index_chronic]{\small{[返順時目]}}
~
\hyperref[sec:index_scriptual]{\small{[返順卷目]}}
~
\hyperref[sec:aTaCCC_F45U]{\small{> > > NEXT SERMON > > >}}
\newline
\newline
$^{1}$感恩我們又可以再一次透過聖經的真理.
去思想一些神學的題目.
我們今天晚上會探討關於前千禧年那個證據.
他們所持的理由.
都盼望神給我們能夠明白.
特別我們想到我們眼見到的地球.
已經有兩場差不多是跨越了地域的戰爭.
其實都沒有停止過.
在俄烏或者在以巴.
現在去到黎巴嫩.
甚至乎去到伊朗.
我們看到這些情況.
基本上沒有辦法有人可以駕馭.
我們很盼望天父給我們見到.
這個世界是需要有主的救恩.
和主親自來管治.
才可以真正解決那些問題.
我們都盼望基督徒.
我們不要失去了對永恆.
對主的再來的盼望.
這個熱切的渴望.
給我們在很有限的人生.
為永恆去打算.
為神的國度去傳福音.
真是將人帶進神的愛和救恩裡面.
我們今天晚上會看前線的一些經文.
願神都帶領幫助我們.
亦都給我們往後.
我們要思想一個關乎到宗教改革.
其實一個都幾大幾重要的時代.
都真是扭轉了基督教和天主教的關係.
和對真理的重新的堅持等等.
這些都是很影響後世.
我們過去那500年基督教的發展.
願神都給我們大家透過這些簡單的講解.
絕對不可能是神學院的級數.
詳細的30個小時的課程.
但是給我們深入淺出.
大概知道什麼叫做宗教改革.
給我們又有了一個印象.

$^{41}$或者宗教改革帶來的那些成果.
給我們真是有一個持定的心靈.
今天晚上我們等候在神的面前.
這樣感恩禱告奉耶穌基督的名.
阿們.
是的 正如我想講的就是.
毛遷希年很簡單的理論就是.
否定《啟示錄》的20章.
因為他們認為全本聖經只有那裡.
講關於有千希年.
而《啟示錄》是充滿了很多的異象.
很多的一些象徵.
就不應該20章按字面去理解.
就去否定了有實質的地球上的一千年.
是用calendar這樣去數去計算.
而且那個希年的內涵.
就和後千是不同的.
後千就是說我們一直計算.
由主耶穌升天.
大使命神的國度在地上管治.
只會越來越好.
不會越來越差.
這樣就叫做希年了.
但是《啟示錄》的20章.
和前千希年的理論的主張.
其實某程度上真的很難達到.
按我們人的理解.
所以你說可以很高檔次.
因為真的要講一千年的那個.
你真的要揭月曆揭一千年.
那一千年是耶穌在地球.
和那些的聖徒在地上.
所以這個是按我們今天很難去理解.
但是我們前千的主張就是.
那段經文是你不應該不按字面解.
這個是其中一個很重要的argument.
另外因為無論後千和無千希年的理論.
他們全部都沒有處理到舊約很多的預言.
這個我已經找到了.
其實當然不是我找到.

$^{81}$是聖經一直都浮現給我們.
我所有的手指1 2 3 4 5 6 7 8 9 10 11.
至少有11段經文.
我們要處理.
講到又不是天堂.
又不是現在的地上.
但是講的是明顯將來的一個世界.
這些經文要處理了.
如果不處理的話.
你說已經過去了.
就不會發生了.
那就變得太容易否定了先知他們的預言.
但是不約而同地.
講到有一個很特別的日子.
在遙遠的將來.
不是現在的世界充滿了戰亂.
更加不是將來永恆的天堂.
而是好像有一個地方將來會有太平盛世.
所以我們就要進入這些經文裡面.
某程度上前千不是單單透過啟示了的20章.
所以其實大家聽了這4堂.
你已經比其他的弟兄姊妹多了一些知識.
就是原來這個世界有神學家主張.
有千禧年或者沒有的千禧年.
或者是後千禧年.
但是無論如何都是主的再來的一個時間.
我們要揭幾段經文.
跟大家去看看.
有些是比較生.
大家很少去揭的.
西番雅書.
全部舊約的先知書.
大家可能連聽都沒聽過.
又是要背一下聖經的目錄.
荷爾摩羅那尼紅哈番該亞瑪.
哈該書之前就是西番雅書.
西番雅書的第三章.
可能我們讀多一點給大家概念.
往往去到先知書的後半段.
一般先知書就明顯神給這些先知.

$^{121}$是有當時代要來預言當時的國家.
或者周邊的國家.
他們的光景.
靈性的光景.
所謂靈性光景.
就是和誰做一個標準去做比較.
一定是和摩西的律法書.
神立約的子民.
那個屬靈的光景做一個比較.
或者一些比較好的時代.
猶太人想到的一定是大衛.
所羅門的王朝.
這些時代他們遵行神的旨意.
國家就會興旺.
所以所羅門皇后期.
和他之後的皇帝.
大部分都不是政治局勢不好.
而是和神的靈性關係不好.
先知就會出來勉勵,指責,控訴等等.
所以先知書大概是這樣的層次.
不會是無的放矢.
走出來自述自話.
一定是有那個差遣.
或者有那個呼召.
去到百姓當中.
或者本身就是生活在百姓當中.
這樣來講解.
所以每個先知都有不同的時代的特性.
所以猶太人的聖經.
他們不是分大小先知.
沒有大小之分.
先知書只有前先知書和後先知書.
這個大家要記得.
所以十二小先知是我們自己叫的.
是因為七十事易本的翻譯.
就變成這樣.
這些小小的書卷其實不算小的.
其實他們的書卷.
比較內容的寓言是小.
絕對不是小先知.

$^{161}$因為以賽亞是大先知.
就沒有這樣的分法.
是前先知和後先知.
一定是按皇國的歷史做一個界分.
所以後先知書大部分是差不多亡國.
耶利米等等.
所以耶利米哀歌就不是先知書.
耶利米哀歌是詩篇裡面的經文.
對不起.
耶利米哀歌是被歸納去普通的信件.
所以大家要記得.
猶太人他們分聖經不是這樣分的.
我們就分了大先知和小先知.
大先知是賽耶哀結旦.
但易理書原來不是先知書.
但易理書是放在普通的書卷裡面.
所以你會留意到這樣理解就很不同了.
他們就真的會按經文.
按寓言去作解釋.
西班牙書第三章.
先講了一些當時的光景.
然後就會突然加插了一些好的希臘的世界.
第一節開始.
這輩亦污穢欺壓的城有禍了.
當然是講耶路撒冷.
他不聽從命令.
不領受訓誨.
不依靠耶和華.
不親近他的神.
他中間的首領是咆哮的獅子.
他的審判官是晚上的豺狼.
一點食物也不留到早晨.
晚上的豺狼.
即是聯群結隊.
你遇到這些狼群.
基本上你都沒得靜.
晚上出現.
一點食物也不留到早晨.
似乎審判官是講那些狼吞虎嚥的酒席.
飯局.

$^{201}$喝的可能是貴州茅台.
你可以想像到那些的光景.
他的先知是虛浮詭詐.
虛浮詭詐即是說.
講的東西可能是平安.
不會去責備.
不會講事實.
不會講真話.
講一些大家喜歡聽的虛浮.
講一些大家都覺得很舒服的.
回來都說阮明你是最好的.
不斷這樣講的時候.
不需要去責備.
先知就容易講.
他的祭司攝讀聖所.
強解律法.
你想想.
首領審判官.
即是有實權的.
先知本來就是要有神的啟示.
那個寓言.
去對當時講出.
此時此道的.
那個實際的光景.
祭司就應該辦理那個憲制.
要根據律法去憲制.
什麼時候不根據律法呢.
就是百姓都求求其其.
你都照收那些祭物的.
那就叫強解律法.
耶和華在他中間.
是公義的.
斷不作非義的事.
每早晨顯明他的公義.
無日不言.
此時不義的人不知羞恥.
所以三章五節.
這些是靈修很重要的經文.
因為原來神一直在他的百姓當中.
是他自己秉公行義.

$^{241}$他一定不會作一些非義的事.
所以我們能夠有神的.
每早晨的訓誨.
第六節.
我耶和華已經除滅列國的民.
他們的城樓毀壞.
我是他們的街道荒涼.
以至無人經過.
他們的城邑毀滅.
以至無人也無居民.
我說:你只要敬畏我.
領受訓誨.
如此,你的住處不至照我所指定的除滅.
只是你們從早起來.
就在一切事上敗壞自己.
這些就是後先知書的主題.
就是告訴他們.
如果你們肯悔改.
你們是可以改變的.
這些是前先知書的主流的信息.
但後先知書就說.
如果你們改變.
也無法取消.
收回神的審判.
這是時代的問題.
所以你讀耶利米書.
你會看到一個好像完全沒有希望的國家.
透過耶利米不斷的哭喊.
說明這個國家將會被擄到巴比倫的光景.
即使他們有一些憤青會.
也改變不了他們的結局.
這就告訴我們.
現在講出他們內心國家的問題.
第八節.
耶和華說:你們要等候我.
直到我興起努力的日子.
因為我已定義朝聚列國.
聚集列邦.
將我的老奴.
就是我的列奴都傾在他們身上.

$^{281}$我的憤怒如火.
必消滅全地.
留意第九第十節.
突然間又有一個不同的光景.
那時我必使到萬民用清潔的言語.
好求告我耶和華的名.
同心合意的事奉我.
祈禱我的.
就是我所分散的民.
必從古實河外.
給我獻供物.
你就可以想像到.
這是一個美好的光景.
到底那個日是哪一日呢?.
第十一節.
當那日.
你必不因你一切得罪我的事.
自覺羞愧.
因為那時.
我必從你中間除掉.
驚誇高傲之輩.
你也不在於我的聖山狂傲.
我卻要在你中間.
留下困苦貧寒的民.
他們必投靠我耶和華的名.
以色列所盛夏的人.
必不作罪孽.
不說謊言.
口中也沒有詭詐的舌頭.
而且吃喝躺臥.
無人驚嚇.
錫安的民應當歌唱.
以色列應當歡呼.
耶路撒冷的民應當滿心.
歡喜快樂.
耶和華已經除去你的刑罰.
輾轉等等.
一路說著美好的光景.
你說這裡就比較沒有那麼明顯.
說到一個很特別的好日子.

$^{321}$但是已經有一個好像悔改之後的復興.
而且留意第九節就是.
萬民用清潔的言語來求告耶和華的名.
而且萬民同心的合意去事奉神.
所以你可以想像到這一段聖經.
如果你放在千禧年的角度去理解.
其實就很合理.
但是我們從來都還沒見到.
以色列到今天.
都有這樣的光景在地上.
來有這樣的歡樂平安和太平的日子.
所以你可以說一定是遙遠的將來.
但是這個遙遠的將來.
是怎樣來發生和應驗呢?.
所以很多千禧年的主張的人.
就會將這些經文在千禧年裡面發生.
就可以這樣理解.
如果這些都不行.
我們又再多揭幾段.
我們看看尼加書.
尼加四章.
以摩訶拿尼.
和以摩訶拿尼洪哈番該亞瑪.
即是俄巴底亞書約拿書之後.
尼加書大家讀過第五章.
就是彌賽亞降臨.
在伯利恆第五章的這一段.
就是主耶穌的降生.
這裡將來有一位從里那裡出來.
在以色列當中作掌權.
他的根源是從耿古.
從太初就由伯利恆而發他這個小城.
在猶大的諸城當中為小.
那些文士.
希律王問他基督生在哪裡.
他們就立即引述到尼加書第五章第二節.
但是我們今天看的不是這一段.
我們看的是四章一到第四節這一段.
第三章是說一個人血饅頭的耶路撒冷.
第四章是突然說一個很美麗的耶路撒冷.

$^{361}$所以就有一個很強烈的對比.
這個對比我稍微讀一讀給大家聽.
三章尼加書.
我說亞國的首領以色列的官長.
你們要聽你們不當知道公平嗎.
當然會知道不過不肯.
你們互善好惡.
互就是那個動詞.
從人身上剝皮.
從人骨頭上剔肉.
吃我們的肉.
剝他們的皮.
打折他們的骨頭.
分成筷子將要下鍋.
就是打火鍋.
由丈夫中的肉.
到了遭災的時候.
這些人必哀求耶和華.
他卻不應允他們.
那時他必照他們所行的惡事.
向他們掩面.
就是說到處都是貪官污吏.
耶路撒冷是這樣的光景.
輪到使我們走差路的先知.
先知本來要說真話.
先知有一個標準.
祭司嘴裡當傳知識.
人也應該從他口中知道律法.
這是馬拉基書二章的經文.
先知應該讓百姓走神的路.
輪到使我的民走差路的先知.
我們都要很小心.
所以港台不可以隨便讓一些不太信任的牧者去說.
你不知道他在說什麼.
當值祝福的童工有責任把關.
如果他說錯就要問他.
你剛才這樣說是不是真的有這個意思.
其實剛才很累.
說錯了耶穌不是神.
你要很禮貌地問他.

$^{401}$他說是耶穌不是神 是真的.
你要當眾譴責.
因為聚會還沒完.
祝福的牧者就要說.
這位牧師說的是錯的.
不正確的.
聖經不是這樣說的.
他也繼續說耶穌不是神.
要當眾譴責.
因為他在說異端.
當然很少這麼戲劇性的.
怎麼會這樣.
當然不會.
但是我們把關的責任是真的.
真的要把關.
因為論道使我民走差路的先知.
他們牙齒有所著.
他們就呼喊.
果然我們時時說的道都是平安了.
從來都不痛不癢.
其實我們都要問.
我們是不是使我們民走差路的先知.
因為那個訊息永遠都是這個主題.
真是好了.
這樣才是真是好了.
不用悔改.
大家不回祈禱會沒問題.
大家不回團體也沒問題.
不奉獻也沒問題.
這樣就真是使我們民走差路.
聖經的要求我們盡量減低它.
某程度上也不是全道.
不過我們在這裡.
大家開心地唱聖詩就好了.
不是真正神的國度的指紋.
我完全沒有意思指桑罵槐.
完全沒有這些意思.
而是聖經真的應該要教我們.
什麼時候神要我們講一些安慰的話.
我們就要講安慰的話.

$^{441}$什麼時候神要我們講一些免禮和指正的話.
我們就不可以避免不講.
怕得失人就不如不要做全道先知.
當然我們不是為得失人而得失.
而是為了神的道.
而我們講一些說話.
有些人不喜歡聽我們都要講.
這就是那個意思.
所以永遠他們的主題.
某程度上就是.
如果宮廷那些.
在耶利米年代.
那些叫擦鞋仔先知.
就會不斷講平安鳥平安鳥.
這裡果然尼加斯的時候也是這樣.
不過凡是不供給他們吃的.
他們就預備攻擊.
耶和華如此說.
你們必遭遇黑夜.
以致不見異象.
又必遭遇昏暗幽暗.
以致不能夠占卜.
日頭必向你們沉落.
白晝變為黑暗.
先見必抱攜.
占卜的必蒙羞.
都必烏著嘴唇.
就是沒有話可以講.
因為神不再應運他們.
其實到了一個地步.
他們已經沒有神的道去講.
已經沒有感動去講.
已經不想講.
最好不要講.
不是大家爭先恐後.
表達謙卑.
而是沒有道去講.
這是一個牧者最大的可悲.
是沒有道可講.
抄回十年前.

$^{481}$二十年前的講章.
去講.
反正現在電腦隨便搜尋.
就馬上找到.
找不到就叫保護找程序表.
曾經講過哪個題目.
講過哪段經文.
拿出來再講.
古代就是這樣.
沒有話可以講.
神要掩著他們的嘴.
接著.
尼迦就很不同.
你看尼迦的內心.
至於我.
我夕耶和華的靈.
滿有力量公平才能.
可以向雅各說明他的過犯.
很勇敢.
每個人都不敢講.
甚至乎也不肯講.
因為和利益掛帥.
和自己的面子掛帥.
怕得失人.
但尼迦就說.
不是啊.
我夕耶和華的靈.
我就滿有力量公平才能.
這就是剛剛的對比.
假和真的才知道.
向以色列指出他的罪惡.
所以神的靈原來是感動我們.
所以神的靈原來是感動尼迦.
去說什麼呢.
滿有力量公平才能的道.
就是講出.
雅各以色列的實際光景.
你做醫生.
你是肺科的.
拿著一張肺片.

$^{521}$沒事沒事.
那個陰影沒事的.
你回去就行了.
你還在吸煙嗎.
明明那個肺都花了.
你說沒事.
你可以想像到.
你去看這樣的醫生.
你是否期望那個醫生這樣跟你講.
明明你的肺都花了.
他就告訴你沒事.
所以這裡就是.
我們往往期望找一些黃綠醫生.
來判斷.
我們大家的.
那個屬靈的光景.
這個其實是很可憐的.
所以有時物極必反.
傳道人很有神的恩典.
反而有時他罵人.
人們反而會去找這些道聽.
原因就是他肯講真話.
這個很重要.
第九節.
我們還沒講第四章的千禧年.
是在講強烈的對比.
第九節.
雅各加的首領以色列加的官長.
當聽我的話.
你們厭惡公平.
正常的首領.
就是要秉公行事.
在一切的事上.
他們就屈枉正直.
以人血.
我就是講這個人血饅頭.
以人血建立錫安.
以罪孽建造耶路撒冷.
也就是說有很多的不公不義.
很多的冤案.

$^{561}$很多的貪贓枉法.
首領為了賄賂就行審判.
祭司為了顧價.
這個顧字真的非常好.
神的僕人成為受顧的.
米飯班主的一個手下.
講的是米飯班主想聽.
喜歡聽的.
祭司為顧價施憤悔.
先知為銀錢行占卜.
每句都很押韻.
先知,祭司,首領.
包括軍王.
他們卻依賴耶和華說.
耶和華不是在我們中間.
災禍必不臨到我們.
所以因你們的緣故.
錫安必被耕種成一塊田.
所以還是要依賴耶和華講.
耶和華在我們中間.
大家不用擔心.
必定不會有災禍.
先知這樣的預言.
這樣的安慰.
所以這只是一個假象.
不斷給假象的百姓.
但實際上.
神說錫安必被耕種成一塊田.
耶路撒冷必變為亂堆.
這殿的山必將叢林的高處.
如果你去過聖殿山.
你就知道那個宏偉.
那個輝煌.
但變成了一個叢林.
這就是尼加先知.
他講真話描述真正的情況.
不過突然間四章.
有一個很強烈的對比.
講到一個很美的耶路撒冷.
只能夠說.

$^{601}$如果不是天希年.
都很難發生.
這四章.
這段是很重要的天希年的經文.
聯合國都拿了來做他們的座右銘.
四章的第三節.
所以你去聯合國紐約.
你會看到這節金句.
末後的日子.
耶和華殿的山必建立.
超乎珠山.
高舉過於萬靈.
萬民都要留歸這山.
當然不是說.
突然地殼變動再升高一點.
好像現在大家知道.
珠穆朗瑪峰升高了.
當然他的溫室效應.
雪地越來越融.
但是因為真的有地殼變動.
當然是很精準的量度.
但是耶路撒冷就不是突然升高了.
而是在所有的萬民當中.
他成為一個最重要的宗教中心點.
萬民都要留歸這山.
必有許多國的民前往說.
來吧.
我們登耶和華的山.
奔雅各臣的殿.
主必將他的道教訓我們.
我們也要行他的路.
所以這些經文其實放在新約.
就是我們敬拜神.
要有這樣的心情.
我們要奔耶和華的山.
我們要得到神的訓誨.
他要指引我們的道路.
因為訓誨必定出於石安.
耶和華的言語必出於耶路撒冷.
所以每一個禮拜帶著一個期待.

$^{641}$去聽港台的訊息.
是重要的.
有這一份等待.
神有說話要跟我說.
接著第三節.
耶路撒冷.
他必在多國的民中心.
施行審判.
為遠方強盛的國.
斷定是非.
他們要將刀打成泥頭.
把槍打成鐮刀.
這國不舉刀攻擊那國.
他們也不再學習戰事.
不需要再打仗了.
不需要再空寒了.
想想要用什麼戰術.
製造什麼戰爭的工具.
彈藥.
人人都要坐在自己的葡萄樹下.
無花果樹下.
無人驚嚇.
這是最重要的.
描述那個光景.
大家安居樂業不止.
還要毫無驚嚇.
大家是盛評的時代.
這是萬軍之耶和華親口說的.
萬民各自為之盛評.
這些可能大家比較少有機會去講和去讀.
第十七節.
看啊.
我做新天新地.
從前的事不再被紀念.
也不再追想了.
你們當因我所做的.
永遠歡喜快樂.
因我做耶路撒冷為人所喜.
做其中的居民.
為人所樂.

$^{681}$我必因耶路撒冷歡喜.
因我的百姓快樂.
其中必不再聽見哭泣的聲音.
以及哀號的聲音.
第二十節.
其中必沒有數日.
夭忙的嬰孩.
也沒有受數不滿的老者.
因為百歲死.
仍算孩童.
有百歲死的罪人.
算被咒詛.
他們要建造房屋.
自己居住.
栽種葡萄園是其中的果子.
他們建造的.
人不得住諸如此類.
一直這樣說.
然後到第二十四節.
二十五節.
又再一次這樣說.
這些曾經在十一章都出現過.
第二十四節.
他們尚未求告.
我就英文了.
正說話的時候我就垂聽.
接著二十五節.
豺狼必與羊羔同食.
獅子必吃草.
塵土必作蛇的食物.
在我聖山的遍處.
這一切的野獸.
有殺傷力的.
兇猛的.
都不會再傷人.
不會再害物.
這是耶和華說的.
如果你真的要按字面去理解.
或者為什麼你不按字面去理解呢.
如果按字面去理解.

$^{721}$就解不通了.
如果不是擺一個.
千禧年的世界.
這些預言.
就變成一個笑話.
有那麼遠就那麼遠.
豺狼見到羊.
怎麼會一起在這裡找東西吃呢.
羊本身就是豺狼的食物.
怎麼會一起同食呢.
獅子還吃草.
跟牛一樣.
牛我們知道.
這是神造的特性.
大自然界也有根本性的改變.
變成吃草.
將來我們不用吃東西了.
喝水就飽了.
當然神可以做.
通常我們沒有水就不能生存.
但喝水飽.
這是開玩笑的.
有情可以喝水飽.
這是預場片.
今天.
我忘記了是星島日報報導.
原來錢.
是男女分手.
最大的原因.
金錢.
所以有情喝水飽.
原來是假的騙人.
所以這裡告訴我們.
如果有這樣的自然界改變.
如果不是千禧年.
我就想不到要怎樣.
天堂怎麼會有野獸.
也不會.
在目前的地球.
我們任何地方.

$^{761}$無論是非洲也好.
在香港.
好像還有一隻老虎.
在冰頭花園.
我忘記了.
都是已經被困到純了.
但純了還純了.
你走進去.
牠會咬你.
我們知道大自然神造的.
賦予的.
自然界的.
生態系統.
還沒有改變.
如果改變.
如果按字面應驗的經文改變.
就應該放在千禧年裡.
這就是.
前千的理論.
我覺得是合理的.
因為如果不是一種.
憧憬.
不是一種烏托邦.
以賽亞.
不是說一些誇張的東西.
其實我誇張的.
其實不可能的.
就不應該說得這麼實在.
但一直提到.
在神聖山的遍處.
這一切.
本來有兇猛的.
獸性.
是不會再傷害.
這是耶和華說.
所以.
這些經文是等待.
遙遠的應驗.
沒有人知道何時是千禧年的開始.
因為千禧年的開始.

$^{801}$如果前千的理論.
是首先結集了.
七年大災難.
有大災難的降臨.
就是馬太福音24章.
我們上星期看的.
主耶穌的預言.
何時降臨的.
前奏.
是會這樣的.
所以如果過了那些災難之後.
主耶穌第二次回來.
開展千禧年.
魔鬼撒旦又被拘禁.
在無底坑.
這些事情改變.
我覺得是很合理的.
就是說.
我們有機會見識到.
將來有這樣的世界.
就不是一個冒險樂園.
是真真正正的千禧年的國度.
是的.
這裡就是.
我們看了.
我們不如看以賽亞書還有幾段.
我們看第二章.
第二章的第一節.
阿摩斯的兒子.
以賽亞得到默示.
論到猶大和耶路撒冷.
末後的日子.
耶和華殿的山必建立.
超乎珠山.
高舉.
臥於萬陵.
萬民都要留歸這山.
不斷在說這個訊息.
不過絕對不是在.
他那個時代.

$^{841}$其實是令到那些百姓.
流離失所.
絕對不是萬民.
都要留向.
耶路撒冷.
但是將來.
末後的日子就會這樣.
必有許多國的.
文前妄說.
來吧 我們登耶和華的山.
賓雅各臣的殿.
好似曾相識.
剛才我們看的.
離家書就已經說了.
所以離家是有一點.
和以賽亞.
重疊的時間.
主必將他的道教.
訓我們.
我們也要行他的路.
因為訓誨必出於石安.
耶和華的言語.
必出於耶路撒冷.
宗教改革.
最重要的改革.
就是回到聖經.
你可以說有1000年.
大概1000年.
天主教已經將聖經.
封閉了.
只能夠神職人員.
才可以講解.
和去中讀.
是用拉丁文.
當時如果馬丁路德.
就是德國人.
他是一個本身.
受過足夠的訓練的律師.
所以他的思辨能力.
他的文學.

$^{881}$他的語言能力.
都是很強.
至少他懂得希臘文.
和拉丁文.
因為他自己也翻譯了.
很多卷的新約聖經.
然後他還有很多著識書.
而且他用德文.
去翻譯.
將拉丁文翻譯成德文.
就是說.
一般的德國人.
都看得明白那些文章.
才會帶來復興.
才會帶來宗教的改革.
如果他還是用拉丁文.
那就是說.
一小撮精英分子才會.
就不是一般人會.
所以宗教改革.
其中一個最重要的.
當然首先是針對贖罪卷.
這個大家要知道.
但其實不是的.
是回到人人可讀聖經.
就好像周星馳的搞笑片.
人人都可以有功練.
那些我不記得是那套戲了.
是不是破壞王還是什麼.
這就告訴我們.
將來千禧年.
當然.
我對號入座也好.
千禧年的光景.
就不是說.
每天吃自助餐.
遊山玩水.
不會有惡魔不會有豺狼.
相信不是這樣.
相信是有傳遞的神的.

$^{921}$同在和訓誨.
因為留歸耶和的山.
都是很想去尋求神的道.
所以其實我們每個禮拜.
你自己有空.
你想一想.
你走回去教會.
你去敬拜神.
其實這件事本身是很偉大的.
我們在地球上的人.
其實很多事好做.
很多事都重要.
但我們用一個多小時.
去朝聖.
敬拜神.
其實這件事是很偉大的.
我們要讓這件事真的偉大.
你真的整裝待發.
準備好你自己的新心靈.
穿好一點.
不是為了炫耀.
是為了禮貌.
我真的不是很接受.
穿涼鞋.
穿短褲去敬拜神.
方便一會去旅行.
一會去郊遊.
不想走路.
又要穿恤衫又要穿皮鞋.
所以明仔有時很激氣.
上午要靈祥.
下午又要做團契.
變成他穿西裝打領帶.
做主席.
他很煩.
家裡又遠又要出去元朗.
穿著整套等下午的團契.
有時反而覺得有困難.
但他又要穿.
因為他知道敬拜是對神.

$^{961}$他要做好自己的本份.
那個榜樣.
作為一個無論是靈祥.
或者是一個主席.
這些麻煩.
有時要看我們內心的態度.
當然我們不是說要很吹毛求疵.
但其實有志於內的內心.
相信我們這一代的人.
明白不同的場合.
穿不同的衣服.
是一個尊重.
其實你從這個尊重去看.
不要覺得是挑戰.
你從尊重.
你就會感覺到這件事.
其實是沒有抵觸的.
尊重的場合.
譬如有些高級餐廳.
你穿短褲泳衣.
根本是不准你進去.
這是基本上的禮貌.
所以.
如果我們見到千禧年.
或者末後的日子.
我不想引導大家.
一定要接受這是千禧年.
但那個日子.
就是以賽亞同境的.
將來的美好世界.
百萬民.
去留歸這個山.
尋求耶和華的訓誨.
這是可以肯定的.
所以.
很明顯他們有了神的道之後.
自然性情也會改變.
不再私算打仗.
因為神的道已經是.
私行審判.

$^{1001}$為許多國民.
斷定了是非.
就是大家都有公平.
就不需要再爭.
現在以巴的世仇.
巴勒斯坦佔據了.
以色列的故土.
因為他曾經被趕逐.
去到天下.
當他復國.
1948年.
回到以色列.
整個阿拉伯世界就佔據了.
他就要不斷爭.
自己神聖的.
神選民的土地.
這就是麻煩所在.
就是磨心.
就是這個問題.
大家都覺得不公平.
問題就是.
人們已經在.
差不多第七世紀.
已經在這裡住.
但問題就是.
第一次世界大戰之後.
英國和美國.
准許和鼓勵.
支持.
以色列回到復國.
就帶來整個阿拉伯世界.
對以色列的.
針對和攻擊.
所以這個問題.
將來就不會有了.
因為神的判斷.
審判.
是公平的.
為了許多國家已經斷定了.
是非了.

$^{1041}$所以他們能夠做的就是將刀打成泥頭.
把槍打成鐮刀.
這國.
不舉刀攻擊那國.
他們也不再學習戰事.
那國家.
內霸我們在耶和華的.
光明當中行走.
將來.
那個幕後的光景.
這個就是.
以塞翁的書和尼加書.
所描繪的.
太平的.
希臘.
可以再看多一兩段.
第九章也有.
第九章就是.
以馬來尼.
第九章.
我們可以稍為讀前面一點.
譬如我們很熟悉的就是第六節這一段.
描述.
當將這個嬰孩出現的時候.
它就會帶來一個很戲劇性的.
國家的扭轉.
似乎這個嬰孩.
在這個時期.
馬上應驗.
以塞瓦的第二個兒子.
這個.
我忘記了是不是斯瓦瓦術.
因有一嬰孩.
為我們而生.
有一子賜給我們.
當然記住.
新約的應驗.
肯定是講到主耶穌.
政權必擔在他的肩頭上.
他明知.

$^{1081}$他必擔在他的肩頭上.
他明知.
他必擔在他的肩頭上.
他明稱為奇妙策士.
全能的神.
永在的父和平的君.
這樣一個整全的.
主耶穌的稱呼.
他的政權與平安.
必加增無窮.
他必在大位的寶座上.
治理他的國.
以公平公義.
使國堅定穩固.
從今直到永遠.
萬君之耶和的熱心.
必成就者事.
就相當肯定.
這個將來的掌權者.
就會帶來一個這樣的統治.
所以這段經文.
可以說是真正主耶穌在地上的管治.
九章的第六第七節.
所以不只是.
彌賽亞在伯利恆降生的耶穌.
嬰孩耶穌.
是講到這一段的第七節.
是講到耶穌在地上的管治.
第一次降臨.
根本就.
不可能管治.
是被人殘殺.
但是第二次降臨.
就會有這個.
權力的施政.
他在大位的寶座上.
來治理他自己的國家.
用公平公義.
使耶路撒冷.
這個人.

$^{1121}$在地上的管治.
是完全不可能的.
使耶路撒冷.
猶太將來的國度.
堅定穩固.
所以這些經文.
放在任何一個現在的時代.
的以色列都還未能夠.
發生的.
主耶穌更加未回來做.
統治者.
所以他的第二次降臨.
就是猶太人等待的彌賽亞.
的降臨.
其實就是前千的.
追求者.
他的理論將來就是.
主在地上的執政.
這個就是.
千禧年的主張.
接著我們可以.
看第十一章.
都是異創的書.
十一章.
稍為後一點.
我們可以看多一點前面.
因為又是講到.
主耶穌的統治.
從耶西的本.
必發一條.
這個本.
原文是那個.
等.
那個木字沒有了上面的.
一棟.
只剩下下面那個不是筆字.
那個是等.
樹等.
就是巨大的樹.
斬了之後.

$^{1161}$就剩下那個等連著地上.
那個就是樹等.
就是打等.
從他更新的.
之子必結果實.
所以耶西的本.
在耶西的家族.
出現的.
這個.
似乎是.
表面看.
好像很脆弱的.
一個支條.
但是因為有神的靈.
第二節.
耶和華的靈必住在他身上.
就是使他有智慧.
和聰明的靈.
所以一個牧者.
如果和神近.
是有屬靈的智慧.
和悟性.
牧者.
未必是.
博覽群書.
但是牧者應該.
和神的關係.
是會帶來一種.
屬靈的敏銳.
智慧 聰明.
所以大家要近一些傳道人.
會得到一些亮光.
和恩典.
謀略和能力的靈.
知識和敬畏耶和的靈.
所以.
任何一位主學老師.
領袖必須要.
追求具備.
其實簡單就是與神相近.

$^{1201}$就會有神的靈.
同在帶來的.
原來可以這樣分析.
有一種智慧.
有一種聰明.
有一種謀略.
即是他的判斷.
一些銳志.
未必是經驗.
當然不限於經驗.
包含了經驗.
有能力.
能力就是.
事情的果斷.
等等.
勇敢.
當然這一切的背後.
整體的內心.
是一個敬畏神的心.
第三節.
他必以敬畏耶和華為樂.
行審判.
就不憑眼見.
斷是非.
也不憑耳聞.
所以.
眼見也不是實在.
聽也不要偏聽.
要公平地聽.
特別我們要做一個決定的時候.
卻要以公義審判貧窮人.
貧窮人就不是完全對.
舊約往往告訴我們.
一般來說.
貧窮人受欺壓.
但是原來公義的審判.
一樣不要只是向有錢人.
窮人也要審.
因為窮人也可以有鬼擇.
所以聖經說.

$^{1241}$貧窮人.
是要公義審判.
貧窮人.
有鬼擇.
所以聖經說的公義審判.
往往不是只是為審判有錢人.
反而變成不公義.
不要屈望正直.
以正直.
沒錯.
以正直去判斷世上的謙卑人.
以口中的仗擊打世界.
以嘴裡的氣.
殺戮惡人.
公義必當他的腰帶.
信實.
必當他脅下的帶子.
就再出現這個光景.
這次就連同.
公義管治的日子.
來講到.
世界的秩序.
恢復.
到底.
是不是墮落之前的樂園的光景呢?.
猛獸原來不兇猛的.
是墮落之後.
亞當夏娃的沉淪之後.
野獸.
動物.
變成兇猛的.
聖經沒有交代.
我們不知道.
我們看到的.
自出生之後就發現.
這個世界有分家禽.
家畜.
可以養的.
但有些是不可以養的.
就算熊貓原來也是兇猛的.

$^{1281}$雖然養到肥了.
但其實也很兇猛.
更何況.
大自然兇猛的獅子.
柴狼.
會令到他們更兇猛.
所以在第十一章和第六十五章.
不斷出現這些光景.
似乎不是一種象徵.
似乎是一個真是有管治.
而且是真正.
這位彌賽亞.
耶西的子孫.
在大衛的寶座上.
去管治的.
那個光景.
還要再描述.
連那個大自然界.
也有一個很戲劇性的扭轉.
可以去到一個.
所有的事情.
都是去到一個.
完全的恢復.
這樣.
不過再加多一樣東西就是.
第九節.
這是十一章的第九節.
在我的聖山的遍處.
這一切都不傷人不害物.
因為認識而來和華的知識.
要充滿遍地.
好像水充滿洋海一般.
又是和整個的.
那個國度裡面.
最重要的.
就不是你去到.
獅子可以和你玩.
去吃草.
不是這樣.
而是因為那個更新.

$^{1321}$是有神的知識.
主的道.
遍佈全地.
這樣的一個.
管治的.
好的國度.
我不敢說是千禧年.
但是如果我們將這些經文.
放在千禧年.
那個物質的千禧年.
現在任何一個日子.
你都不會找到地上.
這樣的光景.
天堂根本不需要這樣的光景.
不需要這樣的描述.
所以這些不是天堂裡面的經文.
不需要在永恆裡面去應驗.
這些經文.
但是在地上完全找不到.
有哪一年.
有幾時會這樣.
有的只能夠在千禧年裡面.
這些日子會發生.
所以.
又不是真的誇張.
真的有那個含義在其中.
我們再揭多一段.
就是25章的.
都是二次世界書.
25章.
二次世界書真的講了很多.
第25章的.
我們就第一次開始看.
其實嚴格來說是第八和第九節.
第25章.
二次世界書.
(念二次世界書).
你作貧窮人的保障 作困乏人急難中的保障.
作暴風的這個地方 作被炎熱的陰涼.
你要壓制外邦人的喧嘩 保障乾燥地的熱氣下落.

$^{1361}$禁止強暴人的哀歌 保障熱氣被雲影消化.
是提到那個的解救啦 第六節.
在這個山上 萬軍之耶和華必為萬民用肥金.
即肥美的 可能當時是肉類比較少.
有肉食就是一個豐收 還有是一個很重大的日子.
即是在山區就是這樣 特別是少數的民族.
有肉食是很矜貴 即是用肥金設擺筵席.
用陳酒 即不是一些新釀的酒 還有滿水的肥金.
即骨髓滿了 即是豐盛的筵席.
並澄清的陳酒 即是那些酒很靚的.
用很靚的水去釀 設擺筵席.
他猶必在這山上 除滅遮蓋萬民的墨.
和遮蔽萬國蒙面的柏子 即是那些修辭.
他已經吞滅死亡 直到永遠.
主耶和華必擦去國人臉上的眼淚.
又除掉普天下他百姓的羞辱.
因為這是耶和華說的.
到那日人必說 看啊 這是我們的神.
我們素來等候他 他必拯救我們.
這是耶和華 我們素來等候他.
我們必因他的救恩 歡喜快樂.
這是第六到第九節.
是會有神預備的豐盛的筵席.
到了那天 我們一路去.
即是當時的人 將來的人.
是會這樣來等候.
在神的歡恩 救恩裡面來快樂.
來享受豐富的筵席.
這樣的描述 又是再次說到.
舊約的先知所預言.
他們所看到的將來會有豐盛的筵席.
我們套入新約的啟示錄.
就會有這樣的筵席.
接著我們可以揭開還有一兩段.
哈該書第二章 哈該書是尾二的一卷.
紅哈番蓋亞瑪第三卷 尾三的一卷.
哈該書 哈該書只有三章.
撒迦利亞書之前.
哈該書就是重建聖殿的時候的作品.
就是索羅巴伯和大祭司約書亞.

$^{1401}$第一次他們回歸的時候.
哈該都跟回去重建聖殿.
立了根基之後 因為政治的干預.
停了15年 後來神的靈感動哈該.
和感動當時的撒迦利亞兩位先知.
他們就一起起來勉勵百姓要建造.
不要怕敵人 所以輾轉.
用了不久的時間 聖殿就建好了.
到了第二章 講到將來的聖殿.
這個又是要放在經文裡面.
對不起 我搞錯了 應該是哈巴谷書.
應該是第二章的哈巴谷書才對.
不是哈該書 蓋就是蓋亞瑪的蓋.
第二章哈巴谷書.
第二章我讀給大家聽 又是一個對比.
很強烈的對比.
因為哈巴谷有些不是很甘心.
神興起那些迦勒底人來攻擊耶路撒冷.
他們就看不過眼 為什麼會這樣呢?.
他們就立在望樓上 其實是看神怎麼做.
看你怎麼來解救 用一些惡人.
去攻擊我們耶路撒冷的百姓.
所以他們做這個望樓不是去禱告.
而是去看神怎麼去收拾他自己定的規矩.
他說漢耶和對我說什麼話.
我可用什麼話向他訴冤.
他對我說 將這默示明明的寫在板上.
使得讀的人容易讀.
因為這默示有一定的日期.
快要應驗 並不虛荒.
雖然延遲 還要等候.
因為必然臨到 不再遲延了.
第四節 迦勒底人自高自大 心不正直.
為二人恩順得生.
這一句聖經 哈巴菊沒想到會這麼受歡迎.
在新約 只是這句聖經 保羅也印了三次.
迦勒泰書 羅馬書 希伯來書.
希伯來書未必是保羅寫的.
但新約聖經印了三次 為二人恩順得生.
這句聖經也影響了宗教改革.

$^{1441}$馬丁路德很重要的整個想法.
扭轉整個當時的宗教體系.
不靠行為 靠信心.
二人不是因行為得救 是因順得生.
是哈巴菊發明的.
是他得到神的啟示.
當時的意思是.
神興起惡人來懲罰惡人.
第五節 迦勒底人因酒鬼詐狂傲 不住在家中.
擴充心慾 好像陰間.
他如死不能知足.
聚集萬國 堆積萬民 都歸自己.
這些國的民豈不都要提起.
詩歌並俗語 幾次他說.
和哉 迦勒底人 你爭天不屬自己的財物.
多多取人的當頭.
即那些抵押品.
要到幾時為止呢.
咬傷你的 豈不忽然起來.
妖害你的 豈不興起.
你就作他們的老物嗎.
因你搶奪許多的國.
其實是嘲笑這些迦勒底人.
他們那些的暴行.
但是神就興起他們.
第八節 因你搶奪許多的國.
殺人流血 向國內的城.
並城中一切居民施行強暴.
所以各國剩下的民必搶奪你.
為本家積蓄不義之財.
在高處踏窩.
指望免災的有禍了.
你圖謀剪除多國的民.
犯了罪.
使到你的家蒙羞.
自害幾命.
牆裡的石頭必呼叫.
房內的棟樑必應聲.
以人血建城.
以罪孽立一的有禍了.

$^{1481}$眾民所勞碌得來的被火焚燒.
列國由勞佛而得的歸於虛空.
不都事出於萬軍之耶和華麼.
但是要留意.
突然間第十四節.
認識耶和華榮耀的知識.
要充滿遍地.
好像水充滿洋海一般.
這就告訴我們.
會有這樣的光景.
但是在目前.
哈巴菊所預言看到的.
就是迦勒底人被神興起.
其實他們本質是很邪惡的.
但是去對付國內耶路撒冷犯罪的人.
這就令哈巴菊覺得不思其解.
一位公義的神為何會興起.
容許罪惡犯罪的人.
會對付自己的國民呢.
但是有信心的人.
這段聖經大概的意思就是.
二人因信就得生了.
就是你憑信去相信.
神的工作是不會不義的.
神的作為是有祂的時間.
是有祂的目的.
甚至有祂的手段的.
但是如果我們認識神.
榮耀的知識.
我們自己就會得到那個生.
生就是在神面前有信心看到局勢.
我們看到這些短暫的過去之後.
就會有真正的.
充滿全地的神的知識.
榮耀的知識.
和像洋海一樣被整個大地.
都充滿了神的管治帶來的知識.
榮耀.
和一定是說管治帶來的和平.
所以這段聖經雖然偶然出現.

$^{1521}$但一定是遙指不是當時哈巴菊年代的光景.
這已經是很後期的階段的先知審判.
和預言的責備.
一度讀了這麼多.
千禧年是說災難之後的地球上的太平一千年.
這是千禧年要說的內涵.
無千禧年是說耶穌還沒回來之前.
耶穌升天之後.
就會有最後的耶穌回來之前的一千年.
在地球上的太平盛世.
越來越趨向太平.
但記住是耶穌沒有回來.
但地球上已經有一千年會好的.
這叫後千.
但無千是什麼都沒有.
總之就是耶穌第二次回來就是了.
地球是OK的.
但千禧年前派就說不是.
有大災難.
耶穌升天之後.
教會的時代月趨是黑暗.
有災難.
以色列也是關鍵.
會被萬民攻打.
接著耶穌第二次回來.
就腳踏橄欖山.
接著就開始一千年的希年.
這就是千禧年前派不斷說的.
一千年的好日子.
在地球上發生的好日子.
再加上啟示錄.
正正就是因為我們相信啟示錄的寓言.
魔鬼擾亂天下的.
他被捆綁囚禁在無底坑.
好像功夫熊貓那個.
我又忘記名字了.
那個惡魔.
那個不知道叫什麼名字.
我要想一下.
功夫熊貓的那個.

$^{1561}$被囚禁了.
被千年鎖鏈鎖住.
不知為什麼無意中被解了出來.
又出來為禍.
這裡告訴我們.
可以再引多兩段的.
詩篇的72篇.
72篇應該是所羅門的詩.
所以這些是幾千年前的作品.
他們寓言的美好光景.
如果新約有這樣的寓言.
有千禧年.
其實我們這樣信.
也不是天馬行空.
耶穌再回來.
去改變那個宇宙萬物.
也不是沒有可能.
第一次降臨做BB的.
那些人.
他們的心是有可能.
去改變這個宇宙萬物.
他們的心是有可能.
第一次降臨做BB的.
耶穌.
都可以走在海上.
可以五餅二魚.
可以叫拉薩路復活.
他都可以稍微行駛.
屬天的那個能力.
要改變整個.
宇宙.
整個地球.
又為什麼不可能呢.
如果我們這樣理解.
就覺得千禧年不一定是一個.
好像天方夜譚的洪水猛獸.
或者是一個.
很無厘頭的神學.
相反.
我自己看了這麼多經文.

$^{1601}$我都仍然.
到今天那幾十年.
信了主.
我都認為我自己是.
千禧年前派的擁護者.
當然我們見不見到.
是否會進入就是另一個問題.
一定會進入.
因為我們是復活回來的.
我們是在生的.
才會改變.
我們死了的將來.
在哪裡呢.
可能是和合石.
在那裡復活.
會和耶穌一起.
在地上一千年了.
千禧年了.
我們將來會有一個榮耀的身體在千禧年.
所以是所有的聖徒.
就有一千年在地上.
但無千禧年和後千.
我們是直接回天家.
我們是有一千年之後才回天家.
所以分別就是這樣.
七十二篇.
詩篇.
所羅門的詩.
我稍微讀一讀很快.
神下求你將判斷的權並賜給王.
將公義賜給王的兒子.
他要按公義審判你的民.
按公平審判你的困苦人.
大山小山都要因公義.
使民得享平安.
他必為民中的困苦人申冤.
拯救窮乏之輩.
壓碎了欺壓人的.
太陽還存.
月亮還在.

$^{1641}$人要敬畏你.
直到萬代.
這個好像有詩歌唱的.
他必降臨象雨.
降在已割的草地.
如甘霖滋潤大田地.
在他的日子.
二人要發旺.
大有平安.
保障月亮長存.
他要執掌權並.
從這海直到那海.
從大河直到地極.
住在曠野的.
必在他面前下拜.
他的仇敵必要添土.
他施和海島的王要進貢.
士巴和西巴的王要獻禮物.
士巴女王就是向所羅門獻禮物.
就是好像有人.
說是那個.
阿丁灣出口.
說是那個.
那個叫什麼.
那個叫什麼地方.
也門.
就是士巴女王的原居地.
就是去到很遠.
你想想.
當時沒有車沒有船.
去到耶路撒冷.
諸王都要叩拜他.
萬國都要事奉他.
因為.
窮乏人.
呼求的時候他要答救.
沒有人幫助的困苦人.
他也要答救.
他要連率行寒和.
窮乏的人.

$^{1681}$拯救窮苦人的性命.
他要救贖他們脫離欺壓.
和強暴.
他們的血在他眼中.
看為寶貴.
他們要存活.
供給他.
人要常常為他禱告.
終日稱頌他.
在地的山頂上.
五谷必然茂盛.
所結的谷實.
要響動如尼巴嫩的樹林.
城裡的人.
要發旺如地上的草.
他的名要傳到永遠.
要流傳如日之久.
人要因他蒙福.
萬國要稱他有福.
獨行其事的耶和華.
以色列的神應當稱頌.
他榮耀的名也當稱頌.
直到永遠.
願他的榮耀充滿全地.
阿們.
很多人說這篇詩篇.
在所羅門的管治下.
未曾有過.
一個.
太平盛世的管治.
這裡說的是一個.
非常美好的.
一個光景.
一個國度的平安和和平.
為所羅門王祈禱.
但是.
這裡告訴我們.
這是他自己對於管治的期待.
所以有些人將這篇詩篇.
歸入所羅門將來的後裔.

$^{1721}$就是這位主耶穌.
在地上千禧年的管治日子.
跟大家多揭一段.
因為已經到十了.
我們可能要看新約了.
舊約我不再揭了.
其實撒迦利亞書還有一段.
我不揭了.
要看哥林多前書.
第十五章.
復活的次序.
保羅在什麼情況下.
講主耶穌的復活呢.
是在說哥林多人不相信死了會復活.
第十五章.
首先第一節.
第三節.
給大家有些印象.
弟兄們.
我如今十四章處理了.
在教會裡的恩賜的大小.
十四章特別強調.
講方言和先知講道的對比.
和分別.
有什麼好處.
和應該怎樣去追求.
講了背後的價值觀.
在公共崇拜裡.
公共聚集裡.
我們應該羨慕什麼呢.
應該羨慕追求先知的講道.
每個人都聽得明白.
每個人都受勸.
不是去講一些自己明白.
自己享受.
用廣東話.
自己high就開心了.
所以崇拜自己high.
是一個很低級的崇拜.
你不要以為這樣就厲害了.

$^{1761}$不是.
是應該眾人都有阿們.
眾人都是聽明白.
有一個決心去進行.
這是一個最優質的崇拜.
這是公共崇拜十四章裡說的.
尤其是先知的靈.
是馴服先知.
很重要的.
什麼叫先知的靈馴服先知.
就是說你不可以在公共崇拜裡.
你自己很high.
現在很手舞足蹈.
現在被聖靈感動.
我要做什麼.
出位的事.
不應該這樣.
這個靈馴服先知.
神不是叫人混亂.
是叫人安靜.
所以我們就應該要留意這些.
很重要的價值觀.
突然十五章有一個很大的責備.
這次不是討論哪些比較好.
不是說港獨和講方言.
沒錯.
不要禁止講方言.
但是更加應該要羨慕港獨.
不是這樣的對比.
是一個很重的責備.
就是他們不信死了的人會復活.
第十五章.
弟兄們,我如今把先前所傳給你們的福音告訴你們知道.
所以什麼是福音呢.
阿保羅就定義了.
這個福音你們都領受了.
又靠著站立不住.
所以是一直持守.
我們所信.
我們所信才能得到永生.

$^{1801}$不是一次得救.
就永遠得救.
是一次的相信.
一直的相信.
所以是要持守的.
要站立得住.
並且你們又不是徒然的相信.
能以持守我所傳給你們的.
講了兩次持守.
就必這福音得救.
那究竟是什麼福音呢.
第三節.
我當日所領受.
又傳給你們的.
第一就是基督照聖經所說.
為我們的罪死了.
而且埋葬了.
又照聖經所說.
第三天復活了.
那《小路信經》就跟足了.
完全沒問題了.
並且顯給基法漢.
然後顯給十二使徒漢.
後來一時顯給五百多弟兄漢.
五百多人.
五百多人是什麼概念呢.
禮堂三百人.
三百個位.
加上輔堂.
加上所有的房間.
都是四百多人.
可能四百三.
都不夠五百.
同一時間見到耶穌.
那你可以想像得到.
是很震撼的一件事.
五百人見到耶穌復活.
其中一大半到如今還在.
寫哥林多前書的時候.
保羅都遇見過.

$^{1841}$這五百個其中那些.
親眼見過耶穌的弟兄姊妹.
你可以想像得到.
是多麼的一件特別的事.
卻也有已經垂了的.
即是死了.
第七節.
以後顯給雅各漢.
再顯給眾使徒漢.
到末了.
也顯給我漢.
我如同未到產期而生的人一般.
原因是他逼迫耶穌.
所以他講到自己這麼渺小.
不配稱為使徒.
因為我從前逼迫神的教會.
即是說他都配得到耶穌的顯現.
他本來不配.
根本不應該向他顯現之餘.
應該懲罰這個保羅的.
但是保羅成為了可以認識見到耶穌的一個人.
所以他的生命有一個很大的改變.
所以我們都求神給我們.
都是大家認真想學聖經的弟兄姊妹.
其實我們應該羨慕耶穌.
你讓我見一下你.
你用你的方式讓我見一下你.
我們這樣的祈禱是沒有錯的.
我們見到主的榮耀.
我們的生命會改變.
如果我們只停留在信念.
一些道理.
有時道理就還道理.
生活就還生活.
但是如果我們想生活和道理有一個結合.
我們反而應該有這樣的祈禱.
這些不是妄求.
不是妄想的不拉肩的祈禱.
而是你應該見到神的榮耀.
應該在心靈裡面.

$^{1881}$所以我們應該要學習靈修.
每天用一點時間自己安靜.
默想一下神.
讀聖經.
有一個渴慕神.
今天我很想見你的面.
你用你的方式讓我顯現.
讓我認識你.
可以是大自然.
可以是超自然.
可以是你的生活.
但是給你一種新的眼光.
新的心靈的感動.
都可以的.
所以我認識一個傳道人.
他一路上班.
突然間坐地鐵.
見到很多人下地獄.
其實每個人都只是坐地鐵.
下地獄坐車.
每天都是這樣.
他那次不同.
他裡面有很大的感動.
他就一直去裝備等候祈禱.
令到他自己熱心.
越來越印證他.
是神呼召他做傳道.
所以原來.
可以這樣去祈禱.
因為這個不是為了.
自己的興趣.
而是為了和神有更好的關係.
所以保羅說.
我去到幕後一個不配的使徒.
他都向我顯現.
即是說原來我們都可以.
這樣去渴慕.
認識神.
當然了.
整段不是說誰見到誰見不到.

$^{1921}$而是第十二節這一段.
既然全基督是從死裡復活.
為什麼在你們中間.
有人說.
沒有死人復活呢.
如果沒有死人復活.
基督也沒有復活.
因為基督是一個人.
在物質的.
受限的那個秩序裡面.
他這樣死.
即是說他進入我們人類的時空.
是他自己製造的.
那個秩序.
他要尊重這個秩序.
他在這個秩序底下.
被人誣告 被人殺害.
但是正正就是.
他的神能.
超越了他自己製造的這個秩序.
他可以復活.
他可以升天.
而且是一個永活的身體的升天.
這個就是信仰.
這個就是我們基督徒.
不能夠沒有的.
那個凝信.
每個星期我們崇拜的凝信.
就是相信.
主耶穌是.
從死裡面復活.
那一位再來的君王.
那所以.
如果沒有死人復活.
基督就沒有復活.
如果基督沒有復活.
我們所傳的就是謊言.
你們所信的.
也是謊言.
並且明顯.

$^{1961}$我們是為神.
妄作見證.
我們見證神.
是叫基督復活的.
如果死人沒有復活.
神也沒有叫基督復活.
因為死人如果不復活.
基督也就沒有復活.
基督如果沒有復活.
你們的信便是陶然第二次說.
你們仍在罪裡.
因為基督是擔待我們的罪死.
就是在基督裡面.
睡了的人也要滅亡.
所以寶來.
這裡是說得很清楚.
寶來這裡是說得很重.
很重的真理.
的那個責備.
我們如果靠基督.
只是今生有指望.
就是回來教會.
我們有排排舞.
我們有拉筋班.
還有高球.
那你想想你的人生.
在錦繡堂.
這麼多東西玩.
如果我們今生只有這些指望.
而我們根本不信基督會復活.
那就比眾人更可憐.
其實就不可憐.
你這麼多活動.
不過.
和其他人一樣.
完全都是為自己的罪.
面對神的將來審判.
我們根本就沒有信仰可言.
所以如果教會沒有信仰.
就很大件事.

$^{2001}$因為我們藉著那些.
只不過是一個媒介.
我們將那些媒介變成為終極的目的.
就很大件事.
媒介是帶我們去看.
真正的目的.
就是基督.
所以教會是基督.
教會不是活動.
我們要搞清楚.
教會不是回來唱聖詩.
就叫做是教會.
唱聖詩是為了領導弟兄姊妹.
一起敬拜神.
所以我們千萬不要將那些本末保置.
團契是為了.
不是我們一群人很開心地吃東西.
是我們一群人.
彼此相交.
更加團結.
更加去傳福音.
事奉神.
第二十節.
但基督已經從死裡復活.
就是這一段.
成為了睡了之人初熟的果子.
死既然是因為一人而來.
死人復活.
也都是因為一人而來.
死人復活.
也是因為一人而來.
在亞當裡.
眾人都死了.
照樣.
in the same manner.
在基督裡.
眾人也都要復活.
在說眾人.
不是一個半個.
第23節.

$^{2041}$但.
這裡開始.
但各人是按著自己的次序復活.
初熟的果子.
就是基督.
他首先復活.
他現在在天上.
以後在他來的時候.
留意一下.
那些熟基督的.
不熟基督的還沒復活.
再後末期到了.
是不是說千禧年完結呢.
那時.
基督既將一切執政掌權.
有能都毀滅.
就把國交予父神.
因為基督必作王.
等神.
把一切的仇敵都放在他的腳下.
去到末了.
所毀滅的仇敵就死了.
這個就是回應.
啟示了的.
20章.
那個白色大寶座.
是過了一千年之後的.
那個魔鬼.
被釋放.
接著去到白色大寶座.
死就要交出.
其中的死人.
死亡都要被關進陰間.
有沒有留意20章那裡.
現在這裡臨前.
保羅是說那個次序.
是有一段時間.
基督在地上掌權.
接著去到盡末後.
死都要交出.

$^{2081}$都要被制服.
盡末了.
所毀滅的仇敵.
就是死了.
因為經上記著說.
服在他的腳下.
這是詩篇.
萬物都服了他.
明顯那叫萬物服他的.
不在其來了.
萬物既服了他.
那時子也要自己服那叫萬物服他的.
就是這位造物主.
叫神在萬物之上.
為萬物之主.
這段聖經.
第23到28節.
其實我們應該理解為.
阿爸的死.
其實我們應該理解為.
保羅說的復活次序的千禧年.
這樣理解.
會更加準確.
不只是說一個復活.
有復活那麼簡單.
因為他說的按次序復活.
是有它的考究.
所以如果這段聖經.
fit在啟示錄的20章.
其實就剛剛解得通.
那些疑點.
所以這段是一個很強烈的.
對千禧年的支持.
時間不夠.
還有兩段.
因為喉嚨開始有點乾.
還有兩段.
可能我下星期再補充.
因為那裡也有些關鍵.
就是啟示錄的13章.

$^{2121}$說到地獸海獸.
邪惡的三位一體.
和20章的關係.
下一次我再補充一點.
我們下星期.
就會正式開始.
這個神學.
就是說.
宗教改革.
由1517年.
應該去到1536年.
其實有幾個重要的向度.
深入淺出.
我不是很熟悉.
但也研究過一點.
和大家有些分享.
可能環繞著.
加爾文和馬丁路德.
他有些前因.
有些再早一點的神學家.
和一些殉道者.
比如約翰·胡斯.
可能大家沒有聽過.
又有當時的慈雲尼.
麥蘭頓.
威克里夫.
大家可能聽過.
這些就是那些前人.
奠好根基.
馬丁路德.
當然是他自己身受其害.
他自己去.
神學的扭轉.
才帶來一個.
整個時代的改變.
所以.
時勢做英雄.
也是英雄.
帶來一個新的時勢.
一個時代.

$^{2161}$所以馬丁路德肯定是.
那個靈魂人物.
對他的人生.
我們也要一點點了解.
他因為懼怕神.
做傳道.
到了晚年.
他一生守獨身.
因為他是天主教的神父.
但是晚年.
真的津津樂道.
晚年就娶老婆.
和一個修女結婚.
你可以想像到.
對他來說是一個.
履行宗教改革.
那些附帶的條件.
或者後遺症.
因為這麼多人跟了你.
離開天主教.
很多修女沒有人要.
馬丁路德就娶一個修女.
做太太.
一來神父可以結婚.
神職人員可以結婚.
二來就是他要照顧這些.
孤兒寡婦.
因為天主教.
不是成為他們的任備之後.
需要有.
叫做還俗.
其實就沒有還俗的意思.
而是為了.
真的要你.
帶來的改變.
要向那些平民百姓.
負責任.
這個就很津津樂道.
很感動.
因為可以想像到.

$^{2201}$宗教改革會帶來.
很多人根本的改變.
所以如果我們.
信耶穌 而他真的有很大的改變.
而他面對一些.
困難 比如被家人.
趕走.
面對一些.
例如在澳門工作的.
我真的不再在賭場工作了.
我要找一些很低收入的工作.
教會都要明白.
幫助 諒解.
支持等等.
這些就是信仰.
帶來的社會的改變.
這些其實挺有趣的.
挺值得我們大家一起去思想.
或者反過來.
最後我們自己.
有沒有帶來.
我們自己的家 或者.
我們的社會有改變呢.
多了你信耶穌 有沒有改變.
那個社會呢.
至少我們多一個人.
沒有那麼暴戾.
多一個人沒有那麼.
怨天尤人.
多了讚美.
這些都是.
我們應該要想的.
就是怎樣能夠造福.
那個時代的人.
這次我們用了.
四次講的神學.
這個末世論.
給大家有一個這樣的眼光.
開一開啟 原來這個世界.
有人相信是沒有千禧年的.

$^{2241}$有人相信是有後千的.
但是前千呢.
原來是難相信的.
因為主耶穌又要再回來地球.
又要在地球開一個.
一千年的王國.
有些人和他在一起.
這個就是聖經二十章.
啟示錄的經文的字面意思.
看看你相不相信.
或者你有沒有更強烈的.
聖經的支持.
你認為一定是不會有千禧年的.
只是有.
耶穌第一次降臨.
第二次降臨.
就是開始新天新地.
絕對沒有說錯對的.
因為沒有百分之一百的.
除非你不相信耶穌再回來.
那這個就是百分之一百.
我們一定要留意這個.
我們一起祈禱.
感恩天父給我們用了四堂.
去講完這個末世論.
當然還有一些細節.
我們都可以再補充.
都能夠透過我們在.
每一堂的一個半小時.
都將很多的經文.
和他們背後的主張.
都簡單地講述了.
盼望都給我們對於主宰的.
有一個很根本的渴慕.
我們的人生.
無論怎樣摧殘.
如果沒有了復活.
都是徒然的.
我們很盼望我們能夠.
在很有限的.

$^{2281}$我們人類地上的生活.
我們去為主的福音而活.
我們很盼望.
我們能夠成為.
可以主所用的人.
去幫助帶領人.
進入神的救恩.
多謝您給我們今天晚上的安靜.
默想神的話語.
我們這樣禱告.
奉耶穌基督的名.
阿們.
基督的名.
\newpage



\section{}
\label{sec:aTaCCC_F45U}
\textbf{2024-10-10 周四晚查經班《改革之路:過去與未來.(1)》 講員:曾裔貴牧師}
\newline
\newline
連結: \href{https://youtube.com/watch?v=aTaCCC-F45U}{\texttt{https://youtube.com/watch?v=aTaCCC-F45U}} ~~~~ 語音日期: 2024-10-12
\newline
\newline
\hyperref[sec:H4W4IN0LZRY]{\small{< < < PREV SERMON < < <}}
~
\hyperref[sec:index_chronic]{\small{[返順時目]}}
~
\hyperref[sec:index_scriptual]{\small{[返順卷目]}}
~
\hyperref[sec:Sb_mbdX9Sa4]{\small{> > > NEXT SERMON > > >}}
\newline
\newline
$^{1}$我們一起有簡單的禱告.
又一次來到神面前來到有查經班的時間.
我們盼望我們研讀這個宗教改革.
馬丁路德和嘉爾文他們那個中世紀時代.
盼望我們只是一個很基礎的形式.
我們不可能像神學院那樣做一些大的專文.
而是我們能夠可能就這樣聽一下.
又反思一下和我們現代有什麼關係.
因為已經是500年前的事.
福音派都是由這樣而開始.
一直到今天一直都是需要那個改革.
在我們大家在教會的帶領裡面的.
盼望神給我們明白什麼是宗教改革.
以致我們大家都能夠很簡單的掌握到.
也都給我們有趣味的去學習.
我們仰望禱告奉耶穌基督的名 阿們.
好 很抱歉 大家要靜音.
因為怕影響課堂.
我就幫大家做了靜音.
那就開一個新題目.
可能有些人未必有興趣.
或者都未必會覺得容易.
但是我盡量很簡單.
當是趣味性的給大家去認識.
就當是大家聽一下教會的歷史就可以了.
所以就不要太過擔心.
不用做功課的 不用寫文的.
就給大家知道一下來龍去脈.
就當是一個了解.
當然其中重要的就是.
因為我用的題目今早都想了一段時間的.
都問了一些同學的意見.
宗教改革這麼多東西講.
用什麼題目好呢?.
改革的原因和過去.
和對將來的影響.
就盼望這個都是大家一起去.
有個機會去碰一下教會的歷史.
500年前的.
一定是相當遙遠.

$^{41}$但是給我們都有一點點的輪廓.
就幫助我們對今天我們去讀的時間.
我們又會有一些的了解.
會有很多的資料給大家.
大家當是興趣有空看一下.
沒空就隨便他.
那就無所謂的.
當然你自己再有盡心的研讀.
就更加好.
但是上個星期答應了.
我們要處理最後一個問題.
就是啟示錄的13,19章.
怎樣來印證20章的災前被提.
這個是最後的那個point.
都想給大家一個完整.
因為13,19章是講到海獸和地獸.
就是撒旦和那些假的宗教.
就是邪惡的三位一體.
去到20章就再一次講.
哪些人沒有拜到獸像.
和向獸來跪拜的.
就可以復活.
可以首先復活.
所以那裡就間接印證了.
是有大災難.
就不是災後的被提.
更加不是那個所謂的後遷.
和毛遷希年.
是講那個社會的世界是越趨完美.
然後主耶穌回來.
就是說基督的掌權.
是令到整個世界越趨美好的.
但是似乎啟示錄告訴我們就不是這樣.
所以我要幫大家翻這兩段經文.
稍後我們就會進入宗教改革的探討.
我們先看這13章的啟示錄.
或者不要了.
看20章比較好.
因為20章就講明了.
哪些人是第一次復活.

$^{81}$死人那些未信的,不信的是不能復活的.
這個是前天最重要的重點.
所以我自己到現在.
我還是相信有分開兩次的主耶穌降臨.
就是第一次祂再回來.
就是第二次祂再回來.
就是接所有信的人被提復活.
但是那些死了的,不信的是不能復活的.
要等待第二次的白色大寶座再復活的.
這個是啟示錄20章那個微妙的.
經文的那個世界字面的意思.
如果你不信字面的意思.
那你就是後遷和無遷的了.
無遷欺憐的了.
但是如果你信字面的意思.
那你就是前遷的了.
就是有遷欺憐的了.
應該這麼說.
我們看回第20章的第4節.
因為第3節就講了.
撒旦是被捆綁在無底坑.
一位天使就可以將他捆綁的了.
所以原來他所謂的權力不是無限的.
他不是神來的.
一位天使就可以將他捆綁的了.
一千年.
前遷可以說是幾根據字面.
所以你不根據字面是沒有問題的.
不過你要合理化你自己不根據字面.
你要解釋到為什麼你不根據字面.
那你要有一套很好的完整的解釋.
那你才能令到我去信服.
但是現在這裡就字面告訴我們.
為什麼不可以字面呢?.
近代的人就覺得不需要按字面去理解.
那就不會有遷欺憐了.
那個爭論就是這樣.
留意一下第4節.
我等一下會將大家帶上去第13章.
就是之前約翰講過.

$^{121}$有將來的日子.
基督和那些假使徒會出現.
那些假的宗教.
就是拜那個666.
第4節.
我又看見幾個寶座也有坐在上面的.
並有審判的權柄賜給他們.
我又看見那些因為給耶穌作見證.
接著留意.
並為神之道被斬者的靈魂.
和那沒有拜過獸與獸匠.
也沒有在額上和手上.
受過他印記之人的靈魂.
他們都復活了.
接著留意.
是與基督一同作王一千年.
復活.
接著就在地上與基督作王一千年.
哪些人呢?.
什麼條件呢?.
哪些人會復活呢?.
就是很明顯講到.
在災難的時候.
沒有拜,獸和沒有受印記.
不肯受印記.
不是沒有受.
是不肯受印記的人.
就是忠於信仰的基督徒.
那就是.
他們當然都死了.
因為耶穌還沒回來.
都是按照著神給你的肉身的年紀.
那你就離開世界了.
但是呢.
你在你的人生.
是沒有拜,獸和受印記的.
接著這些人就復活了.
與基督一同作王一千年.
記住.
接著再一次強調.

$^{161}$這只不過是那個理論.
其實不是我的理論.
是極多人對這個前線的看法.
留意一下後面那裡.
這裡提到這是頭一次的復活.
其餘的死人.
還沒有復活.
直到一千年完了.
在頭一次復活的.
一定是因為與主同在.
因為在再生的時候.
是沒有拜過獸像.
和認主是基督的.
也沒有受印記的.
所以是有福的.
他們有份於第一次的復活.
就是有福了.
聖潔了.
接著還要強調.
第二次的死.
就不是第二次的復活.
因為他們已經復活了.
第二次的死.
就是像夏文說的.
白色大寶座的.
那幫第二次復活.
一復活就馬上死的人.
當然是審他之後死的.
所以是有兩次復活.
而中間的時間.
無長無短就是一千年.
所以前天.
其實那個理據是強過無遷和後遷的.
而且是講明了.
是忠於信仰沒有拜獸像的.
所以現在我們就要回頭.
就是第十三章.
這些我們曾經說啟示錄的時候.
已經說過.
所以大家上YouTube.

$^{201}$找回啟示錄去聽.
就會比較完整一些.
十三章.
因為我們很快的跳一跳.
我們因為想帶出二十章的重點.
就是約翰之前已經說過有大災難.
有為信仰而受逼迫的信仰群體.
教會要面對迪基督將來的欺凌.
但是原來有一幫人.
是忠於信仰不跟隨的.
十三章.
我又看見一個獸從海中上來.
就是所謂的海獸.
有十角七頭.
在十角上戴著十個冠冕.
七頭上有蝕毒的名號.
我所看見的獸形狀像炮.
腳像紅的腳.
口像獅子的口.
相當怪獸.
那龍將自己的能力.
座位和大權並都吸了他.
所以這是邪惡的三位一體.
這個龍就代表著聖父聖子聖靈.
所以現在這個龍就是魔鬼.
就將他的權力全部移交給海獸.
就是迪基督.
留意一下.
第三節.
我看見獸的七頭中.
有一個似乎受了死傷.
那死傷卻醫好了.
全地的人都稀奇.
跟從那獸.
其實描述的是一個人.
所以這裡告訴我們.
他的形態.
他背後的.
你可以說他裡面的靈.
一定是和邪惡和魔鬼拉上關係的.

$^{241}$接著這裡告訴我們.
有一個似乎受了死傷.
那死傷卻醫好了.
全地的人都稀奇.
跟從那獸.
又拜那龍.
因為他將自己的權力吸了獸.
也拜獸說.
誰能比這獸.
誰能與他交戰呢?.
所以是一個聯合.
不過是一個邪惡的聯合.
等於我們敬拜主耶穌.
就是敬拜父了.
所以魔鬼都會自己搞一個邪惡的三位一體.
接著第五節.
「又賜給他說誇大褻瀆華的口.
又有權並賜給他」.
就是這個海獸.
「可以任意而行四十二個月三年半」.
也是啟示錄經常這樣提的.
「獸就開口向神說褻瀆的話.
褻瀆神的名並他的帳幕.
以及那些住在天上的」.
即是對造物主的褻瀆.
咒罵.
以及任意的去罵神.
「又任憑他與聖徒爭戰.
並且得勝.
也把權並賜給他.
制服各族,各民,各方,各國.
凡住在地上.
明志從創世以來.
沒有記載被殺之羔羊生命賊上的人.
都要拜他」.
很特別的.
因為第八節這裡.
去到二十章的後半段.
講到那個白色大寶座.
就是再講那個生命賊.

$^{281}$明志沒有記載生命賊.
就要被拎入火壺.
又是再次提醒.
原來歷世歷代.
因為不只是某個時間.
歷世歷代都有這個大的原則.
就是魔鬼都要帶人去褻瀆神.
否認神 否認祂的存在.
否認祂是那位絕對的榮耀者.
我們要敬拜祂 敬畏祂.
和我們一會兒晚一點.
我們講的中國改革.
當時他們是有一個很強烈的信仰.
當然那個信仰已經去到一個.
很混亂的地步.
怎樣去撥開那個混亂呢?.
就是燒烤已經那些炭堆滿了灰.
你要撥開那些灰.
那些炭火才能再旺盛.
其中一個人就是馬丁路德.
這個我們稍後再講.
這個時候很嚴肅.
凡是在地上名字在創世以來.
沒有記載被殺之羔羊生命冊的人.
都要拜祂.
凡有意的就應當聽.
擄掠人的必被擄掠.
用刀殺人的必被刀殺.
聖徒的忍耐和信心就是在此.
你可以想像得到.
你表面就是拜魔鬼.
實際上原來是令到你的內心.
應該這樣說 反過來.
可能會更加好一點.
就是你表面上就以為自己的信仰.
是你自己可以去選擇.
可以去操控.
所以我不會這麼盲目的拜真神.
但其實背後你自己不知不覺間.
其實你是被魔鬼迷惑.

$^{321}$你對神的絕對和權威和榮耀.
你根本就沒有那個認識.
這樣不會使你變成中立國 中性.
反而令你更加偏向魔鬼的迷惑裡面.
越來越對真理是一個模糊.
讀完十一節下面這一段.
我又看見另有一個獸.
這個就是地獸.
這裡說有兩隻角.
好像羊羔的 所以是假的才知道.
說話很像龍.
他在頭一個獸面前.
就是在海獸面前.
施行頭一個獸所有的權柄.
並且叫地和住在地上的人.
都拜這個死傷醫好的頭一個獸.
而且又能夠行大騎士.
甚至在人面前.
叫火從天降在地上.
他因此給他權柄.
在獸面前能夠行騎士.
就迷惑住在地上的人.
做這麼多事其實都是想迷惑.
所以是三位一體的.
對真理的那個否認.
和要迷惑人不要信真神.
不要信真理.
不要信這個世界有審判.
不要信我們的生命要由神去主宰.
這些妖言惑眾的東西.
其實歷史歷代都有的.
很多有些名人都提出自己的價值觀.
有人信的.
譬如美國Tom Cruise.
他是信科學教的.
很怪的宗教都信的.
他不是一個沒有知識的人.
很厲害.
躋身在整個美國最上流的階層的人.
錢多到.

$^{361}$但是信科學教.
那你就可以想像是極之迷信的一件事.
已經是其中一個例子.
他公開這樣說.
你可以想像到.
假先知就是要帶人去敬拜魔鬼.
用一些.
可能是一些似是而非.
因為他能夠.
他說話好像龍一樣.
說話好像龍.
像魔鬼.
有兩隻角好像羊羔.
像羊羔的其實是魔鬼.
你可以想像到.
這樣的形態的出現.
叫我們非常小心.
似是而非.
會帶來一個使我們無伶兩可.
接著繼續說.
這裡告訴我們第十四節.
他因此給他權並在獸面前能夠行其事.
就迷惑住在地上的人說.
要給那受刀傷還活著的獸.
作個像.
就是幫這個海獸.
其實基督做了一個像.
接著下面就說.
要拜這個像.
第十五節.
又有權並賜給他.
叫獸像有生氣.
並且能夠說話.
那幅相也會說話.
又叫所有不拜獸像的人.
都被殺害.
你就可以想像到.
去到一個大災難.
真是要為信仰而犧牲.
第十六節.

$^{401}$他又叫眾人.
無論大小貧富自主為奴.
都在右手或在額上.
受一個印記.
以前我們小時候紋身.
我肯定被我媽媽和我爸打到飛起.
因為那時候紋身一定是黑社會.
和黑社會有關.
想做一隻龍.
一隻老虎.
現在不同了.
現在什麼圖案都有.
現在是普遍到.
年輕女孩的腿.
都印上了.
我真是很有興趣.
是貼水紙.
還是真的紋上去.
其實很痛的.
是終身的.
但是將來.
一點都不會覺得是一件噁心的事.
反過來.
受印記.
就不能買賣.
除了受印記.
有了壽命.
或者有壽命數目的.
都不得作買賣.
現在上內地沒有支付寶都麻煩.
麻煩而已.
不是不能買賣.
是麻煩而已.
你沒有WeChat都不行.
就是一路方便.
現在講的是生和死的問題.
是你要拜壽.
和要有他的印記.
那個問題.
所以是怎樣會發生呢.

$^{441}$這個我們不是太知道.
第十八節.
在這裡有智慧凡有聰明.
可以算計壽的數目.
因為這是人的數目.
他的數目是666.
我們小時候.
我們都覺得很恐怖.
原來那時候不懂得看.
原來是說那個迪基督成長.
那個迪基督怎樣打他都死不了.
誰發現他都會被殺.
死於非命.
真的是恐怖片.
但是根據這個十三章.
去拍的電影.
當然是電影.
不過我們想說.
啟示錄是一路預言將來的大災難.
十三章肯定是宗教上的大災難.
也可以反過來.
宗教上的堅貞.
和背棄宗教信仰的人的分野.
不是沒有手尾巾.
手尾就去到二十章.
就有手尾巾.
就是拜過壽像的人.
就沒有第一次復活.
就是絞豬毛利粉.
因為你根本就不在絞豬的箱裡面.
絞豬的箱裡面的全部是守神的道.
信主的人.
堅守不拜壽像的人.
第二次復活就有利粉了.
因為那些基督徒已經是與主作王一千年在地上.
第二次復活.
我們剛才還沒讀第二次.
我們現在讀第二次那裡.
因為很明顯是追究十三章那些跟隨魔鬼的人.
二十章後面那段.

$^{481}$第十一節.
二十章的十一節.
當我們從第十三章一跳過來.
我們就終於明白這個問題了.
第十一節.
我又看見一個白色的大寶座.
與坐在上面的.
當然就是我們的造物主真神.
從他面前天地都逃避.
這是最後的末日了.
是沒有任何的躲藏的地方了.
再無可見之處.
所有東西都要在他的面前.
我又看見死了的人.
無論大小都站在寶座前.
按卷展開了.
並且另有一卷展開.
就是生命冊.
死了的人.
所以那個按卷.
肯定就是每個人生前所做的.
要自己現在來到去.
因為是審判.
要自己來到去供出.
自己來到去承擔.
就是說他們一路在人生.
都不依靠羔羊的寶血.
去洗乾淨自己的罪.
那就要靠自己面對了.
某程度上就是馬丁路德.
他對於自己的那個發言.
很明顯的.
他說死了的人.
都憑著這些按卷所記載.
照他們所行的受審判.
所以於是.
因為是最後的末日.
於是海交出其中的死人.
死亡和陰間都要交出其中的死人.
死了的都要被交出來.

$^{521}$他們都照個人的所行.
受審判.
死亡和陰間.
因為已經不需要再存在.
就被泯在火壺裡了.
亦都是魔鬼和假先知.
和獸的那個地方.
這是第20章的第10節.
我們剛才沒有讀到的.
不過現在所有不跟隨神的人.
都要去到魔鬼的那個地方.
就是第10節.
那迷惑他們的魔鬼.
被泯在流亡的火壺裡.
就是獸和假先知所在的地方.
因為獸和假先知.
邪惡的另外那兩位.
19章已經說了.
一早就被泯了.
因為魔鬼要將他拘禁在無底坑一千年.
不是拘禁在火壺.
是拘禁在無底坑.
所以一千年過了.
才要釋放魔鬼出來.
再將他泯入火壺.
這樣的層次.
都很戲劇性.
如果我們拍一套電影.
這樣都很有趣.
現在就有無數的人.
生前不信主.
在他們的人生都敬拜魔鬼.
他們就要面對.
他們自己一生的惡行.
即是沒有主耶穌救贖的惡行.
死亡和陰間.
亦都被泯在火壺裡.
這個火壺就是第二次的死.
每個人只能夠一次死.
但現在是復活之後再死.

$^{561}$是第二次的死.
所有的罪人要復活.
然後被泯入火壺.
就是第二次的死.
所以都幾恐怖.
第十五節.
若有人明志沒記在生命冊上.
他就被泯在火壺裡.
所以弟兄姊妹.
前天的理論.
其實不只是解字面這麼簡單.
是因為十三章已經預言將來.
在宗教信仰上有大迫伯.
有大災難.
魔鬼會將他所有的權.
交付給代理人.
就是迪基督和假先知.
使這個邪惡的地上的三位一體.
在地球上橫行四十二個月.
三年半.
這樣就出現了信仰的大迫伯.
所以跟隨的人.
一定是面對很大的困難.
所以這裡就告訴我們.
將來原來是要審判.
要清算.
要計算這條數.
所以那些以為一生拜偶像.
拜假神 拜魔鬼.
以為可以心安理得.
原來不是.
第一次的復活呼召是沒有他們份的.
叫不到你的名字.
因為帖撒羅尼加前書四章.
主耶穌必定親自從天降臨.
有呼叫的聲音.
天使掌的聲音.
應該有號角聲.
所有的信徒.
應該說在生的.

$^{601}$首先身體改變復活.
然後死了的.
第二批復活.
都是會升到天上.
這個時候那些死了.
但是不信的人就不能復活了.
這個就是我這樣的理解.
其實不是我.
是前天的理論的理解.
所以基於這些經文的好幾段.
我們之前處理過.
應該是上個星期.
處理過好大量的舊約的預言.
講到那些日子.
又不是天堂的日子.
又不是現在地上的日子.
似乎就好像是.
千禧年的一些光景.
一輪半找的一些場景.
所以希望我能夠游說你相信前天.
遲些國民慈牧師.
他在10月31日.
他也開了一六堂的神學講座.
他最後那堂第六堂.
就是講千禧年.
我希望他是前天.
我不知道.
應該上一代多數都是前天.
這個是我今天的補充.
關於千禧年.
因為我們之前講過.
關於千禧年的思想.
現在大家有沒有一點點的問題.
不過可能大家都不方便問了.
可以在那個群組裡.
如果你有分享 有問題.
我們可以慢慢看了之後.
我再去了解一下.
然後我再去解答.
或者私下再WhatsApp我.

$^{641}$也可以的 沒問題的.
或者你有一些見解.
這裡未必夠時間去分享.
你可以公開分享.
又或者你可以是.
其他的渠道去分享都可以的.
如果大家沒有問題.
這個是最後的補充.
然後我們就正式開始.
其實不會太長.
我都希望是.
我看看分享了沒有.
請問是不是已經看到了.
那個PowerPoint.
看到的.
那就行了.
我已經分享了.
宗教改革之路.
都是想了一段時間.
過去和未來.
過去一定是講那個歷史的宗教改革.
那個來龍去脈.
未來就是我們現在都仍然是在改革之路上.
雖然我們不能夠假設.
天主教他們會最終有一個和我們完全的融合.
擁抱我們的教義.
但是呢.
到今天仍然還是有一定的分歧.
而那個分歧一定是由馬丁路德那個時代開始.
到今天仍然是沒有一個完全的融合.
所以大家要知道這一點.
當然撇開了天主教和基督教的分離.
我們自己.
基督教裡面.
我們自己有什麼需要改革呢.
宗教改革有什麼的持守.
他們的教義.
他們為什麼要畏著當時的信念.
要來到去.
甚至乎真的要付上代價.

$^{681}$就是生死的代價.
比如說馬丁路德是周圍被那些教皇.
因為他們都有當時的德國的一些皇帝.
來到去捉拿的.
所以有一些的城邦的一些領袖.
都是一些諸侯的王.
來保護馬丁路德.
所以不只是文字的對話.
大家彼此不同立場.
沒有可能的.
因為涉及到整個天主教的根本利益.
你說我們錯了.
就是說我們是異端的.
這是一件很大的事情.
所以他們當時的持守.
我們就要問.
持守什麼呢.
這是當然的.
這些持守.
是不是今天我們都已經忘記了.
淡化了呢.
這些就是過去與未來.
我們就應該去問.
為什麼我們是福音派呢.
我們就要不斷反思這件事.
其實這樣已經有很多東西要說.
所以一定是深入淺出的.
所以如果你去德國.
你去馬丁路德的故鄉.
你就會見到這個真人的銅像.
就是500年前.
馬丁路德就是這樣的樣子.
你可以去到見到的.
這些網上隨便都找得到的.
給大家來一個小小的了解.
這個就是一個引言的開始.
我們要開始來看.
最重要的就是在1517年.
馬丁路德.
他就是1483年11月10日出生的.

$^{721}$就是說他34歲的時候.
34歲.
居然間.
他本身就是威登堡的大學裡面的教授.
同時間他也是威登堡地區主教的神職人員.
當然他不是那個地方的主教.
34歲的那一年.
1517年的10月31日.
真的是萬聖節.
不知道是不是巧合還是怎樣.
在萬聖節那天.
他就自己私下.
就是這個95條的輪桿.
這樣釘在了威登堡教堂的門口.
不得由自可.
因為.
當然當時他沒有用德文.
用拉丁文來釘.
因為當時的普遍官方宗教的語言.
無論港獨都是用拉丁文的.
所以一般的平民百姓是聽不明白的.
聽不明白都不會解釋給你聽的.
等於錦宣每個星期港獨都是用英文的.
你聽不明白就說對不起.
是你的問題不是教會的問題.
當時他們是這樣想的.
教廷.
因為覺得拉丁文才是唯一的中文語言.
那個聖經.
後來.
當然.
這個95條釘出來之後.
就帶來很多的討論.
因為其實這個95條.
本身不是一個罪狀.
是想引發一個討論.
所以是一個論綱.
是一個叫做physic.
是一個我提出一些我的見解.
大家怎樣去理解這些事情呢?.

$^{761}$當然在這個95條裡面.
其中有幾個大的方向.
教皇不是無誤的.
教皇所宣判的赦罪.
因為他們當時已經是很.
根深蒂固的.
教皇一宣告你無罪.
你就是無罪的.
不過後來因為很實在.
因為他們要建教堂.
沒有錢.
既然教皇可以宣判誰無罪就無罪.
就要想到一個辦法.
就是有一個叫做贖罪券.
indulgence.
現在好像是一元歐羅.
買了一張贖罪券在德國.
就會發售的.
一元歐羅.
就被人看到500年前.
德國人就是這樣的問題.
主要是處理教皇不是無誤.
還有那個赦免其實是無效的.
另外就是.
當然我等下會再講.
所謂的無效就是.
如果那個人不是根據基督的赦罪.
而只是教皇的宣判是不行的.
那就是不應該有教皇成為.
或者主教成為那個代理人.
因為當時的告解都是很普遍的.
所以教皇認為你要怎樣怎樣做呢.
那你就要跟隨的了.
或者我再提前一點.
在馬丁路德的年代.
人是怕死到一個地步是很嚴重的.
所以教會提出的那個煉獄.
是令到很多人很懼怕的.
馬丁路德之前的大概.
其實都差不多100年了.

$^{801}$整個歐洲有一個病大家知不知道.
是死了差不多三分之一的歐洲人.
黑死病.
就是鼠疫.
由亞洲傳去歐洲的.
那個黑死病是死了很多人.
所以歐洲是差不多好像沒有人一樣.
是很恐怖的.
是很怕死亡的.
而且更加有些人將他掛勾.
是神對我們這個時代的一個審判.
所以人就更加懼怕.
所以因為這樣的緣故.
那個煉獄的觀念帶來.
使得那個又是很實在.
在地上有些不治之症不斷的蔓延.
找不到那個解決辦法.
所以你可以想像到.
那個對死亡的意識是非常之強烈的.
所以一定是有很多前因後果.
導致到這些方案提出來.
就是會帶來那個越來越.
變相的一個整個教義的扭曲.
贖罪券起初不是這個意思.
贖罪券起初是不用錢的.
但是後來真的因為打仗.
很多的年年的損耗.
再加上要起教堂和那個修葺.
所以最好的辦法就是販賣贖罪券.
主要不是為自己.
不是為在生的人.
贖罪券的精義那個含義.
是為了已經死去的親人.
你買贖罪券.
那個親人在煉獄裡面.
就可以少受一些災難.
那些刑罰.
這樣的原因.
所以當時的時代.
是深信不疑的.

$^{841}$那些人可以說文字還沒開.
很難用我們500年後的今天.
我們這樣去理解教廷的權威和宣判.
今天聽到.
如果300人準時坐在那裡聽.
很熱切地聽牧者講道.
我想從來都沒有試過.
要不就遲到 要不就早走.
要不就睡覺 要不就在那裡胡思亂想.
要不就看程序表.
其實是看那些人 其他後面的人奉獻.
不是真的想關心甚麼.
很多很多不同的想法.
因為那個權威已經不是像500年前的.
這樣的宗教那個權威.
當然我們不是要追求這種權威.
正正就是要宗教改革.
就是不應該這樣權威.
就不應該人的想法.
人製造出來的傳統.
就凌駕了聖經.
這個是95條最重要的精神.
馬丁路德他一直都捍衛自己的信念.
就是說除非你們在聖經裡面.
找到我有甚麼跟聖經的主張是錯的.
我就一定會承認.
如果不是的話 我要捍衛我自己的立場.
這個是關乎我自己的良心.
面對神 求主幫助.
這個是他的名言.
我是不會修改我的主張的.
所以這個95條.
第一個就是對教皇的權威.
因為他不可以再判人無罪.
你要買他的贖罪券.
贖罪券買了之後.
當然是給了錢教廷.
印發一張贖罪券給你.
那你已經離開世界的那些親人.
就會在煉獄減少他受苦的年期.

$^{881}$這樣的意思.
等於好像超渡一樣.
就是這樣的一個情況.
你可以想像到.
後來馬丁路德很實在.
他上到羅馬.
見到更大的教堂.
他想著會更加的敬虔親近神的主教.
發現他們很多的問題.
令到馬丁路德自己都有一個很大的失望.
所以這裡對於當時教皇的權威.
已經開始提出一個夠膽的質疑.
這個是95條裡面.
大家可以回去再看一看其中這一點.
當然整個議題.
95條裡面其實主要是想講一個point.
就是到底贖罪券是不是真的合乎聖經呢.
因為他們有一個package.
主教宣判哪些人你再生某些罪.
這些罪就可以透過買贖罪券減輕.
但有些死罪是不可以的.
所以一般人就不怕永生.
因為你在主教裡面.
你在教廷裡面就講明了你一定有永生.
你離開了教會就沒有永生.
但是有永生是一件事.
你在地上做了一些不好的事.
那些事耶穌的寶血就不會洗這些.
你要在煉獄受刑.
這樣就變成原來有兩個關卡.
你過了一關離開世界之後.
所有人就去到煉獄.
然後煉獄有些人很長有些人很短.
那些聖人因為他們當時都會拜聖人.
那些聖人就更加快去天堂.
還有就有敬拜瑪利亞.
就是聖母.
另外有一些很得道的.
他們不是叫高僧.
得道的一些聖人.

$^{921}$他們有些功德賦.
這些全部是一個package.
就是那些功德多到根本不需要在煉獄之餘.
還可以分給其他在煉獄裡面的人.
這樣的package就令到當時代的歐洲人很害怕.
還有個個爭著買.
出現了一個很大的後果.
就是買了贖罪券就肆無忌憚的犯罪了.
原來我可以買贖罪券.
等於我們70年代80年代.
媽媽的年代.
回大陸去百貨公司買那些代用券.
就可以上到廣州提取那個電視機.
就不用自己托上去.
類似這樣的代用券.
這些就是當時他們好像為著.
其實背後的動機不是想減少人對神的犯罪的意圖.
而是相反本身的動機就是想要百姓奉獻來建教堂.
這樣的問題就會變成整個教會的教義.
就在人的傳統裡面不斷地改變.
就導致到其實真是越來越人云亦云.
慢慢就變成有一個積重難返.
所以馬丁路德你可以說他列舉了95條.
其實只不過是幾個重要的綱領.
質疑教皇有沒有這樣的權威.
還有哪個人去定哪些罪可以赦.
哪些罪不可以赦.
哪個人去定呢.
另外為什麼會給一些錢.
就說在煉獄裡面就可以超脫呢.
誰去證實這件事呢.
誰去說有煉獄呢.
這些要推翻上去.
就是天主教本身的教義.
這個就是很大的一個問題.
不過要記得.
馬丁路德本人其實真是一個神父.
他為什麼信耶穌.
和他怎樣信耶穌呢.
他爸爸要他做律師.

$^{961}$所以他的思辨能力是很強的.
很厲害的.
所以他宗教改革之後.
他要逃亡歐洲.
他就在一個秘密地方.
做什麼呢.
翻譯德文聖經.
整本德文聖經翻譯出來.
可以給他的同胞德國人去看.
因為當時只有拉丁文的聖經.
所以後來因為16世紀.
印刷術在歐洲已經很流行的時候.
所以基於原來自己可以讀聖經.
就有很多神興起的人.
例如最早期的丁道爾.
再早一點的馬丁路德之前的.
已經開始質疑教皇的.
已經被燒死的約翰胡斯.
這些殉道者已經引發歐洲的啟蒙.
兩類的啟蒙.
個人自由的啟蒙.
和宗教信仰的啟蒙.
就在那個時候就來引發了.
所以馬丁路德絕對不是一個原創者.
是很多人曾經一點一滴.
反對教皇的誤誤.
開始質疑信仰是不是真的這樣呢.
這樣來引發宗教的改革.
真的直捲全個歐洲.
就是因為這個95條.
夠膽釘在威登堡教堂的門口.
來列舉.
其實這個列舉不是罪狀.
是一個論綱.
大家來討論一下.
其實是否適合聖經呢.
其實是可以有討論的空間的.
有沒有人可以給我一個更好的論述呢.
所以他沒有意義要攻擊教皇.
不過提出聖經可能不是這樣理解.

$^{1001}$當然在1517年10月31日.
一發不可收拾的就是.
這個用拉丁文寫的95條.
大量印發在整個歐洲.
當然是德國先行.
就來傳播.
一發不可收拾.
居然所謂龍裡雞作反.
你自己都是一個教士.
你現在質疑教廷所提出的這些贖罪券.
贖罪券背後就是教皇.
說他有這樣的權利去赦罪.
所以你買我的贖罪券.
你的親人就可以離開煉獄了.
很有名的有一個主教在羅馬.
有名到一個地步.
後世的人個個都嘲笑.
只要你的錢放進去奉獻箱.
你的親人就會立刻從煉獄離開了.
很出名的.
變成這麼噁心的犯賣贖罪券.
你可以想像到.
教會要建教堂.
只要你奉獻超過10萬元.
你的名字就響噹噹了.
類似這樣的意思.
你可以想像到.
都有人信和害怕.
剛才說的黑死病.
和煉獄這個概念.
其實不是迷惑.
是令一代人.
一直都對神敬而遠之.
再加上原來每一次崇拜的彌撒.
很神秘地只會讀拉丁文.
是不會解釋.
是不會有翻譯.
所以一般的人真的變成平民百姓.
每次真的去看彌撒.
根本不會明白主教說什麼.

$^{1041}$也不會用例子來解釋.
去讓真理使你自己更加能夠明白.
是沒有的.
所以你就可以想像到.
我多貼後面的幾頁.
因為我做了一些筆記.
今天未必能講完.
你想想馬丁路德的時候.
他是威登堡.
這個PowerPoint我會給大家的.
大家不用擔心.
整個PowerPoint給你.
我會陸續這樣做.
因為這些無聖經.
我這樣說很空泛.
所以我只能夠做PowerPoint.
給大家一個大小小的標題.
希望大家有個輪廓.
你可以看到.
馬丁路德當時在修道院.
學習拉丁文.
原文的希臘文.
他身軀比較強.
在大學裡面教授聖經.
你想想他1513年至15年.
他已經傳過150篇詩篇.
逐個星期來教授.
然後去到1515年.
他真的要講羅馬書.
又是verse by verse這樣講.
講了一年.
所以他自己有機會去背熟羅馬書.
戲劇性就是.
大家聽過了.
他自己感受到.
他說他自己是一個.
其實是人都知道.
連他的朋友.
他的同行都知道.
他是一個很好的人.

$^{1081}$但是他說他自己要解決內心的罪惡.
那種的罪疚感.
他說他自己聽從主教講.
你自己做一些克服.
他就跪上去威登堡教堂.
一路跪.
那些聖經就閃入他的內心.
突然間他背到羅馬書的一章17節.
大家應該熟悉的了.
我不以福音為始.
這福音本是神的大能.
要求一切相信的.
先是猶太人後是希利利人.
因為神的義.
正在這福音上顯明.
這義是本於信而至於信.
如經上所記.
義人必因信得生.
整個人突然間.
我從來沒有聽過這樣的經文.
突然間閃入他的心.
他自己真的連夜重看羅馬書.
終於他發現.
原來義人不是靠自己的義行去得救.
這個我們等一下會看.
他早期的神學觀.
就是契約論.
就是神和地球的人納了一個約.
人只需要滿足他自己的德行.
就可以得到神的恩典.
問題是天主教在這裡.
不斷講到人是做不到的.
必須要有那個救贖.
但那個救贖不是基督唯一救贖.
所以這個就是馬丁路德.
福音派和天主教最根本的分歧.
就是天主教說.
基督加上教皇的傳統.
即教會的傳統.
就是救贖.

$^{1121}$所以這裡這樣的家藏疊屋.
就可以有很多事情發生.
例如你要買贖罪券.
例如你要向教皇告解.
向當地的主教告解.
他要你罰你什麼 你要做什麼.
你都要做.
我又沒試過告解.
所以我不可以亂說.
如果有些人去過天主教告解.
我莊重的.
我們不是稀少老麻的.
因為天主教徒是很敬虔的.
如果有過這些經驗.
後來你又改了成為基督徒.
可以說一下你的經驗.
是怎樣的一種心情.
因為我們一信就是恩信稱義.
我們就沒有這件事.
我們是直接向神祈禱的.
我們就不需要找張牧師.
不需要找朱牧師.
請你代我向神認罪.
罰你跑五個圈的湖.
沒有的 我們福音派就是.
恩信就是稱義的.
所以馬丁路德就是過來人.
他從來沒有想過.
原來他自己可以直接和神.
有這樣的對話.
有這樣的稱義.
為義人只需要恩信.
基督耶穌就得生.
所以這個恩信.
就令他徹夜再重新看羅馬書.
真的發現整本羅馬書.
從來都是講恩信稱義.
不是講靠行為的.
連第九到十一章的猶太人.
為什麼不得救.

$^{1161}$都是他們不肯承認耶穌.
都是講信心的.
所以馬丁路德就真的重新一個發現.
到他真的一路去了解.
希望將他自己所得到的這個亮光.
希望找不同的主教.
他所敬重的去講.
都被人很自然的.
揭露教廷的根本的實學問題.
怎麼可能容許你呢.
所以都是禁止他去講的.
所以去到一個地步.
他只能夠是連夜將這個95條.
當然沒有這樣的紙.
沒有這麼好的印刷.
釘在威登堡教堂.
希望引發那個討論.
希望引發大家一個反思.
一個神學的討論.
所以他不是以下犯上.
去挑戰教皇的權威.
但是你做了這個行動.
你說教皇可能真的不是無誤.
教皇沒有權去赦免人的罪.
透過賣淑罪券.
還有那個信仰原來只是信心就可以了.
不需要加上.
這個最重要了.
就是牽動了整個教廷的權威.
就是教廷一直積集下來的傳統.
原來和舊恩是不需要拉上關係的.
你做這件事.
你是一個很敬虔的天主教徒是沒有問題.
但是你成為一個和信仰的根本真理拉上關係.
那就是一個很大的問題了.
就是說你是令基督的那個救贖.
無條件的救贖.
就會慢慢不知不覺被淡化了.
那些人就會追求很多的傳統.
所以大家知道天主教一定有七個聖禮.

$^{1201}$而基督教我們都是慈雲里的那些.
後來的嘉爾文的神學主張.
我們只有兩個聖禮.
洗禮和聖餐.
兩個聖禮.
這些浸禮班一定要認識的.
我們只有兩個聖禮.
我們認為基督徒一生一定要做的.
洗禮一次而已.
聖餐就是一生.
只有兩個聖禮.
我們不會拜嬰兒洗禮.
我們更加不需要神父最後的沒油的臨終的附油禮.
當然之後中間的有些婚禮.
或者其它的東西.
我們就看為是一個本身的禮儀.
未去到神聖的一個必須的.
因為有些人是獨生的.
是不需要看那個是一個絕對的聖禮.
這些就是天主教的傳統.
他們提升到一個必須性的問題.
所以令到聖經的恩訓清明.
越來越蒙上了塵埃.
人不能夠再憑自己和神直接的對話.
去得到那個心靈的赦免.
根本問題就是天主教是耶穌基督.
加上教廷的傳統才等於得救.
這樣就有一個很大的問題.
我們回到95條的開始的部分.
主要就是救贖罪券的意義.
和它帶出的效果其實是沒有效果.
是沒有效的.
馬丁路德見到更多人.
真的有恃無恐的拿著贖罪券去犯罪.
因為不需要害怕.
一來就是教廷說我有永生.
那我還怕什麼呢?.
那個贖罪券就是我不需要去煉肉.
那我又怕什麼呢?.
不怕之後就在地上犯罪了.

$^{1241}$馬丁路德就是親眼目擊到這樣的情況.
他就覺得原來那個好心做壞事.
因為你淡化了那個真理.
人不再依靠耶穌.
根本不會覺得犯罪是得罪神.
其實就是一個很嚴重的問題.
所以當然我們不可能逐句去區分.
大致上幾個範疇贖罪券的那個意義.
如果它沒有意義.
犯受的人或者有那個權威去印製贖罪券.
就變成別有用心了.
還有煉肉都是沒有人知道的.
聖經沒有說得這麼詳細的.
基督教就很撇脫的.
所有基督教都不相信有煉肉的.
這個應該都沒有錯的.
所以我們是直上天堂的.
信主的一次相信 一直相信 必然得救.
這個就是聖經的應許.
不是哪個教皇 不是哪個傳達人說.
更加不是哪個教會說.
是聖經很清楚的應許.
一次相信 一直相信.
就是到死都不否認耶穌.
依靠主耶穌就得救了.
所以臨終信主.
他那一刻真心相信.
我們看不到他的心 看不到他的行為.
但是我們相信他的承諾.
我們就都收貨了.
他肯跟我眨眼 我都幫他洗禮了.
我們不是神 我們不知道.
那個善意我們都要接受.
除非他真的怎麼都不肯點頭.
怎麼都不在捉客傳達人的手.
為他祈禱都不停的轉頭.
那我們就沒有辦法了.
就算我們很愛的弟兄姊妹.
你幫我爸爸洗禮吧.
不行 因為他真的沒有承認.

$^{1281}$我們都不能夠當他是 不行.
其實那個要求已經低到捉住一手.
眨一眨眼 就是這樣了.
或者點了點頭 但是他都不肯做.
那就不行了.
所以其實是一次相信 一直相信.
就必然的得救.
這些人就一定可以上天堂了.
還有基督教好的地方就是.
是沒有煉獄的.
不過是有基督將來的審判.
那個就涉及到希伯來斯的問題.
就不是煉獄 是得到賞賜.
還是受到責備的問題.
所以也都等於我們那個前千.
死了就是死了 等待復活.
一復活就去到千禧年了.
所以這個就是我們暫時我能夠理解到的.
那個信仰 那個信念.
所以放心 我不會賣贖罪券給你.
你也都不需要去煉獄.
你不需要擔心在煉獄裡面受苦.
好像說那個彈丁.
當年就是當時.
他是等於邪惡的代表.
彈丁都信有煉.
你想想可以想像中世紀.
教廷的權威和他們的教義.
對當時代的人的影響是很深遠的.
是使很多人都很怕死亡.
馬丁路德就是他怎樣信主呢?.
爸爸要他做律師受訓練.
接著有一次閃電 很大雨.
當時都是比較簡陋.
可能地方大 幾乎劈死他.
他就在當時向神.
他自己向神 因為找不到神父.
向神認罪 祈禱.
就說如果你不想我死.
我就奉獻自己做神父.

$^{1321}$接著就 應該不是好天的.
就是沒事了.
他就履行他的承諾.
爸爸就很生氣 履行他的承諾就進修道院.
接受神學的裝備.
後來就真是如願在威登堡那裡.
先做一個教授聖經.
和在教堂裡面來參與一些服事.
這個就是馬丁路德.
你可以說是他的前半生.
所以其實他得天獨厚.
就是他在教會裡面工作.
就有讀經的權利.
因為普通人是不可以家裡有聖經.
有都沒有用 因為你不會拉丁文的.
你拿著一本聖經都沒有意思的.
我給你一本希伯來文在家裡.
當然你可以供奉它.
你都不會那些字母.
是沒有意思的 看不明白的.
當時根本沒有大量的翻譯.
根本還沒有英文的聖經.
英文的第一本的音定版.
英王的音定版King James Version.
都是忘記了164幾年 我先回去查一下.
你可以查到 是很後期才翻譯的.
是因為馬丁路德他們既然打開了這個門.
就大量不同的人就翻譯了.
所以馬丁路德就第一個人翻譯德文聖經.
將拉丁文翻譯成德文的聖經.
就是說希望他的同胞都能夠自己讀到聖經.
那你就可以想像得到 整個歐洲.
其實都是需要有科技的配合的.
如果沒有印刷的 如果沒有人民主義的抬頭.
當時也有一個叫做文藝卿.
和人民主義的卿 就是每個人.
開始不再那麼懼怕教廷了.
因為發現教廷裡面原來不同地方又自己立教皇.
同一個時間有三個教皇.
那誰才是權威 誰才是真正的領袖呢?.

$^{1361}$就是說香港有一個特首 接著有一個民間特首.
接著他說自己又是特首 那怎麼辦呢?.
那我們跟誰呢?.
問題就是當時有幾個教皇.
所以你就引發人家會想的.
真的還那麼蠢 要來怕你.
那就要開始問 為什麼你們教廷自己都這麼文員呢?.
所以這些文字這樣開了.
可以導致當這個蓬打開.
人就開始可以去挑戰教廷.
是不是真的絕對沒有錯呢?.
所以神就興起了馬丁路德.
他自己是一個很敬虔 很懼怕神的人.
他已經有了這個自己的內心.
和很敏銳罪惡的一個人.
但是重點在這裡 他跟足了教會的傳統去做.
都解決不了自己內心裡面的罪疚.
就是說那個良心仍然時時都覺得自己離開神很遠.
就是說我已經用了教會要我做的辦法.
我都不行 你想我怎麼樣呢?.
他就開始要問有沒有其它的辦法呢?.
或者說其實原來這個辦法從來都不是.
聖經開了這個念頭就很重要了.
就引發整個宗教的改革.
他說原來有一幫人一而再的掩著聖經不告訴我們.
原來他不是想存心掩著聖經.
原來他是自己有一套的東西.
導致到聖經被人忽視了.
和聖經的真理不能被揭開.
原來就是這樣的意思了.
所以回到一個教士 他自己也是這樣身歷.
自己在那個成長裡面有過這些掙扎.
而那個掙扎不是根據教廷給的方法來做到.
就是說一定是有另外的辦法了.
另外的辦法就是恩信稱義了.
所以我們今天就太容易了.
但是對當時代的人 懼怕死亡.
懼怕神的當時代的歐洲人.
一定是要根據教廷所賦予的傳統來活的.
所以基督的白白舊恩就慢慢被淡化了.

$^{1401}$不再被人看為是那個唯獨了.
所以相信馬丁路德首先揭開的就是那個唯獨稱義這個問題.
去到加爾文就再將他發揚光大.
加爾文是後期的人 是瑞士人.
後期當然他應該是有法國血統.
法文和德文都是很精通的.
神都興起這些這麼厲害的神學家去影響整個歐洲.
加爾文就更加提到五個唯獨.
唯獨聖經 唯獨恩典 唯獨信心.
唯獨榮耀神 還有一個.
我經常都差一個說漏了.
唯獨恩典 唯獨信心 唯獨聖經 唯獨榮耀神.
還有一個我忘記了 我回去再找給大家聽.
我自己都親自教過.
我寫過那些早期的栽培部的主要學的.
我回去再找.
恩典信心 唯獨基督 唯獨聖經.
最重要的唯獨基督.
唯有基督是我們每個人的義.
所以其實一覽不止的.
你不信聖經根本說不出來.
聖經裡面所說的.
關乎到我們人的一定是那個信心.
那個信心和誰結合呢?.
一定是和基督.
所以又是那個唯獨就是恩典.
因為基督伯伯給我們的.
所以全部是一個套餐來的.
人的存在 你發現自己得救.
其實在命運當中.
其實我們是神的兒女.
所以我們整個人生.
由我們信主 我們離開這個世界.
唯獨榮耀神.
所以是五個唯獨.
但是馬丁路德首先揭開的就是.
唯獨清義的恩典.
所以就是這個95條論綱.
那個最重要的精神.
我講了這個95條.

$^{1441}$小小的花絮.
我們下一次再講那個神學.
這個原來我自己再看的時候.
是令我很感動的.
因為馬丁路德1517年.
輾轉他就要面對很多的教廷的審判.
被逐出教廷.
就是說你離開了教廷.
你沒有救恩了.
馬丁路德的內心就是.
我在主耶穌裡面就是救恩.
我不需要擔心這件事.
但是在那個感情上.
他真的是一個神父.
他真的是一個教士.
所以其實他一生都沒有想過要這樣脫離教廷.
不過是教廷要驅逐他.
他沒有辦法.
但是1517年之後引發了一些情況.
就是原來有很多修道院的人.
當然有男的女的.
有11個女的修女逃亡離開教廷.
因為會被處死的.
因為你離開是宗教的叛亂.
所以是偷偷地離開.
當然你離開之後.
你能夠嫁了人.
你又正式否認天主教.
那也沒有你扶.
因為當時也有其他一些君主.
是很反對天主教.
是很支持馬丁路德這個宗教改革.
其實背後是利用宗教改革.
發動一些政治上的利益.
所以後來馬丁路德既然引發.
你們都可以向天主教反對.
有些農民被貴族一路欺凌.
就有農民的戰爭.
種地的農民.
不是耳聾的農民.

$^{1481}$農民的戰爭就帶來整個歐洲很大的戰亂.
但反過來 馬丁路德又看這些平民百姓.
這樣的起義又不對.
他又支持貴族鎮壓.
有很多的農民是這樣死的.
所以後期路德他自己對這些社會的議題.
他自己都是舉棋不定.
因為他主要是在宗教上.
他沒有想到背後有這樣的政治問題.
因為宗教改革而衍生出來.
所以他都要面對要抓這條數.
有很多的農民被殘殺.
所以這個是農民的戰爭.
但是說回那件奇妙的事.
有一些修女跟隨宗教改革逃亡離開修道院.
如果你沒有人收留.
就一定會被抓 又會死的.
所以馬丁路德就盡力安排了那十個修女.
嫁給了好人家.
嫁剩下這個沒有人要.
這個叫做嘉達娜林娜.
Cathina 當然是荷蘭人.
Borla 博拉.
當然當時是這樣畫的.
沒有好相機.
不會高清 是一些人畫出來的.
其實是很感動的.
1525年我們都還沒有出生.
如果不是我們參加這個婚禮.
真的是世紀婚禮.
我們經常說的世紀婚禮.
就是一些皇室.
教士晚年去娶修女.
去還獨做他的太太.
他們的理念就是.
天主教一直所說的是不適合聖經.
聖經從來沒有禁止.
聖經就是說他們.
馬丁路德要為自己所說的真理.
來去身體力行.

$^{1521}$最終沒有人要他去跟這個太太.
就是跟這個修女結婚.
所以真的是世紀婚禮.
其實是很令人感動的.
因為路德就是為他自己的信仰.
來付出一個這樣的實踐.
似乎他的婚姻都不錯.
他有四五個兒子 都挺好的.
都令到他專心的去寫作.
還有一個賢內助去幫他處理很多日常的事情.
似乎馬丁路德都是多病的.
似乎他的神經需要是相當嚴重的.
他晚年的時候.
如果你去到他的故居.
你會看到有些墨水.
他特意留在那裡.
因為他不知道是做夢還是什麼.
因為那個壓力真的很大.
他說他一直在寫東西的時候.
突然間看到魔鬼.
控訴他.
他就拿起那個墨水瓶.
就扔過去.
拒絕他這個控訴.
這個是他自己說的.
所以到底是神經衰弱還是有異象呢?.
無論如何那個墨水瓶就扔在牆上.
就濺到整個地上都是.
我猜可能是神經衰弱.
因為太嚴重的攻擊.
他寫的一首聖詩是很出名的.
就是「上與士我堅固」.
是很繪著人的.
可想而知當時那個壓力.
他們要整個歐洲逃亡.
就是為了一個恩信清儀的信念.
一個信仰.
是基於聖經的一個信仰.
而來到這樣的捍衛自己的信念.
不是一種執著 不是一種固執.

$^{1561}$不是一種死要面.
而是如果這個是聖經的真理.
我就用我的人生付上那個代價.
來捍衛這個恩信清儀.
所以我們今天是在享受宗教改革.
至少帶給我們這個.
基本上沒有人不懂的.
就算你沒有去上過訓練班.
你都會跟人說信耶穌就行了.
但那個時候是沒有人敢夠膽說的.
因為是基督加上教廷的傳統.
有很多很多的傳統.
所以晚年的時候.
他能夠跟這個修女結婚.
都是一時的佳話.
或者說他都真的看到她的可憐.
都願意收留她.
他一直都想著獨身的.
不結婚的.
專心來事奉的.
所以今天我們在座或者在線上的.
一定是一百年多已經退休.
有沒有真的承諾自己為神.
某些事情你要堅持呢.
有你自己的參與呢.
是真心參與的.
因為沒有人能夠逼到.
大家現在這個時間.
你都不會說你不做就不行.
但是有沒有用一種愛神的心去參與呢.
例如我經常都說.
星期三其實好看的那些劇集.
越看越好看的.
但是你能不能夠為神.
逢星期三你半個小時.
公禱會的時間才祈禱.
八點半才開上線.
現在那個門檻只是按一下那個Zoom.
都不用你搭車回來教會.
那半個小時的公禱.

$^{1601}$都是一個對神的承諾.
其實都是一個參與.
其實似乎又要哄又要哀.
其實人家當時為了這個恩信清義.
真的要付上一個生命的代價.
這個就是馬丁路德給大家一點點的認識.
最後他就1546年離開世界.
這個是馬丁路德的人生.
他很勤力的.
他的逃亡是很安靜地.
不斷地寫作和翻譯聖經.
當然應該有其它的教學等等.
所以這個就是馬丁路德.
很快的給大家一點點印象.
我們下個星期就會正式講他的神學.
首先就是他早期的契約論的神學.
最重要的這個沃木斯會議.
就是正式斷定馬丁路德.
在天主眼裡面.
他是徹底的異端.
被驅逐離開教廷.
當然他就成為了宗教改革的最重要的推手.
我帶了兩本書上來.
是不是大家要看.
我要開鏡頭才能看到.
應該是我要開鏡頭才能看到.
現在我就停止分享.
我就給大家看看這些書.
這兩本就是上下冊的.
嘉爾文的神學.
就是基督教要義.
就是Calvin將他畢生的精力.
就寫了這兩本.
當然有其它的信.
其它的墓會.
嘉爾文也是翻譯聖經等等.
但是這個基督教要義.
就很有層次地.
將整個基督教的要理都寫出來.
就成為了你改革完.

$^{1641}$你都要有一些神學的課本.
才能讓跟隨你的人.
知道信什麼.
所以是需要不只是叫口號.
我們要改革.
而是需要有一些真理的.
聖經的解讀.
來做一個支撐.
我是讀神學的時候.
先看看 是2011年買的.
不是.
應該是.
我在讀的時候就看神學院的書.
後來都是方便墓會自己買一套.
我記得了.
因為當時按木.
教會就不可以收現金.
每個團契認頭.
這套書我很想買的.
500多塊.
有團契認頭.
我就可以買了.
因為這些是實文.
就可以歸給那個墓者.
就買了這套.
所以是2011年買的.
我記得了.
是中文來的 是翻譯來的.
英文是一本.
這套書當然是難看.
因為很枯燥的神學的東西.
嘉爾文是一個很有條理的人.
據聞他當時是瑞士一個很重要的領袖.
大家不會看的.
偶然間你揭開就無所謂.
因為你一打開就想睡著了.
很悶的.
但是對於我們研究神學的人就需要去了解.
那個時代是有很多的教義問答出現的.
去幫助很多當時還不是很識字的歐洲人.

$^{1681}$他們可以有啟蒙對真理的一個認識.
當然有大量的印刷聖經.
開始有不同的翻譯聖經.
德文 法文 後來又有英文.
就會再多不同的語言的翻譯.
所以就引發那個宗教改革.
就不只是叫了個口號.
而是有很多的真理的研讀來支持.
所以也開始有很多的解經書會出現.
所以馬丁路德自己也寫了很多解經書.
就是這樣.
今天我們就很簡單的開始了一個類似的introduction.
所以大家看到就是我是用馬丁路德做那個切入點.
之前的就不能夠再講太多了.
因為很多的那些資料.
但是馬丁路德這裡95條做一個起點.
就講一講那個來龍去脈.
和他自己早期的主張.
和後來的一些演變.
就給大家有一些的概念.
概念就可以了 聽一下就可以了.
盼望最重要的我想是最後那一堂.
我還不知道我要用多少堂.
但是最後那一堂我們一定要問.
我們今天享受宗教改革之後.
我們自己是不是要再改革呢?.
教會的恩順清義是不是已經荒廢了呢?.
或者恩順清義是不是慢慢就變成廉價呢?.
這些我們都要問的問題.
一起禱告.
感恩天父給我們能夠透過一個很有興趣.
趣味性的一個研究.
我們絕對不是做一些大的論文.
或者在神學裡面去看那些原典.
那些的原著.
我們只不過是就著很有限的材料.
我們對馬丁路德那一代的人.
他們提出的一些宗教改革.
我們去認識.
都是一個純粹是很簡單的淺談.

$^{1721}$盼望主都祝福弟兄姊妹.
能夠對於500年前的歷史.
我們有一些的興趣.
一起去學習.
也都更加重要就是.
怎樣引發我們對於信仰的堅持.
對於信仰的那份的單純.
那份的誓死的捍衛真理.
願主幫助我們.
我們仰望感恩禱告.
奉耶穌基督的名.
阿們.
\newpage



\section{}
\label{sec:Sb_mbdX9Sa4}
\textbf{2024-10-24 周四晚查經班《改革之路(2):馬丁路德的十架神學》 講員:曾裔貴牧師}
\newline
\newline
連結: \href{https://youtube.com/watch?v=Sb-mbdX9Sa4}{\texttt{https://youtube.com/watch?v=Sb-mbdX9Sa4}} ~~~~ 語音日期: 2024-10-25
\newline
\newline
\hyperref[sec:aTaCCC_F45U]{\small{< < < PREV SERMON < < <}}
~
\hyperref[sec:index_chronic]{\small{[返順時目]}}
~
\hyperref[sec:index_scriptual]{\small{[返順卷目]}}
~
\hyperref[sec:ziJGxjL2HKY]{\small{> > > NEXT SERMON > > >}}
\newline
\newline
$^{1}$好 弟兄姊妹 我們有一次再有機會去研讀.
馬丁路德的一些關於神學.
恩順清二 我們可以開始一起祈禱.
我們在晚上又能夠有一段安靜的時間.
我們去仔細的去思想宗教改革這個課題.
這個不是一個神學研究這麼簡單.
是帶來整個歐洲.
馬丁路德成為那個起始的人.
帶來差不多是很多的天主教徒.
決心的跟隨馬丁路德一幫的改革者.
離開天主教的教皇和教區的轄制.
開始他們基督教的聚會.
當然一直蔓延到德國不同的地方.
去到甚至乎英國 去到北歐等等.
一直的宗派的發展.
後來5月花號就去到美國.
然後也都去到20世紀.
這個普世的宣教運動.
去到中國,印度,非洲.
都是由宗教改革這個開始的.
所以對於恩順清二.
給我們能夠有一個機會.
這樣來到大家靜心的去思考.
去看下當時他們的處境.
也都盼望我們能夠去到展望未來的時候.
我們有沒有一些的缺失.
是不是需要再改革.
還是我們有些東西是我們需要去保存.
前人給我們的寶貴的那個屬靈的寶藏.
叫我們能夠珍而重之.
還有一個真的要去傳揚福音的心智.
還有那個行動.
今天晚上我們第二講的.
馬丁路德的這個恩順清二.
特別是關乎到信心的三個層面.
都給我們能夠掌握得清楚,明白.
多謝主給我們能夠有一幫弟兄姊妹.
很願意的來到去.
好像奉獻自己的時間.
去聽一些可能要神學院才會教的課程.

$^{41}$幫助我們能夠很趣味性去掌握,去了解.
我們這樣禱告.
奉耶穌基督的名 阿們.
有個簡單的請求.
因為上個禮拜看到在那個錄音的留言裡面.
有弟兄姊妹就說.
可能那個現場的雜音.
可能就是線上大家沒有關閉到那個講話的那個.
所以大家就靜音.
除非你真的需要發言.
那你可以舉手.
那你就開了 如果不是呢.
靜音就會方便一些.
我們在講的時候.
就不會有一些雜音騷擾其他的弟兄姊妹.
或者再聽的時候.
他們聽到就會比較明顯的.
所以大家就幫幫忙來到去處理.
好 那我們開始第二講.
上次很簡單講了.
馬丁路德的那個背景.
生平.
甚至他在晚年的時候.
15,16年就和他的太太結婚.
當然他的太太是一個修女.
所以是幾世紀婚姻的.
就不是查理斯和戴安娜這些皇室的婚姻.
是一個神父.
當時他就叫做一個.
一個類似是祭司.
司鐸這一類的天主教的名稱.
他不是一個主教.
馬丁路德還沒有升到他做主教.
但其實在他的轄區.
威登堡他都要管理很多教堂的一些行政事務.
因為他都是一個相當出色.
有神學博士學位的.
很能幹的一位神學家.
但同時間他自己真的是一個修士.
你想一下一個修士.

$^{81}$在晚年的時候.
接近都年紀大了.
和一個修女結婚.
體現的不只是可以結婚這麼簡單.
是體現你打開了那個救恩的門.
叫人不需要懼怕.
離開教廷離開教皇的轄區.
那個轄館你都可以有救恩.
這個是天方夜譚的對於中世紀的人來說.
因為教區就是他們的命脈.
他們的生存.
你想一下他們要去買贖罪券等等.
所以和我們今天已經騙著是教會.
教堂其實很難這樣做一個類比.
當時那些修女真的是無家可歸.
還有十個都嫁得出.
唯一一個嫁不出.
馬的路德就收留她成為夫婦.
到底他們的情況是怎樣呢?.
將來如果我有機會去看這一類的書.
我不知道有沒有他的自傳.
和大家再分享一下這樣的婚姻.
其實是很有趣的.
正正式式他自己作為一個修士.
他自己以身作則.
來到這樣的年紀.
一方面娶一個修女.
另一方面有婚姻生活.
婚姻的生活都是相當重要的.
如果一個好的太太.
可能是牧者.
特別是馬丁路德.
那個事奉有一個非常大的後盾.
一個在家庭裡面的模樣.
會對他有更大的安全感.
因為馬丁路德由1517年之後.
其實應該嚴格說是1521年的.
最後的大會決裂的時候.
他就不斷被教廷追殺.
所以當時一些皇室的人.

$^{121}$就將馬丁路德收藏.
所以你可以想像到.
馬丁路德要面對的代價都是很大的.
這個就簡單.
我們上一個星期講了.
所以今天如果大家不夠精神.
你看不到我 看不到你.
你可以睡覺沒問題.
因為我做的時間.
真的很神學的.
真的很關乎他怎樣理解恩信的青意.
我自己一直準備的時間.
那個時代其實每個人都很怕死.
尤其是我們上個星期講過.
黑死病整個歐洲死了差不多四分一人.
你可以想像到那個嚴重性.
對於死亡這麼近.
大又死小又死.
你想像一個這麼先進的歐洲文化.
當時已經是文藝復興等等的一個世代.
其實很崇尚對自由.
對那些原著的追尋.
很講那個哲學等等.
一個黑死病從亞洲.
其實應該是從蒙古.
一直傳過來的時候.
是沒有救的絕症.
很多人死.
所以很多時候瀰漫著死亡的詩歌 陰影等等.
其實是令到歐洲人很害怕.
所以馬丁路德其實過了一百年.
但是對於死亡這件事.
歐洲人就體會得很深.
所以你可以想像到.
就算是一個平民百姓.
或者是一個飽學之士.
一樣是面對死亡.
他們會覺得很近.
固然馬丁路德他自己的見證.
就是他遇到一次雷電交加的一個時間.

$^{161}$他幾乎被雷劈死.
他就懇切向神祈禱.
如果你不劈死我.
我就要奉獻自己做一個修士.
果然真的沒事.
他就履行這個承諾.
本來他爸爸要求他做一個法律的律師.
他真的修養過 讀過這些法律.
所以對於他那個思辨的能力.
其實都很強的.
輾轉就走神學的路.
進入神學院進修.
我開一些powerpoint給大家看看.
因為這樣說.
《武聖經》必須要用powerpoint.
才能夠大家掌握到馬丁路德的思路.
馬丁路德的關於贖罪卷和95條.
這是1517年10月31日.
大家一定要記住萬聖節的日子.
1517年 其實已經超過了500年.
在2017年的時候.
500年香港的神學界.
對於宗教改革有很多研討會.
2017年的時候已經發生了.
如是者已經過了7年.
已經有7年的時間過了.
所以我們都需要.
這個是舊的powerpoint還是新的powerpoint.
我看看 我那個powerpoint.
這個是舊的 上個禮拜的.
對不起 我應該有一個新的powerpoint.
這個 今天做的.
所以我奇怪 為什麼是新的.
我再從頭再說.
我之後整合之後.
我會最後給大家這個powerpoint.
大家稍為有一點印象.
反正我們都不會有機會.
在其他主要學會教.
大家可以當是一個趣味性來聽一聽.

$^{201}$有一些很重要的日子.
大家要留意.
1512年10月21日.
就是他神學博士學位的畢業.
在威登堡大學裡面.
但要記住.
他其實同一時間.
他已經是一個修士.
他已經在教廷裡面.
來做一些牧養.
主要是教聖經.
所以由1512年一直開始.
去到1517年.
其中是1515年.
很記住.
1515年就是他自己一直去教聖經.
早期教詩篇.
後來1515年.
你猜他教什麼書呢.
就是教羅馬書.
我好像做過一個powerpoint.
做了一頁.
講到他的時間.
我先跳一跳.
我等下再給大家看相關的經文.
這些是他自己的自述.
這就是這一頁.
我先看這裡.
他教的聖經.
是導致他直接自己去.
被聖經來挑戰.
他的內心.
所以就令到他.
有很大的發現.
就是他發現.
這個恩順清二的神的義.
到底.
1513至1515年.
教詩篇.
然後用了兩年時間教.

$^{241}$應該很仔細地教.
用了超過一年的時間.
1515至1516年.
所以他在1514年的年底.
就有重大的心靈自己的發現.
就是對於神的義.
他有新的理解.
一個insight.
直接在他自己的內心.
所以你會看到.
1517年.
不是突然自己寫95條出來.
控訴教廷.
而是他自己先有一個很大的復興.
到底什麼是信.
他有三個層次.
可能大家其實已經擁有了.
不過我們等一下再仔細去界定.
哪三個層次.
缺一不可.
因為是有互為關聯的.
所以你看到這個時間軸.
你就會看到.
1516,1517就是迦拉太書.
羅馬書的姐妹書信就是迦拉太書.
都是講恩順清二.
迦拉太書主要就是說.
靠律法是不可以清二的.
包括猶太人都不行.
所以甚至乎.
因為迦拉太書第三章就是說.
當時的弟兄姊妹.
準備上受割禮.
來覺得自己.
感覺上不會這麼噁心.
給很多人來挑戰.
做多無壞.
但是阿保羅就說.
如果你這樣去做.
而你失去了對基督的唯一的.

$^{281}$簡單的赤裸裸的信靠.
那就是同福音為敵.
如果全割禮的人.
就更加.
阿保羅第五章的迦拉太書.
恨不得這樣的人.
如果全割禮.
就將自己割絕.
割禮就是生殖器官的包皮.
男性.
保羅就罵這些全割禮的猶太人.
五章應該是第一節還是第二節.
你可以想像得到.
保羅看恩順清二.
是看得很嚴肅.
所以馬丁路德.
所以他成為一個修士.
你可以想像得到.
普通人是沒有機會.
不可以自己擁有聖經.
不可以自己家裡私藏聖經.
讀聖經.
更加不可以解經.
因為根本不懂.
因為當時馬丁路德之前.
其實只有很少量的翻譯本的聖經.
中文就更加不用說了.
去到二十世紀才有.
還要等多五六百年才有.
你可以想像得到.
當時的德國人.
是沒有自己的德文聖經.
所以我們要看當時的背景去理解.
聖經只有拉丁文.
當然有希臘文.
當然有希伯來文.
因為本身就是希臘文和希伯來文.
但是在整個中世紀.
以拉丁文作為官方講解聖經的語言.
所以一定要一些很厲害的教廷裡的專家.

$^{321}$主教 修士.
一定是每個人都很厲害.
讀書的人才可以做得到.
所以你想想他們沒有意思.
不過客觀地壟斷了解釋聖經的權利.
就變成了平信徒沒有機會.
沒有能力自己去研經.
所以馬丁路德成為修士.
既是讀神學.
既是他自己是一個法律的律師訓練.
所以他自己更加有身份去研讀和教導聖經.
所以在這樣的過程裡.
由1513至1516年.
我相信是1514年年底.
或者1515年那段時間.
就是他自己最大的心靈啟發.
所以這些就是聖經帶來的直接的confrontation.
就是衝擊我們人的良心.
面對聖經等於面對神.
所以為什麼中國改革成為了這麼重要的.
似乎是二十世紀近尾.
美國開始主日學的運動.
有一個很好的牧者.
不是牧者 是一個在美國的教會的姐妹.
見到很多小朋友.
星期天到處在街上玩.
不回教會.
她覺得很大的觸動.
她就邀請他們回來教他們上主日學.
就開始一個叫做主日學運動.
就遍及全世界的教會.
差不多基督教是獨有的.
星期天個個無論大人無論小朋友都要上學.
就是主日學.
成為了聖經普及教育的開始.
或者是那個基礎.
這個你可以想像得到.
肯定是從宗教改革對聖經的重要看待而出現的.
所以你會看到.
首先是馬丁路德他自己.

$^{361}$就有這個得天獨厚的機會去研讀聖經.
使他能夠在講解的過程.
他自己得到生命的啟迪.
從而讓人看到原來中世紀教廷.
那些守傳統的月來月.
嫁床疊屋.
使整個信仰本來很簡單的恩信清義.
就變得遙遠複雜.
很多的手續層層關卡.
所以大家都不會覺得不以為然的買贖罪券是很理所當然.
教會沒錢就不如賣贖罪券.
一舉兩得一雞兩美.
既有錢去收集教堂.
又有可以幫到平信徒.
好像預先解決了很多他們懼怕的問題.
就是死去的親人在煉獄.
怎樣可以解除沒那麼多的苦難呢.
就賣贖罪券.
因為那些聖人死了的聖人.
是有過多的功德用不完.
就可以賣給那些在地上的親人.
那些親人買了之後.
就令那些在煉獄裡死去的親人減少那些痛苦.
所以就有這些古怪的想法.
疊加在了救恩的恩信清義裡面.
其實就是源於馬丁路德開始自己去嚴經.
他起初是對羅馬書的神的義.
神的義是在第一章至第四章.
羅馬書是出現很多次的.
是不斷在說神的義.
所以早期他說自己是很憎恨讀羅馬書的.
為什麼保羅會這樣說呢.
就說這是這麼寶貴的恩典.
神的義他起初的定義.
我們回頭看.
因為我做了一個PowerPoint.
就有很多引述了他自己的.
當然這是翻譯本.
我們不是看那些原著.
所以如果大家要讀神學.

$^{401}$嚴格來說是要看原著的.
但是我們現在是主一學.
我們就不需要看原著.
我們看一些翻譯.
有人將他自傳的東西.
他在他的《釋經書》裡面說的東西.
做了一個緝錄.
所以我只是引述這些情況.
我們回頭看.
從頭開始.
這裡很清楚.
這個肯定是對歐洲人一個很大的衝擊.
死亡的陰影濃鑿.
大家懼怕死亡.
那就對了.
教廷只有一個.
只有一個天主教.
當然他們不是叫天主教.
叫大公教會.
就是普世的.
所以這個就叫一聖而公.
但是小道信經說的.
是說所有屬主的一聖而公的教會.
天主教就說是.
彼得就是我們教皇.
一直傳下來.
所以天主教就是唯一的教會.
是不會容許有異端的存在.
你有稍微和教皇不同的意見.
你就是異端的.
你是異端的話.
當時是有異端裁判所的.
異端裁判所是抓你去審.
審完就要燒你.
如果你不肯改變.
你不肯跟隨教廷的那個說法.
你就要受到那個懲治.
聽過一個叫約翰胡斯沒有?.
比馬丁路還要早.
都幾十年了.

$^{441}$他就是真的來到去.
新鮮事卒被燒死.
因為他拒絕教皇的絕對的無誤.
所以已經有前人成為了宗教改革的火種.
所以馬丁路德不是一個人.
他自己突然間的醒悟.
有很多前人開始意識到.
教廷已經湧腫到一個地步.
出現了很多很多的問題.
所以神會興起一些人.
當然一定是一些.
可能是他們本身有機會接觸教廷.
接觸聖經的人.
所以在馬丁路德之前.
已經開始有翻譯聖經這回事.
也是因為印刷術傳到歐洲.
所以科技原來不是洪水猛獸.
科技是帶來宗教改革的.
好像是撥黃的火的那把扇.
所以當時的科技是帶來一個大復興的原因.
就是每個人都有聖經.
當然一定是需要有這些神學家.
有這些精通希臘文.
希伯來文的傳道人.
馬丁路德.
還有其他的.
那個叫做丁道爾.
現在有一個聖經翻譯協會就是丁道爾.
這個人本身就是那個時代.
翻譯聖經的一位神學家.
你可以想像到.
比如嘉爾文也是翻譯聖經.
所以在馬丁路德這個改革之後.
第一件事就是他們知道.
如果我們講完.
大家仍然是憑感覺.
信仰的改革很快就會消亡.
要怎樣能夠令到這個改革有根基呢.
就是人人手上有聖經.
自己有聖經.

$^{481}$因為你離開教廷是沒有用的.
你自己家裡都不讀經.
人人對你有什麼要求呢.
是帶不了大規模的復興的.
那個改革的.
所以他們要做的.
就不再是對抗而對抗.
而是他們回到聖經的解釋.
所以講解聖經成為宗教改革之後.
雨後春筍.
不斷有這些研習班.
某程度上就是由馬丁路德開始的.
加爾文就更加在瑞士.
整個瑞士都受到加爾文的神學影響.
上星期我給大家看的基督教要義.
如果大家真的感興趣.
可以去書室買的.
也不是便宜的.
但是可以了解一下中世紀.
他們是怎樣很仔細的.
將一些信仰問答.
很具體很細節的寫出來.
一條一條的.
好像法律書一樣.
其實分得很清楚.
什麼叫信.
信徒的人生應該怎樣.
很詳細的.
其實某程度上都很好看的.
他分了類等等.
所以對於死亡.
就是中世紀.
其實不只是馬丁路德.
就是對於整個世代的人.
他們的那個懼怕.
懼怕死原來是一個重要的神學發現.
今天不是太多人怕死.
我們不怕.
但是我們知道它存在.
我們可能用不同的方法去麻醉自己.

$^{521}$永遠不會發生在自己身上.
但是對於馬丁路德來說.
無論是客觀地歐洲人有黑死病帶來的那個懼怕.
到馬丁路德自己.
其實他一直對自己面對一個難題.
就是他沒有辦法接受神的意義.
這個就是他的突破.
他不接受到.
因為他自己做不到.
我們看看他自己個人的自述.
他怎樣對自己有一個認識呢.
所以.
可以的 拿走它.
好 再接著.
他自己的自述.
我是一個好的修士.
嚴守我的會規.
應該就是奧古斯丁的修會.
以至如果有一個修士.
可以因著修道操練而進入天國的話.
我就是那個修士.
你可以想像到.
他很坦誠地面對自己.
按照著應該做的教規.
例如五點鐘要起身.
有靈修 有神操.
天主教叫神操.
不是馬的神操.
是早晨起來去靈性操練.
如果你去天主教的神學院去做退修.
你就會起身.
你會看到那些修女.
我們在香港那邊.
應該是博富林那邊.
那裡有一間修院.
我們曾經去過那裡退修.
真的 五點多還是六點.
那些修女就齊了.
在那裡大家獨經祈禱.
他們的生活就是這樣.

$^{561}$所以這是其中一個.
我在修道院所有的同伴.
都會肯定地告訴我們.
這件事如果要靠行為.
我就是那個可以上天堂的修士.
不過他要面對的是他的良心.
但我的良心卻不會這樣肯定.
我時常懷疑.
並且這樣說.
你所作的並不完善.
你的懊悔都不足夠.
你在認罪中還有遺漏.
我越想試圖藉著人的傳統.
所以這就是天主教.
去到馬丁路德的時代.
那個傳統疊加上去的各式各樣的東西.
已經是多到一個地步.
做的人都不知道為了什麼.
很多手續.
所以你可以想像到.
那個慢慢慢慢就變成一個枷鎖.
失去了真正的自由.
所以那個時代.
很多不信的稱為人民主義的人.
開始醒覺.
因為教廷是唯一的.
教廷說的話就是絕對的.
教廷要求每一個信眾要做的.
其實越來越複雜.
所以崇尚個人自由.
其實是越來越強烈的.
越來越同一時間湧現.
所以我們今天只是從宗教這個層面.
去看宗教改革.
但其實那個問題.
在民間帶來一場叫做三十年的戰爭.
是那些農民起義.
但馬丁路德因為他一直都在教廷的氛圍裡成長.
他就覺得反對教廷.
你這樣暴亂是錯的.

$^{601}$他是贊成鎮壓農民的.
所以帶來很大的人對馬丁路德都有批評.
所以我們看到其實在那個時代.
正正因為教廷以為自己是拿著天國的弱勢.
這個要很小心.
要不就是神.
要不就是邪教.
講完別人不聽你講.
其實那個責任反而是你講完就行了.
就不用你負那個責任.
你講完個個都聽你講.
你真的要很小心.
可能你真的變成邪教的教主.
所以當然我們期望我們講完.
我們自己有實踐.
別人不實踐不進行.
他自己要面對聖經對他的要求.
但我們都沒有權去強迫人.
這個世紀其實就是強迫人必須守教廷的教規.
他們都不是講真理的教規.
是傳統.
我認為這樣就行了.
我認為我們星期去伯特利.
不去伯特利就一定要被定罪.
慢慢成為一個傳統.
好像變成一個約定俗成.
不是這樣就不行了.
我們那個傳統本來是很好的.
大家一起有聯合崇拜是很好的.
但慢慢慢慢大家都不理背後的爭議.
只覺得那個做法就已經當是真理.
很容易就有很多這樣的情況.
所以我們其實要很小心.
相信當時一定是很多很多.
馬丁路德是做完這些傳統.
做到他自己和其他人看他都是無可指責.
但那些是外面表面的功夫.
不是講內心的態度.
他說如果是這樣的時候.
我是足夠去滿足那個要求.

$^{641}$但是軟弱與煩惱的良心.
我就每天越發發現他的不確定.
和我內心的軟弱煩惱.
看到自己不能夠滿足拯救的先決條件.
這個就是先後的次序.
所以馬丁路德也有人對他很嚴厲的批評.
叫人不需要靠行為得救.
這就說令到整個世界.
當時的世界變成無政府狀態.
說不需要行為.
馬丁路德在這一方面被人攻擊.
尤其是那些人民主義者對他的攻擊是很大的.
但是如果我們真正理解路德講的靠行為.
不靠行為的背後分別.
其實很重要的是我們要了解三個層次的信心.
就是馬丁路德他整個人的發言.
剛才我一開始已經講了.
神的異在早期路德沒有那個insight之前.
他是很懼怕.
因為神的異的先決條件就是.
你要做到才可以給你.
你做不到就不給你了.
就是說如果我們這個自由.
得救在神的面前的一個救恩是絕對的.
但是你要得到這個自由.
就要一千萬 假設.
但是你和我都沒有這一千萬.
你有一千萬你就可以得到這個救恩.
問題是每個人都沒有.
問題就是馬丁路德很憎恨神的異.
因為你定一個標準但是沒有人達到.
那怎麼辦呢.
所以他一直卡在神的異的情況下.
他一直都很想了解.
和很討厭為何保羅講到神的異.
在這個救恩上顯明.
他一直都理解不到.
所以理解不到他就有責任要教聖經.
那就神工作了.
就讓馬丁路德終於接觸聖經.

$^{681}$終於思前想後不斷去研讀的時候.
背熟很多聖經.
聖靈就工作了.
讓他看到那一千萬其實就是信心.
其實就是每個人都可以得到的信心.
不是靠信心去換取.
信心只不過是我們每個人都可以具備的.
那個真正的先決條件.
其實是神已經為人完全的預備好.
所以不是我們要靠那個信心去換取那個救恩.
所以等一下我們就會講信心的三個層次.
馬丁路德他自己的真正的發言.
他就帶來這三個層次.
和教廷所講的傳統和要求.
簡直是天淵之別.
原來可以這麼簡單.
和每個人都可以清楚自己和神的關係.
而立即得到這份平安福音和自由.
心靈的自由.
所以其實和我們某程度上今天的初信班.
我們一直這樣講.
都是前人500年前的人已經思考過.
發現真是唯有.
所以這個唯獨是很重要的.
所以嘉爾文就將這個唯獨的概念.
背後的精神.
將它再延伸.
其實有五大方面.
其實都是馬丁路德的東西.
不過再有系統地分成五個唯獨.
唯獨信心.
唯獨基督.
唯獨聖經.
唯獨恩典.
唯獨榮耀神.
五個唯獨.
就是我們人生裡面.
差不多是我們人生信仰的根基.
有這五樣東西.
我們就一定可以上到天堂.

$^{721}$在人生必定會有這份心靈的平安.
所以馬丁路德在這個時候.
在1514年之前.
他還是一直發現自己沒有辦法.
來達到神的義的標準要求.
神的義已經定了.
沒有辦法改.
但是怎樣可以達到呢.
他一直想怎樣才能達至呢.
但是達不到.
傳統 教廷的傳統.
他已經達到了.
但是還沒有真正有心靈的平安.
到這裡大家有沒有什麼問題.
大家都應該明白和不會太快.
他問的問題就是.
既然這位神是這麼有恩慈的.
我怎樣才能找到呢.
但是在1514年之前.
他認為神是很嚴厲的.
對人很要求高的.
我已經放在這裡了.
神的義是不變的.
你得到就可以得著平安.
不過你要達到.
你要找到達至.
去得到神的義的方法.
那個交換.
所以馬丁路德就說.
當然他那時候說的德國的貨幣.
多少錢.
但是從來沒有人可以有這樣的錢去換這個救恩.
一直都是說.
覺得神的義放在這裡.
成為定罪世人的一個方法.
所以他覺得神是一個殘忍的神.
很嚴厲的神.
怎樣可以找到一位恩慈的神.
令到我得著平安和自由呢.
所以轉捩點.

$^{761}$他的生命的改變.
就源於一方面對死的懼怕.
一方面想有心靈的平安.
他已經開始根據教廷給的指引.
去活每一日的生活.
但是最終良心面對神.
知道有神的存在.
面對神.
就是他內心的那個交戰.
這個交戰怎樣能夠得到平復.
得到那個突然的豁然的開朗呢.
就進入信心的問題.
其實他1512年已經.
剛才我都說了.
他已經開始有那個教席.
有那個修士的身份.
他已經可以駕馭希臘文.
希伯來文.
和可以用拉丁文.
或者用德文去講學.
教那些平民百姓.
來教聖經.
所以他自己有聖經的研讀.
是帶來一個最大的更新改變.
所以很擔心的.
今天我們喜歡談論聖經的教會.
研讀聖經的教會是少之又少.
雖然我沒有一個統計.
但是我們今天可能崇尚音樂.
多過崇尚聖經.
我們來聽聖詩.
我們就覺得很舒服.
一聽聖經.
有時很舒服.
不過睡覺.
睡覺本來不是錯的.
但是習慣性的睡覺.
回家又不研讀.
那是問題.
你睡覺沒問題.

$^{801}$你回家很認真研讀.
或者你平日很認真去走聖經.
是沒問題的.
但問題是.
你平日又是不讀聖經.
來到禮拜堂.
最需要講聖經的你又不讀.
甚至說長一點.
有沒有搞錯?說這麼長的.
妨礙我吃飯.
我們從來沒有聽過人說.
去演唱會.
喂!不准唱這麼長.
又延期.
不准延期.
我要趕著走.
給多貴的錢去聽演唱會.
我們覺得意猶未盡.
延期再延期.
打了蛇燈.
那個情緒.
那個情緒去到頂點.
但是很可惜.
我們不能夠這一代人.
來到禮拜堂.
應該是義正嚴詞.
應該是最基本.
禮拜堂就是禮拜神.
敬拜神.
去聆聽上主的教誨.
教導.
應該就是.
既然現在已經有這樣的傳統.
有這樣的設定.
但是那個設定好像已經失效.
大家已經不是這麼看重聖經.
如果聽到神學生講道.
就更加.
更加不是很想聽.
或者一些你不喜歡聽的傳道人.

$^{841}$你更加不想聽.
你責怪他講得不好.
但是從來自己不需要付任何代價.
我好像說遠了.
其實這個是路德他自己對聖經的發言.
他自己有機會去研讀聖經.
是讓聖經打開他的心.
他自己面對聖經.
面對神的時候.
他真的有一個很大的發言.
而這個發言是令他要付很多的代價.
他也是在所不惜.
因為真是一個insight.
他要和其他人分享.
所以他由始至終.
是沒有想過.
他根本不想和教廷決裂.
他是一個很認真的修士.
他只是想帶來一個更新.
不是一個革命.
不是一個改革.
所以教廷是他的同路人.
不是說不上就大家決裂.
就不是的.
這個從來都不是馬丁路德的.
其他人呢 比如慈雲里.
就比較激進一點.
可能是了.
因為人家沒有這個包袱.
比如嘉爾文.
他都不是修士.
沒有這個包袱.
所以他們不需要受到教廷的養育.
任避成長的關係.
來左右.
但馬丁路德就真是一個修士.
他很敬重他的主教等等.
所以每一次.
其實他都是很婉轉地來講.
這個是聖經.

$^{881}$我是沒有辦法收回.
所以如果你要這樣來做的時候.
我只能夠說.
這是他很重要的宣言.
求主幫助我.
他的清二觀早期就是這個.
你要去justify你自己.
進入和神的正確關係.
你要在神眼中成為二人.
所以神的二.
就是一個標準.
一個每個人必須要達到的.
但問題就是.
我們沒有辦法達到神的標準.
同一時間.
個人主義的意識已經抬頭.
不是一個集體性.
很多人.
當時他們那個時代.
有一樣東西是很特別的.
就是開始追尋原典.
對於古代一些很出名的作家.
或者其他重要的著作.
他們會去研讀.
所以那個時候.
就有這些對於原典的追尋.
所以馬丁路德.
其實不是一個徹底的革命家.
他突然間的發現.
其實是將前人已經建造好的.
他自己再去更新.
所以馬丁路德.
其實是受奧古斯丁的神學.
一個很大的啟發和影響.
所以這些是我們要留意.
所以在這裡的時候.
他早期的教義.
其實都很清楚.
人如果不根據神的意義去生活.
其實就不能與神建立正確的關係.

$^{921}$和好的關係.
和去到神的天家.
他很清楚.
今天我們初訓班探索班的人.
第一堂就要問這個問題.
你與神的關係是怎樣呢.
探索班的人就說.
神呀 誰是神呀.
我們就說我們這一位做物主.
他如果經過了 接受了.
就進入第二個問題.
就是我怎樣可以與神的關係呢.
馬丁路德不斷在問這些問題.
等於幫我們建構一個很清晰的初訓班的信徒.
他要經過的心靈的歷程.
繼續下去.
我先喝一口水.
這個就是神的意義.
另一種解釋.
沒有人可以逃避這個契約.
就是說神已經定了他的義.
等於與人已經定了.
我既然創造萬物.
我就定了一個標準.
這個標準就在馬丁路德來說是一個契約.
如果你做不到.
你就要面對神對你的審判.
所以這個就叫做清義.
但問題就是靠什麼清義呢.
這個就是宗教改革的精髓.
我們今天當然知道.
因信稱義.
已經成為很濫用.
但今天我們就有時間.
我們就去思想這個信的三個層面.
可能大家已經有了.
不過彼此再去溫習.
最低的先決條件.
就是盡你一切所能.
知道的 未知道的.

$^{961}$你去滿足.
但問題是馬丁路德已經說了.
我已經盡我自己去做.
人家都說我是一個好的修士.
如果靠這樣就可以上天堂.
我就是那個可以說出來.
就是我可以上天堂的了.
但馬丁路德說.
我自己看回我自己的內心.
我的良心.
我的軟弱.
我的懊悔.
我的苦惱.
根本未解決.
就是說他裡面.
根本還有一個對於罪的那種敏銳的纏繞.
所以這個契約論就變成了.
在路德的神學裡面早期的問題.
這裡就開始進入這個很重要的信心的突破.
就是原來清二是在說那個信心的特質.
三個層面的信心.
所以他講的恩信清二.
是憑著恩典恩信清二.
重點就是剛才那個例子.
馬丁路德自己說.
清二是絕對的標準.
誰有一千萬.
你就可以去換取那個清二了.
但是馬丁路德說.
沒有一個人可以有這個.
你可以說交換的條件.
他就停在那裡.
僻在那裡.
但是原來.
當後來他看羅馬書.
他就發現原來根本神的義.
神自己已經預備了一千萬.
每個人都可以拿到的.
這個就是什麼呢.
就是恩典.

$^{1001}$我們今天五百年後的人就很簡單.
我們不斷講恩典恩典.
講得很簡單.
但是對當時的人換取.
要守住很多的傳統.
用行為換取救恩.
今天其實等我一會再講.
馬丁路德的發現就是.
行為是恩典的產物.
是恩典因為信.
一個對神的恩典的感恩的結果.
就是好行為.
我們今天聽就很簡單.
但是當時他們就混亂了.
必須要做好才能換取神的義.
但是沒有人可以做好.
在這裡就僻住了.
不單僻住.
還要很多的繁文縟節.
買贖罪券.
守很多的教廷的要求.
你要怎樣來告解.
你要怎樣來做很多的聖禮等等.
甚至你去到教堂.
基本上用拉丁文的彌撒.
只不過你是一個觀望者.
你好像都沒有資格去到再深入的參與.
這樣就令到信仰和人的生活.
其實有一個很大的距離的脫節.
甚至中間還要加上一些中介者.
就是教廷對你不同的人生.
都有很多的要求.
所以令到一個普通的平民百姓.
根本就沒有機會再深究自己的內心.
和神的關係.
所以這個清二是by grace.
透過信心.
這樣的演繹.
清楚的去界定了原來.
那個中介每個人都可以得到的.

$^{1041}$其實就是神預備的恩典.
其實這個恩典再去仔細分.
其實就是基督耶穌.
就是主耶穌的成就的救恩.
所以我們所謂的憑著恩典.
就是我們進入和主耶穌的生命的聯合.
就是路德的發言.
就是他看二弗所書五章.
講到基督和教會的聯合.
就是進入這個婚姻.
永恆的婚姻的約裡面.
去享用基督他完美的人生.
給我們的救恩.
生命的改變.
所以這裡就開始進入信心的定義.
我們信甚麼呢.
信甚麼這個才是重要的.
清二的根基.
甚麼叫個人的參照呢.
甚麼是關乎信靠上帝的應許呢.
甚麼是信徒和基督的聯合呢.
這就是信心三個特質.
或者說可能用層次.
其實就不應該有一二三的.
不過這三樣東西就結合為真正的得救信心.
第一 參照的意思.
馬丁路德認為.
要一個人認同有耶穌的存在.
福音書的客觀事實是不難的.
但是要一個人對於他個人自己認同.
基督是為他而生 為他而死.
這就是所謂的信心的前設的特質.
我先讀這裡.
1515年的前後.
他的清二的神學經歷了決定性的改變.
某些人不明白為何唯獨憑信清二的原因.
在於他們不明白信心到底是甚麼.
我們開始講這三個層次.
路德對信心的了解.
第一 信心是個人的參照點.

$^{1081}$而不是純粹歷史性的.
意思很簡單.
你認同福音書的細節.
這就不是清二的信心.
即是說我知道有耶穌的存在.
是呀 耶穌是很偉大的 是一個聖人.
那與我呢? 啊 跟我沒有關係.
今天香港也有很多人聽過關於耶穌的東西.
小時候 或者他讀中學 讀聖經科.
他知道關於耶穌的東西.
說他在海上走 說他五餅二魚.
我認同有這個事實.
但與我無關.
這就不是個人的參照點.
真正的清二的信心.
第一方面就是.
福音書說的主耶穌.
祂不只是為某個人而死.
這裡就很清楚了.
唯一稱得上是基督徒的信心的.
第一個方面.
就是你要無條件的相信.
基督不只是為彼得.
或者其他的聖人而生而死的那一位.
祂是為你這個人來生來死.
即是說祂是愛你.
每個在座的 善上的.
或者將來聽講學的.
這就是救恩信心與神的結合的第一個起點.
就是這個得救的信心.
所以原來得救的信心是這麼重要的.
我們在初信班.
第一堂是邀請人向神自己決志信主.
如果他不肯.
其實說不上信心.
因為你認同福音書說的耶穌.
與你不接受祂.
可以是兩回事.
所以原來路德他真的有認真讀經.
他就真的發現.

$^{1121}$只需要我們每個人聽明白福音書所說的耶穌.
而打開心靈.
讓祂進入我的生命.
接受祂真的愛我這個罪人.
知道我這個人的存在.
那個內心仍然愛.
這樣才是基督教.
這樣才是神學.
仍然愛我們這個人.
這樣你就打開心靈.
決志相信祂.
向祂認罪悔改.
這就是第一個開始.
這個大家應該沒有問題.
第二.
信靠.
信心是信靠.
即是說這個信靠原來是一個有動作的行動.
所以保羅的羅馬書其實有這樣暗示.
是因信而致於信.
是信心的行動.
我信有這位耶穌.
我信祂愛我.
不過問題就在這裡.
原來另外一個解釋關於信心.
這個信靠.
他用的例子都相當好.
信心不單止是一隻船.
當時他們是歐洲.
要去哪些地方.
一定要坐大帆船.
當時的航海學都應該開始蓬勃.
開始慢慢發現新大陸等等.
特別是荷蘭人他們的航海很厲害.
你要去一個地方.
一個很遠的目的地.
你知道必須要上船.
你要相信那隻船.
但問題就是.
你關乎那隻船.

$^{1161}$你要踏上去.
它才載你去彼岸.
所以信靠是說整個生命交托在這隻船裡面.
這個就是信心的信靠的行動.
所以其實有很多人.
如果他沒有這樣的行動.
他的信就停留在一個知識層面.
當然我們不能夠百分百說他不得救.
但其實要自己負心自問.
我們這個信靠.
是不是整個生命已經信靠給主呢.
而在這個信靠裡面.
我們就要附上一個很具體的行動.
就算每個星期你回教會.
都是一個具體的行動.
你知道神喜歡你.
你要回教會.
你要敬拜神.
你要和弟兄姊妹有相交.
有成長.
有對真理一個比較全面的形式掌握.
這些都是我們信靠.
這個不是一個行為的交換.
是信心本身的那個本質.
就是要信靠.
所以馬丁路德這樣的界定.
其實是挺好的.
不是客觀認為.
有耶穌存在.
不過我和他沒有關係.
我不需要他.
那你根本就不得救.
你不需要他.
你靠自己.
靠自己就是不行.
所以當一個人.
他認同耶穌.
他又接耶穌.
接著那個信心就是.
立即去到那個信靠的行動.

$^{1201}$所以大家要問一問.
你信了耶穌.
有沒有持續的信靠呢.
其實真的有時屢見不鮮.
其實真的很想訪問他.
你洗完禮不回教會.
你想怎樣呀.
你是不是很享受那個洗禮呀.
享受沒問題.
但是你要追求.
你要繼續說的是信靠神.
一生的信靠.
是要你踏上這個生命之旅.
所以馬丁路德就是.
這樣很清楚的說了.
那個不只是一個客觀的認同.
而是對一個歷史的認同是沒有用的.
是一個你自己個人的關注點.
就是你已經開始對這個恩典.
你已經開始明白.
你就覺得你需要這個恩典.
所以原來神的義.
那個絕對標準.
原來只不過是神的禮物.
那個恩典就可以立即fulfill到.
你可以想像得到.
是多麼容易得救呢.
第三個.
信靠上帝的應許.
剛才說的信靠就是上帝的應許.
這是第二個解釋.
同樣是第二點.
信靠的對象.
對路德來說.
是對生命的一種永遠不偏離的.
信靠神的觀念.
第三個就很清楚了.
這個就是宗教改革最重要的部分.
信心就是把整個基督徒.
你聯合於基督.

$^{1241}$這個union和基督的聯合.
馬丁路德就不說太多了.
他只是發揮發展他的十誡神學.
如果沒有十誡神學.
沒有那個贖罪.
根本就不會有清二.
他就停在那裡.
因為可能他已經年紀大了.
後來的加爾文.
就在唯獨基督和唯獨信心這幾方面.
就發展得非常淋漓盡致.
所以有不同的神學家.
來將整個宗教改革再發揚光大.
建構在真理.
尤其是加爾文引很多聖經.
所以基督教要義.
大家其實應該看的.
因為他每一個教義.
是引很多聖經的.
其實這個值得我們去參照的地方.
特別是教主一學.
其實基督教要義都挺好的.
就是這三樣東西.
信心不是一個只是認同某些東西.
是你自己信心的一個參照點.
不是對歷史的一個認同.
那個是幫不了你的.
信心.
如果是真的信心.
一定是會信靠神的應許.
神的應許就是我已經在聖經的應許裡.
賜下了這份禮物.
這一個白白的恩典.
這個恩典就是會帶來你和基督耶穌的完全的聯合.
這樣就是整個套餐的信心的定義.
所以這樣的恩信清義.
憑著恩典我們去透過信心.
所以信心不是那個行為.
信心是神給我們的一份禮物.
那個禮物就是恩典.

$^{1281}$我們就是透過信心去到神的面前.
看到基督耶穌就是那個救贖.
我們就得著基督.
我們就得著神的義.
所以神的義就在這個福音上顯明出來.
一章十七節.
接著再有一個經文就是為義人恩信得生.
引述的是哈巴谷書的二章.
保羅在羅馬書.
令到馬丁路德就真的完全的豁然開朗.
他重新再讀羅馬書.
發現神的義原來只能靠基督耶穌才可以得著.
原來信心帶來的就是被神清義的人.
他很自然就會有好的行為流露出來.
所以不是叫人不守行為.
而是行為不可以換取救恩.
信了主耶穌自然就會有好行為.
這個就和天主教完全的不同.
天主教就是說信心加你的行為就是救恩.
馬丁路德和中國改革就是說.
行為不是你去換取.
就是信心恩信清義的結果.
你有和基督的聯合.
你就會有新生命.
就是你的好行為.
就是說你信了主耶穌.
不斷在聖靈在基督裡面.
你自然會流露神改變你的生命之後的好行為.
所以好行為不是換取的交易.
好行為是信心帶來的結果.
信心就是與基督聯合.
信靠和聯合.
基督就是神在人類裡面的救恩的禮物.
因為人不可以自己清義.
不可以得著神的義.
必須要透過基督我們就可以得到.
比如哥林多後書第五章近尾.
神使那無罪的替我們成為罪.
好叫我們在他耶穌裡面成為義.
本來就是這個意思.

$^{1321}$斬釘截鐵的說了.
沒有人可以用行為換取救恩.
信了耶穌有了聖靈有了基督.
我們自然就會有好的行為.
這個就是信心的三個重要層面.
這個也是我們剛才看過.
很早就揭開了.
聖經帶給馬丁路德的改變.
那個影響.
就是他研讀羅馬書的前後.
使他自己終於明白了.
不是行為是信心.
信心本身不是行為.
信心是看到那份禮物已經擺在那裡.
就是恩典的禮物.
就是基督耶穌.
所以信心只是連接和基督的聯合.
所以信徒的人生是不斷思想耶穌.
和主耶穌有愛的關係.
這個就真的需要每天的累積.
每天的成長.
就會在基督裡面.
我們就會享用到基督的神學名詞.
基督的美善.
他的美麗.
就會在我們的生命不自覺地流露.
很自然地流露.
流露不了.
其實要回頭再問一問.
就是說你根本就很少和基督有關係.
很生疏 很陌生.
一個很久沒見的朋友.
連名字都忘記了.
就是因為太久沒見.
如果我們很久都沒見.
所以其實真的.
每個禮拜上教會的人都不少.
當然很少 相對香港.
但是上教會的人和基督陌生.
也都真的不奇怪.

$^{1361}$可能靈唱的 唱詩歌的.
他可能和基督很陌生.
因為他已經變成回來只做這個部分.
所以很小心.
我回來做私事.
但是可以和基督沒有關係.
這些是很危險的動作.
其實應該是彰顯基督的榮美 愛等等.
如果我們回來對基督很陌生.
根本從來沒有關係 沒有聯合.
其實他的事奉一定會很生硬.
很容易就會流露自己的個性 自我中心.
當然絕對不是教大家.
甚至煽動大家來批判.
不是這個意思.
而是其實我們每個人信了主耶穌.
真正的信心不是客觀認同.
是主觀認定我是被基督愛的對象.
是很重要的一個對象.
基督愛我這個罪人.
所以我打開我的心.
讓他進入成為我的生命之主.
這就是信靠神的應許.
和我和基督有一個聯合的信心.
在這個過程.
自然就是恩典流露到我們的生命.
有神的意在我們裡面.
不是地位上.
甚至是主觀我們人生.
不斷經歷神的意在我們裡面的洗滌.
所以與基督聯合絕對不是一個無關重要的要求.
差不多是一個有關生死的要求.
與基督聯合.
你還說自己可以上天堂.
其實是說不上的.
所以每一個基督徒應該就是流露基督.
多少可以是時間的問題.
但是必須是神已經給了我們的吩咐.
95條就是終於.
其實某程度上.

$^{1401}$剛才我們也說了.
一個修士受天主教的薰陶教育.
飲感的奶成長影響.
一般人都是.
馬丁路德就更加是.
怕死 被雷劈 祈禱 奉獻自己.
接著做一個修士.
又讀到書 成為一個神學博士.
又開始在教廷裡面安排教區去做牧養.
做一些行政的工作.
又可能要私餐等等.
等等.
對主教都很敬重.
不斷問主教一些問題.
在他的自傳.
接著一直自己教聖經.
發現與教廷講的傳統.
有一個很大的分別.
其中就是他發現.
一來賣贖罪券的背後動機是極之錯.
因為教廷沒有錢.
還要大家爭來建很漂亮的教堂.
收緝也是其次.
要大家鬥不同的輝煌.
那就變成好大喜功.
那就變成要找一些巧立名目.
就要找一些賣贖罪券.
背後還要是一個錯誤的神學.
你又怎知道你的親人在煉獄裡面.
賣了贖罪券.
他就可以減少那些苦楚呢.
又何來你會得知到.
這種代用券.
就可以解除死了的人的罪的刑罰呢.
天主教分了兩種.
必死的罪靠耶穌.
不是必死的罪.
小的罪.
就要靠贖罪券.
其實不是靠贖罪券.

$^{1441}$其實是靠煉獄.
有煉獄.
將一些小的罪.
譬如貪心.
經常說謊.
就要在煉獄裡面.
誰去判多少年呢.
就沒有人知道.
有這樣的教導.
因為是天主教自己發明出來的煉獄.
就是天堂的前身.
就是你死了之後.
是一個中間的地方.
好像佛教的無間道.
有一個地方就叫無間道.
真是一個地方.
未去到十八層地獄.
天主教都有這樣的觀念.
引兩段聖經.
《馬太福音》第五章近尾.
你和你的仇敵和好.
如果你不和好.
他要拉你去見官.
耶穌說我實在告訴你們.
一文錢沒有還清.
你斷不能從那裡出來.
天主教就說這個就是煉獄了.
就是那個煉獄的前身.
就是那個經文的關鍵.
似乎好像對.
因為只有他們去解釋聖經.
他認為是就是.
所以九十五條.
主要不是抗爭.
不是改革的革命.
是希望提出一些見解.
引發討論.
因為當時原來都是一個習慣.
你有什麼想討論.
你就釘出來給大家看到.

$^{1481}$然後大家有一個理性的辯論.
他期望是一個理性的辯論.
但是因為你一涉及到教廷的可能有錯.
教廷是不能輸的.
教皇是無誤的.
已經有一個金剛圈綁著自己.
所以我很開心.
有弟兄姊妹們給我一些意見.
感恩 原來我沒有想過.
所以不是教皇無誤.
我們是有一點點的誤.
你們認為不對.
我也會有一個從善如流去改.
因為你能夠說我做得再好一點.
其實是很感恩.
是可以更加的徵一求徵.
所以應該牧者是這樣.
越是連之內的.
其實越難勸.
越難聽人勸.
很想勸你都要鼓起勇氣.
其實不應該.
大家是弟兄姊妹.
不過崗位上大家不同.
我們好像已經慢慢忘記了這樣的身份.
好像慢慢忘記了.
我們中國人當然尊師重道是很對的.
但是忘記了.
其實大家是不同崗位的服侍.
大家都是弟兄姊妹.
往往背後的是非.
不是面前的善意的勸導.
充斥在教會裡面.
其實是很可惜.
所以贖罪卷不是偷偷的釘在那裡.
來嘲諷 挖苦教廷.
一個以下犯上的羞事.
完全不是.
其實提過了很多的對話.
但是都沒有得到任何一個善意的回應.

$^{1521}$所以其中對於贖罪卷的議題.
有很多提到.
其他都是比較次要的.
所以馬丁路德看贖罪卷是一個很大的敗筆.
錯得很厲害.
因為很多人甚至對於犯罪.
已經有一種肆無忌憚.
買贖罪卷可以肆意的犯罪.
那個贖罪卷變成更加促使人犯罪.
不懼怕神的犯罪.
是一個很大的後果.
所以根本不是聖經對於罪人要親近神的要求.
所以95條某程度上.
我們大家都需要來深思一下.
逐條讀一讀.
和我們今天有沒有什麼關連呢.
當然主要是當時代的一個文獻.
引發宗教改革的一個原先不是想這樣決裂.
但是教廷不肯改變.
就只能夠成為決裂.
所以馬丁路德不想這樣做.
他都沒有辦法.
某程度上我自己個人就這樣認為.
神既然在馬丁路德的內心.
給他一個這樣的亮光.
有一個生命的改變.
其實就期望他和他的同路人.
有相同的神學.
不是為反而反的.
一定有的 這麼大的運動.
所以慈雲就很決絕和天主教完全勢不兩立.
很激進的.
但是馬丁路德是很溫和的.
所以你看到他的聖參觀是同質說.
慈雲就是我們今天說的.
就是紀念說.
原來當時對聖參都有三個看法.
天主教就是變質說.
馬丁路德就是同質說.
慈雲就是我們今天說的.

$^{1561}$我們吃的是餅 喝的是葡萄汁.
不過是代表 象徵 代表基督的身體.
所以絕對不是基督的身體.
這個餅和葡萄汁就變成基督的身體.
變成基督的血.
天主教就咬定一定會變的.
每一次的彌撒就是基督臨在.
再一次好像再死一樣.
馬丁路德就是同質說.
就溫和一點.
既不是紀念說.
但又不是變質.
兩個中間叫做同質.
相近一點的和天主教那個.
最激進的就是慈雲里的.
後代福音派的.
必須是用這個.
因為某程度看聖經都是這個紀念說.
比如哥林多前書11章提到.
應該是第11章.
保羅說的聖餐.
他們有一個混亂.
哥林多教會.
但他都是紀念主的身體.
今天差不多說到這裡.
叫做馬丁路德的神學.
最重要就是對於行為的滿足.
那個滿足.
是必須透過憑著恩典.
透過信心.
表面上我們今天已經很明白.
每個教會都是這樣說恩信稱義.
沒有人說你不行.
你的行為不夠.
你不可以得救.
很少人會這樣說.
但是我們又要小心.
我們慢慢會不會那個信本身的定義.
或者那個本質.
我們又有一點淡化了.

$^{1601}$因為信心本身不是客觀認同一樣東西.
而是主觀將自己的生命打開.
然後信靠神的應許.
和基督的聯合.
這樣才是信心.
這個我們真的需要去捍衛.
我們必須要一生.
我們個人的辭令.
和教導其他人.
都要這樣去看.
所以這個是很重要.
還有幾分鐘.
希望大家能夠掌握到.
這個我們上一次都有機會看過.
那個真正決裂的.
就是這個叫做沃木斯的會議.
1521年正正式式的決裂.
就是要他收回這個恩信稱義.
收回對教廷講贖罪券的錯誤.
和收回人不需要靠行為就可以得救.
憑著相信基督就可以得救.
不需要任何的行為.
但是馬丁路德就是說.
我的著作已經流通在市面了.
我是不會收回的.
個人對我的攻擊.
除非有聖經的證明.
攻擊基督教除非改變.
但是我都不會收回.
不接受傳統為權威.
只接受聖經.
這個也是重要的分水嶺.
因為教廷已經有很多的傳統.
成為了在基督和人中間.
成為了一個好像中介.
你必須要透過這個中介.
告訴你基督教是這樣的.
你要滿足所有才可以得救.
所以我們很少的.
初信班就浸禮班.

$^{1641}$都是想告訴你聖經的真理.
基本的真理.
講教會的體制講一堂.
其實沒有其他東西.
講你的權利.
你可以洗禮.
我們一生都要守聖餐.
就是這麼多.
沒有其他東西.
根本就是將聖經最核心的.
用幾堂就講給你知.
幾乎就是你的信心.
是否信靠神和基督聯合.
是就可以洗禮了.
沒有其他東西.
洗禮之後.
你有這個信心.
你自然就會流露那個好行為.
所以這裡就是很大的落差.
馬丁路德不接受傳統.
教廷的傳統.
作為權威.
甚至乎凌駕聖經.
記得我們沒有東西可以凌駕聖經.
任何的傳統都是好的.
但如果傳統凌駕了聖經.
我們就要小心和求改變.
所以宗教改革.
是否要再改革呢.
是否要溫和地改革呢.
這些都是我們要問的問題.
當然他不接受.
不肯收回.
他自己講了的.
他在著作裡面已經流通的.
馬丁路德就說.
這個就是我的立場.
除非你們在聖經裡面找到我的問題.
對我個人的攻擊.
我是不會介意.

$^{1681}$這是我的立場.
求主幫助.
某程度上都挺囂張的.
沒有反過來的.
整個教廷.
會勸你 逼你.
你就是一個人.
我站在這個立場.
我是不會改變的.
因為我已經看清楚.
這是聖經講的得救信心.
不靠個人的行為.
成為全眾千古.
我的立場.
上主是我堅固保障.
那首聖詩就是馬丁路德寫的.
很激昂的宗教改革的聖歌.
你可以想像得到.
當時他面對的壓力有多大.
一個人這樣發言.
有小撮的.
當然也有些教廷以外的.
當時的分封的一些國家的國王.
那些元首.
那些諸侯.
因為都很不滿教廷.
都不會認同.
那個權威.
可以延續在個人身上.
所以有些都已經不拜天主教了.
就暗中去支持馬丁路德.
所以他都是有一定的人的幫助.
他才不是一個人赤手空拳.
所以那些君主.
都將馬丁路德收藏.
帶去其他地方等等.
就變成安全地不受到攻擊.
這就是1521年的正式的決裂.
可以說是天主教和基督教的分開.
因為之前還是不斷有書信來往.

$^{1721}$邀請他來開一些聽證會.
要求他解釋書信裡面講的東西.
但是到了這個1521年.
17到21 也就是四年.
正正式式馬丁路德就要離開.
不再以修士天主教的身份.
是以基督徒 神學家 傳代人的身份.
帶領那個時代的宗教改革.
就是這樣.
下一堂還有馬丁路德一些其他的.
小小的枝節可以講一講.
接著很快就講一講慈雲尼.
馬丁路德有一個好朋友叫麥蘭頓.
影響力比較低.
都可能講一點點.
但是更加重點一定要講的是嘉爾文.
因為他才是更加奠定.
宗教改革更重要的神學根基.
就是五個唯獨.
記得數手指.
唯獨信心 唯獨基督 唯獨聖經 唯獨恩典.
最後就是唯獨榮耀神.
對宗教改革那些人.
特別是小朋友.
他們有些教義問答.
人最重要是什麼.
第一個就是榮耀神.
就是你的人生你的生命是榮耀神.
所以你是一個耕種的.
你是一個醫生 你是一個律師.
你是一個傳道者.
總之每個人你的人生就是要榮耀神.
這個是最終極的.
所以人家已經發明了這樣的教理問答.
所以就帶來整個歐洲的宗教改革的復興.
這個復興都經過了幾百年了.
去到二十世紀了.
才開始發現原來世界還有很多地方還沒有信耶穌.
根本從來沒有接觸過基督教.
他們就醒覺了.

$^{1761}$歐洲人就向這幾個大的陸地去傳道了.
所以美國 英國他們就真的輸出傳道人了.
所以宣道會也是當時大復興的產物.
就是孫迅博士在紐約.
紐約他在那個叫回膜教會.
基本上那個時候大復興差不多每天都有報道會等等.
奮興會等等.
這樣來鼓勵更多年輕人去不同的地方去傳道.
就真的有一大班人去非洲 去印度 去中國的時候.
所以從1517到二十世紀.
都有三四百年還沒有這個發現.
這個醒覺要向普世傳福音.
還在歐洲裡面研讀神學.
所以又是一個要小心.
當然當時還沒有那麼多發明.
所以歐洲還有很多地方 比如北歐各方面.
輾轉發現新大陸就是美國.
所以17世紀也有開始去到美國大陸.
所以這個衛斯理約翰經常由英國去美國港島.
開下令會等等.
他們大西洋的穿梭是很頻繁的.
所以慢慢慢慢他們穩定了.
開始神就感動.
開始另一種大復興就是普世宣教的運動.
所以現在中國 印度也不行了 完全封閉了.
非洲就比較多的 比較文明的就順著.
都有基督徒 但也不是全部.
但是中國現在一路輾轉經過了.
其實那個香港歷史 中國教會的事.
二十年代去到中國輾轉經過了各種的內戰.
要全部撤退.
這個也是一個歷史 我們也要來了解.
今晚就講到這裡.
因為我的part point(部分)也去到這裡.
我們一個簡單的祈禱.
多謝主給我們很透徹的看到馬丁路德.
很認真看待他自己的良心.
對神的義的追尋.
他知道自己永遠都沒有辦法達到神的標準.
多謝你讓他確立了信心.

$^{1801}$不是一個客觀的認同.
而是主觀的自己進入那個同神的關係.
以致他能夠信靠神的應許.
同主這份神賜下的寶貴的禮物.
就是救贖主耶穌 是我們的生命.
是我們的一切.
我們能夠在主裡面.
很自然地聖靈就會感動帶領我們.
去活出一個新生命和好的行為.
這個不是換取救恩的條件.
是救恩帶給我們很自然的.
對神感恩的生命流露.
多謝你讓馬丁路德很透徹地分析信心.
使我們能夠明白.
今天很多教會的初信班 浸禮班.
都是很清楚地講這個信心的定義.
多謝主讓這件事沒有失傳.
但是我們很盼望今天的信徒.
對自己的得救是很寶貴.
是看為很重要的生命的禮物.
以致信仰不再是一個裝飾.
不再是一個例行公事.
而是很多香港的信徒.
看信仰是至死有關的.
我們能夠為主作見證.
是一件我們光榮的事.
求主幫助.
多謝你讓我們今天晚上有時間來分享.
禱告奉耶穌基督的名.
阿們.
\newpage



\section{}
\label{sec:ziJGxjL2HKY}
\textbf{2024-11-07 周四晚查經班《改革之路(3)》 講員:曾裔貴牧師}
\newline
\newline
連結: \href{https://youtube.com/watch?v=ziJGxjL2HKY}{\texttt{https://youtube.com/watch?v=ziJGxjL2HKY}} ~~~~ 語音日期: 2024-11-09
\newline
\newline
\hyperref[sec:Sb_mbdX9Sa4]{\small{< < < PREV SERMON < < <}}
~
\hyperref[sec:index_chronic]{\small{[返順時目]}}
~
\hyperref[sec:index_scriptual]{\small{[返順卷目]}}
~
\hyperref[sec:jvEmZbSGJ_A]{\small{> > > NEXT SERMON > > >}}
\newline
\newline
$^{1}$我們一起祈禱.
再一次我們去到第三堂的.
馬丁路德的神學研究.
我們很盼望.
宗教改革.
雖然真的已經超過了500年.
1517年的10月31日.
我們盼望我們能夠看到.
現在改革帶給我們教會.
對真理的尋索.
對自由.
對恩信,清義,同神的關係.
這些的恢復.
已經是落實在我們幾百年的教會裡面.
但是我們很盼望.
那個背後的改革的精神.
是繼續在每一個信徒的內心.
來到去有一個強烈的一個呼求.
使到每一個信徒對於自己的信仰.
都能夠有那個認真.
和去認識.
願神帶領我們.
也都讓我們在今晚.
能夠試一下從另一個角度.
去看宗教改革.
有沒有什麼後遺症.
或者是一些的影響.
願神都施恩給我們.
我們這樣仰望禱告.
奉耶穌基督的名.
阿們.
其實都是想勾起大家一些的思想.
一些回憶.
馬丁路德其實從來都沒有想過.
自己要革命.
要脫離教會.
教廷 當時叫教廷.
就不是一個小小的.
好像我們今天的不同的教會這樣.
馬丁路德自己本身是一個修士.

$^{41}$決心的來到去.
因為那個雷擊事件.
他自己來到去奉獻.
成為不再做律師.
賣光那些法律的書.
然後來到去將自己完全奉獻.
但是其實家人是反對的.
因為那個收入明顯是相差很遠.
決心來到去做一個修士.
很殷勤的來到去.
完成了他的大學學位之後.
後來的神學博士.
他就可以在教會裡面.
無論是施行一些禮儀.
尤其是教書 教主一學.
教聖經.
所以他就有機會去接觸.
當時仍然不是家家戶戶都有聖經.
准讀聖經.
和根本沒有聖經.
沒有自己的德文的聖經.
所以對馬丁路德來說.
他自己本身拉丁文也很純熟.
所以他有機會自己去接觸聖經的時候.
就好像親自去研讀經卷.
就不是靠當時一些不同的教廷裡面的.
一些神父或者一些主教來講授.
所以他自己是直接讓神的話語.
在他的生命裡面流通浸淫.
這是其中一點.
之前第二還是第一堂講過.
當時代在馬丁路德同時代.
其實很大部份的人未必是完全信耶穌的.
有一個很大的運動.
就叫做人民主義的復興.
就是開始對自己的存在意識.
即內心個人的訴求.
要純上個人自己可以去做一些決定的情況.
是較中世紀時代來得更加的迫切.
因為又經過了早期十二,十三世紀的時候.

$^{81}$文藝復興的時候.
開始有一個很強烈的訴求.
因為不再是單一來自教廷的教導.
所以他們要追尋一些古典的典籍.
來做一些研究.
和希望從那些典籍裡面.
找到每個人自己怎樣的一個個人的取向.
特別他們強調的自由.
那個freedom的情況是很強烈的.
不同地方 這個是很大的課題.
不同地方有不同的人民主義者.
相當有學問的大學教授.
來做這些的原典.
某程度上就是復興古籍的那些.
無論思路那些人.
他們以往的年代的思辨的能力.
辯論 修辭學等等.
所以馬丁路德也都是適逢在這樣的時代.
他自己都是受這些風氣的影響.
所以馬丁路德自己.
其實是極之受奧古斯丁的學派的影響.
他直接於在奧古斯丁的神學裡面.
再他自己去建構上去.
那個對於自由意志.
對於和神的那個關係等等.
所以是有一些基礎的.
是有一些時代的那個氛圍的.
有一個這樣的沉澱了的.
當然有一個最迫切.
最靠近的就是教廷的壓迫.
對人民的心靈自由的壓迫.
已經去到好像人的內心已經受不了.
必須要尋求一個出路.
使到不能夠繼續這樣下去.
這個就是可以說是馬丁路德.
適逢在這樣的時代裡面.
今天我就沒有時間說.
這個人民主義.
其中有一個很出名的.
在荷蘭阿姆斯特丹.

$^{121}$當時叫做路特丹.
一個很出名的人民主義的領袖.
叫做伊拉斯姆.
很出名的.
和馬丁路德都是同路人.
但是他就極之反對路德的改革.
和離開教廷.
脫離教會.
接著就自立門戶.
自立基督教.
本來好朋友就變成了陌路人.
這個也是值得我們大家去做一些研究.
伊拉斯姆是相當出名的人民主義者.
所以他們也是同時間.
他們那個追求的典籍.
研讀古代的文獻.
尤其是對聖經的尋索.
所以某程度上.
導致他們要將聖經翻譯成不同國家的語言.
那個文字.
也是來自於人民主義者.
他們對於古籍的興趣.
那個的重拾.
某程度上可能那個時代.
沒有一些有份量的典籍.
來支持他們內心的空虛.
所以有這些情況.
這個是其中一樣.
就是說人民主義對文學.
尤其是古典的很多年前的文學.
他們有興趣.
聖經肯定是最古老.
最具修辭的文學.
所以對人民主義者來說.
研讀聖經不是一個困難的.
不過不是每個人民主義者.
都是朝向那個心靈的尋索.
去信主.
去跟隨馬丁路德的悔改.
因信清義的範疇.

$^{161}$那條路去走.
所以夾在這樣的情況下.
這個就是馬丁路德.
他自己內心的思潮.
是衝擊他的.
所以他也是具備那些修辭學.
因為他本身是一個律師.
他具備這些分析和思辨的能力.
他能夠將一些很繁複的東西.
可以很簡化地來講解.
馬丁路德的講道.
都是相當動聽.
對於當時代的人.
他好像天天都講道.
中國改革之後.
不斷地教學和講道.
這些似乎都是當時.
這些巨人他們的必需.
因為太多的人.
趨向於宗教改革.
崇尚自由.
崇尚自由分成兩派.
大家可能已經分到.
人民主義者追求的是個人的自由.
個人如何去找到自己.
對於自己身份的尋索.
但是馬丁路德.
一直停留在聖經的世界.
神學的世界裡面的自由.
所以某程度上.
是有神的話語.
去捍衛著馬丁路德的心靈.
他不會受到當時思潮的影響.
脫離了對於神學的追尋.
所以他理解的改革.
只不過是在神學上釐清.
什麼是在神的面前的義行.
而這個義行.
他窮他一生都清楚.
不是自己去換取的.

$^{201}$沒有人可以去換取救恩.
只能夠是神的義.
藉著我們的相信.
賜給我們.
我們上星期.
即是前星期我們說過.
這一切都是恩典.
所以恩典不是一個交換.
不是一個功德.
不是一個工作.
是我們沒有辦法做到的.
是神的義.
白白地加給我們信的人.
信神的人.
所以你可以說.
馬丁路德他面對這麼大的.
那個時代的一個思潮.
他沒有受到影響.
而一個好朋友.
伊拉斯姆成為他的一個敵人.
其實是很可惜的.
即是說大家本來都是改革的.
但是伊拉斯姆.
這個我下星期會再補充多一點.
給大家一些powerpoint.
因為今天我做的是另一個層面的.
那個時代的衝擊.
對馬丁路德.
對那個時代的信徒的衝擊.
所以你可以想像到.
本來是好朋友.
大家都是走在改革之路的.
但是一個就朝向人民主義.
即是人本主義.
人到底是存在在那個時代.
我們是不是還要受教廷的操控呢.
還是不是要買贖罪券呢.
還是不是要其他其他的東西呢.
所以原來改革.
都可以有幾條路走的.

$^{241}$我今天就會講另外一個.
就是更嚴重的.
就是一個戰爭的改革.
我們先看看這個powerpoint.
這個就是馬丁路德一生.
他從開始.
到他的離世.
他只是很想釐清.
到底什麼樣的人.
可以在神的救恩裡面呢.
所以觸動他的.
就是贖罪券.
當時已經去到失控的.
一種這樣的販賣.
販賣到一個地步.
簡直是將救恩完全扭曲了.
為什麼神職人員會這樣.
他們有些代理人.
教廷的主教.
每一個地方的主教.
就會委派誰誰去賣.
你去到希敦.
你去到斯敦.
你就去河北.
去賣贖罪券.
逐家逐戶地賣.
誰給的錢多呢.
那些錢不是給我們的.
是給教會來興建教堂.
修葺教堂.
所以名義是很好的.
很堂而皇之.
但是大家可能有bonus.
或者有回佣.
我不知道.
各自就會有不同的賣贖罪券的促銷.
如果你給多些捐獻放進去奉獻箱.
你在煉獄的親人.
他就可以快些脫離火刑.
記住煉獄永遠不是那些沉淪去火壺的罪.

$^{281}$是每個人自己犯的小罪.
可以饒恕的罪.
不過就要在煉獄裡面受一些火刑.
這是天主教當時代的宣稱.
所以永恆要下火壺的罪.
只能靠耶穌的救恩.
但是煉獄仍然需要一個信徒去面對.
所以就出現了這個中介贖罪券的手法.
一個這樣的手段.
使賣贖罪券的人以為自己好像已經完成了神的赦罪.
所以更加嚴重的就是激發了那些人有贖罪券在手.
更加的任意的犯罪.
就令到馬丁路德感覺這件事一來是很嚴重.
二來是很褻瀆神的救恩.
所以他是看得很清楚自己在神的聖經.
特別是羅馬書的一章十七節.
「這信是本於信,以至於信,正如經上所記,為二人因信得生」.
當他重複看整本羅馬書的時候.
發現內心裡面不知不覺間.
因為自己進入了信心裡面.
不知不覺間他自己一直在內心裡面的罪疚感.
那種自己沒辦法解決自己內心的罪的那個好像捆綁.
一下子如釋重負.
所以他更加再去詳細讀羅馬書的時候.
他就更加確信這個因信稱義.
神的義就是因信而得到這個客觀的被動的稱義.
就是神主動將這個救恩賜給人.
所以對馬丁路德來說.
永遠他停留的是在神學上的那個勸說.
那個更新.
那個希望的改變.
所以嚴格來說.
他一生都不覺得自己是改革.
是革命的.
不過你這樣想.
剛才我們說了.
有些溫和的就是尋找那個典籍.
希望透過古代的文獻.
來釋放自己的信仰.
個人的權利.

$^{321}$自由.
這個就是人民主義的路線.
但其實在人民主義.
在馬丁路德的神學的路線.
還有一條路線.
我們今天必須要說.
就是革命.
戰爭.
但這一場的戰爭.
是令到德國的那些農民.
死了10至15萬人.
在幾年間被鎮壓.
所以那個困難和那個悲哀的地方.
我們一路看下去的時候.
就會知道那個情況.
我們再簡單去看看.
馬丁路德.
看完這兩個方案就夠了.
我就引述了.
我有帶來的.
那個95條.
我用一張紙讀給大家看.
因為現在有些人名遮住了.
我看不到.
第21條.
他說.
賣贖罪券的人.
其實那些贖罪券的那些人.
那些代理人其實沒有罪的.
是叫他去賣的人.
他們就有罪了.
他說贖罪券叫賣者宣稱.
教宗的贖罪券.
能夠使人免除一切的刑罰.
並且得救.
完全是荒謬的.
就是由教宗那裡發起.
各地的主教.
來到去.
可能有quota.

$^{361}$你要賣多少張.
教區要賣多少.
就要不斷的去推銷這些贖罪券.
去到一個地步.
是極之荒謬的陳述.
所以說.
這樣就可以免除一切的刑罰.
甚至可以得救.
離開那個煉獄.
這個是.
馬丁路德看為是一個很荒謬的謬誤的事.
這個是第21條.
留意一下第51條.
基督徒.
我看下一頁.
我拿走這個人名.
這樣大家就看得清楚一點.
可能大家會看得好一點.
如果不是就.
因為你們沒有預備.
第51條.
基督徒你要知道.
教宗願意並且應當自己.
他自己出錢.
去補償那些被贖罪券叫賣者.
騙了的那些錢的人.
就是說.
教宗他有這個責任.
你曾經叫人買贖罪券.
但是馬丁路德當然不只是這裡.
是說95條.
都是說贖罪券的功效.
他說如果既然真的沒有效沒有用.
教宗就有責任.
要將所有的.
其實也不是叫做鵝片.
賣了的那些錢.
全部歸還給所有的信徒.
因為是沒有效的.
不過他甚至乎這樣說.

$^{401}$因為當時他們有一個很重要的原因.
為什麼要賣贖罪券.
就是教廷沒有錢.
沒有錢去收緝和興建教堂.
尤其是當時的聖彼得大教堂.
在羅馬.
所以你可以想像得到.
馬丁路德就說.
教宗有責任.
甚至乎要賣了整個聖彼得教堂.
也在所不惜.
所以你可以想像得到.
95條在威登堡大學的門口.
釘了晚上.
接著第二天.
當時又沒有什麼狗仔隊傳媒.
但是印刷商就將這95條大量印刷.
接著大量的分發給其他的地方.
本來是一個很小的威登堡.
後來10天左右.
整個歐洲已經知道這95條.
居然有一個修士.
膽敢地挑戰教廷.
膽敢地說贖罪券是毫無功效.
你可以想像得到.
那個壓力.
那個挑戰有多大.
所以你可以說.
馬丁路德只是停留在這個壓力下.
他只是在神學上.
希望大家彼此辯論.
他期望辯論.
理性的辯論.
到底贖罪券是否真的像教廷.
所聲稱的有效.
是從哪裡聖經引伸而來的呢.
還是只是一種.
將它變成一種更加奴役的一般平民百姓呢.
所以你回去再.
可能你真的印了一份出來.

$^{441}$有空就去看看.
到底贖罪券對當時的人有什麼影響.
今天我們未必.
會再有人.
有這種錯誤的想法.
來到教會一定要做某些事.
我才可以上天堂.
我們也不會再這樣傳.
但是如果我們的信仰不是一個很純正的.
透過認識聖經.
透過認識主耶穌.
透過信心和神的結合.
我們都會有我們自己.
由傳統慢慢成為一種枷鎖.
甚至是一種無形的枷鎖.
這個是我們需要避免的.
因為中世紀已經去到這樣的情況.
所以似乎去到黑暗的一個景況.
一點光都感覺份外的光明.
所以馬丁路德的95條.
就帶來一個很根本性的.
對教廷的挑戰.
這裡要小心.
因為馬丁路德你自己本人.
是完全沒有這個想法.
要挑戰 要背叛.
和要推翻教廷.
從來到死他都沒有這個想法.
但是其他人已經在思潮.
已經在這種種子的土壤.
就帶來可以是很嚴重.
看你是朝哪個路線走.
剛才我說了第一種就是.
人民主義只是希望回到典籍.
回到古代 回到聖經怎麼說.
所以也引發正面的就是.
聖經翻譯成為一個很蓬勃.
但是他們只是從個人.
在那個時代.
應該有自己的個人自由.

$^{481}$那個方向.
他們就不是朝向神學的方向.
但是我今天想和大家.
比較重點要說的第三條路線.
就是這個農民戰爭.
革命帶來的 宗教改革帶來的衝擊.
就是如果連一個修士.
都可以夠膽挑戰教廷.
對教皇的不敬.
甚至乎聖彼得教堂.
是整個天主教的中心地方.
都說可以隨便賣了它 燒了它.
這樣的說話帶來的震撼.
就是說原來任何一個人.
你都可以不再懼怕教皇對你的信仰.
對你的靈性 那個判決 那個定罪.
這樣的時候.
就開放了一個很大的空間.
如果有些別有用心的人.
他可能很多年都受著壓迫的.
受著不公的.
他這個不公是很真實的.
這樣就會帶來.
他可能會借助了這個宗教改革.
而導致到那個社會的動盪.
果然就是在1524年.
1517年10月31號的萬聖節.
就是宗教改革.
就是應該95條的釘在了威登堡教堂.
就是大學裡面.
就引發整個歐洲那半年.
整個歐洲就天翻地覆了.
但是去到1525年.
就出現了這個激進的運動的領袖.
這個文次爾.
因為那個德文我不懂得讀.
我只知道他的翻譯.
都是這個翻譯.
大部分的書都是用這個翻譯.
這個人原來和馬丁路德一起.

$^{521}$在宗教改革上這樣同路人.
所以剛才的伊拉斯姆和文次爾.
都是同路人.
甚至乎他是一個傳道者.
一個主教.
他因為已經離開了天主教.
他主持完聖餐之後.
1524年的3月24號.
是有一些書的記載.
他主持完聖餐之後.
當時有一些叫做田農.
做農夫的那些人.
因為當時代中世紀.
就是馬丁路德的時代.
有一些是貴族.
特別是那些君王.
就是那些很多不同的分封的王.
就有他們的轄區.
有他們自己的.
屬於他們自己的那個.
又不是叫教區.
是他們自己擁有的地方.
就細集的了.
所以永遠都是貴族.
他們就會給錢一些很窮的人.
幫他們耕種.
所以又會有一些.
永遠都是勞.
等於美國年代的黑奴.
那個階層已經分了.
是不能夠某程度超越的.
可以給你基本的溫飽.
但是他們用另外的方法.
去壓制那些田農.
例如你有好的收成.
那些收成就不是屬於你.
全部屬於你的.
因為你在我的轄區裡面.
做好了你的事.
其實你要納稅的.

$^{561}$因為我要有譬如軍隊.
我要有其他的公共設施.
我要來維持.
所以你只不過在那裡為我耕種.
其實那塊田就是我們的.
是那些貴族的.
所以那些農民.
永遠就是這樣不斷地做.
但是實際上.
可能只是夠溫飽.
而且再加上.
還有來自教廷的壓迫.
又要他們買贖罪券.
又有貴族對他們的徵稅.
又有其他的.
很多不同的壓迫.
所以導致在那個時代.
其實貧富懸殊.
已經是相當之普遍的了.
所以這個情況就是.
這個文次義.
他就跟隨宗教改革.
但是相當激進的改革派的領袖.
所以一拍即合.
那些田農.
當時有50個的田農.
就草擬了12條的條款.
寫好了12條的條款.
我等下和大家逐條去看一看.
你就知道當時代的壓迫.
去到什麼樣的情況.
所以在這個3月24號.
似乎是主日.
主持完那個聖餐之後.
這個文次義就發動農民戰爭.
就是要去推翻教廷和那些貴族.
所以當時已經是侵略了一些貴族的地方.
尤其是燒教堂.
和佔據一些教廷的產業.
這樣的一個戰爭.

$^{601}$一發不可收拾.
在整個德國就這樣蔓延下去.
在這個時候.
因為馬丁路德是宗教改革的領袖.
他又是在那個時代.
他被迫.
因為是你自己帶來的對教廷的不信任.
你是鼓動很多的平民百姓.
跟隨你離開教廷.
不再在天主教裡面去聚會.
而是你自己另立門戶.
有基督教.
路德派開始自己的聚會.
守聖餐.
和有自己的教導體系.
例如教聖經等等.
所以引發的.
不原來那麼單純的所有的信徒.
都是趨向於那個恩信清義這個部分.
原來很特別的.
當看文獻的資料.
我今天預備的時候.
這一班所謂的革命的人.
其中有50個的代表.
當然他們還會繼續鼓動.
因為太多人受壓迫很久了.
一直都是受壓迫.
你看看那些內容12條款.
你就會知道全部跟他們的生死存亡有關.
逐樣去看看.
你就知道.
為什麼這50個人.
在這個文賜爾這位.
都是傳道人.
宗教改革的一個領袖.
不過後來那個運動完結.
他就被捉拿和被處死.
這個文賜爾是馬丁路德的好朋友.
你就可以想像到.
他也是改革.

$^{641}$他走一個極端的要徹底的改革的路線.
所以我們很快就要反思.
那個改革到底要停留在哪一個界線呢.
我們先看回這裡.
給大家一點印象.
當是一個興趣大家聽一聽就算了.
很特別的.
就是我們要有自己的牧者的委任.
那意味著一直都是教廷派人來的.
去到這個牧區我派人來.
那個人是不是他們喜歡呢.
都不能說的.
全部是要接收的.
所以他們第一個就是說.
我們要選立屬於和我們認為適合我們的牧者.
在我們的牧區裡面.
所以嚴格來說.
這個運動所有的農民都說自己是基督徒.
第十二條的條款.
就是說如果以上的這十一條.
是有違反聖經的.
我們就會願意撤銷.
就是說他們一直在這個運動.
聲稱自己很清楚是神的聖經的教導.
這樣來制定.
還有要那些貴族那些地主.
一定要執行.
要符合就是說認同.
簽這個大家彼此的條款.
當然正正就是不行.
第二個就是在我們的轄區裡面.
那個十一奉獻是不可以歸給地主.
是應該我們自己教會.
我們自己的牧區.
自己去運用.
用來給我們自己的那個牧區裡面的需要.
例如給一些新津那些傳道人.
或者那些的需要.
例如那些的燈油火蠟.
或者其他的關於傳福音的工作.

$^{681}$這樣一個叫十二條款.
第三就是我們只會馴服.
由選舉產生出來的統治者.
就是那些貴族.
那些的領袖.
就是整個國家的領袖.
是要由選舉出來的.
你任命我們我們是不接受的.
這樣開始由個人的那個自由.
去到國家層面的地上的政治.
他都想干預.
其實這些是很微妙的.
就是等於衝擊著整個政治的那個任命體系.
不是我們自己選出來的.
我們就不會馴服了.
這樣就是留下了一個很大的一個空隙.
一個尾巴.
第四似乎就是那些地主.
一路以來對那些農民的壓迫.
水的資源.
和不准他們在田地以外.
有一些的收入.
有一些的補充.
甚至乎那些水源.
都不准他們隨便去開採.
不准他們捕獵.
不准他們去做一些其他的.
種地以外的收入.
所以這個就真的直接到.
去到令到那些農民.
是受盡很大的剝削和壓迫.
所以這些某程度是一個很極致的反擊.
所以一會兒我就會說馬丁路德.
最慘就是那個牧者就很慘了.
兩邊都不討好.
因為你說了12個條款.
有一條不聽.
我們就立即革命.
那個革命是立即去佔據那些教堂.
去火燒一些物資.

$^{721}$來將那些地主迫害等等.
接著就說森林的那些木材.
當時德國都很大.
很多大的森林.
他們可以去砍伐那些樹木.
作為那個材料.
是不准的.
所以他們說要有一個開採樹木的權利.
是這些農民可以這樣做的.
第六就是地主對於那些農民的奴役.
你做完你自己應該做的事.
之後你就要幫我們地主做我們要你做的.
這個就變成了無盡的剝削等等.
變成真的會令到那些農民.
積壓了太多太久的那些不滿.
和那些層層的剝削.
所以導致到有這樣的一個事件.
這個事件某程度其實不關馬丁路德的事.
因為你的改革.
是灌注了一個對教廷的不信任.
這個問題就變成了這些農民和文次兒.
都繼續拿來更加激進的改革.
更加激進的要推翻整個地主的霸凌.
所以你可以說馬丁路德是不是要負一點點責任呢.
等一下我們又再提.
第七個 私人財產不可以被地主剝削.
因為層層的稅項.
他們見到你收成好.
他就剝削你每年交不同的考納名目的稅項.
所以令到那些田戶的收入是足襟見爪.
真的是很慘的.
當然那個問題就是那個田租.
不斷的來加租.
來令到那些田戶受盡那種剝削.
還有新的法律全部都是favor那些田主.
那些貴族.
導致那些田戶受盡一個在法律的保障下的不公平.
你可以想像到對他們的影響.
要充公屬於社區的土地.
不可以再給那些地主霸佔.

$^{761}$地主有自己的軍隊.
他有自己的勢力.
所以有些本來是屬於公家的地方.
他們都拿來自己用.
所以現在這個革命.
這個十二條款就是說.
要將這些公用的地方充公了.
不再屬於那些地主.
這個遺產稅我不知道有多嚴重.
就放在第十一條裡面.
要廢除這個遺產稅.
就是說是不是一個人死了.
他擁有的東西.
都要由那些地主和貴族去徵收.
所以可能真的導致層層的剝削.
相信其實全部大部分.
如果是很講到那個靈性的.
有的不過是那個自由.
其實有一個很嚴重的就是.
一直以來的公平的問題.
怎樣去擺平這個公平呢.
在這樣的怨氣底下.
似乎就凸顯了這個公平.
和那個貧富懸殊.
就用什麼手段來達至一個.
在農戶眼中的公平呢.
就是要將地主的擁有.
就要立即無條件地輸送給這些農戶.
不過第十二條就很有趣.
他說以上的全部是適合聖經的.
我們是基於聖經.
我們是以信徒這個角度層面.
來提出十二條款的.
所以就要馬丁路德拆這個蛋.
因為馬丁路德已經是當時的.
整個德國的領袖.
大家不要用今天的眼光去看當時.
今天任何一個教會的領袖.
都不可以在香港有什麼話事權.
但是馬丁路德和嘉爾文那些就不是.

$^{801}$當時的德國和當時的瑞士.
他們兩個都是有很強的話事權.
所以馬丁路德現在就要面對這樣的戰爭.
他們已經侵佔了教堂.
焚燒一些產業.
又要殺害一些貴族.
那怎麼解決呢.
馬丁路德怎麼去面對.
說這些人自己是跟聖經來做.
又說是長期受到剝削呢.
馬丁路德作為一個領袖.
你就沒辦法不面對.
首先那個戰爭已經開始了.
你怎麼去解決呢.
馬丁路德就寫了一篇文章.
一個叫和平的訓誡.
如果大家拿來看一下都好.
其實很清楚.
兩邊都不好.
但其實他兩邊都要譴責.
首先就是任何暴力的抗爭.
一定不是聖經.
一定違反聖經.
我和大家揭開一兩段經文.
他都是引這些經文的.
羅馬書的12章.
其實在社會事件的時候.
這些經文都是被引述過的.
因為教會始終都要表態.
面對這些這麼大的議題.
12章羅馬書.
或者在這一大段裡面.
由16節開始.
因為下面都有出現.
幾重要的大方向的原則.
16節要彼此同心.
這個相當的說是在教會裡面.
和不同的弟兄姊妹不同的看法.
都要志氣高大.
個人的膨脹.

$^{841}$自己個人的意識.
都要苦就卑微的人.
或者卑微的事.
苦就即是你自己明明有理.
有那個sense.
有那個智慧.
但你願意去用一種很謙卑的心態.
去和他在那個水平去思考.
引導他進入更高的思考範圍和水平.
這些叫做苦就.
不是一句說了他.
你懂什麼?.
這些就不是叫做苦就.
這些就是志氣高大.
不要以為自己聰明.
一個相對的說法.
可能那個人的角度去看事物.
可能是很有限.
那個有限可能有point都不一定.
所以都不要這麼快自作聰明.
就覺得一概否定別人.
總是可能有一個point他對的.
不要以惡報惡.
如果遇到一些人對你不公平.
甚至惡待你的時候.
不要以惡報惡.
眾人以為美的事要留心去做.
Pay attention 認真去思考.
這些美事.
不要整天把惡放在自己的心.
有時候有些人.
不知道你有沒有見過.
整天都很生氣.
整天都好像殺了他全家一樣.
就是說他的內心經常以惡報惡.
經常都想用惡去保護自己.
這樣其實不健康.
保羅說.
不要之餘還要時時想一些美善的事.
你的內心就會美善.

$^{881}$你就會能夠時時有一種善良.
所以我們不要想著.
只是我不對付你.
其實不是.
我們內心要有善.
所以就會去到十八節.
若是能行.
總要盡力與眾人和睦.
因為你自己內心有了那個善.
有了那種愛.
你就有力量.
就會多走一步.
如果沒有愛和善.
我們就不單止不走多一步.
甚至我們退後一步.
有時候突然間那些念頭.
這些魔鬼最容易.
就是會給我們一個念頭.
這個人真是壞到沒得救了.
都是經常得罪我的.
經常都想和我作對.
你經常這樣想.
你就會有一個內心的報復心態.
和一個防禦.
你這個防禦.
開始去到一個越來越不健康的那種.
隨時想別人怎樣害你.
想多了.
其實自己就會有這種報復心態.
所以你一定不能夠達到.
那個美帶來的那個行善.
就是與眾人和睦.
所以我們自己一定要保持自己的善良.
常常想事情.
不要從別人一個壞的角度.
或者壞的動機去想人.
一定要自己首先存一個善良的動機去想人.
但是我們要防.
但是千萬不要變成一種.
一個用自己的.

$^{921}$好像天天起來就武裝自己.
沒有人可以害到我的.
我們這樣做得多.
我們不知不覺就戴了假面具.
和穿了一套好像防禦的衣服.
我們慢慢失去了什麼呢.
就是那個善良.
十九節.
親愛的弟兄不要自己伸冤.
寧開讓步聽憑主路.
因為經常記著主說.
伸冤在我.
我必報應.
這就是馬丁路德對這些農民戰爭的直接譴責.
就是和平訓誡裡面.
引述的羅馬書這段經文.
記清楚不要自己伸冤.
寧可讓步聽憑主路.
因為經常記著主說.
伸冤在我.
我必報應.
所以你的仇敵若餓了就給他吃.
若渴了就給他喝.
因為你這樣行.
就是把炭火堆在他的頭上.
你不可為惡所勝.
反要以善勝惡.
相當清楚.
引完《箴言》的經文之後.
在第21節的一個總結.
仍然都是不要自己伸冤這個課題.
那個的大方向.
就是說一定要以善.
就是自己內心的善.
或者說自己成為了基督徒之後.
神的愛.
神的美善.
充滿你的心.
和有永恆盼望的基督徒.
這樣的善.

$^{961}$勝過地上的惡.
無論惡人.
或者是惡士.
這樣的一個大方向原則.
還有包含了第13章.
所以第12和第13章.
一拼去看.
你就更加具體地看到.
不可能不是說.
這種面對政權.
我們那種.
到底是一個奈約.
還是真正的勇氣.
這個勇氣不是用暴力的勇氣.
是用愛.
用善的勇氣.
因為在上有權的人.
人當順服.
因為沒有權並不是出於神的.
訴諸於上帝的.
凡掌權的都是上帝所命的.
所以抗拒掌權.
就是抗拒上帝的命.
抗拒的必自取刑罰.
等等.
所以你會看到.
一直都是在說.
那個秩序.
說一個基督徒.
處身在一個地上的秩序.
我們是要遵守的.
我們不可以違抗的.
但是.
第12章那裡就是.
包括了心靈裡面的那種勇氣.
就是善良的勇氣.
去面對那些不公.
那些無論是迫害.
或者是那些惡事.
或者惡人.

$^{1001}$所以這個和平分界.
是直接相當嚴厲的.
所以又是招來那些農民.
對馬丁路德徹底的抗議.
抗拒.
不幫他們.
所以令到整個農民戰爭.
馬丁路德是支持貴族去鎮壓這些暴亂的.
所以為什麼死了十幾萬人.
就是強烈的鎮壓那些農民.
所以馬丁路德他.
給我們看到的歷史的問題.
就是他兩邊都勸說.
對於農民的勸說就是.
要他們立即停止.
不能夠用暴力.
還有他們所提的訴求.
有些是合理.
但是有更多是涉及到.
是由地上的法庭.
地上的那個法律去分配的.
就不應該是由教會.
是由一些神職人員去支持和去分配.
例如那些屬地.
例如那些十一奉獻等等.
就不應該是由教會去干預那些不同的地區.
他們那個已經制定好的法律.
所以馬丁路德對於農民的革命.
他們表面就說.
我根據聖經的.
我們是想要那個公平.
因為那個壓迫是很多年了.
是很結構性的壓迫.
對於我們的干預.
我們那個被剝削.
貧富懸殊的那個這麼嚴重.
當時真的是這樣的.
所以馬丁路德.
你自己本身是出於一個貴族家庭.
因為他的爸爸其實是開一些礦場.

$^{1041}$他爸爸都是一些有錢的.
但是其實有些人就這樣批評馬丁路德.
根本不理農民的那個疾苦.
但其實馬丁路德自己.
曾經都沒落過.
他自己都會有他們自己的一些田地.
他都要這樣來做.
所以其實馬丁路德不是不了解農民的困苦.
不過仍然是說.
你用激進的革命的暴力的手段.
是絕對聖經禁止.
所以馬丁路德是很義正潛移.
贊成地主和貴族.
強烈的盡快的鎮壓這些農民.
所以在幾年裡面.
十萬至十五萬的農民被鎮壓而死.
他的好朋友文次義被捉拿,逮捕.
然後被處死.
所以都是一個當時的幾大的宗教改革的污點.
某程度上就不關路德事的.
他一直都在說.
我只是想在神學上做一個更新.
撥亂反正.
沒有想過要脫離教廷.
是教廷趕他離開.
甚至要追殺他.
他才能夠逼不斷在諸侯王的保護下.
來隱藏.
所以這樣的脫離.
就不是馬丁路德帶頭鼓勵人們脫離教廷.
所以這樣的情況.
對於農民戰爭.
他們就假借宗教改革.
假借原來教廷是腐敗不堪.
假設這些貴族用宗教信仰.
其實當時整個國家都是用宗教信仰.
沒有人是不信的.
所以這樣的層層的壓榨.
就令這些農民必須要作出一個抵抗.
所以就帶來.

$^{1081}$他們只能夠用訴諸於最激進的手段.
重奪這些屬於他們的財產.
自由.
訴諸於一個真正公平的社會.
所以這些就很複雜了.
因為涉及到個人的利益.
和受到另一個階層的壓迫.
要這樣來拉勻他.
這些只能夠.
要不就用暴力.
要不就協商.
但是暴力也不行 協商也不行.
只能夠接受社會本身存在不公平.
只能夠是那些貴族.
他們自己良心發現.
自己看到自己的問題.
所以剛才我說傳道者就是兩邊不討好.
馬丁路德不是單方面偏幫貴族.
馬丁路德說那些農民不是你的敵人.
你們是神的真理的敵人.
這就令到那些貴族嚇一跳.
因為表面上你叫我們鎮壓那些農民.
因為他們在用暴力對付那些貴族.
但是現在反過來.
馬丁路德又說有某些條款.
你們是神的僕人.
因為你們去壓迫他們.
用盡種種的手段來抽重稅.
又侵佔一些公用的土地.
又禁止這些天然資源.
讓其他人可以來享用.
你們其實就是導致這個戰爭的始作俑者.
就是這些地主和貴族.
所以又導致馬丁路德對他們.
又是一個很無情的譴責.
說這一句已經足夠了.
你們不是那些農民的敵人.
你們這些地主貴族是神的真理的敵人.
但是他們有沒有悔改呢.
當然就難了.

$^{1121}$因為涉及到經濟和很切身的利益.
因為那些的封地是世襲的.
如果給了其他人.
我們就什麼都沒有了.
所以那個情況是一個複雜的.
所以這個農民戰爭又不是沒有影響的.
因為慢慢過後的幾十年.
需要慢慢的時間.
大家都開始在這些這麼嚴重的衝突底下.
整個德國就已經去到內亂.
在這麼大的影響底下.
大家都開始有些沉澱.
大家都有一些反思反省.
慢慢就引進了更大的一些公平.
或者一些的分配.
所以都不是沒有影響的.
始終人都不是鐵板一塊的.
地主或者貴族都開始有一些反思.
是不是真的鎮壓得太緊要.
剝削得太緊要.
導致這些農民有這樣的壓迫底下的抗爭.
不過用的手段是一種極之暴力殘忍的手段.
去殘殺,暗殺.
所以導致到這個不愉快事件.
你可以說是宗教改革.
另外一個的煽惱.
人訴諸於自己內在的公平.
分配要均勻.
誰去做那個分配者.
所以這些就變成一個極大的矛盾.
所以馬丁路德就想置身於事外.
他說這些涉及到有地上的法律.
地上已經有的分配制度.
不是一下子就可以推倒重來.
不是一下子最好的標準就立即要做到.
所以他就很清楚.
有另外一封信寫給當時的貴族和革命者.
他說牧者扮演的角色.
是在神學上做一個教導.
做一個真理的詮釋.

$^{1161}$和一個勸說.
所以牧者不能扮演分配者.
來做調停人.
甚至扮演直接將整件事.
例如那些十一奉獻.
例如那些牧者的委派.
如果由牧者去介入.
就會帶來一種角色的重疊和衝突.
所以馬丁路德就很想避開這些事情.
所以最終都是由一個不愉快的結果收場.
就是用軍隊 用武力.
去鎮壓這些平民百姓的田戶.
導致在這個衝突裡.
我忘記是不是有13年.
最終是10至15萬.
國務司說是15萬.
有些記載說可能是10萬多點.
不是小事.
因為當時是因為宗教改革那個時代.
而導致馬丁路德還是如日方天的.
宗教改革的勢頭的時候.
突然間有這樣的戰爭.
真的會帶來很大的傷害.
所以這個又給大家有一個小小的認識了解一下.
對於當時的這個情況.
尤其是那個農民的訴求.
我們是根據聖經.
我們是沒有違反到聖經.
我們要用這樣的手段.
都是聖經賦予給我們.
這個在馬丁路德來說是沒有可能的事.
怎麼可能你會訴諸於暴力.
來做這些事情呢.
他還有引述《馬阿泰福音》第五章.
我很快讀給大家聽.
其實大家都應該很熟的了.
但是不要小看這些經文.
在這些這麼大的事件裡面.
耶穌基督的教訓就是最直接的那個真理.
如果不是其實就很難了.

$^{1201}$第38節.
你們聽見有話說.
以眼還眼 以牙還牙.
只是我告訴你們.
不要與惡人作對.
有人打你的右臉.
連左臉也轉過來由他打.
有人想要告你.
要拿你的女衣.
連外衣也由他拿去.
有人強迫你走一里路.
你就和他走二里.
有求你的就給他.
有向你借太的不可推辭.
你們聽見有話說.
當愛你的淪陷.
恨你的仇敵.
只是我告訴你們.
要愛你們的仇敵.
為了迫不得你們的禱告.
諸如此類.
一直這樣下去.
你就可以想像到.
主耶穌的教導登山寶訓裡面.
說的是用愛去饒恕.
用愛去化解很多的爭執.
是一個大的脈絡.
是一個大的方向的內心的心態.
所以對馬丁路德來說.
這些是不在話下的.
他很難想像.
他自己也完全不敢去得罪那些教廷.
他自己從來沒有想過.
要違反教廷的教規.
他只是提出辯論.
但是居然有些填護.
集體的.
由一些激進的領袖的鼓動下.
都還沒有談好.
就已經放火.

$^{1241}$就已經去毀壞當時的社會的正常運作.
這些是暴徒的所為.
所以馬丁路德說.
那些填護其實根本就不是信耶穌.
根本就不是根據聖經.
根本就是在這裡做一個暴力的抗爭.
無社會無政府主義的秩序.
是沒有辦法接受的.
對馬丁路德來說.
所以他引述的這兩段聖經.
其實是足以對當時的填護.
是一個譴責.
也就等於給了權力.
這些地主和貴族.
用軍隊去鎮壓這些填護.
所以這場戰爭.
嚴格來說不應該叫做宗教戰爭.
但是在歷史裡面.
因為涉及到宗教改革.
就變成了一個這樣的農民戰爭.
涉及到信仰.
因為他們都說.
我們要找自己的傳道人.
我們轄區裡面的奉獻.
不可以歸給這些地主.
諸如此類就好像涉及教會.
宗教的層面.
但其實馬丁路德看得清楚.
不涉及到宗教的.
是假借宗教的名義.
來將內心的所謂公平不滿.
來去極化了它.
變成了一種激烈的抗爭.
所以對填主和貴族來說.
就是有馬丁路德的開綠燈.
鎮壓那些平民.
就變成有理可據.
因為是馬丁路德自己引述.
他們不合適聖經.
所以馬丁路德又受到譴責.

$^{1281}$被後人.
完全不站在那些填戶的一邊.
去提出抗爭.
所以兩面都是難做的.
因為你是一個宗教領袖.
你要自己擺一個持平的位置.
兩邊都不幫其實是很難的.
因為大家都有自己很強烈的利益衝突.
所以你不站在我那邊.
你就是敵對的了.
那些地主就是.
你幫那些農民.
等同於你就是要幫我們革命.
所以對馬丁路德來說.
他自己都要受這些痛苦.
因為看著十萬人有多死.
這樣被鎮壓而死.
其實你都看到那種痛.
那個矛盾等等.
你可以想像得到.
社會事件的時候.
在早期的時候.
佔領中環的時候.
張牧師都叫那些青年的同工.
我是成年牧區.
那時候我剛來沒多久.
叫青年的牧者都要出去陪那些弟兄姊妹.
當時很和平的示威.
佔領中環的時候.
但後來輾轉又有一些比較激進的抗爭.
更加要留意那些年輕人情況怎麼樣.
你作為青年牧者.
作為當時的堂主任.
你都要面對這些問題.
你就要怎樣去處理.
你一定要擺在自己一個要超然的位置.
但是兩邊的人都因為自己的訴求這麼強烈.
和這麼對錯.
講中間話的人永遠都不被討好.
你去分析 你去持平.

$^{1321}$只會覺得你沒有立場.
不肯去體恤那些困難.
完全置身於事外.
這樣的情況是沒有辦法.
但是我很感恩.
馬丁路德真的仍然站在自己神學的牧者.
他將那個真理去釐清.
至於怎樣去審判.
因為他認為有兩個角度.
這個加爾文都是一樣的.
奧古斯丁也是一樣.
有神的角度 有地上的角度.
地上的角度有他的法律體系.
有他的不完善.
但是仍然是為著一個社會的和平.
可以提出不同的意見去做一個改善.
但就不是用一種激進的改革.
所以這個也是我自己的一路的信念.
這樣來在教會或者政治的.
都有人問我 當然不是昨天.
昨天知道了.
一兩星期前有人問我.
你支持共和黨還是支持民主黨.
我只說了一句.
兩個我都不知道哪個好.
有一個不是美國人.
所以很難做一個賣身.
但是其實由2019年一直到我們疫情之前.
那些賣身賣到的地步.
是逼我表態的人多的時.
一不是在一個立場表態.
就立即被人譴責.
無論WhatsApp 當面說你.
你這樣偏幫一定是你這樣.
什麼都關我事.
入我數.
不要緊的 既然是一個.
在這樣的社會的形態.
你也看到有些結構性的不公.
但也看到有些手段不合適聖經.

$^{1361}$你只能夠將這件事盡量用一種愛.
一種提出問題的結構性去改變.
但不是人人都願意這樣立即去做一個更新.
所以重要的就是我們有機會去看看.
我今晚嘗試找到那些翻譯.
但不是很多 因為中文不多.
那些翻譯和平分界的馬丁路德那篇文章.
還有另外一篇對於譴責的暴力.
他都有提出他的意見.
我嘗試回去找一找.
和平分界是有的.
所以我們今天就要簡單做一點.
這一方面的總結.
三條路線.
馬丁路德始終站在.
他只是扮演著一個神學家.
一個傳道人.
一個對真理要走回在聖經裡面.
他由始至終只是強調.
人的真正公平的身份.
是你自己看到在神的救恩裡面.
你所獲得的客觀的意義.
就是你自己主動的憑信心去領受神的恩典.
在馬丁路德一句名言.
無論你是一個黑丐或者你是一個貴族.
你在神的眼中.
你都可以守聖餐.
你都可以得到救恩.
這個已經是將公平發揮到最極致.
但其實沒有影響到社會結構.
因為那個說自己是黑丐的.
永遠都是黑丐的.
那個貴族仍然享受著地上的富裕.
這個就變成說了等於沒有說.
但其實不是的.
當如果兩個人去到神的面前.
他如果領受馬丁路德的教導.
他仍然是一個黑丐.
但他在神的眼中.
在神的救恩.

$^{1401}$在神的義裡面.
他和那個貴族一樣的身份.
一樣在神眼中的寶貴的兒女.
這個就是神學的重要地方.
就是馬丁路德.
始終他自己能夠站穩.
不會偏去訴諸於人民的主義.
對於公平 自由 個人的身份.
要奪取回.
不要被教廷再操控.
當時就是有這樣的思潮.
所以有些人就吃著這個宗教改革.
就發動了戰爭.
就說要激進的重奪.
屬於我們自己的利益 財產.
或者再一些抽象的就是自由.
就是我們自己可以當家作主.
這些一直發展.
其實你就不難發現.
世界的潮流就是這樣走向.
慢慢就去到很多國家出現的問題.
所以我們都求神幫助.
如果我們在這個角度去看.
馬丁路德的這個宗教改革.
平衡的原來有些人都是吃著改革.
不過他的改革原來是為自己地上的身份.
來捍衛革命.
本來是很高情操的.
是有一個意識發現自己的存在原來這麼重要.
但是慢慢不知不覺就離開了.
神的救恩的恩訊清兒.
慢慢就對於更重要的永恆的.
在神眼中的身份就不那麼在意了.
因為地上是實在 更加現實.
永恆就是我死了才有的.
所以說來都太遙遠了.
地上現在的壓迫.
所以你都不難想像.
在南美洲 中美洲.
繼續都有革命神學.

$^{1441}$真的要推翻國家.
你就可以想像到.
如果我們一擺錯一點點那個位置.
再加上一些領袖的.
一些他的所謂的角度的神學.
這樣來去演繹.
就會帶著一班人走向一種極端化.
所以我們都很需要求神給我們去反思.
所以馬丁路德對地主對貴族的譴責.
其實是給他們一個很大的反思.
其實那些農民不是你的敵人.
你們是神的敵人.
因為你們離開了聖經.
你們去壓制那些農民.
如果那些貴族有一個反思.
有一個行善的心.
有的可以去分給那些窮人.
其實不會導致一種完全的決裂.
要用流血收場.
正正就是多年的壓迫.
導致民不聊生的對抗 革命.
才會出現這樣的鎮壓.
如果他們有一個這樣的反思.
常常有一個這樣的反思.
自己都是在神眼中.
是一個黑蓋.
其實可能更加可憐的.
如果沒有救恩.
這樣的心態.
其實某程度上是很強的一種自我的意識.
對自我的了解.
待人接物 做人處事.
就會朝向聖經的教導去做.
就會帶來至少不會再用一種敵視的眼光.
覺得這些人是搞事者.
是阻礙社會發展的一些廢物.
而自己是有責任將神給你的那些祝福.
那些物資可以分配給有需要的窮人.
這樣帶來.
如果這樣化整為零.

$^{1481}$整個國家就會健康地兩下合二為一.
現在很多時候那種貧富懸殊.
大家都捍衛自己自身的利益.
所以你可以說.
馬丁路德其實有提到.
兩方面都是需要悔改.
兩方面 兩邊的陣營.
都是需要在神的真理裡面.
來復歸於一.
所以這個和平的訓誡這篇文章.
當時是很識切對那個時代的一個模樣.
我今天.
主要集中就是想和大家.
看這個第三條路線的改革.
希望大家都多了一個角度去看宗教改革.
原來不是這麼簡單.
不只是離開了教廷.
原來有些人在民間.
那些德國人開始有一個很強烈的自我意識之後.
一膨脹.
一受到一些人的鼓動.
就會帶來革命.
一定是流血收場的.
因為大家都不容許對方的對話.
只會帶來一個激進的暴力.
所以馬丁路德一生在這一點上.
真的被人這樣詬病.
但某程度上又真的與他無關.
他一直都處於神學的改革上.
希望尋求教廷的一個更新和悔改.
所以給大家.
可能平時你們沒有什麼機會去聽.
這些在神學院裡面教的歷史課題.
有一個這樣的概念.
但是擺在我們現在的社會.
我們都會看到各種的不公平.
昨天說有倫爾的不公平對待.
又說椅子不遇他.
這兩天經常聽到明仔說.
早上在大堂坐車出元朗.

$^{1521}$其實很繁忙的.
每個人都趕上班.
車內幾個站都已經滿了.
要開加班車.
他說每天最繁忙的七點多的時間.
都有倫爾的乘客.
全車就慘了.
因為要叫人們退後.
甚至下車.
然後開地板讓他上.
我不是歧視.
我只是將這件事描述.
要讓那個輪椅慢慢上了.
你那裡浪費了多少時間.
每個人都趕著九點上班.
住在元朗.
不是在旺角.
坐車去西鐵站.
再轉西鐵出九龍.
可能那個坐輪椅的都要上班.
是他的權利.
是沒有問題.
不過對於很多坐巴士.
擠巴士的人.
那個就帶來.
在他們的心裡面就是一個滋擾.
就是一個很麻煩.
司機就沒有所謂.
我都是開車.
我就停下車.
叫乘客麻煩退後.
退不到.
擠到全都擠滿了.
就請車頭的下車.
然後我就開門的地板.
讓輪椅客慢慢上車.
接著那些人上車.
那裡又阻礙了一段時間.
所以有時那個到底是歧視.
還是公平.

$^{1561}$怎樣去處理.
我都沒有辦法.
我都沒有一個答案.
因為當你在真正的處境裡面.
你就知道其實很多很多的事情.
我們真的要用在神裡面的愛.
包容去接納的.
彼此都不要站在自己.
那個很強的自己的立場.
勢不兩立.
這樣只會帶來更加激化了那件事.
所以有時怎樣去更好的處理.
這些都是很高的藝術.
所以透過今天我們講的.
宗教改革三個路線.
馬丁路德的路線.
下星期我們就會講.
這個伊拉斯姆的人民主義者的路線.
還有今天我們講的.
集中講的就是.
這個文次以帶動的激進的革命.
流血收場的.
這樣的十二條款的訴求.
表面上就好像很簡單.
但是帶來的就是.
那些地主覺得.
自己的利益受到最大的侵害.
就只會.
當然是假借馬丁路德的授權.
鎮壓這些農民.
就真的有很多的死亡.
這個是德國歷史裡面一個很大的污點.
宗教改革帶來的污點.
所以求主都讓我們很實在去處身處地.
實際設身處地去想一下.
在我們自己面對的每日裡面.
有沒有一些結構性的不公平呢.
有沒有一些事情.
我們可以繼續在我們自己內心的積蓄的善.
來去成惡呢.

$^{1601}$這個大原則我們就要去求主讓我們學習.
今天我就短一點說到這裡.
我們一起祈禱.
感恩天父讓我們今天都可以將這個農民戰爭.
講得清楚.
來龍去脈.
還有那個處理的手法.
我們看到這些都是很可惜.
其實每一個時代都有這些戰爭.
都是我們不能夠理解.
甚至有些可以避免.
但是我們想到當事人.
總是會用自己的角度去捍衛和去看那些事情.
導致那種勢不兩立帶來的就是聖靈塗炭.
其實每一個時代我們都會見到.
在今天我們見到國家與國家之間.
從來都沒有停止過這些戰爭.
好像以巴或者是俄烏.
或者是北韓的南韓的.
這些內部的內戰.
或者是長期的對峙.
也都有些非洲的部落.
很多的內戰的問題.
貪污的問題等等.
天父我們很盼望.
我們在真理的道路上面去追求.
雖然我們的生活是相當簡單和平順.
但是我們很盼望主給我們成為一個和平使者.
我們為很多事情去禱告.
也給我們看到我們是生在香港一個這麼有福的地方.
我們真的從來都沒有經歷過戰爭.
或者是嚴重的大規模的結構性的一些剝削.
但是我們懇求你.
給我們既然成為神的兒女.
我們就有責任去將主的愛 真理去傳遞.
去使更多人去認識.
以致我們就可以將人帶進神的救贖恩典裡面.
因為唯有在神裡面.
我們才得到真正我們自己那個存在的身份.
在神眼中我們是一個寶貴的兒女這個身份.

$^{1641}$願你幫助我們.
多謝你讓我們大家可以這樣去學習.
都是一個趣味性的對宗教歷史.
對宗教改革.
對馬丁路德.
對遲些我們講的嘉爾文.
或者慈雲尼等等.
去有一點點的了解.
幫助我們能夠在我們的生活.
我們最終都要問一個問題.
宗教改革的過去和將來.
我們應該怎樣走向願神來去施恩.
再一次多謝主讓我們一起學習聖經的真理.
奉耶穌基督的名.
阿們.
\newpage



\section{}
\label{sec:jvEmZbSGJ_A}
\textbf{2024-11-14 周四晚查經班《改革之路(4):伊拉斯姆》 講員:曾裔貴牧師}
\newline
\newline
連結: \href{https://youtube.com/watch?v=jvEmZbSGJ_A}{\texttt{https://youtube.com/watch?v=jvEmZbSGJ\_A}} ~~~~ 語音日期: 2024-11-15
\newline
\newline
\hyperref[sec:ziJGxjL2HKY]{\small{< < < PREV SERMON < < <}}
~
\hyperref[sec:index_chronic]{\small{[返順時目]}}
~
\hyperref[sec:index_scriptual]{\small{[返順卷目]}}
~
\hyperref[sec:DzI9N0JOL0s]{\small{> > > NEXT SERMON > > >}}
\newline
\newline
$^{1}$和大家一起祈禱.
天父愛尼,雖然天氣不是那麼穩定.
一起在不同地方聚集.
也可能有些人在稍後透過一些紀錄.
YouTube的一些紀錄來聽信息.
我們很盼望能夠透過這樣來保存.
可以使得更多弟兄姊妹.
甚至來自不同國家的弟兄姊妹.
能夠聽到一些簡單的宗教改革的信息.
可能引發大家有不同的思想或者言讀.
使這個運動不是一個已經過去了的.
而是我們不斷可以在信仰上有那個屬靈的更新.
因為所謂的宗教改革.
其實就是生命的更新的運動.
一個心靈的重整.
不過是一個整個當時歐洲的宗教信仰.
生命的重整.
所以也盼望神給我們.
既然已經來到我們中國人的社區.
來到香港.
我們華人有我們很獨特的傳統.
盼望我們去掉那些跟傳統不是最直接的關係.
而是信仰的核心.
我們能夠把持.
尤其是當初最重要的恩信,清義.
這個客觀,被動的義.
是神主動將救恩藉著愛子賜給我們的.
給我們每一個時代.
每一個信的人牢牢地抓住.
我們真的賴以得救的.
都不是我們的功勞.
不是我們的好行為.
是神無條件的救恩.
幫助我們能夠明白.
我們今天晚上會說的人物.
都給我們能夠對他有一些的認識.
可能其實都不是可能.
是已經淹沒在歷史的長河.
但是多謝主給我們有一些的書籍.
尤其是宗教歷史的書籍.

$^{41}$我們可以透過當時代的人.
對他們的描述,記錄.
我們後世的人能夠知道.
當然有更多的基督徒.
更多的可能沒有那麼出名的.
已經是我們不知道的了.
因為沒有記錄到.
幫助我們能夠認識.
也給我們有思想,有反省.
我們怎樣能夠留給後世有那個記憶.
有那個我們的擺上為神國度的福音.
因為每一個時代都是那麼短暫.
我們的人生都是那麼有限.
怎樣留給後世.
真是值得我們大家一起去思想的.
我們這樣仰望,禱告.
奉耶穌基督的名.
阿們.
相當重要就是.
今晚提的第四講的.
這位宗教改革的早期.
早期的改革家.
叫做伊拉斯姆.
他是一個荷蘭的六特丹人.
他終其一生.
只是70歲壽命.
70歲.
但是他做的事.
他的著作.
他帶來的復興的種子.
如果用一些雞蛋和雞.
他就是孵化.
好像生了一隻蛋.
那個蛋就是改革.
但是真正使那個蛋能夠孵化.
就是馬丁路德.
大家可以想像到.
他和馬丁路德的關係.
是相當的亦師亦友.
但後來就變成敵人.

$^{81}$亦都是他越來越遠離宗教改革.
是挺特別的地方.
但當然.
對這個人絕對不是一個壞人.
來攻擊馬丁路德.
或者當時的改革.
而是他採取的.
他自己的路線.
和他的看法.
特別是一本著作.
這個自由意志的部分.
馬丁路德就馬上寫了一篇文章.
去反駁.
差不多是他的老師.
就是這個伊拉斯姆.
可想而知.
那個最大的爭論點.
就是在那個恩順清二.
和那個自由意志衍生的恩順清二.
他簡單如果說這個伊拉斯姆.
他最重要的神學.
關乎到和信仰有關的.
就是那個合作.
就是和神合作.
就不是神自己唯一給人的恩典.
這個就是馬丁路德.
看為很嚴重的錯誤.
和他有很多的書信的挑戰.
好像是不戰.
所以後來都不能夠有那個彼此的和好.
當然不是說成為仇人.
而是不再有往來.
還有他就已經在後期.
宗教改革他已經是.
未去到禁止.
未去到極大的反對.
但是都已經是絕對不支持.
所以這個就是一個很簡單的縮影.
對於伊拉斯姆.
我們可以去朝這個角度.

$^{121}$因為很簡單.
我們說馬丁路德.
宗教改革.
其實環繞著三個人.
當然還有很多其他重要的人物.
但是都不要太多.
比較重要的就是馬丁路德.
慈雲里.
當然更重要的是嘉爾文.
因為嘉爾文確立的神學.
是更加的穩固.
奠定了宗教改革的福音派.
往後的那個進路.
就帶來17世紀那個普世的憤興運動.
普世憤興運動.
即大復興運動.
進入18,19世紀.
這個復興運動帶來的就是美國.
其實歐洲都有的.
英國.
美國是最主要的.
特別是愛德華茲.
帶來的就是.
由那個憤興運動帶來的.
就是那個普世宣教運動.
所以正正就是這個普世宣教運動.
才奠定了1900年在中國的大復興.
由義和團之後.
清朝末期.
一直到二,三十年代.
就是中國的基督教最蓬勃.
不要算秦朝.
景教那些.
那些不算的.
就算明朝.
利馬竇那些.
尼華那些天主教.
其實那些都不算的.
雖然他在上層.
和那些清朝的皇帝.

$^{161}$明朝的皇帝.
都有一定的相交.
甚至封了他做一些職位.
但是那些不可以叫做.
整個中國的大復興的一個福音運動.
是由20世紀的20年代開始.
中國就大復興了.
就有一班.
比如王明道.
倪拓星.
王治.
還有其他的一些重要的領袖.
輾轉後來就來到香港等等.
所以其實是幾個大的浪潮.
所以宗教改革.
是很重要的那個土壤.
去引發美洲和歐洲的大復興運動.
或者叫做奮興運動.
在英國就很重要.
約翰·衛斯理和他的哥哥查理·衛斯理.
但是我們再推翻一點.
我們第一次說的就是馬丁·路德.
他適逢的就是.
差不多有接近100年前.
其實輾轉都還有的.
因為伊拉斯姆的爸爸.
都是死於那些鼠疫的.
當時的西班牙的.
歐洲的那個很嚴重的世紀的疫症.
死了三分之一的歐洲人.
大家都束手無策.
都不知道怎麼辦.
甚至有些宗教人士說.
是神興起的懲罰.
對於歐洲的宗教沒落的懲罰.
那個黑死病.
也是14世紀.
都很嚴重的就是土耳其的崛起.
又是佔據了很多歐洲的地方.
大家今天知道的伊斯坦堡.

$^{201}$土耳其的首都.
或者當時的地方.
都是被土耳其人佔據了.
當然就是會引進伊斯蘭教.
所以這兩股勢力.
帶來歐洲的人都有很大的波動.
死亡是一件很實在的事.
如果看回中世紀的名畫.
都有這樣的對死神的描述.
歐洲人是多麼的懼怕死亡.
所以賣贖罪券.
變成一個有市場的銷售.
也是透過很多層層的渠道.
所以馬丁路德主要看到贖罪券帶來的負面影響.
賣完之後.
教廷原來只是為了錢.
但是普羅百姓.
又不是為了不犯罪.
是為了更加安心的去犯罪.
因為有贖罪券在手.
所以整個神學扭曲了.
所以在這樣的天時地利.
馬丁路德自己在他的成長.
也是非常之懼怕死亡.
他又承諾自己成為一個修士.
去回應他自己許下的承諾.
如果你這一次的閃電不死.
他就不再做律師.
而是做一個修士.
輾轉他就在威登堡學習.
後來有機會當時的主教派他去羅馬.
一方面是朝聖.
一方面在教廷裡面給他認識等等.
去完回來他的心靈就更加忐忑.
因為本來想著是一個潔淨之旅.
誠信之旅.
但見到的是很多扭曲了的教義.
和帶來的行為問題.
所以整個土壤必須要改革.
問題在這裡.

$^{241}$改革不只是一個人.
也有不同的人有不同的角度.
上星期我們說過一些田農.
種地的農民.
始終到了高壓手段的壓迫.
那些田主就徵他們的重稅.
又要很多不同的考納名目.
使那些農民終於要去真正的起義.
甚至乎要奪權取回自己的土地等等.
我們上星期說的21條.
終於去到一個爆發點.
但問題是到底宗教改革.
是否帶來他們有這種思想呢?.
原來教廷都可以去推翻的.
教廷都可以挑戰的.
甚至是馬丁路德帶頭去挑戰.
所以這個農民戰爭.
真是帶來.
也都是假借宗教改革之名.
實際上是一個暴力的抗爭.
所以就帶來當時的聖靈塗炭.
差不多有多少人這樣被殺害呢?.
一方面馬丁路德.
上星期我們說過.
他是禁止和絕對不贊成用暴力.
所以在這樣的情況下.
最終農民的戰爭.
是帶給馬丁路德都是有一些負評.
因為居然不支持那些農民的基本訴求.
但其實不是的.
馬丁路德是贊成他們提出的那些不公義.
提出要取回一些他們的尊嚴.
他們自己的一些生計等等.
但是用這麼激烈的手段.
真是一個差不多沒有彎轉.
讓當時的一些諸侯.
用強力的更大的武力.
去鎮壓這些農民的戰爭.
所以在宗教改革的時候.
有十多年的戰爭.

$^{281}$所以在德國.
這個都是帶來一個挺負面的對於宗教改革.
但是同一時間.
今天我們要說的這個人物.
就剛剛相反了.
他是更早看到教廷的問題.
但是他採取的手段.
和那個改革的路線.
比馬丁路德還要溫和.
他只能夠在教內.
透過教育去改變人.
就不可以對教廷.
對教皇帶來的侮辱和挑戰.
所以這個也可以這麼說.
是第三條的路線.
比馬丁路德的那個.
某程度上是勢成水火.
其實是教皇說.
你如果不收回這個95條.
你就要被逐出教廷.
就是教會以外.
就是離開天主教.
離開天主教你就沒有任避.
你就沒有Blessing.
你就沒有救恩.
這個是當時的神學.
馬丁路德就說.
我是迫於無奈的.
我是沒有辦法收回的.
除非你們用聖經告訴我我有什麼錯.
否則我會為自己所做的負責任.
求主幫助.
他就沒有收回95條.
也沒有收回他其他的著作.
所以自此和天主教1525年的.
這個沃木斯會議.
從此就正正式式的決裂.
馬丁路德就被看為是異教徒.
就被逐出天主教.
他也帶著那些信眾離開天主教.

$^{321}$當然是有一些諸侯收留他.
這個佛德列的一個德國的諸侯.
保護支持和收留馬丁路德.
餘下的時間就是教學和著作等等.
翻譯等等.
所以這個今天我們要講的人物.
在這個伊拉斯姆的眼中.
馬丁路德是帶來教皇的顏面盡失.
因為等於你挑戰他的權威.
等於你可以在他的頭的冠冕上.
去壓著他的冠冕.
和令到很多的教士.
主教教宗以下的一些修士.
我不知道香港怎樣分.
我其實研究一下是好的.
那個層級要幾多年才可以升到做師德.
升到做主教.
接著可以去到成為紅衣主教.
紅衣主教才可以選教宗.
香港天主教書記.
就是現在陳日君.
我忘記了他的名字.
他們去到成為紅衣主教.
他們就可以有權去到羅馬的梵蒂岡.
去選立教宗.
當然是要現任的教宗.
要不就是離任.
通常是死了才會沒有.
他們就會開會,祈禱.
如果我沒有記錯.
就是當你看到那些白煙出來.
就是選到教宗的了.
我先回去查一查.
我不敢肯定是白煙還是其它的煙.
這個是他們選教宗的情況.
所以大家如果熟悉的.
不妨留言在群組.
告訴我天主教他們怎樣升天的階梯.
我們回來說說伊拉斯姆.
是第三種的改革.

$^{361}$絕對不是用武力.
絕對不是像馬丁路德那樣.
馬丁路德那個很重要的立場.
他的立場就是.
除非你在聖經的真理找到我有什麼錯.
不然我是不會收回我對教廷提出的指控的.
求主幫助我.
這是馬丁路德的名言.
但是伊拉斯姆就是第三種.
他就放下了聖經的真理.
他可以有一點妥協.
千萬不要帶來分裂.
要在教會內部進行教育.
所以當時稱為人民主義者.
其實就是一些現代教育家的前身.
所以伊拉斯姆在另一個角度.
你可以說他是一個宗教教育家.
其實更加是現代教育之父.
他出了很多兒童的教科書.
教那些兒童的.
所以他的成就是相當大.
但是因為我們今天是用宗教改革這個大課題.
去看這些當時代的人物.
他就妥協了對教廷違反真理的譴責.
他最早當然不只是他.
再早一些就是波希米亞的約翰胡斯.
當時翻譯聖經的人.
比如丁道爾.
或者有其他一些出名的歐洲的人.
這些宗教改革家.
馬丁路德都還沒有出現.
就已經去指斥教廷的一些謬誤,錯誤.
但是得到的就是.
因為當時是有宗教裁判的.
所以約翰胡斯是比較早期被行火刑.
來燒死的.
所以這個就是你可以想像得到.
要麼就是很清晰的指出教廷的錯誤.
要麼就是要承受帶來的後果.
所以雖然馬丁路德一生沒有面對真正的死亡威脅.

$^{401}$但是他要遠離威登堡.
他要在不同的地方.
德國來逃避那個追捕.
因為教皇都立了那個追緝令.
要抓拿馬丁路德.
但是伊拉斯姆他絕對不是說.
好像叫做兩邊做人那樣.
他也不是說想保留自己.
怕時而是他的意念.
這個就看大家怎麼理解.
他覺得就算他是錯.
但也要在裡面做改革.
等一下我們有一個PowerPoint.
就很清楚反映到他的內心.
他的主張和策略.
所以我們就當是一個趣味性.
對這些可能大家聽都沒聽過的人.
但其實在那個時代是相當出名的.
我一邊研讀他的時候.
找到一些資料.
現代的荷蘭人.
他們所認識的其中五個歷史人物.
最重要的.
伊拉斯姆是佔據頭五位.
在荷蘭人的內心.
他是非常出名的一個歷史名人.
所以他在歷史裡面.
留過非常重要的人生的註腳.
譬如他是第一個翻譯拉丁文的聖經.
叫做《武加大》的譯本.
將他翻譯成為新約的希臘文聖經.
所以馬丁路德.
或者一些當時翻譯聖經的人.
都會參考伊拉斯姆這一本新約希臘文聖經.
他也做了很多其他的教材.
讓人可以作為參考.
雖然他是70歲的人生.
但他是一個博學.
和周遊歐洲列國的一位很出色的.
又是演說家.

$^{441}$又是教育家.
主要的教育是他相當多的參與.
和很大的人生都是以教育作為他的服事.
很重要的是.
他其中就是說.
他要在教內.
所謂教就是天主教.
他不會參與基督教的改革.
因為他已經不贊成.
後期他已經不贊成馬丁路德脫離天主教.
所以說教內的改革.
他認為要在天主教裡面.
開設平信徒的主日學.
只要那個文字可以提升.
整體一路來有成長.
自然就會有改變.
所以這個救恩合作論.
就是他和馬丁路德最大的分野.
馬丁路德認為.
人是不可以在救恩.
信主這件事.
有自己的參與.
自己是無能為力的.
所以他和伊拉斯姆說的自由意志.
剛剛有一個對立.
所以我們可以逐個部分.
和大家聊聊.
因為說到馬丁路德.
你不能不提伊拉斯姆.
是很早期的好朋友.
伊拉斯姆制定的一些想法.
和他的成就.
很多的書和翻譯.
聖經的翻譯.
是對宗教改革帶來非常大的直接影響.
因為如果不是伊拉斯姆寫了這麼多著作.
流通在整個歐洲.
就已經借助印刷術.
當時馬丁路德都還沒開始改革.
都還沒正式的.

$^{481}$他自己在後期.
做一個翻譯者.
翻譯不同的聖經.
來成為德文.
伊拉斯姆的著作已經印刷了很多版.
借助了當時印刷術在歐洲的流通.
所以神都用這些社會和科技.
說到伊拉斯姆.
必須要說一件事.
因為十四世紀.
剛才我們說了.
黑死病令很多人懼怕.
但同一時間.
其實在歐洲湧現的.
就是文藝復興.
有可能對於教廷.
一直封閉文字.
封閉人可以自己讀聖經.
不提倡讀經.
不提倡要識字的一個反動.
所以文藝復興.
很多有識有學問的人.
希望回到古代著作的原典.
所以原典的研讀.
非常盛行.
所以他們回到聖經的原典.
意味著那個只看為唯一權威的拉丁文的武家大譯本.
這個耶尤米.
主後的第三第四世紀的一個教父.
他的通俗的拉丁文聖經.
成為權威比希伯來文更厲害.
更厲害的比希臘文.
大家不看希伯來文和希臘文.
只看拉丁文的聖經.
但那本聖經是很異譯的聖經.
裡面其實有些錯誤.
但天主教就不理那麼多.
羅馬是用拉丁文做官方語言.
用拉丁文用了差不多一千年.
大家不認識希伯來文和希臘文.

$^{521}$當然有些很厲害的人一定會認識.
但那時根本不會翻譯其他的聖經.
自己國家語言的聖經.
所以德文聖經是馬丁路德開始翻譯的.
大家可以想像到一直以來.
由教父年代主後的三百多年.
一直到1517年之前.
都不會有大量的翻譯聖經.
因為真的要人手,抄寫,不會有印刷術.
就變成聖經永遠停留在拉丁文的世界.
不是每個人都能讀得明白.
原來讀得明白的教義背後.
都只是有些人滲入了意思去解讀聖經.
所以就帶來中世紀不是每個人都識字.
更加不會是學者.
但神又很奇妙地興起了一小撮的人.
一些可能是當時大學的教授.
或者是其他的聖經學者.
但伊拉斯穆就是其中一個很重要的.
當時的殿堂級的人物.
他本身其實是一個私生子.
他的父親都好像是教廷裡面的一些修士.
他的母親似乎是一個醫生的女兒.
有些報道這樣說.
但是他的父親很早就死了.
剛才我講過了.
有些記載說他的父親是黑死病死的.
所以其實很靠近的.
在歐洲人來說死亡是很恐怖的.
因為是絕症來的.
是沒得醫的.
所以他們對死亡是一點都不會陌生的.
不會覺得很遙遠的.
所以說要買贖罪券是很有市場的.
人們覺得搞笑是不會的.
因為每個人都怕死.
死了之後在煉獄裡怎麼辦呢?.
當時的歐洲人是很擔心的.
但是我們回來講這個伊拉斯穆.
他是一個人民主義者的意思.

$^{561}$就是因為他也是很聰明.
他有機會在某個程度上叫教會學校.
但是當時所謂教會學校就不是今天的Marylock.
不是今天的呂拔萃.
不是今天的其他那些.
或者那個宗教味道已經很淡了.
只是學術高.
但是他是在修院裡讀的教導.
一定是相當嚴格.
他也是一個很出色的教育家.
他的學問是相當高的.
高到一個地步.
他能夠自己去研讀希伯來文,希臘文.
好像是差不多20歲左右.
他已經是能夠精通的.
你可以想像得到.
這樣的人可以有那種學問,學養.
後來他因為受到文藝復興.
復興就是回頭.
回到最good old days的時候.
研究那些原典.
一方面不僅僅限於聖經.
文藝復興就是古代有名的著作.
所以對奧古斯丁的熱潮.
就是在中世紀這個時代再出現.
所以難怪馬丁路德對於自由意志.
馬丁路德也是採納奧古斯丁的定調.
大約的教義真理.
和伊拉斯姆是有很明顯的不同.
所以這種回到古代.
回到聖經成為了當時的熱潮.
就很寶貴了.
因為正正就是一種物極必反.
正正就是當時拉丁文的通俗的武家大譯本.
是帶來整個中世紀對賢經的興趣完全缺乏.
沒有人羨慕或者認識.
學無有這些賢習班,查經.
是沒有的.
只是你去到聽當時的.
可能是主教.

$^{601}$如果是大一點的教堂.
但是有些小的就是修士.
就會讀拉丁文的聖經給你聽.
接著就不解釋.
所以這樣的彌撒.
如同你遠遠地望.
是不會令到你的智力,智慧,靈性.
有一個跟神的接觸.
這個是問題所在.
所以伊拉斯穆就在這個熱潮下.
他就回到賢讀聖經.
他也主張.
這個就是他開始建構他自己的人生.
因此他覺得.
也是建構他的神學.
他認為神給人的自由意志.
是做什麼都可以的.
什麼都可以自己參與下去.
所以是神人合作論.
就是這樣在神學上.
他跟馬丁路德的直接衝突.
因為他寫這個自由意志.
馬丁路德立即寫了一篇文章.
反對駁斥伊拉斯穆的自由意志.
挺有趣的.
如果我們在那個時代出現.
見到兩位名人這樣對話.
其實如果是真理是對的.
不是不對的.
但是我們等一下有幾個重點.
看看怎麼處理.
剛才我說的文藝復興.
帶來當時的文人.
就是有色之士.
不一定是信耶穌.
不過伊拉斯穆是敬虔的成長的.
在教廷裡面.
後來他是做過修士的.
但是當他一路周遊列國.
也看到教廷很多問題.

$^{641}$他申請還俗.
後來以一個人民主義的教育家.
成為他人生的志業.
當然他也寫了很多書.
寫了很多書.
他的書是很受歡迎的.
其中一本叫基督徵兵手冊.
如果我們今天能夠找到.
我這兩天不斷看他的資料.
我有興趣這本基督徵兵手冊.
等一下告訴你.
正正就是他說.
怎麼帶來復興.
但是又不會帶來決裂.
就是需要訓練基督徵兵.
基督徵兵就是所有的平信徒.
要對神的話.
真理有認識.
在生活裡面表現出基督徵兵的品格.
你就可以改變教廷裡面的腐敗.
完全是人為性的參與.
不需要你行禮儀.
甚至嚴重到你不需要回教會.
進去教廷搜彌撒.
你都可以做好.
當然其中一點是很對的.
就是回到讀聖經.
他的主張是影響了馬丁路德.
影響了加爾文.
影響了當時就算溫和激進的.
都認為聖經的翻譯.
是必須去到不同的語言.
這是完全沒有意義的.
起初引發這一點.
就是伊拉斯姆他自己的洞見.
他自己看到根源.
當然他是從教育入手.
所以某程度上一直推到今天.
就是新派去神跡化.
去其他的東西.

$^{681}$只需要你做好人就行了.
信仰是帶給你做好人.
你不需要有對耶穌的神跡.
對這些很難接受的超自然.
你不接受是沒有問題的.
只需要因為聖經.
你做好人就行了.
這就是新派的核心.
但是福音派就不是.
我們沒有人可以做好人.
我們是需要救恩.
改變我們的生命.
在這個新生命裡面.
去活出一個神期望的新生活.
新生命.
所以那個不是合作的真理.
是一個唯獨恩典的真理.
所以唯獨和合作是勢不兩立的.
但是伊拉斯姆就認為是一個合作論.
他強調恩典的.
但是人需要做好.
那就是神給人的恩典.
所以剛才提到這幾個元素.
文藝復興帶來的對古典的研讀.
那個興趣.
而且那個興趣更重要是回到聖經.
因為都提過古代的文獻.
重要的文獻.
相對於聖經出土的文獻.
無論是手抄本.
是很懸殊的比例.
所以文藝復興.
其實都是回到最重要的研讀.
就是聖經.
因為大量流傳的手抄本.
是很懸殊的跟其他普通的作品.
所以這個人民主義的教育.
其實他們得到那個知識之後.
就是要用知識去改變那個時代.
那個時代就是宗教改革.

$^{721}$就是教廷的腐敗的時代.
伊拉斯姆不是不見了.
他見到的.
我給大家看了一個PowerPoint.
這個就是剛才說的基督徵兵.
有一個作品是挺有趣的.
而且是非常受歡迎的.
我很快給大家看看.
這個叫做讚美蠢材.
俞恆讚.
是用很簡單的文字.
來稱讚那些傻人.
其實原來是諷刺宗教改革的時候.
那些教士.
那些修士和教廷的主教.
用一些嬉笑怒罵.
但不低俗的.
不是很卑劣的.
不是的.
用一些很嬉笑怒罵.
但不低俗的.
這樣來諷刺.
居然是大受歡迎.
我試試找找.
看看香港圖書館有沒有這類翻譯的書.
大家可以再研讀一下.
都挺有趣的.
好像當時印了很多重印.
這個基督徵兵印了23版.
在歐洲.
全片成歐洲.
所以你可以想像到.
伊拉斯姆在他的教育理念.
或者在他的改革理念.
是要人平順圖.
裝備起來.
做好自己.
自然就會帶來改革.
他就是一種.
可以說是從下而上.

$^{761}$不挑戰教皇權威下.
來做溫和的改革.
當然其中他就是.
一定涉及到他的神學信念.
就是他認為自由意志.
人是可以自己掌控的.
人沒有用到的地步.
是一無是處的.
你可以決定你自己.
要做好的.
不需要講救恩.
你都可以做好的.
這個就開始有點危險了.
也激發了馬丁路德對他.
這個自由意志這篇文章的譴責.
因為這篇文章.
大家越走越遠.
根本就已經分道揚鑣.
他就不再在宗教改革上.
有任何的發言.
參與和合作.
對馬丁路德的評擊.
說他的作品.
都是很粗魯的野蠻.
這個是歷史給我們看到的.
當時的情況.
他很早就死了.
馬丁路德都還沒有死.
他已經離開世界.
他應該是146幾.
至到1536年.
當時都是宗教改革最蓬勃的時候.
這個就是伊拉斯姆.
這個樣子我想都是當時的情況.
可能古代的人.
他們沒有那麼多其他的娛樂.
他們整個人生就是在聖經研讀.
因為當時都沒有電燈.
只用油燈.
但是好像沒有影響到他們每天用的時間.

$^{801}$對於翻譯.
對於寫作.
對於到處巡迴的講道.
好像沒有太受影響.
為什麼他們會用我們現代人那麼多的時間.
可以有那麼大的成就呢?.
這個就是很難回答的問題.
我們只是稍微看一看手機.
WhatsApp已經搞了你半個小時.
你還是用一個玩偶來回覆.
都可能搞了你幾個字.
所以現代人很碎片化.
那個思想很難集中.
因為你一集中看東西.
那些WhatsApp就已經幾十個了.
有時有這樣的難處.
但是沒辦法 每個時代都不同.
但是我們回頭看這些先驅.
某程度上是他很早就已經洞見了教廷的腐敗.
但是他選取的路線.
可能受著他自己的成長.
受著他自己受的教育的影響.
他認為必須是溫和的.
所以他反對馬丁路德.
這樣挑戰教皇的權威.
所以你可以想像到.
都埋下了大家不和的種子.
今天我們要講的.
他是一個很重要的人民主義者.
是北歐荷蘭很出名的.
路特丹出生的這位伊拉斯姆.
當時的熱潮就是文藝改革.
回到文藝復興 回到古代的文學.
所以回到聖經.
這樣就對了 除非你最怕你不讀.
你讀神的話語就會改變你的.
所以到今天仍然無論香港.
我不知道現在歐美是怎樣.
香港的閻經運動還是強的.
每年的培靈會.

$^{841}$不同的教會都有自己的培靈會.
都是強的.
至少每個星期.
無論其他的節目.
應該這樣說.
無論其他的崇拜裡面的項目.
多重都好.
沒有一間教會是不講聖經的.
如果有間教會不講聖經.
他自己都生存不了.
任何一間教會崇拜.
必定會講聖經的.
當然可以是九個字.
或者有些可能是十五分鐘.
一般都很少十五分鐘的.
起碼25至到大概45分鐘.
如果一個小時的講聖經.
現在的人都不行了.
太長了.
在崇拜來說.
所以想像到.
由宗教改革.
回到聖經.
即是十四世紀.
來到現在21世紀.
在這五百多年.
聖經仍然佔據著.
信徒的心靈.
是相當重的地位.
所以這是一件很好的事.
當然一個運動.
慢慢成為了.
譬如必須要經過的.
感宣就比較明顯一點.
初信班,浸禮班.
個個都知道要堅持.
多難,多忙都要上.
但是浸完.
重擔就好像.
卸除了.

$^{881}$可以稍為鬆一鬆.
變成.
初信浸禮.
只不過是為了浸禮.
我才要上.
其實他不是真的愛慕聖經.
或者不知道.
原來回到聖經的重要.
這一點就已經肯定是.
宗教改革的過去和未來.
我們要去反思.
回到聖經.
因為今晚我們正正借著.
伊拉斯姆.
他本身就在當時代的文藝復興.
回到聖經.
他是受益人.
他很受到聖經的影響.
他真是有這樣的能力.
他自己學會了拉丁文之後.
懂得看武加大譯本.
然後再去對照.
新舊約聖經.
希臘文的原典.
他就發現.
原來武加大譯本.
拉丁文的通俗聖經.
耶尤米這位教父.
原來翻譯到真的有很多的錯誤.
甚至乎是很有意義的錯.
問題就在這裡.
天主教看武加大譯本是權威.
是唯一在教會裡面的譯本.
其他都不准看.
也不給你有機會看.
當時也不會有印聖經的時代.
所以你可以想像到.
天主教只用一個版本.
單一的武加大譯本.
作為他們的信仰權威.

$^{921}$就帶來很大的錯誤.
所以他慢慢被一些傳統.
一路牽著走.
越走越遠.
所以天主教就越來越離開.
聖經的教義.
加了很多人為性的東西.
所以.
一方面你可以說.
你要怪責耶尤米.
但是在他第三世紀.
第四世紀的時候.
你還想他怎樣呢?.
一個人能夠將新舊約聖經翻譯成拉丁文.
你還想他怎樣呢?.
什麼都沒有的年代.
沒有聖經沒有其他的參考書.
已經盡了他人生的力.
印了一本拉丁文的聖經.
因為羅馬當時是拉丁文的.
所以.
到了十四,十五世紀.
這個伊拉斯姆才發現了這個問題.
而他有能力.
將這個新舊約聖經.
來將它整理好之後.
就出了一本.
他自己的希臘文.
他自己的希臘文的聖經.
所以你可以想像到.
是帶來.
宗教改革.
一個很重要的.
對翻譯.
對聖經這兩個重點的重視.
所以.
用廣東話就不能輸了.
因為聖經.
被重視.
不單止被重視.

$^{961}$不是用來做裝飾.
而是要平信徒.
你要有自己文字的聖經.
所以德國的聖經.
自然是馬丁路德做的.
法國和瑞士.
自然是加爾文去翻譯聖經.
有其他的人去翻譯.
所以後來的英皇音定本.
在英國就被印出來.
所以你就可以想像到.
那個翻譯聖經就成為了.
宗教改革.
必然會出現的後果.
就是天主教要綁住那本聖經.
單一的拉丁文.
但是基督教離開了教廷.
他們什麼都不做.
就是做翻譯聖經.
就是去到不同的地方.
去講道.
所以由十五世紀.
一直到十七世紀.
十八世紀.
都是講道為基督教的.
主要傳道的職事.
如果我們有機會.
回到時空穿梭.
回到十七世紀的.
這個.
愛德華茲.
他一篇講道.
據文的歷史記載.
其實他照著講章讀.
他是一個大近視的傳道人.
他一邊讀著講章.
那篇講道就叫做.
當一個罪人在憤怒的神手中.
他在美國的講道.
就說當時大復興的時候.

$^{1001}$那些人要抓住教堂裡面的椅子.
怕掉下去的地獄火湖.
這樣的信主.
你可以想像到.
宗教改革.
就是透過這些教育家.
伊拉斯姆就很重視聖經.
他有能力去翻譯希臘文聖經.
成為了那個藍本.
讓其他人可以用來做一個參考.
根據希臘文聖經.
翻譯成.
首先就是德文.
就是馬丁路德的德文聖經.
所以他也是一個翻譯家.
更重要的是他是一個推銷者.
他將他自己的神學.
透過著作.
來寫.
剛才說的讚美那些愚蠢人.
其實就是挖苦那些教廷的領袖.
不過就用一種很通俗的手法.
個個都看得明白的.
這樣來寫.
就令到那個銷量很大.
當然他有很嚴肅的.
就是他要支持.
他覺得人是可以改變的.
這是教育理念.
其實現代的教育.
就是認為人透過教育增長知識.
是可以有改變的.
這個某程度是對的.
絕對是的.
不過我們現代人也都知道.
越聰明的人.
越有學問越聰明的人.
越容易被人欺騙.
越聰明的人.
越有學問越聰明的人.

$^{1041}$如果他裡面沒有神生命.
其實他壞的可以更嚴重.
就是那些犯罪聰明的人.
他的犯罪的腦袋.
是更加嚴重.
我想無論是犯罪心理學.
或者那些從事法律的人.
或者警務處的人.
都更加知道這些.
越聰明的人.
如果他裡面沒有神的生命.
其實那個知識.
不單止幫不了他.
反而更加很掩飾地犯罪.
所以教育是可以一部份改變到.
但這個就是.
伊拉斯姆自己.
他一個很強烈的信念.
所以很多人說他是現代教育之父.
這個伊拉斯姆.
他提倡的教育.
是非常之重要.
但很有趣的是.
他是一個修士.
去到一個地步.
差不多任命他成為當地的主教.
但因為他周遊列國.
又去英國.
又在北歐一帶.
他看到的東西.
是宗教不可以真正改變到人的.
要教育在人的裡面.
所以他那本基督徵兵手冊.
是相當暢銷.
他印了23版.
那個原因也可以這樣說.
所以最終.
當教皇想任命他成為當地的一個主教.
他自己當時是修士.
他的決心.

$^{1081}$不單不接任主教.
他要還俗.
不再做修士.
給修士的銜頭.
教廷.
正正式式.
來解除修士.
即是可以結婚.
簡單來說就是離開了教會.
但馬丁路德.
不是打算離開教會.
但被教皇驅逐走.
所以兩者是完全不同的.
所以你會看到.
他一生對教廷.
都是採取一種溫柔.
希望用教育改變平順圖.
不是去挑戰教皇的權威.
這一點和馬丁路德.
是有一個直接的不相同.
所以他很不贊成.
馬丁路德這樣以下犯上.
針對教皇來提出的95條.
但是.
他也不贊成.
教皇的這個主意.
他也不贊成.
馬丁路德提出的95條.
但馬丁路德提出的.
不是針對教皇.
是針對教義本身.
所以他不是挑戰那個人.
而是為真理去捍衛.
所以怎樣分呢.
大家真的要.
在你自己人生神秘.
你什麼崗位.
你要怎樣捍衛.
譬如作為牧者.
我們要捍衛的一定是真理.

$^{1121}$福音的完整性.
不是你個人的野心.
你個人的偏見.
不是應該捍衛這些.
是要捍衛真理本身.
這個就一定要注意.
所以絕對不應該有仇恨.
絕對不應該有一個偏合的心.
為反而反.
這個就已經是大錯特錯.
作為一個領導人.
絕對錯.
如果你自己的心偏合.
看什麼都是從你的格局去看.
你自己的心.
看什麼都是從你的格局去看.
你的教會就會很偏合.
就會很有問題.
所以牧者.
必須是為真理才發聲.
如果不是真理.
那些就沒有所謂.
如果是為個人就更加沒有所謂.
千萬不要為爭論而爭論.
或者說為自己的面子.
更加不值.
因為人都是有限的.
所以你可以說從這個角度.
伊拉斯牧.
沒有挑戰教廷.
他就能夠在教廷裡面.
維持一個好的關係.
他的理念就是.
希望平信徒.
所以他那本書很受歡迎.
希望平信徒.
得到生命成長.
自然自己又會有改變.
但是這個.
是不是真的可以呢.

$^{1161}$不接受基督.
就可以改變.
這是做好人的邏輯.
所以馬丁路德.
對他這個主張.
是極之.
當然我不要說反感.
是極之不贊成.
其實起初的時候.
兩個是好朋友.
而且伊拉斯牧.
某程度還是馬丁路德的前輩.
馬丁路德.
都用他的新約希臘文聖經.
作為他翻譯.
德文聖經的一個藍本.
其中一本.
你可以想像到.
那種關係是很深的.
但是當伊拉斯牧寫了這篇文章.
論自由意志.
馬丁路德看到之後.
立即就反駁.
錯的.
人是不會有絕對的自由意志.
在救恩上面.
根本我們是無助.
我們是完全一無是處.
對一個教育家來說.
這是一個很大的錯誤.
對一個教育家來說.
等於一盤冷水潑下去.
教育是完全沒有任何作用.
所以難怪伊拉斯牧.
有這種.
對馬丁路德的.
不是仇恨.
是很重的說話.
他的作品.
馬丁路德的作品.

$^{1201}$是野蠻的.
是粗俗的.
對一個人民主義者.
對一個這麼有教義的.
對一個這麼有教養.
這麼有學養.
希望透過教育.
可以改變那個時代.
和馬丁路德的一個徹底性的.
恩順清義.
在神的眼前.
沒有任何的異可言.
必須是一個.
你要主動.
獲得神的恩典.
你才可以得到這個救贖.
在馬丁路德來說.
是背道而馳.
伊拉斯牧.
所以大家有一點點趣味的了解.
那個時代.
一方面如果我們再高一點.
因為我們已經經過了500年.
我們高一點去看.
神都用不同的人.
正正反映了.
那件事未蓋棺定論.
大家一定有自己的性格.
有自己的成長.
受到的影響.
或者分逃.
採取的路線.
自然就會有些激進一點.
上個星期那些農民的戰爭.
簡直是用革命.
簡直是到處去燒教堂.
你去說教廷錯.
我們多走一步.
燒教堂.
掠奪那些地.

$^{1241}$自己田地去分.
你不肯接受我們的21條.
我們就已經是.
沒有轉圜轉了.
就像革命.
大家死了就發動戰爭.
那麼多人死了.
馬丁路德也禁止這些戰爭.
但是伊拉斯牧.
也不贊成.
馬丁路德的改教.
離開教廷.
所以有三條路線.
神都用不同的人.
在那個時代裡.
最終是否令到教廷.
有一個改變.
有的.
開了一個天特會議.
但是有些最基本.
最核心的.
都是不能夠變.
因為一變就不是天主教.
就是說教皇是無誤.
現在我們就宣佈教皇無誤.
那就不能再變成天主教了.
等於你變成.
立刻變成基督教.
所以天特會議.
只不過是更加確定他們那些東西.
不過是one time.
做一點點的改變.
所以也等於.
徹底的.
由1517年的.
一個.
里程碑的.
10月31號的萬聖節.
就奠定了.
天主教和基督教.

$^{1281}$政績慢慢的.
慢慢的決裂.
慢慢的.
各自去做回.
那個時代的歐洲人.
的心靈.
所以基督教.
當然就要採納兩個聖禮.
因為天主教有七個聖禮.
基督教就兩個.
禮和聖餐.
是最重要的.
所以.
其他就沒有了.
大家就是百花齊放.
你有什麼領受.
你就會去.
治理你的教會.
所以那個時候也都開始.
產生了不同的.
宗派開始的出現.
好的就是.
化整為零去到歐洲.
不同的地方.
當時也都是航海業.
和當時的.
拓展那些.
陸地的地方.
當時的航海.
和殖民.
所以當時的西班牙.
當時的.
荷蘭等等.
英國.
開拓新大陸.
英國就更多人.
去到新大陸就是美國.
和其他的地方.
所以在民間.
開始也都等於.

$^{1321}$帝國主義的.
早期的文雅.
當你發現那個地方.
文化的知識和你相差那麼遠.
你就去征服它.
大量的資源.
拿回自己的國家.
就開始蓬勃了.
十九世紀的.
普世的.
殖民主義或者叫帝國主義.
的問題.
所以也都有我們清朝末期的.
很多的.
八國聯軍等等的.
戰役等等.
你可以想像到.
都是源於基督教.
和天主教既然已經分開.
基督教就一定要.
自己開拓不同的領域.
向不同的人.
傳福音.
所以某程度上是神祝福他們.
透過這個分裂.
神都可以用這些基督教的.
不同的宗派去到不同的地方.
就開始有.
不同的宣教組織.
這個當然是.
十九世紀的末期.
才大量湧現.
那個時候還沒有.
最多是去.
環越大西洋.
去港島.
還沒有去到這麼大規模的.
普通的弟兄姊妹.
都參與那個宣教運動.
就是.

$^{1361}$漂洋過海的宣教運動.
當時還沒有的.
在十六世紀,十七世紀都沒有的.
有的是那些殖民主義.
所以這裡就.
透過伊拉斯姆.
給我們總結了.
他是某程度上.
是那個始作者.
創始人.
對聖經的.
研讀的濃厚興趣.
和他做好一些教材.
使到翻譯聖經.
成為容易.
也是他提倡的.
教育理念.
使得更多人.
要盡快.
認識聖經.
實踐聖經.
這個就是伊拉斯姆.
我們又看看另外一個PowerPoint.
我嘗試找一找這本書.
看看Amazon有沒有.
十五世紀的作品.
五百年前的作品.
現在還有沒有呢.
在當時影響著歐洲的天主教徒.
平信徒.
要讀經.
要生命成長.
教會自然就會改革.
就是這麼簡單的理念.
就真是.
希望可以透過教育.
去改變人的心.
但是呢.
其實如果夾雜著.
應該說夾在一起.

$^{1401}$馬丁路德的恩順清二.
來做一個改變.
相信那個效果會更好.
如果你將兩樣東西分開.
單是從只是教育改變一個人.
其實一般社會.
現在譬如美國大學生.
大把的.
歐洲英國有讀書的人.
難道每個人都是文盲嗎.
但是那些罪案.
是不會因為教育而減少的.
全世界都是.
所以如果單純只是教育.
其實就.
不會改變的.
當然這本基督徵兵.
是從聖經的角度.
入手的.
所以伊拉斯姆絕對不是反對.
讀聖經改變.
而是他覺得.
在改革裡面.
需要用溫和的手段.
需要合作的手法.
就是和教廷.
一個合作.
所以這本基督徵兵手冊.
是為那些天主教的評述圖寫的.
據聞.
印了23版.
不斷帶來更新.
其中就是1515年.
距離1517.
還有兩年.
這一版的第三版.
據文歷史記載.
這一版的更新.
每一次翻印.
就會有增幅.

$^{1441}$增加一些篇幅.
第三版.
據聞就是.
很受當時的影響.
很影響.
馬丁路德的改革.
所以.
不可以說伊拉斯姆.
他做這個舉動是錯的.
建立評述圖的心靈.
其實是一個相當重要的工作.
所以.
真的奉勸.
我們如果讀了很多聖經.
我們要找機會去用.
我們用.
我們自己就會先受益.
然後又幫到一些人.
甚至乎.
如果我們又不是.
有這麼忙碌的.
教會工作.
不是全時間.
教主一學.
其實不只是教.
而是你要成為他的mentor.
真的可能是幾個.
要看著他陪伴他成長.
這樣的.
運動.
其實那個成果.
是很大.
所以今天我們要提倡的.
就是一個門訓.
基督的門訓.
是很重要的.
門徒訓練.
這樣的生命影響生命.
是很重要的.
所以如果不是的話.

$^{1481}$我們就很難去建立下一代.
他們對神有一個很.
強烈的毀身的.
那種情操.
如果我們用一些市場學的分析.
其實任何一個的時代.
都是想找自己的骨幹分子.
強於.
共產黨的.
共青團.
少青團.
由小朋友就要讓他知道共產黨.
共產精神的偉大.
那個重要.
去建立他的心靈.
廣告也是.
都是要從小就要讓他知道.
什麼是好.
譬如大學生.
既然是未來可以花錢的.
就要在大學的校園.
無條件給他用卡.
還有很多禮物送的.
你肯要就行了.
你未畢業不要緊.
先給你.
你就享受一下.
讀書都可以簽帳.
其實他不知道簽完帳原來沒錢還.
一樣要給這麼多的利息.
你就知道那個後果了.
所以.
其實不同的時代.
不同的組織.
都是想找自己的骨幹.
所以基督徵兵搜查.
背後的精神是很好的.
是真的希望在平信徒當中.
做到那個.
他自己對聖經的認識.

$^{1521}$對神的認識.
他自己有一個強烈的.
要活著為耶穌.
這樣的心靈.
其實就化整為靈的.
這樣的改革.
某程度是好的.
不過馬丁路德都提出.
這樣的論點.
這樣的建議.
不過立即因為.
那個感覺是.
教皇有很大的問題.
那95條裡面.
是帶來.
教廷不能夠接受.
所以只能夠.
分道揚鑣的決裂.
所以.
伊拉斯姆他的想法.
是好的.
但最終其實又不真正.
改變到天主教.
反過來馬丁路德的.
這麼強烈的訴求.
帶來的分裂.
其實又湧現了.
基督教的成長.
即往後其實.
由15,16世紀.
一直到公元的.
2000年其實都是基督教.
的蓬勃的黃金時期.
天主教都沒有.
那個蓬勃.
所以你可以想像到.
如果以全教的事業.
回到聖經.
和盡快培養.
信主的人.

$^{1561}$有基督徵兵守冊的精神.
這個兩者結合.
就是.
最重要.
所以任何一個教會.
如果平信徒運動做得好.
門徒訓練做得好.
就會有很扎實的下一代.
去承擔福音的.
傳遞和建立.
所以.
這兩樣東西是伊拉斯姆和馬丁路德.
如果結合.
就很好.
但這個基督徵兵.
我仍然未看.
不知道有沒有英文版的翻譯.
因為我不懂拉丁文.
所以沒得看.
你可以想像到.
他是希望這樣的改革.
但這樣的改革.
其實根本就沒有觸碰到.
核心的問題.
就是馬丁路德觸碰到的.
核心問題.
就是恩信清義這個問題.
伊拉斯姆就不觸碰.
未必他不敢觸碰.
而是他覺得不應該這樣觸碰.
不應該用這種方法去觸碰.
他就認為.
他見到那個腐敗.
但他就認為.
是以這樣的教育.
去改變人心.
這個就是.
大家採取的.
第三條的改革路線.
若教會整體.

$^{1601}$這個就是基督徵兵手冊的重點.
大家上網搜索一下.
你都會大概看到.
有些中文的寫關於.
宗教歷史的都會有這些記載.
若教會整體.
回歸到教父.
和聖經著作.
因為他自己就是文藝復興.
要回到原典.
回到教父.
那就是中世紀不讀聖經.
也不准人讀聖經.
不准普通的信徒讀聖經.
所以其實大家是對神的話語一知半解.
對神的話語一知半解.
即是對神一知半解.
神父告訴你神是怎樣的.
那你就是這麼多.
所以叫你買贖罪券.
就可以贖罪.
那就買吧.
就沒有想那麼多.
所以就不可以用今天.
那麼多的研習班.
有培靈會.
有不同的傳道人都可以講道.
女傳道都可以講道.
以前是沒有這樣的.
你可以想像到.
那種的研讀.
是真的多元化.
真的各有各不同的領受.
你聽不明白.
誰在講.
但另一個你聽就聽得明白.
鄭牧師講得很清楚.
曾牧師呢.
我都不知道他在講什麼.
不要緊.

$^{1641}$每個牧者就有.
不同的演繹.
神的話語總是.
豐富在不同的人.
的領受.
所以這是好的.
基督徵兵手冊.
如果教會整體.
他就是想平信徒運動.
不要理那些教廷.
不要理那些腐敗.
整體我們回到教父年代.
回到聖經的著作.
回到讀聖經.
就可以帶來革新.
恆常閱讀主聖經的.
平信徒就可以帶來.
敬虔的那條鎖匙.
就是解開.
那些腐敗.
但這是.
相信是一個很強的.
個人心願.
但惰性.
罪性.
那個問題.
只會很大的.
影響.
但有一個很重要的.
一點就是.
由一路都不准讀.
到你可以讀.
這是一個零.
和一的突破.
其實早期.
剛剛耶穌的年代.
使徒的年代.
教父的年代.
是一個大逼迫時期.
大家對神的話語.

$^{1681}$對成為基督徒.
的寶貴是沒有意義的.
但自從君士坦丁.
成為基督教的.
開放之後.
人就不再覺得.
讀聖經和生命有一個掛勾.
所以.
如果有人說.
不用你讀了.
他又覺得挺好的.
變成.
起初都沒有什麼.
但慢慢慢慢.
就變成聖經不可以.
隨便給一般人讀.
只能用拉丁文.
讀.
變成聖經和.
普通信徒中間.
有一個中介人.
中介人.
兩邊都不會告訴你.
實際情況是怎樣.
令到平信徒.
蒙在鼓裡.
但聖經很想人認識.
就被中介人攔了.
所以中介的角色.
越來越沉重.
越來越變成.
一個枷鎖.
所以變成必須有.
改革.
你就可以想像到.
為什麼一定需要.
有這個變革.
但伊拉斯穆就是一個.
溫和的.
在平信徒當中.

$^{1721}$做的一個革新.
這個就是.
你可以叫他和.
馬丁路德.
分別的地方.
他就在.
這個武加大譯本.
拉丁文.
應該叫做Rogge.
武加大很出名的.
因為當時.
只有這一本聖經.
流傳.
因為當時是不會有人.
翻譯其他文字的聖經.
所以第一本.
翻譯聖經就是這一本.
武加大譯本.
把希伯來文和希臘文.
翻譯成.
拉丁文的聖經.
然後停了.
不再有其他國家的語言文字的翻譯.
因為大家都不看重翻譯.
所以到了.
15世紀,16世紀.
才正式有.
再一次的蓬貝的翻譯.
聖經的翻譯.
還有對於原典的追尋.
對於希臘文.
的研讀.
對於希伯來文的研讀.
所以這個也孕育了.
後來的發展.
他有一些作品.
其實很多的.
我只選取其他.
比較重要的.
剛才說的讚美蠢材.

$^{1761}$The Place of Foley.
用很通俗的.
讀得明的.
我忘記是不是拉丁文.
來對於.
教廷的.
不是挖苦的.
用一些很有趣的趣味性.
來講他們的問題.
讀起來.
大家大快人心.
引發了很多人.
爭相去看這本書.
所以引發.
基督徵兵手冊更加.
出名.
更加暢銷.
更加流通.
某程度上也是伊拉斯姆.
他的神學.
的灌注.
剛才提過.
在平信徒當中.
做一條約士.
扭開.
教廷的腐敗.
的門.
帶來溫和的.
格心.
不是革命.
也不是沒辦法談的決裂.
所以他是一個最溫和的.
所以他覺得.
不應該譴責教皇.
所以馬丁路德.
在那個.
講這番說話.
這個伊拉斯姆.
其實他不贊成.
後來他寫的這篇論字.

$^{1801}$猶豫字.
就真是他的神學.
人是可以參與.
在神的救恩上.
合作論.
馬丁路德很快.
一見到.
馬上反駁.
因為牽動了.
他的恩信的青衣.
牽動了他.
人是沒辦法靠自己的.
也不能夠和神合作.
是神.
將救恩給我們.
連那個信的動力.
都是神在我們裡面去做的.
所以那種唯獨.
和可以參與.
又是零和一的距離.
所以這個論自由.
自由意志.
我要去搜索一下.
找找這篇文章.
應該可能有的.
因為我們有些基督教文藝的.
古典文學.
一些著作.
教父的著作.
我們在神學院也有讀.
我不知道這一篇還有沒有.
可以讀一讀.
馬丁路德的反對自由意志.
應該肯定有.
變成某程度.
基督教改革來到.
第三條路線.
被排走.
因為只能夠是主流的.
馬丁路德的決裂.

$^{1841}$離開.
基督教就還基督教.
天主教就說.
這些分離的弟兄.
在耶穌的救恩.
陽圈以外的弟兄.
快點回來.
他們永遠的陽圈是沒有問題的.
你分離了.
你們就有問題.
所以你們是分離的弟兄.
天主教時常這樣稱呼我們基督徒.
這是歷史帶下來的.
當然.
現在到了現代.
就不會爭到你死我活.
不過仍然在教義上.
我們都不可以再.
仍然說自己是天主教徒.
這是事實.
所以你想一下總結.
伊拉克斯姆.
一個這麼出色的教育家.
修士.
翻譯聖經的始祖.
很多的元素.
根本就是宗教改革.
根本就是他.
甚至是元祖.
差不多是元老.
但是在處理.
真理上.
他的採取.
就是溫和的合作.
不要挑戰教廷.
在教廷以下.
做那個.
改變.
人心的改變.
所以不徹底.

$^{1881}$馬丁路德.
是他的後者.
又不可以直接叫他.
學生.
是好朋友.
但是因為這一點.
馬丁路德採取的就是.
要不你們改變.
收回那個贖罪券.
承認是錯誤的.
和恩順.
如果不是.
我是不會收回.
我已經說了.
除非你們在聖經裡找到.
我有什麼錯誤.
如果不是.
我就不會收回求主幫助.
這一番言論.
就是馬丁路德採取的.
這樣的.
不是決裂的.
他想討論的.
但是天主教不肯改變.
他只能夠站在.
他的捍衛真理.
這一點.
我相信大家都可能.
明白到的.
馬丁路德不是為反對而反對.
不是為挑戰教皇.
而反對.
他一個人生.
他自己一直以來.
他都看自己是修士.
看自己是教會內部的.
一個勸說.
不是一個挑戰者.
不是一個以下犯上的麻煩友.
所以馬丁路德.

$^{1921}$一生都是以自己是修士為榮.
反而相反.
剛才說了.
伊拉斯姆當他周遊列國之後.
發現他更有興趣.
就是教育.
所以他就放棄.
做修士和做一個地方的主教.
那個任命.
就將這個修士的.
那個終身的.
承諾.
一個誓約.
誓言.
交回給教廷.
我要離開.
我不再做修士了.
因為關係好.
都接納他這個請求.
他就真的還俗.
成為一個正常的.
一般的教育家.
神學家.
所以.
兩者.
我們今天第四堂就說第三條路線.
溫和的改革.
馬丁路德.
就不是.
很激進的改革.
那些農民的戰爭.
就假借.
宗教改革.
因為開了文字.
教廷都是有錯.
和那些田主.
都是有貪婪.
和掠奪.
我們想拿回自己那份.
但是用的就是.

$^{1961}$立即的一種革命.
恐怖主義.
燒教堂.
等等.
變成三條路線.
我們今天.
在主流的基督教.
一定是以馬丁路德.
這個唯獨.
真理的捍衛.
作為主流的.
改革.
所以我們在今天.
我們未去到一個這麼.
嚴重的改革.
改革的思想的.
一些種子.
其實我們應該都要具備.
例如我們.
看到伊拉斯姆.
他在平信徒當中的.
那個論說.
希望透過他的文字.
去改變很多.
信徒的心靈.
回去讀聖經.
給聖經影響.
這些明顯不是保羅的書信裡面的.
教導來的嗎.
中世紀還要強調這些東西.
正正就是這麼多年被教廷.
綁住了.
大家不知道.
聖經有這樣的教導.
就出現了這樣的落差.
雖然回去.
用這些文字.
去引領他們回去讀聖經.
其實是這個意思.
所以伊拉斯姆都扮演著.

$^{2001}$一個文字工作者的角色.
用文字去影響.
那一代人對聖經的重視.
和對.
實踐聖經的訴求.
所以.
這個是.
我覺得可取.
也是宗教改革.
不需要再有什麼改之前.
仍然我們.
每一個教會都要.
平情而論去想一想.
我們有沒有再推動.
那個平信徒讀經的.
那個熱誠.
那個真的看到聖經是寶貴.
而他.
很想受聖經影響.
這樣去.
去做呢.
當然自由意志那裡.
我都不贊成伊拉斯姆.
這樣說的可以和神合作.
慢慢就.
不知不覺淡化了基督的.
唯一救恩.
唯獨基督這一部分.
就會有點牽強.
所以後來嘉爾文.
再將他總結了.
嘉爾文是後代的.
看到你們兩個競爭.
在說自由意志.
反自由意志.
嘉爾文是後代的人.
他就看到了.
那五個唯獨就很重要.
唯獨恩典 唯獨聖經.
唯獨基督 唯獨信心.

$^{2041}$唯獨榮耀神.
這五個唯獨.
就更加奠定了.
這個.
宗教改革的.
真理的.
立場.
一定要有真理.
要有一個學說.
是站得穩的.
你那個運動.
提出運動的人都死了.
你那個運動繼續下去.
不會變質.
就是這五個唯獨的捍衛.
這五個唯獨不斷被傳遞.
就引發十七世紀的大復興.
所以全部事情.
都是神一直一直在做.
我們一定要相信.
地上有神學家在做這些事.
但是平衡的世界.
天上神是管理著.
所以你會看到.
當人大量翻譯聖經.
翻譯完之後要讀.
大量讀聖經 講聖經.
十七世紀就是大復興的開始.
大復興之後.
人就發現原來自己.
得到救恩沒有用.
還有這麼多非洲人.
印度人 中國人.
全部是野蠻人.
不是很學問的.
你快點去跟他們傳道.
就不是去做那個叫做帝國主義.
但是沒辦法.
因為帝國主義太方便了.
你坐戰艦去中國.

$^{2081}$比你自己坐木船快.
有些宣教士就坐戰艦.
因為他們是國家的傳道人.
很自然就坐戰艦.
去廣州.
去那些地方.
你是有軍隊護航的.
當地的人就會反對.
所以早期是有這些.
梁惠不齊的傳道人.
來到不同的國家.
做這些殖民.
表面是傳教.
但其實都是殖民的主義.
所以我們回來看這裡.
今晚我們看這個人.
其實是很出名的.
他就是在.
這個人其實是很出名的.
一個教育家,神學家.
和他對早期的宗教改革.
他留下的那些種子.
馬丁路德使他能夠孵化.
能夠生長.
所以如果看.
伊拉斯姆.
其實他的貢獻都不會少得過.
馬丁路德.
所以每一個時代.
我們總有我們人的.
神給我們的恩賜.
我們的極限.
如果我們用盡.
在一個完全沒有包袱.
沒有成見.
神給你有什麼恩賜.
你在那個時代用了它.
就會祝福後世.
那個功過就由後世的人去評.
等於我們今天講回頭.

$^{2121}$我都不贊成這個自由意志.
講到這個可以合作.
你吃什麼餐.
今天叉燒包好吃.
明天你吃蛋治.
那就沒什麼所謂.
這是神給我們的自由意志.
但救恩的部分.
伊拉斯姆可能多了一些.
就會帶來.
一般人就慢慢趨向.
我自己決定了.
那個人本主義就是今天的現代社會.
歐洲人.
其實那個個人自由.
就無限放大了.
其實會不會受這些影響呢.
是一個我們值得去研究的一個大課題.
今晚我就和大家講到這裡.
我們一起祈禱.
多謝主給我們.
都是趣味性的.
將宗教改革的三條路線.
我們很簡單.
透過幾個人物就講了.
多謝你給我們看到.
伊拉斯姆這位這麼出色的教育家.
神學家.
他帶給那個時代.
對聖經.
對實踐聖經的重視.
和他自己.
是一個溫和的人.
不希望挑戰教廷.
對教廷一種侮辱.
反而他很想改變.
很多平信徒的內心.
雖然他都不贊成.
馬丁路德的離開.
甚至後來.

$^{2161}$他成為了宗教改革的一個陌生的觀望者.
不再是一個.
參與者.
但我們知道.
每一個時代.
每一個人都有他的恩賜.
有他的限制.
有每個人自己的盲點.
或者包袱.
我們很盼望.
我們在不同的時代.
我們在我們自己現在身處的教會時代.
我們懇求主.
神給我們有多少的學問.
知識.
恩賜.
盡量在神的裡面.
用愛.
來幫助其他人.
不要建立我們自己的野心.
我們自己這麼小的.
所謂王國.
反而我們真的很想.
讓更多人認識主.
得到救恩.
願神幫助.
多謝你們.
讓我們又可以有一個半小時.
一齊聽聽這些古代人物.
作為一個趣味性的言擇.
但他們留下的.
很多他們重要的想法.
其實到今天.
還是這麼重要.
願主幫助.
下星期我們會正式講.
加爾文的時候.
幫助我們能夠在神的恩典裡.
學習.
獻上我們的感恩禱告.

$^{2201}$奉耶穌基督的名.
阿們.
\newpage



\section{}
\label{sec:DzI9N0JOL0s}
\textbf{2024-12-12 周四晚查經班《改革之路(5):加爾文(一)》 講員:曾裔貴牧師}
\newline
\newline
連結: \href{https://youtube.com/watch?v=DzI9N0JOL0s}{\texttt{https://youtube.com/watch?v=DzI9N0JOL0s}} ~~~~ 語音日期: 2024-12-13
\newline
\newline
\hyperref[sec:jvEmZbSGJ_A]{\small{< < < PREV SERMON < < <}}
~
\hyperref[sec:index_chronic]{\small{[返順時目]}}
~
\hyperref[sec:index_scriptual]{\small{[返順卷目]}}
~
\hyperref[sec:4EM8nFC2gtc]{\small{> > > NEXT SERMON > > >}}
\newline
\newline
$^{1}$我們才一起祈禱.
很感恩天父給我們今天晚上.
再一次來到這個課題.
宗教改革 我們去到第五講.
盼望我們由嘉爾文.
都是已經去到第二代的宗教改革的.
神學家 聖經的學者.
他們都是帶動那個社會的改變.
也都讓我們看到第二代.
盡心在聖經的研究.
在一些神學的真理上面.
是很有系統的將整本聖經.
怎樣來到在社會上面去實踐.
來到很有系統的寫出來.
這些都是很重要的基石.
成為宗教改革一直的發展.
很重要的一個藍本.
嘉爾文只不過是55歲就離開世界.
可以想像得到他的人生.
由信主 應該說由他接納宗教改革.
加入基督教 馬丁路德的行列.
可能根本只有二,三十年的時間.
就是這樣寫了句著的基督教要義.
不單只是一個神學.
還有很多的《釋經書》和翻譯的聖經.
可以想像得到他那個勤奮.
那個被神用的人生.
盼望我們透過這幾堂.
去講嘉爾文的時候.
我們都可以掌握到這位神的僕人.
他一些值得我們去效法的優點.
求主幫助.
我們今天晚上在你面前等候.
奉耶穌基督的名 阿們.
好 我都當大家一路看到那個PowerPoint.
我們又開始寶血的.
上一次因為忘記錄音.
所以今天要重新講多一次.
第一講就是今晚希望能夠講到.
整個嘉爾文的生平.

$^{41}$他為什麼會跟隨宗教改革.
和他在前人奠基好之後.
他有什麼的成果.
和在前人的成果之上.
他又再有貢獻呢.
我們要看一看對於他的生平.
他在1509年出生.
1564年離開世界.
其實嚴格來說他是法國人.
在這個叫洛安的地方.
我先看看這個PowerPoint.
在這個地方是巴黎.
我今天剛剛搜尋過這個地圖.
巴黎以北的地方叫洛安.
現在這座建築物就是他的博物館.
所以如果大家有機會去.
都會參觀到嘉爾文的博物館.
應該是很小的地方.
但是你想一下.
400年後的人都對他的事跡有一個懷念.
就當是讓大家知道一點點這些現況.
但是他真正的起家是在瑞士.
他一直以來都在瑞士來事奉.
這個就是嘉爾文他那個小小的生平.
我們再看下去的時候.
我們上一次講過.
宗教改革帶來的.
我們先回頭.
第一次我們講了馬丁路德.
他是一個德國的修士.
本身就是天主教大公教會裡面的一個牧者.
所以從來沒有想過他要去做什麼的革命.
或者改革.
這個是馬丁路德的內心那個初心.
他從來都沒有想過.
我要這個人像拍到我現在要怎麼死.
淑英 你要將那個.
現在我正在搞.
好 其實將你自己那個譬如靜音.
還有那個停止視訊.

$^{81}$你自己停止視訊.
這樣都可以了.
對不起 我們先回到上面.
我要退回這個 不要阻礙.
馬丁路德他看到當時.
正值就是教廷委任了一個人.
去周圍在德國去販賣那個贖罪券.
還有他看到教廷這樣的做法.
其實是差不多令到民不聊生.
很多的田農耕種農田的人.
其實他們是最底層的.
他們是沒有自己的田地.
要租田.
有很多層 譬如最高的就是那些貴族.
宗教人士.
下面就有些分侯.
有些騎士.
所以他們很多級.
各級都有自己的封地的.
所以你可以想像到.
最底層的就是那些農民.
所以當宗教他們要建教廷.
他們要建很宏偉的教堂.
或者需要修葺.
他們什麼都不想了.
就是要來叫平民百姓.
不是捐錢 不是自願.
而是透過一個神學的扭曲.
就是說死去的人.
他們現在在陰間.
只要你們買贖罪券.
他們在陰間的受苦的年期.
就可以減低了.
這樣的一個做法.
就是說你可以看到.
原來一直以來.
教廷是壟斷了解釋聖經的權.
和扭曲了聖經和人和神之間.
原來是有一個中介.
一定要透過中介.

$^{121}$你才可以和神有交往.
所以你自己直接是不可以和神有任何的交往.
必須要透過這些神職人員.
所以我們看到那個精髓.
宗教改革其實就是馬丁路德.
很想和教廷有對話.
指出這個贖罪券本身是有很大的錯誤的.
和聖經的恩順清義.
其實是沒有關係之餘.
反而會阻礙了人對於贖罪的觀念.
所以由這樣而產生.
就是1517年10月31日.
那個威登堡的95條.
因為當時的印刷術我們說過.
印刷術已經是很普遍在歐洲開始流行.
很容易就將馬丁路德的原稿印刷散佈.
所以很多人都可以看到.
因為是用德文來寫的.
很多人都可以看到這個95條.
由一個這麼普通的教士.
來越權挑戰教廷的權威.
這個就是宗教改革最後一根稻草.
就是馬丁路德終於來引發整個德國.
對宗教的腐敗來到一個抗爭.
但不是一種很決裂性的抗爭.
所以這樣的緣故.
各地信賴支持馬丁路德的.
無論是宗教人士.
或者是平民百姓.
都會透過這樣.
來跟教廷.
跟當地不同的地區.
來提出決裂.
決裂的有幾樣東西.
你不再向教廷納稅.
原來十分一個人的人頭稅.
是要交給教廷的.
另外還有其他不同的貴族.
收取的地方稅.
田稅等等.

$^{161}$所以有很多的名目.
所以令到平民百姓.
其實是幾乎不堪言.
所以都去到一種.
積壓去到很大的壓迫.
真的要去到一個爆發點.
所以馬丁路德.
就很自然地.
被無論是有善意.
或者是有惡意的人.
來利用.
所以你會看到.
到了2015年27年.
大概十年後.
就有農民戰爭.
這個很嚴重的.
德國的田農.
集體大規模的.
當然是有一些騎士的階層.
變動來向當時的貴族.
向教廷來做一個抗爭.
最簡單就是不再給稅.
燒教堂.
和驅逐神職人員.
他們自己去納自己的傳道人.
或者是一些的牧者等等.
所以這樣的情況.
就變成完全那個勢頭.
是一發不可收拾的.
德國就再慢慢地.
有其他的地方.
例如我做了一些筆記.
例如在蘇黎世.
瑞士.
有這個慈雲里.
這個慈雲里都很重要的.
因為這個PowerPoint的.
1529年的馬保會議.
就因為聖餐的問題.
就是慈雲里和馬丁路德.

$^{201}$兩個人之間的.
談不攏.
就決裂了.
就是教廷裡面.
已經是腐化了.
但是居然間宗教改革裡面.
為了這些的聖餐問題.
本來就不需要這麼堅持的.
但是都必須這樣堅持.
只會帶來就是.
宗教改革就很.
你可以說是很多的山頭.
不單止是這樣.
接著還有其他的.
例如在.
我不說其他的.
密蘭頓在其他的地方.
密蘭頓本來早期想說的.
但是都可能不說了.
因為他的影響力沒有.
馬丁路德.
沒有這個伊拉斯.
沒有這個嘉爾文.
這麼大的影響力.
所以都間接地說一說就算了.
不集中說他了.
所以你會看到.
在當時馬丁路德.
做了這個行動之後.
各地的傳道人.
或者他們本身是一些主教.
或者是一些教廷的人.
紛紛地就要求獨立了.
所以就很簡單.
幾樣東西.
立即宣佈.
和教廷來斷絕關係.
首先就是不再向教廷.
立十分一歲.
不再接受教廷委派的傳道人.

$^{241}$接著就要自己選立.
自己按目的.
這些傳道或者是牧者.
還有.
你可以想像到.
當然有些激烈一些的.
就是會將那些教產焚燒.
不給教廷.
令到教廷處於一個很嚴厲的境況.
各地有不同程度的反抗.
法國就是極之保守的.
仍然對教廷是擁戴的.
所以不是每個地方.
譬如德國就相對比較多.
對教廷有抗議.
所以不同的地方都有不同的問題.
但是當我們再去看嘉爾文的時候.
我們就留意到.
在這個時間.
就是1529年.
這個馬保會議.
就著這個聖餐問題.
這個上一次我們都提過了.
馬丁路德的同質說.
再說一次.
三個其實應該是四個.
天主教到今天還是一樣.
變質說.
變質就是說.
聖餐的餅和聖餐的葡萄酒.
是變成耶穌的肉.
耶穌的血.
是神聖化了.
是有神蹟在那裡的.
但是問題就是.
明明是沒有變.
你怎麼騙自己有變呢.
所以他們就不理那麼多.
真的試過有人吃到.
咬到是耶穌的肉.

$^{281}$這樣就變成明明是沒有神蹟發生.
你又要愚弄那些信眾.
說有神蹟.
說那個就是聖餐餅.
就是耶穌的肉.
所以馬丁路德.
天主教的一個教士.
他自己都不想完全否定.
他就說這個叫同質.
相同的質地.
耶穌的肉.
沒錯在現場.
他沒有變到.
那個餅沒有變到耶穌的肉.
就是說他那裡就去神蹟化了.
但是那個意義是完全一樣的.
你吃下去的時候.
你等同於吃了耶穌的肉和血.
所以在餅和杯的角度.
都是和天主教一樣的變質.
不過你現在看到的是同質.
是那個餅沒有變到肉.
那個杯的葡萄酒沒有變到耶穌的血.
但是是相同的質地.
這樣又變成好像兩者拿個中間點.
因為第三個就是慈雲里強調的.
什麼都不是.
我們只是紀念耶穌的十價犧牲.
那個餅那個杯的意義是沒有用的.
我們紀念的是耶穌的十價犧牲就夠了.
這個都是近代福音派比較多去採納的.
我們應該福音派不會採納變質和同質的.
是慈雲里的紀念說.
但是嘉爾文就第四個了.
嘉爾文就說沒錯.
那個紀念都是對的.
但是不要忘記聖靈在每一次聚會.
聖餐的聚會的臨在.
是會帶來那個聖餐餅和聖餐的杯.
是會有一種對我們基督徒.

$^{321}$就是洗禮的基督徒.
是有一個屬靈的模樣.
和那個臨在帶來的生命的一個活化.
所以仍然都是很重要.
是要來到去在聖餐裡面.
來到讓聖靈在我們眾人裡面的工作.
所以嘉爾文是加上聖靈的同在的.
所以你可以說是一種靈性上的變質.
那又不可以全錯.
所以我都傾向第三和第四的結合.
一定是紀念的.
但是聖靈的工作的確我們要強調.
和我們要等候.
要安靜.
這樣就不會變成行禮如儀的一個活動.
因為我作為牧者.
尤其是在感受的時候.
在台上有時候真的會見到.
弟兄姊妹未必很專心在聖餐的聚會.
其實一個月一次.
一年只有12次.
你想想10年才是120次.
其實不是很多.
人生不是很多的聖餐.
所以這個都是我們應該要很認真.
所以嘉爾文提到的聖靈工作.
其實都要留意.
為什麼要說1529馬堡會議呢.
其實是令到嘉爾文激發他去想.
宗教改革到底所為何事.
和讓他有機會去做一個沉澱.
他去思想到底宗教改革.
是不是真的靠譜的呢.
是不是真的可以帶到我們去看聖經的真理呢.
所以這個馬堡會議的決裂.
是很可惜的.
是馬丁路德和慈雲李正式的決裂.
真是大家是咬到耳瘟面紅耳綠.
這樣來到去大家分開.
你有你做.

$^{361}$你在蘇黎世.
我就在德國.
這樣大家做.
所以你又會令到教廷.
嘲笑基督教.
你們走了.
但是為了這些這麼小的事.
都來到去.
因為這些不是絕對的真理.
這樣來到去大家分開.
這個就是馬堡會議帶來的一個小小的漣漪影響.
對嘉爾文產生一個.
你可以說好奇.
也可以說讓他正式有機會去思考他人生的走向.
我們又回來看這個.
剛才我這樣說了.
應該剛才我按下去就會好一點.
變則說.
就是這個.
我連名字都難讀.
Trans就是轉變.
Substantiation.
Con就是共同的意思.
譬如Consensus.
就是共同的意見.
紀念說就是慈雲你所主張的.
大家記得變則和合則.
或者叫同則.
其實相近就是那個餅和杯變還是不變.
這個就是馬堡會議.
一個會議裡面只是討論聖餐.
你可以想像當時的人.
為了這些事情都看得這麼緊要.
雖然我們還沒有說宗教改革的過去和現在.
將來.
我們其實都要思考.
如果一些的真理.
我們都應該要多一點去思考裡面的細節.
和到底哪一樣東西是更加合乎聖經.
這個都是值得我們去參考.

$^{401}$但是去到一個爭論.
去到一個決裂.
這個有可能都是.
因為當時真的是一個風頭火勢.
馬丁路德17年到現在29年.
其實都是11,12年左右的時間.
可能那個情緒都很繃緊.
既要面對教廷的決裂之後的追殺.
那個驅逐.
和當時教廷結合的國王.
那些君主.
對馬丁路德的追究等等.
所以可能都是一個很繃緊的裡面.
對於自己的同路人.
慈雲里.
都好像大家都有一種這樣的情緒的糾結.
不知不覺都看對方是敵人.
所以原來我們那個情緒.
都可能需要面對.
在一個大運動裡面.
都需要用一些空間去沉澱.
所以你會看到第二代的嘉爾文.
就沒有這個問題.
一來避開了第一批的宗教改革的人.
那種面對教廷的抗爭.
需要用很多的時間去想一下.
有什麼辯解的真理.
和怎樣逐點去反駁教廷對他們的指控等等.
所以不知不覺.
自己都進入了一個抗爭的那種的.
那種的反抗裡面.
就連自己的同路人.
都變成抗爭的對象.
其實含默聖經的確.
我們今天四百年後再去想.
當然我們都認為.
應該是以紀念加上嘉爾文的聖靈的學說是合理.
但是不至於我們要這樣去變成一個抗爭.
一個對立.
這個要留意.

$^{441}$我們要看到嘉爾文正正式式加入改革運動的行列.
我又要說一下他的人生.
因為他本身出身在一個富裕的家庭.
他的爸爸原來就是當時在法國一個很大的教會裡面.
主教的一個助手.
所以無論收入.
無論學識.
所以十二歲那年.
嘉爾文已經在法國一些出名的學校讀書.
他開始要修讀的已經是讀哲學.
讀那些辯證學.
對很多事物的辯證.
當時是一些人民主義的學科.
很重要的.
對你的思辯能力很強.
在那段時間嘉爾文有很多的機會.
因為教廷的學校有很多機會去讀教廷的一些著作.
尤其是第一代的神學家奧古斯丁.
很多很多很重要的作品.
神學的作品.
嘉爾文是有機會專心去研讀的.
所以他對神學的很多觀念都是來自奧古斯丁的.
我們提過的基督教要義.
關乎到三一神論.
其實嘉爾文都是已經採納了奧古斯丁的論點.
譬如三一真神.
三位一體.
這些嘉爾文都不再加以有什麼新的見解等等.
我們想像得到受奧古斯丁的神學影響的嘉爾文.
還有就是有其他的一些.
當時特別是中世紀的時候.
有一位叫做伯納多.
Brenner.
這個法文我不認識了.
那位神學家也是出名的.
在中世紀的著作.
那個神學.
影響著嘉爾文.
所以嘉爾文本身是一個法國人.
他對於古代的古典的原典.

$^{481}$古代的原典.
那個涉獵和研讀.
他是第一手的資料.
再加上原來當時.
教廷一定是那些修士才有機會讀的希伯來文.
希臘文.
因為要讀聖經.
所以某程度上.
你在教廷裡面受教育.
你就有機會接觸.
這些不要小看他還沒有加入基督教之前.
原來教廷給他的優厚的栽培.
塑造他將來在宗教改革裡面.
他是一個非常重要的.
一個著述的神學家.
因為他小時候已經開始讀希伯來文和希臘文.
所以對於他將來的聖經翻譯.
和他研讀聖經.
是有很大的幫助的.
所以他不單止是寫基督教要義.
他寫下到今天流傳的.
《釋經》書都有很多本.
你可以想像到.
嘉爾文在他童年的時候.
受的宗教的教育.
是對他將來投身這個改革.
是有非常大的幫助.
無論是人民的學科.
一些的修辭.
思辨.
辯證學.
哲學.
加上有神學.
奧古斯丁.
白納多.
再加上聖經的原文.
希伯來文.
希臘文.
差不多他已經讀了一個神學的學位.
是真的認真地讀.

$^{521}$而且是從12歲開始就已經開始讀了.
不過很有趣地發現.
因為他在法國.
剛才不是說了.
法國是教廷一個很保守.
你可以說是很.
如果用我們香港說的.
就是很深藍的.
就是擁戴教廷的.
所以爸爸又發現.
原來做傳道.
做修士.
將來那個收入.
不是那麼好.
同時也對基督教.
到處去燒教堂.
又是產生反感.
所以在這樣的情況下.
爸爸就改變初衷.
不准嘉爾文去走那條修士的路.
好像馬丁路德那樣.
而是強迫他要立刻離開.
就要讀法律.
所以根據那個記載.
嘉爾文是一個很順服.
很聽話的兒子.
所以就是這樣的緣故.
他就不再走神學的路.
但是他已經有一定的裝備.
對於聖經.
特別是他自己能夠看原文的聖經.
就不是看那個拉丁文版本的武加大譯本.
所以嘉爾文就有很多的機會.
他自己去思考.
和自己去研判.
特別他好像是坐享其成.
因為馬丁路德在1517年.
來到去宗教改革.
一路走下來.
然後他又看到有些農民發動戰爭.

$^{561}$又看到那個混亂.
又看到教廷和那些君主鎮壓那些平民百姓.
嘉爾文其實是看著這一切的發生.
但是他就在大後方.
在法國的出生地.
對不起.
他已經送去巴黎讀書.
在巴黎最優質的學校裡面.
就是說在教廷裡面.
他受到某程度的保護.
還沒有加入改革的行列.
你看到這個部分就知道了.
去到1533年.
其實他一路涉獵奧古斯丁.
一路看馬丁路德.
一路看一些其他的著作.
更重要的是.
所以他日後對於聖經的研究.
和他自己神學其中很重要的主張.
就是一定要回到聖經的真理.
所以一定是他自己在那個運動裡面.
他不只是跟風.
更加不是情緒的一個波動.
更加不會是一個人魂逆魂的.
這樣一窩蜂的去跟一個熱潮.
而是他有機會在巴黎去沉澱.
在聖經的原文裡面.
他自己去認真的研讀.
他有一個判斷.
就是原來馬丁路德這一班人.
不是為反而反的.
他們真的有他們的立場.
是真的有聖經的真理的支持.
所以來去改變的.
那個導火線.
為什麼嘉爾文都要被迫離開巴黎.
離開法國.
去到瑞士呢.
就是這個事件了.
就是他的好朋友.

$^{601}$這個巴黎大學的校長.
托克利克勒的這位好朋友.
被邀請成為巴黎大學的校長.
在那個就職的演詞裡面.
托克校長就引述宗教改革.
一位比較溫和的.
就是我們上一次講的那個人物.
就是伊拉斯姆.
這位都很偉大的一個教育家.
的一篇文章.
就引述.
但是伊拉斯姆本身就是.
宗教改革的其中一個重要人物.
不過他仍然強調.
要在教廷裡面.
來做那個對話和改革.
就不要像馬丁路德那樣脫離.
分裂那樣.
所以托克都沒有想過.
這個就職演詞帶來.
一個這麼大的迴響.
因為你在教廷大公教會裡面.
一個重要的巴黎大學裡面.
來發表一個宗教改革的演詞.
所以可想而知.
那個後果那個影響.
所以固然托克受到牽連.
被來解職.
上一次我都講過.
因為這一篇的演詞.
嘉爾文是有份幫手撰稿的.
所以嘉爾文更加受牽連.
他就要被出走.
他就必須離開.
因為那個時候.
你可以想像得到.
那種的決裂.
1517到1533.
再加上有13年的農民戰爭.
其實已經是定性了.

$^{641}$是革命是要推翻整個.
教廷和國家.
所以是勢不兩立的.
誰支持宗教改革.
誰就是反動的叛亂分子.
所以很正常地.
嘉爾文只能去一些.
支持宗教改革的地方.
所以他就.
他真的不去德國.
其實他一直都不在德國.
他就去了瑞士.
將來瑞士的宗教改革.
嘉爾文就成為最重要的推手.
嘉爾文去了瑞士之後.
就帶來一個很大的.
你可以說那個地方的一個很大的發展.
這個時候.
我們簡單地講了他的生平之後.
相信他最重要的.
就是要定義什麼是宗教改革.
什麼導致到和教廷不能夠不分開.
還有沒有什麼更加善意的.
一些真理.
我們可以大家對話.
只要那個對話是貼近聖經的.
相信那個分裂是可以慢慢地減少.
甚至乎會合一.
那就必須要大家都有個基礎.
教廷定性了自己的.
譬如他們開過天特會議.
定性了自己.
譬如教皇無誤.
贖罪券是沒有錯.
等等.
其他的.
還要有那個聖禮.
和其他的一些他們自己衍生出來的.
那些約章條例等等.
但是基督教.

$^{681}$當時大家群雄.
不同的程度的分開.
只需要你和教廷決裂.
只是你不再納粹.
不再接受他們委派的主教.
還有再激進一點的就是.
將那些教產焚燒.
這樣的一個新的基督教.
這樣是不能夠帶來真正的.
一個教會的慢慢的建立的.
因為當這些情緒過了之後.
到底人信什麼呢.
那些弟兄姊妹離開教廷之後.
跟隨這一班的領袖.
我跟隨你們.
到底我們信什麼呢.
是很需要有一個真理的基礎.
來作為一個基石.
使到宗教改革可以繼續蔓延下去.
而且這個蔓延.
不是再出於一種感性情緒.
甚至是為反對而反對.
這樣的不合邏輯的.
這樣的一個煽動.
不單止是沒有意思.
而且只會帶來更多人誤解了.
所以我們第二堂已經說了.
農民戰爭就是吃著馬丁路德的改革.
你們的教廷都可以這樣有人去挑戰.
為什麼我們長年受壓迫.
我們不可以挑戰呢.
所以某程度上他們的訴求.
那些農民的訴求完全沒有錯.
不過用的方法.
用的手段是不合適.
和是聖經禁止的.
你想想殺人放火.
燒教堂.
這樣的掠奪.
這樣的暴徒式的.

$^{721}$所以馬丁路德由一開始.
對他們的同情.
到後來公開的譴責.
甚至乎站在那些農民對立的一面.
你可以想像到.
那個改革帶來的實在太過宏大了.
所以大量的人訴諸於那種情緒.
特別可能是長期受到教廷的各種壓制.
精神思想的壓制.
當突然間這樣釋放的時候.
就可能會帶來那種反抗.
是會很不理性.
訴諸於情緒.
甚至乎訴諸於心底的報復.
只會帶來更大的破壞.
所以在這個時候.
其實很早了.
他1533年走到瑞士.
1535年已經有第一版的基督教要義.
你可以想像到.
這些教義問答.
是宗教改革走了差不多20年了.
是必須要有的真理的基礎.
不是每個人都能夠讀得懂聖經.
聖經也未能夠大量的分發.
當然已經有了.
但是如果有簡單的.
逐個的關於神的事情.
所以基督教要義是很有系統的系統神學.
第一刻就是神.
神是怎樣的.
我們要怎樣面對神.
神對人有什麼要求.
各樣東西是很清楚的.
好像一個手冊.
你想找關於基督.
你可以去那裡.
關於聖餐.
每樣東西都有的.
所以如果我們逼自己每天看一段.

$^{761}$其實也很有幫助的.
現在在我辦公室的桌上.
我嘗試強迫自己每天都看一段.
反正買了兩大冊.
你都不看就浪費了.
有需要才來找.
所以其實是很淺白的.
真是一個教義的問答.
好像一問一答.
所以某程度上挺好看的.
因為已經翻譯成中文.
真的相當容易.
因為用法文寫.
然後經過很多很多的再版.
更正.
但大部分其實沒有什麼怎麼變的.
你可以想像得到嘉爾文那個神學觀.
是很受奧古斯丁的影響的.
教義問答.
剛才說的就是.
必須在宗教改革差不多20至20多年的時候.
神就要興起一些神學家.
有系統的寫.
為什麼我們要改革.
為什麼宗教改革是事在必行.
而且是會延續.
就算我們這一代一代二代的傳道者.
宗教改革的神學家.
相繼過去了.
但那個運動是會繼續延續下去.
那個運動能夠延續下去.
首先就是神學觀合乎聖經.
另外就是嘉爾文很重視的聖經權威.
真的可以延續下去17世紀.
17世紀就去到英國的大復興.
英國沒有受到太大的影響.
當然是聖公會.
因為已經和天主教決裂了.
因為是國家國王的決裂.
因為教廷不准英國的國王離婚.

$^{801}$所以國王說我理你都傻.
我脫離教廷.
自己開個聖公會.
變成他很早就離開了.
但在教義上.
英國沒有真正跟隨宗教改革.
只不過是政治上的改革.
不是宗教的改革.
但無論如何.
嘉爾文的影響力都去到英國.
就是聖經的權威.
已經被下一代的人.
看為很重要.
所以一來那個聖經翻譯.
已經去到差不多整個歐洲.
不同語言的翻譯都已經湧現了.
所以每一個國家裡的人.
都可以讀聖經.
可以去教會.
可以得到聖經的教導.
就帶來開始在人的心靈有那個醞釀.
有那個神的聖靈.
透過聖經的被尊重高舉.
而帶來一個復興.
所以是有一個方程式的.
所以我們要看到這些.
我們今天的福音派.
譬如我們的培靈研經會.
其實都是在這個方程式裡面.
每年為甚麼十天去研讀聖經.
是真的會帶來有信徒的復興呢.
是真的因為神的道在人的心.
所以我們看到基督教要義.
就有好幾個特質.
要來去強調.
首先就是.
你想一下剛才我提過.
其實我們一直說.
如果有些立體圖就好了.
因為嘉爾文要面對的就是.

$^{841}$譬如馬丁路德他怎樣延續.
他支持的.
怎樣延續馬丁路德的中國改革呢.
這是其中一樣東西.
另外就是.
他怎樣去駁斥那些極端的改革呢.
例如松洗派.
很激進的.
例如教廷.
教廷不是沒有了的.
到今天都存在的.
怎樣面對教廷的政治.
宗教的壓迫呢.
所以嘉爾文.
他安身在瑞士.
就讓他能夠有足夠的空間.
去沉澱.
去寫基督教要義.
既是幫助宗教改革.
有一套很合真理的神學.
這是很重要的.
任何的運動.
如果沒有真理基礎.
或者叫做法理基礎.
那個運動就一窩蜂完了.
就散了.
是不會有繼續下去的.
第二就是.
面對教廷的那個逼迫.
激進派的幫倒忙.
怎樣能夠去撥亂反正呢.
如果我是一個持平的.
我都很想了解.
到底你們宗教信什麼的.
為什麼要和教廷決裂.
我都很想了解你們信什麼的.
為什麼你們要燒教堂.
為什麼你們一定要人家重洗呢.
所以這個基督教要義.
當時的重要性.

$^{881}$不在話下.
必須要讓一般人都能夠看得明白.
所以真的很簡單.
很淺白的講解每一個教義.
讓人感受到原來自己可以同神相交.
我想我們今天沒有這個問題.
我們的問題就是我們自己不肯同神相交.
不是別人逼我們不准同神相交.
要經過主教.
要經過其他的手續.
這又是另一個今天的時代的極端.
我們是一個不祈禱的世代.
但是當時其實是你不可以自己祈禱.
或者你根本不敢自己祈禱.
因為已經是經年累月教導你.
你不要隨便讀聖經.
聖經是要由教士用拉丁文讀的.
你祈禱你要來到教堂.
你要找神父.
你要告解.
因為神父就是那個中介人等等諸如此類.
很多很多他們的一層一層的.
好像架床疊屋的那些繁文縟節.
所以令到一般人.
我只顧著謀生.
我是一般又沒有什麼學問的人.
我怎麼會理你那麼多.
我總之禮拜.
我穿一套漂亮的衣服.
乾淨的衣服.
去教堂望彌撒.
就當作是宗教的生活.
其他你不要煩我.
你要我做什麼我就做什麼.
要交稅.
要買贖罪券.
那我就照買.
就是這樣.
所以基督教要義的出現.
首先就是指出大公教會的重大的錯誤.

$^{921}$重要的是基督教要義絕大部分是根據聖經的真理的啟示來寫出來的.
所以看基督教要義.
都某程度等同於你對聖經有一種同步的認識.
甚至乎是嘉爾文自己都有一定他自己的詮釋.
另外就是改革的性質.
就是要駁斥這些激進的從洗派.
那個激進.
那個帶來社會的不安和動盪.
和對教廷的一個直接的挑戰.
直接的一種的不單止架空.
甚至乎是用一些的武力.
要使教廷好像要瓦解一樣.
這個就是激進的從洗派.
當然更重要的重點就是基督教要義裡面.
就是要確立神的道是一切所有人信仰的基石.
聖經的權威.
你可以想像到當時家家戶戶都沒有聖經.
沒有人翻譯聖經.
沒有人夠膽翻譯聖經.
識原文的只有教廷裡面的讀書的教士.
是可以來去研讀.
平民百姓.
我講說讀自己的文字的聖經.
根本沒有翻譯.
所以聖經是從來都沒有試過這樣來去被傳遞的.
所以某程度由零開始.
因為有零零星星的馬丁路德之前的一些前輩.
開始膽敢地翻譯聖經.
有了這一點的火就一直延續.
所以其中的伊拉斯姆就是將這個武加大譯本翻譯成新的希臘文.
所以就帶來很大的一個時代的一個劇變.
當然配合印刷術傳到歐洲.
所以人人都可以擁有聖經不再是神話.
不再是大罪.
甚至成為一個風潮.
甚至成為一個人人可讀經的一個運動.
所以結集這樣的時候.
加上有加爾文的這個推動.
對聖經的權威的重視.
天主教都是重視聖經的.

$^{961}$不過是他們要默許的重視.
要由他們的講解.
要在教廷裡面去宣讀的拉丁文的聖經.
才叫做重視.
那就是說一般的平民百姓完全聽不明白你說什麼.
完全不知道聖經說什麼.
那其實就不是重視聖經.
其實你是捆綁著聖經不讓平民百姓知道神的心意.
這個是很大的罪.
某程度就是神使用這些改革.
叫神的話語被釋放出來.
讓所有當時的歐洲的人都有機會.
終於可以讀回自己文字的聖經.
這樣帶來的就是下一代真正的宗教的大復興.
你可以想像到加爾文一生最大的影響力.
就是對聖經權威的高舉.
我們先看還有一點時間.
加爾文神學大家應該稍為都有機會聽過.
雖然加爾文絕對不是荷蘭人.
更加他沒有去過荷蘭做大量的牧羊工作.
但為什麼荷蘭的神學會用這個鬱金香.
看為是加爾文的神學呢.
正是因為這個預定論對於加爾文的改革中.
有很大的影響力.
所以荷蘭的改革中.
就真的奉加爾文為他們的一個很重要的神學家.
他們後來演變了所謂的突獵.
這個鬱金香神學.
是加爾文主義的一個精髓.
總共有五樣.
他們選了第一個字母.
全然的墮落.
Total Depravity.
講到人是不可以自己拯救自己.
必須透過神的救贖.
人是全然無可取的.
所以對中世紀,十六世紀,十七世紀的人.
一點都沒有困難.
沒有人會覺得自己是無神主義.
沒有人會覺得自己不需要神.

$^{1001}$所以這個全然墮落.
對很多人來說.
無論是加爾文的.
或者是亞米紐斯的神學.
都是接納的.
但是加爾文就更加絕對.
因為他強調的是那個揀選.
所以這個全然墮落.
是沒有影響到人可以得救.
你不會因為你的墮落.
去到一個地步神就拒絕你.
不要你這樣的一個情況.
所以當然要五樣東西結合起來.
才是真正的加爾文主義.
如果你只是要一兩個.
就不是加爾文主義的精髓所在.
大家要留意.
這些當然大家聽完就算了.
因為今天的教會比較少再講.
這樣的五個向導.
反而加爾文的五個唯獨.
反而多人會去接納.
這個就是無條件的揀選.
有另外一派的神學.
我就沒有機會做那個powerpoint.
大家聽一下名字就算了.
因為我們在神學院.
真的會被人問.
你是加爾文主義還是亞米紐斯主義.
特別我們上第三季的神學的時候.
就必須要自己做一個表態.
你是屬於加爾文主義.
還是屬於亞米紐斯.
因為亞米紐斯就剛剛跟加爾文有相反的.
就是人是有限度的揀選的.
還有人可以自己去有選擇的.
就是說你可以有選擇不得救的.
對於加爾文那個剛剛的一個對比.
這些太複雜了.
將來如果有機會再講.

$^{1041}$今天我們主要是講加爾文.
一定要懂得後人對他的這個主義的一個詮釋.
當然你不可以說100\%都是加爾文自己本身.
因為他已經死了.
他就已經這樣去講.
沒有什麼修改.
但是綜合來說.
突裂的精神就是.
人是傳言墮落.
但是神有一個無條件的揀選.
因為神願意萬人得救.
明白真道.
不願意有一個人沉淪.
乃願意人人得到那個拯救.
這是提摩太前書的二章一節的經文這樣講的.
所以在加爾文的揀選裡面.
是無條件的.
沒有因為你是什麼人.
就有限度的揀選你.
是不會的.
第三就是有限的袋淑.
這個Limited Atonement.
下面一點.
有限的袋淑就是說.
我們講過了.
耶穌基督的救恩.
一次相信.
但是.
必須是要每個人.
在這個信心裡面.
要堅持到底.
這樣才可以是真正的信.
這個精神在加爾文的主義裡面.
也都反映的.
所以某程度上.
你不是一信.
就一了百了.
就必然得救.
不是加爾文主義的那個精神.
而是你必須保持自己在神的愛.

$^{1081}$在那個信心裡面.
對自己的得救和耶穌的關係.
是一直要保持.
這個叫做有限的袋淑.
無法抗拒的恩典.
就是說.
如果人認識耶穌的救贖和大愛.
你是不應該.
抗拒主耶穌的救恩.
神賜的這份恩典.
你可以想像到.
和上文講的.
其實有些表面的衝突.
其實沒有衝突的.
人的墮落.
是不會自己去尋找神的.
已經失去了尋找神的那種能力.
已經沒有了神的形象.
在人的心靈裡面.
但是那個大前提就是.
有一位他者.
就是我們人的以外的一個存活者.
這位造物主.
他認識我們.
他揀選我們.
甚至乎應該叫做預定我們.
在這個救恩裡面.
所以當我們一回應這個呼喚的時候.
我們就在這個無法抗拒的恩典裡面.
來接受這個呼召.
這個呼召在嘉義民的眼中看來.
就是無法抗拒的.
所以有些沒有被揀選的人.
就一直藐視這份恩典.
你想一下.
那個恩典本身已經是很美的禮物.
是令你無法抗拒的.
你想一下多麼美.
但是大量的人.
藐視耶穌.

$^{1121}$嘲笑十字架.
嘲笑馬槽.
甚至乎看耶穌的救恩.
是和他作對.
因為你要說我是罪人.
我相信你的耶穌.
你現在這樣批評我是罪人.
他們都會有這樣的一種抗拒.
但是.
所以很奇妙的就是.
全然墮落的人.
居然會看.
當他的信心得到了.
當他有聖靈工作.
打開他的心靈.
他會看到耶穌基督.
是他無法抗拒的一份恩典.
所以.
這個肯定就是背後一個無條件的揀選.
因為如果有條件的揀選.
我也不會羨慕耶穌.
我也不是在基督教家庭長大.
怎麼會羨慕這個洋教呢.
怎麼會得到這個救恩呢.
所以冥冥中是有對我這個人的揀選.
我很堅信嘉義民的這個觀點.
所以對於圖利普.
我沒有什麼問題.
我又表態我是嘉義民主義.
因為在神學院真的爭論.
老師要我們做功課.
我就已經說了我是嘉義民主義.
其實有很多是亞米老師的.
挺好玩的.
同學之間幾百年前的事.
但是我們自己做神學的時候.
都要有一個表態.
哪一方是強一些的理據等等.
對於我個人來說.
這個無條件揀選.

$^{1161}$肯定是塑造我這個人的人生.
很感動的.
冥冥中人海.
剛才我在元朗走了一個小時.
買東西吃東西.
慢慢慢步走.
在教育路很多人.
準備今晚的講道.
即是默想.
這麼多人.
很偉大.
想起嘉義民55歲的人生.
寫了這麼多書.
留到我們今天都可以看到.
他那本基督教要義.
我的桌上有中文版.
連我喜歡的時候.
我姐姐幫我改名叫嘉義民.
Calvin.
我姐姐是摩門教的.
當時被摩門教帶回教會.
中學的時候.
學校是英文書院.
英文老師就說.
每個同學回家要改一個英文名.
上課的時候就用英文名.
我什麼都不懂.
回家問姐姐.
姐姐說叫Calvin.
我都不知道是什麼.
叫自己做嘉義民.
所以讀神學的時候.
一來被人笑.
嘩 你叫自己做嘉義民.
有沒有搞錯.
我不關我的事.
真是冥冥中揀選.
接著再加上對他的神學.
尤其是那五個唯獨.
我自己在主義學都教過.

$^{1201}$都做過一些筆記.
真的非常敬重這位神的僕人.
所以剛才我在元朗這樣走的時候.
慢慢走就會思想 會默想.
所以鼓勵大家在越多人的地方.
安全地慢走.
用腦去想想.
是會想到一些東西的.
我感動到幾乎想要一點點眼淚.
當然天氣又冷又少.
黑 人們看不到.
嘩 人生原來最悲哀就是.
你認識神.
我自己當時的默想.
人生最大的悲哀.
當然可以有很多程度的.
例如有些人.
他的身體不健全.
他的思想不健全.
這些是一些悲哀的.
是一些人生的困難的.
我都遇過這些人.
對他們特別需要有一種憐憫.
但是神對這些人.
是不會特別苛求的.
甚至會特別給他們有恩典的.
但是有些人就是很厲害.
很有學問.
甚至乎又認識神.
但是最大的悲哀就是.
你的人生不被神用.
這樣就真的是人生的最大悲哀了.
因為你空有你所有的恩賜.
就是神給了你這麼多的恩賜.
你一點都不用.
而且你用你的人生.
去做一些不能夠存留後世的事情.
我覺得就是最大的悲哀.
因為嘉義民留給我們的財產.
我們今天去研讀他們的神學.

$^{1241}$都很得幫助.
你可以想像得到.
如果他一直都不寫他的文字.
一直只是停留在教廷.
多一事不如少一事.
不要決裂.
那就不會有這個聖經的人物.
這個神學的人物.
讓我們可以看到.
宗教改革到今天.
都需要繼續延續下去.
所以我當時的默想.
就是覺得如果能夠.
在有限的人生被神用.
那個用不一定要很大.
總之是神安排你一個很小的崗位.
重要就是你忠心做好.
人家就算給你一個不是很好的崗位.
你都做好它.
用心去做好它.
專一的做好它.
就算多不好的崗位.
都可以帶來一些改變.
因為用心就會有改變.
剛剛突然間下午又忽發奇想.
快到年尾了.
叫弟兄姊妹寫一些感恩和願望.
在淘寶淘了一棵樹.
仿真樹.
叫弟兄姊妹弄上去.
突然間想到.
不是什麼的.
在講台上.
讓弟兄姊妹看到.
大家的感恩.
大家的願望.
我們呈獻給神.
大家就為著前面教會的願望.
一起去努力.
帶來一個很好的彼此的啟發.

$^{1281}$不單純只是弄一棵聖誕樹.
亮燈都有意義.
讓人知道我們這裡是有神的光.
但是那是弟兄姊妹的心願.
所以有時候每件事.
你想一下.
又會有很多靈感.
所以這裡就是.
為什麼呢?.
聖徒的永無保守.
這個就真的是嘉爾文的精髓.
真正的一次得救永遠得救.
是不對的.
是一次得救.
一直相信.
就必然得救.
這就是嘉爾文的精髓.
神學的精髓.
這個Preservatism of the Saints.
就是聖徒永無保守.
對當時代的十六世紀.
是至關重要的.
因為沒有了教廷的保護傘.
教廷說你離開了我們教會.
我們才是唯一的教會.
你們是分離了的弟兄.
你不回來.
你是沒有救恩的.
嘉爾文說不是的.
聖經說在主耶穌的救恩裡面.
就是救恩.
在主耶穌的裡面.
就是永無保守.
《約翰福音》十章二十八節.
我父賜給我的羊.
必定會到我這裡來.
到我這裡來.
我必定會使得牠得救.
我與父緣為一.
父賜給我大.

$^{1321}$到我這裡來.
我必定會保護牠.
所以是一個雙重的拯救.
你可以想像到.
或者我讀整句的原文給大家看.
《約翰福音》十章二十八節.
就更加準確地讓大家看到.
嘉爾文的神學精髓.
十章《約翰福音》.
當時是祭司長和那些人.
對耶穌的對話的攻擊.
第二十七節.
我的羊聽我的聲音.
我也認識牠們.
牠們也跟著我.
我又賜給牠們永生.
所以你留意到.
是一次相信.
一直相信.
因為我的羊聽我的聲音.
我也認識牠們.
牠們也跟著我.
一直保持這個信仰關係.
接著二十八節.
要識背.
因為這個真的很重要.
我又賜給牠們永生.
牠們永不滅亡.
誰也不能從我手裡把牠們奪去.
我父把羊賜給我.
牠比萬有都大.
誰也不能從我父手裡把牠們奪去.
這個就是雙重的保障.
因為耶穌會保守.
天父也會保守.
所以這裡就是十章的二十七至三十節.
希望大家都識背.
這個就是聖徒永蒙保守的精髓.
所以你會看到.
這是一個盡情的.

$^{1361}$你可以說是一個很完滿的.
對人性的了解.
在救恩裡面.
人為什麼要救恩呢.
因為是傳言的敗壞.
但是現在.
在神的啟示裡面.
是有一個無條件.
任何人都可以是神所愛.
只要你願意接受.
你就可以得到了.
這個是有限的救贖.
就是說如果人拒絕.
就不可以得到神的救恩.
所以神的願意.
不等於人人都可以順理成章的.
自動就轉帳.
就成為基督徒.
是不可以的.
他必須要自己有一個選擇.
打開心靈去相信.
所以當這個相信的時候.
一定是受著一份無法抗拒的恩典.
就是耶穌十誡的大愛.
來感動.
這樣來信的.
一個信主的人.
他要一生看到一件事.
不是他去保守自己的救恩.
是神的恩典.
聖靈的工作.
去令到他蒙保守.
所以他只需要活在.
與主耶穌的生命的關係裡面.
聖徒就永蒙保守了.
這就是加爾文主義的五個重點.
當然在基督教要義裡面.
就會充斥著這些教義.
今天我要做的筆記.
就來到這裡了.

$^{1401}$大概我們講了一點點的生平.
講了基督教要義.
帶來的社會當時的運動的穩定.
加爾盟約很多的神學.
透過聖經的真理去做詮釋.
讓當時很多未必識字的平民百姓.
都有機會去讀這些神學.
很簡單的明白.
再加上後期開始翻譯聖經.
讓聖經無論是法文.
德文他就沒有做.
讓當時的人都有機會自己讀聖經.
再強調加爾文和馬丁路德.
都是經常講道的人.
就是說講道的普及化和宣講.
是基督教一個很重要的核心工作.
相對於教廷每個禮拜的彌撒.
好像在唸南無阿.
當然這樣有一點貶低.
是不應該的.
始終他們都是神職人員.
這樣用拉丁文去讀.
唱一些這樣的經文.
是不會令到人有共鳴感.
是不會令到朝聖去望彌撒的信眾.
得到任何的牧養.
這樣久而久之.
教廷和信眾中間.
是有一個很大的鴻溝.
這個某程度我們都要留意的.
因為我的兒子讀聖公會.
他中學的時候.
學校規定一年有兩次上下學期.
都要去聖公會.
所有學生都要去聖公會.
去禮拜的.
所以不約而同收到的feedback.
就是根本不知道那些主教在做什麼.
只是看到很多儀式.
又要他們跪.

$^{1441}$你完全沒有去解釋.
給那些學生知道.
又要那些學生.
人家根本還沒信.
只不過因為讀聖公會.
白約翰就要去參加你的聚會.
你不講解.
和你的千篇一律.
行禮如儀.
對一個年輕人.
對一個少年人.
只會覺得是一個厭煩.
只會覺得是一個包袱.
只會覺得完全和他們沒有關係.
好像受壓迫.
一定要參加.
這樣就帶來對宗教有一個誤解.
這個陰影.
其實影響著他們的成長.
所以在中學讀基督教學校的人.
長大後很多都不信耶穌.
可能那個底不是打得好.
就會對宗教有一個誤解.
看到的可能只是禮儀的部分.
沒有看到禮儀背後有一套神學.
神學其實是一個很活的有機體.
去承載我們的人生.
當然長大之後.
經過人生.
經過歲月.
經過不同的事情.
經過不同的朋友.
有些人也會信主.
有些人面對人生.
也會開始反思.
這個世界是不是真的有神.
也可能他們小時候的種子.
開始發芽.
這些我們全部不知道.
但是這個也是我們教會需要反省的.

$^{1481}$基督教教會要反省的.
如果大部分的信眾.
每個星期過來崇拜.
都不知道發生什麼事.
都沒有任何的所謂得著.
沒有所謂的一種共鳴.
哇 奇怪.
為什麼今天的講完.
講中我內心的困難.
疑惑的問題.
我祈禱.
神居然透過他解答我.
或者今天我坐下.
這首詩歌就解開我的結.
這個共鳴感.
未必需要說出來.
而是聖靈運行在聚會當中的那個連繫.
所以我們既然是聖靈的教會.
我們無論是牧者.
信徒.
一定要相信這種超自然的整合.
奇妙地.
我完全不知道在在誰有什麼問題.
我只是按感動.
甚至是根據牧者給我的指派.
我今天要講這一段經文.
要講這個題目.
我就是這樣講.
我從來沒有想過.
誰誰誰有這樣的.
沒有的.
這樣這個就是聖靈的工作.
這個就是聚會的神聖的契合和共鳴.
如果沒有共鳴.
不行的.
因為小時候回來玩.
在兒童主義學是沒有問題的.
中學開始有一點點的拉鋸了.
開始睡眠不足.
很多事情要做.

$^{1521}$溫書.
各方面.
慢慢又變成一個.
好像被強迫去做.
到了再大一點.
去到大學.
大部分的同學朋友.
都是不信耶穌.
那就更加是一個.
好像少數派.
那個壓迫感就更加大了.
加上如果回到教會.
都是完全未同著立.
見到很多周圍的哥哥姐姐或者叔伯.
都是這樣.
每個星期一樣坐那個位.
一樣這樣去聚會.
聚會也是一樣.
大概都猜到了.
哪個時間說什麼.
到時哪個報告.
全部大家都完全有心理準備.
沒有任何驚喜.
沒有任何聖靈工作的火花.
這個就等於教廷.
使很多平民百姓.
根本不知道什麼叫信神.
根本不敢和神有溝通.
這個就是宗教改革的導火線.
所以我們都要留意.
我們的詩歌.
我們的港道.
我們整個聚會的設計.
是不是真的讓每一個人.
都能夠來到神的面前.
他自己可以安靜下來.
敬拜神.
回應神的教導.
還有他有一個走出去.
立志要改變的動力.

$^{1561}$還有他感受到自己.
是蒙神保守的一個信眾.
他繼續有光彩.
面向他自己的難處.
所以牧者不是幫他解答了.
在公司的困難.
而是再次賦予他這個身份.
他有力量靠著主去面對.
同樣都是星期一要面對的困難.
但是他已經有了那個力量.
所以這個應該就是崇拜.
帶給那個同在.
和那種共鳴感.
所以港道有時真是有這種效應.
和這種祝福.
你真的不知道哪個人聽完.
那個帶來的生命的改變.
所以很盼望.
透過我們研讀嘉爾文.
今天是第一講.
下一個星期我們就會進入.
他的一些神學裡面.
我們先講五個唯獨.
和他在瑞士帶來的社會改革.
也都有一些爭論的.
因為當他做了瑞士的城邦的領袖之後.
又開始有他們的內部的問題.
又會鎮壓一些反對嘉爾文主義的人.
又是所謂行私刑.
又是一個好像五十步笑八步.
這個我們都值得去反省.
今天我們就講到這裡.
早一點就完結了.
再一次多謝主.
我們能夠透過一個半小時.
差不多的時間.
我們去將嘉爾文的生平.
他的影響.
特別是他的神學.
他的宗教的改革.

$^{1601}$尤其是在聖經的權威的高舉上.
他帶來的影響.
多謝你讓他早年在教廷的栽培下.
他有機會涉獵大量的研讀奧古斯丁的作品.
和有很多哲學辯證學的訓練.
再加上他自己很喜歡看希伯來文和希利文的聖經.
使他能夠有相當的造詣.
來成為宗教改革的一個神學論述的一個作家.
和一個神學家的一個重要的推手.
多謝主.
你是這樣來讓宗教改革不會失控.
也不會回頭走回頭路.
更加不是為反而反.
多謝你透過基督教去給當時的歐洲.
有一個很大的復興.
回到聖經.
回到研讀一些舊約的著作.
和興起這班年青的傳道者.
來勇敢地宣講福音的真理.
求天父給我們今天.
雖然我們不是要做什麼革命家.
我們懇求主讓我們看到.
我們需要在一個時代裡面.
看到這個時代的癥結.
有什麼是影響著.
無論是明顯的.
或者是很隱晦地影響著教會.
或者是一個時代的.
盼望主都給我們能夠有這一份的眼光.
更加不是為反而反.
而是我們要將一些錯的.
是不對的.
真理的.
已經可能是大家都習以為常的.
我們能夠有效.
有智慧.
很客觀地將它提出來.
提上去一些議程.
成為大家都關注的.
好像我們說.

$^{1641}$今天很大部分的信徒.
其實都不會回到祈禱會.
甚至乎自己的祈禱生活.
都是比較少的.
天父我們發現.
這個都不是一個個別的問題.
忙碌的社會.
真的會帶來我們成為一個很少祈禱的人.
但是我們懇求主讓我們能夠看到.
正正就是因為這樣.
我們更加需要來呼籲.
願神幫助我們.
這個只是其中一個例子.
還有可能是我們根本就沒有興趣去讀經.
自己用心去研讀神的話語的重要.
願神真的幫助我們.
讓我們能夠很簡單地將嘉爾文.
先講了他的生平和學術.
願主讓我們在前面.
還有一兩堂的時間.
我們再認真看他的神學裡面的精要.
願神幫助大家.
奉耶穌基督的名.
阿們.
好 謝謝各位.
\newpage



\section{}
\label{sec:4EM8nFC2gtc}
\textbf{2025-01-02 周四晚查經班《改革之路(6):加爾文(二)》 講員:曾裔貴牧師}
\newline
\newline
連結: \href{https://youtube.com/watch?v=4EM8nFC2gtc}{\texttt{https://youtube.com/watch?v=4EM8nFC2gtc}} ~~~~ 語音日期: 2025-01-03
\newline
\newline
\hyperref[sec:DzI9N0JOL0s]{\small{< < < PREV SERMON < < <}}
~
\hyperref[sec:index_chronic]{\small{[返順時目]}}
~
\hyperref[sec:index_scriptual]{\small{[返順卷目]}}
~
\hyperref[sec:3Sqlgs98HDM]{\small{> > > NEXT SERMON > > >}}
\newline
\newline
$^{1}$好 感恩我們能夠透過停了一兩個禮拜之後.
又再恢復今天我們這個宗教改革的查經.
盼望大家都能夠透過這個分享.
去認識關於這個16世紀的宗教改革.
一些重要的發展.
和一些我們值得去關注的.
想問一問在線上的看到這個畫面.
就是這個PowerPoint的了.
如果看到的可以給我手勢就行了.
我們開始一起祈禱.
再一次感恩天父給我們恢復.
我們今天晚上的查經.
我們在這個比較大的教會歷史的題目.
宗教改革的過去的內容,發展.
還有我們真的很想去很簡單的思想一下.
宗教改革到今天已經是500年後.
我們看到基督教一直透過普世的憤興運動.
普世的宣教運動.
去到非洲,去到我們的中國,去到印度.
當然有其它的地方.
但是我們想到.
都真是想到福音一直這樣行進.
去到一些從來都沒有福音到的地方.
都是透過這個天主教內部的一個再一次的反省.
而出現的這個改革.
所以我們都盼望神給我們在今天的教會.
雖然我們在香港都有不同的宗派.
但是都給我們秉承宗教改革這些的先賢.
他們的那個內心對神.
特別是對神的道的重視的態度.
給我們能夠有所學習,有所反思.
我們今天晚上會講加爾文.
求主給我們對他的一些事蹟.
他推行的改革有一定的掌握.
又給我們能夠在其中去做一個思想.
我們有什麼可以效法.
或者有什麼我們要給面子,求神幫助.
我們這樣仰望禱告.
奉耶穌基督的名 阿們.
好 這是第二講.

$^{41}$我先要看回那個錄音 等一等.
可能我還沒有開.
OK 因為好像沒有電.
現在行了 沒問題.
我們裝一點重溫.
他的加爾文其實法國人.
他年紀不是很大.
他55歲就離世了.
他1509年出世.
1564年就離開世界.
當然是在瑞士的日內瓦的地方.
輾轉你可以想像到他1534至1535年.
就已經奠定了他的基督教要義的第一版的面世.
當然就是整個歐洲普及.
已經奠定了他在宗教改革的神學家的身份.
和更加讓宗教改革有理有據.
有真理的奠基.
和讓宗教改革不會百花齊放.
不會不同人有不同的演繹.
他這個建樹是相當重要的.
因為改革開了第一槍的馬丁路德.
他從來都沒有準備任何的文件.
他只是很想和內部的天主教友對話.
裡面有一些東西不合聖經.
特別是贖罪券.
買來只是去圍著教會的修葺.
那個維修.
變成一個基金.
但是背後是帶來很多的平信徒.
買了贖罪券之後.
只會更加肆無忌憚地去犯罪.
所以更加遠離了神和那個赦罪本身的本質.
馬丁路德看到這些情況就去做.
所以他本身不是有一套很完整的神學.
來做一個改革的.
所以只是好像被逼.
因為他釘了95條在威登堡的教堂.
就引發輿論.
整個歐洲廣泛的對天主教的.
真的有人突然間提出.

$^{81}$很多人就會附和.
因為壓迫了很多年.
每個地方又要向天主教的教廷交稅.
還有很多種種其他的繁文縟節等等.
所以宗教對當時的中世紀.
是一個壓迫.
是一個枷鎖.
多過是一個心靈的釋放.
多過是一個.
大家一家人很想去禮拜堂敬拜神.
所以你可以想像到.
如果馬丁路德一直這樣推行的時候.
沒有一個神學的.
一個很有力.
又基於聖經.
那個改革慢慢慢慢就會淡化了.
因為那個一鼓作氣的熱心.
慢慢都會冷卻的.
很容易又會被天主教.
因為人家答應過千年的系統.
很容易又會將你打垮了.
所以神就興起了.
有加爾文這位本身讀法律的一位年輕人.
他在1534,35年左右.
就寫了這個基督教要儀.
而且還是非常淺白的.
很容易看的.
很有條理的.
分了不同的類別的.
什麼是神.
我們信的神是怎樣.
三一神是我們的信仰.
變成對一個普通的信徒.
只要他稍為識字.
都能夠透過這些這麼有系統的編排.
來找到.
譬如關於祈禱的.
他很詳細講.
我看了有幾個章節.
有四個重要的部分.

$^{121}$關於祈禱等等.
這些是很有條理.
很簡單易明的.
一點都不深奧的.
這樣就奠定了.
當時的宗教改革.
是有一個很重要的神學基礎.
很重要的一個清楚的.
根據聖經的論述的教義.
所以他的奠基.
成為下一代的門徒.
他的門生.
他的學生.
把基督教的宗教改革.
更加發揚光大到整個歐洲.
其中一個最出名的學生.
就是來自蘇格蘭的John Cox.
約翰·洛克斯.
也是帶來英國大復興的神學家.
所以基督教要義對當時代是很重要的.
所以你可以想像到.
他二十多歲.
他已經寫了這本巨著.
過了二十多年.
他再去改版.
其實都是大體上.
可能加一些的潤色.
加一些的解釋.
整體的神學其實沒有變的.
再加上很大部份的神官.
都是採納奧古斯丁.
和中世紀的一些神學家.
他沒有自己的新的原創.
或者沒有自己的石破天驚的一些改變.
所以就令到整個基督教的道統.
那個神學一脈相承.
當然他就針對當時的.
所以你會看到.
這個改革運動.
上一次我們講過了.

$^{161}$就是他的好朋友.
這位叫做Kirk.
本來就是去到法國巴黎大學裡面做校長.
但是他的就職禮的演說.
就是支持宗教改革.
整個法國就是天主教的勢力.
所以他的演說.
就帶來天主教正正式式.
就和這個Kirk.
包括Kirk的好朋友就是加爾文.
加以逮捕.
所以他們就被逼.
要離開法國.
輾轉去到這個.
似乎是瑞士的地方.
巴塞爾的地方.
在那裡避難.
在那裡因為瑞士.
又是支持宗教改革的.
所以一個地方一個城邦.
是有不同的.
對宗教改革的支持.
或者是他的熱心等等的.
所以在這樣的情況下.
他就有很大的空間.
因為他不是馬丁路德那些.
最原始的抗爭.
所以他是看著一些前輩.
那個對天主教的背棄.
那個的決裂.
那個的離開.
所以他就有機會.
這個年輕人.
他不是在一個大的漩渦裡面.
所以他有足夠的空間.
他去沉澱他自己.
所受的天主教的教育.
因為他在巴黎.
都是一個很出色的學生.
在巴黎大學裡面.

$^{201}$都有很多機會來涉獵.
很多的經卷.
很多的神學家的一些文章.
一些著作.
特別是古典的著作.
所以嘉爾文就有很多的.
他的足夠的空間.
因為他不會被人來批鬥.
所以他就一直去研究整個基督教.
的一些神學.
所以還有一些古代的著作.
再加上馬丁路德提出的這些論點.
當然最終一定是在理性上.
他是同意和認同這個宗教改革的.
就不只是一種情緒的激進.
我們上一次都有機會提過.
宗教改革有三條路線的.
馬丁路德只是很想在教內.
在宗教內部產生對話.
不是想決裂的.
不過是被迫.
天主教說你不收回.
你就要被逐出教廷的了.
只能夠馬丁路德.
就唯有他黯然的.
不想的.
他是一個神父.
就離開天主教.
但是當他引發這個改革之後.
就有些激進的從洗派.
就覺得什麼都要反.
就不單純只是贖罪券那麼簡單.
而是整套的教義.
整套對人民的壓迫.
都是要改變的.
包括當時的聖禮觀.
聖餐的問題是最嚴重的.
就帶來有一個派別.
就是從洗派.
特別是慈雲里.

$^{241}$都是同時代的神學家.
宗教改革.
就帶領一個很激進的從洗派.
當然還有普通的平民百姓.
因為上次都講過.
那些城邦有不同地方的稅收.
有不同的諸侯.
有不同的一些騎士.
都是由教廷.
由一些當時的世俗的皇帝.
分派,分一些的地等等.
就好像不同的諸侯那樣.
就變成真正那些平民百姓的種田.
做那些農業的.
就受盡各種不同的稅收的剝削.
所以他們變成一個好像受壓迫的階層.
而這些是最多的.
所以德國因為馬丁路德的這樣改革.
既然你是一個神父.
你都可以和教廷有一個決裂.
我們受壓迫的農民.
就更加要提出組織他們那個.
大家的團結.
就來帶來一個真是徹底的戰爭.
直接是燒教堂.
驅逐那些神父聖職人員離開.
不再向教廷交稅.
還有那個叫十二條款.
就講到明我們要列了這十二條.
你不遵守.
不去滿足.
我們就抗議.
那個抗議是一個戰爭來的.
所以馬丁路德是看著這些農民戰爭.
這些務農的田農.
這樣來到去武裝的革命.
是產生戰爭.
所以馬丁路德無奈地.
他是支持政府要鎮壓那些農民.
大概有十三年.

$^{281}$有差不多十萬的德國的農民.
被鎮壓而戰死.
所以你想想在馬丁路德的時代.
開始的宗教改革.
原來你溫和.
人家是不會溫和的.
在宗教上有從洗派.
在世俗上有農民戰爭.
所以其實都要面對這個改革.
到底第二代怎樣將那個改革.
定性在教義上.
這個就很重要了.
所以這個時候.
基督教要以這本書.
成為很重要的信仰綱領.
當時沒有足夠的時間.
他們去翻譯聖經.
大量的去講道,講解.
所以同步進行中.
一方面要面對.
天主教從來都不妥協的.
從來都是基督教.
是一路一路的激進的.
脫離自己的教廷的.
但是失控的宗教改革.
就有很激進的從洗派.
要在整個信仰上.
全部撥亂反正.
天主教什麼都是錯的.
這是其中一個激進.
另外世俗的農民戰爭.
就是長期受壓迫的群體.
要借宗教改革來做一個大幅度的推翻.
就變成整個歐洲.
其實在一個很動盪的情況下.
在這個時候.
1535年.
基督教要義出版了.
就是嘉爾文自己一個人.
整理整套的神學教義.

$^{321}$大部分是系統神學.
但是他也印了很多聖經.
我這樣翻譯一下.
我那兩本中文版的翻譯.
真的他每一個的題目.
都有聖經的支持他的教義.
所以挺好看的.
原來發現英文版.
在香港的圖書館是有的.
在網上都可以下載到英文版的.
變成大家看中文版就難了.
就應該沒有了.
就只是在教會的書室可以買到.
但是很貴的大家不要買了.
因為一套都要五百多塊.
買了大家又不看就不要浪費這些錢了.
其中很重要的.
剛才我說了.
這些教義問答.
就是要指出到底大公教會錯在什麼地方.
所以教義問答.
他們真的很特別.
因為嘉爾文或者馬正路德也是.
但是嘉爾文就更加.
我稍後再講到.
真正在日內瓦帶來的.
真正的一個地上的天國.
這是嘉爾文的心願.
教義問答.
原來是要在小朋友的小學.
我不知道他小學多少歲.
就要開始這些教義問答.
就是給小朋友一問一答.
問一個問題.
他們要懂得回答 懂得背.
這樣開啟神學的文字.
原來宗教改革.
是在那個時候就已經看重小朋友的心靈.
從小就要灌輸正確的神學教義.
這個我們真的要想的.

$^{361}$所以兒童的主要學.
我們不需要說.
我們一定要很好的power point(力度).
那個是技巧而已.
不是說技巧不重要.
而是內容才重要.
所以如果我們怕小朋友.
他們聽這些悶的真理.
行不行的.
其實不要想他行不行.
而是我們要有趣味地.
講到他們明白為止.
用不同的方法去包裝.
譬如用一些圖畫.
或者要power point(力度).
但是那個還是技巧.
重要的就是.
真是很硬朗的教義問答.
是要給他們從小就要認識.
原來這樣的宗教教育.
原來在中世紀的時候.
已經是正正式式地開始了.
所以沒錯.
兒童真是每一個時代的很重要的主人翁.
真的 如果教會.
這些是公認的.
18歲以下是不准去教會的.
在內地.
這個是不合《聖經》的.
《聖經》是告訴我們.
每個人都要清楚自己的信仰.
當然每個人的智力.
每個人的能夠接收.
是很不同的.
是有強有弱的.
有些人真是.
講深奧一點他都不能夠掌握的了.
但有些人是要很理性地去分析的.
每個人的構造不同.
但是大體來說.

$^{401}$宗教教育.
必須是從小就要開始.
有條理地教他們這些重要的《聖經》真理.
所以如果我們不教.
就錯失了宗教改革.
這麼體重這些兒童的教理問答.
所以原來是很重要的.
所以最重要是要指出大公教會的謬誤.
指出不是想推倒它.
人家也有別人的傳承.
有別人的敬殿.
但是如果那個道理本身是有問題.
我們指出那個道理的錯誤.
免得別人信錯了真理.
就不是根據《聖經》來相信神.
所以那個責任是指出那個問題.
然後大家帶來一個理性的討論.
和最終尋找那個真理.
所以激進的從洗派.
就在這裡有過之而無不及了.
所以就令到那個改革的性質混亂了.
這個改革就變相好像為反而反.
什麼都是錯的.
至少人家都是一千年.
由使徒時代.
由主侯的317年.
那個君士坦丁的悔改.
讓基督教可以自由的傳道,輾轉.
都是有一個時代.
接近一千年的時代.
都是承載著這麼多人的心靈.
不能夠因為現在16世紀.
出現了這些賣贖罪券.
就差不多好像是一個的錯誤.
當然裡面有其它的系統的問題的錯誤.
但是不因為這樣.
你就要推翻整個的改革.
或者別人天主教的謬誤.
所以變相嘉義文在這個改革的性質上.
又再一次定性回到聖經的真理的尋索.

$^{441}$就是說等同於不認同.
和駁斥激進的從洗派.
對什麼都要反.
就是一個過度的性質.
當然更重要的就是要將聖經的真理.
透過這個基督教要義再一次著重.
或者我下一次帶兩本書回來給大家看.
引一些chapter我最近看了的.
因為其實我買了很久都沒有看過.
有做功課才看的.
現在逼自己每天看一點點.
累積了一些chapter.
挺好看的.
它講得很有條理的.
翻譯成中文.
由第一版是法文.
接著輾轉去到英文的翻譯.
接著英文再轉中文.
都經過了源遠流長.
但是都仍然很有條理.
真的它已經分了類.
所以很容易看的.
某程度上講得挺好和挺深入的.
有些教義講得挺深入的.
比如信,悔改.
接著就駁斥那些當時世俗對於悔改的定義的人的錯謬.
它真的講得很仔細的.
還有時它自律行間會引同時代的一些神學家.
他們有些錯誤.
或者有些人的經典.
它會引述.
豐富了我們對古代的著作.
少少的涉獵.
所以挺有趣的.
要慢慢坐下來看.
但是相信最重要就是.
整個基督教要義.
最重要的不是為反而反.
而是想將有條理的基督教到底是信什麼的.
最終是要想確立聖經的權威.

$^{481}$所以系統神學都是要服侍聖經的.
系統神學是在聖經裡面讀.
聖經讀出來.
神是三位一體的.
所以系統神學就是要將聖經講的三位一體講出來.
所以系統神學是必須要在聖經神學之下的.
所以兩者是應該結合.
不是有抵觸的.
所以如果有系統神學說聖經不是這樣.
那就一定是系統神學錯.
就不是聖經錯.
所以相信基督教要義最重要.
嘉爾文想帶出的.
或者他自己的內心.
為什麼認同宗教改革.
不只是認同馬丁路德.
而是他認同馬丁路德背後的真理.
不過還沒有一個很詳細的論述的真理.
所以嘉爾文很重要的地方就是承接改革.
將改革放在一個客觀理性的道統裡面.
那個真理的基礎裡面去進行改革.
所以神的道就是一切的.
無論信仰或者所謂的宗教改革的重要的基石.
那個Milestone(基礎).
變成了如果你不是這樣的時候.
那就要反思了.
為什麼你要這麼激進.
要燒人家的教堂.
要驅逐那些神職人員.
要來到那些農民.
要用武力去搶掠那些教堂.
來殺那些神職人員.
那你就可以想象得到.
這個肯定不是聖經.
或者肯定不是宗教改革馬丁路德的初心.
所以在這個時候過了1517 1535都是十幾年.
這一本基督教要儀.
差不多是當時的巨著.
你可以想象得到.
我看梁家倫院長的書.

$^{521}$見到我的院長.
他說如果在這個世界.
只剩下兩本神學書.
要你選擇.
他說不用選擇.
基督教要儀一定要存在.
這樣評價這本書.
有兩冊的.
其實真的是挺重要的著作.
你想一下一個人.
當時是16世紀.
不會有什麼很大量的藏書.
更加不會有今天的AI等等.
可以寫出一本這樣的巨著.
都不簡單.
當時只不過是二十多歲.
加爾文原來是一個體弱多病的人.
好像一個弱質書生.
但是不因為這樣就影響了他的成就.
所以這是基督教要儀.
對宗教改革帶來的.
更加重要的基礎.
不是為反而反.
不是為改而改.
而是為著真理.
能夠在每一個人的心靈.
來產生與神相交的作用的改變.
所以某程度當然是一種撥亂反正.
但是這個撥亂.
又不是好像全盤的都是錯的.
而是將一些真理再一次重新提出回來.
所以當時的特別是關乎到.
那個聖參觀.
馬丁路德都與他自己的好朋友決裂.
宗教改革裡面自己都有分裂.
但是對加爾文來說.
可以提出那個真理.
但是就不需要有一種這麼激進的決裂.
所以這個我們就是.
你可以想像這個就是他.

$^{561}$一個很重要的對當時歐洲的教會.
的一個建樹.
的一個貢獻.
而當時他只不過是二十多歲的一個神學家.
就有這樣的成就.
是非比尋常.
神怎樣用這些的牧者.
這些的傳道人.
但是好景不常.
就是他去到這個巴塞爾.
接著輾轉就去到日內瓦.
與他的好朋友.
其實這個當時已經在日內瓦.
已經是.
應該這麼說.
我看一些著作裡面就是說.
日內瓦在.
這個我已經寫了.
1936年的5月21日.
他們已經是整個的日內瓦的一個城邦.
就是說他自己本身就是一個國家.
當然這個國家很小.
一個城邦.
未與其他周邊的地方合併.
在16世紀的時候.
大概3萬人口的日內瓦的城邦.
有他們自己的議會.
就是類似政府.
他們是公民的投票.
一致贊成要脫離天主教.
歸順基督教的.
所以這個也都受著.
這一位神學家法也納.
我不懂得讀.
這些可能是法文的拼音.
Ferrer.
那個影響.
輾轉.
因為加爾文要逃難.
就去到日內瓦.

$^{601}$法也納就知道.
當時1535年的基督教要義已經出版.
所以因為這本書.
而帶來整個歐洲的影響.
所以就很自然地.
法也納邀請加爾文.
去跟他一起進行在日內瓦瑞士的宗教改革.
當時挺戲劇性的.
法也納就挺直接的去挑戰加爾文.
他說這個是神在我心裡面的啟示 指引.
如果你不來日內瓦.
你是會面對神的咒詛的.
因為你違反神透過我向你下達的命令.
這個指引.
加爾文自然害怕.
因為神的指引.
他輾轉就跟法也納.
就去到日內瓦.
或者真是在日內瓦.
1936年至1938年.
短短的大概三年.
就在日內瓦進行宗教改革.
正如剛才所說.
他們規定.
因為整個.
當時我們很難想像.
整個日內瓦.
差不多三萬的市民.
全部是基督徒.
全部都是脫離天主教.
要歸順基督教.
整個心靈都準備去聆聽.
牧者這些改革家的教導.
所以這個就變成.
開始我們要思想.
16世紀的政教.
不是分離.
是政教的合一.
為什麼三年這麼短.
加爾文他300年之後去了哪裡.

$^{641}$他1564年才死的.
還有很長的時間.
二十多年.
他去了哪裡.
你可以想像到.
原來頭三年的改革.
慢慢的.
那些所謂全城邦的日內瓦的居民.
都是基督徒.
都受不了有這些規範.
是宗教對國民的規範.
你可以想像得到.
只是這些這麼簡單.
聖餐.
人人都要遵守.
每個月一次.
還有如果你不守.
除非你是患重病.
否則你就缺席.
你就要被審查.
為什麼你不守聖餐.
就是說.
原來你個人的自由.
我們很難用今天的思想.
今天的世界.
去想一下16世紀500年前.
他們是怎麼理解個人的自由.
個人的選擇.
剛才提過了.
兒童教育.
就是說.
每個小朋友必須要上主一學.
必須要學會這些要你問答.
就是說一問一答.
你不懂得回答.
又是一個可能要處分的.
還有崇拜.
當然嘉爾文也很看重.
崇拜唱的詩歌是有規範的.
不可以唱一些世俗的詩歌.

$^{681}$必須要唱詩篇.
這些只不過是宗教上.
還沒有去到其他的.
更加重要的一些.
關乎到民生那些.
不過.
但是這三年的改革.
就已經帶來城邦裡面的議會.
發現有很大的反抗.
有很多的日內瓦居民.
覺得很專制.
是很壓迫的.
所以.
最終那個輿論.
帶來就是.
法也納和嘉爾文.
被驅逐離開日內瓦.
去到史特拉斯堡.
靠近德國和法國邊界的地方.
今天是屬於法國的.
但是當時是屬於德國的.
這個史特拉斯堡也是很漂亮的地方.
就在那裡逃難.
在這個史特拉斯堡.
或者當時叫做斯塔斯堡.
但是名字英文就完全一樣了.
Strasbourg.
去到這個地方.
就遇見了這個出色的布塞爾.
也是他一個很好的朋友.
因為布塞爾已經在斯塔斯堡.
做了一段很長的牧羊工作.
還有這個地方.
德國這個史特拉斯堡.
對嘉爾文非常信任.
愛戴.
擁護.
尊敬.
所以.
在這個史特拉斯堡的三至四年.

$^{721}$他又是很好的沉澱.
因為雖然日內瓦不要你.
驅逐你.
但是在德國的斯塔斯堡.
嘉爾文就在這個時候.
開始寫很多的《釋經》書了.
現在他的《釋經》書.
在網上是可以隨便下載的.
嘉爾文的《釋經》書.
500年前的羅馬書的翻譯.
不同的經卷的翻譯.
都有他的注釋了.
是的.
那麼.
這段時間對嘉爾文來說.
是一個很大的保護.
心靈的釋放.
因為這個小小的地方.
對他的接納.
他可以來到更加穩定的.
這樣去寫一些著作.
做《聖經》的翻譯.
還有寫不同的《釋經》書.
所以也再修訂基督教要義.
不是大幅的修訂.
是蕴飾一些裡面的內容.
所以更加奠定他是一個多產的神學家.
神學家和著作的作家的一個地位.
所以在這幾年.
1938至1941年.
就更加讓嘉爾文在德國這個地方.
有更好的研讀.
著作.
講道.
這個都是.
中國的鬼格.
包括馬丁路德.
包括嘉爾文.
包括布塞爾.
他們都是一些講道家.

$^{761}$一個星期要講很多堂道.
所以其實很習慣的.
歐洲人他們當時.
不是今天.
當時他們是很多地方.
當時也沒有那麼多的錄音記錄.
所以他們都是一定要去到現場.
去到不同的鄉鎮去講道.
所以這個也可能是當時巡迴報道的需要.
其中我看的著作裡面.
很寶貴的就是.
都對嘉爾文有一個很大的正面的影響.
就是他在1540年的8月.
就和這位女士.
這位中文的翻譯.
就是伊德萊德·德布爾.
因為應該是荷蘭人.
非常敬虔的一位婦女.
她本身是重婚的.
因為她的丈夫.
都是從洗派的聖徒.
可能是早死.
所以她是一個寡居的寡婦.
布塞爾撮合他們兩個的認識.
大家就結婚.
1540年.
1506年.
也就是說42歲結婚.
也有20多年.
沒有啊.
9年.
因為他的太太在1549年離世了.
所以有9年的婚姻.
但是原來這9年.
對嘉爾文的幫助.
無論是家庭的溫暖.
讓他能夠更加的穩固.
來到無後顧之憂的去發展宗教改革.
等等.
不過記載就是說.

$^{801}$好幾個小朋友.
很小的時候就已經夭折了.
不知道是當時的醫學.
還是什麼原因.
還是營養等等.
對嘉爾文都是一個很大的打擊.
但是夫婦的婚姻.
是令到嘉爾文很得到扶持.
那個幫助.
我想穆姐的婚姻.
都是很重要的一個支持.
所以這裡是小小的題外話.
因為在這個情況的時候.
日內瓦的居民.
剛才說.
日內瓦在頭三年.
嘉爾文推行的宗教改革.
其實都是溫和的.
不是極端的.
只是想有一些的規範.
你既然說自己是信耶穌.
有一些的規範.
在宗教上面有一些的紀律.
這些你不可以說錯的.
就是要兒子.
星期無論多忙都要回教會.
其實很自然就應該要這樣的.
可能對某些人來說.
可能過於的一個嚴苛.
其實雖然表面上個個都是基督徒.
但是我們上一次.
對上一次都有提過.
馬丁路德那個時代.
他們的宗教改革.
原來有一些.
因為馬丁路德的年代.
是有一個運動叫做.
重返那些古代著作的文藝復興.
那些要追尋的原典.
原來有很多這些的.

$^{841}$當時的有識之士.
可能是大學教授.
可能本身是很有學問的人.
其實他們本身對信仰不是太感興趣.
他們對古代的著作會很有興趣.
但是當時.
因為那個壓迫帶來的.
已經是很反感的了.
不過沒有辦法用一己之力.
來對抗教廷.
或者是對抗當時的一些政府等等.
但是當有宗教改革之後.
這些有識之士.
其實他們是自由派的.
他們稱為人民主義者.
Humanist(人道主義者).
他們是很著重人權.
很著重個人的自主權.
這些對當時中世紀是一個很大的.
很前衛的.
但是正正因為改革.
就注入了這個元素.
你可以隔走那個神職人員.
你可以不聽教皇的話.
就是說.
原來背後有很大的空間.
讓這些開始本身已經崇尚個人自由的.
這些自由主義者.
或者這些人民主義者.
就有更大的議價的能力了.
所以也都等於.
在整個政府的運作.
市議會裡面.
我們集中講加爾文.
就只能夠集中講日內瓦當時的情況.
所以在第一次的日內瓦的改革.
就受到很大的.
那些市議會裡面的人的阻撓.
因為那些人本身.
就是崇尚人民主義個人自由的.

$^{881}$我們不要天主教的那個核制.
反過來你找了一個加爾文來.
還更加給我們這麼多規範.
這樣不行 那樣我要跟著做.
那就變成很大的.
就是沒有了一個壓迫.
多了另外一個的箝制.
他們就受不了了.
所以就在那個議會裡面.
醞釀那個罷免.
就最終是將法爾納.
和這個加爾文驅逐離開日內瓦.
對加爾文來講可能是一個打擊.
但是剛才講了.
他去到斯特拉斯堡.
其實是一個很好的沉澱.
重新更加的透徹的去.
讀聖經 研究 還有寫書.
讓他的神學更加的穩固.
所以有一個大事發生了.
就是去到1541年.
當時市議會.
輾轉那個宗教改革的那一派.
就真正得勢.
真正聯合所有的那些市民.
一致好像重金禮聘那樣.
邀請加爾文回來日內瓦.
還要承諾我們一定會跟著你那個宗教改革的步伐.
來進行的.
那封信據聞加爾文看到之後.
都流眼淚.
因為很沉澈 懇切.
還有打動他的心.
還有讓他看到.
這個都是神在他的生命裡面.
給他的一個的召命.
因為其實他在斯特拉斯堡.
已經是很安穩的了.
當地的人是很愛戴這個家庭 加爾文.
還有給他很大的空間.

$^{921}$雖然那個地方是小的.
但是現在日內瓦的市議會.
發出這一份的邀請信.
我就引了這一句.
Post Tenebris.
就是黑暗已經過去.
現在是光明.
一個議會這樣講.
是拉丁文來的.
給加爾文知道.
我們已經是等待 虛位以待.
你再回來帶領我們做宗教改革.
所以1541年的9月13日.
加爾文就真的舉家回去日內瓦.
帶來一個第二次真正的改革.
由41年一直到他離開世界的1556年.
他就正式的離開世界.
這裡有20多年的時間.
都是帶來一個很大的.
使日內瓦奠定了.
宗教改革的發源地.
真正在 我所指的是.
一個很穩固的神學基礎的發源地.
就不是馬丁路德的.
德國的威登堡提出的.
對贖罪卷或者教廷一些問題的.
提出的討論.
而是真真正正一套的基督教神學.
在日內瓦這個地方.
透過加爾文.
透過基督教要義.
還有這個都很重要.
日內瓦的教會憲章.
你准我回來.
之前那些.
你們是不是真的接受呢.
所以他在1941年開始.
就正式草擬這個日內瓦教會憲章.
有幾個重要的部分我們要留心.
將教會.

$^{961}$因為他既然說.
整個的日內瓦都是信耶穌.
都是基督徒.
所以不能夠.
政教分離.
政治.
即日常的生活.
人民的活動.
和信仰.
是不可以分開的.
所以.
你就見到有兩個體制.
整個市議會.
是由.
兩種 應該是四個類別.
評述圖的長老執事.
聖職的牧師和教師.
教師可能是傳道人.
教聖經 不是教一般的學術 學科.
這樣來共同的組成的.
那個市議會裡面.
有12個的議員.
每個星期要聚在一起.
開會談一些關於宗教.
關於民生的事情.
制定很多不同的法規.
律法等等.
市議會裡面的長老執事.
負責去執行那些紀律.
但是關乎到信仰.
嘉爾文就訂定了.
這個長老制.
就是說一定是牧師帶領整個議會.
關乎神的事情.
市議會的長老都沒有權過問的.
所以無論講章 著作.
他們都不能夠去否決的.
特別是講章的講道.
這裡就給我們見到.
嘉爾文就在這個體制裡面.

$^{1001}$灌注了表面上是議會的評訊圖.
他們被選立出來的.
來做這個執行紀律.
治理.
你可以想像得到.
我們今天的唐董會.
是整個感受.
當感受吧.
當感受五千戶都是信主的.
整個唐董會就管理整個五千戶的.
感受花園的居民.
接著呢.
平日他們就有這個巡查.
牧者就要去探訪.
發現原來A段某一個家庭.
他們不單止不去崇拜.
在家裡看電視.
還有那些佛教的書.
還有專佛在那裡.
這個時候.
長老執事就去了.
就發現不是邀請.
是拘捕他們.
來審問.
為什麼你這樣.
就給他們機會.
做一個悔改.
悔改之後.
就來到去.
沒事了.
如果不悔改.
三次警告.
就逐出教會.
不可以守聖餐.
這樣的一個情況.
這個就是.
政教合一的日內瓦當時的情況.
所以.
你可以想像得到.
其實都很難的.

$^{1041}$那個牧養.
牧者就一定要定期.
探訪每一個家庭.
了解他們的信仰.
他們的情況.
有什麼困難.
這個牧養.
原來某程度.
又伴隨著監視.
因為發現有什麼問題.
就將那個個案.
轉交市議會去執行.
那個社會紀律的.
所以.
其實都幾.
在我們今天來看.
好像幾有趣.
幾難以想像.
但是.
某個程度.
今天我們不可能.
因為整個世俗的運作.
已經是去到個人的主義.
和個人的選擇.
教會已經不可能.
好像中世紀那樣.
但是某些元素.
那個觀察.
那個牧養.
那個探訪.
那個了解.
那個信仰的溫度計.
幫他來做一個量度.
問題是你怎樣將這個東西.
將這個所謂個案.
帶回去.
我們今天就不會有市議會.
也不會唐董會談.
哪個人.
他們沒有來守聖餐.

$^{1081}$我們要罰多少錢.
應該不會這樣.
但是.
在這個過程裡面.
怎樣在發現.
然後牧養.
使他能夠回轉呢?.
至少嘉義文是採納了一個.
人是要管的.
你是信主耶穌都要管的.
我們是.
他那個信念是這樣的.
你認不認同呢?.
可以再談的.
我們是神揀選的牧人.
我既然被你們去尊重.
委派在日內瓦.
你們又願意被這個制度去管理.
所以這個制度.
我們就設立了.
你們要遵行.
因為這個憲章來的.
所以這個教會組織法.
確立之後.
我們的牧羊.
就發現有違規.
就不是當沒事了.
這個本身是對的.
不當沒事是好的.
是每個人都要重視和神的關係.
問題是你怎樣將他帶回神那裡.
嘉義文就用這個方法.
我們做了我們牧羊的責任.
我們就將你的個案轉交市議會.
因為市議會還有一個叫做宗教裁判法庭的.
那個叫Constance.
Constance.
那個就真的是一個仲裁.
就真的去判定.
你是要面對什麼的刑罰.

$^{1121}$這樣就變成了很均真的.
你既然願意接受整個日內瓦.
都是有這個憲章.
這個是教會的組織法.
你又承認你自己是神的兒女.
現在牧者去牧羊你.
發現了你有一些不合聖經的東西.
那你要改變.
但是這個改變.
原來是具體的一個行動.
嘉義文認為是要管.
要管才帶來那個成長.
那個改變的.
所以就有同步的長老執事的.
這個市議會的宗教的裁決.
所以這個社會紀律.
是伴隨著牧羊雙軌並行.
所以如果當時的人都是接受的.
大體上接受的.
當然一定有個別的一些人.
是會很反感.
是會動員帶來挑戰.
這個難免的.
每個人的思想構造.
對紀律的執行.
有些人是很馴服的.
最好你告訴我應該怎樣做.
有些人就不是.
你說一他就說為何要一呢.
這樣不同的想法.
所以你會看到.
其實由那個情況.
1947年一直到55年.
都有這十幾年.
輾轉有幾個我找到的.
其實有四個是重要的.
反對得很厲害.
挑戰日內瓦的市議會.
矛頭是挑戰嘉義文.
這個自由思想家.

$^{1161}$人民主義者.
不懂得讀他的法文.
古埃.
他的情況就是挑戰嘉義文.
接著在他家裡找到一些叛亂.
要推翻市議會.
針對嘉義文來毀謗.
寫一些文章等等.
發現他的文章裡面有些叛國.
所以最終你看到那個畫面.
1547年的7月26日.
這個古埃被市議會執行死刑.
要斬頭.
你可以想像到原來嘉義文頭十年.
他的地位都不是那麼穩固.
他推行的改革大體上大部分人都會遵行.
但是會有一些很極端的人.
崇尚個人的自主自由.
例如這個老師.
他本身是在日內瓦的神學家.
但是又是挑戰嘉義文的神學.
我先看回我的筆記.
又是對聖餐的一些討論.
對嘉義文的預定論的否定.
覺得嘉義文講得太絕對太極端了.
來挑戰寫這個文章去抵擋嘉義文.
所以又是經過市議會的裁決.
被禁止守聖餐驅逐離開教會.
這個就是嘉義文的政教合一的.
當時的日內瓦的城邦.
你可以說是負面新聞.
其實必然的了.
這麼多年二十多年的改革.
只是三至四個人夠膽挑戰.
其實一定不會那麼少.
不過又不可以說當時的民意還沒開啟.
而當時的人都是願意這樣去集體的負責.
去遵從一些條例契約.
所以就會有這樣的小撮的人去試圖挑戰.
你可以想像到嘉義文經過了第一次年輕的時候.

$^{1201}$急進的改革帶來的就是整個議會對他的否定.
很多的普通的日內瓦居民.
都認為嘉義文太過年少氣盛.
就驅逐他離開.
但是這次重回來就心思熟慮了.
就更加成熟了.
就更加做得比較客觀了.
持平了.
但是在這個持平裡面.
他都要確立自己的權威.
所以這個老師.
其實市議會起初是反對嘉義文的.
嘉義文就說我寧願死.
寧願離開日內瓦.
我都不會收回對這個老師.
就是中間這一位神學家的禁制.
一定要驅逐他離開教會.
一定不准他守聖餐.
因為他的行徑是背離了這個宗教改革和聖經的真理.
所以你可以想像到嘉義文是不知不覺間.
樹立了一個很專橫的體制.
因為他要保護整個日內瓦的居民.
或者整個宗教改革的核心.
就是有一個地上的理想國就是日內瓦.
就是這個地方雖然人很少.
但是所有人都是要遵從基督的教義.
某程度上是很好的.
如果整個錦繡堂,整個錦繡花園.
當五千戶都是信主的.
大家一定要守聖經的契約.
一定要跟隨神的僕人的指導來生活.
某程度上是很好的.
不過當時有這樣的人挑戰嘉義文的論述.
他的主張定了的法規.
例如那些稅項,借貸的利率.
全部是議會有規範的.
是一個好像不是自由經濟.
是一個規劃經濟的社會.
最嚴重的就是這個西班牙人.
翻譯應該叫什麼名字呢?中文叫塞維塔斯.

$^{1241}$這個比較嚴重.
就是反對三一論.
神不應該是三一真神.
聖父,聖子,聖靈是不對的.
這樣寫書來反對嘉義文.
反對這些.
這些基本上是不能改的.
由奧古斯丁的系統神學已經確立了.
已經不需要再討論,再爭論的.
所以那個問題的嚴重性.
就是他始終不肯收回自己的言論.
又變成一個異端的妖言惑眾.
宗教裁判所不只是天主教.
基督教,日內瓦都有它自己的宗教裁判所.
所以諷刺地一個這麼溫和的宗教改革家.
嘉義文一個這麼讓人敬重的.
都要執行將塞維塔斯用火燒死.
是對待異端的刑罰.
當時代在嘉義文的宗教改革,基督教裡面.
所以這個塞維塔斯的死.
更加引發人,其他歐洲對嘉義文的否定.
和對他的預定論的挑戰.
這個就是後話.
是下一代的人將他的學說整理成Tulip的問題的爭論.
但是對整個這二十多年的宗教改革.
其實是最後的由55年至到1564年是較為順暢.
頭十多年由1541年回去到1555年.
其實都有很多不是沙石.
是很反感的.
有些是神學家.
有些大部份是一些自由主義者.
即是自由派的人民主義者.
只是崇尚個人的權利,自由的.
亦都有學問的.
會挑戰你們的說法.
聖經這樣說.
但是根據他們的理性.
不應該是這樣規範.
帶來的就是在這個規範底下.
是會有反動和壓迫的反抗.

$^{1281}$所以有些是溫和的.
驅逐出教會.
不准收聖餐.
有些是叛亂就殺頭.
但是神學上的異端就要用火刑.
和天主教一模一樣.
這個我都覺得是匪夷所思.
有了絕對的權力.
因為加爾文1555年之後.
真的有了絕對的權力.
整個城邦都信服於教會組織法.
所以當有人要挑戰.
其實他有這樣的異端.
當然我們今天最多可以做的.
就是寫一篇文章.
去反駁這個A君說沒有三一論.
神是獨一的.
沒有三位的.
沒有聖父,聖子,聖靈的.
只能夠用神學去澄清,去講解.
甚至你用很重的話.
再不悔改會落地獄.
但是你不能夠.
跟手就用一個國家的法律.
去將他抓去坐牢.
或者要行死刑.
早期的摩門教.
這個約瑟史密夫.
的確是被教會定性是異端.
政府就跟手去抓他.
所以他們要逃亡去到鹽湖城.
逃避.
當然他們有其他的問題.
特別多妻的問題.
當時的美國政府是禁止的.
某程度上都是大部分.
當時都是基督徒.
所以政教如果合一.
不知不覺間權力的錯配.
或者模糊.

$^{1321}$就容易導致這些聖靈塗炭.
今天我們不會這樣做.
但是我們的形式.
是不是溫和一些.
還是更加激勵一些.
激進一些.
譬如我們的文章對一個人的定性.
或者我們一些流言蜚語.
對一個人背後的責罵.
是不是更嚴重過一個火刑呢.
這些其實.
不過是不同的處境.
不同的形式的權力.
這個也是值得我們去想的.
當一個人被眾人定性為異端.
我們要怎樣幫他回頭呢.
他始終都不想改變他自己的想法.
他都這樣認定.
有權力的嘉義民就是透過議會.
不是他.
是透過議會.
去將這個人公然的行使火刑.
綁在木柴裡面.
生人,鑑生,燒死.
就不是印度教的.
將遺體盡快燒死.
在恆河裡面灑灰.
是真正一個生人,鑑生的燒死.
因為他是異端.
在那個時代中世紀.
因為政教合一.
就會有這樣的權力.
但是今天.
當教會沒有行使這個權的時候.
而這個權已經收回在地上的政府.
地上的政府也不會再干涉你們的教義.
是異端還是不是異端.
是你們自己教內的討論.
這個討論.
我們當然有一定的權.

$^{1361}$我們會在其他的教會都說.
這個人已經在我們的判斷,討論,對話裡面.
他不肯收回自己的三一論.
反對有聖父,聖子,聖靈.
我們就會讓其他教會都知道.
這個人的論述是錯的.
不能夠接受的.
他是異端.
That's it 是這麼多.
你不可以再有任何的事做.
反而是希望和他多對話.
給他機會去悔改或者回轉.
否則只會慢慢成為陌路人.
可以做的就是這麼多.
所以這些值得我們去思想.
我們五百年後再看這些宗教改革.
很想將他內心的理想國.
在日內瓦.
他不是想做一個試驗場.
他真的很認真地.
讓這個日內瓦成為一個敬拜神的中心.
其他地方他管不了.
譬如法國.
天主教還是很強的勢力.
法國一直都很敵視他自己的同胞嘉爾文.
所以他一世都沒有再回法國了.
在瑞士,日內瓦,中路.
在那裡作為一個當地的牧師.
他不是一個長老.
但是他就帶領著12個長老.
開每周的議會.
談一些關乎宗教的事情.
但是無論如何.
你可以想像到.
在政教合一的16世紀.
當時代的人.
他們看這些事情.
要麼就要跟隨基督教.
要麼就要被驅逐.
或者.

$^{1401}$當時天主教還是沒有改變的.
他內部都沒有真正的重整.
但是世俗的人民主義.
只會越來越高漲.
越來越去到一個.
所以法國的大革命.
就是人民主義的抬頭.
就讓法國的皇帝.
6月14日.
送上了斷頭台.
因為法國根本就沒有做宗教改革.
所以就是人民主義的抬頭之後.
就帶來的反動和革命.
所以宗教改革的餘波.
都在民間不信的人當中.
其實產生一個很大的漣漪的影響.
不過某程度在聖經規範底下的.
加爾文的宗教改革.
就算是真正的溫和.
在溫和裡面.
都被人詬病.
原來加爾文都真的命令.
執行這個火刑.
去燒死這個.
塞維塔斯這位朋友.
這位西班牙人.
所以你可以想像得到.
在整個歐洲的宗教改革.
我們看有些地方是失控的.
的確是失控的.
但是在上主 造物主的管理底下.
容讓人有不同的自由意志.
不同的看法.
但是有其中一派.
在這個聖經的框架底下進行的改革.
就一直延續到今天基督教的發展.
都是受這一派的影響.
就是加爾文的基督教要義底下的宗教改革.
那個的影響.
所以某程度我們能夠這樣.

$^{1441}$今天香港的基督教.
是秉承上一代西教士去到國內.
1949年全世界的宣教士.
在任何省市的.
所有都要撤退來到香港.
香港成為整個宣教基地.
百花齊放的.
尤其是美南浸信會的宣教士.
很敬虔的.
改革中的派別.
直接影響著香港的60,70年代.
50年代,60年代的基督教發展.
因為香港最早期的聖公會.
那些大的教會.
就是英國的國教早期來的.
所以跟香港本地那些很窮的居民.
1949年逃難來香港的.
完全無法募養.
所以是需要一個新的宗教的改革.
1949年變色之後.
解放大陸之後.
所有宣教士被迫要來到香港.
神就用這班宣教士.
引發整個香港的大復興.
就是福音的爆炸.
跟聖公會.
跟傳統的在香港已經過百年的西方教會.
完全做不到.
因為大家的身份背景都完全不同.
所以你可以想像得到.
嘉爾文這個溫和的改革.
其實是慢慢的一個很重要的.
一個穩固的樹幹.
慢慢的一直在生長.
所以下一代約翰·洛克斯.
就帶來英國蘇格蘭的大復興.
一直有條根.
接著洛克斯之後.
就是去到這個衛斯理·約翰.
和當時代的那些改革.

$^{1481}$所以一直延伸下去.
英國的這些敬虔主義.
就因為五月花號.
要逃離聖公會的迫迫.
就去到新的世界美國.
又是這樣爆發一個復興.
接著就去到十九世紀.
去到不同的黑暗大陸.
還有大量的那些原始的人未信主.
所以就一九零幾年前後.
西教士湧進中國.
帶來中國的福音工作.
但是不成功的.
因為那個迫迫是很嚴重的.
就去到一九二幾.
經過了非基督教運動之後.
就開始去到接近五四運動的時候.
就真的有一個大爆發的那個信主.
所以回過頭來.
就是溫和的宗教改革.
加爾文對世界的影響.
是由他這本基督教要義產生.
是由這個人.
一個原創的.
將一些系統神學.
很有條理地寫給.
因為當時原創是沒有這些作品的.
這樣來寫給那個時代.
作為一個模樣.
所以就帶來一個很大的影響.
影響後世.
後世當然一定是要慢慢.
不可以再政教合一.
你不可以有聖經.
然後有把劍這樣來執行.
這樣是不行的.
所以慢慢慢慢就分開.
所以教會今天.
我們只能夠用神的真理.
肯定地勸說.

$^{1521}$不聽的就是不聽神的道.
將來面對神的審判.
只能夠這樣恐嚇.
不可能用劍.
你今天禮拜天不見你.
明天就找議會來抓你回去審查.
就沒有這件事.
變成我們就要這樣去看.
但是其實那個貢獻是.
由基督教要義.
由加爾文他的神學.
奠定了宗教改革.
到今天還是有這條線牽連著.
所以可能我最後用一堂.
就將那個發展.
譬如加爾文最重要的學生.
帶來的貢獻就是John Gluck.
是很重要的一個報道家.
奮卿家 神學家.
帶來蘇格蘭一個很大的改變.
所以你可以想像到.
一代一代的人受著.
宗教改革的影響.
和那個幫助.
今天我就講到這裡.
下一次就再聚焦在.
加爾文提出的他們的神學.
特別是五個唯獨.
和我們都簡單看看.
Tulip大比當時.
因為長老中受著這個影響.
我們一起簡單禱告.
感恩天父給我們.
透過思想加爾文.
我們好像聽一些趣味故事.
很有趣.
又覺得很奇怪.
為什麼基督教都會執行死刑.
去將一些異端.
用火來燒死.

$^{1561}$我們很難回到古代.
去看這些事情.
可能古代的人覺得很理所當然.
因為政和教是沒有分開的.
我們懇求主幫助我們.
讓我們能夠看到.
其實牧者對信徒的了解.
禱告 代禱 慰問和關心.
這個是神給教會的一個很重要的恩賜.
和聖職.
同時也給香港的教會.
不同的牧者.
都能夠履行這一份神聖的職份.
真是去牧養.
將人帶到神的國度.
使人對永恆有意識.
自己的選擇.
在神的感動裡面進行.
但是當我們看到有這樣的瑕疵.
被後世的人來救病的嘉爾文.
又給我們有警惕.
當我們真是有所謂的權的時候.
幫助我們用愛.
用恩典去服侍.
將人帶到主耶穌裡面.
願神幫助.
我們知道保羅都曾經說要趕一些.
在哥林多前書第五章.
取繼母的人要趕他出教會.
不准他吃聖餐.
在當時都有這樣的紀律執行.
其實教會都有的.
盼望神都給我們看到這些事情.
我們一起在主的恩典裡面.
去守望不同的弟兄姊妹.
願神幫助我們.
再一次我們獻上感恩.
奉耶穌基督的命.
阿們.
\newpage



\section{}
\label{sec:3Sqlgs98HDM}
\textbf{2025-01-09 周四晚查經班《改革之路(7):加爾文(三)》 講員:曾裔貴牧師}
\newline
\newline
連結: \href{https://youtube.com/watch?v=3Sqlgs98HDM}{\texttt{https://youtube.com/watch?v=3Sqlgs98HDM}} ~~~~ 語音日期: 2025-01-10
\newline
\newline
\hyperref[sec:4EM8nFC2gtc]{\small{< < < PREV SERMON < < <}}
~
\hyperref[sec:index_chronic]{\small{[返順時目]}}
~
\hyperref[sec:index_scriptual]{\small{[返順卷目]}}
~
\hyperref[sec:Dn07Lw4F37U]{\small{> > > NEXT SERMON > > >}}
\newline
\newline
$^{1}$好 我們一起又祈禱.
多謝天父讓我們又可以來到第七講.
關於這個嘉爾文.
今晚我們會想思想一下他的神學的主張.
和對當時這個宗教改革所帶來的影響.
我們都盼望我們吸取到裡面一些很重要的原則.
和方向使我們能夠對屬靈的事情有一個認識.
盼望我們都能夠在整個過去宗教改革.
和一路的發展.
我們一方面明白了一些歷史之後.
我們自己都能夠秉承著前人他們一些的思想.
他們的主張來去學習.
和我們真是有一個對自己的信仰的深化的認識.
我們今天晚上都等候在神面前.
我們這樣禱告奉耶穌基督的名.
阿們.
上一次我們提到這個宗教的組織法.
和因為它排除了一些反對的人之後.
很自然地就在日內瓦這個城邦.
或者當時的國家瑞士這裡.
就來到可以很順利開展它一些神學的主張.
它從來都沒有想過要將天堂,神國.
要付諸於實行在瑞士,日內瓦.
日內瓦就成為了人間的天堂.
它沒有這樣想的.
不過一直看一些不同的文章和一些著作.
特別它自己的基督教要義.
它提到地上的政府是有一個很重要的責任.
特別因為當時日內瓦全部都是基督教信主.
他們也是在1541年的時候.
極力地邀請加爾文回去帶領他們.
大概三萬多人.
你可以想像得到.
既然大家都承諾要遵守聖經.
所以就在整個日內瓦就制定了這些公民的約章.
和在社會上去實踐基督教的教義真理.
所以很自然地它會有很多地方.
它會有一個管理.
稍後我們會說.
會分成兩級的教會的管理.

$^{41}$就是由牧師和傳道人.
地上的治理就是由長老和這些的執事.
這樣去分開去治理的.
當然和我們今天一定不同.
因為我們不是整個錦繡 整個元朗是基督教.
我們只能夠在個別的教會裡面去實踐這個治理.
所以它現在就擴大了整個日內瓦這個地方.
成為了好像一個大的教會那樣.
在整個城邦裡面去推行教會的一些聖經的教導.
所以連那些比如借貸要多少利息.
他們都定了不可以高過某一個數.
還有如果是借給一些窮人.
那個利息又不可以好像借給一些做生意的人那樣.
政府當時嘉爾文和他們那個日內瓦的長老團.
都有制定這些稅制的責任和權威.
所以你可以想像得到等同於整個城邦.
都是要實踐基督教的教義.
所以我們上一次最後講的就是.
有一位神學家其實某程度上和嘉爾文是熟的.
英文就是Swertis 塞維塔斯.
他原先在法國都被天主教驅逐.
說他不相信三一神論.
就是聖父聖子聖靈同尊同榮同等.
今天我們都大概知道.
我們甚至乎可能有些錯信的都不知道什麼叫三一論.
你不會因為這樣就要將他驅逐出教會.
但是當時法國的天主教就將他驅逐.
因為他不肯否認他不信三一神論.
就是三位一體.
所以輾轉輾轉他就在歐洲流亡.
因為我再看回他的資料.
後來最終就在日內瓦被人發現了他出現在日內瓦.
在嘉爾文那個核區那個牧區那個國家.
嘉爾文的那個宗教法宗教治理.
人都知道他是不相信三一論的.
或者否定三一論的.
就是不相信三位一體的.
所以就被當然好像是捉拿.
來去召見.
嘉爾文給他一個機會.

$^{81}$你否定你所主張的.
你承認三一論.
承認聖父聖子聖靈.
其實有多難呢.
但是他又忠於他自己那個信仰神學.
怎麼都不肯承認.
當時宗教法庭.
因為日內瓦都有宗教法庭.
天主教也有.
日內瓦自己也有.
其實只是長老去執行哪些人.
例如他不守聖餐.
無緣無故不守聖餐.
不去教會.
長老就要代表教會.
來對他們嘉爾的審問.
懲罰.
但是這次很嚴重的神學偏差.
所以嘉爾文就給他機會去悔改.
他堅決不肯否定他自己的信念.
所以嘉爾文就和那個宗教法庭說.
日內瓦的宗教法庭.
不是天主教.
是宗教裁判處.
就叫他收回他不肯.
嘉爾文就說.
現在他需要受懲罰.
不過用一些比較文明一點.
溫柔一點的懲罰.
意味著最高的懲罰就是用火刑燒死.
天主教就一直都做的.
沒有異議的.
因為誰不信他們的神學.
就要被燒死.
但是現在連基督教.
連宗教改革的嘉爾文的神學.
因為在日內瓦他們是得到那個權.
最後問他.
他說我都不肯否定.
我堅決不相信有三一論.

$^{121}$嘉爾文就請求長老團不要行火刑.
但最終長老都否決了嘉爾文的建議.
將這個塞維塔斯行火刑燒死.
這件事就成為日後攻擊嘉爾文主義的人.
後世的神學家,歷史學家.
對嘉爾文的一個批評.
因為你自己反對天主教.
就將你們基督徒捉拿去行火刑,審判.
但現在你們基督教.
因為這些似乎這麼嚴重的神學立場.
就對一個神學家去做一些.
不是基於一些殺人放火.
就行這些火刑.
這件事就成為了歷史的污點.
到底是因為教會和社會.
是一個緊密的地上的理想國.
還是怎麼樣的情況.
而導致有這些實權去做這樣的刑罰呢?.
這個我們後世都要想的.
今天我們就覺得是天馬行空.
很匪夷所思的.
怎麼可能教會會行使這個權呢?.
已經沒有了.
大家要留意一下.
如果大家有那個開了聲呢.
可能你要關一關.
不然就怕影響其他人聽.
那個一路的分享.
這個就是上一次我們講過的.
現在我們集中講回嘉爾文神學.
幾樣東西是很重要的.
一方面他一定是很著重神學.
神學尤其是神學和聖經的結合.
就是說系統神學和聖經神學的融為一體.
我今天就忘記帶那兩本基督教要義.
但是你看基督教要義.
一來就是分類的.
系統神學的分類.
由神觀一路去到社會的治理都有的.
因為他真的是一個城邦.

$^{161}$所以他是有關於社會的秩序的.
所以挺好看的.
挺容易看的.
但是你會看到.
穿插在他的神學裡面.
是有大量的聖經來支撐他的主張的.
所以你可以想像一個人.
用了二十多年的時間.
他就已經將一個這麼有系統的神學.
來到鉅細無遺地寫出來.
然後分類.
還有賦予很多聖經的經文的支持.
所以你看到某一個topic.
他就會有經文去支持的.
所以你可以想像到.
嘉爾文對聖經的通透.
那個認識.
是很深入的.
當然說到關於聖靈.
這個就是嘉爾文的聖參觀.
就在那裡可以看到.
如果可以的話.
我們就進入聖經的神學裡面.
嘉爾文的神學.
有哪些重要的主張呢.
大家知道.
馬丁路德是神父.
是天主教本身的神父.
所以他的宗教改革.
純粹.
這是第一批.
馬丁路德純粹是很想教廷.
不要再用贖罪券.
去影響了平民百姓.
那個錯誤的對贖罪的觀念.
和其他教廷內部的一些腐化.
有一些的改變.
所以馬丁路德一直不是想著要革命的.
一直想著希望教廷就著他列舉了95條的政綱.
來大家討論辯論.

$^{201}$理性的辯論.
但是天主教是鐵板一塊.
還有當時的印刷術.
當時已經很多人是不滿意天主教的.
所以就借馬丁路德來一個很好的契機.
就將他95條.
你釘在那裡.
我們就撕下來.
大量印給歐洲的不同的地方.
所以就引發那個革命.
引發大家大量的人.
開始對天主教產生不滿.
和最終經過神學家的一路的教導.
他們就脫離天主教.
所以馬丁路德從來沒有想過在一些崇拜.
禮儀.
聖禮上要來有一些改變.
所以相反上一次我們說過.
嘉爾文是第二批第二代的神學家.
和這些宗教改革.
所以他有很多的時間.
看著前人是怎樣來革命改革.
而他自己有一個很大的時間.
他在年輕的階段.
其實他在法國的巴黎大學讀神學.
他又接觸很多古典的.
比如奧古斯丁的典籍.
中世紀時代的.
比如托馬斯.
托馬斯.
糟了,我忘記了神學家的名字.
這樣來到有機會去讀.
再看馬丁路德這一班第一代的神學家.
怎樣來跟天主教的論爭.
所以他自己又沉澱了.
更重要的是嘉爾文有很多機會.
在他年輕的時候.
他已經對希伯來文相當熟悉.
所以他不需要有掙扎.
要脫離天主教.

$^{241}$他本身就很認同基督教的這個改革.
所以他就在根本性上.
看到天主教的崇拜和勝利.
是有很大的落差問題.
所以你可以想像得到.
在他的一個.
多麼的站身.
因為你反對別人.
那到底你有什麼自己可以建立呢.
他就在這幾方面.
真的可以說是很原創性.
馬丁路德都沒有這樣的.
大家的成長都不同.
其中就是重視港壇.
你想想在當時的市議會.
日內瓦的市議會.
他們通過了1549年.
1549年其實是一個分水嶺.
因為之前在他.
由1541到1549.
其實還有很多反對的聲音.
在日內瓦.
上一次我們講過.
不單是天主教是他的敵人.
還有一些人民主義.
所謂自由派.
那些人其實很崇尚個人的選擇權.
個人的自主權.
所以那些人未必是信耶穌的.
或者未必是那麼commit(承諾).
來去很想成為那個改革的.
但是無論如何.
他們一定不會喜歡天主教.
所以他們就在兩者中間.
所以他們當在日內瓦的時候.
很多時候就成為嘉爾文的一些敵人.
他推行的東西是非常之嚴格.
規定了.
例如你不去崇拜.
不去聖餐.

$^{281}$就派長老去調查.
為什麼你不去聖餐.
為什麼你不去崇拜.
就帶來其實很大的壓力.
所以有些人民主義的人.
有些甚至乎都是有在市議會裡面的.
都極力反對嘉爾文的改革.
所以去到1549年.
那些起初.
這個就是我們上星期講過的.
我不懂得讀那個名字.
是類似法國的名字.
這個人最終被遊街示眾.
因為反對嘉爾文.
後來因為不斷攻擊他.
後來可能找到他有些問題.
要將他去遊街示眾.
好像背著牌來打壓.
你想像得到這樣的情況.
當然可能是有他的長老團.
最後得勢支持嘉爾文.
所以這裡告訴我們.
原來當時市議會.
是一直都通過了.
是牧者.
這裡有分兩種的.
牧師和傳道人.
牧師就是要固定的.
來探望每一個家庭.
了解家庭的屬靈狀況.
和有什麼需要.
如果發現到原來家裡藏著.
看邪術的書.
或者其他的一些不合真理的東西.
就要告訴長老團.
市議會就要來調查.
要執行 要來跟進.
所以你可以說.
在日內瓦的生活其實都很嚴格.
所以不是每個人都能夠接受到.

$^{321}$但最終嘉爾文終於有這樣的主張.
這個和他之後的神學很有關連.
或者這個點大家記住.
是影響著日後16世紀.
無論是蘇格蘭的大復興.
英國之後就去到美洲.
因為發現新大陸.
當然就是五月花號和輾轉.
英國的清教徒就移民去美國.
就帶來17世紀的美國的大復興.
這裡就是一個很重要的.
你可以說是一個種子.
他認為人就算是信耶穌.
都需要對神的真理明白.
明白之後要實踐.
要在生活裡面去應用.
這樣才能夠全面的.
那個屬靈的生命會成長.
所以不是一種只是掛口號式的一個基督徒.
而是牧者是有那個責任.
將神的真理來去後設的講解.
所以你看到第三個下面那裡.
這個講道唯一的合法的方式.
就是按那個系統的《釋經》講道.
即是有系統的續章續卷這樣去講解.
就不要一個topic這樣.
來到這裡抽一折那裡抽一折.
那就是他們不能夠接受的.
這個是嘉爾文的想法.
所以在當時他在這個市議會.
已經得到那個絕對的勢力之後.
他就正式推行牧者必須每日講道.
可能不同的地方.
大大小小的講壇.
或者去到家庭裡面去講道.
所以很忙碌的.
星期早午晚不用講了.
就一定是會有.
其他每日都在不同的地方來講道.
重視講壇就由這一刻開始.

$^{361}$在日內瓦成為一個示範單位.
輾轉去到歐洲.
又是在這裡出現.
那你會看到.
你很快就已經和天主教當時做比較.
因為稍後我們就講天主教的崇拜觀.
這個聖禮其實對於聖經.
近乎是不會探討的.
近乎是不會有那些講經的時間和聚會的.
所以變成很多的天主教徒.
其實對聖經完全是不認識的.
所以在這樣比較之下.
嘉爾文的這個第二代的宗教改革.
就尊重聖經.
體重講壇.
和來到去將.
那個信念就是.
神在地上.
這個就和稍後的天主教的那個傳承系統.
就是由教皇領受了那個權柄.
傳遞給每一個世代的那個主教.
大主教.
跟著有這個權.
再分受下去.
但是嘉爾文就不是這樣的.
完全將這件事打破了.
只需要那個牧者.
是清楚神的呼召.
他就是在地上.
是代表神.
不過是在真理的教道上代表神.
不是那個權就是成為了教皇.
兩者是不同的.
就是說牧者仍然是一個普通的人.
但是當他領受了這個職份.
他就在這個職份上.
就要將神的道.
講解明白給他的會眾知道.
如果那個會眾是有一些東西不認識.
而將來面對神.

$^{401}$他真的不認識沒有聽過.
那麼那個牧者就要有責任.
就受到神的那個責罰了.
那麼這個就是嘉爾文當時.
他的信念.
所以他是.
其實你想一下他多忙啊.
還要翻譯聖經.
還要每天講道.
還要去探訪不同家庭.
還要長老團的十二個長老和牧者.
是要每個禮拜開會.
談一下那些日內瓦的公務.
今次.
今年我們要加多少稅.
或者那個利率要調整多少.
他們要開會.
你想一下.
好像我們要去開行政會議.
我們要開立法會會議這樣.
談這些九巴加多少錢.
船加多少.
的士加多少價跳標.
多少錢跳標.
要談這些東西.
當時因為整個城邦都是基督徒.
整個城邦都是基督教.
所以連這些事情.
都要成為他們要談的事項.
所以你可以想像得到.
但是有一樣東西就是很重要的種子.
就是十六,十七世紀的敬虔主義的開始.
就是因為從這一刻開始.
不同天主教.
講壇是很重要.
講神的道.
很認真地.
有系統地將神的道講明白.
譬如我們剛剛講完末世論.
講完災前,災中.

$^{441}$或者前千禧年,後千禧年.
或者沒有前千禧年是怎樣.
我們至少都揭過一些聖經.
知道聖經大概有這樣的講解.
或者我們講啟示錄.
或者我們講羅馬書.
我們叫做有系統地看過整本書.
大概它的主題.
或者每一章大概講什麼.
我們都有一些零零星星的印象.
但是嘉爾文就在這裡.
成為了一個很重要的里程碑.
能夠將聖經這個權威.
不只是放在這裡給人看.
而是要深入,淺出地講解,明白那個真理.
所以這件事到今天的華人教會.
到今天的福音派.
講道都仍然是崇拜的中心.
沒得假的.
所以如果崇拜有一個小時.
如果只有五分鐘是講道.
在嘉爾文的眼裡.
你基本上根本就不是崇拜.
簡直是喧賓奪主.
所以你可以想像得到.
當時他們就是看重這個講道.
所以接著還要是有系統的釋經講道.
所以他認為唯一的方言.
就是對希臘文,希伯來文的掌握.
將經文有系統地講解,清楚,明白.
你想想不再用拉丁文.
用自己能夠明白的語言.
然後更重要就是將原文講解清楚.
給回眾知道.
當然他們就分了兩個很不同的界別.
天主教就是教士.
就是這些priests.
就是神職人員.
我一會兒會講天主教那裡.
特別是提到其中一個聖禮.

$^{481}$就是講到有三個很重要的神職的職份.
就是主教,司德和那個.
我遺漏了,我一會兒再揭開.
下一頁應該有的.
但是嘉爾文就清楚了.
我今天因為預備到這裡.
我就揭開基督教要義的呼召.
他講得很詳細.
Pastor一定要有神的呼召.
而且那個呼召只不過是神話語的管家.
將神的真理去講解.
所以那個職份不是一個好像天主教的.
因為當時就天主教來做改革.
所以就真的某程度將那個聖職.
將那個主教,教皇這種權威.
來拉低.
跟普通人一樣.
不過是特別有神的呼召去做這個牧養.
傳福音和來教導.
如果是傳道者就很清楚了.
是一個教師,一個teacher.
不過是教聖經的真理.
所以牧者當然也有一個重要的聖禮.
就是聖餐和洗禮.
這些都是牧者要做的.
還有主持婚禮和主持喪禮等等.
所以你會看到.
在他的聖禮觀裡面.
那個最重要的核心.
都是神的道.
所以嘉爾文他自己整個心靈.
他自己的神學都是建基於聖經神學.
而且他一直都是一脈相承.
秉承著奧古斯丁的神學觀.
所以他沒有什麼大的改變.
所以那個神論,基督論,救恩論和聖經論.
都是秉承著過去的那些神學家.
已經有的基礎來去建構上去.
所以你可以說嘉爾文唯一.
最重要的就是對聖經的重視.

$^{521}$還有一個很重要的信念.
就是牧者的存在.
就是要在你的轄區.
轄區就是今天就是教會.
我們不可能在元朗大街叫人守聖經.
你不對,這樣不行,這樣我們不可能.
因為已經分了兩樣東西.
不過我們就要在我們那個轄區的教會.
讓弟兄姊妹明白聖經.
和遵行聖經的真理.
那他自然就會經歷神的同在,恩典.
這個第一方面在崇拜的聖禮上的更改.
我們開始看天主教.
因為天主教的聖禮觀就會是很不同的.
你可以想像得到否認彌撒.
因為彌撒原來不單只是聖餐.
彌撒原來是有為死去的人祈禱的.
他們覺得是彼此的一個協助,幫助.
但是因為他們又賦予一些他們自己的信念.
或者天主教自己的一些神學.
所以很多時候會呼喚聖母.
或者一些逝去的聖人來祈禱.
所以彌撒某程度是某一種的超度.
如果用我們香港的想法.
就有這種含義.
所以你可以說.
加爾文對這件事是極之否定的.
完全沒有這件事.
他認為公共的崇拜只有一個很重要的目的.
一個教會的存在.
就是要讓信主的人明白聖經的真理.
明白神的心意.
和在生活裡面去進行神的心意.
那就是一個基督徒成長的階梯.
所以他最重視的就是兩件事.
神的道的宣講.
有系統地宣講.
還有聖餐的執行.
所以只有兩個.
你可以說是明顯的兩個聖禮觀.

$^{561}$就是這兩個.
但是天主教就不是了.
七聖禮.
可能大家都聽過了.
聖洗聖事.
其實都是英文Baptism.
原來這裡有考究的.
因為這裡表面上和我們一樣洗禮.
原來不是.
Baby一出生都要洗.
但是後來我就找了近代香港.
不是香港.
好像是美國的天主教一些洗禮的申請表.
去到五歲.
你就要和他行堅振禮.
第二個已經講了.
這個confirmation 堅振.
所以不是要到十八歲.
原來Baby洗禮之後.
五歲就要開始.
因為開始懂事.
五歲就要申請入那些學道班.
聽明那些聖經.
然後才confirm他是堅振.
所以原來頭這三個是一個package.
洗禮.
堅振.
和所謂聖體聖事.
我不懂拼音.
其實就是聖餐.
不過天主教的變質說的聖餐.
這個是一個package.
三個.
所以如果你只是做了洗禮.
然後從來你長大之後.
你都不confirm.
其實你的生命不完整.
他們認為你是不完整的.
還有你如果不守聖餐.
你就更加經歷不了神.

$^{601}$當然這些某程度是對的.
我們都認同的.
不過我們就不是這樣的方式.
這個就不只是普通的成人.
成人當然簡單.
你一定要去初信班.
不少過一年.
我看過他們的申請表.
不少過一年.
原來都很嚴格.
成人都可以洗禮的.
當然有些成人信是跟進天主教的.
但更重要.
天主教當時更加重要的.
就是一定要BB馬上洗禮.
因為怕他們夭折.
就不可以了.
因為洗禮加上這個堅振禮.
就已經洗去你的原罪了.
他們有這個信念.
但是我們認為.
洗禮不是這樣的意思.
我們是跟救主耶穌的生命認同.
我們在祂的救恩裡.
我們就得到聖靈.
我們的生命就其實改變了.
但我們仍然保留我們這個原罪.
我們沒有改變.
但不過我們在神的聖靈裡.
我們是不斷成長.
不斷成聖的.
所以這個套餐.
所以他們一定要確認.
確認本身不是錯的.
我們每一次聚會.
都應該對自己的信仰有確認.
所以他們現在不需要這樣說.
要多少歲才要堅振.
其實BB五歲之後.
就要開始對自己的信仰有基本的認識.

$^{641}$但是我們就沒有這個分拆.
所以我們為什麼一定要十六歲.
當然不知道這是不是因為當時.
香港因為區聯會定的.
當時這個十六歲是開始.
接近成年的階段.
所以一個比較籠統的計算.
沒有科學根據.
可能現在十二歲都可以.
不過區聯會仍然是十六歲.
我們就跟著這樣.
其實沒有一個劃定.
十二歲就不可以信耶穌.
又不需要這樣想.
但是我們只能夠是一個方程式.
不用過三關.
我們只是洗禮.
我們確認其實每天.
你自己都要更新自己的靈明.
這個是對的.
但就不是用一個儀式去做一個堅振.
但是聖餐我們是可以告訴我們.
我們每一次.
大家一班人尤其是在聚會裡面.
所以嘉爾文.
我很認同這件事.
嘉爾文認為如果你自己守聖餐.
其實是罪的.
因為這都是教會的聖公.
教會的聚集.
所以不可以說.
今天戚戌徒根很有mood(氣氛).
自己守聖餐.
叫明仔一起在客廳裡面守聖餐.
其實嚴格是不行的.
嘉爾文認為是不可以的.
一定要在教會的setting裡面做的.
是由牧者去帶領做的.
這個是嘉爾文的.
某程度今天.

$^{681}$我都不覺得有太多人.
夠膽自己經常這樣做.
我覺得少.
當然有些私人的我不知道.
背後什麼原因我不知道.
反過來我們為了一些牧羊的原因.
我們去到一些老人院.
因為他已經去不到教會.
我們就有責任去派餐.
那個仍然是教會的一個延伸.
但其實家裡就不算是教會的延伸.
至少嘉爾文是這樣理解.
所以如果你在那個時代.
就捉到你自己守聖餐.
那就要去宗教法庭.
那就又要調查.
所以當時是不會有人這樣做的.
所以嘉爾文就將這件事簡化了.
只有洗禮和聖餐.
才是兩個重要的聖禮.
還有崇拜絕對不是一個彌撒.
絕對是我們一起敬拜神.
所以詩歌都是很少.
不是太多.
重點仍然是那個牧者領受神的道.
他代表神去講道.
將神的道講解.
所以那一刻是嚴肅.
因為他是代表著神去講解那個真理.
然後大家一起學習.
一起在生活裡面去經歷.
所以這裡是將七個濃縮成兩個.
我們今日福音派都只是這兩個聖禮.
算為神聖的聖禮.
洗禮和聖餐.
這裡我們之前講過有幾個不同的.
其實應該有四個關於聖餐的.
我們下一頁就會再講.
我先講七個聖禮.
剛才三個是一個套餐.

$^{721}$記住洗禮 堅振禮.
小朋友和大人都要堅振.
還有聖餐.
即是說如果你跟著那個步驟次序.
你洗禮 跟著你行了堅振禮.
跟著你又固定的去守聖餐.
聖餐即是彌撒.
你去這樣做.
你就會經歷神的生命的聯合和好.
這個某程度上和我們都是一樣.
不過在那個做法上.
有時會令到天主教的教徒會有一些誤解.
以為在那些個人的行為.
多於主耶穌的十個犧牲.
這個就reconciliation和好.
就明顯了.
就是透過與神和好.
我們一向都在做的.
我們有什麼事想起.
就去省察 向神祈禱 向神認罪.
不過福音派或者宗教改革.
是強調個人和神的關係.
但這裡其實到今天都是.
就是原來告解.
為什麼你要去向神父告解呢.
因為原來.
這裡就和他們那個教皇的權威有關連.
因為彼得就是第一任教皇.
是神已經將那個權給了教皇彼得.
彼得也都將這個權一直傳遞到今天.
都是由主教領受了這個權柄.
而主教就分給每一個地方的司德.
所以那個神父.
就是代表著那個地方成為某程度的主教.
所以你麻煩去元朗的天主教堂.
來做告解 扮告解.
當然你可以自己.
但更好就是你要告訴.
是你處理好了 你已經完成了.
你就去告訴神父.

$^{761}$你做過什麼 越詳細就越好.
講你自己的.
你的心裡想的事.
你實際做過的事.
昨天不經意在百佳.
感受百佳.
見到有一個.
有一個有包煙.
拿了 袋了.
然後又忘記給錢.
然後想了想 還是不要給了.
這麼噁心 這麼難看.
那就不拿回去了.
你知道之後不對了.
你就認錯了.
你回去告訴神父.
我做了這件事.
這些還是簡單的.
因為這些很小的東西.
其實是有點貪心.
或者不方便.
但有些可能思想 心靈的東西.
其實你怎麼告訴一個人呢.
但是他認為教皇.
分了那個權給神父.
所以這個和好.
不單止是人和人的和好.
都是透過你看看那個主教.
來跟他說.
主教就有權了.
你這件事罰你什麼.
有時候大家看電影都會見過.
罰你要背多少節金句.
玫瑰經 念多少遍.
類似這樣.
本來那個意思.
和好是好的.
但是有時經由了人的代理.
慢慢不知不覺.
當然我們說回中世紀時代.

$^{801}$是馬丁路德和嘉爾文的時代.
你想想神父的權是多麼的大.
所以這裡如果打破了.
其實帶來的.
其實是很大的民間的釋放.
很大的將那個代理權.
再一次回到耶穌基督那裡.
所以每一個信徒.
宗教改革.
他們回到與神的相交.
和主耶穌的生命的聯合.
所以很多時候.
其實真是的.
我們要這樣信.
這種神跡性的聖靈工作.
很多時候在聖餐的聚會.
你安靜 慢慢沉澱.
真的去默想 去思想.
真的想起自己做過一些.
譬如說過一些不好聽的話.
有些不好的思想念頭.
或者和哪個人衝突過.
其實那一刻就是你沉澱.
你就是去到主耶穌的.
雖然不是告解亭.
其實是去到主耶穌的私恩寶座.
尋求那個赦免.
那個和好.
那個聯合.
那個恢復.
然後剛剛好.
走到去跟著跟勵士牧師拿聖餐.
不是他 他在台上.
應該是鄧全道拿那個聖餐餅.
跟著楊姑娘拿那個聖餐杯.
其實那個時候.
你走的時候.
你已經搞定了你自己的心靈.
那個認罪.
所以這種和好.

$^{841}$某程度上不要向人.
是要向神.
因為這個其實.
人是很微妙的事.
因為那個神父聽著聽著.
可能就變成一個更是非.
更麻煩的人.
所以就未必是最好的.
正正因為他們就有這個錯誤的神學.
由教皇傳遞給每一個地方的主教.
而主教因為他只是一個人.
他變成要分給其他的師德.
那就變成當地的神父.
就有這個責任.
這個是聖事.
是一個聖禮.
七個聖禮其中一個.
就是說你如果不做.
你是不完整的.
你是天主教的生命.
所以都令到很多人都怕這件事.
所以有時候你會看電影.
就是那些黑社會黑幫.
去告解.
你很難想像有這樣的事情.
接著就是這個病人的附尤.
現在原來他們改了.
不需要臨終趕去都可以.
原來可以.
他們可能學了基督教都不一定.
因為他們都是引雅各書的.
很特別的.
但是他們慢慢就擁有了.
原來應該是1962年.
我今天看了一些文件.
所以我自己都有update了.
他們新的教義.
原來是1962年的梵蒂岡第二次的大會.
在梵蒂岡開的.
是全世界的天主教開的.

$^{881}$他們終於定性了.
不再需要在臨終去抹油了.
就是說平時探病都可以抹油了.
甚至乎如果有很特殊的情況.
信徒都可以抹油.
不用等神父趕去醫院.
趕去病床.
這個某程度我們現在都會是這樣.
我們不是假設.
是試過的.
楊姑娘打電話來.
現在我在醫院.
剛好有個病人一定要洗禮.
現在來流了.
牧師我可不可以幫他洗禮.
因為這麼特殊的情況.
我都去不了.
我們都談過的.
教會談過的.
是容許的.
區聯會都准許的.
所以他告訴我們.
我立刻授權給楊姑娘.
她可以幫他洗禮.
這是合法的.
所以這個都是某程度.
是不是他們學了我們.
天主教學了我們呢.
因為當時如果沒有神父趕去醫院.
沒有了這個臨終的抹油.
似乎是有一個很嚴重的後果.
但是現在他們沒有了.
改了.
所以某程度都從善如流.
當然第六個所謂的Holy Order.
和你們沒有關係.
他們說的是在眾信徒裡面.
他們有某程度的儀式.
所以不是經常行的.
其實就是按立一些主教.

$^{921}$師道和一些.
我先看一看.
執事.
他們這個叫做Holy Order.
某程度是我們弟兄姊妹去讀神學.
他將來成為傳道人.
他畢業了.
我們去參觀觀禮.
接著他入職了.
後來輾轉可能五年後.
經過查問.
他可以按目了.
類似這些.
所以和基督教都近似的.
不過我們就不會看為七聖禮這麼嚴重的東西.
不過這個和平信徒沒有關聯.
最後婚姻都看為一個聖禮.
所以七個聖禮.
是天主教一生一定要做的.
要堅守的.
但是基督教只有兩個聖禮.
我們現在看關於聖餐.
因為剛才說的.
剛才說的就是嬰兒洗禮.
見陣禮.
和聖餐.
那個package(套餐).
是要你不斷去做.
不斷去做的.
不過記住.
和我們一樣.
洗禮和見陣禮.
一生只用走一次.
但是聖餐就是每個禮拜.
但是聖餐.
他們仍然都是這樣.
應該就沒有變了.
仍然都是這個transubstantiation.
變則說.
就是說主耶穌的身體.

$^{961}$神聖的臨在.
那個餅和杯.
所以是變了質地.
所以餅就是肉.
那個杯葡萄酒.
是葡萄酒.
就是變成血.
所以是餅.
肉和血的一個聖餐.
明明是沒有變.
何來變呢.
所以馬丁路德就改了一點點.
是那個叫will presence(真實臨在).
一個真實的臨在.
所以那個餅.
那個杯.
已經賦予了一個真實的臨在.
就是主耶穌.
真實的臨在在那個餅和杯.
所以你看到的都是那個餅.
某程度上就改了一點點.
人性化一點.
因為天主教就不理那麼多.
總之你現在咬的餅就是肉.
你不覺得是也是.
你明明說是餅.
你騙我.
都是肉.
那個杯就是血.
但是馬丁路德就知道很多人反對.
明明都不是.
原來是will presence(真實臨在).
真實的臨在.
所以你吃的時候.
那個是肉.
那個是血.
不過那個質地就沒有變.
是同質.
所以這個同質.
consubstantiation(共存性).

$^{1001}$所以你可以這麼說.
兩者兼收.
又是變質.
又是那個餅和那個杯.
那嘉爾文和慈雲尼就更加清清楚楚.
慈雲尼我們今天不說了.
慈雲尼就是我們今天說的紀念說.
就是說我們不要理那個餅和杯.
我們紀念主基督在十誡上犧牲這個歷史的事實.
所以我們只不過借這個餅和這個杯.
一班人一班人很認真地聚集.
思想耶穌的救恩.
所以我們只是借那個餅和那個杯.
作為一個紀念的事情.
所以不是那個餅那個杯給我們有屬靈的生命意義.
而是這個聚會.
耶穌基督叫我們去回想救恩.
而帶來的那個生命意義.
嘉爾文在這一點上有第四個主張.
這個就叫做spiritual presence(靈性臨在).
他就加了一樣東西.
我們回頭看他的神學裡面.
其中這裡神學聖經和聖靈.
聖靈就在這裡顯出他的重要性.
就是這裡的重要性.
嘉爾文這一點很重要.
就是他說這個紀念.
沒錯.
那個餅和那個杯的確是一個表記之物.
只是我們去回想.
帶我們回想耶穌設立聖餐的時候.
和他在十字架上的犧牲.
但是嘉爾文說不是啊.
再加一點點.
當然不是變質.
當然也不是同質.
不是一個real presence(真實臨在).
real就是真實的臨在.
所以這裡又是一個創建.
是spiritual presence(靈性臨在).

$^{1041}$spiritual presence就沒得假了.
我們就應該要崇尚這件事.
所以我自己是很認真在聚會裡面.
等候聖靈的spiritual的恩典和臨在.
不是耶穌的親自真實臨在.
反而是聖靈在聚會.
在我們這班人.
無論那班人是一百人.
一千人.
或者是譬如晚堂只有五十人.
最多的時候就八十人.
八十人和三百人是一樣的.
那班人在這一刻對十字架的救恩的默想.
那個虔誠守聖餐.
這個spiritual presence.
嘉爾文這一點上就是一個很大的突破.
突破了那個變質.
突破了馬丁路德的那個真實的臨在.
更加慈雲就比較很撇脫.
因為他要否定天主教.
總之就是什麼都不用想了.
不要想那個餅.
不要想那個杯.
純粹是紀念就行了.
嘉爾文就是紀念裡面.
原來不只是靠我們的腦.
百般明想東西.
原來是有聖靈的同在.
所以聖餐有聖靈的同在.
所以某程度上就不可以私人自己在屠宰.
自己在守聖餐.
每晚自己守.
這樣是沒有用的.
不會令到你吃得多就會有幫助的.
反過來你去到一個聚會.
你要用那個聖靈.
即我們裡面的聖靈.
去感覺一下那個聚會.
你是不是感覺到一份平安.
感覺到一份寧靜.

$^{1081}$感覺到一份聖靈的同在.
帶給我們的靈感.
這些就大家一起學.
不是我們要神秘地追求.
什麼時候來呀.
到了沒有呀.
不用這樣想的.
是你整個人很心靈平安.
當然.
牧者說我們現在要開始要審查自己.
是需要的.
如果不是.
聖靈的臨在都不關你的事.
因為你根本沒有審查你自己.
所以.
這個紀念說再加spiritual presence.
這個就是嘉爾文對聖靈的理解.
所以這個我覺得是四個學說裡面.
最好的一個.
如果不是.
其實你沒有理由硬說變質.
又不可以是一個real presence.
其實你換湯不換藥.
你又怕得罪天主教.
但你又沒有自己一個的創見.
你的同質又變得不是真的對的.
所以又不需要像硬邦邦.
不用理那個餅那個杯.
不需要這樣想的.
所以呢.
這裡就讓我們可以去看回.
嘉爾文他對於聖經.
對於聖靈的期待.
對於聖餐的重視.
很正常的.
因為當時他們宗教改革.
正正就是要抗衡.
什麼是天主教的錯誤.
我們要避免.
聖餐肯定是其中一個最大的衝突點.

$^{1121}$所以他們一定要銳意去想.
到底聖經怎樣理解聖餐呢.
所以.
慈雲你不可以說他錯的.
一概否定.
不是變質.
但是又變成純粹紀念.
靠人又是不好的.
但是嘉爾文就加上了聖靈的同在.
這個是一個很大的突破.
是對於每一個時代.
我們守聖餐都認為是這麼重要的原因.
次數就不一定要跟足每個星期.
按實際的情況.
反而是那班人對聖餐的意義的了解.
掌握神學.
掌握沒有用的不是知識層面.
是心靈的層面.
所以.
這個就是嘉爾文的神學.
我還沒有說其他.
還有將這套聖經神學之後的演繹.
例如唯獨聖經.
五個唯獨.
還有他對於教會觀.
對於社會的治理.
對於文化.
對於學術.
他有他的見解.
我今天不夠時間寫了筆記.
因為參考的資料都比較多.
我們回頭看.
他在1541年正式回到日內瓦.
在這個市議會裡面.
制定了日內瓦的教會憲章.
記住.
因為整個日內瓦都是基督徒.
三萬多人.
所以整個日內瓦就是一個大教會.
所以這個就是大教會的憲章.

$^{1161}$就是我們的會章.
你要這樣遵守.
這個就是市議會就暫定了.
有他的條文.
這張椅子如果你去日內瓦.
你都可以參觀他坐過的這張椅子.
Calvin's Chair.
嘉爾文坐過的那張椅子.
這個當然成為景點.
在日內瓦.
我在網上很容易就找到.
所以如果你經常去瑞士.
不妨去.
教會組織.
就是剛才講過.
牧師和教師.
就是傳道人.
長老和執事.
原來是要制定.
這些的組織法之後.
他們是負責去執行.
由牧師去牧養.
去探訪.
去講道.
牧師才可以講道.
長老執事是不可以講道.
是普通的信徒.
當然是資深.
所以他們就會有這樣分了這兩個的界別.
所以牧師和教師.
是負責整個日內瓦.
當時的教會的牧養和聖靈的工作.
神聖的工作.
或者和關於教導聖經的真理的工作.
長老和執事.
其中執事.
對的 跟聖經的.
比特前書有.
哥林多前書有.
執事主要是扶貧.

$^{1201}$幫助貧窮人.
是做那些.
分發那些的詞彙.
可能當時有些地方是窮.
因為日內瓦都大.
有些是貴族.
有些真是比較窮苦的人.
所以教會.
既然整個地方都是教會.
整個日內瓦都是教會.
所以教會的長老和執事.
就要扮演著這個角色.
長老有多大的程度.
就是處理那些民間的.
宗教法的事情的審裁.
所以我們可以透過這個.
我們可以了解.
這個教會組織法.
成為十六世紀的長老中.
蘇格蘭的長老中.
首先就採納.
成為了他們教會的憲章.
所以長老中就是長老自會.
所以就是牧者和傳道人.
帶領教會.
所以就不是.
平信徒帶領教會.
所以不是行政主導.
是神的聖經的真理主導.
所以是很不同的.
所以長老制是由這個組織法.
那裡衍生出來的.
所以如果我們下一次有機會說.
這個蘇格蘭的神學家.
這個洛克斯.
是加爾文的一個學生.
他就將這一套教會組織法.
和加爾文的神學.
整個精髓帶回蘇格蘭.
帶來蘇格蘭的大復興.

$^{1241}$所以從蘇格蘭又輾轉去到英國.
英國輾轉移民潮.
就去到美洲大陸.
就帶來17世紀的那個.
敬虔主義的事情.
那個復興.
所以你看到是由加爾文.
這個神學.
對聖經的重視.
開始和天主教就越來越分水嶺.
因為天主教仍然要保持.
他自己的聖禮.
他仍然都是用他自己的彌撒.
他就很少會這樣逐卷逐卷的說聖經.
所以基督教就由這樣開始.
不斷興起不同的牧者.
傳道人.
來去教導聖經.
所以很多的講習班.
很多的不同的聚會.
都是以講道為那個中心的重點.
所以你可以想像得到.
這個就是加爾文他.
那個帶來的一個很大的改變.
這個就是我們上個星期都提過的.
市議會就是長老制.
長老有12個是選出來的.
民間的領袖.
但是他們開會必須是要有那個牧師帶領著.
當時當然就是加爾文.
每個禮拜他們會開會.
談一些整個城邦日內瓦的一些事務.
當然之前會有祈禱.
然後也都會有一些牧師.
認為有些事情要他們跟進.
就會這樣做.
這個就是那個市議會的組成.
所以不是一個民間組織.
是長老制.
是由牧師和長老.

$^{1281}$共同來召集開的.
所以剛才講了.
長老執事是under這個市議會.
來做一些民間的.
例如州制窮人.
例如下個月或者下一年.
那些稅要多少.
那些借貸只能去到多少的利息.
這些他們鉅細無遺都要跟進都要理的.
因為只有一個市議會去管理整個地方.
所以都是一個.
當時是好還是不好呢.
教會間接.
好像連這些基本的民生都要跟進.
都要去處理.
都要去決定.
好又不好.
所以都會很繁瑣.
都會有很多瑣碎的事務.
甚至我們講過.
那個不信三一論的蔡唯塔斯.
就是這樣被火營燒死.
就是因為市議會.
根據牧師嘉爾文所講.
他是經過勸喻都不肯悔改.
所以市議會的長老就執行那個火營.
那你就可以想像得到.
這樣算不算是政教完全合一呢.
表面上是的.
但實際嘉爾文他認為其實不是的.
他只是希望教會扮演著教導的角色.
要教導平民百姓對神的真理的認識.
但教完之後.
如果你不付諸於那個執行去調查觀察他的成長.
其實這樣教完都沒有用的.
他認為是要輔助在那個實際生活裡面的觀察.
所以我們不知道當時他們有多緊密的探訪.
家訪.
現在這個家訪似乎又好像有入屋搜查令.
一定要入屋的.

$^{1321}$不是說沒空很忙.
你不要進來我家很亂.
今天我們就比較少.
但那個時候是沒有這些問題的.
他們是講明牧師是要定期探訪每一個會友的家庭.
了解他們的屬靈的狀況.
還有其實很緊密.
他們如果不去崇拜不去聖餐.
其實就好像那個名單就已經要去到那個勸喻.
和了解他們的狀況.
以致不需要去到市議會去討論.
但似乎當時他們都是一個這麼緊密的城邦群體.
所以長老執事是執行市議會的社會紀律的.
所以其實為什麼有些人民主義者這麼憤怒.
這麼不喜歡嘉爾文.
就是什麼都管 什麼都關你事.
說明人民主義就是不想被天主教壓迫這麼厲害.
想有自己個人的選擇.
譬如現在美國 今天都在講.
如果你今天覺得自己是女人.
你就可以進女廁.
現在那個性別的混亂.
他認為不是混亂 是我個人的自主權.
所以很難想像中世紀什麼都要有監控.
你很難想像我們今天有很多的空間.
純粹是我們個人聽完那堂道.
我怎樣去思想 我怎樣去實踐.
是我個人的成長.
某程度教會都沒有這麼多空間去了解和去干預.
除非你很想和牧者分享.
否則很少人說你今天聽完堂道之後.
你沒有執行 你沒有實踐.
今次你要見一下市議會.
沒有這樣的東西.
不會堂東會見一下你.
為什麼你今天聽完堂道完全沒有反應.
沒有的 任何教會都沒有.
強如恩福那些都沒有.
反而聽完堂道就說這個傳統人講得好 講得不好.
馬上寫那些紙 放在桌面上.

$^{1361}$紙 所謂的非紙.
就投訴那個傳統人說這堂道講得很差.
不知道說什麼 就寫紙.
反而是平順圖去平頭品足.
但是當時就不是這樣了.
當時牧者就負責目樣之後.
要跟進目樣探訪.
發現有違規就要帶到市議會.
執行那個社會紀律.
你想一下都很害怕.
所以牧者是分開的.
牧者是負責目樣探訪.
最主要是教導真理.
當然不知道實際在整體的裡面.
其實他都是在講十幾年.
因為去到1563還是1561年.
他已經死了.
所以由1541到1563 二十幾年.
其實還要是1549年後.
才可以全權在整個日內瓦.
有一個管理的權.
因為之前1541到1549.
都有很多反對的人.
對他的攻擊 對他的針對.
所以你可以想像得到.
在這個憲章裡面.
都一定會產生一些張力.
但是中世紀的時代.
可能那些文字.
大家都覺得是集體負責制.
就沒有我們今天這些的想法.
他們覺得很自然.
有人教我們 有人看著我們.
我們一起成長.
禮拜年一起 所有東西都停了.
要去教會.
他們覺得很正常.
不會覺得有什麼問題.
但是慢慢慢慢衍生的.
那個時代一直轉變.

$^{1401}$是有不同的.
這個的確是那個時代有轉變的.
要喝水了.
應該我今天預備的PowerPoint就完了.
主要是想講教會組織法.
嘉義民將七個聖禮.
幾有創意地認為只是洗禮.
要保留的.
還有聖餐.
聖餐其中更加否定辯則 同則.
即是連他的師父馬丁路德.
雖然不是直接受教於他.
但是他都有那個.
你可以說有那個吉士.
有那個insight.
都連同則都否定.
同時代的.
他說的是spiritual presence.
聖靈的臨在.
在聚會當中.
即是說每一個信徒.
你都可以經歷主耶穌的十誡.
不是一件歷史那麼簡單.
不是只是一個拼命去想.
而是有聖靈在聚會當中.
這是很大的一個創建.
不是那麼容易想到的.
這個就是嘉義民的神學.
其中就是要.
其實某程度是關於牧羊.
因為那麼多的天主教徒離開了天主教.
成為基督徒.
跟中國改革成為基督徒.
到底我們信什麼呢.
這是很大的問題.
我們要信什麼呢.
什麼我們要保留呢.
什麼是不對的.
我們不要再信.
不要再持守.

$^{1441}$所以我們看就很簡單.
但是當時其實那個意義.
壓力是很大的.
因為天主教仍然健在.
法國仍然是整個國家是天主教的勢力.
所以其實沒有宗教改革.
有的都是很零星的.
嘉義民自己都是由法國走了去日內瓦的.
所以其實面對的是整個歐洲.
不是每個人都贊成基督教的.
但是當然有更多人是不滿意天主教.
不滿意不等於他想有一套東西你守住.
他只是不滿意.
他想要自由.
所以你又要擺平這些自由主義者.
其實就必須要有一套是根據聖經的神學.
所以你可以想像是很偉大的人.
嘉義民是很偉大的.
將一套很有系統的神學.
去教導當時離教者.
這些離開天主教的人.
有一套很穩固.
還要是根據聖經的神學.
的一個教義的主張.
然後確立.
就給大家去學了.
既然現在有一套神學.
有一套系統.
所以那個延習就重要了.
所以嘉義民就很強調.
小時候就開始要有延習班了.
所以每一個小朋友的那個靈性.
越早接觸聖經.
越早接觸神就越好了.
就越在單純裡面來成長就越好了.
不然就越大有時候就更加難.
所以這個你可以想像得到.
嘉義民對主日學的重視.
雖然當時不是叫主日學.
對於小朋友對於那些教義問答的重視.

$^{1481}$所以嘉義民不要以為只有基督教要義.
他有很多那些要理問答的.
他寫這些書是駁斥那些異端.
詳細講解那些神學背後的意義.
這些又是要將很多很深奧.
一般人不是很明白的那些神學.
用很簡單的語言講解.
連小朋友看完都很簡單明白.
所以就有教理問答.
一問一答.
小朋友回到家就是問.
接著你就要答了.
所以小時候的小朋友.
自小就已經深化了對聖經神學的認識.
這些很多時候長大之後.
就可以被神用的不同的傳道人.
或者是將來的一些教會的領袖等等.
所以嘉義民為我們奠定了一套很簡單.
但是很重要的神學.
有一套神學.
我們不是為了反對天主教而離開.
我們是為了行這些真理.
而我們不是我們想離開.
而是我們要行這些聖經真理.
而天主教不認同.
所以我們要離開的原因.
那就變成不是為了激情.
不是為了情緒的反對.
是為了真理而持守.
兩者又是不同的.
所以如果教會是越多弟兄姊妹對聖經有認識.
有明白.
教會自然會成長和復興.
他會去到不同的地方.
都會將神的道傳遞給那個地方的人.
今天我就講到這裡.
都差不多了.
我預備了的還沒有整理的嘉義民另外的神學.
其中提到的教會觀.
社會的治理的到底是不是政教分離還是政教合一.

$^{1521}$還有一些教育文化.
當然還有很重要的五個唯獨.
還有他的預定論.
就是那個Tulip那五樣東西.
我們下個星期就將餘下的都講了.
如果我試想一想.
可以再講多一點.
就是嘉義民之後的十六世紀.
十七世紀頭的蘇格蘭的大復興.
英國大復興.
甚至乎應該講美國的大復興.
有幾個重要人物.
例如愛德華茲.
這個Jonathan Evers.
愛德華茲.
還有就是衛斯理約翰.
還有就是.
有些想不到了.
不過遲一點再講.
大概就變成一堂到兩堂.
我們就完結了宗教改革這幾個重要人物.
其他的沒那麼重要就未必講了.
當然希望最後一堂我們講過.
過去和未來.
我們又要講一下宗教改革.
在我們現在這個時代.
我們應該怎樣理解.
還有我們有什麼出路呢.
希望大家一起帶出一些point.
大家一起去思想.
我們一起禱告.
因天父給我們都有一小時多一點時間.
我們去將嘉爾文的神學做一個整理.
多謝你讓他能夠在天主教的壓力.
在人民主義者.
或者一些城邦的君主.
對他的攻擊.
他都能夠很有系統.
很安靜地去整理一套神學.
系統的神學.

$^{1561}$和結合在聖經的神學裡面.
以致讓當時的信徒.
不只是為了情緒跟風離開天主教.
而是知道要跟隨聖經的真道.
來相信主耶穌.
今天我們都能夠可以由基督教要義.
流傳到今天我們可以去研讀.
我們很盼望這位神學家.
雖然已經是500年前的人.
都給我們能夠有一份新鮮感.
有一份看到.
他是怎樣在他有生之年.
差不多都是20多年.
最有精力的人生的青年階段.
來服侍寫書,講道.
制定教會,組織法.
和在日內瓦的瑞士.
帶來的一個大復興等等.
都給我們看到他的貢獻.
短短的50多歲.
就帶來這麼大的歐洲的影響.
一方面我們想到他對聖靈的重視.
對聖經的重視和尊重.
對講道,講壇的重視.
都是很重要的一些復興的種子.
埋下了在歐洲的教會.
我們很盼望.
其實我們今天都不斷在做的.
但都求主給各教會.
我們能夠看重.
不要只有外表.
外在的外學.
而是那個心靈.
那個很重要的精髓.
都能夠把握到.
願神幫助.
我們再一次獻上感恩.
我們這樣祈禱.
奉耶穌基督的名.
阿們.

\newpage



\section{}
\label{sec:Dn07Lw4F37U}
\textbf{2025-01-16 周四晚查經班《改革之路(8):加爾文(四)》 講員:曾裔貴牧師}
\newline
\newline
連結: \href{https://youtube.com/watch?v=Dn07Lw4F37U}{\texttt{https://youtube.com/watch?v=Dn07Lw4F37U}} ~~~~ 語音日期: 2025-01-17
\newline
\newline
\hyperref[sec:3Sqlgs98HDM]{\small{< < < PREV SERMON < < <}}
~
\hyperref[sec:index_chronic]{\small{[返順時目]}}
~
\hyperref[sec:index_scriptual]{\small{[返順卷目]}}
~
\hyperref[sec:5Lw6EDNBsNI]{\small{> > > NEXT SERMON > > >}}
\newline
\newline
$^{1}$多謝天父給我們一直都有機會.
去將宗教改革.
相當重要的兩位神學家.
他們的人生,他們的主張.
和當時他們為著要將信徒.
帶進聖經的真理.
他們高舉的是聖經的重要性.
和唯獨聖經在那個時候.
就很認真地被確立.
而圍繞著的是那個唯獨信心.
和透過基督這個完整的.
無條件的救贖.
我們很盼望能夠.
承傳這些前人.
他們在非常時期.
他們所帶進的那個對真理的尋索.
以致不是為了要反對天主教.
而是要將信徒.
帶回神的真理和救恩裡面.
這些是非常難得的.
所以我們都很盼望.
其實一開始我們都去思考.
宗教改革過去和未來.
我們一直都用了八堂的時間.
去看過去的歷史.
和種種的神學的主張.
但都盼望我們能夠.
可以有一些空間去想一下.
在21世紀的香港.
當我們去跟隨主的時候.
有什麼我們需要去堅持去探索.
和去一直去傳承.
使我們的下一代.
或者有更多人認識真理.
不會被一個時代扭曲了.
譬如現在我們看到的個人主義.
個人那種絕對的無上自由的選擇權.
已經在歐美是非常嚴重.
他們在性別的差異已經完全沒有界線.
我們很盼望我們在我們這個時代.

$^{41}$特別是消費主義已經是無孔不入的.
在我們每一個層面來滲透.
我們想到其實在中世紀時代.
根本沒有所謂這些消費主義.
我們想到時代的確是有不同.
但都盼望我們能夠.
看到有哪些是我們需要去堅持的真理.
有哪些可能是相對的.
我們未必需要這麼執著的.
願神都幫助我們.
在我們最後一兩堂的神學課程.
都能夠想一想.
好像一個盤點那樣.
幫助我們對信仰更加的認真.
更加的對一些要認真的真理.
我們無論如何都要捍衛.
願神幫助.
我們今天晚上在你面前等候.
禱告奉耶穌基督的名.
阿們.
一直都說第二代的嘉爾文.
就是真正將整個宗教改革.
不再是為了反對天主教.
而去做種種的舉措.
而是真的要沉澱下來.
確立為什麼需要脫離天主教.
固然那個現實就是.
天主教已經不要你了.
你要走就能走.
要不然你就回來.
必須要跟足我們天主教的.
全套的綱領.
全套神學的範圍.
否則你根本就不能在教會.
大公教會裡面聚會立足.
所以是全盤的一個套餐.
所以嘉爾文的重要性.
和他的開創性.
就是不再沉溺在恐懼的對抗.
而是真的要大進.

$^{81}$到底為什麼要改革.
這個是最重要的.
所以是帶信徒回到聖經這個大方向.
回到對神的唯一的救贖的專注.
這些就是真的很開創性.
讓當時的信徒有一個很大的心靈.
所以引進了一次得救.
永遠得救的預定論.
後人當然將他再去到極端.
極化.
你沒有被揀選自然就不得救.
變成會給人一種錯誤的理解.
你們嘉爾文主義這麼霸道.
我沒有被揀選.
不關我的事.
神不揀我.
那你就說我注定要滅亡.
其實那個意思就不是這樣的.
我看回一些著作.
一些對嘉爾文神學的描述.
有些宗教.
特別是教會歷史的書.
其實預定論.
其實在當時的提出.
是基於一份很人道的模樣.
因為天主教就很清楚.
給了一個最後通牒.
你離開教會就沒有救恩.
沒有救恩就是說你會沉淪.
去到永恆的火湖.
那嘉爾文和他們的同路人.
要帶領這幫人離開了教會.
其實很多是一些一般的普通平民百姓.
那他哪有空理你這麼多的神學.
你給我一個信心.
到底我現在跟你們這幫神學家.
離開天主教.
跟隨你們去信基督教.
那我有什麼保障呢.
這個就是在那個模樣上.

$^{121}$嘉爾文真的要回到聖經那裡.
看到神的那個預定.
當然他有很多聖經的支持.
我剛才來的時候.
很趕就印了一份筆記.
就去到前面我們看到.
嘉爾文怎樣看經濟.
怎樣看一個信徒.
信了主耶穌.
既然現在是預定論.
神預定了我們得救.
我們是被揀選.
那到底我們在人生.
應該怎樣來看自己呢.
那裡我就直接.
因為我真的看基督教要義.
就整個章節關於經濟.
我就印了出來.
我要完了.
今晚我才把PDF給大家.
很好看的.
很容易看的.
還有幾個信念是很寶貴.
很重要.
因為我現在沒有帶那個.
我自己印刷的那本.
就給不了大家.
我又保留不了給大家看.
因為現在屏幕就在PowerPoint上.
我稍後就會把PDF給大家.
大家回去可以再看一下.
嘉爾文怎樣看一個.
蒙了主耶穌拯救的人.
他的人生應該怎樣呢.
其實有幾個信念.
就是很簡單的信念.
就是要過一個很.
Humble.
簡樸的生活.
和盡量把神給你的恩賜.

$^{161}$來去發揮.
還有很重要.
要照顧那些有需要的人.
就是說要有一個憐憫的心腸.
大致上是這樣.
因為他有好幾個Point.
我就簡單的總結了.
我們稍後會講到很多點.
因為我稍後要逐點和大家看.
整套的嘉爾文的神學.
我們就開始進入.
其實有些枯燥的.
如果你累的話.
先喝一口水和洗個臉.
因為都要多些心力.
三個重點.
我們上一堂講過了.
真的很重要的.
嘉爾文整套的神學.
完全是建立在聖經神學裡面.
所以如果你真的買本嘉爾文的基督教要義.
其實你看到.
就是系統神學的將神.
聖經啟示的神是怎樣的.
他的屬性.
他的位格.
等等.
接著就是基督論.
接著就是聖靈論.
接著就是教會.
當然夾雜的是救恩和人論.
他都有講的.
所以真是完完本本的.
一本系統神學的書.
他一個人寫的基督教要義.
但是更寶貴的就是.
他是基於聖經神學.
來確立他自己的信仰和綱領.
來跟隨他在日內瓦的三萬多人.
那個人口.

$^{201}$當然其他人.
有些在外面的都會跟隨的.
因為開始印刷術.
都已經是蓬勃了.
所以這一套的神學.
就成為整個歐洲.
一路發揚光大的嘉爾文神學.
所以聖經神學是嘉爾文整個神學的精髓.
這個就真是很寶貴了.
我們華人教會.
都受嘉爾文神學很大影響的.
所以我們對簡選論.
對於長老字會.
長老鐘.
我們一點都不陌生的.
我們華人教會大部分都是以長老制造那個字會的.
我們上一堂課講過.
他當時在日內瓦已經奠基了.
牧師和傳道.
就是治理教會的.
或者那個神的事情的.
他要探訪.
要講道.
要做聖禮.
這個他們都會這樣做.
但是一般世俗的事.
就會有長老和會有那個執事.
來負責.
來帶領那個社會的秩序.
因為嘉爾文認為.
神安排了有一個社會.
就不可以不管的.
因為人是必須要教育.
所以正面是教育.
不聽才是要管.
你可以說是教育的某程度的雛形.
因為當時都沒有普及教育.
神學院基本上後來就成為了很多大學的前身.
所以這些都是一個宗教改革的發展.
所以嘉爾文都很著重那個大學教育的.

$^{241}$或者普及的宗教和一般的教育.
所以神學和聖經.
其實是在嘉爾文裡面.
是完全無縫的結合的.
你看基督教要義.
大量的聖經是他引述的.
所以有些人說他對聖經.
簡直是去到一種近乎瘋狂的愛慕和尊重.
這個你可以說對嘉爾文的了解.
就是他真的很認定聖經的無誤.
聖經的權威.
一切的東西都是要根據聖經來執行.
聖靈比較重要的是什麼呢?.
他有一個章節的.
但是更重要的是上個星期我們說過.
就是在聖餐裡面.
聖靈是嘉爾文獨到之處.
不再是變質.
更加不是同質.
馬丁路德就是兩樣都是.
又是變質.
但是其實是沒有變的.
主耶穌臨在那些餅和杯.
所以某程度上都是質地有一點點共同.
你吃的是餅.
喝的是葡萄酒.
但是其實意義上就等於主耶穌的血.
但是嘉爾文.
第三個就是慈雲尼.
就說一概不是.
只是紀念.
不需要想得那麼神聖.
紀念耶穌十下的犧牲.
但是嘉爾文再加多一樣東西.
在紀念裡面.
每一次聚集的聖餐.
這個我們真的要學習.
要經歷.
因為現在你和我都是每個月.
或者每個星期.

$^{281}$我們真的在做.
所以我們要有一個等候.
讓聖靈在我們的聖餐裡面充滿.
這個是嘉爾文主義.
對聖靈這一方面他很著重.
所以我們就.
又不是變質.
又不是同質.
又不是只是紀念.
而是有聖靈臨在.
那就是說聖靈代表著主耶穌臨在.
在我們的聚會裡面.
在我們一班的聖徒.
領聖餐的時候.
所以盡量克制自己.
不要胡思亂想.
不要想著待會午餐吃什麼好呢.
那些東西都可以吃的.
但是聖餐聖靈的臨在.
就是我們應該要這樣學悟.
所以這個也是嘉爾文神學其中一個.
你可以說是獨到之處.
這個我們上個星期講過了.
我們進入到每一個point.
我和大家逐點去分析.
剛才提過了.
聖經神學就是整個嘉爾文神學最重要的核心.
就是一定是對聖經的啟示.
所以根本分不開的.
沒有聖經的啟示.
沒有神的啟示.
根本就不會有66卷的聖經.
所以根本就是一種不正自明的.
因為有神直接的啟示.
所以有聖經.
所以有我們今天的66卷的聖經.
所以嘉爾文就是說.
這個啟示是超自然的.
是神已經透過不同的時代.
給了我們有祂的話語.

$^{321}$所以聖經的神學是絕對的權威.
教會存在.
我們要根據聖經來帶領.
牧羊教導信徒.
直到主耶穌回來.
所以這個對於華人教會是沒有困難的.
不過如果近代開始對聖經的批判.
那就是另外的話題了.
因為已經將聖經可能有錯的那種觀念.
如果滲得太嚴重.
有錯字是正常的.
因為執字凹凸錯很正常.
我們看出有些字不同.
甚至有不同的抄本.
那些是人為抄寫過程的錯誤.
是很合理的.
但是在整個我們說的默示裡面.
神的啟示裡面是沒有錯的.
所以那個原典.
第一個的原稿是沒有錯的.
但我們已經不是要崇拜.
去拜那個原典.
因為已經失傳了.
已經找不到了.
已經幾千年沒有了.
但是那個首抄的過程.
是有少少的手誤.
但是不足以令到這個啟示來到被挑戰.
甚至我們要否定整個啟示.
這個是絕對不行的.
因為這個涉及到信仰的問題.
就是我們對聖經絕對權威的捍衛.
嘉義文就是應該和我們華人神學.
有一脈相承.
我們都是秉承這個聖經神學的權威性.
這個真的是很重要的.
我絕對不是所謂古老.
而是不能討價還價的.
什麼時候哪個傳道人.
他說我不信聖經是神所默示.

$^{361}$那個傳道人是立即要被解僱的.
因為他根本已經是離經半道.
所以這個是不能討價還價的.
這是絕對性的真理.
應該說福音派絕對看為聖經是無誤的.
這樣才有東西可以談的.
如果你說聖經這件事那裡不行那裡不行.
那就沒得說了.
你不知道哪裡對哪裡錯.
那就帶來一個很大的問題.
這個要留意了.
因為近代的學者.
未必真的相信聖經絕對的無誤.
已經有些教會.
有些宗派是未必再這樣捍衛的了.
我聽郭文慈牧師在播神的神學講座.
他都沒有點名.
有些神學院已經將他自己的聖經無誤修改了.
即是神學宣言都修改了.
你可以想像得到.
這些絕對不是我們要追求的古老等等.
所以我們回來集中就是.
嘉爾文神學一定是聖經神學做那個骨幹的.
就是因為有聖經有這個啟示.
所以就有他的神學的出現.
所以他的基督教要義就是等於聖經神學.
有系統地將聖經神學逐個部分去詮釋.
有系統地將神是怎樣的.
基督,聖靈,教會,人論和為什麼要有救恩.
諸如此類就是這樣的聖經神學.
那麼關於啟示.
不如我直接讀一段我在基督教要義裡面所抄出來的.
應該是將他的說法說出來.
關乎啟示我覺得這個很精彩.
一切的起點為什麼會有基督教呢?.
就是因為神直接向人類啟示他自己.
我直接讀了.
因為這樣是更加原汁原味.
當然是由拉丁文轉為我們現在翻譯成中文.
當然那個遞減是有一些.

$^{401}$但是都不足以將人類的原著的意思完全扭曲.
一定不會.
所以我讀給大家聽.
關乎啟示這個topic.
神所有的啟示的目的.
是為了人的幸福.
但是因為人的罪.
創造的啟示的結果.
反而使人的罪更加深.
因為已經否定這位造物主的存在.
否定他向人類直接啟示他自己.
他是絕對的榮耀.
全能.
他自己的存在.
但是當他要這樣來到.
他好像書專降貴.
啟示給我們知道.
這一切的啟示的目的背後.
都是為了人類最大的幸福.
但是因為人已經有罪了.
所以對神的啟示.
對神要介入我們人類的社會.
我們就有很大的抗拒.
所以我很快就想到.
我們那時候查過《約翰福音》.
一章第七節.
他到自己的地方來.
自己的人都不接待他.
凡接待他的.
去到十二節.
他就似他們權並成為神的兒子.
你可以想像到.
耶穌成為一個可以見得到的人.
真是有血有肉.
在嬰兒一路長大的人.
在以色列長大的.
是神的啟示的最終極.
就是要帶進那個幸福.
就是將人帶到神那裡.
但是人的罪.

$^{441}$看不到這位耶穌.
看到的就是拿撒勒木匠的兒子.
看到的就是他是一個很潛望的人.
所以這個福氣.
這個幸福.
變相成為更深的拒絕神的罪.
所以這個就是嘉義文對啟示的.
原本最重要的給我們的啟發.
原來神一切所做.
就是要使我們人類得著那個幸福.
就是當他和我們接觸.
我們就知道原來人存在在這個社會.
在這個地球.
原來是有目的的.
原來是有他的計劃的.
我們在這個計劃裡面.
一路這樣來到去.
倚靠著聖靈.
倚靠著神.
我們就會得到那個幸福感.
得到那個幸福.
所以很簡單啟示.
在嘉義文的那個理解.
就是為著人的幸福.
所以如果你真的抄下它.
或者越來越可能你很想買這本書來看.
英文版是圖書館有的.
政府圖書館是可以借的.
基督教要義.
甚至英文版在網上某些平台.
是可以免費下載的.
如果你英文看到的就沒有問題.
或者你索性就去基督教書室.
買這一套基督教要義.
很好看的.
不難看的.
很有條理的.
你想找哪個題目就可以找到.
譬如我今天講到經濟.
我馬上搜尋那個題目.

$^{481}$就讓我看到.
有些東西.
我整份直接印給大家.
只不過我剛才來的時候.
忘記了將PDF放在群組.
我稍後回去就會放的.
那個經濟論.
其實對我們今天都是有幫助的.
都真的有很好的提醒.
信了主耶穌.
你更加應該將神給你的恩賜.
給你的恩典.
要用得其所.
這就是對人不要貪婪的警告.
身高兩準.
送機票去泰國.
學問什麼都不用.
那些肯定是騙局.
都要去的.
不要怪別人殘暴.
這些是常識.
所以你可以想像到.
我們如果人信了主.
開始問自己.
我存在.
神對我有什麼要求呢.
我去到經濟那邊再說.
不然越說越遠了.
啟示就是大概這個部分.
很重要的.
是神直接介入我們人類.
透過聖經.
透過歷世歷代.
先是向猶太人顯現.
然後一直那個救恩.
就是這樣進入新約.
所以主耶穌的降生.
就是整個啟示的極致.
就是遇見耶穌.
就是人類進入真正的幸福.

$^{521}$所以真的不難發現.
很多信了主耶穌的人.
他的生命能夠重整.
很多吸毒.
或者他以前很壞.
一信了主耶穌.
他的生命改變.
第一件事就是他自己感受到幸福.
而他一直追求.
幸福不是禍延.
他的幸福是會讓他下一代.
深深感受到神的恩典.
所以如果我們第二代.
甚至第三代.
都是很敬虔愛主.
就很肯定這個家庭有幸福.
真的很肯定.
因為這個是很事實.
你可以想像得到.
不是錢.
不是那個家族的東西.
而是那個家族的對神恩典的愛.
毛耳中看到.
講福堂的一個講道.
他們有一個家族.
已經去到第三第四代.
如果我沒有記錯.
好像是田家丙那個家族.
去講見證.
講到第三第四代.
都是這樣愛主.
印象我有些模糊.
我要找回那些資料才行.
有些少印象.
但是感恩的是.
整個家族十幾人上台.
講述由那個父輩開始.
怎樣領受恩典.
現在繼續.
其實應該是他的女兒.

$^{561}$女兒已經很大年紀了.
已經有他的孫子.
都上台.
一齊上台講述神的恩典.
所以得到啟示.
或者神為什麼要啟示我們呢.
就是要將幸福帶給人.
所以在當時.
跟隨嘉爾文的人.
就會明白.
自己原來不需要懼怕神.
只需要在神的裡面.
在他的護理.
這個稍後我會講.
在神的護理裡面生活.
盡你自己的能力.
去做你人生神給你的工作.
你就會得到那份幸福.
最終當然就是去到永恆裡面.
享受永恆的福樂.
神觀沒有什麼特別.
因為跟奧古斯丁完全一樣.
神的獨一.
神的全能.
反過來.
他們就少去分析神的屬性.
反過來.
這裡就特別這兩點.
是很重要的.
就是神的主權.
和神的護理.
護理就是說.
他創造了宇宙萬物.
到今天.
仍然用他的權能來管理.
這叫做護理.
保護管理.
就是說.
這個世界從來沒有失控.
這是很重要的神學信念.

$^{601}$基於他絕對的主權.
所有事情都要在他的命令.
他的旨意.
他的智慧下發生.
然後既發生.
同時間又護理.
這樣的上帝觀.
就令到我們.
上帝觀就令到當時的中世紀.
離開天主教.
又會有一個很大的.
所以那個預定論.
其實是離不開神的主權.
神的護理.
神對人的美好的心意.
就是要將幸福.
給到每一個相信他.
認識他的人.
因為罪是令到我們拒絕神.
令到我們不相信這個世界有做物主.
讓我們以為.
世界是沒有神的.
所以人沒有了.
這個神的管理.
表面的幸福.
就是自己營造的.
就不是神.
所賜的.
期望的.
真是會給的幸福.
所以我拆開.
回頭.
因為當你說到啟示.
這個幸福感的時候.
怎樣可以獲得呢.
又是一男子的了.
嘉義文就很相信聖經的教導.
怎樣可以得到呢.
就是信徒必須經常.
透過禱告.

$^{641}$與神接觸.
你就可以認識.
明白神在你生命.
在教會的帶領.
這樣才可以.
來到去很深的感受到.
這一份的幸福.
否則就變得很機械化.
我相信耶穌了.
洗禮了.
我就不再回教會.
我就不再祈禱.
但你要給我幸福.
這樣是不對的.
因為你根本沒有在那個軌道上.
與神一起.
所以原來在那個祈禱的部分.
嘉義文的基督教要義.
有好幾個題目.
教祈禱.
所以如果我們不懂祈禱.
我有空借給大家.
你們可以看看那個祈禱的題目.
因為很多我印不了.
幾十頁.
很鉅細無遺的教祈禱.
原來就是怎樣接觸神.
獲得那個真正的幸福.
神的幸福已經給了.
透過啟示.
但怎樣可以獲得呢.
一定是透過對聖經的認識.
更重要就是透過祈禱.
所以教會.
責無旁貸是要鼓勵弟兄姊妹.
祈禱.
祈禱不是唸咒語.
祈禱是與神.
帶著愛的生命相交.
和認識祂.

$^{681}$所以就拆開了說.
祈禱.
不過我們仍然說.
上帝觀的時候.
就不是我們用顯微鏡.
去分析神是怎樣的一位神.
永遠分析不到.
所以今天不同的神學.
學派就對神.
譬如三一論是要爭論.
其實爭論完之後.
都沒有答案的.
因為只能根據聖經的啟示.
告訴我們.
我們目前有限的智慧.
相信是.
三一真神.
三位一體的真神.
就是這麼多.
你不需要再去有其他的辯論.
反而就是說.
我們要認定神這個絕對的主權.
這個是嘉爾文神學.
你說他霸道嗎.
其實是基於那個模樣.
也都基於聖經的啟示.
神是有祂的主權.
有些經文.
可以和大家揭一揭.
不過我沒有聖經在.
我就用手機.
譬如羅馬書第九章.
就已經是這個教導.
保羅自己就說.
這個神的主權.
我很快讀一讀給大家聽.
所以.
基督教要熱好看的地方.
就是很多聖經.
真是很多聖經.

$^{721}$我們很難想像.
一個55歲就死的人.
他寫了很多的.
聖經注釋.
又翻譯聖經.
還要是.
每天講道.
還要在日內瓦開會.
處理那些民生的事.
其實那些時間.
和他們不會有夜晚的夜燈.
他們是用油燈.
更加不會有電燈.
更加不會有空調.
上網.
沒有這些東西.
為什麼可以做到這麼多東西.
都很難想像.
九章羅馬書.
我們講過的.
第六節.
這不是說.
神的話落了空.
因為從以色列生的.
不都是以色列人.
而亞伯拉罕的後裔.
都作他的兒女.
唯獨從以色列生的.
才要稱為你的後裔.
這就是說.
肉身所生的兒女.
不是神的兒女.
唯獨那應許的兒女.
才算是後裔.
因為所應許的話是這樣說.
到明年.
這時候我要來.
撒拉必生一個兒子.
本來一直在引用創世記十七章.
還有利伯加.

$^{761}$一個人就是從我們的祖宗.
以撒懷了人.
雙子還沒有生下來.
善惡還沒有做出來.
真的沒有 兩個都還沒有出世.
此因要顯明.
上帝揀選人的旨意.
不在乎人的行為.
乃在乎召人的主.
神就對利伯加說.
這些已經是創世記的經文.
將來大的要服侍小的.
正如經上所記.
那國是我所愛的.
以素是我所護的.
這樣我們可說什麼呢.
難道神有什麼不公平嗎.
斷乎沒有.
因他對摩西說.
我要憐憫誰就憐憫誰.
要恩待誰就恩待誰.
對摩西說.
是什麼時間呢.
出埃及記對付法老的十災的時候.
據此看來.
這不在乎那定義的.
也不在乎那奔跑的.
即是人的熱心.
自己的行為.
去換取神的恩典.
祝福都不行.
只在乎法憐憫的上帝.
神.
因為經上有話向法老說.
我要憐憫誰就憐憫誰.
我將你興起來.
特要在你身上彰顯我的權能.
並要使到我的名.
全遍天下.
如此看來.

$^{801}$神要憐憫誰就憐憫誰.
要叫誰剛硬就叫誰剛硬.
絕對的主權.
摩西面對法老.
法老是很多阻撓.
很多的攻擊.
很多他以為自己可以駕馭摩西奴隸的民族.
說得這麼誇張.
說我要帶我的族民離開埃及.
回去應許之地.
神應許我們的地方.
神是誰來的.
是我們埃及的神.
還是哪位神.
是哪位.
神透過摩西.
其實摩西是很軟弱很怕事.
拿著棍子.
就當自己是聖子.
果然真是.
神就用神能.
使得棍子都能夠行神跡.
後來摩西.
慢慢去到第六災之後.
就不用靠阿倫.
自己越來越純熟.
越來越有信心.
等於我們做傳道.
越講就越有信心.
經過了很多磨練.
摩西也是.
慢慢他根本不用.
阿倫幫他.
成為一個背後.
他自己已經有信心.
行在法老面前.
因為他知道神有他的權柄.
交付給他.
那一刻.
他能夠行十災.

$^{841}$的神跡.
即是蛙失凝蓄.
蒼薄黃黑.
水滅血.
最後就是滅長子.
都是摩西能夠行使的.
現在保羅就告訴我們.
就是想.
爭論這個問題.
神有絕對的主權.
所以嘉義文的神學精髓.
預定論的精髓.
就是建立在神的絕對主權.
所以嘉義文的神觀.
就不是這樣.
有空在桌子上討論.
什麼叫三一論.
什麼叫無所不在.
無所不能.
中世紀的時候.
有些神學家.
特別是當時.
十二十三世紀的時候.
他們有些神學是很誇張的.
神能不能夠做一個石頭.
他自己舉不起.
討論一頓.
做物主是無所不能的.
他能不能做一個石頭.
他自己舉不起.
這些簡直是.
荒謬的神學討論.
因為根本沒有意思.
根本沒有答案.
神怎麼會.
這樣限制他自己.
然後就說天使.
在一個很小的.
一個的.
手指裡面.

$^{881}$有多少天使在跳舞呢.
討論這些做什麼呢.
這些討論完是沒有意思的.
所以對嘉義文來說.
他的神觀.
不是討論神的那些東西.
而是功能性.
怎樣在他那個時代.
去提出.
牧養他的會友.
現在那個會友.
那班人.
是被天主教趕走.
離開了天主教.
面對的是.
生和死的.
關乎到靈性的問題.
其實面對的壓力.
是很大的.
所以嘉義文的神觀.
絕對是一個.
牧養的.
出發點.
所以這個絕對的主權.
如果我們在神絕對的主權底下.
有人說.
你是被踢出教會.
你是不會有救恩.
但嘉義文就告訴我們.
不是的.
我們相信的這位絕對的主.
這位造物主.
已經在耶穌的救恩裡面.
十字架揀選了我們.
他有主權.
去保護我們.
所以主權和神的護理.
是不能夠分割的.
所以這個神的護理.
不是神.

$^{921}$這個是後期.
哲學一直一直討論.
開始討論.
我們都找不到祂的存在.
我們看到這個宇宙.
是不是祂創造了這個宇宙之後.
祂就自己走了不再出現了.
所以這個宇宙就成為了上年的.
祂自己自動.
來自動運轉.
是不是這樣呢.
這個就是現在的哲學大概的意思.
但是聖經就說.
不是的 祂到今天.
都藉著祂的愛子.
統管著萬有.
希伯來書第一章.
這些經文大家應該都認識的了.
所以我就說那個精髓.
就是嘉爾文他對聖經的尊重.
認定聖經的權威.
第一節.
一章一節.
希伯來書.
神記載古時.
古時就是遠古.
當夏娃的年代.
藉著眾先知.
多方不同的方法.
不同的方法.
來去曉喻列祖.
列祖就是我們的祖先亞伯拉罕.
以撒 雅各 或者是阿當.
或者是以洛.
或者是羅亞.
等等的先祖.
不過留意第二節.
就在這末世.
藉著他兒子曉喻我們.
不再藉地上的.

$^{961}$神興起的先知.
眾先知.
Many的.
這些Professor.
而是透過最後末世.
藉著他的愛子.
直接透過愛子.
曉喻人類.
他的真實.
接著又早已立他為.
承受萬有的.
創立世界以前.
造物主義已經選立了耶穌.
成為那個承受萬有的.
也曾藉著他.
創造諸世界.
如何創造呢.
你想想耶穌未降生之前.
與父神的.
同在.
在永恆裡面的.
那種三一的同在.
是我們不能夠明白的.
但是聖經已經告訴我們.
藉著耶穌創造諸世界.
見得到的.
見不到的屬靈世界.
接著留意第三節.
他是神榮耀所發的光輝.
是神本體的真相.
我們讀完也未必明白.
這些字眼.
但是他就是完完全全.
將神反映出來的.
那一位的愛子.
那一位的耶穌基督.
是神本體的真相.
所以我們見到耶穌.
其實就見到神了.
接著留意.

$^{1001}$常常用他權能的命令.
托住萬有.
上帝的護理.
就在這裡.
很清楚的說了.
從來沒有殺手.
不理人間的疾苦.
不理世界的混亂.
容許.
不等於他.
撤退.
或者闊佬難理.
這些是現代的哲學家.
不斷去推敲.
不斷在爭論.
不斷在懷疑.
帶來地上的人.
只會繼續懷疑神的存在.
神的愛.
但不是的.
聖經告訴我們.
主耶穌就是神的顯現.
他就是這樣.
常常用他權能的命令.
去托住萬有.
到下面.
第三節.
他洗淨了人的罪.
就坐在高天至大者的右邊.
他所承受的名.
既比天使的名更尊貴.
就遠超過天使.
因為當時.
在希伯來書的年代.
他們會拜天使.
天上的使者.
耶穌只是比天使高一級.
但不是的.
從來都不是.
耶穌就是那一位造物的主宰.

$^{1041}$所以絕對不是其中一位天使.
所以這裡是講到.
主權和護理.
神學就叫做Privilege.
其實就是.
保護管理.
就是說.
祂的主權.
祂的旨意.
從來都沒有離開.
祂創造的宇宙萬物.
一直都在祂的旨意裡面.
進行.
這些.
我想真的只能夠用.
信仰去接受.
因為現在世界的哲學.
已經對這些.
簡直是鄙視.
什麼絕對主權.
我都不知道.
現在從來都不需要.
祂的存在來保護我.
很多哲學家.
無神的.
很大言不慚.
但是記得我們信主.
特別我們嘉爾文主義的嘉爾文神學.
那個精髓就是.
相信神絕對的主權.
就是說我們每個人.
都真的要面對祂.
相信祂的護理.
我們要記得.
有些時候我們需要忍耐.
等候.
不要這麼快.
每樣東西都要立刻超說.
按個按鈕.
就要讓我看到.

$^{1081}$有審判.
什麼都這麼快.
就好像違反了神的護理.
神有祂的計劃.
有祂的時間.
我們就是要對神的主權.
尊重和馴服.
在祂的護理底下.
我們過一個.
享受與神同在的.
幸福.
這個就是上帝觀.
嘉爾文比較務實.
不是討論一些不拉肩的神學.
討論完都沒有用.
為什麼要這樣做呢.
人觀就更加和絕對主權.
有直接的關係.
就是告訴我們.
人的徹底的無助.
不可以靠自己.
去得到神的救恩.
所以必須是.
神主動介入.
所以.
人論.
D和E又是一個套件.
是根本上無法分割的.
某程度上我們今天.
呼喚派都應該還是很看重這件事.
人的徹底的無助.
需要.
外在的救恩.
所以基督論.
就是告訴我們.
主耶穌無條件.
神的愛子.
甘願為我們犧牲.
帶我們進入.
對神的認識.

$^{1121}$和透過.
神的啟示.
使我們知道原來主耶穌.
不是普通人.
祂是神的獨一兒子.
是神的計劃.
又要讀一些經文了.
《加拿大書》四章.
講主耶穌的降生.
我們很多時候.
看路加福音的描述.
馬槽.
聖誕節.
但是《加拿大書》又講.
降生的原因.
不知道大家有沒有聽過.
四章.
我說那承受產業.
雖然是.
全業的主人.
但為孩童的時候.
卻與奴僕毫無分別.
當然和上文.
保羅的比喻.
有關聯.
就是那些管家對那些未能夠領受.
繼承產業的.
小朋友.
他有那個全權的對他們.
教化.
比如當時英國有些管家.
就是這樣.
或者當時希臘人.
羅馬都有這些管家制度.
保羅就用當時坊間的.
這些例子做一個比喻.
是講律法.
對猶太人的規範.
前提是講.
這個point(點).

$^{1161}$他就想告訴我們.
當耶穌基督來到.
這個世界降生.
就徹底改變了這個律法.
那個用處.
和背後的精髓.
那個精神.
所以這裡.
大家要對上文有一點認識.
你就會明白.
他這裡所講的東西.
這裡告訴我們.
他講的那個比喻.
管家很嚴格管束.
他的主人.
應該叫少主人.
因為那個主人.
將他的兒子.
交付給這個管家.
這個管家.
某程度就有權.
去對這個少主人.
有很嚴苛的管教.
為了要提升他將來.
成為那個.
全業的主人.
第一節.
我說那承受產業的.
雖然是全業的主人.
即是孩童的時候.
即是少主人.
但為孩童的時候.
卻與奴僕毫無分別.
我是甚麼的.
為甚麼你這樣對我.
管家說你爸爸叫我管你的.
你現在.
你拜託八點鐘就坐好.
溫書不准上網.
不准甚麼.

$^{1201}$他在那個感覺上.
就好像我一個奴隸.
我不是奴隸.
我是我爸爸的兒子.
不過在那個運作上.
他就受制於那個管家.
這裡是一個比喻.
保羅就告訴我們.
耶穌基督未來之前.
我們眾人都在律法之下.
第二節.
乃在師父和管家的手下.
直等他父親預定的時候.
來到.
例如十八歲.
他可以繼承爸爸的產業.
他繼承了.
那麼大了.
管家就不再管束他.
他只是請管家回來工作.
大概是這個意思.
第三節.
我們為孩童的時候.
受管於世俗小學的.
也是如此.
這個又是可圈可點.
保羅都說到.
他們這一班猶太人.
在主耶穌未來之前.
我們都是一樣.
我們循規蹈矩.
守律法.
按時去什麼節期.
要獻祭等等.
安息日.
什麼都不准做.
各樣的事.
我們都是受制於世俗小學之下.
也是如此.
好像上文說的.

$^{1241}$像少主人一樣.
受制於管家.
一樣的意思.
但要留意第四節.
很大的轉捩點在這裡.
及至時候滿足.
即是耶穌降生在馬槽.
上帝神就差遣他的兒子.
為女子所生.
瑪利亞.
生耶穌.
且生在律法以下.
這就是聖誕節.
不是馬槽歌.
是神學的討論.
甚至將所有人類.
在神的舊約律法.
帶離開.
因為在主耶穌的.
救恩裡面.
我們立即得到釋放.
而且將聖靈賜給我們.
我們不是變成無法無天.
而是我們能夠超越了.
律法表面的要求.
因為聖靈會帶我們.
愛神愛人的心.
所以這裡告訴我們.
事後滿足.
神差遣他的兒子.
為女子所生.
且生在律法以下.
留意第五節.
要把律法以下的人贖出來.
叫我們得著兒子的名份.
即是已經成為一個.
合法的繼承人.
我們就可以享受.
神給我們的產業.
好像以撒一樣.

$^{1281}$每個人都要明白.
首先你要知道.
什麼是令我們在律法下的.
那個捆綁.
那個枷鎖.
但是原來主耶穌的降生.
是帶來一個徹底的.
整個時代的扭轉.
甚至將所有人做不到的事情.
在主耶穌裡面.
我們都可以慢慢的.
朝向這個目標.
這樣去達到.
成為一個得贖的兒子.
我們能夠成為.
神兒子的名份.
身份位份.
這個不只是.
嘉爾文是這樣理解的.
不只是一個地位上的兒子.
所以後面說的.
救恩就有關連.
我們看下面.
和基督論之後.
就是這裡的.
再細分到基督論這裡.
這個就是基督論.
救恩就會說.
被動的兒和主動的兒.
的小小的分別.
這裡我可能說多一點.
這個兒子的名份.
在保羅的理解裡面.
不只是.
只是那個地位上.
一個.
好像給你一個身份證.
新的身份.
不是的.
在保羅的裡面.

$^{1321}$不斷的在我們的.
行為生活裡面.
去體現.
這個兒.
你能夠感覺到.
所以這個就是嘉爾文.
其中救恩論和路德.
馬丁路德有些不同的地方.
對不起,先喝口水.
有點乾了.
(音量注意).
這裡我全部是直接抄基督教要義的.
關乎到基督論.
主耶穌.
加泰書這裡大家要留意.
為什麼保羅要這麼著重.
完全都不浪漫.
路加說的馬槽.
馬尼亞大肚.
去到伯里行.
客店沒有地方.
馬槽那裡生.
很狼狽,有天使向他顯現.
諸如此類,有天君.
唱歌,為什麼保羅.
完全不說那些東西.
他知道耶穌是這樣降生.
肯定,但是為什麼他會說.
及至時候滿足.
神就差顯.
他自己的兒子.
為女子所生,不過留意就是.
要生在律法之下.
因為要救贖.
所有在律法之下.
在救藥神的命定下的人.
要離開律法的咒詛.
成為神的.
得到神的兒子的名分.
所以保羅所探討的.

$^{1361}$不是否定那個浪漫.
而是告訴我們更加務實.
基督帶來的.
整個人類的釋放.
成為神的兒女的身份.
所以成為神的兒女的身份.
就會在我們的生活中.
一會兒我會說救恩.
那個被動的義和主動的義.
我們基督論.
加爾文就很清楚.
三樣東西.
先知,祭司,君王.
基督教要義也有.
我也是直接抄的.
向世人宣告神的心意.
就是主耶穌來這個世界.
祂就是神差顯而來的先知.
祭司肯定就是贖罪.
以祂自己的生命作為祭身.
為人類獻上贖罪祭.
希伯來書就非常多的筆墨.
去形容這個祭司和祭身.
就是耶穌.
第八.
第七,第八,第九章.
就是在說主耶穌就是完美的祭身.
又是那一位的大祭司.
所以加爾文都是根據聖經.
多少是希伯來書.
或者是以賽亞書的53章.
耶穌也是君王.
君王就是統管萬有.
在未有世界之前.
祂已經創造出世界.
來到這個世界.
為我們預備了救恩之後.
祂現在仍然回到天上.
也是再來的君王.
這個是.

$^{1401}$四重福音.
孫迅博士.
都是根據.
這個加爾文神學.
因為孫迅博士是長老中的.
他在加拿大出生的時候.
他父母是回長老會的.
長老中其實就是加爾文後人的分支.
所以孫迅博士.
為什麼宣道會是長老中的.
因為孫迅博士是長老中的.
傳道人.
愛德華王子島.
很小的一個加拿大的地方.
多倫多一點.
後來輾轉去到紐約.
在紐約真是發揚光大.
宣道會的一路的發展.
但是宣道會很重要的四重福音.
就是在宣道會的四重福音.
在宣道會的四重福音.
就是在宣道會的四重福音.
向世人宣告上帝的心意.
我們都是代表著神.
因為我們已經得到救恩.
我們已經被主耶穌帶我們脫離了律法的咒詛.
我們現在就要告訴世人.
用我們的方法.
用我們有限的人生.
向世人宣告上帝的心意.
所以我們都是先知.
我們不是為人贖罪.
我們更加不是祭神.
所以祭司很重要的工作就是祈禱.
我們要將人帶進救恩裡.
所以我們需要為人祈禱.
我們需要報道會祈禱.
我們需要真正的將一些.
我都是很後期才學.
真的有些不同.

$^{1441}$有本祈禱簿.
有些朋友我幾年前為他祈禱.
有名字 有時間.
給自己一個責任感.
給自己原來的上心.
不要小看這些.
記得有意識地再為他祈禱.
因為去年我在港福堂講完道.
執事會的副主席就走上來找我.
問我可不可以為我爸爸媽媽祈禱.
我真的寫下了他們兩個名字.
所以偶然間看到那本祈禱簿.
我看到他們兩個名字 又算是祈禱.
變成你有這個心力.
祈禱就會很立體.
我建議大家.
真的找一本簿.
不要撕下來.
要認真寫下.
最好有記事卡.
和他的進展.
這樣上心.
我覺得神肯定會聽我們的祈禱.
因為這是.
我們認定神.
是想施恩給我們.
而我們成為祭司.
就是我們能夠.
將人帶進神的榮耀裡面.
我們一定有這個位份.
所以我們不要小看.
這些表面很老土的東西.
其實不是.
真的有一本簿是你自己個人和神.
好像一個爭論.
好像一個願望清單.
其實原來是挺過癮的.
真的有空撿一下.
你發現真的挺好的.
這個就是祭司的位份.

$^{1481}$我們不是主耶穌.
我們不可能為人贖罪.
我們更加不是.
那個大祭司.
但我們是可以參與在禱告裡面.
君王就更加是.
統攝萬有.
所以我們真的要有.
這個這樣的位份.
讓人看到.
我們在主的救恩裡面的.
那個幸福.
我們不是要去做一個.
控制其他人.
而是告訴人知道.
我們有萬王之王的.
再來的盼望.
所以我們對世界絕對不會絕望.
我們會有.
看法.
但不會對這個世界的.
沒了有一個絕望.
我們是.
前千的人.
我們.
可能後千的人.
但我們絕對不是無千.
但是我們知道這個世界.
將來是會有主再來.
這件事.
所以其實是一個很大的.
對自己.
對其他人的人生的盼望.
是有一個很好的.
一個模樣.
所以嘉爾文的基督論.
一定是先知祭司君王.
這三樣東西.
但在這三樣東西裡面.
大家看看這個基督論.

$^{1521}$他有專門的章節.
講先知祭司君王.
他列舉了很多東西.
大家聽完我講之後.
如果有興趣可以回去看一看.
舊俗論.
即舊恩.
一定是建立在.
路德 馬丁路德.
的唯獨恩典.
唯獨信心.
我們回頭講1517年的事情.
10月31日的萬聖節.
馬丁路德.
其實他前後.
他自己都知道.
這位神逃避不了.
曾經閃電.
他幾乎死了.
後來就決心做神職人員.
不再做律師.
真的進了修道院.
一路一路追求.
但仍然覺得自己.
和神的距離很遠.
去找主教.
去做一些保屬.
去做一些告解.
都令到他的內心.
不感覺到真正的平安和釋放.
所以.
他一路.
潛移默化的聖經.
還未能夠發揮作用.
因為還是天主教的.
那個的舊俗觀念.
總是要人做一些事.
一個共同的.
舊俗.
不能夠徹底的看到.

$^{1561}$一切都是恩典.
一切都是憑信心.
去領取這些恩典.
我想那個時代的人是不敢這樣想的.
我什麼都不用做.
神的恩典就夠了.
當然.
一直直至到.
神的時間到.
讓他.
突然間.
一路背誦聖經的時候.
好像搞笑.
他一路跪上.
威登堡教堂的時候.
神的道就在他內心.
因為他是一個神父.
他有權在教會裡讀聖經.
讀聖經.
變成他能夠對羅馬書.
相當熟悉.
真是.
一章的十七節.
就在他內心.
徘徊.
激盪他.
為義人.
因信.
得生.
他就真的在那個時候.
再回去研讀羅馬書.
不斷看到保羅講的.
憑信可以獲得神的.
白白救恩.
看到唯有基督的犧牲.
才是真正的.
贖價.
所以對路德來說.
唯獨恩典和唯獨信心.
是他整個的信仰.

$^{1601}$是他整個的.
使他夠膽.
提出九十五條.
和最終要被迫.
離開教廷.
的那個back up.
如果不是恩典.
如果不是信心.
他自己怎麼夠膽一個人挑戰整個教廷呢.
還要帶著一幫人.
離開教廷.
整個天主教的所有的.
核區.
不能再守聖餐.
不能再去彌撒.
信主的人.
現在由這一刻開始.
誰誰.
你不可以再回錦繡堂.
你住在錦繡你不可以去錦繡堂.
你不可以再去教會.
唯有錦繡堂才是教會.
中世紀的時候就是這樣.
所以馬的路德一定是對.
這個唯獨恩典.
有一個很大的.
開竅.
所以他的神學一定是.
恩典和信心.
那嘉爾文怎麼將他發揚光大呢.
這裡就.
看一看這裡.
舊恩論嘉爾文就認為.
不單純是地位上.
的那個工作.
所以他認同路德說.
的人是無助.
必須神的那個介入.
所以這個被動的義.
其實就是神將一個我們本來沒有的.

$^{1641}$我們本來就是得罪神.
敵黨神的.
但是就將主耶穌.
給了我們信的人.
看到這份恩典.
就得到這個被動的義.
就是神給我們的.
不是我們自己的.
但是嘉爾文就不單止停在這裡.
他就再加上那個主動的義.
就是說被動的義.
你明白了.
你擁有了.
其實就是救恩.
但是原來嘉爾文另外一個精髓和路德不同.
在舊屬論裡面.
就是人的主動的義.
就是說原來.
這個我們都蠻多.
這些的普遍.
我自己眼見的.
洗完禮就不回教會的人.
就沒有了主動的義.
當然我不是神.
我就不能夠追蹤他每天.
他的每年.
他還有沒有去教會.
我自己有限的眼光.
我眼見的.
有些都不再回教會了.
洗完禮赤裸裸地都不再回教會的人.
都還有.
其實就失去了你那個主動的義.
原來嘉爾文的那個舊屬論裡面.
是這樣的雙倍.
被動的義我們經歷了.
有了.
其實就是憑信心得到的恩典.
救恩.
但是原來嘉爾文強調的就是.

$^{1681}$我們這個人信了主之後.
都要成長的.
都要不斷被教育.
不斷要在我們的生活裡面.
經驗神的.
主動的義.
就是說我們每天.
用普羅的意思.
一天生前一天.
我們應該比起我們十年.
二十年前對神的認識.
對自己的認識.
應該是有這個主動的義.
其實就是說一個靈命的成長.
就是說.
不是只是一種交易.
我洗了禮了我就完成了.
就不是畢業.
而是.
繼續生命的成長.
在嘉爾文的論述裡面.
就叫做主動的義.
所以這兩樣東西.
看你用什麼說話.
其實就是等於成聖.
這個意思.
所以這個成聖觀.
跟路德是有些不同的地方.
就是在這裡.
我們要記得這個主動的義.
聖經是說很多的.
哥林多後書七章.
我突然想到的經文.
臨後七章.
哥林多後書.
六章尾就是在說.
讀多一點.
六章尾.
十四節.
你們和不信的.

$^{1721}$緣不相配.
不要同父恩厄.
義和不義有什麼相交呢.
光明和黑暗有什麼相通呢.
基督和比列.
就是魔鬼.
或者撒旦.
有什麼相和呢.
信主和不信主有什麼相干呢.
上帝的電.
和偶像有什麼相同呢.
因為我們是永生神的電.
就如神曾經說.
我要在他們中間居住.
在他們中間來往.
這段經文是利未記.
還是文掃記.
我忘記了.
我回去搜索一下.
我要作他們的神.
他們要作我的子民.
又說你們務要從他們中間出來.
與他們分別.
不要沾不見剩的物.
我就收納你們.
本來印了兩段舊約.
一段如果我沒有記錯是文掃記.
還是利未記.
一段其實這一段可能是.
二卓亞書我回去搜索給大家看.
就是十七節.
所以第十八節.
你們作你們的父.
你們要作我的兒女.
這是全能的主說的.
第七章.
就是將這個總結.
親愛的弟兄.
我們既有這等應許.
就是上文說的.

$^{1761}$我們是和那些不再信神的地上的人.
不可以再成為.
朋友.
來一起.
做那些不好的事.
他說.
我們既然有這些應許.
神要收納我們.
全能的主要收納我們.
所以這是應許.
就當潔淨自己.
除去身體靈魂一切的污穢.
身體靈魂一切的污穢.
敬畏神.
得以成聖.
所以嘉爾盟約它整理了之後.
就是他所謂的主動的兒.
就是我們必須要不斷.
在信了主之後的成聖.
有意識地成聖.
成為主動的兒.
是我們自己要做的.
但不是換取救恩.
因為救恩是被動的兒給了我們.
分不分得到這一點.
其實就是我們先得到.
這個被動的兒.
神的兒.
加給我們這不義的人.
但是信了主之後.
我們就要追求成聖的主動的兒.
路德就是.
我們都知道.
他對聖經的.
尊重.
看到恩典和信心的.
那個並行.
這個是馬丁路德.
他就沒有再提出另外的了.
但是嘉爾文.

$^{1801}$這裡就真的最重要了.
唯獨恩典 唯獨信心 唯獨聖經.
全部沒有問題.
但是沒有了這一點.
嘉爾文看人.
人就是徹底的無助.
需要神的救贖.
救恩.
無條件的.
為我們犧牲.
將我們帶進神對神的認識的幸福.
接著.
這裡就是嘉爾文最精彩的論述.
就是這裡.
首先是他創的.
其實保羅早就說了.
腓立比書三章.
都說.
無論生或死.
總要基督在我身上造成顯大.
保羅的金句.
其實保羅早就說了.
不過.
如果能夠成為一個神學的論述.
就是嘉爾文某程度的首創.
我將整個拉丁文都說出來.
Solio de Gloria.
Gloria我們知道.
就是那個榮耀.
有時候你會看到中世紀的作品.
譬如.
這個.
一些很重要的神學作品.
譬如.
彌賽亞.
那個神曲.
這個.
我忘記作者的名字了.
他們的.
手稿.

$^{1841}$全部寫完之後.
都有這個唯獨榮耀神.
就是說我這麼出色.
這麼原創的作品.
最終都要歸給神.
神得到一切的榮耀.
所以不要當是我們自己的口號.
我們要當是我們人生.
我們做每一件事.
的目標.
我們做每一件事.
呈獻給造物主.
他得到一切的榮耀.
我們有我們的成就.
都是應該歸給神得到榮耀.
所以我們有成就.
是對的沒有問題.
不過要歸給神.
真是唯獨神得到榮耀.
或者人.
這個就是嘉爾文神學.
或者後期的神學家.
都很重要就是.
很重要的一個.
最單純的目標.
就是要榮耀神.
在我們的人生榮耀神.
所以再仔細一點.
就是在我們的職場榮耀神.
在我們的家庭榮耀神.
某程度上.
真是說了最重要.
我們人得到最大的.
那個福樂.
如果不是.
人就很可怕的了.
有資源的就會用他的資源.
去可能自肥.
為所欲為.
或者其他的做惡事.

$^{1881}$沒有資源的只會每天就.
唉聲嘆氣我的人都存在.
都不知道為了什麼.
就變成浪費了很多的時間.
去尋問自己.
怎麼可能有答案呢.
但是神學給我們的.
比如中世紀給我們的.
就是很簡單原來我們人是有價值.
原來我們是得到神的救恩.
祂的啟示.
在我們進入祂的幸福裡面.
祂為我們預備的幸福裡面.
然後我們就在.
這個幸福裡面.
有我們主動的意義去成長.
其實最終.
就是這個唯獨榮耀神.
所以我們都要.
每天起床.
你自己都跟自己說.
現在的心理學就說.
你自己可以的看著鏡子.
然後你就去做推銷.
這些心理學好像說笑.
我今天起床就算我今天有病.
我的病都是要榮耀神的.
這樣其實不是.
accusing神.
其實是讓我們裡面的潛能.
神給我們的發揮到極致.
因為我不是為自己.
我真是為神的榮耀.
這個是成聖觀.
嘉義民的那個地方.
看看還有沒有時間.
預定論這裡.
開始要說了.
最後這裡.
大家都知道.

$^{1921}$真是不是要.
挖苦那些沒有被揀選的人.
搞錯了.
以為.
你們說了.
那些沒有被揀選的.
注定一定是滅亡.
因為這個是無條件揀選.
但你就是沒有揀選我.
我就注定滅亡.
其實真是會錯意.
或者無辜.
因為是為了那些離開教廷.
大公教會.
天主教的信徒.
現在跟隨改革.
成為了.
基督徒.
徹徹底底不可以再去離散.
不可以再去領聖餐.
教會說.
你們不回來.
你就沒有救恩.
現在嘉義民就將他整理成一套.
幾標準的神學.
人的傳言敗壞.
這個人論很清楚.
所以需要一個無條件的.
神啟示的一個救恩.
就是主耶穌.
唯一的人類的救贖.
所以這個是無條件的揀選.
神沒有因為哪個人.
但因為有限度的救贖.
就不是神.
不能夠傳言救每一個人.
而是.
不是他不能夠.
是只不過人自己的限制.
所以真是會有一個.

$^{1961}$這個是他的神學的那個package.
所以不可以自己抽開一點.
去辯論的.
是因為傳言敗壞.
是因為神的愛的無條件.
這樣來揀選.
但是有些人就argue.
為什麼有些人就沒有被揀選呢.
這個我們答不到的.
因為只能夠在這樣的.
這個邏輯裡面.
來引申這個.
預定論.
記住整個package.
是為著那些離開了.
教廷的信徒.
他們的內心的平安.
既然我現在是跟隨耶穌.
跟隨你們神學家講的.
耶穌的救恩.
那就是說.
我已經是清楚自己是被揀選.
無條件被揀選的了.
那所以呢.
這個有限的救贖.
就不是我們這班人.
我們被揀選.
神一定會救我們到底的.
但有些人是.
不要.
那就變成某程度上.
就是神沒有揀選他們.
但是這個是後人將這些.
去到極致就帶來很大的爭論.
但是整體來講.
其實是互相的補足的.
那個敗壞.
不可以自己救自己.
神透過啟示.
透過主耶穌.

$^{2001}$救我們.
而且是無條件的.
那我們是白白得到這個揀選的了.
每個人都可以的.
但是有些人是拒絕.
那就不能夠.
每一個人都得到救恩.
應該是.
德國後來的神學家就說.
神透過不同的方法.
最終所有人都得救.
那這個就肯定跟嘉爾文神學.
有一個衝突的了.
因為一定是有些人是不行的.
所以不是神.
用另外的方法.
透過不同的宗教救他們.
是不可以的.
我自己是嘉爾文主義的.
嘉爾文神學的.
我在神羅院的時候.
都辯論的.
我仍然相信嘉爾文神學的.
有限度的救贖.
還有後面就說.
既然是這個.
有限度的救贖.
這個不可拒絕的恩典.
既然神揀選了我們.
我們沒有人可以拒絕的.
這個是一個.
又是一個一男子的.
整個package(組合).
所以.
這個無條件的揀選.
背後就是一次得救.
永遠得救的那個演繹.
就是.
信徒永遠蒙神的保守.
真是被揀選的.

$^{2041}$真是基於.
目樣的考慮.
是對那個時代.
所以後來的人.
這樣將他演繹.
就有些不公允.
因為跟隨嘉爾文的人.
就是要問一個問題.
我們是不是真的可以.
得救上天堂.
嘉爾文就說.
你們被揀選 而且是無條件被揀選.
被神的.
那個預定.
所以你們在這個救恩.
在主耶穌裡面.
一定蒙神的保守.
所以後來有些人就.
比較客觀些去修正.
就是一次得救.
不是永遠得救.
是一次得救 一直相信.
就必然得救.
這樣才是精髓.
不然就很機械性.
就是一次得救.
跟著可以拜拜.
就永遠得救.
令到很多人就產生一個救病.
和一個論斷的懷疑.
就變成失去了.
對那個.
目樣上的安慰.
和那個.
令到信徒.
在中世紀.
跟隨改革的時候.
這種惶恐.
這種怕著真的沒有.
救恩的失落.

$^{2081}$所以是對他們.
一個目樣上的安慰.
還有一兩個.
還沒有說完下面.
已經47了.
都是簡單來說.
我能不能夠現在.
稍微離開那個.
版面.
就將那個檔案.
扔給大家.
我應該放上去.
算了 放下吧 我回去再做.
教會論.
很實際.
當時是三萬多人的.
日內瓦.
他的自區的牧區.
整個瑞士.
跟加爾文主義.
其實不是叫加爾文主義.
是跟加爾文的目樣.
願意接受他的管治.
他發展的教會論.
其實很豐富的.
不單只是恆常的目樣.
原來.
必須要教導聖經的真理.
是要讓信徒.
得到聖經的教導.
充分的理解聖經.
這樣.
這是教育.
和要陪伴成長.
所以很重要.
就是教會存在.
就是要將神的真道.
似乎神.
設立教會.
就像一個機制.

$^{2121}$讓神的旨意.
聖經.
讓信徒一代一代.
能夠充分的理解神的.
心意.
所以你想想.
只是這個系統神學.
加爾文神學.
如果我們能夠充分的了解.
每一個題目.
某程度對我們來說.
有些事是不會過時的.
神的救恩和神的存在.
這些不會過時的.
我們能夠明白之後.
就會在我們自己的時代.
如何與神同行.
所以教會論很重要.
精髓.
是存在在地上.
作為神.
教導信徒.
的一個地方.
當然透過牧者.
所以牧師和傳道才可以講道.
來將這些真理.
有系統地.
教導.
教導完之後.
要觀察和陪伴牧師的成長.
其實牧師在做的每天工作.
就是這些.
牧養 教導.
譬如探病.
今天突然又要去探病.
那個姐妹.
失勢.
放針.
聽到媽媽進了醫院.
中風.

$^{2161}$跟著告訴我她去了屯門醫院.
她自己走去屯門醫院.
誰知道原來在天水圍醫院.
她沒有看WhatsApp.
妹妹告訴她.
她一直以為是屯門醫院.
去到之後我差不多到了.
下車了.
跟著告訴我天水圍醫院才是真的.
又要坐車.
牧養和探訪.
都是牧者要做的事.
在牧神無主的時候.
又安慰牠.
又給牠.
一個信心等等.
牧者就是要做這些.
其實就是將神的.
恩典 真理.
透過不同的方式.
去給.
弟兄姊妹能夠.
接收 更重要的是.
建立一個穩定.
機制.
為什麼我們要固定.
有主一學 為什麼有兒童崇拜.
為什麼有兒童的主一學.
為什麼要按不同的程度.
一級一級的上.
就是要建立一個穩定的機制.
不是一種即興式的.
隨意式的.
這樣就不行的.
因為這個堂子.
大家是穩定地回來.
去得到教育和牧養的陪伴.
所以這個嘉爾文.
他對教會存在.
當然他當時.

$^{2201}$整個日內瓦就是教會.
就是那個國家.
這個跟社會觀.
有些相近.
有些不同的就是.
因為始終.
不同的人的不同步伐.
信仰.
所以是需要觀.
社會就涉及到很多的.
行政.
不同的體制.
與時並進.
我們未必可以參考的.
我們根本不會 除非我們整個感受區.
都有幾萬人.
整個感受就是要由.
感受堂治理的.
勵士牧師就是那個.
那個主要的領導人.
長老應該是斯勞永富弟兄.
兩個共同管理.
這個感受的社區.
這個就是當時的日內瓦.
的那個場景.
但今天沒有可能任何一個地方都不會.
所以那個社會觀.
我們有一點點了解.
就已經夠了.
因為時代已經不同了.
經濟觀就不是.
經濟觀這裡.
剛才我們說了一點.
一個信徒必須在信主之後.
認識自己.
原本是一無所有.
原本本來是神的眼中.
是敵黨.
是敗壞.
是有罪 拒絕 啟示.

$^{2241}$拒絕神的存在.
到現在我們明白了.
原來這個經濟觀是帶給人.
一個最高的.
那個他的身份帶來的.
自己自我意識.
的那個存在.
就是成為神的兒女.
之後呢.
在生活裡面去榮耀神.
所以其中就一定是.
那個物質上面.
必須要簡樸.
盡量.
積蓄然後去.
周濟有需要的人.
那個同理心和那個愛心.
那麼詳細.
我今晚給的那個檔案.
大家就可以看了.
四至五頁的.
你可以抽一些題目來看.
應該今天我們.
因為還有一兩個.
我就沒有時間再做了.
就是關乎到文化和教育.
嘉爾文也有一些見解.
所以他很著重.
大學教育就是其中一個重點.
這個就是.
某程度將嘉爾文神學.
當然後人.
看的就是那個.
預定論.
但是其實還有.
其他很重要的.
那個.
就是因為聖經神學.
而衍生出來的.
他的那些.

$^{2281}$有系統的.
那個對神對人對救恩.
那個的辭定.
是我們今天福音派.
差不多大部份都是這樣.
所以對我們來說.
不陌生.
但是如果你對其他的宗派.
就很否定嘉爾文這一套.
因為覺得很壓迫.
覺得好像人一無是處.
這個就.
當然另外的宗派另外的神學.
但是我們嘉爾文神學.
福音派就一定是.
很強的人的無助.
神的白白的救恩.
和人在救恩裡面.
有一個新的身份.
有那個力量.
所以又不可以.
變成個人的救恩.
我們應該是一個整體去看.
個人我們自己的存在.
怎樣來.
活在神面前.
今晚我就講到這裡.
我們還有.
一堂就是過年.
我又很想.
過年之後.
可能會停一次.
但是.
我更加想多一堂至兩堂.
講一講.
嘉爾文神學.
帶來的.
即時的影響.
例如蘇格蘭的大復興.
就是嘉爾文神學的.

$^{2321}$直接的影響.
因為他的學生就將這套東西.
帶回蘇格蘭.
就帶來大復興.
後來就影響了英國.
當然真的承諾了.
我們大家一起要去探討一下.
過去發生了.
但是.
到今天我們怎樣理解.
宗教改革呢.
就不只是嘉爾文.
宗教改革來到今天.
我們的精髓.
我們需要去傳承.
需要去捍衛.
這個都希望我們能夠一起去思想.
所以可能會多一至兩堂.
給大家.
都有空間去.
看看宗教改革.
和我們的關係.
所以到時.
我再看看情況.
但是就要開始問.
回到聖經.
一來.
那個預備不用那麼多.
因為我要整理筆記.
看書.
用很多時間.
都是要回到聖經.
不是說時間不用那麼多.
而是相對我都相對熟.
我就不需要太多的筆記.
要準備.
大家要問.
新約還是舊約.
哪一卷書大家是很想的.
一方面看看我的時間.

$^{2361}$夠不夠準備.
一方面看不太熟的經卷.
我就需要很多時間.
看一些筆記等等.
但如果一些我喜歡和熟悉的.
就相對容易一些.
譬如希伯來書我們肯定沒說過.
肯定是重要的書卷.
曾經想過.
會不會因為那時羅馬書.
我們完全沒有錄音.
是第一卷的.
我們的.
星期四晚的查經班.
羅馬書我們從來沒有錄過音.
所以差不多我們.
有一年.
的講道.
都浪費了.
有沒有需要我們.
重新講一次羅馬書.
給一些從來沒聽過的人.
都很有系統的了解.
保羅的整個救恩的.
進程.
當然有些人聽完.
可能就未必會再來.
有些少限制.
譬如希伯來書.
舊約.
大家想不想看民數記.
或者利未記.
都是很好看的經文.
可能大家從來都不敢看的利未記.
其實都有很多東西值得我們思想.
有好幾卷這些書卷.
但你可以在WhatsApp.
私下給我.
給我了解你們想有什麼意見.
但無論.

$^{2401}$下一次.
我們完了終界改革之後.
就要回到書卷那裡.
有系統的研讀.
我們一起禱告.
感恩再一次我們能夠.
很有系統的整理了.
加爾文神學.
當然我們都是一輪半找.
去將一些重點.
去涉獵.
我們一輪不在那個時代.
很難能夠回到原著.
但多謝主讓我們看到.
加爾文對聖經的.
尊重,捍衛.
和著重.
我相信絕對不是那個時代.
沒有人否定聖經.
沒有人來抗議.
和挑戰聖經.
但是我們看到.
這位神學家.
因為他對聖經的尊重.
和捍衛.
和實踐的時候.
就帶來下一代.
宗教的復興.
來帶來.
其他國家.
蘇格蘭,英國.
後來輾轉去到美國.
都帶來一個很大的.
宗教上的.
大爆發.
很多人信主.
後來去到一個.
普世的宣教運動.
盼望我們都能夠.
透過馬丁路德.

$^{2441}$透過加爾文.
我們都能夠再一次.
思想宗教改革.
帶給我們的福氣.
對真理的認識.
對救恩的明白.
以致我們都能夠.
一起在神的道上.
去追求去成長.
再一次多謝主讓我們一起.
有這段時間的學習.
感恩奉耶穌基督的名.
\newpage



\section{}
\label{sec:5Lw6EDNBsNI}
\textbf{2025-02-06 周四晚查經班《改革之路(9):諾克斯(一)》 講員:曾裔貴牧師}
\newline
\newline
連結: \href{https://youtube.com/watch?v=5Lw6EDNBsNI}{\texttt{https://youtube.com/watch?v=5Lw6EDNBsNI}} ~~~~ 語音日期: 2025-02-07
\newline
\newline
\hyperref[sec:Dn07Lw4F37U]{\small{< < < PREV SERMON < < <}}
~
\hyperref[sec:index_chronic]{\small{[返順時目]}}
~
\hyperref[sec:index_scriptual]{\small{[返順卷目]}}
~
\hyperref[sec:q7QZXuHzcds]{\small{> > > NEXT SERMON > > >}}
\newline
\newline
$^{1}$再一次感恩天父讓我們在晚上再看這個宗教改革.
我們能夠從嘉義民的學生.
這位洛克斯來到去.
從蘇格蘭的改革來到去看宗教改革帶來的歐洲的影響.
還有一些值得我們去深思的事情.
盼望今晚我們講完這個大綱之後.
又進入一個大家可以去思想.
去反思宗教改革對我們大家的影響.
我們都用了第九堂了.
盼望我們可以無論馬丁路德.
這個伊拉斯姆和嘉義民.
再加上他的一個重要的學生.
將整個嘉義民主義長老中的長老制傳到整個英國.
使我們能夠明白.
我們仰望在天父的手中.
奉耶穌基督的名.
阿們.
都真的要停了.
因為希望用十堂.
今天是第九堂.
其實主要環繞著三個宗教改革.
一定要認識的重要的人物.
當然之前提過有持運理.
或者再遠一點的.
有丁道義的翻譯聖經.
這些是一些比較支援宗教改革.
但是主要的宗教改革.
一定是由馬丁路德.
他的九十五條.
引發一個宗教討論.
而帶來天主教的教廷.
逼他要收回九十五條.
他只是希望有對話.
希望不是想推翻.
或者要來攻擊教廷.
或者他所敬重的一些神職人員.
但是沒辦法.
因為勢不兩立.
他只能夠帶領跟隨他的人.
對改革.

$^{41}$回到聖經的原本.
來做一個更新.
所以由一個溫和的.
引發了整個思潮.
我們第二,第三堂都說過.
既然教廷可以被挑戰.
而且是由神學方面.
由本身是神職人員的馬丁路德.
去挑戰教皇的權威.
那就是說.
中世紀的文字就因為這樣而打開了.
是一發不可收拾的.
如何能夠保持這個.
是一個穩定,溫和.
和不是革命.
改革和革命是兩碼子的事.
所以你可以想像到.
馬丁路德,加爾文.
和今晚我們說的洛克斯.
他們都是朝著一個很重要.
回到聖經這個大主題.
來進行改革.
所以你可以想像到.
民間的一發不可收拾的革命.
只會帶來一個拆毀性.
玉石俱焚.
甚至是聖靈塗炭.
我們看過馬丁路德同時期的農民戰爭.
那些田農自己制定了一些條款.
要強迫這些貴族和當時的君王.
要來接納.
不再納十分之一歲給教廷.
他們的教區.
要驅逐這些神職人員.
要來燒教堂.
要拿回自己一些教產等等.
所以最終都是朝著政權的逼迫.
鎮壓.
你可以想像到.
其實是幾個路線在進行.

$^{81}$一直都沒有停止過.
所以不是一個靜態的在神學上.
或者在神學院裡面.
或者在當時的修道院裡面.
大家就束膝夜談.
大家來看看誰比較適合就寫下.
不是的.
是一個關乎到生和死.
關乎到很多人到底是要跟隨宗教改革的路線.
還是要跟隨天主教教廷的規範.
和千萬不要離開教廷.
今晚我們會涉獵到這個部分.
就是法國.
所以我今天下午會看多些資料.
越畫越深.
畫到1789年.
其實我們在說16世紀.
到1789年18世紀的開始.
17世紀末期.
就是法國的大革命.
1789年的7月14日.
大家應該都熟悉的.
7月14日本身就是法國的國慶.
為甚麼用7月14日做國慶呢.
就是1789年的當時的平民.
發動大規模的革命.
要推翻當時的路易十六.
法國的君主.
很出名的.
佔據了法國的巴士的監獄.
然後真的帶來聖靈塗炭.
路易十六被人推上斷頭台.
殺死了.
我們讀中學歷史都讀過.
今晚我們會說.
在整個宗教改革裡面.
法國是完全沒有受影響.
沒有受影響即是很大件事.
即是一直無論如何都站在天主教封閉裡面.
沒錯一定是基督教.

$^{121}$一定是加爾文 馬丁路德你們的錯.
所以法國是一直鎮壓那些基督徒.
所以剛才已經說了.
改革其實是一個不斷變遷的思潮.
有宗教改革的神學的修正.
有神職人員即是牧者.
做一個帶領.
希望回到聖經.
從來沒有想過要推翻哪一個政權或哪一個人.
但在民間有一些長期受到壓迫的田農.
這些低下階層.
在一個這麼大的文字突然間開了之後.
即是說原來教廷是可以挑戰的.
因為借助了舉格家他們的歧視.
原來教廷是有一些需要修正的.
而他們不肯修正.
造就了民間對教廷產生的那種仇恨.
所以一直積累.
最終爆發到臨界點.
真是帶來一個真正的革命.
要流血的革命.
所以這個就是法國堅持不做宗教改革.
而帶來的後果.
當然在一個沒有神的.
特別是所謂民意民主.
他們強調的是民主.
因為中世紀是一個封閉的社會.
受壓制是被天主教的教廷,神學.
他們的講解的壟斷來壓制自由.
所以在這些民意一開啟的時候.
而沒有神學.
沒有敬虔的領袖帶人回到聖經.
那個恐怖就是.
人只會為著很多切身的利益.
沒有任何道德倫理可言.
只需要達到眾人認為的目標.
就會做出對當時的君主.
當時的神職人員.
天主教神職人員進行所謂大報復.
所以由1789到1790年代.

$^{161}$甚至18世紀開始.
最終拿破崙.
我們所認識的拿破崙.
真是帶來法國的統一.
所以有差不多十多年.
接近二十年.
法國走上流血的革命之路.
所以不要以為很簡單.
我們看到宗教改革.
他們一心只希望.
將人帶回神的真理.
對神的敬畏.
沒有想過要為革命而革命.
所以這是一個很根本的思想.
為什麼我要用洛克斯來作討論呢?.
一方面是到底嘉爾文在瑞士這麼偉大.
這麼有影響力.
瑞士和當時的歐洲.
德國有這麼大的改革.
但其他歐洲的地方呢?.
或者神學的改革能否蔓延到其他地方.
帶來一些變化呢?.
歷史都是這樣發展的.
因為這位洛克斯.
相當重要是嘉爾文一個非常.
不可以說是入室弟子.
是一個非常好的學生.
差不多是將嘉爾文神學的精髓.
放在整個蘇格蘭作一個改革.
又帶來一個整個蘇格蘭國家的大復興.
所以這個值得我們透過.
在另外一個地方.
另外一個時代的神學家.
其實都是差不多那個時代.
因為嘉爾文在156幾年死亡.
他也是在157幾年死亡.
所以是同時代的.
但是對於洛克斯來說.
他在年輕的時候.
未認識嘉爾文之前.

$^{201}$所受的苦是相當大的.
所以今晚有機會.
如果英國有興趣.
你又可以聽一下.
因為他做了很多功課.
因為有很多資料.
只是問深層都問了很多.
要他回答的東西.
才做了很多功課.
所以如果你對英國蘇格蘭的歷史有興趣.
今晚可以和大家交換一下資料.
我們今晚說的是洛克斯John Locke.
KNOX.
很出名的.
蘇格蘭現在還記不記得他呢.
不知道.
但是他帶來蘇格蘭的更正教改革.
是非常之重大的.
剛才我開始之前.
就訂了一個蘇格蘭的信仰宣言.
就是由洛克斯和他的同伴所制定的.
有25條的.
大家可以回去看.
不長的.
差不多是將嘉爾文神學的精髓.
都原汁原味地.
嘉爾的在他們的地方去應用.
我們了解這個改革之前.
其實我們就要了解.
在當時的英格蘭.
英倫三島.
其實是一個體系的國家.
大家可能不知道聖安德魯斯.
其實就是耶穌基督的門徒.
安德烈.
就是彼得的哥哥.
安德烈.
據聞他將福音帶到蘇格蘭.
所以聖安德魯斯已經成為今天蘇格蘭國家的守護神.
當然有些可能是天主教的迷信的東西.

$^{241}$就看這個人成為了可以供奉.
可以去祈禱的一個守護神.
但是無論如何.
聖安德魯斯你問蘇格蘭人.
沒有人不認識的.
安德烈.
所以蘇格蘭的聖安德烈日.
就在11月30日.
蘇格蘭的國旗.
做了很多功課.
今天的英國國旗.
不是這個淺藍色.
起初是淺藍色的.
不是深藍色的.
斜的十字架.
安德烈據聞.
當時他被殺害之前.
要釘十字架.
他就請求不要直接釘.
請你幫我打斜釘.
所以那些舉止獸就接納他的請求.
所以是一個斜的十字架.
是蘇格蘭的國旗.
這些冷知識.
紀念安德烈.
是他們國家的守護神.
所以每年的聖安德烈日.
都是在11月30日.
所以如果你剛好去到.
體驗一下聖安德烈日有什麼紀念.
這個就是11月30日.
但是其實因為耶穌的時代.
是主後的一百年.
現在我們說的是一五幾年.
所以一直蘇格蘭.
都沒有真真正正的被基督教的改變.
一直都是服從整個大公教會.
由317年君士坦丁開始.
所謂大公教會.
Catholic就是天主教的大公教會.

$^{281}$所以整個歐洲都是天主教的統治.
去到馬丁路德的年代.
就開始所謂有更正教.
我們能不能回到.
不要只是跟人的宗教傳統.
而我們回到跟隨聖經呢.
這樣的思想就帶來宗教改革.
所以真正帶來蘇格蘭的大復興.
其實就是洛克斯.
所以今天的蘇格蘭人.
對洛克斯應該都未必陌生.
但對聖安德烈就肯定.
所以他們的大學.
聖安德烈大學在蘇格蘭就相當出名.
又是紀念安德烈.
我們說的洛克斯.
1514年出生.
相當出色的一個年青人.
所以他1536年.
在聖安德烈大學畢業.
又是相當敬虔的一個人.
很有趣.
1543年的時候.
他就真的讀到基督教要義.
而自己向神悔改.
所以基督教要義.
其實我都說了.
我桌上有中文版.
已經翻譯了很久.
相當好看.
因為很容易看.
很有條理.
將神論,基督論,聖靈論.
和一些信徒.
應該要注意的.
例如關於祈禱.
很詳盡的關於祈禱.
有很多方面的.
所以兩大冊的巨著.
洛克斯在認識加爾文之前.

$^{321}$他就拜讀了這本基督教要義.
而他自己悔改.
所以更新的種子.
就在他的內心裡.
他本身不是神職人員.
他本身只是一個很普通的人.
當然父母應該有一定的學問.
不是一些田農等等.
很低下階層,不識字的人.
不然都不能夠進入這麼出色的大學.
所以他說他50歲.
看那些資料.
我自己都覺得不太好意思.
50歲才學希伯來文.
不懂的.
學到一個地步.
他可以翻譯.
將希伯來文翻譯.
可能就是英文.
可能是蘇格蘭文.
你可以想像到.
他們這些人.
為什麼可以這麼專注.
為什麼可以沒有什麼科技的幫助.
連電燈都不多的油燈年代.
他們可以這麼專心去做一件事.
最後因為天主教對他正式的拘捕.
他就要逃亡.
在這個過程有相當長的時間.
由1547年開始.
一直到1554年.
他最終要去到日內瓦逃亡.
原因是什麼呢?.
因為某程度他真的生不逢時.
因為當時有三個瑪利.
大家聽過一個叫做Buddy Mary.
很殘暴的瑪利皇后.
瑪利一世.
蘇格蘭有兩個瑪利皇后.
英格蘭又有一個Buddy Mary.

$^{361}$即血瑪利.
據聞將基督徒神職人員.
那些傳道人有三千人去殘殺.
導致有一個名字叫做Buddy Mary.
所以你可以想像到.
我們今天說的洛克斯.
他在認識加爾文之前.
在逃亡到日內瓦之前.
在英國.
特別是蘇格蘭.
因為他在蘇格蘭出生.
他成長.
在這本書之外.
他認識了基督教要義.
認識了我們要回到神的真理.
在他內心有這個要改革.
開始他認同更正教.
認同加爾文他們的宗教改革.
或者馬丁路德的宗教改革.
所以導致他在英格蘭.
特別是蘇格蘭.
受到迫害.
我們開始說他的生平.
他的年青時代.
受到迫害相當長時間.
當時蘇格蘭一位基督教領袖.
就是George Whistler.
因為蘇格蘭文和英文有點不同.
是Isater.
他在1546年殉道.
被天主教的Python.
Python主教來到.
有宗教裁判所.
判斷基督教是異端.
要行刑.
所以Isater當時被天主教的Python主教迫害.
因為洛克斯非常認同更正教的宗教改革.
因為基督教要義在他內心已經產生影響.
所以我們透過周邊的人.
再看加爾文的神學.

$^{401}$對當時歐洲的影響.
作為一個小小的輻射.
讓我們看看神如何透過馬丁路德.
第二代的人就是加爾文.
加爾文很重要的角色.
工作的功用就是.
令到改革一定要在神學和真理的軌道上運行.
如果失去軌道.
就變成革命.
變成無日無之的.
我覺得你錯 我覺得你做得不對.
你長期壓迫我們農民.
我們要抗爭.
無日無之的 只會你死我亡.
但是他們這班神學家.
所以加爾文制定的基督教要義.
和一些信仰宣言.
令到整個歐洲的人.
那個革命不是為了革命.
是為了回到神的真理規範裡.
希望和神有第一手的接觸.
不需要透過代理人.
因為中世紀代理人的制度和觀念.
已經根深蒂固.
沒有人敢自己向神祈禱.
也不懂 也不會教你.
所以你可以說是一個很嶄新的想法.
人可以直接向神祈禱.
對我們今天21世紀當然沒有這個問題.
我們的問題是我們很少祈禱.
不是我們怕自己向神祈禱.
沒有人回教會會覺得.
可以啊 我自己向神祈禱.
很少會這樣.
但是中世紀就真是代理人制度.
所以在這樣的情況下.
這個洛克斯.
他沒有認識加爾文.
只是看他的書.
但是他認同基督教的改革.

$^{441}$所以他就站在了.
當時的伊沙特領袖的一邊.
自然就是一個和一齊的革命.
交織的一定是壓迫.
這個時候你可以想像.
1546年被殺.
當時的洛克斯在蘇格蘭.
已經有相當的影響力.
據聞他的講道和著作.
都很受蘇格蘭人的歡迎.
所以當時有很多人對他的推崇.
甚至乎來抬舉他.
成為當時他們的大領人的領袖.
他在眾人的鼓勵和推舉下.
他就義不容辭.
他都認同中國的改革.
所以他在1547年.
成為了蘇格蘭教會的一位領袖.
在這個時候.
我們就要講三個瑪麗.
先講他太太.
這個瑪朱麗.
大家知道這個應該都是.
有可能是法國人.
不過我要先查一下.
洛克斯是蘇格蘭人.
所以他一生.
雖然在日內瓦有非常好的學習.
和加爾文成為終身的好朋友.
所謂終身其實只有幾年.
那時候大家差不多死了.
所以在日內瓦相當受到保護.
但是很寶貴的就是.
他始終是流著蘇格蘭人的血.
他覺得他還是要回去蘇格蘭.
去帶領教會.
這些我想是牧者的情操.
不會逃亡之後.
就安於自己個人的安全.
就不回去.

$^{481}$所以起碼他的太太.
都不是太長時間離世.
所以有甜蜜的好的時光.
其實都很短暫.
所以你可以想像到.
這個洛克斯.
他一生所受的迫害.
多過馬丁路德.
馬丁路德都沒有什麼被迫害.
加爾文就更加少.
只不過是言語上.
起初日內瓦不歡迎他.
但是當他第二次回到日內瓦之後.
差不多是統領著整個.
日內瓦的國家的城邦.
甚至乎他自己加爾文.
都定了某個人的罪.
要行火刑.
變相加爾文自己.
都被人日後的後代來詬病他.
他自己都做一些對人的.
又不是迫害.
而是那個人決心不肯悔改.
在聖餐那件事情上.
怎樣都不肯接受正統的.
加爾文他們的教導.
後來被定性為異端.
就要死的.
所以這個也可以說是意味著.
馬丁路德,加爾文.
他們是沒有受過肉身的傷害.
但是對洛克斯來說.
由他認識了基督教的改革.
看了基督教要義.
而他自己向神悔改之後.
就很認同剛才說的那位.
艾沙特這位蘇格蘭的領袖.
所以在這樣的情況下.
洛克斯就成為了蘇格蘭的領袖.
由這一刻開始.

$^{521}$1547年開始.
很無奈地當時有幾個.
三個Mary.
是極端擁戴天主教的.
而且不約而同.
當時歐洲其實很.
以前的歐洲一體化比現在更加厲害.
她本身就是法國人.
這個瑪麗德吉斯.
十六世紀的一個女王.
蘇格蘭.
為什麼她會是蘇格蘭呢.
其實她是法國人.
因為她改嫁之後.
嫁了蘇格蘭的查理五世.
蘇格蘭的國王.
對不起.
詹姆斯五世.
她本身是法國人.
輾轉.
她起初嫁第一任的丈夫.
就是這個朗格維.
後來再改嫁.
嫁了詹姆斯五世.
所以非常非常擁戴天主教.
又有權力.
不知為何當時女性的社會地位這麼高.
女性都可以做女王.
帶來對基督徒.
對基督教一個大規模的迫害.
殺害.
所以就捉拿了洛克斯.
有19個月歷史記載.
捉了他去法國.
做那些奴隸.
那些奴隸在船艙.
我不知道是否有幾層.
類似這樣.
在艙底裡面做那些爬船.
奴役了19個月.

$^{561}$後來才輾轉放了他.
由馬麗·德吉斯.
這位法國女子.
後來嫁了詹姆斯五世.
成為了蘇格蘭女王.
正值蘇格蘭宗教改革.
正正在1547年.
洛克斯接任了.
這個艾莎特領袖的位置.
所以首當其衝.
洛克斯也受到這幾位女士的逼迫.
輾轉之後.
就是他的女兒.
馬麗·德吉斯一世.
名字比較繞口.
其實是法文.
馬麗·德吉斯.
當時的蘇格蘭.
女王是她的第一個女王.
因為往後也有.
例如她的女兒馬麗·德吉斯一世.
也是蘇格蘭的女王.
她和英格蘭的馬麗·德吉斯皇后.
有親戚關係.
就是她的表姑.
因為那個族譜相當長.
一來不熟.
特別是中世紀.
所以你可以想像到.
這些都是當時的構想圖.
因為當時沒有拍攝相機.
就一直傳下來.
大概是這個樣子.
你可以想像到.
在這樣的情況下.
她的女兒馬麗·德吉斯一世.
同樣在蘇格蘭.
逼迫洛克斯.
再加上這個更加嚴重.
這個樣子也有些邪惡.

$^{601}$血腥馬麗·德吉斯.
不兼Mary.
大家可能聽過.
在1553年登基.
她是都鐸王朝的繼承人.
所以她登基之後.
大家可以猜到.
1553年她登基.
所以洛克斯是最後一個基督徒.
要逃亡到日內瓦.
在這樣的情況下.
這三個女王.
蘇格蘭和英格蘭.
不兼Mary是英格蘭的.
好像有一套電影叫做《瑪麗皇后》.
這幾年的.
所以值得看.
可以對當時的教廷.
對當時的皇宮.
或者他們的親戚關係了解.
因為他們做了很多資料搜集.
可能值得看.
因為我找資料的時候.
也看到有這些電影介紹.
所以可能我有空.
我再找回這一類的資料.
可能再遲些補充給大家知道.
輾轉這三個女士.
對基督徒蘇格蘭和英格蘭的迫害.
是非常之大的.
是令到.
要補充一點.
英格蘭和蘇格蘭都是.
他們的父輩.
有些是同情基督教的.
同情宗教改革.
但是去到這三位女士.
極端擁戴天主教.
不知什麼原因.
當然是想鞏固自己的權力.

$^{641}$所以在他們幾個的有生之年.
是大力鎮壓基督徒.
不知什麼原因.
殺了三千個居民.
三千個教徒和神職人員.
所以有一個叫血腥瑪麗的稱號.
所以你可以想像到.
在英格蘭或者蘇格蘭.
宣教士或者是Big Brother的種子.
都滲透在整個英倫三島的大地.
有很多人受這些迫害.
這樣來逃亡.
最終好像要回到法國天主教的統治.
因為剛才一開始說.
法國從根本沒有改革.
沒有做到宗教改革.
只是站在天主教的既得利益和立場的立論裡.
所以整個國家繼續是君主的專制.
教廷的封閉,腐敗.
對農民窮人的壓制沒有改變過.
但整個歐洲已經開了民意.
整個歐洲 例如德國開始.
日內瓦,瑞士,英格蘭,蘇格蘭.
已經開始了這些.
到底是否要跟隨基督教.
還是個人的自由自主.
已經產生了一個很大的漣漪.
但法國仍然沒有任何的改變.
居然出現了這三個女王.
對基督徒大力地在他們有生之年.
統治的年代進行鎮壓.
所以你可以想像到.
法國的下場不歸向真神.
又繼續在這樣的故步自封的封閉裡.
最終人的心靈是一直在互動的.
不會說現在甚麼都封閉了.
是不會的.
一定會有這些宗教改革的浪潮.
一定會影響不同的歐洲地方.
你說其他亞洲根本沒有那麼多的資源去.

$^{681}$不足為奇.
但歐洲是不會的.
當時是一體化.
是不會分得那麼緊要.
所以我自己其實很慨嘆.
我今天看的資料越看越恐怖.
再過多二百年.
就是法國真正的流血大革命.
1789年的大革命.
那些皇帝最終那麼腐敗.
那麼昏庸.
和令到國家的債台高築.
因為又打仗.
又支持美國的獨立戰爭.
所以就真是龍門正正式式.
發動一個大規模的流血的革命.
所以整個國家以為不需要歸向基督教.
不需要回到宗教改革的步伐裡.
其實是不會有好的結果.
因為你們自己天主教的神學氛圍.
神學的主張.
不能夠令到自己的國民.
有任何和神的關係的更新.
復興.
人家基督教的宗教改革.
是一路將人有系統地帶回去.
看如何和神建立一個真正的關係.
你沒有這樣做的時候.
是不會免疫的.
會讓這個問題更加嚴重的.
透過人性的敗壞爆發出來.
所以路易十六第一個示範單位.
被人斬頭.
在斷頭台斬頭.
巴士的監獄.
1789年的7月14日.
就是法國的國慶.
就是用那天血腥的那天.
作為他們的國慶.
你可以想像到法國人是多麼崇尚.

$^{721}$沒有什麼可以管束我們.
我們要的是民主自由.
當時他們的宣言就是這個.
所以這三位女士到底是什麼原因.
為什麼這麼支持天主教.
為什麼完全否定基督教的更新.
完全不接納.
還要加上用他們的權力去鎮壓基督徒.
這些神職人員.
這是一個很大的課題.
可能要問問蝶式他們知不知道.
真的很多材料很多資料.
但是幾位女士在他們有生之年.
所謂有生之年其實是說.
那個統治是10年至15年左右.
完全是聚焦地攻擊基督教.
捍衛的是天主教.
是非常之敬虔.
而更加諷刺的是.
他們這幾位的父輩.
有些是很同情基督教的.
容許在英格蘭和蘇格蘭.
在某程度是基督教的更新.
都可以來到發生的.
但是經過幾位的鎮壓之後.
肯定的就是在諾克斯要逃亡之後.
整個蘇格蘭和英格蘭.
在血腥馬力的統治底下.
其實就一直完全傾向天主教的不改變.
但是其實好景不常.
人民對血腥馬力的鎮壓是極之反感的.
最終下場都是很慘的.
所以是很短命的.
幾位都是四十多歲就離世.
當時的科技和醫療.
因為當時的人都是貴族.
如果是皇室其實是很腐敗的.
很荒淫的.
他們的婚姻都有很多問題.
所以不知道是不是有很多背後的原因.

$^{761}$但是起碼那十幾年對基督教.
在英格蘭和蘇格蘭.
在他們在任期間.
是帶來很大的迫害的.
但是我們一講到那個鏡頭.
就是1554年的時間.
他的繼任.
就是洛克斯要逃亡的時候.
1553年的10月1日.
他成為了當時的亨利百世的女兒的繼任.
做英格蘭的女王.
這些歷史對英國人來說應該相當熟悉的了.
但是和宗教完全拉上關係.
他是極之對天主教是擁戴的.
所以不可能容納基督徒存在在英格蘭.
在蘇格蘭.
所以洛克斯從蘇格蘭走到英格蘭.
因為另外一個他的表姪女瑪麗一世.
和他的媽媽.
都是對蘇格蘭的宗教改革有迫害.
所以輾轉洛克斯走到英格蘭之後.
誰知遇上女王殘暴的血腥瑪麗.
有三千個基督徒和神職人員受到迫害.
所以在這個時候1554年.
奠定了洛克斯真的要去到當時.
差不多是對基督教最友善的.
就是加爾文現在能夠在瑞士的日內瓦.
作為整個城邦的領袖.
神學的領袖.
當時的政府的共同執政領袖.
所以洛克斯去到之後.
一來就是原著的作者.
基督教要義.
他一直都冒名很想去認識這位加爾文.
所以一見如故就成為好朋友.
再在那裡受教於加爾文的教導.
所以又是兼修並蓄.
他很快就要將加爾文的神學整理好之後.
就希望引進回去蘇格蘭.
因為他的心一直都想回去蘇格蘭.

$^{801}$帶來宗教的改革.
所以這個就變成了由瑞士的日內瓦.
由德國,瑞士到現在.
就去到英倫三島.
大家知道是離開了歐洲大陸.
帶來蘇格蘭.
首先是蘇格蘭.
一個大復興.
一個改變.
當然那幾位女王相繼五幾年都死了.
所以就有機會在蘇格蘭進行宗教改革的更新.
其中很重要的就是.
整個蘇格蘭必須要有長老制度.
其實不是刻意想要有長老制度.
長老制度的好處就是.
牧者和當時的長老.
被認為是靈命好的長老.
共同治理地方和教會.
必須要神的僕人當然領袖.
當時就沒有分教會和城邦.
因為加爾文連當時日內瓦的日常公務.
他都有份去參與.
這是當時的問題.
現在就沒有可能.
沒理由突然要去行政會議開會.
或者去立法會開會.
沒有這件事.
這是一路的演變.
當時中世紀由加爾文首先制定的長老制度.
就是這個精髓.
屬靈的事必須要牧者做帶領.
當然要用聖經.
要接受人民的監察.
和根據聖經的規範等等.
重要的是成立Church of Scotland.
其中重點就是天主教的主教制.
絕對的管治.
不容許任何挑戰.
這一點是他們真正的拆毀.
讓人民更多的透明去參與一些改革.

$^{841}$所以反對主教制.
不是刻意想引進長老制.
而是想帶進不是教會絕對的沒有錯.
沒有人可以挑戰教會的權威.
這是一個問題.
引進之後就帶來一個天翻地覆.
人民開始發現原來是有傾有講.
有依據標準就是聖經的真理.
就變成不是僵化的主教制.
主教的神父就是社區的代理人.
是由教廷派來的.
整個社區的牧區.
必須要每人都奉獻十分之一的田產收入.
由教區領取到梵蒂岡等等.
不需要這樣.
所以這可以說是對當時.
今天我們有這麼多不同的體制和宗派.
我們反而覺得這麼誇張.
但其實當時是一個很大的引進.
就是讓人民一起回到神的管治.
就算是神設立的僕人 長老.
都可以被監察 檢視.
這是一個很透明的 從來沒有試過的.
所以對中世紀的教會來說.
是一個非常大的突破.
對天主教當時一定是一個很大的反彈.
是不可能的.
這差不多是不能夠挑戰的權威.
現在你挑戰了.
天主教還有什麼勢力可講呢.
每個人都可以監察.
這樣就難了.
所以這就是天主教一直不想.
也不能夠變更的原因.
就是根本的權威受到挑戰.
其中就是一個叫使徒統帥.
Apostolic Accession.
我回去找英文給大家.
使徒統帥就是說.
由教皇 第一任的教皇就是使徒彼得.

$^{881}$一直傳遞下來.
每一個時代的教皇 教宗.
那些主教都是無誤的.
所以不能夠容許有平民百姓.
可以挑戰教皇的權威.
他頒布的教令等等.
這就是正正就是95條和嘉爾文的神學.
要想去更新的地方.
所以成立了蘇格蘭的國家教會.
尤其是要長老制的教會.
就是革新的一個很重要的元素.
後來我找到他們的25條.
就是諾克斯很快在1561年.
制定了25條蘇格蘭信仰宣言.
在長老中的教會 今天也找到.
在長老制的教會也可以找到25條的條文.
為什麼這麼重要呢.
普及教育其實他們.
所以他們已經去到第三代了.
由馬丁路德開始發現聖經的真理.
開始發現二人恩信得生.
突然間的聖靈的感動和對馬丁路德的感召.
使他發現要重新研讀羅馬書.
再看其他的書信.
他有信心 內心有一些渣拿.
透過聖經來對話.
所以到了第二代的加爾文.
他們除了《終極改革》之外.
翻譯聖經和寫聖經的著色書.
都是他們很大的事奉人生.
希望在有生之年留下他們對聖經的理解.
所以他們很自然地就讀希伯來文.
本來拉丁文是他們的語言.
或者法文 加爾文是法文.
蘇格蘭的就是英文.
他們有自己衍生的語言.
但始終他們要對原文有認識.
所以要做一些翻譯和學習這些語言.
所以幾位領袖都是精通希伯來文和希臘文.
所以是好的 應該要學的 應該要讀的.

$^{921}$香港的傳道人認識原文的不多.
半筒水的更多.
但精通的真的不多.
一定要在神學院裡面的老師.
他們會在這方面.
所以是分開了的.
但中世紀就沒有這麼規範.
他們又是革命家 所謂的宗教改革.
改革家 又是神學家.
要制定一些宣言.
更加他們是聖經翻譯.
所以可以想像這些人很奇特的.
這麼短的人生 五十多歲.
做了這麼多事.
今天是不是我們不要這麼多聚會.
不要這麼多活動.
又搞踩單車 又搞這些.
你就做不到你自己應該做的事.
人家真的做了很多這些底層的工作.
所以使宗教改革.
真的定性了在神學上的更新.
信徒的心靈的更新.
所以這個普及教育.
意味著當時大量的人.
我沒有數字.
可能超過一半 甚至超過一半以上.
都是目不識丁的貧窮人.
你講什麼我都信 我都不知道你講什麼.
聖經講什麼我都沒有興趣.
或者我不知道 不知道怎麼讀.
所以你可以想像到.
既然他們有一個這樣的信念.
由馬丁路德開始.
原來根本的問題就是.
天主教不讓人讀聖經.
不怕人讀聖經會成為異端.
實際上他們自己是異端.
他們不准人讀聖經.
變成了一種愚民的政策.
使你既不好又普及.

$^{961}$但原來馬丁路德不知不覺就發現.
原來他自己在聖經裡面才找到真理.
所以經過加爾文的教教要義.
將聖經有系統地系統神學.
結合聖經神學.
來做了很有條理的兩本巨著之後.
其實只是一本.
中文翻譯了兩本.
成為當時歐洲一個很重要的寶藏.
所以隨之而來的就是聖經的翻譯.
著識書的大量撰寫.
沒有這個能力的人就自己去學習.
既然神給了你有這個能力.
你有這個學問.
你就好好去學習那些語言.
幫我的子民去做翻譯聖經.
所以有法文 德文的聖經.
英文聖經在161幾年.
就有King James的英皇音定版出現.
差不多是第一本英文的翻譯.
所以你可以想像到.
他們不只是為改革.
不只是為了挑戰天主教的權威.
而是真的要講什麼是基督教.
基督教信什麼.
這些最根本的問題.
他們清楚了.
需要讓聖經自己講說話.
讓聖經的真理自己和百姓對話.
一定是透過中間人.
就是這些領袖.
他們就要將聖經去做翻譯.
讓普通人可以讀聖經.
讀聖經之前一定要識字.
所以在普及教育之後.
下一代就開始有很多宣言.
比如剛才說的蘇格蘭的信仰宣言.
接著蘇格蘭信仰宣言之後.
就是英格蘭的威斯敏斯特教條.
Westminster的那些短的問答.

$^{1001}$是一問一答的.
即是你上了這一學就問你怎樣怎樣.
你就要答.
可能當時的小朋友就很開心.
可以有這些對答.
所以不斷有普及教育之後.
就是為了人人可以讀聖經.
就開啟了真正的文字.
所以你可以想像到.
由十六世紀一直到英格蘭.
在歐洲做這個聖經的教導.
其實就深化了他們的靈命.
所以去到啟蒙運動的後期.
歐洲的教會就沒有太受影響.
但剛才我們說的法國.
一直都沒有受基督教.
或者拒絕基督教的更新.
一直都只是擁護天主教.
人民的心靈其實就和聖經和神相當遙遠.
所以後果就不堪設想.
再過多二百年.
英國有一個大復興.
有一個清教徒的運動.
大家聽過的五月花號.
因為新大陸已經發言.
美國已經獨立.
大量的基督徒因為有些逼迫.
就走去美洲都是長老宗的教會主力.
所以五月花號的美國之旅.
又叫做長老宗的遷徙.
你可以想像到.
加爾文和洛克斯帶來的宗教普及.
就是開始人人自己感讀聖經.
感同神祈禱對話.
這個其實差不多像我下一堂要講的.
或者我透過講了第九堂.
一直準備.
真的有些東西浮現了出來.
是我們今天的香港教會.
或者普世教會.

$^{1041}$要從十六世紀的.
在這麼艱難的氛圍.
還要是自己教內的一個.
似乎是風波.
愛人如果不信神的人看.
你們教內打來打去.
其實不是.
其實是積壓了十五世紀的.
越來越錯謀的.
那種對神的聖經的誤解.
一下子的.
透過神興起的馬丁路德.
做一個更新.
而這個更新.
不是要玉石俱焚.
而是要將整個的時代.
帶回聖經的研讀.
所以唯獨聖經是當時很大的課題.
所以你看到那個氛圍還未具備.
你怎麼說唯獨聖經是口號呢.
所以他們做的.
就是在自己的牧區蘇格蘭.
做一個普及的教育.
所以我覺得這些是很偉大的.
看到那個目標.
而慢慢從第一步走到目標.
我想一個神學家.
或者一個教會領袖.
應該要有這個前瞻性.
就是要帶領你的會眾.
去到自己可以讀聖經的田地.
那個地步.
就是推動那個教育.
讓大家能夠先掌握文字.
識字.
開始對聖經有認識.
所以教會開多些的查經班.
開多些的.
如果大家基本上.
香港已經沒有人不識字了.

$^{1081}$但是我們要怎樣令弟兄姊妹.
知道聖經的重要性.
和與神有對話.
這是最重要的.
所以在這樣的情況下.
蘇格蘭也是同時間.
加爾文神學的精髓.
讀基督教要義.
有其中一個章節.
就是講信徒的倫理.
我就總結了.
其實這個就是.
洛克斯在蘇格蘭要推動的.
新教的倫理相當簡單.
一定要做一個簡樸的人生.
剛才說我們人大概有82至83年的壽命.
平均.
不是每個人都一樣.
除一除.
原來大概有3萬天.
快點大家現在算一算.
還有多少天.
我們怎樣用我們人生的每一天呢.
不多不少.
有些人可能90歲就會多一些.
應該不會去到4萬這麼多.
那就要算了.
所以你可以說這三個重點.
一定不會錯.
勤奮地工作.
但又勤奮地工作之外.
自然你會有成果.
你是會節儉的.
還有你個人的崇高的.
對神在你生命裡的照明.
是要有意識的.
不是說我是我自己.
我怎樣就怎樣.
有些人勤奮.
但勤奮完之後.

$^{1121}$他就大吃大喝.
瘋狂地玩.
星期五就晚機去.
星期一才回來.
每個星期如是.
他很勤奮地工作.
不過他只是享受他認為的人生.
新教倫理.
我們有信仰的基督徒.
就是節儉.
勤奮.
在神安排的照明.
恩賜裡面去服侍.
所以這三萬天.
不是我們自己怎樣就怎樣.
你可以說是人性一個很高的情操.
可以將我們最好的做好.
做好一定有很多回報.
認定這些回報是神給你的恩賜.
是神給你的資源.
你去投放在神要你去做的地方.
這樣就很不同了.
所以為什麼新教倫理.
很自然就衍生了.
中世紀後期的資本家.
就是他們怎樣可以知道神的揀選.
一定是在他身上.
就是他根據領袖的教導.
我要一生節儉.
不會亂輝煌.
然後我勤奮地工作.
有目的地為神工作.
有個人責任.
自然就積累了很多的資本.
成為了歐洲當時中世紀.
打後就有這些資本家的出現.
因為他們有新教倫理.
我怎樣知道自己沒有被神揀選.
如果我很殷勤地去工作.
為神工作.

$^{1161}$我就清楚認定自己是被神揀選.
因為這是加爾文的揀選論中.
很重要的精髓.
所以其實影響著當時代.
蘇格蘭,英格蘭和歐洲的基督教氛圍裡.
尤其是長老宗教會的影響力就更加大.
當然今天有些不算太強.
但在十六世紀.
長老宗教會直接影響著歐洲的思潮.
所以清教徒都是由長老宗教會一直衍生.
不過不再說宗教改革.
因為清教徒運動已經有很大的課題.
又涉及到維斯利約翰.
到了十七世紀.
然後就到了美國的基督教興起.
時間不夠.
因為要用很多功夫.
看很多材料.
不是我的強項.
所以下兩次完結.
下星期就要回到聖經.
否則我要用很多時間.
自己的工作做得不夠多就不好.
你可以想像到.
我們要問神的話語在我們心靈裡.
有沒有更新的力量.
我們有沒有受神的話語的感動.
我們不要說挑戰.
不要說對我們的刺痛.
感動有沒有讓主的道在我們內心感動我們.
我們聽到有沒有刺痛.
令我們有感動.
而令我們人生有方向.
有責任和願意配合教會的方向.
一起在人生包括我們自己的家庭都有方向.
這些是新教倫理.
所以某程度上.
如果教會裡有很多弟兄姊妹.
知道甚麼是新教倫理.
其實他們就會重新檢視人生的時間和次序.

$^{1201}$記住.
你想和不想都是三萬天.
有些人還要少一點.
例如大S48歲.
沒有想過自己會突然間.
呼吸不力就要離開這個世界.
所以沒有人可以知道預測.
但新教倫理一定是我們人生的一個目標.
我們可以令自己有這個照命.
所以為何嘉爾文他們的影響力這麼大.
是賦予平民百姓一個崇高的照命感召.
令你自己知道我要為誰而活.
當然很重要的是.
我引述給大家的蘇格蘭信仰宣言.
這25條成為了整個蘇格蘭復興的一個很重要的種子.
使他們決心.
記住.
天主教的影響力在歐洲有些地方還是很大的.
在英格蘭.
Buddy Mary和後來的繼任人.
一樣是對天主教擁戴.
不是真正徹底的平棄.
所以你要做這些信仰宣言.
其實只是在你認同的教會的人士底下去做.
所以沒有真正帶來整個歐洲.
每個人都有改變.
所以願者就去選擇.
神的恩典藉著幾個人.
帶來歐洲一個新的氣象.
但不是整個歐洲的人都接納.
至少法國人後來興起很多哲學家.
帶來的影響性是非常之嚴重.
隨便舉一兩個.
法國的北卡爾.
他們質疑人生.
某程度上是不希望神的管治在個人身上.
他們可能有一些數據.
天主教這麼腐敗.
這麼封閉.
沒有帶給我們一個更新.

$^{1241}$但他們又不接受基督教的改革.
所以在兩下都不要之下.
信仰增兇.
憑自己的良知或思想.
其實很危險.
人本主義.
我的存在是一切.
不要再相信迷信的.
天主教說宗教信仰對我們的規範.
煉獄,永生.
所以不信的浪潮會直捲.
人只會追求自由,民主,平等.
法國大革命的宣言就是這樣.
就是這樣的宣言.
只會用任何手段.
為了達到人民自由民主.
這個手段都可以合法.
所以基督教從馬丁路德開始.
一直到現在的諾克斯.
都是用他們的生命.
不是對抗政權.
而是用他們的生命去印證.
回到聖經才是人的最終重要出路.
所以你可以看到.
在未到100年.
由1517年馬丁路德的時候.
到16世紀.
159幾年.
其實都未到100年.
已經帶來一個很大的更新.
所以原來更新是很重要的.
天主教沒有更新.
只會被整個時代的巨輪.
平除在後面.
他們自己只會繼續內部的僵化.
但是一般的平民百姓.
如果不歸附於中國改革的.
回到聖經 回到神的真理.
又會落後了.
變成他們只會以自己的思潮.

$^{1281}$去凌駕聖經.
表面上以為很好.
但實際上其實更加遠離了聖經.
既討厭天主教那些僵化的制度.
告解 七聖禮等等.
亦不敬拜 亦不怕教廷.
但是不回到神的話語的根本.
又出現了一個絕對的心靈真空.
所以我想今天牧者最擔心的就是.
信仰的真空.
很多基督徒不知道自己信什麼.
是一個很大的要面對的問題.
好像一張單從來都沒有找數.
是一個很嚴重的問題.
所以如果我們不找不同的渠道去聽聖經.
可以有很多渠道.
YouTube有很多出名的牧者.
你都一概不聽.
然後你就說自己的牧者不好聽.
你去到神面前 神都不會收貨的.
你不聽曾牧師 你可以聽其他的.
你去聽陳牧師 你可以聽另外的牧師.
不只是單一的渠道.
重點是將你的心靈 將我的心靈帶回聖經.
這個才是重要.
回到聖經的根本就是.
與神對話 有不斷更新的心靈.
這個才是真正的信仰.
你一定要有你自己的新教的倫理.
你既然成為新教徒 成為基督徒.
你要有你自己在神面前的倫理.
我們每個人都一定有個人生目標.
有些人就是這樣吃 有些人就是這樣玩.
這都是目標.
但是我們基督徒的倫理是怎樣呢.
所以今天透過洛克斯.
都將他前半段青年的人生受盡迫得.
但都沒有放棄自己是蘇格蘭人.
有那個照明要將很重要的真理帶回蘇格蘭.
所以他要制定蘇格蘭信仰宣言.

$^{1321}$不過有看過一些報告.
第一篇寫的文章是極之尖銳.
暴露天主教的錯謀.
和當時地方的君王的敗壞.
看看有沒有機會拿到原著和大家分享.
相當尖銳.
你可以說是招致天主教和擁戴天主教的當時君主.
對他的無差別的迫害.
但無論如何他的心都是懷念著蘇格蘭人.
所以在日內瓦那麼好的那幾年.
他的書信告訴他的朋友.
在日內瓦和嘉爾文的對話和學習.
是相當得益.
那幾年是他人生中最快樂的人生.
但最終他還是選擇回到蘇格蘭.
我想真的.
牧者真的要有一個很清晰的照明.
在哪個牧區神安排你.
所以其實要保持這份常常保持的牧養.
事奉的初心.
我2009年畢業.
現在已經是2025年.
已經15年多了.
真的離開神學院.
當年教的東西完全沒有用.
但那個初心.
在神學院裡面訓練的.
和一班同學一起.
即是畜生夜談.
上課.
祈禱會.
那份初心要保持.
所以剛剛過年就約一班同學.
一起吃自助餐.
假期.
再講一下大家的近況.
其中很重要的就是.
大家如果現在還是有心事奉.
就是那份初心不改.
那邊有些同學就有些變化.

$^{1361}$不再牧會.
一定有他們的一些困難.
所以牧者的初心是很重要的.
所以大家必須為牧者.
尤其是事奉了很久的牧者祈禱.
要保持剛剛畢業的初心.
真的好像甚麼都不懂.
但其實懂的.
但那個心要保持單純.
保持那份好像童真一樣.
不要變成老油條.
其實很重要.
其實需要頂尖的祈禱.
和需要頂尖的支持.
互相彼此激勵.
所以提醒絕對不是不可以.
是可以.
提醒不是要責備挑剔.
而是同心的同行.
兩者是不同的.
所以盼望大家千萬不要介意.
如果你知道牧者有甚麼需要可以提點.
彼此互相激勵.
其實這個才是真正最好的同行和代禱.
今天我預備的材料就來到這裡.
下星期就真的一個總結.
大家可以想一下這個課題.
在過去的九堂.
大家又可以有甚麼領受.
我都很難叫大家在咪高峰上說話.
我可能組織幾個點.
大家就著這幾個點可以再思想.
日後你自己再反思.
我們今晚就到這裡一起禱告.
感恩天父給我們透過蘇格蘭的改革家.
諾克斯這位神的僕人.
他早年因為認同基督教的更新.
受盡迫害.
有19個月要成為一個被法國奴役的階下囚.
過盡一些很慘無人道的生活.

$^{1401}$但是感恩天父讓他去到日內瓦.
認識加爾文.
能夠真的可以在整個加爾文神學裡面.
他獲得很多的靈感和指導.
使他能夠在蘇格蘭和他的同伴.
制定蘇格蘭的信仰宣言.
很盼望透過我們在過去的九堂.
馬丁路德,伊拉斯姆和加爾文.
和今晚我們說的諾克斯.
都給我們有一定的啟發.
我們真的需要回到神的話語.
能夠和神有很好的溝通.
我們能夠上上沐浴在主的恩典裡面.
不只是一種個人享受.
而是讓我們清楚自己的人生的召命.
可以在有限的人生去為主而活.
求主幫助我們仰望,感恩,禱告.
奉耶穌基督的名.
阿們.
\newpage



\section{}
\label{sec:q7QZXuHzcds}
\textbf{2025-02-13 周四晚查經班《改革之路(10)完結篇:未來》 講員:曾裔貴牧師}
\newline
\newline
連結: \href{https://youtube.com/watch?v=q7QZXuHzcds}{\texttt{https://youtube.com/watch?v=q7QZXuHzcds}} ~~~~ 語音日期: 2025-02-14
\newline
\newline
\hyperref[sec:5Lw6EDNBsNI]{\small{< < < PREV SERMON < < <}}
~
\hyperref[sec:index_chronic]{\small{[返順時目]}}
~
\hyperref[sec:index_scriptual]{\small{[返順卷目]}}
~
\hyperref[sec:mCrcbeplxo4]{\small{> > > NEXT SERMON > > >}}
\newline
\newline
$^{1}$邀請大家一起同心禱告.
感恩天父讓我們再次到第十堂.
來看宗教改革.
他們的人生,主張,神學論述.
帶出一個很嶄新,脫離了一個很僵化的宗教體系.
其實一點都不簡單.
因為整個歐洲的復原.
天主教已經超過一千年.
在整個歐洲的大陸的潛移文化和傳統的運作.
其實已經在整個歐洲盤踞了.
差不多是神的代言.
但是我們又看到居然有人看到一些問題.
來針對和扭轉.
又帶來一個新的格局.
基督教從此開始在歐洲大陸傳道.
傳講一個回到聖經的真理的福音.
一直我們看到到十六,十七世紀.
又有另外一波的大復興在英國,美國的出現.
以至復興的浪潮直接催生普世宣教的運動.
使中國,印度,非洲,東南亞.
這些從來都沒有福音.
比較整全和有規劃的傳福音的運動正式開始.
我們看到神都是這樣按著一個時代一個時代.
疊加上去這些時宮的福音能力.
我們看到當人回到聖經的時候.
真的會有一個很大的改變.
所以我們都盼望今晚算是最後一堂.
關於我們一個很淺談的宗教改革.
幫助我們能夠有好幾個大方向.
我們都能夠再一次去反思一下.
我們有沒有失落.
還是我們需要再繼續在那個改革裡面.
在一種溫和的.
在一個群體裡面不斷地去更新.
求神幫助.
我們這樣仰望祈禱.
奉耶穌基督的名.
阿們.
多謝大家用了九堂的時間.
一起來看這個宗教改革.

$^{41}$先看看這裡.
它那個分享畫面.
分享了.
我看看怎樣開大一點的畫面.
為甚麼它不懂得動.
還是要…….
好了 行了.
我們有說過……又不走的?.
分享……是…….
先看看 為甚麼不動呢?.
應該要去的 是的.
可以開這個 行了.
對不起 有些技術問題.
我們有機會看過比較完整的.
馬丁路德 加爾文.
和上一次我們說的約翰·洛克斯.
因為他直接受加爾文影響.
和又帶來蘇格蘭和英國.
某程度上還沒有去到大復興.
就是將聖經和蘇格蘭的信仰宣言.
去落實了.
使得往後的牧者.
即是他們下一代的牧者.
就能夠有一個很重要的真理的基礎.
去推動教會的革新和復興.
我們也說過一個叫做伊拉斯姆.
也是一個很重要的人民主義的教育家.
差不多是教育之父.
他是第一個將拉丁文的聖經.
翻譯成新約的希臘文.
所以又是帶來翻譯的過程.
其實在他之前已經有的.
比如比較久的威克利夫.
威克利夫現在是一個機構.
就是聖經翻譯.
所以你可以想像到.
由13,14世紀.
或者14世紀就比較蓬勃一點.
到15世紀.
當然16世紀就是宗教改革.

$^{81}$即是中世紀的時候.
所以神興起這些算是有色之士.
或者在教會裡面.
有一定的參與的人.
他們馬丁路德是神父.
即是修士.
其他那些只不過是一些很擅長讀書.
接近大學的校長.
例如加爾文就是法國大學的校長.
洛克斯和他的好朋友.
也是蘇格蘭大學的校長.
馬丁路德本身是讀法律.
後來做傳道.
成為修士.
有另外一些其他的人.
例如約翰·胡斯.
比較早期的殉道者.
差不多影響著馬丁路德的一個很重要的人物.
相隔100年了.
約翰·胡斯.
他是在捷克.
當時叫波希米亞.
捷克大學的校長.
所以你可以看到.
都似乎有些脈絡.
神興起這些有學問的人.
至少他們能夠讀聖經.
能夠讀明聖經.
還有機會去看.
可能不只是拉丁文.
其實大部分人都懂拉丁文.
因為那時是官方語言.
他們會來看拉丁文.
當然他們更加輾轉地將自己的國家.
例如胡斯是捷克人.
他會將聖經翻譯成英文.
維克利夫翻譯成英文.
加爾文更加將法文翻譯出來.
馬丁路德翻譯成德文.
所以你可以想像到.

$^{121}$我們要留意.
神在一些時代裡.
當然絕對不是說.
神只會興起一些有色之士.
這個說法未必是成立的.
因為任何人都可以被神使用.
那個人是否能夠被神使用.
為什麼在這個時代.
會有這麼顯著的人物.
可能當時真的有這個需要.
第二個重點.
有一點端倪.
因為當時在十六世紀.
其實應該是橫跨十五,十六世紀.
因為約翰·胡斯.
其實很早就已經反對贖罪券.
可能是由外圍去反對教廷.
效力就差不多是以下犯上.
外人要論斷.
要想推倒整個教廷.
是很難去想像的.
所以只是成為了炮灰.
約翰·胡斯在1415年被判斷為異端.
要被火燒死.
當眾在很多人面前.
就像槍斃一樣.
要群眾去看.
所以胡斯的死的殉道.
又引起很多人.
胡斯的學說.
引起很多有學問的基督徒.
有些反思.
是不是這個人真是異端呢?.
還是天主教已經落到一個這樣的境況.
已經是到了一種權力.
到了一種絕對的腐化呢?.
隨便一個人.
很善意的提出他的論點.
說贖罪券是不應該的.
但為什麼又會遭遇被公開的殺害呢?.

$^{161}$還是用宗教的名義來殺害呢?.
所以引起很多人的反思.
但當然未能結合一股的.
具體的神學的勢力.
去論述如何能夠有聖經的支持.
天主教是否應該大家坐下來對話.
那時候還沒有這樣的氣候.
所以一直醞釀.
我們看到.
終於又回到贖罪券.
馬丁路德的《論綱》第95條.
其實主要集中贖罪券的利弊.
其實沒有利.
只是有弊處.
贖罪券的大概意思.
下面也說了.
信徒那個時候.
擔心死後去哪裡.
14世紀.
整個歐洲的黑死病.
瘟疫的鼠疫.
有很多年青力壯的人無端端死.
因為那時候不知道什麼疫苗.
也不懂得.
只希望等到那些疫情自己完了.
有抵抗力就算了.
所以突然無緣無故死了.
所以當時的畫都是魔鬼.
好像歐文使者那樣去抓那些人.
等於我們中國人那樣的地獄使者.
你可以想像得到.
當時的歐洲人其實很怕死.
更加怕死了之後去哪裡.
因為大家要記得.
天主教說的是.
你在今生有罪.
即是現在你再生有罪.
你就要向神父.
一定要向神父告解.
告解即是承認你做過什麼.

$^{201}$有碗就要說碗.
有碟就要說碟.
這樣來告解.
告解完之後.
你就要有些貢獻.
不是叫貢獻.
叫捐獻.
就是contribute(貢獻).
當時天主教都是建教堂.
修葺和維護.
不同的牧區.
又會爭誰建得高些.
漂亮些等等.
所以不知不覺間.
在互相競爭.
在在需要的就是.
教友的十一的捐獻.
所以最好的辦法就是.
賣贖罪券.
贖罪券有兩個好處.
第一就是既然告解完.
你今生有些罪.
現在神父代表神赦免了你.
你要有些貢獻.
所以救贖罪券.
更重要的是.
第三個部分就是.
煉獄.
這是拉丁文.
我特意找了.
英文是atory.
是purgatory嗎?.
是專有名詞.
煉獄.
又不是天堂.
又不是火湖.
但是信徒.
要去到宗教的地方.
做一些保贖.
到底有多少時間呢?.

$^{241}$說不完的.
你的罪多些.
就要在煉獄內一些.
在天主教的時期.
不知什麼原因.
因理成變.
只要你賣贖罪券.
你賣得多.
他還沒死.
你賣贖罪券的人還沒死.
你的親人在煉獄的時間.
就可以減少.
他可以提早去天堂.
這些當時的論述.
一方面.
似乎可以令到.
怕死的歐洲人.
可以有一個保障.
我買了.
就有保障.
現在早些購票.
早些了優惠.
做了這些手續.
我就可以有一個保證.
我的親人就可以減少痛苦.
等於燒一些醫治給他.
就變成了.
還要是由.
主教神父.
牧區的領袖.
簽罰的.
他說有效.
我當然相信他.
所以.
不知不覺間.
他們這種狀況.
令到整個歐洲的人.
因為當時只有天主教.
他們很單一.
他們整個牧羊的傳道策略.

$^{281}$只是每個星期日.
崇拜彌撒.
沒有其他東西.
彌撒之外.
他們的教會.
就不會跳讚美操.
不會去排排舞.
不會踩單車.
會做什麼呢?.
就是要告解.
所以你找教會.
就是這些東西.
主持你的婚禮.
你的嬰孩的英鞋洗禮.
其他就沒有其他東西.
所以對整個歐洲的生活.
他們要面對的困難.
天主教好像都一概不理.
當然你去找告解的時候.
他們可能都會分享.
會給你一些牧羊.
安慰.
但是整體來說.
信仰對歐洲人來說.
當時其實是一知半解.
他們真的以為買了這張券.
就可以得到一種超度.
所以這個就會帶來.
在馬丁路德他自己親眼目擊.
就是贖罪券.
帶來的就是更加肆無忌憚.
對罪的敏感.
不強.
因為買了之後.
越買得多.
那個人就越覺得自己有一種自保.
有一種.
好像很.
有保障.
就變成他更加肆無忌憚.

$^{321}$來犯罪.
所以其實又會變成.
一些有錢的人.
可以更加在.
罪的敏感度上.
窮的人買不起.
又會被排擠歧視.
就變成那件事.
本來不是.
神父主教.
來卸人的罪.
慢慢不知不覺.
就建立了一種.
這樣的代理人的制度.
就帶來.
有很多衍生出來的問題.
所以.
你可以這樣說.
如果我們要去想.
我們要去反思整個宗教改革.
首先就是.
其實不是馬丁路德開始.
他原創的.
其實是由13世紀末期.
到14世紀.
其實已經有人看到.
為什麼其他普通人.
不可以.
自己讀聖經.
所以有這個威克利夫.
他是1320至1384年.
的一位.
很出色的一位學者.
他很主張.
一般人.
要有自己文字的聖經.
所以他翻譯了.
第一本英文的聖經.
當然後來還有其他.
聖經的譯本翻譯.

$^{361}$所以已經有這些.
之前的先賢.
來提出.
天主教的一些問題.
希望得到一個很妥善的對話.
更新.
但是招致的只是.
單方面被評為異端.
所以.
威克利夫在1415年.
為什麼會在這個時候.
已經死了很久.
已經死了30年.
被判為異端.
燒他的書.
把他的屍體遺骸.
拿出來燒.
好像鞭屍一樣.
為什麼呢?.
因為當時1415年.
約翰胡斯被燒死.
被判為異端.
所以原來.
你指出天主教教廷的問題.
你是會有這樣的下場.
但是其實.
在整個體系裡.
他們有一個很大的缺失.
整個歐洲.
不可以說是神.
藉著瘟疫.
藉著鼠疫.
不過是一個自然的.
這麼嚴重的.
大規模的會死人的.
會集體死很多人的.
瘟疫.
等於你可以想像.
如果去非洲.
那些伊波拉病毒.

$^{401}$突然間人們就出血死的.
是很恐怖的.
歐洲當時就是這樣.
但是教會不能夠回應.
百姓的內心驚恐.
信仰本身就是.
牧養人的心靈.
但是不行.
但是又不知道.
有什麼辦法.
因為信徒不懂得.
和你看聖經.
看神學的東西.
所以只會落在一種.
僵化的買贖罪券.
告解.
但是實際上.
信仰沒有真正的更新.
所以.
你可以想像到.
去到馬丁路德的出現.
帶來的.
在這個缺失裡面.
有幾樣東西.
是馬丁路德.
他在教內.
神的工作.
在自己教廷裡面.
一個有心的年輕人.
一個經歷過.
上主聖靈更新的年輕人.
他自己.
很認真.
探求自己的內心.
發現自己.
根本沒有力量.
去面對自己的軟弱.
但是當時教廷.
不斷宣傳.
他們只需要成為主教.

$^{441}$就會自然成為聖品人.
其他的信徒.
只不過是一些平民的信徒.
將聖人.
和普通的平民.
有一個階層.
所以.
這些聖品人.
這些主教.
他們主持的聖禮.
禮儀才有效.
所以.
在彌撒裡面.
或是說你必須去.
望彌撒.
就是要領受.
築聖過的那些聖體.
所以.
在這樣的代理人.
不知不覺的.
在這樣的代理人的制度下.
其實這種缺失就是.
百姓,平民.
得不到心靈的牧羊.
懼怕死亡.
又解釋不了.
給到一個真理.
他們面對死亡.
之後去哪裡.
給不到答案.
只是告訴你.
你在煉獄可以減少一些痛苦.
就是買贖罪券.
就是這麼多.
其他的時間.
一切的聖經教導都很欠奉.
崇拜.
中世紀的時代.
就是將拉丁文每次去讀.
不解釋.

$^{481}$也沒有任何的教導.
當然現在的天主教肯定不會這樣.
因為基督教給到他們的啟示.
那個成功.
他們一定會有一定的改變.
現在甚至在YouTube上.
都有一些神父.
會將一些教理文答去解釋.
去給教徒知道.
實際背後的意義.
或者和基督教有什麼不同.
那個距離.
有多遠.
都有解釋.
所以他們開始意識到.
在牧羊上.
要牧養到那個時代的信徒.
尤其是心靈的牧羊.
如果沒有.
就有很大的問題.
所以這一點.
稍後浮現了.
在牧羊上.
今天的宗教改革.
過去就是這樣.
在一件這樣的事情上.
由95條.
引發的討論.
那個對話.
當然不是完美.
而是一定要一分為二.
等於分裂.
但是基督教的發展.
是一路朝向一個.
越來越趨向真理的部分.
天主教就在他自己.
因為他一定要捍衛.
他們的教統.
不能夠.
當時的主教.

$^{521}$特別是教皇的無誤.
主教坐在他的位置.
說的話全部不會有錯誤.
等等.
在這些時候.
他們都是為著那個傳統.
或者那個教統.
他們不能夠有一個很真誠.
大家平等的對話.
所以教會的權威.
和教會的權力.
受到挑戰.
因為馬丁路德是透過聖經.
和他們展開對話.
希望他們來正視.
贖罪券帶來的問題.
和到底信徒今生有什麼把握.
和來生.
如何知道自己.
在神的救恩裡面.
但是天主教在那個時候.
就沒有這個能力.
所以他們就沒有.
這個能力.
去看神的救恩裡面.
但是天主教在那個時候.
是欠奉.
所以你可以想像到.
最大的死亡威脅.
天主教都不能夠.
給一個具體的.
模樣上的答案.
所以看到天主教的缺失.
正正就是.
這裡就開始出現.
一個很大的.
信仰的轉移.
很有point地.
你可以想像.
馬丁路德這種原創性.

$^{561}$其實.
陸陸續續.
以前的人都已經提出.
聖經要大家讀.
不只是教廷,神父.
拉丁文.
讀給人聽.
然後又不去解釋,講解.
沒有研習班,這樣不行的.
所以你可以想像到.
那個模樣上的到位.
就是提出.
很簡單的口號.
因為不是當時的歐洲人都懂得.
提出很簡單的口號.
但是在這些口號的背後.
其實是回到.
尋求聖經.
那個研讀的解釋.
所以什麼叫.
唯讀聖經呢?.
就不是拿聖經來供奉.
他們是很具體地.
既然馬丁路德是德國人.
他就要為自己的同胞.
把希臘文.
把希伯來文.
因為是沒有的.
保羅只是寫了希臘文.
舊約是希伯來文.
所以.
一直都沒有德文.
所以既然你現在有.
那個復興.
你有那個眼光.
看到那個需要.
那個需要就是.
人人都要自己讀聖經.
認識神.
最重要的背後就是.

$^{601}$打破了代理人制度.
不是某些階層.
壟斷了解釋權.
所以一定要開放.
一般人都要自己讀聖經.
或者說要明白聖經.
那個過程.
就不是一步到位.
馬丁路德就要用他的功夫.
翻譯.
德文聖經.
翻譯完之後.
就要一邊翻譯一邊教導.
把德文的聖經.
成為每個主人.
用的自己的語言.
就可以.
大家有聖經了.
有印刷的聖經.
就開始多的.
研習班就會出現.
就會興起不同的主學老師.
那時候還沒有叫主一學.
查經的.
解經的.
就會湧現.
在教學上補足了.
天主教一直以來.
沒有去解釋聖經.
他壟斷了解釋權.
但是不去解釋聖經.
這個落差就很大了.
大家都習以為常.
所以有這個唯讀聖經.
這個口號.
就很重要.
很新鮮.
帶人去.
原來我是一個很普通的人.
我耕田的.

$^{641}$也可以聽聖經.
在家裡可以有聖經.
你可以想像到.
新的氣派就出現了.
整個模樣的策略.
都已經不同了.
當然.
馬丁路德更深的看到.
是恩典托著.
他這個罪人.
所以他是一個很重要的口號.
同時是聖人.
同時是罪人.
有生之年都是.
不能改變的.
我們的罪根.
一來不是罪根.
我們的罪性.
是伴隨著我們這個人的身體.
當然是說.
身心靈.
不是肉體的軀殼.
是整個身心靈.
我們都是同時是聖人.
但同時都是罪人.
即是說.
《迦拉特書》五章.
我們如果隨從肉體.
就會犯罪.
隨從聖靈.
就會生命平安.
是很清楚的.
但天主教一直的講解.
聖品的制度已經定規了.
執行聖禮的時候.
他們是沒有錯.
甚至乎.
他們是近乎完美的.
這種教導.
但很奇妙地.

$^{681}$當事人馬丁路德.
他自己.
完全.
很謙卑地.
或者很坦誠地.
他自己.
已經是很認真地去追求神.
他都發現.
自己的內心.
還是有很多.
未完善的地方.
他很坦誠地.
跟他的信眾.
和那個時代.
透過他的文字.
來講解.
這就是.
這個嶄新,新鮮.
這個模樣的轉移.
就是有一個人.
願意我們今天說的.
落地.
願意告訴周圍的人.
我的職份是神父.
但我的生命.
和你都一樣是.
普通的人,普通的信徒.
這是在當時.
是很嶄新的.
是很石破天驚的.
沒有一個神父.
說自己.
和一般人一樣.
你可以想像到.
那個階層的對立.
由馬丁路德.
貫注的.
這樣的唯獨恩典.
憑信心去得救.
就帶來一個.

$^{721}$很徹底的.
一個人.
他只不過在職份上.
有責任去講解聖經.
但他的生命.
他的生活.
和普通的人完全一樣.
在神的面前是完全一樣.
所以也有這個.
很嶄新的.
信徒皆祭司.
每一個信徒.
包括神父.
聖職人員.
都是在神的.
揀選裡面是一個祭司.
祭司就是說.
你代表人.
帶到神的面前.
當然舊約就是帶牛羊.
是一個袋屬.
但是新約的祭司.
就是將人帶到神的面前.
當然就是.
為他祈禱.
引領他.
慢慢認識真理.
還有你要在你的人生.
要清楚自己的.
身份,照明.
不是渾渾噩噩的.
這個叫信徒皆祭司.
每一個人都有神.
所給的照明和恩賜.
來發揮.
所以.
今天我們要思想.
絕對不是信徒.
皆牧師.
是祭司不是牧師.

$^{761}$就是說職份.
是真的需要神揀選.
才可以做的.
如果今天.
我們說信徒皆牧師.
Alex你讀完神學.
你就可以上台講道.
只不過是你的口才好.
還是我的口才好.
這樣就有問題了.
信徒皆祭司.
是想強調每個人.
都是在神的面前有一個照明.
但是神給的.
不同的職份.
教會的職份.
的確不是馬丁路德所說的.
這一種.
每個人都可以事奉.
每個人都可以講道.
有時候有些混淆.
會令教會的聖職.
慢慢平淡化.
一般人.
稍為讀過神學.
晚間的可以做傳道.
可以上台講道.
馬丁路德.
絕對不是這個意思.
尤其是當時的講座.
是要將聖俗.
的對立.
拉回平等.
馬丁路德認定自己.
職份上有神的揀選.
但是生命上.
和一般的信徒完全一樣.
所以同時是一個罪人.
同時他又是一個聖人.
這才是一個最大的突破.

$^{801}$因為要將聖品的制度.
帶來一個改變.
所以.
我們第一個要反思的地方.
說了大堆這些東西.
就是.
我們要反思.
就是.
馬丁路德.
包括之前的.
特別是影響著.
馬丁路德的約翰胡斯.
他們已經看到.
天主教已經積雜了很多.
在牧羊上的缺失.
那個僵化.
不能夠帶來.
那個時代的一個.
信徒.
在這麼大的恐懼下.
是真心的恐懼.
是一個傳福音的契機.
更加要讓他看到.
永生的實在.
地上有很多東西.
雖然控制不了.
還沒有袁國勇出來.
發明什麼疫苗.
沒有這些時代出現.
但是信徒.
任何一個時代.
都因為聖經的預言.
啟示和耶穌的救恩.
是應該有生命的平安.
但是天主教.
既然是獨家代理.
給不了那個時代.
有平安.
就是有一個很大的.
牧羊上的偏差.

$^{841}$所以我們就要問.
17世紀.
神又興起美國.
英國的大復興.
去到19世紀.
他們意識到.
我們只是歐洲大陸.
美洲有福音.
但是整個的.
大洲.
印度,中國,非洲.
完全沒有福音.
那怎麼辦呢.
那些人是排著死亡的.
落火湖的.
所以他們就真的.
不同的宗派.
意識到這件事.
就真的大規模的發展.
年輕人.
要出國.
去這些.
落後的地方.
所以.
很多經典的問題.
在印度.
有一些地方.
聽到福音.
為什麼你們這麼遲才來.
我們的父母.
沒聽到福音就已經死了.
為什麼你們這麼遲才來.
這些是當時的宗教.
去到印度,中國.
那些人這樣說的.
你可以想像到.
是一個很大的.
你可以想像到.
是一個時代的神.
在用福音.

$^{881}$在帶進不同的時宮.
來到21世紀.
我們.
有什麼要做呢.
教會.
不是個別的教會.
不是個別的制度.
我們已經衍生了.
現在比較大陸的.
有主教制.
聖公會,中華基督教會.
區制.
有浸信會的一些信徒的.
回眾制.
每個月都有月會.
回眾來投票.
等等.
也有我們的長老制.
例如宣道會.
就是長老制.
就像一個.
代議政制.
一班的執事.
或者是堂委.
堂董被選立出來.
代表著回眾.
來處理.
教會一些行政.
或者是聖公的事.
當然有牧者做帶領.
主要是這幾個制度.
所以現在問題是.
我們透過.
在這十六世紀的宗教改革.
我們看了大概.
五十多六十年.
馬丁路德,加爾文.
和約翰洛克斯.
他們這幾位.
比較主要的領袖.

$^{921}$橫跨的.
由德國,法國,瑞士.
去到英國.
蘇格蘭.
等等.
當然.
北歐那邊也有.
東歐的捷克也有.
怎樣去理解.
這個模樣.
現在來到我們21世紀.
我們有什麼要.
回去注意.
要留心.
我們在這個模樣上.
能否意識到現在的時代.
不是個別的.
個人.
而是整體.
那個時代的思潮.
那個時代.
現在欠缺什麼.
在信徒的心靈裡.
或者在一些未信的人的心靈裡.
共同.
懼怕什麼.
因為福音就是要回應.
人心靈的懼怕.
那個恐懼.
怎樣去處理呢?.
這個就是一個.
很大的問號.
很大的一個答卷.
是教會.
共同.
要去回應的.
也是上主.
讓我們看到.
所以我們需要祈禱.
為一個時代禱告.

$^{961}$這個時代欠缺什麼呢?.
似乎現在.
醫藥不欠缺.
社會的公平.
或者是民主不欠缺.
當然一定有些落差.
有些強有些弱等等.
每一個制度都有它的缺失.
有它的好.
但問題是.
人的內心深處.
到底.
未有福音.
他懼怕什麼呢?.
例如工作的壓力.
社會的不穩.
各方面的事情.
我們能不能夠透過福音.
福音是沒有變過的.
去回應呢?.
這個我想是透過宗教改革.
幾位先賢.
他們所做的.
對那個時代.
有一個清晰的把脈.
而去做了改變.
當然就是就天主教的缺陷.
而做出的修訂.
例如你們不強調聖經.
就出現問題了.
所以這班的領袖.
很重要的.
不是為了要革命.
不是為了要反對天主教.
而是要給予.
那個模樣.
原來欠缺聖經的教導.
所以首先就是要翻譯聖經.
和要講解聖經.
有查經班.

$^{1001}$有一些.
的識經書.
由馬丁路德.
加爾文都大量地寫.
所以你可以想像到.
第一個重點.
我覺得我們要思想的.
就是怎樣.
共同地去回應.
那個時代的心靈.
這個就好像很空泛一樣.
但是.
必須是要面對.
一個時代.
到底他們欠缺什麼.
這個時候需要.
很冷靜,很觀察.
和很多對話.
去了解.
去看到那個心靈的需要.
這是第一個.
在這一大段.
我自己一直.
在這九堂裡.
都很浮現.
在我自己的心靈裡.
牧者.
不只是策劃一些活動.
而是要為一個時代.
做一個把脈者.
當然有不同的角色.
有些人在前線.
做一些.
很賣身的工作.
有些人可能是.
默想.
做一些寫作.
在神學院裡面.
教導神學生.
有大家不同的崗位.

$^{1041}$在那個時代裡面.
但是每一位牧者.
包括教會的弟兄姊妹.
都要問一個問題.
這個時代.
那個懼怕是什麼.
福音怎樣能夠有效.
去填補這個懼怕.
這個人心的空虛.
那個空洞.
這個就是.
每一個時代的牧者.
必須要問的問題.
這個就是我和大家.
做一個.
首先第一點的分享.
宗教改革的未來.
我們必須要有一個.
為時代把脈的.
這樣的心靈.
所以真的要多祈禱.
真的不只是為很少的事祈禱.
是要為一個時代守望.
為整個時代的需要.
心靈的需要去守望.
剛才提過.
神興起的伊拉斯姆.
是一個相當出色的教育家.
雖然他不認同宗教改革.
脫離天主教.
他很同情天主教.
不是說他們不需要改.
但也不要這樣分裂地改.
所以他對馬丁路德.
後來我們之前說過.
本來是好朋友.
後來都是少少的脈絡人.
因為他繼續在天主教裡面.
進行對話.
當然天主教不會有什麼大改變.

$^{1081}$但是他自己仍然.
由他把拉丁文的.
叫做《武加大譯本》.
即是Rogue的拉丁文版本.
翻譯成新約的希臘文.
由伊拉斯姆這位教育家.
這位人民主義者.
他的人品方面都不錯.
很推崇教育.
也就是說.
令到歐洲很快就有普及教育.
都是這些先驅人物.
雖然很多歐洲人他們不識字.
但也有這班人.
不只是寫讀聖經.
也讓普及教育.
慢慢滲透在歐洲不同的地方.
所以就有這些種子去做.
所以今天似乎香港.
或者較為民主的地方.
文盲應該少.
不過也不要一概而論.
美國雖然很民主很自由.
但是識字的人也未必很多.
不是每個人都很精英才能進大學.
你說那些很出色的斯丹福.
麻省理工.
或者其他的加州柏萊.
這些很精英才能讀.
不是真的很多美國人讀書.
只是說一個這麼自由的國家.
香港也是一樣.
可能每個人都識字.
但是關乎到聖經的全面認識.
我也不敢說100\%.
可能很多現在我們在信主的弟兄姊妹.
我們未必真的能夠.
很基本的聖經我們能夠掌握.
我們可能都是很片段式.
很金句式.

$^{1121}$甚至很多東西未必懂得.
所以這是基礎教育的問題.
要爭取.
科技越來越發達.
手機的APP可能有十多個版本的聖經.
已經很感恩它肯裝了.
不肯裝的大有人在也不奇怪.
寧願用來玩撲克牌.
玩玩偶像.
也不會看聖經.
開車當然不行.
坐車也不會看.
可以想像到.
以前普及要透過印刷.
要文字.
要識字.
現在已經不需要.
已經可以是手機這麼方便的流動裝置.
但不等同於聖經更加普及.
這個缺失.
其實真正問題的所在.
不是沒有人講聖經.
而是在信徒的內心.
好像已經對聖經有一種.
不知不覺的厭倦.
這種厭倦可能是厭惡.
可能是根本沒有興趣.
這是一個時代的反應.
不是個別某間教會.
而是普遍的教會.
所以你可以想像到.
以前是由零.
什麼都沒有.
聖經是綁在教廷的.
教堂的神父崇拜.
彌撒先讀.
在家裡找到聖經.
是要被抓的.
很極端.
但現在這麼普及.

$^{1161}$印刷版可能家裡也未必有.
但是教會現在普及的時代.
不鼓勵你拿著硬體本.
甚至很麻煩.
不要擺出來了.
因為普及已經解決了.
包括洪恩也是.
大家都不會寫聖經了.
都不感覺需要寫聖經了.
可能現在另一個層面.
我們不知不覺.
表面以為我們每個人都懂得勝經.
其實可能是由教會間接地.
不知不覺地.
令到信徒不喜歡讀經.
這樣的間接.
要怎樣來堵塞.
要怎樣平民化.
大家對聖經有真心的愛慕.
興趣.
不是為了滿足哪些人的指引.
好像現在叫你速讀聖經.
現在是否正在讀羅馬書.
好像很不好意思地說.
你速讀一下吧.
好像很怕得失你.
你再跟我說.
我下星期就不來了.
當然好像開玩笑.
但是你可以想像到.
那種如飢如渴.
其實是欠奉.
即是說一個時代已經去到.
有一種像聖經的腸胃炎.
如果是.
你吃什麼都不行.
一定要清除.
你才可以身體有一個不再排斥的復原.
現在我們就要問.
聖經的吸收.

$^{1201}$那個愛慕.
是否有了一種不知不覺的排斥呢.
而且是一個時代的排斥呢.
我們就要認真去祈禱.
甚至是否是一個時代我們需要認罪.
是什麼原因令到信徒不喜歡讀聖經.
信了主也不喜歡.
不是說未信耶穌.
信了也不喜歡讀.
那就變成我們真的要問.
因為務求要聖經在很多信徒的內心.
他真的明白自己所信.
所以不是只是一種頭腦上的認知.
一個認同.
而是他要在聖經的教導裡.
他自己的生活要有受聖經的影響.
這個第二方面其實和馬丁路德的主張完全一樣.
就是透過伊拉斯姆.
讓我們看到他們這一班人這麼緊張.
信徒要有自己的文字的聖經.
他們就這樣做這個普及的教育.
相信到了這個情況下.
有很多普通的信徒.
開始有機會拿起自己第一本的聖經.
德文的聖經.
法文的聖經.
這樣來去中讀的時候.
相信那個震撼性是很大的.
因為從頭到尾.
你們很多的祖宗世代.
都沒有機會自己家裡有一本聖經.
現在突然有一本聖經.
在自己的家裡.
其實是一個很大的震撼的事.
當然第二代的加爾文.
就不只是一種抗爭性的思潮.
因為馬丁路德沒有辦法.
因為他就是這樣提出.
就是這樣受到迫霸.
受到排擠.

$^{1241}$他無論如何.
他自己都要被迫分裂.
離開教廷.
成為帶領者去建立基督教.
但是加爾文第二代就不是了.
他們這班人很沉著的.
用真理.
用他們的言讀.
用他們對古教父的信任.
很多的教義.
已經不需要再改變了.
這樣來將他有系統地寫成了這本基督教要義.
就帶來宗教改革更加系統化.
更加合情合理.
合真理.
這樣來繼續推崇.
去開展傳道.
所以不再是傳一個問題.
譬如說馬丁路德的跟隨者.
他們去傳.
他們未必知道這麼多東西的時候.
他們只是仍然在情緒化之下.
我的老師.
我們的這位牧長.
受到迫霸.
他說的我們相信.
所以我們現在要反對.
我們要離開天主教.
這些還是慢慢不知不覺.
你越傳就越亂.
越來越變成情緒化.
越來越為反而反地去傳.
本來分裂就更加形成真正的分裂.
但是有了基督教要義.
是很有條理地說.
到底基督教為什麼我們不是天主教呢?.
為什麼我們被迫要離開.
但我們仍然希望有對話.
因為聖經的真理.
是有系統地被很條理化地整理好.

$^{1281}$普通人都可以讀得明白的.
稍後我都引述了兩段給大家看看.
所以其實就不是再為反而反了.
第二代是最不容易的.
因為你必須要有一個真理去確立.
為什麼你和天主教不同.
所以別人的錯就不等於你對.
別人的錯是你要找到之後.
你要有你自己的立論.
那個立論在聖經裡面.
這樣才算合理.
就不可以說你錯了.
所以我都要抵擋你.
這樣就沒有意思了.
這樣是不理性的.
真是為反而反.
所以其實是填補天主教教導的不足.
所以這個就應該隨著我們信主內.
我們個人應該由初信甚麼都不懂.
只是一股熱誠.
慢慢去到十年二十年.
開始你應該成為第二代.
你自己信仰的更新.
開始對一些真理有融會貫通.
有一些人生閱歷.
你應該進入你自己的第二代.
對真理比較全面的掌握.
對神又有經驗.
又有真理的確具.
這樣自己建立自己的第二代.
第三代.
基督教這樣的教會.
就不斷累積.
越信主內越有經驗.
越有那個不是單純的一股熱心的信徒.
這個中層的信主超過二十年.
開始有小朋友有家庭有穩定工作.
就是教會很重要的骨幹.
如果這一代.
即是他個人的第二代.

$^{1321}$經過十年二十年.
都不思長進.
聖經我從來沒有讀過一次.
不過我是教會的執事.
這樣就大件事了.
他仍然停留在他自己初信的一知半解.
拿著一支金句就做執事了.
可能在社會上很厲害.
但在不同的教會是有人在的.
真是很大件事.
如果他做到終身濟長老.
完全對聖經沒有興趣.
每次回教會都會睡覺.
但做執事坐在那裡就有很多事要做.
你想想那個教會是否一個福氣呢?.
我沒有批評任何一個人.
我只是把現象說出來.
所以是需要有第二代的嘉爾文.
整理整套系統的神學.
所以如果信主內.
而自己從來都不讀經.
教會還會覺得這個人更適合做執事.
因為信主內.
真是慘的.
是一件慘事.
不是針對哪一個人.
而是他根本沒有長進.
他自己仍然停留在初信階段.
他怎可能承擔教會這麼重要的真理解釋呢?.
所以你看看嘉爾文的基督教要義.
是多麼重要的.
把馬丁路德這一股熱誠.
轉化成真理.
去模樣那個時代的歐洲人.
你可以想像到.
這就是帶來的復興.
因為不是每個人都是蒙召做全部的.
但是每個信徒當時是必須要認識.
自己在神面前的身份.
嘉爾文的後來的人.

$^{1361}$當然是把他極端化.
當時我們一而再說過很多次.
無條件的揀選.
是指天主教的恐嚇.
你離開我們教廷.
你就沒有救恩了.
所以基於這個模樣.
嘉爾文很清楚.
透過聖經告訴他們.
你們現在脫離天主教.
我們是在神的救恩裡面.
是無條件被揀選的.
所以這個連消帶打就是.
你是被神愛.
被神認識.
被神揀選.
所以每一個基督徒原來都是神的計劃.
是有呼召在你的人生.
你可以想像到.
這個在模樣上帶來的就是.
每一個信徒.
你有你自己的人生的責任.
我引述了兩段.
在中間第五百多頁之後.
我忘記了二十多三十頁.
是講信徒的生活.
盼望永生.
我引述了一點.
盼望永生就會產生正直的行為.
所以你可以想像到.
就是要有自己的人生照明.
清楚自己的夢照.
清楚自己的永生.
你是盼望永生的.
你就在地上.
我們不應放縱.
當然是中文的翻譯.
就是加爾文的基督教要義第卷二那裡.
我們不應放縱地使用這些福分.
福分就是神給我們享用的萬物.

$^{1401}$或者貪心地追求錢財.
而是在神對我們的呼召上.
忠心盡本分.
所以加爾文給到的那個時代就是.
你是有照明的.
你要發揮你的恩賜.
你要在你的人生生活的重點就是盼望永生.
所以盼望永生產生的是正直的行為.
接著還有.
照明的人生.
他說.
節儉世上的財產都是神所交託給我們的.
我不印這麼多.
如果他們保持這樣的節制.
他們將會在敬虔上大大地長進.
所以基督教要義對那個時代的影響.
就是大家都知道自己是離開天主教.
但是我們是有照明.
而且是神對我們的照明.
在耶穌裡面.
我得到神的救恩.
我現在發現原來神在我的生命裡面是有計劃.
那我就去發展我的人生.
但是同時間.
這個發展原來是平衡的.
而且是根據聖經的.
所以你既要渴望神給你的財富.
是神給你去享用的.
這是第一段所說的.
這一段是這麼說的.
我讀多一點.
B. 最有效約束自己的方式.
是藐視今世和默想天上的永生.
從此就尋出兩個規則.
我們用細物的時候.
好像不用細物.
因為上文提到.
一切萬物都是由神所賜.
而且是給我們去享用的.
但是原來是要有制約的.

$^{1441}$這個制約就是.
我們的人生不要貪婪.
不要在我們一種自私的裡面.
去享用我們自己的.
沒有節制地去享用.
相反.
一個對自己的人生有認知.
是神的照明的人生.
一個基督徒.
他就會有這種.
某程度上是一些聖經的心法.
進入了內心.
因為聖經主人是有教導的.
其實保羅也有說.
用細物好像不用細物.
這樣的一個心態.
這樣的一個人生價值觀的座右銘.
所以保持這種平衡.
既有照明.
知道我們所做的是為神而做.
既有人生的清楚的哲理.
你就會在這個平衡的人生.
發展出一種.
你知道什麼才是最優先.
什麼不是那麼重要.
這樣的時候.
就很不同了.
你就會有這種.
就是嘉爾文之後的.
清教徒的這一班人.
帶來的教會.
那個社會的改變.
所以這裡就告訴我們.
第二,第三個我們要思想的.
第一個就是.
在牧養上對那個時代.
我們要有一個把脈.
那個時代有什麼恐懼.
有什麼缺陷.
有什麼需要.

$^{1481}$我們要去回應.
是要透過真理去回應.
而牧養他們的心靈.
譬如美國的越戰.
一代的年輕人.
很大的失去了方向.
開始有嬉皮士的文化吸毒.
濫交的文化出現.
所以那個時代的教會.
就要回應他們的問題.
譬如我們回到49年.
100萬內地的難民.
蜂擁來香港.
住在牧屋區.
沒有教育.
沒有屋住.
很多的問題.
教會那個時候.
中環的聖公會.
大的教會.
不能夠一下子.
馬上回應到.
那個社會的需要.
因為他們根本就是在香港.
中環最好的地段.
他們習慣的是.
歐洲西人的信徒的聚會.
和本土的.
特別是這些難民更加切.
根本不能有任何的切入點.
但是神怎麼用呢.
就是全中國所有宣教士.
要立即撤退來香港.
這些宣教士各自的宗派.
立即就開教會.
回應難民的工作.
因為根本有些難民.
陪伴這些宣教士.
逃難來香港.
就已經立即做這些難民事工.

$^{1521}$聖公會.
不好意思.
沒有這樣開名.
根本回應不到.
你可以想像到.
有時是很急的時代.
有些是那個時代的內心的狀態.
很需要一些很有智慧的牧者.
神學家.
去highlight(強調)到那個問題.
然後讓不同的教會.
做那個處理.
靠近.
或者是透過講道去牧羊.
但是你可以想像到.
這個就是牧羊的轉移.
教會就必須要不斷在那個改革的心靈裡面.
要回應到那個時代.
如果已經落後了.
我們就要有一個這樣的轉變.
這個就是牧羊上的一個轉移.
另外就是剛才講的教育.
包括聖經的教育.
我們怎樣去保濕那個欠缺呢.
那個的信徒.
他現在那個時代.
現在可能大家很緊張的就是.
每天起床就問deep seat(深層座位)一些問題.
問到他不懂回答就最好了.
我想現在每個年青人.
最喜歡.
這個沒有問題的短時間.
但是仍然問一個問題就是.
這個時代的信徒.
老中青.
為什麼不讀經.
為什麼對讀經是沒有胃口.
如果作為一個父母.
看到一個小朋友不斷的嘔吐.
不斷的上吐下瀉.

$^{1561}$你都不知道應該餵什麼給他吃.
你還繼續帶他出來吃自助餐.
你根本就沒有可能了.
所以到底現在的情況是不是已經改善了呢.
他已經止了嘔吐.
現在是不是需要有些竹筍呢.
還是現在需要喝一些補抗力呢.
還是他開始需要有一點點的營養呢.
這個就是時代對那個狀況的把脈.
也是宗教改革的未來.
我們必須要去面對.
我們感孫和洪恩都很感恩.
每年都有聖經,獨經日營.
都多人參加的.
大家都很想聽.
但是那個延續.
是不是應該建立一個自己讀經的胃口呢.
而且這個胃口.
是能夠成為一個時代的人.
都很想去得到的呢.
譬如深圳現在這麼多美食.
星期一上去深圳.
就看到很多老人家這麼辛苦.
走很多路.
都要上去吃東西.
你可以想像到.
那個是胃口的問題.
是一個趨勢的胃口的問題.
所以我們要問的.
來到21世紀.
什麼令到信徒那個屬靈的胃口.
讀聖經的.
真的自己愛讀聖經的胃口.
怎樣能夠建立呢.
這個是第二方面.
透過伊拉斯夢.
透過加爾文.
透過他們所說的基督教要義.
我們要問第二個問題.
第三個問題.

$^{1601}$沒辦法逃避的了.
要改革就要犧牲的了.
就要有付出的了.
就要有一代的人.
他們面對的死亡.
都在所不惜的了.
我就稍微列舉一些出名的.
我們上個星期已經說過.
這個洛克斯.
給血腥瑪麗.
畢竟瑪麗皇后.
成三千人基督徒.
這樣被他殺了.
你可以想像得到.
那個時代的恐怖.
還要是他是天主教徒.
不是那些邪教.
或者是其他的宗教.
自己聲稱自己是天主教徒.
殺基督徒.
那你可以想像得到.
這個是我們要面對的代價.
21世紀.
文明是我們一定希望的.
需要的.
但是基督徒一定要有心理準備.
至少要有一些.
真是有領導性的人.
他們要有這個心理準備.
是要受苦或者殉道.
約翰胡斯被火燒死.
威克利夫他沒有被燒死.
但是他死了之後三十年.
就是約翰胡斯被燒死的時候.
順便柴還沒有用完.
就拿威克利夫出來燒.
燒了他的書.
燒了他的屍體.
你可以想像得到.
天主教要替天行道.

$^{1641}$在那個時代就什麼都不管了.
你反對就要被攻擊了.
現在約翰洛克斯.
早年還沒有去日內瓦.
還沒有去瑞士的時候.
就受到蘇格蘭的瑪麗皇后.
英國的血腥瑪麗的逼迫.
被關在船艙底.
做了十年的奴隸.
你可以想像得到那種痛苦.
本人約翰.
寫天路歷程的很出名的作家.
他是1628至1688年的人.
都是60歲.
但是在他最壯年的時候.
就坐了12年牢.
在12年裡面.
他就在牢裡面寫天路歷程的第一版.
所以你可以看到.
有很多先聖先賢.
為一個真確的信仰要付代價.
所以21世紀.
我們不鼓勵一些無端白事的無謂犧牲.
但是我們的老宗.
有沒有這個仇苦的殉道呢?.
基督教是否已經不再講血腥.
不再講殉道.
不再講仇苦.
稍為一些東西硬一點.
人們都不可以聽了.
你這樣講他們會很反感.
我們就要問.
這是否我們今天需要思考的.
對未來的基督教.
未來的教會的發展.
是否我們已經將仇苦和基督教.
有一個分道揚鑣呢?.
現代人不應該講仇苦的了.
應該講怎樣得平安.
這些本來不是錯的.

$^{1681}$但如果真的要為真理而受苦.
或者要站在真理的一邊.
那又如何呢?.
我們有沒有為了真理而要站穩.
無論代價有多大我們都願意呢?.
這就是我今天想到的三個大方向.
和大家總結.
透過我們一直看.
當然用了九堂去看過去.
但用一堂我們用一個很簡單的總結.
思想一下我們怎樣去承擔福音.
務求要令到那個時代的人.
例如現在有一大班的基督徒.
去了歐美,英國,加拿大,溫哥華.
多倫多比較少.
其他歐洲地方應該都比較少.
那怎樣去祝福當地的人.
不要只是為了生存找到教會就算了.
需要有一個宏觀去看.
既然神帶領一大班人.
是有信仰的,有使命的.
去到當地了.
那怎樣令到當地華人.
可能當地的歐洲人.
他們得到神的祝福呢?.
這就是一個時代不斷轉變,流動的時候.
我們需要有這些反思.
所以都想用這一堂.
來和大家做一個反思.
透過過去的前人.
他們留下的腳zoung .
他們的服侍.
他們已經服侍了他們那個時代的人.
那個時代也真是跟隨宗教改革.
而帶來一個大復興.
而帶來一個新的現象.
基督教一直很有力的在前進.
所以我們從來不需要擔心.
基督教會沒有了.
主耶穌回來後才沒有了.

$^{1721}$所以我們不需要擔心.
因為有聖靈,有神會帶領著教會.
但是我們那個時代.
能不能回應到那個時代的需要呢?.
這就是我們大家要一起去問的問題.
所以不要再停留在一些很低層次的事情.
那些無關痛癢的事情不要浪費時間.
反而回到教會.
怎樣集合大家的力量去傳福音.
回應那個時代的需要.
才是我們要認真去想的問題.
不要回到教會就想起那些雞毛蒜皮的事情.
又不開心的一天.
那個崇拜完全是為了那個雞毛蒜皮的.
將它這樣擺放到亂了.
整堂就不開心了.
浪費時間.
然後又走了.
然後又完全忘記了.
第二個星期不斷累積.
52個星期 52件不開心的事情.
回來教會就不是這樣的.
要想這些宏觀的事情.
我們對這個時代有什麼虧欠.
這個時代我們有什麼可以令他們的恐懼.
說都說不到的恐懼.
怎樣能夠透過福音讓他們得到呢?.
是教會要不斷想的問題.
如果不是的話.
神可能就會將我們擺放.
使用其他人.
其他的教會去做這個工作.
好了 完了.
因為最後的這篇PowerPoint.
那就感恩了.
我盡了力.
在我自己有限的時間.
來盡力找資料.
當然真的是淺談.
下個星期就回歸聖經.

$^{1761}$我就沒那麼辛苦了.
起碼有聖經.
我盡力了.
我不敢保證.
盡力都用PowerPoint.
還有會有分段.
去到重點的就會講長一點.
絕對不可能一堂講完一章的.
因為羅馬書實在太多東西要解釋.
所以希望有幾個進路.
在不斷講解上下文的時候.
就要教大家解經.
盡量學習怎樣去解釋聖經.
上下文的Context的解經.
還有一些當時的背景.
盡量都來講解.
務求有趣味地.
和深入淺出地.
大家都能夠讀完一卷的羅馬書.
希望大家留一些記錄.
最好就做一些筆記.
我就盡量.
可以的話我就做PowerPoint.
還有會有些輔助的筆記等等.
盼望大家一起學習.
這樣每個星期都能夠迫到我有機會.
認真去研讀.
是辛苦的.
因為有其他工作.
但是樂趣是很大的.
所以星期四晚回家.
就很舒服很開心很輕鬆.
因為星期五上午就沒那麼忙.
心情就會很舒暢.
因為做了自己很喜歡做的事.
又講解完.
有一班弟兄姊妹又很想繼續學習.
滿足是很大的.
一起在神的話語.
在宗教改革裡學習.

$^{1801}$剛才有弟兄姊妹說.
17世紀的衛斯理約翰.
亞米紐斯.
其實都很值得講.
16世紀有一個很重要的清教徒運動.
其實是另外一層的大復興.
所以將來有機會.
我們又間歇地講一下.
都好的.
其實19世紀的普世宣教運動.
其實不可能不講.
都要了解.
所以一涉及到普世宣教運動.
肯定又要講由1900年的義和團事件.
到現代的中國教會事.
都需要有時我們有少少的認識.
就會知道就會好.
今晚就停在這裡.
多謝大家.
感恩我們又可以透過一堂訊息.
彼此勉勵.
我們真的需要牧養一個時代.
任何教會都需要牧養一個時代.
那個牧養不只是講一些活動.
而是那個時代的心靈需要.
中世紀他們那種懼怕死亡.
和煉獄的神學.
扭曲了聖經.
令到信徒更加懼怕.
但是多謝主.
馬丁路德講的簡選.
無條件的簡選.
和白白的恩訊清義.
真是回應到整個時代的懼怕.
而且又帶給信徒.
透過嘉爾文.
有這個人生照明.
以致很多的信徒.
信了主之後.
為主而活.

$^{1841}$在他們的生活.
在他們的職場.
來發展和神的關係的信仰.
因為信仰已經落實到家庭.
落實到信徒的內心.
落實到他們自己的生活裡面.
不再是天主教的時候.
在每個星期的彌撒.
他們只是遠遠觀望.
沒有任何的參與.
更加不會有任何的獲益.
我們盼望.
透過這些反思.
都給我們一班弟兄姊妹.
對這個信息的了解之後.
又有一個自己個人.
或者一個小組的大家的反思.
願神幫助.
我們見到聖經的普及.
就是由宗教改革開始.
不是以前的人不需要聖經.
而是以前因為有錯誤的代理制度.
令到聖經被封閉了.
多謝主透過馬丁路德他們.
這個屈克利夫他們.
打斷了這些錯誤的迷思.
使到原來人人都可以有機會讀聖經.
這個普及的運動.
我們又更加見到.
每一個時代都有信導者.
都有為信仰.
站在真理的行列.
求主幫助.
其實主耶穌,保羅,彼得.
他們全部都是這樣殉道的.
我們看到中世紀的時代.
黑暗的天主教.
反而成為迫害其他.
走真理的人.
這種諷刺.

$^{1881}$求主憐憫.
使到教會不要只是以為.
我們在做一些合神的事.
而我們自己不是在真理裡面.
求神幫助我們.
今晚我們有這幾方面.
其實很大方向性的反思.
更加盼望信徒每一個我們自己.
都在主面前有一個很敏銳的心靈.
下星期我們要回歸羅馬書的研讀.
求主幫助我們.
讓我們真的去愛慕.
一卷這麼重要的書卷.
我們將每一章去講解的時候.
都令到我們透過保羅的有系統的講解.
我們認識一個很整全的福音.
我們獻上感恩.
禱告奉耶穌基督的名.
阿們.
\newpage



\section{羅馬書}
\label{sec:mCrcbeplxo4}
\textbf{2025-02-20 周四晚查經班《羅馬書(1)引言:保羅其人其事》 講員:曾裔貴牧師}
\newline
\newline
連結: \href{https://youtube.com/watch?v=mCrcbeplxo4}{\texttt{https://youtube.com/watch?v=mCrcbeplxo4}} ~~~~ 語音日期: 2025-02-22
\newline
\newline
\hyperref[sec:q7QZXuHzcds]{\small{< < < PREV SERMON < < <}}
~
\hyperref[sec:index_chronic]{\small{[返順時目]}}
~
\hyperref[sec:index_scriptual]{\small{[返順卷目]}}
~
\hyperref[sec:TdP7rHPK3XI]{\small{> > > NEXT SERMON > > >}}
\newline
\newline
$^{1}$我們可以一齊開始今晚.
我們會研讀羅馬書.
都很盼望神給我們大家有這份學務.
認識這本書.
在其中發掘到一些屬靈的寶藏.
盼望神帶領我們.
這本書都是相當長的.
16章.
求主給我們大家有系統的研讀.
去明白我們這樣仰望,感恩,禱告.
奉耶穌基督的名.
阿們.
很多謝各位.
我們一齊有機會去看這本書.
算是有系統地跟大家去研讀.
我將它放大.
為何好像放不到大呢?.
先看看.
是否要收起它才放大呢?.
不太清楚為何放不到大.
與這件事無關.
是有點奇怪的.
一到這個螢幕.
可能要關掉這個.
真的不太清楚.
總之你們看到就算了.
讀羅馬書首先要知道.
為何保羅要寫一卷這麼大的書卷.
拿聖經.
大家在家中就方便.
預備聖經在自己的桌上.
因為我們都會揭得比較多聖經.
為何我不能放大這個畫面.
奇怪.
不太清楚.
隨它吧.
大家如果看到就算了.
不要理會.
要問這本書寫作的動機是甚麼.
我們要拿捏.

$^{41}$因為一卷這麼大的書卷.
裡面肯定有很多分題.
都是重要的題目.
例如關乎恩訊清二.
由第三章一直到第八章.
都是講恩訊清二.
恩訊清二帶來基督徒的成聖.
三到八章肯定是這個主題.
但問題是這是否整卷羅馬書的中心訊息呢?.
或者寫這本書的目的呢?.
因為一卷書.
如果是書信.
肯定是有一個主題的.
所以我們要找出主題是講甚麼.
例如第九到十一章.
明顯是講到猶太人在救恩裡面.
他們扮演甚麼角色.
猶太人在耶穌未來之前.
他們是怎樣看待自己的身份.
耶穌來世是一個徹徹底底的改變.
從舊約過渡到新約.
猶太人在耶穌降生之後.
當然更重要的是他們離世復活升天之後.
保羅怎樣理解當代猶太的同胞.
他們現在在救恩裡面.
他們在救恩的門外.
還是在救恩裡面呢?.
這就是第九到十一章來探討的.
第十二章.
大家亦都很熟悉的.
要將身體獻上.
在生活上人與人之間的倫理關係.
又怎樣去銜接整卷書的主題呢?.
所以如果要拼湊起來.
我們便要問整卷羅馬書是講甚麼呢?.
我們讀完整卷書.
必須要自己去問.
到底這卷書保羅是想他講甚麼呢?.
首先我們都要了解.
當我們讀這卷書.

$^{81}$既然是這位使徒.
他寫這封書信給羅馬教會.
我們都要對保羅.
他自己當然是神.
透過聖經讓我們知道的保羅.
他是一個怎樣的人.
又是未信耶穌之前.
信了耶穌之後.
和現在他成為了耶穌的使徒.
他又有甚麼心路歷程的轉變呢?.
我們如果能夠掌握得越多他的資料.
我們亦會對智慧行間的羅馬書.
有更多的脈絡了解.
因為羅馬書是保羅時代.
二千年前寫的書卷.
所以如果我們稍為對當時羅馬帝國.
有一些歷史方面的了解.
又可能會幫助我們.
尤其是對保羅的用字.
我們又可以了解.
所以讀一卷書.
最好能夠用不斷發問的問題.
對大家有幫助.
我們可以先來看看.
到底保羅.
很快的.
因為有些弟兄姊妹可能是初信.
未必太熟悉.
這裡簡單說一點.
雖然和羅馬書直接的內文沒有關係.
但是我們先稍為了解一點.
如果大家很熟悉保羅的生平.
其實不要緊.
大家當聽一聽便可以了.
我們看看《小道行傳》的七章.
今天我沒有帶大聖經.
所以很難看.
因為細本已經超越了眼睛的老化.
七章的五十八節.
七章是《小道行傳》.

$^{121}$有一位很重要的執事.
司提反洵道.
大家留意七章五十八節.
聖經作者路加醫生.
記載了一個伏筆.
告訴我們一個年輕人.
將來他會很重點地介紹這個人.
如何遇到耶穌.
所以現在先說他未遇見耶穌之前.
他是一個怎樣的人呢?.
七章五十八節.
司提反洵道到城外.
用石頭打他.
作見證的人將他們的衣裳.
即是行刑的人.
放在一個名叫素羅的青年的腳前.
他們正用石頭打司提反的時候.
他呼求說.
主耶穌求你接納我的靈魂.
然後他跪下來大聲喊著.
主不要將這罪歸於他們.
說了這話就長眠了.
路加醫生記載了一件很恐怖的洵道事件.
有一大群人.
每個人拿著石頭去扔司提反.
將他扔死.
但司提反未死之前.
他見到天上開了耶穌.
來跟他說.
不是耶穌跟他說.
主耶穌求你接納我的靈魂.
這是素羅.
他沒有扔.
但他站在群眾附近.
看守他們的衣服.
我們可以思想一下.
他望著一個說自己跟隨耶穌的司提反.
為他作見證.
現在被人用石頭扔死.
但到底保羅未信耶穌前就叫素羅.

$^{161}$素羅是一個很大的名字.
因為保羅是汴雅文支派出現了一個素羅王.
所以保羅自己又叫素羅.
其實是一個很大的支派名字.
我們看看第八章這一段.
或者第八章第一節開始.
「素羅也贊同處死他」.
從那一天開始.
其實八章一節應該歸納到七章五十九節.
因為當時聖經是沒有分段的.
沒有章節的.
所以其實是一個連續的故事.
「素羅贊同處死司提反」.
接著從那天開始.
耶路撒冷的教會.
正正式式遭受一個明顯的公開的逼迫.
除了使徒之外.
眾門徒都分散到猶太和撒瑪利亞的各處.
有些虔誠的人將司提反埋葬.
而且為他大大的哀哭.
留意第三節.
「素羅卻殘害教會」.
「哀家哀戶的進去」.
「拉著男女關在監裡」.
這就是素羅未信耶穌之前.
他所遇到的事情.
他所可以說是生在一個大時代.
因為見到很多他的.
可能是前輩.
或者是一些很勇敢的猶太人.
就要來清算這些基督徒.
來攻擊這些說自己跟隨拿撒勒人的基督徒.
所以素羅就在這個時代底下.
他是附和贊成要來處死和捉拿這些基督徒.
來坐牢的.
所以一定是跟他的信仰.
猶太教的信仰.
那份熱誠有一個很直接的關係.
所以我們先有這個背景.
我們就可以知道原來素羅都是相當之火熱.

$^{201}$不過他的熱心.
只是仍然在猶太教裡面的教育底下的熱心.
接著我們就跟大家揭開第九章.
《使徒行傳》的第九章.
因為記載完第八章.
有一些奇遇.
在該撒尼亞.
其實應該是先從撒瑪利亞開始.
作者路加就很有系統.
就說猶太的信徒.
因為碧巴的緣故.
他們就分散在猶太各處.
還有在撒瑪利亞.
在撒瑪利亞傳道.
在撒瑪利亞就遇到一件事.
就是這個肥利帶了一個埃提瓦伯的太監信主等等.
然後就記載肥利.
這個人將來在二十二章,二十三章.
就會再與保羅有機會再有相遇.
因為肥利有幾個女兒都是先知.
最後就是肥利第八章.
第八章的三十九節.
他們從水裡上來.
就是菲利與太監洗禮.
在曠野立即幫他洗禮.
因為這個太監在很遠的地方上來守這個逾越節.
但是.
對不起.
應該是五旬節.
守完之後.
他就要回去.
在路上就遇見肥利.
其實是肥利被聖靈差遣.
去帶他太監信主.
所以當他信主洗完禮.
肥利就突然離開.
三十九節.
他們從水裡上來.
主得靈把肥利提了去.
太監再也看不見他了.

$^{241}$就歡歡喜喜地上路.
即是回到自己的地方.
後來有人在阿索多遇見肥利.
阿索多是比較遠的地方.
歐洲的地方.
他走遍那地方.
在各城宣揚福音.
一直去到蓋薩尼亞.
即是告訴我們.
作者又是很巧妙地告訴我們.
那些傳福音的信徒.
因為碧巴從耶路撒冷.
輾轉到猶太地方.
去到撒瑪利亞.
後來肥利作為耶路撒冷教會的領袖.
是七個執事裡的其中一個.
去到蓋薩尼亞.
在蓋薩尼亞似乎定居在那裡.
因為將來第三次傳道.
上到蓋薩尼亞的時候.
就去到一個叫傳福音的肥利的家裡.
來到他相交住了幾天.
應該是23還是22章.
肥利已經不再是執事的身份那麼簡單.
已經變成一個叫傳福音的肥利.
是作者路加告訴我們的.
是一個很寶貴的名字.
因為保羅已經傳道很多年了.
現在在他鄉遇到固茲.
其實當時是想殺他.
因為斯提反也殺了.
其中一個執事也殺了.
不僅殺了肥利.
誰知道肥利跑去撒瑪利亞.
殺不了他 抓不了他.
但是經過很多年之後.
大家各自在不同的神安排的工場.
保羅去到蓋薩尼亞.
其實和大家揭曉一下.
反正揭到這裡了.

$^{281}$我看看是22還是23章.
如果沒有記錯可能是23章.
就是第三次的旅行報道.
上耶路撒冷之前.
保羅就在那裡.
應該沒有那麼遠了.
可能是22章.
對不起 是21章才對.
這段聖經也很特別.
很值得我們看看.
也是和保羅和肥利有關.
第七節.
我們從推羅行完航程.
來到了多利買.
問候那裡的弟兄.
和他們同住了一天.
第二天我們離開那裡.
來到蓋薩尼亞.
就進了傳福音的肥利家裡.
和他同住.
他是那七個執事裡的一個.
七個執事.
第一個殉道就是斯提反.
當時沒有相機.
他家裡有張相.
是七個執事就職來的時候拍的.
你想想保羅見到斯提反.
是他親手.
當然他沒有扔石頭.
不過是看守著那些扔石頭的人.
那些衣裳.
還有聖經告訴我們.
他很贊成要處死斯提反.
現在在他鄉.
遇見了很多年之後.
仍然在神的宮牆裡服侍的肥利.
聖經這樣描述他.
路加醫生.
在這個時候跟隨保羅去傳道.
就說他是七個執事裡的一個.

$^{321}$他有四個女兒.
都是還沒有出嫁的.
都是會說禦言.
都是一起事奉神.
所以我們在那裡多住了幾天.
過了幾天之後.
有一個先知阿加布.
從猶太那裡下來.
就說要像一個先知救藥的時候.
有時候會像做一些大龍鳳一樣.
實物教材.
好像做一套話劇一樣.
捆綁著.
拿了保羅的腰帶.
放在自己的手腳.
聖靈這樣說.
猶太人在耶路撒冷.
都要這樣捆綁腰帶的主人.
就是保羅.
所以就把他交給外邦人.
所以第十二節.
我們路加醫生.
可能是肥利.
我們聽見這些話.
就跟當地的人.
苦勸保羅不要上耶路撒冷去.
於是保羅回答.
你們為什麼這樣痛哭.
使我心碎呢?.
我為主耶穌的命.
不但被人捆綁.
就是死在耶路撒冷.
也是願意的.
第二次提願意.
之前就很願意很贊成斯提反遜度.
看著可能拍手掌.
可能該死.
這個人跟一個這樣的人.
拿撒勒人.
跟他說他是救世主.

$^{361}$但現在保羅自己.
反而叫那些跟他一起服事的人.
不要攔阻聖靈的工作在我的身上.
我要在耶路撒冷被捆綁.
這是聖靈的工作.
所以你們不要令我心碎.
令我難堪.
我為主耶穌的命.
不但願意被捆綁.
就算死在耶路撒冷.
他也願意.
這些豪情壯語.
對弗利,路加醫生.
對其他四個女兒.
相信那一晚聚會的教會.
充滿了聖靈同在的勇氣.
那份對主耶穌的救恩的效忠.
你可以想像到.
培靈會可以在任何地方.
只要有兩三個人奉主的名.
大家為神的國度去守望.
那裡就有聖靈有主的同在.
保羅現在甘心樂意.
這樣擺上自己的生命.
這裡有一個很大的功課.
人家用了十多年.
二十年.
應該沒有二十年.
人家由一個敵黨神.
敵黨基督.
甚至有行動捉基督徒.
到現在願意自己的生命.
有一個很大的毀身.
所以二十年.
就看我們大家在主耶穌的救恩裡面.
我們的成長.
說多不多說少不少.
在香港的教會.
二十年大有人在.
二十年之後.

$^{401}$連看聖經都不感興趣.
大有人在.
一來到教會.
現在手機隨身.
除了站在這裡.
背著小道信經之後.
坐下來可能未必是看聖經.
坐在這裡看手機.
看WhatsApp.
都大有人在.
雖然我們沒有監控.
但是其實.
二十年.
由一個敵黨神的人.
變成為我為主耶穌的命.
就算死在耶路撒冷.
都是情願.
肯定背後不是那個人自己有什麼造作.
肯定是神的恩典.
讓保羅慢慢看到.
這個永恆生命的實在.
所以這裡是值得我們想一想.
歸主前後的保羅.
他的生命的成長.
所以這一晚的對話.
其實很美麗.
雖然當時我們沒有相機.
也沒有高清來錄下.
但是透過字面讓我們聯想.
我們已經看到.
一個就是為主.
由耶路撒冷.
輾轉去到.
撒馬利亞.
輾轉再去到蓋撒尼亞定居.
而他四個女兒.
都是講預言的.
在預言的角度裡事奉的.
這樣的家庭.
即是傳福音的家庭.

$^{441}$所以是很美麗的一件事.
再加上保羅路過.
知道肥利在這裡.
就進去住.
你想想那些新仇舊怨.
在耶穌裡面根本沒有.
肥利又認識到.
始終他一定會知道.
保羅是搞亂了天下的.
到處被人抓的.
肥利看著他的生命的轉變.
相信肥利為保羅祈禱.
有很多的禱告.
因為肥利自己就定居在.
蓋撒尼亞去做福音的工作.
所以現在路過.
保羅上耶路撒冷之前.
就有機會來見見.
這一位的前輩.
可能大家差不多年紀.
大家都是猶太人.
大家曾經在耶路撒冷.
都有過一個短暫的居留.
保羅就是當時的拉比.
但是肥利呢.
我們就不太知道他做什麼工作.
只知道他被選出來做執事的.
七個裡面其中一個.
在這樣的情況下.
保羅就有這個機會.
去見見曾經在耶路撒冷.
第一批為主耶穌.
作見證的門徒.
接著我們就回到第九章.
大家不要介意.
因為我怕有些初訊不是太熟悉.
反正都是有系統的研究.
都講一講.
給一些不是太熟悉的弟兄姊妹.
有些少的輪廓.

$^{481}$對我們進入羅馬書.
又會有一定的幫助.
第九章四途行傳.
掃羅不斷用威嚇.
凶悍的口氣.
向主的門徒說話.
他去見大祭司.
要求發信給大馬士革的國會堂.
如果找到信奉者道的人.
無論男女都准他捆綁.
帶到耶路撒冷.
當然是帶回來審問.
用他們的宗教法庭來審問定罪.
無論男女.
第三節.
掃羅在途中.
即是他已經拿了文書.
拿了大祭司的搜查令.
逮捕的文書.
接著.
掃羅在途中.
將到大馬士革的時候.
忽然有一道光從天上下來.
四面照射著他.
他就突然在地.
聽見有聲音對他說.
「掃羅!掃羅!你為什麼迫害我?」.
如果在二十六章.
保羅在說這件事的時候.
是用希伯來話跟他說的.
所以他聽到的聲音是希伯來話.
很特別的.
造物主可以用任何的語言.
用你保羅的希伯來話.
當然保羅是精通希臘文.
羅馬官方的語言就是拉丁文.
但用希伯來話跟他說.
聽見有聲音對他說.
「掃羅!掃羅!你為什麼迫害我?」.
他說.

$^{521}$「主啊!你是誰?」.
這個主不知道是誰的聲音.
是一個很一般性的稱呼.
第五節下半.
主說:我就是你所迫害的耶穌.
起來進城去.
你應該做的事.
必有人告訴你.
同行的人站在那裡.
說不出話來.
因為他們聽見聲音.
卻看不見人.
他們看不到大光裡面的東西.
但其實在二十二章提到.
其實他們不知道那個聲音的意思.
但掃羅是知道的.
所以很特別的.
同在那個地方.
不是說太遠聽不清楚.
是他們聽不到那些聲音的內容.
但唯有掃羅聽到.
是很特別的.
大家都看見大光.
第八節.
掃羅從地上起來.
睜開眼睛.
因為很強光.
很自然地要閉上眼睛.
但他睜開眼睛的時候.
突然失明.
竟不能看見什麼.
有人拉他的手.
領他進了大馬士革.
當然有一段路程.
作者很簡潔地告訴我們.
不在話下.
可能有人一直拉著他.
因為他突然間失明.
看不見東西.
所以當進入大馬士革之後.

$^{561}$聖經這樣形容.
第九節.
他三天什麼都看不見.
也不吃也不喝.
對他來說是一個很震撼的訊息.
他為什麼不吃不喝.
是什麼原因.
是內心的痛悔.
還是非常憤怒.
非常疑惑.
那個光的聲音告訴他.
我就是你所迫害的耶穌.
二十六章他自己在描述的時候.
這裡是路加醫生的記載.
二十六章他跟千夫長說的時候.
二十六章跟千夫長.
和那時的喜伯來人說的時候.
他就說.
我就是你所迫害的拿撒勒人耶穌.
路加就很簡單.
就說是耶穌.
所以拿撒勒人耶穌跟我說話.
現在我一起來.
睜開眼就瞎了.
接著有人拉著我.
讓這個.
記住他是一個很有朝氣.
拿著大祭司的文書.
準備去捉拿那些基督徒.
似乎是已經知道在哪裡.
不是去到才去找.
才去搜捕.
大馬士革的人類似猶太人.
已經點上了.
在哪個地方就有這些基督徒.
你來去捉就行了.
所以全部標籤了.
但是現在真正被人點上的是保羅自己.
現在瞎了.
所以他三日三夜不吃不喝.

$^{601}$是很值得我們去想一想.
他在做什麼呢.
耶穌有沒有再跟他顯現呢.
在這三天裡.
似乎沒有記載.
是讓他自己不斷的去思前想後.
一個對猶太教的聖經.
無論是Torah.
整個他們的經卷.
舊約.
是相當純熟.
認識的這位保羅.
現在有三天要沉澱自己.
那個生命.
自己在走的人生的方向.
是錯了.
因為除非那個光是在耍他.
是一個現在的電話騙案.
那當然沒有問題.
那些網絡的化身.
那當然沒有問題.
但是現在古代哪有這麼先進呢.
那個光說話他又聽到.
又對號入座是耶穌.
那就不是開玩笑了.
所以保羅他要用這幾天時間.
都完全沒有興趣吃東西的意思.
來讓自己去反思.
當然過了三天了.
第十節.
那時在大馬士革有一個門徒.
名叫阿拿尼亞.
主在異象中對他說.
阿拿尼亞他說主我在這裡.
主對他說.
起來往那叫直街.
Straight Road的路去.
在猶大的嘉里.
在這個地方原來是猶大.
很普遍的猶太人的名字.

$^{641}$在猶大嘉里去找一個大數人.
大數人.
名叫素羅.
他正在禱告.
三日三夜.
他有很多時間不斷去反思.
他向誰祈禱呢.
肯定不會向耶穌祈禱.
因為他還沒有認識耶穌.
肯定不會向耶穌祈禱.
還是很多掙扎的.
但是他在禱告向神祈禱.
他們猶太人向神禱告.
所以其實想想這些細節都很有趣.
一個還沒有信耶穌的保羅.
他現在瞎了.
三日三夜又不吃又不喝.
現在一直在祈禱.
不過是向耶和華向神祈禱.
他不敢向耶穌祈禱.
因為他根本不知道耶穌到底是誰.
第十二節.
在異象當中看見了一個人名叫阿拿尼亞.
進來為他按手.
讓他能夠再看得見.
阿拿尼亞當然要有一個確認.
因為不是開玩笑.
那個人說明是大數人 素羅.
其實在大馬士革的基督徒.
全部都知道這個消息.
大家都已經躲起來了.
現在你叫我曝光.
叫我去找這個素羅.
所以第十二節.
阿拿尼亞回答.
主啊 我聽見許多人講道這個人.
說他怎樣在耶路撒冷多多輔待你的聖徒.
並且他在這裡有從祭司長得來的權.
並要捆綁一切求告你名的人.
第十五節.

$^{681}$主對他說.
你只管去 他是我所揀選的器名.
要在外邦人 君王和以色列人面前宣揚我的名.
我也要指示他.
為我的名必須受許多苦難.
這節聖經十五節.
也是其中相當重要.
是保羅後期常常引述的.
他自己的這個呼召.
所以保羅信耶穌.
從信那一刻就有了這個呼召.
而且這個呼召.
連一個普通的平信徒阿拿尼亞也知道.
耶穌對他說.
你去吧.
因為這個人不是普通人.
你一按手.
他就會變成使徒.
或者一個普通人按手在素羅身上.
因為有神的默示.
神就告訴阿拿尼亞.
那個人 沒錯.
現在是迫害教會的.
他有權柄.
用錯了在黑暗的勢力裡面.
要捉拿基督徒.
但是當你按手之後.
他就是我揀選的器名.
他就是成為將來改變世界的使徒保羅.
所以原來我們在教會裡面.
信主內.
就算我們不是使徒.
其實我們也要留意哪些人是可造之材.
其實保羅肯定不是可造之材.
他是可奏可坐.
因為迫害教會.
但是你不要管.
神的工作 聖靈改變一個人.
有時會脫胎換骨.
我們想都沒有想過.

$^{721}$他可以有這麼大的成就.
所以我們在教會裡面.
應該越安靜地默默地在神裡面祈禱.
神啊 你感動我.
哪些人是我需要按手的.
我要和他分享的.
讓他將來可以被神大用的.
這樣的人就是教會的屬靈的寶貴的前輩.
和屬靈人.
因為那個眼光是看下一代.
是看其他人可以怎樣被神用.
所以真的.
當然最好就是教小朋友.
因為你看著他成長.
你就可以不斷地將他潛移默化.
在神的愛.
在神的國度.
他將來成長 被神用.
你就肯定有你那一份.
但是其他的地方.
所以教會其實很多事情要做的.
你坐在那裡崇拜.
不是看手機.
專注敬拜神.
而是在不同的時間.
心裡面安靜地去祈禱.
這個人感覺他的內心有聖靈.
可以被神用.
你就為他更認真的祈禱.
所以亞拿利亞就成為了一個按手洗到保羅.
成為被神大用的開始.
所以在這裡我們可以看到.
你可以想像到.
將來保羅信主之後的傳道.
他對亞拿利亞的恩情.
是有很深的記憶的.
因為一個人 普通的信徒.
不怕死.
聽神的吩咐.
冒險來找我.

$^{761}$都不知道自己是怎樣的.
肯定知道的就是.
他已經收到風聲.
他是抓基督徒的.
我亞拿利亞就是基督徒.
我進去就是被他抓了.
但現在很明顯是有聖靈.
透過耶穌跟他說.
亞拿利亞.
所以普通的信徒.
一樣可以經歷聖靈的工作.
他就真是.
十七節亞拿利亞就去了.
進入那加.
把手按在訴羅身上說.
訴羅弟兄在你來的路上.
向你顯現的主.
你怎知道的呢?.
正正就是聖靈告訴亞拿利亞就知道了.
所以屬靈人參透萬事.
就不是說笑的.
是真的.
向你顯現的主.
就是耶穌打發我來.
叫你能夠再看得見.
又被聖靈充滿.
所以訴羅的眼睛.
按手訴羅的眼睛上.
立刻好像有輪一般的東西掉下來.
當然不是一些污糟的東西.
肯定是神蹟.
他就能夠再看見.
於是他起來受了洗.
可能是灑水.
或者可能有水缸.
跳下去洗禮.
受浸.
總之就是一開眼.
立刻受洗.
立刻成為主的兒女.

$^{801}$就是說他願意跟隨亞拿利亞所說的耶穌.
不再是拿撒勒人的目像的耶穌.
而是死而復活的救世主.
就是猶太人所等候的彌賽亞.
所以訴羅現在是完全釋放過來.
不只是眼睛看到.
是生命的釋放.
看到這位的耶穌.
我碧巴的耶穌.
不單止沒有懲罰我.
現在還大用我.
成為了當他受了洗之後.
十九節真的說他立刻吃飯.
吃過飯體力就恢復.
三天三夜沒有吃東西.
差不多了.
三天都不喝水.
都差不多很累了.
但是不會死的.
但是現在整個人就健壯.
身心靈健壯.
生命有一個重整.
一個大的手術.
一個洗禮.
洗去他的從前.
現在他是新造的人.
這些全部是在保羅的書信裡面說的.
他自己就是一個新造的人.
他自己擁有的.
舊有的猶太教的.
拉比的訓練的身份.
學問.
人際網絡.
這些全部不會忘記的.
但是他現在有的是新的一個概念.
就是對耶穌的教會.
耶穌的國度.
這個亞拿利亞未必是猶太人.
而且可能是外邦人.
你就可以想像得到.

$^{841}$這件事更加將猶太人.
外邦人.
來到去.
完全在主耶穌裡面合一.
又是保羅在往後的書信.
很著力要說的.
在主裡面.
不再分猶太人.
希利尼人.
自主的 圍爐的.
有學問的.
或者那些化外人.
是男是女.
全部沒有再分的了.
加拉太書的三章28,29節.
在主裡面.
都完全是合一的.
不再分那個種族等等.
你可以想像得到.
這個洗禮.
這個按手.
對保羅來說是一個很大的生命扭轉.
在這一刻.
真是來到去成為主耶穌的一個門徒.
不過不單止是門徒.
第十九節下半節.
「素羅和大馬士革的門徒.
一起住了些日子.
立刻在各會堂裡.
傳揚耶穌.
說他是神的兒子」.
立刻 幾天.
就在大馬士革.
普爾奇一帶.
敘利亞 對不起.
來到去傳揚耶穌是神的兒子.
猶太人今天還是有這個困難.
耶穌真有其人.
但是他不是神的兒子.
不是我們的彌賽亞.

$^{881}$但是保羅已經這樣來應訊了.
所以不是普通的應訊這麼簡單.
凡是聽見的人都很驚訝.
在耶路撒冷殘害.
求告明鏡.
豈不是就是這個保羅嗎.
為什麼他現在在這裡.
要這樣做呢.
豈不是要下來大馬士革.
要捆綁那些基督徒.
帶去祭司那裡嗎.
第二十二節.
但訴羅越發有能力.
搏倒住在大馬士革的猶太人.
證明耶穌是基督.
開始要做一個護教的辯論.
要跟猶太人辯論.
那就開正保羅自己的本行.
因為他自己就是猶太教的一個拉比.
聖經是相當純熟的.
相當認識的.
但現在要告訴自己的同胞.
猶太人.
要跟他們辯論.
要告訴他們.
我們現在對這位耶穌.
那撒勒人耶穌.
我們要有一個新的站身的角度去看他.
因為他由死復活.
他就是神的兒子.
一個人能夠死而復活.
肯定是有神跡在他身上.
但這個神跡不是想誇耀一個特殊的事件.
而是要印證耶穌的身份.
本來就是生命的主.
就是神的兒子.
所以當然這樣的傳道.
就會招致猶太人的攻擊.
所以輾轉保羅就要離開大馬士革.
第23節.

$^{921}$過了一些日子.
猶太人商議要殺掃羅.
但他們的計謀被掃羅知道.
他們就夜在城門守候要殺他.
他的門徒就在夜間.
用糠紙把他.
即是那些螺.
從城牆上來到去.
聚了下去.
掃羅到了耶路撒冷.
想與門徒結交.
大家卻要怕他.
不信他是門徒.
因為他不是去猶太教那裡.
他去門徒那裡.
所以他去錯地方.
人家就以為他想來捉拿我們.
所以就不和他交往.
第27節.
出現了一個巴拿巴.
只有巴拿巴接待他.
領他去見使徒.
把他在路上怎麼看見主.
主怎麼向他說話.
他在大馬士革怎麼奉耶穌的名放膽傳道.
都述說出來.
即是說很多人一見到保羅.
就已經不肯聽他講.
他悔改的經歷.
遇到主耶穌的生命的扭轉.
但巴拿巴就給他空間去說.
所以他怎麼怎麼怎麼.
三個怎麼.
就是遇見耶穌.
遇見大光.
遇見阿拿尼亞.
然後他怎麼傳道.
奉主的名傳道.
巴拿巴就有這種生命的特質.
給空間這個掃羅.

$^{961}$講述自己的見證.
所以聽他這樣說.
看到他內心有生命的主.
在他內心有聖靈.
所以深信不疑.
把他帶到使徒的面前.
讓使徒們再聽他的得救見證.
然後由使徒.
因為當時很有權威.
因為耶路撒冷是使徒來看守住.
去到見使徒的時候.
如果使徒們都認為.
保羅真是這樣信主的經歷.
就容易讓其他的門徒.
都能夠接納保羅這個生命的改變.
所以巴拿巴又是一個.
掃羅的信主初期事奉的.
一個很重要的伯樂.
剛才我們說.
我們要成為一個成熟的基督徒.
在教會裡留意.
神叫你何時按手.
現在我們又要學習巴拿巴.
他又是一個能夠.
看到人內心有聖靈的引導工作的人.
你又成為他的伯樂.
可能在他最需要你去引薦的時候.
你去講一句.
就會得到使徒們的信任.
本來不應該講得這麼詳細.
我能夠在錦繡堂.
在2004年的時候.
當時張美眉牧師剛好在.
孫德堂做堂主任.
剛剛在藉堂的一年左右.
知道我被我自己原先的教會.
因為一些原因解僱離職.
他就叫鄧全和當時的一位姑娘.
伍姑娘上來探我.
了解一下什麼情況.

$^{1001}$原來不是什麼問題.
是一些看法上大家的不相同.
其實不需要這樣.
張牧師知道我是傳道人.
就跟我說.
你進神學院吧.
我就說當時的播道神學院不收我.
因為沒有教會的推薦信.
張牧師就跟鄧傳道說.
就帶了一封口信給我.
就說你可以來我們孫德堂聚會.
如果你聚會完一年之後.
就是2005年.
你仍然有這個負擔做傳道.
我就一定會幫你寫推薦信.
因為他覺得很不公平.
現在是想讀神學.
從頭開始.
為什麼不肯幫你的同工寫信進神學院呢?.
進神學院而已.
又不是叫教會支持我.
當然是一年之後.
我在孫德堂聚會.
張牧師真的幫我寫推薦信.
我就可以在2005年進入建道神學院.
所以輾轉有這條尾巴.
原來這條債也要還的.
畢業之後.
我又來探望張牧師.
張牧師第一句就說.
不如你來錦繡堂吧.
事奉吧.
輾轉就開始在錦繡服侍.
所以其實巴拿巴對保羅這個信任.
是很重要的.
因為這個被引薦到見使徒.
是一個很重要的.
公開的印證保羅有神的托付.
呼召和差遣.
所以這裡就告訴我們.

$^{1041}$到了耶路撒冷.
門徒不敢和他結交.
其實是一個很大的挫折.
門徒都不信任你.
你做傳道.
你的生命改變.
都不信任你.
那怎麼辦呢.
所以第28節.
於是訴羅在耶路撒冷.
和門徒出入來往.
奉主的名放膽傳道.
並和說希臘話的猶太人.
來講論辯駁.
他們卻想法子要殺他.
因為那些猶太人.
是不能夠接受這個事實.
你講得多明白都沒有用.
他們就是不接受.
這個拿撒勒人釘十字架的.
說是彌賽亞.
是很大的一個關他們過不了的.
所以他們就要迫害保羅.
想殺保羅.
第30節.
弟兄們知道了.
就帶他下蓋撒尼亞.
送他回到大宋.
他自己的故鄉.
輾轉.
你可以想像到.
保羅就回到大宋.
在那裡沉潛.
等待.
所以什麼時候呢.
就是第11章的第25節.
因為安提安有一個大復興.
所以.
那個消息已經傳到耶路撒冷.
耶路撒冷的使徒.

$^{1081}$就猜險巴拿巴.
去看看為什麼會有這樣的事.
所以巴拿巴就去到27節.
當那些日子.
對不起.
應該講短一點.
第22節開始.
這風聲傳到耶路撒冷教會人的耳中.
他們就打發巴拿巴到安提安去.
他到了那裡.
看見神所賜的恩就歡喜.
勸勉眾人要立定心志.
恆久靠主.
這巴拿巴原是個好人.
滿有聖靈和信心.
於是有許多人歸服了主.
巴拿巴其實不是他的真名.
他的真名叫約瑟.
巴拿巴是一個花名.
巴拿巴其實叫勸慰之子.
巴是誰的兒子.
巴拿巴.
譬如核子巴底買就是底買的兒子.
巴就是巴拿巴.
所以我們照譯叫巴拿巴.
拿巴這個希伯來文的字根.
就是先知的那個字.
Nabi的那個字根.
所以勸慰之子就是舊約先知的角色.
所以巴拿巴其實.
人們說他是一個很能夠安慰人的人.
這個意思.
其實他是住在居比路.
就是塞浦路斯人.
他是利未人.
名字叫約瑟.
Joseph就是他的真名.
所以路加告訴我們.
因為他的出現有許多人歸服了主.
第25節.

$^{1121}$他又往大掃去找掃羅.
找著了就帶他到安提阿去.
他們終有一年和教會一同聚集.
教導了許多人.
門徒稱為基督徒.
從這個地方安提阿開始.
巴拿巴就看到很多人歸主.
就覺得自己做不完.
因為當時是安提阿.
是很近大掃.
所以就去大掃城找保羅.
保羅一直都在大掃等候.
神的時間還沒有到.
他不會到處去傳道.
所以他還不知道自己在哪裡是他的工場.
和應該肩負甚麼崗位.
一直都在大掃.
但是巴拿巴記得這個掃羅.
就去找他.
真的找到了.
和他一齊來到同宮.
在安提阿這個地方.
來教導許多人.
這些人最終更加認真成為主的門徒.
所以就叫基督徒.
跟隨基督的人.
就是這兩位神的僕人.
在安提阿所做.
安提阿由十三章開始.
成為了宣教的中心.
十三章的第一節.
在這一年裡.
跳了十二章.
因為是講彼得的經歷.
和希律的被殺.
十三章.
在安提阿的教會.
有幾位先知和教師.
所以教會裡面應該有使徒.
有先知.

$^{1161}$有傳福音的.
有牧師和教師.
是有不同的職份的.
所以不是單一的.
只是講方言的.
所以教會應該有不同的成熟的基督徒.
按照著神給的恩賜.
在教會裡面去做不同的牧羊的工作.
先知當然是將神的預言去講.
當時因為沒有聖經.
教會將真理更加堅固.
教師有系統地和教徒講解舊約等等.
當時巴拿巴和稱為彌潔的西面.
鼓勵內人路求.
與希律分封王一起長大的馬念.
在王室長大的馬念.
希律好像是一起成長的王室的人.
信主耶穌.
在安提安事奉.
和素羅.
總共有六個.
他們在事奉主禁食的時候.
領袖們一起尋求神的心意.
一起禱告禁食的時候.
聖靈說要為我分派巴拿巴和素羅.
去做我召他們做的工作.
教會安提安是他們兩兄弟建立的.
或者說更加堅固.
但現在將這兩位重要的領袖.
抽出來去做外邦的傳道.
即是說神的時間到了.
純粹是神的沉潛和安靜等候的時間到了.
現在要大派用場備神.
使用他成為外邦人的使徒.
這樣開始.
有聖靈的默許.
默示.
呼召.
而且是在眾人一群領袖祈禱的時候.
聖靈開聲.

$^{1201}$要分派巴拿巴和素羅去做這個工作.
第三節.
於是他們禁食禱告.
給巴拿巴和素羅按手.
然後派遣他們走了.
這裡有一個差遣.
很認同.
不會扣留著兩兄弟.
不讓他們走.
因為這麼幫得上忙.
很尊重聖靈的主權.
經過禱告之後.
兩人都認為自己要離開安提阿的安書區.
可能是一個很大的教會.
很多人尊重.
聚會很開心.
團聚很快樂.
但是現在神說要去外邦.
要為我的名受很多苦難.
這是對保羅說的.
在大馬士革的時候.
由阿拿尼亞說出來的.
但是現在保羅要離開安書區.
要去外邦.
有一位弟兄巴拿巴和他一起.
按手禱告.
然後他們就出去.
這就是聖靈的差遣.
很快.
對不起我要先喝口水.
雖然很快.
因為想解釋得透徹一點.
給大家知道.
大家不要介意.
如果信得不久.
給他們有些少輪廓.
今晚會給大家看回部分.
大家看回那些經文.
如果你不是太熟悉.
再看回.

$^{1241}$就會幫助你.
增強你對使徒保羅的認識.
對你日後看書信.
肯定有大幫助.
因為一個人.
不可能就這樣.
懸空了.
突然出現在這個世界.
他有他的成長.
有他的背景.
有他的文化.
有他的氛圍.
但是神是會用得著這些的.
又不能動這個東西.
看看他怎樣去.
奇奇怪怪的.
真的有點複雜.
先回頭.
有了這個可能又.
我不知道怎樣關這個.
就去那個.
再分享.
稍等.
不好意思.
是不是還沒有分享這個畫面.
可以了.
分享了.
可能在這裡放大.
這樣就行了.
OK.
不好意思.
因為有時和Zoom一起.
會有這樣的情況.
我們看羅馬書.
中心思想.
恩順之義是其中一個.
重要的中心思想.
的確講這個課題.
是講得很透徹的.
講到哪裡呢?.

$^{1281}$講到第八章.
恩順之義帶來的那個福樂.
和宇宙觀.
永生的那個眼光.
完全開啟.
對神的國度.
是有一個真知的.
曲見的形式.
是由三章的二十一節開始.
去到八章的結尾.
都是講恩順.
清義帶給一個基督徒.
整個人生的眼界的開啟.
當然也都講到有福音.
你可以用福音貫穿整本書.
但是你沒有這個那麼好.
就是保羅寫這一封信.
其實是想告訴羅馬的教會.
我雖然還沒有去到你們中間.
但是我很想快點去到你們那裡.
然後和你們有一個交往之後.
我就要將福音帶到歐洲的南面西班牙.
這個就是保羅寫這封信的原因.
我和大家讀一兩段經文.
我們為什麼能夠說這個主題.
更加貼切過剛才提到的.
那兩個大的分題.
一章第五節羅馬書.
如果大家寫得慢也不要介意.
因為有時候也要寫相當多經文.
第五節.
我們從他蒙恩受了使徒的積分.
為他的名在萬國中.
使人恩順而信服.
信而信服.
這個信心是要進度服從神.
服從這個真理.
這個信是有一個叫come into obedience.
是很特別的.
不是只信口的宣認.

$^{1321}$有十章九節有說.
口裡承認.
心裡相信.
就必定會得救.
我們又揭開第十五章.
頭尾保羅也說到他自己的心境.
對羅馬教會的內心.
十五章.
我們看看第十六節.
當然這些經文我們往後會詳細再提.
現在我想點出羅馬書的寫作動機和主題.
整卷書的中心.
就是萬民需要得到神的福音.
保羅被猜遣要將這個福音傳到萬國之中.
第十六節.
使我為外邦人作基督耶穌的僕役.
即是奴隸.
希臘文就是Dolos.
作神福音的祭司.
福音的祭司.
祭司就是獻祭.
即是福音的道種進入微信的朋友的心靈.
將祂帶到神那裡.
福音的祭司使我所獻上的外邦人.
因著聖靈成為聖潔可蒙神的月立.
獻祭就是這樣.
獻祭羊羔.
首先要帶到祭司那裡.
要驗明正身.
是否沒有殘疾.
有殘疾就不可以獻.
在黑暗時期.
撒母耳的時期.
他們的獻祭是有殘疾都可以獻.
因為不理會那麼多.
但真正的獻祭是沒有殘疾.
當祂按手後.
那些祭參被燒.
被神月立.
因為是輕香的火祭.

$^{1361}$所以保羅自己扮演著新約的祭司.
他真的獻祭.
將外邦人呈獻給神.
很高尚的對自己的描述.
這樣看自己很寶貴.
原來外邦人可以獻祭給神.
要帶他們去認識耶穌的福音美好.
而且是輕香被神月立.
所以我們帶一個信主.
得到神的月立.
我們便完成了在地上福音祭司的功能.
大家要這樣理解自己.
使祂獻上的地神月立.
所以關乎到神的事.
我們不再看了.
當然後面不可以再揭開太多.
否則我們便會變得遠了一點.
就是這樣.
這本書保羅很想寫給羅馬教會.
是一間非常有規模的教會.
何時開始的呢?.
不是保羅建立的.
第二章的《四道行傳》說.
在五旬節有很多天下萬國的猶太人.
上去守五旬節.
似乎是一個很大的節日.
有遠道在羅馬的猶太人.
都上來守五旬節.
上來守五旬節只是一天.
這麼遠上來.
如果是過年.
我可以坐三天.
在內地也有十天或七天.
上來守五旬節只有一天.
但是五旬節第二章有聖靈降臨.
有彼得的講道有三千人信主.
有羅馬的猶太人信主.
回去建立教會.
不是再建立會堂.
羅馬的教會就是這樣開始.

$^{1401}$所以現在保羅知道羅馬的教會.
何時寫這封信呢?.
我們又看看第十五章.
在二佛所的地方寫的.
我先翻開來看看.
好寫的地方是聖經的版面不同了.
所以要找也有點難.
因為用慣了舊的和合本.
版面一翻就翻到.
十六章這一段.
十六章的二十一節.
大概就是在這個地方寫這封信.
我的同宮提摩泰和我的親戚路求.
耶孫索西巴德向你們問安.
我這代筆寫信的德提在主裡向你們問安.
有人幫保羅做代筆.
保羅口述他們就寫.
那接待我也接待全教會的該由.
向你們問安.
城裡面財務官(管銀庫).
以拉督和弟兄闊土向你們問安.
留意這裡.
以拉督很多時候被我們知道.
他就在二佛所.
所以提到城裡面的財務官.
以拉督似乎是保羅在這個時候.
第三次的旅行報到.
他現在就在二佛所這個地方寫這封信.
當時還看他的那個智慧行間.
還未被捉拿.
即是說他在耶路撒冷之前.
在二佛所寫這封羅馬書.
本來他去完耶路撒冷.
交付了那些獻金之後.
他就會路過羅馬再去西班牙傳道.
大概他想這樣的行程.
當然他一上到耶路撒冷.
就已經被人捉拿.
一路輾轉第二屆去到羅馬.
在阿基帕王的面前受審.

$^{1441}$這就是大概這本書告訴我們.
保羅他寫這本書的目的.
就是很想將這個整傳的福音.
給羅馬的教會知道.
重要的不是只是講內容.
而是羅馬的教會.
都需要支持宣教這個事公.
要萬國萬民信主.
因為萬民都是在神的就坐底下.
這是三章的二十節.
萬民都需要福音.
猶太人同樣需要福音.
所以教會的存在.
就是要使萬民歸主.
這就是保羅寫這本書.
更具體的就是羅馬教會.
請你們支持我的宣教事公.
因為是神託付我成為使徒.
我要將福音傳給外邦人.
因為我是外邦人的使徒.
你們是一個有社會脈絡的教會.
更加需要扮演這個角色.
要支持我的宣教事公去到地極.
這本書我們可以將每一章拼入圖畫裡.
還有十五分鐘.
當然因為有在線和現場.
大家難以在當下發問.
因為時間不夠.
我預備了很多材料.
更好在群組裡面發問.
或者你都尷尬.
私下WhatsApp我又可以.
我收集了一些問題.
我會在下一次WhatsApp回答.
希望大家能夠教學彼此的激勵.
記得沒有任何問題是愚蠢的.
你不懂就需要問.
是不重要的.
千萬不要介意一些問題.
其他人會取笑你.

$^{1481}$千萬不要這樣理解.
問問題是一個成長的必須功課.
所以大家一起學習.
我經常說.
如果問到一些問題我不懂回答.
我會找資料.
我會了解去消化.
來解答.
你開啟了我.
原來我是不懂的.
所以我是很開心的.
所以千萬不要介意.
有問題不懂就不問就算了.
這樣就不好了.
這個我們剛才已經讀了.
保羅對自己身份的認識.
外邦人的伯役.
耶穌的伯役.
很尊貴的保羅.
能夠把自己的身份降得這麼低.
由一個猶太人拉給.
一個很有社會地位的議士.
在猶太人的身份裡面.
很尊貴.
教法師.
法利賽人.
對自己的要求是無可指責.
就是律法的義無可指責.
但是在主的眼中看.
他是一個有罪的人.
所以當主耶穌改變了他.
他很甘心樂意.
接受這個身份.
所以在保羅的書信.
他經常看自己是一個伯役.
所以如果教會想問題少.
每個人都看自己.
是主耶穌的伯役就沒有問題了.
每個人都看自己是服事主的.
還有什麼問題呢.

$^{1521}$你做完事這麼辛苦.
人家不領情.
還在笑你.
或者單打你.
其實算了.
因為你做了.
耶穌滿意就行了.
我們要看到的.
當然更好是.
人又會讚賞.
又看到耶穌天上對你說.
我的僕人.
我看到你的忠心.
但是很多時候我們就看不到.
耶穌的稱讚.
看到的是那些人的黑面.
我們就黑面.
我們就不能有保羅的心境.
他的事奉能夠堅韌.
就是因為他真的看自己是一個僕人.
看自己只需要所求於管家的.
是要他有忠心.
耶穌是這麼說的.
保羅自己一生忠誠對主就已經足夠.
有沒有回報.
或者即時的回報.
他當然期望.
但是沒有也沒有辦法.
但是不要緊.
耶穌看到了.
所以他成為福音的祭司.
這些身份他交替地說.
在不同的書卷都會這樣說.
但是其中卜役的身份.
是他每一本書差不多都說的.
所以一個使徒絕對不是坐在高位.
我初來感受我自己都不是很習慣.
有些位是領袖坐.
當然是上台方便.
但是有些教會真的有些位是長老坐的.

$^{1561}$當然是當值的長老.
或者等等.
沒有對錯的.
不過千萬不要有一個位.
我們以為那個位很尊貴.
是某個人才可以坐的.
千萬不要這樣.
為了方便.
上台容易.
事奉那裡又可以簡單.
那就沒有問題.
但是如果有某些位置.
是某個人才可以坐.
你坐下去就一定會被人罵.
那就千萬不要這樣.
唯有主耶穌才配得榮耀.
我們不過是做僕人.
大家真的要平衡這些東西.
這幾樣東西.
基督耶穌的僕役.
剛才我們說了.
耶穌基督的僕人.
就是Boros這個字的奴隸.
當時有奴隸制度.
保羅看自己是一個奴隸.
是甘心樂意成為耶穌的奴隸.
做奴隸是為了傳福音.
所以他看到自己有夢照.
從信那一刻.
就已經知道自己夢照成為使徒.
我們每個人都清楚到自己肯定.
其實那個時間要多久呢?.
每個人都不同了.
但至少你必須知道你的人生.
神擺你在什麼位置.
用你的恩賜.
你在那個崗位裡盡忠.
一定要盡快發現.
你要發現自己的恩賜.
神安排你的照明在哪裡.

$^{1601}$所以他很快就知道自己是使徒.
使徒就是.
我們揭開第一章第一節.
因為那個描述都很特別.
他看到耶穌的僕人保羅夢照為使徒.
奉派傳神的福音.
使徒什麼都不准做.
唯一一生要做的.
就是特派去傳神的福音.
使徒一定要在教會裡.
所以彼得他們每個人都要走.
每個人都受逼迫要到處去.
使徒就留守在耶路撒冷.
成為一個地標.
要殺就殺.
要捉就捉.
我們是耶穌的跟隨者.
我們是耶穌的使徒.
我們是十一使徒.
是耶穌選擇我們的.
我們要留守在耶路撒冷.
和這個福音同生同死.
所以使徒就真的只有一個使命.
就是傳福音.
如果你是開店鋪.
問你賣什麼.
我只賣牛雜.
什麼都不賣.
你整間店鋪都知道.
你的牛雜一定很好吃.
很好味.
只是一樣.
真的應該要這樣做.
所以使徒的人生存在.
就是要告訴別人.
耶穌就是救主.
是他們見過的這位耶穌.
就是人類的救主.
當然他們在教會裡面.
就會設立不同的門徒.

$^{1641}$分派不同的崗位和恩賜.
去發掘.
去發現其他的門徒.
他們的成長.
他們的恩賜.
然後就繼續將這個福音傳承.
但是使徒保羅就很清楚.
他要奉神的差派去傳福音.
接著就是福音的祭司.
將外邦人帶到神的面前.
還要獻祭.
即是那個人完全歸服主.
這樣都很好的.
很重要.
我們的責任就是.
要將每一個人帶到耶穌那裡.
他能夠越來越認識主.
他之後能夠有美好的事奉.
很多年前我在舊教會.
我自己成長的母會.
有一個小弟兄.
信主.
後期信的.
都比信主大.
但一直看他的追求.
成長.
後來現在.
甚至他現在過來幫忙.
紅恩堂教主一學.
他現在在原浸聚會的弟兄.
他四月又會再過來教.
看著一個人信主.
慢慢他的改變.
現在可以成為同工.
是很開心的一件事.
其實我們應該就是這樣.
看到後輩一路信主一路成長.
所以他有時還叫我做「銳貴哥」.
因為我們以前福音堂.
就是「哥」,「姐」這樣叫.

$^{1681}$其實差不多這麼大.
他今年退休了.
但他還叫我做「銳貴哥」.
很搞笑.
你看到這些後輩現在能夠事奉.
他預備的主很認真的.
真的看到神有不同的工作.
所以我們都是一個福音的祭司.
將人呈獻給神.
留意的是小童行傳九章.
耶穌說的.
這個人現在好像凶神惡煞.
其實他是我所揀選的器皿.
神的器皿.
保羅在提摩太后書提到.
在大戶人家有不同的器皿.
有金器,銀器,木器,瓦器.
即是不同的工用.
然後怎樣用呢?.
人若自竭,脫離悲戰的時.
必作貴重的器皿.
器皿最主要是乾淨.
其實沒有所謂金銀不銅鐵.
最主要是乾淨.
不乾淨就不能用.
我們想像得到.
保羅很體重.
叫提摩太后不要捲入時時非非的紛爭裡.
浪費自己的心力.
而是要潔淨自己.
成為潔淨的器皿.
合乎主用.
所以保羅自己也聽得很清楚.
耶穌揀選了他成為他的器皿.
所以保羅一生都保持自己對主的純貞,純潔.
事奉的專注,專心.
我們說了這麼多的引言.
我們看看保羅的一些心境.
去到第十五章.
跟羅馬教會說.

$^{1721}$除了基督.
因為上文說他是福音的祭司.
獻上外邦人可蒙月立.
接著說他自己很謙虛地.
除了基督藉我做的事情.
我什麼都不敢提.
只提他藉我的言語作為用神之歧視的能力.
使外邦人信服.
所以他一直都看自己是外邦人的使徒.
其他事情他不誇耀.
但只提到神在他生命裡.
給他的機遇,能力.
使外邦人信服主.
其實很不容易.
要帶一個人信耶穌.
我們未必有神的歧視.
但我們需要聖靈的力量.
愛的力量.
忍耐的力量.
看到他的靈魂失散的焦急.
這份力量.
我們需要有這份心.
有意識地為人禱告.
有目的地為人祈禱.
希望他被神改變.
使外邦人能夠信服.
這就是保羅的人生的寫照.
不是大鑼大鼓地做.
如果要為神.
他一點都不介意有神跡,有能力.
所以他的衣服都可以治病.
那條毒蛇咬他的手.
他一甩就甩開了.
那些土人在馬爾他島上岸的時候.
這是《使徒行傳》的28章.
那些土人的一個酋長.
那個領袖叫做布伯留.
以為他一定會死.
誰知道他不死.
每個人都拜他.

$^{1761}$保羅說你千萬不要.
我是普通人.
就跟他們講福音.
所以馬爾他島這個地方.
每年都有聖保羅的日子.
紀念保羅.
所以馬爾他原來很早就已經信主了.
可能不夠時間講後面.
我們要分下一次再講.
在羅馬書有三個很重要的祈禱.
我們先講完這個就完了.
第九節和第十節.
我在他兒子的福音上.
用心靈所事奉的神.
可以見證我怎樣不住地提到你們.
在我的禱告中常常懇求.
或許照神的旨意.
最終我能夠毫無阻礙地往你們那裡去.
又是自己的心願.
又是在禱告裡面.
來懇求神開啟他的心意.
他的旨意.
他很想探望羅馬的教會.
將他自己的一些屬靈的恩賜跟他們分享.
和分享他們的信心.
這就是保羅在禱告上的服事.
為羅馬教會去禱告.
這是第一個祈禱.
保羅在開場第一章的時候.
先講述他自己寫作的原因.
為何他很想來看望你們羅馬教會.
他很想尋求神的旨意.
經過他們當中.
跟他們分享.
然後再去其他歐洲的地方傳道.
因為耶路撒冷是彼得的墓區.
是其他使徒的墓區.
亞西亞也是.
所以保羅是要進入歐洲大陸.
向歐洲人,外邦人,希臘人來信主.

$^{1801}$來事奉.
所以你可以想像.
保羅看到的很重要的就是羅馬教會.
不如我們就停在這裡.
九時了.
再一次多謝主.
我們開始了羅馬書的研讀.
盼望今天我們看保羅的一個.
信主前後的生命改變.
都成為我們彼此的一個思想.
我們自己的信主前後.
我們自己的蒙召.
我們自己的被主.
因此的選召.
到底是否已經到位.
是否已經我們是埋班.
在神的聖靈引導下.
清楚我們自己.
要怎樣來在有限的人生.
在教會的崗位去服侍.
求神感動引導我們.
再一次給我們看見.
神你的作為.
一個誤解了舊約.
以為外邦人是不應該信主.
以為猶太人是必然可以得救.
以為拿撒勒人耶穌不是彌賽亞.
這樣的誤解成為他那種怒火.
要來追擊基督徒.
多謝主.
不是像我們人那樣想.
要將他制止.
而是神改變他.
讓他有一個新的生命.
有一個好像他自己所說的.
二章.
心靈的割禮.
不是肉身的割禮.
成為神的兒女.
多謝主.

$^{1841}$幫保羅做了心靈的手術.
使他得到永恆生命.
得到這麼尊貴的使徒的職份.
更加神使用他的學問.
他寫下的書信.
成為我們後世的人.
可以去研讀.
可以對真理有認識.
我們往後還有很多課.
看羅馬書.
求主讓我們能夠有.
這個渴慕的心靈.
去認識.
去明白.
我們這樣祈禱.
奉耶穌基督的名.
阿們.
\newpage



\section{羅馬書}
\label{sec:TdP7rHPK3XI}
\textbf{2025-03-06 周四晚查經班《羅馬書(2):分題:神的義(一)》 講員:曾裔貴牧師}
\newline
\newline
連結: \href{https://youtube.com/watch?v=TdP7rHPK3XI}{\texttt{https://youtube.com/watch?v=TdP7rHPK3XI}} ~~~~ 語音日期: 2025-03-08
\newline
\newline
\hyperref[sec:mCrcbeplxo4]{\small{< < < PREV SERMON < < <}}
~
\hyperref[sec:index_chronic]{\small{[返順時目]}}
~
\hyperref[sec:index_scriptual]{\small{[返順卷目]}}
~
\hyperref[sec:WjqHgH3CzNw]{\small{> > > NEXT SERMON > > >}}
\newline
\newline
$^{1}$很開心又再一次.
我們可以有機會仔細去研讀羅馬書.
大家都看到現在的PowerPoint.
我們都嘗試用PowerPoint.
來幫大家在記憶上方便一些.
也提一提如果你開了麥克風.
麻煩你也關一下 因為影響到錄音.
申訴申請提醒和提交文件.
Jenny也要把那個.
是的 沒錯 麻煩你.
我們有一個簡單的祈禱.
天父我們第二次開始講羅馬書.
很盼望我們真的有系統地去研讀.
能夠了解這本書的整體.
無論是中心思想.
或者是每一個的分段.
大的大綱等等.
更加讓我們在讀入那些經文的時候.
我們自己也透過保羅這一份熱心.
傳福音 我們都有感動.
還有在我們人生.
找到我們自己在神賜給我們的那個照明.
我們來成為在神眼中合用的器皿.
可以被主使用 求主幫助.
多謝主給我們有在線的空間.
使得天氣寒冷其實也沒有太大影響.
我們很盼望成為了錄音之後.
在網絡上又可以祝福不同的國家.
不同的弟兄姊妹.
我們更盼望更多的初信者.
能夠對羅馬書有掌握.
有認識 我們仰望禱告.
奉耶穌基督的名 阿們.
大家看到那個分題.
已經開始和大家組織.
因為由一章開始.
一直到四章的二十四節.
就已經說了一個很重要的主題.
就是神的義是如何運作.
神的義給哪些人.

$^{41}$這個大課題.
還有為什麼猶太人.
他們是神的舊約的子民.
立約的 亞伯拉罕的子孫.
到底神的義.
他們在保羅的時代.
他們是怎樣呢?.
是不是在神的救恩以外呢?.
所以其實神的義.
是穿梭在這本書裡面.
一到第四章就試探討這個問題.
當然第五到第八章.
就是說到一個信了主耶穌的外邦人.
他可以得到什麼樣的福氣.
和這個福氣怎樣在他的生命裡面運作.
就去到第八章.
接著第九章.
真的又要再探討一個問題.
就是神的義.
怎樣和猶太人有什麼關係呢?.
所以又去到第九到十一章.
所以你看到保羅.
一直都在說這個很重要的題目.
所以這個就是我們要看到那個分段.
看到那個總綱和那個中心思想.
當然我們上一堂講過.
這本書不只是說恩訊清義.
不只是說我們怎樣可以得到神的義.
而是一群在神裡面被清義的基督徒.
怎樣去回應大夫召傳福音給萬民.
這個才是羅馬書真正的中心思想.
所以大家一路讀.
來去掌握就會好些.
我們開始進入經文.
有一兩個勸勉.
大家都長大了成人了.
真的不要小看讀聖經這個課題.
其實嚴格來說就不應該再要勸了.
不是一個小朋友.
不是一個小學生.

$^{81}$大家真的要買一本自己.
伴隨你一生的聖經.
你寫下多少東西是你自己寫的.
當然絕對不是否定現在電話的APP.
那些是不在話下的.
我自己的手機也有很多不同版本的聖經.
但是完全沒有抵消.
我們需要有一本大字版的聖經.
在研經的過程隨便揭開.
都能夠方便你去記熟經文的位置.
和背誦經文.
這些是很重要的.
當然不是因為你不用做傳道就不需要這樣.
其實不是的.
如果我們能夠有一本聖經.
隨時放在桌面這樣去研讀.
其實是一個很大的益處.
這個就題外話.
大家彼此這樣提醒.
所以如果你在家裡方便又有桌子.
麻煩你自己準備好.
攤開它.
最好有一本筆記本.
自己記一些重點.
我盡量錄完音之後.
現在AI方面.
我就可以將講道轉成文字.
方便大家做一個小小的總結.
但是其實聽了一個半小時.
有些都不會漏了.
這是看大家的時間.
所以真的要有一本大聖經.
不要拿電話在那裡跟著我.
你來不及的.
要揭聖經.
最好你自己也揭一下.
這樣會幫助你增強你對神的話語的認識.
越來越有一種熟絡.
還有就是大家要不斷看.
那節聖經的上下文是怎樣理解.

$^{121}$這就是更重要的解經的大原則.
方便大家將來不會讀錯聖經.
最好有本記事本自己記下重點.
我們開始吧.
要背熟這兩節聖經.
羅馬書的一章五節.
保羅寫這本書的原因.
第十五章他很想羅馬的教會.
意大利的弟兄姊妹.
支持他將福音傳到西班牙.
這就是大前題.
他寫這本書.
當然他鋪排的.
是相當有系統地去講.
神的義怎樣加給不義的人.
三章的二十三節.
就是整個段落的總結.
萬民都是不義.
所以萬民都需要神的福音.
神的福音到底是什麼呢?.
就是神的義加給那些不義的人.
這就是福音.
那怎樣運作呢?.
當然就是藉著基督耶穌的流血.
使我們可以得到神的拯救.
或者說神的義加在我們不義的人身上.
這就算為義了.
這就是整個神學的運作.
我們繼續.
上星期我們不夠時間講完.
在整個文安裡面.
我們先讀讀.
一到第七節.
大家稍微聽一聽.
我讀一讀.
耶穌基督的僕人保羅奉召為使徒.
特派傳神的福音.
他的僕人是使徒.
我們的僕人只是做聖徒.
搞清楚了.

$^{161}$因為後面就馬上講了.
他是特派傳神的福音.
這福音是神從前.
直眾先知在聖經上所應許的.
論道他兒子我主耶穌基督.
按肉體說.
即是按肉身來說.
耶穌是來過這個世界.
做一個有血有肉的人.
按肉體說.
是從大胃後裔生的.
按聖善的靈說.
因從死裡復活.
以大能顯明是神的兒子.
我們要問.
為什麼一定要大胃後裔.
為什麼不是亞伯拉罕後裔.
為什麼不是摩西後裔.
摩西是阿倫支派.
利美支派.
大胃才是猶大支派.
所以這就牽連神在創世記的應許.
猶大的後裔.
和去到大胃的子孫.
今天如果時間夠.
我都引述了很多舊約.
才講到大胃後裔的應許.
所以保羅拉比.
這位仕徒.
他一講耶穌肉身.
按照肉身來說.
他是大胃後人.
即是有皇族血統.
有皇族身份.
位份.
等等.
當然更是應許.
第四節.
按聖善的靈說.
因從死裡復活.

$^{201}$以大能顯明是神的兒子.
特別提到死而復活的能力.
更加按聖靈的意思.
Spiritually speaking的意思.
這樣來理解主耶穌.
祂這個神人異性的身份.
在萬國之中.
第五節.
我們從他受了恩惠.
並使徒的積分.
第二節講了.
使徒的積分.
這裡我們可以有機會再探討一下.
積分和恩賜兩者之間.
是否有對立呢?.
是否有矛盾呢?.
為何這樣說呢?.
正正因為我出身於福音堂.
很小的時候.
初信主的時候.
教會不應該有牧師.
牧師是恩賜.
不是積分.
所以牧師是驕傲的.
不可以安立牧師.
到現在還是這樣.
你去福音堂聚會.
是沒有牧師的.
你問這位傳道人你是甚麼牧師.
他很害怕.
我們不是牧師.
我們沒有牧師.
所以我們去到這些經文.
就當作是一個趣味性.
或者是一些香港教會的冷知識.
給大家有機會了解一下.
有些人這麼認真去想這些問題.
我們稍微知道一下也好.
甚至他們都認為是根據聖經.
所以不要覺得別人是洪水萬壽.

$^{241}$但當我們了解多了.
所以我們自己都安立了做牧師.
你可以想像到.
有些事我們要更加融會貫通.
更加要明白.
有些事不是絕對.
當然我們跟從經文去看這個問題.
使徒的積分是做甚麼的呢.
在萬國之中.
叫人為他的名信服真道.
本來就沒有地球儀.
也不會有GPS.
更加不會有這樣的系統.
知道這個世界的版圖.
他們當時只能夠靠用步行.
或者是他們地中海坐船都快的.
都方便的.
馬車.
根本不會有飛機這回事.
更加在他們那個時代.
就是初期教會的年代.
地中海稱為大海.
大海其實他們很多時候會.
看為是一些鬼怪.
所以啟示錄有一個海獸.
在海中上來的獸.
你可以想像到.
對古代的人不要小看.
因為真的行船跑馬都三分險.
是很危險的.
所以保羅體的萬國.
一定沒有我們今天的概念.
我們有了精準的全球定位.
你去任何的地方.
只要你有定位系統.
你就不需要擔心迷失等等.
但是在保羅的眼中.
他已經知道整個地球.
都是造物主所創造.
總之不論是甚麼人.

$^{281}$他是甚麼民族.
都是在救恩的涵蓋範圍.
所以在萬國之中.
叫人為他的兒子的名去信服真道.
第六節.
其中也有你們這蒙召屬耶穌基督的人.
再一次提到蒙召.
絕對分得很清楚.
保羅他們是蒙召作為使徒.
我們包括意大利人.
羅馬教會.
後世的教會.
我們所有的人都是有一個蒙召的.
所以你和我信耶穌.
其實是冥冥當中神將我們選召的.
我們是蒙召作耶穌基督的兒女.
祂的人.
和我們是蒙召作聖徒.
這個在羅馬書也有出現.
第七節.
我寫信給你們在羅馬為神所愛奉召.
這裡提到奉召作聖徒.
屬於耶穌基督原來就是聖徒.
屬於神的人.
有了神的生命就是一個聖徒.
這個一定是身份地位上.
已經立刻具備.
一信耶穌基督就具備.
但要怎樣在我們人生.
不斷的自潔.
不斷的成長.
不斷對真理的認識.
和在生活裡面.
讓我們的生命能夠流露出那份聖潔.
這個不單沒有抵銷地位.
反而在我們的生活.
在有限的人生.
更加反照這個召命.
所以信耶穌應該越來越.
對自己的罪有敏銳.

$^{321}$有敏感.
信了耶穌應該生命會有改變.
因為蒙召是作聖徒.
所以這裡告訴我們.
願因平安從我們的父神.
並主耶穌基督歸於你們.
我們先讀七節.
開始說.
在文安裡面.
一到第六節.
已經提到兩次他自己的積分.
第七節打後.
一到第十五節.
就是他對羅馬教會的文安.
和對他們的認識.
其中更重要的是.
第九節為他們禱告.
我做了一些筆記.
讓大家來思想一下.
使徒.
普羅的蒙召.
作耶穌基督的使徒.
一章一節.
提到使徒的積分.
我們又來看看羅馬書的十一章.
大家可以看看螢光幕.
在線上的一定是看著螢光幕.
十一章的第十三節.
普羅自己親口承認.
他作使徒.
這個使徒原來有一個特別的群體.
就是外邦人的使徒.
即非猶太人的.
就是普羅的對象.
神沒有要求他向猶太人傳道.
當然不是說猶太人.
他就不目瞪.
不好意思你是猶太人.
拜託你找彼得.
不要找我.

$^{361}$也不會這麼奇怪.
不過他在人生大概二十多年.
由信主之後.
他主要的服事.
就是不斷向外宣教.
不斷進入歐洲大陸.
一路由東歐殺到南歐.
即是多倫多.
再去到西歐那邊.
所以你可以想像到.
普羅的一生就在地中海一帶.
當然他每次都會回耶路撒冷述職.
回耶路撒冷去交付事工等等.
似乎和彼得他們有些連繫.
有些相交.
當然絕對不是耶路撒冷.
就是總會.
他們是分堂.
一路都不獨立的分堂.
就不是這樣.
將來說到這些類似的體制.
因為到少藤傳十五章.
都有些少這樣的味道.
如果我們碰到那些經文.
我就會和大家提一提.
十一章這裡.
十三節.
我對你們外邦人說這番說話.
因我是外邦人的使徒.
所以敬重我的職份.
或者可以激動我骨肉之親發憤.
好救他們一些人.
若他們被凋棄.
天下就得與神和好.
天下是說萬國的外邦人.
他們被收納.
豈不是死而復生.
接著用了一個比喻.
其實是兩個.
就是麵粉和樹枝條等等.

$^{401}$做了一個比喻.
這裡我們就不詳細說了.
保羅在整本書裡面.
一而再提到後尾.
都提到他自己有這個職份.
他是蒙召作使徒.
即是說給我們有一份特別的.
不同的地方.
就是神呼召他.
我們剛才說過.
一章六節.
一呼召就成為聖徒.
保羅也不例外.
他也成為神的兒女.
不是因為他一生守律法.
是因為他得到神的義.
這個我們稍後一章九節開始.
我們就會說神的義.
其實應該是到十七節才說.
但是保羅有多一重的蒙召.
就是蒙召作使徒.
這個不是每個人都可以有.
大家又要認識某些東西.
所以開始要做一些powerpoint.
給大家有機會認識一下.
這位Mike Bickel.
在美國應該是北卡羅來納州.
第四波的靈因運動的領袖.
他就去宣傳他有使徒的恩賜.
和他有使徒的神秘經歷.
他就宣傳新的寓言.
和他認為有新的使徒.
所以他有新的使徒運動.
他在一九九幾年運作到現在.
在美國有一個24小時全年無休的祈禱會.
叫做International Prayer IPG CO.
任何時間你去都可以參加他們的祈禱會.
24小時不會停的.
不會年初一就會停的.
沒有這回事的.

$^{441}$這樣的運作.
不過他現在已經沒有再操作祈禱的事工了.
這位第四波靈因運動說自己是使徒的Mike Bickel.
因為我在神學院的時候.
做過他的一個功課.
關乎到靈因運動.
搜尋到他的一些事件.
他認為有新的使徒.
這就是我們今天要說的.
保羅這一班人.
是很特殊有神的呼召.
來有使徒的職份.
因為為了在耶穌基督升天之後.
需要有一班人去傳道.
所以保羅就被納入成為十二使徒裡面的其中一員.
所以當然你說是十三個還是十二個呢?.
因為馬提亞保選之後.
還是不是可以叫做使徒呢?.
所以這個就很難去界定.
但無論如何保羅是一個很獨一無二.
和彼得他們是並列的使徒.
而且是由主耶穌升天之後.
直接的呼召.
接著還有沒有使徒呢?.
又很特別的.
巴拿巴又可以叫做使徒.
所以慢慢是不是有些比較.
比較獨一無二的那種職份呢?.
這個就要看大家怎樣去理解.
遲些我們去到那些位置.
說到巴拿巴可能我又會給大家一些概念.
但無論如何.
現在第四波的靈因運動.
的確是有新使徒運動.
但在我們傳統的福音派.
我們就不接受有新的使徒.
因為使徒這個職份是獨一無二的.
在那個時代.
一來有幾個重要的原因.
聖經的正典.

$^{481}$我指的是新約的聖經正典.
根本就未寫好.
未完成的.
所以還是口述的.
神的直接啟示.
所以是需要有這些使徒.
需要有這些先知.
他們得到神的靈感.
好像舊日的先知那樣.
這樣作預言的.
以及得到那個啟示.
來寫聖經.
還有那個權威的.
所以現在聖經已經完備了.
所以我們傳統福音派.
我們不會接受有新的使徒.
當然剛才我說了.
為甚麼IPGO這個組織.
是它自己創辦的.
就將它淘汰了.
後來揭發了近年.
它有一些戀童.
和一些私德.
被趕出了那個組織.
那個組織仍然在運作.
但這個Mike就已經不行了.
我想說的是我們沒有新使徒.
那個職份已經close了.
使徒那一代的人一離開就沒有了.
開始我們說福音.
保羅很快就說.
這位奉召作使徒的保羅.
是特別被分派去傳神的福音.
所以使徒是為福音作見證的.
大家有沒有印象.
使徒行傳第一章.
保選亞提亞的時候.
其實有兩個準則.
在那個保選裡面.
既然我們讀到關於使徒.

$^{521}$大家重溫一下聖經.
使徒行傳的第一章.
所以我們多讀一點.
一章是十五節.
因為尤大死於非命自殺.
所以十二沒有了一個.
對他們來說就好像不太完整.
不知道神有什麼特別安排.
又沒有告訴他們.
將來我會找保羅.
不要搞那麼多事.
從來沒有這樣的事.
十五節.
那時有許多人聚集.
約有一百二十人.
彼得就在弟兄中間站起來.
說弟兄們.
聖靈值大位的口.
在聖經上所預言.
令人捉拿耶穌的猶太.
這話是必須應驗的.
他本來列在我們訴中.
本來他是使徒.
的確耶穌是枯照的.
不是臥底.
也不是耶穌看錯的.
並且在使徒的職任上得了一份.
這人用他作惡的功架.
買了一塊田.
以後身子扑倒或是腐倒.
土腹崩裂.
腸子都流出來.
因為吊頸死了.
繩斷了.
第十九節.
住在耶路撒冷的眾人都知道這事.
所以按著他們那裡的話.
「吸那塊田」.
起名叫「阿甲大馬」.
翻譯出來就是血田的意思.

$^{561}$因為詩篇上寫著.
彼得引聖經69篇的咒詛詩.
「願他的住處變為方牆.
無人在內居住」.
原來是說猶太的.
終極的應驗就在猶大身上.
然後.
「願別人得他的積分」.
這裡就是詩篇109篇八節.
第21節就是那個結果.
所以主耶穌在我們中間始終出入的時候.
就是他出來傳道的時候.
就是從約翰施洗起.
因為主耶穌是由阿司西約翰印證出來.
幫助耶穌施洗.
耶穌就正正式式.
經過40天的試探之後.
就正式來到在加利利一帶傳道.
開始呼召彼得,雅各,約翰,安德烈等等.
這些漁夫.
所以由耶穌始終出入開始.
直到主離開我們被接上升的日子為止.
必須從那常與我們作伴的人當中.
即是有很多人跟隨耶穌.
很多人都是跟隨耶穌.
但是始終都繼續跟隨.
沒有離開.
所剩無幾.
因為現在在那個壁壩裡面.
很多人已經走了.
變節了.
很多人就要釘死耶穌.
怎會從施洗約翰開始.
那個時候就多人跟隨.
因為還沒有事發生.
耶穌還沒有被人覺得是褻瀆神.
有好處.
有行神蹟.
講道又好聽.
每個人都跟隨.

$^{601}$但是後來就未必那麼多人跟隨.
所以現在定了一個準則.
就是由施洗約翰開始傳道.
去到主耶穌升天.
即是這三年半.
始終都與門徒一同出入.
跟隨耶穌的跟隨者.
在他們中間就要選出.
立一位與我們同作耶穌復活的見證.
即是兩個準則.
選出來是要見證耶穌的復活.
就是使徒的位份.
所以在這裡給大家了解一下.
蒙召作使徒.
其實是不合資格的.
他從來都沒有跟隨耶穌.
但現在就是為耶穌復活作見證.
羅馬書就是為了這位復活的主耶穌作見證.
但為什麼他沒有這個條件.
他又能夠得呢?.
因為是耶穌直接呼召他.
因為是耶穌在大馬士革的路上呼召他.
所以大家有沒有印象.
他在大馬士革被人捉拿的時候.
他走上耶路撒冷.
是沒有人相信他.
當然不是不相信這個人存在.
而是知道這個人本身是一個拉比.
是一個很敵黨基督教的一個人.
所以是巴拿巴接待他.
聽他的故事.
他的得救見證.
冒險將他引進給使徒彼得和約翰他們.
勉為其難他們兩位就接待保羅.
聽他講的見證.
才正正式式將他納入使徒的行列.
所以你可以想像到.
保羅原先真的不能夠成為在使徒的行列裡面.
即是說在早期還不是很透明的時候.
保羅都會給人一個懷疑.

$^{641}$他的身份等等.
這個就是使徒.
很重要的.
使徒做什麼呢?.
原來保羅很清楚告訴我們.
是特派全神的福音.
為耶穌復活作見證.
或者說見證耶穌基督是死而復活的救贖主.
所以那個時代就需要有這個積分.
就是使徒的積分.
有時候真的忍不住.
講到積分我們就耳不所書.
對不起.
我很多時候會講到某一個課題.
我就會馬上引很多聖經.
所以大家為什麼要做筆記呢?.
就是我怕我一帶遠了.
大家就回不了頭.
其實很重要的.
因為使徒既然是早期耶穌親自設立的一個積分.
那就是說在保羅他們這一班使徒裡面.
他們一直去事奉,去牧養.
保羅就有更多的啟示之後.
他就更加清楚了.
初期教會原來有些奠基的人.
這些人是那個時代才有的.
過了就未必需要有的了.
因為教會已經設立了.
最重要的是有了聖經的導統.
有了新舊約傳書已經成為了信仰的絕對標準.
就不再需要有這些得到啟示的人.
所以為什麼我們說我們不贊成.
有第四波的新使徒運動.
可以去訓練人們做使徒.
其實是很難講得通的.
所以我們就要看回.
保羅自己親自講的.
在耳不所書的四章第十節.
他一直在講原來我們每個人蒙神的恩典.
得到各樣的恩賜.

$^{681}$是因為什麼呢?.
是因為我們這一位主死而復活.
或者說他降世來到這個世界.
接著他為我們死.
然後為我們復活.
他引聖經就是詩篇的68篇.
我們可以多看一點點.
第七節.
四章的第七節.
這是耳不所書四章第七節.
如果初上的弟兄姊妹.
你揭聖經盲也不要介意.
有時間可以聽回.
或者你自己做一些筆記.
四章七節.
我們各人蒙恩.
都是照基督所量給各人的恩賜.
我們每個人都蒙自己特別的恩典.
因此這裡未必是講事奉上的那種.
反而是在救恩上.
我們每個人都有自己的神娘的份.
有些人很難信主.
有些人自小就在教會.
在基督教家庭第三代第四代信主.
所以每個人得到的份都不同.
是神娘給我們的.
保羅說.
他引舊約聖經詩篇.
第八節.
所以經常說.
他升上高天的時候.
努力了受敵.
將各樣的恩賜賞給人.
這是詩篇68篇的經文.
第九節.
保羅做了一點點的詮釋.
「既說升上,豈不是先降在地下?.
那降下的,就是遠升諸天之上.
要充滿萬有的」.
原來理解為詩篇的意思是.

$^{721}$耶穌基督.
他升到天上.
其實他先降下來.
在地上.
降到一個最低的.
最低的就去到陰間.
然後又死.
然後復活.
然後去到天上永活.
就是我們的主耶穌.
是充滿萬有的那位復活的主.
所以他得到這樣的位份之後.
留意第十一節.
直接連繫主的死而復活.
將各種不同的恩賜賞賜給人.
第十一節.
「他所賜的有使徒,有先知,有傳福音的.
有牧師和教師.
為要成全聖徒.
各盡其職.
建立基督的身體」.
這就是關鍵.
原來在保羅的理解裡.
教會能夠達至最終極的成長.
就是「同歸於一」.
認識神的兒子的「同歸於一」.
需要有不同恩賜.
不同職份的人來效力.
所以他就列舉了一些很重要的職份.
就是使徒,先知.
使徒是傳福音.
先知是講預言.
在《使徒行傳》有好幾個地方提到有先知.
甚至乎肥利.
耶路撒冷的七位執事.
其中一個肥利.
他四個女兒都是先知.
女先知來的.
講預言的.
所以那個時候是有這些職份的.

$^{761}$是有這些職份設立在教會.
和他們在職份裡面有這樣的恩賜.
可以來行使去運作的.
這裡說有傳福音.
有牧師和教師.
所以職能就是承傳聖徒.
各盡其職.
聖徒有不同的職份.
有不同的在教會裡面的服事.
使大家各按各位去服事.
這裡關乎恩賜.
羅馬書有講12章.
哥林多前書有講很多.
因為哥林多前書講這麼多.
又是有另外的問題.
就是哥林多教會爭誰的恩賜大.
誰的恩賜小.
但羅馬書是正面的.
去列舉不同的恩賜.
12章就會列舉一大堆.
那又那些恩賜的列表呢?.
如果我們去到12章.
我們又一起去比對一下.
整個保羅書信裡面講關於恩賜.
是挺好玩的.
如果你有興趣.
你自己可以先做好這些功課.
但這裡我們要問一個問題.
到底使徒先知,傳福音,牧師,教師.
是否只是恩賜.
不可以成為職份呢?.
這就是爭論點.
我初期信主當然是這樣理解.
因為長者這樣教我們.
但是再久一點.
我們不需要這樣局限.
那個只是恩賜.
不應該在教會裡面設立成為職份等等.
當然他們的爭論就是.
現在都沒有使徒了.

$^{801}$現在都沒有先知了.
為什麼你還要有傳福音.
為什麼教會要有牧師呢?.
你在教會有沒有聽過.
這個傳福音的職份呢?沒有.
做個人工作,做陪談就有.
怎會有一個職份叫做傳福音呢?.
教師也有.
有些教會.
中華基督教會就是宣教師.
是否要這樣極端呢?.
其實兩者做一個平衡.
的確沒有意義的.
使徒先知有特殊的時代性.
已經沒有了這個恩賜.
沒有了這個職份.
也不等於代表我們在教會裡面.
設立牧者,牧師,教師.
就等於是一個洪水猛獸.
不過呢.
這又是另一個我們要小心的地方.
因為有了這個職份.
而沒有了這個恩賜.
而沒有了這個職份.
包含的性質.
又變成一個虛銜的虛位.
這又是一個我們要小心的地方.
譬如有些牧者.
在教會裡面.
只有一個具名.
好像虛有其表的身份.
而達不到對弟兄姊妹的模樣.
承傳聖徒的一生的服侍.
又能夠令大家覺盡奇蹟.
有時我們要小心.
其實兩者的融合就是最好的.
既有這個職份.
同時這個職份賦予他的恩賜.
神絡亦不絕地給予恩賜.
他不斷地建立自己.

$^{841}$越來越成長,成熟.
譬如我2015年和禪木一起按木.
現在2025年.
已經十年了.
明顯十年前的按木和現在.
做了這麼多關於牧養的事.
是不同了.
所以應該朝著恩賜和職份.
應該是相輔相成.
這就是教會的祝福.
因為設立職份不是升級.
有些人說按木就是升級.
其實不應該這樣看的.
這樣理解的.
讓人發現具備了牧養的恩賜和成熟.
所以我們去理解這些經文.
我們說關於使徒,蒙召.
保羅原來在牧養裡邊.
一路事奉.
神給他越來越多啟示.
他發現原來一個教會.
不是只能夠有使徒.
而是有很多不同的職位.
職份的人來豐富神的國度.
其實最終每一個職份都是為著承傳聖徒.
很重要.
每一個教會.
所有的人不是你做你的事.
而是我做我的事.
這個山頭那個山頭不是這樣的.
是每個人得到他獨有的神給他的恩賜.
一起去承傳聖徒.
去建立基督的教會.
使眾人在基督裡邊同歸於一.
你可以想像到.
不同的恩賜是應該在一個領導下運用得很好.
這個永遠都是每一個時代的教會的挑戰.
每一個時代的教會都面對.
有時沒有恩賜 有時有很多恩賜.
不過各有各做.

$^{881}$這樣的話 整體上都不健康.
我們回頭看.
保羅看自己這個使徒的職份.
記住 蒙召要傳神國的福音.
開始來到這裡.
我有這張圖給大家看看.
保羅就是生在該撒奧古士督的時代.
當然他出來傳道之後.
奧古士督已經離開世界.
輾轉就到了尼祿王.
不過不要忘記.
奧古士督很重要的一點就是.
他打生打死.
掃平了群雄的割據.
成為了大一統的羅馬帝國.
所以稱為羅馬的聖平的200年.
就是由奧古士督開始.
所以為了要紀念這個聖平.
其實就是福音.
鑄幣.
我給大家看看這些幣值.
鑄了一個幣.
所以這個Pasvomena是拉丁文.
就是羅馬的聖平.
歷史學家都告訴我們.
由主前27年.
這個奧古士督.
一直到主後的180年.
就由奧古士督開始.
因為打贏了屋大維.
掃平了諸侯王.
將整個羅馬帝國統一.
然後修條堡路.
最重要就是一上任就鑄了幣.
紀念他的功績.
所以就告訴我們.
開始對宗教有一個寬容政策.
所以保羅都在享受這些最惠國待遇.
在享受道路網.
貨幣的統一.

$^{921}$和紀念統一的聖平.
所以可以說這是一個很敏感的字眼.
因為現在耶穌告訴人.
只要信耶穌.
你就會有好日子.
你就會有福音.
就是直接挑戰羅馬當時的政權.
要鋪排出來的盛世.
誰給你的?.
耶穌?當然不是.
當然是我奧古士督.
我所信的羅馬的眾神.
才帶來和平.
所以保羅的罪名肯定是.
他講的是另外一個福音.
《少唐傳》第17章.
直接講這件事.
第17章的第一節.
就是挑戰所謂的福音.
保羅和西拉經過暗非波尼.
阿波羅尼亞來到帖撒羅尼加.
在那裡有猶太人的會堂.
保羅照他素常的規矩進去.
路加醫生很仔細記錄這些事情.
保羅不忘本.
剛剛才說他是外邦人的使徒.
神要他去外邦人.
但是他每一去到一個地方.
他首先就去猶太人的會堂傳道.
猶太人就不需要跟他們辯論.
一連三個安息日.
安息日就是猶太人讀聖經的時間.
本著聖經.
整個新約所說的聖經.
不是我們現在所看的新約聖經.
本著聖經就是舊約聖經.
大家要有一點點冷知識.
因為在使徒的時代.
根本沒有聖經的原備出現.
第三節講解陳明.

$^{961}$基督必須受害從死裡復活.
又說我所傳語你們的這位耶穌就是基督.
對猶太人來說.
這位耶穌就是拿撒勒這個木匠的兒子.
就是我們的彌賽亞.
猶太人一聽就明白.
基督就是彌賽亞.
就是舊約所應許的拯救者.
大衛的後裔.
大衛的子孫.
就是這個意思.
第四節.
黑門中間有些人聽了勸.
服從保羅和西拉.
並有許多乾淨的希利尼人.
尊貴的婦女都不少.
但是那些不信的猶太人心裡疾到.
招聚市井匪類.
搭伙成群.
聳動合成的人.
闖進耶穌的家.
要將保羅西拉帶到百姓那裡.
找不著他們.
就把耶穌和幾個弟兄.
拉到地方官那裡去.
記住是一個政治的罪名.
喊叫說.
「那搗亂天下也到這裡來」.
「耶穌收留他們」.
「這些人違背該殺的命令」.
「說另外有一個王耶穌」.
你可以想像到是一個政治的罪名.
聳動合成的人.
喬治·萊尼加城.
我們都有去過這個地方.
相當美的地中海地方.
是亞歷山大大帝紀念他.
為他的妹妹建造的城.
而且不單止這樣.
去到奧古斯都和屋大維打仗的時候.

$^{1001}$劫殺羅尼加城的人.
就壓在奧古斯都身上.
他會贏.
所以當統一之後.
劫殺羅尼加城就有一個很大的待遇.
就是他們的軍糧是不用給政府.
有駐軍.
有羅馬的駐軍.
所以劫殺羅尼加是享有一個很特殊的.
因為壓中了奧古斯都.
而來到有這個最惟國待遇.
所以相當敏感.
所以當地人只需要列舉這個罪名.
保羅在傳另外一個王.
是耶穌.
當然扭曲了這件事.
耶穌那個做王是說將來回來.
不是現在要來招集他的會眾.
來推翻羅馬政權.
所以這樣的罪名已經夠死.
因為直接衝擊到當時地方官.
一直享有的福利和待遇.
是很嚴重的指控.
所以你可以想像到.
保羅基於這樣的情況.
他在劫殺羅尼加只是大概四個星期.
或者最多一個月至兩個月.
他就立即要離開.
都不夠時間和他們分享解釋.
他就要離開.
一來不想教會受到無謂牽連.
受到很大的迫迫.
輾轉他就去到比利亞.
後來再輾轉就去到哥林多.
先去雅典.
接著就去到哥林多.
所以劫殺羅尼加他就不是逗留很久.
你可以知道這個就是關乎到福音.
所以保羅一直說的是什麼福音呢?.
所以信耶穌我們在保羅的年代說的福音.

$^{1041}$其實是相當之帶來敏感的.
因為帝國剛剛在宣傳.
我賜給你們眾百姓的平安.
即是不用打仗.
因為已經打了.
你可以想像到還要鑄幣.
來紀念這個太平盛世.
太平盛世是羅馬皇帝給你的.
保羅說不是.
我們的平安.
我們的福音是耶穌這位死而復活.
不是經過年年的戰爭.
來攻城掠地.
是用他的生命來犧牲.
賜給世人的平安.
而且不單是普通的福音.
是說神的義.
在福音裡顯明.
所以兩碼子的平安.
所以在保羅的年代.
其實是一個很劃時代的開始.
你想想在羅馬帝國這麼多年的戰爭.
希臘大帝.
希臘的王國瓦解之後.
一分為四.
其實是年年的戰爭.
沒有停止過.
但是到了奧古斯都.
的確是一個劃時代.
里程碑的開始.
那個時候正正是主耶穌出來傳道.
所以使徒們他們說的福音.
我們今天當然很容易.
我們看福音書.
我們會明白.
但是當我們回到那個時代.
你要信耶穌.
得到這個福音.
其實就要有一個很大的信心.
一個考驗.

$^{1081}$給大家有一個了解.
福音是聖經上所應許.
保羅一而再告訴我們.
不是他的原創.
保羅告訴我們.
他自己也是一個猶太人.
是熟讀舊約聖經的一個拉比.
是一個教導律法的律法師.
他是一個在加瑪列門下受訓的.
很嚴謹在祖宗的律法為人的一個保羅.
但是他讓我們知道.
原來舊約聖經不是這麼簡單.
原來一直已經是神.
我們讀回那些經文.
第二節.
這福音是神從前直種先知在聖經上所應許.
這裡就有些講究了.
保羅看舊約的先知書.
是神的應許.
即是神在舊約興起過眾多的先知.
來成為代言人.
Prophet這個字.
原文的那比這個字.
就是代言的意思.
就是先知.
或者有時叫先見.
Seers.
或者是Prophet這個字.
又可以有其他的意思.
即是先人所知.
即是有一個預知的能力.
不是那些微卜先知.
而是神先向他有一些啟示.
再將啟示告訴民眾.
例如大衛的先知.
就是那位拿丹先知.
大家如果記得.
大衛做了一件錯事.
神先跟拿丹說.
然後拿丹去找大衛.

$^{1121}$很多時候在舊約就會有.
保羅告訴我們.
這個福音不是他發明的.
這個福音原來一直是神在舊約的時候.
就已經直中先知.
在聖經上所應許.
先知有其他.
應該這樣說.
在舊約裡面有兩類的先知.
有些是行動型的先知.
例如以利亞,以利沙.
他們沒有什麼特別的著作.
他們是行動型的.
在那個時代裡面成為一個力挽狂瀾的人.
所謂力挽狂瀾就是.
將百姓的心帶回神的真理.
應該說是帶回摩西的律法書.
那個要求的水準裡面.
這些就是行動型的先知.
Working Prophets.
但有些是Wording.
是寫書注書納說的先知.
例如以賽亞就是先知.
寫他的以賽亞書.
耶利米更加是有他的書.
所以是有中先知.
不過這裡說是在聖經上所應許.
即是說神原來在啟示的過程.
當然曾經有直接對人的面對面說話.
例如阿當 羅亞.
直接面對面.
甚至阿巴拉罕一而再見到天使.
見到天使跟他說話.
但後來就有先知的學校.
甚至有先知是注書納說的.
將那些社會現況說出來.
盡量的列舉.
還有針對好像有些講章.
針對那個時代去做一些鞭策.
做一些責備.

$^{1161}$做一些勉勵警告.
又帶進將來的一些預言和應許等等.
所以你可以想像到.
神已經不是再面對面的默示一些人.
而是透過寫作.
來讓他的真理流傳到後世.
慢慢地成為了舊約的經卷.
我們今天的舊約.
即是現在我們的和合本.
有39卷.
大家應該隨便識背.
這個聖經.
但原來嚴格來說.
希伯來的聖經.
即是舊約希伯來民的聖經.
或者在主耶穌的時代.
舊約的希伯來民.
他們的聖經的分段.
只有24卷.
不是39卷.
因為他們只有三個類別.
Torah就是摩西五經.
即是創出利民生.
第二個類別就是先知的類別.
先知書.
有前先知書和後先知書.
另外的11卷.
就是普通的著作.
應該叫做雜項的雜類.
一般的寫作.
大家留意一下.
螢光幕可以看到.
這個就是摩西五經.
Torah.
就是教訓 教導.
這個就是拉比他們定了的名稱.
接著就是這個Proverbs.
這個Levithim.
就是先知的意思.
又分開前先知和後先知.

$^{1201}$原來有《孫雅記時詩記》.
是先知書.
《撒慕爾記》沒有分上下.
《列王記》也沒有分上下.
所以前先知書就是四卷.
後先知就是四卷.
大家留意一下.
以賽亞 蔡爺 爺利.
蔡爺.
如果是中文的和合本.
大先知就是蔡爺 埃 傑 但 爾利.
但是大家留意一下.
原來希伯來文的先知書.
那個分段是比我們的和合本好.
因為耶利米埃哥就不是先知書.
是耶利米為自己的國家.
快要滅亡的那種流淚悔改.
那種無助.
是自述的一種痛苦.
所以耶利米埃哥在猶太人的拉比眼中.
不是先知書.
所以是很合理的.
以賽傑書 但 爾利書.
其實嚴格來說不是先知書.
其實是掀起文學.
是像啟示錄的型.
那些題材的書卷.
所以十二小先知書.
人家就沒有這樣叫.
是十二先知書 是一卷書.
所以前先知書 後先知書.
接著就是普通的著作.
就是雜行 詩歌書.
詩篇《箴言》 《約伯記》.
《約伯記》是詩歌書.
另外有五個小卷.
就是《雅歌》 《路得記》.
所以是耶利米埃哥.
《爾斯帖》和《傳道書》.
就是五個小卷.

$^{1241}$這五個小卷很特別的.
就是在節期會讀的.
《路得記》就好像是在駐彭捷.
就會讀這些書卷.
歷史書就三卷.
但以理書 以色列和尼西馬是一卷書.
所以所有的希伯來文的舊約聖經.
最後的卷書其實是歷代志.
不是馬拉基書.
我們的分法和猶太人是不同的.
所以歷代志很特別.
耶穌曾經講過 預言過 提過.
在歷代志記載過.
有一個人在電和電中間.
就是耶和耶大的兒子.
會控訴以色列人.
說你們流這些血.
要報應在你們的子孫身上.
在講耶和耶大的兒子.
控訴當時的約阿司殺了他的舅父.
這樣的情況.
耶穌就引述歷代志記載.
所以我們看這個著作.
看希伯來文的聖經.
就有這樣的分野.
當然嚴格來說不是什麼大的關係.
影響我們對聖經的理解.
不過希伯來文的聖經就有這樣的分法.
24卷 我們的就是39卷.
當然猶太人就不會有新約.
所以他們就沒有這個問題.
所以根本猶太人就不會有.
福音書之後的經卷.
羅馬書更加不用說.
所以如果將羅馬書變成希伯來文.
因為全部是希臘文寫的.
你可以想像到.
如果猶太人能夠讀到保羅的希伯來文羅馬書.
其實是一件很偉大的事.
如果我們能夠將新約變成希伯來文.

$^{1281}$來派給這些猶太人.
如果他們肯讀.
那個震撼性就很大.
但是猶太人就整體而言.
到今天他們都不會讀新約聖經.
是基督徒 基督教的聖經.
不是他們猶太人的聖經.
所以要記得一件事.
所有保羅所說的聖經都是舊約.
譬如提摩特後書三章十六節.
聖經都是神所默示的.
於教訓督責使人歸正.
教導人學義都是有益的.
叫屬神的人得以完全預備行各樣善事.
所說的聖經和提摩特說.
你留意提摩特.
聖經是神所默示的.
你不要隨便動搖.
一直說的都是舊約的希伯來文的三分法二十四卷的聖經.
大家稍為了解一下這些冷知識就可以了.
開始進入一個很重要的討論.
大衛的後裔 保羅一開始說到福音.
福音就是耶穌基督.
第二節 福音是神從前直眾先知.
在聖經上所應許.
再簡約就是說.
這個福音就是神在遠古的時候已經應許.
就是說從亞當開始.
神就讓這個福音越來越運作.
透過猶太人這個民族.
一直這樣應許.
最終的成就就是在耶穌基督身上.
所以福音再簡化就是耶穌基督.
這位神的兒子.
所以很快就說.
到他兒子我主耶穌基督.
按照肉體說是從大衛後裔生.
很快就已經濃縮了最重要.
主耶穌一定要是大衛的後裔.
主耶穌就是要承受神和大衛的應許.

$^{1321}$我就做了一些筆記.
三部義記很多段.
夠不夠時間揭開呢?.
我們逐一看吧.
大家一定要轉入這些經文.
沒辦法 舊約是有些沉悶.
但是也要的了.
因為保來一開始就告訴我們.
福音不是他發明的.
福音是神應許的.
福音在歷代已經透過先知.
眾先知來寫下.
神為一個應許.
一個默示將來就會應驗了.
那怎樣應驗呢?.
我們說神在歷史裡工作.
神從來不是突然出現.
你猜不到的.
是一直在歷史.
容讓歷史的發展.
本來不應該立王.
但是他們堅持要立王.
好像其他列國一樣.
齊齊齊齊有自己的制度.
當然會好一些.
誰知道一立了王之後.
神就說明了.
那些王不僅不會帶你更好.
因為你給了他這個職位.
他就會成為家族的壟斷.
更加會欺壓 會奴役.
因為你有了這個制度.
你就要供養他.
等於現在英皇的皇室.
以前的產物.
你就要繼續養著這班貴族.
皇室也是這樣.
所以神早就說了不要立王.
但是他們堅持要這樣.
神就容讓他們這樣做.

$^{1361}$容讓歸容讓.
神的救恩的永恆計劃.
在人的有限和不完全的制度下.
繼續這樣進行 繼續運作.
神從來沒有因為人的有限知識.
來取消他的救恩的運作.
這是一個很重要的發現.
今天我們面對現在這個時空.
我們面對現在這個政權.
我們面對現在這個生態.
我們不要以為神已經離開了.
神可以用我們現在的有限.
我們現在的社會制度.
來繼續運作他的教會.
我們不是脫離那個時空.
脫離我們自己的現實.
神可以在我們的現實裡.
創造他自己的救恩.
所以我們要記得.
我們不是離開現實.
但是不要變成現實的主人.
現實的奴隸.
沒辦法了 現實就是這樣.
沒得救了 沒辦法了.
我們就變成兩樣東西都不對.
所以神怎樣在舊約裡面運作呢?.
明明這些制度是神說明不好的.
撒慕爾記的下.
我們站一站這些經文.
第七章應該是相當重要的了.
第七章的十二節開始.
其實我已經精簡了.
因為這裡是很重要的.
大衛有一個很強烈的心.
他自己住在一個很漂亮的皇宮.
突然之間.
不是感懷身世.
突然間他有一個.
好像忘記了詞彙應該怎樣說.
突然間他有一個很大的感觸.

$^{1401}$我就住在一個這麼漂亮的皇宮.
神的約會在一個這麼簡單.
這麼落後.
這麼隨便的帳幕.
怎能行呢?.
所以他和先知說.
我要為神建一個聖殿居所.
這個就是大概的意思了.
第七章的第一節.
王住在自己的宮中.
耶和華使他安寢.
不被四圍的仇敵擾亂.
那時王對先知拿丹說.
看啊 我住在香柏木的宮中.
神的約會反在萬子裡.
即仍然在示羅的地方.
怎能行呢?.
拿丹就對王說.
你可以照你的心意而行.
因為耶和華與你同在.
這是他個人的感恩.
因為王這麼有神心.
自己也將心比己.
有沒有搞錯呢?.
我們現在這麼好.
為甚麼呢?.
拿丹就說.
好啊 你的心意多好啊.
神祝福你啊.
他又沒有問神.
第四節.
當夜神不是和大衛說.
神就和拿丹說.
當夜耶和華的話臨到拿丹.
你去告訴我僕人大衛.
耶和華如此說.
是當時一個很重要的神脈事的證據.
耶和華如此說.
你豈可見做奠宇給我居住呢?.
自從我領以色列人出埃及.

$^{1441}$直到今天.
我從未曾住過奠宇.
常在會幕和帳幕中行走.
凡我和以色列人所走的地方.
我何曾向以色列一個支派的事司.
就是我吩咐牧羊我民以色列的說.
你們為何不給我建造香柏木的奠宇呢?.
現在你要告訴我僕人大衛說.
萬君知耶和華如此說.
我從羊圈中將你召來.
你不再跟從羊群.
立你作我民以色列的君.
你無論往哪裡去.
我常與你同在.
剪除你的一切仇敵.
我必使你得大名.
好像世上大大有名的人一樣.
我必為我民以色列選定一個地方.
栽培他們.
使他們住自己的地方.
不再遷移.
兇惡之子也不像從前.
妖害他們.
並不像我明事司治理我民以色列的時候一樣.
因為大衛之前有450年是事司時代.
大家要對這些稍微有些印象就會好.
否則你對先知說的話.
或者神說的話.
你就會有些不知所謂的.
然後說.
我必使你安寢.
即是很穩定.
不被一切仇敵擾亂.
並且我耶和華應許你.
必為你建立家室.
你受數滿足.
與你列祖同睡.
我必使你的後裔.
接續你的位.
就是這裡.

$^{1481}$我也必堅定他的國.
他必為我的名建造奠宇.
當然是在說所羅門.
我必堅定他的國位直到永遠.
我要作他的父.
他要作我的子.
他若犯了罪.
我必用人的象責打他.
用人的鞭責罰他.
但我的慈愛仍不離開他.
像離開在你面前所廢棄的素羅一樣.
也很不公平.
素羅做了一件事.
他就絕子絕孫.
沒有神的祝福.
但大衛說.
你這樣為我建造奠宇.
你擔心我的居所.
我不僅擔心你的居所.
我還要祝福你的永遠後代.
你想想這個福氣真是大到不得了.
大衛其實當時是有感而發.
不好意思.
自己住得這麼漂亮.
也幫藥櫃裝修一下也好.
可能是這樣想的.
那個心都是出於真誠的.
但神說.
祝福我的居所.
我要祝福你的家和你的後代.
這個投資真是幾千倍.
你可以想像到大衛從來沒有想過.
他還不知道.
因為是晚上.
阿拿丹在自己家.
神跟他說.
先掛了他兩巴掌.
你為什麼這麼快答應了他呢?.
你都沒有問過我.
但不要緊.

$^{1521}$神也真是滿意大衛這個心願.
所以我們也應該沒有什麼錯.
我們自己不知道我們的後代是怎樣的.
但如果我們自己敬虔.
我們自己熱心愛主.
後代不知不覺會蒙福的.
神會祝福我們的後代.
從大衛這個事件.
原來那個回報是千倍萬倍.
你可以想像到是多麼的厲害.
這就是第一個應許.
大衛有心要為神見面.
其實他不好意思自己住在皇宮.
若貴還在帳幕.
就是這麼多.
他將這份心意跟阿拿丹說.
拿丹當晚被神教訓完之後.
拿丹就將這個默示告訴大衛.
所以大衛聽到之後.
十四節.
於是大衛進去坐在耶和華面前.
即是去到會幕.
當然他不可以去到會幕裡的致聖所.
因為他不是祭司.
主耶和華,我是誰呢?.
我的家算什麼呢?.
你竟使我到這地步呢?.
主耶和華,這在你眼中還看為少.
又應許你僕人的家置於苟遠.
主耶和華,這豈是人所常遇的事嗎?.
主耶和華,我還有何言可以對你說呢?.
因為你知道你的僕人.
你行這大事使得僕人知道.
是因你所應許的話.
也是照你的心意.
主耶和華,你本為大.
照我們耳中所聽見.
沒有可比你的.
除你以外再無神.
世上有何民能比你的民以色列呢?.

$^{1561}$你從埃及救贖他們.
作自己的子民.
又在你贖出來的民面前.
恆大而可畏的事.
就是埃及的法老.
驅逐列邦人和他們的神.
顯出你的大名.
你曾建立你的民以色列.
作你的子民.
直到永遠.
你耶和華也作了他們的神.
耶和華,你所應許僕人和僕人家的話.
求你堅定直到永遠.
因為真的感到無比的祝福.
為你永遠的後代.
都有神的慈愛.
管教裡面也有慈愛.
是不會廢棄大衛家的子孫.
所以這個心願.
換來神的永恆的祝福.
所以照你所說的而行.
願以色列的神.
對不起.
願人永遠尊你的名為大.
萬君子耶和華是治理以色列的神.
這樣你僕人大衛的家.
必在你面前建立.
萬君子耶和華以色列的神.
因你啟示你的僕人說.
我必為你建立家室.
所以僕人大膽.
向你如此祈禱.
真是一個很大的福氣.
這些祈禱.
大家可能真的少讀舊約.
這些祈禱就是蒙福的祈禱.
大衛其實他沒有想過自己.
會有這麼大的祝福.
神的祝福好像完全沒有任何容量.
可以傾倒在他和他的家室.

$^{1601}$即往後的家族.
你可以想像到.
這裡還沒有很明顯說到後人.
就是耶穌.
但是新約的作者.
例如加拉太書.
例如羅馬書這裡.
都告訴我們大衛一定有一個後人.
是神所設立的受告者.
所以還沒有明顯.
我們又再揭露.
新約的89篇.
我已經做了筆記.
大家回去看這個part point.
大家自己再溫習.
可能會再豐富一些.
對大衛的子孫.
大衛的後裔這個概念.
猶太人就會很明白.
大衛將來有一位子孫.
就是那位彌賽亞.
89篇.
第三節開始.
這篇很長的.
這是以色列人的耳探的一個訓誨詩.
我要歌唱耶和的慈愛至永遠.
我要用口將你的信實傳於萬代.
因我曾說你的慈悲必建立到永遠.
你的信實必建立在天上.
我與我所揀選的人納了約.
向我的僕人大衛起了誓.
我要建立你的後裔至永遠.
要建立你的寶座至萬代.
接著跳一跳20節.
19篇.
當時你在異象中曉喻你的聖文說.
我以把救助之力.
加在那有能者的身上.
我高舉那從民中所揀選的.
第20節.

$^{1641}$我尋得我的僕人大衛.
用我的聖膏告他.
我的手必使他堅立.
我的膀臂也必堅固他.
仇敵必不勒索他.
凶惡之子也不苦害他.
我要在他面前打碎他的敵人.
擊殺那恨他的人.
只是我的信實和我的慈愛.
要與他同在.
因我的名他的國必被高舉.
我要使他的左手伸到海上.
右手伸到河上.
接著他要稱呼我說.
你是我的父.
是我的神.
是拯救我的盆石.
一直這樣下去.
到後面不能讀太多.
35節.
我一次指著自己的聖潔氣勢.
我決不向大衛說謊.
他的後裔要傳到永遠.
他的寶座在我面前.
如日之行一般.
猶如月亮永遠堅立.
如天上確實的見證.
這樣去建立他的寶座.
他的皇位.
他的後裔.
你可以想像到.
這個就是呼之欲出的.
一直朝著歷史的發展.
才繼續確定這個應許.
這個將來的對象.
就是那一位彌賽亞耶穌.
可能我們這次不夠時間再看這些經文.
很快和大家揭一揭.
在《使徒行傳》.
在另外的人的講道裡面.

$^{1681}$如何引進大衛的子孫.
就是這位主耶穌呢?.
13章.
其實這章聖經是保羅講的.
其實是這篇信息.
保羅在居比路.
即是塞浦路斯傳完道之後.
輾轉進入土耳其一帶.
在比西底的安提阿.
然後再返回上去.
經過龐非利亞的別家.
然後上耶路撒冷.
第16節.
在安提阿.
這不是敘利亞的安提阿.
是比西底的安提阿.
聽一聽這段經文.
相當有趣.
講到大衛的子孫就是耶穌.
第16節.
保羅站起來舉手說.
以色列人和一切敬畏神的人請聽.
這以色列民的神揀選了我們祖宗.
當民居埃及的時候.
抬舉他們.
用大能的手領他們出來.
又在曠野容忍他們.
約有四十年漂流曠野.
摩西第19節.
記滅了迦南地的七族的人.
迦南七族即是當地的迦南人.
就把那地分給他們為業.
此後給他們設立事司.
約有四百五十年.
大約總共十三至十四位事司.
橫跨四百多年.
都不短的.
直到先知撒母耳的時候.
就是一個劃時代的開始.
一個小朋友聽到神的默示.

$^{1721}$一直都沒有默示.
但神的時代開始就是先知的時代.
撒母耳就是一個象徵性的人物.
由小到大都是神面前的先知.
第21節.
後來他們求一個王.
神就把汴雅敏茲派中基氏的兒子.
素羅給他們作王四十年.
既廢了素羅.
就選立大衛作他們的王.
所以一而再說到素羅.
都是說被廢的.
都幾悲慘的.
一點都不容易.
是神廢你的.
就不是你過不了probation.
過不了試用期.
是神每次都是廢了他.
既廢了素羅.
就選立大衛作他們的王.
又為他作見證.
說我尋得耶西的兒子大衛.
這是撒母耳記上的經文.
他是合我心意的人.
凡事要遵行我的旨意.
第23節.
從這人的後裔中.
神已經照著所應許.
為以色列人立了一位救主.
就是耶穌.
你看保羅一脈相承地說.
他在《小童行傳》中已經說到.
第一次傳道.
在羅馬書之前.
他已經很清楚.
將舊約我們都這樣讀的歷史.
他已經幫我們做了詮釋.
從這個人的後裔.
就立了一個就是耶穌.
所以大家要記得.

$^{1761}$保羅一而再說的福音.
他終於相信.
大衛應許的子孫.
就是耶穌.
對猶太人來說.
如果了解了.
消化了這個概念.
整個民族就應該完全降服.
在神的恩典的應許.
沒有廢棄的應許.
終於兌現在耶穌身上.
所以耶穌絕對不是猶太人的咒詛.
是神的應許的祝福.
已經臨到.
所以在哭場繼續.
這樣叩頭的拉比.
那些猶太人.
其實很可惜.
因為他們等待彌賽亞降臨.
其實彌賽亞就是耶穌.
已經降臨了.
所以這個就是.
我們看到.
如果新約真的成為希伯來文.
給猶太人看.
其實就很震撼.
盼望神使用.
因為我們錦書家.
都印了一批新約的希伯來文.
送給猶太人.
我們上一次去聖地的時候.
就是將這批聖經.
交付給以色列的聖經公會.
這樣來去印地.
希望真的有猶太人看到新約.
他們就會生命很大的改變.
今天我們就先說到這裡.
一起禱告.
再一次感恩.
天父讓我們開始.

$^{1801}$去看羅馬書.
我們的進度一定會很緩慢.
因為我們有很多.
舊約的經文.
要大家一起研讀.
很盼望主給我們.
有研讀的興趣.
那份心力.
我們一起去承載神的話語.
以致將來我們.
能夠將我們所學,所識.
所融會貫通的.
和其他以後信主的人分享.
讓神的道.
可以更加廣傳.
讓我們能夠有聖經.
有《仕途保羅》.
得到神的啟示.
寫在聖經裡面.
也流傳到今天.
我們有我們自己中國人的和合本.
來去研讀.
盼望神的話語.
去祝福更加多人.
更加盼望新生代.
學會神的話語.
以致我們能夠.
不會變成一個.
在教會但不懂聖經的.
一個群體.
我們這樣仰望禱告.
奉耶穌基督的名.
\newpage



\section{羅馬書}
\label{sec:WjqHgH3CzNw}
\textbf{2025-03-13 周四晚查經班《羅馬書(3):分題:神的義(二)》 講員:曾裔貴牧師}
\newline
\newline
連結: \href{https://youtube.com/watch?v=WjqHgH3CzNw}{\texttt{https://youtube.com/watch?v=WjqHgH3CzNw}} ~~~~ 語音日期: 2025-03-14
\newline
\newline
\hyperref[sec:TdP7rHPK3XI]{\small{< < < PREV SERMON < < <}}
~
\hyperref[sec:index_chronic]{\small{[返順時目]}}
~
\hyperref[sec:index_scriptual]{\small{[返順卷目]}}
~
\hyperref[sec:8LIvofDzP5M]{\small{> > > NEXT SERMON > > >}}
\newline
\newline
$^{1}$感恩天父讓我們可以繼續在羅馬書盡心去研讀.
很盼望我們能夠不會遺漏一些重要的真理.
以及我們能夠明白保羅整個寫作的動機.
以及如何去講解這個福音.
讓我們能夠掌握.
使得我們一方面.
我們自己的屬靈的生命.
都得到那個慰養建立.
更加讓我們能夠在整卷書.
都很想每一個教會.
每一個時代的教會.
都要向萬民.
未得之民傳福音.
去見證耶穌的救恩.
求主幫助我們.
我們去到這個二十一世紀.
我們仍然能夠有羅馬書.
流傳到今天.
都是神自己保守這些聖經的經卷.
幫助我們用心靈.
去發掘裡面的重要的部份.
使得我們能夠掌握整卷書.
那個重要的脈絡.
使得我們真的可以有效.
有能力去將福音帶給有需要的人.
我們今天晚上等候在你的面前.
也給弟兄姊妹在家可以安靜.
甚至乎可以做筆記.
使得我們的研讀.
更加有一份興趣.
和有聖靈的感動.
我們愛慕神的話語.
尤其是羅馬書這卷這麼重要的書卷.
我們這樣禱告奉耶穌基督的名.
阿們.
第三講的羅馬書.
仍然是第一章的引言.
上一個星期我們講到.
很重要保羅看自己.
很清楚第一節第五節.

$^{41}$他也清楚自己是被神呼召去做使徒.
記得福音派.
至少福音派.
我們是拒絕承認這個世界還有新的使徒.
因為使徒這個職份已經過去了.
他們是耶穌的見證人.
他們在未寫聖經之前.
當然是新約聖經.
他們是介乎見過復活的耶穌.
降世的肉身的耶穌.
和他們是為這位復活的主耶穌.
作他們的生命見證.
當然見證不是只是那個Verbal.
講述.
是他們有那個很特殊的恩賜.
來記載聖經.
寫下一些書信.
這些是很特殊性的.
福音派到今天仍然沒有改變.
就是我們不承認.
還有新的使徒運動.
新的先知的運動.
因為不需要了.
因為已經有傳遍的新舊約.
已經完備的聖經.
不需要再有新的啟示.
我們是在已經有的聖經裡面.
去得到那個.
不只是靈感.
而是得到那個很明確的.
聖經的教導.
和我們跟隨著聖經的教導.
而去傳福音等等.
所以我們在這一章裡面.
那個引言部分.
保羅就很強調他這個使徒的身份.
使徒做什麼的呢?.
特別差派去傳福音.
問什麼是福音呢?.
福音就是這位降世成為人的耶穌.

$^{81}$他的按肉身來說.
即是說.
如果他是按照著舊約聖經的預言.
他就是大衛的子孫.
即是說.
這個救恩不是今時今日才發生.
是神已經在創立世界以前.
這些就是已不所書的經文.
創立世界以前.
即是未有這個宇宙.
已經有神的救恩計劃.
你想想多麼偉大呢?.
我們人類多麼渺小呢?.
我們如果拒絕有做物主呢?.
其實我們的所謂的智慧.
其實是很有限的.
所以保羅就看到.
福音就是講耶穌基督.
這位降世為人的耶穌.
和他是由死復活的.
所以我們今天傳的.
是一位復活的耶穌.
絕對不是一位很好人.
很有禮貌.
很有學識的一個拉比.
是沒有用的.
救不了我們的.
所以我們講的福音.
就是一位曾經流血.
在十誡上犧牲.
而且沒有東西可以拘禁他.
連死亡都不能夠拘禁耶穌.
所以他能夠由死.
第三天復活.
而且他還要在他升天之前.
頒下了這個大的使命.
歷世歷代的教會.
都要向萬民傳福音.
所以保羅就在第一個引言.
應該是第七節.

$^{121}$就講到他自己的身份.
講到他自己作為使徒.
是一個呼召.
所以這就是我們上個星期.
跟大家講的.
引述了舊約的經文.
我們上個星期只是碰到.
為什麼不動?.
我們上個星期只是去到.
關於撒慕爾記下.
他想為神建造聖殿.
因為他自己已經建好了自己的皇宮.
住在皇宮裡面.
見到這個約櫃.
在那個tong 房的意思.
其實那個時候是叫做回望.
回望是在時羅的.
我就住在皇宮.
二千呎的神的約櫃.
就在那個tong 房.
那怎麼行呢?.
所以他就很想為神建殿.
這個心意.
神就閱立他了.
起初是拿丹先知說.
你有這個心當然好.
神一定閱立的.
死了.
下班回家.
神就和拿丹說.
你告訴我的僕人.
你為我建造聖殿是不行的.
不過我會為你建造你的家室.
所以這樣的投資.
發自真心對神的愛的那個服事.
會令到自己的世世代代的子孫蒙福.
所以你想想.
為你的子女祈禱.
為他奉獻給神.
世世代代祂會祝福你的下一代.

$^{161}$下幾代.
你已經不在了.
已經回天家了.
你的下幾代都因為你這個.
今日對神的愛的那個奉獻.
而得到那個一路世代的延伸.
我們揭開詩篇89篇.
大家應該看到那個PowerPoint.
應該看不到我們.
看到PowerPoint.
我們看PowerPoint的經文.
我都盡量用PowerPoint.
方便大家做一些記錄.
每次完結後.
我就會上傳到網上給大家群組.
大家可以自己再參考.
詩篇的.
這些叫做明顯.
這幾段一定是說到大衛的後裔.
就是耶穌.
和他的身份.
詩篇的89篇.
所以這些叫做君王司.
我們差不多是.
進入猶太人的思想體系裡面.
去看他們的彌賽亞是怎樣的呢?.
不過他們停在了舊約的世界.
不能夠跳出他們的成天.
彌賽亞一定是君王.
一定是榮耀.
一定是來到這個世界.
拯救我們的.
誰知道耶穌就是彌賽亞.
不過是一個木匠的兒子.
是一個受盡人的靈慾的一個普通人.
所以就不在他們的思想的期望裡面.
所以要置耶穌於死地.
也都到今天都不肯接受.
我們新約說的.
或者我們基督教說的.

$^{201}$耶穌就是基督.
89篇.
以斯拉人的詩篇.
這個以探.
我要歌唱耶和華的慈愛直到永遠.
我要用口將你的信實傳於萬代.
因我曾說.
你的慈悲必建立到永遠.
你的信實必建立在天上.
留意第三節.
我與我所揀選的人.
哪個人呢?當然是大衛.
納了約.
向我的僕人大衛起了誓.
其實是回應撒母耳記下的七章的應許.
納的永約.
所以後代的詩人都會回應那件事.
第四節.
我要建立你的後裔直到永遠.
要建立你的寶座直到萬代.
所以大衛是一個很聰明的人.
他真心為神.
想到神的聖殿.
他很想盡心盡力為神建造聖殿.
但神說你幫我建造聖殿.
不用了.
我幫你建造你的後裔.
你的家室.
所以這就是大衛.
你可以說是他的聰明.
或者是他真心對神的愛.
接著我們不看那麼多了.
我跳到第20節.
我尋得我的僕人大衛.
用我的聖告他.
告就是動詞.
我的手必使他堅立.
我的傍臂也必堅固他.
仇敵必不勒索他.
凶惡之子也不苦害他.

$^{241}$我要在他面前打碎他的敵人.
擊殺那恨他的人.
只是我的信實和我的慈愛.
要與他同在.
因我的名.
他的國必被高舉.
指到君王.
因為那時大家對聖殿的擺設.
要有少少印象.
因為猶太人他們應該很熟悉.
因為小時候就已經讀利未記.
那些藥櫃會有一個角這樣翹起.
這些就是那個舊恩的杯的角.
很多時候就這樣用.
可能以前那些喜歡這樣挑剔.
但是很多時候就會說到.
那個角代表著.
你可以說有時說的德性.
有時說的君王那方面.
第25節.
我要使他的左手伸到海上.
右手伸到河上.
他要稱呼我說.
你是我的父.
是我的神.
是拯救我的盤石.
我也要立他為長子.
為世上最高的君王.
我要為他傳流我的慈愛.
直到永遠.
我與他立的約必要堅定.
我也要使他的後裔傳到永遠.
使他的寶座如天之久.
倘若他的子孫離棄我的律法.
不照我的典章行.
背棄我的律例.
不遵守我的誡命.
我就要用杖責罰他們的過犯.
用鞭責罰他們的罪孽.
只是我必不將我的慈愛傳言收回.

$^{281}$也必不叫我的信實廢棄.
我必不背棄我的約.
也不改變我口中所出的.
我一次指著自己的聖潔起誓.
我決不向大衛說謊.
他的後裔要傳到永遠.
他的寶座在我面前.
如日之行一般.
猶如月亮永遠建立.
如天上確實的見證.
你可以想像到.
作者,詩人.
當然我們不知道當時的寫作背景.
以色列人.
一定是他們在王國的時代.
後期的了.
但是再再反映.
大衛的王朝.
大衛的後裔.
是顯赫的.
將來會有一個很榮耀.
掌管萬有的君王興起.
還有神與大衛立了這個不改變的恩典約.
而且是永約.
所以你會看到這些經文.
就是這樣反映到.
大衛的投資.
一元變成了過百倍,過千倍的回報.
神要為他的家人,家族.
為他的寶座.
一路延伸下去.
你可以想像到.
這就是大衛的子孫.
大衛的恩約.
所以又叫大衛之約.
我們揭一揭去《爾賽亞書》.
有兩段很重要的經文.
第一段是《爾賽亞書》.
《爾賽亞書》的第九和第十一章.
也是有相當的代表性.

$^{321}$這個時候應該是約坦的時候.
烏西亞約坦.
未到希西加的.
第九章.
可能不看那麼多.
看後一點.
主要是第六節.
九章的六節.
因有一嬰孩為我們而生.
有一子賜給我們.
政權必擔在他的肩頭上.
他名稱為奇妙賊士.
全能的神.
永在的父和平的君.
他的政權與平安.
必加增無窮.
政權通常追隨著暴力.
戰爭,掠奪.
但這個政權.
伴隨著的是平安的政權.
很少人平安地發展政權.
神的兒子的國度.
就是以平安去建立他的政權.
蚵烏的戰爭.
就是誰贏就掠奪他的土地.
簽和約各方面.
但這個政權.
是以平安.
麻煩請姐妹要mute.
不然在線聽到會很騷擾.
以公平公義.
使國堅定穩固.
從今直到永遠.
萬軍之耶和華的熱心.
必成就這事.
這是大衛的子孫.
這一段當然是在說一個寓言.
因為當時.
以賽亞是主前的.
七八年世界的人.

$^{361}$她歷經四代王朝.
剛剛烏西雅駕崩的時候.
她就開始出來做傳道.
做先知.
因為她進入聖殿.
見到異象.
神向她顯現.
她就呼召她.
她就做先知.
其實她在王室長大的.
以賽亞.
接著到她的兒子約坦.
阿哈斯.
接著到希西加.
根據猶太人的傳統.
以賽亞被嗎哪西.
用鋸鋸死的.
將她放在木的樹幹裡面.
挖通樹之後.
就像一個尖陽三寶一樣.
用鋸鋸死她.
這是以賽亞.
所以應該是五朝的先知.
這裡所說的嬰孩.
有可能是第八章的寓言.
以賽亞的太太會生一個兒子.
要幫她改名.
叫做馬克爾.
沙拉勒.
烏斯巴斯.
名字很長.
所以可能這個嬰孩.
是指她的太太.
生的兒子.
但你可以想像到.
這段經文.
都很肯定是指.
我們這位和平之君.
這位拿撒拉人耶穌.
這就是大衛的寶座.

$^{401}$強調要坐在大衛的寶座.
他的政權是這樣延伸的.
這就是神.
透過歷世歷代的先知.
做一個這樣的寓言.
現在我們仍然在說第三節.
按照肉身來說.
耶穌基督是大衛的後裔.
這樣的意思.
即是說保羅.
一來他是對舊約非常精通.
所以他一說到福音.
說到主耶穌.
不是空降憑空出現在這個世界.
他是按照著.
神在舊約的應許.
一直的發展.
他終於到了時候滿足.
他降世來到這個地上.
生在猶太人的地方.
所以不是偶然的.
是神的時間和旨意.
即是說.
他要在猶太人的地方出現.
即是說.
一定是大衛的家.
就是伯利行.
所以不是偶然的.
所有事情是按照著.
那個應許的應驗的.
所以大衛的寶座.
就是主耶穌要來承繼.
接著十一章.
差不多全章都是說到這位彌賽亞.
十一章.
從耶西的本.
必發一條.
從他更新的子子.
必結果實.
耶和華的靈必住在他身上.

$^{441}$耶西的本.
大衛的父親就是耶西.
必發一條.
耶和華的靈就在他身上.
就是寫到他有智慧和聰明的靈.
謀略和能力的靈.
知識和敬畏耶和華的靈.
所以我們如果追求聖靈的充滿.
就算是舊約也已經是一個啟發.
一定是說到很全面的.
對真理的渴慕.
認識知識的層面.
理性的層面.
感性的層面.
和我們整個人.
是受聖靈的感動.
引導.
馴服.
那個知識.
有些人.
你聽他說話.
完全不拉筋.
你就知道他一定不會有聖靈.
那種知識的靈.
你讀經讀到自己都亂了.
那就肯定是.
因為他的智慧有問題.
智商有問題.
但更多的人.
他如果炒在一起.
根本不懂理順.
整個舊恩的原委.
那就肯定.
他在聖靈的工作上.
一定有些混亂.
你看舊約.
已經能夠這樣界分.
這麼全面的一個人的知識.
情感.
意志.

$^{481}$都是馴服在聖靈底下.
有謀略.
有能力.
有知識.
更重要的是敬畏耶和華的靈.
他必以敬畏耶和華為樂.
行審判就不憑眼見.
這個就是最大的教會.
有時的問題就是.
行審判.
每個人都覺得自己可以行審判.
因為根本沒有人可以這樣做.
最慘的是行審判.
不憑理性.
不憑事實.
不憑真理.
憑眼見.
人家說你就是.
原來是怎樣怎樣的.
聽完就當是.
很多時候教會的很多問題紛爭.
就變成行審判.
因為每個人都說話不需要成本.
就帶來教會.
有時就是最多是非的地方.
當然其他的社團都會有.
有人的地方就會有.
但其實最需要的就是我們學習.
這一位的彌賽亞.
這一位的受高者.
這一位預言將來會出現的主.
他真的有審判的權.
但他不憑眼見.
憑事實.
憑公理.
憑真理去審判.
斷是非都不憑耳聞.
所以你看看.
眼見耳聞.
很多時候是道聽途說的是非.

$^{521}$所以你看到審判和是非.
是非黑白的是非.
不是說是非的那個是非.
卻以公義審判貧窮人.
以正直判斷世上的謙卑人.
以口中的仗擊打世界.
以嘴裡的氣殺辱惡人.
公義必當他的腰帶信實.
必當他脅下的帶子.
這樣的一位將來的掌權者.
大衛的根.
你可以想像得到.
這就是對主耶穌的稱讚.
或者他受託.
就是在這樣的王權底下.
而且是有神的靈.
一直帶領著他.
所以十一章大.
其實如果大家可以背了它也好.
因為你看看.
耶和華的靈在他身上.
是什麼靈呢?.
令他有智慧.
有謀略.
有聰明.
有知識.
更重要.
一個人有了這些.
最重要的底層.
就是敬畏神.
對神有敬畏.
你運用一切的知識謀略.
都是在神的愛神.
敬畏神底下運用.
這就是敬虔.
這是我們一生要追求的敬虔.
我們不要追求成為審判者.
我們應該追求成為一個敬畏神.
愛人的人.
所以當我們真是神安排我們.

$^{561}$給我們一點點的管理.
我們就要很小心去行使.
所以你看看.
行審判不憑眼見.
斷是非不憑耳聞.
你眼所見.
耳所聽的.
都要經過查證.
即是fact check.
那你就可以想像得到.
這是多麼重要呢.
因為人很容易聽完就當是.
我們繼續下去.
還有最後兩段.
都是相當有代表性的.
這位耶和華.
我們的公義.
這是耶利米蘇23章.
23章很長的.
大家揭讀一點.
第23章.
關乎到大衛的.
那個應許的.
我讀多一點點.
就是第一節開始.
耶利米蘇23章.
第一節.
耶和華說.
那些殘害趕散我草場之羊的牧人有禍了.
即是那些惡瀑.
耶和華以色列的神.
斥責那些牧羊他百姓的牧人.
如此說.
你們趕散我的羊群.
我有看顧他們.
我必討你們這行惡的罪.
這是耶和華說的.
我要將我羊群中所餘剩的.
從我趕他們到的各國來.
招聚出來.

$^{601}$令他們歸回本圈.
他們也必生養眾多.
我必設立照管他們的牧人.
牧養他們.
他們不再懼怕羊.
就是最怕驚嚇的.
因為牠沒有任何保護自己的能力.
羊是所有動物裡面.
雞也懂得叮你.
貓也懂得咬你.
狗更加會咬你.
但是羊又不會咬.
又不會咬你.
只會叫.
甚麼都做不到.
但是牠最怕驚嚇.
他們不再懼怕.
不再驚惶.
也不缺少一個.
這是耶和華說的.
為甚麼會這樣呢.
因為有這位好的牧人.
第五節.
耶和華說.
日子將到.
我要給大衛興起一個公義的苗裔.
他必掌王權.
這些經文要背的.
真的說到這位彌賽亞.
將來的出現.
大衛的子孫.
他必掌王權.
行事有智慧.
在地上施行公平與公義.
真的很難的.
一個有管理身份的人.
最難就是按公義公平.
在地上施行.
接著說.
在他的日子.

$^{641}$猶大必得救.
以色列也安然居住.
他的名必稱為耶和華.
我們的義.
這裡已經遙指.
保羅已經深深體會到.
神的義藉著耶穌基督.
賜給那些不義的人.
因為這位受高者.
這位大衛的子孫.
他有一個名稱.
耶和華我們的義.
可以想像到.
整個以色列.
一直背棄神.
一直面對很多外敵.
他們自己內部的問題.
但是有這位公義的主.
會掌管.
當他掌權的時候.
猶大必得救.
以色列會安然居住.
安然居住.
很多時候是說到.
那個好日子.
沒有戰爭了.
他們就在葡萄樹下.
無花果樹下.
大家很暢快地安然居住.
就是先知書對那個.
美好世界的憧憬.
就不會再驚嚇了.
相信你也想到了.
中秋節.
家家戶戶.
走在湖邊的.
過來外面的廣場.
在那裡很開心.
將燈結彩.
安然居住.

$^{681}$家家戶戶都這樣將燈結彩.
你想想疫情的時候.
整條街.
整個湖邊.
一個人都沒有.
那就不是安然居住了.
是一個大家都驚嚇.
碰完之後.
死了.
馬上要什麼.
那就不是安然居住.
所以這位主的真正將來的掌權.
就會帶來真正的永恆的平安.
所以現在以巴都還在面對那種戰爭.
所以猶太人其實很害怕.
雖然他們這麼厲害的高科技.
但其實他們都很擔心.
有時坐巴士.
就會有自殺式的襲擊.
因為人家就是用一個人的性命.
換取你們這麼多猶太人的危機.
那個懼怕.
那個心慌膽跳.
那你能想到.
一定不是靠科技.
靠人.
就可以得到真正的平安.
就是這位在地上.
他施行公平和公義的日子.
所以這個是將來.
遙指到主耶穌彌賽亞在地上的掌權.
最後37章的以西結書.
所以其實你想想.
剛才保羅都說得很清楚.
第二節已經說了.
這個福音是何時開始的?.
不是保羅發明的.
是神直著眾先知在聖經上所應許.
一章第二節羅馬書.
所以現在我們好像和保羅一樣.

$^{721}$看回舊約的經文.
其實原來已經充斥在舊約先知書裡的語言.
37章以西結書.
這些比較明顯.
你肯定應該要熟悉一下.
知道一下.
你就可以了解這些語言.
37章全章是說以色列會復活.
枯骨.
本來已經死了.
但是神的靈在他們身上.
突然間那些枯骨變成大軍.
復活.
好像一隊隊的兵一樣.
這就是以西結見到的異象.
我們看的是24,25節.
不說這麼多了.
否則我們第一章可能要說十堂也不行.
所以我們只能夠點出大衛的子孫.
給大家有個概念.
原來舊約已經對大衛的子孫有這樣的應許.
第24,25節.
我的僕人大衛必作他們的國.
就是這些復活的大軍.
眾民必歸一個牧人.
他們必順從我的典章.
謹守遵行我的律例.
他們必住在我次級僕人那國的地上.
就是你們列祖所住之地.
他們和他們的子孫.
並子孫的子孫.
都永遠住在那裡.
我的僕人大衛必作他們的王.
直到永遠.
並且我要與他們立平安的約.
不是平安紙.
平安的約.
作為永約.
我也要將他們安置在本地.
使他們的人數增多.

$^{761}$又在他們中間設立我的聖所.
直到永遠.
我的居所必在他們中間.
我要作他們的神.
他們要作我的子民.
我的聖所在以色列人中間.
直到永遠.
外邦人就必知道.
我是叫以色列成為聖的耶和華.
你就可以想而知.
將來如果這些經文.
真是在地上應驗.
就是所謂的千禧年猶太人的應驗.
主在地上的一千年.
這就是對大衛將來的君王的應許.
和他會治理整個的以色列.
特別說的地圖就是以色列的地圖.
就是亞國神彼亞國的應許之地.
所以我們將第三節講了.
我們再揭到羅馬書.
我們上個星期還沒有講完.
我讀讀第一章羅馬書.
一到第七節這一段.
就是使徒他知道自己的人生.
一信耶穌.
很特別的.
一信耶穌就立即蒙召.
成為使徒.
不是普通的門徒.
他成為特別的使徒.
好像一個特選的.
其他的都是呼召.
耶穌在加利利海呼召的門徒.
有些是打魚.
有些是做稅利.
接著一大群人.
很多的就是門徒.
一起跟耶穌上山.
他去哪裡就跟他.
但是耶穌在山上真真正正有一次.

$^{801}$知道神的時間到了.
他就徹夜禱告.
接著就在山上.
下山就選了這麼多的門徒.
當中選了12個.
就成為了使徒.
所以是有特別的使命.
但是很特別的就是.
在保羅信主的時候.
神耶穌又親自欽點他成為使徒.
現在保羅就說回.
這個呼召的原委.
他清楚自己的人生.
成為一個這樣的職份.
就是使徒.
而且是去傳耶穌基督的福音.
他就開始.
他不是拉彼上身.
而是他很有系統地告訴外邦人.
告訴羅馬教會.
福音是什麼呢?.
如果耶穌就是福音的中心.
那個核心.
那麼耶穌是誰呢?.
原來在舊約.
神已經有應許.
藉著猶太人.
藉著亞當的後裔.
一直延續.
這就是說是全人類的了.
所以在聖經上.
我們記得是舊約聖經.
保羅所說的聖經.
都是舊約聖經.
因為保羅死後.
都沒有一本聖經的書.
成為他們的正典.
因為正典是到了主後百多二百年.
甚至到三百多年.
才27卷的新約.

$^{841}$成為真正的Canon.
就是第二,第三代的教父.
來把它收緝成為了27卷.
所以只能夠是這27卷.
其他的就變成次經.
所以我們要記得.
聖經就是舊約.
希伯來文的聖經.
當然拆開來說.
耶穌的時代.
保羅的時代.
已經有七十時.
譯本的希臘文的舊約.
七十時譯本的舊約.
就不是希伯來文.
是用希臘文寫的舊約.
這裡我們要留意.
「論到他兒子我主耶穌基督.
按欲體說是從大位後裔生的.
按聖善的靈說.
因從死裡復活.
以大能顯明是神的兒子」.
我讀完那幾節.
接著我們就說復活.
「我們從他受了恩典的福氣」.
接著再第二次說.
「並使徒的積分.
在萬國之中.
叫人為他的名信服真道」.
所以一定要人.
他要歸服耶穌的聖名.
要呼求耶穌.
奉耶穌的名去決志.
他才成為神的兒女.
第十章九節.
聖職輩.
「你以口裡認耶穌為主.
心裡信神.
叫他從死裡復活就必得救」.
這裡保羅很清楚.

$^{881}$一定要呼求耶穌的名字.
所以耶穌的名字真的有能力.
所以我們如果真的要趕鬼.
我們奉耶穌的名.
記得不是帶一本紅色聖經去扔那些鬼.
是呼求耶穌的名字.
奉耶穌的名字趕這些鬼離開.
或者這些靈體.
這些靈界離開那個人.
所以耶穌的名字是非常有力的.
所以這裡告訴我們.
他是以大能顯明.
他是神的兒子.
受了恩惠使徒的積分.
萬國之中叫人為他的名信服真道.
信服.
Obedience.
所以不只是認同.
認同和Obedience是兩個層次.
你說的耶穌很好.
我們星期四的長者團契.
有一兩個長者已經回來了很多年.
你們的耶穌真的很好.
每個星期都來.
亞瑟蒙斯你說得很好聽.
但是你信吧.
我還沒有得罪.
他認同你說的話是好的.
他自己不認同.
他要come into obedience.
他不願意去信服這個真道.
服這個道.
所以信耶穌.
記得不是認同你的福音.
我認同你的教會的人很好.
我們聽聖詩.
令我們很安靜.
事實是很多人很紛亂.
工作很忙碌.
或者有時他們的生活有些混亂.

$^{921}$來到教會一坐下聽聖詩是好的.
但是還沒有去到Obedience的真道.
那可能是第一個階段.
但是他要進入那個救恩.
這裡告訴我們.
是要信服真道.
所以保羅說我要在萬國當中.
叫人為耶穌的名字.
信服真道.
一定要為耶穌的名字.
不是為誰.
這個老師最好了.
我跟他.
這個老師不回教會.
那我就跟他走.
為他的名字.
不是信服真道.
只是為了和他有好的關係.
所以其實是不好的.
我們記得我們是將人帶進救恩.
帶進耶穌的聖潔的生命裡面.
所以也有羅馬教會的份.
其中也有你們這蒙召.
屬耶穌基督的人.
所以保羅看是蒙召的.
每一個信耶穌的人.
是得到神的恩典的選召.
主義是揀選的.
所以由開始回到家爾文主義.
還是亞米紐斯主義.
將來有人有興趣.
將來我們完了羅馬書.
如果還有沒有被提及.
我就可以試一下說一下.
都挺有趣的.
到底是神的揀選.
還是人的自由意志.
其實兩者沒有什麼衝突.
因為將來我們去到第九章第十章.
第九章就是揀選.

$^{961}$第十章就是自由意志.
所以保羅其實沒有衝突的.
我們讀神學的人.
就會將這些歷史曾經發生過的事.
我們都要學習 都要了解的.
各有各的理論.
但是這裡保羅很多時候都說.
信徒能夠成為基督徒.
是蒙召成為神的兒女.
屬於耶穌基督的人.
你想想多寶貴.
我們是屬於耶穌基督的人.
第七節.
我寫信給你們在羅馬.
為神所愛奉召.
記住不是奉召作使徒.
奉召作聖徒.
聖人這個字.
因為不可以叫聖人.
因為寫中文的聖人很重.
通常是那些已經擺在神台.
孔子這些聖人.
你一信耶穌就變成聖人.
這件事很大.
令中國人很抗拒.
所以就說聖徒.
其實英文都是聖的字.
所以我們要留意.
在救恩在聖靈裡面.
我們就是聖徒.
接著就是一個的祝願.
願因平安從我們的父神.
並主耶穌基督歸於你們.
所以去到一些重要的經文.
我就希望幫大家停一停.
做得豐富一點.
有些不需要太多解釋的.
我們整段解釋完就算了.
除非大家有問題.
那你就可以在群組裡再問.

$^{1001}$一來可以快一點.
二來可以繼續掌握經文的脈絡.
所以第一章第一到第七節.
就是保羅講述.
他有這個身份.
託付他要向萬國的人.
證明耶穌.
唯有耶穌才是人類唯一的出路.
就是救主.
要come into obedience.
那個真道.
所以保羅就有這個責任.
他現在在這個時候.
在哥林多可能他已經知道.
有些時間他一定要.
pay a visit 去羅馬.
所以他就想寫一封很完整的信.
去講解神的義.
怎樣可以加給不義的人.
更重要的是.
一到第十三十四章十五章.
這三章聖經就是教會內部的倫理.
原來猶太人和外邦人.
都為近一些很小的事.
都開始有一些爭論.
例如食祭偶像之物.
例如安息人.
外邦人就說天天都可以.
但猶太人是受安息人.
所以你會見到.
保羅在內文裡面.
都開始告訴羅馬教會.
真正能夠解決這些問題的.
就是認識我們大家在神的救恩裡面.
我們就不再分彼此了.
應該大家一起.
只有一個共同的目標.
就是向萬民傳福音.
所以羅馬教會.
必須對保羅的宣教支持.

$^{1041}$這樣的意思.
所以我們去到第十二章後.
開始內部教會的倫理.
我們就會再去碰那些經文.
所以這一到七節.
就是保羅講述.
他作為使徒.
他有這個托付.
他就要將這個福音講解明白.
所以在這個時候.
讓羅馬教會知道.
一個整傳的.
完備的救恩.
到底是怎樣的運作.
所以可以令.
差不多等於是讓更加多的人.
可以成為一個傳福音的.
一個好像福音班的導師.
這樣.
這個是一到七節.
第八節去到十七節.
就是保羅不是拖大腳.
他是很想很想.
不斷祈禱去羅馬教會.
將他自己一些屬靈的恩賜.
和彼此的分享.
就同得安慰.
但是他說他都等候神的時間.
他可能覺得現在的時間到了.
他就會很快去到他們中間.
但是在智慧行間.
他自己極力的禱告.
但是同時間願意馴服神的.
那個時間.
那個安排.
兩者好像很衝突的.
你是你自己祈禱就祈禱.
是你的心.
心所想你就做.
但是保羅不是的.

$^{1081}$很平衡的.
他極力的禱告.
但是希望按照著神的旨意.
終必得平安的道路.
這個是第十節的.
在禱告之間常常的懇求的.
或者照神的旨意.
所以不是隨便這樣說的.
一個獨者能夠又很認真策劃.
但是同時間等候的.
最重要的.
不是自己喜歡做什麼就做什麼.
是神的旨意.
終能夠有平坦的道路.
往你們那裡去.
所以你想想.
保羅在智慧行間.
都流露出他個人.
在神面前的那種與神親密的.
那種肅靈的敏銳.
我想我們牧者最難的就是.
譬如我們要等候祈禱講道.
我們帶領教會的方向.
很多事情我們會姿態.
最終是不是神要我們這樣做呢?.
其實要不斷去揣摩.
最難的就是你要很微妙地.
很敏銳地.
神的時間和旨意.
這個才是最不容易的.
所以這裡就告訴我們.
保羅就有這樣的初心.
所以他在這個引言的部分.
和羅馬教會初步講述.
他為什麼要寫這封信.
我們今天還有一點時間.
處理復活的部分.
因為耶穌基督.
他按照著聖善的靈說第四節.
因從死裡復活.

$^{1121}$他來這個世界降生在馬槽.
馬尼亞這位童貞女.
生的這位主耶穌.
仍然是一個嬰兒.
一路每分每秒一路成長.
和我們人是完全一樣.
唯一不同的就是從來沒有犯罪.
這個就是耶穌了.
他整個人是沒有犯罪.
這樣成長.
不過是一個很困難的環境.
兩歲就要逃難去埃及.
希律死了就回來.
回來都不敢回到伯利恆.
就上了拿撒勒.
在拿撒勒長大.
拿撒勒是當時的駐軍.
羅馬駐軍的地方.
所以有很多外邦人出入.
所以稍微有良知的猶太人.
有自己對律法的要求的猶太人.
就不會去拿撒勒.
我們小時候經常聽電視.
聽父母說.
千萬不要去那些紅番區.
那些紅番區就是秀茂坪.
藍田那些.
就是經常斬人的葵涌那些.
你就可以想像到.
某些區你小時候一聽.
就知道那些地方是三山五嶽的.
在當時的年代.
拿撒勒就是一個外邦人經常出入的地方.
你會沾染的.
那些猶太人.
應該不是猶太人.
是那些拉比去完市集.
回家都要洗手.
要行一個潔淨禮.
為什麼?.

$^{1161}$因為他們不是怕有洛伊病毒.
是怕沾染了外邦人的污穢.
這樣的想法.
所以耶穌基督.
他是大衛的後裔所生.
但是第四節.
就好像很濃縮的.
將主耶穌的一生就說了.
他後半段的人生.
就更重要了.
就是因為他的復活.
四章有說的.
因為他的復活就拯救我們每一個人.
所以他不可以只是死.
死只是一次就擔了我們的罪.
但是他的復活.
使得我們永遠在他的生命裡.
我們得到那個生命.
那個復活.
所以有講究的.
四章本來就會再解釋這件事.
所以你看到四章已經有這麼多經文說復活.
我已經全部列舉了.
羅馬書講關於耶穌的復活的經文.
我們一定會去到那些經文.
就會詳細解釋.
這裡告訴我們.
這裡就是.
這個part1大家不用抄了.
因為我晚上就會扔給大家.
大家自己儲起來.
開一個folder裝起來.
慢慢的有空就去研讀.
復活節長假.
你就讀一下那些經文.
都希望復活節長假.
當然要上深圳.
要玩一下.
但是都很希望大家.
自己一口氣坐下.

$^{1201}$有杯飲品.
自己一口氣讀出聲.
讀整卷羅馬書.
大概要三個小時.
讀完整十六章.
要讓你的腦海裡.
不斷有經文.
跟你說話,對話.
不然就變得不熟悉.
很難將它串連在.
整個脈絡的思想.
因為這本書.
詩辯的論證.
是相當緊密的.
相當細緻緊密.
所以如果我們能夠追隨著.
保羅的思想進程.
我們就能夠.
更加掌握羅馬書的訊息.
聖善的靈說.
從死裡復活.
你可以想像到.
有哪一個人.
能夠從死裡復活呢?.
但是初期教會.
所有的使徒.
在主耶穌復活之後.
不再懼怕說這個真理.
在任何人.
任何地方處境.
都說耶穌是復活的救主.
這樣去推銷基督教.
基督教最大的賣點.
就是耶穌是復活的救主.
永活的救主.
所以如果我們今天.
教會已經被.
整個世界同化了.
人就會將我們打殘.
只需要告訴你.

$^{1241}$你不要相信復活.
你這麼傻的.
誰會復活?.
人怎麼可以復活?.
人死怎麼可以復生?.
所以記住.
在最早期的使徒.
每個人都說耶穌是復活的.
而且還要為復活的見證.
而犧牲生命.
所以你可以想像到.
基督教如果沒有復活.
那一定不是基督教.
一定不是救恩.
這裡已經引用了這麼多段的經文.
因為羅馬書.
我們都會繼續講.
所以就先不講這麼多.
到時我們揭露經文.
再進入裡面.
四章.
大家今晚看完PowerPoint.
回家自己先揭露.
六章又講救恩復活.
和我們有什麼關係呢?.
七章,八章,十章.
所以已經全部列舉了.
反而很想和大家.
去看一看哥林多前書.
菲勒比書和.
帖撒羅尼迦加前書.
暫時這三段.
其實有些我都不再引用了.
但後書又有引述.
我們揭一揭.
哥林多前書.
十五章的前書.
全章都是爭論.
為什麼哥林多人.
不相信有死人復活.

$^{1281}$你想想.
可以在書信裡.
去責備教會.
不相信復活.
你可以想像到當時.
使徒們傳道.
傳一位復活的耶穌.
所以我們真的要記住.
我們傳道,我們聚會.
我們被人見到的教會.
應該是一位復活的.
耶穌的教會.
不是死氣沉沉.
每個人回來都頭暈暈.
聚會完之後大家都頭暈暈.
那你的復活的耶穌去哪了?.
你們好像沒有經歷過.
耶穌的復活.
所以我們真的有時要反思.
教會有時聚會.
形態久了.
習以為常.
慢慢我們都忽略了.
其實基督教最大的賣點就是.
一位復活的主和我們一起.
十五章.
第一節到第三節.
接著下面那些我不讀了.
因為不夠時間了.
因為今天要開始講神的異議.
弟兄們.
我如今把先前所傳給你們的福音.
告訴你們知道.
當時.
當時保羅帶.
哥林多教會信主.
在四道行傳的第十六章.
第十六章.
對不起是十八章.
不是十六章.

$^{1321}$十六章是馬其頓二章.
十八章去到哥林多就帶了他們信主.
本來受迫害.
他想要走的.
但是晚上耶穌在二章向保羅顯現.
告訴他不要怕.
只管說有我與你同在.
必定沒有人下手害你.
然後就叫他.
繼續這樣去停留.
保羅就在那裡住了一年半.
在哥林多教會.
現在他聽聞.
哥林多人.
不知道受了什麼影響.
就不信有死人會復活.
弟兄們我如今把先前所傳給你們的福音.
告訴你們知道.
這福音你們也領受了.
又靠著站立得穩.
並且你們若不是信主.
並且你們若不是徒然相信.
能以持守.
我所傳給你們的.
就必因這福音得救.
所以要持守.
要hold fast.
舊的英文.
要拿著.
持定著.
持定著什麼是福音呢.
接著就保羅說.
第三節.
我當日所領受.
又傳給你們的.
第一.
就是基督照聖經所說.
有沒有留意.
又是舊約聖經的應許.
為我們的罪死了.

$^{1361}$而且埋葬了.
中文的翻譯其實挺好的.
好像一個past tense.
一個完成了的.
為我們的罪死了.
而且埋葬在墳墓.
在一個婦人的墳墓.
在亞利馬太的.
藥室.
這是.
《耳塞書》53章的.
預言.
53章第十節.
死的時候.
與財主同葬.
《耳塞書》這樣說.
這裡告訴我們.
而且埋葬了.
又照聖經.
又照聖經所說.
原來所有東西都有.
好像我們認為有一個劇本.
其實不是劇本.
是神的預言.
和神的定義.
要耶穌這樣死.
又照聖經所說.
第三天復活了.
並且顯給基法漢.
然後顯給.
十二使徒漢.
後來一時顯給.
五百多弟兄漢.
其中一大半到如今還在.
卻也有已經垂了的.
也就是一段很長的時間.
那些為耶穌.
作見證的弟兄姊妹.
五百多人.
還有一大半還沒死.

$^{1401}$還在民間.
作見證.
為耶穌的福音作見證.
在羅馬帝國作見證.
即是說你預計.
一定會受逼迫.
當時寫這書的時候.
就是尼祿王的時候.
所以那個死.
似乎幾多都不是.
壽終正寢.
是受逼迫而死.
但是他們仍然.
為這個見證.
付代價.
第七節以後顯給.
雅各漢.
這個就是耶穌肉身的弟弟.
因為使徒雅各.
在《使徒行傳》十二章.
已經殉道了.
被希律王殺了.
所以這個雅各是.
耶穌肉身的弟弟.
然後再顯給眾使徒漢.
第八節留意.
沒了也顯給我漢.
即是說.
大馬士革的事件.
保羅去大馬士革的時候.
就遇見大光.
就是耶穌.
向他講說話.
然後呼召他入城.
然後就是亞拿寧.
向他按手.
所以沒了也顯給我漢.
保羅看自己.
我如同未到產期.
而生的人一般.

$^{1441}$七星仔.
我原是使徒中最少.
不配稱為使徒.
因為我從前迫伯臣的教會.
不是要謙卑.
而是講一個實在的真話.
即是認識自己.
一個認識自己.
是多麼蒙恩的人.
就很賣力.
用心去事奉.
所以這裡是要.
吃力不討好的.
因為你要講明真理.
你就會得罪哥林多有些人.
所以一個使徒.
你追求的是神的榮耀.
真理的彰顯.
因為你怕人不敢講真話.
不敢講真話的人.
容易到不得了.
沒人說你啞的.
你不出聲.
但出聲講真話.
才是真的先知.
舊約經常都是這樣.
所以保羅說.
我好像一個未到產期而生的人.
但他也要為這個真理作見證.
其實是因為他自己.
從前迫伯臣的教會.
所以然而.
我今天成了何等的人.
其實是蒙神的恩才成.
不是自己的努力.
不是自己的能幹.
智商被人強.
有很多的條件.
優越.
就不是這樣.

$^{1481}$深深看到的是神的恩.
才成為今天這樣的人.
而且成為使徒.
絕對不是一個.
很榮耀.
到處都被人稱讚.
到處都會招致猶太人的反擊.
羅馬教會.
羅馬皇帝政權的迫伯.
所以其實是內外受敵.
不過既然福音這麼寶貴.
他就很願意去做.
所以這是蒙神的恩典才成.
並且他所賜我的恩.
不是徒然的.
我被眾使徒格外勞苦.
真是一個真正被恩典感動的人.
所以教會需要的是一個一個.
被神恩典感動.
他就很願意為主而活.
為主而做事.
還需要叫他迫他.
他根本很想參與事奉.
這原不是我.
不是我特別被人勒.
那是神的恩.
第二次第三次講恩典.
與我同在.
不單止是我 眾使徒.
我們這樣傳 你們都是這樣信.
其實你們以前是這樣信的.
為甚麼現在你們不是這樣信.
開始辯論了.
這就是十二節.
所以我今天很想給大家看到一個roadmap.
保羅的書信.
差不多捲捲都講主耶穌的復活.
再闊些就是.
所有在生的使徒.
他們寫的書信.

$^{1521}$他們的傳道都試探討復活.
例如彼得前書.
三章.
這個時候保羅已經死了很久.
三章應該是十八節.
這個時候就開始有組織的迫害教會.
彼得就叫他們學校耶穌.
要用行義.
不要行惡.
來面對那些苦難.
我們看看第十三節.
因為上文那裡我們不引.
因為他引三十四篇的詩篇.
第十三節.
然後一直到十八節.
你們若是熱心行善.
有誰害你們呢?.
你們就是為義受苦.
也是有福的.
開始木羊頂子妹在苦難的時代.
教會應該怎樣自處呢?.
唯一的就是在迫害我們的人的面前行善.
也是有福的.
不要怕人的威嚇.
也不要驚慌.
專注基督為聖.
所以你們裡面對神有敬畏.
是最大的.
你的膽量也會最大.
你面對的事情.
那些難就不是那麼難.
因為你裡面有神.
你專注基督為聖.
有人問你們心中盼望的緣由.
就要常作準備.
以溫柔敬畏的心回答各人.
不是以惡報惡.
以牙還牙.
你再惡過他.
不是這樣.

$^{1561}$不是用你的理去駁斥他.
而是以一個溫柔敬畏的心靈去回答.
存在一個無愧的良心.
叫你們在何事上被毀旁.
即是開始白熱化.
教會受到有組織的針對.
就在何事上.
可以叫那些誣賴你們的人.
在基督裡面有很憤恨的人.
自覺羞愧.
所以是以德,品德,德行去服人.
讓人看到你們這些基督徒.
所謂何事.
為何你們會躲起來守聖餐呢?.
是否吃人肉呢?.
還是有甚麼秘密組織呢?.
是否有甚麼秘密組織呢?.
讓他們知道其實不是的.
我們是紀念耶穌十日的犧牲.
我們在大庭廣眾.
在羅馬的生活.
譬如他們會敬拜君王.
你們不做.
別人就一定會逼迫你.
但是你要讓他們看到.
我們是敬拜一位獨一真神.
不是故意要和社會或君王為敵.
然後告訴我們.
第十七節.
神的旨意.
即是如果迫害是別人刻意製造.
你已經進入了那個時代.
這個你要看為是神的旨意.
不是一個你可以逃避的警方.
神的旨意若是叫你們因行而受苦.
總強如因行惡受苦.
然後就開始引述基督.
因基督也曾一次為罪受苦.
一次就是十字架上的一次.
但他一生背負著被人羞辱的罵名.

$^{1601}$所以一次和一生我們都要記住.
所以耶穌有個名字.
《二乘書》53章第三還是第五節.
他稱為「憂患之子」.
Man of sorrows.
他的人生充滿憂患.
你想想多麼不容易.
他的人生就是這樣.
但主耶穌更加嚴重的是一次為罪受苦.
就是「疑的代替不疑的」.
為要引我們到神面前.
按照肉體來說.
他和保羅所說的完全一樣.
他被致死.
按照靈性說.
他復活了.
接著是19節.
這段是相當難解的.
是新約的難解經文.
他值者靈.
他就是耶穌.
曾經去傳道給那些在監獄裡的靈聽.
所以那個煉獄觀.
其中一段就是這裡.
天主教的煉獄.
就是這樣.
但這段其實不應該解為煉獄.
第20節.
就是那從前在羅亞預備方舟.
神容忍等待的時候.
不順從的人.
在監獄裡的靈.
其實就是從前.
在羅亞的時代.
那些不順從的人.
神容忍等待.
他們進方舟.
但他們是笑羅亞.
怎麼會有洪水.
怎麼會造這麼大的船.

$^{1641}$我們要進去.
在裡面就不准出來.
是神容忍等待.
羅亞造了多久.
就笑了這麼久.
因為從來都沒有大淵的泉源打開.
就有洪水來.
所以我們要留意.
第19節.
耶穌.
藉著靈.
曾經傳遞給那些在監獄裡.
即是說現在在監獄裡的靈聽.
那些監獄裡的靈.
就是從前羅亞的時代.
沒有進入方舟的人.
接著就是下面的20字的中間.
當時進入方舟.
有多少人呢?.
8個.
因為羅亞一家.
和他三個兒子.
閃,含,瓦忽.
和他三個媳婦.
總共8個人.
當然還有很多動物.
這裡說藉著水得救的人不多.
只有8個人.
沒錯 只有8個.
數字給得都對了.
但這裡提到第21節.
藉著水得救.
說明是藉著進入方舟得救.
怎會藉著水得救呢?.
這麼奇怪的阿彼得論點.
接著繼續解釋那個水.
他說「藉水所表明的水禮.
現在藉著耶穌基督的復活.
也拯救你們.
而這洗禮不在乎除掉肉體的污穢.

$^{1681}$只求在神的面前有無愧的良心」.
接著22節說一個現在在彼得的時代肯定的.
他說「耶穌已經進入天堂了.
在神的右邊.
眾天使和有權柄的.
並有能力的都服從了耶穌」.
這裡其實想說什麼呢?.
其實想說第17節.
神的旨意如果叫我們.
因行善受苦.
總強如因行惡受苦.
基督就是那個行善受苦的榜樣.
祂現在在哪裡呢?.
耶穌已經進入天堂了.
在一切的權力的中心裡.
祂已經掌權了.
但祂的一生其實是憂患之子.
所以我們在人生.
我們面對憂患和苦難.
我們是不知道那個結局.
我們不知道什麼時候有出路.
什麼時候完結.
不用看我們.
彼得教我們看耶穌.
耶穌就是這樣走完祂的一生.
現在在天堂了.
當然中間那段.
因為今天不想拆解.
就先不解釋了.
免得再拆解.
所以那個水其實不是說洪水.
那個水是在說審判.
那個是耶穌基督好像被大水淹沒.
受了那個刑罰.
受了我們人類面對神的那個咒詛.
那個憤怒.
所以這個是一個審判的水.
所以我們是藉著這個水得救.
彼得的意思就不是說.
其實他們沒有進過方舟.

$^{1721}$不是.
進方舟才可以得救.
但現在這裡是象徵.
借羅亞的這個洪水.
所以是主耶穌替我們在十誡上.
忍受那些咒詛,痛苦.
所以祂是受了那個審判.
我們就藉著耶穌.
所以耶穌就是我們的方舟了.
我們就在耶穌的生命裡面.
我們得救了.
所以彼得的意思是這樣的意思.
而且受苦有益的原因.
耶穌已經得到神的抬舉.
升高成為在永恆宇宙萬物當中的掌權者.
而且還是由死復活.
將來可能我們見到的是.
戴著釘痕的手的耶穌.
這樣的掌皇權.
這樣受苦的主耶穌.
所以其實那個重點.
不要想那麼複雜.
第十七,十八節.
就是神的旨意.
如果要我們在人生中有苦難.
你就要學習.
好像耶穌那樣.
經驗那個困難.
有神同在的.
有主安排的.
我們要認定.
總是有神的計劃的.
不要那麼快只是用我們自己的水平.
我們自己有限的眼光.
去看我們自己現在的處境.
你的處境好的當然開心.
但很多時候正正就是我們有時不開心.
不好的處境.
我們就不知道怎麼辦.
這段聖經已經很清楚告訴我們.

$^{1761}$我們可以學習耶穌.
我要忘記喝一口水.
所以這個水是在說耶穌替我們受審判.
所以嚴格來說.
這段經文耶穌就是那個方舟.
我們其實在耶穌裡面.
避開了神的就坐審判.
所以彼得也借用這個羅亞事件.
一個靈異的解經.
但整個思路是說我們要行善.
不要行惡.
如果神的旨意真的要我們行義受苦.
我們都要接受.
國共內戰的時候.
很多基督徒祈禱共產黨快要倒.
但不是的.
輸了很多基督徒去了台灣等等.
那些戰爭.
這是歷史.
我們不是要抵毀甚麼政權.
而是告訴我們神的旨意.
如果要這樣安排.
都要進入社會的時代.
我們永遠只有一個的向度.
就是一定要行義.
如果有苦難的我們都要面對.
因為當時尼祿歷史說.
他自己不喜歡舊的羅馬城.
他就洗還手.
找一些黑社會燒了舊的羅馬城.
給元老院對他有一個問責.
要他交代.
他就賴基督徒做了這些反抗的事.
所以就嫁禍給基督徒.
就變成大規模合法地去捉拿基督徒.
所以在那個時代你沒有選擇.
彼得就要被刀釘十字架.
這裡就是彼得前書這一段.
我說連使徒彼得在他的事奉裡面.
在這個苦難的對教會的模樣.

$^{1801}$他都提到我們一定要看到主耶穌是復活.
而且現在在天上掌握真正的永恆的王權.
一切有能力的都要服從他.
但在他的人生.
其實他是一個好像寂寂無聞.
到處都被人嘲笑的一個木匠的兒子.
但現在就在這樣的永恆裡面.
所以我們要有這種好像一個平衡世界.
我們實實在在活在香港.
我們實實在在需要每天要吃東西.
要工作 有家庭.
但我們要記住我們有一個永恆的天堂.
是同時間發生的.
所以如果基督徒沒有了對永恆的盼望.
那就很慘了.
你又要每個星期回教會.
那些人不用回教會.
睡到十一點.
其實我也很想的.
睡到十一點.
接著去喝茶.
如果熟的指一下鼻子.
可以進去的.
不用排隊的.
其實很爽的.
這樣的人生.
接著看一場戲.
接著下午去哪裡.
不如去遊船河.
其實很開心的.
一個不信耶穌的普通人.
但他就只能夠是這樣.
但如果我們看到有永恆.
有天堂.
就很不同了.
耶穌已經進入天堂.
這是22節.
很想告訴在苦難中的基督徒.
要看到有永恆.
彼得的意思就是這樣.

$^{1841}$有兩段反而重要.
就是《肥勒比書》三章.
我們可能看完這幾段.
我們就要試一試.
我們用一點時間.
講神的「二」的段落.
《肥勒比書》的第三章.
這些是保羅書信.
所以我也不引其他使徒的書信.
譬如約翰的書信等等.
剛才彼得那裡和苦難有關.
我就稍為引一引.
三章第十節開始.
或者這一段.
保羅對自己的自我的辯解.
講他自己的身份.
我們可能從第五節讀起.
經文稍為長一點.
對他所說的由死復活有些幫助.
第五節.
我第八天受割禮.
真正是一個很敬虔的猶太家庭.
才能夠這樣做.
第八天就受割禮.
不是每個人都做到.
在外邦那些其實做不到.
所以他們在猶太人的眼中.
在拉比的眼中.
他們就不夠虔誠.
保羅在第八天也沒有選擇.
是他的父母幫他行割禮.
你可以想像到他真是一個敬虔的家庭.
我是以色列族.
卞雅敏茲派的人.
原來保羅是卞雅敏茲派的人.
卞雅敏茲派最出色的就是有君王.
這個索羅.
都是一個大族.
是希伯來人所生的希伯來人.
有些不是.

$^{1881}$有些是父親是外邦人.
母親是猶太人也有.
比如提摩泰就是.
他的外祖母就是羅爾.
他的母親就是有利基.
是猶太人.
但是他父親是希臘人.
所以保羅就幫提摩泰行割禮.
即是說他原來小時候都沒有行割禮.
因為保羅要帶提摩泰去猶太人的地方傳福音.
你這個半堂番都沒有割禮.
你現在教我們?.
我們才是猶太人.
所以保羅都很懂得禮儀.
很知道那些人的心結.
所以強制幫提摩泰行割禮.
那個時候他已經是青年人.
因為他要做傳道.
神有揀選他.
所以保羅就知道.
你是一半的猶太血統.
所以你都要行割禮.
反過來提摩泰完全是外邦人.
完全不是竹昇仔.
完全是外邦人.
所以保羅是沒有幫提摩泰行割禮.
所以你看到保羅就能夠拿捏到這些律法.
他也很明白人性.
對猶太人他也會有一種遷就.
我們看下面.
這裡說所生的希伯來人.
就律法來說我是法利賽人.
即是嚴守律法.
差不多去到一個律法主義.
保羅按照著明文的律法.
似乎猶太人就說.
將它濃縮到變成613條.
即是說他做足了.
就熱心來說.
我是迫不得已教會.

$^{1921}$既然他嚴守律法.
他見到有些猶太人歸信了基督教.
不再信猶太教.
他就按他的熱心.
他就有責任去抓這些人.
要將他們定罪.
所以他迫不得已教會.
就律法上的意義來說.
我是無可指責.
即是說按照著表面的律法.
那個要求.
所以這個到第七章.
保羅就說到律法和心靈的關係.
所以他真的是過來人.
就可以說到.
他按照著律法表面的意思.
他是無可指責的.
即是說我做足給你看.
但是這裡告訴我們.
第七節.
只是我先前以為與我有益的.
我現在因基督都當作有損.
一個失敗者.
不但如此.
我也將萬事當作有損.
因我以認識我主基督耶穌為至寶.
我為他已經刁棄萬事.
看作糞土.
為要得著基督.
我不知道我是否初信的時候讀到這段經文.
我扔掉了那些LV.
給了我姐姐.
姐姐的眼睛都發光了.
你不要的LV是真的LV.
不是A貨那些.
她開心得跳起來.
可能我讀錯了聖經.
保羅這裡就不是我們這麼決絕.
好像二分法那樣.
只是屬靈不要世界.

$^{1961}$當然好像開玩笑.
其實保羅的意思.
其實他看到什麼才是真正永恆的價值.
所以他不是扔掉猶太人的書.
那些律法書.
不再讓自己什麼記憶都沒有了.
完全不是.
他有了新約有了聖靈.
他能夠運用舊約他認識的背景.
熟讀的背熟的經文.
那些每一卷書的脈絡.
他能夠將它化作新約的教導.
所以不是這樣丟棄萬事.
我們是在心智上改變我們的價值觀.
所以這個有損是我們和基督做一個比較.
他將萬事當作有損.
重點是因為認識我主基督為至寶.
所以這裡是真的很重要.
人生如果我們什麼都有.
各樣都好.
這是神給人的福氣.
但是他看不到耶穌基督是他的至寶.
沒有揀選沒有相信.
沒有去承認基督是他的救主.
一切都是徒然的.
一切都是.
雖然表面上他的人生很好.
很順暢.
但是給一個在主耶穌的救恩裡面.
而人生未必很順利的人.
相信大家都想到.
基督耶穌就是那個人的至寶.
其他東西都是短暫和過眼雲煙.
不知不覺我們都.
不是我們.
我都已經有落油卡了.
你可以想像得到.
由初信讀這些經文.
到現在這樣走了幾十年的基督徒的路.
你可以想像到.

$^{2001}$基督是至寶這個才是最重要的.
其他東西真的會過去.
所以保羅他真的很明白這件事.
來勉勵菲律賓教會.
他自己.
然後不單止.
其實我們要說的就是下面那裡.
我為他已經刁棄萬事.
看作糞土.
即是他不理人怎麼看他.
為要得著基督.
並且得以在他裡面.
不得是有自己因律法而得的義.
這裡真是可圈可點.
不是他做好自己.
而是在基督裡面.
那個新生命.
讓他能夠達至.
在服事.
在他人生.
他不需要刻意去守什麼律法.
但他整個人.
就已經能夠有信基督的義.
就是在神面前.
而來的義.
信從神而來的義.
即是說.
他的裡裡外外.
是被神所閱納.
稱許.
如果外在律法的義.
就在人面前.
你好像是無可指責.
你可以想像.
保羅已經提升了自己一個很高的層次.
其實馬丁路德就是這樣.
他很虔誠.
想做一個天主教的神父.
因為他許了一個諾言.
因為那次閃電.

$^{2041}$電他不死.
他就很害怕.
就祈禱.
他就要奉獻自己.
但是越奉獻.
他就越發現自己.
根本沒有真正的律法的義.
所以他最終都是因為羅馬書.
他就明白.
原來義不是自己去做到的.
為義人因信得生.
一章十六,十七節.
就使他重新啟發.
再研讀羅馬書.
所以保羅其實在《肥牛理說》已經說了.
他自己都不是因律法而得的義.
而是信基督的義.
就是在基督裡面.
他就能夠得到從神而來的義.
這就是一章羅馬書.
由一章到第四章.
就是講神的義.
加給什麼人呢?.
什麼人需要神的義呢?.
保羅在這裡就已經作了一個.
他自己後來在監獄裡面寫的監獄書信.
講述他自己的人生.
他的足跡.
他的成長.
接著第十節.
使我認識基督.
曉得他復活的大能.
並且曉得和他一同受苦.
現在他在坐牢.
效法他的死.
就是說我們為義.
為真理而受苦.
這樣來效法他的死.
或者我也得以從死裡復活.
這就是保羅現在的座右銘.

$^{2081}$他那個事奉人生的格言.
他已經找到秘訣了.
他已經找到他的事奉.
所以你們監禁我沒有用的.
因為我整個生命是屬於耶穌的.
耶穌要我坐牢.
那我就坐牢.
我學效耶穌的死.
但是我也可以和他一樣.
從死裡復活.
他已經有的是.
信基督的義.
就是因信從神加給我的義.
所以保羅真的可以在神面前.
隨時可以見主.
如果神的時間到.
他要見主.
他隨時可以見主.
很坦然地見主.
所以這裡說的復活.
保羅就是在事奉裡面.
在現在這麼大的艱難開始的時候.
甚至很快就不能陪伴教會了.
最大的就是他不能陪伴提摩太.
看著他一路成長.
不要緊.
那個是神的工人.
所以保羅就說.
你打完仗了.
到你打了.
他都信任神揀選的工人.
難道我看著你一路服侍嗎?.
難道我一定要一百歲才可以死嗎?.
不需要了.
神的時間到.
保羅就交付棒給提摩太.
要信任神的工人.
不是你拿著不放不交棒.
是要提摩太.
你是神的工人.

$^{2121}$所以你要記住.
神會在你生命裡面帶領你.
我算得了什麼呢?.
我完了我的人生.
我要去見主.
從今以後就有公義的官面為我存留.
即是你有你的人生.
你有你的牧場.
你有你將來的關難.
但更重要的是.
你和我一樣有那位主.
這是很重要的.
所以保羅就在這裡勉勵.
菲律賓教會的弟兄姊妹.
說他自己的事奉人生.
這就是復活的大能.
所以這裡的復活大能.
直接應用在事奉裡面的能力.
不僅是說基督的復活這麼簡單.
這裡是整段聖經.
他講述他自己的事奉.
就是有一種復活的大能.
在他的內心裡面.
再看貼撒羅尼加前書的四章.
都要完結了.
因為又變成下星期我們才說.
神的義那個大課題.
一到四章都是神的義這個主題.
神的愛要到第五章才說.
貼撒羅尼加前書四章.
我們講過的末世論都提過了.
是有一個物質的復活.
是講述信徒.
死了的信徒復活.
四章要很快讀了.
第十三節.
四章十三節貼撒羅尼加前書.
論到睡了的人.
即是死了的.
我們不願意弟兄們不知道.

$^{2161}$恐怕你們憂傷.
像那些沒有指望的人一樣.
這裡提到那個盼望.
我們若信耶穌死而復活.
乃已經在耶穌裡睡了的人.
即是死了的人.
神必定會將他.
他就是死了的人.
與耶穌一同帶來.
死了即是完全沒有了.
完全沒有了都可以帶回來.
你可以想像到.
神將來的能力.
叫所有存在在地球的人.
信耶穌的人會復活.
你想想有多少.
你可以想像到.
整個鑽石山.
整個和合石.
很多信主的人.
那些沒有那麼多.
最多應該是博夫林.
博夫林那些全部都是信耶穌的.
博夫林全部一起復活.
你幻想一下那種美麗.
那種震撼.
當然你想想.
整個香港有那麼多墳場.
都復活了.
基督徒.
你可以想像到.
還有些根本沒有墳墓.
有些死無葬身之地.
已經沒有人認識.
但是所有的在主耶穌裡面都會復活.
這個還不是盼望.
而且保羅是講得很實在的.
第十五節.
「我們現在照主的話告訴你們一件事.
我們這活著還存留到主降臨的人」.

$^{2201}$所以是第二次主耶穌的榮耀降臨.
不是第一次的降生.
有呼叫的聲音和天使掌的聲音.
對不起.
我可能讀早了.
「我們這活著還存留到主降臨的人.
斷不能在那已經睡了的人之先.
因為主必親自從天降臨.
有呼叫的聲音和天使掌的聲音.
又有神的號吹響.
那在基督裡面死了的人」.
這裡提到次序.
「必定會先復活」.
死了的先復活.
即是主回來的那一刻.
死了的首先復活.
第十七節.
「以後我們這活著還存留的人.
必和他們一同被提到雲裡.
在空中與主相遇.
這樣我們就要和主」.
當然和那些復活了的人.
不只是他們和主.
一定是所有復活的人.
包括天使.
永遠的同在.
「所以你們當用這些話.
彼此勸諱」.
所以肯定是自勉的意思.
是安慰的說話.
是對死去的人的一個指望.
所以我們要承認.
是有復活的盼望.
所以一定是基督首先復活了.
祂現在在天上.
將來祂會降臨.
第二次的降臨.
所以應該是福音派.
所有福音派.
除非是新派.

$^{2241}$不信有主耶穌的再來.
福音派無論是災前災中等等.
都是相信復活的.
但是千禧年,毛千禧年.
都相信聖徒要復活.
主耶穌會再回來.
是肉身的再回來.
一個榮耀的降臨.
所以我們其實是在說.
《羅馬書》一章四節.
不過很快幫大家引述一下.
保羅書信裡面.
對於復活的運用.
他所理解的闊度.
首先這裡只是說到.
耶穌的福音.
祂的身份.
祂是來過這個世界的.
這位造物主來過這個世界.
成為人的.
祂是大衛的後裔.
即是說祂是照著.
神在大衛的約的應許.
但是實在按照聖靈來說.
主耶穌從死人當中復活.
而且是永活.
所以顯明祂是神的兒子.
即是說是造物主的.
這位耶穌的身份.
我們說到一章第四節.
下一頁給大家看看.
不過這裡是禱告.
保羅的書信有三次的祈禱.
我都印了.
我們下一次再看.
包括向外邦人祈禱.
十五章他有禱告.
接著開始真的進入神的義這個課題.
Witness of God.
在保羅的書信出現的次數等等.

$^{2281}$我下星期會詳細和大家看.
大家記住.
由第一章第十八節開始.
一直說到四章尾.
都是說神的義.
是一個大的主題.
當然我會再整理下星期的PowerPoint.
再和大家講述.
我們一個簡單的禱告.
感恩天父給我們一個半小時.
我們可以將羅馬書.
將保羅的行文.
他的內心和他的呼召.
大家一起去學習.
我們很盼望經文成為我們的感動.
我們在人生裡面.
雖然我們絕對不可能成為.
另外一個時代的使徒.
但我們的確是蒙召成為耶穌基督的門徒.
聖徒.
盼望我們能夠在主的國度裡面.
我們是有份的.
我們是彰顯神賦予的大能.
我們能夠在地上.
為著主的國度.
我們去行義.
告訴別人不是我們因為律法的義.
是順基督而來的義.
是從神而來的.
加給我們順主的人的義.
以致我們能夠在主面前被稱許.
被閱立.
我們再一次獻上感恩.
保守我們大家一起用信心去研讀.
因為很長的書卷.
十六章.
有時真的會很難去掌握.
盼望我們能夠一起去學習.
這樣感恩禱告.
奉耶穌基督的名.

\newpage



\section{羅馬書}
\label{sec:8LIvofDzP5M}
\textbf{2025-03-20 周四晚查經班《羅馬書(4):分題:神的義(三)》 講員:曾裔貴牧師}
\newline
\newline
連結: \href{https://youtube.com/watch?v=8LIvofDzP5M}{\texttt{https://youtube.com/watch?v=8LIvofDzP5M}} ~~~~ 語音日期: 2025-03-21
\newline
\newline
\hyperref[sec:WjqHgH3CzNw]{\small{< < < PREV SERMON < < <}}
~
\hyperref[sec:index_chronic]{\small{[返順時目]}}
~
\hyperref[sec:index_scriptual]{\small{[返順卷目]}}
~
\hyperref[sec:nWNWvRiUB3k]{\small{> > > NEXT SERMON > > >}}
\newline
\newline
$^{1}$感恩天父給我們今天晚上.
可以去到第四堂的羅馬書.
我們都朝著一個去對這本書可以掌握.
基本能夠可以對於每一個分段.
每一個重要的神學.
我們都一起去探討.
尤其是保羅在新約的羅馬書.
是引述相當多的舊約經文.
都盼望神給我們在去到這些經文的時候.
我們又回頭去看一下.
當時是怎樣來講述.
又讓我們在舊約,新約穿梭的時候.
不單止不會迷失.
反而能夠豐富我們對於保羅.
他的引述背後的原因.
使我們能夠更加知道.
我們雖然在新約的世界.
但其實這個救恩已經是在萬世之前.
神已經計劃要藉著你的愛字.
來賜給我們.
求主幫助.
帶領我們開始進入這些重要.
和可能都頗抽象的神學觀念的時候.
又讓我們明白.
我們這樣仰望禱告.
奉耶穌基督的名.
阿們.
讓大家知道一下.
我們在研讀裡面.
當然我們是有.
現在我都盡量有PowerPoint.
如果時間許可.
大家可以看著.
相對容易些去掌握.
因為也引用不同的經文.
但始終有時突發性想起一些經文.
就不是在我預備的時候有引述.
所以大家如果能夠在家裡.
方便有本記錄簿.
記下.

$^{41}$是方便大家去研讀.
另外就盡量.
如果晚上我可以做到的.
就是將錄音轉為一個很簡潔的文字分段分析.
很短的可能幾版紙.
整篇的錄音.
但更好的就是原汁原味.
你再聽一次.
就會掌握多些.
當然大家的時間每個人都不同.
最好就是.
我們比較少的就是課堂的討論.
但現在用Zoom比較難.
所以盡量爭取.
你可以在群組裡問.
或者你可以私下在WhatsApp問我.
我可以盡量就那個經文.
有些你未必跟得到.
未必明白.
未必掌握的.
你可以問.
我盡量因應你的程度去解釋.
你就可以繼續.
不要越積越深.
你都不明白.
就變成沒有興趣.
我都是想培養大家讀經的興趣.
但多一點要求的.
就是上下文的解經.
就必須要學習.
要掌握.
將來你讀其他的書卷.
特別是新約的書卷.
那個上下文就相當重要.
因為一定是有它的背景.
有它的論點.
所以我們都要有這樣的學習.
不要只是即興性的.
讀到一節聖經就覺得很重要.
金句是重要.

$^{81}$但要看它放在上下文的整體思路是怎樣.
大家學習.
我們去到第四講.
都是一個分題.
就是神的義.
其實是由一章的十六節.
一直到三章的二十節.
就是神的義是怎樣運作.
當然三章二十一節.
這是下一節的講道.
不再詳細了.
如果分段是一章十六.
一直到四章尾.
都是講神的義.
你可以看到這是一個很大的課題.
我們去掌握這本書.
我們第一章.
那七節我們上星期都算做講了.
今天應該集中是八到十五.
如果可以就講十六,十七節.
就會更加圓滿.
再溫習一點.
一到第七節重點是保羅.
知道自己信主.
以及已經被召.
蒙召.
他的蒙召是很特別的.
既是蒙召成為基督徒.
成為耶穌的門徒.
更加重要的是.
他成為了使徒.
一個特選的.
本來他沒有份的.
他自己都描述.
他自己是不配成為使徒的.
因為他逼迫神的教會.
他要去抵擋神的.
所以神的工作.
恩典有時會蓋過我們軟弱.
和無知的罪惡.

$^{121}$這就是使徒的蒙召.
但是他很清楚他的蒙召.
大家不要介意.
如果還沒有描述就儘量描述.
否則就會影響到其他弟兄姊妹.
大家自己安排.
他很清楚自己的使徒.
就是傳福音.
最重要是向外邦人傳福音.
他很清晰.
彼得·約翰是向猶太人作使徒.
他就要去到很遠的地方.
成為了外邦人的使徒.
所以連他的整個成長.
他的學識.
他的背景.
原本是沒有受訓去幫助外邦人的.
但是神可以這樣將他改變.
成為一個外邦人的使徒.
在他整個人生的每一個一輪半找的生活.
他的細節.
都是想到外邦人.
所以你可以想像到.
完全沒有猶太教的驕傲.
那股情緒.
一定是不斷對外邦人有一個擁抱.
很不容易的.
你可以想像到.
使徒神在他生命裡面.
做了一個很大的手術.
他才能夠不再有自己猶太人的驕傲和優越感.
他能夠對外邦人完全不會看不起.
反而會對外邦人神交給他人生的工作.
當然一到第七節.
都探討到福音到底是什麼呢?.
什麼是福音呢?.
有些人說禿頭也是福音.
生髮不同的美容也是福音.
很多濫用了真正的好消息.
我們說過奧古斯都.

$^{161}$第一任的羅馬皇帝.
他帶來整個羅馬.
因為打仗贏了.
屋大維被他打敗了.
所以他帶來整個羅馬帝國.
當然當時還沒有那麼大.
是一個羅馬的承平.
所以他帶來的國民的福音.
就是羅馬的盛世.
然後鑄幣.
不知道有沒有放煙花.
鑄幣來彰顯他的國力.
和皇帝賜給那些人的一個承平.
所以當時福音這個字.
都是有一定的敏感的.
但是保羅和使徒們.
他們很清楚福音就是耶穌的降生.
大衛的兒子成為人.
在十誡犧牲和死而復活.
這個就是福音.
所以要記得這一點.
這就是一到七節.
上個星期已經講了.
給大家一點印象.
第八到十五節.
大概的內容.
我們先讀讀第八到十五節.
讀聖經也是重要的.
我用舊的和合本.
不是修訂版.
因為我這本大字聖經比較舊.
第八節開始.
第一.
我靠著耶穌基督.
因你們的信德.
信心的德行.
英文很簡單.
都是他的信.
他們的信德全遍了天下.
我在他兒子福音上.

$^{201}$就是耶穌的福音上.
神的兒子.
用心靈所事奉的神.
可以見證我怎樣不住地提到你們.
就是神為他作見證.
見證什麼呢?.
就是他一直對羅馬教會.
有一個負擔.
很想去到他們當中.
在禱告之間常常懇求.
或者照神的旨意.
終能得平坦的道路.
往你們那裡去.
因為我切切地想見你們.
要把些屬靈的恩賜分給你們.
使你們可以堅固.
接著十二節.
就是大家的見面.
大家彼此分享的時候.
或者分受保羅的恩賜的時候.
這樣我在你們中間.
如果他真的去到的時候.
因你與我彼此的信心.
就可以同得安慰.
就是說每一次的見面.
使徒們都帶給他們.
有信心的成長.
信心的建立.
接著就得到真正的安慰.
第十三節.
弟兄們.
我不願意你們不知道.
我屢次定意往你們那裡去.
要在你們中間得些果子.
如同在其餘的外邦人中一樣.
只是到如今仍有阻隔.
上文就說照神的旨意.
這是第十節.
或者照神的旨意.
不是照他自己的心思.

$^{241}$他自己的獨斷獨行.
無論如何都要去.
但他都是等候神的旨意.
能夠有一個平坦的道路.
可以去到他當中.
接著這裡就說.
只是到如今都有阻隔.
但是對他內心的感動.
那個熱切.
那個迫切.
是沒有減少.
就是說困難沒有減少.
保羅要向未去過的外邦的教會.
或者地方去傳福音.
這是他的生命的責任.
所以既是等待神的旨意.
能夠有一個平坦的道路.
還有那些難阻能夠除開.
但是在另一個角度.
他的內心有神為他作見證.
他是很焦急.
很想很切切.
將福音神的恩賜.
與羅馬的教會分享.
使彼此得安慰.
他們的信心又可以得到見證.
這是一個很成熟的牧者.
很想將他自己所有的.
分享給其他的教會.
尤其是外邦的教會.
所以保羅是兩方面的.
不是他自己祈禱一個地步.
要做什麼就做.
其實也不是.
他也有一個平衡.
既等候神的旨意.
有一個平坦的道路.
還有沒有難阻.
但是其實他內心是切切想去.
想見.

$^{281}$這樣的平衡其實很好.
很多事情我們不要強求.
神的時間,神的開路.
神的印證.
這些我們也要想的.
我們要看的.
但是如果只是站在那裡祈禱.
只是神的印證.
怎樣去印證呢?.
我們又不去策劃.
我們又不去等候.
又不去祈禱.
是不行的.
所以兩者其實沒有衝突.
但是怎樣將這兩樣東西調教到平衡呢?.
每個人都在神的感動裡面.
有不同的步伐.
重點我想也是.
那個心意應該要先有.
否則我們就很難說這是神的旨意.
突然想起尼希米.
大家認識嗎?.
尼希米他做什麼工作?.
他是政府的高官.
不過是波斯帝國的政府.
皇帝阿俠隋老身邊的酒政.
他是在外邦做酒政.
有一天他見到有些猶太人.
在以色列來到波斯.
就問他現在家鄉怎樣.
告訴他現在城牆破落.
城門都被火焚燒.
尼希米聽到之後.
悲哀幾日禁食祈禱.
一直這樣禱告.
到了第二章.
突然間皇帝發現他有什麼問題.
為什麼悶悶不樂.
不開心.
做酒政.

$^{321}$現在不是做南和許冠文.
現在不是做破地獄.
那些人經常哭.
如果去安息聚會.
如果笑.
經常笑.
我也不可以經常笑.
去到安息禮拜.
你要表達沉重的心情.
和哀哭的人同哭.
所以做酒政.
第一個要求就是要笑.
但是皇帝發現他不笑.
就問他為什麼不笑.
他說我祖宗的神的地方現在很破落.
現在很慘.
皇帝就問他想要什麼.
讓我回去重建聖殿.
重建我列祖的地方.
這樣.
皇帝就准許他.
那裡祈禱了多久.
三個月.
不是一分鐘.
因為在二章一節說了.
是基斯留月.
他在之前的一章那裡.
是另一個月份.
所以已經差了三個月.
你可以想像到.
原來尼希米貴為高官的酒政.
他也這樣去禱告.
默默的等候.
聽到這個悲哀的自己祖國的破落.
他就有一個很大的感動.
很想回去參與重建.
一邊等候一邊祈禱.
所以保羅也是這樣.
所以剛才這個尼希米就是一個例子.
保羅就是一樣這樣.

$^{361}$很迫切的禱告.
很想將他這個整傳的.
對主耶穌的救恩的原備的形式.
來跟羅馬的教會分享.
所以同時他也在等待神的時間到.
神的旨意來了.
他就去.
所以這裡是兩個很平衡的舉措.
接著就是十四節.
無論是希利尼人或者希臘人.
化外人就是barbarian.
就是野蠻的人.
當時在歐洲的北面.
是有這些未開化的人.
相對於羅馬人他們是未開化的.
所以很殘暴的.
可能到處去打家劫舍的.
聰明人不用說了.
很有知識的.
因為希臘也有很多哲學家.
很有知識的人.
愚蠢的人.
一般的平民百姓.
接著保羅說.
「我都欠他們的債.
所以情願盡我的力量.
將福音也傳給你們在羅馬的人」.
他不是說那些上文的四類人.
是說羅馬你們的人.
就是說你們也有聰明的.
窮乏的.
因為當時有奴隸.
也有不同的民族在羅馬.
這麼重要的首都.
所以保羅也欠羅馬人的債.
保羅也看自己對羅馬有虧欠.
他要將福音傳給他們.
第十六,十七節讀完.
「我不以福音為恥」.
回應上文說的.

$^{401}$他既然是欠債的人.
又看到他自己有托付.
是試圖.
他自己要傳整傳的福音.
這個福音.
當上個星期大家有溫習.
是神已經在舊約的時代.
透過眾仙支已經來啟示.
已經這樣默示出來.
所以保羅不以福音為恥.
不只是現在這個時代.
因為當時羅馬帝國.
肯定福音是不受歡迎的.
真正的福音市面上賣的.
就是奧古士都.
他帶來的整個羅馬的承平.
整個世界復歸於統一.
在奧古士都的統治下.
宗教都有一個寬容的容許.
某程度上是福音.
某程度上不需要再打仗.
都是一個福音.
但是保羅不會因為當時的思潮.
淡化了真正的福音.
就是耶穌基督.
但是你可以想一想.
當我們去到將來的十六節.
你可以想一想.
他一說到這個羞恥.
一般人怎樣理解耶穌.
猶太人看他是沒用的.
你說自己是神的兒子.
那你為甚麼會死呢?.
為甚麼你會在十字架上呢?.
十字架是被咒詛的.
所以這個就是猶太人看的耶穌.
外邦人更加是求那個哲理.
你都沒有學問.
一個人替我死了有甚麼了不起.
要很特別的哲學.

$^{441}$才能夠讓那些希臘人覺得很厲害.
所以後期保羅再過了一兩個世紀.
那些洛斯底主義就出現了.
保羅時代還沒有.
約翰差不多死的時候.
開始有那些洛斯底主義.
但是保羅時代就不是這樣.
但是你可以想像到.
突然間說「不以為恥」.
即是說在當時代.
一定是一個這樣的思潮衝擊.
信耶穌絕對不高檔.
很低俗.
出來還可能受到欺壓.
所以保羅就這樣跟羅馬人說.
我不會看我現在這個傳福音.
是一個羞恥.
我用一個屬靈的眼光去看.
是神的大能可以拯救一切相信的罪人.
保羅自己就是一個猶太的罪人.
本有猶太教的學問.
對舊約很認識.
但是他在神的眼中.
他是一個罪人.
但是大能的神.
透過福音改變保羅.
完全扭轉他對舊約的誤解.
對神的要求的誤會.
因為九章到十一章.
就正正式式說.
猶太人到底掉落在哪裡呢?.
他們整個民族到底在哪裡掉落呢?.
就是以律法為義.
不肯接受神的義.
十章四節還是五節呢?.
我們遲些會碰到這些經文.
你可以想像到.
保羅一直追求的都是律法的義.
所以他現在更加明白.
這個大能的神.

$^{481}$透過福音.
可以將他整個錯誤的神學觀.
徹底扭轉.
現在明白福音才可以改變.
所有世界的人.
不論他是猶太人還是希利利人.
即是泛指一切的人類.
所以提到猶太人和希利利人.
即是猶太人以外就是希利利人.
就是外邦人的意思.
是一個統稱.
所以不只是這兩個民族.
猶太人和希利利人.
不是這個意思.
羅馬人就不夠了.
一提到猶太人和希利利人.
保羅就變成一個統稱.
所以他告訴我們.
神的大能.
透過福音去拯救相信的人.
首先是猶太人.
然後是希利利人.
十七節就說了那個「因為」.
神的「而」在福音上顯明.
所以神的「而」.
在沒有「而」的人類的生命裡面.
藉著福音的訊息.
藉著甚麼途徑呢?.
藉著人的信.
就可以得到神的「而」.
所以三章二十一節開始就說.
神的「而」在律法以外顯明.
這裡先說了一個標題.
神的「而」在福音上顯明.
三章二十一節就是.
如今神的「而」在律法以外顯明出來.
所以是一個首尾呼應.
很快讀讀第三章.
三章二十一節.
你會看到整個氣氛就是這樣.

$^{521}$但如今神的「而」在律法以外.
已經顯明出來.
有律法和先知為證.
一提到律法和先知.
即統稱是舊約的意思.
不只是律法書和先知書.
不是這個意思.
律法和先知是保羅自己的用語.
怎樣運作呢?.
神的「而」二十二節.
「因信耶穌基督」.
即是福音的內容.
「加比一切相信的人並沒有分別」.
很清楚地將神的「而」.
怎樣運作在不義的罪人的生命裡面呢?.
即是神不再在律法裡面尋找人的「而」.
人亦不能透過律法達到神的「而」.
所以一定是透過主耶穌的救恩.
神的「而」就白白地加給那些信的人.
所以你看到第十七節和三章二十一節.
剛剛一個開括弧,關括弧的.
中間十八到三章二十節.
這麼長的篇章說了人類都是罪人.
萬民都需要「因信」「清義」.
都需要在律法以外尋求神的「而」.
神的「而」哪些人可以得到呢?.
任何人只需要信就可以得到.
這個分段很清楚.
所以第四章又是很簡單的一章.
原來舊約已經有「因信」「清義」.
大衛和亞伯拉罕是保羅第四章引的例子.
原來舊約已經有這些名人,偉人「因信」「清義」.
大衛保羅引的是詩篇三十二篇.
「得赦免其罪」「遮蓋其惡」的人就有福了.
「神不以此人為罪」就有福了.
大衛的懺悔詩.
亞伯拉罕就說到他在無可指望的時候.
信神就成為有指望的人了.
所以第四章告訴我們.
舊約已經有「因信」「清義」.

$^{561}$不是保羅發明的.
其實就是要和猶太人說.
他們一直都不肯接受耶穌.
信這位耶穌就可以得到神的義.
他要立的是自己的義.
所以得不到神的義.
這個分段遲些我會再和大家做的時候.
大家要牢記.
現在輕輕的說一說.
我們現在說第一章的內容.
一到第七節就是保羅的一個自述.
他自己的身份.
他自己的託付.
他就傳福音給外邦人.
第八到第十五節就說到.
他自己的內心是很愛外邦人.
很想透過禱告去前行.
去將福音都傳給羅馬的教會.
就要等候神的旨意.
神的時間有一個平坦的道路.
難阻又可以除開.
他也切切同時間.
他的內心很熱切地為羅馬教會弟兄姊妹去禱告.
很想去到他們中間和他們彼此分享.
彼此的信心得到建立的時候.
大家都得到神的安慰.
另外這一段經文也值得我們再想一想.
尤其是重點我想說的是.
十三節.
「我不願意你們不知道.
我屢次定意往你們那裡去.
要在你們中間得奢果子」.
留意一下這個描述.
「如同在其餘的外邦人當中一樣」.
如同.
即是我向貼薩諾尼加.
菲臘比,馬其頓一帶,被推尼.
或者是那個.
還有一個地方.
比利亞那些人很殷勤讀經.

$^{601}$例如雅典,哥林多.
很多這些外邦人本來已經去了.
他現在寫這封信的時間就在哥林多.
所以很多歐洲的地方他已經去了.
即是東歐的地方.
一直沿著海岸線一直下來.
現在差不多是第三次旅途.
土耳其的地方.
例如安提阿又會去.
又有外邦人 很多外邦人.
所以保羅說.
「我也很想將福音傳給你們羅馬的教會的弟兄姊妹」.
「如同在其餘的外邦人當中一樣」.
這個「一樣」.
好像我在其他外邦人當中.
所以保羅真的沒有一個偏心.
哪些地方的教會大一些.
我就要用心一點.
有些地方沒有那麼富庶.
就不要放那麼多精力.
有時我們的策略性.
好像有時很容易受環境影響.
當然我們作為弟兄姊妹.
我們最直接的就是.
在教會千萬不要有一些個人的平頭品足.
這個有錢一點 我貼近一點.
這個窮一點 我坐遠一點.
千萬不要這樣.
現在香港比較少.
很少會這樣.
偶然間也會有.
以前好像也有一個例子.
偶然間說過.
我的妹妹去了傳宣崇拜聚會.
去了一段時間.
教會要直行天水圍.
他們經過等候.
兩夫妻一起.
跟隨牧師.
那時大概五十多人.

$^{641}$去天水圍直行.
傳宣是一個差不多一千人的教會.
七八百人 沒一千人那麼多.
直行天水圍.
你想想當時天行村.
是二十多年前.
剛剛開的時間.
真是很多內地的新移民.
大家看過天水圍的日與夜.
大概看到那個劇情.
老夫少妻.
被人哄了下來 嫁了下來.
發現老公很窮 很老.
真是這樣的.
她聚會了一段時間.
好像不夠一年.
跟她先生說.
不如我們轉去傳宣.
她先生問她.
為什麼 我們來直行很有心.
你都答應了.
她說 最近坐在她旁邊的崇拜姐妹.
在撩鼻屎.
很噁心 受不了.
我們轉教會了.
我妹妹又是大銀行經理的秘書.
習慣了穿得漂亮就上班.
很不習慣.
她先生跟她說.
我們來有這個負擔.
要幫這些基層的教會.
當然這個可能她說了一下.
後來都沒有實行.
有時我們都有這個準備.
尚且有時都會有這些.
總會有不同的比較.
真是很不同的.
譬如現在說在基層的地方.
洪水橋真是不容易的.
有些是很基層的.

$^{681}$有個別的肢體.
有些朋友都有情緒病.
有時要處理這些事.
其實都不同的.
所以保羅是真的有分的.
你看他的十四節.
無論是欺劣人,犯外人,聰明人,愚蠢人.
但他都提到.
我都欠他們每一個類別的債.
所以真是他一視同仁.
這個真是很不容易.
我想我們一生都要面對.
不要用我們自己的尺去量度別人.
我們都給別人一個空間.
我們都有一個同理心.
我們能夠有富裕.
有各方面餘暇的時間.
可以沒有那麼緊張,沒有那麼加班.
不需要做到乾火性.
有些人真是日以繼夜地做.
子女沒有人理.
這樣的情況.
我們應該感恩.
我們應該對自己有的恩典.
神給我們的祝福.
我們看為一個恩典.
當然不是我們每天都要施捨.
不是要這樣做.
而是至少我們的內心像保羅一樣.
我都欠他們的債.
他真是會分辨這些人.
他真是會看到不同的類別.
說得太高深都沒有用.
他都不識字.
譬如現在我星期四上午的長者團契.
他們要求要加碼.
我本來是隔個星期四的.
因為我要預備.
晚上要查經.
時間不夠有時做powerpoint.

$^{721}$他們覺得很開心.
又不是太深.
他們很想每個星期都有.
但是你說得太深都沒有用.
因為每次說一節聖經.
簡單講解.
用很多不同的例子.
他們才會明白.
因為真是又老又不懂讀書.
真是不能夠.
你看他.
好像說了很抽象的東西.
所以.
保羅就有這種.
屬靈的情操.
真是看.
福音去到不同的人.
有不同的傳達方式.
但始終如一.
他對每一個民族.
都有這樣的一視同仁.
我傳給你羅馬的人.
如同.
其餘的外邦人一樣.
這個一樣.
就是反映他.
內心裡面的態度.
那個.
對別人真是不同.
我很感恩.
真是那些牧者自己講.
因為我去到紅欣之後.
都是請一些.
某程度現在基督教裡面.
都很難請的牧者.
其實因為錦宣過去已經認識.
他們來.
很多是第一次來洪水橋.
這麼大的兒子.
第一次來洪水橋.

$^{761}$港渡.
所以覺得很新鮮.
有時吃飯招呼他們問開.
其實原來他們不介意.
就在這裡.
去大教會.
去小型的教會.
基層的教會.
如果時間就到.
就很願意來講.
這個我想是每個牧者.
香港的牧者都有這樣的情操.
當然.
開始約的時候.
他看到自己的行事曆.
明明可以的都說不行.
那就另計了.
一般都不會這樣的.
但原來真的.
來到這些基層的教會.
基層的社區.
好像他們在大教會.
這樣講道.
例如這個星期林誠信牧師.
你想想金鐘的廣福堂.
你想想有多少.
出色的人.
他的教會有多大.
但他知道我們.
的邀請又很樂意.
這個星期日又過來洪水橋.
來講.
崇拜的信息.
你可以想像到.
這個牧者很重要.
不要選擇別人來傳.
不要選擇教會來服侍.
你可以想像到.
我屢次往你們那裡去.
要在你們中間得些果子.

$^{801}$如同在其餘的外邦人中一樣.
即是說.
不是因為你們羅馬是大的教會.
你們有很多的資源.
所以我怎樣都願意.
怎樣都會來你們當中.
不是.
如同好像其他的外邦人.
當中的教會一樣.
一樣真的很不容易.
我們實在.
真的需要.
在教會.
我們真的要學習.
真的每個人.
有不同的人生崗位.
人生的.
神秘的恩賜.
每個人都不同.
但來到神的家.
大家是一個很重要的身份.
大家都是神的兒女.
早信主.
遲信主的弟兄姊妹.
所以.
不要分那些工作.
你的階層.
諸如此類.
我們很重要.
首先我們的心.
千萬不要有一種.
還是那樣世俗的觀念.
如果是.
如果有.
其實很難成長.
你很難.
看得起其他人.
神的恩賜.
往往不會祝福這樣的人.
因為你的心已經有了一個.

$^{841}$錯誤的觀念.
所以我們記得我們一定是.
分享.
將我們有的去分享.
看自己是一個欠債者.
肯定是一個.
很重要的內心的.
價值觀.
所以我想著密一點.
大家可以勉勵分享.
否則.
我們都很麻煩.
一個人如果他.
時時都有這種優越感.
很多事情他就高不成低不就.
低的事情.
你找我做有什麼搞錯.
什麼人什麼身份.
這樣就不好.
這個就是欠債.
始終如一.
對於不同的人.
傳給羅馬的人.
我們就看看.
有沒有時間看看.
我們先來說說.
保羅的祈禱.
其實這裡.
一到九和十節.
我們讀過了.
這是保羅對於羅馬教會.
他切切的禱告.
另外在羅馬書還有.
兩段祈禱.
第九章.
猶太人恩信稱義.
第九章第一節.
弟兄們我心裡所渴慕.
和向神所求的.
是要以色列人得救.

$^{881}$這是第九章.
我們可以稍微揭一揭.
保羅在第九章.
開始講述.
為什麼猶太人.
會得不到神的義.
他們在哪裡跌倒.
就是因為他們.
只是追求自己.
是猶太人是神的選民.
他追求的是律法的義.
所以得不到神的義.
第九章.
保羅一章的語氣.
由講猶太人.
講外邦人.
轉去講猶太人的時候.
他就有這個.
內心裡的禱告.
九章一節.
我在基督裡說真話.
並不謊言.
有我良心被聖靈感動.
給我作見證.
我是大有憂愁.
心裡時常傷痛.
為我弟兄.
我骨肉之親.
就是自己被咒助.
與基督分離.
我也是願意.
剛剛第八章才講完.
沒有東西可以使我們與基督的愛隔絕.
這裡就說.
我願意自己被咒助.
甚至乎咒助到一個地步.
與基督來分離.
我也是願意.
願意什麼呢?.
就是願意猶太人得著神的義.

$^{921}$接著就開始講.
這個.
講到他們是怎樣.
我們遲些去到最後再說.
對不起.
剛才我引錯了經文.
應該是十章第一節.
十章.
弟兄們我心裡所願.
向神所求.
是要以色列人得救.
我可以證明.
他們向神有熱心.
但不是按照著真知識.
因為不知道.
開始提醒了.
不知道神的義.
我們外邦人就知道.
耶穌就是神的義的具體彰顯.
耶穌的死而復活.
這個就是神的義.
所以就要立自己的義.
猶太人.
不服神的義.
所以律法的總結.
就是基督.
使凡信他的人.
都得著義.
所以你會看到神的義.
一直穿插.
去到第四章講完.
暫時停一停.
但講到猶太人的時候.
再一次他們欠缺的.
就是他們沒有做那些行為.
他們有憲制.
有做其他的儀式.
但是他們沒有信耶穌.
就得不到神的義.
所以一提到.

$^{961}$羅馬書講神的義.
一定涉及到關於.
福音和耶穌.
生命的擺上.
有關連.
這裡提到猶太人的義.
保羅很想他們得到.
保羅就為.
猶太人.
切切的禱告.
他心裡面渴慕.
他甚至乎自己很想.
如果有需要他受咒詛.
他都願意.
對自己的同胞.
有這一份熱切.
他自己是過來人.
他深知.
繼續這樣去.
這些叫做.
甚至乎不是自圓其說.
自己活在.
自己的安書區裡面.
的猶太人.
以為自己就是亞伯拉罕的子孫.
真是肉身.
真是行割禮.
真是有很多不同的儀式.
真是吃那些.
不會吃那些不潔之物.
真是做了很多很多的儀式.
在自己的安書區裡面.
但是整個安書區.
原來是律法的義.
是得不到神的義.
這樣的情況.
保羅就情願.
自己被咒詛.
都想將他們帶離開這個安書區.
就是接受神的義.

$^{1001}$所以.
這個都是我們去到那些經文.
我們再看看怎樣給錢免禮.
跟著就是第十五章.
十五章.
看完這一節.
關乎到保羅的祈禱.
就是這三段.
在整個羅馬書裡面.
比較大的祈禱.
十五章.
我不讀這麼多了.
可能我們要講第十六十七節.
十五章的三十節.
弟兄們.
我直著我們主耶穌基督.
又直著聖靈的愛.
勸你們與我一同竭力.
為我祈求神.
叫我脫離在猶太不順從的人.
因為他要上耶路撒冷.
似乎是帶那些鵑蟀上去.
也叫我為耶路撒冷所辦的鵑蟀.
可蒙聖徒月立.
並叫我順著神的旨意.
歡歡喜喜.
到你們那裡.
很清楚.
他寫這封信.
就是很想去羅馬之前.
很具體地將神的義.
即是福音.
和聖靈的義.
和聖靈的愛.
而這個福音.
是怎樣的?.
寫得很詳細.
希望外邦人.
更加得到.
神各方面的祝福.

$^{1041}$所以他現在在這裡.
接近結尾了.
因為十六章是文安.
結尾的時候.
他最後邀請.
羅馬教會.
為他禱告.
而且竭力地一齊.
希望.
猶太.
因為他深深知道.
回到猶太的地方.
即是以色列.
要面對死硬派的猶太教.
那些知識份子.
大部分是由安提阿.
回到去的.
這樣來對保羅的攻擊.
無中生有的攻擊.
果然.
我們看《小道行傳》22章.
保羅在聖殿的範圍.
因為那些猶太人.
都上來守節.
在聖殿.
見到保羅.
聖殿外面.
見到保羅.
和年輕人.
特羅非摩.
在聖殿的範圍.
以為保羅.
帶外邦人進入聖殿.
這樣是不行的.
外邦人是不可以進入的.
只可以在外院.
所以就鼓動當時的人.
暴動.
捉拿保羅.
所以保羅就被.

$^{1081}$由那時候開始正式被捉了.
去到大祭司那裡審問.
但是.
輾轉.
因為他是猶太人.
他就.
對不起.
因為他是羅馬公民.
所以那些祭司長都懼怕.
當時的千夫長.
在以色列的.
軍兵.
就怕.
受牽連.
就問保羅.
我是羅馬公民.
我要上訴去蓋薩.
去到意大利.
所以他們就不敢動他.
就將他遞解.
一路去羅馬.
首先就送到蓋薩尼亞.
這是聖經22,23章開始.
就開始講那個情況.
但是沒辦法.
但是神的計劃旨意.
都是要保羅帶著鎖鏈.
去羅馬.
但是在他還沒去之前.
他很想.
藉著我主耶穌基督.
又藉著聖靈的愛.
勸你們與我一同竭力.
聖靈是有愛的.
所以如果聖靈充滿.
應該感受到神的愛.
聖靈充滿.
如果你是不由自主地在那裡.
又害怕.
又震驚.

$^{1121}$靈恩派那些就是.
狂笑,狂哭.
這個是不是聖靈的愛呢?.
似乎不是.
因為很有sense地.
很有那種對猶太人.
既是.
他們是很難信主.
但是保羅仍然.
好像都要上去.
將外邦人的.
外邦人的捐項.
記住是那些獻金.
是給猶太人的選民.
的基督徒.
來到外邦的教會.
不是錢有多少的問題.
是那個心意的問題.
外邦人現在.
信了主耶穌.
是猶太人所拒絕的福音.
外邦人現在回饋.
耶路撒冷的聖徒.
其實是很大件事的.
就是神的工作.
已經藉著外邦人.
來到去.
很想去敲醒.
猶太人.
他們那個固執.
他們那個偏見.
你想想外邦人本來是.
你和外邦人是拒絕來往的.
猶太人是.
不可以和外邦人來往的.
如果來往完之後.
你去市集買完東西.
你回家要行一個潔淨之禮.
就不是洗個澡.
是行潔淨之禮.

$^{1161}$這樣來表示.
外邦的污穢.
是你沒有份沒有沾染.
不是回來消毒.
不是回來因為有流感.
不是這些.
你可以想像得到.
是多麼的自掘.
一個鴻溝.
拒絕外邦人.
所以外邦人.
由保羅和西拉.
應該提摩泰也有.
帶鵑航上去耶路撒冷.
所以保羅說.
你們與我同歉歷.
為我去祈求神.
叫我脫離在猶太不順從的人.
絕對不是那個祈禱.
不應驗.
而是神有祂自己的計劃.
神要透過猶太人的逼迫.
使得保羅一路帶著鎖鏈去羅馬.
但不等於這個祈禱不成就.
然後又說.
也叫我為耶路撒冷所辦的捐項.
我蒙聖徒月.
即是讓猶太的聖徒.
感受到這一份愛.
是外邦弟兄姊妹的愛.
所以真的.
如果我們曾經有些人.
看我們都不可能信耶穌.
如果我們現在信了耶穌.
所以我很懷念.
因為小學就沒有印象了.
一點印象都沒有.
但中學還有很多印象.
還有些老師的印象.
但如果我的老師現在都八十多歲了.

$^{1201}$可能不止了.
所以都有些印象.
但是從來都沒有機會.
能夠和他們說.
我信了耶穌.
生命改變.
我記得我都提過.
我在中二的時候.
有一個老師叫Miss Lee教英文的.
後來回到加拿大移民.
她教了一段時間.
因為我哥哥是高班一點的.
不知道是中四還是中五.
我哥哥因為是很出名的.
踢校隊.
經常都要出去比賽.
所以很受歡迎.
很出名的.
某程度上都很乖的.
Miss Lee因為她教我中二的時候.
就認識我.
知道我哥哥是阿富.
我是阿貴.
她認識我哥哥.
都很疼我哥哥.
我中二的時候已經在學校很頑皮.
很學壞.
有一次在課堂.
在課堂不是小息.
在課堂大聲.
因為坐得很後.
很頑皮的坐得很後.
就大聲罵Miss Lee.
用粗口.
Miss Lee說沒見過人這麼粗魯.
算了.
但我想在她內心是很傷的.
因為她很疼我哥哥.
認識我哥哥.
接著就沒有了.

$^{1241}$一路輾轉.
那件事一直到現在我腦海還是這麼深刻.
所以.
這個印象.
就在Miss Lee的情況.
因為我那時十幾歲.
她已經30歲.
有多了.
可能都已經八九十歲了.
所以你可以想像到.
這裡告訴我們.
可蒙越納.
外邦人是猶太人這麼憎恨的.
在.
律法裡面.
是無份無關的.
現在外邦人居然聽到耶路撒冷.
有饑荒.
集那些錢奉獻給猶太人.
的聖徒.
所以透過保羅帶去.
就更加促成.
猶太和外邦人.
在主耶穌裡面的和好和合一.
所以保羅的書信.
已不鎖書.
菲勒比書.
歌羅西書不斷說.
在基督裡面合而為一.
不再分任何的民族.
所以.
這個可蒙越納.
某程度上是一種.
我們信主生命改變的見證.
一個挽回.
一個讓人看到.
神的大能的改變.
所以很可惜.
我見不到無份的老師.
現在有個.

$^{1281}$我哥的同學.
和一班同學.
現在來了紅恩聚會.
試過一次的世紀晚餐.
我哥和我.
和那兩位同學.
真是同班的.
很搞笑.
所以.
如何能夠讓我們的生命.
和一些以往.
可能看死我們的.
可能我們得罪的.
現在看到我們的生命改變.
所以外邦人現在奉獻給耶路撒冷的窮人.
其實是一個很大的見證.
是一個很大的.
神的恩典的一個印證.
其實從來都沒有界線的.
是舊約以為有界線.
其實在舊恩裡面.
從來都不會有界線.
只是信和不信的界線.
人類其實沒有界線的.
所以.
外邦人奉獻給耶路撒冷的窮人.
真是帶來一個.
在主裡面的一個合一的見證.
其實很感動的.
很寶貴的.
其實保羅體.
這個這麼集合性的.
整個歐洲的教會.
馬其頓宗教會.
菲拉比,帖撒羅尼迦.
哥林多.
他們這樣奉獻.
對耶路撒冷的窮人這些聖徒.
信了主的.
是一個很大的感動.

$^{1321}$這些是外邦人的錢.
我們現在用他們的錢去買糧食.
你可以想像到.
這個就是對保羅.
對外邦.
對耶路撒冷的窮人.
是一個三贏的祝福.
更大的就是神的榮耀彰顯.
所以.
第32節.
「並叫我順著神的旨意.
歡歡喜喜,到你們那裡.
與你們同得安息」.
一同得到.
這份主裡面平安的.
那份的.
平安,平靜.
所以.
保羅的內心.
是很有那種的胸襟.
很有那種的平安.
因為.
救恩令到保羅看的東西.
更加高,更加闊.
看到的是宇宙觀.
不只是他自己猶太人.
這些民族主義.
所以.
我們千萬不要有太大.
的一些我們個人的.
一些主義.
譬如民主,共和.
或者怎樣不同的世界.
都有不同的意思.
當然你說運動.
你喜歡哪一隊?喜歡利物浦?.
喜歡曼聯?這些就當踢球.
沒所謂.
當是一種平時的消閒.
沒所謂.

$^{1361}$但如果你去到極端.
又是一個很大的問題.
所以保羅就希望外邦羅馬教會.
和他一起.
竭力的去祈求神.
所以.
在羅馬書有這三個重要的.
祈禱.
我先和大家.
分享一下.
我稍為喝口水.
我們要進入第十六.
十七節.
一到十五節.
大家有沒有問題想問?.
當然你問我都看不到.
就當.
當大家是甚麼.
大家自己再溫一溫.
暫時.
還是保羅.
他自己的.
心願.
引言為甚麼我要寫這封信.
很想去羅馬.
和講述他自己的.
使徒身份.
他要傳福音.
他很熱切.
為羅馬教會祈禱.
很想去到他們中間.
十五章再一次回應.
他很想.
平安去到他們中間.
一同的同德的.
安息.
在主裡面的一份.
平靜安穩.
但是整個過程.
《少林寺傳說》告訴我們.

$^{1401}$一點都不平靜.
甚至沉船.
甚至那些水手想殺了.
囚犯逃走.
保羅就說如果他們逃走.
船就不可以了.
不准他們逃走.
只扔掉貨物.
到海裡.
平安擱淺之後.
上岸.
上岸去了馬爾他島.
Malta那個地方.
當時.
那些.
荒島.
當時應該不是荒島.
在保羅的年代.
那島主叫布伯留.
看到保羅.
上岸.
突然看到天氣冷.
他們生火.
一條蛇走出來.
毒蛇咬住保羅.
保羅一跳就跳到毒蛇.
那島主布伯留.
就說一定死了.
那條蛇咬.
但沒事很正常.
有個神蹟.
那些人就拜保羅.
當然不可以拜.
所以.
馬爾他每年都有一個聖保羅日.
紀念保羅.
帶這些島民.
信主.
雖然一點都不平靜.
去到羅馬.

$^{1441}$但是.
保羅的內心.
原來一直都有基督的平安.
去支撐他.
一直有基督的平安.
在他的內心.
所以我們不能夠.
改變世界的起起伏伏.
不可以改變.
不同時代的思潮.
但是.
我們必須要.
自己的內心.
有真正的平安.
去統治我們.
去管理著我們.
使我們處變不驚.
這也是重要的學習.
透過保羅的第一章.
和第十五章.
我們都感受到.
他的熱切.
他的平安.
他對羅馬教會.
很想分享他的恩賜.
給他們.
很想和他們同得.
安慰.
同得安息.
一章第十二節.
和剛才的第十五章.
那段經文.
背後.
所以我們看了.
保羅的心願.
很想去羅馬.
所以寫這封信.
先讓他們知道.
他的目的.
和更加想.

$^{1481}$將這個救恩.
神的義.
是怎樣能夠.
透過人的信.
就能得到呢?.
我們開始講.
第十六,十七節.
你可以將十六節.
歸到第十八節.
之後的段落.
也是對的.
或者第十六節歸到.
第十四,十五節也是對的.
因為一直都在講.
他不偏袒任何一個民族.
無論他是.
希臘人.
有哲學,有修養.
有學識的.
有化外人.
野蠻的,未開化的.
聰明的,愚蠢的.
這樣的人.
他都欠他們的債.
情願傳福音給他們.
所以這裡一直都在講.
他看這個福音是榮耀.
所以他就說.
我不以福音是羞恥.
這樣的原因.
所以歸十五節.
也是對的.
接著就解釋這個福音.
為什麼保羅不看為羞恥呢?.
他不感覺是羞恥.
反而感覺是榮耀呢?.
因為.
可以將任何一個地上的人.
無論是有學識,沒有學識.
貴族或平民.

$^{1521}$等等奴隸.
都可以帶他.
進入耶穌的救贖裡面.
就可以得到救恩.
成為神的兒女.
得著神的義.
所以.
這個就是神的大能.
可以改變地上的人.
所以保羅因為這樣的緣故.
他不覺得是羞恥.
不等於他沒有被羞辱.
兩件事是不同的.
他被羞辱.
但他完全不覺得羞恥.
這個就是榮耀.
所以我們今天的基督徒.
我也不反對心理學.
見父道我也不反對.
但是我們的內心.
那個質素.
那個堅韌夠不夠呢?.
是不是每次都說.
我又不知為何不開心.
很大件事.
經常多愁善感就不用做了.
我們這一代的基督徒.
是不是經常都多愁善感呢?.
是不是經常都.
無無故故地.
又說很不開心.
沒有悶不做事.
這樣是不行的.
所以一定要以福音為榮耀.
不代表我們不會受羞辱.
不代表我們一帆風順.
不代表我們.
甚麼都很平安.
甚麼都遷就你.
去到總站車是等你的.

$^{1561}$去到西鐵又要等你的.
你一進車就有位給你坐的.
個個都不敢坐的.
不可能是這樣.
可能更多的時候.
例如昨天74號X.
撞過去.
現在嚴重危殆.
完全沒有想過.
就算你怎樣捉實都沒有用.
可能是司機睡著了.
可能是其他因素.
你沒有選擇.
就算你多麼精神捉住椅子.
他都會撞過去.
都不關你的事.
所以我們在這個世界.
一定會面對不同的人生.
例如不信的人對基督教的攻擊.
現在是基督教裡面的大復興時期.
還是基督教裡面的軟弱時期.
等等.
這些都不是我們可以控制的.
我們信耶穌.
是我們個人去選擇明白了福音.
但是我們處於哪個時代呢?.
這是神的安排.
不是我們去選擇的.
沒得選擇的.
等於我出生在九龍塘小西湖.
1997號.
是沒得選擇的.
你就是在網區長大.
你出了地鐵就要走了.
你沒得選擇的.
是你安排在那裡.
八歲的時候.
那個收養我的娘太.
媽媽不肯就告吹了.
如果不是的話.

$^{1601}$我現在在加拿大多倫多.
我現在可能是多倫多大學退休教授.
但是沒有人知道的.
神安排了的.
定了的.
是你沒得選擇的.
就是去到你門口.
13個兄弟姐妹.
只要第9個阿貴.
已經定了的.
但是媽媽都說不行.
13個沒有了一個.
都是不行的.
我就不知道.
知道了就馬上撿包袱走.
但是不知道.
冥冥中我們回想.
就是神有他的管理.
我就一定要在.
這樣的母屋區長大.
在大亨東那段時間.
出入.
見到那些人開片.
斬人 黑社會.
等等.
接著輾轉就要搬.
就搬了去巴富街小學.
在窩兜窩兜都出名的.
接著就去了新發書院.
太太知道的.
都出名的 但是中一就踢出校了.
就沒得上了.
接著就派去無範.
九龍塘 隔壁的.
整個名校區最差的那間.
英文中學.
輾轉輾轉就在學校.
站在主任室門口的.
常客.
一直畢業.

$^{1641}$你沒想過突然會信主改變.
所以真是.
我想最大的.
重要就是保羅這一句.
不以福音為恥.
真是看到自己為什麼會信耶穌.
這個福音原來是一個很大的能力.
這個能力是.
讓你感受到神對你的愛.
一種無條件.
犧牲的愛.
在你的生命裡面.
讓你.
可以醫治自己的過去.
很多人.
到今天還不肯接受自己.
接受自己的原生家庭.
你也不肯接受自己.
不肯饒恕其他人.
其實是不需要的.
當然你家裡很順利.
你不需要饒恕.
但是很多人.
一直帶著這種仇恨.
自己的出現.
自己的成長.
所以不以福音為恥.
就是他內心.
看這個救恩是榮耀.
所以從來沒有說.
信耶穌.
就會很順利.
保羅的時代正是尼祿.
最當權的時代.
也不是奧古斯都.
奧古斯都已經死了.
然後就是格魯德都要趕走猶太人.
《小唐人傳》第18章.
第四任的羅馬皇帝.
將所有在羅馬帝國裡面的猶太人.

$^{1681}$全部要趕走.
所以阿居拉和伯基拉兩夫婦.
就由羅馬去到哥林多.
遇到保羅.
就和保羅同住.
一起織帳棚.
就是這樣的緣故.
所以猶太人.
信徒 基督徒.
一樣在那個時代.
受著羅馬帝國的迫害.
所以保羅.
他這樣說.
當時福音絕對不是一個高檔的娛樂.
是你信耶穌.
是有相當的代價的.
是要受很多猶太人對你的迫害.
如果你是猶太信徒.
如果你是外邦的.
為什麼你不拜蓋薩.
為什麼你不認蓋薩是神.
為什麼你不去參與公共崇拜.
敬拜那些偶像.
譬如那些巴特隆神殿.
眾神廟.
在羅馬的.
為什麼你不做這些.
為什麼你不拜那個阿底米女神.
很多這些.
譬如保羅說.
你在哥林多.
在市場買的肉.
你直接吃吧.
但是你去偶像的廟拜就不行了.
其實信耶穌.
是會讓人感覺到.
你會被針對.
所以保羅在這裡的時候.
很快就已經講述.
他自己一個很堅強的內心.

$^{1721}$就是他感受到.
神的大能.
在他內心.
至少先改變他這個猶太人.
他這個以為很厲害的拉比.
他以為還是為神.
去發熱心的一個法利賽人.
神改變他.
拯救他.
他曾經害過那些信徒.
無論在耶路撒冷.
在大馬士革.
他看著斯提反.
他在那裡笑.
斯提反殉道.
他在那裡笑.
他可能也有機會扔石.
保羅幫那些人看守那些衣服.
那時候他還有一個.
你可以幻想一下.
如果拍一套話劇.
他是站在那裡的一個少年人.
你想想他的眼神那種奸笑.
那種幸災樂禍.
對斯提反.
被人扔死.
又說見到天.
他見不到.
這個人還說見到天開.
差不多死了就說見到天開.
問題是保羅見不到.
他要去到第九章.
在大馬士革的路上.
他才見到.
都不是每個人都見到.
是他被神.
光照,選照.
才見到.
所以斯提反的殉道只有斯提反見到.
其他人見不到.

$^{1761}$見不到大光.
見不到天開.
見不到耶穌站起來.
接施提反.
所以不以福音為恥.
都是我們一生.
都要學習的.
弟兄姊妹.
今天我們沒有一個肉身的逼迫.
我們在工作地方.
都很少人會說.
你信耶穌.
我就會迫害你.
我就不加你工資.
我們很少這樣的情況.
一般就算他們不信.
基督教不會有太大的敵視.
我們好像慢慢忘記了.
基督教原來不用受苦.
基督教原來要受苦.
好像從來沒有想過要受苦.
但是其實.
只是第一節已經說.
我不以福音為恥.
就是說當時福音絕對不是一個榮耀.
在一般人的眼裡.
在政權的眼裡.
羅馬的政權的眼裡.
從來都不是一個榮耀.
但是.
作為牧者.
作為使徒.
保羅就很清晰的講述.
他很清楚.
他為什麼今天成為使徒.
就是因為他看到福音的榮耀.
所以他很願意.
很賣力的將福音傳給.
不同類別的人.
他裡面內心.

$^{1801}$是再沒有一些對其他人的偏見.
或者是一些誤解.
或者是一些偏見.
或者是一些誤解.
這樣.
所以福音是.
神的大能.
這個大能未必是一種神蹟.
往往是.
讓一個人.
明白真理.
感受到愛.
感受到這份神人的復和.
而帶來有力量.
他去行出.
和好.
和人和好.
首先和自己和好.
明白神在你的生命裡有計劃.
真的.
我經常回望.
很開心的.
原來童年是這麼開心.
這麼有計劃.
很讚嘆神.
讓我在一個.
這麼窮苦的世界成長.
我真的會回望.
感受到原來是有安排.
有計劃.
這樣你看自己.
其實是健康很多.
不需要那麼多抑鬱.
思前想後.
而是冥冥中有安排.
所以.
這就是神的大能.
其中這裡.
不是說不可以有神蹟.
可以的.

$^{1841}$但神的大能往往是.
你感受到福音的榮耀.
然後你下力.
去叫人帶人.
去進福音裡面.
所以就要拯救一切相信的人.
要拯救一切相信的人.
就去到第十章.
保羅說如果沒有奉差遣.
怎麼可能去傳道呢.
如果沒有傳道的.
怎麼能夠聽見呢.
不聽見怎麼能夠信呢.
如果不信.
怎麼能夠得救呢.
所以一直都在說.
那個差遣.
所以這裡告訴我們.
當你感受到神的大能.
神的醫治.
神的福音的榮耀.
你就會有這份心.
很想將福音盡力的.
去拯救一切相信的人.
所以.
我們的責任就是.
先去和自己和好.
讓福音改變我們的生命.
改變我們的價值觀.
使我們能夠.
很想.
將耶穌告訴人.
讓人看到.
我們的改變.
真是福音令我們改變.
接著就是第十七節.
這個「因為」.
因為神的義.
在這個福音上顯明出來.
所以這個福音.

$^{1881}$其實是在說耶穌受咒助.
耶穌.
為我們成為罪.
所以我們在.
祂的裡面.
就成為了神的義.
這是哥林多後書.
第五章.
近尾那裡.
神使無罪的耶穌.
替我們成為罪.
好叫我們在神的裡面.
就得到神的義.
所以這裡告訴我們.
神的義.
神的義就在福音上顯明.
所以不在福音上.
就顯明不了.
所以你不信主耶穌.
你不會預見到神.
也不會稀罕神的義.
就一直納自己的義.
一直都覺得自己好.
就不能夠.
感受到.
神的義的重要.
迫切.
神的義.
所以第十八節.
就講到神的憤怒.
因為你不緊張.
不稀罕神的義.
你就不會感覺到神的憤怒.
在天上顯明.
所以息息相關.
一直都是一個.
很重要的.
討論的邏輯.
因為神的義.
包含的就是.

$^{1921}$首先我們都沒有神的義.
我們自然就在神的憤怒下.
但是我們信了主的人.
第五章就告訴我們.
我們就在神的喜悅裡面.
寵愛裡面.
我們既因信稱義.
就藉著我們的主耶穌基督.
得以與神相和.
我們又因著他.
因信進入.
現在所站的.
聖地.
我們就在神的喜悅裡面.
藉著我們的主耶穌基督.
得以與神相和.
我們就在現在所站的.
恩寵之中.
第五章第二節.
羅馬書.
所以沒有神的義.
就在神的憤怒下.
有神的義的人.
因信稱義的人.
得了神的義的人.
就在神的.
Blessing.
不是Blessing 還是恩寵裡面.
第五章第二節.
你可以想像到.
這個神的大能的奇妙.
第五章第二節.
我們又藉著他.
因信得進入.
現在所站的恩典中.
其實英文就是.
恩寵的字.
並且歡歡喜喜.
盼望神的榮耀.
一直都在說一個.

$^{1961}$很平安喜悅的.
一個狀態.
一個心靈的狀態.
不只是心理學.
是心靈的狀態.
與神已經接通了.
天地線.
即是說你的Wi-Fi.
現在已經去到5G.
不是很慢速的.
是很強的Wi-Fi.
可以去到神那裡.
現在站的恩寵.
你有沒有這樣想過.
在救恩裡面.
在神的眼中.
你現在站在一個恩寵.
你見過雨場片.
公庭電視.
那些太子女.
那些皇帝女.
隨時可以去到.
皇帝身邊.
很開心.
或者那些皇后.
得寵的皇后.
不是那些打入冷宮的.
皇后就是得到皇帝的恩寵.
或者那些很重要的大神.
就是在皇帝身邊的紅人.
現在保羅告訴我們.
原來我們外邦人.
我們可以這樣.
把自己固定在神的恩寵裡面.
還要是歡欣喜.
盼望神的榮耀.
這就是神的義.
延伸在現實的地球的基督徒.
回到天堂就甚麼都不用了.
現在在地上.

$^{2001}$在這個苦難的世界.
基督徒就是怎樣去.
在神的義下.
我們的生活的狀況.
第五章到第六章.
就講這個課題.
五章先講那個.
恩寵的平安.
第六章就探討一個課題.
不准犯罪.
第六章其實更重要.
就是不准犯罪.
不准犯罪.
第六章其實更重要.
原來基督徒.
是可以不再犯罪.
是這個意思.
我們去到這些神學再討論.
但是現在我們仍然在講第十七節.
神的義.
原來是很廣泛.
是包含很多東西.
不只是卸下你的罪.
你走吧不再刑罰你.
現在釋放你了.
釋放完之後怎麼辦呢?.
救恩不是這樣.
救恩就是說釋放在神的眼中.
我們已經得到神的義.
其實就是耶穌的寶血救贖.
但是我們現在.
在神的面前生活.
有生之年.
我們就是在神的面前.
怎樣來生活.
所以一直延伸.
都是去到.
是不是要換電話呢?.
第十七節.
神的義.

$^{2041}$是怎樣運作呢?.
留意很重要.
顯明出來.
這義是.
本於信.
以至於信.
英文就很清楚了.
無論NASB,King James.
NIV.
比較權威的版本.
都是從信仰.
由那個開始的相信.
由你第一次.
有人跟你陪談.
你舉手決志信耶穌.
從信仰的開始.
一直到你的一生.
你的生命結束.
都是仍然是那個信.
從信仰.
到信仰.
這樣來得到信心.
所以.
還有幾分鐘.
(時間到).
這裡說了.
接著呢.
接著保羅引述.
如經上所記.
二人必因信得生.
哈巴谷書二章四節.
在整本新約.
有三次出現這一節的聖經.
在整本新約.
有三次出現這一節的聖經.
這裡是保羅第一次引.
加拉太書也是他引.
所以希伯來書.
到底是不是保羅寫的呢.
三段.

$^{2081}$不過很特別.
二人因信得生.
羅馬書是說二人.
如何得到神的義.
因信.
得到神的義.
但加拉太書是說有信心.
三章十一至十四節.
希伯來書更加有趣.
就是說.
生命的問題.
得到永生的人.
在多大的苦難中.
他都不會失腳跌倒.
所以一節聖經.
哈巴谷書二章四節.
二人因信得生.
羅馬書是二人.
加拉太書是因信.
希伯來書是得生.
所以下個星期.
我們會詳細講.
這個經文.
今天到這裡.
我們一起祈禱.
多謝天父.
簡單地將第八節.
到第十七節.
都簡單講解完.
盼望弟兄姊妹都能夠掌握.
一方面對保羅的熱心.
很想去羅馬.
將他的恩賜.
將福音.
跟他們一起分享.
得到神的恩典.
和安慰.
繼續讓我們看到.
開始進入神的義.
這個神學課題.

$^{2121}$都讓我們能夠充分地明白.
下個星期.
我們會由十八節開始.
講到三章的二十節.
律法的義.
或者個人的.
以為自己的行為.
其實是得不到.
和拒絕神的義.
求主讓我們能夠明白.
當我們由三章二十一節.
開始講述.
神的義在律法以外.
這樣顯明的時候.
又讓我們能夠掌握得清楚.
我們這樣感恩禱告.
奉耶穌基督的名.
阿們.
\newpage



\section{羅馬書}
\label{sec:nWNWvRiUB3k}
\textbf{2025-03-27 周四晚查經班《羅馬書(5):分題:神的義(四)》 講員:曾裔貴牧師}
\newline
\newline
連結: \href{https://youtube.com/watch?v=nWNWvRiUB3k}{\texttt{https://youtube.com/watch?v=nWNWvRiUB3k}} ~~~~ 語音日期: 2025-03-29
\newline
\newline
\hyperref[sec:8LIvofDzP5M]{\small{< < < PREV SERMON < < <}}
~
\hyperref[sec:index_chronic]{\small{[返順時目]}}
~
\hyperref[sec:index_scriptual]{\small{[返順卷目]}}
~
\hyperref[sec:7C1Uqd4_AYQ]{\small{> > > NEXT SERMON > > >}}
\newline
\newline
$^{1}$感恩我們再一次在晚上.
我們可以有恢復黃馬書的《查經班》.
求天父讓我們能夠對經文有認識.
也讓我們能夠對保羅所說的整傳的福音.
這個神的義是怎樣運作.
還有我們一直這樣去研習的時候.
我們自己都是深受神的話語的激勵,感動.
還有我們真的可以在人生裡面.
去為主的真理,為福音來效力.
求主幫助,仰望,感恩,禱告.
奉耶穌基督的名 阿們.
好 我們先開這個檔案.
我們來到第五講.
其實一直都在說一個大的課題.
就是神的義.
神的義其實由第1章16節.
嚴格來說應該延伸到第4章結尾.
當然比較重要的關鍵的鑰句.
就是3章21節這個段落的轉變.
所以我們可以逐段去看.
一個大的成卷書,一個總的主題.
肯定是保羅很想羅馬教會來支持他.
對這個世界,對萬民的宣教.
即是傳福音.
因為唯有福音可以改變和拯救萬民.
所以羅馬的教會要明白.
第一章很明顯就是一個引言.
就是很想羅馬教會認識保羅試圖的身份.
所以這個位份他是神特選的使徒.
特派他去傳福音.
而且特別是向外邦人傳福音.
這些大家在過去透過《四道行傳》.
或者其他的書信都有提到.
在羅馬書一開始就開宗明義說.
他是特派去傳神國的福音.
接著就說那個福音到底是什麼內容呢?.
就是神在舊約透過眾仙支已經有默示.
有那個啟示.
直至耶穌基督的降生.
就讓這個福音正正式式成為人類的盼望.

$^{41}$福音其實就是講到主耶穌的降生,受死.
祂的復活和祂的永活.
所以我們要傳揚的就是這位.
叫人得到生命的耶穌基督.
接著保羅就提到他自己和羅馬教會的關係.
就是他很想切切地去到他們中間.
要將他自己一些的屬靈的恩賜.
和他們去分享.
使得他們見到保羅的時候.
又可以同得信心和安慰.
接著他們又可以繼續支持保羅.
向西班牙傳道宣教的使命.
所以保羅就一直切務地去.
但是有難阻.
所以他又等候神的時間.
希望能夠有一個平坦的道路.
去到羅馬.
這個是保羅他想講的原因.
就是第十五節.
因為他是一個欠債者.
他很想要還這個福音的債.
無論是什麼人.
他都欠他們的債.
情願都來將這個福音都傳給他們.
所以這個第一章的所謂引言.
就是表白了保羅自己的身份.
他的使命.
和很特別的就是寫這封信給羅馬教會.
就很想他們知道他自己這一份的心意.
接著他突然間第十六,十七節.
就大進了一個救恩.
只要人相信.
就可以獲得神的義.
就是說似乎萬民唯有這個方法.
才可以改變到人性,人心.
到底人性和人心是怎樣的呢?.
在神的眼中是什麼狀況呢?.
所以十六,十七節.
首先是給了一個要解決.
一個不完整的世界的人性.

$^{81}$就是罪性.
首先就講了答案.
就是第十六,十七節.
我不以福音為恥的原因.
因為他知道福音是可以改變萬民的這一份能力.
這一份神的大能.
可以拯救一切相信的人.
一切在不義,在神的憤怒底下的罪人.
因信就可以得著這一份的救恩.
就是神的義.
所以他就很清楚看到.
唯有萬民因信稱義.
才可以解決人類的罪.
所以這裡就讓我們看到.
保羅很快在第十六,十七節.
就已經提出人類的答案.
那個命運.
那個解救之道.
所以十七節就是一個很重要的鑰節.
我們看看第一.
大家看到了.
很簡單的分段.
福音的趙明的保羅.
十六,十七節.
就是信心和神的義.
那個兩者的關係.
人必須要擁有神的義.
才可以得到救恩.
免去神的審判.
那個憤怒,那個刑罰.
但是怎樣可以得到呢?.
就是因信而得的義.
這個就是十七節.
因為神的義正在福音上顯明出來.
稍後我給大家看一個powerpoint.
在三章二十一節.
也是同樣出現.
神的義在律法以外顯明出來.
現在這裡就說.
神的義在福音上.

$^{121}$就是耶穌的救恩上.
來顯明出來.
這個義是本於信.
而致於信.
就是由開始信.
一直整個人生的過程.
都是要持續地信.
去到見主之前.
都是持續地相信.
這樣就可以得到神的義了.
所以是本於而致於.
大家要掌握這些很重要的真理.
所以不是一個交換式的.
我信一次.
以後不需要再信了.
其實這個信心是不斷地成長.
不斷在神裡面去更新.
這樣的因信而致於信.
就在這個義裡面有份了.
如經上所記.
「二人必因信得生」.
接著就是.
這裡第一章的大綱.
萬民欠缺神的義.
今天我們會重點講18到32節.
這個段落.
我們還沒講之前.
再溫習一下16,17節.
這裡大家看到.
神的大能就是拯救人類的計劃.
有足夠的力量可以改變信的人.
無論信的人.
未信之前.
是怎樣的一個罪人.
或者活在罪中.
無望的.
但是.
神拯救的大能.
都有足夠的力量去改變人的生命.
就是說.

$^{161}$與世界任何的思潮.
任何的政治體系.
任何的教育.
都是不能夠改變人心的.
唯有神的大能.
才可以改變人心.
所以保羅看到.
唯有福音可以去更新舊有的生命.
得到神的義.
神的義就是福音的核心.
所以主耶穌的死而復活.
人這樣相信主耶穌.
就可以得到神的義.
這是幫大家重溫一點.
所以神的義.
必須要通過信心.
是那個人的信主耶穌.
嚴格來說是信靠主耶穌.
在這個信心上就會顯明神的義.
所以如果我們真的這樣信.
就會立即獲得神的義.
神的義可以很闊的.
那個領域很大.
第五章就說.
神的義.
信的人會有甚麼的生命改變呢?.
這就是原來救恩.
會使我們的生命.
能夠在地上.
會帶來一個很嶄新的.
那個的.
不要叫內容.
那個的生命的表達.
所以第五章第三節開始.
一直到第五節.
就說到一個基督徒.
可以在神的義的包圍下.
那個生活是會有很大的不同的.
例如第五章的第三節.
這個羅馬書.

$^{201}$不但如此.
就是在患難中.
也是歡歡喜喜的.
因為知道患難生忍耐.
忍耐生老練.
老練生盼望.
盼望不至於羞恥.
因為神的愛.
藉著所賜給我們的聖靈.
就囂觀在我們的心裡.
這就是一個有義有信的門徒.
他面對世界種種的困難,問題.
他都能夠有一個新的力量.
就是神的愛在他內心的囂觀.
就可以有力量去面對.
所以你其實可以想像到.
保羅是在說.
有一班新的人.
新人類.
就是在主耶穌的救恩裡面的新人類.
仍然活在羅馬世界的那個世代.
到處都是偶像.
和帝國都要他們去拜君王是神.
但是這一班新人類.
跟隨耶穌的.
就可以因為這個信心.
就得到一個新的生命的更新.
就可以在地上活出一個.
完全不一樣的人生.
所以這個是由信開始.
因信而致於信.
「以致於」即是不斷透過信心.
去經歷上主的幫助.
那個加力.
這樣去面對人生不同的際遇.
我們開始在神的義這裡.
大家看看這個PowerPoint.
這一頁都頗重要.
上一次我們就停在這裡.
我們不夠時間講完.

$^{241}$和大家稍微揭一揭.
三段聖經.
新約都是講到哈巴谷書二章四節.
引述這段聖經.
但三次都不同的重點.
一章的十七節就明顯.
剛才我們講了.
義人因信得生.
這裡告訴我們.
因為神的義在這個福音上顯明.
怎樣顯明呢?.
就是這個義是本於信.
而致於信.
如經上所記.
義人必因信得生.
特別強調.
哪些人可以成為義人呢?.
就是信主的人.
因信主的人.
就可以得著成為一個義人.
或者說神的義就加在這些信人的生命裡面.
雖然和我們沒有什麼關聯的加拉太書.
但和保羅的羅馬書是有相當大的關聯的.
我們可以稍微揭一揭.
這個加拉太書第三章這一段.
三章由第十一節開始.
第三章是一個很重要的論點.
來駁斥加拉太教會.
如果他們要靠自己的好行為.
其中的好行為就是要受割禮.
因為外邦人沒有受割禮.
被猶太的信徒來迷惑.
說你們是二等的信徒.
因為你們沒有受割禮.
可能會帶來一種淡化了基督的救贖.
唯一的重要.
所以保羅知道這個危機.
就來責備加拉太教會.
他們好像畫蛇添足那樣.
加多一些東西.

$^{281}$以為做漏了.
其實不需要的.
信主耶穌就已經立即獲得那份信心.
那份神的義.
我們來看看保羅如何在這裡.
引用《阿巴菊書》二章四節.
引用的重點.
到底是在義人那裡.
在信心那裡.
還是另外的情況呢.
我們來看看第十一節這一段.
第十節再詳細一點.
「凡以行律法為本的.
都是被咒詛的.
因為經上記著.
凡不常照律法書上所記一切之事去行的.
就被咒詛」.
這是《申命記》的經文.
第十一節.
「沒有一個人靠著律法.
在神面前稱義.
這是明顯的.
因為經上說」.
即是《阿巴菊書》二章.
「義人必因信得生」.
第十二節.
「律法原不本乎信.
只說行這些事的.
就必因此活著」.
第十三節.
有一個對比.
「基督既為我們受了咒詛.
或者成為了咒詛.
救贖出我們脫離律法的咒詛.
因為經上記著.
凡掛在木頭上.
都是被咒詛的」.
這是《申命記》的經文.
二十一章二十三節.
第十四節.

$^{321}$「這便叫亞伯拉罕的福.
因基督耶穌可以臨到外邦人.
使我們因信得著所應許的聖靈」.
這裡的重點提到.
在這段聖經.
是和那個行律法的對比.
所以我們不需要守律法.
守不到.
也沒有人能夠守到.
保羅說.
因為基督為我們成為了咒詛.
就將我們帶離開律法的規範.
那個咒詛.
使我們因信主耶穌.
就可以得到這一份最寶貴的禮物.
就是應許的聖靈.
即是神自己住在我們裡面.
所以這一段聖經.
保羅的引述.
特別強調.
只要信.
不是靠手行為.
不是靠自己的功德換取.
是信主耶穌.
因為主耶穌已經成為了咒詛.
我們信就得到這一份應許的聖靈.
雖然他不是強調神的義.
但也是另一個角度去講救恩的恩典.
就是我們得到應許的聖靈.
住在我們的內心.
即是說我們已經脫離了手行為.
靠律法.
來得到神的那個稱許.
所以第二段保羅引述.
是強調信心的信靠.
不要靠行為.
就是加拉太書這裡的論點就是這樣.
羅馬書那裡特別強調.
那些人為什麼可以得到神的義呢?.
就是因為在神的運作裡面.

$^{361}$我們是信.
就立即成為了義人的身份和地位.
第十章的希伯來書.
就是新約第三段最後一段引述.
是哈巴谷書的.
也都揭露了.
雖然不是保羅寫.
不知道是誰寫的.
希伯來書.
但是這裡更加明顯.
是說到一個有信心.
相信主耶穌的人.
他應該是在信心裡面.
帶來他的生命是活的.
在地上是一個很活潑的.
屬靈生命的人.
第十章.
我們可以從第37節開始看.
可能要讀多一點.
第32這段開始.
當然想上文更加詳細.
不過我們只想引述關於.
這個「義人因信得生」這一段.
大部分人都說未必是保羅寫.
就當是不知名的作者寫.
第32節.
你們要追念往日.
蒙了光照以後.
所忍受大征戰的各樣苦難.
一面被毀旁.
即是說這些初代信主的弟兄姊妹.
面對無論是帝國羅馬的逼迫.
如果他們是猶太人.
受盡猶太教的冷嘲熱諷的叛教.
又或者是當地的人.
拜偶像的.
又會受到這些牽連.
影響.
一面被毀旁.
遭患難.

$^{401}$成了氣境.
即是說被人嘲笑.
挖苦.
被人排擠.
這就叫成了氣境.
好像一台戲一樣.
一面陪伴那些受這樣苦難的人.
因為你們體恤了那些被困所的人.
並且你們的家業被人搶去.
也甘心忍受.
知道自己有更美長存的家業.
所以你們不可丟棄勇敢的心.
全這樣的心必得大賞賜.
這裡提到我們要堅持到底.
不要看短暫的眼前.
那個苦楚.
那個困難.
因為所有歷代的信徒.
都是這樣去忍受這些逼迫的.
是的.
重點是不要丟棄勇敢的心.
全這樣心的人.
必定會得到將來的大賞賜.
接著第36節.
你們必須忍耐.
使你們行完了神的旨意.
就可以得著所應許的.
因為還有一點點時候.
那要來的就來.
並不遲然.
只是二人必因信得生.
再次引述哈巴谷書的二章.
這次連第三節也引述了.
「他若退後,我心裡就不喜歡他」.
接著結論就是第39節.
我們即是面對困局的弟兄姊妹.
我們卻不是退後入沉淪的那等人.
退後即是改變信仰.
因為太艱難而放棄信仰的人.
希伯來西亞的作者說.

$^{441}$我們不是這樣的層次的人.
我們是怎樣的呢?.
我們不是退後入沉淪的人.
乃是有信心以至靈魂得救的人.
即是特別強調.
二人因信得生的生命力.
有力量面對種種困難.
就是因為他內心有神的意義.
因為他相信將來會有永恆的上次結局.
與現在短暫的相差很遙遠.
所以他在地上,在苦難中.
印證他的生命是活的.
所以第39節是一個總結.
接著第11章大家都很熟悉.
就是引述舊約歷代的信徒.
在困難中怎樣展現他們的生命力.
就是這些信心偉人.
所以其實一點都不是偉.
只是有信心.
所以第39節就是說.
我們不是那些退後入沉淪的人.
是有信心.
即是被列入信心的行列.
現在寫這本書的時候.
他們是面對著現在的苦楚.
作者說我們其實是有信心的行列的人.
我們將來就會靈魂得救.
而且將來還會有上次的.
然後作者列舉一些已經得到上次的信心的人.
逐個數.
由阿當,以諾.
一直這樣數.
羅亞.
接著一直到亞伯拉罕.
一直這樣數下去.
就是說我們是不是那些信心行列的人.
現在面對的困難.
我們都要堅持.
這就是哈巴谷書.
全卷新約裡面.

$^{481}$全部新約的27卷書裡面.
只是引述了三次.
所以如果你說第三卷喜伯來書都是保羅引.
可能也對也說不定.
透過這個內政.
當然這樣比較牽強.
但是你看到運用哈巴谷書.
都有這麼巧妙地將二人.
和恩信.
和得生三樣東西.
都在三卷書裡面.
淋漓盡致地講解.
我們回到羅馬書第一章這段經文.
留意神的兒在這裡.
很快就做了一個首尾呼應.
保羅都很巧妙地.
將這個首尾呼應神的兒.
加給了什麼人呢?.
你可以說一開始17節.
就是那個答案.
萬民只能夠靠恩信稱義.
得到神的兒.
第三章的21節.
就說神的兒是怎樣的運作呢?.
中間有一個大的括弧裡面是什麼呢?.
就是萬民都是活在罪中.
一章18到三章的20節.
一大段的聖經.
是在講列舉人類種種的罪行.
種種的不義.
最嚴重的就是否認神的神聖.
不肯榮耀神.
所以透過這個括弧裡面的種種的指控.
我們基督徒只需要剛剛相反.
就可以成為神的很大的祝福的新人類.
就是說我們一信主.
就必須盡快學習怎樣去榮耀神.
因為一開始就提到這些不乾不義的人.
明知道神的永能神聖.
但是故意的去否認.

$^{521}$去漠視神的榮耀.
然後自己的心靈就昏暗去拜各種各樣的偶像.
或者更重要的就是榮耀自己.
就變成更加遠離神的救恩.
和遠離神的榮耀,管治和牧養.
所以你看到這個括弧.
阿保羅多巧妙地將萬民都需要稱義這件事.
一開始就已經提出這個標題.
因為神的義就在福音上顯明.
接著講了一大段之後.
三章二十節大家懂不懂得背?.
我們斷定了.
所有人都是在罪當中.
沒有辦法得到自己的拯救.
所以凡有血氣.
即萬民,各族.
無論是希利尼人,法外人,聰明人,愚處人,猶太人,外邦人.
凡有血氣的沒有一個因行律法能夠在神的面前得到神的義.
為什麼呢?.
因為律法本是叫人知罪.
即神不單是憤怒.
神還用他的律法來規範地上的人.
所以我們不能夠靠自己的行律法.
就可以在神的面前稱義.
因為律法只是叫我們知道我們自己是罪人.
今次第二十一節很明顯有一個答案的轉向.
就是那個三括弧.
你看看如今.
保羅所說的如今.
就是基督的出現在這個世界救恩的成就.
如今神的義不是現在保羅寫信的時候的如今.
但如今神的義在律法以外.
所以不是在律法以內.
在律法以外.
因為人已經不能夠透過行律法得到神的義了.
所以神不是要將律法廢棄.
而是神在另外一個恩典裡面.
來將義賜給信的人.
所以這就告訴我們.
在律法的以外已經顯明了.

$^{561}$而且有律法和先知為證.
律法和先知即是舊約聖經.
有舊約聖經的啟示.
即是說神雖然這麼嚴苛.
但其實祂有恩典.
有憐憫的.
祂會將祂的愛賜給信的人.
去得到這份救恩.
所以神如何運作呢?.
這就告訴我們.
22節.
「就是神的義」.
有沒有留意?.
是很清楚的一個開始.
1章17節和3章21節.
就是神的義.
「因信耶穌基督,加給一切相信的人」.
然後再強調.
「並沒有分別」.
即是說不會大細超.
每個人不論什麼民族.
只要你接上信主耶穌的救恩軌道.
你就立即可以獲得神的義.
就是這麼簡單.
就是神在人類的社會.
在一個已經背棄了神存在的地球.
神自己將救恩.
由舊約已經應許的救贖.
透過耶穌基督去施行普世的救贖.
當然普世是要信.
萬民必須是因信才可以稱義.
不是耶穌死了.
每個人都可以得救.
不是的.
必須要信的人才可以得到.
這一份神的義.
所以你看到.
這裡有一個大的括弧.
開引號就是1章17節.
告訴我們.

$^{601}$人只需要認識神的義.
相信就獲得神的義.
接著一直說.
萬民到底欠缺什麼.
而得不到神的義呢?.
保羅用了三個具體的種類.
就是否認神存在.
表面好像無神論.
但其實保羅以前的年代.
沒有人是無神論的.
沒有人說神不存在.
因為很多的拜偶像.
不同的各樣的東西.
當然近代有唯物主義.
否認神的存在.
應該是七十八世紀.
開始無神論.
或者說神死了的神學.
以前的保羅的世代.
是不會有人是無神論的.
整個世代都是知道有神的存在.
不過是拜什麼神.
這裡讓我們看到.
保羅很工整地.
將恩順清義.
萬民可以得到生命改變.
就是說解決人類罪惡.
唯一的答案就是恩順清義.
什麼人需要恩順清義呢?.
就是我們今天要開始說的.
一章十八節這一段.
就是神的憤怒.
我預備的時候.
又順便反正羅馬書.
都是說神的憤怒.
說了兩次.
原來在保羅書信裡面.
全本新約有五次說神的憤怒.
我們可以揭開.
以忽所書的第五章.

$^{641}$和歌羅西書.
讓大家感受一下.
保羅看神的憤怒.
是怎麼樣的一件事呢?.
我們先看回頭.
看十八節的羅馬書.
五章的六節.
以忽所書五章六節.
大家無論在線的或者現場的.
都要學學揭聖經.
最好就是家裡有一本大一點的聖經.
和有一本筆記本.
記下來.
當然我會將部分給大家.
保羅的書信裡面有五次說神的憤怒.
說得很重的.
第六節.
不要被人虛浮的話欺控.
因這些事神的憤怒必定臨到那些背逆的人.
就是上文和下文所說的這些事情.
原來神的憤怒會臨到背逆的人.
背逆之子.
就是說他的生命充滿了對神的叛逆.
這個神的憤怒就會臨到.
接著歌羅西書三章.
當然大家回去看上下文會更加感受到那些.
我只是想引述給大家知道.
原來神的憤怒保羅也引述了五次.
第三章第五節.
所以要致死你們在地上的肢體.
就如淫亂污穢邪情惡慾和貪婪.
貪婪就與拜偶像一樣.
因這些事神的憤怒必臨到那背逆之子.
就是說很相近的經文.
不過歌羅西書再列舉一些內心裡面.
一個背逆的人的內心寫照.
這些X光照不到.
磁力共振或者是PET scan都照不到.
唯有神的道.
將人性的內心裡面的這些層次都描繪得很細緻.

$^{681}$其實都是人性來的.
另外就是貼薩萊尼加前書一章.
所以在保羅的事奉裡面.
他傳遞神的愛 神的義.
但是他同樣都說.
在他的時代都有很多敵黨神的人.
他都說這些人就是背逆的人.
這些人就是神的憤怒會臨到他們身上的人.
一章第十節.
貼薩萊尼加前書.
這裡很明顯是在說迫伯保羅的猶太人.
可能我們讀多一點.
一章第六節開始.
貼薩萊尼加前書一章第六節.
並且你們在大能之中.
說貼薩萊尼加的教會.
因為當時保羅在那裡四個星期傳道.
他突然間緊急要離開貼薩萊尼加.
他去了雅典之後.
很擔心貼薩萊尼加的弟兄姊妹.
會不會因為迫伯而有變動的心靈.
有軟弱.
所以就打發提摩泰回去.
保羅因為各種原因.
他要離開.
而且離開得很急.
不能夠有很詳細的解釋.
一來就怕貼薩萊尼加的教會有誤會.
以為他因為一些的迫害.
看到那些環境不順利就走了.
保羅不是這樣的人.
但是要寫這封信解釋給他們知道.
現在保羅就說這裡.
並且你們在大能之中.
蒙了聖靈所賜的喜樂.
領受真道.
就效法我們.
也效法了主.
即是稱讚貼薩萊尼加的教會.
甚至你們作了馬其頓.

$^{721}$和阿該亞所有信主之人的榜樣.
阿該亞就是哥林多的省會.
那個省的首都.
阿該亞省所有信主的人的榜樣.
留意第八節.
因為主的道從你們那裡.
已經傳揚出來.
你們向神的信心.
不但在馬其頓和阿該亞.
就是在各處也都傳開了.
所以不用我們說什麼話.
因為他們自己已經報名.
我們是怎樣進到你們那裡.
你們是怎樣離棄偶像歸向神.
要服侍這位又真又活的神.
等候他兒子從天降臨.
即是說貼薩萊尼加的弟兄姊妹.
是信得很清楚.
他們離開了拜偶像.
歸向真神.
等候神的兒子的第二次降臨.
就是他從死裡復活.
這位救我們脫離將來憤怒的耶穌.
將來的憤怒.
即是將來有一個末日的大審判.
信主的人就可以逃避將來的.
極大的神的憤怒.
所以保羅就這樣理解神的憤怒.
在一些地上背逆的人.
不肯悔改的人.
將來面對神的審判.
公義的審判.
所以我們稍微闊一點去看保羅的書信.
就會讓我們那個範圍.
那個眼界.
淑寧的眼光.
又會再闊一點.
保羅信將來是有末日的憤怒審判的.
所以我們記住.
我們信主都是逃避將來的審判滅亡.

$^{761}$而進入永恆的那個真道.
真理的天堂.
所以其實我們今天我們信主.
都真的要繼續在這個信裡面.
我們不要有偏差.
這個信裡面其實是.
包含著敬畏和榮耀神的.
所以如果我們只是說交換.
信啊信啊.
這個信就不是聖經所說的信.
是你自己演繹的信.
這個就要小心.
因為你看到這些經文.
保羅所說的信不是這麼廉價的.
好像你不信都不好意思.
不是的.
保羅所說的信是包含著.
對神的憤怒的懼怕.
對神的榮耀的尊崇.
和對神的旨意的實踐.
對神的愛的擁有的改變.
是一個很豐富的生命的套餐.
不只是一種像講價的.
這麼高要求的不信.
低一點我就信.
就不是這樣.
我們要小心.
不要自己以為的信就是.
認為這樣就是信.
我們留意這裡開始.
幫大家做一個圖像.
一張十七和三張二十一.
是一個很重要的括弧.
就是說神的兒就是這樣加給信的人.
而且是在律法以外.
就是說人絕對沒辦法靠自己.
可以獲得的.
哪些人才可以得到呢?.
就是認識自己在神憤怒底下的人.
沒辦法自救.

$^{801}$更加不是你個人.
根本我們一生來到這個世界.
就活在罪中.
我們就在神的憤怒底下.
但是只要相信.
我們就可以獲得這一份無條件的救恩.
所以你看到.
保羅就很刻意將十八節到三張二十節.
作為一個很重要的段落.
來去羅馬的教會去展述.
差不多已經涵蓋了第二張.
有一些很有道德行為的人.
和神的選民猶太人.
第一張就是否認神的榮耀.
而活在自己的虛妄的人.
差不多是涵蓋了整個地上所有的人類.
所以你看到.
這個結構.
一張十八節到三十二節.
內容就是萬民都欠缺神的義.
不能夠靠自己獲得神的義.
第二張就是我們下一個星期會講.
外表的義不是神稱許的義.
第三張一到二十節.
作為一個總結句.
所有的經文都需要在律法以外的恩典.
就是主耶穌的救贖.
我們開始一張十八節這一段的經文.
簡單地幫大家分析了.
一張第一到第十七節.
尤其是最重要的十七節這一節的經文.
和三張二十一節.
這個很重要的豁兀的神的義.
如何顯明.
一張十八節開始.
我們先稍微讀一讀經文.
第十八節開始.
原來神的憤怒從天上顯明.
在一切不乾不義的人身上.
即是說保羅講完十七節.

$^{841}$人類唯一的解救.
得到神的義的方法.
就是透過信主.
接著就開始講.
人類現在在神的義以外.
其實就是在神的憤怒.
因為神的義也包含了他的公義的憤怒.
所以神的義不是我們要就行不要就行.
是必須我們要恩順清義.
去得到神的義的這份恩典.
因為保羅告訴我們.
人類現在的光景.
人民都是在罪惡底下.
或者說在神的義的以外.
是哪些人呢?.
就是那些行不義.
阻擋真理的人.
第十九節.
神的事情.
人所能夠知道的.
原顯明在人心裡.
因為神已經給他們顯明.
去到這裡.
開始講創世.
講神的這位造物主.
是可以讓人認識尋求.
去到第八章第十一節.
又再講創造.
又再講亞當的墮落.
又再講萬民在虛空底下.
切忘可以得到解救.
更加完整地將人類面對的困局.
去剖析得很淋漓盡致.
這裡就開始告訴我們.
人原先是可以被設計,被創造.
是認識造物主的.
但是什麼原因令我們不認識呢?.
接著保羅就解釋了.
第二十節.
原先是這樣.

$^{881}$但實際就不是這樣.
自從造天地以來.
神的永能和神性.
是明明可知的.
雖然眼不能見.
明顯不是肉眼見到神.
肉眼見到也不會相信.
祂的永能,祂的神性.
是要透過我們內心的靈去認識.
去俯伏,去尋求.
最重要的是去榮耀我們這一位主.
我們現在信了主.
我們不在神的憤怒下.
我們可以探求,認識,尋找.
這位永能和神性的主.
可以對祂的存在.
在我們生命裡的帶領.
其實我們就可以去尋找了.
所以我們記得.
不要活在一個沒有祈禱的日常生活.
基督徒可能回到禮拜堂.
才跟著牧者,公土.
或者在主日崇拜之後的牧土.
才去祈禱.
其餘時間根本不想要祈禱.
不覺得自己需要祈禱.
其實保羅告訴我們.
原來是可以認識,可以尋找的.
不過是罪惡令人否認和拒絕他的榮耀.
如果我們信了主.
即是剛剛相反.
我們信了主.
只需要我們每天學習榮耀.
尋求祂.
所以其實每天的靈修安靜.
這個不是粗言穢語.
就是我們每天認定神是掌管萬有的主.
是我們的生命的主.
我們每天去尋求祂.
我們是可以尋求到.

$^{921}$但這裡告訴我們.
原來可以尋找到神的人.
居然是不可以的.
藉著所做之物.
就可以了解祂的存在.
叫人無可推諉.
這個相信不是一種知識的層面.
是心靈上對神的榮耀敬拜.
俯伏的跟隨.
對祂的敬畏.
欠奉.
否定.
拒絕.
所以因為.
他們雖然知道神.
似乎是保羅時代的讀者群體.
或者在保羅以前的世代.
同樣也有這一類人.
這裡說.
雖然知道神.
卻不當作神榮耀他.
所以絕對不是無神論.
是知道神存在.
故意拒絕.
即是說存一個很邪闢的心拒絕.
用自己的方法去尋找.
要降低神的永恆.
那個神性.
就用一些降級神的屬靈的價值.
這裡告訴我們.
他們無知的心就昏暗了.
第22節.
自稱為聰明.
反成了愚拙.
將不能扭壞之神的榮耀.
一直在說神的榮耀.
即是說.
祂本質的存在.
做物主本來應該得到的尊崇.
人就拒絕去敬拜祂.

$^{961}$拒絕承認祂是掌管萬有.
掌管生命的做物主在.
所以自稱為聰明.
其實就反成了愚拙.
即是說.
故意想祂的不存在.
不要干涉我們的個人空間.
大家知道.
過年那些親朋好友.
姨媽姑姐.
阿明仔.
何時結婚?.
這些個人空間.
個人的空間.
很冒犯.
現代人就很明白這些東西.
還沒結婚?.
賺多少錢一個月?.
做什麼工作?.
這些是很冒犯的個人空間.
我們現代人就大家明白了.
我們不會討論這些東西.
但是過年那些姨媽姑姐.
因為一年見一次.
一定問的.
明明問過年.
每年都問一次.
現在在讀什麼書?.
其實已經做了很多年了.
都還是這樣問的.
這些就叫個人空間.
所以其實這裡的思念變為虛妄.
無知的深昏暗.
其實真的不是慢慢變得愚蠢了.
而是不想神侵入每個人自己的個人空間.
希望以為這樣.
就可以讓神的榮耀.
不存在我們自己的生命裡面.
所以就成為了一個愚拙的人.
將不能夠扭壞的神的榮耀.

$^{1001}$就將它變成了偶像.
彷彿扭壞的人.
飛禽.
走獸.
昆蟲的樣式.
就是說.
將這些沒有殺傷力的.
成為神.
你可以駕馭它.
你可以操控它.
還要它保護你.
保佑你.
明明它不行的.
你都要它這樣做.
來讓自己的內心好過一點.
我有拜神的.
不過不是拜獨一的榮耀神.
這樣的描繪.
都是很生動的反映.
保羅時代.
拜眾多的神.
譬如亞底米女神.
我們在爾弗所見到的柱子.
那個遺址.
就是將人的無限的想法.
活現在偶像裡面.
譬如.
赫爾瑪斯.
大家都知道是那個名牌.
赫爾瑪斯那個女神.
其實是傳遞信息.
希臘和羅馬都有很多這些假神.
所以你會看到.
保羅時代.
是有很多這些拜不同的神的.
那些信仰,宗教.
其實就是要讓人.
有空間去管理自己.
不要被神介入.
所以就拒絕神的存在.

$^{1041}$所以第24節就開始.
就說到.
這裡很特別.
保羅用了三次的任憑.
任憑就是Let go.
你喜歡怎樣就怎樣.
就是說.
讓人的無知.
繼續無限的發展.
這個好像表面上.
有很多個人空間.
因為神都不再干預.
但其實實際上.
這個是一個刑罰.
是一個令你自己活在自己的虛荒.
虛妄.
和那個無知底下.
拜的神子.
幫不了那個內心的潔淨.
和得到真正的與神和好的平安.
所以第五章.
立即就說.
聖經帶來的第一個祝福.
就是與神和好.
所以一直都沒有說和好的.
一直都說.
從第18節開始.
一直都說.
要神保持距離.
離開我的心靈越遠越好.
不要和我有關係.
其實是想自己不和神有關係.
猶太人一樣.
所以第二章也說猶太人.
你可以想像到.
那個任憑.
你會看到那個下降.
我們說過.
我們看到.
那個問題.

$^{1081}$是和罪物不可分.
不論你是外邦人.
猶太人.
有沒有道德的.
或者只是享樂主義者.
每個人都是和罪有份.
普世的福音就是要解決.
罪這個普世的問題.
唯一的答案.
就是萬民恩信.
可以得到與神的和好.
所以保羅是在解決.
人類種種的罪惡.
我做了一些筆記.
拒絕看到神的存在.
不榮耀祂.
是一個主動的.
否定祂的榮耀.
不榮耀祂.
所以弟兄姊妹.
不要說我們不懂祈禱.
是你與神之間的生命的愛的關係.
所以我們只需要榮耀神.
我們就會建立到這一份親密和好的關係.
是你與神之間的和好關係.
每天都這樣繼續進行.
所以真正的背影.
就是不榮耀神.
明知道祂是眾物主.
不榮耀祂.
不榮耀祂之後.
在心靈裡面.
變得虛妄.
變得有恃無恐.
有恃無恐.
這樣去做一些神喜悅的事.
當然要找一些東西.
來滿足自己內心的空虛.
去拜偶像.
有三個孕鵬.

$^{1121}$我們站在24節.
這個PowerPoint有的.
回家就會給大家.
其實不用拍了.
24節的解構就是.
所以神孕鵬他們.
請著心裡的情慾.
行污穢的事.
以致彼此玷污自己的身體.
他們將神的真實.
變為虛謊.
繼續這樣不斷說.
拒絕祂的榮耀.
祂的真實.
祂的真理.
出現了.
去敬拜.
事奉受造之物.
敬拜然後事奉.
是兩個字.
拜完還要服侍.
你可以想像到.
那種的落力參與.
不過投射的.
是在這些假神裡面.
滿足自己內心的私慾.
接著就不敬奉那造物的主.
主乃是可稱重的.
直到永遠阿們.
稍微揭一揭羅馬書.
又有三次說到.
直到永遠阿們.
羅馬書十五章.
羅馬書十六章都有這樣說.
接著保羅書信.
還有其他的經文也是這樣說.
所以保羅說到一些激動的地方.
他就說我們人類.
為甚麼會這麼可憐呢.
將神的真實變為虛方.

$^{1161}$虛方即是說.
都不知道祂是否存在.
當祂不要存在.
接著就拜一些根本不存在的.
自己製造出來的木頭.
去敬拜這些受造之物.
不敬奉造物的主宰.
接著保羅就說.
我們的主是可稱重的.
直到永遠阿們.
所以我們大家都學習.
每天起床祈禱.
都好像保羅那樣.
從心裡面對造物主的讚歎.
有日能夠有生命起身.
可以去刷牙.
準備一日的工作.
或者今天你是去迪士尼的.
或者今天你是上飛機去看櫻花的.
都是一個感恩.
今天你很忙的工作.
或者有些很艱難的事務.
都是一個感恩.
但是對這位造物主.
祂是可稱重的.
直到永遠阿們.
保羅寫著寫著都很激動.
這麼多人原來.
這麼拒絕神.
保羅就知道.
自己的存在.
對造物主的榮耀.
所以原來很健康.
我們讀一下羅馬書.
保羅很有感情.
說到那些拜偶像的人.
這樣來降低神的榮耀.
神的真實.
很看不過眼.
那些是偶像.

$^{1201}$居然我們將神的榮耀降格.
不去敬奉這位造物主.
主乃是可稱重.
直到永遠阿們.
這樣來再一次.
其實不需要這樣說.
都這樣說的時候.
你就可以想像到.
這位神對救恩改變萬民.
充滿了激動.
充滿了熱情.
充滿了欠債的心靈.
所以弟兄姊妹.
我們真的要慢慢學習.
看到原來人的信主.
是最大最寶貴.
所以無論你擇假意也好.
真心也好.
報道會真的要參與.
真的要祈禱.
真的要守望.
去爭戰.
因為真的要搶救靈魂.
不知道撒旦在聚會裡.
如何令到來聽的人.
充滿了疑團.
充滿了對講員的藐視.
等等.
但如果你的祈禱.
捆綁著撒旦.
讓他的心安靜.
客觀去聆聽.
可能那一刻就改變了.
他得到生命的拯救.
所以要令到那些.
不信有榮耀.
不信有神的.
真實的人.
其實他們是在拜他們自己的.
受造之物.

$^{1241}$自己製造的.
神祇是假的.
是得不到真正永恆的福樂.
但是.
主是可稱頌.
直到永遠的.
這是第一個任憑.
心靈拒絕真理.
真實.
榮耀.
去敬拜.
受造之物.
這是第一個任憑.
接著第26節就第二個任憑.
到第28節就有第三個任憑.
因此又有另外一種任憑.
神任憑他們.
放縱可羞恥的情慾.
他們的女人.
把純性的用處.
變為逆性的用處.
男人也是如此.
是在說.
male和female.
是男性和女性.
就不只是說.
男人,女人.
是說人類的結構.
只有男性,女性.
就是兩個性別.
就不會.
這樣去混亂.
是不同的功用.
所以這裡明顯說.
那個所謂逆性.
就是違反.
創造的定律.
在保羅的年代.
如果大家看文獻.
歷史都知道.

$^{1281}$或者大家有看過.
《亞歷山大》那套電影.
很多年前播放的.
都是三級電影.
因為.
在說亞歷山大的年代.
亞歷山大大帝的年代.
就是希臘末期.
然後到羅馬帝國.
那些皇帝.
大部分是戀童的.
就是和男人.
有性行為.
歷史都.
印證了這些事實.
你可以看到.
保羅在他的時代.
其實他都清楚.
那些帝王的.
雙方拜德.
他都說到.
人類因為.
有意將神的真實.
改變為虛方.
去滿足自己.
那個.
有平安 有渣拿.
又是受他操控的.
所以一定是拜偶像.
帶來的就是.
連人性本來的創造.
的秩序都違反.
違反創造的.
倫理秩序.
就不是神的期望.
所以這裡就告訴我們.
女性本來在婚姻以來的性生活.
就是神的配合安排的.
已經違反了.
這種倫理.

$^{1321}$創造.
所以導致保羅就列舉了一些.
其中一項.
不是唯一的罪大惡極.
因為下面都有很多.
離開神之後的人性.
扭曲.
不協調.
內心的惡毒.
說虛假的話.
等等.
所以不可以.
只針對同性戀的罪.
就說是罪大惡極.
審判.
其實下面一樣.
所以這裡整個大段落.
一章十八到三章二十節.
萬民活在神的憤怒下.
因為萬民被罪操控.
在罪的操控下.
人是故意不懂得去榮耀神.
這是問題的所在.
所以我們不要特別高興.
一章25至26兩節.
特別是27節.
男人都是這樣.
棄了女人純性的容處.
欲火攻心.
彼此貪戀男和男.
恆可羞恥的事.
就在自己身上.
受者妄為.
當得的報應.
福音派.
特別早期.
都很強烈針對同性戀的罪.
很嚴重的去指出.
去定罪.
這是其中一個例子.

$^{1361}$所以這裡不是唯一一樣.
保羅在整個羅馬書這段經文.
說到罪人.
這就是面對神極大的憤怒.
因為還有下面的.
這是其中第二個的任憑.
即是說.
那個下級.
首先是否定神的榮耀存在.
真實.
然後製造偶像.
製造偶像之後.
任意的.
操縱偶像.
要對自己有任何的利益.
包括自己有一個很大的個人空間.
誰都不可以侵入.
這樣的空間.
製造了人自己去設計很多.
違反了創造倫理的生活行為.
這就讓我們看到.
不要只針對同性戀.
第28節.
他們既然故意不認識神.
已經是第三次說了.
是故意的.
是否定神的存在.
所以基督徒其實很簡單.
我們每天有意識地在我們的思想.
在我們實際的生活.
我們認定神.
掌管著我們的生命.
和有一個美善的計劃.
我們要去尋找.
這是第12章的第一節.
差不多是整個1-11章的總結句.
所以弟兄們.
我以神的慈悲.
勸你們.
將身體獻上.

$^{1401}$當作活祭.
有意識地去敬拜神.
就是一路回應.
基督徒得到這麼大的恩典.
我們要有意識地去榮耀神.
在我們的生活裡.
承認祂的主權.
將身體當作活祭.
獻上給神.
是聖潔.
是神所喜悅的.
我們這樣事奉.
不是事奉那些偶像.
是神所喜悅的.
每天主動在神的喜悅裡.
去榮耀祂.
所以基督徒嚴格來說.
很容易得到神賜的滿足.
所以你會看到這一段.
其實這樣的描繪.
人類只會自相殘殺.
人只會不斷自私下.
要不斷擴大自己的空間.
沒有任何的.
連神都不可以進入.
更何況其他人.
只會不斷地打打殺殺.
第28節說.
神突然故意不認識.
神就任憑他們.
第3個任憑.
傳一個邪闢的心.
即是說那個心靈開始出現了.
沒有了那種美善.
沒有了那種和睦.
沒有了那種和諧.
出現了一種對各樣的事情.
都是行不合理的事.
邪闢的心.
然後就裝滿了各樣不義.

$^{1441}$心靈.
記得量不到多深.
那個心不是用尺去量的.
所以裝滿了.
還沒有滿滿.
裝滿了什麼呢?.
各樣的不義,邪惡,貪婪,惡毒.
滿心都是嫉妒.
整個心滿滿都是嫉妒.
你想想.
如果有一個人包了這麼多東西.
那個人是怎樣的樣子呢?.
去看看上是怎樣的.
先生祝福你.
你長命百歲.
你可以賺很多錢.
你的事業線很厲害.
但是.
這個人可能是神的憤怒下.
面對將來審判的人.
因為心靈是完全量度不到的.
這裡告訴我們.
滿心都是嫉妒,兇殺.
憎恨,詭詐,毒恨.
又是殘毀.
所以你有毒恨,貪婪,邪惡.
自然你就會不喜歡其他人比你好.
你要「殘毀」.
「殘」即是「慚嘴」.
即是很喜歡說人壞話.
「殘毀」即是要詆毀.
你不是說人的好話.
你知不知道是哪個牧師多好.
多敬虔.
我們很少這樣說.
我們很少稱讚敬重的牧者.
或是敬重的弟兄姊妹.
很少人這樣說.
但我們很多時候聽到的是「殘毀」.
你想想用一些「滋滋滋」的說話.

$^{1481}$去毀憾人.
接著就是背後說人.
怨恨神.
又是武漫人.
狂傲,自誇.
所以內心裡的自我中心.
自我保護.
自我的行仗.
「獵造惡事」.
違背父母.
無知,背弱.
無親情.
不憐憫人.
他們須知道.
神判定行這事的人.
是當死的.
然而他們不但自己去行.
還喜歡別人去行.
即是說.
有這種奇怪的傳染性.
相信疫情之後.
每個人戴口罩.
已經很有意識地.
明明都沒有甚麼.
在戶外都戴著口罩.
30度都戴著口罩.
好像都怕那個傳染.
但我們比較少怕這些傳染.
保羅告訴我們.
如果離開了神的榮耀的人.
就很喜歡和這些人來往.
不但自己去行.
還喜歡別人去行.
所以原來我們的朋友.
我們的敬虔的弟兄姊妹.
如果我們不是彼此有互相扶持.
大家學習去敬虔的道人.
敬畏神.
只會落在這些光景.
其實就是第三個淫貧裡面.

$^{1521}$就會充滿了內心.
沒有受到聖靈的管理.
就結出不了聖靈的果子.
就會有這些惡果.
是心靈的滋生.
因為我們是一個不斷活著.
心思是不斷活著的人.
除非死了那天.
所以這些事情.
是不斷進入我們的心靈.
或者是來自魔鬼.
更多是我們已經否認神的榮耀.
而慢慢就裝滿了這些不好的.
內心的惡果.
所以真的很需要求神幫助我們.
我們去看羅馬書.
《恩訊稱義》的時間.
保羅告訴我們.
原來《恩訊稱義》是這麼寶貴.
是這麼有神的大能.
將這些活在種種的罪惡昏暗的人.
都可以拯救出來.
而且不是靠自己改好.
不是靠你行律法.
是在律法以外顯明出來.
就是因信主耶穌.
就得到這份神的義.
就得到聖靈進入我們的內心.
就會讓我們慢慢看到神的榮耀.
所以原來我們每天對神的榮耀.
有一份好像每天自己起身的朝聖.
我們生命就一直有成長.
所以其實很奇妙的.
救恩真的改變了人.
改變了很大的.
所以是很大的一份恩典榮耀.
所以三個的任憑.
24到32節.
這麼短的經文裡面.
就已經不斷地講到.

$^{1561}$我們是不會有一個中立的心靈.
我又不拜偶像.
不過我又不拜神的.
但我的心是很乾淨的.
似乎應該不行了.
我們離開了神之後.
故意否定祂.
神的任憑就是我們去敬奉受造之物.
將有限的.
不是神的.
當做是神.
最終其實骨子裡是我們要操縱我們自己的心.
是我們要操縱這些的神.
要祝福我這個人.
但26,27節就是.
慢慢地放縱我們內心的情慾.
心靈會更加地邪闢.
真的每樣東西都看不順眼.
內心裡面自己不斷地湧現一些對人的不好.
自己又不開心.
又控制不了.
就一直擴大這個範圍.
喜歡自己走.
更加喜歡其他人和你一樣走.
其實就是神的任憑.
所以我們都要很小心.
所以神的憤怒.
不一定是地裂開的.
我問過大家一個問題.
很早的第一堂的時間.
在舊約有沒有哪些是神的憤怒呢?.
例如洪水滅世.
就是一個很大的憤怒.
就是神那個.
我怕大家都不熟.
創世記第六章.
第一節.
當人在世上多起來又生女兒的時候.
神的兒子們看見人的女子美貌.
就隨意挑選取來為妻.

$^{1601}$耶和華說.
人既屬乎血氣.
我的靈就不永遠住在他裡面.
然而他的日子.
還可以到一百二十年.
那時候有偉人在地上.
後來神的兒子們和人的女子交合生子.
那就是上古英武有名的人.
耶和華見人在地上罪惡很大.
終日所思想的盡都是惡.
耶和華就後悔做人在地上.
心中憂傷.
這是一種描述.
對神的那個.
不是他的後悔.
是他慨嘆人為何會失落得這麼嚴重.
來拒絕神的榮耀.
來將自己降格成為一個好像沒有神的受造物.
所以神就有一個刑罰.
就是要讓洪水來毀滅世代.
第十一節.
世間在神面前敗壞.
地上滿了強暴.
見事敗壞了.
所以你會看到保羅說的《羅馬書》.
其實那個藍本都是根源於他對創世記的了解.
離開了神之後.
充滿了強暴.
人.
即是說世界都充滿了不安.
很多的動亂.
人和人之間的詭詐.
虛謊.
毀謗.
各方面的東西.
所以你留意到.
見到人敗壞.
凡有血氣的在地上都敗壞了行為.
神就對羅馬說.
凡有血氣的人.

$^{1641}$打的盡頭已經來到我面前了.
因為地上滿了他們的強暴.
我要把他們和地一併毀滅.
這就是神的憤怒.
還有其他的在舊約都有的.
例如可拉一黨的人.
來去窒道和攻擊摩西亞倫.
神就叫可拉的叛黨.
其實是利未人來的.
其實都是祭司的家族.
就叫他們聚集在一起.
在利未記的十六章.
如果沒有記錯.
民數記十六章.
就全個地裂開.
那些可拉的叛黨.
全部掉到地下.
然後地合起來.
這是一個很大的神蹟.
在曠野的神蹟.
就是一個這樣的憤怒.
所以保羅在新約.
沒有說到一些很即時的.
所以我們不要直接套用.
愛滋病就是神對同性戀的憤怒的審判.
可能是的.
但是我們又不需要這麼快的對號入座.
因為這裡提到的任憑.
是人性已經去到一個腐敗.
不能夠自己自我救贖.
神的任憑是讓人看到自己的需要.
所以福音是解決人.
不能夠自己拯救自己的答案.
那個真正的出路.
所以我們在傳福音的過程.
我忘記了聽新穆斯.
我們在那個傳道裡面.
怎樣可以製造這個list.
令到他感受到自己的需要.
面對神的憤怒和任憑.

$^{1681}$他怎樣能夠進入那個救恩裡面呢?.
每個人都有不同.
他到底在哪一個光景的任憑底下呢?.
這個需要某程度上的友誼的報道.
帶來他的抗拒心沒有這麼嚴重.
沒有這麼大.
但也不代表他會很順從.
那個好朋友的善意的鼓勵.
但總是我們要相信.
這個福音是神的大能.
可以改變那些信的人.
所以我們只能夠用一些很原始的.
為這些人守望,禱告.
在一路的困難裡面都陪伴他面對.
讓他感受到這份真誠.
一份他不應該去抗拒的友誼.
讓他看到原來信耶穌是一件這樣的事.
可能他就會有一個改變.
這個很高的藝術.
是我們整個教會所有人都要一起參與.
萬民都需要福音.
所以我們一起去參與.
我們可能在不同的部份去做那個工作.
最終是聖靈讓他進入那個真理.
他看到真光.
看到自己的需要.
神光照他.
他有一個悔改的心.
他去尋求十章九節.
口裡承認.
心裡相信.
就可以得救.
凡求告主明十章十三節.
就必得救了.
所以一定要進到那個求告裡面.
所以這裡很寶貴.
就是講到萬民需要恩信清義.
在這一段經文裡面.
今天就差不多了.
我們講到這裡.

$^{1721}$所以基督徒的逆轉.
就是每天學習怎樣去榮耀神.
讓神在我們的生活介入我們裡面.
每天學習主事可稱重的.
直到永遠.
阿們.
每天我們都能夠.
當然不是念口往.
真是我們在我們的生命裡面.
看到.
保羅很寶貴.
原來這個世界是敬拜受造之物.
不敬奉造物的主.
但保羅現在將這位造物的主.
他的救恩.
有系統地寫出來.
所以他內心裡面.
是從心裡認定.
主事可稱重的.
直到永遠.
阿們.
保羅就是這樣在他的人生裡去事奉.
所以他不會有一種憤世嫉俗.
尼祿就快要逼迫他了.
就快要殺他了.
他都不會因為這樣.
來否定造物主在他生命裡面的帶領.
和他盡他的人生.
去將這位造物的主.
去描繪出來.
讓後世的人.
讓羅馬當時修信的讀者.
都認識.
原來我們投放在永恆裡面.
為萬民信主去祈禱守望.
永遠是最大最崇高的.
我們的人生.
所以信了主之後.
不是說不應該有我們個人的空間.
個人的休息.

$^{1761}$有時旅遊.
但總的來說.
我們是要大家一起參與.
在教會的傳福音的行列裡面.
否則我們又慢慢.
不知不覺.
疏離了教會的模樣.
和尋求神榮耀.
所以基督徒剛剛的逆轉.
我們不要學這些.
故意不認識和否定他的存在.
這樣的行徑.
今天就到這裡.
我們一起祈禱.
再一次感恩.
天父給我們都有系統.
將十八到三十二節.
就是萬民活在罪和否定神的存在.
這樣的生活光景.
天父幫助我們.
讓我們已經信主的.
特別我們看和去研讀羅馬書的弟兄姊妹.
有初信有信主內的.
有已經參與很多年事奉的.
我們一起去學習.
每天真是有意識地.
我們去敬拜這位創造的主宰.
認定你的榮耀和你的管治.
使我們能夠在生活裡.
不再讓憂愁.
讓重擔霸佔我們的心.
而忘記了凡事交託.
謝恩禱告.
沿主幫助我們.
可以算是有系統地.
將十七到一章的結尾去講解.
還有第二章第三章.
去到二十節的.
講有道德的人和神的選民的猶太人.
為什麼都是活在罪中.

$^{1801}$求主讓我們能夠在下兩次去聽的時候.
大家一起去反思.
我們這樣仰望禱告.
奉耶穌基督的名.
阿們.
\newpage



\section{羅馬書}
\label{sec:7C1Uqd4_AYQ}
\textbf{2025-04-03 周四晚查經班《羅馬書(6):分題:神的義(五)》 講員:曾裔貴牧師}
\newline
\newline
連結: \href{https://youtube.com/watch?v=7C1Uqd4-AYQ}{\texttt{https://youtube.com/watch?v=7C1Uqd4-AYQ}} ~~~~ 語音日期: 2025-04-05
\newline
\newline
\hyperref[sec:nWNWvRiUB3k]{\small{< < < PREV SERMON < < <}}
~
\hyperref[sec:index_chronic]{\small{[返順時目]}}
~
\hyperref[sec:index_scriptual]{\small{[返順卷目]}}
~
\hyperref[sec:JU6MXStfi9g]{\small{> > > NEXT SERMON > > >}}
\newline
\newline
$^{1}$一起同心祈禱.
感恩天父再一次給我們有機會.
可以在羅馬書上深入研讀.
盼望我們能夠對每一段的講解.
上下文的處理.
和保羅要帶出的神學教訓.
我們都能夠掌握.
更加重要的是.
我們可以透過聖經的啟示.
去認識同神的關係.
以及如何在我們有生之年.
成為一個活祭.
呈獻給基督.
在地上為主作見證.
求主幫助帶領我們.
有些艱深難明的.
求你都讓我們大家一起同心.
一起願意虛心地學習.
一起思想,討論.
我們仰望感恩禱告.
奉耶穌基督的名.
阿們.
好,我先開大Power Point.
來到第六講.
其實我們才真正進入神的義的討論.
你可以想像到一個很大的課題.
如果你再用比較闊一點.
其實是由一章的十八節.
其實應該延伸到第五章的結尾.
都是在講這個分題.
就是神的義.
當然第五章有兩個大段落.
第一個段落一到第十一節.
應該是講關於我們信了主.
恩順清義的人.
享受甚麼身份帶來的福樂.
但突然間第十二節.
去到第五章的結束.
突然間講述了.
阿當與耶穌基督的關係的對比.

$^{41}$是很特別的經文.
所以到時我們再去到那些經文再探討.
所以我們是由第十八節一章.
一直在講神的義的分題.
當然總的寫作動機.
就是保羅很想羅馬教會.
特別是一個首都的教會.
能夠支持他們向萬民宣教.
很感恩昨晚我們三堂的兄姊妹同工.
一起在線.
有Chris 余建峰弟兄.
和Mustafa這一位伊朗的宣教士.
因為他已經被囚禁了.
他的太太在線上分享.
他們這幾個月的遭遇等等.
都是一件很美麗的事.
我們一起在線上透過Zoom.
可以一起連線.
可以聽到消息.
在土耳其的現場.
他們透過Zoom分享他們的近況.
是很感恩的.
都是我們對宣教的事公的支持.
所以保羅很想羅馬教會.
所以第十四十五章.
就是教會的必須合一.
限制很大.
因為有猶太人.
有外邦人.
有來自不同的想法.
不同的一些傳統.
有守聖日.
守安息日.
有些說不用守的.
有些是吃一般的食物.
有些猶太人不潔不可以吃.
有很多限制.
但是在更大的神的角度.
宣教的事業.
這一切都要放下.

$^{81}$所以保羅說一定要教會要合一.
才能夠支持宣教士保羅.
他們去傳福音.
所以你可以看到.
宣教是一個最大的羅馬書課題.
是一個總綱.
在這個總綱下.
我們再思想怎樣信耶穌.
為何要信耶穌.
甚麼人是信不到耶穌.
有很多這些討論.
在這個救恩裡面.
我們怎樣看自己舊有的舊人呢.
第六,第七章也有涵蓋.
第八章也是釋放.
生命,聖靈的律.
住在我們裡面.
所以每一個環節.
都是緊扣著福音要向萬民宣講.
所以我們千萬不要在這麼大的書卷裡迷失.
所以大家一定要多讀.
多對聖經的經文掌握.
串通它.
你就可以慢慢浮現一個很重要的.
整個脈絡的圖畫.
我們回來看這個分段.
一個小的分段就是.
一章十八至三章二十節.
就是萬民在神的憤怒之中.
必須透過耶穌.
才可以得到神的義.
或者叫做「恩順清義」.
所以由一章十八節.
到三章二十節.
都是在講神的憤怒.
在類似三類的人當中.
我們上個星期提過.
一章十八節全章.
是講一些故意以為自己是無神的人.
因為拒絕神.

$^{121}$不承認祂的存在.
不尊崇祂等等.
所以神就有三個淫膨.
這是我們上個星期也講過.
作為一個印證.
這些以為沒有神.
活在一個不信神的世界.
心靈裡面的人.
他們不斷降低自己的靈性.
出賣自己給偶像.
去拜偶像等等.
所以這也是萬民的寫照.
很難不拜偶像.
本身又很好.
第二章剛剛也試探討這個問題.
真的不是去拜偶像的人怎麼辦呢?.
所以第二章我們今天就涵蓋.
在當時希羅的世界.
希臘羅馬的世界.
一些哲學家.
一些哲學的導師.
都是相當流行.
其中有一派叫斯多瓦.
或者斯多葛.
很崇尚道德.
很崇尚這個世界是有理性的.
是存在於有結構的.
當然不是有違格的神的創造.
所以他們覺得是這個理想的世界.
只要人做好就可以了.
也是靠自己的自然主義.
所以第二章本來就是針對這樣的人.
不過留意我們今早有姐妹問到.
我覺得是很好的.
先給大家有一點點的範圍.
當然大家都忙.
基本上都沒有時間預備.
如果你有時間是一個福氣.
你躺在桌面上開杯咖啡.
去讀一下羅馬書.

$^{161}$很靜心去讀.
是很有幫助的.
譬如你在Pacific.
拿一本聖經去看.
明顯的.
都不用太認真看.
都留意到一章18到一章的結尾.
這32節不斷提他們他們他們.
說到眾生所有的人.
但是突然間去到第二章.
全章是針對你個人的單數.
這樣去責備保羅.
這又是一個結構上的轉變.
所以我們又要留意.
保羅這種結構的轉變.
若有所指.
又變成一個神學的指涉.
留意現在你看到的Power Point.
三章21節是明顯的一個很重要的轉變.
就是如今萬民虧欠神的.
而在神的陣路下的人.
已經被驗證.
是沒有辦法逃避神的審判.
但是現在什麼是福音呢?.
就是三章21節很清楚的一個轉向.
所以由三章的21節.
嚴格應該是五章.
不過去到四章結尾.
結構性就會更加強.
因為第五章有另外的講解.
是告訴我們什麼人可以得到神的義.
萬民都是欠缺神的義.
萬民都是在神的憤怒下.
所以三章21節是一個很重要的轉向.
盼望大家對結構上的Contrast(相反).
大家能夠拿捏到.
你就會知道對自己的心靈.
都是一個很大的釋放和牧養.
因為你知道原來沒有人可以靠自己.
也不需要靠自己.

$^{201}$也不要逼自己.
做不到的.
萬民都在罪惡裡面.
是做不到的.
只需要將你的生命信心.
依靠投入在耶穌的救恩裡面.
你就立即可以獲得神的義.
所以你看這個轉向.
分水嶺是這麼嚴肅.
所以這個就是保羅運用的辯論技巧.
一直帶進救恩的真理.
所以首先講到神的義.
首先講到神的憤怒.
所以神的義包含他的憤怒.
所以不可以隨便稱人為義.
一定是有一位義者.
無條件替我們成為罪.
我們才可以在神的裡面.
得到神的義.
所以這是很重要.
我們要不斷去了解.
我們可以再仔細去看每一個段落.
我們看了18到32節.
神的義如何通過信心去顯明呢?.
萬民都欠缺神的義.
第二章就是今天我們要講的.
外表的.
表面以為自己有義的.
其實是一種自義.
這種自義.
第二章的有趣地方.
不斷講那個「你」.
其實是一個虛構人物.
只是若有所指.
即是說「泛指」.
這樣的行徑的人.
都不得到神的義.
到真正講到.
很指明道性的某些類別的人.
就去到例如第九章.

$^{241}$就立即指向猶太人.
他們種種的問題.
就真的會整個民族這樣講.
所以第二章其實.
好像保羅在一個假想一個人.
在他面前這樣來講.
即是說某人.
例如.
差不多等於對當時代.
一些以為可以靠這樣的神學.
靠這樣的哲學.
哲理來得到.
以為自己可以去到將來的世界.
或者得到他們認為的神的認許.
就不可以了.
第三章就更加.
去到一個總結句.
三章的十九二十節.
所以我們已經斷定了.
沒有人可以在神的面前稱義.
沒有人可以靠律法.
在神的面前稱義.
因為律法本來叫人知道自己的罪.
好像一個顯影劑.
將我們本來的面相.
顯露出來.
所以它是不會帶到我們去稱義的.
這是律法的功用.
所以大家看到這個大概的分段.
所以現在仍然是在神的憤怒下的人.
不同的表現.
第二章就很明顯是講一些.
不需要神的人.
真真正正以為自己是無神的人.
他們就靠自己.
以為自己可以得到.
這種良善.
道德規範.
就能夠令到自己.
好像找錯了平安之境.

$^{281}$就可以得到神的蔭秘.
其實他們根本就沒有神的觀念.
這裡又要講一點點.
因為我其中一個引述.
因為當時保羅的時代.
是一個比較出名的哲理運動.
一個這樣的信念.
一套哲學.
就叫做斯多瓦的學說.
你上網可以搜尋一下斯多瓦.
就很詳細的教義.
要人們信奉這個哲學.
要怎樣為人.
可以看《史道行傳》.
因為保羅在雅典.
就真的有這樣的人跟他辯論.
《史道行傳》的十七章.
就因為他在帖撒羅尼迦加.
要立即.
因為一些的迫發的緣故.
他沒辦法.
要立即離開帖撒羅尼迦加.
輾轉就去到.
突然間說到口上邊.
去到比利亞.
後來就去到雅典.
在雅典的阿勒巴古那個地方.
一個半小山丘.
我們去過了.
前年我們去希臘之旅.
那個阿勒巴古.
是被人辯論的地方.
很適合一個像會議室的地方去辯論.
我們看看這段經文.
16到34節.
我們先看這些引文.
然後我們進入第二章的經文.
就是那個單數的「你」.
是說哪一類的人呢?.
有三個「你」.

$^{321}$第16節開始.
保羅在雅典等候的時候.
等候的是誰呢?.
就是提摩泰和西拉.
叫他們回去帖撒羅尼迦加.
因為似乎是針對保羅的歹徒.
有些猶太人買一些地方土匪.
想逼迫保羅趕走他.
保羅見到這個情況.
他就自己先離開帖撒羅尼迦加.
第16節.
保羅在雅典等候他們的時候.
看見滿城都是偶像.
心裡著急.
於是在會堂裡與猶太人和前進的人.
並每日在市上所遇見的人辯論.
市上就是希臘下面的市集.
那裡相當出名.
相當漂亮的地方.
再上一點點.
就是巴特隆神殿的位置.
即是今天的位置.
當日是非常鼎盛的.
那個地方是很出名的.
希臘那裡.
第18節.
還有以比古羅和斯多瓦兩門學士與他爭論.
保羅與這些哲學家爭論.
無論是神官.
無論是怎樣得到神的救恩.
斯多瓦其中一個就是.
你靠好行為就可以得救.
他們承認這個世界有秩序.
有理性.
可以透過你的個人修養.
道德的修養.
就可以得到國度的進入.
他們的人生相當高尚.
第一章不拜神.
拜偶像.

$^{361}$拜完偶像之後.
就有很多生命道德的降低.
第二章不是這樣的人.
所以保羅其實很涵蓋他所認識的人類.
可以想像羅馬書這裡是多豐富地探索了.
不同的人對救恩觀.
對存在在這個世界的想法.
都涵蓋了.
但是三章的十九二十節.
我們已經斷定了.
世人沒有人可以在神面前稱義.
沒有人可以在神律法面前稱義.
所以一除定音已經說了.
所以第二章和第一章.
是相當不同的人生的表現.
如果用我們簡單的第一章的人.
就是甚麼都不要管.
我們是最大.
總之令到我們的身體.
歡愉,開心.
每天大吃大喝.
玩到盡就行了.
還有那些人對你不好.
你就特別要針對他.
一種利己主義.
不是利他主義.
第二章的人.
這裡我們再看看.
就和他爭論了.
說保羅在這裡的講道.
都有少少的含義.
就是將人以為的世界觀.
人以為的神觀.
都簡單的講出來.
所以保羅都用這些.
當時的哲學家.
相信的東西.
作為一個題材.
作為一個切入點.
令到他要帶出的.

$^{401}$當然不是和他們一樣的結論.
保羅要帶出的就是說.
你們不知道的那位神.
你都敬畏.
我就告訴你.
我們信的是一位怎樣的神.
這樣就帶入一個關注.
特別因為既然這些.
這麼有學問的人.
聽到保羅不是在胡言亂語.
不是在訴諸一種情感.
是一種很理性的分析和辯論.
他們會聽.
聽可能會說你講的東西.
我們從來沒有想過.
可能有些insight.
這是一個切入點.
我們可以看看.
相信大家都可能相當熟悉的.
這一段經文.
就是未得之神.
未認識之神的一個講道.
整篇的福音講章.
就放在17章這裡.
下面就說.
「與他爭論有的說.
這胡言亂語的」.
要說什麼呢?.
就是說他們否定保羅講的.
他們聽不明白.
又沒有聽過.
就以為保羅胡言亂語.
但不是的.
「有的說他似乎是傳說外邦鬼神.
這說是因保羅傳講耶穌和復活的道」.
就是說復活對根本無論什麼哲學都好.
是超脫不了的.
你怎麼可能一個人死了.
無論你什麼身份都好.
怎麼能夠復活呢?.

$^{441}$怎麼能夠永活呢?.
但保羅或者基督教.
當時的基督教就是.
一直傳遞的就是.
一位已經死了而復活的主耶穌.
所以一直就是胡言亂語.
要麼就說不是.
他傳說外邦鬼神.
第19節.
「他們就把他帶到阿勒巴古」.
是一個當時辯論的地方.
說你所說的這個新道.
我們也可以知道嗎?.
有些人沒有聽過.
有些人就覺得.
哇 很特別.
你這個見解.
可不可以再詳細的去展述一下呢?.
如果引起人注意.
也是一件好事.
是會真的帶來一個辯論.
尤其是你去到一些有識之士.
有時你能夠讓他們都看到.
我們不是一些目不識丁的人.
來帶進一個大家共同都關心.
但又好奇的討論課題.
其實都是重要的.
接著下面的意思就是20節.
「因為你有些奇怪的事傳到我們耳中」.
我們願意知道.
這些事情是甚麼意思呢?.
「雅典人和住在那裡的客人都不顧別的事.
只將新聞說說聽聽」.
這是當時的人.
對於哲學的思想.
新的思潮.
他們不是一概抱一個拒絕.
而是他們很想了解.
這是一個好的探求的基礎.
但問題是最終達至的是怎樣的神觀呢?.

$^{481}$是否能夠帶人去到救恩耶穌的裡面呢?.
這才是那個課題.
譬如大學的地方.
特別讀哲學的學系.
真的有很多討論.
但最終都是告訴人.
人可能要靠自己等等.
不能帶進救恩對人的生命的釋放.
留意第22節.
「保羅站在阿勒巴古當中說.
「眾位雅典人啊.
我看你們凡事很敬畏鬼神.
遇見一座壇.
上面寫著『微色之神』.
你們所不認識而敬拜的.
我現在告訴你們』」.
保羅就是從他們好奇.
這個探求真理.
對真神,對鬼神的好奇.
和很想來到去.
不要錯漏了.
拜少了.
這樣這個本來是一個很好的態度.
但是『微色之神』都是一個偶像.
但保羅現在要傳遞的.
在這一點上.
大家都共同承認.
大家承認有人類以外的神明.
但是哪一位才是真正的神明呢?.
保羅藉著這個契機.
因為真的見到你們有一個這樣的像.
保羅不是要傳遞一個新的偶像.
而是透過這個偶像.
告訴你帶你進入永恆的這位造物主.
這個可以說是一個很高明的.
引起關注的切入點.
傳福音的一個技巧.
譬如大家熟悉的《約翰福音》第四章.
那些門徒出去買食物.
在這個.

$^{521}$死人的地方.
在這個薩瑪利亞.
有一個薩瑪利亞的婦人.
中午出來打水.
耶穌就問她要水喝.
那個婦人就很懼怕.
也很詫異.
為什麼你會問一個薩瑪利亞的女人要水呢?.
引起她的關注.
就帶進那個救恩給這個婦人.
所以傳福音的技巧原來都是這麼重要.
不是天方夜譚.
不是說一些無關痛癢的事.
而是跟這個婦人有關的事.
現在保羅跟這些辯論的人有關的好消息.
就是你們所拜的微色之神.
我告訴你.
就是造物主.
這位耶穌基督.
然後告訴我們.
你們所不認識而敬拜的.
我現在告訴你們.
保羅對自己的信仰相當肯定.
不是無伶兩可.
不是自己都好像知不知之.
又不怕又不敢說.
等等.
第24節.
創造宇宙和其中萬物的神.
既是天地的主.
就不著人手所做的事.
也不用人手服事.
好像缺少什麼.
自己倒將生命氣息萬物賜給萬人.
他從一本(由一個人).
做出萬族的人.
住在全地上.
並且預先定準他們的年限和所住的疆界.
所以我們不是宿命論.
是神的揀選論.

$^{561}$所以我們每樣東西某程度都是整定的.
這不是我們中國人迷信的命運.
所以要看獎.
問問神婆.
不是這樣.
而是神定了我們的人生.
定了我們不同的.
我們要去的階段.
再準確些.
每個人都有自己的跑道.
有些人是短跑的.
有些人是長跑的.
所以我們每個人都要有自己對神的認定.
祂給我們的總是有祂的揀選和祝福.
保羅自己也這樣印證.
定準了他們的年限和所住的疆界.
要叫他們尋求神.
或者可以揣摩而得.
其實祂離我們各人不遠.
我們生活,動作,存留都在乎祂.
這種又看不到.
但是又在祂的掌管下.
這種掌管又容許我們有絕對的自由意志.
這是一個很特別的想法.
對這些哲學家來說是很新穎的.
你是自己決定自己一路活在這個世界成長.
但是原來一個平衡的世界.
還要掌管著萬有.
這樣才是基督教真正對真理的解釋.
對人生最美麗的解釋.
這裡告訴我們.
然後就引他們當時的哲學思想.
就如你們作詩的友人說.
我們也是他所生的.
這裡是引甚麼呢?.
我不太清楚.
可能要找一些資料才知道.
可能是當時一些古籍很出名的哲學家.
例如柏拉圖,亞歷士多德等等這些人.
第29節.

$^{601}$我們既是神所生的.
就不當以為神的神聖.
像人用手藝,心思所雕刻的金銀石.
世人蒙昧無知的時候.
神並不鑑察.
如今卻吩咐各處的人都要悔改.
所以是層層進入了救恩的核心裡面.
你以為與神無關嗎?.
其實我們每個人的生活動作,存留.
都是在乎祂的.
祂定準了我們的人生.
不只是一個你有你我有我.
我掌控我自己關你甚麼事?.
為甚麼你要掌控我?.
不是的.
原來表面上我們可以自己揣摩,揀選.
但人要得到神的,認識神的機能.
其實已經很弱化.
如果不是神的介入聖靈工作.
人其實自己是尋找不到神的.
但這裡告訴我們.
祂已經安排了,定了時間.
要各處的人都悔改.
即是福音的傳揚.
福音就是要帶進神的能力.
神的慈愛,神的義.
要讓那些聽明白,相信的人.
開啟他們與神溝通的機能.
尋找神的機能.
來得到與神的接觸.
因為祂已經定了日子.
要藉著祂所設立的人.
就是耶穌.
按公義審判天下.
並且叫他.
從死裡復活.
給萬人作可信的憑據.
一個非常完整的救恩.
說得很清楚.
由一個未識之神.

$^{641}$切入到造物主.
原來與每個人都息息相關.
我們去尋索他.
其實必須要透過.
在神啟示的這位耶穌裡面.
我們要悔改.
我們要去到祂的救恩裡面.
這樣就可以與神.
恢復這份靈性的交往.
認識和尋求祂.
所以一個很完整的救恩.
這個時間應該是呼召了.
在座的.
其實不是在座.
他們是站在那裡的.
的朋友.
不論你是什麼哲學家.
或者是什麼背景.
現在你要藉著耶穌.
祂所設立的人.
因為祂已經由死復活.
來證明神已經要將這份救恩.
賜給每一個願意謙卑信的人.
所以你記得第32節.
就是一個呼召.
眾人聽見從死裡復活.
就有譏笑他.
又有人說.
我們再聽你說.
於是保羅從他們當中出去了.
有幾個人貼近他.
信了主.
其中有阿略巴古的官.
即當時這個地方的地方官.
好像類似是管會堂.
叫做丟尼修.
和有一個婦人名叫大瑪麗.
還有別的人都一同信主.
即是說這些多麼顯赫.
多麼有社會地位的人.

$^{681}$聽明白了就信了保羅說的道.
在這個廣場的事公.
得到神的救恩.
保羅藉著一個偶像.
來帶進耶穌神設立的救主.
來得救.
所以當我們看到這些經文.
保羅的時候.
所以保羅一點都不難.
他很有意思地.
不只是說那些本身很壞.
拜偶像的人.
道德行為不好的人.
原來保羅針對的是.
當時很流行的文化的哲學思想.
即是人靠自己.
去做好自己的行為.
就以為自己可以很理直氣壯地.
就可以去上天堂.
這個就是斯多瓦或斯多葛派的翻譯.
就是美德是人生最高的目標.
幸福來自一個合乎理性的生活.
所以相當高尚.
相當有道德水準.
但是是一個沒有神國度的地上人生.
進不到永恆的天家.
為什麼呢?.
因為這裡只是點出了他們的問題.
二章一到第十二節.
就會點出這些有道德的人.
他們的限制.
但是真正的結論.
其實是在三章二十節打好.
所以我們慢慢逐步逐步.
所以剛才我說了.
從第一章的福數的他們.
他們.
到第二章.
幾次都是說這三類的人.
有道德.

$^{721}$有倫理的人.
當然包括有哲學.
因為有哲學才會有道德.
因為有一套做人態度.
做人的價值觀才會有道德.
沒有這個價值觀就不會影響你的人.
很簡單.
如果你不是有政治的信念.
你都不會變成某些人的行為.
一定是有這個信念先行.
所以原來信仰是影響著我們每個人的人生.
所以我們有時候回教會.
慢慢變成習慣不為意.
其實真正是因為神學.
因為宗教的信仰.
導致我們有這樣的行為.
所以這些哲學家一樣.
所以保羅其實很認.
是不是有些好朋友都是這些人.
所以保羅在寫《羅馬書》的時候.
就想起這一類的人.
其實他們需要恩信稱義.
他們不能夠靠自己的道德.
靠自己的理想.
對這個世界的認知.
來得救.
我們都一而再說過.
保羅處身在一個很黃金的羅馬帝國.
剛剛開始.
就是這個阿古斯都.
回到剛才還沒說到這裡.
阿古斯都的人生.
是在說羅馬的承平.
你在這個帝國底下的隱蔽.
你就有平安.
這個Roman Peace.
你過一個非常太平盛世的好日子.
就是說你可以很容易地做生意.
你可以很好的到處去探索.
就是說一個新的很平安的世界出現.

$^{761}$但是保羅探求的不是那個帝國.
不是要衝擊這個帝國.
而是很想告訴當時的人知道.
真正的平安是來自與神的和好.
真正的平安第五章就說.
真正的平安是來自你在神的義裡.
被稱為義的人身份的轉變.
才來得到真正的.
可以行神期待的道德的義行.
第六章和第七章就會探討這個問題.
我們內心的那個罪性.
怎樣去處理它呢.
其實斯多瓦的哲學都在面對.
他們都是崇尚著這種理性.
這種完美的道德.
但是勝在不了神和耶穌復活的救恩.
所以大家要留意保羅在第二章.
我們說了這麼多的引言.
當我們進入經文裡面.
我們就可以逐段去掌握.
今早有姐妹說.
希望能夠列出新約保羅羅馬書引述的舊約.
我想想也不錯.
不過我希望和大家一起.
完成十六章之後.
我們大概可以總結到有多少的舊約.
保羅引述的.
以賽亞書是引很多節的.
詩篇,耶利米書,創世記.
譬如一章的十八節.
我又多說一點.
關乎到保羅引述的舊約.
有些是明顯經上記著說.
那段經文本來是舊約說的.
但是我現在不是拿來用.
而是應驗在現在我討論的神學裡面.
就是這樣應驗.
經文可以這樣應驗.
例如十四,十五章.
保羅也這樣理解舊約聖經.

$^{801}$那些就複雜了 我們先去到經文.
有些未必是直接引用.
上星期我們說過直接引用的.
哈巴谷書二章四節.
直接引用的.
就如經上所說.
「義人必因順得生」.
直接引用的.
所以他提到的題目.
他很清楚知道經文,舊約討論過這個問題.
真的很熟悉.
沒有AI,沒有電腦.
能夠這樣熟悉的這麼厲害.
真的是一個拉比.
但是他現在寫新約的聖經.
可以馬上引述舊約.
根本不會有參考書在他面前.
舊約的聖經是卷軸.
很重,很大.
你可以想像到保羅對舊約那種純熟.
和對整卷書的上文下理.
那種了解.
不只是聰明這麼簡單.
是神給他的特別的一個恩賜.
為著要寫聖經成為後世的.
新約權威的經卷.
另外有些就是間接引述.
例如一章十八字說什麼呢?.
自從造天地以來.
神的永能和神聖是明明可知的.
雖是眼不能見.
但藉著所造之物.
就可以曉得.
叫人無可推諉.
是引述哪一段經文呢?.
《創世記》第一章第二章.
起初神創造天地.
告訴我們造物主是這樣的形態.
因為保羅的年代.
對神的存在.

$^{841}$都有很多穿鑿附會.
很多哲學家.
都有他們自己的鬼神學說.
例如宙斯和其他神在天上打架.
所以很清楚.
造天地是造物主.
創造宇宙的真神.
所以在羅馬書第一章十八字.
已經有一個間接的引述.
對創世記第二章的認定.
所以如果今天我們有些人說.
特別是新派的人說.
創世記只是一個傳說.
不是一個事實的根據.
你想想保羅的創世論.
已經是站不住腳.
即是說原來我們都要留意.
使徒們他們怎樣理解舊約.
今天第二章有一節聖經.
詩篇的62篇.
又是直接的引述.
62篇已經是第十二節.
我們到先開始講述.
第二章的第一節開始.
我們又要讀聖經.
希望我們每一次.
我做了Power Point.
但其實要進入經文.
因為開始要逐段逐段解.
我們希望大家熟悉段落之後.
每一次我們能夠找到一節的金句去背.
是多重要的金句.
我們要找到來中讀.
去背誦它.
第二章大家聽一聽這一段.
今天我們只能夠講到第十二節.
「你者論斷人的」.
開始若有所指.
誰論斷呢?.
說的是哪一類人呢?.

$^{881}$你大概明白這個意思.
無論你是誰.
可能是不同階層.
有些可能是做官.
有些可能是學者.
有些可能是很有修養的人.
唯一的就是.
那個口和心.
是對沒有他那麼好的人.
一個論斷.
一個審裁.
接著說.
無論你是誰也無可推諉.
你在什麼事上論斷人.
就是說你在什麼事情.
他做了什麼不好.
你去論斷他.
接著說.
就在什麼事上定自己的罪.
這個真的挺特別.
為什麼呢?.
原因就是.
因你者論斷人的.
自己所行卻和別人一樣.
被論斷的人一樣.
我看一些參考書的解釋.
別人呢.
我再看原文.
原來保羅刻意用「別人」.
我們中文看不到那個分別.
原來是另一個類別.
就是說你去論斷那個人.
保羅很重的說話就是.
你看他是異類.
就是說.
不同你那個身份.
那個一類的人.
就是說.
這個看不起對方的嚴重性.
保羅用原文的字.

$^{921}$就可以反映出來.
我們中文看不到這裡.
原來是另一類.
就是說.
你那個眼角高到一個地步.
看他不起之餘.
還要看他是另一類.
我們中國人以前早段時間.
香港大陸撐.
現在沒有人這樣說.
我一聽.
因為有阿晴撐那套電視劇.
《望中人》.
在裡面偷渡下來的.
這些不務正業.
好食爛飛的人.
就是阿晴.
我們就將這個阿晴.
本來這個字是沒有的.
一個電視的人物.
阿晴就代表著這些異類.
和我們香港人.
現在沒有了.
以前應該是80年代的時候.
大家一聽阿晴.
就知道我們看不起.
一見到這些人.
就算他穿得很漂亮.
我們看他一看.
就覺得他是另類的.
這個另類就是保羅所說的.
自己所行卻和別人.
是另類的人一樣.
保羅說話是很重的.
其實大家都是人.
他做得沒那麼好.
你就看不起他.
但是有時候口不說.
心說.
眼說.

$^{961}$就看他是一個另類.
所以就論斷.
不過第二節.
我們知道這樣行的人.
神必照真理審判他.
原來照的是真理.
那把尺是真理.
是神的真理.
不是A和B.
或者我們香港人和阿燦的尺度.
那把尺當然.
我們又不會蹲下坐.
現在在地鐵.
很少見到人再蹲下.
但是80年代的時候.
偷渡下來的人.
在香港很容易蹲下.
等巴士的時候.
是很噁心的.
現在就很少了.
現在真的很少了.
沒有了.
我們真的一看到他蹲下.
我們覺得很不順眼.
因為我們已經覺得自己很有文明.
很有修養.
再加上吐痰.
應該都少了.
除非那些爛仔才會.
一般都不會.
但是你可以想像到.
原來保羅告訴我們.
神的真理去審判每一個人.
按那個標準.
是神的真理的標準.
所以保羅在這裡.
其實都指著一大班.
表面有修養.
有哲學.
有好的行為的人.

$^{1001}$但是這樣的情況.
相對於上文講的種種罪.
已經是從裡面孵化出來.
裡面先是否定神.
不容許祂干預.
不要祂的尊榮.
而自己降歸去拜偶像.
拜偶像就帶來內心沒有底線.
就去犯種種的罪.
但是第二章沒有講到什麼罪.
只是講到對人的論斷.
即是說那個批評.
近乎是一種自己已經.
沒有需要再受到評價.
或者那個判斷.
那個審判.
所以這個都是一個.
我們需要去求助幫助.
真的每個人都需要面對這種危險.
第二章第一第二節.
第三節.
你論斷行者事的人.
自己所行卻和別人一樣.
第二節再提.
因為第一節已經講了.
你在什麼事上論斷人.
自己所行卻和別人一樣.
第三節又再講一次.
自己所行卻和別人一樣.
你以為能夠逃脫神的審判嗎?.
即是原來做物主是監察人心.
這是基督教很重要的一個真理和事實.
這裡告訴我們.
不可以逃避神的審判.
第四節.
還是你藐視他豐富的恩慈.
寬容忍耐.
不曉得他的恩慈是令你悔改.
你竟任著你剛硬不悔改的心.
第五節就是神要審判的第二個致意者.

$^{1041}$就是心已經變得麻木剛硬.
第四節是很重要的出路.
神有豐富的恩慈,寬容,忍耐.
就是對一個不斷從事.
越來越將別人踩低而抬高自己.
表面的德行的人.
又是能做到的人.
就有一種寬容.
不是寬容被論斷的人.
是寬容被論斷的人.
能夠在他的恩慈,寬容,忍耐.
都是神另一種屬性.
不是憤怒.
是一種包容,等候.
等候悔改.
等候我們心靈的一種發現.
自己在神面前的真正光景.
真正的我.
所以我們信主.
我們能夠回教會.
是很寶貴.
就是不斷有聖靈的提醒.
不斷有聖靈在我們內心的教導.
所以原來是恩慈,神的寬容,忍耐.
是領我們悔改.
回到祂的面前.
所以第二個神要審判的.
自然的人.
就會有一個剛硬的心.
不悔改的剛硬的心.
即是心靈.
很明顯的例子.
就去到第九章.
講法老就是這樣.
法老就是將那個心.
每次的神蹟他都悔改.
但是神蹟一過了.
哇 我知道錯了.
整個尼羅河都是血.
摩西舉象.

$^{1081}$那些魚死光了.
發臭.
那些大神來說不行了.
那些魚死光了和臭了.
法老就悔改.
出埃及記第七章.
一而再講了很多次.
當那些災一過了.
法老就將心一硬.
即是心硬.
即是心不再害怕了.
沒事了.
可能是那些術士.
是摩西在搞出來的.
再找一些更厲害的術士.
來攻擊以色列人.
第二個災又來了.
如是者終於去到第十災.
這個滅掌日.
法老就低頭了.
低頭完之後.
死光了那些兒子之後.
死兒子之後.
再找一些馬車.
去追逃了去.
準備出埃及的以色列人.
否則都不會有紅海的追兵死光了.
被紅海覆滅了.
即是原來死長子.
都沒有真正的悔改.
法老.
可以想像得到.
那個剛硬就在這裡.
更加淋漓盡致的形容出來.
所以我們求主讓我們真的能夠.
要不斷有一個柔軟的心靈.
柔軟不是一個沒有知識的心.
而是對自己一個很客觀的.
一個真正的.
在神面前的.

$^{1121}$坦開.
闖開的心靈.
讓神的聖靈去光照.
所以.
早晚的.
個人的檢討和祈禱.
當然大家臨睡.
我不知道有些人會失眠.
一睡都要一個小時才睡得著.
那你就更加可以反省了.
但有些人一睡就立刻睡著了.
那就.
都要珍惜那幾十秒.
即是和神有一個很美麗的分享交通.
其實是訓練自己的心靈柔軟.
真的不需要用言語都可以.
那種和神之間.
打開的心.
向祂說你的狀況.
今天的際遇等等.
讓我們心靈能夠柔軟的一個.
很重要的訓練.
所以第二個就是剛硬的心.
自然是會帶來剛硬的.
沒有道德的人.
他真的知道自己很差.
他有時容易去悔改.
但有道德的人.
難一點去悔改.
就是他有一個剛硬的心.
即是他自己已經覺得自己很好.
不需要不稀罕你的救恩.
所以就.
如果不悔改.
就為自己積蓄憤怒.
以致神震怒.
顯他公義審判的日子來到.
即是說.
如果你一直都靠自己.
沒有主耶穌.

$^{1161}$那信念強到一個地步.
有些人真的很強.
突然想起.
其實Terence.
鄧.
現在忘記了.
阿紫堅.
他爸爸鄧家寶.
在圍村.
他要辦很重要的圍村事件.
很難信主.
這麼多年.
真的非常抗拒福音.
但去到差不多晚年.
他去年離開世界.
很奇妙.
真的看到他的生命轉變.
一來他的癌症都嚴重.
自己知道自己的情況.
但真的看到他的心靈的柔軟改變.
所以我們真的有不同的人.
那個心靈有時就是那種剛硬.
其實一直在積蓄神的震怒.
到公義的審判.
但在這個日子未到之前.
不要忘記.
神仍然有豐富的恩慈.
寬容和忍耐.
所以我們不要灰心.
繼續將我們的救恩.
有愛心地傳講.
有溫度地去傳講.
繼續在友誼的關係上.
很深刻的是.
有一次他已經痛得很厲害.
在天主教醫院.
過去祈禱.
我問他為何要祈禱.
他說好啊.
他說祈禱很痛.

$^{1201}$他真的痛得很厲害.
祈禱完之後.
我說世伯鄧先生.
不如相信耶穌吧.
他說那還未必.
他很堅決.
一樣歸一樣.
我現在很痛.
你說我祈禱.
我沒有問題.
他閉上眼睛祈禱.
接著我就說不如相信耶穌吧.
他說那還未必.
手又插滿了東西.
再下一次.
再過了兩個星期.
他又真的很不舒服.
那次真的有些改變.
我說不如相信耶穌吧.
他身體已經很弱了.
他就有改變.
他又點頭.
我說我看不到你點頭.
他戴著氧氣罩.
我說不如你相信.
你舉起手指.
他真的這樣舉起手指.
手指包滿了東西.
其實很大隻.
真的用力這樣舉起.
真的立刻幫他洗禮.
那晚就離開世界了.
你可以看到.
那個生命的轉變.
都看到神的恩慈.
寬容忍耐.
所以這裡告訴我們.
第二章說的有道德的人.
他們有這些偽善的偽裝.
有心靈的剛硬.

$^{1241}$原來越有自己道德一套.
道德標準的人.
那個心是相當硬的.
很難去改變他自己的心靈.
所以我相信心靈剛硬.
必須用聖靈.
用愛.
用祈禱.
我猜眼淚可以融化剛硬的心.
最親的人對他的眼淚.
第六節.
他必照各人的行為.
報應各人.
第七節.
凡行深行善.
行深.
是consistent地行深.
保羅用這個字的可圈可點.
即是說.
哪可能有人.
包括保羅自己.
由始至終.
由出生到你認識這些哲理.
到你採納這些哲理.
到你認同這些哲理.
到你實踐這些哲理.
是哪有可能呢.
行深這樣做.
行深這樣行善.
所以根本沒有人可以做到.
凡行深行善.
尋求榮耀尊貴.
和不能扭壞之福.
即是說.
對自己的心靈.
對身體以外的榮耀.
去探求.
即是對做物主的探求.
對真善美的探求.
那哪有人可以這樣做到呢.

$^{1281}$既是行深這樣行善.
同時間.
不只是為地上的短暫的.
這樣來生活.
還要是為榮耀尊貴.
和不能扭壞之福.
神是怎樣去報應呢.
就是用永生報應.
這樣的人.
所以這裡不是叫我們靠行為得救.
是在討論.
自以為自己是足夠的.
但其實是不足夠.
義者以為自己是足夠.
其實根本就不足夠.
必須是要這樣行深的行善.
所以我們信主耶穌.
這個才是我們人生的標準.
我們真的要行深的.
在主面前.
那個.
分別為聖的奉獻的生活.
還有我們的追求.
不是說我們不緊張物質.
我相信.
我們有很多東西.
我們不緊張物質.
我不應該說這些的.
因為已經去到病態.
我認識一個人.
我不說名字.
也不說他的身份背景.
免得太敏感.
不斷被人借錢.
但家裡有很多.
幾百件恤衫.
我有十件.
我都很煩.
是方便出汗.
夏天換.

$^{1321}$是幾百件.
但到處被人借錢.
是一種這樣的病態.
如果我們只是為了地上的物質.
忘記了尋求榮耀.
尊貴.
人生的意義.
就很有限.
所以這裡告訴我們.
原來生命以外.
是有榮耀,尊貴,美善.
我們要追尋的.
是不會扭壞的.
福氣是有的.
我們要追尋的.
是不會扭壞的.
福樂.
所以保羅用的字眼.
用的對真理的探索.
你們這些有道德的人.
有沒有這個尋求呢?.
當然沒有.
所以基督教.
如果我們能夠.
每個弟兄姊妹.
看到這個永恆的榮耀.
就不會再太介意.
地上短暫的得如失.
這是很寶貴的經文.
第六和第七節.
是很重要的.
因為四節.
四,五和八節,九節.
是一個對比.
三文治的中間第七節.
是最重要的.
所以.
治意的人會在這些爭拗,論斷裡.
第八節開始.
唯有結黨.

$^{1361}$不順從真理.
反順從不義的.
就以憤怒,怒恨,報應他們.
將患難,困苦.
加給一切作惡的人.
剛才說的.
將榮耀,尊貴加給這些人.
永生報應給這些尋求榮耀的人.
但是如果我們一心作惡等等.
神就會用他的憤怒,怒恨,報應他們.
將患難,困苦.
加給這些作惡的人.
不過這裡特別留意.
開始說先是猶太人.
後是希利利人.
第十節就剛剛這個對比.
先是猶太人.
後是希利利人.
因為神不偏待人.
神不偏待人又出現過.
在這本書重複出現過.
神不偏待人.
應該是第十章.
十三節.
你可以想像到.
這裡立即引進猶太人,希利利人.
因為下面第十六節開始.
對不起.
十七節開始.
正正式式說猶太人.
猶太人是另外一類.
因為真是在神的律法下生活的人.
連羅馬帝國的人都知道.
這群人有自己獨特的生活形態.
所以保羅特別分一段出來.
是說他自己的同胞.
不是要映射在自己身上.
而是要說他自己的同胞.
所以保羅特別分一段出來.
是說他自己的同胞.

$^{1401}$不是要映射在自己身上.
而是更加希望神憐憫猶太人.
他們更加難接受耶穌是救主.
他們的律法的意義.
讓他們更加難去承認自己的需要.
所以這裡開始引進先是猶太人.
後是希利利人.
和上文第一章說的有些不同.
這裡說的是.
原來猶太人和希利利人.
一樣有在神面前偽裝的善.
以為自己做到的那個自然.
所以這一段聖經.
真的給我們很清楚的一幅圖畫.
就算我們道德上沒有其他犯罪的人這麼嚴重.
但是心靈始終逃避不了神的那個意義.
不了神的那個監察.
那個等候的寬容.
那個的恩慈.
在這個等候裡面.
我們就應該有一個柔軟的心靈.
有一個敞開的心靈.
奔跑去到神的救恩耶穌裡面.
所以每一個決志信耶穌的人.
其實就是我們的靈魂.
突然間有一個全速前進的跑去耶穌的救恩裡面.
向祂祈禱.
希望祂進入我們的心靈.
我們雖然看不到這個圖畫.
但是我們信主內又事奉的人.
我們就知道其實是他的心靈打開.
他的肛硬的保護膜已經被拆開.
聖靈的工作讓他整個人被光照.
看到自己的心靈去到主的救恩裡面.
所以原來聖靈可以融化最肛硬的心.
很多人的心如果在聖靈的工作下.
他就會看到自己內心的罪.
所以這裡告訴我們.
所以這裡保羅在這一段有一個很重要的對比.
如果我們一生在神的恩慈以外偽裝自己.

$^{1441}$那就很可惜.
亦都很浪費我們自己的時間.
因為你的偽裝就不去追求榮耀尊貴.
不能扭壞的福.
追求的只是那些作惡等等.
結黨,不信從真理,信從不義.
神就會報應我們.
將患難困苦.
即是我們一直追尋的東西已經令我們不開心.
最終更加不開心的就是.
神就用患難困苦加給我們這些作惡的人.
但相反我們如果很願意打開我們的心.
去追求神的.
尋求的是榮耀,行善,尊貴.
神就會用永生去報應我們.
你看看一個疊加的相括弧.
中間那個最重要的就是第七節.
凡是行心,行善,尋求榮耀,尊貴.
和不能扭壞之福.
就以永生報應他們.
這些是金句中的金句.
我們很少人有看羅馬書的習慣.
就不知道.
因為它夾雜著中間左右三文治的兩塊麵包.
就是一些人尋求神.
一些人尋求自己的自義.
未必是被人看到的罪等等.
最近有一個新聞.
很少事的.
有個阿太買菜.
見到地上有些錢.
腳踩一踩.
沒人看見.
撿起來.
放在袋子裡.
可惜天眼被抓了.
所以你可以看到.
報紙就這樣拍出來.
還有片段.
你可以想像到.

$^{1481}$這麼小的.
一下子的貪心.
表面上沒怎麼傷害人.
我小時候.
中學還是大學.
還是出來工作.
撿了一百塊.
拿去元朗警署.
我忘記後來是怎樣了.
因為你不撿.
你不撿.
我那時信了耶穌.
你不撿.
又被人撿了.
就沒有了.
但如果你撿了.
拿去警署.
其實找不到人認領.
聽聞說如果三個月沒人認領.
就是你的.
我好像記得.
好像都有去拿.
有這樣的措施.
但我們回頭來說.
第四節到第十節.
是一個很強烈的對比.
我們的人生.
到底是以自然去安頓我們的心.
還是以尋求神.
作為我們生命不斷的遊園.
第十節又再一次說.
卻將榮耀尊貴平安.
說了兩次.
加給行善的人.
但這裡讓我們看到.
有這樣的情況.
這是第二章第六到第十一節.
就是看我們實際人生的行為.
我們是怎樣為人.
因為一定是神學影響我們為人.

$^{1521}$如果我們的神學觀是很正確.
對這位造物主我們有一份敬愛.
敬畏.
我們就會在我們能夠認知的範圍.
做好我們自己.
至少現在我們能夠渴慕.
想認識多些聖經的真理.
這已經是在和路軌.
尋求榮耀尊貴.
和不能扭壞之福.
你明白了.
你就會有些東西通了你自己的思想.
你就會有些東西.
不知不覺你沒那麼執著.
你會慢慢對永恆.
你會有一個渴慕.
你會有一個積極.
在地上怎樣能夠裝備自己.
今天有位姐妹.
在紅一很多年了.
突然間說我近來有些東西想和你分享.
又很怕和我分享.
原來她說她想全職做傳道.
我聽到之後.
我心裡面說阿們.
這樣就可以想像得到.
我一聽到她說.
我已經說阿們了.
她一向在教會的服侍.
擺上的參與.
都已經不知不覺.
她自己都感覺到自己具備這種感動.
當然還未詳細和她談.
還要見她的先生.
但是這是很寶貴的.
我們的人生.
可能我現在真的去到60歲.
真的會不同了心境.
有種夕陽觀的想法.
這個夕陽觀不是無奈.

$^{1561}$而是很快就過頭了.
永恆.
有一種對永恆的渴慕.
是積極的.
不是負面的.
不是很累.
就是保羅這樣說.
所以保羅雖然說這班人.
其實他自己在滲透.
他自己對榮耀尊貴的渴慕.
說了兩次.
一個對比就是說.
你追求那些東西.
首先你追求就令你不開心.
你作惡.
你去犯罪.
你去肛硬自己.
你得到的報應就是.
神就會用那些各方面的來報應你.
但是尋求榮耀的人.
其實心靈是很豁達的.
你參與服侍.
不問什麼.
只為榮耀神.
其實是很開心的.
你能夠這樣做的時候.
你整個動機被神監察和傳順.
其實你尋求的過程已經是很開心的了.
神就將榮耀,尊貴,平安.
加給這些行善的人了.
所以不是靠行善得救.
也不是靠行為得救.
不過要留意.
為什麼我們不是靠行善.
靠行為得救呢?.
因為有一位聖者.
他的義行.
他的義.
已經加給我們這些不義的人了.
所以是我們的主耶穌.

$^{1601}$用了他的人生.
成為了我們的義.
所以我們就不需要再靠我們的義.
靠我們的行為了.
所以我們就不是靠行為得救了.
但是我們信了主.
有了神的義.
我們就會在生活裡面.
有一個渴慕榮耀,尊貴的心靈.
而這個心靈.
是會慢慢影響我們的生活為人的.
我真的已經從信耶穌開始.
完全沒有這個意欲要說髒話.
我以前說髒話說得很厲害的.
是完全沒有的.
沒有這個意欲.
不是因為那個環境.
不是因為做了傳道不需要說.
或者不准說.
不是的.
是心裡面根本沒有這種意欲.
包括吸煙也是.
完全沒有.
你可以想像到.
我們原來去追求神的國度.
所以.
《馬太福音》第六章33節.
你們要先求神的國和他的義.
這些東西都要加給你們了.
就是說如果我們只是求地上的東西.
本來不是錯事.
但是基督徒有神國度的.
只是求地上的事.
而對神的國度完全沒有興趣.
完全沒有興趣.
這樣你最終求不到自己地上的事.
因為你的心很亂.
你會做錯很多決定.
你會有很多很多不同的動機,手段.
求神的國度的人.

$^{1641}$過程是聖靈調整著我們的內心.
我們的動機.
我們做事的手段.
這樣求.
表面上我們求神的國度.
但其實地上的事.
很多時候神就會有很奇妙的引導.
就會去到類似水到渠成.
所以不要說笑.
自義的人只會看到自己利己主義.
但是保羅告訴我們.
真正的義是恆心的恆善.
尋求榮耀,尊貴和不能扭壞之福.
而章七節成為今晚給大家的金句.
其實第十節也是.
是一個強烈的對比.
所以保羅的修辭,神學的精密.
其實很值得我們逐段去欣賞.
所以第六到十一節是說兩種行為.
所以耶穌說在登山變象.
有些人是路是闊的,路是大的.
得著的人多.
那是滅亡的路.
有些路是少的.
得著的人少了.
這是人生的路.
兩條的門路,兩種的果子.
這是第六章在《馬太福音》的經文.
對不起,我要先喝口水.
我要先跟大家引一引詩篇的62篇.
大衛的詩篇.
說到那個報應.
保羅怎樣引用呢?.
今天第二章就是第二節.
第二節關於新約引用舊約的經文.
62篇.
跟經文差不多一樣.
應該是第十二節.
當然我們可以用一點時間.
因為今天我只是預備到第十六節.

$^{1681}$所以還有15分鐘的時間.
我們看看第62篇的詩篇.
主要我們看的是第十二節.
最後一節的.
我們看看大衛怎樣理解一個在神保守的人.
大衛的詩.
照耶道頓的作法.
交予靈長.
我的心默默無聲.
專等候神.
我的救恩是從他而來.
唯獨他是我的盤石.
我的拯救.
他是我的高台.
我必不肯動搖.
你們大家攻擊一人.
把他毀壞.
如同毀壞歪斜的牆.
張倒的壁.
要到幾時呢?.
他們彼此相倚.
專要從他的尊位上把他推下.
他們喜愛方話.
口需祝福.
心卻就坐.
我的心啊.
你當默默無聲.
是主題曲.
專等候神.
即是說一個敬虔的希伯來人.
他怎樣在這樣的嘲諷.
就坐下.
他專心等候神的裁決.
審判.
唯獨他是我的盤石.
我的拯救.
他是我的高台.
我必不動搖.
已經第二次說.
好像主題曲再唱一次.

$^{1721}$我的拯救.
我的榮耀.
都在乎神.
我力量的盤石.
我的避難所.
都在乎神.
因為他在猶大的曠野.
真是有很多巨石.
真是有很多山洞.
所以就地取材.
靈感就是來自神.
穩固.
對古代的人.
這些巨石就是穩固.
就是保護.
尤其是四十多度.
在這樣的炎熱.
在山洞裡面.
是很陰涼的陰秘.
你可以想像到.
大自然這麼大的落差.
神的盤石.
那種穩如盤石.
對作詩歌的人.
有很大的保障.
今天就是.
如果你經常上網.
你是受到欺凌的.
在學校很多人霸凌你.
或者在公司.
有人排擠你.
每個人去吃飯都不叫你.
在這樣的情況下.
現代人就叫輿論.
壓力下.
神就是你的出路.
神就是你的避難所.
神就是你可能在廁所裡.
自己偷偷的飲泣的場景.
類似這樣.

$^{1761}$所以每個時代都有那個避難所.
因為你想想.
網上的輿論.
那些公審很嚴重.
起底啊.
很恐怖的.
有些人為了這樣而自殺.
所以可以看到大衛.
在當時就看到神.
是他的種種的保護.
不過重要的是講夏文華.
你們眾民.
何時倚靠他.
開始有些教導了.
在他面前傾心吐意.
神是我們的避難所.
就是這個詩的主題.
下流人真是虛空.
上流人也是虛假.
不同你什麼類別的人.
你是低下階層的.
平民百姓的.
一些牧羊的.
或者你是上層的人.
做官的.
等等.
放在神的天平.
第二章的字.
咦.
那類人.
越來越有所指.
因為待會第十二節.
跟第十節完全一樣的引述.
所以現在整個詩的題材.
一路大進.
你不要小看自己.
神才是真正的天平.
去稱讚你.
放在天平裡.
你不要小看自己.

$^{1801}$放在天平裡.
就必浮起.
他們一共比空氣還輕.
當然是一個比喻.
在神的法馬.
在神的標準下.
你以為自己自以為義.
你以為自己比其他人好.
其他人好.
即是你會論斷.
這樣的人.
其實在神的眼中.
比空氣還輕.
不要仗勢欺人.
論斷的人就會.
也不要因搶奪而驕傲.
若財保家增.
不要放在心上.
即是說.
大衛觀察時事.
又見到社會有這樣的現象.
他也是一個聰明人.
也看到這個時代.
很多當時的一些人.
就是這樣.
譬如約克.
或者當時的.
做瑞靈的將軍.
等等.
第十一節.
神說了一次兩次.
我都聽見了.
就是能力都屬乎神.
主也慈愛.
也是屬乎你.
突然間說了這句.
因為你照著各人所行的報應他.
這個慨嘆.
是很對的.
一切都是神的.

$^{1841}$所以不要爭.
不是不要上進.
不要用手段去爭.
在神裡面.
去渴慕神的國度.
為神的國度.
不要計較.
真的不要和神計較.
神也不和我們的罪.
我們的永死計較.
我們和神去計較.
這麼少的東西.
其實很蠢.
所以一切都是在乎神.
大衛說.
所以他說慈愛也是屬於你的.
如果不是你的憐憫.
我也是被審判的行列.
因為你照著各人所行的報應他.
所以這句說話才是最精簡的.
原來整篇詩篇的氣氛.
都是第二章的羅馬書.
保羅說得沒有大衛那麼有詩意.
大衛好像越說越灰心.
其實不是.
一切都在乎你的.
你照著各人的行為.
我們今天看不到.
誰是最準確的判斷.
看不到不要緊.
原來一切神都會照著各人所行的報應.
這就是第二章的羅馬書.
說了兩次.
這就是第二章說那些.
有詩意的人的內心表徵.
這就是62篇.
全篇裡面的氣氛.
所以神是我們的避難所.
為什麼要避難所?.
經常都躲在廁所哭.

$^{1881}$因為受壓.
因為蘇羅逼迫他.
他的兒子阿沙隆逼迫他.
追殺他.
或者他自己面對很多困難.
但是他終於知道神是他的倚靠.
所以要看緊張的就是神的判斷裁決.
所以他必照各人的行為.
報應各人.
第二章六節羅馬書.
還有第八節第九節都有這個氣氛.
所以這個報應就是62篇的詩篇.
是由保羅引述的.
詩篇是62篇的12節.
但是原來全篇都是在第二章的氣氛裡面.
所以如果平衡地看.
又覺得很有味道.
保羅就沒有那麼詩意.
因為他在寫的是教義.
他現在不是用詩歌的題材.
不是將個人的心境去觀察社會.
和他自己個人的遭遇.
無限放大情感.
寫這些是神學.
所以是硬邦邦的報應.
但是原來引述的是詩篇的62篇.
保羅是相當熟悉詩篇裡面的意境.
背後的真理.
你看多麼寶貴.
第二章六節就是保羅引用的詩篇62篇.
他必照各人的行為報應各人.
原來舊約已經有了.
現在保羅更加確定將來一定有一個公義永恆的審判.
所以我們人生信了主.
其實是一個很大的生命的改變.
如果再不去追求榮耀,尊貴和不能扭壞之福.
其實就很可惜.
真的希望我們作為牧者.
我完全沒有介意不是要針對哪個人.
我真的見到有人回來教會崇拜.

$^{1921}$根本不稀罕.
很渴慕聽到.
很想唱聖詩敬拜神.
很想聽聽我回來是否有份於這個教會.
不是啊.
那你就很可惜了.
當然說笑了.
你不要回來睡到11點.
然後你就去喝茶.
得福也好.
脆乎.
食物也好吃.
你何必這麼早回來.
回來生氣的.
你做什麼呢?.
回來崇拜是敬拜神.
不是敬拜人.
每個星期都是這樣.
不是尋求榮耀,尊貴.
所以這些是牧者要教導的.
是要彼此分享,彼此勉勵的.
絕對不需要篤責.
因為每個人都是這樣.
我們因為我們研經就能夠看到.
原來神期望我們人生是榮耀,尊貴.
得著的是不能扭壞的福.
盼望這成為我們彼此互相勸勉提醒.
所以第6到第11節是說不同行為的人.
不過這裡是說在智異裡面的利己主義的人的人生.
那種尋求.
但保羅說真正的人生是應該尋求永恆的.
所以當然第5章基督徒的生命有了清義之後.
就應該是這樣生活.
不過這裡不是說基督徒.
是說在神陣路下的人.
這個類別的人就是自以為義的人.
就會有這樣的行徑.
所以某程度上都相當透徹地將人性剖析出來.
還有一兩個power point還沒說完.
可能要留待下一次.

$^{1961}$還有我們可以第12節那段我們下一次再說.
因為時間都不夠了.
抱歉的下個星期因為我11,12號.
有三堂都挺難預備的道.
因為要跟鄧全道進去青年會小學.
這次輪到我講.
講小學.
你真的想我死.
怎樣跟小朋友講道呢?.
我都希望10號那晚停一停.
星期五就可以有個預備講道.
12號我就要去以樂團契講道.
所以都盼望我們停一次.
但是不停的就是大家要背金句.
還有要找到.
所以這個list.
要將所有列出來.
大概有60節.
羅馬書引述舊約.
大概60節.
所以我們一起慢慢去高呼.
現在已經有了.
一章17節.
二人必因信得生.
哈巴谷書二章四節.
還有今天我們講的二章六節.
就是引述詩篇的62篇12節.
所以不要只是這樣引述.
而是將經文比對.
就會得到原來新約引述舊約.
不是這樣引述.
原來整個氣氛都放在經文裡面.
你就會覺得原來.
使徒們這麼尊重聖經.
舊約的聖經.
所以經常這樣記載.
他們就是這樣運用.
所以運用到很淋漓盡致的.
講那些治理人的內心.
就是大衛那篇詩裡面.

$^{2001}$所浮現的那些所謂下流上流的人.
其實在神的天秤底下.
比空氣還要輕.
所以我們真的要不斷在神裡面.
看自己的擁有.
和自己現在做的事.
是怎樣能夠榮耀神.
所以大家一起豐富那個清單.
第二章後面再加有其他舊約的經文.
未必是直接引述.
但那個氣氛一定有.
例如割禮.
創世記十七章就是割禮的起源.
這些就是使徒們.
猶太人他們很熟悉的.
這些舊約的事件.
割禮到底是甚麼一個屬靈的含義呢?.
新約有其他保羅書信有提及.
歌羅西書有提及.
真的割禮是怎樣的.
所以大家一起去豐富.
我們對保羅書信的認識.
今天就說到這裡.
讓我們在《查經班》.
可以看到保羅的.
很重要的關於真理的教導.
我們不需要.
也不應該尋求自義.
和去硬心.
而是我們能夠有一份.
在神面前敞開柔軟的心靈.
來尋求榮耀,尊貴.
和不能扭壞的福.
願主幫助我們.
再一次帶領我們能夠在前面的復活節.
我們可以在不同的聚會.
去經歷上主.
你的模樣.
和我們在一起服事裡面.
將神的救恩.

$^{2041}$很有力地講解.
接下來我們在紅因20日.
也有福音主日.
求你使用見證的嘉賓的見證.
能夠真的帶來很多人的生命.
有一個反思.
來一齊去決戰信主.
多謝主讓我們能夠有服事.
有言經的機會.
奉主耶穌基督的名.
阿們.
\newpage



\section{羅馬書}
\label{sec:JU6MXStfi9g}
\textbf{2025-04-17 周四晚查經班《羅馬書(7):分題:神的義(六)》 講員:曾裔貴牧師}
\newline
\newline
連結: \href{https://youtube.com/watch?v=JU6MXStfi9g}{\texttt{https://youtube.com/watch?v=JU6MXStfi9g}} ~~~~ 語音日期: 2025-04-19
\newline
\newline
\hyperref[sec:7C1Uqd4_AYQ]{\small{< < < PREV SERMON < < <}}
~
\hyperref[sec:index_chronic]{\small{[返順時目]}}
~
\hyperref[sec:index_scriptual]{\small{[返順卷目]}}
~
\hyperref[sec:b2w9fc63m3c]{\small{> > > NEXT SERMON > > >}}
\newline
\newline
$^{1}$感恩再一次我們在晚上去研讀羅馬書.
我們很盼望能夠對每一段特別重要的分段.
能夠掌握對裡面的.
仕途堡羅想要帶出的訊息.
那個神學的脈絡.
我們能夠慢慢地去掌握.
以致當我們自己去再研讀其他經卷的時候.
我們都能夠有把握對經文的上下文可以來了解.
羅馬書是引述很多舊約經卷.
盼望我們一路看的時候.
又可以掌握到堡羅他對舊約的引用.
讓我們都能夠看到聖經的一致性.
舊約,新約不是分割的.
而是一個整全的重要的神的啟示的訊息.
求神帶領我們.
多謝你讓我們能夠在壽呼節之前.
一起去思想主耶穌十字架的救恩.
帶給我們人類的拯救.
很盼望我們在《查經班》裡面.
我們深入對整個羅馬書的恩訊,清義.
是有很基本,握要的認識.
使我們能夠對全福音有信心,有把握.
我們一起仰望等候奉耶穌基督的名.
阿們.
好,我們繼續第七講.
仍然停留在一個重要的段落.
也是堡羅在他講恩訊,清義的救恩的時候.
他提到萬民其實是規缺神的義.
萬民都需要恩訊,清義.
在一章的十六,十七節.
相當要緊就是將福音.
什麼是福音呢?.
來到他講了很重要的答案.
就是必須要恩訊以至於信的信心.
連上耶穌基督的義.
我們就可以得到神的義.
接著一章十八節.
我們上兩個星期開始看.
提到萬民都是活在神的憤怒之中.
因為我們都是虧欠神的義.

$^{41}$所以我們自己怎樣進行.
都不能夠符合滿足神的義標準,要求.
所以堡羅由一章十八節.
一直到三章二十節.
就很清楚講有三類的人.
都是虧欠神的義.
這三類人等於全民.
所以結論就是萬民都需要恩訊,清義.
萬民都是在神的憤怒之下.
所以三章二十一節.
又是剛剛頭尾呼應的答案.
中間就是需要救恩的原因.
後尾就是救恩的答案.
就是恩訊藉著耶穌的寶血.
所以三章二十一節.
就是另外一個更詳細的.
講到怎樣可以在神的義裡得到永生.
所以我們一直在這幾堂集中講.
堡羅列舉萬民在神面前實際真正的光景.
所以我們留意到一章十八節到結尾.
很多時候是講他們,你們等等的眾數.
但到了第二章.
一而再出現了你這個假設性的人物.
即是說某一個類別的人的代表.
其實就是泛指著在這樣的形態生活的人.
其實就是在這種狀況,類別底下的人.
這些人在神的眼中又是怎樣的光景呢?.
堡羅列又是很仔細地探討.
所以我們又可以回到.
我們上一次講到二章十一節.
特別強調神是不偏待人的.
所以幫大家透過這裡的講解.
我們在組織的分段.
一章如果大體上來說.
一章(沒聲音).
沒聲音的意思是甚麼?.
對不起 是收不到嗎?.
是不是聽不到我說話?.
聽到的 聽到的.
OK OK.

$^{81}$我們繼續.
很重要的是我們在這個分段裡面.
要掌握到堡羅想要大進的救恩.
當然大進救恩不只是解釋救恩.
而是很想羅馬的教會.
支持他在宣教事業裡面.
同心地向萬民,未得之民傳福音.
這是這本羅馬書重要的寫作目的.
不只是講恩信清義.
這是很重要的.
但亦很清楚讓我們看到.
堡羅既然是神所呼召.
來被選召成為使徒.
就是要特派傳神的福音.
這是他很清楚自己人生的召命.
他將這件事跟羅馬的教會分享.
所以他很想去到羅馬.
將他所得著的恩典.
跟羅馬的教會分享.
所以我們在總綱裡面.
一至十五節.
你可以說是堡羅的夢召,事奉.
接著十六節很快就說.
他不以福音為恥.
福音是神的大能.
要救一切相信的人.
所以這個能力.
在我們作為傳福音的人.
我們記得要讓人.
將他帶進去.
他自己看到自己生命的需要.
他去求告主名.
他一求告就可以得救.
所以這是一個大能的彰顯.
我們能夠叫人去求告主.
這是一個很大的能力的彰顯.
所以不是我們的能力.
是神的能力.
是聖靈的能力.
在不信,拒絕,承認,拒絕,接受.

$^{121}$有神,有真理的世界,世人.
其實都是蒙昧,無知的.
但是神的作為就可以這樣來彰顯.
所以一章十六,十七節.
就是恩順,清義,萬民可以得救的保證.
那個答案.
接著一章十八節.
一直到三章二十節.
就是那個內容.
什麼人需要福音.
為什麼需要福音.
因為萬民都是在神的義之外.
在神的憤怒裡面.
這是人的光景.
所以其實我們越信主就越發現.
三個類別.
我們差不多都有了.
當然我們絕對不是肉身的猶太人.
但是有時候那種狀況.
都有些很相似的地方.
我們繼續下去.
大家看著這個重點.
要留意這些分段.
其實都是重要的分段.
接著我們看回上一次我們還沒講完的.
我們看到第二章的一到十一節.
似乎真的不是叫映射.
真的好像就著一些當時相當有哲理.
也都有理想.
有他們的生活上面.
某程度是他們認為的.
憑著自己的心靈.
自己的理想.
來去活.
就可以得著那個比其他人更好的將來的永生.
但是保羅斷言說了.
這些人其實一直不能夠每分每秒地行神的要求.
所以二章的十三節.
就是凡是沒有律法犯了罪.
也必不按律法滅亡.

$^{161}$凡在律法以下犯了罪.
也必按律法受審判.
十三節就是這樣說.
原來在神的面前.
不是聽律法的遺疑.
是行律法的清義.
所以沒有律法的外邦人.
如果他能夠做到.
這當然是一個假設性的問題.
是沒有人可以做到.
所以我們記得不要將這一節聖經.
無限放大.
抽離了上下文.
上下文是說這些有道德的人.
有理想的人.
有自己的哲理的人.
甚至推銷這些哲理的人.
以為行自己的善就可以得了.
但是保羅說.
是的 如果你能夠行就是了.
所以十三節這裡的意思.
不是再一次反過來說.
原來保羅也想我們靠行為得救.
我們就看上下文看得不夠透徹.
因為上下文也在說.
那些道德家都是需要恩順清義.
所以這裡不是叫我們靠行為得救.
意思是沒有人可以這樣做.
但是其實福音.
因為後面十六節.
保羅說這一切都是照著.
保羅所傳的福音所言.
我們等一下解釋這一節聖經.
就是說如果不是有主耶穌一生.
祂由出世到祂離開世界.
祂的兒在神的面前被越辣.
祂真的在神的面前做好了祂的一生.
祂的兒就加在我們身上.
或者是遮蓋我們的罪惡.
所以是有一個人行了兒.

$^{201}$行了善.
遵行了神的律法.
最終和遵行神的旨意.
在十誡上流血犧牲.
所以就不是我們要靠行為得救.
所以千萬不要將二章十三節.
作為抵觸保羅所說的恩信稱義.
從來都是恩信稱義.
從來都是我們沒有能力.
靠自己行神的律法要求.
不能夠靠行為得到神的喜悅.
神的兒是在律法以外賜給信的人.
所以是因為有二者耶穌.
替我們成為了兒.
我們才可以得到這個救恩.
所以多麼可笑.
二章這一班人.
有所謂自己的道德.
有自己所謂的完美.
其實是論斷其他人.
和只是會有一種徒有外表的偽善.
有這個偽善.
就會令自己的心靈很剛硬.
不肯悔改.
不肯讓神的聖靈軟化和感動.
有時要到很後期才有大切大悟的改變.
在實際的行為上.
根本符合不了自己一直說的行義和行善.
所以你可以想像到.
神的審判一樣會臨到這些自以為是的人身上.
很多解經家都認為.
保羅針對當時的時代.
有哲學家推銷靠善行和理想去得救.
但保羅看得很清楚.
他自己也是這樣追求.
但他自己也發現.
因為認識了耶穌基督之後.
他自以為是的過去行為.
在主耶穌面前不值一提.
甚至看為糞土,污穢.

$^{241}$這是保羅自己生命的經歷.
這是二章一到十一節.
我們上星期看了.
也想和大家看看.
因為保羅提到神偏待人.
我在新約也幫大家做了一些功課.
聖經也有很多講.
至少我找了四段.
我們先稍為揭一揭這裡.
和保羅講的這裡有些關連.
保羅強調.
在整個萬民都在罪惡.
在神的憤怒.
在需要救恩底下.
特別強調神不會偏待任何一個人.
所以沒有人可以靠身份.
沒有人可以靠自義.
沒有人可以靠自己的理性,理想.
就可以得救.
這個不偏待人.
更加反映只要在救恩,在耶穌裡面.
而得救絕對不是一種廉價.
不是一種便宜.
而是耶穌負上極大的生命代價.
所以我們要看看這個不偏待人.
聖經如何看待這個課題.
或者當時幾位使徒.
他們如何引述這些經文.
我們揭一揭這裡的使徒行傳.
《哥尼留信主》.
兩段都是彼得和保羅講的.
比較特別的第十章《哥尼留信主》這件事.
大家如果在座有些不太熟悉使徒行傳.
大家不要介意我讀多一點.
讓大家知道背景.
希望這些歷史書大家能夠掌握.
因為其實很好看的使徒行傳.
你掌握到大概知道裡面.
路加醫生想要帶出的重要真理.
拆開剛剛上個星期.

$^{281}$利斯基穆斯講使徒行傳.
極度地極.
很透徹地很淺白地.
將使徒行傳做一個很握要的分析.
和傳講.
非常具體.
將路加醫生想要帶進的舊恩.
聖靈如何在那個時代工作.
如何引導教會.
引導使徒們的行進.
在舊恩的傳播機遇裡面.
這樣講出來的.
所以相當值得聽的.
大家可以在群組找到那個連結.
很短的半個小時.
而且是很精煉的.
利斯基穆斯的講道.
十章這裡說.
在該塞尼亞有一個人名叫戈尼劉.
特別強調他的身份.
他的社會地位.
和他心靈的狀態.
這裡說其實本來不是很匹配的.
因為又當警察.
又做到這麼高級.
做到指揮官.
白夫長即是有100個士兵在他的統管下.
意大利型的白夫長.
該塞尼亞是一個重羅馬駐軍的地方.
大家看《四道行傳》都見過.
保羅在該塞尼亞.
當時有千夫長將他救了.
才免於死亡.
因為猶太的暴徒要置他於死地.
該塞尼亞意大利的白夫長.
他是一個虔誠人.
他和他全家都是敬畏神.
相信敬畏神是我們猶太人耶和華.
多多的州制百姓常常禱告神.
又州制又禱告.

$^{321}$其實很難的.
因為他們是武官不是文官.
意大利型的外邦人.
強調他的國籍.
但他有一個很敬虔的對神的心.
還有具體的行動.
所以第三節.
有一天約在新初.
新初是下午三時.
他在異象中.
明明看見神的一個使者進去.
到他那裡.
見到光明的使者.
是白天的.
下午三時.
不是那些鬼怪.
半夜的夜談.
說鬼故事.
好像某個明星有陰陽眼.
經常見到那些鬼.
令他很害怕.
但不說是誰了.
施明.
白夫長戈尼留.
明明看見神的使者跟他說.
一定很害怕.
因為靈界是天使.
驚怕說.
主啊 什麼事呢.
天使說你的禱告和你的州制.
達到神面前已經無紀念.
即是說他做的善事.
都被神紀念.
當然.
這些善事做完.
他都還未得救.
他不是因為做善事得救.
是因為之後他接待彼得.
而彼得說要他們全家信耶穌.
他們才得救.

$^{361}$所以分得很清楚.
不過蒙神的垂聽.
蒙神的紀念.
就是他這個人很值得被神拯救.
神要向他顯現.
所以.
這些特殊的案例.
因為今天未必需要.
這樣去一個好人.
一個經常奉獻的人.
就可以見到使者.
不應該是這樣的.
很少這樣的情況.
要告訴我們.
他當然很害怕.
第五節天使告訴他.
現在你當打發人.
往約柏去.
稱呼彼得的西門.
即是西門彼得.
即是這位使徒.
他住在海邊的一個蕭皮匠.
西門的家裡.
他的房子在海邊.
當時沒有Google map.
告訴他.
去哪個房間.
哪個地方.
去找這個.
稱為西門彼得.
而且住在哪個區.
等等.
第七節向他說話的天使去後.
突然消失.
哥尼留就叫了他兩個家人.
和上上事後.
他的一些虔誠兵.
來到去.
連士兵都很虔誠.
就將這件事述說給他們知道.

$^{401}$就打發他們往約柏去.
跟著.
那件事就另一個鏡頭.
就是彼得.
在約柏的地方.
即是這個.
蕭皮匠.
即是那些牛皮.
的地方.
做那些皮革.
這個西門的家裡.
作客.
短暫逗留.
好像旅館一樣.
似乎這個西門都是信著.
彼得就去到這個地方.
九章結尾的時候.
彼得就周圍去.
探望弟兄姊妹.
在哥尼留.
很奇妙的就是.
神原來是移動的.
彼得在哥尼留.
就不用回到耶路撒冷.
可以在任何地方.
神都可以向哥尼留.
向彼得顯現.
這次就另一個角度.
就是.
神.
讓彼得見到異象.
而且這次不是讓天使.
一個不認識的人.
向一個未信的人.
顯現.
而是神向他的僕人彼得顯現.
吩咐他 命令他.
要接待外邦人.
給他一個異象.
才引起他的關注.

$^{441}$他覺得很奇怪.
為什麼會做夢.
看到這些污糟的東西.
在猶太人眼中的污糟東西.
不是我們說的那些鬼怪.
不乾淨的東西.
這些是不可以吃的.
不潔之物.
所以他沒有理由會做這些夢.
沒有理由他會肚餓到.
突然間很想吃魚蝦蟹.
猶太人自小不敢.
也不會這樣想.
因為知道那些是不准吃的.
他似乎真的沒有吃過.
有四足的走獸.
昆蟲.
天上的飛鳥.
接著就是第十三節.
有聲音向他說.
彼得起來.
載了吃.
彼得說.
主啊 這是不可的.
凡是俗物和不潔淨的物.
我從來都沒有吃過.
這是很真實的.
第二次有聲音再和他說.
神所潔淨的.
你不可以再當為.
俗物.
這樣一個異象.
給他一些思考.
為什麼會突然間做這樣的夢.
為什麼會突然間.
有這樣的念頭.
還有這樣的聲音告訴他.
要起來載了吃.
不過第二次就再具體一點.
神已經潔淨了.

$^{481}$你不可以再看為是足而不潔.
但是怎麼厲害都想不到.
很快就會有.
外邦人來找他.
就是說猶太人.
和外邦人.
真的沒有任何的來往.
所以令到他.
有一個很大的詫異.
不敢默默地.
隨便地.
接受這個邀請.
所以.
輾轉一路發展下去.
哥尼留的僕人已經訪問到.
西門的家.
第十八節.
喊著問有沒有稱呼彼得的西門.
住在這裡呢.
第十九節.
彼得還思想那異象的時候.
聖靈向他說.
有三個人來找你.
起來下去.
和他們同往.
不要疑惑.
因為是我猜他們來的.
聖靈.
路加醫生很快.
將兩件事完全.
將它連繫.
不是一個誤會.
不是哥尼留.
很想巴結彼得.
完全不是.
是聖靈在居中.
撮合.
哥尼留不知道是甚麼事.
因為沒有告訴他.
只是叫他找這個人.

$^{521}$似乎是大人物.
沒甚麼所謂.
由哥西尼亞去到.
西門的家.
但是.
彼得就有疑惑.
因為要他做一些.
非比尋常.
一般都不會做的事.
就是要接觸外邦人.
不過.
他應該會開始有些印象.
神所潔淨的.
你不要再看為不潔.
這個都是.
頗重要.
第十章很想告訴我們.
原來救恩.
臨到你認為不潔的.
外邦人.
神已經拯救了.
哥尼留.
所以日後任何的外邦人.
信主都不需要.
再有這些聖靈的.
matching(配合).
一個至關重要的.
真理就是.
神潔淨了的.
我們自以為清高的人.
彼得.
千萬不要再看為逐而不潔.
我們要在神的心腸.
去看.
其他和我們不同類別的人.
那種的.
心境.
所以第十章路加醫生想說出.
這件事.
輾轉輾轉對話後.

$^{561}$第三十四節.
我們先看看.
少少註腳.
就是三十節.
哥尼留見到彼得後.
彼得問你找我來做什麼呢?.
哥尼留說.
前四日在這個時候.
都是新初.
我在家中守著新初的禱告.
即是說原來我們自己.
固定有些時間祈禱.
都是蒙神的祝福.
我們信徒.
我們自己固定.
可能是早上,中午.
或者晚上等等.
都是頗重要的.
我都說過.
我們有本祈禱簿.
是很重要的.
寫下一些名字認真的祈禱.
原來有寫過.
感恩,他現在回到菲律賓.
我們在星期三晚期禱.
提到一位菲律賓的一位姐姐.
有一次.
因為有英國的姐妹叫我去探望.
去到屯門醫院的時候.
她已經昏迷了.
昏睡了.
但完全不懂得和醫生.
醫護對話.
但輾轉.
我們祈禱了一段時間.
我們都放下了.
但上星期英國的姐妹就說.
她的情況已經可以.
對話.
醫生評估過她身體的狀況.

$^{601}$已經可以.
回到菲律賓.
所以在12號.
今個月的12號.
早幾天就回到菲律賓.
回到菲律賓之前.
我都有機會去探望她.
她坐在那裡很有感情.
坐在那裡.
都有聊天.
我都問了她和她拍了照.
真是很看到神的作為.
所以弟兄姊妹.
我們要學習祈禱.
哥尼留未信耶穌.
他守著.
新初的禱告.
他有定時的.
祈禱.
當然我們不要因為祈禱.
要等候天使.
向你顯現.
敬拜禱告.
他就說忽然之間.
有一個人穿著光明的衣裳.
站在我面前說.
哥尼留你的禱告.
以蒙誰聽.
你的州際達到神面前.
以蒙紀念.
你當打發人往約柏去.
請那稱呼彼得的西門來.
他住在海邊一個燒皮像.
西門的家裡.
所以我立時打發人去.
請你你來了很好.
現今我們.
哥尼留就叫了他的家人.
和他的親屬士兵.
一個準備好的報道會.

$^{641}$不過還沒有講完.
因為講完還沒有到.
那些會眾全部是哥尼留請的.
哥尼留自己都還沒有信耶穌.
所以是一個很戲劇性的聚會.
所以彼得其實沒有預備講章.
因為他根本沒有想過要做什麼.
所以這個時候就說.
所以我立時打發人去.
叫你來.
現在我們都在這裡.
現今我們都在神的面前.
要聽主所吩咐你的一切.
接著34節就是我們想引的經文.
彼得就開口說.
我真看出神是不偏待人.
原來各國中那敬畏主行義的人.
都為主所悅納.
神藉著耶穌基督.
傳和平的福音.
將這道次給以色列人.
這話在約翰宣傳.
朕例以後.
從加利利起.
傳遍了猶太.
以聖靈和能力.
告拿撒勒人耶穌.
這都是你們知道的.
他周流四方行善事.
醫好凡被魔鬼壓制的人.
因為神與他.
他就是耶穌.
同在他在猶太人之地.
並耶路撒冷所行的一切事.
有我們作見證.
他們警罷.
他們就是猶太人.
警罷他掛在木頭上殺了.
第三日.
神叫他復活.

$^{681}$顯現出來.
不是顯現給眾人看.
乃是顯現給神.
如先所揀選.
為他作見證的人看.
就是我們這些.
在他從死裡復活以後.
和他同吃同喝的人.
他吩咐我們全部給眾人.
證明他是神所立定的.
要作審判活人死人的主.
眾先知也為他作見證說.
凡信他.
的人必因他的名.
得蒙赦罪.
彼得還在說的時候.
聖靈就已經呼召.
降在一切聽道的人身上.
那些奉割禮.
和彼得同來的信徒.
見聖靈的恩賜.
也囂在外邦人身上.
就都稀奇了.
因聽見.
全部人.
說謊言.
稱讚神為大.
這個哥利留他們.
於是彼得就說.
這些人既受了聖靈.
和我們一樣.
誰能夠禁止用水.
吸他們的屍浸呢.
就吩咐奉耶穌基督的名.
吸他們的屍浸.
馬上洗禮.
完全不用問心事.
不知道有多少人.
這裡可能都超過幾十人.
就再和他們住了幾天.

$^{721}$就說.
就再和他們住了幾天.
當然這件事是有後著.
因為耶路撒冷.
聽到彼得進去找外邦人.
又帶他們信耶穌洗禮.
就很大件事.
就要調查.
審問彼得.
為什麼要這樣做.
和外邦人來往.
彼得其實帶了六個弟兄.
都是信徒.
猶太人.
有七個人.
見證著聖靈奇妙的工作.
所以34節.
彼得所說.
他很清楚.
其實他之前很疑惑.
應不應該接見外邦人.
外邦人還是.
羅馬的白夫長.
意大利人.
我和他們來往.
可不可以呢.
果然是不可以.
因為回到上元耶路撒冷.
十一章就是解釋.
所以我們也要學習.
彼得在此.
發自內心的讚歎.
神是不偏待人.
原來這個哥尼留.
不知道他在哪裡有的神學.
讓他成為一個敬虔的白夫長.
他自己在家裡.
守著禱告.
有周濟窮人.
有行善.

$^{761}$等等.
原來神已經閱立了這個人.
所以我們要從這個角度去看.
福音的偉大.
福音的無遠佛界.
對不同的人種.
對不同民族的拯救.
盼望我們對於福音.
能夠臨到不同民族.
不知何時我們成為修葛者.
因為神是不偏待人的.
可能是等待我們.
感動向人說.
聖靈在他的生命工作.
所以這是很值得我們振奮的地方.
在彼得面前這麼簡單的說道就順著.
第一個就在這裡.
接著我們揭一揭.
二不所書六章.
這五段都是保羅用的.
羅馬書固然在剛才說的那裡.
二不所書.
這是很重要的倫理教導.
六章.
由第五節讀起.
因為在教會裡面.
我們面對世界生活.
六章五節開始.
你們作僕人的要懼怕戰經.
用誠實的心聽從你們肉身的主人.
保羅沒有刻意將那些倫理.
不公平不公義的奴隸制度.
用福音或群眾的壓力.
去立即取消.
保羅沒有這樣做.
因為當時那個奴隸.
普遍成為了社會的勞動力.
如果輕易地這樣做.
等於基督教很快就會倒下.
羅馬帝國也不會放過你.

$^{801}$所以不是用強權.
輿論.
壓力去推倒羅馬帝國.
這個絕對不完善.
但保羅在這裡講回實際的情況.
你看到之後你會發現.
保羅不是服從帝國的.
歌功頌德.
他在說原來在神裡面有更高的標準.
只要大家願意這樣做.
問題反而未必是.
立即要將奴隸制度廢除.
原來我們作僕人.
要好像聽從基督一般.
要像聽從基督一般.
聽從肉身的主人.
不是要為帝國效力.
每個人都有不同的.
在神面前的身份.
那個位份.
那個人生的不同階級.
第六節.
不要只在眼前事奉.
像是土人喜歡.
要像基督的僕人.
從心裡遵行神的旨意.
基督的僕人.
原來我們不是只是做了基本的.
九點至五點放工.
探訪?明天才開始.
快要離樓?不要緊,明天才開始.
我九點至五點.
這樣是不行的.
所以僕者不是九點至五點.
我們要記住這一點.
像是土人喜歡.
今心事奉.
好像服侍主.
不像服侍人.
因為曉得各人所行的善事.

$^{841}$不論是為奴.
這很清楚了.
奴隸是自由人.
或者有一定的身份.
社會地位,財富.
的一些主人.
這樣.
都必按所行的得主的賞賜.
第九節.
你們作主人的待僕人.
也是一理.
不要威嚇他們.
因為知道他們和你們.
同有一位主在天上.
他並不偏待人.
這裡是說.
倫理關係的將心比己.
而且不單是地上的將心比己.
即是你對我好,我對你好.
而是一直我們要認定.
我們有一位主在天上.
他有他的主權.
他有他的分配.
他安排我們人生不同的崗位.
那個恩賜.
我們要在那個馴服下.
發揮到淋漓盡致.
一個敬虔.
愛主的奴隸.
一個敬虔.
服侍主的有錢人.
主人.
這是一件很美的事.
不自覺地.
那些不完善的制度.
都自己瓦解.
所以.
一段經文.
保羅不是不知道.
現在是奴隸制度的強權壓迫.

$^{881}$但是保羅更加知道.
不是隨便.
用他的輿論.
用他的號召.
就可以將帝國的不完善的制度.
馬上推倒.
他從來不是這樣.
這些都是智慧.
這些都是保羅對世情的那種透析.
和智慧的運用.
聖經的教導.
所以重點是.
大家都知道.
有一個主.
同一個主.
所以這裡很寶貴.
神不偏待人.
如果你是奴隸.
如果你是主人.
按照自己的形態.
對準主.
去對待其他人.
所以怎樣對待別人.
都是一個生命的功課.
接著到歌羅西書.
最後這一段.
就是歌羅西書之後.
三章同樣都是倫理的教導.
我們可以多寫一點.
由十二節開始.
三章.
可能會太長.
還是不要了.
短一點.
由十八節開始.
三章.
十二節.
短一點.
由十八節開始.
三章十八節.

$^{921}$你們作妻子.
當信服自己的丈夫.
這在主裡面是相宜的.
你們作丈夫.
要愛你們的妻子.
不可輔待他們.
某程度反映當時社會的.
女性只是一個物件.
是男性的附屬品.
沒有社會地位.
或者有些人過份了.
變成女權.
未必是對的.
所以將當時的.
主流的社會階層.
說出來.
但是這裡夾雜了.
丈夫要有愛.
妻子要有信服.
這樣來匹配.
又成為和諧家庭的典範.
又成為可以打破.
當時男性主導.
帶領女性的.
女性的.
男性主導.
帶來的.
這裡提到是輔待.
對女性可以.
隨便的呼喝.
毆打.
現在印度還是這樣的.
因為他們的種姓制度.
女性的社會地位的.
低落.
那種難以想像.
嫁妝不夠就燒死.
已經不人道了.
但是.
古代.

$^{961}$不會少於這些.
新約的出現.
保羅對於.
當時的信徒的.
教訓能夠很淋漓盡致.
既沒有.
將當時的.
不完善的.
已經積雜了這麼多年的制度.
或者那些傳統.
一下子將它推倒.
而是讓人的.
生命因為主耶穌.
而帶來的.
新的展現.
那種力量.
然後作兒女的.
要凡事聽從父母.
因為這是主所喜悅.
你們作父親.
不要惹兒女的氣.
恐怕他們失了志氣.
即是.
隨隨便便呼喝.
罵一句到尾.
粗口.
一巴掌就掌摑過去.
在街上.
有時看到我也受不了.
對一個小朋友已經在哭.
媽媽可能太激氣.
又未相信耶穌.
可能相信也不知道.
一巴掌就掌摑過去.
那種難受.
那種侮辱.
他自己不知道.
他已經很火遮眼.
對那個小朋友的傷害.
是會很深.

$^{1001}$那種仇恨.
激發他內裡的.
胡亂的想法.
其實是很慘的.
所以這裡說.
惹兒女的氣.
恐怕他們失了志氣.
我們小時候就好一點.
因為重男輕女.
男的怎麼錯都不會罵.
女的就不行了.
真的粗口事奉.
罵粗口.
所以.
你可以想像到.
二千年前已經教我們.
作為父親.
他不是你的擁有.
他是神給你的一個產業.
他是神託付給你的.
你只不過是一個信託人.
你是一個信託人.
信託給你.
不是給你去操控.
你逼他去讀醫.
他喜歡讀其他的,喜歡讀藝術.
你就是要他去做醫生.
很多這些的錯配.
這個就是.
惹兒女的氣的後果.
然後就作僕人的.
要凡事聽從你們肉身的主人.
不要只在眼前事奉.
像是土人喜歡.
為什麼又再說一次呢.
因為兩本書是寫給不同的地方.
哥羅西教會和爾弗所.
相距都很遠的.
所以.
保羅在這裡教導.

$^{1041}$就是說當時的社會的急切需要.
其實都大概是主僕關係.
家庭的關係.
等等的.
所以如果能夠.
這些倫理都能夠梳理得好.
一個家庭.
都是這樣愛主敬權.
羅馬帝國就會被推倒.
那些不完善的制度.
就會得以瓦解.
所以你可以想像到.
保羅都是很清楚.
講完.
他就說.
講完那些教義之後.
在實際的生活的倫理.
就是這樣去勉勵.
信了主的弟兄姊妹.
當時在歐洲.
保羅的年代.
有一些斯巴達人.
是很強悍的.
是很凶神惡煞的.
是還不講理的.
所以.
男士隨時就會突然暴露.
一點都不奇怪.
中國人都會.
我們上一代已經是.
隨時會暴露的.
很暴烈.
很憤怒.
所以如果一個男士.
信了主耶穌.
而對他的太太.
對他的小朋友.
對他的僕人.
是很溫文有禮.
是很有同理心.

$^{1081}$那就是福音的德性.
是福音的偉大的改變.
所以保羅看重的是.
每個每個家庭的生命改變.
很多人.
他們的生命有了一份同理心.
而是來自基督的救恩.
所以.
這裡說23節.
無論作什麼都要從心裡面作.
像是給主作.
不是給人作.
因你們知道從主那裡.
必得著基業為賞賜.
即是我們一生.
做在主裡面的.
不要問那個回報.
因為將來就會有基業作賞賜.
我們一心.
為主去做.
不過留意最後那句.
你們所事奉的乃是主基督.
那行不義.
必受不義的報應.
主並不偏待人.
突然間加了這個負面的這一句.
就抵銷了.
所有你所做的事.
即是說如果我們一心.
以為主不知而行不義.
將來就會有一個大的審判.
的報應.
一直說是做那麼多.
本來都很仔細.
即是他都很有.
那種社會的sense.
都是不外乎.
工作的場景.
家庭的場景.
倫理的.

$^{1121}$掌有的場景.
人類的組合.
就是那麼多.
你只需要這樣教導.
重要的邏輯.
原則.
就能夠讓.
基督的精兵在不同的.
家庭裡面來孕育.
大家都是為.
主而活.
這樣就會帶來.
羅馬帝國慢慢.
就會有改變.
好像那些白蟻日積月累.
吃著吃著.
整間屋就會倒塌.
當然我們不是白蟻.
我們是基督的僕人.
所以最後那句就很恐怖.
不要以為.
說笑的.
保羅突然間這樣說.
走不義的.
說那麼長都是正面的.
突然間加這一句.
他很肯定將來會有審判.
所以我們真的.
每天都要靠主的救恩.
來面對.
我們自己做得不夠好的地方.
我們要.
不要走精面.
不要以為神不知道.
因為他並不偏待人.
保羅再一次說這句話.
羅馬書說過.
以弗所書說過.
現在歌羅西書又繼續說.
都是在監獄裡面寫的信.

$^{1161}$在林中又說.
彼得前書記一章.
就是說到.
對永恆的盼望.
因為在苦難裡面.
他有一個信.
他說.
盼望.
因為在苦難裡面.
教會已經進入苦難.
如果用一個稍微的時代性.
耶穌離開世界.
當然有逼迫.
但後來羅馬帝國的逼迫.
不是每天都逼迫.
也有一段時間是太平盛世.
所以容許你們去傳道.
等等.
就是救恩傳播的一個很蓬勃的時期.
但是奈何.
到了尼祿開始.
尼祿王開始.
那個苦難.
又越來越明顯.
白熱化.
水面化.
不是暗地裡.
是公然的.
針對基督教.
想將他嫁禍給他.
使自己的.
政治目的.
能夠達到.
尼祿王據聞.
他很想建一個新的.
一個.
宮殿.
一個新的羅馬城.
舊的城.
還未壞.

$^{1201}$放火燒了它.
然後嫁禍給基督徒.
所以.
可以轉移視線.
將矛頭指向基督教.
去迫害.
所以保羅在尼祿下.
他就真的要去.
坐監殉道.
被斬頭.
在這裡告訴我們.
在彼得要寫這卷書的時候.
就是黑暗時期的開始.
他要勉勵.
在苦難裡面的信徒.
教會.
要怎樣去等候.
永恆的盼望.
第十三節.
一章十三節.
所以要約束你們的心.
謹慎自守.
專心盼望.
耶穌基督顯現的時候.
所帶來給你們的恩.
專心的盼望.
即是說我們有一個很強大.
忍受苦難的力量.
就是.
專心盼望.
主耶穌的降臨.
祂的顯現.
這就是在多困難下的信徒.
內心的光明.
就是因為我們相信.
主耶穌回來了.
十四節.
順命的兒女.
不要效法從前.
蒙昧無知的時候.

$^{1241}$放縱私慾的樣子.
那照你們的既是聖潔.
你們在一切所行的事上.
也要聖潔.
因為經上記著說.
這是利未記.
你們要聖潔.
因為我是聖潔的.
接著第十七節.
你們既稱那不偏待人.
按各人行為審判.
人的主為父.
就當傳敬畏的心.
到你們在世寄居的日子.
知道你們得贖.
即將來身體的解除.
去到永恆.
脫去你們祖宗所存留.
虛妄的行為.
不是憑著能壞的金銀等物.
乃是憑著基督的寶血.
如同無瑕疵.
無玷污.
羔羊之血.
下面的我們不再說了.
是想一路帶進.
對主的永恆的救恩.
祂再回來的盼望.
去.
不是抵消.
是去忍耐.
地上現在進入.
黑暗的苦難時期.
成為很大的力量.
彼得就說.
這位審判人.
不偏待人的.
是我們永恆的父.
是我們天上的父.
所以我們要帶著.

$^{1281}$這樣的敬畏的心靈.
在世寄居的日子.
去渡過.
在世寄居.
一切都會過去.
甚麼都不會帶走.
我們活在當下.
和在主的面前.
全世界的.
所有的.
在主的面前.
傳一個敬畏的心靈.
仍然帶著.
平安的享用.
享受每一日的.
生活的機遇.
不同的.
我們人生的責任.
等等.
剛才我們吃飯的時候.
65歲後.
不要只顧著工作.
開始要去體驗一下.
做一些平時.
人生的責任.
做不到的事.
例如去一些地方旅行.
去享受勞苦工作的成果.
不要等到80歲.
走都走不到就不能去.
這是一個享受.
一個敬畏的心靈.
去渡過.
我們寄居的日子.
所以寄居的日子.
不是等於一定要去修道院.
等等.
職場都可以是我們的修道院.
帶著一份享受.
神給我們的.

$^{1321}$祝福的感恩.
都是一個寄居的日子.
當然每天都是為玩而玩.
要問自己.
是否真的看這個世界是寄居呢?.
如果是寄居.
我們做的事情.
心態是不同的.
因為我們無論去多好玩的地方.
都會找一個地方.
天家就是我們永恆的家.
地上就是短暫的家.
所以我們要記得.
地上不是我們可以.
完全擁有的.
我們看了這裡.
幫大家.
稍為將這些重要的經文.
在新約裡面.
大家去掀開.
豐富一下使徒們.
如何引用這些經文.
他們對做物主.
這麼重要的.
他的屬性.
他將來要行使的審判.
我們要帶著敬畏.
信了主耶穌.
不是一了百了.
洗禮就不是大了.
我們是要傳敬畏的心.
我們要有一個同理心.
對不同的家庭的家人.
工作的同事.
社會上對不同的人.
我們要有一份.
不偏待人的心.
這樣敬虔的成長.
所以有時很多基督徒.
被人感覺到.

$^{1361}$又是很嘎嘎的.
對人又沒有禮貌.
又惡.
其實就反映不了我們.
懼怕神.
反映不了我們有一個敬畏的心.
知道他是不會偏待人的.
那一位主.
所以這樣揭曉.
自己也有豐富的領受.
我們繼續下去.
《羅馬書》第二章.
11至13節.
其實一直都是說.
外表的道德.
理想的.
以為自己達到的行為.
都是在神的面前.
是不可救的.
這是第二類的人.
就是這些.
似乎在當時.
有一定道德水平的.
道德家.
哲學家.
甚至勸導人.
要這樣效法他們的生活.
這樣的形態.
但是在保羅的眼裡.
如果不是.
在救恩的.
拯救也不行.
因為第14節.
沒有律法的外邦人.
如果能夠順照著.
他本性行律法上的事.
雖然他們表面上.
沒有律法.
自己就是自己的律法.
但是奈何13節.

$^{1401}$沒有人能夠做到.
因為不是說.
聽律法的來為義.
是要實踐.
行律法為義.
所以沒有人能夠做到.
15節說.
「這是顯出律法的功用.
刻在他們心裡.
他們是非之心同作見證.
並且他們的思念.
互相較量.
或以為是或以為非」.
到底第15節是甚麼意思呢?.
擺在整個上文.
就是道德家.
他們有一套.
自己認為的標準.
但是實際上.
在我們內心.
互相較量的判斷.
其實我們.
都是應該.
苦心自問.
做不到神的標準.
所以沒有人.
能夠達到.
接著就是.
就在神.
接耶穌基督審判人.
隱蔽時的日子.
照著我的福音所言.
相當難解.
甚麼會突然帶出福音呢?.
福音的確是在說審判.
福音的基礎就是.
耶穌基督被審訊.
被釘死.
被定罪.
就是這樣.

$^{1441}$帶出他一生的義行.
就可以.
不是轉嫁.
是完全能夠給我們.
得到神的義.
而且.
福音是叫人.
現在今生要信.
否則.
將來神就會透過.
耶穌基督.
來審判世人.
甚至乎是.
每個人隱蔽的事.
所以.
沒有人可以靠.
自己的義行.
道德標準.
來得救.
福音不是叫人靠自己的義.
去得救.
福音是藉著耶穌基督的審判.
去得救.
所以主耶穌先為我們受審判.
是他的一生義行.
滿足了神的要求.
我們就可以得到神的拯救.
所以再回來.
都是在保羅所說的福音裡面.
所以福音不是廉價的.
你要吧.
不要吧,我真是很可惜.
福音是一個控訴.
福音是一個審判.
福音是一個要人改變.
是一個要人.
趕快悔改.
承認自己的罪.
在救恩耶穌的愛裡面.
進入與神和好的關係.

$^{1481}$所以.
不要單是說.
福音是很好的,很有愛的.
你有病嗎?沒有問題.
祈禱神一定赦免醫治你的.
福音不是一個萬靈丹.
福音不是叫我們.
信了耶穌.
你的小朋友讀書又有學位.
將來又怎樣好.
福音不是單是這樣.
那些全部是副產品.
是生命的改變.
帶來的很多和諧.
福音其實是說.
主耶穌為我們.
進入死亡.
他接受了.
無情的審判.
為我們承擔.
十字架上.
我們一生的罪惡.
世人的罪惡.
所以保羅在這裡.
說的是.
內心的拯救.
沒有人可以達到.
所以是藉著.
耶穌基督審判人.
隱蔽時的日子.
照著他的福音所言.
所以保羅傳的福音.
是一個要人悔改.
承認.
而拯救的救恩.
所以我們千萬不要以為.
我們將救恩減低.
不要太殺傷力.
什麼都可以.
有耶穌就什麼都可以.

$^{1521}$沒有問題.
當他發現有問題的時候.
他就覺得福音是騙他.
是假的.
福音其實不是這樣.
這就是第一類.
當然上文一至十八節.
到結尾的.
三十二節.
是說那些.
否認神的人.
他們走的種種惡行.
罪惡.
那一類的人.
第二類就是.
很有道德.
很有他們自己的.
理想的人.
但是達不到.
神的意義水平.
先停一停.
喝一口水.
接著我們就進入猶太人.
我們要問的問題.
就是.
非常獨特的民族.
救恩都是.
由耶路撒冷.
由耶穌開始.
由使徒們開始.
由保羅自己.
更加鹹淡水.
穿梭在舊約的世界.
的法利賽人.
接著戲劇性.
見到耶穌.
接著成為外邦人的使徒.
現在他要寫一封信.
充分反映.
萬民都需要救恩的時候.

$^{1561}$他怎樣針對自己的民族呢.
他有沒有偏見呢.
還是持平呢.
我們就要進入.
這一段.
他對猶太人的批判.
當然我們要問的問題.
就是.
因為在第一章十六十七節.
和二章的第九節.
都不斷在說.
猶太人.
都不斷在說.
這個概念.
先是猶太人.
後是希利利人.
這個概念.
大家要學會一個新約的.
常用的詞彙.
就是泛指所有的人類.
的意思.
猶太人和非猶太人.
用當時希臘.
當時的主要文明.
文化.
希利利人就是泛指.
所有的猶太.
非猶太種族的人.
就是非希利利人.
所以.
你讀《小徒行傳》.
就會見得多.
但保羅的《羅馬書》都有的.
已經有兩次了.
一章十六和二章的第九節.
已經有了.
但是來到第十七節.
這裡就是一個獨特的身份.
就是單單指.
猶太民族.

$^{1601}$這個民族.
又是用一個單數.
好像一個.
舉例.
這個人.
在這樣的形態下.
反照所有.
現在在耶路撒冷.
或者在周邊不同的地方.
居住的猶太人.
都是這樣的想法.
這樣的.
防護盾.
對舊恩的那種抗拒.
那種保護膜.
的心態.
所以保羅就要.
拆解他們的困難.
原來律法.
是他們的保障.
是他們的一個保護盾.
保羅就在這裡.
列舉.
第十七節.
我們要看看將這一段.
我先讀一讀.
因為我們上次沒有讀過.
十一到第十六節.
大家如果有問題.
我不再重複解釋.
有問題想不通邏輯.
就在群組私下.
WhatsApp我.
我盡量去解釋.
我們進入第三個類別.
就是一定要說猶太人.
因為他們.
到今天都還認為自己.
是亞伯拉罕的子孫.
是大衛的子孫.

$^{1641}$是律法之子.
是有神納若的一班人.
這個獨特性.
如何將他理解呢.
這是一個大課題.
去到第九到第十一章.
保羅要特別將猶太人.
為什麼不信耶穌.
很詳細地去拆解.
所以在這裡.
這個段落.
仍然想在大段落裡.
是什麼段落呢.
就是一章十八.
至三章的二十節.
問耶穌.
到三章的二十節.
萬民都規缺神的義.
萬民都需要.
恩信稱義.
這個大課題.
所以二章的十七節.
到二十九節.
都是在說猶太人.
都規缺神的義.
猶太人都需要.
救恩.
你稱為猶太人.
這是十七節.
又倚靠律法.
誇口.
既從律法中受了教訓.
就曉得神的旨意.
也能夠分別是非.
原文就是喜愛美好的事.
又深信自己是.
給核子領路.
黑暗中人的光.
是蠢笨人的師父.
是小孩子的先生.

$^{1681}$在律法上.
是有知識和真理的模範.
你既是教導別人.
還不教導自己嗎?.
你講說人不可偷竊.
自己還偷竊嗎?.
你說人不可姦淫.
自己還姦淫嗎?.
你厭惡偶像.
自己還偷竊妙中之物嗎?.
你指著律法誇口.
自己倒犯律法.
玷污神嗎?.
神的名在外邦人中.
因你們受了褻瀆.
正如經上所記.
這裡我們稍後會揭曉.
二賽亞書的五十二章.
就是被擄時期.
他們的民族.
落得很痛苦的審判.
第25節.
你若是行律法.
國禮固然於你有益.
若是犯律法.
你的國禮就算不得國禮.
所以那未受國禮的.
若是犯了律法.
所以那未受國禮的.
若遵守律法的條例.
他雖然未受國禮.
豈不算是有國禮嗎?.
而且那本來未受國禮的.
若能夠全守律法.
豈不要審判.
你這些有疑問.
疑問即是表面的.
那個的律法.
例如他們守安息日.
等等.

$^{1721}$豈不要審判.
你這有疑問的.
和國禮.
竟犯律法的人.
因為外面作猶太人的.
不是真猶太人.
外面肉身的國禮.
也不是真國禮.
唯有裡面作的.
才是真猶太人.
真國禮也是心裡的.
在乎靈.
不在乎疑問.
這人的稱讚.
不是從人來的.
乃是從神來的.
將這段聖經.
做了兩個.
很簡單的分段.
17到29節.
他們有.
律法的教訓.
但違反了律法的教訓.
即是說.
很明顯十誡.
他們都是肆無忌憚.
去破戒.
去犯這些.
違反神的律法.
的心意.
的事情.
特別強調.
保羅在這裡引用了幾個例子.
原來我看一些.
的.
解經家的.
引述.
原來在歷史裡.
是發生過.
在主後的19年.

$^{1761}$曾經有.
當時是該察第二任.
就是提比劉.
約翰出世的時候.
提比劉.
就是奧古斯都.
的接任人.
接著有些.
猶太人因為偷了.
聖殿的錢.
被.
州制的窮人.
的錢被提比劉.
驅逐離開羅馬.
有這樣的案例.
記載在.
約瑟夫的.
猶太的史記裡.
你可以想像到.
保羅在這裡不是.
隨便說的.
是有一些案例的根據.
保羅在引述的一些案例.
真是有這樣的情況.
在外邦人當中.
神的名.
被褻瀆.
因為猶太人的緣故.
所以保羅在這裡.
強調.
如果17到24節.
這一段他們是有.
依靠律法.
誇口在律法裡.
受了教訓.
但是能夠.
說就不能夠實際.
在生活.
生命裡面去實踐.
經常性地實踐.

$^{1801}$即是說.
是有一種.
頗掩耳盜鈴.
自己都已經慢慢麻木了.
這樣做.
所以.
可能他真是.
以一個假設性.
指出某一個人.
其實那個人未必存在.
但亦都是說.
這樣的狀況.
那麼有道德律.
那麼有神的律法.
摩西五經十誡.
但是都不能夠.
因為這樣而可以.
逃避神將來的審判.
因為.
他們都是.
在神的面前.
甚至乎連神的名字.
都因為他們自己.
因為他們自己.
這些獨特的.
行事為人.
猶太人.
而引起外邦人.
特別是羅馬當時的帝國.
給他們有一定的信仰.
宗教上的.
自由.
自己宗教的禮儀.
但是.
在一般.
人的標準裡面.
都達不到.
帝國的要求.
都被驅逐.
可想而知.

$^{1841}$在神的律法.
就更加達不到.
保羅想帶出這樣的含義.
即是說.
就算有律法.
似乎就好悲哀.
律法有什麼用呢.
這章就會講解.
律法的功用.
猶太人他們帶出.
律法的功用.
但是他們自己也未能夠在這一段.
做到神的要求.
連基本人.
對他們的期望.
都達不到.
我們看看24節.
引述的是.
以賽亞書那段經文.
52章.
之後會再引述.
10章講到.
報福音傳喜訊.
腳中何等皆美.
那段經文.
52章其實就是受苦的.
尼賽亞段落的開始.
52章.
錫安娜.
興起.
披上你的能力.
聖成耶路撒冷.
穿上你華美的衣服.
因為從今以後.
未受割禮.
不潔淨的.
必不再進入你中間耶路撒冷.
當然那個城市是死物.
不會跳起.
一定是講裡面的人.

$^{1881}$就是說.
神會將那個聖潔.
將那個榮耀.
帶回去這個.
現在已經是被老的.
那個世代.
將來會有這種踴躍.
這種榮耀.
他們可以.
抖下塵土.
起來坐在位上.
錫安被老的居民.
要解開你頸項的鎖鏈.
也就是說.
他們都是因為.
他們自己的罪惡.
受到神的刑罰.
就是被老.
所以在這個被老的光景.
這裡有一個很好的盼望.
第三節.
耶和華如此說.
無價被賣.
也必無銀被贖.
主耶和華如此說.
豈先我的百姓下到埃及.
在那裡寄居.
又有亞述人無故欺壓他們.
耶和華說.
我的百姓既是無價被老去.
如今我在這裡.
作什麼呢.
耶和華說.
乞濟他們的人呼叫.
我的名整天受褻瀆.
所以我的百姓.
必知道我的名.
到那日他們必知道.
說者說的就是我.
看那是我.

$^{1921}$接著就帶出一個.
很重要的信息.
就是要揚聲.
報福音報佳音.
第七節.
這裡下半段的.
二丈亞書52章.
就在羅馬書第十章出現.
所以保羅就將這兩個段落分開.
在第二章.
這裡說的猶太人.
整個民族.
我看到圖就是.
在神面前的羞恥.
52章.
一到七節.
整個氣氛.
所以看保羅的引述.
整個氣氛就是.
在被老底下.
在巴比倫人面前.
你就是成為一個羞恥.
成為一個笑談.
說你的神在哪裡.
你們就是這樣.
將神的名字.
被人嘲笑.
聖殿的器皿.
在致聖所的器皿.
去了哪裡.
就放在巴比倫的寶物庫.
後來.
以色列回來.
所羅巴巴.
就拿回這些器皿.
聖潔的器皿.
回到耶路撒冷.
來供奉.
即是說連聖殿.
連耶路撒冷.

$^{1961}$連聖城都被外邦人踐踏.
猶太人要被老.
所以保羅將整個.
被老時期.
在二章這裡用.
他來論述.
猶太人.
其實你不要因為.
自己是猶太人就很得意.
其實你在外邦人眼中.
其實是一個shame.
保羅的意思就是這樣.
還有真的有些案例.
就是提比留的.
19年.
將這些猶太人驅逐出境.
所以.
你可以想像得到.
保羅.
不是幸災樂禍.
因為他自己也是猶太人.
他想說什麼呢.
以為有律法.
就可以成為一個保護網.
成為一個.
免疫.
是不可能的.
保羅都把自己放在.
罪人的行列.
保羅都看自己在基督的面前.
其實他也是.
要主的救恩.
他才可以得救.
不是因為他的學識.
不是因為他的法律地位.
不是因為他是一個法理錯人.
他就可以不需要.
而得救.
他需要主耶穌的拯救.
他需要亞拿利亞幫他洗禮.

$^{2001}$他成為基督徒.
成為主耶穌的.
憐憫的對象.
所以我們看這一段經文.
就是.
違反了律法的教訓.
我們看25到29節.
他們根本.
特別強調割禮.
就是說.
保羅在這裡.
突破了一件事.
猶太人獨特的身份.
絕對不是護身符.
羅馬帝國.
都是很懷柔的.
容許他們有.
自己的宗教的自治.
所以你要行割禮.
你看為這麼重要的第八天.
行割禮.
羅馬人是容許的.
不會禁止的.
因為這件事關係到滅絕.
如果我們在神面前不行割禮.
比殺了我們更嚴重.
所以羅馬人.
不和他們計較這些.
所以割禮成為了.
猶太人獨特的身份.
當然有趣的.
就是如果某些皇帝.
很支持割禮.
有很多外邦人.
都要行割禮.
因為有好處.
很多不信的人都知道.
讓子女進基督教學校是好的.
所以都很想他們.
進基督教學校有學分.

$^{2041}$又會有好的教導.
類似這些.
但是第二段.
保羅說.
雖然割禮是這麼明顯的.
猶太人被分別出來的一個記號.
一個立約的記號.
但是你們有.
割禮的愛在.
但是沒有.
割禮背後的.
真正的精神.
就是心靈裡.
進行神的旨意.
歌羅西書.
有一段這樣講割禮的.
不緊要.
今晚我回去找給大家.
我們先講了這段.
接回來羅馬書第二章.
最後這一段.
其實都.
很簡單容易明白.
不可以像.
守律法.
有獨特的.
生活形態的猶太人.
為什麼不講以色列人.
以色列人是地方的民族.
猶太人是被老之後的身份.
因為猶大人.
被老.
會認定自己的猶大族身份.
所以用這個.
作為他們的普遍.
在被老之後.
的獨特的.
身份來自居.
所以就算帝國的.
更替他們都是用.

$^{2081}$這個身份來自居.
希伯來就是在講.
阿伯拉鄉的時代.
在大河那邊.
可能是整個民族.
的神.
選照出來的希伯來人.
所以.
去到保羅的年代很少再用希伯來人.
一般是用猶太人.
以色列就講.
國家的地理的狀況.
還有另外一個.
就更少人用.
就是叫黎凡德人.
就是.
Levant那個字.
就是黎巴嫩.
一帶的地方的.
中東人.
再大一點的區域.
所以猶太人都屬於.
黎凡人.
或者再.
近代一點的就叫黎凡德人.
所以猶太人.
在保羅的時代.
是一個獨特的.
羅馬帝國.
底下的民族.
有自己堅持的宗教的.
法律.
有自己堅持的獨特的身份.
就是他們要行割禮.
他們要穿自己的.
特別的衣服.
他們要留鬍子.
他們有很多.
傳統的禮儀.
他們要這樣去做.

$^{2121}$來形成自己獨特的身份.
他認為這個獨特的身份.
就可以帶進.
永恆的救恩裡面.
或者神的喜悅裡面.
但是.
保羅在這裡.
都差不多完全打破了.
不是靠.
獨特的身份.
而是獨特的一個.
傳統信仰的.
表現.
而是心靈裡面.
所以這裡告訴我們.
他當然在說.
未受割禮的如果能夠遵守.
就不是說那些未受割禮.
就等於可以做到.
可能有的 例如哥利留沒有受割禮.
但是.
守著一個敬虔的生活.
但是那些是很特別的.
例子.
可以作為一個普遍性.
但是這裡再一次告訴我們.
第28節.
外面作猶太人.
不是真猶太人.
其實是真猶太人.
真是身份的血統.
是真的.
不過保羅在說的是.
心靈裡面對神.
原來都不認識割禮的意義.
是要帶進.
與神的相交的聖潔.
所以外面肉身的割禮.
肉身的割禮.
真是第八天.

$^{2161}$一個男孩割包皮的割禮.
但是這個.
是外在.
不是真的割禮.
第29節.
唯有裡面作才是真猶太人.
真的割禮.
是在乎心靈的.
不是在乎移民的.
人的稱讚.
就不是從人來.
原來受割禮是會加分的.
如果皇帝很抱猶太民族.
受割禮就變成高人一等.
所以.
這人的稱讚.
不是從人而來.
是從神而來.
所以我們.
回教會成為教徒.
18至21.
四天的浪漫之旅.
開始了.
不是北上.
是南上.
不是北上.
是壽苦節.
復活節.
唱聖詩等等.
在錦繡的居民.
看我們錦繡堂.
這一班人.
穿著詩袍.
敬拜朝見神.
有一個這麼獨特的生活形態.
年年都如是.
很多節日都是這樣.
他們是很嚮往.
對的.
他們是很清善的.

$^{2201}$因為有一個獨特的文化.
獨特的形態的生活.
不過要記住.
如果我們真的從心裡.
去認識.
我們所做的.
而我們做在神的裡面.
唯有裡面.
才是真猶太人.
真割禮都是從心裡.
在乎靈.
不在乎疑問.
所以我們得到的稱讚.
是從神而來.
所以我們不是否定.
做這些儀式.
做這些宗教禮儀的錯誤.
應該做.
但是原來做的.
背後是要得到神的喜悅.
以及.
不是要得到人的稱許.
是要從心裡.
去做這些事情.
所以有時.
做著做著.
會進入真理.
進入心靈與神的交契裡面.
那種相交.
所以如果施班.
當然是亂說.
原來這裡穿完袍.
很開心.
大家合一.
一揭開袍就走了.
平時是吵架的.
沒有來往的.
就不是做給神看.
如果人都發現.
對信仰就更加有打折扣.

$^{2241}$所以你看到這裡.
保羅不是想批判這個民族.
因為他自己也是猶太人.
他不是想要論斷這個民族.
而是以事論事.
這個民族既然有這麼獨特的宗教文化.
傳統.
甚至帝國都容許你的存在.
有這種自知式的宗教禮儀.
但是這些宗教禮儀.
不能讓你裡面的神.
讓人認識.
其實就是違反了律法的教訓.
即是犯罪.
違反了和神立真正的約的背後.
所以最後要揭一段經文.
都是相當重要.
耶利米書的32章.
以及以西結書的36章.
相當重要.
在乎靈不在乎疑問.
保羅就是在這裡透過耶利米書的啟示去領受.
所以保羅說的心靈的覺慮就在這裡.
32章耶利米書.
原來我剛才弄錯了.
是31章.
耶利米書的31章.
不是32章.
是31章的33節.
因為希伯來書有引述.
31章.
第31節.
耶和說:日子將到,我要與以色列家和猶大家另立新約.
不像我拉著他們祖宗的手.
令他們出埃及地的時候.
與他們所立的約.
我雖作他們的丈夫.
他們卻背了我的約.
這是耶和說的.
到底那個約是怎樣呢?.

$^{2281}$第33節.
耶和說:那些日子以後.
即是說一切的審判過了之後.
被老歸回.
接著有耶穌的降生之後.
那些日子以後.
我與以色列家所立的約.
乃是這樣.
就是那個約的條款.
我要將我的律法放在他們裡面.
寫在他們心上.
我要作他們的神.
他們要作我的子民.
他們各人不再教導自己的鄰舍.
和自己的弟兄說.
你該認識耶和說.
因為他們從最小的到最大的都必認識我.
我要赦免他們的罪孽.
不再紀念他們的罪惡.
這是耶和說的.
就是新約.
就是神將他的律法.
寫在信徒的心靈裡面.
相信這就是聖靈的工作.
讓我們從心裡面對神的愛.
和敬畏.
還要認定自己是立約的子民.
所以我們能夠最終認識耶和說.
不需要人教導.
當然我們不要否定.
我們要上查經班.
我們要學習.
不需要教導的意思就是.
心裡面已經有和神的相交.
那種默契.
那種愛的關係.
所以這是一個很重要的應許.
新約希伯來書引述這段經文.
就說聖靈在我們裡面工作.
所以和保羅第二章說的心靈的覺慮.

$^{2321}$差不多是異曲同工之妙.
保羅都是拿了這些靈感.
這些啟示.
來講到真正的神子民.
是心裡面的覺慮.
即是心裡面行神的旨意.
與神相交.
願意聖潔合神心意的生活的聖徒.
這才是真正的覺慮.
所以不要以為保羅自己突然發明出來.
其實舊約已經有這樣的啟示.
應許將來神的子民.
可以從心裡面.
所以一個成熟的教會.
就是很多弟兄姊妹.
從心裡面愛神.
敬畏神.
行神的旨意.
樂意去傳福音分享.
教會越成熟.
就應該越來越多這樣的弟兄姊妹.
我們一起禱告.
感恩我們將剛才第十七到二十九節.
簡單的做一個分享.
和來分析.
我們其實在幾個類別裡面.
我們都是被證明有罪的人.
多謝你將我們拯救.
帶進耶穌的救恩裡面.
如果不是主耶穌的兒行.
我們根本不能夠靠自己去得救.
多謝你讓我們能夠在主耶穌裡面.
得著的是神的義.
幫助我們能夠認清.
我們心裡面是在真正的覺慮.
我們是在聖靈的感動和呼召下.
來到地上成為主的見證人.
去見證這個救恩帶給我們的生命改變.
求主幫助我們這樣感恩禱告.
奉耶穌基督的名.

$^{2361}$阿們.
\newpage



\section{羅馬書}
\label{sec:b2w9fc63m3c}
\textbf{2025-04-24 周四晚查經班《羅馬書(8):分題:神的義(七)》 講員:曾裔貴牧師}
\newline
\newline
連結: \href{https://youtube.com/watch?v=b2w9fc63m3c}{\texttt{https://youtube.com/watch?v=b2w9fc63m3c}} ~~~~ 語音日期: 2025-04-25
\newline
\newline
\hyperref[sec:JU6MXStfi9g]{\small{< < < PREV SERMON < < <}}
~
\hyperref[sec:index_chronic]{\small{[返順時目]}}
~
\hyperref[sec:index_scriptual]{\small{[返順卷目]}}
~
\hyperref[sec:DvUFUpNSEFw]{\small{> > > NEXT SERMON > > >}}
\newline
\newline
$^{1}$我們一起同心禱告.
再一次感恩.
天父給我們晚上.
可以安靜的.
可以來到在線.
利用科技的方便.
我們可以來到同步.
去研讀神的話語.
羅馬書是一卷.
相當有結構.
相當重要將教義.
來到去詮釋的書卷.
幫助我們能夠.
進入保羅當時的.
世界的場景.
在羅馬帝國.
這個尼祿王.
作王的時代.
我們能夠掌握.
當時的背景.
他們的信仰.
在這個政治的.
氛圍帝國底下.
是怎樣來運作.
幫助我們能夠.
從保羅所.
帶給我們的視野.
對永恆的.
認識.
尤其是和神的關係.
建立的認識.
來掌握.
幫助我們這些永恆不變的.
能夠我們每一個.
的個人弟兄姊妹.
能夠蒙福.
能夠神的話語.
可以藉著不同的方式.
去傳揚開去.
使更加多人.

$^{41}$認識與神的關係.
今天晚上我們.
在你的面前.
我們這樣禱告.
奉耶穌基督的名.
阿們.
我們都差不多用了一堂多一點.
去講到第二章的.
羅馬書.
猶太人.
為什麼保羅要用一個特別的.
題目來講猶太人.
都需要救恩.
都需要.
不是靠律法.
而得救.
是靠律法.
他們已經是由阿巴拉罕.
祖宗的年代.
一直已經這樣來運作.
但是保羅現在要講的.
就是不是有另外一個救恩.
只是給外邦人.
而是一到第三章.
都是想告訴我們.
萬民都需要.
神的救恩.
萬民都需要恩順.
所以保羅先講了.
第一章.
外邦.
種種的藐視.
否認.
故意拒絕神的存在.
而帶來的.
種種的罪惡.
到第二章.
就來將一些.
可能表面很有教養.
修養道德.

$^{81}$或者有哲理哲學的人.
他們雖然.
不是像一些大奸大惡.
這樣的犯罪.
但他們的自義.
拒絕神.
都同樣在審判的行列.
去到第三類.
就是上個星期我們講的.
就是猶太人.
有他們的罪惡.
而且這個罪惡.
更容易被他們的.
宗教的外衣.
所包裹住.
因為他們有行割禮.
他們有每個星期的安息日.
有他們每個季度.
不同的節日.
的憲制.
所以他們.
好像當時的帝國.
當然不是每個皇帝.
都容許.
都給他們有一個.
宗教上的自治.
自由.
所以猶太人就有他們自己一套.
的生活方式.
或者一些食物.
哪些是乾淨可以吃.
哪些是摩西.
認為不潔不准吃.
他們都是這樣遵守.
所以很自然就會有一個.
自己一派.
自己一個.
在當時的帝國.
底下的一個.
和其他人是不相同.

$^{121}$這樣的不相同.
又會帶來.
一個分別.
自視過高.
自義.
這些都是在這樣的.
氛圍.
底下的群體.
是會獨有的問題.
所以我們上個禮拜.
提到二章的29節.
關於.
割禮.
我們再讀一次.
我們也做了一些資料.
上個禮拜不夠時間講.
和揭不出割禮.
在新約出現了很多次.
不過據象徵.
意義的.
屬靈的含義.
保羅在歌羅西書和.
菲勒比書等等.
包括《少奴行傳》.
都有提到.
關於割禮屬靈的含義.
稍後我會和大家揭一揭.
我們再回到.
整個羅馬書.
最重要的大綱.
既然福音是關乎萬民.
既然萬民.
都需要恩信稱義.
萬民都需要.
所有的民族.
都需要恩信稱義.
都需要神的救恩.
都是在神的憤怒審判下.
所以萬民.
都需要得到.

$^{161}$聽聞福音.
所以宣教肯定是.
保羅寫這本書.
最重要的目的.
所以救恩論,基督論.
聖靈論或者教會.
或者倫理關係.
那些全部要在.
一個大的Umbrella(圖案)下.
就是羅馬的教會.
必須支持.
保羅去推動.
福音工作.
其中就是要去到.
很遠的歐洲.
最南面的.
西班牙.
來宣教.
保羅很想羅馬教會.
有這樣的形式.
所以如果.
大家都對宣教有這樣的.
眼光.
有這個胸襟.
有這個神的.
真正的帝國神的國度.
這個胸襟的時候.
其他東西就很次要.
否則我們就會放大.
內部的矛盾.
所以很難怪.
說到倫理關係.
就出現了第十三章.
一些和帝國的.
相處.
應不應該納粹等等.
要不要馴服在上掌權.
第十四章就是.
內部更嚴重的.
宗教的矛盾.

$^{201}$應不應該守安息日.
有沒有一些食物是不可以吃的.
是不潔的.
就帶來.
猶太人和外邦人.
就會有這種.
可能是彼此.
觀望 但可能更嚴重.
就會引發.
一種衝突.
一門別類.
你一派 我一派.
是很可惜的.
很難能夠.
齊心合一.
去做宣教的工作.
因為大家都不同心.
所以在這個消耗戰裡面.
魔鬼都不用猜險.
我們出去.
我們都不會出去.
向外邦宣教.
所以這是最重要的.
中心思想.
神的義.
在第一章.
你可以說.
由十八節開始.
一直到四章結尾.
當然嚴格來說.
到五章都可以.
不過更多的.
關注.
一章的十八節至三章的二十節.
就是將.
罪的本質說出來.
將神的憤怒.
為何他會憤怒.
說出來.
就是萬民欠缺神的義.

$^{241}$哪些人可以得到神的義呢?.
由三章的二十一.
一直到四章的結尾.
就是神的義.
這份白白的恩典.
賜給.
相信主耶穌的人.
所以.
這個分段.
也是相當重要.
你看看這個簡單的分段.
萬民欠缺神的義.
一章十八節.
直到三十二節.
第二章.
就說到有兩類人.
表面有德行的人.
外表有義.
人與人之間相對的義.
不是真正神能夠稱許.
能夠閱立的義.
因為不是神本身的義.
他的本質的義.
第三章.
就說到.
一到二十節.
也是我們今天重點會說.
就是萬民.
都是被定罪在.
律法底下.
律法會使那些罪.
顯明出來.
所以透過律法.
我們知道什麼是罪.
我們有提過猶太人.
他們獨有的罪惡.
就是因為他們本身的身份.
是神設立這個民族.
由亞伯拉罕被呼召.
離開了吾爾.

$^{281}$加勒底的吾爾.
即今天的伊拉克.
一直到迦南地.
神賦予他們這個身份.
神要求他們在.
創世記十七章.
他們行割禮.
所以他們和其他民族.
是有不同的.
所以他們是一個很獨特的身份.
但問題是.
雖然他們有.
律法.
表面的律法.
表面他們都是很循規.
道矩的守一些節日.
守一些食物的條例.
或者.
有一些表面的.
和外邦人刻意不同的.
例如他們不會隨便.
去剃那些鬍子.
和他們的頭髮.
等等.
但是真正重要的.
心靈的.
同神的.
道德的.
分條.
有關係的.
心靈他們就未能夠.
完全滿足.
和滿足.
所以二章的十七.
這個我們上星期已經講過了.
就是他們違反了.
律法的教訓.
我們都引述過.
在.
蓋薩提比留的年間.

$^{321}$真的有猶太人.
偷那些善款.
即是穿櫃筒底.
被.
羅馬帝國驅逐.
離開意大利.
令到神的名.
被羞辱.
這個就是.
保羅很可能.
都知道當時的情況.
這樣去列舉.
這些例子是真有其事.
所以.
帶來的就是.
神的名.
被外邦的人.
或者軍王.
嘲笑和攻擊.
二十五到二十九節.
他們本來.
就是與神納了約.
一個有.
神的約的約民.
但是他們都.
違反了神的約.
他們表面上說我們受國禮.
我們因為我們是神的.
兒子.
即是律法之子.
可能又拆開一點.
律法之子.
是大概猶太的男孩.
十二至十三歲.
他們.
那個不是國禮.
因為國禮是第八天行的.
即是剛剛出世.
八天就要受國禮.
男孩.

$^{361}$但是律法之子.
有可能是.
十三歲.
有些人可能是十二歲.
他們正式行.
律法之禮.
是一個成人禮.
由他們.
行了這個禮之後.
即是說.
你自己是一個成年人.
你不需要再由爸爸.
去保護你.
即是說你應該要明白.
神與你立約.
的關係.
所以如果你犯罪.
就是你自己個人去承擔.
神對你的罪.
那個追究.
但是十二歲以前.
所做錯的事.
家庭有責任去幫他.
來到去.
承受.
所以這個叫律法之子.
大家記得在路加福音.
的第二章.
耶穌就是這樣.
去耶路撒冷.
聽拉比的教導.
一邊聽一邊問.
就是一個這樣的儀式.
所以你可以想像到.
猶太人.
他們有他們自己的一套.
法規.
宗教的條例.
但是保羅.
爭論的是.

$^{401}$他們兩樣東西都做不到.
所以他們雖然有律法.
但是都不可以.
在律法底下.
證明他們是.
得到神的義.
他們都需要神的救恩.
所以一定要在.
神的律法以外.
的恩典的約.
來將他們拯救.
這些是上個星期我們講了.
我就幫大家來做一點.
的重溫.
所以第29節.
你可以說.
差不多我們所有人.
都需要在.
這樣的屬靈的.
氛圍底下.
去看自己和神的關係.
一定要心靈裡面.
來作真的猶太人.
或者說.
作真的基督徒.
或者說作真的錦宣人.
紅恩人.
嘉恩人.
我們才能夠是真正神的兒女.
所以你不要拿著證書.
這個就是.
我是屬於神的人.
然後你在證書.
掛在家裡.
最當眼的地方.
你邊年邊月洗禮.
但是你的生活你的為人.
你的靈命和你的洗禮.
是完全不符合.
你就落在這個.

$^{441}$這樣的審判.
甚至是奏詐底下.
唯有裡面作.
才是真猶太人.
真覺禮也是心裡的.
在乎靈不在乎疑問.
不是說疑問不重要.
如果你沒有.
心靈而只有外表的疑問.
那就是極之可惜和虛假.
所以如果你只是在乎外表.
在錦繡堂洗禮.
真是不得了.
在湖邊.
天鵝相伴.
然後拍照.
漂亮絕不得了.
外在疑問.
很可惜.
下星期拍照也沒興趣看.
你完全沒有.
靈命的支撐.
你外面的感動.
其實就不得了.
真實了.
所以最終會變成.
記憶的淡忘.
基督徒的身份.
不能夠一直持久.
你越來越.
可以在心靈裡面.
與神的救恩帶來的相交.
帶來的相愛.
帶來的和好.
這是第五章說的.
你就不會享受到,經驗到.
所以洗禮這個儀式.
不是不重要.
如果只是你強調這個儀式.
而沒有內心的.

$^{481}$靈命.
裡面的.
與神的相交.
外面的儀式.
不單止沒有幫助你.
反而嚴重地說.
可以害了你.
令你與教會越來越陌路.
這是很可惜的.
洗禮是一個真正.
進入與神相交的開始.
不是畢業禮.
洗禮不是畢業禮.
洗禮是進入與神.
更加緊密的.
相交的開始.
所以這人的稱讚.
不是從人而來.
是要從神而來.
神稱許我們內心的靈命.
所以我們要看一兩段.
關於割禮的經文.
我們先看.
歌羅西書2章.
11到第13節.
保羅如何引述割禮.
當然是一種.
責備口吻.
歌羅西書.
如果我這本.
舊的和合本.
是新約的280頁.
差不多.
大家去掀一掀.
如果是電話.
相對容易.
如果是大本聖經.
我建議大家買一本大聖經.
在家做研究.
看讀經.

$^{521}$大隻些.
方便些.
掀聖經又容易些.
2章.
歌羅西書.
由第6節讀起.
段落比較詳盡些.
你們既然接受了主基督耶穌.
就當遵祂而行.
在祂裡面.
生根建造.
信心堅固.
正如你們所領受的教訓.
第8節.
你們要謹慎.
恐怕.
有人用他的理學.
和虛空的妄言.
不照著基督.
乃照人間的遺傳.
和世上的小學.
就把你們擄去.
因為.
神本性一切的豐盛.
都有形有體的.
居住在基督裡面.
即是說叫我們千萬不要離開主耶穌.
去聽一些虛無飄渺的學說.
一些似是而非的.
教導.
因為神已經將他一切的豐盛.
藉著耶穌.
給了我們全人類.
有形有體的.
居住在基督裡面.
第10節.
你們在他裡面.
在耶穌裡面.
也得了豐盛.
他是各樣執政掌權者的元首.

$^{561}$你們在他裡面.
也受了不是人類的福.
而是人類的福.
也受了不是人手所行的國禮.
乃是基督使你們脫去肉體情慾的國禮.
你們既受洗.
與他一同埋葬.
也就在此.
與他一同復活.
都因信那叫他從死裡復活.
神的公用.
你們從前在過犯和未受國禮的肉體中死了.
神赦免了你們一切的過犯.
便叫你們與基督一同活過來.
這裡就告訴我們.
保羅就藉著國禮.
這樣的舊約儀式禮儀.
因為女性就不需要行國禮.
那是不是只有弟兄才需要這樣做呢?.
這裡的意思當然不是這樣.
既然我們已經在主的救恩裡面得到豐盛.
主已經藉著他的聖靈去重生了我們.
進入我們的生命.
洗乾淨我們的罪.
我們就已經好像做了一個心靈的手術.
就是將我們整個人.
透過洗禮.
我們已經好像浸入水裡.
與主耶穌同埋葬.
然後我們又再上水.
與主耶穌同復活.
所以就有這個復活的大能.
在我們每一個弟兄姊妹的心靈裡面.
所以我們要不斷經驗.
耶穌是怎樣死的.
為什麼他會死.
我們要不斷認同.
主的死是和我有關.
因為是擔當我的罪.
所以我既然現在耶穌復活.

$^{601}$即是我的過去.
神已經看為是.
已經放在耶穌的生命裡面.
來帶進審判和帶進刑罰.
所以主耶穌已經成為了我的代表.
所以這個就是第十一節所說的.
基督使你們脫去肉體情慾的割禮.
所以不是物質的割禮.
所以男女都需要的就是.
聖靈在我們生命裡面的進入.
那個掌管.
那個感動.
聖靈就在我們每一個弟兄姊妹的內心.
所以一信主.
第一件事就應該感受到自己是一個罪孽深重的人.
要向神認罪.
你的罪行就會像告白那樣一筆勾銷.
因為主的寶血已經洗乾淨.
和他已經擔當了我們的罪.
所以這個就叫做肉體情慾的割禮.
已經割了.
所以我們要認同這個十字架的死.
其實是你和我應該釘上去的.
但是就算我們釘上去也沒有用.
因為我們不能夠再解決我們自己一生要行善的問題.
所以基督的死就是關乎我了.
基督的復活也是和我有份.
所以耶穌的復活現在在天上.
保羅在《歌羅西書》也說.
其實現在我們和主耶穌一同在天上.
這個當然是說原來我們有兩個的世界.
一個是屬靈的世界.
所以我們可以隨時向耶穌禱告.
其實就進到天上的致勝所.
我們同時間實實在在我們是活在這個地上的.
我們實實在在是要面對我們的生老病死.
萬事就是好事壞事.
我們想的事我們期望發生的事.
和我們意料之外的事.
都是會淋到我們的人生的.

$^{641}$但是這裡告訴我們.
基督已經為我們做了這個國禮.
所以記得我們從前在未信主之前.
我們在過犯和未受國禮的肉體中死了.
所以我們必須要認同.
我們要認同我們是一個罪人.
是恩典將我們拯救.
我們現在在耶穌的裡面.
我們舊有的過去已經解決了.
所以我們是一個新造的人.
這個就是保羅在《歌羅西書》所說.
意味著當時可能有一些學說.
有一些似是而非的哲學.
可能夾雜一些神秘主義.
就說你信基督不夠的.
你們要追求有一種屬靈的.
深奧的神秘的經歷.
你們才會成為高級.
成為一種靈命的突破.
其實不需要這樣的.
每一個基督徒.
必須是對主的救恩有全盤的了解.
掌握和相信.
這樣信靠.
你就會經歷這個國禮.
記得是身體情慾的國禮.
所以這是一個頗重要的屬靈的概念.
今天我們在這個聖靈的國禮裡.
我們就是新的屬天的以色列人.
我們就是屬靈的亞伯拉罕的子孫.
加勒太書三章有這樣說.
三章的十九節有這樣提到.
所以我們都是亞伯拉罕的子孫.
不過不是肉身的.
因為耶穌我們透過聖靈.
我們就是屬天的亞伯拉罕的子孫.
也就是說都和神立了這個約.
透過聖餐我們就有這個約和神立了.
所以我們是永遠屬於神的.
所以千萬不要被那些錯謀的異端.

$^{681}$來勾引離開神的兒子耶穌基督.
最近都有弟兄跟我說.
現在在洪水橋都幾多那些.
又不是摩門教.
摩門教固然有.
就是守望台每個星期六日都站在天橋下面.
不是天橋 隧道那裡.
來派他們的傳單.
所以我們的弟兄姊妹.
其實如果稍為聽了就很容易受影響.
因為我們自己本身對真理不是很全面的掌握.
所以我都盡快會開一個異端的分辨.
給受洗的弟兄姊妹去上.
希望他們能夠對這些真理有一個更加全面的掌握.
和這些異端到底他們有什麼錯謀.
我們有稍微少少的了解就會好些.
我都會盡快開.
另外肥勒比書三章 翻前一點.
前一卷書就是三章的肥勒比書.
三章由第一節開始.
弟兄們 我還有話說.
你們要靠主喜樂.
我把這話再寫給你們.
於我並不為難.
於你們卻是妥當.
又是叫他們留意一些似是而非的錯誤道理.
應當防備犬類.
防備作惡的.
防備妄自行割的.
妄自行割即是說.
有些可能極之熱心的.
又是一知半解的猶太人的信徒.
勉強那些外邦人.
你未行割禮.
你不夠聖潔.
你沒有我們猶太人那麼.
跟足神的律法的要求.
我們也是信耶穌的.
不過你不行割禮.
你就不夠我們那麼好.

$^{721}$類似這樣.
保羅看這些就是一知半解的半桶水的.
妄自行割.
又或者當時如果帝國政治.
對猶太人有一個favor.
你如果想得到這樣的favor.
你就要不得我都行割禮.
我這麼大個都行.
證明我真的有這樣的熱心.
那可能會有這樣的妄自行割.
所以一定要跟當時的國情.
社會的背景來相掛鉤.
所以保羅可能很意識到.
教會都可能會受這些輿論的影響.
所以保羅說不是你外表要扮一個猶太人.
而是你要跟舊約的形式去生活.
可能很快地說一個.
大家聽過安息日會沒有?.
他們的主張是很好的.
有一些激烈的事我們不做不吃的.
例如我們不喝咖啡.
這些沒所謂的.
我也不喝因為我會敏感.
但是他們說的就是那些食物帶來神喜悅.
那是過火了.
接著他們認為崇拜是應該跟猶太人.
安息日星期五的日落計起來去做崇拜.
做禮拜.
所以安息日會是必須星期六崇拜.
我們感受也有的.
感受是不是安息日會呢?.
當然不是.
因為我們有些人禮拜要開工.
所以都夠一定的人數.
有這樣的需要.
所以我們就開星期六晚的崇拜.
是分流.
不是因為要守安息日.
所以是兩樣東西不同的.
所以不同的思潮.

$^{761}$那些的思想.
那些的教導.
我們要很小心.
所以不是我們每一樣東西都要這樣跟.
就叫做合神的心意.
保羅這裡的說法就是說.
要防備這些人.
因為第三節.
希望大家都背熟它.
都是重要的經文.
因為真受割禮乃是我們者以神的靈敬拜.
在基督耶穌裡面誇口.
不靠著肉體誇口的人.
接著保羅就說.
如果要誇口.
第四節.
其實我也可以誇口.
接著他就立即說.
我就是第八日受割禮.
所以很明顯說.
這裡第三節說的.
我們是以神的靈敬拜的真受割禮者.
是以神的靈去敬拜.
哪些人可以敬拜神呢?.
就是信主.
認識神.
有聖靈內在我們心裡面的人.
才可以敬拜神.
所以我們一定要記得.
保羅說的.
其實他自己剛剛就說.
我第八天受割禮是根正苗紅的猶太人.
但是他完全不誇耀這些東西.
他現在向外邦人傳福音.
完全不會說舊約的背景.
他自己有什麼優越.
他說他自己是迫害神的教會.
他比其他任何一個使徒還要低微.
還要該死.
但是他現在告訴我們.

$^{801}$每個人都在神面前.
是徹徹底底的平等的.
只需要你去到神的面前.
是用心靈去敬拜神.
所以某程度教會.
信仰是帶給每一個民族.
每一個時代.
都真真正正的在神裡面的平等.
回到教會大家真的是弟兄姊妹.
大家真的去敬拜神.
所以保羅這裡說了一個很重要的秘訣.
真受割禮是以神的靈敬拜.
不誇耀肉體的人.
才得到那個敬拜.
所以你不要回來.
每個星期回來.
一回來就說.
我真是很蒙福.
每個星期回到這麼漂亮的教堂.
我坐下來之後.
我回到去就很有福氣.
不是教堂令你很有福氣.
是神的恩典令我們每一個人得到真正的福氣.
所以我們不是去敬拜教堂.
我們不是來到教堂.
向教堂來敬拜.
教堂你真是偉大.
我敬拜你.
接著我就去會所喝茶.
不是.
我們是向神敬拜.
來到神的面前.
所以你可以想像得到.
在舊約.
神給猶太人最大的羞恥是什麼.
就是聖殿被拆毀.
被擄到巴比倫.
那些聖殿裡面所有的聖潔的神聖的獻祭的器皿.
全部被帶到巴比倫成為戰利品.
是要去到70年之後.

$^{841}$所羅巴巴得到神的開放出路.
感動波斯王古烈.
將那些器皿帶上來重建耶路撒冷.
所以你可以想像得到.
如果我們搞錯了敬拜對象.
拜個石頭都不拜神.
那就很可惜了.
所以第三節是一個很重要的屬靈秘訣.
讀多一次.
因為真受割禮的乃是我們.
這以神的靈敬拜.
在基督耶穌裡誇口.
不靠著肉體來誇口.
讚美.
服侍耶穌.
敬拜神.
這就是神的靈的敬拜.
所以大家每個主日.
真的要有一個很認真的心靈.
以靈以誠敬拜神.
沒有人提醒你.
提醒你可能是一個罪人.
得失了大家.
叫大家安靜一下.
有時候好像很不好意思.
其實每個人真的要安靜自己來敬拜神.
來準備自己的心靈.
向神發出真誠的敬拜.
當然你可以帶著很多困難.
憂愁.
有很多需要.
是沒有問題的.
正正就是這樣.
才是一個真人.
但是帶到來的時候.
就跟神說.
這些事情我真的解決不了.
我真的很煩惱.
但是我現在要敬拜你.
我放下你幫我來解決.

$^{881}$安靜我的心靈.
我完善崇拜.
我才去想.
現在不想了.
我們要這樣學習.
要這樣操練.
否則我們就不是真正去敬拜神了.
第十六章的《四道行傳》.
我們讀完這一段就算了.
第十六章的第三節.
這就是保羅第二次的宣教旅程.
保羅來到特比.
猶道路斯德.
在那裡.
這是土耳其的內陸地方.
在那裡有一個門徒.
名叫提摩泰.
是信主猶太婦人的兒子.
她的爸爸是希利尼人.
特別強調.
她的爸爸是外邦人.
媽媽是猶太婦女.
母親是羅爾.
對不起是友尼基.
是她媽媽.
她的外祖母才是羅爾.
這個家庭特別強調.
路加醫生特別強調提摩泰.
就是說保羅見到這個年輕人.
一定是有些東西不同的.
我們再讀下去.
第二節.
路斯德和以哥廉的弟兄都稱讚他.
突然間多了一個以哥廉.
以哥廉是那個省.
以哥廉是新界西的元朗和屯門洪水橋都稱讚的弟兄.
大概是這個意思.
以哥廉的弟兄都稱讚提摩泰.
第三節.
保羅要帶他同去宣教.

$^{921}$只因那些地方的猶太人都知道他父親是希利尼人.
是什麼意思呢.
那裡的一帶 以哥廉 路斯德 特比.
一帶的地方猶太會堂的人.
都認識提摩泰的家庭.
認識提摩泰的父親.
沒有名字的.
是希利尼人.
所以那是怎樣呢.
提摩泰有一半的血統是猶太人.
現在保羅帶他去傳道.
見那些會堂的世叔伯.
所以猶太人還有這種的一種的戒條.
還有一種這樣的迷思.
就是說這個人自己都不守律法.
帶來可能會有論斷.
對提摩泰的傳道又會打了一個折扣.
所以這個無關是救恩.
而是關乎到傳統和一些對長輩的.
他們未能夠一下子消化的一種緩和.
這樣是有一種禮儀.
是有一種識大體.
是有一種保羅識人情世故的一種智慧.
傳道就會較為方便.
那是怎樣呢.
就不是勉強提摩泰行國禮.
而是幫這個年青人未受國禮的.
就來受國禮.
這意味著這個手術都需要一個星期左右.
我不知道當時做手術的祭司.
行國禮的祭司.
醫護的常識.
那些包紮.
那些洗傷口的藥水.
是不是夠乾淨.
我相信都應該是很專業的.
所以年青的提摩泰去傳道之前.
就要行國禮.
一個年青人要割包皮其實都幾痛的.
保羅就幫他做了.

$^{961}$記住這裡是無關得救.
所以保羅是講了背景.
就是令到提摩泰將來傳道.
不會被人拖後腿.
譬如你說你自己是猶太的信徒.
自己原來都沒有行國禮.
自己都類似其身不正.
其實講大了.
其實不是.
因為已經沒有需要了.
外邦人是沒有需要行國禮.
但這個就是一半猶太人.
一半希臘裔.
為了傳道的方便.
就幫他行國禮.
這個就是保羅的智慧了.
所以原來國禮猶太人做當然固然好.
但原來無關在神面前的救恩和靈命.
所以記得今天我們不是去哪間教會.
大教會,小教會.
令到你的靈命特別好.
而是你在那個教會裡面.
得到被栽培.
全面掌握聖經真理.
和有參與事奉的那份熱心.
那份擺上.
是這樣去蒙福.
所以不關你在哪間大教會.
每個星期都出名的神學院院長都沒有用的.
如果你自己都不去追求.
自己都不好好去成長.
不將自己的生命擺上給神.
去讓神裝備.
給神去用.
這樣其實都幫不了多少.
所以.
在教會越久.
越難去更新.
你好像說什麼道我都聽過.
我更加.

$^{1001}$你未講我都懂得怎樣講給你聽.
沒錯.
不反對的.
的確是.
但是那個道有沒有進入到我們的心靈呢.
原來每間教會都一樣.
有些人回來只是做事.
他很熱心做事.
他沒有想過聽道.
沒有想過安靜坐在那裡.
有些更加諷刺.
他做完事就走了.
回來開了個檔之後.
他就走了.
他不崇拜了.
很有趣.
這樣就是一個問題.
應該是要有牧者去勸導.
去循循善誘.
我們不是要請廉價勞工.
也不需要我們是要回來敬拜神.
如果事奉令到你.
回來做完事就走了.
那真的是一個很大的問題.
到底是教會的問題.
去不用錢去用弟兄姊妹.
還是弟兄姊妹自己不知不覺有一個錯覺.
有一個誤解.
聚會的目的呢.
那你就可以想像到.
如果我們經年累月是這樣.
自己的靈命不成長.
你就不要怪任何一個人.
必須是要自己面對神.
在我們生命裡面不斷的重整.
不斷的去更新.
這個就是割禮的部分.
所以和二章的二十九節.
是完全一樣的.
真受割禮.

$^{1041}$是我們這些以神的靈敬拜.
所以外面肉身的割禮.
保羅說不是真割禮.
這樣就挑起很多猶太人的反感.
這一句說話已經帶來.
很噁心.
很衝擊猶太人的律法主義的觀念.
你叫人不要受割禮.
你說割禮沒用.
你說我們猶太人一文不值.
類似這樣.
其實保羅的意思.
不是否定肉身應該行的禮儀.
但是如果只有外在的禮儀.
沒有了裡面的真誠.
和靈明與神的相交關係.
那就是一個真正的問題.
下一頁.
割禮這裡有沒有問題?.
如果沒有我們就繼續下去.
猶太人的身份.
我們看完第二章之後.
接著就會有一個.
其實大家應該浮現的問題.
就是那是不是猶太人一無是處呢?.
神的律法和割禮和這些藥.
是不是完全沒有任何的功效呢?.
保羅又在假想有這樣的針對他.
告訴人什麼都不用做了.
有耶穌就行了.
恩典就行了.
是不是這樣就令到猶太人有這種困惑.
或者讓一些刻意挑剔保羅的言論.
講解福音的人.
會有機可乘.
好像要挑起猶太人和外邦人.
因為神的藥的矛盾.
所以保羅在第三章的第一到第五節.
他就開始講了一些關於猶太人.
其實不是一無是處的.

$^{1081}$第三章我們讀一讀經文.
第三章的第一節.
這樣說來猶太人有什麼長處.
割禮有什麼益處呢?.
第二節.
凡事大有好處.
第一是神的聖言交託他們.
即便有不信的.
這又何妨呢?.
難道他們的不信就廢掉神的信嗎?.
斷乎不能.
不如說神是真實的.
人都是虛放的.
如經上所記.
再一次引述舊約.
這裡就是詩篇的五十一篇第四節.
是大衛向神認罪的詩篇.
你責備人的時候顯為公義.
被人議論的時候可以得勝.
保羅就將一些當時可能他想到.
可能真是有其人挑剔保羅的傳道的.
這些人在挑骨頭.
所以保羅就假設.
以為他這樣說.
那些人就會說.
來論斷批評保羅.
說到猶太人一文不值.
說到神的約好像完全沒用.
保羅就告訴我們.
其實他不是這個意思.
一方面.
神的猶太人.
神的子民.
和割禮的約.
有什麼好處.
凡事大有好處.
第一就是.
神的聖言交託他們.
這真是很奇妙.
很特別.

$^{1121}$猶太人.
他們真是一直傳承著.
抄寫神的摩西律法書.
到今天.
所以在舊約神容許.
未有印刷術.
未有photocopy之前.
神容許這個民族.
抄寫神的聖言.
舊約的啟示.
能夠原本本流傳到我們今天.
我們手頭上的舊約聖經.
是十三世紀的馬索拉抄本.
即較為權威的和合本.
但是這個馬索拉抄本.
已經是極之.
來到去符合.
舊約的當時.
在主前大概.
兩世紀至一世紀.
那些抄本.
那個準確度.
原因就是死海古卷的出土.
讓我們發現.
原來我們今天最權威的.
最近的.
流傳今天最多的抄本.
這個馬索拉的抄本.
是和一世紀的死海古卷.
差不多90\%的完全一樣.
你可以想像到那個抄本的嚴謹度.
所以神的聖言.
交託給猶太人.
猶太人就忠於他們這個民族的身份.
捍衛著神的書卷.
來到一路流傳.
讓後世的人都可以中毒.
這個就是.
記住是一路每個時代.
每個時代來到過去.

$^{1161}$因為神容許.
沒有那些高科技發明.
不會說主前的2000年.
就已經有互聯網.
就已經有影印.
就已經有電子版.
就已經有相機高清.
等等.
是沒有的.
什麼都沒有的.
只有一支筆.
還要是那些蒲草的籽.
是要這樣流通.
還要是那些羊皮.
這樣來保存.
這樣流通.
流傳.
那你可以想像到.
神的聖言交託給猶太人.
他們去抄寫.
所以今天新舊約傳書.
又有電子版.
又有這麼多不同版本的APP.
如果我們不讀神的聖言.
其實就很可惜.
當然神都會感動和興起.
有很多聖經學者.
有些宗派很嚴謹讀經.
所以一個教會的復興.
絕對不可能沒有人愛慕神的道.
所以我們錦宣.
現在我們紅恩都有.
我們每年都很想有讀經日營.
那個初心就是這樣.
都是張牧師一直這樣孕育.
每年都要有讀經日營.
就是很想喚起弟兄姊妹的心.
我們專注一起去聆聽一些.
準備優質的經卷的研讀.
這樣來提起我們以後自己去研讀的方法.

$^{1201}$那個興趣.
那個投入等等.
所以神的聖言交託給猶太人.
神的聖言有沒有交託給我們錦宣家.
紅恩家呢?.
我們真的要問的.
我們必須對神的道.
要有一個掌握,認識.
以致我們能夠對神的真理.
有一個了解.
所以其實我自己是很賞心.
如果我們屬靈的信徒領袖.
根本不讀經.
根本不知道什麼是聖經.
根本不知道什麼是聖經的真理.
只是一套專業,智慧,大學畢業.
很有學問.
沒錯,不否定的.
但是更重要的是.
監督既然是神的管家.
必須要有健全的,真道的裝備.
來教導其他人認識神.
所以如果我們真的搞錯了.
選立領袖的方法.
你不要怪教會一直軟弱.
因為我們什麼叫做屬靈.
我們自己都不懂.
那你如何告訴別人這是教會的成長呢?.
量的成長是從來不否定的.
但是如果量的成長不是伴隨著質量.
至少一班領袖的質量.
是要有神的聖言.
這個差不多是肯定的.
所以這個大家就要彼此守望祈禱.
所以保羅受到挑剔.
保羅一定會知道有些人會對他.
突然轉變他自己的跑道身份.
他成為外邦人的使徒.
去講一個好像這麼廉價的信耶穌福音.
一定會受到很多的挑戰.

$^{1241}$逼迫和背後的閒言閒語.
在骨頭裡面挑釁.
所以你可以想像到.
保羅預計到會有這樣的情況.
是不是猶太人一無是處呢?.
但是保羅就說不是啊.
神的聖言交託給他們.
就是說交託給你.
你們只是守著那個羊皮卷.
你們自己不去中毒.
不去遵守.
不去反思和神的關係.
只是有一套徒有外表的猶太人的外殼.
沒有了真正的因為神的聖言而改變你的生命的轉化.
這個問題不關神的事.
更加不關神的聖言的事.
是猶太人自己本身.
所以他就立即說.
即便有不信的又有何妨呢?.
難道他們的不信就廢掉了神的信實.
神的真實嗎?.
神的不變嗎?.
神的信實就是那個不變.
斷乎不能.
不如說神是那位永恆不變的神.
即是說他的應許.
他的約是不變的.
猶太人變.
猶太人當歸他自己的身份.
但是神的信實.
應許亞伯拉罕的子孫永遠可以蒙福.
是不變的.
神自己是不變的.
所以保羅的一個爭論.
是很重要的.
告訴我們.
就算我們見到有不同的人.
我們過去很尊敬他.
而發現他有很多的限制.
很多的瑕疵.

$^{1281}$我們採取包容和接納.
因為大家都是人.
都是罪人.
我們不需要高舉人過高.
但千萬不要反過來.
踩一個人踩到過低.
大家都是人.
所以這個世界不應該有聖品人.
當然我們不拜天主教的情況.
我們也不需要去論斷人家教宗怎樣怎樣.
但是我們自己記得.
人是虛荒.
但是神是信實的.
這就是第三節所說.
即便有不信又何妨呢.
難道他們的不信就廢掉了神的信實嗎.
難道你看到哪個人的軟弱.
你就否定你的信仰嗎.
不需要這樣做的.
我們應該更加包容.
禱告.
我們應該為那個人.
為那個教會祈禱.
而不是加入論斷.
網上去公審.
不需要這樣的.
大家都是神的兒女.
我們當然如果有機會.
可以好言相勸.
他不聽的.
我們盡了我們的責任.
如果認識的.
如果不認識的.
就更加不應該這樣往下判斷.
但是記得.
保羅這裡告訴我們.
就算人的不信.
都不應該廢掉神的信實.
所以我們不要因為人而離開教會.
不要因為我們見到人.

$^{1321}$好像我們自己特別熟悉一點.
這個人就是這樣.
沒得救了.
不需要這樣的.
我們應該帶著一個恐懼.
但是我們要記得.
神是信實的.
不會因為人的虛荒.
人的不信.
就廢掉神的真實.
就是這個意思.
接著就說.
人都是虛荒的.
不如說.
人都是虛荒的.
神是真實的.
那就是.
反過來說.
所以.
你不要看人.
人根本就是虛荒的.
人根本就是不能夠做到100分的.
是永遠都做不到的.
所以永遠都是在恩典裡面.
去倚靠神的.
所以.
這個就是告訴我們.
永遠要轉眼仰望耶穌.
我很喜歡一節聖經的.
對我自己有很大的幫助.
希伯來書的十三章第八節.
大家一定要懂得背.
如果你是參與事奉的.
就更加要銘記在心靈裡面.
我們事奉的人.
是有很多的口訣.
好像乾坤大羅兒的.
波斯文的口訣.
張無忌練完之後.
就可以功力猛進.

$^{1361}$聖經都有很多對於事奉者的重要的口訣.
其實聖經的金句.
十三章.
很重要的.
第七節.
十三章的第七節.
作者在這裡.
我相信都可能是保羅.
是我相信的.
因為很多人都不認為是保羅寫的.
當然作者沒有說他自己是誰.
所以就變成2000年的原案.
回到天家就要問.
誰寫希伯來書的.
請你走出來.
讓我們開估.
到底是巴拿巴還是保羅.
因為都有人說是巴拿巴寫的.
第七節.
從前引導你們.
傳神之道給你們的人.
即是傳道者.
從前即是他已經死了.
你們要賞念.
賞念他們.
效法他們的信心.
留心看他們為人的結局.
賞念.
效法.
留心看.
即是說很重點的去注意這些人的人生為人.
接著第八節.
十三章的第八節.
就是耶穌基督昨日今日.
一直到永遠都是一樣.
你看人.
他帶領過你.
你思念他.
你留心看他.
耶穌基督才是昨日今日.

$^{1401}$一直到永遠都不改變.
即是說人.
昨日你見他.
嘩 很正.
很勁.
很好.
很尊敬.
今日再看他.
不外如是.
和他一起做同工.
都差不多.
甚至某些事比他好.
有一次.
我這樣說好像不是很好.
不過都當是笑話算了.
我們在神學院的時候.
有個神學院的老師就說.
嘩 去完洗手間.
看著一個比我年紀大的教授.
去完洗手間不洗手.
即是當作是開玩笑.
當然他完全沒有開名.
即是有時.
這些當是茶餘飯後.
這樣說笑.
但其實我們如果認真去看人.
我們有時會有過去的印象.
和現在的大家的遺憤.
和值得我們再去無全清的.
總是會有每個人的限制和瑕疵.
但是十三章第八節就很重要.
是我們一個很重要的事奉的基礎.
無論人怎樣.
曾經在我們生命裡面.
留下很大的我們對他的敬重.
但是人始終是人.
或者你可能有一天已經超越了他.
但是不要緊的.
耶穌基督昨日,今日.
一直到永遠都是一樣.

$^{1441}$即是說永遠我們要注目仰望耶穌.
就會有力量繼續這樣孤單前行.
我們應該要擔負起的.
神安排給我們的責任.
好像那些電影,電視.
那些臥底警察.
經常都要忍辱負重.
我們年輕的時候有一套叫《獵鷹》.
劉德華就給太太,給全家人都看他.
因為他已經走進黑社會.
但是他又口難言.
因為他是臥底.
有時候這種心情.
事奉者就一定會有.
真的不是拍戲.
有時候你做一些決定.
真是千呼所指.
但是不要緊的.
你對準神.
你知道這件事一定是對整個群體有益的.
你就一定要這樣去做.
不做才容易.
做才難.
這個時候你要記得.
耶穌基督昨日,今日.
一直到永遠都是一樣.
所以千萬不要做一個犬儒.
對人怎樣都看不起.
個個都看不起.
你最厲害的.
完全不要這樣.
反過來.
你覺得有些曾經幫過你.
現在他都可能被你超越了.
你要記得.
他都是人.
他都會過去.
等於我們一樣會過去.
所以耶穌基督才真正是不改變.
值得我們投放我們整個生命.

$^{1481}$去服侍他.
去被他去使用.
所以送給大家這一節.
真是很重要的口訣.
是很重要.
尤其是應該在你自己很失落的時候.
你就背下這節金句.
讓聖靈提醒你.
你看誰?.
你看耶穌.
你看他.
他有他的跑道.
他有他的工場.
神都收納他.
你是他誰?.
你是他的主人?.
神都收納了他.
那你為什麼還不看你自己的工場.
你自己的位份.
你自己的神安排你的人生呢?.
所以這個是很重要的.
所以神的聖言交託給猶太人.
他們忠於他們自己的聖言.
不過只不過是他們虛方做不到.
引述的是《詩篇》51篇第四節.
不詳細讀.
因為大衛很真誠的懺悔.
拿丹先知醫來到去責備他之後.
他就去到神面前向神認罪.
所以他就很清楚.
不是我們的義.
是神的義光照在我的心靈裡面.
所以第四節就是《詩篇》.
我每一次出現了舊約.
我已經做了一個table.
所以遲些就來不及就訂給大家.
那你就可以知道.
保羅在第三章已經引了差不多九節的聖經.
舊約的《詩篇》和《爾賽亞書》.
只是現在這一段都有很多經文保羅引述的.

$^{1521}$第五節.
我且照著人的常話說.
我們的不義若顯出神的義來.
我們可以怎麼說呢.
即是說有功勞.
原來我們犯罪.
神要給我們功勞.
因為我們的犯罪.
我們的不義.
可以顯出神的義.
即是說如果我們違規駕車.
運輸署也要給我們一些獎金.
因為令到運輸署的公義可以彰顯.
這些簡直是荒謬的.
是不可能的事來的.
你違規就是違規.
你的違規可以令到神的義彰顯出來.
這個是完全說不通的一個謬論來的.
保羅就這樣提出.
可能真的有這樣的人.
有這樣荒謬的想法.
所以神是不應該降奴.
因為我們的不義是顯出神的義.
如果沒有我們的不義.
神怎麼可以顯出他的義呢.
真是這麼荒謬的事.
所以第六節就說.
端乎不是.
如果這樣的時候.
神怎能夠審判世界呢.
如果神的真實.
因為我的虛方.
越發顯出他的榮耀.
為什麼我還受審判.
好像罪人呢.
就是說.
為什麼還要審判我.
我令到你的義彰顯.
其實我是有功勞的.
等於猶大出賣耶穌.

$^{1561}$那耶穌才會被出賣.
他才能夠成為救主.
其實兩件事是沒有關係的.
沒錯.
猶大真是被預定要出賣耶穌.
但不等於這樣他就有功勞.
他的確是出於他自己的貪心.
懼怕.
想和耶穌劃清界線.
想在最後的耶穌的死亡.
他都要拿一些油水.
從來沒有想過.
要為神的義.
來付上代價.
他一心想著犯罪.
因為一心已經認定.
耶穌對他來說.
是沒有任何的利用價值.
所以有時候神都容許我們.
在人生萬事都效力.
有時候讓我們都有一些困難.
考驗我們的信心.
每次都順利.
每天都稱心如意.
有求必應.
這樣的信仰遲早有一次.
很嚴重的崩潰.
所以有很多人.
有些人生的低谷.
第一時間他要否認的就是神.
其實在低谷.
在幽暗.
更加要在黑暗裡面.
抓住神.
不要離棄.
這才是真正的信仰.
所以保羅說斷乎不是.
如果這樣的時候.
神怎能夠審判世界呢?.
神要和一些人去討價還價.

$^{1601}$好像要買通一些人.
這是不可能的事.
神本身就是真實.
人自己去犯罪.
人自己去背棄神.
人自己不去根據神的意義去生活.
就是必須要面對神的懲罰和審判.
所以這裡告訴我們.
第七節.
若神的真實因我的虛方月法顯出他的榮耀.
為甚麼我還受審判.
很像罪人.
即是將一件事扭曲了.
即是本末倒置了.
將罪還要怪責給神.
神是你令我犯罪.
神是我犯罪.
應該你給我一些獎勵.
因為令到神的義可以彰顯出來.
這簡直是一個難以想像的廢話.
保羅也想到可能有這麼極端的人.
會這樣想事.
所以一切的罪.
不是在於知識層面.
是心靈裡面對神的背棄.
由一章十八節到三章八節.
都已經很清楚說.
是心靈與神雙離開而帶來的罪行.
因為人有自己的自由意志.
選擇的是犯罪.
不是選擇敬畏 跟隨和愛神.
第八節.
為甚麼不說我們可以作惡而成善.
即是我做錯事.
神應該要獎勵我.
因為我作惡可以帶出善.
這就是絕對不可能的事.
因為神根本有權制定法規.
道德律和每個人要面對神的要求.
因為他是做物主.

$^{1641}$是他定的規條.
所以後面就說.
這是毀破我們的人說我們有這話.
所以這裡可能暗示.
有些人其實刻意去挑剔保羅.
說一個這麼漣漪的福音.
信就行了.
不明白過終背後的.
其實是需要基督耶穌的犧牲.
痛苦的寫明.
被神離棄的那種痛苦.
那種絕望.
以為信就行了.
帶來這種好像廉價的福音的論述.
所以保羅說.
是這些人試圖去毀破我們有這樣的說話.
這些人其實更加應該被定罪.
因為不僅不謙卑.
反省自己在神面前的不義.
反而挑剔保羅這群傳道者.
說的恩,信,清,義.
帶每一個人進入救恩與神和好的.
這麼重要的福音傳遞.
所以如果我們在這樣的人生.
在教會專挑剔人.
而自己不思長進.
在神的救恩裡面.
要很小心.
不是說笑的.
因為你自己延遲了自己與神的相交和好.
那個生命關係.
你就不斷用挑剔.
作為你自己自然的材料.
因為你在這樣的狀況下.
永遠覺得自己是好.
因為你不好就不會挑剔人.
所以會帶來一個很嚴重的.
沒有了自知反思的能力的狀態.
所以千萬不要有這種喜歡在教會.
眼睛不斷留意哪個人的不好.

$^{1681}$這是很嚴重的.
第三章一到八節已經說了.
保羅想到有人會挑剔保羅.
說到猶太人好像一無是處.
好像要否定猶太這個民族.
好像保羅自己就厲害了.
自己拉比現在轉了跑道.
成為外邦人的使徒.
特派去傳神的福音.
所以很可能對他的福音.
有時帶著一種偏見誤解.
這樣來到好像炒在一起.
對保羅有系統的講解救恩.
一概的否定.
這個很容易就會出現.
所以記得我們喜歡一個人.
對他的瑕疵都要了解.
我們不喜歡一個人.
對他的優點千萬不要一筆勾銷.
人就是這樣.
喜歡一個人.
他的缺點你都很大的包容.
但你不喜歡一個人.
他有多好的長處優點.
你當看不到甚至乎在挑剔.
這樣我們永遠不會成長.
所以我們求主幫助我們.
能夠像保羅這樣.
他在自己猶太人的身份去講.
猶太人一樣需要救恩.
因為他們根本不能夠照著神的律法.
去生活滿足神的義.
所以他都將自己放在宏觀論述裡.
萬民不論什麼民族背景.
都是欠缺神的義.
這就是一章十八節到三章八節.
對於罪的論證.
三章九節到二十節.
一路去講.
這次將一個人的內心淋漓盡致地描述.

$^{1721}$這裡引用了大約五次的舊約聖經.
在這一段裡面.
三章九節開始.
答九.
我們很快看完這一段.
稍為聽一聽.
先喝一口水.
有少許乾.
希望我們今天能夠看完三章九到二十節.
有些難度的.
因為還有三個字.
第九節開始.
我們比他們強嗎?我們是哪些人呢?.
「缺不是的」.
因我們已經證明猶太人和希利利人都在罪惡之下.
即是所有的人類的意思.
就如經上所記.
沒有二人連一個也沒有.
沒有明白的沒有尋求神的.
到時偏離正路一同變為無用.
沒有行善的連一個也沒有.
引述的是詩篇的三段聖經.
詩篇的十四篇第一至三節.
詩篇的五十三篇一至三節.
另外應該有一個.
我忘記了是第十篇還是甚麼.
第十三節.
他們的喉嚨是敞開的墳墓.
他們用舌頭弄鬼扎.
這是詩篇的第五篇.
嘴唇裡有蟺蛇.
蟺蛇即是一些很毒的小蛇的毒氣.
詩篇的一百四十篇.
滿口都是咒罵苦毒.
這是第十篇的詩篇.
殺人流血.
他們的腳飛跑.
所經過的路變成殘害暴虐的事.
平安的路他們未曾知道.
這就是以創亞書的五十九章.

$^{1761}$接著十八節.
他們眼中不怕神.
詩篇的三十六篇.
接著就是第十九節.
這裡引了多少段經文.
一二三四五六七.
已經有七節的詩篇和以創亞書.
只是這一段對於罪惡罪人的內心的描述.
很特別.
由心一直去到手腳.
因為是走路.
接著是口.
接著是頭.
接著是眼.
眼中不怕神.
即是說.
全個人.
所有民族的整個人.
都是在罪惡底下.
當時羅馬帝國是一個極之強盛.
和鼓吹羅馬的盛平.
盛平即是說太平盛世.
條條大路通羅馬.
即是說.
不會有盜賊.
不會有土匪.
不會有突然間會有些人抓了你.
即是說.
是很安全.
可以四通八達到處去.
帶來那個貿易.
帶來那個大家人和人之間的交往.
所以.
那些無論是文化的交流.
貿易的交流.
當然不會有關稅.
等等.
是塑造當時的帝國的那個太平盛世.
所以.
保羅.

$^{1801}$這兩章聖經講的罪的論述.
其實.
他不是想刻意打對台.
而是將人.
無論你帝國認為.
多好的一個制度.
當時羅馬的法律制度是相當好的.
其實今天的法律.
特別是普通法.
多少是根據羅馬人他們制定的律法.
那個的背景.
慢慢孕育的.
所以羅馬人.
是對於法律的制定.
是非常之嚴謹的.
有他們一個這樣的.
這個民族的特性.
所以.
表面上.
太平盛世的帝國.
那個文化的多樣性的包容.
那個交流.
促進很多的不同的貿易.
或者文化上彼此的激勵.
各方面.
但是.
在神的救恩底下.
萬民仍然有他的限制.
有他的欠缺.
就算多好的制度.
都是不能夠與神有直接的那個相交.
和那個的復和.
必須是要透過耶穌基督.
這個救恩.
才能帶來人整體的.
可以與神和好.
那就是說意味著.
如果人一個一個的人.
一個國家一個國家.
這樣來悔改信主.

$^{1841}$就會帶來在制度上.
只會更加的完善.
只會更加的帶來平等.
只會更加的帶來人.
可以克制自己內心的慾望.
帶來那個大同的一個人和人之間.
再沒有猜忌.
再沒有任何的敵對.
那所以呢.
不是帝國.
不是任何一個制度.
去改變人類的命運.
人類其實就按照著一章十八.
去到三章的八節.
其實應該是三章二十節.
所有的世界的人.
只是步向神的審判的.
將來的永恆的刑罰.
就是火狐了.
所以保羅不是用一種敵對的心態.
去和地上的政權來挑戰.
而是他講述人性的限制.
所以沒有任何一個制度.
可以將人來承載.
改變和轉化.
必須是要透過恩順清二的這份恩典.
為二人恩順得生.
那個生命就會改變.
所以最大的動力是源於.
裡面的心靈的改變.
是有神的靈來改變的一個國禮.
所以這裡.
真是每一個基督徒都可以講這個見證.
其實每一個信耶穌洗禮的基督徒.
都要講他自己獨一無二的見證.
不懂得講.
不敢講那些是技術性問題.
可以幫他的.
但必須要告訴別人.
為什麼你信耶穌.

$^{1881}$為什麼你會可以信到耶穌.
這些是每個人都獨一無二的.
你不可以找人幫你代講的.
代講是可以.
因為你已經講了.
但是這樣是不夠傳神.
所以每個人都應該有這樣的機會.
去講他自己的得救見證.
這個是很重要的.
家人來聽都是想聽.
誰的見證.
我不是想聽你的見證.
我想聽我的兒子的見證.
就是類似這樣.
所以其他一切都是技術性問題.
時間不夠那些是技術性.
不要緊的.
那些看怎樣調整.
最重要的那件事就真是最重要.
所以第三章我們要記得.
這裡告訴我們.
保羅去形容一個人.
從裡面到外面.
都是不敬畏神.
而來犯罪.
而且很恐怖的一個人.
如果我們犯了這些人.
沒有義人.
偏離了正路.
然後就一直這樣來到.
他的喉嚨是敞開的墳墓.
喉嚨是敞開的墳墓.
我們見到墳墓我們都害怕.
因為是髒東西.
髒東西就是不潔的意思.
不是那些鬼 不是那些髒東西.
就是這樣.
但是是一個敞開的墳墓.
就是發出一些惡臭.
發出一些惡毒.

$^{1921}$果然是.
舌頭是弄鬼炸.
就是說一嘴兩舌.
背後說人壞話.
各方面的東西.
嘴唇裡面有毀蛇的毒氣.
比范鏟頭還毒.
頭上有毀蛇的毒氣.
是嘴唇.
整個的心靈.
由舌頭 由嘴 由口.
這樣來講出來.
其實舊約都很恐怖.
這樣描述罪人.
相信我們自己都有.
過去.
其實我不知道.
中二都講大了.
很小的時候就從聽髒話.
到家裡都講髒話.
到處講髒話.
已經是習慣了.
在母區每個人都講.
沒有人不講.
所以每個人都沒有文化.
我們以為沒有文化才會講.
其實沒有關係.
不關文化的事.
是內心的問題.
所以聖經已經告訴我們.
我們整個的身心靈.
都是這樣壞透了.
殺人流血.
腳是飛跑.
走的路.
是殘害暴虐的事.
去到哪裡就做盡壞事.
平安的路.
簡直已經忘記了.
眼中不怕神.

$^{1961}$所以.
因為目中無神.
就會帶來目中無人.
這樣的比較.
驕傲.
仇恨 嫉妒.
很多這些衍生出來的問題.
都是源於.
每個人的眼中沒有神.
所以我們就求主.
讓我們自己有.
對神存在的敬畏.
大概上就是這樣.
所以第十九.
就是總結.
我們曉得.
律法上的說話.
這個我們剛才講了.
全人的犯罪.
從內至外.
身心靈 手口 腳 頭 眼睛.
你想想保羅.
多麼將人的人性.
很全神地.
還要引用七節的舊約.
在一段這麼小的經文.
引用七節詩篇的舊約.
來作為描述.
所以詩篇某程度.
不是詩人多愁善感的創作.
在引述的時候.
是肯定了詩篇有神的默示.
所講的不是一種創作的靈感.
不存在.
而是真是人性裡面的問題.
但是最後的總結就是這裡.
面臨的是集體的毀滅.
罪其實就是真正人性和世界的本質.
本相.
罪惡就是現在人類的社會.

$^{2001}$我們很擔心核戰.
核子輻射.
但我們很少人會擔心罪惡對人類的毀滅.
我們很少人會這樣留意.
所以基督教.
如果我們不去責備罪惡.
去宣認人需要救恩.
其實根本不需要再存在了.
因為世界是不容許人講罪惡.
不容許人講.
人人都需要救恩.
所以這是我們需要面對人類集體.
面對的這個再不悔改.
就要面對神的審判.
所以我們要努力去傳福音.
將神的名.
要來宣告.
叫人去求告悔改.
得著神的永生.
所以最後的總結就是19,20節.
我們曉得律法上的話.
都是對律法以下之人說的.
好塞住國人的口.
叫普世的人都服在神審判之下.
所以這個就是那個結果.
有沒有留意.
就是我們曾經講過.
一張十八節.
是那個.
對不起.
一張十六,十七節.
是人類的出路.
就是答案.
我不以福音為恥.
這福音本是神的大能.
要救一切相信的.
先是猶太人.
後是希利尼人.
你留意第三章第二十節.
19,20節這裡.

$^{2041}$又再一次講.
所有的人類都面對在神的就坐下.
審判下.
所以那個答案就是.
唯有因信清義.
才可以帶來那個拯救.
所以一張十八.
到三張二十節.
就是那個頭尾的中間的內容.
就是人類目前在神面前的光景.
但是神一早透過保羅.
講了那個答案.
只需要因信就得到生命.
所以三張二十一節.
就是另外一個答案.
就是現在在耶穌裡面.
因基督耶穌的救贖.
就白白清義.
所以一張十六.
到三張二十一.
就是那個三文字.
中間就是人性的描述.
罪性,罪行等等.
今晚我們就到這裡了.
天父我們將我們今天看的三張一到二十節.
和一張十八節到三張二十節.
這個對全人類的描述.
就算是一個這麼好的帝國制度底下.
人都不可以靠自己.
去到神的面前清義.
多謝天父.
因為福音就是神的大能.
可以拯救任何一個相信的人.
但是需要是相信.
所以求主讓我們能夠努力.
將這個要人決志相信的福音.
帶給人去認識.
去得救.
多謝你讓我們能夠有這個時間.
將這一段這麼重要的經文.

$^{2081}$做一個分段.
和做一個簡單的詮釋.
盼望弟兄姊妹能夠掌握.
當我們一路繼續講.
這個恩順清義的發展的時候.
我們更加能夠明白和掌握.
獻上感恩禱告.
奉耶穌基督的名.
阿們.
\newpage



\section{羅馬書}
\label{sec:DvUFUpNSEFw}
\textbf{2025-05-08 周四晚查經班《羅馬書(9):分題:神的義(八)》 講員:曾裔貴牧師}
\newline
\newline
連結: \href{https://youtube.com/watch?v=DvUFUpNSEFw}{\texttt{https://youtube.com/watch?v=DvUFUpNSEFw}} ~~~~ 語音日期: 2025-05-09
\newline
\newline
\hyperref[sec:b2w9fc63m3c]{\small{< < < PREV SERMON < < <}}
~
\hyperref[sec:index_chronic]{\small{[返順時目]}}
~
\hyperref[sec:index_scriptual]{\small{[返順卷目]}}
~
\hyperref[sec:BtALIL64hKM]{\small{> > > NEXT SERMON > > >}}
\newline
\newline
$^{1}$各位大英姐妹晚安.
我們一起同心禱告.
讓我們今天晚上.
羅馬書可以繼續的來查考,講解.
一同的研讀.
我們很盼望.
無論是現場,在線,聆聽的,參與的.
還有當放上去社交媒體或者是YouTube的頻道.
能夠造就更加祝福更加多的人.
去對神的話語有一個基本的框架的明白.
對恩信稱義,對救恩,對神的義.
對神為什麼要將恩典賜給我們人類.
賜給我們犯罪的人類.
我們能夠去探索,去明白.
以致我們跟神的關係.
因為這樣而真的來到有一個恢復.
我們晚上在你面前等候,禱告.
奉耶穌基督的名.
阿們.
也有大英姐妹今晚同步了.
七點半.
在恩福堂又見到今天晚上的一個陪靈會,獻身會.
有大英姐妹又出去了聽.
盼望神也去祝福不同的時宮.
不同的聚會.
都能夠去叫神的兒女得到生命的建立.
我們去看看今天我們的範圍.
因為我們上一個星期.
上一次我們已經處理了.
一章十八至三章二十節.
就是那個題目.
就是萬民都被證明是有罪.
是有不同的形式.
不同的層次的犯罪.
是規缺了神的榮耀.
和即是欠缺神的義.
所以真正的解救的辦法.
就是三章的二十一節.
來告訴我們.
到底神的兒.

$^{41}$是怎樣用什麼渠道.
用什麼方式去賜給人呢?.
這裡就是.
今晚可能我們只能夠處理到三章的三十一節.
結尾.
四章可能說一點點.
但是看看我們的進度.
但是如果一個稍微大的段落.
去到五章第一節之前.
就是由三章的二十一.
你可以延伸到四章的二十五節.
都是想告訴我們.
什麼人可以得到神的義.
或者說神的義.
透過什麼方式.
透過什麼管道.
去輸送給人呢?.
所以.
涉及一個問題就是.
保羅很想在律法以下的猶太信徒.
都要明白.
其實.
恩信稱義.
絕對不是一些全新.
嶄新.
以及違背祖宗律法的一些事情.
這是至關重要的.
所以.
你很快就會看到.
四章突然間會說舊約的亞伯拉罕.
和大衛兩個重要的人物.
就是告訴我們.
其實他們不只是想說舊約.
而是告訴我們.
在舊約.
在神所閱立.
神所拯救的人.
他們都是靠恩信稱義.
來得到神的恩典的.
是這樣跟神立約的.

$^{81}$所以.
你可以想像得到.
保羅的辯論辯證.
是很細密的.
就讓我們看到.
原來萬民都必須恩信稱義.
無論是舊約的人.
就是耶穌還沒降生.
還沒來這個世界.
讓人看到他是耶穌救主的時候的人.
他們都是一樣恩信得到神的那個恩典的.
在耶穌降生之後.
就是現在的新約.
就是使徒們.
就是保羅開始傳這個新道.
這個就更加是恩信稱義.
所以.
保羅雖然說他當時代.
這個嶄新的.
叫萬民可以得救的福音.
但是原來.
福音在舊約.
已經是神透過信心.
去將人拯救.
這就告訴我們.
我們就要留心.
這個恩信稱義這麼重要的課題.
什麼人可以得到神的義.
我們先讀一讀經文.
就是三章的二十一節.
去到三十一節.
這裡告訴我們.
但如今.
神的義在律法以外.
已經顯明出來.
有律法和先知為證.
就是神的義.
因信耶穌基督.
加給一切相信的人.
並沒有分別.

$^{121}$因為世人都犯了罪.
虧缺了神的榮耀.
如今卻蒙神的恩典.
因基督耶穌的救贖.
就白白的稱義.
神設立耶穌.
作挽回制.
是憑著耶穌的血.
藉著人的信.
要顯明神的義.
因為祂用忍耐的心.
寬容人先時所犯的罪.
這個「他」就是神.
好在今時顯明祂的義.
使人知道祂自己為義.
也稱信耶穌.
就是祂的愛子的人為義.
既是這樣.
那裡能誇口呢?.
這很明顯是針對那些.
自以為義.
或者有律法的猶太人.
以為可以憑著自己.
肉身一些宗教的傳統.
以為可以得到神的救恩.
得到神的閱立.
其實是不行的.
所以是沒有人可以誇口的.
是沒有可誇的.
用什麼方法沒有呢?.
是不是用立功之法呢?.
立功就是靠自己的行為.
靠自己以為可以計算.
得到神祝福的恩典的自己的行為.
或者好行為.
其實是不行的.
「不是乃用信主之法.
所以我們看定了.
人稱而是因著信.
不在乎遵行律法.

$^{161}$難道神只作猶太人的神麼?.
不也是作外邦人的神麼?」.
再一次將神創造的做物主.
萬民都是他所創造.
萬民都要向神負責.
所以不是猶太教專利的福音.
更加不是外邦人來拒絕.
猶太人可以得救.
所以這裡你可以想像得到.
我們必須自己一個群體.
我們是新的群體.
再沒有我們自己的背景.
再沒有我們自己的強烈的傳統.
我們自己的族群.
而我們應該要在新的神創造裡面.
我們去看我們是一個.
成為神新創造的一個群體.
所以不再分猶太人.
不再分外邦人.
所以這裡告訴我們.
這個救恩是關乎到每一個信的人.
所以不再局限在那個種族.
局限在他的宗教傳統.
傳承等等體系裡面.
這裡告訴我們.
神不是個別猶太人的神.
更加不只是外邦人的神.
他同樣都是作外邦人的神.
這是第29節.
所以你可以想像到保羅.
他自己有一個很大的豁達的心胸.
因為救恩將他整個人撐大這個心胸.
他自己的背景是一個極之.
對猶太教來捍衛.
就是說他的心胸只是停留在.
對猶太人這個群體.
他有負擔.
其他人就一概無關.
跟他的信念.
跟他所傳的那個誠信.

$^{201}$是跟外邦人沒有分的.
你想想那個心胸是多麼狹窄.
但是救恩耶穌將他完全扭轉改變.
所以我們信耶穌.
如果我們還是停留在一種很狹隘的宗族.
我們自己背的背景.
其實就很可惜了.
救恩不能使我們的心胸擴大.
第30節.
神既然是一位獨有的真神.
他就要因信稱那受割禮的為義.
受割禮就是猶太人這個群體.
也要因信稱那未受割禮的為義.
即是凡是外邦人.
這樣我們因信廢了律法嗎.
斷乎不是.
更是堅固律法.
最後31節.
突然說了一個命題.
說了一個假設.
到底這個命題是否成立的呢.
即是說.
保羅也假設有人會挑戰他.
說一個信耶穌就可以一了百了的福音.
是否即是廢掉了神的舊約系統.
整個律法呢.
這是一個很大的課題.
一個爭論.
你如何能夠處理猶太人.
那種可能是善意的誤解呢.
如何能夠處理呢.
所以.
三章31節.
其實緊接的就是.
如何因信稱而沒有廢掉律法呢.
而且是堅固律法.
即是說.
讓這個律法原原本本的心意.
真正的神的期望.
能夠淋漓盡致的.

$^{241}$這樣去彰顯.
就是第四章的論證.
即是他提了一個命題.
是否廢掉律法.
因信稱而.
是否廢掉了神舊約.
透過摩西頒布給猶太人的律法.
那些神與猶太人納的約.
那些的應許.
是否所有東西都完全廢棄呢.
還是堅固呢.
保羅就很清楚很斷言說.
是堅固律法.
第四章就是.
堅固律法的那個論證.
他提出就是.
原來亞伯拉罕是舊約的人民.
在他未受割禮之前.
因為17章的創世記.
他才受割禮.
但是15章.
神已經稱他為異.
所以第四章.
是很想說一個很重要的題目.
就是一個論證.
保羅現在就要論證.
因信稱而.
不單止沒有廢棄律法.
而且是堅固律法.
怎樣堅固法呢.
原來舊約都是因信稱而.
就是說萬民都是需要因信稱而.
這個就是很重要.
對保羅來處理他那個時代.
整個群體.
整個猶太的族群.
對他的不信任.
對他的懷疑.
對他所說的新的道.
那個保留.

$^{281}$保羅如果能夠解釋得清楚.
是有Point的.
是有舊約的證據的.
那就令到至少不要有一種.
先入為主的一班猶太教的人.
他們可能一下子改觀.
但是至少.
有一些可能是中立.
可能是一些還是抱懷疑的.
透過這樣的解經.
就能夠讓更多的猶太人都可以歸主.
所以.
福音不是硬來的.
福音不是橫斑的.
福音是.
原來在舊約裡面.
已經有這樣的因信稱而.
所以我們.
好像保羅那樣.
神就大大的使用他.
可以將舊約.
能夠這樣來重新去解讀.
其實.
就是說過去他那個解讀.
他那個意思.
那個分析.
原來是在他猶太教.
那個俠義裡面.
但是當主耶穌向他的顯現.
以及他認識.
耶穌就是這位彌賽亞的時候.
他重新看舊約.
他就會看通了.
所以這個.
當然固然是有神的默示在保羅的心靈裡面.
所以.
我們就可以透過這一段經文.
就是三章的二十一.
到四章的二十五節.
神的義是賜給一切相信的人.

$^{321}$舊約都是因信稱而.
新約我們就更加是.
我們就沒有這個困難.
我們可以再看下去.
這個如今這個字.
就不是說一個時間.
而是說耶穌基督來到這個世界.
成就舊恩之後.
那個扭轉.
那個使在律法底下.
在罪惡底下的人.
因為現在基督的降臨.
犧牲.
完成舊恩的復活升天之後.
神的義.
就可以在律法以外.
去彰顯和賜給信的人.
那個運作過程.
我們就可以看看.
Papon裡面.
有沒有留意.
我們很早期已經說過.
一章的十六,十七節.
就已經說了福音是神的大能.
要救一切相信的.
先是猶太人.
後是希利尼人.
因為二人因信得生.
好像已經說了答案.
還沒說哪些人需要因信得生.
哪些人看到自己的需要.
那個欠缺.
要去到神的面前.
靠著信心.
得著這個永生.
去到三章的二十一,二十二節.
你看到這個首尾的呼應.
好像一個括弧一樣.
神的義.
只會加給哪一類人呢?.

$^{361}$就是一切相信的人.
所以.
我們好像在這裡已經找到一個很重要的重點了.
保羅在頭尾.
都說了那個答案.
所以一章十八.
到三章二十節.
就是我們上兩堂所說的.
人類都在神的憤怒.
就坐底下.
必須透過一個另外的救贖.
將人拯救.
離開他自己那個犯罪的沉淪的光景.
所以這裡.
是我們很重要的.
保羅給我們的這樣的括弧.
是的.
所以你會看到.
頭尾都是那個答案.
就是人只需要信.
就可以得到永生.
接著二十二節開始.
就是說那個信的過程.
以及到底在這個信裡面包含什麼呢?.
大家看一看這個關鍵字.
這個因信耶穌基督.
強調其實不是信心本身.
這個指的信心.
不是得救的原因.
或者是那個條件.
信.
更加不是相信一套的教條.
教義.
而是信.
原來在這裡說的.
是對主耶穌的一個傳言的.
認識祂之後.
傳言對祂的毀身.
就是說.
整個的生命投向耶穌.

$^{401}$在耶穌的裡面.
跟祂這樣來去.
同行人生的路.
這裡告訴我們.
第二十二節.
就是神的兒.
因信耶穌基督.
加給一切相信的人.
並沒有分別.
所以這個信.
不是神的信實.
不是真理.
信實.
是在說那個信靠.
所以這裡很重要.
是動詞的一個行動.
信靠.
透過耶穌基督.
所以是因信耶穌基督.
所以信心不是那個條件.
更加不是換取救恩的動力.
而是我們要信耶穌基督.
所以不是信教條.
這裡說有律法.
和先知為證.
其實這裡很籠統.
就是說.
舊約已經有蘊含著.
救恩的啟示.
以及神一直救贖的歷史的發展.
如何來允許.
揀選亞伯拉罕這個民族.
一直將救恩推進.
一直使他的愛子.
要生在猶太這個族群.
這樣來孕育救恩.
或者說一直都在他的範圍裡面.
讓救恩.
讓人類差不多息息相關.
差不多是每一個人都能夠隨手.

$^{441}$就可以來相信.
所以你會留意到.
在第十章有一段經文.
是挺特別的.
第十章的第五節開始.
十章的第五節.
所以有律法為證.
大概是這個含義.
摩西寫著說.
人若行那出於律法的義.
就必因此活著.
這是利未記的經文.
第六節.
唯有出於信心的義如此說.
你不要心裡說.
誰要升到天上去.
就是要領下基督來.
這個所謂誰要升到天上.
是摩西臨終之前在申命記的三十章所說.
就告訴他們.
你們不要硬著頸看.
你們要入迦南了.
我就快要離開這個世界.
我要脫離地上的這個網羅了.
你們這個新生的世代.
你們現在要入迦南.
因為你們的刑罰.
那個妖行在曠野的時間到了.
你們要跟隨若舒雅入迦南了.
所以就說.
誰去帶誰去導我們.
其實摩西就說.
他好像有一種.
其實不是假設.
好像是一種對他們的.
又不是一個挖苦.
而是想告訴他們.
你不要以為要上天下地.
說要去到很遠去尋道.
其實不是.

$^{481}$這個道就在我們的心裡.
這個道就是隨手可得的.
這裡說.
就是要領下基督來.
誰要下到陰間去.
就是要領基督從死裡上來.
他到底怎麼說呢.
保羅引述的就是摩西的三十章的遺言.
他說.
這道離你不遠.
正在你口裡.
在你心裡.
就是這個意思.
保羅借摩西的這樣的告誡.
你們不要上天下海.
因為當時猶太人.
他們寧願拜那些很難拜的偶像.
都不寧願相信摩西的律法.
跟隨神的帶領.
所以就是說.
捨近逃遠.
但是現在保羅就借這一段歷史.
來勉勵.
什麼叫做恩順清義呢?.
其實.
這個道就是在你的口裡.
在你的心裡.
就是我們所傳.
信主的道.
這就是十章九節的來源.
你若口裡認耶穌為主.
心裡信神.
叫他從死裡復活.
就必得救.
所以.
原來真的有律法為證的.
三章的二十二節.
保羅說的這段.
其實真的有舊約.
很多很多聖經的根據.

$^{521}$就是說原來舊約.
都已經在說恩順清義了.
舊約都已經在說.
這個道.
就是這位造物主.
隨時可以被我們去呼求.
我們去求告.
就可以得著這個救恩了.
所以.
我們真的多去學習求告.
多去學習.
跟神有一個很親密的.
屬於我們自己的.
現在說的me time(獨處時間).
就是放下小朋友.
自己有一個寧靜的地方.
可能是看一場戲.
可能是吃一頓飯.
是屬於可能自己.
可能是夫婦兩個人.
這樣的me time(獨處時間).
這些都重要的.
但是.
更重要的就是我們每天都要發展.
有一個跟神之間.
什麼都不要做.
只是閉上眼睛.
安靜去敬拜祂.
有一種這樣的親愛的默契.
一種甜蜜的生命關係.
繼續下去.
所以信心的行動和本質.
都沒有辦法拯救人類.
史奴保羅所說的這個福音.
就是使得人恩信基督.
立即得到這個信基督.
投靠基督.
所以用信靠.
是更加貼切的.
將那個動作反映出來.

$^{561}$所以我們每一個信耶穌的人.
是一個信靠這位救主.
而得到神的義作為回報.
神的義.
住進我們的心靈裡面.
這個不只是一個免除了刑罰.
所以神的義包含的.
不僅僅是一個好像法律地位.
沒有罪了.
就是法庭.
看到一些大奸大惡的人.
用各種的方法.
因為刑事法是要舉證的.
控方要舉證.
辯方就答辯.
是沒有責任去提出任何的東西.
你找到我有什麼證據.
如果證據是毫無疑點的.
就會被定罪入罪.
所以如果入不了罪.
法官可能還很生氣.
雖然我現在找不到你有什麼錯.
但是現在基於法律的定義.
那個邏輯.
我只能夠宣判你沒有罪.
那真的犯了罪的人.
就會被判離開法庭.
就無罪了.
不可以再抓他了.
在那件事情上.
因為沒有證據.
沒有任何的疑點.
就要將他釋放了.
但是法庭的釋放.
是不會說那個傷害和好.
和恩賜的.
所以我們理解的這個恩訊清義.
有別於只是宣判無罪.
宣判無罪是一個基礎.
在這個基礎上.

$^{601}$我們就是透過信主耶穌.
和神建立一個很親密的神人關係.
這個也是第五章.
開始會比較淋漓盡致地去講解.
這個恩訊清義帶來的一個新群體的祝福.
所以我們回到第三章.
我做了一點點的分析.
三章的二十二節.
恩訊耶穌基督加給一切相信的人.
並沒有分別.
出現了好幾次了.
在這一大段經文裡面.
一章的十六節.
要求的是一切相信的人.
三章的二十二節.
是加給一切相信的人.
你看到那個詞彙.
那個字眼.
是這麼強調一切相信的人.
就是說在保羅這個救恩的宣講裡面.
再沒有哪些人特別強一些.
哪一個民族特別蒙恩一些.
是再沒有了.
是已經好像拉平了.
只有信的人.
就成為了一個新的族類.
新的群體.
或者在已不所知就叫做教會.
叫做基督耶穌的身體.
這一班人.
所以這裡告訴我們.
在救恩.
在保羅所傳的福音裡面.
其實很早已經告訴我們.
後來就講宣教.
其實在這裡已經告訴我們.
福音是令到真真正正人類的合一.
因為人類就是在罪惡之下.
三章九節.
你留意下這個圖表的中間.

$^{641}$三章九節.
猶太人和希利利人.
都在罪惡之下.
所以你可以想像得到.
保羅由第一章開始講.
對於神的疑.
哪些人可以獲得.
由神的憤怒講起.
要追討.
要來去刑罰.
虧欠他而榮耀的人類.
所以其實人類在神的眼中.
只是一個無法自救的沉淪的群體.
但是現在神透過保羅.
傳遞這個拯救人類的福音.
就告訴我們.
原來神一直期望要救一切相信的人.
一切不論是什麼民族.
都可以獲得這個恩順清義的福音.
所以你可以想像到.
教會有了與生俱來的責任.
向萬民宣講恩順清義.
向萬民傳講一個人類可以合一的福音.
就是我們無論是什麼民族.
什麼人犯了什麼罪.
都可以在神的救贖裡面.
透過主耶穌.
就可以帶來合一.
帶來永生和清義.
再下去下面我們繼續看下去的時候.
你可以想像得到.
在三章的二十三節.
這些告訴我們世人都犯了罪.
虧缺了神的榮耀.
這裡就很清楚了.
其實是一個撮要.
保羅在一章十八節.
一直到三章二十節.
講的無論是拒絕承認有神.
在各方面的無論宗教道德.

$^{681}$來到每方月下敗壞的人類.
或者自命為有教養.
有道德規範.
有崇高理想的人.
都是陷落在那種的自然.
都是對神的一種背影拒絕.
所以三章的二十三節.
是一個很濃縮的一個總結.
對人類的命運的一個宣判.
就是我們虧缺了做物主的榮耀.
我們在這位創造的主的面前.
失去了我們自己本來應該有的.
神期望的形象.
樣式.
和那種的美德.
我們是欠缺了的.
所以我們可以想像到一個沒有神的人.
他的內心是和神不能溝通的.
他不能祈禱.
他不能認識神.
所以可以有一個生命的分別.
一個就算不識字.
但是他信神信得很清楚.
很單純.
很簡單.
在主耶穌救恩裡面的弟兄姊妹.
就算他的背景.
他與生俱來的條件很有限制.
但是他的內心可以和神很通的.
例如一些可能是有智障的.
信主的弟兄姊妹.
他們未必懂得表達.
也表達不了.
我遇過一個.
很多年了.
因為我教會有一個弟兄.
他很有心.
因為他自己的哥哥.
是有一些智力的問題.
和有些精神的問題.

$^{721}$有一次.
其實每年過年他就去探.
就在博愛醫院隔離那些.
很簡陋的村屋的護老院.
其實沒有什麼照顧.
任由他們每天坐在那裡百無聊賴.
那次的弟兄和另一個弟兄.
那時候是年初二.
就去到.
見到很多很多的院友.
在那裡真的百無聊賴.
過年也沒有家人來探.
那個弟兄就戴了一個口琴.
因為他很喜歡吹口琴.
就和這個弟兄一起.
和那些院友聊天.
吹聖詩給他們聽.
很感動.
又和他們簡單地說一些福音信息.
就燃點了這個弟兄覺得.
這個是一個很好的機會.
我們去探這些沒有人探.
甚至連教會也未必會去探的.
很多這些院舍.
輾轉真的成為了一個很大的事工.
引起很多不同教會的弟兄姊妹.
當知道這個事工之後.
就來支持.
一起去參加這個探訪.
做到很大.
做到要聯絡一些旅遊巴.
由九龍,香港分幾個地方.
有些去新界.
有些去九龍.
這樣去出隊.
這樣探望這些院舍.
其中我親自去到一些很嚴重弱智的院舍.
那些就是政府營運的.
很先進.
很多儀器.

$^{761}$很多設施.
很多工作人員.
其中有一個就是中度弱智.
他講不了話.
他真的每天在床上.
他終身都是這樣的狀態.
但是他聽完福音和聖詩.
他很有感動.
當然他講不了.
很大的反應.
輾轉有人幫他做一些記錄.
大概是這樣.
雖然他說不了話.
給他有這個救恩.
大概的意思是這樣.
你可以想像到.
這樣的生命.
似乎近乎是完全沒有用.
如果我們不是因為對人類.
有一種善心,愛心.
如果在一些比較極權的國家.
這些沒有生產力的人.
其實就要被淘汰.
但是因為基督教的精神.
每個人都有他自己本身的價值.
就算表面沒有價值.
其實他都是一個有價值的人.
所以政府仍然提供這些住宿.
差不多是終身永久.
因為他根本不可能有康復的可能.
但是在他的心靈裡.
你如果能夠有這樣的眼光.
你見到基督在他生命裡面.
基督仍然愛這樣的人.
愛一個輕度弱智的人.
因為他恩信稱義.
所以你可以看到.
救恩的偉大就是這樣臨到每一個不同的人類.
信就可以獲得這個新生.
信就成為了一個.

$^{801}$在神的救贖計劃裡面.
新的群體.
所以我們既然有了這樣的生命.
我們的責任就是.
帶來的.
對內必須是我們要合一.
對外我們必須要讓人.
帶人去到恩信稱義.
帶人去聽明白救恩和他的關係.
他就可以恩信稱義.
我們繼續看下去.
這章24節.
這裡告訴我們原本我們萬民都是虧缺神的榮耀.
無一倖免.
沒有人是不需要救恩.
但是如今.
當然這個如今也是在說基督的出現.
如今卻蒙神的恩典.
因基督耶穌的救贖.
就白白的稱義.
所以是神的恩典的禮物.
就是賜下他的愛子耶穌.
這份恩典的禮物.
是只需要我們去接受.
透過信心.
我們和主耶穌有這個生命的關係.
立即我們的生命就可以得到一個.
嶄新的一個.
新的.
你完全沒有想過的新生命的那種力量.
就是有些人是立即戒煙.
立即脫離那些不光彩的事情.
譬如賭錢的賭人.
立即就沒有了.
很簡單就是.
因為這一份恩典的禮物.
一份新的生命.
進入了那個信的人的內心.
所以這裡告訴我們.
如今卻蒙神的恩典.

$^{841}$記住是因基督耶穌的救贖.
基督耶穌是怎樣救贖我們呢?.
就是他在十字架上為我們白白的犧牲.
記住是一個流血的宗教.
基督教是一個流血的宗教.
如果沒有流血.
我們是不得救的.
基督教是一個很血腥的宗教.
就是基督親自為我們在十字架上犧牲流血.
是他的兒代替我們的不兒.
本來在十字架上是應該我們每一個罪人承擔神的審判.
但現在基督已經代替了我們.
所以因基督耶穌的救贖.
就白白的清兒.
白白是神的一個免費的禮物.
是無條件的.
而且這一份的禮物是無價的.
記得這個無價和無條件是一個很大的關連.
白白不是代表很廉價.
白白不是代表我們可以去藐視.
白白是因為背後有一個無價的恩典.
是這樣為我們付出.
所以第25節.
接著就是提到神的寬容.
神設立耶穌作挽回祭.
挽回祭就是救藥獻祭的一個概念.
那個人要親自按手在祭參.
而祭司驗證過那個祭參之後.
就幫他來獻祭.
挽回了他的罪.
所以這裡特別強調耶穌是作為挽回祭.
好像耶穌基督就是這樣.
在神的面前替我們去代求.
替我們去承受那個罪的刑罰.
所以這個是一個挽回祭.
所以你會看到是一個流血的祭.
耶穌為我們流出他的寶血.
所以是憑著耶穌的血.
藉著人的信.
就是說和這件事直接相連.

$^{881}$十字架就是我們怎樣能夠去到神那裡呢.
就是透過耶穌在十字架上的犧牲流血.
我們在耶穌裡面就立即和神和好.
我們又可以在這個救恩裡面.
和任何一個人都可以和好.
這個就是憑著耶穌的血.
藉著我們的信.
就可以顯明這是神的義.
不是我們的功勞.
不是我們的發言.
是神的義彰顯在我們的生命裡面.
為什麼呢.
就是這裡了.
就是第25節.
因為他用忍耐的心去寬容人先時所犯的罪.
原來一直都是神的忍耐.
神的慈悲.
神的憐憫.
神對我們人類犯罪的人類.
虧欠他 虧欠他榮耀的人類.
來忍耐.
這個忍耐到底有多麼的好大呢.
對我們每一個人.
他有他的時間.
他有他的計劃.
未到時間好像未能夠出手.
所以我們都不要灰心.
不要抱一個絕望.
以為神不能夠改變一些我們認為改變不了的人.
所以我們要抱有這樣的學習.
等候的心靈.
既然神都可以忍耐.
我們就沒有理由這麼快就灰心放棄.
他用忍耐的心去寬容.
所以忍耐帶來的就是.
好像孖公仔一樣.
忍耐寬容.
是同一樣的情操.
造物主原來有這樣的屬性裡面的一種忍耐寬容.
造物主有愛 有公義 有憤怒 有聖潔.

$^{921}$有他的真實 有他的全知 有他的全能.
更加原來可以這麼寬容的.
是這麼忍耐的.
這就是一個神的屬性.
所以我們人類才未至於完全被毀滅.
所以很多人慢慢隨著人生不同的階段.
都會有對自己的改觀.
都會有很多自己的性情的轉變.
這正正證明原來造物主讓時間去沖淡了人對他的誤解.
對他的背逆.
所以慢慢有些人慢慢明白了.
自然會信主.
因為這也是寬容忍耐的運作過程.
所以你可以想像到.
神怎樣對我們自己的寬容忍耐.
我們未信主之前.
都有拒絕 都有嘲笑.
都有說基督教是洋教.
是一些我們完全不懂.
有時還要詆毀耶穌.
但這就是寬容忍耐.
原來造物主是有意識地.
所以是怎樣運作呢.
我們真的不太知道.
我們都叫做忍耐我們的小朋友.
他們有時做錯事.
有時還以為自己很得意.
其實我們心裡都知道.
其實小朋友你都不要怪他.
但你可以想像到.
造物主對我們人類的寬容忍耐.
我自己都是暴風少年.
十三,四歲的時間.
真的好像不可一世.
好像仗著自己有少少的牛力.
就在學校.
但有時那些老師見到.
這些那麼反叛的青少年.
有時搖搖頭就算了.
一來都叫都叫不住.

$^{961}$叫都沒有用.
越叫就越反抗.
所以由得你.
但其實在一個有智慧.
有社會經驗.
有成熟的老師來說.
其實看著你.
知道你一直走向那條路.
只會越來越學壞.
將來沒有一個社會上的.
很好的工作.
很好的身份.
人家有智慧的人都會看得到.
等於現在我們看到.
可能在鐵水圍裡面的一些同黨.
我們都會見到.
你跟他鬥都沒有用的.
因為他現在這段時間.
就是這樣衝.
就是這樣到處撞.
好像一些盲牛一樣.
完全沒有方向.
可能躲起來就是吸毒.
太空油那樣.
各方面的東西.
就是受著整個思潮的影響.
以為自己很厲害.
一班人出來很威風.
但其實你會看到.
那條路就是這樣.
一直朝向的.
那種自我毀滅.
自我放棄.
但神是不會放棄罪人.
他用寬容忍耐的心.
去寬容人先時所犯的罪.
所以這裡.
我們要多些回望.
我們自己的過去.
多些去想想.

$^{1001}$我們未信主之前.
是多麼不好.
被我們心裡謙卑.
因為在教會內.
慢慢有時我們就會聚焦.
在一些很少很無謂的事.
有時就生氣一頓.
生氣人又生氣一頓.
生氣事又生氣一頓.
每個星期回來都生氣一頓.
這樣又走了.
如果是這樣思想的運作.
其實就是不斷地.
掏空自己的心靈.
其實我們要回去.
看回自己.
如何寬容忍耐.
我們先時所犯的罪.
我們才能夠有力量.
看到原來我們今天的改變.
不是因為我背過甚麼玫瑰經.
不是因為我背甚麼聖經.
而是因為主耶穌在我的生命裡.
不斷改變我.
我每一次去敬拜神.
我都有新的那種.
新鮮的.
無論是領受.
還是生命的建立.
是因為這樣.
所以我們作為牧者.
我自己.
我最看不過眼這件事.
看不過眼的意思是.
很慨嘆.
你回來是敬拜神.
回來不是為了.
某些人很少的那件事.
那張椅子怎樣擺.
那些花要怎樣插.

$^{1041}$其實那些很小事.
那些根本不值得.
不值得浪費時間.
應該是我們怎樣去謀求.
我們使到萬民.
能夠恩順.
清義.
所以忍耐寬容的心.
對我們先時所犯的罪.
好在今時.
去顯明他的義.
顯明他的義.
就是藉著耶穌的出現.
在我們的生命裡.
我們就看到原來我們是不義的.
靠賴耶穌.
我們就立即得到神的義.
所以原來.
這個義從來不是我們自己有的.
我們現在更加印證.
是神的義.
加在我們這些相信的人身上.
所以這個就叫.
顯明他的義.
使人知道他自己為義.
也稱信耶穌的人為義.
他自己為義.
即是說.
他從來沒有用一些.
隨便的方法.
就將人算數.
不追究你了.
這些叫濫權.
主耶穌真的為我們人類.
釘死在十架.
面對神的.
自重的刑罰.
所謂的苦杯.
但他都甘願.
去承擔.

$^{1081}$所以.
以至今日我們能夠認定.
神是義.
我們因為在耶穌裡面.
他能夠稱我們為義.
所以.
我們被主耶穌包圍住.
所以神看我們.
是因為耶穌的寶血.
就不再看到我們是一個.
憎恨和虧欠的罪人.
他看到的就是.
在他愛子的救恩裡面.
寶血的救贖裡面.
那個人.
就是一個義人.
所以他能夠.
稱信耶穌的人為義.
第27節.
現在就做一個小小的總結.
保羅.
既是這樣.
那裡能誇厚呢.
沒有可誇的了.
在這個裡面.
既然這裡告訴我們.
稱信耶穌的人為義.
其實再加上一件事.
就是帶來.
我們與神和好.
所以.
我們與神和好.
不是只是判你無罪.
請你快點離開.
我與你無關.
這不是聖經.
神說的恩順清義.
或神的義的定義.
這裡說神的義.
加給我們.

$^{1121}$使我們能夠透過耶穌基督.
這位神的兒子的禮物.
我們與神和好了.
所以這個救恩.
是帶給我們與神的關係的建立.
所以.
這裡告訴我們.
不是只是.
抹掉了我們的罪.
跟著就沒有任何的生命關係.
這個絕對不是基督教.
下面.
這裡告訴我們.
很重要.
我們在27,28節裡面.
既是這樣.
那裡能誇厚呢.
沒有可誇的了.
用何法沒有的呢.
是用立功之法嗎.
不是乃用信主之法.
所以我們看定了.
人清義.
是因著信.
不在乎.
遵行律法.
這個就是27,28節.
不可以誇厚.
信告訴我們.
沒有人.
去在神面前.
透過任何的方式.
可以換取到救恩.
神稱我們為義.
必須是透過耶穌的犧牲.
祂所流出的寶血.
使我們能夠被稱義.
因為神容忍我們.
直至到救恩成就.
我們就得著這個.

$^{1161}$拯救.
就是神的義.
加給不義的人.
所以保羅在這裡特別強調.
既是這樣.
那裡能誇厚呢.
他的答案就很清楚了.
沒有什麼可以誇厚.
用什麼方法.
來不可以誇厚呢.
是不是用立功之法呢.
當然.
這裡暗示著.
仍然有一些人.
拒絕.
去透過這麼簡單的.
恩訊稱義.
去獲得神的救恩.
仍然很想靠賴自己的行為.
可能自己的明星.
或者是自己的家人.
可能自己的明星.
名譽,理想.
或者自己的德行.
或者自己的宗教熱誠.
宗教的熱心.
來希望可以換取神的救恩.
但是.
這裡告訴我們.
不是用立功之法.
而是用信主之法.
是用相信.
就可以獲得.
所以這個是.
我們必須學習.
謙卑在神的面前.
任何東西都是.
信心帶來的.
恩典的禮物.
所以我們要謙卑.

$^{1201}$在神的面前.
另外呢.
第29節和30節.
這裡又好像帶進一個.
既然神是一位.
他又要將恩訊稱義.
給猶太人,給外邦人.
就是想暗示著.
其實這個救恩.
肯定是帶來.
一個神人的和好.
更加是人和人的和好.
我可以先揭一揭.
已不鎖書.
第三章的經文.
可能第二章也有了.
第二章的第十一節.
對不起.
十四節可能也有了.
不如我們讀多一遍.
第一章第十一節.
第二章的十一節.
已不鎖書.
二章的十一節.
所以你們應當紀念.
你們從前.
已不鎖書的讀者.
你們從前按照肉身來說.
是外邦人.
他們的種族.
膚色,血統.
是稱為沒受割禮的.
他們的血統.
是稱為沒受割禮的.
不是猶太人.
這名原是那些憑人手.
在肉身上稱為受割禮之人所起的.
即是你們是未受割禮的.
即是猶太族群的.
猶太教的人.

$^{1241}$去看外邦人的.
一個這樣的稱呼.
第十二節.
那時你們與基督無關.
在以色列國民以外.
在所應許的諸約上是局外人.
特別強調無關局外.
神的國民以外.
並且活在世上沒有指望.
所指的是在猶太人眼中沒有指望.
在猶太人的理念,宗教裡面.
你們是沒有我們的耶和華.
沒有我們的立約.
沒有我們各方面的禮儀.
更加沒有我們這位的神.
第十三節.
你們從前遠離神的人.
因為沒有耶穌.
如今卻在基督耶穌裡.
靠著祂的血已經得親近了.
所以原來基督的血將我們拉近.
跟神,跟人的關係.
這首先是說跟神的關係.
但第十四節是說跟人的關係.
因祂使我們和睦.
原文是說.
耶穌就是我們的和睦.
祂就是我們的和睦.
將兩下合而為一.
猶太人和外邦人.
拆毀了中間隔斷的牆.
有什麼阻礙呢?.
是否宗教的禮儀阻礙了兩個群體的合一呢?.
所以保羅用一個比喻去說.
一個很象徵性的.
有牆阻隔住.
其實這個牆就是.
律法帶來的執著.
固執.
執迷不悟.

$^{1281}$的一種間隔.
第十五節.
而且以自己的身體.
就是耶穌廢掉冤仇.
就是十字架上.
就是那既在律法上的規條.
沒錯了.
猶太教神的律法.
猶太人讀了之後.
以為我們唯我獨尊.
是一神的.
敬拜耶和華的.
其他人如果不進入猶太教.
就被排擠在.
猶太教的舊恩以外.
所以呢.
就有這些隔斷的規條.
就將人.
帶來好像有.
等級.
不同的優越.
這樣來規範了.
所以現在耶穌就是做什麼呢?.
他就是以自己的身體.
去廢掉了冤仇.
在人和人之間.
原來宗教.
原來這些阻隔.
是成為了冤仇.
即是說路死不相往來.
你走你的路.
我走我的路.
但是呢.
這些說要將兩下.
藉著自己.
造成一個新人.
新的創造.
新的群體.
在聖經裡面.
所有信主的人.

$^{1321}$加入一個新的群體.
就是基督耶穌的身體.
不再分什麼種族.
不再分什麼宗教.
接著說.
第十六節.
或者先讀下面.
十五節的下半節.
造成這個新人之後.
如此便成就了和睦.
在這個新人裡面的和諧.
合一.
第十六節.
既在十字架上廢了冤仇.
便藉著十字架.
使兩下歸為一體.
就是這裡了.
與神和好.
所以保羅很重要.
他一直傳講這個福音.
不是想讓猶太人和外邦人.
更加決裂.
更加你走你的路.
大家成為陌路人.
保羅就是成為一個.
復和之子.
成為一個使人和睦的人.
基督已經成就了.
這個抗世的救恩.
將通向神的路.
打通了.
他自己就是那條路.
他自己就是那個真理.
和生命.
所以他更加希望跟隨他的人.
促進這個人.
和人之間的合一.
在救恩裡面.
才能夠合一.
在救恩裡面才能夠合一.

$^{1361}$這是一個我們享用救恩帶來的成果.
合一不是變成我們成為一個會所的會員.
合一是帶給我們.
向萬民傳福音的動力.
我們在這個基礎底下.
向人宣講福音的大能.
向人招手.
帶他去到救恩耶穌裡面.
所以很多時候我們將合一扭曲成為很開心的團契生活.
每個星期見到自己很喜歡的人.
我也一而再講過.
有些人慢慢慢慢只喜歡團體.
不喜歡教會崇拜.
我也很懷疑這樣是不是信耶穌.
或者他自己演繹的信耶穌.
我也很公開地講.
我見過這樣的人.
如果你聽到這些信息.
千萬不要介意.
我們扭曲了你對耶穌的關係.
那個合一是帶給我們和所有信的人成為一個新的族群.
一個新的創造.
我們就在這個基礎上面一起去傳福音.
所以我們不是有一班我很開心的人.
我只是和這班人一生一世.
他們做什麼我就做什麼.
做什麼就是大家聊聊天.
開心一下.
唱一唱聖詩.
接著就這麼多了.
其他人想加入就說不行了.
這個權是不准人加入的.
其實某程度上不應該是這樣的.
我們是應該.
這些告訴我們.
保羅的胸襟.
保羅的國度觀.
是使兩下合而為一.
耶穌已經拆毀了.
用他自己的身體拆毀了中間隔斷的牆.

$^{1401}$用他的身體廢掉了冤仇.
所以在這個基礎上.
我們真的可以成為神的兒女的群體.
就是要合一了.
所以三章29-30節.
就是告訴我們.
我們可以成為合一的群體.
因為你看看我怎麼說.
第29節.
難道神只是拯救猶太人.
只是作猶太人的神.
保羅自己斷言說.
不是啊,不也是作外邦人的神嗎?.
是的,也作外邦人的神.
神既是一位.
所以救恩是一主一信一神.
所以我們一個救恩.
帶來的應該是一個合一的群體.
是的,也作外邦人的神.
神既是一位.
他就要因信稱那受割禮的為義.
也要因信稱那未受割禮的為義.
所以這個就是很清楚的合一的要求.
神也是這樣拯救萬民.
沒錯,不同的時代.
神有他在舊約裡面的時候.
那個救恩的工作.
但是最終透過基督耶穌的出現.
就是要使到萬民因信稱義.
記得不是自己信自己.
信完,使完就好像拜拜一樣.
不是的.
應該就是我們加入基督耶穌的教會裡面.
在這個合一裡面.
不斷地壯大這個基礎.
我們就成為了這個合一的見證人.
接著就是邀請未信的人進入這個愛的團契.
與神與人的相交.
所以這個是很重要的.
我也是有來自福音堂的背景.

$^{1441}$在初信的時候.
我相信屬靈的長輩從來沒有想過.
是會這樣塑造到我們這一班年輕人.
越是熱心的.
就越對自己的教會的傳統.
就越是非常的那種.
好像原教旨那樣.
就不知不覺在經年累月的運作.
就會看一些其他教會不是這樣運作.
不是這樣聽道,不是這樣嚴謹.
就覺得好像是洪水猛獸.
好像不是信主那樣.
那種合一是不知不覺地形成的.
是一個極端.
繼續下去就會變成一種排他主義.
就是自己我們這個群體就可以得救.
其他的就不行.
這個其實是一個胸襟的問題.
對真理不夠全面的問題.
所以保羅是很全面.
講救恩不是只是講個人的.
是講不斷成為了一個群體.
猶太人和外邦人都在神.
因為神是那一位獨一的神.
他的救恩都是只是透過耶穌基督.
所以一個救恩.
一個宗教的教會.
就是神期望地上的教會.
要在合一的基礎底下.
去傳揚神的福音.
所以這裡告訴我們.
你可以想像得到.
三章27-28節.
個人或教會必須謙卑無可誇耀.
一生保持這份謙遜.
一生都看到是恩典托住我們.
能夠事奉.
能夠成為神的兒女.
是恩典的托住.
所以無可誇.

$^{1481}$到我們在教會裡面.
我們的心態就是.
必須尋求合一.
這是很重要的.
最後三章30節.
神既是一位.
他就要恩順稱那受國禮的為義.
也要恩順稱那未受國禮的為義.
為什麼呢.
第31節.
保羅就開始在這裡.
突然間.
有一個轉向.
就是他很闊.
擔心在律法底下的猶太人.
如何理解這個信心.
可以叫任何人都得救.
都得到神的稱義和恩典.
是否廢掉了我們猶太人.
一直以來的信條和律法呢.
這是一個很大的指控.
保羅可以有這麼細密的思想.
想到猶太人的關注.
猶太人對他的猜忌.
他提出一個命題.
到底恩順稱義.
這個信心的義.
有沒有抵觸了律法.
廢掉了律法.
還是如何能夠讓律法.
在這個平衡底下.
更加得到啟示.
更加不單止沒有前後的矛盾.
舊約和新約有抵觸.
這些都是一些很重要的課題.
由保羅在這裡.
他講完這個恩順稱義的運作之後.
接著就講到底.
這個恩順稱義有沒有違背.
律法本身的頒佈和精神呢.

$^{1521}$所以第31節.
是一個很重要的命題.
這樣我們.
保羅一直都在講.
不是他自己.
是我們傳福音的一班使徒.
我們或者被神呼召的當時代的傳道人.
我們恩順廢了律法.
代表那些猶太人.
提出的這個疑問.
這種指控.
到底猶太教和基督教.
怎樣可以找到一個銜接點.
能夠讓大家.
可以這樣在這個信裡合一呢.
所以是一個必須要回答的問題.
所以到底是神啟示給保羅.
他有這個思想.
還是他自己都想到.
他就是那個過來人.
這個就是一個非常寶貴的問題.
我們都未必會這樣想.
我們都未必會想到.
其他人會這樣擔心.
但保羅就能夠在這裡想到這件事.
答案就是.
很簡單.
斷乎不是.
那你為甚麼斷乎不是呢.
你怎知道是斷乎不是呢.
你有甚麼證據是斷乎不是呢.
你說你堅固律法.
是怎樣堅固呢.
甚麼的堅固呢.
所以第四章就是.
舊約堅固了恩順清義.
沒有抵觸到律法.
是堅固律法.
所以保羅就引述.
亞伯拉罕的信.

$^{1561}$亞伯拉罕都是恩順清義.
舊約的人如果有信心的.
都是在神面前恩順清義.
你可以想像到.
保羅的思想真的很周到.
這樣解答對猶太人.
真是一個如釋重負.
因為真是基於聖經創世記.
好幾段的大經文.
來作出恩順清義.
沒有廢掉律法的指控.
恩順清義是更加堅固律法.
今天和大家看的三章二十一到三十一節.
就是恩順清義是怎樣運作.
羅馬書三章三十節.
都很想每一個禮拜.
我們都抽一節來背了它.
都是成為我們大家一個很大的福氣.
神是獨一的.
祂將舊恩賜給每一個人.
所以祂會使到恩順這個信.
給那些受割禮的猶太人.
即是說有一定背景的人.
他們都可以恩順得到舊恩.
亦都可以恩順給那些未受割禮的人.
所以到我們去傳.
我們就要很了解不同的人.
來自甚麼不同的背景.
甚麼的困難.
甚麼阻礙他信主.
我們都要了解.
譬如他不同的年紀.
他對事物的接觸又會有不同.
對一個飽歷滄桑的人.
有社會的閱歷的人.
可能你這樣說.
他都未必會相信你.
但亦都不等於我們用一些很簡單的方法.
他就不會不信.
人生總是有很多的事情的.

$^{1601}$有時在教會裡面.
我們牧者有時都會知.
知我們就不會去張揚的.
有些人有故事的.
那些故事都不方便跟人說的.
但你會見到.
神的恩典就在弟兄姊妹的身上.
真是很奇妙.
所以神可以叫有故事的人.
或一般從小就在教會長大的人.
都可以一齊恩順清二.
這就正正是神是一位.
所以這節聖經可以有很大的.
給我們的領受.
我們自己先領受了.
我就是未受割禮就被稱為二.
我們是外邦犯罪的人.
稱為二.
我就是有過一個暴風少年的階段.
有五年左右迷失.
來到去賭錢.
離家出走.
不務正業.
有五年是完全荒廢了.
那個人生.
還有一些真是不方便說.
這樣的人生.
但神的寬容.
忍耐的心.
讓我能夠成為他的兒女.
第二年就已經在建道.
神學院的斜路.
有聲音跟我說要做傳道.
那個時候是1988年的八月.
第一次去參加夏令營.
那個時候是一些剛剛信主的夏令營.
那晚聽完道.
心已經很大的感動.
很想奉獻.
不知道是什麼.

$^{1641}$不知道什麼叫傳道.
傳道人又很害怕.
不要那麼緊張.
不要那麼衝動.
先冷靜下來.
人生的路就這樣走到今天.
所以寬容忍耐的心.
神在我們生命裡面.
很盼望今天給這一節聖經成為大家的一個護面.
下個星期.
今天就不再說了.
都可以休息一下喉嚨.
辛苦的都有點不舒服.
下個星期我們就第四章.
舊約如何恩順清二呢?.
以及恩順清二帶來什麼與神和好的福氣呢?.
我們就停在這裡.
天父我們感恩.
因為你的愛.
藉著愛子耶穌.
讓我們知道你自己為義.
亦都是稱信耶穌的人.
我們得著這一份的義.
是無條件的.
是無法誇口的.
是神的寬容忍耐.
多謝你.
雖然我們不知道.
做物主是什麼形象.
或者是什麼模樣.
但是我們知道.
你就是那一位有真理.
有人意.
有聖潔.
更加我們知道.
你是有寬容忍耐的心靈的真神.
我們每一個信主.
都是你自己的傑作.
都是你自己手的奇妙的工作.
願你幫助我們.

$^{1681}$越聽羅馬書.
越知道要合一.
要傳福音給萬民.
因為求告主明的.
才可以得救.
所以幫助我們.
不要介意我們帶人去基督那裡.
不要介意會被人嘲笑.
甚至乎侮辱.
但是我們知道.
我們值得去付這些代價.
願主施恩.
現在在恩福堂的.
這個.
見道的大會.
盼望神都使用.
葉博士的呼召.
真的有年輕人.
或者不同年紀的.
他們清楚自己在神面前的呼召.
回應上主你的Calling.
來做傳道.
去接受裝備.
去拯救其他未信的朋友.
和家人等等.
我們這樣感恩抱告.
奉耶穌基督的名.
阿們.
\newpage



\section{羅馬書}
\label{sec:BtALIL64hKM}
\textbf{2025-05-15 周四晚查經班《羅馬書(10):分題:舊約中的因信稱義(一)》 講員:曾裔貴牧師}
\newline
\newline
連結: \href{https://youtube.com/watch?v=BtALIL64hKM}{\texttt{https://youtube.com/watch?v=BtALIL64hKM}} ~~~~ 語音日期: 2025-05-16
\newline
\newline
\hyperref[sec:DvUFUpNSEFw]{\small{< < < PREV SERMON < < <}}
~
\hyperref[sec:index_chronic]{\small{[返順時目]}}
~
\hyperref[sec:index_scriptual]{\small{[返順卷目]}}
~
\hyperref[sec:QTpL9O5GR3E]{\small{> > > NEXT SERMON > > >}}
\newline
\newline
$^{1}$各位晚安 我們一起有個祈禱.
感恩我們再一次來聚集.
去研讀《羅馬書》.
我們去到第四章的經文.
特別是關乎到亞伯拉罕.
在舊約的聖經人物.
他是怎樣能夠堅固律法.
以及和這個恩信清義有什麼關係.
都盼望我們能夠透過保羅他的展述.
能夠掌握,能夠認識.
以致我們能夠在整個神學的框架.
不會有一種錯誤的觀念.
舊約就是律法.
新約才是恩典.
我們很盼望神給我們認識.
整個聖經的那個救恩的歷史.
以及神的救恩的那個計劃.
我們等候在你的面前.
我們會處理到相關的舊約的經文.
盼望神給我們能夠有一個明白的心去認識.
我們這樣感恩禱告.
奉耶穌基督的名.
阿們.
上一個星期我們提過.
這個恩信清義.
其實只是一至四章的經文.
當然哪些人需要清義呢?.
就是萬民都欠缺神的榮耀,神的義.
所以其實每一個人都必須來信主.
藉著耶穌才能得到神的救恩.
所以恩信清義是說每一個人都需要去到神的面前.
得到這個白白的恩典.
但是在整個的羅馬書.
反而有一個很重要的主題.
既然萬民都需要救恩.
萬民都需要恩信清義.
以及不再在保羅的神學裡面.
其實在這個救恩底下.
根本就沒有分猶太人.
或者是希臘人,外邦人.

$^{41}$那個不同的等級.
不同的階級,不同的民族.
所以在救恩底下.
萬民都是平等的.
所以信了耶穌.
這個就是往後我們就見到.
一方面我們需要向萬民宣教.
但是更加重要的就是教會必須合一.
不要管是甚麼民族.
不要管是甚麼不同的意見.
都需要在那個救恩底下.
來去謀求那個一切的同心.
重點仍然是要希望神的福音去到地極.
所以羅馬書就是這個大的主題.
反而就不是恩信清義是一個大的主題.
當然每一個在教會裡面的人都是恩信清義.
不然根本就不可能在神的教會裡面.
能夠有新的生命.
大家彼此的相愛,彼此的同心去服侍.
所以我們第三章稍為有些總結.
就是在這個恩信清義底下.
就是說神的義在律法以外來彰顯出來.
就是恩信耶穌基督是白白的清義.
所以是必須讓人看到不是自己的功勞.
去換取救恩,去得到神的義.
所以所有的人在救恩的下面.
都需要謙卑,承認自己的無助.
而得到神的救恩是白白的,是無條件的.
我們就可以成為神的兒女.
所以這個救恩帶給我們.
其實第五章就會詳細講.
其實是建立了關係上的恢復.
與神的關係的一個恢復.
就不只是一個宣判無罪.
那個很負面.
只是不再追究我們的罪.
那個是一個很基礎的開始.
但其實是真真正正在關係上的一個建立.
所以三章二十九,三十節.
也都是再一次.

$^{81}$其實不止這裡.
是一直在整卷書都流露出.
萬民都需要合一的.
在救恩底下.
神看我們不再分我們是什麼人.
其實第四章都有的.
所以我們又要有一個總結.
就是三章二十九至三十節.
是告訴我們.
我們需要在這位獨一的神面前.
我們必須學習一個合一的功課.
合一的真理.
接著我們就來做這個開始的總結.
之前的最後.
就是三章的三十一節.
因為保羅突然間帶出了一個很重要的一個命題.
到底恩順清二.
這個這麼嶄新的一個真理.
在保羅的年代.
你想一下是二千年前.
當時猶太教還是在耶路撒冷有憲制.
有那些不同的節期.
猶太人仍然看自己是神選民的身份等等.
所以到底保羅怎樣去處理.
舊約對每一個猶太人的要求呢?.
或者舊約.
如果我們看得不清楚.
好像是規範了外邦人必須要透過加入猶太教.
才可以得到神的眷顧.
所以突然間保羅傳出一個.
似乎好像撇開了所有的舊約.
就可以透過耶穌.
就可以恩順清二.
甚至進入和神的關係上.
這個對猶太人來說.
尤其是一些拉比,法利賽人.
對他們來說是一個很大的侮辱.
一個冒犯.
到底保羅是不是一個離經叛道.
離開猶太教.

$^{121}$叫人不要守律法的一個人呢?.
所以保羅也假設.
或者他也意識到.
他也需要處理猶太人的不滿.
猶太人可能是善意的誤解.
怎樣來疏離.
以致真是有根有據.
又合乎聖經這樣來推的一個恩順清二.
如果能夠處理得好.
那是帶來一個很大的人類的突破.
就不會再分猶太教以外的人.
這是舊約一直都是這樣運作的模式.
所以他就反而自己去提出這個問題.
到底這個恩順清二是不是廢掉了律法呢?.
當然不是.
但是恩順清二是怎樣去兼顧律法呢?.
所以你可以看到在三章二十一節.
這個如今神的意.
在律法以外顯明.
到了四章的第一節開始.
就好像一個很重要的分水嶺.
保羅都為了讀者.
不是一個好像不清不楚的一個講述.
而導致信了主之後.
可能還會被猶太教的人逼迫.
或者影響或者誤導.
你們不加入猶太教.
你們不守國禮.
你們就不夠純正了.
沒有我們猶太人那麼純正.
那就會帶來恩順清二好像要靠一些行為這樣.
所以無論如何.
保羅他都是很有意識地要告訴我們.
其實跟舊約不僅沒有抵觸.
原來保羅再重新去演繹舊約的時候.
當然不是一套嶄新的演繹.
而是舊約原來一直都有信心.
恩典這些精髓在其中.
所以第四章就立即有一個很重要的聖經人物.
就是亞伯拉罕.

$^{161}$所以保羅要解釋到底恩順清二是如何兼顧律法.
很自然地就必須要回到舊約.
保羅採取哪一個聖經的人物.
是足以讓就算是猶太人或者外邦人.
都能夠心悅誠服地看到.
原來舊約都有恩順清二.
完全是可以很有理性地.
很合理地 很有point地.
來解除猶太人的敵視.
所以如果我們能夠在辯論一些真理的時候.
能夠很平心靜氣.
很有真理的基礎.
有神的智慧這樣來解說.
就不只是一句到尾.
令到那些懷疑的人只會更加懷疑.
更加不滿.
但如果我們能夠有point.
所以你看到第四章是一個很重要的補充.
因為第五章就再一次回到.
那個恩順清二不只是宣判你無罪.
恩順清二是帶來同神關係的一個恢復.
在這個關係上我們怎樣在人生裡面.
與神同工 與神同在.
所以第四章是一個很重要的補充.
就是這個命題.
怎樣兼顧律法.
怎樣uphold舊約的律法.
原來舊約不單止我們沒有違背.
原來你們看錯了.
以為舊約要你做什麼什麼.
其實不是 原來真正在神面前的舊約的聖徒.
他們都是憑信心在神面前的.
那就很不同了.
原來恩順清二從來不是暫生的.
是一路都是恩順清二的.
一個很重要的解說.
所以保羅將亞伯拉罕帶出來.
我們要問這個問題了.
初期教會 修信羅馬的教會.
或者當時周圍的地方 歐洲的地方.

$^{201}$信主的弟兄姊妹.
他們很多不是猶太人.
或者可能有些地方是有猶太教的會堂.
但是到底亞伯拉罕在這些外邦人的心靈裡面.
是一個陌生人.
還是大家已經有一定的概念.
對亞伯拉罕有這個認識呢?.
如果不是從頭開始講過.
好像是需要費很多的唇舌.
還有很多的功夫.
但是你看到第四章.
其實是很深入地.
假設那些讀者.
已經是對《創世記》15到17章.
是已經很療癒指掌的.
我不知道在線的,現場的.
你們對《創世記》亞伯拉罕的故事.
由他離開吳爾到他一直成為信心之父.
你有多少的掌握.
他的年歲,他的家庭的遭遇有多少呢?.
就是說保羅似乎假設.
羅馬教會是很認識亞伯拉罕的人生的.
就是說亞伯拉罕不是一個陌生人.
從來都不在外邦人的世界裡面出現.
就是說原來在教會有很多關於這些舊約的人物的討論.
這個時候對保羅的講解就相對容易很多了.
當然我們今天沒有證據.
沒有一個人告訴我們.
亞伯拉罕我們很熟悉的.
我們沒有這些證據.
但是我們只能夠在經文裡面去假設.
保羅的探討其實是很深入的.
很多時候他有他的見地.
待會我會跟大家一起去讀的時候.
他怎樣去運用《創世記》.
那個經文的由人有餘的這樣遊走.
還有連那個引述都告訴我們.
保羅都是相當透徹的了解整個亞伯拉罕的人生的.
所以他是非常運用自如的.
我們就感覺到透過第四章讓我們看到.

$^{241}$原來舊約的表面的經文.
我們以為很熟悉.
但是原來我們再進一個再深度的一個認識.
或者有計劃地去這樣再重新去研讀.
又會有一些新的聖經舊約的時候.
沒有這樣去反映出來的.
看到保羅無論神給他的智慧,啟示.
還有他對舊約的純屬的運用.
都是值得我們去借鏡和學習.
我們就在這裡第三章帶出這個問題.
是保羅點出來要處理猶太教一直以來的困難的.
但是唯一的答案就是第四章這裡.
原來舊約就是已經有恩順清二了.
那怎樣有恩順清二呢?.
哪個聖經的人物.
舊約就是說一定是一些故事.
一些歷史人物.
他們是怎樣反映他們的人生是恩順清二呢?.
所以就是第四章我們開始進入今天的經文.
這個PowerPoint我用了很多時間.
今天下午做的了.
想到我頭都快爆了.
雖然是一頁而已.
但是裡面包含了保羅很多的心思.
他很多巧妙的討論.
我們先回到經文第四章.
由一到十.
十二節先.
這一段已經足以告訴我們.
舊約的人的得救.
其實真的是憑信心.
不是憑遵循律法的.
那就看看保羅怎樣去討論.
怎樣去帶出他的論點.
命題就說了.
恩順清二怎樣去堅固律法.
不是廢掉律法.
那怎樣堅固呢?.
就是說如果舊約也有人是恩順清二.
那就是堅固律法了.

$^{281}$不僅沒有抵觸.
原來我們今天新約是一文雙成.
神都是看每個人的信心.
去得到神的恩典的.
第四章由第一節開始.
如此說來.
我們的祖宗亞伯拉罕.
由頭到尾都沒有叫亞伯蘭.
這些巧妙之處.
還有他引的經文.
十五章第六節.
他都用亞伯拉罕.
其實舊約十五章.
如果熟的時候.
當時亞伯拉罕叫甚麼名字?.
叫亞伯蘭.
所以第十五章的創世記第六節.
亞伯蘭信耶和華.
耶和華就以此為他的兒.
但是保羅引十五章.
他不是就這樣平鋪直敘地引.
是他自己加了注釋.
就是亞伯拉罕信實.
這就算為他的兒.
就是說十七章.
十七章是過了十三年了.
亞伯蘭是八十六歲.
繼續他的大兒子以實瑪利.
接著他的小兒子都還沒有出生.
就是這個以實還沒有出生.
後來去到十七章.
由十六章的最後一節.
亞伯蘭那時候是八十六歲.
接著十七章.
就突然間亞伯蘭九十九歲了.
耶和華向他顯現.
接著就幫他改名.
所以亞伯拉罕信神.
其實他已經是九十九歲的時候才改名的.
十幾年之後才改名的.

$^{321}$但是保羅好像不是很熟悉那個歷史.
其實不是這個意思.
其實保羅一直是以亞伯拉罕.
就是看亞伯拉罕.
是一個一直得著神的信心.
憑信心來活在神面前的.
一個屬靈的領袖,長輩.
所以你會看到這些巧妙.
就不要小看.
以為亞伯拉罕引錯聖經.
對不起.
以為保羅引錯了聖經.
其實不是的.
他已經加入了他自己的神學判斷.
他已經加入了他自己對亞伯蘭這個人.
由舊變新的信心的歷程.
你可以想像得到.
保羅就是真的要將亞伯拉罕帶出來.
讓羅馬的教會知道.
其實你們不需要懼怕猶太教對你們的威脅.
對你們的藐視,攻擊.
因為舊約都已經有恩順清義.
舊約的聖經人物.
亞伯拉罕他都是恩順清義的.
第四章由第二節.
他說倘若亞伯拉罕是因行為清義.
就有可誇.
就有可誇就是回到三章那裡.
這個恩順清義是無可誇的.
就是再一次保羅用一個一連貫性的詞彙.
去將舊約,新約去無縫的接合.
就是說就算亞伯拉罕他信神.
都不是因為他有好的行為.
而是他真是在無助的時候.
相信神的應許.
所以神就看他這樣的信心.
來去稱他為義.
稱義他這個信教對神的應許.
就是他的兒子以撒的出現.
這樣來得到神的應許.

$^{361}$就是說他成為萬國之父.
亞伯拉罕的意思.
希伯來文的意思.
還有就是他的兒子以撒.
成為真正的繼任人.
繼承他的產業.
這裡告訴我們並無可誇.
接著第三節經常說什麼呢?.
說亞伯拉罕信神.
這就算為他的義.
這是《創世記》第十五章第六節.
我們讀完經文十二節之後.
我們就回到《創世記》.
先看那些引述的經文.
還有當時的背景.
我們再回來看.
你就知道保羅他的引述.
是真的很巧妙.
第四節.
他講了一個分析的標準.
什麼叫做恩典.
什麼叫做靠行為.
什麼叫做你自己賺回來的.
這個就無法誇耀.
第四節.
作功的得功價不算恩典.
就是是應份的.
我付出了我的勞力.
付出了我的智慧.
給一個公司,一個機構.
或者等等的.
譬如在田裡的工作.
我拿回我的功價很合理.
這個交換而已.
不是什麼恩典.
不是什麼憐憫.
但是是該得的.
第五節.
唯有不作功的.
只信稱罪人為義的神.

$^{401}$就是沒有靠任何的功德.
沒有靠任何的行為.
其實若有所指的.
猶太人仍然存在一個割禮.
是一個符號.
是一個標記.
是一個夢神喜悅的儀式.
一個禮儀.
就是說我們猶太人.
跟你們外邦人不同的地方.
因為我們有割禮.
你們外邦人是沒有的.
所以這個似乎是一個符號.
所以你去看到第十二節.
保羅是刻意打破了.
這個絕對不是夢神恩典的記號.
這個只不過是.
亞伯拉罕的一個.
未受割禮的時間已經信的信心.
但是當他受割禮.
只不過是為了那些.
沒有這樣做而信的人的一個印證.
一個記號.
待會我們去第十二節繼續說.
接著告訴我們第五節.
他說.
「只信稱罪人為義的神.
他的信就算為義.
正如大衛稱那在行為以外.
夢神算為義的人是有福的.
他說」.
就是詩篇的三十二篇.
「得赦免其過.
遮蓋其罪的.
這人是有福的.
主不算為有罪的」.
舊約就是耶和華.
「耶和華不算為有罪的.
這人是有福的」.
就是大衛一個自述.

$^{441}$在詩篇裡面.
他向神的認罪.
神赦免他.
他寫這篇詩篇.
很想表達他對神的感恩.
神寬恕他.
所以這個不是作弓.
這個是白白的.
無條件的一個恩典.
所以保羅想要帶進的.
就是這個信心的義.
然後第九節.
「如此看來」.
就是一個小小的總結.
他說這個福.
這個恩信稱義的福.
憑信心而得到神的恩典的福氣.
是不是只是單單加給那些受割禮.
就是猶太人.
「不也是加給那未受割禮的人嗎?」.
然後他要解釋.
到底是還是不是.
是只是加給受割禮的人.
還是加給未受割禮的人.
那就要留意這裡了.
因為保羅引述的十五章的《創世記》.
是亞伯蘭當時未受割禮.
就已經信耶和華.
耶和華已經向他顯現.
告訴他.
你的後裔要像天上的星.
海裡面的沙.
那麼多的後裔.
亞伯蘭在那個時候.
他沒有做任何事.
他好像張開了嘴巴.
都啞口無言的.
來接受神的這個應許.
他就相信.
當然神就看到他有這個信心.

$^{481}$就不是只是一種厭口枉.
而是他看到他內心的信心.
就稱他這個信心為義.
就是說他已經得著神對他的應許.
雖然還未兌現.
他的小兒子以撒也還未出生.
十五章的時候是完全未出生的.
但是神好像已經看到亞伯蘭有這個信心.
而得到神的喜悅.
那麼亞伯蘭繼續說.
是不是只是加給受了割禮的猶太人呢?.
其他的未受割禮的外邦人就沒有份呢?.
那麼亞伯蘭就引述了.
第九節的下半節.
因我們所說.
「亞伯拉罕的信就算為他的義」.
是怎麼算呢?.
意思就是說.
在什麼情況下算呢?.
如果保羅在這裡突然間凸顯割禮.
又要將亞伯蘭他未受割禮來去信呢?.
這就是一個很重要的論點.
就是讓當時似乎在保羅的年間.
叫人,勉強人受割禮.
加入猶太教.
就是一個很重要的切入點.
就是蒙神的恩典的標記.
一個記號.
但是保羅在這裡.
就好像想淡化割禮.
他不是反對猶太人行割禮.
因為猶太人基於傳統.
他們的民族受割禮是沒有問題的.
而也不要因為受割禮.
就變成得天獨厚.
變成有特殊的身份.
特別的權利.
那就是錯誤了.
所以變相外邦人也不需要刻意去受割禮.
才可以得到神的信心的恩典.

$^{521}$所以現在他就要強調.
是怎麼算呢?.
亞伯拉罕是怎麼被算為義呢?.
是不是因為割禮他被算為義呢?.
如果是的話就困難了.
就好像有了案例.
後世的人都需要這樣來跟隨.
但是我們看到亞伯拉罕.
他的被算為義就不是割禮.
而是單單他信耶和華的應許.
所以我們待會去到四章二十五節.
保羅就做一個最重要最重要的人類的總結.
就是耶穌出現在這個世界.
耶穌被交付給那些敵人.
就是為了我們的罪惡.
祂的復活就是為了我們能夠清義.
那就很清楚了.
等於就是亞伯拉罕什麼都沒有做.
神現在有一個應許.
你信我.
你的子孫就好像天上的星一樣.
亞伯拉罕什麼都沒有做.
就接受了造物主的承諾.
這個應許.
那麼他就立即被稱為義了.
所以原來我們的恩順清義.
我們的那個應許.
就是神已經將他的愛子交付了.
將他的愛子復活了.
我們就可以得到這個恩順清義.
所以又是回到神的義在律法以外.
這樣顯明了.
所以三章二十一和四章二十五.
又是一個頭尾的神應許了給信的人.
的這個白白永生的福氣.
但是要解釋.
就是到底恩順清義是怎麼兼顧律法呢?.
所以這個第十節怎麼算呢?.
我們就要來留意了.
是不是在他受割禮的時候呢?.

$^{561}$是在他還沒受割禮的時候呢?.
到底是還是不是呢?.
答案就是第十節的下半節了.
不是在受割禮的時候.
乃是在還沒受割禮的時候就恩順清義了.
這個是猶太人接不接受呢?.
沒錯.
表面上的確是.
但是我們一路再看下去的時候.
因為下面由第十七節到二十一節.
就很詳細將《創世記》的十七章.
做一個很濃縮的解釋.
阿普羅.
就是說他很想帶我們回到《創世記》.
回到好像時光隧道那樣.
回到亞伯蘭.
他怎麼經歷神帶領他的人生.
怎麼能夠一路的成長呢?.
這個是比我們在今天.
當然我們今天就沒有面對猶太人對我們諸多的挑剔.
說我們不受割禮.
我們根本已經過了那些年日了.
但是在保羅的年代.
外邦人信主.
其實面對的四面八方的嘲笑,攻擊.
羅馬帝國要拜多神的.
你不拜.
你要拜一神.
你要拜耶穌基督.
這是一個攻擊.
但是猶太教又說你不受割禮.
你不守節旗.
你不跟我們一樣.
你又不是我們神的選民.
又變成另一種嘲諷,侮辱.
他未必會逼迫你.
但是會對外邦人信主.
會有一種歧視.
會有一種誤解.
你怎麼分薄我們的福呢?.

$^{601}$你有什麼身份去分我們在神面前的福呢?.
所以這段比較重要.
去拆解.
其實不需要兩個民族有這種對立面的.
而是我們回到舊約.
保羅說.
其實連我們這麼有信心的這位亞伯拉罕.
稱為信心之父的這位人物.
原來他也有他的軟弱.
有他的成長階段.
他得到神的應許的時候.
根本就未受割禮.
不是刻意要淡化.
或者保羅要扭曲那些經文.
而是根本那些經文的人生的順序.
他根本都未受割禮.
所以割禮不是一個獲取恩典的記號.
這個就很肯定了.
所以拆開一件事.
回到信仰我們的基督教.
如果聖經沒有堅持的傳統.
是我們自己慢慢約定俗成的一個傳統.
就不要看為一個絕對的真理.
去灌輸給任何一個時代的神的兒女.
因為聖經沒有應許過的東西.
我們就不要看為好像是祖宗家法不可以改.
可以是一個傳統.
可以大家都覺得這樣做會比較舒服.
不等於是絕對.
現在問題就是猶太人就看割禮是絕對的.
你不受割禮你就不可以得救.
我們有割禮的記號.
我們就是蒙神的恩典和喜悅.
保羅就等於很間接很委婉地去拆解了.
從來割禮都不是一個記號.
是記號不過不是一個唯一的蒙福的記號.
真正蒙福的記號就是信心在神的應許上面.
所以這裡告訴我們第十一節.
並且這裡就說記號了.
很快我想去下一個部分.

$^{641}$因為我將英文都帶出來.
這裡就是第十一節.
是很重要的經文.
我先看英文.
這裡告訴我們他獲得割禮的記號.
原來是一個銷.
英文NASB是一個銷.
是一個印記.
埃及法老約瑟未從監獄出來.
解釋了七個方年就快來了.
現在是七個豐年.
現在就要去掘米去囤積那些田地.
到方年來的時候.
就可以開倉派米.
不然大家都會同歸於盡.
法老聽完之後.
你解了我的夢.
然後就將他手的戒指給了約瑟.
這個戒指就是古代的人.
一個像文書的蓋印.
所以是一個銷.
一個印記.
所以這裡英文的NASB都是用銷.
這個不是用一個銷這麼簡單.
所以這裡中文沒有英文講得這麼好.
他說這個割禮的印記.
原來就好像一個錄印.
一個標記.
一個印鑑.
是一個有信心的人的印記.
就是在他未受割禮的時候.
我highlight了.
Uncircumcised.
未受割禮的時候有這個印記.
他就可以成為了.
Father of all who believe.
without being circumcised.
就是說他能夠成為未受割禮的人.
一個指標示範的單位.
我們看回十一字的經文.

$^{681}$並且他受了割禮的記號.
在他未受割禮的時候.
欣順清議.
果然就是印證這個字.
印證這個字就是那個字.
就是那個好像一些印鑑.
好像我也講過在《異筆所書》.
受了一個印記.
就是古代他們要買大批的木材.
那些買家或者買手.
就會去到那個木材的地方.
就看到這批木材我要了.
你什麼時候就斬好之後.
通常就扔到河裡浮下來.
免及那個運輸.
因為那些牧會浮.
那他們怎麼做呢?.
這批貨是屬於我的.
我現在就用我的公司.
以前的木材.
用一個印.
好像我們做月餅那些印.
大銅月餅,恆香月餅.
就蓋上一個印.
就證明這批貨我將來就會要了.
就是那個印記.
所以這裡保羅告訴我們.
亞伯拉罕他在未受割禮的時候.
就已經欣順清議的一個印證.
一個銷.
他就已經欣順清議了.
他就能夠作為一切未受割禮.
而只是相信神的人.
他就是他們的父.
就是說他是所有信神的人的父.
就是說他自己第一個信.
神就好像在他的生命裡面.
給他一個印記.
一個印證.
你就是一個好像示範單位.

$^{721}$就成為日後都是這樣.
他就是一個憑信心.
相信神的應許的人.
不管那些人是猶太人.
還是外邦人.
都可以得著神的應許.
這個就是保羅想要去解釋的.
所以我們回來這裡告訴我們.
其實舊約裡面.
我們的這位亞伯拉罕.
他能夠得到神的救恩.
同樣都是憑信心.
不是憑律法.
保羅在這裡的解釋.
其實是很重要的.
接著他不僅是說未受割禮的人.
實際上他們也算是為異.
留意一下第十二節.
保羅其實是想說萬民都是靠恩典得救.
在保羅的救恩裡面.
這個福音其實是給任何一個.
開放自己心靈的人.
就可以得到的.
絕對沒有你的背景.
你的種族.
你的文化.
可能有一些幫助.
但不是先決的條件.
先決條件是神看我們有沒有信心.
所以十二節.
「又作受割禮之人的父.
就是那些不但受割禮」.
保羅多加一句.
受了割禮而你沒有亞伯拉罕的信心.
你只是一個猶太教的人.
你沒有在神面前有真正真實的信心.
你也不可以因為你是猶太教.
受割禮就可以得到神的應許.
因為保羅真的很有智慧.
並且你受了割禮.

$^{761}$但是你也要按照我們的祖宗亞伯拉罕.
他未受割禮而信的信仰的宗旨去走.
你才稱為有信心的人.
他講的是猶太教的割禮分子.
你可以想像到保羅的救恩.
原來涵蓋的神的愛.
神的救恩涵蓋的就是.
就算你沒有受割禮.
從來都不知道甚麼是割禮.
但是你知道有耶穌的應許.
你相信你就可以得救了.
很多穆斯林的群體.
如果他們根本沒有聽過其他的東西.
他只是憑他的內心相信耶穌.
他就可以得到永遠的救恩.
但是猶太人就要多一件事.
他如果只是從小長大.
知道有猶太教.
照著那些規矩,禮儀去生活.
但是他一直都沒有信心.
他也不行.
因為要根據好像亞伯拉罕在他人生.
未受割禮的時候.
那個信心的生活.
這樣去依靠神.
是那個宗旨.
所以17節就開始講亞伯蘭的宗旨.
他的信仰宗旨是怎樣的呢?.
作者保羅就列舉了好幾件事.
尤其是他憑信心去等候以實的出現.
所以來到這裡一到第十二節.
首先第一個很重要的論點.
就是其實舊約.
當然以亞伯拉罕做的單位.
已經是恩信稱義了.
我們現在就要回到看《創世記》.
跟大家看回那個歷史.
怎樣去組織保羅所說的經文.
那個信心之旅.
就是亞伯拉罕.

$^{801}$或者亞伯蘭的信心之旅.
我們從第十五章可以看起.
十四章就是講到一場戰爭.
就是因為羅德他的侄子.
被所多瑪王去捉拿了.
當然當時是五王的.
七王和五王的.
四王和五王的戰爭.
所以這場戰爭.
所多瑪王應該大敗了.
後來亞伯蘭就去召集他自己318個壯丁.
去打贏了那場仗.
也拿回了財物.
輾轉就出現了一個叫做麥基駛德的祭司.
為亞伯蘭祝福.
接著輾轉.
第十五章.
這時以後.
作者告訴我們.
又過了一段時間.
亞伯蘭那個時候.
他還是住在萬里象樹的地方.
接著告訴我們.
耶和華在異象中有話對亞伯蘭說.
亞伯蘭你不要懼怕.
可能他怕人尋仇.
因為他打贏了那些四王.
接著.
你不要懼怕.
我是你的盾牌.
必大大的賞賜你.
亞伯蘭說.
主耶和華.
我既無子.
你還賜我什麼呢?.
並且要承受我家業的大馬色人.
以利以謝.
就是他的僕人.
第三節.
亞伯蘭又說.

$^{841}$你沒有給我兒子.
那生在我家中的人.
就是我的後製.
很清楚與神有這樣的對話.
有這樣的交換.
就是說不用這麼麻煩了.
我的僕人就當是我的收養的繼承人.
後製就是繼承人的意思.
第四節.
耶和華又有話對他說.
這人.
就是以利以謝.
這人必不成為你的後製.
你本身所生的.
本身就是由他自己.
就是與太太撒萊所生的.
才可以成為他的後製.
才成為你的後製.
第五節.
留意一下.
於是就領他走到外邊.
怎麼領呢?.
是一種榮光.
舊約的時候神怎麼顯現的呢?.
十七章就有兩個天使.
怎麼來向亞伯蘭顯現.
是不是牽著他的手走出去呢?.
怎麼跟他說話呢?.
總之亞伯蘭是可以接到這些信息的.
不是報夢的.
真的走出去的.
是日光日白.
所以神可以在日光日白向來顯現的.
是的.
然後就叫他看了.
數算眾星.
白天應該看不到眾星的.
能夠數得過來.
可能就是晚上也不一定了.
亞伯蘭就留意第五節.

$^{881}$對不起.
第五節還沒讀完.
又對他說.
你的後裔將要如此了.
提出了這麼大的應許.
不是信口開河.
是神跟亞伯蘭說.
而且是根據上下文.
肯定不是告訴你.
所有的星球都是你的了.
對亞伯蘭來說.
根本就沒去過外太空.
根本就不知道.
連太空館也沒去過.
根本不知道有多遙遠.
因為古代的人怎麼可能有這樣的技術呢?.
但是.
看到眾星肯定比我們今天多.
因為沒有那麼多污染.
在晚上的星空.
是非常壯觀壯麗的.
但是很清楚的一個應許.
肯定這些星是跟後裔有關的.
因為上文是在說.
你不能生.
你所生的.
才是你的後裔繼承人.
亞伯蘭就說不要這麼煩了.
我的中心的僕人.
我過世之後.
就全部給他就行了.
他是我繼承人就行了.
都是屬於我的.
全部給他就行了.
神就說不行.
必須是要你身所生的.
才是後裔.
接著就真的帶他去看戶外.
去觀星.
告訴他們.

$^{921}$你的後裔將要如此.
就是一個這樣的應許.
聖經就在這裡.
留意第六節.
保羅說是亞伯拉罕信耶和華.
不是亞伯蘭信耶和華.
好像保羅引錯聖經一樣.
其實不是.
保羅已經在他一直寫的時候.
就已經認定亞伯拉罕.
已經是信心之父的人物了.
保羅他是後世的人.
當然知道亞伯蘭一直的成長.
過了十幾年之後.
他都是這樣跌跌碰碰的.
經歷這個信心的考驗.
輾轉他就真的神幫他改名.
但是保羅好像已經立即幫他改名.
保羅就說亞伯拉罕信神.
耶和華就以此為他的兒.
所以你可以看到保羅那種.
這麼巧妙地.
來將他自己的神學.
灌注在這些經文裡面.
但是在這個時候.
在十五章的時候.
亞伯蘭是還沒改名的.
要十三年之後才改名.
輾轉神就跟他說.
我要跟你立這個約.
你是我將這個地賜給你為業的.
然後就告訴他.
你要怎樣在這個約的儀式.
怎樣去做.
就教他要將一隻三年的母牛.
一隻三年的母山羊.
一隻三年的公綿羊.
單溝,初鴿就拿來.
準備獻祭.
所以每樣都劈開.

$^{961}$分成兩半.
就是死了的.
一半對著另一半擺.
就是這樣平擺.
就是說準備開壇作法的意思.
就是獻祭.
但是只有鳥沒有分開.
只是那些牛羊,公綿羊就分開了.
就是劈開了.
那些單溝就整隻了.
那些初鴿就整隻了.
沒有劈開.
然後第十一節.
有飼鳥下來.
落在那死畜的肉上.
就是說過了一段時間.
阿巴蘭就把牠嚇飛了.
就是走過來.
那些蒼蠅拍開牠們.
因為牠走過來想吃獻祭的死畜.
接著是十二節.
日頭正落的時候.
就是黃昏.
阿巴蘭因為太久了.
沉沉的睡了.
不過留意.
忽然間有驚人的大黑暗.
落在牠身上.
就是說神降臨.
耶和華就對阿巴蘭說.
你要的確知道.
你的後裔必寄居別人的地.
又服侍那地.
當然是埃及.
那地的人.
那地的人要輔待他們四百年.
原來以色列人要寄居埃及四百年.
其實是做奴隸.
因為要被輔待.
已經預言說了.

$^{1001}$接著.
並且他們所要服侍的那國.
我要懲罰.
後來他們就是以色列人.
就是你的後裔.
必帶著許多財物從那裡出來.
但你要享大壽數.
127.
阿巴蘭人.
平平安安歸到你列祖那裡.
他的列祖就是阿當.
以洛.
以羅氏.
馬土撒拉這些先做.
羅亞這些.
因為阿巴蘭人是第一個被神由吾爾帶離開.
去到迦勒地.
接著你留意一下.
歸到列祖那裡被人埋葬.
到了第四代.
就是你的第四代.
阿巴拉罕的第四代就是.
阿巴拉罕,以撒,瓦割.
接著就是約瑟勒.
第四代.
那四代之後第五代就會出了.
所以他們很快就要.
不過留意這裡.
他說你們就會歸回到這個地方.
就是迦勒地.
為什麼呢?.
原來因為阿摩利人的罪孽.
還沒有滿盈.
就是說原來神的計劃是很雙向的.
原來他們就被老.
不是被老.
他們自己下去埃及.
去吃米.
後來就在埃及那裡一直生活.
因為祖宗約瑟得到王的喜悅.

$^{1041}$所以整個民族下了埃及.
接著就一直都回不來了.
接著原來最終的原因.
是我們都沒有想過.
神現在就很早告訴阿巴蘭.
因為阿摩利人的罪孽還沒有滿盈.
就是說他們為什麼回不來迦勒呢?.
原來迦勒地的原居民.
還沒有罪大惡極.
還沒有得到神最終的審判.
所以以色列人還要等待四代.
在埃及地.
這些就是超出了我們對歷史的新的了解.
原來神已經很早告訴阿巴蘭.
你的後裔將來會這樣這樣.
將來那個世界會這樣.
出現這樣的情況.
很多事情表面到時就會自然地演變.
但是原來背後.
原來神是正在積蓄阿摩利人的罪孽.
當地迦勒人原居民.
他們拜偶像.
他們是很得罪神.
到了一個真的滿盈了.
神就要刑罰了.
神怎麼刑罰呢?.
原來是透過亞伯拉罕的後裔.
約瑟的子孫.
他們回到埃及.
當然是摩西帶領他們的民族離開埃及.
來到迦勒地.
或者約瑟帶領他們回到迦勒地.
原來已經是在四百年前.
神已經跟阿巴蘭說了.
但是現在這個重點就是.
這個應許就是阿巴蘭還沒有改名.
中國人就很重要.
改名就是轉運.
我姐姐也改名了.
但是阿巴蘭還沒有改名.

$^{1081}$就是他這個時候.
還是一個很普通的人.
但是到了17章.
神的應許的時間到了.
他要生一個兒子就是以撒了.
神就說我要將你這個舊的名.
阿巴蘭改為阿巴拉罕.
這個名字不是因為轉運.
是因為神的應許的時間到了.
你應該具備信心之父這個名字.
但是在15章.
是完全沒有這個信心之父的名字的.
但是保羅已經用阿巴拉罕.
這樣去定性了.
阿巴蘭其實已經因為他的信心.
來得到神的應許.
不是因為國禮.
因為17章才行國禮.
整個希伯來人才行國禮.
所以如果我們熟悉這個歷史.
我們就知道保羅為什麼這麼振振有詞.
恩信稱義其實是肩負律法.
不是否定律法.
更加不是將律法架空.
原來舊約也已經有恩信稱義了.
舊約的阿巴蘭.
都是因為他對神的應許的信心.
其實信心有多大我們也不知道.
但是神就看到已經足夠了.
就稱阿巴拉罕為二了.
當然我們覺得蠻恐怖的.
立約的儀式.
第17節.
日落天黑.
神的出現很特別.
有無煙的爐.
並燒著的火把.
從那些玉塊中經過.
神這樣經過被阿巴蘭看到.
當那日.

$^{1121}$耶和華與阿巴蘭立約說.
我已經賜給你的後裔.
從埃及河直到伯拉大河之地.
就是基尼人,基尼洗人,甲摩尼人.
亨人,比利洗人,尼弗欽人.
阿摩尼人,迦南人,甲加薩人.
耶布斯人之地.
整個迦南都是阿巴拉罕的.
神跟他納了這個約.
就在這個時候.
這個儀式跟他立約就是當那日.
納了這個約.
這個民族就這樣蒙福.
是因為阿巴蘭的信心.
因為阿巴蘭信耶和華.
這樣來得到.
接著十六章就回到日常生活.
十六章憲制完.
納了約.
阿巴蘭的妻子撒萊不給他生兒女.
撒萊有一個侍女名叫夏甲.
是埃及人.
這個時候.
這個夏甲原來當時在十一章.
阿巴蘭因為饑荒落了埃及.
在埃及怕他自己的太太撒萊.
被法老王抓了她做妃嬪.
就跟撒萊說.
現在我們去埃及了.
我也不知道當地的王是怎樣的.
你就當是我的妹妹.
其實也可能是同父異母等等.
某程度上也不是絕對的謊言.
但也是一個欺騙.
果然法老看到撒萊真的想娶她為妻.
果然就真的將她.
不是抓的.
而是用很多聘禮拿了.
後來神就向法老報夢.
你娶撒萊.

$^{1161}$你想死嗎?.
你快點歸還給阿巴蘭.
否則你的家族也會有事.
法老就很害怕.
就叫阿巴蘭來.
阿巴蘭就說我很害怕這個地方的人.
所以我就認我的太太做妹妹.
她就變成了一個大寇.
是不是大寇?.
哥哥嘛.
大寇.
就拿了很多聘禮.
不是聘禮.
是那些禮物.
送走阿巴蘭和撒萊.
下鴿是其中一個禮物.
就是帶她上去迦勒地.
所以現在就下鴿.
說得很明顯她是埃及人.
妖后那些埃及人.
撒萊對阿巴蘭說.
耶和華使得我不能生育.
求你和我的侍女同房.
或者我可以因為她而得到孩子.
阿巴蘭就聽從了撒萊的話.
於是阿巴蘭的妻子撒萊.
張侍女埃及人下鴿.
給了丈夫為妾.
那時阿巴蘭在迦勒已經住了十年了.
就是八十幾歲.
八十五歲左右.
阿巴蘭與下鴿同房.
下鴿懷了孕.
就是說.
原來不關阿巴蘭的事.
生不了就是撒萊不能生.
就不是阿巴蘭不能生.
現在下鴿一同房就有得生了.
就懷孕了.
她見自己有人.

$^{1201}$就小看她的主母.
是不是可能真是因自己的兒子.
將來繼承的兒子.
都不知道.
都沒有去做DNA化驗.
也沒有照超聲波.
怎知道是兒子.
可能是女兒.
都不知道.
但是她已經開始.
因為真是自己身懷六甲.
就開始對主母.
這個撒萊.
開始有一些不同的眼光.
第五節.
撒萊對阿巴蘭說.
我因你受屈.
我將我的侍女放在你的懷中.
她見自己有了人.
就小看我了.
怨謠在你我中間判斷.
阿巴蘭對撒萊說.
侍女在你手下.
你可以隨意待她.
撒萊輔待她.
她就從撒萊面前逃走.
我們可以想像到.
撒萊怎樣去輔待夏甲.
古代的人怎樣去輔待那些妾.
是用言語.
還是具體的行動.
關她去柴房.
不讓她吃東西.
要輔到一個地步.
她要從主母的家逃跑.
然後跑到曠野.
因為耶和的使者.
在曠野的樞爾的路上.
似乎想回外家.
因為回到埃及.

$^{1241}$一直在曠野.
輔到一個地步.
一個有了身孕的婦人.
痛苦到要逃跑.
相當痛苦.
因為其實古代的人.
尤其是一個下人.
一個妓女.
一個要奴隸.
被人送走的妹子.
其實根本就沒有社會地位.
根本就沒有人的尊嚴.
但是她都要苦到一個地步.
寧願逃跑.
我們可以想像到.
撒萊原來都不是善男信女.
撒萊都有些手段.
對付夏甲好像綽綽有餘.
聖經就沒有描述.
她有多苦.
她的手段是多惡毒.
都不理她.
又身懷六甲.
還有肚子裡的嬰兒.
是亞伯蘭.
不是其他的僕人.
你可以想像到.
有時那種失控.
都不理那麼多.
撒萊就固然在這裡.
應該要受到神的譴責.
但是.
某程度上神又沒有反對.
因為夏甲所生的.
都不可以繼承亞伯蘭的產業.
還要是撒萊自己的肚子裡生的.
才可以繼承神的產業.
所以好像默許做這件事.
我們一路不說那麼詳細.
接著我們突然跳到第十五節.

$^{1281}$就是講夏甲聽了天使的話.
一來就是.
她太辛苦太受屈.
她走而已.
其實她沒有全盤計劃.
沒有收拾乾糧.
指南針.
手提電話回埃及.
因為一個女流之輩.
還大了肚子.
你怎樣能夠.
頂得住在曠野的生活.
所以她只想吐口苦水.
突然迷了路.
天使就向她顯現.
叫她回去.
要伏在她的主母的腳下.
就是說你回去.
主母的氣都過了.
都會讓你生的.
還有你的兒子.
某程度上因為你.
因為亞伯蘭.
他都是無福的.
天使這樣安慰.
所以天使都很重要.
都救了兩個生命.
因為如果不是.
夏甲可能被野獸擔心.
那就沒有了.
但是你可以想像到.
有時候那口氣.
都需要有一個空間.
給夏甲他自己又可以明白.
是喔.
我非理性地逃跑.
都是不行的.
回去可能就是.
改變她的態度.
你不要再小看主母.

$^{1321}$其實你要自己.
麻煩你照鏡子.
你不是因為你的兒子.
你就水鬼升成王.
這樣是不行的.
她開始明白.
降低她的氣焰.
當然主母殺來.
就知道原來自己都過份了.
就容許.
既然都是丈夫的骨肉.
就容許.
這個妾侍.
都在家裡.
生了這個二十晚.
沒錯.
生了這個二十晚.
十五節.
後來夏甲.
給亞伯蘭生了一個兒子.
亞伯蘭給他起名叫.
二十晚.
夏甲給亞伯蘭生二十晚.
的時候.
亞伯蘭年八十六歲.
八十五歲.
夏甲就被.
跟亞伯蘭.
來到同房.
是阿薩萊的一個.
獻祭.
但是這告訴我們.
當時這件事.
叫做一個小小的風波.
暫時可以完滿的解決.
那個大妻子.
終於可以.
順產二十晚.
這個時候.
亞伯蘭已經是八十六歲.

$^{1361}$留意十七章.
有十三至十四年.
一個真空期.
在一節聖經裡面.
作者.
如果真是摩西.
好像十三年十四年.
是完全沒有任何的記錄.
讓我們有一個空間去聯想.
二十晚.
他們一路長大.
在亞伯蘭的家裡.
有兩個男士.
當然不算亞伯蘭.
有一個叫以利爾謝.
那個老的僕人.
那個已經講到明不是你的繼承人.
不可以的.
後來獻祭.
由我的侍女所生的.
歸入我的名下.
亞伯蘭就說.
好啊 也對.
因為真是由我所生的.
是我和夏甲所生的.
或者我和由我一個女士所生的.
都屬於我的骨肉.
可以是神的.
應許的產業.
他可以繼承.
但是.
以實瑪利出世.
一路由嬰兒長大.
然後滿月走.
然後八日宴.
然後一路成長.
越來越大了.
十三歲都不小了.
在這個時候.
以實瑪利是十三歲.

$^{1401}$但是聖經是完全沒有神的.
再介入.
再顯現.
也沒有告訴我們.
亞伯蘭每天弄兒為樂.
撒萊很滿足.
什麼叫姊妹.
相安無事.
大家快樂融洽.
那些雨場面的大結局.
亞伯蘭就成為.
齊人之福.
又有兩個老婆又有一個兒子.
就可以繼承產業了.
十三年就這樣過了.
突然之間十七章.
亞伯蘭年九十九歲的時候.
顯現很晴天霹靂.
告訴我們.
神又有些新的東西.
要亞伯蘭去遵守.
就是說那十三年的天倫之樂.
是廢的.
不是天倫之樂不對.
而是那個應許的兒子.
你想多了.
以實瑪利不是繼承人.
必須要撒萊.
你可以想像到.
其實亞伯蘭不斷經歷.
越來越深的信心功課.
越來越不容易.
因為十三年已經定了.
是吧?.
是要和他一起工作的.
但是突然之間.
晴天霹靂神就再顯現.
第一次開始.
亞伯蘭年九十九歲的時候.
耶和華向他顯現.

$^{1441}$對他說.
我是全能的神.
你當在我面前作完全人.
我就與你立約.
就納了.
十五章就納了.
說應許了我的子孫.
原來那個信心的考驗.
越來越深入.
所以信主的內.
我們的信心應該越來越深入.
越來越堅固.
作完全人.
我要與你立約.
使你的後裔極其煩多.
是吧?沒錯.
以實瑪利就是那個繼承人.
第三節.
亞伯蘭苦服在地.
神又對他說.
我與你立約.
從此以後你的名字.
不要再叫亞伯蘭了.
要叫亞伯拉罕就是多國之父.
因為我與納你.
作多國之父.
一直都沒有碰到實瑪利.
以為可以過關.
我必使你的後裔極其煩多.
國度從你而納.
君王從你而出.
這段.
就是.
加泰書三章十九節.
那個應許的言子就是耶穌.
就是這段聖經的語言.
我要與你.
並你的世世代代的後裔.
建立我的約.
作永遠的約.

$^{1481}$是要作你和你後裔的神.
我要將你現在寄居的迦南地.
全地.
賜給你和你的後裔.
永遠為業.
我也必作他們的神.
很重要.
不只是地圖.
地圖之上的.
這一位的.
主神.
所以我們今天說.
這個教會.
不只是一個物業.
還有神的同在.
我應許每一次你們的聚會.
我會同在.
這樣才重要.
不是那個硬件.
而是神成為他們的神.
第九節.
神又對亞伯拉罕說.
你和你的後裔必世世代代.
遵守我的約.
你們所有的男子都要.
受國禮.
在這個時候過了十三年之後.
才正式頒布這個.
國禮的約.
這就是我與你.
並你的後裔所立的約的記號.
是你們所當遵守的.
你們都要受這個國禮.
這是我與你們.
立約的證據.
也是那個紀號.
所以你看到.
保羅引述經文的時候.
他是非常.
充分了解了.

$^{1521}$整個亞伯蘭的成長.
信心的歷程.
的軌跡.
所以.
他去講解.
恩順清二在舊約.
原來在亞伯蘭.
他還沒有受國禮之前.
這麼久之前.
已經經歷了.
他曾經給他的應許.
而他相信.
就可以得到這個證據.
所以第十二節.
你們世世代代的男子.
無論家裡生.
是你們後裔之外.
用銀紙在外人買的.
即是那些奴隸.
生下來第八天都要受國禮.
你家裡生的.
和你用銀紙買的.
必須受國禮.
這是神頒布的命令.
是一個條款.
但是原來.
神的這些頒布.
是沒有抵觸到.
他第十五章的時候.
同樣都是這位做民主.
來去應許亞伯蘭.
同樣都是受益人.
就是亞伯蘭 就是亞伯拉罕.
所以恩順清二.
在舊約已經有了.
不是.
行國禮才可以得到.
這個恩順清二的福氣.
恩典是源於.
相信神的.

$^{1561}$應許.
所以亞伯蘭.
就是相信神的應許.
我們今天後世.
我們是屬靈的.
亞伯拉罕的子孫.
我們都是恩順.
相信神的應許.
就是他的愛子被交付人.
來為我們的罪.
他復活.
是為了我們能夠.
在神面前.
建立神人的和好關係.
所以你可以想像到.
保羅的論點.
是非常清晰.
印證舊約的日.
同樣都是恩順清二.
這個就是.
為什麼恩順清二.
不是廢棄律法.
是更加堅固律法.
所以當代的猶太人.
仍然堅持受國禮.
受國禮才是無福的證據.
是錯誤的.
猶太人可以受國禮.
但是國禮不是.
蒙神的.
祝福的條件.
必須是要.
因為猶太人有這個傳統.
受國禮是沒有問題.
但是你都要根據亞伯拉罕的.
信心的軌跡去相信神.
你才能夠無福.
所以猶太人.
受國禮是不會有特別的優待.
外邦人.

$^{1601}$不受國禮更加不會是一個欠缺.
是一個很重要的.
將兩個民族拉平.
在恩典底下.
大家都是透過恩典.
去蒙神的祝福.
所以有時我們都要.
不斷迅速地.
其實要苦心自問.
我們慢慢因為信心.
因為恩典而改變了.
我們不要慢慢就變成.
有一種優越感.
看不起那些未信的人.
他們的行為.
他們的想法.
我們慢慢就有一種.
癡之以備.
或者不知不覺我們自己有一種優越感.
就忘記了.
其實我們都是蒙恩.
在恩典裡面神改變我們.
神賜福給我們.
我們能夠家庭幸福.
各方面都能夠順暢.
我們走的人生路的軌跡.
又能夠走在.
那個二路上面.
但是我們要記得.
我們要回復那個初心.
未信主之前.
我們都是一樣這樣敵擋神.
在神的憤怒底下.
招致神的憤怒的人.
我們就會.
有一份胸襟.
去接待.
接納一些未得救恩的人.
他們那些錯誤的行徑.
例如說髒話.

$^{1641}$例如一些不好的行為.
我們當然不是縱容.
但是我們要明白.
他未有恩典.
在他們的生命裡面.
自然就會有這些種種的錯誤.
對神的那種敵擋.
或者嘲笑.
所以這裡告訴我們.
留意一下.
我們下面這裡不到.
十五節.
神又對亞伯拉罕說.
你的妻子撒萊.
不可以再叫撒萊.
她的名要叫撒拉.
我必賜福給她.
也要使你從她.
得一個兒子.
我要賜福給她.
她也要作多國之母.
就是撒拉的意思.
必有百姓的君王.
從她而出.
再一次去肯定.
以撒的後裔將來.
是坐在這個皇座上.
就是我們的救主.
耶穌基督.
第十七節.
亞伯拉罕就俯伏在地.
喜笑心裡說.
一百歲的人還能得孩子嗎?.
撒拉已經九十歲.
大十年了.
還能夠生養嗎?.
亞伯拉罕對神說.
但願以實馬利.
不是以利以謝.
以利以謝已經肯定不是.

$^{1681}$已經是飛了出去.
局外了.
十五章還在說.
但願以利以謝.
在你面前無福.
現在就不是了.
現在是她的大兒子.
以實馬利.
但願.
什麼意思呢?.
就是信心不夠.
不要這麼複雜.
這個兒子已經萬事足.
夠了.
撒拉其實不應該再生了.
古代就沒有高齡產婦的概念.
神就像有一個玩笑.
九十歲了.
真的九十歲了.
實際上不應該再生了.
現在是一個反建議.
亞伯拉罕就告訴神.
但願以實馬利也夠了.
在你面前活.
神就說.
不是啊.
你妻子撒萊要給你生一個兒子.
你要幫她改名叫以撒.
我要與她.
堅持生一個兒子.
我要與她堅定所立.
和你的約.
就是你要傳遞你的恩典.
你的約給你的兒子.
哪個兒子?.
以撒,還沒出生的以撒.
堅定這個約.
就是說意味著.
神再一次這麼說.
再一次確認這個約的時候.

$^{1721}$就是雙重確認.
意味著.
亞伯拉罕.
信神這件事.
不過這個信.
是戰戰兢兢.
難度高了很多.
因為明明有個大兒子在.
難度高很多.
太太就已經.
十年前就告訴我.
現在九十歲叫她生.
我都不好意思.
但是.
神就說.
這就是考驗你的信心.
要再進一層級.
所以我們不要再.
相信那些見證.
我們要不斷有一個新的經歷.
實在我們生命裡面的.
大靈.
在這個時候你可以想像到.
第二十節.
至於以實瑪利.
我也應允彌.
我必賜福給他.
使他昌盛的其繁多.
他必生十二個族長.
我也要使他成為大國.
這是很大的安慰.
作為一個父親.
你想怎樣?.
你說他沒有繼承權.
你也要給他一個好的下場.
好的晚年.
神就說我應允彌.
他也就是你的兒子.
我也應允彌.
使他成為一個大國.

$^{1761}$他就會.
留意第21節.
到明年這時節.
撒拉必給你新耳神.
我要與他堅定所立的約.
神和亞伯拉罕說完了.
就離開他上升去了.
即是暗示.
作者暗示了一個.
很重要的訊息.
即是15章第6節.
亞伯蘭信耶和華.
耶和華就以此為他的義.
17章第23節.
我們應該要假設.
亞伯拉罕在這個時候.
神說了兩次.
我要與你未出世.
都未行訪.
的兒子.
明年才出現.
神就說已經說了兩次.
我要與這個未出世的兒子.
堅立.
與他的約.
所以第23節.
和第22節中間.
我們應該要假設.
或者應該要肯定.
這句說話.
就是亞伯拉罕再一次.
信耶和華.
祝福他.
再一次信耶和華.
耶和華就以此.
為他的義.
大家明不明白這個邏輯.
就是這個很重要.
亞伯拉罕在這個時候.
是另一種信心.

$^{1801}$對神的信靠.
很不容易的.
因為你現在好像在一個降鮮上.
我就100歲.
太太就90歲.
但現在是要我們兩個老人家.
生的兒子.
才可以繼承.
但神已經說了兩次.
這個人就叫做以撒.
明年就會出現.
其實是多有趣.
多令人憧憬.
亞伯拉罕知道之後.
第23節證明.
我的假設是真的.
因為上文說了一個命令.
就是神向亞伯拉罕顯現之後.
幫他改名.
之後他自己就叫亞伯拉罕.
然後有一個命令.
就是你和你所有的男丁.
要立即行割禮.
所以23節.
正當那天.
神走了之後.
離開他上升之後.
正當那天.
亞伯拉罕就遵著神的命令.
給他的兒子以實瑪利.
和家裡的一切男子.
無論在家裡.
很清楚沒有計錯數.
所有的男子.
總之和他有關的.
用銀紙買的.
都行割禮.
第24節.
亞伯拉罕受割禮的時候.
年99歲.

$^{1841}$也是呼應17章所說.
神在這裡向他顯現的時候.
他的兒子以實瑪利受割禮的時候.
就是13歲.
接著第26節.
正當那天.
亞伯拉罕和他兒子以實瑪利一同受割禮.
家裡所有的人.
無論在家裡生的.
用銀紙從外人買的.
也都一同受割禮.
要再加多一個.
很重要的訊息.
在這個時候.
亞伯拉罕.
其實他是不是再一次.
堅定他的信.
來相信.
神會在我這個殘區.
我太太的殘區.
是不是應該成立呢?.
今天.
今天的90歲.
和100歲.
和亞伯拉罕的年代.
100歲有沒有不同呢?.
我們又沒有證據.
不知道.
他自己都說自己是100歲.
怎麼可以生呢?.
大概不會相差很遠.
你可以想像得到.
我們肯定可以加這個.
就是在這個時候.
既然他遵守神的命令.
去行割禮.
也就是說.
他都應該根據.
神已經兩次說.
明年這個時候.

$^{1881}$10月懷胎.
你的太太就會生一個叫做耳實.
這個耳實才是.
真正的繼承人.
有一個很大的神蹟.
亞伯拉罕正在期待.
所以這裡暗示.
雖然作者沒有說.
亞伯拉罕信神.
會給他一個兒子.
大家有沒有留意.
第四章17節.
到21節.
保羅就很平鋪直敘.
將亞伯拉罕的心路歷程.
這一刻的心路歷程.
就說出來了.
我們回到第四章的17節.
所以.
剛才17章的創世記.
22,23節.
是應該加一個.
我們自己加下去都沒有錯.
就是亞伯拉罕在這個時候.
再一次要.
相信神的應許.
他明年會有這個兒子.
這個不是天方夜譚.
不是自己兩腦.
在無事生事.
而是真的從神的應許.
而且還是兩次說.
我會和你的兒子.
堅定約的.
就是說這個兒子一定會出現.
如果不出現就很大件事了.
神好像都收不到科.
肯定是有一個兒子.
好像都收不到科.
肯定不是.

$^{1921}$所以我們透過這個經文的組織.
我們看看第17節開始.
很奇妙的.
保羅告訴我們.
這段經文真的這樣發生.
第17節第4章.
第13節我不看了.
我們看第17節.
亞伯拉罕所相信的.
是叫死人復活.
使沒有變為有的神.
證明他自己當時.
已經很清楚.
他們兩個是不可以生的.
撒萊一直都不可以生的.
所以他看自己是死人.
在生育上.
所以他在這個時候.
他的信心.
原來可以成為一個新約的演繹.
就是叫死人復活.
使沒有變為有的神.
就是剛才說的22,23節.
那個心路歷程.
他內心的那個.
你這樣說我一定要信.
你說我將來一定.
不是將來.
明年就一定有這個耳塞.
我一定要信.
所以就等同於.
由沒有變為有.
由死人可以復活.
這樣的信心.
那就是我們的信心.
我們就是透過耶穌的復活.
我們去相信神的應許.
所以保羅就將.
亞伯拉罕的內心世界.
用真理.

$^{1961}$來描繪出來.
好像畫了一個圖畫.
所以很巧妙地將.
舊約的精髓.
保羅用新約的言語.
來講述出來.
接著留意一下.
使沒有變為有的神.
他在主面前.
作我們世人的父.
就是信心之父.
如經上所記.
我已經立你作多國的父.
第17章第5節.
他在無可指望的時候.
因信.
因信亦有指望.
就是他的太太.
可以因為他.
而生這個.
應許的兒子.
耳實.
所以原來亞伯拉罕在這個時候.
再一次承諾他自己.
那個心靈.
在神的應許上.
比當年更加難.
但是他信.
神就.
以此為他的義.
所以他現在能夠真是叫做.
多國之父.
這個就是.
無可指望的時候.
因信亦有指望.
聖詩就是憑信.
阿伯拉就有這個信心.
接著第19節.
他將近百歲的時候.
就是神應許給他的時間.

$^{2001}$雖然想到自己的身體.
如同已死.
但生育已斷絕.
他的信心.
還是不軟弱.
並且仰望神的應許.
總會有因不信.
心裡起疑惑.
反倒因信心裡得堅固.
將榮耀歸給神.
保羅.
將整段聖經.
最重要的17章.
描繪出來.
就是記得信心是建立在.
不是我們自己的幻想.
是建立在神的.
兩次的應許.
你的太太.
明年這個時候.
必定會幫你生一個兒子.
你一定要幫他改名叫做以撒.
我就要跟他.
堅定那個約.
那你可以想像得到.
原來阿伯拉漢.
那個時候.
他看到自己的身體如同已死.
即是不可以生的.
但是他的信心不軟弱.
因信不得堅固.
不會起疑惑.
因信心裡就堅固.
更重要就是.
可以說Hallelujah.
敬拜神.
將榮耀歸給神.
神仍然眷顧他們兩老.
神仍然要在他們的家族.
衍生救世主.

$^{2041}$你可以想像得到.
保羅真的很淋漓盡致.
將整個阿伯拉漢.
他在信心的成長.
原來我們可以做一個總結.
15到17章.
原來亞伯蘭.
一直都是信靠神.
一直都是憑信.
得到神的應許.
神的喜悅.
即是我們今天信耶穌的弟兄姊妹.
我們一直在我們的生活.
去經歷神的帶領.
有時帶我們越走越難的.
那個信心的旅程.
但是不要灰心.
好像阿伯拉漢那樣.
神的應許一定會兌現.
這就是我們.
看到這一段經文.
更加告訴我們.
4章17到21節.
第二個重點.
第二個重要的重點.
信心和恩典.
是阿伯拉漢成長的經歷.
一切都是恩典.
亞伯蘭不斷說.
以利以謝.
不如以實馬利.
十三歲了.
神話不可以.
所以其實是恩典的祝福.
信心的成長.
亦都是我們今天.
信耶穌恩信稱義的.
一個很重要的基礎.
我們可以回到第13節.
稍為解釋一下.

$^{2081}$因為神應許.
阿伯拉漢和他的後裔.
必得承受世界不是因律法.
是因信而得的義.
若是屬乎律法的人.
才可以得為後裔.
信就歸於虛空.
即是說做一個分析.
如果只是那些受割禮.
猶太教猶太人.
才可以成為後裔.
那就變成.
只是局限了恩典.
是透過行為.
透過一些猶太教.
他們的傳統才可以得.
神就不是用這個系統.
而是用信心.
信靠他的愛子.
這個系統.
即是說普世的人.
都可以透過主耶穌.
得到神的救恩.
因為如果只是這些人.
信就沒有必要.
因為你行入猶太教.
你受割禮就可以了.
根本就不需要信心.
所以信就歸於虛空.
應許都會廢棄.
但第15節.
律法是惹動憤怒.
即是說如果你要靠自己的好行為.
你永遠達不到神的標準.
永遠不能夠透過律法.
透過自己做好人.
就可以得到神的.
祝福和恩典.
因為律法.
本身就是告訴你你是一個怎樣的人.

$^{2121}$他不可以幫你.
去做到神的要求.
所以這裡有一個.
這樣的解釋.
去到第七章就會再詳細地描繪.
律法的功用.
所以這裡就一個.
基本的分析.
保羅不是想暗示要這樣做.
反而在第17至21節.
保羅就告訴我們.
原來亞伯拉罕.
都是舊約的人物.
但他的成長.
是不斷憑信心.
去經歷神的應許.
經歷神的恩典.
那就是真正.
舊約信心的人.
人生的路.
就不是律法.
所以這裡告訴我們.
第16節也有一個總結.
所以人得為後至.
是本乎信.
因此就屬乎恩.
就是說舊約的人都是.
憑信得到神的.
應許,這就是恩典.
叫應許定然.
歸給一切的後裔.
不單止是屬乎.
律法的後裔,即猶太人.
也歸給效法.
亞伯拉罕的.
信心的軌跡的人.
即是外邦人.
所以第17節就開始說.
亞伯拉罕,我們反過來說.
就讓我們知道.

$^{2161}$亞伯拉罕的人生其實都是.
信心和恩典的偏激.
而成長的一個人.
那就是和我們今天信耶穌的弟兄姊妹.
是完全一樣.
舊約的人都是.
信清而律.
的意思就是這樣.
我們還有一個時間.
9點鐘了,我們可以留待.
下一個星期.
因為21到25節.
我們還沒有詳細解釋.
因為25節是一個很重要的總結句.
就是人類的命運.
必須是要像亞伯拉罕那樣.
來憑著神的.
應許,將自己的.
信心投射到.
耶穌那裡,因為耶穌是為我們.
傳言的,擺上的.
神應許的禮物.
如果我們將我們的信心.
放在主耶穌裡面.
我們就得著.
已經被消除了.
信心的建立就是.
與神的關係的復和.
就是我們今天大概.
將這一段.
一到第二十節.
就很簡單地解釋了.
還有一兩個重點.
都差不多了.
應該都是大概的意思.
我們可以看看第二.
不過25節我們還沒有講.
我們先不背了.
我們下一個星期再看.
我們一個祈禱.

$^{2201}$多謝天父.
今天我們都很快地將第四章的羅馬書.
特別是同創世記的15到17章.
亞伯蘭的成長的旅程.
其實就是一個信心.
經歷神,相信神.
應許的旅程.
保羅要強調的.
亞伯拉罕雖然是舊約的人.
甚至乎他都是後來受割禮.
但是割禮絕對不是一個得救的憑據.
反而是一個信心的.
相信的印證.
他來到去因為信神.
神跟他立約.
得到神的祝福.
多謝你讓我們今天透過主耶穌.
可以跟你有永生的關係.
我們可以隨時去到天上的致聖所.
向你禱告.
向你去求問.
還有你的聖靈的帶領.
很盼望這些經文的講解.
成為很多弟兄姊妹.
可能初初讀羅馬書.
能夠明白.
能夠掌握.
所以我們大家都建立.
我們跟神之間很重要的.
信心的基礎.
還有繼續在信心上.
因信而至於信的成長.
我們感恩禱告.
奉耶穌基督的名.
\newpage



\section{羅馬書}
\label{sec:QTpL9O5GR3E}
\textbf{2025-05-22 周四晚查經班《羅馬書(11):分題:稱義帶來與神復和》 講員:曾裔貴牧師}
\newline
\newline
連結: \href{https://youtube.com/watch?v=QTpL9O5GR3E}{\texttt{https://youtube.com/watch?v=QTpL9O5GR3E}} ~~~~ 語音日期: 2025-05-23
\newline
\newline
\hyperref[sec:BtALIL64hKM]{\small{< < < PREV SERMON < < <}}
~
\hyperref[sec:index_chronic]{\small{[返順時目]}}
~
\hyperref[sec:index_scriptual]{\small{[返順卷目]}}
~
\hyperref[sec:23gUJZhHg2o]{\small{> > > NEXT SERMON > > >}}
\newline
\newline
$^{1}$各位晚安 我們一起有個祈禱.
感恩我們再一次來聚集.
去研讀《羅馬書》.
我們去到第四章的經文.
特別是關乎到亞伯拉罕.
在舊約的聖經人物.
他是怎樣能夠堅固律法.
以及和這個恩信清義有什麼關係.
都盼望我們能夠透過保羅他的展述.
能夠掌握,能夠認識.
以致我們能夠在整個神學的框架.
不會有一種錯誤的觀念.
舊約就是律法.
新約才是恩典.
我們很盼望神給我們認識.
整個聖經的那個救恩的歷史.
和神的救恩的那個計劃.
我們等候在你的面前.
我們會處理到相關的舊約的經文.
盼望神給我們能夠有一個明白的心去認識.
我們這樣感恩禱告.
奉耶穌基督的名.
阿們.
上個星期我們提過.
這個恩信清義.
其實只是一至四章的經文.
當然哪些人需要清義呢?.
就是萬民都欠缺神的榮耀,神的義.
所以其實每一個人都必須來信主.
藉著耶穌才能得到神的救恩.
所以恩信清義是說.
每一個人都需要去到神的面前.
得到這個白白的恩典.
但是在整個的羅馬書.
反而有一個很重要的主題.
既然萬民都需要救恩.
萬民都需要恩信清義.
以及不再在保羅的神學裡面.
其實在這個救恩底下.
根本就沒有分猶太人.

$^{41}$或者是希臘人,外邦人.
那個不同的等級.
不同的階級,不同的民族.
所以在救恩底下.
萬民都是平等的.
所以信了耶穌.
這個就是往後我們就看到.
一方面我們需要向萬民宣教.
但是更加重要就是教會必須合一.
不管是什麼民族.
不管是什麼不同的意見.
都需要在那個救恩底下.
來去謀求那個一切的同心.
重點仍然是要興旺神的福音去到地極.
所以羅馬書就是這個大的主題.
反而不是恩信清義是一個大的主題.
當然每一個在教會裡面的人都是恩信清義.
否則根本不可能在神的教會裡面.
能夠這樣來到有新的生命.
大家彼此的相愛,彼此的同心去服侍.
所以我們第三章稍微有些總結.
就是在這個恩信清義底下.
就是說神的義在律法以外來彰顯出來.
就是恩信耶穌基督就白白的清義.
所以是必須讓人看到不是自己的功勞.
去換取救恩去得到神的義.
所以所有的人在救恩的下面.
都需要謙卑承認自己的無助.
而得到神的救恩是白白的,是無條件的.
我們就可以成為神的兒女.
所以這個救恩帶給我們.
其實第五章就會詳細講.
其實是建立了關係上的恢復.
和神的關係的一個恢復.
就不只是一個宣判無罪.
那個很負面.
只是不再追究我們的罪.
那個是一個很基礎的開始.
但其實是真真正正.
是一個關係上的建立.

$^{81}$所以三章二十九,三十節.
其實不止這裡.
是一直在整卷書都流露出.
萬民都需要合一的.
在救恩底下.
神看我們不再分我們是什麼人.
其實第四章都有的.
所以我們要有一個總結.
就是三章二十九至三十節.
是告訴我們.
我們需要在這位獨一的神面前.
我們必須學習一個合一的功課.
合一的真理.
接著我們就來做這個開始的總結.
之前的最後就是三章的三十一節.
因為保羅突然間帶出了一個很重要的一個命題.
就是在恩順清二.
這個這麼嶄新的一個真理.
在保羅的年代.
你想想是二千年前.
當時猶太教還是在耶路撒冷有憲制.
有不同的節日.
猶太人仍然看自己是神選民的身份等等.
所以到底保羅怎樣去處理.
舊約對每一個猶太人的要求.
或者舊約如果我們看得不清楚.
好像是規範了.
外邦人必須要透過加入猶太教.
才可以得到神的眷顧.
所以突然間保羅傳出一個.
似乎好像撇開了所有的舊約.
就可以透過耶穌.
就可以恩順清二.
進入和神的關係上.
這個對猶太人來說.
尤其是一些拉比,法利賽人.
對他們來說是一個很大的侮辱.
一個的冒犯.
到底保羅是不是一個離經半道.
離開猶太教.

$^{121}$叫人不要守律法的一個人呢?.
所以必須保羅都假設.
或者他都意識到.
他都需要處理猶太人的不滿.
猶太人的善意的誤解.
怎樣來梳理.
以致真是有根有據.
又合乎聖經這樣去推的一個恩順清二.
如果能夠處理得好.
是帶來一個很大的人類的突破.
就不會再分猶太教.
猶太教以外的人.
這是一直都是這樣運作的模式.
所以他就反而自己提出這個問題.
到底這個恩順清二是不是廢掉了律法呢?.
當然不是.
但是恩順清二是怎樣去兼顧律法呢?.
所以你可以說看到在這個三章二十一節.
這個如今神的義.
在律法以外顯明.
到四章的第一節開始.
就好像一個很重要的分水嶺.
保羅都為了讀者.
不是一個好像不清不楚的講述.
而導致信了主之後.
可能還會被猶太教的人逼迫.
或者影響或者誤導.
你們不加入猶太教.
你們不守國禮.
你們就不夠純正.
沒有我們猶太人那麼純正.
就會帶來恩順清二好像要靠一些行為.
所以無論如何.
保羅他都是很有意識地要告訴我們.
其實與舊約不僅沒有抵觸.
原來保羅再重新去演繹舊約的時候.
當然不是一套嶄新的演繹.
而是舊約原來一直都有信心.
恩典這些精髓在其中.
所以第四章就立即有一個很重要的聖經人物.

$^{161}$就是亞伯拉罕.
所以保羅要解釋到底恩順清二是如何兼顧律法.
很自然地就必須要回到舊約.
保羅採取哪一個聖經人物.
是足以讓就算是猶太人或者外邦人.
都能夠心悅誠服地看到.
原來舊約都有恩順清二.
這樣就完全可以很有理性地.
很合理地.
很有意識地解除猶太人的敵視.
所以如果我們能夠在辯論一些真理的時候.
能夠很平心靜氣.
很有真理的基礎.
有神的智慧這樣去解說.
就不只是一句話.
到最後就令到那些懷疑的人.
只會更加懷疑,更加不滿.
但如果我們能夠有重點.
所以你看到第四章是一個很重要的補充.
因為第五章就再一次回到.
那個恩順清二不只是宣判你無罪.
恩順清二是帶來同神關係的一個恢復.
在這個關係上我們怎樣在人生裡.
與神同工和與神同在.
所以第四章是一個很重要的補充.
就是這個命題.
就是怎樣兼顧律法.
怎樣是保持舊約的律法.
原來舊約不單止我們沒有違背.
原來你們看錯了.
以為舊約要你做甚麼.
其實不是.
原來真正在神面前的舊約的聖徒.
他們都是憑信心在神面前的.
這樣就很不同了.
原來恩順清二從來不是站身的.
是一路都是恩順清二的.
這是一個很重要的解說.
所以保羅將亞伯拉罕帶出來.
我們要問這個問題.

$^{201}$初期教會.
修信的羅馬教會.
當時周圍的地方.
歐洲的地方.
信主的弟兄姊妹.
他們很多不是猶太人.
或者可能有些地方是有猶太教的會堂.
但是到底亞伯拉罕.
在這些外邦人的心靈裡面.
是一個陌生人.
還是大家已經有一定的概念.
對亞伯拉罕有認識.
如果不是由頭開始說過.
好像需要費很多的唇舌.
和很多的功夫.
但是《列卷》第四章.
其實是很深入地.
假設那些讀者.
已經對《創世記》十五到十七章.
是已經很瞭如指掌.
我不知道在線的.
現場的.
你們對《創世記》亞伯拉罕的故事.
由他離開吾爾.
到他一直成為信心之父.
你有多少掌握他的年歲.
他的家庭的遭遇.
有多少呢?.
就是說保羅似乎假設.
羅馬教會是很認識亞伯拉罕的人生.
就是說亞伯拉罕不是一個陌生人.
從來都不在外邦人的世界裡面出現.
就是說原來在教會有很多.
關於這些舊約的人物的討論.
這個時候對保羅的講解.
相對容易很多.
當然我們今天沒有證據.
沒有一個人告訴我們.
是亞伯拉罕我們很熟悉的.
我們沒有這些證據.

$^{241}$但是我們只能夠在經文裡面去假設.
保羅的探討其實是很深入的.
很多時候他有他的見地.
一會兒我會和大家一起去讀的時候.
他怎樣去運用《創世記》.
那個經文的猶人有餘的這樣遊走.
還有那個引述都告訴我們.
保羅都是相當透徹了解整個亞伯拉罕的人生.
所以他是非常運用自如的.
我們就感覺到.
透過第四章讓我們看到原來舊約表面的經文.
我們以為很熟悉.
但是原來我們再進一個再深度的認識.
或者有計劃地再重新去研讀.
又會有一些新的聖經舊約的時候.
沒有這樣去反映出來的.
看到保羅無論神給他的智慧.
那個啟示.
還有他對舊約的純屬運用.
都是值得我們去借鏡和學習.
我們就在這裡第三章帶出這個問題.
我們會有一個很好的回應.
因為我們會有一個很好的回應.
我們會有一個很好的回應.
第四節 他說了一個分析的標準.
什麼叫做恩典 什麼叫做靠行為.
什麼叫做你自己賺回來的.
這個就不能夠誇耀.
第四節 作功的得功價不算恩典.
就是是應份的.
我付出了我的勞力.
付出了我的智慧.
給一個公司,一個機構.
或者等等的.
一個團隊裡面的工作.
我拿回我的功價 很合理.
這個交換而已 不是什麼恩典.
不是什麼憐憫.
但是是該得的.
第五節 唯有不作功的.

$^{281}$只信稱罪人為義的神.
就是說沒有靠任何的功德.
沒有靠任何的行為.
其實若有所指.
猶太人仍然存在一個國禮.
是一個符號.
是一個標記.
是一個夢神喜悅的儀式.
一個禮儀.
就是說我們猶太人.
跟你們外邦人不同的地方.
因為我們有國禮.
你們外邦人是沒有的.
所以這個似乎是一個符號.
所以你看到第十二節.
保羅是特意打破了.
這個絕對不是夢神恩典的記號.
這個只不過是亞伯拉罕的一個.
未受國禮的時間已經信的信心.
但是當他受國禮.
只不過是為了那些沒有這樣做而信的人的一個印證.
一個記號.
待會我們第十二節繼續說.
然後他就告訴我們.
第五節這裡.
只信稱罪人為義的神.
他的信就算為義.
正如有一個例子.
大衛稱那在行為以外.
夢神算為義的人是有福的.
他說就是詩篇的三十二篇.
得赦免其過.
遮蓋其罪的.
這人是有福的.
主不算為有罪的.
舊約就是耶和華.
耶和華不算為有罪的.
這人是有福的.
就是大衛一個自述.
在詩篇裡面.

$^{321}$他就說神的認罪.
神赦免他.
他寫這篇詩篇.
很想表達他對神的感恩.
神寬恕他.
所以這個不是作弓.
這個是白白的.
無條件的一個恩典.
所以保羅想要帶進的.
就是這個信心的義.
然後第九節.
如此看來.
一個小小的總結.
他說這個福.
這個恩信稱義的福.
憑信心而得到神的恩典的福氣.
是不是只是單單加給受國禮.
即猶太人呢.
不也是加給那未受國禮的人嗎?.
然後他要解釋.
到底是還是不是呢?.
是加給.
只是加給受國禮的人.
還是加給未受國禮的人呢?.
那就要留意這裡了.
因為保羅他引述的十五章的創世記.
是亞伯蘭當時未受國禮.
就已經信耶和華.
耶和華已經向他顯現.
告訴他.
你的後裔要好像天上的星.
海裡的沙.
這麼多的後裔.
亞伯蘭在那個時候.
他沒有做任何事.
他好像張開了嘴巴.
都啞口無言的.
來接受神的這個應許.
他就相信.
當然神就看到他有這個信心.

$^{361}$就不是只是一種厭口枉.
是他看到他內心的信心.
就稱他這個信心為義.
就是說他已經得著神對他的應許.
雖然還未兌現.
他的小兒子以撒丁還未出生.
十五章的時候是完全未出生.
但是神已經好像看到亞伯蘭.
有這個信心.
而得到神的喜悅.
亞伯蘭繼續說.
是不是只是加給受了國禮的猶太人呢?.
其他那些未受國禮的外邦人就沒有份呢?.
亞伯蘭就引述.
第九節的下半節.
因我們所說.
亞伯拉罕的信就算為他的義.
是怎麼算呢?.
意思就是說在什麼情況下算呢?.
如果保羅在這裡突然間凸顯國禮.
又要將亞伯蘭他未受國禮來信呢?.
這就是一個很重要的論點.
就是讓當時似乎在保羅的年間.
叫人,勉強人受國禮.
加入猶太教.
就是一個很重要的切入點.
就是蒙神的恩典的標記,一個記號.
但是保羅在這裡就好像想淡化國禮.
他不是反對猶太人行國禮.
因為猶太人基於那個傳統.
他們的民族受國禮是沒有問題的.
也不要因為受國禮就變成得天獨厚.
變成有特殊的身份,特別的權利.
那就是錯誤了.
所以變相外邦人也不需要刻意去受國禮.
才可以得到神的信心的恩典.
所以現在他就要強調.
是怎麼算呢?.
亞伯拉罕是怎麼被算為義呢?.
是不是因為國禮他被算為義呢?.

$^{401}$如果是的話就困難了.
就好像有了案例.
後世的人都需要這樣去跟隨.
但是我們看到亞伯拉罕.
他被算為義不是國禮.
而是單單他信耶和華的應許.
所以我們一會兒去到四章二十五節.
保羅就做一個最重要的人類的總結.
就是耶穌出現在這個世界.
耶穌被交付給那些敵人.
就是為了我們的罪惡.
祂的復活就是為了我們能夠清義.
這樣就很清楚了.
等於亞伯拉罕什麼都沒有做.
神現在有一個應許.
你信我 你的子孫就像天上的星一樣.
亞伯拉罕什麼都沒有做.
就接受了做物主的承諾.
這個應許.
他就立即被稱為義了.
所以原來我們的恩順清義.
我們的那個應許.
就是神已經將他的愛子交付了.
將他的愛子復活了.
我們就可以得到這個恩順清義.
所以又是回到神的義在律法以外.
這樣去顯明.
所以三章二十一和四章二十五.
又是一個頭尾的.
神應許了給信的人.
的這個白白永生的福氣.
但是要解釋.
到底恩順清義是怎樣兼顧律法呢?.
所以這個第十節怎麼算呢?.
我們就要來留意了.
是不是在他受割禮的時候呢?.
是在他未受割禮的時候呢?.
到底是還是不是呢?.
答案就是第十節的下半節.
不是在受割禮的時候.

$^{441}$乃是在未受割禮的時候就恩順清義了.
這個是猶太人接不接受呢?.
沒錯 表面上的確是.
但是我們一路再看下去的時候.
因為下面由第十七節到二十一節.
就很詳細將創世記的十七章.
做一個很濃縮的解釋.
阿普羅.
就是說他很想帶我們回到創世記.
回到好像時光隧道那樣.
回到亞伯蘭.
他怎樣經歷神帶領他的人生.
怎樣能夠一路成長呢?.
這個是重要給我們在今天.
當然我們今天就沒有面對猶太人對我們諸多的挑剔.
說我們不受割禮.
我們根本已經過了那些年日.
但是在保羅的年代.
愛邦人信主其實面對的四面八方的嘲笑.
攻擊.
羅馬帝國要拜多神的.
你不拜 你要拜一神.
你要拜耶穌基督.
那是一個攻擊.
但是猶太教又說你不受割禮.
你不守節旗.
你不跟我們一樣.
你又不是我們神的選民.
又變成另外一種嘲諷.
侮辱.
他未必會逼迫你.
但是會對愛邦人信主.
會有一種歧視.
會有一種誤解.
你怎樣分薄我們的福.
你有什麼身份去分我們在神面前的福呢?.
所以這段比較重要.
去拆解.
其實不需要兩個民族有這種對立面.
而是我們回到舊約.

$^{481}$保羅說.
其實連我們這麼有信心的亞伯拉罕.
稱為信心之父的這位人物.
原來他也有他的軟弱.
有他的成長階段.
他得到神的應許的時候.
他根本就未受割禮.
不是刻意要淡化.
或者保羅要扭曲那些經文.
而是根本那些經文的人生的順序.
他根本都未受割禮.
所以割禮不是一個獲取恩典的記號.
這個就很肯定了.
所以拆開一件事.
回來信仰我們的基督教.
如果聖經沒有堅持的那些傳統.
是我們自己慢慢約定俗成的一個傳統.
就不要看為一個絕對的真理.
去灌輸給任何一個時代的神的兒女.
因為聖經沒有應許過的東西.
我們就不要看為好像是祖宗家法不可以改.
可以是一個傳統.
可以大家都覺得這樣做會比較舒服.
不等於是絕對.
現在問題就是猶太人就看割禮是絕對.
你不受割禮你就不可以得救.
我們有割禮的記號.
我們就是蒙神的恩典和喜悅.
保羅就等於很間接很委婉地去拆解了.
從來割禮都不是一個記號.
是記號不過不是一個唯一無福的記號.
真正無福的記號就是信心在神的應許上面.
所以這裡告訴我們第十一節.
並且這裡就說記號了.
很快我想去下一個部分.
因為我將英文都帶出來.
這裡就是第十一節.
是很重要的經文.
我先看英文.
這裡告訴我們.

$^{521}$他獲得割禮的記號.
原來是一個銀.
英文NASB是一個銀.
是一個印記.
埃及法老約瑟未從監獄出來.
解釋了七個方年.
快要來了.
現在是七個豐年.
現在要去採米去囤積田地.
到豐年來的時候.
就可以開倉派米.
不然大家都會抱著死.
法老聽完之後.
你解了我的夢.
然後就將他手的戒指給了約瑟.
這個戒指就是古代的人.
一個好像文書蓋印.
所以是一個銷.
一個印記.
英文NASB都是用銷.
不是用銷這麼簡單.
所以這裡的中文沒有英文說得那麼好.
他說這個割禮的印記.
原來就像一個錄印.
一個標記.
一個印鑑.
是一個有信心的人的印記.
就是在他未受割禮的時候.
我highlight了.
Uncircumcised.
未受割禮的時候有這個印記.
他就可以成為了.
Father of all who believe without being circumcised.
就是說他能夠成為未受割禮的人的一個指標示範的單位.
我們看回11字的經文.
並且他受了割禮的記號.
作他未受割禮的時候.
因順稱義.
果然就是印證這個字.
印證這個字就是那個字.

$^{561}$就是那個好像一些印鑑.
好像我也說過在《異筆所書》.
受了一個印記.
就是古代他們要買大批的木材.
那些買家或者買手.
就會去到那個木材的地方.
就看到這批木材我要了.
你什麼時候就斬好之後.
通常就扔到河裡浮下來.
免及那個運輸.
因為那些牧會浮.
那他們怎樣呢?.
這批貨是屬於我的.
我現在就用我的公司.
以前就是那些木材.
用一個印.
好像我們做月餅.
那些印.
大同月餅.
恆鄉月餅.
就蓋上印.
就證明這批貨我將來就會要了.
就是那個印記.
所以這裡呢.
保羅告訴我們.
亞伯拉罕他未受割禮的時候.
就已經因順稱義的一個印證.
一個銷.
就已經因順稱義了.
他就能夠作為一切未受割禮.
而只是相信神的人.
他就是他們的父.
就是說他是所有信神的人的父.
就是說他自己第一個信.
神就好像在他的生命裡面.
給一個印記他.
一個印證.
你就是一個好像是示範單位.
就成為日後都是這樣.
憑信心相信神的應許的人.

$^{601}$不管那些人是猶太人還是外邦人.
都可以得著神的應許.
這個就是保羅想要去解釋的.
所以我們回到這裡告訴我們.
其實舊約裡面.
我們的這一位亞伯拉罕.
他能夠得到神的那個救恩.
同樣都是憑信心.
不是憑律法.
所以保羅在這裡的解釋.
其實是很重要.
接著他不單只是說未受割禮的人.
實際上他們都算為異.
留意一下第十二節.
保羅其實真的想說萬民都是靠恩典得救.
在保羅的救恩裡面.
這個福音其實是給任何一個.
開放自己心靈的人就可以得的.
絕對沒有你的背景.
你的種族.
你的文化.
可能有一些幫助.
但不是先決的條件.
先決條件是神看我們有沒有那個信心.
所以十二節.
「又作受割禮之人的父.
就是那些不但受割禮」.
保羅多加一句.
受了割禮而你沒有亞伯拉罕的信心.
你只是一個猶太教的人.
你沒有在神面前有真正真實的信心.
你都不可以因為你是猶太教.
受割禮就可以得到神的應許.
因為保羅真的很有智慧.
並且你受了割禮.
但你都要按我們的祖宗亞伯拉罕.
他未受割禮而信的信仰的宗旨去行.
你才稱為有信心的人.
他講的是猶太教的割禮分子.
你可以想像到.

$^{641}$保羅的救恩原來涵蓋的.
原來神的愛,神的救恩涵蓋的.
就是就算你沒有受割禮.
從來都不知道什麼是割禮.
但你知道有耶穌的這個應許.
你相信你就可以得救了.
很多穆斯林的群體.
如果他們根本沒有聽過其他的東西.
他只是憑他的內心相信耶穌.
他就可以得到永遠的救恩.
但猶太人就要多一件事.
他如果只是從小長大.
知道有猶太教.
照足那些規矩,禮儀去生活.
但他一直都沒有信心.
這樣他都不行.
他要根據好像亞伯拉罕在他人生.
未受割禮的時候.
那個信心的生活.
這樣去依靠神.
是那個宗旨.
所以17節就開始講亞伯蘭的宗旨.
他的信仰宗旨是怎樣的呢?.
作者保羅就列舉了好幾樣東西.
尤其是他憑信心去等候以實的出現.
所以來到這裡一到第十二節.
首先第一個很重要的論點.
就是其實舊約.
當然以亞伯拉罕做的單位.
已經是恩信稱義.
我們現在就要回到創世記.
和大家看回那個歷史.
怎樣去組織保羅所講的經文.
那個信心之旅.
就是亞伯拉罕.
或者亞伯蘭的信心之旅.
我們從第十五章可以看起.
十四章就是講到那個一場戰爭.
因為羅德他的侄子.
被這個所多瑪王去捉拿了.

$^{681}$當然當時是五王.
四王和五王的戰爭.
所以這場戰爭.
所多瑪王就應該大敗了.
後來亞伯蘭就去召集他自己318個壯丁.
去打贏了那場仗.
也拿回了那些財物.
輾轉就出現了一個叫麥基駛德的祭司.
為亞伯蘭祝福.
輾轉第十五章.
這些事以後.
作者告訴我們又過了一段時間.
亞伯蘭那個時候還是住在萬里匠樹的地方.
然後告訴我們.
耶和華在異象中有話對亞伯蘭說.
亞伯蘭你不要懼怕.
可能他怕人尋仇.
因為他打贏了那些四王.
然後你不要懼怕.
我是你的盾牌.
必大大的賞賜你.
亞伯蘭說.
主耶和華.
我記無子.
你還賜我什麼呢?.
並且要承受我家業的大馬色人.
以利以謝.
這是他的僕人.
第三節.
亞伯蘭又說.
你沒有給我兒子.
那生在我家中的人.
就是我的後子.
很清楚跟神有這樣的對話.
有這樣的交換.
不用那麼麻煩.
我的僕人就當是我的修養的繼承人.
後子就是繼承人的意思.
第四節.
耶和華又有話對他說.

$^{721}$這人.
就是以利以謝.
這人必不成為你的後子.
你本身所生的.
本身就是由他自己.
就是跟太太薩萊所生的.
才可以成為他的後子.
才成為你的後子.
第五節.
留意一下.
於是就領他走到外邊.
怎麼領呢?.
是一種榮光.
就是神在救約的時候.
神怎麼顯現呢?.
十七章就有兩個天使.
怎麼來向亞伯蘭顯現.
是不是牽著他的手走出去呢?.
怎麼跟他說話呢?.
總之亞伯蘭是可以得到這些信息的.
不是報夢的.
真的走出去的.
日光日白.
所以神可以在日光日白向他顯現.
然後就叫他看.
數算眾星.
日間應該看不到眾星的.
能數得過來嗎?.
可能就是晚上也不頂了.
亞伯蘭就.
留意第五節.
第五節還沒讀完.
又對他說.
你的後裔將要如此.
提出了一個這麼大的應許.
不是順口開河.
是神跟亞伯蘭說.
而且根據上下文.
肯定不是告訴你.
所有的星空星球都是你的.

$^{761}$對亞伯蘭來說.
根本就沒去過外太空.
根本就不知道.
連太空館都沒去過.
根本不知道有多遙遠.
因為古代的人.
怎麼可能有這樣的技術呢?.
但是.
看到眾星就肯定比我們多.
因為沒有那麼多污染.
在晚上的星空.
是非常壯觀壯麗.
但是.
很清楚的一個offer.
肯定這些星是跟後裔有關.
因為上文是說.
你不能生.
你所生的.
才是你的後裔繼承人.
亞伯蘭就說.
不要那麼煩了.
我的終生僕人.
我pass away之後.
就全部給他就行了.
他是我繼承人就行了.
都是屬於我的.
全部給他就行了.
神就說不行.
必須是要你新所生的.
才是後裔.
接著就真的帶他去看戶外.
去觀星.
告訴他.
你的後裔將要如此.
就是一個這樣的應許.
聖經就在這裡.
留意第六節.
保羅說是亞伯拉罕信耶和華.
不是亞伯蘭信耶和華.
好像保羅引錯聖經.

$^{801}$其實不是.
保羅已經在他一直寫的時候.
就已經認定.
亞伯拉罕已經是信心之父的人物.
保羅他是後世的人.
當然知道亞伯蘭一直的成長.
過了十幾年之後.
他都是這樣跌跌碰碰的.
經歷這個信心的考驗.
輾轉他就真的神幫改名.
保羅好像已經立即幫他改名.
保羅就說.
亞伯拉罕信神.
耶和華就以此為他的兒.
所以你可以看到保羅那種.
這麼巧妙地.
來將他自己的神學.
灌注在這些經文裡面.
但是在這個時候.
在十五章的時候.
亞伯蘭是還沒改名的.
要十三年之後才改名的.
輾轉神就跟他說.
我要跟你立這個約.
你是我將這個地賜給你為業的.
然後就告訴他.
你要怎樣去在這個約的儀式.
怎樣去做.
就教他要將一隻三年的母牛.
一隻三年的母山羊.
一隻三年的公綿羊.
單溝,初甲就拿來準備獻祭.
所以每樣都劈開.
分成兩半.
就是死了的.
一半對著另外一半擺.
就是這樣平擺.
就是說準備開壇作法的意思.
就是獻祭.
但是只有鳥沒有擺開.

$^{841}$只是那些牛羊,公綿羊就擺開了.
就是劈開了.
那些單溝就整隻了.
那些初甲也是整隻了.
沒有劈開.
然後就是十一節.
有飼鳥下來.
落在那死畜的肉上.
就是說過了一段時間.
亞伯蘭就把牠嚇飛了.
就是跑過來.
那些蒼蠅拍開牠們.
因為牠跑過來.
想吃那些獻祭的死畜.
然後就是十二節.
日頭正落的時候.
就是黃昏.
亞伯蘭因為太久了.
沉沉的睡了.
不過留意.
忽然之間有驚人的大黑暗.
落在他身上.
就是說神降臨.
耶和華就對亞伯蘭說.
你要的確知道.
你的後裔必寄居別人的地.
又服侍那地.
當然是埃及.
那地的人.
那地的人要輔代他們四百年.
原來以色列人要寄居在埃及四百年.
其實是做奴隸.
因為要被輔代.
已經預言了.
說了.
並且他們所要服侍的那國.
我要懲罰.
後來他們.
就是以色列人.
就是你的後裔.

$^{881}$必帶著許多財物.
從那裡出來.
但你要享大壽數.
127.
亞伯蘭說.
平平安安歸到你列祖那裡.
他的列祖就是.
亞伯拉罕.
以洛.
以羅氏.
馬土撒拉.
這些先做.
羅亞.
因為亞伯拉罕是第一個被神.
由吳爾帶離開.
去到迦林地.
然後你留意一下.
歸到列祖那裡.
被人埋葬.
到了第四代.
就是你的第四代.
亞伯拉罕的第四代就是.
亞伯拉罕以撒.
瓦割.
然後就是約瑟勒.
第四代.
過了第四代.
第五代就會出.
所以他們很快就要.
不過留意這裡.
他說你們就會歸回到這個地方.
就是迦林地.
為什麼呢.
原來因為阿摩利人的罪孽.
還沒有滿盈.
就是說原來神的計劃是很雙向的.
原來他們就被老.
不是被老.
他們自己下去埃及.
去吃米.

$^{921}$後來就在埃及那裡一直生活.
因為祖宗約瑟得到王的喜悅.
所以整個民族下了埃及.
然後就一直都回不來了.
然後原來最終的原因.
是我們都沒有想過.
神現在就很早告訴亞伯拉罕.
因為阿摩利人的罪孽還沒有滿盈.
就是說他們為什麼回不來迦林.
原來迦林地的原居民.
還沒有最大惡極.
還沒有得到神最終的審判.
所以以色列人還在等待四代.
在埃及地.
這些就是超出了我們對歷史的新的了解.
原來神已經很早告訴亞伯拉罕.
你的後裔將來會這樣這樣.
將來那個世界會這樣.
出現這樣的情況.
很多事情表面到時就會自然的演變.
但原來背後.
原來神是正在積蓄阿摩利人的罪孽.
當地迦林人原居民他們拜偶像.
他們是很得罪神.
到了一個真是滿盈了.
神就要刑罰了.
神怎麼刑罰呢.
原來就透過亞伯拉罕的後裔.
約瑟的子孫.
他們回到埃及.
當然摩西帶領他們的民族離開埃及來到迦林地.
或者約瑟帶領他們回到迦林地.
原來已經在四百年前.
神已經跟亞伯拉罕說了.
但現在這個重點就是.
這個應許就是亞伯拉罕還沒有改名.
中國人就很重要了.
改名就是轉運.
我姐姐都改名了.
但是亞伯蘭還沒有改名.

$^{961}$這個時候他還是一個很普通的人.
但是到了十七章.
神的應許時間到了.
他要生一個兒子就是以撒.
神就說我要將你這個舊的名.
亞伯蘭改為亞伯拉罕.
這個名不是因為轉運.
是因為神的應許時間到了.
你應該具備信心之父這個名字.
但是在十五章完全沒有這個信心之父的名字.
但是阿保羅已經用亞伯拉罕來定性了.
亞伯蘭其實已經因為他的信心.
來得到神的應許.
不是因為國禮.
因為十七章才行國禮.
整個希伯來人才行國禮.
所以如果我們熟悉歷史.
我們就知道阿保羅為什麼這麼振振有詞.
恩信稱義其實是肩負律法.
不是否定律法.
更加不是將律法架空.
原來舊約都已經有恩信稱義了.
舊約的亞伯蘭.
都是因為他對神的應許的信心.
其實那個信心有多大呢?.
我們都不知道.
但是神就看到已經足夠了.
就稱亞伯拉罕為義了.
當然我們覺得立約的儀式是很恐怖的.
第十七節.
日落天黑.
不料神的出現很特別.
有無煙的爐.
並燒著的火把.
從那些玉塊中經過.
神這樣經過讓亞伯蘭看到.
當那日耶和華與亞伯蘭立約說.
我已經賜給你的後裔.
從埃及河直到伯拉大河之地.
就是基尼人,基尼洗人,甲摩尼人,.

$^{1001}$黑人,比利洗人,尼弗欽人,.
阿摩尼人,迦南人,甲加薩人,.
耶布斯人之地.
整個迦南都是亞伯拉罕的.
神與他納了約.
在這個時候.
這個儀式與他立約就是當那日.
納了這個約.
這個民族就是這樣蒙福.
是因為亞伯蘭的信心.
因為亞伯蘭信耶和華.
這樣來得到.
接著十六章就回到日常生活.
十六章憲制完立了約.
亞伯蘭的妻子撒萊不給他生兒女.
撒萊有一個侍女名叫夏甲.
是埃及人.
這個時候夏甲原來當時在十一章.
亞伯蘭因為饑荒落了埃及.
在埃及怕他自己的太太撒萊.
被法老王抓了她做妃嬪.
就跟撒萊說.
現在我們去埃及了.
我也不知道當地的王是怎樣的.
你就當是我的妹妹.
其實也可能是同輩異母等等.
某程度上不是絕對的謊言.
但也是一個欺騙.
果然法老見到撒萊真的想娶她為妻.
果然把她不是抓的.
而是用很多聘禮就拿了.
後來神就向法老報夢.
你娶撒萊你想死.
你快點歸還給亞伯蘭.
否則你的家族也有事.
法老就很害怕就叫亞伯蘭來.
亞伯蘭就說我很害怕這個地方的人.
所以我就認我的太太做妹妹.
她就變成了一個大舅.
是不是大舅啊?.

$^{1041}$哥哥嘛 大舅.
就拿了很多的聘禮.
不是聘禮 是禮物.
送走亞伯蘭和撒萊.
下鴿是其中一個禮物.
就是帶她上去迦勒地.
所以現在就下鴿.
那明顯是埃及人.
埃及妖后那些埃及人.
撒萊對亞伯蘭說.
耶和華使我不能夠生育.
求你和我的侍女同房.
或者我可以因為她而得到孩子.
亞伯蘭就聽從了撒萊的話.
於是亞伯蘭的妻子撒萊.
張侍女 埃及人下鴿.
給了丈夫為妾.
那時亞伯蘭在迦勒已經住了十年了.
就是八十幾歲.
八十五歲左右.
亞伯蘭與下鴿同房.
下鴿就懷了孕.
那就是說原來不關亞伯蘭的事.
生不了就是撒萊不能生.
就不是亞伯蘭不能生.
現在下鴿一同房就能生了.
就懷孕了.
他見自己有孕就小看他的主母.
是不是可能真是因自己的兒子.
將來繼承的兒子呢?.
都不知道 都沒有去做DNA化驗.
也沒有照超聲波 怎知道是兒子呢?.
可能是女兒來的呢?都不知道.
但是他已經開始.
因為真是自己身懷六甲.
就開始對主母撒萊.
開始有些不同的眼光.
第五節 撒萊對亞伯蘭說.
我因你受屈.
我將我的侍女放在你的懷中.

$^{1081}$他見自己有了孕就小看我了.
怨謠說在你我中間判斷.
亞伯蘭對撒萊說.
侍女在你手下.
你可以隨意待她.
撒萊輔待她.
她就從撒萊面前逃走.
我們可以想像到.
撒萊怎樣去輔待夏甲呢?.
古代的人怎樣去輔待那些妾呢?.
是用言語呢?.
還是有具體的行動呢?.
關她去柴房不讓她吃東西呢?.
還是怎樣呢?.
要輔到那個地步.
他要從主母的家逃跑.
然後跑到曠野.
因為耶和的使者.
在曠野的舒爾的路上.
似乎想回外家.
因為回到埃及.
一直在曠野.
就是輔到一個地步.
一個有了身孕的婦人.
痛苦到要逃跑.
就是說相當痛苦.
因為其實古代的人.
尤其是你是一個下人.
一個婢女.
做一個奴隸.
被人送走的一個妹子.
其實你根本就沒有社會地位.
根本就沒有人的尊嚴.
但是她都要苦到一個地步.
寧願逃跑.
你可以想像到.
撒萊原來都不是善男信女.
撒萊都有些手段.
對付夏甲.
好像都綽綽有餘.

$^{1121}$聖經就沒有怎樣描述.
她有多苦.
還有那些手段是怎樣的惡毒.
都不理她又身懷六甲.
還有肚子裡的嬰兒是亞伯蘭.
不是其他的僕人.
你可以想像到.
有時候那種失控.
都不理那麼多.
撒萊就固然這裡.
應該是要受到神的譴責.
但是某程度上.
神又沒有反對.
因為夏甲所生的.
都不可以繼承亞伯蘭的產業.
還要是撒萊自己肚子裡生的.
才可以繼承神的產業.
所以好像默許做這件事.
我們一直不講那麼詳細.
然後我們突然間跳到第十五節.
就是講夏甲聽了天使的話.
就是她太辛苦太受委屈.
她走而已.
其實她沒有全盤計劃.
沒有收拾乾糧,指南針,手提電話.
回去埃及.
因為一個女流之輩還大了肚子.
你怎樣能夠定得順在曠野的生活呢?.
所以她只是想吐口苦水.
突然間迷了路.
天使就向她顯現.
告訴她要回去.
要伏在她的主母的腳下.
就是說你回去.
主母的氣都過了.
都會讓你生的.
還有你的兒子某程度上.
因為你,因為亞伯蘭.
他都是無福的.
這個英雄天使這樣安慰.

$^{1161}$所以天使都很重要.
都救了兩個生命.
因為如果不是.
夏甲可能真的被野獸擔心了.
那就沒有了.
但是你可以想像到.
有時候那口氣都需要有一個空間.
讓夏甲她自己又可以明白.
是喔.
非理性地逃跑都是不行的.
回到家裡.
可能就是改變她的態度.
你不要再小看主母了.
其實你要自己.
麻煩你照鏡子.
你是小妹而已.
你不是因為你的兒子.
你就水鬼升城王這樣.
這樣不行的.
她開始明白.
降低她的那種氣焰.
當然主母殺來.
就知道原來自己都過分了.
就容許.
既然都是丈夫的一個骨肉.
就容許這個妾侍.
都在家裡這樣生了這個二十晚.
沒錯 生了這個二十晚.
十五節.
後來夏甲給亞伯蘭生了一個兒子.
亞伯蘭給他起名叫二十晚.
夏甲給亞伯蘭生二十晚的時候.
亞伯蘭年八十六歲.
那就是說的對.
八十五歲.
夏甲就被跟亞伯蘭同房.
是阿實來的一個獻祭.
但是這就告訴我們.
當時這件事叫做一個小小的風波.
暫時可以完滿的解決.

$^{1201}$那個大妻子終於可以順產二十晚.
這個時候亞伯蘭已經是八十六歲.
不過留意就是十七章.
有十三至十四年.
一個真空期.
在一節聖經裡面.
作者如果真的是摩西.
好像那十三年十四年.
是完全沒有任何的記錄.
讓我們有個空間去聯想.
二十晚一直長大.
在亞伯蘭的家裡有兩個男士.
當然不算亞伯蘭.
有一個叫做以利爾謝.
那個老的僕人.
那個已經說明不是你的繼承人.
不行的.
後來獻祭.
由我的侍女所生的歸入我的名下.
亞伯蘭說好啊.
因為真的由我所生的.
是我和夏甲所生的.
或者我和任何一個女士所生的.
都屬於我的骨肉.
可以是神的應許的產業.
他可以繼承.
但是.
以實瑪利出世.
由嬰兒長大.
然後滿月走.
然後八日宴.
然後一直成長.
越來越大了.
十三歲都不小了.
在這個時候.
以實瑪利是十三歲了.
但是聖經是完全沒有神的.
再介入.
再顯現.
也沒有告訴我們.

$^{1241}$亞伯蘭每天弄而為樂.
撒萊很滿足.
夏甲是兩個.
是不是叫姐妹.
相安無事.
大家快樂融洽.
那些雨場面的大結局.
亞伯蘭就成為齊人之福.
又有兩個老婆.
又有一個兒子.
就可以繼承產業了.
十三年就這樣過了.
突然之間.
十七章.
亞伯蘭年九十九歲的時候.
顯然是晴天霹靂.
告訴我們.
神又有些新的東西.
要亞伯蘭去遵守.
就是說.
那十三年的天倫之樂是廢的.
不是天倫之樂不對.
而是應許的兒子.
你想多了.
以實瑪利不是繼承人.
必須要撒萊.
你可以想像到.
其實亞伯蘭不斷經歷.
越來越深的信心功課.
越來越不容易.
因為十三年已經定了.
兒子已經牛高馬大了.
已經和他一起工作了.
但是突然之間.
晴天霹靂.
神就再顯現了.
第一節開始了.
亞伯蘭年九十九歲的時候.
耶和華向他顯現.
對他說.

$^{1281}$我是全能的神.
你當在我面前.
作完全人.
我就與你立約.
就納了.
十五章就納了.
說應許了我的子孫.
好像天上的星一樣多.
我相信了.
原來那個信心的考驗.
越來越深入了.
所以信主越來越久.
我們的信心應該越來越深入.
越來越堅固.
作完全人.
我要與你立約.
使你的後裔極其繁多.
沒錯.
其實瑪利就是那個繼承人.
第三節.
亞伯蘭苦服在地.
神又對他說.
我與你立約.
你要作多國的父.
從此以後.
你的名字不要再叫亞伯蘭了.
要叫亞伯拉罕.
就是多國之父了.
因為我以納你作多國之父.
一直都沒有碰到以實瑪利.
以為可以過關.
我必使你的後裔極其繁多.
國度從你而納.
君王從你而出.
這段就是.
加泰書三章十九節.
那個應許的兒子就是耶穌了.
就是這一段聖經的預言.
我要與你並你的世世代代的後裔.
建立我的約.

$^{1321}$作永遠的約.
是要作你和你後裔的神.
我要將你現在寄居的迦南地.
全地.
賜給你和你的後裔.
永遠為業.
我也必作他們的神.
很重要.
不只是那個地圖.
而是那個地圖之上的.
這一位的眾日主.
神.
很重要.
所以如果我們今天說.
這個教會.
不單只是一個物業.
還有神的同在.
我應許每一次你們的聚會.
我會同在.
這樣才是最重要的.
不只是那個硬件.
而是神成為他們的神.
第九節.
神又對亞伯拉罕說.
你和你的後裔.
世世代代遵守我的約.
你們所有的男子.
都要受國禮.
在這個時候.
過了十三年之後.
才正式頒布這個國禮的約.
這就是我與你.
並你的後裔.
所立的約的記號.
是你們所當遵守的.
你們都要受這個國禮.
這是我與你們立約的證據.
也是那個紀號.
所以你看到保羅.
他引述經文的時候.

$^{1361}$他是非常充分了解了.
整個亞伯拉罕的成長.
信心的歷程.
那個軌跡.
所以他去講解.
因順稱義在舊約.
原來在亞伯拉罕.
他還沒有受國禮之前.
這麼久之前.
已經經歷了神給他的應許.
而他相信.
就可以得到這個證據了.
所以第十二節.
你們世世代代的男子.
無論家裡生的.
是你們後裔之外.
用銀紙在外人買的.
即是那些奴隸.
生下來第八天.
都要受國禮.
你家裡生的.
以及你用銀紙買的.
必須受國禮.
這些是神頒布的命令.
是一個條款.
但是原來神的這些頒布.
是沒有抵觸到他第十五章的時候.
同樣都是這位做民主.
來去應許亞伯蘭.
同樣都是那個受益人.
就是亞伯蘭.
就是亞伯拉罕.
所以.
因順稱義在舊約已經有了.
不是行國禮.
才可以得到這個.
因順稱義的福氣.
恩典是源於相信神的應許.
所以亞伯蘭就是相信神的應許.
我們今日後世.

$^{1401}$我們是屬靈的亞伯拉罕的子孫.
我們都是因順相信神的應許.
就是他的愛子被交付人.
來為我們的罪.
他復活是為了我們能夠稱義.
在神面前來建立神人的和好關係.
所以你可以想像到.
保羅這個論點是非常清晰.
印證舊約的日.
同樣都是因順稱義.
這個就是為什麼.
因順稱義不是廢棄律法.
是更加堅固律法.
所以當代的猶太人.
仍然堅持受國禮.
才是無福的證據.
是錯誤的.
猶太人可以受國禮.
但是國禮不是蒙神的祝福的條件.
必須是要.
因為猶太人有這個傳統.
受國禮是沒有問題.
但是你都要根據亞伯拉罕的信心軌跡.
去相信神.
你才能夠無福.
所以猶太人受國禮.
是不會有特別的優待.
外邦人不受國禮.
更加不會是一個欠缺.
是一個很重要的.
將兩個民族拉平.
在恩典底下.
大家都是透過恩典去蒙神的祝福.
所以有時我們都要不斷迅速來.
其實要苦心自問.
我們慢慢因為信心.
因為恩典而改變了.
我們不要慢慢就變成有一種優越感.
看不起那些未信的人.
他們的行為.

$^{1441}$他們的想法.
我們慢慢就有一種.
癡之以備.
或者不知不覺.
我們自己有一種優越感.
就忘記了.
其實我們都是蒙恩.
在恩典裡面.
神改變我們.
神賜福給我們.
我們能夠家庭幸福.
各方面都能夠順暢.
我們走的人生路的軌跡.
又能夠走在那個二路上.
我們要記得.
我們要回復那個初心.
未信主之前.
我們都是一樣這樣抵擋神.
在神的憤怒底下.
那些招致神的憤怒的人.
我們就會有一份胸襟.
去接待,接納一些未得救恩的人.
他們那些錯誤的行徑.
例如說髒話.
例如一些不好的行為.
當然不是縱容.
但是我們要明白.
他未有恩典在他們的生命裡面.
自然就會有這些種種的錯誤.
對神的那種抵擋或者嘲笑.
所以這裡告訴我們.
留意下第.
我們下面這裡不到.
十五節.
神又對阿伯拉函說.
你的妻子撒萊.
不可以再叫撒萊了.
神要叫撒萊.
我必賜福給她.
也要使你從她得一個兒子.

$^{1481}$我要賜福給她.
她也要作多國之母.
就是撒萊的意思.
必有百姓的君王從她而出.
再一次去肯定.
以撒的後裔將來坐在皇座上.
就是我們的救主耶穌基督.
第十七節.
阿伯拉函就俯伏在地.
起笑心裡說.
一百歲的人還能得孩子嗎?.
撒萊已經九十歲.
大十年了.
還能夠生養嗎?.
阿伯拉函對神說.
但願以實馬利.
不是以利以謝.
以利以謝已經肯定不是了.
已經是飛了出去局外了.
因為他是僕人.
十五章還在說.
但願以利以謝在你面前無福.
現在就不是了.
現在是她的大兒子.
以實馬利了.
但願.
那是什麼意思呢?.
就是信心不夠.
就是說不要這麼複雜了.
這個兒子已經萬事足夠了.
撒萊其實不應該再生了.
古代就沒有那些高齡產婦的概念.
神也好像有這種玩笑.
九十歲了.
真的九十歲了.
實際上你真的不應該再生了.
現在是一個反建議.
阿伯拉函告訴神.
但願以實馬利也足夠了.
在你面前活.

$^{1521}$神就說不是了.
你妻子撒萊要給你生一個兒子.
你要幫她改名叫以撒.
我要與她堅定所立和你的約.
就是你要傳遞你的恩典和約給你的兒子.
哪個兒子呢?.
以撒.
未出生的以撒.
堅定這個約.
就是說意味著神再一次這麼說.
再一次確認這個約的時候.
就是雙重確認.
意味著阿伯拉函信神這件事.
不過這個信是戰戰兢兢.
因為難度高了很多.
因為現在明明有個大兒子.
難度就高很多了.
太太就已經.
你十年前告訴我.
現在你九十歲叫她生.
我都不好意思了.
但是神就說.
這個就是考驗你的信心.
要再進一層級.
所以我們不要再是初信的時候的見證.
我們要不斷有一個新的經歷.
在我們生命裡面的帶領.
在這個時候.
你可以想像得到.
第二十節.
至於以實瑪利.
我也應允彌.
我必賜福給他.
使他昌盛的其繁多.
他必生十二個族長.
我也要使他成為大國.
這個是很大的安慰.
作為一個父親.
既然你帶他來這個世界.
那你想怎樣呢?.

$^{1561}$你說他沒有繼承權.
那你也要給他一個好的下場.
好的晚年.
神就說我應許.
他 以實瑪利.
因為都是你的兒子.
我都應許他成為一個大國.
他就會甚麼呢?.
留意第21節.
到明年這時節.
撒拉必給你生耳實.
我要與他堅定所立的約.
神和亞伯拉罕說完了話.
就離開他上升去了.
即是說暗示.
作者暗示了一個很重要的訊息.
即是15章的第6節.
亞伯蘭信耶和華.
耶和華就以此為他的義.
即是說17章的第23節.
我們應該要假設.
亞伯拉罕在這個時候.
神說了兩次.
我要與你未出世.
都未行房的兒子.
明年才出現的.
神就說已經說了兩次.
我要與這個未出世的兒子.
建立他與他的約.
所以第23節和第22節中間.
我們應該要假設.
或者應該要肯定.
這句說話就是.
亞伯拉罕再一次信耶和華.
耶和華就以此為他的義.
大家明不明白這個邏輯.
就是這個很重要.
亞伯拉罕在這個時候.
是另一種信心對神的信靠.
很不容易的.

$^{1601}$因為你現在好像在一個幹線上.
我就100歲.
太太就90歲.
但現在是要我們兩個老人家.
生的兒子才可以繼承.
但神已經說了兩次.
這個人就叫做以撒.
明年就會出現的.
其實是多有趣.
多令人憧憬.
亞伯拉罕知道之後.
第23節證明了我的假設是真的.
因為上文說了一個命令.
就是神向亞伯拉罕顯現之後.
就幫他改名字.
之後他自己就以後叫亞伯拉罕.
然後有一個命令就是.
你和你所有的男丁.
要立即行割禮.
第23節正當那一天.
就是神走了之後.
離開他上升之後.
正當那一天.
亞伯拉罕就遵著神的命令.
給他的兒子以實瑪利.
和家裡的一切男子.
無論是在家裡.
很清楚沒有計錯數.
就是所有的男子.
總之跟他有關的.
用銀子買的.
都行割禮.
第24節.
亞伯拉罕受割禮的時候.
年99歲.
也是呼應17章所說的.
神在這裡向他顯現的時候.
他的兒子以實瑪利受割禮的時候.
就是13歲.
接著第26節.

$^{1641}$正當那一天.
亞伯拉罕和他兒子以實瑪利.
一同受了割禮.
家裡所有的人.
無論是在家裡生的.
是用銀子從外人買的.
也都一同受了割禮.
要再加多一個很重要的訊息.
在這個時候.
亞伯拉罕其實他是不是再一次.
堅定他的信.
來相信神會在我這個殘區.
我的太太的殘區.
是不是應該成立呢?.
今天的90歲和100歲.
和亞伯拉罕的年代.
100歲有沒有不同呢?.
我們又沒有什麼證據不知道.
他自己也說自己是100歲.
怎麼可以生呢?.
大概不會相差很遠.
你可以想像到.
我們肯定可以加這一句.
在這個時候.
既然他遵守神的命令去行割禮.
也就是說他都應該根據.
神已經兩次說.
明年這個時候.
10月懷胎.
你的太太就會生一個叫做耳塞.
這個耳塞才是真正的繼承人.
這樣就有一個很大的神蹟.
亞伯拉罕正在期待.
所以其實這裡就是暗示.
雖然作者沒有說.
亞伯拉罕信神會給他一個兒子.
大家有沒有留意.
第四章17節到21節.
保羅就很平鋪直敘.
把亞伯拉罕的心路歷程.

$^{1681}$這一刻的心路歷程就說出來了.
我們回到第四章的17節.
所以剛才17章的創世記.
22,23節是應該加一個.
我們自己加下去也沒有錯.
就是亞伯拉罕在這個時候.
再一次要俯伏相信神的應許.
他明年會有這個兒子.
這個不是天方夜譚.
不是自己兩腦在無事生事.
而是真的從神的應許.
而且還是兩次說.
我會跟你的兒子堅定約的.
就是說這個兒子一定會出現.
如果不出現就很大件事了.
神好像也收不到科.
肯定不是.
所以我們透過這個經文的組織.
我們看看第17節開始.
很奇妙的.
保萊告訴我們.
這段經文真的這樣發生.
第17節第4章.
第13節我不看了.
我們看第17節.
亞伯拉罕所相信的.
是叫死人復活.
使母變為有的神.
證明他自己當時已經很清楚.
他們兩個是不可以生的.
撒萊一直都不可以生的.
所以他看自己是一個死人.
在生育上.
所以他在這個時候.
他的信心原來可以成為.
一個新約的演繹.
就是叫死人復活.
使母變為有的神.
就是剛才說的22,23節的心路歷程.
他內心的那個.

$^{1721}$你這樣說我一定要相信.
你說我將來一定.
不是將來 明年就一定有這個耳塞.
我一定要相信.
所以就等同於.
好像由母變有.
由死人可以復活.
這樣的信心.
那就是我們的信心.
我們就是透過耶穌的復活.
我們去相信神的應許.
所以保羅就將亞伯拉罕的內心世界.
用真理來描繪出來.
好像畫了一個圖畫.
所以很巧妙地.
將舊約的精髓.
保羅用新約的言語.
來講述出來.
然後留意一下.
使到母變為有的神.
他在主面前.
作我們世人的父.
就是信心之父.
如經上所記.
我已經立你作多國的父.
17章第5節.
然後第18節.
他在無可指望的時候.
因信 留意一下.
因信亦有指望.
就是他的太太.
可以因為他而生的應許子.
以臣.
所以原來亞伯拉罕.
在這個時候.
再一次承諾他自己.
那個心靈在臣的應許上.
比當年更加難.
但是他信.
臣就以此為他的義.

$^{1761}$所以他現在真的能夠叫做多國之父.
這個就是無可指望的時候.
因信亦有指望.
有一首聖詩就是憑信.
阿伯拉就有這個信心.
然後第19節.
他將近百歲的時候.
就是神應許給他的時間.
雖然想到自己的身體如同死.
撒拉的生育已經斷絕.
他的信心還是不軟弱.
並且仰望神的應許.
總沒有因不信心裡起疑惑.
反倒因信心裡得堅固.
將榮耀歸給神.
阿保羅將整段聖經.
那個最重要的17章描繪出來.
就是記得信心是建立在.
不是我們自己的幻想.
是建立在神的兩次的應許.
你的太太明年這個時候.
必定會幫你生一個兒子.
你一定要幫她改名叫做以撒.
我就要跟她堅定約.
你可以想像到.
原來亞伯拉罕那個時候.
他看到自己的身體如同死.
即是不可以生的.
但是他的信心不軟弱.
因信不得堅固.
不會起疑惑.
因信心裡就堅固.
更重要的是可以說.
Hallelujah敬拜神.
將榮耀歸給神.
神仍然眷顧他們兩老.
神仍然要在他們的家族.
衍生救世主.
你可以想像到.
阿保羅真的很淋漓盡致.

$^{1801}$將整個亞伯拉罕.
他在信心的成長.
原來我們可以做一個總結.
15到17章原來亞伯蘭.
一直都是信靠神.
一直都是憑信.
得到神的應許.
神的喜悅.
即是我們今天信耶穌的弟兄姊妹.
我們一直在我們的生活.
去經歷神的帶領.
有時帶我們越走越難的.
信心的旅程.
但是不要灰心.
好像亞伯拉罕一樣.
神的應許一定會兌現.
這就是我們看到這段經文.
更加告訴我們.
就是四章十七到二十一節.
第二個重點.
第二個重要的重點就是.
信心和恩典是亞伯拉罕成長的經歷.
一切都是恩典.
亞伯蘭不斷說.
以利以謝.
不如以實馬利.
十三歲了.
神話不可以.
所以其實是恩典的祝福.
信心的成長.
也是我們今天.
信耶穌.
恩信稱義的一個很重要的基礎.
我們可以回到第十三節這段.
稍微解釋一下.
因為神應許.
亞伯拉罕和他的後裔.
必得承受世界不是因律法.
是因信而得的義.
若是屬乎律法的人.

$^{1841}$才可以得為後治.
信就歸於虛空.
即是說.
做一個分析.
如果只是受割禮.
猶太教.
猶太人.
才可以成為後治.
變成只局限了恩典是透過行為.
透過一些猶太教他們的傳統.
才可以得.
神就不是用這個系統.
神就是用一個信心.
信靠他的愛子.
這個系統.
即是說普世的人.
都可以透過主耶穌.
就可以得到神的救恩.
因為如果只是這些人.
信就沒有必要.
因為你走進猶太教.
你受割禮就可以了.
根本就不需要信心.
所以說.
信就歸於虛空.
應許都會廢棄.
但是第十五節.
因為律法是惹動憤怒.
即是說如果你要靠自己的好行為.
你永遠達不到神的標準.
永遠不能夠透過律法.
透過自己做好人.
就可以得到神的祝福和恩典.
因為律法本身就是告訴你你是一個怎樣的人.
它不可以幫你去做到神的要求.
所以這裡有一個這樣的解釋.
去到第七章.
就會再詳細地描繪律法的功用.
所以這裡就一個基本的分析.
保羅不是想暗示要這樣做.

$^{1881}$而是反而在第十七到二十一節.
保羅就告訴我們.
原來亞伯拉罕也是舊約的人物.
但是他的成長是不斷憑信心去經歷神的應許.
經歷神的恩典.
這就是真正舊約信心的人生的路.
就不是律法.
所以這裡告訴我們.
所以第十六節也有一個總結.
所以人得為後製是本乎信.
因此就屬乎恩.
就是說舊約的人都是憑信得到神的應許.
這就是恩典.
叫應許定然歸給一切的後裔.
不單止是屬乎律法的後裔.
即猶太人.
也是歸給效法亞伯拉罕的信心的軌跡的人.
即是外邦人.
所以第十七節就開始說.
我們反過來說.
就讓我們知道.
亞伯拉罕的人生其實都是信心和恩典的偏激.
而成長的一個人.
即是和我們今天信耶穌的弟兄姊妹.
是完全一樣.
舊約的人都是因信稱義.
意思就是這樣.
我們還有一個時間.
九點鐘.
我們可以留待下一個星期.
因為二十一到二十五節.
我們還沒有詳細解釋.
因為二十五節是一個很重要的總結句.
就是人類的命運.
必須要像亞伯拉罕那樣.
來憑著神的應許.
將自己的信心投射到耶穌身上.
因為耶穌是為我們傳言的.
帶上的神應許的禮物.
如果我們將我們的信心放在主耶穌身上.

$^{1921}$我們就得著.
罪已經被消除.
信心的建立就是與神的關係的復和.
我們今天大概將這一段.
一到第二十節.
很簡單地解釋了.
還有一兩個關鍵字.
都差不多了.
應該都是大概的意思.
我們可以看看第二.
不過二十五節我們還沒有講.
我們先不背.
我們下個星期再看.
我們一個祈禱.
多謝天父.
今天我們很快將第四章的羅馬書.
特別是同創世記的十五到十七章.
亞伯蘭成長的旅程.
其實就是一個信心經歷神.
相信神應許的旅程.
保羅要強調的.
亞伯拉罕雖然是舊約的人.
甚至乎他後來受割禮.
但是割禮絕對不是一個得救的憑據.
反而是一個信心的.
相信的印證.
他來到這裡.
因為信神.
神與他立約.
得到神的祝福.
多謝你們今天讓我們透過主耶穌.
可以與你有永生的關係.
我們可以隨時去到天上的致聖所.
向你禱告.
向你求問.
以及有你的聖靈的帶領.
很盼望這些經文的講解.
成為很多弟兄姊妹.
可能初初讀羅馬書.
能夠明白.

$^{1961}$能夠掌握.
以致我們大家都建立.
我們與神之間很重要的信心的基礎.
以及繼續在信心上.
因信而至於信的成長.
我們感恩禱告.
奉耶穌基督的名.
阿們.
\newpage



\section{羅馬書}
\label{sec:23gUJZhHg2o}
\textbf{2025-05-22 周四晚查經班《羅馬書(11):分題:稱義帶來與神復和》 講員:曾裔貴牧師 NEW}
\newline
\newline
連結: \href{https://youtube.com/watch?v=23gUJZhHg2o}{\texttt{https://youtube.com/watch?v=23gUJZhHg2o}} ~~~~ 語音日期: 2025-05-24
\newline
\newline
\hyperref[sec:QTpL9O5GR3E]{\small{< < < PREV SERMON < < <}}
~
\hyperref[sec:index_chronic]{\small{[返順時目]}}
~
\hyperref[sec:index_scriptual]{\small{[返順卷目]}}
~
\hyperref[sec:r2uz84DdreY]{\small{> > > NEXT SERMON > > >}}
\newline
\newline
$^{1}$我們一起同心禱告.
多謝天父的恩典.
我們在晚上再次又恢復.
可以來到去查考《羅馬書》第四章.
很盼望這一份恩順清儀的恩典.
我們能夠牢牢地掌握.
以及認識我們可以透過主耶穌.
與神,與人可以發展一份和好的關係.
願神幫助帶領我們.
我們晚上不等待在你的面前.
我們這樣禱告.
奉耶穌基督的名.
阿們.
感恩我們再一次來到去.
開始我們晚上的查經.
上一次我們差不多講到四章的十七節.
可以跟大家重溫這個段落.
然後我們今天要處理就是.
因為加拉太書的第三章.
保羅就從一個實際的案例.
就是當時加拉太教會.
已經有一些接近異端.
來衝擊教會的信心,信徒.
加拉太教會的信徒.
他們那個恩順清儀的信心.
就隔閘了要行律法.
其中特別是行割禮.
所以是很直接對抗著.
保羅所講的舊約已經有恩順清儀.
亞伯拉罕是一個明顯得到神祝福.
而是憑信心去經歷.
恩典的一個重要的聖經舊約人物.
還要是未有律法頒佈.
連摩西也要很多年後才出現.
所以是一個很重要的.
其實是聖經自己本身的一個印證.
就是保羅發現.
所以不是保羅去發明的.
所以我們今天完結了這一段十七到二十二節.
二十五節可以去完.

$^{41}$我們就會處理第三章的加拉太書.
最後簡單做一個總結.
四章這一段經文的總結.
保羅為什麼要講亞伯拉罕的清義.
當然上個星期我們講過了.
就是堅固律法.
不是否定律法.
更加不是叫人以為是不需要做任何事情.
就可以得到神的救恩.
其實是希望能夠讓猶太人.
不要繼續帶著一份誤解,執著,敵視.
來抗拒傳福音和信心的清義.
所以很重要第三章的加拉太書.
我們就溫習一點.
仍然是這個大段落.
我們已經差不多處理了這兩個大段落.
這個分段是相當重要的.
一章的一到十七節就是一個重要的引言.
特別是十六,十七節.
就是講了信心帶來的.
就是跟神的關係和好.
這個福音是神的大能.
要救一切相信的.
猶太人,外邦人.
都可以恩信稱義的.
所以萬民都是透過信主.
來得到神的義.
講了答案.
但是十八節到三章二十節.
就講了需要恩信稱義的原因.
就是萬民欠缺神的榮耀.
虧缺神的公義.
所以萬民都是在神的憤怒.
咒詛之下必須趕快接受.
在律法以外的救恩.
就是神藉著他的愛子.
來為人贖罪.
這樣就可以得到無條件的.
憑信心去得到的義.
所以三章二十一到四章的二十五節.

$^{81}$有一個很重要的段落去講解.
怎樣可以得到這個義呢?.
就是在律法以外.
藉著耶穌基督的犧牲.
所以是神賜給每一個相信的人.
當然了,保羅一而再說.
先是猶太人,後是希利利人.
就是說其實每個人都是平等的.
在造物主這位神的眼中.
每個人都是需要信就行了.
牽連了可能就是令到猶太人.
有一個很大的誤解.
就是你們將摩西的律法.
我們一直都虔誠地遵守的律法.
來到完全的廢掉嗎?.
怎樣能夠讓我們猶太人.
能夠有一個滿意的解答呢?.
到底律法的功用又是怎樣呢?.
這些都是往後經文會處理的.
關於律法的意義等等.
待會我們三章的迦太書.
都有提到律法有什麼用呢?.
所以我們就要在這個分段上.
大家都應該可以很簡單地拿捏到.
就變成不會一個十六章的經文.
我們就很籠統.
你能夠找到很重要的分段的段落的轉折.
這些轉折是關鍵的.
譬如「但如今」.
就是說說完一個需要救恩的原因.
就說那個答案.
所以一章十六十七.
和三章的二十一至二十五.
就是一個頭尾的呼應.
就是哪些人需要救恩.
什麼方式去得到救恩.
其實就是我們一直說的.
信就可以得著神的救恩.
信就是和這位主有一個生命的聯合.
就是一個救恩.

$^{121}$所以我們就要記住這些重要的分段.
接著我們再簡單分一點.
上個星期我們就說了這三件事.
就是四章一到第十六節.
就是割禮.
其實保羅很強調.
割禮是在《創世記》的十七章.
亞伯蘭神幫他改名叫做亞伯拉罕.
然後就要他們所有的男丁.
行一個割禮.
象徵著這個民族.
當然以亞伯拉罕做中心.
就和神有立約的記號.
是屬神的.
但是保羅很清楚告訴我們.
亞伯拉罕信神.
那是《創世記》的十五章.
所以其實是在很多年後.
其實是差不多十三年後才行割禮.
所以保羅強調割禮.
只不過是內心相信神的一個記號.
其實內心的相信.
亞伯拉罕早就已經具備.
而神就稱他為義.
所以我們看看四章的十一節.
讀這些經文就給大家有少少的牢固掌握.
因為大家未必會經常去讀.
經常去看.
你不能夠重複這樣看.
很快就會淡忘.
希望大家都能夠牢記.
十一節就說.
並且他受了割禮的記號.
留意一下保羅的描述.
作他未受割禮的時候.
因順清儀的印證.
叫他作一切未受割禮.
而順之人的乎.
就是所有沒有受過割禮的.
無論是猶太人沒有受過割禮.

$^{161}$或者外邦人根本就不需要有割禮.
總之是順就可以得著神的稱義.
這個就是亞伯拉罕他自己.
作了一個示範.
神不是因為他做一些行為.
一些律法.
來得到神的稱義.
而是神應許給他.
他相信神的應許.
所以這個信心.
使他能夠連上.
紅神的這個.
清義的祝福.
這個恩典.
所以對我們外邦人來說.
我們就是效法亞伯拉罕這個信心.
所以我們是未受割禮而順的人.
亞伯拉罕就是我們一個屬靈的父親.
接著保羅還不是停在那裡.
他就說使我們順民義.
第十二節就繼續說.
他說有作受割禮之人的父.
就是猶太人.
可能自小就根據傳統的律法.
家教拉底的教導.
他就循規蹈矩這樣來到去.
因為他是一個猶太人.
他是亞伯拉罕的肉身的子孫.
他很自然地就要受割禮.
但是保羅就在這裡已經很清楚了.
他說有作受割禮之人的父.
就是那些不但受割禮.
不單止受肉身的割禮.
而且是要按我們祖宗亞伯拉罕.
他未受割禮而順的宗旨去行的人.
就是說仍然不是建立在那個行為.
那個肉身的某些做了什麼儀式.
來量度你是不是得到神的.
稱義的那個條件.
相反你行割禮.

$^{201}$你身為一個猶太人.
你不需要廢棄那件事.
但是那件事不能夠令到你.
和其他外邦人有不同.
有一個差異.
唯一的就是你行完割禮之後.
你的人生都要憑信心去相信.
神的應許.
神給人類的救贖.
就是耶穌基督了.
所以猶太人欠缺的.
就是他們要看到他們的彌賽亞就是耶穌了.
所以仍然都是要好像亞伯拉罕.
未受割禮的時候那個信心去行.
所以必須要這樣去帶出那個信心的意義.
當然我們也說過.
亞伯拉罕在他的人生.
由創世記15章到17章.
當然18章之後也是他成長的人生歷練.
我們看到的就是神不斷在他的家庭.
訓練他讓他經歷神的信心.
信實.
更重要的是神的應許.
一直將他的應許.
越來越具體地落實在亞伯拉罕的身體.
在他的家庭.
尤其是在他的生命裡面.
所以神好像有計劃地讓亞伯拉罕.
所以垂垂老矣.
連日的那個時間的洗禮.
那個虛度.
原來是一個恩典的歷練過程.
好像一些純酒.
越釀越放得久就越純.
越能夠體驗到那種生命的成熟.
所以亞伯拉罕要去到100歲.
其實應該是99歲.
他才真正經歷神的應許的兌現.
所以一直以來的以實.
以利以謝.

$^{241}$以實馬利.
都不是真正的真命天子.
但原來在這個過程.
原來是亞伯拉罕.
最寶貴就是他的人生在跌宕裡面成長.
和經歷恩典.
這個才是人生信仰的路最寶貴.
就是那個人自己本身是有成長.
如果沒有成長是嚴重的.
回教會20年30年.
依然顧我的人其實大有人在.
我不想說太多.
變成不是時非.
是一個觀察.
往往有時回來事奉的.
特別是施事.
有時不自覺.
不為意.
慢慢忙於事務.
有時根本沒有想過原來自己是敬拜者.
其實是要彌此勸勉提醒.
不是批判.
更加不是譴責.
是要很有意識地.
我回來不是給誰面子.
我去做些事.
我是傳奉獻袋.
我派程序表.
我要開門.
我要甚麼.
不是這些.
這些不是不重要.
是重要的.
是需要有人去做的.
但是每個人回來都是敬拜者.
有時觀察真的開始說到.
在後面打牙交談.
這不是生命成長.
偶然有些特別的事要處理.
關乎到那個聚會.

$^{281}$是沒有問題.
但是有意識地完結了基本的運作.
當講完很認真地讀經.
在後面的人在聊天打牙交談.
這就不是小事.
這是反映我們原來有意識地回來.
只是參與一個我們很開心.
見到我們自己很開心的那幾個朋友.
來聊天.
來做一些事.
本身的本質是嚴重的.
因為那一刻是所有人要平息以待.
神現在要向我們說話.
你可以想像得到是多麼嚴肅.
雖然我說到這件事好像誇張了.
每個星期都聽到.
你為何要搞到這麼嚴肅.
現在頒聖旨嗎?.
背景本來的意義是這樣的.
神的僕人領受了.
他盡他的能力.
盡他想的題目去準備好.
無論經文,例子,powerpoint.
都預備好.
他要怎樣演繹.
他要盡他的能力.
去讓回眾捕捉到.
得到那個信息.
這個信息是神透過他和那個人說的.
這個性質就是神性.
我們福音派就是相信這個東西.
聖靈藉著那個人.
來和那個回眾說話.
那個說話的人還不知道原來是和某個人說的.
我們要相信這個東西.
但如果成為一個結構性的打牙交.
這件事就嚴重了.
嚴重不是你影響了聚會.
其實沒有人影響到聚會.
沒有人會影響到聚會.

$^{321}$影響的是你那個人.
你經年累月都沒有經歷神恩典的生命成長.
這個就是.
如果再更加見微知著.
就更加嚴重.
所以,在《論中》就要提醒.
我們大家是回來參與敬拜的人.
講道的人都是安靜地敬拜.
參與一起事奉的.
不是坐在上面就等於我可以置身於事外.
你們唱歌,我就在上面.
時間到了我就站起來說話.
不是的,每個人都是一起敬拜的.
所以我們要很深化,很認真去想想.
我們設計這些聚會背後的目標.
和本身的意義.
所以伯拉罕,神一直不給他兒子.
原來就是他經歷這個信心.
對神應許的信靠.
無可倚靠的時候.
因信仍有倚靠.
就是這位伯拉罕.
所以第二方面.
應該是第四章的13到16節.
就是這一段.
當然17節就再描述多一點.
就是講他在17章創世記那裡的.
很繪形繪聲的描述.
保羅可以將伯拉罕的心路歷程.
剖析得很透徹.
讓當時羅馬修信的讀者去了解.
原來伯拉罕就是這樣考驗信心.
而相信神的應許.
所以最終他就在主的面前.
作我們世人的父.
再一次講全球化.
再一次講萬民都是神的兒女.
所以萬民都需要福音.
不是猶太人才需要福音.
不是一些發達國家的人才需要福音.

$^{361}$是萬民都需要福音.
如果有些地方是因為各種的政治,宗教,歷史.
導致那些地方是沒有福音的.
例如穆斯林的群體.
我們基督徒就有責任想不同的辦法.
例如透過衛星,營商宣教,醫療宣教.
讓穆斯林能夠認識主.
因為伯拉罕就是在主面前作我們世人的父.
所以萬民都需要福音.
我們繼續透過這些重溫.
幫大家更加牢固掌握伯拉罕的信.
堅定堅固舊約的律法.
是沒有抵觸的.
而且比律法還要早430年.
伯拉罕就已經經歷神的應許.
所以神的應許是用信心去信靠.
在這個信靠的過程.
其實是恩典伴隨我們的生命成長.
所以信耶穌不是你簽了一個紙.
我信了將來就在天堂見.
洗了禮就拿著證書.
就不是這樣.
是我們生命因信而致於信.
不斷在信心裡經歷生命與神與人的成長.
當然第五章就特別強調和好,復和.
這都是事實.
所以伯拉罕就是那一位.
相信神能夠叫死人復活.
第十七節我想再讀多一點.
再給大家有更深刻的印象.
看著這個powerpoint.
看第十七節.
這裡就是保羅再一次將創世記十七章.
比較有深度地.
好像他看到伯拉罕當時是這樣想的.
這樣去描繪.
留意一下第十七節.
伯拉罕所相信的是這位叫死人復活.
洗到沒有變為有的神.
當然他當時不知道誰叫耶穌.

$^{401}$他根本沒有聽過.
或者不知道將來會有一位神的兒子.
真的會死而第三天復活.
他相信的是他自己的身體.
他和太太薩拉的身體已經垂垂老矣.
是不可能孕育新的生命.
是這樣的.
相信自己好像死的人一樣.
是不可能有新的生命.
可以透過兩個老人家來產生.
伯拉罕所相信的是這位叫死人復活.
使無變為有的神.
他在主面前.
作我們世人的父.
如經上所記.
我已經立你作多國的父.
就是創世記十七章.
應該是第五節.
跟著他在無可指望的時候.
因信仍有指望.
就得以作多國的父.
正如先前所說.
你的後裔將要如此.
保羅就是這樣前後的兌換經文.
就引述創世記十五章第五節.
就是神帶他出到營外.
叫他望天上的眾星和地上的沙的時候.
就應許他你的後裔.
好像天上的星和地上的沙一樣.
伯拉罕就相信耶和華.
耶和華就以此為他的義.
保羅我們上星期都特別點出.
保羅好像人錯聖經一樣.
第四章的第三節.
經上記著說.
伯拉罕信神.
肯定是人錯聖經.
保羅你都不知道歷史的亂來.
正正就是不是亂來.
正正就是他還沒改名.

$^{441}$保羅在這裡幫他改名.
就是說保羅用經文.
用了一個純屬的程度.
就是在神的眼中.
亞伯蘭仍然是亞伯蘭.
但是他信的那一刻.
其實已經是多國之父了.
已經是亞伯拉罕了.
所以保羅這樣人聖經.
真是可圈可點.
不是死人那樣.
依書直說.
而是他真的能夠運用.
亞伯拉罕是信心之父.
在他還沒受割禮的時候.
還沒改名的時候.
他已經信神的應許.
在一個好像自己沒可能會信的那種心境.
都來相信神.
這一位能夠改變萬事的神.
所以這個信心是得到.
一定稱他為義.
所以當我們看回這裡的時候.
就是第十九節.
接著他說他將近八歲.
再一次跳回十七章那裡.
將近八歲的時候.
雖然想到自己的身體如同已死.
當然包括了.
生育是一男一女的事.
撒拉的生育已經斷絕.
其實斷絕了很多年.
他的信心還是不軟弱.
所以不是一種強硬的迷信.
是神的應許.
清清楚楚.
你們夫婦要生的兒子.
已經改了名字.
就叫做爾薩.
明年這個時候.

$^{481}$你們就可以得到這個應許的兒子.
所以亞伯拉罕做了盤點.
不是一種阿Q精神.
真的做了盤點.
我九十九歲了.
撒拉九十歲了.
真的沒有了生育的能力.
但是他唯一要在他兩個人的信心以外.
去相信的.
就是這位可以改變人生命的神.
這位可以叫死人復活的神.
所以他的信心是軟弱的.
所以我們不要否定自己的信心的軟弱.
但是有一個外力的介入.
就是神的應許.
所以我們求主讓我們能夠.
在主耶穌裡面.
一方面不要看死自己.
一方面在神裡面有無限的可能.
所以第二十節.
並且仰望神的應許.
總會有因不信心裡起疑惑.
反倒因信心裡得堅固.
將榮耀歸給神.
所以就帶來一個這樣的循環.
清楚自己不是自己逼自己相信.
而是很明顯清楚神的應許.
如果要兌現.
我的信心大和小是沒有關係的.
所以既然有這麼清楚的應許.
我就信靠神.
所以信心是投靠在神那裡.
就帶來一種改變.
所以不是一種不斷自己的幻想.
不斷自己的心理作用.
而是投靠在神的應許裡面.
所以最終如果那件事成就.
就不是哪個人的信心特別大特別少.
而是神的應許的兌現.
所以榮耀永遠不是歸給大有信心的人.

$^{521}$變成我們有時說靈恩的神人,神醫.
所以不是的.
世界不應該有神人,神醫.
是應該信靠神.
所以榮耀歸給神.
所以你看到保羅很清楚說的.
心裡面得堅固.
那件事兌現就將榮耀歸給神.
並且滿心相信神所應許的必定能夠作成.
所以這就算為他的兒.
所以保羅是不斷遊走在十五,十七章的創世記.
和亞伯拉罕他最重大的考驗.
當然這個考驗完之後.
二十二章又有另外一個大考驗.
就是請你將你所愛的以撒.
獻在摩尼亞山上.
所以不斷有人生不同的考驗.
但是在這裡亞伯拉罕在找神的應許.
所以我們今天我們基督徒.
我們有一個客觀的事實.
耶穌基督神的兒子真的來過這個地球.
生在以色列.
所以我今天的PowerPoint都給大家.
那些苦路大家看一下.
主耶穌真的走過這樣的階梯.
苦路來到這個世界.
這個就是神對人類的應許.
只要透過耶穌.
我們就可以進入和神的那個復和.
那個關係的重新的建立.
就會帶來種種的生命的豐盛.
福樂,平安.
五章第一一講就馬上講平安.
我們與神相和就立即帶來那份心靈的平安.
所以我們要記得亞伯拉罕都在找應許.
我們今天信徒都在找神.
透過新約使徒們給我們的啟示.
我們在找那些的應許.
所以有多少的應許呢?.
我們在找越是讀聖經的人.

$^{561}$越是深信聖經的人.
就會越感受到神的應許的兌現.
例如我們真的相信我們的禱告.
是可以成就神的工作的.
我們用信心不斷去為人祈禱.
是真的神會聽祈禱的.
我們雄恩也是.
有位姐妹的女兒.
三十多歲的.
三個星期前突然找我找得很急.
要我一定要趕去威爾斯的ICU.
因為她有肺炎.
可能感染了一些病毒.
很嚴重很危急.
因為呼吸不了.
去到已經插喉.
我們都沒什麼可以做了.
在她耳邊控得很緊.
因為她根本完全不能夠說話.
見到人.
好像近乎昏迷一樣.
她在耳邊祈禱.
回到教會.
星期三.
我們都祈禱.
交託給神.
第二個星期是三.
再去探望她.
還是在威爾斯.
又是在ICU.
但是好一點了.
但還是插著一些喉等等.
又是要靠呼吸.
但是感覺開始有些眼神.
但是說不了話.
都是閉上眼睛.
很累.
可能用了藥.
第三個星期.
上個星期.

$^{601}$就轉去上水的北區醫院.
開始有一點反應.
開始可以坐在那裡.
但是說不了話.
因為她一直都不認識我.
我兩次去.
她都是分在威爾斯的ICU.
我在北區的ICU去探望她.
她坐在那裡.
可能我又趕回夜晚的祈禱會.
很快叫護士開門進去.
她看到陌生人衝過去.
她很害怕.
她的反應很害怕.
我也很尷尬.
因為她不知道我是誰.
我表明我的來意.
我是你媽媽教會的牧者.
為你祈禱.
她也說好.
祈禱.
祈完我就馬上走了.
因為很尷尬.
她好像不知道我是誰.
真的看到神的工作.
剛剛昨晚.
第四個星期.
她已經開始可以吃糊了.
開始不需要再用呼吸機.
還有說話.
逐個字逐個字用力說出來.
看到神真的有很大的作為.
譬如我們之前有Chris的野地牧人.
那位伊朗被囚禁的宣教士.
在土耳其.
Mustafa.
我們祈禱了一年.
一直都有為他祈禱.
真的很振奮.
在監獄最終能夠拿到居留權.

$^{641}$難民的身份釋放.
就可以去加拿大.
是一件很振奮的事.
所以最近也感受到.
大家一起禱告的力量.
所以我比較老套.
我仍然覺得.
星期三晚上的公禱會.
我們不是要集氣.
而是大家認定.
我們分別半個小時.
可能45分鐘.
來給神.
大家同心的禱告.
神會聽祈禱.
重點是我們同心去信靠神.
所以亞伯拉罕.
他掌握到神的應許.
用他的信心去信靠.
然後就是他的生命裡.
在神的恩典的伴隨.
這個伴隨肯定是今天我們第四章結尾.
第五章頭會說的.
就是聖靈的介入.
在每一個信徒的心靈裡.
所以亞伯拉罕能夠成長.
肯定就是聖靈在他內心的工作.
當然舊約就不是那麼明顯.
但新約就肯定是這樣的.
包括保羅待會在加勒太書三章.
都是這樣說的.
所以我們看第四章.
保羅不是要和猶太教對抗.
而是很想讓猶太人知道.
其實不需要掙扎.
救恩是白白的.
在舊約已經有恩信清儀了.
在舊約神已經設計了.
人其實是不信.
人是信心軟弱.

$^{681}$人也不可以靠自己去得到神的救恩.
或者說人靠自己.
在神以外自己可以成就救恩.
是不可能的.
因為罪已經令到人扭曲了.
已經沒有力量行善了.
所以如果猶太人還堅持在摩西的律法.
其實他也不明白.
整個神的創造和救贖計劃的進程.
所以第四章的重要性就不言而喻.
保羅真的很需要告訴猶太人.
其實他們欠缺的是信主的義.
不是他們的律法行得不足夠.
這就是我們上星期說的一個簡單的總結.
所以當伯拉罕信死人復活.
使母變為有的實人的時候.
他就是在我們人類後人.
在神的設計.
他就是我們信心之父.
所以我們也要求主讓我們能夠在我們成長.
某程度教會信主越來越大的弟兄姊妹.
應該是信心堅強.
信心能夠透過很多的經歷.
來勉勵到新信主的朋友.
那就成為大家不斷去建立彼此的信心.
應該是這樣的一個成長.
否則就有很大的問題.
教會如果越大人越多.
信主越久的人信心越更加軟弱.
那就是一件不公平的事.
教會根本沒有成長.
這是一個大的問題.
第四章我們差不多到這裡說完.
不過第23章開始就很重要.
因為保羅再將我們帶回我們自己.
如何透過耶穌來成為亞伯拉罕的子孫.
這個就是保羅在第四章的重點.
就是到底新約的恩信青兒.
其實是如何去兼顧律法.
沒有抵觸到律法本來的設計.

$^{721}$保羅就告訴我們.
拆解了亞伯拉罕的信心.
他就是在找神的應許.
就是他的兒子.
未出生的兒子.
就是神應許給他和撒拉的.
而且是蒙祝福的後裔.
當然一路說.
加拿大書說其實那個後裔是耶穌基督.
三章應該是19節.
這個就是我們今天和亞伯拉罕的信是完全一樣.
所以他相信神就被神稱為義.
這個就是兼顧律法的原因.
不是否定律法.
不過人是不能夠靠自己去行律法.
就可以得到神的應許.
所以不需要繼續堅持在律法底下.
但是在律法以外.
神已經做了功了.
神已經付了贖價了.
所以只需要透過相信神的應許.
就可以被神稱義了.
這裡是新漢語譯本.
我覺得不錯.
再一次讓大家一起重溫這一節聖經.
其實就是四章十一節.
他又領受了國禮為標記.
作為他在未受國禮的時候.
因順稱義的印證.
使他就是亞伯拉罕.
作一切沒有受國禮而相信的人之父.
好讓他們我們這些沒有受國禮的人.
和亞伯拉罕一樣.
這樣去信就可以被算為義了.
所以我們未必需要有這個印記.
但是我們今天不是用國禮.
今天保羅就很斬釘截鐵.
每一個基督徒信主就受了聖靈為印記.
所以很清楚的.
以不鎖書的一章.

$^{761}$一章我忘記了是十七還是第幾字.
就受了這個聖靈為印記.
印記這個字就特別意思了.
就真的好像一個訂金的意思.
在以不鎖書就是這個意思.
所以這裡保羅很想強調.
未受國禮的亞伯拉罕.
和我們完全一樣.
這樣去經歷神的呼召.
帶領.
接著就是次下應許.
相信.
經歷成長和應許的兌現.
所以我們留意一下.
我們看看第四章的23至25節.
23至25節.
算為他兒的這句話.
算為他兒就是第15章第5節創世記.
不是單是為他寫.
保羅好像很大膽.
將舊約重新去解釋.
當然不是這個意思.
而是他能夠看到.
如果我們這樣跟隨亞伯拉罕.
這樣去相信神的應許.
我們就會被稱為義.
所以保羅在這裡有一個這樣的演繹.
就是算為他兒這句話.
不只是為他亞伯拉罕寫的.
可見聖經原來真的可以超越了時代.
第24節.
也是為我們將來得算為義之人寫的.
就是我們因信聽聞這個福音.
我們就信我們的主.
就可以稱為義.
就是我們這些因信稱義的人.
就是我們這信神.
使我們的主耶穌從死裡復活的人.
這就是神的應許.
神無條件的救贖.

$^{801}$神賜下的和好的關係.
就是藉著他的愛子.
所以基督教的那種感人之處.
那種偉大之處.
不是用壓力.
不是用威嚇.
是用愛的關係的復活.
這個關係的復活就是神的兒子.
親自來為人死.
所以基督教最偉大就是.
不斷傳遞這份有力量的愛.
這個愛不是廉價的.
是帶來關係的復活.
所以這裡就告訴我們.
就是我們這些信神.
使我們的主耶穌從死裡復活的人.
就是這樣被算為義.
就是說和亞伯拉罕信以撒的出現.
是完全一樣的.
是完全異曲同工的.
完全吻合的.
所以原來我們今天.
都是一樣亞伯拉罕的子孫.
是因信而成為了亞伯拉罕的子孫的人.
接著第25節.
這裡特別補充一次.
人類必須要面對的就是.
這個神解救人類免除罪惡和審判的方法.
就是耶穌被交給人.
是為我們的過犯復活.
是為我們稱義.
或者有另外的翻譯.
耶穌是為我們的過犯交付了.
交付真的完全好像拿出去一樣.
一種行刑的犧牲.
你可以想像到.
真是肝腸寸斷.
一個父親要犧牲他自己的兒子.
亞伯拉罕在22章的創世記.
又經歷這種生離死別的考驗.

$^{841}$將他所愛的以撒.
是真的需要量度的.
真的要他命的.
要親手殺的.
這樣來交付出去.
所以這種交付是一個很大的信心的行動.
主耶穌就是這樣被交付.
交付是去刑場.
交付是去到死.
好像有誰說過.
好像有團體說過.
有一位都不算很熟大式的.
一位舉世知名的教授.
港大的.
早一年退休.
現在退休很年輕.
他現在已經轉去.
因為他在港大的講座教授退休了.
甚至是一些重要的.
大的一個項目的主管.
大學生,同事.
很多人歡送他.
因為一個最後的授課.
很多人說你不需要退休.
我們很珍惜.
他說他真的每一堂課.
人們是不會走堂來聽他講.
他在他的教會講見證.
講道,分享.
他說.
他回家.
晚上回家.
結束了聚會.
人們很開心.
很捨不得來歡送他.
讚他這麼多年的教學.
很多的圖你滿門.
是一個舉世知名的學者.
他說他回到家大哭.
他這樣說.

$^{881}$他的退休得到這麼多人的肯定.
這麼多人對他的尊崇.
他將榮耀歸神.
因為是一位很謙卑的弟兄.
他不是傳道人.
但他很謙卑.
也在教會講道.
我自己也有很大的感動.
他說他想起耶穌.
他的退休在十字架上.
他的一生造福這麼多人.
被唾罵.
被嘲笑.
被辱罵.
他說他算得了甚麼呢?.
教學是他的興趣.
也是他終生的職業.
他能夠得到一些的信用.
一些的榮耀.
是神的恩典.
但他想起主耶穌的退休.
就是這樣被人殘害至死.
你可以想像到.
那個落差有多大.
這種感動是很感動人的.
真的很感動.
可以想像到.
耶穌被交付出去.
是為了我們的過犯.
你怎會為一個有錯誤的人.
來交付你自己的愛子呢?.
所以這些不是小朋友的好消息.
是我們所有人類.
必須要重新去想想.
我們自己的存在價值.
耶穌這樣被交付.
我們應該如何在人生.
盡力為主去擺相呢?.
然後他的復活.
使我們能夠清怡.

$^{921}$所以是一個死而復活的福音.
所以是不斷持續關係的.
建立的復和的福音.
所以不是一不料.
按個按鈕.
我洗完禮就走了.
我不再回你的教會.
我不再回你的基督教了.
這樣的人其實根本不相信耶穌.
嚴格來說.
他洗了禮也是假的.
因為他根本錯誤地理解信仰.
信仰是不斷地說.
我們在這個應許神的兒子裡.
不斷地生命的建立和成長.
不斷地和神和人建立那個關係.
不斷地和好.
所以求神幫助我們.
四章我們就講完.
這個時候我們就要處理加拉太書的經文.
先喝口水.
第三章是加拉太書.
肯定就是在羅馬書接近那個時間.
處理教會內部問題.
因為開始有一些異端.
想來影響信徒.
本來信徒很好的.
憑信心入門又得到聖靈的工作.
外邦的教會.
但是突然之間感覺自己好像有些疑拉.
不是一個因信而帶來的滿足和喜樂.
好像被打沉了.
就是因為你們沒有像我們猶太人那樣.
信耶穌包括我們受割禮.
我們是神.
亞伯拉罕的子孫.
這樣的一種誘惑.
或者一種此事而非.
會令基督徒處於一種困惑裡面.
當然一定是傳遞到保羅的耳邊.

$^{961}$保羅知道.
又意識到如果再讓這種.
此事而非的教導去深化去影響.
就變成一種毒素.
來摧毀信心繼續的成長.
所以保羅意識到這個問題.
就寫這封信.
我們可以收集的範圍.
第三章全章開始就是一個很重的譴責.
就是對這些迦拉太信主的弟兄姊妹.
用比較重的言辭.
所以是針對一件事.
不是針對人.
無知的迦拉太人.
耶穌基督釘十字架.
已經活化在你們眼前.
誰又迷惑了你們呢?.
我只要問你們這一件.
你們受了聖靈.
是因行律法呢?.
是因聽信福音呢?.
所以舊約不等於沒有聖靈.
新約也不等於我們撇開了律法.
所以沒有抵觸的.
律法和信心都不應該有一個強烈的異分.
不是舊約就是律法.
新約就是信心.
不應該這樣想的.
但相反.
信心帶來可以在內心運作.
和遵行神的旨意.
肯定是聖靈在心靈裡面.
在每一個信弟兄姊妹的心靈裡面.
工作的那個.
不是動力.
是那個引導.
是聖靈這位主.
在我們大家心裡面的工作.
所以你們受了聖靈.
是因行律法而得到.

$^{1001}$還是因為聽信福音呢?.
肯定答案就是一個切問.
一種挑戰.
其實就是透過信福音而得到聖靈.
這個已經是普遍的事件.
不是那些怪異的行徑.
那種聖靈降臨.
聖靈充滿.
就不是初期教會那些.
而是一個普遍的事實.
所有的基督徒.
信耶穌是救主.
就立即獲得勝靈.
接著第三節.
「你們既靠聖靈入門.
如今還靠肉身承傳」.
即是說.
沒有否定他們是信了主.
有聖靈的.
但那個行徑.
開始像是一個.
自己不斷要透過做好.
做好.
這個不做不行.
那個一定要做.
這樣的做法.
這種心態.
就落在一種.
好像自我控訴.
自我去完善的救恩.
就不是一直在聖靈的引導下.
用愛用信心.
這樣來經歷人生的不同階段.
兩者是完全不同的.
結果.
那條路也不同.
這裡說.
「如今還靠肉身承傳」.
你們是不是這樣無知呢?.
第四節.

$^{1041}$「你們受苦如此之多.
都是徒然的魔」.
似乎就是初信的時候.
他們可能經歷了.
信耶穌帶來的代價.
例如羅馬的君王.
對一些.
如果決心信主.
整個城邦都信主.
而不再君王敬拜.
就會有一定的考驗.
可能是一些逼迫.
對於猶太人.
也會有更大的考驗.
就好像你否定了守律法.
和脫離猶太教.
所以都相同的考驗.
所以對於保羅來說.
跟隨這個福音的人.
在保羅的認知裡面.
一定會伴隨著困難的.
所以你們受苦如此之多.
是不是徒然的呢?.
難道果真是徒然嗎?.
第五節.
再一次提到.
「那似給你們聖靈.
又在你們中間行義能」.
是因為你們行律法.
還是因為你們聽信福音呢?.
正如亞伯拉罕信實.
「這就算為他的義」.
再一次引用十五章的《創世記》.
第五,第六節.
第七節.
「所以你們要知道.
那以信為本的人.
就是亞伯拉罕的子孫.
以信為本」.
即是說.

$^{1081}$我們用廣東話就很明白了.
「跌埋心水去信」.
「一心一意去信」.
不是靠律法.
不是要用律法去完事.
這樣就會明白很多了.
「以信為本的人.
就是亞伯拉罕的子孫」.
第八節留意這裡.
「並且聖經既然預先看明.
神要叫外邦人.
因信稱義」.
這是保羅自己的演繹.
「就早已傳福音給亞伯拉罕說.
萬國必因你得福」.
這段聖經.
保羅是引用《創世記》的十二章.
第三節.
是亞伯蘭離開梧爾.
去迦南地.
有神的呼召.
帶他離開本地附加.
有這個應許.
「我要使你成為萬國的祝福」.
所以保羅引用這節聖經.
即是說.
在他改名.
或是應許他的兒子之前.
還要再早.
所以保羅說.
「聖經既然預先看明.
神要叫外邦人.
因信稱義」.
所以使萬國都會蒙福.
這麼早期.
那時是七十五歲.
離開梧爾.
輾轉去迦南地.
到底他在梧爾時被宣告.
還是在迦南再宣告呢?.

$^{1121}$似乎是再早一點.
在梧爾的時候.
這裡讓我們看到.
原來聖經在舊約.
已經有神的計劃.
要萬國都因信稱義得福.
所以首先就用亞伯拉罕的人生.
作為一個示範.
所以後人.
只要跟隨亞伯拉罕的信心.
就可以經歷神的應許.
和神的祝福.
所以神就透過亞伯拉罕.
去祝福萬民.
當然一直再仔細說.
神是透過子孫去祝福萬民.
就是耶穌基督.
第九節.
「可見乃以信為本的人.
和有信心的亞伯拉罕.
一同得福」.
可以看到不斷強調.
我們必須要有信心.
不是守律法的信心.
是信靠神.
應許的信心.
這樣來在我們的生活.
這樣依靠神.
就跟有信心的亞伯拉罕.
同樣的得福.
第十節.
接著做一個對比.
就是針對那些傳假道的.
可能是猶太的信徒.
「凡以行律法為本.
都是被咒詛.
因為經常記著.
凡是不常照律法書.
所記一切之事去行.
就被咒詛」.

$^{1161}$這是摩西在他差不多離世.
在《申命記》.
應該是27章說的.
就是說這一切的律法書.
就是神跟你們這個民族立的約.
我就要走了.
我就要見神了.
你們就要去跟隨這些律法.
你們就要在這個約下.
去做你們應該要做的了.
所以律法不是說我信.
就是你要實踐去實行.
也告訴我們.
正正就是人不可以靠自己.
來實行神的律法.
所以是需要神的.
應許和恩典.
去將我們帶進信心的意義裡面.
第十一節.
「沒有一個人靠著律法.
在神的面前稱義」.
這是明顯的.
因為經常說.
「二人必因信得生」.
哈柏谷書我們引過.
一章的十七節就引過.
哈柏谷書的二章四節.
二人是因信得生.
所以這一段.
是特別強調那個信.
就可以得到這個稱義.
第十二節.
「律法源不本乎信.
只說行者些事的.
就必因此活著」.
這段是利未記的經文.
第十三節.
「基督既為我們受了咒作」.
或者說成為咒作.
「救贖出我們脫離律法的咒作.

$^{1201}$因為經常記著.
凡掛在木頭上.
都是被咒作的」.
申命記的經文.
第十四節.
「這便叫亞伯拉罕的福.
因基督耶穌」.
所以那個關係就是.
亞伯拉罕的這個子孫.
他的信心的那個兒行.
就是被交付到死亡.
來拯救.
幫我們脫離我們的過犯.
所以變成亞伯拉罕的福.
因基督耶穌.
就可以臨到外邦人.
其實要拆解.
這些全律法主義的.
可能是猶太信徒.
可能是一些領袖.
來逼令加拿大人.
要回頭.
不要只是信.
要做一些行為.
其中要顯著的就是國禮.
很大個人都要去國禮.
表明你們是屬於神的.
但是扭曲了國禮本身的含義.
以及會令到這些.
本來用信心入門的人.
開始出現混淆.
以為耶穌的救恩不是白白.
以為耶穌的救恩不夠徹底.
去將人帶進神的救恩的愛裡面.
那就是大件事.
所以這裡就必須.
我們要得到亞伯拉罕的福.
就是神允許亞伯拉罕的.
他的遺孫將來的後裔可以蒙福.
所以那個管道.

$^{1241}$必須是透過基督耶穌的犧牲.
所以亞伯拉罕的福.
因基督耶穌.
當然就是因基督耶穌的犧牲.
就可以臨到外邦人.
使得我們可以能夠恩信.
接著這裡繼續說.
再一次補充.
恩信得著所應許的聖靈.
信心的入門.
就是藉著基督的救恩.
十價的犧牲.
我們就得著所應許的聖靈.
聖靈進駐我們的內心.
會使我們知罪.
會使我們對自己有一個更加認識.
更加有力量去行.
神要求我們行的.
所以這個就變相.
神的律法就在我們內心.
給我們力量去行.
遵行神的律法.
所以全律法都是包在.
愛人如己.
這一句話.
其實是加里大書的文章.
所以你可以想像到.
原來沒有抵觸律法的.
不過是給我們一個新的生命力量.
去行神的律法.
所以這裡很重要.
從頭開始說.
我們是憑信心入門得到聖靈.
現在再一次告訴我們.
叫亞伯拉罕的福.
因基督耶穌可以臨到外邦人.
使得我們恩信就得著應許的聖靈.
因為聖靈是這麼寶貴.
這麼重要.
在整個救恩的運作裡面.

$^{1281}$聖靈進入我們每一個人的生命裡面.
我們就等同於和亞伯拉罕一樣.
這樣去過我們信心的人生.
不同階段有不同的應許.
有不同的兌現.
有不同的我們生命的眼界成長.
我們作為牧者.
人生的閱歷如果換成甚麼都咒詛咒罵.
眼界就停留在一個很窄的範圍.
因為你衝破不了自己的苦毒.
但如果人生的歷練.
有很多可能是眼淚.
是令你更加成長.
胸襟更加大.
對人,對事物的眼光更加遼闊.
這就是你的眼淚經驗不單沒有白費.
而且成為了成長的種子.
所以仍然是聖靈在我們裡面的工作.
所以一個成長基督徒.
心胸越來越大.
眼光對人,對神的愛越來越真誠.
信心越來越大.
信心不是要走過紅海.
要這個人按手掣立即站起來.
不是這樣的.
信心伴隨著一種很平安的內心.
很有一種讓人感受到一種.
又不能說是靈氣逼人.
但總是有一種氣質.
人感覺很舒服的.
這個人會有散發聖靈的工作的氣質.
所以我們要記得.
保羅在這裡告訴我們.
他是特別針對那些全港要受割禮.
才能夠補足基督的救贖.
那就大件事了.
肯定是接近異端的.
所以原來神的方法不是我們脫離律法.
是把聖靈給我們.
讓我們的生命在聖靈裡面.

$^{1321}$一直在愛裡面去進行神的律法.
所以全國禮的人就以為做了這件事就行了.
其實不是這樣做的.
所以這裡和我們去了解亞伯拉罕.
也是相當重要的一個指標.
在這裡的亞伯拉罕.
是以信為本.
一心相信神.
成為了我們的信心之父.
我們效法他的信心來得到.
還有一樣東西就是神賜給亞伯拉罕的福.
現在是因著基督耶穌的犧牲.
就在我們外邦人身上.
這個福其中在今生運作中.
我們要很有意識地.
所以我們不要被那些靈恩派壟斷了.
和搞砸了聖靈的工作.
所有基督徒這裡很肯定的.
因信就得到應許的聖靈在我們心裡面.
是每一個基督徒肯定獲得的.
所以是常態來的.
不是特殊的.
不是突然間有些神人,神醫.
不需要的.
每個基督徒都要有聖靈.
是應許了給我們的.
我們只要是因著基督耶穌.
我們就不知不覺已經成為了亞伯拉罕的蒙福後裔.
這個福首先體現的是聖靈住在我們內心.
還有一點我們再讀.
後面的太詳細了,我就不讀了.
我看看時間夠不夠.
另外的部分都預了.
還有最後一個部分.
這裡我們說了.
這個比較重要.
這裡稍為重要一點.
我們解釋第十五節開始.
神的應許和律法到底有什麼關係或是抵觸呢?.
留意第十五節.

$^{1361}$保羅強調這是一個契約.
神定立了,包括人定立了.
都不能隨便違約.
何況是神和亞伯拉罕定立了呢?.
弟兄們,我且照著人的常話說.
即是一個很常見的習俗.
羅馬都是一個定契約比較重要的國家.
我們今天更加明白.
簽訂了的約是不能隨便改的.
就算多兒戲都不行.
都是叫違約.
弟兄們,我且照著人的常話說.
雖然是人的文約.
若已經立定了.
就沒有能廢棄或加珍.
即是說,他用一般性來討論.
給加拉太人有一種明白.
其實你們都明白的.
你們信了主,難道神還要收回這些東西.
要你們重新來過,當然不會這樣.
所以這裡說的是亞伯拉罕.
所以第十六節.
所應許的原事向亞伯拉罕和他的子孫說的.
神並不時說眾子孫.
指著許多人乃是說你那一個子孫.
指著一個人就是基督.
對不起,不是十九節.
是十六節才對.
眾子孫的那個獨特的子孫.
原來亞伯拉罕將來那個獨特的子孫.
就是耶穌基督.
所以萬國恩利的子孫蒙福.
就是指著耶穌基督.
在十二章第三節.
創世記.
神已經有救贖世人的計劃.
就是藉著他的愛子.
就是亞伯拉罕將來很多世代之後的孫子.
就是耶穌.
所以這裡說已經訂立了契約是不可以更改的.

$^{1401}$所以就是這位基督使世人能夠蒙福.
所以保羅現在就在說.
第十七節.
一來就是勉勵加拿大人不要受影響.
但其實更加用真理去駁斥那些傳這些所謂律法的人.
他們要留心.
其實神的聖經都不是這樣.
你不要扭曲了聖經.
十七節.
「我是這麼說」.
「神預先所立的約」.
就是和亞伯拉罕所訂立的約.
「萬國恩他得福」.
「萬國恩他的子孫得福」.
怕大家都不太熟悉.
不如我們記起十二章和十三章的《創世記》.
十二章.
「耶和華對阿伯難說」.
十二章第一節.
《創世記》.
果然在吾爾就已經跟他說.
「你要離開本地本族父家」.
「往我所要指示你的地去」.
「我必叫你成為大國」.
「我必賜福給你」.
「叫你名為大人」.
「你也要叫別人得福」.
「為你祝福的」.
「我必賜福於他」.
「那就助你的」.
「我必就助他」.
留意一下.
「地上的萬族都要因你得福」.
接著留意第十三章.
要揭後一點.
或者這裡不用了.
十四節開始.
十三章的十四節.
羅德離開了亞伯蘭之後.
「耶和華對阿伯難說」.

$^{1441}$「從你所在的地方」.
「凡你所看見的一切地」.
「我都要賜給你和你的後裔」.
「直到永遠」.
「我又要使到你的後裔」.
「如同地上的塵沙那樣多」.
「如果能夠數算地上的塵沙」.
「才能夠數算你的後裔」.
「你起來」.
「就中環走遍全地」.
他就在希伯倫的萬里象樹那裡足行.
在保羅眼中看見的後裔.
他再認真詳細地經過了十字架.
十字架的救恩和聖靈在他內心不斷地洗滌.
他告訴我們.
這個眾子孫不是說到哪一個.
而是說到耶穌基督.
這個就是《加拉太書》三章的十六節經文.
我們繼續看在筆記下面.
「我時這麼說」.
「神預先所立的約不能避那四百三十年以後的律法廢掉」.
計算就是亞伯拉罕.
對不起.
應該是由雅各落了埃及之後.
一直計算四百三十年.
就是他們要重新返上迦勒地.
在埃及.
才有摩西的出現.
才有神和摩西在西乃山訂立的西乃山之約.
即是舊約.
所以保羅就比較籠統地說.
就是四百三十年摩西的出現.
這個民族的起源.
來去立了這個舊約.
就是以安息日作為他們那個單位.
那個立約的記號.
所以這裡說.
神的應許是早於摩西都未出現.
雅各都未落埃及已經訂立了.
就是雅各的爺爺已經訂立了這個和神的契約.

$^{1481}$就是恩順清儀的契約.
所以不會因為四百三十年以後的律法.
去廢掉這個應許.
叫這個應許歸於虛空的.
第十八節.
因為承受產業若本乎律法.
就不本乎應許.
但我們可以肯定.
神是憑著應許把產業賜給亞伯拉罕的.
這樣說來.
律法為什麼有呢.
原是為過犯天常.
就是說.
律法只是會將人的需要的問題去指出來.
律法沒有那個能力.
讓那個人他看到自己的問題.
而他能夠改善.
這個律法的限制就是這樣.
所以是會將人帶到信心的那個層面.
那個限制就是將他帶到神的兒子舊恩那裡.
透過信心來引向那個應許.
當然就是聖靈.
就是現在我們在新約的代理人.
在我們生命裡面帶我們去看到我們人生的需要.
我們的限制.
然後賜給我們力量去完成神的計劃.
和進行神的旨意.
所以這裡告訴我們.
律法不是沒有功用的.
這樣說來.
律法為什麼有呢.
原是為過犯天常.
就是說點出我們的限制.
包括律法自己本身的限制.
就是給不到力量我們去做神要我們做的.
這個就是律法本來的意思就是這樣.
如果我們這樣去運用律法.
就恰如其分了.
但如果我們認為律法可以叫你做你就做到.
那就是說我們根本就沒有力量.

$^{1521}$亦都誤解了律法本來的限制.
律法其實是將我們帶到基督那裡.
使我們能夠在基督那裡得到力量.
來恩信清義.
和得到應許的聖靈.
去實踐神在我們生命裡面的救贖計劃.
所以你留意一下.
保羅這裡是真是幾可圈可點.
他現在在這裡真是幾有獨特的眼光.
來到他在他透過舊約新約這個臨界點.
一個整個民族猶太教在一個這麼大的包袱.
應該是否全部扔掉.
還是重新接受一個新的東西.
保羅沒有將兩件事二分了.
原來律法只不過是讓我們看到我們的情況.
就是說原來第十九節.
我再讀得清楚一點.
這樣說來律法是為甚麼有的呢?.
原是為過犯添賞的.
就是說律法其實它最終極能夠做的.
就是使我們知道我們的過犯.
它沒有力量令我們能夠可以.
因為明白了律法就不會再有過犯.
其實是不可能的.
所以這裡說.
是等候那蒙應許的子孫來到.
並且是藉天使經中保之手設立的.
但中保本不是為一面作的.
神卻是一位.
這裡有些竅口有些困難的.
就是告訴我們.
律法一直它的時代性和它的限制.
就是只會點出那個人在律法面前的本相.
其實就要帶你去看那個蒙應許的子孫.
在他那裡才能得到信心的意義.
和那個力量去行義.
就是這個意思.
所以原來律法又不是隨隨便便的.
因為如果我們這樣看上下文.
是在說律法的設立.

$^{1561}$這樣表面的理解就簡單一點.
並且是藉天使經中保之手設立.
經天使中保.
相信是摩西在西奈山神跡聖地上山40天.
畫一個法板.
神親自用指頭畫.
給他帶著兩個石板.
不是iPad.
石板帶下山.
頒布給以色列人.
所以是一個很嚴肅的.
你可以說是一個基督未來之前的過渡性過渡期.
是很嚴肅的.
經天使透過摩西去設立.
即是神親自頒布來設立.
雖然是他的過渡性.
但這個過渡性都是神允許的嚴肅.
因為那個子孫還未來到.
即是說神的救贖的兌現的時間未到.
人還在舊約的體制下.
所以做著那些憲制.
做著那些守安息日.
是無可厚非的.
來規範著人.
但是律法不可能繼續這樣來存在.
因為基督已經來了.
救恩已經完成了.
不可以再停留在律法的憲制.
以為可以解除我們的罪.
因為那個時候已經到了.
所以保羅說的現在如今.
是神的救恩計劃的開始.
這是一個相當重要的分水嶺.
所以我們需要有這些在舊約,新約.
介乎耶穌未成就救恩和耶穌成就救恩之後.
那麼重要的歷史時刻的這班人.
就是這些使徒.
就要幫我們梳理律法到底還有沒有功效.
在舊約耶穌未來之前.
律法有什麼功效.

$^{1601}$耶穌已經來了.
律法現在給我們有什麼功效呢?.
那就要很清楚地梳理.
就令到我們很清楚.
愛邦人我們恩順清儀.
聖靈會在我們裡面.
叫我們在愛裡面去行.
因為全律法都包在愛人如己這句說話裡面.
我們要用愛去待人.
就是這個意思.
就告訴我們神既然是一位.
就告訴我們律法有它的功用.
但是耶穌基督來.
律法的功用已經慢慢地減退.
因為聖靈會補足律法的限制.
本來律法就是叫人做好.
跟隨神.
但是律法沒有那個能力.
叫那個知道律法的人行到神的要求.
律法繼續存在.
但是聖靈進入信徒的內心.
就有足夠的力量.
如果我們順從.
就可以引向神的應許.
得著神祝福亞伯拉罕的福.
就臨到每一個信的人身上.
所以這段聖經.
講到關於律法.
其實主要是針對.
當時加勒太人過份地熱心去傳.
一知半解地去講.
要受割禮.
才能夠更加優越.
才能夠更加蒙神恩的祝福.
過份了.
淡化了耶穌的救恩.
所以一來不需要擔心.
因為信耶穌得到應許的聖靈.
聖靈就是真正的代理人.
在我們每一個信徒的內心.

$^{1641}$引導我們在神的要求下生活.
所以聖靈才是真正的新約代理人.
所以我們要傳遞聖靈的感動.
聖靈的工作.
在每一個信徒的內心.
如果我們傳遞的只是律法.
只是傳統.
就會越來越沉重.
變成一個很重的法典.
其實有些教會都很循規蹈矩.
慢慢要守.
不知為何要這樣守.
就變成不健康.
所以有時.
例如明仔讀白約翰.
聖公會的學校.
一年兩次要去聖公會的教堂.
所有學生都要去.
每年規定的.
但你對一個未有信仰.
未得救的中學生.
他去到見到的.
又不去解釋的儀式.
就未童絕肋不得已.
就更加對你這些儀式.
這些傳統.
就更加由疑惑變反感.
就更加令到他信主的信心.
未能夠堅固建立.
因為這個落差太大.
見到你這班人.
尤其是老人家在做這些事.
在跪下聽祈禱.
假如這班年青人.
這班少年人.
見到那些老人家又睡覺又聊天.
就更加反感.
本來是好事.
去教堂一年兩次.
但對一個未有信心.

$^{1681}$又是一知半解的中學生.
你叫他去看這些儀式.
其實如果你不經過很仔細的傳譯.
和性靈工作.
其實見到這些宗教的禮儀.
只會更大的誤解.
誤解帶來的就是反感.
更加覺得宗教就是這樣死咕咕.
就是這樣一成不變.
他根本完全不知道背後的傳統.
其實是需要解釋.
所以性靈就會在我們的心靈裡.
不斷讓我們透過認識真理.
讓我們能夠更加的豁達.
更加的信心增長.
所以我們如果強迫人做一些宗教儀式.
而我們不去解釋.
又不帶他去看到基督的愛.
那些宗教儀式其實是有毒素的.
是害他的.
是令他對基督教更加的卻步.
因為他見到的基督教.
只是這一套的傳統.
他見不到基督在這個傳統裡面.
見到這班人在這裡遊走.
在這些傳統裡面.
其實就是很大的問題.
其實第二代第三代.
他們回教會就是這樣.
他們就看到一班成年人.
很熱心的在那些傳統裡面.
每個禮拜這樣.
但是他們在這裡成長.
他們的衝擊很大.
需要的是基督在他們心裡.
需要聖靈在他們裡面進駐.
教導 引領.
他們一直長大.
他們的信心不斷的建立.
經過很多考驗.

$^{1721}$在大學裡面被人取笑.
等等各樣的東西.
他們就更加牢固.
他們再參與這些傳統.
他們就非常明白.
做這些傳統背後的意義.
原來是這麼真實的.
例如聖餐.
為什麼要這麼嚴肅.
大家一起慢慢走出去.
領聖餐.
原來背後是有這麼嚴肅的神學的.
對一個初信未信.
是很不明白的.
但是他越經歷了.
他就會越體會.
這些傳統原來真的有.
他承載的那個恩典.
那個愛.
那個信心在裡面.
所以很多東西是經過要詮釋.
所以重點是.
保羅不斷強調.
新約很簡單.
你透過耶穌.
你就能夠.
當然是他的交付.
你就可以和亞伯拉罕.
成為信心的連結.
就是他的子孫.
你就會得到相應的福.
那個福是什麼呢?.
應許的勝利.
進入我們每一個信徒的心靈.
這個就是.
亞伯拉罕的福.
就會一直延續下來.
所以基督就是.
那個成就.
神應許的代理人.

$^{1761}$他當然就是.
為我們死的代理人.
為我們復活.
在神面前有這個關係的.
復活的代理人.
更加是我們生命成長.
聖靈在我們裡面的代理人.
因為耶穌自己說的.
《約翰福音》十四章.
是不是第十六節呢?.
我要求父.
另外賜給你們一位保衛師.
就是另外一位.
和我一樣的.
因為我要走了,我要回天家了.
叫他永遠與你們同在.
就是真理的勝利.
就會在我們內心.
進駐的了.
所以透過姦太書.
這兩個大段落.
就讓我們看到.
稍微總結一點.
亞伯拉罕在這裡.
強調是以信為本.
他從來沒有走過律法.
而且.
基督就是那個代理人.
為我們成就了.
那個咒詛.
救我們脫離了.
律法的咒詛.
其實每個人都是面對.
律法對我們的追討.
但是基督已經為我們付了代價.
所以我們其實已經被救脫離.
律法.
被救脫離律法.
獲得勝靈.
所以在聖靈的新生命裡面.

$^{1801}$我們去遵行神的旨意.
所以不是無法無天的.
基督釋放了我們.
叫我們得以自由.
姦太書五章一節.
所以要站立得穩.
不要再被奴僕的騙去脅制.
我們得到自由.
但是這個自由就不是.
我們自己的自由.
這個自由就不是叫我們放縱.
我們.
大致上.
我將第四章的羅馬書.
和第三章的姦太書.
就做了一個完整的.
一個的詮釋.
就讓大家對保羅.
所理解的亞伯拉罕.
和他的行徑.
有一個再多一點.
充分的了解.
但是期望大家.
你自己要再由創世記十二章.
去到十八章.
亞伯拉罕這一段.
大概有二十五年.
這一段的人生.
都要首先.
療愈止長.
然後你再讀第四章.
因為這個某程度是.
加插下去.
因順清義.
不單止沒有脫離律法.
沒有違背律法.
而且是堅固神的律法.
不過這個堅固.
就是亞伯拉罕原來.
都是因順清義.

$^{1841}$這樣告訴我們.
其實舊約已經有因順清義了.
其實我們一直都是在.
這個因順清義的.
那個救恩的行列裡面.
所以我們不需要突然之間.
我們走去跟隨一些異端.
守一些的行為.
這樣來獲得.
那個救恩.
所以一切的敬拜.
都是隨心去做.
按真理去敬拜神的.
所以每個基督徒.
是有意識來教堂.
當然現在在線.
有些是疫情.
有些要上班.
要輪班.
是有意識地準備聚會.
所以真的不要說笑.
聚會前真的要有自己的安靜.
暫時真的不要說話.
不要聊天.
安靜自己的心靈.
是要有意識.
你要出席.
敬拜神.
是天大的事.
我們如果經常.
這樣訓練自己.
就很有這種.
的恩典.
因為你的心.
準備好.
等候神跟你的對話.
顯現.
聖靈的工作.
所以我們就拆開.
我們以往說過的馬丁路德.

$^{1881}$加爾文.
加爾文的聖餐論.
也是紀念.
但是聖靈在.
聖餐的聚會.
臨在聚會當中.
令守聖餐的人.
經歷到基督.
他的救贖.
所以不只是紀念.
紀念就很乾燥.
每年我們都要紀念.
抗戰紀念日.
沒錯的.
但是.
聖餐就不只是紀念.
加爾文就是.
更加走前一步.
否定了不是肉身.
不是變質.
更加不是同質.
不需要.
但就不是慈雲所說的.
只是紀念.
記憶.
不是這樣.
是紀念裡面有聖靈的運行.
在聚會當中.
所以大家麻煩.
真的要醒醒定定.
很認真.
有意識去守聖餐.
有意識去敬拜神.
是我們一生人生的成長.
很重要的一個.
操練和學習.
記得是對神.
我們去朝見神.
敬拜朝見神.
我們要停了.

$^{1921}$因為九點鐘了.
亦都是我今天剛剛說完.
第四章的羅馬書.
和迦勒納書的補充.
因為相近的時間.
保羅用亞伯拉罕.
去捍衛.
這個恩順清義.
和律法的關係.
第五章我們下個星期.
就會進入.
那個恩順清義.
那個關係上的.
復和帶來的那個祝福.
似乎第五章.
有六大方面的祝福.
我們可以一起去討論.
和去思考.
我和大家有個祈禱.
感恩我們能夠.
在第四章的羅馬書.
充分地掌握.
恩順清義.
在舊約已經存在.
而且不單止沒有抵觸到律法.
反而是堅固律法.
藉著使徒保羅.
在當時羅馬的教會.
寫這一封信.
很想告訴萬民.
不單止是羅馬的教會.
或者是猶太人的信徒.
都要在這個救恩裡面.
看到自己.
其實是和亞伯拉罕有關係.
就是因為這個信心.
而來到有這個生命的結連.
雖然都是幾千年前.
舊約的人物.
是主前的三千多年.

$^{1961}$但是這個信心.
一直運作在我們每一代的基督徒.
在在讓我們看到.
福音是關乎萬民的.
保羅不斷地說.
先是猶太人.
後是希利尼人.
所以很盼望主讓我們參與.
不同的方式的參與.
在普世宣教的運動.
我們不是叫口號.
而是我們很認真.
為一些在未信主的.
這些未得之民.
來為他們禱告.
或者我們能夠透過有限度的參與.
去無論是金錢.
或者是參與一些的短訊.
去體驗一下.
外邦未信的群體.
他們的需要.
願神幫助我們.
首先讓我們在教會裡面.
有一個合一的心靈.
大家一起向著神的國度的福音.
這樣來活著.
我們這樣仰望感恩禱告.
奉耶穌基督的名.
阿們.
\newpage



\section{羅馬書}
\label{sec:r2uz84DdreY}
\textbf{2025-06-05 周四晚查經班《羅馬書(12):分題:因信稱義的福》 講員:曾裔貴牧師}
\newline
\newline
連結: \href{https://youtube.com/watch?v=r2uz84DdreY}{\texttt{https://youtube.com/watch?v=r2uz84DdreY}} ~~~~ 語音日期: 2025-06-06
\newline
\newline
\hyperref[sec:23gUJZhHg2o]{\small{< < < PREV SERMON < < <}}
~
\hyperref[sec:index_chronic]{\small{[返順時目]}}
~
\hyperref[sec:index_scriptual]{\small{[返順卷目]}}
~
\hyperref[sec:G8aDRsOGytw]{\small{> > > NEXT SERMON > > >}}
\newline
\newline
$^{1}$好 和大英姐妹一起同心禱告.
再一次多謝天父.
晚上我們又能夠安靜地去研讀羅馬書.
多謝您讓我們能夠有一卷.
這麼有系統去講解福音的書卷.
我們很盼望藉著舍圖保羅.
讓我們看到神國的規模.
那個福氣.
以及我們能夠進入這一份這麼大的恩典.
我們去掌握,我們去享用.
更加重要的是我們能夠.
好像保羅所講的.
要有一個欠福音的債.
這樣的心情.
以致他盡力去將神的恩典與人分享.
我們很盼望主給我們有不同的機會.
有不同的學習.
在神安排我們的人生.
來很有力地,很專注地去傳講.
認識我們的人.
得到福音的好處.
今晚我們仰望在你的手中.
這樣禱告奉耶穌基督的名.
阿們.
簡單的一個再幫大家做一點回顧.
在經文裡面.
我們當然收集範圍.
由三章的二十一節.
就是神的義.
是甚麼人可以得到呢?.
就是賜給一切相信的人.
所以是說萬民都可以恩信稱義.
萬民都可以得到福音.
如果不是.
萬民都在神的陣路.
虧欠了,虧缺了神的榮耀.
但是我們今天會集中在第五章這一段.
當然第四章我們用了兩節.
講到其實舊約已經有恩信稱義.
這個運作.

$^{41}$這個概念.
以及神的安排.
保羅很巧妙地將亞伯拉罕.
他未改名之前.
已經神來呼召他.
給他有這種眼光.
就是他相信神.
他就可以得到他暫時未得到的.
因為第十五章的《創世記》.
他當時是完全未能夠生他自己的兒子的.
連這個以實瑪利都未出世的.
神就已經叫他去相信神的應許.
他相信神就以此.
就是以他這樣的信心.
這個行動為義.
所以保羅很巧妙地告訴我們.
恩信稱義的亞伯拉罕.
是在未受割禮之前的.
所以不關乎到亞伯拉罕他的為人.
他的好的行為.
這是羅馬書要告訴我們.
現在唯獨我們要效法亞伯拉罕的信心.
就是我們相信神.
這位能夠使死人復活的神.
我們就是相信耶穌從死復活.
將我們的罪帶進他自己的審判.
然後耶穌的復活.
就帶給我們能夠稱義.
這是第四章的結束.
所以我們看這段經文的時候.
第五章其實是否在說另外的東西呢?.
其實不是.
因為保羅用了一個希臘文的連接詞.
中文都已經反映了.
就是「所以」這個字.
所以我們的恩信稱義.
所以再一次告訴我們.
同三章的二十一節.
保羅所說的.
現在神在律法的以外.

$^{81}$來將信心賜給人.
就是信的人.
應該說信靠耶穌的人.
就可以白白的稱義.
可以得著神的義.
不是人自己換取的所謂公義.
而是神的義.
藉著耶穌的犧牲賜給信的人.
所以再帶我們看三章二十一節.
到四章二十五節.
這個大的段落.
就是神的義如何在地上運作呢?.
我們今天就會將第五章.
可能我們只能夠去到第十一節.
因為十二節開始.
有一個很強烈的對比.
就是舊的一個人.
就是阿當.
和新的人類的代表.
就是主耶穌.
做一個很大的對比.
就是恩典和審判做一個這樣的對比.
這個可能我們要下一節.
甚至下兩節才有機會講到.
所以我們今天看第五章.
我們就來留意.
上一星期也提過.
大概可以在這十一節聖經裡面.
我們可以找到六個很重要的福氣.
這一段經文是講.
「恩順清義在耶穌救恩裡面.
逃避了神的審判的基督徒.
得到什麼福氣呢?」.
如果今天晚一點有時間.
其實要將這一段.
和羅馬書十二章.
提到的各種在教會裡面的不同的恩賜的配搭.
來做一個簡單的對比.
所以基督徒的福氣和你獲得的恩賜.
某程度上是兩樣東西.

$^{121}$福氣是你個人與神的關係復和.
帶來的新的生命.
那種流露.
但是恩賜就是特別給某一類的人.
不同的恩賜.
來在教會裡面.
為著使教會群體能夠成長.
被牧養.
得到恩賜帶來的祝福.
所以兩者可能未必有太大的分野.
不過要留心的.
可能待會會講一點.
就是我們如果沒有了第五章的福氣.
在我們的生命裡面.
而我們只是以為擁有那些恩賜.
而我們去用那些恩賜.
這樣的人就要很小心了.
因為沒有這些福氣的人去運用恩賜.
可能只會不斷帶來的就是增敬.
炫耀自己的才幹.
還有未必是有心很想放下身段去服侍.
因為恩賜唯一的用途就是去服侍教會.
使教會得到益處.
當然教會就不是說建築物.
是那些人去得到益處的.
所以服侍背後一定是要謙卑.
一定是覺得自己是神的愛的人.
很想去用自己的恩典.
神給的恩賜去祝福其他人的.
所以如果沒有了這個態度.
你擁有的恩賜只會帶來競爭.
只會帶來炫耀.
你越厲害某些東西.
你越炫耀.
其實只會帶來更多的增敬.
你很厲害嗎?.
我用得著聽你的嗎?.
其實就不是聖經所說的真正的恩賜.
所以當說到今天這十一節的恩典的時候.
我們要記得.

$^{161}$是我們整個生命從一個罪人.
在神的憤怒之下的罪人.
而進入一個跟神關係復和的那種福氣的一份恩典.
這是比恩賜來得更加重要.
所以如果沒有了這樣的福氣.
享受這樣的福氣.
而我們去運用我們的恩賜.
那就要小心了.
其實就等於坊間的.
大家你爭我奪的那種競爭.
就不是聖經所說的恩賜本身的運用.
這些就是題外話.
這段聖經本身沒有叫我們這樣去聯想.
沒有這樣去跟別人競爭的.
這段聖經很寶貴的就是.
從和好開始進入一個更加終極的和好.
所以也是頭尾呼應的.
和好這個詞.
Peace英文是反映不了的.
應該希伯來文這個Salom.
就是那種關係上的復和.
我一邊說下去就會再跟大家去理解.
第五章我們沒有打PARPON出來.
可能我用經文去讀一讀.
因為十一節其實也挺長的.
做PARPON也做幾頁.
我等一下會逐點跟大家去做一個思想探討.
我先看看第一到第十一節.
也盼望大家如果能夠背.
盡量背.
因為這些經文是相當寶貴.
背了它對你自己的內心也有很大的福氣.
我們既恩信清義.
就藉著我們的主耶穌基督.
得與神相和.
我們又藉著他.
恩信得進入.
得進入就是第二個福氣.
得進入現在所站的這恩典中.
現在的站在恩典.

$^{201}$而且這個得以進入.
是由一個過去式一直延伸到現在的.
就是說一個信主的人.
與神一個隔絕.
進入到現在在恩典裡面.
還有這個字.
等一下我會說一點.
有些解經家就說這個字.
是一些臣子.
進見皇帝的那個字.
就是入宮面聖的意思.
如果我們中國人就明白了.
看很多宮廷的劇.
什麼時候入宮面聖呢?.
就是進入這個字.
當時希利文會用這個字.
告訴我們現在我們站在恩典中.
當然就是在神同在的福氣.
叫做恩典.
並且歡歡喜喜.
歡歡喜喜中文就翻譯不了這個意思.
應該加上歡歡喜喜的誇耀.
這個盼望.
所以你看到英文.
是Rejoice和Bose兩者合在一起的希利文.
所以是歡歡喜喜的盼望.
其實應該是歡歡喜喜的誇耀.
誇耀是誇耀.
就是誇口這個字.
所以等一下我們會說第二節.
誇口原來在保羅的書上說了很多.
似乎唯一這裡.
是很正面的叫我們要誇口.
誇口我們現在在這個盼望裡面.
什麼盼望呢?.
就是神的榮耀.
神的榮耀.
第八章又會再詳細說了.
所以你可以想象到.
原先我們是虧缺神的榮耀.

$^{241}$原先第一章說.
我們是在神的憤怒之下.
而且我們與神的榮耀有很大的虧欠.
但是現在我們基督徒.
因為我們站在這個恩典裡面.
我們可以誇耀.
我們可以誇口.
誇耀就變成很正面.
我們可以很誇口.
我們有一個日子.
就是盼望神的榮耀的彰顯.
第三節.
不僅如此.
就是在患難當中.
也是歡歡喜喜的誇耀.
原來可以在患難裡面.
向患難誇聖.
這些是中文的不同翻譯.
其實都是誇口這個字.
我們可以在患難裡面.
誇聖.
為什麼呢?.
因為知道患難生忍耐.
忍耐生老練.
老練生盼望.
盼望不至於羞恥.
所以第四個.
就是在患難裡面.
我們基督徒有一個很大的力量.
就是基督徒才有.
所以我們就要操練這些.
保羅告訴我們我們要操練.
為什麼我們可以歡歡喜喜地誇耀.
患難也打不倒我們.
因為知道患難就生出忍耐.
忍耐又會帶來更加有一種歷練.
老練.
英文也是experience.
就是有經驗.
不會很陌生.

$^{281}$老練生盼望.
又是回到原來一個很大.
似乎絕望的基督徒.
因為患難.
似乎絕望的.
但是原來患難會使我們生出更加真確的盼望.
更加是屬於我們在患難裡面的人.
才會有的盼望.
你可以想像得到.
所以我們基督徒為什麼不需要懼怕苦難呢?.
因為當然沒有人會求苦難的.
但是在苦難裡面.
原來終極會使我們有更加真確的盼望.
這個就是保羅告訴我們.
第五節說盼望不至於羞恥.
為什麼我們可以不需要有一種這樣的慚愧.
有一種這樣的羞恥帶來的可能是痛苦,絕望.
為什麼不需要呢?.
接著就是第五個福氣.
就是充滿了神的愛.
那怎麼來得著呢?.
因為賜給我們的聖靈.
將神的愛囂貫在我們心裡.
就是說我們常常有活水供應.
而這個活水就是神的愛囂貫在我們的心裡.
所以是在心靈裡面的.
不斷有活水的囂貫.
就是神的愛來囂貫.
這個囂貫讓我們在患難裡面.
不至於羞愧,羞恥,絕望.
這些全部都是福氣.
你想想連保羅也說了.
連患難都不會打倒一個基督徒.
一個在神的恩典裡面復和的基督徒.
這樣的時候.
就是說沒有東西可以讓我們產生一種悲氣神.
在恩典裡面的基督徒.
所以在恩典裡面去享受.
在教會與神相交的恩賜的運用.
才是真正的福事.

$^{321}$如果我們回來想演戲.
想做一些事情.
很快就會被人打沉.
別人對你的閒言閒語.
對你的諸多挑剔.
對你做的事情.
你用良心別人當作狗吠.
很快就會被人打沉.
因為你想回來在教會裡面做你想做的事情.
但是如果我們在這六方面能夠有滿足.
我們在這份愛的包圍下.
神同在的平安.
在我們的心靈裡面.
這樣的基督徒慢慢去運用恩賜.
你的恩賜就會祝福很多人.
因為你很知道自己都是無愛的人.
所以很多時候你看到一些人做一下.
突然間他會全低下.
原因可能就是他忘記了.
其實在恩典裡面去用恩賜服侍.
所以兩者是不能夠分割的.
而且如果沒有了前者.
今天我們說的這十一節的六個大的方向.
一個生命的那種豐滿.
而我們只是回來以為做一下事情.
演一下戲來用恩賜.
那就很慘了.
你那個付出帶來的收穫.
往往就是令到我們很挫敗的.
很不開心的.
甚至乎再嚴重一點.
連神都生氣了.
為什麼你的教會的人會這樣?.
那麼就連神都不要了.
就不回教會.
這樣是一個很可惜的.
出發點是很好.
但是忘記了.
我們就在神同在的恩典底下.
去服侍其他人.

$^{361}$還有多和少.
那個量是顯出還是隱藏.
這些反而不重要.
這些是神的安排.
最後的就是第六節.
「盼望不至於羞恥.
因為所恃給我們的聖靈.
將神的愛囂貫在我們的心裡」.
其實應該去到第十一節.
才是最終極和好的那份以神為樂.
接著保羅就講解一下.
這麼大的福氣怎麼運作.
第六節.
「因我們還軟弱的時候.
基督就按鎖定的日期.
為罪人死.
為義人死是少有的」.
因為你是義人.
為什麼為你死呢?.
為仁人死.
為了好的.
很有善,有人義的人死呢?.
或者有時候會的.
但是第八節.
「唯有基督在我們還作罪人的時候.
為我們死.
神的愛就在此.
向我們顯明了.
還是做罪人的時候.
就已經為我們死了」.
所以這一份才是神的真正偉大的神聖的愛.
明知而故意去死.
就是神的愛.
明知我們是敵對背叛和藐視神的人.
已經為我們死.
這真是一份很大的奇妙的愛.
第九節.
「現在我們既靠著他的血清怡」.
再一次講耶穌基督.
這位聖子為我們流血.

$^{401}$使得我們得到清怡.
「更要藉著他免去神的憤怒」.
因為這就是和好的描述.
第十節.
「因為我們作仇敵的時候.
且藉著神兒子的死.
得與神和好.
既以和好.
就更要因他的生」.
就是耶穌的復活.
來得救.
「不但如此.
我們既藉著我主耶穌基督.
得與神和好.
也就藉著他以神為樂」.
不是以神為一個取笑的對象.
以神為樂就是.
你想起與神的這種這麼甜蜜的.
在耶穌的愛的救恩裡面的交往.
帶來的就是我們想起神.
我們不會覺得陌生了.
我們會覺得很開心.
很平安.
很歡樂.
很有一種的滿足.
很有一種令到你的內心.
是心心靈的那種暢快.
例如我們很喜歡吃某一類的食物.
而讓你能夠享受到.
是最優質的.
你感覺到你的肚子裡面.
是一種很大的滿足.
就是那個口味.
最好味的那個狀態.
進入了你的口腹.
你會感覺到那一餐很豐富.
很滿足.
還有你很好的朋友一起去享受.
我們就覺得人生就是應該是這樣.
現在就是以神為樂.

$^{441}$就是我們進入常常在這樣的狀態裡面.
所以這個是一個終極.
我們能夠在一個和好裡面.
不斷去發展與神的一個這樣的關係.
能夠以神為樂.
這個就是一個很大的終極.
很多的信徒.
很多的其他宗教.
與神可能一年一次.
可能又用什麼方式.
戰戰兢兢.
舊約就是那些人.
大家聽過以利亞的年代.
那些拜巴力的先知.
是真先知.
巴力的先知.
不是假先知.
他們叫醒神.
因為要比較.
就要降水下來燒那些獻祭.
但是一直都沒有下雨.
以利亞就說.
你的神睡了覺了.
你快點做多些動作.
他們拿刀插自己.
好像玩神打一樣.
你可以想像到.
外邦的宗教.
他不會說到與神有這樣的一個相交.
當時保羅正正就是在羅馬帝國.
當時他們拜的神.
都是敬而遠之的.
所以保羅能夠說到一個基督徒.
其實是外邦人.
根本與耶和華無關.
但是他可以告訴我們.
保羅邀請外邦的信徒.
羅馬的信徒.
今天香港的信徒.
我們都在這個真理下.

$^{481}$我們在救恩裡面.
藉著耶穌.
我們不僅脫離了敵對憤怒的狀態.
而是在一種很大的歡樂裡面.
神是我們的滿足.
神是我們的喜樂.
神是我們的平安.
神是我們的一切.
這樣的一份心靈的心境.
所以大家要一起學習.
十一召裡面就是這樣的狀況.
就是一個在信心裡面得救的基督徒.
應該擁有的.
這個肯定是universal的.
無分國籍.
無分年齡.
無分性別.
無分階層.
你應該在救恩裡面.
都要享受.
由與神復和開始.
到終極.
神是你的喜樂.
這個肯定是每個人都要經歷的.
就沒有所謂的.
我不懂希臘文.
怎麼有這個學習呢?.
無關的.
只要你是有性靈.
神就可以在我們的生命裡面.
不斷的滋生祂那份恩典在我們內心.
令我們有這份滿足的喜樂.
好像大衛詩篇的十八篇一樣.
在神的面前.
祂有滿足的喜樂.
第十二節就突然間去到另外一個.
似乎又好像沒有什麼關連的.
為什麼突然間會說阿當呢?.
說阿當聖經新約只有兩段.
有沒有人知道另外一段是什麼?.

$^{521}$大家好少開咪.
不知道大家知不知道.
另外一段是說阿當的.
說得很詳細的.
就是哥林多前書十五章.
說到耶穌就是復活的.
末後的阿當.
當然這裡是在說.
舊有的人類帶來的審判.
刑罰.
死亡.
新的救恩的開始.
也是三章二十一節.
但如今的很重要的概念.
但如今神的義.
在律法的意外.
已經顯明了.
就是因信基督耶穌.
就可以白白稱義了.
所以來到第五章的結束.
保羅就將我們帶回去看.
耶穌基督帶來什麼更新呢?.
原來就是整個舊有體系.
人類的舊有體系.
阿當吃了分別善惡的樹果.
帶來的咒詛.
神沒有取消咒詛.
不過神興起了另一個恩典.
就是他的愛子親自來到這個世界.
將人帶進了救恩永恆的福樂裡面.
所以我們今天在主耶穌裡面.
我們不僅勝過了物質的死亡.
連我們的罪性.
我們都可以慢慢在救恩.
在聖靈裡面.
慢慢我們可以克勝我們的罪惡.
來得到神的更大的.
比我們整個新心靈的救贖的福氣.
所以原來這一段又不是多餘的.
這樣的比較.

$^{561}$又是有它的重要的存在的.
所以原來也等於告訴猶太人.
根本這個民族還沒有出現.
人類的代表就是阿當.
第一個人.
第一個人已經失敗了.
所有的人類都在咒詛之下.
但是現在在耶穌裡面.
所有的不管你是猶太人.
是外邦人.
都可以在救恩耶穌裡面.
有一個所謂的新的開始.
所以這個就是大概保羅在第五章的佈局.
第六章就多去探討罪性這個問題.
我們回來看第五章.
開始看我們現在這個部分.
仔細一點看.
在這個和好裡面.
很重要的.
因為猶太人的和好的概念.
和當時羅馬人看的和好的概念.
是兩樣東西.
我看看到沒有.
很快先看這個.
接著我們回頭看.
這裡要留意.
在羅馬的時代.
他們這個私恩的制度.
特別他們很強調.
很推崇的.
就是公元27年.
到大概主要180年.
連後世的歷史學家.
都是這樣稱為羅馬的承平的承世.
100年的承世.
在這個承世之下.
沒有城邦的戰爭.
大大小小的戰爭.
已經差不多沒有了.
帝國復歸於統一.

$^{601}$有好的羅馬的法律制度.
有好的貨幣.
好的貿易.
平安的道路網的通商.
和有不同的文化等等.
所以在歷史學家都說.
這是羅馬的承平的太平承世.
所以是帝國.
是帝王.
是該殺賜給人的平安.
沒有戰爭的平和的平安.
這就是羅馬人說的平安.
但是在聖經.
舊約猶太人.
他們說的那個salum.
那個民安.
那個平安.
不是說只是沒有病沒有痛的平安.
是說一個關係上的大家的一份平和.
一份大家的親密的平安.
所以聖經保羅告訴我們.
留意第一章第七節.
就說這個平安.
這個都很重要.
保羅似乎他的民安.
反映他的內心一些對真理.
對神的認識.
值得我們說.
我們很多時候都見慣了.
保羅的民安通常是公式化的.
願因平安.
從神我們的父.
和主耶穌基督歸於你們.
但是在羅馬書這裡就分外.
來得重要了.
因為第五章他說的平安.
正是根據第一章第七節這裡.
這樣來告訴我們.
就是這個和好.
所以第一章第六第七節.

$^{641}$他說我們從他.
就是從這位主.
這位耶穌基督受了恩惠.
並使徒的積分.
他自己都是得到恩惠.
Grace.
恩典.
都是這個字.
和使徒的積分.
所以兩者是要平衡的.
先有恩惠.
然後才有積分的恩賜的服侍.
所以如果我們只是回來要做事.
那你就真的搞錯了.
從來連保羅自己.
都不是說我要在教會裡面.
怎樣大展拳腳.
其實他自己個人.
蒙受了神的恩惠.
然後當然他有特別的積分就是.
在他人生.
神不僅不再追究他的罪.
反而賜給他有使徒的身份.
沒有跟耶穌都有使徒的積分.
那你可想而知他有一個特選的使徒.
真是很特別的一個人.
當然這麼特別一定會經歷很多的困難.
考驗,痛苦.
但這就是神給他人生的這個份.
所以我們如果透過這些.
我們大家一起去到這些人生的階段.
應該要想一下神給你的份是什麼.
你要忠心,你要專注.
去發展神給你的那個位份.
有些人就是專注做兒童的.
他專注.
他的專注慢慢都被人看到他的恩賜的.
就是他一看到小朋友就好像很貪心一樣.
所以你看到一個主學導師.
看到小朋友就板著臉.

$^{681}$他肯定不是那個份.
是被迫要這樣做的.
小朋友會感覺到.
雖然好像特別針對小學那樣.
其實不是的.
比較明顯的.
所以我們應該去到這樣的時分的人生.
其實不要騙自己.
什麼是神給你的份.
忠心去做,樂意去做,盡力去做.
因為你不做.
神會興起其他人.
你忠心去做.
你又會做得很開心,很喜樂,很滿足.
然後別人因為你的服侍.
他又會得到那份的祝福.
所以這個就是份.
積份.
不是每個人都是一樣的.
恩典就是每個人都是一樣的.
然後在萬國之中.
叫人為他的名信服真道.
這些真的對一個傳道人很大的勉勵.
神學生畢業.
通常就是要勉勵神學生.
我們的心要胸懷普世.
胸懷普世不等於你真的普世.
可能你在一個很小的教會裡面.
幾十年就是這樣度過你的人生.
但是仍然不要小看.
我們的心要胸懷普世.
在萬國之中叫人.
為他的名字信服真道.
所以我們是將人帶進耶穌的名下.
去歸服真道.
記得千萬不是歸服我們.
是歸服真道.
認識主耶穌.
第六節.
其中也有你們.

$^{721}$者蒙召.
屬耶穌基督的人.
每一個鼎姐妹都是蒙召的.
都是神在命運當中.
人類眾多人當中揀選我們的.
所以我們成為了屬於耶穌基督的人.
我們的人生應該歸屬耶穌.
我們應該要很滿足.
很誇口.
我們是歸屬於耶穌的.
我們是真的很想告訴別人.
我們是屬於耶穌的人.
第七節就是講保羅的自己.
我寫信給你們在羅馬.
為神所愛.
奉召作聖徒的眾人.
為神所愛.
我們又要苦心自問.
我們的信主.
是否真的感受到為神所愛.
我們是奉召作耶穌基督的聖徒的人.
神所愛.
然後感覺到自己被蒙召.
接著保羅留意一下保羅所說的話.
他說願因為平安.
這是保羅的祝願.
從我們的父神.
並主耶穌基督歸於你們.
這就告訴我們.
因為平安.
就像孿生兄弟一樣.
當然很少是平安因為的.
通常是因為平安的.
就是說神的恩典.
神的救恩.
臨到你的身上.
平安自然會跟隨你的.
所以保羅在每一次的問安.
都一定是從神那裡而來的.
因為平安的.

$^{761}$所以是有因為有平安的.
平安是在說和好.
和平的這個Salom.
新約很多時候都是用Salom這個字.
就是救約的那個字.
所以我們要看看保羅說的平安.
蒼西基有好幾段經文.
就告訴我們平安是怎樣用的.
我們可能要看幾段.
可能先看37章.
記得這是約瑟的故事.
37章開始.
就是那個采依事件.
37章我們不看那麼多了.
第12節.
就是約瑟被賣到埃及.
約瑟的哥哥們往事劍去.
放他們父親的藥.
以色列對約瑟說.
你哥哥們不是在事劍放羊毛.
你來 我要打發你往他們那裡去.
約瑟說 我在這裡.
14節就是這裡了.
以色列說你去.
以色列就是亞國.
以色列說你去看看你哥哥們.
平安不平安.
群洋平安不平安.
就回來報信給我.
這個平安不平安.
其實就是和好那個字.
就是SALOM那個字.
就是去到猶太人的關係上的一種的問安.
就不是說好像很富氣.
爸爸叫我來問候你.
問完就馬上走了.
就不是這樣.
就是一種希望兄弟們.
不要因為上一次的采依事件.
大家反目成仇.

$^{801}$反而就是.
以色列就是打發她的最愛的兒子.
約瑟就帶平安.
爸爸的平安去到那些現在正在放羊的.
在事劍的那個10個哥哥.
來希望有回這一份和好的關係.
所以大家要知道.
舊約這樣用平安這個字.
就不是說一聲就算了.
可以我們再看多一兩段.
譬如43章.
另外15章那裡就不需要看了.
去到43章.
40多章你們知道就是.
約瑟準備和他的哥哥相認了.
43章可以看後面.
因為怕時間不夠.
由26節開始.
約瑟來到家裡.
他們就把手中的禮物拿進屋去給他.
又俯伏在地向他下拜.
約瑟問他們好.
又問你們的父親.
你們的父親平安.
對不起.
你們的父親就是你們所說的那老人家平安嗎?.
他還在嗎?.
接著他們回答說.
你僕人我們的父親平安他還在.
表面中文看不到.
就是在說一份和好.
一份關係上的文安.
這樣的意思.
所以在《創世記》出現好幾段.
都是Salam這個字.
那種合意是很深的.
猶太人真的有感受到這種關係上.
所以他們親咀文安.
按舊約的傳統都很合適的.
真的是一份關係.

$^{841}$所以有時候你看到中東人.
都真的這樣見面.
大家抱著輕輕的.
來一個文安親咀.
他們可能已經是很習以為常.
中東的人.
但是我們亞洲人就比較含蓄.
西方就算有.
都不是這樣的情況.
所以我們回到舊約這些經文的應用.
到保羅用的這個Salam.
就不是當時的世界.
那種的沒有事沒有幹.
沒有戰爭.
有帝國的庇蔭.
就是平安.
沒錯 那個當然是一個好的.
一個太平的盛世.
絕對不是戰爭.
但是猶太人.
或者保羅用的平安.
是牽涉到關係上的一個復和.
所以我們回來看第五章.
這個和哪一個的復和呢?.
肯定是告訴我們.
我們恩順清二.
帶給我們的.
就是藉著我們的主耶穌.
同神有這個Salam.
與神Salam.
就是同神有這個關係上的恢復.
所以這個恢復.
是恩順以至於信.
是不斷的可以去到終極的.
以神為樂.
同神有一個很密切的愛的關係.
就是好像水與膠肉的一種狀況.
所以這些是我們應該經年累月.
每天早上.
當然不一定要早上.

$^{881}$享受你自己獨處.
親近神.
閉上眼睛.
與神聊天.
表達你對神的愛.
表達你對神的渴慕.
每天這樣操練.
不是去練功.
不是去修練.
是每天有一個時間.
是屬於你與神的.
這樣一個很密切.
很溫馨的一種愛的關係.
神人的神性的愛的關係.
我們就會慢慢進入這個和好裡面.
沒有罪的刑罰已經是一個解決了憤怒的.
但是不是停留在那裡.
所以仍然不斷的好像當作是開玩笑.
不要洗禮就不再回教會.
那個開始你明白了.
在契約底下.
我們已經與神立了一個約.
我們是在這個地上屬於神的.
我們是神的產業.
我們要怎樣來到地上榮耀神.
但是我們個人應該發展與神的關係.
我們每天沒有人強迫我們.
我們很想去到神的恩典的寶座.
去享受祂的同在.
這個就是我們的默契.
所以每個人都應該要發展.
就算你是在月球上的太空站的太空人.
你更加可以去發展與神的相交.
其實你離開天堂更加近.
我們在地球都很密集.
當然是開玩笑的 沒有地域限制.
是心靈的限制.
所以我們閉上眼睛.
只不過是代表我們更加專注.
在家裡你更加專注.

$^{921}$或者你午餐的時候.
在你自己的座位吃完飯.
你有五分鐘大家自己安靜在神面前.
總之是屬於你與神之間的一個很好的關係.
所以這裡是特別保羅告訴我們.
好像神邀請我們記得進入這個相和的關係.
就在神的同在下.
去與祂有這份平安的相交.
這是第一個福氣.
就是和好.
進入了這個關係裡面.
第二節保羅好像告訴我們.
原來可以從從進入.
層層的進深.
告訴我們.
現在我們不再遙遠.
所以基督徒.
你對神的稱呼.
你不要騙自己.
如果你不敢閉上眼睛向神祈禱.
你甚至不享受這個祈禱.
你是陌生的.
好像你完全是口吃一樣.
這是一個問題.
因為這段經文.
是邀請每一個基督徒在和好的情況下.
還要進入.
得以進入.
就是說.
本來是離開.
現在是慢慢進入到神同在的恩典裡面.
就是說.
上上要我們閉上眼睛.
進入天堂.
在神的同在下.
這樣的享受.
所以你怎麼可能.
還不懂得去祈禱.
你不敢自己與神說話.
叫神的名字.

$^{961}$稱呼祂是你的天父.
如果你都不敢這樣祈禱.
證明有些事情需要去克服.
這告訴我們.
得以進入.
剛才說過.
這個字本身.
《解經》告訴我們.
是有一些神伯進見皇帝的字.
如果是這樣的時候.
保羅刻意用一些入宮綿盛的措辭.
來做他的修飾.
就告訴我們.
等於真的與大相關.
就告訴我們.
如同我們本來是沒有資格的.
現在能夠隨時去到神的恩典的寶座.
所以私恩寶座是我們可以隨時進入的.
去到神的面前.
所以我們真的要去學習操練.
我們在神面前的禱告的服事.
以及我們要常常去到恩典裡面.
在恩典裡面.
我們就會感受到我們是蒙恩的人.
這就是第二節.
得以進入.
就是第二個福氣.
當然.
其實是一個大的救恩的福氣.
不過可以這樣來到.
保羅將我們這個大的福氣.
來到分得很細緻地告訴我們.
我們在什麼狀況.
全部在說的是屬靈的福氣.
是生命的福氣.
是新的生命的福氣.
這個新的生命的福氣.
足以去令到我們有力量抗衡地上.
因為我們還是活在地上的.
我們還是在舊有的亞當體系底下生活的人.

$^{1001}$除非有一天.
就是我們安息主懷.
我們就進入真正的新生命的永恆.
我們就不需要再對抗舊有的.
無論是罪.
世界的困苦.
我們就不需要再面對.
所以今天我們就是在恩典中去面對地上.
形形色色的苦難.
有些是肉身的苦難.
有些是心靈的苦難.
有些是心靈肉身的苦難都有.
每個人都有不同.
但是無論如何.
我們要留意.
今天我們得以進入.
是第二個福氣.
Gained access.
可以進入到神的恩典裡面.
是要站在恩典裡面.
保羅在說的形容.
「現在所站」.
強調.
現在就已經可以站在恩典裡面.
即是說.
就算你現在在甚麼光景.
你都可以.
所以一個轉念.
我們就可以在神的恩典裡面.
來到去.
這裡提到歡歡喜喜盼望就是一個誇耀.
歡喜地誇口.
歡喜地誇耀.
是帶著歡喜的.
不是得戚 不是一種自誇.
驕傲 不是那些.
這我們明白.
但這個就是在恩典裡面.
誇耀神.
而且是在說歡歡喜喜的誇耀甚麼呢?.

$^{1041}$神的榮耀.
誇耀神的榮耀.
這個就是歡歡喜喜地誇耀.
神的榮耀.
即是說本來.
我揭開之前.
怕大家有時會忘記了.
第一章第十八節.
原來神的憤怒從天上顯明.
在一切不奇不異的人身上.
就是那些行不義.
阻擋真理的人.
神的憤怒.
然後去到第三章.
我們這班人本來是在.
祈求神的恩典裡面.
但現在.
我們這班人.
本來是怎樣的呢?.
三章的第二十節.
所以凡是有血氣的.
所有人.
地球的人.
會有一個因行律法.
能夠在神的面前稱義.
因為律法.
本是叫人知罪.
到第七章再強調.
再詳細講解.
律法的功用.
第七章會再講的.
二十一節.
但如今神的義.
在律法以外.
已經顯明出來.
有律法和先知為證.
就是神的義.
因信耶穌基督.
加給一切相信的人.
並沒有分別.

$^{1081}$留意第二十三節.
因為世人都犯了罪.
留意一下.
我們是規缺了神的榮耀.
因為我們的罪.
使我們得罪了神.
神不可以在我們受眾物身上.
得到榮耀.
我們不單止沒有去尊榮他.
反而我們規缺了神本來的榮耀.
他創造的人類.
不應該這樣自相殘殺.
這麼多的默辱.
遠離神.
但是在這個規缺下.
保羅很快告訴我們.
我們可以歡喜的盼望神的榮耀.
神的榮耀就是我們可以跨口.
可以去等待.
而且是在一個歡喜的心靈裡.
去等待,去盼望的一個福氣.
所以這是第三個福氣.
接著.
不單止這樣.
不但如此.
我們也要去意識到.
人類.
無論他自己的時代.
或者他自己的切身處境.
或者舊約.
新約的人類.
在罪的底下.
在死亡的底下.
在神的救助底下.
其實我們是無奈.
是無望的.
但是不單止這樣.
這個恩順清義帶來的福氣.
就是不單止我們不需要再無望.
不需要再痛苦.

$^{1121}$不需要再有這種壓抑.
不單止這樣.
在患難當中.
開始引進一個很現實的處境.
就是所有的人.
都一定會有大或小的患難.
原來患難神都容許我們經歷.
就是使我們更加成長.
基督徒.
我們不是要去追求患難.
但是患難是必須每個人都會面對.
但是神讓我們在患難裡面.
份外感受到.
他同在帶來的福氣.
留意第三節.
不單止這樣.
就算在患難當中的人.
基督徒也是歡歡喜喜的.
培肅這些經文.
保羅說在患難當中.
我們可以歡歡喜喜.
可以歡歡喜喜地誇耀.
可以誇耀我們在患難當中.
都可以有一種滿足.
為甚麼呢?他的答案就是.
因為患難只會使我們更加忍耐.
在患難當中的基督徒.
經過這樣的訓練.
經過這樣的人生的閱歷.
是會使我們更加有一份耐忍.
是一種慢慢的對事物的透徹的認識.
突然想起.
楊牧師.
他曾經患了癌症.
他在癌症期間.
他說他更加深體會神學.
更加深體會神.
更加深體會人的渺小無助.
這告訴我們照他兒子的模樣.
使他的兒子在許多弟兄中作長子.

$^{1161}$就是我們的元首.
預先所定下的人.
又照他們來.
我們是被神選照.
所照來的人.
又稱他們為兒.
我們成為神稱兒的人.
稱為兒的人又叫他們都得到榮耀.
所以你見到是一個盡心的關係.
聖靈由我們不懂祈禱.
不敢和不肯祈禱.
帶我們進入成為神的榮耀.
叫我們得到榮耀.
就是這樣的一個狀況.
你可以想像到.
在保羅的羅馬書講的.
救恩的恩信稱義.
聖日三一的童宮.
是緊密沒有間斷的.
主耶穌來這個世界.
成就了和平.
成就了救贖.
但原來聖靈.
成為另一個保衛師.
另外一個主耶穌.
在我們裡面.
長遠的在我們裡面.
不斷為我們有生之年.
在我們裡面.
不斷為我們有生之年.
在我們裡面.
不斷為我們有生之年.
在我們裡面.
不斷為我們有生之年.
在我們裡面.
不斷為我們有生之年.
在我們裡面.
不斷為我們有生之年.
在我們裡面.
不斷為我們有生之年.

$^{1201}$在我們裡面.
不斷為我們有生之年.
在我們裡面.
不斷為我們有生之年.
在我們裡面.
不斷為我們有生之年.
在我們裡面.
不斷為我們有生之年.
在我們裡面.
不斷為我們有生之年.
在我們裡面.
不斷為我們有生之年.
在我們裡面.
不斷為我們有生之年.
在我們裡面.
不斷為我們有生之年.
在我們裡面.
不斷為我們有生之年.
在我們裡面.
不斷為我們有生之年.
在我們裡面.
不斷為我們有生之年.
在我們裡面.
不斷為我們有生之年.
在我們裡面.
不斷為我們有生之年.
在我們裡面.
不斷為我們有生之年.
在我們裡面.
不斷為我們有生之年.
在我們裡面.
不斷為我們有生之年.
在我們裡面.
不斷為我們有生之年.
在我們裡面.
不斷為我們有生之年.
在我們裡面.
不斷為我們有生之年.
在我們裡面.
不斷為我們有生之年.

$^{1241}$在我們裡面.
不斷為我們有生之年.
在我們裡面.
不斷為我們有生之年.
在我們裡面.
不斷為我們有生之年.
在我們裡面.
不斷為我們有生之年.
在我們裡面.
不斷為我們有生之年.
在我們裡面.
不斷為我們有生之年.
在我們裡面.
不斷為我們有生之年.
在我們裡面.
不斷為我們有生之年.
在我們裡面.
不斷為我們有生之年.
在我們裡面.
不斷為我們有生之年.
在我們裡面.
不斷為我們有生之年.
在我們裡面.
不斷為我們有生之年.
在我們裡面.
不斷為我們有生之年.
在我們裡面.
不斷為我們有生之年.
在我們裡面.
不斷為我們有生之年.
在我們裡面.
不斷為我們有生之年.
所以保羅說如果是這樣的時候.
還有甚麼可以說的?.
神這樣為我們幫助.
這個Hepa其實應該是.
for us,他站在我們這邊.
神如果站在我們這邊.
誰還可以抵擋我們?.
神既然不愛惜.

$^{1281}$祂自己的兒子.
為我們眾人都捨去.
豈不是把萬物.
和祂和祂自己.
神自己一同白白賜給我們嗎?.
誰可以控告神.
所揀選的人呢?.
有神稱他們為異.
誰能夠定他們的罪呢?.
有基督耶穌已經死了.
而且從死裡復活.
現今在神的右邊.
也替我們祈求.
這麼詳盡描述耶穌.
由死到復活.
到現在在天上.
神的寶座的右邊.
正在為每一個基督徒祈禱.
原來聖神的三一.
聖父,聖子,聖靈.
是如此緊密地配合.
為聖徒的生命成長.
來效力的.
聖靈就在我們地上的基督徒.
效力使我們最終得到神的榮耀.
然後耶穌.
為我們不被控訴.
來禱告.
為我們的生命.
能夠不斷在神的裡面.
得到那個豐富來禱告.
因為他的意思就是.
誰可以定我們的罪.
保羅的答案就是.
基督耶穌為我們死.
償付了那個罪債.
然後他死裡復活.
就是現在活著為聖徒禱告.
這個就是聖經.
希伯來書也有說.

$^{1321}$所以羅馬書這裡.
先這樣再說耶穌現在的身份.
那個職份在天上.
原來是一直在為.
聖徒的不被控訴來禱告.
你可以想像得到.
聖父,聖子,聖靈對我們的那個祝福.
就是讓我們的生命成長.
所以今晚你什麼都聽不到.
你要聽到就是.
聖父,聖子,聖靈.
三一真神為我們每一個信主的人.
預備了一個套餐.
就是要讓我們得榮耀的套餐.
我們的生命要是榮耀的.
所以真的.
先看看誰先走快一點.
最後一口氣.
你應該跟所有人告別.
多謝大家.
我現在要去見主了.
因為祂在我生命裡面.
令我得到榮耀.
現在我就進入那個榮耀.
享受神的榮耀.
這樣跟大家說再見.
就是一個地上榮耀神的基督徒.
當然是痛的.
當然戴著面罩.
戴著氧氣罩這樣掙扎的.
但是真的可以這樣誇耀自己的勝利.
這個就是救恩帶給我們的生命的豐富.
經文是這樣告訴我們.
經文是一而再地說這些真理給我們知道.
所以就是救恩帶給我們的生命的豐富.
將來我們第八章再說.
第五章這裡就是聖靈囂冠神的愛.
在我們裡面成為了一個中介人.
如果你熟而不所書.
聖靈更加是我們將來在天堂.

$^{1361}$得到基業的一個憑據.
憑據那個字原文就是訂金.
神在地上.
在我們每一個聖徒的生命裡面.
祂已經給了我們一個承諾.
你將來來到我的天家.
你會有產業的.
而且祂用聖靈作為訂金給了我們.
就是說我將來一定會買贖你.
回到天堂.
所以我們每一個基督徒.
我們是肯定將來進入那個永恆的天家.
死亡是我們離開那個物質的世界.
就進入永恆.
那個屬靈的真正的新的世界.
神的家.
那裡享受那個永恆的福樂.
原來我們本身就是神的產業.
神已經買了.
祂現在下了訂.
聖靈就是那個訂金.
這是二不所書一章那裡告訴我們的.
所以我們能夠有聖靈就是充滿了愛.
充滿了的是神的愛.
所以神的愛是肯定要驅逐那些黑暗的.
神的愛肯定會驅逐那些仇恨的.
神的愛就是會讓我們學習彼此相愛的.
所以我們學習不要挑剔.
學習不要對人總是有一些苛刻.
我們要有愛包容.
我們學習從另對方的角度去看事物.
我們就會不知不覺學習到有愛.
記得要學這樣學的.
教會很容易就給了很少的東西.
很多人就說我要據理力爭.
其實有時候所謂據理力爭.
是我們自己的內心的那些個性.
脾氣.
舊有的舊的自我.
其實是不必要的.

$^{1401}$終極的就是剛才說的.
以神為樂.
我們不如讀那段.
還有一點時間.
我們讀後面那段.
第六節這一段.
因為我們還是軟弱的時候.
第八章剛才二十六節也是說軟弱的時候.
不過這裡就說是我們得救恩的時候.
那裡是我們連去賀武神的旨意的時候都沒有.
但是現在這個是進入救恩的門的基督徒.
因為我們還是軟弱.
還是敵黨神的時候.
基督就是神差遣而來的愛子.
按照所定的日期.
記得.
保羅很著重這件事.
加泰書四章三節.
記不記得.
及至時候滿足.
神就差遣他的兒子.
為女子瑪利亞所生.
而且生在律法之下.
他能夠將我們在律法之下的人拯救出來.
巴黎大學四章三還是四節.
所以原來在神的眼中.
有時候日期的.
基督的降生.
基督的犧牲.
是按照著神所定的日期.
為罪人死的.
記得是為世上每一個罪人.
就是說你只要承認你是一個罪人.
你向神呼求就可以得救了.
這就是信耶穌的最基本.
所以是為罪人死.
所以每一個基督徒.
應該看到自己的罪.
是多麼得罪神.
但是又多麼蒙了神的愛.

$^{1441}$基督為我們犧牲.
所以保羅就再描述.
就是說似乎針對羅馬人.
因為當時羅馬人有一個恩師的制度.
就是私恩主.
大家有沒有聽過Patreon.
Patreon.
現在如果YouTube.
你要入會做會員.
那些就叫做Patreon.
這些就是羅馬人的制度.
就是有些叫做私恩主的制度.
因為羅馬人有很多奴隸.
那些無論地主或者富有的人.
他就會有很多奴隸.
那就是他們的Patreon.
那些奴隸就是他們的一個.
我忘記了另外一個名字.
我要想一下.
這些Patreon的制度.
就是羅馬書十六章.
那個叫做菲比.
就是私恩主.
就是保羅的私恩主.
留意一下保羅怎樣稱呼這個菲比.
英文就是Patreon那個字.
第一節.
我對你們舉薦我們的姊妹菲比.
她是堅革理教會的一個女執事.
請你們為主接待她.
合乎聖徒的體統.
她在何時上要你們幫助.
你們就幫助她.
留意一下這裡.
因她數來幫助許多人.
也幫助了我.
這個幫助這個字就是Patreon會員制的那個字.
就是她是保羅的一個恩主.
我記得了Patreon和Kaien.
就是如果你是一個奴隸.

$^{1481}$你有一個恩主.
他會給你好處.
那個好處有時未必是金錢.
他成為你的恩主.
就是這個Patreon.
你就要怎樣呢?.
羅馬當時的制度就是.
如果當他想要得一些功績.
你就要幫他.
例如你要幫他投票.
要有一些利益輸送.
是當時的一個制度.
保羅就用這個.
提到原來保羅的恩主.
都是這個菲比.
他幫助過許多人.
所以就叫其他的羅馬教會.
將這個菲比.
他遲些會來你們羅馬教會.
你們要按聖徒的體統去接待他.
因為他是很多人.
在很多福音工作上的恩主.
那個Patreon.
所以我們就要學習這件事.
保羅不是想強調.
這種因果對照的利益輸送.
而是保羅很想告訴我們.
我們今天最大的恩主就是耶穌.
我們就是他的Client.
我們就是在神的無條件的救贖裡.
得到這個任避.
所以我們如果在古代羅馬.
那些人能夠成為了哪些人的Client.
他都尚且要為那個恩主.
私恩的主去效力.
要做很多可能未必是他心意想做的事.
但是最大的私恩主就是耶穌.
他現在為了我們去犧牲.
當然在羅馬的世界是不可能這樣發生.
所以保羅說為二人死是少有.

$^{1521}$為好的人死或者有感激的.
但是保羅說不是這樣的輸送制度.
一個這樣的Patreon和Client的制度.
而是我們從來都沒有想過人間都不會有的.
就是唯有基督在我們還作罪人的時候.
就為我們死了.
所以沒有可能一個私恩的主.
為一個Client本來完全不值一舊的.
這樣為他犧牲.
羅馬的臣子都不會這樣做的.
只會用來供奉.
但是現在告訴我們主耶穌無條件為罪人犧牲.
是用他的生命為罪人犧牲.
所以我們既然靠著他的血清義.
所以血帶來的真正的清義.
洗乾淨我們的生命.
我們的心靈.
再沒有罪的追究和控訴.
所以基督徒要有平安.
就是用寶血面對這個控訴.
寶血帶給我們潔淨.
寶血帶給我們去遮蓋我們的羞恥.
我們犯罪的時候那種不好.
已經是透過血清義了.
這個是在神的面前.
就好像古代猶太人憲制.
在寺廟所彈血.
其實很髒的.
你想一下你的血是這樣彈的.
在家裡在廚房你斬一隻雞.
不斷地彈.
你先生就說你幹什麼.
很髒的.
至聖所就是這樣彈血.
但原來他們覺得這個髒.
在神面前就是清義的遮蓋.
所以這些概念不是猶太人.
我們未必明白.
但是現在這裡很清楚了.
主的寶血洗乾淨我們的罪.

$^{1561}$是應該連罪咎感.
羞恥感都應該要除去.
更要靠著他.
免去神的憤怒.
為什麼呢?.
因為我們作酬藉的時候.
借著神兒子的死和好.
既然和好就更要.
因為他現在是復活.
現在他還是在神面前.
他復生了.
我們就得救了.
所以這個得救.
應該就是有一種很強大的新生命的力量.
在地上展現的一種新生命.
基督徒不應該好像要充電.
要打強心針.
而是我們應該有一種常常慢慢.
感受到救恩帶給我們的生機,生氣.
所以第十一節就是終極了.
不但如此.
所以說完這裡就剛剛完結了.
不但如此,我們的藉著我主耶穌基督.
得與神和好.
就藉著他以神為樂.
終極.
由一個和好進入.
去到與神有一個很緊密的相交.
愛神.
神愛我,我愛神.
甚至乎不為什麼.
沒有什麼福氣.
我都仍然愛神.
有困難,更加愛神.
更加倚靠神.
更加思想一下.
耶穌給我受的難.
更加多十倍,千倍.
但耶穌都願意為我這個罪人死.
赤裸裸知道我這個人的本性.

$^{1601}$他都願意為我釘十字架.
這份愛是無懈可擊的.
所以我們想想十字架的主.
為我們犧牲.
我們就會感受到.
神的愛源源不絕地進入我們的內心.
今晚就說到這裡.
我們一起禱告.
再一次感恩.
天父給我們能夠很仔細地.
去將這六個神給我們的救恩.
帶來的生命的福氣.
與神的交往.
帶來的生命力.
那個在患難裡面.
都仍然展現一種很大的盼望.
甚至保羅說.
是歡歡喜喜地誇耀.
幫助我們,誇耀我們的困難.
誇耀的是我們的救主.
你的愛的力量.
誇耀的是我們有永恆的盼望.
誇耀的是我們可以在神裡面.
我們站在恩典裡面.
享受你的同在.
享受你那個環繞我們這一份的力量.
愛的那個包圍.
幫助我們發展我們那個口對你的頌讚.
不是一種好像不拉肩的一些讚美詞.
而是我們有經歷的.
是心靈裡面的一種讚美.
對神的榮耀的讚美.
尤其是做物主.
根本就是與我們人類不相同的.
因為你是永恆的.
你是無始無終的.
你是自有永有的.
讓我們對你有認識.
以致我們的信仰是更加堅穩.
更加的扎實.

$^{1641}$多謝主聽我們感恩禱告.
奉耶穌基督的名.
\newpage



\section{羅馬書}
\label{sec:G8aDRsOGytw}
\textbf{2025-06-12 周四晚查經班《羅馬書(13):分題:克勝罪惡的新生命》 講員:曾裔貴牧師}
\newline
\newline
連結: \href{https://youtube.com/watch?v=G8aDRsOGytw}{\texttt{https://youtube.com/watch?v=G8aDRsOGytw}} ~~~~ 語音日期: 2025-06-13
\newline
\newline
\hyperref[sec:r2uz84DdreY]{\small{< < < PREV SERMON < < <}}
~
\hyperref[sec:index_chronic]{\small{[返順時目]}}
~
\hyperref[sec:index_scriptual]{\small{[返順卷目]}}
~
\hyperref[sec:7LZOyo5Zt3c]{\small{> > > NEXT SERMON > > >}}
\newline
\newline
$^{1}$我們一起同心禱告.
再一次感恩 我們可以再去研讀《羅馬書》.
幫助我們能夠一路進心去到第五章的時候.
都差不多去到經文的心臟很重要的部分.
幫助我們能夠明白 清晰地拿捏.
以致我們知道整個救恩的發展.
和神是怎樣來籌劃這個救恩.
使到全人類可以得福 求主幫助帶領我們.
我們晚上都可能會有疲倦的.
加添我們身體心靈的力量.
我們這樣仰望禱告 奉耶穌基督的名.
阿們.
第五章上個星期都很快地講了.
六個在基督裡面新的生命.
每一個基督徒都應該可以體驗的新的人生.
當然就是聖靈進駐在我們裡面.
帶來的福樂和平安.
我們上個星期都有提過.
第五章這裡講的一到十一節的福氣.
和第十二章講的恩賜上面的配搭.
是兩樣東西來的.
這一份的恩典是每一個基督徒.
都應該共同地擁有.
和一生都要追求.
這個和我們是否服事得好與不好是沒有關係的.
十二章那裡我們領受了恩典.
我們在這麼大的救恩裡面.
我們應該要把身體獻上.
成為一個活祭來給神使用.
那裡就是神會把不同的恩賜給個別的基督徒.
按你所得的恩賜去忠心地服事.
所以兩者是不等同的.
所以大家要留意.
這是第五章.
對不起 還有一點感冒還沒清理好.
今天講的第十二到二十一節.
某個程度上是相當重要的.
因為我們一開始講福音.
講到一開始第十六,十七節的一章.
就是講了那個結果.

$^{41}$就是哪些人可以得到神的恩典呢?.
這個福音是神的大能.
要救一切相信的.
首先是猶太人 後來就是希利利人.
因為神的意義在這個信心上.
來彰顯這個信是因信而至於信.
這個結果來的.
就是說信的人就可以得到那個的清義.
但是十八節開始就講.
哪些人需要救恩呢?.
就是萬民都在神的憤怒底下.
在神的咒詛底下.
就一路列舉世人的罪.
有一定的道德的人的罪.
猶太人的罪.
還有神就已經宣判.
普世的人都服在神的律法底下.
都是被定為有罪.
但是還沒有解釋到萬民為什麼會有罪呢?.
為什麼萬民都需要救恩呢?.
所以今天第五章的十二節.
保羅不單止好像第四章那樣.
把亞伯拉罕救約的這位重要的人物.
他有什麼特別呢?.
就是原來在救約的時候.
沒有律法.
沒有這個西乃的立約.
沒有摩西的出現.
已經有恩信清義.
這個這樣的教訓了.
原來亞伯拉罕都是在神面前.
不是因為他的行為.
是因為他的信心.
相信神而得到清義.
所以原來救約就已經有這樣的行使信心.
而得到神恩典的人了.
所以救恩是給每一個民族.
不只是給猶太人.
更加不是因為你是亞伯拉罕的子孫.
受割禮.

$^{81}$就可以自然的成為神的兒女.
這個就在第四章.
已經很清楚的交代了.
接著第五章就說.
一個在救恩裡面的人.
可以享受的六方面的福氣.
很快再和大家溫習一下.
然後十二節開始.
就是很重要的重點.
保羅帶我們去看人類的起源.
就是當這個人.
原來是人類的一個總代表.
所以原來我們就明白.
為什麼萬民服在神的憤怒底下.
因為原來.
保羅告訴我們.
阿當的犯罪.
罪就進入到人類裡面.
就帶來全人類都要面對.
那個就坐.
那個神的審判.
所以怎樣可以帶我們離開這個就坐呢?.
當然你沒有出過世.
你沒有存在這個地球.
那你就不用就坐.
根本就沒有你存在就什麼坐呢?.
但現在說的就是.
所有眾生在這個地球出現的人.
都要面對阿當這個就坐.
那怎樣可以解除呢?.
所以今天這一段就很重要.
真的給我們拆解了.
人的最根本的問題.
好像我們人類現在越來越先進的科技.
已經可以描繪一個基因圖譜.
不過還沒有開始可以去分析.
每一種病的致病的原因.
怎樣能夠找到那個密碼呢?.
那個DNA就可以配對一些藥物.
或者一些治療等等.

$^{121}$但還沒有完全的.
來做一個詳細的分析.
但是人類靈性的問題.
本來已經在神的啟示裡面.
或者不知道是不是他自己再參照舊約整體.
而得到的一個這樣的智慧.
他能夠說出.
阿當在神的面前.
他的失落.
和一個末後的阿當.
帶來的一個新的生命的開展.
這就是五章的十二到二十一節.
要給整個羅馬書.
一個很大的.
一個從上面神的角度去看.
人間的那種困難.
所以如果這一段經文.
在當時讀經的受眾.
羅馬的教會.
他們去看的時候.
他們明白了.
他們知道五章的一到十一節.
是帶來整個世界.
真正的一個超越.
因為他們正正就是處身在羅馬帝國的那個承平.
那個售賣的平安.
讓人感受到.
你在帝國的保護下.
你就會有平安.
你就會有小的戰爭.
你就可以去到不同的地方去做貿易.
不怕那些盜賊.
有各方面的保護.
但其實這個只不過是一種很人間的.
很稍遜即逝的.
也都不是真正幫到人的.
但唯有在耶穌裡面的這個福音.
才是保羅所告訴我們的.
真正的平安的福音.
所以如果我們能夠明白到.

$^{161}$這個這麼重要的中心思想.
對我們去理解福音.
就會有很大的感動和得著.
我跟大家再很快的去溫習上個星期.
因為上個星期都講得比較快.
大家可以去看回那個PowerPoint.
和看回這個經文.
我們收集的範圍.
三章二十一節就是.
神的義賜給每一個相信的人.
因為每個人都需要福音.
萬民都可以恩信稱義.
萬民都是在神的憤怒底下.
三章二十一節就是說.
萬民只需要相信耶穌.
就可以得著神的義.
這個是三章二十一節.
至四章二十五節.
包括連舊約的亞伯拉罕.
這麼有信心.
這麼有在神面前的地位.
他都是恩信稱義.
是創世記的十五章第五節.
那時候根本就沒有受割禮.
還沒有成為神的割禮的約的子民.
更加沒有律法.
但是亞伯拉罕已經得到神的稱義.
得到神的應許.
到了第五章.
就是稱義帶給我們進入一個新的.
與神的復和關係.
這個就真的很重要.
上個星期都提過了.
羅馬的帝國君王.
不斷去宣傳自己是神.
不斷叫人去膜拜羅馬的神子.
宣稱可以與神有關係.
但是唯有保羅所傳遞的耶穌.
是真的將我們帶進與神復和的關係裡面.
所以首先就是我們能夠與神和好.

$^{201}$一個稱義的基督徒.
第一件事就是立即感受到.
那個關係上面.
所以這個是Salam.
Salam就是關係上面的復和.
不是只是一份感覺上的平安.
是關係上面的一份和好.
是沒有隔膜.
在神的裡面.
接著不單只是一個.
我們好像在地上感受到的一份和好.
而是你能夠上上.
可以獲得進入.
Gain access.
是進見君王這樣的一個狀態.
所以上上要自己學習.
有時是脈搏.
坐巴士的時候.
坐車的時候自己脈搏.
就立即自己去到天上的致勝所.
去享受那份進入的恩典.
進入和站在那個恩典裡面的這種福氣.
上上要這樣去體驗這種平安.
不是只是我們一個人來去胡言亂語.
而是真的你感受到.
你好像進到天堂的一個進入.
這是第二個福氣.
其實都是一個套餐.
將這個描繪.
而且是歡歡喜喜地.
去誇耀你自己現在擁有的.
同神復和的這種歡樂.
這種心靈的暢快.
所以是歡歡喜喜地誇耀.
誇耀是真的Bose.
真的好像去自誇的那個誇耀.
所以原來誇耀神.
就不會帶來負面的.
我們如果自誇.
經常稱讚自己很厲害.

$^{241}$稱讚自己怎樣怎樣.
我們聽得多我們都覺得很討厭.
可能說的人自己不覺得.
但是原來我們聖經教我們不要自誇是對的.
但原來保羅在這裡.
真的用誇耀這個自誇這個字.
歡歡喜喜地誇耀神.
在我們生命裡面的工作.
帶給我們的福樂.
所以你可以想像得到.
一個基督徒的新的生命.
就是在神裡面.
神的聖靈的進駐.
在我們內心裡面.
帶給我們的平安.
帶給我們能夠.
夠膽.
這個是要學的.
夠膽上上進到神的面前.
你想想如果你是一個地上的.
當時羅馬帝國.
你怎麼可能隨時去見那個皇帝呢.
他如果不讓你來.
你都不能進去.
大家去過英國的那個白金漢宮.
有些御林軍.
我們遊客.
就是很遠在大閘望進去.
你進不去的.
除非他有特別的日子.
你可以近一點.
如果不是.
你想爬進去.
你就馬上被人抓了.
所以不是每個人都容易.
可以坦然的進入.
那個重要的位置.
但是現在.
神告訴我們.
透過保羅.

$^{281}$我們可以坦然的進入.
你想想如果保羅想起了.
致勝所.
基本上.
表面上是沒有設防的.
人人都可以衝進去的.
但是沒有猶太人敢衝進去致勝所.
因為裡面是神的居所.
連祭司.
大祭司.
每年只能夠進入一次.
那個致勝所.
在裡面其實黑漆漆的.
在那裡彈血七次就要走了.
還有一條繩子綁住那個祭司.
有些鈴鐺.
那個祭司一路進去的時候.
如果一路很久都沒有鈴鐺響.
證明他已經死了.
可能睡著了.
死了在裡面.
那些外面的利未人.
就要拉那條繩子.
將屍體拉出來.
據聞猶太人就是這樣的傳統.
不是隨時可以進入致勝所.
但是現在保羅說.
所有在竇恩裡面的基督徒.
所有.
無分貴族.
無分你是猶太人.
還是你是外邦人.
都可以隨時進入這個恩典.
是可以獲得進入.
就是說你有一個許可證.
大家去過機場.
如果你的證件是最高級別的.
任何地方你都可以去.
這樣都挺特別的.
因為前年我們去參觀過機場的禁區.

$^{321}$帶一些小朋友.
真的有個證你可以進去.
當然是有些職員帶著.
你都覺得很特別.
沒來過這些地方.
你可以想像得到.
我們基督徒很少去說這些.
好像表面上這麼神秘的東西.
去分享我們在神的恩典裡面.
在神的榮耀裡面的那種享受.
那種心靈的無限的平安.
這些就是我們需要去學.
這些是我們需要.
要很沉靜地.
在我們自己的個人的一個靈修.
安靜的時候.
就這樣來到去.
你不是拼命地這樣祈禱.
祈十分鐘都不停.
都未必需要這樣的.
而是你很安靜地閉上眼睛.
享受與神的獨處的這種進入.
所以弟兄姊妹.
這個不是個別的恩賜.
這是普遍的救恩帶來的福樂.
我們真的要這樣去經驗.
所以第三方面.
這個歡歡喜喜.
原來連帶下面.
為什麼能夠誇耀呢?.
原來在患難當中.
我們都可以歡歡喜喜地誇耀.
就是說.
救恩不只是輸打贏要的.
救恩帶給我們就是.
在任何的情況.
因為有神.
我們就已經有這個力量.
有這個本身的心滿意足.
所以我們應該練就一個.

$^{361}$常常有主同在的福樂平安.
不要一步十寒.
不要突然間很厲害.
厲害到一個地步.
好像靈恩一樣.
又不要突然間枯乾到.
好像你對信仰完全感覺乏味.
又不可以的.
我們要慢慢地進入那種常態.
所以在患難當中都有力量.
第五方面.
就是有充滿了神的愛.
我們上一次說過了.
第五章的第五節.
羅馬書的五章五節.
為什麼我們能夠在患難中不至於羞恥.
仍然有盼望呢?.
第五節就是告訴我們.
因為所賜給我們的聖靈.
將神的愛囂貫在我們的心裡.
聖靈就住在我們裡面.
上星期我們看過其他關於聖靈的工作.
第八章我們看過了.
聖靈就住在我們裡面.
真是神賜給我們每一個信的人.
聖靈就會將神的愛囂貫在我們的心裡.
在我們面對患難的時候.
在我們面對盼望.
感覺可能會失望的時候.
祂會給我們這一份愛.
神的愛.
充滿在我們的心.
使我們有力量地去盼望.
所以第四和第五又是一個相連的.
然後就解釋.
哪些人可以得救呢?.
其實就是那些真是無助的人.
神不會只是救那些公義有義的人.
有義的人就不需要救的了.
這個就真的涉及到上個星期我們講過.

$^{401}$就是羅馬人的那個回報系統.
就是說我是一個家主.
很有錢,很有社會地位的.
我養了一幫人.
我就會供應他們的必需品.
一些物資,金錢.
到我有需要的時候.
例如我競選議員.
例如我要做一些公職.
這些食客就要為我效力的了.
就有一種這樣的.
好像傳書的一種利益輸送的關係.
但是保羅說.
基督徒就不是這一種的.
是超越這一種的了.
是一種愛裡面的關係.
所以你看到主耶穌.
不是為這些人死.
不是為一些好的道德的人死.
基督是徹徹底底的.
不是為了回報.
為了一些價值的交換來去死.
是為無助的罪人來去死.
所以是和羅馬當時的那種.
Honor and shame(榮譽與憐憫).
那種這樣的私恩受恩的制度.
是一個徹底的完全是兩樣東西來的.
那你就可以想像得到.
我對你好不是希望有回報的.
但是羅馬當時就是這樣的回報制度了.
你就可以想像得到.
這個就是聖經保羅說的.
真正的為罪人無條件犧牲的愛.
就是耶穌十字架的愛.
所以對保羅來說.
他都覺得是一件很偉大的事來的.
因為保羅都沒有辦法想像.
神這樣為罪人而死.
他更深的體會.
他自己原來是一直都.

$^{441}$似乎自己很有聖經的知識.
舊日的聖經知識.
更加是教導律法的一個拉比.
原來他自己根本不認識.
神這一份無條件犧牲為罪人死的愛.
所以對保羅來說.
他自己都是一個很大的驚嘆.
因為他處身在那個時代.
都是一種利益輸送交換的制度底下的人.
所以他對亞彌西姆.
對這個肥利門.
他都跟肥利門說.
其實你是虧欠我.
叫肥利門不要對亞彌西姆這個奴隸.
有過分的要求.
在社會制度上是可以的.
奴隸偷了你的錢走了去.
你有權抓他回來.
你有權去刑罰他.
但是現在他成為了.
在監獄裡信了主的弟兄.
你就要用基督的愛的這個新的制度.
來對待這個弟兄.
不是說他的錯你不需要理會.
而是他仍然要向你道歉.
還有可能需要還債.
保羅說我替他還這個債.
但是你可以想像就是.
肥利門就會明白.
原來在基督的救恩裡面.
超越了地上這種回報.
你對我好我對你好的一個機制.
所以我們看第八節.
五章八節.
唯有基督在我們還作罪人的時候為我們死.
神的愛就在此向我們顯明.
真正的愛就是無條件替罪人死.
那個罪人根本自己都不知道自己的問題.
也都不覺得需要神的愛和救恩.
但是神已經這樣為他死了.

$^{481}$這個就是聖經說的.
無條件的一份禮物.
就不是地上交換的一個制度.
這個就是基督徒為什麼要愛人如己的原因.
我們有時候星期四上午.
都會有祈禱會.
然後才有團契.
有些家長都有分享.
就說他有些微信的朋友.
很特別就是.
為什麼你們基督徒會這麼關心我們.
他有些朋友都知道他患了重病.
都沒有去關心.
沒有去問候.
姐妹就很坦誠地說.
因為是基督的愛.
我們就願意打個電話.
問候一下.
甚至有些姐妹很有心陪他看醫生等等.
這些就是慢慢學習.
展現一份無條件的愛.
我相信任何一個人.
無論多麼的硬心.
無條件的愛.
去到他的身上.
他慢慢發現.
他真的會受感動.
所以保羅就是被這份無條件的愛感動.
所以他見到有些為二人死.
是少有的.
為人人死.
或者有感觸的.
這些就是當時的社會制度.
我對你好.
我期望你將來對我有回報.
這個是為人人死.
有道德的人.
或者是一些很有社會地位的人.
是好人來的.
你為他死.

$^{521}$你為他付出.
一定是會有這樣的事情.
這樣的動作.
但是第八節就真是福音的核心.
基督在我們仍然是罪人的時候.
在他仍然知道上十字架的終極的痛苦.
是極之不容易的.
但是基督也願意去享受這個福輩.
第九節.
現在我們既靠著他的血清而更要藉著他免去神的憤怒.
再一次說一章十八節.
我們未信主的人.
所有人都在神的震怒之下.
因為我們是規缺神的榮耀.
現在我們已經免去了神的憤怒.
我們已經透過主基督的血.
清而免去了憤怒.
這個不只是一了百了.
而是後續就是再一次保羅說的.
一個更終極的和好.
更終極的就是以神為樂.
這個將來的.
你可以說將來的盼望的終極.
就是去到神面前的快樂.
這告訴我們.
第十節.
因為我們作仇敵的時候.
且藉著神兒子的死.
得與神和好.
既已和好.
就更要因他的身得救了.
更加因為主耶穌現在在神面前.
因為他現在還是活的.
這樣來得救.
不是赦免了你.
你就離開.
不要再見到我.
就不是這樣.
神期望我們一生由信主.
不斷在他的裡面.

$^{561}$所以我們藉著耶穌基督得救.
第十一節就是最後這一節.
不單止是這樣.
我們藉著我主耶穌基督.
這一位活的主.
與神和好之後.
我們現在又再一次藉著耶穌基督.
有這份終極的和好.
和好就是以神為樂.
或者說等候神的榮耀.
我們不是規絕.
現在我們是等候神榮耀的降臨.
我們可以以神.
作為我們心靈最大的喜樂.
所以神的整個位置.
填滿我們的內心.
所以真正的心靈的空虛.
是讓神去填滿.
所以基督徒在這六方面的福樂裡面.
真的會塑造我們整全的人生.
不是說物質的軀殼.
肉身的軀殼.
你的近視不會一下子.
因為你信了主耶穌.
你會得到一個即刻改善.
神也不希望這樣.
你不會突然間變高了.
或者等等.
而是我們整個的心靈.
是在一份平安.
福樂.
誇耀裡面.
愛的力量裡面.
去令到我們整個的內心的價值觀.
不斷的更新.
裡面的人越來越有力量.
越來越像耶穌.
那種真理,仁義,聖潔等等.
所以一個基督徒.
應該展現這六方面的新的經歷.

$^{601}$是救恩.
輕易帶來的新的經歷.
就是新的生命.
也可能是以弗所書.
歌羅西書所描繪的新人.
所以這就是羅馬書的中心思想告訴我們.
原先是在神的憤怒.
咒詛.
虧缺神榮耀.
現在每一個信主進入救恩的人.
就可以享用這六方面的福樂.
帶來新的經歷.
接著保羅就說.
這次不是在說救約哪個人物.
是在說人類的始祖.
亞當.
所以又是很間接地告訴我們.
保羅對救約的理解.
保羅對聖經的權威的深信不疑.
他不會去問一些我們認為不需要問的問題.
例如亞當之前是什麼人.
原始人又是怎樣.
等等.
恐龍又是怎樣.
保羅也不會再理解這些問題.
也不會再探討那些問題.
會不會是外星人.
保羅不會去理解這些問題.
他很清楚聖經的權威告訴我們.
神做的第一個人在地球上.
就是亞當.
在亞當的肋骨做的就是夏娃.
神可以做很多很多的人.
但祂只做了亞當和夏娃.
就可以繁衍後代.
就是我們今天透過創世記所知道的人類的歷史.
那個起源.
所以大概是五千至六千年.
就是人類的歷史.
所以帶我們去看這個開始的時候.

$^{641}$有一兩點相當重要.
告訴我們為什麼神的憤怒要淋到地上的眾生.
保羅在第十二至二十一節.
告訴我們一個原因.
我們要了解人類共有的DNA.
原來罪是透過一個人的過犯.
就是亞當吃了分別善惡的薯果.
違反了神對他的命令.
他就犯罪了.
因為是人類的始祖.
罪就從他進入到世界後裔.
我們要先讀讀這段經文.
其實如果我們不是逐節去看.
整段就很簡單.
其中比較難解的可能就是第十五節.
或者是第十二節這一段.
十五節就最難解.
但是整段十二到二十一節.
其實是很清楚很簡單.
將亞當帶來人類的定罪.
那個悲慘的結局.
和耶穌他在一個信服裡面.
信服神的旨意.
成就的救恩帶給新的人類.
就是在救恩裡面的新人類.
就可以得到各種的釋放和平安.
就是一個這樣的對比.
一個這樣的類比.
所以在這一段經文.
耶穌又稱為末後的亞當.
我們先讀讀這段經文.
第十二節.
這就如罪是從一人入了世界.
這個一人雖然我們這時候還不知道是誰.
肯定下文我們就清楚.
因為整段經文都是說亞當和耶穌.
就是亞當帶進世界的罪.
死又是從罪來的.
雖然我們今天見慣不怪.
見到慣了.

$^{681}$但原來保羅將我們人類的命運.
去做一個很清楚的神學詮釋.
讓我們知道原來為什麼會死.
自然啦 新老病死.
原來不是的.
死是從罪來的.
死是臨到眾人.
而且下面繼續說.
就算我們不是和亞當犯同一樣的罪.
我們都是在死亡的權勢下.
要面對死亡.
所以你可以想像到.
亞當一次違反神的命令.
罪就因為這樣.
這個背影就進入世界.
死就因為罪.
而令到我們人要面對.
活在一個咒詛底下.
原來我們是有壽命.
我們是有一個日子.
剛剛收到何柱國先生.
75歲因為肺癌離世.
算是香港近十幾二十年.
一個比較出名的人.
每個人都有他的壽命.
這裡告訴我們.
保羅不是解釋什麼原因的病.
而是說人類原本.
如果亞當不犯罪.
罪就不會進入世界.
進入人類.
死亡也不會借罪去刑罰地上的人.
就是說原先的設計.
人是不會死的.
是直接可以享受在地上的永生.
但是這裡就已經很清楚.
是回不了頭的.
我們在歷史之後的人.
我們只能夠看到亞當.
真的這樣違反神的命令.

$^{721}$來面對審判.
就是那個咒詛.
死就臨到眾人.
因為眾人都犯了罪.
固然是每個人有罪性.
是會犯罪.
到底是不是犯和亞當一樣的罪呢?.
保羅就很清楚說不是.
就算不是這樣.
就更加證明.
死已經是掌控了人類.
就算沒有犯同一樣罪的人.
都同樣在權勢底下.
要被命運就是.
被死亡,被罪惡.
牢牢的監控著.
特別強調.
第十三節.
沒有律法之先.
罪已經在世上.
但沒有律法.
罪也不算罪.
大概的意思就是說.
因為下面第十四節就說摩西.
摩西是神叫他頒佈十誡.
和十誡所衍生的不同的法例.
所以這裡說的意思大概就是說.
當摩西的時候.
還沒有正式的法律.
但是已經有不同的人.
在地上犯不同種類的罪.
但沒有明文的法例.
那些律法.
就不可以點明你是什麼罪.
你是犯了哪一條罪.
這樣的一個意思.
但不等於人沒有犯罪.
但是重點反而是想說.
就是下面第十四節那裡.
然而從當道摩西.

$^{761}$雖然從當道摩西.
都沒有明文的神頒佈的那些法規.
那些律法.
那些條例.
但是明顯就是.
一個一個都會死.
從當道摩西.
有哪個人不用死呢.
只有一個.
就是那個爾諾.
然後就沒有了.
包括阿當都死了.
930歲.
摩西又死了.
120歲.
由阿當到摩西.
每個人都死了.
那就是說那個邏輯就是.
每個人都是犯罪而死.
沒有人不是犯罪而他不死.
但是那時候還沒有律法.
那你就可以想像得到.
原來罪的權勢已經藉著阿當進入到人類裡面.
就是在那個咒詛下.
就是說.
正如保羅自己一章十八節所說的.
神的憤怒從天上顯明.
在那些不乾不易的人身上.
所以現在保羅回頭說人類的開始的時候.
告訴我們.
為什麼神要追討地上他所做的人的罪的審判呢?.
原因就是因為罪已經透過阿當進入到人類裡面.
這個就是一個已經注定了不能改變的結局.
如果透過只是人去改變.
是沒有辦法改變的.
但是神安排另外一個人來.
就是耶穌.
來將我們拯救.
那就是一個新的開始.
保羅就是很想說.

$^{801}$救人救的人類的命運結局就是死.
新的人類的命運結局就是永生.
你記得第21節就是這樣.
21節就說.
就如罪作王叫人死.
罪藉著阿當作王叫人死.
照樣恩典藉著耶穌的義作王.
叫人因我們的主耶穌基督得永生.
舊有的體系.
在阿當裡面.
每個人都要面對就在死亡.
但是現在在耶穌這個.
清二的救恩的寶血裡面.
我們就可以得到永生.
這個對比.
一個這樣來襯托.
就是說現在已經超越了摩西的律法.
已經超越了猶太人.
他們自始是亞伯拉罕的子孫.
現在是說人類.
你可以想像保羅的南逃.
那個普世人類的救恩.
已經不再是小小的猶太民族而熄滅了.
現在是說整個人類.
所有的人類.
他不知道原來這個世界有這麼大的.
他們那時候沒有飛機.
更加不會有互聯網.
更加不會有一些很快速的大洋的洋船.
就是說他如果要走去中東.
也要走很久.
不要說要去到我們中國.
根本就不知道.
在羅馬帝國的時候.
是我們很繁盛的.
應該是漢朝那些.
你可以想像得到.
在中國有幾億人.
都沒有聽聞有神.
但是保羅已經告訴我們.

$^{841}$人類一開始就已經失敗了.
一開始就已經要註定審判和滅亡了.
但是神就是要為人類預備一個救恩.
所以保羅這一段經文是很重要的.
將我們帶回去.
為什麼神要憤怒地上的眾生呢?.
為什麼萬民必須要恩信稱義呢?.
為什麼唯有在救恩裡面才能夠成為新的人類呢?.
所以上文的六個福樂.
如果你不是基督徒.
你沒有聖靈.
你是享用不到的.
你可以分析基督教.
但是你是不會享用到神進入心靈的那份平安.
那份膽量.
那份誇耀.
因為必須要有神.
才能夠讓我們有這份滿足.
所以第十四節.
正正告訴我們.
律法雖然還沒有頒布.
因為要到摩西的時候才頒布.
那你可以看到救恩的歷史的一路進程.
沒有律法.
但是人類都一樣在死亡的籠罩操控下.
面對各種不同的死亡.
那怎麼去解決呢?.
所以你會看到.
哪些不與阿當犯一樣罪過的.
也在他的權下.
就是罪的權下.
不過保羅告訴我們.
這節聖經就很特別了.
那乃是那以後要來之人.
就是耶穌.
因為整個段落都是將兩個人討論.
那以後要來之人的遇象.
遇象就是這個Tablet.
我給大家看看希臘文.
這個就是上一次我們說的.

$^{881}$在羅馬時代就是這個私恩制度.
上個星期本來我跟著回去就深入了.
那個私恩制度.
整個深入的很詳細.
我又忘記給大家看.
今晚我再給大家看.
什麼叫私恩制度呢?.
你稍微用深入打一打.
就已經整個大段.
幾頁出來.
大概就是你給我恩惠.
我是一個有錢人.
我給你恩惠.
將來你就要回報我.
這個就是羅馬的私恩制度.
其實香港一樣.
有時候你去茶樓找位.
大時大節.
就叫執務司幫我買張卷.
這就是私恩制度.
大概就是這樣.
一個彼此一個大家好像有個關係的.
但你留意到保羅正正在說的就是.
基督是無條件的.
為沒有任何回報的人去死的.
沒有東西可以回報給基督的.
他都願意為這些人死.
我們繼續下去下面.
這些我們講過.
我們不詳細講了.
這裡也提到.
就是希臘門這個字.
阿當.
我打錯了字.
阿當就是人類的代表.
神是這樣理解.
他是第一個人.
本來是不會死的.
他是可以有永生的.
但是當他違反神的禁令.

$^{921}$不是那個果是毒果.
是那個背役.
令到他因為一次的過犯.
這裡是一個人的過犯.
眾人都死了.
所以因為他的過犯就帶來罪和死亡.
第十四節尾這裡.
他說.
以後要來之人的遇象.
遇象這個字.
希臘文本身就是這個字.
不過用了英文的串法.
就是一個街模.
好像一個帽子.
所以阿當和耶穌有相類.
就是他是人類的舊有制度的代表.
因為他第一個人.
本來在神的構造裡面.
是不會死的.
其實他可以吃生命樹的果子.
但是他沒有吃.
他吃的是蛇.
引誘他吃分別善惡樹的果子.
因為他以為.
吃了就像神一樣的智慧.
不過他有那個智慧是沒有用的.
因為他有罪.
他可以知道什麼是善惡.
但是他沒有力量.
不能夠去不行惡.
這個就是.
保羅說的這個遇象.
很明顯.
在這段經文.
是很想將阿當所做的.
和基督所做的.
帶來不同的後果.
那個結論.
那個的.
複落.

$^{961}$我們就可以.
從第十六節開始去看.
這一段的比較.
所以.
保羅就很清楚說.
如果是這樣的時候.
一個人的過犯.
就及不上.
神的恩賜.
為什麼這樣說呢.
就是說阿當一次的.
違反神的命令.
就是他個人的過犯.
基督一次的.
順服神的旨意.
釘十字架.
的恩賜.
然後告訴我們.
如果一個人的過犯.
眾人都死了.
何況是神的恩典呢.
與乃因耶穌基督.
一個人恩典中的賞賜.
豈不更加倍的淋到眾人.
一個人他做了這件事.
都尚且令到全人類.
他的後裔.
都要夢想這個.
被罪惡.
死亡.
操縱.
無奈悲慘的結局.
何況是神.
主動的.
他的計劃裡面.
要使到.
在他的.
兒子的救恩裡面.
神主動要做的.
這件事.

$^{1001}$豈不是更加倍的淋到.
每一個信主的人呢.
阿當是.
好像在他.
表面上是他一個很強的意識.
去做了這件事.
但無奈地他這個行動.
就帶來.
整個人類就面對死亡.
但現在神有意識地.
有計劃地.
要將他的恩賜.
賜給信耶穌的人.
那就更加是神.
加倍的將.
這份恩典去賜給.
所有信的人.
那你又可以想像得到.
保羅看人的.
過犯面對神的審判.
如果人.
主動投向耶穌.
神賜給.
這樣的人的福樂.
那個美好.
是無限.
保羅很想帶出.
那個很懸殊的對比.
所以,頂子們我們要.
感受一下救恩帶給.
我們的榮耀.
救恩帶給我們每一個基督徒.
其實是共有的.
是一個集體的.
所以基督教是應該.
不是夕陽工業.
基督教不應該讓人感覺到.
是沒有生命力.
好像支離破碎.
不應該是這樣的.

$^{1041}$基督教是展現一份.
永恆的平安.
跟著告訴我們.
因為一個人的犯罪.
不斷說一個人的犯罪.
一個人的過犯.
一個人的背役.
一個人的恩典.
一個人的信服的.
擺上等等.
好像做一個很強烈的對比.
阿當所做的.
和基督所做的做一個對比.
不斷用不同的.
角度,不同的詞彙.
去修飾.
阿當多麼背役.
主耶穌多麼信服.
帶來的是南轅北轍.
那個結果.
那個結局.
十六字我們再看.
一個人的犯罪.
就被定罪了.
因為神說了是不能收回的.
阿當.
方便大家.
我們回到第三章的措施記.
稍微看一看那段經文.
相信.
保羅都是.
欣著這段經文.
他很清楚神的啟示.
摩西寫這段經文的時候.
都很清楚.
三章的十四節.
創世記第三章.
耶和華神對蛇說.
你既作了這事.
就必受咒助.

$^{1081}$比一切的生畜.
野獸更甚.
你必用土耳行走.
終生吃土.
我又要叫你和女人.
彼此為仇.
你的後裔和女人的後裔.
也彼此為仇.
女人的後裔要傷你的頭.
你要傷她的腳筋.
又對女人說.
我必多多加增你懷胎的苦楚.
你生產兒女.
必多受苦楚.
你必戀慕你丈夫.
你丈夫必管轄你.
十七節.
又對阿當說.
你既聽從妻子的話.
吃了我所吩咐你.
不可吃的那樹上的果子.
比較詳細去描述阿當.
所以.
雖然是.
阿夏娃.
叫他吃.
但始終主動都是阿當自己吃.
你吃了我吩咐你.
不可吃的.
就是這個遺逆.
地必為你的緣故受咒助.
即是說以往是沒有受咒助的.
一個.
很容易耕種.
很容易出產.
你必終生勞苦.
才能從地裡得吃的.
地必給你長出荊棘和窒泥來.
你也要吃田間的菜蔬.
你必汗流滿面.

$^{1121}$才得糊口.
直到你歸了土.
因為你是從土而出的.
你本是塵土.
仍要歸於塵土.
這是一個很無奈.
很悲哀的結局.
本來是塵土.
但神的生氣.
吹入阿當裡面.
他是有靈的活人.
這簡直是神跡.
一堆泥怎會變成一個人呢.
好像那些.
那些.
那些療齋.
神跡那樣.
傳統故事.
但不是的 這是聖經的啟示.
阿當是由泥土變成人的.
那麼.
但現在神就宣判.
那個結局那個咒助.
就是他要歸回塵土.
所以真的會.
人就會土葬化骨.
回到大自然.
這個就是結局.
就是塵土.
原來不需要是塵土.
本來是永恆的人的永生.
現在就沒有了.
所以要在基督裡面.
再有一個新的創造.
我們才可以有.
這個永生的福樂.
這個就是那個咒助.
這個就是保羅.
他就著這個神的很清楚的啟示.
來講述.

$^{1161}$在阿當這一次的過犯裡面.
死就作了王.
這個就是第十六節.
原來審判是由一個人而定罪.
就是剛才我們所讀的經文.
三章十七節.
因此乃是由許多過犯而稱義.
因此就是神藉著耶穌基督的一個救恩.
給我們這些犯了很多罪的人.
都可以白白的去死.
所以我們就要.
在這個過犯裡面.
我們要有一個新的啟示.
讓這些犯罪的人.
都可以白白的享受神的恩賜.
就是永生.
所以.
你可以說很不值.
阿當只是吃了一次禁果.
違反神的命令.
但是.
就已經立即.
宣判他的那個.
咒助結局.
由幾多歲開始宣判呢.
聖經沒有詳細講.
但是他也有.
九百三十歲.
在這個地上.
最終他也要歸回塵土.
但是.
我們雖然有很多的過犯.
很多未信之前的背影.
但是神就給我們.
這個恩賜.
就是眾多的罪.
在救恩 在耶穌裡面.
就立即得到那個拯救.
稱義.
還有去有這份的永生.

$^{1201}$是說真的.
一個地上的結局.
和永恆的生命的對比.
所以真的有永生.
有天堂.
是真的一個我們將來.
可以進入的地方.
到底那個靈魂.
是要多少秒才能去到天堂呢.
或者靈界根本就沒有距離.
不應該用我們物質的頭腦.
去想那個距離.
根本就可能已經是.
都不是想.
可能就已經到了.
所以我們已經和靈界.
相差太遠.
我們的認知.
所以我們不知道的.
但是無論如何.
這裡是說地上的死亡就坐.
和耶穌帶給我們的永生.
這個對比.
再繼續.
十七節.
如果一個人的過犯.
死作亡.
何況受洪恩.
又蒙所恥之義.
豈不更要因耶穌基督.
一個人.
我們可以在生命裡面作亡嗎.
意思就是.
一個人一次的過犯.
死就籠罩整個人類.
操控著整個人類.
是人都怕死.
懼怕.
和要面對死亡.
沒有人可以倖免.

$^{1241}$是很絕望.
其實.
稍微訪問一下都知道.
沒有人說.
不怕死亡.
但是你可以想像到.
保羅這裡的描述.
用這麼強烈的形容詞.
修昔.
生命作亡.
基督徒.
生命作亡.
所以.
原因是因為我們領受了.
洪厚的恩典.
洪恩.
是.
abundance.
是很巨大的恩典.
和蒙了所賜的.
清義.
以至保羅說.
豈不是更加可以因耶穌.
我們可以在.
生命裡面作亡嗎.
不再受死亡的威脅.
所以.
不應該再懼怕死亡.
無懼死亡的威脅.
因為連死亡.
我們已經看穿了.
他不過如此.
令我們.
懼怕不知去哪裡.
現在我們知道自己去哪裡.
基督徒就應該.
當然痛是有的.
有知覺.
沒理由我們不痛.
那些是一個過程.

$^{1281}$但是我們知道的是結局.
所以.
我們就應該今天.
不斷要在我們所領受的.
洪恩和賜給我們義的.
裡面.
我們要有這樣的行為.
就是我們要在.
生命裡面自己做亡.
當然是和神.
和主耶穌去做亡.
就是說我們.
不再需要.
好像我們舊有的人.
的那個個性.
那樣.
有些人說.
我是這樣的.
我未信之前都是這樣的.
我沒有人說.
這樣不對.
你表現你的真性情.
但是在基督裡面.
這段經文就正正告訴我們.
可以在.
不斷領受的.
洪厚的恩典裡面.
我們就可以由.
嘗試.
和真的慢慢.
在生命裡面去做亡.
就是說我們不再要服從.
我們舊有的脾氣.
舊有的.
那些價值觀.
我們可以有新的價值觀.
我們可以愛人如己.
我們可以來學習.
無私的對人的寬恕.
接納.

$^{1321}$就是在生命裡面做亡.
這就是真的要學.
這些真是很高級的基督徒的生命.
才會這樣想.
如果不是你做初一我做十五.
信不信耶穌都是這樣.
不信耶穌的人都會這樣.
有時不信耶穌的人反而不會這樣.
變了我們就不是在生命裡面做亡.
這裡說一個人過犯.
阿當的過犯.
死就藉著這樣去籠罩.
操控著人類.
現在耶穌.
帶來的菁兒.
給我們洪厚的恩典.
不斷對我們的.
接納對我們的寬恕.
對我們的愛.
使我們整個人的生命.
有一種很大的感動.
帶來的就是.
我們可以在生命裡面.
自己作亡.
不再被.
我們所謂的價值觀.
操控我們.
這就是第十七節的意思.
所以今晚我回去再.
因為我今天開了好幾個會.
不夠時間再做.
再仔細一點的PowerPoint.
我回去再列舉給大家.
就是下一堂的PowerPoint.
這裡不斷做一個.
這樣的描述.
這樣的對比.
舊的一次就永不超生.
一次的過犯.
死就藉著.

$^{1361}$阿當這個過犯.
操控人類.
但現在藉著基督的一次的馴服.
他的這樣的.
無條件的.
對我們的菁兒.
帶給我們.
可以在生命中.
我們自己作亡.
十八節.
如此說來.
因一次的過犯.
眾人都被定罪.
照樣對比.
一次的義行.
一次過犯一次義行.
好清楚了.
不斷在說舊的代表.
阿當.
新的代表.
這個人是背役.
那個人就是信徒.
這個就是耶穌.
他對神的一個義行.
眾人怎樣?.
就被稱義得生命.
不斷強調生命.
那個Life.
說的是那個活力.
基督徒的.
新生命的活力.
是有活力的基督徒.
今天回來開會.
屯元天.
很多的牧者.
因為我們兩個月見一次面.
一起分享.
祈禱.
好幾個牧者.
因為上一次去不到.

$^{1401}$因為星期四我就有福音團契.
他說見到曾牧師.
為甚麼你越來越年青.
越來越光彩.
你很忙的.
你紅恩現在又做大學.
我還不知道.
類似這樣.
希望我們能夠讓人感受到.
保羅不斷說.
生命作王.
這裡被稱義得生命.
不再被定罪.
所以我們在.
耶穌的義行裡.
所以我們不斷要想想.
主耶穌他無條件的犧牲.
帶給我們的.
我們的那個稱義.
第十九節.
好像怕我們不夠描述得透徹.
他不斷用不同的.
角度去描繪.
阿當所做的.
耶穌所做的.
他這個不是偶然的過犯.
是他內心全心背逆神.
這樣五逆.
蛇那個引誘.
其實有多致命呢.
只是說一句.
吃下那些.
分別善惡的水果.
你就可以像神一樣有智慧.
這樣去勸他.
不是恐嚇他.
不是那條蛇頭去咬他一口.
是引誘他.
但是保羅說.
這個原來是背逆.

$^{1441}$因一個人的背逆.
眾人成為罪人.
不斷描繪不同角度.
又是照樣.
一個人的順從.
成為義.
不只是這樣宣判為義.
而是我們本身就得到神的義.
我們在義的地位裡面.
所以保羅看這個義的地位是很高的.
相信他在寫這段經文的時候.
他在想的是全人類.
不只是猶太人.
他在想的.
不是一個儀式的義.
是在基督裡面.
圓滿救恩的義.
可以站在神面前的義.
所以這個是很重要.
所以不斷在說背逆和順從.
大家可以試一下做一個對比.
做一個table.
第20節.
律法本是外天.
這個就是回應.
沒有律法的時候眾人都死了.
因為死了就立不到眾人.
就算不是犯同一樣的罪.
就算沒有法律.
條文不能夠點出.
什麼罪是什麼罪.
但是人的死就證明.
罪已經在人的裡面.
現在這裡告訴我們.
律法本是外天.
摩西頒佈律法.
叫過犯去顯多.
只是罪在那裡顯多.
恩典就更顯多.
就如罪作王叫人死.

$^{1481}$再次回應第12節.
照樣.
恩典藉著義作王.
叫人.
因我們的主耶穌基督.
得永生.
不斷想告訴我們.
在耶穌的救恩裡.
帶給我們的.
就是這個永生.
帶給我們的就是.
我們不再需要在死亡.
在罪.
那個權勢底下.
這裡第6章.
到第8章.
就詳細講.
這個.
新的那個律.
在我們裡面.
勝過罪.
和死亡的那個律.
這個就是第6章就要.
緊接的探討.
叫我們可以.
勝過罪惡.
因為現在我們仍然.
是處在舊人的.
舊有的體系.
舊有的生命.
舊人的操控.
和我們現在得著的.
新的生命同樣在我們.
這個人的內心.
所以必須是要再進一步.
不只是在那個.
神學上.
理論上.
我們是進入.
那個新的生命.

$^{1521}$除非我們這個人死了.
一個新的當然沒有問題.
但現在同樣這個信主之前的人.
和信主之後那個人.
在那個心靈裡面.
就真是有一個新舊的.
交戰.
亦都是第6章第7章要探討.
的那個課題.
怎樣保羅你可以.
給一些具體的實例.
一些重要的.
神學我們能夠.
可以知道.
我們.
可以在義裡面.
去作王.
我們怎樣能夠藉著恩典.
藉著義去作王.
怎樣去.
剋聖我們自己.
舊有的生命.
這裡就是保羅點出了.
那個重點之後.
先做了一個對比.
如果你不在耶穌裡面.
你只能夠在阿當裡面.
因為你出世就是阿當的後裔.
你就是在那個就坐.
和.
那個死亡的.
操控底下.
只能夠等待你的人生.
生老病死.
但在基督裡面.
不同了沒錯.
生老病死這個定律.
仍然未解除.
因為神的就在大地.
但是神給我們一個新的.

$^{1561}$律就是.
聖靈住在我們裡面.
和基督的義行已經將我們.
帶進一個.
可以作王.
來去管束我們自己的生命.
那個新的人類.
所以要怎樣能夠具體實踐.
第六章就交代了.
第六章就告訴我們.
原來我們可以.
在基督裡面.
不再犯罪.
這就是第六章要交代的問題.
第五章這一段.
是告訴我們人類.
那個原先的.
創造.
但是一個過犯一個背役.
一個.
犯罪.
就帶來全人類都在.
那個就坐底下.
但現在基督帶給我們.
新的生命.
新的活力.
這就是第五章.
所以這段經文是很重要的.
一個解釋.
一章十八到三章二十節.
那個罪惡.
就坐.
就是那個原委.
所以第五章的一到十一節.
就是一個新人類.
基督徒所有基督徒.
都應該要盡快.
享有的福樂.
就是與神和好.
帶來的那個生命的.

$^{1601}$那個福樂.
大家有沒有什麼問題想問呢.
今天可能我講快了一點.
都.
還有一些時間.
第六章我還沒有預備.
那個power point.
如果大家.
很想來到去再.
看得深一點.
這裡的對比其實就.
很呃要.
但很有透視.
就是將人類.
那個的.
結局.
講了出來.
就是說共有的DNA.
無論你是什麼民族.
都在阿當的.
那個死亡.
就坐底下.
但只需要在基督的救恩.
裡面的人.
就立即成為新的人類.
可以在生命.
當中來到去作王.
所謂的作王.
就是可以超越.
自己的舊有價值觀.
在神的國度裡面.
有份有用.
這樣的一個.
新人類.
我們是很需要透過這一段.
基督和阿當的.
對比.
來到去有這樣的.
對生命的一種透視.
不要小看.

$^{1641}$這麼簡短的經文.
沒有的人就是沒有的.
他不斷去搜尋人類的.
意義他找不到的.
如果他不回去聖經這一段經文.
是找不到人類的意義的.
找到的可能就是很多的鬥爭.
很多的無奈.
很多的所謂的命運.
但是這段告訴我們.
我們可以超越了.
所謂的命運.
我們在基督裡面就可以立即.
有這個永恆生命.
和地上的平安.
和在生命裡面作王.
我漏了.
我問大家的問題.
新約還有另外一段經文.
講阿當的.
我就用另外一個角度去講.
就是哥林多前書的.
十五章.
趁還有幾個字的時間.
和大家聊聊這一段哥林多前書.
保羅.
對阿當的理解.
就是這兩個大段落.
當然在保羅的年間.
也有一些次經.
是講阿當的.
但是那些通常都是一些錯誤的說法.
哥林多前書.
十五章.
明顯是講.
關於復活的事.
詳細我們就不看那麼多.
由第35節.
看起.
這裡就特別提到.

$^{1681}$這裡就明顯說.
說耶穌.
就是末後的阿當.
又是將.
首先和末後作一個對比.
35節.
當然上文就是在講.
要爭論的就是.
基督徒會復活.
我們只要是很真心.
相信耶穌基督.
是降生.
為我們釘十字架.
和他死而復活.
這個就是福音.
我們如果信靠這個福音.
我們就一定會有將來的復活.
這個是保羅告訴我們.
當時的人.
最怕以為沒有.
所以保羅就很斬釘截鐵.
只要你信這個福音.
以這個福音來.
繼續來信靠.
就是一次得救.
一直相信.
就必然的得救.
一直相信這個信.
因信而至於信.
不斷的憑信心.
去信靠神.
就一定可以得到永恆的復活.
35節.
然後他就用另一個角度.
去講關於復活的身體.
那個形體.
這裡又有一些.
給我們看到將來的復活.
就是這樣.
第35節.

$^{1721}$或者有人問.
假設有些人會這樣想.
他就去提出.
一個這樣的討論.
帶出他的觀點.
有人問.
死人怎樣復活.
帶著什麼身體來呢.
無知的人.
你所種的.
若不死就不能生.
並且你所種的.
不是那將來的形體.
不過是子粒.
即如麥子.
或是別樣的菊.
但神隨自己的意思.
給他一個形體.
並叫各等子粒.
各有自己的形體.
凡肉體各有不同.
各有不同的形體.
即是我們站立一個人的身體.
一個模樣.
瘦又是一樣.
又是一個不同的模型.
鳥又是一樣.
魚又是一樣.
有天上的形體.
也有地上的形體.
但天上形體的榮光是一樣.
地上形體的榮光又是一樣.
日有日的榮光.
月有月的榮光.
日有日的榮光.
月有月的榮光.
星有星的榮光.
這星和那星的榮光.
也有分別.
他一直在說.

$^{1761}$一些我們能夠理解.
明白的一些形態.
地上神確從其類.
做了人.
做了地上的生畜.
生畜可以分野獸.
可以分家禽.
可以分鳥.
可以在天空上飛.
在海上游.
不同的形態.
當然還有樹木.
他這樣來形容.
然後他又形容.
天上有不同的.
形體的可能是天使.
因為他們都是讀舊約.
會見過天使.
包括保羅自己也見過耶穌.
有不同階層的天使像.
有加伯列.
有米加勒.
有其他的天使.
告訴我們.
他們都有不同的榮光.
有不同的.
甚至大自然.
日月星辰.
都有不同的榮光.
這樣來描述.
每樣東西都是神賦予他的存在.
讓我們能夠了解.
接著第42節.
有些是我們暫時未能夠了解.
因為我們根本沒有見過.
死人復活.
但是保羅.
他們這班的使徒.
肯定知道.
耶穌的榮耀的復活身體.

$^{1801}$因為已經向不同的人.
顯現過.
第42節.
死人復活都是這樣的.
所中的是必扭壞的.
復活的是不扭壞的.
所中的是羞辱的.
復活的是榮耀的.
所中的是軟弱的.
復活的是強壯的.
所中的是血氣的身體.
復活的是靈性的身體.
如果有血氣的身體.
我們這個肉身的人.
神賦予我們存在在這個地球.
降生在這個地上.
就是一個形體.
規定了我們的眼耳口鼻.
規定了我們的膚色.
我們的種族.
這是神去做的.
沒有選擇的.
所以我們要去.
做的了.
沒有一個人知道自己.
投胎去了哪裡.
本來你是香港.
你去了非洲.
沒有這些東西的.
那些科幻片就會有.
這個就是神安排了我們.
我們發現.
這個原來是一個血氣的身體.
就是說真的會生病.
真的會在咒詛下.
我們會有絕望感.
我們會有開心不開心.
等等.
但更重要的是我們會過去.
因為是扭壞的.

$^{1841}$但現在保羅很想告訴我們.
將來.
必然會有一個復活之後.
靈性的身體.
如果有血氣身體.
必定就會有靈性的身體.
經上也是.
這樣記著說.
引述創世記的第一章.
首先的人.
當成了有靈的活人.
創世記.
對不起這是第二章才對.
第二章第七節.
然後末後的阿當.
就是耶穌成了叫人.
活的靈.
但屬靈的不在先.
屬血氣的在先.
以後才有屬靈.
就是先有人類的.
始早他的墮落.
他的犯罪.
人就在這個血氣底下.
會生老病死.
屬靈的來了.
就是耶穌.
第一節他們所說的.
有死人會復活.
所以現在不是屬靈的先有.
後來才是屬血氣.
是剛好相反.
屬血氣的先有.
我們體會過地上的人.
原來有時是這麼虛空.
有時是這麼無奈.
有時是這麼悲哀.
但是現在.
有了屬靈的身體.
所以我們應該期盼.

$^{1881}$永恆的降臨.
應該期盼.
我們離開.
血肉之軀之後的永恆.
第四十六節.
屬血氣的在先.
以後才有屬靈.
頭一個人就是阿當.
出於地屬土.
第二個人就是耶穌.
出於天.
由天上來這個世界.
那屬土的怎樣.
所以凡屬土的也怎樣.
不是金木水火土.
屬於泥土.
即是說是灰塵.
是Nothing.
是無有的阿當.
神給他一個這麼燦爛.
這麼美麗的軀殼.
但是他不珍惜.
他背逆神.
就受咒作.
所以我們有這個血氣的身體.
就承接了.
阿當的咒作.
和這個血氣的身體.
但是.
保羅告訴我們.
第二個人我們今天相信的耶穌.
是屬天的.
永恆的主降世成為人.
成為那個救贖主.
成為那個死囚.
為我們.
帶來新的生命.
新的永恆的身體.
所以我們信主.
暫時地上的身體.

$^{1921}$還沒有變化.
而是離開世界.
就會得到屬靈的身體.
所以屬土的怎樣.
記住不是我們那個.
金木水火土.
那個屬於什麼.
屬於什麼年.
牛啊.
虎啊.
還是土啊.
屬於什麼生肖.
是說.
只是地上.
塵土.
我們真的能夠.
越信主內.
越謙卑.
不反對我們.
要有各樣的知識.
但一切的知識背後.
我們都看到一切都是恩典.
使我們能夠得到.
這一切的生命.
所以知識.
和我們謙卑是沒有抵觸.
應該越不斷有不同的知識.
但一切都要服在謙卑.
在神面前.
底下.
記住你和我不過是塵土.
真是塵土.
所以.
我們不斷帶著.
這種無奈感去服事.
就很有建設性.
因為.
做多少都是神的恩典和榮耀.
不能誇口.
我們是屬土.

$^{1961}$所以屬天的怎樣呢.
這就是我們大膽的盼望.
既是屬地的無奈.
謙卑.
同時間.
要擁有這種屬天的怎樣.
我們凡是屬天的.
就會怎樣.
如果耶穌現在在天堂.
我們還怕什麼.
我們去不到天堂呢.
我們要表達出.
我們是要活在這個世界.
我們是要生存.
但是我們不是屬於這個世界的人.
我們是一個.
短暫的過客旅客.
但是要用心用力去活在這個世界.
不要成為一個逃避主義.
什麼都隨隨便便.
是不應該的.
但是記住.
凡屬天的也就這樣.
是因為基督.
已經將這個永恆的救恩.
給了我們.
所以我們將來一定.
和耶穌一起.
永恆裡面得到這個福樂.
所以我們既有.
屬土的形狀.
即是地上的軀殼.
所以是將來.
必定會有屬天的形狀.
所以就是地上.
首先的阿當.
和末後的阿當.
耶穌的對比.
阿當承受了.
軟弱和無奈的軀殼.

$^{2001}$因為他的背影.
神對他的刑罰.
造就了人類.
面對很多.
因為死亡的威脅.
帶來的很多悲歡離合.
但是在耶穌裡面.
不是說我們要逃避.
或者我們不能夠.
抗拒這些東西.
而是就算我們同樣面對這些困境.
我們有這個永恆的盼望.
和這個屬天的福樂.
就是耶穌賜給我們.
所以第五十節就是這個.
弟兄們.
我告訴你們.
血肉之體不能夠承受神的國.
必扭壞的身體.
都不能夠承受.
我如今將一件.
奧秘的事告訴你們.
我們不是都要睡覺.
乃是都要改變.
就是主在內.
就在一煞時.
斬眼之間.
浩同末次吹響的時候.
因為浩同要響.
死人就要復活成為不死.
因為浩同要響.
死人就要復活成為不死.
因為浩同要響.
死人就要復活成為不死.
死人就要復活成為不扭壞.
我們都要改變.
這個不扭壞的.
總要變成不扭壞的.
必死的這個軀殼.
總要變成不死的.

$^{2041}$必定扭壞的.
既然間變成不扭壞.
這個必死的既變成不死.
重複這樣講兩次.
證明保羅在這裡.
真的很著緊告訴我們.
舊有的扭壞的.
是可以將來.
那個浩同吹響.
主再回來.
就會有一個驚人的.
大逆轉.
就是所有信主的死人.
都會有一個不扭壞的身體.
就會進到永恆.
甚至連死亡都被吞滅.
第五章就說.
死亡就是.
籠罩著亞當的後裔.
現在我們只是在.
價值觀.
生命裡面有一個作王.
不是我們整個的身體.
去改變了不死.
我們一樣會死.
但是將來那個日子.
連死亡都被吞噬.
那就是.
那個真正的永恆的.
永遠舊恩的降臨.
所以今晚.
透過這兩段關於亞當的描述.
給我們其實有一個.
很大的釋放.
至少我們不應該再懼怕死亡.
至少我們應該.
勇敢的過好我們每一日.
為福音去效力.
將人帶進.
新的生命裡面的.

$^{2081}$那個經歷.
就是與神.
復和帶來的那個福樂.
就是我們人生.
大家一起努力的照明.
努力的大人認識福音.
脫離亞當的咒詛.
就是我們.
每個人一起可以在神.
國度裡面參與的.
那個服事.
今晚和大家講到這裡.
喉嚨有點乾一起祈禱.
再一次感恩.
天父給我們能夠透過五章.
頭那段落.
講到六方面的.
新的生命的.
福樂經歷.
與主的同在.
帶給我們新心靈的.
那種的輕省.
那種的.
不應該叫超脫.
但是一份很大的心靈的平安.
更加讓我們看到.
保羅給我們有很.
重要的透視.
人類的始祖的犯罪.
帶來的那個背役.
帶來的那個罪的籠罩.
死亡的威脅.
令到我們人類.
都不知道原來在這樣的.
救助底下.
多謝主在救恩在耶穌裡面.
是你將我們帶進.
這個永生.
這個永恆的盼望.
願你得到一切的榮耀.

$^{2121}$更加讓我們能夠享用.
救恩帶給我們的平安.
福樂和我們一班基督徒.
因為這個是集體的福音.
我們一起努力.
去傳揚福音.
讓更多人認識主.
我們的救贖.
我們這樣感恩禱告.
奉耶穌基督的名.
阿們.
第六章下個禮拜.
要探討的就是.
字面的意思大家可以先讀第六章全章.
可能讀多幾次.
由不可犯罪.
不准犯罪.
到反過來就是.
可以靠主不犯罪.
多謝大家.
\newpage



\section{羅馬書}
\label{sec:7LZOyo5Zt3c}
\textbf{2025-06-19 周四晚查經班《羅馬書(14):分題:不可犯罪?可不犯罪》 講員:曾裔貴牧師}
\newline
\newline
連結: \href{https://youtube.com/watch?v=7LZOyo5Zt3c}{\texttt{https://youtube.com/watch?v=7LZOyo5Zt3c}} ~~~~ 語音日期: 2025-06-20
\newline
\newline
\hyperref[sec:G8aDRsOGytw]{\small{< < < PREV SERMON < < <}}
~
\hyperref[sec:index_chronic]{\small{[返順時目]}}
~
\hyperref[sec:index_scriptual]{\small{[返順卷目]}}
~
\hyperref[sec:wh3ZqIs7b68]{\small{> > > NEXT SERMON > > >}}
\newline
\newline
$^{1}$好 各位聽眾姐妹 我們一起同心同道.
感恩再一次我們開始《羅馬書》的閱讀.
盼望我們去到第14章.
關乎到第5,6兩章.
我們很盼望神在我們中間帶領我們.
讓我們能夠認識你自己的真理.
盼望第6章的經文.
我們能夠可以充分的掌握.
還有我們真的去實踐保羅的教導.
願主幫助帶領我們.
我們這樣感恩禱告.
奉耶穌基督的名 阿們.
好 再一次我們盡量.
各位姐妹如果在線.
因為我現在沒有開到那個畫面.
如果在線你們在聽.
就記得靜音.
那就方便我們大家一起靜心的聆聽.
那就開始.
上個禮拜我們有提過.
關乎到第5章.
第2個段落12到第21節.
很明顯就是這個首先的阿當.
就是人類的始祖.
和末後的阿當耶穌基督.
兩者或者兩個人.
他們兩個都是各自一個代表.
就是舊有的人類的那個族類.
和在舊恩耶穌被差遣來這個世界.
開展舊恩之後.
就是引進一個新的人類.
新的生命.
所以保羅在這裡講了.
六個在主裡面同神復和之後的.
那個新的生命.
進入一個新的經歷的時候.
可以享用.
我可能聽不懂.
(記者: 因為現在是星期三).
(記者: 星期三收到的).

$^{41}$不好意思 如果要靜音先 不要介意.
我經常都碰不到那個位置 很煩.
是 先靜一下.
是 好的 謝謝.
我們繼續回到這個講道.
保羅很巧妙地講了.
那個新的生命.
同神復和的新的開始.
就是每一個都有一個新的開始.
就是由信主與主耶穌結連.
接著就是有聖靈內住.
其實是同時間的 不是有先後次序的.
就是與神復和.
我們明白了之後.
其實我們立即開展了.
永恆的生命就是與神的同在.
所以這個新的生命.
是由主耶穌的犧牲復活.
所帶來給我們的.
但是保羅在這個時候.
他帶回我們看那個根源的問題.
就是為什麼需要救恩.
為什麼需要救主來這個世界.
就是人類的始祖.
阿當他那個違背神的命令.
而他的背影就帶來.
因為他是一個代表.
所以所有在他以後出現的人.
都是同樣在什麼的權勢呢?.
就是罪藉著這個背影來掌權.
死亡又因為罪的緣故.
就被神宣告.
所有的人都要面對肉身的死亡.
包括將來在永恆裡面.
面對神的審判就是永死.
所以你可以想像得到.
這是一個神必須要處理的救贖.
所以保羅真的是.
再一次帶我們去創世記.
不是帶我們去看《亞伯拉罕》.

$^{81}$是帶我們去看人類的那個開始.
既然從那個源頭就已經出現了那個問題.
他告訴我們有一個新的源頭.
就是新的耶穌.
就是新的阿當就出現了.
所以你看到那個分段.
其實真的很想告訴我們.
沒有其它的方法可以將人去拯救.
就是間接告訴我們.
等於是對當時羅馬帝國所宣稱的.
帝國帶來的平安.
帶來的福樂.
帶來的人類的幸福.
那些只是地上的表面的.
甚至嚴格來說是假象.
因為其實只不過是少了戰爭的意思.
是這個強權的帝國奧古斯都.
來做統治者.
但其實他一樣是很暴虐地虐待非我族類的.
但是新人類就不同了.
主耶穌是用救恩,用愛.
將人帶進那個新的國度.
而且是跟創造的主有這個相和.
所以你會看到.
保羅是在做一個對比的.
這個對比其實不複雜的.
不難去看到的.
我們看回上個星期不夠時間.
我幫大家做一個簡單的分析.
中間綠色這個部分.
就是講到那個起源.
那個行動.
當然講兩個代表.
阿當.
舊的.
左邊這個就是舊的阿當.
右邊就是新的阿當.
就是耶穌.
所以.
他做了這個行動帶來什麼結果呢?.

$^{121}$這個結果又對人類有什麼影響呢?.
以及神因為這個結果有什麼宣判.
有什麼審判呢?.
以及到底誰是擁有了統治呢?.
其實很簡單.
就將這個段落.
12到21節.
很清楚地講了.
從左邊開始.
他就是人類的始祖.
最終是從一個人進入了世界.
所以他是人類的始祖.
他一次而已.
吃了分別善惡樹的果.
其實就是背逆神.
不聽,不服從,不信從神.
他一次過範而已.
不要因為他930歲.
很多次加起來.
不是.
他第一次不聽神的命令.
聽蛇的引誘.
他一次過範.
就是他的行動.
因為他這個過範.
就進入了世界.
就是說.
進入了地上的世界.
其實不用拍了.
我今晚都會有這個PowerPoint的.
大家可以看回.
所以.
所有的人.
就是人類.
所有的人類.
都成為罪人.
就算不同的多一樣的罪.
都成為罪人.
都在罪的權勢下.
所以第六章.

$^{161}$就帶我們去看.
新的人類.
是勝過罪惡權勢.
所以我們找到那個方程式.
第六章就講那個.
你可以說那個被決.
所以.
這是一個結局.
是已經成為事實.
無法改變的.
信不信耶穌.
我們都是成為了罪人.
在罪的權勢下.
所以.
我們只能夠因為阿當的緣故.
我們不信耶穌的.
就被神定罪了.
這就是審判.
當然了.
死是藉著罪去作王.
所謂作王.
就是說.
人人都要經過死亡.
但是進入的就是永恆的黑暗.
所以這就是阿當帶來的.
一個這樣對人類的禍患的影響.
所以.
我們看到.
現在保羅想講的救恩的對比就是.
告訴我們.
這位不是阿當的子孫.
他是神的兒子.
所以他是降世.
由童貞女瑪利亞降生.
他是神的兒子.
是.
我先拿走這個.
不要妨礙我們.
放在上面.
那麼.

$^{201}$他一次的兒行.
就是十字架上.
為我們無條件的.
深知痛苦的無助.
被神氣絕.
被神刑罰.
他要擔待人類的罪.
這是一次的耶穌釘十字架的兒行.
那麼.
神的恩典.
藉著耶穌.
就可以臨到所有相信的人.
所以是臨到相信的人.
不是所有地球的人.
是相信耶穌的人.
才可以得到這份恩典.
所以.
眾人就因為耶穌的救恩.
因為信而帶來清義.
這就是耶穌所做帶來的影響.
這就是第五章告訴我們.
神有什麼宣判呢?.
就因為這個人的緣故.
恩典現在可以稱義作王了.
就是說.
現在在神的裡面.
永遠.
當他還沒回來之前.
永遠還可以有月立.
有拯救.
有恩典.
不是一個審判.
不是一個定罪.
所以.
你看到第一章的十八節就是說.
一開始保羅告訴我們.
其實我們現在人類的光景.
因為阿當的緣故.
我們現在已經在神的憤怒之下.
列舉各種各樣的敵黨神.

$^{241}$無論拜偶像.
或者是一些道德上的罪惡.
是被定罪的.
但是現在.
因為在耶穌裡面.
一個人的義行.
就帶給我們得到這個清義.
而且是恩典.
來藉著義作王.
所以不再是死作王.
是恩典.
藉著耶穌的義行.
來作王.
就是說.
恩典現在是統管著我們的生命.
地上各方面的眾生.
所以.
我們現在為什麼能夠在神裡面.
永遠還有恩典.
所以我們千萬不要看自己是絕望的人.
不要看自己是無助的人.
譬如今早我說的.
禿鷹團契.
我說這個撒瑪利亞的婦人.
耶穌點出她的內心的真相.
其實是要挽回.
是要接納.
是要寬恕她.
是要讓她生命有盼望.
當她感受到這一份愛和盼望的時候.
她扔下那些水罐.
衝回城裡.
告訴城裡的人.
你們快點出來看.
因為有一個人將我數來所做的事情告訴我.
而且這個話不是指責.
不是批判.
不是來去論斷.
而是接納.
而是寬容.

$^{281}$所以這個婦人.
裡面有一份很大的生命力.
就是盼望.
她邀請城裡的人出來.
他們來到見到耶穌.
就邀請耶穌在城裡面.
在撒瑪利亞城住了兩晚.
後來就跟這個婦人說.
現在我們信.
不僅僅是因為你的勸導.
現在是因為我們見到.
聽到耶穌.
祂說的真理.
我們現在相信了.
更加知道.
祂就是救世主了.
所以這個婦人在很大的恩典的接納裡面.
是耶穌點出了她的真相.
耶穌說.
你給我水喝吧.
這個婦人說.
先生你沒有打水的器具.
你怎會有什麼呢?.
這個井是我們的祖宗留給我們的.
耶穌就說.
如果你知道我是誰.
你早就求我給你一杯水.
這個婦人就說.
好啊.
先生你給我喝這杯水吧.
耶穌就說.
你叫你的丈夫來.
這個婦人第一句就說.
我沒有丈夫.
耶穌怎麼說呢?.
你已經有五個丈夫了.
你現在家裡的都不是你的丈夫了.
就是六個.
這個婦人就很吃驚.
為什麼她會知道呢?.

$^{321}$馬上說第二件事.
我們知道你們猶太人拜神在耶路撒冷.
我們撒馬利亞人敬拜神就在基尼森山.
我們知道我們所拜的.
你們所拜的你們知不知道呢?.
避開那個話題.
不要再碰我的內心.
但是耶穌不放棄.
就跟她說.
跟你說話的就是那個彌賽亞.
這個婦人不是因為被人奚落.
不是因為被人點出她的罪惡.
而是耶穌讓她看到她自己的內心.
心靈的痛,需要.
而且是給她有光,有盼望.
是有愛的一個恩典的接納.
所以你看到耶穌帶來的是恩典作王.
是恩典讓人稱義.
是恩典讓人可以在神裡面.
得到生命盼望和愛.
所以在亞當裡面.
逃避罪的蔓延.
最帶來的就是更加互相敵視.
藏起自己.
不敢承認自己是怎樣.
很多很多的問題.
所以因為罪和死.
來籠罩著人類.
保羅讓我們帶我們回去.
現在連那個源頭都可以去更新.
是有一個新的源頭就是基督.
就是我們的主耶穌.
或者叫新的亞當也沒錯的.
所以這一段.
真的將我們帶回去.
那個根源怎樣解決的問題.
所以現在我們就能夠從第六章開始.
這裡大家沒有問題.
我們就從第六章開始.
去看回這一段聖經.

$^{361}$我太激動了.
我慢慢說吧.
因為都怕喉嚨不舒服.
這一段先用回舊的和合本.
大家先聽一聽.
但是我今天想跟大家去看的.
就是新漢語的譯本.
因為它照原文的直譯.
有些部份是比舊的和合本.
相對會好一點.
當然既然是直譯.
就不會說文法等等.
所以有時可能會有些掉漆.
但是某些字眼.
它的翻譯可以作為一個參考.
都是權威的新的漢語.
第六章大家先聽一聽這一段.
第五章如果大家還有什麼問題.
想搞清楚它.
可以在此之下再問我.
但是無論如何.
舊有的體系.
舊有的根源.
最帶來的.
已經在耶穌裡面過去了.
新的源頭開始了.
我們在信主裡面.
我們現在已經被納入.
在新的族類.
就是新亞當之下的人了.
所以叫做新造的人.
保羅經常說的.
在《以弗所書》.
《哥林多後書》五章的十七節.
若有人在基督裡面.
他就是新創造.
過去的事已經過去了.
剛才那個例子.
撒瑪利亞的婦人.
她知道耶穌點出她.

$^{401}$其實是寬恕她.
使她完全可以沒有了她的過去.
她可以由新的開始.
這樣來做一個人.
第六章羅馬書.
它說這樣.
這個是Dare for.
就是說.
所以也有一個這樣的結果.
這個結果就是上文說的.
現在保羅說得很重.
說得很多的就是.
恩典廣泛地作王.
就是說.
剛才我也說了.
神閱納我們每一個罪人.
在救恩裡面.
我們是無條件獲得寬恕.
而且恩典是不斷的恩上加恩.
可能會衍生在.
保羅也要信徒要很小心.
特別是羅馬的信徒.
不要將救恩.
成為了一種.
可以在地上成為一個安殊的避難所.
來過一些只是開開心心.
平平順順的生活.
因為整本羅馬書是一個很激動的.
萬民都需要福音.
萬民都是要恩順清義.
萬民都是在神的就坐下.
所以如果沒有人去傳福音.
萬民聽不到福音.
就在那個就坐下滅亡.
所以我們得了救恩.
在恩典裡面.
不是叫我們停留在恩典裡面.
而是我們要將這個恩典.
成為一個你可以說是一個武器.
一個福音的機動的去傳開.

$^{441}$所以仍然都是那個主題.
羅馬教會要好像生生勝勝那樣.
跟隨保羅向西班牙宣教.
或者支持保羅和他們的同工.
向歐洲其他地方.
沒有福音到的地方去宣教.
所以寫這本書本來的目的.
就是鼓勵羅馬教會.
要有這樣的裝備.
一起來向外宣教,傳福音.
所以恩典藉著二作王.
就不等於我們繼續停留在恩典裡面.
重重複複的犯罪.
過一些我們過去的生活.
所以第六章就是新的生命已經有了.
第六章就是說我們應該在這個新的生命裡面.
怎樣有新的生活呢?.
就會是這個六章.
可能先去到第十四節.
因為第十五節.
因為我一直都沒有時間預備.
或者我們根本不夠時間講到第十五節.
我們看第一到第十四節.
大家靜心聽一聽.
有些字眼都挺神學的.
我們就學習進入裡面的討論.
保羅其實都很想告訴我們.
我們今天信的耶穌是復活的.
我們怎樣在祂的復活裡面.
有份重新開始這個新的概念的生活呢?.
就是你下個月就要去美國了.
或者加拿大.
美國現在沒有那麼多人喜歡.
去加拿大或者去英國.
你開始要想一下那邊的生活.
他幾個月就會很冷.
你開始要這樣想的了.
這些是準備很正常的.
現在我們已經進入了基督的新生命.
我們要告訴自己.

$^{481}$我們已經不是當年的舊有的族類了.
我們是基督的新的族類了.
所以第六章第一大段落.
一到第十四節.
就是想告訴我們.
在新的族類裡面怎樣有新的生活.
一定是從神學開始.
就是說要通過我們自己的心態.
然後在我們實際的生活就去行動了.
就去表現出來了.
第一節開始.
這樣怎麼說呢?.
就是回應上文.
第21節.
21節就是說就如罪作王叫人死.
同樣地.
恩典現在不再被罪來操控人了.
而是恩典.
不然就操縱了.
恩典是讓信的人得著清義.
而且恩典是藉著義來包圍著信的人.
所以恩典是藉著義作王.
叫人因我們的主耶穌基督得永生.
在主耶穌裡面.
我們一直在恩典裡面進入那個永恆.
就是說我們的人生充滿了恩典.
恩典包圍著我們.
但是恩典包圍著我們.
千萬不要理解錯了就可以任意的犯罪.
所以保羅他可能真的意識到.
羅馬教會可能會有誤解.
我經常用一個例子.
我早上剃鬚就用剃鬚膏的.
用鬚刨.
就不是用那些電鬚刨.
就是那些刀片.
這樣刮的.
有時候你突然間想著另一件事.
要刮鬚的.
因為你不小心.

$^{521}$它真的是刀.
你刮鬚你就要塗藥膏.
你想一下有沒有人笨到一個地步.
突然間有人送一瓶藥膏給我.
這瓶是哪裡出品的.
哪裡的藥膏最好.
是不是澳洲還是不知道.
塗了就什麼都沒事了.
我會不會傻到每天都用刀片刺自己.
弄傷自己塗藥膏呢?.
所以第一節就是這個意思.
應該我們只是一個防衛.
所以犯罪不應該成為一個正常.
信耶穌犯罪就不再是正常.
就是我偶然間不經意弄傷了自己.
很不舒服.
很乾.
流血.
但是我沒理由塗了藥膏之後.
我每天都割一割自己.
再塗那種藥膏.
證明那種藥膏很有效.
我們覺得是沒有可能的事情.
但是信耶穌.
保羅第六章就想告訴我們這個問題.
他說這樣怎麼說呢?.
我們可以凝在罪中.
叫恩典顯幫忙.
果然是在說恩典.
你自己說的.
恩典要作王的.
叫我們可以因為主耶穌基督得永生.
我們就真的好像坐花轎上天堂一樣.
我們不需要擔心.
反正都有了永恆的保障.
我們還需要擔心嗎?.
他反而想說的是.
在我們仍然在生的時候.
我們怎麼能夠勝過罪.
那個權勢.

$^{561}$以及我們不再是他的奴隸.
這個就是新的概念.
新的生命裡面的.
新的族類的那個概念.
我們逐個點去看.
其實都列出來了.
有五個點的了.
我們逐個去看.
當然那個答案就是.
他自己問.
其實已經預計了是這個答案的了.
保羅.
他好像加重了語氣.
其實是可不可以的呢?.
我們繼續可以犯罪.
那恩典就繼續顯多了.
保羅已經立刻回答了.
絕對不可以.
Absolutely not.
完全不行.
接著他就解釋.
所以不僅是不准.
不僅是不行.
不僅是通常我們教兒女.
不准這麼做.
不要這樣.
他可能未必明白.
所以保羅立即解釋.
我們在罪上死了的人.
新概念了.
以前就是罪作王叫人死.
這是五章的十二節.
我們回頭看十二節.
幫大家一點點總結.
然後我們再看第六章.
第十二節五章.
所以有關聯的.
保羅讓我們知道.
為什麼我們是阿當的族類.
就如罪是從一個人進入了世界.

$^{601}$死也是從罪來的.
於是死就臨到眾人.
因為眾人都犯了罪.
就是說.
你一出生.
你就在罪的權勢下.
你一出生.
你就已經在罪上死了.
死了的意思就是說.
不是真的肉身的死.
而是我們已經對罪.
是沒有能力去抗衡的了.
我們已經在他的權勢下.
也是後面所說.
所以我們繼續讀下去.
第三節.
豈不知我們者受浸歸入.
或者受洗也可以.
歸入基督耶穌的人.
是受洗歸入他的死嗎?.
突然間引進一個.
似乎當時初期教會.
也是大家人盡皆知的禮儀.
就是洗禮這個禮儀.
洗禮這個禮儀當然是一個浸禮.
這個浸水.
就是告訴我們.
我們好像與主耶穌有這樣的生命聯合.
當然我們後世的人都知道.
是必須你裡面先有信主.
得到聖靈.
你裡面已經與主耶穌有這樣的聯合.
那麼我們的水禮就來得有意義了.
我們不是藉著水禮才得救.
才與主耶穌聯合.
所以先有真理的認順.
然後才有具體的行動來接受.
所以這個洗禮.
某程度上當時他們也是跟隨猶太人的洗禮.
例如施洗約翰的洗禮.

$^{641}$就是叫人悔改.
跟隨將來的彌賽亞.
施洗約翰就在約旦河外洗禮.
幫人.
猶太人去找他.
就叫他洗禮.
所以約翰也說.
我是不可以告訴你.
使你上天堂的.
你必須要相信.
我以後要來的.
那個耶穌.
彌賽亞.
你才可以得救的.
我幫他解鞋也不配的.
所以猶太人也明白的.
純粹是一個基本表面的悔改而已.
所以.
原來新約的時候.
初期教會.
已經有洗禮這個儀式了.
已經有與主耶穌聯合.
他們進入水裡面.
然後從水裡面上來.
代表著一個新的生活的樣式.
所以現在這裡.
保羅說到.
我們不再屬於舊有的體系.
就用洗禮這個禮儀.
來做一個比喻.
因為稍後再用奴隸的比喻.
現在就用浸禮這個比喻.
來告訴我們.
我們就是這樣與主耶穌浸在水裡面.
代表著耶穌的犧牲.
稍後再引一段經文.
我先讀完.
第四節.
所以我們藉著洗禮歸入死.
和他人同埋葬.

$^{681}$原是叫我們一舉一動有新生的樣式.
新生命的樣式.
要讓我們有這個思維.
我們是新造的人.
我們要有新生的樣式.
像基督藉著父的榮耀.
從死裡復活一樣.
好像什麼呢?.
基督復活的形象.
藉著父的榮耀.
他來到去復活一樣.
所以耶穌的復活.
就等於我們上來之後.
我們就有這種新生的樣式.
但我們還沒有耶穌的復活.
因為耶穌的復活是很超自然的.
他能夠繼續穿牆過壁.
繼續隨時可以出現在任何地方.
已經好像不受時空建築物的阻礙等等.
你就可以想像到.
將來我們復活的形體是很超自然的.
告訴我們.
這裡很明顯就是叫我們效法耶穌.
要和他有一個完全的認同.
完全的聯合.
所以你可以說.
信耶穌到了另一個階段.
就是我們要認同耶穌所做的一切.
其實我們信耶穌的時候已經認同了.
我們認同祂是我的救主.
我才會信舉手的.
所以現在保羅針對的是.
罪的權勢這個部分.
針對的是新的族類.
怎樣能夠抗衡地上.
我們仍然有舊有的身體和罪性.
怎樣能夠在有限的人生.
去抗衡這個罪的性情呢?.
這就是這一段聖經.
所以保羅說不可犯罪.

$^{721}$但是其實這一段聖經.
可以用四個字來做一個很簡單的分析和總結.
就是可不犯罪.
可以不犯罪.
就是這個段落裡面很重要要告訴我們.
當然了.
續點的秘訣就是這裡有五個大點.
第五節.
我們若在他死的形狀上與他聯合.
也要在他復活的形狀上與他聯合.
一直都在說要完全認同耶穌.
祂所做代表我.
我現在做就代表著祂的生命在我裡面.
所以絕對不是我們沒有了耶穌.
我們自己不斷可以改好.
是不可能的.
但是在耶穌裡面.
每個人都有他的可能.
這就是救恩的奇妙.
救恩的偉大.
第五節說.
也要在他復活的形狀上與他聯合.
第六節.
因為知道我們的舊人.
這個翻譯我待會用新的漢語來換.
因為知道我們的舊人.
和他同釘十字架.
使罪身而絕.
叫我們不再作罪的勞瀑.
開始說那個勞瀑的關係了.
就是說我們是一個自由的人.
關於奴隸.
羅馬帝國就很盛行.
差不多七成至八成的公民都是奴隸.
所以你不要想到奴隸就是戴上鎖鏈.
被皮鞭不斷打.
其實是一些自由人.
甚至有些是很有學問的.
所以在那個制度下.
假如有一個人買了那個奴隸.

$^{761}$釋放他.
他就變成一個釋放了的自由人.
那個主人就不可以再勞役他.
就算見到他都不可以再勞役他.
這個就是奴隸的制度.
當時的概念.
法律的地位.
所以保羅這裡很想告訴我們.
我們的確生而為阿當的後人.
阿當的族類.
但是因為在法律的層面.
在神的真理之下.
我們已經成為基督的聯合的人了.
所以基督是怎樣的呢?.
我們就在地上.
就可以不再受罪的勞役了.
所以仍然是罪作王.
現在就是罪不可以再勞役我們.
這個就是第六節.
告訴我們.
叫我們因為罪身滅絕了.
罪身滅絕六章的第六節.
漢語曆本.
要知道我們的舊人.
已經和他一同釘在十字架上.
留意一下.
罪的身體已經失效了.
不是滅絕 是失效.
叫我們不再作罪的勞服.
為什麼呢?.
因為死了的人已經擺脫了罪.
所以我覺得這個曆本.
是更加貼切原文.
舊的王合本也有它的意思.
但是因為罪身的滅絕.
到底是罪惡的身體.
何為滅絕呢?.
到底罪身是什麼意思呢?.
這個也是很值得在這段經文裡面要考究的詞彙.
到底是不是舊有的某些身體的一些不好的壞習慣.

$^{801}$等等.
滅絕了.
就是不再有效了.
失效了.
不再需要做那件事了.
可能有些人會是這樣.
待會我們先解讀.
因為我想未必可以解讀到第六節這麼長.
已經不夠時間解讀了.
第七節開始.
因為已死的人是脫離了罪.
我們若是與基督同死.
就信必與他同活.
記得這是將來.
Future tense.
第九節.
因為知道基督既從死裡復活.
就不再死.
死也不再作他的主了.
他死是向罪死了.
只有一次.
就是十字架.
他活是向神活著.
這裡告訴我們.
這樣就是一個Therefore.
你們向罪.
也當看自己是死的.
這個看字.
基本上原文已經不是叫你看.
叫你看的那個字.
這個原文的字.
是叫你Consider.
是叫你Reckon.
Reckon就是你要去計算.
你要去計算.
一加一等於多少?.
這是計算.
就是等於二.
是叫你Consider.
就是你要思想.

$^{841}$所以不是看一看.
後面也是這個意思.
第十一節.
是的.
沒錯.
第十一節.
我讀給大家聽.
其實我給了大家.
他說這樣.
你們也要看自己向罪是死的.
在基督耶穌裡.
向神卻是活的.
後面有解釋的.
這個就是他已經說了Reckon.
Consider這個字.
所以這個看.
不是叫你這樣看一看.
是告訴你.
拜託你.
耶穌是怎樣向罪死.
我們今天也可以效法耶穌.
我們也看自己.
在罪的面前.
我們是死了的.
在神的面前.
我們也可以效法耶穌.
我們可以看自己是活的.
所以這個看自己.
其實不是看一看.
是Consider.
是Reckon這個意思.
是要這樣來計算自己.
是已經在耶穌裡面.
我們可以向舊有的罪.
我們死了.
新的.
就是在神面前.
我們是活的.
第十二節.
所以不要用罪在你們必死的身上作惡.

$^{881}$果然就是這個意思.
如果我們這樣認同.
耶穌的死和復活.
他的能力.
在他的救贖裡面.
神的能力.
現在我們就是要看自己和耶穌完全一樣.
耶穌是在罪的面前.
罪是不可以克制的.
因為他一心一意要行神的旨意.
在神的面前.
他現在是復活的.
是看死的面前.
所以當我們來到這裡.
第十二節的總結.
很明顯.
真是抱歉.
弟兄姊妹.
如果你還沒有關那個.
就關一下.
謝謝.
我們繼續來看我們今天的經文.
所以其實重點就是.
不可以再用罪.
因為舊有的時候.
罪藉著.
當進入了世界.
死又藉著罪作惡.
所以罪和死的律.
這是第七章所說的.
在我們裡面.
所以現在.
保羅說.
原來可以有一個新的律抗衡.
就是你不可以再用罪.
在我們這個必死的人身上作惡.
使我們繼續順從.
身子的私慾.
也不要將你們的肢體.
就是我們的手腳器官.

$^{921}$獻給罪作不義的器具.
這裡說的是獻.
就是主動做了那個動作.
主動的意思就是說.
我們先有裡面的心思.
意念.
接著有行動.
這就是獻.
所以不要將我們自己.
現在新的族類的人.
我們各方面的器官.
器官當然包括心思意念.
不只是手腳.
獻給罪.
作不義的器具.
保要將從死裡面復活的人.
留意一下.
原來基督徒.
如果我們不主動.
不斷呈獻自己的生命給神.
我們就會留於過去的壞習慣.
所以要將自己獻給神.
所以不是全道人才獻給神.
是每一個新的族類.
在基督裡面的族類.
我們都可以獻自己給神.
並且將肢體.
果然當你有了決心獻給神的主動.
自然你的身體.
就會隨著你的意念去做.
我今天還是不想上茶經班了.
我想怎樣怎樣.
但是突然間你覺得.
不是啊.
裡面有感動.
我還是要回來聽.
你的意念就強過其他的.
比如有些人說.
我們去哪裡玩呢?.
當然玩不是錯的.

$^{961}$不過純粹是一個例子.
這未必是這麼貼切的例子.
可能極端一點就是.
有人繼續叫你去打牌.
你要回崇拜.
那你去打牌還是回崇拜呢?.
可能初信階段還不是很清楚.
你繼續去做.
繼續打牌的時候.
又說是非,又說髒話.
激氣的時候.
但是後來你知道其實不應該.
信了耶穌應該不要這樣了.
所以有一次.
突然間你裡面又有很大的感動.
就是我今天一定要回崇拜.
一定要來敬拜神.
你開始有了這個意念.
你就有一種感動.
你開始會穿衣服.
你選擇穿什麼衣服能夠得體的.
榮耀神的.
回去敬拜的.
你開始有各方面的準備了.
所以你會看到.
原來寶來真的與大相關的.
原來真的有知於內就有行動了.
當你有了將自己獻給神.
接著很簡單的.
就將肢體.
就是你自己的身體.
作為器具獻給神.
你自己就會將自己的身體整理好.
獻給神.
好像一個呈獻給神的器具.
一個義的器具.
譬如我今天回去.
是要做施事的.
我回去就準備好自己.
我來到去.

$^{1001}$來準備跟隨施事長的指引.
一起禱告.
帶著一個笑容.
來服侍今天回來的人.
派了一個程序表.
可能你那天是很不開心的.
但是你知道.
你今天已經呈獻給神了.
你就將你的肢體.
作義的器具獻給神.
罪必不能作你們的主.
因為你已經在律法之下.
乃在恩典之下.
不只是一種公式化的.
受著一個好像舊有的觀念.
舊有的想法.
來做這個行動.
而是在恩典之中做.
所以所有甘心樂意獻給神的.
都是明白在恩典之中.
神將我們無條件的拯救.
我們能夠有今天這樣的為人.
這樣的蒙恩.
所以我們不斷主動地呈獻.
將自己呈獻.
很自然你的身體.
你的肢體就會做義的器具來獻給神.
因為我們是在恩典當中.
恩典管束著我們的人生.
所以第一個段落.
其實就是正正就是在耶穌裡面.
怎樣去開始我們新的族類.
新的生活呢?.
大前提肯定就是.
怎樣讓這一班基督徒.
能夠過一個在神眼中的公義的生活呢?.
在神的眼中生命的那種綻放的光彩呢?.
有別於當時羅馬帝國.
他們那種的醉生夢死.
去擄掠不同的敵國.

$^{1041}$搶了人家的物資.
殺光那些士兵等等.
然後回到大肆的放縱慶祝.
就代表著帝國的和平.
沒有再有戰爭了.
甚至那些貿易,通商.
都能夠條條大路通過羅馬了.
但是這些只是地上由帝國.
由強權,征服帶來的所謂和平.
所以我們上個星期說的那個Women Peace.
羅馬的承平.
100年了.
歷史學家告訴我們.
羅馬帝國最承平的時候.
就由奧古斯都開始.
去到大概150多年的時候.
就是那個承平.
就是太平盛世.
是羅馬的承平.
但是耶穌正正告訴我們.
在祂的國度裡面.
是真正的平安.
真正的福樂.
所以保羅肯定要帶我們去回到醉這個問題了.
帶我們去看回人類根源的問題.
所以一來就是神將我們所有人都已經滅絕了.
一個新天新地.
那就沒問題了.
但是現在我們仍然是在這個地上.
我們是有了那個靈魂的復活.
永生.
但是還不是真正整個的.
身心靈的肉身的復活.
還是將來的.
所以既然我們都是亞當的族類.
死和罪仍然是在這個世界.
他仍然是掌權在一些未信的人.
和信的人身上.
現在保羅就教我們.
怎樣有另外一個律.

$^{1081}$可以不再在罪這個權勢下呢?.
所以就是保羅在第六章這裡.
就帶領我們一起去探討這個課題了.
所以就是我們今天要開始講的.
先用這個喝一口水吧.
不用抄下來.
我今晚就會給大家這個PowerPoint的了.
開始講這個歸屬聯合.
保羅就一個很寶貴的例子.
就是說當時代洗禮是一件很榮耀的事.
每一個基督徒要排除萬難.
羅馬帝國 猶太教.
或者你是其他的外邦宗教.
都有他們的儀式.
入會的儀式.
但是基督教就是告訴我們.
這個洗禮原來是有一個很重要的真理.
就是耶穌替我們進入那個永恆的死亡.
使得我們能夠因為跟祂完全的聯合.
就可以得到這個新生了.
所以是新生命的樣式.
就是在第六章不斷要探討的這個課題.
新的樣式就是我們可以不再在罪的權勢.
不再是祂的奴隸.
剛才我說了.
在洗禮有另外一段經文.
都相當重要.
可以跟這裡有一個呼應.
就是《馬可福音》第十章38節.
就是耶穌臨死之前.
當時門徒就跟耶穌說.
35節開始.
第十章的35節.
就是其中兩個門徒很想跟主耶穌.
同坐在天上的寶座.
西比泰的兒子雅各,約翰.
兩兄弟進前來.
後來雅各就殉道了.
約翰也是最年老的.
被流放到拔摩海島.

$^{1121}$在地中海一個很偏僻的海島.
對耶穌說.
夫子我們無論求你什麼.
願你給我們作.
又沒有說是什麼.
所以都沒有那麼快答應.
要求什麼.
要我給你們作什麼呢?.
第37節.
居然兩兄弟都敢開口這樣說.
到底是不惜好歹還是怎麼樣呢?.
到底那個心路歷程是怎麼樣呢?.
我事奉這麼久都沒有人不告訴我.
找我 曾牧師.
我們兩兄弟都一定要做陀董.
一定要坐在哪個位置.
我要做副主席.
我要做師父.
暫時沒有人敢找我.
很感恩.
一來都不是這樣找的.
當然是開玩笑的.
我們說的是.
果然是耶穌真的聽到這樣的要求.
很搞笑的.
第37節.
他們說.
次我們在你的榮耀裡面.
不是現在.
榮耀裡面.
將來.
一個坐在你右邊.
一個坐在你左邊.
果然是最重要的兩個位置都給了他們.
兩兄弟.
其實是不是很搞笑呢?.
到底他們在說什麼呢?.
其他門徒會死的嗎?.
其他門徒都會抗議的.
為什麼兩兄弟會突然間這麼飆出來.

$^{1161}$成為了一個前列分子呢?.
不過38節這段經文很值得我們去反思.
耶穌一方面.
可能心裡面都會笑.
但是耶穌帶出一個很重要的真理.
就是那個認同.
就是那個聯合.
而且是那個苦難上的聯合.
留意第38節.
你們很想在榮耀裡面坐在我的左右.
第38節耶穌說.
果然.
雖然他是半戴嘲諷.
你們不知道所求的是什麼.
其實怎麼會不知道呢?.
就是想坐在左右.
接著耶穌說什麼呢?.
他不是說將來的榮耀.
是說很快就要面對的那個苦杯.
那個死亡.
是由死亡進入榮耀.
由犧牲進入德性.
是由痛苦進入將來的永恆的福祿.
所以現在耶穌跟他們說的.
你現在就要認同我這個苦杯了.
所以第38節.
耶穌說.
你們不知道所求的是什麼.
我所喝的杯你們能喝嗎?.
我所受的浸你們能受嗎?.
他們說.
我們能.
耶穌說.
我所喝的杯你們也要喝.
我所受的洗你們也要受.
只是坐在我的左右.
不是我可以賜的.
乃是為誰預備.
就賜給誰.
當然十個門徒聽到的話.

$^{1201}$就非常老老雅各和若寒了.
耶穌就叫他們來.
對他們說.
你們知道外邦人有尊為君王的.
自利他們.
有大臣操權滿足他們.
只是在你們中間不是這樣.
你們中間誰願為大.
就必作你們的用人.
在你們中間誰願為首.
就必作眾人的服人.
因為人子來.
並不是要受人的服侍.
乃是要服侍人.
並且要捨命.
作多人的屬.
所以耶穌帶他們看的是.
你們很想認同.
你們必須要認同.
因為你是我的跟隨者.
是我的門徒.
這裡就帶出了一個很深的真理.
保羅正正就是在說這個認同.
但耶穌藉著兩個門徒.
他們很失禮的尋求.
沒錯.
可能將來你們得到榮耀.
神真的安排你們兄弟倆坐在那兩個位置.
但現在我不可以代表神.
但現在我可以代表著.
我們這個世界要進入死亡.
去將人類拯救.
所以你們是我的門徒.
你們都要認同這個苦杯.
都要認同這個洗禮.
所以這個洗禮在這段經文.
是在說進入那個苦難裡面.
不是普通的洗禮.
所以耶穌說你們都要接受.
所以這裡的認同.

$^{1241}$就等於是羅馬書提到的聯合.
就是保羅所告訴我們.
我們就要跟耶穌有這個完全的聯合.
其實已經聯合了.
我們已經透過洗禮做了這個動作.
不過在那個教義真理上.
我們當面對罪的時候.
我們就要這樣來看自己.
這個新的身份.
所以這裡真的是我們一起去學習.
所以保羅用一個當時全國初期教會.
共同的那個禮儀.
每一個信主的弟兄姊妹.
無論多困難.
他們甚麼背景.
例如外邦的那些錯誤宗教的.
例如拜阿波羅.
拜這個亞底米.
或者其它的宗教的那些偶像.
他們棄假歸真.
都要一個代價.
或者他們在羅馬帝國底下.
如果要跟隨耶穌.
可能都有一定的代價.
如果猶太教就更加了.
等於叛教的.
所以他們願意接受那個洗禮.
代表著他們已經是完全明白.
跟耶穌聯合了.
所以保羅很快就引這個洗禮的例子.
去帶出那個聯合的真理的意思.
所以這裡值得告訴我們.
我們回來看這個功課.
就是保羅首先第三節說.
我們受浸歸入基督耶穌的人.
就是受浸歸入他的死.
也都是剛才馬漢福音主耶穌所說的.
你們不知道你們求什麼.
但是我所受的那個苦.
你們都要受.

$^{1281}$我所面對的那個洗.
你們都要面對.
就是說很快了.
因為當時是受難周.
很快我要去到被捉拿.
被審問.
被刑罰.
被釘十字架.
其實你們都要這樣認同.
我這個救贖.
所以這個洗禮本身就是一個完全的認同.
完全的聯合.
所以保羅在這裡就告訴我們.
第四節就是這裡了.
所以我們藉著洗禮歸入死.
我們開始看這個新漢語了.
第四節.
所以我們藉著洗禮歸入死.
和他一同埋葬.
和耶穌一同埋葬.
是要我們行事為人.
有些不同了.
行事為人.
所以是在說我們新的人生.
那個生活.
是要按照生命的新樣式.
那就很清楚了.
新的樣式.
就是要和耶穌這樣的新生命.
已經在我們裡面了.
已經在每一個信徒的內心了.
聖靈.
所以我們只是配合聖靈在我們裡面的引導.
我們就可以在行事為人.
記得是按照生命的新樣式.
這樣來生活.
正如第四節的下半節.
正如基督藉著父的榮耀.
從死人當中復活一樣.
主耶穌就是這樣從死人當中復活.

$^{1321}$是藉著什麼呢?.
藉著父神的榮耀.
藉著做物主.
創造的權柄,能力和榮耀.
這樣來復活的.
所以耶穌的復活是來得很自然的.
第三天神定規了祂要復活.
那麼耶穌就從墳墓出來.
所以也不用說誰吹噓祂.
不需要了.
神的話語就是那個權能.
所以既然神在基督裡面這樣看我們.
我們也要這樣看自己.
我們已經是藉著洗禮認同了基督擔待了我的罪.
我承認祂是我的救主.
是我的生命的主.
所以這個認同就告訴我們.
原來要在新的生活裡面.
表現出新的樣式.
有別於我們過去舊有的生命.
所以這裡就告訴我們.
其實我們的行事為人.
新的概念,新的族類,新的樣式.
以前打牌的現在不打了.
那隻手我們就有決心地為神去工作了.
有些人打慣了打牌的時候手癢.
玩十三章,啤牌.
或者有些心癮突然間就會坐船去澳門.
信了主耶穌.
他就知道原來有一個新的生命.
但是那些也會有隱憂的.
他的手真的會手癢的.
但是他就會記住.
保羅告訴我們.
當我們跟主耶穌完全認同了.
以及在新的族類的生命裡面.
我們就要在行事為人.
按照著生命的新樣式去做.
也就像基督在賦的榮耀.
由死復活一樣.

$^{1361}$也就是說我們都可以藉著賦的榮耀.
當然就是聖靈的力量在我們內心.
我們就可以過一個新的樣式了.
第五節.
因為既然已經與他聯合.
像他那樣死了.
我們在耶穌的十字架的犧牲.
我們認同了.
祂是為我而死.
就是說我原先本來要在十字架死的.
但是耶穌擔待了我.
所以既然我們有這樣的認同.
因為已經與他聯合了.
像他那樣死了.
接著就是第五節的下半節.
就告訴我們.
那麼我們也必像他那樣復活了.
這是第五節的羅馬書告訴我們的.
我們有一個學習耶穌這樣的死.
和祂的復活.
所以我們現在.
這個復活是將來的.
身體的復活是將來的.
但是現在我們在地上.
死的聯合已經發生了.
但是復活的聯合在我們身上.
其實真正的實現.
就是我們要有一個新生命的樣式.
新的樣式來生活.
而且這個新的樣式.
是在神的能力榮耀下.
這樣來做的.
有沒有一些例子.
能夠讓大家可以比較具體明白.
比如剛才說的那個替數的例子.
或者是我們習慣了做一些事情是神喜悅的.
信了主耶穌之後.
你會有一種力量.
不如我說一些靈界的那些.
真人真事.

$^{1401}$信主的時候有一個很大的真戰.
它裡面的靈體,靈界.
是禁止它承認耶穌.
甚至操控它.
但是當它信了主.
當然很多人為了祈禱.
它接受了耶穌.
它終於肯和終於能夠呼叫耶穌是救主.
還有幫它趕走的人.
就斥令裡面的靈體要離開.
那個靈體是必須離開的.
但是在這個時候就有一個關鍵.
就是這個人如果他不接受耶穌在他的生命.
在他的靈裡面.
在他的內心.
那個靈體是會再回來的.
因為他已經過去一段很長時間.
依附了這個人.
或者這個人也開放了他的心靈.
讓那些靈體進入.
所以在這個時候的關鍵就是要斬斷那個連接.
首先就是如果他不肯信耶穌.
幫他一次 下次也會繼續有被鬼附的.
所以第一時間是要這個被鬼附的人.
趕快承認耶穌是他的救主.
接受耶穌進入他的內心.
有主耶穌在心靈裡面.
魔鬼就不敢再來騷擾了.
記住是騷擾 不敢再進入.
但是除非他自己繼續開放.
讓邪靈進來.
所以這個時候.
他會用一些不同的手段.
去騷擾那個曾經被鬼附的信徒.
可能用一些恐嚇.
製造一些意外等等.
但是那個被鬼附的人.
必須有一個堅定.
承認自己是新的生命.
屬於耶穌.

$^{1441}$要有聖靈在裡面的力量.
來定義抵抗魔鬼.
不再相信他 不再需要他.
甚至是缺絕的.
他是不可以再來.
慢慢地 那些鬼就永遠不會再回來.
就會走了.
這些是有些趕鬼的人寫的書的一些個案.
真人真事的個案.
告訴我們知道.
這些比較特別的例子.
但是一般人.
我們有時候有些不同的習慣.
所以第十一節就會說.
我們那些過去的習慣.
就要開始一個新的改變了.
所以一定牽涉到我們整個人.
所以我們要深思意念裡面.
要不斷地去更新.
才是說我們整個人不斷地成長.
所以我們聽完貝兒的金句.
我們沒有去認真地思想.
具體地落實實踐.
所以保羅他強調的.
過去我們怎樣獻給罪.
現在我們要主動地獻給神.
所以這個獻本身.
就是一個很主動對信仰的尋求的活出.
這個就是很明顯.
你會見到一個有心行神的妒忌 弟兄姊妹.
她的成長就很明顯了.
一個不想再行.
她就停在那裡了.
她可能有很多聖經的知識.
但是她不能夠實踐到.
那個新的生命帶來的新的樣式.
所以這裡就告訴我們.
這個聯合帶來的新的樣式.
這個就是第五節.
也是我們第一個要說的.

$^{1481}$歸入基督的死.
第二個就是與基督一同的復活.
記得現在我們在洗禮裡面.
我們進入那個死亡.
我們上水是一個新的樣式.
不是一個復活 未是一個復活.
真正的復活是將來我們離開世界.
真正的永恆復活是身體.
所以這個是將來 記得.
所以這裡說的新生的樣式.
就不是現在我們個個好像不在這個地球.
就不是 我們仍然是腳踏在這個地球的.
我們仍然是舊有的人.
但是怎樣來更新我們呢?.
第六,第七 一直這樣說下去的.
所以我們要記得.
這個是尚未完全實現的復活.
我們現在是在靈魂裡面.
生命裡面.
與耶穌完全的認同.
祂的死與我是完全有關.
祂的復活帶給我們是新生的樣式.
記得這件事.
所以將來才是永恆的復活.
第三個就是什麼呢?.
就是被基督釋放.
當然就是罪的釋放.
釋放就是說不再是奴隸 是自由人.
第六,第七節.
我們看漢語的譯本.
第六節 要知道我們的舊人.
已經和他一同釘在十字架上.
整個舊有的生命.
我們這個本來屬於亞當的人.
在亞當的體系底下的舊人類.
就是幾個舊人.
當然這個是我們的個體.
我們每一個的個體.
但是在一個集體來講.
亞當的族類的人.

$^{1521}$我們都是在這個舊人底下.
所以現在這個舊人.
我們認同了耶穌的十字架.
我們就在十字架裡面被釘死了.
使到罪身失效了.
它就不是滅絕.
是失效 罪身失效.
到底這個罪身有不同的解釋.
就是說罪惡的身體.
就是說可能就是那個觀念.
就是說可能是其他的一些舊有的事情.
似乎這裡和上下文的意思.
可以擺脫罪 可以有新的生活.
就是說可能就是哥林多後書那裡.
較為貼近這個罪身的意思.
就是第五章的哥林多後書.
這樣去理解罪身.
可能會相對那個神學的衝突沒那麼大.
就是第十 十六節開始.
或者可能我們從第十一節.
保羅怎樣理解他們現在這個積分.
和他現在做這麼多事情.
似乎好像帶來有些人不是太明白.
但是保羅說他就是要為主.
這樣去得到哥林多人的心靈.
第十一節 五章的十一節 哥林多後書.
我們既知道主是可畏的.
所以勸人 但我們在神的面前是顯明的.
就是說我們所做的.
是神監察和神知道的.
不是要得誰人的心.
更加不是要操縱誰人.
他說盼望在你們的良心裡面.
也是顯明 也就是說不顯明.
有些人誤解了保羅的動機.
甚至對他所做的帶來一個負面的批判.
甚至打對台.
所以保羅就說了這麼婉轉的話.
希望他們明白.
我們在神面前顯明的.

$^{1561}$你們不顯明 也就是說證明你們心裡有刺.
你們可能受別人唆擺.
所以導致我們好像不顯明的意思.
保羅說其實不是的.
第十二節 我們不是再向你們舉戰我們自己.
我再要這樣來不斷抬舉.
說自己多偉大.
保羅說你根本不需要這樣說的.
乃是叫你們因我們有可誇之處.
其實很想你們能夠因為認識了我們所做的.
你們有些誇耀自己的地方.
就是你看到我們完全是為你們的緣故.
為福音的緣故的這樣.
好對那憑外貌 不憑內心誇口的人有言可答.
就是說有些人只是靠那個外在的.
可能是身份 可能是社會地位.
可能在教會 可能在教會裡面.
是哪個的支派 哪一個派別的人.
太容易了 分黨 分門 結黨.
保羅說其實不需要靠那些的.
所以我們很希望你們能夠有東西可以誇耀.
就是誇耀我們是神的僕人.
為神的福音所做的.
當他們明白之後.
他們就會有一種這樣的諒解和接納和同心.
所以第十三節保羅就說.
當然他好像舉例子那樣.
你們看到我們好像發瘋似的去傳福音.
保羅說 我們若果癲狂.
就是說背後蘊含著很多.
他們哥林多人對保羅的誤解.
以為他們好像瘋癲了那樣.
我們若果癲狂是為神.
若果謹守是為你們.
就是說如果我們是很有sense的.
做事做每一樣東西都很步步為營.
循規蹈矩.
其實就是不要讓人有誤解.
不要讓人以為我們有特權.
這樣的謹慎,謹守.

$^{1601}$其實我們根本不需要.
我們是神的僕人.
我們是使徒.
所以第十四節.
這就是大家都很熟悉的.
原來基督的愛激勵我們.
因我們想一個日記替眾人死.
眾人就都死了.
並且他替眾人死.
是叫那些活著的人不再為自己活.
乃為替他們死而復活的主活.
你可以想像到.
你已經完全歸入主的死裡面.
我們現在只是一心要為主而活.
所以這就是一個很清楚的聯合的意思.
留意第十六節.
所以我們從今以後.
不再憑著外貌認人.
高林多人開始比較哪些強勢的權威使徒.
哪些領袖是哪些強勢的.
不斷在憑外貌.
憑那個表面的來誇口,誇耀.
開始追星 類似這些.
保羅說其實這些是背後的價值觀是不行的.
不可以這樣的.
這樣帶來只會分化和比較.
所以保羅說我不再憑著外貌認人了.
雖然曾經憑著外貌認過基督.
憑著外貌就是說.
他起初看這個拿撒勒人.
以為他只不過是一個木匠的兒子.
為什麼這麼大膽說自己是彌賽亞呢?.
就是連保羅自己還沒信的時候.
他都會這樣看耶穌基督.
何況你們在我們這些領袖.
很容易就會對號入座.
很容易就會喜歡誰就追到死.
不喜歡誰就生氣到死.
這樣是很大的對教會的傷害.
就不應該是這樣的.

$^{1641}$所以保羅說如今我們這一幫的使徒.
都不會這樣認他了.
當然真正的承認的就是.
基督的愛現在是激勵著我們去做所有的聖功.
是基督的愛.
所以你們在一個冷的心場.
在我們的愛就以為是發瘋.
就以為是在壟斷.
在霸佔哥林多教會.
其實不是的.
是基督的愛激勵我們去做.
我們根本沒有任何的意思.
要來牢籠你們的心.
但正是因為有其他的領袖.
進入了哥林多教會.
來做這些的分化.
做這些在他們面前來詆毀保羅.
這樣的說話.
就令到哥林多教會憑外貌.
不再憑內心.
或者憑聖力.
所以保羅在這裡告訴我們.
終於出現了這一節.
我們今天想討論的就是.
若有人提到罪身的問題.
若有人在基督裡.
他就是新造的人.
舊事 即是過去他所做的不好的習慣.
曾經犯過的罪.
過去 在基督裡是新的.
因為神已經饒恕,赦免,接納和淘沒.
所以舊的事.
就是在阿當族類之下的舊人.
所做的不好的事.
神不約立的事.
是已經過去了.
一切都變成新的了.
因為一切都是出於神.
神藉著基督使我們和他和好.
所以這個罪身.

$^{1681}$有可能是在說過去舊有某些的壞習慣.
某些的罪惡.
這樣的罪身.
所以罪身已經失效了.
不再成為我們的基本了.
所以這裡告訴我們.
罪身和下面說的.
因為要知道我們的舊人.
已經和他和耶穌一同釘在十字架上.
使到罪身已經失效.
所以是失效 不是滅絕.
失效的意思就更加具體一些.
就是不再是這樣了.
以致我們不需要再作罪的奴僕了.
這就是前因後果.
就是轉成上下文的意思.
舊有的因為已經進入了新的族類裡面.
由於我們認同耶穌的十字架.
祂的死與我們完全聯合有關.
所以那個舊人已經釘在十字架上了.
舊有的價值觀.
舊有的犯罪的傾向.
犯罪的意欲.
全部已經釘在十字架上了.
所以罪身已經失效了.
就是說可能我們部份曾經犯罪的身體.
一些舊有的牢籠著我們的.
已經失效了.
失效不代表完全滅絕.
就是說如果我們不在下面所說的.
要繼續將自己主動的獻給神.
這個失效可能慢慢又會恢復它的運作.
繼續蠢蠢欲動的拉我們回去.
其實你是阿當的族類.
你不要看自己是耶穌的兒女.
這麼多掣肘,這麼多規條,這麼多麻煩.
為什麼你不回來舊有的享樂.
那種的釋放呢?.
這個就是價值觀的問題了.
所以罪身失效.

$^{1721}$有可能就是我們要告訴我們.
舊有的整個價值觀.
都要這樣看自己在十字架上.
這個是要這樣認同.
然後那個結果就是叫我們.
這個六章的下半節.
對不起.
第六節的下半節.
就是叫我們不再作罪的勞瀑.
重點是不需要再作罪的勞瀑.
因為罪身已經失效.
過去的壞習慣.
過去的舊有價值觀.
人一世民一世.
你對我做初一.
我不對你做十五.
我就枉為人了.
很多這些價值觀.
我們被塑造要去敢愛敢恨.
新的族類就應該要有新的價值觀了.
所以記住是失效.
罪身失效.
以致我們不再作罪的勞瀑.
這個就是保羅想帶出的第四方面.
釋放.
被基督律師釋放.
第七節.
這個第三節不是第四節.
因為死了的人已經擺脫了罪.
因為死了的人.
所以罪身失效.
和我們與基督舊的人釘在十字架.
就是我們已經是死了的人.
擺脫了罪的基本.
就是說罪仍然在這個世界.
因為耶穌還沒有回來.
還沒有再回來.
神的永恆救恩還沒有完成.
地上是有罪.
我們都在一個罪惡的世界.

$^{1761}$但是.
所以千萬不要搞錯.
以為教堂是天堂.
因為那是物件中物.
我們這班人才是教會.
我們這班人在神的救恩裡面的認同.
我們才會有新生命.
所以你來到教堂.
你來到一些基督教機構.
那裡一樣是罪的地方.
所以有些人去完基督教機構工作.
就永遠撒手露頭.
原來是一樣的.
以為世界的人很詐.
你如我詐.
去基督教機構.
大家都是基督徒就很好了.
不然有些人真的很好.
所以千萬不要一概而論.
不過發現其實我有罪性.
那些基督徒的同事一樣有罪性.
和其他微信的機構是完全一樣的.
所以這次說的這個聯合.
是我們每一個的個人.
本來是阿當的族裡.
因為我們已經和主耶穌有死的聯合.
所以我們一定要這樣認同.
不斷這樣認同.
我們就是銅釘在十字架.
以致罪身失效.
那些過去的舊有的壞習慣.
舊有的事已經過去了.
不再和我有關了.
在神的眼中和我沒有關係了.
所以我們不需要再作罪的牢獄了.
原因就是死了.
擺脫了罪.
死了不是因為我們釘十字架.
是耶穌替我釘十字架.
所以這個聯合必須要這樣認同.

$^{1801}$所以耶穌就是這樣將我們釋放.
因為祂的犧牲.
祂的復活.
就將我們釋放出來.
我們就能夠.
因為已經死了的人.
脫離罪和脫離罪的勞瀑的身份了.
第八節就是將來必定與基督同活.
還有十分鐘.
我們可以多講一個.
第八節.
我們若與基督同死.
就相信也必定與他同活.
因為知道基督既已從死人中復活.
再也不會死.
死再也不能克制他.
他死是向罪死了.
只有一次.
他活是向神而活.
這就是第八到第十節.
我們要與基督同死.
重複地講這個很重要的聯合的概念.
不斷地講我們要歸入基督的死裡.
與他復活.
在他的新的生命裡面.
有我們自己新的身份帶來的新的樣式.
在這個新的樣式裡面.
我們原來是被基督釋放的.
我們不再是罪的勞瀑.
我們的罪身.
舊有的無論觀念也好.
壞習慣也好.
是已經失效了.
即是說已經過了期了.
不再有效了.
我們原來這樣的狀況.
就是已經脫離了罪的核制.
和不再是罪的勞瀑.
我們不需要再是勞瀑.
被主人不斷地監視,操控,奴役.

$^{1841}$這個奴隸本來他的身份是這樣.
但現在我們是基督的兒女.
是屬於真理的勞瀑.
是屬於神的勞瀑.
是公義.
下面第十一節說是義的勞瀑.
我們不再是罪舊有的主人的勞瀑了.
用奴隸這個比喻去說我們新的生命.
很全神的.
如果我們真是一個奴隸.
舊有的主人如果是很苛刻的.
很嚴厲的.
是很不斷地強迫我們的.
我們是痛苦萬分.
但現在有一個新的主人將我們買熟了.
我們屬於這個新主人.
舊有的主人就再沒有任何法律的地位.
去再操控我們.
阿貴 做事已經不行了.
因為我已經是屬於新的主人了.
不再是舊的主人了.
要這樣不斷地.
不是洗自己腦.
而是在這個真理下.
不斷地提醒自己.
我是屬於基督的新的主人.
我是和主耶穌有這個聯合.
祂已經將我釋放了.
所以保羅經常用初期教會.
那個時代的政治.
一定不是贊同有奴隸.
但是你沒有理由.
你傳福音第一時間就打破當時的奴隸制度.
那就差不多全國大亂了.
因為七成到八成都是奴隸.
那你怎麼打破它呢?.
那些事情是誰做的呢?.
羅馬帝國.
那根本就瓦解了你基督教.
所以保羅是用一個真理的教導.

$^{1881}$大家是平等的.
在神的裡面.
神現在將我們釋放.
不斷講那些奴隸.
這些字眼.
那你講得多.
其實可能你說你想煽動.
什麼什麼釋放.
又說奴隸.
當時真的有奴隸的.
那你可以想像到保羅.
不是介意別人怎麼看他.
他是將那些大家都明白的那些概念.
套用在這些永恆的真理裡面.
讓我們看到.
原來我們是舊有的體系.
當體系裡面聚的奴隸.
死將我們籠罩著.
但現在基督將我們釋放.
將我們買熟.
將我們拯救.
我們現在是義的奴僕.
是神的奴僕.
是愛的奴僕.
是基督的奴僕.
是其他人的弟兄.
這樣的一個新的概念.
那我們的生命就很尊貴了.
因為這個主人永遠是對我們充滿恩典的.
這樣來讓我們在生命裡面不斷的成長的.
這樣的福音.
就不再是一些很抽離的.
一些的頭腦知識.
而是一個新的生命.
等待我去活出這個樣式.
所以這個就是聯合帶來的.
所以第八節又是講這個聯合.
若與基督同死.
就相信必定的同活.
因為知道基督祂自己就是這樣死了.

$^{1921}$和復活了.
再不需要死了.
死都不能夠核祭祂了.
所以祂的死就是向罪.
一次死了.
祂活就是在神面前向神而活.
所以我們將來必定會同主耶穌.
有這個永遠的復活.
但現在我們就要認同主祂的死.
在我們生命裡面.
這樣我們就能夠在將來.
必定可以活在神的面前.
當然在講現在.
現在就是罪和死的權勢.
在我們這個復活的生命裡面.
它不可以再去搞擾.
記住是失效了.
這個失效是要我們主動.
這個就是第十一到第十四節.
下個星期我們再講.
然後再進入第七章.
保羅又再講一個在罪.
核祭下的那個人.
是怎樣的光景呢?.
和那個釋放.
然後就講到聖靈.
進入那個人的心裡面的釋放.
就是和這一段經文有一個息息相關的.
所以第十一到十四節.
就是我們主動將自己呈現.
就帶來那個聯合的體驗.
新的樣式的那種經歷.
所以你可以說第十一到十四節.
雖然是第五點.
但其實是最重要的實踐.
就是在耶穌裡面新生命的實踐.
最重要就是這樣看自己有這個可能性.
所以不要看自己是怎樣體驗贖活.
其實不需要.
我們已經認同主耶穌這樣的犧牲.

$^{1961}$我們的新心靈已經不斷有一種更新的了.
今天就講到這裡.
我們一起感恩禱告.
多謝天父.
雖然我們重重複複講一些很抽象的教義真理.
但是簡單來說.
保羅很想告訴我們.
我們已經因為基督新的阿當.
我們成為了新的族類.
新的生命.
我們需要在愛,在恩典裡面.
去告訴自己.
我們是一個集體的新人類.
我們是有新的生命.
我們是有新的樣式.
我們每一個都需要將自己呈現給神.
以致讓恩典在我們的內心來作旺.
不再讓罪和死來去克制我們.
盼望基督的犧牲.
我們在洗禮的時候.
是已經完全的認同和聯合.
願神幫助我們.
在這個段落裡面.
我們能夠掌握.
不斷重重複複告訴我們.
舊有的那個個性.
舊有的那個價值觀.
舊有的無論過去的事.
罪惡.
都在神的愛子裡面.
是已經成為新造的人.
一切都過去了.
不再被追究.
願主幫助帶領我們.
讓我們能夠明白之後.
我們每一個都一起去實踐.
這個新生命的樣式.
願主幫助.
我們這樣感恩禱告.
奉耶穌基督的名.

$^{2001}$阿們.
記得不可犯罪是那個標題.
但是整段的聖經是告訴我們.
在新生命的新樣式裡面.
是可以不再犯罪.
這個段落裡面很清楚.
要帶進的那個教訓.
我們就停在這裡.
多謝大家.
\newpage



\section{羅馬書}
\label{sec:wh3ZqIs7b68}
\textbf{2025-06-26 周四晚查經班《羅馬書(15):分題:不可犯罪?可不犯罪(二)》 講員:曾裔貴牧師}
\newline
\newline
連結: \href{https://youtube.com/watch?v=wh3ZqIs7b68}{\texttt{https://youtube.com/watch?v=wh3ZqIs7b68}} ~~~~ 語音日期: 2025-06-27
\newline
\newline
\hyperref[sec:7LZOyo5Zt3c]{\small{< < < PREV SERMON < < <}}
~
\hyperref[sec:index_chronic]{\small{[返順時目]}}
~
\hyperref[sec:index_scriptual]{\small{[返順卷目]}}
~
\hyperref[sec:PTM8SSdjvVE]{\small{> > > NEXT SERMON > > >}}
\newline
\newline
$^{1}$各位我們一起同心禱告.
多謝天父讓今天晚上.
我們再一次可以查《羅馬書》第六章.
幫助我們能夠對這兩個段落.
可以掌握保羅想要探討的.
不可犯罪.
以及不可以濫用恩典.
也讓我們能夠看到.
到底怎樣可以靠主得勝.
在地上,在有限的人生.
過一個土神喜悅.
以及將自己奉獻給神的生活.
求主帶領我們.
這些都不是說要怎樣去引修.
怎樣去過一些被世的生活.
而是在我們的日常生活裡面.
這是我們基督徒正常.
應該是正常的生活.
正常的基督徒生活.
就是要將自己獻給神的.
求主帶領我們看這些經文.
裡面都有相當關鍵的一些字詞.
我們能夠掌握.
雖然我們不懂原文的希臘文.
但是讓我們能夠透過不同的翻譯館.
都來去了解,明白.
我們再一次仰望,禱告,求主施恩.
奉耶穌基督的名 阿們.
上一個星期我們講到第十一節.
十一節我們還沒有開始講.
尤其是第十一節到十四節.
是很重的一個概念.
就是必須要主動獻上自己給神.
這個似乎是正常基督徒的生活.
必須的起點.
成為基督徒之後.
我們憑信心,靠恩典.
去得到神的救恩.
但是原來在正常的基督徒的人生.
因為我們已經歸屬了給神.

$^{41}$所以保羅常常都告訴我們.
我們要記住洗禮所象徵的意義.
就是由死復活.
你好 歡迎你.
隨便坐.
我們今天多了一位年輕人.
有些在線的 有些在現場的.
當然我們之後可以在網上的YouTube.
有我們專屬的頻道.
有羅馬書 由第一講到第十五講.
到第十四講 大家又可以再重溫.
另外的書卷都有的.
現在我都算是習慣.
可以盡量趕到做PowerPoint.
所以晚上回家之後.
都會把PowerPoint給大家可以就著那個講道.
來看一些經文或者一些重點.
大家來留意.
我們要看第十一節的時間.
我們看到第六章的十一節.
上個星期 今天應該可以完到第六章.
但是我們要重溫那個很重要的上下文的段落.
第四章結尾就已經告訴我們.
其實在舊約已經有恩信稱義.
恩信稱義絕對不是保羅的發明.
其實舊約的人有信心的亞伯拉罕.
在他未受割禮之前.
他已經憑信心得到神的.
對他的信的行為來去稱義.
就是他相信神 應許他.
雖然他當時未有子孫 未有兒子.
但是神說你的子孫將來.
好像天上的星 海裡面的沙那麼多.
亞伯拉罕信神.
神就以此為他的義.
就應許要將福氣給亞伯拉罕和他的後裔.
所以保羅都在第四章告訴我們.
恩信稱義不只是給某一些人.
例如猶太人 他自命是神的子民.
他們自認自己有割禮.

$^{81}$有律法等等 有摩西的律法書等等.
但是保羅說不是.
恩信稱義是萬民的.
每一個人都需要神的救恩.
恩信稱義是給每一個信的人.
以及萬民都是虧欠神的榮耀.
萬民都是在神的憤怒下.
所以每一個人都需要恩典.
然後講完第四章之後.
第五章的第一個段落.
就告訴我們原來我們恩信稱義.
其實是進入與神的復和關係.
與這位造物的主.
來到有一個生命的連繫的關係.
所以這個連繫 這個關係.
就帶來各種的生命的福樂.
我們立即會享受與神復和的平安.
站在神的那個.
應許的恩典.
嚴格來說 我們是去到神面前.
站在那個恩典裡面.
以及我們就算在患難.
都可以歡欣喜喜的誇耀.
因為患難只是讓我們更加的.
那個生命可以成長,成熟.
所以保羅告訴我們.
在我們信主成為基督徒之後.
很重要,很正常,不是特殊.
不是一些神秘主義.
而是我們應該每天發展.
與神相和的生活.
第五章頭的段落.
一到第十一節.
就盡量保羅都說了.
我們怎樣能夠以神為樂.
怎樣能夠在神的裡面.
等候那個真正的應許的日子.
接著到第二個段落.
五章十二節.
一直到結尾的二十一節.

$^{121}$就告訴我們.
原來人類共有的那個最大的問題.
就是因為我們是阿當的子孫.
人類的始祖.
因為人的犯罪.
被神來刑罰.
罪是從阿當入了人類.
因為阿當就是舊有人類的總代表.
所以阿當的背影.
就導致神的審判對每一個人.
所以我們一生來這個世界.
保羅就帶我們去看一個.
最重要的人類共有的DNA.
就是我們在神的.
因為阿當的背影.
就在神的奏座底下.
我們面對罪.
成為我們的核制.
罪會帶我們面對將來的審判的結局.
就是死亡.
所以阿當帶來的.
是我們好像沒有辦法.
你一生下來在這個世界.
物質的世界.
我們就要面對這個審判.
但是在這一段裡面.
保羅告訴我們.
他帶我們去看.
所有的東西.
最基本的就是阿當的出現.
在這個世界.
第一個人在這個世界.
他第一次的背影.
人類已經不能回頭了.
但是既然造物主是愛世上的人.
他就已經設計了.
預備了人類是可以得救.
所以是藉著耶穌.
就帶來一個扭轉.
怎樣扭轉呢?.

$^{161}$聖經不如我們讀一點.
回到第五章這一段.
第十七節.
他說如果因為一個人.
就是阿當的過犯.
阿當的過犯就是.
吃了不是禁果.
是分別善惡樹的果.
其實就是背逆神.
不聽神的話.
因為阿當的一次過犯.
是一次而已.
一個人一次的過犯.
死就因這個人來作王了.
死是因為阿當的犯罪而來作王.
因為神的咒詛.
所有人類都要面對死亡.
這是一個定律.
現在在地球上這個定律.
就是所有人都要面對死亡.
除非很特殊的.
有兩個人.
有沒有印象?.
以洛和以利亞.
兩個人不需要死.
連耶穌都要面對死亡.
直到待會.
因為他要代替.
擔當我們人類的罪.
所以你留意第十七節的下半節.
是很強烈的對比.
我想第十四講的那個.
PowerPoint就做了那個對比了.
上一次的PowerPoint.
第十七節說.
死就因為這個人作了王.
接著.
何況那些受洪恩.
又蒙所似之義的.
得到恩典.

$^{201}$以及得到神的義.
豈不是更要因耶穌基督.
一個人在生命當中作王.
就不是死作王.
就不是罪來操控人.
因為第六章會講罪的問題.
這裡告訴我們.
現在有一個很宏大.
浩大的恩典.
就是耶穌來這個世界.
他在生命中當王.
耶穌的義行.
他對神的信服.
使得每一個在他裡面.
在他的救恩裡面.
在他的生命中的基督徒.
就可以在義中當王了.
保羅很想告訴我們.
這才是基督徒正常的生活.
如果我們信了耶穌.
在耶穌裡面.
繼續在罪惡裡面.
才是不正常基督徒生活.
這就是一個概念.
我們要開始學習.
要有了.
所以這個對比.
很強烈.
十八節繼續有對比.
他說如此說來.
就是說.
既然如此.
到底當帶來的人類有什麼審判.
耶穌基督稱為新的亞當.
或者叫末後的亞當.
新的人類的代表.
又帶來什麼平反或者更新呢?.
十八節就說.
因一次的過犯.
眾人都被定罪.

$^{241}$剛才是說一個人.
現在是說一次.
分別善惡旋的果.
亞當違背實.
聽從魔鬼那條蛇的引誘.
吃了那個果子.
其實不是吃.
使他中毒.
是他違背神的命令.
不聽神的吩咐.
而做出這個背逆的行為.
所以.
第十八節就說.
一次的過犯.
一次.
眾人都被定罪.
就是所有的人類.
往後衍生的人類.
所以很快你會留意到.
在《創世記》第三章.
第四章之後.
就馬上說.
亞當的兩個兒子.
戈隱和阿伯.
兩個人都是有罪的.
不過.
戈隱殺自己的弟弟.
那個罪就更加嚴重.
這裡告訴我們.
下面那裡這麼說.
他說照樣.
因一次的疑行.
就是耶穌的這個公義的.
為人類的罪釘十字架.
對不起.
有一點點乾.
我先喝口水.
不好意思.
大家要留意.
耶穌所做的是一個疑行.

$^{281}$在保羅的眼中.
主耶穌沒有罪的.
祂是很明顯.
必須是無罪.
祂才可以擔待我們人類的罪.
所以祂是一次的疑行.
又是一次.
使眾人也被稱義得生命.
被神稱義.
得著這個新的生命.
在我們裡面.
當然就是聖靈進入我們的內心.
給我們有新的力量.
給我們有新的生命.
所以每一個洗禮信主的人.
都得著這個新生命.
接著第19節.
又再用另一個角度.
講述這個對比.
因為一個人的背影.
眾人成為罪人.
照樣不斷這樣說.
就是說一個類比.
亞當所做的.
帶來人的審判.
但是耶穌基督現在所做的.
又帶來人的一個很大的更新.
他說一個人的背影.
眾人成為罪人.
第19節的下半節.
照樣一個人的信徒.
一個信服神.
定義上十字架.
本來有神聖的耶穌基督.
但是他取了一個奴僕的形象.
有人的樣式.
任由人對他的嘲笑,侮辱,攻擊.
來欺凌耶穌.
都願意走上十字架.
來被人釘死.

$^{321}$背負著種種羞辱,羞恥.
這樣來受盡痛苦.
這樣來死.
所以這是一個很徹底的對神的信徒.
所以這個信徒.
就帶來什麼呢?.
眾人成為義.
在神的眼中.
我們得著神的義.
不是我們自己的義.
不是我們自己換取回來的義.
是神的義.
透過耶穌.
透過耶穌這個義行.
加給每一個信的人.
所以我們在神的眼中.
已經不再被定罪.
我們在神的眼中.
嚴格來說.
再沒有這個定罪的控訴.
只有站在神的恩典裡面.
所以我們如果能夠經常默想主.
上上到天上神的面前.
成為一個待討者.
讚美者.
我們的靈命就會不斷地有成長.
不斷地有更新.
我們會與神有一個很近的生命.
以致我們的思想.
我們的價值觀.
就不再太停留在生活的種種糾纏裡面.
例如很多煩惱的地上的事等等.
我們可以慢慢去梳理它.
不會變成糾纏不清.
就會變成很多的煩惱.
這裡告訴我們.
因為一個人的順從.
眾人都成為異.
所以這裡告訴我們.
下面那一段我們不看那麼多.

$^{361}$這裡告訴我們第21節.
就如罪作王叫人死.
照讓恩典也藉著義作王.
叫人因我們的主耶穌基督.
或者說在主耶穌基督裡面.
就是永生.
在主耶穌基督裡面.
不再面對肉身的死亡.
以及將來的永恆的沉淪的死亡.
這裡告訴我們.
其實這裡說照讓.
就是說恩典現在是藉著耶穌的這個義行.
來作王.
恩典現在可以普遍地.
臨到每一個信的人.
在主裡面.
或者說倚靠主耶穌的人.
就可以得著這個恩典.
而在恩典當中.
我們來到去得永生.
在這個恩典裡面.
我們是藉著義作王.
所以不是恩典作王.
恩典是藉著義作王.
所以其實跟第六章是有息息相關的.
因為如果我們只是覺得.
我們是一個得到.
獲得恩典.
而我們的生命.
我們的生活根本沒有改變.
根本跟神沒有任何關係.
信了耶穌也是很若即若離.
回不回教會也完全沒有問題.
是否事奉也完全不介意.
這樣你就是恩典濫用者.
這裡告訴我們.
保祿很強調.
正常基督徒就是跟神的關係復和.
恩典我們是享用.
但是藉著義作王.

$^{401}$所以這裡告訴我們.
叫人因我們的主耶穌基督而得永生.
在主的救恩裡面得永生.
所以第五章.
保祿就將我們的焦點.
帶回人類第一次墮落失敗.
然後帶我們看到.
我們既然在主裡面.
當時羅馬教會在保祿很想他們有這種眼光.
因為他們都面對帝國當時的尼祿王.
對教會的迫害.
也有猶太人在羅馬的教會.
大家是共野一勞的.
猶太人有時候如果他們還有一個.
未全面的神學觀.
就會有一種排斥.
對外邦羅馬教會的信徒.
是會有一個排斥.
到底如何理解這些意大利人.
這些外邦人.
我們是不是還要守他們這些猶太人的規矩呢?.
譬如到安息日.
要不要聚會崇拜呢?.
這些問題其實是一些很簡單的倫理問題.
但是如果解決了神學的問題.
其他就很次要了.
所以到了十四章才再探討.
到底哪一天比較好呢?.
到底哪些東西可以吃.
哪些東西是不潔的.
哪些東西是聖潔的呢?.
這些就變成了倫理的問題了.
但是如果我們處理得不好.
當然倫理的問題也可以帶來教會的分裂.
但是現在保祿告訴我們.
回到人類最起初.
就是說任何的後代.
任何的民族.
都要回到那個起源.
就是那個源頭.

$^{441}$就是阿當的背影.
阿當的離開神.
阿當的犯罪.
現在神就不再在救恩以外.
來追究我們了.
如果我們在神的預備的救恩裡面.
在耶穌的裡面.
我們就得到神給我們一個新的更新.
生命的改變.
所以不斷在說作王.
不過是恩典藉著義作王.
就是說真正的神的恩典.
救恩的恩典.
這個恩典會帶來我們越來越明白.
以及接近神的義.
當然這個接近不是說我們還沒清義.
還沒得救.
是我們越來越在生活裡面.
有神給我們的義.
在我們的生活能夠反映我們的義行.
這個就是第六章.
有兩個問題.
保羅都可能會想到.
弟兄姊妹羅馬教會.
他們有可能會被濫用.
也可能有一些批評保羅說的恩順正義.
說得這麼簡單這麼隨便.
就是說會衍生了.
你說的恩典是無限量的.
人都可以這麼信就行了.
是不是會變成一種濫用.
就是在恩典裡面我們犯罪也不需要懼怕.
不需要害怕.
上個星期我們已經處理了第一個問題.
第二個問題就在第十五節.
是另一個問題.
所以你看到很工整.
第六章有兩個假設性的問題.
雖然是假設性.
保羅已經很清楚.

$^{481}$他提出.
他設問了一個問題出來.
這個問題.
他可能預計有一些批評舊恩的人.
恩順正義的人.
會想出來.
來挑戰教會.
挑戰保羅.
又或者可能真的有一些人.
他們誤解了自己得到的舊恩是這麼棒.
所以繼續在舊有的生活形態裡面.
例如什麼呢?.
當時在羅馬教會.
你想想你在帝國的教會裡面.
尼祿要拜偶像.
自己是神.
你們要向君王媒拜.
如果你相信了耶穌.
那些是一般的習俗.
不要緊.
就照拜.
祂是人而已.
我們就照拜祂.
如果不拜就麻煩.
這些問題.
可能我們已經誤解了.
以為在恩典裡面什麼都可以再做.
所以第一個問題.
因為今天新同學.
再簡單地很快重溫一個段落.
但是要聽比較完整的.
可能看之前放在網上的YouTube.
就會比較完整一點.
不過抱歉就是每一個星期.
我們必須要用比較長的時間.
就是一個半小時.
來逐段很仔細地梳理上下文.
以及理解重點的經文.
所以每一堂都一個半小時.
大家就可以在家慢慢聽.

$^{521}$第六章.
一個問題而已.
保羅想到.
可能他是深思熟慮.
又或者是保羅很擔心人誤解恩典的定義.
自己去演繹恩典.
對他這個人的生命.
生活的要求.
不是聖經保羅所期望.
或者神所期望.
所以他立即在第六章.
已經立刻提出一個假設的問題.
我們可不可以仍然在罪惡當中.
叫恩典顯多麼?.
答案根本就是直接了當.
斷乎不可.
中文的語氣和原文都很貼切.
是一個很強烈的否定句.
加重了語氣.
你連想都不要想.
這樣的意思.
我們會不會在我們的生活中.
我們根本慣性了.
慣性犯我們自己喜歡的罪.
慣性的慣犯呢?.
保羅說.
是會的.
真的有的.
但是.
如何去跟著這個問題呢?.
接著他用了十一節經文.
來解釋.
我們必須要站在一個新的身份地位裡面.
就是透過洗禮.
這個儀式.
我們已經很清楚.
向神.
向其他人.
去宣告.
當然更加重要是向自己的心靈.

$^{561}$第六章這裡主要是向自己的心靈.
去宣告你是一個新人.
你是一個新做的人.
所以你的生活.
是應該過一個新做的生活.
的意思.
所以都不是向人.
不是向神.
神已經知道我們信主了.
有性靈.
都不是向魔鬼.
這裡完全沒有說魔鬼.
其實是向自己宣告.
是要讓自己明白了.
你是新做的人.
我們可以看回這一段.
很快就將我們上個星期說過的.
來做一些總結.
我做了上一個星期的PowerPoint.
跟大家先看一看.
這裡有一點.
大家看回PowerPoint.
在網上應該都看到這個PowerPoint.
剛才我們提過.
五章一到第十一節.
就是同神復和的恩順清兒.
是帶來我們因為與神的相和.
而有新的生命.
十二節說回新的生命的來源.
就是因為我們已經成為新的亞當的族類.
所以我們說的是一個集體的.
不是個別的基督徒.
第六章.
個別不同的基督徒.
如何去看自己是新做的人.
這樣的一個集體.
其實與當時的羅馬帝國.
所說的羅馬的平安.
羅馬的承平.
羅馬帝國給你的福祿.

$^{601}$其實可能不是錯的.
他搞好了道路網.
全國復歸統一.
帝國帶來的就是.
你去到任何地方做生意.
都不再懼怕那些盜賊.
因為羅馬帝國很強大.
奧古斯都已經剷平了城邦等等.
不會再有其他反對的勢力.
那些反對的勢力已經變成奴隸了.
成為帝國之下被人用的工具.
但是這些都是物質上.
給我們表面的平安.
但是保羅現在推銷的福音.
是真真正正人類的新的開始.
就是所有在耶穌救恩裡面的基督徒.
就成為新的族類了.
所以就享受到.
耶穌由死復活這個新生命.
帶給我們的新的力量.
所以第六章就開始講這個新的生活.
應該嚴格來說.
第六章我的分題就是.
不可以再犯罪.
因為兩段.
第一節和第十五節.
都很清楚講了不可以再犯罪.
但是其實第六章全章.
我們如果將這兩個問題.
再重新去看.
其實第六章應該不是一個負面的不可犯罪.
應該是正常的基督徒.
可以不犯罪.
這個就是第六章.
你想想這是多大的震撼.
人類歷史歷代都是在罪惡之下.
最帶給我們的就是與神隔離.
最帶給我們的就是死.
藉著罪作王.
每個人都要面對死亡.

$^{641}$還有罪就會帶來人與人之間的疏離.
這些是人類解決不了的.
完全沒辦法.
可以透過制度去解決的.
但是第六章.
你應該得到一個結論就是.
可以不犯罪.
當然那個大前提就是.
我們要過正常基督徒生活.
什麼是正常呢?.
第六章就告訴我們.
首先是很負面的.
不可以犯罪.
但是它說的理由其實很簡單.
上個星期我們第十一節還沒正式說.
現在我們稍微將第六章再重複一點.
上個星期的PowerPoint還有一點點在這裡.
你可以看到.
在耶穌裡面的新生活.
這個正常基督徒生活.
原來是我們在救恩裡面.
其實我們因為浸泥.
這是大家共有的一個這樣的禮儀.
一個這樣的開始.
就歸入基督耶穌的死.
我們以往很少人會說.
我是歸入了亞當的死亡的.
我們很少人會這樣理解的.
但是我們歸入了耶穌的死.
所以原來耶穌的死.
是帶我們的罪進入審判裡面.
耶穌擔待了.
但是我們的主耶穌.
是第三天由死亡裡面復活.
而且這個復活的大能.
是神使他復活.
所以這是一章第四節.
羅馬書.
因從死裡復活.
以大能顯明他是神的兒子.

$^{681}$耶穌一章四節羅馬書.
所以現在保羅訴諸正常基督徒生活.
就是你立即信了主之後.
你要看自己是和主耶穌同死.
其實是他死我們沒有死.
是他死了.
是他進入陰間.
是他帶我們的罪.
去到神的審判.
他已經擔待了.
神因為耶穌就赦免了我們所有人的罪.
然後.
耶穌已經復活了.
是我們人類的代表.
以及現在在天堂上.
為我們禱告.
為我們預備將來的居所.
當然我們不要用物質去想.
到底要多大呢?.
Winnie分到一千呎.
我分到五百呎.
我們不要用這個物質去想.
而是將來有許多.
因為在永恆裡面是一個浩瀚的靈界.
何來需要有尺寸呢?.
所以我們不要用我們人間的物質.
香港人人住tong 房.
我們不要用這些去想.
而是做物主.
現在就是為我們預備將來的永恆.
在那裡不會再有任何痛苦了.
所以今天正常基督徒生活.
原來就是對基督所做的.
我們要完全融合下去.
完全看自己.
就是這樣來與主同釘十字架.
這是《加拉太書》第五章.
我已經與基督同釘十字架.
現在活著的不再是我了.
當然那是二章二十節.

$^{721}$《加拉太書》.
現在活著的不再是我了.
是基督在我裡面活著了.
並且我如今在肉身活著.
是因為信神的兒子而活了.
他時愛我為我捨己了.
完全是二谷同工.
完全是在說對真理的凝訊.
我們可以想象到.
原來真正的基督徒生活.
是要讀神學的.
是要對這些神學.
要掌握深信不疑.
以及在生活裡面實踐.
你每天起來告訴自己.
我是新做的人.
通常哪些人才會這樣呢?.
做完大手術.
甚至被人死了.
心臟停了.
然後離開了物理的身體.
去到一個很光的地方.
然後突然間回來.
又說得出醫生搶救你.
曾經這樣的人.
真是新做的人.
他真的完全發現自己曾經在地上做的事情.
不再重要了.
真的要為神作見證才是重要的.
真是大有其人.
但是我們不需要去到那些極端的時刻.
反而現在保羅是全面性的說.
不可犯罪.
是成為基督徒的正常生活.
但是要怎麼去不可呢?.
所以這裡告訴我們.
第五節很重要.
與基督一同的復活.
是被基督釋放.
以及將來.

$^{761}$真真正正將來我們會與基督一同活在天上.
所以.
第十一節就是具體.
正常基督徒生活的.
上個星期我也說過.
那個獻上的問題.
原來保羅拆解了.
我們要看第十一節.
所以上面的我不再詳細說了.
因為上個星期我們已經很詳細解釋了.
你可以上網看那個PowerPoint就有了.
還有每個星期我都會將今晚的PowerPoint.
放在群組的.
所以問群組的人就可以拿到的.
或者你手上給我一個電話.
我就可以加你進去群組.
那你就可以收到群組裡面一些分享.
或者一些查經資料等等.
非定期就會有的了.
但是PowerPoint現在每個星期我都會放上去的了.
所以大家就不需要現在寫筆記.
反而大家一起進入那個討論裡面.
這次我們就看第十一節了.
十一節.
六章的十一節.
這個又是一個很重要的實踐.
原來正常基督徒生活是不斷獻上.
還有如果用香港人都明白的就是.
你到底跟誰大腦.
你跟了罪做大腦呢?.
這就是我們過去的問題人生了.
現在阿保羅不斷強調.
我們要跟隨主耶穌.
獻上給主.
將我們的生命.
這就是第十一節開始了.
他說這樣.
第十節的總結.
他死就是耶穌.
向罪死了.

$^{801}$一次.
所以第十節是一次而已.
所以.
拆開來說就是天主教的彌撒.
每個星期都是耶穌再釘十字架.
就不應該這樣.
因為所有新約告訴我們.
耶穌是一次的彌行.
一次的十字架的犧牲.
就已經解除了人類的罪了.
所以我們每一次.
是相信耶穌基督在十架的當年.
當日為我們釘十字架的事件.
我們這個十架的事件.
發生了.
不需要再發生了.
我們就可以立即跟神復和的關係了.
所以每一次如果我們犯罪.
我們再一次回到耶穌的救恩裡面.
求祂赦免.
我們就立即又立即清義.
立即得到神的赦罪.
這就告訴我們.
他死是向罪死了.
只有一次.
就是溫誠伯拉.
就是一次性而已.
不會有第二次了.
不需要有第二次了.
有第二次全部都是異端.
接著他活.
就是當他現在說的復活.
是向神活著.
留意開始.
伯拉就要我們應用了.
正常基督徒生活就是要應用了.
他說這樣.
就是同樣.
你們就是我們跟隨耶穌的人.
你們向罪也當看自己是死的.

$^{841}$向神在基督耶穌裡.
卻當看自己是活的.
在耶穌基督裡.
這是一個很重要的保羅書信.
尤其是羅馬書的一個很重要的概念.
我們不斷要看自己是在基督耶穌裡.
記住是生生不息的.
由我們信著的那一刻.
我們就已經在基督耶穌裡.
基本上你是不可能離開基督的.
除非你完全否定你的信仰.
否定到一個地步.
你要拒絕這件事.
其實是很難的.
因為你的良心或者你自己也不敢這麼做.
甚至有些人去到失憶.
老退化.
也很少會這樣.
很多人老到一個地步.
他什麼都不記得.
他都記得自己是信主.
剛剛跟一個98歲.
在台灣回來.
住在錦田的老人院.
又是透過有人介紹.
就要幫他洗了.
因為很大年紀了.
他已經聽不到東西了.
要很大聲在他的耳邊.
他的兒子就替我幫他.
好像翻譯一樣.
其實是大聲說話.
我說不用了.
我都夠大聲的跟他說.
說他信耶穌.
他真的完全很懵.
那個懵不是貶義.
是他真的什麼都不記得.
但是說到信耶穌.
他很開心.

$^{881}$說他自己小時候已經是信耶穌了.
在台灣已經是信了很多年了.
日於年表的說他自己很開心.
說他信耶穌.
就說了這麼多東西.
但其實他什麼都不記得了.
又年紀大了.
你就可以想像到.
他在主的救恩裡面.
相信一定有很多經歷.
但是也不方便.
因為他年紀太大了.
說不出話來.
就幫他做了一個洗禮.
所以這裡告訴我們.
很重要的.
我們要這樣看自己.
原來真正的正常基督徒生活.
其實回到第六章的問題.
不可以犯罪.
但其實嚴格來說.
保羅是在說正面的.
真正正常基督徒就是.
可以不再犯罪了.
我們要看到自己原來.
我們已經跟主耶穌完全聯合.
所以在基督耶穌裡.
在羅馬書的概念是很重要的.
我們在主基督耶穌裡面.
就是永生.
我們現在在主基督耶穌裡面.
我們就可以.
我們可以用義去作王.
就是說我們可以過一個.
清義的生活.
這個就是保羅告訴我們的.
秘訣就是這裡.
所以你看到第十一節的下半節.
要看自己是活的.
這裡就開始我們第十一節.

$^{921}$就是今晚我們要說的.
我們看看這段.
新的翻譯.
你們也要看自己.
向罪是死.
在基督耶穌裡面.
向實向善活著.
原文的意思英文也有.
NASB就是Consider.
King James Version就是Reckon.
Reckon大家就知道.
是數學的計算.
就是說你要算一個賬.
你花了十塊.
你就沒理由花九塊.
你就要在你的計算機裡面.
花十塊.
那就叫做Reckon.
算了.
這個就是算術方面的詞彙.
在另外一個.
我蓋了一個希臘文.
第十一節.
這個我高光了.
就是要這樣理解.
我就叫AI幫我做一個直接翻譯.
透過這個希臘文直接翻譯.
他說這樣.
你們也要認為自己.
認為不是漢字.
漢字我們以為.
就這樣看一看.
其實你要這樣認定自己.
對罪是死的.
卻在基督耶穌裡面.
對神是活的.
記住這個是生活的常態.
成為基督徒的生活的常態.
你要這樣認為.
你要這樣Consider你自己.

$^{961}$你要這樣去Reckon自己.
相信我想我們那個類比一定沒那麼完整.
但是其實也不是很難的.
現在我去了一些基層的社區洪水橋.
開始現在有些中產.
或者新搬進來的那些階層.
但是那些基層的小朋友.
其實你觀察他.
你也看到他有時候真的沒有那種眼界.
他真的沒有見識過.
一直都住在公共屋村.
沒有貶義的.
不過在他的生活圈子裡面.
他未必感受到一些那種態度.
待人接物.
或者沒有被人教育.
跟一些有教養.
出入過一些可能比較乾淨的環境.
有Caste的地方的人.
那個Behaviour是完全不同的.
我們作為牧者我們會觀察到的.
完全沒有貶義的.
不是說躋身在漂亮的地方就一定好.
那些就不好.
完全不是的.
而是我們說的就是.
他的Behaviour就是因為在每天的生活裡面.
他已經重覆重覆.
這樣來跟自己說.
例如他叫阿甲.
阿甲你是很有錢的.
你出入可以去會所的.
你出入有車載你的.
阿甲你記住你是不同其他人的.
其實他不用這樣問自己的.
但是他每天的生活.
他出入的地方.
就已經讓他不知不覺.
他就認為自己是這樣.
所以這個認為.

$^{1001}$這個Recognition.
其實就是正常基督徒生活的那個思維.
而且這個思維不是我們自己在想當然.
而是不斷要我們對應.
耶穌已經是向罪死了.
耶穌現在已經是向神活了.
已經是活著在天上了.
我們就是要將這個本來是耶穌的.
現在因為我跟耶穌聯合.
我們就要怎樣成為正常的基督徒不犯罪呢?.
是不准犯罪是負面.
其實是可以不犯罪的基督徒.
原來就是你要認為自己.
當然了.
躋身在豪宅.
你是含著金鎖匙出生.
你已經很明白你小時候長大.
你已經跟別人與別不同了.
你已經不斷別人告訴你.
少主人你不用幫我做事的.
我請回來是服侍你的.
綁鞋帶都不用你綁的.
我幫你綁的.
可以有這樣的人的.
當然很少數.
但是他已經知道自己與別不同了.
但是現在正常基督徒就是.
任何每一個人.
不同階層.
不管你是外邦人,猶太人.
因為現在基督就是神的拯救代表.
已經為我們向罪死了.
我們所以現在能夠看自己.
是向罪都死了.
我們在基督耶穌裡面向神.
我們都可以活著了.
要這樣行動.
就需要在生活裡面.
常常帶著基督這一份信心.
與主一起面對我們生活的人生.

$^{1041}$所以正常基督徒.
就是在生活的每一個層面.
來體驗這一節聖經.
保羅對我們的宣告.
而且是一個事實的宣告.
不是我們自己一廂情願.
這就是十一節告訴我們.
我們要向罪看自己死了.
在神的面前.
在基督耶穌裡面.
卻當看自己是活的.
就是說現在我們是一個活的生命.
接著第十二節.
思考就是一個基督徒的人生.
原來是一個成長的人生.
是一個有方向.
是一個向神活著的.
我們要計算.
這個計算就在這裡出現了.
就是清二是神已經告訴我們.
他已經犧牲他的兒子.
審判他的兒子在十字架.
所以耶穌的犧牲.
就帶來我們得到清二.
但是現在十一節是告訴我們.
基督徒就是要在這個義裡面.
怎樣計算自己.
當然這個計算就是第十二節.
就是開始講到獻上.
留意第十二節.
所以不要容罪.
這個是以人法.
罪成為了一個惡魔.
一個曾經對我們很大的威脅.
壓迫和管束我們的.
第十二節.
罪容罪在你們必死的身上作王.
使你們順從身子的私慾.
意味著第十二節告訴我們.
罪曾經在我們未信耶穌之前.

$^{1081}$就在我們的生命裡面.
身體裡面作王.
以人法罪是一個背役.
對神的背役.
罪不是一個物體.
一個有位格的.
將它以人化了.
罪你不可以再讓過去的惡的王.
再去操控我們.
使我們這個必死的身體.
即是容易犯罪的身體.
到了一個階段當然會死亡.
但是必死的身體.
即是容易犯罪的身體.
繼續這樣被他操控.
操控我們就會再傾向.
順從那個私慾.
就會犯罪.
所以原來是很微妙的.
好像轉眼之間.
現在就是保羅要強調.
所以從第十三節連貫來看.
我們要跟另外一個大佬.
我們要跟從耶穌.
十三節.
也不要將你們的肢體獻給罪.
作不義的器具.
剛剛的對比.
說: 禱要將從死裡復活的人.
將自己獻給神.
正常基督徒的生活.
就是選擇哪個大佬.
跟隨哪個大佬.
以及怎樣獻上自己.
原來每天都是一個獻上.
對主耶穌基督的一個效忠的獻上.
而且是自己的身體.
這裡說: 要將自己的肢體.
作義的器具.
不再是不義的器具.

$^{1121}$現在是義的器具.
獻給神.
這樣罪就不能再作我們的主.
因為我們不在律法之下.
一直都在說在恩典之下.
我們來獻上.
所以這個獻上.
是一個甘心樂意的奉獻.
這是一個神喜悅.
這是保羅勸告每一個基督徒.
所以十二章的總結就是.
所以弟兄們.
我以神的慈悲勸你們.
將身體獻上.
當作活祭.
是聖潔的.
原來我們成為基督徒.
在神的眼中.
我們如果真心的獻上.
是聖潔的獻上.
是神所閱立的獻上.
所以第十二節.
十三節就已經說.
在這裡已經很快說.
所以不是到十二章才說奉獻.
第六章已經說了.
是我們每個人要有一個意識.
我不再跟隨罪.
做他的奴隸.
我要將我的身體.
獻給罪.
做犯罪的不義的器具.
現在我是要有意識地.
Consider.
Reckon.
我要認為自己.
是向罪死了.
向神.
在基督耶穌裡面活.
我要這樣獻上自己給神.

$^{1161}$原來我們的身體.
可以被神大用.
以前是不義的器具.
現在在神眼中.
每一個有聖靈.
有主救恩的基督徒.
每一個的身體.
你的意念.
你的心思.
全部是神體為聖結的.
我們可以獻上給神的.
不獻就笨了.
你的生命不獻給神.
原來這裡第十二節.
十三節是一個對比.
我們不獻給神.
罪又會再來成為我們的王.
克制我們.
我們又會將我們的身體.
不知不覺又會成為不義的器具.
所以這裡是說.
不義和公義.
獻給罪還是獻給神.
來做奉獻.
說的是生命的奉獻.
所以留意第十三節.
其實是一個負面的.
如果我們不獻給神.
對不起.
十四節.
罪必不能作你們的主.
就是說.
如果我們不獻給神.
不主動的認定.
認為自己已經死.
以及在神面前活.
我們就會繼續被舊有的罪性.
包括罪這個王.
繼續這樣拖著我們走.
所以很多基督徒.

$^{1201}$過著一個不快樂的生活.
因為可能不斷在罪裡面糾纏.
或者跟人之間的矛盾.
又解除不了.
所以我們不知不覺.
就將自己的肢體.
我們的肢體.
來獻給罪.
成為不義的器具.
例如你的腦想的全部都是很負面的.
想想誰恨你.
你恨他.
你就會整天在你的腦袋裡面.
充斥著很多的憤怒.
埋怨.
嫉妒.
紛爭.
甚至惡毒.
陰毒.
很多在我們的腦海裡面.
這些是腦海裡面沒有人看到的.
是神看到的.
不過慢慢你整個面容都不開心.
人人都看到.
因為你已經很少去祈禱.
很多的去想別人不好.
你就將你的器具獻給不義了.
但是現在保羅就告訴我們.
原來正常基督徒.
是可以獻給神.
作公義的器具.
公義的器具就是達到可以行義.
可以活出義的生活.
所以以前你不行的現在可以.
可以來到去常常感恩.
常常從別人的好處出發.
想到別人的好.
來到感恩.
聽到別人的好.
你就可以成為一個更加豁達的人.

$^{1241}$原因就是因為有主的救恩在我們裡面.
所以第十二節和十三節.
真是一個很重要的秘訣.
但是保羅甚至要我們必須盡快實踐出來.
如何不用罪在我們必死的身上作王呢?.
必須要將自己的肢體獻給實.
這就是十二,十三兩節的對比.
不要將你的身體獻給罪.
再一次獻給罪.
作不義的器具.
我們的身體是一個器具.
可以被神用.
所以上了主喜歡唱聖詩.
參加師班.
雖然水準都不是那麼好.
但是重點是我的心很想獻.
盡力去訓練.
獻最好的給神.
這是其中一樣東西.
我們的聲音可以成為獻給神的一個頌讚.
你的內心可以多去想到別人.
為別人祈禱.
我教過了.
大家預備自己一本心寫的簿.
寫下哪些人你為他祈禱.
即是說你很著緊.
你久不久就再揭開神.
他還未成功.
你真的幫他吧.
其實提醒你要祈禱.
提醒自己要祈禱.
不要下排輕輕.
這樣日子有功.
真的會有很大的變化.
其實真正的變化是你的心靈的變化.
是你開始看到祈禱的重要.
看到祈禱的重大功用.
所以原來我們的身體.
我們的不同的器官.
不同的器官.

$^{1281}$可以作為義的工具獻給神.
所以正常基督徒不准犯罪是負面的.
正常的基督徒生活.
就是不斷呈獻自己.
就是可以不犯罪.
這是第一個段落.
保羅解釋的第一章.
第一節的六章第一節.
我們可以仍然在罪中.
是不可以的.
即是說我們現在可以不在罪中.
是正面的意思.
以至恩典就更加的凸顯.
其實是恩典托住我們這個人的生命.
我們可以在神的救恩裡面.
我們可以不犯罪.
所以記住不是躲在深山野嶺裡面.
自己獨自修行.
修行也沒有用.
可能你很苦悶.
你去打水.
你很不開心.
天天都想著這個世界等等.
所以是在生活裡面過我們正常基督徒生活.
是在生活的場景過正常基督徒的生活.
所以你就可以想像一下.
你的時間如何運用.
你的金錢如何運用.
運用在永恆裡面.
運用在你能夠成為義的工具裡面.
如何為神去用.
所以每個人都有自己那把尺.
不過就要根據聖經.
那個大邏輯原則就是在這裡.
不要再獻給罪.
這個大佬是跟不過的.
我們曾經跟過的.
跟過就不斷犯罪.
但是現在我們跟的這個新的大佬.
就是最好的主耶穌.

$^{1321}$就是神的兒子.
所以我們在耶穌裡面.
我們就有平安.
我們可以將我們的肢體獻給神.
所以很快說這個獻.
一定大家都很熟悉的就是第十二章.
十二章是一個新的段落的開始.
因為第十一章是整個的頌讚.
保羅對舊恩的敬拜.
然後第十二章就是說.
所以弟兄們.
既然我們領受了這麼大的恩典.
一切都是由神那裡而發出.
一切都是因為他.
因為他無限的愛.
來憐憫我們這些罪人.
所以一切都是應該要本於他.
依靠他.
歸於他.
一切的榮耀都是要歸給他.
直到永遠.
阿們.
突然間那個轉變.
就是第十二章.
保羅開始要探討.
在羅馬的教會.
基督徒的生活倫理.
實際應該怎樣運作.
人與人之間應該怎樣運作呢?.
所以首先就是.
在第一時間說就是.
每一個基督徒正常的生活.
就是要奉獻自己給神.
所以弟兄們.
我以神的慈悲勸你們.
不是神的命令.
不是神的憤怒.
不是神的要求.
是神的慈悲.
來勸你們.

$^{1361}$勸就是說.
你可以不做的.
你不做就是你自己的損失.
不是神有任何的損失.
獻上的人.
主動獻上的人.
就會得到這些.
一般的基督徒有的更大的平安福祿.
將身體獻上.
就是說整個人.
身心靈都獻上給神.
不是分割了.
這半天是神的.
當然的確是.
星期天的上午.
你還要想其他東西.
你真的不是敬拜神了.
走到教會.
坐在禮拜堂.
你還想其他東西.
你還在想.
一坐下就想.
待會吃什麼.
一個半小時你還要想.
你真的不是呈獻了.
但是保羅說的就是.
整個生命是呈獻給神的.
因為我們本來就是神將我們拯救的.
將我們的身體.
真的是身體.
將我們的身體獻上.
所以不是一種.
好像抽象的.
心裡面自己想就行了.
我心想了就行了.
其實不是的.
在你的整個生活裡面.
是有具體的行動表現.
他說將身體獻上.
當作活祭.

$^{1401}$接著保羅還要勉勵信徒.
羅馬的教會信徒.
你這樣獻上.
你這個具體的行動.
由心靈開始.
到生活.
到實際的場景的實踐.
是聖潔的.
是神所喜悅的.
你不用擔心.
你這樣獻上.
神是不要你.
不稀罕你.
神是不喜歡你.
不會的.
曾經有素羅的獻祭.
神就不喜悅.
大衛的獻祭.
神就喜悅.
大家可以想像得到.
這裡告訴我們.
每個基督徒.
神都喜悅.
因為你這樣獻上.
在基督裡面的獻上.
已經是聖潔.
是神認為聖潔.
是神閱立的一個獻上.
所以你們如此事奉.
這個事奉.
英文是Stand.
站在這裡.
你這樣準備去事奉.
就已經是神所喜悅的.
而且保羅說.
其實是理所當然的.
我們不是奉旨的.
我們能夠服侍是理所當然.
是神要求我們每一個基督徒都要這麼做.
所以我們回到.

$^{1441}$當然我們先不解釋五章.
將來我們先去到.
很久都還沒到.
還有五章.
我們盡量去解釋.
所以第一個段落.
第一節已經說了.
不可犯罪.
不可以因為恩典.
讓你繼續肆無忌憚地犯罪.
不可以.
但是去到第十二十三節.
其實告訴我們.
其實我們每天呈獻自己的人.
其實就可以不犯罪了.
所以很正面.
是很積極的.
所以這個就是呈獻的問題.
這裡真的要具體實踐.
我們基督徒沒有任何的秘訣,捷徑.
就是每天都奉獻自己.
將自己放在神的旨意裡面.
神的旨意有時候.
你要可能跟某個人聊天.
或者你要跟某個人幫他面對一些困難等等.
這個就是神怎樣為我們人生去安排.
但是我們的責任就是.
每天我們要看自己.
正常的基督徒就是一個奉獻給神.
將身體奉獻給神的基督徒.
所以我們的心思意念.
我們的口.
我們所聽的.
都不可以再是一些神不喜悅的東西了.
第二個問題.
十五節.
十三節我還有多加一個花樣.
你想想.
神一開始的震怒是對哪些人呢?.
就是那些神不喜悅.

$^{1481}$而他自己喜悅的生活.
所以神就任憑這些人.
隨著他心裡面的私慾.
行一些污穢的事.
這個就是一張在神的憤怒下的那些人.
但是現在六章十三節.
基督徒我們離開了舊有生命的行列之後.
你看看這個對比.
其實是很首尾呼應.
沒有神的人.
他以為自己做什麼.
自己負責.
自己心裡面想做什麼就做什麼.
的確是的.
很多人現在不信耶穌.
這個就是他們認為我喜歡做什麼就做什麼.
不需要有任何的規範.
原來真正基督徒的規範.
只不過是神的生命的律裡面.
道德的律裡面.
神期望的真正的自由的生活.
就是我們所謂的自由.
不是去放縱.
而是在神的裡面.
將我們的肢體.
將我們的BODY.
我們的器官.
來呈獻給神.
是主動的呈獻.
所以這個主動的呈獻.
要成為我們正常基督徒的思想.
我們的認為.
就是你要行為.
像一個很高級的人生.
你是神的子女.
在這個地球上.
我們真正的永恆歸宿是天堂.
神才是我們屬靈.
這個創造主.
我們的父親.

$^{1521}$所以我們現在要在地上.
代表著神來活的時候.
我們的生活.
就應該反映天上的那種榮耀.
我們有這種.
讓人感覺到.
所以牧者走出來.
如果別人都感覺不到你有平安.
你有喜樂.
這也是一個問題.
因為你正正就是.
透過聖經你得到的真理.
應該就是首先你自己要活出.
當然你可以為一些事情很煩惱.
要禱告.
有些很嚴重的事.
需要大家認真甚至禁食祈禱.
這些是對的.
但是平常.
沒有什麼特別的事情的時候.
你是不是真的有一份平安.
有一份發自內心的平安.
對永恆的盼望.
和有信心.
有信念.
和舉止投足.
你都能夠讓人看到.
這個人是有神的.
有些先知.
因為今天早上講福音傳契.
就講以利沙.
他是神人以利沙.
人們看到.
那些人是不相信的.
我看到這個人是有聖神的靈的神人.
就是讓人感覺到有一種氣質.
這個人就是跟隨神的人.
所以你看到這個對比是很強烈的.
沒有神的.
那些以為自己喜歡怎樣做就怎樣做的生活.

$^{1561}$但是有神的正常的基督徒.
就是保要像從死人當中復活過來的人一樣.
將自己獻給神.
獻給神是一個很大的一個運圖.
但是實際每天的生活.
就把你們自己的身體的肢體.
獻給神.
作義的器具.
記得我們已經是一個義的器具.
我們的身體,我們的手腳.
已經是應該要來到分別為勝出來.
作義的器具了.
所以記住仍然是生兒養女.
有生活有不同的需要的興趣.
需要有你的嗜好.
有你的生活的不同的品質.
但是記住都是獻給神的義的工具.
這個就是六章十三節.
和一章保羅所說的那些敵黨神的人.
大家的心態完全是兩樣東西.
我們以前不信耶穌.
不知道有耶穌救恩.
我們都是這樣過自己的生活.
但是現在我們有神在我們心裡面.
我們現在就是要不斷呈獻.
所以六章的一到十四節.
可以不犯罪的正常基督徒.
就是不斷奉獻自己給神的基督徒.
這是第一個段落.
我看看我有沒有過時.
已經達八了.
我們看第二段落十五節.
另一個問題.
我也是做了小小的總結.
雖然有些操作.
不好意思.
到底我們跟誰呢.
我們保羅說千萬不要再跟錯人.
跟錯主人.
當時他用奴隸制度.

$^{1601}$奴隸就是羅馬帝國一個很普遍的社會制度.
因為他們抓了很多俘虜回來.
難道就全部殺了他嗎?.
這些全部是勞動力.
所以很多的俘虜都變成為奴隸.
有奴隸制度.
所以你不要想到美國的黑奴.
虐待他死.
不要這樣想.
因為全部是很重要的勞動力.
有些奴隸是做律師.
做管家.
或者有其他的工作.
當然最低等的就是做很粗重的事情.
都有的.
有不同等級的奴隸.
但是仍然是奴隸主.
所以那個制度是保羅的時代.
所以他不斷有時候都會用奴隸的制度告訴我們.
因為聽的人一聽就明白.
如果奴隸被一個主人釋放了.
他就不再屬於那個主人.
就算那個主人是很殘暴的.
他都不需要再懼怕他.
因為那個制度賦予他可以有自由.
這個就是當時的奴隸制度.
所以第六章告訴我們.
不要再跟隨這個主人.
因為他是不好的.
他是很殘暴的.
他是會操縱我們的.
我們現在要跟隨的就是基督了.
所以罪和基督變成一個擬人化的對比.
其實罪是沒有違格的.
罪不是一個神聖的主.
第二個問題就是.
既然我們不可以再犯罪.
根本的問題就是生活方式要改變.
要從看自己.
這個看自己就是認為自己.

$^{1641}$計算自己.
Recognize.
Consider.
這樣來看自己的時候.
我要過一個成為神的子女的基督徒生活.
所以是生活方式的問題.
所以如果再問我可不可以吸煙.
我可不可以來到馬會.
那時候我進馬場.
頭柱站很簡單.
穿著白色校服.
讀新法的時候.
後來讀模範.
牛津道出名的班五的時候.
我們撕掉校章就能進頭柱站了.
在頭柱站擔著一瓶煙.
拿一份報紙.
那時候沒有平板的.
沒有電話的.
那個時代.
拿一份報紙.
馬報.
跟其他的叔叔在那裡聽跑馬.
買毒癮,煉癮.
有時候試過跟同學進馬場.
我唯一沒試過進香港的麻將館.
因為不敢進.
因為是三教九流.
不敢進.
但是打牌,麻將那些一定有.
你可以想像得到.
是生活方式的問題.
舊有的就是原來我們不知道.
罪是我們的主人.
但是基督徒.
我們要很快.
保羅維利一定要不斷認為自己.
是基督將我們拯救了的.
我們現在是.
基督是我們的主.

$^{1681}$我們是他的奴隸.
我們要跟隨基督成為新做的人.
所以其實是生活方式的轉變.
這個就是正常基督徒.
可以不犯罪的根本的.
其實不是秘訣.
是每個基督徒都要這樣的.
所以不是特選的基督徒才要這樣做.
不是因為你要奉獻做傳道才要這樣做.
是每個基督徒都要這樣理解.
自己現在已經歸屬了耶穌.
這樣的意思.
所以洗禮原來是一個很重要.
當時那個你可以說一個.
又不是過關.
又不是突變.
而是普遍的一個禮儀.
代表著,象徵著.
你已經進入了死亡.
你已經由死復活.
這樣的一個新的開始.
所以正常基督徒.
新的開始就是.
跟隨耶穌.
他是好的.
以及生活的方式.
要一個新的轉變.
第二個問題可能我點出之後.
就要完了.
因為都答了九了.
第十五節.
第二個問題.
那時候再一次說斷乎不可.
就是問什麼問題呢?.
就是這卻怎麼樣呢?.
因為上文說.
我們已經不在律法之下.
即不是規範的.
我們在恩典底下去做這一切.
所以整個人生.

$^{1721}$由信主到我們離開世界.
都是在恩典底下.
所以我們要記得.
既然現在是恩典.
即是說就算我們做得不對.
做得不夠.
可以繼續重新開始.
這個就是在恩典底下.
所以我們就更加要把握我們的時間.
我們要盡快認識.
保羅告訴我們.
原來要這樣看自己的.
你要考慮自己是.
同主耶穌同死.
同主耶穌同復活.
這樣的正常.
就會在我們的生活中反映.
不過.
保羅又想到另外一個.
可能會有一些知識份子攻擊他的.
攻擊他.
說恩典這麼普遍.
這麼廉價.
就會帶來一些人濫用.
放縱.
第15節的問題就是這樣.
第6章的第二個問題.
他說這卻怎麼樣呢?.
我們在恩典之下.
不在律法底下.
是不是我們就可以犯罪呢?.
好像和第一個問題差不多.
第一個問題就是.
我們可不可以仍然在罪惡當中.
仍然在罪惡當中呢?.
就是說.
我們堅持舊有生活的模式和形態.
和那個的思維.
就是說.
我完全習慣舊有的.

$^{1761}$我不改的了.
基督對我沒有任何的影響.
我不是一個新做的人.
這是第一個問題暗示的原因.
但是第二個問題就不是這樣了.
沒這麼嚴重了.
就是說.
偶然間犯罪.
因為恩典.
使得我們可以偶然間犯罪.
就是偶然間.
回到舊有的生活.
和第一個就有點不同.
這裡告訴我們.
仍然是斷乎不可.
就是說.
保羅也可能會設想到.
有羅馬的頂子妹.
或者在這個救恩底下.
到底他犯的是猶太人.
羅馬人.
還是其他特定的人.
或者甚至是挑戰恩典的普及的知識份子.
都不重要了.
重要的是保羅又帶出了一個很重要的方法.
同樣這次其實和上文差不多的.
不過很強調.
我們要獻上自己給主耶穌.
成為主耶穌基督的一個奴隸.
這裡整個段落都講到奴隸制度的負面和益處.
可能很快地將第六章先讀一讀.
我們下個星期再認真看這一段.
我也做了一點初步的分析.
就是在恩典之下.
所以可以犯罪.
這是新的翻譯.
AI幫我做了一個直接的翻譯.
用希臘文直接的翻譯.
在恩典之下.
所以我們可以犯罪了.

$^{1801}$保羅不行.
斷途不可.
因為我們不再是罪的奴隸.
我們現在是神的奴隸.
或者是義的奴隸.
所以這裡涉及到罪和義.
我們看自己是奴隸.
但是到底是罪的奴隸.
還是神的奴隸呢?.
記得不要看在負面.
因為現代人看奴隸.
就是很負面,虐待,壓迫,黑工.
我們不要這樣看.
記得當時的奴隸是很普遍的.
差不多有三分之一的人口都是奴隸.
在羅馬帝國.
所以其實是很高的階級.
所以千萬不要用我們今天的現代觀念.
一說到奴隸就覺得很負面.
其實不是.
所以保羅這樣敢說.
如果有基督徒.
你到底是想繼續成為罪的奴隸.
還是成為基督的奴隸,神的奴隸呢?.
這就是原來可以幫助我們.
可以不犯罪的一個重要的出路.
其實我們是否相信,明白也好.
其實我們一出生就是罪的奴隸.
但是成為了基督徒.
我們就必須盡快思考自己.
原來已經加入了成為愛我們的主,神的奴隸.
記住不是負面的奴隸.
一個異的奴隸.
我們可以在我們信主之後的人生.
過一個異的奴隸的生活.
甚至我們信從這個異的人生.
我們得到的結局就是永生.
信從罪的結局就是死.
一個很強烈的對比.
所以你看到就是.

$^{1841}$本來有三個課題.
罪和異就在這一段裡面.
不信服不是明顯說的.
但是信服就很清楚了.
死和永生也是很清楚的說的.
是兩個的結局.
所以第六章這一段.
我很快先讀一讀.
第十六節.
為什麼「鑽福不可」呢?.
為什麼我們在恩典底下不可以再犯罪呢?.
第十六節.
「豈不曉得你們獻上自己作奴僕」.
這就是上文說過的.
信從誰就作誰的奴僕.
或作罪的奴僕以至於死.
或作順命的奴僕以至成異.
留意第十七節很特別的.
這裡說到一個道理的模範.
那個forms of teaching.
那個規範.
那個規模.
那個規模是甚麼來的呢?.
其實就是整全的福音.
我下一個part 1就有了.
一個完整的福音.
就帶來我們這個規模.
就會帶給我們一個順從神的心.
這個順從神.
我們就變相成為了異的奴僕.
不得到釋放了.
我再讀下去.
第十六節.
「豈不曉得你們獻上自己作奴僕」.
信從誰就作誰的奴僕.
或作罪的奴僕以至於死.
或作順命的奴僕以至於成異.
因為你們從前.
每一個人都是的.
未信耶穌.

$^{1881}$雖然作罪的奴僕.
現今信耶穌之後.
卻從心裡順服了.
所傳給你們道理的模範.
為甚麼要這麼詳細做一個形容呢?.
從心裡順服.
傳給你們的這個救恩.
這個道理的模範.
一個教訓.
一個形式的教訓.
即是說一個規模的.
即是說一個很整全完整的福音.
主耶穌的來這個世界救恩.
主耶穌的死.
主耶穌的復活.
和我們的關係.
我們如何和祂聯合.
原來是一個形式的教訓.
人類面對的現在的心靈的問題.
終極的問題.
不只是生老病死.
是有永生.
是有永死.
是有神的審判.
是有神的賞賜.
所以這個是形式的教訓.
還有保羅特別強調.
如今如果我們想過正常基督徒的生活.
是要從心裡順服.
福音原來是要順而順服.
順完你不服.
只有口講.
你就經歷不到基督在你生命裡的釋放.
是從心裡.
不是口講.
是從心裡順服.
所傳給你們道理的無法.
所以是一個完整的形式的教訓.
是一個完整的架構.
所以基督徒一間健全的教會.

$^{1921}$如果沒有兒童主義學.
如果沒有初信班.
沒有浸禮班.
沒有成長班.
那個教會就很大有限.
就算你的性學多厲害.
就算你的人講道理多厲害.
什麼都沒有了.
沒有那個配套了.
沒有人會去主義學的.
是不會開主義學的.
那個教會就基本上沒有用了.
因為保羅這一句已經很清楚了.
賢今從前就是罪的奴隸.
賢今是為什麼你們會成為一個新做的人.
為什麼你們現在能夠這樣奉獻給神.
是因為從心裡順服了.
所傳給你們道理的無法.
所以保羅為什麼這麼緊張提摩太後書的二章.
第二章.
第二章的第一節.
提摩太後書.
聽一聽.
我兒啊.
你要在基督耶穌的恩典上剛強起來.
你在許多見證人面前.
聽見我所教訓的.
親耳聽到.
在現場.
你要將這個教訓.
你自己要來去.
接受.
明白.
融會貫通.
接著你要交托.
又要交托那忠心能教導別人的人.
到底有多少人呢.
交托給一些人是忠心的.
那些忠心的人又可以再教導其他人的.
那些人.

$^{1961}$即是說.
好幾個的外圈.
即是說.
使徒保羅教會第二代的年輕的跟隨者.
提摩太太 提多 路加 馬可等等.
這些人繼續傳道.
他就會繼續去招生授徒.
訓練一些自己認為忠心的.
看他有忠心的 又肯學習的.
使他能夠成長.
那些人就成為第三代 第四代的傳道人.
繼續去傳.
繼續將這個道傳開去.
來到今天 二千年了.
所以基督的真理.
二千年前已經告訴我們.
要不斷開班授徒.
沒有神蹟.
不需要一定要有神蹟.
只需要將一個很重要的完整的真理.
有系統的 用心的 用禱告的.
去教導下一代.
自然那個道就可以一直傳遞開去.
神的道就可以發揚光大.
所以這裡告訴我們.
所以十七節真是今晚的金句.
我再讀一次十七節.
因為你們從前.
雖然作罪的奴僕.
然今卻從心裡信服了.
所傳給你們道理的無法.
就是說現在你看自己是神的奴僕.
因為你已經聽明白了這個真道.
而且是從心裡信服.
幾個動作的.
心裡接收.
所以很多基督徒.
雖然好像很空泛 好像很誇張.
但是不帶一個心來聽道.
來崇拜.

$^{2001}$不帶一個人來.
不帶一個心來.
其實是一件很榮耀的事.
你禮拜日可以放下你任何的工作.
和你整個人生的身份.
你是大機構的CEO.
是很偉大.
很好 很有成就.
但是你一走進教堂.
你是錦繡堂的一個弟兄.
你是元朗區宣道會的其中一個弟兄.
你走進去.
你去敬拜神 朝見神.
有時很多人已經忘記了這個動作.
他走進去教會.
帶著的還是自己的思想.
自己的成見.
自己的看法.
自己的心思意念.
他已經沒有帶一個心去敬拜神了.
何來會在敬拜裡面得到神跟你說話呢?.
所以是從心裡面馴服.
所傳給你們.
你想想這幾樣東西.
心馴服.
聽道.
肯謙卑聽道.
肯聽.
不過不謙卑.
裡面只是你這樣說.
我不會聽你說的.
不過很溫柔地.
我不會聽你說的.
有些人就索性不聽你說.
來都不來.
但是我們就是要從心裡面馴服.
所以別人講道.
我們要聽.
所以你不要以為.
神學生講道.

$^{2041}$你就可以不用聽.
其實不是的.
那個人是代表著神去講道.
我們無論是多麼資深的牧者.
堂主任.
其實那個人代表著神講道.
他那一刻不是神學生.
他那一刻就是神的僕人.
你要聽他講道.
麻煩看著.
幫我傳開去.
做私事的.
不單是做私事.
做完就在後面聊天.
我是不介意說.
你是失去了從心裡面馴服.
所傳給你的道.
你不要等高永謙來你才聽.
李思敬來你就坐下來聽.
然後第二個人來你就在聊天.
其實不要.
我們每一次是來朝見神.
聽神.
透過神的僕人跟我們講道.
誰站在台上就是神的僕人.
我們就要聽他.
神透過他和你和我講道.
這樣你才會真正明白.
17節這一句聖經的重要.
你們羅馬教會.
是從心裡面馴服了.
傳給你們道理的部分.
所以保羅用了這麼詳細的.
去展述新的生命.
其實正常基督徒的生活.
就是要這樣理解自己.
第18節.
你們既從罪裡面得了釋放.
過去式.
因為在耶穌裡面.

$^{2081}$就作了義的奴僕.
我因你們肉體的軟弱.
因為我們還是不斷.
你想想那個奴隸.
你可以想像一下當時的奴隸.
有人將他釋放了.
買斷了.
但是舊的主人是很兇的.
還在的.
看到他那個僕人都害怕.
那個主人會不會再抓我回去呢?.
但其實不需要了.
因為他已經屬於新的主人了.
在羅馬的制度下.
他已經是屬於新的主人.
因為已經買斷了.
但是在心靈上.
他仍然存在.
他不是死了.
所以這裡本來就說.
我因為你們肉體的軟弱.
當然包括自己.
就照人的常話對你們說.
就是說.
我是用一個好像說人話的意思.
就是說我現在不是說一些很超脫.
很虛無的東西.
是說一些很實在的.
你們從前怎樣將肢體.
獻給不潔不化作奴僕.
以至於不化.
就是做了這些不好的事情.
自然就會帶來神的憤怒.
得罪神.
但是現今再一次一個對比現今.
現今也要照樣將肢體獻給不潔不化作奴僕.
所以其實就是回應上文的13,14節.
也是真真正正在第六章.
為什麼可以連續說兩次.
斷乎不可不可以犯罪.

$^{2121}$但原來不是斷乎不可.
是非常可以我們看到.
只要我們認真獻上自己.
我們就可以過正常的基督徒生活.
就是不斷呈獻自己給神.
過一個義的生活.
以至於誠信.
所以真正的誠信.
是我們人生的閱歷.
我們對神的經歷和認識.
我們就可以不斷在這個呈獻下.
我們就可以以至於誠信.
第20節.
因為你們作罪之奴僕的時候就不避而約束.
你們現今所看為羞恥的事.
當日就是過去的時候.
有什麼的果子呢?.
再一次說那個果子.
大家記得勝靈所記的果子.
加拉太書五章應該是19節.
人愛喜樂和平.
忍耐恩慈良善.
信實溫柔節制.
九個果子.
聖靈所結的果子.
這裡是罪帶來的果子.
就是羞恥.
結局就是死.
不斷說那個果子.
是有具體的行動的.
是有具體的開花結果的.
你犯罪.
自然你這個人.
其實是你越想很多不開心的事.
你真的有之於內.
亦有之於外的.
其實是很慘的.
還有上一代的牧者說.
聽一個人彈琴.
那些詩琴的壓力很大.

$^{2161}$聽到他的靈性.
他可能這段時間都沒有靈修.
那個詩琴聽到的.
不知道怎麼聽.
千萬不要這樣去量度.
禮拜的詩琴.
可能人家真的很累.
是不方便告訴你的.
但是其實要為他祈禱.
如果你聽到他彈琴.
是很緊張的.
或者沒有了那種心靈的釋放.
其實可能內心有些東西.
其實我們就需要祈禱.
或者今天這個講者很累.
他講話好像有點重膽.
其實你默默的要為他祈禱.
因為可能他有很多事情要面對.
所以成熟的基督徒.
就應該這樣來從同理心去看待其他人.
這裡告訴我們.
但現今你們既從罪裡面得了釋放.
作了神的奴僕.
接著就立刻說了.
就有誠的果子.
你想想是兩個對比.
我們有什麼果子呢?.
我們人心結出什麼果子呢?.
這些果子是真的傳到永恆的.
我們的錢財,我們的財物.
只是放我們今生用的.
帶不到將來永恆的.
但是誠的果子是屬於你自己.
永遠跟隨你的.
所以這裡是一個很重要的.
這是一個生與死的結局.
如果關乎生與死.
在第十章的羅馬書.
阿保羅有提到申命記三十章.
摩西那個生與死的呼喚.

$^{2201}$這裡也差不多有了.
不過十章就更加詳細.
九點零六分.
很快的揭一揭吧.
三十章的申命記.
很快的揭一揭.
因為下個星期也會說的.
但是去到第十章就更加一定要說.
因為保羅直接有這段經文.
三十章.
我看看後面.
因為我的手很乾.
十五節沒錯了.
第十章羅馬書.
如果今晚大家可以看.
十五節這裡.
生與福,死與和.
陳明在你面前.
吩咐你愛耶和華你的神.
遵行他的道.
謹守他的誡命,律例,典章.
使你可以存活.
人數增多.
耶和華你的神.
就必在你所要進去等等的事情上.
來賜福給你.
這是生死和福.
第十一節.
更加是保羅在第十章所引述的.
回到第十一節.
三十章的申命記.
我今日所吩咐你的誡命.
不是你難行的.
也不是離你遠的.
不是在天上.
使你說誰替我們上天取下來.
保羅就立刻加一句.
就是基督從天上來.
使我們聽見可以遵行.
第十三節.

$^{2241}$也不是在海外.
使你說誰替我們過海取了來.
使我們聽見可以遵行.
即是那些人在挑戰摩西人.
摩西就在陳明.
其實不是你們在托詞.
說你讓我看看.
我就會信了.
你告訴我哪裡有經書有神的道.
就算多遙遠的海外.
即是地中海以外的地方.
因為當時根本去不到.
因為那些輪船是木船.
很危險的.
我們都會去取回來的.
摩西說不是的.
就是第十四節.
這話或這道卻離你甚近.
就在你口中在你心裡.
使你可以遵行.
就是我們所傳的信主之道.
這就是羅馬書第十章.
保羅直接引述.
《申命記》的三十章.
十一到十四節.
第六章已經有這樣的含義.
生死和福.
我們到底是要呈獻自己給罪.
繼續做罪的奴隸.
因為我們跟了另一個主人.
跟了另一個大佬.
就是耶穌基督成為義的器具.
今晚我們就停在這裡.
我們一起禱告.
感恩天父給我們能夠.
有機會將第六章.
仔細地去分析.
去學習.
多謝主.
讓我們看到.

$^{2281}$這不是甚麼神秘的主義.
這是每一個基督徒都需要學習的.
就是要跟隨新的主人.
過新的生活.
不再是一個故意.
或者是習慣犯罪的生活.
而是我們能夠開始.
過一個能夠結義的.
果子的生活.
求主幫助.
讓我們能夠有一個很敏銳的心靈.
去跟隨主的聖靈.
這樣在生活裡面.
去等候成為義的一個奴僕.
我們的器皿.
我們的身體.
我們的每一個器官.
都能夠被主提醒和使用.
再一次我們多謝主.
給我們今晚的聚會.
奉耶穌基督的名.
\newpage



\section{羅馬書}
\label{sec:PTM8SSdjvVE}
\textbf{2025-07-03 周四晚查經班《羅馬書(16):分題:恩㤟與律法》 講員:曾裔貴牧師}
\newline
\newline
連結: \href{https://youtube.com/watch?v=PTM8SSdjvVE}{\texttt{https://youtube.com/watch?v=PTM8SSdjvVE}} ~~~~ 語音日期: 2025-07-04
\newline
\newline
\hyperref[sec:wh3ZqIs7b68]{\small{< < < PREV SERMON < < <}}
~
\hyperref[sec:index_chronic]{\small{[返順時目]}}
~
\hyperref[sec:index_scriptual]{\small{[返順卷目]}}
~
\hyperref[sec:HoQGEbdeqSk]{\small{> > > NEXT SERMON > > >}}
\newline
\newline
$^{1}$我們一起禱告.
多謝天父讓我們能夠在晚上.
繼續來到你的面前去研讀《羅馬書》.
盼望我們能夠從第五到第七章.
掌握我們是怎樣去呈獻自己給神.
成為神這位愛我們的神的奴僕.
我們來到去獻給義 作義的奴僕.
以致於能夠我們成聖 願神帶領我們.
第七章開始會涵蓋關乎到律法.
和我們基督徒信了主之後的關係.
願你都帶領我們明白.
我們仰望感恩禱告 奉耶穌基督的名 阿們.
麻煩大家跟進著 因為我一直在錄音.
我不知道實際是不是錄到今晚.
大家說是聽到就可以 但我這裡是沒有聲音的.
有人發言 應該是我來的.
我們繼續下去 第五章頭的段落就是.
信了主在救恩裡面 就是與神恢復生命的關係.
有神的性情.
第八章會進入聖靈工作的部分.
享受與神同在帶來的生命的平安 福樂.
和在神的恩典裡面 生命的地位的改變.
下半段的第五章 就不再是舊有阿當的族類.
我們是新的族類 就是新的阿當.
末後的阿當 第二個的人 耶穌的裡面.
我們能夠得著這個永生.
就是說 我們本來就是舊有創造的墮落的人類.
但現在在救恩裡面 耶穌帶我們看到.
我們是新的人類 是一個新的創造.
是一個新的國度 不是當時羅馬帝國的.
那個太平盛世 地上的盛世的國度.
而是在神建造的永恆的國度.
這個國度是包含猶太人 外邦人 所有的民族.
這個就是帶我們好像超越了舊有的創造的失敗墮落.
以致我們在耶穌裡面 得著這個新的生命.
那五章結尾就是在說罪在哪裡顯多呢?.
恩典就更加顯多.
叫我們在耶穌裡面 得著這個新生命的永生.
第六章有兩個問題 我們上個禮拜都有機會探討過.
我們可以犯罪 因為恩典是會繼續顯多.

$^{41}$那保羅說不行的 因為我們已經向罪死了.
藉著洗禮這個儀式 我們是已經向罪死 向神活的.
所以我們要這樣理解 我們這個新的生命.
在耶穌裡面 我們就有這樣的開始.
要這樣去看 自己是已經在罪的面前.
我們以往沒有耶穌之前 未在救恩裡面.
我們是不能夠向罪死.
但是現在在神耶穌的救恩裡面.
我們就可以向罪死 而且是應該要這樣理解自己.
向罪死 甚至六章的十一節就是.
應該要這樣去認同 應該要這樣去考慮自己.
要在罪的面前 我們是一個死了的人.
因為是耶穌將我們買熟 將我們帶進死亡.
又使我們能夠有這個新生命復活.
所以我們向罪就應該這樣理解自己是死了.
所以這個就不是不准 不可以禁止犯罪.
而是我們要在新的生命裡面.
一到第十四節 應該一個結論就是.
我們可以開始去嘗試可以不犯罪的新生命.
接著第十五節又是另一個問題.
我們是不是可以要在那個舊有的裡面.
我們是不是要繼續讓那個罪繼續去顯多.
保羅說不行的 因為我們是已經獻了給神.
我們藉著耶穌 我們不再是罪的奴隸了.
我們應該是獻自己給義 獻給神.
這樣來過一個好像愛的奴隸的生活.
當然這個是保羅的一個比喻.
就著當時的社會的現象 社會的文化.
社會的體制來講一些羅馬的教會能夠明白的道理.
以往我們是獻給罪 以致帶來的是不法.
帶來的結局就是死.
但是現在我們是要將自己呈獻給神.
好像一個神的奴隸一樣 一個奴僕一樣來獻給義.
我們上一個星期都有機會很快的看了下半段的問題.
所以你可以看到整個第六章都環繞著耶穌的死.
我們要進入認同耶穌的生命 新生命.
我們要來去擁有 也就好像已經告訴我們.
間接將第七章 很重要的七章的二十四節.
我真是苦啊 誰能夠救我脫離那個取死的身體呢?.
當然就是主耶穌的死,埋葬,復活和永活.

$^{81}$祂這個新生命當然是藉著聖靈.
已經進入我們的內心就可以得著了.
所以其實第六章都是很想告訴我們.
我們要有那個意志方面 才是從真理有這個認識.
原來我們洗禮就已經是好像一個這樣具體的行動.
來認同耶穌的死和我的生命有關.
祂已經帶著我的罪進入審判了.
我又在耶穌裡面得著這個新生命了.
所以第六章其實兩個問題.
都是很想讓基督徒有一個新的動力.
就是原來信了耶穌可以不再犯罪.
可以在新的生命裡面體驗這種奇妙的變化.
那個的改變和歸屬一位愛我們的真神.
那種的釋放,那種的真的自由.
本來奴隸就不應該有自由.
但是一個愛我們的主人.
這樣來愛我們的主人就是這位神.
這樣的主人就是一份愛的呈現.
就是一份我們刻意自己去選擇.
要走的人生的路.
這樣獻自己給義,獻自己給神.
好像一個神的奴僕那樣.
六章二十二節就是那個對比.
罪的結局就是死.
唯有神的恩賜在主基督耶穌裡面就是永生.
這個就是我們不再作罪的奴隸.
是要作神的奴僕的那個對比,那個生命的改變.
所以你可以想像得到.
由第五章一直帶進我們去看.
其實白白恩典到底是不是可以駕馭地上的道德規律呢?.
白白的恩典如果不是會令人更加放縱,更加犯罪.
那麼白白的恩典和神的律法有什麼不同,有什麼比較呢?.
保羅就是一個很心思細密的人.
因為肯定猶太人會問一個問題.
如果你們說的白白恩典可以令人不需要再面對罪的審判.
甚至有新的生命動力.
這個肯定是道德一直提升.
一直合乎神的心意的道德的提升.
那麼律法還有什麼用呢?.
律法到底在救恩裡面有什麼用呢?.

$^{121}$所以保羅似乎必須要回答.
可能是猶太人的質問.
可能是一種不理解的,一種善意的提問.
又或者是更具侵略性的.
就是對保羅的福音的一種冷言冷語.
因為同一時間猶太教又在灌輸守律法.
安息日要行國禮.
要做不同的大小的儀節,儀文的事情.
同一時間在社會,在世界,在世俗.
又說要拜地上的神才能得到各種的福樂.
又是要做很多不同的儀式.
或者甚至一些比較極端的行為.
宗教的行為才能夠換取到心靈的平安.
現在保羅是一個這麼新奇的.
講的好消息,福音.
那到底因順稱義,白白恩典.
和律法有什麼的違背或者是衝突.
還是說現在根本不需要再有律法了.
這些必須是保羅在一個這麼深思細密的邏輯裡面.
他需要去讓無論猶太人,無論反對福音的人.
都要明明白白地看到.
其實最終都是告訴我們在耶穌裡面就是真正的出路.
在主耶穌的裡面不斷地來和主走人生的路.
就會有一個更加超越律法的道德律在我們的內心.
這應該就是第八章所說的生命,聖靈的律.
在基督耶穌裡面釋放了我們.
這是八章的第三節.
所以其實一直是引進.
不是不斷地講第二件事.
而是在講福音,耶穌的新生命進入了我們裡面.
就有一種這麼大的一個改變.
而且是整個心靈的價值觀改變.
而且是一個很清楚是歸向神,歸依神.
跟隨神的一個人生.
這樣就可以將第五到第七章.
就是一個很完整的關乎到處理罪.
關乎到處理我們那個歸屬的問題.
我們原來信了主是屬於神,屬於耶穌的.
還有就是我們要告訴未信的人知道.
我們不是逼自己繼續做一些宗教的行為.

$^{161}$不是猶太教變成基督教.
就要守基督教的戒律.
不是這樣 而是要讓我們看到.
同樣都是在耶穌的新生命裡面.
自然就來到有一個這樣的.
應該說是加勒太書第五章,第六章那裡講的.
全律法都包在愛人如己這句話裡面.
稍後我們有機會可以揭一下加勒太書那裡.
所以簡單幫大家做這個第五到第七章.
很簡單的分段和邏輯的轉變.
很重要的問題就是第七章我們今天要去處理.
這個就是上一次我們講了一點.
可以再加多一點的.
就是保羅明顯在第十五到二十三節.
第二個問題在恩典之下.
是不是所以可以犯罪呢?.
保羅說同樣都是不行的.
不應該是這樣 因為我們不再是罪的奴隸.
就是說不信的人,不認識主耶穌的人.
其實沒辦法 原來一出生就是罪的奴隸.
就是不斷在罪的奴役底下.
第七章更加恐怖.
保羅說罪原來藉著律法去操控我們.
罪不是直接來干預我們的生命.
罪是藉著律法在我們裡面來禁閉我們.
所以原來那個出路.
都不是我們只是不斷抗拒不犯罪.
原來更重要的出路就是我們要認定.
我們有新的主人歸屬於耶穌.
這個就是生命力在我們裡面.
所以才能夠可以不斷享用與神同在的平安.
還有你感受到在我們的生命是有永生這件事.
就不是一種好像死氣沉沉的.
一種罪帶來的結局的死亡.
所以你可以看到第六章.
就有這樣的一個很強烈的對比.
主要三方面的罪和義的對比.
不馴服的和馴服的那個強烈的對比.
那個結果當然不馴服帶來的就是死.
死就是可能在今生就是好幾件事與神隔絕.

$^{201}$最終我們慢慢模糊了我們的人生方向.
沒有了對永生的渴慕和服侍.
這些都是死的狀態.
但是永生就不同了.
我們知道自己是順從義的人生.
我們是神的奴僕.
奴隸都可以的.
我們是義的奴僕.
這樣在人生不是被人強迫去就範.
而是我們選擇跟隨這一位愛我們的主.
這樣來做他的奴僕.
所以這個是心甘樂意的一個擺上.
是一個甘心樂意神喜悅的擺上.
所以十二章最後那個都是一個dare for的總結.
就是在這麼偉大的救恩底下.
應該每個基督徒.
保羅說應該每一個都是甘心樂意的擺上自己.
好像一個活生生的祭來呈獻給神.
活祭,凡祭是要很講究.
整個洗乾淨.
裡裡外外五臟六腑洗乾淨.
然後擺好原狀.
整個推進去爐裡燒.
那個銅祭壇.
就是說大家是很清楚.
基督徒就是最怕沒有意識.
哎呀 隨便啦.
無所謂啦.
都是坐著坐著.
不要這麼高要求.
那就糟糕了.
那就類似了.
我們基督徒現在這一代.
會不會已經沒有了那種很強烈.
要為神而活的意識呢?.
這個相信是第五到第七章.
一個很重要的秘訣.
就是主動將自己發現.
原來做神的奴隸才是最好的.
原來做義的奴隸才是最好的.

$^{241}$這樣的基督徒人生.
就真是可以不知不覺.
我們就不再去想犯罪.
已經慢慢在他那個力量的引誘底下.
已經不再受操控了.
這樣就是第六章和第七章的精髓所在.
我們可以一路收窄了.
所以保羅是用很多的篇幅.
不斷說我們要有意識地呈獻給神.
offer我們自己.
向神呈獻.
是有意識的.
所以我們的奉獻.
我們的聚會.
牽涉的是內心的態度.
和那個主動性.
這個主動是源於.
救恩對我們生命的改變.
這樣的意識.
這樣有意識地去裝備自己.
這樣就是基督的精兵大軍.
就不是一盤散沙的.
就算你一萬人又如何呢?.
個個都是一個半小時的基督徒.
接著那六天半.
就是自己的人生.
獻給罪 獻給不義.
那十萬人都沒有用.
不能夠將救恩那種力量.
在人間帶出那種影響力.
所以這裡告訴我們.
第六章多重講的就是那個奉獻自己.
就零零碎碎幫大家再做一點點.
第五到第六章的那個分段.
和那個大的邏輯.
保羅的思想.
他要討論的.
果然是如果我們基督徒.
是一幫仍然和羅馬帝國的.
未信的 拜偶像的.

$^{281}$尼祿手下的那些人.
的一樣生活的形態.
那基督教何來影響這個社會呢?.
基督教何來會被人羨慕.
逼迫他都不會否認呢?.
你就想到有一幫人就是.
有一種很獨特的生活.
而這個獨特生活不是一些怪異的宗教行為.
最近有十個家庭.
已經走了上山避地震.
避難.
你可以想像得到.
如果我們錦繡堂.
哇 慘了 七月日本大地震.
不單止不要去日本.
我們現在要起碼一半的錦繡堂的弟兄姊妹.
我們現在快點上大帽山.
曾牧師就買了那些求生包.
每個人一包 免費的 不收錢.
別人說我們邪教.
真的有人響應.
那這些古怪的行為.
你會很吸引到那些傳媒的關注.
因為個個都不發現 而你發現.
那你有什麼根據呢?.
那你有什麼動識天機呢?.
如果不是這樣的怪異.
而仍然都是在正常的生活裡面.
所以第六章.
應該是正常基督徒生活.
就是獻情給神.
向罪死 向神活.
的正常基督徒的生活.
所以如果在羅馬帝國.
保羅期望羅馬教會.
他們如果過正常的生活.
就是一幫有影響力的信徒.
是影響當時羅馬帝國.
不需要再拜偶像.
也不需要靠你自己的行為.

$^{321}$而是靠伯伯的恩典.
而這個伯伯的恩典.
不僅不會令到人放縱犯罪.
因為第六章就試探討這個問題.
反而我們是有道德規範的.
這個道德規範.
就不是逼自己去走律法.
而是原來.
當然這個第八章就開估了.
就是告訴我們.
原來我們有了這個新的生命在我們裡面.
所以你看到一個進路.
保羅是想建立一幫.
有影響力的羅馬的基督徒.
影響力帶來什麼?.
就是要宣教到萬民.
就是萬民都唯有信主.
才可以避開神的憤怒.
就是唯有福音才可以叫人得到.
與神的和好.
所以就算所謂的誠信.
都不是我們個人自命清高那種.
不吃人間煙火.
那就不是講聖經保羅講的誠信.
那個是有些可能偏激的宗教行為.
例如我們有一段時間我自己.
福音堂.
真是我們以不看戲.
以做每一件事都很認真.
作為我們那個持守.
本來的態度不是錯的.
心態都是對的.
但有時慢慢就會矯枉過正.
說要簡樸.
簡樸到一個年輕人剛剛信主.
又用開LV的.
用開一些名牌的.
突然間穿白恤衫藍褶褲.
在國貨公司元朗還有的那時.
穿那些衣服這樣就叫敬虔了.

$^{361}$有一個這樣的模式的.
本來做的人不是錯的.
但跟著做的人可能就誤解了.
以為你有這個形態.
就等於敬虔等於屬靈.
但保羅從來不是講那些外貌.
外在.
裡面打通了與神關係的呈獻自己給神.
這樣的態度.
是要不斷經歷.
主耶穌在我們生命裡的帶領和感動.
所以最終都是要處理這個問題.
恩典與律法有什麼關係呢?.
律法在舊約神頒佈的猶太人.
還是這樣遵守的.
對於我們新信主的基督徒.
在保羅在使徒他們的帶領下.
信主的又要面對羅馬帝國.
又要面對猶太教.
到底我們跟隨.
信這位主耶穌.
我們還需不需要守律法呢?.
律法對我們基督徒.
有沒有規範呢?.
這些又是必須要處理的問題.
所以你看到加拿大書會處理這些問題.
哥林多前書又處理這些問題.
所以其實都是大的議題.
如果不是那我猶太教都可以.
不用再跟一個信一個基督教.
所以你會看到.
初期教會會爭論.
要不要守國禮.
要不要守安息日.
所以我們就要問這個問題.
恩典和誠信是如何過渡呢?.
白白的恩典.
如何能夠防止道德上的濫用呢?.
白白的恩典即是我犯罪都不怕.
白白的恩典.

$^{401}$到底如何去稱平衡呢?.
這些就是第六章.
保羅立即要處理的課題.
我們繼續下去.
對不起.
先回到正題.
我們要思考.
已經有這個恩典.
律法還要不要有呢?.
是不是那個時代已經過了.
我們不需要再守律法呢?.
所以也是羅馬書第七章.
主力要探討的.
那律法如果需要的.
它的功用又是什麼呢?.
有沒有它的極限呢?.
第七章就是處理這個問題.
所以如果我們能夠第五到第七章.
我們都能夠掌握保羅的中心思想和脈絡.
其實我們都好像找到一個.
我們可以開始學習不犯罪.
不是閉上修煉.
不是勤軍大羅儀.
只能夠去到第三級.
之前的教主.
我忘記了他的名字.
都是去到第五級.
張無忌因為他有九陽神功.
就去到第七級.
不是這樣練.
而是在正常基督徒生活裡.
不犯罪才是真正的正常基督徒生活.
這個發現就不是開玩笑.
原來任何基督徒.
都不需要有什麼的弊典.
不需要有任何的刻意的苦修.
正常的生活裡.
和耶穌一起.
就是正常基督徒生活.
就是可以慢慢地不犯罪.

$^{441}$這樣就不是秘密.
是一個重要的生命成長的邏輯.
所以我們到第七章.
當然必須要處理律法和我們基督徒的關係.
可能大家偶然會聽到.
有人這樣說.
立志為善 由得我.
此事行出來 由不得我.
這個七章十七節.
保羅就這樣說.
我們可以慢慢去探討一下.
到底他在未信之前有這樣的生吟.
這樣的痛苦.
還是他想反映在這樣的狀態下.
罪的律和想向善的律交戰下的困難.
是怎樣呢?.
當然有不同的解經家.
會有不同的理解.
我們純粹從一個.
我們在整個宏觀.
第七章處理律法這個課題.
來思想.
我們第七章.
跟大家先來到.
快一點 我們先讀這一章.
因為我今天沒有帶漢語聖經譯本.
我就用舊的和合本.
不是新修訂版 舊的和合本.
有些字眼可能大家去到某些位.
就要再用其他的版本去看一看.
第七章.
弟兄們.
我現在好像是另一個段落的討論.
討論了.
罪不應該成為恆常化.
罪不應該是基督徒正常的生活.
不犯罪才是正常的基督徒生活.
你要記住.
不犯罪才是我們正常的生活.
那就很不同了.

$^{481}$有了這個概念.
我們記住我們去實踐.
就不是騙自己.
而是我們去實踐.
我們認定耶穌在我生命裡有這樣的力量.
就會經驗到不犯罪.
所以現在就有另一個段落的討論.
弟兄們 我現在對明白律法的人說.
你們豈不曉得.
律法管人是在活著的時候.
再一次從最重要的根基.
就是我們已經是死了.
在耶穌裡面死了的人.
這樣來到.
六章已經說了我們在罪上死了.
現在這裡告訴我們.
我們已經是和耶穌同死.
要這樣看自己.
這樣才是我們正常.
我們是因為救恩 因為耶穌.
我們已經是死了.
所以律法就不能夠再追究.
追討一個死了的人.
所以這個就是第一到第六節.
那個中心的思想.
你們豈不曉得律法管人.
就是對人的約束.
對人的要求.
是在活著.
就是那個人仍然在生的時候.
就如女人有了丈夫.
就是她有婚約.
丈夫還活著.
就被律法約束.
就是那個婚姻的盟約.
丈夫如果死了.
就是說那個婚姻的約就解除了.
就脫離了丈夫的律法.
就是說.
我想當時羅馬帝國也是這樣的.

$^{521}$雖然他們也有多妻.
甚至乎初期.
也不是初期教會.
在希臘時代.
據聞這個亞歷山大大帝.
都是同性戀的.
大家好像看過亞歷山大那套電影.
拍成電影.
有一幕是說.
他有一些戀童.
有一些男妓.
跟他同床.
亞歷山大大帝.
所以那個時代.
希羅時代.
絕對不是聖經說的一夫一妻.
但是這裡保羅也告訴我們.
他說.
丈夫死了就脫離了丈夫的律法.
肯定是一個舉例.
第三節.
所以丈夫活著.
他若歸於別人.
就是他的太太便叫淫婦.
丈夫如果死了.
他就脫離了丈夫的律法.
就是婚姻的規範.
在神面前的二人的盟約.
雖然歸於.
對不起.
丈夫死了.
他就脫離了丈夫的律法.
雖然歸於別人.
也不是淫婦.
我的弟兄們.
開始透過這個例子.
一個比喻.
一個簡單的比喻.
帶進了我的弟兄們.
第四節.

$^{561}$我現在要講那個道理.
我的弟兄們.
這樣說來.
你們藉著基督的身體.
因為是基督死了.
在律法上也是死了.
叫你們歸於別人.
因為變成我們是新婦.
如果在那個例子.
在那個比喻裡面.
我們是嫁了給耶穌.
太太.
那個女人.
因為耶穌已經為我們死了.
我們就是歸於別人.
我們就不再在這個.
這樣的律法的規範下.
當然.
我們還是歸給耶穌.
因為耶穌是復活了.
叫你們歸於別人.
就是歸於那從死裡復活的.
當然就是耶穌.
叫我們結果子給神.
這裡真的已經講了重點.
我們認定.
我們是歸屬耶穌.
我們是歸屬神.
原來是真的可以.
結果子.
我們結神喜悅的生命果子.
正如六章.
我們在上個禮拜.
都曾經點出過.
這裡告訴我們.
六章的第二十節.
應該是十七節開始.
因為你們從前.
雖然作罪的奴僕.
然今卻從心裡信服.

$^{601}$所傳給你們.
道理的模範.
你們既從罪裡得了釋放.
就作了義的奴僕.
我因你們肉體的軟弱.
就照人的常話.
對你們說.
你們從前怎樣.
將肢體獻給不潔不法.
作奴僕.
以至於不法.
現今正常基督徒.
生活應該怎樣.
也要將自己.
獻給義作奴僕.
以至於成聖.
後面就講.
那個果子.
22節就是那個總結.
現今.
一個真的認定.
自己和主耶穌完全的.
聯合的基督徒.
正常的生活是怎樣.
現今你們既從.
罪裡得了釋放.
作了神的奴隸.
就有成聖的果子.
結局就是永生.
真的會有成聖的果子.
真的可以.
有一個.
好像很無形.
裡面有一個新生命.
規範我們.
向神而活.
而結出神.
認可的成聖果子.
而且最終的結局.
是永生.

$^{641}$六章都已經這樣講.
所以七章繼續講.
我們歸於別人.
我們就能夠.
結果子給神.
歸於從死復活的.
耶穌.
我們就能夠結果子給神.
不再在律法規範.
反而能夠結到果子給神.
你可以想像.
多麼奇妙.
似乎保羅.
一會兒就會講.
很強烈的對比.
就是.
如果你繼續活在.
罪的引誘下.
舊有的.
性情下.
我們不刻意去過正常.
基督徒生活.
我們就會在一種困擾.
在罪直著律法.
來干擾我們.
但是.
我們只要回到耶穌那裡.
因為第七章的二十五節.
靠著我們的.
主耶穌.
就能夠脫離罪.
脫離律法.
來到去.
我們可以結果子給神.
所以.
不斷地將我們帶回去.
看.
現在每一個信主耶穌的基督徒.
你現在的生活才是正常.
歸向神.

$^{681}$獻給公義.
獻給神.
跟隨耶穌去生活.
這樣才是真正的正常基督徒.
正常基督徒.
就是會結果子給神.
其實是很簡單.
接著我們下去.
第五節.
因為我們屬肉體的時候.
就是說仍然.
未在救恩.
的規範底下.
那因律法而生的惡肉.
這裡有些考究.
原來罪.
是會藉著律法.
來要求我們.
而律法的要求是神的要求.
我們根本就沒有力量.
面對律法的要求.
而最終.
其實是我們舊有的軟弱的肉體.
來到去透過律法.
反而我們很想走.
但是其實.
滋生了惡的肉念.
這裡說.
因律法而生的惡肉.
就在我們肢體中發動.
以致結成死亡的果子.
對比.
死了的.
歸於耶穌的.
就結耶穌神的果子.
我們.
活在罪底下.
罪就藉著律法.
來滋生.
很多的惡念.

$^{721}$惡肉.
這個肉念是惡的.
是神喜悅的.
是邪惡的.
滋生的就是.
死亡的果子.
死亡就是說.
最直接的就是.
與神的疏離.
其實與神疏離.
未必馬上.
與教會疏離.
可能.
現在我們地鋪.
招牌有些白泥.
外面還好地地的.
但是上兩個禮拜.
找人來打開.
裡面已經充滿了白泥.
裡面已經爛了.
你可以想像得到.
這個告訴我們.
我們結出死亡的果子之前.
我們其實.
首先就是與神隔絕了.
與神越來越疏遠.
沒興趣再談論神的事.
但是可能.
每個禮拜還在教堂.
有一個責任.
又很擔心.
人怎麼看你.
一個月都沒來教會.
你去了哪裡.
有時候都不想解釋.
但是其實裡面已經是空洞.
其實你是在結出死亡.
遠離神的果子.
所以那個仍然的解決之道.
是要看到.

$^{761}$立即回到耶穌那裡.
在耶穌裡面.
去看自己正常的基督徒的生命.
你就會有.
那個生活的力量.
去結神喜悅的果子.
兩個果子.
所以是.
很大的懸殊的分別.
所以這裡.
普羅加進了.
一件事.
原來律法.
暗示不關律法事.
就討論關於律法的功用.
這裡就告訴我們.
首先我們必須認識.
我們是已經.
和主耶穌同死的.
律法不可以.
對一個死的人再有任何的.
追討.
罪都不可以再透過律法.
來使我們.
滋生一些惡慾.
第六節.
但我們既然在.
捆我們的律法上死了.
普羅用這麼.
傳神的字眼.
來形容.
捆綁著我們的律法.
如果我們是猶太教.
的中堅分子.
普羅自己就是法利賽人.
他憤外體會.
在律法底下的.
各種的疑問.
對他的每日的.
規範要求.

$^{801}$他做到叫做勉為其難.
做到,但更多的人.
未必真的做到.
那種做不到的枷鎖.
就像普羅這裡的.
很傳神的形容.
是捆綁著.
我們心靈的律法.
但是普羅說.
我們已經在.
無論他捆綁我們多大.
多深.
由何時開始.
但是如果一個.
死了的人.
多大的.
那個律法的要求.
都不再有效.
我們在這個律法上死了.
所以現今.
留意這些字眼.
不斷說現今.
就脫離了律法.
叫我們服事主.
是按照著.
心靈的新樣.
或者聖靈的新樣.
第八章就說.
不再按照著.
移民的舊樣.
按照著一個表面的.
守律法的一個形態.
很好的外表.
的一個.
可能是猶太教的人.
可能是一些有道德規律的人.
他們有他們自己的.
基本的道德的.
主張.
或者是那個生活的要求形態.

$^{841}$但是普羅說.
仍然都是死了.
看自己是在律法.
在罪的面前.
死了.
所以第六章就是.
罪就不可以再作我們的主人.
第七章就是說.
罪不可以再藉著神的律法.
來對我們有規範.
有一個捆綁.
是令到我們.
捆綁著.
我們一個被緊逼的人生.
只會不斷帶來的.
就是好像更多的.
惡慾.
因為捆綁而滋生出來了.
下面就有一個例子.
就是說不可貪心.
第十屆的第十屆.
下面普羅就說.
其實真正的根源.
原來是內在的.
罪性不是律法.
不關律法事.
律法從來都是.
聖潔公義良善.
這個是七章.
第.
十二節.
所以你看普羅這個對比.
層層的拆解.
原來最終是我們這個人.
是要認識.
自己是要死.
在耶穌的救恩裡面.
其實已經死了.
我們是歸於愛我們的主.
我們的神.

$^{881}$和有意識地呈現自己.
成為一個活祭給神.
那麼我們的生活.
就會進入正常基督徒生活.
很少人今天說正常基督徒.
更加很少人.
知道.
什麼才是正常基督徒的生活.
我們很多時候.
就是每個禮拜回教會.
每個月有奉獻.
有手聖餐 有參與事奉.
甚至乎去短宣.
甚至乎讀一些神學課程.
那就是正常基督徒的生活.
那些是基督徒.
參與這些活動的形態.
是希望幫助我們.
更加豐富對神.
對自己事奉的認識.
但是原來真正的.
保羅理解的正常基督徒.
從來都只是.
認識耶穌.
在我們的生命裡面.
祂所為我們擺上的.
就是祂的降生.
為我們犧牲.
為我們向罪死.
向神活.
而且這個生命已經給了我們.
每一個基督徒.
所以正常基督徒.
當然剛才說的全部都是對的.
都是正常的 但是真正的正常.
原來就是.
我們每一天認定我們.
是一個新生命的人.
我們就活在神的面前.
這樣的形態.

$^{921}$的思考.
就是正常基督徒的生活.
第七節開始.
就是律法到底有什麼用呢?.
是不是應該廢棄呢?.
其實所以不關律法事.
我們就來到看到.
原來保羅在這裡.
不是叫我們不理律法.
而是根本如果我們是正常基督徒.
有耶穌的新生命.
有那個復活的能力.
我們就.
不需要再理會律法.
那個對我們的.
無論過去 現在.
那個的影響.
因為我們有更大的影響.
就是耶穌的復活大能.
影響著我們的人生觀.
影響著我們的生活.
所以其實.
我不知道.
其他外國人.
中國人男士就最難說.
對不起的.
越大年紀就越難說對不起.
很多時候像魯鳳.
必然這樣.
但是一個正常基督徒就應該.
越來越能夠做錯.
說對不起.
真心的說對不起.
人家幫了你一點點.
懷緬別人的恩典.
在你身上.
你要不斷有一種感恩的心靈.
感恩的行動.
這樣的一種.
傳遞的愛.

$^{961}$的交換.
正常基督徒就會這樣想.
所以根本不需要理會什麼民族.
什麼大男人主義.
因為就會做的.
我們要讓基督徒的復活能力.
在我們裡面.
就是七章.
其實都是說這個主題.
接著就說.
其實保羅都要告訴我們.
猶太教都會質疑.
那律法.
還有什麼用呢.
保羅就.
都是一個猶太的拉比.
他自己都是全律法的.
法理賽人.
他自己最有資格去說.
他自己舊有生命.
和現在在耶穌裡面.
一個傳道者.
叫人來歸順.
耶穌的福音.
所以他就.
將這個律法原原本本.
本來的.
意思和現在.
在舊恩裡面再去看律法.
相信有了.
不同的眼光.
第七節.
他說這樣.
我先喝一口水.
相當精彩的.
我們.
如果不是透過羅馬書.
其實我們可能從來都不會想.
律法和我們外邦人.
有什麼關係呢.

$^{1001}$和猶太教有關.
我們外邦人根本不需要再去憲制.
不需要再做其他的安息日.
不需要做其他東西.
其實這個是倫理關係.
去到第十四章就要處理.
這些不同的文化的.
倫理關係.
所以教會.
從來都不懼怕多元.
不懼怕不同的民族.
共野一路.
懼怕的就是我們.
不貫穿我們對真理的.
全面的認識.
所以律法現在在這裡.
是有一個定位.
這個點就是因為耶穌已經.
成就舊恩.
律法是應該怎樣理解.
由摩西頒佈律法.
到現在耶穌已經成就舊恩.
那現在往後.
基督教怎樣走呢.
是易拉命的.
猶太教的分支.
還是根本不需要再有.
猶太教是一個.
全新的人類的.
唯一舊恩呢.
所以必須要處理.
律法有什麼用.
有什麼限制.
第七節開始.
這樣.
我們可說什麼呢.
律法是罪嗎.
斷乎不可.
也是一個這樣的設問.
就是說預計.

$^{1041}$有人會這樣問的.
律法是不是罪.
是不是因為律法令我們犯罪.
我們沒有了律法就不怕.
其實就不是.
接著下面就說.
只是非因律法.
我就不知何為罪.
非律法說.
開始說十大.
這是出埃及的二十章.
十七節.
不可起貪心.
十屆的第十屆.
不可以貪戀別人的僕人.
那些財物.
家庭.
那些婦女.
等等.
這裡說.
不可起貪心.
我就不知何為貪心.
原來律法是有一個.
對我們心靈的定義.
我們裡面的.
舊人.
的那種.
慾念.
那種和神.
的要求很大距離.
的那種罪性.
律法可以.
幫我們很透析地.
讓我們看到.
原來我們是有.
神所不起悅的.
這些舊有的.
惡慾.
神所憎惡的.
這些舊人的.

$^{1081}$心思意念行為.
所以這裡說.
不可起貪心.
我就不知何為貪心.
然而罪趁著機會.
就藉著誡命.
叫諸般的貪心.
在我們裡頭發動.
因為沒有律法.
罪是死的.
意思就是說.
我沒有一個這樣的.
顯影劑.
將我們顯露我們裡面的.
諸般的貪婪.
這裡就告訴我們.
不可貪心就是.
神對我們人類的要求.
就是說神已經知道.
人類是有貪婪的.
生命.
舊有的生命.
所以有了這個.
律法在我們裡面.
我們知道了清楚了.
也定義了什麼才叫貪婪.
這個時候.
開始.
你去.
現在我早上有時會和媽媽喝茶.
吃那些腸粉.
就叫她給多一點點甜醬.
有時因為熟.
她拿了三碟來.
那是不是貪心呢.
又好像是.
但是.
是不是這樣去定義.
這個數量呢.
保羅這裡告訴我們.

$^{1121}$只是for example.
不可貪心.
這樣就定義了.
但是保羅在這裡很巧妙地告訴我們.
原來罪.
是會藉著.
諸般的貪心.
這個不關律法的事.
諸般的貪心是我們裡面的性情.
所以罪是.
趁著機會.
藉著律法的禁止.
反而會挑動我們.
有諸般的貪心.
這樣來到.
在我們裡面發動.
所以他想在這裡定義.
其實從來不關.
律法事.
律法是忠誠的.
律法是神頒佈的.
律法是公允的.
律法是定義那些人所做的行為.
當然是由心思開始.
是神所不起悅的.
再做是不能夠承受神的國.
所以.
罪就會趁著.
本來就應該沒有.
但是我們知道在世界.
墮落了的世界.
罪當然是魔鬼.
藉著罪.
來操控地球上的人類.
角色的形態.
所以已經在.
保羅這個描述裡面.
罪好像已經是已人化了.
是有為格.
是有目的.

$^{1161}$是有那種能力去挑動.
我們地上的人.
去範神的律法.
所以律法就不是成為.
一個磨心.
成為一個幫兇.
律法從來是忠誠.
這個我們要留意.
我們要牢記.
律法從來是忠誠的.
令到我們貪心.
就用這個例子.
保羅.
是我們裡面本身就有.
諸般的貪心.
是罪.
裡面的那個慾念.
挑動我們去犯這個律法.
就是不可.
變成我們就更貪心.
這個就是.
保羅想告訴我們.
所以第六章已經說了.
我們如果不呈獻.
自己給神.
就繼續在罪的權勢底下.
所以第六章已經定義了.
罪是一直.
不會放過我們.
和我們爭戰的.
我們去離開神的.
但是我們要認定.
正常就是我們歸神歸耶穌.
我們不再歸向罪了.
這樣的呈獻.
所以第七章仍然告訴我們.
罪更加具體.
就是原來罪不是.
引誘我們去貪心.
而是藉著律法.

$^{1201}$去挑動我們.
諸般的在內心裡面的貪心.
來去犯罪.
這裡告訴我們.
表面上沒有律法.
罪是死.
因為我們沒有那個.
律法在我們裡面的.
那個要求.
或者我們未被這樣挑動.
但是當有律法.
第九節我以前沒有律法的時候.
是活著.
但是誡命來到.
罪又活了我就死了.
這個很巧妙地.
到底那個運作是怎樣呢.
相信每個人都不同.
每個人那個對於.
神的誡命的.
那個能力.
那個持守.
那個軟弱受引誘.
都很不同的.
但是始終那個形態就是.
我們表面是.
很好.
很想跟隨神.
譬如猶太教的人.
個個都很想行律法.
個個都很想去合神心意.
的獻祭.
都能夠希望跟隨神.
這樣去生活.
我想任何基督徒信耶穌都是這樣的.
但是要找回這個.
根源.
保羅就說原來.
表面上沒有律法.
我們未知道之前.

$^{1241}$我們是自由的.
保羅說我們是活的.
但是當誡命來到.
就不是誡命.
活.
是罪又活.
是罪在我們裡面蠢蠢欲動的.
挑動我們舊有.
那個.
保羅說的舊人.
來到去.
不要理那些律法.
我們去貪婪.
挑起我們去.
刻意有目的地.
抗衡神這個律法.
就是不可貪心.
所以保羅就說我就死了.
所以似乎這一大段.
保羅是很想.
用一種.
共用共有的一班人.
所有基督徒.
所有人在律法底下.
不要理他是猶太教也好.
明白律法的.
外邦人可能他有.
律法的規範的道德.
或者更加.
具體的.
就是信了主耶穌的基督徒.
同樣都是在這個.
這樣的狀態底下.
我們就要很巧妙地.
認識其實.
不是律法對我們.
這個無情的要求.
而是罪.
藉著我們認識了.
這個律法就挑動.

$^{1281}$我們去挑戰.
這個律法的.
要求就是不可貪婪.
所以保羅說.
既然罪有活我就死.
那就是說保羅曾經.
都挺貪婪的.
如果他真是以.
他自己作為一個例子的話.
當然似乎就不應該說.
他個別地引用自己這個例子.
他說我就死了.
貪婪可以是.
貪那些人對他的.
稱讚的 譬如他是一個.
這麼崇高的拉比.
他是一個教法師.
他說他在律法的面前是無可指責.
可能這個是他貪婪.
的一種表現.
未必是在錢上.
或者是.
社會地位的.
無論如何第十節.
本來叫人活的誡命.
就是神本來期望.
每個人在這個世界生活.
都不要貪不屬於你自己的東西.
這個本來就是一個.
很大同的世界理想.
但是好像往往.
事與願違.
原來.
本來叫人活的誡命.
反倒叫我死.
表面上就是誡命.
使我們跌倒.
這個死不是死亡.
是跌倒.
不能夠在那個律法面前.

$^{1321}$我們可以得到神的喜悅.
就是說我們裡面.
有些可能已經做出來的.
例如那些貪官污吏.
過千萬.
等等.
又要雙開.
等等.
還要上到去那些.
正什麼位.
正部位.
那你就可以想像到.
其實已經不是第一天.
是很多年了.
由49年到現在.
當然不是代表外國沒有.
但是如果這麼平常化.
證明就是說.
不可貪心.
已經正正就像這段經文.
套用在任何人身上.
都好像是這麼.
全神地告訴我們.
他找不到那個辦法.
那個出路.
去制止這個問題.
成為已經是嚴重的社會問題.
那個原因.
就是一來不是.
根據聖經.
給我們的那個教導.
那個引進的出路.
甚至連這個定義.
都不肯接受.
不肯接受有罪在每一個人的.
生命裡面.
一路拉住我們.
只不過是藉著律法.
來挑動我們.
保羅這裡給我們看到.

$^{1361}$所以我們好像在探究.
人類的進化史.
那個成長.
現在保羅告訴我們.
新的人類.
不再需要舊有的.
那個形態.
不再需要局限在舊有的.
那種敗壞裡面.
怎樣能夠有出路呢.
層層地進深.
第十一節.
因為罪趁著機會.
就藉著戒命.
引誘我並且殺了我.
所以我覺得這個不一定是.
他個人.
是一個應該是集體.
所有人都.
共有的經驗.
當然或多或少.
或重或輕.
不是每個人都一樣.
那你就可以想像得到.
那個描述就很全神地.
將一個人.
由他很想不犯罪.
到他無法不犯罪.
那個心路歷程.
透過文字去描繪出來.
就是這樣去.
告訴我們.
他就是這樣.
所以不應該怪那個律法.
要怪的就是裡面的舊人.
沒有處理到.
不懂怎樣處理.
也沒有從羅馬書第六章講.
其實死了.
我們的法律地位.

$^{1401}$不需要再向罪活.
我們是向神而活.
所以第十一節.
罪趁著機會.
就藉著戒命.
引誘我並且殺了我.
這樣看來.
律法是聖潔.
戒命也是聖潔.
公義良善的.
做一個這樣的結論.
這樣看來.
一直批評的是罪.
在我們裡面.
從來都是與神對立的.
所以我們在人生.
何時我們回到.
不獻自己給神.
而獻給罪呢.
我們就會受到.
律法的捆綁.
繼續這樣去犯罪.
所以我們要留心.
保羅給我們的出路.
很清楚.
那個定義是很清晰.
律法和戒命是聖潔.
公義良善.
因為是由神.
對我們人類的規範.
是道德.
原來人類是必須面對.
自己的道德.
一生的道德.
在神面前.
所以救恩的足稱.
偉大就在這裡.
每個人要面對神.
對我們人類.
由出世之後到死之前.

$^{1441}$那個道德的標準.
所以誰可以站在.
永活的神面前.
所以救恩的偉大.
就在這裡.
每個人都是罪人.
在神的道德追討的規範裡.
現在就告訴我們.
所以我們明了.
與律法無關.
我們就知道.
我們不是逼自己去行律法.
我們有另一個新的出路.
就是在主的救恩.
在聖靈.
在主的力量裡.
我們就可以.
去遵行神的要求.
十三節.
既然如此.
那良善的是叫我死嗎?.
即是律法.
良善的.
叫我死的原來是罪.
即是叫我們跌倒.
離開神.
叫我們不去行神的.
律法誡命的原因.
是罪叫我們去跌倒.
所以.
沒錯了.
這個是十三節.
但罪藉著那良善的.
叫我死.
就顯出真是罪.
叫罪因著誡命.
良善和下面誡命.
是同義詞.
罪因著誡命.
更顯出是惡極了.

$^{1481}$將罪.
好像很繪影繪聲.
一個惡魔.
一個怪獸.
在我們裡面.
所以我們的內心.
如果我們有一段時間.
不歸向神.
罪就會在我們裡面.
好像慢慢.
由單細胞.
變成一路膨脹.
一路越來越發大.
即是越來越.
去到一種.
任何人都見到.
你這一輪真是犯罪.
犯得很厲害.
你可以想像到.
真是顯出它的惡極.
十四節.
我們原曉得律法是屬乎靈的.
但我是屬乎肉體的.
因為是從神而來的.
那個道德標準.
所以一定是屬靈的.
神期望我們每一個人.
都要達到祂期望的.
我們那個人生.
所以.
有些人間地獄.
人渣等等的描述.
即是.
惡魔 禽獸.
這樣的描述.
就真是很全神地.
神的律.
本來就是聖潔的.
是屬靈的.
但我們是屬乎肉體的.

$^{1521}$你就可以想像這種對立.
原來救恩可以將一個屬乎肉體.
禽獸不如的人.
可以轉變成.
正常基督徒.
就是屬靈.
原來每個基督徒.
都應該這樣看自己的正常.
這樣的正常.
就是成為羅馬帝國.
以後的世界.
的一個宣教力量.
就不是怪異的.
死守某些戒條.
我們可以想像得到.
原來是一個這麼正面的.
救恩.
帶給我們不再屬乎肉體.
而是屬靈.
十五節.
因為我所作的我自己不明白.
我所願意的我並不作.
就是他所願意.
就是很想行神的律法旨意.
神的旨意.
要他來進行神的.
道德誡命.
但是他所願意都沒有用.
因為他沒有力量去做.
如果我們不要再考究.
到底是微信之前.
還是在律法之下的.
保羅快利塞人的時代.
因為他信了主之後.
還有軟弱.
如果我們不要這麼仔細去考究.
是信主前後的.
這種苦況.
如果一直都看一個很正面.
將罪.

$^{1561}$將我們這個舊有.
生命和新生命.
的一個對比的比喻.
這樣的.
一個分析.
反而就不需要陷落.
到底我們要執著.
保羅在信主之後.
還有軟弱.
還是微信之前的那個光景.
因為.
聖經學者就會在這裡.
就會有一些.
不同的理論.
如果我們整體來看.
其實我們就未必需要給.
他這個很個人.
的經驗.
因為似乎一直講下去.
上下文都是共有的.
經驗.
不因為你是一個.
傳道人.
不因為你是一個普通的信徒.
你就沒有這些經驗.
其實應該每一個人.
只要他有一個心想.
一個律法.
他就已經有這個掙扎.
這個掙扎.
當然信了主有聖靈.
那個就再.
有一個更加明顯的.
那個催逼.
那個感動等等.
我們將保羅這個.
的掙扎表面就是我我我.
但其實我覺得是一個.
共有的基督徒的經驗.
甚至乎.

$^{1601}$在猶太教底下.
都一樣有這樣的經驗.
若我所作的是.
我所不願意.
我就答應律法是善的.
就是說我很想做到.
不貪婪.
但是我沒有能力去不貪婪.
我只能夠.
答應.
認同.
律法本身是良善的.
只是我裡面.
沒有這個力量去行良善.
這樣大概的意思.
所以這個答應律法是善.
就好像.
越來越以人化.
將他的內心描繪出來.
就是有心無力.
做不到.
既是這樣.
十七節就不是我作.
而是住在我裡頭的罪作.
它不是.
委過於罪.
不是不關我的事.
我這樣做.
是不由自主的.
其實不是這個意思.
就變成.
扭曲了那個真理.
它的意思就是說.
我根本沒有行善的力量.
而罪在我裡面.
由始至終都不斷的.
去引誘我們.
去抵擋.
違背神.
好像阿當不住.

$^{1641}$不住的受住.
蛇或罪的引誘.
我們每一個人.
都不斷有這個.
這樣的律在我們裡面.
來挑戰.
我們要行善的.
這個可能.
所以保羅在這裡.
某程度真的很全神地.
剖析了我們人性裡面的.
運作.
其實沒有.
一些儀器.
可以將心靈道德.
照出來.
但透過真理的聖經.
保羅這樣描述.
真是將一個人.
他什麼時候落在罪裡面.
受到壓迫.
律法的枷鎖.
到他真的要行犯罪的行為.
然後就離開了神.
得罪了神.
在一個循環的律法裡面.
就這樣來到.
越來越陷落在罪惡裡面.
但是基督徒.
當然這裡未講.
七章後面就會講.
所以第十八節.
我也知道在我裡頭.
就是我肉體.
之中沒有良善.
因為立志為善.
由得我.
只是行出來由不得我.
故此我所願意的善.
我犯不作.

$^{1681}$我所不願意的惡.
我倒去作.
又是其中一段的描述句裡面.
所以這些不是.
名言來的.
不是格言來的.
不是說基督徒是這麼軟弱的.
是沒得救的.
你不要再逼我做好.
我是這樣的.
其實我們完全解錯了第七章的羅馬書.
是不斷去描述.
我們裡面有一個.
要接受我們有一個罪的律.
在我們裡面.
那個舊人就算信了耶穌.
他也沒有死.
因為罪本身就不是一個物體.
不是你割了他就等於.
麻煩順手金峰幫我割了他.
沒有這件事.
罪是一個敵對神.
的一種.
近乎有位格在保羅的描述裡面.
但他本身是沒有位格的.
但是這種敵對.
已經與生俱來.
在阿當的舊有的.
那個族類裡面.
是與生俱來有的.
我們沒得選擇了.
因為阿當就是這樣背逆神.
罪就從一個人.
阿當進入了世人世界.
所以我們只能夠.
離開世界.
將來就沒有罪了.
所以離開世界第一件事.
就不是你馬上輕飄飄.
而是你沒有了罪.

$^{1721}$不再犯罪.
你想想那種快樂.
這種快樂應該大於.
你痛十幾秒.
十幾天的.
癌症,癌末那種的痛苦.
那種的絕望.
但是你想想.
一直說從第五到第七章.
是說我們仍然在.
罪底下的.
我們基督徒現在要記住.
真正的基督徒.
是要拒絕.
不是逼自己拒絕.
而是法理上我們.
要告訴自己.
我們不再在你的權勢.
在耶穌裡面.
但是真正的不在那個權勢.
就是死了就沒有了.
那個罪的律就沒有了.
罪的律就沒有了.
這件事是非同小可的.
你想想.
我們將來一離開世界.
罪的律就馬上離開.
那種的狀態是怎樣呢?.
簡直是完美到不得了.
所以現在保羅他一個.
那麼敬虔的聖徒.
這位拉比.
這位與神.
同工的使徒.
他現在都再一次.
告訴我們.
在他的有生之年.
他都是在.
罪的利誘.
逼迫.

$^{1761}$那種的.
無奈底下.
但是保羅很早就找到這個秘訣.
這個真理.
就是他歸向耶穌.
他將自己完全奉獻給耶穌.
為耶穌而活.
這樣就可以.
來抗衡這個罪的.
枷鎖.
的綑綁.
所以記得17到第19節.
絕對不是我們.
不成長的.
格言.
由不得我的.
所以.
不要這麼高要求.
因為在我裡面有一個律.
這是一個描述.
舊有的人.
每個人都在這個律底下.
但是基督徒.
就是在耶穌這個愛的律底下.
救恩.
偉大救恩的律底下.
我們來到去活.
所以表面上就是他.
故意地這樣剖析.
我本想.
做行善的.
神的旨意的.
但是無奈裡面那個律太強了.
令到我們.
我想願意做的就反不做了.
我所不願意做的.
我就倒去做了.
好像.
不關我事.
身不由己.

$^{1801}$其實不是.
正正就是關我事.
真實存在.
就算信了耶穌.
就算有道德的律.
但原來保羅說不用擔心.
你承認他的存在是必須的.
但原來要承認.
另一個新的律.
在我們裡面承認中.
在我們裡面.
我們願意.
就可以被他進佔.
當然用佔就不對.
是他會進入.
來與我們同在.
陪伴幫助我們.
去過.
神喜悅的人生.
所以其實五到第七章.
是很正的.
拆解心靈問題.
某程度.
都涵蓋心靈的.
醫治.
某程度對一些心靈特別脆弱的人.
時時都活在自己的.
自我控訴裡面的人.
這裡都已經是一個很大的出路.
承認有舊人存在.
不要緊 我承認你的存在.
但是我更加要.
在舊因裡面承認.
復活的耶穌在我裡面.
這樣的力量就真的無窮.
無窮的力量.
所以第20節.
就約我去作所不願意.
我剛才說了.
第21節就出現了.

$^{1841}$我覺得有個律.
開始說一個.
好像一個定律.
一個定律這樣.
每個人都應該.
有相同的經歷.
每個人都是一樣.
每一個生而為.
亞當後人的.
舊有的族類.
你都要盡快發現.
你有這個律.
這樣就很不同了.
因為下面說另外一個律.
生命,性靈的律.
八章應該是第三節.
沒錯.
第二節和第三節的對比.
所以你可以想像到.
現在引進的是.
人人都犯罪的律.
我覺得有個律.
就是我願意為善的時候.
便有惡與我同在.
因為按照著.
其實我裡面的意思.
我是喜歡神的律.
但我覺得.
肢體中另有個律.
和我心中的律交戰.
把我擄去.
叫我服從那肢體中犯罪的律.
大家有時看那些劇集.
久不久便有魔鬼黑色.
白色的天使.
很小的.
在你的肩膀上.
來慫恿你.
某程度上都是保羅.
這樣的說法的一種擬人化.

$^{1881}$都不可以說他錯.
因為保羅說的.
更加人人都是這樣叫律.
即是你在香港.
你現在在錦繡的203.
跳到地下.
在廣場.
你會很痛 因為二樓不是很高.
不會跌你不死的.
不過很痛 除非你的頭衝下去.
當然會更嚴重.
那你在二樓的錦繡.
和我在元朗大堂的二樓.
和我在嘉湖山莊的二樓.
跳到地下 是不是會不同呢?.
應該是一樣.
因為是有個律.
這是物理的律.
但現在說的生命的律.
就是亞當後人.
必然會有的律.
就是我們要行善.
為善的時候.
就會攔阻我們.
挑動我們去犯罪.
所以保羅說.
不是他的發言.
而是神的真理的宣告.
告訴我們.
為什麼萬民都需要救恩.
而不是需要教你做好人.
而是告訴我們.
我們每個人都在亞當的.
那個生而為罪人的這個律底下.
所以這裡就告訴我們.
我們越早發言.
就越能夠盡早享受.
救恩帶給我們的那個平安喜樂.
和那個生命的釋放.
所以對不起 第22節.

$^{1921}$按照著我裡面的意思.
我喜歡神的律.
把我擄去.
因為交戰.
叫我服從那肢體中犯罪的律.
似乎這裡想告訴我們.
沒有主耶穌愛在的律.
愛在的那個愛的救恩.
在我們裡面支撐.
幫助 陪伴我們的時候.
我們每次都好像打敗一樣.
但現在就告訴我們.
阿寶萊就用24節來描述.
他這種有心無力.
很想 但是沒有力量去走出神的要求.
他說 我真是苦啊.
誰能夠救我脫離這個取死的身體呢?.
抱歉 我忘了按powerpoint.
所以一直讀那些經文.
都差不多在說這些的意思.
怎樣能夠從律法轉向歸屬基督呢?.
就是我們已經在罪裡面死了.
律法其實是揭示罪的本質.
律法是公義良善的.
但是18至25節.
就是在這個無力行善掙扎底下.
怎樣的出路呢?.
這裡就告訴我們.
取死的身體.
這個字挺特別的.
我聽一些解經家都挺權威的.
就很清楚說.
當時羅馬帝國懲罰那些死囚.
尤其是叛國那些.
用什麼來懲罰呢?.
找一些屍體綁住他.
跟那個屍體.
一直這樣從早到晚的生活.
那些屍蟲慢慢咬死那個活人.
這樣的聯合.

$^{1961}$就是說保羅告訴我們.
這個取死的身體就是我們舊有的罪.
那個舊人.
在我們裡面一直跟我們蠶食我們.
一直給我們不斷的養分.
就是不要行神的旨意.
就是要行神不喜歡的事.
那就是罪的律.
好像就是取死的身體.
如果保羅真的是這個意思.
真是這樣來描述就很恐怖.
就是說原來我們傾向去罪.
原來是不應該的.
因為我們已經是屬於耶穌的了.
我們已經不再在這個取死身體被他綑綁的了.
應該就是我們是一個自由釋放的.
跟隨耶穌的.
所以今晚我自己的領受最大的就是.
我們原來可以過正常不犯罪的基督徒生活.
我們是可以不再被這個取死的身體.
就是舊人來綑綁著我們一起.
所以25節.
保羅好像很戲劇性.
好像你在這裡將大結局濃縮了半分鐘.
第25節就是出路了.
第7章最後這一句.
記住.
那個取死的身體仍然在我身體裡面.
這個當然是比喻.
這個舊人仍然在我們裡面.
這個聖情仍然在我們裡面.
是一個律.
但是原來有外在的幫助.
原來有一個永久常存的幫助.
就在我們的面前.
在我們的心裡.
就是這位主耶穌.
我們去呼喚他.
他就立即成為我們的力量.
可以抗衡這個罪的律.

$^{2001}$第25節.
果然最終又是回到.
就是和主耶穌的生命關係.
靠著我們的主耶穌.
就能夠脫離這個取死的身體.
這樣看來.
我以內心馴服神的律.
我肉體卻馴服罪的律了.
記住不是各自表述.
好了算了.
這一部分就是說.
馴服神的律.
另一個部分就是.
馴服罪的律.
其實不是這個意思.
就是說很清楚.
裡面的舊人.
是一生都存在.
不會離開的了.
所以維斯利約翰17世紀.
那個叫拔罪根.
second blessing.
就是說你信了主.
你要經過一個考驗.
當然有他們的規範等等.
我用最簡單的邏輯.
讓大家知道.
曾經在歷世歷代.
有一些聖徒有這樣的主張.
由那一刻開始.
你的罪根就拔掉了.
你由那一刻開始.
你就不會再犯罪了.
信主那一刻就叫first blessing.
由你不再犯罪開始.
就是second blessing.
這個就肯定不適合羅馬書了.
因為保羅說.
舊有的律.
一直存在一生存在.

$^{2041}$只要離開世界.
才不再存在.
但是原來可以不犯罪的原因.
不是拔罪根.
而是我們有生命聖靈的律的主耶穌.
在我們裡面成為我們的力量幫助.
這個才是真正不犯罪的正常.
所以不是突然間去到某個階段.
因為衛斯理耀漢的確他個人的經歷.
就是曾經他體驗過這種無盡的快樂.
感恩.
和從此他整個生命的脫胎換骨.
他集合了就以為.
這個就是那個blessing.
這個也是早期1900年.
第一波的靈恩運動.
就以這個作為藍本.
就是他們講方言.
就是那個second blessing.
就好像有一種過關的形式.
其實保羅這裡從來不是說過關.
你突然間會有一個領悟.
你從此就不再犯罪了.
不是這樣的.
保羅說.
從來我們如果不依靠耶穌.
都會犯罪.
但是都會犯罪.
就不是正常基督徒的生活.
不犯罪在耶穌裡面呈現自己.
不犯罪才是正常的基督徒.
這個就是五到第七章.
保羅很想給羅馬的信徒.
有一個這樣的理解.
今天剛剛好就講完第七章.
一個段落的完結.
重重複複.
我們要肯定.
罪的律不是我們做過什麼壞事.
是我們是亞當的族類.

$^{2081}$我們一出生就必然.
罪就伴隨著我們.
直到我們死了.
但是沒得解決.
因為不信耶穌.
死了之後有審判.
但是我們作為基督徒.
我們就不應該憧憬將來永生.
今生就算了無奈了.
因為有罪在我們裡面.
我們都只能夠是某一小撮人.
他們享受這種奧秘的blessing.
我們是一般的平凡的平信徒.
就不應該享用這種second blessing.
這個本身就是錯的邏輯.
保羅就在這裡.
每一個基督徒.
記得正常是可以不犯罪.
因為不是我們沒有罪性.
更加不是我們有能力.
靠自己不犯罪.
而是我們是一班呈獻了給神的.
新生的新人類.
亞當的.
另外一個亞當的族類.
我們是要認定自己是新的族類.
所以second blessing不是一種懸殊.
是應該是普遍每一個基督徒都可以享受.
可以這樣進入的.
這個就是第七章.
比較仔細講律法的功用.
因為都九點零二分.
我們就今晚停在這裡.
抱歉 因為下個禮拜四.
我們要開洪恩的發展組會議.
就不了其他時間.
所以他們又要改星期四.
所以就要停一次.
我們下個禮拜四晚的茶經班.
除非大家很想將它改為星期五.

$^{2121}$那我都OK的.
看大家的意見.
如果不是就要停一次.
看大家想不想進度可以快一點.
因為我們還有大半沒看的.
現在暫時都是16講.
我們先祈禱.
感恩.
天父讓我們發現不是特別的屬靈基督徒.
才可以不犯罪.
而是每一個基督徒都應該過一個正常.
在耶穌的救恩裡面.
因為歸屬了給新的主人耶穌基督.
和我們是呈獻自己給神.
我們就可以過一個神喜悅不犯罪的生活.
而且這個動力是來自聖靈.
是來自生命聖靈的律.
盼望我們能夠將整個大的脈絡.
越來越清晰地去拆解.
分析和具體的應用.
使我們每一個弟兄姊妹.
都能夠得著這份生命平安的福樂.
更加我們能夠看到.
我們何時離開主.
其實何時就在罪的律底下.
再一次被挑動去犯罪.
所以盼望主給我們一個警醒.
一個感恩的心靈.
再一次多謝你讓我們能夠逐段.
逐段有系統地去解釋羅馬書.
盼望主的聖靈透過這本書.
讓我們更深地認識.
萬民真是需要福音.
真是唯有福音去拯救萬民.
願主幫助我們大家一起.
在很有限的人生和資源.
齊心地努力.
興旺神的福音.
叫多人信主.
我們這樣感恩禱告.

$^{2161}$奉耶穌基督的名.
阿們.
\newpage



\section{羅馬書}
\label{sec:HoQGEbdeqSk}
\textbf{2025-07-24 周四晚查經班《羅馬書(17):分題:恩㤟與律法》 講員:曾裔貴牧師}
\newline
\newline
連結: \href{https://youtube.com/watch?v=HoQGEbdeqSk}{\texttt{https://youtube.com/watch?v=HoQGEbdeqSk}} ~~~~ 語音日期: 2025-07-26
\newline
\newline
\hyperref[sec:PTM8SSdjvVE]{\small{< < < PREV SERMON < < <}}
~
\hyperref[sec:index_chronic]{\small{[返順時目]}}
~
\hyperref[sec:index_scriptual]{\small{[返順卷目]}}
~
\hyperref[sec:mDmug5anuis]{\small{> > > NEXT SERMON > > >}}
\newline
\newline
$^{1}$我們一起祈禱.
再一次感恩天父讓我們恢復.
晚上周四晚的茶經.
我們開始看到羅馬書第八章.
聖靈和恩典的運作.
如何能夠全力在我們每一個信徒的心靈裡.
來運行.
來對抗在我們內心裡的罪和死的律.
這是一個驚人的發現.
是很小的信徒都必須要這樣來認識.
以致能夠在神的面前.
原來我們今生是可以過一個不犯罪.
選擇歸向神.
來向神奉獻的人生.
願神帶領我們明白裡面的經文.
特別是舍圖保羅是很平鋪直敘的.
將救恩的偉大在人心裡面的工作.
來講出來.
讓我們能夠從其中得到很大的啟發.
多謝主原來正常的基督徒生活.
就是一個可以不犯罪.
可以歸神為聖的生活.
願主幫助我們這樣祈禱.
奉耶穌基督的名.
阿們.
大家由第五章一直開始.
你都已經看到保羅塑造一個新的人類.
就是在基督裡面的新的族類.
如果哥林多後書的五章十七節.
如果有人在基督裡面就是新的創造.
新造的人就是New Creation.
新的創造.
舊的事情都過去.
全部都變成新的.
保羅就是告訴我們.
我們在基督裡面享受一個恩典的地位.
其實就是本來不配蒙恩的.
本來是與神為敵的.
但是現在能夠在救恩裡面.
我們就站在一個蒙恩典的地位.

$^{41}$而且我們是神的兒女這個身份.
所以是要持續保持這份與神的和好.
因為得來不易.
是耶穌為我們在十誡上所成就的.
所以我們要很主動的.
來到上上與神保持一份和好.
相愛的關係.
這是羅馬書第五章很想帶給我們知道.
接著第五章的第二個段落.
十二到二十一節.
就是很想說.
本來我們是舊有創造的罪人.
阿當的族類.
因為阿當的背影.
阿當的犯罪.
阿當這樣來去違逆神的心意.
罪就進入了人類的社會.
罪就在人的內心.
所以變成與神俱來的某程度就是所謂的如罪.
現在我們在救恩裡面.
神沒有將這個罪性徹底的根除.
但是也讓我們人在地上.
看到我們如果不依靠神.
我們都會繼續回到我們舊有的罪的性情.
罪的主宰或者叫引誘.
我們就會跌落在神的殯路,刑罰裡面.
所以這是第六章.
本來很快就已經探討.
我揭開一些PowerPoint.
幫大家重溫一點.
可能有些印象.
都是以前的一些PowerPoint.
給大家再看一些舊有.
幫大家有些溫習.
所以第六章很重要.
真正的基督徒是應該過一個新的生活.
就是一個我們在恩典底下.
.
各位抱歉.
如果是頂展會.

$^{81}$如果是賽仙.
麻煩將那個靜音就行了.
不好意思.
因為變成不需要影響其他人.
我們繼續說.
在主的救恩裡面.
我們就是享用一個白白的恩典.
所以不可以作為一個犯罪的藉口.
因為在我們內心裡面.
其實那個舊有的性情.
我們要記得在羅馬書給我們的啟示.
神從來沒有將我們舊有的罪.
所謂罪根拔除的.
也沒有任何的方法可以拔罪根的.
我們那個罪的性情.
是必須透過聖靈在我們裡面的引導.
感動.
帶領.
我們就可以不再在罪的奴役和引誘底下.
而是我們越來越傾向.
在神的兒女這個新的生命的生活.
所以羅馬書六章就是很想說.
本來是強迫我們不准犯罪.
要我們不可以濫用恩典.
但其實你看看羅馬書第六章.
其實整章的含義就是說.
原來基督徒當你意識到原來這份這麼大的恩典.
使我們成為神的兒女之後.
原來我們是可以選擇去跟隨一個新的.
很好的一個主人.
這個選擇權就是隨著我們的靈命的成熟.
所以不是有什麼新的啟示.
新的突破.
石破天驚.
而是我們越來越認識我們的人性.
越來越認識神的愛.
在我們生命裡面的不斷的介入.
那種愛的憐憫.
使我們很樂意成為一個誠信的.
一個基督徒.

$^{121}$抱歉 弟兄姊妹不要介意.
真的要靜音.
否則會影響其他的弟兄姊妹就不太好.
不要介意.
否則有時候聲音會影響到別人的聽力.
當然如果弟兄姊妹你想發問是可以的.
不過要舉手讓我看到.
第六章一個很重要的發現就是.
我們原來是可以選擇不犯罪.
這個體驗我想是前所未有的.
因為我們不信耶穌是不可能不犯罪的.
但是信了耶穌.
越來越認識我們神兒女的身份.
原來我們可以選擇不犯罪.
當然第八章就說具體的運作.
所以第六章就似乎是禁止.
是警告是一種很負面的.
不可犯罪.
但是更加正面的應該就是可以不犯罪.
人就是喜歡勸勉.
人很不喜歡強迫你做.
禁止你做.
越是禁止越是強迫.
有時候就越是想犯規.
但是如果我們慢慢發現.
原來不是的.
原來我們是真的可以靠著主的恩典.
我們是可以不犯罪的.
所以有幾個重要的詞語.
就是我們要主動的獻上自己.
第六章應該是第十七節.
那裡告訴我們.
我們一定要主動的獻上.
十六,十七節.
豈不曉得你們獻上自己作奴僕.
順從誰就作誰的奴僕?.
或作罪的奴僕.
以至於死.
或者作順命的奴僕.
以至於成異.

$^{161}$所以保羅很強烈的勸勉羅馬弟兄姊妹.
必須真正看到自己一個新的生命.
就是過一個正常奉獻給神的生活.
這就是基督徒的生活.
應該再正確一點.
這就是正確的基督徒生活.
正常的基督徒生活.
就是不斷呈獻自己.
就是說那個心是湧入到一個地步.
什麼時候神需要我做一個奉獻.
我就願意奉獻了.
因為主權是在神那裡的.
這樣一個不犯罪的心靈.
就是一個不犯罪的生活.
所以這裡就是.
其實某程度上真的是人類一個大的發現.
因為任何地上的政體.
不同的政治主張.
都希望自己的國家是國泰民安.
大家知法守法.
大家有規律.
大家彼此相愛的生活.
整個社會不需要有監獄.
不需要有任何的日出而作.
夜不閉戶.
但是這些好像同境的一種幻想.
從來都沒有出現過在人類的社會.
因為人的罪性.
所以只能夠不斷製造,建立不同的規範.
來讓犯罪的人的成本要很高.
但是從來都解決不了人性裡面的.
無論敵對,仇恨,貪婪等等的問題.
這就是罪性的使然.
但是原來保羅他到底有什麼特別之處.
神給他有這麼大的一種智慧,啟示.
他能夠很輕描淡寫.
慢慢平鋪直敘地.
將人性的問題去剖析.
而且提供到一個合乎神真理的答案.
就是不斷看到自己神兒女寶貴的身份.

$^{201}$當你越珍惜自己.
很珍貴自己這個新的身份的時候.
你就在這個新的身份裡面.
不斷有一個主動的呈獻.
主動的獻上自己.
這樣的一個奴僕的人生.
其實就是神兒女蒙福,蒙愛的人生.
很矛盾的.
做奴隸何來有自由?.
何來有愛呢?.
但是往往越是奉獻的人.
就越是感受到神愛的偉大.
所以我們很多時候羨慕一些宣教士,牧者傳道.
只不過因為在他的工作,他的崗位.
他不斷都必須要走前一步.
不斷地體驗神的工作.
所以我們往往忽視了.
以為我來到崇拜,我來到聚會.
我能夠做的就是給一個半小時,兩個小時神.
其餘的時間往往我們忽視了.
其實我們在其他的時間.
如果你都這樣呈獻自己.
其實你一樣體驗到那種神兒女自由的身份.
所以求神都幫助我們.
第七章就是.
保羅仍然很想來探討.
有實際的問題我們提出過.
因為猶太人會繼續質問你們基督徒.
為什麼做到一個好像不需要守猶太教的律法.
摩西十誡衍生的法律.
猶太教的那個律法.
到底你們是憑什麼可以不需要.
好像我們猶太人一樣有宗教的規範等等呢?.
所以第七章必須緊接第六章.
這個自由是用什麼來換取呢?.
這個呈獻和律法又有什麼關連呢?.
第七章又將一個好像在內心很想行神律法.
所以未必是猶太人.
當然這章聖經就引發不同的解經家有不同的爭論.
當然是討論不一定是爭的.

$^{241}$到底是保羅在未信主之前的掙扎.
還是信了主之後.
他的一個內心的矛盾.
有了聖靈的保羅.
他來到有這樣的掙扎.
還是他根本未信耶穌之前.
他自己就在律法底下.
這樣爭很難找到一個正確絕對的答案.
因為保羅沒有強調他自己是在未信主前.
還是信了主之後仍然有這樣的掙扎.
所以其實對任何一個人都對的.
每一個基督徒.
我們什麼時候不願意讓神的靈在我們裡面.
因為這就是第八章的解釋答案.
所以什麼時候我們不願意讓聖靈.
引導我們的內心.
我們就會落在罪對我們的引誘.
就會令我們好像落在一個未信主前的狀態.
這樣跟神有一個敵對的關係上的疏遠.
所以這裡我們就要留意.
當然第七章原本的意思.
是很想說到底律法有什麼功用呢?.
律法有什麼好?.
但它也有什麼限制呢?.
這個就是保羅很想讓猶太人要停了.
沒有耶穌的出現.
猶太教仍然是受著摩西的律法的規範.
你可以說是籠罩著他們這個民族.
是正確的 是神的律法.
但是當耶穌基督來了這個世界救恩完成.
那麼律法是不是應該有它的時限.
有它的終結.
有它的限制要過渡到歸屬基督的轉向呢?.
這個就是很重要的一個轉向.
所以第七章有它的過渡性.
有它對可能是未信的人有這個要求.
但是對猶太教一些不認識基督的人.
或者不肯將自己的心帶到基督裡面的猶太教的人.
這一章也是一個很重要的提醒.
就是說如果他們沒有基督.

$^{281}$沒有將自己歸屬基督.
其實律法對他們來說.
只會是不斷的壓迫他們的心靈.
他們沒有真正成為蒙福自由神的兒女.
他們是亞伯拉罕的子孫.
他們是每個星期有去會堂去獻祭.
去向神祈禱.
但是他們的心靈仍然是受著律法.
對他們的內心的壓迫.
因為剛才說過.
罪的性情唯有到了死亡那天才真正結束的.
罪是會藉著律法來引誘人.
令人不能夠實行神的要求.
所以不是律法本身的問題.
是我們的內心不能夠靠自己去實行神的律法.
所以罪就在那裡潛入.
非法的潛入我們的內心.
其實從來都是在我們的內心的.
但是就好像隱藏在律法底下.
以為是我們自己要實行律法.
但實際上我們就不能夠靠自己本身的人性.
來喜歡,遵行,樂意去守神的誡命.
這就是第七章 應該是第24節.
將那份無力感用一個圖像化來揭示出來.
我們落在一種完全不能夠在律法面前去抵抗罪.
原因是因為律法本身只是揭示我們罪的存在.
它沒有能力,它很間接地不能夠帶我們去勝過罪的引誘.
不能夠讓我們能夠在神的面前.
成為一個蒙福,自由的神的兒女.
那怎麼可以得到呢?.
第八章就立即解答.
原來生命,聖靈的律在我們信主的一刻.
就已經進入我們的內心.
所以一個信耶穌,愛主的弟兄姊妹.
跟一個很老練,成熟,被神用的工人完全是一樣的.
都沒有多到因為她突然間拿到一本秘笈.
她就修煉 所以她就能夠突飛猛進.
從來都沒有秘笈的.
原來不多不少就是聖靈在我們內心的工作.
那就是一件很天大的真正的秘密了.

$^{321}$原來每個人都領受聖靈.
一信耶穌都領受聖靈.
比如我這兩天 星期二.
聖經班24個小朋友.
老師問誰要信耶穌舉手決戰.
說一起舉手.
老師就叫牧師為這二十多位小朋友祈禱.
每個人閉上眼睛跟著說.
不用說我的名字 說回小朋友自己的名字.
每個人說自己的名字跟著說.
如果我們真心.
如果我們真是有一個屬靈的透視.
我們就看到每一個小朋友一信耶穌.
他接了聖靈的 他不知道的.
聖靈就住在他的內心 他真心信的.
可能比你和我在報道會決戰更加真心的.
他不懂的 但是如果我們真是這樣相信.
聖靈就在每一個信主的人的心裡面.
就進入工作了.
所以真是不多不少的.
每個人要經歷不犯罪的生活.
唯有就是要看聖靈已經在我們裡面.
成為我們可以抗衡不犯罪的那個力量了.
所以原來不是律法令我們不犯罪.
是讓我們知道我們是沒有力去靠自己去勝過罪.
所以應該是24字在這裡.
我們看看這個PowerPoint.
律法和犯罪的張力.
再將第七章這裡稍微再分點來說.
律法是揭示罪的存在.
因為非律法我們就不知道何為罪.
神的要求 但是沒有辦法使我們能夠勝過罪.
勝過就是說我們能夠滿足神對律法的要求.
在我們生命裡面的.
那意味著我們知道罪的存在.
但是沒有辦法勝過它.
怎樣去勝過它呢?.
又仍然是那個原因.
就是我們是神兒女.
我們就可以勝過罪.

$^{361}$這個就是正常的基督徒生活了.
真正的正常就是.
原來我們真的可以不犯罪.
有條話告訴我們.
律法本身本來已經一而再強調.
是聖潔的 是公義的 是良善的.
所以不是律法激發我們犯罪.
是罪藉著律法激發我們心中對神的背影.
意味著原來.
本來也給我們一個很重要的啟示.
原來我們要搞清楚.
原來有神的律的就是聖靈.
也有我們內心那個渴慕行善的那個心.
就是我們裡面的.
本來我們的本意是很想愛神.
行神的旨意的.
但是無奈地.
我們是亞當的族類.
我們有罪的律在我們每一個生而為人的身上.
就在我們裡面.
所以有幾個律在我們內心來交戰.
所以保羅就在第七章的二十四節.
就很明顯地告訴我們.
我們如果什麼時候違背神的命令呢?.
保羅就提到.
他說我這個人真是苦.
這是漢語《聖經》逆背.
誰能夠救我脫離這個必死的身體呢?.
我查過一些文獻.
甚至我搜尋AI.
都是有這樣的解釋.
就是說當時羅馬帝國.
對待一些嚴重的死囚.
很簡單.
那種殘酷性.
是你難以想象的.
就是古代的酷刑.
就將一個死囚.
你殺了那個人.
就用那個人綁住你.

$^{401}$那個是死人.
跟你在一起.
每天你在監獄裡面.
他綁住你在你的身體.
那些屍蟲是慢慢進入你裡面.
感染你.
咬死你.
這樣你慢慢死.
你殺了那個人.
人家一秒鐘就死了.
但是你不知道要多少天.
你才慢慢被侵害而死.
這些死亡是很恐怖的.
你每天就綁住一個死人在你的身上.
保羅就說.
這個就是取死的身體.
所以如果我們能夠.
如果這樣的典故.
是真的保羅來引用羅馬帝國的酷刑.
那就難以想象.
聽羅馬書的意大利的基督徒.
他們一定知道當時有這樣的行刑.
就是在清朝.
去到五門斬首.
那些人一定要出去看.
是要拿那些大隻佬.
特別是胸口有毛的.
拿大關刀斬在頭上.
這樣脫了頭.
是要人家個個侍眾給人看的.
就是說.
代表著如果你不聽話.
犯罪.
可能叛國.
可能其他的事情.
你都是這樣的下場.
個個都害怕就不敢做了.
所以保羅在這裡描述.
原來罪性的恐怖,可怕.
是想基督徒看到這個問題.

$^{441}$也再一次告訴我們.
原來主耶穌是為一些這麼大罪的罪人釘十字架的.
所以我們信耶穌.
應該越來越看到自己的本質.
不在恩典裡面的人的本質.
是一無是處的.
只是在神的憤怒下.
我們本來就是這樣的一個取死的身體.
就是說如果沒有了救恩.
沒有了聖靈.
沒有了我們時常去倚靠聖靈.
我們的智慧.
我們的思考.
有時候會偏向於鬼詐.
極端.
有時候憤怒.
狹隘.
這些就是告訴我們.
我們什麼時候不倚靠聖靈.
我們就會落在這種取死的身體.
你可以想像得到.
是伴隨著我們的人生.
直到離開世界的時候.
才沒有了這個取死的身體.
所以你可以想像得到.
離開教會.
某程度上是對的.
離開了教會就沒有救恩了.
當然這句話的原初的解釋就是.
天主教因為基督徒.
因為馬丁路德的改革.
就說教會以外沒有救恩.
就是說你分離的弟兄.
你快點回來天主教.
否則就沒有人可以將那個祝福給你了.
他們的教會就是自己的天主教.
在天主教以外就不會有救恩了.
但是我們當然知道不是這個問題.
但是如果我們用一個在這個角度.
在七章二十四至八章一節.

$^{481}$在這個轉接句裡面.
保羅是說得很誇張地.
一個這樣的圖像化的.
用一個死囚的痛苦的死亡.
來到誰可以解除他呢?.
有一個人願意解除他.
就是主耶穌了.
所以第二十五節.
所以這些轉接句.
大家有時候如果能夠很熟悉.
就能夠幫助我們看到.
原來在保羅的眼中.
罪那種徹底的.
不可能再和他有一種.
好像我們以往未信主之前那樣.
我們應該不斷將自己.
看自己是一個新的生命的人.
所以第二十五節.
藉著我們的主耶穌基督就能夠了.
能夠救我們脫離取死的身體.
就是說原來罪的性情.
我們知道它存在.
但是我們用另外一個律.
來幫助我們行神的旨意.
那就是我們過正常生活的一個基督徒了.
所以這樣看來.
一方面我自己在理智上.
是服侍神的律.
但另一方面.
肉體的律.
是不斷我們裡面的舊人.
就會服侍罪的律.
這個矛盾就不是代表著.
一個分裂的人格.
不是這個意思.
而是保羅清楚.
何時他傾向於這個肉體.
他就會受罪的律的試探.
引誘,操控.
所以基督徒能夠每個星期敬拜神.

$^{521}$有鼎子們彼此相交扶持.
這是其中很大的動力.
使我們能夠不斷對自己身份的奉獻.
當然去到第八章.
就明顯是一個很重要的生命轉化.
就是真的進入在第五到第八章.
正常基督徒生活的最高峰.
其實不應該叫最高峰.
就是平白地告訴我們.
我們是神的兒女.
這個身份不是用錢.
不是用我們自己的功德去換取回來的.
好像我們在說福音.
其實就是告訴我們.
本來我們在神眼中是敵對.
是一無是處.
而且我們擁有的只是取死的身體.
但是聖靈進入我們的內心.
居然不斷地去將我們聖化改變我們.
使我們能夠.
特別是八章28字那段經文.
很快要讀一讀.
沒有辦法講到那裡.
八章28一直到後面的段落.
保羅告訴我們.
原來神已經在我們的生命裡面.
預定了我們要由信主進入救恩.
去到一個能夠得到神榮耀的身份.
那個兒女的身份.
我下次要帶一個夾子.
帶一個衣夾來.
因為冷氣吹到聖經不斷地動.
阻礙了.
八章的28節.
希望大家喜歡這段經文.
看到這段經文的重要.
因為保羅在上文.
我們遲些再講.
因為眾生在墮落裡面.
是嘆息勞苦.

$^{561}$但是當救恩神的救恩成就.
基督就帶給這個新的人類.
有一個新的重要的身份.
所以留意第28節開始.
我們曉得萬事都互相效力.
叫愛神的人得益處.
就是按他旨意被召的人.
開始講我們能夠在萬民當中被神選召.
29節.
因為他如先所知道的人.
就預先定下效法他兒子的模樣.
使他兒子在許多弟兄中作長子.
就是說再一次強化我們神兒女的身份.
而且這個身份.
必須要和主耶穌這位長子的模樣.
好像一個倒模的模樣.
這樣匹配.
這樣的相像.
這就是說神期望我們每一個基督徒.
都應該得榮耀進入天家.
這裡告訴我們第30節.
預先定下的人又召他們.
就是信耶穌.
所召來的人又稱他們為兒.
稱兒就是不再追究了.
沒有罪了.
以及開始過一個誠信的生活.
這個就是稱兒的套餐.
所稱為兒的人又叫他們得榮耀.
這起碼就是第五章第十一節那段經文的.
另一個角度的解釋.
就是由呼召到得榮耀.
完全是只是因為聖靈在我們內心的參與.
我們人生的工程.
所以不是有新發現.
是每一個信耶穌得到聖靈的基督徒.
讓聖靈在我們內心不斷的教導引領我們.
所以這裡不斷告訴我們.
神賜下的聖靈就在我們裡面引導我們.
見證我們能夠是神的兒女.

$^{601}$就是說讓我們有一個很實際的.
很實際的見證.
我們走出來讓人見到.
我們真的等候神國度見證神存在的一個基督徒.
我們會為我們代求.
他的代求令到我們發現.
原來我們有一個很寶貴的身份.
使命就是我們可以成為一個為萬民祈禱的人.
很寶貴的.
你能夠運籌為額.
只是在家裡知天下事.
你就可以為不同的領袖去禱告.
譬如為普京,為特朗普.
有時為一些重大的事情我們去守望禱告.
就在我們的智慧,我們的見識.
我們的心靈裡面成為一個代求者.
所以原來我們不斷在人生最大最寶貴的.
就是進入神國度的那個稿庫.
來祈禱,守望.
更加重要的是使我們活出.
我們是神兒女的身份.
其實第五到第八章.
不斷用不同的角度去講這個身份.
我們是神的子女這個身份.
而且是要和主耶穌一樣的榮耀的身份.
你可以想像得到我們得到的救恩是多麼的寶貴.
原來每個基督徒都應該得到這樣的身份.
當然在事奉上,恩賜的運用上.
就真的有不同的角色.
有不同的神安排的位置.
那就真的是神的旨意和祂的主權.
但是同一個生命.
一個正常的基督徒就是一個不犯罪,愛神,有自由.
來越來越感受到聖靈帶領的那種寶貴.
這個就是第八章.
在第五,第六章的累積裡面.
第七章到第八章.
就是真的講到底是怎樣運作的呢?.
所以第八章很自然地就是整章開始講聖靈.
在一個基督徒裡面的運作.

$^{641}$羅馬書出現最多的聖靈就是第八章.
第五章只是第六節出現那一次.
賜給我們的聖靈.
將神的愛囂觀在我們的生命裡面.
使我們能夠在患難當中仍然可以堅持.
在患難裡面仍然歡歡喜喜地等待.
因為有聖靈囂觀神的愛.
聖靈就是將神的愛囂觀.
讓我們不斷有盼望.
不斷活著我們像耶穌.
這就是聖靈的工作.
所以我們就應該開始看第八章.
我們要先把經文補讀.
很顯然第一節就是回應第七章24,25節.
保羅很圖像化地對聚聖的描繪.
如果我們真的沒有看錯.
那個必死的人必死的身體.
就是一個死人我們綁住了在囚犯的身體上.
和他共存亡.
我們都想像得到.
聚聖和我們的人生都是共存亡的.
由出世到死.
每個人無論他是政界的領袖.
或者是普通的凡夫走卒.
聚聖都在每個人的生命裡面.
伴隨一生的了.
所以如果真的這樣伴隨.
保羅就認為如果不趕快信主.
接受聖靈其實就沒得救了.
我們現在從基督徒的角度去看.
我們是有聖靈的.
我們是可以在聖靈裡面自由的.
我們怎樣去經驗到這種基督徒身份的榮耀呢?.
首先就是第一個.
一至四節.
我先讀一讀.
這個餘金就不是一個時間的餘金.
在保羅裡面就是一個神學性的宣告.
就是說本來我們是聚聖.
綑綁著我們必死的身體.

$^{681}$是真是苦的.
不能夠不犯罪的苦.
但是現在在救恩.
耶穌來了這個世界.
成就救恩的餘金.
所以保羅這樣說.
就真是一個神學性的邏輯.
將舊有的生命和新生命.
做一個很強烈的對比.
所以保羅就說餘金.
所以我們又要問.
我們有沒有我們自己的餘金呢?.
我們是不是真的很寶貴.
救恩帶給我們的新生命呢?.
我們信了主.
我們是不是很珍貴.
我們能夠親近神.
敬拜神.
還有祈禱的.
不是責任.
而是權柄呢?.
所以這個餘金是很寶貴的一個連接詞.
餘金那些在基督耶穌裡的就不定罪了.
因為似生命聖靈的律.
在基督耶穌裡釋放了我.
使我脫離罪和死的律了.
因為罪和死就是孿生兄弟了.
我們得罪神.
神的判定.
神的審判刑罰就是我們會死.
當然有兩種.
真是肉身的死亡.
也有我們與神隔絕的那種失去自由的疏離.
這也是一種靈性上的狀態.
用我們能夠明白的很沒mood.
很不想回教會.
做基督徒有什麼好.
諸如此類的感覺沉悶.
感覺沒有生趣的基督徒的生活.
我們以為這樣是正常.

$^{721}$其實這樣是不正常的.
真正的正常是很開心很活潑.
是很有聖靈在我們內心裡面.
讓我們真的看到.
意識到原來我真的可以不犯罪.
勝過我的自己.
勝過我的自我.
這就是聖靈在我們裡面的工作.
所以這裡告訴我們.
原來在耶穌裡.
聖靈的工作表面上好像很隱藏.
我們不會感覺到聖靈.
令我們感覺祂的存在.
而是讓我們感覺我們是在基督裡.
所以這是一個很重要的代名.
代功一種這樣的代理.
聖靈是隱藏在我們的內心.
好像你不知道祂存在這樣.
祂讓你不斷看到耶穌愛你.
耶穌是真實.
耶穌是有能力.
耶穌是在我們內心.
和我們一起.
以致我們在基督裡面的人.
其實就是在聖靈裡面的人.
這裡告訴我們.
使我們脫離罪和死的律.
這個脫離不是拔罪根.
是使我們有一個新的傾向.
就是我們行神旨意那種暢快.
那個準備是準備就緒的心情.
基督徒準備就緒.
呈獻自己給神.
是神掌管我們的人生.
雖然講得這麼戲劇性.
沒有具體的例子.
要你向東你就向西.
是不是真的這麼戲劇性.
有時這個就不是違反我們自己天然的.
神給我們的智慧.

$^{761}$我們的知識.
我們每天的生活的形態.
但有時有一些的感動.
有時你很難想像.
總是有一些這樣的安排等等.
當然我們不應該說每一樣東西都是巧合.
也不需要訴諸於一種.
好像甚麼都是聽裡面的聲音.
也不一定要這樣的.
往往有時我們去順著我們自己每天的決策.
每天能夠分析的事情.
去做一些我們的安排.
但不等於這些安排神是不會用的.
但偶然間真的會有聖靈引導我們的心.
叫我們做某件事.
例如我之前也說過.
其實是李思敬牧師他自己親口說.
他有一次上台.
說到其實已經預備好講章的了.
準備五分鐘就上台了.
已經主席和在領唱的了.
心裡面真的突然間有感動.
你不要說這篇信息.
已經預備好了.
怎麼辦呢?.
你說另外一篇的信息.
他掙扎了五分鐘就上台了.
就說神要他講的那篇信息.
沒有講章.
因為是臨時五分鐘叫他講的.
完全不知道為甚麼要這樣講.
講這篇信息.
作為牧者我們當然平時都會有不同的例子.
在心靈裡面.
或者那些講章.
那些講道的上下文.
我們本身都會認識.
他講出來.
他說講完之後都沒有東西.
那就完了.

$^{801}$但是有人就將這篇道給中大的院長.
什麼醫生.
糟了 我忘記了.
SARS的時候.
我突然間忘記了名字了.
拿這篇講道給他聽.
聽完之後.
他就清楚.
因為他自己當時都是有一個等候.
中大要他做院長.
他都是一個等候.
聽了李斯敬牧師這篇道之後.
他就清楚了.
神要他接任中大的院長.
甚麼姚啊?.
沈祖堯.
沈祖堯教授.
沈祖堯院長.
很特別的.
聖靈感動李斯敬牧師.
要他講那篇信息.
不要講其他的一篇.
好像我沒有記錯.
就是若書雅記那篇信息.
要他講.
講完之後.
有人就將這篇道給了沈祖堯教授聽.
聽完之後.
他說他清楚.
他當時都是有一個等候.
到底我應不應該接任這個院長的職位呢?.
他聽完之後.
清楚神要他做這個決定.
所以你可以想像到.
聖靈可以這樣運籌為額.
所以正常基督徒.
就是應該不斷有聖靈.
一路引導我們做一些.
祂要我們做的事情.
所以其實一個基督徒.

$^{841}$一個正常的基督徒的生活.
就是一個準備好去服侍.
去聆聽神帶領的基督徒.
是要有一個準備.
是要有一個每天都奉獻自己的一個基督徒的人生.
那就是一個正常的生活.
繼續下去.
第三節.
律法既因肉體軟弱.
有所不能行的.
這是第七章的描述.
的一個總結.
神就差遣自己的兒子.
成為罪身的形狀.
作了贖罪祭.
就是主耶穌成為一個無罪.
替我們成為罪的一個死囚.
使我們能夠脫離罪的刑罰.
就是神的咒詛.
也使我們能夠得著內著的聖靈在我們的內心.
神就差遣自己的兒子成為罪身的形狀.
真是有一個肉體的主耶穌.
真是在地上行走過的三十多年的人.
祂不是一個靈體.
不是一個鬼魂.
因為初期教會有這些異端.
就說耶穌沒有身體的.
祂是一個靈體.
其實不是的.
這些不是還影說.
耶穌是真是有血有肉的一個人.
成為罪身的形狀.
作了贖罪祭.
很明顯.
在肉體中定了罪案.
就是說一個本來不可能不犯罪的.
當的罪的後人.
怎麼可能這個人可以不犯罪呢?.
耶穌就是這樣.
祂就定了罪案.

$^{881}$就是說祂定罪案.
在肉體中.
在祂的肉身.
祂能夠可以行神的旨意.
行到一個地步.
神說我要你去釘十字架.
耶穌都願意順服這樣去做.
當然我們沒有一個人.
一信耶穌就馬上像耶穌那樣.
有的可能是精神上有一點問題.
因為很少人可以馬上這樣做.
都有很多的掙扎.
包括保羅自己也是.
但是我們已經有聖靈.
所以羅馬書五到第八章.
完全不是奧秘.
是每個基督徒都有聖靈.
聖靈就可以帶領我們進入榮耀的人生.
正常基督徒的生活.
可以靠著聖靈不犯罪.
可以喜歡神的旨意.
因為我們的救主.
他已經以他的肉身去定了罪的罪案.
就是說罪的引誘.
罪的罪性.
不可以令到耶穌就範.
所以耶穌能夠完成神的救恩.
所以這個救恩就是我們每一個.
無論是一個小朋友.
或者是一個很認真信主的成年人.
當然小朋友也認真的.
都是這樣得到一個整位的聖靈在我們內心.
每一個都是得到整位的聖靈住在我們內心.
所以無所謂高級一點,低級一點的聖徒.
只有我們願意讓我們自己.
不斷呈獻自己給神.
所以第六章的重點就是.
奉獻自己給神.
所以第四節就說.
使律法的義成就在我們這不隨從肉體.

$^{921}$只隨從聖靈的人身上.
這個就是一到第四節.
不定罪的自由.
就是我們要將自己隨從聖靈.
即是說奉獻自己讓聖靈帶領我們的人生.
接著就繼續解釋.
在我們內心裡面他怎樣引導我們呢?.
本來就做一個對比.
因為隨從和體貼肉體的人就體貼肉體的事.
隨從聖靈的人就體貼聖靈的事.
這個其實都是第七章倒轉的描述.
第七章就是說我們很想行神道的時候.
就有罪在我們裡面來引誘我們.
試探我們甚至殺了我們.
但是這裡告訴我們.
現在就從另一個角度.
怎樣讓恩典在我們內心裡面不斷滋生呢?.
這就是告訴我們原來我們現在不需要再隨從肉體.
我們現在可以隨從聖靈了.
第六節體貼肉體就是死.
體貼聖靈就是生命平安.
原來體貼肉體是與神為仇.
因為不服神的律法.
也是不能服.
而且屬肉體的人不能得神的喜歡.
這是很詳細,很細緻地描繪一個人.
如果他繼續走自己的路.
以自己為中心.
不以聖靈,不以神為中心.
不以神的旨意為他的喜悅,為他的中心.
那就只會繼續讓舊有的.
其實從來都在的.
舊有的罪性.
就會繼續操縱我們這個心靈.
所以一個基督徒其實就是心靈裡面.
我們揀選主.
或者揀選自己的心意.
第七節就是原來體貼肉體就是與神為仇.
因為不服神的律法.
也是不能服.

$^{961}$而且屬肉體的人不能得神的喜歡.
第九節如果神的靈住在你們心裡.
你們就不屬肉體乃屬聖靈了.
清楚地宣告我們已經屬於聖靈.
因為由接受救恩那一刻.
就是成為屬於聖靈的人.
不是靈因運動的.
要第一波,第二波,第三波.
要講方言.
要能夠行神蹟.
要有一種特殊的.
好像第二波的進入我們的心.
令到我們的罪根拔掉.
這些是很容易的.
都很想大家追求一個很聖潔的生活.
但是仍然都是那個問題.
不是迫自己過聖潔生活.
而是要清楚我們已經是屬於聖靈.
不再屬於肉體.
要向自己宣告我們是屬於聖靈.
屬於聖靈就需要隨從聖靈的生活.
這是第十節的中間.
第九節.
「人若沒有基督的靈,就不是屬基督的,.
基督若在你們心裡」.
從另一個角度講聖靈的工作.
其實聖靈的內著.
就是等於基督的愛,基督的救恩.
基督自己的內著.
基督如果在你們的心裡.
身體就因罪而死.
心靈卻因義而活.
這裡我們要少少解釋.
上文就是說.
我們有神的靈住在我們的心.
我們就不再歸屬於肉體了.
就是說我們要在我們的內心.
有一個好像要宣告獨立的.
好像如果我們是殖民地的一個地方.
現在那個殖民地的政府.

$^{1001}$因為各種原因.
最終他們要撤退了.
下半期他們的國家要交出主權了.
那個殖民地的政府.
新政府成立了.
有自己的新的國家的旗幟了.
好像回歸等等.
你可以想像得到.
因為我們香港都是回歸給祖國.
我們就感受過97年的.
應該是6月30日和7月1日的轉折.
就是真真正正宣示獨立了.
所以保羅在這裡告訴我們.
我們要有這一刻宣示我們的獨立.
神的靈住在我們的心裡.
這個是讚身的.
從來我們未信耶穌.
都沒有神的靈在我們的心靈裡面的.
但是由信主那一刻.
我們就知道原來我們不再屬於肉體了.
我們是屬神的了.
我們是屬基督的了.
其實是同義詞.
保羅是交替的用.
當我們這樣屬基督的時候.
我們就身體因罪而死.
這裡稍為解釋.
這個身體怎樣因罪而死呢?.
信了耶穌就因罪而死.
那我們就死了.
這個意思相信不是這個意思.
上下文都是在說舊有的肉體的生命.
我們要理解.
好像《六章》這樣說.
我們藉著基督的犧牲.
其實我們向罪已經是死了.
我們現在是有一個新的生命.
我們是向神活的生命.
不過我們這個肉身的身體.
仍然還沒有到一個時候.

$^{1041}$神要它結束在地上的生活.
在地上的時間.
所以就不是這個肉身的身體因罪而死.
是我們裡面的那個舊人.
是因罪而死了.
基督已經將我們拯救了.
所以我們要宣示這個主權的獨立.
我們是屬神了.
所以這是一個很重要的開始.
對基督徒來說.
我們不是要給人家一些枷鎖.
說你上初信班,浸禮班.
你浸禮班之後你繼續上裝備課程.
成長課程.
好像繼續給人有一種壓迫感.
哇 原來上教會有這麼多壓迫感.
其實是希望讓人不斷對你自己的主權.
屬神的這個主權.
越來越掌握清楚裡面的真理.
所以好像我們讓人感覺到.
行了 我浸完我就夠了.
我不用再成長了 不需要裝備了.
我只是來崇拜就行了.
這是一個錯誤的想法和一個教導.
所以應該這麼說.
我們要讓弟兄姊妹去到一個階段.
原來舊恩帶給我們的生命.
舊有的生命已經好像隨隨的主權已經失去了.
新的主權已經慢慢地升起了.
就是基督的主權在我們裡面.
就是我們新的主人成為了我們的主宰.
這樣的說法就很不同了.
我們就要制定新的律法了.
我們都有我們自己的基本法.
有我們自己的香港或者一些獨立的國家.
有他們自己國家的律法 很正常的.
所以這裡告訴我們原來我們不斷聽多些真理.
聽多些人講聖經的真理.
我們就更加認識神的計劃.
那就會更加豐富我們這個主權的穩定.

$^{1081}$這裡就告訴我們.
身體因罪而死 心靈因義而活.
這就是這個意思了.
我們的心靈越是體現到這個主權的真實 偉大 自由.
我們的內心就會奔放了.
我們的內心就因義而活.
因為這個主權是一個公義的主權.
是舊恩耶穌換取回來的主權.
就已經能夠很理直氣壯地.
不需要再犯罪了.
因為耶穌的新生命在我裡面.
仍然透過第八章講的.
聖靈的律在我們的內心.
所以基督徒是活的.
基督徒的心靈是永日公義和活潑的.
這就是告訴我們.
又是真的不斷 好像很誇張.
我們每天將我們自己體現一下.
基督國度的那個旗幟.
舊恩的元帥主耶穌.
他已經是那位耶和華升旗.
他來掌管帶領著我們.
任避著我們.
我們就在這個新的國度裡面.
我們就是他的子民了.
這就是基督徒的生活.
第十一節 然而叫耶穌從死裡復活者的靈.
若住在你們心裡.
那叫基督耶穌從死裡復活的.
也必藉著住在你們心裡的聖靈.
使你們必死的身體又再活過來.
再一次加重那個語氣.
耶穌的復活是因為神的能力.
神的聖靈進入耶穌裡面.
這種是我們不能夠去分析.
不能夠用科學去解釋的能力.
不是物質界的能力.
是屬靈界的能力.
(聽不清).
開電掣 沒電.

$^{1121}$稍等一下 現在沒電了.
是不是這裡鬆了掣呢?.
(聽不清).
應該可以的.
電路沒電.
不知道為什麼現在這個掣不成功.
為什麼不開上去呢?.
(聽不清).
可能是線鬆了.
可能這個掣鬆了.
這個沒電了.
這個線呢?.
(聽不清).
這個掣可以了嗎?.
(聽不清).
原來有些乾坤.
繼續吧.
保羅再一次告訴我們.
仍然是說基督復活的大能在我們裡面.
所以其實基督復活的大能.
其實就是聖靈一樣的能力在我們裡面.
叫耶穌復活的靈在耶穌裡面.
所以你可以想像.
死了能夠復活.
能夠是一種新的復活.
不是拉薩路那種剛剛死了.
翻生的仍然是肉體.
地球肉身的一種復活.
已經是不得了的了.
但現在耶穌是一個站生到.
地球上從來沒有的.
以後就會有了.
因為當耶穌再回來.
這種能力就會在所有的教徒裡面發生了.
但現在保羅想告訴我們.
他體現到的就是這個奧秘.
這個神的能力.
在耶穌身上發生的.
令到他永活的.
這一種的能力.

$^{1161}$不是甚麼鬼怪神秘主義.
而是很實實在在.
眾人都能夠看到耶穌的身體.
耶穌可以隨時出現在那四十天裡面.
任何地方的.
這樣的奇妙.
是人類的社會沒有的.
不再受時空限制.
不再受地理物質的限制.
這樣的能力在耶穌身上.
就已經很自然地體現了.
現在再一次說.
然而叫耶穌從死裡復活者的靈.
躍駐在你們心裡.
這個躍當然不是假設.
是想告訴我們.
我們要這樣去經驗.
因為上文已經很清楚.
生命,聖靈的律.
已經藉著我們在耶穌裡面.
已經進駐在我們的內心了.
所以再一次告訴我們.
如果你體驗到這位這麼大的能力的聖靈.
能夠曾經在耶穌的死亡裡面.
進駐第三天復活.
能夠繼續活著.
能夠在他們眾門徒面前升天.
慢慢地駕雲升上去.
這種能力.
不是拍那些用吊威也的一種掩眼法.
而是實實在在我們見到一個超大的神蹟.
保羅現在告訴我們.
就是這位聖靈.
就是這位使耶穌復活的聖靈.
就在我們裡面.
所以祂既然進入我們的內心.
使我們能夠在這個復活裡面.
我們必死的身體.
必定又會再活過來.
必定再活過來.

$^{1201}$可以告訴我們.
如果在這裡上下文處理著的.
就是不犯罪的這個自由.
在這裡上下文明顯告訴我們.
就不是我們死了的人.
我們去到現在.
哥倫神閣三樓四樓.
都有很多我們的弟兄姊妹離開世界.
安息主懷的.
有時走過這個這個.
記得.
你沒理由突然去到叫一下他的名字.
你現在復活.
因為未到那個時間就不應該這樣復活.
所以這裡的意思.
不是說那個身體的復活.
是說我們不犯罪的那個自由.
在我們認定聖靈已經在我們裡面.
就有這種能力去帶領著我們的人生.
所以第十二節就真是這個意思.
所以 Therefore.
這樣看來.
他這裡就不是說.
你要快些儲備那個能量.
夠了100\%了.
那你就可以復活了.
不是的 是叫我們清楚再一次認定.
其實我們沒有再欠肉體的債了.
就是說我們不需要再懼怕這個舊有的人.
對我們的那個無理的欺騙.
我們不需要再欠肉體的債.
去順從肉體活著.
這個意思我們現在本來就應該是新的族類.
是順從聖靈而活.
因為基督的力已經進駐我們的內心.
所以很明顯第十二節.
保羅再一次好像婦後婆心那樣說.
要清楚自己是有新的主權.
是獨立了 是榮耀的.
我們就不需要再被肉體支配我們了.

$^{1241}$我們就能夠順著聖靈而活了.
所以第十三節又再一次說了.
如果你們繼續順從肉體活著.
就必定要死.
就是那個靈性上與神的關係隔絕.
但是如果靠著聖靈.
致死身體的惡行.
致死就是說向他有一個宣告.
我獨立了 我有新的主權了.
我是歸屬了這個新的主人耶穌了.
我是屬於主耶穌.
我是不需要再讓舊有的性情去霸佔我的內心了.
所以我們靠著聖靈致死身體的惡行.
就必定會活著了.
因為凡被神的靈引導.
都是神的兒子.
表面上就好像很簡單.
但是更加一而再.
很想強調其實我們有沒有神的靈引導.
如果我們信主我們都是神的兒子.
為什麼再強調要有神的聖靈引導.
我們才是神的兒子呢?.
為什麼要這樣說呢?.
是不是沒有神的聖靈引導就不是神的兒子呢?.
相信也是和上下文有關.
就是更加體現我們不再在肉體的支配裡面.
我們有一個新的主人帶領我們.
這個主人他要在我們心靈裡面.
他不是用一種佔據,威嚇來威逼我們.
而是要引導我們.
所以這個引導是可圈可點的.
保羅用的字眼.
好像聖靈是很害羞的那樣.
你聽祂的話.
你讓祂引導你.
就是說你可以不讓祂引導.
你可以拒絕祂的引導.
這樣的時候就讓我們看到.
罪,舊有的性情是強迫我們.
但是聖靈是用愛去引導我們.

$^{1281}$就是說甘心樂意讓聖靈在我們內心帶領我們.
這個時候我們就可以看到一個很大的分別.
就是好像《馬太福音》第五章.
有些好像很矛盾的那些百福.
使人和睦的人有福了.
稱為神的兒子.
其實你不用人和睦.
你經常在教會裡面.
有些人經常去挑骨頭的.
他都是神的兒子.
他信耶穌就是成為神的兒女.
到底耶穌這樣說.
保羅又這樣說.
我不讓聖靈引導我是不是神的兒子呢?.
如果你要這樣爭辯的.
這種哲學的假設性的辯論.
其實就沒有意思了.
因為這裡說的是.
讓聖靈釋放自由的神的兒子.
一直都在說的是.
抵抗肉體的那個方法.
這裡一直都在說的是.
讓聖靈.
他好像很默默地引導我們.
而我們願意順從的時候.
這裡就寶貴了.
保羅說那個的結果就是.
你自己知道你是神的兒子.
其實慢慢就是讓人都知道你是神的兒子.
這樣的一種的順從.
所以我們不要用一種反過來的.
好像特意和保羅唱對台那樣.
我就是不讓人引導.
但我是神的兒子.
其實就不是這段經文的意思.
而是告訴我們.
更加體驗我們可以克性罪惡.
不再在肉體的引誘下的神的兒子.
就是我們讓神的聖靈引導我們.
這個引導就是一定朝向.

$^{1321}$和神的兒女身份匹配的人生的旨意.
那些的帶領.
所以你可以想像到.
一定朝向那個真理.
愛主,愛人.
傳福音.
一定朝向這樣的事件來生活.
所以保羅就告訴我們.
凡是被神的靈引導.
都是神的兒子.
就是做神喜悅的事情.
神的兒子.
那個心靈應該就是.
我們都有幾十年信主的經驗.
那個心靈不是你突然間越來越大.
而是那個對真理的深化,內化.
和你的心胸對神的國度.
應該就是越來越有那種的眼界.
越來越擴闊那種的眼光.
那個人越來越謙虛.
越來越很想為神做一些事情.
不需要問為什麼的報酬.
你都願意去做.
這樣就是神的兒子.
所以這個不是說我相信耶穌就是神的兒子.
為什麼我要被聖靈引導.
我才是神的兒子.
其實這樣的質問.
他已經是很離開.
和聖經的真理.
保羅這位上下文的教導.
還有一段距離.
因為根本就不是討論一個.
要怎樣才是神的兒子.
而是他想說出一個很寶貴的生命的流露.
就是我們讓聖靈引導.
不讓肉體引導.
不讓肉體來引誘.
來壓迫,威逼我們.
我們就是神的兒子了.

$^{1361}$很自然.
一定是說到阿爸父這個問題.
你們所受的就不再是奴隸的心了.
是害怕的奴隸的心.
所受的是兒子的心.
因此我們呼叫阿爸父.
這是一個共同的見證.
所以內心的引導.
聖靈的引導.
我們是神的兒子.
即是說我們已經脫離了罪.
脫離了肉體的引誘,威逼.
我們已經來到可以內心.
一直配合神在我們生命裡面的帶領.
那個引導,那個感動.
所以很自然地.
我們是很樂意稱呼造物主.
阿爸父神.
而神更加在天上可以說阿們.
我們真是祂的兒女.
這個是複合字.
亞蘭文的複合字.
阿爸父神.
阿爸真是亞蘭文.
可以想像得到是非常親切.
能夠這樣教我們.
在聖靈裡面.
我們不單純是一個很機械性的神的兒子.
而是親切到一個地步.
我們的內心深處是寶貴.
我們有這位愛我們的父神.
能夠這樣這麼親暱的阿爸父神.
所以不是我們廣東話的老爸.
不是這樣.
有時這樣是過分了.
以為這樣叫親暱.
其實不是這種的親暱.
不是說我們天上的爸爸.
也不需要這樣.
但這裡就真是在說生命的親暱.

$^{1401}$真是告訴我們.
我們和神近到一個地步.
我們因為在聖靈的引導下.
已經不再在肉體的支配.
而是讓聖靈的愛引導我們.
很容易就一定達至這個內心的感受.
似乎應該都可以感受得到.
保羅的內心.
他帶領外邦人稱呼做民主.
猶太人不可以妄稱神的名字.
這個是十誡裡面的第三誡.
你妄稱神的名字會死的.
耶和華.
你看猶太的文獻.
如果你讀神學.
很有趣的.
那個God的字.
GOD沒有一個O.
只有一個hyphen.
好像他們不敢讀神的名字.
不敢妄稱耶和華.
所以保羅引導外邦人,羅馬人.
今日香港人.
我們稱呼做民主.
是阿巴庫神.
是非常非常的親密.
所以我們其實每天都在享用這種的親密.
享用這種跟他的關係.
所以這個就是聖靈在我們裡面的引導.
所以第十六誡.
聖靈與我們的心同正.
共同的見證.
我們是神的兒女.
既是兒女便是後子.
後子即是可以繼承產業.
這個意思.
就是神的後子.
繼承神將來給我們的產業.
而且是與基督一同繼承神的產業.
如果我們和他同受苦.

$^{1441}$也必和他一同得榮耀.
在這個段落裡面就停頓了.
然後就不再說我們的內心.
是說整個受造物的.
由亞當開始說原初的創造.
可以想像到.
一到第十七誡.
是告訴我們聖靈引進我們.
要我們看到這份能力.
是叫基督復活的能力.
在我們的內心.
這個能力不是讓我們行雲駕霧.
不是讓我們做一些超自然的事.
要行神蹟.
要講方言.
要怎樣怎樣.
不是.
是要我們能夠意識到.
我們裡面的主權已經移交給基督.
是一個愛的移交.
我們要讓這個主權進駐在我們內心.
是一份愛的進駐.
聖靈引導我們.
引導我們看到.
神一個很大的人類的功情.
就是使到本來罪人.
能夠在救恩裡面成為神的兒女.
而且要享受將來得到產業的這個位份.
這個名份.
和這個福氣.
所以救恩帶給我們.
不是我現在只是禮拜日回教會.
我是一個基督徒.
我要在人面前謝罰.
保羅很想帶我們看到.
真正的那個時間軸.
由信主到得榮耀.
是神的後製.
要看到我們能夠有將來永恆的產業.
我相信在香港.

$^{1481}$很多基督徒如果他越來越疏濫回教會.
你問他信不信有天堂.
信不信他死後會去到永恆.
可能他都要想一下.
因為和他的生活太遙遠.
距離太遙遠.
如果他的生活沒有什麼病痛.
沒有什麼困難.
工作又很順暢.
做的生意又很穩定.
可能他根本不會想永恆.
但是在保羅的《羅馬書》.
他們是在面對苦難的.
尼祿王的逼迫其實越來越白熱化.
但是保羅在這裡告訴我們.
他們可能很多都是奴隸.
沒有什麼社會的地位.
再加上猶太教另外一種的夾擊.
以為你們違背了摩西的律法.
所以信又不是.
不信又不是.
這樣的羅馬外邦的政權的迫迫.
猶太教的誤解.
針對冷嘲熱諷.
對初期教會的信徒來說.
其實是一個很大的張力.
這個張力保羅怎樣去理解呢.
不要緊的 回教會就什麼都沒事了.
信了主就行了.
保羅都不是這樣.
一些這樣的安慰性.
好像完全講完等於沒有講.
保羅正是人性裡面.
導致我們離開神的那個實體.
就是罪在這個世界.
因為阿當的背影.
罪在地球上.
你用任何的方法.
都不能夠好像聖經這個方法最好.
因為唯有神的聖靈.

$^{1521}$在我們每一個基督徒裡面進駐.
我們有新的生命.
新的智慧眼光.
這樣去看自己.
所以基督徵兵就是這樣形成.
我們每一個都是成為聖傑的基督徒.
每一個都是由救恩帶給我們進入.
和聖靈的相交.
聖靈的引導.
由意識我們是神的兒女.
這個名份,這個身份.
我們在地上我們是有新的身份.
就是神的兒女.
神的後製的兒女.
的一個身份的時候.
我們就能夠.
也是保羅一直都講的.
就是我們要接觸那個救恩.
就是要將福音傳到地極.
宣教.
去到西班牙.
去到其他將來的國家.
因為航海的技術一直的發展.
基督的救恩一直的傳開.
最終是去到中國,印度,非洲,大陸.
你可以想像得到.
這個就是保羅的意思.
所以保羅給我們人類.
是有新的創造.
就是在基督裡面的新創造.
在這個新創造裡面.
就是聖靈扮演著主導性的角色.
所以我們要記得.
今晚我們要認定.
我們是有一個新的主權.
我們移交了給基督.
聖靈已經藉著基督進駐.
而且是愛的進駐.
不是強迫的威嚇.
是愛的進駐在我們裡面.

$^{1561}$它會扮演著那個引導.
一直這樣來感動.
將神的愛囂觀在我們的內心.
所以透過聖靈.
我們就能夠越來越認識.
我們是神的兒女.
所以聖靈和我們的心同樣.
這樣來呼叫阿爸父神.
對神這位創造的主.
我們有這份新的關係.
從來都沒有.
但是因為基督的緣故.
我們有這份新的關係.
所以有多少弟兄姊妹.
從來都沒有這樣稱呼神呢?.
今晚你睡覺多叫幾下吧.
真是因為基督的緣故.
我們稱呼阿爸父神.
是很親暱的.
不需要尷尬.
不需要起雞皮疙瘩.
所以不是那些俗話的老爸.
不是用這些這樣的.
而是整個神學的概念.
阿爸父神.
這個概念是聖靈和我們的內心.
共同的印證.
我們是神的兒女的身份.
所以後面又再補充那個榮耀.
所以這裡告訴我們.
我們雖然一同受苦.
在地上.
剛才也說了.
猶太教對新興的基督教的誤解.
保羅就是那個令這件事惡化的人.
因為他是猶太教的拉比.
他帶領外邦人.
不需要守律法.
來信耶穌就行.
所以得到的那種誤解.

$^{1601}$那種攻擊是難以想像的.
但是你又帶領人一個新宗教.
舊的羅馬的宗教.
希臘的宗教全部不要了.
因為全部是假的.
那就更加激怒政權.
尼祿王.
要來迫害這個基督教.
保羅也是那個始作俑者.
但是保羅在這裡給我們看到的就是.
人類一個很大的啟發.
就是正常的基督徒.
是可以犯罪.
可以在主聖靈的帶領下.
看到自己是神兒女的身份.
今晚就講到第十七節這裡.
因為下一個段落就是創造的墮落.
帶來眾生都要嘆息勞苦.
再一次帶回來看.
人類的救贖原來都是唯有藉著基督.
整個受造物的盼望.
都是等待基督的再來的.
所以盼望我們能夠透過這些經文.
不斷告訴我們.
真正的基督徒是不犯罪的喜樂.
不犯罪的自由的基督徒.
才是正常的基督徒生活.
我們一起有一個簡單的禱告.
我們的阿爸父神.
我們很少這樣重複地稱呼你.
多謝你藉著保羅的教導.
聖經的啟示.
我們無論是甚麼民族.
我們都可以稱呼我們的造物主宰.
我們這位耶和華真神.
你是我們的阿爸父神.
多謝主.
你是創造的主.
其實整個宇宙萬物.
我們根本不知道有多浩大.

$^{1641}$那種偉大,那種浩瀚,那種奇妙.
但是造物主就是這樣.
拆遣你的愛子.
為我們人類釘十字架.
而且是極痛苦的死亡.
換取的我們與主相和.
我們享用在救恩裡面.
帶給我們的恩典.
帶給我們的新的身份.
幫助我們今晚聽到這一堂的訊息.
就是我們新的主權.
已經移交給基督.
我們是屬於基督的人.
幫助我們能夠真的看到.
這是多麼寶貴的真理.
聽我們的感恩禱告.
奉耶穌基督的名.
阿們.
\newpage



\section{羅馬書}
\label{sec:mDmug5anuis}
\textbf{2025-07-31 周四晚查經班《羅馬書(18):分題:聖靈帶動萬物復興》 講員:曾裔貴牧師}
\newline
\newline
連結: \href{https://youtube.com/watch?v=mDmug5anuis}{\texttt{https://youtube.com/watch?v=mDmug5anuis}} ~~~~ 語音日期: 2025-08-02
\newline
\newline
\hyperref[sec:HoQGEbdeqSk]{\small{< < < PREV SERMON < < <}}
~
\hyperref[sec:index_chronic]{\small{[返順時目]}}
~
\hyperref[sec:index_scriptual]{\small{[返順卷目]}}
~
\hyperref[sec:HwCOo8EiKp0]{\small{> > > NEXT SERMON > > >}}
\newline
\newline
$^{1}$好 弟兄姊妹我們一同在晚上.
查經班再去禱告.
多謝天父讓我們能夠在香港有不同的渠道.
去學習聖經的真理.
明天開始就會有港九的培靈年經會.
我們很盼望神去感動不同年代的弟兄姊妹.
尤其是我們更加盼望有新一代的年輕人.
他們能夠學習聖經的真理.
不要覺得是沉悶 不要覺得是湧場.
反而是每一堂都有主聖靈的感動.
讓香港的教會能夠有復興.
我們一起去明白真理.
還有其實香港的信徒都不是真的佔大多數的.
在整個的社會裡面.
其實福音來到我們中間已經是好幾個年代了.
所以都盼望香港既然我們有這麼多的資源.
讓我們能夠真的看到我們能夠有更多人.
如果蒙福得到神的救恩.
還有起來去裝備來去被主聖靈去感動.
還有真的去傳福音給我們身邊的人.
使更多人得到神的恩典.
我們將十天的培靈會仰望在你的手中.
也都讓我們能夠有系統地慢慢去讀羅馬書.
來到第八章了 雖然我們不是用原文識經.
也都沒有這樣的能力去駕馭這些經文.
但是透過很準確的翻譯.
都給我們已經差不多是九成我們掌握到.
仕途保羅他一直的神學.
還有他在整個的上下文的討論裡面.
帶進的是一個很重要的神學觀.
就是我們的日論 我們的內心.
還有神怎樣藉著耶穌基督去救贖.
還有聖靈怎樣在我們內心進駐.
成為我們的保衛師 一直這樣帶領著教會.
我們今晚去到第八章的重要的中心思想.
真是讓我們看到整章聖經講聖靈是相當透徹的.
比起其他的書卷來得更加透徹.
還有來得更加自然 就是聖靈的律.
已經在我們每一個信徒的內心.
所以我們盼望主讓我們能夠慢慢去分段.

$^{41}$逐段去看的時候 我們又找到一些.
我們可能過去忽略了的那些提醒和教導.
我們仰望等候在你的面前.
我們這樣祈禱 奉耶穌基督的名.
阿們.
開始我們一些上一課我們講過.
第八章肯定是和第六,第七章的關聯性.
第六章是很想告訴信徒.
在救恩的這麼浩瀚的恩典裡面.
不應該甚至是禁止犯罪.
不可犯罪這個主題是相當明顯的.
但是當第六章你再看得深入一點.
其實是要帶進我們有新的生命.
因為第六章第三節已經講了.
我們藉著和主耶穌的聯合.
我們就已經歸入死.
一同埋葬 跟著有一個新生命的復活的樣式.
其實開始暗示.
當然到了第八章就很詳細講.
這個新生命的運作.
所以第六章那個結論.
就不是一個很負面的不可犯罪.
而是我們如果能夠一直在基督徒的生命的成長.
成熟裡面.
一個正常基督徒.
是應該說可以再犯罪了.
這個可以不犯罪和不可犯罪.
是兩個層次來的.
是兩樣東西.
一個小朋友看到一顆糖.
他不斷想吃.
但是到長大之後他明白.
他可以吃.
有需要的時候.
吃一顆糖.
譬如行山.
可以增加一些力量.
和他貪吃是兩樣東西來的.
所以應該是基督徒的生命成熟.
我們不再是這件事可以做.

$^{81}$這件事不准做.
這件事我偷偷地做.
應該就是一個成熟的生命.
第七章就再一次告訴我們.
好像一個屬靈的儀器.
描繪到底.
當然我們如果將那個討論.
在《神學家》裡面.
到底是保羅信主之前.
還是信主之後的掙扎.
我們很難有一個正確的答案.
因為保羅根本沒有告訴我們.
這個就是我未信之前的掙扎.
如何來律法.
透過罪,透過律法.
在我們裡面來引誘我們.
保羅沒有說.
而是在律法底下.
我們如果沒有做一個新的生命力量.
去引導,去幫助我們.
我們一定會面對那個犯罪的傾向.
以及那個罪透過律法來引誘我們.
第七章就我們上一次說過.
就是告訴我們.
保羅用了一個很重要的連接詞.
第七章的二十四節.
他說我這個人真是苦.
誰能夠救我脫離這個必死的身體呢?.
或者舊的和諧館叫做取死的身體呢?.
我搜尋過一些不同的解經家.
一些著色書.
都告訴我們.
這個取死的身體.
極有可能就是羅馬當時最殘酷的一個刑罰.
就是一個很大的罪犯.
就是那些死囚.
他們用很殘酷的刑罰.
就是將一個死人.
綁在這個罪犯的身體上.
就是說和他完全的聯合.

$^{121}$那個死人.
慢慢那些腐臭的蟲等等.
就會咬那個活的人.
這樣你可以想像到那種死法.
比那個死人更加慘.
那個死人可能被人一刀就死了.
但是那個活的人.
那個懲罰就是被綁在死人裡面.
你可以想像得到.
如果保羅是用這個當時.
大家能夠聽得明白的這樣的一個例子.
那就可以想像得到.
原來罪在我們舊有的人身上.
就是這樣的一種狀態.
讓我們根本上動彈不得.
還有在一種很大的痛苦裡面.
既不能夠行神的旨意.
沒有力量去行神的旨意.
但是同時間我們自己.
又因為罪帶來各種的後遺症.
人際關係的疏離.
或者自己犯罪之後的種種的刑罰等等.
你可以想像得到.
24節就應該是一個在罪惡裡面痛苦掙扎的人.
偶然間有時候都要去探監.
有些等待判監的.
早幾個月了 最近我都沒有去了.
就是你沒有辦法.
因為已經是當場抓到販毒.
那你可以想像得到.
一位的男士都跟我差不多年紀.
是一位肢體的親戚.
去赤柱去探望有什麼可以做呢?.
就跟他祈禱 希望他信主.
他當下都表示信的.
但是你感受到一直都在等待上庭判決等等.
因為已經初步自己都承認了罪.
因為當場抓到販毒.
你就幫人帶.
你可以想像得到那種罪帶來的後果.

$^{161}$那個刑罰 雖然香港沒有死刑.
但是你都可以想像得到.
他的人生就是這樣 基本上沒有的.
都不少這樣的情況.
這樣的人有時候在香港或者其它地方的監獄.
但是這裡說的是我們生命裡面的.
囚禁著我們 就是罪這個性情.
在羅馬書很想探討這件事.
是每一個人 生而為人的人.
都有這個罪的性情在我們裡面.
現在整個羅馬書就是告訴我們.
一個在神救恩裡面的新造的人.
就是用什麼新的救恩.
帶我們離開這個罪的死亡和罪的律呢?.
所以你看到七章二十四節.
和七章二十五節的對比是很強烈的.
我這個人是多麼的痛苦呢?.
二十五節就立即告訴我們.
就能夠脫離這個取死的身體.
接著保羅就做了一個這樣的總結.
就是說我現在有一個新的律去支配我.
去帶領我 去愛我 引導我.
我怎樣去信服祂 而達至那個自由呢?.
那個救恩帶給我的榮耀.
那個用我們講的是升級呢?.
我們的階級都可以提升呢?.
信耶穌真是的.
你真的會越來越斯文啊.
大方啊 禮拜天回教會啊.
真是每個人如果你真的順從聖靈呢.
真是會有這個提升的.
慢慢的原來子女讀書又變得厲害了.
在家裡都不再講髒話 不再賭馬.
不再吸煙喝酒 不再去花天酒地的時候.
多點時間照顧子女.
子女出生就一起帶回教會.
明顯你是整個階層的改變的.
不過這些就是那個救恩的副產品.
真正的那個救恩是.
保羅這裡所針對的是每個信主的人的內心.

$^{201}$有了一個新的動力 就是聖靈的律.
所以你看到24和25節那個強烈的對比.
就已經是一個很強烈的對比.
就是說保羅是很巧妙地帶進那個.
咒詛底下的人類舊有的亞當的族類.
是沒有任何的辦法可以解決的.
最終將來還要面對神的審判.
但是在耶穌這一位新的亞當的來臨.
拯救我們透過信進入這個救恩裡面.
就要不要只是得救 要看在我們的人生.
怎樣去讓聖靈去.
任何一節第八章都找不到支配這個字的.
你一路找到的就是 等一下我會再跟大家分析.
引導,貼近,隨從.
這些某程度是支配的.
不過這個支配是要那個人選擇去跟隨的.
同罪就真的支配了.
你不能不選擇的.
你都不到你有選擇的.
所以保羅用的字眼都是不斷告訴我們.
聖靈是不會禁止我們.
聖靈是不會強迫我們.
但是聖靈會感動我們.
讓我們越來越長大.
我們是很主動地.
所謂我們中國人說的生性.
長大了 生性了.
我弟弟的兒子就是這樣.
他小學五年級.
因為太太 就是媽媽離開世界.
他在天蘇圍的小學是很壞的.
沒有人理會.
爸爸就在日本工作.
三個月才回來一次.
媽媽就患重病.
他就很壞 很頑皮.
就是想引人注意.
他根本上就是突然間的突變.
媽媽離開世界之後就帶他去日本.
他一去就已經跟日本人打架了.

$^{241}$因為進屋又穿一雙鞋.
去廁所又換另一雙鞋.
然後每件事都要一起做的.
他完全不習慣.
但是我這次過去.
其實很多年沒有見他了.
已經是一個很斯文.
很有紀律的一個小朋友.
一個青年人.
又有教養 有禮貌.
當然一去到某些地方.
他跟爸爸很快就衝突了.
因為他的模仿是很少的.
因為只有他自己跟爸爸在日本.
還有他的後母媽媽.
所以成長真的是有一個改變.
但是這裡說的是.
我們信了主之後.
盡快要知道.
所以羅馬書為什麼是這麼重要的一本書.
就是救恩我們很多時候以為信了耶穌.
我們就會自動成長.
但其實不是的.
我們必須對保羅這個神學.
有一個充分全面的掌握.
這裡用第八章.
來將第六章,第七章的不准犯罪.
還有罪在我們裡面是不會離開的.
到死才會解除那個咒詛的.
所以在我們離開世界之前.
罪的律和聖靈的律.
我們是並存的.
所以第八章的重要就是.
要我們成為一個成長的基督徒.
就是要開始貼近選擇跟隨聖靈的感動引導.
我們就能夠不再在那個律法的咒詛.
和罪的引誘下.
所以保羅就告訴我們.
這樣看來.
一方面我自己在理智上服事神的律.

$^{281}$另一方面肉體卻服事罪的律.
就是說是完全並存的.
於是我們自己是否願意成為一個成熟的基督徒呢?.
每樣都能夠隨從聖靈在我們內心的帶領引導.
這就是七章的二十五節.
接著很快就第八章.
就是「餘金」.
我們讀讀這部經文.
「餘金」.
保羅一而再用了好幾次「餘金」.
就是說不是時間.
而是在救恩的發展裡面.
現在在基督耶穌裡面的人.
所以就不是現在這個時刻的時間.
而是每一個在耶穌的救恩裡面的人.
都必須要認識.
我們現在在主的救恩裡面.
第一個就是不再定罪的自由.
或者說一個釋放的恩典.
來在我們裡面的運作過程.
就是生命,聖靈的律.
第二節.
在基督耶穌裡面將我們釋放.
使我們脫離罪和死的律.
所以不是罪根拔掉了.
而是罪的性情仍然在我們每一個基督徒的內心.
但是聖靈.
這位是生命的聖靈.
在我們裡面.
藉著耶穌將我們釋放了.
所以相信這個就是那個氛圍了.
是那個全面對真理的認知和認同.
和完全在我們的生活.
投入在這樣的氛圍裡面.
我們就會感受到這種這樣的生命.
所以如果你去一些很有恩典.
有聖靈同在的聚會.
你會感受到一份平安.
一份內心的那份溫暖.
感受到弟兄姊妹之間.

$^{321}$就算沒有說話.
都有一份力量給予互相的支援.
這個就是聖靈臨在聚會裡面.
所以我們經常在這樣的氛圍.
或者和這樣的弟兄姊妹相交.
能夠大家都是為神的國度這樣服侍.
你就會感受到生命聖靈的律.
那種活潑,那種玲瓏浮特的同在.
相信大家我們感酸.
有時候有些大型的音樂會.
敬拜的聚會.
有時候大家很投入地去帶領那些聖詩.
你又加入去唱.
你是一心一意去敬拜神.
包括你自己.
包括那些回眾都感受到神的同在.
這個就是.
保羅就很想告訴我們.
就是因為聖靈的工作在信徒的內心.
所以這裡告訴我們.
這裡第一件事就是說.
第八章就肯定很想告訴我們.
不再探討罪的問題.
而是透過聖靈.
怎樣可以幫我們去解除罪的性情.
過去的那個引誘.
當然第六章曾經都說過.
保羅說.
你從前是怎樣將自己獻給不法不義做奴隸.
現在你要做一個新的奴隸.
就是要將你自己的身體獻給主.
做主的奴隸.
這樣又是一個整個內心生命的轉.
去另一個主人的問題.
所以來到這裡保羅告訴我們.
生命聖靈的律.
在我們信主的那一刻已經注進我們的心靈.
現在它的運作.
首先是要讓我們感受到不再定罪.
和不再犯罪的自由.

$^{361}$因為它已經釋放了我們.
所以這個釋放.
大家不要以為你就好像一了百了.
其實是一生的戰鬥.
因為下面一直繼續會說.
你必須要學習隨從,體貼,聽取引導.
和不願有這種共證的心靈.
我們是神兒女的生命.
不斷你要感受到.
所以其實如果你在崇拜.
首先問一問在座今日或者最近.
你覺不覺得自己是神的兒女呢?.
你有沒有這種感受到的福氣呢?.
可能這個這麼簡單的問題.
不是每個人都敢回答你.
其實不是的.
我信耶穌.
理性上我就是神的兒女.
有洗禮證的.
但是那個生命的關係.
你能夠很真誠地.
一閉上眼睛就能夠稱呼神是你的阿爸父.
這個是聖靈的印證.
所以這個是關係上的共證.
所以其實一直到已經大概有17方面的.
在這一章由第一節到17節.
這裡已經有7個方面.
18節是明顯另外一個段落.
本來我今晚要講的.
就是聖靈帶動萬物的復興.
因為18節是要再一次帶我們去看.
保羅再帶我們去看.
墮落了的人類世界.
怎樣得到那個改變呢?.
所以保羅的視野是在講創造.
是在講人類的墮落.
一來他有舊約.
他是舊約的拉比.
他真的對舊約是相當的通透.
在舊恩裡面.

$^{401}$他就能夠揮灑自如的.
將神的救贖計劃.
和人類的墮落來結合.
在舊恩在耶穌裡面.
他現在告訴我們.
原來現在是萬物都嘆息的時候.
連基督徒22節都是心裡面嘆息的.
不過我們信了耶穌.
有第三種的嘆息.
這應該就是24節.
聖靈的嘆息.
我們稍後再講.
這是第八章頭的17節.
我們上一次也講過了.
不定罪.
聖靈會引導我們大家的心靈.
要多點去敏銳.
聖靈在我們裡面的教導.
引領.
讓我們能夠學習隨從聖靈.
在我們裡面的內心的帶領.
以及讓我們感受到.
我們是神兒子的名分.
和羨慕永恆.
有我們自己的產業.
因為我們每一個都是後製.
是有繼承產業的權利.
將來是會有的.
不僅是《羅馬書》.
《爾弗所書》一章也講了.
大家如果不熟悉的話.
有時會錯過這些重要的經文.
《爾弗所書》的第一章.
明顯是講天上的福氣的第七樣.
第十四篇.
他說這位聖靈.
《爾弗所書》第一章.
新約的古羅的書信.
十四節.
他說這位聖靈.

$^{441}$是我們得到基業的憑據.
原文就是那個「志」這個字.
直等到神之民被贖.
使到他的榮耀得著稱讚.
他要令到我們能夠得到將來的產業.
以至他得到榮耀.
而且是我們去稱讚他.
古羅就很強調.
我們有這個印記.
就是聖靈.
這個印記的意思.
就是一個訂金.
這個「志」原文就是訂金.
這個憑據.
就是說.
你們將來一定可以得到天上的屬靈的基業.
是神樂意賜給我們.
擔心我們神會失信.
所以神就要保證.
所以他就要下訂.
下一個很大的訂.
就是三位一體的聖靈.
給我們每一個.
其實不用了.
神是不會失信的.
他是公義.
他是信實的.
你為什麼需要用聖靈做一個訂金呢?.
聖靈本身就是神.
將神給我們做訂金.
古羅的意思就是這樣.
這個應許的聖靈作為印記.
這位聖靈是我們得到基業的憑據.
就是說將來我們就可以得到基業了.
所以我們應該要有這種這樣的.
不是奢望.
更加不是一種幻想.
是一種渴望.
將來我們要得到天上的基業.
你可以想像得到.

$^{481}$而且還是有聖靈在我們的內心作為一個訂金.
所以一直去到神的子民.
得到救贖.
就是耶穌回來了.
我們被贖了.
我們就稱頌這位主的榮耀.
在這一大段經文.
榮耀得到稱讚.
是這一段頌歌的主題曲.
因為講了四次.
在這一段經文裡面.
例如第六節.
使得他的榮耀得到稱讚.
接著下面又有講的.
第十二節.
使得他的榮耀可以得到稱讚.
接著就是我們剛才看的經文.
你可以想像得到.
這裡保羅給我們看到的.
就是聖靈其中一個這麼重要的工作.
所以原來祂來到我們的心靈.
不要以為我們只是得到的憑據.
我們有聖靈.
就證明我們是神的兒女.
原來不單止是這樣.
神給我們的是永恆的基業.
頌這個夾子很重要.
因為那些風.
終於我拿了個衣領回來.
夾住那些紙.
因為又薄.
你可想像得到.
我們在一到第十七節第八章.
我們就逐樣去看看.
不定罪的自由.
就是釋放.
就是說第六第七章.
我們可能不同階段的人生.
比如說初信主.
信了主之後慢慢再被捆綁.

$^{521}$再被引誘.
再落在一些各樣的試探.
那種捆綁.
那種痛苦.
是信了主和未信耶穌.
是有一點點不相同的.
因為我們信了耶穌.
再重蹈覆轍.
那個從來沒有信過耶穌的人.
他們會習以為常的.
那叫做罪終之樂是不相同的.
所以第八章.
保羅很強調就是.
聖靈的進入.
應該讓我們越來越釋放.
越來越感受到.
救恩帶給我們不是限制.
是自由.
如果我們還停留在.
為什麼不可以吸煙.
為什麼我不可以買張六合彩.
我又不是貪心.
十塊而已.
我做做善事也不行.
很難說逐樣去計算是否可行.
染佛也不行.
紋身也不准.
這麼多限制.
如信耶穌.
很多時候就會為了這些事情.
但是整體來說.
信耶穌應該你自己的內心.
充分感受到一份釋放.
一份自由.
而且這個自由是不定罪的自由.
而不是放任自己的自由.
你會很有意識.
你去到禮拜堂.
你自己會知道我要安靜.
準備心靈敬拜神.

$^{561}$這不是一個規律.
不是感受堂硬性.
你坐下來聊天.
你就不好了.
而不是這樣的.
我們應該慢慢.
不是要挑戰權威.
而是我們要進入.
真的同神同在帶來的生命自由.
所以這裡應該是表達.
一個生命成長的表現.
某程度上這17節.
就是一個生命成長.
在救恩的愛裡面.
越來越感受到的救恩的套餐.
你買了這東西.
裡面就有這些內容.
你就可以用了.
你想想.
在這17節裡面.
我們用了多少救恩的內容.
我們是不是感受到不再定罪的自由呢?.
我和我弟弟都說.
我和他差不多一樣的開始.
我們那個青少年階段.
突然之間.
在我未信耶穌之前的前一個月.
有一天突然間起來.
我不吸煙了.
一天吸了兩包煙.
突然間起來我不吸了.
接著我就接近信耶穌了.
就信耶穌.
由那一刻開始.
從來沒有再吸煙了.
完全什麼戒煙糖.
那些完全沒有.
突然間我就不再吸煙了.
就已經感受到.
其實接近信耶穌了.

$^{601}$還沒有信.
已經有這種自由了.
已經不吸了.
沒有那個癮了.
當然現在聞到煙味都很討厭.
會咳嗽,會發作.
你可以想像得到.
這是很奇妙的生命力.
所以這一大段經文.
是說一個很正常基督徒的生命力.
所涵蓋的內容.
所以保羅不是要和大家逐點去挑戰.
那時候都不會有煙.
那時候都不會有其他的紋身.
可能有的.
但是問題就是.
那時候也有很多拜偶像.
有很多更難的挑戰.
但是保羅面對的.
現在要否釋給我們知道的就是.
他有神給他的啟示.
他能夠將在我們裡面的聖靈.
這麼具體地否釋給我們知道.
叫我們跟隨著去做.
我們就會進入到.
與聖靈同在帶來的救恩的榮耀.
就是我們這個人本身.
真的會脫胎換骨地去成長.
每個人都有他不同的時間.
來去有他的生命的成長和改變.
所以作為基督徒的老師,教師.
我們要有耐性去引導每個基督徒.
他都有自己的不同的時間表.
神在他的生命裡面的工作.
所以有些人很快 有些人很慢.
不一而足的.
我回望跟我一起成長在團契的弟兄姊妹.
都有不同的階段.
包括我自己.
所以我當時是一個很突變性的信主改變.

$^{641}$跟他們是很不同的.
所以他們某程度上都好像發光地看著我.
接著很快就奉獻傳道.
在教會讀神學等等.
已經跟他們走了一條不同的路.
但其實再回望.
其實一樣都是生命的成長.
如果不成長.
有一段階段我就很律法主義.
看不起那些走不了聖經的人.
很快就是一信耶穌.
大概一年都沒有.
在太子長榮大廈.
12樓就是種子書室.
很深刻的.
剛初信很雄心壯志.
跟一個弟兄在12樓看完書買完書.
一起下樓.
那個時候我們上的是很傳統的福音堂.
很嚴格的.
我們就覺得自己是跟足了做的.
在長榮大廈對出太子道過馬路.
紅燈說我們是禁忌.
不可以過的.
站著不動的.
我看著弟兄這樣走.
我的心就說.
怎麼搞的.
剛剛才跟你一起坐電梯.
這是我剛剛初信的時候.
你其實表面很公義.
其實你就是很看不起人.
你慢慢回望再看.
你已經知道原來聖靈的工作不是這樣的.
不是你內心對人的批判.
沒錯.
你是很公義.
你是不衝紅燈.
但是不代表你真的成長.
這是初信的階段.

$^{681}$你就很快就有因為律法而帶來的這種的恣意.
這種恣意為是的驕傲.
雖然是很少很低層次的驕傲.
但是你發展上去.
其實就是律法主義了.
你再發展上去.
你慢慢就變成很犬儒.
甚至內心是很苦毒的批判.
其實是不好的.
問題是還要覺得不對.
不是自己錯.
那裡更慘.
那裡更嚴重.
你就變成了.
其實根本就不是第八章羅馬書所說的真正的自由.
所以你看到.
這裡告訴我們.
一到第四節.
一種真正不定罪的自由.
有解答的.
是因為生命,聖靈的律.
在基督耶穌裡面釋放了我們.
就是說我們是認同主耶穌的愛.
祂的犧牲.
祂的十誡的救恩.
祂的由死復活的能力.
在我們裡面.
那個能力就是什麼呢?.
不是電.
不是一些神秘的力量.
很清楚告訴我們.
是聖善的靈.
已經在我們裡面了.
聖善的靈就會將救恩的成果.
不斷在我們裡面感動.
來引導我們.
所以聖餐.
如果我看到有人哭.
不是裝模作樣的哭.
是真的感動的哭.

$^{721}$甚至他走出來.
他自己不由自主地.
有聖靈感動地哭.
其實就是聖靈工作.
所以每個人應該自己安靜.
去到神的面前.
在一個這麼神聖的.
我再講原來.
原來有四個解釋的.
聖餐的理論.
我們不是相信變質說.
就是天主教的.
也不是馬丁路德的同質說.
因為都是換湯不換藥而已.
因為他不敢完全脫離天主教.
我們表面就是嘉爾文的紀念說.
但是.
對不起.
慈雲里才是紀念說.
但是應該是四個結合紀念說的.
嘉爾文說.
我們仍然是那個餅.
是那個杯.
和慈雲里一樣的解釋.
但是在那個聚會裡面.
有聖靈的臨在.
餅沒有變.
杯也沒有變.
是我們經驗到聖靈臨在.
在那個聖餐裡面.
所以嘉爾文是第四個學說.
這個聖餐.
所以我們要記住.
我們是讓我們的心靈.
感受聖靈的同在.
帶給我們的感動.
對主耶穌的救恩的愛的一個回應.
所以如果能夠.
哼起一些詩歌.
你自己內心去感動.

$^{761}$或者大會安排的樂的聖詩.
你去唱.
或者詩班的頌唱.
你都要投入進去.
我經常都說笑.
我們在台上.
尤其是感受這麼明顯.
看到人們.
走快一點.
很累.
你感覺到他好像是這樣.
當然我沒有權去批判.
但是覺得.
啊.
一直這樣.
每個人都有不同的生態.
每個人真的是很神聖的時刻.
是你自己跟神.
跟主耶穌的生命的聯繫.
雖然好像開玩笑.
但是其實那一刻.
你想想一個月只有一次.
一年只有十二次.
有時候去旅行又沒有再一次.
那你一生又有多少年呢?.
當你每個星期都有.
你也是最多一個月四次而已.
其實是很少的.
都不珍惜.
都不緊張.
那就可惜了.
不是罪.
但是可惜而已.
如果真的可以讓你.
又不是我們可以像草芬那樣.
又增加你的什麼.
又不是這樣.
不應該是這樣.
而是真的是一個愛的時刻.
跟主耶穌的敬拜.

$^{801}$是的.
無論如何.
將來我們在永恆.
都會吃這個新的筵席.
但是地上是一個禦常.
所以.
相信這一大段經文.
是我們要認同耶穌基督的救恩.
帶給我們的生命的感動.
如果不是基督的死.
沒有一個地上的人.
可以脫離神的審判.
因為那個就坐已經發生了.
我們每一個人都要面對生老病死.
還有為什麼會死呢?.
因為罪在我們裡面.
與神已經隔絕.
我們已經與神隔離.
所以.
救恩的開始是.
原來神的憤怒.
從天上臨到地上不揭不移的人身上.
但是.
最終告訴我們.
這樣的人.
神都願意拆遣他的兒子.
為我們犧牲十字架.
所以.
這就是讓我們可以脫離罪.
以及死的律的原因.
就是基督的愛已經盡了.
藉著聖靈.
進入我們大家的心靈.
所以.
我們如果只是靠自己.
是不能夠去行神的要求的.
但是現在.
第四節就說.
律法的義.
就成就在我們這些隨從肉體.

$^{841}$自己隨從聖靈的人身上.
再一次告訴我們.
能夠讓我們不再捆綁.
是因為我們有新的主人.
第六章的意思.
這裡再一次告訴我們.
我們跟隨聖靈去行事為人.
跟隨.
原來要跟到怎樣呢?.
保羅玉靈的另一個字眼就是第五節的.
體貼聖靈的事.
不要體貼肉體的事.
這裡每個時代都有不同.
方式也有不同.
現在的青少年被太空遊影響.
我們都完全不知道是什麼.
因為不是我們那個年代.
我們的年代嚴重的就是白粉.
連小K也沒有.
在我的那一區.
很多跟我一起成長的年輕人.
吸毒,吃白粉.
真的進出監獄.
變成盜竊.
他的人生就毀了.
體貼.
這裡說是完全貼近聖靈.
所以這裡是多麼寶貴.
我們能夠有一位愛我們的主.
帶領我們的人生.
我讀中學的學校都是Band 5.
都是最差的飛車學校.
經歷過同學斬人,開騙,講數.
有些是跟黑社會大佬.
即是大佬在門口等放學等等.
你可以想像得到你體貼誰.
我試過一次跟同學去一個佛堂放學.
發現他見到大佬和大佬的女朋友.
在那個佛堂裡面.
原來都挺惡的.

$^{881}$後來輾轉 因為家裡很嚴.
太晚回家又會罵.
因為我哥哥沒什麼多.
哥哥姐姐多.
變成令我徘徊在差不多可以進入黑社會的階段.
但最終都進不了.
因為家裡真的很嚴.
所以見到同學後來真的跟了大佬.
覺得沒辦法了.
因為住在那些屋村.
父母又沒有人管.
變相跟他的分別.
只是他進入了黑社會我就沒有了.
這個就是我中學階段.
如果回望相信是有神的保守.
各方面的.
但這裡告訴我們.
不是說個別哪些人特別慶幸.
特別有福.
而是每一個在主救恩裡面的弟兄姊妹.
在線上的或者之後聽錄音的弟兄姊妹.
或者在不同地方比如荷蘭,嘉欣.
你們教會聽的.
我們要記得每一個地方.
每一個時代都不同.
但聖靈扮演著新主人的角色.
是永遠不會停止.
一直到我們回到天家為止.
他都要做這個大靈.
而且他是歡喜.
這樣在我們的心裡面.
來引導我們.
所以這個就是說.
如果我們跟團.
我們最早期開始去旅行.
一定不會自己去.
一定是跟旅行團.
當然你去得多就發現很悶.
指定的地方.
沒有自由.

$^{921}$但是你要記得有一個領導.
你跟著他.
你就可以聽到很多東西.
學到很多東西.
這裡就是我們永遠的.
內著的這一位新的主人.
就是聖靈.
藉著基督的救恩.
就在我們的內心裡面.
所以第四,第五到第十七節.
就是聖靈內著之後.
一直保留著不同的角度去說.
我們要聽取他的引導.
我們就能夠得到剛才說的自由.
而且我們能夠自死身體的惡行.
這就是第十三節.
你們若順從肉體活著.
必定會死.
就靠著聖靈自死身體的惡行.
必要活著.
就是說有一個幫助者.
永遠了解,明白我們的內心.
我們只要跟隨他.
來去順從他.
我們就能夠自死身體的惡行.
就是說.
譬如用第七章保羅所說的.
貪心.
他在聖靈裡面.
他可以勝過罪的人有.
可能後面會更詳細.
這就告訴我們聖靈的律.
上一次我們說了.
聖靈的律和罪和死的律的對照.
一個對立.
然後我們再不是奴隸.
我們是神兒女的身份.
所以這是一個轉化.
所以正常基督徒.
是要不斷地越來越肯定你是神兒女.

$^{961}$你越來越能夠感受到神兒女的氛圍和氣質.
流露在你的生活裡面.
這是神的兒女.
不是神經兒女.
是神的兒女.
我們能夠舉手投足.
不是古怪,不是神秘.
是一種很自然的待人接民.
其實我很開心.
在洪水橋.
那些街坊,石橋,柏宇那些.
新來的那個紅姐.
還沒來就已經說曾先生.
那些街坊有時候特意開門.
比較遠一點.
有個老人家準備出門.
我都開門等他出門.
慢慢又建立了這種鄰舍關係.
慢慢一點.
而且不是特意被逼做的.
而是很開心地去做.
你走過單車路.
跟大嬸打聲招呼.
人家慢慢又會跟你有認識.
現在開了一間寵物醫院.
又跟那些職員很熟.
他們就在我們後巷吸煙.
見到我就說:曾牧師.
立刻把煙收起來.
其實不用的.
你又不信耶穌.
但是他們知道是教會.
牧者見到他吸煙就不太好.
其實他們要放電.
放鬆.
但是希望我們能夠舉手投足.
讓人們看到我們內心的盼望和喜樂.
所以很開心.
可以不斷每天都見到很多街坊不認識的朋友.
這就是由奴僕的身份.

$^{1001}$轉化成為神兒女的身份.
所以我們有什麼機會讓我們的年青人.
我們的老人家.
感受到自己是神的兒女呢?.
相信聖經的教導.
和不同的方式.
讓他去承載這個新生命.
其實是重中之重.
如果當人來到教會.
都不再覺得自己是神的兒女.
只是覺得我們有一個很好的氛圍.
大家有好吃的.
喝的.
聊天的.
是一般的事情.
不再探討聖經真理.
那就慢慢會失衡.
慢慢就會失去那種觸覺.
慢慢就變成一個類似俱樂部的聯誼.
其實是不健康的.
因為那只是救恩帶來的副產品.
不是真正救恩的內容.
真正的內容是我們由奴僕的身份.
變成神兒女的身份.
所以你可以想像到.
保羅在這一大段經文.
不斷想說.
聖靈已經負責起領導我們心靈的工作.
如果我們一生都不讓聖靈去引導.
那為什麼難怪我們慢慢回教會.
你和著了.
慢慢我們根本不稀罕.
對人不夠的那種冷淡.
根本就不再嚮往一些陪靈的聚會.
希望我們的靈明可以不斷地提醒.
不單純是充電.
是有一種共鳴的提醒.
不再嚮往有日得救的快樂.
原因就是我們根本都已經沒有讓聖靈.
或者從來都不讓聖靈.

$^{1041}$或者不知道要讓聖靈來帶領和管理我們的人生.
所以你可以想像到第八章是多麼重要的一段經文.
而且是由第六,第七章.
一路這樣去講解.
一路慢慢推進到那個核心.
怎樣能夠不犯罪.
不可犯罪.
變成可以選擇犯罪.
不再在罪的那種游論底下.
就是因為第八章講的聖靈.
負起這個領導我們心靈的那個重任.
那個工作.
下面第三段那裡就還沒有時間講了.
又要再等一堂.
另一個角度.
釋放剛才講了.
再仔細一點就是.
神的律法的義.
我們是要行神的道.
有力量去使我們承傳律法的義.
某程度上就是耶利米蘇的三十二章還是三十六章.
那裡講的將來神的道.
寫在我們的心裡.
那段經文的應驗.
就在第八章裡已經解釋和應驗了.
例如希伯來書的八章.
就是引耶利米蘇的.
整段引.
第十節.
八章的十節.
當耶穌基督成就了大祭司的工作之後.
他說那些日子以後.
我與以色列所立的約乃是者一樣.
我要將我的律法放在他們裡面.
寫在他們心上.
我要作他們的神.
他們要作我的子民.
這起碼就是《羅馬書》八章講的福氣.
聖靈與我們的心同正.
我們是神的兒女.

$^{1081}$起碼我們有律法的義在我們的心裡面.
因為我們隨從聖靈.
現在這個預言在耶利米蘇的三十一章.
不是三十六章.
他們要作我的子民.
第十一節.
他們不用各人教導自己的鄉鄰和自己的弟兄.
而是說你該認識主.
因為他們從最小的到最大的都必認識我.
我要寬恕他們的不義.
不再紀念他們的罪天.
記說新約.
這個就是新約裡面的內容.
原來舊約耶利米蘇三十一章.
已經有對將來的應許.
耶利米蘇就是亡國時期的作品.
警告已經沒有辦法再回頭了.
神一定會刑罰.
所以耶利米蘇的宣告就是要投降.
不要再抗拒.
去到巴比倫.
七十年後就會有回歸.
但是當時的以色列人不聽.
最後的王就是西底加.
覺得可以.
因為當時有些賊才知道.
不用怕.
我們可以抵抗巴比倫.
最終耶利米蘇的預言應驗.
那個說話.
但是原來在這樣的亡國裡面.
都預見將來.
遙遠地望見將來.
會有一個新約出現.
現在新約的作者就引述耶利米蘇的說話.
整段差不多就是羅馬書第八章.
那些福氣.
不再定罪.
有那個心靈的感動.
我們是神的子女.

$^{1121}$我們是祂的子民.
我們有教導.
我們可以認識主.
而且我們從最小到最大都認識神.
神會寬恕我們的不義.
不再紀念我們的罪惡.
你可以想像得到.
全部就是第八章的羅馬書的描述.
所以保羅某程度上.
他對舊約的認識.
神很早期.
耶利米蘇就是主前的七百多年.
六零幾.
亡國的階段.
已經望到六百年後.
耶穌基督來成就舊恩.
就會有這個效果.
只要在舊恩裡面.
我們就能夠有律法的義.
可以成就.
我們就能夠是神的兒女子民.
即是說.
就算我們不是肉身的亞伯拉罕的子孫.
我們都已經透過耶穌.
我們就是屬靈的神的兒女.
某程度上.
加泰書已經說了.
三章.
我們就已經是這個.
因為舊恩帶給我們的.
我們就是亞伯拉罕的屬靈的子孫.
所以我們想像得到.
這個就是新約.
回到羅馬書八章.
這裡說.
律法的義成全.
就是這個意思.
我們能夠可以透過聖靈一路引導.
我們就可以有力量去成全.
去遵行神的律法的義.

$^{1161}$當然就是心思意念的更新.
因為我們已經可以在神裡面.
在第五章.
我們一路隨從體貼.
就是生命的平安.
所以我們的心思意念.
是被聖靈去更新.
接著.
再下面還有12到14節.
就強調我們有身份的.
對身份的認識的確認.
因為14節.
因為凡是被神的靈引導.
都是神的兒子.
已經不再糾纏在我們.
在肉體的遵行.
和要來隨從肉身的肉體的凡序.
而是我們已經在神兒子的生命裡面.
有這種引導.
即是對自己的屬靈身份的確認.
確具.
我們是神的兒子.
因為我們有神的靈的引導.
馬太福音耶穌的百福.
使人和睦的人必稱為神的兒子.
所以如果用另一個角度.
屬靈身份的確認.
我們隨從聖靈.
被神的靈引導.
被神的靈引導不是去做古怪的事.
突然間講方言三天都不停.
不是這樣的.
更多的是使人和睦.
因為耶穌說的.
使人和睦的人有福了.
必稱為神的兒子.
所以保羅的說話就是.
神的靈引導我們去做甚麼呢?.
一定是做一些和相愛.
帶來和平.

$^{1201}$帶來合一.
帶來全福音.
一起有一個心智去做.
我們就是神的兒子.
所以絕對不是古怪的.
所以我們一聽聖靈就覺得是靈因.
古怪.
叫你做一些很特殊的事.
那些你要小心.
可能個別的人都會做.
但我們不能說是普遍性.
真正的普遍性.
和聖經真理的引導.
一定是和群體的合一.
相愛.
促進相愛.
肯定是聖靈引導的.
聖靈不可能引導你去做一些挑撥.
離間.
窒道.
紛爭.
不可能的.
絕對不可能的.
連你自己的心都知道是不可能的.
但你可以想像到.
這裡就告訴我們身份的確認.
我們是神的兒子.
你想想有幾個寶貴.
地球上有八十億人.
其實很難叫人信耶穌.
很難告訴他們.
耶穌是真的.
是復活的.
是全能的.
是永恆的.
我們今天已經信以為真.
以及已經這樣行使著這種力量.
在我們的生活.
但原來在這個世界.
和我們一起生活的人.

$^{1241}$是有更多比我們厲害的人.
完全藐視這種能力.
但他自己也不能改變自己.
那你想想多可惜.
原來死而復活原來這麼難去信.
但我們今天已經得到了.
這個就是身份的確認.
在聖靈裡面.
我們越來越對自己身份的尊貴確認.
其實是很尊貴的.
接著就是親密關係.
在心靈的平安.
第十五節是你們所受的.
不再是奴隸的失落.
一直繼續害怕.
我們沒見過奴隸.
我們有生之年都沒見過奴隸.
在我們以前住的地方.
現在也是.
錦鹽的大江埔村.
原來有一個家庭.
後來我們知道.
原來他們是移民來的.
原來他們的祖先是真是奴隸.
挺有趣的.
都認識了這個家庭.
某些我姐姐和他們都熟的.
他們的祖先真是奴隸.
不知為何後來來了香港.
可以想像得到.
你們在那個時代.
如果是奴隸的.
那種懼怕,那種痛苦.
我們沒有這個感受.
因為香港這麼自由.
各樣東西都越來越趨向平等.
不會有剝削.
不會有一些嚴刑的.
對工人的毒打.
是沒有這件事的.

$^{1281}$也都基督教的氛圍.
香港都趨向平等.
個別可能還是很苛刻的僱主.
但都很少.
因為一有就馬上被警察抓.
馬上禁止.
但你想一下整個時代.
你是奴隸.
你就見怪不怪了.
但本來能夠講到那個精髓.
如果我們繼續是奴隸的心.
就會害怕.
就會沒有那種心靈的平安.
現在我們受的是兒子的心靈.
所以這裡說我們呼叫.
阿爸,父神,聖靈和我們的心.
同證我們是神的兒女.
而且是一個可以繼承將來產業的繼承人.
後詞的意思就是一個繼承人.
所以你可以想像到.
這是一個親密關係之下的有福的人生.
就是因為聖靈在我們裡面.
不斷讓我們自己對自己的存在.
那種感恩,確信,確認帶來的平安.
不再是懼怕的奴隸.
而是兒子.
兒子就是說繼承產業.
兒子就是有法定的地位.
有那個氛圍.
你就可以來享受這個家庭屬於你的各樣的東西.
所以如果神國度是這麼浩瀚的.
我們有沒有行使到禱告的權柄呢?.
我們可以捆綁一些神要我們捆綁的重大的事呢?.
我們有沒有為一些宣教士在危險的地方.
我們為他們去釋放.
可以來到去傳講福音的那個能力呢?.
這些是權柄來的.
我們在這個家裡面就行使著這個權柄.
比如明仔回家.
我問他他問我.

$^{1321}$我可不可以開雪櫃吃點東西.
整個家是你的.
你還問我可不可以開雪櫃吃東西.
是不是?.
我們覺得匪夷所思的.
但是保羅告訴我們.
我們是基督徒.
現在聖靈和我們有這個親密關係.
神國度的那個浩瀚的豐富.
是要我們去知取的.
將來去天堂.
每個人都是神在天上的兒女.
還知取什麼呢?.
都沒有爭戰了.
今生就是.
現在你就要寶貴這個平安的身份.
有多少人每天活在那種恐懼之下.
現在星島.
我喜歡看星島.
你看一個新聞.
它現在有一個連接.
你一直下去.
全部都是事實.
幫你全部列出來.
上個星期再早一點.
再早一點.
一些是老人家.
一些是年輕人.
列出全部的新聞.
因為它有全部的字的概念.
所以你一看這個新聞.
你一掃下去.
它就會將你列出來歸類了.
那你可以想像到.
原來每天有很多人活在恐懼裡面.
這一段經文.
是聖靈在一個成熟的基督徒裡面.
其實應該叫正常的基督徒裡面.
就有這份親密關係帶來的平安.
所以我們能夠稱呼阿爸父神.

$^{1361}$造物主就是我們天上的父神.
造物主就是將來有一個計劃.
所有的人都要面對審判.
世界是有結束的.
而這個結束不是一個滅亡.
是一個更新的新的世界.
後面就會說了.
所以如果我們自己信了耶穌.
自己都沒有那份平安.
又何來去祝福人呢?.
我們自己有這份平安.
是不是我們應該和別人分享呢?.
是不是我們應該.
不是要黑銷.
而是讓人感受到.
其實信耶穌不是一個宗教.
而是一個生命的一個皈依.
和神造物主有一個生命的聯繫.
你看不到神.
我又看不到神.
但為什麼我的生命會有平安呢?.
我一樣和你一樣面對不同的困難.
為什麼我會有平安呢?.
就不是因為我比你厲害.
而是因為我和這位造物主有接觸.
接觸就是聖靈一直的引導.
所以不單止是我們有平安的生命.
更加我們是有繼承人.
所以不是開玩笑的.
你家裡的子女.
你不會告訴他們.
請你先問我.
你要這樣做這樣做.
你進屋都要先問我.
都不可能是這樣.
你就會知道.
這個是家.
家就是所有的東西都是你的了.
都是屬於你的了.
都可以公平分享的了.

$^{1401}$你可以想像得到.
我們又要做一個問卷調查.
崇拜頂子妹.
你知不知道你將來是繼承造物主的產業的.
我想大部份人都未必聽過.
但這個保羅是講了很實在的.
如果我們是兒女便是後製.
這裡不是說因為你得性.
因為你帶了幾個人信耶穌.
如果我們是兒女便是後製.
就是神的後製.
和基督同作後製.
不言不告的角度去講.
我們是後製.
是和耶穌基督同作後製.
所以我們在地上.
自然就是要學孝耶穌.
祂是一個被拒絕.
被人不認識.
知道都不會接受祂的一個救主.
所以我們去介紹.
我們去傳揚祂.
我們一樣會面對這個苦楚的.
不過保羅已經講了.
地上的是短暫的.
因為永恆的才是真正的榮耀.
所以大進十八節.
就是一個新的段落的開始.
所以事實上.
這個是漢語聖經來的.
我認為現今所受的苦.
與將來要向我們顯現的榮耀相比.
根本就算不得十惡.
所以將來的榮耀有多少呢?.
我們根本不能夠用我們現在的肉身的感官.
那個智慧去了解明白.
可能是一種類比.
因為我們知道墮落是什麼.
當的墮落導致我們人類墮落.
我們又知道生老病死的痛苦.

$^{1441}$我們每天都見到都發生.
所以我們已經有心理準備.
我們都知道什麼是苦.
我們傳福音有時候都會被人笑.
如果你在極權的地方會被人打.
被人抓去坐牢.
我們知道地上這個已經在發生的墮落的世界.
受咒詛的世界那個痛苦.
但是保羅要帶我們去看的就是.
在這個永恆將來要顯現在我們當中的榮耀.
我們要有一個這樣的憧憬.
其實這個叫做渴望和盼望.
根本就不算什麼.
保羅在講這番話的時候.
他自己都是身心靈都是受壓.
他面對哥林多教會.
因為在哥林多寫了這封羅馬書.
面對哥林多教會對他的誤解.
又會有不同的手段去針對保羅.
對他的傳道.
對他的為人.
有很多不同的攻擊.
還沒有到要抓他坐牢.
因為那個是彭波特後書人生的階段.
最後的階段.
在這個時候他還是可以自由的傳道.
還沒有被當時的猶太人去說他踐踏聖殿.
來抓他.
現在還是很自由.
但是面對教會種種不同的需要.
他也是內心有他的那種擔憂.
有他那種為基督的教會.
他也很想感同身受的那種承擔.
就是保門基督患難的缺欠.
歌羅西書應該是第二章的經文.
他也很想保門基督患難的缺欠.
這樣的心智.
所以他這樣講的含金量就很高.
他說我認為現今所受的苦.
他現在在傳道裡面.

$^{1481}$又不是猶太教的拉比.
面對外邦的信徒.
要帶他們不要再拜戴安娜.
就是在義不所的守護神.
那個低迷雕.
還是忘記了名字了.
你想一下會招致外邦人.
特別是拜偶像的那些產業的人.
來對他的攻擊.
這個應該是少林寺十七章.
對不起 少林寺二十三章那裡.
提到那幫人對他的攻擊.
低迷女神對不起.
就是一個戴安娜.
就是在義不所入口的神像.
或者一條柱子.
你可以想像當時的幾個民衛.
本來就在那個時代.
雖然兩邊都討好的.
猶太教又會攻擊他.
帶猶太人不要守律法.
外邦人又不要拜偶像.
不要拜君王.
就變成了真是有一個很大的.
對他的那種壓迫.
所以他所描述的苦.
一定是很具體.
一定是為基督而受的困苦.
但是他是很正面的去看這些人生.
短暫的苦楚是實在的.
是真的會有各種的傷痛的.
哥林德後書十一章.
他就描述了很多.
他自己既然你說我不是使徒.
我就告訴你.
我為基督所受過的一些患難.
他不是要誇耀.
他說我好像一個愚昧的人.
來講我的一些需要.
哥林德後書十一章的二十一節.

$^{1521}$我說者說是羞辱自己.
保障我們從前是軟弱.
然而人在何事上勇敢.
我說句愚妄話.
我也勇敢.
他們是希伯來人嗎?.
說的是那些假的使徒.
我也是.
他們是以色列人嗎?.
我也是.
他們是亞伯拉罕的後裔嗎?.
我也是.
就是說這樣的誇耀自己的.
幾根正苗紅的猶太的拉比.
就是他的對手.
所以保羅就再說了.
他們是基督的僕人嗎?.
我說句狂話.
我更加是.
我比他們多受勞苦.
多下監牢.
受鞭打是過重.
冒死是屢次都有的.
被猶太人鞭打五次.
每次四十下.
減去一下.
因為他是猶太人的一個拉比.
被棍打了三次.
被石頭打了一次.
遇著損壞三次.
一晝一夜在深海裡.
又屢次行遠路.
遭江河的危險.
盜賊的危險.
同族的危險.
外邦人的危險.
城裡的危險.
曠野的危險.
海中的危險.
賈弟兄的危險.

$^{1561}$受勞碌.
受困苦.
多次不得睡.
又飢又渴.
多次不得食.
受寒冷.
赤身暴體.
除了這些外面的事.
就是肉身的事.
還有為宗教會掛心的事.
天天壓在我身上.
有誰軟弱我不軟弱呢?.
有誰跌倒我不焦急呢?.
是在跟那些假使徒對話.
你說我不是使徒.
我就要好像誇耀一些.
讓你們知道.
他不是想去炫耀.
這些是不能炫耀的.
但是我們這樣去看完.
我們現在看的羅馬書.
它在說的現在的苦.
所以是短暫.
越來越有這個心境.
我們要.
因為現在經常在說那些明星.
什麼優雅地老去.
久而久之就說葉紅.
一頭銀髮.
葉善文.
很優雅地不染了.
由得他.
我自然.
不需要再染色.
其實基督徒也是.
我們應該在不同的階段.
來接受自己.
如果地上不順主的人.
都願意這樣接受.
我們也要這樣接受.

$^{1601}$神在我們生命裡面.
不同的安排.
例如喪偶.
例如我們要面對突然間的打擊.
我們也不要因為這樣而倒下.
當然人一定有他的情緒.
需要有同行者.
需要有陪伴.
走出那個痛苦.
是正常的.
因為我們真的是人.
不是一個石頭.
但是我們要記得.
地上始終是短暫.
所以我們要接受.
神在我們的生命裡面.
不同階段的帶領.
所以突然間有時面對一些事情.
看我們當下是否已經整個信仰崩潰.
有感受.
有困難.
有哀告.
是應該的.
是正常的.
但是哀告完之後.
你放棄信仰.
這就是我們的神學不夠全面.
我們在好景的時候.
我們沒有裝備一個整全.
一個全面的神學觀.
這一節已經告訴我們.
真的要將永恆和現在的短暫相比.
現在你不可能覺得是風平浪靜.
沒有事發生.
現在是我們越投入在那個服事.
越投入在基督的國度.
一定是會有苦楚的.
就算我們不是.
人生一定是有不同的悲歡離合.
不同的生老病死.

$^{1641}$所以我們要接受.
有時看那些老的牧者.
他們能夠很安然地接受自己面對不同的困難.
甚至接受自己不能再行動自如的港度.
其實這些是很成熟的生命的流露.
人生已經覺得自己可以被忘記.
是最高超的一個境界.
你是很樂意地被忘記.
不要霸佔自己的位置.
你要讓身來接續神的工作.
所以這就是要對這一節聖經的一種很重要的透析.
但是原來下面不單止是這樣.
保羅帶我們看的是.
原來整個宇宙萬物在就坐底下.
現在保羅在這裡有三個嘆息.
已經45分了.
看不到那裡了.
我們可以進入這裡一點點.
這裡就是我們今天講的主題.
「聖靈在萬物復興中的參與」.
肯定就是在說基督徒扮演著這個主導的角色.
保羅很有趣的.
受造之物其實是沒有靈魂的.
花草,樹木,動物何來他們會嘆息呢?.
但是保羅用了這麼美麗的擬人法.
其實就不是天馬行空的.
因為他帶我們回去創世記第三章的就坐.
原來萬物都因為人類的犯罪而受就坐.
在整個的就坐底下.
原來受造之物就服在這種虛空底下.
就是沒有辦法可以改變.
神的時間來到.
所有的東西都要服從在這個定序的定律底下.
直到真真正正的復興.
就是說還沒墮落之前的樂園的神創造.
我們不知道.
我們是阿當生出來的後代.
所以我們沒有阿當那個又有樂園.
又經驗過墮落的失樂園.
所以阿當就是唯一一個不同的夏娃.

$^{1681}$但是我們就不行了.
因為我們已經是墮落的人.
所以我們現在應該要另一個渴慕.
因為這個是保羅說的.
渴慕將來恢復沒有就坐的榮耀樂園.
所以又是一個問題.
又要問崇拜的弟兄姊妹.
你知不知道將來是一個沒有墮落的世界呢?.
就是我們回到阿當沒有墮落之前的那種與神相會的境界.
將來在天堂就是了.
似乎如果神告訴我們.
我們離世就已經立即與基督同在了.
我們的靈魂就已經去到樂園了.
所以其實應該是越來越渴慕.
進入樂園的那種法理律.
不是自尋短見.
一定是按照神給我們的時間去活.
我們的每一個日子.
這裡就告訴我們.
19,20節保羅說了這個引言之後.
其實是連繫上林的.
所以他給我們出路的.
在聖靈的引導下.
我們在苦難裡面其實是與神同工.
是與基督同工的.
這個是17節.
所以必定將來與他一起得榮耀.
所以將來基督回來.
來接教會.
來呼喚我們復活的時候.
我們就得榮耀.
真真正正全人的得榮耀.
那個新的宿命的身體進入永恆裡面.
所以今生我們是應該在短暫裡面.
我們的每一個重大.
應該都是與主有關連.
接著18節就帶進新的段落.
就是一個很大的保羅說的範圍.
就是整個受造物現在的狀態.
只要主回來就會有一個徹底的改變.

$^{1721}$所以我們如果有這個眼界.
我們就不只是個人的得救這麼簡單.
而是整個受造物真正解決的一個救恩.
夠不夠時間再講一點點呢?.
其實這裡就立即因為19節和20節.
就已經在講那個墮落.
保羅就已經.
雖然他沒有直接引用《創世記》.
就已經告訴我們.
《創世》的墮落是一個不能夠逆轉的.
在當時未有新約之前.
是不能夠逆轉的一個咒詛.
當人犯罪之後.
接著就是三章的14節.
耶和華神對蛇說.
你既作了這事就必受咒詛.
比一切的生畜野獸更甚.
你必用土止行走終生吃土.
我又要叫你和女人彼此為仇.
你的後裔和女人的後裔也彼此為仇.
女人的後裔要傷你的頭.
你要傷他的腳筋.
又對女人說.
逐個咒詛.
就是夏娃.
我必多多加增你懷胎的苦楚.
你生產兒女必多受苦楚.
你必戀慕你丈夫.
你丈夫必管轄你.
又對阿當說.
你既聽從妻子的話.
吃了我所吩咐你不可吃的那樹上的果子.
分別是樂樹的果子.
地必為你的緣故受咒詛.
所以受咒之民.
伏在虛空之下.
似乎保羅就是在這裡.
其實地沒有犯罪的.
不關他的事.
是阿當犯罪.

$^{1761}$如果地會說話.
你說他也關我什麼事.
但沒辦法.
因為整個大地是為人類設計的.
人類是生活在這個地上.
所以你是配角.
主角的結局就關你事.
所以大地也受咒詛.
所以將來我們回到一個沒有受咒詛的榮耀.
其實比起我們想像最美的地方.
更要美的不知道多少萬倍.
沒有咒詛的地方.
你可以想像到那種美麗.
因為我們已經是暴露的人.
我們不知道在哪裡.
時時提醒我們有地震.
地球原來地殼在變動.
原來是很不穩定的.
原來氣候現在有些是逆轉不了.
因為氣候的暖化.
因為地球本身可以自然調節溫度.
海洋的溫度原來有自己的修復機制.
但人類不斷破壞.
可能就逆轉不了.
墮落越來越明顯.
現在普羅表告訴我們.
未墮落之前就是將來的榮耀世界.
所以你看到.
你必終生農苦才能從地裡得吃的.
地必給你長出荊棘和質泥.
這些不是實創造的.
是現在才出現的.
野草是自己出生的.
不用人種的.
你也要吃田間的菜蔬.
你必汗流滿面才得糊口.
直到你歸了土.
因為你是從土而出.
你本是塵土.
仍要歸於塵土.

$^{1801}$即是說原先本來如果吃生命樹的果.
是不需要死亡的.
現在告訴你.
你要歸於塵土.
即是你本來是不值一文的塵土.
現在將來你要歸於塵土.
不過也要歸於930歲的命.
沒有說夏娃多少歲.
女性大.
男性香港就是85和80.
平均.
通常女性是比較長命的.
長命3歲還是5歲.
是不是夏娃就936歲.
我不知道.
我沒有記載.
只是用男系的歲數來記載家譜.
即是上古古代的人.
原來是有一個結束.
這個奏作就是說.
人是有結束的一天.
而且一代一代這樣結束.
很偶然的.
第七代的當日孫爾諾.
就不需要成為塵土.
他直接去到神那裡.
好像都被人感受到.
原來神的能力可以超越了.
那個奏作的死亡.
所以今天我們在耶穌的救恩裡.
肯定可以勝過那個罪和死亡.
對我們與神隔絕的這種關係.
不再是一種敵對.
而是能夠和好的關係.
所以現在我們已經在享用一個聖靈.
給我們的能力.
講到十九和二十節.
下次就要在十八節這一段重新再講.
我就停在這裡.
我們一起有個簡單的禱告.

$^{1841}$感恩天父給我們能夠再一次.
透過羅馬書我們繼續去深化.
第一個段落.
一到第十七節.
那個聖靈在一個正常基督徒的內心.
那個引導,感動.
讓我們去跟隨和體貼聖靈.
和成為聖靈的新的僕人.
是真正愛的一個跟隨.
盼望主甦醒我們的內心.
讓我們看到不是地上的生活不實在.
我們每天都面對生活.
但是我們有聖靈的感動和帶領.
我們面對不同的人事物.
讓我們能夠帶著一份這樣的愛.
生命的激情.
我們在地上去讓人感受到.
就算我們面對困難.
那種快樂都仍然是沒有失去的.
那份平安依然是好像在我們的內心做主導.
願主幫助帶領我們.
因為明天就真是新一年的培靈會.
很盼望主真是藉著幾位的僕人.
去復興香港的教會.
透過現在網絡的便利.
海外甚至國內不同的華人的地方.
都能夠透過主的僕人的教導.
去得到那個啟發.
得到那個生命的更新的力量.
奉耶穌基督的名.
阿們.
\newpage



\section{羅馬書}
\label{sec:HwCOo8EiKp0}
\textbf{2025-08-14 周四晚查經班《羅馬書(19):分題:聖靈帶動萬物復興二》 講員:曾裔貴牧師}
\newline
\newline
連結: \href{https://youtube.com/watch?v=HwCOo8EiKp0}{\texttt{https://youtube.com/watch?v=HwCOo8EiKp0}} ~~~~ 語音日期: 2025-08-18
\newline
\newline
\hyperref[sec:mDmug5anuis]{\small{< < < PREV SERMON < < <}}
~
\hyperref[sec:index_chronic]{\small{[返順時目]}}
~
\hyperref[sec:index_scriptual]{\small{[返順卷目]}}
~
\hyperref[sec:3_mKpJn_jtA]{\small{> > > NEXT SERMON > > >}}
\newline
\newline
$^{1}$好 電影節目晚安.
開始我們今晚的茶經班.
我們都一起同心禱告.
多謝天父 我們雖然停了一個星期.
大家都能夠在不同的學習.
透過不同的牧者.
在培靈賢經會去聽一些憤興的信息.
讓我們每年香港的信徒.
或者華人的信徒.
都能夠透過神靈的僕人.
在多方面去教導我們.
我們很盼望這一份屬靈的渴望.
是激發我們去認識聖經.
認識我們與神的關係.
認識我們在人生.
其實我們可以用一個很尊貴的身份.
神的兒女的身份去服侍身邊的人.
求主幫助.
也都讓我們能夠成為一個君尊的祭司.
是一個禱告的守望者.
也都為一些很重要的事情.
我們去同心禱告守望.
我們今天晚上我們看第八章.
一段相當令人振奮的經文.
讓我們看到在保羅的那個世界.
他所看到的榮耀.
那個將來的榮耀.
所涵蓋的範圍是多麼的大.
叫我們能夠真的看到.
以致我們不只是著眼於很短暫的.
我們個人的信仰的開心與不開心.
而是我們能夠看到.
神的國度要一直這樣發展.
要成就,要超越舊有的世界.
那個墮落的世界的那個對神.
以及對人的那個影響.
願神帶領我們.
我們仰望等候在你的面前.
禱告奉耶穌基督的名.
阿們.

$^{41}$很感恩我們雖然停了一次.
但是經文都給我們有一個延續.
今天我們上一個星期.
就是上一次我們就講到.
關乎到第八章.
特別我們停到第17節.
接著18節之後就沒有再詳細去講.
我們可以再一次回到第八章.
很重要就是在這17節告訴我們.
我們在主耶穌的救恩裡面.
我們得到什麼身份.
這個身份是如何確立我們與神的關係.
例如聖靈進駐在我們的內心.
就會令到我們可以開始學習.
勝過罪和死的律.
就是說原來我們一直不只是得到永生的天堂.
一個好像是將來的世界.
而是我們這個在地上由出世到死的.
很實實在在的有血有肉的這個人.
其實是有靈魂的.
是有裡面的那個舊有的那個性情.
是與神的國度.
與神的性情是有一個很大的距離的.
所以你可以說今生的不救.
由我們信主接受救恩.
我們就立即成為神的兒女.
聖靈在那一刻已經進駐.
其實就是不斷更新我們舊有的性情.
要與同神的性情有份.
所以你可以說基督教的圖造工程.
就是令到每一個基督徒意識自己有罪性.
意識罪性在耶穌的復活裡面.
已經失去他本來的力量.
當然每個人有不同的步伐.
有些信主立即意識到自己的生命的改變.
有些是慢慢的循序漸進.
但是又是不斷在《羅馬書》的第五章到第八章這一段.
就是不同的角度告訴我們.
我們要克性罪惡的辦法.
就不是我們去死守那個律法.

$^{81}$也不是要像猶太人那樣.
跳進那個律法的世界裡面.
而是讓我們看到由信主那一刻.
我們已經在神眼中成為尊貴的兒女.
我們可以.
後面那裡不是說了嗎?.
我們可以稱呼聖靈與我們的心靈.
一同的見證,印證我們是神的兒女.
既然是神的兒女.
我們就稱呼神做阿爸父.
其實我們犯了十屆摩西的律法裡面.
第三屆的.
不可以妄稱耶和華的名.
妄稱神的名在舊約是要被治死的.
不可以稱呼神的名.
但是我們居然今天稱呼阿爸父神.
這是一個亞蘭文的複合字.
就是很親暱的意思.
所以是告訴我們可以大膽地.
去稱呼這個身份.
我們就是神的兒女.
在這個地上.
所以其實是很尊貴的身份.
所以你可以說第五,六到第八章的十六節.
都很想告訴我們.
我們要在聖靈裡面.
克性罪惡.
以及要宣告罪惡已經不再是我們的主人.
舊有的聖情在基督裡面已經被自伏了.
我們現在.
雖然它仍然存在我們舊有的肉身的生命裡面.
但是我們已經在新的聖靈的管理.
帶領下來生活了.
所以隨著聖靈而行.
在聖靈的引導.
在這位新的主人愛的一直的引導下.
我們可以不斷地經驗神兒女這個身份.
或者說讓我們的心靈越來越活躍.
越來越敏銳.
越來越與神的關係是一種融合.

$^{121}$我剛剛在星期天就說了有聖靈有靈魂.
在圖書館發現了一本書.
就是叫做《After》.
這本書的中文名就是一個精神科醫生.
用了五十年行醫.
去找了很多很多無數的個案.
都是告訴他.
他寫成一本書是很專業的研究報告.
五十年他遇見的個案.
或者他訪問的個案.
是有靈界的.
很多人死過翻生.
全部是死過翻生.
說他們去到一個很光明.
很聖潔,很平靜的地方.
那個地方居然還見到那些已經逝世的親人.
當然他們是回到所謂的陽間.
就是地球.
回到自己的身體.
講述他們臨死之前的經歷.
不約而同他們的生命都有改變.
當然他們一樣繼續生活.
但後來又會慢慢離開世界.
但是他們的經歷.
不約而同寫成一本書.
但問題是奈何這本書.
原來我現在在第32個排隊.
找不到.
因為只有幾本.
我現在排到第32個.
所以暫時也沒有看了.
其實第八章已經講了靈界的真實.
第八章的羅馬書已經很清楚告訴我們.
我們是很正面去看罪.
看我們這個舊有的性情.
我們不再虧欠它.
因為耶穌已經幫我們付了代價.
我們現在要學習依靠耶穌.
由我們仍然未成熟.
仍然偶然間會失敗,會犯罪.

$^{161}$到我們可以慢慢.
真的能夠越來越宣告.
我們現在活著.
就好像耶穌一樣.
所以這是我們每一個弟兄姊妹.
必須要去認識的聖經的真理.
所以這裡羅馬書很有系統地.
一點都不會掩飾每個基督徒.
你怎樣可以成聖呢?.
什麼方法可以令你成聖呢?.
所以如果我們找對一個渠道.
在這條跑道上這樣走.
始終我們會達到終點.
就是在聖靈裡面.
藉著耶穌這個新生命.
我們就能夠可以剋聖我們舊有的性情.
其實要再講得清楚一點就是那個自我中心.
那個Sub-Center.
基督不以愛作為中心的人的價值觀.
其實世界的人都知道的.
他找人工作,看人面試.
都找一些品格好一點,良善一點.
就算不是信教的機構.
都是希望這樣的.
沒有人希望找一些很詭詐.
說話很玄幻,避重就輕這樣的人.
你都不會信任他在你的公司裡面.
去做這些的崗位.
現在基督徒我們就已經知道.
原來神在我們這個救恩裡面.
就是希望每一個基督徒.
都能夠認定自己是神兒女這個身份.
自己能夠在這個很有限的.
我們每個人神定了的人生.
來活著像耶穌.
所以第八章的頭這個大段落.
就是不斷再將第六章.
不同的角度來講.
怎樣可以剋聖罪惡呢?.
就是隨從聖靈.

$^{201}$就是要讓聖靈在我們的內心不斷的引導我們.
所以真的離不開我們要安靜讀經.
跟神聊天或者叫默想.
或者叫我們靈修.
或者我們思考神.
或者我們讓神的真理.
比較全面一點的神學去幫我們.
因為其實我們每一個人信不信耶穌.
總有一套我們做人的哲學的.
某程度上這個就是神學了.
每個人都有自己的神學的.
就是地球的人不信的人.
人不為己就天誅地滅.
這樣已經成為哲學,哲理了.
要改變他這個想法其實是很難的.
因為不斷不同的世界的事物.
都好像告訴你一定要這樣才能生存.
你就可以想像到.
要一個基督徒慢慢改變.
是以愛為出發點.
在基督的愛裡面願意去服侍.
願意用愛去對待其他人.
其實是難之又難的.
但是因為在聖靈裡面.
我們就有這個力量了.
所以其實保羅都是很想講.
那個正確的神學觀給我們每一個基督徒知道.
特別是當時羅馬教會.
他們既然是一班在帝國裡面的一班信主的人.
來自不同的外邦的.
但是如果他們有這個這麼具體的神學的思維.
這樣的理解.
他們將來在教會的治理.
以及傳福音去到不同的其他的國家.
例如西班牙就更加的重要.
因為很快就去到第十三,十四章.
就已經出現了羅馬教會的弟兄姊妹.
就是外邦人和信了主的猶太人.
就為了一些表面上好像很重要的東西.
其實聽起來就不是這麼重要.

$^{241}$安息日是不是真的這麼重要呢?.
是不是一定要禮拜崇拜才重要呢?.
那些祭過偶像.
或者是那些安息日.
我們是不是有些東西不應該吃呢?.
很多東西就會成為了大家衝突的那個點.
但是如果放在救恩,放在永恆,放在宣教.
這些就是次要了.
這些就是大家都要放下.
大家按照自己的良心來去堅持就行了.
所以這個就是第十四章之後.
就有這樣的教導了.
所以第八章我們曾經在上個星期講過的.
就是我們怎樣能夠轉化每個弟兄姊妹的屬靈生命.
讓他們知道自己是在恩典裡面.
神已經將聖靈在我們信著的那一刻.
已經進駐在我們內心.
祂是一路引導我們.
一路在我們裡面.
讓我們認識基督的愛.
讓我們知道我們在救恩裡面的穩妥.
以致我們慢慢可以扮演一個代求者的角色.
代求不是這麼簡單.
你揣摩神的心意.
你在神的心意裡面去為世界.
為其他的人去禱告.
為那個地方.
你在那裡聚會的地方去把脈,去祈禱.
用神的眼光.
就是說你在神裡面禱告.
這個就不是小事了.
我們本來是一個罪人.
本來求的是我們自己的自身,健康.
我們自己的財富等等.
但是現在你是提升到你在神的角度裡面.
成為一個守望者.
所以這些一定要操練的.
一定要你慢慢去發揮.
由發掘,發展.
到這樣來發揮.

$^{281}$就是一個代求者.
所以如果在這個世界.
有一千,一萬個.
是貼近神的心意去守望.
去為世界禱告.
整個世界都會有很多的不同.
所以.
盼望我們能夠真的看到.
原來我們屬靈的身份.
是這麼具體,這麼真實.
所以我們要活出神異女這個身份.
這個是上個星期我們講了一個比較重要的部分.
所以你看回這個機密就是.
我們在這個簡單的分段裡面.
不再定罪.
已經是自由了.
其實這是一個開始.
就是說我們有這個不再定罪.
但是去到更遠的八章的三十.
應該是八章的三十四字這一段.
或者早一點.
是講到我們沒有東西可以隔絕我們與神的愛.
沒有東西可以隔絕.
本來不再定罪已經是很偉大了.
對一個人來說.
但是現在基督徒應該享受的就是.
不隔絕的愛在我們的生命裡面.
不隔絕神的愛.
不會有任何的東西來隔絕.
由今天我們.
等一下我們會開始講第十八節.
萬物在那個痛苦裡面.
因為現在還有一個就坐的地球.
昨天不知道大家有沒有看新聞.
這個也不算是頭版.
在內地有一個小朋友六歲.
很有趣的樣子很乖很純的.
因為玩水.
那些淡水.
無意中可能是喝了還是怎麼樣.

$^{321}$還是有傷口.
那些吸血蟲進去腦子.
一個月ICU救不了.
就這樣死了.
六歲大.
你可以想像得到.
那個家庭的打擊.
那個傷害.
那個內心的撕裂.
是你難以想像的.
基本上可能有生之年都未必能夠填補.
你可以想像得到這些是完全不是人為性.
或者是有一些的罪帶來的.
本來就是一個這樣的世界.
這個是一個從來都告訴我們.
還不是真正永恆復興的世界.
所以我們有罪.
我們有死亡.
我們有些可能是人為性的.
例如你被人搶劫.
你抗拒被人刺了一刀.
那你就要入院.
可能你就毀容.
或者有其它東西.
這些是在一個墮落了的地球世界是會發生的.
所以就算我們相信了.
都一樣告訴我們是這樣的.
這就是18節到第26節告訴我們的.
接著告訴我們.
從這一大段.
由8章1到第16節.
都是想告訴我們.
我們在聖靈的帶領下.
我們一直在它的引導下.
是對自己神兒子名分這個身份.
是不斷有一種確立的確立.
確立是很重要的.
你用這種氣質在地上生活.
所以為什麼有時候你會看到.
有些人就算他很坎坷.

$^{361}$很困難.
但是當他清晰自己是神兒女這個身份.
他不會覺得那些苦是真的苦.
或者那些苦是不會令他意志消沉的.
那你就可以想像到.
裡面的這個神學的價值觀.
是多麼大的能力.
當然不信的人就會反過來說.
洗腦.
教會洗這些人腦.
令他們好像在地上不覺得是一個受苦.
但是如果那個人是活在一種尊嚴.
一種真理.
一種愛之下.
這個就不是洗腦了.
這個就不是喝雙氧水.
就不是一些古怪的神學了.
是他真的意識到自己是神的子女.
神是愛他.
環境他無法改變.
因為神定了一些東西.
但是環境不會令到他內心受影響.
這個就是正確的神學觀.
或者說八章一到第十六節.
救恩帶給我們心靈的光芒光彩.
就是聖靈.
如果沒有聖靈.
根本我們怎麼能夠有力量.
自己騙自己嗎?.
你怎麼有力量面對地上的困難呢?.
所以一定是聖靈在這一章裡面.
應該是第八章才開始正式出現聖靈.
大量提到聖靈在信徒內心的工作.
第五章的第六節.
只是說了一次神的愛.
藉著聖靈.
好像淋水一樣囂冠我們.
但是第八章就很詳細地說.
祂的引導.
祂是我們新的主人.

$^{401}$我們已經不再是罪的奴隸.
但是我們是一個自由人.
我們就要服在聖靈的帶領下.
我們就能夠活出這一份新生命的新生活.
我們繼續說.
第五到第八章.
其實就是正常基督徒的生活.
不正常就是我們繼續犯罪.
正常的基督徒生活.
就是開始可以學習犯罪.
因為聖靈在我們的內心.
別人罵你.
你可以不需要反駁他.
因為你裡面有愛,有寬容,有包容.
偶然間你會很生氣.
你會頂嘴.
但是慢慢你會發現.
我們裡面的聖靈會提醒我們.
學習溫柔.
學習忍耐.
學習愛.
不要以惡還惡.
這就是聖靈在我們裡面的進駐.
帶來的生命的改變.
所以這也是上一次我們比較仔細的說了.
一到第十六節.
特別最後他提到的後字.
就是將來可以分受神的產業.
與基督一同成為後子.
接著十七節就留了一個連接下文的.
就是因為十八節打後.
一直都是在說萬物受苦.
萬物在就坐下的虛空.
虛空就是說.
他本來的形態不是這樣.
現在在這樣的情況.
這是就坐的地球.
所以有天災,有人禍是正常的.
真正的永恆的榮耀.
是主耶穌再回來.

$^{441}$所以保羅的神學是有苦難的.
如果神學只有成功.
只有開心,得滿足.
每個星期你都很享受唱歌.
接著你的人生只能夠容許恩典,祝福.
其他一切就不要.
那麼你的神學就很扭曲了.
只是一些民間宗教.
沒有人會求得黃大仙,求觀音.
說拜託你給我一點災難吧.
沒有人這樣求的.
從來都不會有的.
但是基督徒我們不是要求.
而是神告訴我們.
這個還不是真正得到榮耀的救贖.
所以我們處於這個地上.
一定是有這個困難的.
但是十七節就是那個連接前了.
我們再讀一次就是十七節.
從這裡開始說起.
「幾是兒女便是後製,就是神的後製,.
和基督同作後製.」.
說了三次告訴我們可以繼承神的產業.
與基督一起繼承神的產業.
我們是一個不是普通的兒女.
是神收養我們之後.
拯救我們之後.
這個收養就是adoption.
就是神藉著他的愛子的犧牲.
將我們拯救.
然後他要賦予我們.
將來我們可以承受神的基業.
到底是怎樣的呢?.
是否將來永恆的屬靈世界.
都是物質的流放呢?.
我們不是太知道了.
說的不是太多了.
但是一而再告訴我們我們可以繼承.
所以其實我們就應該要裝備我們自己的生命.
將來我們帶著這個永恆的生命.

$^{481}$去到神的面前領受祂所賜給我們.
某程度上是我們自己那份的產業.
但是既然我們仍然在地上.
不過留意了.
後面是這樣說的.
「如果我們和他一同受苦」.
就是說耶穌.
一定是說耶穌在地上所做的.
和他現在在救恩所延續的影響.
就是說我們每一個信耶穌的.
都是跟隨耶穌在地上等候神的國度的.
等候神的國度.
我們同樣要面對敵黨神的人.
他們在地上仍然有他們生存的機會,權利.
因為神都是愛每一個人的.
所以既然現在是這樣的時候.
我們就延續基督的苦難在地上的見證.
就是說這個受苦只能夠在地上的.
將來離開世界就沒有苦了.
天堂是不會受苦的了.
所以今天是一同受苦.
所以如果我們今天不肯認基督.
不肯和基督來受苦.
那麼我們的救恩就停留在成功的神學了.
就沒有了一個苦難神學的重要一環了.
所以保羅在這裡就告訴我們.
這個就是一個連接詞.
如果我們在地上和基督一同的受苦.
接著果然就是了.
將來必定會一同得榮耀.
所以榮耀是在將來.
受苦是在現在.
這個就是基督徒在現世我們需要面對.
所以我們就不會覺得太稀奇.
為什麼在一些中東的地方.
他們受的迫害是這麼嚴重.
受很大的苦楚.
剛剛聽到Chris他們在祈禱會提出.
現在伊朗極需要聖經.
他們已經開始付運.

$^{521}$開始印刷.
土耳其是合法可以印聖經.
他們現在準備好.
到時候就會呼籲教會.
我們可以奉獻金錢.
來支持這些聖經的派發.
可以想像到伊斯蘭世界是會禁止聖經的傳遞.
所以如果你在那些世界.
在那些國家做基督徒.
你就是和基督一同的.
甚至是肉身要受苦.
不過保羅在這裡仍然告訴我們.
我們是有一個將來的榮耀.
將來是與耶穌一同得榮耀的.
所以再簡單做一個總結.
聖靈的律和罪.
死的律.
一個對照.
要我們有這一套神學的裝備.
信了耶穌.
我們要認定.
我們要有很強烈的意識.
聖靈已經在我的心裡面.
在我們眾人的心裡面.
來做一件事.
不是哄我們開心.
而是告訴我們.
我們就算有困難.
有很大的痛苦.
都可以靠主喜樂.
因為聖靈的律.
釋放了我們.
脫離了罪.
脫離了自我.
所以我們不再是一個奴僕的名分.
我們應該是神兒女的身份.
所以如果我們用這個身份來活在這個地上.
就是說我們不是雙重人格.
而是我們是一個平衡的時空.
我們既是活生生.

$^{561}$我們真的踏在這個地上.
但是我們整個的心靈.
我們的神學.
我們是活在天上的.
跟主每天這樣來起居生活的.
這樣的基督徒.
就是正常的基督徒生活.
是有主跟我們同行的.
每天的生活.
這個就是正常基督徒生活.
就是第五到第八章.
保羅不斷想引進聖靈.
就會帶我們看到這個身份.
所以到最後.
當然就是我們稍後會說的.
我們這個很重要.
救恩帶給我們.
由一個本來我們是完全不認識神.
我們被神來預知.
祂來到我們萬民當中將我們選照.
使我們能夠得救清二.
最終我們是得榮耀.
你可以想像.
這個就是每一個基督徒.
我們的人生.
應該進入永恆.
是帶著榮耀的這個身份.
進入永恆的.
所以這個就是保羅很想讓每一個時代的基督徒.
當讀羅馬書.
有系統地認識聖靈工作的時候.
就是聖靈進駐.
讓我們能夠克性罪惡.
克性我們舊有的性情的自我中心.
讓基督在我們生命裡面來作主.
這個就是第八章的一到第十七節.
所以苦難神學是重要的一環.
某程度苦難會令我們更加有一種警醒.
苦難會令我們更加沉澱.
苦難會令我們更加不是低頭.

$^{601}$而是我們會靜下來去思想人生.
我們越是願意思想人生.
越是在主的裡面.
讓祂來陪伴我們思想人生.
就有出路了.
所以原來苦難是不能夠欠缺.
對我們的生命的成長和成熟的.
所以告訴我們.
現今所受的苦與將來要向我們顯現的榮耀.
去相比做一個天平的比較.
不算不得什麼.
所以第十八節就是回應第十七節的經文的意思.
保羅說「與主同苦,與主同得榮耀」.
接著保羅就立即說了.
就開始用一個新的段落來說.
本來我們現在就已經有苦了.
保羅可能想起自己面對的.
在全道試圖這個身份.
要去建立教會的時候.
所帶來的那個逼迫,誤會,攻擊等等.
被人排擠,被人不信任.
但是保羅說現在的苦楚.
這個涉及到不只是人世艱難,搵食艱難這麼簡單.
而是相信他以試圖的身份.
在這個時候一定是提升到.
撲票又撲不到.
謝霆鋒的票買不到.
想看利物浦又買不到.
我相信不是這些表面的生活上的不方便.
而是真的在全道的過程.
四面受敵的這種尷尬.
這種內心的痛.
猶太人看你是叛國.
外邦人看你不拜偶像.
不承認君王.
又會帶來很多政治上的誤解等等.
各方都不討好.
更加要講到基督徒.
已經不應該再用舊有的價值觀去生活.
羅馬人有很多道德上的問題.

$^{641}$在保羅的時代.
有時候我們看電影.
將那些古代都列出來.
你都可以看到.
譬如亞歷山大大帝都原來是同性戀的.
你可以想像得到.
當時的南記那個時代都有這樣的情況.
羅馬的世界更加是.
你可以想像得到.
保羅既要講真理.
更加會不斷的去指出世界社會的罪惡.
但同時間他自己在不同的群體當中.
都會面對很多的困難.
所以保羅說.
結接到這裡.
他很實在.
不是一種抽象的,客觀的.
世界都是這樣的.
是在講他自己.
他現在面對的苦楚.
但如果比起將來的榮耀.
就不足介意.
保羅很清楚他自己會帶著榮耀的身份.
得冠冕,得榮耀,得稱讚的身份.
這樣進入永恆.
所以他只是五十多歲被斬頭.
何來榮耀可言呢?.
可能比耶穌更長二十多年.
受的苦更加多.
還要去很多地方.
整個當時的歐洲.
他都差不多去了.
當然當時就沒那麼麻煩.
因為地中海坐船很快,很方便.
所以其實又不是想像中要走那麼多路.
但無論如何.
你在那些地方.
你去到歐洲大陸.
那些價值觀是很抵擋真理的.
猶太人一直都追著他.

$^{681}$要來審問他.
但是保羅看到的是將來的榮耀.
所以以一個活在二千年前的保羅.
他能夠看到將來永恆的榮耀.
可以想像得到是多麼大的一個力量.
心靈的推動力.
推動他去傳道.
都不會覺得苦.
是會有那種疲倦.
但不會是苦.
所以我們要學習.
真的一定要在神的裡面.
有這種很強的屬靈的價值觀.
所以苦難神學.
必須是我們重要的一環.
我在未來的24日.
穆都安排了我回來講早課,午課.
我都定了題目.
信仰的知能.
就是苦難神學.
必須是我們在信主的人生.
必須是要一個重要的一環.
不要沒有了這一環.
否則我們的神學就會很扭曲.
所以一有一些困難.
我們就會有些人崩潰,會倒塌.
就是因為這個原因.
19節開始.
剛才我給了大家一個work檔案.
就是非定期增加保羅引咎約.
我們很久都沒有放出來.
因為由第四章到第八章.
保羅都沒有引咎約.
這裡引的都是一個邏輯的引.
那個虛空的奏座.
其實就是上一次我們講的.
《創世記》三章那段經文的.
因為阿當違背神的命令而來受奏座.
所以地就伏在虛空.
萬物都同受苦.

$^{721}$因為人的緣故.
他們很被動地來受苦.
好像躺著都中槍.
我們香港人講的一個諺語.
可以想像到.
這個創造本來是完美的.
神認為是甚好的.
但是因為人類的罪惡.
地受奏座.
阿當受奏座.
夏娃受奏座.
蛇受奏座.
人類受奏座.
這個原先的創造不是這樣.
所以我們都很肯定.
十八節這一段.
保羅就再一次告訴我們.
不是個人那麼簡單.
原來整個世界.
由阿當的犯罪.
是已經墮落.
是已經離開了神.
所以現在在這個情況下.
神其實是恢復創世之前的秩序.
屬靈的秩序的完美.
以及要引進基督帶來整個宇宙萬物.
即是說神的.
無論是屬靈界.
物質界的創造.
恢復未墮落之前的榮耀.
所以在《羅馬書》裡面.
一講到苦難.
保羅就立即.
不是個人的方不方便那麼簡單.
而是再一次.
好像說笑一樣.
今天黑雨有兩個小時.
即是7點50分.
去到11點10分.
都有八九十.

$^{761}$當是三個小時.
三個小時都有人罵.
罵到天龍台不得了.
有三個小時你就喝一杯茶.
不要那麼緊張.
都可以的.
我們有一點點不方便.
就已經有很多怨言.
但是保羅告訴我們.
那個真正的不方便.
是墮落了.
是就坐著.
我們就明白了.
我們一生下來這個世界.
其實.
雖然我們沒有統計.
不是真的個個都生得出.
有些人難產.
根本來不了這個世界.
所以其實很多東西.
我們都帶著感恩.
就算不信人他都會感恩.
但是我們基督徒的感恩.
就更加是因為我們.
更加認識在救恩裡面.
現在我們是新的創造.
所以我們就是.
在神的救恩計劃裡面.
和神來同工.
恢復這個創造墮落的光景.
去到神的榮耀.
所以其實應該說哪裡有基督徒.
哪裡應該有開心.
有平安,有盼望.
有一份的喜樂.
感染其他人.
我們基督徒應該是這樣.
所以這一段.
真的不可以沒有.
就是告訴我們世界的墮落.

$^{801}$怎樣來處理呢?.
所以這個第十八節.
留意一下是苦難.
和榮耀.
一個強烈的對比.
地上真的有苦難.
記得不是那些佛家的哲理.
是在說.
在創世記的事實.
萬物受咒狀.
所以我們要記得.
我們今天是迎向.
對這個永恆的.
再臨的.
國度的盼望.
所以這是我們基督徒.
共有的一份.
這樣的等候.
和這樣的一切的同工與神.
所以來到這裡的時候.
我們就進入第十九節開始.
聖靈在萬物.
同得榮耀之中.
祂有什麼作為呢?.
首先就是要讓我們知道.
今天我們在地上.
最重要的原動力.
或者可以說是秘訣.
就是必須與主同活.
祂如果在地上.
暫時祂的身份.
祂的救恩的工作.
是被排擠.
被抵擋.
跟隨祂的人.
我們是要面對這些苦楚的.
所以我們要預了.
成為基督徒.
其實就是魔鬼敵對的陣營的.
另外一面.

$^{841}$就是屬於聖靈的.
所以我們要記得.
我們真的需要站在真理這一邊.
接著保羅就告訴我們.
帶我們回到第十八節.
一直到第二十五節.
我們先讀一讀這段經文.
其實大家應該在不同的聚會.
可能有機會聽過.
「我想現在的苦楚」.
我用舊的和合本.
「若比起將來要顯於我們的榮耀.
就不足介意了.
受造之物.
切忘等候神的眾子顯出來」.
神的眾子就是每一個將來在天上.
回來的基督徒.
當然包含了天使.
第二十節.
「因為受造之物伏在虛空之下」.
這是墮落之後的警方現實.
那個事實.
「不是自己願意.
乃是因那叫他如此」.
就是神令萬物受咒詛.
本來萬物是不會表達的.
他不會說很辛苦.
突然嘆氣.
沒有這回事.
普羅有一種擬人法告訴我們.
唯有人能夠感受到痛苦.
感受到那種無奈.
或者悲歡離合的嘆息.
但普羅的擬人法告訴我們.
雖然萬物在虛空之下.
他來到等候二十一節.
「但受造之物仍然指望脫離敗壞的痕跡」.
敗壞的痕跡是神的咒詛.
「不享神而累自由的榮耀」.
我們現在已經有在聖靈內住的自由.

$^{881}$但是去到永恆.
沒有了我們這個肉身.
離開這個地球.
記住如果那些死過翻生的人說.
可以有一個靈體在神的愛裡面.
將來那個更加榮耀的自由就更加浩大了.
其中一個個案就是剛才我說的那本書.
第一個個案就是令到這個精神科醫生.
原來當時他剛剛在醫學院畢業.
他就要馬上.
因為當時他做Houseman.
就要在醫院的辦公室休息.
半夜突然間收到急電話.
他的呼吸器響起.
正在吃東西的.
那些人是外國人.
就知道整條領帶都是.
馬上趕去.
然後就搶救那個病人.
那個病人是自殺不遂.
就昏迷了.
經過很多的儀器的診斷.
就沒有了.
就睡在那裡病床.
然後又有另外一個床位的病人.
輾轉就和這個剛剛畢業的醫生.
就出去去另外一個房間.
就談一些他們的身體醫學的問題.
就離開了那個病房.
第二天早上那個醫生.
因為要回去寫報告.
開始要看那些病人.
去到那個昨晚已經昏迷了的病人的床邊.
醒了 沒事.
和那個醫生說話.
告訴醫生.
昨晚你和我的同房.
在另外一個房間說話.
我聽到.
醫生就震驚到不得了.

$^{921}$和那個醫生說.
你的領帶很髒.
這樣說.
他是一個受科學訓練的精神科醫生.
剛剛畢業.
他說他相信的是科學.
他不相信有靈魂靈界.
令到他不能夠不相信.
怎麼會我們在另外一個房間說話.
因為他的靈魂飄去.
聽到而已.
根本沒有人會和他說話.
是半夜來的.
他怎麼可能會這樣呢.
就開始去研究.
由他做醫生開始.
一直到他退休50年.
研究這個個案.
到底是不是有靈魂.
他的答案就是人是有靈魂的.
我就還沒有拿到這本書.
還沒有看其他的個案.
這個就很明顯第一個個案.
就令到他由一個懷疑.
到他真的知道.
不是一個意識.
而是實在的存在.
是一個世界是永恆.
是光明的.
好像一個隧道去到一個光明的地方.
所以要告訴我們.
我們要記得.
我們一定要告訴我們.
我們基督徒其實一點都不會覺得詫異.
因為第八章不斷說.
我們有心靈的世界.
舊有的人是屬於罪.
是被罪核制的.
我們現在的聖靈在我們裡面.
是帶我們進入永恆.

$^{961}$令到我們不斷修正舊有的價值觀.
能夠有神的愛.
有神的性情.
這樣的一個基督徒的正常生活.
所以正常基督徒是不犯罪的.
正常基督徒的生活.
就是一個愛神愛人的基督徒.
所以求主幫助我們.
能夠對保羅的羅馬書.
能夠有非常大的認識和渴慕.
所以我們記得這裡就是說.
神兒女自由的榮耀.
真是永恆的歸宿.
自由的榮耀.
第二節.
我們知道一切受造之物.
一同嘆息勞苦.
直到如今.
再一次說那個就坐的起源.
為什麼萬物會受就坐.
為什麼地球會受就坐.
為什麼魔鬼還會在地上去迷惑人.
為什麼還有這麼多的人.
拒絕神的真理.
拒絕承認神是他們的主宰.
所以原因就是萬物受就坐.
第23次.
不但如此.
就是我們有聖靈的基督徒.
聖靈初結果子.
我們已經嘗了天恩神愛的滋味.
感受到永恆的真實.
我們都一樣在這個強烈的改變.
不可以用人力去改變的就坐下.
所以我們一樣會嘆息.
所以這段出現了三次嘆息.
第八章.
很快給大家看看.
沒錯了.
22節開始.

$^{1001}$保羅對於嘆息的理解.
首先他說.
受造之物萬物受就坐.
在嘆息下.
然後信徒.
收集範圍.
每個人都在就坐下.
但現在我們有聖靈初結果子的23節.
信徒都是自己心裡嘆息.
等候得著兒子的名分.
乃是我們的身體得屬.
所以記得.
我們在地上.
我們不要天真地.
信了耶穌就一了百了.
一定就會立即改變.
不要受那些影響.
成功神學就是說.
我們不要理那些負面的.
我們不應該有負面的.
我們不應該有不開心的.
我們應該要成功.
應該要很開心.
我們是神的子女.
他只是說了一半.
保羅的神學裡面.
是有苦難這個部分的.
而且那個部分.
還要是萬物創世記的就坐.
那個真理的眼界.
所以如果我們沒有了這個.
萬物現在仍然在受就坐下.
而我們這一班人.
就零零碎碎.
一信了耶穌就一了百了.
就開心到底.
就可以見主.
我們變成一些自己騙自己的一班人.
我們活在一個烏托邦.
不存在的世界.

$^{1041}$如果這樣的神學就很恐怖了.
所以就不斷地.
喝水喝水就會好了.
你怎樣怎樣就會好了.
一切的苦難來到.
你都要避開它.
你都用一些錯誤的理解去解釋.
那就大件事了.
但是如果我們很正面去看.
就不擔心了.
因為保羅說不是現在.
是在我們身體得悟的兒女自由的榮耀.
所以信徒是嘆息.
一點都不奇怪.
所以苦難神學.
其實應該是我們一個.
很重要的一個真理.
所以在苦難裡面怎樣自處呢.
才是真正的神學的認識.
在苦難裡面.
我都不會覺得苦.
我能夠仍然在主裡靠主.
不要苦.
這個就是神學了.
正面的神學了.
我們就不會逃避.
甚至乎抗拒埋怨神.
給我們有苦難.
二十四節.
我們得救是在乎盼望.
只是所見的盼望不是盼望.
誰還盼望他所見的呢.
意思大概就是這個意思.
人是不會因為他所見的.
還去等候盼望.
而是我們是等待未真正全面出現的.
因為根據上下文.
是在說榮耀的自由.
身體的得悟.
那個盼望.

$^{1081}$所以還在地上.
我們就等待.
所以真的很期待.
應該沒有人這麼期待.
很期待快點閉上眼.
進入永恆是怎樣的呢?.
天使怎樣引路我去到見到主耶穌呢?.
那個稱讚呢?.
應該這樣期待.
但不是一種身無可戀.
自殺傾向.
不是這樣的.
所以那個盼望應該是這樣.
等候.
應該說.
今天我已經忠於神要我做的.
完成了.
明天如果還有的.
我就繼續忠心地去做.
不是希望突然變成揚名立萬.
而是忠心地做好我明天的工作.
神的託付.
這樣慢慢累積這種在神的角度裡的參與.
我也開始想將來退休要做什麼.
開始有一點眉目.
有一點個人的想法.
很想怎樣做什麼.
有機會和你們個別分享.
但是我們記住.
眼見的盼望就不是盼望.
所以我們對於我們有一些獎勵.
例如我的兒子升職加薪.
讓上司賞識.
去到一些很好的機會.
甚至補他去讀書.
這些是開心的事我們不需要粉飾.
我們不應該開心不應該什麼.
也不是的.
但是真真正正的永恆的盼望.
就是真正的榮耀.

$^{1121}$所以如果我們能夠在地上.
例如我們做事有上司的賞識.
被人稱讚.
而得到業界的認同.
這些是一些榮耀歸榮要給神.
是神給我們的忠心.
是神給我們知道.
我們要做一個有忠心的管家.
做好我們自己的本份.
這是一個很好的見證.
但是保羅在這裡仍然告訴我們.
一方面我們不要抗拒地上.
有時候有些很好的美食.
或者有人對我們的工作的欣賞.
或者我們見到我們自己愛的人.
得到一些榮耀.
這些是神給我們的賞賜.
我們不要抗拒.
我們不要像佛教徒那樣.
覺得世俗一切都是虛空.
不需要這樣想的.
不過我們真真正正要看到.
問題是我們更多著眼於那些.
以為是終點.
我們少了對保羅這一份.
其實真正身體的得贖.
才是真正將來主耶穌對我們的稱讚.
這才是最重要的.
如果將來他不這樣稱讚我們.
那才是真正的大事.
所以我們看羅馬書.
應該怎樣活我們的正常基督徒的生活.
那才是大事.
我們要探望將來耶穌對我們的稱許.
所以我們記得.
真正的稱讚不是短暫的目前.
那些真的不是錯的.
但如果那是一個終極.
就是我們已經夠了.
我們就有一個問題了.

$^{1161}$因為我們會擺錯了方向.
擺錯了焦點.
所以我們盼望.
我們繼續在這個二十五節.
但我們如果盼望那所不見的.
就是將來的榮耀.
就必定會忍耐等候.
所以今生我們是一個信徒.
又是一個有盼望的信徒.
又是一個忍耐等候盼望的信徒.
所以忍耐盼望.
一定是我們對永恆的價值觀.
心靈的準備.
如果我們問一些人.
原來他從來沒有想過.
要等待這個永恆.
原來基督徒正常的生活.
原來有很大的對真理.
根本從來不懂.
不懂就只能為自己認為最好的去活.
所以他的子女.
進了Band One 進了名校.
他就興高采烈.
進不了Band Three.
整個世界就倒塌了.
本來是有困難.
但是不要將焦點放在終極.
終極就是盼望未見未出現的永恆.
所以是不斷由少少少少累積的.
所以我們基督徒.
是應該不斷對真理.
對自己處於這個世界.
越來越認識.
我們就能夠達至保羅所說.
接著26節就告訴我們.
聖靈就在我們當中帶我們去看盼望.
所以第26節到30節.
就是聖靈會在他的同在參與.
在我們的生命裡面.
其中就是禱告.

$^{1201}$讓我們對永恆不斷有很強烈的意識.
所以到26節這一段.
況且我們的軟弱有聖靈幫助.
我們本不曉得當怎樣禱告.
只是聖靈親自用說不出來的嘆息.
第三次講嘆息.
萬物嘆息.
信了主有聖靈的信徒都會嘆息.
因為是一個受造的世界.
原來聖靈都有他的嘆息.
你可以想像到.
保羅刻意營造.
讓我們看到原來我們在主裡面.
我們都需要由這份嘆息.
成為我們生命的動力.
所以這裡指示聖靈親自用不能說得出來的嘆息.
替我們去禱告.
即是聖靈嘆息的祈禱.
在我們的生命裡面.
然後告訴我們監察人心的就是神.
曉得勝靈的意思.
即是曉得勝靈的嘆息.
曉得勝靈在我們心裡面嘆息的禱告.
你看多複雜.
接著說因為聖靈照著神旨意替聖徒祈求.
很感恩.
原來神賜給我們的聖靈.
永遠站在神那一邊.
希望我們的生命能夠成長.
而且永遠站在我們這一邊.
為我們去禱告.
不會放棄.
不會將我們排擠.
更加不會覺得我們無可救藥.
來將我們投閒置散.
而是聖靈會在我們裡面.
按照著神在我們生命裡面的旨意去禱告.
所以如果我們讀到這些經文.
千萬不要斷章取義.
是在說我們要盼望永恆地上的成長的困苦.

$^{1241}$或者我們在事奉裡面面對的困難.
那種困境.
那種困難.
我們原來不需要放棄.
因為有聖靈在我們裡面.
會持續不斷.
而且是按照著神的旨意替我們祈禱.
所以我們就能夠達到.
所以28節就出現了.
我們曉得萬事互相效力.
萬事互相效力.
就是因為聖靈在我們裡面的歎息禱告.
叫愛神的人得益處.
就是按他旨意被召的人.
哪些人得到益處呢?.
愛神的人.
神選召的人.
在神國裡面有聖靈禱告的基督徒.
我們就會能夠看到.
意識到萬事互相在效力.
為被祈禱的人效力.
或者為教會的時代效力.
因為在裡面說的.
保羅在收集的範圍.
更加準確地說是哪些人呢?.
第29節.
因為他預先所知道的人.
就是他預知被選召的.
就預先定下了效法他兒子的模樣.
即是學像基督,持基督.
這就是我們正常的基督徒生活.
預先知道的.
他就定下了每一個都要效法他兒子的模樣.
使他的兒子在所有被拯救的許多弟兄中.
當然包含了姐妹.
所有被主聖靈拯救的神的兒女.
耶穌就成為了長子.
即是說他有這樣的位份,形象.
和神安排的尊貴身份.
他就在所有弟兄當中作長子.

$^{1281}$然後30節.
預先所定下的人.
就呼召他們成為基督徒.
呼召的人就被稱議.
被稱議的人就叫他們得到榮耀.
所以是由預先知道,選召,信主,稱議,得救.
然後來得到榮耀.
這就是我們基督徒的人生冥冥當中.
神在萬民.
所以已不所知的時候我們都很難明白.
為甚麼保羅說他在萬民當中創立世界以前.
在基督裡面已經將我們選立呢?.
是不可能的.
保羅怎會知道這些東西呢?.
整個宇宙未出現.
神在未有創造之前已經預定了我們.
選立我們了.
所以這一定是冥冥中我們能夠信主的預定論.
告訴我們神預先已經選召我們.
然後讓我們出現在這個世界.
去到某年某月某日我們就信主成為基督徒.
記得那是起點.
在我們救恩開始是起點.
在神的角度裡面.
他已經很早很早創立世界以前.
就是他的起義.
但是來到我們這裡.
由信主那一刻就是一個起點.
但是記得.
保羅不斷說的是榮耀才是我們的天花板.
我們的司令.
所以我們信了耶穌.
洗禮不是我們的天花板.
8月24日就洗禮了.
洗完禮就很多人說沒空了.
不成長了.
不參加了.
慢慢可能他只是維持崇拜.
其實不足夠.
當然每個人真的有他的生活.

$^{1321}$有他的忙碌.
但一定要自己有意識地調整.
你用什麼方法都不要緊.
你聽什麼人.
總之令你的生命成長成熟就行了.
要令你聽得明白.
和你願意實踐.
那才是真正的牧者的牧羊.
你聽了很多東西.
你都不願意走出.
你不能體驗到.
你自己實踐過程的困難.
因為別人說完是別人的經驗.
別人說的經驗是希望你聽完.
你受感動.
你自己開始走.
你才體會到原來做基督徒.
要面對世界是這麼難的.
所以記住我們是.
「清二得榮耀」.
是一個成長的階梯.
所以今天的基督教.
只是說開心,成功,快樂,平安.
那就是一個未完整的神學.
應該包含了.
由「清二」.
其實是由自己為什麼被耶穌選召.
到你意識自己得救.
很清楚你自己有一刻得救.
你是神國度的兒女.
到你一路去成長.
神安排你的人生.
不同的場景.
你去體驗.
行真理的成長.
到你真的將來得到榮耀.
這就是我們基督徒的一個人生.
這就是《羅馬書》第八章.
不斷說的就是.
正常基督徒的生活.

$^{1361}$就是一個「清二得榮耀」的過程.
所以聖靈的代禱.
和萬事你意識到的效力.
是有絕對的關係.
所以聖靈是令到我們.
能夠不斷在祂的祈禱裡.
來去符合朝向神的旨意.
這樣來活.
所以其實是很奇妙的發言之旅.
自己已經三十多年由信主到現在.
傳道.
因為今年我們區聯會每年都會給一些.
五年,十年,十五年,二十年的牧者.
宣道給的牧者.
我們每年都會有一個晚宴.
頒個老人獎給你.
我們今年就是十五年.
剛剛今天填了表格.
一定要出席.
拿那個阿加利膠的獎牌.
當然不是金牌.
是一個膠牌.
五年,十年,十五年才會有的.
當然二十年,三十年就更加資深.
我就是十五年了.
所以在地上走這條路.
一條信心的,服侍的路.
其實都是很奇妙的.
能夠讀神學.
能夠將神的話語去傳遞.
自己去研究,看書.
深入淺出去講解.
現在要用科技,用PowerPoint去做.
不斷去消化.
怎樣能夠講得更加明白.
整件事是很奇妙的.
在地球上有一個職業是傳道人.
你可以想像得到.
這個就是神的一個不同的計劃.
在一個人身上.

$^{1401}$所以第八章這一大段.
二十六到三十節.
很明顯講到聖靈.
肯定在每一個信徒的內心.
為你能夠符合.
能夠行神的旨意去禱告.
所以其實是很奇妙的.
所以我們可以沉澱我們自己.
因為我們已經有足夠人生經驗了.
不再是二十多歲.
剛剛初出茅廬.
氣沖沖地去遇到.
我們已經沉澱了.
我們是怎樣的人.
我們的品格.
到底我們是一個很快,很慢.
慢條斯理,很急躁的人.
其實我們已經很清楚了.
我們應該沉澱.
我們在聖靈裡面去思想.
聖靈一直為我這個人禱告.
讓我能夠體驗.
由清二到榮耀.
其實是一個很大的.
真正的平安之旅.
一個旅程.
這樣就是讓我們.
所以為甚麼去到這個年紀的基督徒.
憤外對培靈會有感動.
因為你真的對人生很明白.
對很多事情.
你已經算是某些.
很多背後的哲理.
都已經很認識了.
當然不是食鹽多過食米那些說話.
而是真的對人生都有一種味道.
我們品嘗了.
我們發現原來在聖靈裡面才是最好的.
才是真真正正.
他完全了解,認識我這個人.

$^{1441}$為我這個人禱告.
而且全力要我能夠在神的愛裡面.
在神的國度裡面.
有誰會這樣愛我們呢?.
真是唯有聖靈.
所以原來這一段經文.
是要告訴我們.
原來在神的創造計劃.
是有我們份的.
你可以想像.
譬如我突然間想到.
新約聖經有一個人.
耶穌說他是滅亡之子.
一聽就知道是猶大.
他本來.
耶穌說.
他本來列在.
應該是施提反說的.
他本來列在十二使徒的遺訓裡面.
但是耶穌說.
他從起初就已經墜落了.
他是滅亡之子.
這個是.
《約翰福音》的十四章的經文.
還是十三章?我回去找找.
除了那個滅亡之子.
沒有一個失落.
你可以想像.
對不起,是《約翰福音》的十七章.
滅亡之子.
他的出現就是去到永恆的沉淪.
猶大是跟隨耶穌.
你可以想像.
更多在世界出現到死.
無論有什麼豐功偉績.
如果他不是在耶穌.
在聖靈的救恩裡面.
對自己有這個強烈的.
清義的意識.
成長,愛耶穌,愛教會.

$^{1481}$愛傳福音.
這個意識其實就這樣.
枉過了人生.
因為那些時間.
好像長江的滾滾流水.
將人一代一代的沖走.
我們很快就會過去.
但是如果我們意識.
這一段經文26到30節.
原來聖靈.
由我們信的那一刻已經在我們裡面.
就是這樣不斷持續.
為我們禱告.
萬事.
為我們這個日效力.
我們就會感受到.
我們真是.
越信主越有平安.
沒有人可以搶奪.
我們對這個盼望的.
永恆的來臨.
是很實在.
會有這個把握.
因為.
眼見的就不是盼望.
誰會盼望那些見到的呢?.
但是我們如果盼望.
那個未見的.
就必定會忍耐等候.
所以基督徒的美德.
忍耐等候.
原來是很好的美德.
雖然很多人說我吃東西又快.
做事又快.
但其實我是忍耐等候.
我們真是有這樣的.
屬靈的價值觀.
成為了我們一些很重要的神學的.
意識.
所以其實讀羅馬書.

$^{1521}$是可以有很準確的神學.
對地上我們明白.
對永恆我們更加知道.
是會出現.
所以如果在一個.
現在失序的世界.
失去秩序的世界.
我們就真是扮演著.
代求者的.
這個位份.
所以我經常都呼籲頂子妹.
我們星期三晚的公土會.
無論你是現在在屋村的一個婆婆.
或者你在哪個地方.
在線.
我們就是神的兒女.
我們就代表著神.
為這個世界去把脈.
好像特朗普.
因為有姐妹提.
星期六.
這個重要的高峰會.
特朗普和普京的對話.
閉門的.
有姐妹提出在公土會祈禱.
當然很好.
我們真是扮演著.
這樣的代求者的.
這個位份.
那半個小時其實是很重要.
我們在行駛著.
我們神的兒女的這個位份.
所以我們要記得.
這裡就是說我們在這一段經文裡面.
最顯眼當然是.
八章二十八節.
本來認定原來萬事.
真是會互相效力.
其實不是那些很意外.
感恩啊我今天上不到那班飛機.

$^{1561}$我不用死啊.
對的不要輕或錯的.
但如果我們太強調.
我們是不信的人.
而他們犧牲了自己的家人.
本來他們是無辜的.
你的基督徒神保佑你.
我們就全部死了.
所以我們只能夠說.
我們盡量不要太過強烈的去感恩.
應該說.
我們要記住萬事互相效力.
背後有什麼的作為.
我們真的不知道.
因為我們不是神.
所以我們反而要為那些.
無辜.
不幸的我們去祈禱.
我們不能夠說.
那個上不到飛機的.
我們就要不斷的為他感恩.
不是錯的感恩.
但有更多人是不幸的.
我們應該要留意.
這些這樣的.
崩口人忌崩口滿.
某程度上都應該是.
但是我們基督徒.
就應該不要幸災樂禍.
但反而我們要有一種悲天憤人的心.
對其他人的苦難.
我們應該要有一種感同身受.
所以最後這一段.
你看看這個部分.
最後.
在基督裡面.
由開始很強烈地.
在基督裡面.
不再定罪了.
我們是被釋放.

$^{1601}$那個釋放就是當時的奴隸制度.
跟著第七章.
說的是取死身體.
如果真的那些.
行死刑的人.
有一條屍體綁住他.
慢慢被那些蟲蛀死他.
那取死的身體.
就是罪性在我們裡面.
我們現在正常.
就是我們在基督裡面.
不再定罪.
沒有了這個取死身體.
而是我們能夠自由.
我們釋放我們有新的主人.
就是聖靈.
愛我們為我們祈禱的聖靈.
所以最後這一段.
保羅的讚歌.
31到39節.
相信很多基督徒都很喜歡.
最高峰.
就是保羅好像大包圍.
在基督裡面.
沒有東西可以隔絕.
我們跟神之間.
這一份親密的愛.
神愛我們.
是永遠不會改變.
就是說我們要有一個很正確的神學.
就算苦難.
就算我已經在苦難的氛圍.
那個時間裡面.
但是我雖然苦.
但是我不會感覺絕望.
因為在基督裡面.
沒有東西可以隔絕我們跟神的.
就是頭尾第一節.
和31節.
這樣的一個括弧.

$^{1641}$一定是中間.
整章裡面.
聖靈在我們生命裡面.
的那個.
那個結合.
在我們心裡面.
生命裡面的那個.
的介入.
的幫助.
所以第31節.
可能讀一讀我們就差不多完結.
既是這樣.
就是說又是一個.
上文說的.
原來每一個基督徒.
神安排我們是要得榮耀的.
如果是這樣的時候.
那個計劃原來是這樣.
才是真正滿足的.
那就是怎樣呢?.
就是說神如果幫助我們.
神如果站在我們這邊.
英文和原文.
不是幫助.
不是只是祂過來幫助.
而是一直站在我們這邊.
God is for us.
一直都for我們的時候.
誰能夠抵擋我們呢?.
神既不愛惜自己的兒子.
為我們眾人捨了.
豈不也把萬物和他.
就是耶穌呢?.
一同白白賜給我們.
當然這裡說的萬物.
就不是令我們飛黃騰達.
不是令我們一世無憂.
好像我們拜觀音.
這樣的民間宗教.
所謂的萬物.

$^{1681}$就是要讓我們能夠.
在萬事的效力裡面.
我們的生命的成長成熟.
還有來清楚自己的恩賜.
我們的生命的成長.
還有來清楚自己的恩賜.
位份,崗位.
真是的.
要發現自己在什麼崗位.
神正在用你的地方.
去被神用.
如果你一直都沒有自己的崗位.
沒有自己的最愛.
神給你的恩賜.
你在那裡發揮是自如的.
就是你一個田徑100米的選手.
叫你去奧運游泳.
100米是不可能的.
因為都不是那個場景.
所以每一個基督徒.
應該你慢慢發現.
自己神安排你哪一方面.
你的崗位,你的恩賜.
就在那裡盡心盡力地去服侍.
這個就是我們要這樣去做.
神既然不愛惜自己的兒子.
就是他獨生的兒子.
就是在創造的永恆裡面.
神的這位耶穌基督的聖子.
所以我們很難去明白這些奧秘.
為什麼神要犧牲他自己的兒子呢?.
聖子本身就是神聖的神.
居然由天上來到地上,地球.
住在以色列.
在那裡出生,在那裡成長,在那裡傳道.
在那裡被自己的同胞釘十字架.
但是這位是造物的主耶穌.
所以神不愛惜他自己的兒子.
為我們眾人寫了.
「豈不也把萬物和他一同百百的次給我們?」.

$^{1721}$這個不是一個設文.
是一個告訴我們.
現在就是這樣發生的了.
所以我們要記得.
神一定不吝惜對我們的生命的成長.
因為救恩都是他的愛子成就.
不是天使,不是普通人.
是他的愛子犧牲的血成就的.
所以對於我們的生命的成長.
其實我們要有這個信心.
他一定會這樣來幫助我們.
所以保羅很認識自己.
好像擁有萬物給他去使用的權利.
因為他是使徒.
因為他是神的使徒.
所以他在哥林多後書.
他說他自己有三層天的經歷.
真的好像他的靈魂去到那個樂園一樣.
他自己這樣描述.
十二章第九節那段經文.
第八節哥林多後書的十二章.
你可以想像得到.
原來他也不是死了,翻生.
他真的有三層天的經歷.
第三十三節.
誰可以控告神所揀選的人呢?.
這樣提出的一個疑問.
其實就有答案了.
有神稱他們為異.
即是說不能控告了.
因為神已經揀選了我們.
揀選了我們之後.
是要令我們是清而得救的基督徒.
很有意識地知道自己是得救.
自己是屬於神.
已經是進入永恆行列的神異類.
這樣就叫做清而了.
誰可以控告呢?.
沒有人可以控告.
因為神已經將我們清而了.

$^{1761}$接著三十四節.
誰能夠定他們的罪呢?.
是否真的可以呢?.
是的.
但是基督耶穌已經袋贖了.
他已經死了.
而且是由死裡復活.
所以是一個套餐.
降生,死,埋葬,復活和永活.
是我們救恩的一個必需.
所以我們今天信的是一位復活了的耶穌.
在天上面待求的耶穌.
所以我們的救恩就有力量.
就是一種屬靈的力量.
因為我們不是信一個死了的教主.
我們是相信一位復活的耶穌基督.
所以我們在傳揚的就是一位復活的耶穌基督.
接著說:言今在神的右邊也替我們祈求.
這就是聖經很清楚地告訴我們.
現在耶穌就在天上的永恆.
神的寶在右邊為我們祈禱.
接著第三十五節.
誰能使我們與基督的愛隔絕呢?.
接著一長串的假設的苦難.
很有趣的是.
有沒有東西可以使我們與基督的愛隔絕呢?.
打麻將,跑馬,吃自助餐,去北海道.
或者其他地方買一部超級跑車,法拉利.
很少說這些開心的事情.
全部說的苦難.
因為我們從來沒有一個苦難神學在我們的意識.
所以我們認為一定要開心.
一定要滿足.
我們才能與神的愛隔絕.
但是保羅就告訴我們.
難度的他在數.
例如難度是患難.
即是各式各樣的患難.
困苦.
困苦可以是他試過要絕食,禁食.

$^{1801}$又被人囚禁.
被人誤會來打.
在菲律賓監獄被玉卒帶領信主.
以為那些人都走了.
後來那個玉卒就打他.
因為他之前是那個神婆.
可以講一些預言的.
後來就失效了.
那個主人就告他.
然後就抓了他.
那些禁卒就打他.
打了很多棍子.
打了保羅和西拉.
在菲律賓監獄.
所以這些是困苦.
有逼迫,有飢餓,有赤身怒體.
有危險,有刀劍.
然後他真的引聖經三十六節.
如經上所記.
他終於引回這些舊約的聖經.
就是詩篇的四十四篇.
我們為你的緣故終日被殺.
人看我們如張載的藥.
所以你想想那個意境.
張載的藥.
代載的羔羊.
記住我是羔羊.
我們不是宰狗.
宰狗就是和你死過.
你抓一隻貓.
貓都會抓你.
羊是從來不掙扎的.
張載的羊.
所以基督徒有時候我們要記得.
如果有些苦難是我們要鳴定.
我們就接受.
從容的接受.
其實要學習從容都不是容易.
但是詩人已經說.
我們為你的緣故終日被殺.

$^{1841}$就是說他暴露在那種危險裡面.
就是一種可能沒有真理的氛圍裡面.
他要面對那些人對他的針對迫迫.
人看我們如張載的藥.
所以都是可圈可點.
我們基督徒應該都是.
不是盲目的迷途羔羊.
或者跟風的迷羊.
我們是好像張載的羊.
我們都不介意.
如果為福音.
我們為神.
因為有神在我們這一邊.
所以《保羅人》真是44篇的22節.
大家回去可以看回這篇詩篇.
詩人在什麼意境裡面.
他講自己是張載的羊呢?.
但是言而喻.
靠著愛我們的主.
在這一切.
即是上文所講的例子的這些苦難.
在這一切的事上已經得勝有餘.
所以我們記得不是沒有了這些苦難.
而是在這些苦難裡面.
我們靠著可以在那個困境裡面.
仍然有盼望.
仍然有平安.
仍然有喜樂.
叫做得勝有餘.
就不是要擺過.
不要.
苦難這張牌我不要.
就不是這樣.
不能不要的.
有時在教會裡面都會被人誤會.
記得我們要在主裡面有這份信心.
所以得勝有餘.
就不是沒有苦難.
而是在苦難當中.
我們有喜樂.

$^{1881}$有平安.
有主的同在.
和有主的稱讚.
因為我深信.
無論是死是生是天是.
是掌權是有能是現在是將來.
是高處是低處.
是別的受造之物.
都不能叫我們與神的愛隔絕.
這愛是在我們的主基督耶穌裡.
第一節就已經講了.
在基督耶穌裡不再定罪.
現在第31節.
在基督耶穌裡沒有任何的東西.
可以使我們與神的愛隔絕.
所以一開始就反映了.
這個很重要正常基督徒的生活.
我們今晚就講到這裡.
我們一起祈禱.
再一次我們多謝主.
讓我們都算是用了好幾堂.
將第五到第八章.
正常的基督徒生活.
保羅所告訴我們的.
我們盡力去講.
就是要我們在主的救恩裡面.
克制罪惡.
克性我們的自我中心.
讓主讓聖靈來帶領居守衛.
成為真正的屬靈的自由的一個基督徒.
盼望我們能夠在生活裡面.
去開始實踐.
和更加重要去發現.
去發掘我們的恩賜.
我們在神的眼中的定位.
分派的安排的崗位.
去盡心盡力的服侍.
我們這樣感恩禱告.
奉耶穌基督的名.
阿們.

\newpage



\section{羅馬書}
\label{sec:3_mKpJn_jtA}
\textbf{2025-08-21 周四晚查經班《羅馬書(20):分題:神的選民猶太人可得救嗎?》 講員:曾裔貴牧師}
\newline
\newline
連結: \href{https://youtube.com/watch?v=3-mKpJn-jtA}{\texttt{https://youtube.com/watch?v=3-mKpJn-jtA}} ~~~~ 語音日期: 2025-08-24
\newline
\newline
\hyperref[sec:HwCOo8EiKp0]{\small{< < < PREV SERMON < < <}}
~
\hyperref[sec:index_chronic]{\small{[返順時目]}}
~
\hyperref[sec:index_scriptual]{\small{[返順卷目]}}
~
\hyperref[sec:MB59D1BA0Dc]{\small{> > > NEXT SERMON > > >}}
\newline
\newline
$^{1}$我感恩今晚又再一次去查羅馬書.
今天會開第九章這個課題.
我們一起先祈禱.
天父我們感恩.
讓我們都能夠平安的過了一個星期.
我們又可以能夠在工作上來服侍你.
榮耀你.
也都讓我們能夠在救恩裡面.
我們不斷在主的愛和恩典底下.
這樣來生活.
我們很盼望我們在經文一直進行的時候.
保羅的思路成為我們思考的模式.
他現在突然轉向面對一大群猶太人.
其實大部份是不得救.
應該怎樣處理他們的神學.
都是讓我們能夠有一些發現.
盼望透過這三章的九至十一章.
都讓我們能夠再一次從更加宏觀的去看.
這個恩信清義和神的揀選.
和神定義的私恩到底那個銜接是怎樣.
求你帶領我們.
因為我們都有很多舊約的背景.
盼望我們都能夠有一個掌握.
使我們更加明白.
我們在你面前等候禱告.
奉耶穌基督的名 阿們.
八章的結尾.
我先開小一點的這個方案.
我們上一個星期就很清楚的說了.
第八章保羅告訴我們.
其實他沒有特別針對猶太人或者是外邦人的.
而是在聖靈的一路的工作底下.
人很自然就會看到在耶穌裡面.
祂的救恩 祂的犧牲.
就帶給我們有這個得救的清義.
再加上聖靈會令到我們感受到.
我們和神的關係 一份的親密.
我們多數是說外邦人.
就可以來去稱呼神做阿花姑.
對猶太人來說是一件很大亦不道的事.

$^{41}$因為十誡的第三誡就已經禁止.
我們不可以妄稱神的名.
妄稱神的名就會被刑罰的.
所以保羅這樣來去呼喚弟兄姊妹.
外邦的弟兄姊妹.
可以來去這樣稱呼叫做阿爸父神.
這樣的阿難民的說話.
再加上我們在聖靈的一路的引導底下.
我們明白世界的世情和人.
生活在這個世界.
是真的會有生老病死的困苦.
就是說信了主.
都不要以為是順順利利,平平安安,上天堂.
所指的不是心靈的平安.
而是生活要每一個環節都是盡善盡美.
都是要相安無事.
都是要什麼有什麼.
這樣就不是真正的信仰.
保羅說主耶穌在這個世界都是受苦.
祂將來才得著真正的榮耀.
當然祂因為已經完成救恩.
祂就已經回到上天上.
但是今天七十多億人裡面.
有多少人會來去尊耶穌為主呢?.
很少.
只有基督徒.
或者一些可能是東正教或者天主教等等.
但是絕大部分人根本就不認識耶穌.
只會當祂是一個歷史出現的人.
所以不會當祂是神的.
你可以想像到.
保羅在這裡很決絕地告訴我們.
最後的經文就是說.
我們在主裡面.
沒有東西可以使我們和神的關係斷絕的.
我們就能夠有一個很穩妥的把握.
我們就是屬於主的人.
就是說.
苦難,困難只是讓我們生命更加成長而已.
不會使我們和神的愛隔絕的.

$^{81}$這是很重要的一個神學.
但是突然之間.
保羅就很想告訴我們.
如果不是因順清而這條路.
這條救恩的路去信主的人.
無論他是什麼人.
都不可以得救.
都不能夠得救.
所以你可以說.
似乎是在說另外一個課題.
探討猶太人為什麼信不了耶穌.
但是其實一直在說的是.
救恩到底是怎樣的呢?.
什麼是救恩呢?.
怎樣能夠得救.
得救之後.
怎樣過一個正常的基督徒的生活呢?.
但是來到第九章.
一到第五節.
其實是一個很好的銜接.
因為第八章都已經說了.
沒有東西,沒有人.
沒有任何的事物.
在一個跟隨耶穌的人身上.
可以將他帶離開神的愛.
就是說如果我們很堅定的在耶穌裡面.
就沒有東西會使我們失去了神的愛和同在.
就是說永恆的救恩.
但是第九章.
第一到第五節.
保羅就將他自己心靈的感受.
放在第八章的結尾這裡.
他說.
可能如果為了其他人的救恩的緣故.
他自己願意與基督分離.
很特別的.
第一節.
我在基督裡說真話.
並不妨言.
這是舊的和合本的翻譯.

$^{121}$有我良心被聖靈感動.
給我作見證.
我是大有憂愁.
心裡時常的傷痛.
為我弟兄.
我骨肉之親.
就是猶太的同胞.
就是自己被咒撞.
與基督分離.
我也是願意.
剛剛才說沒有東西可以使我們與基督的愛隔絕.
沒有東西使我們能夠.
與基督的愛隔絕.
但是這裡告訴我們.
如果要猶太人得救.
如果要我被咒詛.
如果要我保羅與基督分開.
他說他也願意.
你可以想像到這樣的落差.
保羅不是說一些.
好像說了也沒說過的話.
言語的藝術.
保羅不是這個意思.
而是他真心願意.
猶太人.
當然在整個三章裡面.
猶太人不要再用自己的辦法去得救.
不要以為自己的身份可以得救.
身份不行.
所做的不行.
就帶來只會步向滅亡的光景.
所以保羅就說.
他願意與基督分離.
自己被咒詛.
使得他的弟兄.
他的骨肉之親.
來信主.
或者說明白救恩救贖的神學.
他就可以恩信稱義了.
但是就是說.

$^{161}$猶太人現在不肯恩信稱義.
這就是整個三章.
九至十一章.
保羅想要來講述的.
接著保羅就說到.
什麼讓他們不肯.
或者不認為需要恩信稱義呢?.
原因就是.
第四節開始.
「他們是以色列人.
那神兒子的名分.
榮耀.
諸若.
榮耀就是神的榮耀.
向這個民族不斷的顯現.
例如在出埃及的時候.
致勝所是有光.
曾經試過摩西.
離開了那個會幕出來的時候.
他的面容發光.
那些人怕見到他.
那些人想了辦法.
就是戴一個小手巾.
那時候沒有口罩.
沒有豬嘴的口罩.
就是用一個小手巾.
遮住整個面.
因為他有神的榮光.
從那個會幕出來.
這是摩西.
所以有神的榮耀.
有諸若.
有不同的若.
神與以色列人.
納了不同的若.
那這些若.
一直都是這樣對著.
他們的後代子孫.
來去應許和祝福的.
這就是普羅在說的.

$^{201}$他想到的諸若可能就是.
西賴山的若.
神與大衛所納的若.
與亞伯拉罕所納的若.
堅定與以撒所納的若.
接著後代的人也有不同的應許.
例如耶利米.
神也與他納若.
70年之後.
這個被擄的民族.
也會回來.
耶利米你可能見不到那天.
但是他的應許.
神的若是不會廢棄的.
你可以想像到.
在普羅寫這封信的時間.
他很清楚.
猶太人是一群與神納若的子民.
接著有律法.
有禮儀.
就是摩西的五經.
就是有其他十誡衍生的一些條例.
接著還有禮儀.
就是說他們有一個規模.
那個規模就是.
安息日.
他們就要大家一起去聖殿.
或者當時的舊若的時代.
就是未有聖殿之前的會幕.
大家都要圍聚那些男丁.
這樣去守節.
去朝見神.
去敬拜神.
另外又定了一些節日.
節日就是代表著休息.
無論你住在哪裡.
都要上耶路撒冷守節的.
這些就是各種的禮儀.
猶太人就是世世代代遵守的.
很小的第八天就要行割禮了.

$^{241}$猶太人他們會有這些禮儀.
來代表著不斷在他們生活的每一個層面.
都是跟神有接觸.
有關聯的.
所以如果小朋友從小在教會長大.
他參加不同的聚會.
兒童崇拜.
兒童主日學.
夏天有不同的營會.
他的思維不斷讓基督教的思想.
神學.
聖經.
詩歌.
來在他的心靈的價值觀塑造.
是很重要的一個成長的元素.
就是有神觀念的小朋友.
跟一個沒有神觀念的.
無論基層或者富貴的.
都是不同的.
有神.
他有一個敬畏的心.
這個就是很重要的.
這些就是透過不同的禮儀去塑造的.
所以每個星期一起背《主禱文》.
小朋友一起去唱一些聖詩.
或者唱關於奉獻的詩歌.
都是一個慢慢潛移默化.
對他的生命的成長.
所以如果你不帶小朋友回禮拜堂.
回到主日學.
他就沒有了這個成長的.
因為不到他去選擇.
他就被強迫地獲得這種觀感.
和這樣的教育.
這些禮儀不知不覺就在他的內心.
慢慢地成長.
聖靈就會用他認識的神學.
當然是一輪半抓地不斷地加增上去.
他就會不斷地越來越認識.
猶太人一樣.

$^{281}$一出生他就在這個氛圍底下.
這些食物是不准吃的.
那個房屋有發霉.
就要找祭司來檢查.
接著生鱔.
因為現在用那些je 喱 je 喱頭.
有時候中途會有頭瘡的.
查一些東西就沒事了.
在舊約你就要找祭司了.
你是什麼的鱔.
如果是生的滿了.
就要被禁閉七天.
七天之後就要找祭司再驗.
這個利未記的第十三章十四章就有教導了.
所以在整個生活的每一個部分.
來跟神有關係.
所以保羅都肯定了.
這些事情其實就應該更加讓一個猶太人.
更加認識神的存在.
神期望他的子民有信心.
但是很可惜.
他們有了這些應許.
保羅說全部都是他們的.
但是第五節.
列祖就是他們的祖宗.
列祖就是肉身的先祖.
就是亞伯拉罕 以撒 雅各 約瑟.
然後約瑟之後就會有其他的.
例如以伯年 馬來西.
一直下去.
當然有猶太派那邊的族譜.
列祖就是他們的祖宗.
按照肉體來說.
肉身來說.
基督都是由猶太人所出的.
是約瑟的兒子.
這個木匠約瑟的兒子.
馬尼亞所生的.
真的生在這個地球.
真的來這個世界做人的.

$^{321}$這位基督耶穌降世.
說要拯救罪人.
保羅說可信可佩服.
原來連救主都是由猶太人所生出來的.
所以按道理來說.
這是尼塞那在舊約的應許.
神應許將來會有救世主.
尼塞那.
所以猶太人如果熟讀舊約的聖經.
按道理應該有一個時候.
等待尼塞那.
尤其是以賽亞書的七章.
必有童女,懷孕,生子.
人要稱他的名字.
叫做「以馬內利」.
七章的十四節以賽亞書.
所以耶穌基督都是由猶太人這個民族所出來的.
當然保羅再一次說到關於耶穌的肉身.
他立即補充.
這位耶穌不僅是地上這個短暫的身份.
其實他真正的身份.
他是萬有之上永遠可稱頌的神.
真是不得了.
說說耶穌降生的位置與猶太人有關.
即是說你們應該更加近水流台.
你應該更加可以得救.
即是說保羅的意思是.
他自己也有一個反省.
即是說我們這些熟悉舊約.
一直在這樣的氛圍成長的民族.
居然與救主耶穌隔開很遠.
甚至本來已經來到我們當中.
好像《約翰福音》一章十二節.
他到自己的地方來.
即是來到這個地球.
自己的人倒不接待他.
凡是接待他的.
就是相信他明的人.
耶穌基督這個名字就有權力.
有權柄使我們能夠成為神的兒女.

$^{361}$所以你可以想像到.
保羅其實很慨嘆.
其實我們這個民族是很應該來信主的.
即是說如果你讀黃人的.
你居然考到尾.
居然完全DSE科科都拿U.
你讀黃人的.
你現在不是天水圍那些Band 5的學校.
除非你的智力本身有問題.
如果不是 合格根本就不是問題.
進大學根本就沒有難度.
除非你說要全部拿五星星.
一般來說就是這樣.
你可以想像到.
保羅真的很痛心.
所以他現在很用心地.
將他所認識的真理.
他現在的眼睛開了.
看到原來如果我們再繼續在猶太教.
這個救恩的門路.
這樣的身份認同.
就不能夠得救了.
所以後面就正式陸續解釋.
原來是有在救恩裡面的主權.
揀選和神的旨意.
這就令到我們好像覺得.
跟之前說的不同.
是人信就行了.
原來在神的另一個角度裡面.
在救恩的角度裡面.
原來不得救的人.
原來又是有他們自己的問題之餘.
原來又牽涉到神的揀選.
所以去到第九十章.
很平衡地第九章講揀選.
第十章就講自由意志的相信.
兩者其實完全沒有衝突.
在保羅的神學裡面.
在我們人就要有一個很大的相對.
是揀選就不用我們信.

$^{401}$是我們信就不是揀選.
我們有時候就會有這樣的矛盾.
但是在神在保羅的眼中其實沒有.
不過保羅用了五節聖經.
就將他內心的傷痛.
由他自己都很硬心的去攻擊逼迫基督徒.
逼迫耶穌.
逼迫相信跟隨耶穌的人.
到現在他很想成為一個代言人.
用心希望拉到一個只有一個的同胞上岸.
救恩的岸.
好像現在就是洪水.
你在那個水裡面.
你遲早是會被淹沒.
但是現在他能夠救一個只有一個.
所以他就將整個的神學.
仍然都是因順清義.
被一些先入為主的神學去攔阻.
而這個先入為主的神學.
就是猶太人好像一個示範那樣.
所以我們也都要求主幫助.
我們信了主.
我們也要求主給我們一個謙卑的心靈.
他看到一切都是恩典的時候.
我們就要憤外的恩勤.
去幫助未信.
不知道或者故意不因順清義的人.
來得到神的憐憫和救恩.
因為我們信了主.
再希望叫人信.
其實是很難的.
因為不信的人.
他看不到突然之間那個聖靈的感動.
那個部份.
不是他聽不明白.
而是沒有那個聖靈的感動.
那個跳躍式的相信.
那個因順清義.
有時候真的需要某一些的時機.
某一些突發的事件.

$^{441}$就會令到那個人.
有一個好像火花.
好像心靈裡面遇見聖靈.
所以這個.
就是我們要去做.
這個就是第一到第五節.
保羅告訴我們.
他的內心的傷痛.
猶太人歸主.
猶太人是不是真的可以得救呢?.
答案是.
沒錯是的.
但是不是因為他的血統.
不是他的身份.
與生俱來的猶太人的身份.
而是因為信心.
即是說猶太人.
原來一直拿著聖經.
是沒有聖經.
帶來的得救的信心.
即是說讀舊約的聖經.
他們固然拒絕新約.
讀舊約的聖經.
是讀不到信心.
所以就不能夠因順清義.
也都不容讓聖靈.
在他們的內心.
軟化改變他們.
這個就是讀經的障礙.
一打開舊約的聖經.
他們的卷軸.
就已經.
心靈被蒙蔽.
讀的是不斷.
猶太教的.
氛圍裡面的.
那些手續.
那些律法等等.
不能夠帶到他們.
跳躍的去信靠主.

$^{481}$我們基督徒.
新約我們就.
沒有了那個包袱.
但我們都有我們小心的地方.
就是我們太容易信.
我們又變了不珍惜.
我們因順清義.
帶給我們的那個跳躍.
如果不是主恩典的憐憫.
我們根本就不能夠.
跳躍地信了主.
而得到這份的救恩.
所以我們常常保持.
這份所謂初心.
就是這個原因.
第六節開始.
我們就小小的分析.
第一到第五節.
保羅為他的同胞憂傷.
自己被咒助.
他不是.
講一些言語的藝術.
明明剛剛第八章.
才說沒有任何東西.
可以使我們在耶穌裡面.
與神的愛隔絕.
保羅說.
如果為了未信的同胞.
我願意被咒助.
與基督分離.
我都願意.
所以不是說.
他真的很想這樣做.
就是說.
他自己成為那個橋樑.
來將耶穌.
與這個.
失落了的.
那個環.
能夠扣上.

$^{521}$讓他的同胞.
可以踏上這條恩信的路.
去遇見耶穌.
所以.
如果我們為著耶穌.
好像有些大家給錢的免禮.
為著耶穌.
我們不要計較.
為著耶穌.
我們可以奉獻一切.
當然我們要有理智.
但是如果真是為著別人的信主.
我們多走一步.
學習多走一步.
我們的靈命.
就會在這樣的過程.
是會有突飛猛進.
我們不要停留在一個頭腦的知識.
是呀 我都很想帶人信主.
但是都沒有具體的行動.
就不是保羅所說的.
他願意為同胞.
來憂傷.
所以他要排除很多的誤會.
嘲笑.
被猶太的同胞.
來逼迫.
甚至乎.
來誣告他.
叫人不要遵守.
摩西的條規.
這是聖行傳的十五章.
這樣說的.
所以保羅.
是很為難的.
是很痛的.
被自己的同胞這樣.
推下去.
其實他是一個很虔誠.
守神誡命的人.

$^{561}$他在第三次在以弗所開船.
回到耶路撒冷.
上到耶路撒冷.
還沒到耶路撒冷之前.
在該撒尼亞.
他有怨在身.
就是拿西爾人的怨.
他自己就完了那個怨.
之後才剃頭.
你可以想像.
保羅其實他自己是從小.
是一個很循規蹈矩的人.
但他終於發現.
不是循規蹈矩才可以得救.
是因為耶穌的救恩.
你信才可以得救.
其實好像很簡單.
但猶太人就是不肯接受.
就是不能夠.
透過一個這麼簡單的.
死囚木匠的兒子.
信他.
我就可以完全.
和神立即恢復.
這樣的關係.
所以對外邦人是容易.
但對猶太人就太難了.
在屋牆.
這個.
西牆.
鎖羅門的位置.
那個聖殿.
很巨大的石.
他們會將那些紙.
神的許願.
塞進那些縫隙裡面.
很多那些紙等等.
用來祭祀.
不斷的叩頭.
不斷的表達對神的敬虔.

$^{601}$一定要戴頂帽.
遮住頭.
不可以讓神看到他們的頭髮.
等等.
這個猶太人.
記住 只有男丁才可以去.
女士是不能去的.
你可以想像到.
這個就是猶太人.
他們的虔誠.
帶到他們去不了神那裡.
接著第六.
到十三節.
我們可以先讀一讀這段經文.
保羅就去解釋.
他們的原因.
也可以這樣說.
就是神揀選人得救的重要原則.
「這不是說.
神的華落了空.
因為從以色列生的.
不都是以色列人.
肉身來說一定是.
那個血源.
尤其是亞伯拉罕的後裔.
就都作神的.
他的兒女」.
然後他引聖經.
這個錯世紀21章.
唯獨從以色列生的.
才要稱為你的後裔.
為甚麼從以色列生的呢?.
因為亞伯拉罕有另外一個兒子.
是由下格所生的.
叫做.
以.
以實瑪利.
這樣來到去.
以實.
年長13歲.

$^{641}$這樣呢?.
但是神就跟亞伯拉罕說.
這個大兒子不算數.
一定是要由你.
跟你的太太.
撒拉.
所生的.
那個應許的兒子.
才是神的兒子.
得到神的祝福和那個應許.
所以這裡告訴我們.
八節.
「欲身所生的兒女.
不是神的兒女」.
即是說.
就算你是有猶太的護照.
有猶太的血統.
甚至你是亞伯拉罕.
欲身的兒女.
都不是每個人都得到這個應許.
所以這裡告訴我們.
這裡也開始.
要令到猶太人非常反感.
就是說原來血緣.
不可以使你上天堂.
原來你是甚麼身份.
都不可以使到你.
能夠去到神那裡.
必須唯一的方法.
我們外邦人就知道了.
就是唯有我們在耶穌的.
愛和救恩裡面.
要告訴我們.
「唯獨那應許的兒女.
才算是後裔」.
第九節.
「因為所應許的話是這樣說.
到明年」.
這是17章的創世記.
「到明年這時候我要來.

$^{681}$撒拉必生一個兒子」.
這是創世記.
第十八章第十節.
「不但如此」.
就是說保羅說這個應許.
說這個真正的永恆.
屬天的身份.
原來這個救恩.
是一直一直發展下去的.
不只是亞伯拉罕的兒子爾薩.
現在連到利伯加.
就是爾薩的太太.
還有利伯加.
「既從一個人.
就是從我們的祖宗爾薩懷了人」.
雙子 因為他是雙子.
生的是爾素和雅各.
雙子 即是雙子.
還沒生下來的時候.
善惡都還沒作出來.
雙子出來.
誰好誰不好.
是沒有人會知道的.
善惡都還沒作出來.
就是說他們一直長大.
他們是龍還是鳳.
他們是甚麼修為.
我們說三歲還是八十.
但是保羅的意思.
他不是否定這些東西.
而是當時到底是爾素還是雅各呢.
沒有人知道誰是神揀選的.
是應許的兒子.
這裡告訴我們.
善惡都還沒作出來.
只因要顯明神揀選人的旨意.
這裡很重要.
揀選的旨意.
不在乎人的行為.
乃在乎召人的主.

$^{721}$表面上這句話.
引述的就是爾素的人生.
和雅各的人生.
雅各其實是奸狡一點的.
陰濕一點的.
深沉一點的.
很公於心計的.
爾素的人就好像很隨和.
很無所謂的.
兩個都是罪人.
但是在爸爸的眼中.
其實喜歡爾素多於喜歡雅各.
當然媽媽就喜歡雅各.
因為她時時學煮東西吃.
在家裡不出去打獵等等.
大家的性格都不同.
神不是看兩個人的行為.
是看神對誰揀選的旨意.
很奇妙地.
神揀選了第十二節.
神就對尼伯加說.
已經是第三次引聖經了.
在第九章.
所以九到十一章是大量引咎約的.
因為保羅現在在這三章裡面.
主要針對的就是.
很想和猶太人探討一個問題.
為什麼我們猶太人.
不可以順利地坐直通車.
直達天堂.
必須要下車.
必須要信主耶穌.
你一定要謙卑地承認.
這個木匠的兒子.
他就是永遠的神.
你要對他的真正身份.
謙卑,俯伏.
猶太人就做不到這件事了.
我不可以拜另一個神.
再沒辦法想像.

$^{761}$耶穌是我們的神.
是我們的彌賽亞.
所以這個很大的重擔.
使猶太人都好像呼吸不了.
但這個不是神.
要使他們這樣沉淪.
而是很想告訴他們.
保羅在這裡你可以想像得到.
他現在在新約成為了使徒之後.
他現在再去消化整個救恩.
他現在從一個猶太人的拉比.
的一個博士的身份.
和有救恩的使命的神的使徒.
來到綜合了之後.
再去看回舊約.
重新解釋舊約.
保羅可以得到一個這樣的結論.
就是告訴我們原來這個揀選.
早就已經有了.
在亞伯拉罕.
他的兒子以撒和以實瑪利.
神揀了以撒.
然後在以撒的兒子.
還要是孖仔.
孖仔揀誰呢?.
揀了的就是第十二節.
真是可圈可點.
保羅人的經文.
神就對尼伯加說.
尼伯加說就是還沒出生.
我們回到創世記才25章.
創世記25章.
雙子都還沒生下來.
25章.
因為這本聖經比較大.
所以很重.
所以他的福斯章有時候會跌.
25章說到以撒.
以素和雅各出世.
第十九節.

$^{801}$亞伯拉罕的兒子以撒的後代記在下面.
亞伯拉罕生以撒.
以撒娶尼伯加為妻.
那時候是40歲.
尼伯加就是巴旦亞蘭.
亞蘭人彼土里的女兒.
是亞蘭人拉班的妹妹.
以撒因他妻子不生育.
就為他祈求耶和華.
耶和華英文他的祈求.
就是以撒祈禱.
他的妻子尼伯加懷了孕.
孩子們在他福中.
孩子們當時沒有超聲波.
但是有雙胞胎.
有一對不知道是男還是女.
孩子們在他福中彼此相憎.
他就說若是這樣.
我為什麼活著呢.
可能不斷胎動.
又踢媽媽.
他就去求問耶和華.
尼伯加求問耶和華.
第二條23節.
保羅引這節聖經.
耶和華對他說.
兩國在你福來.
兩族要從你身上出來.
即將來他們的成就.
這族必強於那族.
留意這一節.
將來大的要服侍小的.
到頭來就是.
他們兩個都不知道.
可能不知道有沒有超聲波.
大的要服侍小的.
誰先不要出.
尼伯加就沒得選.
也不到他選.
他第一個出來就推他進去.

$^{841}$他也不知道誰好誰乖.
所以這樣說.
對一個大老婆都很困擾.
對一個快要生的大老婆.
將來大的.
你會生兩個孩子.
大的要服侍小的.
意味著小的就是應許的兒子.
這是第12節羅馬書.
保羅所引述的.
《創世記》25章.
第二代的揀選.
首先是亞伯拉罕的兒子.
不選擇13歲的以實瑪利.
選擇和晚年.
和撒拉所生的兒子.
就是這個以撒.
晚年生兒子.
其實都高危.
又沒得產前檢查.
又沒得教你怎樣呼氣.
但古代的人就好像很容易生.
當然一定有難產.
例如拉傑也是生在弱勢的時候難產死的.
所以你可以想像到.
到第二代了.
以撒自己又是有這個困擾.
原來神說兩個兒子.
我只是選擇其中一個.
就是小的.
在這個媽胎裡面.
選擇其中一個成為接續神的應許的子孫.
即是說應許的兒子.
就會有得救的信心.
就會有對神的信心.
是和一般的猶太人.
自己的同胞是不相同的.
這就告訴我們.
這就是我們需要學習的.
所以回到.

$^{881}$我稍為收窄範圍的一個例子.
一個的勸勉.
就是在眾多的基督徒當中.
如果我們只是將我們的信心.
放在一個教會.
放在一些人.
或者放在一些我們自己喜歡的團體.
要留心.
因為猶太人都是這樣.
他們大部份人都是信神.
每個星期都是去聖殿的.
但神告訴我們.
蒙應許的兒女.
就有這個和神的關係的信心.
才能夠接續神的祝福.
所以是以撒的後代.
雅各的後代.
約瑟猶大的後代.
才可以一路承受那個祝福.
大家都是同一個時代出生成長.
就沒有了那個信心.
你可以想像得到.
這是我們今天都要留意.
如果我們地球上.
有沒有十億人的基督徒呢?.
如果有的我們都要小心.
在十億個裡面.
我們一定要自己有和神.
直接的信心的關係.
否則就會好像變了.
就是第十三節所說的.
正如經上所記.
保羅引述的不是《創世記》.
是引《馬拉基書》.
雅各是我所愛.
即是說他們已經過了他們的人生.
已經是蓋棺定論.
小兒子雅各是我所愛.
是我所揀選.
是我的旨意.

$^{921}$他始終能夠接續到那個信心的行列.
成為信心的偉人.
但是以素蓋棺定論.
是我所恨惡的一個人.
嘩!你這樣說真的很坎坷.
又讓我們看到眼白白.
你大家聽到神這樣的應許.
近乎有些咒詛的成分.
因為大的將來一定要服侍小的.
大的那個命運.
原來隨著他的成長.
原來和神是有一種距離的.
雅各表面上是一個很壞的人.
但是原來最終他和神是有很密切的關係.
雅各是我所愛.
保羅還要是引《馬拉基書》的經文.
不是引《創世記》的經文.
就更加加重那個語氣.
和他對這個揀選的堅持.
保羅你就想像得到.
我們認為不公平.
不可能這樣的.
大家都是這樣在媽媽的肚子裡出來.
在你父親的肚子裡.
都是爸爸以撒的兒子.
以撒是蒙應許的.
以撒的爸爸就是亞伯拉罕.
那以撒的子女是不是都應該蒙應許呢?.
應該似乎是.
但是神在兩個裡面再有篩選.
再有那個揀選.
再有那個獨一無二的旨意.
就是要選擇雅各.
就是要愛雅各.
是我定義要這樣做的.
誰可以挑戰我呢?.
這個就是我們在這樣的經文裡面.
看到神這樣的計劃和心意.
所以其實我們真的要留意.
我們不要看見有些基督徒很順利.

$^{961}$很平安.
我們就羨慕或者嫉妒不值得.
我們有我們自己人生的路.
最重要我們要保持自己在神的揀選裡面.
我們必須要讓自己常常在神的愛裡面.
所以是一次的得救.
一直的相信.
才是必然的得救.
來到第十四節.
他說這樣我們可說什麼呢?.
神有沒有什麼不公平呢?.
果然真的提到公平的問題了.
因為好像在這樣的教導下.
是非常的偏心.
是非常的針對耳塞瑪利.
為什麼大家都是爸爸的兒子.
為什麼他們得到這個屬天的祝福呢?.
大家記住耳塞瑪利不是沒有地上的恩典.
耳塞也有耳塞分給他的一些地上的財物.
有句很經典的.
就是耳塞打獵.
因為爸爸已經老眼昏花.
我吃了你的野味.
我就為你祝福了.
你就可以得到神的祝福了.
或者神在我生命裡面完結.
將這些福轉嫁到你的戶口.
耳塞就興高采烈出去打獵.
那一刻媽媽李百嘉就聽到.
就馬上自己去弄一些野味.
很快就預備好.
就拿進去給爸爸.
爸爸聞到那些香氣.
又摸摸它.
真的全身都是毛.
其實是用山羊皮包著.
讓爸爸不要有錯覺.
或者有一個發現.
就盡量幫牙角.
吃完那些野味之後.

$^{1001}$盡量幫牙角祝福.
又揭一揭.
《痛世紀》的27章.
還是29章.
我要揭一揭.
大家稍微看看.
可能27.
我先把這本聖經翻下來.
因為這個站常常動.
耳塞祝福牙角.
沒錯27章.
第18節.
27章第18節.
「牙角對他父親那裡說.
我父親他說我在這裡.
我兒你是誰.
當然是阿爸爸.
牙角對他父親說.
我是你的長子耳掃.
我已照你的吩咐做了」.
即是做了野味.
「你起來坐著.
吃我的野味.
普吸我的祝福.
耳塞對他兒子說.
我兒你如何走得這麼快呢?」.
他也有些懷疑.
「他說因為耶和華你的神.
使我遇見好機會得著的」.
連說謊都要拿神來做甚麼.
第二節.
「耳塞對牙角說.
我兒你近前來.
你真是我的兒子耳掃不是.
牙角就靠近他父親耳塞.
耳塞摸著他說.
聲音是牙角的聲音.
手卻是耳掃的手」.
即是半信半疑.
「耳塞便不出他來.

$^{1041}$因為他身上有毛.
像他哥哥耳掃的手一樣.
就給他祝福」.
接著就說那個福氣.
「又說你真是我兒子耳掃嗎?.
他說我是.
耳塞說你弟給我.
我很吃.
我兒子的野味.
給你祝福.
牙角就給他.
他便吃了又拿走給他.
他也學了他父親耳塞對他說.
我兒.
你上前來與我親嘴.
他就上前」.
即是這些是當時他們的文化.
與父親親嘴.
即父親好像都大能.
即有他的大限.
「他父親一聞.
他衣服上的香氣.
就給他祝福說」.
留意一下這個祝福「我兒的香氣.
如同耶和華似福之田地.
的香氣一樣.
願神賜你天上的甘露.
地上的肥土.
並許多五穀生酒.
願多民事奉你.
多國跪拜你.
願你作你弟兄的主.
你母親的兒子向你跪拜.
凡咒助你的願他受咒助.
為你祝福的願他蒙福」.
盡他的.
即是.
一個老人家.
一個家長一個族長.
一個支派的領袖.

$^{1081}$盡他的.
Blessing.
來為這個兒子.
祝福是屬天的Blessing.
留意下面才是重要.
「以實.
為雅各祝福元北.
雅各從他父親那裡才出來.
他哥哥爾素正.
打獵回來也作了.
美味」.
即是「野味」「拿來給他父親.
說請父親起來.
吃你兒子的野味.
好給我祝福」.
留意這段經文32字.
「他父親以實對他說.
你是誰?他說.
我是你的長子爾素.
以實就大大的震驚.
說你未來之先.
是誰得了野味.
拿來給我呢?.
爾素驚吃了.
為他祝福.
他將來也必蒙福.
爾素聽了他父親的話.
就放聲痛哭.
說我父啊.
求你也為我祝福」.
這些是次等的祝福.
「以實」.
很清楚.
他做了這個.
錯誤的決定之後.
原來那些福氣.
是直接流到雅各身上.
「以實」不可以說.
吐口水說過.
他知道神.

$^{1121}$我們用一個比較俗的.
但很傳神的.
就是神推鬼推.
使得「以實」祝福了雅各.
其實他不想的.
但不想神也要祝福他.
不過留意這裡.
他在這裡終於明白.
神的揀選.
他現在不敢.
違逆神的旨意.
他現在終於知道.
原來神起初.
和他的太太李百嘉所說的.
是真的神會成就.
所以他現在.
作為家長.
他才真正在這一刻悔改.
他現在不敢.
明明這個真是爾素.
明明這個真是.
他所愛的.
很重的.
那個兒子.
他現在有什麼不敢做.
他不敢幫他去祝福.
屬天的那個身份.
所以這裡.
是看到「以實」.
終於明白.
和肯接受神的旨意.
所以他祝福的.
只是田地上的.
那些成長.
你可以想像得到.
他是不敢.
去越過了神的旨意.
他終於在這裡.
肯低頭.
一直都是違反.

$^{1161}$違抗神的命令.
一直都要和神鬥氣.
我都喜歡大兒子.
我希望我的大兒子蒙福.
但是神說不是.
李百嘉所生的第二個兒子.
其實出池可能幾分鐘.
就是雅各.
才是神揀選的那個兒子.
你可以想像得到.
經文的精彩之處.
和神的僕人.
「以實」.
他去到這個階段.
去到自己晚年.
才真是大徹大悟.
不敢再得罪神.
希伯來書有說這件事.
希伯來書有說.
「以實」.
祝福.
我們看看第十一章.
有說到「以實」.
他這一刻.
你可以想像得到.
你是最愛的那個大仗.
他就在那裡.
你是最愛的那個大仗.
現在陰差陽錯.
被小兒子騙了.
如果我們現代人.
誰理他.
我現在照樣祝福我的大仗.
再祝福多一些都可以的.
都是口說出來的.
還要誠心所願地說都可以.
雅各那些不算數.
廢棄他再周初他.
正常人就應該這樣.
但是「以實」就是不正常.

$^{1201}$他這個不正常.
就是他終於成為一個屬靈人.
終於知道神的旨意.
他不敢再違抗.
抗拒神的旨意.
很多年在那裡兩個角力.
其實是同神角力.
留意「以實」這裡.
第八節開始.
不過很長的第八節.
我們就看.
再後一點.
第十七節.
是不是第十七節呢.
應該是的.
這個是希伯來書十一章.
信心的榜樣.
應該說不同的時代.
不同的人.
怎樣在他自己的生活.
運用那個信心去信靠神.
所以不同的人.
所以沒有所謂毀人.
不同的時代的人.
怎樣去信靠神.
有些在獅子坑裡面信靠神.
有些就在生仔裡面信靠神.
第十七節.
亞伯拉罕因著信被試驗的時候.
就把以實獻上.
這便是那歡喜領受應許的.
將自己的獨生兒子獻上.
輪到這個兒子.
就是以實.
曾有話說.
從以實生的才要稱為你的後裔.
他以為神還能叫人從死裡復活.
他也彷彿從死中得回他的兒子來.
留意第二十節.
作者怎樣描述以實的信心呢?.

$^{1241}$以實因著信.
就指著將來的事.
給雅各以素祝福.
一句說話就說出.
以實他有對神的敬畏.
信心.
就是將次序調轉.
作者也告訴我們.
真正蒙福的是雅各.
以素只是得到一般地上的福氣.
父親對他的愛心的回饋.
不是神的祝福.
所以原來真正的信心就是.
不敢在自己喜歡的事情上.
去做超過了神的旨意.
這就是很高的屬靈境界.
以實原來一生都愛這個大兒子.
去扶植這個大兒子.
做那個蒙應許的兒子.
可惜怎樣扶呢?.
最終原來神的旨意.
沒有人可以來抵抗.
真是陰差陽錯.
雅各又先走了.
出世又先走了.
拿過長子的祝福就先走了.
就騙了父親.
肯定是欺騙的.
是謊言.
但是神的工作.
就不會停留在那裡.
神是會有辦法去訓練這個雅各的.
所以在這個時候.
當我們這樣總結這些經文.
很自然就會和保羅一樣.
保羅都很想引導我們.
或者猶太人去想.
一定去到這裡的時候.
就說到底.
這樣就是神很不公平.

$^{1281}$十四節果然就是問了.
這樣我們可以怎樣說呢?.
難道神有不公平?.
因為上行已經說了.
對以實,瑪利不公平.
對以素不公平.
為什麼給他不給我呢?.
我也是父親的兒子.
這樣你就可以想像到.
還有雅各不應該得到這麼好的祝福.
再加上就是.
以實連最後他肯為兒子祝福.
祝福完之後不敢再祝福給以素.
是不敢的.
因為他已經很清楚.
神攔阻他.
神的旨意攔阻他再這樣做.
如果這樣做.
他整個的靈命就跌入谷底了.
就是說有雙重的祝福.
那就是等於你是神.
神說你要祝福給小孩.
就是小孩得到那個福氣.
真是在這樣的擺弄下.
他騙了祝福.
但是最終是走在神的計劃之上.
所以以實他明白之後.
真是不敢再做一些他以為可以做的事情.
他過去一直想扶植以素的.
他現在不敢了.
所以使得以素很害怕.
大哭.
因為爸爸為什麼你沒有福給我呢?.
為什麼我不可以成為那個應許的兒子呢?.
他終於知道原來應許的兒子這個身份.
是天大的一個福氣.
他以往都是很輕視的.
所以這裡就告訴我們.
這裡就立刻帶出一個很嚴肅的問題.
神很不公平.

$^{1321}$那保羅怎麼去拆解呢?.
原來他講到這裡的時候.
他真的會想到神這樣做.
是會讓人有一個救病.
你真的會很不公平.
保羅就立刻回答.
斷乎沒有.
就是說是很簡單節竭.
神是不需要有這個.
被人放在天平上來審判.
來去論斷他有不公道.
為什麼呢?.
第十五節.
《聖經》出埃及記.
因他對摩西說.
我要憐憫誰就憐憫誰.
要恩待誰就恩待誰.
《出埃及記》33章第19節.
當時摩西求民.
不如叫我回埃及吧.
那幫百姓都是這樣來得罪你.
我真是有一個磨心.
氣力不討好.
不如你不要找我吧.
神就憐憫他.
就講了這一番話.
如果大家不介意就揭一揭吧.
怕大家都不熟悉了.
《出埃及記》33章.
都要放下《聖經》.
因為這個下次要拿另一個好一點的.
不停地墮.
跟Yoyo說一聲.
33章.
本事已經重了.
承受不了.
太薄了.
想用口水幫他.
最方便.
33章.

$^{1361}$第17節.
因為當時.
摩西已經很心力交瘁.
那幫人時時都敵黨他.
其實不是敵黨他.
是要敵黨神.
17節.
耶和華對摩西說.
你這所求的我也要行.
因為你在我眼前無了因.
並且我按你的名認識你.
摩西說.
求你顯出你的榮耀給我看.
當時摩西真是想.
所謂劈炮不撈.
太辛苦了.
面對這一幫百姓.
不斷地去催逼他.
攻擊他.
埋怨他.
他很難去應付200萬人.
60萬的藍丁.
當時還沒有咒詛.
還沒有進入迦南之前的刑罰.
所以現在是剛剛在西乃山三個月.
《出埃及記》這段經文.
離開了埃及三個月.
三個月.
那個記憶還是很留在目的.
已經忘記自己是奴隸了.
現在覺得自己很興奮.
這樣要行那樣要行.
所以就很多埋怨.
不過留意第19節.
耶和華說.
我要顯我一切的恩賜.
在你面前經過.
宣告我的名.
我要恩代誰就恩代誰.
要憐憫誰就憐憫誰.

$^{1401}$又說.
你不能看見我的面.
因為人見我的面.
不能存活.
所以神就將他擺在盆石上.
遮住他.
神的榮耀.
向那個盆石經過.
摩西見到神的榮耀.
這個真的很神跡的奇妙.
一個人能夠見到神的榮耀.
見到神的悔之門.
這是第23節.
所以我們回來這裡的時候.
保羅就引這節聖經.
來告訴我們.
神是絕對不需要被人去論斷.
是否在救恩裡面.
公平與不公平.
因為這就關乎到.
簡選和公平.
這個邏輯的相對.
神的簡選.
神的旨意.
是不需要因為那個人本身.
這就和公平沒有關係.
因為神是造物主.
祂創造人類.
人類本應每一個都不應該得到神的拯救.
因為每個人都是罪人.
但是神現在就將他的愛和憐憫.
淋到一部分的人.
而這一部分的人.
始終是有信心.
與神的救恩去吻合.
銜接.
而得到生命改變.
但是未改變之前.
是出於神的簡選和憐憫.
所以與公平沒有關係.

$^{1441}$現在沒有這個邏輯.
在神的角度裡面.
這樣去看.
好像如果我有一些物件.
那個物件如果能夠說話.
我喜歡這個電話.
與另一個電話.
如果兩個電話都能說話.
為何你特別喜歡我這個電話.
而不喜歡我.
甚至將我這個電話賣了.
這都是我的主權.
是我的簡選.
是沒有公平可言的.
與公平無關.
因為兩個電話都是我買回來的.
都是屬於我的.
所以你可以想像到.
人是爸爸生的.
親生的媽媽.
媽媽到幾分鐘.
就不能簡選.
這就真的很不公平.
但是連這個部份.
神都說他不需要解釋.
不需要辯護.
我是不公平.
因為這就是簡選的問題.
神是要定義愛,雅,郭.
恨,惡,義,素.
完全令到我們今天的邏輯.
是過不了的.
所以猶太人再擴大民族的範圍.
神定義愛這些拜偶像的外邦人.
當時羅馬人,希臘人.
是拜那些假神的.
那個二不所的廟的那個柱.
那個大安娜.
那個叫做.
忘記了名字.

$^{1481}$那個守護神.
整個二不所都要拜的.
你想想這麼多的人.
以前是拜偶像的.
是交軌的.
現在大量的信耶穌.
恩信就清二.
對猶太人來說.
他們從來都沒有一個對神的敬畏.
但是現在就被簡選.
是很難接受的.
對猶太人來說.
或者對耳塞瑪利.
是很難接受這個所謂簡選.
對今天的外邦人.
對今天我們未信的一些地上的人.
你們基督徒.
你說信耶穌.
神就簡選你上天堂.
我們現在不信.
未明白.
只不過是有一種理性的分析.
不肯信.
你就說我們要下地獄.
神沒有簡選我們.
就令到這些未信的人.
就有很大的反感.
對基督教這個簡選論.
是很大的拒絕.
你越講.
他就越反感.
我平生做人.
對得起天地良心.
我是沒有被簡選.
我就要下火壺地獄.
你們這班人基督徒.
就說被簡選.
如果就這樣去看.
就這樣去分析.
就這樣去辯論.

$^{1521}$我們基督徒都好像站不住了.
但是保羅就很理直氣壯地.
他以一個拉比的身份過來人.
再引述舊約的聖經.
來拆解這個所謂公平不公平的問題.
神是沒有不公平.
因為神要憐憫哪一個.
就要憐憫哪一個.
這個就真的回到神絕對的簡選裡面.
就涉及到這個簡選的背後.
原來神是對某一些特別的愛.
以致簡選了他.
而這些簡選的這樣的恩典.
是會塑造那個人.
是會慢慢越來越因為信而信服.
而有一個愛神的心.
來願意回饋救恩的.
所以我們記得.
如果我們真的信主.
其實在萬人當中.
神定義憐憫簡選我們.
恩待我們.
因為十一章保羅就警告外邦人.
其實我們是天生的野橄欖.
現在蒙了一個這麼大的恩典.
我們就應該感恩.
所以第十六節.
據此看來.
不在乎那個定義.
也不在乎奔跑的.
即是在說耳塑.
在說耳實馬利.
或者照著這樣的邏輯的人.
我自己對得起自己.
我為甚麼要稀罕你的簡選.
你不選就不要選.
我有一個人.
我有我自己自主的靈魂.
我可以主宰我自己的人生.
但他不知道其實.

$^{1561}$靈魂不滅.
將來要去到永恆的主的面前審判.
所以始終是創造和受造.
兩個的天淵之別.
我們是受造和做物主.
那個的相距是很不同的.
所以我們今天帶著這一份感恩的心.
去看簡選就相當容易了.
因為聖靈已經改變了我們.
我們已經有了一個生命的扭轉.
我們已經由遠我們的心靈.
認識耶穌不是只是木匠的兒子.
而是祂是在永遠之上的神.
我們就已經有了這個降伏.
所以神的簡選只有一個條件.
就是被簡選的人.
去到耶穌的面前.
承認祂是你生命的主.
所以我們的敬虔是每天承認.
耶穌在我的生命裡面有祂的大吏.
很多人突然聽到自己的含情.
就晴天霹靂.
當然好好地基督徒.
慢慢地會沉澱下來.
安靜了 祈禱了.
然後他會承認.
耶穌基督你有你的主權.
你現在要在這個階段.
需要面對癌症.
我願意馴服.
敬虔的基督徒就要經過這個過程.
當然會有困難 會痛 會辛苦.
會嘔吐等等.
但是承認這個主權.
所以簡選和那個人慢慢降伏.
對神有一個愛和敬畏.
原來真是相輔相成.
所以被簡選的人.
就不要以為不用擔心了.
坐著花轎上天堂.

$^{1601}$其實原來是沒有關係的.
那個簡選和那個人越來越慢慢.
心靈的謙卑.
對主 對人生 對神的認識.
會越來越在他的內心有一種深化.
所以做基督徒應該越來越告訴別人.
是呀 神本來不應該簡選我的.
其實你這麼好 神應該簡選你的.
不過我的不好 神都來寬恕.
所以我得到這個生命之後.
你看看我的人生.
慢慢地去到在主的面前謙卑降伏.
帶著感恩的心.
你就不行了.
因為你根本沒有降伏你的心靈.
你以為你自己就是神.
你就沒有降伏你自己.
所以其實表面上就好像.
為什麼不公平 你選你不選我.
但其實不是的 兩個都不應該選的.
但當選了之後.
你看看我的人生就是這樣.
慢慢地改變 是我真心樂意的改變.
這個就是簡選帶來的那個人的信而信服.
所以原來神的憐憫.
是會令到那些在還的人.
心靈會柔軟的.
所以希望救恩令到我們越來越柔軟.
我們應該越來越喜歡報道會.
不是去看誰舉手誰什麼.
而是在一個聚會裡面參與.
你是可以在祈禱 在爭戰.
為一些的牧者去守望禱告.
一場的肅靈爭戰.
還有你看到有人信主.
你是真心的替他高興.
他的生命改變 他舉手決戰.
雖然那條路還很漫長.
但是你真心的為整個聚會.
神得到榮耀 你去感恩.

$^{1641}$應該帶著這樣的眼光.
來到聚會參與.
不要說關我事 今天沒有朋友.
我不回這個教會 不回這一課了.
就不要這樣.
應該真是一起參與一場爭戰.
所以神的憐憫.
原來帶給那個受憐憫的人.
會有生命改變.
所以你會看到熱誠.
所以你會看到這個雅各的生命改變.
雅各的改變就是32章.
那個彼路伊拉的事件.
接著就是第16節.
他此一看來不在乎定義.
不在乎奔跑 只在乎發憐憫的神.
所以原來憐憫是可以令一個人的生命改變.
神的慈悲憐憫.
因為經常有話向法老說.
不單止是這樣.
在出埃及的時候.
法老是定義的剛硬.
一而再神憐憫他.
他就一而再的十災.
即是十次的剛硬.
災禍一止息了.
他就立即剛硬.
法老說好了.
要讓以色列人走.
求你叫那些災離開.
十災.
首先就是最後的滅長子.
首先就是蛙塞形.
就是蒼蠅.
蓄就是那些生蓄的瘟疫.
接著就是獨蒼.
薄就是下了冰雹.
黃就是蝗蟲.
黑就是全地都黑了.
最後就是滅長子.

$^{1681}$第一個就是水變血.
就是十災.
每一次災禍一減去.
法老就立即變剛硬.
所以他的心靈越來越剛硬.
所以下面就這樣說.
第十七節.
因為經常有話向法老說.
我將你興起來.
特要在你身上彰顯我的權能.
並要使我的名全遍天下.
如此看來.
臣要憐憫誰就憐憫誰.
要叫誰剛硬就叫誰剛硬.
保羅這樣一個重要的總結句.
就好像將雅各和以素比較.
法老是經過十次神的憐憫.
透過摩西給他的教訓.
災禍一個大於一個.
而且有人說.
十災是當時埃及拜的十個神.
拜尼羅河.
拜其他自然界的東西.
所以神就透過這樣的災難.
這些神蹟.
教訓法老要敬拜真神.
不要再拜這些偶像.
但法老真的將他的心把心一橫.
這樣來到不斷.
沒事了.
就立即逼迫摩西.
再加重那些刑罰.
他們要做專家.
要加重那些刑罰.
不斷神就透過摩西.
要讓法老看到.
原來真正傳遞的造物主.
就是耶和華.
就是神.
所以他和神作對.

$^{1721}$最終也是死於紅海的淹沒.
所以神就令到法老心裡剛硬.
出現好幾次的.
在《出埃及記》第七到第十一章.
出現了好幾次.
神令到法老剛硬.
後來摩西的描述.
創世記的作者告訴我們.
法老自己將心剛硬.
慢慢慢慢.
神好像令到他的心硬.
但慢慢.
法老自己的心越來越硬.
所以.
神的教育.
沒有令到沒有被揀選的法老.
心靈柔軟.
心靈只會越來越剛硬.
但如果一個被揀選的人.
在救恩底下.
心靈是越來越柔軟.
對造物主.
對神.
對永恆的認識.
因為主要使自己的名全遍天下.
所以我們應該在生活的每一個層面.
不斷去發現上主的同在.
恩典和祝福.
所以不要擔心我們的人生.
有主的同在.
所以第十九節.
我忘記了.
所以這裡是很大的課題.
神的主權和人應該要負的責任.
所以這裡告訴我們.
繼續在這一章裡.
不斷去強化揀選.
不斷告訴我們.
外邦人其實一樣是不應該得救的.
不是那麼容易的.

$^{1761}$但是神有一個普遍的恩典.
就是要施行給每一個信主心靈柔軟的人.
第十九節.
這樣你必對我說.
他為什麼還指責人呢?.
有誰抗拒他的旨意呢?.
你這個人.
保羅再提出另外一個假設的問題.
就是說知道人心.
保羅你都說了.
神祂最偉大的.
祂要選擇誰就選擇誰.
開始軟皮蛇.
開始好像沒有東西可以反駁.
你想怎樣就怎樣.
你喜歡.
這些是非常負氣的相反的說話.
但實際上就是很不服神的這個揀選.
很不公平.
為什麼不選擇我們猶太人呢?.
所以保羅就告訴我們.
為什麼你還要指責人呢?.
誰抗拒他的旨意呢?.
其實這些是相反的說話.
就是一些相反語.
帶著有骨的相反語.
接著保羅就很重.
好像自己和自己對話一樣.
其實就是和一個民族.
對自己的身份有恃無恐.
這個都是一個信仰的大忌.
當我們有恃無恐.
我們就會慢慢對信心的信仰.
慢慢有時候我們就會.
沒有了初信的時候那份單純.
所以保羅的說話就很重.
你這個人你是誰.
竟敢向神來博取.
受造之物.
豈能對造他的說.

$^{1801}$你為什麼這樣做我呢?.
這個就是創造和受造的徹底的.
天淵之別的分別.
所以神造我們是有自由意志.
是有生命的存在的個體.
是活在這裡有知識.
我們有領悟力.
我們是懂得分析,懂得思考.
但唯一我們的思考就是.
不能夠去思考這位造物主的榮耀.
那份偉大,那份的全能.
所以我們用人間的邏輯和哲學.
來和神有這樣的針鋒相對.
甚至有一種輕佻.
他喜歡怎樣就怎樣.
他說這樣就這樣.
他創造他就選擇.
你不選擇我我就這樣.
這些就是很晦氣的說話.
忘記了受造.
就是說神永遠是做那個永遠的那個莊.
我們永遠都是在他創造的旨意底下而存在.
我們就很難接受.
哇 原來他對我有生殺之權.
我們很難去想這些東西的.
我們只會想可能特朗普就會有生殺之權.
可能那個和約簽不成.
我們烏克蘭人就慘了.
可能有很多他們的擔心.
但是很少人在地上會擔心神的時間到我們到了.
我們就要被收回.
很少這樣的情況.
所以第二節.
保羅再引一些當時的一些容易理解的例子.
堯將難道沒有權柄.
從一團泥裡拿一塊作成貴重的器皿.
又拿一塊作成卑賤的器皿.
所以一而再要強化這個揀選.
倘若神要顯明他的憤怒.
彰顯他的權能.

$^{1841}$就多多忍耐寬容.
那可怒預備遭毀滅的器皿.
又要將他豐盛的榮耀.
彰顯在那蒙憐憫.
早預備得榮耀的器皿.
就是我們信主的.
所以都是神豐盛的憐憫.
令到我們終於可以信主.
而去到那個信而信服.
對神有一個敬畏.
認識他是至聖的造物主.
但是那些被神擊打.
來毀滅的器皿.
神就仍然用多方的寬容忍耐.
來忍耐這些仍然在地上.
對神不斷藐視.
不斷挑戰.
不斷說不要神存在.
我決定我自己人生的主宰.
這就是神多多忍耐寬容.
原來這些人要就坐沉淪.
神都用多多的寬容忍耐.
但對於蒙神的揀選憐憫.
神是用豐盛的榮耀.
將他最豐盛的榮耀.
當然就是十字架救恩.
耶穌的大愛.
耶穌來這個世界的救贖.
就是神豐盛的榮耀.
彰顯在我們的生命裡.
所以我們就能夠成為了.
得榮耀的器皿.
所以你想想一個天淵之別.
真是分道揚鑣.
我們繼續蒙在神的恩典底下.
所以記得我們一生都要.
在神這份豐盛的榮耀底下生活.
我們就會成為榮耀的器皿.
第24節.
這個器皿就是我們被神所召.

$^{1881}$說得很清楚.
不單止在猶太人當中揀.
可以在外邦人當中揀.
有什麼不行?.
保羅是這個意思.
然後他又引《河西那書》.
可以想像到.
保羅一而再在第九章不斷講揀選.
不斷透過舊約的例子,個案.
告訴我們原來我們蒙恩.
真是揀選的.
不是我們自己去選擇的.
不過神的榮耀,恩典,揀選.
來到我們身上.
我們就變成一個.
很願意放開自己的心靈.
打開自己的心靈.
來接受舊恩.
因為豐盛的榮耀.
就是顯在耶穌的生命裡面.
舊恩裡面.
後面看看時間夠不夠講完.
應該不行了.
還有一大段.
我再看一些powerpoint.
所以30到33節就是最後的總結.
猶太人為什麼被排擠在舊恩的門外.
一直不肯運用信心.
一直不能對主耶穌有一個全新的看法.
一直憑外貌肉身去看耶穌.
不能憑心靈去接受耶穌.
後面我們繼續告訴我們.
大概這樣的情況.
神的選民和信靠彌賽亞是一個很大的抵觸.
就是猶太人不能夠承認.
耶穌就是他們的彌賽亞.
所以本來大家都是同樣蒙到恩典的.
外邦人和猶太人都可以得的.
但是外邦人就很容易地.
因為主耶穌的愛.

$^{1921}$也是初期教會的使徒們.
冒險在困難將舊恩去傳.
很多人相信.
當然也有大量的人不相信.
例如尼祿這些皇室的人員根本就不會相信.
還是逼迫基督徒.
但是當蒙恩的人.
可能是猶太人.
但是更多在第一世紀是外邦人.
因信而清二.
所以只有一條路就是.
蒙救恩的路就是因信而清二的路.
所以猶太人不是他們有什麼問題.
是他們不肯接受耶穌這份信心的得救.
所以保羅就強調神的絕對主權.
在第九章這裡.
所以第九章是一個很精彩的經文.
我們可以去了解一下.
神是不會不公平的.
神是有絕對的主權的.
在救恩裡面.
肉身的血統還是信心的信靠.
猶太人一直用信心的肉身的血統.
所以排擠了這個因信而清二.
但是神的揀選和憐憫.
就臨到那些可以相信而生命得到改變的.
當然當時就是外邦人.
其實初期教會有很多祭司都聽了這個道.
《士徒行傳》第六章第七章.
都信奉這個道.
可能怕就不敢立刻離開猶太教.
但是重點就是我們和猶太人不相同的地方.
猶太人以為自己是被神揀選.
但是我們從來都不知道.
我們原來在萬民當中神揀選我們.
是我們信了主.
慢慢來到去謙卑.
在神的救恩底下.
認識了耶穌是我們寶貴的救主之後.
我們就來到去不斷的認識.

$^{1961}$透過羅馬書.
保羅就用另一個角度告訴我們.
原來是出於神那個不改變的揀選.
所以我們是被揀選的人.
其實我們真的要帶著一份迷迭行的心.
是冥冥當中有那個主宰.
所以我們是非常之有福的.
原先本來沒有份的.
但是現在我們成為無福的人.
這個就是在第九章這一段.
我們需要來到去學習.
所以真正的神的指紋.
是說憑信心接受神的揀選.
和神的珠背的憐憫.
豐盛的榮耀在我們生命裡面.
所以猶太人其實就是被基督.
他們自己的這個盆石來叛變.
這個是33節所講的.
下幕我們透過這個段落.
繼續來到去學習.
今天可能就停在第24節這裡.
佛西亞書那裡就下一次.
我們才再回來看.
保羅怎樣再繼續強化他那個揀選的神學.
因為下個星期28號都想停一次.
因為剛剛那天是我生日.
就停一次.
但是去到9月.
我們就要去到9月5號.
因為9月4號開始.
我們就有12個弟兄姊妹去參加.
在播道神學院和陳淑娟牧師所推行的.
第二屆的行在社區的福音運動和神學課程.
到時有7個教會在元朗區.
其實應該是元朗屯門區.
7間教會一起去上課.
到時會有一些神學教導.
和會有彼此的分組交流等等.
所以要將我們這個茶經班改為下個星期.
9月5號.

$^{2001}$下一次.
因為下個星期28號本來都有課的.
但是我想停一次.
因為今天是正日生日.
大家抱歉.
我們今天就停在這裡.
天父求你.
祝福我們認識你豐盛的榮耀.
這一份的揀選是堅定不移.
在萬世以前.
好像保羅的邏輯.
神的啟示.
讓他知道神的揀選是不會有後悔.
而且是堅定不移.
所以我們很盼望.
我們的責任不是滿足於在神的揀選.
而是我們在神的愛底下.
來不斷的生命的更新.
我們在主的恩典底下.
不斷的去學習.
去服事.
去回饋.
還有神給我們的光偉.
恩賜.
這樣來盡心的去服事其他人.
帶其他人去認識神的救恩.
使他們可以得救.
求主幫助我們.
也給我們做一個忠心有見識的管家.
在地上來生活.
求主幫助.
我們為了將來的9月5日開始.
我們有新的時間上課.
盼望主保守.
讓更多弟兄姊妹都能夠抽到時間.
或者之後聽一些錄音.
一起去學習聖經的真理.
我們人說都有很多的限制.
盼望主的聖靈.
都在不同的弟兄姊妹的心靈裡面.

$^{2041}$讓他們能夠融會貫通明白.
也為了我們在下個4日的星期四.
我們正式同恩堂.
和部分的感恩的弟兄姊妹.
一起去參加這個課程.
盼望神去祝福我們.
使我們能夠更深的認識.
這個服事.
多謝你聽我們感恩禱告.
奉耶穌基督的名.
阿們.
\newpage



\section{}
\label{sec:MB59D1BA0Dc}
\textbf{宣道會錦繡堂2025年8月30日晚堂崇拜直播}
\newline
\newline
連結: \href{https://youtube.com/watch?v=MB59D1BA0Dc}{\texttt{https://youtube.com/watch?v=MB59D1BA0Dc}} ~~~~ 語音日期: 2025-08-30
\newline
\newline
\hyperref[sec:3_mKpJn_jtA]{\small{< < < PREV SERMON < < <}}
~
\hyperref[sec:index_chronic]{\small{[返順時目]}}
~
\hyperref[sec:index_scriptual]{\small{[返順卷目]}}
~
\hyperref[sec:3TbpwI2Yugs]{\small{> > > NEXT SERMON > > >}}
\newline
\newline
$^{1}$接下來是宣導的時間.
請聽我將經文讀出.
經文是記載在詩篇十三篇一至六節.
詩篇十三篇一至六節.
葉和華,你忘記我要到幾時呢?.
要到永遠嗎?.
你轉念不顧我要到幾時呢?.
我心裡愁善,終日受苦.
要到幾時呢?.
我的仇敵升高壓制我要到幾時呢?.
葉和華,我的神求你看顧我,應運我.
求你使我眼目明亮,免得我沉睡至死.
免得我的仇敵說我姓廖他.
免得我的敵人在我動搖的時候喜樂.
但我倚靠你的慈愛,我的心因你的救恩快樂.
我要向葉和華歌唱,因他厚厚地恩待我.
經文讀筆.
主,今天有李富城牧師在場當中靜悼.
李牧師是弘恩基督教學院.
基督教士工高級文憑課程總監.
今天的講題是「要到幾時呢?」.
我們恭請李牧師.
相信剛才主席弟兄讀這段經文的時候.
「要到幾時呢?」這個片語.
很明顯是整篇詩篇最搶我們眼球的一句說話.
我選了這幅相.
不知道大家在人生中有沒有經歷過這種狀態.
「要到幾時呢?」.
疫情的日子.
局長要出來經常說的就是.
多少日限聚令.
7日,14日,21日.
如果你在限聚令期間.
有一些人生重要的事情要辦.
例如你要結婚.
例如你有兒女要結婚.
訂了教堂,訂了酒席.
就看著疫情的限聚令.
你就會在想「要到幾時呢?」.
我有一位弟兄.

$^{41}$也是傳道人.
跟我非常相熟.
他的年紀跟我女兒是相同的.
所以我就當他是兒子.
在疫情期間他結婚.
他租了在長沙灣.
一間浸信會的禮拜堂.
那裡可以坐400人.
比貴堂還大一點.
當日的限聚令令到他.
這個婚禮可以繼續舉行.
不過當時的要求就是.
只可以坐20人.
男家,女家一人10個.
我因為負責領唱.
才可以分到位置進去.
當日有一個舅父不適性.
就自己走來,即是沒有預位.
被幹事趕你出去.
因為這樣是違規.
在人生的重要關頭.
你如果遇到一些不順意的時候.
你會問「要到幾時呢?」.
人生會面對病患.
面對失意.
會面對失業,失學,失戀.
喪妻,喪夫.
甚至有時候你會遇到一些.
你認為是無妄之災的時候.
你可能會對上帝發出這個問題.
為何你容許他出現呢?.
而且還不斷尾.
要到幾時呢?.
在士徒行傳有一個畫面.
大家看一看這幅圖.
你會想起甚麼呢?.
主角是史提反.
這幅圖表達.
史提反因為捍衛真道.
引來猶太的狂熱分子.

$^{81}$要即時對他進行一個屍刑.
用石頭將他扔死.
這就像是史提反記載的第一個殉道者.
你想像一下.
如果當時是史提反的同行者.
或者是他的肢體.
或者是他的門生.
你看到他被攻擊.
而且屍刑是用石頭.
大家想著不是一槍.
李富城信不信耶穌阿信.
一槍就回天堂.
這樣就容易解決.
用石頭將你慢慢折磨.
不斷身體受摧殘.
流血.
炎炎日色.
如果你望著這個殉道的場景.
你作為史提反的朋友.
弟兄姊妹.
你會不會心裡有一個反應?.
神呀.
這種的逼迫.
要到何時呢?.
而聖經很幽默記載.
當史提反殉道之後.
聖經跟著不速一轉.
就是教會大遭逼迫.
大家不知道有沒有唱過一首詩歌.
叫做《教會根基歌》.
裡面有兩句歌詞是這樣唱的.
它說信徒警醒.
爭相問.
黑夜到底有多長?.
有時無論是個別的信徒.
或者是教會的群體.
有時落在一個逆境.
一個黑洞的時候.
特別你看不到出路的日子.
那個黑夜的時間.

$^{121}$並不是用物理的時間.
逐秒逐秒計算的.
是你看不到一個.
沒有出路的隧道深谷.
那種的沮喪.
2016年12月14日.
是我太太.
離留在馬來醫院的日子.
那一晚主診醫生說.
家人自稱.
你們可以不用離開醫院.
我和我女兒就在醫院.
守候了一個通宵.
那一晚這個黑夜.
在我人生來說.
並不是在滴答滴答.
計算一個物理的時間.
是一個人生進入了一個黑洞.
到第二天中午.
上帝把我太太接走了.
其實我心裡面都有發出抱怨.
其實我和上帝的祈禱.
都不是很要求過分的.
我覺得.
當時我女兒已經談婚論嫁.
我和上帝求什麼呢.
就是讓我太太見到女兒.
成婚.
在我們中國人的傳統裡.
媽媽的都叫做安樂.
不過上帝都沒有答應我這個祈禱.
我太太在2016年底過身.
我女兒在2018年才結婚.
詩人大衛同樣在詩篇第十三篇.
發出了這個吶喊.
在這個詩篇裡面.
很工整可以分做三個要點.
第一.
在這個詩篇裡面.
詩人大衛是對神發出了一個.

$^{161}$真誠的抱怨.
大家留意這個抱怨.
我們用詩歌的題材.
我們讀得有味道一點.
我們讀出那個意境出來好不好.
你看到紅色的要到何時呢.
你想想要用甚麼語氣.
甚麼語態來讀呢.
我們一起讀一次這段經文.
你們讀黑色的字.
我讀紅色的字.
然後反過來.
你們開始.
要到何時呢.
反過來.
又說你忘記我.
要到永遠嗎.
你轉念不顧我.
要到何時呢.
我心裡愁善終日愁苦.
要到何時呢.
我的仇敵聲高壓制我.
要到何時呢.
那就是詩人的心情很生氣.
詩人的心情很想.
如果上帝在他面前以人法.
就去搖他.
到何時呢.
在這兩節經文裡面.
讓我們看到詩人抱怨的原因是甚麼.
那你說當然是第二節了.
因為他說我心裡愁善終日愁苦.
我的仇敵聲高壓制我.
我心裡不斷這樣想.
終日愁苦即是我愁腸百結.
我想多想都想不通.
為甚麼我的仇敵在上風.
我在下風.
他不斷壓制著我.
我是被壓者.

$^{201}$這個現象當然是個原因.
不過我看為這是一個表層的原因.
更加深層的原因在第一節.
他跟神說你忘記我.
要到永遠.
這個表達是一個非常尖銳的表達.
用廣東話.
即是一輩子.
然後他說你轉念不過.
對猶太人來說.
耶和華是向他們養念.
是代表上帝的賜福.
在民數記六章二十六節裡說.
願耶和華向你養念.
致你平安.
這個耶和華轉念不過.
即是廣東話.
嘔吐在臉上不看你.
施恩覺得這個上帝在哪裡呢.
他覺得上帝漠視他.
離棄他.
置他於不利.
所以他心裡有一個很濃烈的抱怨.
一輩子.
其實在舊的聖經裡面.
因為個人的苦難對神發出真誠的抱怨.
在智慧文學在約伯記裡面.
都是一段很清楚的經文.
其中有幾節是很深刻的.
約伯說我厭惡自己的性命.
任由我講我的苦情.
因為我心裡面苦惱.
我要講.
第十八節他說.
你為何使我出母胎呢.
甚願我當時氣絕.
我希望沒有眼睛看見我.
然後第十九節.
第十九節那裡說.
這樣就如從未有過我.

$^{241}$我一出母胎就被送入墳墓.
我的日子不是短少嗎.
求你停止.
求你放過我.
使我稍得其樂.
你見到約伯和上帝的抗辯.
或者不明白.
去到一個就坐自己出生.
我不出世好過.
第二十節那裡說.
求你停止.
我們現在很流行說.
Stop.
不玩了.
不和你玩.
上帝不要再玩我了.
我寧願沒有出過世.
不要搞我.
有些弟兄姊妹有時在信仰的地獄都會說.
我不信耶穌好過信耶穌.
我不跟從你好過跟從你.
為何搞到我現在這樣的田地呢.
另外.
先知亦都因為社群的苦難.
對神亦都發出真誠的抱怨.
譬如哈巴谷書.
我們試試讀這三節經文開始.
耶和華.
我呼求你不應還.
要到幾時呢.
我向你呼喊暴力.
你還不拯救.
你為何使我看見罪孽.
你為何坐視姦惡呢.
毀滅和兇暴在我面前.
爭執和紛爭不斷發生.
因此律法無效.
公理從未彰顯.
惡人為困二人.
所以公理遭受扭曲.

$^{281}$惡勢力猖獗.
無法無天.
惡人當道.
無枉管.
無神管.
先知和上帝好像對部宮廷一樣.
在法庭上.
他作為控方.
他要控告上帝.
不管.
你的公義已經蕩然無存.
無論是詩歌.
無論是智慧書.
無論是先知書.
我們都看見很清晰的.
一種人對上帝的抗衡.
抱怨.
不過請留意.
上帝是接納我們這種真誠的抱怨.
因為天父對我們的關愛.
我們可以傾心吐意的跟祂說.
就是祈禱.
甚至是對神發出這種真誠的抱怨.
有時候我們在教會上.
一段時間我們的祈禱文.
都頗八股的.
你覺不覺得.
我們祈禱有套標準.
起承轉合.
作為公眾聚會.
即是公共崇拜.
領土員有準備.
我是非常欣賞.
不過如果你個人和上帝的心交.
你和祂在密室裡的聊天.
你的靈修.
你的不明白.
我相信我們不需要太工整.
有套祈禱的方程式.
你有碗就說碗.

$^{321}$你有碟就說碟.
你不需要再在上帝面前裝模作樣.
我想我們不需要寫篇亮麗的土文.
才可以和神.
將你心裡的抱怨來表達.
舉一個例子.
假如有一個小朋友.
他今天上學.
快要開學了.
上課的時候.
有些同學頑皮.
很吵.
老師罵了幾次都不聽.
於是老師就殺一儆百.
李富城不要再頑皮.
李富城就撓著頭.
我都沒有頑皮.
老師就睜大眼睛看龍.
其實是曾睿貴頑皮.
他就說李富城頑皮.
就抓了我.
我想抗辯的時候.
老師還橫蠻無禮.
你還說.
我覺得投降.
於是我很委屈.
我很無理地被老師責備.
我很不開心.
我覺得很冤屈.
如果我今年還要跟這個老師.
我今年不知道有沒有運氣走了.
他帶著一個很委屈的心情.
我回到家.
爸爸下班回來.
我想問這個小朋友.
即是李富城.
要不要寫一篇文章跟爸爸說.
我想爸爸一走進來.
我就抱著他的腳.
然後就哭冤.

$^{361}$然後就說.
黃老師冤枉我.
不是我.
是另外一個同學.
又是看錯我.
我想申辯.
他還說你再說就罰多一點.
於是我不敢說.
不過我真的冤屈.
爸爸.
而女對父親.
發出一個真誠的反應.
盡訴心中情.
大衛,約伯,哈巴菊.
就是這樣.
將她跟上帝的關係.
這種抱怨.
其實是將她跟上帝拉近.
由抱怨會進入另外一個信仰的境界.
叫做對神懇切的呼求.
我們一起讀第三,第四節.
開始.
耀我的上帝.
求你看顧我.
應允我.
你使我眼目明亮.
免得我沉睡至死.
免得我的仇敵說.
我勝了他.
免得我的敵人.
在我動搖的時候起落.
由抱怨會進入另外一個層次.
就是求神你出手.
去呼求那個小朋友被老師冤枉.
他一定會嘮叨他爸爸.
你明天打電話去學校.
跟老師解釋一下.
不是我頑皮.
大家留意這個呼求的對象.
呼求的對象是耶和華.

$^{401}$我的上帝.
有什麼特別呢.
耶和華是猶太人對上帝一個客體的稱呼.
我對神的一個稱呼.
陳一華牧師就是陳一華牧師.
但是後面那個是猶太人信仰者.
跟上帝的親暱的關係.
我的上帝.
這是一個立約的關係.
剛才我說小女兒2018年結婚.
那晚擺酒的時候.
新郎新娘都要上台說一下致謝的話.
到我女兒說的時候.
這個是她小時候對我的暱稱.
爸爸是她小時候對我的暱稱.
我很久沒聽過了.
她長大之後都不叫我.
不用怎麼稱呼了.
最近一次叫我爸爸.
是因為她生日.
我發訊息祝她生日快樂.
她有兩個小孩子.
我希望兒女聽教聽話.
她就多謝爸爸.
再聽一次.
在婚宴上她說.
我的爸爸.
這個稱呼馬上令我現場失控.
就在我坐的主家席上拿了紙巾.
因為將我跟她的關係升級到一個.
不能甩掉父女的關係.
她不是當我一個客體的李富城爸爸.
是一個很親暱的稱呼.
仍然當神是她的神.
這裡有三個祈求.
一個是看顧我.
看顧我是回應前面祂抱怨的.
你轉念不看我.
應允我.
是前面抱怨的.

$^{441}$好像上帝沉默.
上帝不出聲.
上帝不知道去了哪裡.
求神來回應她.
應允她.
使我眼目明亮.
一個解釋就是她哭得太厲害了.
所以哭得眼睛不舒服.
不過我更加接受一個比較屬靈的解釋.
是她的屬靈的眼睛掉進一個黑洞裡面.
她跟上帝仍然有一個不能夠理解的距離.
如果我們看這個屬靈的眼目.
很多時候就指向我們的信心.
我們也有一句經文說.
我們是憑信不憑眼見.
在逆境裡面仍然能夠跟上帝接連的唯一途徑.
是你的信心.
而不是你今天祈願土.
明天就來一個戲劇性的轉變.
我那個年輕的童工.
他也為他的婚禮祈禱.
限聚令不要只有20個人.
因為他有很多朋友.
預計要預約600人的堂.
也預計有400人來.
人生一次而已.
他會不會向神祈禱.
神啊你令到限聚令.
突然間變成可以很多人聚集.
但現實沒有.
但是有沒有打擊他的信心呢.
沒有.
他繼續結婚之後事奉主.
他今年年初生了一個B.
不過B早產27週就出世了.
是很危險.
他第二次經歷一個非常大的逆境.
感恩B受到很好的醫護照顧.
上帝的恩守.
B已經出院了.

$^{481}$現在回家了.
最近這位弟兄按目.
按目例裡B第一次抱出來見人.
大家都為他非常非常的感恩.
肅靈的眼睛我們說信心.
其中在列王記下有一個故事.
可能大家還記得.
亞蘭王大軍壓境.
要捉一個神人叫做以利沙.
因為他能夠去令到以色列人對抗亞蘭王.
令到亞蘭王沒辦法攻陷猶太人.
他們來捉以利沙的時候.
以利沙的僕人看到這麼多的敵人.
很害怕.
這是人之常情.
正常.
但是以利沙卻看到有千萬的天君與他同在.
所以他不怕.
在列王記下六章十六節有句這樣的記載.
我們一起讀這節經文開始.
神人說不要懼怕.
因與我們同在的比與他們同在的更多.
那個僕人心裡聽完這句話.
一定說以利沙你瘋了.
因為對方就千軍萬馬.
我們這邊就是道.
以利沙僕人ray.
me都沒有.
怎麼會你說與我們同在的比與他們同在的更多呢.
從眼見到的根本是見不到更多.
在當時那一刻.
當然聖經記載以利沙求神開那僕人的眼.
信心的眼睛見到千萬的天君.
弟兄姊妹.
當你面對人生的黑洞逆境.
不明白無妄之災的時候.
你的信心是你唯一與上帝繼續可以連接的.
一個很重要很重要的一個元素.
當你向上帝用信心去經歷的時候.
有時候上帝會為你開一個你認為很滿意的出路.

$^{521}$有時候上帝沒有為你開一個你認為很滿意的出路.
不過是上帝認為對你最好的出路.
我舞會明年一百週年.
那就是經歷過三年零八個月.
在我們舞會的文獻裡面記載.
在日治的時代.
教會都走了很多人.
這個都是很正常的.
要逃命.
但是教會亦都感恩.
每個主日都仍然維持有聚集.
當然很少人.
有一次有一批的日本兵操入來我們的教會.
大家都很害怕.
不過很奇妙地.
原來在這個會中裡面有位姊妹.
根據文獻記載.
她叫做趙朱竹子姊妹.
她竟然站起來用日文跟那班日本兵對話.
沒有人知道她會說日文.
她站起來跟他們對話.
自然那些日軍覺得這個人會說日文.
就善待了沒有去衝擊到我們當時舞會的聚集.
記載了在我們舞會的文獻當中.
成為了我們很多時候說上帝奇妙恩手的一個案例.
不過亦都很多時候上帝沒有按我們的劇本.
來填寫劇場.
我求主讓我太太多活兩年命.
都沒有.
當時的姊妹我們看教會二千年的歷史.
有一句名言說是容易但很痛苦的.
叫做殉道士的血是教會的種子.
當殉道即是迫伯尼的時候.
是很致命的時候.
當時你一定會覺得是很困難.
上帝的劇本就讓一些弟兄姊妹將生命擺上.
繼續成就這個天國擴展的奇妙計劃.
三個免得我沉睡至死.
這個很可能就是如果上帝不介入的話.
可能我的信心受不了.

$^{561}$我就可能會真的離棄神了.
第二是免得我的仇敵說我聖了他.
即是你的神在哪.
幸災樂禍.
然後我的敵人在我動搖的時候就起落了.
弟兄姊妹.
當去到這個情況的時候.
其實私人面對的迫伯都還未解決.
但他已經進入了第三個狀態.
就是對神發出信靠的讚歌.
我們一起讀第五和第六節.
但我倚靠你的慈愛.
我的心因你的救恩快樂.
我要向耶和華歌唱.
因他厚厚地恩待我.
第六節是很多聖詠團或者音樂事奉者很喜歡的.
我要向耶和華歌唱.
因他厚恩來待我.
不過請留意.
其實在這篇詩篇裡面.
私人轉到這個信靠的讚歌這個位置.
並不是他已經經歷到那些壓制.
在那個時刻離開了.
其實在這些叫哀歌的詩歌的結構裡面.
通常就是遇到困難.
私人的掙扎.
最後憑信.
見機在上帝的信實裡面.
而唱出這些讚歌.
大家留意這裡的經文.
個「但」字是一個非常重要的轉折詞.
是從掙扎到歡呼的一個轉捩點.
轉捩點.
倚靠甚麼.
倚靠是神的慈愛.
即是說他的信心不是掛勾他見到的變化.
是掛勾在上帝的屬性.
神的慈愛.
快樂歌唱是一個未完成的時式.
即是說他還未經歷到這個後因.

$^{601}$他已經唱了.
不過平衡位在個「但」字那裡.
個「但」字是一個完成的時式.
那即是甚麼意思.
意思是說我們的信心一部份.
真的見不到是憑信心去仰望上帝的慈愛屬性.
所以我唱得出這首詩歌.
但同時間你看看你的人生記錄.
我們的check record.
你經歷過上帝無數的後因.
你不要突然忘記了.
所以我們的信心一方面是建基於.
過去上帝的奇妙作為.
在我們的身上.
明年我舞會一百週年.
一定會述說很多上帝的信實.
在我舞會很具體的事情上.
這成為了我們面對未來的一個很大的激勵.
和一個具體的明證.
剛才我說我太太在2016年因為cancer離開世界.
其實在她年輕28歲的時候.
她的甲狀腺已經發現有cancer.
當時醫護做了手術控制了.
當時我聽到醫生朋友告訴我.
我太太cancer.
我很腳軟.
因為我女兒剛剛出世.
怎麼搞啊.
如果我太太有什麼事.
在那一次上帝按他的心意.
讓我們經歷他慈愛的手對我太太的醫治.
以致我太太活多了三十年.
這成為了當時我太太離開的時候.
模擬在中國人的觀念.
五歲而已.
比較年輕.
我們不是很舒服.
不過既然上帝已經給我太太多活三十年.
這個check record這麼清楚.
就讓我和我女兒.

$^{641}$抹完眼淚之後.
繼續走上帝要我們走的路.
這個姐妹叫梁世珠姑娘.
梁世珠傳道.
她是我舞會.
大概比我小六年左右的後輩.
在舞會差六年都算是後輩.
差一代了.
她曾經做過我女兒的兒童團契的導師.
我女兒小時候.
弄她的牛角那個女孩是我女兒小時候.
她人生四十的時候.
蒙召要做宣教的事奉.
她是我們舞會差委會的副主席.
她一心求主開路.
她去到建道神學院.
接受跨文化的訓練.
到她最後那年已經和舞會談.
會不會她畢業的時候.
舞會的差權將她安排到柬埔寨.
去實踐她所領受差權的使命.
也都是我們當時舞會的一個差權策略.
因為我們已經去了.
大概有接近一百人用不同的方式.
去過看柬埔寨.
覺得那個是我們教會應該要發展的宣教工場.
既然自己有姐妹受了訓練.
又配合到大家情投意合.
即是異象相同.
所以大家都沒有簽約.
不過有默契.
她畢業的時候.
我們就差遣她去柬埔寨宣教開放.
她也都有這樣的心理準備.
不過在她讀神學.
最後那一年.
她突然之間患了這個病.
有沒有醫護人員讀得出.
這裡總共有25個字母.
她讀過很多次給我聽.

$^{681}$我也想出來講道講見證的時候.
我講的字我沒辦法讀得到.
那就不如記住林子祥.
簡稱叫LAM.
中文也很長.
叫做肺淋巴管平滑肌增生症.
即是肺隨時會爆肺.
她怎樣發現呢.
她在健道突然之間很不舒服.
去到長洲的醫院急叫直升機去東區醫院.
僅僅救了一條命.
然後就發現了這件事.
是非常之困難.
所以她要休學.
當然她最後也完成了她的課程.
不過醫生跟她說.
梁小姐你的肺這樣的情況.
你以後有一樣事你不可以做的.
叫搭飛機.
那她怎樣去選店呢.
還有不單止搭飛機.
她要拖著一個氧氣.
即是要聞貓鬚.
她講道的時候有個氧氣在下面.
她有不斷的起起伏伏.
她媽媽未信主.
她媽媽跟她說.
你信耶穌信到這樣嗎.
她作為一個傳道人.
她不會用神學的知識回答這個問題.
她很難過.
試過有一次因為病引起很多併發的反應.
她真的跟神說.
你不如接我走.
因為她怎樣可以跟著她這個靈呢.
只有一個方法.
就是緩肺.
她沒有辦法用藥物可以跟著.
很感恩在2021年2月2號.
她突然發通知給很多關愛她的弟兄姊妹.

$^{721}$她說為了祈禱因為她排到去緩肺.
一有新鮮的肺就叫你進去.
所以她好像生孩子一樣.
經常都預備訂袋東西.
一召就馬上去.
她當時在我們舞會裡面做兒童區的傳道人.
她就拍了一條短片說我現在馬上進去.
主她的肺換了.
她現在是唱福音粵曲和唱古典的.
她現在可以享受上帝給她歌唱的事奉.
她是很大聲的.
所以她做兒童侍工基本上不用拿咪高峰.
她經常覺得我回來了.
她憑信經歷了一個非常大的帝國.
甚至被媽媽質疑她.
你信耶穌信到這樣.
信什麼呢.
要到何時呢.
世界的哀歌是一個人在絕望的時候.
孤零零地一面唱一面流淚.
走上樓梯從一路開始往上走越走越高.
極其的悲痛.
可能一時想不開就從高處跳下來.
要到何時呢.
聖經的哀歌不是這樣的.
詩歌中間會出現信靠甚至是讚美.
苦難裡面仍然有生命力和盼望.
就如這位弟兄Nik.
他天生無手無腳.
他的見證一樣.
哀歌所表達的不是一種絕望無助的悲哀.
而是在悲哀掙扎的心境裡面.
仍然產生盼望勝過苦境的生命力.
七年前我去到弘欣開始辦課程.
招生當然不容易.
因為都不知道我們是什麼.
都有人來讀.
我們就見其中黃色圈的姊妹.
她當時來面試的時候.
發生了這樣的插曲.

$^{761}$我們就問中姐妹.
她在一間教會做幹事.
我們就說你做幹事了.
你都叫事奉主了.
你現在想來讀這個課程.
想是怎樣呢.
有什麼期望呢.
她說沒有的.
我覺得我想好好去事奉主.
我覺得你的課程對.
我就來讀.
說著說著的時候.
因為我以前是校長.
又八卦人家的子女在哪裡讀書.
你的小朋友在哪裡讀書.
她說她的兒子在某間中學.
女兒當時大概十一歲.
我說哪間小學.
我想我都認識吧.
然後她說在明愛.
當時我一陣子.
在我認識裡面.
明愛沒有辦小學.
我說哪間明愛.
她說不是啊牧師.
不是明愛小學.
在明愛醫院.
我和當時拍檔一起見的牧師.
兩個打一打.
眼色一擲一擲.
好像碰到別人的傷痕.
不太想再說.
不過她很主動的說.
牧師是這樣的.
其實我女兒出世沒多久.
就變成了植物人.
她在明愛醫院已經住了十年.
當時我們都不知道怎樣回應.
但是她仍然帶著很堅定的信心.
她說我會定時去探望她.

$^{801}$有時候要由明愛去戴口環.
做一些事的時候.
我要陪著她.
這位姐妹經歷了這麼大的困難.
但是她仍然恃至裝備去事奉主.
她的女兒到今天還在明愛.
過了幾年.
她今年九月會回來銜接.
讀我們的學士課程.
她仍然很笑容.
很有信心去經歷.
這根本沒有變到她女兒的狀況.
她仍然憑信唱出讚歌.
她是一個慈善領袖.
很多時候大詩歌都很踴躍.
但是我每一次見到她很踴躍.
我一來就知道她的背景.
我心裡面就對她有很大的肯定和欣賞.
甚至應該說我要向她學習.
要到什麼時候呢.
開始是對神真誠的抱怨.
然後去對神懇切的呼求.
對神唱出順口的讚歌.
弟兄姊妹.
不知道你今天在你人生的境遇.
又或者你的家庭.
或者你的家人.
或者你的社群.
你有沒有一些.
你都在問要到什麼時候的掙扎.
不過透過今天這個訊息.
透過私人的掙扎和她對上帝的信靠.
我們彼此的鼓勵.
一起來去仰望這位為我們信心創始成中的主.
想跟大家唱這首短詩作為一個回應.
\newpage



\section{羅馬書}
\label{sec:3TbpwI2Yugs}
\textbf{2025-09-05 周五晚查經班《羅馬書(21):分題:信心的義與福音的普世性》 講員:曾裔貴牧師}
\newline
\newline
連結: \href{https://youtube.com/watch?v=3TbpwI2Yugs}{\texttt{https://youtube.com/watch?v=3TbpwI2Yugs}} ~~~~ 語音日期: 2025-09-07
\newline
\newline
\hyperref[sec:MB59D1BA0Dc]{\small{< < < PREV SERMON < < <}}
~
\hyperref[sec:index_chronic]{\small{[返順時目]}}
~
\hyperref[sec:index_scriptual]{\small{[返順卷目]}}
~
\hyperref[sec:1GDpbDYTYf0]{\small{> > > NEXT SERMON > > >}}
\newline
\newline
$^{1}$我們一起同心祈禱.
天父我們很感恩.
我們今天星期五晚能夠一起在教會.
在線上或者在頂節會之後聽回錄音.
在不同的地方學習神的話語.
今晚已經是第21講.
第10章的羅馬書.
盼望我們能夠掌握這一點書的主要主題.
以及明白保羅一直要講神的救恩.
是怎樣讓人可以得到.
求主保守我們能夠明白.
也讓我們有這份心得到救恩之後.
都很想和其他人分享.
叫更多人透過耶穌的救贖.
來得到永恆的拯救.
我們今天晚上等候在你的面前.
我們這樣感恩禱告.
奉耶穌基督的名.
阿們.
多謝各位姐妹.
我們改了週五晚.
因為我們週四晚.
現在逢星期四.
都去到一段比較長時間.
我們都要上課到神學院.
在元朗的生命堂借用那個地方.
因為我們有九間在元朗屯門.
和天水圍教會.
加入了一個行在社區福音的課程.
就是陳淑娟博士和愛護神學院合辦的課程.
其實很想推動不同地區的教會.
對自己的社區有認識.
以及找到一些適合的方法.
和教會的資源.
或者和那個社區的需要有關連的.
來在整體的配對.
和動員教會裡面的弟兄姊妹.
在不同的崗位一起去服事.
就會帶來那個福音.
就會不只是一兩個譬如牧者.

$^{41}$或者有心的弟兄姊妹去做.
而是全個教會都一起去擁抱那個福音.
對那個社區有那個關懷.
更加要走出教會去接觸人群.
所以這個課程有它那種很重要的機動性.
第一屆都真的有些教會自己盤點.
自己一起更多的禱告.
那個教會是有一個很大的增長.
當然我們都不是為了增長才去讀這個課程.
而是很想將福音有效傳給社區裡面的人.
所以如果每一個地方的教會.
都正視自己的社區.
可能真的傳香港福音的一個祝福.
整個社會,整個香港的這個地區.
所以我們改了星期五晚就查經班.
就盼望大家都不要介意.
反而對我來說.
因為星期四上午我要預備早上的福音團體.
要預備.
免得有時那個時間如果能夠鬆一點.
我反而那個預備看的資料還可能會好一點.
所以周五晚也有它的好處.
不過現在暫時因為那個課程的緣故.
我們來到講第十課.
第十章的二十一講.
上一次我們前一個星期.
我們講到第九章.
那章很強調神的主權的揀選.
尤其是在猶太人.
保羅看到一個很重要的中心.
為什麼猶太人好像是不得救呢?.
所以神從來沒有保證.
只要你是生下來的猶太人.
你就無論如何.
不管那麼多你都可以得救.
完全違反了恩順清義.
但是原來在保羅一直對創世記的理解.
神給他的啟示.
原來這個救恩.
原來是神在無條件的愛裡面揀選的.

$^{81}$沒有人可以異議說.
為什麼神不揀我這樣.
因為人當然有自己的責任.
所以猶太人還是在背影狀態.
但是第十章就很特別了.
反而講到人的責任就講了很多.
所以要順序的方法.
就是要透過耶穌.
所以第一個階段.
我們看看哈邦.
第一個階段就提到猶太人.
他們不是要以神的義.
來作為他們真正的知識.
反而他們還要立自己的義.
這個就是說他們根本不是按照神提供的方法.
提供的救贖來與神和好.
想用自己.
以為自己一直就已經承繼了波西的律法.
就想用律法來表現自己是神的選民.
就是這樣有那個渲染關係.
有那個傳遞下來的律法書.
或者那些律法帶來的藥.
就會帶來一個很大的落差.
以及對神的真知識就有一個誤解.
所以我們就在第十章就發現.
保羅不再講神揀選的主權.
反而第十章很強調.
神已經為人預備了.
這位降世來到這個世界的耶穌.
他印的是《申命記》的三十章.
講那個可以認識神的門路.
波西就是三十章那裡講的.
但是現在保羅將那段經文.
做一個新的詮釋.
就是講到耶穌基督降到陰間.
升到天上的.
所以不需要到處去找.
上天下地去找.
苦苦地去尋求.
而是那個道就在你口中.

$^{121}$在你的心裡.
就是我們所講的信主之道.
所以這裡就是讓我們看到.
保羅很強調.
就是如果人找到這個知識.
明白這個道.
而他決心去跟隨.
他就可以得到這個永恆的真道.
所以第十章相當強調.
不單是猶太人.
連普世的人都需要自己.
來開放心靈去做.
表面上就是人應該要負的責任.
我們可以逐段去看一看.
我們先觀察一下.
九章神的揀選.
這個重要的部分.
第九章就講了很多.
保羅為他的同胞憂傷.
十章也有.
所以十章的第一段.
和十一章也還有.
保羅一開始轉的新的段落.
他也提到自己跟他的同胞.
猶太人的血緣關係.
那份心靈的互相的緊扣.
特別是在救恩的失落上.
保羅就很焦急.
還有他很憂傷.
所以第六章開始很快就講.
神揀選的原則.
是基於一個無條件的.
因為耳塑和這個啞角.
其實大家都還沒出生.
兩個小朋友還在媽媽的肚子裡面.
根本沒有任何的行為.
誰比較好.
誰沒那麼好.
根本沒有的時候.
神就已經預言給媽媽李伯嘉.

$^{161}$將來大的耳塑要目視小的弟弟啞角.
所以保羅也強調.
不是那個行為.
是在乎神的揀選.
而這個揀選是一個愛無條件的揀選.
所以我們現在回想.
無論在座或者在線.
大家都是已經得救.
我們很清楚.
其實我們不是繞著手去信主的.
我們是感受到一份的愛.
一份的接納.
我們越信主久就越明白.
無條件的揀選.
我們做的其實是很有限的.
但如果我們不進了這一步.
神的工作 聖靈的工作.
神的揀選就不能夠產生那個作用.
但其實歸根結底.
是神對我們罪人的揀選.
所以揀選就肯定涉及到主權的了.
肯定是涉及到神在祂的預知.
在祂的傳能.
在祂的愛裡面來揀選了我們.
以致我們終於能夠得救.
在這個世界出現.
我們能夠得到神的呼喚.
我們就回應神的救贖的選照.
所以我們今天應該讚嘆.
是神無條件揀選我們.
這一份是白白的恩典.
我們應該常常懷著這一份感恩的心.
也不應該反過來說.
活該你死了.
神沒有揀選你.
你又要落地獄了.
你是沒得救的了.
我們就不應該這樣.
應該是我們其實和他一樣.
不同層次,不同方式的罪.

$^{201}$不同方式的罪人.
我們應該帶著一份憐憫的心.
好像保羅所說的.
他為他的同胞憂傷.
我們不是比被信的人好.
我們其實是一樣的.
當然罪是有嚴重的.
有些是沒那麼嚴重.
但是在神的眼中.
我們本身的本質就是徹徹底底的罪人.
我們可能生在一個比較有教養的條件的環境.
我們的表現不是窮凶極惡.
不是一些勢不膽敢.
不是一些走古惑,做壞事.
我們未必是這樣的.
整個的性格,整個的氣質都不是這樣的.
但是在神眼中看我們.
那個本質完全一樣都是罪人.
所以我們能夠信著.
是神那個無條件的揀選.
所以我們要感恩,常常懷著這一份這樣的心情.
以致我們看其他人.
我們就不是看他有沒有揀選了.
因為我們不是神.
我們不知道他有沒有被神揀選的.
我們的責任就是第十章有說.
我們就要認可另外一個叫做.
就是報福音傳喜訊的人.
他們的腳肘就是皆美了.
這個就是第十章.
無論是信的人或者說未信的人要負的責任.
就是看清楚耶穌基督其實就在你的口裡.
在你的心裡的.
但是信了的人我們就要看到.
如果沒有聽見福音.
那些人是不會信的.
所以連同保羅第十章說.
其實需要很多的傳道者.
傳道的人.
我們去說這個皆美福音的人.

$^{241}$那些人的腳肘都是皆美.
所以你可以想像到.
保羅多麼重視宣教.
所以第九章我們來看的時候.
這個就是揀選涉及到神的主權.
這個就很抽象了.
很神學了.
甚至是很奧秘了.
甚至是令到那些不信的人.
很鄙視我們.
越說得重越說得嚴重.
越說得很決絕.
是啊你肯定沒得救了.
因為神沒有揀你.
但是這個說話我們替神說是不公允的.
因為神從來沒有告訴我們.
誰是A君不得救.
B君得救.
我們從來不知道的.
所以我們是人.
我們是信了耶穌.
我們的責任就是要傳福音.
不是去宣佈.
對不起誰你是沒有被揀選.
所以我不會向你傳福音的.
我們就搞錯了.
不是這樣.
所以我們記住第九章.
保羅想展現猶太人被拒絕.
其實不是神拒絕他們.
是他們一直以自己亞伯拉罕後裔血緣的身份.
一直守著那些律法.
以為這樣就是合標準.
反而還帶著一個鄙視外邦人的心.
即是驕傲.
因為他有一套東西.
每個星期就是這樣做.
一套宗教的行為.
這套宗教的行為.
令我很安心.

$^{281}$很穩定.
我一定可以上天堂.
做好不做可以不舒服.
你就可以想像到.
其實他們不是在拜神.
拜這位造物主.
是在拜禮賢的宗教.
猶太教的憲制.
食物的條例等等.
所以第九章給我們有一個概念.
但第十章就開始講到.
人應該要負的責任.
我們在第十章開始.
我再看第九章三十節.
以色列人被叛島.
就是保羅引述的.
以賽亞書的二十八章.
就是這個絆腳石.
當然就是夏文.
保羅再用另一個角度去引述.
就是申命記三十章.
去到不知哪個地方去尋道.
這個道不用你去尋.
在我們每一個人.
生而為人類.
其實真的舉頭三尺.
有神明這句說話.
有它的道理.
但保羅講的.
每一個人只要他呼求這位耶穌.
就可以得救.
十章我們稍為讀這段.
到第十一節.
這段特別強調.
就是責任他們沒有負.
其中就是因為他們還擁有.
或者拿著一些知識.
拒絕神的救恩.
提供給他們的方法.
第十章第一節.

$^{321}$弟兄妹.
我心裡所願的.
向神所求的.
是要以色列人得救.
我可以證明他們向神有熱心.
但不是按著真知識.
所以這就是異端出現問題.
非常熱心.
逢星期六的早上.
在我們洪水橋的隧道.
有兩個人站著.
好像不動了.
他不會跟你說話.
除非你跟他說話.
他就在那裡放個書攤.
在那裡派傳單.
他們每個星期都要這樣做.
我又沒有刻意去留意他們.
或者去挑戰他們.
我不知道是不是每個星期都是那兩個.
總有一男一女.
穿著襯衫.
西褲.
很斯文地站著.
派守望台的傳單.
很有報道的熱心.
但他們不是按照聖經真理的真知識.
因為他們傳遞的耶穌.
不是聖經所說的耶穌.
他傳遞的福音.
是有那節扣老曲的.
耶穌不是獨一的真神.
所以我們要留意.
特別是針對猶太人.
猶太人有全套的律法規模.
有對律法的認識.
甚至他們都認為律法是被小朋友引路的.
是那些蠢笨人的師傅.
是那些外邦人的光.
但是他們都只是有自己那一套的知識.

$^{361}$所以不是真確帶到他們去信耶穌的知識.
所以這個真的要拆開.
我們真的要不斷更新我們的神學.
我們不要有自己那一套的觀念.
我們這樣做就行了.
我都說過笑.
當然不應該是說笑的.
有些人去拿聖餐.
目睽睽看不到.
拿多一兩塊拿回家.
好像當那個是聖餅.
回到家裡給那些沒有回教會的.
他吃就會神聖了.
就會無福了.
其實那個餅不會令你無福的.
這個就是神學的知識不足夠.
因為他覺得在聚會很感動.
在那個聚會很有靈受.
還有牧師祝謝了那個聖餐餅.
又在那一刻有種特別的意義.
就是耶穌的靈臨在聖餐的聚會.
其實不是聖餐餅.
是整個聚會在眾弟兄姊妹的心靈裡面.
所以嘉義文說的.
聖靈的臨在是比餅變質重要的.
所以我們應該要拜的就是.
聖靈在整個聖餐裡面運行.
去感動我們對主耶穌的敬拜.
那個就是臨場.
所以如果你不是臨場.
你回到家裡.
阿仔這個餅是曾牧師祝聖過的.
快點吃啊.
其實沒有用的.
如果越來越迷信.
那個熱心就變成一個錯的知識.
吃了那個餅分分鐘有危險.
因為不是那個餅被人詛咒了.
也不是有什麼髒.
而是因為你錯誤了那個理解.

$^{401}$有時候更加會更加迷信中的迷信.
這是不好的.
所以我們一定要搞清楚.
這個只是其中一個例子.
說的是全套的質理.
我們自己都要不斷的更新.
所以不要變成一種質誤.
一種基督教自己的迷信.
譬如某一個傳統是好的.
擺這棵花這樣擺就一定是最好的.
別人動了就說他死過.
然後你就說不要亂搞.
一代一代傳下去.
就變成那棵花擺的位置.
就變成一種迷信.
這樣就會帶來更多的爭鬥.
更多的問題.
所以這個我們要記得.
那個熱心是要配合真理的知識.
本來這個是很平衡的.
然後開始說什麼是假的知識呢?.
或者錯誤的信主的知識呢?.
就是第三節了.
因為不知道神的異.
其實神的異就是藉著耶穌基督.
來拯救三章21節.
三章23,24節.
不知道神的異有兩個解釋.
是故意的不知道.
因為福音已經傳給猶太人了.
由耶穌離世之後.
很多的使徒們已經在耶路撒冷傳道.
為耶穌作見證.
所以當保羅寫這段經文的時候.
寫給外邦人的時候.
就不是猶太人不知道.
是故意的不知道.
另外一個就是堅持自己那套.
不要神那套.
覺得不應該這麼簡單.

$^{441}$耶穌.
也就是說連帶的就是拒絕耶穌這個彌賽亞身份.
因為他一直都咬著耶穌只是那個木匠的兒子.
這個就不是真的知識.
這個要怎樣打破一個很大的靈魂的手術呢?.
保羅就是一個示範單位.
他自己就是故意不知道.
還有否定耶穌的身份.
所以當他去大馬士革.
敘利亞差不多到城了.
就在路上耶穌向他顯現.
光照他.
令他三天三夜都看不到.
那麼大的心靈手術.
就是耶穌直接進入保羅的心靈裡面.
告訴他你逼迫的那些拿撒勒人.
或者那些傳拿撒勒人是耶穌的.
拿撒勒人是耶穌是尼塞馬的.
那個就是我了.
保羅就突然間晴天霹靂.
所以這個突然間灌注給他的真知識.
就將他整個舊有的那個所謂城堡.
那個觀念.
那個他自己的防衛.
完全的崩潰.
他要重新整理他自己的神學.
所以神都容許他有三天三夜禁食.
眼瞎了.
讓他在家裡.
在那個地方去左思右想.
去想想自己為甚麼會這樣呢?.
那個復興就來了.
那個叫阿拿尼亞的弟兄就來找他.
說是耶穌差遣他來.
叫他可以看見.
正式呼召他做傳道.
所以這個真知識.
第三節保羅自己去回顧解釋.
他自己就是過來人.
他有對迫害耶穌的熱心.

$^{481}$但沒有對神的熱的真知識.
第三節因為不知道神的義.
想要立自己的義.
就不服神的義了.
自己的義就是靠自己的行律法的好行為.
哪些的食物如果是利命計劃.
不乾淨不潔的.
他們就不吃了.
所以那些肉類的處理.
是所有肉都要放血.
不可以帶血吃的.
你可以想像得到.
他們做到加零一.
守這些律法.
他不理那個時代.
不理那個時代的轉變.
都要繼續這樣做的時候.
有時就變成矯枉和正了.
所以這個就是保羅很痛心.
和他們專注在那些很微小的事情上.
而他們沒有注重的.
如何能夠與神建立真正生命的關係.
就是要獲得神的義.
如何獲得呢?.
我們揭開三章.
幫大家溫習一下.
三章的23節羅馬書.
保羅早就說了.
開宗明傑已經說了.
人是虧缺神的榮耀.
人是在神的憤怒下.
但現在神在救恩.
在律達以外.
就已經顯明了.
顯明了什麼呢?.
顯明了他的義.
三章的第21節開始.
但如今神的義在律法以外.
已經顯明出來了.
有律法和先知為正.

$^{521}$即是說原來一直基於舊約的律法書.
不是全新沒有人認識的.
是一直都在舊約有這樣教導.
有這樣預言.
有這樣講到的.
所以現在第22節.
就是神的義.
因信耶穌基督.
加湊一切相信的人.
並沒有分別.
因為世人都是犯罪.
即是不可以再靠自己了.
虧缺了的是神的榮耀.
但現在蒙神的恩典.
因基督耶穌.
所以基督耶穌就是救贖的源頭.
也是神的義最重要的指標.
所以剛才我們還沒看第4節.
律法的總結就是基督.
所以原來律法是帶給我們.
只是知罪之後.
就要去到耶穌.
這個救恩的源頭裡面.
去得著這個救贖.
所以當我們虧缺了神的榮耀之後.
如今卻蒙神的恩典.
因基督耶穌的救贖.
就可以白白獲得神的義.
獲得神的義.
也等於神不再看我們是罪人.
將我們的身份轉變.
所以這就叫做白白的清義.
即是宣告無罪.
以及得著神的公義.
成為神的兒女.
即是身份的立時轉變.
不再是神的敵人.
不再是普通的跟神無關的人.
現在成為了神家裡面的人.
即是神的兒女.

$^{561}$而且神的兒女還要是被清義.
不再有罪的一個義人.
所以原來一直都在說.
神的義就是要透過耶穌基督.
來得到.
所以當他們要得到的是律法的義的時候.
他們就不會再服神的義了.
所以十章第三節.
不知道神的義.
所以要立自己的義.
就不服神的義了.
所以你看到第四節立即就說律法.
律法不會帶給人有神的義.
律法的總結.
即是說律法真正的功用.
其實就是帶我們去認識基督.
使得相信基督的人.
都得到這份神的義.
所以其實是相當吻合.
第三章已經說得很清楚.
來到第十章再說猶太人不得救.
他們不肯負責任的時候.
不如這個責任怎麼負呢?.
第五節開始就說了.
保羅立即就說.
這個信主的道.
這個所謂的行為.
其實都不是我們能夠負的.
這個信已經在我們的心裡.
在我們的口裡.
只要我們呼喚.
我們就可以得到.
問題就是去到後面.
保羅邀請很多羅馬的信徒.
必須要做一個傳道者.
因為真是的.
報福音傳喜信.
十章尾應該是十三節.
腳中是皆尾.
你想不想你的腳中皆尾呢?.

$^{601}$你走的路都是無福的.
好像滴滿了滋味的街道道路.
腳中我們人生.
其實都是在走一個人生的旅程.
但是我們的旅程.
有些人一路走.
走到滅亡.
個人滅亡.
自己也帶不到什麼.
可以去到永恆.
帶到的不是好的品格.
不是好的名聲.
但是有人在.
那些惡人死.
表面可以很風光.
很多人為他哭.
但是他自己是一個罪人.
他一生犯罪.
做盡壞事.
你可想像得到.
他的腳中.
不單止不是皆尾.
是邪惡.
聖經保羅堂的腳中皆尾.
這些就是我們真的要羨慕.
要學習.
如果我們信徒主乘二十年.
都不肯去參與教會的傳福音.
或者自己都去參加傳福音.
其實你的腳中.
只是你在走自己的路.
我不能夠說你是在走犯罪的路.
但是肯定沒有保羅所描述的那麼好.
所以就開始說了.
律法的總結.
就是帶我們去認識.
自己的不濟.
自己的無助.
律法反而不斷光照我們那個無助.
帶我們去看到.

$^{641}$信心的意義就是耶穌基督.
使到凡人相信.
因為看到自己的無助.
這個就是信心的呼喚.
信心的呼求.
呼求誰呢?.
就是信這位主耶穌.
就是這位尼塞馬.
所以很適合神的意義.
就是透過耶穌基督.
就可以得著.
就可以得救.
就可以身份的轉變.
所以我們要知道.
信祂的都得著.
這個意義是神的意義.
因為上文一直都在說.
包括三章都一直在說.
他們不服神的意義.
好了.
怎樣可以得到這個信心呢?.
律法讓我們看到自己的罪.
但是如果你沒有一個正確的方法.
就是說你的鏡頭如果是模糊.
你都不能夠讓你去找到那個信的道.
所以保羅在這裡很特別.
他突然間引了利未記和申命記兩段經文.
第五章開始.
摩西寫著說.
人若行雅出於律法的義.
就必因此活著.
利未記的十八章.
不如我們來看一看.
希望大家能夠對舊約的時候.
現在都可以看到一些印象.
崇拜的時候都比較難說這些經文.
當然不是說這些經文不重要.
十八章.
耶和華對摩西說.
你曉遇以色列人說.

$^{681}$我是耶和華你們的神.
你們從前住的埃及地.
那裡人的行為你們不可效法.
我要你們道的迦林地.
過去未來.
過去就是埃及.
將來就是迦林.
因為現在利未記就是在西乃曠野的一個月.
出了埃及的一個月的時間.
應該是第二年的一個月.
這個就是利未記的經文.
然後就說.
第四節.
你們要遵我的典章.
守我的律例.
按錢而行.
我是耶和華你們的神.
所以第五節.
所以你們要守我的律例典章.
人如果遵行.
就必定因此活著.
我是耶和華.
是相對於迦林地.
埃及地那些邪風惡俗.
可能拜偶像.
交軌.
巫術.
很多不同的人亂.
但是聖經波西就說.
你們進去之後.
因為十戒.
十戒所衍生的戒條.
將會陸陸續續出現.
就是說要行一個新的道德觀.
就是跟隨神.
因為他們已經是神.
在迦林.
應該叫做在曠野.
這個西乃曠野.
跟他們納了約.

$^{721}$他們就是從這一刻開始.
是神的選民.
他們已經做好了那個螢幕.
已經第四十章.
那個回望已經安好了.
開始有些擺設.
所以利未記是接續著.
第二年的第一個月.
蒸奪的時候.
摩西就開始設立.
祭司的職位.
開始選召大祭司.
首任的大祭司就是亞倫.
他來高末他.
在他能夠正式就職典禮之後.
就開始幫以色列人贖罪.
所以在曠野.
第一年就已經開始有憲制的事.
所以相對於過去的埃及.
和將來的迦林.
他們不是神的選民.
他們是沒有神的戒律.
他們只是隨心所欲.
按照自己去定的.
可能是一些很惡劣的律法.
例如塔利班這樣才算.
還是很原始的.
還是很極權的.
還是做一些很沒有人道的一些律法.
猶太人他們進入將來的迦林地.
他們要守的是新的道德律法.
神的國度的律法.
所以聖經就說.
人如果遵行.
即是說按照神的真理去生活.
用心靈去生活.
你就能夠活著.
得到神的喜悅.
和因為這樣而帶來的信心.
這個就是摩西在第五節這樣說.

$^{761}$其實應該說保羅引述摩西十八章的領備記所說.
現在就開始說.
哪裡找到這個信心的義呢?.
第六節開始就說.
唯有出於信心的義如此說.
相對於律法的義.
其實十八章的第五節已經說得很清楚.
很清楚是神的法則.
神的要求下的屬靈生活.
猶太人他們就會得到.
找到信心的門路.
即是說一直過一個神喜悅的人生.
有這樣的人.
例如在事先的時代.
有那些事先.
例如鵝駝獵.
例如有其他一些的人.
他們都是可以帶領那個時代.
回到神那裡.
例如以利亞.
又是在阿哈耶駒別.
叫全國都拜偶像.
拜阿斯塔羅.
拜巴力.
但是以利亞就是要他們不要生存良異.
要拜要敬畏.
要跟隨耶和華.
所以不同的時代.
他們都找到自己律法的義.
但現在保羅做這個對比就更加強烈.
先引了裡面十八章提到的那個律法的義.
即是人如果真心去做.
都尚且找到.
但保羅說現在不用那麼麻煩.
因為那個舊恩已經出現了.
耶穌已經降生了.
所以他就開始說第六節.
其實這裡幾可圈可點的.
因為三十章的申命記.
就是摩西臨終的時候說.

$^{801}$你們不用到處去找.
因為我快要離開這個世界.
你們怎麼辦呢?.
其實摩西就說這個道就在你們的口裡.
在你們的心裡.
問題是你們想不想要不要得道.
我們就用新約保羅引述.
保羅將這段經文重新演繹.
就很可圈可點了.
第六節.
唯有出於信心的義如此說.
你不要心裡說.
誰要升到天上去呢?.
就是要領下基督來.
誰要下到陰間去呢?.
就是要領基督從死裡上來.
即是說保羅這裡很清楚.
他對耶穌的認識.
耶穌是誰來的?.
耶穌是從天上永恆的居所.
神所住的地方.
來到我們的地球.
地上來到這個拿撒勒.
來到這個伯利恆降生.
基督現在在哪裡呢?.
他是從死,復活,永活.
回到他本來的地方.
天堂.
所以原來耶穌基督就是這樣的一位救贖主.
即是說連帶他救贖的人生.
他所做的才是整全的救恩.
不只是信耶穌這個人,這個名字.
很多人都說你說的耶穌是很好的.
很影響到很多基督徒,很多教友都很好.
不過那個是耶穌.
跟我無關.
但是耶穌和他的十字架.
才是保羅所傳的一個神的義.
因為只能夠主耶穌和他的十字架.
才可以將我們拯救.

$^{841}$所以你看到這裡.
信心的義就是相信耶穌由天而降.
相信耶穌由死復活.
即是為我們釘十字架的由死復活.
以及永活將來審判活人死人.
所以這段聖經的源著在哪裡呢?.
為什麼這樣說呢?.
這樣說有什麼目的呢?.
保羅怎樣將它轉化成為信主的道呢?.
我們要回到《申命記》的30章.
因為我的舊的串珠.
有很多的經文.
會立刻告訴你.
這段聖經其實原來在舊約哪裡出現的.
所以也不是我的記性特別好.
當然某些經文都是.
但是絕對不可能是每一段我都記得.
所以大家如果想讀聖經熟悉.
就要買一本這樣的聖經.
手機肯定很難幫你找到.
但是大本的聖經.
有串珠就可以立刻參考.
舊約哪裡說過呢?.
我就記住它.
當然記住也不夠.
就是兩段聖經.
新約的作者.
為什麼這樣引用.
這樣引用之後有什麼目的呢?.
近幾年石榴園就有這一科.
新約引述舊約.
一個專門的學科.
在聖經研究裡面.
其實不是說近幾年才有.
而是開始有人會寫這一類的書.
很厚的.
很大本的.
外國已經有.
我好像給大家看過了.
中文翻譯了兩大本的.

$^{881}$就是將新約引述舊約的.
每一段都解得很透徹.
我們揭開《申命記》.
這些功夫如果你喜歡.
其實當你解通了.
你會有很大的滿足感.
因為如果解完.
其實你不知道要說什麼.
為什麼無端端會引用這些經文呢?.
什麼上天下海.
第30章.
因為第30章是說聽話就無風.
不聽話就就坐.
是說最後的叮嚀.
摩西他快要死了.
他很快就要在尼泊爾告別人世了.
但是他仍然帶著一個語重心長的長者.
看到這班後輩.
20歲以外的.
在抗日剛剛出生.
上一輩已經死了.
他們現在長大到20歲.
因為過了差不多40年.
這個時候他們就是20歲以外.
即是大約60歲以下的人.
就是現在這班新生的世代.
所以摩西就非常語重心長.
最後第15節.
30章的15節.
《申命記》.
他說漢啊.
我今日將生與福.
死與禍.
第一個措辭.
所以寶萊人這段經文.
差不多是遺足.
人用這段經文.
摩西當時的情詞迫切.
來告誡自己的下一代.
你們一定要入到迦林.

$^{921}$你們一定要無福.
你們不聽話.
是要就坐的.
所以寶萊就說到救恩的重要.
怎樣可以得到信主的道呢.
留意這段經文的重點.
15節.
生與福死與禍.
神明在你面前.
吩咐你愛耶和華你的神.
遵行他的道.
謹守他的誡命律例典章.
使你可以存活.
人數增多.
耶和華你的神就必在你所要進去.
得為業的地上.
賜福與你.
彷若你心裡偏離.
不肯聽從.
卻被勾引去敬拜.
事奉別神.
我今日明明告訴你們.
你們必要滅亡.
在你過躍但河.
進去得為業的地上.
你的日子.
必不長久.
我今日呼天喚地.
可以想像到那個氣氛.
那個措辭.
那個態度.
我呼天喚地向你作見證.
我將生死.
禍福.
神明在你面前.
所以你要揀選生命.
使你和你的後裔.
都得存活.
且愛耶和華.
你的神.

$^{961}$聽從他的話.
專靠他.
因為他是你的生命.
你的日子長久.
也在乎他.
這樣你便可以在耶和華.
向你列祖.
亞伯拉罕.
以實瓦各.
起實應許.
所賜的地.
那裡來永居.
講了利害關係.
信者就蒙福.
拒絕摩西的警告.
不聽話.
不守神的誡命.
或者說不愛這位耶和華.
只會註定.
一生欺負.
受咒罪.
但仍然問那個問題.
信主之道在哪裡.
我們未讀的.
就是11節這一段.
回到11節.
我今日所吩咐你的誡命.
不是你難行的.
也不是離你遠的.
不是在天上.
使你說.
誰替我們上天取下來.
取經.
取道.
使我們聽見可以遵行.
其實一直都在講.
這個律法書.
就是神頒佈.
你只要根據你的心靈.
去敬拜神.

$^{1001}$去跟隨神.
你就能夠蒙福.
所以道就是你一直認識神的真理的時候.
你就會具備.
這份愛.
和信.
摩西就說.
這麼好的道.
我去找.
你告訴我.
我一定要去找.
天上我都上去.
誰替我取下來.
我們就會遵行.
這就像是一種切文.
像是Ralph他們那樣.
摩西說不需要.
不是在天上很遙遠.
你們拿不到.
拿不到不關我的事.
不是.
下面還有.
下面的第十二節.
十三節.
也不是在海外.
使你說.
海外就是說當時他們對地中海.
他們看為.
是一個很大的海洋.
海洋在古代的人.
有鬼怪.
有惡獸.
海怪.
所以要航海.
其實也不容易.
如果在海外.
我們怎樣去拿.
很危險的.
我們哪裡找.
大海茫茫.

$^{1041}$摩西就說誰替我們過海取下來.
使我們聽見.
我們就會相信.
其實根本不是這個意思.
摩西說.
我已經將生死禍福.
在你們面前.
你們想聽就一定會得到.
所以根本不用上天下海.
這樣去說.
第十四節.
就是那個總結.
這話卻離你甚近.
留意了.
就在你口中.
在你心裡.
使你可以遵行.
保羅引述十一至十四節.
就是這段羅馬書的十章第五節.
保羅引述這段.
誰去升天上.
替我們找到道.
就是基督從天而降.
保羅這樣告訴我們.
就是領下基督來.
誰領下基督來.
當然就是神的旨意.
差遣耶穌來這個世界.
所以原來第十章.
這裡的救恩.
原來就是神差遣他的愛子.
這個就是道.
所以猶太人說.
我們不知道什麼是神的義.
我們不知道什麼是道.
其實連摩西在三十章的申命記.
都已經說了.
當然當時還沒有基督降世.
但是律法的總結就是基督.
就是說如果你真的用心去愛神.

$^{1081}$讀神的誡明的律法.
你是可以認識自己的無助.
需要救恩.
你就會看到基督的原備救贖.
不過不是用名字叫耶穌.
因為耶穌還沒有出世.
還沒有降世.
但是保羅這裡就很清楚.
和羅馬的弟兄姊妹說.
其實原來舊約已經有神的差遣.
現在基督就是被領下到地上來.
就是說這個道.
耶穌基督就是神的道.
顯然在地上讓人可以認識.
接著.
誰要升到天上去.
就是領下基督來.
誰要下到陰間.
剛才說的是海外.
更加清楚.
保羅來將一些本來不可能的事.
告訴我們已經成就了.
原來誰下到陰間呢.
就是基督從死裡上來.
復活.
就是告訴我們.
我們不要再有任何藉口.
你給我看耶穌.
我就會相信.
耶穌已經降世成為人.
也已經由死復活.
沒有人可以在死亡裡復活.
沒有人可以被陰間拘禁.
可以來復活.
耶穌已經完成了.
所以這個就是.
到底摩西怎麼說.
保羅就引述了.
這道離你不遠.
正在你口裡.

$^{1121}$在你心裡.
心靈記三十章.
他直接引述.
接著保羅就將這句說話.
開始演繹了.
這裡就告訴我們.
就是我們.
留意到嗎.
就是我們所傳信主的道.
這裡真的很厲害.
我之前都沒有這樣的發現.
今天再看就有一個很大的發現.
原來第十章.
不單呼召外邦人.
不要小看猶太人的歸主.
雖然他們現在是拒絕神.
但是原來傳這個字.
是很突出的.
如果沒有人去傳.
這個道就不能進入到人的心靈裡.
就不能夠在人的口.
在人的心裡.
去呼求.
所以就是保羅說的.
我們所傳信主的道.
所以宣教士用他的遠方.
真的有他的原因.
有他值得我們去敬重.
支持的地方.
今天大家都看到.
《野地牧人》出了一些資料.
就算我們去不到.
我們都應該用心靈.
用金錢.
用行動.
去支持這些福音工作.
保羅說如果沒有人傳.
人們是聽不到這個道.
就不能夠在他們的口中.
在他們的心裡.

$^{1161}$轉化成為他們向神的求告.
所以這就告訴我們.
就是我們所傳的.
信主的道.
信耶穌的道.
需要有人傳.
所以你可以想像到.
當有人傳.
當是傳信耶穌的道.
接著就.
即是到第九章.
你就是那個聽到道的人.
你就考慮.
以耶穌為主.
心裡信神.
叫他從死裡復活.
就必定得到.
所以我們要令到人.
去到一個他自己向神求告.
所以個人的談道傳福音.
一定要他自己去.
跟你祈禱.
向你.
來到去.
跟著你向神祈禱.
然後他自己承認是罪人.
呼求耶穌的救恩.
那就立即.
聖靈當然他是真心祈禱.
聖靈就進入他的心.
你就得到一個果子.
成為神養.
所以第九章.
十章九節.
很重要.
不只是教我們.
去傳福音用.
而是在告訴我們.
神的義.
已經透過耶穌降世.

$^{1201}$即是神差遣.
耶穌降世.
成就了救恩之後.
神的義已經具備在.
現在地球上.
充斥著在教會.
我們去傳道的時候.
所以誰人.
聽明白了神的義.
他就可以得救.
所以教會是傳遞.
神的義的地方.
不過更多是傳遞.
是非的地方.
因為太密集太熟.
變成我們有時會掉以輕心.
慢慢.
我們不是集中用心.
去傳揚.
這個義.
而是傳揚我們自己的見解.
那就.
難怪有時福音的大能.
不能夠.
大量彰顯在教會裡.
所以我們很需要求神幫助.
所以一說到.
我們所傳的.
信主的道.
立即就引到第九節.
就是說.
當這個道在你的口中.
你去祈求.
你就得救了.
你考慮你的誠意.
心裡面相信.
即是說.
真的只需要你.
宣認耶穌是你的救主.
是你的生命的主宰.

$^{1241}$你就立即得救了.
你想想.
這個這麼戲劇性的過程.
是怎樣發生的.
有沒有什麼步驟.
然後立即說.
就告訴我們.
第十節.
因為人心裡相信.
就可以清義.
口裡承認就可以得救.
其實清義承認.
清義和得救是兩個.
詞彙其實不同的角度.
看那個人的.
同神關係的復和.
所以口和心.
是要一致的.
所以不只是.
口認心不信.
兩者是要共同的一致.
所以我們不要.
操之過急.
總之他講了就行了.
但一定要他的內心.
那一刻是很坦誠.
所以.
陪談兩個人.
有時都要很安靜的環境.
讓他的心.
認真去想這件事.
然後留意到.
第十一節.
經常說凡信他的人.
必不至於羞愧.
二賽馬書的二十八章.
就是講盆石.
我在石安放了一塊盆石.
信的人必不羞愧.
或是必不動搖.

$^{1281}$也是一樣.
保羅不要以為他亂引聖經.
原來立即就在講.
他的信心.
必不羞愧.
其實就是建立在耶穌這個大盆石.
所以你帶人信耶穌.
去到某個階段.
你要交給耶穌.
你不要成為他的主人.
你做Mentor沒問題.
但你不要操控他.
操縱他的信主.
你的罪要比他更大.
因為他要在盆石上.
穩固的信心.
才必不動搖.
所以.
保羅突然這樣說.
信耶穌的人.
必定不會羞愧.
其實就是二十八章.
不如跟大家再引述.
二十八章這一篇.
二賽馬說.
在石安有一個盆石.
是穩固的根基.
信的人必不至於羞愧.
二十八章.
這些經文其實都重要的要熟悉.
二十八章的十四節.
所以.
你們這些洩漫的人.
二賽馬書在聖經中間的部分.
因為實揭了.
因為他會說六章.
就是.
「合本住在耶路撒冷者百姓的」.
就是那些官長.
要聽耶和華的說.

$^{1321}$「你們曾說」.
都很誇張.
「我們曾經與死亡納入.
與恩奸結盟」.
就是沒有人可以動他們.
就好像內地.
有時候那些.
那些叫什麼.
就是.
去找那些.
那些吳盆的官.
我忘記了.
那些名字.
「敵軍如水漲漫經過.
的時候.
必不留不到我們」.
就是他們真的很有恃無恐.
敵人來到.
就像水漲.
洪水.
都不會淹沒我們.
「因我們以謊言為被所謂.
以虛假藏身」.
就像隱形了.
不需要被人找到.
甚至不會被人找到假象.
不過第十六節.
耳上就這樣說.
「所以主耶和華如此說.
看啊.
我在石安放一塊石頭.
作為根基.
是試驗過的石頭.
是穩固根基.
寶貴的防角石.
信靠的人必不著急」.
你們可以想像得到.
這一節聖經的重要性.
就在舊約.
以賽亞時代.

$^{1361}$已經預言.
彌賽亞就是人類的穩固根基.
所以救恩.
是要建立在.
神兒子的生命上.
所以我們作為.
傳福音傳報的人.
我們要記得.
一生要將人帶去.
耶穌的救恩.
讓他找到.
耶穌穩固的盆石.
我可以用這樣的例子.
其實每代都是這樣的.
所以我們.
有些警惕.
我初信的時候.
很多我和我一起.
回教會的年青人.
和比我更厲害的大學畢業.
大學生.
五六個根的同學.
回來我們教會的團契.
很熟悉.
比我回更久.
那時候我中學.
勉強畢業.
這班大學生.
很有知識很有學問.
回到團契.
但慢慢地.
開始工作各方面.
不再回教會.
完全不再信.
完全不信耶穌.
原來他沒有跟隨在耶穌的救恩裡.
只是在團契.
只是在開心.
在關係.
但當有少少變化.

$^{1401}$你出來工作.
和信耶穌是兩件事無關.
應該更加依靠耶穌.
連信仰都不要.
我真的約過一個同學.
他後來做信仰中心.
剛剛入職.
和他聊天.
問他為何不回教會.
好好的弟兄.
但他都不再信.
甚麼原因令到連耶穌都不信.
即是說.
根基建立在哪裡.
建立在好好玩的團契生活.
每個星期有周會.
每個月有茶點.
有生日會.
這些本來不是錯事.
是讓我們更加依靠神.
但這裡就告訴我們.
原來二次世界已經說了.
真正的穩固.
我們心靈的牢固.
是要在這個盤石.
在這個防角石.
在這個被試驗過的石頭上.
這個石頭肯定是指耶穌.
所以保羅告訴我們.
人凡相信的人.
必不至於羞愧.
所以是相信耶穌.
所以我們要記得.
教會不是很好.
我們很有愛心.
我們有團契.
我們有愛言.
我們有很多不同的活動.
我們有跳讚美操.
很好的教會.

$^{1441}$但最終我們帶不到他去.
找到耶穌.
是他穩固的盤石.
他始終就不穩固.
他的信心不穩固.
所以這裡說的凡信.
他的人.
必不至於羞愧.
那個大前提是.
基督是他心靈的盤石.
是他的安穩的避難所.
所以弟兄姊妹.
我們真的要在耶穌裡.
建立我們的信仰.
教會是一個.
提供給我們認識耶穌的地方.
所以後面就是說.
我們必須要讓人.
找到耶穌.
第十二節.
猶太人和希臘人並沒有分別.
因為眾人同有一位主.
他也厚待一切.
就是猶太人.
希臘人.
所有的人類.
希臘人只是一個集體的稱呼.
保羅的內心.
總之地上.
當時不是那麼多的飛機.
到處去.
又有AI.
不知道有多少人種.
保羅.
當他知道所有的人種.
假如他知道有中國人.
中國人也沒有例外.
凡求告耶穌.
就必定得到.
意思就是這樣.

$^{1481}$這裡更加告訴我們.
原來這個信心的意義.
原來這個道.
在口裡在心裡.
是說所有地球上的人類.
哪些人需要救恩.
就是未聽聞耶穌救恩的人.
不知道耶穌來這個世界的目的的人.
所以十二節已經說了.
求告主的人.
必定會得到.
這是第十三節.
神會厚待求告的人.
厚待.
所以初信的人.
特別在祈禱裡有經驗.
有經歷.
即是單純對祈禱很敏銳的人.
上星期我的神學院同學.
其實我的學子.
高我兩班.
他說了一個見證.
祈禱的見證.
他說他讀短宣中心.
的時候沒什麼錢.
有一天.
很感懷身世.
吃早餐.
他說他祈禱.
他很想.
其實很想.
吃一個很好的早餐.
可能全餐.
即是有蛋.
有通粉.
有香腸.
他不敢出外.
好像很勉強神.
突然有個心願.
神說可以祈禱.

$^{1521}$他真的祈禱.
神給我一份早餐.
他開始回神學院.
在沙田.
回到去.
見到有一個神學院同學.
拿著兩大袋的早餐.
匆匆忙忙.
進去餐廳去吃.
可能是幫人擺.
他就.
酸溜溜.
自己沒有早餐吃.
因為沒有錢.
很想吃早餐.
剛剛祈禱.
走進去餐廳.
突然神學院同學.
跟他說.
阿晶晶.
今天神感動我.
買了一份早餐.
是給你的.
我買的時候神感動我要買給你.
阿晶晶整個人大眼睛.
眼都大了.
剛剛才祈禱完.
其實沒有串通.
不是裝模作樣.
突然想哭.
神聽到這樣.
在他未祈禱之前.
已經感動那個姐妹.
去買一份早餐.
平日不會這樣買.
原來有些人的心靈單純.
可以這樣祈禱.
很奇妙.
大家可以聽上星期.
蒙恩堂的廣告.

$^{1561}$很好.
這裡告訴我們.
神會厚待一切.
求告祂的名的人.
厚待.
我們會厚切豬扒.
厚多士.
厚豐厚的獎賞.
但我們好像.
不太善報.
神的厚待.
求告.
所以我們如果愈來愈喜歡祈禱.
就祝福你.
因為神會厚待.
祈禱的人.
我們多去為人祈禱.
將人的心.
將人的名字掛在祈禱裡.
將他帶到施恩寶座.
他一定會蒙福得救.
所以有一個例子.
真人真事.
大衛李雲斯敦.
在非洲.
他是一個探險家.
有一個很厲害的傳道人.
去非洲傳道.
學那些非洲話.
然後進去中心.
有一次.
有隻獅子撲過來.
和他搏鬥.
他就用盡力.
打了他之後.
剛好有槍.
打死他.
他很慶幸.
能夠這樣解救.
因為那個時候.

$^{1601}$沒有IDD.
沒有WhatsApp.
要寄信.
要好幾個月才回到英國.
他的教會.
告訴他這個危急的事.
但後來他收到信.
原來當晚.
在英國.
有一個弟兄姊妹的家庭.
突然間裡面.
有很大的.
催促.
要他晚上立即起床.
為David Stone.
李雲斯敦在他的自傳這樣說.
來去祈禱.
但不知發生什麼事.
他心裡突然間.
立即起床去祈禱.
原來計算時間.
就是當時獅子和他搏鬥的時候.
李雲斯敦這樣說.
所以原來可以月陽.
根本沒有WhatsApp.
根本沒有IDD的年代.
神已經可以.
叫一個人祈禱.
不是因為收到WhatsApp.
收到WhatsApp.
可能只是給他祈禱.
但真心的祈禱.
神感動.
所以第十三節.
相當寶貴的弟兄姊妹.
要人明白.
神的義就是基督.
要人知道基督在這個世界.
是愛罪人.
要將人的生命改變.

$^{1641}$所以我們今天的責任就是.
原來每一個民族都可以.
信主 所以是沒有分別.
因為.
眾人同有一位.
做民主在天上.
所以不要以為.
只有猶太人才可以得救.
或者某些人才可以得救.
原來每個人人人都可以得救.
問題就是.
右邊的人我們還沒有去.
跟他們傳他不得救.
所以我們的言語.
我們的行動.
都要表現我們是屬於基督.
我們是基督的使者.
我們要令到他.
心悅誠服的.
找到原來口中.
原來心裡就可以.
向神祈禱.
向神求這份恩典.
所以當他一求.
我很喜歡第十三節.
大家都一起背熟他.
因為凡求告主名.
就必定得救.
原因就是.
他會厚待一切求告的人.
我們將這個.
呼喚告訴他們.
其實不是我們更加幸福嗎?.
我們將這個呼喚.
告訴那個不懂祈禱.
從來都不祈禱的人.
的內心.
他自己向神呼求.
或者我們帶他.
向神第一次的呼求.

$^{1681}$你們可以想像得到.
那個厚待.
是留意想像的.
我們內心那份滿足.
因為是我們將耶穌告訴他們.
當然我們絕對不應該.
成為以後他的操縱者.
我一而再要彼此勉勵.
做傳道.
做導師.
我們千萬不要操縱.
別人的靈魂.
你要他做甚麼他就做甚麼.
這樣的傀儡絕對不應該.
我們某個階段.
我們要退後.
我們要離開.
他要自己成長.
或者他要跟靈魂的指導.
這是應該的成長.
如果一生都跟著你.
他不懂得轉別的彎.
他就不懂得轉彎.
這個意思就是說.
我們要有這一份.
對福音的熱心.
這個就是差遣了.
這是廣義的差遣.
應該每個人都應該要有的.
應該要在那個救恩上效力的.
你連習廣台講的福音信息.
你都可以是背後祈禱.
或者你派傳單.
或者你有蘋果子.
你的朋友.
你的邀請.
諸如此類我們都要這樣做.
這個就是那個差遣.
所以你不要說沒有逢差遣.
是應該我們每一個都有不同程度的差遣.

$^{1721}$我們就要來將那個道告訴他.
那個傳可以是很闊的.
可以是真的帶他來教會.
然後就介紹一些領子給他認識.
為他祈禱.
坐在他旁邊陪伴他.
很有善的.
或者完了聚會還是新朋友.
還是不太熟悉教會的.
你和他一起吃東西.
讓他認識一下其他的教友等等.
這些都是差遣.
這些都是傳道的.
可能很多的步驟的其中一個步驟.
所以都是要這樣走上去.
所以這裡告訴我們.
如果有逢差遣他就會去傳道.
記住是傳道.
最終他得救都是因為耶穌基督.
聽明白認識耶穌基督.
當他去傳道.
有人會聽明白.
聽明白他就會信.
信他就會祈禱.
五部曲.
保羅覺得很清楚.
所以在他講這個過程.
這個賢兒.
即是說他都保妃.
救恩是普世的.
但是賢兒如果不傳.
即是說他就和羅馬教會說.
你們這麼多年來.
你們和羅馬教會說.
你們這麼有資源.
你們這麼有眼光.
你們是帝國的教會.
你們這麼有福氣.
你們是外邦人信主.
信得這麼好.

$^{1761}$如果你們只是為了那些無無畏畏的事.
十四章就講那些無無畏畏的事.
就是猶太人說星期六安息日.
我們外邦人就禮拜日安息日.
就爭論這些事.
猶太人說這些是不可以吃.
外邦人就說這些可以吃.
保羅說.
各人要義字堅定.
要堅.
你的義字要堅.
守日的人為主守.
不食的人為主不食.
都是榮耀神.
保羅說那些事.
好像要不要爭論.
在教會裡面.
最重要的就是.
合一和福音.
所以現在見到就是.
在這裡有一個.
羅馬教會.
如果猶太人被信主.
羅馬教會就開始要.
祈禱.
蒙猜險.
然後就有意識地去傳道.
這個就涉及到你邀請.
你講什麼.
福音的重點.
你跟一個小朋友講福音.
跟一個大學教授講福音.
都是不相同的.
福音是一樣.
但是傳的方法.
和內容.
都是不相同的.
你怎樣引進他.
願意謙卑.
認識福音.

$^{1801}$都是不相同.
一個婦女.
一個家庭主婦.
一個高級的管理人員.
那個福音的進路都不同.
一個福音.
但怎樣講那個福音.
都是不相同的.
所以羅馬教會就是要奉猜險.
去祈禱.
明白了福音的寶貴.
人去求告就可以得救.
得到神的厚待.
他就會去祈禱.
然後就會有感動被猜險.
他就會裝備.
就會去傳.
傳就會有人聽.
聽就會信.
信就會祈禱.
最終你帶他去祈禱.
記住帶他去基督穩固的盤石.
是他的生命的救恩.
的來源.
開始是帶到你面前.
是沒有問題的.
你要餵奶.
嬰兒出生.
你都可以說你躺在這裡.
我現在走了我生完你我沒事做.
都不行的.
人類就比較長.
沒有十二歲你都不可以放下他.
開始就長大.
中學了.
不聽你講了不跟你了.
對的應該就是這樣.
初信主的.
有些很快有些很慢.
記住傳福音的人.

$^{1841}$我們只不過是一個代理人.
將他帶到耶穌那裡.
我們絕對不要操縱.
任何一個人的靈魂.
否則.
罪會比異端更加大.
這裡告訴我們.
就是那五部曲.
所以.
最終保羅想說.
我很喜歡的經文.
十五節.
無逢差遣不能傳道.
如經上所記.
以創阿書五十二章.
報福音傳喜信彼人.
貪門的腳中.
何等皆靈.
五十二章我們來揭一揭.
以創阿書.
就是五十三章的序幕.
是講這個.
講的詩詩如寒.
但是可以是其他人.
五十二章.
第五十二章第一節.
因為神的好訊息.
已經傳開了.
就需要有傳的人.
石安興起.
披上你的能力.
聖成耶路撒冷穿上你皇帝的衣服.
因為從今以後.
未受國禮不見證的.
必不再進入你中間.
即是他們要準備被老歸回.
耶路撒冷要抖下尋土.
起來坐在位上.
石安被老的居民.
要解開你頸鎖的鎖鏈.

$^{1881}$即是被老已經完結了.
神的審判.
懲罰已經完結了.
現在是無福了.
第三節.
耶和華如此說.
你們是無價被賣.
也必無銀被贖.
即是他們現在可以.
來到翻身.
站起來.
主耶和華如此說第四節.
起先我的百姓下到埃及.
在那裡寄居.
又有亞述人.
無顧欺壓他們.
耶和華說.
我的百姓既是無價被老去.
如今我在這裡做什麼呢?.
耶和華說.
克制他們的人呼叫.
我的名整天受褻瀆.
所以我的百姓.
必知道我的名.
到那一天他們必知道.
說者說的就是我.
看了是我.
即是神現在親自去呼喚他們.
轉去他們的苦難.
因為神的懲罰時間已經滿了.
所以以前.
他們對神有誤解.
也有那些欺壓的人.
對神的褻瀆.
所以他們聽不到神的呼聲.
那個呼喚.
那個拯救.
但現在要大力揚聲.
留意第七節.
那裡又是直接引述.

$^{1921}$那報佳音傳平安.
報好信傳救恩的.
對石安說.
你的神作王了.
這人的腳登山.
因為上石安山.
因為大家去的聖地.
就會有這個概念.
在西羅萬池.
原來是最低位.
所以那些祭司在西羅萬池.
在鑄坪節.
擔水.
一直上聖殿.
是很斜的大斜路.
所以我們去過西羅萬池.
我們就知道.
原來是很高.
很急速的.
祭司就要浩浩蕩蕩.
擔著兩擔水.
就上去堵聖殿的祭壇.
所以那個就是.
鑄坪節最興高采烈的日子.
所以登山.
猶太人就會明白.
就是上聖殿.
這人的腳登山何等皆美.
保羅就引了這一節聖經.
52章第七節.
報告神月立猶太人.
現在保羅說.
報告耶穌基督的恩訊稱義.
人明白了耶穌.
就會求告神.
他就會得到永恆的救恩.
所以那個福氣更加大.
所以記住.
如果我們已經在做這樣的工作.
記住我們是非常蒙福的.

$^{1961}$我們的腳zoung 是皆美的.
神體的皆美.
是這種心靈的滿足.
和將來永恆的賞賜.
其實更加是你帶到那個人的生命改變.
他的下一代的生命改變.
其實這些就是價值連城的生命.
比錢更加重要.
將人的生命帶到神的面前.
所以一個教會的老穆者.
就是那個教會的寶.
因為看著第二代第三代這樣出世.
一路這樣牧養.
一路這樣看著他們成長.
這些就是福氣皆美的腳zoung .
所以十三第十五節.
大家都背熟它.
真的看到原來全福音不可以這麼輕散.
原來可以這麼蒙福.
記住演唱阿書說的.
那個神拯救的日子到了.
懲罰的日子已經過去了.
就可以得到神的月.
現在保羅告訴我們.
他直接引述沒有傳到.
所以如果肯去傳的人.
那些人就會蒙福了.
因為他十四節說延延.
然後他呼籲.
如果報福音的人.
腳zoung 就皆美.
當然十六節我們下一次再說.
因為時間不足夠了.
下一次再說這一段.
因為仍然猶太人都是以他的良心.
來不肯接受.
那這一節就先說到這裡.
我們一起同心圖告.
今天我們能夠解第十章人的責任.
和神的解選不單沒有抵觸.

$^{2001}$而是告訴我們神的選擇.
就是差顯你的愛子.
無條件為我們人類釘十字架.
而且已經是具體的證據.
主耶穌由死復活.
從高天來到世間.
從陰間升到天上.
回到主你原本的地方.
我們很盼望我們能夠明白.
絕對不深奧的救贖的恩典的福音.
讓更多人能夠明白.
可能有些人覺得是愚蠢.
但是在信的人就覺得是神的大能.
我們盼望主作供.
讓我們今天晚上.
無論之後聽到查經的弟兄姊妹.
在不同的地方.
我們都能夠成為傳福音的使者.
我們能夠將人帶到禱告的根本裡面.
他去禱告,求主的接納和拯救.
以致他們得到神厚待的救恩.
幫助我們能夠成為一個傳福音的教會.
和一班傳福音的人.
以致我們的路徑.
是兌換了基督的恩典的自由.
我們的路徑是無雄福的.
是皆美的.
是在神眼中來去蒙主使用的.
所以我們盼望你在我們錦書家.
嘉欣堂,洪欣家.
我們一齊來去同心的服侍.
現在我們有新田部屋.
有簡約公屋,洪福村.
我們都是集中現在去做.
還有青年會小學.
區內洪水橋都有兩間中學.
我們現在都有接觸.
盼望主都使用我們.
叫我們在教會多傳主的愛和穩固的根基.
少說的是是非,無謂的行為.

$^{2041}$願主幫助我們這樣感恩禱告.
奉耶穌基督的名.
阿們.
多謝各位.
記得下星期再會.
再見.
\newpage



\section{羅馬書}
\label{sec:1GDpbDYTYf0}
\textbf{2025-09-12 周五晚查經班《羅馬書(22):福音:人的選擇與後果》 講員:曾裔貴牧師}
\newline
\newline
連結: \href{https://youtube.com/watch?v=1GDpbDYTYf0}{\texttt{https://youtube.com/watch?v=1GDpbDYTYf0}} ~~~~ 語音日期: 2025-09-13
\newline
\newline
\hyperref[sec:3TbpwI2Yugs]{\small{< < < PREV SERMON < < <}}
~
\hyperref[sec:index_chronic]{\small{[返順時目]}}
~
\hyperref[sec:index_scriptual]{\small{[返順卷目]}}
~
\hyperref[sec:ZdF__6Sl2No]{\small{> > > NEXT SERMON > > >}}
\newline
\newline
$^{1}$各位姐妹我們一起同心禱告.
我們再一次感恩天父.
我們在晚上可以一齊來到賢讀神的華園.
我們很盼望藉著保羅他很用心.
以及語重心長.
這樣來探討什麼是福音.
怎樣可以得到福音.
以及如果得不到福音會有什麼後果.
都很透徹地告訴我們.
多謝你將這個這麼寶貴極大的英雄.
藉著聖經的真理啟示我們每一個時代的人.
多謝你也讓我們由西教士來到中國.
來到很多年的年日.
來翻譯了從希伯來文.
希臘文翻譯了和合本的中文聖經.
已經差不多有200年的時間了.
多謝你我們能夠一代一代的華人.
透過中文聖經去得到教導.
而且現在在21世紀都有不少的神的僕人.
已經開始能夠掌握很多的願望.
以及在神學院裡面都有很多的教導.
很盼望我們能夠真的看到過去很多的困難.
無論外在的環境.
敵視,攻擊福音的.
或者內在普遍我們的教育水平很低.
到現在我們都已經大部份人.
可能由中學,高中甚至大學的學位.
能夠掌握不同的語言.
但是我們很盼望真正愛慕神的話語.
這個不單止是有學術水平修養這麼簡單.
而是需要有靈命.
需要有一個屬靈的胃口.
需要有一個性靈更新的新生命.
才會對神的話語有這一份的學悟的.
所以我們盼望神幫助.
也都讓我們不要看我們自己不行.
盡我們的能力在神恩典裡面去學習.
而總是有我們自己那份的祝福.
因為我們打開神的話語去研讀.
願你使用我們在一直研讀每一節的聖經的時候.

$^{41}$都能夠有很多啟發,個人的反思.
這樣來一起學習.
我們這樣仰望,感恩,禱告.
奉耶穌基督的名 阿們.
好 我們上一個星期是講到第十五節.
十章的第十五節.
跟大家再來溫習一下.
有幾個重要的題目.
就是一直都在講.
因為第九章開始.
要探討的就是以色列人現在是不得救了.
因為他們是非常拒絕.
到底什麼是救人的途徑.
是信心的信靠.
還是仍然靠那個律法的意義.
律法的遵守呢?.
這個第九章就已經相當詳細的探討了.
稍微的分段就可以給大家看看.
這個第九章我們已經上下落地講過了.
保羅為他自己的同胞憂傷.
十章,十一章都還有這樣講的.
每一個大段落的轉向.
他都提到他自己對同胞的憂心.
但是第十一章是正面的.
所以衛邦人你不要這麼得意.
當神的恩典來到的時候.
以色列全家會得救的.
所以都不是絕望的.
第六到第十三節.
就是第九章我們上星期都講過了.
九章和十章好像剛剛一個很強烈的天秤的平衡.
神的揀選在乎人的行為.
但是十章又很在乎人的選擇.
人對福音的態度.
十章我們從十六節開始看.
你就會看得很清楚.
保羅引了好幾段的舊約.
來到講到原來不是聽不到.
是故意的不聽.
這是兩件事.

$^{81}$可以很清楚十章是很強調個人的責任的.
不信就不是神沒有揀選.
是他拒絕福音的選召,呼籲.
所以兩者的平衡.
就不會給我們一種錯覺.
他選擇我們就不用做自己的責任.
一定會相信的.
這些又是一個推搪的邏輯.
所以我們要記得.
但是在原來很多的事情上.
的而且確神的主權.
他引用摩西和法老十宰的比試.
法老是一而再.
那些災禍一放鬆.
他就將那個心來硬.
把心一橫.
應該把心一硬.
法老就是這樣.
因為恃著自己有權力.
恃著自己還有那個能力.
更加可能就是訴諸於他們拜的.
他們自己的埃及的假神等等.
所以你可以想像得到.
這個是神的主權.
來將法老成為一個示範單位.
示範一個人輾轉拒絕.
可能有一點聲音的雜音.
大家不要介意.
如果有些人還沒有關咪.
就要關一下咪.
好 多謝.
我們繼續這個講道.
所以以色列人最終.
保羅給我們的結論.
他們沒有去運用神給他們的信心.
到第十章.
我們又來看看.
這個第十章很重要的發言.
保羅他想強調的.
仍然是信心而不是律法.

$^{121}$所以是主耶穌被神差遣來到這個世界.
去到第十三到十五節.
真的告訴我們.
如果沒有傳道的人.
人聽不到福音.
聽不到福音.
人就不能夠信.
或者不能夠透過福音.
燃點他內心的渴望.
那個需要.
他就看到自己.
原來在神的面前.
是有那個罪.
他需要得到神的拯救.
這個是一個很重要的.
聖明光照.
但是都是需要有傳道者.
有意識地帶著神的道.
神的愛.
來呼籲人去悔改.
所以如果我們真的有這一份報福音.
傳喜訊的心靈.
心態.
這樣的人.
腳中是何等佳理.
自己都會忘福.
提安太后書第二章.
保羅就說.
叫提安太要留意.
我們是基督耶穌的徵兵.
凡在軍中當兵.
就不將誓無前身.
就是很專注在自己的神的召命.
然後.
人在牆上比武.
非按規矩.
就不能夠得冠冕.
那個牆上可能是運動場.
可能是當時的馬拉松.
或者其它的競技.

$^{161}$有很多不同的競技.
但是要按規矩.
就是說要按那個.
設計體育運動的規則.
來去進行.
就是踢球你不能用手.
一用手就犯規了.
這個就是按規矩.
就不能夠得冠冕.
提安太可能在他的時代.
有很多壞了心術的傳道者.
或者有其它目的的傳道人.
是為了要得利來去傳.
或者甚至乎在《肥沃表書》有提到.
有些人是因為保羅坐牢.
就故意去激他.
不是想傳福音.
是想激保羅.
現在我們去傳了.
我們要做給你看了.
現在你坐牢了.
保羅就說他不介意他的安康.
或者是真心或者是假意又如何呢?.
基督被傳開.
他就歡喜了.
就是說就算你們不認同保羅.
有什麼所謂呢?.
你們去傳.
有人信.
那就是神的恩賢.
我得不得到那個高舉.
被人讚賞又何妨呢?.
這個又是一個很大的胸襟.
接著就是第三個.
保羅有第三個血面.
應該是儀章的.
我忘記了第四,五,六節.
勞力的祿夫.
你當先得糧食.
這個也很重要.

$^{201}$一個傳道者.
或者一個真心想服侍神的人.
一定要專注的勞力.
勞力就是說你自己背後為一個人的禱告.
你是想方設法的.
希望他能夠慢慢進入福音的胃口.
很多可能的御公.
你跟他不是一句信耶穌.
不信就先這樣.
這樣就很難的.
不是說不行.
但是一般都不行.
可能是建立了一個友誼,信任,關係.
這些就好像要存錢一樣.
每次親一些,每次親一些.
我們做一個勞力的祿夫.
就是說如果我們在世界上.
我們尚且都專注我們自己喜歡的人事物.
有些人喜歡音樂.
有些人喜歡看運動.
有些人喜歡其它的投資或者其它的東西.
有些人喜歡車的研究.
或者研究相機.
這些全部很忠誠的.
一些嗜好,一些興趣.
這些沒有壞事的.
但是如果我們能夠專注在那個程度上.
勞力的祿夫.
用自己的勞力.
可能是汗水.
可能是眼淚.
來去囂觀一個人的生命.
這樣的勞力.
你當他自己先得糧食.
或者他自己會得到神的豐厚恩典.
其實那個過程已經令到他自己.
不斷從一個感恩的心.
不斷從一個自己.
都會將自己放在神的恩典裡面.
這樣來做一個服侍的人.

$^{241}$得到的功勞,得到的感恩.
其實榮耀是歸給神的.
因為每個人都是這麼軟弱的.
所以更加重要的就是.
福音不單止是猶太人的尊美.
而是他們反而拒絕十章後面一會兒會說的.
反而福音是一個普世性.
你看看很巧妙的.
保羅怎樣引幾段舊約的聖經.
其實表面真的很不吉利.
引詩篇的十九篇第四節.
引五十三篇的第一節.
又引六十五篇的二賽亞書等等.
其實表面跟那個福音沒有什麼關係.
但是保羅就能夠游刃有餘地引述那些經文.
而且那些經文本身又真的不是離經叛道的.
跟上下文沒有關係的.
你可以想像保羅.
他那種東西能夠將舊約的精華.
放在他自己的福音的宣講.
來講述到底猶太人現在在救恩那裡.
他是什麼定位呢?.
怎樣將他帶回去救恩的軌跡呢?.
這個九至十一章.
就真的給我們有很大的眼光.
視野.
更加重要就是神透過保羅給我們的啟示.
所以很重要的.
這十章的九節.
保羅告訴我們.
福音其實就在我們的口裡.
在我們的心裡.
上個星期都跟大家讀過《申命記》的三十章.
摩西林中跟這些以色列人講.
你不要以為要認識神.
要上天下海.
要去很遠的地方尋道,去問道.
好像那些人要踏破鐵鞋那樣.
要找冰山雪蓮這樣的心情.
其實是不需要的.

$^{281}$這個道就在我們的口裡面.
在我們的心裡.
當時是在講摩西的律法書十誡.
摩西死了.
留下的就是神頒布的十誡.
和一路衍生的神的各方面的話語.
但是保羅講的這個宣認.
就是一路在講神差遣他的愛子.
彌賽亞耶穌來成為人類的救贖.
所以我們一宣認耶穌.
就立即可以得救.
所以這裡不僅是一隻重要的金句.
而是告訴我們整個的時代.
是一個里程碑的開始.
就是主耶穌降生在馬槽.
長大,傳道.
被人殺害.
第三天復活.
在地上四十天.
升天.
現在坐在神寶座的右邊.
將來要再回來.
一個整全的福音.
就是我們要宣認.
一宣認.
一令一個人聽明了這個套餐.
是耶穌一生所做的.
和他本身的榮耀的身份屬性.
那個人就立即得救.
所以是一件很寶貴.
很美麗的事.
我們在人生能夠盡量搶救一個得一個.
所以我們一定要加入傳福音的行列.
去做一個福音的傳道者.
接著我們就下去.
真正的問題就是.
猶太人缺乏的是真的知識.
所謂真的知識.
其實在道育一種就有的.
不過他們一直都用十誡.

$^{321}$用那些生活的條例.
或者一些猶太人的習俗.
一些拉比的傳說.
做了很多的移民.
很多的傳統.
這樣來到去慢慢以為.
這樣就是他們的一個護身符.
一個保護罩.
拒絕一個好像很嶄新的.
突然間的新的道理.
就是有一個人說自己是神.
你就要相信他.
起初不相信是很合理.
但是慢慢慢慢發現.
原來這位耶穌.
必須是要你用信心去.
將他接納在心裡.
認信.
所以這個就是真的知識.
猶太人有他們自己的防衛機制.
就是他們一直都是有律法的人.
這個律法還要是來自神的摩西的十誡.
還有在西瀨山上已經立了一個榮約.
他們是神的子民.
所以不是因為這樣.
令到他們叛變.
而是仍然是好像保羅引述的.
《以賽亞書》的28章.
仍然是因為這個被驗證過的石頭.
寶貴的防國石.
耶穌基督使得他們叛變.
所以他們不是叛變在律法.
他們的知識.
是叛變在耶穌這一位真神的兒子.
彌賽亞的身上.
所以我們信主不斷需要更新我們自己的神學.
我們自己都要不斷與時俱進.
那個真理的更新.
真理是不變的.
但是我們個人的成見.

$^{361}$就有越來越牢固的偏見.
這個不知不覺會令到我們.
與真理有時真的背道而馳.
我們自己都不以為意的.
所以每一代的人都要反省.
每一代的基督徒都需要去追尋那個真理.
真理從來都沒有變過.
只不過是我們接近真理的那個路.
有時就越來越遠.
我們自己的喜好.
我們自己的執著.
我們自己的成見.
有些人一回到教會.
就好像很多的恩怨情仇就開始了.
完全連最基本回來教會.
就是真心的開放自己.
與神與人的相交,相愛.
不應該有任何的嚴規.
有嚴規,有誤會.
就應該尋求和睦一心的追趕.
這樣的返教會的目的.
每個返教會的人都知道的.
但是慢慢就原來我們自己那一套的知識.
就原來與真理可以是這麼遙遠的.
這樣真的很可惜的.
你可以想像得到.
有時神在教會的情況.
其實是一樣的.
猶太人一代一代都是這樣守著那些律法.
都是這樣拒絕福音.
甚至將耶穌彌賽馬置於死地.
他們都要這樣做.
還以為自己是守著神的戒律.
可以想像得到.
原來這個真知識真的多重要.
我們需要不斷的要來與時俱進.
讓神的道在我們大家的心靈裡面.
不斷的更新我們.
所以我們很重要的責任就是要人人去到基督那裡.
這些全部是上個星期講的.

$^{401}$我們可以來做一個很簡單的總結.
今天我們會從第16節開始看起.
猶太人沒有真的知識.
是沒有得救的真的知識.
其實就是我們衛邦人很清楚的.
告訴別人耶穌就是你的救主.
就是你的生命的主.
你相信祂 你就可以得到新生命.
與神和好.
本來我們衛邦人根本不知道哪一位神.
我們拜都拜錯了.
沒有人會拜一位他不懂得又敬畏的神.
和不斷尋找祂 是沒有的.
但是透過聖經我們就認識了.
我們就知道了.
原來做民主一直呼喚我們 尋找我們.
所以我們有了這個真知識.
我們就可以來得救了.
這個就是猶太人沒有的真知識.
我們衛邦人很早就已經有了.
所以這些告訴我們.
所以保羅都告訴猶太人.
其實律法的總規.
是要大人看到自己的不足.
自己的無助.
自己在神面前的審判.
公義的審判的絕望.
是告訴我們.
其實是要讓我們看到不是守律法.
不是以為守律法就可以解決.
而是帶我們去看到那個律法的總結.
其實就是基督.
帶我們看到其實基督就是那一位.
帶給我們可以完善.
我們在神面前的新生命的那一位.
所以保羅很強調.
律法的總規 律法的總結就是基督.
在《利害經文》.
這個羅馬書的14章都提到的.
律法的總規都是愛.

$^{441}$所以你會看到.
兩者是從不同的角度去講那件事的.
所以律法的總結.
如果是我們可以超越.
我們可以自大 可以驕傲.
我們就連信耶穌都信錯了.
應該讓我們看到.
這一位真是無與倫比的基督耶穌.
是我們一生的發現的至寶.
是我們越來越在祂面前.
是要謙卑 是要俯伏.
是要傾盡我們的生命.
有限的人生去愛祂.
去跟隨祂.
去以祂為我們的主.
這樣的基督徒就是看清楚.
律法的總結是帶我們去到耶穌那裡.
雖然我們沒有律法的頒布.
守十戒.
但是不代表我們沒有這個.
來對我們自身在地上.
我們去成為基督徒的限制.
所以不是繼續做自己的喜歡.
以為自己做好了就行了.
相信在座或者在線上的弟兄姊妹.
都不會認為我每個星期準時上崇拜.
只是上崇拜 我就很好了 在神面前.
相信那個是一個給我們的詛咒.
其實是要告訴我們知道.
這個整個星期的其中這一點.
是要擴散到在我們的生活.
盡量每一個層面.
我們都能夠得到神對我們的生命的認可.
和喜悅 這個才是最重要.
所以你可以想像 律法的總結.
不只是猶太人才需要.
其實我們基督徒也需要.
所以如果這些事情帶你去不了基督那裡.
無論是律法 無論是崇拜.
無論是那個教堂 那個教會.

$^{481}$其實都不是真正發揮到真正的屋陽的作用.
所以我們真的要將人帶到基督那裡.
這個是不斷的老生常談.
但是往往有時候我們就以這些禮儀,傳統.
做慣了 行禮如儀.
就以為我們已經親近了神.
那個一定是對的途徑.
但是那個不是終點.
那個是一個途徑帶你去親近主.
以基督為你的生命的主.
如果達不到那個目的.
那個途徑還是很遙遠的.
就是你長跑.
你不可以說我衝一段就等於我已經完成了.
你根本都沒有去到終點.
所以你就一定要明白.
所以很小的福音.
就是第三段我們就會集中看最後這六節聖經.
我們要先讀一讀聖經.
我們從第十章九節開始.
這裡說:你若考慮認耶穌為主.
心裡信神.
叫他從死裡復活.
就必得救.
這裡就是回應上文.
到底這個「道」在哪裡?.
也就是等於對某些猶太人.
在這裡用一些藉口.
你告訴我救恩 我就會相信.
是不是這個耶穌?還是怎麼樣?.
保羅就是要這樣來跟他們說.
其實我們的摩西都已經告訴我們.
在聖經裡面都已經可以尋找到真道.
尋找到神的生命.
更何況現在神自己的愛子.
來親自去宣告這個救贖.
所以一宣認耶穌就立即得救.
所以我們幫人去宣認耶穌是他的生命的主.
是他的救恩的主.
這個真是一個很大的.

$^{521}$我們變成一個助診使.
屬靈的助診使.
帶人來到這個世界.
所以是一個屬靈的世界的生命.
所以將來那種無限的那種的發現.
有很多人是我們曾經跟他說.
有意無意或者我們的生命的影響.
跟他說他信耶穌.
或者是他小孩的時候.
因為這一句話.
他後來這顆種子萌芽長大.
這些是將來在天堂.
很多這些想不到的那種的恩典和賞賜.
就會發生.
所以千萬不要吝嗇我們能夠在地上有限的年日.
我們成為一個福音的使者.
因為人口裡承認.
心裡相信就必定的得救.
因為他已經跟耶穌生命的主有一個的結連.
所以很重要的.
我們繼續讀下去.
第十節.
因為人心裡相信.
就可以輕易.
口裡承認.
就可以得救.
經常說.
凡信他的人.
必不至於羞愧.
這個就是《耳塞孟子》講的盆石.
28章的16節那裡.
我在石安.
安放一塊石頭.
是試驗過的石頭.
是穩固的根基.
信靠的人.
就必不至於羞愧.
28章的16節.
所以保羅告訴我們.
基督耶穌就是現在人類唯一的盆石.

$^{561}$所以在主耶穌裡面.
我們心靈應該是安穩的.
安穩在耶穌的救贖.
恩典裡面.
在他整個的榮耀.
神聖的屬性裡面.
就是我們人生最大的依靠.
是的.
第十一節.
第十二節.
猶太人和希利利人並沒有分別.
因為眾人同有一位主.
他也厚待一切求告他的人.
再一次.
保羅已經將他的福音.
完全沒有局限性了.
不會偏幫猶太人.
也不會因為怪責猶太人.
就好像只是偏待外邦人.
神要他做外邦人的使徒.
但是他對猶太人的愛.
從來都沒有停過.
因為就是他自己的寫照.
就是他自己得著的救恩.
神給他一個新生命.
所以他就更加對自己的同胞.
有這一份的著急.
有這一份的迫切的愛.
所以也要問一下我們自己.
我們有沒有普世性呢?.
我們會不會偶然間的時候激氣呢?.
對別人敏震呢?.
這個偶然間你悔改.
但是如果你是結構性.
那真是不恭喜了.
如果我們的心胸狹窄到一個地步.
我只是對我自己喜歡的人.
這個就絕對不是救恩.
就是說你跟救恩的邏輯相對太遙遠了.
保羅自己這樣說.

$^{601}$這裡不停在《來話書》出現好幾次了.
猶太人和外邦人都沒有分別.
原來保羅立刻上到神的神光.
神的體人類.
在這麼高空的體人類.
哇 不是啊 這個是美國首富啊.
這個不是啊 是非洲的奴隸啊.
沒錯 在人類的社會.
是有這麼懸殊的分別.
你想一下 只是香港你住在壽臣山.
你住在壽茂坪 已經是兩回事了.
那個生活 那個一代人.
他的兒女 那個將來的成長.
都受著環境影響.
都很大的改變.
你試一下現在天水圍這麼多的太空遊.
那些什麼伊索偉達.
那個學名忘記了.
你可以想像得到.
在一些比較有學問,有知識.
又有一定的金錢富裕的家庭.
肯定不會這麼容易就接觸這些壞的事.
壞的事,壞的人.
那他的人生都會有改變.
但是我們記得 我們信徒.
要有保羅這種普世的眼光.
並沒有分別.
福音是向萬民傳的.
我們不可能對每一個民族.
我們都有特別的感情.
我們又會各種的語言 除非是神.
一般都不可能的.
但是不代表我們的心是這麼狹窄的.
是這麼刻意的.
是這麼偏心和這麼小器的.
這個是我們個人的問題.
舊恩從來沒有這個東西的.
是沒有任何的門檻的.
哪個民族特別優越一點是沒有這個東西的.
是不是也有頂梓滅未山的那個靜音.

$^{641}$大家不要介意 一定要靜音了它.
否則影響其他人就不是那麼好.
你可以想像得到.
猶太人和希臘人並沒有分別.
因為眾人同有一位主.
原因就是這個 神愛世人.
不偏待任何一個人.
所以這個就是真是舊恩的寶貴之處.
真是神給每一個民族.
某程度上都是公平的.
所以現在如果我們參與宣教.
跳出我們自己同胞.
廣東人,香港人,中國人.
其實都好的 都給我們有眼光.
有新的視野 有些能界.
或者那個心靈的胸襟大一點.
也都要為宣教士祈禱.
不是每個人都可以做到這樣的工作.
譬如我們認識的穎兒.
一家人去向索馬里群體.
譬如陳姑娘向尼泊爾.
或者現在阿克里斯向阿拉伯世界.
這些穆斯林的世界來傳福音.
埃及,土耳其 甚至乎亞洲大疆.
亞美尼亞這些 伊朗這樣.
是很困難的 但是更加需要為他們祈禱.
給他們有一份普世的心靈.
所以中人,全人類同有一位主.
我們大家都是向著神的.
所以如果在我們的心靈裡面.
捆綁著一個人 每次一祈禱.
我們就想起他的壞事.
想起對他的仇恨.
對不起了 你真的祈禱不了.
因為在神的眼中 他愛世人的.
多大罪的人 好像保羅.
神都給他那個救恩 使他改變的.
所以我們更加應該為.
我們自己不喜歡的人要祈禱.
這個告訴我們 他會厚待一切.

$^{681}$求告他的人 這裡立即說到.
怎樣可以求告呢? 就是福音.
就是聽明福音 他就會求告.
所以我們作為基督徒.
我們現在得到救恩.
我們就應該更多祈禱.
現在高文寶都不斷增加人.
真的耐不足見.
最近又有小偉的祈禱.
又立即祈禱 原來挺好玩的.
自己在神面前 你自己要認真.
原來記下那些人命的祈禱.
其實是讓我們自己更加感受到.
原來我們是一個福音的使者.
是一個祈禱的勇士 一個守望者.
神一定會祝福的.
因為求告主名就必得到.
為人代禱的人.
肯定是神來興起和保守.
所以盼望 香港現在普遍.
不只是香港 我想世界.
祈禱有時候可能只局限於在星期日的崇拜的公禱.
其實不健康 應該我們有一個中國.
看到不同的事情 我們就有一個祈告的心去祈告.
當然每個人的時間 他的程度都有不同.
但是很重要的就是求主給我們祈禱的生命長大.
這個禱告就更加重要了.
我們裡面那個祈禱的生命是不斷有成長.
就是說你的胸襟 對事物的研判越來越成熟.
13節 求告主名必得到.
約珥書的2章32節.
保羅引述的舊約聖經.
然後第14節 然後就懸疑了.
我們上個星期都有提到.
人未曾信 怎麼能夠求呢?.
就好像再一次告訴我們.
原來那個責任不只是神的揀選.
人有責任 所以神既然將這位主耶穌.
保衛唯一的救世主 透過聖經的啟示.
讓我們認識了 我們就要一個時代的教會.

$^{721}$牧者信徒就要承擔起全陽的責任了.
所以這是一個很好的美麗的循環.
我們自己因為聽明白別人跟我們講的福音.
我們信主 我們一直去裝備去長大.
然後我們明白聖經的真理.
越來越清楚告訴我們.
原來每一個教會 每一個時代的教會.
都有它全道的責任的.
所以信道是從聽道而來的.
就是說那個教會自己本身.
要明神的道的規劃,規模.
是怎麼樣能夠信主,得救.
然後它就要化整為靈的教導弟兄姊妹.
一起去做福音工作.
這樣就是信道從聽道而來.
而這個聽道就是我們這一幫已經有神的生命.
有神的道的弟兄姊妹.
所以每個星期的聚集.
是告訴我們我們在神面前做著神喜悅的事.
告訴我們 還沒聽見怎麼能夠信.
所以你看到這個循環.
沒有全道的怎麼能夠.
沒有奉差遣 又怎麼能夠去傳呢?.
所以這個雷體的就好像我們是必須有差遣.
但是這個不是特定講的一些全道者宣教士.
是在講每一個基督徒我們明了救恩.
知道人的生命的去處.
就是永生或者是滅亡的永死.
我們就有這樣的意象.
有這個責任我們就要被差遣.
求主差遣 這樣去做.
當然特定的是有特別的差遣.
但是一般教會如果我們回來.
只是為了我們開心 我們來去敬拜.
我們沒想過要向陌生人傳福音.
很麻煩 很多手續 很浪費我們的資源.
每個星期唱聖詩開心完.
大家回去自己的工作 自己的家庭.
就可以了 不用那麼麻煩.
那些不信的人有這麼多的問題.

$^{761}$這個就不是純正的教會存在在這個地上的必須.
不信的人唱歌都很開心的.
像卡拉OK或者一些夜總會的唱歌.
正經的那些都開心的.
一班人一起去唱一些原曲都開心的.
問題就是聖經告訴我們.
我們不應該只是內部的開心.
我們是要讓人聽到神的道.
就是基督祂來去求告耶穌.
所以這裡告訴我們.
就是14節打後到15節.
若會有逢差年怎能傳道.
如經上所記報福音傳喜訊的人.
他們的腳踮何等皆美.
很盼望我們因為看到福音的寶貴.
我們去主動求神差遣.
就是我們將自己放在福音的宣講裡面.
我們就已經被差遣了.
所以不是特別的人才可以被差遣.
我們這樣來擺上自己.
我們就是能夠成為皆美腳中的人.
因為我們走的是將寶貴讓人得到新生命.
可以生命改變的福音來傳遞給人.
那個傳遞的人本身就已經是皆美了.
保羅說這樣的人的腳踮是何等皆美.
在神眼中 不只是在人的眼中.
在人的眼中如果那個時代是不容許你的.
你就要死了 你的腳踮就要綁著那個腳踮坐牢.
保羅很快就要在羅馬坐牢.
所以腳踮的皆美不是每一個時代都是這樣那麼順暢.
好像坐滿了看得到的五萬人一樣.
五萬人舉手 當然不可以說沒有 神要做.
但是有時更多的就是不斷被人嘲笑,被人拒絕.
有時會應難生 有時會很灰心.
有時甚至有些人單單打打.
但是記住傳福音在神眼中是何等的皆美.
腳踮是何等的皆美.
我們要記住要在神的眼中看我們自己的福事.
所以已經等於杜絕了那些沽名釣譽了.
一次,兩次都沒有馬上回報的 不做了.

$^{801}$那個動機就已經不純全了.
真正純全的動機就是一次,兩次,三次都不氣餒.
因為沽名釣譽不斷不氣餒的 沒有可能的.
一定是他內心的動機是很純全的.
他就真的看到很希望將福音讓人能夠得到.
我這個人被人誤解有什麼所謂的.
他得到福音就行了 不要緊的.
我有沒有讚賞 不要緊的.
他跟耶穌不就是應該是這樣.
他必興旺 我必衰微.
這個是約翰福音三章的33節.
施洗約翰 他的僕人來問.
夫子 從前跟你在藥膳行.
你幫他施洗的那一位耶穌.
現在的人都跟了他.
約翰就說 如果不是從天上賜.
沒有人能夠得到什麼.
我曾經說過 我不是基督.
你們可以為我作見證的.
娶新婦的就是新郎.
新郎的朋友看到新郎來 開心.
那就很合理的 是應該的.
所以人們去找新郎 那個新郎的朋友被冷落.
是應該被冷落的 是應該被忘記的.
所以這個就是他必興旺 我必衰微.
傳道者最高的境界就是.
完成一切被人忘記.
甚至被人誤會都不要緊的.
你終身做好你自己神要你做的.
問心無愧 神看你的內心.
知道你已經盡了你的力了.
他必興旺 我必衰微.
所以最高境界就是被忘記.
那就是一個很重要的修為.
應該叫修養 我們不是修行.
是一個我們個人的品格.
傳道人的品格就是被忘記是開心的.
因為基督被人高舉 被人尊崇.
所以你可以想像得到.
但是第十六節就開始我們今天的講道.

$^{841}$告訴我們第十六節.
「只是人沒有都聽從福音.
因為以塞瓦說.
主啊 我們所傳的有誰信呢?.
可見信道是從聽道來的」.
再一次引述剛才的循環.
不過這次加了一點東西.
「聽道是從基督的話來的.
但我說人沒有聽見嗎?.
誠然聽見了 自問自答」.
就是說當時的人一定有一個這樣的護身符.
我們沒有聽過.
我們不知道什麼可以使我們得救.
故意的好像掩耳盜鈴一樣.
故意的去按著心.
其實不是沒有聽 是不肯信.
這個就是問題的關鍵所在.
是拒絕這個純傳的福音.
所以保留在這裡 我還讀了一點.
不是太多經文的.
接著他引的聖經就很特別了.
詩篇的19篇.
「朱天術說神的榮耀.
窮倉就傳揚他的手段.
即是這一日到那一日等等等等」.
就是說那個大自然的創造.
是普遍地一直存在.
例如我們可以計算到.
8月15農曆是全年.
我們看到那個月球是最圓最光的.
每年都可以這樣計算到的.
這個就是「朱天術說神的榮耀」.
窮倉傳揚神的手段.
就是詩人告訴我們.
大自然如果我們認真去尋索.
是可以讓我們看到有造物主.
雖然不知道是哪一位.
所以後來就是說神的精煉.
從宇宙窮倉的浩瀚.
那個人盡皆見 人盡皆知.

$^{881}$那個這樣的設計.
你可以量度 你可以去計算.
有很複雜的數學 物理等等.
但是真正要讓人認識那個內心的真神.
就是透過聖經的神的導覽.
就是第十九篇詩篇的大概主題.
所以保羅這個人就是引第十九篇的第四節.
這裡說「他們的聲音傳遍天下」.
這個他們就是在說宇宙的星神.
太陽,月亮,其他的星系.
我們想到的例如獵戶座.
或者九大行星,宇宙的銀河系,太陽系等等.
還有外太空的銀河系以外的.
我們都不認識的那些星系.
他們的聲音傳遍天下.
就是說從實創造開始.
一直到今天都是這樣生生不息地告訴我們.
這個宇宙是有設計者.
是有美麗的創造者來存在的.
接著就是「他們的言語傳遍地極」.
特別強調就是無處不能夠認識這一位的創造主.
就是無處不知 沒有一個地方是隱瞞.
是沒有了對這一位造物主的那種謙卑的認識.
保羅這樣的人就很想告訴我們.
其實以色列人就千萬不要用一個托詞.
就說我們沒有聽見 我們不知道哪裡有.
其實就是拒絕去聽 聽到而拒絕.
我不知道大家有沒有這種感覺.
因為每個星期都有報告.
報告有時候有兩三次.
重要的都有兩三次報告.
報告完之後問一些大紀元的知識分子.
他們聽不到的 因為跟他沒有什麼關係.
而且又聽不到 是一個報告的時間.
好像他就可以放鬆的.
很多人講完報告完.
他是完全聽不到任何的細節的.
當然可能跟他有直接關係的 他就會留意到.
當然我這個是開玩笑的.
絕對不是說要拒絕去聽.

$^{921}$但是這裡提到以色列人.
他們是聽見知道耶穌已經在那個時代.
被人高舉 被人冒生命危險都要傳遞.
都要來講述 耶穌就是我的尼塞俄.
就是我的救主 無論你怎樣對待我的生命.
我都不能夠否認 耶穌是我生命的主.
其實是很清楚的一個信息.
但是猶太人 保羅就說拒絕相信.
拒絕去那個信帶來的生命改變.
所以原來聽從信心跟選擇.
是完全拉上關係的一個很重要的.
生命轉變的不同的步驟.
我下一篇就應該有的了.
我們先讀下面.
第十九節 保羅再說.
我再說 果然是對的 以色列人不知道嗎?.
就是那個問題了 其實他們一直都在問.
保羅你說得這麼偉大 哪裡啊?.
告訴我吧 我們相信吧.
其實真的是一個托詞 真的是一個敷衍.
真的是一個很賴皮的一種心情.
很多人都是這樣的 離家會.
你的耶穌給我飯吃嗎?.
很多的誤解帶來對這位宗主嘲笑侮辱.
這些都在神的憐憫下沒有計較.
使得它慢慢的改變.
所以你可以想像到.
以色列人就很浮現到這種輕佻.
不介意 不在乎 不懼怕的心情,心態.
接著告訴我們 保羅說我再說.
再引用《約瑟經》說.
以色列人不知道嗎?先由摩西說.
我要用那不成子民的惹動你們的憤恨.
我要用那無知的民.
不成子民,無知的民.
又是一個好像重疊的描述.
就是這些人不濟的.
在你們猶太人的眼中.
他們既不是神的選民.
但是神都會興起他們.

$^{961}$使他們可以得到那個恩典.
去觸動你們的怒氣.
就是說好像我特意去特別的祝福這些.
本來不配的 你自認是配的人.
你們就會被觸怒了.
以致其實希望他們悔改.
不是好像要特意去打壓他們.
所以如果我們真的看到一些事情.
是我們要去反思.
為什麼別人比我們更加謙卑,更加蒙恩.
那個教會會更加的蒙福.
我們就不要嫉妒 更加要反思.
所以每一次每個人,每個教會.
反思是最重要的.
因為這裡已經說了.
我特意要興起這些不是子民的.
這些無知的朋友.
不是說他們真的不識字.
而是表面以為在神的知識裡面的.
這些猶太人 自命優越的人.
所以有時候你看到那些驕傲的人.
他們那種輕佻對福音的態度那麼淺.
你們這幫人在那裡聚集.
又浪費時間 又浪費金錢.
很輕佻的去看.
當然有時候我們自己表現得.
比別人有夠窄的地方.
比別人覺得我們有挑剔的地方.
這個我們要接受,要承認.
但是其實雖然我們是這麼有限.
但是每次我們聚集.
有一種神秘力量的.
頂替我們一定要承認的.
我們是要和永恆,和做物主,和聖靈相交的.
所以我們要記得.
我們去到一個神的名字聚集的地方.
總有神聖的光輝在那個地方的.
所以你不要說笑.
你不要輕佻.
帶著一個虔敬的心.

$^{1001}$你就會得著的了.
真的是一個福氣來的.
你自己不斷成長 你自己不知道的.
你成為了教弟,姐妹.
你如果是一個成熟的人.
你坐在一個聚會裡面.
其實你不是應該去探討那個港道有什麼瑕疵.
而是你為那個港道者來祈禱.
神啊 你使用這位僕人.
讓他能夠祝福在這個現場的會眾.
或者他的信息成為罐頭之後.
當其他人開罐頭去聽這個錄音的時候.
又可以得到祝福.
你這樣的祈禱.
其實代表著你一個胸襟很大.
還有你是一個感覺聖靈在你生命在現場的人.
因為你不是一個旁觀者.
你是參與去敬拜神的人.
所以原來你崇拜都有很多東西學的.
不是沒有東西學的.
所以我們大家一起記住.
原來神是特意興起這些無不成子民的.
例如那些亞蘭人.
上個星期我們講到就是乃蠻.
亞蘭國其實是迫害以色列的.
將軍來到以色列.
就應該立刻抓他做人質 還救他.
以利沙有沒有搞錯.
但是神的工作就是.
興起那些無知的民 不成子民來去圍告他.
就是要觸動猶太人的怒氣.
這是《申命記》的32章第21節.
接著最後兩節.
又有以賽亞放膽說.
要放膽 先知要放膽.
就是那個時代很拒絕先知的預言.
警告 勸勉.
又有以賽亞放膽說.
會有尋找我的 我叫他們預見.
會有訪問我的 我向他們顯現.

$^{1041}$這個肯定就是神不斷在不同的人生.
來到去 透過不同的傳道者.
去讓我們認識 讓我們在人生有一個偶遇.
例如我的老師 或者我的同學.
其實人生走在一起 真的冥冥中有神安排.
雖然是父母幫我們選學校.
有得選 其實沒得選.
因為你派位也不好 一定是那個編定的.
冥冥中都總是有神的安排.
這個就告訴我們 尤其是舊一年.
神 那些人 就是提到外邦人.
沒有尋找 沒有訪問.
可以預見神向他們顯現.
相反猶太人 第21節.
至於以色列人 所有人很相對.
就是那些外邦人白白的得到被預見.
被神尋找 但是以色列人.
我整天伸手招呼那背影頂嘴的白色.
自己的選民 自認為神選民的猶太人.
神來到去伸手招呼那些背影頂嘴的.
怎樣伸手呢?他們都不肯接.
因為他們是有一個內心的背影.
對神的頂嘴 全神的來反應.
這個也是《耳塞》書的65章.
第一和第二節剛剛的相對.
我讀了這段經文之後.
我們就要verse by verse 逐節去看了.
16節這裡 彌賽亞被自己的子孫拒絕.
就是保羅引述的53篇的第一節.
就是彌賽亞受苦.
這位受苦的僕人的篇章.
跟大家一起看這段經文.
《耳塞》書 先喝口水.
第53章的第一節這一段.
這一段聖經都是保羅的favor.
引述過不止一次.
包括新約的作者無論是彼得,馬太.
都有引述過《耳塞》書的53章.
希望大家能夠背熟這一章聖經.
都是相當重要的.

$^{1081}$不過記住 今天的猶太人.
仍然死硬的他們的口就是.
是 這個受苦的僕人不是耶穌.
是我們這個民族 猶太人.
所以我們是受苦的民族.
我們是受苦的僕人.
但是新約的作者已經很清楚定性.
這一章聖經的寓言.
就是直接指到耶穌基督的.
第53章第一節.
「我們所傳的有誰信?」.
就是保羅所引述的.
「耶和華的膀臂向誰張開?」.
就是顯露 接著就說這個僕人.
他在耶和華面前生長如亂林.
像根出於乾地.
就是說他的成長相當坎坷 相當無助.
絕對不是含金的出世.
很多很多的困難 但是就好像證明.
他更加能夠成為人類的救主.
更加能夠到落地.
有一份很深的同理心.
所以如果我們能夠體會一些民間的疾苦.
其實都是基督徒應該有的神的那個性情.
這裡說「他無階型美容」.
「我們看見他的時候也無美貌 使我們羨慕他」.
所以這個可以想像得到.
在舊約不止一次.
無論是約瑟 原來約瑟是盡美的等等.
原來這個阿沙隆是很盡美的.
他的頭髮每年剪一兩次 重二千四百瓦.
就是說非常之.
他的頭髮是令他可以誇耀自己的美麗的美麗.
頭髮多的 在古代.
你可以想像到聖經時常都形容這些人能性.
但是這位救主神給他取的一個外在的形狀.
或者可能先天後天.
在一個勞苦的家庭 勞苦大眾的家庭成長.
做粗重事 體會很大的痛苦等等.
然後很自然 不只是他長得不好看.

$^{1121}$他是被藐視的 被人厭棄.
傳道者如果真的問心無愧.
被人藐視 被人厭棄.
就很貼近耶穌的成長和事奉.
如果一個牧者 所有人都稱讚你.
一定有的 是非常好的 感恩的.
如果那個稱讚是好的.
只不過是你根本沒有做到.
作為傳道應該做的那個責任.
那就不是很公平了.
你只是一個很好的先生.
其實你未必是傳神的道.
因為傳神的道肯定不是每個人都喜歡的.
如果那些人聽完不喜歡而對你藐視.
那就與你無關 是那個道的問題.
如果你是個人的 當然應該要改變.
應該要悔改.
但是耶穌肯定是沒有犯任何罪.
被藐視 被人厭棄.
所以其實我們作為傳道既然選了這條路.
你應該這個是常態.
被藐視 被厭棄應該是常態.
不要覺得很特別.
我要為自己據理力爭.
據理力爭就不要做傳道了 真的不要做.
因為本身就是這樣的.
多受痛苦 常經憂患.
你可以想像得到聖經的描述.
三言兩語已經告訴我們.
他的人生多麼不容易.
但是他就是這樣一步一步很堅實的.
走上十字架的路.
所以耶穌的腳踵是何等的加起來.
我們如果跟隨耶穌 我們都應該是這樣.
所以我們真的不要追求好像世人那樣.
說我要不斷被人高舉 被人讚賞.
誰不讚賞我 誰不支持我 我就是我的人.
那你不如自己開一個政黨就好了.
不如你就不要做傳道了.
因為真正我們的主都是這樣的人生的結局.

$^{1161}$好像被世人去掩面不看.
我們也不尊重他.
就是沒有人來對他有一個真正的生命的認識.
就是這樣 保羅告訴我們.
所以保羅看這段經文.
當他信了主再去看.
那個生命那個流淚不在話下了.
這一位的救贖主原來就是這樣犧牲.
所以保羅就是這樣一生願意向他的同胞.
來告訴他們 我們的彌賽亞.
不是我們這個民族 是耶穌 是神的兒子.
是瑪利亞的兒子.
是我們見過拒絕他的那一位耶穌.
所以保羅就一定看到這53章.
一定是受很大的感動.
因為他就是一個拉比.
他是很熟悉這些經文的.
所以對一個猶太人.
如果他能夠真是放開那個心懷.
去認定耶穌是主.
他們再看舊約就不得了.
真是不得了的.
因為他們一直受教.
要拒絕耶穌是他們的彌賽亞.
他們看到就是戴著一副有色眼鏡.
或者蒙了一層沙 看不到.
所以我們新約的人 我們外邦人.
我們也都相對沒有了那些包袱.
我們就沒有猶太人那麼深的那種歷史.
那種的眼淚 那種的偏見.
但是我們變相 如果我們不珍惜.
我們就很輕易就以為誰手可得.
我們又會比較沒那麼重視.
這些經文帶給我們的生命的感動.
你想想一個被藐視的人.
被拒絕的人.
常經憂患的人.
是憂患之子的這一位耶穌.
Man of sorrows.
sorrows是眾多的sorrows.

$^{1201}$我們一次的打擊 我們都頂不住.
是眾多眾多的打擊.
成為憂患之子.
你可以想像得到.
一個人服侍主 甘心樂意.
肯定是因為他想通了這位救主的美麗.
他的人生的美德.
以致我們願意放下我們自己有幸的.
可能我們本來可以享受的.
都願意放下 去為主 去行禮.
我們不是代表我們不應該有一點點休息.
運動 就不是這個意思.
而是我們整個人生.
只是為了運動 為了玩 為了旅遊.
真的要負心自問.
是不是距離我們對主耶穌的認識相差很遠呢?.
甚至越來越遠.
信仰成為一個宗教生活.
其實就很遠了.
因為我們的救主.
原來是一直被人嘲笑.
這樣走上十字架.
才成就了救恩.
所以保羅就告訴我們.
只是人沒有聽從.
這個就不是沒有傳.
是視而不見.
根本不認為他是你配的救主.
這個就變成了一種拒絕.
在以賽亞的年代.
拒絕這一位的彌賽亞.
每一個時代都有對耶穌的拒絕.
更加希望不是教會自己的內部的紛爭.
使一些不信的人拒絕.
其實是不好的.
所以我們真的要記住.
教會還是未得榮耀的一個地上的群體.
所以我們要戰戰兢兢.
帶著十字架的價值觀.
這樣在地上生活.

$^{1241}$否則我們就慢慢變成一個貴族.
變成一個享樂.
變成一個我們內部為偉的.
很多東西競爭.
你的言語 我的立場.
大家就會有很多的問題.
就會有很多的爭拗.
如果大家都是看到主耶穌這樣的犧牲.
大家都願意放下很多的誤會,成見.
那些不值得提.
完全不需要提的.
大家的心就是為福音.
還需要介意那些小事嗎?.
你可以想像得到.
如果我們真的已經沒有那種心靈.
這個是需要大家一起反思的.
不需要譴責任何一個人的.
每個人都有不同的時間.
那種對福音的拒絕.
不同的形式的拒絕.
懶惰都是一個拒絕福音.
所以我們都是撲朔迷離.
我們有這一份的恩恆.
印了第53章一節之後.
保羅就說.
可見信道是從聽道而來.
聽道是從基督的話.
這裡有幾樣東西.
之前那裡沒有提到的.
人如果沒有聽能就不能夠信.
沒有奉差遣就沒有人去傳.
所以報福音傳喜訊的人.
腳中皆未.
但是保羅這裡就聚焦.
聽是聽基督的話.
當然不是基督直接跟我們講話.
基督的話.
應該就是整個救恩的原本.
就是一個很整全的救恩.
所以為什麼我們能夠定性.

$^{1281}$五安門教是異端.
或者守望台是異端.
就是根本不是講一個整全的福音.
主耶穌祂那個神聖的身份.
和祂在地上完成的救贖的工作.
那個整全性.
所以原來信道是從聽道而來.
聽道相信是講聽明白.
這個救恩的原委.
所以我們的責任就是令到人喜歡聽道.
他願意去聆聽.
而且是一個開放的心靈去聆聽.
所以是真的要我們做了很多前期的.
囂觀,關係的建立.
一份很深厚的友誼.
而帶給他的信人.
以致他很想聽一下.
為什麼你會這樣信耶穌.
所以信道是從聽道而來.
聽道是從基督的話而來.
應該是這樣的一個邏輯.
猶太人就是聽而不相信.
但是外邦人是聽而相信.
所以真的是那個結果.
就是那個問題了.
信心的得救的信心.
是聽見,聽得明白.
而且是聽基督的話語.
就是說不是一個人間的哲學.
不是一個你的做人道理.
不是這個教會特別開心.
團契特別融洽.
那些是副產品.
真正的得救的信心.
是聽見以致透過基督的話.
這個真理,這個救恩的宣講.
就會建立那個得救的信心.
所以得救的信心.
必須是要透過他很認真聽明白.
一個整全救恩.

$^{1321}$以致他能夠打開心靈.
接受耶穌進入他的心裡面.
成為得救的那個信心.
所以這些告訴我們.
保羅就很清楚.
聽道是從基督的話.
我們教會初信班的課程.
真的是將很有系統的.
將耶穌的屋降生到上十字架.
是一個完整的課程.
十二堂來講明白.
其實每間教會的初信班都是大同小異.
將最精髓,最重要的救恩的原委.
人是罪人.
人需要面對神的審判.
面對將來永生.
靈魂離開這個身體去哪裡.
現在在地上.
什麼是我們需要去尋找.
神給我們空間,有機會.
這個就是一個整全的福音.
要讓他看到這樣的一個發言.
以致他就真的有一個.
不只是好奇.
而是他必須急切要去尋找這個答案.
所以信心是等於聽見家常基督的話.
所以保羅都相當清楚告訴我們.
尤其是當時羅馬書.
整本的新約還未完備.
還在陸續寫.
根本還沒有彼得前後書.
彼得還沒有坐牢.
還沒有去到人生最後階段.
連約翰寫約翰福音都還沒有.
都是主後的八十多年才開始寫約翰福音.
接著約翰一二三書.
他已經是晚年的一個長老.
已不鎖的長老.
其他的侍奴都離世了.
約翰是最後的了.

$^{1361}$後來甚至被放逐去到白魔海島.
領受了異象,寫啟示錄等等.
所以在使徒們的內心.
基督的話.
就是說真理是透過這些成文的記錄.
來給我們後世的人.
有真理的教導和基礎的.
所以如果我們只是道聽途說.
從來不讀經.
其實我們會大打折扣.
我們那個舊恩都會大打折扣.
所以異端邪教一定有它自己的經書.
不要聖經.
不要聖經本身已經是異端了.
是真的.
所以你沒有聖經一定是異端的了.
所以我們求主幫助.
教會不會希望一代一代.
其實現在科技這麼先進.
反而我們讀經的那種文能更嚴重了.
其實每個手提電話都可以安裝很多的Apps.
都有不同的版本的聖經.
但是讀的人真的很少.
這個不健康的.
因為基督的話就是得救的最重要的根據.
沒有基督的話.
那個聽只是人的理解.
甚至乎是人的哲學.
所以那些教主如果是厲害的.
洗腦式的 譬如這個雙氧水.
你可以想像到他不斷在推銷自己那一節經文的領受.
其實只是借那一節經文.
然後就說他那些發言.
他的心得.
他的所謂的領受.
然後售賣.
慢慢地將人帶到異端裡面.
聽他那個人的見解.
所以這些絕對不是基督的話.
所以我們真的要裝備自己.

$^{1401}$讀經不是為了什麼.
是為了我們很想認識神.
所以仍然是那一句.
《提摩特後書2章》.
勞力的勞夫.
提摩泰雖然你不懂得耕田.
但是你是神的工人.
勞力就會得糧食.
你就會得到自己對神的認識.
這些是人都搶不到的.
這些是生命的道.
在我們內心.
我們自己先來養活自己的心靈.
所以勞力的勞夫.
你當先得糧食.
我者說你要先賞.
因為凡事主必給你聰明.
就是說你自己要舉一反三.
為什麼我會這樣說這三個例子呢?.
提摩泰你要有這三種的警告.
你是一個軍人.
在軍中當兵就不要給世無前身.
就是說既然你已經是承諾了自己給神.
你就要忠心自治的為福音.
你整個生命不斷投入去為福音.
你就整個人就變成福音了.
所以這個就是在軍中當兵.
不將世無前身.
就不是說不可以結婚.
不可以生孩子.
不可以吃自助餐.
不是這些.
這些就是世無.
而是做的那些事情.
你仍然那個心都是為神去做的.
不將世無前身.
接著人就再從上比武.
你要體會自己是一個嚴格的運動員.
你是要fight lawfully.
你要按規矩來希望得冠冕.

$^{1441}$所以如果我們用一些錯誤的手段.
去得到人去建立教會.
最終都是會散的.
因為你正在用你的手段.
所以人若在牆上比武.
非按規矩就不能夠得冠冕.
這些都是傳道者要銘記的重要價值觀.
保羅可能見到他的同伴.
有些是跌倒.
底馬灘那現今的世界.
就離棄了我.
往跳薩羅尼加去.
格勒斯去哪裡.
銅像亞歷山大多多的駭我.
主邊報應他.
你又要防備他.
所以同一個時代.
總是有跟隨主到底.
或者有些人半路中途又會有一些出差子.
便捷等等.
所以我們要記得.
但是更寶貴的就是勞力的農夫.
提安泰 我再敦促你.
你自己要好好地勤奮地去為福音.
去養生.
生命的生命來服侍主.
所以這三個的身份.
是一個傳道者都需要這樣去反思.
所以你見到大企業那些成功的人士.
他們是非常專注在他們的事業裡面.
而最終有一點點的成就.
即是地上的企業都尚且是這樣.
何況我們作為牧者傳道.
豈不是應該比其他人更加的專注.
更加的用心呢?.
因為是要將人帶進永恆的救恩.
所以你們說我們知道.
這裡就是保羅說「但我說人沒有聽見」.
誠然聽見.
其實就是很想說.

$^{1481}$再一次說第十八節.
沒有聽見其實聽見.
想來引述就是說.
猶太人 以色列人.
其實他們是故意地拒絕福音.
福音已經廣泛的被傳播.
所以引述第十九章原來很精進.
保羅為甚麼可以這麼有那種智慧.
可以突然間 那個時候沒有電腦.
也沒有AI 沒東西可以問.
連一個可以按鍵的機器都沒有.
靠腦去想的.
他就馬上想到.
是的 詩人曾經說過 應該是大衛.
說宇宙窮倉是印證了造物主的存在.
是沒有一個地方是看不到的.
只要你很平生而論 你是會看得到的.
但是 我看看時間是否夠了.
因為我看不到那個錶.
OK 還有一點點時間.
所以他在第十九章的第四節.
整篇詩篇的確是在說.
從那個普遍的啟示.
去到聖經的律法書的.
不如先讀一讀吧.
我都怕大家都不懂這些經文.
都是這些很美麗的詩.
如果能夠背熟它都好.
十九篇.
《朱天述說神的榮耀》.
大衛的詩 沒錯了.
窮倉傳揚他的手段.
大衛也有很多時間接觸大自然.
當時應該是他看不到的那些滿天星.
比現在我們在香港高樓大廈遮住的星系多.
他說「這日到那日」.
就是說那個時間的普遍.
發出言語.
「這夜到那夜傳出知識」.
是在說宇宙 其實哪裡有聲音的呢?.

$^{1521}$我們收到的可能都是電波.
電池等等.
但是詩人就是用這種優美的描述.
在說大自然說話的.
神要讓大自然向人類說話.
那種結構的精美的精密.
種種的安排的佈置.
是不斷在傳遞一個訊息.
就是做物主的偉大豪漢.
傳言的美麗等等.
接著第三節.
「無言無語果然也無聲音可聽」.
但是他的量大通遍天下.
他的言語傳到地極.
不僅是一直在說.
而且沒有一個地方.
所想的行規就是任何的地方都被傳遞清楚.
但是現在保羅就借用了這麼優美的詩篇.
去講述其實福音同樣都已經被傳遞了.
如果宇宙這個很貴格的存在.
都是這樣印證著有神.
這位美麗的神存在.
主耶穌的降生傳遞這個救恩.
就更加告訴我們這位做物主.
是多麼愛我們犯罪的世界的世人.
所以保羅從這裡真的很美麗的描述.
但是當然是一個負面的.
因為「成言聽見」.
人沒有聽見.
你們以色列人沒有聽過.
聽不明白.
其實是聽見的.
所以這樣其實是一個嘲笑.
一種挖苦.
對以色列人的這種視而不見的那種硬心.
是一種責備.
詩篇都已經講了.
這麼多年前.
大衛的年代主前的大概700年多一點的時間.
對不起.

$^{1561}$主前的1000年多一點的時間.
因為700年已經是北角的被擄了.
你可以想像得到.
大衛的年代都已經有這種啟示.
有這種的發言.
現在就借用這一節來講述.
福音更加的浪離盡至.
當然不是那個領域去到中國.
去到其他的南美洲.
就不是講這個領域.
因為還是局限於人類那個可以去航海.
或者那個科技還沒有那麼發達.
還沒有那麼普遍.
都是靠一雙腳這樣走.
但是講的就是福音的傳遞.
或者什麼是福音呢?.
其實在保羅的年代已經是相當清楚了.
什麼是福音已經是這樣明智了.
所以每一個時代的教會都要捍衛.
福音的那個唯一的傳遞.
這個是很重要的.
我們教會存在都應該是這樣.
告訴我們知道.
接著就是第19節.
我們看下面.
19節.
所以人在救恩面前的責任.
就不是說沒有.
就不是說我聽不到.
是聽得見的.
所以是聽見之後的那個回應.
如果沒有回應.
其實就是等於沒有真正的聽見.
最後我們都看完它了.
第20節.
我們下一次就可以集中開始講第11章.
當然還有一節沒講都不要緊的.
接著保羅再一次.
第19節.
我再問以色列人不知道嗎?.

$^{1601}$先有摩西說.
我要以不成國的.
這個是新和本.
激起你們的嫉妒.
我要以愚頑的國.
惹起你們的法路.
這樣來到好像是刻意激猶太人.
其實是要使他們有一個反思.
要使他們有一個悔改.
不要以為仰尊處優.
以為自己是坐定了六.
是亞伯拉罕的子孫.
我們有聖殿.
保羅的年代還有聖殿的.
保羅在以弗所最後一次的宣教.
他回到耶路撒冷.
大鵑巷上去.
他是去聖殿找祭司幫他行潔淨的禮.
因為他許了一個男婿二人的願.
就要行潔淨的禮.
所以他仍然自己都是循規蹈矩.
但是因為有人見到特羅非摩這個外邦人.
在聖殿的附近出入.
以為保羅帶這個外邦人進入聖殿.
因為純粹是一個觀察見到.
所以這個誤會就帶來要立即捉拿保羅.
因為本身他們在安提摩.
對著保羅已經很生氣了.
已經是恨之入骨了.
上到來還沒有證據.
你還不是帶外邦人來糟蹋我們的聖殿.
這樣的無妄之災莫須有罪名.
就是這樣捉拿保羅.
但是保羅去聖殿敬拜.
還有聖殿.
是主後的70年.
聖殿才被提多將軍移為平地.
你可以想像得到猶太人那種宗教的自信心.
宗教的那個堡壘.
使得他們看不到實的舊恩.

$^{1641}$就沒有了這個真的得救之聲.
所以他們就會嫉妒.
他們就會發怒.
為什麼外邦人說自己可以無恩.
或者自己可以親近神呢?.
親近神要去聖殿的.
要祭司驗祭的.
為什麼外邦人可以這樣做呢?.
所以保羅就引述《申命記》的32章第21節.
用不成指紋做了幾分鐘.
不如多揭這一段經文.
還是已經夠了.
再說這一段.
就是《申命記》的32章.
我們就完結了.
我們下一次再從這裡開始.
到底他上下文想說什麼呢?.
32章其實就是.
這本事因為手很難揭.
32章的第21節.
我揭了.
摩西的歌.
然後我們不說那麼多了.
從15節打後讀起.
耶穌倫漸漸肥胖,粗壯.
這個原來就是開始.
越來越以宗教成為他們的一個護身符.
光潤,踢跳,奔跑.
便離棄做他的神.
興旱救他的盆石.
敬拜別神.
觸動神的憤恨.
行可憎惡的事.
惹了他的怒氣.
所祭祀的鬼魔並非真實.
乃是素不形式的神.
是近來新興的.
是你們早所不畏懼的.
你興奮新來的盆石.
忘記產來的神.

$^{1681}$是摩西的郎中的詩歌.
好像數以色列人一樣.
或者預言他們將來沒有領袖帶領.
他們就會不斷地.
其實內心從來都沒有得救過.
出埃及只是看熱鬧.
看到摩西大主很興高采烈.
當有少許挫折便馬上埋怨摩西.
這些典型的跟隨的.
好像羊群一樣.
19節.
耶和華看見他的魚類惹動他.
就厭惡他們說.
我要向他們掩面.
看他們的結局如何.
他們本是極乖嬌的族類.
心中沒有誠實的魚類.
留意21節.
就是這一段聖經.
他們以那不算為神的觸動我的憤恨.
以虛無的神.
即是回應上文那裡.
惹了我的怒氣.
然後留意.
我都要以眼瞎眼.
我都要激你們這些以色列人.
你這樣激我.
令我憤怒.
我都這樣激你.
我又要以那不成子民的觸動他們的憤恨.
以愚昧的國民惹了他們的怒氣.
所以因為舊約的版本跟新約的版本不同.
所以字眼上是有些不同的.
但意義是完全一樣的.
因為舊約馬索拉差本.
新約他們用七十字譯管.
所以是有些許的字眼是不同的.
所以你不要以為聖經.
為什麼不是跟舊約一樣呢?.
是不一樣的.

$^{1721}$因為版本不同的.
變成我們真的看到不同的翻譯的意思.
但意義完全一樣.
現在神就是要興起這些不成子民.
這些無知的民.
其實就是好像讓本來不可能得救的外邦人.
都可以得到.
記得不是這樣字面.
什麼看我們無知.
看我們是不成子民.
我們很有民族的尊嚴.
你現在神是這樣看我們外邦人.
記得是一個修辭上面.
這邊是《申命記》摩西的歌.
所以不要因為這樣.
突然間你就有一些這樣的誤解聖經.
貶低了外邦人就不是了.
反而就是說猶太人這樣拜假神.
神就要來去興起一些本來不應該得到的.
來去激怒他們.
希望他們有一個知錯反思悔改.
所以今晚大家.
我想我們都要不斷反思.
我們信主.
我們是不是真的繼續蒙恩.
繼續蒙神的祝福呢?.
或者我們已經有一個宗教的很大的堡壘.
你不怕嗎?.
我們在錦繡堂這個牌頭.
多大一個湖邊的教堂.
拿獎的多厲害.
說神不祝福沒可能.
不過記住.
猶太人的聖殿.
在以賽馬的年代有.
考勞多元被擄就沒有了.
接著回去重建.
但是到了耶穌的年代.
到了保羅的年代.
還有聖殿.

$^{1761}$但是耶穌已經預言了.
將來一塊石頭都不在另外一塊石頭上面.
這個急迫輝煌的聖殿被夷為平地.
因為它已經成為了一個賊窩.
宗教已經成為了一個方便.
有宗教的禮儀.
沒有了宗教背後帶人去到神面前的那個信心.
所以求主幫助我們能夠.
我們一群聽《曹經》的弟兄姊妹.
我們是有得救的信心.
我們真是被神無條件地揀選的.
帶著一份感恩的心來服侍.
我們一起禱告.
感恩我們能夠有系統地將這幾節聖經.
和上晚文的關係一起分享勉勵.
求主幫助.
人說的是很悠閒.
很盼望神主的聖名.
來作供在我們的心裡面.
讓我們能夠看到.
保羅其實是很語重心長.
勸勉他們的同胞.
來有一個反思,悔改.
還有真是有得救的信心.
就是要聽明白基督的話.
來得救.
求主幫助.
多謝你讓我們能夠有救恩.
有福恩去傳.
很盼望更多人來認識主的救贖恩典.
也都讓我們在教會.
真是看到什麼是最最重要.
福恩的傳開.
人得到神的救恩才是最重要.
人和人之間更加不應該提升到.
一種好像你死過亡那樣的.
錯誤的,扭曲的一種心態.
教會根本就不應該存在這樣的思想.
所以盼望主的聖名幫助我們.
每一次都有份警醒的心.

$^{1801}$來到你的面前.
為你感恩禱告.
奉耶穌基督的名.
阿們.
多謝各位的參加.
祝福大家.
\newpage



\section{羅馬書}
\label{sec:ZdF__6Sl2No}
\textbf{2025-09-26 周五晚查經班《羅馬書(23):分題:神選民與恩典—神的不變救恩》 講員:曾裔貴牧師}
\newline
\newline
連結: \href{https://youtube.com/watch?v=ZdF-_6Sl2No}{\texttt{https://youtube.com/watch?v=ZdF-\_6Sl2No}} ~~~~ 語音日期: 2025-09-27
\newline
\newline
\hyperref[sec:1GDpbDYTYf0]{\small{< < < PREV SERMON < < <}}
~
\hyperref[sec:index_chronic]{\small{[返順時目]}}
~
\hyperref[sec:index_scriptual]{\small{[返順卷目]}}
~
\hyperref[sec:bw6FVf5f_n0]{\small{> > > NEXT SERMON > > >}}
\newline
\newline
$^{1}$第一節會晚安.
我們一起再有機會去看羅馬書.
我們先有個祈禱.
很感恩天父讓我們再一次恢復羅馬書的言經.
我們一直在第九,第十章.
看到保羅他的思想.
對以色列人他不會有一個拒絕.
反而就是將神的計劃和人的限制.
人的失敗基於人自己的自以為是.
以為自己是神選民.
就會有一種必然.
好像是誰手可得那樣.
必然的會得到救恩.
多謝主讓我們能夠透過仕途保羅的教導.
還有在第十一章他有一個很大的奧秘.
都讓我們揭開對於神永恆計劃的了行形式.
做物主是這麼偉大.
是一位前愛,真理,信實的主.
讓我們能夠認識你.
尤其是透過聖經的啟示.
讓我們知道你在人類的救贖計劃裡面.
那個規模,那個計劃.
我們今天晚上都會接到去.
十一章這些經文的時候.
都讓我們能夠一個這樣的了解.
無論是透過引述舊樣的經文.
都讓我們認識保羅的內心.
他當時寫作的時候那個內心的慨嘆.
還有那個對神偉大的救恩的讚嘆.
都讓我們能夠認識.
讓我們都有一個很熱的心情.
能夠真的感受到我們原來蒙到恩典.
是多麼的難得,寶貴.
讓我們能夠珍惜.
還有一個在人生的責任.
就是要傳福音.
讓人聽明白聖經真道.
我們仰望,感恩,禱告.
奉耶穌基督的名 阿們.
再一次我們回來看第十一章的羅馬書.

$^{41}$因為都停了一次.
都想和大家再重溫第九到第十章.
試圖保羅他的思想.
讓我們能夠掌握得比較緊貼一點.
第23章我用的題目.
其實都和第十一章有直接關聯.
就是神的選民和恩典的那個關係.
當然更加重要就是整個十一章.
想告訴我們神的不變.
所以在救恩上面.
他的計劃對以色列人的得救.
他都不會說好像沒有他們的份一樣.
因為保羅的那個比喻就說到.
猶太人,以色列人在救恩.
一個很重要的概念.
就是他們是橄欖樹.
所以反而外邦人能夠蒙恩.
保羅看其實我們是那些野橄欖枝.
所以如果野橄欖枝都尚且蒙了這麼大的恩典.
神的救贖.
何況那個橄欖樹.
神自己親自去栽種的.
豈不是更加的應該蒙恩呢?.
所以這樣理解就好像救恩從來都沒有失敗過.
我們可以進入那個經文的討論.
讓大家再去看第九十章那個大綱那個脈絡.
保羅他整個思路的架構.
在第九十章.
當然他大前提就是.
他一開始就概括.
為甚麼以色列人不得救呢?.
是基於甚麼原因?.
是否神已經撇棄了他的子民呢?.
為甚麼現在外邦人.
不需要經過很多的儀式.
進入猶太教.
就可以白白的清義得救呢?.
而且那個得救是很真實的.
經歷到聖靈的同在.
聖靈的更生.

$^{81}$例如第五章提到的.
還有那個活培的盼望.
將來身體會得贖.
這些都是外邦人一直準備迎接的.
將來那個永恆救恩.
但是猶太人有沒有墳呢?.
是否有另外一個救贖呢?.
這些都是保羅在當寫完第八章之後.
他都是用了三章的聖經.
來去探討為甚麼現在猶太人這麼背影.
怎樣去面對將來的審判.
神的救恩有沒有愚昧到.
猶太人都可以得救呢?.
所以又讓我們能夠透過保羅在這三章.
很獨有對猶太人的那個救恩的情況.
又讓我們認識.
當然單純是認識保羅他的思想的周密這麼簡單.
而是神透過保羅.
他自己都是一個過來人.
他寫出這些的教導.
就讓我們看到.
原來背後神真的有他自己從來不改變的計劃.
雖然人好像一直攔阻著神向萬民的私恩.
阻撓 但是沒有影響到神自己一直的作為.
以及推進救恩的拯救.
所以我們有些重大的一些脈絡.
一些架構 一些保羅給我們看到的救恩的大綱.
我們就要思想了.
我們現在就站在外邦人蒙恩的那個立場.
我們去看原來在整個救恩裡面.
神沒有虧欠任何的人 包括猶太人.
因為就算在猶太人當中.
原來不是每一個都是蒙神的揀選的.
造物主是有他的主權,權柄來揀選人的.
所以第九章已經很清楚講了救恩.
原來在神的永恆的計劃裡面.
他揀選了一些蒙恩得救的人.
去到第十章 原來這個揀選.
和人自己要負的責任是完全沒有抵觸的.
所以對我們來說就很矛盾了.

$^{121}$揀選了就不需要負責任 負責任就不是揀選.
但是原來兩者是一個很和諧的調整.
一個這樣的和諧的合一.
導致這個人既是神的揀選.
同時間也是他對福音,對耶穌基督.
對他聽明白了救恩關乎耶穌的釘十字架.
而他作出的那個信心的回應.
就是成為了他得到信心的意義的那個標準.
所以你就可以想象得到救恩在神這麼偉大的計劃.
和他的智慧裡面 因為十一章的二十六節.
特別強調神的智慧和知識.
這個是非常深奧的 在保羅的眼中.
所以救恩就是反映著神的那個智慧.
不單純是能力 不單純是權柄.
而是神的那個智慧 神的那個知識.
神怎樣能夠讓一切的事情.
在他的那個作為裡面一直這樣開始.
所以原來神對整個宇宙的救贖計劃.
是從來都沒有改變的.
我們可以來到他 稍微看看.
你可以說神的揀選很自由地.
神一向都是自由的 在他的那個旨意裡面.
他的揀選不是憑人的行為.
或者是他的血緣 他的血統.
或者他是哪一個的子孫.
而是在乎神的那個揀選 很特別的.
在起初揀選的時候.
都不代表著那個人本身就是大有信心.
或者那個人本身他的各方面的道德.
對神的心意是很出類拔萃.
例如雅各就已經是一個這樣的例子.
反映到其實他真的不是這樣來對神.
好像他的父輩以實 或者是亞伯拉罕那樣.
是神不斷將他修剪.
導致到他原來最終都能夠.
被神揀選的行列裡面的這位雅各.
所以這裡就已經反映到神的揀選.
尤其是透過舊約 很多聖經的人物.
與神接觸 在他們的人生去認識神 跟隨神.
又讓我們看到 原來真的不是他們的行為.

$^{161}$更加不是他們的血統.
所以基督教肯定就不是靠律法.
律法只是一個標準告訴你 你應該這樣做.
不代表律法你行完.
你就能夠好像達到那個信心的標準.
所以是必須憑信心去接受.
所謂的信心就是口裡認耶穌為主.
心裡面相信神已經使到耶穌基督復活.
這個所謂的口認心信的接受.
這個就是信心 就是將我們整個的生命.
投向這位釘十字架的耶穌基督.
就是得著這個榮幸救恩的信心的意義.
所以都很清楚 第九 第十章告訴我們.
基督教不是靠律法 是憑信心的意義.
愛邦人的基督教 原來都是神的救恩的計劃.
所以其實那個是愛邦人 蒙恩.
其實是猶太人對神的背影.
堅持硬頸要行律法.
要以自己是亞伯拉罕的子孫後裔.
這樣來用這種身份 血統來換取救恩.
其實都是行為的一種.
但是一講到行為.
就沒有人可以滿足.
貼近神的標準 心意.
就是神對公義 對罪的要求.
所以想像得到 猶太人拒絕福音.
就是他們自己的硬心.
是他們自己的堅持自己是義.
而拒絕神的義.
所以有時候我們都要手停一停 要反省.
我們慢慢順主來.
會不會堅持我們自己的傳統呢?.
有時候連真理我們都來不知不覺的去抗拒.
因為猶太人就是日積月累.
在律法的分途底下.
在很多舊約那些的戒條.
拉比的演繹 教導底下.
慢慢越來越牢固的宗教化了自己的身份.
自己這個民族 慢慢就習以為常.
這一種習非誠實的這樣的義.

$^{201}$就難阻住了.
當尼采娜出現來歡迎自己就是神的救贖.
但是猶太人就已經不能夠用他們信心的義.
來迎信耶穌了.
因為他們已經很牢固的有因行律法而來的義.
因這個身份血統而來的義.
但是這些的義在神的眼中其實是不足夠的.
所以本來一個很好的傳統.
是要讓人更加有機會去認識耶穌的.
反而成為了一個絆腳石.
好像從小回基督教學校.
每個星期有早會 有老師祈禱.
有一起去借飯.
但是原來變成了一些完全沒有靈魂的.
每個星期學校的一些禮儀.
慢慢畢業 基督教成為了一個可能是婦女.
可能是一個對他們的一個誤解 難阻.
以為基督教不過是做一些儀式.
我根本不會接受這樣的這一套.
尤其是因為那些張力就是老師.
有時候如果那個過程不是那麼公允.
或者有時候有一些不好的 大家的待遇.
就會遷怒於基督教 其實不關基督教的事.
可能是人自己學得沒有那麼好.
而導致很多的學生讀基督教學校.
對教會 對學校 對基督教.
都會有一定的反感 抗拒.
原因其實可能就是他接觸的基督教.
只是很外表的儀式 傳統.
所以這個就帶給我們.
猶太人就是這樣落在這種宗教法之下的.
這樣的拒絕福音.
福音本來讓人真是帶著一個完全打開的心靈.
無助地去到耶穌那裡得到憐憫.
所以猶太人就是拒絕憐憫的一個民族.
所以這些雖然是跟猶太人講.
其實我們成為教會的人 信主內在.
尤其是我們連讀羅馬書.
我們更加都需要去反思.
我們自己現在那個信仰.

$^{241}$是不是已經越來越系統化了.
越來越傳統化了 甚至越來越宗教化了.
以致我們這樣做 這樣做這些儀式.
但是沒有真正跟耶穌基督有生命的相交.
來到教會都沒有真正想著去親近神 擊敗神.
這些有時候我們自己每個星期.
做出來的小小的那些行為.
其實不知不覺就會被一些熟人能夠看到.
當然不是要譴責 而是真的彼此的勉勵.
彼此的提醒.
我們記住是來到敬拜神.
我們真的要帶著一個空虛的心靈.
來到神的面前 朝拜祂.
得到祂的安慰 祂的加利.
這樣的一個信仰人生.
所以不是宗教化 本來宗教本身不是錯.
但是慢慢就是有了那個儀式.
沒有了那個實際的對神的敬拜.
所以這個部分告訴我們.
神在祂的救恩裡面的主權.
就是真是神憑祂的旨意 憑祂的自由.
憑祂的權柄 憑祂的那個揀選.
來將人選照的.
絕對不是任何一個人的行為.
所以我們能夠順處.
大家來自不同的背景都能夠順處.
而且應該不斷那個心靈越來越柔軟.
對別人越來越有一種的感恩.
是因為救恩帶給我們的生命改變.
如果我們覺得是必然的.
是真是 好像是魯奉的.
那就很大件事了.
慢慢我們就未必會再珍惜.
我們覺得必然是我的.
我們必然是要得救的.
祂必然是要拯救我的.
這樣的想法.
就已經和救恩完全離開了那個標準那個行列了.
所以大家要非常的小心.
我們不要變成另外一種基督教的猶太人.

$^{281}$接著 這裡告訴我們.
第十章就很明顯講到我們必須要去傳.
因為必須是要人自己打開心靈.
去明白那個救恩那個真理.
他就可以得救.
所以那個得救的信心.
或者那個信心的意義加給什麼人呢?.
就是聽道 聽明白神的道.
而願意打開心靈向耶穌呼求的人.
求告耶穌的人就必然的得救.
不再分猶太人 不再分外邦人.
不再分任何一個民族.
不再分你是貴族 還是男是女.
無論地上怎樣將人分開不同的等級.
但是在救恩的面前.
所有人都必須聽明白聖經真理.
所以傳道的重要性就是這裡了.
所以教會存在這個世界.
不只是我們在一個教堂裡面.
來滿足我們的無論聽道.
無論去敬拜 無論成樂的昌頌.
或者在一個教堂裡面去禱告.
全部都是對的.
但是這一切都是想告訴我們.
我們的存在就是為了要出去傳道.
我們的聚集就是為了要出去傳道.
所以我們的聚集應該是一個裝備.
應該是一個身份的確認.
應該是一個在我們這個地方.
我們深深感受到我們是無福的.
我們是一幫真是有幸福感的人.
但是不是終點.
如果我們一幫人很開心.
我們的子女讀書很厲害 很乖.
入名校 將來做事很有成就.
全部都是對的 但是那個不是終點.
那個是一個過程.
既然我們已經感受到這份幸福感.
其實是來自救恩 來自耶穌基督.
將我們拯救 將我們揀選.

$^{321}$所以第十章很重的責任.
就是我們必須 教會必須要去傳福音.
將耶穌基督有系統地 有計劃.
有方法地去告訴人.
使得那些人可以得到 聽明白.
原來耶穌和我們的生命有關.
他願意相信 他一求告就可以得到了.
所以信心是來自聽到.
聽到是來自有人去傳道.
所以教會一個相當重要的責任.
無論本地的傳道 海外的傳道.
或者是支持海外的傳道等等.
用不同的方式 或者禱告 或者是金錢.
或者是猜派等等 我們想得到的方法.
無論如何 我們就是要一起去傳揚福音.
所以教會如果都不合一.
這就是第十二章開始就講這些問題了.
如果猶太人和外邦人.
還圍著那些吃什麼的食物.
還守不守安息人 看得很大.
大到要分開 其實可惜到不得了.
因為最重要的就是大家共同的看到.
大家的分歧 所謂分歧.
大家有不同的背景 大家有不同的民族.
但是在救恩底下 大家可以繼續.
外邦人的生活方式 猶太人的生活方式.
大家彼此的尊重已經不需要提升到不夠的層面.
反而就是大家在這樣的不相同底下.
合一的去傳道 合一的去支持史努保祿.
和他們的宣教士 傳道者 教會.
一起很有系統地去對外去傳福音.
抗衡的是羅馬帝國的明顯的逼迫.
或者暗中的一些可能是道理.
或者一些哲學的影響.
因為洛斯底主義在第二世紀就比較濃厚了.
在保羅的年間都還沒有嚴重地影響教會.
很多那些哲學思想是會帶來影響.
所以教會的存在很重要的就是.
真的要將福音講明白.
所以這個真是我們每一位讀者.

$^{361}$必須要怎樣能夠用最好的方法.
去集合大家去為福音來努力.
所以在福音底下是有很多不同的崗位的.
有些人可能是接單張.
因為今天和昨天有姐妹回來.
接了那些福音單張.
明天就放進那些福袋裡 明天就去派.
接單張這個舉動本身都已經很神聖.
因為如果假如真的有一些朋友,小朋友.
看到那些單張又明白.
他自己又真是心靈裡面很單純地向耶穌求告.
那就得救了 所以你可以想像得到.
原來每一個崗位在救恩.
在耶穌的大愛,拯救底下.
都是這麼重要,這麼寶貴,這麼值得尊重.
所以大家一起努力.
這是第十章全章裡面給我們很重要的大綱.
所以你想想救恩神的揀選.
和人的責任是調和的 是沒有衝突的.
求告主名 所以這裡告訴我們.
並沒有分別猶太人和外邦人.
同有一位主在天上.
祂會厚待一切求告主名的人.
十章的十三節.
所以救恩是向所有的人類開啟的.
所有的人都需要救恩.
不同的階層,不同的背景都需要救恩.
所以盼望我們大家都找到我們人生裡面.
我們自己有負擔的一個.
可能是你凝結了這個群體是我很有負擔的.
譬如有些人 譬如陳靜娟姑娘.
她很有負擔在尼泊爾這個族群的.
整個的心思,意念都放在民族的信主底下.
有些是泰國的 例如我們在紅一堂.
高淑芬姑娘也是對泰國人非常有負擔的.
另外有其他的牧者同工.
譬如金峰醫生.
對阿拉伯人信主是很有負擔的.
所以大家都有不同的那個群體.
在你的心靈裡面.

$^{401}$神很想放在你和我的心靈裡面.
我的同學何世傑博士他也是讀神學的時候.
他說他也是對阿拉伯人有負擔.
很想帶他們這個族群有機會信主等等.
所以每一個的同工,牧者.
總是在他的門罩裡面.
有一個神給他的那個心靈的感動.
所以我們大家都一起.
就算我們是普通的一般的信徒.
雖然我們有我們的工作.
但是我們的心靈裡面.
如果對某一個階層我們是有負擔的.
不妨在那裡繼續去禱告.
尋求神的心意,帶領.
譬如有牧者特別是對小朋友是很有負擔.
很有興趣.
慢慢他不知不覺發展到他自己的恩賜.
他整個人的個性.
都能夠被小朋友接納,歡迎等等.
所以如果救恩是向所有人開放.
我們既然是蒙了神救恩的人.
我們真的要想一下在我們人生.
有哪個族群,有哪一類的人.
我們是特別有負擔的.
我們就多點留意,擺上禱告.
還有我們裝備我們自己.
怎樣能夠有神為我們開路去配合.
使我們能夠帶到這些.
我們心裡面感動的人來去信主.
所以保羅我們等一下十一章會講.
他終於他的經文.
「大心碎小抽」?誰找你來的?.
老闆不夠錢就在那裡碰水.
不好意思 如果還沒有關掉.
就要靜音 大家不要介意.
你可以想像到保羅他說他是外邦人.
當然是企業企業 十一章二十六節.
本來我們等一下才講的.
我們看看這段經文.
保羅說他是外邦人的使徒.

$^{441}$對不起 應該是十三節才對.
我對你們外邦人講這些話.
因我是外邦人的使徒.
所以敬重我的職份.
你可以想像到保羅是盡力.
令到這些外邦人的祖先.
對自己的職份有一個敬重.
這個敬重包括他整個人.
化身成為一個不再仇恨外邦人.
不再聞到別人外邦人的氣味.
就已經有拒絕 有反感.
有那種身體最誠實的態度.
有那種身體最誠實的那種的厭惡.
那種的苗嘴 已經沒有了.
很難的 一個拉比.
能夠整個人化做一個外邦人的使徒.
所以十一章的十三節.
也是很值得我們大家來消化一下.
我們是不是已經能夠.
好像保羅那樣道成肉身呢?.
主耶穌由榮耀的天上來到世界.
真是娶了一個人.
見到他沒有任何一個DNA.
是神聖的一個存在.
而是只會人見到的.
如果沒有信心.
見到的就是一個木匠的兒子.
一個在拿撒勒.
一個很窮乏不起眼的地方.
成長的一個街童這樣的身份.
一個很老到的家庭.
來長大的一個普通的猶太人.
所以主耶穌他那個內心的榮耀.
當然透過神蹟那些人會見到.
但是對主的身份.
真是不是所有人都見到.
是真是有信心的人才能夠認識.
例如有一個白夫長.
他的僕人差不多死了.
因為他是一個很好的官.

$^{481}$又幫百姓.
又幫猶太人捐錢.
收集會堂.
所以當地的一些長老幫他去求耶穌.
那麼耶穌就肯去他家.
當這個白夫長聽到之後.
這是《路加福音》第七章.
聽到之後就跟那些人說.
我不敢當 耶穌來我家不敢當.
他只要講一句話.
我的僕人一定會好的.
那個時代是沒有任何通訊.
沒有任何電話 沒有任何WhatsApp.
沒有任何的互聯網.
講一句話 那種遙控是很誇張的.
你可不可以想像得到.
這個白夫長是深深感受自己的不配.
但是又同時間認定耶穌有這樣的權柄.
因為他自己是白夫長.
他講了一句話.
我在這樣的權勢底下.
我叫那些人離開去做事.
那些人就要立即去的了.
叫那些人馬上回來.
那些人就要回來了.
就是那些叫僕人 那些叫兵丁.
所以他很感受到.
耶穌有那個權柄 一句話就可以了.
所以當耶穌聽到這個白夫長這樣講的時候.
就稱讚這個白夫長.
他說我在以色列當中.
沒有見過有這樣的信心的一個人.
就是說連他跟隨著他的十二使徒.
都沒有這樣的信心.
這個白夫長完全是外邦人.
都有這樣的信心.
這樣就可以想像得到.
這個就是使徒保羅真正的道成肉身了.
他說我尊重我的職份.
神既然不准我只是向猶太人遊說傳道.

$^{521}$而是要我去那麼遠的向外邦人傳道.
所以保羅就忠於敬重他的職份.
首先喜歡 首先自己想一下.
既然我是外邦人的使徒.
我是不應該去強調哪些飲食是比較重要的.
猶太人有些東西是不能吃的.
保羅能不能去到這些外邦人當中.
那些帶著血的食物 那些肉類.
他會不會吃呢?應該是不會的.
但是當他自己一個人跟猶太人一起的時候.
他就立即很尊重猶太人的那個教導傳統.
又不會令到猶太人自己難堪.
他又會令到他在神的面前.
有他自己的謹慎,謹守.
這種這樣的道成肉身.
沒有將自己的利益放到最大.
這個就是使徒的情操.
就是對舊恩那個心悅誠服的愛.
和一個的馴服.
所以真是的 牧者我們都應該要這樣學習.
突然間想起一段英文.
《哥林多侯書》應該是第十章.
因為在跟那些假使徒做一個對比的時候.
保羅就再一次說他們作為使徒.
這次是對哥林多教會.
這樣來說他這個身份.
他怎樣理解自己使徒的身份.
還有他又期望弟兄姊妹哥林多教會.
都要這樣理解使徒的身份.
有一段這樣的經文來跟大家彼此分享一下.
我稍微看一看 找到了.
十章的第七節.
如果大家不熟悉上限文.
我讀多一點.
都跟保羅說的使徒身份的尊重有關.
羅馬書就是說他敬重他的職責.
所以他盡力令到外邦人臨離盡似的信主.
是因為耶穌基督 不是因為保羅.
不是因為他的口才 不是因為他的才學.
不是因為他的身份背景.

$^{561}$是因為救恩這個寶貴的信息.
所以真的是很重要.
我們不要誇耀自己的才學.
自己的身份 自己的所謂拋書包.
本來我們是應該裝備的.
讀多一點書 通透是必須的.
但是不要拿來拋.
第十章這樣說.
他說我保羅 這個是哥林多後書.
開始針對那些假使徒迷惑哥林多教會的.
我保羅就是與你們見面的時候是謙卑的.
我不在你們那裡的時候.
向你們是勇敢的.
如今親自藉著基督的溫柔和平勸你們.
有人以為我是憑著血氣行事.
我也是以為必須用勇敢待這等人.
求你們不要叫我.
在你們那裡的時候有這樣的勇敢.
因為我們雖然在血氣中行事.
卻不憑著血氣爭戰.
我們爭戰的兵器本不是屬血氣的.
乃是在神面前有能力.
可以攻破堅固的刑壘.
將各樣的計謀.
各樣難阻人認識神的那些持高之事.
一概攻破了.
又將人所有的心意奪回.
使他都順服基督.
並且我已經預備好了.
等你們十分順服的時候.
要責罰那一切不順服的人.
大義的意思.
保羅講他的內心.
其實他是很溫柔,很和平的.
但是好像現在因為要跪著哥林多教會.
不要偏向於受那個迷惑,受那個影響.
而甚至乎離開舊恩.
保羅就要表現他的勇敢.
不過是在神面前.
圍著福音而勇敢.

$^{601}$有這樣的權柄的意思.
所以令到哥林多人相當困惑.
為什麼保羅他的言語.
他寫信的時候.
和他實際來到的時候又會有不同呢?.
因為保羅去到哥林多.
又見到本來要悔改的已經改變了.
所以他就可以收回他的那種權柄.
他的威嚴等等.
所以又更加給人一個誤解.
講就這麼偉大,天下無敵.
實際來到做又好像有心無力那樣.
是不是我們這些廣東的俗語呢?.
完全不是的.
保羅就是他可以隨時調整他的內心.
因為他重點是要帶領那些人.
不要偏於異端的邪惡.
反而就是當他們有改變,有迴轉.
保羅的心就立即有一個改變.
這個就是保羅.
你可以說他的那種達人性.
不會好像死在一邊的.
針對那個人就不針對到他死.
人家改變你都不肯接受.
有時候就好像不知不覺.
我們會有這種偏見.
對一個人的偏濕.
不知不覺就看死了.
他怎麼變呢?.
總之在你的眼裡面都總是不夠的.
但是到了第七節就要留意到這裡了.
我想講的就是這裡了.
保羅很想告訴哥林多教會.
他們是怎樣的使徒.
這個是一個氣質.
這個是一個價值觀.
這個是一個影響著保羅做人的一個很重要的價值觀.
所以這些真是很寶貴事奉的態度.
終生的受用的態度.
所以如果我們不是這樣的所謂內功深發.

$^{641}$有時候不知不覺就會有一個受迷惑,受影響.
留意第七節開始了.
你們是不是在眼前.
倘若有人自信是屬基督的.
他要再想想他如何屬基督.
我們也是如何屬基督的.
主賜給我們權柄是要造就你們.
並不是要敗壞你們.
我就算為這權柄稍微誇口也不至於慚愧.
我說者說免得你們以為我寫信是要威嚇你們.
因為有人說他的信又沉重又厲害.
及至見面卻是氣貌不揚言語粗俗的.
所以就是剛才我所說的.
保羅因為知道哥林多教會的人.
有一些離開了真道.
所以保羅知道之後.
他就要寫信去警告這樣的人要悔改,要改變.
但是當他收到.
因為一段長時間才收到信.
或者有人跟他說.
他就真是一收到之後見到人在神面前改變.
他的心就馬上牛軟馬上接納.
所以就被那些挑剔保羅,攻擊保羅的人.
有機可乘就說他的信就是這樣這樣.
見面又另一種了.
就是說原來真是無力無權的這樣的情況.
不過大家留意第七節和第八節.
保羅就好像做了一個初步的分開,分界.
就是將這些假使徒的行徑.
和保羅一幫人他們的真使徒的行徑.
先做一個概括的說法.
接著去到第十章打開和十一章.
就有相當長的保羅說我們作為使徒.
我們不誇任何東西.
只誇我們為主所面對的.
為福音所受的苦.
這些就是他們誇.
就是說保羅不是故意這樣為誇而誇.
而是當時的那些假使徒.
誇耀的就是自己的人脈,財學,資源,金錢等等.

$^{681}$所以就在不斷的以為.
這樣才是真正的使徒.
所以保羅在這裡很想哥林多教會去明白.
就是很用心良苦地去告訴他們.
不過有一句很顯然很重要的.
他們要想一想他們是怎樣屬基督.
我們一幫使徒我們又是如何屬基督.
這句話真是可圈可點.
保羅雖然這裡我就不可以再詳細.
因為我們現在不是看哥林多的後書.
保羅就說他是怎樣屬基督.
保羅在這裡做了一個.
給哥林多人自己去做一個判斷.
他們是怎樣的去屬基督.
而我們一幫的使徒.
當然包括保羅自己.
我們又是怎樣的屬基督.
這個就真是值得我們大家一起去想一想.
我們在一個時代裡面.
教會的牧者是應該被人看到的.
是懦弱,怕事 還是不食人間煙火.
沒有陪伴 只會面對各種的困難.
那些只是平安的訊息.
鬍子就仰處 又不敢指出教會的問題.
各方面等等.
因為我們陪伴著有錢的人.
大家一起開心,平安,快樂.
不需要擔心 神一定會保守我們.
各方面都平安.
因為是怎樣屬基督呢?.
每一個時代,每一個地方都不同.
所以真的要求主讓我們能夠在神的恩典裡面.
我們要看準耶穌基督為我們最高的準繩的標準.
以及我們跟隨使徒們.
他們是怎樣能夠為那個福音的真理.
在那個時代的變質.
來讓更多的人可以聽明白.
聽完不明白,不相信.
我們沒有辦法 我們盡力講明.
但是聽完不僅不明白 還要拒絕,還要攻擊.

$^{721}$這個就是我們要為那個福音講明.
那個福音的代價就是要受苦,受逼迫.
所以每一代人都需要來告訴會眾.
我們是如何屬基督的.
那些人又是怎樣屬基督的.
那些就是開始帶哥林多教會離開真理的假使徒.
所以做了一個分別.
所以牧者需要弟兄姊妹的禱告.
我們是如何屬基督.
這個是一個很大的課題.
既然現在是終身蒙照 成為神的僕人.
我們就要不斷要問自己.
我是如何屬基督呢?.
基督教是如何屬基督呢?.
這些真的是相當重要.
是不是有苦難的時候.
牧者就第一個走了去移民.
我沒有說移民是錯的.
就是走了去離開群洋.
撇下群洋就走了去.
那是不是屬基督的牧者應該這樣呢?.
當然這樣是太過簡單化了.
就不能夠將一些很複雜的問題去梳理.
所以絕對沒有意思說移民就是錯.
離開工場去別的地方是錯 不是這個意思.
而是我是說那個大方向.
牧者是一個時代的牧者是如何屬基督.
這個是最大的課題.
要讓回眾看到我們這一代的牧者.
真的有他的氣節 有他們的堅持.
有他們的價值.
他們真的跟隨福音 傳遞福音.
讓人明白 讓人聽得明白.
讓一個可能完全不識字的婆婆.
都能夠很有系統的講明白.
什麼叫福音 為什麼說自己是信耶穌.
如果都能夠這樣做到.
那就是我們的責任 我們是屬主.
既然神託付我們去傳福音.
我們就要講明那個福音的奧秘.

$^{761}$所以透過這兩段跟大家來想一想.
保羅他很尊重自己外邦人使徒的這個積分.
積分不是一個身份這麼簡單.
他可以不做 不能逼的.
但是他被託付 由信主 接著差遣.
接著神透過亞拿尼亞告訴他.
你要去外邦來傳道.
向君王來傳遞這個恩訊清言.
所以保羅慢慢知道 起初又抗拒.
又覺得自己打爛陣.
看到外邦人已經想殺他了.
現在利用他向外邦人傳道.
這個心路歷程都是保羅自己.
一定是思前想後很多十字架的愛.
很多的眼淚 使他終於豁然開朗.
終於愛上外邦人 向外邦人去傳道.
他的同宮 現在和他一起的同宮.
很多都是外邦人 包括提倭泰.
他的屬靈的遺址.
都是一半有希臘人的身份 那個血統.
所以譬如提倭 肯定是外邦人.
你可以想象到 真的有他對自己的積分的尊重.
尊重是很重要的.
有沒有我們成為神的兒女.
我們在教會的服侍.
很尊重神給我們的身份和積分.
我們很賣力去做那個聖功.
這些是尊重 當我們尊重 敬重我們這個身份.
神都會敬重我們 神都會祝福我們.
所以我們都真的需要在這些價值觀.
這些意識形態 要自己想清楚.
不要只是好像不是稀罕那樣.
做又可以 不做又可以那樣.
一種很敷衍的心態 很敬重.
好 我們看第十一章.
稍微讀一讀經文 對不起 我想要一口水.
我們繼續下去 第十一章.
還有一些時間 我們可以看第十一章.
再看一段經文 我們稍微讀一讀.
十章我們當差不多講完了.

$^{801}$因為有一點我們還沒有講.
我們也不再講了.
普羅他引述的幾段的舊約經文.
稍後我會將整個清單.
羅馬書出現過的舊約.
新約 普羅人的舊約.
一路陸續這樣添加上去的.
第十章很多舊約的經文.
十一章也都不少.
第十一章的第一節.
「我且說神棄絕了他的白成魔 斷乎沒有.
因為我也是以色列人.
亞伯拉罕的後裔 屬辨那敏茲派」.
他是掃羅 辨那敏茲派都是在猶大.
南面的南國 猶大和辨那敏.
都是相當大的一個的茲派.
「辨那敏茲派 神並沒有棄絕.
他如先所知道的白成」.
強調的是他預定 預知道的白成.
即是一個概念在《義昌書》有的.
就是那些叫做愚民 聖愚之民.
就是普羅眼中的愚之的白成.
就是說絕大部份的猶太人都是拒絕尼塞馬.
拒絕福音 但原來神為他自己選了一些愚民.
聖愚之民就會得救了.
所以普羅這裡強調的不是普通的猶太人.
神沒有拒絕 棄絕他自己預先所知道的白成.
當然我們不知道.
所以我們不知道誰有沒有被神揀選.
我們知道的就是那個人聽完福音有沒有信主.
這個就是我們能夠做的.
所以感恩就是神沒有讓我們知道.
誰是被揀選 誰必然會得救.
是沒有人知道的.
所以那個重要性就是.
我們成為教會在地上就是傳福音給萬民.
接著說你們豈不曉得經常輪到以利亞.
是怎麼說的呢?以利亞可以很多東西.
他出來事奉的人生.
由他在基列一個提斯比利亞.

$^{841}$由他出現在阿哈的面前.
在北面的耶穌的地 加密等等的地方.
加密山上面 這樣去跟阿哈.
跟他的假的先知 其實是假神先知.
是真先知 不過是假的 拜偶像的.
拜巴黎 拜亞斯塔羅等等的.
這樣去對抗的.
但是保羅在這裡說的那個被揀選的少數的漁民.
就特別highlight以利亞求死的那一幕.
他不想再做先知.
甚至覺得自己所做的一敗塗地 因為沒有人.
因為他求雨 整個以色列豁然的興奮.
因為那些拜巴黎的都求雨.
甚至用刀刺自己都沒有東西看.
反而以利亞的祭物就被天火燒盡.
十月立了 接著沒多久就有大雨.
就是已經三年六個月沒有下在地上的了.
所以按道理就是以色列人應該很自然的.
找到在我們中間有大先知興起了.
我們現在要跟隨以利亞入城.
去殺了阿哈 殺了這個耶洗別.
接著就擁戴南國的造王.
但是很可惜不是那回事.
這個耶洗別聽到這樣的消息.
就要追殺以利亞.
所以以利亞就很懼怕.
也都真是有一個很大的心靈的落差.
就是所做的沒有成果.
所以他看自己是一無是處的失敗.
以為神猜起他都是浪費.
有時去到那種突然間心靈的疲倦.
那種崩潰 什麼都想得出 什麼都說得出.
所以都要小心 因為太疲倦了以利亞.
就一路走 大家都熟的了.
走到一個羅藤樹下求死.
後來神就向他顯現.
不過留意一下保羅人的這段經文.
第二節的中間那裡.
他在神面前怎樣控告以色列人說.
主啊 他們殺了你的先知.

$^{881}$拆了你的祭壇 只剩下我一個人.
他們說要尋殺我的命.
當然就是耶洗別和宮裡面的敵對勢力.
阿哈只是一個傀儡 什麼都做不了.
又怕自己的皇后等等.
經文就是在列王記上的19章.
就是行完大的神蹟.
殺死了那850個巴黎先知.
但是以色列人表面興奮完之後.
但是耶洗別是無動於衷.
而且是要下追殺令 要殺以利亞.
所以以利亞就在神面前有一點點的訴苦.
或者他內心是一種很大的黑暗軟弱.
但是第四節 神的回話是怎麼說的呢?.
他說 我為自己留下七千人.
是未曾向巴黎屈膝的.
這是19章的18節 就是神親自向以利亞顯現.
讓他知道神仍然在這裡.
你現在心靈的軟弱是絕對可理解的.
因為你的落差太大了.
你所做和你回報是不成正比例.
我想我們人都是這樣.
太看自己的得與失 一次就成敗.
一次的講道就要決定性.
好像不是慢慢的去用苦口婆心去勸導,去教導.
希望一次就要所有人都要悔改.
所有人都要馬上歸向耶和華.
耶洗別馬上要來供奉以利亞等等.
但是不是這樣的時候 就開始忘記了.
以為自己只是你慣狂瀾.
也看不到身邊的同路人跟隨者.
所以又看到因為這個阿哈的牧人.
私下收養的,收留的.
在山洞裡面的先知真是被殺了.
所以在表面那一層.
就是看到以利亞講的是事實.
是真確的一個社會.
現在的靈性,宗教的迫害的開始.
其實一直從來都是開始的.
但是問題是告訴他們.

$^{921}$神要打開以利亞的心靈 讓他們看到.
他們真的未必看到這七千人是哪些人.
其實可能是很普通的以色列人.
可能不是什麼先知.
只不過是散居在不同的地方.
有他們自己的因為信心,信靠神.
而守住自己的那個靈性.
那個不變節的信仰.
未必是一些努力抗衡,洗滅的人.
所以有可能是在皇宮裡面的.
不起眼的牧人士女.
不是以利亞所知道的.
七千不是少 是相當大的數目.
你想想在加密山上面850個先知.
這樣同時間求雨.
又用刀來插自己的儀式.
希望喚醒巴力快點醒.
需要下雨了.
但是以利亞還嘲笑這些先知.
但是當以利亞稍微祈禱.
天就下雨來淋濕了水溝.
然後點火 記住不是石油.
石油根本點都點不著.
石油是要提煉的.
所以不是那些天然的在地裡面的石油.
所以你可以想像得到.
都是一個很大的神蹟.
突顯了以利亞的神是真神 是喚神來的.
但是以利亞自己第二天面對的.
就是一個很大的伺服帶來的挫敗.
其實不應該這樣看是挫敗.
你盡了自己的力向阿哈表明要歸向神.
向以色列人不要再生存疑義 歸向神.
其實就已經完成了自己的工作.
但是要和自己的明星.
和自己得到的讚賞回報.
來掛勾的時候 太多就慘了.
原來巴力不是真神來的.
原來我們一直都拜錯了.
現在你幫我澄清了 以利亞多謝你.

$^{961}$本來這樣的稱讚.
是以利亞都可能很想聽到的.
但是卻不是這樣的情況下.
以利亞就懼怕 甚至擔心自己的安危.
生命的威脅 這樣來逃走.
但是神就和以利亞說.
有七千人是未曾向巴力屈膝.
未曾向巴力親嘴.
所以這裡告訴我們.
保羅引述這一節聖經.
是講到那個揀選的愚民.
神有他那個在救恩裡面的揀選.
所以原來以色列人.
是有一些已經是分別為聖.
被神揀選了的愚民是得救的.
這裡告訴我們.
所以第五節 如今也是這樣.
如今就是說 我現在講救恩.
猶太人是不是真的被拒絕.
這是第一節 不是.
因為保羅說如今都是這樣.
所以我們外邦人就不要學以利亞.
看到很多的猶太人都拒絕.
我們就以為神的恩典.
是不能夠臨到他的亞伯拉罕的子孫.
我們又再忽視了 再忘記了.
第九章已經講了 神揀選了一些人.
來去有信心 這樣來得獎.
第十章也都講了.
當人如果肯很坦誠面對福音.
聽明白福音就可以得救.
所以可能是我們外邦人.
未盡力向猶太人傳道.
我們未盡這個責任.
告訴猶太人尼塞就是耶穌.
所以就導致我們找不到這些.
被神揀選的漁民出來.
所以向猶太人傳道.
又不要認為是天馬行空.
或者是一些在現在這兩年政治的因素.

$^{1001}$我們就恨猶太人恨到入骨入心.
我們更加要告訴他們.
你們是離開了彌賽亞.
是一個非常危險的情況.
但是保祿告訴我們.
神從來沒有放棄猶太人.
照著神的揀選.
第五節 如今也是這樣.
照著揀選的恩典.
原來揀選是恩典來的.
原來在永恆裡面.
神是冥冥當中揀了一些人的.
而這些人是最終可以信主得救的.
所以我們信主.
其實我們是信主得救.
才知道自己是被揀選.
不要搞亂次序.
千萬不要說.
是啊 我出世神揀了我.
十八代之前揀了我.
知道我曾荺貴會信主.
千萬不要搞錯.
是我們明白了救恩.
是有人向我們傳福音.
我們願意相信.
我們現在回頭看.
原來冥冥當中神揀了我們.
是多麼大的福氣.
這個是揀選的恩典.
對於那些拒絕福音的人.
他們很討厭這句話.
揀選 也就是與我無關.
我做錯事也不需要負責.
其實這樣是很得罪神.
揀選的恩典就是說.
還有所留的餘數.
可言就是那個餘民的概念.
第六節 既是出於恩典.
就不在乎行為.
原來揀選的恩典.

$^{1041}$都不是因為他有好的行為揀選.
這個就更加超過了我們的理性.
能夠明白的了.
我們揀一定是找最好的.
但是神的揀選.
既然是恩典.
他的揀選令我們覺得是匪夷所思.
不能夠理解.
但最終我們只能夠讚嘆.
神的偉大作為.
所以不在乎行為.
不然恩典就不是恩典了.
在這個揀選裡面.
保羅仍然都強調.
整個揀選都是恩典.
都是白白的.
所以如果真的在乎有一點的行為.
就不是恩典了.
這個真是非常重要的一個邏輯.
而揀選的背後都是恩典來的.
所以沒有被揀選.
就絕對不能拒絕.
也不能來挑戰神的智慧和偏心.
因為是恩典來的.
意味著我們能夠明白.
舊因是白白.
我們能夠聽福音.
能夠平等.
所以不信耶穌沒有被神揀選.
是真的咎由自取的.
所以這個矛盾要到天上永恆.
才能夠平衡調和.
其實我們自己由信主那一刻.
越信主內.
越發現我們的軟弱.
我們對神的背影.
憤慨感受到.
是恩典托住我們.
是神無條件的愛.
來將我們改變.

$^{1081}$將我們拯救.
所以當我們真是很心願成福.
對舊因的那個讚嘆的時候.
這個就是神無條件的揀選.
這個就是在永恆裡面.
神為什麼揀了我.
而沒有揀另外的人呢?.
這就是一個很寶貴的解釋.
接著繼續下去.
第七節.
這事怎麼樣呢?.
保羅就在這裡有這個解釋.
他說以色列人所求的.
他們會有得著.
他所求的就是.
他們是理當理直氣壯.
是神的選民.
是得著彌賽亞向他們的顯現.
他們帶著彌賽亞建立榮耀的國度.
神就是他們的神.
他們每日每個星期的敬拜.
都是對準耶和華的.
所以以色列人所求的沒有得著.
接著.
為了無揀選的人得著了.
其餘的就成了還敢不化的.
如經上所記.
神給他們昏迷的心.
以眼睛不能看見.
以耳朵不能聽見.
直到今日.
保羅在這一段引述兩節.
《新榮記》的29章和《以塞馬書》的29章.
連續講兩節.
來到講到如果還敢不化.
即是變成石頭.
變成化石.
基督徒的化石是有的.
如果還敢不化.
就是眼睛不能看見.

$^{1121}$耳朵不能聽見.
即是拒絕看見.
拒絕聽見.
拒絕改變.
那就是還敢不化.
仍然認為自己是神選民.
亞伯拉罕的子孫.
是理應,理當,得到神的恩典.
是神的選民.
是沒得改變的.
沒人可以改變的.
這樣的還敢不化.
當然那個原因很簡單.
就是被基督這個活石所判斷.
拒絕基督.
拒絕承認基督.
慢慢就會變成還敢不化.
不僅沒有被揀選.
而且跟福音是越來越.
所謂離學未勒.
所謂是距離越來越遠.
而且是還敢不化.
即是心靈硬化.
心靈硬化就是最嚴重的.
屬靈的那個警況.
即是已經是聽不進任何的真理.
甚至以為自己擁有真理.
這樣的內心.
所以絕對要避免.
絕對要令到心靈柔軟在神面前.
絕對真是要學習.
帶著一份上上感恩的心.
是被揀選的愚民.
當然一直在說的是猶太人.
沒有被棄絕這個上文的.
也是跟第十章拉上一個關係的.
所以我們要記得.
這裡是在說猶太人.
那些還敢不化的大部份.
是讀著拉比的通俄的那些律法書.

$^{1161}$守著六個誡命.
帶著祭物去獻祭.
羅馬書的年代還有獻祭.
你可以繼續想像就是.
這樣的人一直走到聖殿.
是還敢不化.
是拒絕福音.
譬如到耶穌的年代.
在會堂裡面中讀.
他們慢慢慢慢開始很喜歡耶穌.
去蜂擁的去聽他.
但慢慢慢慢見到人喜歡耶穌.
又見到耶穌偶然間很似乏安息日.
那些門徒在安息日就摘些墨水等等.
又在安息日就醫病.
開始就因為宗教上的條例.
還有過不了自己的傳統.
就來針對耶穌了.
挑剔耶穌了.
後來就甚至針對又挑剔.
變成.
所以對不起 又要稍微關一下.
大家不要介意.
如果真的有開了麥克風.
就稍微關一下.
不要影響其他領導姐妹.
不好意思 謝謝大家.
我們又再回去繼續.
這裡告訴我們.
大部份的以色列人是環境不化的.
但仍然告訴我們.
沒有人是注定滅亡的.
救恩的偉大就是.
他們離開世界.
他都可能會改變順處的.
所以我們不要有一種灰心.
去拒絕猶太人的順處.
當然外邦人更加.
雖然沒有這些包袱.
我們更加有其它的東西.

$^{1201}$綁住我們的心靈.
金錢,社會地位.
不同的人世間的哲理.
都會影響外邦人.
拜偶像 或者拜自己.
無神思想的唯物.
這樣都是一些攔阻我們.
對真神的謙卑,降服的.
所以我們是過來人.
我們更加要為人的內心的驕傲的祈禱.
往往不是什麼敵人.
往往是內心的驕傲.
魔鬼的蒙騙.
所以變成環境不化.
猶太人的環境不化.
就是用宗教來成為自己的護身符.
黨節牌.
自己築起的堡壘.
外邦人就是吃喝玩樂.
有今生沒來世.
這樣的價值觀.
只是營營役役地工作,謀生.
享受 忘記了永恆的準備.
好像那個無知財主.
不斷為自己積蓄.
沒有為自己的靈魂來打算.
沒有為永恆來準備.
雖然是一個比喻.
但耶穌說我們要小心.
人的生命不在乎家道豐富.
耶穌告訴我們就是這樣.
就是說原來也挺拒絕福音.
那個價值觀原來.
我們一生只想地上的東西.
就挺拒絕那個福音.
所以這裡就告訴我們.
一直直到今天.
直到保羅寫信的時候.
整個猶太民族是環境不化.
所以都挺痛心.

$^{1241}$在保羅的眼中去看.
更加要使到.
好像我們能夠.
將那些已經被人點了穴的那些解穴.
令到他能夠順處.
所以如果那個講道的人.
例如保羅.
對外邦人有這份熱誠.
如果對自己的猶太人.
都不稀罕.
你信也好 不信也好.
都不關我的事.
如果那個傳道者沒有了那種心靈的焦急.
其實都已經是一個問題.
因為正正就是他已經被人點穴.
成為環境不化.
你要怎樣能夠用方法.
令到他那個環境的化石.
可以來點頭呢.
就是一份很深厚的愛去改變.
所以牧者最可貴的就是.
不斷有一份柔軟的心.
對群羊的愛和牧羊.
如果只是打一份工.
本來打一份工不是錯的 應份的.
但是有時候.
每次都遷就著遷就著.
要有些東西要講.
要表態 又不出聲.
擇不知 明知不好心.
久而久之.
就真的變成打一份工.
沒有了那個真正的心靈的黑.
人家會感受得到的.
所以這裡就讓我們看到.
本來看到現在有大部份的猶太人.
在那個時代環境不化.
他就更加的焦急.
所以第九節.
大衛也說.

$^{1281}$願他們的延籍變為網絡.
變為機構.
變為絆腳石.
作他們的報應.
願理事常關下他們的腰.
這個是詩篇69篇.
彌賽亞的詩篇.
彌賽亞受苦的詩篇.
那裡是講迪基督.
即是猶太.
69篇全篇是講就坐那個人.
所以你可以想像到.
保羅在這裡的69篇.
是講到大部份的猶太人.
現在是背影是環境是不化的.
你可以想像得到.
原來神是整天招呼這些背影頂嘴的百姓的.
第十一節.
可能我們要看多一點就要完了.
未必講得完.
因為都差不多69.
我們看第十一節.
我且說 他們失腳是要他們跌倒.
斷乎不是 反倒因他們的過失.
這個就是第二個段落.
就是在這個table裡面的11到24節.
外邦人得救.
這個所謂的激勵或者激怒.
其實是嫉妒.
是使得猶太人很不值得.
為什麼外邦人會這麼忘恩.
為什麼他們不守猶太教任何的條例.
守任何的日子.
為什麼外邦人這麼忘恩.
或者認為自己這麼忘恩.
要跟我們的信仰的真實拉上關係.
其實是要刺激.
那個英文是要激勵.
其實是激動.
其實是要使得他們嫉妒.

$^{1321}$Jealous.
他們有一種很嫉妒外邦人.
這樣的意思.
他說.
反倒因為猶太人的過失.
就是整體來說.
救恩就臨到外邦人.
要激動.
不是啦.
舊的和諧本就是用這個字.
這個激動就使到他Jealous.
使到他嫉妒.
就是做出一些事情.
好像神特別祝福外邦人.
蒙福.
得著神的恩典.
得著一個生命的改變.
但是猶太人.
繼續停滯不前.
宗教變成未同象立.
這樣的一種恥笑比賬.
就窒暴了.
原來神又用這樣的方法.
去刺激猶太人.
所以呢.
如果我們見到.
教會.
有些弟兄姊妹突然復興.
而我們不抵得呢.
不抵得就大家一起做.
不抵得你又不做呢.
這就叫做被激動.
這是使你Jealous.
使你嫉妒.
這個嫉妒.
其實是使你有一個反思.
不是反感.
所以有時教會就是這樣.
希望拉低她.
希望她不要那麼熱心.

$^{1361}$不會使你的冷淡.
被別人發現.
但原來不是.
神就是要激動猶太人.
引致他們窒暴.
來反思自己.
離開神.
原來是很遠了.
所以這個價值觀.
我們很想這樣想.
原來真的是這樣的.
所以我們為了如果有人.
真是他的靈性.
被神感動興起.
我們應該感恩之餘.
我們應該投入去獻祭.
我們應該投入去協助.
去支持.
去守望.
就不是被激動.
去Jealous.
所以我們要小心.
原來神是會真的這樣.
令到我們.
這個反思是暴露我們自己的.
靈性的.
環境和墮落.
是墮後了.
所以這裡就告訴我們.
是要激動他們.
去發奮.
是開始有一個.
從反思而來的復興.
所以.
原來神要這樣做事.
某程度屬靈的原則.
都是這樣做事的.
如果我們見到一些教會.
整個教會都是很童心.
很單純地希望福恩.

$^{1401}$不覺得辛苦.
覺得很榮譽.
很榮幸.
這個就是一個很大的祝福.
我們在星期四晚.
參加恆在社區的福音.
現在才是.
聖經的基礎.
但是聖經的基礎.
已經是很厚很實.
是非常合乎聖經的.
整體的真理教導.
和九間教會.
其實大部份都是.
成年和長者.
但是見到大家.
是很渴慕.
參與.
改變那個社區.
自己教會附近的社區.
希望大家.
童心去做.
不是希望鬥誰多人數一點.
誰做得厲害一點.
也不是.
大家盡力去做.
原因就是被那個運動激勵.
激勵而發奮.
所以神的作為.
是激勵那個人.
原來Jealous.
可以是正面的.
是發奮.
但是往往的Jealous.
我們因為那個心靈已經冷淡.
那些熱的.
復興.
來到我們身上.
我們就起反感.
起抗拒.

$^{1441}$和會攻擊.
會拖後腿.
希望他不要那麼熱心.
就不要暴露.
我們的罪惡和冷淡.
所以神的作為.
是要透過.
外邦人的蒙恩.
去激動猶太人.
激動他們.
就是.
以色列人.
來去發奮.
來去熱心的復興.
接著.
第十二節.
保羅在這裡說了一個好像對聯那樣.
若他們.
即是.
若他們的過失.
為天下的負足.
他們的缺乏.
為外邦人的負足.
何況他們的豐滿呢.
所以原來神.
不是想猶太人.
就是這樣一沉不起的.
猶太人現在這樣拒絕.
一種誤解彌賽亞.
所以彌賽亞成為.
他們這個民族的絆腳石.
但是.
使得.
這樣的錯誤.
使得天下的外邦人.
來得著負足.
所以他們的缺乏.
為外邦人的負足.
何況將來.
他們這個民族真是.

$^{1481}$當然是說那些被揀選的.
漁民.
來去復興.
豈不是整個宇宙.
回到一個和諧.
合一 敬拜真神.
的那種歡樂.
那種讚頌呢.
何況他們的豐滿呢.
他們就是猶太人.
就是說他們終於.
向神來去.
悔改.
認耶穌基督為主.
他們心裡相信.
耶穌基督是由死.
復活的那一位.
神的兒子.
是他們的彌賽亞救主.
所以何況.
當他們來到去.
接受主之後的.
整個宇宙的.
救恩的圓滿.
所以現在好像還欠一個.
黑數 猶太人.
現在是大部份拒絕.
福音 所以.
他們現在都要面對一個.
國家 民族.
很多的災難仍然.
纏繞著他們.
當然我們不應該將.
政治.
時事放入聖經裡面.
但是原來.
將來他們的豐滿.
而且後面那一段.
今天不夠時間講了 保羅說.
將來.

$^{1521}$以色列就會真正的.
豐滿的時候 整個宇宙.
就復歸於一了 所以.
現在真是只欠猶太人.
未真正聽聞.
耶穌是他們的彌賽亞.
外邦人 當然還有很多.
外邦人 例如阿拉伯世界.
他們沒有機會聽 受各種.
因素限制 但是.
我們現在就是有責任.
要向不同的民族.
宣告 主耶穌就是.
彌賽亞 今天我們就先講到.
第十二節 十三節.
之後我們再繼續看下去.
請我們一起同心禱告.
感恩天父給我們.
在第十一章.
都有機會比較仔細的.
去探討保羅內申.
對猶太人.
信主這一份的感動.
多謝你讓他能夠.
不是要.
刻意的針對.
自己的同胞 反而.
他懷著一份.
很悲天憫人的心.
很想他們的同胞.
能夠早日的歸住.
但是他們都知道.
是有揀選的餘數.
來被串流.
但是大部份都是.
這樣來還銀 不化.
拒絕福音 拒絕耶穌基督.
我們仰望主.
讓我們透過一路看保羅.
他的思想架構的時候.

$^{1561}$更加認識.
更加認定.
福音的寶貴.
福音改變了保羅.
福音也改變了不同時代的.
傳道者 信徒.
讓每一個時代.
我們都能夠去.
見證 捍衛.
我們都一路看下去的時候.
發現原來.
以色列.
他們現在是被質鬥.
他們其實.
是需要發奮.
很盼望神幫助我們.
當我們有時看到一些.
很不喜歡的新聞.
但是都讓我們知道.
這個都是一個過程.
可能我們不是很明白.
背後的格局.
但是求主讓我們能夠為猶太人.
發奮禱告 為猶太人.
他們悔改禱告.
讓他們能夠終有一天.
可以認識彌賽亞.
禱告 完全幫助.
多謝你我們能夠有.
可以這樣的仔細.
去看到羅馬書.
這些平時我們都可能很少.
會去碰的課題.
可能真的和我們未必有太直接的生活關係.
但是那個大邏輯.
都是我們能夠.
要明白 我們是被神.
揀選的神的兒女.
我們現在都滲透在.
不同的教會.

$^{1601}$不同的地區 我們很盼望主.
讓我們看到這一份.
莫大的恩典.
不是我們的行為.
使得我們能夠帶著感恩.
這樣去回饋.
救恩帶給我們的.
那個要求 就是我們努力.
有系統的 同心的.
去傳福音 如果教會.
有不合一的 求主憐憫.
讓弟兄姊妹.
看到哪些才是真正重要的事情.
以致為最重要的事情.
同心的擺上.
我們這樣禱告.
奉耶穌基督的名.
阿們.
\newpage



\section{羅馬書}
\label{sec:bw6FVf5f_n0}
\textbf{2025-10-10 周五晚查經班《羅馬書(24):分題:神國救恩歷史的智慧》 講員:曾裔貴牧師}
\newline
\newline
連結: \href{https://youtube.com/watch?v=bw6FVf5f_n0}{\texttt{https://youtube.com/watch?v=bw6FVf5f\_n0}} ~~~~ 語音日期: 2025-10-11
\newline
\newline
\hyperref[sec:ZdF__6Sl2No]{\small{< < < PREV SERMON < < <}}
~
\hyperref[sec:index_chronic]{\small{[返順時目]}}
~
\hyperref[sec:index_scriptual]{\small{[返順卷目]}}
~
\hyperref[sec:FAt1alHLvdg]{\small{> > > NEXT SERMON > > >}}
\newline
\newline
$^{1}$天父我們感恩再次恢復《羅馬書》的十一章的研讀.
他相當詳盡地去講.
外邦人如何可以透過恩信耶穌.
得到這個清義.
以及在這個清義之後.
如何能夠在我們的舊有生命.
不斷有聖靈的工作.
讓我們能夠脫離我們舊有的自我.
進入神的豐盛生命.
因為是神預先選召我們.
讓我們要在人生學校耶穌基督的榮耀的樣式.
我們過去有相當長的時間.
開始看第九到十一章.
今晚我們會有一個總結.
我們真的很盼望猶太人他們現在的光景.
雖然我們很難用政治.
用他們在過去的兩年對巴勒斯坦,哈馬斯.
或者在加沙地帶的平民百姓的轟炸.
來做一個很終極的判斷.
因為是超過了我們的了解.
但是我們都看到在聖經保羅給我們的啟示.
我們知道有一個很大的奧秘.
盼望今晚我們能夠去梳理一下這個奧秘.
以及以色列這個民族將來的光景.
我們等候在主的面前.
求你一路帶領我們.
很盼望能夠.
雖然有些很難每個星期都能夠聚在一起.
但是當聽回記錄的時候.
又給大家有一個啟發.
可能大家安靜的去聆聽.
以及在各個分段裡面.
去有一個很好的掌握.
我們仰望在主的手中.
這樣感恩禱告.
奉耶穌基督的名 阿們.
我們可以先跟大家溫習一下.
因為上個星期我們停了一次.
我們的演唱會都很感恩.
今晚我給的分題.

$^{41}$是第24講的分題.
相當於是十一章十三節打後.
尤其是最後的保羅的頌歌.
大家又要問一下這個問題.
十一章最後的段落的頌歌.
和第八章的頌歌又有什麼不同呢?.
如果是在這麼短的章節裡面.
出現了兩個保羅很有情感的.
一個對神的重責.
當然是靠文字.
但是我們都應該要有一點分析.
第八章特別強調就是神的愛.
那個永遠不改變.
對信徒的保守.
第十一章尤其是強調神的.
救恩的歷史一直的發展.
神的計劃.
超乎我們人們想到的智慧.
祂的策略.
這個就是十一章比較重點.
對這方面去頌讚神.
這方面的偉大,屬性.
最終也是一切的榮耀.
都要歸給這位的真神.
我們稍微有一點重溫.
上兩個星期我們講過的.
因為這幾章聖經是一個整體.
原因就是以色列人.
為什麼他們不得救.
至少在保羅的時代.
在保羅現在寫羅馬書的時候.
就不是普遍他們對福音有一個認同.
接受等等.
所以保羅就必須要處理一個問題.
為什麼外邦傳福音.
相對是這麼多人信主.
來歸服神.
但是為什麼猶太人.
無論在耶路撒冷的.
在外邦的.

$^{81}$為什麼他們都是這樣拒絕呢?.
羅馬書還不是保羅最後的人生.
都是中段的時間.
特別是已經第二次的宣教旅程.
稍微完結.
就應該在哥林多寫羅馬書.
所以其實已經有相當多的印證.
就是保羅去到外邦宣教.
其實首先是進猶太人的會堂去傳道的.
但是往往拒絕迫迫保羅的.
就是會堂裡面的一些猶太人.
或者一些可能是有拉比背景.
訓練的這些人.
就會對保羅有一些攻擊.
後來當第三次保羅的旅程完結.
上耶路撒冷.
最針對他的一方面就是在安提摩上去.
守節旗的一些猶太教的中堅分子.
但是結集的.
一直要控告.
一直要上訴去到該寨.
來去控告保羅的.
就是猶太一些的拉比.
一些他們安排的專業的律法家.
來去對付保羅的.
所以你可以想像得到.
保羅他現在靜下來.
寫完了外邦人在神的救恩計劃裡面.
來到因為神的愛永遠不改變.
人就可以透過愛神的愛.
當然神的愛就是犧牲他的愛子.
這樣來得救.
而且是永遠無神的保守.
但是他必須要問一個問題.
或者要幫助外邦.
特別是羅馬教會的弟兄姊妹.
都要面對一個很赤裸裸的現實就是.
猶太人到底現在在救恩裡面.
他們處在什麼的光景呢?.
或者神有沒有一些的心意或者計劃呢?.

$^{121}$果然就有了.
所以保羅就在第十一章.
就真的將神這個智慧的計劃.
或者我們叫救恩的歷史的計劃.
來到很詳盡的講出來.
所以如果我們看猶太人的得救.
不要只是看一兩個年代.
更加重要的就是.
我們要在保羅的屬靈的視角.
去看猶太人的將來.
整個民族的結局.
這樣就會健康一些.
安全一些.
因為如果我們只是從時事.
從歷史這樣去看下去.
有時候我們就會離開了.
保羅那個屬靈的視角.
我絕對說不可以.
但是我覺得保險一些就是.
因為我們福音派是堅信聖經的啟示.
還有這個默示.
所以如果由保羅給我們的奧秘.
看到原來將來是一個很亮麗的結局.
我們就有這份的安心.
所以希望大家都從這樣的理解.
這樣去看這三張獨立.
猶太人的過去,現在和將來.
就會相當好了.
所以我們可以稍微溫一溫第九到第十張.
第九張強調的是神的主權.
就是說為什麼說猶太人可以得到神的憐憫,救恩.
為什麼有這麼多的猶太人拒絕呢?.
所以保羅就要強調.
不是每一個猶太人自然的就可以順利.
就是這樣去成為神的兒女.
必須是被神揀選的人.
才可以是神的兒女.
保羅從亞伯拉罕一直說起.
所以他很強調第九張.
神的主權和揀選是不會改變的.

$^{161}$所以必須是在這個前提下去理解猶太人.
所以不是每一個都可以順主.
但是很奇妙.
他第十張就不再說太多揀選.
反而說人對福音的態度,責任.
以及他只需要有一個很簡單的信心.
就是口形深信.
就已經得救了.
第十張是在跟猶太人說的.
不是跟那些初信班,決志信耶穌那些說的.
因為整個脈絡是在說猶太人怎麼可以得救呢?.
原來第十張又強調.
原來不只是揀選就好像兇了一樣.
他沒有揀過我就不信了.
其實不是的.
揀選是在神的永恆裡面.
祂的預知,祂的權柄,祂的主權.
但是第十張特別強調.
就是人要面對福音.
其實應該是面對耶穌基督.
怎麼去回應耶穌基督這個救恩.
這樣打開心靈呢?.
因為耶穌是從天上來到這個世界.
引述的生命的三十章.
摩西跟那些以色列人說.
你告訴我哪裡有道.
我就去尋找.
是不是在天上?.
是不是在海外?.
保羅就引述摩西的申命記那裡.
不是在天上很遙遠的地方.
我們找不到的.
而是基督從天上來到世界.
而且主耶穌從死裡面復活.
所以我們只需要口應心信.
就可以得救.
神是不會偏袒人的.
凡是求告的.
就必定可以得救.
這就是第十章.

$^{201}$使徒保羅強調的.
人的責任.
所以這裡告訴我們.
神對以色列的計劃.
救恩的普世性.
凡求告主名.
十章的十三節.
這些探望大家都要背熟.
凡是求告.
就是不管你是什麼背景.
不管你是什麼民族.
不管你是什麼身份.
你在世界可能是很卑微.
很軟弱.
或者是很顯赫.
都是需要求告主名.
就可以得救了.
所以我們都很渴望能夠發揮.
作為一個傳福音的人.
這是第九章的總結.
神是自由的去揀選人.
可以看看螢幕.
相當大的電腦.
不是憑人的行為.
或者血統.
這是跟猶太人.
第九章特別強調.
不是從以撒辛的.
都是以色列人.
得救不是靠律法.
是憑信心接受.
所以愛邦人得救.
是神計劃的一部分.
第十一章就講得很詳細了.
猶太人的拒絕福音.
並非神的話語落空.
是猶太人自己拒絕神的救恩.
所以記住.
揀選主權在神那裡.
我們在地上的人是不知道.

$^{241}$誰有被揀選.
誰沒有被揀選.
我們是不知道的.
我們知道的就是.
向他講一個福音的訊息.
他聽明白.
他的心受感動.
他求告.
他就可以得救了.
所以我們要知道的就是.
他有沒有將自己的心靈打開.
他自己的責任去相信耶穌.
所以他不知道有沒有被揀選.
我們也不知道.
他的對象有沒有被揀選.
所以我們不要以為他現在信得很開心.
我們就一定以為他是揀選.
可能他之後就會離開的.
但也不要因為有些人表面很抗拒.
但注定看死他.
一定沒有被揀選.
這就不是我們那個層面去看.
我們的層面就是記得.
既然信心一定是來自聽到.
所以一定是需要有傳到.
因為沒有傳到就不能夠讓人去聽到.
所以是需要這樣的發生的.
是有一個這樣的相互的關係的.
所以教會的存在就是將這個福音去傳.
當然用什麼方法.
用什麼技巧.
用什麼方式.
那是不同的教會.
不同的時代的責任.
我們必須在不同的時代.
要去把握現在這個時代.
現在這個群體.
他們怎樣去聽福音.
所以有時候用電話.
現在這個時代.

$^{281}$都是一個傳福音的媒介.
但是一樣東西很重要的就是.
一定要聽明白什麼叫做信耶穌.
他一定要自己心靈打開.
求告主名.
那就可以得救了.
使得人求告主名.
十章就很強調.
就是教會的責任.
報福音.
傳喜訊的人.
他們的腳中何等皆美.
所以我們服侍主傳福音的人.
腳中皆美.
去到將這個寶貴的福音.
去跟他們說.
所以這是一個很重要的信息.
好像在這個星期.
因為張董介紹我認識一個家庭.
是很大的困擾.
那個家庭的心靈.
那個媽媽的心靈的需要.
很緊急.
要我立刻去趕鬼.
了解之後.
其實一定不是靈界的問題.
是心靈落在一個很大的困境.
感恩今天他肯去跟女兒.
女兒在錦繡工作的.
去找廖醫生.
得到很大的安慰幫助.
多謝教會.
這些都是一個渠道.
讓人聽明福音.
可能有很多東西.
我們需要有時.
可能先在他的生活上.
先在他的需要上幫助他.
對教會有好感.
慢慢可以來聽道.

$^{321}$所以最終也不是因為他喜歡教會的音樂.
喜歡教會的朋友.
喜歡團契.
大家吃飯.
那些全部可能是一個準備.
最終必須是.
信心一定是來自聽明白神的道.
他必須是聽而相信耶穌.
然後他就求告.
跟他傳福音的人.
一起去禱告.
他去決志承認耶穌是他的主.
這樣就得救了.
所以這是我們的責任.
教會的責任就是要傳揚這個真道.
這個三節大家都應該熟悉的.
凡是求告主明的就必定得救.
所以救恩是向所有的人開啟.
你可以想像到.
如果你要去阿拉伯世界.
你要懂阿拉伯文.
你要可能用那些技巧.
不同的方法.
這個就是我們必須要在每個時代.
都有每個時代的傳福音的方式.
例如現在有些營商的宣教.
就去到那個地方.
表面就不是一個福音堂.
不是一個教會.
但是透過辦一些的商業.
大家的結識來往.
慢慢建立友誼關係.
就可以帶他順處.
所以有這些迂迴.
但是問題是那個重點.
就是要讓每一個人聽到福音耶穌的救贖.
這些外邦人一重溫.
我們看這個第十一章.
保羅一開始第十一章.
就說了一個很重要的題目.

$^{361}$就是其實最終神有沒有放棄以色列人呢?.
保羅很清楚說是沒有的.
一到第十節就已經幫大家做了一些分段.
是沒有放棄以色列人.
但是現在以色列人好像被撇棄.
其實是他們自己拒絕耶穌.
拒絕承認耶穌是他們的彌賽亞.
一直這樣來對號入座.
彌賽亞應該是很榮耀的君王.
但是耶穌一點都不榮耀.
不是君王.
只是一個死囚.
所以他們這樣對號入座.
就拒絕承認耶穌是他們的君王.
是他們的彌賽亞.
所以他們就被撇棄.
在這個時候.
激勵這個詞十一到二十四節.
其實是嫉妒.
是嫉妒這個詞.
所以是讓猶太人感受到那個嫉妒.
為什麼外邦人可以這麼無恩.
為什麼我們猶太人每個星期去會堂裡面.
敬拜保羅的年代還是有聖殿的.
在耶路撒冷有聖殿的.
他們去到聖殿敬拜神.
為什麼沒有外邦人這樣的福氣蒙恩呢?.
當然他們不承認外邦人的恩典是神的恩典.
只是一群人在奉耶穌的名聚會.
說敬拜我們的真神.
其實是很諷刺的.
但是保羅就真的以一個過來人.
他是由那個黑暗猶太教跳出來.
現在能夠跟外邦人有這樣的來往.
他是很深的感受到.
外邦人是得著一個屬天的平安.
永恆的救恩的那個福樂的.
所以保羅他看外邦人.
是真的蒙了神的恩典的.
所以他就在一個這樣的淡水.

$^{401}$他就很失痛很悲哀.
猶太人現在在一種這樣的質度的光景.
來有一種敵視保羅.
敵視這些外邦的教會.
所以這個其實是醞釀12章打後.
去到15章.
保羅要處理在教會裡面的那個律理關係.
尤其是大家不同民族.
大家不同的文化.
大家來自有不同的背景的信仰.
共野一爐.
如果不處理好那個關係.
就有很多的衝突.
所以13章開始就是處理這些衝突.
你可以想像到.
現在外邦人就蒙這麼大的恩典.
所以原來是神的計劃.
去逼到猶太人要對自己的信仰有一個反思.
這個就是從那個質度開始.
去反思到底我們的民族怎麼辦呢?.
如果這樣想的時候.
他們就會近乎得救了.
因為他們真的要問.
如果不是我們的祖宗.
如果不是我們的血統.
如果不是我們宗教的行為.
還有什麼呢?.
唯有信耶穌.
真的要降服在耶穌的救恩裡面.
所以這裡就是保羅一直的思路.
接著我們要講這個.
我先跟大家讀經文.
11章這一段.
也是有那個比喻.
保羅也很巧妙地用橄欖樹這個比喻.
我跟大家稍後會引述耶利米蘇.
有一段經文也很重要.
11章之前有讀過.
我們可能從11節這個段落讀起來.
第11節.

$^{441}$「我且說他們」就是猶太人.
「他們失腳是要他們跌倒嗎?.
斷乎不是」.
所以已經保羅很清楚.
不是一個終結.
不是說他們已經永不超生了.
現在是暫時的失腳.
是暫時的被撇棄.
是暫時他們因為這樣的這個疾度.
他們就要來不斷的品行.
外邦人可以無恩.
但是他們自己就一直這樣來到去.
在那個救恩的以外.
來到去有這種疾度.
接著留意一下.
「盡乎不是.
反倒因為他們的過失.
救恩便臨到外邦人.
要激動他們發奮.
若他們的過失為天下的富足.
他們的缺乏為外邦人的富足.
何況他們的豐滿」.
這個所謂豐滿就是.
他們將來真的整個民族.
來迎接彌賽亞的那種雀躍歡欣.
我們外邦人就沒有這種強烈的感受.
因為我們一個個都是這樣相信耶穌.
認定耶穌是主.
你不認就不會去教會不信.
所以每一個都是這樣認.
但是猶太人呢.
他們是全家的.
如果他們整個民族.
一起宣認耶穌是救主.
那就是整個宇宙的一個豐滿.
就是一個很完整的救恩的成就.
所以保羅用的字眼都很考究的.
現在他們的過失.
現在他們的缺乏.
就令到外邦人天下都富足.

$^{481}$何況將來當猶太人整個民族歸主.
承認耶穌就是彌賽亞.
那你可以想像到.
那天的日子的那種奇妙美麗.
所以現在有一個猶太人信主.
都是大件事的.
當然他要面對.
如果他在以色列.
整個家族是要跟他出殯.
當他死了.
猶太人.
現在是這樣的.
在外國 比如英國.
那些相對沒有那麼嚴重.
但是猶太人信主.
一個已經那麼嚴重了.
但是將來如果他全部.
整個民族一起歸向耶穌.
承認他是真真正正的彌賽亞.
保羅看是一個豐滿.
整個宇宙的一個圓滿的救贖.
所以我們要在這個大前提.
這樣去理解猶太人.
所以現在保羅完全不適合政治的.
不適合現在猶太人那種惡行的.
最近 應該是近年.
有一位神學家.
就問在猶太 在以色列的一個拉比.
現在這樣的情況.
怎樣呢?.
猶太人現在這樣來到去狂轟濫襲哈馬斯.
巴勒斯坦的加沙地帶的人.
死了超過百萬的難民.
這樣怎麼辦呢?.
超過五萬人死.
其實很多在亞利根本找不到的.
那怎麼辦呢?.
那個拉比流著眼淚這樣說.
最近的 近期的.
即是說可能幾年前.

$^{521}$就說可能我們猶太人.
需要再一次被擄多一次.
就是給神再一次被擄.
這樣流著眼淚跟神學家說.
基督教的神學家.
就真的在以色列問他.
這個拉比這樣回答.
相當智慧的智者.
承擔這個罪要面對.
所以可能還有另一次的.
大規模的以色列被擄.
再一次好像被拋到天下.
都不出奇的.
因為現在不同的國家.
歐洲的都承認巴勒斯坦要立國.
所以可能再要快一點的還打.
可能就更加不會幫以色列了.
所以今天就簽訂了第一階段的和約.
就不再進侵,戰爭等等.
就是說現在要盡快釋放那些俘虜.
但是我們看到的就是.
將來保羅有一個很奇妙的預言.
應該說是一個奧秘.
就是以色列將來整個民族的歸主的豐滿.
第十三節.
接著就回來跟我們說.
所以你會看到這個過程.
現在猶太人是這樣被撇棄.
他們現在處於一個質度的裡面.
他們自己以為自己享受著的.
是與生俱來的民族的血統.
那個血緣,那個福氣.
但是原來保羅說其實已經停了.
你們現在是沒有的.
救恩現在是給外邦人的.
當然不是說個別不可以有.
保羅自己就是說他這麼背影.
他都信了主了.
他都能夠被警察向外邦人傳道.
所以現在保羅就跟那幫外邦人說.

$^{561}$留意了 跟外邦人就是一個集體的說.
當然不是某個時代.
因為他們根本就不認識我們.
他們都不知道香港.
他們不知道2000年後有一個叫中國.
因為那個時候的地理根本不會知道我們有.
在遠東原來羅馬帝國的時代.
就是已經是漢朝的時候.
他們怎麼會知道有一個這麼大的地方.
有這麼多人在那裡聚集呢?.
第十三節.
我對你們外邦人說者說.
因我是外邦人的使徒.
上個禮拜我們就說到這裡.
所以敬重我的職份.
他是完全沒有任何的對外邦人的.
那種以往的歧視.
那種以往的看不起.
那種的優越感.
一個使徒能夠所謂洗盡原話.
就是完全沒有.
這些對那些風塵裡恥的說法.
就是保羅是完全沒有對其他的民族看不起.
我們都覺得有些汗顏.
現在可能好一些.
我們都不敢再得失那些內地來的自由行.
就是希望他們花錢.
希望他們不要窮遊.
都不敢踢行李箱.
因為經濟不好.
如果我們風光的時候.
又踢人家的行李箱.
又走去笑人家等等.
就是個別的行為.
當然我們不應該鼓勵.
諸如此類.
但是保羅說他敬重他自己的職份.
就是一個使徒成熟到這樣.
他知道自己的民族現在是被神撇棄.
但是他盡心盡意的去服侍外邦人.

$^{601}$為外邦人禱告.
盡力帶他們信主.
讓外邦人蒙受現在神的大恩.
所以作為牧者.
我想我們都是要反思的.
我們有沒有對某一類人特別喜歡呢?.
因為簡單的,聽話的,乖的.
我們就好一些.
一些不聽話的.
我們就差一些.
有時候這樣的差別對待.
其實都不好的.
更何況其他的民族.
保羅就是這樣的敬重他的職份.
榮耀他的職份.
接著保羅反而就是說.
或者可以激動我骨肉之親法群.
第二次說法群.
真的要有一個激動他們的心靈.
要來到有一個很大的反省.
信仰原來真的要不住的反省.
猶太人在保羅的年代.
他們還好像在吃血統的紅利.
還不斷以為自己.
亞伯拉罕的子孫.
為什麼我們每個星期六.
安息日.
就是星期五的六點黃昏.
到第二天的星期六晚的黃昏.
才叫做安息日.
因為他們算黃昏開始算起的.
每個星期去安息日.
去聖殿.
我們不得到神的救恩.
你怎麼說都說不完.
說不完就是不對數.
但是保羅很想讓他們知道.
原來其實神沒有改變他的計劃.
是人自己接不上神的心意計劃.
所以這些都是我們要不住.

$^{641}$無論個人或者教會.
我們都要一起反省.
否則我們一直以為.
自己在吃那些信仰的紅利.
其實就很嚴重.
所以記住.
這裡告訴我們.
十四節.
可以激動我骨肉之親法群.
好救他們一些人.
保羅的年間.
救一個就一個.
誰拉起肯信主.
就得到神的救贖.
所以你會看到保羅的同行的同工.
很少猶太人.
很少猶太的年輕人.
所以提摩那一代都很易數.
他是猶太裔.
但是希臘人.
他爸爸是希臘人.
他媽媽是羅爾和有利基.
是猶太的勁權的婦女.
所以不多猶太人.
這樣來跟隨保羅.
但是都有的.
不過你告訴我們.
第十五節.
若他們被丟棄.
猶太人.
天下就得與神和好.
他們被收納.
豈不是死而復生.
接著再用另一個比喻.
所以的新棉.
新棉不是那些公仔棉.
是棉粉.
很幼細的棉粉.
所以上的新棉.
若是聖結.

$^{681}$全團也就聖結了.
樹根若是聖結.
樹枝也就聖結了.
這裡用了一個概念.
跟大家讀一讀民數記.
如果沒有搞錯.
應該是19章還是18章.
這段經文.
獻的棉 新棉.
因為那個是關於獻祭的問題.
我看看下一個部分.
有沒有記錄下來.
好像有記錄下來.
大家先看一下.
我看看另外一個記錄了.
不要緊.
跟大家記錄這段經文.
給我一點時間找一找.
可能記錄在另外一邊.
關於獻祭的問題.
這個舉掣.
很快就有了.
跟大家記錄.
15章的17節這段.
民數記15章17節.
耶和華對摩西說.
你曉如以色列人說.
你們到了我所領你們進去的那地.
那地的糧食.
就要把舉掣獻給耶和華.
舉掣的意思.
定義就是.
有一大塊的收成.
先拿其中一個小部分.
那個部分就變成一個舉掣.
將它變成一個.
一堆的東西.
例如有水果.
就將.
一籃的水果.

$^{721}$那一籃的水果就叫舉掣.
整個收成的地方.
可以是很大的果園.
果園裡面那一籃的水果.
就叫舉掣.
這裡特別強調什麼呢.
留意一下.
第19節.
就要把舉掣獻給耶和華.
你們要用初熟的物質磨麵.
最幼細的小麥.
來磨成麵粉.
那時候沒有磨的.
可能用木棍等等.
做餅.
當舉掣奉獻.
就是那個小麥的餅.
因為用這個幼細的麵粉.
最幼細的麵粉.
當然加些油.
來搓成一個餅.
這個餅就成為.
舉掣獻給神.
所以就代表著.
這個大的豐收.
小麥的大豐收.
這個餅就象徵.
聖潔的獻給神的一個舉掣.
自己未吃之先.
要先獻了這個舉掣.
所以為什麼神打工剛順著呢.
第一個月的出糧.
就奉獻給神.
那個新鮮的一個舉掣.
就是說神一路祝福你.
以後的工作.
第一份工第一份第一個月的出糧.
就奉獻給神.
當然那個月的十份一.
是一個很重要的象徵.

$^{761}$這個象徵就是說.
所有的農作物的收成.
這個舉掣就是再一次將榮耀歸還給神.
如果不是神祝福.
根本上是不能夠有收成的.
所以舉掣這個動議.
保羅就再一次在十一章那裡用了.
我先說完.
你們舉上.
好像舉和場的舉掣一樣.
你們世世代代.
要用初熟的物質磨面.
當舉掣獻給耶和華.
所以.
如果我們真的有這一份感恩的心.
記得.
一切都是來自神的恩典和祝福.
所以我們.
每樣東西.
帶著這份感恩的心的舉掣.
來將榮耀歸還給神.
是神所配得的.
所以保羅這裡就有一個概念.
我們要留意這裡.
十一章這裡就提到.
因為上下文說的就是.
現在是被撇棄.
所以天下.
即是外邦人.
就能夠與神和好.
將來如果他們被收納.
豈不是死而復生.
這個民族的死而復生嗎?.
這是十一章的十五節.
接著第十六節.
所獻上的新面.
若是聖潔.
全團.
也就聖潔了.
又是說一小撮.

$^{801}$和一個整個和場.
這樣來做比較.
這裡所獻的新面.
不是說物質的面粉.
不是小麥的餅.
不是那個舉掣.
因為這裡要與下文再讀在一起.
就是在說.
這裡說.
樹根若是聖潔.
樹枝也就聖潔了.
說那個舉掣的小麥餅.
和一棵樹的樹根.
是聖潔.
即是根源的部分.
是聖潔.
長大的樹枝也是聖潔.
到底這個小麥餅.
和這個樹根.
是代表什麼呢?.
如果你看下面的.
橄欖樹的比喻.
就相當清楚.
就是在古代.
神揀選.
亞伯蘭.
離開吾爾.
成為一個民族.
就是以色列.
即是說亞伯拉罕,以撒,雅各.
約瑟.
這一班人.
一直都是神揀選的器皿.
去到摩西終於成為了.
在西賴山.
神頒布十誡.
這個民族永遠歸神.
所以這裡說的.
心面.
這裡說的樹根.

$^{841}$是說蒙神揀選的.
舊約的.
這些聖徒.
這些是有信心.
與神立約.
或者說應該是神.
和他們立約的.
這班有信心的.
領袖列祖.
所以現在就是.
保羅稍後就會說.
根源的問題.
所以現在這裡先有一個比喻.
就是說為什麼這個民族.
是沒有被撇棄呢.
暫時撇棄.
不是終極的撇棄.
因為原來神的計劃.
是從來都沒有停止.
要祝福給以色列人.
原來呢.
由開始亞伯蘭.
由他的兒子以撒.
一直到雅各,約瑟.
然後去到其他的有信心的先祖.
所以難怪.
創世記.
用這樣的敘述的技巧.
將列祖.
古代的這些信心的人.
他們的人生的事跡.
盡量最重要的部分.
記下了.
所以創世記都是四個古人的故事.
就是亞伯蘭.
當然有阿當.
但是記的很少.
主要是由亞伯蘭開始.
就很詳細地記載.
什麼時候結婚.

$^{881}$去哪裡.
在希伯倫又怎樣.
然後去埃及又營祭作會.
很多生活的事情.
去記載.
原來就是很想告訴我們.
這個民族.
原來他們的先祖.
是有信心跟隨神.
不是他們有好行為.
是神令到他們.
有信心.
得到神一直的眷顧.
一直這樣來牧養.
所以他們過去了.
他們的人生之後.
其實現在就成為.
位於天上的信心.
的行列的人.
不是偉人.
是有信心而得到.
神的救恩的人.
所以保羅在這裡.
他現在內心寫到這裡的時候.
他感受到一份.
很大的恩典.
就是原來.
為什麼我們這個民族沒有被撇棄.
原來神.
一直跟我們的.
列祖立約.
是神沒有遺棄他的約.
不是人.
盡了自己的本分.
所以這裡說的新面.
跟樹根.
其實就是說亞伯拉罕.
跟亞伯拉罕一樣.
信心的人.
他們就是這樣在神的.

$^{921}$國度裡面.
而在他們那個時代.
短短的人生.
來經歷和跟隨.
這一位信實的主.
和這一位榮耀的主.
所以我們又要.
回頭看回我們自己.
雖然我們不是猶太人.
我們也不是肉身的血統.
但是加拿大書保羅.
也提到我們是.
屬靈的子孫.
我們是有信心的.
亞伯拉罕的子孫.
這樣說的保羅.
加拿大書有三章.
可以翻譯一下.
所以原來那個樹根.
和那個新面.
都跟我們有關.
所以原來.
其實這個真的可以直接.
形容敬虔的.
家庭的後代.
如果一代一代的敬虔.
真的感受到.
神的一路的祝福.
所以我們真的有生之年.
真的要為子女祈禱.
為子女的將來.
你見不到的了.
因為你會離開.
你還未出生.
你要為他們祈禱.
為那些未出現的都祈禱.
一路去到三代四代.
我之前也說過.
偶然間聽到.
港福堂.

$^{961}$有一個相當大的家族.
那個家長還在.
然後那些子孫上台講見證.
就是說.
自己的爺爺媽媽是怎樣敬虔.
他們是這樣家族.
一路傳承.
幾十人這樣坐在那裡講見證.
其實是一個很大的福氣.
所以現在保羅告訴我們.
其實不是保羅.
原來是亞伯拉罕.
原來是這些列祖他們對神有信心.
神的應許.
神的救恩計劃是不改變的.
所以現在保羅能夠.
這樣寫這些關於.
我們這個民族將來怎樣.
將來就一定是良例.
在保羅的奧秘裡面.
所以當我們回頭看.
剛才還沒讀完.
關於太師這一段.
三章.
如果沒有記錯應該是十九節.
這些突然間這樣閃過來的經文.
講到這些.
其實你就要有一個聯想.
就是因為多用.
我才能夠馬上聯想到.
因為都在講血統.
講列祖的信心.
保羅在這裡.
很特別強調.
我們雖然不是血統上.
是亞伯拉罕的子孫.
但原來在信心的層面.
我們就是亞伯拉罕的.
屬靈的子孫.
即是說我們都可以得到.

$^{1001}$亞伯拉罕所得到的.
應許的祝福.
三章.
我看看要不要讀多一點.
三章可能讀多一點.
第一節開始.
三章的第一節.
加拉太書.
無知的加拉太人.
耶穌基督經十字架.
已經活化在你們眼前.
誰又迷惑了你們呢.
我只要問你們一件事.
你們受聖靈.
是否因為你們行律法呢.
還是你們聽信福音呢.
當然就是聽信福音.
第三節.
既然你們都是靠聖靈入門.
即是說.
得到聖靈在內心的內著.
為何現在還要靠肉身去承傳呢.
即是說.
聖靈在我們裡面.
就算我現在未能夠成熟.
成為好像耶穌的基督徒.
但都不是要反過來.
靠行為去換取功德.
那個意思就是這樣.
所以.
未完全成熟.
不等於不是一個基督徒.
的完整.
因為一個洋娃娃.
她永遠都不會說話.
但一個剛出生的嬰兒.
就算現在不會說話.
她一定將來就會說話.
因為她是一個人的生命.
所以.

$^{1041}$我們要記得古來這裡的.
那種很激氣的一種責備.
他們現在好像.
寫本續末.
去跟那些人說.
你們要行律法.
就是受國禮.
或者其他的規條.
你們是不是這麼無知.
第四節.
你們受這麼多的苦.
即是遺福恩.
是不是徒然.
連續這樣說.
賜給你們聖靈.
在你們中間行義能.
是不是因為你們行律法.
即是他們經歷了神蹟奇事.
經歷聖靈的牧養.
其實是因為什麼呢.
當然再一次說就是聽信福音.
即是原來福音.
我們聽明白.
不斷去思考耶穌.
和你的生命關係.
可以經驗聖靈.
還可以經驗聖靈的工作.
所以我們牧者.
如果不斷去跟別人說.
信耶穌去為人祈禱.
有時有些出人意表的事.
發生我們都不知道.
因為就是這樣.
如果神要透過我去祈禱.
去醫好他 是馬上可以的.
所以我們.
如果我們不去運用聖靈.
不去倚靠聖靈.
不去思想耶穌的大愛.
我們就會枯乾.

$^{1081}$我們就會慢慢覺得.
耶穌都不是值得我們.
這樣去讚美.
因為這裡就是很想.
告訴我們.
他們其實根本不需要再做其他事.
只需要在聖靈裡面.
去思想神的大愛.
第六節.
正如亞伯拉罕信神.
這就算為他的義.
這個就是十五章的第六節.
亞伯拉罕當時.
因為沒有兒子.
神告訴他.
來看看天上的星 海裡的沙.
亞伯蘭就信.
神說你將來的子孫.
就好像這麼多.
亞伯蘭就信.
神就以此為他的義.
不留意第七節.
所以你們要知道.
乃以信為本的人.
就是亞伯拉罕的子孫.
並且聖經既然已經.
預先看明.
神要叫外邦人因信稱義.
就早已傳福音.
給亞伯拉罕說.
萬國都必因你得福.
可見乃以信為本的人.
和有信心的亞伯拉罕.
一同得福.
記得了.
以信為本.
我們一生要信主.
為我們的本分.
為我們的本業.
為我們的本源.

$^{1121}$就能夠和有信心的亞伯拉罕.
一同得福.
而且不是地圖的福.
而是永恆的.
屬天的.
屬靈的福樂.
這是相當重要的.
所以記住.
我們福氣的來源.
必須來自.
有信心的.
和神的關係.
就是和亞伯拉罕一樣.
所以在這六個層面.
我們就是亞伯拉罕的子孫.
那裡我不再詳細說了.
後面那裡.
我本來要再揭的.
不揭了就回到羅馬書這裡.
所以這裡告訴我們.
那個本源的問題.
所以如果由我們開始信主.
我們要切切在我們有幸的人生.
將最好的屬靈的福氣.
就是禱告.
為下一代禱告.
要他們愛主.
跟隨主.
我們見不到第三代不要緊.
我們為第三代未出現的祈禱.
希望他們將來.
效法我們的先祖.
這樣的信心.
所以這裡就是保羅告訴我們.
因為新年是聖節.
新年是聖節.
因為新年是聖節.
所以全團都聖節.
因為樹根是聖節.
所以樹枝都聖節.

$^{1161}$保羅這樣來描述.
接著第十七節.
開始說這個比喻.
如果有幾根.
外邦人是野橄欖枝.
猶太人.
就是這棵橄欖樹.
原原本本的橄欖樹.
要揭經文先.
後面的經文就是.
耶利米書的十一章.
沒錯耶利米書的十一章.
相當重要.
差不多就是.
為何保羅.
他不用其他樹.
就用橄欖樹呢.
十一章的耶利米書.
都是耶利米的哀求.
十一章.
又是要讀長一點.
十四節開始.
所以你不要.
為這百姓祈禱.
不要為他們呼求禱告.
因為他們遭難.
向我哀求的時候.
我必定不應允.
已經去到.
將他們被擄都不行了.
就算他們祈禱都不能改變.
神定了.
因為在早期的.
先知的時代.
他們祈禱.
悔改神的刑罰.
就不會臨到.
但現在去到耶利米.
最末後的時候.
神的憤怒已經去到極處.

$^{1201}$所以跟耶利米說.
你不要祈禱.
祈禱都沒有用.
我都要將這班人被擄到巴比倫.
所以你就可以想像到.
已經是.
完全被神.
好像.
被遺棄了.
這個孕憑.
相當嚴重.
我記得中二的時候.
英文科的老師.
那時候很貪玩.
很頑皮.
那個老師.
跟我哥哥很熟.
教我哥哥.
我哥哥是高班中四還是中五.
因為我哥哥.
都是一些荒頭等.
因為他是校隊.
足球隊.
拿很多獎.
特別給學校看重.
他們不用上課.
有很多特殊優惠.
隨時可以去踢球練球.
很多老師.
都認識我哥哥.
叫大富.
因為他踢後裔.
後來認識我.
知道我是富的弟弟.
對我都很客氣.
但我很賭蛋.
很頑皮.
老師.
對不起.
那時候.

$^{1241}$我哪會怕.
我不理你.
任憑你.
其實這樣放任.
才是最大的悲哀.
因為佛下還是有一個可以悔改的機會.
所以耶利米蘇.
這裡是說他們已經不能再悔改.
都不會再收了.
聽他們禱告.
我偏不應允.
我所親愛的.
既行許多淫亂.
即是他們拜偶像拜得很嚴重.
性慾也來了你.
即是獻那些濟生.
都不會得到神的願立.
你在我念中作甚麼呢.
你作惡就起立.
從前耶和華給你.
起名叫青鑒冷樹.
又華美又潔好果子.
如今他用洪陽之聲.
即是那些軍隊.
要來攻擊耶路撒冷.
點火的在旗上.
之子也被截斷.
原來栽培你的萬軍之耶和華.
已經說要降和攻擊你.
是因以色列家和猶大家行惡.
向巴黎燒香.
惹我發怒.
是自然而然.
是因為我.
是天下的天子.
我才會被人家.
惹我發怒.
是自作自受.
這裡就很清楚.
原先這個民族是親自.

$^{1281}$設立成為一個.
鑒冷樹.
可以來出鑒冷.
鑒冷可以榨油.
又可以來滋潤.
但是.
不能夠.
來成為一個.
華美又潔好果子.
的一個民族.
所以神就要.
來將他們.
被擄到他方.
是一個很悲慘的.
刑罰.
但這裡告訴我們.
原來鑒冷樹是整個民族.
的一個象徵.
是神親自設立.
這個民族.
成為一個.
可以潔果子.
來榨油.
吃或者.
熔抹.
特別是中東那些乾燥的地方.
需要油來滋潤.
所以.
你可以想像到.
現在他們成為了野的鑒冷.
成為了要被神.
懲罰的鑒冷.
我們回到王亞宣這裡.
我們看最後這一段聖經.
他說若有幾根枝子.
被截下來.
所以不是整棵樹斬了.
是將那棵樹的.
枝子被截下來.
你這個野鑒冷.

$^{1321}$得摺在其中.
一同得著.
鑒冷根的肥汁.
然後你就.
不可向舊的枝子.
即是被截下來的枝子.
來誇口.
若是誇口.
當知道.
不善理托著根.
可能我們有一些.
恩怨.
蒙神的祝福.
我們就不要看不起那些被神截肢的.
無論猶太人.
或者其他.
已經再對耶穌.
沒有興趣的人.
我們都不要這樣.
因為保羅強調一件事.
一切都是恩典.
一切都是神.
在每一個人的生命裡面.
的那個作為.
所以.
是根托著那些樹枝.
是那棵樹的.
本源.
去托著那些樹枝.
能夠生長.
能夠有那個質象.
能夠茂盛.
能夠最終是結果子.
等等.
都是來自那棵樹.
沒有被枯死.
那棵樹仍然是神自己親自栽種.
這個是很重要.
告訴我們.
我們要知道.

$^{1361}$我們是神的肥質.
所以我們就不可以向舊知子誇口.
若是誇口.
當知道不是你托著根.
那是根托著你.
保羅一直用一些很擬人法的說法.
首先他已經說了.
這個樹根.
他就不是特指其他的樹.
是橄欖樹.
特指本來是神親自的設立.
這個民族.
是應該出果子.
是高柔.
成為很多人的享用.
那些橄欖的果實.
又可以吃.
但是現在.
他們被截下來的.
團支.
猶太人.
仍然在神的舊恩裡面.
因為說那個樹根.
就是說那個舊恩.
一直由亞伯拉罕開始.
這個救贖.
其實神一直在栽培.
從來沒有撇棄.
地球上的人.
尤其是猶太人.
所以我們.
記得信主.
我們都應該將猶太人.
放在我們的禱告裡面.
因為現在照著保羅的視野.
神給他的.
宿命的奧秘的眼光.
猶太人.
仍然在.
舊恩裡面.

$^{1401}$不過現在在這個本樹裡面.
這個橄欖樹裡面.
神的舊恩計劃裡面.
是被截下來的.
團支.
所以這個團支.
現在就沒有生命了.
因為斷了 離開了本源 沒有生命了.
現在神很奇妙地.
這是保羅的比喻.
很有智慧.
又很有想像力.
保羅.
我們這些是野橄欖枝.
摺上去.
原來橄欖可以這樣.
來生的.
所以沒有違反天然的.
植物學的.
智慧.
和常識.
我們是野橄欖枝.
被摺上本樹.
所以我們就沒有這麼大的恩典.
為什麼我們有生命力.
為什麼我們.
外邦人會知道有永生.
為什麼我們可以這麼肯定.
我們離開這個世界.
靈魂就去到神那裡.
我們得到這個永恆的.
舊恩 就是因為.
我們現在被摺上了.
那個本源.
那個樹根.
那個.
中心就是神自己.
所以.
在神裡面才有舊恩.
離開了神就沒有舊恩.

$^{1441}$再次告訴我們.
我們不可以誇后.
第十九節.
保羅又好像有些假設.
是不是他又意識到.
有些外邦人多看不起猶太人.
多拒絕猶太人.
所以保羅現在.
又再一次一個假設.
一種好像有些嘲諷的口吻.
李若雪.
那枝枝被摺下來.
是德衛叫我摺上.
是不是反映著.
某一類的人.
在一個不成熟的裡面.
以為猶太人是終究被撇棄.
外邦人沒有很大的恩典.
表面上現在是的.
但是原來在保羅的內心.
第二十節.
不錯.
可能你說得對的.
你表面看到.
他們因為不信.
所以被摺下來.
你因為信所以立得住.
那就不可以自高反要懼怕.
原來信而信服.
敬畏主.
信心帶來的.
是一份信而信服神.
是一份敬畏的心靈.
所以信.
是應該帶來一份謙卑.
降服.
和對自己所擁有的.
這樣的心情.
如果我們有一份自高.
以為是必然的.

$^{1481}$其實我們立即成為了猶太人.
那個錯誤的想法.
就是看上去.
自己是有血統.
有血緣的亞伯拉罕子孫.
那就錯了.
其實一切都是因為信.
所以猶太人.
現在被摺下來.
保羅說得很清楚.
是因為不信.
我們能夠摺上去.
就是摺了這個支條.
在這個本樹上.
是因為信.
所以.
必然的.
所以.
這裡再一次告訴我們.
信帶來的生命.
就是不可以自高.
因為一切是恩典.
不是靠行為.
能夠做到什麼都是恩典.
是的.
當如果一切都是恩典.
我們就應該有責任.
盡力做好神交付給我們的東西.
盡力就行了.
最重要是盡志.
就行了.
好之後.
而且我們反而要懼怕.
因為要帶著一份敬畏的心靈.
因為第21節就說了.
神既然不愛惜原來的資質.
也必不愛惜你.
可見神的恩慈和嚴厲.
兩方面.
對那些跌倒的是嚴厲.

$^{1521}$就是猶太人.
以自己血統.
身份.
肉身的身份.
或者有做多少禮儀.
這些.
神對這些跌倒的就是嚴厲.
向你是有恩錢.
這個你.
就是在主的救恩裡面.
被聖靈.
來光照.
感動和充滿.
在主耶穌現在救恩裡面的弟兄姊妹.
我們一直帶著這份感恩的心靈.
得到甚麼呢?.
神的恩慈.
一班人是嚴厲.
另一班有信心的外邦人.
就是神的恩慈.
所以這真是告訴我們一個很大的對比.
所以有生之年.
弟兄姊妹.
無論你現在是在外國的移民.
或者繼續在香港.
總之神給我們有生之年.
帶著信心.
和別人分享.
一切都是恩典.
我們為主應該不要太有保留.
為神的國度.
我們盡力去服侍.
讓人享受和我們一樣的恩典.
所以慢慢千萬不要因為慣了一年.
又是不知不覺.
就好像理直氣壯.
順理成章.
一切都是神應該給我們的.
因為連猶太人都不應該得到.
來被神撇棄.

$^{1561}$因為他們的不信.
所以我們要記得.
原來慣了在享用恩典裡.
信心要同時間成長.
否則我們的信心一滯後.
我們就不知不覺.
以為可以繼續坐享其成.
就出現了問題.
所以一切都是在信裡面.
所以第十.
接著第二十二節.
所謂只要你長久在他的恩前.
否則你都會被砍下來.
第二十三節.
而且他們若不是長久不信.
亦要被接上.
神能夠把他們重新接上.
這就給了我們一個很大的空間.
讓我們去反思.
讓我們去看到這個奇妙.
如果猶太人不是長久不信.
即是偶然歷世歷代總有一些.
慢慢有機會有信心信主.
是很難的.
向猶太人傳方是很難的.
但偶然間都會有一個叛國信.
他們就不是長久不信了.
他們就會重新被接上.
何況神有能力.
能夠把他們重新接上.
所以我們要帶著這份盼望.
為猶太人祈禱.
禁書間我們都印了猶太人的希伯來文聖經.
希望將來他們不知甚麼原因.
看到聖經就信主.
記住是希伯來文的新約.
希伯來文本身就沒有新約.
只有舊約.
現在我們禁書家.
房大兄他們印了希伯來文的新約聖經.

$^{1601}$即是他們看到羅馬書.
用希伯來文看到羅馬書.
知道原來求告主.
哪位主呢?.
就是耶穌.
就是歷史的那位耶穌去祈禱.
他們就得救了.
神能夠使他們重新被接上.
所以我們都要為猶太人祈禱.
真的希望他們能夠得到耶穌的救恩.
第24節.
你是從天生的耶駕覽上.
撼下來的.
尚且逆著聖德.
接在好駕覽上.
當然你不要有一種感到冒犯.
保羅為何這樣看待外邦人.
是天生的耶駕覽.
這樣撼下來.
這是一個比喻裡面的描述.
保羅從來都沒有對外邦人有鄙視.
看不起.
我們都說過他尊重自己的職責.
情願更加多外邦人能夠信主.
這是第13節的經文.
所以我們不要因為這樣的描述.
而說我們天生的耶駕覽.
相對猶太人.
其實他們一直都有十戒.
一直都不敢再拜偶像.
他們很專一地對耶和華.
不過在他們一直的認知.
就以為彌賽亞應該是榮耀的救主.
所以他們就看不到耶穌.
這位受苦的僕人.
就是他們等待的彌賽亞.
所以他們其實都是在一個本事上的支條.
但是我們相對於猶太民族.
我們就好像天生的耶駕覽樹.
好像很賤生.

$^{1641}$其實不是.
其實我們都是有文化.
有不同的知識.
不過我們對做物主的獨一真神.
我們從來都不是這麼單一地得到.
所以某程度我們就是這種天生的耶駕覽.
因為我們不是一開始就知道十戒.
知道這位耶和華十.
近乎沒有一個小朋友出生就是這樣.
所以我們應該要很珍惜.
我們的下一代.
我們的孫子孫女.
他們從出生就已經在媽媽的肚子裡.
就是基督徒.
我們要很寶貴這些的本事.
不是天生的耶駕覽.
我們應該這樣看.
我們要真是為他們多祈禱.
希望他們的勁情.
不只是留於一種表面.
而是心靈裡認識我們的祖父母.
我們的父輩.
他們是怎樣跟隨神.
憑信心跟隨神.
所以這個是很重要.
所以我們尚且逆著聖德.
接在普甲那上.
即是我們信到耶穌.
冥冥中我們能夠恩信清而不鬥.
可能我們是第一個家族裡的成員.
但我們有責任繼續為其他家人順處去祈禱.
因為我們這麼難.
都可以接上這個好格欖上.
何況當然這裡的比較是說猶太人將來的歸處.
何況這個本事的之子.
要接在本事上.
就是圍著下面的奧秘鋪路.
就是普羅所說的這個奧秘.
第25節.
弟兄們,我不願意你們不知道這奧秘.

$^{1681}$恐怕你們自以為聰明.
就是以色列人有幾分是硬心.
等到外邦人的數目天門.
於是以色列全家都要得救.
如經上所記.
必有一位救主從石安出來.
要消除那國家的一切罪惡.
這段經文在《二賽亞書》.
跟大家見見59章.
《二賽亞書》.
一向都說舊熟的.
不如我們看第15節這裡.
雖然誠實少見離惡的人反成略物.
那時耶和華漢見沒有公平.
甚不喜悅.
所以他親自出現.
他見無人拯救無人代求.
甚為詫異.
就用自己的傍臂施行拯救.
以公義扶持自己.
他以公義為規格.
以拯救為頭盔.
以報仇為衣服.
以熱心為外袍.
他必按人的行為施暴.
擄擄他的敵人.
第18節就是《羅馬書十》第二章提過.
他按個人的行為報應各人.
擄擄他的敵人.
報復他的仇敵.
向眾海島施行報應.
如此人從日落之處.
必敬畏耶和華的名.
從日出之地也必敬畏他的榮耀.
因為仇敵好像急流的河水.
從來是耶和華之氣所驅逐的.
接著留意第20節.
就是《保羅引述》這一節.
必有一位救贖主來到石安.
那國族中轉離過法的人那裡.

$^{1721}$這是耶和華說的.
耶和華說:至於我與他們所立的約乃是這樣.
我加給你的靈.
全給你的話.
必不離你的口.
也不離你後裔.
與你後裔之後裔的口.
從今直到永遠.
這是耶和華說的.
在以賽馬的年代.
已經告訴我們.
有一個時期.
這位救贖主會再次回到石安.
但是《保羅》這樣引述的過程.
似乎按的是字面的應驗.
因為以賽馬這一段引述的.
是說將來地老歸回那裡.
但是《保羅》現在這樣直接的引述.
是相當具體的告訴我們.
將來真是有一個日子.
以色列全家悔改歸向這位再降臨.
即第二次再降臨的彌賽亞耶穌.
這樣的意思.
因為這裡上文很清楚.
現在有幾分硬心.
等到一個時期.
外邦人的數目.
即是外邦人得救的數目天門了.
到底是怎樣計算呢.
香港佔多少比例呢.
全世界有這麼多的國家.
中國人多少比例呢.
代表多少呢.
美國人代表多少呢.
《保羅》沒有詳細說.
總之就是外邦人的人數.
得救的人數天門.
以色列就全家都要得救.
怎樣得救呢.
就好像經上所記.

$^{1761}$那位救主再一次降臨.
到石安山.
使那國家除去一切的罪孽.
即悔改完全歸向這位彌賽亞.
所以這裡很間接地說.
千禧年是有的.
耶穌會回來地上作王.
然後整個以色列就悔改.
在時序上.
在經文上引申出來.
耶穌真的再第二次回來.
如果結合新的經文.
就是我們這些已經去天堂的人.
應該去到樂園.
就會和耶穌一起回來.
在墳墓裡面首先復活.
升天 身體地頭.
然後我們就在天上.
有天使長的呼叫.
在地上死了的人.
復活有自己的身體.
然後就去到天上.
就開始千禧年.
最後猶太人就全家悔改.
第27節.
「就算我除去他們罪的時候.
這就是我與他們所立的約」.
所以那個將來的約是除罪的約.
猶太人是全家悔改歸向彌賽亞.
但現在就是整個民族拒絕彌賽亞.
所以那個分水嶺就是.
外邦人為什麼這麼蒙恩呢?.
就是因為我們在救恩裡面.
所以我們一定要告訴人.
人必須在有生之年.
要在救恩的愛.
就是耶穌的十誡救恩裡面.
第28節.
「就著福音來講.
他們為你們的緣故是仇敵」.

$^{1801}$就著….
可能是在移動車子.
可能是移動車子.
好像是684.
684在裡面.
在裡面.
他們練歌.
對不起.
因為外面已經移動車子了.
大家不要介意.
這裡就說了.
再一次說仇敵和無愛的對比和關係.
所以很難來進至.
將猶太人現在被折下來的.
在神的救恩以外的光景.
無愛成為神的仇敵.
但現在我們接支接上去的本源.
就是說信心的行列裡面.
一直延續亞伯拉罕神的救恩計劃.
以致我們認識耶穌.
求救得救.
現在就在神的愛裡面.
是一個很重要的分水嶺.
所以就著福音來講.
我們都是恩信請言的福音.
他們猶太人為了你們的緣故.
你們就是外邦人.
當然是蒙恩信主的外邦人.
的緣故.
他們就是仇敵.
就著簡選來說.
簡選就是說神開始的救恩計劃.
在他的預知.
在他的永恆裡面.
去拯救有信心的人.
那些人就可以得救了.
例如亞伯拉罕.
例如彭亞國.
所以就著簡選的角度來講.
他們為哪些人呢.

$^{1841}$為列祖的緣故.
他們是無愛的.
就是說原來保羅是很有信心和盼望的.
因為神是在列祖的時候.
才有這個簡選.
是堅定不移的簡選.
不會因為一代一代的猶太人.
批士 拒絕.
來攻擊彌賽亞.
神就好像將他們打入永不超生.
原來不是.
神自己的計劃是從來都不改變的.
所以這個就是我們人的一個福氣.
和那個盼望.
所以就著簡選來講.
這個斟酌的字眼是很重要的.
就著福音.
猶太人就是一個敵神的.
就著簡選.
因為亞伯拉罕這幫人的無恩和有信心.
他們就是無愛了.
就是猶太人原來都在愛的行列裡面.
所以現在還是這樣被撇棄.
那就是說是有盼望.
第二十九節.
因為神的恩賜和選召是沒有後悔的.
就是說神自己是不會後悔.
他定了這樣的計劃.
這個就是一個很大的智慧.
第三十節.
你們從前不信服神.
如今因為他們的不信服.
你們倒蒙了連戌.
這樣他們也是不信服.
叫他們恩著施給你們的連戌.
現在也就蒙連戌.
但是他的意思其實也很簡單.
不過用不同的角度.
不同的身份去講.
我們以前本來就是不信服神的.

$^{1881}$本來就不認識他.
但是真的因為不知道什麼原因就去信主.
譬如我也講過.
真的從來沒有想過.
在1987年的12月24日那天.
明明都看完電影.
英雄本色.
回家突然間裡面的聲音就說.
其實是我弟弟一個多月前告訴我的.
在他的小學.
就是教會的小學的聖誕節的晚會.
有一個所謂的報道會.
是給家長的.
叫我去聽.
是一個多月前叫我.
我完全沒有放在心裡面.
但是不知道什麼原因.
裡面的聲音就說.
你去聽吧.
我從來沒有聽過這些聲音的.
從來不知道有這樣的聲音.
但是就知道要衝去.
去到已經是一半.
已經枯燥了.
然後就舉手決戰.
你可以想像到.
這些經文其實真的告訴我們.
恩典難道.
有時候真的很戲劇性.
所以我們千萬不要以為.
一些人永遠不會信主.
因為當時我弟弟阿昌.
他的同學.
讀中六的同學.
因為開始回教會.
有時候他就來探我.
我在家裡拿著一支煙.
跟人家聊天.
笑他們.
他們沒有想過我會信耶穌的.

$^{1921}$所以整班人和那個傳道人.
看到我這麼出現.
開心到跳起來.
然後問我剛才聽完怎麼樣.
我都沒有什麼期望.
你走上來幹什麼.
坐到最後.
已經滿了人.
他現在說我不知道.
我剛才舉手決戰.
舉手舉手.
他很開心.
立刻分組.
信什麼我怎麼知道信什麼.
總之就是剛才舉手決戰.
這樣再說說.
就回家了.
回到家.
那個時候已經被人家發覺.
我很喜歡看.
完全看得明白.
從創世記一直看.
從來沒有試過.
你可以想像到那個奇妙.
那個釋放.
所以我們從前都是不信服神.
如今因為他們的不信服.
所以我們現在蒙了連戌.
這樣這個魔鬼.
我們都要將心比己.
古羅說.
我們現在他們是整個民族被撇棄.
所以我們其實都很生氣.
猶太人用這麼高科技.
這麼精準的炸彈去炸哈馬斯.
炸這些加沙的居民的時候.
我們都要為他們的背影.
為他們現在對神的那種否定去祈禱.
希望他們都能夠蒙神的憐憫.
因為我們都能夠尚且可以蒙連戌.

$^{1961}$所以因為師給我們的連戌.
所以他們都應該可以蒙連戌.
因為神將眾人都圈在不信服之中.
特意要連戌眾人.
即是說一切都是來自神的恩典.
白白的憐憫.
每一個民族.
每一個個人.
每一個家庭都是.
所以我們要不斷維持在這一份的蒙恩的.
不單只是情感上的感受.
而是在整個意識裡面.
一切都是神的恩典.
所以我們千萬不要這麼容易意氣用事.
稍為有些不開心.
就說不回教會了.
不事奉了.
就不要這樣看.
因為一切都是神的連戌.
開始有這個最後的頌讚.
保羅引述了兩段聖經.
《異象》書的四十章和《約伯記》.
這個重歌.
其中和第八章不同的.
第八章特別強調神永恆不變的愛.
在屬於他的兒女身上.
這裡特別強調神的智慧.
在整個救恩計劃裡面的部署.
安排和設計等等.
心災就是很偉大.
這個意思.
用文言中文.
很偉大.
就是讚嘆神的作為.
以一個拉比的心情.
以一個外邦人的仕途.
就是百感交集的保羅.
來讚嘆神的豐富智慧和知識.
豐富的智慧和知識.
你可以想像到.

$^{2001}$這個地球那種精密.
那種設計.
那種巧妙.
我們都很難以去形容,想像.
譬如我們剛剛對上的風球.
在太平洋只是一些天氣的原因形成.
都可以帶來這麼嚴重的破壞.
整個社區的人都要嚴陣以待.
感恩沒有吹到一隻一隻.
否則都很麻煩.
一個最先進的社會.
都可以有很大的破壞.
樹木都可以連根拔起.
不過是很自然的炎著.
太平洋因為那些暖流.
那些溫度越來越高.
所以那個旋轉的氣候的變化.
就帶來形成風球越來越多.
越來越嚴重.
你可以想像到.
這些純粹是小小的大自然的自然炎著.
我們人類都受不了.
何況做民主的豐富智慧和知識.
所以我們要對神永遠帶著一份渴慕.
認識祂.
還有祂有祂的主權在我們的生命裡面.
所以祂的判斷.
所以神原來不是一個沒有位格的存在.
祂不是一個氣體.
祂不是一個電力.
祂是一個有位格的判斷.
何其難測.
祂的蹤跡何其難尋.
誰知道主的心.
本來是直接引述約伯記的.
約伯記是他得到恩典之後對神的讚嘆.
誰做過他的謀事.
誰先給他償還.
一切都是否定據.
他提出了那個問題.

$^{2041}$答案就是顯而易見.
一定不是.
一定是實自己.
所以第36節就是《羅馬書》11章.
在說救恩最後這一節.
然後就說倫理的關係.
丙子會在教會裡面的關係.
第36節.
因為萬有都事本於他.
依靠他 貴於他.
就是說救恩.
因為一直都在說救恩這個大題目.
救恩的起始都是神的主權.
神的預定.
神的救贖.
開始透過主耶穌的計劃降世等等.
所以一切都是本於他.
依靠他.
就是說一直那個救恩的發展.
都是神一直的經營.
一直用他的智慧.
用他的判斷.
就是說我們能夠信耶穌.
我們幾時幾日信耶穌.
都是神自己的計劃作為.
而且是歸於他.
就是一切的榮耀.
那個救恩.
將來得救的人數.
一定會是天門在天上.
是不會有些人神預定了.
對不起你預定你的事.
我是不相信的.
我是一生都不會相信的.
就應該不會有這樣的人.
將來又上不了天堂.
又在地獄裡面.
然後又再一次這樣去說神.
是沒有這樣的人的.
因為神的預定一定可以順處的.

$^{2081}$只不過是我們不知道誰是神預定.
我們要記得.
等等.
所以要願榮耀歸給他.
直到永遠.
願榮耀歸給神.
在這個救恩裡面.
所以我們越認識救恩.
人信主的榮耀.
人信主的困難.
人對神的背影.
拒絕.
嘲笑人信了主之後的謙卑.
降伏.
信心增長.
愛主.
服侍主.
例如Chris.
余建峰弟兄.
他們現在有個團隊.
快月底他們就去埃及了.
他們住的是相對簡陋的.
基本的.
還是為一些當地的信徒領袖去籌款的.
所以你看到日弟兄姊妹這樣為福音擺上.
你可以想像得到.
一切的榮耀都是神自己的作為.
你又反過來看到.
有些人一離開了主.
那種嘴臉.
那種硬心.
那種對神.
一看到神的憤怒.
一說教會.
不要再說教會了.
再說教會就沒有事做了.
出去吃東西沒問題.
有時候你都會很痛心.
為什麼他回教會.
慢慢就會變成這麼硬呢.

$^{2121}$兩個極端.
所以我們仍然都要等候神.
所以願榮耀歸給他.
直到明月.
都不夠時間揭經文.
所以看看二創亞書這一段.
40章這一段.
二創亞書.
40章的13節二創亞書.
當然二創亞對以色列人的質問.
又是一些好像撒文的方式.
好像一個審判.
不多了.
只能看一點.
40章的第十節開始.
他必像牧人牧養自己一樣.
用薄肥聚雜羊羔在圍中.
慢慢引渡那羽羊小羊的.
誰曾用手心量諸水.
其實提出這些根本沒可能得到.
一點都不需要複雜的問題.
一定是神.
誰能夠用手苦口.
這樣的手苦口.
來量蒼天呢.
用星斗去乘大地的尋圖呢.
用秤去秤那個山嶺呢.
用天秤去秤江嶺呢.
用一些這樣的道量行.
去跟這個宇宙做一個對比.
即是說根本沒可能.
我們就應該要宏觀一點看神的作為.
第十三節就是這段.
誰曾猜度耶和華的心.
或者作過他的謀士去指教他.
他與誰相議.
誰教導他.
誰將公平的路指示他.
又將知識教訓他.
將通達的道路指示他.

$^{2161}$看萬民都像水桶的一滴滴.
又算如天秤上的微塵.
你可以想像到.
作者爾塞亞對於神的偉大作為.
與人的渺小.
做一個很強烈的對比.
所以保羅就選取十三節.
誰做過他的那個謀士.
來給他.
他要還呢.
一切都是.
所以最終保羅用一個最後的總結.
因為物有都是本於他.
依靠他.
歸於他.
榮耀歸給他.
這個都是我們信主.
我們現在每日維持一份愛主的心.
都是神的遺棄.
最終我們人生都要完結了.
我們要歸榮耀給神.
所以我們一定要到最後.
我們的人生無憾歸榮耀給神.
不是甚麼特別的良藝.
但是問心無愧的盡了我們信主.
以及盡量傳福音的責任.
至少下一代都是敬虔的.
將這個遺產交付給下一代.
就是他們都是繼續承繼.
這個虔極對主的心.
所以都奉勸大家.
建議不是奉勸.
將這一節36節成為我們人生的座右銘.
真人是很寶貴的.
萬有都是本於他.
依靠他.
歸於他.
月榮耀歸給他.
直到永遠.
阿們.

$^{2201}$也都應該是我們每一個人完結的時候.
如果有人幫我讀.
Alex還沒死我就不死了.
幫我出來讀這一句.
當年曾牧師說.
如果有人幫我讀這一句.
他見證著.
曾牧師一生就是在這一句話裡面.
在這一句金句裡面.
這樣去活.
盡量在這句金句裡面去活.
其實是很美麗的.
因為你一切都依靠神.
一切榮耀都是歸給他.
而且是本於他.
如果不是他給你的.
你根本就不會有.
所以榮耀是應該歸給神.
一直都是屬於神.
人生就很平衡.
還有一份平安.
和真正的心靈的福樂.
我們要完整的十一章.
這個救恩.
由第三章打後.
去到十一章.
這個大段落.
算是比較完整的.
逐段逐段的說了.
十二章開始一直到十五章.
就是內部的conflict.
如何能夠在一個屬靈的.
主要是愛.
在基督的層面去看大家不同.
不同可以是一些分歧.
不同可以是大家不同的做法.
如何能夠有智慧的.
讓教會可以是一種.
大家同有一個心智向福音.
這個就是大前提.

$^{2241}$就是大家都要為宣教.
來放下自己.
一起普告.
感恩 天父讓我們在晚上.
又可以用一個半小時.
將第十一章最後這一段.
說得清楚.
多謝你讓我們能夠現在.
我們天生的耶格蘭.
能夠被接上.
在這個本算.
在神的恩典.
神的救恩計劃裡面.
我們很盼望香港人.
或者我們宣道會.
或者我們錦繡堂.
紅恩堂 佳恩堂等等.
宣德都是.
我們現在有四間教會.
我們都是仍然在主的恩典裡面.
我們每一個弟兄姊妹.
都是很寶貴.
我們在恩典裡面.
因為我們的信.
所以我們蒙恩.
我們可以在生活.
在信仰上仍然去持守.
特別是這個星期一.
就是我們的會友大會.
已經是第四十五週年.
天父幫助我們.
我們一直秉承著前人.
我們的前輩.
他們當年如何來建立.
在信心裡面.
來有這個錦繡堂的出現.
亦都是一直秉持著.
宣道會的宣教精神.
我們有印尼語.
菲律賓語.

$^{2281}$我們紅恩都有泰國的肢體等等.
我們很盼望天父讓我們看到.
原來我們在這麼大的救恩的脈絡裡面.
我們是在享受神給我們在主救恩裡面的福樂.
所以我們為離開了的土共.
被折壓的.
我們要流淚.
要為他們去擔心.
要為他們土共.
希望他們重新可以被接上.
尤其是猶太人這個民族.
盼望主憐憫他們.
讓我們都能夠為他們土共.
我們不知道何時是外邦人的數目天門.
何時是以色列全家歸主.
但我們盼望我們都擺上這個心力.
特別是《希伯來文》的《新約聖經》.
似乎是我們第一批的中國人.
來出版印刷.
盼望送到猶太,以色列的教會.
有很多猶太人.
尤其是年青一代的猶太人.
有機會去讀,去信.
他們來求告主.
你的聖名.
來得救.
我們仍然盼望你.
多謝你讓我們完成了一個多大的創舉.
一到第十一章.
我們都是有系統地去研讀.
願主幫助我們.
好像第八章.
尤其是普羅的讚歎.
就是沒有東西可以與神的愛隔絕.
沒有東西可以使我們信徒與神的愛隔絕.
第十一章讚歎的是神的豐富榮耀的知識.
願主得到一切的榮耀.
奉耶穌基督的名.
阿們.
\newpage



\section{羅馬書}
\label{sec:FAt1alHLvdg}
\textbf{2025-10-17 周五晚查經班《羅馬書(25):分題:從奧秘到信徒的生活》 講員:曾裔貴牧師}
\newline
\newline
連結: \href{https://youtube.com/watch?v=FAt1alHLvdg}{\texttt{https://youtube.com/watch?v=FAt1alHLvdg}} ~~~~ 語音日期: 2025-10-19
\newline
\newline
\hyperref[sec:bw6FVf5f_n0]{\small{< < < PREV SERMON < < <}}
~
\hyperref[sec:index_chronic]{\small{[返順時目]}}
~
\hyperref[sec:index_scriptual]{\small{[返順卷目]}}
~
\hyperref[sec:yNx82XHOkbk]{\small{> > > NEXT SERMON > > >}}
\newline
\newline
$^{1}$多謝天父讓我們再一次可以來到他在萬常.
我們有系統地去研讀羅馬書.
多謝你讓我們能夠從第一章.
講到神的憤怒臨到不祈不移的人.
來到我們今天要講的第十二章.
我們看到我們每一個基督徒.
原來當我們的內心有了更新之後.
有聖靈的進駐.
我們每一個基督徒可以去測驗.
神在我們生命裡面的旨意.
這件事真是很奇妙.
是非常之偉大.
原來神是一個一個地拯救我們.
幫助我們能夠在很多轉折的經文.
我們多停留和去思想.
以致有些可能我們平時在崇拜.
都未必能夠夠時間詳細講的.
我們可以有系統地去研讀和明白.
使到保羅他內心想交代,教導的.
叫我們都能夠掌握,明白.
以致激發我們的心靈.
有一份的服侍,對主的效忠.
我們今天晚上等候在你的面前.
亦都盼望這些錄製了的錄音.
都成為不同的兄弟姐妹.
在不同的時間 當他們聆聽的時候.
都得到造就 我們仰望著.
奉耶穌基督的名 阿彌陀佛.
好 我們開一個大一點的PowerPoint.
我們之前一直都在說.
我先開那個PowerPoint.
今天我想給一個題目.
就是上一個星期我們講過.
保羅說他知道有一個奧秘.
就是當外邦人的得救數目添換之後.
以色列整個民族全家要得救.
當然我們在第九章已經講到.
不是肉身的以色列人就是等於得救的.
必須是蒙神揀選.
才可以讓他得到這個救恩.

$^{41}$所以以色列全家.
就不是每一個的以色列人.
所以大家要留意這件事.
所以由奧秘 接著去到第十一章.
我們上個星期有提到.
就是一個很大的重站.
對於主 特別是對於他的一切.
救恩上的那個計劃的智慧.
和第八章講到的神的愛.
又是有一點點相同的.
或者說更加補充那個豐富.
這個就是我們在上個星期提到.
關乎到奧秘和整個救恩帶給保羅.
好像呼喚我們整個宇宙.
都要重站我們這一位的主.
他的奇妙,他的智慧,他的知識等等.
但是十二章我們今天會講.
多麼寶貴.
保羅突然間改變那個思路.
以及將一個由神學那麼高的.
神的偉大的智慧.
突然間收窄了範圍.
就是每一個得救的信徒.
就要做一件事.
或者說必須要這樣理解自己.
由信耶穌的一刻.
應該就是這樣理解自己.
所以這個就是讓我們看到.
得救原來是一個開始.
和神的生命有這個連繫.
所以和之前的第三章.
或者第五章,第六章面對罪的問題.
面對內心的救人.
其實都是一脈相承的.
我們很快看回上一個星期我們講到.
第十一章保羅的思路.
神是沒有放棄以色列人.
然後他也說以色列人現在是暫時被撇棄.
但是在神的救恩的智慧計劃裡面.
猶太人按照亞伯拉罕他們的先祖.

$^{81}$其實就是這一棵樹.
就是那棵橄欖樹.
所以是神自己親自去栽植的.
所以這一棵樹的樹根.
那個樹等都還沒有連根拔起.
所以有部分的以色列人.
當然各時代的以色列人.
被神撇棄了.
不代表神放棄了以色列人.
然後我們又看到在這個過程.
保羅很巧妙地用橄欖樹的枝子.
本樹的枝子被摺下來.
然後有一些野橄欖的外邦人.
就可以被摺上這個本樹.
作為一個這樣的比喻.
就是說現在普遍當主耶穌的救恩.
由使徒們在耶路撒冷一直傳開之後.
外邦人就很容易因順清賢.
明白了耶穌的救恩.
接受了主耶穌在十誡上的犧牲代屬.
聖靈就會在人的心裡面有感動 有工作.
所以你會看到有很多的外邦人.
不是猶太人的外邦人陸續的順主.
原來有一個目的.
就是要激動 要使猶太人嫉妒.
以致他們去反思為什麼自己現在在救恩的門外.
所以這個就是一個這樣的進程.
當然25至32節.
就是保羅說到的奧秘以色列全家將要得救.
這個就是第十一章的思路.
接著我們可以看下面.
所以你會看到我們可能要先讀聖經.
將很重要的反差給大家有一個這樣的理解.
就算我們不需要懂得希臘文.
當然懂得更加好.
我們已經感受到 中文我們已經感受到.
那個由神學到信徒的生活的一個轉折.
我們先看看經文.
十一章我們再看這個續集.
第33節這個續集.

$^{121}$他說心災很偉大.
十豐富的智慧和知識.
這個以賽亞書的40章.
和約伯記的41章.
他的判斷何其難測.
他的蹤跡何其難尋.
誰知道主的心.
誰作過他的謀事.
誰是先給了他.
使得後來償還.
就是說每個人都可以知道神的心.
這是第十二章.
第二節就已經要求信徒.
要察驗什麼是神.
在我們個人或者在教會.
那個生命裡面的那個旨意.
這裡說沒有人可以察驗神的偉大.
神的智慧.
在這個救恩裡面那個奇妙.
但是十二章立即就告訴我們.
每個信徒都要察驗.
所以你看見保羅那種豐富.
那種奇妙的他得到的啟示.
然後告訴我們誰是先給了他.
使得後來償還.
很多很多的切文句.
其實答案是顯而易見的.
就是說一切都是神自己.
使得救恩來開展.
那個過程就是救主來這個世界.
或者救若的時候就是醞釀來趁.
要降生.
接著主耶穌就是這樣來降生.
保羅就見到復活的耶穌.
接著現在在羅馬書的時期.
保羅寫的時候.
就是救恩正是被傳於外邦的時候.
所以你可以想像到.
神就是這樣來讓我們見到他自己的作為.
第36節.

$^{161}$因為萬有多事本於他.
倚靠他 歸於他.
願榮耀歸給他.
直到永遠 阿們.
這個就是保羅來到最後最宏大的一個重擇.
所以越是認識救恩.
和自己生命關係的基督徒.
每一個時代都要重擇著.
都要從生命裡面向神發出這一份呼喊.
這一份敬拜.
因為榮耀真的要歸給他.
因為萬有都是本於他.
救恩是神用他最奇妙最偉大的智慧.
和那個你可以說是最大的犧牲.
就是他的愛子為我們人類釘十字架.
所以保羅在這裡的時候.
再一次將一個敬拜來發出.
雖然好像是他個人在寫希利文的時候.
自己在台頭向神發出敬拜.
但是都已經足夠我們看到.
保羅神學的豐富.
對救恩的澎湃的認識.
突然間第十二章.
很明顯就是一個轉折.
是一個新的段落.
你怎麼看都看到是一個段落的開始.
但是立刻有一個所以在這裡.
翁這個希利文的連接詞.
他說所以弟兄們.
我以神的慈悲勸你們.
我用修定版的翻譯.
我用修定版.
對不起不是修定版.
是《新漢語譯本》的第一節.
他說所以弟兄們.
我以神的憐憫勸你們.
要把身體作為活祭獻上.
就是這個身體.
是聖潔的.
是神所喜歡的.

$^{201}$這裡相當明顯讓我們看到.
在神學的頌讚敬拜神之後.
就開始保羅探討在教義.
神學的教義.
去到實際的生活的實踐.
生活的實踐.
第三節就說那個平台.
必須在教會裡面.
當然是那些人.
不是那個教堂.
這些概念大家都要知道.
在這樣的情況下.
那個實踐首先就是每一個信徒.
必須真的要有最大前提的裝備.
所以參加事奉的弟兄姊妹.
如果是被洗禮,被得救.
其實根本就不是叫事奉.
根本就說不上事奉神.
還有根本就不知道自己原來是活祭.
根本就不是活祭.
是在神眼中的憤怒.
因為第一章已經說了.
你可以想像.
當我們回頭看這裡的時候.
保羅就告訴我們.
在這個神學的轉折裡面.
他很快將對這位做物主的頌讚.
轉移在信徒的生活層面裡面.
首先對自己裡面的人.
生命的轉變.
要有一個很強的意識.
所以讀羅馬書.
很應該看到自己.
一邊讀的時候.
看到自己原來不再在神的憤怒之下.
而是活在神的愛.
第八章是.
而是活在神的智慧的揀選之下.
第九章也是.
一邊讀到這裡的時候.

$^{241}$告訴我們.
我們求告主明的.
就是說跟神有一個生命關係.
我們開始透過祈禱.
跟祂有一個接觸.
這樣的基督徒.
接著保羅就告訴我們.
原來我們個人.
在神的救恩裡面.
我們立即成為了一個生活上的活祭.
所以不是我們的頭腦認同福音.
是整個人的裡裡外外.
我們在神的救恩裡面.
我們是整個生命被救贖.
是從裡面.
我們認識自己現在的新的身份.
新的價值觀.
第九節到二十一節.
就是新的價值觀.
人與人之間相處的價值觀.
三到八節.
就是在教會這個平台.
在這個場景.
怎樣能夠運用神給我們的恩賜.
你可以想像到.
既是有救恩.
清楚自己與神的關係.
是無礙的兒女.
是可以隨時禱告.
因為聖靈與我們的心同征.
我們是神的兒女.
我們可以去到神的面前.
接著告訴我們.
同時間我們有神給我們的恩賜.
這是第三到第八節.
其中一個比較大的描述.
當然關乎到恩賜.
就有很多經文.
特別是前十二到第十四章.
是詳細地做恩賜的調查表.

$^{281}$來講.
講這麼多.
講這麼詳盡.
一定是有問題.
才會講這麼詳盡.
因為那裡當然是另一個話題.
是哥林多教會.
他們爭誰的恩賜最厲害.
最厲害.
最顯眼.
所以保羅再將恩賜的本來目的講出來.
但是在羅馬書這裡的恩賜.
是相當正面的.
告訴我們.
每一個基督徒.
首先要對救恩有一個降服.
敬拜.
頌讚.
你知道自己與神的關係.
接著就開始探討.
我們在神的救恩裡面.
我們獲得什麼新身份呢?.
這樣的時候.
可能今晚只能夠講第一第二節.
我們還沒讀完經文.
不如我想跟大家讀第三節.
我用舊的和合本.
因為我預備了第一到第二節.
但是後面那段可能大家先看看.
我們可能看到.
不如看完這一章.
第三節開始.
我憑著所賜我的恩.
對你們各人說.
不如看自己過於所當看的.
要照著神所分給我們各人信心的大小.
看得合乎中度.
合乎中度的是信心.
有多大的信心領受.
神給你有多少工作的範圍.

$^{321}$這是保羅告訴我們.
合乎中度.
修昔的詞句是相當美麗的.
合乎中度.
不要過多也不要過少.
是一個很美麗的教導.
第四節.
正如我們一個生子.
如何合乎中度呢?.
首先就是那個心態.
要認定自己是在基督裡面其中一個器官.
這就很重要了.
不是在基督以外的自己獨立一個人.
所以有時候在教會裡面.
有些人一不喜歡.
他就自己脫離了.
不再服侍.
這是很可惜的.
因為我們是在基督裡面.
記住第四節.
正如我們一個生子上有很少的肢體.
不同的器官.
肢體也不都是一樣的容處.
很清楚的.
有不同的器官發揮過作用.
所以保羅就這樣來描述.
我們者許多人在基督裡成為一身.
一而再出現在基督裡.
這個生子.
這個很完整的身體.
當然是裡裡外外不同的器官.
成為一身.
然後互相聯絡.
作肢體.
也是如此.
按我們所得的恩賜.
各有不同.
就是說.
我們大家都是活祭的生命.
大家都要察驗神的喜悅的旨意.

$^{361}$在我們生命裡面.
但是我們同時間.
都是在基督這個身體裡面.
來去被安插在神的幫助.
教導下的一個肢體.
所以我們就要察驗.
我們有什麼恩賜.
我們怎樣來去用.
所以你看到保羅是非常平衡的.
教導信徒.
要有一個重點就是.
信了耶穌.
很簡單來說.
信了耶穌.
是必須有行動的.
在生活裡面的服侍和成長.
第六節.
按我們所得的恩賜.
各有不同.
或者說預言的.
就當照著信心的程度說預言.
當時有這些這樣的預言.
這些作預言的先知.
例如11章的《聖體行傳》.
有一個叫阿加布.
就預言.
他拿了保羅的皮帶.
就說將來這個皮帶的主人.
就會被綑綁.
帶上耶路撒冷.
這些就是講預言.
第七節.
或作執事.
執事是那個部門.
就是做一些服侍在教會裡面.
可能是一些開會.
有些事情要商討.
例如當時還有繁文公用.
需要預備一些飯吃的.
需要有不同的策劃.

$^{401}$和編排.
這些都需要有一個系統的.
所以執事可能就是做這些.
很快的 保羅提到.
在《提多書》提到.
要訂一些寡婦的州制的系統.
有些是真寡婦.
有些年輕寡婦就不要在教會了.
快點去嫁人.
有些真的是年老的.
無依無靠的.
這些就是真寡婦.
就要列在寡婦的名冊.
按時教會就公認.
所以這些都需要有一點智慧,策劃.
有那個嚴肅性的弟兄姊妹.
來去做.
有一個持平等等.
這些其實都是一般的智慧.
這裡告訴我們.
如果作執事就當專一執事.
或作教導的就當專一教導.
就是大家都在自己的專注裡面.
不要和其他人比較.
但是專注在自己所做的時候.
但是那個重點是.
圍著整個基督的身體的.
不是自己做惡了.
說自己霸了.
不是這樣的.
而是專注在那個事工上.
盡心竭力.
接著第八節.
或作勸化的就當專一勸化.
施捨的就當誠實.
自理的就當殷勤.
憐憫人的就當甘心.
這裡是很清楚的將一些大陸的恩賜.
這樣來列舉.
和哥林多前書12章到14章講的那些恩賜.

$^{441}$有些是不相同的.
或者沒有數得那麼清楚,那麼詳細的.
因為這裡不是想這樣來得那麼詳細.
而是很想告訴我們.
我們有一個平台.
在神學了解了自己的得救之後.
我們就要立即看到.
整個生命是神已經喜悅.
是神愛的生命.
等下我們回頭講第一節.
接著這個被神愛的生命.
就要開始測驗了.
你想一下.
剛剛初信連什麼叫信耶穌都不是很清楚的.
保羅已經說我們要測驗神的旨意.
通常這些是每到讀神學才會想的東西.
但是原來羅馬書是非常好及性的.
給每一個信徒都要知道.
原來自己在救恩裡面.
神這樣愛我們,這樣看我們.
說要我們測驗.
剛剛我們11章才講完.
誰會知道神的心.
誰會幫他想了.
借給他償還.
沒有可能的事.
但是現在每一個基督徒.
我們都要每時每刻測驗.
接著第九節.
我們也要讀了.
因為這一段就是人際關係的一個相當重要的教導.
如果舊的《和學本》有一個標題.
基督徒的生活守則.
這樣也相當對.
他說愛人不可虛假.
惡要厭惡.
善要親近.
愛弟兄要彼此親熱.
恭敬人要彼此推讓.
恩情不可濫度.

$^{481}$相當有那個真言的氣味.
那個格式.
好像一對一對那樣.
保羅自己也是一個拉尾.
對舊樣也是相當的通透認識.
所以一個真言的世界就是.
一個齊齊齊在神國治理的猶太民族.
他們正常應該就是這樣.
當然正常是這樣.
但是以色列人應該從來都不會這麼齊齊齊.
有這樣的情況.
所以有和國的歷史.
有真言這些的教導.
作為一個平衡.
神期望的心意.
一個正常的.
有神管理的國家.
世界應該就是這樣.
所以這裡告訴我們.
教會的生活很多的原則.
我們因為教會不是一個人.
就是一群人.
然後第十三章打後.
就不只是一群相同文化的人.
是有猶太人.
是有外邦人.
就是不斷將那些複雜性層面讓我們知道.
但是我們每一個都一定要根深蒂固的認識.
自己就是蒙神舊恩的.
活祭的一個弟兄姊妹.
這個大前提.
當我們這樣能夠明白.
就相當寶貴.
我們可以學習這些新的價值觀.
我們逐樣去做.
就會得到神使我們的生命會很大的成長.
生命會豐富.
還有就會好像其他經文講的結果子.
將來我們會得到神給我們的公義冠冕.
我們再看多一點.

$^{521}$反正都是讀這張聖經.
因為很多這些轉折.
我們都需要詳細講清楚.
第十二節.
在寂寞當中要起來.
在患難當中要忍耐.
禱告要行切.
因為就是等待的階段.
聖徒缺乏要幫補.
這個彼此相愛的好撒瑪尼加人.
所以保羅都認定這件事.
講到周濟窮人.
昨晚我們那個行在社區的福音.
因為我們整個課程就是在講.
我們教會 每間教會 每個社區的教會.
怎樣去做這些社區的工作.
所以幫助窮人.
這個肯定是新疆有多大的教導.
所以這裡告訴我們.
聖徒缺乏要幫補.
客要一味的款待.
逼迫你們的.
要給他們祝福.
只要祝福 不可咒詛.
特別強調.
不要只是口說完.
真的要祝福.
與喜樂的人要同樂.
與哀慨的人要同哭.
這個同理心.
我想大家都會明白.
要彼此同心.
不要志氣高大.
倒要苦就卑微的人.
或者卑微的事情.
當我們的內心.
如果能夠恰如其分的看自己.
我們就能夠.
所謂游刃有餘.
對不同的人 對不同的事物.

$^{561}$不會一個態度.
一種自己的眼光去看.
其他人 其他的事.
不要自以為聰明.
不要以惡報惡.
再一次說.
眾人以為你的事要留心去做.
若是能行.
總要盡力與眾人和睦.
不斷說和睦.
不斷說盡力.
因為有些人可能根本不想和你和睦.
但是我們作為一個活祭的基督徒.
就必須履行在人生裡面.
真的不能夠有仇人.
或者有仇恨.
第十九節.
稱愛的弟兄不要自己失淵.
就是說有些人不僅不和睦.
還繼續加害於你.
無論暗中的 明顯的.
保羅就面對著那些猶太人的拉比.
對他的逼迫.
現在在羅馬書的時候就沒有那麼明顯的.
針對 但是就沒有真的去追殺他.
當他第三次上耶路撒冷的時候.
那些拉比就直接要去告狀.
一直這樣來到要置他於死地.
所以保羅真的面對著那些假弟兄.
面對在耶路撒冷的猶太教的這些忠奸分子.
特別是安提哥上去首節的五循節.
應該是二十二還是二十三章的《聖行傳》.
一見到保羅就二話不說.
以為他帶著那個叫做特洛菲摩.
在聖殿範圍.
但是特洛菲摩是外邦人 可以進入聖殿.
其實保羅沒有帶他進入聖殿.
保羅自己去聖殿還願.
所以一見到就已經火遮眼.
無論怎樣的捏造都馬上出現了.

$^{601}$就抓住保羅和特洛菲摩.
你就可以想象到保羅說.
原來真的是若是難行.
總要與眾人和睦.
不過接著就說親愛的弟兄.
不要自己伸冤.
寧可讓步 聽憑主路.
主路就是主的審判.
有些人是不會拍實.
我們要繼續認識我們在人生.
我們不去結怨之外.
我們被別人去針對和被別人去陷害.
我們不要自己伸冤.
當然我們可以伸冤.
但是最終的結果我們交託給神.
保羅自己也是被斬頭.
昨天我才認識了一個新的概念.
從來沒有聽過的.
原來後世 即是現代的畫像.
保羅原來是拿著一把劍的.
我以為他隨時會和別人打.
原來後世的畫家 敬勤的畫家.
全部都告訴我們.
那把劍原來就是象徵著保羅被斬頭而死.
所以原來無論在現在聖伯多祿的大教堂.
有保羅的雕像都是拿著一把長劍.
我那時候從來都不知道.
昨晚老師說了我們才知道.
保羅都是拿著那把劍.
所以和羅馬孫有抵觸.
其實不是的.
保羅就是這樣 聽憑主路.
不要自己伸冤.
相當不容易的.
我想他想起主耶穌.
在十字架上.
受害不說威嚇的話.
被罵不還口.
就是耶穌基督的默默無聲.
真的不容易的.

$^{641}$為自己伸冤.
為自己來解釋.
而且那些是莫須要的那種.
對主耶穌的這種侮辱.
耶穌基督是一句都不出聲.
帶到希臘王的時候.
真的是默默無聲.
希臘王說你來講.
你拜我 我就可以釋放你.
耶穌基督根本就是不屑跟這個人這樣講.
所以保羅也是秉承著.
救恩耶穌基督的十字架上的大愛.
他現在這個時候.
保羅也是傳道的中期.
都沒有去到後期.
他知道自己差不多離世的殉道.
但是現在他已經有這樣的教導.
作為使徒.
都真的給我們看到.
他真的是身體力行.
因為在《四輪傳》十八章.
那些哥林多人.
或者說應該是哥林多的猶太人.
打到他半死.
但是他都仍然願意留守在哥林多的地方.
抱歉 弟兄姊妹如果是開了.
忘記關了.
就記得關掉那個麥克風.
可以安靜.
就不要騷擾其他的弟兄姊妹.
麻煩關掉.
然後就告訴他們.
要關一下那個.
不好意思.
要靜音.
麻煩.
OK 好.
我們要開回我們的經文.
在那邊.
開不開就行.

$^{681}$行 OK 好.
又要吹一吹.
是不是要縮回這個才行.
有時候很少用這個.
這個要縮回才行.
行得通.
我們繼續.
聽憑主路因為經上記著說.
引述《新靈記》三十二章.
主說失冤在我 我必報應.
這段經文十四章再出現.
《新靈記》三十二章三十五節.
所以你的仇敵若餓了就給他吃.
若渴了就給他喝.
因為你這樣行.
就是把炭火堆在他的頭上.
《真言》的二十五章.
本來很熟悉舊約.
知道神看著整個宇宙.
整個宇宙多麼浩大.
多麼浩瀚.
他都能夠看著.
用他的權能的命令.
希伯來書一章.
托著萬有.
所以每一個基督徒.
每一個有的冤屈.
人所共知的冤屈.
或者人都不知道的冤屈.
神都看著的.
所以如果我們真的有這個信心.
聽憑主路.
我們要學習.
當然要相對.
我們令到人受冤屈.
我們要記住.
我們要尋求和睦.
否則就不行了.
一定會被審判.
這裡告訴我們.

$^{721}$你不可為惡所成.
都要以善成惡.
所以就是那個善.
在我們內心裡面不斷的滋生.
我想這是人生的過程.
真的不同階段.
但是我們讀《羅馬書》.
其實真的相當寶貴.
我們每一個弟兄姊妹.
如果能夠學習《羅馬書》的教導.
真的在我們的生活裡面.
有很大的生命突破.
我們今天集中要講的是第一,第二節.
這個轉折句的重點.
這告訴我們神起源.
就是在這個轉折裡面.
神學的.
從神的憐憫.
帶來的新生活.
就是活祭的獻上.
這裡是.
保羅特別強調.
用神的慈悲憐憫.
來勸每一個基督徒.
得了救恩的基督徒.
無論是神的預先的揀選.
還是我們自己明了救恩.
而來求告的人.
都是在神的眼中.
在這個救恩底下.
其實我們是得到神的憐憫.
你看看有解經家告訴我們.
兩段經文原來是極度相似的.
不過那個相似.
剛剛一個強烈的對比.
我們翻到第一章.
我們很早很早期講的.
和我們今天看第十二章.
一個很強烈的對比.
救恩的原委.

$^{761}$就是神的憤怒要遵土.
在地上已經敗壞了生命的人.
已經是神任憑他們.
我們看到保羅一直講的時候.
是不斷講神那種放棄.
那種對地上的人.
已經是沒有任何的憐憫可言.
第十八節.
原來這是第一章.
原來神的憤怒.
從天上顯明.
在一切不權不義的人身上.
就是那些行不義.
阻擋真理的人.
你想一下只是這一節.
就已經反映了.
我們離開神非常遠.
其實從一出生.
我們就已經和神有一個這樣的敵對.
沒有人知道自己都不知道.
但是原來一直成長.
就是不斷的犯罪.
因為我們內心裡面.
聖經告訴我們.
內心裡面我們有這個罪性.
所以我們就一直這樣來到.
敵黨真理 阻擋真理.
而且過著一些沒有義的生活.
所以接著是十九節.
神的事情.
人所能夠知道的.
原本可以顯明在人的心裡.
但是因為已經墮落了的人.
就不能夠去明白.
去認識和去尋找神.
所以一直告訴我們.
在這個強烈的反差.
就是神的憤怒.
要臨到沒有救恩.
不願意接受救恩的人.

$^{801}$那個敗壞裡面.
這個就是神的憤怒.
但是你想一下.
去到救恩房裡.
那個忠心.
神的偉大.
神的慈愛.
基督裡面的救贖.
接著告訴我們.
現在我們就是在神的憐憫下.
他的慈悲憐憫下.
我們現在就來到去成為.
神喜悅的一個活祭.
就是說本來已經敗壞了.
但是現在在心靈裡面.
因為聖靈進駐救恩.
使我們明白.
我們是一個在神的救恩的愛.
被拯救的人.
原來立即成為神愛的人.
不再是內心敗壞.
而是我們裡面有聖靈.
我們可以同時溝通.
我們知道聖靈和我們的心.
共同證明.
共同印證.
我們是神的兒女.
我們是神的後持的兒女.
接著就告訴我們.
這個新生命.
原來在神的眼中.
每一個信了主耶穌的.
必須是活祭.
記住不是讀神學才是活祭.
不是要奉獻做傳道才是活祭.
那是另外一個差遣.
是另外一個揀選.
不是每個人都有的揀選.
但是每一個信主的弟兄姊妹.
都是由憤怒之子.

$^{841}$成為神喜愛的活祭.
是很清楚的.
在這個神學的轉向.
告訴我們原來留下我們在地上.
由信主到尼西.
是有一個責任.
就是要成為代言人.
就是我們是救恩的活祭.
我們本來是敗壞.
是不欠不義.
阻擋真理的人.
但是現在我們是尋求明白真理的人.
我們要尋求不只是客觀的真理.
明白神在我們生命裡的旨意.
而且要察明察.
要試驗.
要不斷的去了解.
不是只是看一看就算了.
是留心的認真去察驗.
很仔細的東西才叫察驗.
就是我們人生神有一個計劃.
我們要察驗甚麼是祂的旨意.
而且這個旨意是純存神喜愛的.
要在我們生命裡達到的.
普羅用這樣第一節的形容詞.
形容這麼多東西.
所以今天如果我們要告訴別人.
原來你真正的存在.
是一個神的救恩的傑作.
是一個很大的祝福.
人就不會將自己的人生虛度.
因為太多人知道自己原來可以有救恩.
他的人生只能夠為自己而活.
為自己而活還有救友的性情的價值觀.
所以普羅就跟著第九節來說.
新的價值觀應該是怎樣.
還有那個平台應該是在那裡.
所以你可以看到原來救恩可以帶給我們這麼多的發揮.
這麼多的東西我們可以在神裡繼續發現.
發展和發揮.

$^{881}$這個就是對比.
所以我做了這個簡單的比較.
神喜悅的兒女和沒有救恩底下的人.
是任憑他們.
你想怎樣就怎樣.
你最終是敗壞.
你不要以為你最終是敗壞.
你人一世人一世過了你的人生就算了.
不是的.
將來要面對審判.
你的人生已經壞了.
但是還要審判.
但是信了耶穌.
本來壞了的人生.
現在可以重見光明.
現在可以成為一個神喜愛的活祭.
活祭就是說.
特別是這段題的活祭.
就是繁祭.
所以我們稍微看看普羅那種很美麗的比喻.
因為他是猶太的拉比.
他就很明白什麼叫做獻祭.
所以我們如果一路看下去.
我們就發現其實是裡裡外外.
所以我們要有聖經的真理.
所以神學也幫了我們很多.
我們要有聖經的基礎.
我們需要有內心的根深.
我們需要不斷對自己的新的身份.
要有一個好像認證的形式.
你就會慢慢在這種新的身份裡面.
你的行為就會越來越像耶穌.
這個就是很奇妙地.
神喜悅的一個在他的憐憫下的活祭.
不再是憤怒.
不再是任憑.
不再是好像來到.
由得我們繼續這樣越來越離開真理的那種混亂.
而是我們是越來越在神的喜悅裡面.
來得到那個開心,快樂,平安.

$^{921}$所以你可以想像到.
一個小朋友如果從小在家裡.
很受到家裡人的愛,寵愛.
在那個保護,安全感.
一直這樣成長.
你就會眼看的.
你會看到那個小朋友.
無論樣貌,無論他的笑容.
各方面身體的成長.
和一些在屋村沒有人理的.
好像街童那樣.
他經常受人欺凌的.
他就會不斷武裝自己.
越來越眼神.
各方面都充滿了一種的嫉妒,仇恨,保護.
你就會分辨到.
其實一看就看得到.
所以基督徒有沒有心靈的平安呢?.
其實就是反映我們自己有沒有知道.
原來我們一信主就是一個活祭.
所以是神喜悅的活祭.
和一章開始說到.
在沒有救恩之前的人.
那個生命的敗壞.
自己都不斷有很多對人,對實的敵視.
那種眼光,那種內心.
就是這樣一個很強烈的對比.
我們繼續下去.
不單止是信徒要理解救恩.
而是要在生命裡面一個回應.
對救恩的回應.
所以在簡選裡面.
我們要過一個活祭的生活.
這個活祭的生活.
就是從神的憐憫.
作為一個動力.
我們來學習.
不是告訴我們律法對我們的要求.
而是用神的慈悲去感動信徒.
我們要成為一個活祭.

$^{961}$我和大家看看利弊記.
關乎到憲法文祭.
表面上好像是一些祭身.
利弊記應該是第三章.
利弊記.
應該是第一章已經有了.
耶和華 高度有五個祭.
以分類來說.
凡祭,素祭,贖罪祭,贖天祭,平安祭.
五個祭.
耶和華從會幕中呼叫摩西對他說.
你曉遇以色列人說.
你們中間若有人獻供物給耶和華.
要從牛群羊群中獻生畜為供物.
接著就開始說.
他的供物如果是以牛為凡祭.
要怎樣獻呢?.
我們要留心.
因為保羅現在用凡祭.
這個獻祭的過程.
說我們這個信徒.
每一個信徒的內心.
我們這個人對自己的認識.
這樣的一個全人的關注.
第三節開始.
他的供物若以牛為凡祭.
就要在會幕門口獻一隻沒有殘疾的公牛.
可以在耶和華面前蒙月立.
他要按手在凡祭身的頭上.
凡祭便蒙月立.
為他贖罪.
他要在耶和華面前載公牛.
就是那個獻祭者.
阿倫子孫作祭司的.
要奉上血.
把血灑在會幕門口.
壇的周圍的銅壇.
即是會幕外面的銅祭壇.
第七節.
祭司阿倫的子孫.

$^{1001}$要把火放在壇上.
把柴擺在火上.
阿倫子孫作祭司的.
要擺玉塊和頭.
就是那個祭身.
並滋油擺在壇上.
火的柴上.
首先要有那個凡祭壇.
要燒著柴火.
這是祭司的工作.
或者叫利未人的工作.
然後那個祭身就搞定了.
按手.
當然那個手續就是.
要由祭司驗證有沒有殘疾.
無論是明顯的.
或者是暗病的.
那就要驗證一下.
沒有了就按手.
然後就要斬那個祭身.
然後要.
阿倫子孫作祭司的.
這個就是八節.
要擺玉塊和頭.
並滋油擺在壇上火的柴上.
即是整個代表著.
因為玉塊和頭.
和滋油都擺在壇上.
火的柴上.
因為凡祭每天都要獻.
所以那個銅祭壇.
慢慢就叫做凡祭壇.
然後又要.
祭司就要把一切全燒在壇上.
當作凡祭獻與耶和華為輕香的火祭.
雖然全隻要燒完.
但是其實不是的.
因為去到第七章.
告訴我們凡祭身的那塊皮.
即是那些祭身的皮.

$^{1041}$整塊的牛皮.
是歸給祭司作為下欄的.
祭司就可以有皮外套穿.
用皮可以做衣服.
是歸給祭司的.
所以每次要跟一次.
就有整隻的牛,羊,羊皮.
作為下欄的祭司.
然後我們再看第十節.
不過就不再說了.
因為這裡告訴我們.
剛剛留意這裡.
應該是第九節這裡.
第九節這裡就很特別.
因為你既然整隻去燒.
其實沒有所謂的.
你斬不斬都是整隻燒的.
你說我只是想燒某一隻腳.
我就只斬一隻腳.
但現在不是.
凡祭的臟腑和腳的腿.
要用水洗.
祭司要把一切全燒在壇上.
原來那個意思.
就是說原來獻之前.
沒錯是全隻.
由頭和肉和脂肪.
都要整隻獻上.
但是原來獻之前.
要洗那些臟腑.
臟腑即是內臟.
即是說切開之後.
都要洗乾淨.
其實都是燒的了.
這麼猛烈的火.
什麼都燒了還用洗嗎?.
即是說原來裡裡外外都要聖潔.
起初已經驗證了.
第一層已經是祭司.
測驗那個祭身有沒有暗病.

$^{1081}$測驗完之後可以獻祭了.
那個獻祭者就按手.
即是說這個祭身就代表著我奉獻給神.
這樣做了這個手術之後.
我想應該是利未人幫手的.
因為猶太人自己應該不敢這樣殺身.
應該利未人.
即是劊子手差不多.
很純熟的了.
開始去屠宰.
在這個過程要求一件事.
所以保羅告訴我們.
原來我們身心靈都是在神眼中.
成為基督徒之後是聖潔.
神體是他喜悅的兒女.
所以這裡告訴我們.
原來獻梵祭是很奇妙的.
在這五個祭裡面.
大家可能真的不是太熟悉.
這五個祭裡面唯有梵祭要洗乾淨那些內臟.
贖罪祭那些不用的.
如果輕輕的贖罪祭.
只是拿一點點象徵式的獻就行了.
那個祭身.
贖罪祭的祭身.
有些都是可以祭司成為下臨的.
除非是全民族的贖罪祭.
就不准吃.
什麼都不准吃.
平安祭當然不用說了.
平安祭就是神和人的共享祭.
即是祭司可以吃.
獻祭者可以吃.
唯有平安祭.
獻祭的人都可以吃的.
現在梵祭除了其他的祭.
只有梵祭是要洗那些臟腑.
各位我們信了耶穌.
有沒有想過自己要追求生命成長呢?.
有沒有想過自己要察驗神的美好旨意呢?.

$^{1121}$原來舊約利未記第一章.
已經告訴我們.
獻梵祭是要洗乾淨才整隻扔下去.
不是那隻牛一扔進大焚化爐裡面.
不是這樣的.
我都說過我們以前家裡是養豬的.
有時有些大規模的豬瘟.
那些疫情.
人沒事的.
是豬才有事的.
很快幾天就整個豬林都染了.
因為人不吃東西.
整隻豬都弱了.
那些叫做口蹄病和腳蹄.
那些腳會出那些包.
應該是發炎.
就是癢癢的.
自己走不了.
奄奄一息的扔死.
因為那些病毒.
我們拿去.
人們就這樣燒.
整隻要毀滅.
就是這樣.
但是聖經告訴我們.
獻梵祭原來是有這樣的講究.
這個塘祭壇很大.
其實扔下去就燒得著.
但是神要求梵祭要洗乾淨裡面的器官.
然後擺好.
要整隻擺好.
才去燒.
這樣的講究.
其實豈不是像保羅所說的.
要我們心意的更新.
要我們能夠認識自己.
有這樣的新生命嗎?.
所以信徒的心意的更新.
豈不是就是不斷的潔淨呢?.
豈不是就是不斷.

$^{1161}$仍然活在這個世界.
不再屬於這個世界.
耶穌十七章的約翰福音就說了.
我們是在這個世界.
但不再屬於這個世界.
所以更新基督徒的活祭.
首先就是要認定.
我們不再和這個世界有合一.
說的是生命和價值觀的層面.
不是說我們今天還要像保羅的年代.
穿衣服穿古裝.
不是那樣東西.
是我們不斷與時俱進.
但那個價值觀.
世界的價值觀是沒有神的.
是沒有永恆的.
是沒有所謂審判的.
但我們的價值觀就是.
我們是神拯救的活祭.
原先我們都是從舊有的世界.
神憤怒的那個世界.
被拯救出來.
所以我們是離開了那個世俗.
進入神喜悅的那個世界裡面.
所以我們就是立即一定要看自己.
信了耶穌.
很少人會看自己是活祭.
只看自己是一個信徒.
當然這些不是錯的.
不過我們要快些讓頂子妹看到.
其實你信耶穌.
是神揀選你.
是神愛你.
是神用重價將你拯救.
現在你在神的眼中.
因為聖靈已經進入我們內心.
我們已經與神有關係.
所以我們現在就是神的一個活祭.
在地上一個神憐憫的一個見證人.
見證神的憐憫.

$^{1201}$為什麼我本來是憤怒的神的敵人.
現在我們能夠改變我們的生命.
成為神的喜愛的活祭.
所以是要活出神的旨意.
我們繼續再看下去.
剛才我們說了.
下面這裡說.
這是你們理當獻上的敬拜.
很寶貴.
任何一個信了耶穌的基督徒.
他立即會向神敬拜.
打通了與神溝通的管道.
所以每一個基督徒.
就算是一個老人家.
一個小朋友.
都要學習敬拜神.
當然不是說你要懂唱歌.
要有八級鋼琴.
不是每個人都有這樣的恩賜.
但現在說的不是恩賜.
是生命我們與神要有這樣的管道.
對祂有頌讚敬拜.
他說我們這個人的這樣的生活的想法.
就是一個敬拜.
我們每一個人成為基督徒之後.
成為活祭之後.
我們這個人就是一個敬拜.
就是讓人看到我們認識這位造物主.
所以神的憐憫勸我們.
將我們的身體作為活祭.
獻上.
現在我們要主動.
原來神學帶來的就是生活的主動.
就是將自己呈獻給神.
看自己是一個凡祭.
因為我們是人.
我們不是獻祭物.
所以我們這個祭物就是活的.
是一直等待我們神接我們回去.
或者我們完成人生的旅程.

$^{1241}$我們壽終正寢.
我們安然離世.
剛剛說禹堂的媽媽就安息處懷.
她是一個90多歲的老人家.
那時候她星期五.
因為那時候星期五有茶經班.
在大家樓吃飯.
有時候偶爾看到我.
就知道有時候沒有吃飯.
就很緊張.
晚上6點多就在大家樓等我.
一定要看著我吃飯.
幫我買賣.
這些愛主的老人家.
很敬虔的這樣的人生.
剛剛安息處懷.
我們知道原來每一個信徒.
都是在神眼中.
救恩帶給我們的活祭.
所以我們如果已經是的時候.
我們就要更加的呈獻自己.
將自己的身體.
即是說不再是我們自己的了.
我們是離開了永恆死亡的行列.
我們已經來到在神的裡面.
來得到榮生.
所以我們整個生命應該是屬於神的.
所以大前提要想想.
神你想我怎樣呢?.
這個就是權柄的問題.
那個是主次的問題.
即是說神真的有權去用我們的人生.
所以這個就是我們要學習馴服.
這個就是相關.
我們首先認定自己是屬於神的了.
用其他經文就是.
神在萬民當中將我們買熟了.
將我們拯救了.
拯救完之後.
不是你仍然顧我的自己.

$^{1281}$不是說我偶然間回教會.
我喜歡就回.
不喜歡我就不回.
這樣就不是一個活祭了.
我們要盡快讓弟子們知道.
我們是屬於神的.
這個屬於神原來是一件很美麗的事.
因為在神眼中.
祂在救恩看我們.
在祂的憐憫下看我們.
是一個聖潔的.
是祂所喜歡的一個活祭.
而且這個是我們理所當然地理解自己的生命.
而我們要獻上給神的一個敬拜.
我們這個人對神的感知.
就是一個敬拜.
當然在教會裡面有另外的服事.
那些是恩賜上的發揮.
每一個人無論你是甚麼身份.
你就是一個活祭.
你就是一個敬拜者.
你就是一個神喜悅的敬拜者.
你就是一個神喜悅的聖潔的敬拜者.
一層一層地告訴我們.
我們要盡快在生活裡面.
對自己有這樣的認識.
剛才說了.
我們本來就是不會敬拜神.
不會慰藥神.
是不會認識尋求真理.
我們是敵黨真理.
是左黨真理.
是不甜 是不義的.
但現在我們成為了救恩底下的一個活祭之後.
我們就是神的憐憫的一個活生生的一個活人.
所以我們是見證神的憐憫.
見證神在我們這些罪人身上的憐憫.
所以我們應該要充滿了這份神給我們這麼大的憐憫.
我們內心就不要鐵石心腸.
我們內心就要柔軟.

$^{1321}$我們對別人開始會有一份的愛 憐憫.
當然配合了那個智慧.
就不會突然間排隊.
有個人走過來問你.
先生有沒有五百塊.
因為我掉了個錢包 你可不可以借給我.
我把它還給你.
我們要有智慧去跟其他人交往.
譬如聽電話 什麼什麼.
我們要小心 因為太多騙案.
但是我們都要學習神的憐憫.
這裡就告訴我們.
本來這裡告訴我們的就是.
我們在神的慈悲憐憫下.
祂勸我們要有這樣的形式.
勸我們要有這樣的形式.
這個就是真理的透視.
告訴我們原來我們有新的身份.
就是神的愛的兒女.
是一個活祭的生命.
所以活祭是必須我們自己主動的呈獻給神.
我們本來是不會敬拜.
是不會榮耀神的.
現在我們要榮耀神.
我們又要將人都帶到神的救恩裡面.
接著就是那個察言.
這個全人奉獻的敬拜.
就是在外在的生活裡面的選擇.
活祭和世界.
外在生活的選擇.
就開始講第二節了.
就是第二節的經文.
告訴我們原來我們之前是一個拒絕.
頌讚神的一個人.
對神的智慧.
神的救恩.
我們都是拒絕頌讚.
我們在眼中是一個憤怒.
敵黨阻擋真理.
但是現在我們就是一個活祭.

$^{1361}$這就是十二章的第一節.
一定要當自己是一個活祭.
來呈獻給神.
而這個活祭在神的眼中.
我們就是已經蒙喜悅了.
所以我們要不斷察言.
不斷在生命裡面去成長.
比去年或者十年前.
更加認識人生.
更加願意感恩神.
帶領我來到這個階段.
感恩神給我現在擁有的.
我可以服侍的.
不斷有感恩.
不斷有走前一步的獻上.
這就是聖見神喜悅的活祭.
不要十年前和現在完全沒有改變.
喜歡就回到大學會.
喜歡就去奉獻.
喜歡就什麼.
那就不行了.
就是你沒有成長.
沒有不斷的去認識自己.
這是一個活生生神喜悅的祭.
所以一定要和自己的過去做一個比較.
和自己的過去做一個比較.
不要和別人.
和自己.
你就會不斷看到自己的成長.
接著就是第二節.
告訴我們不要效法這個世界.
反而要藉著心意更新.
又是漢語四邊營館.
讓自己被改變過來.
不要效法.
記住神在墮落的世界.
將我們拯救了出來的一個活祭.
所以我們不能夠再和這個世界有一個效法.
反而要藉著心意的更新.
裡面的生命改變.

$^{1401}$讓自己被改變過來之後.
使我們可以察現什麼是神的旨意.
他很詳細的做了一個描述.
漢語營館也不錯.
什麼是美善.
神喜歡.
完全的善.
很特別的翻譯.
舊的文學館就沒有這麼多.
舊的文學館就是說.
不要效法這個世界.
只要心意更新而變化.
心意更新變化.
就是說藉著心意的更新.
心意.
心靈和意志.
這樣的一個.
裡面的內心的思考模式.
內心的思考價值觀.
那個mindset要改變.
然後要讓自己被改變過來.
自己也感受到自己有新的看法.
這樣的生命.
使你們可以去察現.
以前是抗拒.
是阻擋神的旨意.
神的旨意我們根本不會明白的.
根本不相信有這個造物主存在的.
有神的存在的.
現在我們已經去到一種.
這麼立地進至的.
首先裡面被改變.
首先知道自己不再屬於這個世界.
現在當我們自己被改變過來之後.
使我們現在可以察現.
什麼是神在我生命裡面的旨意.
什麼是美善的.
是神喜歡的.
是完全的善.
什麼事是神想我做的.

$^{1441}$是很美善的.
就是說我做這些美麗的事情.
是會帶給自己也有益處的.
做完之後.
神是更加喜歡的.
而且這件事是會帶來一個很完全的結果的.
我們參與了其中的人.
首先裡面自己個人的生命得到改變.
然後我們開始去想神國度的事.
神要在地上做祂自己的工作.
祂要找見證人為祂而活.
證明這個世界有神.
有愛.
有美善.
有審判.
有短暫的人生.
是很短暫.
是很痛苦.
但是有永恆.
所以閉上眼其實是進入永恆的光明的國度.
所以神就要每一個得救的基督徒.
就要先有這樣的裝備.
先有這樣的裡裡外外的更新.
然後就開始可以去察驗了.
所以又是三部曲.
信主.
認識自己的生命.
在這個世界的存在.
然後我們在這個世界原來是有目的地.
被留在這個救恩的裡面.
我們就要察驗神在我生命裡面.
在大家的生命裡面.
當然發生的場景一定是教會.
甚少有人信了主.
決心離開教會.
然後他就沒有了第二步.
而他有第三步.
他就不斷去察驗神在他生命裡面的旨意.
察驗不到.
因為他都沒有客觀的真理的印證.

$^{1481}$完全是他自己主觀的一刻.
沒有了一個平台.
他怎麼可能發揮到神要在生命裡面的工作呢.
是不可能的.
所以一定是在教會這個場景.
讓我們去發展.
讓我們不斷去發掘.
之後最後就是一個發揮.
這個就是神的旨意.
所以就算我們在外面有一些服侍的公貞.
都一定要掛鉤在教會.
我們這個場景.
如果你說我出去做了.
做很多的工作.
為神而做了.
但我不回教會了.
對不起了.
真的不是神的旨意要這樣做.
所以其實真的會令到我們.
有一個更加全面的對自己的教會生活.
是有一個掌握的.
大家試一下.
如果你找到一個人.
他不再回教會了.
但他做事又好又會.
每個人都見到他是神的旨意的代言人.
如果你找到這個人.
你告訴我.
我也很想認識一下他.
是啊 我不回教會了.
我永遠都不會再回教會了.
教會很麻煩的.
接著我是做了神的旨意的工作.
基本上是近乎零.
是不可能的.
所以接著告訴我們.
原來我們發掘測驗.
記住是篩測神的旨意.
是要測驗.
是要認真仔細準備好自己的生命.

$^{1521}$來去真是願意行神的旨意.
這樣的叫做測驗.
這裡就相當寶貴清楚.
兩節聖經已經說了.
在救恩裡面.
我們是一個活祭的生命.
就是說.
祭身是必須呈獻給這位偉大的造物主.
我們的生命如果沒有以神為中心.
必定我們人生是會黯淡的.
不以神為中心的人生.
不可能是一個榮耀的生命.
當然記住.
那個榮耀的生命.
不代表每個人都是很風光很美麗.
有些人的生命.
真是很多人.
香港很多這些院舍.
有些是低度的弱智.
中度的弱智.
因為最近有弟兄.
他的哥哥離世.
在院舍裡面.
由這樣開始十幾年.
這個弟兄就是這樣去不同的院舍.
每個月去報道.
跟這些人講福音.
他有負擔退休之後.
要在這些低度中度弱智的人傳福音.
你想想那種愛有多大呢?.
這是一個車房的大老闆.
你可以想像到.
那種活濟的生命.
就是這樣來將自己的人生燃燒.
是很奇妙的.
所以我們在這兩節聖經告訴我們.
必須在神的憐憫下.
看到自己原來我們奉獻給神.
整個生命要為神而活.
原來人生最大最美好的一個人生.

$^{1561}$這樣用一定是會用你的時間.
用你的精神.
用你的心力.
有時用你的眼淚.
用你的錢.
但是最好.
人生這樣是最有意義.
是最滿足.
為主而活的人生.
那個回報.
儘管這樣說.
那個回報是最大的滿足.
真的.
那個弟兄來過我們江中講見證.
可以想像到一個大老闆.
這麼有心.
你想想誰會去探望這些弱智的人呢?.
他就很想去.
很想怎樣講得明白給他們聽.
那個救恩.
這樣一個人.
所以你可以想像到.
如果我們真的很快認定.
自己是一個活祭的生命.
我們就會想像.
神在你的生命裡有什麼大靈呢?.
這樣想的.
這樣看自己是神有什麼大靈呢?.
原來是將你自己整個人.
發揮到最高的.
你那個人生的目標.
和我們那個崇高.
發揮到最好.
就是神要在你生命裡有什麼大靈.
所以我們要藉著這個心意的更生.
所以我們要在主的裡面.
看自己是一個更生的生命.
讓自己來不斷地改變.
當然這個改變.
就是透過見到不同的人.

$^{1601}$我們上星期洪恩.
如果大家有時間.
上網去洪恩堂聽一聽.
我們上星期的崇拜.
我們也保留了見證嘉賓.
這位姐妹.
她現在是傳道人.
在觀塘的教會.
但她原來坐了監四年.
一點都不簡單.
她的內心.
從小說到她的過程.
但是神的生命改變她.
現在能夠一直讀神學.
她現在在教會裡面讀書.
是很感恩的.
很感人 很動聽.
這就是我們要不斷地讓自己被改變.
所以不同的見證.
不同的講道.
對我們的生命的牧養.
對我們的更生.
對我們不斷認識神的旨意.
我們開始可以.
因為我們這個人裡外都已經認識.
自己是一個新的生命之後.
這個生命是屬於神的.
我們就要察驗什麼是神的旨意.
所以不空凡的.
每個人都有不同.
神會帶領我們的人生.
所以讓一位最有智慧.
最愛我們的主.
帶領我們的人生.
這就是神的旨意.
在我們生命裡面.
或者我們有兒女.
不斷要祈禱.
希望他自己開始學習.
他去尋求神.

$^{1641}$認識自己神兒女這個身份.
小時候從小就知道.
因為一定要回教會.
一定要回兒童主義學.
一定知道.
但是那個知道只是頭腦上的知道.
要盡快讓他發現.
自己是蒙神愛的兒女.
他開始自己去學習.
察驗他的人生.
你不用擔心.
因為神就開始在他的生命裡面.
好像進駐.
好像開始有溝通.
開始有祈禱.
他就自己懂得去學習.
神在他生命裡面的帶領.
我們就可以放心.
甚至可以放手.
這就是我們作為父母.
必須給到最好的屬靈遺產.
就是我們自己去察驗.
他們開始有他們自己的人生的察驗.
而且要想想.
這個察驗不是去想做什麼壞事.
而是想想什麼是美善的.
神喜歡的.
是完美的事情.
你想想如果我們的兒女.
每天除了讀書之外.
想想神在他生命裡面的旨意.
你想想我們還用得著擔心他那麼多嗎?.
他自己就會有一個很大的愛的動力.
他會做神給他力量去做.
還有他知道神要他做的事情.
所以你會看到有些弟兄姊妹.
當他慢慢被神更新改變.
應該沒有那麼明顯.
譬如我們的弟兄.
余建峰弟兄Chris.

$^{1681}$他十月底就要帶隊去土耳其.
接著去埃及的阿里山泰.
這次我不可以加入他們的團.
本來我也很想去了解一下.
那個野地牧人的事公.
就是看著他這兩年生命的艱辛.
就不斷去察驗.
也帶領弟兄姊妹去察驗神的旨意.
保羅是很寶貴的.
告訴我們原先我們看都不用想.
我們是神憤怒的人.
我們抗拒真理.
阻擋真理.
不乾不義.
不斷的降級.
但是現在保羅說.
我們可以開始去想一下.
什麼是美善呢?.
什麼是神喜歡的呢?.
什麼是完美的事呢?.
這樣的神的旨意.
我們去每天這樣去察驗.
所以今天我們這兩節聖經.
其實相當寶貴.
在這個世界羅馬的鬥獸場.
保羅也在這裡.
還有很多的聖徒.
保羅離開世界之後.
離陸王之後.
再有更大的規模的大巴士.
就在這個羅馬的鬥獸場裡面.
一個個的.
我還沒有去過.
有機會去也好.
有很多的閘.
是關著那些奴隸.
放那些獅子出來.
那些奴隸就不斷走.
那些獅子就追來咬.
在上面的那些貴族.

$^{1721}$帝王就看著.
這些高級娛樂.
原來就是殘忍地.
將那些奴隸撕裂.
所以我們保羅說.
基督徒要在這個世界裡面.
活出一個語別.
分別的.
對不起,那個分詞是錯的.
分別為性的生命.
一個的生活.
這個就是我們.
基督徒必須要這樣學習.
心意的更新,變化.
本來我們誰會知道主的心呢?.
十一章近尾那裡.
現在我們從一個認知到一個行動.
就是我們可以去測驗主的心意.
主的心意原來是告訴我們.
一定是美善.
一定是神喜愛.
一定是完美的善.
在我們生命裡面來去落實.
我們做了傳道很多年.
很多的艱難.
有時人與人之間的很多的事情.
但是真的回望這二十多年.
整個生命不斷的成長.
有時做到一些自己都很滿足的事情.
就幫到人.
特別近幾年很多時候.
突然之間要面對的喪禮.
不認識的人找教會.
我們去幫.
真的帶來自己有一份滿足.
幫到這樣困難的家庭.
面對那些哀傷.
所以原先是我們不知道.
拒絕知道神的事情.
神的存在.

$^{1761}$到現在一個基督徒.
保羅說我們可以更新.
可以去測驗.
可以有行動上對神的旨意的落實.
在我們的生活裡面.
所以就是十二章神學的轉折.
就好像先告訴我們.
要有這樣的準備.
接著第三節就說教會生活的一個身體.
要記住永遠是一個身體.
基督耶穌.
我們每一個人在這個世界.
嚴格來說我們有不同的宗派.
但是在神的裡面.
我們是一個的身體.
大家都是弟兄姊妹.
不要因為一些派別.
我們有比較.
或者有一些彼此的敵視.
這樣就不行的.
自己教會就更加是同一個宗派.
更加要有一個學習.
相交相愛.
同心的服侍.
我們就應該要這樣去成長.
今天就我預備了這麼多.
就說到這裡.
和大家有一個簡單的討論.
再一次我們感恩.
因為由十一章的結尾.
保羅告訴我們.
其實沒有人知道神這麼偉大的計劃.
我們根本不能夠明白的.
但是奇妙地.
這個不是收錢上面的裝飾.
而是真的讓我們看到.
原來救恩帶給我們認識神.
和神可以有相交.
我們可以來到去.
生命改變.

$^{1801}$我們是一個很實實在在的活祭的生命.
多謝主給我們有食靈.
多謝主給我們能夠可以去合法地.
隨時的去到施恩補助.
敬拜你.
尋求你.
來滿足你的同在.
帶給我們內心的甜蜜.
和那個的安慰.
這個活祭不單只是自己一個人.
而是我們不斷去和頂智慧一起.
去測驗神在我們生命裡的旨意帶來.
而且保羅還要用三個形容詞.
去描述這個測驗.
就是這個旨意是美善的.
是神喜悅的.
而且是完美的一個旨意.
元氣幫助.
我們準備好.
就在第三到第八節.
下個星期.
我們就會看看教會的生活.
教會不同的恩賜弟兄姊妹.
怎樣能夠互相的欣賞.
互相的配搭.
互相的服侍.
使基督的身體得到潔淨.
多謝主聽我們感恩禱告.
奉耶穌基督的名.
阿們.
\newpage



\section{羅馬書}
\label{sec:yNx82XHOkbk}
\textbf{2025-10-31 周五晚查經班《羅馬書(26):分題:教會是發掘,發展及發揮恩賜的場景》 講員:曾裔貴牧師 1}
\newline
\newline
連結: \href{https://youtube.com/watch?v=yNx82XHOkbk}{\texttt{https://youtube.com/watch?v=yNx82XHOkbk}} ~~~~ 語音日期: 2025-11-03
\newline
\newline
\hyperref[sec:FAt1alHLvdg]{\small{< < < PREV SERMON < < <}}
~
\hyperref[sec:index_chronic]{\small{[返順時目]}}
~
\hyperref[sec:index_scriptual]{\small{[返順卷目]}}
~
\hyperref[sec:B5vFf1eCUBA]{\small{> > > NEXT SERMON > > >}}
\newline
\newline
$^{1}$我們同心抱告.
感恩天父讓我們晚上再一次有機會.
一齊去研讀羅馬書.
也謝謝你讓我們能夠透過現在現今的媒體.
我們可以儲存這些錄音.
可以讓弟兄姊妹在不同的時間.
不同的地方都可以去聆聽.
來去查考神的話語.
我們很盼望神的息靈.
祝福你自己的說話.
當我們來到詳細的.
盡量來還原保羅的意思的時候.
我們盼望能夠建立在那個時代的歷史背景.
還有當時寫作的需要和動機.
也讓我們感受到.
使得保羅對於教會.
對於外邦的教會那份的愛.
還有對於猶太人.
他們不能夠得到福音的那個焦急.
都讓我們能夠感受到.
也讓我們有一個很成熟的心靈.
當我們一路去思考的時候.
我們自己都能夠將我們整個的生命.
都放在主耶穌的救恩裡面.
求你帶領我們.
今天晚上我們等候在你的面前.
其實剛剛在感受整個的社區.
好像瀰漫著一份很開心.
很雀躍的心情.
就是很多的家長帶小朋友.
穿那些萬聖節的衣服裝飾.
可能已經完全沒有那些迷信的觀念.
沒有那些對於外邦.
特別是外國人那種對鬼的觀念.
但是我們都盼望.
不單止只是一個節日.
而是讓更加多人去思想屬靈的生命.
尤其是小朋友.
自小可能就是耳喧目諫.
他們很需要有神的救恩.

$^{41}$透過盼望是幼稚園.
他們一路升學.
遇到的是很好的基督徒老師.
將他們的心靈來改變.
去分逃.
求你帶領我們將這一切都揚忘在你的手中.
我們這樣揚忘禱告.
奉耶穌基督的名.
阿們.
再一次感恩我們又能夠.
可以有時間再去看羅馬書.
我們上一個星期都很快講到.
十二章有三個重要部分.
就是十二章的第一節.
就是很清楚告訴我們.
保羅在十一章他講完對上主.
一切都是從神而來的救恩.
是他發明.
是他去建立.
是他去經營和成就.
所以每一個在救恩裡面的弟兄姊妹.
這裡是很收窄的範圍.
十二章很肯定是對每一個信了主.
領略過救恩的弟兄姊妹.
就應該有一個這樣的呼籲.
所以這個呼籲都是一個很理所當然的擺上.
就是第一個部分.
一到第二節.
就是我們的生命從此.
不再是一個在神眼中的.
就坐墳墓的一個器皿.
而是我們是聖潔蒙神愛的一個器皿.
整個的生命應該怎樣呢?.
這裡就是一個很清楚.
就是要將自己作為一個活祭獻給神.
這是第一個部分.
一到第二節.
很快十二章的三到八節.
就立即講有一個很特定的場景.
這個場景都不用解釋了.

$^{81}$肯定就是在教會那裡來發生.
就是在教會裡面.
怎樣去在這個合一裡面.
來成為一個群體.
群體自然就是要彼此在恩賜上去發揮.
這個發揮過程就是身體的一個成長.
所以第三到第八節.
是緊接著和這個活祭個人的生命奉獻.
和一個整體群體的生活.
是有密切息息相關的關係.
接著到第九節.
就推而及之.
就是我們整個的基督徒.
我們應該怎樣來到在不單只是教會裡面.
更加擴展到我們整個的生命.
應該就是在我們每一個方面.
生活的每一個方面.
就是我們可能在外面社會的買賣.
我們的家庭.
我們和人的接觸.
當時如果有些人是有奴隸.
變成一個新的觀念.
他是羅馬的貴族.
信了主耶穌.
怎樣能夠處理自己的奴隸呢?.
這些又是非常重要的倫理的關係.
所以本來就很有格言式的.
由九節一直到第二十一節.
就去講到方方面面.
這個基督徒應該怎樣.
在整個帝國底下去有那個表現.
所以十二章都很重要的.
將一個基督徒整個人生.
好像濃縮在第十二章這裡去講.
十三十四章就已經有些具體的案例.
應該怎樣去處理一些重要的.
具爭議的議題.
所以這裡就很重要的告訴我們.
我們要從神學的救恩開始.
去到我們的生活.

$^{121}$怎樣去在我們的生活裡面.
參與在神的救恩的工程底下.
所以信了主.
其實是救恩工程的開始.
我們就應該要在神的救贖裡面.
有這個學習.
如果用另一個角度.
這些上一個課都講過了.
我們個人的靈命.
怎樣去成為一個活化了的活祭.
跟著就是我們在教會的生活.
恩賜方面的運用.
我們應該怎樣去處理.
跟著就是愛和人際關係.
那個重要的元素.
在我們內心怎樣去灌注.
這些都是十二章很簡單的分段.
看似簡單.
但真正要實踐的時候.
是相當不容易.
所以你看到第十四章,十五章.
保羅用了兩張聖經.
只是講在教會裡面的人際關係.
講關於安息人應不應該守.
是不是只應該守主人.
講到一些潔淨的食物.
猶太人認為潔淨的食物.
我們外邦人是不是應該這樣去跟隨呢.
這些是沒有大是大非的.
對猶太人和外邦人.
但在羅馬書當涉及到這些生活倫理的處境.
保羅用了兩張聖經來講解.
當然他先講了一些很重要的大原則和方向.
給信徒一個跟隨.
然後就講那些案例.
所以我們遲些去到那些經文.
又可以一起去探討.
不過回到這張聖經.
我們今天集中在第三節.
這個符中道.

$^{161}$就大家一起來思想.
這個道是不知道什麼問題.
可能是網絡問題.
因為Canva是用網絡做的.
我們上一次可能再跟大家傳一點.
因為我們停了一個星期.
好像一個首尾呼應的.
救恩的開始.
保羅立即訴諸於一張十八字.
就是岑任憑那些他極之憤怒的.
已經敗壞了自己生命的人.
所以神的憤怒是從天上顯明.
在那些不乾不易的人身上.
就是說在整個地上.
由我們出生到死.
如果我們沒有主耶穌的挽回.
就是那個救恩.
我們就不能夠面對神永恆的審判.
但是當我們在救恩裡面.
其實保羅這個結論就很清楚了.
我們是一群在神眼中喜悅的.
是可以成為聖潔的活祭的.
這樣來獻給神的.
你可以想像這個首尾呼應.
就是整個救恩包圍著我們每一個人.
整個人在主的救贖裡面.
那個憤怒.
神已經用他的愛詞來清除了.
所以我們今天就在這個救贖裡面.
所以我們現在不再是神憤怒的器皿.
現在我們應該每天活在神的喜悅裡面.
我們要不斷獻上自己.
怎樣獻呢?用什麼方法呢?.
具體的場景在哪裡呢?.
所以你會想到.
保羅是層層地來到.
一路推進.
說了我們現在在神的整個眼光底下.
我們是可以成為一個聖潔的器皿.
我們上一次說過.

$^{201}$一個墳祭是要將它洗乾淨之後.
將它切成一塊一塊.
那些內臟都要洗乾淨.
然後擺好那個形狀.
然後就推到祭壇裡面來奉獻.
這個就是一個祭身的獻祭.
但是神透過保羅告訴我們.
我們不是一些邪教.
是我們每一個人.
我們得到救恩.
聖靈就洗乾淨我們的心.
所以我們是從裡面成為一個潔淨的人.
我們就可以被神使用.
大家今晚有看新聞.
安德魯王子正式被自己的哥哥查理斯.
正式是整個英國的皇室.
褫奪他一切的名銜.
又要勒令他立即搬離皇宮.
似乎是哥哥讓他過世.
他住一些普通的民宅.
那個費用不是由皇室給.
是由他哥哥幫他給.
就是說可能根本不會做事.
沒有收入已經不再是王子了.
但是他連他自己應該有的尊榮.
身份都全部被褫奪.
大家如果一直有追蹤這個新聞就知道.
連寫那個自傳的那個女受害人.
最後都死了自己.
你可以想像得到這件事所牽連的問題.
是非常之大的.
你可以想像得到原來在人類的社會.
不論你是什麼階層.
背後原來有很多的黑暗.
很多的敗壞.
很多的罪惡.
最近有位弟兄他做了三十多天的陪審團.
因為現在已經定了可以講那些案情.
他是親身和其他的陪審團.
聽了整個過程.

$^{241}$最後是七個四男三女的陪審團.
一致裁定那四個被告.
全部都是罪有應得.
一致裁定是有罪的.
你可以想像得到原來在很多的社會裡面.
不論你是什麼年歲.
那些是同黨的犯罪.
如果你看新聞.
因為已經公佈出來就可以講那些案情.
你可以想像得到很多的破碎家庭.
他們做了很多.
因為沒有大人去管.
父母又沒空管.
父母又要謀生.
低收入家庭.
很多很多你可以想像得到那種困難.
但是原來在皇室這麼尊貴.
這麼榮耀的身份.
一樣會有這樣的敗德的事.
你可以想像得到.
在羅馬書告訴我們.
原來在救恩裡面.
我們原本的憤怒已經被更換成為.
神喜悅的一個活祭.
所以如果我們信了主.
而不把自己獻上給神.
好像我們覺得是應份得到救恩.
應該要救我們.
我們其實從來都不明白.
救恩那個真正的含義.
和我們對自己的生命的那種了解.
所以很感恩我們一直看羅馬書.
來到這裡的時候.
保羅告訴我們.
我們開始在救恩裡面.
我們要有具體的回憶的行動.
這個行動是需要我們自己有一個很強烈的選擇.
等下和大家稍微看看第一第二節的原文.
這個理所當然.
其實應該說是一個很理性.

$^{281}$很深思熟慮的一個決定.
通常一個很成熟的人.
作出一個很慎重.
經過深思熟慮的決定.
就是保羅所說的理所當然的敬拜.
所以這裡大家可以一起思想一下.
這些經文我們上一次都提過.
這裡告訴我們.
這個是漢語易門.
所以弟兄們.
我以神的憐憫勸你們.
要把身體作為活祭獻上.
所以記住我們整個的生命.
是先由那個見證.
繁殖就是先洗乾淨.
驗證過是聖潔.
就是沒有殘疾.
這樣才可以獻上.
接著切開了.
切開了即是說.
而裡面的內臟都要洗乾淨.
即是裡裡外外.
整個的身心靈.
我們的思想.
我們的行為.
我們的心靈.
都是聖潔了.
這樣來成為一個繁殖.
記住不是我們一了百了這樣獻一次.
而是終身我們是歸神的.
所以我們不斷不斷來呈獻上.
給神來神所喜悅的一個.
這裡告訴我們.
是神所喜歡的.
是聖潔的.
用幾個去形容詞.
而且保羅最後這裡告訴我們.
這個理所當然是禮當獻上的敬拜.
這個禮當.
其實英文也有的.

$^{321}$有些日本就是rational.
rational就是很理智地去思考.
即是我們中國人說的.
心思熟慮.
即是經過深深的思想.
很成熟的考慮.
即是從各方面這樣考慮.
那個陪審員說.
他們經過了退庭.
大家是很認真地.
不同的角度去討論那個案情.
他說很精彩.
因為大家都是很持平.
不會因為那幾個同黨.
很多案底的.
連那個受害人本身.
都是同黨的朋友.
所以其實是很複雜的.
但是他們都是很理性地.
大家去討論那個案情.
不會因為這些是慣犯.
即犯就好像有一個先入為主的成見.
所以這個叫做心思熟慮.
所以保羅告訴我們.
我們成為末制.
我們要憤外地給其他的基督徒熱心.
其實應該每個基督徒都熱心.
如果我們突然之間.
有感動有熱心.
其實是經過了心思熟慮.
不是曇花一現.
不是偶然間.
偶然間我們能夠慢慢成為一個.
習慣成為一個我們內心的感動.
這個就是理性地的心思熟慮.
這個就是理當獻上的一個對神的敬拜.
所以我們人生的事奉.
其實就是一個敬拜.
整個人生就是一個敬拜.
當然每一次的崇拜.

$^{361}$又是一個很現場式的敬拜神.
這個我們上次看過了.
很期望信徒開始要有一個回應.
就是說救恩是要迫使我們.
這個強迫是神的言文的一個感動.
一個催迫.
我們要有一個生命上的回應.
然後就開始告訴我們.
我們在這個生命裡面.
真的是以奉獻作為一個起點.
愛是我們很重要的原動力.
所以第五章說的第七節.
神的愛 聖靈藉著神的愛.
囂冠我們在我們的心裡.
所以我們能夠好像不斷有一種力量.
從裡面而來的力量.
當然這個就是第四節後面說的.
以天秘恩賜去服侍.
作為一個成長的方式.
使我們活出了和福音相稱的新的生命.
大概這個意思.
重複這樣說.
我們是整個人 全人的獻上.
就是說我們的知識要增長.
我們的情感又要增長.
我們的意志 決心又要不斷的成熟.
不是一時三刻.
是我們整個生命.
投向終生至此對神的毀身.
就是說如果我們越來越成熟.
神感動我們要做甚麼呢?.
我們就要回應那個感動.
去領受 去接受 和去擺上.
最快就告訴我們.
就是要在教會的場景.
來發揮我們的恩賜 彼此的服侍.
這就是救恩的成果帶給我們基督徒.
有這樣的一個參與.
這個我再重溫.
今天我們集中會看的可能第二就是三節.

$^{401}$第四節都未必夠時間再看.
我們先看第二節.
不要效法這個時代.
反而要藉著心意的根深.
讓自己被改變過來.
使你們可以測驗甚麼是神的旨意.
即甚麼是那美善的.
就是將這個旨意來做一個詳盡的形容.
神的旨意所以不是對我們來說是苛刻.
是強迫.
因為保羅的意思是美善的.
是很佳美的.
聽下去都令人感覺到是很美的整件事.
是神喜歡的.
而且是很完美 完全的事情.
就是當我們走在神的旨意裡面.
每個人的旨意都不同.
神安排的 例如今早我和謝牧師.
主辦了一個安息禮拜.
一位弟兄教他幫他洗禮.
他其實好幾十年都要他本身自己.
因為有情緒精神的問題.
來不及出入精神病院等等.
那你敢說是怎樣叫神的旨意呢?.
那怎樣叫做美善呢?.
怎樣叫做我們人都不覺得是一件很美的事.
怎樣是一個完美 是神喜歡呢?.
我只能夠就著我過去大概二十年.
我實際的年期就不是太清楚.
大概十幾二十年.
我對這位弟兄的認識.
因為他而引發了在基督教裡面.
是有一個很大的事公.
而這個事公根本上都不會見報.
根本上都不會有人留意.
因為剛才已經說了.
這個弟兄是長期有精神的問題.
要住在院舍.
基本上是沒有出過院舍.
要回家裡面去住.

$^{441}$由年青住到他剛剛過世.
就一直都在院舍裡面.
不在院舍.
其實可能就是家裡人陪他出來喝茶.
然後要回到院舍.
甚至乎再嚴重一點就要入院等等.
為甚麼我會這樣說呢?.
因為今天假裝安息禮拜的時間.
就讓我們見到這件事.
就是大概十年.
我不敢確定那個日期.
大概這個時間.
你想一下這麼長的時間.
就是因為他的弟弟信了主很熱心.
當年就是20年前.
當年這個弟兄因為他很熱心信了主.
就很想跟傳福音.
有一天過年的年初三.
就跟另外一位我們自己的教會的弟兄.
來到他進來.
他就住在博愛的醫院附近的老人院.
那些老人院是很簡陋.
基本上沒有人理會.
當然有吃東西,有衣服穿,有床位.
但是其他就不用想了.
欠缺的了.
每天就坐在老人院裡面.
好像一個活動的.
又不叫監獄.
活動的地方.
他們就是這樣.
三層的原居民的丁屋.
三層就是這樣生活的.
那一次20年前.
當這個弟兄跟另外一個弟兄一起來探望他哥哥的時候.
那個弟兄就喜歡吹口琴.
就見到很多的院友.
年初三沒有家人來探望.
說的是幾十人.
他就開始吹口琴.

$^{481}$吹那些聖詩.
沒有帶歌譜.
什麼都沒有.
純粹他記得的.
幾十個男男女女就聚在一起聽.
這個弟兄就馬上跟他們說耶穌的愛.
有人舉手勸志.
這個原因.
這個弟兄.
莫弟兄.
見到這樣的情況之下.
心就更加受感動.
更加每個月就進來探望他哥哥.
慢慢就萌生了一個.
由他的哥哥有這樣的精神困擾.
到發現原來不是.
不只是博愛附近的這麼簡陋的老人院.
有這些人的需要.
原來新界,元朗,上水,九龍,香港.
都有很多這樣的需要.
由他哥哥開始.
就引發了他跟這位弟兄.
每個月大家去崇拜.
就開車進來.
慢慢神再感動他們.
就成為了一個很寶貴的時光.
就是全城的時期.
香港,九龍,新界三部大旅遊車.
他們買那些物資.
買那些食物,餅乾和其他的東西.
來給這些院舍的院友.
跟他們講完福音,唱完聖詩.
就派禮物.
那些小朋友,那些院友.
開心到一個地步.
就是有人關心.
我自己有一段時間讀神學的過程期間.
我也有幫忙港島.
連明仔我的兒子.
也有一段時間他讀初中的時候.

$^{521}$也有跟著去崇拜.
去探望這些院舍.
你可以想像得到.
是由這一位剛剛今天安息禮拜的弟兄.
而引發一個基督教的時光.
我自己今天早上在港島的時候.
我就想到神感動我.
這個是不是神的旨意.
透過一個看似不可能有恩典.
是完美的人生.
而引發了有一班弟兄姊妹.
而且是平信徒的運動.
沒有傳到人參與.
他們是這樣不斷有這份感動.
弟兄出錢,出力.
後來甚至乎要成立一個機構.
就是聯合福音團契.
而且因為開始有越來越多弟兄姊妹認識.
有弟兄姊妹是做社署的.
幫忙去問社署.
申請那些資金.
拿到政府的資金.
來幫助這些院舍.
所以我曾經讀神學的期間.
也做過福音團契的一個董事.
後來因為入職我就沒有再做.
可以想像到.
由一個我們看似.
他的人生都是這樣.
受精神困擾.
沒有明天的一個弟兄.
因為他的弟弟.
看到這個一點點的需要和意象.
就引發了一個很大的福音工作.
我們曾經去過沙田.
看到那個國權.
那個唱歌的那個.
蔡國權.
後來起初是很憤怒.
他不來聽報道會.

$^{561}$不來聽福音聚會.
慢慢地他的心又有些動.
因為他自己不能行動.
坐那個輪椅.
電動輪椅來的.
自己過來聽.
據聞他都有舉手決字上樹.
你可以想像到.
我們看神的旨意.
和神真正的旨意.
有時那個闊度.
那種榮耀.
是很難去理解.
所以真是這麼巧.
今早我和謝牧師去辦這個安息禮拜.
感覺到那個離世的弟兄.
是很有那份安詳榮耀.
一方面就是他已經脫離了地上的牢苦.
精神的困擾再不影響他.
那些困難.
身體的疾病再不影響他.
今日特別是化妝師.
化得很漂亮.
那個面容很持長.
很有平安.
真是很大的一份恩典.
所以保羅告訴我們.
我們必須看自己是一個被神喜悅.
拯救的活祭.
我們就會將自己的生命.
擺在神的手中.
來完成神在我們生命裡面的旨意.
我們過完人生.
我們一定會認定.
真是美善.
是神所喜歡.
是完全的一件事.
所以盼望大家都能夠理解.
很快就去到第三節.
立即說原來這個活祭.

$^{601}$不是一個個人的世界.
個人的自我更加中心.
以前我都不理人.
信了主更加不需要理.
我得到救算術.
就不是的.
保羅很快就已經將我們的場景.
擺在第三到第八節.
肯定是教會才能夠去發掘我們的恩賜.
發揮我們的恩賜.
當然要發展.
然後要發揮我們的恩賜.
那個作用就是.
保羅用了一個很重要的比喻.
就是基督的教會.
就是一個身體.
不同的器官.
都是共同的努力.
為了這個人越來越成長成熟.
所以真是很巧妙地.
將每一個來自猶太人的背景.
外邦人的背景.
或者有些是奴隸的身份.
或者有些是當時女性.
男男女女.
主人奴隸.
或者是不同的民族.
都可以在教會這個場景裡面.
不是一個教堂.
是教會.
我們就是神的家.
我們就是基督耶穌的身體.
和器官.
這樣來發展的.
所以活祭絕對不是逃離這個世界.
不是避開這個世界.
所以避世主義.
不是基督教的一個出路.
真正的出路就是不效法世界的價值觀.
而我們是進入世界裡面.

$^{641}$所以大概的意思.
應該就是教會是一個.
我們一班活祭的弟兄姊妹.
的一個訓練的地方.
我們在教會裡面.
由信主到我們受教.
我們去學習聖經.
無論這個學習可以多方面的.
可以是個人讀經.
可以聽講道.
可以有不同的.
在我們前面有經驗的.
屬靈的領袖.
屬靈的長輩.
給我們的新教,現教.
這個是訓練的場景.
所以可以自己一個人躲起來讀聖經的.
一定是會變異端.
一定是會慢慢冷淡的.
所以在教會裡面.
一齊的.
是一個這樣的.
絕對不是絨毛.
絕對是一個身體.
一個是不斷可以成長的身體裡面.
我們來去發展.
所以教會不可以只是同一個階層的一班人.
教會更加不可以只是年青人.
只是老人家.
是不可以的.
教會是應該來自不同階層.
不同的背景.
不同的年齡.
這樣大家彼此去補足.
彼此去豐富.
彼此去幫助.
這樣我們就能夠在神的安排裡面去成長.
接著不是僅限於自己的教會.
而是我們需要.
這個有其他經文的印證.

$^{681}$我們需要向世界出發.
告訴教會以外的人知道.
世界是需要救恩.
萬民需要拯救.
萬民需要福音.
萬民需要耶穌.
所以是我們的責任要這樣做.
所以你可以想像到.
是一個很豐富的救贖的計劃.
不是我們舉一手決志.
一了百了就不再去成長.
這個不是基督徒救恩帶給我們的.
我們是要不斷地生命的成長.
以及在神安排的場景.
去發掘我們自己的恩賜.
以及要發展.
最終是要發揮.
一說到這裡的時候.
我們就會發現.
保羅就很緊張合乎中道.
很緊張第三節說的.
要看自己應該怎樣看.
不用看了.
我們是活祭.
我們是聖潔.
是神喜愛.
那就大聲地去參與示福.
但是保羅有其他經文是平衡的.
例如提摩太后書.
她說初入教的.
千萬不要讓他做監督.
恐怕甚麼.
其實你想到的.
初入教就大熱得發熱身.
就恐怕他自高自大.
就陷入魔鬼的網羅.
即是說魔鬼變成魔鬼.
就是因為他自高自大.
保羅說叫提摩太后.
你不要這麼快見到人做.

$^{721}$很標青.
很熱心.
剛剛信主一個月就做.
就這樣來到不斷讓他做.
而且還是很出名的大明星.
就找他來做一些很重要的崗位.
很快就要收拾這個層局.
因為當有些問題出現.
那個輿論.
教會以外的人嘲笑.
你要收拾這個層局都不容易.
要讓他經驗了一些考驗.
一些挫折.
一些他自己生命都需要調整的心態.
等等.
所以你看看.
不要效法這個世界.
就是這樣.
大概上其實都猜到.
保羅的內心.
為甚麼不要這麼效法這個世界.
不是應該人家這樣做.
我們要跟著跟著這樣做嗎.
保羅說人性.
保羅很清楚.
其實自我中心.
是必須要先有一個對付.
用甚麼對付呢.
學習謙卑.
學習合一.
不要一個自我的中心.
要從對方的角度去想那件事.
跟著很容易就會有世俗的增進比較.
羅馬帝國是一樣的.
人性是一樣的增進.
在尼祿身邊的大臣.
一樣會盡殘言.
所以尼祿自己封建.
一定是他身邊的人討好他.
巴結他.

$^{761}$叫他跟元老院對抗.
所以尼祿就想到.
如果將舊城燒掉.
建一個新的城.
元老院不喜歡他大肆揮霍.
那些人慫恿他.
可以嫁禍給基督徒.
這樣就行了.
所以這些就是增進比較.
跟神的角度裡面的彼此服侍.
所以根本就是完全兩個世界.
兩個角度.
是兩種價值觀.
所以初入教.
如果我們都沒有這種智慧.
作為帶領的人.
見到人馬上熱心.
馬上給他很重的事奉.
一方面就害死他.
因為他都未有機會學習那個謙卑.
學習那個聆聽.
學習那個服侍.
他就已經背負著很重的事奉.
其實不是一件好事.
所以能夠坐在一個領袖的崗位.
其實是經過很多年月日的跌跌碰碰.
很多長輩的議題面明.
很多的經驗.
很多聖靈的教導.
很多不同的考驗.
這樣才能夠達至這種成就.
這樣的一個服侍的心態.
心靈.
所以.
阿哥牧師上個星期的講到的其中一點.
真的是講阿們的.
絕對不是我們每一個傳道人.
各樣的恩賜都有.
各樣的才華,才幹都有.
阿哥牧師往往講.

$^{801}$很多有才華,有恩賜.
有各方面的成就的人.
重點是有很多時候是不被處用的.
以致神要用一些其他的忠心的牧者.
肯肩負那個重擔的牧者.
來去成就神自己的工作.
所以我們服侍的人.
深深體會.
越事奉得久.
越深深體會.
必須要謙卑.
必須要不斷地學習.
不斷地去看到自己的不配.
能夠參與神國度的事奉.
就是這樣.
所以這個彼此服侍的心.
是真的從心底裡面.
看別人比自己強.
要這樣互相欣賞.
互相競爭.
不是比較.
這些初入教的未必馬上領略到.
不是說一定所有都不行.
但大部份都會出現這個問題.
所以太快地坐在高位.
很顯眼的位置.
其實令到他暴露一個很大的事態.
就是心很容易輕飄飄.
很容易將榮耀歸給自己.
忘記了一切都是神的工作和恩典.
這個叫利己主義.
做甚麼.
甚至去到一個地步.
包裝了原來最終都是想自己.
那個我得到人尊重.
得到榮耀.
就不是真正的事奉的態度.
愛和恩典.
才是我們做一切事的背後目的和動機.
否則當我們有利己主義自我中心.

$^{841}$那種明爭暗鬥就很嚴重.
所以保羅告訴我們.
不要效法這個世界.
一個活祭已經超然地屬於神國度的人.
但主耶穌又將我們放回一個場景.
就是人與人之間很緊密的教會生活.
就是一個很好的磨練的地方.
因為只有自己一個人.
你不需要磨練.
但當我們一起的時候.
一班人.
這就是一個很好的學習的機會.
所以第二節就告訴我們.
千萬不要再效法這個世界.
記住不是離開這個世界.
而是要進入這個世界.
但當我們進入之前.
我們是有一個教會生活的.
所以如果我們今天你聽到人說.
我們做回教會有沒有落後.
有沒有老土一點.
聽網上都已經給面子了.
其實不需要回教會了.
我自己就可以敬拜神了.
你聽這些言論.
如果你懂的你就笑一笑.
因為如果使徒保羅.
在當時面對這麼嚴重的尼祿王的逼迫.
他們都不是傻的.
他們都要堅持在教會裡面.
這個群體的.
其實剛剛開始的教會.
基督教都要這樣去凝聚.
大家彼此之間的關係的建立.
使徒們他們給我們的教訓.
正正就是告訴我們.
神在地上.
他還沒回來完拯救贖的計劃之前.
他已經容許教會伴隨著他的工作.
聖靈的工作.

$^{881}$在教會裡面治理一班人.
就是屬神的兒女.
所以誰人如果說自己得到領受.
他可以離開教會.
獨處生活.
將來結神.
這個簡直就是異端.
因為沒有可能的.
有時不出一年兩年.
他已經不再信耶穌.
已經沒有那份感動了.
所以必須要在教會這個場景.
來聚會成長.
服侍.
這個就是保羅告訴我們.
我們要在地上的神的國度.
所以今天神的國度就是教育他.
我們又繼續下去.
活祭和恩賜.
這裡又讓我們知道.
這個活祭不單純是個人的.
沒錯.
第十二章的第一節.
是說每一個得救的基督徒.
我們要看自己的生命.
是神給我們.
透過主耶穌洗乾淨.
我們現在個人.
必須對神有這份委身.
這樣擺上.
原來第三到第八節.
就是告訴我們.
如果嚴格來說.
提到兩次身體.
我們看一看下面的卡卡班.
十二章的這裡告訴我們.
將身體獻上.
第一節中間.
將身體獻上.
這個身體就是每一個基督徒的生命.

$^{921}$個人的生命.
但到第三節這裡.
留意這裡.
我憑著所賜我的恩.
對你們各人說.
不要看自己過於所當看的.
要照著神所分給各人信心的大小.
看得合乎中度.
接著就立刻再說那個身子.
正如我們一個身子上.
當然這是一個比喻.
有好些肢體.
肢體也不都是一樣的用處.
留意第五節.
我們者許多人在基督裡成為一身.
互相聯絡作肢體.
也是如此.
我們在基督裡成為一個身體.
所以第五節說的一個身體.
再不是個別基督徒.
而是每一個個別的基督徒.
集合成為耶穌基督的身體.
所以很清楚.
就是個人的得救.
要將個人的得救.
放在整個基督教.
整個耶穌的救恩裡面.
我們才成為真正的基督徒.
如果我們個別.
當然沒問題.
個別舉手決志.
甚至在家裡自己祈禱.
信耶穌是沒問題.
但是祈完禱.
決志就要開始在教會生活.
就要開始在神安排你的教會裡面.
有聚會.
有敬拜.
有學習事奉.
和生命的成長.

$^{961}$否則我們只是個別的活祭.
而我們從來沒想過.
要將自己放在基督的身體.
在這個身體裡面.
一個有機的身體.
就是會不斷成長的身體裡面.
這樣來成為一員.
而在神的聖靈的安排下.
祂會將恩賜.
將祂的旨意.
放在大家的心靈裡面.
以致當我們去發掘.
當我們去發展的時候.
最終我們就能夠.
可以在神安排的場景.
去發揮我們的恩賜.
這樣就叫做一個合一的基督的身體.
所以是由個別的個人的得救.
進到教會的生活裡面.
所以獨行俠的基督徒.
注定的只會失敗.
我不可以說注定滅亡.
因為信主求告主名就必定會得救.
但是記得不是一次的求告.
是一生的相信.
是一直的相信才可以得救.
不是一次得救就永遠得救.
是一次的得救.
一直的相信.
這樣才是得救.
所以教會就是令到我們的新聖靈.
不斷有一團火.
不斷去認識自己是活祭的生命.
不斷知道自己原來.
是有一個很美麗的人生.
就是在我們的人生裡面.
包括我們的成長.
興趣,工作,家庭,婚姻.
我們要走在神的旨意裡面.
而神有祂的美善.

$^{1001}$有祂的帶領.
這樣的人生才算得上是活祭.
所以很快的就說到一個集體性.
當他要說到這個集體性的時候.
我要和大家看看希臘文.
第二節的合乎中道.
或者現在第三節.
原文的意思.
這個我已經不是很會讀了.
內契形容詞的那個修尾.
KN.
用英文去做希臘文的符號.
怎樣收走這個.
這個怎樣拿走呢.
為什麼整天都在動.
這個希臘文.
第一個肯定是和那個logos.
那個理性有關.
那個logos那個詞根.
不過是形容詞.
後面這個partia.
就是那個關乎敬拜.
兩者合謂.
就是合理的敬拜.
就是說在神的慈悲的面前.
將自己獻上最合和七和.
最自然的一個回憶.
就是在整個神這麼偉大的慈悲.
憐憫底下.
我們真的認識自己原來本來就是神的憤怒之子.
我們本來就是在神的眼中.
我們已經是不填不二的人了.
我們現在能夠看到.
神這樣收納我們.
喜悅我們.
其實是因為耶穌的緣故.
但是這個這樣的結論就很清楚了.
理性的敬拜就是.
我們應該很滿足於.
我們能夠被神拯救.

$^{1041}$能夠被神使用.
我們能夠服侍.
所以將自己獻上的人.
其實就是最合和七.
最深思熟慮.
是最自然的一個回憶.
如果不是這樣回憶.
就變成了一種.
好像我們得到恩典.
就離開.
不曉得感恩的人.
好像那十個大麻風的人.
當他們醫好之後.
他們找完祭司.
其實都沒有找祭司.
走到路上已經好了.
當他們見到自己好了.
可能找完祭司之後.
只有那個萬大麻風.
但是那個撒瑪利亞人.
就是說本來已經慘了.
再加上他的身份.
就是撒瑪利亞的身份.
就更加被其他人歧視.
但是他仍然帶著這種過去.
那種舊有的.
那種看自己的無助.
看自己的卑微的身份.
反而看自己是一個.
蒙主耶穌恩典.
醫治的一個人.
誰知道原來耶穌不單止.
醫治他的身體.
連他的過去.
被人歧視的民族身份.
耶穌都醫治了他.
他是神的兒女.
可以平平安安地離開.
其他的猶太人是慘的.
有大麻風.

$^{1081}$但是當他好了之後.
他自己就是猶太人.
再看不起那個撒瑪利亞人.
可能已經立刻離開.
跟那個撒瑪利亞人來往.
因為已經康復了.
這裡告訴我們.
原來一個願意擺上自己的基督徒.
原來是經過深思熟慮.
很合乎邏輯.
很合乎那個理性.
深思熟慮的理性.
很自然地回應神對我們人生的呼召.
因為這是神的慈悲.
所以希臘文.
因為我都找了AI幫忙.
請你幫我把希臘文用中文翻譯給我聽.
做了這一段解釋.
所以現在看希臘文.
可以很簡單.
不如你去看希臘文聖經.
網上有這些軟件.
但那時候我們讀神學.
我們要買一個希臘文的.
其實不單止希臘文.
叫做原文聖經的軟件.
那套軟件二千多元.
當年.
十幾年前.
已經那時候很貴了.
是神學生先童.
你們基本上不用用它.
你要將這段經文.
人家就會幫你分析.
做出來.
所以現在AI真的很方便.
但是這裡的解釋就很清楚了.
活祭的意味就是告訴我們.
這句希臘文的意思.
所謂理所當然的敬拜.

$^{1121}$就是在和諧本這樣翻譯.
所以敬拜不再是獻上動物.
而是獻上自己.
獻上自己就是說.
出於理性屬靈.
全心全意地回應神的大愛.
所以盼望我們都不斷看自己.
是一個這麼尊貴的身份.
就是一個被神拯救的一個人.
而這個人可以成為一個獻給神的活祭.
所以我們不斷去參與事奉.
其實在神的眼中.
在人的眼中.
都是一份恩典.
輕香.
當然那個過程是有疾疾的.
所以為什麼我們不能初入教.
馬上有很吃重的事奉呢?.
就是因為我們都很容易會自我中心.
是神不斷地會調教我們.
聖父會不斷地在我們內心提醒我們.
有不同的挫折.
使我們更加地願意獻給自己.
所以面對挫折.
不是退出不事奉.
而是思想反思.
有什麼做得不好.
就有些修正.
更加擺上的事奉.
這樣才是一個活祭.
所以每次我們都洗乾淨.
用神的愛.
用主耶穌的寶血.
洗乾淨我們的動機.
我們的內心.
再參與事奉.
這個就是一個活祭.
絕對不是一有挫折.
一有人的論斷批評.
就立即什麼都不做.

$^{1161}$這是合理的.
世界就是這樣.
因為我是義務.
義務就是我不做就不做.
不要這樣想.
應該想你是神眼中的活祭.
是獻給神的.
所以不是隨便就可以收回就算了.
反而是真是一個理性的.
理所當然.
是經過邏輯的心思熟慮.
這樣叫做理所當然.
就是這個字.
Logos的詞的詞根.
我們繼續下去.
所以活祭不是一個孤立的靈修.
而是進入教會的身體的服侍.
每個人都有不同的恩賜.
是神按信心所分配.
目的不是為了建造整體.
而是恩賜的使用.
也不是炫耀或自我的實現.
而是謙卑地在身體裡.
發揮功能彼此的配得.
榮耀神.
所以不是建造我自己這個人.
而是我們這個人.
擺在整個身體裡.
不是要炫耀自己.
而是盡我的才幹.
盡我的能力.
使整個身體發揮功能.
彼此的配得.
所以整個身體的器官.
是互相的協助.
互相的協調.
當你眼見到一件事.
那隻手就會有行動去做.
是一個大腦驅使的行動.
你可以想像到.

$^{1201}$是非常完美的結合.
正常嚴格來說.
一個越是對人體了解的醫生.
或者無論外科內科也好.
精神科也好.
都應該對人體的結構.
不斷的讚嘆.
因為他們是專家.
應該更加了解.
那個結構的精美.
每一個裡面的機能的配合.
是天衣無縫的.
整個人體.
所以稍微有些出問題.
無論是精神方面的.
器官方面的.
先天的後天的.
已經要我們盡辦法去解構等等.
但如果一個正常的人的成長.
每一個器官的互相協調.
是很奇妙.
是很精美的.
所以你可以想像到.
神透過保羅告訴我們.
教會就是一個耶穌基督的身體.
這個概念太少的兄弟會懂得.
所以他只是不斷的看自己.
是教會的一個會員.
他不是看自己是耶穌基督的會員.
這樣就帶來.
有時就不知不覺.
就會有自己的看法.
價值觀.
自己的自我的中心.
或者利己的主義等等.
很容易就會引發很快很快的衝突.
但是耶穌基督的身體.
就是一個這樣.
所以這段經文的分析都非常重要.
因此不是去炫耀.

$^{1241}$不是要自我的實現.
例如我彈琴很厲害.
我要在教會做私琴.
嘩好啊真是好啊.
但原來他做私琴.
彈得很巧妙.
登峰造極.
原來是想自我的實現.
這樣就不對了.
他整個人坐在那裡準備崇拜.
是要令到回眾.
是要感到回眾敬拜神.
他是代表著好像天線去彈琴.
這樣而已.
是要引導回眾.
一起用歌聲.
用心靈去敬拜神.
所以他只是其中一個小小的器皿.
但當他沒有人知道.
他自己內心的自我的實現.
這樣就糟糕了.
因為他在炫耀自己的技巧.
這樣就不對了.
等於港獨的人一樣.
只是口沒遮攔地說他的口才.
這樣就不是一篇港獨.
那篇港獨是要感動人心.
提醒.
是喔.
這個提醒是適合我現在面對的困難.
或者試探.
這個提醒就是剛剛滲透到他.
向我講說話.
所以港獨不是一個分享.
港獨崇拜的港獨是一個正道.
是神透過那個傳道人去見證神的道.
可能神只是希望某一個人聽的一篇道文.
所以這裡告訴我們.
這個就是先看了很重要的價值觀.
保羅不是太在乎那些其他的恩賜.

$^{1281}$反而擁有恩賜的人.
有什麼價值觀呢.
才是最重要的.
所以開始他就講那個合乎中道的概念.
一切原來都是為著要榮耀神.
這樣來告訴我們.
我們繼續下去.
我們開始講第三節這個合乎中道.
又是漢語聖經逆叛.
希望給大家多些角度去看.
這一節經文的意思.
舊的和諧館就是第三節.
我憑著所賜給我的恩.
對你們各人說.
不要看自己過於所當看.
要照著神所分給各人.
信心的大小.
看得合乎中道.
就是說保羅的認知裡面.
肯定有些人信心大的.
有些人信心軟弱一點的.
甚至關乎倫理的.
那些信心保羅都有提的.
在哥納多前書就明顯.
有幾段經文提到.
有些人的良心還是軟弱.
他們就不敢去吃那些製藥醬之物.
所以如果那個傳道者.
或者很成熟的基督徒.
他自己去廟裡面坐靜.
就會令到那些信心弱小的弟兄.
就被叛奪了.
所以保羅說如果我吃肉.
叫弟兄跌倒.
我就永遠都不再吃了.
不是不吃那些普通的肉.
是那些製過偶像的肉.
最近在網上看一些很有趣.
很解壓的那些好笑的.
那些什麼臥底旅行團.

$^{1321}$我發現了一件事是很恐怖的.
就是很好笑的.
你放鬆地看.
是很有趣的.
但每次進酒店.
第一個動作你知道是做什麼嗎?.
有看的你就知道了.
那兩個主持人.
進到酒店.
拿著鎖匙.
要敲門.
這些動作是告訴我們.
我們相信靈界是有髒東西.
對不起我們.
你麻煩你走開吧.
沒有這件事的.
你相信耶穌當然我們沒得說.
但如果你告訴我.
我們是相信耶穌.
麻煩你走開吧.
我根本不相信他們會搞到我們基督徒.
千萬不要這樣.
你這樣的時候很嚴重的.
就是說不好意思.
我來你們的地方.
你麻煩不要搞我.
今晚你讓我睡一覺不好.
沒有這件事的.
你回去看看.
這樣的動作是代表著.
你已經怕了靈界.
其實不需要.
當然他沒有耶穌.
他們有這種迷信.
那種迷信已經塞入到.
香港人的那種價值觀.
你可以說多麼嚴重.
所以基督徒不需要.
我們甚至要告訴他們.
我們是沒有任何禁忌的.

$^{1361}$我們要去旅遊.
旅遊是要付錢的.
理直氣壯的.
不需要做這件事的.
不需要開燈.
不需要什麼方位.
什麼尾房.
沒有這件事的.
我們是基督徒.
我們是神的兒女.
所以這裡告訴我們.
我藉著神賜給.
這個就是漢語譯管.
我藉著神賜我的恩典.
保羅所領受的恩典.
是極其的浩大.
一個本來是拉給.
一個在社會上.
猶太教裡面的一個.
很尊貴的一個議事.
一個能夠去行審判的一個議事.
即是投票權.
是不是可以去懲罰一些人.
在宗教的條例上.
保羅他能夠信主.
能夠本來是逼迫基督徒.
嘲笑侮辱耶穌的.
他可以領受神既往不計的那種罪惡.
來賜他恩典.
還賜他一個職份.
而且明顯是使徒的職份.
這樣來對你們羅馬人說.
不要對自己的評價過高.
即是意味著.
合乎中度.
即是不要過高.
更加不要過小.
超過所應該的評價.
所以這個超過所應該的評價.
即是說按部就班.

$^{1401}$循序漸進.
我們每個人得到的都是聖靈.
都是一份的救恩.
都是耶穌為我釘十字架.
每個人都一樣.
原來在教會裡面的生活.
就開始不一樣了.
就是神會在不同的弟兄姊妹裡面.
賜下那些恩賜給每個人.
所以每個人有些就明顯些.
有些就需要運用大些的信心.
但是原來每個人整體都是神拯救的兒女.
神用寵家.
他的愛子的寶血來買熟的兒女.
但是在那個實際彼此的服侍上.
就有不同的那個.
因為信心.
因為神給的那個恩賜.
就會有高低之分.
那個高低不是一種道德上的.
即是有比較.
而是神安排我們不同的人生.
要告訴我們.
不要對自己評價過高.
即是說不要再從自己的角度.
去看自己的價值.
要在神聖靈的帶領感動底下.
去看神給我們應該有對自己的評價.
即是說不再用自己去量自己.
是要在神的真理去量度自己.
這樣就完全不同了價值觀.
在神眼中看自己.
在永恆裡面看事奉.
幾十年就這樣很快過去了.
我不斷提我是1987年12月24日.
平生第一次去一個家長小學晚會.
不是為我預備.
根本沒有這樣的準備去這個報道會.
是臨時這樣裡面有個感動.
叫我衝上去.

$^{1441}$衝上去濕透了.
然後就已經是一半的枯燥了.
舉手就成為了我由那個時刻到今天.
就真的是一個人生三十多年.
就是這樣順著的人生三十多年就這樣過去了.
所以在永恆裡面看自己的評價.
這樣看自己能夠參與事奉.
是忙的 是重責任的.
最好不要自己再做這些重的責任.
希望快點退休.
找一間沒有人認識我的教會.
去做一些很小的弟兄的角色.
這樣很享受敬拜.
就是我很想的心願.
但是在永恆裡面看.
能夠參與事奉是很尊貴.
其實你不參與事奉才是不健康的.
不理性的 不心思術類.
但就不是從自己去評價自己.
而是從永恆 從神的眼裡面去看自己.
所以你這樣坐在這個崗位去服侍.
冥冥中總是有神的帶領感動.
所以倒要按照神的良比各人的信心.
所以是神用他的智慧 永恆的計劃.
他這樣度身訂做給我們.
這樣給我們真正為著.
不是想做好那件事那麼簡單.
是做那件事的人自己會成長 成熟.
我都真的坦誠分享.
張牧離開2017年9月.
他轉去建道教書.
我就要承擔一個很重的堂長的崗位.
但是7年過去了.
我現在回想.
我經常說 張牧你不要找我了.
我不是有這樣的喜歡.
不過好像真的 我當時都這樣說.
好像臨危受命那樣.
但是7年過去了.
我坐在這個位置對自己的生命.

$^{1481}$對自己的成長.
現在是祝福雄恩.
我有很多不無盡的創意.
那種曾經有的經驗.
有的各方面的人物的建立.
可以令到我現在可以祝福.
在雄恩這個小小的教會.
我們的分堂.
我覺得是神的準備.
訓練了我7年.
而離開在雄恩的發展.
給我肩負很大的重擔.
令我真的呼吸不了那麼辛苦的壓力.
但是很感恩.
真的是神良美國人的信心.
所以要忠肯地評價自己.
所以每件事過去了之後.
是很大的一個祝福.
如果我沒有坐在那麼吃重的位置.
開始經歷社會事件.
開始經歷疫情.
我會有很多的困難.
但是可以真的令到自己的生命成長.
成熟.
能夠服侍.
是一個很大的祝福.
所以不是做到多少那個重要.
是做事的人才得到成長.
才是重要.
所以有些長輩.
以前很多年前提醒你不做.
神就將你bypass了.
會祝福其他的人.
他做到.
他將你本來得到的福氣.
你就被人拿走了.
所以我們不是要爭那些事奉.
這個就是剛才說的價值觀.
但是當神給我們.
其實是困難的.

$^{1521}$當你參與事奉.
但是忠心地.
來到好像剛才所說的一個傻子.
去參與去服侍.
其實默默地這樣做.
真正的祝福就是神會訓練我們的生命成長.
會使我們更加看到.
原來服侍是那麼尊貴,那麼寶貴.
又說一個很簡單的小小的例子.
上個星期六.
去探我們洪恩一位姐妹.
因為神有些發言.
本來她已經在去機場.
去韓國旅行.
去機場.
突然間痛得不得了.
緊急要入急症.
就進了北區醫院.
輾轉幾天.
又忙 沒空.
到星期六去探她.
原來她已經和隔壁的床位.
原來後面發言.
只不過是小小發言的酵素高.
就要慢慢打點滴.
自己就會好.
就在醫院裡面.
就沒事了.
知道和其他器官有問題.
只需要慢慢調理.
繼續在醫院裡面.
隔壁床一位女士.
年輕很多的.
她是有癌症末期的.
她和我們的姐妹說.
你可不可以為我祈禱.
我很痛 我很辛苦.
姐妹真是神安排.
去不了旅行.
就在醫院裡面.

$^{1561}$就和她簡單講福音.
跟著在WhatsApp裡面說了.
這位女士肯信耶穌.
就為她祈禱.
我看到 但我都忙.
星期六不太忙.
又有其他的活動.
上午就有幼稚園的典禮.
剛剛有個空檔.
未到下午兩點半.
我們有一個福曲的22週年的慶典.
剛剛有一個空檔.
由北區醫院可以12點探病.
因為屯門現在不准.
剛剛有個空檔.
就在錦繡開車過去.
就見到這位女士.
就說 女士不如你這麼清楚地信耶穌.
我現在和你做一個更加認真的決戰.
然後立刻幫你洗禮.
她說好啊.
她很感恩.
之後WhatsApp我.
她家人在這裡.
見到我立刻幫她洗禮.
家人說好啊.
因為是不能再回頭的絕症.
她自己知道.
而且痛得很厲害.
所以很多事情.
神給我們的工作.
特別我們作為牧者.
去到這些階段.
真的將以往幾十年.
電影報告成長.
越來越有經驗.
自己原來是在神的恩典裡面.
被你養大的.
所以現在.
時報不會再介意人怎樣看.

$^{1601}$總之就是能夠做到.
就是神的祝福和恩典.
我現在又帶了一個女士.
真是素面謀面的.
很快就能夠一句立刻幫她洗禮.
她已經轉去百寶利的靈陽中心.
很快就要離開.
遲一些過兩天.
我就會將那個洗禮證帶給她.
因為已經做好了那個洗禮證.
她是我們神的兒女.
她很快就會見主.
但是你可以想像得到.
我們如果不參與事奉.
是很輕鬆的.
是很好的自己的人生.
當然每個人真的有不同.
但是如果神真的給你一個服侍.
就要忠於所託的.
去盡心做好它.
現在紅恩的弟兄姊妹.
真的很愛惜我.
叫我不要那麼快.
不要衝.
不斷地從這裡做到那裡.
他們一直跟著我.
他們很開心.
他們覺得牧師兄.
麻煩你休息一下.
但是.
總是內心有很大的感動和開心.
譬如昨天.
又見到一群老人家.
好幾個從來沒有去過香港.
石澳,赤柱.
不如秋天.
我們團契去一次石澳.
探我的同學的教會.
拍照.
然後再去赤柱.

$^{1641}$吃點東西.
在赤柱逛一會兒.
老人家開心到不得了.
但是他們很累.
因為走不了那麼多.
開車我更加累.
回來整個人都散了.
但是內心有很大的滿足.
所以要在永恆裡面.
去看自己的時況.
不要自己評價自己.
如果我自己評價自己.
我寧願不做傳道.
如果中國銀行可以再投胎.
就不要再傳道了.
千萬不要傳道.
吃力不討好.
但是如果在神的永恆的眼光去看.
既然現在是神驗證.
是神揀選.
是神感動.
那就盡力去做好它.
我每一次講我自己個人的分享.
他講張牧師真的是我自己一個很好的恩師.
他不是給我一個陷阱.
是真的給我一個很好的祝福.
讓我能夠成長.
對自己對事奉有更深的體會.
所以真的.
如果我們有事奉.
不要輕易的放棄.
應該要更加咬緊牙關.
參與服侍.
我不知道聽這個經文的荷蘭.
或者外國弟兄姊妹.
在教會裡面會不會面對一些不同的挫折.
每個教會都有的.
但是真的要合乎宗道.
所以這裡說.
忠肯的評價自己.

$^{1681}$記得是在禱要按照著神.
給我們個人信心的那個大小.
忠肯的評價.
所以不要自己評價.
我們一想到思前想後.
做什麼我要這麼辛苦.
做什麼我是義工.
做什麼我要被你這樣仇氣.
一做什麼就不想做.
往往很少看神的恩典.
神的感動.
所以我自己都講過.
我很喜歡寫下每一個的婚禮.
和安息禮拜.
特別是安息禮拜.
每一次安息禮拜我都會做一些記錄.
很多年我看回.
有時候很感恩這個家庭.
曾經在他最困難的時候.
我們教會怎樣幫過他.
繼續為這些家人禱告.
原來是很開心的一件事.
你會感覺到一份滿足感.
你能夠幫助一些困難的家庭等等.
所以我們要在這個合乎宗道裡面.
我又再查過希臘文.
這個修辭.
意思就是謹慎地思想.
自制地去判斷.
剛才所講的很理直氣壯.
很經過深思熟慮去了解.
是一種清醒謙卑.
合宜的自我的形式.
和剛才的理解差不多都是一樣.
不要過於自己所當下.
但在這裡講的合乎宗道.
就是有這樣的含義.
所以告訴我們.
我們原來一個事奉.
不是隨便自己喜歡怎樣就怎樣.

$^{1721}$不開心就不做.
不是這樣的.
而是謹慎地去思想.
是從整個宏觀.
這個事奉本身.
如果我能夠做得好.
是可以榮耀到神.
可以祝福到一些家庭.
我吃力一點.
辛苦一點.
其實是沒有所謂的.
慢慢地.
當你不再介意那些個人的得與失.
就更加從容.
就更加地去輕鬆治癒.
因為你看的是事奉.
和看的是神的工作.
所以我們要經過了.
現在沒有人掘井.
我們以前在家裡錦田.
我們要自己掘井.
粉嶺也是.
掘那些大口井.
頭十多呎.
你掘完泥.
掘完石.
它都沒有水.
再掘再掘.
慢慢就有些泉.
有些水出來.
地下水.
所以告訴我們.
原來參與事奉這個合乎宗道.
都是不斷要自我的反思.
不斷謹慎的思想.
強調是一種很清醒的.
很謙卑的.
很俠義的自我的形式.
就是合乎宗道.
就是我們能夠繼續.

$^{1761}$可以事奉的原動力.
所以求主讓我們找到我們的恩賜.
找到神要我們去服侍的那個場景.
那個對象.
以致我們就擺上我們自己.
神啊.
我是一個你眼中的活祭.
我現在就這樣奉獻給你.
在你安排的那個場景.
我就會忠心的去服侍.
就是這樣的意思.
所以剛才提過.
第一到第四節.
其實應該是第五節.
身體的獻上是個人的.
每一個基督徒都是個人的一個活祭.
明白救恩.
被神拯救.
身體.
心靈都是屬於神.
我們再獻給神.
是很理性.
很俠義的.
但是我們是要加入一個大的群體裡面.
就是基督的身體裡面.
我們要將自己放在這個身體裡面.
是其中一員.
所以不是一個教堂那麼簡單.
是整個的教會.
所以這樣的概念就更加偉大.
因為不再是錦繡堂.
不再是蒙恩堂.
不再是香港的教會.
元朗區的教會.
在神的眼中.
我自己就是古往今來.
每一個重生基督徒的地方.
就是我的教會.
這樣理解的時候.
就會很不同.

$^{1801}$所以也不會眼紅.
那個教會突然間那麼服氣.
我們可以參考.
但是不要需要眼紅.
因為他也是為神服氣.
會服侍.
我繼續做好我自己安排的.
可能是一個鄉村的教會.
細小的教會.
一個可能老齡化的教會.
好好地去做好我自己的部分.
去中心地服侍神的群羊.
大家都有大家的中心.
所以如果這樣看的時候.
神的教會根本就從來不需要競爭.
反而互相的祝福.
互相的合作.
才是那個重點.
所以現在我們同恩堂很感恩.
就是我們可以和生命堂有一個這樣的聯繫.
有一個這樣一起去做一些時工的服侍.
大家都不分彼此.
這樣分享那些資源.
是一件很美麗的事.
所以是一個群體的建造.
今天大概就講了這麼多.
剛剛好55分.
我們今天就先停在這裡.
今天講了第三.
第四節都還沒正式開始講恩賜上.
多樣性的恩賜方面.
這個可能要下一個星期.
但是要講一講.
因為下個星期.
我星期天晚上就要飛土耳其的伊斯坦堡.
因為我會參與SAT7.
大家可能都聽過SAT7.
就是金峰醫生他們的組織.
就是用人造衛星.
向回教徒的世界傳福音.

$^{1841}$我們會去開會.
還有我們有另外一個教育部.
筆書的.
我們會參觀他們土耳其的製作中心.
所以都很值得去開銀介.
我就要下個星期才回來.
所以星期五晚的茶經要暫停.
所以盼望大家為我祈禱.
為SAT7的國際性會議祈禱.
希望這個時光能夠讓更多阿拉伯世界的地方的人.
回教徒來去信主.
同一時間都要為Chris.
他們在這個野地牧人.
他們現在就在勒班縣.
遲一些他們就會去到亞歷山派.
那裡繼續去傳福音.
希望大家都在土報裡紀念.
多謝大家一起土報.
肯定我們感恩.
我們今天學習了很重要的真理,秘訣和原則.
是每一個基督徒都是一個活祭.
都是一個在神眼中.
其實是神已經喜悅的人.
所以求主幫助我們.
要從這個成為我們的起點.
在神的眼中的喜悅.
去參與事奉.
絕對不是神憤怒的器皿.
根本就不可能去事奉.
多謝你讓我們能夠在耶穌裡面.
我們成為你的兒女.
而且不只是成為一個身份的改變.
而是神很想我們每一個的兒女.
參與在教會裡面服事.
多謝你透過保羅給我們有很重要的教導.
就是我們要合乎忠道.
是要很認真地看自己.
需要很有神的恩典.
從神的角度去看我們所領受的信心.
那個大小和所得到的恩賜.

$^{1881}$盼望主幫助我們.
因為其實每個人都有自己的才幹.
有自己的背景.
有自己的恩賜.
有自己的優勢.
有自己不同的方面.
但是我們很盼望.
我們千萬不要在世俗的洗腦底下.
我們用一種很負面的眼光去看事奉.
反而我們不是去爭去競爭去比較.
反而我們是很渴望.
能夠成為神喜悅的器皿.
去成為別人的幫助和祝福.
求主幫助我們.
裝備我們.
也願意我們自己.
在神的手中.
被神裝備.
被聖靈去更新.
使我們能夠在不同的地方.
為主作美好的見證.
願你幫助.
也能夠甦醒我們大家的心靈.
一起來去想.
用心的.
用思想去想.
我們有什麼地方可以被神大用的.
以致我們不枉費我們的人生.
來做一些可能在神行中的次要善情.
與神幫助.
我們這樣感恩禱告.
奉耶穌基督的名.
阿們.
\newpage



\section{羅馬書}
\label{sec:B5vFf1eCUBA}
\textbf{2025-11-14 周五晚查經班《羅馬書(27):分題:合乎中道的自我認識與恩賜的正確運用》 講員:曾裔貴牧師}
\newline
\newline
連結: \href{https://youtube.com/watch?v=B5vFf1eCUBA}{\texttt{https://youtube.com/watch?v=B5vFf1eCUBA}} ~~~~ 語音日期: 2025-11-15
\newline
\newline
\hyperref[sec:yNx82XHOkbk]{\small{< < < PREV SERMON < < <}}
~
\hyperref[sec:index_chronic]{\small{[返順時目]}}
~
\hyperref[sec:index_scriptual]{\small{[返順卷目]}}
~
\hyperref[sec:0TRElWqGJ4E]{\small{> > > NEXT SERMON > > >}}
\newline
\newline
$^{1}$我們一起同心祈禱.
天父我們再次感恩.
因為我們又可以在晚上.
又可以一起無論在線上,在現場.
我們一起去炒聖經.
也都當成為了錄音之後.
又可以在不同的地方.
一些姐妹都可以來聆聽.
我們都盼望這樣的解釋.
將保羅的原意我們盡量去還原的時候.
我們能夠掌握到保羅他對於福音的真言.
對於每一個信主的弟兄姊妹.
明白了救恩的弟兄姊妹.
怎樣能夠在地上神安排的教會.
即稱為基督的身體的地方.
在不同的地方的教會.
來各自去作一個有用,合用的肢體.
成為一個被神可以用的器皿.
因為這個是你所謂的奉獻.
因為我們的生命都是主用寶血拯救,買贖的.
求主幫助我們都能夠在第三到第八節.
當我們今天有機會去探討恩賜的時候.
我們都盡量能夠去吸收保羅所教會我們的.
我們仰望等候在你的面前.
我們這樣祈禱奉耶穌基督的名.
好 我們.
甚麼事?為甚麼不動?.
可以了.
我今天給的題目就是.
都是跟第三到第八節有相當高的重點的關聯.
就是怎樣能夠自我的認識.
這個認識特別是在一個場景.
保羅在說我們每一個都是一個活祭.
就是說我們每一個信了主耶穌.
被主的寶血拯救的神的兒女.
必須有一個這樣的對神的虧欠的心.
就是如果我們不奉獻我們自己.
我們就不能夠嚴格來說.
算得上是神期望我們的身份.
跟我們在地上信主之後的生活來得相似.

$^{41}$所以每一個基督徒絕對不能夠說.
我信了主 我洗了禮.
我仍然都是過我自己的生活.
是有一些牽強的.
甚至我自己都講了不止一次.
當然我近這十多年在感受的範圍.
真的洗完禮之後 簡直是不回教會.
見到你都不會覺得不好意思 沒有回教會.
當然每個人我們沒有辦法去深究.
他有什麼背後的原因 背後的困難.
但是如果那個意識是跟我們今天.
在《羅馬書》這個範圍所講的.
那就講不過去了.
可能他根本不知道原來神對我們有這麼高.
這麼榮耀的一個期盼.
就是我們很想我們每一個都是一個活祭.
這個是理所當然的 是應份的.
是我們虧欠神的一個這樣的生命.
我們應該是理應 我們盡我們的人生.
在我們有限的人生來奉獻自己.
這個就是合乎中道的自我認識.
在這個活祭的生命的大前提下.
我們就立即告訴我們知道.
這個活祭 這是理所當然的 是否敬拜.
原來發生在我們這個人.
怎樣將自己fit進去第四到第八節所講的.
保羅所給我們的一個很淑女的一個比喻.
就是一個身體有不同的器官.
這個是相當簡單.
但是更加深奧的就是第五節告訴我們.
基督都是這樣.
原來我們基督徒有一個淑女很深的奧秘.
就是第五節所講的.
我們許多的基督徒已經被主拯救了.
這裡告訴我們知道.
在基督裡都是成為了一個身體.
都是一個身體裡面.
然後是各自的作之體.
元宗的元朗作之體.
.

$^{81}$OK 謝謝 我們回去 怎樣縮小畫面?.
我們回去這個落下書的畫面.
又是這樣 有時就是這些彈出來.
怎樣可以這個畫面?.
.
是 是的 行 行 OK.
我們就回去這裡來看.
我先讀回剛才那一段.
第五節告訴我們知道.
我們許多人在基督裡成為一身.
互相聯絡作之體 也是如此.
就是說跟保羅在第三節所講的.
正如我們一個身體有不同的器官.
有不同的肢體.
我用了一個翻譯.
第五節我跟大家看看.
可能那個意思會更加再闊一點.
那個含義 大家可以在這裡告訴我們.
新的翻譯 一個原文的直譯.
照樣 照樣就是跟剛才上文所講的那個比喻.
我們在一個身體裡面.
是有不同的器官.
不同的器官各自的功用.
都是為了整個人每天的生活.
每天的思想 每天的情感.
每天所處的不同的抉擇.
不同的意志的決斷等等.
是令到整個人有那個力量去做任何的決定.
例如手,腳 腳要走路去做一件事.
手要做一件事.
完全是被這個身體的意志去操控的.
但是現在這個第五節告訴我們知道.
保羅就說教會就是這樣.
教會必須要這樣理解.
每一個弟兄姊妹要這樣理解.
自己在教會裡面的那個功用.
這裡告訴我們知道.
他說照樣我們這許多的人.
就是許多得救的基督徒.
被主拯救 在主的愛底下.

$^{121}$來得著救恩的人.
我們應該要獻上自己.
獻上自己之後 就要這樣理解.
我們在一個教會的整體裡面.
要有這個整體觀.
在基督裡是一個身體.
這裡比剛才第五節舊的和盤來得更加直接.
在基督裡面是一個的身體.
而且是各日 每一個基督徒.
互為彼此的肢體.
互為彼此就是說.
原來大家都是彼此有一個互相幫助,協調.
而且是一致和諧.
這樣來使基督的身體不斷地增長.
不斷地在生命裡面.
更加的無論是靈活.
更加的恰如其分地.
在不同的地方發揮那個功用等等.
這個就是保羅的那個行義.
所以他在第五節之後.
就開始列舉有不同的恩賜.
在第六節開始就說.
按我們所得的恩賜.
所以我們等一下就簡單地分析一下.
我們現在先回到這裡.
我們從第三節開始.
這個是漢語聖經的本.
我藉著神賜我的恩典.
對你們各人說.
不要對自己評價過高.
超過所應有的評價.
都要按著神量給各人的信心.
忠肯地評價自己.
評價自己在神所給我們的信心有多大,有多小.
忠肯地評價自己.
得到神給我們在我們的生命裡面.
有什麼的用處.
而這個用處呢.
那個領受的過程就是信心.
所以基督徒那個信心是要不斷地增長的.

$^{161}$信心就是說我們要跟神的關係.
跟神給我們在聖經的應許.
我們要透過信心來領受.
來明白.
先是從明白.
聽明白.
然後我們就知道.
原來神在這些事情有這些應許的.
我們就透過這些應許.
我們就去實踐.
我們就去知取.
這個信心就會不斷地增加的了.
信心是可以成長的.
信心是可以不斷地有增進的.
所以這個就是跟我們那個所謂的生命.
是有一個息息相關的.
所以我們要.
應該說.
用英文就是.
in terms of truth.
就是說我們要不斷在真理裡面.
來去建立自己.
信心才會增長的.
信心是不能夠.
你坐在那裡.
你不去認識神.
你不去學習.
去運用.
那麼信心是不可以增長的.
甚至乎.
如果我們不用.
那個信心還可能會慢慢地倒退.
那麼那個就出現問題了.
到時候我們就跟神.
是有一個很大的隔膜.
距離.
這是《記得後書》一章.
那裡告訴我們.
原來是要在我們的屬靈生命上.
來不斷地去學習.

$^{201}$去知取.
以及去行洗.
或者叫做運用.
那麼我們的信心就能夠增長.
所以這裡告訴我們.
原來是神給我們的.
這個糧原來可以增加的.
有的可以再加給他.
使得他再多一點.
所以這裡.
跟主耶穌說的比喻裡面的那些僕人.
其實沒有抵觸.
因為不是說一次過糧了.
就不再增加了.
也不會減少了.
也不會收回了.
就不是了.
是每個基督徒.
如果大家.
用我們大家明白的.
大家都是由初信班.
來自不同的背景.
甚至乎可能來自一樣的報道會.
假設有五個人信主舉手.
五個人都願意參加初信班.
大家都是聽.
當然那個部份.
都開始是一致的了.
但是慢慢的.
那個朋友.
那個弟兄姊妹.
繼續透過他在初信班裡面.
或者他自己開始去讀聖經.
去明白之後.
他自己繼續這樣去.
透過可能是讀金句.
可能是他很恨慕讀聖經.
不斷地去明白神的心意.
包括透過聽道.
不同的牧者的教導.

$^{241}$不斷地去增長的時候.
就會帶給他對神.
在他生命裡面的帶領.
那個給的恩典和恩賜.
就會越來越了解了.
當他開始.
他有時候是他自己.
已經開始有一些這樣的.
客觀的他自己的.
內心裡面的感動.
例如可能看到教會有某些崗位.
他開始頗有負擔.
頗有興趣去參加.
但是原來更多.
保羅告訴我們.
原來那個最大前提.
是神已經給我們每一個基督徒量了.
必須是基督徒.
所以真的某程度上信了主.
在教會裡面有些重要的崗位.
必須是洗禮.
他始終認識自己是神的兒女.
才能夠明白他所事奉的背後的目的.
以及事奉過程裡面有時候的委屈.
就能夠有神的恩典去補足.
當然你透過長輩,牧者.
弟兄姊妹的勸勉,安慰.
那些委屈就可以過渡了.
否則極容易會因為做完之後的負評.
吃力不討好.
就會令人很消沉,很沮喪.
因為完全是出於一份義務的心.
出於一份愛心.
但是換來諸多的批評.
是人都受不了的.
所以浸了禮的弟兄姊妹.
就知道自己是無恩的基督徒.
就知道是神的恩典.
是神安排給我們人生的崗位和恩賜.
所以不要小看這些價值觀.

$^{281}$這樣的價值觀是令到參與事奉的弟兄姊妹.
一方面就如第三節這裡所說.
原來在事奉之前.
我們真的要幫助弟兄姊妹.
對自己的評價,對自己的認識.
這裡說了 這裡已經好幾次提到評價.
對自己有一個不是平頭稟竹.
對自己的這裡所指的評價.
原來涉及到我們是怎樣的個性.
我們的過去的原生家庭.
到底塑造我們這個人的性格.
那個性向是怎樣的呢?.
是很衝動,很容易一句話就爆炸.
因為很柔弱的 什麼都不動的.
兩者都是需要調整的.
所以原來那個性向和那個崗位.
那個事奉某程度上都有掛勾的.
在正常情況下.
有些崗位真的不適合某些人做的.
一做就會帶來那個崗位的事奉都被他搞壞了.
可能不是不關他的事.
所以原來在這個恰到好處.
和合乎中道的事奉之前.
首先就是每個基督徒的弟兄姊妹.
一方面要認識自己是蒙恩的一個基督徒.
是恩典使我們能夠.
可以參與在基督的身體裡面的事奉.
這個就涉及到真的要將自己的生命調整.
這個不是關乎到你的才幹的問題.
是生命的問題.
所以如果調整得不好.
其實又是很多的衝突的.
因為我們一開始就是用我們自己的才幹.
用我們自己的才華.
用我們自己一腔的熱誠,熱心.
一衝過去可能那種被挫敗的銳氣.
會令到那個人有時回不了頭.
他不會再事奉 永遠不想再事奉.
原因就是可能在起步的階段.
我們沒有留意到保羅在第三節.

$^{321}$每個人都是蒙這麼大的恩典.
每個人都是在神的勸勉下.
要將自己奉獻給神作為活祭.
但是保羅告訴我們.
當他準備要講到教會.
要彼此來服侍去建立基督的身體之前.
第三節就是一個很重要的大前提.
保羅說他也是在神的裡面領受了這個恩典.
又或者不知道是不是他已經參與了一段時間的事奉.
他都看過有些人跌倒.
有些人在神的國度裡面.
在教會,在傳福音裡面去參與,去服侍.
不知道他內心還是不是.
因為現在第二次的旅行報導.
完了第二次 他現在在哥林多.
還是不是想起第一次的時候.
和巴拿巴一起去塞浦路斯傳道的時候.
遇見馬可 因為太辛苦.
就離開了他們在柏弗那個地方.
就自己走了 不再和他們兩個一起去傳道.
輾轉保羅和巴拿巴就回到唐非尼亞的別家.
然後在那裡坐船 就回到安提摩來述職.
這是第一次的報導.
所以他會不會在寫這篇經文的時候.
馬可在他的腦海裡面.
他是不是自視過高呢?.
還是他只是很想跟著保羅和巴拿巴.
導致他突然間不辭而別的一個舉措.
令到保羅第二次報導的時候.
都不要再帶著馬可去了.
和巴拿巴就有些分道揚鑣.
所以這裡原來告訴我們.
多清楚是要參與神國裡事奉的弟兄姊妹.
必須首先要做一個評估表.
所以教會是應該要幫弟兄姊妹做一個評估表.
當然在那個生態上.
真的每個崗位都沒有什麼人做.
所以有人肯做 就發光了那些牧者的眼睛.
你肯做啊 你肯做啊 真是開心啊.
但是先不要開心.

$^{361}$又可能害了他 因為他還沒有清楚自己.
還沒有做過一個自我評估.
還有牧者對他的評估.
如果我們沒有這樣認真去做.
而帶來的就是他未必是中肯的評價自己.
過高的評價自己.
這裡就已經肯定說超過所應有的評價.
對自己 一個初信主的弟兄姊妹.
無論她以往有多少的才華才幹.
在屬靈的國度裡面.
在一個新的屬靈生命裡面.
的而且確只是一個屬靈的初生輕孩.
這個我們真的要很小心 很留意.
我自己有時候不知不覺都會有一種錯覺.
或者真的需要去不斷地提醒.
有時候很容易會找一些名人.
他一初信 他本身是一個名人信主.
我們就覺得千載難逢.
簡直是萬年一遇這樣的奇人.
還不快點把他放在港台上講見證.
吸納起碼有200人來聽.
坐滿了整個紅恩堂.
或者甚至是錦繡堂這樣的效應.
總之不要理他.
但我們剛剛犯了第三節的.
這個保羅很嚴肅的一個勸勉.
未必是警告.
也不應該說絕對沒有例外.
保羅自己就是一個例外.
一信主立即被呼召.
不過又要留意一下.
如果我們稍微熟悉了這行傳.
保羅 我們先看一看這行傳的11章.
即是看回前一卷書.
應該是11章這裡.
對不起 應該是第9章.
9章的26節.
9章的16行傳 26節.
這裡說 素羅就是保羅.
他在大馬色信了主.

$^{401}$因為他在大馬色受到迫害.
不如多揭一點 就是23節.
過了一段日子.
猶太人商議要殺素羅.
其實他一直在安提奧.
一直都在大馬色.
商議要殺素羅.
但是他們的計謀被素羅知道了.
他們就在晚上在城門守候.
要殺他.
他的門徒就在夜間.
用一些竹籮 用一些糠紙.
把他從城牆上墜下去.
在哥羅德後書的11章.
保羅記回這件事.
說這是一個極大的恥辱.
就是說他本來打算去大馬色.
拿那些文書去找那些祭司長.
誰知道現在就反過來被猶太人追殺.
甚至包括大馬色的軍兵.
都要守住城門.
不讓他出入.
就是說你沒有通行證.
就不知道了.
就看是誰保羅.
當時不知道有沒有證件.
應該就沒有了.
就要捉拿他.
門徒就用竹籮把他放進城牆外面.
他就走了.
然後26節這樣說.
素羅到了耶路撒冷.
想與門徒結交.
他們卻都怕他.
不信他是門徒.
唯有巴拿巴接待他.
第一次認識巴拿巴.
巴拿巴不是使徒.
但是巴拿巴在那個時候.
已經在耶路撒冷生活.

$^{441}$與使徒們出入,來往,相交和事奉.
就帶他去見使徒.
把他在路上怎麼向見主.
主怎麼向他說話.
就是說保羅就說了他的得救見證給巴拿巴知道.
巴拿巴就成為了一個橋樑.
因為使徒敢見他.
怕他.
因為使徒們是耶路撒冷人.
耶路撒冷人就很清楚保羅是誰.
所以知道是保羅.
但是不敢相信他是信主.
信得這麼好.
信得這麼戲劇性.
所以現在就要由巴拿巴做一個人界.
去見使徒.
把他怎麼在路上.
就是大馬式的路上.
有大光光照他.
他就瞎了.
接著有人帶他去大馬士革的直街的一個地方.
在那裡居住.
三天三夜都禁食.
不吃東西.
很多的思前想後.
因為你瞎了什麼都做不了.
又上不了網.
又什麼都做不了.
又去不了街.
又看不到電影.
什麼都做不了.
瞎了.
不用瞎了.
我們如果整個感受停電.
我們什麼都做不了.
暴動一次.
上不了網.
沒有冷氣我們都受不了.
但是保羅就三天三夜.
其實他是思前想後的懊悔.

$^{481}$因為沒辦法想象.
自己本來是一個殺人兇手.
準備去抓基督徒.
被這位叫做耶穌的.
在那個異象裡面的大光.
跟他對話.
用希伯來話跟他說.
我就是你所逼迫的耶穌.
用希伯來話跟他稱呼.
不是用希文.
用希伯來話跟他說.
接著主怎樣將他拯救.
來改變他.
所以保羅現在跟巴拿巴說這個得救見證.
巴拿巴就深信不疑的接納這個弟兄.
接著就說這是一個很奇特的見證.
得救的見證.
所以一定要去找使徒.
直接找使徒.
所以使徒們就知道之後.
他在大馬式怎樣奉主耶穌的名.
放膽傳道.
都述說出來.
就是說保羅已經不管其他人怎麼看他.
他就很大的內心的性命感動.
他就自己在大馬式格.
到處去跟別人說.
當然可能就是跟猶太人說.
證明耶穌是基督.
接著就是28節.
於是訴羅在耶路撒冷.
和門徒出入來往.
奉主的名放膽傳道.
並與說希利尼語的猶太人.
講論辯駁.
他們卻想發旨要殺他.
這裡告訴我們.
27和28節之間.
使徒們聽了弟兄的得救見證.
這個弟兄就是以往的拉斐.

$^{521}$一個很有社會地位的猶太人的議事.
代議事.
上議會的一個議事.
再加上他現在有這麼大的戲劇性的生命改變.
所以他們就接納.
也准許他在耶路撒冷.
跟門徒出入來往.
就是說保羅開始天天可以抹頂.
可以跟門徒一起查經.
來分享他在神裡面的形式.
因為他是一個舊約的權威的拉斐.
但是這樣就沒有空間.
他去沉澱整理自己往後更深.
更大的信心的事復.
所以開始有些逼迫.
就是這個原因.
29節他奉主的名放膽去傳道.
而且就是說.
他開始跟那些不信的猶太人.
不只是跟猶太人的信徒來往.
不只是跟使徒們來往.
給自己分享查經.
然後他開始知道自己內心的感動和照明.
就是要向那些未信的猶太人.
去傳講耶穌彌賽亞.
不僅如此.
還要跟那些 這裡告訴我們.
還要跟希臘尼旺的猶太人講論辯駁.
講希臘尼旺就是從外地回來的.
就是說不是在耶路撒冷的猶太人.
就是說有其他的從外地上來的猶太人.
他們就不是住在耶路撒冷的.
保羅就是認識希臘尼旺.
就用希臘尼旺跟他們講道.
跟他們分享 跟他們去查經.
跟他們去討論關於耶穌的事.
當然了 這個講論辯駁.
帶來的就是這些從外地回來的.
就是說我們上來就是想守節.
我們就是想守猶太教.

$^{561}$你居然阻礙我們去祭拜我們的神.
說耶穌是神.
所以這個就是很大的對他們的得罪.
所以他們就想發旨要殺他.
如果保羅是唯有他是特殊的例外.
那就沒有辦法了.
但是更多的在現代的21世紀.
我們似乎好像從來都不敢面對這一次聖經.
我們講聖經 講給別人知道.
我們很怕別人怎麼看我們.
所以我們永遠都不會去到保羅這裡.
他們要想發旨殺保羅.
保羅其實是不是橫斑棒呢?.
不講道理呢?.
你們每個人都要聽我的 不管那麼多.
用權威 用他的職位.
用他過去的官階.
相信不是的.
相信保羅是用很謙卑的心來用舊約.
因為當時沒有新約的 對不起.
沒有新約的.
當時是用舊約來跟他們講論辯駁.
保羅是引舊約聖經的.
所以真的駁不了.
保羅一通了舊約的寓言.
就是來塞亞耶穌.
所以在這樣的辯駁裡面.
他們又好像不單只是普通的禮貌地婉拒.
就算 似乎不是禮貌地婉拒.
禮貌地婉拒 不聽那也沒辦法.
但是似乎到了一個階段.
他們是要表態.
你們一定要相信 如果不相信是有後果的.
你們等候的彌賽亞.
將來會審判你們的.
這樣的說話措辭.
才會勾起這幫人.
這幫講希臘話的人.
對保羅的敵視仇恨.
才想要殺他.

$^{601}$如果是很輕輕的 不相信了.
不要緊 不要緊 大家討論而已.
就不會有這樣的情況.
所以我想我們今天欠缺的.
不是我們沒有福音單張.
不是我們沒有報道會.
而是我們很少有弟兄姊妹.
願意講那個聖經真理.
用我們自己的生命代價.
去為這個真理去作見證.
我們很少真的有一個心.
就是我們對那個人的信主是很迫切.
你不相信 你將來落火湖.
我是有責任的.
這樣的心情.
就會好像保羅第28章所講的.
他真的要面對殺生之禍.
講這個福音.
就是說講得很恆切.
講得很杜獄.
那些人是真的感受到威脅.
那個威脅是生命生死.
永生永死的威脅.
所以才對保羅不客氣.
就不是說他很沉.
說他很煩等等.
所以這個也是保羅神給他的一個.
你可以說是他事奉的必須要這樣去做.
但是苦心自問我們信了主這麼多年了.
我們一提到要跟人講福音.
是不是我們就立刻避開.
我們就盡量低調.
我們盡量不知道.
這個也是一個我們大家要小心.
我們這個時代欠缺的就是.
原意為主去付代價.
我不敢說要付死亡的代價.
至少我們付代價.
被人嘲笑.
被人侮辱.

$^{641}$我們願不願意呢?.
沒有什麼殺生之禍.
但是我們願不願意這樣.
真的為了福音的緣故.
來受一些這樣的.
被人這樣針對呢?.
保羅就願意.
不過當然了.
這裡給我們看到《聖行傳》.
勞駕醫生記載.
都挺持平的.
剛剛才在大馬士革.
剛剛才上耶路撒冷.
剛剛才開始出來傳道保羅.
就是說其實這些的迫害.
這些的障礙.
這些的挫折.
其實是令到他更深的得著.
就是第三節所講的.
那個信心可以更加的進深的.
一些的原動力.
一些的必需的元素.
所以千萬不要以為我們面對一些的逆境.
一些的扶貧.
一些的人對我們的誤解.
我們就馬上退.
我們所謂的潔身自愛.
我是一個義工.
誰都不能說我.
一說我我就馬上走了.
我馬上退的.
其實這樣就很可惜了.
因為原來這裡告訴我們.
原來我們可以不斷的增長我們自己的信心.
不斷的在無論別人讚賞.
別人對我們的誤解.
都可以使得我們有成長的.
所以這裡告訴我們.
保羅暫時就不在耶路撒冷了.
這裡告訴我們很特別的.

$^{681}$是送走了保羅.
這個可以說是一個很大的智慧.
就是說保羅現在暫時還不可以在耶路撒冷.
受到萬人的尊敬.
和使徒們.
其實他真的是耶穌特別挑選的使徒.
應該是第十三個使徒.
其實他的身份是相當平起平坐的.
和十二個使徒.
其實是十一個.
因為猶大死了.
和十一個使徒.
真的平起平坐.
出入耶路撒冷是沒有問題的.
不過留意一下.
三十次.
這些經文我們平時比較少留意.
弟兄們知道了.
這個弟兄們到底是不是使徒呢?.
因為是耶路撒冷的門徒.
相信一定是使徒們.
就送他下到該撒尼那.
然後再打發他回大掃.
大掃就是他出生的地方.
回到大掃.
似乎是回到他的家鄉.
在那裡做什麼呢?沒有說過.
接著到第十一章.
我們就知道了.
我們站在第十一章.
所以我們先看這裡.
再看第十一章就會更加準確些.
這個要看後面那裡.
第十九節.
十一章的十九節.
剛剛是彼得他帶了哥尼留一家信主.
他就由該撒尼那回到耶路撒冷.
一回去就遇到一個很大的風波.
就是說使徒彼得.
都一樣會被一些不解真理的人.

$^{721}$他所做過的事情.
而來到這個落差.
那時候不能拍片.
沒有YouTube.
不能夠有圖有真相.
他就要去解釋.
為什麼你向外邦人傳福音.
以及你是猶太人.
為什麼你要進入外邦人的家.
跟他們同食同住.
你可以想像原來初期教會.
那種觀念還是這麼大的局限.
那種狹隘的民族思想.
受著他們的宗教律法的規範.
原來框得這麼緊.
就算是彼得自己見到異象.
神親自跟他說話.
他自己都半信半疑.
後來帶了全家的哥利留的所有信主.
又見到聖靈的降臨.
但是耶路撒冷完全不知道這些事情.
知道的已經傳了回去耶路撒冷.
彼得居然夠膽向外邦人傳道.
當然傳道是可以.
但是進入外邦人的家.
就像我們今天看的匪夷所思.
沒辦法想像.
不知道多開心.
有外國人來探望我們.
有猶太人來住在我們家.
我們都不知道多開心.
如果有阿拉伯人來.
更加開心.
如果大家都是信主的.
但是原來初期教會.
強如彼得這麼崇高.
他自己本人都要經過一番掙扎.
所以他也不能怪耶路撒冷.
可能連使徒們都對他有這樣的咎責.
那個風波最終都平息了.

$^{761}$一方面彼得也很有智慧.
原來他不再衝動.
他帶了六個弟兄一起去哥利留家.
不是他自己一個去.
原來有見證人也很重要.
有時候做一些事情.
我們真的需要有一些謹慎.
千萬不要過急.
原來莫彼得也稍微有些印象.
我自己都受不了.
耶路撒冷的弟兄會不會誤解了我.
來跟外邦人同住.
玷污自己呢.
就帶了六個弟兄一起去.
所以有六個見證人幫彼得解話.
神的國度.
耶穌完成了救恩升天之後.
再沒有分任何的民族.
沒有任何的民族是特別聖潔.
有些民族是特別污穢.
已經沒有了.
主耶穌已經將這一切完全歸於一.
這些在當時.
保羅還沒有寫這些經文.
還沒有這些經文之前.
就有這些處境.
原來他們要克服.
猶太人自己信了主.
都一樣要克服這些宗教帶來的負面.
這些錯誤的觀念.
原來這些觀念要慢慢地去洗除.
所以原來神的救恩.
讓人有自由意志.
都可以讓這些在舊約背景的猶太人.
又見過耶穌的.
又信了主的.
又經歷了聖靈很大的恩典的.
又見過耶穌復活.
又和耶穌一起生活過的.
一樣原來舊約的觀念都還沒有一下子清除.

$^{801}$所以我們小時候一直在香港長大的.
很多的價值觀.
很多的觀念.
或者迷信的觀念.
我們都要學習.
在基督裡面來學習.
我都講過.
有時候晚上YouTube會看到臥底旅行團.
挺好笑的.
但是我又發現了一件事.
阿正和阿….
忘記了什麼名字.
每次進酒店一定敲門.
這些觀念肯定是信的人.
他們知道他們目前香港人的心態.
進去一個陌生的地方.
不好意思我現在進來住了.
謝謝.
不好意思.
我們基督徒就一定不要做了.
除非你說有人中午來收拾房間.
他進來敲門是沒有問題的.
因為人和人.
但是現在很明顯.
你是進去你給了錢的房間.
你用門卡.
之後應該拆下去.
然後開燈.
燈自己開.
那房間就是你的了.
因為給了錢的.
為什麼我要敲門.
不好意思我現在來住這間房.
這個是我們自己迷信.
我們自己還有這些殘存觀念.
所以彼得就要面對猶太人.
居然說你不行.
你去哥利留家裡.
那裡是外邦人.
他們是白夫長.

$^{841}$其實人家都很敬虔.
人家經歷了聖靈的降臨.
他們是完全懼怕.
等待聽彼得說.
怎樣可以得救.
所以你想想原來.
在耶路撒冷的猶太人.
還以為自己在做一些很偉大的事.
其實原來那個心胸.
還是沒有經歷聖靈工作的修剪.
還是停留在耶穌還沒有降生之前.
舊約的背景.
忘記了耶穌帶來一切都更新.
耶穌帶來的一切都更新.
但是人還跟得上.
所以聖經就記載這些事件.
告訴我們.
給我們自己的後世基督徒.
你是香港人.
你就要有香港人那些錯誤的觀念.
那就要更新了.
如果你是一個韓國人.
或者你是日本人.
你有你自己的傳統觀念.
你都要去更新.
所以這就告訴我們.
彼得用了很多功夫才勉強.
第十八節.
我還沒讀下面之前.
讀這裡.
十八節.
眾人聽見這話.
就不言與子歸榮耀語神說.
這讓看來神也賜恩給愛邦人.
叫他們悔改得生命了.
你可以想像到.
彼得用了很多的唇舌.
解釋整個事件.
他們才真的肯接受.
好像他們比神還大.

$^{881}$接待愛邦人本來是不應該的.
但是現在神接待我們沒辦法.
類似好像這樣.
其實這些不知不覺.
在我們大家的生活裡面.
總有一些我們要學習的.
我們要真是.
我們自己的生命應該要.
來遠來遠有一種救恩帶給我們的全面.
救恩帶給我們的生命的成長.
這些跟信心有直接關係.
接著就安提摩教會.
我們留意一下.
因為主要我想講保羅.
不是講彼得.
第十九節.
那些因斯提反的士.
遭患難四散的門徒.
直走到菲利基和居比路.
居比路就是第一次去旅行.
去宣教報道的地方.
就是塞浦路斯那個島.
地中海那個島.
就是巴林巴的家鄉.
巴林巴就在那裡出生.
被安提摩.
接著他們不向別人講道.
只向猶太人講.
但是內中有居比路和古里內人.
古里內人就是北非那地方.
卡門到了安提摩之後.
也向希利尼人傳講主耶穌.
所以路加醫生就告訴我們.
因為斯提反第七章尾.
他就殉道.
斯提反就是耶路撒冷的執事.
在那裡來到殉道之後.
那些門徒因為大兵不.
那些門徒到處走.
走到不同地方.

$^{921}$不單止是避難.
的確是避難的.
不過避難之後.
他們有更大的內心.
就是去到不同的地方.
向當地的人去傳福音.
所以就說了.
一方面可能有些人不認識希利文.
就只向講白萊話的猶太人傳道.
但是原來有些.
這裡說第20節.
內中有居比路和古里內人.
他們到了安提摩之後.
都向在安提摩的.
講希利尼話的那個世字古卷.
就是說向說希利尼話的猶太人傳講主耶穌.
可能兩者都可以.
就是向希利尼人傳道.
也都是向講希利話的猶太人傳道.
可能他們走到那裡之後.
也有些人是能夠.
可以用希利文溝通.
所以也很合理很正常.
所以就是說那個涵蓋面.
不單止只是耶路撒冷的猶太人.
或者旅居在安提摩.
就是小拉西亞地方的土耳其中部的地方的人.
而是一直發散出去.
就有現在講希利話的人.
他們都是福音的對象.
那時候聖經告訴我們.
在第21節主與他們同在.
信而歸主的人就很多了.
本來這個是大兵八走難的.
但是走難帶著的.
不要只是一種惶恐.
是的要去很多部署策劃.
也對的但是不要失去了.
那個福音的照明.
不能夠在耶路撒冷暫時去傳福音.

$^{961}$就不代表不去其他地方傳福音.
所以就是了.
移民去英國去加拿大的.
無論是多倫多,溫哥華.
或者英國的去倫敦,曼徹斯特等等.
或者其他再遠一點的中部等等.
有各樣的準備.
但是記得福音的照明一定要帶著去.
這樣就能夠可以再經歷.
在不同地方的特別的恩典.
而且是特別豐厚的恩典.
這裡告訴我們.
他們都面對一個很大的逼迫.
必須要走.
走是忠誠的.
但是去到之後.
當地他們就要馬上安頓好之後.
第一時間他們要集合去傳福音.
這裡告訴我們.
主與他們同在.
當然不止同在.
就是同工.
所以信而歸主的人就很多.
這裡是安提摩.
在土耳其.
所以當地的土耳其.
只有1\%都沒有的基督徒.
其實是相當少的.
當時的土耳其是很多人信主的.
然後這個風聲傳到耶路撒冷教會的人中.
又說耶路撒冷有大兵包.
因為那些人都走了.
可能受牽連的人都走了.
可能他們要抓的是那些執事.
或者是和斯大林有比較密切的.
因為斯大林說的道是激怒了猶太人.
所以那些人都走了之後.
耶路撒冷的使徒就好像沒有受到牽連.
也都相對慢慢平靜了.
所以在耶路撒冷又繼續可以有正常的教會運作.

$^{1001}$但是這裡告訴我們.
這個風聲就傳回到耶路撒冷教會的人的耳中.
他們就打發巴拿巴出去.
為什麼不打發保羅呢?.
因為在這個時候.
保羅還是在大掃這個地方.
這裡是過了一段日子.
接著巴拿巴就一直去到安提安.
23次他到了安提安之後.
看見神所賜的恩.
就歡喜.
就勸勉眾人.
就是那些新信主那些.
立定心志.
行久靠主.
巴拿巴就由耶路撒冷所猜探的特使.
就是要看一下.
現在安提安為什麼這個風聲說他們有大復興.
有很多人信主.
很多人歸復主.
所以他們就要去看.
就看到那裡很大的神的恩典.
巴拿巴看到的就是一班新信主的弟兄姊妹.
是很充滿了神的救恩帶給他們的喜樂.
平安.
還有追求的心.
還有大家彼此的連結.
相交.
擺餅.
聚會等等.
所以巴拿巴就勸勉他們.
就叫他們立定心志.
要行久的.
不要一時三刻.
不要現在因為沒有冰巴.
沒有困難.
你們就很忠心地跟隨著.
但是當有冰巴來的時候.
你就沒有了行久的心去靠主.
所以原來是需要這個信心.

$^{1041}$這個信仰.
是一定要堅持到底才夠的.
接著24節就是這裡.
就是保羅的出現.
24節.
這巴拿巴原是個好人.
被聲明充滿路加醫生的評價.
帶有信心.
於是有許多人歸服了主.
又是一個牧者的出現.
就帶來教會的復興.
更加深層的復興.
因為他的感染力.
巴拿巴的感染力.
他是一個好人.
被聲明充滿.
好人必然的.
為什麼要這樣描述呢.
原是個好人.
但是這裡提到的就是.
他不是一個很好的先生.
一個君子這麼簡單.
巴拿巴的強項.
就是一個叫做勸慰者.
《小塔堅實四章》那裡這樣說的.
他是對人有很大的安慰.
勸練.
所以他是Man of Comforter巴拿巴.
巴拿巴本來就是花名.
其實他的真名叫做約瑟.
是在塞浦路斯的一個藥室.
所以他不是叫做巴拿巴.
巴拿巴是一個花名.
巴拿巴.
巴就是BAR.
是猶太的.
希伯來話的.
誰的爸爸.
這樣的意思.
對不起.

$^{1081}$誰的兒子.
巴拿巴.
他是拿巴的兒子.
所以叫做巴拿巴.
福音書有時候叫做巴土賣.
就是土賣的.
巴底賣.
那個黑蓋.
盲眼那個.
黑蓋巴底賣.
其實就是底賣的兒子的意思.
所以巴拿巴就是拿巴的兒子.
拿巴其實就是原文.
舊約原文.
那個先知的那個子君的那個譯音.
所以巴拿巴都是一個在神的角度裡面.
他是一個很能夠血滅人的人.
所以第一個是他接納保羅的.
因為保羅那些人都怕他是一個臥底.
厭仔.
想回去顛覆耶路撒冷的教會.
所以這裡告訴我們.
他被聖人充滿大有信心.
所以教會因為他的信心.
因為他的見識.
因為他的聖人的充滿.
所以令到更加多人.
不只是信主.
是歸服主.
歸服主就是說.
願意將生命都來交給神.
好像巴拿巴那樣.
來到他一生的.
來到他讓神帶領.
這樣的歸服主.
是一個很重要的根基.
很重要的基礎.
教會需要一些很重要的基礎的人.
他們已經將生命奉獻給神了.
很願意無論什麼環境.

$^{1121}$我和教會一起面對的.
這樣的中流砥柱.
就是要歸服主.
所以今天教會的問題.
有時候不是不夠人.
不是不多人.
而是不是太多人歸服主.
不是太多人願意將自己的生命交託給神.
在神的角度裡面.
很希望有份.
很希望參與.
無論有富貴.
無論有患難.
都一同面對.
這樣的基督徒.
就算是稀罕了.
所以有時候用一些苦難.
就可以練成基督徒的生命.
有時候神容許一些困難.
是會使我們更加練成我們的生命.
就是說如果你有面對逼迫.
你都真的願意為主受苦.
就真的證明你是歸服主了.
所以這裡說歸服主.
是說他們不怕付代價.
不怕別人對他們嘲笑.
如果假如耶路撒冷再有人逼迫我們.
他們都預了這個代價.
這就叫歸服主.
所以盼望我們真的在人生裡面.
學習有這樣的追求.
當然人多不是我們不喜歡.
而是量和質一定要有一個平衡.
如果沒有.
就變成一個很大的隱憂.
接著25節巴拿巴馬上想起一個人.
因為看到這麼多人歸服主.
他覺得一個人都不夠.
其實這個時候.
他應該馬上打報告.

$^{1161}$當時不知道是不是非格傳書.
如果有WhatsApp就好了.
馬上就收到.
馬上就要通知耶路撒冷.
因為他是一個特派的特使.
他應該馬上告訴耶路撒冷.
但是因為安提阿離開大塑很近.
在土耳其.
他不回到耶路撒冷.
他突然間想起一個人.
就是第25節.
他又往大塑去找塑羅.
找雀鳥.
就帶他到安提阿去.
他們足有一年的功夫.
和教會一同聚集.
教訓了許多人.
門徒稱為基督徒.
是從安提阿起手.
你可以想像到.
到底保羅他沒有遇到這個.
其實是由他開始的.
他看到斯大林被殺.
被人用石頭扔死.
他是喜歡這件事的.
他是看守著那些扔石頭的人.
那些劊子手那些衣服的.
接著他就輾轉.
他就去大馬士革去抓基督徒.
就是說這段時間.
這一班走下來.
安提坊的人.
其實就是從耶路撒冷那邊走下來.
其實是和保羅有關係.
保羅當年也逼迫他們.
但是你可以想像到.
巴拿巴他沒有回耶路撒冷報告.
沒有回去找救兵.
找一些厲害的使徒來.
一起去做這個工作.

$^{1201}$然後他想起在土耳其.
有一個他親愛的弟兄.
到底在這段期間.
保羅在耶路撒冷被人.
那個人其中可能有巴拿巴.
就將他帶一些蜆送回大宋.
在那裡安靜等候.
這個就是保羅的那段真空期.
接著就發生很多事了.
彼得就下了去哥薩尼亞.
接著就有這個腓力.
還有那個大壁巴等等.
接著就有撒馬尼亞的事件等等.
接著作者就再巧妙地.
透過安提拉的大復興.
記載了巴拿巴.
這次是第二次引薦保羅.
參與那個侍服.
所以都挺特別.
沒有回耶路撒冷找救兵.
反而自己靈機一蹤.
想起自己親愛的弟兄.
曾經被這樣拯救的弟兄.
現在怎麼樣呢?.
現在他的生命是否具備.
那個大大復善的時間呢?.
但是不管了.
先去找這個保羅.
就去大宋.
意味著巴拿巴都有機會.
曾經都去探望過保羅.
和他都有一些機會相交.
不過似乎神的時間還沒到.
保羅仍然在大宋等候.
當然那個等候意味著.
一樣會傳道.
一樣在教會裡面聚會.
一樣他都會按他的內心的照明.
他會向外邦人傳道.
不過只是留在大宋.

$^{1241}$盧家醫生沒有記載這一筆.
但是無論如何.
巴拿巴很快就去大宋.
一想起就去大宋.
去找保羅.
找到了.
他們兩個人.
當然找到之後就帶他去安提拉.
讓他看到這些過去.
猶太人受理.
和那些猶太的忠奸分子.
猶太教的那些人.
逼迫的基督徒.
現在走到安提拉.
聚集在安提拉.
所以本來是逼迫人的保羅.
現在在安提拉成為這些被迫害人的牧者.
又是一個很戲劇性的神的工作.
神帶給當時安提拉教會.
一個這樣的美麗.
兩個傳道者在安提拉的教會.
這些告訴我們.
用了一年的時間.
在教會和他們一起聚集.
當然有查經.
有講道.
查的是舊約聖經.
這是保羅的強項.
然後就教訓了很多人.
然後告訴我們門徒稱為基督徒.
是從這個地方開始.
在這個地方開始.
就是說這個地方去到一個地步.
大家非常之感覺榮幸.
我是基督耶穌的門徒.
就是基督徒.
我是跟隨耶穌的一個人.
在街上被人看到.
原來你是基督徒.
我很感到榮幸.

$^{1281}$都很想跟他說信耶穌.
這是一個榮幸.
昨晚我們參加的課程.
大角咀現在越來越基督化.
因為陳淑娟牧師帶來的街的店鋪等等.
都感受到這個教會的教友.
是真的有一份活力.
一份光彩.
一份服侍人的心.
大概是這個意思.
安提摩見到有些人越來越有活力.
越來越有內心的喜悅.
平安.
就叫這班人跟隨耶穌的門徒.
就是基督徒.
這個意思就是說.
不介意跟耶穌的過去拉上關係.
因為耶穌的過去就是一個死囚.
說自己是神的一個人.
但是用基督徒的眼光.
就是屬靈的眼光.
就是他真的從天上來這個世界.
拯救我們.
代替我們的罪的救主.
就跟耶穌拉上這樣的關係.
人家看你跟他有關係.
他是死囚.
是的 你們誤會了.
他是死囚.
但是他是從上天來替我們罪人付代價的死囚.
所以以自己跟耶穌有份有關為榮的.
這樣才叫基督徒.
所以今天的香港.
不要說這麼遠.
說回自己的香港.
我們有多少百分比.
都有差不多.
其實好像是4-5\%.
全香港的基督徒.
已經算上天主教等等.

$^{1321}$不同的宗派.
其實很少.
但是我們能夠每一個都很大的榮譽感.
我是跟耶穌這個過去.
完全拉上關係.
我不會怕醜.
我還是很榮耀.
他真的為我的罪死的.
他真的是神的兒子.
所以如果他是普通的死囚.
當然我就會感覺羞恥.
你跟葉繼歡是有關係的.
你也可能會羞恥.
好像歐陽炳強.
他家人太太跟他離婚.
因為他是第一宗紙盒藏屍案.
你想想.
如果你跟一個死囚是有關連的.
你真的是他太太.
或者你是他子女.
你都想跟他割席.
他做出這樣的喪心病狂的事.
假如真的是他做的.
你可以想像到.
現在我們就應該這樣理解.
這個人被認為是死囚的耶穌.
我們基督徒就知道.
他是由死復活.
以大能顯明他是神的兒子.
《羅馬書一章》四節的這位救主.
所以如果我們稱為基督徒.
就是說意味著.
我們要知道他是死而復活的救主.
而且是永活將來要再回來審判的救贖主.
所以現在他還沒回來之前.
世人看的就是他的死是羞恥.
我們基督徒看的.
他的死就是一個榮耀.
所以十字架是一個羞恥的記號.
但也是一個榮耀的記號.

$^{1361}$就是這個意思.
所以就是這樣.
將他們的生命更加的堅固.
每個人在教會裡面都看自己成為神兒女的一份榮耀.
所以這份榮耀感.
如果我們能夠散發在整個教會.
大家都感覺到回來教會是一份很大的福氣.
能夠事奉神是一個很大的恩典.
這樣就全民皆兵了.
這樣就很不同了.
很多時候就是我們還帶著我們自己.
可能很優越的身份回來教會.
那心態就很錯了.
因為回到教會就是.
我是英佛坐在這裡.
你講道是英佛講得好聽,好笑.
我是要很滿足我的頭腦,我的思想.
稍為悶一點,稍為沒有勵志.
我就做什麼呢?.
你們唱的詩歌一定要很完美.
當然準確.
當然我們要有好的準備是應該的.
因為是一個聚集來的.
試過的東西一定要有準備.
不應該隨便.
但是我講的是每種個別的心態.
我們回來是敬拜朝見神.
頂尖師傅已經盡了他的力事奉.
所以如果我們一用自己的心態.
沒有用基督徒的心態去看.
我們就會有時在這個落差裡面.
我們就會有一個真的來源不切的心情.
那個這樣的錯誤扭曲的心態.
如果一直這樣坐大就更嚴重了.
如果這個心態都沒有根本來到他自服了.
根本沒有以自己是基督徒為榮耀.
你見他是一個初信者.
還不斷加一些事務給他.
越加就越高,越加就越重.
加到他成為堂董,成為主席,成為什麼什麼.

$^{1401}$我沒有任何的貶義.
但是如果他內心根本沒有真正.
在自己是神兒女的一個尊貴身份.
你越加給他就越是一個對他的傷害.
對教會的傷害.
因為他慢慢就繼續用他自己的價值觀.
來治理教會.
這是一個很大的問題.
說了這麼長其實就是第三節.
黑如保羅所告訴我們.
千萬不要衝之過急.
帶領的牧者要去評價能夠事務的人.
首先千萬不要歧視.
每個人都是神良信心給他.
一定有他的價值,有他的事負可以參與.
也可以這麼說.
如果我們根本沒有這麼多事負給人.
他都發揮不了.
這是一個問題的錯配.
這是我們的責任.
但是更大的就是.
每一個在主救恩裡面被拯救的新生命的活祭.
必須首先要過這一關.
就是對自己有一個評價.
對自己參與事負的一個評價.
就是不要看自己過高.
超過了所應該有的評價.
就是說他一定要自己的價值觀.
信耶穌本身就好像一隻駱駝穿過真蘭.
就是說你真的要放下你自己過去的價值觀.
你要接受一個死囚.
他是永恆的救主.
不是說你突然之間要穿破衣服回來教會.
不是這件事 不是那個外表.
是你的生命有一個洗滌.
就是神.
而我這個罪人他都肯拯救.
就是說一切都是恩典.
那個時候我們就應該在恩典裡面去評價自己.
所以不是說你要將你所有的才幹全部扔掉.

$^{1441}$不是這個意思.
神可以再用你的才幹.
但是那個用的過程.
是要被神的愛和恩典去改變.
以致你再用你的才幹的時候.
已經不是發揮你自己的才華.
而是要貢獻給教會的整個群體.
去成為祝福.
所以兩者可能只差一線.
那個線就是你怎麼看自己.
如果看自己是英勳這樣做.
那就會帶來只會越來越有一種互相的評價.
互相的一個可能是明顯的.
來到有攻擊.
或者暗中的就會有批評.
就會帶來教會的一個分化.
所以如果我們這裡做得不好.
保羅的教導是這麼奇妙的.
每一個參與事奉的人.
一定要認定自己是活祭.
然後才要做一個自我評估.
所以那個評估書.
不只是說你有什麼恩賜這麼簡單.
而是要對我們自己的生命.
是要有一個重新的洗禮.
就是我們成為神無條件拯救的罪人.
現在成為了他的兒女.
這個其實要更多的層面.
所以保羅現在在大肅出來.
那裡有多長時間呢?.
我們就不知道了.
在大肅一個這麼小的地方.
他可以做什麼呢?.
但有更多的空間就是.
再沒有在耶路撒冷的光芒.
那些的射燈射著他.
再沒有一個剛剛初上的保羅.
大肆的宣傳周圍去講見證.
沒有了 因為他先不用去做.
他就要回到自己的大肅城.

$^{1481}$一個很小的土耳其中部的地方.
在那裡安靜.
原來這一種沉澱.
這一種對自己的追求.
重新要整理舊約.
是不是真的這麼說呢?.
為什麼我以前從來都不知道呢?.
為什麼我自問這麼熟悉舊約的一個拉比.
為什麼原來舊約講的東西是離錯的呢?.
我完全是一知半解的呢.
因為他沒有聖靈.
因為他根本沒有新的聖靈.
根本他都沒有經過主耶穌的救恩.
所以看的舊約.
就只能是字面的舊約.
再加上猶太人那種驕傲.
那種錯誤的理解.
以致他以為神都要求他們.
永遠要跟外邦人分化.
所以你想一下.
原來這個沉澱的保羅.
在大數的評價自己.
是相當相當的重要.
我們有多少人初信.
就有自己的大數成呢?.
有自己在初信班.
我什麼都不懂的.
我很想先聽多一點.
我上主耶學.
我完了朕禮班.
更多人如果朕禮班以為會飛了.
以為什麼都不用再聽了.
以為什麼都不用再學了.
查經不用了.
我星期天回來崇拜就行了.
崇拜是那半個小時講一段經文.
整體的聖經根本就是覆無倫國.
你怎麼可能靠每個星期.
你自己不去讀經.
你靠那半個小時.

$^{1521}$你就說我熟讀了聖經.
然後保羅這麼熟悉舊約.
他都要在大數.
神將他放在大數.
其實是人將他放在大數.
在那裡沉潛.
在那裡沉澱自己的生命.
再一次思想.
為什麼神要這樣拯救我.
現在還要差遣我.
要我向外邦的外邦人傳道.
還要去到很遠他都要去.
所以這個沉澱.
是給他將來更大的和祥的工作.
練就他的信心.
所以我經常都講.
初入教的不可作監督.
是保羅自己講的.
在晚年的時候.
跟提安泰講.
提安泰因為當時在爾忽所教會.
他要治理一個很大的城的教會.
因為保羅留下他在爾忽所.
他要去馬其頓.
保羅第一次放監的時候.
所以保羅跟提安泰講.
你不要這麼快.
見到一個人稍為熱身.
就好像發了光一樣.
什麼都給他做.
馬上找他做監督千萬不要.
初入教的千萬不要做監督.
很簡單的.
根本下面就已經講了.
恐怕自高自大.
就陷入魔鬼的網絡.
原來越劈重的工作.
大力的工作.
如果沒有經過生命的鍛鍊.
聖靈的工作.

$^{1561}$和一些的磨練.
其實就真的會.
落得一個驕傲的收場.
有時候那些驕傲.
是一發不可收拾的.
甚至乎他行使的職權.
來顛覆你的教會.
也有.
令到你整個價值觀都錯了.
所以我們要留心.
聖經真的給我們.
很真實的教導.
如果我們不聽.
其實就會帶來一個很大的禍患.
所以在一個很正面.
每一個人參與事務之前.
原來要對自己有一個很忠肯的評價.
所以合乎忠道.
就是一個很忠肯的.
不是主觀的.
很客觀對自己的一個認識.
所以原來參與事務.
對自己認識都很重要的.
動機.
所以要上一些.
事奉學習班.
是應該的.
我在紅恩也開過一個.
事奉學習班有四課.
其實講的不是那些技巧.
是為什麼要事奉.
什麼人才可以事奉.
其實都是很想再一次灌輸.
其實沒有人有資格事奉.
能夠事奉是恩典.
所以我們要帶著一份謙卑的心靈.
去參與事奉.
要看自己.
是否真的恰如其分.
以致我們還要記住.

$^{1601}$是在教會裡面合一的事奉.
不是特別誰大一點.
強一點.
是不同的崗位.
不同的恩賜的配搭.
所以對自己的謙卑.
和彼此的配搭.
就帶來對事奉帶來的一個祝福.
就不是一個影響.
我們就在下面.
另外一個翻譯.
我算直著.
所刺及我的恩典.
對你們中間.
每一位說.
不要高看自己.
超過所當看的.
而要以清醒的心.
來看自己.
客觀的.
真是認真看自己的.
來看.
看完之後.
你對自己有一個這樣的認識.
都不夠的.
你應該讓一些牧者.
有你的了解.
然後按他所給你的.
做一個這樣的客觀正反.
來做一個這樣的調和.
這樣就算是一個更加客觀的.
因為有時自己都未必太認識自己.
當有了一個清醒的心思.
來讓其他.
特別是牧者.
一定是對我們的關心.
對我們成為基督徒.
會有了解.
其實更多.
牧者會為我們祈禱.

$^{1641}$在祈禱裡.
有時他會知道.
這個弟兄姊妹.
是適合這樣去做的.
這個弟兄姊妹.
其實未適合造傳.
有時這種感動.
是來自神給我們祈禱裡面.
有的帶領.
有時這種是.
你可以說我這麼主觀.
他可以馬上做.
他就不行.
你是不是偏心.
他有神的心腸.
有那種肅靈的眼光.
要相信.
這是很難說得明白的.
我們牧者可能就未知.
所以有時我們在感受的時候.
我們開會.
10月,11月.
我們會談新一年的事務崗位.
我們牧者會談.
現在張國豪的位置.
誰誰.
我們會提出.
誰暫時不要住.
另外他還是.
這方面比較幼齡.
先不要參與.
我們會拿出來談.
其實不是針對他.
可能他暫時不適合住.
我們真的這樣談.
所以這裡就告訴我們.
原來.
事務是要清醒.
原來.
一個很榮耀的身份.

$^{1681}$能夠成為主的.
使用的僕人.
真是一個師事.
也是一個僕人.
其他的.
原來在事務之前.
一定要用一個這樣的態度.
清醒的心思.
來看自己.
所以.
很多如果一腔熱誠.
而受到挫折.
受到一些嫌疑.
一些誤會.
他就立刻有一種挫敗.
這個挫敗可能不關他的事.
但可能是有關他的事.
因為他真是被.
經過這樣的洗禮.
他自己原來.
不懂得黃碼書.
不知道原來我們領受了恩賜之後.
我們應該要給我們.
信心的不同階段的成長.
我們就參與不同階段的事奉.
還有在事奉之前.
自己是很明白.
我能夠事奉.
是恩典.
我們是很樂意去做.
做不到不夠能力.
信心還沒有去到那個階段.
那個不是問題.
就會慢慢自己放下.
最重要就是.
不會傷害自己裡面的生命.
因為一參與事奉.
一定會有.
不同的人對你的.
平頭品足.

$^{1721}$這樣的平頭品足.
本來不關你的事.
經過這樣的洗禮.
認真去思考事奉的.
應該有的準備和態度.
這個就是我們實在說得太少了.
有時我們不知不覺.
令到弟兄姊妹.
一開始參與.
很快就慢慢.
有時找一些不同的原因.
慢慢退卻.
原因就是因為.
太多負面或者誤會.
會令到義務參與的弟兄姊妹.
有一種內心的受傷.
所以原來.
先有這樣的心靈的準備.
清醒的心思去看自己.
神竟然感動我.
又有牧者邀請我.
即是說很客觀地.
我現在應該要做一些事.
在教會裡面.
由牧者帶我們去安排去事奉.
這樣一步一步.
可能由淺的.
自立的信心.
慢慢越來越.
因為一路事奉.
客觀地周圍的弟兄姊妹.
都會有一種心靈的.
一種印證.
他真是很適合這個崗位事奉.
他真是參與做得很好.
很忠心.
很默默地做.
不多是非.
這樣的時候.
其實就更加印證.

$^{1761}$你的信心一路增加.
越做越感動.
即是說神又可以再.
加多信心給你.
你再讓更多人感受到.
你的事奉.
是帶來教會的聚會.
或者群體.
是一個福氣.
譬如.
用一些例子.
譬如亞歷斯.
一直都在兒童的工作中.
如果其他的弟兄姊妹見到.
這麼大的男人在做小朋友.
小朋友?.
不行啊.
因為大家這麼多年都.
印證到愛小朋友.
很喜歡小朋友.
很想跟他們講.
福音的故事.
很想帶小朋友進入.
耶穌的救恩的國度裡面.
連大家同居都見得到.
這就是信心不斷地增長.
不斷地讓神去量.
信心加進去.
就變成這個人.
無論他多大多大.
大家都知道他對小朋友.
就是有負擔的了.
根本就是膽子.
你可以想像到.
這就是我們在教會裡.
人家都會發現神在你生命裡面的.
那個恩賜和大人.
譬如有些弟兄.
教初信班.
真是很令那些未信的朋友.

$^{1801}$慢慢明白.
進入那個救恩.
信主裡面.
很印證到.
他的恩賜是來自神.
加給他.
一直這樣去做.
你可以想像到.
這裡主要特別對.
剛剛成為基督徒的弟兄姊妹.
在羅馬的教會.
本來不是很認識每一個人.
但是這些中肯的.
血緬教導.
其實是相當重要.
每一個人開始之前.
我們就要來做.
一個這樣的警示.
一方面教會就要有這個機制.
初信.
真是不可以讓他做太吃蟲的工作.
一信主.
都未夠半年.
洗禮都未夠半年.
馬上做司事長.
馬上做主席.
不是不行的 可能可以的.
但是如果機制不健全.
原來是誰喜歡找就找.
那就不好了.
可能他很好的.
更多就是.
保羅所說.
一洗禮半年馬上做主席.
馬上做司事長.
其實可能對他的生命.
是會有一個影響.
是那個工作.
是那個時服帶給他一個挫敗.
一個影響.

$^{1841}$那就會很嚴重.
因為可能會受傷.
因為可能會自高自大.
這裡就告訴我們.
很容易就會看自己.
不要說笑.
要情願.
要謹慎.
因為一旦事奉了.
有時就很難叫他不要再做.
很多時候是他受挫敗.
受傷才不做.
這個時候其實是一個影響.
對他個人是一個影響.
對教會又是一個影響.
如果他自高自大沒有受傷.
那就更慘.
受傷要跟他療傷.
又是很大的事.
有很大的反感.
對教會的時服.
甚至對神也有誤解.
所以兩者都是不行的.
你可以告訴我們.
這個就是要以清醒的心思.
來先思考.
甚麼是時服.
我是否真的可以在神.
給我的恩典.
信心裡去參與.
他是一個很好的基督徒.
但是他需要去調整.
所以.
初信.
有一個階段的沉澱.
學習.
沉澱是整理自己的神學.
對神學更加的穩固.
對主更加的認識.
不用擔心.

$^{1881}$因為他還有很多年去參與.
所以如果只是開心.
很開心地去做.
很快就會.
這種熱誠會減弱.
這個時候就會出現.
很多的問題.
所以這個量度.
就是要按照神.
分給各人信心的度量.
來看自己.
也就是說.
要在神的眼中.
去看自己能夠成為基督徒.
這個起點.
來去看我們過去的家庭.
對我的影響.
我現在的性格.
我的性向.
我是一個暴躁的人.
我是一個很冷漠的人.
我是一個很怕事的人.
如果是這個崗位.
我是否合適呢.
我有足夠的準備.
無論知識的準備.
心靈的準備.
去做這個崗位.
我們很多時候都沒有去想.
一聽到有人叫.
就做了.
這就是一個最好的部署.
都是這一句.
第三節.
合乎中道的欺凌原民.
我們上一個星期已經說過.
都是謹慎.
自制的判斷.
一個很清醒的.
英文.

$^{1921}$姆努的書信.
就是一個很深的mind.
很深.
好像很清脆.
理智上.
是很客觀的.
很安靜的.
很誠實的.
去面對自己.
不要只是一種熱辣辣的做.
就不行了.
清醒背後就要謙卑.
和很合宜的.
來自我的形式.
我們仍然是這樣說.
第三節.
所以信心的大小.
中間才是重要.
教會中謙卑和自卑的分別.
又是有不同的.
又是有相同的.
所以大家要留意一下.
這一點.
我們在教會裡.
必須要學習謙卑.
要真是在神的眼中去看自己.
我自己也是不斷地學.
我們教會初信.
我們連受浸都要排小椅.
沒有兩年.
都不想幫你洗禮.
你沒有讀過一次聖經.
你都不應該來告訴我們.
想洗禮.
我們這麼嚴格的教會.
當然有它的缺陷.
有它的困難.
有些人根本不再洗禮.
但是那個裝備.
特別剛剛洗禮之後.

$^{1961}$其實我們在元朗堂教會.
都不是很大.
都沒有資格去做一些重要的崗位.
都是不斷地去學習.
所以過程的階段.
有足夠的空間去聽.
讀者的教導.
透過聖經的真理的教導.
讓自己知道.
事奉的那個事奉官.
事奉的態度.
是有準備好的.
變成一洗禮就甚麼都可以做.
這樣就很大件事.
真的會帶來彼此的傷害.
如果他是驕傲.
就會傷害教會.
傷害其他人.
如果他裡面沒有準備好.
沒有一個很充足的信心.
他就會受傷.
慢慢地由熱切的心.
變成一個冷卻的心.
其實很多教會都是這樣.
所以你看到這個圖表.
其實真的要告訴我們.
很清醒.
很謹慎地去思考.
我們到底有多大的信心.
一定是由事奉裡面.
不斷地去累積.
這個恩賜給我們的信心.
所以一定是從參與事奉.
所以由那裡開始.
參與事奉.
強如保羅都要在大掃.
去沉澱.
才成為真的舉世.
影響世界的大仕徒.
他都要有.

$^{2001}$我不知道有多少年.
聖經沒有記載有多少年.
他就在大掃這個地方.
等巴拿巴來找他.
他都沒有周圍的傳道.
他就在這裡.
好像神給他一個安靜等候的時間.
使徒保祿見到耶穌.
尚且要等候.
所以不要以為洗禮就好像很偉大.
一定是偉大.
救恩使我們成為神的兒女.
我這裡說的是事奉這個範疇.
所以作為領袖.
不要見到剛洗禮.
就像發了達一樣.
什麼都給他做.
不是不做.
先沉澱.
他要清醒地去思考自己的生命.
然後那個信心.
神量給他.
因為一定見到的.
不會走寶的.
神量給他的信心.
那個恩賜.
慢慢就被人發現.
慢慢他就會開始流露.
這個恩賜.
他對不同的群體.
他有那個事奉的心.
這個是神.
因為是神量給每一個人.
我們一個土壤.
他栽培.
那個土壤自然會生長.
因為裡面有種子.
就是神量給他的恩賜.
信心的種子.
所以我們很多時候.

$^{2041}$比神快去找人事奉.
我們就出事了.
我們就會帶來自己招致.
有時候要收拾困局.
都難.
因為始終是人.
始終有時你不准我去事奉.
你有沒有搞錯.
或者你叫他不要介意別人怎樣說你.
那些是誤會.
他都不再理你.
不再事奉.
不再熱心.
兩者都是相安.
源於第三節.
我們沒有很認真跟隨保羅.
記住身體的現象.
不是惡了.
活祭.
是要在群體裡面去建造.
這個活祭.
不是你自己拜拜.
我洗了禮了 我拿著一張證書.
將來見吧 沒有這件事.
除非你快死了.
最近這個禮拜.
三個禮拜前.
我不知道.
有沒有在這裡說過.
我們教會有位姊妹.
洪恩堂教會的姊妹.
和另一個家庭.
在飛機上.
搭飛機去韓國玩.
但在去機場途中.
肚子痛.
痛得很厲害 要急症.
進了北區醫院.
輾轉安頓好.
檢查過了.

$^{2081}$原來都是一些腎的指數高.
要禁食.
打點滴就可以了.
但要過幾天.
接著安頓了之後.
發訊息給我們.
安頓了之後 還是痛.
沒辦法了.
然後隔壁床.
有一位婦女.
年紀不大的.
50歲 鬆一點.
見到我們的姊妹.
你可不可以為我祈禱.
因為我很痛.
原來這個女士是癌症的末期.
見到我們的姊妹.
當然很痛.
都跟她一起.
叫她順序.
她原言.
馬上WhatsApp.
我們的群組.
牧師你快點來 她要順序.
她現在很辛苦.
輾轉 上一次.
剛剛開放日.
幼稚園.
正式的.
奠基.
45週年的.
啟動禮.
拍完照 剛剛有一點時間.
因為下午有福曲.
我就趕去北區醫院.
見到那個女士.
她已經說了兩句.
她已經很肯定地相信.
馬上幫她洗禮.
這個是三個星期前.

$^{2121}$接著剪輯了再探望她.
上個星期因為要去土耳其.
我去之前.
星期天晚上上機.
我上機之前下午.
在教會拿了洗禮證.
親自送到她的手中.
她整個人很感動.
因為她拿著洗禮證.
因為她知道自己.
很快就要離開世大.
接著就說.
因為我要去哪裡.
一個星期就探不到她.
我答應了 星期二就回來探望她.
我星期二.
剛過去星期二就去探望她.
已經是一個皮包骨的一個人.
但她見到我之後.
將她夾住的.
那些.
不知道是什麼.
心臟夾住的那些線.
就拔走 兩隻手抓住我.
很感動.
很硬.
其實她已經.
好像集中營那樣身體的樣子.
抓住我.
接著就放開.
告訴我.
她有信心.
我跟她祈禱.
但今天.
她姐姐.
醫院打來.
要我馬上趕她.
剛才來之前都說要祈禱.
所以你就說這是一個活祭.
一個身體的獻上.

$^{2161}$我想說.
就是說.
如果我們一信主就已經要離開世界.
那當然.
真的不需要想太多.
但是更多更多的.
99.9\%的.
都是信了主之後一定要看自己的條制.
一定要將自己獻上.
但是獻上去哪裡呢.
就是教會的群體的建造.
所以第四節立即就說身體.
第五節就說.
基督都是這樣.
基督的教會都是這樣.
所以我們首先對自己要評價之後.
還要將自己評價好的這個生命.
擺在教會這個場景.
不是自己.
是教會的.
教會的.
場景.
不是自己一個繼續自己的生活.
是要在教會裡面.
一齊這樣來到去.
同尊同榮.
同生同死.
這樣的心態的擺上.
的擺上.
所以.
這裡是很重要的.
原來九點鐘.
我只是說到第四節.
其實我今天相當累的.
本來我打算第二節.
對不起 一起祈禱.
肩負我們感恩.
讓我們能夠在今天晚上.
我們可以來到去看到.
保羅的教導.

$^{2201}$其實保羅自己.
都是這樣來到去經驗過.
他可以是即刻就去到全世界.
去報道.
因為有神給他的授權.
那個差遣.
但是我們看到的就是.
耶路撒冷的弟兄.
將他遣散.
去大宋.
在他自己一個無名的小城.
在那裡沉澱等候.
其實他不需要等候.
但是他等待巴勒巴.
神的印證時間.
一方面是大復興的安提安.
需要有同工.
原來神的作為.
真是奇妙.
這位保羅.
不是被投行置省.
而是沉澱的一個基督徒.
在那裡等候.
等候巴勒巴的召喚.
即是神的時間到.
可以被大用的保羅.
求你讓我們在第三節.
十二章.
真是經常要有一種.
恰如其分.
很理性.
這樣來評價自己.
能夠在神的國度裡.
這個葡萄園參與.
事奉是一個很大的福氣.
願主幫助.
我們今晚就.
只能夠講到第三節.
願神來使用.
我們這樣感恩禱告.

$^{2241}$奉耶穌基督的名.
阿們.
\newpage



\section{羅馬書}
\label{sec:0TRElWqGJ4E}
\textbf{2025-11-21 周五晚查經班《羅馬書(28):分題:恩賜是要來用的,在教會的合一場景中運用》 講員:曾裔貴牧師}
\newline
\newline
連結: \href{https://youtube.com/watch?v=0TRElWqGJ4E}{\texttt{https://youtube.com/watch?v=0TRElWqGJ4E}} ~~~~ 語音日期: 2025-11-29
\newline
\newline
\hyperref[sec:B5vFf1eCUBA]{\small{< < < PREV SERMON < < <}}
~
\hyperref[sec:index_chronic]{\small{[返順時目]}}
~
\hyperref[sec:index_scriptual]{\small{[返順卷目]}}
~
\hyperref[sec:WdwoyEOQVUg]{\small{> > > NEXT SERMON > > >}}
\newline
\newline
$^{1}$我們一起同心禱告.
多謝天父我們在晚上再一次來到羅馬書12章第4到第8節.
多方面恩賜和在合一教會的基督的身體裡.
來發展,發揮,盼望神來帶領我們.
讓我們能夠知道這段聖經和整個12章的關係.
也讓我們透過這樣發展我們自己在神面前所領受.
伯伯所得著的恩賜來彼此服侍.
求主帶領今天晚上我們仰望在你的手中.
我們這樣禱告奉耶穌基督的名.
阿們.
第28講羅馬書.
上一個星期我們只是講第3和第4節.
也講到保羅他看自己恰如其分.
神入了一個很特別的經歷.
他曾經很快就在大馬斯加周圍傳道.
後來他認為一個屈辱.
就是大馬斯王把守著城門不讓他出入.
門徒就將他用鑼,筲箕來扔在後面的城牆.
這樣可以暫時解除那個捆綁的壩.
就上到耶路撒冷.
去到耶路撒冷不久又受到猶太教的信徒的誤解和迫害.
門徒又將他送回大肅.
似乎那個地方就不是那麼激烈的反對.
也對耶穌的福音未必那麼認識等等.
所以大肅曾經發生過甚麼我們就不知道了.
聖經是完全沒有記載的.
但告訴我們這是保羅的故鄉.
他一直在等待直至安提摩有大復興.
所以巴拿巴就想起這個地方需要有聖經教師.
有恩賜,有夢召的神的聖徒或者使徒.
來幫他一起服侍.
巴拿巴就立即不是回去耶路撒冷拯救援.
而是去到大肅找保羅.
即是說他一直和保羅保持著.
有可能是宣信的往來.
有可能是知道大家的近況.
去到大肅就找到保羅.
所以你可以想像到保羅在他未真正.
在一個大的宣教旅程之前.
他自己一個那麼有特別的拯救,救恩的經歷的保羅.

$^{41}$有那麼大的託付.
神要他向外邦人傳道.
一個猶太教的領袖.
熟讀舊約的解舊約的一個拉比.
能夠去到外邦傳道.
即便是這樣的託付,積分,恩賜.
還有這樣的才幹,這樣的社會地位.
都需要尋潛在大肅.
都需要來到密集的等候神的時間的道.
來到才可以被神大用.
所以千萬不要小看神的栽培時間.
否則太快有時就不太好.
剛剛昨天說福音序會.
就說到在米顛曠野.
摩西有四十年來牧養他的岳父葉蒂羅的羊群.
有一天那時已經八十歲.
他就見到那些野外的荊棘著火.
又燒不盡.
為什麼會著火又不會燒壞呢?.
這個好奇的摩西就去看.
輾轉再一次讓他經歷神的呼召.
蹉跎了四十年才再呼召他.
這些時間和這些數學都是我們現代人.
未必太能夠明白,接受這樣的蹉跎.
但無論如何保羅在大數有這樣的沉潛.
我們繼續回來看第四到第八節.
這裡為什麼會有亂碼.
我先讀這段經文.
暫時我喜歡用漢語聖經譯本.
新的譯本.
因為它是照原文的直譯.
要告訴我們.
「這就正如我們一個身體上有許多肢體.
各肢體不都有一樣的功能.
照樣我們這許多人.
在基督裡是一個身體.
而且每一個都是彼此的肢體」.
這句是有特別強調那個含義.
再加重了那個解釋.
我們每一個不是獨立的為自己而活.

$^{81}$每一個都是彼此的肢體.
就是說在這個身體裡.
我們好像是被拯救.
放在基督的身體裡.
雖然在地上.
我們很肯定我們在不同的堂會.
在不同的地區.
在不同的國家.
在不同的民族.
我們肯定是有不同的.
但是原來在屬靈的眼光.
在屬靈的角度裡.
我們要這樣看到自己.
是彼此的肢體.
是對方的肢體.
是這個身體的肢體.
這樣的翻譯就來得更加的嚴肅和重要.
因為往往我們要麼看自己很大.
看不起其他人.
要麼我們很容易會自卑.
就會覺得自己沒用.
應該在這個身體裡有任何的參與功用.
但是保羅這裡說不是.
我們正如一個身體上有很多肢體.
照樣.
就是第五節這裡.
Therefore基督的身體都是這樣.
我們很多的基督徒.
在基督裡面.
就是一個的身體.
所以身體的教導比喻.
是相當重要.
我們要有這個很強的意識.
每一個是彼此的肢體.
然後後面就說按照著所賜給我們的恩典.
當然恩典和恩賜.
其實沒怎麼分得那麼仔細.
恩典和恩賜.
其實都挺重疊的.
當然恩賜就可以再去到一些功用上的.

$^{121}$那個實際的應用等等.
然後就說我們各有不同的恩賜.
記住是按照著神所賜給我們的恩典.
以致我們有不同的恩賜.
或者是說預言的.
就要按照著神的預言.
那個信心去說.
就是說原來有程度的大小.
或者是服事的.
就要專一的服事.
強調那個專一性.
專注.
一心一意.
就是說.
原來服事的背後有一個態度.
那個態度要正確.
保羅強調要專一.
就是說.
不只是那個功用那麼簡單.
而是你如何運用你的功用呢?.
在教會裡面.
就是強調的就是服事.
原來很容易會分心的.
原來很容易就會有一些錯誤的想法.
但是當我們專一.
專一就是當時適合的態度.
如果你是清楚神給你這個崗位.
有這樣的恩賜.
你就要專一去做.
專一就帶來你會不斷優化那件事.
你會不斷讓自己在做的過程.
你自己都不會不過且過.
而是你要很放心下去.
相信每一樣東西放一個心下去.
都會很不同的.
就是你要去喝茶.
那個事憶都好像完全不存在的.
叫完都忘記了.
本來那壺茶都不是你們喝的.
是保理的.

$^{161}$拿壺壽比來.
你就會覺得有沒有搞錯.
為什麼這麼不專業.
我們都會覺得很難受.
我們覺得沒辦法接受.
但是你想一下.
如果我們專一的服事.
我們如何放一個心下去做好那件事呢?.
這些就是很不同.
原來保羅.
等一下我們會說的.
因為等一下我會跟大家做一個.
可能平時你沒什麼機會這樣做的.
就是《新約聖經》.
只是說保羅.
我不說彼得和其他的使徒.
只是保羅有三卷書講恩賜的.
哥林多前書的十二至十四章.
已不鎖書的第四章.
都提到恩賜的.
所以等一下就會有幾個部分.
就和大家做這三卷書的恩賜的比較.
是.
就變成我們會多了一點對恩賜這方面的教導.
這裡是專一.
教導也是要專一.
去勸勉.
又要專一.
施予的.
施予的就是你去幫人.
這個施予.
原來就是實心實意.
就是說你裡面的心意.
是要切實的.
誠實的.
是要從心而施捨的.
這樣的一種施予.
也是提升了我們去周濟的時候.
去幫助人的時候.
有心靈的態度.

$^{201}$似乎每一個的恩賜.
保羅都涉及到你內心的動機和態度.
這樣就很不一樣了.
我們從來叫人做事.
我們只是說你這一方面具備.
你這方面又可以有八級鋼琴.
你又可以這樣.
當然了.
事奉之前.
我們會講到事奉的裝備.
事奉的態度.
正正就是這裡了.
正正就是保羅在列舉這裡還有多少個恩賜呢?.
就有七樣的恩賜.
這七樣恩賜.
等一下和其他的裡面有什麼不同呢?.
在羅馬書和哥林多前書又有什麼不同呢?.
我們等一下就可以仔細去看了.
但是這裡已經發現了.
不斷提到專一.
不斷提到按照著你的信心的程度.
不斷提到要實心實意.
然後治理.
治理就是去辦事.
要殷勤.
殷勤就是你很賣心賣力.
用盡你的功夫時間.
這樣去做那件事情.
這叫做治理.
然後憐憫別人.
就要甘心樂意.
出於一份很樂意的心靈.
因為這個也很明白.
因為原來有些人你憐憫.
未必有一個很快的回報.
甚至有時候又會帶出一個反效果.
但是如果你的心靈不是真心樂意.
就會有時候受傷.
變成玻璃心.
對兩者之間都會受到傷害.

$^{241}$所以原來保羅強調的.
在第四到第八節.
反而強調的是態度本身.
剛才已經說了.
殷賜是神按照著祂給我們的恩典.
賜給我們的.
賜給我們之後.
原來殷賜是要用的.
下一個節要說了.
雖然不同.
但是源自同一位神.
我們有時候真的很難去肯定.
和認定我們所得到的不同的恩賜.
都是來自同一位神.
往往有時候就是我的先天.
我的才幹.
我的成長.
我的優越.
你就沒有我這樣.
很多時候我們在那個事奉裡面.
帶來一些誤解,衝突.
我們就很難去認同.
但實際上真的是從神那裡而來的.
首先就是從我們得到救恩.
得到救恩之後.
神會將祂給我們有不同的人.
有不同的恩賜.
而且這個恩賜是必須要發揮在教會這個場景.
才能夠讓這個恩賜.
在多樣化之下.
可以成為一個合一.
這種很美妙的和諧.
所以合一帶來的是教會和諧的成長.
就是那個身體很健康地一直長大.
所以目的呢.
一切恩賜的發揮.
都是為了建造基督的身體.
所以那個是終極的目標.
所以你的存在.
當然有不同的崗位.

$^{281}$會有些顯眼一些.
有些比較隱藏一些.
但無論顯眼還是隱藏的.
最終其實是要榮耀這位基督.
要榮耀他的教會.
不是榮耀那個人.
不是那個人要得到榮耀.
那個人盡了力.
他自己本身就可以得到神的獎賞.
神的稱讚.
但不是那個人得到榮耀.
大家一定要搞清楚這個大前提.
否則我們遲早一定會碰釘.
一定會後退.
一定會受傷.
有時那個傷就是自己隱身退.
或者有時不是自己受傷.
你就會傷害其他人.
所以兩者都是傷害.
是內傷還是向外去攻擊其他人.
所以記得一切是源於.
所以不要小看這些總結.
很重要是要造就基督的教會.
所以保羅避免了比較.
所以根本就沒有所謂使徒的信心講預言.
所以講預言就厲害一些.
大一些.
受神的恩典會多一些.
其實就不應該這樣看.
所以每一個崗位都有它獨特的恩賜的作用.
所以不是比較和競爭.
反而要忠於自己的角色.
發揮到在神面前.
問心無愧的臨離進至.
就變成你做好你自己的管式.
相信在其他的教導邏輯裡面.
就是那些忠心的僕人.
有的要再加給他.
沒有的要從他本身有的搶奪.
那個一千銀子的又懶又惡的僕人.

$^{321}$他本來有的 但他拿著捏著不放.
找個洞藏起來.
就變成了那個是比喻裡面的羞羞愧愧.
相反肯去做肯去投資的.
肯去做買賣的.
一千或者五千就會賺到五千.
二千的又賺到二千.
就是說他不斷去用主人給他的財幹.
他就越能夠發揮可能在另外的層面.
不斷可以被神用.
和角度被神用.
所以必須要忠於自己的角色.
角色就是神的託付.
就是神安排的人生的光維.
希望大家已經信主一段長時間.
當然如果有些是新信主.
而聽到這些經文的很感恩.
證明大家有一個看慕的心.
無論新的舊的.
要盡快在你自己的生命裡面.
認定神給你有什麼負擔.
你就為這個負擔去不斷的努力.
不斷的琢磨.
不斷的祈禱.
以致你能夠發揮得很淋漓盡致.
忠於自己的人生的角色.
去為神為基督的身體而活.
我們繼續下去.
先繼續分析一下.
多樣性就是肢體不同.
功能不同.
但是同屬於一個身體.
所以不同的人.
是有不同的因此的配搭的.
所以教會裡面是應該多樣性的.
有些人是音樂.
有些人是講道.
有些人是默默的做事.
有些人是做一些安慰,血免.
有些人是特別喜歡去做一些施捨.

$^{361}$關心人等等的.
所以是真的有不同的.
在這個不同裡面.
大家要學習同屬於基督.
同為基督的僕人士類.
這樣來按照著神給自己的感動.
或者教會在客觀的印證底下.
安排你或者我一些不同的崗位.
我們就中心地去做好這些崗位.
這樣就明要實了.
所以當然我們可以有自己的選擇.
但是在這個選擇底下.
最好能夠學習在服侍之前.
大前提應該是學習馴服.
對長輩的教導,安排.
有一個馴服.
以及對於那個制度我們都要有一個尊重.
我們繼續下去.
這些都是我們一直都講過的.
這裡我們應該要留意.
剛才講到的恩賜.
是有行動的指引的.
就是說是要做出來的.
那個不是一個頭銜.
不是一個職位.
就給你有一個這樣的冠名.
以致你不是去做事的.
那就不行了.
其實每一個恩賜都是有一個行動的付出的.
那個付出當然有些人真的要落手落腳去跑.
去做,去安排等等.
那些很明顯被人看到他真的在動.
但是有些就算他未必是在動.
他可能是心靈裡面有這樣的裝備.
都是一個服侍.
所以我們要記得.
絕對恩賜不是一個頭銜加給你的.
一個虛銜,一個顧問.
不是這樣的.
恩賜加給我們就是有行動性的一個指標.

$^{401}$是服侍的一個責任.
我們如果明白了.
我們就知道保羅在十二章的《羅馬書》.
很強調教會的弟兄姊妹.
記住教會的弟兄姊妹.
第一節是說我們要將自己獻給神的活祭.
是聖潔的活祭.
是神喜愛,喜悅的活祭.
但是在這個喜悅裡面.
原來就開始要尋找我們的恩賜.
尋索甚麼是神.
要我們去有的崗位去服侍.
所以人生必須有一至兩個的崗位.
是你擺得很重的.
是你由信主一直都有興趣.
例如教導聖經.
我自己就一直由信主那一刻開始.
就很重很重的負擔.
音樂完全沒有.
但是當感恩.
因為教會這樣唱.
你跟著唱.
多不少你都要懂一些.
或者有時候會走音.
但是至少都跟得到.
但是真的不是我的恩賜.
但是變成我們要一直發現.
神在我們的心靈裡面.
是有甚麼是會帶給我們.
好像真是如雨得水一樣的.
如雨得水.
這樣來到很順暢的.
不會令到你很噁心的.
做起來你是很辛苦.
拉頭上樹的.
當然有另外一些未必是你經常做的.
但是你都要參與去做.
所以應該至少有人生一樣的恩賜.
神給了你.
你自己要去發掘.

$^{441}$不會沒有.
你要發掘.
你是神的兒女.
神一定會將一個恩賜給我們每一個弟兄姊妹的.
我們必須記住這是真理.
當然給了我們不去運用.
不去發掘.
不去發展.
不去發揮.
那就可惜了.
那就浪費了我們人生.
就會有這樣的浪費了我們自己.
神給我們的服侍的機會.
因為服侍如果是榮耀基督的生態.
服侍的人自己都會得到滿足的.
得到一份榮耀的安樂的.
所以一定會的.
因為我們是為民主之主而做事.
所以記得不是一個頭銜.
在教會不要要這些頭銜.
我們要的是實際做事的崗位.
那些的責任.
可以說一下.
預言是第一個.
說第二項預言.
按信心的程度說預言.
就是說和謙卑.
避免自高又是拉上關係了.
其他的恩賜.
在其他六個的恩賜裡面.
都是強調那個心態的準備.
你的動機是什麼.
保羅體重的.
你的動機要正確.
那你做起來.
由你生疏.
到慢慢純熟.
這個過程.
是會有一個心態.
來護衛著我們的內心.

$^{481}$以致我們做的時候.
是不會偏離了那個事公.
那個事奉.
背後的原因.
就是心態強於那個成果.
所以基督教的事奉.
就不是按那個結果做目標.
過程比結果重要.
有些人用心做.
未必做得很顯眼.
但是他只要忠心.
他的動機純全.
也要知道神安排他.
要做這個崗位.
而且帶著一份和其他的弟兄姊妹.
合一的心去做.
這樣的動機純全.
他做的恩賜.
就會得到神的喜悅.
當然記住那個喜悅.
不代表會有一下子.
立即的那個結果.
會有很多可以見得到.
量度得到的那個果效.
所以那個內心的動機.
和實際的事奉的果效.
未必需要是相輔相成.
因為有時候你在一些很小的教會.
你能夠盡量在你人生做好你自己.
未必那個教會會有成長的.
甚至乎是一直走下坡的萎縮.
這個未必是關事奉者的參與的事.
不因為他而帶來.
一個教會由小教會變成一個極大的教會.
也不一定是這樣.
當然有時候真的會透過這樣.
教會慢慢會增長.
也都是相當多這樣的例子.
但是記得不是按那個結果.
去看那個事奉的人.

$^{521}$他的動機是否真的純全.
所以這個必須是要苦心自問.
常常謙卑在神面前等候.
來測試自己的內心.
是否已經變了那個心呢?.
心態是否已經有一些雜質不純全呢?.
保羅體重的那個動機.
是否已經有了一個世俗的價值觀.
來取代了呢?.
繼續都是這樣做.
但是做的時候已經不是這樣的態度呢?.
這個就要留意了.
所以你看到有時候提摩太前後書.
保羅相當強調提摩太和提多.
他們作為年輕的傳道人.
他們要追求的.
不要那麼快去想那個高速增長.
是想一下你做的動機是什麼.
站一站提摩太前書.
因為今天想起.
其實預備等遲一點那個講道的.
站一站提摩太前書第六章.
第六章的提摩太前書.
因為我站的是羅盤.
頗重要的.
第十一節.
提摩太前書的六章十一節.
如果你知道上文.
可能我不夠時間放上文.
你就知道這段聖經的重要性.
要提摩太非常要留意他自己的起初的宏照.
起初的呼召.
十一節.
但你是屬實的人.
你就要逃避這些事.
什麼是這些事呢?.
上文就是說.
那些假的使徒.
或者那些傳義教的人.
已經慢慢因為他內心貪財.

$^{561}$他們自己落在了一個敗壞和滅亡裡面.
所以有人貪戀錢財.
就被引誘離開了真道.
因為以前他是在真道裡面的.
但是慢慢就偏離了.
離開了.
被引誘離開了真道.
就用許多的仇苦.
將自己像萬箭穿心那樣刺透自己.
仇苦刺透自己的心靈.
所以第十一節.
保羅立即說完那些壞榜樣之後.
立即就叫提安泰.
你是屬實的人.
你都要小心.
如果你那個心被一些東西引誘離開了真道.
你一樣會被那些仇苦刺透.
所以你就要逃避上文所說的這些事情.
就是那些貪財.
或者讓世界的價值觀進入你的心靈.
令到你感覺不是的.
原來傳道人戴著勞力士.
走出來講道真的有信心一點.
不知道什麼價值觀進入了你的心呢.
那就麻煩了.
你就發現原來你越來越追求.
越來越極致.
一套西裝就不要上大陸做了.
可能二萬多一套會出來好看一點.
那你層次又不同了.
本來大陸一千多一套就很漂亮了.
但是你要追求二萬多一套西裝.
才能夠配得起你在港台出來的效果.
我不知道有沒有這樣想的 牧者.
希望沒有.
那你留意了 保羅提提安泰.
你要逃避這些事情.
那就是說原來提安泰都不是絕緣的.
保羅不在他身邊.
他現在在一個已不鎖的大教會服侍.

$^{601}$會有很多不同的思潮的.
會有很多不同的價值觀.
來邀請他.
來讓提安泰可能會不知不覺.
喜歡了這些世界的價值觀.
然後原來.
怎樣令自己常常心靈都有一種的猶豫.
始終如一的初心呢.
留意了 追求公義 敬潔.
信心 愛心 忍耐 溫柔.
一個傳道人 原來裡面的那個動機.
事奉的心態.
才是常常要力久如新的.
這樣來到什麼.
比如說十二戰.
你是為真道打仗的兵.
你是打美好仗的傳道者.
所以如果我們今天已經.
那些武器已經放下了.
已經在初信的時候就有了.
有那個頭盔 有盔甲.
有用信心那個傳福音的鞋.
預備去傳福音的.
或者有其他的信德銅牌.
有聖靈的寶劍.
全部已經放下了.
因為現在在一個文明社會已經不用打仗了.
保羅說不是的.
你現在仍然是要打美好的仗.
你要持定的是永生.
就是說還沒得到永生.
你都不可以去說我可以退役.
不用再打仗了.
你是那個打仗的將士.
原來一個傳道者.
一生就要這樣帶領弟兄姊妹.
去打這場宿命的美好的仗.
而且是美好的仗.
然後持定永生.
然後就說你是為此而無照的.

$^{641}$也在許多的見證人面前.
已經 曾經作了美好的見證.
那個是起步點.
原來一直都要保持自己有這些追求.
公義敬虔 信心愛心.
就是說整個的心靈裡裡外外.
都是能夠讓神可以去監察.
可以去光照.
可以讓神去清許的一個心靈.
然後說都不夠.
你自己有美好的見證.
你留意一下誰才是最美好.
一定是基督耶穌.
十三節.
我在叫萬物生活的神面前.
並在向本丟彼拉多.
作過美好見證的基督耶穌面前.
去祝福你.
就是為了命令你.
在誰面前命令你呢?.
就是在這位創造的主.
叫萬物來生長的神.
就是說神賜萬物給我們.
如果神要我們在全道.
就是提摩太太.
你要面對有一些貧窮.
有一些困難.
神叫萬物生長.
就是說神會看顧祂自己的僕人.
提摩太太你千萬不要跟隨那些傳異教的人.
他們就是慢慢貪愛錢財.
就是用自己的心態.
用自己的動機.
用自己的方法.
去盡快得到滿足.
那個滿足就是金錢的付足.
金錢其實代表我們就是安全感.
安全感就變成了另外一個神.
拜安全感神.
所以貪財就是要有自己的安全感.

$^{681}$所以保羅就叫提摩太太.
我指著這位叫萬物生長的神.
就是說提摩太太你不用擔心.
你服侍祂.
祂是會知道你的需要的.
叫提摩太太知道.
還有.
還有誰要向祂無愧呢?.
就是基督耶穌.
祂在管彪彼拉多手下受難的基督耶穌.
祂是面對彼拉多對祂的嚴刑拷問.
耶穌都沒有辯解.
都是這樣默默無聲.
都是這樣忍受十字架的苦楚.
來走上十字架.
所以提摩太太.
你的見證有沒有基督那麼難呢?.
有沒有基督那麼痛苦呢?.
所以這樣就真的相形見絀了.
提摩太太就不會將心比己.
為什麼那個傳道人只要放鬆一點.
就可以這麼開心.
這麼活潑.
這麼好玩.
做人做得這麼舒服.
為什麼我做傳道在這個役所面對的.
那些人又不尊重我.
因為我又不是保羅.
我的師父那麼受人尊重.
我現在剛剛出來在役所.
我又不受人尊重.
那些人又看不起我.
那我如果放低一點是不是會好一點呢?.
保羅就意識到.
知道這個是年輕傳道者.
一些其實要超越的門檻.
或者打擊的你要升級.
就必須要勇於面對自己.
因為你真的不是保羅.
因為你真的沒有那些牧羊的.

$^{721}$新孫的那些獵人.
讓人去對你有尊重.
所以你就要用時間去爭取別人的尊重.
所以不能夠快.
所以這個就是提摩太太要來學習.
所以保羅是最大最寶貴的王牌.
所以頂指每人.
如果大家是參與有不同的崗位.
尤其是如果這些信息能夠海外的頂指會聽到.
如果你在教會有不同的事工參與.
如果你覺得已經很灰心.
提摩太太在這個時候其實很灰心的.
她其實因為前孫都已經提到.
你真的有那些胃病.
你不要再喝水了.
可以稍微用點酒.
來去面提摩太太.
這個是五章的經文.
第五章.
第22節.
給人行按守的禮就不可以急促.
不要在別人的罪上有犯.
要保守自己清潔.
因你胃口不清.
屢次患病.
再不要照常喝水.
可以稍微用點酒.
這樣來勸勉.
新聖靈都關注到.
提摩太太其實壓力相當大.
未必身體比較虛弱.
而是信心,歷練,自信心沒有保羅那麼大.
一定會不斷與保羅比較.
保羅說你不要跟我比較.
你要與基督耶穌做比較.
你的見證要能夠對得起耶穌.
耶穌是怎樣來作見證.
他是在本彪彼拉多手下受難的見證.
神指著耶穌來祝福你.
第14節.

$^{761}$你要守這個命令.
就是那個祝福.
「毫不憋惡,無可止則」.
直到什麼人?.
直到耶穌的顯現.
即是說.
要終生到你完成你人生的旅程.
終生到死.
你都要這樣去打美好的仗.
這裡講完了那些被自己的異教影響的人.
和提摩太太受純正的真道影響和呼召.
現在她自己的事奉就要配合這些真道.
所以原來一直對一個事奉的人.
提摩太太的前書.
活生生一個傳道人.
年輕的傳道人.
將他擺在一個大教會裡面面對的各種各樣的.
可能千奇百怪.
可能根本看不起你提摩太太這個人.
但是提摩太太你不需要.
別人怎樣看不起你不要緊.
你自己要對得起基督耶穌.
對得起這一位萬物.
賜你我們享用的神.
我們來去這樣裝備.
所以原來保羅由始至終講恩賜.
他很清楚每個基督徒都會有從神而來的恩賜.
但是他更加擔心更加期望.
用恩賜的人一生要不斷去試驗.
去探問自己的內心的動機.
是不是純全.
是不是偏離了.
所以傳義教的未必他一下子一開始就傳義教的.
這是六章第三節.
如果有人傳義教.
不服從我們主耶穌基督純正的說話.
和合乎真理敬虔的道理.
即是說一個整全的福音.
應該帶來我們有一個很正確的神學觀.
和做人的人生的態度.

$^{801}$那即是什麼呢?.
當然就是第八節所講的.
敬虔加上知足的心.
就是帶來.
那即是說傳義教的人會被自己的教導影響.
原來我們說一些假的似是非的道理.
自己是首先受影響的.
所以這個真的很奇妙.
傳義教的人就會自己慢慢變成異端.
然後他就會逃避那些苦難.
盡量避免.
我不解釋了.
因為這裡不是我們要講的重點.
不過仍然是和傭君恩賜的原則.
是一個完全合一的.
由始至終的保羅.
無論在羅馬書.
在提安泰前書.
叫提安泰.
你不是要發揮恩賜.
盡量如何吸引別人的眼球.
如何讓自己很亮麗.
受人尊重.
最重要就是你自己的內心.
如何裝備.
如何打理好仗.
如何去設定這個永生.
以及你自己的呼召.
默默地在這個呼召裡面.
不斷地進步.
那個進步就是.
追求公義,極權,信心,愛心.
溫柔和忍耐.
彼此是一個很平和的平衡基督徒.
不會突然間很熱心.
突然間冰冷冷漠.
不會是這樣的.
是很平常地成長的.
我們再下去考考哈根.
我們一直在看的是恩賜.

$^{841}$所以這裡告訴我們.
這一段我們今晚看的是四至八節.
這一段經文不是一些理論.
而是具體的一些指導.
幫助當時羅馬的教會.
即是意大利羅馬首都的教會.
在一個多元文化背景之下.
維持教會信徒的合一.
因為十四章,十五章就說.
外邦的信徒和猶太的信徒.
怎樣能夠在文化都有這麼大的差異之下.
可以帶來合一.
因為沒有合一就不會有宣教了.
沒有合一教會就不能夠一起做一個大方向的發展.
所以必須要放下文化的差異.
彼此之間的不同.
那些不同.
原來當我們大家都堅持自己的不同.
專心去做好自己.
而又認定大家是合一.
這個不同是可以堅持的.
那就很奇妙的.
個個都堅.
你又堅,我又堅.
打架的原因就不是各自的堅持.
是不合一的原因.
所以堅持去做好師班的人.
能夠在整個教會裡面是合一.
是沒有問題的.
是可以很調和的.
但是如果你自己又堅持.
但是又看不起自己.
在那個群體裡面的那個合一.
那就出現問題了.
就不是你的堅持出了問題.
是你那個合一的心態出了問題.
就帶來那種的分化.
那種的角力.
那種的都不知道哪個方向才是一起向的.
所以這是一個大家要留意的地方.

$^{881}$亦都這樣說.
在多元文化的教會.
保羅很意識到羅馬教會有這麼多的恩賜.
有這麼多的機會.
有這麼多神給他們的恩典.
所以要怎樣用呢?.
保羅列舉了七樣的恩賜.
很特別的.
待會我們看和哥林多前書的對比.
羅馬書的恩賜是特別不相同的地方.
又可以這樣做一個整理.
首先我們要記得.
我們不要就這樣背完之後.
所有的恩賜只有這七樣.
你搞錯了.
因為保羅根本就沒有告訴你.
只有七樣的恩賜.
因為哥林多前書.
有些是超自然的恩賜.
講方言,翻譯方言.
能夠行二能.
又能夠醫病.
這些是恩賜.
我們想都沒有想過.
但是在羅馬書這裡.
全部沒有超自然恩賜.
全部是可以神給了你.
以及你可以是與生俱來的性格.
接著被神再用.
又可以更加成為更加好的恩賜.
所以特別羅馬書這裡.
不是一些超自然的.
接著.
所以重點剛才講過.
這七樣.
每一樣都有加上一個解釋.
你怎樣用你的語言.
就要按照信心的程度大小去講.
不要講多了.
不要講離開了聖經的真理.

$^{921}$你要做施捨的.
你就要誠實.
你是做勸導,勸教的.
你就要甘心.
等等等等來描述.
就是說重點是態度.
內心的態度.
以及實踐的過程.
接著這裡好像是兩個裡面.
就是關乎言語的.
以及關乎行動的.
這裡是語言.
教導.
勸化.
都是對人的.
是言語方面的類別.
當然行動方面也是一樣對人的.
不過這裡特別是強調用言語.
這樣來服侍的.
語言和教導.
勸化.
都是不同的方式.
令到那個人回到真理.
你語言就是說.
神給你一些特別的領袖.
你要在群體裡面講出來.
就令到大家學習到這個課題.
或者這個道理.
教導就更加有系統地.
將聖經神怎麼看待這件事.
例如我們接下來的30日.
就有「活由祈禱」.
原來紅印堂從來都沒有試過「活由祈禱」.
他們覺得很奇怪.
什麼叫「活由」呢?.
「活由祈禱」用什麼「由」呢?.
「活由祈禱」之後有什麼特別的神蹟呢?.
等等.
後來想了想.
原來我們感受堂先祭.

$^{961}$每年有「活由祈禱」.
紅印堂沒有.
他們真的不懂.
我就整理了一篇文章.
當然收集了宣道會.
過去的牧者寫的.
整理之後.
讓弟兄姊妹能夠明白.
和靈恩有什麼不同呢?.
有一點不同.
因為靈恩強調一定要好的.
不好就是沒有信心.
但是我們宣道會所說的「活由祈禱」.
不是強調一定要好.
而是強調牧羊的同伴.
那個的同心.
那個的扶持.
等等.
所以不一定是立即好的.
這就是我們對「活由祈禱」.
《雅各書》那段聖經的理解.
所以你會看到.
有幾個層次.
就是說不同的行使.
真理.
有預言的.
可能就是用一個講道.
將一個信息講出來.
大家眾人就聽到了.
就去學習了.
這個教導.
例如關乎主在內的.
就去做一次性的教導.
就讓人得到這個概念了.
譬如保羅說.
你們不要以為你們現在.
因為主在內.
那些死了的人首先復活的.
我們還沒死的.
到時候身體變化就結主.

$^{1001}$這些就是預言了.
就是一個在那個可能聚會裡面.
來到去所謂釋放出來的那個信息.
教導就是已經有的那些預言.
有系統地教導.
無論是新信主的.
或者是其他不同的程度.
不同的階層的信徒.
他們有系統地學習.
這些就是教導了.
所以主要是要分班.
是要分不同的程度.
是要按道就班.
可能四堂.
可能是六堂.
這樣來教導一個術聲.
這些就是教導了.
勸化就是教了你之後.
就看看你的實踐.
這樣就再勸導你.
這樣來做了.
所以幾樣東西都是一個包圍.
透過那個服務的人的言語.
將他所理解的真理.
來到去可能很生活性的.
跟一個群體講.
或者是跟個人講.
或者是跟一個小組去講.
這個是言語方面的.
所以教會現在我們都一樣有.
我們有主日港台的講道.
某程度上這個就是預言了.
當然因為新約的書信已經完備了.
就不需要再有突然間.
一個震撼式的爆肚式的預言.
就是我什麼都不知道的.
現在我上台了.
我講道了.
大家等等.
可以了.

$^{1041}$來了.
現在上面正在過了料.
現在開始講預言了.
寧恩派就會這樣.
極端的寧恩派就會這樣.
但是基督教和呼幽派不是這樣的.
就是說全本的聖經已經傳遞了.
傳道者就會按照他在神裡面的感動.
加上他自己的準備.
他自己在書台的汗水.
這樣來準備.
加上他禱告等候這個群體.
這個禮拜他需要什麼信息呢?.
這樣來等候.
接著有了.
自己就好好的整理.
加上分析.
加上例子.
加上生活的見證.
這樣來將那個信息.
好像一個很聚焦的一點的信息.
就打入了弟兄姊妹的心靈裡面.
所以每一次的講道應該都帶來一個的果效.
就是他有什麼的得著改變.
還有他走出去那個聚會之後.
他是有一個心很想.
照著今天的講道.
試一下去實踐.
教導就是已經有的真理.
再去有系統的講明白.
給不同的群體去知道.
勸化就是已經聽過的信息.
你怎樣能夠再勸導他.
要實踐那個信息帶給他的生命.
是了.
那你關懷行動的就簡單了.
這四個.
服侍,施捨,治理,憐憫.
都是很實際的針對那個人.
所以教會不只是停留在聽道的.

$^{1081}$而是人與人之間.
你看到他的需要.
怎樣去幫助他.
施捨不一定是金錢的.
可以是你用時間去聽他說他的困難.
可能是婆媳關係.
可能他的兒子.
剛剛我吃飯就掃WhatsApp.
社區中心就問怎樣處理.
有一個SEN的小朋友.
在課堂裡不斷的像猴子一樣走來走去.
那些老師就已經很辛苦.
就叫他出去.
那個爸爸一直都知道兒子有這樣的情緒.
和那個活躍.
過度的活躍.
在客廳已經又罵又打他.
都沒有用.
就帶來幾方面都受到影響.
那應該怎樣辦呢.
上一堂已經講過了.
但仍然是這樣.
怎樣處理呢.
所以作為那個導師.
你很困苦.
你不知道怎樣辦.
所以這些都是施捨.
這些都是憐憫.
這些就是你怎樣能夠實際面對那件事.
那個行動就是怎樣幫助在那個困境裡面的弟兄姊妹.
所以教會可能只有傳道人.
只有《廣聖經》.
是需要有平日對弟兄姊妹的接觸.
關心.
私底下低調的勸勉.
以及一起同心禱告.
這樣服侍施捨治理憐憫.
就變成了對一個人的身心靈是整全的關心.
又不可以不斷地這樣做.
而沒有講道.

$^{1121}$沒有教導.
這樣是不行的.
所以有講道,有教導.
有勸勉之外.
就需要有人和人之間有溫度的對話.
昨晚我們聽《行在社區的福音》.
陳淑勤牧師給我一個很大的啟發.
他說為什麼教會找不到人事奉.
你叫他.
現在我們弄一個QR Code.
麻煩你自己刊登.
刊登完你就參加做什麼的崗位.
其實我們每年感受堂我們都有的.
列出所有恩賜.
那個效果真的完全一樣.
十多個報名.
還是一向都做的那十多個報名參加.
我現在有一兩個新人會參加.
但原來大部份根本都不會做這件事.
陳牧師給我們一個挑戰.
原來他的教會自己也是這樣.
是逐個弟兄姊妹和他談.
按照著他有的來勸勉他.
挑戰他.
這樣來讓他踏上.
教會聽到陳牧師所說.
就是柯福堂.
由十多個突然間暴增到二百多個.
參加福事.
原來真的要這樣.
所以只是這樣說完.
聽不明白的.
因為大衛的聽.
聽不進自己的.
因為是被動的.
但原來這樣默默的逐個逐個.
在他的需要裡面了解.
因為大部份人.
陳牧師說得很對.
大部份人怕是否.

$^{1161}$不是他不想.
不是他沒有能力.
而是一答應了之後.
原來是由二十歲做到八十歲.
是沒得停的了.
那就慘了.
他就會覺得自己負擔不了這樣的門檻.
跨越不了.
所以有這些錯覺.
所以就很需要有這種這樣的.
在行動方面.
讓弟兄姊妹真的看到.
其實按照你所得到的恩賜,才幹.
以及你的時間.
來去承諾.
所以這些就是需要逐個逐個這樣勸勉.
這樣來去練文.
這樣來去施捨.
這樣來去做一個治理.
就會讓人發現原來我真的可以行的.
原來我真的可以嘗試.
原來嘗試之後行是不要緊的.
這樣的時候.
教會就活潑了.
當教會一活起來的時候.
大家就.
其實不是搶著做的.
而是大家都找到自己在教會這個屬靈的家.
那個自己的慰佛.
我很擔心,很害怕的.
感受變成一個戲院.
我來就是看完這場港獨我就走了.
我繼續舒舒服服我的生活.
我去飲茶.
我去下午打波爾夫球.
我去打乒乓球.
我去看一場戲.
我晚上和太太牽著手去吃飯.
這些是沒有問題的.
不過你更大的人生.

$^{1201}$你應該是在屬靈的家裡面.
找到你自己事奉的光位.
我擔心的就是這個禮堂太舒服了.
那個泡泡太漂亮了.
那個螢幕太好了.
望上十字架令你神魂顛倒.
太美麗了.
你回來就變成一個戲院.
如果這個是戲院就大件事了.
那個是教會.
是神的家.
所以我不知道是不是有二百人.
還是坐在那裡從來都參與事奉.
如果是的話.
我們就是在蹉跎弟兄姊妹的人生時間.
因為弟兄姊妹肯定是被動的.
肯定是有很多錯誤的迷思.
影響了她事奉的心態.
所以為什麼這個行動.
不只是大包圍地講一堂道.
給她一個QR Code.
她就會馬上變成參與熱心事奉.
原來是需要不斷地勸說.
不斷地這樣言論.
不斷地來施捨.
了解她的心靈.
很明白她的內心的困難.
這樣逐個去帶她進入那個事奉裡面.
所以如果我們有二百人從來沒有參加事奉.
從來只是來崇拜.
準時11點15分坐定在這裡.
這樣其實沒有真正的看教會.
是她的宿命的家.
這個肯定是我們牧者未盡的積分.
也肯定是將來神要查問我們.
我們沒有做好這一種施捨,憐憫.
所以今天講到這七樣的恩賜.
其實真的需要全方位.
我們需要有不同的牧者.
不同的性格.

$^{1241}$不同的恩賜的牧者.
去做不斷地彼此的補鐘.
使得最重要是弟兄姊妹.
知道恩賜必須要用在神的家.
他是那個有恩賜的人.
他要用.
所以如果一直坐得久.
其實不健康.
接著每個恩賜後面都有行動的指引.
剛才講過了.
就是要恩勤.
顯示了保羅關心的是服侍的心態.
這個是保羅一而再強調的.
所以其實是一個這麼原創性.
是保羅自己發明出來.
你的外邦人可以參與服侍這位真神.
可以來跟隨耶穌.
跟隨我保羅.
來到新開的新成立的教會.
羅馬教會來服侍.
開始的時候.
保羅強調的不是做做做.
而是強調你的內心的心態.
即是說內功的心法.
所以我們一定要不斷檢視自己.
事奉.
剛才我漏了講得不夠深.
就是在灰心的時候.
記得想起保羅講.
我在這位在.
我讀回這裡.
並在向本雕彼拉多作個美好見證的基督耶穌的面前.
去祝福你.
這個就真的去到很高的命令了.
提安,我祝福你.
我在誰面前祝福你.
我在基督耶穌的面前祝福你.
他什麼時候最偉大的面對苦難.
就是在本雕彼拉多手下的基督耶穌.
這樣的面前祝福你.

$^{1281}$就是說耶穌就是這樣來守住真道.
耶穌就是這樣打美好的仗.
耶穌就是這樣忍受這個秦婦的羞辱.
記不記得那些羅馬的官兵將那些荊荊插在耶穌的頭上.
又吐口水,又打他.
又幫他穿著紫色的袍子.
恭喜猶太人的和蔓.
又拉他出去釘十字架.
所以提安太何時軟弱呢.
當然他可以賞一下保羅.
但是保羅更希望他.
你不要只跟著我.
你要賞的就是賞基督耶穌.
所以事奉者能夠始終都保持這一份事奉的熱誠.
唯一的原因.
那個原動力就是被基督耶穌的愛去推動.
去參與事奉.
這樣就會非常有力量.
繼續下去.
應該就去到這裡.
應該就是我們開始要講這個.
跟其他保羅講的恩賜聖經有什麼關聯呢.
羅馬書應該是早過哥林左前書.
更加早過《已不所書》.
所以羅馬書是強調態度.
以及群體的服侍.
就是說你看到不會有一些超自然的現象.
似乎羅馬書是沒有什麼超自然的教導的.
哥林左前書我就不能夠詳細的跟大家看.
總結了就是處理關於對恩賜的誤解.
那個濫用的問題.
濫用就是哥林左教會有一個錯誤的想法.
以為講方言就是最大的恩賜.
所以就不斷在做恩賜的排名.
哪些是重要一點的.
次要一點的.
但是他們一直強調.
講方言是最重要的.
所以這個迷思.
當然就是有一些不整全的神學.

$^{1321}$或者有人影響.
導致他們會有這樣的錯誤.
所以當他一定要處理這樣的問題的時候.
他就帶出整個恩賜的準方向是什麼呢.
教會有不同的恩賜.
強調愛和在教會裡面有一個秩序.
不是混亂.
不是沒有王管的.
你喜歡怎樣就怎樣.
是不可以的.
所以要留意在聖靈的感動裡面.
原來反而不是混亂.
聖靈的感動是釋放.
內心是有一個開放寬廣.
但是不代表是失控.
但是哥林教會就是有一種失控.
他又突然間講方言.
他又突然間自己胡言亂語地翻譯方言.
各種的情況.
所以保羅在第十四章.
就定了講方言有規範.
只能夠兩三個講.
但是如果講完之後.
沒有人去將這個方言翻譯.
很特別.
他講完之後有另一個有感動的.
可以立刻站起來翻譯.
保羅說如果沒有翻譯.
講方言的就請坐下.
因為沒有人翻譯.
不會幫到其他人.
所以你有方言的恩賜感動要講.
都不是最大.
都請你安坐一下.
是這樣的.
很特別的.
而且又不可以超過兩三個.
就是說要規範它.
就是說就算有感動.
你都可以壓抑你內心的講方言的想法.

$^{1361}$或者是實而來的當下的恩賜.
為什麼呢?.
因為先知的靈是馴服先知.
就是說你自己是可以駕馭你自己的內心.
駕馭不了完全失控.
那他就不是聖靈了.
是邪靈了.
這樣規範.
所以哥林多前書的恩賜.
原來是有解決問題的原因寫出來的.
所以哥林多前書的恩賜.
是有很多超自然的.
可能又讀一點給大家知道.
其實大家自己回去就要做一點功課.
就是臨前的恩賜有哪些呢?.
就像這裡開始寫了.
第27節.
十二章的哥林多前書.
你們就是基督的身子.
並且各自作肢體.
完全和羅馬書所講的完全一致.
一個身體各自有不同的功能.
接著就講恩賜.
第28節.
神在教會所設立的.
第一是使徒.
這裡排列次序.
是排列那些人.
第二是先知.
第三是教師.
其次就是行異能.
等於就是超自然的能力.
異能.
是真的異能.
再次是得到恩賜去醫病.
真的神醫.
不是原來那個.
幫助人的.
治理事的.
說謊言的.

$^{1401}$是真的有講謊言的.
接著29節.
這裡就開始講那個強調追求獨一最大.
那種錯誤.
那種驕傲.
那種爭排位.
那種容不下其他人.
所以保羅在這裡就用一個蛙虎的說話.
就來講29節.
起都是使徒.
起都是先知.
起都是教師.
起都是行異能.
起都是翻謊言的.
這些全部都是恩賜.
但是教會不可以只有一種的恩賜.
而且是要獨大.
就不可以.
你們要切切的求更大的恩賜.
那就是先知的講道.
不過在講之前.
保羅就說.
我現在就要將一個最妙的道.
告訴你們.
就是在教會裡面.
用恩賜的人.
必定需要內心有愛.
愛就是成為了聯絡其他人的.
那個必須的.
那種好像潤滑劑.
那種的條件.
所以就算有多大.
就算整個教會裡面.
個個都可以做使徒.
都沒有用.
因為都沒有愛.
還有教會是不應該單獨一種的恩賜.
是應該多樣化.
是應該跟你不同.
是的.

$^{1441}$所以同工之間真的不同的性格.
是非常好的.
甚至是一定放在一起的.
不能夠一把聲音.
一個人完全一樣的.
這樣才是最美麗的一個調和.
但是大家雖然不同.
但是有愛.
有合一的心.
認定大家是基督的僕人.
基督的身體.
這樣怎樣不同都會帶來合一的.
那這種美麗就會散發出來.
所以這就是臨前要處理的恩賜的問題.
所以你現在有些超自然的.
還有保羅在講恩賜的時候.
好像有些激氣.
有些蛙虎.
什麼都是使徒.
什麼都是先知.
整個教會都是教師.
就是這樣的意思.
這樣就措辭.
你想想多麼羞愧.
可能教會就要追求.
我們是講恩講方言最好的.
保羅應該是這樣.
你保羅整個頭.
整個身體都是手.
千手觀音.
就怪獸.
不可以這樣.
這就是這樣.
已不所說.
反而強調的是那個人本身.
是求領袖的積分.
目的就是這班人成長.
成熟.
是會促進教會的信徒.
能夠在真道上面.

$^{1481}$不斷的來同歸於一.
尤其是四章的第一節.
這個是已不所說.
所以你看到這幾段經文.
已經將羅馬書的功能.
又襯托出來了.
原來保羅講述的恩賜.
是這麼多樣性.
是這麼處境性.
是一定教會的問題.
意味著羅馬教會就沒有這個問題.
不過有這個要留意.
因為是一個多元文化的地方.
就更加需要具備這種心態.
已不所說的第四章.
這樣說.
第四章.
我為主備酬的勸你們.
既然蒙召行事為人.
就當與蒙召的恩典相稱.
凡事謙虛溫柔忍耐.
用愛心互相寬容.
用和平彼此聯絡.
竭力保守聖靈所賜.
合而為一的心.
保守這個心.
是聖靈賜的.
這個心只會帶來合而為一的.
因為當我們追求.
我們是蒙召的時候.
我們就要和其他人.
保守這個心.
原來我們自己就要有這樣的準備.
我們凡事都要這樣.
有這些心靈的價值觀.
可以這樣說.
就是愛裡面包含這些形象.
愛是很空泛的.
但當我們追求凡事謙虛溫柔忍耐的時候.
我們去寬容別人的時候.

$^{1521}$我們用和平和和睦.
這樣去和別人聯絡的時候.
我們就可以保守合而為一的心靈.
然後留意身體又只能有一個.
就是基督的身體.
聖靈是因為我們.
我們可以保守.
我們的身體.
聖靈是因為正如你們的蒙召.
同有的是一個天上永恆的子王.
一位主耶穌.
一個信心.
一個洗禮.
即是說大家都是同樣.
飲於屬天的.
生命的活水.
一位的十.
就是我們這位創造的主宰.
就是眾人的父.
超乎眾人之上.
貫乎眾人之中.
也住在眾人之內.
這樣描述我們的創造主.
無處不在.
無所不能的真神.
這就告訴我們.
我們的眾物主.
是一直和我們每一個人有關.
即是說教會.
我們作為事奉者.
都要有這種信心.
既然他貫乎眾人之中.
教會真的會負責.
我們要信靠這件事.
一邊事奉一邊學習.
兩者應該不會有衝突.
你越參與事奉.
越愛教會.
越緊張教會.
越擔心教會.

$^{1561}$很自然你整個心.
我們牧者就會放在教會裡.
你的人生專注在教會.
有所謂的小小的感受.
你會覺得很自然.
你的人生專注在教會.
有所謂的小小的改動.
有小小的風波.
你就全心全力去防衛.
去保護.
這是人生的積分.
但是不知不覺.
教會就變成了你.
很多以往的老牧者.
有時不知不覺.
由他創立.
一直在教會裡.
變成沒有他不行.
本來應該是對的.
至少一頭半個月.
真的沒有你不行.
當然我們好像說笑.
其實不是的.
其實要嚴肅.
要留意.
保羅一直在說.
這個合一裡面.
我們一個的神.
這位神從過去到將來.
他都是伴乎在我們眾人之中.
他住在眾人之內.
但是他超乎眾人之上.
他超乎眾人之上.
所以我們要記住.
教會是神的.
我們盡了一生的努力.
去保護.
是對的.
但是當我們需要退場.
就真的要撇脫.

$^{1601}$記住教會是神的.
我們要退場.
就退場.
牧者最大最崇高的人生.
就是被忘記.
被忘記.
被人忘記.
是最崇高的.
一個氣節.
一定要被忘記.
你還記不記得.
一百年前.
哪一間最厲害的教會.
的那些牧者.
都可能忘記了.
忘記不要緊.
神會記得他.
神會honor他.
所以被忘記.
因為神.
才是最重要的教會基石.
我們一定要銘記這些.
否則就會有很多困難.
一來我們又不敢扶植新人.
不敢交棒給新人.
又不能夠讓新人.
在一路的成長中.
慢慢在錯誤中成長.
永遠就變成了一個很古怪的教會.
記得完全沒有任何的藥藥所指.
是我們在解聖經.
不過要留意.
我們各人蒙恩.
都是照基督所量給各人的恩賜.
各人的蒙恩.
又一樣.
幾個段落都在說.
恩賜是神給我們的恩典.
我們每一個人蒙恩典.
就會給我們不同的恩賜.

$^{1641}$然後就引述一段聖經.
就是詩篇.
68篇.
他升上高天.
努力了十年.
將各樣的恩賜賞給人.
第九節.
保羅就解釋這個詩篇.
既然間升上.
即是耶穌的復活.
豈不是先降在地下屋.
那麼特別.
降下的豈不是.
受死埋葬.
去到陰間.
降下的誰.
就是遠升天上.
充滿盲油的基督耶穌.
然後.
他升上去之後.
就將各樣的恩賜.
即是耶穌基督完全的.
得勝受敵.
得勝死亡.
當他復活升天之後.
即是神的勝利已經完備了.
不過還未在地上.
但他已經有全面的勝利.
所以他現在在天上.
就將各樣的恩賜.
當然按照著各個時代.
不同的人.
譬如非洲的李文斯.
中國的.
戴德森.
印度的克里威尼.
按照不同的時代.
將這些恩賜.
給不同時代的信徒.
譬如我們今天.

$^{1681}$曾聚貴.
二十多三十年前.
我在福音堂信主.
就將恩賜給我了.
那時候開始信主.
每天在家自己讀聖經.
一年就讀了幾次聖經.
簡直是瘋了.
每天就在家讀聖經.
那時候.
自己都沒辦法想像.
那時候又不會有電話又不會有網絡.
真的一定要拿著一本聖經.
不斷讀不斷讀不斷讀.
這個就.
有不同的神秘的.
機會 恩賜.
接著告訴我們.
留意是第十節.
對不起 第十一節開始.
接著就開始講.
那個恩賜的列舉.
他所賜的.
有使徒.
有先知 是講那些人.
不是講.
那個恩賜本身.
是講使徒 先知.
傳福音的 是那個人.
有牧師和教師.
不是講那個.
教導 所以很不同.
接著原來.
神安排了.
不同的恩賜的.
這些職份 這些可能是.
傳職 可能是義務.
的人.
有這樣的恩賜的人.
留意 他們是為了.

$^{1721}$要承傳聖徒.
各盡其職.
重點是紀錄基督的身體.
以致.
等到.
得以長大成人.
等到我們眾人.
在真道上.
同歸於一.
認識神的兒子.
得以長大成人.
滿有基督長成的.
新娘 所以.
神安排不同的領袖.
按照著.
不同的恩賜.
這些領袖.
各按各職.
在教會裡面.
去服事 當然.
不是一間教會裡面.
必須有使徒.
這裡有人解釋.
可能是由聖經.
耶穌升天之後.
一直到不同的時代.
在初期教會.
有使徒的職份的人.
例如保羅 彼得.
他們真的有超自然的能力.
去做傳道.
當聖經寫好之後.
證典完備之後.
這是福音派對聖經無誤的理解.
就不需要再有使徒的職份.
不需要有使徒的人.
在教會裡面.
所以我們一定不拜靈恩派的第四波.
有新使徒運動.
可以訓練自己成為新的使徒.

$^{1761}$使徒的權威就是.
隨時即興講語言.
神給你的特權威.
你講的東西就是聖經的真理.
你可以說.
如果我們真的喜歡.
都頗過癮.
現在甚麼都不用準備.
我站上台.
那些東西就到了.
我就站上去講語言.
我就是新的使徒.
不是我們福音派的傳統.
不是我們順聖經無誤的.
福音派的傳統.
基督教的傳統.
亦都是相當大的傳統.
所以靈恩肯定是小部份會這樣.
所以現在比較純正的靈恩派.
都不會接受有新的使徒運動.
這大家要留意.
所以可能是不同的時代已經有了.
現在不需要再有使徒.
所以他所設的有使徒有先知.
可能就是歷代未有聖經之前.
因為寫的時候.
未有全本聖經.
所以就會有這樣的教導.
亦都不要抵觸到有過.
今日不需要再有.
因為我們已經有全本.
已經是封頂的新舊約傳書的聖經.
就是唯一的聖經的真理.
神的默示.
給我們的元備.
不需要再有新的啟示.
斟酌的就是使徒和先知.
他們有新的啟示.
甚至他們強調.
可以和聖經不同的新啟示.

$^{1801}$這樣我們就不可能接受.
我們今日做傳道.
我們是將聖經已經有的真理講明白.
這些不叫做新的預言.
的確神給我們當下有這個講道的感動.
但不是新的創造性的預言.
所以這裡就要留意.
這裡告訴我們.
是這班人在教會的領袖.
他們領受了這些積分.
是這些屬靈的成熟.
所以這裡很重要.
成熟帶來的就是教會同歸於一.
成熟的領袖.
可以帶來教會不同的人.
能夠長大成人.
即是不斷有新信主的.
他都越來越成長.
越來越成熟.
越來越接近去接觸.
那些使徒先知.
教師.
牧師等等的積分.
來讓教會在合一裡長大成人.
有基督長成的新娘.
即是大家對真理都很融會貫通.
而且不只是聽.
是有行動的.
不斷去讓更多人感受到基督的恩典.
所以這裡第四章.
強調是這些成熟的領袖.
帶來教會的合一和長大.
所以原來在三卷書.
那種恩賜的描述.
有這樣的多樣性.
很值得我們可以.
很欣賞保羅.
那種他對不同教會的把脈.
對不同的需要.
那種的教導.

$^{1841}$已不鎖教會.
他看為是一間很大的教會.
阿提安是在那裡的.
但現在他坐牢了.
他寫給已不鎖的時候.
他就要強調教會的領袖.
不應該發揮這樣的功能.
去輔助傳道學者.
使教會的弟兄姊妹.
更加要來長大成人.
弟兄姊妹要長大成人.
所以哥林多教會在謊言和肅明恩賜上有爭議.
保羅強調的是愛才是最高的用恩賜的原則.
並教導恩賜.
為著的是要建造教會.
為榮耀自己.
所以他這裡特別提到.
已不鎖書他列出的是積分.
所以不是單純恩賜.
所以這個積分.
就是講到那些成熟的基督徒.
在教會裡面有使徒.
先知.
所以使徒就可以說是我們牧師傳道.
講道.
所以記住我不是分享.
嚴格來說.
上一代的牧者很介意.
什麼分享.
我現在是代表神向你講道.
是正道.
不是分享.
分享就是大家坐下來聊聊天.
所以我那一刻上台.
是代表神講道.
講得好不好是另一個問題.
我真的代表著神.
所以會眾要來聆聽.
所以其實嚴格真的不是分享.
主席報告的時候.

$^{1881}$現在要請曾牧師去正道.
那就對了.
不要請曾牧師現在分享.
我可以不聽你分享.
當然我們不要這樣捉字蝨.
不需要的.
已不鎖的教會.
是有不同積分的人.
他們有具備這些恩賜.
這些恩賜其實就是要.
這樣去合一地.
建立教會成長.
當然是人成長.
是生命成長.
是生命遇到.
好的教師.
好的牧者.
好的血道.
好的施捨.
他能夠得到.
對福音越來越明白.
堅固.
越來越知道自己信什麼.
越來越有信心.
有把握.
得到上天堂.
的功用.
所以已不鎖的強調那些人.
那些領袖.
那種合一的事奉.
但是.
臨前強調.
愛是用恩賜.
最大的原則.
就是一切都是為了.
讓教會被建造.
就是.
今天在整個的教道裡面.
所以羅馬書是.
群體的生活.

$^{1921}$服侍的態度.
很少說超自然的恩賜.
所以並不代表.
羅馬的教會.
沒有超自然的恩賜.
就算保羅他在船難.
二十七章.
《善後行傳》.
他在那個米尼都.
即是今天的馬爾他島.
保羅就上岸之後.
沉船上岸.
扶著木頭上岸.
冬天來的.
有些土人在那裡生火.
有個土人.
遊長.
叫做游士都.
就見到保羅.
拿一條柴.
因為熱.
那些蛇在柴裡走出來.
毒蛇來咬保羅.
保羅一手就甩開那條蛇.
那些土人就以為.
他一定死了.
因為是劇毒來的.
但其實沒有事.
跟他們聊天很治癒.
那些土人完全當下就.
歸服主.
因為沒有可能.
完全接受保羅所講的.
信了主.
即是整個馬爾他島.
就信主.
馬爾他每年有一個叫聖波斯.
就是紀念聖保羅.
所以如果你有機會去馬爾他.
你就會知道.

$^{1961}$有這個故事.
是少林主義27章.
所以不要以為羅馬教會.
是沒有超自然.
所以整本羅馬書.
完全沒有超自然的教導.
哥林多前書就有.
爾忽所書.
都有講.
但問題是羅馬書著重的.
是這個多元文化的教會.
應該在服事上.
有甚麼態度呢.
所以提到服事上.
內心的態度.
例如專一.
專一講了三次.
做的時候要專一.
即是很容易分心.
很容易就會被影響.
很容易就會灰心.
所以就不得專一.
所以教小提琴超自然的.
不是沒有 是有的.
就是今天.
跟大家講的.
剛剛好九點鐘.
我們要完結.
今天我們要講的.
感恩.
我們透過羅馬書第四到第八節.
羅馬書的關乎恩賜.
表面上好像只有七樣.
但我們稍為分段.
就見到在言語方面.
有三樣.
在行動上的服事.
有四樣.
但更加奇妙的.
跟哥林多前書的描述.

$^{2001}$跟爾忽所書的描述.
是完全不相同.
是相當溫和地.
去講述.
參與事奉者.
那個內心的態度.
多謝主給我們有這麼清晰的.
屬靈的眼光.
以及那個.
屬靈的原則.
幫助我們在教會一定要投入.
事奉 以及在事奉的裡面.
一定要不斷去探索.
自己的內心的動機.
更加.
我們不斷要仰望耶穌基督.
是這位曾經.
在半雕彼拉多.
手下受難的基督耶穌.
你是怎樣來承接.
神給你的.
人生的恩賜.
那個服事.
就是要擔待世人的罪.
願主幫助我們.
多謝你聽我們的感恩禱告.
奉耶穌基督的名.
阿們.
\newpage



\section{羅馬書}
\label{sec:WdwoyEOQVUg}
\textbf{2025-11-28 周五晚查經班《羅馬書(29):分題:在群體生活中的自我認知》 講員:曾裔貴牧師}
\newline
\newline
連結: \href{https://youtube.com/watch?v=WdwoyEOQVUg}{\texttt{https://youtube.com/watch?v=WdwoyEOQVUg}} ~~~~ 語音日期: 2025-11-29
\newline
\newline
\hyperref[sec:0TRElWqGJ4E]{\small{< < < PREV SERMON < < <}}
~
\hyperref[sec:index_chronic]{\small{[返順時目]}}
~
\hyperref[sec:index_scriptual]{\small{[返順卷目]}}
~
\hyperref[sec:2B_XWriXVXo]{\small{> > > NEXT SERMON > > >}}
\newline
\newline
$^{1}$一起來去禱告 真是天賦.
我們在晚上 我們都知道牽動著我們全香港.
以至於全世界 一件這麼特大的災難.
來發生在我們的面前.
而且媒體不斷的過去兩天.
不斷的去報告 更新那些資料.
讓我們的心更加痛 因為有百多人.
還有200人還沒有聯絡到.
我們知道裡面甚至很多是老人家.
我們想到突然間那個災難臨到.
很多人都是措手不及.
所以我們很盼望 現在可以做到的.
物資似乎都相當豐富.
金錢的救濟也有不同的機構.
大的財團等等 還有政府來支援.
也有心靈的安慰輔導的社工.
精神科的醫生等等.
心理學家都臨場的去幫助.
天賦我們很盼望 教會在這個時候.
我們成為那個陪伴者 同行者.
我們未必能夠在這個時候.
進入到他們的心靈 這個就是我們平常.
如果能夠是社區的教會.
他們是認識社區 認識教會的時候.
就不會感覺到對教會有陌生.
所以我們都盼望 我們都成為.
社區的一間教會 是居民.
就算他們未信主 都是感覺到教會.
是跟他們有關.
他們現在面對這麼大的.
突然間的慘況 教會是同行的.
所以我們很盼望主.
讓我們有主理的恩典.
很特別 我們今晚就開始講到.
十二章羅馬書最後的第三段.
關乎到一些生活的倫理.
我們個人怎樣去跟其他的教友相處.
還有跟外面未信的人.
當時是一個敵對的羅馬世界.
帝國是對基督徒相當不友善的.

$^{41}$但是一而再叫我們用愛 用善去成惡.
所以相信字面的經文已經是相當難去學的.
所以我們看到斯圖保羅.
他那個淑寧的胸襟 淑寧的眼光.
還有教會要存在.
在當時的羅馬世界是需要忍耐的.
是需要用溫柔的心的.
所以我們很盼望我們不會變成極端.
也都不會變成冷漠 冷淡.
我們盼望我們能夠跟未信的人同行的.
今天晚上我們等候在你的面前.
我們這樣仰望禱告 奉耶穌基督的名 阿們.
好 頂曉會萬安.
我們繼續第三段落的羅馬書的研讀.
我給的題目就是.
因為第三到第八節.
我們都用了差不多兩課的時間.
來跟大家總結了一個恩信青兒被主的恩典.
來拯救弟兄姊妹.
就應該在神的憐憫下.
看到自己是應該奉獻自己給神.
作為一個活祭 是神喜悅聖潔的活祭.
這個時候 第三節就叫我們.
首先不要看自己過高.
要來謙卑的 恰如其分的.
按照著我們靈情的成長.
這樣來參與在神的國度.
很快就提到教會這個群體.
很快就已經告訴我們.
信了耶穌不是一了百了.
不是自己一個人這樣來生活.
而是要加入一個地方的組織.
稱為耶穌基督的身體.
當然那個是一個很屬靈的概念.
就是說很小的每一個基督徒.
就是在基督耶穌的身體裡面.
但是這個身體裡面.
也都代表著我們實實在在.
是要參與一個地方.
一個社群的教會.

$^{81}$所以羅馬的弟兄姊妹.
他們如果是由使徒函主的第二章.
在那個五旬節信主.
當時一些人上去守節的猶太人.
信了主 得到聖靈.
有感度回到羅馬去設立教會.
開始有自己的聚會.
就是在羅馬的一個地方.
不知道是哪一個的省市.
那裡設立了一個聚會點.
那個聚會點就是奉基督耶穌的名字聚集.
好像當時猶太人的會堂那樣.
但是要分別的不是信獨一的耶和華.
而是相信三一真神.
這個就是很清楚的界分.
所以地方的教會.
是當時無論得時不得時.
都是需要來到去行常的聚集.
偶然間使徒函主告訴我們.
在耶路撒冷的時候就是天天的聚集的.
天天的來到去準備.
好像神的感動.
或者神的差遣未到.
後來就用的壩令到他們要分散.
必須要走.
就一直去到第八章講的.
就是去到撒冷亞,龜撒尼亞等等.
然後去到土耳其的安提摩等等.
然後再由使徒保羅和巴拿巴.
正式有聖靈的差遣.
就開始一個宣教的一二三次的旅程.
這樣來到去有系統的.
由教會差派的去宣教.
這裡就是聖經告訴我們.
意味著當時的教會.
剛剛成立的教會.
還不是一個很大的,很嚴峻的.
很公開的一個壁壩的時期.
所以教會就開始在不同的地方.
來設立聚會點.

$^{121}$早些你信了主.
你就會去那個地方.
所以在羅馬.
等一下我們不是等一下.
是下一次.
我們就會看到第十三章打後.
就有猶太人和外邦的希臘人.
或者其他種族的外邦人.
都是在羅馬的教會一起聚會.
那就是另外一種倫理文化.
彼此怎樣來到去.
又守著自己的傳統.
但是又要尊重別人的傳統.
但是又在這兩個傳統底下.
不要抵觸,不要論斷.
不要去爭拗.
反而不要讓真理說沒所謂.
你喜歡怎樣做就怎樣做.
就不是這樣的.
所以這些真是有一個很大的.
保羅的教導我們.
所以我們來到這個群體的教導.
所以第九節一直到第二十一節的.
第三個段落.
就是告訴我們每一個基督徒.
在一個群體的生活裡面.
個人應該怎樣來到對自己的生命.
有一個認知的過程.
就是怎樣和其他人,和外人來到相處,交往.
這個就是第三個段落.
想要探討的.
我們一路進入PARAPAN裡面.
跟大家來到去看看.
很快的轉折點.
說了群體的恩賜彼此.
服侍彼此來到去學習的時間.
接著很快就繼續說到關於.
個人的生命的操守.
那個態度,那個內心.
好像保羅已經意味著.

$^{161}$基督徒是擁有一個新生命.
而這個新生命是有一個很大的神的愛.
這樣來到去在我們的內心.
所以說完第八節.
那個恩賜彼此的協調去用的時間.
我們看第八節.
或者作勸化的.
當專一的勸化.
就是去做一些的取消.
去幫助一些有困難的人.
接著施捨,周濟,幫人.
可能是金錢的.
可能是一些時間的付出的.
或者是一些生命的付出的.
當然說到就是.
似乎金錢上會比較明顯些,比較多些.
就當誠實.
就是說要做在光明當中.
不要沽名釣譽.
不要假道學.
不要只是為了一個好像打卡.
是一種的炫耀.
或者是來去博取所謂的呃Like.
古代人不知道什麼叫呃Like.
不要只是博取別人覺得你是很有熱心的.
就是說要很誠實地.
很低調地來默默地去做.
不需要打鑼打鼓.
這樣叫誠實.
接著治理事情就要殷勤.
就是說任勞任怨.
以及要動聲.
要做一些事情.
一定有紛爭的.
比如耶路撒冷也曾經有過.
希臘的婦人和希伯來的婦人.
因為言語的因素.
言語還是其次.
而可能會覺得有口齒薄彼.
當一些治理的時候.

$^{201}$總是會有人覺得.
是有一些快慢.
或者多少.
或者用廣東話的大細超.
就會很容易.
很快就受眾就會有怨言.
所以當治理的人必須要有一個付出之後.
還要有一種的殷勤.
就可以讓那些事情盡量持平.
客觀.
以及均勻.
不要讓人有一種誤解.
接著憐憫的要甘心.
甘心樂意這樣做.
也某程度上真的反映了那個心態.
因為被憐憫的.
可能又會有很多的反效果.
所以當帶著一個甘心樂意去做.
做的目的.
都已經有了那個心理準備.
就不會這麼快因為可能少許的誤會.
就會自己都受傷.
而自己都退卻.
那就不是這麼好了.
所以第二之間.
說到這個群體之間.
彼此的相交之後.
好像一直都在說教會.
地方教會的setting.
那個內圍的一班信徒.
似乎只是很想告訴我們.
只需要在教會工作就可以了.
在教會不需要在教會外面.
但是第9到第21節.
就正正告訴我們.
有外人的.
而且是有些反對福音的外人的.
而且教會是要面向外人的.
就讓我們看到原來教會.
從保羅的世界神學.

$^{241}$已經是無牆教會.
已經是不要自己去規範了.
一個自己的尼巴.
就沒有了其他的淪陷.
外面的不信的人.
他們的救恩的情況.
怎麼辦呢?.
所以很寶貴.
雖然是短短.
很多好像很格言式的.
很短的精句.
每一句其實當我們要真的去實踐.
從來都不容易的.
所以這些倫理的教導.
一定是在一個很重要的基礎.
就是必須清楚自己是蒙了救恩.
恩順清義.
是在神的憐憫下.
我們是一個很虧欠神恩典的人.
所以我們必須不斷去更新自己的心靈.
我們是一個活祭.
我們在神面前.
我們是呈獻給神.
呈獻給神就是要用在教會上彼此的服事.
教會成為一個平台.
去祝福我們信主的弟兄姊妹.
另外當我們自己有一定的準備.
我們就要祝福外面世界拒絕.
未明白,未認識福音的人.
所以這個就是保羅在第十二章的.
活祭的生命裡面.
帶給我們有一個這樣的世界觀.
新的世界觀.
新的我們基督徒的新生命.
所以很快就說愛一定是從神而來.
所以基督徒就能夠用真誠去愛不可愛的人.
所以我們不能夠控制虛假的心靈.
我們不能夠去改變其他人.
他們用一些虛偽,用一些有機心的交往.
來欺騙我們基督徒的愛心.

$^{281}$所以當我們去處理這些倫理經文的時候.
我們就要知道保羅很想讓弟兄姊妹基督徒.
羅馬的教會要有一個這樣的重要的前設.
不斷說的就是我們自己不能夠有虛假.
我們要厭惡一些惡事,神不喜悅的事.
然後善,善就是美善.
goodness,好的事.
我們就要持守,愛主裡面的弟兄姊妹.
就是說可能跟你不同性格.
可能現在所謂的愛人.
現在有所謂的外向的依人等等.
都是心理學上給我們人有不同的性格,品格.
有不同的做事的方法.
都是不斷告訴我們.
原來真的需要有不同的藝術.
需要對人有不同的認識.
這樣我們就不會撞板.
但是這個告訴我們.
當我們要去實踐這個愛的時候.
自己是真心實意很想去關心,去愛其他人.
但是記得是很容易受傷的.
所以我們自己首先有一個不虛假的愛心.
但是當我們去愛其他人的時候.
我們就要有一種彼此的相親,相愛.
就是說仍然待人以禮.
給人有空間.
這些真的很重要的.
我成長都是在一個很沒有愛心的家庭.
大家一直聽我說過.
13個兄弟姐妹.
出生的時候我已經是排第九.
我還有4個弟弟妹妹.
16個減4個.
我家裡一出生我就12個.
你可以想象得到.
接著我陸續看到我的弟弟妹妹出生.
16個人在一個木屋區.
其實是完全沒有什麼所謂愛心的感覺的.
成長是很冷漠的.
還要是一個木屋區.

$^{321}$很多犯罪,販毒,黑社會等等.
又窮.
社會又窮.
家裡就更加窮.
你可以想象得到.
從小在小學都是這樣掙扎.
中學也是在一個班班的學校.
就是你不兇就會被人欺負.
就是說整個環境.
是令我很壓門的人.
是完全不懂得去憐憫,關心.
對人有一種了解,同理心.
沒有的 完全沒有的.
家裡只會鬥快搶吃飯的.
嫌宋是鬥快的.
最喜歡小朋友就是最喜歡打風.
因為那些問題與我無關.
但是開飯與我無關.
就是開罐頭.
最喜歡開午餐肉.
鬥快還沒裝半碗飯.
馬上就夾了兩塊.
夾多一塊都會被人罵的.
但是在你家裡就沒有了.
你就可以想象得到.
這個就是在這樣的家庭長大.
但是保留給我們的教導.
真的是非常寶貴.
要我們真是這一種的愛.
是有同理心的愛.
不是說你強迫別人做任何事情叫做愛.
而是你要令到對方舒服.
要令到對方不會受傷.
要令到對方感覺到.
這是一種很有禮貌.
很有愛心的一個彼此的付出.
是彼此 是雙方的.
所以單是一句話.
已經是讓我們可以感受到.
這一種這樣的相親相愛.

$^{361}$你可以想象得到.
別人走過叫聲你.
如果你不喜歡他的.
理他也不理.
這樣叫做弟兄姊妹.
你嚇死人了.
在教堂裡面.
大家已經是神的兒女了.
你不是說有些人現在是迫不及待教他的回教徒.
或者是一些人.
不是的.
大家是洗禮的弟兄姊妹.
還要大家一起走去擺餅.
你可以想象得到.
你可以想一下.
你想一下神怎麼看我們.
在聖餐聚會.
在感受這麼遠.
走到外面拿聖餐.
你看到某個弟兄某個姊妹.
然後立刻勾起你那個仇恨的心.
不喜歡他的心.
立刻板著臉.
看都不看他.
然後一直走過去拿聖餐.
然後很開心地吃聖餐.
神怎麼看我們這樣的虛假呢?.
這樣的舉動呢?.
這樣的不願意饒恕的舉措呢?.
如果做物主.
猜顯他的愛子.
這樣無條件的釘十字架.
而我們對一個人.
還要是我們的弟兄姊妹.
一個人.
我們視他為我們的.
好像眼中釘.
最好他不要出現在我面前.
這樣的就不是這段聖經的教導.
可以想像到一對一對的經文.

$^{401}$好像一些牌匾一樣.
愛要彼此的相親相愛.
尊重人.
就是對人有一種敬重.
欣賞.
總有每個人的長處優點.
來彼此有一種這樣的欣賞之外.
彼此一種真心的對別人的一種讚賞.
訴諸於口.
由心靈發出的對別人的欣賞.
相當不容易的.
因為你稍微對一個人有一個意見.
就會累積那種不喜歡.
這種不喜歡不斷地增加.
你就不斷看到他的過失.
這個聖經讓我們看到.
保羅很想我們在因此彼此的服善之餘.
我們都要顧及我們自己內心的心靈.
我們自己的內心的品格.
是不是真的具備愛呢?.
所以整段的第九節.
其實就是講一個有神愛的弟兄姊妹.
如何與其他人相處呢?.
接著我們就做一些分析.
這個告訴我們.
在這三個段落的羅馬書十二章.
可以大概分成幾個段落.
第一到第二節.
就是一個不斷深意的更新.
不斷重新的生命.
而且要認定是屬於神的.
是要向神帶領的.
當然這個不斷讓神帶領.
在我們一路成長.
就會不斷有更新的.
神帶領我們有時候去到一些際遇.
去到一些環境.
去到一些崗位.
去到一些人生不同的階段.
那種精彩.

$^{441}$那種騰訊心.
要讓我們到底站到信仰的哪一個台階,階段呢?.
這一種不是抱高抱遠.
而是神會開啟我們的心靈.
我們的生命會越來越成長,成熟.
這樣的.
所以越來越成熟的基督徒.
就是看自己.
真的很心悅誠服地奉獻給神的活祭.
讓神去帶領自己的人生.
這個就是很重要的基礎.
就是每一個基督徒.
所以保羅如果這樣寫完.
而羅馬教會.
真的這樣來看自己是一個不斷更新的活祭生命.
隨時可以讓神去帶領的.
這樣的時候.
神可以在羅馬教會做很大的工作.
因為整個羅馬帝國都沒有福音.
有的是拜裡面的羅馬的巴特隆神.
獨一的神.
就是他們的宙斯.
已經是整個歐洲相當普遍的了.
你可以想像到.
現在基督徒有一個這麼嶄新的生命觀念.
就是這個人信主的這個人.
就是在神的喜悅裡面的活祭.
這個是不得了的.
每一個基督徒在自己信主.
是一個很大的突破.
你一切的過去.
無論你的原生家庭.
怎樣讓你成為鐵石心腸.
我是一個例子.
真是鐵石心腸.
原來是慢慢經過數十年.
很多的事奉.
很多的接觸人.
你就會慢慢有一份同理心.
慢慢自己又更新.

$^{481}$所以我的事奉帶給我自己的生命.
是一個很大的感恩.
因為真是不斷去學習的時候.
你就會有更新有改變.
原來真的可以改變到的.
真是不知為甚麼.
去到這樣的年紀.
小朋友又很喜歡我.
抓住我不讓我走.
老人家又很喜歡我.
很開心的.
因為裡面就有一個這樣的生命.
以前不會喜歡人的.
黑口黑面.
但現在不是了.
所以這是很重要的基礎.
每一個基督徒在保羅的內心.
你信了耶穌.
你就有一個虧欠神的制.
你必須要將生命獻給神.
要甘心樂意地獻給神.
然後在群體裡面.
要看自己是基督的身體的一個器官.
所以整個的生命.
是要讓基督的身體得到榮耀.
所以我們的貢獻.
是應該整個身體去看的.
所以不要凸顯自己的偉大.
凸顯自己的狂妄.
自己的那種驕傲.
那種的懶惰.
千萬不要.
要慢慢看到自己的功用.
是在整個群體裡面.
發揮那個作用的.
所以每一個都應該有自己的功用.
所以這個就是開始.
第一個階段就是.
自己是活祭.
然後就求神.

$^{521}$神啊.
我有什麼可以被你用呢?.
當然你要裝備.
所以為什麼不讀主一學呢?.
最好又沒有那麼多事奉的時候.
你多回一個聚會.
你不回都可以.
你自己能夠有追求就可以了.
你自己可以看不同的書.
你自己又能夠聽不同的道.
你自己又不願去有進修.
那沒問題的.
主一學只是其中一個教會提供的培訓而已.
你有其他的培訓就可以了.
但問題就是很多時候都沒有.
變成我們自己.
未能夠很快很迅速的去知道.
神怎麼用你的生命.
你沒有去發掘你自己的恩賜.
你沒有去向神祈禱.
因為你越參與教會.
你就越看到那些需要.
這裡需要人.
那裡又需要人.
這裡有這麼大的需要.
那裡又有這麼多的需要.
我們如果只是按我們自己的步伐.
我們就不會有這份感動了.
也不會有這份大靈.
我們一個月這樣過.
很快就12月了.
很快就2026了.
很快我們就去到7680.
很快我們就要見主了.
那我們的人生就這樣過了.
所以如果我們連一個自己的教會.
一個地方的教會.
我們是整個基督身體裡面.
其中一個的肢體.
我們都還沒有學會.

$^{561}$我們都還沒有去發掘自己的恩賜.
那也是蹉跎了我們在人生那個.
可以大大被神用的機會.
因為是真的需要發掘的.
所以教會是一個很重要的平台.
是神設立的平台.
不單是敬拜神.
而是我們要培養弟兄姊妹.
發掘自己的恩賜.
還有一個心態就是.
不是你最大的.
你必須是要跟其他的肢體.
成為一個會合.
這樣來讓整個身體.
是每個星期很正常地運作.
那教會就健全.
教會健全就是一個見證.
就是成為當時羅馬世界一個重要的組織.
去告訴這個世界.
我們這個平台是有信仰.
是可以獲得愛.
是可以經驗神.
可以知道永恆的生命的實在.
就是一套神學.
就是世人需要的神學.
世人有他自己的神學.
就是我要拜羅馬皇帝.
我要拜宙斯.
我要拜其他的神子.
我要來到怎樣等等.
但是基督徒就有一個嶄新的神學觀.
嶄新的世大觀.
所以平台就是要告訴世人.
我們不只是在這裡唱聖詩那麼簡單.
我們這個平台是不斷培育弟兄姊妹.
發掘自己的恩切.
在神的聖靈的賞賜帶領下.
讓我們這個人可以在基督的身體裡面.
越來越清楚自己的角色.
那個功用.

$^{601}$因此提到的功用上面.
是有不同的.
但是都是為了那個身體.
最後就是三個倫理的關係.
因為十三章很快就講到.
跟社會和政治的群體.
當時就是羅馬帝國的關係.
所以現在先講了個人有這樣的意識之後.
教會這個重要的平台.
怎樣能夠發揮得好呢?.
相信也是在恩賜裡面.
在愛裡面彼此的服侍.
然後就是第九到第二十一節.
就是告訴我們這是個人的倫理.
在整個群體裡面.
我們怎樣去活出這個新生命呢?.
所以跟大家先讀讀第九到第二十一節這一段.
很短的.
但是相當寶貴.
每一節都是很值得我們去思想的.
我又做了一點點的分段.
這裡告訴我們.
保羅不是單純地列舉道德的規條.
不是說基督徒你要這樣.
你不准這樣.
我初信的時候是一些很嚴謹的福音堂.
所以很懷念和很感恩的.
因為初信它給你很清晰.
不准這樣 不准那樣.
你做好就好了.
傳道人叫我們走過戲院盡量要繞圈.
不要走過那條街.
就是你要避開.
那些戲院你都不要去.
走那條街都不要.
我們初信什麼都不知道.
又傻.
什麼都聽話.
但是慢慢慢慢發現.
其實這種可以不可以.

$^{641}$隨著我們成長.
就變成一種律法主義.
其實未必真的幫到那個人.
反而這裡就要留意.
不是一種道德的規條.
基督徒不是要死守那些道德規條.
而是要信徒在日常的生活裡面.
來活出福音的果效.
所以其實不應該離開團體.
其實不應該離開世界未信的那些世人.
我們應該和他們一起在那個社會裡面.
但是我們怎樣能夠在那個社會裡面.
而我們能夠活出那個福音的那個果效.
就是剛才我們上次講的.
就是愛的生命.
怎樣能夠具體的讓人感受到呢?.
比如我講一些例子.
就是這兩天的事.
我都有問到宣道會大報堂.
趙家石牧師.
都是曾經和他一起在英國開會認識的.
都邀請過他來感受和講道.
我早兩天就問他的情況.
很快就知道他們的教會.
和當時大報的社區很快第二天.
就是前天的下午就已經組織了聯盟.
和開放教會.
開放了一個他們的類似社區中心.
給那些災民來休息.
很快就已經讓居民.
那些災民知道教會和他們同行.
就是說這個意味著.
他們過去一直都有這個運作.
不是突然之間的熱情.
教會都不認識你.
外人更加不知道你的教會存在.
這樣就來得很噁心.
說平時都不關心我們的需要.
如果突然之間又說出現來關心.
那麼你的居民就會有一種反感和鄙視.

$^{681}$相反如果我們一直都是社區的教會.
沒有缺席.
當然不是說和他們做一些拜偶像.
犯法 不合真理的事.
不是說開放在這裡大家可以買六合彩.
不是這些 而是真的和居民有關的事.
我們都很想他們知道我們是關心的.
這樣的情況.
例如很簡單.
整個感受的社區.
三萬多人 五千戶.
是很擔心那個水管的漏水問題.
那麼是不是我們都應該.
和教會和社區有一個聯繫.
我們不是要偏幫誰.
而是居民的擔心.
我們是不是應該在我們有的資源底下.
我們去例如禱告.
我們來協助去讓這件事能夠透明.
各方面這些教會是可以做的事.
因為我們本身就在這個社區裡面.
所以這裡保羅很想強調的就是.
信徒怎樣在日常生活.
所以9到21節差不多是講日常生活的倫理.
就是基督徒應該怎樣去行動.
我們的品格應該怎樣.
去活出福音的果後.
先讀一讀.
愛人不可虛假.
惡要厭惡 善要親近.
這是舊的和合本.
愛弟兄要彼此親熱.
這個一定是講教內.
恭敬人要彼此推讓.
殷勤不可懶惰.
要心裡火熱.
常常服侍主.
在指望中要喜樂.
在患難中要忍耐.
禱告要行切.

$^{721}$聖徒缺乏要幫補.
客要一味的款待.
就是有客人來.
當時沒有那麼方便的旅館.
或者其他的社會的州制.
需要有一個對其他人的接待.
可能就是家裡的接待.
當然整個社會整個時代不同了.
我們要詮釋.
我們現在比如香港和美國都不同.
大家的處境都不相同.
但是那些原則就是相同.
第十四節開始講教外的逼迫.
就是說我們越做得好.
越有利好見證.
你要預期那個反對就會更加大.
我們應該怎樣去處理這些處境呢?.
逼迫你們的要給他們祝福.
只要祝福不可詛咒.
就作提了兩次.
就是重疊地提.
對希伯來講就是重中之重.
與喜樂的人要同樂.
與哀哭的人要同哭.
這個同理心我們明白了.
所以昨天陳淑娟牧師說.
現在大埔的災民不是需要物資.
是需要你坐在他旁邊陪他.
讓他去講一下他的痛苦.
他的擔心.
這樣就已經是一種同理心了.
所以與喜樂的人同樂.
所以現在這個階段肯定.
不可能你要這樣去開心喜樂.
相反現在這個時候.
是哀哭的時候.
我們就應該要同哭了.
所以很欣賞一些大學生,中學生.
都來自發在網上呼籲.
網上去做一些分流.

$^{761}$還有去到現場.
大家來自不同的社區.
這樣去協助災民.
這些這麼自發性的.
又是這麼有組織的.
就是很寶貴的一種同哭.
所以大家按照大家有的恩賜.
去發揮這一份同哭.
都很重要的.
接著第十六節.
要彼此同心.
不要志氣高大.
倒要苦就卑微的人.
或者卑微的事.
不要自以為聰明.
不要以惡報惡.
眾人以為美的事.
要留心去做.
若是能行.
總要盡力與眾人和睦.
很明顯就是在說.
對著郊外的人.
他們對我們的信仰.
對我們的為人.
是有誤解.
甚至是有挑剔.
甚至是有針對.
我們應該怎樣去.
跟這些人相交.
來相處.
很防守性.
好像很退縮.
但是.
因為正正我們不能夠.
改變別人怎樣想法.
必須是神.
藉著長時間的禱告.
和我們生活為人的.
始終如一的見證.
才可以改變一些.

$^{801}$頑固的.
硬身的人.
所以.
第十八節說.
若是能行.
總要盡力與眾人和睦.
親愛的弟兄.
不要自己伸冤.
寧可讓步.
聽憑主路.
因為經常記著.
申命記三十二章.
主說:伸冤在我.
我必報應.
所以你的仇敵.
若餓了就給他吃.
若渴了就給他渴.
因為你這樣行.
就是把炭火堆在他的頭上.
周原有這樣的教導.
不知道實際.
那個典故是什麼意思.
怎樣將炭火堆在頭上.
炭爐還是一塊炭.
是燒著的炭還是怎樣.
一定是當時古代的典故.
就是說.
只會.
他拿個石頭.
扔你.
你仍然用愛心.
去包容.
去感化.
這樣的意思.
所以第二十一節.
再一次強調.
不可以為惡所勝.
為惡所勝就是.
以眼還眼.
或者好些的就是.

$^{841}$不跟他計較.
但是普羅說.
反要以善勝惡.
這裡大段的經文.
是一串論理.
我們怎樣來自處在神的面前.
或者一個教會.
處在羅馬帝國.
應該怎樣去面對呢.
在耶路撒冷可能很遙遠.
的確相當遙遠.
耶路撒冷發生的事.
未必馬上來到羅馬.
根本就沒有即時通訊.
根本就是需要好幾個月.
需要有人真的將這封信.
帶到來才可以讀到.
所以.
這些論理教導.
就是普羅寫給羅馬教會之後.
他們都要好幾個月收到.
他們就要來揣摩.
學習.
很多可能現在.
處在羅馬帝國.
普羅的年代就是尼祿作王的時候.
一個相當自大.
封編.
未必很有才幹.
但肯定是一個當代的暴君.
又是和他們當時的制度.
羅馬的元老院.
是有相當的張力.
因為尼祿都是.
根據教會歷史.
是一個很大喜功.
很想自己在他人生.
有些建樹.
因為他的祖父.
都是很偉大的君王.

$^{881}$加里居拉.
羅馬的大帝.
所以有時都會.
感覺自己不如自己的先祖.
很想有些建樹.
最容易就是大興土木.
所以他的後期.
就是想建一個新的羅馬城.
當時的元老院.
說他是很好勝.
好大喜功.
禁止他這樣做.
他就很簡單.
用一些詭計.
收買一些人.
可能是黑社會 爛仔.
就燒了那個羅馬城.
舊城.
就說是基督徒.
這一班新開始.
對帝國有威脅的.
不信我們的宙斯的.
要敬拜獨一的神.
要敬拜釘十字架的耶穌的.
每個星期都吃聖餐的.
相當明顯的地步.
連帝王都知道這個組織的存在.
信耶穌的教會.
這個組織的存在在羅馬.
所以他們就很簡單.
嫁禍給基督徒.
所以保羅的後期.
就是去賣了.
他就開始要面對逼迫.
也是由尼祿王所引發的.
這些公開性的逼迫.
所以你可以想像到.
保羅寫這封信的時候.
還不是大規模明顯的逼迫.
但是保羅一傳道.

$^{921}$我們上一次都說過.
保羅一出來傳道.
他自己就已經面對那個反對.
在大馬士革已經有當地的人反對他.
基督徒還沒有了解.
保羅的信仰是否真的.
是不是做一個臥底.
要來密謀.
來將我們一網打盡.
可能慢慢觀望之後.
發現不是的.
保羅真的用心.
用他的性命去傳講.
基督耶穌是我們猶太人.
外邦人的救主.
這樣的一個傳道就引起.
很自然的聖經也是這樣說.
就是那些猶太人對他逼害.
所以他也沒有辦法解釋.
不如避一避.
所以信徒就將他.
可能找一個鑼.
運走他.
將他扔到城牆下面.
因為前門的城門已經關閉了.
所以由城牆的房屋.
將保羅調下去.
他就可以連夜這樣走.
上去耶路撒冷.
也就是說.
都要逃避這樣的危機.
所以保羅對於逼迫.
一點都不陌生.
輾轉上到耶路撒冷.
很感恩巴拿巴信任他.
將他引薦去見使徒.
得到使徒的確認.
又相交.
就可以在耶路撒冷出入去傳道.
所以這個就是保羅在耶路撒冷.

$^{961}$上一次我們說了.
那個情況.
好景不常.
很快耶路撒冷的猶太的信徒.
或者那些猶太教的人.
又開始以為保羅又在傳講.
耶穌是神.
是和這位耶和華神他們敬拜.
是有徹底的抵觸.
又是要來迫害保羅.
所以弟兄們.
所以保羅開始有很好的教會生活.
這位好像入室的使徒.
能夠有戴瑪士教的教會.
耶路撒冷的教會的弟兄姊妹.
對他的慢慢的接納.
然後帶他離開.
叫他回去他自己的家鄉.
基尼加的大塑.
那裡不是那麼重猶太教的氛圍.
也沒有太多羅馬帝國那種嚴謹的政治氛圍.
所以猶太教和帝國的政治.
在大塑就不是窮坑僻壤.
變成沒有那麼大的風頭火勢.
所以保羅就回到大塑.
在那裡好像一直靜默了.
在那裡過一個很平靜的生活.
他的平靜讓他有更多的空間.
未大規模參與事奉之前.
他就有很多的默想.
很多的整合舊約和新約.
所以今天我們很擔心人不事奉.
很擔心弟兄姊妹信得主不夠出名.
一信得好一些就將他擺在最前面.
什麼都給他做.
如果稍微更出名一些.
又好勝一些.
那就更加開心了.
拍手掌了.
這個信主半年就已經可以做主席了.

$^{1001}$已經可以什麼什麼了.
小心哦.
強如使徒保羅.
他都有他的優如期.
他都有他的學習.
他自己回到去.
當然神給他有很大的恩典和祝福.
和聖靈在他心裡面的工作.
他將他自己舊約的經文.
做一個很徹底的反省和評估.
做一個盤點.
然後他就能夠有能力寫新約的十三卷書.
能夠整合舊約.
能夠在講到他自己一些經文的原創性的時候.
至少羅馬書是這樣.
他能夠引用舊約.
是相當的游刃有餘.
這個是經過思考的很詳細.
很多的累積的整合.
所以是很有深度的.
所以今天基督教其實是太沒有深度了.
很多人信了主.
只有開心 空殼.
然後裡面什麼都沒有.
經文都不讀.
根本連讀一次聖經都很難為他.
有時候難聽說一句.
可能是教會很重要的領袖.
你就一定死.
教會最高權的人都是空的.
沒有聖經的.
問他一次聖經.
他說揭都揭不懂.
當然現在不用揭.
現在用APP就可以搜尋.
但問題是腦子裡面什麼都沒有裝聖經.
完全沒有這些概念.
那就很嚴重了.
所以你看保羅他自己這種神給他的.
真的很客觀地體驗過Big Bang.

$^{1041}$經驗過由一信耶穌.
為什麼要受Big Bang.
我以前不用這樣.
我以前去到哪裡就受猶太人的歡迎.
尊重.
我是在社會上相當有社會地位的拉底.
為什麼現在我要為耶穌受這麼多苦.
只是受苦都需要有這樣的學習.
保羅就是這樣經歷過一次一次.
然後去到大宋之後.
根本沒有人再理他.
在大宋怎麼辦呢.
我就在大宋傳道.
明明神說不是的.
我要去到外邦.
我要向君王傳道.
為什麼沒有事情發生.
我繼續在大宋.
時間還沒到.
神的工作.
神的訓練還沒完成.
還需要沉潛.
所以我們今天太少人說的.
生命的沉澱.
一個領袖要經過很多很多的艱難.
這樣成長.
經過了這麼多艱難.
他的內心仍然要事奉神.
這樣就沒有人可以難阻了.
我又說一下.
某程度上是我的Mentor.
是我的恩師.
我2023年12月7日.
2003年.
因為和我自己的部會.
有些意見上不相同.
分開了.
我要重新.
因為我以前沒有讀神學.
也不相信神學的裝備.

$^{1081}$那個價值觀當時是這樣.
對神學對什麼都不知道.
外面教會覺得好像行水萬獸.
我當時真的這樣看事物.
看世界觀,教會觀.
根本看不起什麼神學院.
覺得自己已經可以去做.
輾轉.
我就在孫德堂.
就是我們的分堂.
開始聚會.
因為我的兒子阿平就是讀徐澤林小學.
可以想像到那個關係.
是神配合我們.
和孫德堂的傳道人又認識.
又上來我們家探訪.
就是伍敏儀姑娘.
伍敏儀姑娘.
當時就是張牧師做堂主任.
不認識的 完全不認識的.
但是當時張牧師知道.
我離開了自己的部會.
準備進神學院.
好啊 進神學院.
再重新裝備.
但是很可惜.
自己的部會不肯幫我寫推薦信.
沒有推薦信.
就沒有神學院可以收我.
我是2004年.
我是去播道神學院.
來申請入學.
神學院的院長就約見我.
親自約見我.
他也算是認識的牧者.
就是說.
如果我的舞會的傳道人.
不幫你寫信.
神學院都不會收你.
沒辦法.

$^{1121}$這件事我就.
因為張牧師知道這件事.
就問我怎樣.
我回去告訴他.
神學院這樣說.
教會就說.
我們不會幫你寫的.
就變成了半天調.
我說不要緊.
如果神要我不要做傳道就不做.
張牧師就給了一個.
一個這樣的恩典的.
一條路給我走.
他說這樣好像很不公平.
因為你在教會這麼多年.
又是曾經和教會一起去事奉.
現在是意見不同.
為什麼你會阻礙了我的進修呢?.
因為你不寫.
就等於完全斷了.
張牧師覺得這樣不是很公平.
不如這樣.
你在孫德堂聚會一年.
如果一年之後.
你都仍然有這個感動.
我就會幫你寫個人推薦信.
不是錦繡堂幫我寫.
是張牧師個人.
其實與他無關.
他不認識我的.
就是因為這樣一個.
給我一點光明的路.
我就在孫德堂一路聚會.
等了一年.
我是2005年.
差不多到他們入學的.
開始去呼召營.
開始填申請表.
張牧師都很好.
就找了郭鴻彪牧師.

$^{1161}$我不認識的.
完全不知道是什麼人.
約他來和我吃一頓飯.
因為都是孫德堂.
都是錦繡堂.
約我吃一頓飯.
了解我是什麼人.
了解了之後.
還要我等了一年.
2004年至2005年.
這兩位牧者都幫我寫個人推薦信.
後來再加其他一些我認識的弟兄姊妹等等.
我總共有15封個人推薦信.
但是沒有教會的推薦信.
所以建築神學院特別處理我這個案件.
又了解為什麼等等.
後來還要發現.
他們有其中一位老師.
真的憑信心接納我.
收我做學生.
沒多久開學之後.
在學校居然有一個個人的基金.
給我讀書交學費的基金.
就是弟兄姊妹.
放錢進去.
四年我不用交一毫子學費.
神學院的四年的學費.
看到神一直給我這樣的等候裝備.
所以我時常都回想.
如果不是張美玲牧師一句話.
可能現在我拿著小巴.
或者可能做其他事.
都對的 沒問題的.
我當時已經很豁達.
我是很清楚 如果神不要我做.
因為正正是武堂不想我做.
他不想我做才不幫我寫.
所以不要緊的.
我覺得彼此尊重的.
等等.

$^{1201}$所以真是一句話.
另外有一條路走到今天.
所以2009年畢業.
我不是說過了嗎?.
想起四年都沒探過張牧師.
其實從來都沒去過錦繡堂.
我第一次去錦繡堂.
就拿成績表給他看.
這個是2010年畢業.
在一間教會事奉.
接著差不多去到五月.
就是2010年的五月至六月左右.
就拿成績表和獎狀去探望張牧師.
告訴他現在的事奉.
是一間小的堂會來服侍等等.
諸如此類.
張牧師說不如你來錦繡吧.
就這樣開始了.
我在2010年的11月.
就開始在錦繡服侍到今天.
現在又在分堂繼續服侍.
尋潛有很多眼淚.
有很多個人的遭遇.
每一個都是一個光環.
每一個都是令你更加對事奉著緊.
明白什麼叫事奉.
多難都不會覺得辛苦.
不會覺得是一種虐待.
因為十字架的救恩太大了.
大到一個地步就是.
你能夠服侍主是你的榮耀.
是你一份很大的祝福.
但其實是很不容易的.
所以這些經文.
是真的要我們去沉澱.
是要有經歷.
是要有痛苦.
你才會感受到.
小路普羅的內心的那種胸襟的豁達.
他自己就是經歷過大馬士革.

$^{1241}$經歷過耶路撒冷.
經歷過大掃城.
不知道一年還是兩年.
都跟他無關.
已經被人遺忘.
不過巴拿巴記得他.
就去找他.
就開始他的裝備好了.
可以備受大用.
就開始在安提摩.
終有一年的時間.
跟弟兄姊妹查經.
然後聖靈感動安提摩教會的同工.
聖靈就跟他們說.
要為我分派巴拿巴和普羅去做神要他們做的工.
他們知道之後去哪裡好呢?.
在安提摩教會祈禱完.
巴拿巴說既然我跟你是神夫照的.
不如這樣吧.
不要去大掃了.
去我的故鄉居比路.
就是塞浦路斯島.
所以第一次的旅程.
就是安提摩去到塞浦路斯.
整個島我不知道有多大.
我退休後很想去塞浦路斯.
先講一下我的心願.
我很想去看看塞浦路斯有多大.
普羅第一次大規模的.
正正式式用宣教士的使徒身份.
跟巴拿巴去塞浦路斯.
整個島走遍了.
一路由帕窟去到.
有些名字忘記了.
在那裡上岸.
然後回到安提摩.
就是第一次宣教士的報道.
所以當我們看這些經文.
肯定是普羅經歷過.
然後寫的經文.

$^{1281}$所以他寫這些經文.
是有血有肉的.
他真的受過迫迫.
是真的有人會嘲笑他.
拉比不做.
做一個這麼下賤的人.
去講耶穌.
說是我們猶太人的尼塞那.
是被人單單打打的.
是被人嘲笑迫迫的.
去到一個階段.
甚至乎要用石頭打他.
在第二次旅行報道.
在路斯德就是了.
打到半死.
不如起結吧.
怕大家都不熟.
十四章就已經見到有這個經文.
第十四章第一節開始.
這是第二次旅行報道.
二人就是巴拿巴和普羅.
二人在以戈連同進猶太人的會堂.
在那裡講的.
叫猶太人和希臘人信的很多.
但那些不信從的猶太人.
就聳動外邦人.
叫他們心裡怒恨弟兄.
二人在那裡住了多日.
依靠主放膽講道.
主席他們的手.
施行實質歧視.
證明他們的恩道.
城裡的眾人就分了黨.
有服從猶太人的.
有服從使徒的那時.
外邦人和猶太人.
並他們的官長.
一起湧上來.
要凌辱使徒.
用石頭打他們.

$^{1321}$使徒知道了.
就逃往呂高尼的路斯德.
特比兩個城.
和周圍地方去.
在那裡傳福音.
呂高尼是那個省.
路斯德和特比是那個省裡的城.
廣東省的深圳市.
或者是其他的地方.
廣州市.
為什麼呢?.
到了第二世紀.
呂高尼再不是這個省的地方.
所以這段聖經肯定是第一世紀寫的.
所以路加醫生寫得這麼仔細.
因為他就是那個時代的人.
所以有些文件是印證了.
聖經的真實.
第八節.
路斯德城裡坐著一個兩腳無力的人.
心內就是瘸腿.
從來都沒有走過.
他聽保羅講道之後.
保羅定睛看他.
見他有信心.
可得存於.
大聲說.
在這個對望過程.
聖靈的工作.
聖靈的感動在保羅身上.
也是聖靈的能力.
要他做一個神蹟.
希望引發周邊的群眾.
拜的不是拜路斯德的偶像.
因為之後他們就會拜那個雕師.
要他們敬拜真神.
所以就叫他立即站起來.
所以那個人就立即跳起來.
然後行走.
好像《使徒行傳》.

$^{1361}$彼得在第三章.
叫那個在聖殿門口.
吃的那個黑丐.
是生下來瘸腿的.
從來都沒有走過.
就是說印證著.
保羅真的好像彼得一樣的使徒.
是使徒的權柄.
所以神也給了保羅.
也有這樣的權柄.
斑馬斑就沒有.
很特別.
是保羅才有的.
然後第十一節.
眾人看見保羅所作的.
很正常.
用呂高靈的話.
好像換廣州話.
大聲說.
有神藉著人形.
降臨在我們中間.
所以有這種美麗.
是不合理的.
因為正正就是展現.
神要透過保羅去宣講.
耶穌基督才是真神.
但是原來當時的人是拜那個雕師.
那些的紐像.
所以第十二節.
於是就稱巴留巴為雕師.
稱保羅為希爾米.
希爾米大家都知道.
就是那個叫做馬車的神.
Hermes 愛馬仕.
現在很出名的.
很貴的品牌.
就是Hermes這個.
當時希羅文化就是拜這些假神.
雕師就是一個傳遞訊息的.
因為他說話領首.

$^{1401}$第十三節說話領首.
就是說因為保羅負責代言.
巴留巴就很安靜的.
就出聲.
所以有城外的雕師廟的祭司.
就牽著牛拿著花圈來到門前.
要跟眾人向仕途獻祭.
就是說當了他們是神人.
當了他們是一些有神跡騎士的.
那個廟裡面的顯靈.
所以一定要澄清.
一定要很清楚的撇脫.
我們是凡人.
不過我們是神的僕人.
不是你們雕師的那些領袖.
或者那些聖人等等.
所以第十四節.
巴留巴保羅義士和聽見之後.
就立刻撕開衣裳.
跳進眾人中間.
喊著說.
諸君為甚麼作這事呢.
我們也是人.
聖情和你們一樣.
我們傳福音給你們.
是叫你們離棄這聲虛妄.
歸向那創造天地海.
和其中難物的永生神.
藉這個機會.
有機會講見證.
講福音.
講他們的需要.
你們拜的不夠我們厲害.
當然不是一種挑戰.
而是說我們要敬拜的是.
真真正正創造天地宇宙萬物的真神.
和其中萬物的真神.
留意第十六節.
他在從前的世代.
任憑萬國.

$^{1441}$各行其道.
十七章的四行傳.
去到希臘的瓦點的那個地方.
也是這樣再講這篇道的.
就是和那些拜偶像的人.
好像對號入座那樣告訴他們.
你們拜的是錯了.
我們在拜的是一位.
我們可以去尋問他.
因為是他向我們顯現的.
他說他是任憑萬國各行其道.
然而為自己未嘗不顯出證據.
原來神不斷顯明他的威嚴.
他的憤怒.
他的慈愛.
他的信實.
他的存在.
所以我們很重要對神要有一個這樣的尋求.
有這樣的認識.
就如常思因為從天降雨賞賜豐年.
叫你們飲食飽足滿心喜樂.
原來保羅都沒有拒絕.
那個自然的啟示.
原來整個的地球.
整個的宇宙.
從前到今天都沒有改變到.
原來是神的管理.
那個護理.
一直到今天.
所以我們就看到.
造物主一直這樣去憐憫世人.
所以保羅雖然沒有什麼科學的知識.
但是他很清楚那個世界觀.
很清楚的創造論.
很清楚神是真正的永恆的主宰.
我們一定要為他.
向著他而活的.
這樣的基督徒.
所以保羅在這裡就要攔著那些人.
拜他們是神.

$^{1481}$所以第十八節.
二人說了這些話.
就緊緊攔著眾人不欠濟於他們.
但是有些猶太人從安提安和以哥廉那裡來.
就挑唆眾人.
就用石頭打保羅.
打保羅.
就是用石頭扔他.
曾經保羅就看著那些暴徒扔死才是犯.
保羅就看著那些人的衣服.
他未信耶穌之前的青年的時候.
很氣焰要去抓基督徒的時候.
現在他為主就被石頭打.
以為他死了.
將他拖到城外.
門徒正圍著他.
他就起身走進城去.
好像不怕死那樣.
拖他出去城外.
想著就這樣讓他自生自滅就算了.
你可以想像得到.
被石頭打完再拖他.
不是舒服的紅地毯.
是很崎嶇的石頭路.
怎麼會開心?.
怎麼會容易?.
拖到出去.
他醒了.
他不是馬上走.
是回到城裡去.
去見證他是神的僕人.
剛才的訊息是千真萬確的.
我不是那些招搖撞騙的.
這樣來澄清他自己的身份.
他不是走了去別的地方.
是回到城裡去.
所以我們需要了解.
在羅亞書寫的時候.
保羅已經有這樣的屍體.
有這樣的事奉.

$^{1521}$所以越是要被神用的弟兄姊妹.
越要學習神對你的隱藏.
沉潛,裝備.
還有那個心是真的很想服侍神.
不是為了自己的野心.
所以在開始服侍的時候.
一定會慢慢的.
神會將我們的內心很多的雜質都除乾淨.
這些是需要一些的考驗.
一些的困難.
一些的阻礙.
會讓我們更加的純全的心去服侍.
所以越是經驗過神的恩典.
神的考驗的人.
就越能夠在事奉的路上.
能夠越來越豁達.
那個心胸又夠闊.
還有事奉的心態是正確的.
這些需要很多的心靈的裝備.
所以被隱藏的基督徒有福了.
如果一初信就已經有很多的.
鎂光燈射著.
自己又很心花怒放的人有禍了.
因為很快你就會原形畢露.
你就會表達你那種驕傲,自大,無知.
和不堪一擊的那種困難.
在困難底下很洋雞的不堪一擊.
所以不要以為很開心.
剛剛信主就說甚麼都找你做了.
那不是一個恭喜的事.
所以我們要留心.
這些是生命的功課,生命的服侍.
所以保羅告訴我們.
我組織了一點.
這個就是信仰不是一個抽象.
是具體在生活裡面.
你真的遇到不同的人對你的態度.
你要學習怎樣跟人家交往.
有誤會的怎樣能夠化解.
有真正針對的.

$^{1561}$你怎樣能夠學習忍耐.
用生命,用時間去改變人家對你的看法.
有時百詞莫辯的就莫辯.
你不需要再辯了.
因為怎樣講都沒有用.
所以若是能行.
總要盡力和睦.
盡力了之後都不和睦.
這就不是我們的責任.
內心都要表達.
不要有一種仇恨.
或者有一種自己內心的錯誤價值.
接著就是告訴我們.
在保羅的睿智裡面.
教會肯定是處於一個多元的文化.
因為既然羅馬的教會.
有來自世界不同的民族.
包括當時羅馬帝國.
是這麼強大的軍事的勇力.
他們的文化,他們的法律,他們的建宗.
你看看那個鬥獸場到今天都健在.
你可以想像到羅馬帝國.
他們的版圖,他們的雄心.
他們要征服世界的.
承繼了希臘的文化等等.
可以繼續將這個帝國去到很遠的地方.
所以你可以想像到.
在意大利,在羅馬.
一定是來自不同的國籍的人的聚集.
所以教會就在那裡成為一個很重要的組織.
也是神擺在帝國的一個平台.
所以那些基督徒就有很重大的存在的意義.
就是要印證基督的福音.
所以當然透過宣講.
但是更多就是透過在那個互動裡面.
猶太人對你們的嫉妒,攻擊.
因為你僭越了猶太教他們的真理.
還有你又在帝國底下.
你又不拜雕師,又不拜宙斯的時候.
你又不拜希爾米.

$^{1601}$你又不拜猶太人.
當時迦南的那些假神.
很多的東西你都不要的時候.
所以這個是文化上,真理上的對立.
不等於對人的對立.
這些就是我們需要在教會這個平台學習.
所以很難的地方就是.
因為你很簡單去做一些引修士.
與世無爭或者不關我們事.
你死你事.
我們自己在修道院自己自給自足種下菜.
不問世事那就沒問題了.
你有什麼大的波動都不關我們事.
所以我們就要去學習.
我們要去學習教會的教義.
在教會裡面.
相信第九節到第十三節.
我們可以稍為歸納了.
要彼此的一種接納.
避免的就是我們個人主義太強.
個人的性格好勝失去主宰著其他人.
所以你看到這裡是彼此的推讓.
是很恭敬 是很有禮貌.
是很有同理心.
是很有一種對人的空間.
給予人一個空間去做.
大家有一個決定.
但這個決定不會帶來很大的壓迫感.
這些是大家要學的.
即是要有一種彼此都願意讓步的愛.
在愛裡面.
但十四到二十一節.
外面透過經文給我們感受到.
教會是初期教會 都不是很多年.
就已經要面對很大規模的迫害.
甚至是白熱化的.
真的要面對那種攻擊.
所以你可以想像.
這種大方向原則.
在經文裡面已經讓我們看到.

$^{1641}$基督徒的生活原則.
在這個群體裡面.
這個群體對外面未信的人.
我們應該有什麼倫理守則呢?.
要珍惜.
我們的根源都是從裡面的內心而來.
要愛得真誠.
這個我想我們要學習.
不能假的.
因為如果我們的心不是真而是虛假.
神會監察我們.
剛才說了彼此的尊重.
這個就是十節的經文.
還有要常常保持這份熱忱.
這份對主,對人的熱心.
有時會受傷到一個地步.
去到冷漠.
冷漠變成一個保護.
對自己心靈的一個保護.
但這樣是不健康.
不是神期望我們這樣的.
所以可能要透過不同的方法.
例如向人傾訴.
有人代禱.
有時真的需要去整理一下.
將受傷的向神說.
以致自己不要有一種創傷的後遺症.
反而是釋放自己的心靈.
能夠重新振作.
所以有時可能要退休.
或者真的要找信仰的人禱告.
分享.
保持這份熱心.
相信教會信主來的.
不熱心的是多的事.
不熱心絕對不是每個人自己本身想.
而是積累了很多的誤解,誤會.
言語上的衝突.
被人針對.
受到其他人的排擠.

$^{1681}$這樣你想熱也不行.
所以為何古羅要這麼鄭重再說.
因勤勿難度.
心裡要火熱.
是心靈裡面被神的愛感動.
被主的救贖來改變自己.
所以真正的原動力.
都是回到神那裡去看.
基督耶穌的犧牲.
能夠心裡重新燃點.
當然有些人可以透過詩歌.
或者透過去下山.
總之是適合自己的方法就行了.
令自己重新有熱誠.
事奉.
知道自己是一個活祭.
是神喜愛.
是神悅立的.
這樣就能夠回到起初的熱誠.
這個心裡火熱去服侍主.
記住.
你心裡火熱去服侍主.
是你和神之間的默契契約.
但不代表神不容許其他人.
去將你這個火熱的心淋熄.
不斷淋也不熄.
唯有是基督的愛.
不斷在我們內心感受到一份祥和高深.
我們才能夠有力量.
不是逼自己.
所以不斷需要有安靜.
有默想.
我們自己找到一個.
我們自己可以復興的方法.
每個人都有不同的.
有些人可能是吃東西.
和別人聊天.
他就可以回到那份熱心.
重點仍然是熱心.
是一個淋熄的熱心.

$^{1721}$就是考驗到我們繼續在神裡.
為神而活的決心.
所以你可以這樣說.
第九到第十三節.
就是某程度這幾句.
就已經反映了我們那個心靈.
自己是真誠對神.
和期望對人是真誠.
對人也有一種尊重.
相信每個人都有他自己的空間.
有他自己的崗位.
恩賜 性格.
就是說原來我們對人的認識.
都要有長進.
不要一本通書看到老.
一個小朋友和一個老人家都已經不同了.
應該這樣說.
一個很有學問的大學生.
和一個沒怎麼讀書的年輕人.
大家一樣的年紀.
他們的價值觀都已經不同了.
我們都要這樣來看的時候.
對不同的人.
我們都需要有這種尊重.
仍然是對事.
要有這份熱誠.
十四節開始就處理那個冰霸.
要祝福冰霸的人.
十四節差不多是很肯定地說.
那種難見的冰霸.
甚至你可以馬上心裡想起.
哪個人 什麼名字.
冰霸的你.
但是我們能夠做的.
就是只能夠祝福.
說了兩次.
只要祝福.
不可以咒詛.
你可以想像到.
咒詛是最極端的.

$^{1761}$我看到他死.
我已經對我是這麼慘.
我是要奉神的名來祝他.
保羅說不單止不可以.
算了 我就不再記他了.
不行 還要為他祝福.
所以是另一種極端.
當沒有什麼不可以.
還要為他去祝福.
就是要為他祈禱.
為他去求神.
讓他能夠改變.
他可以蒙恩.
他可以將來被神用.
他的過去成為他.
服侍的一個很大的動力.
保羅自己就是這樣的為人.
他自己就是這樣的過來人.
曾經看著斯提反殉道.
他是很開心.
他是幸災樂禍.
後來他自己去大馬色.
因為長大了.
開始有這個權.
就將他自己的熱心.
要將基督徒有這麼多.
抓上來去受審.
他就好像將來由他去行刑.
這樣石頭的將這些基督徒處死.
就是說他自己都是過來人.
他現在發現.
哪個人哪個教會曾經為保羅祝福呢?.
為他的信主來禱告呢?.
不知道 聖經沒有說.
只知道有一個阿拿尼亞.
聖靈跟他說.
耶穌跟他說.
叫他去找這個素羅.
他也跟耶穌說.
這個人是迫不得教會的.

$^{1801}$你不知道的.
是啊 我知道.
不過就要你去改變他.
你按手在他身上.
所以阿拿尼亞好像沒有為他祈禱.
沒有為他祝福.
因為根本知道這個人是無法改變的.
所以原來保羅自己也是這麼戲劇性.
他現在是教導羅馬教會.
就要求他們在這個倫理上.
首先要有一個心理準備.
基督徒一定會受逼迫.
但是無論那個肉身的.
言語的.
心靈的.
侮辱性的.
單單打打的.
或者用一些手段的.
都只能夠祝福.
所以真的不簡單.
基督徒只能夠祝福.
所以如果我們不是相信有神.
如果我們不是不斷在神裡面.
我們是沒有可能做到的.
因為真的一點都不容易.
所以我們要對於.
如果有些人他們懷著一份憤怒.
仇恨.
我們不要太快去譴責.
因為真的不容易的.
因為真的需要有很大的愛.
才能夠做到.
所以我們要學習.
要為人祈禱.
所以如果是直接針對我們.
當然這個感恩.
因為針對就是說我們有機會去反省.
我們做錯了什麼.
因為如果每個人都喜歡你.
你就不能反省了.

$^{1841}$人是不會為了自己好去反省的.
不好的話.
我們就會有一個反思.
反思就是會帶給我們成長.
帶給我們更新.
所以逼迫你的人.
我們應該要提高一線.
其實他是我們的恩人.
他會逼我們去反省.
反省就是我們會有在神面前悔改的空間.
所以逼迫真的會是一個很好的.
對我們的成長最大的一個祝福.
所以我們經歷了逼迫.
我們就要學習.
所以這些差不多是很重要的真理.
是基督徒.
是我們願意做一個活祭的生命的基督徒.
我們不要做一個很敷衍的基督徒.
我們要做活祭的基督徒.
這些就是我們活祭基督徒的內功心法.
就好像我經常說的張無忌.
因為整個香港太熟了.
一代一代不斷去翻拍.
很多版本.
但是那個故事是一樣的.
所以代代的人都喜歡張無忌.
都知道他的劇情.
講到就容易做例子了.
他裡面就有乾坤大萊兒.
有這些九因真經的神功護體.
所以我很感恩就是我初信的時候.
我的教會是很嚴謹的.
逼迫我們讀聖經.
是使得我們很用心去讀.
這些是內功心法.
這些是神功護體.
因為現在很多基督徒是稀罕讀經.
不懂聖經覺得很合理.
為什麼要讀經?為什麼要跟我一起讀經?.
為什麼要我這麼辛苦讀經?.

$^{1881}$好像他還比你兇.
其實不是的.
我們需要有神功護體.
當然我們講的是例子.
我們需要有這些真理去做一個很重要的價值觀.
你想一下.
如果這個是我們基督徒的價值觀的時候.
我們還要擔心弟兄姊妹將來面對逼迫.
我們還要擔心弟兄姊妹自己遇到個人困境.
不需要了.
你在公司被人逼迫.
不要緊我們為這個人祈禱.
或者如果可以的.
我們因為有選擇我們轉工.
但是以後有些崗位是不能選擇的.
你是必須面對的.
這個時候我們就只能夠為他們祝福.
不要小看這些神功護體.
因為這些格言是塑造我們裡面的DNA.
塑造我們的價值觀.
塑造我們不敢得罪神.
我們真的要學習.
學習不到是不要緊的.
因為還沒有成長.
孩子還沒有走路不要緊.
他一定會走的.
但是現在走不了不要緊.
因為我們真的沒有可能.
我們是人.
所以我們是學習.
祝福不代表我們任你打.
我們也不是這樣.
我們是真的有血有肉的人.
所以這裡就說了兩次.
告訴我們.
就是說你真的會面對逼迫的.
你就要給逼迫你的人祝福.
只要為逼迫你的人祝福.
不可以咒詛.
就是說保羅自己曾經.

$^{1921}$以往都是這樣的過來人.
但是現在他就很感受到.
所以每一句話背後.
都是你明白不明白.
基督的愛那種無條件.
那種的無罪.
替罪人死.
現在我們根本就是可能.
我們自己都有錯.
我們都可能有虧欠.
所以沒有一個基督徒是對的.
那些人完全是壞蛋.
都不是的.
有時候都是經過這樣.
我們自己又可以成長.
所以這裡特別提到兩次的祝福.
是一個給我們一個很重要的.
基督徒的價值觀.
我們的內心的想法.
撒旦教剛好相反.
將這些經文.
因為我都說過了.
在神學院有一個叫.
別異宗教的研究一科.
抽籤抽了我做撒旦教.
怎麼辦呢.
老師還說你要飛去美國.
去找這些撒旦教徒.
我被人用了.
他又說不是撒旦教.
在網上搜尋.
原來這些人.
將聖經是完全剛好相反的.
撒旦教都有他們的十誡.
就是逼迫你.
你要加倍的逼迫他.
加倍的對他.
做一些完全無情的報復.
將一些東西完全扭曲.
將人性不斷的扭曲.

$^{1961}$就是撒旦教的教條.
美國是合法的宗教.
撒旦教.
所以很恐怖.
這個告訴我們.
保羅告訴我們.
這些經文不是一信主.
你馬上可以做到.
而保羅自己都要這樣經歷了很多的困難.
但是他就是這樣去教導信徒.
這個就是我們基督徒的生命的本質.
正正因為我們是被主的愛改變.
因信而清義.
我們在神的眼中.
現在是活生生的一個祭身.
是神喜悅聖潔的祭身.
呈現在世人面前.
我們就是這樣呈現給神的.
所以我們不是沒有感覺.
不是逼迫我們.
完全不會有任何的反應.
我們是有的.
不過神要我們在你仍然有空間可以信主之前.
我們應該加上我們的祝福.
希望你將來可以改變.
可能將來你就想起.
有一個人你曾經迷失的.
現在你就悔改了.
所以我們都是神用的一個器皿.
在這個世界.
我們是有一個祈禱的.
\newpage



\section{羅馬書}
\label{sec:2B_XWriXVXo}
\textbf{2025-12-12 周五晚查經班《羅馬書(30):分題:愛與順服》 講員:曾裔貴牧師}
\newline
\newline
連結: \href{https://youtube.com/watch?v=2B_XWriXVXo}{\texttt{https://youtube.com/watch?v=2B\_XWriXVXo}} ~~~~ 語音日期: 2025-12-13
\newline
\newline
\hyperref[sec:WdwoyEOQVUg]{\small{< < < PREV SERMON < < <}}
~
\hyperref[sec:index_chronic]{\small{[返順時目]}}
~
\hyperref[sec:index_scriptual]{\small{[返順卷目]}}
~
\hyperref[sec:ego6cJukrIE]{\small{> > > NEXT SERMON > > >}}
\newline
\newline
$^{1}$晚上大家一起我們禱告.
天父我們將晚上今晚我們的《查經班》.
我們開始進入第十三章的經文的時候.
讓我們一直追隨著保羅他當時的處境.
對教會處身在羅馬帝國的那個場景.
那個需要面對的政治上的考慮.
都讓我們能夠可以先明白經文.
然後我們能夠在21世紀當我們要應用的時候.
我們都能夠有一些很重要的大方向原則.
來對這些經文的了解.
求主大力幫助我們.
我們仰望禱告奉耶穌基督的名 阿們.
這樣我們就恢復今晚的《查經班》.
大家一直都在看第十二章的時候.
我先開大那個檔案.
要留意 因為我們今天可能要碰到第十三章的經文.
第十二章其實很重要.
整個段落由那個在救恩底下每一個信主的人.
都是一份憐憫的新生命.
我們應該成為一個活祭 奉獻給神.
過一個有意識地為神而活的生活 人生.
接著很快就告訴我們.
我們是一個蒙了恩典 救恩.
還有獲得神給我們的恩賜.
可以在教會這個場景 基督的身體裡面.
有我們的功用 我們要發現我們自己有什麼功用.
我們要在這些功用上要去發展.
接著要不斷去發揮.
令到我們能夠在我們由信主到我們回天家之前.
可以成為一個被主容的人.
我們每一個都是有那個作用的.
接著很快就提到那些恩賜之後.
反而第九節開始.
就是每一個基督徒應該都是一個新的生命.
是有神的愛 性命充滿的一個基督徒.
所以我們那個新的倫理關係.
是應該以愛 以善來作為我們那個生活.
一個很重要的基督徒的守則.
所以不能夠用一些操作 兇惡.
或者是不合真理的手段.

$^{41}$這樣去做我們跟人接觸的那個價值觀.
所以這裡首先是涵蓋在教會裡面.
多重要就是提到愛的那個生命.
是一個不虛假 對別人是以善.
來去相待的一個基督徒的生活.
這個是提了好幾次.
接著保羅也想到.
基督徒不只是在教會裡面.
在一個教堂的建築物裡面.
所以很快就觸及到教外是怎樣去跟別人相處.
很快就提到在教外就會有Big Bang.
所以你要面對第17至18節開始.
就已經在說有Big Bang.
包括14節都已經有提到.
Big Bang你們要給他們中風.
只要中風不可就坐.
你就可以想像到這樣的狀況.
絕大多數都是在教會以外你的生活的那個環境.
一定是你被人知道你是基督徒.
你在羅馬帝國底下.
你重新要去思想你自己的階層.
你的身份 你的不方便.
因為我等一下會說一些帝國.
和當時猶太人的衝突.
那個的張力等等.
所以基督徒都是一樣面對這個問題的.
所以本來就已經很意識到.
不可能只是在教會裡面你要去用善.
去對待其他的弟兄姊妹.
反而你應該將這一份的善.
怎樣能夠在日常的生活裡面.
接觸一些未信的人.
要來去相處.
重點就是提到祝福這個生命的課題.
就是我們有意識地知道那些困難.
針對,逼迫.
很有意識地知道對方不單不喜歡你.
還要可能在政策上針對你.
來攻擊的 你怎樣去面對呢?.
所以成為了基督徒絕對不是一件容易的事.

$^{81}$因為由你受贈開始.
你就要向其他的人.
如果你是一個猶太人.
你去宣告你要跟隨耶穌.
就已經是一個很大的一件事.
猶太的族群或者你的家庭.
是會來將你排擠.
你不再可以在猶太教這樣出遊.
所以這個你已經想像得到.
那個代價是相當大的.
這個是在當時保羅的時代.
信耶穌要付的代價.
但是如果我們跟隨的是這位復活.
為我們犧牲的救主.
這個代價是必須.
也是應該要這樣付出.
因為是一個很榮耀的跟隨.
所以你會看到就是.
從14節之後.
都是提到我們要有一個心理準備.
你要怎樣面對微信.
甚至乎迪士尼福音信仰的人.
我們應該怎樣跟他們相處.
跟他們交往呢?.
很多時候都提到我們內心要堅強.
我們裡面要有一個.
很大的一個屬靈的生命的那個泉源.
就是一個愛的能力.
來幫助那些微信的人.
透過我們那個表現.
不是以牙還牙的表現.
而是用一份的接納.
用一份的溫柔.
用一份的忍耐.
使用這些拒絕和抵擋我們的人.
當時的基督徒.
就是要這樣以善來去食惡.
以及不可以自己伸冤.
伸冤就是有冤.
就是有人陷害.

$^{121}$有人故意的攻擊.
詆毀.
甚至加害於你.
那你就有冤.
但是保羅說不可以自己去伸冤.
當然我們就不要極端到.
在法庭去作見證.
成為一個見證人都不行.
那就犯法了.
不可以.
因為我們要宣誓.
來作證的.
作證如果是刑事的.
法庭的調查.
那你宣誓完之後.
你就必須要講真話的.
就不可以說我不可以宣誓.
我不可以伸冤.
其實就不是這個範疇.
但是關乎到個人.
或者關乎到教會.
就開始要面對這些.
將來都很快了.
其實不是很久了.
保羅寫這封信.
大概是56至57年.
就是主後.
那你想一下.
很快去到68.
不如很快跟大家看看一個部分.
很快給你看一看.
因為我都做了反正.
你可以想像到.
66年就是10年後.
這個的.
尼祿之後.
就掠奪聖殿的財物.
侮辱猶太的信仰.
對不起.
尼祿皇未死之前.

$^{161}$因為他68年.
被元老院來到去.
逼迫他.
他就自殺了.
67年.
這個將軍維斯柏森.
威斯柏森.
就鎮壓猶太人.
就奪回了加利利.
來到他.
本來就是分封的皇.
加利利分封的皇.
就是這個.
糟了.
現在想不起來.
肥力.
68年.
主後就是尼祿皇自殺了.
接著69年.
就是這個維斯柏森.
就成為皇帝了.
繼承了尼祿皇.
接著就委派了.
很出名的.
委派了提多將軍.
正式因為有猶太人的叛亂.
猶太人起義.
起革命.
來推翻羅馬人.
所以提多將軍.
就攻打耶路撒冷.
將聖殿焚燒.
那些叛軍.
很少數.
很多大部分都死了.
有些逃亡.
有些始終都不肯投降的.
就跑到馬薩大.
提多將軍很厲害.
他弄那些土豪.

$^{201}$一路弄上去.
慢慢上去.
就不需要爬上懸崖了.
我們去過那個地方馬薩大.
那裡的人在裡面全部自殺.
不可以讓羅馬人殺.
所以70年.
就整個耶路撒冷完全被攻陷了.
所以猶太人全部分散到外國了.
你想一下.
是10年後而已.
保羅尼西.
大概這封信.
56年到57年寫.
10年後.
政治上已經有這麼大的動蕩了.
所以絕對不是保羅幾人憂天.
而是當我們看第13章的時候.
我們就知道原來保羅真的很有那個.
屬靈的敏銳和智慧.
所以我們看到這個就是.
12章已經給我們一個很大的生命的大原則.
成為了基督徒.
就是因為我們蒙受了神很大的愛.
救恩的無條件.
所以我們已經有聖靈將神的愛.
不斷敲冠在我們的心裡.
所以基督徒就是一群有愛心的人.
所以我們就要面對.
所以這個面對.
愛心就不是等於一種好像被人濫用.
愛心的背後都是有感情的.
都是有真理的.
都是有思考的.
都是有那個理智的.
就是說我們是很有理智地.
面對不公的迫害.
迫不受屈受冤.
我們用什麼態度呢?.
所以基督徒就不是任人打的.

$^{241}$這樣的傻的人.
而是有意識地我們知道.
我們要用愛去改變人生.
一定要用事去聖樂.
所以12章已經幫我們打了底.
基督徒這個群體.
一定是用愛作為我們那個.
如果我們用一個兵器或者一個武器.
愛就是我們那個大前提的手段.
達至問題解決的目標.
所以12章不斷在說.
我們都是要用善去聖樂.
這裡就是在說那些反對一個逆境的環境.
但是始終你要面對的.
就是當時那個社會的制度.
因為保羅不是憑空抽離了他的工作.
他現在生活的那個環境.
他那個時候就是彌祿作王的時候.
早期沒有直接將矛頭針對基督徒.
但是因為基督徒剛剛開始沒多久.
還沒有真正去到帝國的意大利.
成為了一個很主流的勢力.
來讓羅馬皇帝感到威脅.
還在耶路撒冷一些很小的地方.
一直這樣傳播而已.
的確保羅現在還沒有第三次旅行的宣教.
所以你可以想像到.
現在還是很中期.
保羅還是很大力去傳播的時候.
福音還是在土耳其一帶.
歐洲的南面.
還沒有去到很遠的意大利那邊的地方.
所以你可以想像到.
保羅現在正正就是他有意識地.
他要將福音帶到意大利.
然後再要希望透過意大利的首都.
可以將福音透過不同的.
可能他那一代的人.
或者下一代的年輕的傳道者.
來將福音傳到更遠的地方.

$^{281}$在他們力所能及的.
因為當時的航海.
當時的造船都不是那麼先進.
所以他能夠做的.
只能夠在地中海一帶.
離開地中海其實就相當危險了.
因為再去到遠一點的.
就要進入那些大洋.
他們未必有能力去駕馭那些海洋.
所以在地中海一帶.
都希望能夠將福音傳到最遠的西班牙.
所以一定是面對當時的社會.
那個制度,那個狀況的.
保羅知道他自己本身是猶太人.
他自己 我們讀福音書.
耶穌的門徒有一個是粉瑞黨的西門.
粉瑞黨後來越來越激進.
粉瑞黨本身的宗旨已經是激進的了.
就是一定要復辟自己的猶太教.
以及自己的以色列國.
就是要自己的國能夠再獨立.
羅馬人再不可以在我們的地圖.
來去又分封,又掠奪等等.
又駐軍等等.
所以你可以想像到.
粉瑞黨原來都是當時在耶穌的時代.
都已經有一些這樣的對政治和宗教.
拉上關係的這樣的混雜了.
所以你可以想像到.
一定是他們會對於自己的復國.
是有一個很大的期望的.
所以耶穌不能夠讓他們這樣的主張.
他們這樣的這麼好的心靈.
那個抱負來達到.
耶穌居然不是地上拯救我們的彌賽亞.
耶穌只是說自己是神的兒子.
但是最終在十字架上是毫無反抗的能力.
帶不了我們這個國家.
可以來到去有一個革命,一個復國.
所以粉瑞黨在耶穌之後.

$^{321}$他們就不再是很好的基督徒了.
他們就成為了革命的分子了.
所以在當時.
羅馬帝國已經對猶太人.
是已經有一種很嚴重的敵視.
其中一個很簡單的原因就是.
猶太人是拜獨一真神.
不可以在獨一真神以外有任何的敬拜.
基督徒都是在這樣的範疇底下.
來到去成為基督教的.
所以基督徒都是拜獨一的真神的.
不過後來就發展到.
聖經告訴我們.
特別是新約告訴我們.
我們是三一真神.
所以對猶太人來說.
是不能夠拜羅馬的皇帝.
當時君王的敬拜就沒有這麼嚴重的.
但是越演就越嚴重了.
尼祿之後.
他已經後期.
已經需要全世界.
當時他統治的國家的神民.
都要敬拜他.
就是羅馬的皇帝.
就是神.
所以對猶太人來說.
是一件你寧願殺了我都不能做的事.
對基督徒來說.
一樣不能夠敬拜這些地上的人.
作為神.
所以這個已經是足夠是一個理由.
是叛亂的了.
所以不斷開始那個鎮壓就越來越厲害了.
當然有另外一件事就是.
尼祿一個很大起高的人.
似乎又不是很有能力的一個君王.
他就見到舊的羅馬城很殘舊.
他想建一個新的羅馬城.
就需要很多的稅收.

$^{361}$很多的錢.
所以元老院就禁止他這樣做.
他一不開心.
就不知道怎麼辦.
就找一些人來嫁禍.
燒了那個舊的羅馬城.
就好像來到米爾城區那樣.
我先下手為強.
我燒了那個舊城.
就嫁禍給當時的基督徒.
所以元老院就沒辦法了.
他又可以建一個新的城.
就變成那個時候.
就用基督徒來做一個這樣的埋單的對象.
他就理直氣壯.
可以建一個新的羅馬城.
這個是舊約的歷史告訴我們的.
但是其實從耶穌的年代開始.
已經告訴我們知道.
羅馬人他們成為了整個當時的歐洲.
統一了的一個帝國.
所以他們要達至他們的尊榮的極權.
就慢慢推崇那個君王就是那個神.
要拜這個君王就是他們的神.
所以這些君王的敬拜就等同於拜偶像.
所以猶太人都是很清楚.
但問題是這個時候.
保羅的年代.
還不是一個真的全面迫害基督徒的時候.
仍然讓基督徒可以在他們自己的生活裡面.
有他們自己的宗教信仰.
因為就算猶太人都是.
他們可以有自己的聖殿.
他們可以有自己的宗教的審判.
所以彼拉多就跟那些猶太人的拉比.
那些法利塞人說.
他有什麼罪呢?.
我審過沒有罪的.
現在我釋放他了.
你們關乎你們宗教的我就不理了.

$^{401}$原來猶太人是有一個自己內部的宗教的自治權的.
但是關乎到刑事.
關乎到同國家的.
他們就沒有了.
必須要去巡撫彼拉多那裡去審判的.
所以那些猶太的長老就說.
我們沒有殺人的權.
所以就逼彼拉多要結帳.
要將耶穌治罪.
所以你可以想像到.
羅馬人很清楚.
他們會給不同的民族.
有你們自己的宗教上的自由.
那個自治權.
但是關乎到納粹.
關乎到一些羅馬自己的法律制度.
就必須一切的遵守.
所以就會出現了這裡.
你可以想像到保羅.
其實他是很有sense的.
知道基督教還是一個.
既是由猶太教那裡分出來.
和猶太教有一定的關聯.
人家未信的人是不知道你怎樣分的.
你們都是猶太人的樣子.
可能你們就不再慢慢去回堂.
我那時候沒有回堂.
沒有去成年敬拜.
但是對外人來說.
是不懂得分你到底是猶太人.
猶太教還是基督教.
因為基督教還不是一個很全面的.
讓更多的人來認識的.
固然他們一直開始去傳.
所以來到這裡的時候.
你就一定要面對第十三章.
保羅來到這裡.
當他講了這個大的方向.
基督徒我們在一個社會裡面.
我們在外人的眼中.

$^{441}$我們是一群怎樣的人呢?.
保羅已經很清楚.
來告訴我們知道.
這個大的方向就是.
必須以善在我們的生活裡面.
去戰勝邪惡.
包括惡人.
加諸於我們的一些惡行.
一些陰謀.
一些冤屈.
我們都是不可以為惡所勝.
所以禱告就是一個很大的出路.
我們可以向神哀告.
有些事情我們不能夠用一些人用的手段.
來以牙還牙的時候.
禱告就是我們可以訴諸於神.
所以羅馬書有兩段都講了.
不可為惡所勝這兩段.
然後要等候神.
讓神去伸冤.
十四章裡面講的.
所以陸續當面對社會這些.
相交的生活.
在社會上來自處.
或者教會來在那個社會上.
很視乎當時的政策.
羅馬帝國的國策.
如果有些是和真理敬拜真神.
沒有抵觸的.
沒有直接抵觸的.
基督徒就不能夠說.
因為這個政府.
羅馬的政府.
到處去掠奪.
當時真的是這樣的.
到處去殺其他的俘虜.
就是一個很不好的政權.
我們不能夠納粹.
就不行了.
本來就很強調是需要去做的.

$^{481}$我們一路看下去.
你就會看到.
很特別.
因為十二章和十三章.
中間夾雜了一到第七節.
因為第八節.
立即再講愛這個題目.
再講愛這個的本質是怎樣.
所以你可以想像十二章.
當講完那個恩賜.
彼此要用在神裡面給我們的綱位.
恩賜去配搭.
去服侍之後.
就講我們那個.
我們的生命的原則.
我們每一個基督徒.
我們在社會上與人相處.
就是一定要以善.
作為生活的一個標準.
待人的接物.
接著突然間的第一到第七節.
本來表面上就是講納粹.
對政府的信服等等.
與愛有沒有關係呢?.
其實阿波羅更大的考慮就是.
讓基督徒的愛在社會上實踐.
不是讓社會認為.
我們是獨樹一格.
標旗立意.
甚麼都不聽政府的話.
甚麼都來對抗.
我們有自己的社群.
我們有自己的一套一套的律法.
這樣就不合理了.
所以你會看到.
這個納粹這麼具體的.
是羅馬人強制每一個民族.
都必須向羅馬納粹.
因為當時很多的公勞網.
很多的法律制度.

$^{521}$當然還有很多軍王的奢侈.
需要很多的錢.
很多的稅來讓公共的社會.
可以運作.
當然我們今天就很明白.
又是到了交稅的季節.
1月就要交稅了.
我們就很明白.
公共的稅收是要來.
一個政府的運作來管治.
這是相當中性和相當需要的.
所以阿波羅沒有在這裡.
與猶太人有一個相同的想法.
猶太人就認為我們是神的子民.
我們沒有得罪你.
我們不納粹.
而且是不公平的.
我們是亡國的.
所以我們是不會納稅的.
所以就製造了很多的緊張.
很多羅馬人.
看猶太人是一個很大的威脅.
他們又要行國禮等等.
所以對羅馬人來說.
是一個很大的一個反動的民族.
所以本來在這裡的時候.
其實就不是突然間說1到第7節.
只是突然間說政府.
我們要上復.
反而其實是一直脈絡都是在說.
基督徒的生命的特質.
就是函善.
我們用愛來活在教內,教外.
兩個不同的環境.
所以很快就要見到第8節.
又馬上回到說愛的.
我讀一讀給大家聽聽.
13節開始.
在上有權柄的.
人人當順服他.

$^{561}$因為沒有權柄.
不是出於神的.
凡掌權的.
都是神所命的.
所以抗拒掌權的.
就是抗拒神的命.
將這樣的抗拒.
提升到抗拒神.
猶太人就是.
我怕得罪神.
所以我就不向羅馬人.
去納粹和去順服他們.
但保羅剛剛反過來.
不同猶太人這樣看這件事.
他認為一個.
健康和平的管治.
沒有戰亂.
沒有仇殺的社會.
那個秩序是神所命定的.
生活的一個安居樂業的環境.
不是因為那個政權不好.
神就好像要猜險基督徒或者猶太人.
將他們毀滅.
興起一個完全好的政權.
就應該不是這樣.
所以你會想像到.
這個就是保羅眼中告訴我們.
原來地上的政權.
一切的政權都是神所授命的.
就是說在一般的狀況下.
神仍然用人去管理這個地球.
其實羅馬人的法規是最好的.
羅馬人的法律是最.
當時是最完善的.
他們的法律的基礎是相當蓬勃的.
所以現在的現今的法律.
都是很多源於古羅馬的法律的基礎.
當然不斷會有不同的衍生.
因為時代不斷地進步.
但是法律的精神.

$^{601}$法律的基本的邏輯原則.
邏輯都是來自羅馬人 意大利人的.
所以這就告訴我們.
保羅將他現在要基督徒.
他的教會在羅馬的教會.
或者日後歐洲收到這封羅馬書的教會.
都要認定.
原來我們不要有一種錯誤的想法.
他這麼邪惡 我們就要推翻他.
就變成了基本神透過地上政權的護理.
保護 照管大地.
都失去機制 失衡的時候.
就只會帶來更大的混亂.
所以保羅不希望在他有生之年.
或者基督教成為一個不斷戰爭的宗教.
他反而很想什麼才能夠改變.
記得不是保羅看不到他助紂為虐.
不知道羅馬政權的問題 那些的殘忍.
因為在他的子女行間.
哥林多後書第一章.
保羅說我們這幫基督徒讀給大家聽.
你就會知道保羅其實不是不知道.
羅馬帝國那種殘暴的.
一章 所以當我們對保羅當時的背景.
有一個更多的認識.
那是不是幫助我們更好的去有一個反省.
哥林多後書第一章後面.
第二章十四節.
哥林多後書的第二章十四節.
並藉著我們在各處顯揚.
那因認識基督而有的香氣.
在各處去揚基督的香氣.
透過保羅他們去傳道的時候.
讓人認識基督.
而將那個生命改變而有的香氣.
然後留意第十五節.
因為我們在神面前.
無論在得救的人身上或滅亡的人身上.
都有基督輕香之氣.
在這等人 這等就是滅亡的人.

$^{641}$就作了死的香氣 就是拒絕福音.
拒絕保羅所傳的福音.
就有死的香氣叫他死.
在那等 那等就是上人所說的得救的人.
就作了活的香氣叫他活.
他說 這是誰能當得起呢?.
我們不像那些人 其實保羅想說什麼呢?.
保羅想告訴我們一個勝利的海船場景.
羅馬人打勝仗之後.
他們就會將他們所努力的金銀財寶.
以及那些俘虜用那些魚鉤.
鉤著一路回來一路示眾.
兩旁就有很多那些美少女.
意大利那些人在那裡殺那些香花.
兩旁就是海船回來.
那些奧運冠軍回來.
總有一些人夾道歡呼.
我們就不會這樣灑那些花 灑那些花瓣.
當時就是將那些花 那些香花.
灑在那些軍人那裡.
就是說他們回來的時候.
沿途就充滿了香氣 勝利的香氣.
保羅就用這個做一個比喻.
他說 我們讓那些得到的人.
有基督興香之氣.
就叫他們能夠感受到他們的生命的改變.
他們有一種脫胎換骨.
但是有些人拒絕福音的.
那就成為了死的香氣了.
就是說 什麼是死的香氣呢?.
就是那些俘虜回來.
他們也經過那些少女一樣灑那些香花.
但是是死的香氣 為什麼呢?.
因為很快他們回來.
完了這個列隊的遊行之後.
他們就會被處決了 所以是死的香氣.
所以保羅在這裡已經告訴我們.
他不是不知道羅馬人的殘暴軍人的.
他不是不知道當時他處理著.
他正面對的是尼祿王的時候.

$^{681}$所以如果我們在那個時代.
我們就會很明白保羅他寫的這些經文.
所以我們就不要用我們現在的那個心情.
就讀進去聖經了.
反而我們就要透過保羅當時的處境.
當時他想基督教怎樣能夠在這樣的.
好像甲房底下 慢慢的更加的茁壯成長呢?.
保羅很意識到基督教其實還是一種.
雖然一定是有神那個很大的超自然的力量.
但是更多是實實在在.
你怎樣去用神學去牧養弟兄姊妹.
怎樣用一些很重要的價值觀.
加諸在弟兄姊妹的心靈裡面.
有什麼價值我們要持守的.
我們追求的是地上的一個新的王權.
耶穌是主.
這裡已經是一個反動了.
耶穌是主 那麼蓋撒呢?.
蓋撒是誰?.
已經是一個很嚴重的指控了.
所以羅馬書不單止是這樣.
你就可以看到在保羅的書信.
根本就是要處理這些君王敬拜.
我們看看《貼撒羅尼加前書》.
在歌羅西書之後.
如果我沒有記錯 應該是第二章還是第一章.
我先看看.
這裡不是很明顯.
我們可以去看《班仕行傳》的17章就肯定有了.
就明顯是一個指控了.
所以保羅其實是很清楚他自己現在做的事情.
是需要很溫柔地,低調地去讓教教.
慢慢在整個羅馬帝國的監視下去發展.
讓人看到的是一個和平的福音.
不是一個對抗的,革命的一個推翻政權的一個武力.
17章的《貼撒羅尼加前書》.
對不起 是《士徒行傳》17章.
這裡告訴我們.
我們就多讀一點.
保羅西拉經過暗非波尼阿波羅尼亞來到貼撒羅尼加.

$^{721}$就是今天的撒羅尼加.
我們都去過的 很漂亮的一個地中海的城市.
相當舒服的那個環境.
貼撒羅尼亞是亞歷山大大帝幫他的妹妹建的城市.
所以我們在那個廣場.
是有亞歷山大大帝騎馬的英姿的銅像.
所以也是挺出名的一個地方.
保羅就來到這裡.
那裡有猶太人的會堂.
保羅就照他素常的規矩進去.
一連三個方式進去.
本著聖經與他們辯論.
講解陳明基督必須受害.
從死裡復活.
又說我所傳與你們的這位耶穌就是基督.
他們中間有些人聽了保羅的教.
就服從保羅和西拉.
然後有許多乾淨的希利利人.
尊貴的婦女也不少.
來跟隨耶穌 跟隨保羅.
但是留意了.
但那不述的猶太人.
心裡疾度焦醉了些使者匪類.
就是那些土匪或者那些地貓.
搭伙成群 洗黃手.
聳到合成的人闖進耶穌的家.
要將保羅西拉帶到白城那裡.
抓不著他們.
就把耶穌和幾個弟兄拉到地方官那裡.
然後留意他們喊叫什麼.
「攪亂天下的也道者裡來」.
「耶穌收留他們」.
這些人留意了.
違背了該殺的命令.
該殺是什麼命令.
該殺才是主.
然後說他們現在違背該殺的命令.
說有另一個王.
就是另一位主.
來統治我們.

$^{761}$就是耶穌.
你想想這件事.
等同於耶穌和該殺對敵.
跟隨耶穌的人.
原來想推翻該殺.
所以匪徒立了這樣的一個控告.
就很嚴重了.
所以眾人和地方官聽見就驚慌.
那些官為何要害怕呢?.
害怕什麼呢?.
就是害怕耶穌是主這件事.
被羅馬皇帝知道.
就是說現在的地方官原來包庇這班人.
慫恿他們要作王.
另外就是帖撒羅尼迦這個城.
當時因為很多年前.
他們在羅馬剛開始的時候.
有兩個中一還是中二的西史.
我都有聽過.
安東尼和屋大維兩個打亂.
後來奧古士都.
就是這個屋大維就贏了.
所以就成為了羅馬第一個皇帝.
奧古士都.
接著帖撒羅尼迦城的人.
當時就押注在奧古士都身上.
所以當奧古士都做了皇之後.
帖撒羅尼迦城就有一個最為國的待遇.
可以給他們軍事上的免軍糧.
所以現在聽到耶穌做皇.
那些地方官就很害怕.
因為本來我們現在享用的是國家最優的福利.
現在居然有一個人保羅來傳什麼呢?.
耶穌才是主.
我們今天就隨便掛一個招牌.
在左邊也有.
在同路圖書館也有.
耶穌是主.
那些不信的人就說耶穌是豬.
來抵禮.

$^{801}$但是在羅馬的時代.
耶穌是主是要殺頭的.
是要死的.
因為只有該殺才是主.
所以地方官是害怕的.
不是害怕保羅.
而是害怕他們受牽連.
所以可以想像到.
對保羅來說.
他是很清楚.
當時的羅馬帝國是不會歡迎基督教講的獨一真神的.
講耶穌基督是主的.
所以這些都是我們敬拜真神必須要付的代價.
如果被人嘲笑.
如果我們只能夠用一個忍耐寬容的心.
但是甚至到了一個地步.
如果是要逼迫你再相信耶穌是主.
就要受到一些刑罰.
甚至監禁.
我們也只能夠接受.
我們沒辦法去推翻.
耶穌基督才是唯一人類生命的主.
要告訴我們.
當我們看回這些經文的時候.
當大家稍微看了兩段保羅的書信.
我們就已經知道.
保羅其實很清楚.
他現在寫的羅馬書.
那個影響性是相當大的.
是直接影響著那一代的人.
怎樣去想自己的存在.
教會的角色.
在地上的身份.
所以如果保羅用一些很煽動性.
很革命性的.
只會帶來宗教很快就被瓦解.
保羅不是要忍辱負重.
避重就輕.
而是他再更加在他的智慧底下.
他是一個猶太人.

$^{841}$他也是敬拜獨一真神的.
但是在他的智慧底下.
他很清楚.
羅馬書也很清楚.
唯有萬民恩順清延.
不是政治可以叫人恩順清延.
不是一個廉潔的政府.
就可以自動的恩順清延.
萬民必須恩順耶穌基督的寶血救贖.
而來到生命的改變的清延.
所以保羅現在很清楚.
他在地上的身份責任.
他是基督耶穌的一個奴隸.
一位使徒.
是神特派他去傳救人的福音.
所以他很明白.
他要讓教會集中生命力去傳福音.
不是集中在一些.
例如只是關心社會物質的需要.
今次的火災.
或者過去一直有不同的貧窮.
我們都要關心.
但是總體來說.
我們要關心人的整全的.
就是他的身體和心靈是救恩的層面.
所以我們一定不可以偏離.
福音才能夠改變人生.
所以保羅能夠在這樣的處境.
他不是要步步為營.
不是要避重就輕.
而是他能夠說到一件很寶貴的事.
我們先讀一下經文.
之前還沒讀完的.
第三節.
作官的原不是叫行事的懼怕.
乃是叫作惡的懼怕.
說的是一般性.
政府的設立.
就是能夠維持一個基本的秩序.
人與人之間.

$^{881}$有一個和平共處的空間.
這樣來運作.
一路來讓不同的制度.
不同的人可以得到在制度的安全下的保護等等.
所以保羅就說.
你願意不懼怕掌權的.
你只要行事就可得他的稱讚.
因為他是留意了神的用人.
是與你有益的.
你如果作惡卻當懼怕.
因為他不是空風的佩劍.
他是神的用人.
第二次說.
是申冤刑罰作惡.
在一個很正常的狀況下.
神對大地的護理.
就是設立一個完善的政府.
在一個正常的狀況下.
我們都有不同的架構.
有不同的管理.
例如我們去大自然.
去大坑行山.
我們為什麼可以不怕野豬或者野獸或者獅子.
一來就是已經沒有了.
老虎都沒有了.
但是野豬都是有規模地.
有意識地.
有那個漁農處去管理.
所以就不會容讓那個大自然.
或者那個社會.
例如最近的日本.
那些北極熊出沒.
就進去民居去殺那些人.
你可以想像到.
如果一個地方完全沒有一個管理.
人根本就不能生存.
不可以居住.
各樣東西都會亂了.
所以我們看到保羅很清楚.
羅馬帝國的制度相對很完善.

$^{921}$那個法律的條規是很完善的.
所以保羅就用這個很具體的實例.
我們始終信了主.
都要面對在世界的身份角色.
我們在主面前.
我們敬拜的是獨一的真神.
不可以有第二位主在我們的生命裡面.
但是和現在羅馬帝國在地上的管治.
保羅知道奴隸的制度是不對的.
保羅都知道羅馬的政權是很殘暴.
很邪惡.
很多裡面的問題.
但是保羅他不是一個社會改革家.
他是一個傳福音的使徒.
所以他現在在這裡.
你可以說神給他的一年.
給他的啟示.
他來寫羅馬書的時候.
他將他自己在神裡面所領受的.
來寫出來.
我們應該怎麼做呢?.
提了兩次.
這個政府,這個政權是合法的.
他是神的用人.
就是說神透過地上的有規模.
有秩序的制度.
讓不同時代的人.
在其中的生活得到保障.
你就可以很自由地來順處.
得到因為政府的制度包括.
新疆死狀,教育,不斷的醫學提倡.
各方面軍事的保護.
在這個環境裡面.
就不會有外敵的入侵.
所以各方面是帶來一個護理.
一個保護.
在裡面的居民就可以自由地生活.
自由地生活就可以自由地.
羅馬的道路.
當時代已經沒有再需要和其他的城邦打仗.

$^{961}$已經完全統一.
所以他就可以發展道路網.
促進貿易.
當然就是因為軍事的原因.
所以就帶來可以有更多的人.
可以去不同的地方去傳播那個宗教.
傳播當時的哲學.
一些的哲理等等.
但是保羅就利用這些優勢.
來將福音傳到更遠的地方.
所以保羅看政府.
不是他要隱瞞那些壞事.
而是他很知道教會.
不應該再是一個地上的壓力團體.
一個革命的團體.
教會是地上一個在神關懷的社區.
做福音的傳播的一個社區教會.
一個福音的燈台.
所以我們就不要將這些經文無限放大之後.
就會讀錯了保羅的意思.
因為保羅在這裡告訴我們.
並不是他不知道羅馬帝國的腐敗邪惡.
但是他的責任是要讓神去審判.
他的責任是要將罪人帶進舊日裡.
得到那個基督的香氣.
所以保羅就很清楚自己所到的地方.
當他真是講活這位耶穌.
就有基督的香氣留下.
譬如我們現在每個星期四晚.
我們有九間教會去上陳淑娟牧師的課.
個個都聽到真的完全沒有人睏.
因為講很多個案見證.
怎樣叫那些基層的家庭很窮苦的.
生命得到改變.
那些男士失業或者打太太的.
慢慢又有生命的改變.
基督的香氣保羅就在那些信的人身上.
就會有基督的香氣叫他來活.
所以當我們真是深信福音可以改變人心的時候.
真是有一個很大的生命的扭轉.

$^{1001}$譬如我之前都講過.
我們有一個姐妹在醫院認識隔離床一個.
癌症末期的姓林的叫嘉慧已經回到天家了.
早幾天她的姐姐打來找我.
因為我開車聽不到電話.
後來我再打給她.
她說什麼事呢?已經過了兩個禮拜的喪禮.
原來她是多謝教會多謝我們.
他們感覺到很榮耀.
妹妹很安詳地回到天家.
她說她承諾遲些因為她要回大陸一段時間.
她說她回來就會來我們紅印堂聚會.
有時候就是這樣將福音傳開.
你就會看到一些未信的人.
感受到一種基督的香氣.
所以保羅其實很清楚政府是神的用人.
完全不同了.
當你看到他是你的敵人.
你是要無論如何要推翻他.
要跟他作對.
這樣就完全邏輯錯了.
他是神的用人.
當然我稍後可以言配說一點點怎樣形容.
不過可能未必說得到.
因為時間未必夠.
接著我們再到下面.
神的用人是與利有益.
你如果作惡卻當懼怕.
因為他不是空空的配件.
即是他有實權.
制度賦予他有不同的職級.
例如警察,軍人或者法官等等.
所以他們是神的用人.
我講了兩次很特別的.
在正常情況下.
他們都知道要扮演在制度下的角色.
例如他是伸冤的.
是凌罰作惡的.
所以你們必須信服.
講得好像很誇張很直接.

$^{1041}$有沒有搞錯.
羅馬帝國我都要向他信服.
必須信服.
不是吧.
政權這麼邪惡.
我怎樣向他信服.
原來下面才重要.
第五節.
不但因為凌罰.
也是因為良心.
因為良心.
我們的良心已經被基督的愛.
被基督的寶洗乾淨.
我們已經有一個新的生命.
我們呼叫阿爸父.
第八章的經文.
我們已經是活祭獻給神.
所以現在我們是代表著神的國度.
來活在地上.
所以是因為我們的良心的緣故.
我們是要維持這個政府的正常運作.
所以我們有一方面相信神會有審判.
有管治.
但是我們更加相信.
神會用我們每一個基督徒在地上.
做一個奉公守法的公民.
為了我們自己的良心.
我們就要來做.
所以納粹對於羅漢人來說.
是一個很必須要做的.
維持他們的國家機器運作.
但是猶太人就認為.
我是不會向你納稅的.
我們只會向神納.
因為他們要納丁稅.
他們要納性別稅.
人頭稅.
所以猶太人就不納.
他帶來的就是對抗.
帶來的就是鎮壓.

$^{1081}$帶來的就是反動.
但是保羅就很謙卑地讓我們看到.
因為這個良心.
不是因為一個的鎮壓.
我們就要怎樣呢?.
第六節.
你們納糧也為這個緣故.
為了良心.
就是說我們要承認.
神透過地上的人.
去管治這個大地.
所以我們相信神有祂自己的作為.
有祂自己的保守.
就像我們看到的.
有些人就說.
我沒有這個保守.
我沒有這個良心.
我沒有這個信仰.
但是我相信神.
就是在我們這個地方.
神的保守.
就是在我們這個地方.
神的保守.
就是在我們這個地方.
神的保守.
就是在我們這個地方.
神的保守.
就是在我們這個地方.
神的保守.
就是在我們這個地方.
神的保守.
就是在我們這個地方.
神的保守.
就是在我們這個地方.
神的保守.
就是在我們這個地方.
神的保守.
就是在我們這個地方.
神的保守.
本來就是要面對這個政權的.

$^{1121}$他原來是要幫政權去講好話.
來去歌功頌德.
他告訴我們.
是為了我們基督徒的良心.
我們要承認神.
是有那個大地的護理的.
如果他們是神的用人.
說了兩次羅馬政權是神的用人.
你們就要為這個緣故.
必須要這樣做.
因他門是.
再一次說.
因他門是神的差役.
本來這樣說.
羅馬的政權是神的差役.
上上特管者事.
所以就七字.
凡人所當得的就給他.
當得良的給他立良.
當得遂的給他上遂.
當懼怕的懼怕他.
當恭敬的恭敬他.
很清楚.
不是要我們做一個阿諛奉承.
隱惡揚善的一種這樣的東西.
而是讓我們知道我們存在這個社會.
我們就要讓這個社會的制度.
能夠健康的維持.
這是一個最基本的定義.
所以保羅在這裡說了.
好像突然間加插了.
為什麼要馴服政權呢.
為什麼要這麼順從的順民呢.
但其實是一個十二章的再一次的現實.
就是不單止在教會裡面.
我們要有一個愛的生活的原則.
對人真的要有真誠.
不要虛假.
還有面對一些人.
對我們的針對攻擊.

$^{1161}$我們都要以善去承諾.
這是一個很清楚的大原則.
就是我們基督徒的價值觀的改變.
但是保羅格不知道.
整個教會處在那個社會.
是要怎樣來有自己在神眼中的角色呢.
所以就必須是這一段經文的解釋.
接著回頭再一次說愛的原則.
凡事都不要虧欠人.
唯有彼此相愛.
要常常以為虧欠.
就是總是覺得自己愛得不夠.
因為愛人的就完全阻礙律法了.
接著放一個心圖.
「象那不可姦淫不可殺人不可偷渡」.
十誡.
「不可貪婪或辯極戒命」.
就是十誡.
他隨便都舉了.
不過要留意.
這些全部是十誡第六誡以後的.
因為頭四誡是對神的.
除了我之外不可以有別的神.
不可以有上天下地稱為神的.
我們去敬拜的.
不可以妄稱神的名.
第四誡就是不可以.
一定要守安息.
第五誡就是後面父母.
所以第六誡開始到第十誡.
就是講到人與人之間的道德倫理.
接著.
「一切都包在愛人如己」.
這句話之內的.
愛是不會加害於人的.
所以如果我們有這些動機.
有這些陰謀.
我們真的要很小心.
因為已經離開了神.
所以愛就完全阻礙律法了.

$^{1201}$所以你可以看到.
這段聖經包住了一到第七誡.
所以一到第七誡.
就是我在紅色那裡的字.
就是愛的實踐在社會裡面.
我們基督徒都是.
我們不是不知道羅馬的政權.
你們敬拜該神是錯的.
敬拜這個宙斯.
敬拜那些假神是錯的.
我們是不會這樣拜的.
但是我們仍然滿足我們良心.
我們為這個政府能夠健全地運作.
我們去納稅,納糧.
我們去做我們應有的本分.
當然這是最基本的.
不代表我們不可以有一些議政.
或者談論政制,政治等等.
不能夠說我們不可以進入.
那個政府的架構裡面.
去改善那些制度.
我們更加應該做.
所以是沒有抵觸的.
但是那個大前提.
因為良心的緣故.
我們要認定神是讓地上.
在祂們回來之前.
主耶穌第二次回來之前.
地上仍然有一個管治.
是在地上.
這樣來讓地上的人.
可以活得自然.
以及有生存的空間.
如果不是整個社會.
很簡單 如果我們去了一個地方.
比如像南丫堡這麼小的地方.
那裡是沒有警察的.
是沒有政府的.
是沒有法律規定的.
你做什麼都可以的.

$^{1241}$那就大件事了.
可能第二天你也想馬上走.
因為那裡是無政府狀態.
所以是沒有可能的.
所以你會看到.
原來政府這個制度是好事.
是對人類的一個好的傑作.
神容許人就已經有這些.
由可能消息穿落一個城邦.
慢慢成為一個國家.
就是羅馬帝國.
所以這裡就讓我們看到.
保羅絕對是知道自己在說什麼.
他更加希望教會.
一定要看到福音.
才能改變人的生命.
甚至乎嚴格來說.
改變帝國的命運.
所以羅馬帝國不斷的迫害基督徒.
他自己最終都覆亡了.
因為不斷的好大喜功.
不斷的去殘暴不仁.
不斷的得不到民心.
又繼續荒人無道.
就帶來那個國家慢慢的覆亡了.
所以你會看到西羅馬.
東羅馬最後都亡國.
一個這麼大的帝國.
都一樣會變成灰飛煙滅的了.
所以我們看到第8到第16節.
再一次告訴我們那個愛的狀況.
留意第11節.
再者你們曉得.
現今就是該趁早睡醒的時候.
白晝將近.
我們就當脫去暗媒的行為.
戴上光明的兵器.
行事為人要端正.
好仗行在白晝.
不可方便聚酒.

$^{1281}$似乎就是當時羅馬的政權.
那些人.
他們的生活的常態.
但是基督徒就要知道.
我們有一個更大的使命.
更大的那個榮幸.
就是主耶穌回來的日子是相當近了.
就快去到白晝了.
所以我們應該再活在黑夜.
這樣的意思.
這個概念在以弗所書.
歌羅西書都相繼會再出現的.
不可方便聚酒.
好色邪道.
憎激嫉妒.
總要欺負主耶穌基督.
不要為肉體安排.
去放縱私慾.
再一次給當時羅馬教會有一個勸勉.
就是一方面我們要很認真.
做我們的身份角色.
就是要保留那個敵意.
神在這個墮落的世界裡面.
仍然設立秩序.
就是政府.
避免混亂.
混亂帶來的就是很明顯的.
就是要感受停電兩晚.
沒了面子.
警察都不多.
就很亂的了.
有些人就處心出現了.
搶奪的了.
在美國就是這樣.
所以你可以想像到.
這個秩序是帶來人類的安全的.
所以信徒的信服.
不是因為政權本身.
而是因為承認神的那個管治的主權.
所以這個如果我們能夠這樣去理解.

$^{1321}$其實是一個很成熟的表現.
所以和第十二章這個不義惡報惡.
是完全互相呼應的.
信徒不是靠暴力去改變世界.
而是靠愛和善行.
這個是十二章和十三章給我們的結論.
所以保羅寫書的時候.
仍然是尼祿掌權的時候.
不過在他的早期.
就尚算穩定的.
但是當後期就逐漸越來越暴躍.
尼祿.
這個就是羅馬帝國對他的治理下的要求.
一定要尊重羅馬的法律和納粹.
另外就是君王的崇拜.
必須要做的.
任何被俘虜的城邦.
那些地方國家.
都要尊該什為主為王.
所以面對猶太人和基督徒.
因為他們是敬拜獨一真神.
所以就帶來在羅馬帝國的眼中.
是一個很敏感的身份.
所以時時都有張力.
《少林寺十八章》第二節.
猶太人當時是要被排遊.
因為當時格魯迪要驅逐猶太人.
因為他們就是一個反動.
一些主張得罪了帝國.
就要驅逐猶太人離開意大利.
當時就有一對夫婦.
伯基拉和阿居拉.
就是一對夫婦.
阿居拉夫婦.
他們被迫離開意大利.
輾轉就去到哥林多.
遇見保羅.
其實應該是保羅遇見他們.
他們已經在哥林多住下.
他們是即帳棚.

$^{1361}$保羅知道他們在羅馬被趕逐離開.
所以你可以想像到.
在保羅的年間.
猶太人和基督徒的身份.
是非常之敏感的.
因為他們會得到帝國對他們的監控.
監視.
甚至後來的敵視.
因為猶太人一直不肯接受.
君王的敬拜和納粹.
所以大家可以想像到.
基督徒如果承認耶穌是主.
就等同於敬拜奴隸真神.
和凱撒.
即是該撒是主.
形成了一個極大的張力.
所以在這個不妥協的信仰前提下.
避免政治衝突.
不要這麼快因為一些納粹.
因為一些事情.
來到有一個這樣的政治衝突.
是保羅要帶進給教會的智慧.
和一個很重要.
教會存在在地上的一個決策.
因為除了君王敬拜之外.
其他基督徒都需要去接納.
都需要去馴服.
何時就要出現這個問題呢.
很快保羅面對的就是斬後.
被捉拿.
要承認帝國.
就一定不能夠接受.
所以他們是很安然地坐牢.
不能夠馴服政權說.
我們要敬拜凱撒.
當你不馴服的時候.
你就要接受那個不馴服的代價.
那就是坐牢或者有身體的迫害.
或者甚至是死亡.
因為當時基督徒全部在那個.

$^{1401}$由主汗70年打後.
一直到主汗大概200多年.
都是十大迫害的時期.
很多的基督徒殉道.
是被殘害的.
所以就唯有這個.
如果耶穌是主受到威脅.
我們就沒辦法接受和馴服的了.
我們只能夠接受政權加給我們的刑罰.
但是除了這樣之外.
在一般的情況我們是應該要馴服的.
這個不是要特意和帝國或者政權對著幹.
或者好像歌功頌德就不是這個意思.
剛才已經說了這個.
你就可以想像到.
所以保羅的智慧令到教會可以一直繼續存在.
就是一些很重要的價值觀.
很重要的宿命的理念.
認定神仍然管理著大地.
神還沒回來.
政權的興替是神的事情.
我們基督徒教會的存在.
是要將人帶進基督的救恩的香氣裡面.
他的生命改變.
這個是很重要的.
我們就不是一個議席的教會.
參政的教會.
甚至推舉牧師曾牧師去參選.
做什麼就不行的.
我們都知道區議會的規定.
牧師可以做議員的.
不過牧師做議員沒有問題.
麻煩你辭職.
你不可以再做宣道會的牧師.
其他教會可以的.
大宗派是可以的.
可以同時做牧師.
同時做議員都可以的.
如果選上的話.
但是宣道會就不行.

$^{1441}$因為我們政教是要分開的.
這個要很清楚.
所以保羅這裡.
不是要對政教這個課題很詳細地講.
而是很想告訴我們.
在一般的情況下.
政權是神的牧人.
是神的用人.
講了三次.
這樣來告訴我們.
我們應該在一般情況下.
我們是要有我們的信服.
我們不是來做一個對抗.
甚至推翻革命.
所以這裡是很清楚.
猶太人亡國了.
但是基督徒仍然在.
你可以想像到.
只不過是十年二十年的時間.
就已經印證了羅馬書十三章的威力.
基督徒就是聽從士保羅所講.
知道這個世界觀.
原來神仍然掌管.
神母叫我們這樣做的時候.
我們就不應該像猶太人那樣.
做革命推翻羅馬.
最終的下場就是四五年都被毀.
猶太人都要被分散到全世界.
但是基督徒仍然在那個帝國底下.
當然開始就是大壁北.
是用生命去印證那個福音的真實.
所以很渴望我們大家都能夠要學習.
今天送給大家的經文的金句就是.
所以你們必須馴服.
不但是因為神的憤怒.
這個新的羅盤.
也是因著良心.
我們記得我們都是因著我們自己的一個生命的見證.
我們所得到的神給我們一個這麼寶貴的新生命.
我們看到政權的交替是有神的旨意的.

$^{1481}$所以我們就不要用自己的力量.
去做一些可能是真的違反法律的事情.
所以保羅在後面這裡的.
我們可能下個星期再說.
這個第八到第十四節.
今天可能短一點.
因為我的Power Point已經說完了.
我們就停在這裡.
我們一起有個祈禱.
天父我們將今天我們所看的.
由十二章一直到十三章的八節.
我們看到有一個很寶貴的教導.
就是基督徒整個教會存在在地上.
我們的價值觀是以行事.
是要將人帶進救恩裡面.
我們知道地上是有不同的罪.
包括罪人.
我們都不知道現在那個災難的調查.
將來九個月之後會怎麼樣.
雖然我們很多人都認為是人為的犯罪.
但是我們希望那個調查是能夠找出那些原因.
以及帶來那個根治.
但是我們知道我們活在這個地上.
我們真的有意識地知道.
這個黑夜異心.
保羅也形容了人的心.
很多的犯罪,得罪神.
求主幫助我們看到.
白晝真的將近.
包括我們自己也很有限的人生.
來盡力為福音,為主去傳揚.
將更多有需要的人帶進救恩的香氣裡面.
願主幫助.
也讓我們看到保羅的智慧.
保羅的決策.
大理教會安然地渡過了尼祿和的迫害.
以及往後的君王的迫害.
使得整個猶太國都瓦解.
聖殿都被拆毀等等.
天父讓我們看到我們有一個更加重要的任務.

$^{1521}$就是要將主的這個真理.
用我們的生命去傳揚.
一起去為神的福音效力.
我們仰望,感恩,禱告.
服耶穌基督的名.
阿們.
\newpage



\section{羅馬書}
\label{sec:ego6cJukrIE}
\textbf{2025-12-19 周五晚查經班《羅馬書(31):分題:愛與順服(二)》 講員:曾裔貴牧師}
\newline
\newline
連結: \href{https://youtube.com/watch?v=ego6cJukrIE}{\texttt{https://youtube.com/watch?v=ego6cJukrIE}} ~~~~ 語音日期: 2025-12-21
\newline
\newline
\hyperref[sec:2B_XWriXVXo]{\small{< < < PREV SERMON < < <}}
~
\hyperref[sec:index_chronic]{\small{[返順時目]}}
~
\hyperref[sec:index_scriptual]{\small{[返順卷目]}}
~
\hyperref[sec:aePwbTZJwec]{\small{> > > NEXT SERMON > > >}}
\newline
\newline
$^{1}$第一節我們慢慢又可以一起跟大家去查考羅馬書.
我們先一起禱告.
天父我們臨近聖誕.
在這個章林期我們都看到.
雖然世界都有很多的災難.
但讓我們屬主的兒女確信.
我們因為有主耶穌.
並不是我們能夠立刻扭轉世界.
而是在我們很多人的內心裡面重新燃點.
一個真正的盼望.
就是在主耶穌的救恩裡面.
我們才能夠真真正正的可以安穩.
可以來到祂有一個信心的依靠.
求你帶領今晚的聚會.
我們會看十三章.
我們很盼望我們能夠對於士路保羅.
他要跟我們去論證,分享的.
特別他對那個時代的局勢.
把脈是相當準確.
他都是有士路的那個眼光,權威.
使我們能夠透過士路的一些真知灼見.
我們能夠了解原來不同的時代.
在這麼黑暗的羅馬帝國.
同樣我們看到教會.
不單止沒有萎縮.
反而是更加的在困難底下去復活.
裡面一定是有很重要的肅靈的價值觀.
要我們去學習的.
求主帶領我們.
我們將今晚的聚會仰望在祂的手中.
我們這樣感恩禱告.
奉耶穌基督的名.
他們.
好 我們再來看第三十一講的.
今晚仍然這個分題都是.
因為第十三章我們都還沒有完全的完結.
上個星期都有一點探討到一些.
當時的社會背景.
羅馬帝國在這個尼祿王的.
這個是早期的執政的時間.

$^{41}$這樣的一個狀況.
所以我們又要探討.
第十三章在整個的羅馬書裡面.
那個的出現.
為什麼突然間要講到和政權的那個的關係.
要怎樣去周旋在猶太教對基督徒.
對基督教的敵對.
和羅馬帝國這些外面的.
就是當時的政治的氣候.
怎樣能夠去處理那個的環境.
而教會仍然在這樣的狀況底下.
可以繼續的去安穩的傳福音呢.
這個都需要我們大家一起去學習.
我們可以再看一看一點.
我先開這個.
我們可以看回十三章這樣講.
人人都應當信服在上為掌權的.
這個是新的漢語譯本.
因為沒有權柄.
不是出於神.
這些掌權的.
都是神所指派的.
這樣來到好像很一錘定音這樣.
來到去叫.
因為我們講過了.
十二章就是教會的內部.
怎樣看每一個弟兄姊妹.
看自己是神的恩賜.
得到這個永生的生命.
所以每一個弟兄姊妹都要活著.
奉獻給神.
每一個弟兄姊妹都有一個奉獻給神的心.
然後就要在教會裡面.
首先就是要發現自己的定位.
自己在這個信心底下.
聖靈將不同的恩賜給過我們.
我們大家各就各位.
按照著不同的恩賜.
來到在不同的崗位去配搭,事奉.
這樣自然教會會有一個對外的見證.

$^{81}$對內的和睦.
對內的合一敬拜.
這樣的見證是會影響外面的教會的.
或者從開始猶太教.
耶路撒冷的法利賽人.
是一直敵視基督徒,基督教的.
雖然基督教現在已經不再僅僅在耶路撒冷.
而是一直已經去到安提阿.
土耳其的安提阿.
成為一個新的愛邦的福音的據點.
來到有不同的宣教的事供.
是神去編配保羅和巴勒巴.
所以在這樣的狀況下.
基督徒他們內部的穩定.
自然就成為了一個很美好的見證.
這是一個愛和合一的見證.
但是十二章的那個倫理關係.
不僅僅是教內自己的和好.
自己的開心.
而是很容易就會影響到外面的人.
對教會那個.
當然我們不知道當時他們有沒有很大型的戶外報道.
或者開不同的聚會去吸引人們.
似乎當時不是這麼大規模.
但是無論如何.
羅馬的國教.
羅馬的皇帝政權.
一定是會對一些這麼有規模的聚集.
宗教上的聚集.
他們都會有那個留心,敏感.
還有會不會是一個很重大的勢力.
到底那個教義是不是要跟羅馬本身的法制.
那個君王的治理.
有一些的敵黨違背呢?.
這些一定是實在會引起當局的關注的.
所以在這裡的時候我們又看到.
保羅很快就十二章我們講過了.
十二章的最後是在講教會面對不公平的對待.
甚至乎是一些冤屈.
我們的態度都是需要.

$^{121}$好像有一種的忍受,馴服.
還有不可以就坐.
要為那些逼迫的人.
我們要為他們禱告.
使得能夠仍然可以在有限的空間.
來到去聚會.
還要是傳福音.
一直將這個福音去到不同的地方,地域.
所以一直這樣的思路.
似乎感覺到保羅很想讓教會的弟兄姊妹.
尤其是修信的羅馬教會的弟兄姊妹.
有一個很強烈的意識.
就是十三章第一節這裡.
對不起 講話有些不舒服.
因為有一顆很大的飛之在口裡面.
有些可能是太熱的氣.
所以碰到那個位置就有點不舒服.
大家不要介意.
這個時候這裡就好像有一個很重要的宣告.
給弟兄姊妹千萬不要不知不覺地.
訴諸於地上的人為的一種的權柄.
和一種的衝突.
還有一種的生態想法.
尼塞馬很快會降臨.
主耶穌很快會回到地上.
主耶穌會成為那個掌權者.
這個時候有這種情意結.
猶太人肯定有.
所以他們一直都是和羅馬的政權一直對抗.
由主耶穌升天之後.
那個對抗就更加一直的嚴重.
後來那個憤懼黨.
真的成為了反對派革命的分子.
最終羅馬就透過一些高壓.
來將整個革命剷平之後.
就將猶太人都這麼不順服的這些民族.
就將他們分散到天下.
所以這個就是大概由主後30多年.
去到主後的70年.
你可以想像35年左右的時間.

$^{161}$猶太教就因為這樣的錯誤理解彌賽亞降臨.
彌賽亞執政.
彌賽亞推翻當時的政權.
當然就是羅馬.
所以這樣的錯誤的理解之下.
一方面就失望了.
彌賽亞沒有真的這樣降臨.
耶穌這樣說自己彌賽亞.
也是這樣失望.
但是要透過人的那種對自己國家的情意結.
對自己猶太教,猶大國的情意結.
所以也有一幫人,有一代人.
就用這樣的武力的方法.
想來推翻,趕走羅馬的政權.
好像要回頭去看.
在主前的165年的那個頌光節.
就是當時的敘利亞王.
要來使猶太教的信仰添上很多雜質.
所以就激發起當時一個祭司的家族馬加比.
就來鼓動其他的猶太人.
去趕走,推翻這個敘利亞王.
所以可能對一些很忠誠的猶太教的.
尤其是法利賽的這一幫人.
他們就會有一個很大的心結.
就是我們什麼時候可以真正.
有我們自己猶太人的自由,統治呢?.
所以一直都是無論是暗中.
低調到高調的對抗.
所以猶太人是這樣的時候.
就不斷受到那個鎮壓.
但是現在保羅都意識到.
教會始終都要面對.
當政權留意教會的運作.
那個主張,那個教理,那個傳播的時候.
就一定會看到到底是信服政權.
還是反對和敵對當時的合法的羅馬政權呢?.
所以你可以說這一次聖經.
在整個羅馬書裡面.
是一個很大的宏觀的那種對福音.
對神的國度.

$^{201}$那個治理的一個很重要的價值觀.
來給我們看到.
就是說不要隨便去反對.
也不是我們要做一個.
無論多麼殘暴的政權.
我們都要默不作聲.
甚至要舉手附和.
都不是這樣的.
所以你留意到保羅很想給弟兄姊妹知道.
我們要在神的國度去看.
原來地上的那個政權的交替.
是有神的管理.
這個就超越了我們人的那個眼光.
那個的眼界.
因為我們看到的就是.
可能哪個王推翻哪個王.
因為羅馬的皇帝就是這樣.
所以你會看到.
這個加里居拉.
這個凱撒大帝.
又是被人暗殺的.
所以你看那些古裝的羅馬片.
你都會看到.
其實一代一代都是這樣暗殺上位的.
那你就可以想見到.
哪個那裡有神的掌權.
神的管理.
有神的差派等等.
但是保羅他很清楚.
他不是不懂歷史.
他更加不是不懂當時的羅馬的.
君王的內部的那些問題.
那些的殘暴.
還有那些的荒人等等.
所以我們看到.
保羅在第十三章這裡.
接著第八節到第二十節.
給我們有很大的一個學習.
所以我們看到.
這一次聖經真的好像一個延節詞那樣.

$^{241}$將上文講的教會.
我們自己一定要內部的合一.
我們一定要.
每一個人都認定自己是一個.
獻祭的一個祭物.
獻給神.
因為我們的人生.
整個的心態價值觀.
都是要以神作為我們的守衛.
為神而活.
在教會裡面.
我們能夠有一個服侍.
是神安排的.
然後我們就會面對外面的衝擊.
外面的衝擊.
我們自己要有一個很強烈的神學.
我們自己要首先有一個神觀的教會觀.
我們就會能夠處在不同的時代.
不同的政權.
我們都能夠可以站穩.
因為我們清楚.
那個真正的中心點是在神那裡.
就是說如果不是神來容許.
尼祿都不會作王.
如果大家熟悉羅馬的歷史.
你可能真的會看到.
尼祿都不是一個本來他可以坐的獄位.
你可以想像到.
每一樣東西背後總有神的美意安排.
我們福音派是很好的.
我們是相信這些超自然的.
我們不是自由派.
我們真的是相信神掌管的.
所以既然跟保羅這樣說.
他絕對不是一個很俠麗的政治家.
他更加不是一個目不識丁的.
什麼都不懂的一個普通的無學民的小民.
他真的是一個很掌握到時代脈搏的這位使徒.
他自己本身對希羅的文化.
對當時的政權.

$^{281}$他其實是非常掌握的.
甚至他能夠.
你想一下他能夠夠膽.
那種膽量可以在當時的阿基泊王面前.
來講他的基督教見證.
可以想像得到.
他那個氣場.
其實是他的信心.
他在神裡面對自己身份的確認.
那種自信心.
是來自神給他的使徒這個身份的權柄.
所以當他這樣講的時候.
我們作為福音派的弟兄姊妹.
其實我們就不應該再有懷疑.
我們要認定保羅這樣講的時候.
一定有他的智慧.
有他那個要擺平了弟兄姊妹.
不要每一次都訴諸於用人的力量.
去推翻或者抵擋地上的政權.
這個就變成了.
裡面已經包含了政治的東西.
我們盡量不要有那個盲心的去碰.
反而教會就真的是很超然地.
我們是為了福音而生.
我們是為了敬拜神而存在的.
所以教會我們如果很直接.
用教會,用教務去參政.
或者去用一種這樣的結合的力量.
去對政權有任何的.
例如是推翻,敵對.
這樣就給那個空間.
就是我們變成了一種很反動.
就帶來當時的教會,羅馬教會.
就會受到那個迫害的了.
所以保羅告訴我們.
首先一個很重要的定義.
政權它那個合法的身份是來自哪裡呢?.
保羅就用.
記得不是地上的富裕.
保羅的那個代言.

$^{321}$保羅是為神而代言的.
一切的政權是來自神的安排,管理.
當然它的作用就是第二節已經說了.
是神所命定的.
是神所編配,神掌所來到去.
為什麼開不了?.
我先看看.
可能太….
看什麼.
行.
好了.
可能那個很慢.
不要緊.
我們稍微看看那個分段.
上個星期已經說了.
不過今天我們就會再詳細的說.
在保羅他當前的政治下的模樣.
12章9到第21節.
就是一個服侍的群體.
彼此相愛的群體.
一個就算受迫不得.
當然這個迫不得是來自外面.
受到迫不得.
我們仍然都是帶著愛.
這樣來去謙卑地忍耐.
還有訴諸神的審判,申冤.
我們就不要自己來做一個對抗性的申冤.
對抗性的一種以牙還牙的審判.
因為教會就沒有這樣的力量去做這樣的事情.
更加大前提不是有沒有力量這麼簡單.
而是我們整個存在是愛的群體.
是要彰顯十加福音的那份接納,赦免.
所以教會只能夠在這樣的狀況下.
為自己來去申辯.
剛才我們說了13章第8到第10節.
當然後面還有其他.
就是告訴我們愛或者愛人.
愛不同的人才能夠真正承傳律法.
這裡我們遲些再講.
不是守十誡.

$^{361}$不是來去我們好像要信主不夠完全.
未能夠靠自己得救.
不是這個意思.
而是我們這個新的生命.
信了主的新生命.
我們能夠怎樣不斷達至神的生命的完整呢?.
就是一個信的人是會有愛的.
一個信耶穌的人就會有愛.
一個信神的人就會有愛.
一個等候到永恆的人.
我們是有愛在我們的心裡面.
所以這一種行出來的愛心.
就是基督徒的一個很大的特徵.
所以我們很奇妙地.
我們以前就是沒有愛.
但現在我們有.
但這裡很有趣的是.
大家看到很明顯.
一到第七節是包含在愛的裡面.
在社會事件的時期.
這一段聖經一到七節帶來很大的爭論.
也被一些不同的宗派.
或者不同的弟兄姊妹批評得很厲害.
保羅他這麼久.
這樣去助紂為虐.
這樣去教導無論如何要信服.
如果支持保羅這一段經文的教會.
保守一點的教會.
不要說必戒.
這樣就被人罵得七彩.
但是我看過一篇文章.
有一位我都很敬重的教授.
以前的老師.
寫了這篇文章.
正正就告訴我們.
這一段聖經的上下文.
如果你稍微去解一解就知道.
保羅的意思.
不是在說我們做一個無論如何的順民.
而是在說教會無論如何.

$^{401}$都要有愛.
這樣去面對無論教內教外.
都是要展示我們是一個愛的群體.
愛的群體不是等於我們去做一個順民.
不是做一個我們甚麼都看不到的.
一個這樣的人.
去助紂為虐.
絕對不是這個意思.
這個時候我們就留意到.
在中間這一段.
那個含義就真的很清楚.
告訴我們.
教會不要隨便行駛地上自己的人.
甚至乎是保莊,革命這些力量.
這樣來參與在政治上的那種瓜葛.
反而教會存在.
是要讓人看到福音帶來生命改變的.
這個愛的群體.
所以其實沒有離開到.
保羅不是現在飲惡揚善.
或者做鵪鶉.
甚至乎去做羅馬的走狗.
完全不是.
保羅沒有興趣這些政治.
因為看到裡面的鬥爭.
裡面幾個君王的殘暴等等.
保羅很清楚人性.
很清楚任何的制度.
都不是真正帶人去基督那裡.
得到這個愛的生命.
所以保羅來到很想讓教會知道.
在這段經文裡面.
所以他留意下第三節.
他說「作官的原不是叫行善的懼怕.
乃是叫作惡的懼怕」.
意思就很清楚.
政權,制度是要讓社會處於一個相對的穩定.
是要讓犯罪的人被緝拿,被監禁,被刑罰.
讓社會的制度盡量在一種和諧裡面.
當然絕對沒有說制定法律的掌權的.

$^{441}$自己本身都是犯法的.
保羅不是不知道.
而不試探討那個問題.
所以仍然在這樣的情況下.
教會在一個相對穩定的治理下.
可以繼續將神的福音.
真的改變成生命的福音去傳.
這個是教會很重要.
你看到保羅從來都沒有偏離到他的主題.
是在說我們是一班新的創造.
我們應該是在神的眼中.
我們是這個世界的一個活祭.
我們要服侍.
所以不只是教會裡面的恩賜.
而是對外我們都有這份愛的力量去接納,去包容,去改變.
怎樣改變呢?夏文就說了.
所以你會看到是要讓一個地方的社會相對的和平.
因為作官的是要讓那些作惡的懼怕.
不是讓那些行善的,守法的懼怕.
所以後面就說.
提了兩次.
你只要行善就可以得到他的稱讚.
所以教會在正常的狀況下.
我們就是秉公行義的基督徒.
第四節.
因為他是神的用人.
再一次提到政權是神的默許,容許.
甚至冥冥中是有神的安排.
這樣去做.
任何的政權都不是不公義的.
你留意一下委內瑞拉的馬杜羅.
或者美國的特朗普.
甚至乎在早期的烏克蘭的前總統.
他的官邸.
昨天看了一套無線的網上介紹.
烏克蘭前總統.
我忘記名字了.
他的官邸現在可以成為博物館.
因為大到一個地步.
你難以想像.

$^{481}$要坐車走都走不到.
裡面的擺設全部是價值現成的.
一個皇帝,一個總統.
你可以想像得到.
他在烏克蘭的前總統.
他在前總統的官邸.
當懼怕的懼怕他.
當恭敬的恭敬他.
要我們真的從心底裡面去看清楚.
我們要願意信服在那個掌權底下.
來做一個我們是安居樂業.
更重要的是我們要傳福音.
支持保羅的宣教工作.
來將福音由羅馬傳到其他地方.
所以這一段聖經.
是很貼切地告訴我們.
一個彼此相愛有愛的群體.
是能夠讓未信的.
甚至乎一些政權.
都會對教會有一個改觀.
再留在下面.
9至21字.
是之前我們講過.
第8至10字就是我們今天要講的.
愛就是律法的總歸.
這裡我們講了.
權並是來自神.
所以納粹和尊重掌權的.
接著8至14字.
愛是律法的總綱.
基督徒怎樣讓這些政權外面的人.
是敵黨基督的人.
他們看到基督.
藉著光明的生活 八節開始了.
凡事都不可虧欠人 唯有彼此相愛 要常以為虧欠.
因為愛人的 就完全了律法.
像那不可姦淫 不可殺人 不可偷渡 不可貪婪.
或有別的戒律 都包在愛人如己 這一句話之內了.
愛是不加害於人的 所以愛就完全了律法.
保羅告訴我們 我們不單止是一班好像奇特的人.

$^{521}$因為你不拜君王 你不拜那些假神.
你又來到去和猶太教是有分別的.
我們不只是信耶和華 我們是信耶穌基督.
是人類的救贖主 我們信的是三位一體的神.
這樣的一個基督教 這樣就帶來.
猶太教會迫使我們叛變 移端.
因為你拜的是三位神 不是拜耶和華獨一的神.
所以對猶太教 看我們基督教就會有一種很大的痛恨.
但是對於政權 他們一定是會用他們的宗教.
來強制不同的民族 都要敬拜他們的神.
之後就再要 因為要令到整個社會有一個穩定.
君王就是那個神了 所以一切的好處.
都是由於君王的治理 而令到你們有福樂,平安.
所以必須要敬拜每一代的君王.
尼祿都會看自己是一個神.
所以這裡就讓我們看到 教會的存在.
就必須透過一個新的生命的流露.
來讓無論猶太教或者是羅馬政權.
都來看到我們真是有生命的力量.
來成為一個見證 所以第八到第十四節.
很簡單來說 既是讓猶太人看到.
我們不是拜錯了異端的假神.
我們敬拜的三位一體的神.
是讓我們反而在裡面有愛.
而且愛是根據《舊約聖經》說的.
就完全了律法了 律法的總綱,律法的總規.
律法的總的目的是想我們有對神的愛.
對人的愛 所以可以先記起這些經文.
《出埃及記》的20章.
第20章的13節《出埃及記》.
13節開始 20章其實是10節.
不可殺人 不可姦人 不可偷渡.
不可作假見證 陷害人 不可貪戀人的房屋.
也不可貪戀人的妻子 瀑皮 牛奴.
並不貪一切所有的.
這個就是10節的保羅所引述的.
因為全律法都包含了愛人如己之內.
所以我們看完引述了之後.
應該就是《申命記》的另外一段了.
沒錯 《申命記》第5章.

$^{561}$可以稍微看一看《申命記》.
第17節同樣都是這裡.
都是摩西臨終之前.
向新生的世代.
描述當時在西乃山上的律法.
神頒布的十誡.
再一次17節這一段.
同樣都是不可殺人 不可姦人 不可偷渡.
不可作假見證 陷害人 不可貪戀人的妻子.
也不可貪戀人的房屋 田地 瀑皮 牛奴.
並不貪一切所有的.
這樣的經文就是保羅在這裡所引述.
接著他給我們的教訓就相當意義深長了.
因為當他這樣來引述十誡餘下的五誡的時候.
他告訴我們 或者有其他的誡命.
所有這一切都包在愛人如己這句話之內.
正正好像耶穌所說的.
愛是分加害與恩.
這是保羅自己加上去的.
所以愛就完全阻律法.
這個就是律法的總綱 就是一個愛.
其實十章已經說了.
這個告訴我們 保羅讓我們看到.
原來我們在馴服裡面.
我們雖然一到第七節.
好像拆開了講教會要馴服掌權的.
在當時的羅馬政權.
讓我們能夠平安的.
可以做好我們的本分之後.
教會就能夠有力量作見證.
為耶穌的福音作見證.
其中一個見證就是我們在愛裡面.
所以第八和第十節.
凡事都不可以虧欠.
唯有彼此相愛.
常常覺得做得不夠.
第十節 愛是不會加害人.
不加害人 但是你不對人好.
這樣都是不夠完美的.
愛是完全阻律法.

$^{601}$讓律法的精神能夠貫徹.
在我們一群屬主的兒女身上.
去行使出來.
所以這一段聖經真的告訴我們.
如果教會以為自己可以在地上行使武裝.
行使革命 行使武力 行使仇恨.
其實就不明白了保羅的意思.
當然你可以不引用聖經.
你可以不根據聖經做.
你純粹根據地上的.
以牙還牙 以眼還眼 沒問題的.
但是就不可以說是教會.
更加不可以說是保羅的教導.
保羅的教導剛剛相反.
保羅就告訴我們.
他不是不知利祿.
不是一個好的君王.
但是在一個制度底下.
相對穩定的時間.
我們是要做一個順服的兒女.
我們要去讓那個平安.
地上的平安.
可以使我們能夠在傳福音的.
這個偉大的天國的事業.
不受到阻礙.
如果不是 早就已經被人封了教會.
你根本存在不了.
根本何來有見證可言呢.
不是沒有這些日子.
遲一些會有的.
當政權到了一個極端.
到了一個完全的失控.
就是要使教會不存在.
那個不存在只不過是那個物件上不存在.
就是說教堂不存在.
不是那幫信主的人不存在.
信主的人仍然以愛的守則.
來愛迫害的人.
那你就可以想像到.
原來教會從來就不會沒有的.

$^{641}$是那個教堂沒有了.
不是教會沒有了.
因為教會那幫人.
就是有神生命的人.
我們就是這樣來遵照保羅的教導.
我們仍然的學習.
來在主裡面去接納.
去寬恕 去忍耐.
還有更重要的就是.
愛是我們最重要的.
內心的那個力量.
不可以說是一個武力.
不可以說是一個.
我們的一個武裝的本錢.
愛本來就不是武裝.
愛本來就是一個內心的力量.
就是相信有神.
相信有永生.
相信主耶穌是掌權的.
所以這個愛.
不是一種盲目的愛.
是有意識地來愛一些不可愛的人.
所以基督徒真的有一個這樣的力量.
所以如果真是認識聖經的.
我們就不會做出一些跟社會敵對的事情.
這個不是做一個盲目的信徒者.
而是我們有意識地.
知道教會存在在地上的那個功能.
所以我們應該展示的就是愛的接納.
接著告訴我們後面就是.
所以活在愛和光明中.
原來就是一個很大的見證的力量.
真正的基督徒的力量.
就是愛和光明.
留意第十一節.
「再者你們曉得現今就是該趁早睡醒的時候.
因為我們得救.
現今比初衷的時候更勤了.
黑夜與深白就將近.
我們就當脫去暗昧的行為.

$^{681}$帶上光明的兵器」.
肯定是在說我們的生命.
要來讓聖靈的管理.
讓主耶穌的愛在我們的內心.
十誡都已經說叫我們不可以加害人.
我們要來在我們自己的生命裡面.
有意識地來活在主的聖潔生命裡面.
所以你可以看到.
原來我們就是一班白晝之子的見證人.
我們應該過一些光明磊落的生活.
你就可以想像到.
保羅就是這樣說我們.
所以既然白晝將近.
即是主耶穌快要回來了.
所以我們也不是沒有意識地來等待.
和來馴服.
原來我們是有意識地.
來等候主耶穌的榮耀的再臨.
然後怎麼做呢?.
整個教會就應該要這樣教導弟兄姊妹.
我們就當脫去暗媒的行為.
帶上光明的兵器.
現在不斷討論那些圍標.
那些物料.
或者建築公司.
如果我們是基督徒有良心的公司.
我們就應該那個標準更加嚴格.
就不要隨便帶來的.
只會有事就可以很嚴重.
所以其實重點在這裡.
要過得了自己.
所謂過得了自己就是在神的愛底下.
去看這件事.
來到一個.
大家可能聽過法理情.
做一些事情.
當然你做對的.
在法律上你沒有錯的.
在道理上你也對的.
但是在情感上有沒有過分了呢?.

$^{721}$有沒有很苛刻呢?.
有沒有一些好像不近人情呢?.
所以我們都要留意.
比如有些童工突然要拿假.
女兒病了要進醫院.
有時候她沒有假.
我們都應該要通情,答理.
有時候盡量找到其他童工.
補到位就可以解決了.
當然這些是很小的事情.
但是如果我們有時候每次都講法,講理.
而我們不講情.
有時候就好像變成一個很不近人情.
這樣是不對的.
但是你可以想像得到.
保羅告訴我們.
我們要帶上光明的兵器.
所以很敏感.
剛剛才說要馴服.
不要去對抗.
因為猶太教當時.
就是這樣來到有革命份子.
那些憤慮黨.
曾經跟隨耶穌的憤慮黨的西門.
後來成為革命份子.
最終當然是被打解了.
所以一提到兵器.
有點敏感.
不過保羅很清楚.
我們哪有可能用武力.
我們是用光明.
我們是用愛.
那就是說在愛裡面.
我們自己有一份聖潔的生命.
我們有一份很崇高的永恆的盼望.
我們是有很光明的內心.
所以這兩次.
我們都是幫一些未信的朋友.
辦一些安息禮拜.
辦一些聚會.

$^{761}$有時候人家都很緊張的.
我們已經很清楚有共識.
我們在人家沒有問我們之前.
我們已經一開始跟他們說.
什麼叫安息禮拜呢.
我們就說明了.
教會是一毛錢都不收的.
是義務做的.
大家不要擔心.
不需要有什麼想法.
做完了.
你們很想對教會的福音工作有支持.
有奉獻.
是沒有問題的.
可以自由奉獻.
說不自由呢.
你問我多少錢.
我不會告訴你多少錢.
你可以奉獻一元就行.
一百元就行.
一千元就行.
我覺得三五千元都可以的.
我們一定要很清楚的說了.
這就是光明.
這就是讓人看到.
我們真是所做的.
是要讓他知道.
我們不是故意這樣幫你做的.
暗中收你的錢.
我們不做這些事情的.
反過來你不做.
人家安心.
人家很自然地.
曾牧師是奉獻給教會的.
我們要說明.
我們要發修例.
那個可以免稅的修例.
讓你們知道.
有些說不用了.
我們信一個教會.

$^{801}$不是信不信的問題.
一定要這樣做.
所以光明的兵器.
帶上光明的兵器.
行事為人.
要端正普象.
行在白晝.
這些已經帶來羅馬的政權很尷尬.
因為他們正正就是連白晝都獻落.
白晝都謊言敘酒.
給前書長看的.
白晝獻落.
白晝就是做事.
就是很忙.
如果你是奴隸的.
當時奴隸的社會.
根本是很辛苦的.
白晝獻落就是極權的.
那些管理的階層.
很博尚平民的.
是擁有大權的.
你可以想像得到.
如果能夠白晝.
你的生命行在白晝.
就是說可以在陽光.
在神的面前.
我們都能夠這樣去.
讓神去做一個審判和管理.
這裡告訴我們.
接著有些負面的.
就是說當時的人.
普遍都會這樣.
不可以謊言敘酒.
不可以好色邪道.
不可以憎恨嫉妒.
就是為了自己個人的利益.
總欲欺大主耶穌基督.
不以為肉體安排.
去放縱私慾.
當然不是說平安夜去吃個聖誕大餐.

$^{841}$行不行呢?.
是有一個計劃.
為肉體安排.
去放縱.
早期我初信主真的有這種心態.
有這種罪覺感.
但是慢慢成長了.
慢慢成熟了.
首先不要這樣對號入座.
保羅的意思是.
一些持續性的.
很階層性的.
已經慣性的每天的生活.
就是想著怎樣去為自己的肉體去放縱.
就像那些酒林浴池.
當時的羅馬的皇帝.
或者那些元老院裡的人.
所以你可以想像到.
保羅很想我們有一個品牌.
有一個品牌.
有一個我們是披戴著耶穌.
人們看為是一個死囚.
釘十字架的一個.
這麼噁心的一個行刑的人.
但是我們很樂意披戴耶穌.
在這個帝國裡面出入.
披戴著耶穌在帝國出入.
即是人們嘲笑你.
猶太教是敵視你的.
你繼續說耶穌是我們的尼塞那.
是討厭到你死的.
但是你披戴是怎樣披戴呢?.
就是用生命的聖潔.
那個盼望.
那個光明.
來成為兵器.
成為了讓人看到我們是有愛的一班人.
所以這段聖經最終是說愛的結晶.
就是神的兒女.
因為整個的內心已經有意識地.

$^{881}$由信那一刻是為主而活的活祭.
所以如果每一個基督徒是活祭.
就算將來尼祿下半段的執政.
就燒了羅馬城.
然後就嫁禍給基督徒.
所以基督徒就不能再聚會了.
所以沒有了那個教堂聚會.
就不是沒有了那班人敬拜神.
因為那班人已經是活祭.
那班人雖然分散.
但是他們就是帶著一個生命的見證和獻祭.
這樣去讓那些未信的人看到.
我們一班人是有愛,是有光明,是有盼望.
所以這個就真的是教會最大的資產.
當然我們也不要這麼極端.
教堂如果很值錢的.
已經是在一個很中心的社區.
幾千萬的物業.
我們就不要這樣.
但是真正重要的不是那個硬件.
是那班人的心靈是在神的救恩裡面的活祭.
所以保羅在這裡.
你想一下多麼偉大.
不僅不是一種盲目的順從.
做一個逆星蟲.
做一個好像助紂為虐.
多殘暴我們都要服他.
我們要支持他.
我們甚至閉上眼.
不要管那麼多.
要支持.
保羅不是.
保羅夏文其實是在審判這些在上掌權的.
這些邪惡的生活.
但是教會就讓審判權給教會的人.
教會的審判就是用我們的生命去見證.
兩種的生活.
這一種是上不了天堂的.
我們這一種的生活.
就在神的裡面.

$^{921}$我們現在已經知道自己是有永生的了.
因為我們就是活在神面前.
所以保羅其實好像暗啟一樣.
在審判利祿.
在審判在上掌權.
所以你看到就更加肯定.
一到第七節叫我們那個順服.
不是盲目.
是真的有意識地讓一個平安的社區.
社會可以讓福音可以傳開.
所以其實真的很寶貴.
我們再看其他餘下的PowerPoint.
第十四章我是沒有預備的.
還有那個很痛.
裡面那個飛之.
所以利祿的時間.
羅馬帝國的政治環境.
利祿的時期.
他是要規定秩序至上.
剛才說了.
納粹一定要.
他看得很重要.
和要順服君王是必須的.
所以由他之後的君王.
全部都是君王敬拜的.
將自己看為神.
但是對於教會來說.
我們就要留意了.
猶太人的處境和基督徒的處境.
我們就要留心.
基督徒要面對猶太人這種對國家的敵視.
所以我們就要避免.
免得我們就好像猶太教一樣.
好像敵對叛亂的政權.
所以保羅這樣的教導.
都是非常智慧.
保羅能夠.
以他自己猶太人.
他沒有了這種亡國的.
這樣的情義結.

$^{961}$可想而知.
保羅已經看到有永生的神國度.
更加大於地上的猶太國度.
一個拉比.
能夠看通了.
不是地上的國度.
而是永恆的神國度.
才是真正的重要.
所以才能夠教導教會.
教導其他人.
即猶太人.
猶太的信徒.
用一種敵視的態度去對抗.
因為很快就會被鎮壓.
教會就根本不存在了.
所以尼塞拉他期待.
他就要澄清.
神的國度.
這裡很重要.
神的國度不是政治的革命.
這段聖經對一代的年青人.
早幾年的社會事件.
這段聖經帶來他們非常之反感.
保羅是一個.
好像效忠政權.
助紂為虐.
一概都不聽.
其實那個經文根本不是這個意思.
所以保羅告訴我們.
我們期待的尼塞拉.
神的國度.
不是一個地上的政治革命.
所以當年青人聽到這些經文.
就不會做出一些太過不好的事.
當然我們沒有權力判斷.
因為很多東西根本掩蓋著我們.
沒有人知道.
所以權柄是神設立.
我們再做一點總結.
因為我沒有預備十四章下面.

$^{1001}$就是為著有秩序的治理.
不至於帶來混亂.
所以這個是相對每一個政權.
存在的必須的.
是神的智慧安排的.
所以如果有些國家沒有政權.
是沒有政府.
那些地方是更嚴重的暴亂的.
最終人民一定要有一個紀律的.
所以始終都需要有一個制度的.
基督徒的使命就是傳福音.
不搞政治的顛覆.
不參與一些革命.
保羅這裡是很清楚.
我們的兵器就是.
我們能夠好像有一個這樣的.
帶著光明的兵器.
光明就是聖潔的生命.
活出愛和聖潔.
等候基督再來的群體.
不斷告訴人知道永恆的真實.
現在我們從星期四的老人家的團契.
差不多每個星期都告訴他們.
死是不怕的.
因為我們是有永生的.
很多老人家都很明白.
因為每個人都自己有很多病痛.
現在都不怕死的.
有一個瑞文婆婆說我不怕死的.
死了都無所謂.
一離開世界就見主了.
就是很大的那種豁達.
不是每個人都這樣想的.
所以我們是一班等候主再來的群體.
我們要在地上活出愛和聖潔.
所以如果我們不用了保羅的經文.
去濫用了他.
就變成在上有權柄.
人人都要信服.
就變成支持一個專制政權.

$^{1041}$就不是保羅的意思.
保羅是要養一個相對穩定的社會制度.
那個制度是怎樣呢?.
保羅說他不是沒有興趣.
他現在是在教會福音的差遣的角度去看.
我們處在那個社會.
所以我們信服是為了良心的.
是為了英文.
因為神的管理著的.
所以如果我們將這些英文濫用.
就變成真的支持專制政權.
去強迫其他人都必須聽.
那就不對了.
但又不是極端.
所以平衡的觀點就是信服神.
不是信服人.
我們要看這段經文.
最後第五章.
《史獨行傳》.
你聽聽彼得怎樣說.
第五章就是《史獨行傳》.
經文是在29節的.
但我們可能要讀多一點這段經文.
第17節.
可能讀多一點.
但第12節這一段.
對不起我要先喝口水.
因為有點乾.
這裡有一個完整的故事.
接著就會探討第29節.
主耶穌.
使徒的手在民間行了許多神蹟奇事.
他們都同心合意在所羅門的廊下.
其餘的人.
就是聖殿其中一個廊.
其餘的人沒有一個敢貼近他們.
百姓卻尊重他們.
信而歸主的月法真天.
連男帶女很多.
甚至有人將病人抬到街上.

$^{1081}$放在床上.
或者是浴紙上.
只望彼得過來的時候.
或者得他的影兒照在什麼人身上.
還有許多人戴著病人.
和被污穢纏磨的.
從耶路撒冷四圍的成因.
來到全部得到醫治.
一個很大很徹底的神.
透過彼得和使徒.
來做了一個很大的見證.
就是讓人看到耶穌基督的能力.
在這些病人身上.
所有人都起來.
接著第17節.
大祭司和他的一切同仁.
就是撒都該教門的人.
都起來滿心忌恨.
就下手拿著使徒.
收在外監.
但主的使者夜間開了監門.
令他們出來說.
天使對他們說.
你們去站在殿裡.
把這生命的道都講給百姓聽.
就是猶太人在耶路撒冷的時候.
使徒聽了者說.
天亮的時候就進殿裡去教訓人.
因為主的使者將使徒們帶回家.
離開了監獄.
離開了他們囚禁的地方.
所以使徒是完全不怕不懼怕不走.
聽了天使這樣說之後.
天亮就再進聖殿去教訓人.
就是說好像和法利賽人在曬冷一樣.
繼續進去聖殿教訓人.
就令到大祭司更加生氣.
大祭司和他的同人來了.
叫齊工會的人和以色列族的眾長老.
就差人到監裡去.

$^{1121}$要把使徒提出來.
很奇怪.
一邊睡醒以為使徒還在監裡.
準備開一個嚴肅的審判會.
批鬥會.
就開始叫壁丁去監裡.
把使徒提出來.
第22件事很好笑的.
但差人到了監裡.
不見他們在監裡.
就回來稟報說.
我們看見監牢關得極妥當.
看守的人也站在門外.
及自開了門.
裡面一個人都不見.
要麼就蒸發了.
要麼就不知什麼原因.
他是一個有血有肉不是靈魂.
物質的軀殼是怎樣穿過監門的牆.
就好像周星馳的電影.
要蒲肉都要撞爛牆才能走.
外面還有人站在那裡.
全世界都沒有人這樣越獄.
天使帶走彼得.
天使就是靈體可以走.
彼得怎樣走呢?.
不能走.
所以難以想像這種神蹟.
所以他們都不明白.
當他們這樣的時候.
下面就這樣說.
第24節.
守殿官和祭司聽見這話.
心裡就泛濫.
不知這事將來如何.
然後突然有個人來稟報.
你們收在監裡的人.
現在站在殿裡教訓百姓.
已經很奇怪了.
簡直覺得很害怕.

$^{1161}$為什麼?.
那些是不是神人呢?.
為什麼他們可以穿牆過壁呢?.
彼得可以自己出入呢?.
那個兵兵還在外面看守著.
難以想像的就是.
現在他們不是逃走.
還在聖殿裡.
如若無人地繼續教聖經.
教人信耶穌.
這種力量就很大了.
所以留意一下.
在這樣的氣氛下.
彼得說這句話.
順從神不順從人.
是應當的.
所以我們就不要濫用聖經.
隨便別人叫你做事.
我順從神.
我不順從你.
這樣就不對了.
不是這個意思.
留意到經文.
他們很泛濫的時候.
然後第26節.
於是守殿官和警察.
就去帶使徒來.
留意並沒有用強暴.
因為怕百姓用石頭打他們.
因為這些使徒.
已經成為了很超然的.
在百姓的心裡面.
一個很超然的.
一定是遇見過神的人.
所以守殿官都不敢用武力去捉拿.
去威脅.
反而邀請彼得去作見證.
去到大祭司的院裡面.
去法庭去講他們自己的見證.
到底什麼事.

$^{1201}$因為起初你們捉拿彼得.
是因為很多人聽他們講.
關於耶穌的救恩.
所以27節.
「大道了便叫使徒站在公會前.
大祭司問他們說.
我們不是嚴嚴地禁止你們.
不可奉這名教訓人們.
你們倒把你們的道理.
充滿了耶路撒冷.
想要叫這人的血歸到我們身上」.
留意29節.
這節聖經等於平衡了.
什麼叫「馴服在上掌權」.
留意29節.
現在明顯是一個禁令.
禁止彼得和使徒們.
去講一個訊息.
就是嚴嚴地禁止.
說耶穌就是那一位彌賽媽.
我們嚴嚴地禁止你們.
不可以再奉這個人的名.
去教訓其他人.
奉什麼名.
就是奉耶穌的名去醫病.
奉耶穌的名去趕鬼.
奉耶穌的名告訴他們.
你們等候的彌賽媽就是耶穌.
就是我們所信釘十字架的耶穌.
就是我們今天講的完整的福音.
所以禁令很清楚.
留意彼得怎麼講.
彼得和眾使徒回答說.
「信從神,不信從人,是應當的.
你們掛在木頭上,殺害的耶穌」.
所以原來是第30節.
所以不要抽這些字.
又抽來上下文.
什麼都信服神.
我不信服你了.

$^{1241}$這樣不對.
這樣就是你在濫用經文.
因為原來上下文.
彼得就解釋了.
什麼叫信服神.
不信服人.
就是你講的.
你禁止我講這個重要的信息.
這個我不能夠信服你了.
因為這個是必須要信服神.
給我們人類有這個訊息.
就是福音.
所以第30節.
「你們掛在木頭上,殺害的耶穌」.
就是那個木匠.
就是那個死囚.
就是那個被咒詛的人.
就是那個被禁止去傳講的人.
彼得就說.
「我們祖宗的神」.
就是耶和華神.
「已經叫祂復活了.
神且用右手將祂高舉」.
就是已經復活了.
「採神寶座了.
叫祂作君王,作救主.
將悔改的心和赦罪的恩賜給以色列人」.
就是信的人就會得到赦罪,悔改.
成為新的以色列人.
所以第32節.
「我們為這事作見證.
神賜給信從之人的勝利.
也為這事作見證」.
很清楚的上下文.
如果我們真的要講信服掌權.
就是在神的真理下去信服.
如果違反了神的真理.
我們就不能夠信服掌權的.
不信服帶來的就是坐牢.
要殉道.

$^{1281}$所以保羅是殉道.
彼得也是殉道.
所以不信服神的真理.
我們就不能夠不信服了.
就是說我們不能夠信服掌權.
關於神的真理.
如果不是關於神的真理.
我們就願意信服來做.
但是關於真理的時候.
我們不信服不等於不服從.
就是說現在我們要判你坐牢.
你也只能夠接受.
因為在這樣的狀態下.
受到宗教的迫害.
就不是因為你自己無端端反對政權.
而受到牽連.
所以一定是在這樣的情況下.
所以接著工會就聽見極其牢牢.
就要將他們收監.
所以彼得是預了坐牢.
但是就說得很清楚.
這件事我們要為主耶穌作見證.
這件事我們是絕對信服神.
不信服你們的命令.
所以如果我不信服這個命令.
那個後果就是要監禁.
就麻煩悉從尊便.
我們是預了接受監禁.
這樣叫做不信從.
這樣叫做信服掌權.
或者這樣叫做不信服掌權.
關係到真是宗教上神的真理.
這個就是最重要的價值觀.
所以大家要留意這裡.
今天我們怎樣面對這個不公義的政權呢.
就是避免政治的對抗.
專注福音的使命.
除非福音的使命受到挑戰.
否則我們都是在整體的情況下.
仍然在這個社會制度下.

$^{1321}$做一個信服的公民.
我們避免教會進入那個漩渦裡面.
來參與政治的對抗.
這也是當時在社會事件的時候.
我自己的立場.
根據保羅的教導的立場.
所以我是很反對頂智慧要上街.
甚至要用武力禁止.
不應該是這樣.
當然可以和平表達.
但是如果你覺得和平表達是不行的.
政權不聽的話.
那你繼續要再升溫.
人家都會繼續升溫.
那就不會有一個停止的.
這就等於是被一些有用心的人.
用這樣一個對抗.
所以教會是避免參與政治的對抗.
我們的責任是專注福音的使命.
所以你看到彼得是很清楚的.
不信從人 信從神是應當的.
這個應當就是說.
你叫我不說耶穌是不行的意思.
就不是他們做一個不斷去抗議.
對抗這些法利賽人.
或者這些大祭司的人.
因為耶穌都是這樣死在大祭司的手下.
是彼得自己親眼看到的.
難道彼得要做革命去為耶穌討公道嗎?.
耶穌有更大的使命就是.
宗教上為神的救恩.
祂成為那個代罪的羔羊.
所以就不是政治的迫害.
耶穌是甘心樂意成為真正的贖罪制.
所以彼得是沒有一個時代的人.
彼得保羅為耶穌上街去對抗政權.
對抗猶太人 對抗羅馬人.
是沒有的 因為耶穌就是為神來成為贖罪制.
所以盼望大家能夠幫助我們一起在愛.
在光明裡面 在盼望裡面.

$^{1361}$來等待這個終極的神國.
而不是只是我們這個地上的政權.
所以我們是不負面的.
不會做一個完全不積極的.
不是的 我們的積極就是在整體裡面.
我們是一個光明聖潔的群體.
今天就說到這裡了.
又要有些時間休息.
因為有些事.
先跟大家做一個眉頭.
再一次感恩.
我們都叫做第十三章.
都很詳細的講解.
因為他處身的.
正正是一個羅馬敵黨神的政權.
一個對猶太教都是非常之強大鎮壓的一個政權.
但是我們看到教會能夠繼續在保羅的教導下.
來到去安穩地活.
安穩地去傳福音.
當然後期因為是尼祿自己的那種狂妄.
那種好大喜功.
來到去轉嫁給基督徒來受害.
但是整體而言.
基督教仍然是以愛,以光明成為兵器.
這個這麼重大的原則.
求主給我們每一個時代的弟兄姊妹.
都能夠銘記我們在救恩裡面.
其實我們最大的武器,最大的兵器.
就是光明,就是愛,就是永生的盼望.
願主幫助.
就算給他們的教產都被沒收.
但是他們的見證是更加的真實.
有更多的猶太人因為這樣信主.
願你大力幫助.
這個星期二在我們洪恩家.
都有李頂光牧師的福音主人.
求你使用見證加賓.
求你使用我們這個福音聚會.
讓更加多人來信主.
有些我們上個星期接觸過的大埔的災民.

$^{1401}$都搬了來我們附近的駱橋樓.
盼望他們都能夠收到這些傳單.
可以來聽福音信息.
我們這樣仰望禱告奉耶穌基督的名.
阿們.
\newpage



\section{羅馬書}
\label{sec:aePwbTZJwec}
\textbf{2026-01-02 周五晚查經班《羅馬書(32):分題:如何處理教會內「非核心信仰」的差異》 講員:曾裔貴牧師}
\newline
\newline
連結: \href{https://youtube.com/watch?v=aePwbTZJwec}{\texttt{https://youtube.com/watch?v=aePwbTZJwec}} ~~~~ 語音日期: 2026-01-03
\newline
\newline
\hyperref[sec:ego6cJukrIE]{\small{< < < PREV SERMON < < <}}
~
\hyperref[sec:index_chronic]{\small{[返順時目]}}
~
\hyperref[sec:index_scriptual]{\small{[返順卷目]}}
~
\hyperref[sec:nB_6d0PSa6M]{\small{> > > NEXT SERMON > > >}}
\newline
\newline
$^{1}$第二次圍魂灣我們一起祈禱.
再一次我們很感恩過了聖誕節的長假期.
也有不同的特別聚會.
無論是舞堂的平安夜崇拜.
聖誕節的崇拜.
還有我們在雄恩都有來到.
有地區在洪水橋.
每個屋苑的大型的佈街恩.
很盼望神的福音.
藉著不同的教會.
不同的弟兄姊妹.
可以來讓更多的香港人.
或者其他地區的人.
能夠認識救世主的降生.
改變真真正正將人.
無作的罪人來改變.
的一個唯一的機會.
求主幫助.
我們回來看《羅馬書》.
多謝您讓我們又進入第十四章.
在教會的內部.
怎樣處理一些不同意見.
看法甚至乎良心的一些事情.
我們求主幫助.
讓我們能夠透過.
一段本來講這麼重的.
救恩的教義.
好像一些的問理.
一些真理的問答那樣.
由保羅來按照神給他的啟示.
來講出來的.
但是當來到教會的內部.
又讓我們看到在不同的處境.
有不同文化背景的猶太人.
外邦人一起在羅馬教會聚集.
就是說保羅可能都意識到.
不同的文化聚在一起.
一定會有一些的差異的.
這些差異用什麼的手法去平衡.
去能夠和衷共濟的聚會.

$^{41}$這些都是比文化上的差異更加重要.
多謝你讓我們可以一起去商討.
一起去了解保羅的教導.
我們在你面前等候.
求主施恩.
我們這樣感恩禱告.
奉耶穌基督的名.
阿們.
我們就很感恩.
我先開大一點的螢光幕.
如果我們停了雖然一次.
但是十三章,十四章的分佈.
都將愛和警醒.
成為了兩個很重要的主題.
警醒就是帶來我們在地上的盼望.
所以都是環繞著愛心的.
無論在社群,社會.
因為十三章是說到.
我們要對在上的羅馬政權有一個的馴服.
這個馴服就是基於每一個管治都是來自神的.
所以教會不是特別得天獨厚.
和那個社會脫節.
不需要活在地上.
不是的時候.
就是要按正常的情況.
我們就要納粹.
順從那個黑己.
來到他在一般的情況下.
來做一個順命的兒女.
順命的教會.
順命的公民.
這樣來使全福音的運作.
一直都能夠繼續讓更多的人.
包括在上的官員.
羅馬的君王.
能夠有機會去聽聞福音信主.
所以一直這個機制運作.
都是這樣來進行的.
甚至乎到了主後70年.
猶太人受到高壓的統治.

$^{81}$來變成為.
而為平地的整個羅馬當時的提多將軍.
來圍城.
使得整個耶路撒冷血流成河.
接著所有的猶太人.
就要被分散去到歐洲不同的地方.
這個是歷史告訴我們.
就是那個對抗帶來的.
沒有真正帶到的猶太教有一個平安.
反而基督徒.
基督教同樣都是不能倖免的.
都是受逼迫的.
但是基督的信仰.
不因為這個逼迫而停止.
沒錯.
逼迫帶來的就是那些教產.
就是那些教堂.
或者基督徒不能夠正常的聚會.
這個是高壓帶來的.
令到他不斷的受到壓迫.
但是這種化整為零.
去到不同的地方.
有一樣東西是不能夠被壓制的.
就是每一個基督徒內心裡面的那個信仰.
對三一真神的凝訊.
這個是無論如何.
怎樣壓迫都不能夠改變的.
好像我們引述過《使徒行傳》提過.
彼得就說.
信從神.
不信從人.
是應該的.
你們要我們不准信奉耶穌.
這位拿撒勒人耶穌.
是我們的救主.
是我們的尼西亞.
你要我們否認這件事.
是不可能的.
你可以將我們監禁.
但是神是會將我們解救的.

$^{121}$所以你會看到彼得他們就是用一種信念.
信心.
在丁伯底下.
讓福音繼續運作.
所以連加瑪尼也不敢對這班基督徒.
有太多的高壓的箝制.
反而就是說.
如果神真的要將他取替.
自然他自己會瓦解的.
所以如果我們基督徒的信心不夠強.
就算平安.
只會令我們更加放任.
但是相反.
壓迫.
壓力.
迫不及待.
反而令我們更加珍惜.
原來我們是屬神的.
所以.
你會看到十三章是在正常一般情況下.
我們基督徒要在社會上有公民責任.
接著.
很快就去到第十三章的後半段.
就是我們要過一個光明的兒女生活.
大前提就是.
黑夜已深.
白晝將近.
耶穌來這個世界成就救恩.
完成救恩.
聖天.
就是這個世界已經被定罪.
很快神的永恆救贖就要全面展開了.
所以在保羅的眼中.
或者在使徒們的眼中.
主耶穌來到這個世界.
就是末世的開始了.
所以.
主耶穌再來.
就是我們真真正正.
這個永恆的人類歷史的結束了.

$^{161}$所以現在我們就是在耶穌降生之後.
和耶穌再回來之前的這個時代.
就是保羅所很有詩意的說的.
黑夜已深.
就是說可能現在是凌晨的兩三點.
可能現在已經是四點,五點.
就是說很快就到清晨.
很快就有清晨的日光.
就會有神希的破曉.
就是看日出的破曉.
就快會來了.
這樣的意思.
所以保羅就在第十三章的十一節.
引進一個很重要的.
就是愛之餘.
基督徒要有愛心.
無論對社會上的社群.
或者政權.
但是對永恆.
保羅想告訴我們.
有一個盼望.
我從第十三章的十一節讀起之後.
就進入我們第十四章.
講到將那個焦點.
又再一次帶回教會的內部.
怎樣能夠帶來和平的合一.
仍然最終不是說飲飲食食.
不是說一些的節日.
哪些是重要的.
相反更加重要的就是.
能夠有一個和諧,合一.
彼此的尊重,彼此的欣賞.
這樣的大前提下.
最終仍然是帶來的那個.
很小要聽到神的福音.
就是羅馬的教會要支持保羅和使徒們.
他們宣教的工程.
所以不是教會內部這些這麼小的矛盾.
保羅你不要這麼擔心我們.
不是的.

$^{201}$保羅就是要意識到.
最重要就是大家在救恩底下.
來去謀求的那個同心,合一.
以至真正將最重要的資源投放.
就是福音要傳到地極.
傳到更多未得的民族當中.
所以十三章的第十一至十四節.
又是一個很重要的.
每一個基督徒.
沒錯都已經是神救贖的活祭.
應該人生為神而活.
應該來到為到主的榮耀.
這樣來到去生活.
每天都是一個活祭.
這樣來到去呈獻給神.
但是更加重要就是.
我們怎樣能夠在政權底下.
社會,在公民,在人看到我們的範圍.
我們是一班怎樣的人呢?.
所以猶太的猶太教.
和基督徒的新興的基督教.
就有兩個很大的分別.
基督徒好像什麼都盲從.
什麼都信服.
但是猶太人在耶穌升天之後開始.
那一班憤慮黨繼續抗爭.
繼續對抗羅馬政權.
最終就高壓地去剿滅.
你就沒有了你們的存在.
但是基督徒仍然為神作見證.
怎樣作見證呢?.
這裡又很重要了.
原來是一班光明的基督徒.
光明的兒女.
所以保羅開始引進光明和黑暗的分別.
或者說有救恩和沒有救恩的人的生命分別.
所以羅馬人都不是蠢的.
他們過著一些敗壞的生活.
比特前書說白晝燕落就是那些政權的人.
白晝燕落.

$^{241}$那種燕落不是抓一下就走那麼簡單.
是狂歡作樂的燕落.
白晝是工作的時間.
你想一下現在不是放1月1日的假期.
古代的人是很忙碌的.
一定是有很權貴的階層.
才會這樣去享受奴隸的勤奮得來的成果.
所以你可以想像到.
基督徒就要進入一個新的社群.
就是一個光明兒女的社群.
第十一節.
再者,你們曉得,現今就是該趁早睡醒的時候.
因為我們得救,現今比初信的時候更勤了.
甚麼叫得救呢?.
就是再信耶穌,再聽福音,再走出去決戰?.
相信肯定不是.
是一班已經明白了救恩的人.
這個所謂的得救.
即是說距離我們離開這個世界.
去到見耶穌基督.
或者我們的屬靈生命一直成熟,成長.
當然,未必是每一個時代都有見到耶穌回來.
一定是每一個時代.
我們都有我們人生的那個年限.
我們完了,我們就去見主了.
所以我們得救,現今比初信的時候更勤了.
即是說我們應該對於我們整個的生命.
我們自己有一個責任.
就是我們應該怎樣比初信的時候更加成長呢?.
更加對罪的敏銳呢?.
更加對自己的內心有一種敏銳呢?.
更加對神有一種渴慕,追求和愛呢?.
這個就是比初信的時候更勤的意思.
不過留意第十二節.
引進那個盼望.
盼望就是要令到我們真的有一份生命帶來的.
光明兒女的生活.
所以保羅用一個比較負面的.
負面就是說在當時的社會其實相當普遍的.
「黑夜已深,白晝將近.

$^{281}$我們就當脫去暗霉的行為.
帶上光明的兵器.
行事為人要端正.
好像行在白晝.
不可謊言醉酒」.
這個是例子.
就是一些人盡皆知的了.
就是說其實就等同於.
就是用一個很迂迴的手法的口吻.
去譴責當時的權貴.
有特權的人士.
他們那種沒有神的生活的形態.
就是這樣.
一個奴隸怎麼可能會這樣謊言醉酒呢?.
謊言.
荒唐的宴會.
就是吃到飽到不得了.
然後很糟蹋食物等等就叫做謊言.
醉酒.
如果你去過一些窮奢極侈的食物.
可能是很刁鑽的.
很講究的.
吃的東西很昂貴.
但是大量的來糟蹋那些食物的.
而很多的人根本吃都沒得吃的.
你可以想像得到.
那種對比,那種落差.
那個背後的精神.
才是神更加的厭惡.
不是說食物.
神就不准我們吃得稍微好一點.
而是我們如果是一種這樣的全心的謊言.
這種其實是在神的眼中是更加的不好.
不可好色邪道.
不可憎勁嫉妒.
好像將我們人類裡面幾樣東西.
貪圖錢,貪圖女色,貪圖權力.
好像已經否釋了出來.
是人類生大的要害.
這樣說了出來.

$^{321}$就是說在這樣的價值觀背後的.
這種壞的心態底下.
背後所承載的.
有不同形式的包裝的.
比如我不知道是不是真實.
可能有些人已經知道.
原來現在觀塘海濱.
大家不知道有沒有聽過一個新名.
叫Generation Z.
Generation Z就是00後的那些.
就是現在大約20歲的那些.
如果到晚上 尤其是長假期.
是一些很放縱的地區.
還要是明仔跟我說.
有時候開車載他上班的時候.
他說有時候他去看那些留言區就會這樣說.
有時候那些殘廁,公廁.
就會是那些年輕人.
來做一些不道德的事情.
原來觀塘海濱有一個新的名稱就是這樣.
剛剛昨天就有一宗案件.
就是一個13歲的女孩打一個15歲的女孩.
就是說是一個很普遍的情況.
這樣就可以想像到.
保羅要求的基督徒.
一方面不是沒有意識地.
成為一個好像聽命的,盲徒的.
在社會上一個好像手握搏雞之力的公民.
不是的.
保羅想基督徒有意識地成為一個良好的公民.
這個良好是要在生活裡面.
是一個新的群體.
這個新的群體是一個光明兒女的群體.
有的是光明的兵器.
留意一下 保羅用這些話其實是很敏感的.
兵器就是說你是否想推翻羅馬呢?.
但是那個兵器原來是內心的新生命.
是內心有意識地要意識到權力.
色情,金錢這些權力.
對於一個人的腐化.

$^{361}$所以第14節.
因為光明的兒女這個群體.
這一群羅馬教會的弟兄姊妹.
總要披戴耶穌基督.
不要為肉體安排.
去放縱私慾.
你可以想像到.
整個的意識不是每天就想著去玩.
去玩到盡.
而是我們要披戴著耶穌.
我們要讓主耶穌這個福音.
這個新生命.
讓羅馬的政權看到.
為什麼我們要納粹呢?.
是因為我們要活在這個世界.
信服神的旨意.
為什麼我們要信服呢?.
我們就是因為我們是一群.
在社會上為耶穌作見證的兒女.
但是當你要我否認耶穌和福音.
那就不行了.
這個就是極之核心的信仰.
就不能夠信服了.
你可以將我們監禁到羅馬坐牢.
但是你不能夠將我裡面的這個信仰.
這個對永恆真理的盼望.
來去禁閉.
因為你禁閉不了我的內心.
所以你可以想像得到.
羅馬政權所見到的.
是一群猶太教的知識分子.
是一群跟隨保羅的教導的忠心.
在你上面開始過聖潔生活.
有盼望的生活的神的兒女.
光明的兒女.
不是那些中世紀的光明會.
那些秘密組織.
是一群真是光明的.
可能他們的身份仍然是奴隸.
仍然是有猶太的背景.

$^{401}$因為14章就已經說了.
仍然是有一些外邦人來自不同的國家.
有帖撒羅尼迦人.
有特羅非摩的就是帖撒羅尼迦.
有特比人的.
忘記了什麼名字了.
接著就是有不同的地方.
所以你可以想像得到.
每一個基督徒來自不同的背景.
都要披戴主耶穌.
都要告訴別人.
我們是信了耶穌的.
我們是神的活祭.
我們準備好為主耶穌作見證.
甚至乎要犧牲生命.
我們都要為主耶穌作見證.
我們已經準備好.
我們有這個永恆盼望.
因為主耶穌回來的日子已經近了.
所以黑夜已深.
白晝將近.
所以你可以想像得到.
我們一直在等待耶穌的再來.
這就是13章結尾.
就是說這裡提到一個盼望.
盼望帶來的一定是警醒.
所以披戴主耶穌.
原來是一種對我們生命的很清楚的檢視.
不再屬於舊有了.
不再屬於以前了.
不再屬於罪惡.
這是回到第六章的羅馬書所說.
已經將我們自己生命的主權.
主動地交出來.
信服聖靈.
信服耶穌基督.
所以這個披戴是要每人主動的.
這個時候我們就能夠為主去用我們整個人生.
所以現在AI的年代.
我們又要學一下.

$^{441}$我們又要了解一下.
就是幫我們更加有意識地.
我們應該怎樣每天為主而活.
所以這個帶來的盼望警醒.
14章又再說回這個題目.
很特別.
最後第14章.
所以我們可以先停一停.
你可以想像得到.
在保羅的意識裡面.
很明顯就是13章的段落.
13章的一到第七節.
就正正是一個插日的段落.
本來羅馬書是在說教義.
是在說什麼是恩順清義.
什麼是人類唯一的盼望.
什麼是我們在主耶穌裡面.
我們就有這個永恆的.
跟神關係的恢復.
我們就能夠等待那個黎明的日子.
主耶穌再回來的時間.
或者我們離開世界.
與基督同在.
這個是教義永恆救恩的啟示.
透過保羅用他的希臘文.
寫出來的這封書信.
這樣來去展述.
大家讀羅馬書.
你就會跟其他的書信.
保羅的書信有一點不同的就是.
羅馬書是沒有針對教會.
一些錯誤教義.
因為根本他的意思都不是想這麼說.
羅馬書是想說.
什麼是基督教的教義.
什麼是基督徒.
要必須信守的真理.
而且保羅很希望.
他現在準備一路向歐洲大陸.
去傳道的時候.

$^{481}$他就很想先來去知會羅馬的教會.
讓他們知道他的身份.
他的使命.
他是神的使徒.
他現在就是要來去帶領更多人.
進入更多微德之盟.
特別他提的十五章.
講到西班牙那個地方.
就是說當時福音不是很大規模的.
去到任何一個角落.
而是需要有使徒們.
這樣來深入這些地區.
但是更寶貴更重要.
保羅首先很有系統地.
將重要的這些教義真理.
這樣寫下來.
使後世的人讀的時候更加有條理.
更加知道我們在信守的是什麼的福音.
因為眾說紛紜的時候.
什麼才是福音呢?.
福音要怎樣來承載它呢?.
所以在羅馬書就很詳細了.
跟其他書信唯一不同的地方.
羅馬書從來都沒有譴責.
教會出現了對真理的偏差.
或者哪些情況下.
有一些隱憂.
就沒有這樣說的.
反而就是很正面地.
很有系統地去闡述整個福音是什麼.
以及怎樣將這個福音承載.
這就是羅馬書相當重要和清楚的邏輯.
來到這個時候.
保羅拆開了港十三章一到第七節.
就是對當時羅馬的政權.
應該怎樣來教會的自處.
教會推翻它.
為敵.
還是我們仍然是一個公民的社會.
來做一個公民去馴服.

$^{521}$不因為信了耶穌.
成為了光明兒女.
就另外一個群體.
不需要履行遵從任何地上法律.
因為我們有的是天上的法律.
這樣是說不通的.
因為地上的政權未信的人.
仍然是在世.
神容許他們繼續這樣來生活在地上的.
所以教會唯一要做的.
就是以光明兒女.
披戴耶穌這樣的新生命.
來活在地上.
所以保羅拆開了講羅馬政權關係之後.
回到第十四章.
就立即再一次說.
教會內部一些跟真理完全沒有關係的東西.
我們怎樣來處理那個分歧呢?.
這就很重要了.
也再一次重申一件事.
就是說十三章一到第七節.
這樣的插入句.
就顯得分外的重要了.
因為這樣的時候.
因為十二章最後的那段段落的倫理關係.
人與人之間的相處.
應該有新的等格.
十四章再一次說.
要來到怎樣和諧共處的邏輯.
那個重要的精神.
當然就是彼此的接納.
不要論斷.
接納與論斷剛剛有一個對比.
一個相反.
可以想像到十三章這個加插.
是一定有他的邏輯思考.
以及分外顯得重要.
就是保羅也很想教會.
不要跌進一個抗拒政權的思維裡面.
不知道保羅是不是有一種怨憂.

$^{561}$真的看到猶太教殺死了耶穌之後.
繼續似乎那些主戰派.
更加嚴肅的要結集他們猶太人的力量.
要來準備好像來一次大的暴亂那樣.
其實最終就是主教70年完全被瓦解.
所以保羅在這裡的時候.
他真的有這個意識.
教會處在社會.
必須要有這種價值觀來看.
所以不是一種盲從.
但是更加不是一種抗拒.
一種用武力去對抗.
否則只會是一個一定會失敗.
根本不是神的心意.
這個才是更加重要.
十四章很重要.
這一段聖經第一節.
我就將整個希臘文引出來.
因為現在AI真的很方便.
你叫它打出來.
它馬上一秒鐘就出來了.
不用你再翻譯.
我有新的希臘文聖經.
但是已經不可能.
因為現在那些字你想打都打不到.
輸入法打不到.
所以很方便.
其實不難的.
不過已經完全忘記了.
十四章.
信心軟弱的.
你們要接納.
但不要辯論所疑惑的事.
又不准說.
又不准討論.
不要辯論.
你疑惑的就有個疑惑.
是不是這個意思呢?.
當然不是.
重點是接納.

$^{601}$那個大主旨.
那個主題.
是在說要接納.
這裡說的.
不是說不可以討論.
不是說不可以去釐清.
那個事情的背景的來龍去脈.
所以第二節就開始了.
保羅就立即說.
有不同的人.
他們的信心.
他們的良心.
這個不是涉及恩信精明的信心.
是那個人.
他的傳統帶來的良心.
這裡就是問題了.
我用一個例子.
我們在福音堂教會.
相當強烈的一段長時間.
我一信主就已經進入這個系統.
那個系統就告訴我們.
基督徒絕絕對對不准收利是,派利是.
你一信耶穌.
有一個很大的改變.
教會主流的教導.
又看到每個人都不派利是,不說.
一到過年就立即說.
派利是是迷信的.
是會涉及可能交軌的.
有些典故的.
有些原因的.
所以就這樣了.
慢慢變成了一種很嚴重的禁忌.
慘了.
初二回家.
親人都是傳道.
那怎麼辦呢?.
站硬不收利是,不派利是.
那些老人家的親人就很生氣.
後來私底下接充.

$^{641}$私底下就派給明仔.
派給其他的表兄弟妹,表妹.
就不敢告訴父母.
因為是傳道人.
這樣來暗暗底.
接著我們那輩子.
有一樣的子女.
走到過來.
這兩個家庭又不派利是.
很尷尬.
後來慢慢慢慢輾轉.
到我就派.
另外一個家庭又不派.
又帶來一個分歧.
為什麼你們今天有一個親戚有問題.
為什麼你又是傳道人.
他又是傳道人.
為什麼你又可以派.
他又不可以派呢?.
後來我們好像低調地.
大家交換禮物就算了.
可以想像得到.
原來單是在基督教裡面.
都有這樣的分歧.
我說的為什麼會有這個所謂的典故呢?.
因為很多年前.
我們的長輩.
福音堂很出名.
很資深的牧者.
他說他五十年代的時候.
去趕鬼.
趕得很辛苦.
趕了.
所謂辛苦是很難.
那些鬼很兇.
趕了.
接著沒多久.
那個被鬼又出事了.
就一直追問原因.
原來他家裡有一張紅紙.

$^{681}$是曾經不知道是不是那些符號.
所以就因為這樣的緣故.
就有一個結論.
就是說.
原來這些都可能跟教鬼有關的.
都要避忌小心.
慢慢就變成了.
好像一個耳邊蓋全那樣.
凡是過年節日的這些.
紅色的.
例如是年獸.
這樣掛勾.
就一定有迷信的了.
所以基督徒我們絕對不准.
是我們的教會.
其他教會從來沒有聽過這些東西的.
根本很正常.
甚至我們錦書院.
我們曾經印利是封.
叫弟兄姊妹來買.
有金句的.
所以.
你們一向原來.
在這樣的不同的良心.
那個信心程度.
都會有這些的分歧.
而這個分歧.
那個分歧本身.
就可以帶來大家的分化.
大家的不同心.
所以原來這個才是更加需要小心.
是保羅要處理的.
保羅不是叫我們處理了.
就是那些要監食的.
不准吃那些猶太人認為污穢的食物.
就分開兩邊.
就不是的.
所以你可以想像得到.
這一張聖經.
其實如果我們能夠拿準那個原則.

$^{721}$我們又願意學習.
其實是非常之寶貴的.
我們稍微再讀下面這一段.
可能我們先讀到第1到第6節.
有人信.
開始講那個例子.
有人信.
他沒有說哪一個.
有人信百物都可吃.
就是外邦人.
但願若的.
願若是那個良心.
不是他真的信心的願若.
但願若的.
只吃蔬菜.
吃的人.
吃的人就是萬物都可以吃的.
吃的人就不可以輕看那些不吃的人.
不吃的人又不可以論斷.
輕看和論斷.
那種相近的心態.
和剛剛的對立就是.
第一次講的接納.
所以接納和論斷.
剛剛是一個對立面.
是一個對比.
所以我們要怎樣接納到.
跟我們有不同做法的人呢?.
因為神已經收納他了.
不可以去論斷那些萬物都可以吃的人.
信念認為什麼都可以吃的.
神就已經收納了他.
第四次.
「理是順,竟論斷別人的僕人,.
他或站住或跌倒,.
自有他的主人在,.
而且他也必要站住,.
因為主能使他站住.」.
有人看.
另一個例子.

$^{761}$有人看「這日比那日強」.
有人看「日日都是一樣」.
似乎暗示的是猶太人.
他們一到星期五的黃昏六點.
就進入安息日的狀態.
去準備自己去守安息日.
每樣事情都不去做.
希望整個生命的日子裡面.
分別為勝出來給神.
這就是對某一個日子特別的看重.
但是外邦人沒有這些背景.
沒有這些習俗的時候.
他們覺得每天都是一樣的.
但其實慢慢地.
因為他們各種的原因.
到了星期天.
就一樣變成了一個重要的日子.
等於我們基督徒一樣.
所以你會看到.
這裡提到的是.
另一個例子.
剛才說了.
第五次.
有人看「日日都是一樣」.
只是各人心裡要意見堅定.
裡面自己的意志.
裡面自己的良心要很堅定.
為什麼我看這一天重要.
而我沒有強迫別人.
跟我一樣要看這一天重要.
這樣來到各人看自己.
心裡面的意見堅定.
有自己很清楚的.
為什麼要這樣去做.
接著是第六次.
「守日的人是為主守的.
吃的人是為主吃的.
不吃的人是為主不吃的」.
就是說不守日的人.
同樣都是為主不守的.

$^{801}$就是說原來各人的良心.
在做各自對神的敬虔的操作.
這樣的時候.
如果不影響其他人的狀況.
其實就是每個人要自己.
在論述良心.
每個人都要堅定自己的意向.
不是為了討別人的喜悅.
而是為了討神的喜悅.
慢慢隨著那個聚集.
而帶來彼此的同心.
甚至胡子們這裡的分析.
這裡的小結論.
就是說各自做自己的事情.
但是要同心.
要接納.
這是大的公約數.
就是說就算你堅持守自己的日子.
你都能夠要接納,尊重.
以及跟那些不守日子的基督徒.
在更大的福音裡面是合一的.
這是一個很大的邏輯.
很大的價值觀和方向.
原來可以帶來很大的大同和諧.
很重要的一個概念就是.
跟十三章上人所說的.
黑夜以心的警醒有關聯.
為什麼我們黑夜以心要警醒呢?.
因為將來要有審判.
因為將來神要審問我們人生.
所以保祿在第七節開始.
真的引進了審判.
第七節.
我們每一個人為自己活.
就是說守不守日子.
每個人的良心不同.
最終守不守日子.
吃不吃哪些食物.
最終每個人都要去到神面前.
要經過神的審判.

$^{841}$第七節.
我們每一個人為自己活.
因為我們是活祭.
也沒有一個人為自己死.
我們若活著.
是為主而活.
若死了.
無論是壞處.
或者人生的時間到.
譬如病.
意外.
我們都是為主而死.
保祿這樣說.
所以原來保祿看的生死和我們是不同的.
因為人類最後的防線就是自己的命.
最重要的.
吃什麼穿什麼.
最重要是為自己的命.
但是保祿的眼中.
他看我們整個基督徒人生.
原來活是要為主而活的.
就是說如果尼祿不是神的時間要我死.
我也是不會死的.
我也是.
雖然我交在你的手中.
但是如果神的時間還沒到.
我也是不會死的.
這樣的不是一種押響精神.
而是保祿看清楚.
真正的掌權者.
真正的審判者.
是我們的主.
是我們的神.
所以留意第八節的總結.
所以我們說活或死.
總是主的人是屬於主的.
如果你想想.
我們將自己看為一個物件.
就是說.
我的手機是屬於我.

$^{881}$我買了它之後.
那個手機就是我.
我要怎麼用它.
我用到它出神入化也可以.
我將它隨便放下.
不用它也可以.
用到它沒電.
用到它壞掉.
也可以.
但是一個有智慧的人.
一定不會將手機隨隨便便.
明明是智能手機.
明明拍照很漂亮的.
完全不用這些功能.
只是用來打電話.
沒理由這樣的.
所以當我們這樣將自己看為一個物件的時候.
那個主人就有絕對的權.
我們知道.
我們活或死.
總是主的人.
這個比起我們論斷人吃不吃的食物.
守不守的日子.
只是兩個其中的例子而已.
還沒說其他的東西.
人家現在說的就是.
又穿什麼衣服.
又穿什麼顏色.
又這樣配襯.
又這樣等等.
更多更多的平頭等足.
教會裡面的分化.
我們完全沒有想到一個大前提.
我們或活或死.
總是主的人.
就是說我們每個人應該很有意識地.
看自己是光明的兒女.
我們是披戴著福音的見證.
就是耶穌基督.
然後我們的價值觀已經清清楚楚.

$^{921}$我是成為神的器皿了.
或活或死.
總是主的人.
這樣看自己的時候.
其實你有這個淚鏡.
再看其他人.
其實已經不同了.
為什麼不同了呢?.
就是不敢再像以前那樣.
容易的批評.
容易的論斷.
因為我們就是主的人.
所以根本就是九字就是了.
因此.
基督死了又活了.
為要作死人並活人的主.
你這個人.
為什麼論斷弟兄呢?.
本來真的好像來了.
說到我們每個人真正的歸屬是神.
其實是神論斷我們.
批評我們.
當然一定不會.
但是神是清楚我們的生命.
所以我們根本就不會得到神的批評.
因為祂根本是全知的.
我們就是要面對神.
所以當我們這樣面對神的時候.
我們還怎麼會敢隨便.
去對其他人有過多的評論呢?.
有過多的評論帶來的負面的就是論斷呢?.
本來就很重的說話.
好像已經假想羅馬教會已經出現了這樣的問題.
這樣來到很重的一種責備的口吻.
你這個人.
就是說如果你扮演著神的角色.
來到去對別人.
有一種你認為你的判斷才是最準確的.
其他人都是錯的.
這樣你才會論斷的.

$^{961}$為什麼論斷弟兄.
又為什麼興翰弟兄呢?.
又是似乎是孿生兄弟了.
論斷帶來的不是欣賞.
不是尊敬.
不是來接納.
論斷帶來就變成興翰.
很小心.
這些有時候比義專的真理的錯謬更嚴重.
有時候這些是帶來的心靈的損害.
削弱了心靈的原因.
比起那件事本身更加的嚴重.
所以告訴我們.
為什麼興翰弟兄呢?.
因我們都要站在神的台前.
經常寫著.
主說我憑著我的永生起誓.
萬生必向我跪拜.
萬口必向我承認.
這是《自塞亞書》的45章.
他說這樣看來.
一個結論.
就是一到第六節.
教會內部的基督徒聚集.
我們一群光明的兒女.
我們這個社群.
外面是羅馬的政權.
是未信的人.
但是裡面是一群屬神的.
成為神兒女的活祭的基督徒.
我們在教會的生活.
保羅不是說我們全副恩就行了.
什麼都不用管.
而是都要管一下我們大家的心靈.
我們的心靈要不斷讓真理來光照.
就是神要再回來.
更重要的是.
我們準備好自己進入那個光明之前.
我們是要警醒對外.
對那些活在黑夜的人.

$^{1001}$我們要有一個光明的生活.
但是原來還有另外一個邏輯.
一個價值觀就是.
我們是要成為一個接納的群體.
就是對別人和我們不相同的.
我們要有一個接納.
因為重點是.
「失淵在我 審判在處」.
這個很重要的概念.
這裡就是告訴我們.
所以就是第十二節.
這個結論就是相當重要.
「這樣看來.
我們各人必要將自己的事.
在神的面前說明」.
就是說.
就是黑夜以心的概念.
第十四章.
在教會外面的人.
看我們是一群怎樣的基督徒.
我們是一群等候永恆的基督徒.
我們是一群將來要面對神審判的基督徒.
所以我們不能夠再白晝熱落.
好像醉生夢死.
我們是一群有意象.
是有使命的基督徒.
活在世界.
我們是要傳揚福音.
用不同的崗位.
用不同的形式去承載福音.
但是原來教內我們的生活.
我們是一群有意識的人.
我們是一群活祭.
我們是一群已經屬於耶穌的人.
怎樣去展現呢?.
就是對不同的文化.
不同的傳統.
不同的信心的人.
不同的他們本身的文化不同.
而帶來的那個接納.

$^{1041}$這個接納.
是因為我們都是神的兒女.
我們都要面對神將來的.
我們要向祂交代.
這是第十節所說的.
因為我們都要站在神的台前.
什麼叫站在神的台前呢?.
很簡單.
就是要面對審判.
《火羅人聖經》不如我們先揭開.
《宴上書》有45章.
不是叫你站在台前.
是站在基督的審判台前.
就是說將來是要向祂交代.
這是第十三節.
對不起是第二十三節.
我們就看後一點.
我們從第二十節開始看.
45章.
你們從滅國逃脫的人.
似乎預指到將來的被擄歸回.
若同罪疾前來.
那些抬著雕刻木偶.
以賽亞書的45章.
那些抬著雕刻木偶.
禱告不能救人之神的毫無知識.
你們要述說陳明你們的理.
讓他們彼此相議.
誰從古時指明.
誰從上古述說.
不是我耶和華莫.
似乎再一次和猶太人做一個辯論.
再一次將神的傳能.
那個從古的自有永有.
和地上人製造出來的偶像.
做一個辯論做一個比較.
除了我之外再沒有神了.
我是公義的神又是救主.
除了我之外再沒有別神.
地極的人都當仰望我.

$^{1081}$就必得救.
因為我是神.
再沒有別神.
我指著自己起誓.
我口所出的話是憑公義.
並不反悔.
就是這裡開始保羅引述.
萬失必向我跪拜.
萬口必憑我起誓.
人論我說.
公義能力唯獨在乎耶和華.
人都必歸向他.
凡向他發怒的必自蒙羞.
以色列的後裔都必因耶和華得稱為義.
必要誇耀.
意思就相當清楚了.
做物主的偉大.
全能.
自有永有.
他是唯一的創造主.
再沒有第二位了.
所以你拿偶像來跟神做比較.
簡直是侮辱,褻瀆神的.
當時的猶太人.
就是這樣在迦南製造偶像.
所以神就透過以賽亞來譴責.
跟他們辯論.
神是憑公義來建立他自己的威信.
叫做他的存在.
但是將來再沒有這樣的分別了.
只有做物主的存在.
所以到將來.
萬失要在他面前跪拜.
萬口要來承認.
當然就是他的公義.
他的聖潔.
他的審判的權柄.
所以保羅就很精準地將這段45章.
以賽亞書的經文.
歸納成為了末世.

$^{1121}$我們是神的兒女.
我們有一個小心.
因為將來我們是要面對.
基督的台前審判.
所以保羅也很特別.
將以賽亞書舊約的這種概念.
去到新約的時候.
保羅知道將來有一個審判.
就不再是耶和華了.
是基督台前.
好像是一個新的概念.
是耶穌基督坐在寶座上審判.
基督的台前接受一生所做的.
來面對神的光照.
到底是怎樣的情景呢?.
聖經就不詳細再講了.
給了一個概念就是.
將來有一天我們要交代.
其實根本不需要.
做民主是全知的.
我們根本去到神的面前.
根本不需要續約.
由哪一天.
1989年你做過什麼.
不可能這樣來講.
因為古往今來這麼多的人.
怎麼可能會這樣呢?.
所以在這個概念.
就是說其實我們每個人都要向神交仗.
我們每個人的一生.
才是真正有神去審問.
所以其實這裡想帶出的概念就是.
地上我們不要扮演審判者.
地上一些好像這麼小的事情.
大家的做法不相同.
用尊重,用接納,用欣賞.
用一個可能大家未必真的相同.
不相同之下又沒有抵觸信仰.
那當然不是核心信仰.
是一般的不同做法.

$^{1161}$帶著一份這樣的給錢.
互相尊重,接納的情況下.
因為更大的前提.
是我們要在基督的台前.
彰顯,顯露我們自己的工程.
這裡才是重點所在.
所以我們可以看看經文.
剛剛也說了.
十三章一直都說的一個很重要的.
大的總綱.
就是一切以愛取締了那個律法主義.
取締了我們在新的活濟生活下.
我們是以愛,愛神,愛人.
做一個很重要的.
我們整個生命的依歸.
對人有益,有善的.
不要全心與人為仇為敵的.
我們是要愛來作為我們的總綱領的.
所以下面那段剛才也說了.
就是面對黑夜以深.
我們就是保羅所強調.
我們是活在光明中的群體.
因為羅馬政權是看著基督徒的.
真的會看著的.
所以越是光明就趁虛.
反照他們背離了真神的敬拜.
一定有兩個結果.
一是他們信著.
果然是有的.
但也是同樣的.
就是尼祿和以後的君王.
展開所謂十大逼迫時期的羅馬的政權.
就是對基督徒無盡的逼迫.
直到君士坎丁皇帝信主.
那個逼迫才真正的結束,停止.
所以古羅馬有很多的地道.
就是逼迫帶來的產物.
就是要在地下崇拜.
在不同的一些地道裡面.
所以這個也是我們要小心的.

$^{1201}$光明的生活就是一群等候主的基督徒.
警醒就是主在內.
以及要活在光明中的生活.
就是讓未信的教會外面的世界.
他們是有一個見證.
他們需要來悔改.
接著回到教會內部的議題.
就是一些跟信仰不是核心有關的.
真理有關的.
就要來學習一個接納,欣賞.
剛剛上個星期在紅恩.
教了三堂基督教義端.
有守望堂 也就是耶和華見證人.
也都加插了有些人想聽.
聖靈的靈因運動第四波.
最後有講了摩門教.
因為也要做一些功課.
摩門教他們的網頁.
有時候真的很搞笑.
義端的網頁特別漂亮.
特別的工整.
很多的條理.
找一找都不用找得很辛苦.
他自己自問自答.
他說摩門教到底是不是義端呢?.
他在網頁這樣來解答.
不是義端你怎麼解答呢?.
你不用感受堂的網頁上.
感受堂是不是義端呢?.
沒有這些東西的.
他真的這樣回答.
他說有七點.
摩門教不是義端.
你今晚有時間大家上去看看.
是笑死你的.
他不是義端.
他不是把教義跟他辯論.
他說我們摩門教徒是好人.
我們納粹的.
我們什麼的.

$^{1241}$很搞笑.
完全避開了那個教義.
你們有什麼錯謬.
為什麼你們不改變你們的錯謬.
他不回答這些.
所以這些就是核心的教義.
核心的教理.
譬如摩門教很簡單.
耶和華曾經都是由人慢慢演變.
就是人變神.
耶穌真的人.
不過就是神跟一個天使生的耶穌.
他是一個受造的.
所以根本上整個的摩門教.
名字叫做摩門教.
實際上是完全敵基督的.
是非常敵視真理的.
跟聖經是剛剛南轅北轍的.
是徹徹底底的一個很褻瀆神的義端.
所以你可以想像得到.
義端就是這些核心的教義.
出了很嚴重的問題.
不肯改變.
不肯悔改.
就是被教會定性為義端的問題.
所以現在保羅探討的十四章.
完全不是這些核心的教理.
不是這些核心的信仰.
而是一些很簡單的.
大家做法上的一些不相同.
所謂這些不相同的分歧.
這些我們已經講過了.
這是十三章的第八節.
愛是律法的總規.
十四章同樣都是在講.
為什麼我們需要去愛心.
去對待每個人.
尤其是每個人的個人權利.
所以有時候那個論斷.
就是將個人的權利放得太大.

$^{1281}$這個都是要很小心.
所以十四章總括來說.
在這個段落.
一到第十二節.
是在講怎樣去對於我們的自由.
我們要來約束.
我們個人的良心.
要怎樣去將它培養.
在沒有真理底下.
越來越和其他人達致一個越來越有和諧.
有合一這樣的生活的形態.
就是我們在我們自己的自由和良心底下.
那種比例一路一路減少.
不要無限放大自己的自由.
變成為個人的權利.
就好像變成了.
肆無忌憚的在教會製造了.
變成一個帶來的異見.
帶來的一種衝突.
這樣就不行了.
所以原來.
其實每個人都應該有自己的自由.
有自己的良心.
有自己的權利.
但是仍然是那個重點.
十四章.
在保羅的教會內部的大家的生活.
其實就是在運用十二章.
十三章.
這個愛的大總綱.
作為在那個應用上的例子.
就是很實實在在.
保羅就將這個愛.
怎樣應用在教會的場景呢?.
因為一直都在說.
教會是承載著真理和福音的.
神是透過不同的教會.
或者一群聚集的基督徒.
去傳揚耶穌的福音.
所以教會最大的必須是愛心.

$^{1321}$大於個人的權利.
大於我們的自由.
大於我們的良心.
所以在這裡怎樣能夠運用呢?.
原來就是必須要學習接納別人的差異.
必須.
記得從頭到尾.
不是在說那些核心的信仰.
那些是不能討價還價的.
必須是就是不是就不是的了.
但是對於不同的做法.
真的有不同.
所以現在我派利是收利是.
我是完全沒有良心的不安的.
早期還有的.
後來就不需要了.
因為那個原來是一些不同的做法.
不同的傳統.
甚至中國人.
香港人.
根本已經不再看那張紅字.
是不是吉利的了.
當然了.
那些老人家真的還有的.
你過年.
你用一個白的信封包著來派.
那些老人家就真的不肯過你.
真的會當面罵你.
還有這些小小的觀念.
但是那裡都還沒有去到一個地步.
說真的和交軌,迷信.
拉上一個很大的關係.
所以純粹在一個習俗上.
那個又不需要將它放大到一個無限這麼大.
當然不代表我們是不是年初一.
基督徒就可以去廟裡拜.
習俗而已.
那又是另外一種的情況了.
那又是另外一個邏輯了.
所以我們記得在14章說的.

$^{1361}$這些完全是非核心的信仰的一些做法.
就是說可以斟酌的.
有些人這樣做.
是跟罪,跟良心沒有任何的抵觸的.
我們要留心.
就是說我們今天很少人再說.
安生日為什麼你不守啊.
或者某些食物.
為什麼曾牧師你還吃豬紅啊.
又說不可以吃血.
類似這些.
摩門教守望台應該是.
守望台禁止吃血.
而且不單止.
去到一個地步就是.
血裡面有生命.
禁止輸血.
我不知道香港是不是有.
我要先問一下醫生.
美國是有一些儀器.
幫你自己的血循環.
是不可以接受外面的血的.
如果你是守望台的信徒.
你要告訴醫生我是守望台的信徒.
你不准幫我輸血.
我死是我的事.
你不准幫我輸血.
是真有其人.
真有其事的.
是很慘的.
是普通的病.
流血不止的死.
因為不准輸血.
所以那些就是.
關乎真理.
就不是一些非核心.
當然了.
基督徒肯定可以吃豬紅.
因為已經不是我們的文化.
不是我們的習俗.

$^{1401}$連那件事本身當時.
保羅.
對不起.
彼得他們所說的.
都不是因為真理.
都完全和真理無關.
所以我們要記住.
不要將一些我們先入為主的東西.
畢竟異端就是將它放大.
就變成害人不淺.
但是我們自己也要對真理有一個全面的貫通.
那個形式.
這樣來融會貫通之後.
我們就會對真理有越來越真確的了解.
所以這就告訴我們.
下面說的就是第十三節開始.
不可以叫弟兄跌倒.
所以我們就一直在看這個.
告訴我們我們活在這個世界的盼望.
是了.
很重要了.
保羅大進的短後.
那個永生.
那個末世.
那個末世就是說.
將來.
末世一定是和一個掛鉤.
就是將來要面對主耶穌的審判.
所以就叫基督台前的審判.
是了.
所以他就將45章的《以賽亞書》.
這個曼辛要向神跪拜.
曼侯要向祂承認.
要向祂承認祂是神.
原來有跟著的審判的.
就是說你要將你的安寢.
讓審判者來交代.
所以所謂曼辛跪拜.
曼侯承認.
就是十二節.

$^{1441}$有首尾跟的.
就是各人必要將自己的事.
在神的面前說明.
所以保羅就將《以賽亞書》45章.
做一個他很清楚的解釋.
原來曼辛,曼侯要承認.
要來尊崇神.
原來就是我們的人生.
如何為主而活.
如何用神給我們的人生.
所以我們就.
原來十三章的《黑夜耳塞》.
十四章就應用在一個實際的生活的場景.
就是我們在這麼有限的基督徒的教會生活.
我們的人生.
首要要學習.
對別人接納,尊重,欣賞.
合一和同心為福音.
每個人都有不同的價值觀分歧.
你叫我穿紅色.
穿粉紅色.
不如你殺了我吧.
我家是不會有粉紅色的恤衫.
你叫我穿.
就是違反了我整個的良心.
這件事不是罪.
也不是說我絕對不准不做.
但是可能隨著某些成長.
又有些做法不同.
都不出奇的.
所以那些只是個人的喜好.
個人的習慣.
個人的良心.
所以每個人彼此都有一種的欣賞接納.
其實人與人之間發揮的愛的總綱.
愛是不加害於人.
所以愛就是完全了律法.
這個十三章第八節的經文所說.
所以這裡很重要的就是.
保羅原來一直都沒有離開愛這個總綱.

$^{1481}$基督徒的群體.
為什麼愛人都知道我們是一群有愛的頂子妹呢?.
所以其實教會裡面.
更加應該要實踐這個互相接納.
欣賞的這個情操.
這樣的價值觀.
所以保羅也用了很負面的話.
正面就是我們要彼此接納.
大家為福音去效力.
大家真的不要用一些分歧,看法等等.
大前提不要將那些精力.
消耗在這些所謂的內耗上.
反而就是我們大家都要來面對這個世界要結束了.
因為黑夜已深,白晝將近.
保羅告訴我們.
保羅說完這句話已經是二千年後了.
我們主耶穌回來.
是不是一定是再近一點呢?.
肯定是再近過保羅的時候.
保羅也說黑夜已深,白晝將近的時候.
我們就更加需要警醒.
我們自己也需要在我們這麼有限的.
實定我們人生的時間.
匆匆的幾十年過去了.
我們就更加應該要珍惜.
世人都知道每年的倒數.
就是盡快的過去就讓它過去了.
很多人大家都很想展望新的一年會好.
這個心情是對的.
人類的與生俱來的那個生命來的.
所以學無生開始是很好的.
如果沒有的話.
人就根本不可能再生活在世界了.
因為如果這麼悲哀.
每天都活在陰影的話.
就真的死定了.
但是世人只是停留在希望可以好一點.
但是基督徒其實很實實在在的.
保羅說因為我們必要站在基督台前.
將自己的事在神的面前說明.

$^{1521}$這個是肯定會發生的.
這個是我們的盼望.
這個是讓我們警覺.
不要扮演一些我們不知道的事情的救世主.
和那個審判者.
只要我們強烈有這個意識.
我們就不斷會對別人有一種接納.
有一種的欣賞.
至少在接納和欣賞之前.
我們那個問號都不是一個審判的問號.
而是一個可能我不知道他發生了什麼事.
所以我先為對方禱告.
有機會就去問一下.
為什麼你會這樣呢?.
這樣來帶出一種的.
原來是這樣啊.
那就對他有更大的接納.
更大的欣賞.
或者他有他的困難.
就更多的包容.
更多的對他的接納和幫助.
所以原來等候處這個基督台前.
原來是這麼正面的去裝備自己.
不要成為審判者的那個原因.
是讓我們真的看自己是一個.
很有血有肉軟弱的一個人.
其實真的.
我現在.
昨天他們一邊在行山.
還要走鳳凰山.
想我死了當然不去了.
走完回來第二天整條腿抽筋.
因為跨過這麼高的地方.
但是明知道真的沒有以前那種的體力了.
就是拍回當年的照片.
我說當年我也去過這裡拍過.
是當年.
十幾年前.
所以現在是不行的.
隨著人生.

$^{1561}$我們的體力不斷下降的.
我們的精神都不斷下降的.
但是我們在基督裡面的主在來.
和我們要去到基督台前.
應該是越來越清晰.
好像那個焦點越來越強烈.
所以我和我自己初信的時候.
基督徒的那個心態.
經過了二十多年的事奉.
就更加很清楚自己.
真的不會再有一種隨便.
隨意的批判.
論斷.
扮演神的那種審判.
反而多了就是對別人的先諒解.
包容.
體諒.
還有希望能夠明白.
尋求那個形式.
背後的一些原因.
價值觀.
這個就是相當寶貴的.
所以原來我們每一個弟兄姊妹.
必須學習.
寶來真的很好.
他很意識到.
講到多偉大.
有一個很小的福音.
萬民要信主.
但是如果一幫弟兄姊妹.
羅馬的弟兄姊妹.
如果都是四分五裂.
貓合神離.
怎樣有力去傳福音呢?.
要這樣做.
這樣做不行.
這樣做.
要這樣做.
那你就真的.
只是為了討論怎樣做.

$^{1601}$已經是蹉跎了歲月時間.
相反就是.
如果大家都是帶著一種彼此的信脈.
彼此的一種信心的默契.
彼此都已經有一種彼此的同心.
那就很不同.
有些東西是彼此的包容.
我記得.
我開名都不要了.
小玲有時候教的做法.
事奉根本就是彼此的包容.
突然間有些改變.
有什麼所謂呢?.
大家走位.
彼此保衛.
這些是很重要的價值觀.
就是事奉者大家都是同心.
不是看到有些東西突然間改變.
哇.
臉都黑了.
然後他在背後.
講是非.
又在論斷.
這樣做.
突然間改變.
其實不需要.
有時候真的有些東西突然間要改變.
就改變.
改變之後再檢討.
那個改變為什麼呢?.
解釋彼此在一種信仰下的尋求原因.
不是帶著一種本身是敵視的.
去尋找那個罪責.
誰要背負那個責任.
不是這樣的.
教會就是應該要這樣的同心.
所以不要以為14章不重要.
因為保羅的大前提是想告訴我們.
在等候盼望.
黑暗而深的時候.

$^{1641}$我們更多的.
更少的.
是面對邪惡的政權的迫迫.
就是羅馬教會.
其實更多的.
就是在地下大家同心的生活.
所以有時候困難.
會使我們更加凝聚.
太過安殊.
會使大家更加容易分歧.
這個可能是一個普遍的道理.
通常那些爸爸一生做完很多錢.
有很多遺產.
第二代就不理那麼多了.
這份遺囑那份遺囑.
這樣這樣又來搶那些家產.
有時候就是在這樣的情況下.
親兄弟姐妹都會這樣四分五裂.
所以教會14章.
都仍然為什麼要接納.
其實不單純是審判懼怕.
而是愛的正式實踐.
在教會裡面.
如果我們本身以愛做我們的準綱.
根本上我們對別人首先有一種信任.
接納.
在這個信任過程.
我們才尋求大家的分歧.
分歧背後有沒有什麼原因.
然後明白了.
將那個原因慢慢縮小.
更加大家一起學習.
在主裡面同心的服侍.
所以這一段正經其實很寶貴的.
對外的確我們面對的是未信的.
無論社會,國家,未信的人.
我們用什麼態度呢?.
肯定就是我們要光明的兒女.
背負著,披戴的就是耶穌基督.
當然不是穿著T恤.

$^{1681}$不斷要穿著T恤就是耶穌基督.
不是這個意思.
而是我們整個生命走出來的光明.
但是對內就是教會的群體.
我們就應該在這個末世盼望裡面.
彼此接納.
愛就是核心.
然後在這個末世裡面.
我們就在盼望裡面警醒.
因為每個人真正的主是基督.
我們在基督的台前來接受面對他的審判.
非核心信仰.
很清楚了.
保羅就用一個比較大的方向.
就是猶太人的群體.
和外邦人的群體在羅馬的教會.
就是說一定有其他的群體.
不過以這兩個作為一個明顯的主流的群體.
就會有一個很容易在這些做法上的分別.
那個差異應該說.
猶太人就是守安息日.
信了主之後.
還沒定奪一定要禮拜守敬拜神.
所以他們如果認為星期六.
來去敬拜神.
來去敬拜耶穌.
一樣是對的.
所以外邦人根本沒有這個安息日的觀念.
和這樣的包袱的時候.
外邦人只會因應他們的做法.
來去這樣做.
一定是受當時的社會的思潮.
例如君士坦丁信了主之後.
他那時候拜開太陽神.
12月25日.
24日,25日.
他就將這個拜太陽神的日子.
改變成為聖誕節等等.
這些一定是受當時的社會的思潮.
但是在這些不是重要的核心的信仰.

$^{1721}$保羅在傳道的過程.
其實甚至乎是他自己.
叫猶太人信主.
帶外邦人信主.
在喬治亞諾萊加.
在肥牙比.
根本就是猶太人和外邦人.
已經開始在一起聚會.
所以保羅是非常意識到.
這些差異只會去到一個階段.
一定會產生一些張力.
大家會擱持己見.
擱持自己的良心.
而帶來慢慢將這些表面和核心信仰.
沒有直接關係的.
慢慢不知不覺.
就會變成好像核心信仰的價值觀.
這樣去競爭.
這樣去做一個對立.
就帶來教會出現的困難.
所以原來羅馬書.
不是針對哪個教會的錯謬.
異端或者產生的內部的衝突.
反而在沒有這些事情之前.
保羅就列舉了.
我們這個信仰群體.
是來自五湖四海.
來自不同的民族.
必然會產生這個差異.
這個差異根本就是神容許存在的.
在這個容許存在之下.
怎樣辦理的人.
帶來的人.
怎樣讓這個差異成為一個祝福.
不要成為一個對立.
這個就是教會的領袖必須不斷去思考.
怎樣將它帶來和諧相處.
例如聖樂.
這些幾十年都是這樣討論的.
聖樂.

$^{1761}$傳統的聖樂.
和現代的敬拜.
就是用樂隊打樂隊.
唱現代的詩歌.
兩者整個形式.
整個承載都不相同.
怎樣能夠讓教會和諧共處呢?.
紅音還多了福音樂曲.
你可以想像到.
怎樣讓差異能夠和諧共處呢?.
甚至要容許它和諧共處.
要容許它存在在一個教會裡面.
你說不同的教會.
你喜歡這樣就這樣.
你這樣就這樣.
這樣就沒問題了.
但是.
本來的意思不是這樣.
反而就是讓差異縮小到一個地步.
大家不覺得是差異.
反而是各自有一種的欣賞.
各自的精彩.
各自去承載基督的福音.
用它能夠最好準備的狀態.
的那個服侍.
那個獻情.
不是叫演出.
叫獻情.
給上主.
帶領會眾.
用那個方式的敬拜.
例如用福音樂曲.
去承載那個福音.
這樣去使神得到榮耀.
在聚會裡面.
所以.
你可以想像到.
在這樣的情況下.
這是一個真正核心的問題.
就是.

$^{1801}$怎樣帶來大家和諧的.
共同擁有這個福音.
一起在大家.
各自可以繼續堅持的差異裡面.
來服侍.
特別是紅恩.
去年.
凡是第五週.
就是太羽會獻情他們的詩歌.
凡是第三週.
就是紅恩的福音樂曲去獻情.
有些大的節日.
我又繼續叫福音樂曲隊去獻情.
大家各自去承載那個音樂.
弟兄姊妹又覺得很好.
我們去年試過一次.
就是在報道會那天.
福音主人那天.
一般就是著詩棚的帶領會眾唱歌.
我就刻意叫幾位年輕的弟兄姊妹.
用現代的詩歌.
當然因為他們不懂彈結他.
所以就不會用結他.
刻意不准著詩棚.
要用現代的穿襯衫,西褲就行了.
等等.
這樣來帶領會眾唱詩.
就是想引進那個差異.
重點是大家在愛裡面.
不是排擠其他的聖樂.
不是排擠其他的音樂.
你那些不行的福音樂曲不行的.
一定要傳統的.
這個我就禁止了.
這個想法就錯了.
不應該這樣.
反而是大家接受大家的差異.
即不同.
然後大家盡力在崇拜裡面.
帶進那個敬拜.

$^{1841}$帶領會眾進入對神的敬拜.
這樣就希望.
將大家不同的地方.
反而凸顯他的寶貴.
他那種不同的優美.
這樣就會帶來.
大家都學習彼此接納.
所以如果教會甚麼都是一樣.
其實就不最好.
當然這裡說的更加大.
就是文化上.
這些想改都不是一時三刻就可以改.
所以求神幫助我們.
帶領的人要有很強烈的意識.
一定人和人之間有差異.
有不同的看法.
連穿的衣服都有不同的顏色.
那些差異其實很小事.
不要譁眾取寵.
不要堅持自己的自由.
自己的權利.
就行了.
教會要有意識地.
堅持愛是那個總綱.
以及在等候永恆裡面.
警醒.
因為有一天要面對主.
基督臺前的審問.
那個叫.
可能說得很好的.
將自己的事.
在神的面前說明.
不要說說明.
不是一個審判.
但說明的背後.
其實就是你一生.
你有多少愛的總綱的人生.
你有多少對別人.
欣賞,接納,包容,體諒.
不同意見都可以是體諒.

$^{1881}$不是批判.
所以盼望教會.
我們成為一個這樣的群體.
今天就說到這裡.
因為第十三節我還沒準備.
下個星期我們就再進入.
十三節那段.
我們一起祈禱.
多謝天父.
十四章真的很寶貴.
因為是將愛的總綱.
超過了律法的人生的生活.
來到實際的場景.
十四章在教會裡面.
和一些不相同的.
民主的民主的民主人士.
和一些不相同的文化背景的人.
一起生活的.
可以說是一個基督徒的藝術.
一個待人接物.
原來守日.
堅持守日.
原來是可以榮耀神.
不守日又可以榮耀神.
我們很盼望.
我們看到的就是做物主.
給我們人不同的性格.
不同的審美標準.
不同的做事的方法.
但是求主幫助我們.
能夠知道哪些才是真正我們要堅守的.
就是信仰.
就是絕對的恩信.
清義的救恩.
幫助我們能夠知道我們所信.
和有意識地.
活在主的榮光底下.
對別人.
我們代表著主.
其實是為了去愛.

$^{1921}$去接納.
去批判,審判.
願神幫助我們.
我們仰望你.
多謝你給我們這麼寶貴的寶囊.
能夠意識到教會有猶太人.
有外邦人.
有不同的國籍和諧共處.
最重要.
最重要都是要將.
這個福音傳到萬民當中.
求主幫助.
多謝你聽我們感恩禱告.
奉耶穌基督的名.
阿們.
\newpage



\section{羅馬書}
\label{sec:nB_6d0PSa6M}
\textbf{2026-01-09 周五晚查經班《羅馬書(33):分題:如何處理教會內「非核心信仰」的差異(二)》 講員:曾裔貴牧師}
\newline
\newline
連結: \href{https://youtube.com/watch?v=nB_6d0PSa6M}{\texttt{https://youtube.com/watch?v=nB\_6d0PSa6M}} ~~~~ 語音日期: 2026-01-10
\newline
\newline
\hyperref[sec:aePwbTZJwec]{\small{< < < PREV SERMON < < <}}
~
\hyperref[sec:index_chronic]{\small{[返順時目]}}
~
\hyperref[sec:index_scriptual]{\small{[返順卷目]}}
~
\hyperref[sec:pe8biigoai0]{\small{> > > NEXT SERMON > > >}}
\newline
\newline
$^{1}$晚安 各位網友 我們一起禱告.
再一次我們將整個的聚會仰望交託給主.
多謝你招聚我們 透過不同的方式.
我們可以在現場 我們可以在線上.
也可以之後透過YouTube的頻道.
我們可以聽一些錄音.
都給我們在信息裡面 我們有啟發.
特別是在第13,14章 一直到第15章的第一個段落.
都是保羅很想教會在待人接物.
彼此相交,相處上面.
要有基督徒的那個愛心.
我們很盼望主都帶領我們.
因為雖然時代不同 文化背景都不同.
我們都不是猶太的民族.
或者當時歐洲的這些不同的地區的人.
但是我們知道神的話語.
特別是那個大方向原則是不會改變的.
幫助我們每一個地區的教會都來去遵守.
都來去有一個透過啟發,反思.
而帶來的那個具體的落實,執行.
使到教會能夠與時俱進.
去承載不同弟兄姊妹的心靈.
來去認識主和同心的事故主.
我們一起在你面前等候禱告.
奉耶穌基督的名 阿們.
上一個星期我們都差不多講了.
第1到第12節的那個段落.
我們今天就可以再去思想.
或者我們從第13節這個段落開始.
到第23節 就是在這一章裡面.
是有兩個大的分段的.
在第一個分段我們就講過.
第1到第12節就要避免.
不要去壟斷教會裡面一些人所做的事.
如果那些事情不是與真理,.
救恩有關連的 直接的關聯的.
我們就要彼此的尊重,欣賞,接納.
接著去到第13節就剛剛一個對比.
一個負面的不要壟斷.
13到23節剛剛一個進程或者一個對比.

$^{41}$就是要怎樣去具體去挽回弟兄.
以及怎樣可以透過一些飲食.
透過一些文化的差異就可以接納到.
以致那個愛心能夠很具體的落實.
所以這兩段經文都是保羅很有意識地.
很想在一些很實際的.
包括可能保羅自己也知道.
會有這樣的文化差異.
他來舉這些例子很具體的.
就是發生在當時羅馬帝國意大利的教會.
因為有猶太人的聚集.
所以在這麼大的文化差距之下.
能夠這樣來到去有一個教導.
其實殊不簡單 可想而知保羅.
不只是講一些道理.
而是他很具體在每個星期的聚集的生活裡面.
怎樣能夠帶出一份愛的生命.
怎樣能夠彼此的接納.
以致能夠合一呢?.
所以這一章其實都相當精彩的.
不過這一章我們又不能夠抽離上下文.
因為整個的上下文比較遙遠的12章.
就是我們已經要肯定自己是一個活祭的生命.
是神用了很重的代價.
他的愛子為我們犧牲.
所以如果我們就這樣讀14章.
我們就好像覺得我們要在教會裡面做和事佬.
做好人 大家不要生涯不要口.
就變成好像人與人之間.
怎樣來制定一些的律例.
或者一些的守則 圖書館盡量不要吵.
就不是這樣 而是在更大的大前提.
我們對外固然要面對尼祿羅馬帝國.
所以我們在正常的情況下我們要學習馴服.
我們要來到去在那個實際的操作.
例如我們要納稅 我們要奉公守法.
我們要來到去做好這件事.
第14章就是教會是一班無恩的人.
我們都是神的活祭的生命.
奉獻了給神的一個祭身.

$^{81}$我們應該怎樣在我們人與人之間的相處.
明顯是有不同的信主和信心的程度.
都有不相同的時候.
看事物的角度 看事物的情況又有不同.
怎樣可以處理得好一些呢?.
所以這些就是很寶貴.
能夠在具體的生活裡面.
保羅能夠這麼認真去教.
而且教的過程還有好幾個大的原則.
如果我們一直這樣去遵從.
而我們在這樣的生命裡面去成長.
我們的生命就很成熟.
這個是神所喜悅 人所喜愛.
這是第14章最重要的那節聖經.
就是第17節.
因為神的國不在乎飲食.
當然更加不在乎你算哪一天好一些.
安息日好一些 禮拜日好一些.
復活節好一些 重陽節好一些.
不是這些 神的國不在乎這一切.
在乎什麼呢?.
三樣東西 公義 聖潔和在聖靈裡面的喜樂.
在這幾樣上面去服侍神.
服侍人就會得到神的喜悅 人的稱佩.
第14章第17節真的要大家背得很熟.
要很知道 但是今天我們講的第14章.
我們就必須將整個的上下文.
跟大家一起去看.
這樣就非常好.
所以我們怎樣看待彼此.
在本來就提這兩樣.
可能對當時已經是很大的一個大家的矛盾位置.
大家已經開始很顧守自己的日子.
就是安息日.
有些就認為禮拜日大家一起聚集是最好的.
兩者之間這些的不同.
怎樣帶來不同 任讓它不同.
但是這個不同裡面可以有合一.
這個就是一個很寶貴的教導.
所以焦點在第14章上半段不要論斷.

$^{121}$下半段13到23節簡單來說.
就是要很主動的.
用一個很有愛心的思考.
來去接納跟你相同的人所做的事情.
當然這裡講的焦點就是實際具體的操作.
本來就用了很多具體的教導.
我們一起再下去看我們今天這個.
第14章這一段經文.
我先移走這個 先不要阻止.
我們接納的原因.
首先就是要明白我們中間不同的人有差異.
而這個差異不是關乎舊恩.
不是關乎真理.
這樣的差異就容讓它存在.
因為文化和良心上的不相同.
這個不是關乎舊恩真理的問題.
所以這個不是要強迫別人.
必須跟隨你這樣做的.
一定有它的好 有它的不好.
盡量在整體裡面.
從一個整體來看這件事就OK的了.
其他人很具體的.
我們剛剛在星期三開完檢討會.
就是聖誕節我們報佳音.
今年我們知道預計會多人報佳音.
有60多人參加.
你想想如果60多人.
只是排隊都不夠排.
所以我們就分了幾個項目.
有一些是拿咪高峰的.
就是帶領和說話的.
接著有一部分就是負責跟著去唱聖詩的.
另外有一班十多人就編了.
就是跟那些街坊主動邀請他們過來一起唱.
和派禮物.
我們就開了檢討會.
有肢體就說.
你編了我的位置.
只是派禮物.
我很想去報佳音去唱的.

$^{161}$為什麼不行.
有這樣的意見好.
我們就明白.
其實不是的.
其實我們有幾個點.
從大發.
從隧道.
過來這一邊的輕鐵站.
再過去駱橋樓.
再過去幸福村.
其實有幾個點.
其實當初因為可能那個活動的時候.
不是那麼具體.
有些人就遲到了.
變成我們很多時候誤會了.
我好像編了你的位置.
你就不能動了.
你一定要站著.
你都不准唱歌的.
其實不是的.
這些就是一些差異.
這些可能透過一些誤會.
或者一些做法上.
因為教會一個大的活動.
總是有它各方面的考慮和安排的.
所以如果跟你那個想法有一點相同出入.
這些更加可以討論.
可以事前我們多來給意見.
事後我們大家一起檢討.
例如拿咪高峰那個.
講完祝福語.
我們預備唱什麼聖詩.
他就不可以大聲用咪高峰唱.
因為其他人的音被蓋了.
這是很好的建議和檢討.
所以原來我們尚且這些不斷需要有不同的檢討.
何況一些是已經存在了.
在他們未信主之前.
就已經出現的這些差異.
這些不相同.

$^{201}$本來就很具體了.
文化上.
因為猶太人真的要受安息日.
猶太人真的有那些不潔淨.
未放血的食物.
他們是不可以吃的.
這就衍生了一個神學的問題.
要怎麼讓猶太人慢慢明白呢?.
所以這就需要一個時間.
大家的信心程度都相同.
所以在做喜上.
大家一起相處的時候.
就需要有一個藝術.
就需要有一個接納.
就需要有一個包容.
所以這張聖經.
其實真的很寶貴地說到.
如果我們教會的弟兄姊妹.
大家都是神拯救的.
如果我們不含恩.
根本上怎麼能夠將福音帶到那些不信.
拒絕的人的生命裡面呢?.
所以這裡告訴我們.
最重要的.
愛和尊重.
就是勝過一切的爭辯.
有愛,有尊重.
才是真正可以帶來的造就.
對別人生命的造就.
所以我們一直看下去的時候.
畫面好像不太好.
十四章第十二節.
最後那一節說什麼呢?.
他說這樣看來.
我們各人必要將自己的事.
在神面前說明.
連我們尚且將來.
都要面對神的審判.
何況我們現在呢.
現在我們就是更加.

$^{241}$所以第十三節.
立即再補充.
所以我們不可以彼此論斷.
你可以想像得到.
第十二節就是一個總結.
總結一到十二節的經文.
告訴我們.
不要論斷.
最重要.
就是我們自己都要面對神的審判.
某程度上有一個判斷.
當然神就不是論斷人.
但是我們將來都要面對神的交仗.
元文就是這個詞.
我們的仗.
要一切的仗.
要向神呈交.
元文的意思.
因此我們每一個人.
都要向神交上自己的仗.
元文的直譯.
所以保羅就用了第十二節做一個總結.
這句話都相當重.
所以如果我們有意識地.
原諒我們自己都要面對神的審問.
所以我們今天呢.
我們真的沒有權.
扮演一個審判者.
來對不同的事物.
來做很多的批判.
批判背後呢.
如果不是出於一份的愛.
那就不是一個的善意.
更加不是一個可以帶來那個群體.
一個生命的躍進.
一個的更新.
所以十三節他講這個接納之前.
整段是講對別人的接納.
講這個接納之前.
十三節立即就再加多一句的結論.

$^{281}$就是重複第十一節那段經文這樣講的.
第十一節.
十和第十一節.
他說你這個人.
為什麼你論斷弟兄呢?.
為什麼你輕看弟兄呢?.
因為我們將來都要站在神的面前.
站在神的審判台前.
就是四十五章的二創大書.
就是引述將來神要來審判這些扮演偶像的人.
這個告訴我們.
保羅就很清楚引述.
原來將來在永恆.
是有一個的判斷的.
所以我們今天.
我們絕對不應該在地上有不同人之間.
這裡是在教會這個設定.
在大家一起聚會的一群人.
就是大家有關係的.
大家是弟兄姊妹.
大家是神的兒女.
再重一點.
大家都是生命的活祭.
原先是獻給了神.
是神的活祭.
我們就應該在我們實際的每天的生活裡面.
那個緊密的同在的時候.
要有生命的操守.
就是愛和尊重.
接納是我們成長的一個很重要的生命的指標.
十三節到第二十三節.
我先跟大家讀了這一段經文.
然後我們就講這一大段經文.
講到保羅教我們怎樣接納別人.
十三節開始.
我用舊的和盤.
所以我們不可再彼此論斷.
再彼此論斷.
保羅是不是聽到某些人講.
讓他知道羅馬教會有這些問題.

$^{321}$因為在保羅的意識裡面.
其實一直以來.
就算耶路撒冷教會.
安提摩教會.
不同的教會.
哥林多教會.
就說保羅還沒有去.
在這個時候.
他都應該意識到每個人總有一個多大的生命要學習.
要成長.
要對付自己.
就是那個論斷.
再論斷.
所以保羅這裡講的是一個repeat.
一個一而再重複的論斷.
我們已經有了這種慣性的道德.
慣性.
看一件事我們不對.
和我們不相同.
或者我們自己不是很滿意.
我們就會有我們的評斷.
我們有我們對那個人的論斷.
這就是一個多大需要.
大家都要更新學習.
所以我們不可再彼此論斷.
寧可定義誰也不給弟兄放下半腳跌人之物.
我憑著主耶穌確知心信.
在這裡解釋對真理.
對一些事物的更新的形式.
大家留意不是這麼簡單.
保羅在這裡說的是.
其實在主耶穌的救恩申約.
耶穌復活之後.
一切都更新了.
所以可以這樣說.
他說憑著主耶穌確知心信.
凡物本來沒有不潔淨.
唯獨人以為不潔淨.
在他這個人他一路有這個心結.
有這個良心的不安.

$^{361}$有這個對事物不能夠全面的透支.
形式的時候.
在他的心裡仍然是不潔.
這裡好像很具體的說到一個猶太人.
就算信了主.
他也真的覺得有不同的食物.
有不同的對他的玷污.
對他在神面前的身份的重要.
但是保羅告訴他們.
其實在主耶穌的救恩裡.
一切都更新了.
你成為神的兒女.
不是因為你吃了某些乾淨的食物.
即是放了血.
即是沒有帶血的食物.
或者你做過甚麼儀式.
令到你有這個身份.
保羅是很清楚.
斬釘截鐵地這樣說.
其實再沒有任何的事物.
使我們要靠它來得救.
靠它來在神面前.
可以換取我們稍為聖潔一點的身份.
其實不是.
一切都是在耶穌裡面.
在他就不潔淨.
即是說在他的軟弱的良心裡.
他仍然有這種不潔淨.
所以怎樣令到憂愁呢?.
就是第十五節.
你若因食物叫弟兄憂愁.
就不是按照愛人的道理行.
基督已經替他死.
這是愛人的道理.
你不可因你的食物叫他敗壞.
即是說意味著我們有這種信心.
我們有這種對所謂的不潔食物.
已經沒有包袱的時候.
我們就在一個有包袱的人面前.
良心有很大的限制.

$^{401}$很大的罪咎.
來做那件事的時候.
我們去吃.
就帶來一個很大的影響.
不是我們受影響.
是我們這個行動對他的良心受影響.
他的良心受影響.
他就會批評這個吃的人不聖潔.
不聖潔之餘.
還不敬畏神.
因為你的自由.
就令到成為別人的毀謗.
下面就是真的這樣說.
所以我們在這樣的行動.
吃東西的行動.
來去玷污了弟兄.
來去叛奪了弟兄.
這個是罪.
即是說雙方面都是罪.
那個人如果他這樣吃.
又是一個限制.
和他自己的良心有虧欠.
但是我們這樣的行徑.
沒有了愛心的行動.
都是一個很大的傷害.
所以我們信了主耶穌.
要很敏銳我們的心靈.
我昨天上午.
臨去我們星期四上午的福音茶經班.
好端端在星島日報見到一個這樣的新聞的報導.
應該可能是真事.
就是有一個小朋友.
當然他現在已經是大人.
他說他回憶他在小朋友的時候.
和他媽媽去海洋公園.
去看海豚表演.
離遠.
可能是過年還是一些大的日子.
見到兒童節目的其中一個主持人.
叫蓋世寶.

$^{441}$可能是那個戲名.
小朋友很開心.
等於我見到譚玉英我也很開心.
雖然他不認識我.
但他看著他長大的.
他見到這個主持人.
很開心.
他就和媽媽說.
媽媽就說你過去找他.
和他拍照.
我不知道那時候有沒有相機.
小朋友一定是很開心.
誰知道去到的時候.
那個蓋世寶就可能不知道是什麼原因.
就告訴這個小朋友.
你認錯人了小朋友.
我不是蓋世寶.
我是蓋明輝.
蓋明輝是粵劇的豪平.
那個小朋友就被這個.
吃檸檬.
那個小朋友回來和媽媽說.
很不開心.
媽媽就很生氣.
明明你就是那個小朋友節目.
《閃電全真機》的主持人.
你就這樣來耍走這個小朋友.
這個小朋友已經長大了.
這個他也是他的童年陰影.
這個沒有愛心的行動.
可以帶來這麼多年後.
被人這樣清算.
被人這樣說.
是叛變了一個小朋友的心靈.
所以原來可以這麼小的事.
會對我們.
特別是吃檸檬.
特別是一些很傷害到人的眼神.
行動.
會帶來一些人的內心很受傷.

$^{481}$那個只是其中一個例子.
所以在這個原則下.
保羅教我們.
我們要管著的是愛心的行動.
然後他告訴我們.
基督已經為他死.
這是愛心的行動.
造物主來到這個世界做人.
為每一個罪人釘十字架.
造物主就是這樣無條件.
為每一個人這樣死.
如果我們在屬神的兒女.
我們在我們的食物上.
叫他去敗壞.
所謂敗壞就是跌倒.
就是他的良心受虧欠.
受虧損.
或者就是他看到你這樣做的時候.
他就有一個很大的內心的痛苦.
為什麼你也是信耶穌.
為什麼你會這麼不敬虔.
這樣來對神.
為什麼你會這樣做這些事情.
來傷害我們猶太人的良心呢.
所以這告訴我們.
第十五節就是這樣說.
他說不可以因為你的食物叫他敗壞.
不可叫你的善.
你的善就是你的自由可以做這件事.
你有自由去做這件事是沒有問題的.
但是你的自由無限放大.
放大到變成.
你這個善客在另外一個人身上.
就成為了他對你的一個毀謗.
所以沒錯了.
你的自由就帶來他的傷害.
他的傷害就對你所做的這件事的一個毀謗.
那個合一就斷開了.
那個同心就沒有了.
如日積月累這樣.

$^{521}$我是這麼自由的.
我是要這麼做的.
我是要我的自由放大.
因為我們是外邦人.
我們沒有這些包袱的.
我們不是摩西的門徒.
我們喜歡吃就吃的了.
這樣的無限的放大自由.
對一個在那個群體裡面有猶太人.
就有一個很大的傷害.
而這個傷害是有回報的.
就是因為你的善就被人毀謗.
不是被人來感恩.
不是被人將榮耀歸給神.
所以你想想如果我們做的事情.
不是可以將榮耀歸給神.
那真的盡量不要做了.
否則只會帶來受到那個毀謗.
而且是你的善被人毀謗.
所以十七節我們再讀.
因為神的過不在乎吃害.
只在乎公義和平並聖靈中的喜樂.
在這幾樣上服侍基督的.
就為神所喜悅.
又為人所稱許.
這個就是愛心的生命.
愛是聯絡全德.
十三章所講.
所以保羅告訴我們.
一個基督徒能夠在教會裡面和人合一.
和有這份同心.
一定是因為他追求神的國.
什麼是神的國呢?.
不是等候上天堂叫做神的國.
是我們在人生裡面.
反映我們有公義.
我們有和平.
有和睦 有與人為善.
有聖靈裡面的喜樂.
就是說我們的人生是充滿了一份平安.

$^{561}$一份喜樂.
能夠讓人感受到.
我們那種親和力.
我們那種與人合一.
我們那種很想關心其他人.
剛剛就是翻譯之前.
有一位小朋友.
剛剛從幼稚園升上小學一年級.
由九月到現在還沒有笑過.
沒有笑過是很大的事.
媽媽還沒有信耶穌.
來社會中心補習.
因為都認識了他.
暑假都會來參加我們那些聖經班.
媽媽又什麼都不懂.
就去付很多錢.
問問是不是心靈有事.
跟著跟他拜一些符.
都是不開心.
見到他有一個佛珠等等.
同工很好.
仁仁就說.
你不如祈禱吧.
你那些都不行.
不如你寫下祈禱紙.
叫教會牧師幫你祈禱.
真是不開心 真是不笑.
又不是在學校受到欺凌.
就是由幼稚園這個轉變.
去到小學就有一個很不同的階段.
其實很害怕.
因為三個月都不笑.
看著他不笑.
平時回來.
星期六是學K-pop.
在我面前又跳又很開心.
幾個小朋友一起.
就走進來拿火糖吃.
但現在完全不笑.
可以想像到.

$^{601}$教會我們應該怎樣呢.
我們應該將這份聖靈裡面的喜樂.
和睦和公義.
讓他來感受到.
所以人見到的基督徒.
是不是真的追求神的國的基督徒呢.
我們的心靈是不是有公義和平.
聖靈裡面的喜樂呢.
不是狂笑的一種開心.
而是在聖靈裡面帶來的那份平安滿足.
那種的喜悅.
所以這裡說.
在這幾樣上面去服侍基督.
為神所喜悅.
又為人所稱許.
剛才說的是委旁.
現在是稱許.
所以你可以想像到.
當一個基督徒用他的自由.
放大他個人的自由.
就會帶來受到別人的委旁.
即是說我是這樣的.
我喜歡這樣就這樣.
我不會改變的.
這就是我們個人的自由.
你在家裡.
自己一個群體.
你自己一個人完全沒有問題.
但是當我們是一個很大的群體.
而這個群體是基督的國度.
是神的家的時候.
我們就要學習.
追求甚麼呢.
所以古羅真的教得我們很認真.
我們有一班追求神國的弟兄姊妹.
所以我們不要放大我們的自由.
我們自己個人的.
可以用的身份.
我們真正的身份.
就是事奉神.

$^{641}$服侍基督.
和將人帶進基督的恩典裡.
第十九節.
所以我們不要追求和睦的事.
很清楚了.
說完神的國度的弟兄姊妹.
追求的是甚麼呢.
追求的是和睦的事.
和彼此能夠建立得行的事.
即是我們不斷彼此的安慰.
彼此的鼓勵.
用神的道彼此的祝福.
彼此的禱告.
這是很寶貴的.
所以番奇督的弟兄姊妹是特別有福的.
因為她真的很認真.
將她的心靈打開.
將人的需要放在她的內心.
帶到神的面前.
這樣的弟兄姊妹.
其實是很勤神的國度.
能夠很敏銳神聖靈的同在.
所以我們只是勉勵.
鼓勵.
甚至我們很想弟兄姊妹多參加祈禱會.
因為真是在那一刻.
你的心靈就是很近神.
你久而久之.
其實你不知不覺對神的事情.
你會有一種敏銳.
當然不是說你做其他事物就不行.
因為我們始終是活在一個社會.
有不同的工作.
有不同的忙碌.
但是為自己求神.
讓你能夠在那段時間.
分別為聖.
用那段時間去專心的禱告.
陪伴教會其他弟兄姊妹.
一起同心的禱告.

$^{681}$現在那個門檻很低.
你開條鈴.
按下去.
就可以看到現場.
你的心靈就在現場.
一起同心禱告.
其實可以說是真的學習.
那半個小時安靜在神的面前.
這裡告訴我們.
追求的是和睦和健立得行的事.
不可以因為食物毀壞神的功情.
凡物固然聖潔.
剛剛他才說的.
第十四節.
保羅說其實憑著主耶穌.
我確知心信.
沒有東西是不潔的.
但是保羅告訴我們.
凡物固然潔淨.
但有人因食物叫人跌倒.
就是那個用自由去吃的人.
在猶太人面前去吃的人.
就是他的罪.
無論是吃肉還是喝酒.
是什麼別的事.
叫弟兄跌倒.
一概不作才好.
就是說這一段經文.
就是叫我們很善意地主導去接納.
還有我們將那些絆腳的石.
拿走它.
如果是我們令到人絆倒.
我們自己就要去改變.
我們自己就要悔改.
我們自己就要去追求.
我們就要放下我們的自由.
不要放大.
放下我們的自由.
然後我們追求的是基督的事.
服侍基督的事.

$^{721}$那就是服侍人.
我們裡面有愛心的生命.
所以保羅就說.
他一概不作.
如果叫弟兄跌倒.
他就一概不作.
第22節.
你有信心.
就當在神面前守著.
沒錯.
你個人.
你可以將你的自由.
限制在你自己個人的天地.
是沒有問題的.
完全沒有問題.
但是當然不是叫你躲在家裡吸毒.
那些叫放大自由.
不是這個意思.
而是人與人之間相處的時候.
我們盡量放下我們的自由.
不要讓別人受到不必要的一種傷害.
你有信心就當在神面前守著.
人在自己以為可行的事上.
能夠不自責.
就是有福了.
若有回心而吃的.
就必有罪.
因為他吃.
不是出於信心.
所以不是食物本身令他跌倒.
是他帶著一個疑惑.
帶著一個罪疚.
帶著一個他自己的良心有虧欠.
這樣的情況.
不是出於信心.
那就有罪了.
所以不是出於信心的事情.
都是罪.
所以保羅亦講了一個概括的真理準則.
這段經文很寶貴.

$^{761}$其實是回應十四章的一到十二節.
告訴我們要怎樣具體的去接納其他不同的人.
我們一路看下去.
我做了的part 1.
怎樣去應用和反思.
不同的背景.
我們有不同的習慣.
應該要怎樣彼此的接納呢.
我是否因為一些很次要的問題.
而我們去論斷其他的人呢.
大家可以思想一下這兩條的問題.
跟我們在一到十二節有關連的.
接著十三到二十三節.
我們會繼續講.
所以其實今天我們要講的就是.
凡事為主而作.
凡事去合一.
我認識一個弟兄姊妹.
挺好的.
當然這個絕對不是絕對.
沒問題的.
他個人放下他的自由.
他說如果去一班教友的旅行.
因為他自己是做生意的.
很多時候積分可以換到一些商務位.
或者一些比較高一點的檔次的座位.
都是長途的.
他說如果跟一些朋友一起.
頂尖一起.
他就情願陪伴他們坐經濟客位.
這樣可以說他放下他自己.
不是裝模作樣.
不是要特意去收起來.
而是他能夠放下他自己的自由.
一齊去坐經濟客位.
這些其實是一個很大的思考.
他體會不同弟兄姊妹的新聖靈的感受.
這些其實是一個很成熟的表現.
這些很具體的.
這些就告訴我們.

$^{801}$我們是不是要這樣去將我們自己個人的自由.
保羅說你有信心.
你可以在神的面前守住.
無論吃肉或者喝酒等等.
一切的東西其實沒有對錯.
但是問題是如果有疑心而吃.
這就是一個問題.
當然不要反過來.
你見到別人跟你不相同.
見到他坐雙步客位.
你就有一個論斷.
這就不是愛心.
因為每個人都有他自己的東西.
這就不是跟他個人的自由有關連.
所以如果能夠彼此都在一個愛裡面.
這個水平去看其他人跟你不相同.
這個群體就會有一個很美麗的展現的合一和愛.
所以很重要的.
我們這個行動的綱領.
凡事為主.
凡事能夠合一.
學習這三樣東西我們一起.
以致教會或者我們這個群體.
或者我們這個人.
能夠讓人看到基督的生命的尾線.
我們希望的是讓人看到基督的生命的尾線.
我們繼續下去.
所以大家一起去看看.
13到23節.
為什麼我們要小心.
不要成為別人見到的絆腳石.
就是說我們做的事情.
因為我們做的事情背後.
就是因為涉及到我們的自由.
我們常常當我們自己個人的自由.
如果我們這個個人的自由放得很大.
就會不知不覺成為一個絆腳石.
這個絆腳石就會絆跌了其他人.
所以就是這段聖經的一個很寶貴.
原來就是很積極.

$^{841}$不論積極是怎樣.
就是接納別人.
在世界上的亂叛亡裡面.
我們彼此來接納.
我們在1到12節.
就不要將自己放在審判席上.
因為12節說.
我們眾人都要將我們一切的事.
呈現在神的台前.
13到23節.
就是一個很具體的.
就是怎樣能夠接納別人呢.
就是不要把弟兄推到絆腳石那裡.
如果我們簡單做一個宏觀的分析.
保羅首先處理了心態上不要有錯誤.
再處理行為上具體的行動.
不要有錯誤.
所以心態就是我們的內心就是論斷別人.
這樣做和我不相同.
就是不對了.
這個就是論斷.
保羅告訴我們心態會有行動的.
所以心態如果不正確.
行動也不正確.
但當我們的心態正確.
我們也能尊重每個人有他的良心.
有他的文化背景.
這些是次要的問題.
不是我們可以隨便扮演審判者.
如果我們有了這樣的學習.
我們就會對別人有一個接納.
欣賞和包容的時候.
我們就能夠在改變心態之後.
我們的行動和行為都會正確.
這是很重要的.
因為如果我們的神學.
我們只有一套理論.
而這套理論是牢不可破的.
是這樣就是這樣.
這樣就不行了.

$^{881}$我們就會看什麼都不對.
看什麼事物都會看得準確.
不準確就會帶來你的判斷.
就會越來越深沉.
越來越有挑剔.
越來越帶負面情緒.
越來越帶有色眼鏡.
這樣就不對了.
原來你看所有的東西都已經有了偏見.
所以其實源於傲慢.
內心的傲慢.
所以求神幫助我們.
你看保羅用了一章聖經.
講教內處理了和帝王.
和社會的政府的關係.
然後用了一章聖經.
教我們每一個弟兄姊妹.
不是教你做好人.
而是在教會裡面.
我們都是神的請求的贖罪制.
我們是一個活的祭物.
獻給神.
我們如果有這麼偉大的開始.
每一個人都是神喜悅的活祭.
為什麼我們會成為別人的偉旁呢?.
這是保羅在後面的講法.
15節告訴我們.
他會來偉旁你.
不可以叫你的善被人偉旁.
就是說活祭是一件事.
實際的生活是另一件事.
實際的生活.
我們怎樣去處理自己的自由.
就是會令到我們到底.
是讓人看到我們那個生命而帶來喜悅.
對神有一個感恩.
對我們所做的又有一份的欣賞.
不單止不是偉旁.
更加不是一種論奪.
而是一種更加投入在教會裡面的服事.

$^{921}$所以這一章聖經其實很實際之餘.
是相當重要的.
我們繼續讀下去.
所以基督就是我們的主.
所以不能夠論奪的原因.
最大的原因.
記住我們一定要記住.
我們真的要服事主.
我們可以討論.
絕對不是討論就等於論奪.
我們可以開會檢討.
沒問題.
完全有不同意見都好.
我們就去想想.
怎樣下一次的活動做得更加好.
包裝得更加大家都能夠好的崗位.
應該是要這樣的.
所以不同意見不同的看法是沒問題的.
但是真正的論斷的背後.
就是對那件事.
你有你自己那種錯誤的心態.
去對別人的來去定性.
記得我們絕對沒有權.
扮演耶穌這個生命的主的身份.
所以如果我們認定我們只不過是僕人.
我們就會在教會少了很多的煩惱.
少了很多的質鬧.
反而多了更加對別人的接納.
13到23節.
基督既然是主的時候.
基督都為弟兄死.
我們就更加要學習為弟兄去犧牲.
不是放大我們的自由.
犧牲和自由剛剛就是一個對立.
基督就是放下他神兒子的名份.
那個位份.
來到這個世界.
本有神榮耀那個形象.
那個本質.
但是反倒虛己.

$^{961}$所以是放下他的自由.
耶穌基督.
所以教會往往最重要的.
就是我們需要在初歲班成長班.
不斷在教會的講台多教導.
就是我們要學習基督的虛己.
基督的放下.
就是說有意識放下的人.
其實才是成熟的基督徒.
什麼都要拿盡自己那個權利的人.
絕對不是一個成熟.
是幼稚的基督徒.
一直都不成長.
譬如老師班.
我又是亂說話.
沒有任何的指示.
你就要霸佔咪高峰.
一定要站在中間.
可能你未必是最.
每個當然由群體的領袖安排.
但是如果我們的心.
是不斷放大我們自己在群體裡面.
要成為話事人.
要成為那個領袖.
那你現在是最好的事.
求神幫助我們要學習.
為什麼我們英文那裡.
為什麼我們需要愛多過個人的權利自由呢.
所以個人的權利.
是和基督的虛己.
是一個剛剛相對立的事.
求神幫助我們學習.
我們越是信主內.
我們就要學習.
我們很樂意因為基督.
我們放低我們自己.
你堅持自由.
可能就是毀壞了基督所愛的人.
因為神的國不在乎吃喝.
只在乎公義和平.

$^{1001}$並聖靈裡面的喜樂.
所以不是說吃喝是罪是錯.
是相對於這裡說的例子.
這裡是文化之間.
良心之間的信心的程度之間的差異.
那裡是關乎那些飲食.
所以是突止那些不潔的食物.
令到猶太人感覺我吃了這些經過處理的食物.
神就會閱納我們這個人.
但是舊約而已.
新約耶穌來到這個世界.
舊恩完成了.
我們一切的東西已經再不是這樣來運作.
那些只不過是一個禮儀上.
不會因為你吃了.
是會帶給你們有那個聖潔.
譬如現在猶太人.
還是吃那些不潔的食物.
經過那些拉比處理.
甚至乎空運來香港.
但是吃了那些是不是可以上殿堂呢?.
完全不行.
因為唯有基督十字架的舊恩.
才更新我們一切.
更新我們的生命.
所以不可以再用這些飲食.
令到你的身份會有改變.
在神面前蒙恩.
相反就是哪些人不能在神面前蒙恩呢?.
能不能得到神的喜悅呢?.
已經很清楚了.
就是我們有意識地在群體裡面.
去追求神的國度.
神的國度很具體地說出來.
公義和平.
就是說我們是走在神的真理底下生活.
我們來到和睦.
和睦可以是心靈的平安.
可以是散播平安.
可以是成為和平之子.

$^{1041}$可以令到那些紛爭因為你的出現.
可以帶來大家是一個的合一.
這個就是和平了.
就不是叫你坐在那裡出聲都叫和平.
只是沉默和和平是兩回事.
一個和平之子是很喜歡和別人聊天.
將你的歡樂.
將你的喜悅.
將你內心的盼望去和別人分享.
今日很累了.
一早要出去莫大仙.
為一位老人家祈禱.
因為真的未老.
然後年細.
家人都在那裡.
中午突然間又收到電話.
有個弟兄的外公.
這個爺爺又來漏了.
又要去.
本來去不到的.
因為他的父母是拜天道五教.
是不讓教會的牧者去的.
突然間他媽媽.
即是我們的弟兄.
紅恩的弟兄.
他媽媽就說.
叫你教會的牧師來.
爸爸不行.
他老爺不行.
他拜天道五教.
他知道其實他的老爺.
即是我的弟兄的爺爺.
原來很多年前是紅恩的.
很奇妙.
當然後來因為很老.
老人院少了.
是媽媽叫兒子去叫牧師來的.
我剛好在莫大仙.
馬上趕回屯門醫院.
又是兩個離婚的老人家.

$^{1081}$讓那個和牧.
他的爸爸在那裡.
他爸爸也是未信的.
他說不用想.
安息禮拜的了.
多數都不用了.
因為他們真的很迷信.
但是感覺到牧師為他的爸爸祈禱.
他都好像有點感應等等.
所以原來我們追求神的國.
我們的出現.
是帶給那些徬徨.
未信的人.
有一個盼望.
所以剛才下午我安慰那位女士.
都認識她的.
因為暑假很多時候.
每個星期六都過來.
帶小朋友過來跳K-pop.
都認識了.
經常都聊天.
我安慰她.
這個改變要留心.
帶她回來教會.
她去試試.
其實很擔心.
因為真的三個月不笑.
那位小朋友.
又不知道什麼事.
所以我們要追求的不是吃漢.
吃漢就是個人要求.
自己的自由.
當然不是放縱那些.
保羅相對的吃漢.
就是你我行我素.
你自己用盡你自己生活上的自由.
但是這樣不是幫到其他人.
祝福到其他人.
不是說這些是錯.
不是說這些是罪.

$^{1121}$而是相對於你能夠幫到其他人.
所以保羅說更加應該追求的.
在乎的.
在你的生命裡的事實都在乎的.
就是公義和平.
聖靈裡面的喜樂.
在這幾樣上去服侍基督.
就一定是神所喜悅.
十三戰.
我們不要再彼此論斷.
反而要決定這件事情.
不要將弟兄放在半腳石.
我們自己.
或者我們令到他感受到那個冒犯.
或者對他的一種歧視.
白眼等等的負面.
使他跌倒.
十四十五戰.
我們再仔細看.
知識上面.
保羅說食物本來是潔淨.
但是如果我們有愛.
我們就不會隨便將這個知識放大去用.
我們就不會這樣做.
所以保羅告訴我們.
他自己個人憑著主耶穌確知心信.
的確知道有一個真實的真理發生了.
就是耶穌基督將那個救恩完成.
一切都更新.
所以猶太人需要更新.
所以保羅自己就在外邦人的面前.
更新和他們一起飲食.
完全不會有任何的避諱.
但是當他回到耶路撒冷.
他自己又繼續.
仍然是一個猶太人的血統.
他上到聖殿.
他就學習守一個拿西爾人的願.
有些東西他就不吃.
還有他就不剃頭.

$^{1161}$等待那個願的日子滿了.
他就上聖殿.
找祭司幫他剃頭去還願.
所以保羅自己.
能夠很有分寸地.
何時用他的自由.
在他的身份.
他帶著的是兩個的血統.
兩個的身份.
基督的血統.
就是十字架的寶血.
將他改變,拯救,使用他.
但是他自己與生俱來的血統.
就是猶太人受猶太教的昏徒.
在嚴謹的猶太教受教.
是一個拉比.
他要教人猶太教的法規.
但是保羅現在.
他就可以猶日有餘地.
行走在兩個文化中間.
因為他確知.
再不是食物.
再不是甚麼日子.
再不是甚麼事情.
讓我們特別聖潔一點.
是因為基督使我們有這個新生命.
所以保羅說.
我憑著主耶穌確知心信.
凡是凡物本來會有不潔淨.
唯獨人以為不潔淨.
在他就不潔淨了.
所以說原來每個人都有他良心上的軟弱,不安.
當然慢慢在這個對話裡.
就要讓這些猶太人都明白.
舊因已經更新一切.
所以猶太人現在.
你和我們一班外邦人一起飲食.
其實在神的眼中.
這是和諧的合一同心.
不再是你食不潔食物.

$^{1201}$你和外邦人來往.
所以猶太人慢慢都要學習.
這個慢慢過程.
外邦人的信徒.
就要慢慢接納他.
欣賞他.
等候他.
他的良心慢慢的去更新.
所以這個過程.
其實就是接納.
就是欣賞和愛的互相付出.
所以保羅原來教我們.
同不同的文化,差異的人.
有一個愛的共融.
彼此的合一.
要欣賞對方.
所以愛就能夠讓這些事情.
慢慢的更改.
我們繼續下一個部分.
保羅自己的立場.
又確知深信.
食物本身不再有.
任何靈性上的污穢標籤.
不會再有.
因為舊約的確有這樣的分水嶺.
但是那只是為一班人.
預備了.
告訴他們什麼是聖潔.
什麼是污穢.
什麼是神喜悅.
什麼是那個情況.
等於我們的小朋友讀小學.
讀幼稚園.
早上臨上學前.
一定要吃完早餐才可以上學.
如果不吃早餐的小朋友就不乖了.
就不行了.
你的身體就不會健康了.
他們記得.
長大了三十多歲了.

$^{1241}$今早忘記吃早餐.
慘了.
罪大惡極了.
爸爸說不吃早餐.
張開口就這樣.
已經不需要了.
因為他BB的時候.
小朋友的時候.
是有這樣的規矩.
教他這樣的法則.
但是當他長大了.
加完班.
三點多才回家.
根本還沒醒.
然後早上又要去健身.
他已經不同了.
他已經是大人了.
不吃早餐是沒有問題的.
因為他根本就是.
他的工作更加重要.
所以大家明白了.
到了新約.
耶穌完成救恩之後.
一切都更新了.
當然有例外.
比如保羅就罵哥林多人.
你們就是愛邦人.
你們就完全覺得自己沒問題.
走進那些偶像的廟去飲食.
去坐席.
明知道你們是交過軌的.
你們都吃.
保羅就責備這些人.
因為我們是跟神接觸的.
我們是敬拜獨一的神.
我們不能夠跟偶像有份的.
所以這個是例外.
但是現在說的是文化上的.
飲食上的.
很次要的做法上的問題.

$^{1281}$不是剛才說的那些.
真的關乎到靈性.
神的屬性問題的事情.
所以大家要明白.
所以保羅用這裡的時間.
盡量將.
既然這些不是關乎神性的.
神的屬性的事情.
我們每一個基督徒.
應該盡量將我們的自由放低.
不要將我們的自由無限放大.
如果他在認為那是錯的情況下去做.
就變成他在神面前犯罪.
\newpage



\section{羅馬書}
\label{sec:pe8biigoai0}
\textbf{2026-01-16 周五晚查經班《羅馬書(34):分題:擔得起,才配得進神的國》 講員:曾裔貴牧師}
\newline
\newline
連結: \href{https://youtube.com/watch?v=pe8biigoai0}{\texttt{https://youtube.com/watch?v=pe8biigoai0}} ~~~~ 語音日期: 2026-01-17
\newline
\newline
\hyperref[sec:nB_6d0PSa6M]{\small{< < < PREV SERMON < < <}}
~
\hyperref[sec:index_chronic]{\small{[返順時目]}}
~
\hyperref[sec:index_scriptual]{\small{[返順卷目]}}
~
\hyperref[sec:UU4dvSZbobM]{\small{> > > NEXT SERMON > > >}}
\newline
\newline
$^{1}$好 弟兄姊妹晚安 我們一起同心禱告.
天父啊 我們在晚上很感恩.
我們的天氣又暖和了.
我們又可以在這段時間.
還有當錄音放在網上的時候.
又會讓不同地區的弟兄姊妹可以一同學習.
我們真的很盼望我們有系統地研讀羅馬書.
對裡面的經文的那個準的脈絡.
保羅看到 唯有萬民必須因信稱義.
因為萬民都是在罪 在神的憤怒底下.
幫助我們來到第15章.
都是有這幾章聖經教會的生活的倫理.
人與人之間的關係.
還有最終保羅很想福音要傳到不同的地區.
不同的民族.
所以幫助我們在最後這幾章都能夠整理.
尤其是關乎到我們上個星期提到.
相當重要的14章17節.
因為神的國度不是在乎一些外在的.
一些傳統 一些我們看為很重要的飲食.
或者一些的節日等等.
那些傳統是幫我們去敬拜神.
是幫我們的生命去經歷神的.
而我們著重了那些傳統的時候.
正正就好像猶太人他們會來與外邦人有婚飯.
所以求主幫助我們在乎的是公義,和平.
以及聖靈裡面的喜樂.
讓我們一起來同心在地上服侍我們的主.
今天晚上我們再繼續下去.
保羅有甚麼的進心的教導.
都讓我們能夠明白.
我們一起仰望等候禱告.
奉耶穌基督的名 阿們.
好 我們來看看.
我用的題目就是「三十四講」.
擔得起 其實是第15章的第一節.
我們來兼顧的人要擔待那些不兼顧的人.
就是軟弱的人 他們那些軟弱.
我們就能夠一同去建立那個德行.
那個德行當然就是群體一班人.

$^{41}$我們在主裡面一起成長.
在一個合一裡面來服侍.
這樣叫做建立德行.
這個告訴我們很重要.
我們一班等候神國的人.
我們準備要進入神的國度.
絕對不是坐著信了耶穌就可以的.
保羅都告訴我們.
我們眼前在我們每個星期聚會.
都有跟我們不相同民族.
不相同的文化的群體的人.
我們需要跟他們同心.
那怎麼能夠有這樣的學習呢?.
所以這個就告訴我們.
我們這兩章承認是相當重要.
我們很快去溫一溫14章的17節之前.
我們來看回14章 保羅帶出一個話題.
13章我們講過了.
是對當時的社會那個制度.
無論你認為怎樣才是絕對的公平,公義都好.
既然在一個相對是帝國.
剛剛羅馬帝國最強盛的時候.
帶來的真是四海升平.
在這樣的情況下.
無論帝國,帝王他們是敬拜偶像.
或者他們倒喉易思都好.
在那個相對的安全平安底下.
基督徒生在那個時代.
就需要對那個時代有一個的馴服.
那個馴服背後是因為仍然相信.
這個世界是天賦的世界.
是它掌管著的.
所以它容許不同的政權的交替.
好像現在伊朗好像就快要有一個王朝的交替那樣.
我們只能夠祈禱.
我們不應該偏幫任何一邊.
當然對於罪惡,對於邪惡.
我們一定要來譴責.
但是我們求主讓我們看到.
做物主掌管著萬有.

$^{81}$既然它容許羅馬帝國相對地來治理.
當時的城邦,當時歐洲各國的時候.
我們作為當時代的基督徒.
我們就要馴服.
所以保羅是很清楚.
不是對政權的歌功頌德.
而是對神的馴服.
來在那個時代去將福音傳開.
這個才是最要緊的.
因為神設立這個教會.
一班人聚集.
就是要讓那班信了主的基督徒.
在善後門,在當時的領袖帶領下.
來將一個這麼寶貴的福音.
拯救世人,罪人的福音.
要來傳講.
所以教會必須在這樣的安定的狀況下.
來傳揚福音.
所以這個也是教會存在在地上至關重要.
當然大進那個救恩的途徑.
那些就不需要來爭拗.
但是教會的存在必須是因為福音.
因為要讓人認識耶穌.
敬拜這位耶穌.
用什麼方法去帶人去到耶穌那裡.
那個是可以不同的方法.
不同的時代都有不同的方法.
這個是我們要明白的.
所以去到處理了對外政權的那個馴服.
無論用那個例子就是納粹,納良.
要來去遵從.
等等.
但是又告訴我們.
沒錯 我們似乎看到現在那個國度是這麼強大.
好像羅馬帝國是永遠都不會倒下的.
保羅說不是的.
其實黑夜已深了.
白晝將近了.
我們應該脫去暗會的行為.
這些似乎好像在暗示.

$^{121}$帝國他們那個表面很風光.
但是實際背地裡.
他們因為有極權力的掌控.
所以那些貴族當時都是非常的.
那些道德生活都是敗壞的.
但是我們信徒歸入了基督的門下.
我們就要跟隨主.
我們就要在這樣的狀況.
就是黑夜已深 白晝將近.
白晝將近就是主耶穌.
這一位清晨的日光.
將要照臨這個地上.
所以是這樣來告訴我們.
我們需要這樣來警醒等候.
接著就是14章.
保羅又再一次將焦點放回教會.
他就很緊張.
教會必須要達至一個.
不是統一 是合一.
不是要所有聲音都不准出聲.
是要大家有不同的意見.
但是在心靈裡面是合一的.
是在愛裡面合一的.
在神的拯救恩典裡面是合一的.
在我們活祭這個生命底下.
我們是合一的.
這樣來處理大家的那個.
可能是不同的文化.
而帶來的這些沒有真理的分歧.
這些分歧是很次要的.
沒有影響我們大家都是神而類這個絕對身份的.
所以這些是我們要很小心的.
教會不要給一些次要的議題.
變成我們為那些次要的議題來爭論.
來爭論不休 勢不兩立這樣的心情.
用盡最大的那個心情.
連睡覺都不睡.
這樣去想對方是怎樣差 怎樣壞.
這樣去準備用什麼方法來對付他.
而那些純粹是次要的問題.

$^{161}$教會往往就變成一個很大的分裂的原因.
是相當可惜 是相當痛心.
是相當不需要的.
所以這兩張聖經其實真的帶來一個很重要的就是.
當我們在愛裡面接納.
當我們堅固的人.
就是說我們很明白真理已經被釋放.
心靈都是被釋放的人.
我們就要有這種的擔待.
有這種的強壯.
就是一個強人.
強人在地上 在未信的人的世界.
強人就是所有人都要服他.
不服他就要被他打服.
不是這樣的.
強人 等一下我們十五章就會說.
所以我們回來看十四章的時候.
很重要的就是保羅用了兩個例子.
節期的飲食.
和一些特殊的節期.
保羅就來到很明顯.
就是說意味著有一些羅馬的弟兄姊妹.
是猶太的信徒.
輾轉可能就是猶太 耶路撒冷.
輾轉就是去到意大利.
在那裡來生活 定居聚會.
甚至乎可能去到才信主 成為基督徒.
帶來的文化上的差異.
剛剛大家都是在基督教剛剛醞釀的時候.
大家都還沒有空間去思考這些的神學.
怎樣舊約過渡到新約.
保羅的書信都還沒有全面去探討.
到底舊約在耶穌基督降世.
成就舊恩之後.
是不是應該有一些更新呢?.
這些就是往後保羅的書信裡面.
會碰到的問題.
例如《已不所須》的第二章.
提到的他的丁十字架.
就將兩下歸為一體 與神和好.

$^{201}$兩下是猶太人和希利尼人的和好.
在十字架上廢掉了冤仇.
這個冤仇就記載在舊約上的.
有礙於我們的律法.
怎樣十界 怎樣舊約將它剪掉.
將它扔掉 還是將它燒掉呢?.
保羅說不是的.
而是舊約告訴我們的背後的精神.
基督就是那個救主.
祂就是那個贖罪祭.
祂就是那個安息日.
祂就是那個彌賽亞.
祂就是那個大祭司.
祂就是那個聖殿.
這樣就告訴我們.
原來基督的出現.
是一個劃時代的里程碑的開始.
所以在保羅寫這封信的時候.
猶太人 外邦人還未能夠全面拿捏舊約的定位.
以及猶太人已經有數千年的.
守舊約文化的背景.
那個情意結.
應該怎樣去梳理呢?.
這些往後使徒.
他看到這些問題.
已經是關乎神學性的問題.
他就需要去碰這些課題了.
但是現在很實際就是人與人之間的相處.
現在首先要解決生活的倫理問題.
所以那個焦點就是彼此之間.
一方面要接納.
就是說外邦人你就要接納猶太人.
他們不是這樣做.
猶太人就不要論斷外邦人這樣做.
大家就會帶來一個好像彼此的讓步.
其實沒有讓步.
外邦人仍然照樣吃自己的食物.
猶太人照樣守自己的節日.
照樣在一些飲食上不吃等等.
就是大家有大家做.

$^{241}$但是大家在愛裡面一起聚會.
這個是很難的.
如果我們明白主耶穌的救恩.
其實在這個大前提下.
這樣的包圍.
這樣的Umbrella下.
其實是不難的.
但是當我們永遠站在.
我是猶太人.
我是外邦人.
接著我有我的自由.
我外邦人有我們外邦人的自由.
這樣就是我們變成了.
自由本身不是錯的.
但是當我們一濫用了這個自由.
本來不是濫用的.
本來就是你的.
但是當與一群人一起.
大家是奉主的聚會.
那個自由就要增長了.
那個自由就要一個成熟的基督徒.
就會放下自己的自由.
去遷就去接納對方.
去讓對方感受到.
大家是神而屢的這個生命.
這個身份.
這個聚會的那個美麗和寶貴.
所以那個焦點仍然都是.
我們要來學習不要論斷.
以及不要對別人有這樣的否定.
這個告訴我們.
接納的原因就是接納的差異.
彼此的差異.
文化,良心.
這些與舊恩是沒有關聯的.
每個人都有自己的背景.
每個人都有自己的一些想法.
一些的傳統.
說的是大家都在宗教的信仰裡面.
如果不是關乎到一些神學性的.

$^{281}$關乎真理的.
關乎舊恩的.
那就其實很次要了.
當然我們要明白.
什麼才是真正的次要.
什麼才是真正的重要.
我們一定要不斷去反思.
什麼是真理.
什麼是次要的事情.
所以在愛裡面去尊重.
而不是爭辯.
繼續下去.
這裡都是錄得不好.
這是我們上個星期都講過的.
不同的背景,不同的習慣.
我們要彼此的接納.
怎樣接納.
這個就可以大家去討論的.
我們千萬不要因為次要的問題去論辯其他的人.
所以凡事為主.
凡事都是能夠尋求一個合一.
來在主裡面.
來這樣的同心.
接著我們又再講過.
13到23節.
其中的總結句就是.
我們要小心.
我們不要成為別人的絆腳石.
我們去判別了其他的人.
所以我們要很留意這些事情.
在永恆的裡面.
我們盼望的.
就是我們要在主的面前.
14章.
彼此的接納.
就是不要將自己放在一個審判的席位上.
我們成為審判者.
去審判其他的人.
尤其是一些不同的做法.
不同的看法.

$^{321}$不同的意見.
那些不是絕對的.
是相對的.
不同的角度去看這件事.
可能會豐富了結論.
接著第二個段落.
最後一個段落.
就是千萬不要成為別人的絆腳石.
以及不要故意的判別其他的人.
這些告訴我們.
基督才是可以審判的那一位主.
如果祂也不審判我們.
祂也是用永恆的愛愛我們.
祂也是為我們無條件的釘十字架.
這是第十五章的第五節.
第三節是保羅引的.
他說肉麻美人的肉麻都落在我身上.
是保羅引述大衛的詩篇.
彌賽亞的詩69篇的第九節.
所以連基督都願意忍受這些痛苦.
這些肉麻.
何況我們今天不是審判處.
記得我們要學習.
應該我們年紀越大.
越來越多失敗.
越來越多自己的軟弱.
我們就更加應該要學習.
我們沒有權作審判者.
不是不可以判斷.
而是不可以做審判者.
我們是不斷用不同的智慧去思考.
去辨別真理.
辨別什麼對錯.
辨別什麼做得好一點.
什麼食物對我們自己都好一點.
我們不斷在做辨別.
連穿什麼衣服 襯什麼褲子.
某程度都好.
不要讓自己好像從方舟出來一樣.
這樣就不好.

$^{361}$應該是我們大家都有這樣的不斷成長.
但是判斷和論斷.
和我們成為一個審判者.
是完全兩回事.
所以求神幫助我們要懂得分論.
所以相信有愛在心靈裡面.
愛是聯絡全德 保羅說.
所以愛就是不會加害於人.
所以你想想原來一骨子裡面.
我們在判斷的過程.
那個自我就會出現.
就會蠢蠢欲動.
不知不覺就自己扮演那個管理者.
那就千萬不要.
接著我們要記得.
主耶穌他才是那一位主.
他為我們犧牲.
為弟兄來犧牲.
所以我們必須要記得.
我們不要將我們自己個人的權利放大.
來超過我們內心的愛.
所以我們要知道.
十四章 接著保羅就說了最重要.
在這兩章聖經.
其實都不止這兩章.
這句聖經其實涵蓋整個羅馬書.
最終達至的.
為何要恩順清義.
為何我們要與神和好.
為何我們要認識我們本來是神憤怒的兒女.
為何現在我們可以為神的國去效力.
一定是因為救恩將我們改變.
我們現在是活祭.
我們現在在地上.
我們的生命是屬神的.
我們要尋求神的喜悅.
來為神而活.
問題是怎樣為神而活呢?.
是否我們自己走去深山拿著聖經.
一個月讀完一次.

$^{401}$回來告訴勵士蒙聖博士.
我讀完一次聖經一個月.
其他弟兄姊妹都未起步.
是否這樣叫追求神的國呢?.
這是其中一個很小的部分.
是你個人的一個裝備.
但是保羅說得很清楚.
弟兄姊妹的羅馬書太多金句了.
但是這一節聖經.
好像一條很重要的約事.
保羅已經告訴我們.
所謂何事我們成為神的兒女.
一方面就是我們認定.
這個世界不是神的國.
這個世界是準備我們每一個人.
去進入神的國.
所以我們在地上.
雖然我們仍然有一個上命的功能.
但是我們等待的.
永遠是神永恆的國度.
今天早上就去了一個很完美的句號.
那位在北區醫院.
素未謀面的一位女士.
因為癌症的末期.
認識我們的姐妹.
偶然間在隔壁床位.
今天最終都和她一起.
和她家人一起.
一個很簡單的儀式.
就是這樣灼灰.
就完了她一生在地上的人生.
一個很小的石牌.
那些灰和其他的.
在和合石的死者.
都混在一起了.
將來她在永恆裡面.
就會有新的生命,新的身體.
但是現在我已經和她的姐姐,弟弟說.
嘉慧現在她的靈魂.
已經在永恆了.

$^{441}$是不是看著我們笑呢?.
我當然不知道.
我還沒有見到.
我沒有這個機會見到.
但是相信她現在在永恆裡面.
再沒有痛苦,再沒有眼淚.
已經是一種很真正的自由.
這樣來與神同在.
所以保羅在這裡告訴我們.
我們在生的這幾十年.
必須要追求這幾樣.
這幾樣是什麼呢?.
這幾樣真正是我們的生命.
與神一起的一個屬靈生命的成長.
所以我們要為裡面的生命去餵養它.
我們現在都知道.
譬如你來崇拜,來查經班.
你不吃東西你都很累.
做完一天的事你都很辛苦.
你不喝水你又會很渴.
是不是?.
如果你兩天都沒有洗澡.
你開始很癢,不舒服.
我們要養我們自己的身體.
但是我們比較少人去想.
我們成為基督徒.
要養活我們自己裡面的靈命.
我們沒有好好去準備.
原來神的國不是在乎愛在的吃飯.
這個吃飯.
當然保羅不是在說吃自助餐.
保羅在說的.
我們上下文很清楚的.
所以我們為什麼羅馬書.
或者我們讀聖經.
一定要根據上下文去看呢?.
上下文在說的是什麼?.
在說猶太人很緊張傳統.
很緊張哪些東西可以吃.
哪些東西不准吃.

$^{481}$保羅說外邦人完全沒有這些問題.
這樣的對立面.
怎樣去梳理它呢?.
這裡告訴我們.
保羅很想強調.
告訴我們.
17節這裡告訴我們.
我們要學習真正的成長.
就是為神的國度.
我們這個屬靈的生命.
不斷地去成長.
怎樣成長呢?.
原來在乎的是公義和平.
聖靈裡面的喜樂.
這個很抽象.
很神學.
什麼來的?.
怎樣裡面有公義?.
是不是多讀法律的書?.
還是多背十誡?.
等等.
那個是帶我們去明白.
什麼是公義的標準.
那個就是公義.
公義是怎樣來實踐呢?.
接著保羅在第十五章.
告訴我們.
這個神的國度.
是怎樣讓我們一個堅固的人.
我們的屬靈生命.
如果能夠成長成熟的時候.
我們怎樣去擔待那些軟弱的人的軟弱.
這樣就是神國度的彰顯.
就是在我們的生活聚會裡面.
與人的相處.
讓人看到基督.
讓人看到基督是怎樣的一個學習.
他怎樣來將自己擺上.
十五章我們一會兒逐一會看.
我都做了part point.

$^{521}$所以我們回來這裡的時候.
原來十七節這條約時.
真的解開了教會之間的倫理關係.
當大家都很看重神的國度.
很看重我們要得到神的生命在我們裡面.
來去成長.
以致我們在生活裡面.
來讓人看到我們是屬神.
這幾件事是不會錯的.
所謂的在乎公義.
就是我們這個人本身是要正直.
是要適合神的真理.
我們的和平就是我們待人接物.
能夠在星期二.
停車後八時.
在柏宇花園.
八時就很早.
不是,應該是九時.
星期三八時.
七時多,不是八時.
就回去準備.
回教會放下東西.
星期三要進青年會小學.
所以要早點回去放下東西.
就再去.
有一個婦女不認識她.
住在柏宇花園.
差不多我剛放下.
就跑了.
就見到這個婦女.
她說你是牧師嗎?.
我說是的.
然後她說當然了.
完全不知道.
她可能認得.
因為都是附近的街坊.
都是住在柏宇花園.
就很感受到她那份喜樂.
她那份平安.
我不知道她是不是信主.

$^{561}$但是一早這樣打招呼.
又令到我感覺很舒服.
有一種大家好像很和諧的關係.
大家好好的.
每個人早上打招呼等等.
要告訴我們.
我們其實要在這樣的生命裡.
去在乎.
在乎就是你很緊張.
在乎.
以前電視是要錄影的.
以前沒有上網直播.
沒得看就沒了.
所以那時候的技術沒那麼先進.
要錄影的.
所以一到時間就要錄影.
如果錄不到就慘了.
又沒得看直播.
又不知常常看直播的決賽.
慘了.
所以有時候要調鬧鐘.
看世界盃決賽.
在乎.
在乎我們明白.
我們如果在乎.
一定會有很多的心放進去.
所以本來在這裡.
告訴我們.
這是很神學的抽象.
十四章就說了這件事.
所以不是在乎那些傳統.
甚麼要做到最足夠.
所以今天安息日.
哎呀 忘記了安息.
他很不安息.
整天都沒有敬拜神.
耶穌都說了.
人子來.
人子才是安息日的主.
那些法理賽人.

$^{601}$就一直留意著安息日.
看看那些門徒.
看看耶穌有沒有走過那些麥田.
就摘那些麥碎來吃.
那些門徒就有點餓了.
就真的這樣撕了一點來吃.
就很緊張.
抓住那些門徒.
就來問.
為甚麼你的老師安息日.
他們不來跟隨呢.
遵從呢.
摩西說安息日要禁食.
不准甚麼的.
耶穌聽到之後.
就跟這些門徒說.
人子才是安息日的主.
真是好像魚雷貫耳.
好像告訴我們.
保羅告訴我們.
保羅一個猶太人.
這個拉比.
這個一直自己舔手律法的拉比.
他說不在乎吃漢.
他這裡說的吃漢.
是針對著猶太人.
他們很緊張這些傳統.
很怕外邦人會玷污他們.
很怕他們跟這些外邦人一起.
自己會來得罪神.
以致他們就很緊張.
哪些是可以吃的.
哪些是不可以吃的.
這一天很重要的.
一到時間我們馬上要守安息日.
就很大地將心力.
投放在那些物件上.
而不在乎那個心靈.
所以神的國度不在乎吃漢.
記住這個是吃漢.

$^{641}$不是一般的飲食.
是說猶太人遵從那個傳統的十戒.
和成為一個典型的律法之子.
那些的條規.
所以保羅不是叫我們廢棄那些條規.
而是那些條規是帶我們的生命.
去到認識我們的救主.
我們不可能在救主的面前.
不斷埋頭去看那些次要的事情.
很多年前都有聽傳人講一個故事.
相當深刻的一個例子.
我想再講多一次.
和這次聖經也有相當關連的.
有一個小朋友一出生.
他爸爸就不在家了.
因為後來有些戰爭.
他爸爸又回不來了.
他就問媽媽為什麼我沒有爸爸.
媽媽就說不是.
你爸爸去行選.
還有那段時間因為打仗.
所以他在外地回不來了.
於是就過了一段長時間.
有一天他媽媽就說.
這個就是你爸爸.
給他一張照片看.
這個就是你爸爸.
他拿著照片.
我就經常看著照片.
希望爸爸快點回來.
終於有一天.
爸爸真的回來了.
小朋友就不認他.
你不是我爸爸.
我爸爸是這個.
這就很有趣.
我們會不會猶太人死守的傳統.
救贖主尼塞馬來到這個世界.
那些傳統全部是因為耶穌而設立的.
是要將人帶到耶穌那裡的.

$^{681}$但是猶太人對著那個祭物.
對著那個節日.
對著那個時間.
來去敬拜神.
他們忘記了那個神已經站在他們面前.
已經降臨到地上了.
所以我們都不要這麼傻.
好像我們一直看著那張照片.
去等候那位尼塞馬.
等候那位愛我們的主.
所以今天今生我們在乎的是.
公義和平並聖明裡面的喜樂.
在這幾樣上面服侍基督.
為人所愛.
也是神所稱許.
所以告訴我們.
剛才說了.
所以不是外在的行為.
是裡面內在的生命.
重心在這裡.
接著十八節就是告訴我們.
這就是我們上次說完的那裡.
接著我們就開始從這裡引申.
十七節引申到.
擔待別人的實際的實踐應用.
如何將公義和平和聖明裡面的喜樂.
這個神國的價值觀.
來作為第十五章的擔待.
我們從第十八節再看起.
十四章的十八節.
在這幾樣上服侍基督的.
就是上文所說的.
公義和平聖明裡面的喜樂.
就為神所喜悅.
又為人所稱許.
所以作為一個總結.
我們不要追求和睦的事.
所以不要將那個焦點.
放在那些次要的事情去爭拗.
應該要放在生命的合一裡面.

$^{721}$我們能夠同樣成為基督徒.
你是猶太人.
你終於明白耶穌是尼塞拉.
我們外邦人拜訪偶像.
我們做一些犯罪的事.
我們又認識你們的尼塞拉.
我就得到救恩了.
原來大家這樣得救.
你就可以想像到.
保羅說.
所以我們不要追求和睦的事.
和彼此見難不幸的事.
這一句在下文又繼續說.
就是第二節繼續說一次.
所以其實兩章聖經是一個聯繫.
十四章主要是不要去否定別人.
審判別人.
要彼此接納.
第十五章是說要怎樣擔當別人的軟弱.
所以兩者那種相輔相成.
那個重要的鑰節.
重要的鎖匙.
那個關鍵字就是十七節十四章.
就是神國度這個宏圖.
就是在乎我們的生命.
怎樣能夠體現到神國度的公義.
那個和平和睦.
和聖靈裡面的平安喜樂.
所以第十九節保羅說完之後.
不要因為食物毀壞神的功績.
記住那個食物就是大家看的製偶像的食物.
或者敬拜神的.
或者不潔之物.
或者肉帶著血不可以吃等等.
作為他們的爭鬧.
這些毀壞了的神的功績是救恩的部分.
你可以想像到兩個的對立面是很大的.
一些這麼次要的傳統.
相對於神的功績就是使到萬人得救.
神的功績就是要讓每一個人.

$^{761}$來自不同的民族的各民各族各方.
可以聚在神的永恆國度裡面.
怎會計較其他次要的東西呢?.
而且甚至乎那些是有過渡性的.
過渡了就已經沒有用了.
接著就說了.
凡物固然潔淨.
但有人因食物叫人跌倒.
就是他的罪.
這就是罵外邦人.
沒錯 保羅說.
其實神做萬物.
本來就是健康可以吃.
可以給人來享用.
因為是神做的.
除非那些是不可以吃的.
但是如果有人因為食物叫人.
叫其他外邦人.
對不起 叫猶太人跌倒.
那就是他的罪.
這裡就很清楚了.
既是叫猶太人不要論斷.
不要因為食物廢掉神的功績.
就是說猶太人.
你看看外邦人都可以得救成為神的兒女.
不只是亞伯拉罕的子孫.
更加不是因為你是守足律法.
所以成為神的兒女.
都是因為神對列祖的應許.
所以使得部分的猶太人.
都可以因信耶穌得救.
所以這就是神的功績.
所以你不要因為食物廢掉了神的功績.
去毀壞神的功績.
所以我們要小心.
我們看到不同的人信耶穌得救.
我們不要令到他跌倒.
我們不要灌跌他.
他們自己不信是另一個問題.
但是我們不要令到其他人.

$^{801}$因為我們而他不再回教會.
否定救恩 否定耶穌.
這樣是很大罪的.
所以猶太人 普羅都警告他.
普羅自己是守律法的.
上個星期我們引述過.
他有一個拿西爾人的怨在身.
當他上到耶路撒冷之後.
他馬上就去到聖地還願.
不過和他一起的一個叫特勞菲摩.
就被那些猶太人.
在安提摩上來的猶太人見到.
就遷怒普羅.
就抓他去到會堂等等.
來審問他.
所以你可以想像到.
普羅叫猶太人.
你不要毀壞神的工程.
你不要因為食物毀壞神的工程.
這樣是很重的說話.
但是猶太人他們明白了.
原來在主裡面.
真的什麼都是潔淨的.
如果你仍然有這個良心.
有這個文化.
不要緊 你照做.
只要那件事不是要和你信仰掛勾.
加勒太書就開始對付這些問題.
他們說你不行割禮就不得救了.
普羅就在這件事就很嚴重了.
因為救恩永遠只是耶穌的犧牲.
我們相信就得救.
不加任何的附件.
沒有任何的附帶條件.
所以割禮如是.
食物如是.
普羅說猶太人不要叫人.
來毀壞神的工程.
接著普羅在第20節.
又繼續說.

$^{841}$外邦人你們不要以為.
我只是說猶太人.
你們也要小心.
凡物固然潔淨.
如果有人一吃東西.
叫人潔好.
就是他的罪了.
是你的罪.
你在猶太人面前.
肆無忌憚的吃人家.
吃那些猶太人看為不潔的食物.
你這樣吃.
就帶來良心有很大的虧欠.
所以普羅說.
無論吃肉還是喝酒.
是十方別的事.
叫弟兄跌倒.
一概不作才好.
吃肉也叫弟兄跌倒.
喝酒也叫弟兄跌倒.
這個意思就是說.
在一種完全不體諒.
不接納的情況下.
這樣來令對方在那個文化.
在那個良心下.
受盡那種凌辱.
受盡那種欺凌.
而導致關係上的破裂.
彼此在一種不合一.
彼此在一種對立裡面.
這樣是神的審判.
所以這樣的飲食就是罪了.
所以我們不要自己私刑.
去做一個聚餐,聖餐.
如果你私人自己搞一個聖餐.
我們千萬不要去搞這個聖餐.
又悶又不好.
又只有曾牧師帶.
我們去另一個.
我們找人去帶.

$^{881}$我們私下在會所.
我們找最好的餅.
最好的葡萄汁.
用最好的方法去調製.
我們自己開一個.
自己一檔的聖餐.
我們秘密找人去.
千萬不要.
因為我們知道.
我們自己的聖餐是不合一的.
我們會讀到關於耶穌在十架上.
因為縱水要淹沒我.
我陷在深淤泥中.
內心的痛苦.
沒有立腳之地.
我到了深水中.
大水漫過我身.
形容自己像一種水深火熱.
情緒和心靈都很受困.
我因呼求困乏.
喉嚨發乾.
我因等候神.
眼睛失明.
無顧恨我的.
比我頭髮還多.
無理與我為仇.
要把我剪除的.
甚為強盛.
我沒有搶奪的.
要叫我償還.
神啊!我的愚昧.
你原知道.
我的罪軒不能隱瞞.
是大衛自己的自述.
當然這不是說主耶穌.
因為這些經文是告訴我們.
大衛有一種.
人生的經歷際遇.
就好像某個點.
和主彌創下.

$^{921}$有一種相似.
求你叫那等候你的.
不要因我蒙羞.
求你叫那等候你的.
不要因我蒙羞.
以色列的神啊.
求你叫那尋求你的.
不要因我受辱.
因我為你的餘顧受了辱罷.
滿面羞愧.
我的同胞看我為外路人.
我的同胞看我為外邦人.
因我為你的電心裡焦急.
如同火燒.
並且辱罵你人的辱罵.
都落在我身上.
這很特別.
頭一節.
有人觀第二章.
門徒想起耶穌.
入聖殿的時候.
最後一次在儒節上.
聖殿.
驅散那些.
賣牛羊.
賣甲子的.
拿走那些農.
不要讓神的電.
像一個賊窩.
門徒想起耶穌.
想起這節聖經.
引第二天69篇的9節.
因我為你的電心裡焦急.
如同火燒.
引第二章24節.
保羅引的不是這裡.
他引的是辱罵.
都是他都有引.
不過這裡告訴我們.
辱罵你人的辱罵.

$^{961}$都落在我身上.
你可以想像到.
對於保羅來說.
沒錯.
他沒有引前面.
只引下面.
如經上所記.
辱罵你人的辱罵.
都落在我身上.
保羅就引下半節.
約翰福音的作者.
小童約翰.
就引上半節.
基督的那種焦急.
大衛.
似乎那個意思就是.
那時候應該沒有聖殿.
大衛就是.
見到人敬拜神的態度.
那種輕佻.
那種藐視.
而他自己是這麼熱心.
這麼敬虔.
他就有一種焦急.
然後人們就.
就罵他.
我們不知道到底是.
阿撒羅他的兒子.
殺他.
還是素羅追殺他的時候.
兩次逃亡.
的光景.
那個際遇.
但無論如何.
大衛在人生體會過.
成為千夫所指的.
那種屈辱.
那種痛苦.
但是大衛都願意.
為神去忍受.

$^{1001}$這些屈辱.
就是為著.
他自己的同胞.
他的子民.
不要遠離了神.
這個是《詩篇》69篇.
這裡的經文的.
上下文的意思.
真是可圈可點的.
新的作者引經文.
約翰福音就很容易.
因為想起聖殿.
耶穌說的.
聖殿將來要被毀滅.
他見到耶穌這麼熱心入城.
這麼熱心在聖殿裡面.
趕走那些賣賣.
讓那個賣賣要消失.
因為是敬拜神.
不可以這樣去做.
因為背後又有貪污.
又有黑社會收徒弟.
所以就變成了.
門徒後來.
又想起來.
耶穌就是心裡.
焦急為神的殿.
很容易因為門徒想起來.
都沒有神學的關聯.
但保羅不同.
保羅這裡是在說.
他自己的論點.
他一直在說.
我們要在教會裡面彼此接納.
就算不同文化.
就算不同階層.
都要學習彼此相愛.
彼此的來去.
不是一種大家掩飾的.
虛偽的當看不到.

$^{1041}$而是大家在主裡面.
謀求的建立.
德行.
大家都能夠一起成長.
最終我們都是能夠.
為神的真理.
來向人.
向更多人傳福音.
一個合一的心靈.
去傳福音.
所以堅固的人.
保羅一直在說.
就是牧羊論點.
和教導.
他那時候沒有電腦.
他怎會對經文那麼熟.
他怎會只想到.
這一節聖經.
還要引誘的是下半節.
肉麻你人的肉麻.
都落在我身上.
基督也不求自己的喜悅.
基督是為神.
來到這個世界.
是要求神的喜悅.
所以他成為了眾人對他的肉麻.
他一句話都不說.
你想想.
希律王的兵兵.
抓耶穌去到希律王那裡.
希律王是很想見他.
希律王就說.
你走過神跡.
讓我看看.
耶穌是一言不答.
一句話都不說.
很藐視這個皇帝.
根本是傀儡.
真正的皇帝.
真正的審判者.

$^{1081}$是彼拉多.
但是兩個都是軟弱的.
兩個本來都是有仇的.
因為現在耶穌成為彼此的敵對.
大家成為好朋友.
但是耶穌見到這兩個.
表面好像很有力量的君王.
巡撫.
但是其實是一言不答.
你試一下.
人家真的得罪你.
你那個內心.
你都會說一些難聽的話.
你都會罵他.
你都會尋求你自己的公義.
但是耶穌基督.
祂不求自己的喜悅.
祂求誰的喜悅呢?.
祂就是求神的喜悅.
這些就是兼顧者.
這些就是兼顧人.
這些就是屬靈的偉人.
這些就是真正有一個有擔待的人.
那個心態.
他本有神的形象.
那個心態.
耶穌的想法就是這樣.
本有.
我是神來的.
但是他願意取了奴僕的形象.
成為人的樣子.
既有人的樣子.
又是自己的悲悲.
又是保羅自己作出來的.
好像詩歌一樣的題材的.
《肥勒比書》的言章.
五節的經文.
後期保羅對基督的十誡.
越來越認識.
他就完全明白.

$^{1121}$這位兼顧者.
那個生命的想法.
原來是這樣的.
所以原來基督徒為什麼能夠成為兼顧者呢?.
是我們的生命的想法.
不斷的更新.
不斷有主的形象.
不斷我們經歷了很多不同的人生.
我們越來越認識.
唯有這位主.
永不改變.
唯有這位主.
真是值得我們一生去跟隨.
去服侍.
這樣就成為兼顧人.
不是你坐在高位上就是兼顧人.
是呀.
甚至乎那些條例令你很兼顧.
那些開會的議程.
你設計的東西.
都令你很兼顧都可以.
如果要玩那些遊戲規則.
很多人都會玩.
這些政治手段.
或者這些辦公室的開會文化.
甚至乎那些選舉各方面.
你可以暗中有什麼票都可以.
不是這些兼顧.
真正的兼顧.
是當代軟弱者.
基督不尋求自己的喜悅.
真正的兼顧.
是當代軟弱者.
基督不尋求自己的喜悅.
所以這些經文.
本來69篇是講就坐的.
就坐的魔鬼.
魔鬼的化身.
就是猶大.
沿他的筵席.

$^{1161}$成為網路.
成為一個.
捉鳥的機關.
成為一個捉鳥的機關.
猶大就是這樣.
跌落永遠不能悔改的陷阱裡.
被那個陷阱捉了.
這樣.
69篇後面的經文.
本來將這節聖經放在基督裡.
基督就是這樣.
願其他人的就坐.
都落到我身上.
我願意擔待.
所以頂梓櫃.
羅馬書.
來到這裡才是真正的.
中心思想.
原來我們明白救恩之後.
我們要成為一個.
有擔待的人.
這裡才是保羅.
要將他最精髓的.
他為甚麼要周圍去.
傳福音.
向人說耶穌基督.
他好像不介意人怎樣看他.
原來他是有擔待.
他願意學校耶穌.
他就是一個有擔待的人.
所以.
今天我們不是太需要.
太多人.
來做事奉.
我們需要的.
應該相反的說.
今天我們找不到人事奉.
個個都沒空.
我們很擔心.
希望有人吃了維德納.

$^{1201}$其實我們不用擔心這些事.
我們擔心教會.
有沒有擔待的人.
我們的領袖是否有擔待的人.
有擔待就很不同了.
在這個生命裡面.
你就會成為別人的一個幫助和力量.
所以保羅在這裡引述的是耶穌基督.
接著就是四節.
從前所寫的聖經.
當然最直接的.
上文的應驗就是69篇.
就是說.
我引述這一節聖經.
讓你思想一下.
羅馬的教會.
你看看我一直說的救恩.
恩信稱義.
世人都犯了罪.
現在因為基督耶穌的救贖.
我們現在可以白白稱義.
這些救藥的經文.
關乎到救藥的獻祭.
關乎到阿巴拉罕神和他納的藥.
全部都是.
圍著救恩來到這個世界.
現在救恩已經完成.
我保羅就是這樣傳這個救恩.
所以保羅現在在這裡.
做一個很宏觀的告訴我們.
從前所寫的聖經.
當然從前不是他寫的.
是神密士不同的.
歷史上的不同的.
先知來寫成的.
無論創世記.
一直到馬拉基書.
這些救藥聖經.
告訴我們.
當然直接就是69篇.

$^{1241}$當然可以再宏觀一些.
他說從前所寫的聖經.
都是為教訓我們寫的.
原來聖經.
是歷史的記載.
不是讀完之後.
那些過去的不關我們事的.
我們是新約的.
我們不用看舊約的.
那些很悶的猶太人看.
不是的 圍著我們寫的.
保羅這樣說.
是教訓我們寫的.
叫我們因聖經所生的忍耐和安慰.
這裡告訴我們.
Whatever was written.
in earlier times.
was written for our instruction.
是為了要令到我們.
得到那個教訓.
得到那個生命的教育.
來成長.
So that.
through this perseverance.
堅忍的忍耐.
和那個安慰.
encouragement.
由聖經而帶來的.
忍耐和安慰.
以致得著盼望.
原來忍耐和安慰.
忍耐.
就是對別人的包容.
接納.
就是因為我們對聖經的真理.
明白之後.
我們就不會認為.
原來安息日.
我們外邦人都要守了.
我們要好像成為安息日會.

$^{1281}$禮拜日就不准進會.
一定要安息日才進會.
我們就不會這樣了.
因為我們已經明白了聖經的真理.
舊約只不過是一個示範.
告訴我們.
在舊約的世界.
因為救主還沒來這個世界.
需要這樣的規範.
但是當救主來.
那個真正的屬靈生命.
那個永恆真理已經顯現了.
所以我們明白.
凡物都是可以吃.
凡物都是潔淨.
所以猶太人也不需要再拘泥於.
他們自己的傳統.
所以這一切.
都是為了讓我們忍耐和安慰.
原來這裡告訴我們.
基督徒.
原來我們可以帶來.
堅強的對軟弱者的包容.
和等候.
這個就是.
那個忍耐.
弱者對堅強者的信任.
而不是論斷.
這個就是接納.
這個就是一個彼此的同心.
面對誤解.
批評的時候.
效法基督.
不求自己的喜悅.
如果那些人誤會.
都不需要.
為了自己所謂的.
清白.
來去拘泥於真.
甚至反目成仇.

$^{1321}$甚至立即.
從此離開教會.
不回到這個教會.
其實不需要的.
我們要學習在那個誤解裡面.
我們尋求的.
學習的是求基督的喜悅.
所以這是一連串的經文.
是有那個的教導成份.
是保羅告訴我們.
我們必須要學習成為堅固的人.
堅固的人.
就是要擔待.
不堅固人的軟弱.
就是說.
對他有接納.
對他現在未能成長.
都有一個忍耐.
接納背後就是有忍耐.
就是慢慢讓他繼續在.
真道上面.
來去成長.
長大.
這個就是保羅在這段經文裡面.
很重要的意思.
所以原來聖經.
我們一路在神的道裡面.
我們在神的愛裡面.
我們就是可以.
不斷在忍耐.
神對我們的忍耐.
我們能夠成長.
神對我們的安慰.
讓我們能夠得到那個.
心靈的圓滿.
我們的內心.
不會有一種的.
懼怕.
或者有一種不同的.
我不懂怎樣形容.

$^{1361}$我有時要面對一些.
情緒病的朋友.
接觸.
因為現在教會.
有一個小朋友.
幼稚園.
畢業之後.
升上小一.
我上星期好像都有講過.
三個月沒有笑過.
你想想.
是一件很困難的事.
所以神的愛.
就會令到我們的心靈裡面.
有一種的.
平安.
有一種的安慰.
安慰是很重要的.
所以就是這個.
安慰是這個字.
這個Paracletos.
安慰就是保維斯.
那個字.
就是聖靈.
就是保維斯.
所以安慰者.
這裡告訴我們.
十五章四字就是聖靈這個字.
這個安慰斯.
保維斯這個字.
但願似忍耐安慰的神.
叫你們彼此同心.
效法基督耶穌.
一心一口榮耀神.
我們主耶穌基督的父.
保羅.
不單止勸勉.
加上祂的祝禱.
祂的希望.
羅馬教會.

$^{1401}$看到我的這封信.
其實我已經在哥林多.
為你們禱告了.
但願這位似安慰.
忍耐的神.
叫你們.
最終是什麼呢?.
就是彼此同心.
所以堅固的人.
如果我們去不斷學習.
體諒其他人.
不斷為其他人軟弱.
我們就能夠帶來彼此的同心.
因為堅固的人.
他的心智.
就是基督的心智.
就是基督的心靈情操.
所以.
效法基督耶穌.
就能夠一心一口榮耀神.
當我們大家都在主裡.
就能夠一心一口榮耀神.
所以.
真正和諧的歌聲.
是大家在愛裡.
還在唱神的詩歌.
不是唱的時候.
你又想著對方.
對方又想著你.
這樣的不和諧.
其實是貌合神理.
連那些不信人都懂得這樣看.
所以我們.
屬主的元年女.
一定要學習.
生命的成長.
但是這五節.
六節的經文.
已經是相當寶貴.
告訴我們教會的倫理關係.

$^{1441}$必須學習基督的忍耐.
必須學習基督的安慰.
必須學習基督不求自己的喜悅.
這樣的生命.
我們繼續看多一點.
所以一到第六節.
剛才我講的總結.
從不要叛賭的人.
已經提升到一個地步.
就是主動去擔待.
其他人的軟弱.
不叛賭就是一個負面.
我已經沒辦法再忍你了.
但是為了耶穌.
我忍你.
但是對不起.
以後你走你的路.
就不是這樣.
不叛賭但是提升到.
主動去擔待.
這些就是我們剛才講的.
第八節開始了.
主耶穌.
保羅再忍幾段的舊約.
在這四節聖經.
忍了四節的舊約聖經.
我們可以逐節去看.
看我們的時間.
還有多少時間我們可以看下去.
所以再一次一個Den of all.
再一次告訴我們.
同心和接納.
好像基督接納你們一樣.
使到榮耀歸於神.
所以彼此接納才能榮耀歸於神.
才是一個整個的群體.
在愛裡面歸榮耀給神.
所以唱Hallelujah.
唱Hallelujah的時候.
我們的心裡面真的要求主.

$^{1481}$給我們很坦然.
我們記住我們一定要做堅固人.
我們要常常想別人的需要.
我們不知不覺就變成.
對付我們自己的自私.
我們自己的自我中心.
我們在這個時候.
我們就能夠彼此的接納.
所以保羅在這裡.
一而再說的那些說話.
要同心要接納.
然後重點就是.
如同基督接納我們一樣.
如同基督怎樣接納我們呢.
我們大家想一想.
我們是很卑微的.
我們是很虛偽的.
我們是卑微的時候.
我們是犯罪的時候.
我們是無望.
我們都不知道什麼叫永死的時候.
我們更加不會羨慕有天堂.
但是當我們真的明白了.
主耶穌為我們釘十字架的時候.
我們知道原來我們的生命是屬祂.
我們整個人生是祂設計.
是祂安排.
是祂讓我們能夠有今天.
我們生在香港.
我們不是生在格陵蘭島.
我們是香港人.
中國人.
一切有設計的.
所以好像基督這樣接納我們.
所以我們要反過來.
保羅的意思就是說.
給錢接納.
不是說勸你要接納.
你不接納是你的事.
基督都接納你.

$^{1521}$你不接納其他人.
你為了很小的事.
不接納其他人.
你為了他吃的那些祭物.
你就不接納他.
你為了那些外邦人吃的蝦蟹.
你就不接納他.
這樣毀壞神的功情.
你是誰?你是神.
基督都接納我和你.
你是猶太人.
你都沒有手足的律法.
耶穌都愛你.
為你釘十字架 就是保羅.
外邦人更加有基督的接納.
所以保羅就在這裡告訴我們.
使榮耀歸於神.
所以和大英哲會一起.
敬拜神.
衷心的榮耀神.
其實是最美麗的詩歌.
最美麗的畫面.
大家不是貌合神離.
大家是真誠的.
在愛裡合一.
這樣的敬拜.
是真的很美麗.
相信天使會看到.
我們這樣的和諧.
那種敬拜.
記住不是傻乎乎.
甚麼都不懂的.
好像一些有弱智的問題的人.
不是這樣的情況.
而是我們很有意識地.
和我們可能不喜歡的人.
一起和諧的.
唱頌聖詩的敬拜神.
保羅接著.
你看他真是仔細到一個地步.

$^{1561}$他引四節的聖經.
講到神的功情.
救恩的功情.
就是怎樣拯救猶太人.
叫外邦人都無恩.
第八節.
我說.
基督是為神的真理.
作了受割禮人的執事.
即是猶太人.
要證實應許列祖的話.
即是十一章.
因為這一棵.
斬下來的野橄欖.
其實那個本說.
就是亞伯拉罕.
猶太人都沒有.
猶太人都沒有被.
永遠的拒絕.
是因為神對列祖的應許.
即是神是不會廢棄.
祂自己的約.
所以原來基督是為神的真理.
即是神的計劃.
是要將祂自己的子民.
亞伯拉罕的子孫.
是要忘福的.
神要兌現這個應許.
所以使得有部份的猶太人.
一直都有信心.
所以耶穌基督.
就是受割禮人的執事.
保羅很想這樣稱呼自己.
祂都是愛邦人的僕役.
要證實應許列祖的話.
兌現神的承諾.
神的保護.
神的愛.
所以愛邦人.
好像沒有份.

$^{1601}$這就不是.
保羅在這裡第一句聖經.
如經上所記.
因此我要.
對不起.
第九節才是.
這裡沒有引述.
明白了.
第九節.
所以猶太人就是神的應許.
外邦人.
外邦人即是我們.
因為甚麼.
從來沒有應許過外邦人.
今次是因為甚麼.
因為年敏.
所以我們很少看完我們是蒙年敏的罪人.
我們是蒙年敏的聖徒.
我們是蒙年敏的神的兒女.
記住我們一生.
我們是蒙主的年敏.
才能成為基督徒.
我們不是比其他人好.
所以如果我們明白了年敏這個概念.
我們就很感恩.
我們為神不講教.
不計較.
神要我們做的事.
我們就盡力去做.
盡心去做.
因為我們是蒙年敏.
你看過這段片.
有些人被人提攜.
滿心歡喜.
一生為了報恩.
這就是年敏.
所以我們這裡.
保羅用的字眼是增長過的.
我們外邦人.
現在快要到AI的年代.

$^{1641}$快要有機械人.
有那些人工智慧的人.
但是我們唯有人.
我們是罪人.
我們才明白什麼是年敏.
所以無論我們將來怎樣飛黃騰達.
我們在神的國度裡面.
永遠要記住我們是蒙年敏的.
一個因免的人.
是因免改變我們.
所以我們可以榮耀神.
保羅想拉兩班人一起.
你看保羅的心是多麼的廣闊.
多麼的大.
將人拉在一起.
還要不是說兩個人.
是整個世界.
猶太人,外邦人.
即是所有全個地球的人.
這樣來到他.
拉在一起.
是要榮耀神.
他就是一心一口榮耀神.
所以這就是神的救恩計劃.
所以在這個時候.
你想想.
如果猶太人還是那麼計較.
為什麼外邦人那麼不守傳統.
為什麼外邦人吃東西原來那麼骯髒.
那些不潔的食物不能夠吃.
如果他明白神這個工程.
神用了外邦人.
忘他的年敏.
來榮耀神.
所以第九節的經文就引述.
這裡應該是《撒慕爾記》.
或者十八篇的詩篇.
如經上所記.
因此我要在外邦中稱讚你.
歌頌你英名.

$^{1681}$遊說.
再引多一句.
為猶太人引了兩句.
《申命記》的三十二章.
你們泰邦人.
當遇主的百姓.
一同歡樂.
是沒得說的.
《舊約聖經》已經有這樣的教導.
有這樣的預言.
所以猶太人現在在羅馬的教會.
每個星期聚會.
看著那些外邦人.
越看越生氣.
根本就不是聚會.
根本就不是來敬拜神.
給傳統.
給自己的良心來拖住.
保羅說.
舊約都已經有這樣的應許.
舊約已經有這樣的藍圖.
舊約摩西.
未入埃及.
對不起.
未入迦南.
已經有這樣的預言.
你們外邦人.
當遇主的百姓.
猶太人.
一同歡樂.
所以原來神的偉大的功情.
就是要人與人在愛裡合一.
所以教會是更加應該謀求合一的場所.
不是帶來分裂的地方.
很少的東西就分裂.
你說多麼可惜.
可想而知.
對羅馬書.
保羅這樣.
好像哭斷心牆.

$^{1721}$這樣來教導.
用兩張聖經.
彼此接納.
以及要成為兼顧人.
這樣來學習.
時間不容許.
只能夠講到這裡.
因為我們下面還有一個大段落.
保羅再一次講他自己內心的.
他自己就是外邦人的執事.
是福音使者的兼顧人的心態.
所以我們先停在這裡.
一起同心禱告.
天父啊.
我們感受到十四章十七節.
是何等的重要.
求你幫助我們.
我們去追求神的永恆國度.
真的要將公義和平.
以及聖名裡面的喜樂.
成為我們的生命的元素.
讓我們能夠不斷地長大.
原來保羅很想每一個基督徒.
都是一個兼顧的人.
但當然他知道.
有些人是在不同的程度.
所以有些人在他們某些的情況.
都會是軟弱的人.
但是不等於他們不應該在教會.
反而兼顧的人就是要接納包容軟弱的人.
多謝主給我們有很清楚的教導.
多謝主給我們看到.
原來我們一生要追求.
在很短暫我們要見主面之前.
我們要追求成為兼顧的人.
不是要炫耀.
不是要驕傲.
而是真的有基督耶穌的這樣的思考.
就是不尋求自己的喜悅.
要來到我們擔待別人的需要.

$^{1761}$求主幫助.
再一次多謝你.
下個星期我們去看看.
這位學習基督的保羅.
是怎樣將他的兼顧.
他對外邦人需要福音的重擔.
他要來到擔待.
願神幫助我們.
我們這樣禱告奉耶穌基督的名.
阿們.
\newpage



\section{羅馬書}
\label{sec:UU4dvSZbobM}
\textbf{2026-01-23 周五晚查經班《羅馬書(三十五):分題:有盼望,有異象,有擔當的人》 講員:曾裔貴牧師}
\newline
\newline
連結: \href{https://youtube.com/watch?v=UU4dvSZbobM}{\texttt{https://youtube.com/watch?v=UU4dvSZbobM}} ~~~~ 語音日期: 2026-01-24
\newline
\newline
\hyperref[sec:pe8biigoai0]{\small{< < < PREV SERMON < < <}}
~
\hyperref[sec:index_chronic]{\small{[返順時目]}}
~
\hyperref[sec:index_scriptual]{\small{[返順卷目]}}
~
\hyperref[sec:LvE6qlAUC04]{\small{> > > NEXT SERMON > > >}}
\newline
\newline
$^{1}$好 第一節 晚安.
感恩我們又可以有晚上的茶經班.
我們一起先禱告.
天父我們將今天的茶經班仰望您.
我們來到第十五章十三節.
我們這一位主是那一位使人有盼望的神.
幫助我們都能夠真的掌握.
保羅在第十四,十五章.
在教會,多元文化,不同的民族.
怎樣能夠因為福音,因為耶穌基督.
而來到有一個共融.
一切的同心.
可以來到去圍著保羅的宣教神的國度.
這樣來到擴展.
今晚我們再仔細去看.
十三節之後的經文.
求主幫助我們又可以看到.
使徒保羅的一些相當具體的行程的分享.
又讓我們能夠對使徒他的傳道的職涯.
又有一定的了解.
多謝你讓我們有他的分享.
留到今日.
我們在不同的世界的教會都能夠可以讀到.
我們多謝主.
求你帶領我們這樣禱告.
奉耶穌基督的名.
阿們.
十五章很清楚是緊接十四章.
不要批判之餘.
要作為一個堅固的人.
堅固當然就是內心是成熟,是成長.
內心裡面有一個同理心.
是有一個為別人設想.
所以堅固的人就是一個會留心.
會敏銳別人需要.
所以堅固的人就不是自己去很堅.
自己就將自己的自由放得相當大.
而忘記了可能有一些所謂弱勢.
就是相對在信心上.
未能夠徒步這樣去成長的.

$^{41}$總是有一些可能執著.
可能是一些迷思.
有一些包袱.
有一些傳統.
所以保羅,猶太人和外邦人.
他們很明顯的差異.
關乎到一些食物的差異.
一些節旗的差異.
明顯這樣放在一起是一個衝突.
但是正正就是你不能避免.
怎樣能夠在多元的文化體系下.
而能夠謀求到一個合一呢?.
相信就是在心靈上有這種保羅的教導.
就是我們自認為堅固的人.
就要擔待 所以是有擔當.
一個成熟的人是應該在教會裡面有擔當.
是應該看這個是屬靈的家.
看其他弟兄姊妹都是自己的靈魂的骨肉之親.
彼此有一個欣賞,接待.
而不是用一種批評,論斷.
彼此有一種從心靈裡面開始對立.
去到白熱化,去到表面的都滅了.
這樣就不是聖經的教導.
有時相對於羅馬帝國.
因為明刀明槍.
因為說不准信耶穌,不准拜真神.
要拜其他的假神.
相對來說容易一些.
拒絕,不接受 但是自己的弟兄姊妹.
大家都是來自不同的傳統,文化.
就容易有時這些很小的東西.
就引發很大的衝突.
其實真的相對不是真理來講.
是相當少的 所以十四十五章.
都是講一些很次要.
跟真理完全拉不上關係的.
是大家的不同的做法.
對一件事情的那個研判.
那些不影響得救和永恆,永生.
以及對神的敬拜.

$^{81}$所以有時這些才是難堪的.
這些才容易引發那種分裂,爭論.
所以這兩章聖經其實真的很寶貴.
《仕途保羅》講了八章關於救恩.
怎樣能夠獲得.
甚麼人可以白白地得救.
但是去到十四十五章.
就用了兩章來講到.
怎樣能夠對彼此的不同的差異.
有一個接納呢?.
所以很寶貴.
可以更加反面地印證.
《仕途保羅》它不斷地衝去傳福音,宣教.
但是它都會留意教會怎樣能夠處理.
大家在不同的文化.
一層聚會之下.
怎樣能夠有一個彼此的和諧.
那個大原則相信就是這個擔待.
就是有主耶穌的成熟擔當.
他那個擺下自己.
所以我們來到第七節.
重新看第二段落.
我們今天其實集中在十三節打後.
《仕途保羅》在這裡有一個祈禱.
去到十五章的結尾.
又再包一個禱告.
所以《仕途保羅》它都是很想在.
講完一些小的段落之前.
很快就要為弟兄姊妹的合一.
彼此的同心來禱告.
我們看看第七節開始.
我的課堂.
對上一次有些我們要留意的.
這個是四章的十七節.
就是關乎到神的國度.
十四章的十七節.
不是四章.
十四章的十七節.
那個就是因為神的國不是在乎飲食.
是在講那些關乎禮儀.

$^{121}$關乎到猶太人的律例.
那些儀式等等.
所以那個飲食不是普通的自助餐.
或者普通的食物那麼簡單.
而是關乎到這些禮儀.
這些傳統.
這些多年的文化傳統.
在保羅的眼中.
耶穌來了這個世界.
所有東西都要更新.
所有東西都要用一個新的眼光去看.
所以猶太人要慢慢明白.
就算你做足了那些飲食條例.
如果你沒有了耶穌.
沒有了救恩.
一切飲食在神的眼中都是不行的.
所以神的國是在乎內在生命的更新.
所以神的國是在乎公義和平.
聖靈裡面的起來.
要在這幾樣上面去服侍基督.
所以十四章的十七節.
是一個很重要的整全的段落.
然後十五章.
就是如何能夠在實際的生活.
讓這個那麼抽象的神的國.
我們追求在實際的生活聚會裡面.
人與人之間的相處上.
展現神的國度的空襟和內在生命.
所以可以說十四章的十七節.
是非常重要的鑰節.
大家如果有時間真的要背熟.
《上眼文》和分析.
神的國度就是神國度的價值.
基督徒就要擔待有這個責任在我們的生命裡面.
當我們十五章做一個有擔待的人.
我們就會明白.
原來我們正在追求神的國度.
那個屬靈的價值觀.
所以這告訴我們.
神的國重視的是和平.

$^{161}$以及造就別人.
所以強者是有責任不堅持自己的自由.
強者證明就是他對真理應該有一個融會貫通.
應該對人有一個包容,接納.
因為他更深的認識.
就是因為基督接納他這個罪人.
以致我們能夠不斷的認識真理,成長.
以致我們成為一個神眼中喜悅的屬靈人.
我們其中一個很重要的指標.
一個屬靈人,堅固人.
就是不堅持自己的自由.
願意因為別人的緣故來放下.
以致不求自己的喜悅.
求基督得到榮耀.
所以十五章就是這裡的意思.
這些我們上一個星期講了.
幫大家溫習一點.
所以擔待就是不求自己的喜悅.
終極的盼望的榜樣就是學校基督.
我們看看第七節開始.
這個十五章四節.
這些我們講過了.
這些我們提了.
就開始關乎到保羅的祈禱之前.
就是第七節開始.
「所以你們要彼此接納」.
這個就是上文耶穌的榜樣.
因為耶穌都是一心一口榮耀神.
所以我們就要來到.
好像保羅的教導一樣.
彼此能夠接納.
就是說無論自認強者.
還是弱者.
都要彼此接納.
如同基督接納你們一樣.
使到榮耀歸於神.
基督是如何接納我們.
拯救我們.
知道我們的軟弱.
我們的犯罪.

$^{201}$我們對神的軟弱.
我們的拒絕.
但是基督都慢慢引導我們.
明白,福音,得救.
所以我們對其他人.
都要有這種心靈裡面的接納.
這個是很不容易學習的.
因為要從心裡面接納.
首先你就要有一種.
不要有色眼鏡.
不要有自己的先人為主.
更加不要有一個自己很難除去的.
那個成見.
那個對他的觀感.
很難你當作一個新的人.
什麼包袱都沒有.
其實很難的.
所以我們就要不斷在.
基督接納我們的這個思想底下.
這個屬靈的操練底下.
我們就去學習接納其他人.
我們應該要看到.
如果我們有永恆的眼光.
去看到我們自己在神面前.
那個本相,內心.
我們就明白這個接納是多麼偉大.
是多麼無條件.
是多麼愛的深度.
這個時候我們就要用這種愛.
來接納其他人.
當然不是說我們一定要跟他一樣的意見.
可以是很不同的意見.
但是要接納.
所以這個功課是一生要學習.
當然如果我們經常都有這個心態.
我們常常都是被主接納的.
救恩的愛的感動.
我們相對對其他人的接納.
又容易放下.
又容易學得好.

$^{241}$又容易伸出我們的聖靈.
來接納別人.
如果我們一直都從我們自己的角度.
去看那個對錯.
內心那把尺永遠不夠.
基督的接納來得夠寬大.
所以這些都是大家一生要學習的功課.
尤其是我們要有效去傳福音.
去經歷聖靈的同在.
這個就是必須我們內心不要有一個偏見.
不要有一種很重的內心的成見和包袱.
第八節.
我說基督是為神真理.
作了受割禮人的執事.
要證實所應許列祖的話.
並叫外邦人.
這個第八節就是說到基督是擔待神的選民.
為著神的應許能夠兌現.
能夠成就.
因為基督的犧牲.
整個舊約的猶太.
摩西的子孫.
亞伯拉罕的子孫.
以色列屬神的子民.
都能夠被接納被拯救.
所以有敬虔的.
真的有那個屬靈的眼光的猶太人.
就可以得救.
例如保羅就是一個明顯的例子.
一個本來被猶太教.
本來是一個很好的律法.
能夠去經歷神的.
但是誰知道.
當慢慢成為了一個枷鎖.
成為了一個包袱的時候.
連保羅自己都被這個包袱包圍著.
他成為了一個律法的專家.
不過就是去對其他的.
尤其是外邦人.
來攻擊和詆毀這位造物主.

$^{281}$這位耶穌.
所以原來在基督裡面.
他的出現原來是為了成就神.
過去在舊約向列祖的應許.
所以基督就是對猶太人有一個擔待.
有一個接納.
有一個犧牲.
所以這裡保羅就說了.
我們為什麼要學習基督呢?.
正是因為耶穌基督.
是為了神的永恆救恩能夠兌現.
他成為真的用他的生命.
去擔待人類的罪.
去用他的生命.
來使得神的應許.
可以真是生效和兌現.
第九節就說外邦人了.
所以你看到保羅的眼光.
他的心是這麼平衡的.
原來有基督.
對基督有一個這麼全面的認識.
他再不是站在一個有偏見的.
有偏見的立場.
他站在一個使徒的身份.
去看這個救恩帶來的.
對猶太人原來是一個應許的兌現.
是神沒有改變他對猶太人的恩典約.
那個滿足.
但是對外邦人.
保羅告訴我們.
原來神從來沒有偏待猶太人.
而且外邦人都在他的救恩裡面.
所以保羅其實是很詫異神的偉大.
因為保羅自己越來越感受到.
他成為外邦人的使徒的身份的尊貴.
一個拉比.
一個快利塞人.
一個這麼外邦人.
隨隨便便就說自己可以信耶穌.
就可以得到神的永恆救恩.

$^{321}$對保羅來說是很難接受的.
但是當他明白了.
耶穌原來在十字架的救恩.
原來外邦人都有份的.
原來萬民都可以因為耶穌而得救的.
而清義的.
而成為亞伯拉罕的屬靈的子孫.
所以保羅其實很詫異.
很讚嘆神的永恆計劃.
我們讀過的十一章.
第八章都看到保羅對永恆救恩的讚嘆.
對做物主的智慧.
救恩的設計.
耶穌的成就帶來的讚嘆.
留意第九節.
「並叫外邦人因他的憐憫榮耀神」.
這個就是我們外邦人.
保羅一而再用的形容詞.
就是神的諸般憐憫.
使得我們外邦人能夠得救.
所以用這個角度.
猶太人不要自恃自己是神的選民.
都是無因的.
外邦人我們份外感受到.
本來不配.
本來是野橄欖的例子.
我們現在能夠被接上.
封城的那個本源.
就是神的國度裡面.
亞伯拉罕的源頭.
我們就應該要感受到.
所以任何基督徒應該要感受到憐憫.
初信的時候不覺得不要緊.
因為有些人可能慢慢不同的角度.
不同的切入點信主.
有些人可能只是一些例子.
一些見證嘉賓.
一些感動位等等.
但是每個基督徒要自己的生命穩固.
必須感受到神對你對我的憐憫.

$^{361}$我們是蒙憐憫的.
這個蒙憐憫不是施捨.
是神很大永恆的愛.
對我們每一個.
所以每一個基督徒如果越成長.
越感受到自己那個身份的尊貴.
做物主.
創造永恆的真神.
來將我們拯救.
是用他珠官的憐憫.
珠觀的憐憫.
就是說這個憐憫是可以涵蓋任何一個人.
就算我們感覺自己是多麼差.
多麼地底泥也好.
在神的憐憫之下.
我們都可以有一個身份的改換.
所以這個是很重要的.
所以是因為神的憐憫去榮耀神.
接著保羅引聖經.
如經上所記.
因此我要在愛邦中稱讚你.
歌頌你的名.
這個明顯就是大衛他在詩篇的十八篇.
對神的讚頌.
就是說那些不是奉國禮的.
當時他的愛邦人.
其實就是迦南那些的外族.
他沒有想到有中國人.
有美國人.
有歐洲人.
有保羅的時候的希臘人.
因為那時候去不到那麼遠.
但是在永恆的裡面.
神已經感動大衛.
提到將來愛邦人都要來稱讚主.
接著就是十節.
你們愛邦人.
當與主的百姓一同歡樂.
這是摩西他的三十字篇的詩歌.
來講到神的救恩.

$^{401}$在這班選民.
他們準備入迦南.
將來他們成為一個示範單位.
叫萬民都和猶太人一樣的去敬拜神.
一同和神的選民來歡樂.
接著就是十一節.
又說愛邦你們當讚美主.
萬民你們都當頌讚他.
這是117篇詩篇.
接著還有十二節.
又有以賽亞說.
將來有耶西的根.
就是那興起來要治理愛邦的.
愛邦人要仰望他.
不如我們就揭一揭.
以賽亞書的十一章.
這篇是以賽亞的詩歌.
以賽亞書的十一章.
我們可以多看一點這段經文.
也是重要的.
他說從耶西的本.
本字就是樹等.
就是斬了樹幹留下的等.
打等的等.
就是樹等.
木字斬了上面的頭就是等.
樹等必發一條.
從他更新的枝子必結果實.
而和華的靈必住在他身上.
這裡提到的是將來的一位以賽亞.
就是使他有智慧和聰明的靈.
謀略和能力的靈.
知識和敬畏耶和華的靈.
他必以敬畏耶和華為樂.
行審判.
這裡說不憑難見.
斷是非也不憑耳聞.
卻要以公義審判貧窮人.
以正直判斷世上的謙卑人.
以口中的象擊打世界.

$^{441}$以嘴裡的氣殺戮惡人.
當然不是說放飛劍等等那些.
而是說他的教導.
他的公義的言行.
公義必當他的腰帶.
腰帶就是束腰的.
讓一個有力的公義的人走出來.
自然讓人感覺到有威信.
因為他行事是公義正直.
這些我想我們大家都明白的.
這裡告訴我們.
信實必當他脅下的帶子.
脅下的帶子可能就是綁著胸口.
我不知道古代人怎麼穿衣服.
因為都是寬鬆的袍.
所以可能這個描述都是讓他能夠行事為人.
讓人看到他的誠信.
他的公義.
他的真實.
好像最近那個住天津的代辦主任.
在炫耀原來是貪污的.
那就不是誠信真實.
如果是事實的話.
那就真的犯了法.
因為公報員一輪要申報.
上司不知道你做了食用了就不行.
接著我們看下面.
我們可能都要看快一點.
我們看回後面.
保羅所說的關乎到這一段.
外邦人要稱讚就是第十一章的第十節.
我們看看.
第十節.
到那日耶西的根納作萬民的大旗.
外邦人必尋求他.
他安息之所大有榮耀.
當那日主必二次伸手.
救回自己百姓中所餘剩的.
就是在亞述,埃及,巴提羅,古實,.
以蘭,斯拿,哈馬並眾海島所剩下的.

$^{481}$這裡提到將來會有一個日子.
外邦的人要尋求神.
要來看見這位耶西之根.
成為了一個很重要的中心點.
因為萬民的大旗豎起了.
就是說萬民都要留歸管治.
這個君王的那個客管的國度.
這個是萬民將來.
這裡是不是在說將來的千禧年.
耶路撒冷的地上政權呢?.
如果字面上就是了.
就是耶西之根將來就是這樣來到去.
納作萬民要留向,要歸向的那個所在.
這裡再一次出現.
外邦人要尋求,要歡呼.
保羅就直接將這一段改為.
外邦人都要來向主歡呼.
我們回到羅馬書這一段.
外邦人就要仰望神.
所以第十三節就是總結了兩個的民族.
猶太人就是泛指是神的選民.
與神有摩西之約的關係的選民.
無論保羅的年代.
或者耶穌的年代等等.
另外保羅引述了四段的舊約的經文.
來講述原來在舊約已經有萬民可以得救.
萬民在救恩神的計劃裡面的記載.
這樣的寓言.
就是說其實很簡單.
保羅的意思就是說.
大家其實都是蒙恩.
大家都在救恩裡面.
大家都是在耶穌的拯救裡面.
來成為教會,成為弟兄姊妹.
所以就算文化,就算傳統有多大的差異也好.
在救恩裡面.
大家應該看這個是最重要的合一的基礎.
所以能夠合一不是我們去促成的.
是耶穌的救恩使我們已經合一了.
但是當我們去實踐這個合一的時候.

$^{521}$我們需要學習彼此的接納.
這個就是保羅在十三節的祈禱.
第七節的彼此的接納.
所以保羅是很透徹地引聖經.
你想想引了四段.
就是說猶太人你不要裝了.
神在舊約,在詩篇,在申命記,在《爾賽亞書》.
就已經這樣來到去寓言的時候.
就是說神的心胸.
比猶太人,比那些律法的條例更加的寬廣,廣大.
因為律法都是神頒佈的.
不過人自己越來越窄.
以為只是猶太人的專利.
其實就不是.
所以保羅這樣引述這麼大段的經文.
講愛邦人無因.
講愛邦人將來要投靠神.
所以現在猶太人看到的.
可能就是來自不同的歐洲的民族.
有帖撒羅尼迦加,有希臘,馬其頓,有比利亞.
有亞洲的有特比.
提安泰就是特比人.
有路斯德.
這些不同的城市的歐洲,亞洲的這些人.
這樣來到去在羅馬的教會.
輾轉猶太人又在羅馬的教會一起聚會.
我們就分外要學習彼此的接納,彼此的相愛.
所以這是一個很寶貴的教會合一的教導.
倫理的勉勵和關係.
重點仍然不是有相親相愛這麼簡單.
而是大家要在一個合一裡面.
去支持萬民的宣教.
這才是羅馬書最重要的大前提.
接著十三節這個經文.
我們一起要去研究一下裡面的意思.
十三節.
這裡說了.
古羅尼亞四段的《舊約聖經》.
詩篇十八篇.
今晚我會給大家一個重點.

$^{561}$所以大家未必需要抄.
證明愛邦人也要讚美神.
和猶太人同蒙連吻.
所以詩篇有117篇.
以塞爾書的十一章.
都是這樣來印證了神對猶太人的接納.
和對愛邦人的憐憫.
所以來到這裡的時候.
七個第十三節.
就是這個給錢接納.
張倫理的勸勉.
就是這個接納.
奠基在舊恩的歷史裡面.
既然神接納愛邦人.
問題就是.
為什麼猶太人不接納愛邦人呢?.
教會怎麼能夠分裂呢?.
十三節.
就是彼此的同心.
但願使人有盼望的神.
因順將諸般的喜樂平安.
充滿你們的心.
使你們藉著聖靈的能力.
大有盼望.
我們可以.
這裡我引用中文.
接著我們看看這一節聖經.
很重要.
就是原來盼望的源頭.
不是那個環境.
不是我們認為好的環境我們就能夠.
但其實不是.
真正的源頭是來自神.
這裡的希臘文很特別.
小小特別.
因為我已經不太熟悉了.
給大家看看.
盼望之神旁邊的雞腸.
其實是希臘文.
就是告訴我們.

$^{601}$很特別的.
這個O.
這個Fios.
Fios就是神.
下面那個就是盼望的字.
很特別就是.
這一位神.
盼望就是這個意思.
所以盼望之神.
就是這一位神.
盼望.
盼望形容他.
就是說神的屬性.
這次保羅強調.
他不是愛.
他是盼望.
所以嚴格來說.
不是使人有盼望.
而是如果原文就是盼望之神.
因為你們的信.
就將諸般的喜樂平安.
充滿你們的心靈.
所以原來不是但願使人有盼望.
嚴格來說就是盼望的神.
因為你們的信.
就將諸般的喜樂平安.
充滿你們的心.
但原來我們信這一位神.
為什麼會使我們有盼望呢?.
原來他就是盼望.
原來這位神.
盼望是這個意思.
保羅很特別的.
所以原來神的永恆屬性.
可以這麼豐富.
又是這位神.
他是愛的.
所以神就是愛.
《約翰一書》的四章十九節.
這一位神.

$^{641}$信實的是.
其實是很奇妙.
這位造物主.
他的傳言的豐盛.
所以這裡就讓我們有一個這樣的眼光.
原來這一位神.
所以我們永遠在神裡面.
就有盼望.
永遠我們什麼時候很失落.
很低谷.
我們要想回神.
祂就是盼望.
如果我們認識這個希臘文也不錯.
保羅說這位神.
因為是一個Article單數.
這一位主.
這一位神.
Felix就是神.
祂就是盼望.
是這個意思.
不是說我們希望祂將盼望給我們.
所以使人有盼望.
是我們能夠明白.
所以其實神就是盼望.
所以我們的信.
投向祂.
我們就有盼望.
所以其實這樣解釋也是對的.
中文的翻譯也沒錯.
但願使人有盼望的神.
因信將諸般的喜樂平安.
諸般的.
是如何諸般的呢?.
如何多樣化呢?.
如何在我們不同的場合.
我們不同的處境.
使我們有喜樂平安呢?.
大家的處境是不同的.
我們今年的處境.
可能跟十年前的處境已經不同了.

$^{681}$你做大姐仔的時候.
還沒生孩子的處境也不同了.
你現在做奶奶,做婆婆.
那個處境也不同了.
所以人生有不同的處境.
現在我們的子女長大了.
開始凶巢了.
處境也不同了.
你想想小時候.
我們星期一時辰和明仔還在幼稚園.
就坐車由元朗坐到屯門.
過了星期一.
指定的娛樂就是要陪他.
但現在長大了.
你想跟他旅行.
不好意思.
他要自己找朋友去.
你要明白大家不同的處境.
你一定要接受的.
所以原來珠缽的喜樂平安.
原來在神的裡面.
是有珠缽的喜樂平安.
是要用信心去支取的.
所以有時大家安靜在自己.
我不知道哪個是你自己的所謂樹洞.
所謂你的安全的地方.
只有你一個人.
你閉上眼睛.
默想神.
跟神好像聊蜜話一樣.
甚至你有意識地去想起你的童年.
你的成長等等.
跟神好像一直高呼那些以前的日子.
尤其是我們現在年紀大.
小時候那些歷歷在目.
你可以回去.
那時候你沒有神.
現在你有神.
你跟神一起回去.
很甜蜜的.

$^{721}$很讓你的信心能夠增長的.
保羅告訴我們.
當你的信心這樣運作的時候.
當你不急躁.
不急躁地.
這樣慢慢地去沉醉在神的裡面.
得到的.
因為他是盼望的神.
他就會將.
原來盼望是帶來甚麼呢?.
帶來的是喜樂平安.
盼望帶來的是喜樂平安.
就是說你的心靈充滿了一份肯定.
充滿了一份對未來的安全感.
所以喜樂平安充滿在我們的心靈.
所以一定是在心靈裡面.
由然而生的那種喜樂.
那份平安.
以致我們能夠大有盼望.
所以大有盼望的人.
其實就是內心很有那份平安滿足.
接著這裡說.
當我們的信心藉著甚麼呢?.
藉著聖靈的能力.
所以我就做了一點分析.
那個源頭.
盼望的源頭是神.
所以不是環境.
不是環境改變你就有平安.
是的.
表面是的.
但是聖經往往告訴我們.
聖經所應許的就是.
就算你沒有那個環境的改變.
但是已經有平安.
因為神的出人意外的平安.
已經盡著.
那個字是軍隊駐軍.
盡著這個《腓立比書》四章.
神所賜出人意外的平安.

$^{761}$必在基督耶穌裡面保守.
保守就是那個駐軍.
就是解放軍香港的駐軍.
就令到我們.
當然我們沒有打仗.
所以我們不怕.
如果你是一個戰亂的國家.
現在的軍隊很強大的.
來進駐.
你就有平安了.
我們明白了.
所以原來我們親近神.
那個源頭就是平安.
就是盼望.
就是喜樂.
所以環境固然是有幫助的.
但是有時候環境還沒有改變.
我們已經先得到平安.
比如我有個例子就是.
維斯利·約翰.
他那個時候.
他都出名了.
他開始去學習了.
那個時候還沒有全道.
他有一次搭船.
是很大風大浪.
因為那些帆船.
不是今天的那些郵輪.
郵輪那些很強大的巨風.
我們都會怕.
都會懼怕.
但是當時他真的很怕.
但是很奇怪.
他看到一群清教徒.
很安靜的在那裡.
船是搖得很厲害.
看到一群清教徒的弟兄姊妹.
在那裡很平靜的唱聖詩.
為什麼他是基督徒.
為什麼他沒有呢?.

$^{801}$他回到英國之後.
真的認真的去尋求.
後來他又說得到一個第二次祝福.
就是一個滿足的喜樂平安.
當然就是他慢慢找回神.
在他生命裡面的關係等等.
以至於將來.
後來他就真的將自己完全的毀失.
做一個很出名的報導家,傳道人.
這位衛斯理約翰.
後來就創辦了這個衛理公會等等.
就引發了17世紀的大復興.
美國,英國的大復興.
但是他就是沒有經歷到那份平安.
這就告訴我們.
原來十三節是一個應許.
所以這個聖經我們要背熟它.
我們要常常在盼望的神裡面.
來找我們自己的定位.
找我們自己心靈的力量.
以至我們裡面有力量.
我們就得到那個平安.
講到這裡的時候.
又很特別.
昨天我們在辦公室.
跟一些同事行政同工聊天.
說著安息禮拜.
因為我們要準備那些程序表.
說著我們很快要做好自己的.
寫好詩歌,唱什麼歌.
說著說著就說到死亡.
我就說上一次我那本書.
很厚的那本書.
就是那個精神科醫生.
他正式出來做醫生的時候.
那個經歷就說.
有一個女病人服藥過度昏迷.
好像是死了一樣.
第二天早上居然醫生見他醒了.
然後他令醫生感到奇怪.

$^{841}$和震驚的就是.
這個女病人說.
你昨晚跟我的同學.
就是送我來醫院的那個同學.
那個對話.
還有醫生你的胎有茄汁.
這個女病人是昏迷了的.
他能夠說到.
那本書就是醫生他這一次的經歷.
後來成為他50年去研究.
到底生命的死亡是怎麼樣的呢?.
人有沒有靈魂呢?.
是不是真的會經歷到靈界呢?.
我們一直說的時候.
同工很有趣很有興趣.
那本書不行.
因為拿去圖書館續借.
沒辦法了.
原來有很多人跟著借.
我都續借.
就要立刻還給圖書館.
又要重新等.
很多人對靈界很有興趣.
似乎很多人都還沒信耶穌.
我們這樣說完.
誰知道走出去三樓出來的客廳.
有一個我們起初不知道他是信主的.
一位家長帶小朋友來學前的訓練.
他聽到我們說話.
就跟我們說要加入.
是的 我們真的有個經歷.
他的奶奶在屯門醫院100歲的時候.
就說那些家人圍著她.
她已經很老了.
是自然的病.
她說你們不用擔心我了.
我看到光.
我看到有個白衣的人來帶我走.
我看到一個很光的地方.
她很開心地很平靜地離開.

$^{881}$因為她們全家都是信主的.
原來這個婆婆80多歲的時間.
是還沒信的.
原來她是經常看到靈界的.
是那些污穢的靈.
那些靈界就是邪靈那些鬼.
但後來不知道什麼原因.
她信主.
她說她以後就沒有再看到那些東西了.
反而她說有時候會看到這些天堂.
看到這些靈界.
是很安詳很平安地離開.
這樣帶著微笑進入永恆.
這個姐妹就很高興地跟我們分享.
因為我們沒有想過.
她會聽到我們說話.
在房裡面.
所以原來又有很特別的經歷.
我們一直都不知道這位家長原來是信主的.
所以我們看到.
保羅這個應許的禱告.
真的可以兌現.
基督徒.
我們今天在說的這一段經文.
整段的一到第十三節.
都是回應十四章.
神的國度是怎麼運作的呢?.
神的國度的指紋.
是在說公義和平.
聖靈裡面的喜樂.
不是追求外面.
外表.
不是要自己做一百分.
神才會拯救我們.
不是.
而是神給我們有這個新生命.
這個新生命就要流露出來.
就是要將神的救恩的愛.
展現在我們公義的人生.
在我們一個平安的人生.

$^{921}$即是和睦的人生.
和喜樂的人生.
所以我們就要.
很實在的在神的教會就是國度.
所以神的教會就是應該展現合一.
展現愛.
展現接納.
所以來到這裡的時候.
保羅就做了一個小小的總結.
就用一個禱告.
一個祝願.
一個渴慕.
我們能夠經驗.
盼望的神.
給我們的信心.
信心帶來的就是.
我們能夠有珠缽的喜樂.
珠缽的.
令我們滿足的.
令我們能夠感受人生.
不是一種重擔的.
當然會辛苦.
每個階段都會有不同的辛苦.
昨天又有一些不好的新聞.
一對七十多歲的夫婦在酒店.
可能是自殺.
可能是他殺.
因為可能他的病太辛苦.
所以要用類似人為的安樂死.
是一些悲劇.
但是屬主的妮女.
我們記得.
就算我們可能有重病.
都會辛苦.
都會有很多的可能副作用.
很多的問題.
但是心靈肯定是不會辛苦的.
心靈越信主內.
是越有那份真正的平安.
真正的永恆的歸宿的喜樂.

$^{961}$好像那位一百歲的老人家.
是很開心地看到大光.
帶他進入永恆.
所以其實我看得多.
我很羨慕.
我說很慚愧.
我從來都沒有見過這些.
我很羨慕.
鄧英姐妹.
她就神為她開了這個大光的靈界.
但是當然我們透過經文的真理.
使徒保羅告訴我們.
這一位的神.
就是那個盼望的源頭.
使我們一定可以經歷.
途徑原來就是我們的信心.
所以我們基督徒的信心是要不斷增長的.
我們不要停留在一知半解.
算了 隨便吧.
而不是的.
我們的信心是要增長.
對真理有一個認識.
以致我們能夠經驗到聖靈在我們裡面的工作.
所以聖靈的工作不是行神蹟.
不是說一些古怪的謊言.
而是聖靈的能力.
讓我們大有盼望.
讓我們似乎當時都不一定是平安的.
保羅的傳道.
他要面對很多困難.
面對的可能是很實在的.
有些很極端的猶太人.
一直針對他 追殺他.
因為寫完《羅馬書》之後.
他就要上耶路撒冷.
一上耶路撒冷.
就在安提摩上來的那些猶太教的人.
就見到保羅.
見到他和一個外邦的年輕人.
應該是特羅菲柯.

$^{1001}$在耶路撒冷在聖殿外面.
就來鼓動那些群眾.
就捉拿保羅.
所以保羅其實面對的實際的日常生活.
未必是真的有很大的順利的.
所以順利的生活不等於心靈的平安.
有心靈的平安.
不一定是一帆風順.
所以是沒有抵觸的.
我們基督徒有的是心靈的盼望.
盼望帶來的就是喜樂.
是諸官的喜樂.
是諸官的平安.
就是說不同的處境.
那種平安都會不同的.
是令我們真的心滿意足的.
所以我很喜歡保羅這樣的一個小傑的禱告.
就是一個竹猿.
對使徒來說.
他自己就是這樣的心靈.
所以我們這些要背的.
我們能夠大有盼望.
不同的處境.
現在都有多少不同的香港的教會弟兄姊妹.
去了英國,去了加拿大,溫哥華等等.
是這兩個地區比較多.
但是大家可能在新的環境,新的處境.
都相當困難的.
甚至不困難到一個地步.
連回香港都未必可以的.
整個的生活,工作都不同了.
其實新的環境.
有很多的處境都不是我們香港的信徒.
能夠明白的.
所以變成真的需要更加依靠神.
所以十三節.
如果有頂尖的在海外讀到這段經文.
真的很盼望大家要沉澱下來.
找一個你自己可以躲起來.
最好一兩個小時.

$^{1041}$你自己安靜在神的面前.
整個人浸在神的盼望和愛裡面.
去思想自己的困難.
將那些困難帶到神那裡.
會被神的愛,神的盼望包圍著.
當我們完結了這樣的妄想.
保羅就告訴我們.
因為聖靈的能力.
因為我們的信心.
就大有盼望.
而且在那個大有盼望裡面.
心靈就有喜樂,平安了.
所以這裡就是保羅叫我們去學習.
接著我們從第十四節開始看.
開始看就是保羅對自己的一個小小的分享.
就是他現在說了這個大段落.
十四,十五章.
嚴格來說就是一個信息.
就是信徒.
無論你什麼處境,什麼傳統,什麼文化.
都要學習接納.
對別人欣賞.
對自己所做的堅持.
就是彼此的接納.
然後在一個合一的裡面.
來到一層的去為萬民.
得到福音來到效力.
大家在不同的神秘的崗位.
來到去盡心,來到去竭力.
接著保羅在這裡講一講他的使命.
保羅的宣教使命.
十四章我們讀一讀這一段.
他說弟兄們.
我自己也深信你們是滿有良善.
羅馬的弟兄姊妹.
充足了諸般的知識.
在理性上.
也能彼此的血界.
也能夠有這樣的學習.
彼此來到去血緬.

$^{1081}$也不容易.
年紀越大就越難聽別人的血界.
越難聽別人說你這樣不對.
很難.
但我稍微放膽寫信給你們.
好像很委屈.
其實不是的.
保羅真是一個很謙卑的牧者.
這位使徒.
他說要提醒你們的記性.
特別是因為神給我的恩典.
使我能夠成為愛邦人的僕役.
哪個僕役呢?.
基督耶穌的奴隸.
僕役.
作神福音的祭司.
福音來說.
福音是獻上.
當然是獻上愛邦人.
來作為祭司.
祭司.
又是祭司.
又是奴隸.
所獻上的愛邦人.
因著聖靈.
成為聖傑.
可無月立.
強調不是他的powerful.
不是他的知書識禮.
學問淵博.
不是他能夠.
懂AI.
上至天文下至地理.
什麼都懂.
保羅告訴我們.
因著聖靈.
真是不簡單.
一個牧者能夠這樣看自己.
我所做的一切.
微不足道.

$^{1121}$因為聖靈我才能做到一點點的事.
將愛邦人.
成為一個輕輕.
屬靈聖傑的祭.
獻給神.
神是收納的.
舊約獻祭.
我們將來看完這部劇.
就跟大家一起去看.
一定唯一可以蒙神月立的.
就是要照神的吩咐.
要將祭身試驗.
要洗乾淨.
要沒有殘疾.
然後.
就要用輕香的方法.
首先如果你是獻梵祭.
就要將祭洗乾淨.
然後要切成一塊一塊.
內臟要挖出來洗乾淨.
然後放好.
然後.
獻的時候.
再加上酒.
輕香的祭.
即是說那些香氣升上去.
大家燒烤就知道.
聖誕中秋就特別香.
整個Sarandes.
都是燒烤的香.
即是油香.
但不代表聖經說的潔淨.
可能那些豬.
那些肉.
未必是那麼新鮮.
但猶太人就不是.
一定是要驗過的.
祭師是有責任要驗.
那個祭身.
是否乾淨.

$^{1161}$是否沒有殘疾.
才可以獻的.
獻了還要洗乾淨.
獻了還要講究.
你可以想像到.
那個恩爵聖靈.
成為聖潔.
可忘神的月立.
即是保羅.
不是隨隨便便.
做一些.
貓貓呼呼.
而他盡心將外邦人.
無論有.
所以那個不是量的問題.
是質的問題.
他用心.
向外邦人傳道.
無論多困難.
多少猶太人對他的嘲笑.
對他的侮辱.
對他的攻擊.
他清楚自己的遺憤.
清楚自己是神的使徒.
現在記住.
是因為神的恩典.
所以你想想.
單單幾個的形容詞.
已經讓我們看到.
這位謙卑.
牧羊的牧者.
牧羊的心靈.
所以也是的.
越是成熟的牧者.
年輕的時候.
好大喜功.
好像出神學院.
十八般武藝.
要怎樣去耍.
講到個個都.

$^{1201}$說得很好聽.
慢慢他發現.
其實不需要這樣.
應該都是慢慢學習.
是神的恩典.
他能夠站在港台.
去做勉勵.
去做分享.
做教導.
代表神希望能夠.
透過神的道.
去安慰不同的弟兄姊妹.
或者透過神的道.
去指正大家.
不同的人生的方向.
所以你看看普羅.
他這樣講自己.
和這班羅馬的弟兄姊妹.
這樣講.
要提醒他們.
我稍微的放膽.
我不是誇口.
我和你們去分享.
我的生命.
和我為什麼要這樣做傳道.
做使徒.
就要來提醒他們.
寫信給你們.
要提醒你們的記性.
特別因為.
神給我的恩典.
做使徒.
原來普羅.
其實早就不需要去做了.
因為他不配.
因為碧珀神的家.
曾經侮辱.
褻瀆耶穌基督.
是一個.
沒有萬人的人.

$^{1241}$時時講一些很難聽的.
尖酸刻薄的話.
真的用言語.
用行動去攻擊.
那些信耶穌的人.
所以現在.
是神給他的恩典.
所以弟兄姊妹.
是否如果你一直都不覺得.
是恩典的崗位.
其實.
是不可能的.
因為越事奉的人.
越感受到的只是恩典.
給你能夠配.
可以站在不同的崗位去服侍.
所以越是事奉.
而你越是被人感覺到.
你在炫耀.
其實你根本是在用你自己的驕傲.
你自己的自以為是.
是不會幫到人.
祝福到人的.
真正的對事奉的體會.
就是恩典.
所以.
特恩神給我的恩典.
使得我能夠成為.
外邦人的僕役.
為外邦人.
作基督耶穌的僕役.
你想想多麼謙卑.
多麼對自己的.
欺負如其分.
不做不行.
你是奴隸.
你可以想像.
誰會甘心樂意.
說自己是奴隸.
最好就快點釋放.

$^{1281}$快點用錢把我贖了.
怎麼可能要自己做奴隸呢.
但是.
甘心樂意做奴隸.
那就真的沒得救了.
好像我們小時候.
十幾歲的時候.
那些自甘墮落的.
去做魚蛋妹.
去做姑爺仔.
去做妓女.
這樣來賺錢.
給姑爺仔吸毒.
自己還甘心樂意.
這些甘心樂意就是愚蠢.
我們明白了.
但是保羅明白恩典.
明白救恩.
明白為什麼要做使徒.
他說他願意.
做外邦人的奴隸.
僕役.
其實是基督耶穌的奴隸.
僕役.
所以為了外邦人.
他甘願成為基督耶穌的奴隸.
你想想.
謙卑得來又有多尊貴的身份.
所以單是這句話.
對很多年輕.
跟隨保羅的年輕人.
譬如提安泰.
在他身邊.
後來的路加醫生.
或者是.
推基古.
很多年輕人.
看到保羅這樣對自己的評價.
就有很大的感動.
所以.

$^{1321}$今天我們需要.
當然有學問.
保羅大有學問.
希利文也是寶貞通.
所以兩者不應該是抵觸.
而是那個屬靈.
生命對於主的降伏.
因為降伏到成為奴隸.
就不能再降伏了.
因為是最低的.
沒有自由的人.
所以.
你想想.
保羅在當時羅馬世界的奴隸制度.
你想想.
奴隸制度是多麼沒有自由.
多麼沒有身份.
但是保羅願意為主耶穌.
為福音.
甚至為外邦人.
他願意沒有身份.
沒有自由.
沒有那個.
好像尊嚴都沒有了.
原先是一個法理賽人.
是一個猶太常議會的議士.
等等.
所以.
這些經文是有血有肉.
是可歌可泣.
哲學家很簡單.
他只是說說他現在的身份.
就已經能夠令到.
真是服侍主的弟兄姊妹.
包括我們做傳道.
都是一個很大的感動.
這裡告訴我們.
他現在.
在福音這麼榮耀.
他為神.

$^{1361}$這個福音做祭司.
所以全部是.
甘心樂意.
保羅的事奉.
所以弟兄姊妹.
其實我們查經班.
如果能夠在這麼多年.
能夠找到.
幫到一兩個人.
甘心樂意一生奉獻給神.
為神而活.
為神而去.
站在他自己的崗位.
那就是查經班的成功.
因為可以有這麼仔細的.
逐節去分享.
去彼此很坦誠的交通.
去交流.
就是想激發大家.
對神這一位.
做物主的愛.
對耶穌的效忠.
所以他現在.
一直形容他自己.
又是奴隸.
但其實不是.
他是一個祭司.
他對舊約的那種明白.
祭司是很尊貴的.
是全舊約的民族.
最大的那個.
等於現在東正教也是.
或者現在天主教也是.
教皇真是最尊貴.
數十億的信眾.
對你的景仰.
你走出來.
聽你的說話.
本來很明白祭司的尊貴.
但他現在.

$^{1401}$不是回到舊約.
他看他自己是一個傳道人.
他原來是將外邦人.
放進福音裡面.
福音才可以改變.
外邦人的命運.
生命.
所以我們大家.
有沒有這種羨慕.
傳福音呢.
不可以.
其實我傳福音很弱.
我不太喜歡傳福音.
也不敢傳福音.
真是不同的.
我不會很刻意.
直接見到人就說.
投降福音.
但這兩年.
在洪水橋.
很多機會.
無論是前線.
或者策劃等等.
就愈來愈感受到.
其實是很難的.
可以有這種突破.
但本來就已經告訴我們.
他看自己是奴隸的尊貴.
因為是基督耶穌的奴隸.
不是錢的奴隸.
不是罪的奴隸.
他是基督耶穌的奴隸.
所以原來.
我們在主耶穌裡面.
慢慢原來我們自己.
捆綁自己的結.
是不知不覺會解開的.
所以他就是基督耶穌的奴隸.
他就是神的福音的祭司.
祭司獻什麼呢?.

$^{1441}$就是獻外邦人.
什麼獻外邦人呢?.
不是來吃人肉.
行邪教儀式.
是將外邦人帶進福音的真理.
外邦人明白.
恩信稱義.
所以這本羅馬書.
是很有系統地講解.
這個奴隸的真理.
所以這本羅馬書.
是很有系統地講解.
這個恩信稱義.
萬民可以恩信稱義的福音.
所以教初述班的弟兄姊妹.
是非常地蒙福.
因為你經常接觸到一些未信的朋友.
慢慢帶他進入形式救恩.
是很大的福氣.
所以成為一個祭司.
是很尊貴的.
一個大祭司.
是一個大祭司.
是全民族最重要的人.
現在保羅看自己原來在神的福音裡.
他是一個祭司獻上的.
果然是外邦人.
因為他知道自己是外邦人的祭司.
是神給他的使徒的身份.
所以他當然是使徒.
原來使徒.
都是要將福音講解明白.
所以要聽得明白.
其實是不簡單的.
所以不同的受眾.
對福音都有不同的要求.
福音只有一個.
但你能夠講解得他明白.
在他的處境.
他能夠認同和改變他的生命.

$^{1481}$這就是祭司的責任.
這裡告訴我們.
所獻上的外邦人.
因著聖靈成為聖傑.
是可蒙神的月立.
所以必須成為聖傑.
是可蒙神的月立.
所以不要用鬼詐手段.
不要用取巧的辦法.
實事求是.
信就是信.
信是因為耶穌的救恩的愛.
外邦人明白了得救.
所以我們一定要記住.
救恩本身已經是強大的.
救恩本身已經是能夠改變任何一個人.
所以我們不需要害怕.
這些人很有錢.
很難信耶穌.
很嘲笑耶穌.
不一定的.
有錢人都有他們那種生命的改變.
所以我們求是.
讓我們能夠透過這些.
更加看到使徒保羅.
他自我的評估.
這裡不是突然間.
心亂為潮.
突然間有一種.
這樣的.
感懷身世.
不是的.
他也不是在炫耀自己.
而是他很想引導羅馬的信徒.
我講完了倫理關係的合一.
現在我講一下我自己.
因為他差不多要寫完這封信.
他就講一下現在.
我的處境.
我的來歷.

$^{1521}$因為不是每個人在羅馬都認識他.
收這封信的羅馬的信徒.
可能聽聞這個保羅.
未必真的讀過他的書信.
因為當時很多的監獄書信都未寫.
所以保羅就開始介紹他自己.
介紹自己在《救恩》裡面的人生崗位和角色.
很寶貴.
一定選他做使徒.
他現在傳到一段時間.
對自己這個身份的那種認受.
那種熱愛.
對福音的熱誠.
對外邦人的擁抱.
不是格格不入.
你信就信.
你不信就該死.
一個年長的牧者.
那個牧養心就是.
慢慢發現每個人都是這麼珍貴.
每個靈魂都是這麼值得.
放心血,放時間.
所以這裡就讓我們看到.
所以第十七節.
他說論到神的事.
我在基督耶穌裡面有可誇.
論到神要他做的事情.
他就在主裡面.
他就好像誇耀他自己.
說自己是祭司.
然後又說自己是奴隸.
到底你是怎樣的.
保羅就告訴我們.
其實他所謂誇耀.
不過是神要他做的事情.
第十八節.
這裡就一起看.
十八節 除了基督我作的事情.
我甚麼都不敢提.
只能暗示.

$^{1561}$保羅有很多宿命經歷.
保羅的手帕.
使徒行傳十七章.
人們拿他的手帕.
回去治好病人.
人們拿他的手帕.
回去治好病人.
我們認為使徒時代.
因為聖經未傳遍.
未寫好.
正典還未圓滿.
所以神就要透過使徒.
給一些神跡歧視他們.
去傳揚福音的真確.
所以仍然是福音.
不是神跡.
所以不要被神跡影響.
真正的重點是.
神跡是印證福音的大能.
福音的改變.
使徒們能行神跡.
彼得親身經歷過.
自己的家門打開.
冰冰在他左右.
突然昏迷.
彼得自己有天使牽著他.
走出監獄外.
保羅曾經.
將油推鼓.
在三樓.
跌在地上死去.
靈魂都離開了.
保羅抱著他之後.
發現他還有力量.
神就將生命帶回來.
給油推鼓.
你想想.
這個病人繼續上三樓.
繼續講道.
這是《善行經》第20章的經文.

$^{1601}$你可以想像到.
神愛不不感激.
只提他直我言語作為.
果然.
用神跡騎士的能力.
並聖靈的能力.
使得外邦人馴服.
這些就是使徒們獨有的.
那個時代的需要.
那個時代的特別原因.
所以.
我們今天都是同樣.
就算我們沒有.
神跡騎士的能力.
但我們有的是.
聖靈的能力.
我們有的是對福音的擺上.
還有我們有的就是.
基督要我們去做.
我們就有權柄去做.
有這樣的能力.
所以我們能夠.
趕鬼.
如果神要我們行神跡騎士.
我們是能夠行神跡騎士.
所以我們就應該.
要在神裡面.
學習保羅在第十八節裡說.
我們要尊重神.
給我們那個崗位.
那個召名.
那個尊一.
那個擺上.
使得神能夠.
在我們的生命裡面.
叫其他人得到幫助.
益處.
所以這就是.
最終是使得外邦人.
馴服.

$^{1641}$所以神跡不是給保羅炫耀的.
所以他們不會說很悶.
變個神跡出來.
是可以的.
最終是要令到外邦人.
馴服.
所以大家就明白了.
如果神跡是為神跡而神跡.
那就錯了.
為了港方言而港方言就錯了.
是要使得外邦人.
能夠馴服.
馴服神的福音.
馴服耶穌基督的.
偉大的十字架.
救恩.
所以保羅就很明白.
使徒們他們都很明白.
每個行神跡都不會誇耀.
都不會覺得.
我是比其他使徒.
更加好.
因為他們就很明白.
那些根本是會過去.
真正的.
是基督耶穌在人的心靈裡面.
那個生命改變.
而帶來的.
馴服.
所以是.
馴而馴服.
這樣改變了之後.
就會得到神的.
生命和救恩.
所以保羅.
只不過是一個代理人.
今天我們都一樣.
我們都是要.
求主對我們.
藉著我們的言語行為.

$^{1681}$藉著我們能夠依靠.
聖靈的能力.
藉著我們在神裡面.
清楚的照明.
所以照明是很重要的.
清楚神.
要你做什麼.
專注去做.
專心去做.
用心去讓神.
在人的生命裡面得到.
那個建立.
接著十九節.
甚至我從耶路撒冷.
一直轉到以利利古.
就是馬其名.
一帶的地方.
到處傳了.
基督的福音.
保羅不是誇口.
他說他這一帶.
地方他都到處.
將耶穌基督.
讓那些人都知道.
信不信當然是.
人家的責任.
是否拒絕.
是否接受.
保羅的責任就是.
代表著主耶穌.
由耶路撒冷.
一直到外邦的以利利古.
到處.
因為他是外邦人的使徒.
到處傳了.
基督的福音.
將這個救人靈魂的生命.
帶到他未信的朋友中間.
第二十節.
保羅又再一次說.

$^{1721}$他立了志向.
不在基督的名被.
稱過的地方傳福音.
免得建造在.
別人的根基上.
就如經上所記.
未曾聞之他信息的.
將要看見.
未曾聽過的.
將要明白.
雖然我們這樣看.
大家都會明白.
保羅不是特意來分庭抗禮.
玩對立.
你向東我就向西.
因為我和彼得不合.
因為我很憎巴勒巴.
我相信肯定不是這樣.
而是保羅.
他立了一個志向.
他就說很希望.
福音能夠去到.
未得知聞.
未聽聞的.
以及他不做在.
別人的根基上.
即是說他不想去.
做一些神母要求他做.
他來到去疊加上去.
他來到去.
吃別人的便宜.
這樣的一種取巧.
他也不是特意要.
標其立意.
要找苦來生.
不是這樣.
而是保羅都清楚.
神既然給他一些疆界.
是彼得未必做到的.
他就用心去做.

$^{1761}$所以這樣叫不建立.
在別人的根基上.
就不是說.
我去到那間紅印堂.
我就特意去對面那間.
就不是這個意思.
他很明白.
同宮門 使徒門.
大家是很和諧的.
大家有不同的目區.
大家有不同的年齡層.
不同的民族.
但是保羅就很清楚.
神要他.
可能是他的性格.
可能是他的恩賜.
等等.
他不想去到那麼遠.
去向那些彼得的外邦人傳道.
所以這是一個很好的學習.
一個既和不同的牧者有關係.
有愛的彼此的欣賞.
彼此的協作.
但是同時間.
又知道自己有一個學府.
有自己的工場.
有自己的疆界.
這樣的胸襟.
真是一個偉大的人才會有.
就是看到福音是重要的.
所以就不是像炒飯一碟那樣.
人家做就一起做.
就是說好像跟風那樣做.
就不是這樣.
而是保羅已經有他所謂的原創性.
有他所謂的.
他要開闢一些沒有人到的地方.
記住不是要跟人抗爭.
去刻意做給人看.
不是這樣.

$^{1801}$這樣就是鬥氣.
這樣就是小家.
這樣根本就不是一個牧者的情懷.
應該就是真的有些地方.
是必須保羅去的.
所以他去到四大行傳的27章.
他要去到君王的面前.
應該是25,26章.
去到阿基柏王的面前.
來申訴.
所以你可以想像.
本來是一個被人一直嘲笑排擠的.
這位士徒保羅.
居然安排他.
真是一個囚犯的身份.
但是被君王接見.
阿基柏王.
你可以想像.
很諷刺的.
阿基柏王都要驚保羅說的真理.
就是對永恆沒有把握.
但是一個囚犯.
自己對永恆是很有把握.
就算我今天斬頭.
我都對永恆有把握.
因為我就是這位永生神的僕人.
所以保羅這個志向.
都是為大家要學習.
大家有不同的崗位.
不同的托付.
大家忠於自己的托付.
我畢業都15年了.
神學院.
不同的同學開始.
真的不同的崗位.
不同的托付.
不同的工場.
不同的領域.
有些甚至現在已經沒有再做傳道.
做一些自由的工作者等等.

$^{1841}$不同的大家的人生際遇.
但是保羅這個志向.
其實都真的時間會讓我們浮現這些志向.
我們要專注在神托付我們的崗位.
去忠於所托.
剛剛因為星期有同學.
忠於按木.
已經很大年紀了.
在教會30多年了.
忠於在教會長大.
現在忠於按木了.
我們按了10年了.
她現在才按木.
女孩子來的.
做一些兒童工作.
在恩泉堂.
她真的神秘有不同的人生的托付.
所以這裡告訴我們.
保羅她這樣和這些羅馬的信徒分享.
很忍耐.
說她自己是外邦人的一個奴隸.
一個祭司.
獻上的外邦人.
以及說到她自己是怎樣呢.
在她傳的領域已經全傳了.
我想起我們武堂的一位.
很德光旺眾.
已經離開世界回天家很多年了.
在20年代.
上個世紀20年代.
是懷仁教會每人都認識的.
一位很德光旺眾.
很多人尊重她的長者.
她說她那時60年代住在美孚新村.
她說她差不多70多80歲.
每天拿一大疊單張.
寄退休的老牧者.
一位老牧者住在美孚新村.
我一上去她家查經.
她說美孚新村.

$^{1881}$我差不多是這麼多座的居民.
我都派了單張.
一位老牧者.
80歲.
回到教會講完道.
因為身體都不好.
已經年紀大了.
她在教會拿了一疊單張就去派.
這樣在美孚新村.
很多座很多人.
她說我差不多已經派了.
你可想像到這位老牧者.
對永恆的愛那種熱愛.
保羅就是這樣.
所以這些經文其實都是很感動.
不同時代的牧者.
保羅說到處傳了福音.
由耶路撒冷一直到以利利古.
為誇耀.
保羅真的這樣和這些羅馬信徒分享.
因為後面就是鼓勵他們.
一起向萬民宣教.
然後她就說.
已經有人傳的.
我就不希望建立在別人的根基上.
然後引聖經證明.
原來他都不是一個.
好像木空一切.
這麼驕傲的保羅.
又說全了.
又說不會建立在其他人的根基.
即是不稀罕的.
我不需要有同伴.
完全不是這樣.
既是有很多的同路人.
既是有很多的跟隨者.
既是有很多的和他一起的使徒.
他覺得自己是最不配的使徒.
因為他曾經攻擊這位主耶穌.
但是他想說的是.

$^{1921}$神給他的福音的靈力.
他要做盡.
即是未完結.
他都不離開.
這些學習.
真是對一個牧者很重要的生命的教導.
引的聖經就是以賽亞書的52章.
未曾聽聞的可以得到.
這是保羅自己.
你可以說是他內心對福音.
對外邦人的愛的態度.
接著22節.
保羅的行程.
就像他所說的.
22節開始.
我因為多次被攔阻.
總不得到你們那裡去.
但如今我在這一帶地方.
再無可存的地方.
即是上文說的.
而且這好幾年我切心上望到西班牙去.
西班牙不可以這樣說.
歐洲最南面.
當然我們不計其他.
聖保羅陀最南面就是西班牙.
因為再下就是北非.
所以是對的.
整個歐洲是由特羅亞開始.
整個歐洲我們有沒有地圖呢?.
我好像做了一個地圖.
在馬其頓亞西亞.
特羅亞過去對面就是馬其頓.
接著就是菲律賓,帖撒羅尼迦.
北面的就是帖撒羅尼迦.
地中海北面就是希臘.
現在是希臘.
帖撒羅尼迦.
一個相當漂亮的地中海地方.
一直下來.
哥林多都不是最下.

$^{1961}$還要再去南一點.
就是西班牙.
所以可能下個星期.
會給一個再大一點的福原歐洲地圖.
所以為什麼奇怪.
保羅要去西班牙.
就是歐洲最南的地方.
其實再南再下就是北非.
非洲.
所以保羅告訴我們.
他就是這樣.
我切心上望去西班牙.
可以到你們那裡去呢?.
盼望從你們那裡經過.
就是羅馬,意大利.
那裡就是三盤.
如果是南轉那裡告訴我們.
接著說盼望從你們那裡經過.
不見你們.
跟你們彼此的交往.
彼此的分享.
心裡面稍微的滿足.
然後蒙你們的送客.
預先說了這個行程.
預先說我希望將來有一段時間.
我就會直接去到你們那裡.
跟你們有一些的商交.
可能有一些的聚會.
有一些的茶經.
有一些的講道.
接著你們就接納.
接著羅馬教會就.
不是寬送.
是支持保羅的宣教.
就是蒙你們的送客.
就是說你們是認同.
我現在宣教的行程.
神給他的事業等等.
不過現在.
二十五節就說現在.

$^{2001}$保羅他現在就要上耶路撒冷.
留意這裡.
我也預備了最後這兩幅的powerpoint.
但是如今我要到耶路撒冷.
這個就是如今.
將來就是很希望去羅馬.
去意大利.
然後跟你們有一段時間的分享.
交往.
說我的全部策略.
接著你們認同.
你們有警察全部支持.
你們就給我送客.
就是說公開的教會的一個支持.
譬如我們今年都期望.
我們很少我們的金錢.
都是在性質上的意義.
會支持Z7.
會支持士思一家.
會支持Chris的野地牧人.
我們過了同委會.
我們就象徵式.
就是不會有很多錢.
但是一個很重要的就是教會的.
支持他們去承行.
所以蒙你們的送客.
所以教會的公開的接納.
其實很重要.
代表著一個很重要的宿命的.
一個同心.
一個差遣.
所以保羅說.
不過現在我要去耶路撒冷.
我的時間未必能說這麼多.
因為這裡也有些東西要解釋.
我們今晚就停在這裡.
我們從25節下一次再說.
我們一起同心禱告.
再一次我們多謝主.
我們透過這一大段聖經.

$^{2041}$看到保羅的內心.
他對神.
這位盼望的神的認識.
學悟.
以及能夠成為這位做主的僕人.
一個奴隸.
他願意放下他人生的自由.
來為外邦人擺上.
要將外邦人帶進神的救恩的國度.
求你幫助.
讓我們錦宣家.
和一家嘉欣堂的弟兄姊妹.
或者其他.
可能在其他教會裡面.
偶然間上網找到羅馬書.
又剛好找到我們這些訊息的時候.
都成為眾弟兄姊妹.
神國度裡面的兒女.
那個學習.
那個感動.
求主的聖靈去激勵.
我們很盼望.
將我們的人生都放在主的手中.
以致我們能夠蒙主.
閱立.
被主使用.
被主差遣.
在不同的崗位.
為主來盡力.
多謝你.
聽我們感恩禱告.
奉耶穌基督的名.
阿們.
感恩.
多謝大家的參加.
\newpage



\section{羅馬書}
\label{sec:LvE6qlAUC04}
\textbf{2026-02-06 周五晚查經班《羅馬書(三十六):分題:保羅 有擔當的使徒》 講員:曾裔貴牧師}
\newline
\newline
連結: \href{https://youtube.com/watch?v=LvE6qlAUC04}{\texttt{https://youtube.com/watch?v=LvE6qlAUC04}} ~~~~ 語音日期: 2026-02-08
\newline
\newline
\hyperref[sec:UU4dvSZbobM]{\small{< < < PREV SERMON < < <}}
~
\hyperref[sec:index_chronic]{\small{[返順時目]}}
~
\hyperref[sec:index_scriptual]{\small{[返順卷目]}}
~
\hyperref[sec:3rvmIgiASuQ]{\small{> > > NEXT SERMON > > >}}
\newline
\newline
$^{1}$我們可以祈禱了.
天父我們來到第36章.
我們看到在第13章一直到第15章.
都是一些倫理的教導.
特別是一個多元的文化的教會.
保羅是有一個很大的宏觀的教會觀.
因為他知道十字架帶來的.
不再是猶太人,外邦人.
不再是猶太人優先.
外邦人就是神的選民以外的外族.
反而我們看到保羅他自己有一個很大的生命的更新.
好像在我們稍後所說的.
他要去上耶路撒冷.
天父多謝你讓我們看到一位使徒.
是徹徹底底被十字架的愛去改變.
使他完全沒有任何的痕跡.
對其他人有藐視,有歧視.
反而他的愛就像神一樣愛世人.
多謝你讓我們看到他有很多的外邦的同工.
年輕的人是下一輩的門徒.
跟隨他去傳道.
我們在第16章又會有更多的人是保羅問安的.
今天晚上我們再一次去思想.
這一位有擔當的使徒.
他是怎樣在最後結束的文安裡面.
表達他的心境.
他的傳道宣教的生命的計劃.
都讓我們能夠在我們有限的人生.
去掌握神給我們的機遇.
給我們一些福音的使命.
我們再一次仰望主.
感恩禱告奉耶穌基督的名.
Amen.
好 我們繼續下去.
第36章說.
這次我們最後這一大段聖經.
由22到32節.
應該是33 不是32.
是一段很短的經文.
但是我們剛剛在第20節的時間.

$^{41}$就是保羅突然告訴我們.
他已經在周邊的地方.
要傳的道都傳完了.
我們從第15章第20節看起.
他說我立了志向.
不在基督的名被稱過的地方傳福音.
免得建造在別人的根基上.
就如經上所記未曾聞之他信息的.
將要看見未曾聽過的將要明白.
這是《以賽亞書》的52章的經文.
保羅很清楚自己.
就好像這位舊約的以賽亞先知那樣.
他要將神的美好信息.
當然舊約不是那個福音.
就是神的好消息.
就要告訴一些聽不到的人知道.
保羅就很想.
不是標其立意.
而是他很想自己.
來做一些神要他去做的工作.
或者那些群體.
他就不想沾別人其他人的光彩.
也都意味著.
當時大家要有這個意識.
就是原來好像雨後春筍那樣.
是有很多可能是講哲學的.
所以你會看見《少年行傳》有這樣的經文.
《十七章》有一些這樣的人在辯論等等.
當然也都有傳道.
所以你會看見肥立秘書.
有一班人在保羅坐牢的時候.
就很熱心出去傳道.
原來是故意想告訴保羅.
或者跟隨保羅的人知道.
我們沒有保羅還做得比他好.
他現在不能夠傳道坐牢.
我們就去做了.
其實是分庭抗禮.
這樣的心態已經是一定錯的了.
因為你背後的動機.

$^{81}$是一個要來跟別人有一個敵對.
福音其實是合一的.
所以保羅的意思就是說.
那又如何呢?.
或是假意 或是真心.
總之基督被傳開就可以了.
所以我們都要有這樣的眼光和胸襟.
但是保羅這裡就立了一個志向.
就是說如果有些地方已經有足夠的福音涵蓋.
有不同的傳道牧者來做過.
而且是神母要他留下來做的.
那麼保羅就不去這些地方.
去所謂的「心間細作」.
這個是保羅他自己的立志.
他說「免得建造在別人的根基上」.
接著第22節.
我們從今天開始看回到33節這一段.
告訴我們.
「我因多次被攔阻.
總不得到里門那裡去」里門就是羅馬教會.
「但如今在這裡再沒有可傳的地方.
而且這好幾年我徹心賞望.
到西班牙去的時候.
可以到里門那裡盼望從里門那裡經過」.
很合理的.
由意大利一直順著亞得尼亞海.
一直到南部 意大利的南面.
歐洲的最南面應該是葡萄牙.
西班牙就是布羅陀.
接著就是北非那些.
所以很合理的.
保羅是會經過羅馬,意大利.
然後才去西班牙.
就是說他當時的地理世界.
他已經知道西班牙是一個相當富庶的地方.
他都要向西班牙人傳道.
這個時候就告訴我們.
我們要留意第22節.
這個「被攔阻」是什麼「攔阻」呢?.
保羅自己有提過一兩次「攔阻」的.

$^{121}$所以我們可以先看看這個「攔阻」.
為什麼網絡很不順.
「帖撒羅尼迦前書」.
這段是比較早期的.
就是保羅第二次傳道的時候.
開始的時候.
在「帖撒羅尼迦前書」的第二章.
就有提到一個「攔阻」.
第二章這段經文.
第17節.
「帖撒羅尼迦前書」的第17節.
他說:弟兄們.
我們暫時與你們離別.
是面目離別.
心裡卻不離別.
我們極力的想法.
只肯願意見你們的面.
所以我們有意到你們那裡.
就是解釋他為什麼突然間離開「帖撒羅尼迦」.
因為在「士徒行傳」的第17章說.
他在那裡四個安息日講道.
然後突然間當地的人聚集一些市井的匪類.
就要來攻擊保羅.
所以有人就叫保羅快點走.
保羅因為這個原因就立即離開了.
可能連一個分手的機會都不能夠.
就已經立即離開了.
就輾轉去到比利那.
後來就去到安全了.
就去到哥林多的.
所以「帖撒羅尼迦前書」就在哥林多那裡寫.
就是告訴我們.
我們極力想法子.
很願意見你們的面.
所以我們有意到你們那裡.
我保羅有一兩次要去.
接著就留意了.
他說「只是撒旦阻擋了我們.
我們的盼望和喜樂.
並所誇的冠冕是什麼呢?.

$^{161}$豈不是我們主耶穌來的時候.
你們在他面前站立得住嗎?」.
似乎有一點要解釋他突然離開.
是不是因為懼怕艱難逼迫而突然離開.
有一種懼怕.
所以他勉勉強強很想來到去.
比較委屈地和「帖撒羅尼迦」的教會弟兄姊妹說.
其實我們是沒有任何的退縮.
我們期望你們得到我們可以誇耀的冠冕.
你們能夠在主耶穌來的時候.
你們能夠站立得穩 站立得住的.
所以保羅在這裡提到一件事是很特別的.
他有什麼的意識.
他知道是撒旦的攔阻呢?.
是不是敵人 那些反對封印的人.
就是撒旦的攔阻呢?.
還是真的有超自然的介入.
使得保羅清楚這次是有魔鬼的阻攔呢?.
這裡就不太多有這樣的教導.
不過「帖撒羅尼迦」前書.
都有很多這樣的超自然的教導的.
又有主在來 又有天使掌等等.
所以「帖撒羅尼迦」前書.
是不缺少這些靈界的encounter 或者這樣的經歷.
和將來的顯現的.
不過我們就不是在這裡停留.
告訴我們 保羅很緊張他的去留.
很緊張他突然間離開 他要安全了.
他要解釋給弟兄姊妹知道.
免得那個福音因為他突然間離開.
就受到影響 或者打了折扣.
所以保羅就很重視自己的誠信.
這個就是保羅他說是撒旦阻擋.
使得他們不能夠 這次就不是他去哪個地方.
而是他不能夠留久一點.
在「帖撒羅尼迦」這個地方.
這個城真的是非常漂亮的.
今天就是歐洲的希臘的北面叫做薩羅尼加.
很漂亮的地方 那裡有一個很大的像.
就是亞歷山大大帝的一個很威武的銅像.

$^{201}$因為這個城「帖撒羅尼迦」.
就是亞歷山大大帝的父親來建的.
所以是紀念亞歷山大大帝的妹妹.
所以這個城是相當出名的.
保羅不能夠繼續留在「帖撒羅尼迦」.
是有撒旦的攔阻.
我們回頭看剛才我們的經文.
這裡的攔阻 如果我們看上下文.
就不是撒旦 反而是神的攔阻.
因為第22節之前 保羅提到應該是第19,20節.
他說甚至 或者再補上一點點.
第18節 「除了基督藉我作的那些事.
我什麼都不敢提 只提他藉我言語.
作為用神跡歧視的能力 並聖靈的能力.
使得外邦人信服」 接著是第29節.
他說「甚至」 這裡我們上一次也說過了.
「甚至我從耶路撒冷 直轉到以利利古.
到處都已經將基督的福音傳開了」.
就是說再沒有可以傳的地方了.
接著他就有這樣的志向 有些人做過的.
就是說他已經走遍過一大地方了.
似乎是按照著聖靈一直的引導 神的旨意.
尤其是一個大的方向原則.
因為他很清楚自己的一個大方向.
可能大家都知道的 他就是外邦人的使徒.
所以他不會特別留在猶太人的地方.
他會去一個羅馬的城 歐洲的地方.
他會先去猶太人的會堂去傳道.
猶太人不相信 他就先離開 向外邦人傳道.
因為他很愛惜自己的猶太人.
但是他最終的重點.
都是要使得外邦人認識耶穌基督的福音.
這個是他自己很清楚他人生的那個照明.
當然不是說有猶太人問他福音 他不說.
對不起 我是外邦人的使徒.
你找彼得吧 這樣就沒有這樣的事情了.
所以我們要記得 保羅他清楚自己的那個照明.
所以這就告訴我們 當他由耶路撒冷.
一直去到以利利古 雖然我們不是確切知道.
以利利古到底是哪一個確切的地方.

$^{241}$但是似乎涵蓋著一個很大的區域.
他在這裡告訴我們 有一個由他領受福音.
在耶路撒冷使徒們接納他是使徒.
相信認定他是主耶穌的直接揀選.
接著由耶路撒冷開始就去傳道.
輾轉保羅就繼續去到不同的地方.
就是有神一直帶領他.
雖然他在這裡沒有說幾次的宣教旅程.
但是他說了一個很大的 很宏觀的.
他過去在寫這本書的時間.
他做過的那些福音工作 他涵蓋的那些地域.
所以在這裡的時候 他也都說過.
他有一個原則 如果有些地方是已經有.
譬如有神的先知 有譬如耶路撒冷的使徒.
或者其他人在那裡已經做開了.
他就不會插手進去干預 或者去打對台.
他反而有這個志向.
就是要那些還沒有得到福音的地方去傳揚.
所以大紀元22節當你引用了《二丈教書》52章.
講了他自己的心跡之後.
接著他就突然間說 多次的難阻.
不能夠去哪裡呢? 就是去羅馬.
就是說 原來保羅可能有很多次.
有計劃順道去羅馬 但是都是受到難阻.
所以這個難阻 似乎是一個善意的.
是那個優先次序 哪裡需要福音.
保羅就被感動 似乎這個是聖靈對他的難阻.
就來到他感動 他就去到那個地方.
就是說 意味著這個難阻.
就是我們今天大家都要學的 信服聖靈的引導.
追隨內心的那個引導.
其實很多不信的人 一些很有成就的人.
會追隨他內心那個人生的職場.
他的人生的抱負 或者他的終身職業.
或者要讀什麼的學系呢?.
有些人會追隨內心的.
比如有些電視明星 電影明星.
我忘記名字了 都出名的 做什麼電影呢?.
忘記了 《無間道》也有做過的.
那個內地明星 不知道什麼名字的.

$^{281}$可能有些人知道的 我看過他的訪問.
他說他是律師 他是讀研律師的.
但是他的內心原來是喜歡演戲的.
所以排除了家裡人 周圍的朋友的阻撓.
他追隨內心就是去演戲的.
我忘記了什麼名字.
最終今天就成就了一個很出名的內地的電影明星.
相當出名的.
其實我想很多的人都是這樣.
追隨內心的感動.
這樣去成為一個事業.
本來就是看到神一路帶著他的旅程.
那些鹹的 可能是那些味精.
大家聽過利斯頸牧師在我們錦繡堂講道.
他說有一次他已經預備好講章 準備上台了.
上台準備差不多五分鐘就開始主席出來了.
但是內心有很大的一個感動.
其實是催逼 拜託你放下你的講章.
寫好了 要講哪一段經文.
掙扎了一兩三分鐘 這些神的僕人.
就順服上台 悄悄的講章.
就講約書蘭的蒙召.
沒有講章的 真是好像拉天才球一樣就講了.
很順利的講完那一堂道.
有些人就聽得很好.
不經意的就把這堂道給了沈祖堯教授.
就是中大的前校長來聽.
沈祖堯當時原來他就正在等待神.
是不是要他去接任中大的校長.
聽到利斯頸牧師講這堂道之後.
不斷膽怯 神要他同在.
對他好像度身訂做那樣 他都不在現場.
但是有人把這堂道已經成為了錄音了.
成為了一個罐頭 罐頭也原來這麼有益.
可以去幫到沈祖堯教授.
澄清他自己內心的一些思考 未必是掙扎.
就決心承接這個好像臨危的壽命那樣.
就接任了中大的校長 這些是蘭左.
很善意的神的帶領.
你放下然後講 我要你講的道.

$^{321}$其實神教牧人本身就應該是這樣的.
就是如果很完美的 耶和華如此說.
應該舊約那些講道是很簡單的.
因為耶和華如此說不可以加你自己的意思.
這樣講完就可以走了.
因為說什麼時候什麼時候就會有饑荒.
這樣講完就可以走了 好像約拿那樣.
但是當然了 新約就不是這樣.
新約就是我們要在我們人生消化了神的道.
我們每一次都要等候神今天要.
或者這次要我們講什麼信息.
盡我們的能力將那段經文又不可以強解.
然後在《上汗文解經》之後.
就用一些生命生活的例子來演繹那個道.
所以一定有很多的限制.
但是總的來說那個信息本身.
是有神要在當時有一個這樣的傳遞.
那個傳遞的人就是那個牧者.
所以普羅一直就扮演著這樣的被神引導去傳道.
不是說喜歡去哪裡就去哪裡.
所以如果一個墓碑飛驒的牧者.
其實未必是牧者的.
因為牧者應該有一個陀就是神的聖靈一直帶領.
他不是喜歡去哪裡就去哪裡.
所以普羅就被攔了.
這裡他好像講得很負面 其實不是.
《上汗文》是很正面的.
他是順著耶路撒冷 就是由他順著.
由他清清楚楚的被使徒們接納.
到他一直傳道 去到這個以利利古.
最重要就是到處宣揚了基督的福音.
就是說他由他信主一直都秉承著耶穌的吩咐.
差遣來去作這位使徒 這樣來傳道.
所以你可以想象到普羅一直都是受聖靈的引導.
有一段經文是相當戲劇性的.
就是大家都熟悉的《使徒行傳》.
就是那個所謂的馬其頓異象.
一開始信主的人未必很熟悉.
未必很明白 這麼奇怪的.
去哪裡都說不準確 第十六章.

$^{361}$《使徒行傳》的第十六章.
普羅來到特比 又來到路斯德.
在那裡有一個門徒 名叫提摩泰.
是信主之猶太婦人的兒子.
他父親卻是希利尼人.
路斯德和爾戈烈的弟兄都稱讚他.
普羅要帶他同去.
只因那些地方的猶太人.
都知道他父親是希利尼人.
所以那些背景 路加醫生是描述得很清楚.
這次提摩泰要行割禮.
就是有特別的原因.
因為他要向猶太人的叔伯兄弟來傳道.
所以他自己如果都是一個未受割禮的提摩泰.
其實就很難令到那些猶太的同胞.
對提摩泰的道有一個信服.
就是說你自己都沒有做好一個猶太人的身份.
雖然你是半個猶太人.
你的媽媽 你的外婆是猶太人.
但是你爸爸是希利尼人.
所以普羅很有這種感覺.
就幫他行割禮.
一個年輕人 不是猶太裔的純種.
但是普羅要帶他傳道.
都懂得看這些禮儀的.
都懂得看這些別人對他的門徒.
那種平頭等足的.
所以普羅都很敏銳猶太人.
他們的傳統和戒心.
所以普羅就勉強地行割禮.
接著他們經過郭城.
把耶路撒冷使徒和長老所定的條規.
交給門徒遵守.
就是第十五章所說的禁戒偶像的食物和血.
和勒死的生畜和奸人.
這就是普羅把耶路撒冷的條規.
就是那個大公會議的結果的結論.
來頒佈的意思.
在猶太的社群當中.
和在各個教會頒佈.

$^{401}$因為有很多猶太人在不同的地方居住.
就要讓猶太人知道.
不需要再守摩西律法了.
接著就把戒心和門徒遵守.
於是宗教會第五節.
於是宗教會的信心越發堅固.
人數天天加增.
因為澄清了很多的誤會.
外邦人要不要受割禮.
那些祭偶像之物.
或者其他的泛民辱節.
應該怎麼處理.
如果更正了,更新了.
就會有更多人的信心堅固.
就有更多人來信主.
帶來的就是一個復興.
所以真理的辨明.
是帶來很大的祝福和益處.
接著開始講那個很奇妙的聖靈的阻撓.
第六節.
聖靈既然禁止癱瘓.
癱瘓就是保羅和提摩泰.
禁止癱瘓在亞細亞港.
就是不讓他們在土耳其那一帶再繼續停留.
他們就經過佛雷加.
加拉泰一帶地方.
是較為亞洲的北一點.
差不多是高加索地區.
就是那些亞太大疆.
再上一些的地方等等.
差不多再上哈薩克,哈薩克斯坦那些地方.
接著到了美西亞的邊界.
他們想要往秘推尼達.
秘推尼達再高一點.
耶穌的靈卻不許.
他們就越過美西亞.
接著下去特羅亞.
特羅亞就是一個港口.
就是過去對面就是那個菲律賓.
在夜間有異象現與保羅.

$^{441}$有一個馬其頓人站著求他說.
請你過馬其頓來 幫助我們.
馬其頓人是歐洲人.
不是亞洲人.
因為特羅亞是土耳其西面的一個港口.
一個港口城市 特羅亞.
大家讀過《木馬圖成記》的伊木馬.
特羅亞就在那附近.
所以就是那個位置.
你可以想像到.
這個耶穌的靈禁止.
反而耶穌的靈禁止.
夜間又有異象.
就雙重的印證.
保羅就很清楚地說.
我們既看見這異象.
到底是向保羅.
只是言語保羅.
其他人看不到.
保羅既看見這異象.
我們隨即想要往馬其頓去.
以為神召我們.
傳福音給那裡的人聽.
所以他對於耶穌的靈阻止他上美西亞.
保羅和提安蟹沒有發脾氣.
沒有去搞錯 我去傳道你都不准.
到底那個難處是怎樣呢?.
可能只是心靈的一種感動.
一種阻礙 一種不舒服等等.
我很早期讀過大衛·李文斯敦的見證書.
就是這位非洲的國父.
非洲宣教之父.
他是一個探險家.
他是很喜歡去探險和旅行家.
所以整個非洲.
他是第一個去到非洲的內臟,心臟,地帶.
去探索到底有什麼山村,有什麼地理.
他去做繪畫那些地圖.
但他同時間最大的,最重要的職業.
就是他很有負擔向非洲人傳道.

$^{481}$英國的警察會警察他去David Livingstone.
一路去的時間.
有一次有一個很突然的危險.
就是有一隻獅子突然在他後面撲向他那裡.
要來準備殺他 基本上一定會死的.
但是他當時不知道什麼原因.
就轉身 他拿著一支獵槍.
因為他也是獵人.
因為危險 因為真是很大的非洲.
你想一下是多麼的野外的地方.
但是不知道什麼原因.
同一時間在英國有一個家庭.
很敬虔的一個地方的人.
忘記了是爸爸媽媽還是什麼.
突然間晚上很感動 睡不著.
不知道什麼原因 馬上晚上起來祈禱.
禱告裡面說要為在非洲的David Livingstone祈禱.
他祈禱 祈禱到一個地步 平安了.
就睡著了 那件事就過了.
接著他們寫信過去 問什麼原因呢?.
原來很厲害的就是那段時間.
他就是那隻獅子上來準備要攻擊李維斯敦.
所以這些難阻 這樣的情況.
很多時候在宣教的旅程.
神就會做一個好像指揮燈那樣.
當然我們今天可以用WhatsApp.
但是有時候那種感動 聖靈來的感動.
如果我們能夠嘗試有這種敏銳.
不是敏感 是敏銳.
如果我們一直這樣每天活在一種很寧靜的裡面.
當然在做事的 但是實際我們內心有一種敏銳.
追隨聖靈其實是很寶貴的.
有時候真的我們看到一些上一代的神的僕人.
有時候他走過去跟那個朋友講福音.
他就是有感動的 就可以跟他說.
那個人就會相信 很特別的.
所以我們其實要學一下.
我想現代21世紀的基督徒.
我們可能欠缺的未必是做事的那種力量.
是欠缺靜下來的那個安靜.

$^{521}$對神聖靈的敏銳.
所以我們不應該被靈因壟斷了.
或者好像專利了.
一講到聖靈 一講到靈因 就是靈因派那樣.
就不應該是奇奇怪怪的謊言.
就應該是好像保羅一而再聖經提到的那種感動.
引導 攔阻 耶穌的靈禁止.
我們這樣去隨從.
我們的內心很多時候就貼近聖靈.
所以這些其實是很重要的人生操練.
《羅馬書八章》我們之前講過的.
就是體貼聖靈的就是生命平安.
基督的靈如果住在我們裡面.
身體就因罪而死 心靈就因義而活.
如果耶穌的靈住在我們裡面.
那個必死的身體就會因為這樣而活過來.
這裡告訴我們 耶穌的靈禁止保羅.
所以保羅就很清楚.
有那個馬其頓的異象.
他就不是日思夢想.
因為他一直都是想去傳道.
但是去哪個方向呢? 原來不准北上去梅西亞.
他就下到去特羅亞.
下到特羅亞其實都已經有一些訊號了.
他都是去到港口了.
原來保羅都不是蠢的 他不是再回到內陸了.
因為神禁止他在亞細那港島.
所以他就去到特羅亞 似乎是有所等候.
就是說他的心都很想去歐洲了.
很想離開土耳其大陸去歐洲了.
所以他就停留在特羅亞.
那晚就見到異象 就有一個馬其頓人.
所以很多的解經家都說這個馬其頓人.
就是這個勞加醫生.
就去求保羅過來拯救我們的族人.
這個就是保羅他的難作.
我們講多了一點 用多了一點時間.
接著就告訴我們.
但是現在在這一帶再沒有可傳的地方.
而且在這麼多年裡面 他迫切想去羅馬.

$^{561}$終於他現在有一個如願以償.
因為寫羅馬書的時候已經是第三次的旅程完結.
其實現在如果我們再結合《聖誕傳》20章.
就告訴我們 原來我一直都沒有開大電視.
所以等一下會看不到我們的Part1.
告訴我們 其實他一直在這幾次旅行的宣教.
其實已經來到哥林多 已經走了一個大彎.
就是已經兩次在阿西亞 就是土耳其.
第二次就已經過了馬其頓.
這次第三次再探望的時候.
已經再去彌爾忽所 再過來哥林多這裡.
所以這裡告訴我們 保羅現在要去耶路撒冷去.
這個就是我們上次都提過的.
就是在哥林多後書講了.
在馬其頓宗教會極困難的時候.
都因為有一個很清楚的.
由耶路撒冷下來的一個預言.
就是說天下會有大饑荒.
果然去到某個年間就真的出現了.
所以這個先知阿加布作的預言.
所以耶路撒冷就有很大的饑荒.
所以是一個很合適的機會.
保羅就呼籲在外邦的教會信徒.
你們既然領受了神給你們的福氣.
就應該有份於對耶路撒冷教會的窮人的救濟.
所以第25,26節 24節已經讀了.
24節 盼望從你們那裡經過得見你們.
就是講回第23節那個.
現在是他已經各方面他已經做好了.
就是說他的內心現在真的很想.
按照那個心靈的引導就去羅馬了.
但是原來他現在有一個心願.
這個心願好像是一加一等於三.
因為保羅成為一個撮合的人.
讓外邦人和猶太人藉著一些金錢物質的捐贈.
就帶來這個愛的合一.
帶來一個很美麗的和諧.
人和人文化之間的合一.
原因是什麼呢?.
24節講了.

$^{601}$得見你們 先與你們彼此交往.
心裡稍微滿足 然後蒙你們的送幸.
不過接著就說了第25節.
但是如今 這個是新版本.
我們舊的版本就是 但現在我往耶路撒冷去供給聖徒.
因為馬其頓和阿哥拉那人樂意的.
湊出捐項給耶路撒冷聖徒中的窮人.
這個第26節就是說.
讓我們知道原來保羅這次要上去耶路撒冷.
供給聖徒的需要 為什麼呢?.
因為外邦的馬其頓.
馬其頓是一個很大的地方.
就是今天的馬其頓國家.
還有阿哥拉那 阿哥拉那就是.
我上翻譯不行了.
阿哥拉那就是哥林多省.
阿哥拉那省就是哥林多其中一個省會城市.
但是接下來就是北面.
今天還是屬於希臘的.
就是很美麗的地中海地方.
阿哥拉那省就是哥林多的地方.
所以現在是馬其頓省和阿哥拉那省的人.
就樂意湊出一些捐款給耶路撒冷聖徒中的窮人.
所以有時一些天災.
有時我們不想看到的事情.
是會帶來很多美麗的故事.
我們看到保羅想讓這件事有以為之.
因為他自己解釋的 27次就解釋了.
他說這固然是他們樂意.
他們就是馬其頓省和阿哥拉那省的歐洲人.
裡面可能當然有馬其頓人.
有哥林多人 或者有希臘人.
有其他的人 雅典等等.
這裡說 固然是他們樂意.
就是說他們是很受感動.
保羅一呼籲 他們就接受 就響應.
但是留意 其實也算是所欠的債.
因為外邦人既然在他們屬靈的好處上有份.
就當把養生之物供給他們.
讓我辦完這件事 將這個禪禍.

$^{641}$舊的翻譯就有禪禍 新的翻譯就說捐款.
把這禪禍向他們交付明白 他們就是耶路撒冷的.
可能是那些帶領的使徒 或者領袖.
我就要路過你們那裡 往西班牙去.
就是接觸他在第24節的那個心願.
但是他補充了這樣的 現在臨時起義.
要上耶路撒冷 就是帶著一個使命.
這個使命就是外邦人的捐款.
未必是很多的 就上去耶路撒冷.
要來去救濟在耶路撒冷信主的窮人.
就是說可能那個本身的救濟的那個機制.
已經不受這些信了主的基督徒來有份的了.
可能他們被趕出會堂 可能被自己的父母.
如果堅決信耶穌 就要被自己的父母.
據聞就是要幫他出殯 當他死了.
就等於今天的回教徒 如果你信主.
就有權被殺的了 所以可以是很嚴重.
所以一定要由教會去節制這樣的窮人.
所以阿保羅很想將這一件善果的事.
讓外邦人有份 讓一些信了主的弟兄姊妹.
深深的領受 沒錯.
每一個人都是直接從耶穌的救恩得救.
每一個人都是神的愛,憐憫,信主.
但是我們從來都不是直接將臉射向天空就得救的.
我們是透過不同的歷史,不同的人,不同的文化.
我們才成為今天的人 所以神都容許舊約.
一路猶太人醞釀彌賽亞的降生等等.
所以保羅告訴我們 我們外邦人是欠了.
他們屬靈的好處 我們是正在享受.
他們現在失落的屬靈好處的那個份.
所以原來我們本來就沒有屬靈的好處.
但是神其實是愛世人 所以救恩都輪到外邦人了.
所以我們就不要有一種這樣的欺師滅祖.
或者以為我們是獨立於猶太人的.
沒有猶太人我們都可以得救的.
這樣就變成我們有時候忘記了歷史的那個進程.
我們就忘記了 其實的確猶太人.
好像為我們預備彌賽亞和十字架的救恩.
和神的預言的救世的福音.
所以其實我們應該對每一個跟我們講福音.

$^{681}$對我們有恩情的人.
更何況就是猶太人是直接.
他們是這樣將主耶穌去孕育出來.
一路的預言是這樣應驗.
雖然他們不接受 不代表神的應許是落空.
這是11章的羅馬書這樣講.
所以保羅現在反過來就是.
他以一個好像大家長這樣的角色.
他就鼓勵外邦人.
你們看到現在耶路撒冷有貧窮 有饑荒.
他就叫他們聚焦在耶路撒冷的聖徒中的窮人.
這樣來去救濟.
就是說帶上去耶路撒冷的教會.
救濟那些信主的人.
因為國家猶太人自然就會.
救濟他們自己的會堂 他們自己的猶太人.
所以保羅在這裡就不是要分猶太人,外邦人的意思.
而是反而就是想拉近外邦人.
很具體地有一個物質的支援.
猶太人現在的短期的應急的急需.
所以就這樣發動這個感動.
就是說雖然現在保羅就不在羅馬.
其實他沒有要求羅馬的教會 因為還沒有去到.
所以他已經收到馬其頓和阿該那的捐項.
到時帶誰去呢?.
有幾個弟兄和提摩太郎都有.
就上去耶路撒冷.
這裡告訴我們 所以我們求主讓我們能夠看到.
我們真是一個欠債的人.
我們欠了 當然一定是神.
所以不要和神計較 我們欠了神的債.
另外就是我們對猶太人 我們都算是所欠的債.
所以我們就應該將養生之物供給他們.
就是我們能夠有餘的.
真是來做一個感動的奉獻.
來祝福這些在困難裡面的人.
所以救恩是源自猶太人.
外邦人現在是在分享屬靈的恩典.
所以十字架 保羅就真的很透徹地認識.
所以各位心態是很重要的.

$^{721}$如果我們有一些門戶之見.
根本就不會有羅馬書出現.
如果保羅一直都站在自己原原則則本本的.
一個猶太人的拉比.
根本就不會有外邦使徒這樣的.
豐功偉績的人生的事業.
如果他慢慢就是遷怒於猶太人對他的迫害.
他完全站在外邦人那一邊.
也都不會構成他關心耶路撒冷的窮人 猶太人.
所以他整個骨子裡面充滿著的是猶太人的血.
但是他的神學,他的心靈.
是完全被十字架的愛融化.
所以能夠融會貫通的一個牧者.
某程度上是當今世界.
都不用什麼世界領袖.
牧者如果能夠胸襟這麼宏觀.
世界就和平了.
因為你想一下彼此之間沒有那些門檻.
沒有那些人為的那些尼巴.
其實可以帶來一個很大的跨文化的合一.
所以求神幫助我們.
能夠看到保羅這裡絕對不是用心良苦.
其實是很美麗的捐獻.
你想一下耶路撒冷的窮人.
基督徒信主的 現在保羅帶著他們的捐獻上來.
不如我們看看哥羅德侯書那幾個弟兄.
讀給大家聽 因為聖經期間這樣記載.
就是說都很想我們對他們的做法.
記一筆那個紀念.
第九章沒錯 我先看回那個人的名字.
對不起 現在是第七章尾才是.
先看看這裡 第七章的十四節.
如果我若對提多誇獎了你們什麼.
也覺得沒有羞愧.
因我對提多誇獎你們的話成了憎的.
正如我對你們所說的話也都是憎的.
並且提多想起你們眾人的信服.
是怎樣恐懼戰經的接待他.
他愛你們的心腸就越發熱了.
我如今歡喜 能在凡事上為你們放心.

$^{761}$這個就是提多回到保羅面前誇讚哥羅基督教會.
原來對保羅所猜派去的提多是有很大的接納.
也讓提多在他自己很優秀.
似乎好像是患了病 也來對他的接納.
所以提多在保羅面前的誇讚.
或者保羅看到提多的表現.
對哥林多教會也有一份的放心.
接著就開始講那個捐獻.
就是講馬其頓教會 宗教會.
知道耶路撒冷這個大的困難.
所以就來到有一個樂捐的後因.
提多就要開手辦這個慈惠的事.
提多當時就已經在哥林多了.
辦好之後 但是第八章告訴我們.
整年都還沒有辦好 都還有這樣的缺漏.
所以有這樣的催逼 保羅就要寫信.
來對哥林多教會有一種類似的指責.
但是不過是透過其他馬其頓人的慷慨.
來鼓勵他們 激發他們的熱心.
接著去到第九章 第一節開始.
「論道功及聖復的事 我不必寫信給你們.
因為我知道你們樂意的心.
常對馬其頓人誇獎你們.
說阿乖阿仁 預備好了 已經有一年了.
並且你們的熱心激動了許多人」.
接著留意了 「但我打發那幾位弟兄去.
要叫你們照我的話預備妥當.
免得我們在這事上誇獎你們的話就落空了.
萬一有馬其頓人與我同去 見你們沒有預備.
就叫我們所確信 反省要收賄.
你們的收賄就更加不用說了.
因此我想不得不求那幾位弟兄.
先到你們那裡去 把從前所應許的捐資.
預備妥當 就顯出你們所捐的是出於樂意.
不是出於勉強」 一直提到那幾位弟兄.
是不是哥林多教會所認識 所認受.
以及極之有保羅所信任的.
因為的確要帶這些捐項上耶路撒冷.
的確是要講求那個誠信的.
所以這個也是要很留心的.

$^{801}$有時候我們去安息禮拜.
家屬我們就很清楚講明的.
所有的服侍的 我們是一毛錢都不會收.
不需要的 我們是義務做的.
當然那封利是就20元那些.
那無所謂了 一般傳統我們都要的.
但是我們也會加多一句 如果你奉獻給神.
我們是沒有問題的 按感動.
有時候他們會問 麻煩你給我銀碼.
我們說不給的 因為按感動自由奉獻.
是為神的聖光福音.
往往他們就會用信封記錄現金.
方便一點就是 麻煩你在信封寫銀碼.
以及我們可以修條的名字.
那就變成了 如果是封了我們就不拆了.
如果是打開 我們就大家點一點收 就會好一點.
這個就是始終有時候我們怕有時候人的誤會.
或者過失 或者甚至乎應該是試探.
所以最好就是這樣就會好.
以及我們講明了一定會開收條.
我們就拿回來給教會的行政同工處理.
這些這樣的誠信 特別是處理這些的獻金.
就更加要小心 因為我就是圖章收了.
真是沒有人知道的 如果那個人說.
你不用開收條的了 我們信任教會的了.
有時候就有這樣的試探.
所以你要不要收條是你的事.
我們一定開收條.
就是說我一定會將這些錢拿回去給教會.
以及寫到明 譬如是五千塊 在信封上註明五千.
如果是封了的 那我就不需要開了.
所以這幾個弟兄是有保羅信任的.
他們去辦理 保羅說如果這幾個弟兄來到.
見你們什麼都沒有預備.
在現場在那裡嚥嚥噠噠地湊錢.
保羅就覺得是一種羞愧.
你們羞愧就更加不用說了.
因為保羅說 我在馬其頓阿哥亞省的面前.
對不起 在馬其頓省面前.
來誇獎你們阿哥亞人的熱心.

$^{841}$還有不要忘記 阿哥亞到今天還是.
非常富庶的歐洲南部.
北面是比較窮乏的.
所以保羅不是沒有道理的.
哥林多是大都會 又是說那個運河.
就是那裡有一個很特別的地方.
可以通過亞得尼亞海.
直接過去意大利那邊 過去地中海.
不用經過一個大的地方這樣轉折的.
所以哥林多真的有一個很重要的經商貿易的地方.
這裡告訴我們 就是那個印度.
所以保羅現在就說 所以透過這些的捐行.
其實是激發馬其頓人來服侍猶太人.
就是要將十字架的愛帶來合一.
所以都盼望弟兄姊妹.
我們人生都希望能夠怎樣將那些合一促成.
將來一定有一樣東西一定會很憤怒 很審判的.
製造麻煩 製造分裂的人 一定面對審判.
因為神就是愛 神就是合一.
神就是和平的神.
所以製造患難爭端的人 是神所恨惡的.
所以我們人生就是希望怎樣能夠謀求那個合一.
盡力在食靈裡面 以不鎖書四章.
竭力保守聖靈所賜合而為一的心.
一切都是由我們信著.
神就已經給了我們這個合一的靈.
所以我們要竭力的保守.
就是說盡我們的能力.
在教會裡面不要來將教會分裂.
而是將教會拉在一起.
彼此的合一 在愛裡面的合一.
這個就是神一定會祝福的.
那個時空一定會成長的.
人的心一定會有那份和諧的.
所以每一個人 如果我們好像一個黏合劑.
就很重要的了 我們感恩.
洪恩我們有一位姐妹.
她去到哪裡都好像一個黏合劑那樣.
現在去到澳洲了.
在澳洲的教會也是這樣很熱心的服侍.

$^{881}$洪恩星期四就是載老人家去飲茶.
每個星期還要給錢請老人家飲茶.
每個人幾十塊錢就行了 任他們吃.
每個星期都給老人家一個很開心的聚會團契.
這些黏合劑就帶來教會的合一.
你可以想象到十字架本來就是有這樣的功效的.
就是要讓到甚至文化當中的合一.
所以外邦人 一點點錢就可以收買猶太人的心.
猶太人都感動 和外邦人一起祈禱.
不再是你的安息日 我的自由日.
你的不可以吃的食物 我喜歡吃什麼就吃什麼.
沒有這些問題了.
所以其實某程度上已經將第十五章的問題解決了.
就是在十字架的合一裡面.
就是用一點點錢就已經解決了這些很大的傳統.
很大的禮儀 甚至猶太人看為是十界的界條.
都已經在外裡面成為了一個合一.
你想一下有多麼的寶貴呢?.
等我辦完這件事.
我紅色那裡就是一個捐項.
大陸原文就清楚了.
是有一個特別的article(文章).
就是那個善果 這個我已經不懂得讀了.
就是那個果實 它有一個果實.
所以看上下文 果實肯定是在講善的果.
就是說是那些獻金 上文是在講這件事.
所以就不是單純獻金 是良善的果.
就是說 本來看我們這些的 這個就是新的和盤.
把者不捐款 其實就是舊的和盤翻譯得很好的.
就是善果 交付給他們.
我就要路過你們那裡 到西班牙去.
所以弟兄姊妹 我們的獻金是一份善果.
在神眼中 其實中文這個字本身都很好的.
是你內心有善 你才會想起別人的需要.
所以我們常常都有一個善良的心.
有神的感動 我們常常在我們真的有餘裡面.
常常有準備 可以去幫有需要的人.
這些是善果來的.
不信的人都知道自己做善長仁翁是很好的.
所以他們當然看來是積福.

$^{921}$但是我們是坐在神裡面.
我們 保羅看這個是果 善果.
所以不是一個捐獻捐款.
好啊 施捨給你 不是這樣的.
所以原來保羅都用這樣的字 是一個果子.
所以這個果子就一定是捐獻的金錢.
所以金錢是怎樣用的.
如果能夠用得其所 是一個很大的祝福.
接著就應該去到那個民安 保羅的民安的心腸.
他說第29節 我也曉得去的時候.
必帶著基督豐盛的恩典而去.
去的時候 就是他去探望羅馬教會的時候.
他交付了這些的獻金 其實是奉獻.
外邦人的奉獻 可能都有一些手足.
是要換回猶太的貨幣.
來給這些耶路撒冷的窮人.
但是那些是技術性的.
但是保羅看這個這樣的詞.
這一筆收的獻金是一個善果.
所以原來金錢可以不是萬惡的.
金錢原來可以是很有用的.
所以能夠賺到錢的弟兄姊妹.
用在神的裡面是一個很大的祝福.
所以我們要記得 保羅就是這樣理解.
所以他自己既然教人.
就是說他裡面的內心有豐盛基督的恩典.
他隨時都可以去跟外邦人.
將他所得到的真理去祝福他們.
雖然保羅一定不會很有錢.
他可能有足夠的隨行的人員的門徒.
他有足夠的可能基本的飲食等等.
但是他肯定不會是一個很有錢的人.
但是他有的是一幫很忠心的門徒.
跟他視死如歸的去傳道.
這些不是開玩笑的 超過60歲.
如果你都不是一個mentor.
就是很大有限的 沒人喜歡跟你.
已經是一個問題了.
保羅有很多的真是忠心的門徒.
很快讀幾句給大家聽聽.

$^{961}$你就羨慕到不得了.
你看看保羅他有的人.
和他一起的第20章聯定之後.
保羅請門徒來勸勉他們.
就辭別起行.
當時的亂象肯定是以弗所這個地方.
然後走遍那一帶地方.
用許多的話勸勉門徒.
然後來到希利尼在那裡住了三個月.
將要坐船往敘利亞.
猶太人設計要害他.
他就定義要由馬其頓回去.
就是說他仍然堅持要回去.
和他一起在亞細亞去的.
就有比利亞人 北羅斯的兒子 索巴特.
貼薩摩尼加人 阿里達古和西宮都.
還有特比人 戈尤 並提摩泰.
都是特比人 或者是路斯德.
又有亞細亞人 推基古 特羅菲摩.
這些人先走 再去特羅亞等候我們.
你可以想像得到 至少有七個年輕人.
跟隨保羅 還要是大的暴亂之後.
仍然沒有走開 離開保羅.
仍然跟隨保羅 現在是上耶路撒冷.
因為猶太人設計要害他.
他就定義由馬其頓回去.
在他現在這樣的傳道 不是一個亮麗.
不是一個光彩 不是上耶路撒冷.
給一個和平獎勵.
而是有一些極端的猶太教的分子.
要來逼迫保羅 但是年輕人清楚.
這位是神的僕人 愛邦人的使徒.
就是說是我們愛邦人的神.
給我們的一個很大的祝福.
使得我們愛邦很多的地方信主.
現在他們又清楚自己的詔命.
跟隨保羅坐船上.
所以坐船很快的 地中海一過就是什麼.
就上去亞西亞的登陸.
或者是安提拉 來跟隨他.

$^{1001}$意味著保羅他的人生是祝福了很多的人.
成為牧者傳道 跟隨他去接第三代的棒.
這樣來傳道 所以保羅有豐盛的恩典.
其中就是人力的資源.
所以一個成功的牧者 應該有很多的門徒.
接棒人 其實是很重要的.
因為工作那個人可以沒有.
但是那些事工是繼續的.
所以神會興起那些人來去接續的.
所以摩西是有約書亞 以利亞也有以利沙.
當然還有以賽亞 也有他的兒子.
或者其他的人 去到耶穌的時候.
他後期的傳道就訓練12個門徒 親自揀選.
所以保羅現在都有他7個的.
只是剛好跟他在一起.
有些已經分開去到別的地方了.
那你就可以想象得到 剛才提到的20張.
還沒算寫這封信的路加醫生.
一直追隨著保羅 是他好像宣教的一個代筆.
那你一直想象就是 保羅基督豐盛的恩典.
就是他的生命 他燃點著很多那個時代的人.
尤其是外邦人被他燃點 這樣來跟隨他的傳道.
所以這個是保羅的六樣的恩典.
還有他那個服侍宣教策略的深長.
接著一個竹緣了 第三十節.
他的禱告 又很值得我們去看一看.
剛剛還有這一段 16張是一個全新的一張.
所以是要開一個獨立的課來講的.
因為相當重要 保羅的民安裡面.
第三十節.
並叫我順著神的旨意 歡歡喜喜的.
到你們那裡 與你們同得安息.
接著三十三節 就是賜平安的神.
賞和你們眾人同在 阿們.
這樣一段很簡單的祈禱 都已經有幾樣東西.
保羅他是一個很謙卑的來代求.
他說 我藉著我的主耶穌 藉著聖靈的愛.
勸你們要竭力的去祈求 為祂祈求.
一同的竭力 就是說保羅現在不是很膽粗粗.
好像很興奮 現在要上耶路撒冷.

$^{1041}$沒事的 或者我根本不怕裡面的人.
他是說要竭力的祈求 又要為祂去祈求神.
這樣使得他能夠脫離猶太不順從的人的攻擊.
你可以想像得到 保羅感受到.
是不是一個生命的威脅.
還是在他內心的痛.
猶太人真的會有一些硬心 極其的硬心.
保羅自己說的 那個橄欖樹被砍下來.
橄欖樹的樹幹還沒砍斷.
但是現在接上的是野橄欖.
就希望那些橄欖原本的可以發芬.
就是說有幾分硬心的猶太人.
所以保羅是很深的感受到.
原來福音去到這些猶太人當中.
是帶來攻擊 帶來反感 帶來的敵對.
帶來的針對保羅 你欺師滅祖的這樣的罪名.
保羅來到去踐踏神的殿.
帶外邦人進入神的殿.
這個是《大羅飛魔》 這是21章的《聖行傳》.
原來保羅是深深的在他有生之年.
都要面對這些敵對者的.
所以保羅看這些不是他的仇人.
是要祈禱脫離這些人的攻擊.
使他能夠順順利利的帶這些捐項.
而且被聖徒接納 沒人嫌錢腥的.
嫌這些錢不是捐項就會有.
所以要做得很有智慧 要做得很有神的愛.
不是施捨給猶太人 不是現在外邦人厲害了.
是外邦人在猶太人的屬靈的墳上面現在有墳.
所以外邦人看猶太人是自己屬靈的骨肉之親.
所以是這樣的奉獻 所以被聖徒接納.
原來不是用一種施捨的心態.
扔過來的心態 不是這樣 是一種愛的接納.
所以錢只不過是那個途徑的意思.
真正的是在愛裡面的合一.
所以有時候教會 為什麼有時候年輕人.
是很憤怒教會的虛偽呢 就是在唱合一的詩.
有時候他們真的看到教會的上一兩代不合一.
很多時候口唱 心眼都還很憤怒.
是很激氣的 尤其是他真的認識自己的家人.

$^{1081}$譬如有家族的 叫完家族才回來的.
你跟著唱 我們要合一 我們要牽著手.
所以很多中國人 香港人是不喜歡牽手的.
是不喜歡 包括我都不喜歡.
就不是我不想和別人合一.
而是真是這個文化 但是有時候越是這樣牽手.
而沒有心靈的合一 就真是會很虛偽的.
所以保羅說耶路撒冷的聖徒.
越立這些的捐行 就可以想像他們能夠.
從心裡面被主的愛感動 說現在神的救恩.
不單止只是我們來開了眼得救.
現在連救恩去到外邦人 神興起的仕途.
保羅成為那個代言 去到拯救外邦人信主.
現在他們明白救恩了 知道我們是神的兒女.
大家都是亞伯拉罕的子孫 大家同為神一體的.
這樣來彼此的幫助 那個獻金就來得離足珍貴了.
因為是弟兄姊妹奉獻給耶路撒冷的弟兄姊妹.
所以你可以想像到 保羅用的字眼都有這樣斟酌.
可蒙聖徒越立 這個就是第二個祈禱.
第三個祈禱 他都希望照神的旨意.
歡欣喜氣的去到羅馬 神的旨意就是.
上到耶路撒冷 給亞西亞上去的人.
在安提哥上去的猶太人.
見到他 這個應該是22章的《算盒傳》.
稍為我讀給大家聽 所以祈禱似乎不行.
不過神的時間沒辦法了.
第27節 第21章的27節.
「那七日章完 從亞西亞來的猶太人.
看見保羅在奠女 就聳動了眾人下手拿他.
喊叫說以色列人來幫助.
這就是在各處教訓眾人.
糟踐我們百姓和律法 並且地方的.
他又帶著希臘人進奠」 就是這個罪名.
污穢了這個聖地.
因為外邦人是不可以進入聖殿的內室的.
路加醫生就有解釋了.
「這說是因他們曾經看見.
已不所人德羅非摩和保羅在城裡.
以為保羅帶他進了殿.
合城都震動」 就是耶路撒冷.

$^{1121}$所以保羅這兩個祈禱 應該是三個祈禱.
只有一個可以應驗 就是他完了那個捐項.
他自己又有一個拿西人的玉 還了玉之後七天.
可能準備就開始落羅馬.
但是亞西亞的猶太人上來.
就看見保羅就捉了他.
因為看見德羅非摩在聖殿的周圍.
和保羅一起 就以為保羅肆無忌憚.
因為在外面已經很恨他了.
不斷傳那些叫人不要守律法 已經很生氣了.
所以一看見他 所以有時候一看見.
將過去那些積累的誤會.
都一筆這樣發出 就解釋不了了.
明明沒有的 明明只不過在外面看見.
就當作保羅曾經帶德羅非摩進去聖殿.
所以就是百詞莫辯 有時候就是火遮眼.
這些猶太人 所以我們就不要看太表面.
最好就問清楚當事人.
那就會解決很多的誤會 但是本身有成見.
所以這些攻擊的對象真的出現了.
保羅是不是神沒有聽祈禱呢?.
原來神有另外的計劃.
所以保羅就被捉拿 指著他本來想歡歡喜喜.
到羅馬探望弟兄姊妹.
然後由他們送行去西班牙的.
一直被綑綁 是不是歡歡喜喜呢?.
應該不是 不過都是神的旨意.
有可能更高的旨意.
希望他將來要在君王的面前.
作這個見證 所以神的禱告.
聽的程度和方式都不是我們人.
可以隨便的去說不定.
或者說神為什麼不聽我們這樣說呢?.
最後就是33節 有些考究.
平安的神與你們所有人同在.
這個直譯的中文.
就是原文的希臘文的意思.
和諧本也不夠全面.
他說願賜平安的神.
常和你們眾人同在.

$^{1161}$也OK的 眾人.
但是其實是強調所有.
Meta這個字 那個應該是M字.
Meta這個字就是所有的意思.
所以當然中文就很少這樣說 修辭就不漂亮.
常和你們眾人 代表所有人.
但是原文其實是所有人.
但是很特別的.
一而再中文的和諧本都要加個願.
但是其實很多說到平安的神.
從來都沒有祝願的.
就是說平安的神與你們所有人同在.
就是這樣而已.
就沒有說願神 就是希望神等等.
不用願的 古來的意思就是.
平安的神的同在就是平安的了.
其實整句意思就是與你們同在得到平安.
跟著是阿們 得到神賜的平安.
所以中文就很合理的.
願賜平安的神 常和你們眾人同在.
它就強調願賜平安的神.
但是其實原來基督徒的平安.
就是平安的神同在 就是平安了.
就不用賜給你的 神的同在就是平安了.
原來保羅有好幾段聖經都是這樣的.
那個原文都是完全一樣的.
就是《羅馬書》的15章33節.
跟著16章又再一次說平安的神.
16章的20節.
賜平安的神 快要將撒旦踐踏在腳下.
願我主耶穌基督的恩 常和你們同在.
再一次又是這樣說 賜平安的神.
但是在這裡有5段聖經新約.
保羅書信裡面都強調平安的神的同在.
就是平安的意思.
今天不夠時間再揭五段.
因為我怕都差不多時間了.
我們就停在這裡 我們一起祈禱.
再一次我們很感恩.
雖然是一段很短的經文.

$^{1201}$我們都可以來到一起去享受保羅的教導.
了解這位使徒.
他已經充分對十字架的愛.
尤其是合一這個部分.
是帶來人與人 人與神的合一.
充分的了解 明白.
以致在他的生命裡面.
他是成為一個橋樑.
要將猶太人帶進外邦人的心裡.
又將外邦人 當時他的世界就是歐洲.
就是當地的土耳其.
或者是整個歐洲大陸不同的人種.
當然有些很北的北歐 保羅都未必去過.
就是這樣帶到外邦人.
經常可以去耶路撒冷.
與猶太人有相交.
不再需要分開.
可以彼此相交交往 一起吃飯等等.
這些就是十字架的愛帶來的合一.
求你讓我們在香港.
我們真的能夠與不同的教會有合一.
我們能夠與不同的宗派合一.
與不同的弟兄姊妹合一.
我們與不同的堂 晚堂 早堂 午堂.
嘉欣堂 紅欣 孫德堂.
是我們同一個出處的 有合一.
使我們能夠深深被十字架的愛感動.
願你幫助 看到保羅就是這樣.
有一個熱心要來將福音帶進外邦人當中.
他的計劃是要去西班牙.
可想而知 其實當時他已經去了三次宣教.
其實是相當的疲乏和危險的.
但是這位使徒就是願意燃燒他的生命.
為神的福音 雖然我們再見不到.
在少林主之後的西班牙旅程.
願你帶領我們能夠在神的家裡面.
在不同的崗位繼續去孕育下一代.
有更多人信主 讓他成為門徒.
多謝你聽我們的感恩禱告.
奉耶穌基督的名.

$^{1241}$阿們.
\newpage



\section{羅馬書}
\label{sec:3rvmIgiASuQ}
\textbf{2026-02-13 周五晚查經班《羅馬書(三十七):分題:福音的實踐:從教義到生活》 講員:曾裔貴牧師}
\newline
\newline
連結: \href{https://youtube.com/watch?v=3rvmIgiASuQ}{\texttt{https://youtube.com/watch?v=3rvmIgiASuQ}} ~~~~ 語音日期: 2026-02-14
\newline
\newline
\hyperref[sec:LvE6qlAUC04]{\small{< < < PREV SERMON < < <}}
~
\hyperref[sec:index_chronic]{\small{[返順時目]}}
~
\hyperref[sec:index_scriptual]{\small{[返順卷目]}}
~
\hyperref[sec:4pLVUVJIfgU]{\small{> > > NEXT SERMON > > >}}
\newline
\newline
$^{1}$領旨每年晚晚我們可以一起來先有禱告.
天父我們在晚上 我們感恩.
因為我們能夠一路去研讀羅馬書.
雖然我們不是能夠逐節逐節地去看.
但是我們在每一個段落.
我們都能夠盡量去處理上下文.
和斯圖普羅他寫作的那個歷程.
那個動機和那個賞望.
成為我們大家生命的力量.
讓我們看到一些這麼重要的基本的教義.
罪和論法 猶太人和外邦人.
恩信稱義的這個福音.
耶穌基督的白白犧牲等等.
都讓我們能夠看到.
也都讓我們今天晚上.
當我們來到第十六章.
都有差不多26個人名.
是普羅的文安的人物.
裡面所組合的是相當豐富.
我們很快地看一看都知道.
是有不同的階層.
甚至乎有男性有女性.
這些都是對於當時代文安的.
這樣的羞恥這樣的安排.
是很難想像.
可以想像普羅是突破那個時代種種的限制.
所以盼望我們今天.
雖然我們都是2000年後的新約世界的人.
但當我們再去看再去研讀的時候.
讓我們能夠感受到普羅.
他內心因為教義帶給他的生命的釋放.
帶給他對人類那個大同.
那個合一 那個在救恩裡面的平等.
都讓我們牢牢地抓緊.
因為在我們人類的社會.
總是有不同程度的比例的那種歧視.
有時充斥在文化 充斥在人的內心.
盼望主幫助我們.
我們今天晚上再一次來等候在你的面前.
我們這樣祈禱奉耶穌基督的名.

$^{41}$我們上一個星期停在第15章.
最後的段落那裡提到.
普羅說他會去西班牙.
這個是第29節.
他說我又是曉得他去的時候.
必定會帶著基督封聖的恩免去.
提到普羅他的生命.
因為主耶穌的救恩.
和神在他生命裡面的照明.
就是呼召的這個積分.
使他能夠分享神在他生命裡面的種種的工作.
就是普羅的生命的豐富.
可以和其他特別是愛邦的教會的弟兄姊妹分享.
就是分享這個基督封聖的恩典.
所以16章的文案又給我們有很.
真是值得我們去深思.
普羅為什麼會認識這麼多羅馬教會裡面的人.
還有這樣這麼長篇的文案.
有別於其他的書信.
憤外顯出這本書的文案是來得很特別的.
所以我們往後都會嘗試逐段去看.
但是當我們再看的時候.
上一個星期我們講的還不夠詳盡的.
就是提到那個平安.
就是普羅來到他在15章的經文.
就是第30到33節 尤其是33節.
弟兄們 我藉著我們主耶穌基督.
又藉著聖靈的愛.
許你們與我同竭力 為我祈求神.
叫我脫離在猶太不順從的人.
這個是15章他提到.
他要將一些捐項在馬其頓一帶收到的那些獻金.
帶上耶路撒冷 帶上去給耶路撒冷的聖徒.
就是信了主的猶太人.
他們已經是 似乎已經是離開了會堂.
是又在新的教會生活聚集.
但是現在遇到一個饑荒.
就是說傳統的猶太教的那個接濟.
可能這一幫信主的猶太人就沒有份了.
但是當普羅知道之後.

$^{81}$而他又能夠感動馬其頓的教會.
來到他有這樣的一個需要.
有這樣的意象.
馬其頓教會就來到他提集那些獻金.
來到他要透過普羅和一幫弟兄.
就帶上耶路撒冷.
大概就是這樣的背景.
所以普羅就說要為我祈禱.
脫離在猶太不順從的人.
就是那些對猶太教相當執著的.
還是未能夠明白耶穌的救恩.
帶給人類的生命的改變.
所以仍然停留在猶太教崇拜的這些猶太人.
就一定會敵視普羅.
因為他是首當其衝.
勸人不需要再守摩西的律法.
所謂不需要再守摩西的律法.
就是信了主耶穌.
其實已經是完成了,超過了.
超越了摩西的律法書.
已經是直接可以跟神有一個永恆的生命關係.
所以如果猶太教的人不明白.
以為他叫人不要守律法.
其實就是有誤解.
所以普羅只能夠用愛來包容.
但是始終還是要面對這個直接的衝突.
就是當上到耶路撒冷的時候.
好像去到那個宗教的核心的範圍.
普羅就一定是要面對這樣的處境.
所以他現在就呼籲羅馬的教會為他祈禱.
到底那些書信.
因為這封羅馬書是由一個叫菲比帶去羅馬的.
到底那封信去到羅馬.
他們知道這個消息中毒完之後.
然後為普羅祈禱.
還是說普羅都等不及了.
自己都上了耶路撒冷.
那封信都還沒有到.
這個空檔期到底是誰先走呢?.
聖經又沒有交代.

$^{121}$但是無論如何.
普羅都很想羅馬的教會.
他們為普羅的行蹤來禱告.
因為他真的要上耶路撒冷有這個使命.
要帶這些外邦人的獻金.
外邦人能夠現在在神的恩典裡面.
藉著普羅的傳道者.
外邦的使徒來講述了.
外邦人現在都可以白白門恩.
就是這樣拉近和猶太人的那個關係.
所以這些獻金就來得離足珍貴.
因為是由外邦人信了主的外邦人.
來得有感動.
要來回饋猶太人的那個物質的需要.
所以這裡告訴我們知道.
普羅說.
「脫離猶太不順從的人.
也叫我為耶路撒冷所辦的.
捐項可蒙聖徒的月立」.
也就是說讓聖徒不只是看到錢.
所以錢的多少不是最重要.
是那個性質背後的那個意義.
就是外邦人蒙恩來用獻金.
用這些金錢來報答神.
在這些外邦人的生命裡面的那個拯救.
那個救恩源自猶太人.
所以普羅就成為了一個很重要的橋樑.
所以普羅有一個這麼寶貴的信息.
來向外邦人傳講.
叫外邦人可以因信耶穌.
就成為了屬靈的亞伯拉罕的子孫.
也可以在神的國度裡面.
神的約,神的永恆的那個應許是有份.
所以普羅的這個職事是非常的重要.
因為如果他不講.
外邦人沒有聽過這個信息.
外邦人就只能夠拜偶像.
否認神 繼續這樣來犯罪.
面對將來永恆的沉淪.
所以這個就是使徒普羅或者傳道者.

$^{161}$當你有一個很清楚神的呼召的那個職份的時候.
就要盡力去講.
因為那個生命就在那個福音的宣講裡面.
所以「要可蒙聖徒月立.
並叫我順著神的旨意.
歡歡喜喜到你們那裡.
與你們同登安息」.
這個是普羅的祝願.
能夠按照著神的那個旨意.
可以去到羅馬教會跟他們分享.
一同得著神賜的那份安息.
那份的平安.
接著就是「平安的神常與你們眾人同在」.
這個平安的神.
上一次我們就講到這一節 三十二節.
三十三節我們就還沒有講.
在整個的原文裡面.
就不是「願賜平安的神」.
沒有那麼複雜的.
就是平安的神.
這個是對神的這個屬性的形容.
平安的神同在 那個就是平安了.
所以就不是「願賜平安的神」.
神 你在神的裡面就是平安了.
因為祂是平安的神.
所以就是這麼直接的了.
希臘文 或者普羅的寫作.
十六章的二十節又有這樣的稱呼的.
就是羅馬書的十六章二十節.
「賜平安的神 快要將撒旦踐踏在你們的腳下.
願我主耶穌基督的恩 尚和你們同在」.
這個「賜平安的神」原文同樣都是「平安神」.
那就可以想像得到.
和以下書所描繪的「平安之君」是遙相呼應的.
所以有神 如果我們真的認識神.
真的在神的裡面 或者神在我們的裡面.
那個就是平安了.
那個就是我們得到的真正的平安了.
不知道什麼時候是我們最需要平安的呢?.
有沒有哪些弟兄姊妹.

$^{201}$你們經歷過的某些平安呢?.
有些人就說準備要進去做一個大手術.
在那個等候區冷得不得了.
那種冷不僅是冷氣.
而是由心裡面震出來的那種冷.
我想那個時候 我沒有試過.
那個時候就真的需要你要祈禱.
記住要讓平安的神.
籠罩整個冰冷的病床 那個手術室.
神的平安 神同在就是平安.
大家要記住 神的同在就是平安.
因為祂是平安的神.
所以就不是怨賜平安的神.
不需要這樣怨 不需要祝願.
當然中文的修辭也是在反映這件事的.
但是原文真的告訴我們.
平安的神同在就是平安.
所以這個平安就沒有任何東西可以奪取了.
因為沒有東西可以阻礙神的同在的.
甚至連我們的罪 主耶穌都已經為我們擔當了.
所以我們可以隨時.
當然我們不應該說邀請神來到.
因為做物主不需要我們邀請的.
相反就是我們自己的意識.
是我們自己的感官的感覺.
我們有時難不阻神的同在.
其實祂一直都是和我們同在的.
我們要記得這件事 其實很有趣的.
因為最近買了一本書 就是《死亡之後》.
有很多個案就是這個精神科醫生.
他在50年的行醫生涯裡面.
不斷去研究那些瀕臨死亡的人.
靈魂離開身體的那個經歷.
他說有一個年輕人.
他是有一些思覺失調的 都在吃藥的.
但是他居然說有一次有個聲音叫他上去跳樓.
他真的跟著上去 魔鬼說你是沒用的.
你是罪大惡極 你是人渣 你快點死掉吧.
這樣叫他 他聽這個聲音真的跳下去.
但是他說在那一跳之後的那個瞬間.

$^{241}$可能一萬分之一秒.
他說有一個是他聽得很清楚的聲音.
就是神跟他說 告訴他你不是人渣.
你是我所愛的 你是我的兒女.
你不需要懼怕 當然後來搶救 醒了.
後來他將這個經歷就跟這個精神科醫生去講.
不是直接講的 即是輾轉透過不同的人.
來去收集到 而他知道就訪問這個年輕人.
當然是死過活生生 如果不是就講不到了.
他就說有兩個聲音 一個就是魔鬼.
但是另一個更大的聲音就是剛剛我吃東西的時候.
看到的一個個案 所以神的同在就是平安.
所以我們要認定 我們在主耶穌的救恩裡面.
我們是會已經立即得到神的平安.
在我們每一個基督徒的內心.
所以如果某些環境.
例如那些環境真的需要更加憤怒的平安.
當然不是叫你晚上去學人家捉鬼.
走去那些荒山野嶺 那個就不包的.
那個是你自己在挑戰神的 那個就不算的.
你就不要去那些地方 反過來.
有時候你遇到一些很突然間的危險.
神的平安必須要首先來霸佔我們的心靈.
進駐我們的心靈 那個進駐.
等一下《肥立比書》四章就說到這裡了.
我們記不記得《肥立比書》.
幫大家引用了五段關於平安的神進駐.
四章九節《肥立比書》.
或者第八節開始.
弟兄們 我還有未盡的話.
凡事真實 可敬 公義 清潔 可愛 有美名.
又有什麼德行 又有什麼稱讚.
這些事你都要思念 哦 這個很重要.
原來要將我們的心思意念.
貼近神的真理 神的愛.
神的公義 神的同在裡面.
所以你不能夠有神的平安.
但同時你又想一些壞的東西是不行的.
所以這裡保羅讓我們知道.
原來我們在這樣的心靈裡面.

$^{281}$保守自己的思想意念.
在這些聖潔的 乾淨的東西上面.
就會有 這裡說有什麼德行 有什麼稱讚.
這些事你們都要思念.
都要常常去成為我們主流的思想.
時時我們有意識的時候.
我們就想一些健康的 愛人的東西.
祝福人的東西.
這裡就會好像我們鋪平了一條.
好像紅地毯的路那樣.
接著保羅就說.
你們在我身上所學習 所領受 所聽見 所看見.
就是說又是貼近使徒保羅和傳道者.
他們的生命 他們的服侍.
當我們來到去羨慕聖公的時候.
就是說你的思想對神的真理.
你有一個渴望 你有一個這樣的裝備.
就是說我們難免有時候會有些東西.
令到我們很生氣 令到我們很憤怒.
但是記住那個生氣和憤怒不要令到我們犯罪.
我們有這個意識之後.
我們馬上讓主的愛 讓主的十字架的救恩.
再一次充滿我們自己的心靈.
所以我們就能夠再一次有主的平安.
接著就是羨慕聖公.
就是保羅他們所做的我們渴望.
我們未必一定有那個崗位.
但是我們可以來到去認同 禱告.
或者在我們能夠有的資源底下.
比如金錢的奉獻等等 對這些聖公的支持.
就是說讓你的人生都放在那個聖公裡面.
你可以仍然是一個上市公司的CEO.
你可以是一個很忙碌的家庭主婦等等.
但是有份於那個聖公 那是很重要的.
接著這裡告訴我們 賜平安的神.
就必與你們同在.
這裡有另一個角度了 剛才就是說.
神就是平安 祂的同在就是平安.
但是這裡有一些容易達至神同在的那些準備.
原來準備我們的心靈 準備我們的思想.

$^{321}$準備我們在人生裡面投入在神的國度聖公裡面.
當我們是這樣來到去 有這種傾向.
有這種羨慕 有這種裝備的時候.
我們就很快 必定會有保羅所說的平安的神.
就必定與這樣的人同在.
這個大家就會明白了 如果我們讀聖經.
我們祈禱 我們去常常去幫助有需要的人.
我們常常我們的心思意念都為神的國度去考慮.
你就可以肯定一件事.
平安之君 必定會與這樣的人同在.
就是說 這裡有一個更大的一個的應許.
就是說 我們很多時候很想有神的平安.
但是很多時候我們放空了我們自己的心思.
我們想一些很熟悉的東西.
我們根本不羨慕神的聖公.
我們根本沒有興趣 但是我們很想有神的平安.
但是為什麼經常都不行呢?.
整天心都很亂的呢?.
整天都好像不行 看不到這樣的經歷呢?.
原來就是《肥臘彼書》四章九節的地缺.
平安的神必定同在.
但是我們自己才有一個這樣的準備.
就是說 你想一下 如果我們去吃宵夜.
你給我那些材料 那些炭又濕了.
然後那個環境就快去到黃梅天.
濕漉漉的 那個爐又濕漉漉的.
你想一下 剛剛下了一場小雨.
然後又沒有帶那些炭精 只有那些報紙.
也不是不行 但是辛苦很多.
也能透過那個爐的 不過透了很久.
但是如果你想一下 我們萬事都俱備.
快點吃一點就馬上吃.
所以為什麼我們這些姐妹.
這麼容易就會有神給的滿足.
有神給的那個平安呢?.
就是因為第九節.
所以你看一下這裡多麼重要.
平安的神必定與這樣的日同在.
這樣就是說 保羅 其實原來.
絕對不是神的平安藏起來不給某些日.

$^{361}$而是原來神很想每個人都得到平安.
但是原來很多時候 是我們的心已經滿了.
就是很簡單 剛剛吃完一個很好吃的晚餐.
我再給你其它東西吃 你也沒有興趣吃了.
因為你已經滿了.
這些道理我們大家都明白的.
但是真真正正落實去做的時候.
我們就有很多人就未必.
有時候有一種這樣的感受的原因就是.
因為我們沒有陽主的真理.
在我們的內心成為我們的力量.
成為我們的思考.
成為我們人生的一些很重要的價值觀.
所以就會變成一個相輔相成.
但是也是一個結構性的惡性循環.
所以有些人如果不羨慕 就變成惡性循環.
越是想得到神的平安.
大神的平安 好像越來越離開他很遠.
但是相反有些人他很容易就滿足.
很容易就會得到神的愛 感受到神的愛.
他可以很容易就會面對一些事情.
他都能夠很從容的處理.
因為正正就是四章八節和第九節的應許.
他信以為真 他很願意來到去.
這樣將自己放在神的裡面.
所以原來不難的.
原來只需要我們的心思意念 愛主所愛.
常常想念到主的事情.
我們就會必定經歷平安之時.
《貼索羅尼加前書》第五章.
所以這些金句第四章 第八到第九節.
這些要回溯它.
這些是保羅告訴我們 這是他的秘訣.
《貼索羅尼加前書》的第五章.
是講主在來的事情 第五章.
我們可以讀長一點的經文.
第十二節開始.
「弟兄們 我們勸你們敬重.
那在你們中間奴婦的人」.
或者可能是教會領袖.

$^{401}$或者可能要處理一些很棘手的.
可能是神學的爭議的問題.
所以這些奴婦的傳道者.
或者這些信徒領袖.
保羅就說要敬重這樣.
因為當時保羅就離開了《貼索羅尼加》.
就去到希臘那邊.
所以已經是分派給一些領袖.
當地的領袖來治理教會.
所以保羅就說.
沒錯 就是在主裡面治理你們.
勸戒你們的人.
由於他們所作的功.
用愛心格外尊重他們.
敬重和格外的尊重.
對於這些領袖.
他們將主的道來教導,勸戒.
你們也要彼此的和睦.
我們又勸弟兄們.
要警戒不守規矩的人.
勉勵灰心的人.
輔助軟弱的人.
不同的表現.
你就要用不同的心態.
不同的態度來對待不同的人.
所以不能夠用勉勵.
對那些不守規矩的人.
就是助紂為虐.
你就是讓教會成為一個賊窩.
你對於那些灰心的人.
就很重的篤責.
那個灰心就更加沉到底.
更加失去了那個志氣.
他就會更加去自暴自棄.
離開教會,離開神.
所以這個是要勉勵灰心的人.
勉勵就是說你要感受到他的軟弱的原因.
可能他不可以一下子就跳出他的軟弱.
那你要怎樣來用你的安慰勉勵他呢?.
這個勉勵是成為一個動力.

$^{441}$這是很高的人與人之間的藝術.
不可以一本通書就這樣對所有的偽眾.
就是一點都不簡單.
是很高的藝術來的.
要輔助軟弱的人.
他軟弱但是他很想起來.
你怎樣輔助他呢?.
輔助到一個地步.
他最終能夠自己可以獨立獨行的.
他自己能夠可以很穩固的靈性.
甚至乎可以在一個教會裡面.
慢慢甚至乎可以成為領袖.
之前也說過一個例子.
突然間又勾起了.
保羅班德醫生已經離開世界了.
這位很出名的在印度做麻風的.
一個很前衛的醫治的技術.
發言的一位出色的醫生.
他的父母都是在印度很出名的宣教士.
一生就是這樣奉獻自己的生命.
來在印度做這些窮苦的人的福音工作.
所以他的兒子後來已經在英國讀醫.
後來也有感動就要幫那些麻風病人.
輾轉保羅班德醫生就是幫助那些麻風病人.
其中一個令他很深刻.
就是有一個病人真的是因為那些手腳殘缺.
又很難看的面容.
因為那些麻風的病毒入侵.
侵蝕而導致他那個痛感神經失去了.
就不知道痛而很容易會弄傷等等.
後來這個病人信了主.
但是去到教會.
他很想去教會很想守聖餐.
但是在教會裡面的弟兄姊妹.
因為這個病人的出現.
導致他非常大的挫敗.
因為大家都怕他傳染 個個都彈開.
這個很合理的 你聞到這些臭味.
你看到一個人的手腳都畸形.
面容又扭曲 很恐怖的科學怪人.

$^{481}$任何一個人沒有經過詳盡的神愛訓練.
很容易第一個印象就是會害怕.
和會有一種這樣的歧視.
有時候都掩飾不了的歧視.
對這個治好了的麻風病人.
肯定是不會傳染的.
但是製造的那個內心的那種傷害.
比麻風病毒來得更加嚴重.
所以有一段時間都真的很自暴自棄.
都真的懷疑這個信仰的真實.
因為他看不到神的愛.
藉著弟兄姊妹對他的接納.
所以原來一些軟弱的人.
你要怎樣來扶持呢?.
其實是真的很不簡單 很不容易.
相信大家都同意.
感受堂,嘉欣堂,孫德堂和紅欣堂.
都有不同的setting 都有不同的階層.
你想一下紅欣堂那個階層.
有很多基層的 有些是低學歷的.
你想一下在聚會裡面.
有時候那個形態和感受都不相同.
所以如果我們能夠看到這些事情.
弟兄姊妹純粹都有一個很敏銳.
很對於弟兄姊妹的背景.
就是家庭背景,出身.
來到有一個敏銳.
其實就會是一個很大的功課.
因為不可以一下子輕易.
就說我們每個人都一樣一線紅人.
其實不能夠一線紅人的.
因為我們人始終有我們自己過去的背景.
階層的包袱.
本來這些本來都不是太對錯的.
但是當有不同處境.
不同階層的人來到的時候.
就出現了這一種落差帶來的歧視和傷害.
所以要輔助軟弱的人.
真的不是一句就講完就可以了.
因為真是很實際的.

$^{521}$我經常都說笑.
我妹妹早期就是傳宣去天守圍直堂.
接著他們就響應去天守圍直堂.
去到天守圍之後.
坐了沒多久的時間.
她就跟她先生講.
就是跟我的妹夫講.
不如我們回到傳宣.
不要在天守圍了.
她先生就問她什麼事.
她說隔壁有個女士在崇拜的時候撩鼻屎.
她說很噁心.
我妹妹是做秘書的.
在大銀行裡面.
她先生就說我們在基層.
人家的情況是很正常.
當然不是撩鼻屎正常.
人家習慣了這樣.
就是要慢慢一步步改變.
當然好像開玩笑一樣.
慢慢她都明白.
一直都在一段比較長的時間.
都在天守圍那裡去直堂.
去做堂位發展.
所以要輔助軟弱的人.
其實不是容易.
也要向眾人忍耐.
有些人要一直忍耐他.
有些人一直都不長進.
一直那些心態.
一直那些價值觀.
有時候我們都不敢問那些老人家.
過年不准拜偶像.
可能他聽完之後.
他都沒什麼遺憾.
可能初一,十五.
大年初一可能他走去廟裡.
很喜歡拍照.
看不到他在黃大仙爭頭轉向.
有時候他們未必一下子改變.

$^{561}$作為領袖就要向眾人忍耐.
接著下面第十五節.
你們要謹慎.
無論是誰都不可以以惡報惡.
或者彼此相待.
就是說你對我不好.
我對你不好.
或者待眾人常要追求良善.
所以這裡是說內心的愛.
自己不要害人不得已.
還要包容接納其他人.
要等候.
千萬不要用惡去還惡.
反而要用善去對待不同的人.
要常常的喜樂.
不住的禱告.
這個是內心的裝備.
凡事的謝恩.
因為這是神在基督耶穌裡.
向你們所定的旨意.
旨意就是說.
我們不是尋求一些特殊的.
要玲瓏浮特的.
要一祈禱他馬上要答我的.
不是的.
保羅說神的旨意.
其實就是我們這個人的生命.
要常常保持一份的喜樂.
保持一份在神裡面禱告的生命.
常常浸在神裡面的人.
就自然去遵行神的旨意.
接著19節.
他說不要消滅聖靈的感動.
不要藐視先知的講論.
但要凡事察驗.
善美的要持守.
各樣的惡事要禁戒不作.
就是說這裡和聖靈和神的道有關.
就是先知的講道預言有關.
就是說是很平衡的.

$^{601}$你要讓聖靈的充滿.
聖靈的感動.
但是又要讓神的真理平衡我們的心靈.
就是說我們的聖靈感動.
是有真理來做一個平衡的.
就是使我們既有真理.
同時也有感動.
所以不是一套的知識.
我都懂了 不用再說了.
其實那些只是一個知識.
未真正去實踐 未真正去改變我們的生命.
仍然是一套死的知識.
但相反我們有聖靈的感動.
有真理的知識來平衡我們內心的感動.
就來一個很平衡的人生的基督徒.
所以我們就能夠知道怎樣分辨善惡.
接著第23節.
「愈似平安的神 親自使你們全然成聖.
又愈你們的靈與魂與生子得無保守.
在我主耶穌基督降臨的時候.
完全無可指責.
乃照你們的本事信實的.
他必成就者事」.
這裡就再豐富一點.
剛才說了我們要內心.
新生命有主的真理 有主的真實.
一些可羨慕的聖功 我們會去羨慕.
我們將自己放在那個氛圍裡面.
我們自然就會變成神平安的同在.
在我們的內心 在我們的心懷念念.
這裡又再進一步了.
是我們整個的基督徒人生.
就是要追求那個生命的成長,成熟.
就是要讓我們能夠在我們的待人接物.
因為這裡說了好幾樣東西.
同人的關係 就是說有時候會有嫌隙的.
但是千萬不要以惡報惡.
反而就是盡量尋求和睦.
盡量以良事來追求 成為我們的追求.
接著就提到我們要怎樣在神的旨意裡面有份.

$^{641}$就是要不住的禱告 凡常常的謝恩.
常常的喜樂 這樣來成為我們生命的一個.
很重要的每天生活的元素.
接著也都讓我們能夠保持常常被聖靈.
和真理去驅動的一個人.
所以基督徒如果我們都沒有意識.
一離開教會 根本都不知道自己是基督徒.
根本整個星期沒有任何的蛛絲馬跡.
你是和基督有關的 心思固然已經不想了.
實際的行動都沒有了.
不知道有多少弟兄姊妹在公眾面前.
不祈禱 不敢祈禱.
不告訴別人自己是基督徒.
好像那個完全沒有基督的生活的一個星期.
到星期天準備回教會了.
就開始起床了 有意識了.
要回教會早課了 或者午課了 諸如此類.
保羅說不應該是這樣.
應該是自從我們認識基督之後.
就應該讓我們的生活有份於基督.
有份於這個神的國度 那個的善惡 生活.
還有在我們的價值觀裡面.
價值觀就是說我們的心思,意念,最深處想的東西.
我們一路人生的方向.
這一切都讓基督的真理,聖靈在我們內心.
不斷的滋潤著我們.
這種滋潤就成為了我們基督徒的生命的豐富.
那個時候保羅又是一個很大的英雄.
平安之君,平安的神 他一定會親自使我們.
能夠在以上所追求的人生.
來一路漸次的成聖.
就是說一路向那個聖潔來成長.
成為真真正正的離世之前的聖徒.
一路追求我們的良善,我們的公義.
我們內心的聖潔,平安,對別人的愛等等這一切.
是一個真是不斷有成長的一個這樣的生命.
次平安的神必定會成就.
親自的使我們傳言的成聖.
所以這個傳言的成聖.
就是原來我們基督徒的生活.

$^{681}$就是能夠有真理,有愛,有公義,有和平.
所以其實這些經文保羅告訴我們.
平安的神進駐不是你坐在那裡等人到的.
而是我們是很積極的.
是我們的心思想到福音改變人生命的那個雀躍.
就使我們這個追求不是變成一種很空洞的.
或者根本沒有內容,沒有行動的.
但是一直在說的肥勒比書.
在說的《貼索羅尼加前書》是有行動的成長.
而這個行動就是內心裡面.
因為平安的神進駐和來驅使我們.
剛才我獨漏了一點.
四章肥勒比書有一個很重要的關鍵字.
四章這裡就是第七節.
我們剛才獨漏了這一節.
告訴我們神所賜出人意外的平安.
又再一次告訴我們平安的神.
所以出人意外的.
就是說人生一定有不同的意外.
但是我們基督徒如果神的平安在我們裡面.
就算多出人意外.
我們也不會覺得會有一種懼怕,驚恐和失控的內心.
不要小看一些事情.
有時候你聽到晴天霹靂的噩耗.
有時候那個心靈是很需要立即有神的平安的.
人生有時候總有一兩次的.
有時候有些人很慘的.
比如昨天的大嶼山.
一輛的士一個二十多歲的男士.
坐在的士上沒有繫上安全帶.
那個的士直接撞上了工程車.
那個的士司機和另外一個司機只是輕傷.
因為繫上了安全帶.
那個乘客是立即整個人拋在車頭上撞得很嚴重.
內臟就震得爛爛的 這樣離世的.
如果你在家裡聽到這個消息是立即崩潰的.
所以人生總有這些令到我們意外的.
但是要記住保留這些經文就是神賜的平安.
是超過了我們人能夠承擔的.
所以是出人意外的平安.

$^{721}$留意 必在基督耶穌裡.
保守心懷意念.
保守這個字就是那個重要的重點.
保守這個字就是菲律賓的駐軍.
現在我們香港駐軍在哪裡呢?.
石岡最多 其它地方都有的.
主要是石岡或者香港那邊.
駐軍 就是有駐軍就是說你安全.
有安全感 有平安 你不需要害怕.
小朋友 媽媽不在 你哭到死.
別哭 別哭 然後媽媽一回來.
那個平安立即進駐 這些我們都明白的.
在我們成長裡面 我們會知道的.
有什麼令到你突然之間立刻就不再哭的原因.
我們上個星期就去駱橋樓.
來派那些揮春.
那些小朋友很開心 跟著爸爸媽媽來.
然後他們就很醒目的.
完了之後就 曾律師請吃雪糕 去麥當勞了.
然後其中一個小朋友都很熟的.
因為他吃得太多了 我就將他半杯給了另一個小朋友.
他就哭得很厲害 哭到馬上衝回媽媽那裡.
哭到收不到聲音 很簡單的.
馬上就去買多一杯給他.
馬上就不哭了 馬上就笑了.
好像開玩笑一樣 但是小朋友就很明白.
那個進駐 那個駐軍.
那件事令到他突然之間完全感覺到那個平安.
保羅說 平安的神的進駐.
就會馬上保守我們的心懷意念.
心懷意念 令到我們整個人的心靈都崩潰的一些噩耗.
當然我們不希望他聽到 但是人生總有一兩個的.
但是我們要記住 我們是屬神的.
我們要讓主的平安在我們的心懷意念裡面去進駐.
是駐軍那個字.
那個神的那個字 就是保守這個字.
所以中文的翻譯都是比較溫和的.
其實應該說 基督耶穌的同在.
就是成為了那個駐軍.
那就真的告訴我們 這一大段經文都是想說.

$^{761}$平安之君 這一位的神 祂的平安.
因為祂就是平安 所以祂的同在.
就代表我們有平安 所以我們需要平安的.
我們真的需要做物主時保守我們心靈的平安.
讓我們能夠真的有這一份.
這樣的常常都不會失去的平安.
最後就是《貼撒羅尼加號書》第三章第16節.
看看保羅他在寫完前書之後 他又有什麼勸導呢?.
後書第三章 我們就從第六節開始看起.
弟兄們 我們奉主耶穌基督的名.
吩咐你們 凡有弟兄不按規矩而行.
不遵守從我們所受的教訓 就當遠離他.
就是說持續都不聽 不怕警告的這樣的持續.
就不是一次就馬上把他隔走 那就不行了.
一定是持續性的不遵守的.
你們自己原知道 應當怎樣效法我們.
因為我們在你們中間 未嘗不按規矩而行.
也未嘗不踏輕人的範 倒是辛苦勞碌.
晝夜作弓 免得叫你們一人受累.
這並不是因我們會有權柄 就是使徒的權柄.
來傳道 有教會弟兄姊妹的供養等等.
但是保羅就說 但是這並不是因我們會有權柄.
乃是要給你們作榜樣 叫你們效法我們.
就是說保羅刻意這樣做.
是要讓這些信主不久的弟兄姊妹.
不要有一個誤解 不要有一個錯誤的.
對保羅的觀感 對保羅的觀感而帶來打折扣的.
就是對福音的那個認受.
所以就不希望被這些物質的東西影響了.
保羅和弟兄姊妹之間的那個關係.
我們在你們那裡的時候 曾吩咐你們說.
若有人不肯作功 就不可吃飯.
因我們聽說 在你們中間有人不按規矩而行.
什麼功都不作 反倒專管行事.
因為我們太有空 太有空就專管行事.
有空的人就管行事.
其實不是他管的 他也有空就管.
原因為什麼有空呢?.
就是因為他有一個錯誤的利用神學.
就是主在來很快 不要做事了.

$^{801}$要等候主在來.
但問題是他的等候就是成為了一個行人.
講行話的一個人.
就變成保羅都來到去告誡教會要留意.
保羅說的主在來是很近.
但不代表主在來就要大家懶惰做事.
所以好像很矛盾 既然主在來就什麼都不要搞.
但原來不是 主在來是一個事實.
必然會發生 但是我們每天的生活.
就要活在主在來這個永恆裡面.
我們去盡忠我們每天的生活.
這樣才叫等候主在來.
所以等候主在來而不上班,不工作.
保羅就要譴責,責備.
所以原來我們有一個一知半解的神學的觀念.
都很容易會導致一些的謬誤.
所以是一個完整的對神學的全面的掌握才行.
因為他們可以說保羅說的.
主在來很快回來 好像賊一樣.
其實第二章已經說了.
我們就要準備等候他回來了.
不可以再等他了.
不可以再來到他沒有時間的各方面.
但是保羅就說.
反過來好像現在突然間將那個天堂的盼望.
拉回去很現實的生活.
不是啊 等候主在來跟你殷勤的,努力的工作.
是沒有抵觸的.
原來這個平衡的神學也是我們需要學習的.
我們不要過度了某一個信念.
又變成一種自己綁著自己的枷鎖.
第十節 我們在你們那裡的時候曾吩咐你們說.
若有人不肯作供就不可吃飯.
因我們聽說在你們中間有人不按規矩而行.
什麼供都不作 反倒專管行事.
我們靠主耶穌基督吩咐勸戒這樣的人.
都仍然是勸戒的.
要安靜作供吃自己的飯.
弟兄們 你們行事不可喪志.
若有人不聽從我們這信上的話.

$^{841}$要記行他 不和他交往.
叫他自覺羞愧 但不要以他為仇人.
要勸他如弟兄.
就是說有些人有一點誤解了真理.
而行差踏錯.
而有些錯誤的,怪異的信念和行為.
所以要勸戒,要留心.
慢慢引導他離開那個錯誤.
所以不要以他為仇人.
要勸他如弟兄.
就是說他肯改.
慢慢明白了而慢慢改變.
這樣就可以收納了.
但如果那些是錯誤的真理.
怎麼都不肯改變.
怎麼都來傳這些假的錯誤的道理.
甚至遲做一些在教會講閒話的人.
不肯悔改.
那就不是來當他是弟兄的了.
要趕出教會的了.
不過留意第十六節.
平安的主隨時隨時親自給你們平安.
願主常與你們眾人同在.
這裡又多了一些小小的特別的形容.
隨時隨在.
不同的時間,不同的事件.
主的平安可能是大小.
可能是那個的色切性來留在我們的內心.
所以人生真是需要不斷的成長.
我們要經得起風浪.
姜濤應該要經得起風浪.
不是叫你投機.
不是叫你炒金,炒比特幣.
而是我們說的是我們人生要一直成長.
不要一個玻璃心.
不要什麼都沒有平安.
小小事就沒有平安.
又不要膽大妄為.
又不怕神又不對.
而是我們一直隨著我們不同的人生.

$^{881}$所以隨時隨在不同的場合.
不同的一些環境.
你都能夠有神為你而設的那份平安.
這樣來進駐.
仍然都是五段經文.
都是告訴我們.
保羅他很深的體會.
平安的神在他的傳道.
保羅有什麼不同的場合呢?.
在江河的危險,盜賊的危險.
賈弟兄的危險,船壞的危險.
被那些拉比,法利賽人來刑罰40鞭.
減少一鞭.
又去到這個阿里達王的手下.
這個戴瑪士加.
他要連夜被人用那個鑼.
將他扔到城牆外面.
他要離開逃亡 上夜殺難.
很多這些危險.
就是說保羅深深的經歷.
隨時隨不同的時間.
不同的事件.
他都有主給他的平安.
可能有不同的性質.
可能有不同的程度.
原來這麼豐富.
原來這個祝福,這個祝願.
原來可以這麼豐富,這麼宏偉.
原來可以在我們的生命裡面.
可以任何的事.
我們都來到好像有一份很深的經驗.
所以也是這麼說.
教會如果有資深的愛教會.
愛弟兄姊妹的領袖,牧者.
就是那個教會最大的資產和福氣.
真的 他一直在這個教會.
為教會的弟兄姊妹守望,禱告.
忠心為神 為神的道去傳揚.
這樣成為一個好像定海神針.
當然誇張了.

$^{921}$真的好像一個很幼的那份平安.
其實說真的.
你回去見到你自己所愛的敬重的牧者.
其實有很多的問題.
有一半可能已經解決了.
其實不是因為他解決你.
你可能講都還沒有講.
你已經感受到這一位老人家.
這一位的牧者.
他的平安仍然是日於言表的.
這些其實就是每一個教會都需要的.
這樣的年長的老人家.
所以不要小看祖堂.
很多定海神針在這裡.
當然他回來睡覺是另一個問題.
就是老人家累.
但是不要小看教會有的老中青的組合聚會.
因為老人家最喜歡就是見到小朋友.
所以他如果能夠那個心靈很有神的一面的老人家.
他看到的他就會為不同的小朋友,年輕人去禱告.
所以如果年輕人只是說我們不需要這些老人家.
我們只是自己崇拜.
其實他少了很多福氣.
年輕人看年輕人.
就是你玩我 我比較 大家給錢比較.
但是老人家看年輕人是不同的角度的.
是不同的心境的.
所以希望我們在曾經班的弟兄姊妹.
應該都算是年紀大的.
我不知道外國那些聽到的是年輕還是年長.
我們的優勢就是用這種神給我們的眼光.
去看年輕人,看小朋友.
我們要將神的平安進駐在他們的心靈裡面.
就是要為他們祈禱.
希望神將來大用我們下一代,下兩代的年輕人.
他們能夠起來接觸神的福音的棒.
所以這五大聖經是保羅說到.
他可以將他所經歷到的做物主的平安.
散播給不同的外邦的教會,不同的人等等.
但是他自己就是很深的經歷這一份的平安.

$^{961}$所以他在《船難》《少林寺傳》的第27章.
在那個《船難》裡面.
各人都在那裡呼天搶地.
保羅就叫那些千夫長將那些物件全都扔掉.
讓那艘船不要那麼負荷過重.
但問題是保羅說.
始終我們都會撞過去擱淺的.
但是大家不用怕.
因為昨晚我的神跟我說.
那艘船一定會壞的.
一定會沉的.
但是不用擔心.
所有船上的人 如果我沒有記錯.
是270人在船上的全部都會得救.
我先看看《少林寺傳》.
我看看有沒有記錯.
因為太久沒有讀《少林寺傳》了.
第27章的《少林寺傳》.
那個船壞了.
沒錯 這個暴風.
少了一個 是276人.
保羅 我讀一讀給大家聽.
第33節.
這個27章.
《少林寺傳》.
天漸亮的時候.
保羅勸眾人都吃飯.
說你們願望忍餓不吃什麼.
已經第十四天了.
在船上 巨大風浪.
大家走過船沒有?.
去過海洋沒有?.
去過進入太平洋沒有?.
我姐姐的兒子就走過船.
不過他現在已經上岸了.
因為男性很多時候會有一種探險的心情.
不是自己去當然不怕.
問他去過那些大洋沒有?.
他說去過.
他說那些浪打上來.

$^{1001}$基本上有嘔吐 基本上是很恐怖的.
那些真的是滔天巨浪.
所以他告訴我們已經第十四天了.
所以我勸你們吃飯.
這事關乎你們救命的事.
因為你們各人連一根頭髮也不至於損壞.
保羅說了這話就拿著餅.
在眾人面前祝謝.
祝謝神就擦開吃.
於是他們都放下心也就吃了.
所以沒得辦事了.
因為你十四天在船上在大海汪洋在颶風上.
保羅自己作為那個領袖.
他真的有神的平安盡處 很重要.
所以第三十七字.
我們在船上的共有二百七十六人.
他們吃飽了就把船上的物質拋在海裡.
說要叫船輕一點.
所以保羅真是口口聲聲.
他真的經驗到神的平安.
在任何的情況下.
所以隨時隨時不同的時間.
不同的事件 不同的處境.
保羅都學習這些平安.
所以保羅自己不斷.
我相信他不斷將自己的生命.
放在舊約他對神的律法的尊重.
然後他自己將舊約的經文轉化成為新約.
保羅寫書信的這些神學的教義.
所以保羅就是能夠這樣融會貫通.
他自己整個的內心是被主的愛 則愛.
永恆的盼望 這樣來充滿他的內心.
所以他對外人 對猶太人.
都是有一份很大的愛.
所以頂小妹我們信耶穌.
得到永生去天堂固然一定是最大的福氣.
但是原來我們的人生現在有愛 有真理.
所以我都很難想像.
你想一下我十幾二十歲的時候.
其實從七八歲開始已經開始在家裡.

$^{1041}$在木屋區已經粗口囊囊上口了.
但是現在整個心思意念成為基督徒之後.
是根本沒有興趣講粗口的.
聽到這個都覺得噁心.
完全不會有那些粗口的東西在心裡面.
你就可以想像 舊恩已改變我這個人.
一個市井之徒 木屋區長大的流氓.
可以變成一個現在能夠有真理 有愛.
對別人有寬恕 有包容的一個人.
其實都是這個人 都是這樣成長的.
但是如果沒有福音曾經在我的內心.
耶穌一直在我的內心.
我可能比市井更加市井.
因為社會時時都誇耀這些文化.
這些像文化 但是原來不需要.
煙仔由我信主不吸煙之後.
永遠不再吸了 完全不需要再受那個影響.
那個核就是平安的主在每個人的內心.
所以原來這五段經文.
真的可以成就我們的生命.
來到成長和聖潔的.
所以我不應該這樣做比較.
上完《查經班》 不是因為我.
是因為《十日道》 我們大家一起去認真去閱讀.
對大家那個心靈的練成是有幫助的.
你不信你再重新去反思一下自己.
看完《羅馬書》 是不是真的有一個生命的更新.
那個更新真的裡面是有了聖潔.
我們明白了真理.
明白了在耶穌裡面就是白白的救恩.
沒有東西可以使我們和基督的愛隔絕.
這些這麼斬釘截鐵的真理.
其實是我們內心裡面很重要的價值觀.
所以如果你只是一知半解.
勉勉強強上完初信班 洗完禮.
當然我沒有意思說那些經文課堂沒有用.
不是這個意思.
而是你接著停止不再追求.
翻崇拜都偶然地翻.
你的成長 你的生命肯定沒有翻《查經班》.

$^{1081}$浸淫在《十日道》裡面的那種豐富.
這些是肯定的 這些不是做什麼比較.
一定是你越對主的道越有認識.
越渴慕 你的生命不知不覺就會念寂.
所以今天我們最後都要講這五段經文.
就是平安神的同在 因為祂就是平安.
所以祂的同在就是平安.
原文的意思就是這樣.
但原來保羅不是告訴我們.
所以為什麼中國人這麼重.
出入平安是沒有的.
這條鐵圓那裡一樣是艾薩的 一樣經常有炎在家裡.
所以不是那張紙 是神 平安的神在我們的家裡.
所以重點不是指那個牌.
而是我們具體的生命裡面的進行神的旨意.
稍微喝口水.
我們正式進入第十六章.
這是上一次我們講完最後那一節聖經.
到底在我們的信仰裡面.
教義怎樣影響著生活呢?.
保羅在這個文安裡面.
就真的告訴我們 他很想具體的來跟我們講.
這個恩順清義是怎樣能夠很形象化 具體的.
這個Incarnation of Grace 是否叫「退車」?.
可能是的 叫「退車」.
教練呢?.
(教練和教師互動中).
我們講一下羅馬書 看你喜歡的.
(教練和教師互動中).
我們來到第十六章.
首先就是有差不多我數過26個保羅文安的人名.
這個裡面很有考究的.
但是我們來到最後這一章.
首先就要問為什麼保羅要刻意的去問這些人.
這些人來自不同的背景.
我們可以來做一個簡單的一個小小的歸納.
就是神的恩典.
怎樣具體在人與人之間的倫理關係上面.
來去具體的呈現呢?.
所以這個Incarnation of Grace.

$^{1121}$就是好像道成肉身的恩典.
怎樣能夠可以很具體的在人與人之間的相處上面.
例如打破那個階級的分離.
男女當時是有很大的社會階層的落差.
女性在羅馬的世界只不過是男性的附屬品.
所以女性是不能夠出頭的.
不能夠作為領袖的.
你可以想像得到保羅如果這樣的文案被當時的教會的人看到的時候.
是很震撼的一件事.
所以和大家做了小小的組織.
這些是我們以前全部學過的.
恩順清義就是怎樣解決罪.
怎樣能夠將神的恩典白白淋到每一個信的人.
這個就是恩順清義的最基本,最重要的意思.
還有信心和神的十誡律法.
律法是神設立的.
怎樣去在信心裡面和律法的到底是衝突.
還是一個圓滿的來滿足的完成呢?.
這些是我們以往在第五,第六章都講過的經文明.
另外更加重要保羅給我們的概念就是.
在基督耶穌裡面人與人之間過去的那個分離已經打破了.
在基督裡面是合一的.
所以以色列和外邦人再沒有任何的分別了.
因為基督就已經成為了那兩下合一的橋樑.
另外當然保羅也給我們的概念就是.
原來我們有聖靈的愛在我們每一個基督徒的內心.
所以我們在聖靈的愛裡面.
我們大家同有聖靈,同有信心.
所以我們就變成一個新的群體.
好像一個新的創造New Creation這樣的一個群體.
就不是再個別的個人去信耶穌了.
而是每一個人都在救恩裡面成為一個合一的群體.
所以這些是很具體的神學.
但是怎樣將這些具體的神學.
保羅能夠在他與人之間的相處上.
所以你可以說保羅真的有一個很高很強的技巧.
就是透過人的生命的改變.
來去展現剛才我們所列舉的.
小小的羅馬書所講過的這些神學的重要的真理.
所以你想一下 這個福音 十字架的福音.

$^{1161}$是會使每一個人成為了新造的人.
一個新的New Creation.
使一個人可以透過凝訊基督.
進入一個新的群體.
就是稱為神的教會 就是基督的身體.
你可以想像福音令到我們可以在教會裡面.
維護這個純潔和這個合一.
這個就是保羅很想展現.
所以可以這樣說 將十字架的恩典.
活在每一個名字背後的生命裡面.
每一個名字的背後就是一個階層.
一個社會的地位.
一個在當時的時代帶來的被強迫.
你只能夠活在這樣的階層.
例如你出生在羅馬的時代 你是一個奴隸.
你的子弟 你的兒女都是奴隸.
你的身份就是一個奴隸的身份.
所以如果你是外邦人.
你就不能像猶太人那樣在聖殿裡面.
男的女的僱員是貧窮的有錢的僱員.
根本就是已經用那個物質來分了那個階層.
所以你可以想像得到.
保羅透過這些來展現一個很大的突破.
就是保羅對每一個人的名字背後的稱呼.
這樣來帶出他對神學的真理的運用.
具體的呈現在每一個生命裡面.
每一個生命的改變 名字就不能改變.
但是名字背後的那個人 因為信了主.
保羅就提出他們那個生命的改變.
我們開始先看第一第二節.
今天可能只能夠去到這兩節聖經.
我對你們舉薦 我們的姐妹菲比.
她是堅革理教會中的女執事.
十六章的第一第二節.
請你們為主接待她 合乎聖徒的體統.
她在何事上要你們幫助.
留意一下這個幫助這個字.
這個字在當時是一個很重要的制度.
她在何事上要你們幫助 你們就幫助她.
因她素來幫助許多人 也幫助了我.

$^{1201}$出現這個幫助很多次.
到底這個意思是怎樣呢?.
首先看一看 我們很肯定羅馬書這封信.
保羅寫完之後 就透過這個菲比.
來帶到羅馬 第一節已經講得很清楚了.
所以這個菲比 她就是一個女執事.
原來初期教會這個堅革理 我看這個地圖.
哥林多就在這裡 哥林多東面的地方.
就是近一點的那個附近.
那裡就是堅革理一個不是太大的城市.
都是沿岸 地中海的沿岸 堅革理.
哥林多最窄的那個地方就是那個運河.
就是哥林多右上角那裡 最窄的那個就是運河的地方.
通過對面一個對地中海的.
所以那個運河是很重要的.
堅革理的女執事.
就是說保羅肯定她在教會的職份,身份.
她已經在教會裡面作領導了.
所以原來初期教會都有些保羅的時代.
都有些弟兄姊妹 外邦人來的.
都有些弟兄姊妹是這麼長進.
到底她是在金錢上面是很有錢.
而作為那個領袖 例如在菲律賓.
有一個賣紫色布匹的那個女底牙.
那個未必是很多的真理.
但是在她家裡有教會的 這個女底牙.
就是賣紫色布匹 就是說一定是相當貴重的布料.
那裡有提到女底牙.
這裡就提到這個菲比 就是這個教會的一個執事.
就是說一個屬靈的領袖.
她對真理應該都相當有認識的.
留意一下 她帶這封書信去.
因為當時真是唯有寫好那封信.
因為當時沒有印刷術 那封信沿著.
就由那個運輸的人將那封書信帶.
所謂書信是那些蒲草紙 不是今天的紙張.
就運到當地的教會 就給當地的教會的領袖.
當他們聚會的時候.
就會在教會的聚會來宣讀那些保羅的信.
所以她自己都有一定的對羅馬書的首先的了解.

$^{1241}$所以她得到的福氣真是不少.
至少羅馬書是用希臘文寫的.
她又懂得希臘文 這個菲比.
她又要理解整本書.
哇 十六章我們講了這麼多堂.
現在回去問之前的那些 可能大家都不記得了.
但是菲比能夠至少掌握保羅大概那個神學的想法.
因為保羅都會告訴她.
我寫這封信背後的目的等等.
所以在一路朗讀的過程.
菲比有可能有責任要加以解釋的.
另外她就代表著保羅來帶給羅馬教會.
有這樣的授權.
就是說保羅授權菲比帶這封信.
就是說她自己本身在堅吉里.
在保羅的那個傳道的童宮裡面.
她都有一定的權威.
也是保羅信任的.
所以在這樣的推薦信.
這封是一個等於推薦信的.
保羅推薦菲比去到羅馬教會.
就是意味著羅馬的教會要接待.
由堅吉里這個哥林多地方來的女執事.
就是說兩間教會.
就會因為保羅的這封信有一個相交.
有一個合一.
有一個彼此來分享神的工作.
比如堅吉里那裡的工作是怎樣呢?.
哥林多的地方.
現在羅馬 菲比去到羅馬.
她是不是有一些很重要的公務.
要在那裡來完成呢?.
在那裡要來做呢?.
這些都是在《哲理行間》的第一節裡面.
已經有一種日漁言表的來講到.
保羅對於兩處教會的合一的寒蟲.
彼此之間有甚麼宿命的分享呢?.
當然我們不是太知道.
兩間教會的規模.
到底是堅吉里大一些.

$^{1281}$還是羅馬教會大一些呢?.
雖然它不是很遠.
都是寄中海很短的路程.
但是對於當時不是那麼科技發達的傳遞的過程.
始終都需要人用人力.
用腳用坐船.
將那些信息帶到對方的教會.
但是這個充分的神權信任.
就是透過菲比被派去.
由保羅推薦他去.
就讓我們看到保羅對於男女的身份.
是非常之平衡的.
他不會對女性有一個歧視,貶低.
甚至乎來侮辱.
因為古代男性對女性真的沒有那個尊重.
都是會看為一個物件.
就是將女性的邪惡化了.
現代一樣有的.
不過不同的程度,不同的形式.
或者不同的文化裡面.
所以 插一句話.
70年代,80年代.
都是近代宣導會才開始按立女牧師的.
以前都沒有的.
所以 將女牧師都是很早期的.
可能前五個按牧的.
現在當然平常了.
現在當然就是沒有什麼問題了.
播道會都是很後期才按立女牧師的.
你可以想像到.
有時這些的情況.
我們華人教會都是會滯後了一些的.
就是都要盲的.
但是可想而知.
原來男女那個情況.
在聖經的世界和實際現實.
有時都還有些落差的.
何況你不用想了.
不用說就是奴隸的制度.
就算今天沒有奴隸制度的美國.

$^{1321}$一樣對黑人的歧視.
有時Bad Life Matters.
你可以想像到.
要不斷爭取社會的同情,關注.
那些白人的警察見到黑人.
就憤外驚 要壓低他.
有時就會衍生這些槍殺案.
你可以想像到.
連聖經的世界都已經很清楚告訴我們.
我們人類都未必能夠一下子這樣改變.
已經是講21世紀了.
都還有這些情況出現.
所以可想而知.
保羅這一篇長的文案.
待會他會提到的人.
就真的讓我們看到救恩是很實實在在.
在每一個生命的名字的故事背後.
帶來的那個改變.
所以如果菲比能夠這樣自由的去落馬.
相信他一定具備自己足夠的社會的地位.
和那個的資源.
因為當時那個社會不是每個人都可以隨便去一個遠遊.
不是每個人都可以隨便去的.
所以如果能夠有這樣的身份.
由保羅又對他的信任.
而寫這一封的推薦信.
所以背後的含義都相當大.
所以盼望我們都能夠總是對人有那個信任.
對人有一個包容.
對一些如果是初出茅廬的.
能夠給他有一個授權,肯定.
是會帶來那個可能年輕人.
他的生命是會更加趨向快一點的成熟.
人如果什麼都怕.
什麼都不敢給他做.
什麼都不給他做.
不給他碰.
那麼他的成長就會相當緩慢.
慢慢去到一個地步就是什麼都不敢做.
什麼都怕錯.

$^{1361}$什麼都怕老人家不喜歡.
那就變成扼殺了很多這樣的授權帶來的生命成長.
所以作為領袖.
我想我們都要學習這些功課.
因為永遠自己做就一定不會錯的.
錯了都沒有人知道 因為自己做.
但是當你授權.
當你賦權給人去做.
讓他去成長.
其實是一個很大的一個傷盈.
對一個生命的建立都有一個很大的一個祝福.
我們再看多一點.
留意第二節.
請你們為主接待他.
合乎聖徒的體統.
什麼叫合乎聖徒的體統呢?.
就是說去飲茶就不要去大家樂.
去會所.
合乎聖徒的體統.
是什麼意思呢?.
就是說吃好一點.
人家這麼遠來到.
你有沒有理由叫人家去吃大家樂呢?.
李斯勁牧師來 開玩笑的.
講完道 對不起 今天我們招呼你.
我們大家樂 好不好?.
一起排隊.
一般我們都很少會這樣.
有時候覺得這樣怠慢了一些神的僕人.
不是因為階級.
而是那個聖徒的體統.
保羅是不是說這個呢?.
是不是要給他吃餐好一點呢?.
給他去會所吃好一點呢?.
相信不是.
他說他在何事上要你們幫助.
你們就幫助他.
特別提到.
似乎他真的有很重要的公務要辦理.
所以需要當地的羅馬教會.

$^{1401}$給他們一個推薦.
可能是其他的授權.
留意一下.
因他素來幫助許多人.
也幫助了我.
今天就不夠時間詳細講這個.
這個叫做私恩的制度.
那個Patreon那個字.
那個制度.
我們今天不夠時間講了.
只能夠今天先停在這裡.
下個星期就是過年.
所以我們會先停一個星期的.
我們一起重新禱告.
多謝天父.
我們能夠在今晚.
處理了很多段經文關於平安的實.
多謝你.
因為你的同在.
因為我們在主耶穌的救恩裡面.
就立即獲得神的同在.
在我們的生命裡面.
是藉著聖靈住在我們的內心.
多謝你.
我們能夠經驗到你的平安.
很期望每一個基督徒.
都能夠在每天的生活.
在不同的處境.
因為講到明.
在隨時間.
隨事件.
我們都可以有神的平安.
完全幫助我們.
也都讓我們在真理的道上面.
我們能夠不斷地有長進.
我們能夠愛務主的真理.
愛務主的聖公.
這樣來服侍.
再一次我們獻上感恩.
我們下個星期.

$^{1441}$就是下一次.
我們會看聖經的文安的日.
讓我們能夠逐段去研讀.
非比都是一個很不尋常的雷則事.
能夠在保羅的欣賞,授權.
和推薦之下.
去到羅馬.
相信他有非常個人之處.
都是我們下一次開始的時候.
都會再重點地去研讀.
求主幫助.
再一次多謝主.
讓我們今天晚上.
安靜地來聽研讀神的話語.
我們這樣感恩禱告.
奉耶穌基督的名.
阿們.
\newpage



\section{羅馬書}
\label{sec:4pLVUVJIfgU}
\textbf{2026-03-06 周五晚查經班《羅馬書(三十八):分題:使徒保羅的深情問安》 講員:曾裔貴牧師}
\newline
\newline
連結: \href{https://youtube.com/watch?v=4pLVUVJIfgU}{\texttt{https://youtube.com/watch?v=4pLVUVJIfgU}} ~~~~ 語音日期: 2026-03-09
\newline
\newline
\hyperref[sec:3rvmIgiASuQ]{\small{< < < PREV SERMON < < <}}
~
\hyperref[sec:index_chronic]{\small{[返順時目]}}
~
\hyperref[sec:index_scriptual]{\small{[返順卷目]}}
~
\hyperref[sec:oQdnBIyZabY]{\small{> > > NEXT SERMON > > >}}
\newline
\newline
$^{1}$我們同心禱告 天父我們很感恩.
今天晚上我們停了因為過年的原因.
和上個禮拜在菲律賓的短訊.
今晚我們能夠繼續來看最後一章的羅馬書.
我們之前也看過關於菲比這位女執事.
她是保羅的一個贊助人.
可想而知這位菲比.
是相當對福音有熱誠.
我們很盼望主感動我們每一位.
我們能夠真的在我們的崗位.
在我們所身處的地區教會.
來一起同心的興旺福音 願神帶領.
我們會看到一大串的人名.
背後一定有很深的保羅的含義.
她很想逐個逐個去聞安.
當然是她所認識的弟兄姊妹.
求主讓我們能夠在一路很快的.
讀這些人名的經文的時候.
都讓我們內心裡面有一個感動.
就是這些人名在二千年後我們去重讀的時候.
雖然這些人已經是在天家在神的面前.
但是當我們去讀這些人名.
又讓我們能夠好像透過聖靈.
讓我們能夠和他們有一個結連.
在主裡面我們去看看這些聖徒.
在他們的人生怎樣的跟隨.
怎樣的擺上.
怎樣為福音為保羅所傳的道.
來去同心同行.
我們仰望在主的手中.
我們今天晚上的時間.
感恩禱告奉耶穌基督的名.
阿們.
十六章我們上一次說過.
很快就讓我們看到寫這封信的時候.
保羅在哥林多這裡.
他現在要托菲比帶這封信去羅馬.
更加就是要羅馬教會對菲比.
有一個在聖徒裡面的相交之道的接納,接待.
而且特別提到菲比.

$^{41}$她曾經怎樣幫助保羅的福音工作.
是保羅的贊助人.
就是說這個菲比不單止是一個普通的女士.
而是應該有當時的一個社會地位.
也都是有一定的錢.
而且更重要就是認同支持保羅的福音工作.
這個是很不容易很不簡單.
如果她住在哥林多.
因為我們就沒有資料菲比是住在哪一個城.
但是如果她真的住在哥林多.
哥林多的人有好幾派.
有些是很反對保羅的.
因為受了那些異端,假的教師的影響.
就對保羅那個使徒的身份是有一個挑戰的.
但是我們看到的菲比.
這位女執事 堅格理的女執事.
你可以想象.
當然你可以說她可能是住在堅格理.
但是我們不知道實際菲比.
她是怎樣的一個當時的狀況.
不過無論如何.
我們可以看一下保羅在這一大段聖經裡面的文案.
我要先點一點出.
因為我們看第一到第十六節.
很快就會讀完人名.
但是讀完人名之後的十七節.
就立即有一個極大的反差.
就是說原來弟兄姊妹的相交.
是需要這樣彼此的文案.
彼此的祝福,欣賞.
無論男女 原來在保羅的眼中.
男女已經是有一個平等的了.
他不會對一些女性有一些貶低.
或者有一些工作是不出色.
是不值一提 這樣.
不是的 反而我們看回第十七節.
就將十六節的經文.
一到第十六節的經文.
帶來一個好像更加凸顯.
這十六節保羅的文案的寶貴.

$^{81}$那個來得珍貴.
因為十七節就是說什麼呢?.
留意這段聖經.
弟兄們 哪些離間你們.
叫你們跌倒 背負所學之道的人.
我勸你們要留意躲避他們.
要留意的躲避.
不斷在講 在教會裡面就要彼此相交.
十四章,十五章.
就說叫我們不要對別人所做的.
你有挑剔 或者你有誤解.
人家喜歡守日 由得人家自己堅持.
或者你自己覺得什麼都可以吃的.
那你自己來堅持吧.
就不要論斷那些不吃的人 諸如此類.
但是沒有到一個這麼極端的狀況.
就是說要留意的躲避.
留意躲避.
所以這裡就給我們看到一個很大的反差.
也可以這麼說 原來一到第十六節.
保羅是很有意識地.
對每一個他所認識的在羅馬的弟兄姊妹.
是對他們為主所服侍的.
那就更加就凸顯了.
原來教會應該大家在教義,真理的裝備底下.
就應該去到一個地步 就是.
真理不是硬梆梆 這些道理我懂了.
而不是這樣的.
懂了之後 原來就要在生活.
在教會的聚會裡面 人與人之間.
大家應該要找到自己的那個事奉崗位.
大家就要在神的裡面 按照神給你的恩賜.
給你的人生的資源 你就要去盡心盡力的去服侍.
以及彼此的配搭 彼此的欣賞.
彼此的來去帶一份感恩 這樣來服侍神.
所以17節反而就是將前面的16節.
就真的帶出一個很大的一個反照.
所以這個很值得我們大家一起去思想.
第18節 因為這樣的人不服侍我們的主基督.
只服侍自己的肚腹 用花言巧語.

$^{121}$誘惑那些老實人的心.
接著就對羅馬教會一個最後的稱讚.
就是你們的信服已經傳於眾人.
所以我為你們歡喜.
但我願意你們在善上聰明 在惡上愚拙.
似平安的神 快要將撒旦踐踏.
在你們腳下 願我主耶穌基督的恩.
常和你們同在 很清楚告訴我們.
有些人原來是背負了真道.
但是他們又要繼續在教會裡面來做一些搞事.
做一些滲透 做一些分裂.
這樣保羅就說 為什麼要留意的.
躲避這些人 不要被他們的心思意念的價值觀.
那些為人的想法 做事的態度.
或者人生的行徑來影響 甚至受吸引.
這樣來不知不覺就偏離了神的真道.
那就是說 又是回到保羅的感動.
什麼是福音 什麼的恩典將我們改變.
我們應該是在神眼中.
到底我們是一個活祭的生命.
活祭就是說 不到我們自己再去.
有一個很強烈的自決 自主的了.
應該讓神的聖靈帶領我們的人生這樣隨從的.
所以這裡就再一次 好像在文安的背後.
原來是有神學的 原來這個文安.
不是原來好像很隨便 誰誰祝福你.
不是這樣的 保羅是我用的題目就是.
深情的文安 但是原來真正背後的原因.
就是鼓勵羅馬的教會.
既然已經有對真理 對福音的認識.
對自己成為了神兒女 有聖靈在心裡面的.
來這樣的一個新生命 新身份的時候.
我們就要認真的在地上來服侍.
這個服侍就是在傳福音的過程.
就是怎樣來在我們的崗位盡心盡力.
所以17節和18節就是剛剛相反的.
有一個這樣的狀況的人.
其實他們是偏離了神福音的真理的道路.
所以原來人生就算信了主.
都有兩條路我們可以走的.

$^{161}$所以信了主是一個概念來的.
是一個生命的皈依.
但是怎樣讓這個生命成為我們整個人生.
每天的生活呢?.
這就是我們怎樣來信而信服.
怎樣來在福音神的國度裡面有份.
所以回來就是一到第16節.
就是保羅有意識地稱讚每一個人.
他們為主所做 這個就是很實實在在.
他們就是因為明白了福音.
因為也是認同使徒們他們所做的擺上的.
他們就願意這樣來服侍.
接著就問了 原來我們看這16節裡面.
我當然還沒有和大家再分類.
再仔細去看.
因為抱歉今天我差不多大半天上了深圳.
因為有個弟兄的公司開幕.
找很多教務上去禱告.
來去有分享 見證著神的工作.
所以趕回來下午也差不多3,4點.
有很多其它的瑣碎事務.
所以來不及做完整張的PowerPoint.
我做了很少.
所以今天可能還是在那個大的概念上來說.
我們回來就是菲比描述她的時候.
跟手就說在羅馬教會.
所以菲比就一定不是羅馬教會.
她是堅格理的教會的女執事.
再一次告訴我們.
和對他們問安.
所以每一個名字的背後都是一個生命.
裡面是有奴隸 特別如果我們沒有搞錯.
這個老夫和他的母親.
這個應該就是一個奴隸.
甚至有些人就說.
會不會是當時他爸爸幫耶穌背十字架那個呢?.
這個老夫就是他的兒子.
當然我們不是太清楚.
但是那個名字原文.
希臘文是有奴隸的字根的.

$^{201}$有些解經家是這樣說的.
也有一些是自由人.
所謂自由人就是他在社會上的那個身份.
就是有那個你可以說是優越.
還有是有一定的社會地位.
例如菲比就是那個贊助人.
就是說好像她能夠去收養一些奴隸.
成為她名下的一些僕人這樣.
當然不是一定那些贊助人.
大部分都是很邪惡,虐待 都不一定的.
因為都需要幫那個主人去做很多的.
就是日常的勞動的工作.
所以你沒理由天天都虐待他.
虐待到他都皮包骨 那怎麼做事呢?.
所以奴隸不是我們有時近代那些被虐待.
被殺害很少的 所以在羅馬的帝國裡面.
差不多有兩三成聽聞說都是奴隸的身份.
所以是全個國家很重要的勞動力.
所以奴隸某程度都有一定的自由的.
有一定的社會身份.
所以就要看他跟哪一個主人.
如果那個主人本身自己都是一些貴族.
或者一些皇室 那個奴隸就更加有一個.
超越的身份 這些都是當時羅馬教會.
或者保羅的時代 當時的希臘文化.
和羅馬帝國那些的制度 我們都稍微了解一下.
因為保羅在不同的書信.
都有時會用這些當時的羅馬制度的概念.
來講解一些真理 例如現在太書講到.
有些奴僕來去管束那個主人的子女.
就是他們真正可以脫離那個自由成長之前.
都是在那個奴隸手下 或者叫管家.
那個管家本身可能是奴隸來的.
來到他要管束他 讓他能夠一來保護他.
因為怕可能被標心或者等等.
但是更多是要用主人的吩咐.
最嚴格的那個教導 使他能夠成材.
所以這個加勒太書就引進這個概念.
當然在保羅的《稱呼文案》裡面.
這裡說有猶太人 有外邦人 這個不奇怪了.

$^{241}$因為保羅已經非常強烈的意識到.
在他的福音裡面不斷將猶太人,外邦人.
來到去已經在一個合一裡面.
這個還沒有去到以弗所書那些那麼強烈.
在羅馬書裡面已經讓我們深深感受到.
在教會裡面的聚會已經有猶太人,外邦人.
所以一定有他們的文化上的差距.
有他們在一些特別是舊約.
和主耶穌成了舊恩之後.
那個新的身份應該怎樣協調.
應該怎樣來突破.
這些就是保羅在他人生.
尤其是後半段的時候.
對於這個課題 猶太人,外邦人.
到底在耶穌裡面是怎樣的一個新的關係.
新的身份 以弗所書,歌羅西書.
就相當詳細去講解的.
但這裡保羅就用很實在的.
對猶太人的文案 對外邦人的文案.
對一些的男一些的女的文案.
有些人是有錢人 也有些人是貧窮的.
都真的沒有什麼社會地位.
比如剛才講的老夫就是一個可能的奴隸.
這裡就讓我們看到保羅這個文案.
其實就是非常反映到在他的神學裡面.
對不同的人的尊重,欣賞,稱讚.
還有在他的眼裡面.
再沒有高高低低的身份的懸殊.
所以這個是我們骨子裡面需要留心的.
我們對其他人真的需要有這份尊重和欣賞.
就不要有一種不自覺的見高拜見低就踩.
這個是聖經非常之禁止的.
也都近乎我們完全好像一個.
如果真的時時都還是這樣的心態.
那就可能真的和神的國度是沒有分的.
因為神是很禁止我們這樣.
有這一種這樣的心態的.
所以我們求主要幫助我們要了解明白.
所以應該其中有一句話.
就是羅馬書「倒要苦就卑微的人」.

$^{281}$所以保羅就告訴我們.
我們有自己神給我們的身份.
我們自己在人生裡面的成長.
比如我剛才都講了.
小時候住木屋區周山走.
但是經過了很多年社會的進步.
有讀書,有成長,有工作.
各方面經過了數十年之後.
真是那個社會的地位是提升了.
但是提升了 保羅就說提升本身沒有錯的.
但是如果連帶我們那種的驕傲,優越.
鄙視,歧視別人都一起提升 那就是問題所在了.
所以有那個身份,社會地位是沒有問題的.
但是保羅就說「倒要苦就卑微的人」.
卑微的人包括了卑微的事.
所以我們千萬不要看.
有些事情我一定不做的.
有沒有搞錯 要我做.
我是什麼人 你知不知道啊.
就不是這樣的.
我們應該去到在主的救恩裡面.
我們自己的生命都改變的.
我試過的 其實不止一次的.
在錦繡的禮堂有時候看到一些廁紙.
不是特別骯髒的 普通一張的廁紙.
走過都沒有人收拾的.
可能有時候看不到或者不留意.
但是我看到我就一定收拾的.
你說不關我的事 應該叫誰.
現在就是Kitty了 叫Kitty來 叫Kitty收拾.
就是菲律賓的姐妹.
還是說阿冰收拾.
其實你看到 你就收拾起它.
為什麼要分那個事 不是你的身份呢?.
當然每天的工作 我們有事務員.
他幫忙做清潔 要洗廁所.
這個是他每天的工作.
當然沒有理由我的工作跟他對調.
也不應該是這樣的.
但是有時候在那個.

$^{321}$剛好你走過看到.
你就突然間說 我的身份不適合做這件事.
我怎麼都不會收拾的.
其實這樣 保羅的教導就是說.
倒要苦就卑微的人 或者卑微的事.
十二章的羅馬書 這樣給我們看到.
所以保羅的神學真是很落地的.
保羅對福音 對主耶穌從高天來這個世界.
到神祂犧牲祂的愛子 成為了贖罪祭.
我們現在在主耶穌裡面 我們得到永恆救恩.
我們整個的生命都更新改變.
所以我們真的慢慢我們的社會階層提升.
你想一下 我現在有時候都要打領帶.
每個星期天都要打領帶 穿西裝.
肯定是給我的哥哥 我哥哥開巴士的.
我姐姐也是開巴士的.
肯定給他們兩個.
那個社會表面出來的形象高一些.
按一般人的看法.
但是那個心態,心境千萬不要.
這個是很重要的 所以不容易學的.
所以羅馬書給了我們很多倫理的教導.
給了我們很多待人接物的教導.
所以在這個文安裡面.
其實我們直接代入去都已經很感動.
因為我們真的看到保羅.
是對每一個人他都是認識.
當然這個是他所認識的這26個朋友.
他們原來不是教會的什麼大目,領袖.
原來是很卑微的一些接待的奴婦的人.
所以你就可以想象到.
使徒保羅給我們的教導.
這些我們上一次.
就是三個星期前已經講過了.
那就不再詳細講這個部份了.
就是卑鄙就要來去接待他.
要合乎聖徒的體統.
就是說要對他的那個事奉.
他可能有一些重要的事務.
他們必須要按照神的真理來去接待.

$^{361}$就是不要因為他是特別的尊貴.
就特別的來去上賓的款待.
就不是這樣的意思叫做聖徒體統.
而是真的要在主裡面.
大家來彼此的服侍和欣賞這樣的情況.
所以不是因為卑鄙是保羅的贊助人.
不是因為他是一個有錢的女執事.
所以特別的招呼得殷勤.
而是因為他是在神的角度裡面.
來去支持神的僕人保羅所傳的福音.
所以是要這樣的心情去接待這位女執事.
所以原來初期教會都已經在這樣的水平.
去看男女之間的角色,身份.
香港近代 可能大家應該都明白.
我們都是80年代 上世紀的80年代.
才正式可以按納女牧師宣導會.
浸信會可能再早一點 播道會就更加遲了.
這些已經是我們很敬虔的福音派了.
我們都這麼遲疑 對男女的角色.
都有擔心中國人的社會.
還是對好像這個平等的聖職.
還有一個這樣的擔心 怕違反聖經.
這樣的情況 所以譬如我成長的教會.
福音堂女性 當然現在我不知道.
因為離開了很多年 女的姐妹.
就算很有恩賜 星期天的主日講道.
沒有她份的 她可以做主一學,教主一學.
但是禮拜的崇拜講道 福音堂的教會.
女性是不准講道的 你可以想像得到.
原來我們的文化 我們對聖經的理解.
都可以到了今天 21世紀了.
都可以是不同的 但是你可以想像得到.
在保羅的書信裡面 保羅的教導裡面.
其實保羅原來可以是這麼寬的.
對女性的身份的看待.
原來可以是在她的眼中 原來是這麼平等的.
所以我們可以想像一下 我們有什麼要不斷的更新呢?.
或者我們怎樣能夠更加全面對聖經的教導.
我們有一個與時俱進的成長呢?.
這些都是我們在這個文案裡面.

$^{401}$可以彼此提醒的.
我做的PowerPoint在這裡開始.
Patreon就是一個推薦人,贊助人.
在財富,地位,法律的能力上.
提供給受幫助的 比如法律的保護.
有一定的人脈關係,社會的地位.
尤其是推薦信 是很重要的.
如果沒有那個推薦信.
其實你就很難在其他人面前得到那個信任.
或者信納的 這個就是那個受庇護者.
他怎樣去回報這個Patreon呢?.
就是怎樣來回報他的恩惠呢?.
就是對他的主人 或者叫贊助人忠誠.
要公開的支持他 在社會上的聲望,服務等等.
這些在這個羅馬帝國的時候.
就相當普遍的維繫社會的地位的一個制度.
是一個很公然的制度.
某程度上我們中國人的上一代.
我們爸爸的年代 有些人他是很有頭有臉的.
尤其是一些在地區上有一定的身份的.
有很多的所謂的跟班.
那個跟班就大概是一個階級人.
當然這個比較負面 呼呼喝喝.
又或者是有時候做一些不好的事情.
但是羅馬帝國的那個制度.
就會是Patreon和階級人這樣的一個受助人制度.
可能每個社會,每個時代都會有的.
問題是那個形式是去到什麼程度呢?.
這些大家我們知道一下都好的.
保羅其實從來沒有去過羅馬的教會.
當然不是說他沒有去過羅馬.
因為他是一個相當有那個的社會地位.
他自己本身都是因為他爸爸的緣故.
他是羅馬的公民.
但是他也是一個很純正的猶太人.
第八天就行國禮的保羅.
所以他進了羅馬籍的.
但是就是說他有機會去接觸.
比較高的階層的那個社會.
當然在猶太就不在話下.

$^{441}$他就在基里加的大塑那裡成長.
又跟這個加瑪烈一個很出色的拉比.
成為他門下的一個門徒.
一直長大在猶太的社會都有一定的身份.
後來更加有羅馬公民的身份.
就是說他是擁有一些雙重的國籍.
還有在社會上有相當的影響力的一個保羅.
正是這樣的人神就用他,呼召他,改變他.
大大的使用他.
當然這個過程他自己需要有聖靈的工作.
因為他就是一個非常對於自己的猶太教.
有一個的執著,有一個的捍衛,有一個的持守.
所以就令到他越來越是一種的傲慢,偏見,執著.
對於自己的教義好像是滴水不侵.
差不多已經很透徹的解釋了,明白了.
沒有人再可以挑戰他對律法的那個漏洞.
律法應該怎麼執行,怎麼生活.
保羅應該是相當的具體的實踐的.
但是當他從來沒有想過一道的強光照入他的心靈.
當然先照入他的眼睛 肉身就是盲了.
接著其實是照入他這個律法的心靈.
其實是一個好像舊約耶利米蘇32章所說的.
將他的石心變成肉心 變成一個有血有肉.
有愛的一個門徒 基督徒.
因為主耶穌向他顯現.
因為他所認識的耶穌是釘在十字架上.
受咒詛的拿撒勒人而已.
怎麼會可以在天上向他顯現.
用希伯來話跟他說話 告訴他.
蘇羅為什麼你要逼迫我.
你用腳去踢刺是難的.
所以對保羅來說是晴天霹靂.
一個他以為是他已經完全參透了.
當然是佛家的那些說話 也沒有所謂的.
他參透了舊約.
認為他的人生就是一個很出色的.
為律法去捍衛的拉比.
但是當耶穌一句話就令他重新再思考.
到底什麼才是神眼中的信仰呢?.
什麼才是神眼中的一個自由呢?.

$^{481}$什麼才是這位做物主的眼中的公義呢?.
他整個人好像完全崩潰了.
重新完全不認識自己.
這樣就進去大馬士革敘利亞 在一個叫直街那裡.
就在三天三夜 神也挺搞笑的.
不是馬上同步叫阿拿尼亞.
你趕快去 在門口等保羅.
那個時候我想保羅也未必會明白 未必會相信.
就讓他有三天三夜的反思 禁食.
來去思想 到底自己現在怎麼辦呢?.
因為耶穌最後一句話就是.
我就是你所逼迫的耶穌.
起來往直街去.
在那裡就必定會有人指示你應該怎麼做.
保羅就被人拖進去直街裡面的其中的樓房.
在那裡三天三夜的等候.
所以對保羅來講 他現在傳道一段時間了.
應該是第二次的宣教旅程完結了.
準備他要上耶路撒冷.
這個我們十四,十五章講過.
但是羅馬教會對保羅來講 其實還是陌生的.
其實大部分的解經家都認為.
是《使徒行傳》第二章五旬節的時候.
因為有些羅馬居住的猶太人.
上去耶路撒冷守五旬節 信了主 洗了禮.
接著後來輾轉回到原居地羅馬建立教會.
所以那個原先的founder.
就是在耶路撒冷信主的這些猶太人.
回到那裡之後就建立羅馬的教會.
所以就不是由保羅去建立的.
但是輾轉在 因為當時.
大家都應該讀過歷史都知道.
羅馬帝國奧古斯都其中的偉大的成就.
就是做好了那些馬車的寬闊的道路.
那個不僅僅是這樣 而是開始因為統一.
軍事各方面都穩定了 就不會有土匪強盜.
所以變成旅遊就很旺盛了.
就意大利或者歐洲的人就可以到處去.
當然是坐馬車 這樣去不同的地方.
可能是旅遊 可能是經商.

$^{521}$可能是大家 特別是當時的哲學家 講學.
所以很多時候這樣的流彩.
就會有機會彼此的認識.
所以這次我在菲律賓都認識好幾位的教務.
都是很用心的擺上的.
也都有一些是菲律賓當地的信徒領袖.
有些住在很偏遠的海島.
來放下他們的工作 來受培訓.
準備將來他們回到自己的地方來傳福音.
所以很值得去敬重的這些僕人 這些侍女.
他們未必有太高的文化.
但是他們這六天就是裝備對聖經的.
盡快掌握整個聖經的脈絡.
所以那個課程都相當嚴謹的 教導聖經.
接下來九月我去 如果去的時候.
可能我都會幫忙去講一點.
所以我們來到看到保羅.
他對這一幫的人 這裡我應該按錯了.
拿了另外一個圖.
因為今天實在不好意思 我回去會再整理的.
因為今天下午四點多才開始預備.
不夠時間去再整理那些細節的PowerPoint.
現在我們正式和大家再讀從第三節開始.
到第十六節 再給大家思想一下.
我們看這個文安背後的神學.
還有保羅他在整個救恩 整個的真理裡面.
他怎麼看其他的弟兄姊妹呢?.
這個就是主裡面的相交 主裡面的相愛.
第三節開始 因為這個非比我們不再講了.
文伯基拉 阿居拉安這對夫婦.
往往都是伯基拉先行的.
很特別的次序調亂的.
多數人都認為伯基拉就是一個比較行動,領導.
就不是說她壓制著她的丈夫.
而是很多時候她可能比她的丈夫更加的熱心.
所以就有時這樣的出現 這樣女性先行.
似乎未必是一種貶義.
更加不會是保羅對伯基拉特別的有欣賞.
而貶低她的丈夫阿居拉 那就應該不是的.
有什麼不同的地方呢?.

$^{561}$留意一下 他們在基督耶穌裡與我同宮.
也為我的命將自己的頸項致之道愛.
我們可以看看第18章的《使徒行傳》.
就是阿居拉夫婦正式出場的偶遇.
18章《使徒行傳》.
第一節開始 這是以後保羅離開了雅典.
就來到哥林多 遇見一個猶太人名叫阿居拉.
正式開始介紹這個弟兄 這個猶太人.
他生在本都 保羅就是基里加的大宋.
很近的地方 因為格魯雕.
格魯雕就是羅馬的第四位皇帝.
這個格魯雕命猶太人都離開羅馬.
就是當時有一個逼迫 有一個要趕走猶太人的行動.
就是由格魯雕 就是由皇帝來下令要趕走猶太人.
所以他就離開羅馬.
身近帶著妻子阿居拉從意大利來.
來什麼呢? 就是來到哥林多.
保羅就投奔了他們 他們本是製造帳篷為業.
保羅因與他們同業 就和他們同住作供.
每逢安息日 保羅在會堂裡辯論.
勸化猶太人和希利尼人.
保羅就簡單講了這對夫婦的身份 背景.
還有為什麼在哥林多相遇.
他們就落腳在哥林多 他們平時就織帳篷.
就是說他們的工作 因為猶太人有一門手藝.
就是說猶太人他們現在當然是信了主.
但是他們有一個挺好的訓練.
就是都需要有一些技能.
這個織帳篷似乎在當時是挺能夠養活自己的.
保羅也是織帳篷的 原來保羅既是拉比.
又是一個平時他都有這樣的.
這個不是手工藝 是一個很重要的技能.
可能有些人是做鐵匠 或者有些人是做其他的工作.
所以原來保羅自己作為使徒.
他現在就到處去傳福音 神的僕人.
但是他都會這樣織帳篷的.
當然更多保羅就教導.
特別因為哥林多是有他的問題困難.
但是在其他的教會都是會用獻金.
來供養保羅和他的同工的.

$^{601}$但是這裡保羅就給我們有一些這樣的概念.
他因為認識了阿基拉夫婦.
知道他們兩位平時沒有聚會的時間.
他們就會編織帳篷來維生.
可能藉著這樣又可以結識當地的人.
接著保羅就平時跟他們一起同住作供.
一起來到去 我不知道當時居住的條件.
或者是前佈後居等等是怎麼樣.
但是我們就自旅行間給我們看到.
路加醫生的描述 保羅很尊重這對夫婦.
又看到這對夫婦似乎跟保羅有相當的默契.
那個靈性的程度都有一個相近.
不只是一對普通的夫婦.
似乎都是為福音,為神的角度.
輾轉我們去到第19章都再有阿基拉夫婦的.
但是我們不詳細再講阿基拉夫婦.
不過從保羅的文案一開始除了非比之外.
立即就對阿基拉夫婦有這樣的稱讚.
就真的讓我們有一種嘆為觀止.
這對夫婦被趕走離開羅馬.
似乎連他們在羅馬有的產業.
有的財富都一下子被沒收了.
所以神的兒女我們都不要以為是百分之一百的穩定.
我們都要求主給我們有一個.
隨時可以為主而受苦的心智.
因為香港好像從來都對這些事情覺得是不可能.
但是其實在很多國家的信徒.
那些的信徒其實人家真是為了福音,為了信仰.
真是有時連他們應該有的那個基本的生活.
都可以放下的.
所以求主給我們有一些眼光.
有一些人生的機會去認識一下這些的牧者.
這些好像很土豹的神的僕人.
他們甚至乎不是認識很多字.
但是他們那個為主的心就是很熱切.
就是這樣真是願意為福音擺上代價.
都是在所不惜.
所以現在我們就見到這個阿基拉夫婦.
為了保羅的命.
他們可以將自己的頸項截之度外.

$^{641}$很誇張的描述.
就是說如果保羅坐牢.
保羅被人抓.
他們是願意去站在保羅的一邊.
不會說現在見到危險.
我就和保羅割席.
我就來到他有這麼遠走這麼遠.
如果是這樣的時候.
很可惜見異思遷有好處就馬上纏起來.
這樣的居心,這樣的動機.
這樣的人生的氣節其實是很可恥的.
因為真正服侍的是我們的主,我們的神.
所以一句話告訴我們.
致之於道外的阿基拉夫婦.
所以保羅是對他們相當敬重.
這就意味著保羅在這兩次.
第三次還沒有發生.
在這兩次宣教旅程.
真是受了很多的苦.
在路斯德,在以哥廉.
被人用石頭打 打到死.
拖出城外 準備就將他下葬了.
保羅突然醒了 其實是用石頭打.
不是打幾巴掌 蓋幾巴掌算了.
是用石頭打 就是說可能流了血.
爆了肝 但是保羅起來.
那些人都害怕 保羅就直接回到城裡.
你可以想象得到 可能是神蹟.
可能是保羅真的死了.
神令到保羅醒了 復活.
聖經沒有這樣描述.
只是說他醒了就進城了.
繼續去傳道 然後才走.
這是十三章,十四章的經文告訴我們.
但是告訴我們 保羅很多時候.
都是身處在全福音的險境.
但是更加寶貴的就是有同路人.
有弟兄姊妹 阿居拉夫婦不是傳道.
可能只不過是信徒領袖.
但是這樣來與使徒同在.

$^{681}$同一樣的受那個苦.
甚至乎是隨時願意為保羅.
走到那個苦難裡面.
所以 弟兄姊妹 我們人生有沒有這樣的同路人.
有沒有這樣的一起願意為主服事.
為主受苦呢?如果沒有.
那我們要找 或者我們自己成為別人的同路人.
人生短短幾十年 我們真的要開始.
有時候要拋身出來 離開我們的舒適圈.
剛剛我星期 應該不是星期一 星期一回來.
我們星期天那一晚 我們最後的檢討分享.
我都在眾弟姊妹面前講.
我說我從來沒有想過要去菲律賓的.
菲律賓我從來不會去 也都不羨慕去.
也都喜歡去的 我一路上車.
由機場下車去市中心.
進去居住的那個 我們那個短訊中心.
那些又塞車 又討厭的那些交通.
一路在車上都是這樣說的.
但是當我見到那些垃圾村的小朋友.
就不是讓我對這個國家有改觀.
對這個國家完全我沒有改觀到.
只會是更加肯定這個國家是沒得救的.
但是裡面的人是很值得我們要.
擺心擺力的去做 因為那些小朋友天真爛漫.
有個小朋友黏著我 那個女孩子很可愛的.
黏著我 我說抱抱 她馬上讓我抱起.
和我一起到處走 然後有些人幫我們拍照.
我們拍個V 她又跟我拍個V.
是讓我感受到這些小朋友完全.
是好像天使一樣的那種燦爛笑容.
他們是完全無辜的 因為他們還沒離開過那個村.
還不知道外面的世界 比如馬尼拉.
比如菲律賓那種的問題.
那種的結構性的那些制度的腐化等等.
但是他們如果沒有書讀 沒有一個基督教的教育.
來去遮蓋他們 幫助他們.
他們很快上小學就會學壞的了 那些女孩子.
希望盡快有男朋友 就可以脫離自己的垃圾村.
其實不會的 因為當地的男女關係是很普遍的.

$^{721}$那些男的隨隨便便有好幾個女朋友.
還要是那些女朋友是不介意的 很複雜的.
有一個垃圾村的一個我們資助他們的一個家庭.
家裡有10個子女 另外現在在那個學校裡面.
已經是我們資助的有三姐妹.
她們現在在我們面前分享她們的成長.
她們的故事 很流利的英文.
就是在學校裡面學習 三個姐妹 親姐妹.
是同一個媽媽 但是就三個爸爸.
那你可以想象得到 菲律賓就是這麼結構性的了.
就是這麼恐怖 這麼凌亂的了.
所以我們回來的時間 我們要留意就是說.
我們有沒有一個跳出舒適圈的這樣的心情呢?.
在我們人生還可以少許的時間 事奉神的機會.
其實我們都已經不是住在那些地方.
永遠不能翻身的世界 我們都是一個星期.
我們就回到我們自己的香港 我們去幫助.
其實我們都是盡力做到一點點.
幫得一個就得一個 一個改變就改變到一個.
這樣的一個心情 所以我自己是說很大的感動.
真是當下有一個很大的內心的感動.
我一定要將我所見所聞 帶回來香港我的教會.
給我的弟兄姊妹知道.
我們真的要跳出我們的舒適圈.
所以保羅就是有這一位阿居拉夫婦.
其實是萍水相逢的.
其實人家在意大利是做什麼的呢?.
這樣能夠認識真是大家在主裡面.
現在能夠有機會這樣去相遇,同住,同作供.
但是人家阿居拉夫婦為了保羅傳福音.
將自己的頸項置之於道外.
就是說原來我們所做的小小的工作.
是會得到別人的肯定,讚賞.
而且是宗教會保羅 就是說保羅去到不同的教會.
都會去對阿居拉夫婦來稱讚有加.
不是要巴結,不是要抬高抬舉.
而是別人的善行.
別人為了你所做的那些善行.
你來到主裡面去報答.
還有將別人的善功.

$^{761}$為主所做的.
來到其它的教會去分享.
去激勵其它的弟兄姊妹.
所以我真的很希望.
查經班的弟兄姊妹 你們聽到的.
將來有機會透過看YouTube聽到的.
都真的能夠 你一定要找回你自己.
要為主受苦的一些事奉.
要跳出我們自己的舒適圈.
來為神擺上.
因為如果我們不跳出他們的舒適圈.
就沒有事情可以做了.
他啊 誰叫他在菲律賓住嗎?.
誰叫他們那個國家 那些制度這麼差嗎?.
香港我們現在有這麼好的福利制度.
是不會有窮人 是不會餓死人.
正正就是別人會 正正就是那個地方沒有.
正正就是我們是基督徒.
現在我們已經進步到.
真是有一個最基本的社會福利.
所以我們就要幫助還沒有的.
這些 為什麼我會這麼大其中的觸動就是.
我小時候就是這樣的 我們就是這樣.
叫做穿著褲子 就在街上沒有拖鞋就這樣玩.
這樣玩大了 這樣去學說髒話.
接著那些人吸煙 就是這樣.
在家裡比較嚴 所以沒有進入黑社會的意思.
但是那些小朋友其實就是無辜的.
所以我們就不應該讓那個這麼不完善的.
那個制度去殘害這些新生命.
所以很感恩 就是這次都有感動的.
就是跟Chris過去看這個野地牧人.
其實野地牧人就沒有什麼好看的.
就知道是一幫信徒領袖受栽培,受訓練.
但是原來第二天,第三天.
去到這些真是垃圾村.
那就真是一個很大的一個衝擊.
對我的心靈的衝擊.
因為我們從小就是在牧區長大.
是分外了解 對於其他的 譬如我的兒子.

$^{801}$他是在環境改善之後的成長.
可能他覺得很髒 快點走 很臭.
但是我在這樣的環境長大.
一些我都不覺得臭 一些我都不覺得髒.
因為我根本就是這樣長大.
但是問題就是 我們是過了那些.
幾十年的階段之後 看到整個香港.
整個亞洲都有改善 但是菲律賓是沒有的.
就是說神仍然要感動和挑戰屬主的兒女.
要負起這個責任去幫助這些小朋友.
透過教育讓他們的生命得到基督.
他們的生命來有一個動力.
盼望我們是阿居拉夫婦的使徒的同路人.
我們衛著福音 我們有同路人.
不同國籍都可以.
接著第五節 又問在他們家中的教會安.
這裡再一次告訴我們 原來阿居拉夫婦.
成為家庭教會的領袖.
意味著出錢,出力來設立教會.
初期教會有這種的形態.
來有這些小型的聚會 一定有大教堂.
後來當尼祿迫害教會的時候.
他們就要轉到羅馬的地下那些地洞.
所以將來我也有一個心願.
帶弟兄姊妹一起去羅馬 其實不會太貴的.
去看那些地洞的教會 就是使徒時代.
他們在這些地洞裡面聚會的那種情況.
我們感受一下.
土耳其都已經見過那些地下教會.
那些墓穴教會 但是羅馬就另外一種的形態.
又是中環交錯的那些地洞.
某程度上可以想象到.
就是逃避羅馬帝國的士兵的追殺等等.
所以問他們的教會安.
就是說原來他們都負起了一個.
本來打算不讀的了 都讀了第19章.
《使徒行傳》的第19章.
所以阿哥拉夫婦他們不僅是出錢出力.
他們都裝備他們的神學.
他們幫助阿波羅.

$^{841}$請大家來留意這一段經文.
對不起 18章就行了 18章24節.
18章的24節《使徒行傳》.
有一個猶太人名叫阿波羅 來到以弗所.
他生在亞歷山大 是有學問的.
最能講解聖經.
這人已經在主的道上受了教訓.
心裡火熱 將耶穌的事詳細講論教訓人.
不過留意盧家醫生的很仔細的描述.
指示他單曉得約翰的洗禮 單曉得.
就是說他不知道耶穌基督拯救人的那個洗禮.
那個十字架的救恩.
不知道真正的那個靈魂的洗禮.
只是知道一個肉身的禮.
某程度施洗約翰的洗禮就是一個禮儀性.
因為它是在耶穌釘十字架之前的洗禮.
就是預備人的心接受將來的基督 施洗約翰.
所以他那個洗禮 如果你只是受了約翰的洗禮.
你就沒有救恩和救恩無效的.
因為你根本沒有接受耶穌.
你只是預備心 準備要認識將來的彌賽亞.
所以這個是約翰很清楚的那個教導.
他說我不是基督 我是告訴你誰是基督這樣的意思.
所以只是洗了約翰的洗禮是不行的 不能救的.
就是說有些暗示 阿波羅這麼熱心.
又最能夠講解聖經 當然是在講舊約聖經.
又在主的道上受了教訓 那就糟糕了.
最重要的是他不懂 他還沒有 他不通.
就是說每樣東西他都很有辦法去講解 清楚明白.
但是問題就是最重要最關鍵的.
耶穌和你的生命關係是怎樣呢?.
留意了 第26節.
他在會堂裡放膽講道.
伯吉拉亞居拉聽見.
就接他來 將神的道給他講解更加詳細.
他想要往阿該亞去.
阿該亞就是哥林多省.
弟兄們就勉勵他 並寫信 這是戰信.
請門徒接待他 又是需要推薦信.
他到了那裡 多幫助那蒙恩信主的人.

$^{881}$在眾人面前極有能力 搏倒猶太人.
引聖經 證明耶穌是基督.
伯吉拉亞夫婦聽他講道.
覺得有些東西好像可以再給他一些更新.
可以有一些突破.
甚至可能發現他有一些屬靈的.
真理上的小小的偏差.
但是經過了他們再講解詳細明白.
約翰的洗禮 和主耶穌的洗禮.
和主耶穌成就救恩的那個福音.
清楚明白 就真是如融會貫通了.
打通了阿波羅的任督二脈.
真是完全的通了 他就更加有能力了.
更加能夠用他的教導 用他的講道.
來引聖經 證明耶穌是基督.
搏倒了很多的猶太人.
所以原來亞居拉夫婦.
又是認識聖經透徹的信徒領袖.
又是家庭教會的領袖 又是保羅的同伴.
那就是保羅在哥林多的那一年半.
亞居拉夫婦在保羅的身上感受到的.
使徒的不只是權威 而是使徒的心智.
使徒的為人 使徒為主的福音的那份熾熱的心靈.
都感動亞居拉夫婦願意和保羅同工.
願意留守在哥林多 保羅到處去傳道.
他們繼續留守在哥林多的時候.
就認識了這個阿波羅.
就再將阿波羅變成另一個像保羅一樣.
被神大用的傭人.
所以亞居拉夫婦這對信徒領袖.
原來都幫助不同的人進入神的國度被神大用.
所以保羅投奔這對夫婦.
這對夫婦在保羅的內心是很值得敬重的神的僕人.
所以很感恩的.
我們能夠看看這些初期教會的人.
都給我們有很大的感動的.
盼望我們大家都一起去追求.
今晚我就想講到這裡就算了.
我下一次再去看.
我還有一兩個 因為這個拍檔.

$^{921}$連那些人的樣子我都還沒有改變.
所以就不是這麼好意思.
今晚就講到第四,第五節.
我們一起有一個簡單的禱告.
天父啊 我們真的很盼望你的愛感動我們.
被神的福音來改變.
多謝你將我們拯救.
多謝你讓我們在香港信主.
一個相當穩定富庶.
可以自由傳道的一個地方.
其實不是每一個地方都可以自由傳道.
現在菲律賓真的可以很自由.
甚至我們可以隨便去幫助這些兒童.
都可以得到政府的許可.
這裡讓我們看到.
如果我們現在不好好做的時候.
有時候會錯失了這些機會.
願神來幫助 感動我們.
多謝你透過Chris的團隊.
讓我們知道他在不同的地區.
他們在無論菲律賓,中東,埃及.
都是我們自己的弟兄姊妹.
這樣出去將這些有需要的人.
特別現在中東一帶.
特別伊朗的那個痛苦,那個戰爭.
是帶來很多的聖靈塗炭.
我們很盼望主來施恩給伊朗裡面的居民.
可能他們在惶恐裡面,痛苦,懼怕裡面.
但是我們很盼望他們能夠找到福音.
向你去禱告,去尋求真理,願主幫助.
我們很感恩 因為我們能夠透過使徒保羅.
給我們有很深的對初期教會的了解,認識.
我們多謝神 願你再一次在我們大家.
進入就快去到復活節的時候.
我們的內心都是被主受苦的那個痛苦,感動.
和主耶穌是死而復活,永活的主.
叫我們將一切的事情.
都不要在一種絕望的心情去看.
反而在主裡面我們來到去想念.
我們是與主同工的人.

$^{961}$我們仰望,感恩,禱告.
奉耶穌基督的名.
阿們.
\newpage



\section{羅馬書}
\label{sec:oQdnBIyZabY}
\textbf{2026-03-13 周五晚查經班《羅馬書(三十九):分題:來自保羅的問安(二)》 講員:曾裔貴牧師}
\newline
\newline
連結: \href{https://youtube.com/watch?v=oQdnBIyZabY}{\texttt{https://youtube.com/watch?v=oQdnBIyZabY}} ~~~~ 語音日期: 2026-03-16
\newline
\newline
\hyperref[sec:4pLVUVJIfgU]{\small{< < < PREV SERMON < < <}}
~
\hyperref[sec:index_chronic]{\small{[返順時目]}}
~
\hyperref[sec:index_scriptual]{\small{[返順卷目]}}
~
\hyperref[sec:HW8W_v51HY8]{\small{> > > NEXT SERMON > > >}}
\newline
\newline
$^{1}$好 各位晚安 我們一起開始之前同心禱告.
很多謝天父 我們再一次有機會來研讀羅馬書.
有系統地將這本書.
試圖保羅他對羅馬教會的教導.
將真理有系統地去陳述和宣講.
一封書信保羅準備要去羅馬的時間.
就來向當地的教會講述一個這麼重要.
這麼偉大的真理 就是人罪人可以恩順清義.
可以在主耶穌的救贖裡面有份.
而來得到神的無條件的救贖.
我們在這幾次有機會去思想.
試圖保羅對差不多26位弟兄姊妹的文案.
裡面全部是有名有姓的人.
也提到他們一些好像在生活上.
在事奉上的一輪半找的描述.
可以想像這是好朋友的一個文案.
不是一些陌生人.
我們盼望我們在這樣的教導裡面.
有時不經意的陳述.
讓我們看到原來使徒們.
他們都已經這樣將真理落實在生活裡面.
在教導裡面 願神都幫助我們.
今天晚上我們再一次去嘗試.
逐個的聖經人物這樣去看的時候.
都讓我們能夠掌握使徒保羅教導背後的那份心意.
我們等候仰望在主的面前.
我們這樣禱告 奉耶穌基督的名 阿們.
我們算是第二講 其實應該是第三講的.
因為之前很概略的講了菲比這位女執事.
她是那個受助人 或者叫贊助人.
贊助保羅的 就是說是她去支持一些福音工作.
聖經就沒有再提得太多.
堅格利教會的背景.
保羅在那裡有什麼的事奉.
聖經就沒有記載 《聖經行傳》都很少有記載.
所以我們對堅格利教會都不是太認識.
但是無論如何 這封信就是透過菲比.
就帶去給羅馬教會 而且她有些重要的事情.
需要羅馬的教會來推薦.
可能有其它的地方菲比要去.

$^{41}$所以這個就是一個這樣的動機.
我們再看回上一個禮拜.
我們有機會很概略的看了.
我們就要留意 就是上一個禮拜都很快點出了.
一到第十六節 十六章就是告訴我們.
保羅向很多的不同國籍 不同的身份.
不同背景的羅馬的教會弟兄姊妹.
來去祝賀的民安 有些是夫婦的民安.
但是在這個民安的背後.
就是一份很大的熟絡.
親密的在主裡面的關係.
但是講完十六節之後 十七節就分外重要.
因為十七到二十節就提到.
原來在教會裡面有一些人是我們需要留心.
需要去留意 那個留意就是因為.
他們可能有幾個問題是會影響教會的.
所以第十七節開始 他說弟兄們.
那些離間你們 離間是一個問題.
就是說原來人的骨子裡面很容易有分裂的.
很容易就是對一些自己不喜歡的人.
我們就會有一種標籤 慢慢就會將它妖魔化.
慢慢就變成不是我的族類 我們就要來排除.
就是不會再有相交 更加再嚴重一些就是.
刻意去尋找那些問題來攻擊.
這些離間的 接著說叫你們跌倒的.
就是說可能他們的行徑,主張,講道.
會影響了那個對真理的認識,追求,熱心.
跌倒 第三個就是背負所學之道.
更加是刻意的對真理有一個抵觸,違背.
甚至挑戰 甚至傳一些假道.
所以這三樣東西 保羅就要求弟兄姊妹一定要留心.
保羅就說我勸你們要留意躲避他們.
到底是不是保羅聽聞有這樣的人已經在羅馬教會.
還是一般性的呢? 因為講完16節之後.
17節就立即一個語重心長的勸告.
要小心這些人 可以想象得到.
一定是這些人在教會裡面行走.
是有一定的影響力 而那個影響力.
是一個不是正面的 對神,對真理,對人的愛.
是一路一路有增長的 而是會不知不覺.

$^{81}$會令到人受離間的 會令到人跌倒的.
又會令到人背負所學 即是學過的了.
學不懂之外 還要背負 就是背著那個真理.
即是背向那個真理 即是說背道而馳這樣的行徑.
那就變成不是真理的信徒了.
就不是捍衛那個恩信清義了.
所以這些人 保羅就說 如果真是發現這樣的人.
而他是有這樣的存心,有這樣的基心的.
那就要留意,要躲避這樣的人.
所以我們今天回來仔細去看.
在這個整個的民安裡面.
對不起 那個風吹著我 令到有些.
拿了個電視 我看怎樣弄高一點.
不好意思 吹到靜一靜那個位置.
如果有個兜兜著可能好一點.
算了 算了 現在好像有些動了.
好像沒有什麼反應.
我們是有反應的.
要一秒鐘才收到那個訊息.
好了 不要管了 不動了 有得看了.
我們在整體去看整個民安.
我們可以其實有一些歸納的.
因為不單止是對這個呂執事的稱許.
對他那個行徑的一個讚賞.
而且更加原來其他的民安的對象.
全部都是直呼其名.
按他們的名字來稱呼.
有些是猶太的名字 有些是希臘的名字.
就是說反映著有些真是他們是猶太人.
有些當然在羅馬已經有猶太人和外邦人.
來到一起聚會的那個生態.
所以在這個民安裡面.
一方面就是對他們的所作所為來稱讚.
就是說他們的生活.
他們一些可能沒有什麼人留意的地方.
他們保羅都留意到 就稱讚他們.
又或者可能有一些人已經是透過不同的人.
告訴保羅 而保羅知道又認識.
對他們所做的一些生活的行徑有那個稱讚.
當然既然民安一定是保羅所認識的朋友.

$^{121}$我們又可以透過稍後逐個民安的人.
我們就知道那個熟絡的程度.
還有在什麼時候開始去熟絡.
都是很值得我們仔細去看這20多個人的.
更加就是不同階層 不同的身份 不同的背景.
這個就真是要帶到很重要的真理了.
就是告訴我們 喝口水吧 喉嚨還是有點乾.
和大家先揭一揭《加拿大大書》.
三章的《加拿大大書》.
三章的26節開始.
差不多就是這本書的那個中心思想的教義了.
大體來說 《加拿大大書》要針對.
要強烈的反駁的神學的錯誤.
或者一些的讀者 加拿大教會.
他們被人來到去有一個恐嚇.
就是說你信了主耶穌.
如果不是好像我們猶太信徒那樣.
又要受割禮 又要根據摩西的律法來生活.
那麼你們的那個恩信清義.
就不足以使你們好像我們猶太的信徒那樣.
蒙神的那個喜悅.
就是說有另外一個關他們還沒過.
他們就是信入門.
但是需要用他們的那個行為.
這個是保羅的理解.
你做任何的事情在救恩以外要加添的.
這樣的點綴就是一個行為.
如果你說要信再加行為.
那麼保羅的眼中這個就是異端的了.
這個就是會與基督分離的了.
所以整個加拿大書的措辭.
是相當強烈的來就教義去駁斥.
所以是一個護教的書卷.
所以三章這裡要留意了.
第26節開始可以讀長一點的.
所以你們因信基督耶穌都是神的兒子.
你們受洗歸入基督的 或者受浸.
到時披戴基督了 並不分猶太人.
希利尼人 自主的 圍爐的 或男或女.
就是說已經是沒有東西需要再討論的了.

$^{161}$因為保羅已經涵蓋了每一個我們想到的.
在人類的社會那個無論身份,地位,民族.
那個特性 只有男女而已.
不會有第三性的 你就可以想象得到.
這裡告訴我們是一個很徹徹底底的.
來將在耶穌裡面 這個就是唯一的.
不可以叫條件 唯一的分辨.
就是在主耶穌裡面就有那個永遠的確據.
就是都是神的兒女 都是神的兒子.
這樣的概念對猶太人來說.
他們就會有很多的附加上去的.
例如你是不是亞伯拉罕的子孫.
例如你有沒有跟我們猶太教這樣來受割禮.
有沒有跟我們猶太教要進行一些的條規.
這樣疊加上去 這樣加諸在信耶穌的.
無論是猶太人或者是外邦人.
這些都是違反了恩順清二的唯一的信心的.
所以保羅這裡就很清楚的界定.
只要在耶穌裡面 無論過去是什麼分歧.
什麼不同的背景 或者我們的階層.
有多大的懸殊都好.
在主耶穌裡面 立即成為一.
所以這個信心的入門.
就帶給我們一個整個生命的地位的更新的改變.
所以這裡保羅是沒有任何的轉彎餘地的.
他說在基督耶穌裡面都成為一了.
留意29節.
你們的屬父基督就是亞伯拉罕的後裔.
是照著應許承受產業.
好像要將那些猶太教.
一些對舊約有一種錯誤理解.
表面上是沒有錯誤的.
但是當耶穌來了這個世界.
成就舊恩之後 那個理解就不同了.
因為真正的那個贖罪祭.
那個永恆的舊恩就已經完成了.
主耶穌在十誡上完成舊恩.
整個舊約就要更新了.
整個人類的社會只是在耶穌的舊恩裡面.
或者你還沒有在耶穌的舊屬裡面的人.

$^{201}$就只有這兩種分別.
所以在耶穌裡面就算你是外邦人.
所以保羅你想一下多大膽.
在基督耶穌裡面的就是亞伯拉罕的子孫.
而且憑著應許將來可以得到那個應許承受的產業.
這些本來在舊約的時候.
是猶太人獨有的那個福氣.
但是現在在主的舊恩裡面.
人類只有一個的方法.
就是透過耶穌就可以成為神的兒女.
神的兒子 以及將來有永恆的天家.
以及有那個應許承受的產業.
這個就是加拉太書那裡所說的.
所以他用一些實行的原則來反駁那些錯誤的教導.
就是說這個很新的概念.
對保羅來說也是很新的.
哇 明明那個猶太人和外邦人真的是格格不入的.
但是現在突然間這個神學就好像加進來.
我們一些不同的當時的教會,當時的聚集.
說一定要記住這個重要的真理的宣告.
28節說 不再分猶太人,希利利人,自主的.
自主就是那些可能是貴族 可能是一些很有權勢的人.
因為和圍爐和奴隸有一個分別.
或者是男的 或者是女的.
因為男女在那個時代是沒有那個社會的平等的.
有高低之分的.
你可以想像得到.
因為你們在基督耶穌裡面.
立即的成為一個在神眼中的.
只是一個的群體 一個的新的生命.
就是神的兒女.
就是不再分猶太人的神的兒女.
外邦人的神的兒女.
總之不再理會那個民族的了.
都是神的兒女 在主耶穌的救恩.
我們一洗進入這個洗禮裡面.
我們就是神的兒女了.
所以這裡是說要確立在地位上的平等.
在一個好像法定的地位上.
但是這樣的概念只是神學的宣告.

$^{241}$所以我們就要留意.
在保羅的書信裡面.
所以你看到這裡很重要.
保羅 這裡是我印的聖經.
這裡是三章的26-28節.
但是在羅馬書就很特別了.
因為羅馬書在保羅的書信裡面.
已經一直這樣說了.
告訴我們這個神學的基礎.
不再重複這些宣告真理的口號.
而是保羅怎樣示範給我們知道.
在主的救恩裡面.
我們都能夠在實際的生活裡面.
我們正在發生這些變化的.
而是我們怎樣能夠讓那個真理.
已經在我們的生活的層面.
來流露出我們在主裡面的合一呢.
所以你看看第16章的文案.
保羅對每一個有女性的尊重,讚賞.
有一個奴隸 那個奴夫據聞就是奴隸.
來到他這樣對他們的稱呼,稱頌.
對他們的讚賞 你就可以想象到.
加泰書說的.
是告訴每一個錯誤理解真理的猶太人.
或者外邦人一樣.
在救恩裡面.
恩順就進入福氣的合一.
神國度的合一裡面.
所以加泰書告訴我們.
我們是應該要合一.
但是羅馬書這一段的經文.
一到第16節告訴我們.
保羅示範給我們看.
其實已經在這個合一裡面相交的了.
所以保羅就是這樣對每一個不同身份的.
羅馬聚會的朋友.
對他們的感恩,問安,對他們的讚賞.
所以羅馬書更加傾向在那個實際的生活裡面.
那個相交裡面.
大家都已經以弟兄姊妹,神兒女這個身份.

$^{281}$來去有這樣的相交.
就見證那個合一.
是怎樣來去示範給.
所以這個示範就是一個好像示範單位那樣.
示範給人看.
所以外人 未信的人.
對教會的相愛,相交.
都有他們一定的期望的.
有時他們為什麼會憤外的.
有時那種噁心.
就是有時見到教會裡面如果有不和.
見到教會裡面有一些事情.
對一些外人來說.
是會覺得是一種.
你不是教會應該是講愛,講合一.
你說不信的群體,講利益.
講自己的那個放大自己的需要.
來去搶奪.
但是教會不是應該學校耶穌的博愛.
犧牲,捨己,奉獻的精神嗎?.
所以的確是的.
的確聖經就是這樣告訴我們.
所以《加拉太書》是先告訴我們那個原則.
為什麼要合一的原因.
原來我們本來就不是相同的.
大家是猶太人就是猶太人.
外邦人就是外邦人.
你拜開的神和我們是沒有關聯的.
我們拜開的偶像.
我們相信的那些大自然或者無神.
都是我們的事來的.
但是不是的.
造物主差遣他的兒子來這個世界.
要將人類帶進那個救恩裡面.
所以那個合一的那個事實.
因為耶穌已經發生了.
但是發生了.
怎麼能夠令到每一個信的人都要來到去.
學習這個合一,學習這個相交,同心.
去傳福音給未信的人呢?.

$^{321}$這就是羅馬書很寶貴的.
透過這樣的文案就讓我們看到.
原來初期教會就是認定這個真理.
這個在基督裡面的合一這個真理.
保羅就自己去示範.
不是刻意的只是為了文案而文案.
而是在實際的生活裡面.
保羅已經有一個新的價值觀了.
至少他不會對外邦人有歧視,有反感,有憎恨.
他已經看外邦人是他的弟兄姊妹.
他看一些女性是非常的尊重.
我們在神學院應該是二年級的時候做了一份功課.
女性可不可以講道呢?.
《提摩太前書》第二章.
還有應該第六章也有的.
那些人覺得保羅很專環.
應該是第二章的.
我們可以看多一點.
第二章的第八節.
「我願男人無憤怒,無爭論」.
這個《提摩太前書》的二章.
「舉起聖潔的手 隨處禱告」.
「又願女人廉恥,自首 以正派衣裳為裝飾」.
「不以編髮,黃金,珍珠和貴價的衣裳為裝飾」.
「只要有善行 這才與自稱是敬實的女人相宜」.
「女人要沉靜學道 一味的信服」.
「我不許女人講道 也不許她嚇管男人」.
「只要沉靜 因為先做的是阿當 後做的是夏娃」.
「且不是阿當被引誘 乃是女人被引誘 陷在罪裡」.
「然而 女人若常存信心,愛心 又聖潔自首」.
「就必在生產上得救」.
講到生兒育女 懷胎十月.
講到女人在聚會的時候.
提了兩次要沉靜 要學道理.
不許她講道 講道背後就是嚇管男人.
所以是要沉靜 接著講了那個原因.
要沉靜的原因 因為先做的是阿當.
就是說那個是創造的秩序 次序.
後做的是夏娃 接著就要追究.
當初吃的那個分別是鱷,薯,果的是什麼呢?.

$^{361}$而且不是阿當被引誘 是女人被引誘.
接著就陷在罪惡裡 慫恿阿當來受引誘.
所以接著就說15節就是「然而」.
這段經文是挺爭論性的 挺多不同的解釋的.
差不多是有些大的論文都是寫這個題目的.
對華人教會來說 如果我們沒有經過去詮釋.
我自己第一次看 我都認為是這樣的.
女性是不准講道的.
所以有些教會真的是女性是不准講道的.
應該是福音堂的教會.
還是女性都是不可以在主要的崇拜聚會講道的.
所以大家有時候要記得這些傳統.
但是這裡就要問一個問題了.
其實當時是不是因為提摩太前書.
在爾弗所這裡有些很特殊的處境.
而導致保羅好像帶有很強烈的對女性的偏見.
而導致變成一種對後世的教會一個錯誤的示範.
到底整體的保羅書信對女性是怎麼樣的呢?.
所以這個功課原來挺巧妙的 原來要做這份功課的人.
就要對整個保羅書信都要有一個整體的認識.
就是說保羅一貫對女性是怎麼樣的呢?.
在他的行文裡面有沒有這樣的教導.
女性是要被貶低的 女性是不應該做任何事情的.
或者諸如此類的 如果你有這樣的結論.
而再用這個經文去再印證 那就很合理.
這樣的做實學裡面應該是這樣的.
但是如果在整個的保羅書信裡面.
原來一貫都不是對女性有這樣的貶低.
也可以這樣間接印證了 剛才那個假設是對的.
就是提摩太前書是有一個很獨特的處境.
而那個處境日於言表.
保羅要在智女行家教提摩太的時候.
就要告訴他在爾忽所這個地方.
這個時候 有一些可能個別的女性.
特別的出頭 出到一個地步.
是要壓制那些弟兄.
甚至已經為外面其他的教會.
其他地方的教會的人知道.
而導致在那個次序上有一種的錯誤.
有一種的混亂 所以提摩太前書這裡.

$^{401}$傾向就是一個很處境性.
所以現在這個處境性是一個特別的原因.
所以就不是普遍的原因.
不是一個普遍的真理的教導 兩者是不同的.
等於使徒時代禁止吃血.
信了耶穌 外邦人都不准吃血.
因為原因很簡單.
因為當時是有很多的猶太的信徒.
對猶太的信徒來說 看著你吃血.
當然不是喝那些血 是變成豬紅.
變成那些牛的牛是怎樣稱呼的.
有鴨血 牛血沒有的.
豬紅就是豬血 雞紅.
好像我沒聽過牛紅是怎樣稱呼的.
諸如此類 就是說有那個處境.
有那個良心上面.
不要令到猶太人有一種很大的反感.
那就影響了大家的相交.
所以使徒就規定了.
要禁戒的是祭偶像之物和血.
和勒死的牲畜和奸人等等的.
所以保羅告訴我們.
這些處境性隨著那個時代已經不同.
其實就不需要死守了.
所以讀聖經有時候就要留意這些處境.
所以這段經文就告訴我們.
是有它的處境性 我們就不要看得那麼大.
相反 正正就是因為我們讀羅馬書的時候.
正正就讓我們看到.
保羅我們有足夠的經文去印證.
保羅原來對女性的尊重是非常大的.
就是說原來在保羅信了主.
成為神的僕人使徒之後.
他自己因為我不知道聖靈怎樣更新了保羅.
他成為了一個看事物.
差不多貼近真理的那個研判.
就是說他的那個分析.
他的那個神學已經是貼近神的那個心意了.
所以我們看保羅的書信.
差不多他講的一些絕對的真理.

$^{441}$是真的絕對一定會兌現的.
但是處境的就不一定需要了.
比如他說我都希望大家都像他一樣獨生.
不是每個人都可以的 保羅自己也這麼說的.
所以我們要留意什麼是保羅他自己.
已經是因為神給他的啟示.
聖靈在他裡面的工作.
以他來說的這些神學這些真理.
那就是千真萬確的確信了.
但是如果有些處境性的.
我們就要就著那個背景來做一個新的詮釋.
那就會合理一些.
比如我們開始看 菲比我們看了.
我們不會再看這個女性.
保羅是接受這個菲比這個女執事.
對他的支持的.
她是菲比的那個受助者來的.
就是說他接受菲比.
或者可能他的教會或者他的家庭.
對保羅一直傳道的經濟上的支持.
就是說菲比那個身份是相當的尊貴的.
而且保羅這麼有學問,這麼有見識.
還有這麼有神僕人的身份.
都不會介意在這個女性菲比的門下.
因為當時如果你是那個受助人.
你是要對那個資助者,那個贊助人.
是要給他榮譽的.
當然在羅馬帝國的世俗裡面.
這些就是關係的建立.
就是說我給你好處.
將來你有什麼需要.
你就要站在我這一邊.
比如我們能夠明白的就是.
假如如果做議員的.
如果我也得到你很多的幫助.
將來你做議員去選舉.
那我一定會投你 我沒理由投別人.
如果要比較的時候.
因為你真的幫助過我.
那不等於賄選的.

$^{481}$因為我真的覺得你好 那我就投你.
除非有直接的錢銀的利益.
如果不是 這些就是很微妙的.
平衡一個社會裡面的結構.
其實全世界是一樣的 都是會這樣的.
所以很自然的 菲比就是保羅的贊助人.
但是保羅作為使徒.
他不會覺得要女人的錢.
要給她去支持 就覺得很收價.
或者我不應該在她下面.
將來她有什麼好處 我又要給她.
保羅完全沒有這樣想的.
他很認為當時的這個贊助人的制度.
他都沒有反對的.
他都覺得是很合理的.
接著第二對就是這個伯基拉夫婦.
我們可以用一點時間來看這一對夫婦.
第16章的第三節.
問伯基拉和阿居拉安.
記得伯基拉才是姐妹.
阿居拉才是丈夫.
已經不止一次這樣調轉次序的.
先問了太太 再問丈夫.
很特別的 我們比較少這樣的.
他們在基督耶穌裡與我同宮.
也為我的名將自己的頸項致之道愛.
好幾個描述 你就可以想象得到.
這對夫婦對保羅的影響.
還有他們的名聲 他們的服侍.
是多麼得到保羅的感動.
還有保羅對他們的敬重.
我們留意了 伯基拉 應該這樣說.
阿居拉夫婦 他們的足跡.
因為在哥林多認識保羅.
他們本來是住在羅馬的.
也就是說 羅馬的教會是認識阿居拉夫婦.
但是輾轉 因為曾經在早年格魯丁的時候.
有一個羅馬的政府要將所有的猶太人.
從意大利趕走 回去不同的地方.
所以阿居拉夫婦就去了哥林多.

$^{521}$在哥林多就遇見保羅 彼此認識.
平日大家就一起去織帳篷.
猶太人用那些羊毛 山羊毛來織的.
可能是毛毯 或者那些帳篷外面的遮蓋等等.
是一些很專業的手藝.
當然到了安息日 就一起去會堂傳道等等.
一起去聚會 有家庭的教會.
後來輾轉 阿居拉夫婦就去了以弗所.
後來當然也在羅馬服侍.
所以你可以想像得到.
阿居拉夫婦一直都是家庭教會的領袖.
她招呼不同的人.
我們上個星期也提過一點.
這個阿波羅 一個很有恩賜的年輕人.
很有感動要去傳講主耶穌.
但是聖經路加醫生說.
他只是懂得約翰的洗禮.
就是暗示他不知道主耶穌的洗禮.
主耶穌的救恩.
人在主的裡面才真正得到永恆的潔淨.
所以約翰的洗禮.
只不過是叫人預備那個心.
預備一個悔改的心去接受將來的救恩.
所以如果他一直都懂得耶穌的洗禮.
那個十誡犧牲.
那他所傳的只不過是靠自己的好行為.
還沒有得救的一個.
如果再這樣的話就變成異端了.
所以阿居拉夫婦有一次聽到阿波羅講道.
很熱心 很有恩賜 很能夠陳述,講解.
但是他們聽到有些問題.
晚上就接待了阿波羅.
跟著跟他坐定了 跟著再分享.
就是說當時可能真的沒有手抄本那麼多的流通.
就跟阿波羅繼續去講解陳明基督的十誡救恩.
或者說主耶穌為我們受的這個苦楚的洗禮.
使得我們可以得到這個救贖的恩典.
所以阿波羅聽完這對夫婦的教導分析之後.
跟著他的講解更加的細緻.
更加能夠反駁那些猶太人.

$^{561}$跟著引聖經去證明耶穌基督的福音.
可以這麼說 阿居拉夫婦就得了一個年輕人.
阿波羅成為一個將來被大大使用的報導家.
一個周圍去傳道的牧者.
所以阿居拉夫婦就是一直.
她未必是周圍很廣泛的去傳道.
做一些報導的工作 但是默默地在家庭裡面.
作為教會的領袖來去接待.
這個就是阿居拉夫婦 阿波羅怎麼稱讚她呢?.
對他們就是為了阿波羅的性命.
將自己的頸項置之度外.
就是說阿波羅受苦是很合理的 因為他是使徒.
但是你看到阿波羅對這個人的心智.
那個對主耶穌的救恩的服侍 生命都放上了.
你很欣賞這個阿波羅.
但是也需要你們要置之度外 你們的頸項.
而且是夫婦二人同心 是很不容易的.
兩夫婦一起對阿波羅的事奉有一個這樣的支持.
支持到一個地步就是阿波羅受苦 阿波羅受逼迫.
他們都願意跟隨 如果神真的要他們受苦.
他們願意跟隨 這些是很高的生命的情操.
因為面對最大的就是困難 是苦楚.
如果困難苦楚都不會將你的心智去消磨.
那就證明那個跟隨主 服侍的心.
就非常之真確 非常之真實.
這個就是阿基拉夫婦 阿波羅這樣對他們稱讚.
就是說有些記憶 有些事件是真的發生過.
而他們真是願意站在保羅的身邊為福音作見證.
這個就是阿基拉夫婦.
置之度外的這樣擺上自己的頸項.
所以如果我們能夠在教會裡面.
對主 對主的僕人 對主的福音.
我們擺上 我們用生命去捍衛.
甚至神給你有不同的足跡.
在不同的地方來服侍.
你流露出來的見證是讓宗教會都欣賞的.
那就會帶來一個很大的福氣.
因為最近我們都是跟不同的群體.
特別是因為我們這段時間 過年前後.
其實過年之前 聖誕都已經開始接觸.

$^{601}$現在反而少接觸洪福村.
反而落橋樓 成為洪水橋現在比較多群體去接觸.
我也多主動去聯絡接觸.
所以過了一年之後.
有些我不認識他們的人 機構.
有時候打電話來找教會找我.
要談合作 其實我不認識的.
有些人有時候聊天 跟不同的教會聊天.
有時候一起事奉的時候.
好像上個星期六我們去宣明會.
就有一些聚餐給那些落橋樓的火災災民.
有30位 我們弟兄姊妹有4台做義工.
有時候聊天 其他的傳道人其實我也不認識的.
他們有聽過我的名字 聽過洪水橋.
就是洪恩堂曾牧師的名字.
變成好像慢慢有時候你做的事情.
其他教會都認識 都知道.
所以真的 如果我們能夠一心.
在我們人生中心去服侍.
保羅的這句話都挺值得我們去追求的.
當然不是為被人誇而誇.
而是你一心是為主去忠心服侍.
保羅說 不僅僅是我稱讚他們.
就算是外邦的宗教會.
保羅不是要誇大 不是要拍馬屁.
保羅真是真心對這對夫婦的人生.
有一種這樣的欣賞.
連保羅就是一個這麼偉大的使徒.
他看到這對默默的夫婦.
這樣擺上自己的金錢 擺上自己的家庭.
擺上自己的生命 來去服侍主.
所以 如果保羅這句話留下在聖經.
阿居拉夫婦成為.
傳頌於千古的一對神的兒女.
你可以想象到他們離開世界.
短短幾十年 他們已經在天家二千年了.
他們在主耶穌的眼裡面.
他們終身的 走完人生的路的服侍.
將來得到的讚賞和稱讚.
神對這對夫婦的稱讚.

$^{641}$你已經可以感受到那種福氣 那種榮耀.
所以 不要小看使徒對這些平信徒的稱讚.
因為 能夠在羅馬帝國的壓迫底下.
他們也是猶太人 經歷過被政府.
這個格魯雕來去驅逐.
就是說 那些物業都被人搶了.
不會先賣了樓才走的 沒有的了.
說明要排走 就是趕著你走的了.
你是沒有任何東西 你就要馬上離開的了.
不走就有生命危險的了.
等於在印尼曾經排華的時候.
對不起 讓我先賣掉這些田地吧.
完全是被搶奪的了.
你有了一條命 你已經算是這樣的了.
他們就是這樣走了 但是過了一段長時間.
他們的服侍就在保羅的眼裡面.
現在寫羅馬書的時候.
就是說對這對夫婦的服侍有這樣的欣賞.
有這樣的讚美.
所以你想一下 如果人都可以有這樣的稱讚.
何況做物主 對於阿基拉夫婦的讚美.
那個欣賞 還有將來得到的榮耀.
我們再看看下面.
以拜彌陀 這個是希臘文來的.
那個希臘文已經是什麼了.
第四節 啊對不起 第五節.
問我所親愛的以拜彌陀安.
他在亞細亞時歸基督初結的果子.
亞細亞和阿該亞是不同的.
亞細亞這個Asia Minor.
就是凡是土耳其以西的地方.
例如有以弗索 有哥羅西等等.
那幾個的教會 亞細亞.
那怎麼稱呼它呢.
這裡告訴我們它是亞細亞初結的果子.
是歸基督初結的果子.
原來我們看一下保羅很喜歡講果子的.
哥林多前書的十六章 我們又看一看.
當然我們沒有看十五章那裡.
用那個初熟的果子來做一個比喻.

$^{681}$講到我們將來必然會像主耶穌那樣去復活.
主耶穌就是初熟的果子.
講到祂是復活的初熟果子.
我們將來必然會像耶穌那樣復活的.
所以不是一種精神 不是一種幻覺的復活.
我們真的會肉身復活去到天堂的.
不過這裡再一次提醒.
所以保羅講的這些初熟果子.
都是含義著那些初信主的人.
第十六章的十五節 或者十三節.
留意這裡 十六章的哥林多前書.
你們務要警醒 在真道上站立不穩.
要作大丈夫 要剛強.
凡你們所作的 要憑愛心而作.
弟兄們 你們曉得斯提反一家.
是阿該亞初結的果子.
並且他們專以服侍聖徒為念.
這樣的稱呼 一節聖經 已經有很深度的含義.
他們是阿該亞初結的果子 斯提反.
或者使徒斯提法 當然是不同的翻譯.
這裡就告訴我們 阿該亞.
最主要的地方就是哥林多.
所以保羅寫這封信的時候.
再一次提醒弟兄姊妹.
斯提反這個人 當然不是殉道的斯提反.
這個斯提反就是在阿該亞.
保羅初到哥林多的時候.
某程度就是保羅帶的第一個信主的人.
是阿該亞初結的果子.
就是說保羅在哥林多的時候.
其實很多人反對的.
有些人甚至攻擊保羅的.
保羅打算走了 猶太人逼迫他.
但是在晚上 保羅見到異象.
耶穌向他顯現 告訴他 你不需要怕.
只管講 沒人會害你.
因為在這個城 哥林多.
有很多的人 他們是被信主的.
保羅就因為這樣心動.
他就繼續在哥林多留下.

$^{721}$接著有一年半的時間.
所以斯提反應該是很早期.
聽保羅的教導信主 而來到生命改變.
就是阿該亞的初結果子.
這裡告訴我們 斯提反有一樣東西.
他信了主之後就去做的.
並且他們 這裡是斯提反一家話.
全家人都是這樣的事奉的心智.
就是以服侍聖徒為念.
專以服侍聖徒為念.
就是說是非常有愛心的去服侍其他的弟兄姊妹.
這個就是斯提反一家的人.
他們全家是阿該亞初結的果子.
初結果子對猶太人來說一點都不陌生.
就是頭生的和初熟的樹上的果子.
地上的那個收成 那些穀等等.
就要歸神為聖的.
所以舊約就有這樣的教導.
所以保羅都對這些舊約的概念.
其實他不是祭司 他不能夠獻祭.
但是他看這些信了主 由他帶順主的人.
初頭信主的 一直到今天.
保羅都還是記掛他們的生命的成長.
和參與在教會的事奉.
在原浸的弟兄很熟的了.
熟到曾經我讀神學的時候.
他是將他家完全借給我住 住了四年.
沒有簽任何東西 憑一個信字.
他們自己搬去天水圍租地方住.
讓一間房子給我住 給我讀神學.
因為我要賣掉那間房子 讀神學.
四年裡面都沒有任何收入.
他就賣了 我沒得住.
就給我一間房子住 在元朗.
然後他們自己在天水圍租了四年房子.
現在他在原浸已經很長的事奉了.
偶然間有邀請過來洪恩講教主一學.
有時候又會聊聊天 他現在住在Yoho.
Yoho Midtown 初結果子.
是我看著他信主 然後結婚等等.

$^{761}$都算我在元朗的初結果子.
一位很好的弟兄 他很認真預備主一學.
很感恩 跟他家裡 跟我們的關係都很好.
這位弟兄現在在元朗 他曾經試過做理事.
現在他專注訓練一些年輕人.
做一些門徒的訓練.
大家好 一見如故 因為真是我初結的果子.
你可以想像得到 保羅提了兩次.
斯提反一家是阿該亞初結的果子.
但是這個以拜尼圖.
是亞細亞歸於基督的初結果子.
是不是都是保羅在二分所的時候.
帶信主的以拜尼圖呢?.
雖然不是很具體.
但是可能都是隱約暗示在亞細亞的時候.
保羅帶他信主 現在他不再在亞細亞了.
現在在羅馬這樣去生活,聚會和事奉.
雖然沒有提到有什麼對他的稱讚.
是保羅所親愛的這位人士.
這位弟兄 以拜尼圖 是一位男士.
他所親愛的以拜尼圖.
所以弟兄姊妹 如果你初信 或者你有份事奉.
Alex更加明顯 看著小朋友成長.
他們現在繼續參與在初中,高中.
甚至出來工作 又在年輕的團體幫忙.
你看著他們一路成長 一路事奉.
就很不一樣了 這些就是歸實的初結果子.
所以我們要有的 多少都要有的.
如果我們都沒有.
一來你都沒有參與過在別人的生命的建造.
二來就是你自己都若即若離.
沒有在福音工作上擺上 那你自然就沒有了.
所以人生建築 你沒有初結果子.
你沒有結果子 其實這不是一件恭喜的事.
關我什麼事 我都不是全道去做 全道去結.
我都不是全道 何須去結初結果子呢.
不可以這樣說 因為你的生命都會成長,成熟.
你的生命都會影響著一些人.
還有去到一個年紀 其實我們要將我們的經驗去傳承.
最重要的存在 就是要將你的經驗傳遞給其他人.

$^{801}$使他們能夠對主有一個真誠的皈依.
所以我們應該要爭著去做兒童主義學老師.
爭著去做一些我們可以幫助人建立生命的工作.
因為保羅自己都 他不是要沾沾自喜.
其實真的不需要沾沾自喜.
因為根本就是神的工作.
但是如果你看到一個人真的在靈命上面.
你看著他從信主一直這樣成長.
那種喜樂,那種滿足.
那種你曾經在別人的生命裡面.
留過一點點滴的功夫等等.
甚至最好就是有禱告,有眼淚這樣的灌溉.
那個靈命,那個人的果子.
將來就歸於神 歸於你的帳下.
當然我們不是為了勞攏人心來去結果子.
我們是希望將每一個果子歸實的.
因為一切都是實的.
所以兩者之間其實很清楚大家都懂得分的.
我們是要操縱一個人的生命.
他一定要跟我 他成為一個行屍走肉.
因為他有獨立的生命,獨立的恩賜.
獨立的見解 獨立的甚至他長大之後.
已經跟我某些的真理不同了.
某些的見解不同了 那就好了.
因為他已經成長了.
永遠不會停留在我對他的那個教導底下.
超脫不了 所以這兩段我很欣賞的.
保羅說他有亞細亞歸於基督的這個以拜尼圖.
又有這個斯提反的一家人.
是亞該亞初結的果子.
所以弟兄姊妹奉勸大家.
真的不要只是個人的返教會.
我們真的要參與在神的國度裡面的服事.
神會給我們有感動.
神會給我們有這樣的內心的引導來做的.
接著我們就繼續下去了 看看我們夠不夠時間.
接著提到瑪利亞 聖經裡面有六個瑪利亞.
第六節 又問瑪利亞安.
他為你們多受勞苦 只是一句的稱讚.
他為你們多受勞苦.

$^{841}$這個瑪利亞到底是誰呢?.
我們就不是太知道了 因為是保羅所認識.
在這個時候 我們可以透過少許的研讀.
聖經有六個瑪利亞.
我就可能和大家稍微溫習一下.
可能有些大家都不認識的.
我們就當是反正是查經 我們一起去看看.
相信大家不會不認識的了.
就是耶穌的媽媽瑪利亞 就是約瑟的太太.
另外一個叫做莫大拉的瑪利亞.
這個姐妹的出名之處.
就是她未信主之前是被鬼附的.
是住在一個外邦人的地方 馬大拉那個地方.
來到近加利利 可能是加利利海的對岸.
就是說她以往是交鬼也說不定.
或者如果不是交鬼 為什麼會有鬼附.
我們就不知道那個背後的原因.
我們一知道 就是耶穌基督就拯救了她.
幫助她 然後她就成為了一直在耶穌身邊.
跟隨耶穌的一位姐妹.
最後一個階段就是.
她是第一個看到耶穌復活的姐妹.
所以這個莫大拉的瑪利亞.
其實是非常有福氣的.
然後就是有一個叫伯大尼的瑪利亞.
大家可能就這樣聽 可能不太認識.
伯大尼的瑪利亞是誰呢?.
就是拉撒路的姐姐 馬大的妹妹.
就是這個叫伯大尼的瑪利亞.
然後還有兩個 這些可能就更加少了.
因為記載了一段經文而已.
就是一個格羅巴的妻子 瑪利亞.
這個是一個王族的了.
另外就是一個小亞國和約西的媽媽 瑪利亞.
就是這五個人 當然就是第六個.
就是羅馬書 保羅稱讚的這個.
保羅就沒有對這五個的姐妹.
這五個的瑪利亞.
在聖經裡面有書信 或者對話交往.
因為保羅 反而路加醫生.

$^{881}$就反而有機會可能接觸到耶穌的媽媽.
但是反過來 保羅就很少提到.
有提到的就是這個瑪利亞.
這位瑪利亞就是一位多壽奴婦的姐妹.
為了你們多壽奴婦.
到底她是什麼樣的背景,身份.
保羅就沒有詳細教導.
或者說初期教會的羅馬教會.
他們根本就很清楚,很認識這位姐妹.
瑪利亞這個名字是相當普遍的.
一個猶太的名字 相當普遍.
但是大聖經就給我們看到.
在耶穌的年代 有五個的瑪利亞.
和耶穌是有相當的關係的.
反而我們這些大家都應該熟悉的經文.
我們就不太詳細去看.
反而我們看看《使徒函傳》的第一章.
我們翻開前一點的《使徒函傳》第一章.
就提到關乎到這個聚集.
由12節開始 因為耶穌已經升天了.
加利利人就是這一幫的門徒就回到了第12節.
有一座山 名叫「鑒欖山」.
離耶路撒冷不遠 若有安息日可走的路程.
大概就好像傳統大概2000蚤的距離.
安息日的路程 就是說你太長走遠路.
比如假設今天你由錦繡去元朗.
就是安息日的路程才算.
但如果你說不是的 安息日我要在錦繡出去香港.
過了安息日的路程是不行的.
你是犯了安息日的.
所以有時候《使徒函傳》也有這些描述.
因為他們是猶太人.
他們還是在猶太教的那個分途底下.
所以當下他們還是有這些背景的傳統的.
但當然慢慢一定要盡快來分清楚.
在主耶穌成就救恩之後.
這些的禮儀 這些的律法 這些的傳統.
還是不是要遵守呢?遵守到什麼地步呢?.
這些他們一定要在他們有生之年往後傳道.
必須要梳理的 因為很多猶太人信了主.

$^{921}$就會陷在這種困難的兩難裡面.
等於我們有時候也要教導弟兄姊妹.
我們不可以再拜山了.
我們可以帶竹花去紀念.
但是我們不可以帶有一次的安息禮拜.
那就可能沒有講得那麼清楚.
那些家屬在我還沒來到之前.
哇 把那些東西都放在那裡.
他們不是放元寶蠟燭.
是放那些親人生前很喜歡吃的那些水果.
那些東西都放在靈堂裡面.
我說不行啊 我們基督教是不用的.
一支蠟燭紀念就可以了 那些是吃不了的.
他們才會撤走的 可不可以不撤走呢?.
我說不好啊 怕那些不信仰的人誤解了.
不是說那些放在那裡有什麼問題.
而是怕這個傳統 中國人的傳統.
被那些不信仰的人來到基督教的喪禮.
安息禮拜 哇 安息禮拜都要放一些食物在那裡.
你說放在後面給人吃是沒有問題的.
你放在靈堂裡面上了下面 那是不行的.
那就變成了一種錯誤.
我就說不行啊 一定要拿走它.
他們很尷尬的都拿走.
都尊重因為是安息禮拜.
所以我們要留意.
他們還是有安息日的手的禁忌的.
就不可以超過2000爪.
如果那個尺寸真的是這樣的距離.
當下盧家醫生就很仔細的描述.
當下門徒從那裡回耶路撒冷去.
就是說他們就是剛剛見完耶穌升天之後.
就回來了 回到耶路撒冷的樓房.
第十三節 回到耶路撒冷城.
就上了所住的一間樓房.
在那裡有彼得,約翰,雅各,安德烈,.
比利,巴多羅馬,馬太,阿納肥的兒子雅各,.
奮銳黨的西門,和雅各的兒子猶大.
這些人同著幾個婦人.
和耶穌的母親瑪利亞 並耶穌的弟兄.

$^{961}$都同心合意的行切禱告.
這個可以想像得到.
盧家醫生是非常仔細的記載每一個的使徒.
除了這個加勒人猶大 因為他已經背棄了耶穌.
後來就記載他自盡 吊頸死了.
但是這裡這樣去描述.
他們就住在這個樓房上面 在上面來禱告.
然後特別強調這裡同著幾個的婦人.
和耶穌的母親瑪利亞 就是同一些的姐妹.
可能估計有蒙特拉的瑪利亞.
有約西的母親 格羅巴的妻子瑪利亞.
所以好幾個的瑪利亞都在現場.
還有耶穌的弟兄 不是使徒的.
這個雅各書的作者雅各 猶大書的作者等等.
你就可以想象得到 這個這樣的聚集.
這裡說大概有120人 還有很多的門徒.
那些都已經是跟隨者.
仍然誓死效忠耶穌和使徒們.
所以這一次的聚集.
就是說一幫人完全是對復活.
是有一個感覺到很大的震撼的.
他們到耶穌死都仍然不離棄耶穌的.
其實很不簡單的 因為那個壓力很大的.
但是他們都願意繼續.
在這個樓房就有一個復興的開始.
或者說聖靈降臨的開始.
就是這120人來認識了復活的耶穌.
當然還是半信半疑.
但是應該開始他們會內心有很多的交戰.
很多的思前想後.
也都開始他們認定要開始為主耶穌作福音見證.
但是還沒有聖靈 一定沒有那個能力.
沒有那個膽量.
但是這裡我們想說的就是.
使徒告訴我們 馬尼亞耶穌的媽媽.
還在這個行列裡面.
一來就是看著自己的兒子這樣身首異處.
慘死在十架上面 死無全屍這樣.
有很多的傳言說屍體不見了.
但是最終耶穌基督在他們眾人的面前.

$^{1001}$升天 坐著昏 升天 回到天家.
這樣的一個警方.
耶穌的媽媽看著自己的兒子復活升天.
但是知道這一位是自己親的骨肉.
由自己的母腹生出來的.
看著他長大等等的.
但是現在看到的是這一位神的兒子.
就是神人的異性 這樣來升天.
對馬尼亞來說 百感交集的.
又是開心 自己的兒子升到這麼高.
以前是地底泥 被人這樣陷害.
被人這樣去誣告 受盡痛苦 受盡凌辱.
其實出生已經是未婚懷孕了.
一直都帶著很多不好的標籤.
這樣來長大的耶穌.
而且不是靠神蹟 靠突然間飛黃騰達.
是一直每分每秒都是痛苦的長大的耶穌.
所以媽媽墳外感受良多的.
但是看著這一位的兒子升天.
當然不是忠於真相大白.
而是她真的認識這一位自己生的兒子.
其實是神做物主的兒子.
這一位永恆的真神.
所以馬尼亞都挺複雜的.
但是記得 她只不過是120個其中一個的婦女.
所以如果我們真的好像天主教那樣.
將她聖化了 變成聖母馬尼亞.
那就超過了聖經的教導.
她只是其中一個的見證人.
她的見證當然是有血有肉.
因為真的是從她自己的肚腹生出來的耶穌.
但是仍然她是受尊敬.
但是不可以變成一種聖人.
我們去拜一個聖母.
聖母的神學就是 她也是有功德的.
她也可以隨時去到耶穌那裡.
所以我們向聖母祈禱.
聖母就會將我們的禱告轉化.
好像傳訊息出去 傳出去給耶穌.
其實是不需要的 聖經從來沒有這樣教導的.

$^{1041}$聖經叫我們對這個人有一個感恩.
但是從來不應該對任何一個地上的人.
包括保羅 都要將他變成一個聖人.
離開世界就離開世界.
他們不可以將他們的功德.
轉嫁回地上的凡人.
所以我們對馬尼亞有景仰 有感恩.
但是她是一個普通的女性.
所以我們是要直接向耶穌禱告.
千萬不要向聖母馬尼亞祈禱.
這個不是聖經的教導.
但是我們看到在這幾章裡面.
馬太奴家是將那個陳述.
馬尼亞怎樣能夠遇見神的使者.
給她有這個使命.
在千萬人當中的猶太婦女.
未婚的 揀選了她.
所以她是主揀選的一個器皿.
所以很尊貴.
也是神的那個召命臨到這個馬尼亞的身上.
她一定有她很獨特的地方.
連她的丈夫許配給了約瑟.
都是一個這麼公義的 敬虔的一個義人.
就是一直這樣愛著自己的太太.
所以馬尼亞後來都有生.
由約瑟和她兩個生的猶太 雅各.
路加夫林就有記載的.
我忘記了是第七章還是第八章.
就記載馬尼亞在耶穌以外.
有生到好像四個還是五個.
所以原來耶穌有自己的弟弟妹妹.
耶穌是大哥來的.
有時候想一下這些倫理的關係.
想一下我們自己小時候是怎樣的呢?.
耶穌是不是真的好像道貌岸然 又不笑.
硬邦邦的 只是講真理呢?.
還是跟弟弟妹妹都有聊 有講 一起成長.
大幾多年呢?聖經就沒有交代了.
跟著她跟爸爸一起做木工的時候.
又是怎樣的情景呢?.

$^{1081}$那些人對她的愛戴.
就是耶穌基督是一個很神秘的一個人.
但是我們今天所效忠所跟隨的就是這位救贖主.
所以耶穌的媽媽這位馬尼亞.
很值得我們去跟隨.
她就是這樣憑簡單的信心.
被人誤會她都願意生耶穌出來.
來跟約瑟一生盡他們夫婦的責任.
去保護耶穌來成長 成為將來的救贖主.
出來傳道等等 這個就是馬尼亞.
另外一個就是莫大拉的馬尼亞.
所以這個就是馬大拉城.
應該就是外邦的地方.
聖經有記載 在路加 我們看看路加八章這一段經文.
我們可能看完這個就要停一停了.
八章的開始 過了不多日.
耶穌周遊各城 各鄉傳道.
宣講神國的福音 和他同去有十二個門徒.
就是十二使徒 還有被惡鬼所附.
被疾病所累 已經治好的幾個婦女.
特別講其中一個 內宗有稱為「莫大拉的馬尼亞」.
就是說這個人住在莫大拉城的女人.
就叫馬尼亞 因為太普遍的名字了.
所以叫做「莫大拉的馬尼亞」.
曾有七個鬼從她身上趕出來.
接著就提到這些婦女 又有希律的家宰.
就是希律王的家裡御用的廚師 太太.
就是這個撫薩的太太「約阿娜」和「蘇薩娜」.
和「好些別的婦女」.
都是用自己的財物供給耶穌和門徒.
就是說主耶穌都是這樣憑信心去做這個事公傳道.
來跟這些門徒可能一班人住在一起.
去到不同的城 可能去一些客店.
或者有人資助她住 她就留宿.
但是有些起居的生活 都是由這些婦女.
用信心 用愛心去支持主耶穌的傳道.
都是在盧家醫生的記載裡面.
又是很鉅細無遺的記載這些婦女.
有些是有名 有些是尊貴有錢人.
希律的家宰 就是一個膳食的膳長的太太.

$^{1121}$應該都是住在皇宮裡面的.
你可以想像 原來有這樣的背景的人.
都信主 都愛主 跟隨耶穌 服侍耶穌.
支持耶穌的福音工作.
我曾經那個時候在讀神學.
因為有弟兄開車房 他認識一些不同階層的.
有一次有一個傳道人因為有事去不了那個茶經班.
就叫我代 需要有司機來接我.
就要上到半山翠亨村那些會所.
有些婦女 通常是那些什麼夫人那些婦女.
就在那裡茶經班 那個司機就載我.
接著就帶我進去 跟他們講完道.
接著就載我到加魯蓮山 銅鑼灣那裡.
就放下我這樣 很敬虔的.
她在車上就跟我說 她先生是很出名的.
心臟科的醫生 專家這樣.
就是很有錢 但是很敬虔的這些老人家 這些婦女.
所以希律的家宰 這個叫做約阿娜.
所以錢本身是中性的 我們怎麼去服侍主呢?.
這個就很講究我們內心是不是真的看福音這麼重要.
還有地上神給我們的一切的擁有.
都應該是屬於神的 我們用在神的裡面.
那又拆開講 我們這次在菲律賓的這個短訊觀察.
我們看到那裡現在我們基督教可以很大的幫助.
就是做一些助學的工作.
我們如果能夠改變這些小朋友離開那些垃圾村.
去到學校讀書 有基督教背景.
那就差不多已經將他們的命運改變了.
所以現在我們接下來的3月22日.
我們在紅欣就會有一個分享會.
將那些照片 將那些影片.
還有我們所見所聞跟《大紀元》分享.
接著我們就很想鼓勵《大紀元》.
一起去認領一些小朋友.
我們就去支持他們去讀書.
這樣其實就帶給我們很大的一個祝福.
所以可能我們現在就暫定8月底.
希望我們可以8月底就《野地無人》那個時間.
如果8月底的好處就是.
我很想帶一些家庭去 帶一些小朋友去.

$^{1161}$去真的看看那些垃圾村的那種困難,無助,痛苦.
如果能夠真的這樣看到的時候.
我想那個心靈的觸動是會很大.
所以我們預計8月底.
很想辦一次是這次我們紅欣堂帶隊.
會鼓勵《大紀元》節目去 一起去參加.
去了解當地的實際的需要.
你會有很大的感動的.
這個是瑪利亞 我們再看一看 我預備了一個經文.
《馬可》的 因為那個描述都很詳細.
第15章.
雖然和羅馬書這裡的文案是有一點點不相同.
但是既然我們揭開瑪利亞.
聖經都有6個瑪利亞.
我們稍微了解一下都好的.
剛才我們看了8章.
就是說瑪利亞到底她之前有什麼不好運.
她會被鬼附 是不是因為她是交鬼.
是不是因為她天生有什麼家庭背景.
那些聖經就好像沒有興趣交代.
但是由那一刻開始 她的鬼被趕走之後.
就真心實意用她的很單純的心靈.
一直跟隨耶穌 服侍耶穌.
到死都跟隨耶穌的莫大拉的瑪利亞.
所以我們將來會見到這麼多個瑪利亞.
雖然她過去原來她的背景是這麼坎坷.
是這麼黑暗 七隻鬼 鬼符.
但是她整個人生命改變之後.
就走一條光明的服侍主的路.
所以奉勸在座或者線上將來聽到的.
如果你是姊妹 聖經描述很多的姊妹.
她們那條心就是很單純.
一信了主就是愛主 服侍主 沒其它事.
其實這些是很重要的 似乎女性那個專一.
我服侍這麼多年都是見到有些姊妹很專一的.
對主的服侍 這些可能是女性特有的那種特性.
所以都是真的 聖經真的這樣描述給我們見到.
她的生命這麼大的改變 耶穌這樣將她拯救.
她就是這樣一直服侍.
就是整個生命只有耶穌在她的內心.

$^{1201}$鬼在她的內心來影響她 騷擾她.
但是現在耶穌在她的內心 她整個的生命就是這樣服侍.
其實是很輕香的生命 是很寶貴的人生.
就是心有所屬 對於一個女性來說 這個專一.
我相信我們男性不是太了解.
但是對一個女性來說 心有所愛 心有所屬.
心有所事奉 我想你們一定會明白的 比我還要明白.
所以真的大家按照著神給你的感動.
你在得到什麼耶穌的拯救.
你得到耶穌什麼的福氣恩典 你就要還這個恩典.
好像莫大拉的瑪利亞那樣去服侍耶穌.
我們記不記得《馬可》這一段十五章.
十五章的四十節開始 就是耶穌死的時候的那個景況.
她們根本不可能靠近耶穌的身邊.
因為那些男性 那些暴徒 在十字架下面的叫囂.
你是一個普通的婦女 根本就有那麼遠站那麼遠.
等於我們經常說笑 你想想如果我們去唐興拍照.
可能你根本是在最遠最遠的地方.
都不輪到你看到你的大樣子了.
這些不是說要階級 也不是說要歧視.
而是有時實際是這樣的 但是你會看到.
先不要小看 雖然當時那麼多人.
但是最終繼續留下的就是這些婦女.
聖經就很想告訴我們 真誠愛主 完全不會懼怕的.
最終在耶穌身邊的就是這些婦女.
內中有莫大拉的瑪利亞.
又有小雅各和約西的母親瑪利亞.
並有這個撒羅米 就是耶穌在加利利的時候.
跟隨祂 服侍祂的那些人.
就是說 剛才《路加福音》第八章說的這些婦女.
現在來到耶穌在斷氣之後 仍然在耶穌的身邊.
雖然是遠遠的 因為很多的暴徒在十字架下面.
那些兵兵在十字架下面 那些門徒其實都走了.
只剩下約翰和他媽媽在十字架下面.
這些我們叫做二打六的婦女.
其實就遠遠的在山後面等待.
其實她們不是二打六.
她們是真心愛主的 跟隨主的婦女.
可能她們根本沒有機會去講道.
根本沒有人認識她 但是在《路加醫生》的記載裡面.

$^{1241}$她在她的調查裡面 知道有這些婦女.
路加特意寄她們一筆.
這樣來給後世的人知道 就是這些人.
所以路加就很清楚 就是耶穌在加利利的時候.
傳道的時候 跟隨她 服侍她的那些人.
就是這幾個婦女 所以你可以想像得到.
活大拉的瑪麗亞是特別有記名 特別頭上有記名.
就好像不會落人於喉 她是身先士卒的.
來服侍耶穌 不斷感受到自己的生命.
得到潔淨 得到改變的改變.
這個就是活大拉的瑪麗亞.
為什麼我們要寫16章呢?.
過了安息日 活大拉的瑪麗亞.
和雅各的母親瑪麗亞 並殺羅米三個人.
買了香膏 要去告耶穌的身體.
七天的第一天清早 出太陽的時候.
他們三個人就來到墳墓 彼此說.
誰給我們把石頭從墓門滾開呢?.
那石頭原來很大 他們抬頭一看.
卻見石頭已經滾開了 他們進了墳墓.
看見一個少年人 就是這個天使 坐在右邊.
穿著了白袍 就神驚恐了.
那少年人對他們說 不要驚恐.
你們尋找那釘十字架的拿撒拉人耶穌.
他已經復活了 不再這類.
請看安放他的地方.
你們可以去告訴他的門徒 和彼得說.
他在你們耳先 往加利利去.
在那裡你們要見他.
正如他從前所告訴你們的.
他們就出來從墳墓那裡逃跑.
又發抖 又驚奇 什麼也不告訴人.
因為他們害怕.
這裡就告訴我們 穆達拉的瑪利亞.
和另外兩個姐妹.
就是最早去到耶穌的墳墓.
相信主耶穌會復活的這個消息.
而準備要來買香膏 要告末耶穌.
但是已經見到那個實跡發生了.
因為那個大的墓門 那個大石已經滾開了.

$^{1281}$裡面已經有一個天使坐在那裡等待.
其實挺驚恐的 雖然是天亮了.
可能剛剛天亮 清晨.
但是你看到那個墳墓裡面有人坐在那裡.
和你說話 其實無論是天使.
你都會嚇得半死.
所以你就可以想象得到 他們見到一個很大的實跡.
先讓這三個姐妹見到.
反而門徒還是在樓房上關了門 懼怕.
都不敢有任何的行動.
這個又好像反映到.
女性有時候那個感情大於理性.
不去分析 不去懼怕.
總之就是要自己 盡自己最後一點的心意.
去服侍自己的老師 救主耶穌.
你帶些香膏去 會不會受牽連?.
那些門徒就會有很多的思考.
很多的分析 很多的參詳.
但是那些婦女就不想那麼多了.
他們能夠做的 就是用錢買些香膏去告末耶穌.
希望那個屍體不要發臭.
因為已經第三天了.
所以這個也可以告訴我們.
莫迪拉的瑪利亞由她一信主.
就是這樣跟隨耶穌到底.
一直這樣跟隨 致死都跟隨著.
所以你可以想像得到 她對主耶穌的感恩之情.
是多麼的深.
所以我想我們人生如果找到一位.
那麼好的救贖主.
我們用我們的方式 男性的方式就是.
盡我們的心意來服侍.
我們用那種專一,理性.
我們的感性沒有女性那麼多.
但是我們的理性 我們那個擺上.
我們的決心有時候大於女性的.
所以是平衡的.
所以神做人是很奇妙的.
所以又需要男又需要女.
女性 莫迪拉的瑪利亞就是一個表表者.

$^{1321}$她將她的感性去到最重.
就是這樣擺上為主 不管那麼多.
就是這樣 她就見了幾次的神蹟.
所以當人們走了之後 第20章的《約翰福音》.
她坐在墳墓裡哭 耶穌向她顯現.
第一個向她顯現 所以很奇妙的.
她這種近乎麻木 近乎不理智的跟隨.
換來的就是神的兒子.
要不斷讓她經歷這種恩典的醫治.
使得她能夠 她甚至開心到抓住耶穌的腳.
不讓祂走 不讓祂升天.
這樣的一個奇妙的20章的記載.
你可以想像得到 莫迪拉的瑪利亞.
對主耶穌的愛是多麼的深.
所以如果我們感孫家有弟兄姊妹.
或者有姊妹 可能是單身的.
我不知道莫迪拉的瑪利亞是不是單身的.
她如果能夠那個心專一愛主 愛到這樣.
這個是我們感孫的福氣 或者是教會的福氣.
今天有多少個莫迪拉的瑪利亞呢?.
我們可能被很多東西阻礙著我們的心靈.
沒錯的 我們有兒女 有家庭.
有我們自己的個人的想法.
但是我們的心底裡面 到底有沒有主耶穌呢?.
有沒有莫迪拉的瑪利亞這樣擺上事奉呢?.
所以如果我們人生真的沒有東西需要我們牽掛.
我們的心真的要擺在主耶穌裡面.
愛耶穌 服侍耶穌.
甚至乎旁人會覺得我們失去了.
好像一些不是太理智 其實不是的.
我們有時候那個感情完全投向主.
那種擺上又會帶來很大的一些神蹟,恩典.
所以我們今天需要的 當然不是說不用理性.
不用理智 不去分析 當然不是這樣.
但是如果我們過多了 只是很理智.
可能我們的信仰只是一種行禮如儀的功課.
就不是真正的信仰.
真正的信仰是你整個生命擺上去 去服侍.
所以我自己現在真的去到接近退休的這三,五年.
真的越來越感受到這個身份 過去都很害怕.

$^{1361}$但是現在越來越覺得作為牧者.
可以安慰,幫助不同的人 是很榮耀.
是這個身份的榮耀 不是我這個人.
所以感覺到是真的很開心.
實用我這個這麼小恩賜的人.
越來越感受到這個身份多麼的尊貴.
所以如果我們不事奉 其實沒有人會以槍指著你的.
但是如果我們擺上事奉.
得到的只是最大的那種滿足.
以及將來得到的無論神或者人對你的尊敬.
所以如果我們只是一種不冷不熱的服侍.
可能很舒服 你什麼都會記住.
但是如果你每次都跟神計算.
其實你自己就浪費了很多服侍的機會.
帶來的那些福氣 我們一起禱告.
天父啊 我們今天都思想了有不同的人.
特別阿居拉夫婦 一對夫婦一起服侍.
這樣來支持保羅 認同保羅的傳道.
站在保羅的受苦裡面.
我們又見到這個保羅所敬愛.
對他有幫助的這位弟兄.
也是亞細亞初結果子.
我們又想到這位莫大拉的瑪麗亞.
雖然我們保羅未必認識她.
但是在保羅的內心.
總有這位在羅馬的瑪麗亞.
為主多受勞苦 很難想象在保羅的眼裡.
到底多多的勞苦是去到什麼程度.
什麼性質的勞苦 是在教會裡面鞠躬盡瘁.
很單純的服侍 還是其它.
這些保羅沒有詳細的描述.
但是相信對羅馬的教會的弟兄姊妹.
是會認識這個瑪麗亞.
我們見到在《聖經福音書》裡面.
所提的瑪麗亞都是這樣專一的愛主.
興起我們這一代的姐妹.
來用他們的心 最好的去服侍主.
擺上心靈來在神的家裡面.
來做這些默默的工作.
我們盼望主來施恩.

$^{1401}$我們更加需要有阿哥拉夫婦.
這樣的對真理的認識.
又是這麼謙卑的默默的去幫助不同的人.
連阿波羅都是這樣.
在阿哥拉夫婦的教導下.
成為更加有被主用的一個年輕的牧者.
但願我們透過保羅的文案.
給我們有更深的感受.
原來每一個人在保羅的眼裡面.
都是這樣來擺上自己的生命去服侍的.
我們求主施恩.
叫我們都能夠對別人所擺上為主的欣賞.
和帶著一份的感恩.
更加是激勵我們自己.
在主的救恩裡面來到更多的擺上.
多謝主聽我們感恩禱告.
奉耶穌基督的名.
阿們.
\newpage



\section{羅馬書}
\label{sec:HW8W_v51HY8}
\textbf{2026-03-20 周五晚查經班《羅馬書(四十):分題:來自保羅的問安(三)》 講員:曾裔貴牧師}
\newline
\newline
連結: \href{https://youtube.com/watch?v=HW8W-v51HY8}{\texttt{https://youtube.com/watch?v=HW8W-v51HY8}} ~~~~ 語音日期: 2026-03-24
\newline
\newline
\hyperref[sec:oQdnBIyZabY]{\small{< < < PREV SERMON < < <}}
~
\hyperref[sec:index_chronic]{\small{[返順時目]}}
~
\hyperref[sec:index_scriptual]{\small{[返順卷目]}}
~
\hyperref[sec:Ex_WASW50ZA]{\small{> > > NEXT SERMON > > >}}
\newline
\newline
$^{1}$好了 我們一起祈禱.
再一次我們感恩.
晚上我們有一個安靜的時間.
我們可以一起去研讀《羅馬書》.
逐段來思想《述道保羅》給我們的教訓.
多謝主給我們因為這封是一封書信.
去到結尾的時候.
也都真的有當時書信的格式.
對收信的人來去問安,祝福.
我們看到這一段的問安是相當不尋常.
因為出奇的多人.
而且是唯一一卷《羅馬書》.
保羅的書信有這麼大量的問候.
而且都不是一些重要的教會領袖.
有些可能是很普通的弟兄姊妹.
甚至有些名字在希臘文都是一些奴隸.
我們看到《述道保羅》.
都好像真的有一個用意.
透過問安來將一些好像很艱深.
很難明白的一些肅靈的道理.
透過生活,透過當時那些階層的複雜.
以及很明顯那些階級的分明.
以及真的有很多的貧富懸殊.
甚至乎有階級之間的那種競爭.
但是我們看到在基督裡面的這個合一.
這一份大家的平等.
大家在神裡面擺上的那個事奉的生命.
就是這樣令到他感動著保羅.
我們所以真的需要很認真地從保羅的眼光.
去看人與人之間在救恩裡面.
在神將人拯救之後.
尤其是初期教會很明顯是有猶太人.
以及外邦人 而且外邦都分了不同的階層.
一起來到去聚會的時候.
就讓我們真的看到原來初期的教會.
已經是這樣來到去.
經驗到救恩帶來的那個生命的改變.
所以「探望神」都讓我們學習.
其實很多名字跟我們甚至乎.
絲毫沒有任何的聯想關係.

$^{41}$因為都是一些譯音 以及我們都不是.
在那些時代出現的人.
所以可能我們從來都是很快地讀完了.
就算 甚至乎會有一種厭倦.
相當沉悶的這樣的文案 很沉醉等等.
但是「盼望」我們都能夠從其中.
明白薛豪的內心的想法.
求主帶領我們今天晚上.
差不多我們去到這本書的結尾.
都很需要我們思考.
我們查了差不多一年多了.
我們有什麼書卷 或者不同的方式.
繼續大家彼此學習 我們「仰望神」的一點.
這樣禱告奉耶穌基督的名 阿們.
好 感恩 我們能夠來到去再看.
如果在這個分題 就是來自保羅的文案.
當然這個不是開始書信的時候那個文案.
而是在羅馬書 特別是一封書信.
一封當時的格式.
最後都會來去吩咐哪些收信的人.
他們會怎樣 這些都是書信的一些格式.
也都不會例外.
但是我們看到剛才那個祈禱裡面大家都知道.
文案有26位的 有男有女.
是非常不尋常的.
我們有機會再看第17節.
連續上次兩星期都點出了.
去到第17節立即有一個很強烈的.
一個反向的教導.
就是如果第3到第16節.
是很好的在基督裡面的合一的相交.
彼此的來去欣賞.
大家不同的崗位的服侍.
怎樣為神的救恩.
來到去不同的地方 不同的崗位.
不同的人生處境來服侍.
但是去到第17節 去到第20節.
這四節聖經就告訴我們.
有人會全心在教會來搗亂的.
有時候是他與生俱來的天然的性格.

$^{81}$很容易就會和別人作對.
又或者有一些小小的嫌隙.
擺不低又會容易這樣.
但是保羅在這裡講的就是.
有幾樣東西 離間 叫人跌倒.
將所學的真道背棄.
也就是說可能有意無意.
有些是處心積慮要來到去攔阻真理的道.
那就要很小心了.
保羅就是說 叫羅馬的教會.
或者叫大力教會的牧者.
要留意這樣的人.
就是說要守望著教會的真理.
那個的合一 那個的純全.
以致不會有一些像豺狼.
背棄真理的人來搶奪了很多人的心靈.
所以第17節就好像憤外一個反差.
憤外凸顯第三到第16節.
如果加上非比這個女執事.
羅馬教會要用在主裡面接待他.
要合乎聖徒的體統的接待.
這樣憤外顯得這一個大段落的文案.
就是想帶出一些很重要的含義.
就是在主裡面.
我們真的要在真理裡面相交.
在主裡面我們要排除任何的階級.
任何在無論那個社會的主流的民意.
那個輿論是怎麼樣不適合聖經呢?.
我們都要學習跟聖經.
不要跟社會那個思潮 那個輿論.
羅馬帝國是很明顯的.
差不多有七成是奴隸的.
你可以想像得到已經是很普遍的了.
保羅都不是會說.
信了耶穌做了使徒就要馬上要發起.
要完全禁閉任何的奴隸.
因為是不人道的 是不合聖經的真理的.
神要的是平等的.
這些是很高尚的口號和主張.
但是我們任何一個的使徒.

$^{121}$都沒有在他們有心之年.
來將這些東西用一種可能立法,政治.
來推翻它 來撥亂它反正.
反而我們看到的就是在言語之間.
在他們的為人之間.
對不同的身份,不同的民族,不同的階層.
等等的信徒出於很大的愛心的.
互相的平等的相交,接待.
這些就變成了好像在一個.
因為始終是人 始終是一個制度.
有它的結構 如果等於一個主力牆.
你硬要將主力牆拆掉 不漂亮就拆掉.
上面如果有六,七十層樓.
那就很危險了 整座樓就會倒塌了.
所以我們有時候就要看到那個利害關係.
不是那個神不准我們這樣做.
或者保羅不是無視.
說明明是剝削 明明是奴隸 是非人生活.
是一個物件 是一個只會說話的.
不過好像一個機械人一樣.
任由那個主人去逮捕的 為什麼會這樣呢?.
譬如女性沒有地位的 女性她沒有任何的.
甚至在法庭上當時的法律體制下.
女性都不可以去做一個見證人的.
就是找見證人 找男性都不會找女性的.
這是很普遍社會的文化 社會的主流的想法.
所以當我們今天 我們好像在香港,在世界.
我們看到這麼多女性到處走.
完全找工作 各方面是不可能的.
在保羅的年代更加不可能的.
女性是要在家裡面的 是不可能的.
就是她可以這樣自由 自在的生活.
你難以想像 如果保羅他這樣的思潮.
在今天的社會是兩件事差不多是來源不切的.
所以我們就憤外見到保羅.
如果刻意這樣來向不同的人去問安.
那就是說他背後有一個很強烈的神學的想法.
這個神學甚至比地上的.
我們認為是不可能改變的這些政治,這些民生.
這些奴不可破的制度.

$^{161}$原來神學可以慢慢地將這些事情.
慢慢地背叛反正.
慢慢地要在很多人的心裡面.
有了這一種好像共識.
這一種慢慢地滲透了的真理.
就可以指日可待.
那些平等就可以來去發揮了.
所以是一件很奇妙的事.
原來神往往的工作好像很緩慢.
但是非常有效的.
但是我們看到《羅馬書》就是一本很不尋常的書.
居然間一本這麼重要的講真理的書本.
居然間有差不多半章 16章的半章聖經.
因為都不止了 三分二.
因為全章都只有23節.
那你就可以想象到有16節.
就是講到保羅對逐個人的文案.
今天下午也是有些東西在講.
接著逐個名去研究.
因為對我們來說 如果我們不看參考書.
不看其它的輔助的資料.
對我們來說是完全沒有任何意義的一個人.
只是記了一筆這個人是哪個名字.
我們連他是什麼國籍.
保羅都沒有詳細解釋.
連他的詞根 到底他是希臘人 猶太人.
甚至乎是一些拉丁裔的人等等.
我們都不知道的.
但是必須是透過學者這麼多年的研究.
後來有不同的一些文獻的出土.
或者是教父的時代.
他們是怎麼理解保羅對這些.
在教會裡面的信徒的那個文案.
就是說 還是在保羅一代的人.
都可能會認識 但是都不是個個都認識的.
所以我們一直預備的時候.
嘩 講一下這個名字 講一下這個名字.
有時候開始頭暈 每個名字都差不多的譯音.
接著 其實關我們什麼事呢?.
我們在香港這麼多東西要跟.

$^{201}$這麼多東西要面對 這麼多的事故.
就是你一蓋上《聖經》.
這些人根本就馬上淹沒在那個歷史.
在那個文字裡面 跟我們是沒有關係的.
所以可能是一個很宏觀去看的.
宏觀的一個地步就是.
因為學者給我們很多的發現.
有些名字是很肯定是奴隸來的.
而且是普遍的奴隸的名字.
就是我們中國人有一些語場片.
什麼「來福」啊 什麼「緹」啊.
那些類似就是一些下人 約定俗成.
通常叫那些狗仔的名字.
或者叫阿財啊 旺財啊 這樣 大家約定俗成了.
所以 如果使徒們他們來到去.
對這些名字這樣的文案.
而這些名字的背後 相當肯定是奴隸的身份.
而能夠在教會裡面成為弟兄姊妹.
這樣保羅的含義.
就不是想我們記住哪個人 他的什麼名字這麼簡單.
而是保羅是好像不同的類別的人.
都在他文案的那個範圍.
至少在他寫的時候.
一定是他的內心對這些弟兄姊妹.
有份外的一份的感情.
對他們一份的感恩 一份的讚賞.
一份對他們的肯定.
還有在智慧行間 就好像想透過這樣的文案.
將他之前講的一些很艱深的神學的真理.
在基督裡面能夠成為一個好像很抽象的道理.
很實實在在 透過不同的人的這些名字.
就這樣共野一爐 擺在聖經一本.
好像很平原鋪植罪 就是這樣的名字.
就是這樣的設定 就是這樣的一個模式.
問誰安 問誰安.
有些可能會多了一點他們的身份.
或者有些他們所做的事情.
都是一句起兩句止 都不是長篇大論.
因為始終是一封書信.
始終是保羅和這些人的關係 日於言表的.

$^{241}$收信的 因為收信的是羅馬的教會.
是羅馬的教會才會更加明白這一幫人.
等於我現在在雄恩.
如果我要在海外寫一封信.
寫給感受一些我之前跟他們一起服侍的.
譬如栽培部 團契部認識的.
那就很簡單 感孫一弟兄姊妹見到趙宇棠.
這樣才算 趙小玲這樣才算.
我講一兩句就已經是大家都會很明白的了.
小玲在聖幕方面給我們有很大的進步.
啟發 帶我們去敬拜神.
這樣一個描述 感孫一弟兄姊妹個個都明白.
但是對其他的人 可能根本以為不知道是什麼.
就變成我們就會明白 保羅這麼簡單的描述.
但其實對當時羅馬教會已經是很大的一個突破震撼.
首先就是羅馬的教會自己本身有沒有分門別類呢?.
有沒有貴族還貴族生活 奴隸還奴隸呢?.
有沒有我們有一種文化的洗腦 女性沒有地位.
所以在教會都應該有地位.
但是我們看到這個尤尼婭是女的使徒.
那你就要想了 到底她是在使徒的行列當中.
在眾使徒的心裡面 她是有地位.
還是她本身就是使徒呢?.
使徒這個概念 這個積分.
到底是保羅眼中的十二使徒.
還是一個宏觀一點的 闊一點的概念呢?.
這些學者都是做了很多的研究.
所以我們就看到這個不僅是菲比這個人.
其餘的都按照著他們的身份.
他們的階層來去問案.
所以我們其實很感恩.
我們能夠在研讀羅馬書.
有機會我們稍微慢一點去看這些人的名字.
來將整個羅馬書做一個整合的時候.
我們發現原來 特別我上一個星期特別強調.
加拉太書就是很清楚對錯的神學.
就是說如果一有人提出一些異端.
信耶穌要加上行為 那個行為可以是割禮.
可以是其他的一些禮儀來換取救恩.
那就是敵黨神的福音了 那就是錯誤了.

$^{281}$但是羅馬書就不是這樣了.
羅馬書就是很想透過在生活裡面.
將這個合一的平等的這個身份.
來去很平鋪直堆地.
很生活性地來去換出來.
使老人認識 原來真是在耶穌裡面.
我們必須學習 大家都是在神的眼中是合一的.
這個合一就絕對不是每個人都要一樣.
立正一樣 穿的衣服 喜好.
全部要一樣 就不是這樣的意思.
連事奉都要一樣 也都不是這樣的意思.
仍然給我們大家有不同的思想.
有不同的做事的方法 有不同的性格.
但是原來在神的裡面.
一定要謀求的是一個坦誠主裡面的合一.
就是說心靈是要一的.
所以很寶貴 我們有機會再去看.
我上個星期看過了 專輯都給了大家了.
我就不再透露了.
所以是一個示範來的 是一個示範來的.
是一個告訴我們 原來是可以的.
原來在主耶穌裡面.
真是可以與不同的人相交的.
真是可以來到去活出這個合一來的.
加利泰書是應該合一.
羅馬書保羅給我們看到的就是.
我們已經在這個合一裡面.
所以我們要活出來.
我們要在我們的生活.
盡量與不同的弟兄姊妹相交 又交往.
當然你說心交的就不可能每個人都可以.
這個就是我們大家會明白的.
我們上個星期講了 伯吉拉夫婦.
接著我們就看過這個以拜尼圖.
亞細亞初屬的果子.
我們也提過其他的初期果子.
有斯蒂凡一丹 或者斯蒂凡.
是亞干納初期的果子.
剛才這個以拜尼圖就是亞細亞 Asia Minor.
就是以弗所那一帶 哥羅西那一帶的初期的果子.

$^{321}$因為以弗所和哥羅西相當近的.
所以都是叫姐妹食的.
接著我們就講過 上個星期我們最後講到瑪利亞.
為你們多受勞苦的瑪利亞.
瑪利亞很簡單 這個就是在羅馬教會.
一個相當普通的名字的一位姐妹來服侍.
保羅就提到她是為你們多受勞苦的.
這一位的瑪利亞.
在《聖經》裡面總共有六位的女性 瑪利亞.
我們也講過主耶穌的媽媽就是瑪利亞.
活達拉的瑪利亞.
我們上個星期就比較多時間去看這位姐妹.
曾經是有鬼婦 有七隻鬼在她裡面.
鬼婦是很慘的 因為不由自主的.
那些鬼什麼時候附身呢?它就會發作了.
所以它是隨時突然之間.
所以那些鬼婦的人是很累的.
因為那些鬼一來是會令到她突然之間空神落塞.
力大無窮 又或者可以用其它的情況.
很憎恨靈 做出一些連她自己都控制不了自己的事情.
所以這個是真正的鬼婦.
所以不是那些精神病 精神病是另外一種.
但是鬼婦就真是受著那個推控的.
就是被那些靈界的鬼來操縱著那個人的.
是完全不由自主的.
她也都不能夠承認耶穌是救主的 講不了的.
因為鬼令到她不能夠講的 所以一般來說.
如果真是很清楚她是鬼婦 而且有生命危險.
因為可能要她做一些自毀的事情 比如自殺.
那就要立即趕了 立即要奉主的名.
要將這些鬼婦的朋友要趕.
當然必須 她趕了之後 她就必須盡快順主.
有聖靈在她內心 才能夠有力量去抵擋那些魔鬼.
靈界再入她的心.
當然就要問 到底她有什麼原因令到她鬼婦.
比如她是不是自己以前交軌 迷信 占卜.
問米等等那些 她自己開放了那個靈界.
甚至可能有些是上一代一些相當重.
跟那些靈界有關的 那就要求主幫助.
要去清除 了結那些關係.

$^{361}$最主要就是讓聖靈在她的內心.
是有一個很堅定的依靠.
就是說她自己是真的在清醒的時候.
是要立即順主 因為她自己知道有靈界在裡面.
她知道什麼是鬼婦.
所以她更加要快些順主 接受是聖靈.
接受是神 這一位聖潔的神進入她的內心.
所以梅特拉的瑪麗娜很清楚她自己的那個無助.
但是耶穌來到將她拯救.
將七隻鬼在她的生命裡面都拿走了.
她從此就沒有再鬼婦了.
從此就一直跟隨耶穌 服侍耶穌.
供養耶穌基督的.
她有什麼工作 聖經就沒有興趣去交代.
大概可能大家都知道.
因為如果她是沒有家庭 沒有結婚.
其實她沒有什麼經濟能力的.
其實都是相當貧窮的婦女.
不是一些有錢的.
比如在福音書有一個叫做婦薩的妻子.
蘇薩娜 婦薩是希律的家宰.
即是說是御用的廚師.
所以又有錢的婦人.
梅特拉的瑪麗娜可能只是一些好基層.
好像文蜜婆 鵝頸橋那些.
但是她一信得住 就有很大的生命改變.
我們有足夠的聖經給我們去思想這一位的姐妹.
今天我們要看的是伯大靈的瑪麗娜.
這些我們看過了.
我們有兩段經文和大家揭一揭.
伯大靈的瑪麗娜.
我們先看看路加的十章.
其實由第九章的51節開始.
就是路加福音的另外一個序詩.
就是由第九章的51節開始.
就是告訴我們 路加醫生告訴我們.
他寄了九章聖經 很多的神蹟.
耶穌很多不同的工作.
但是現在他應該是進入後的服侍的年代了.
耶穌差不多是大概三年多的傳道.

$^{401}$現在是後半段 接近最後的一年左右的事奉.
路加就答了一個題目 就是他面向耶路撒冷.
向著耶路撒冷上去.
當然所謂面向就不是一兩個星期.
都是一段長的時間.
但是路加就很想用這個標題告訴我們.
耶穌這一位人質最終他做這麼多事情.
所謂已經很出名了.
但是最終其實他是面向那個十字架.
這個是路加告訴我們 這是九章的51節.
耶穌被接上升的日子將到.
他就定義向耶路撒冷去 就是這個標題.
我們看的就是在他一路上去耶路撒冷的時候.
就有一段經文 第38節.
他們走路的時候 耶穌進了一個村莊.
有一個女人名叫瑪大 接她到自己家裡.
她有一個妹子名叫瑪利亞.
在耶穌腳前坐著 聽她的道.
瑪大事後的事多 心裡忙亂 就進前來說.
主啊 我的妹子留下我一個人侍候.
你不在意嗎?請吩咐她來幫助我.
耶穌回答說 瑪大 瑪大 你為許多的事思慮煩擾.
但是不可少的 只有一件.
瑪利亞已經選擇那尚好的福分.
是不能奪去的 一段有時候都令我們很困惑.
我們很容易對號入座 因為耶穌責備瑪大.
稱讚瑪利亞 我們就說.
原來在教會做事就不要做到像瑪大那樣忙亂.
所以瑪大就不要的 反而只是做事不好的.
應該要像瑪利亞那樣沉靜 學道.
專心坐在耶穌的腳前聽道.
我們很容易就會將兩姐妹很快就對號入座.
很快將一個很個別的一次的事件.
就看為瑪利亞永遠坐在那裡聽道.
瑪大永遠就是這樣做事的.
其實這裡只是耶穌剛剛一次去到瑪大 瑪利亞.
當然我們還不知道 在這裡的時候.
聖經還沒有告訴我們.
原來祂還有一個弟弟叫拉撒路.
就是這兩姐妹的弟弟.

$^{441}$但是這段聖經的確耶穌真的這樣來跟瑪大說.
也都是間接 好像都是稱讚瑪利亞.
很安靜的去把握著 耶穌去到自己家裡.
耶穌在這個時候可能真的有一些的教導.
而的確瑪大真的很忙亂的 作為大家姐.
很忙亂的去接待一個貴賓.
一位這麼重要的人物 一位猶太人的拉比.
甚至有能夠行神蹟的.
因為當時很多的風聲傳耶穌是來塞亞的.
所以一個這麼特別的人 可想而知.
如果你在家裡的家庭主婦 有一個很出名的.
是你的子女 大學的教授.
他是諾貝爾獎的 剛好來香港.
接著你的兒子又很榮幸.
邀請到他來感受家裡 來到團年飯吃飯.
如果你是那個家庭主婦.
你是不是很開心 很雀躍.
真的是做很多拿手的菜.
來去款待這個大學教授.
我們不會說 媽媽你不要做了.
團年飯有什麼重要呢?.
這個教授坐在那裡講那些學問才重要.
我們大家一起坐在那裡聽個角度.
其實兩件事是沒有衝突的.
我們很快就變成將馬大.
就標籤他是一個只會做事 不安靜的.
忙亂的 一點都不讓他心靈平靜的.
我們一看到馬尼亞 我希望做馬尼亞.
求主使我成為馬尼亞.
回到教會什麼都不要做.
只想聽到就行了 聽到是對的.
不聽到做事是不行的 是忙亂的.
有時候我們就好像很有趣的.
很搞笑的看了一次的事件.
就說這個馬大就是不對了 不肯屬靈了.
但其實不是的.
一個家庭主婦要接待一個貴賓.
一個很重要的場合.
真是多怕怠慢 多怕我們做不夠.
使這個貴賓享受不了我們中國人的美食.

$^{481}$你就大概想一想 你就不會說 媽媽.
你怎麼這麼馬大的 你怎麼這麼忙亂的.
你怎麼這麼不懂.
最重要的就是安靜 耶穌說的.
一件上好的福份是不能夠奪去的.
馬尼亞已經損傷了.
你還要這麼多事情在思慮煩擾.
那就變成我們太快就一個這樣的小小的記載.
我們就好像斷定了.
馬大的人生就是永遠做事 永遠不用腦.
永遠只是做一個造詞這樣.
所以其實是不需要的.
有時候我們看錯了聖經.
很容易就會太過主觀的對號入座.
只是經文的描述而已.
所以其實耶穌也沒有對馬大有過度的責備.
不過在剛好這樣的情況下.
馬尼亞似乎沒有什麼需要這樣來忙亂的時候.
很自然就安靜的來聽耶穌的教導.
有可能拉薩路在現場.
不過沒有記載他的名字.
這裡就好像告訴我們.
是緊接於這個撒馬尼亞人的故事.
主耶穌就進入了一個小小的村莊叫做伯大利.
在這裡也沒有記載他的名字.
我們再看一看 這次是上限是40幅音.
約翰的十一章.
這篇經文就相當長的描述.
所以對於我們理解伯大利的馬尼亞.
就會多一點點的資料.
十一章的第一節開始 相當長.
有一個患病的人 名叫拉薩路.
住在伯大利 很清楚.
作者約翰來到不斷想讀者.
知道他記載了這個神蹟.
記載了拉薩路復活.
所以他講了那個地點.
是什麼家庭等等的事情.
還有整個的事件,神蹟的前因後果.
來描述 接著他說就是馬尼亞.

$^{521}$和他的姐姐馬大的村莊 就是這個伯大利.
所以路加坤那裡就是伯大利的馬尼亞.
和伯大利的馬大.
第二節 這馬尼亞就是那用香膏抹主.
又用頭髮插他腳的.
患病的拉薩路是他的兄弟.
相當詳盡 原來聖經作者原來這麼細心.
怕我們後代的人 或者收信的讀者.
不知道他描繪哪個馬尼亞 哪一個家庭.
所以他現在就很詳細地來幫我們講述.
就是不會有錯 不會有任何的含糊.
就是那個馬尼亞 這次就要凸顯.
馬尼亞有一件事做了.
就是在耶穌安葬的事情上.
馬尼亞趁著耶穌還沒有死.
因為死人才需要高抹的.
所以死人高抹就很正常.
福音書就說很正常的婦女.
就是帶著香膏去高抹耶穌.
但是耶穌在之前也說過 祂會復活.
所以馬尼亞就可能有一個這樣的理解.
耶穌復活了 我們就高抹不了祂了.
所以這個預先的高抹.
是有相當大的那個香氣.
就是說祂是對復活深信不疑.
而對祂預先的高抹主耶穌.
在一個筵席那裡高抹主耶穌.
所以你可以說祂的信心.
是大到對主耶穌的那個教導.
那個深信不疑.
而來得到耶穌基督真正對祂的稱讚的.
所以現在約翰都要記回這個美事.
就是哪一位馬尼亞.
就是用香膏來高抹耶穌的這一位馬尼亞.
所以有病的那個弟弟就是拉撒路.
第三節 他姊妹兩個就打發人去見耶穌說.
主啊 你所愛的人病了.
那就是幾可圈可點.
那些門徒肯定認識耶穌.
因為耶穌都似乎在福音書的記載.

$^{561}$偶然間耶穌就會回去伯大利過夜.
在他的那個壽蘭州.
所以伯大利的這個家庭.
他們有一些僕人的.
到底他們是不是真的很有錢呢?.
又不是太清楚.
至少馬尼亞如果打破了那個玉瓶.
那個是一個真拿達香膏.
如果就是這個馬尼亞.
那個玉瓶就相當貴重.
還有那些香膏的價值都是相當貴重的.
所以不惜用最好的來高抹耶穌.
還沒死的耶穌 只不過是一個普通的坐席.
在這個法利賽人西門的家庭坐席而已.
其實不需要用香膏來幫他抹的.
找一些僕人的奴隸去洗腳就可以了.
其實不需要用香膏的.
因為香膏是要用在那些房庫.
和死了最重要的那些有錢人的.
那些屍體的房庫的.
所以這樣來到某程度.
猶大賣主那個猶大說.
哇 用32銀子的香膏來抹耶穌.
現在是平常飲宴而已 坐席而已.
不需要用這麼貴的 去洗一下腳就可以了.
所以猶大也不是錯的.
不過他的心其實就不是想這樣.
他是想拿了 不要浪費了.
想拿去變賣.
所以這就是兩者不同的對耶穌的愛.
就是凸顯出來的.
不過這裡告訴我們.
你所愛的人病了.
如果這句話真的照足了原文.
照足了那個語氣.
約翰就很清楚地記載.
那份愛就凸顯了耶穌和這幾姐弟.
這個家庭的關係是非常密切的.
根本不需要說誰是我的所愛.
你說的是哪一個 奴隸.

$^{601}$我是哪一個 是馬尼亞還是馬大.
你說清楚一點 不同做法的.
其實一句話 你所愛的人病了.
這些奴隸其實就不是在顧慮玄虛.
而是凸顯了這個家庭和耶穌的關係.
所以如果我們整個家庭和耶穌的關係.
都是這麼密切就很好了.
那個關係去到一個地步.
根本不需要畫公仔畫出腸的.
很簡單 耶穌已經明白.
已經知道他們的弟弟拉薩路.
有可能是很詳細解釋了.
不過約翰沒有記載也不奇怪.
但是無論如何 耶穌知道一件事.
就是祂所愛的人病了.
接著耶穌就聽到.
這個病不至於死的 乃是為神的榮耀.
叫神的兒子 因此得榮耀.
不僅是神得榮耀.
還要是神的兒子 也就是耶穌自己都得榮耀.
這個病不至於死.
也就是說和生死有關連.
那個榮耀又是和死.
死得榮耀就是復活.
這樣就已經講了.
所以這句話絕對不是一個良薄的人說的.
因為我們要記住.
那些下人那些奴隸來和耶穌說的時候.
是帶著一份信心的 不是一種質問的.
是一份托付 由兩個主人托付他們來找耶穌.
就告訴耶穌.
那耶穌講的這個是一個預言性的.
在這個病了的拉撒路身上將來會有東西看.
而這個有東西看就是叫神的兒子.
因為他後來的遭遇 就會得到榮耀的.
所以我們要帶著這樣的概念.
去讀這個故事就會完整一點.
也都不會有一種誤解.
哇 耶穌很良薄 講這樣的說話.
現在最快的就是你馬上起程.

$^{641}$跟我們回去馬大利 那就是一個問題了.
不過第五節 耶穌掃來愛馬大和他妹子.
並拉撒路 聽見拉撒路病了.
就在所居之地凝住了兩天.
作者真是很好.
將那個故事 很清楚的那個敘事描述給我們知道.
他首先講了那個大前提.
這些大前提 這些畫龍點睛是很重要的.
耶穌是掃來愛這個家庭.
愛馬大 愛馬利亞 愛拉撒路.
是掃來愛 不是偶然一次.
愛 是一直都愛 掃來都愛.
所以他的停留 我們是不應該理解他是不近人情.
在暴露玄虛 在製造緊張 其實就不是.
因為已經講了 這個病是要在他的身上.
彰顯神的榮耀.
然後最終是要叫神的兒子得到榮耀.
得到榮耀 其實就是為了叫人相信有復活.
所以這個是一個序幕.
有真正永恆的復活之先.
有拉撒路 好像一個示範單位.
讓你知道將來會有一個更加大.
更加永恆的復活.
就是死而復活的耶穌.
由陰間死了 在陰間第三天就升天.
在神的大能下來復活.
所以我們如果比較完整的看這個故事.
我們就很清楚不會對耶穌有任何的抱怨的.
但是你是當事人的瑪大馬尼亞.
耶穌現在故意多停兩天都還沒有起程.
有些焦急 而在身邊有些人.
慫恿的來埋怨.
這些慫恿的埋怨是會增加那個主角.
那個當事人那個焦急 那個信心的軟弱.
信心軟弱就會對耶穌開始有抱怨.
對耶穌的應許開始也會有一個這樣的誤解.
所以第八節.
對不起 第七節先.
他停留了兩天之後.
然後對門徒說:我們再往猶太去拜.

$^{681}$接著門徒就說:拉比 猶太人近來要拿石頭打你.
你還往哪裡去?.
意味著耶穌要離開耶路撒冷 上去白大利.
就是離開了那個重要的地方.
的確那是猶太地.
所以當地的猶太人要用石頭打耶穌.
也都很簡單 因為耶穌講的言論.
很多的說話 很多的東西都是很終極的.
好像自己就是最後一次.
他不會再回去這樣的行義的了.
這樣的氣氛的了.
所以當地人聽到耶穌的說話好像相當重.
是若有所指的時候.
一定就是會有一種很有戒心.
對著耶穌也會有一種敵視.
接著是第九節.
耶穌回答說:白日不時有十二小時光.
人在白日走路就不至於跌倒.
因為看見這世上的光.
這個很自然的 當時也都不會有街燈.
晚上如果你沒有這些油燈.
大家聽過《十童裡》那個比喻.
也都不會有電動 也都不會有USB充電等等.
也都不會有一些可以.
日間曬過之後晚上就會發光的那些現在的裝置.
你是晚上就不能夠有活動的了.
只能夠維持在室內很基本的那些團聚.
聊天都不能夠你挑燈夜讀 夜戰這樣不行的了.
所以耶穌就說 如果在黑夜走路就必跌倒.
因為他們有光 當然這裡這麼說的.
只不過是透過一些很技術性的日常生活.
來再帶出一個很重要的含義.
就是在主耶穌裡面永遠都是白日.
永遠都是真光.
所以也都不會有任何的事情來阻礙主耶穌.
祂可以來到去使祂得榮耀.
祂要在拉撒路的身上得榮耀.
這個是最重要的大前提.
所以沒有東西可以阻礙這位真光主耶穌的.
就是耶穌不受那個環境影響的.

$^{721}$如果這裡再用一點點的象徵去看.
就是那個環境不許你去耶穌.
那些人攔阻你.
所以耶穌說 其實祂就是真光.
我不需要在黑夜 現在就是這麼大的攔阻.
我更加要去印證神的榮耀.
神的時間到了.
我要在我的身上.
當然這個時候是透過拉撒路的身上得到榮耀.
得到榮耀最終都是讓人知道.
耶穌基督祂的寓言 祂的犧牲.
祂將來的復活就是跟所有跟祂有關的人是有關的.
所以耶穌就說到自己其實就是白人.
祂自己就是作神的功.
不受任何的攔阻.
就是說現在不合時機.
在耶穌裡面 祂知道這個時間就是神的時間.
祂就可以作為 看為是白人去做功.
所以我們都在我們自己的事務裡面.
我們都要學習耶穌這樣的教導.
有時候那個環境 我們一定要揣摩.
如果是神要我們做 那個環境是不會阻礙的.
如果每次都說黑夜就不要做了.
就是什麼都不要做了.
因為我們不是真光.
那你就搞錯了這段經文了.
耶穌就是想告訴我們 祂知道神的時間.
因為很早已經說了 這個是神要祂去做的.
所以攔阻 門徒很好的勸祂.
現在有年輕人來打你.
耶穌就是告訴他們 白人要做神的工作.
所以攔阻 如果是神的工作 是不可以攔阻的.
第十一次 耶穌說了這話 隨後對他們說.
我們的朋友拉薩路誰了 我去叫醒他.
門徒說 還是很正常的.
我們跟門徒是完全一樣的.
如果現在由重病變死 任何人都失去了信心.
任何人都對耶穌有一種錯誤的誤解.
有兩天時間讓你去到現場 你陪伴也好.
我們有時候會這樣理解.

$^{761}$那是對的 因為我們本身就是凡人.
我們不知道神要透過我去在那裡醫好他.
所以我們應該要去陪伴的.
但是這次是很特殊的.
所以耶穌故意的停留.
表面就令到那些人很焦急.
也都帶出了一個好像被人誤解.
他好像吊高來賣 他好像故弄玄虛.
但是這個就是神的時間.
所以耶穌是在白日做神的工作.
黑夜是不會攔阻祂.
所以他這樣的停留.
的確帶來門徒的誤會.
第十二節 門徒說 主啊他若睡了.
就必好了 如果很安然的睡覺就沒事了.
完全怕不了耶穌說的「睡」.
其實就是耶穌要去行神蹟.
使死了的拉薩路可以復活.
第十三節 耶穌說是指著他死說的.
這些是馬後炮 因為約翰記載了那個故事.
所以耶穌 對不起 約翰在《約翰福音》裡面.
很多時候講那些事後回顧的 事後檢討的.
這件事門徒起初不明白的.
當耶穌復活升天的時候 那些門徒就明白了.
第二章 潔淨聖殿就是這樣.
約翰就很喜歡這樣事後加一些註腳.
就給我們後世的人讀得很簡單很好.
所以耶穌這一句話 其實是指著他死說的.
他們以為是照常的睡了.
第十四節 所以耶穌都很好了.
講了給他們知道 當門徒還是很困惑的時候.
耶穌就明明這樣跟他們說 拉撒路死了.
我們有在那裡就歡喜 這是為你們的緣故.
所以我們理解那個歡喜.
就不是耶穌 簡直就是沒有血性.
簡直還說是朋友 簡直你說自己是神的兒子.
所以南宋的人讀到這些經文就很生氣的.
你不在那裡你就歡喜啊 其實不是這個意思.
是在說耶穌要得到神的榮耀.
要讓門徒認識耶穌他的身份 他的能力.

$^{801}$所以那個歡喜不可以就是有一種這樣的.
好像行災樂禍 就叫歡喜 就不是這個意思.
都講到明 耶穌的歡喜其實是為你們的緣故.
使得他們能夠那個信心得到建立.
得到建立 就是耶穌歡喜.
就不是歡喜他兩天都不過去.
不是歡喜拉撒路死了 他可以去演戲.
做秀 不是的 這個就是耶穌基督的意思.
如今我們可以往他那裡去.
第十六節 多瑪又稱為低土瑪.
就對那同作門徒的說 我們也去和他同死吧.
又似乎真的相當疤苦 對耶穌的這樣的明言.
就是明確的講清楚那個理解.
是為你們的緣故而歡喜.
是為你們的信心增強而歡喜.
不是為了我不去而歡喜.
所以完全是不同的 但是多瑪就講出這樣的說話.
都對耶穌有一種不信任和一種質疑.
第十七節 耶穌到了 到了伯大利.
就知道拉撒路在墳墓裡已經四天了.
兩天不去 兩天再走路 那就總共四天.
伯大利離耶路撒冷不遠 約有六里路.
有好些猶太人來看馬大和馬尼亞.
要為他們的兄弟安慰他們.
因為都知道這個死訊.
馬大聽見耶穌來了 就出去迎接他.
馬尼亞卻仍然坐在家裡.
馬大對耶穌說 主啊.
就是證明馬大是衝到耶穌那裡去迎接耶穌.
馬大的說話就是 主啊 我的兄弟必不死.
就是現在 我也知道你無論向神求什麼 主耶穌求.
神也必賜給你 耶穌就跟他說了.
你兄弟必然復活.
相當令到馬大能以想像 難以相信.
說了一番話 就是說現在你向神求什麼.
神必定會給你的 耶穌就說了.
神必定給我的 就是你的兄弟必然會復活.
馬大就突然間不是太順.
馬大這個24節 馬大說 我知道在末日.
復活的時候 他就是會復活了.

$^{841}$就是給你一個官做 都會復活 每個人都復活.
到時他都復活了.
就是說他現在就會安安靜靜地在墳墓裡.
就是這樣完成了他的人生.
是因為耶穌的延遲 延外之音就是.
耶穌延遲了兩天.
他不能夠得到一個現場的肉身的復活.
但是耶穌開出的那個應許.
就是說他已經有祈禱了.
很清楚神的心意 他故意停留兩天.
就是神的時間未到 因為要得榮耀.
所以令到周圍的人都焦急.
沒有人有這樣的信心.
但是耶穌就是這樣來到 在這個時候.
好像他開出這個應許就是說.
我很清楚神的時間了.
我很清楚神一定會聽我的.
所以你的兄弟 應該這麼說.
稍後就必然復活.
這個稍後不是將來二人復活的時候.
所以末日的時候才要復活 就不是了.
所以馬大的確很難有這樣的即刻的信心.
去信靠耶穌基督.
他是這麼說的.
就是現在你向神求什麼呢?.
神必定會聽的.
耶穌說 神已經聽了.
你的兄弟必然復活了.
所以有時候那個落差.
感恩的神也會給我們人生慢慢的成長.
所以我們真的要跟我們十年前的我.
五年前的我 我們那個信心是應該有增長的.
我們對神的道 對主耶穌.
是應該有一份的心悅誠服的成長的信心.
所以馬大絕對不是在說那些大包圍.
大包圍是人都會說的.
將來會復活的 耶穌也說的.
起碼他這一刻不是挖苦耶穌.
不是來到說耶穌你還行災樂禍.
你兩天沒來 應該就不是這種心態.

$^{881}$他相信他們一家人是很清楚一定有永生.
不容易的.
今天你能不能夠相信你全家都是有永生呢?.
全家都很清楚他們可以上天堂.
瑪利亞 馬大就很清楚.
將來二人復活 末日的時候.
我的兄弟必然會復活.
就是有份於將來永恆的復活.
但是耶穌就說了.
25節 很出名的這一節聖經.
原來就是馬大高低的信心.
耶穌就說了這句永恆的真理.
是我們差不多每一次的安息禮拜的一個宣詔經文.
復活在我 生命也在我.
信我的人雖然死了也必復活.
凡活著信我的人必永遠不死.
你信者運屋.
當然我們只是取了上半段.
下半段其實是一個單日樂.
是一個小小信心的馬大.
和耶穌很清楚神的時間.
就是在這一刻 可能稍後的那一刻.
就是拉薩路必定復活.
必定復活的原因一開始已經說了.
叫神得榮耀.
叫神的兒子因此得榮耀.
當然還要多問一個問題.
得榮耀來做什麼?.
就是要加質門徒的信心.
對於將來永恆的復活是有一個的.
因為他們就是那幫傳道者.
他們自己就是要傳第一個復活的耶穌.
所以這一幫的門徒.
他們就是眼睜睜看著拉薩路跳出來.
他們都是需要經歷這一種成長.
所以這一節聖經25,26節.
耶穌說我就是復活了 我就是生命了.
所以相信我的人雖然死了 必定會復活.
似乎這裡很清楚.
涵蓋的就是在這個地球的.

$^{921}$生老病死的這個定律 這個生命.
就是說肉身的生命在神的時間還沒有完結.
就是這個世界還沒有去到最後.
主耶穌第二次回來之前.
這個世界仍然在這個定律底下.
就是生老病死.
我們突然來到這個世界 是有神的旨意.
我們要生在哪一個家庭 沒得選擇的.
Alice在印尼 沒得選擇的.
我就是在九龍塘的小西湖村 沒得選擇的.
每個人都不同的 後天可以不同.
但是先天就一定是神安排的.
我們信主的必須要認定這件事.
但是後半節就很不同了.
凡是活著信我的人 必永遠不死.
接著耶穌就說 你信者華莫.
這個就是個別了 對馬大當下那一刻.
在這個故事的氛圍裡面.
在這個事件裡面 有一個實在的拉薩路死了.
所以現在耶穌知道他的神的旨意.
要耶穌做一個從來沒有人做過的大神蹟.
就是一個死了的拉薩路 用布都纏住了.
耶穌是叫那些人拿開那個墓門.
就叫拉薩路 其實不需要很大聲.
就是裡面 你想一下有時候我們按那個遙控.
這個遙控就是用紅外線按一下.
有時候我們按兩下 但是原來他接收那個儀器.
要半秒才可以到 這麼近都要半秒.
但是耶穌在那個墳墓 叫裡面的墳墓死了的人.
叫了一聲拉薩路出來 那就出來了.
那個死人就跳出來了 你難以想象.
你今天我們說的那些機械人又打功夫.
在春晚又跳舞 其實用了一整年的訓練.
設定好所有的程序了 所以你以為他自主.
其實一整年前已經開始不斷修改那些程式.
使他做那個動作 做到滾瓜爛熟了.
但是再叫他做其它的 他就不會了.
但是你難以想象 現在第26節.
是耶穌來到跟馬大要提升他的信心.
相信我的人雖然死了 但必定的復活.

$^{961}$所以這句話應該就是即時的那個顯現.
就是稍後當我復活的時候.
拉薩路就會走出來 就會復活了.
接著就問馬大 你相不相信.
凡是活著相信我的人 必定永遠不死.
這個似乎就是另外一個層次了.
就是地上的那個永遠不死.
因為如果是這樣的話 就天下大亂了 因為耶穌死了.
那些門徒全部死了 但是都沒有復活.
不是永遠不死 所以似乎就是在說.
那個生命那個靈性上的那種不死.
就不是活著相信的人.
就可以一直都永遠不需要死 長生不老.
就應該不是物質上 軀殼上.
這個身體的永遠不死 而是在主的裡面.
我們那個生命的力量.
是應該有這種這樣的力量.
你信不信這句話呢?.
所以馬大其實當下他要面對的就是.
到底是復活還是不復活 是沒有東西看這樣.
馬大的意思就是 主啊 你是的.
我相信你是基督 你是神的兒子.
就是那位要流到世界的.
到底是撒耶穌 突然間說這麼正確的神話.
你是基督 你是神的兒子 我相信你.
就是那位要流到世界的.
其實現在問你信不信這句話.
凡換著信我的 就永遠不死.
其實是針對拉撒路 他現在死而復活.
所以馬大是狗不搭白 是太過失儀的懼怕.
而這樣說 好像這麼正確的神話.
你是基督 你是神的兒子.
你是將要流到這個世界的.
其實是信不信呢?.
耶穌問他 你信這話嗎?.
所以這樣的段落的帶來.
約翰這樣來這麼詳細的記載.
其實真是有一個目的 就是要凸顯.
耶穌就是復活 耶穌就是生命.
但是當地在現場 在案發現場的這一幫人.

$^{1001}$其實信不了 因為太難了.
所以今天我們有聖明在我們的生命裡面.
如果我們不相信我們有永恆的生命.
如果我們不斷去操練自己.
對自己的永恆要有一個投資.
要有一個羨慕 要有一個著緊.
其實我們比起馬大這一幫人更加慘.
因為我們這個信仰不是只是給我們去裝飾.
這個是在說 浮臂承載我們進入永恆的準備.
所以我們現在今生要為主而活.
為福音而活就是這麼重要.
因為我們一離開世界 就是無法回頭.
就是馬上去永恆的了.
至少這個是肯定聖經告訴我們的.
這個真確的人類每一個人的命運.
一離開世界就不需要去做任何事情了.
所以趁著大日我們就要做事了.
所以這裡告訴我們 我們不是像馬大那樣.
牛頭不對馬嘴 你是基督 你是神的兒子.
你是將來到世界的 因為這裡告訴我們.
凡活著相信耶穌的就永遠不死.
所以我們和二十五節結合來看.
就是說我們在我們現在有生之年.
我們信主的已經得到一個永恆的救恩.
聖人已經住在我們裡面.
告訴我們我們是有永生.
我們已經是神的兒女.
所以我們現在就要在這個永恆裡面.
為永恆而活 不是地上不要緊.
就是我們能夠再活在地上.
永遠不死的就是我們這個復活的生命.
所以我們是打不死的這個復活的生命.
這樣來為耶穌作見證.
第28節 馬大說了這話就回去暗暗的叫他的妹子.
瑪利亞說 夫子來了 叫你.
瑪利亞聽見了就急忙起來到耶穌那裡去.
似乎耶穌就是在那個露天的近那些.
可能是那個霧的附近.
因為看那些紀錄片都是一個大的霧月.
都是有錢的人才會這樣做的.

$^{1041}$因為如果沒錢就是一亂子的廢墟.
其實是燒垃圾的地方 就像堆填區這樣.
就是沒錢的窮人就是這樣.
但是有錢的就是在墳墓作的.
所以你看到在耶路撒冷.
對著東面就是很多的墳墓.
全部一個方位向著耶路撒冷.
因為一復活他們就向著耶路撒冷 他們就復活.
其實嚴格來說 你向什麼方位都無所謂的.
因為在神那個永恆的真理.
將來所有的人都要復活的了.
如果根據那個次序 主耶穌一回來.
天使像一吹號 所有死了的聖徒就立即復活的了.
所以不是因為你在耶路撒冷.
那個是最貴的那個頭等位 那個是握手的位.
所以就首先復活 就沒有這回事的.
就是你坐飛機 你坐機尾就很長時間都下不去.
就不是這樣 古往今來的 死了的.
完全沒有人認識的都可以復活.
反而沒死的 才是身體的改變.
就被提到天上.
所以這裡馬尼亞就再來到了.
第30節 那時耶穌還沒有進村.
仍在馬大迎接祂的地方.
那些和馬尼亞在家裡安慰祂的猶太人.
見祂急忙起來出去 就跟著祂.
以為祂要往墳墓那裡去哭.
馬尼亞到了耶穌那裡看見祂.
就俯伏在祂腳前說.
主啊 你若早在這裡 我兄弟必不死.
仍然都是這兩天 再加兩天的路程.
所以使得他們有這一種的芥蒂.
也是情有可原的合理.
但是耶穌不需要為了這個遲到.
來解釋和說抱歉.
因為很清楚神的時間 神的原因.
所以耶穌故意的延遲.
是因為神的時間 神的原因.
原因就是神要得榮耀 神的兒子.
要在拉撒路這個死亡上得榮耀.

$^{1081}$所以耶穌不會這樣去解釋.
33節 耶穌看見他哭.
並看見他與他同來的猶太人也哭.
就心裡悲嘆 又甚憂愁.
似乎耶穌是對他們的這一種內心的痛.
有一種的觸動 這個觸動不是發脾氣.
不是因為你們什麼都不信我的.
相信也不是這樣.
耶穌是表達他內心那種的.
與他們那種的心靈的感受.
也是感人的 因為主耶穌清楚自己.
很快就會叫拉撒路復活.
而且是神得榮耀的.
但是他現在在這一刻 因為真的還沒有揭盅.
等於我們看世界盃決賽 如果要護石十二碼.
如果一早知道那個結果的.
如果你真的有這樣的預知能力.
你就不要這麼緊張 一定贏的.
巴西或者法國 因為你知道.
你當然是這樣 但是全世界在看直播的球迷.
都不知道到底射到多少 是那幾秒而已.
那幾秒就是最緊張最重要的.
你一知道那個結果 整件事就完全不同了.
那個心情也完全不同了.
所以耶穌在這裡能夠這樣悲痛.
其實是相當不簡單的.
相當偉大的救主.
「你們把他安放在那裡」.
「他們回答說 請主來看」.
就一直帶著耶穌去到那個墓前.
35節耶穌哭了.
我們未必需要用原文說那種哭法是怎樣.
嘶叫 好像馬的嚎叫聲.
其實未必需要的.
沒理由耶穌突然在眾人面前嚎叫.
其實不需要的 耶穌表達了一份內心的痛.
對於人類面對死亡的那個無助.
那種的無力 那種的絕望.
他來到有這樣的感受 因為他是一個人.
他現在在這一刻 在這個墳墓.

$^{1121}$他有這種的 雖然叫做同理心.
但是肯定比同理心更加強烈的.
所以不是那個動作很大.
好像馬的那樣的嚎叫.
因為也有些人這樣解經 其實不需要這樣想的.
而是你看到那個流程.
耶穌很想將這個生命.
由這一刻灌注到地球,人類.
因為耶穌來這個世界.
最重要的目的就是人類如果沒有救恩.
全部要在永恆的死亡裡面.
是回不了去神的家.
所以死亡是一個非常絕對的敵人.
所以耶穌現在來這個世界就是要剷除這個敵人.
就是用他的生命 用他的死.
來敗壞掌死權的魔鬼和人類.
因為抵擋了神 離開了神.
所以就在魔鬼的掌控下.
所以耶穌就是要面對直向死亡的罪的權勢.
所以耶穌看到人的悲痛,無助.
不僅是見到兩姐妹特別可憐.
而是更加深層的就是神學上.
如果耶穌不去展示復活的能力.
和將來的復活的真實,真理.
人類就沒得救了.
所以現在只是展示一個.
淺嘗一下耶穌就是復活的救主.
所以他的說話就賦予拉薩路.
哇 真的遙控 你想一下.
多恐怖 多厲害.
一句話 拉薩路就可以在墳墓裡面走出來.
很偉大的神蹟來的 真是難以想像的.
這樣遙控法 現在的無人機都很厲害.
飛到這麼遠 準確地投彈.
有了那個軌跡 進入了電腦的程式裡面.
因為那個機裡面有電腦的程式.
就接收了 就去到哪個位置這樣下去.
其實是很偉大的 現在人類的發明.
但是比起耶穌能夠進入一個死亡的人.
他一句話 因為耶穌本身就是復活.

$^{1161}$所以25節就是這麼重要的一個宣告.
所以安息禮拜是真的很好.
一個宣召 是要告訴那些未信的人知道.
不要絕望 我們基督教是有盼望的宗教來的.
復活在耶穌裡面的這樣.
這個告訴他們 三十六節猶太人就說.
你看他愛這個人 愛到多麼的懇切.
當然這個是很合理的這樣的感動.
和這樣理解耶穌 但是其中有人說.
他既然開了核子的瀾 這是第九章.
豈不能夠叫這個人不死.
開始有人這樣對話了.
就是說 原來耶穌之前做的事情.
原來都在人的心裡面有很多不同的想法猜度的.
原來這些狂人的說話 原來才是真正的輿論.
就是說 原來耶穌醫了那個生下來核間的.
那個在盛年的那個人 原來真的會傳開去的.
所以有些人就在那裡指指責責 這樣說.
就是說 意味著可能他們都想.
這次如果耶穌有下一個神蹟.
就真是厲害了.
所以這個只不過是一個對話 不是一個信心.
38節 耶穌由心裡悲嘆 來到墳墓前.
那墳墓是個洞 有一塊石頭擋著.
耶穌說 你們把石頭挪開.
那死人的姐姐 馬大對她說.
主啊 她現在必是臭了.
因為她死了已經四天了.
所以很人之常情.
就是沒有藍燭 耶穌 又知道耶穌有這樣的能力.
又有這樣的慧法 有這樣的說服力.
沒有人敢來嘲笑耶穌 在那裡做什麼騷呢?.
因為耶穌叫那些人把石頭搬開.
這些對死人都挺不敬的 因為你已經告了他.
要等待他的屍體的防護等等.
第40節 這裡我們要留心了.
耶穌說 我不是對你說過.
你若信就必看見神的榮耀.
這個說話其實真是可以不斷的延伸的.
當我們如果對一些事情.

$^{1201}$我們是有所求 認真的去希望.
甚至乎是一個我們的許願.
要做一些事情 為神做一些事情.
這個說話就真是可以應用的.
因為我們如果信 必定會看見神的榮耀.
所以有時候是你要走一步才看見神的榮耀的.
因為你看見那件事就發生了 其實就不需要你信的.
所以一定是信先行的 信然後有行動.
那就會看見神的榮耀 尤其是宗教.
更加要這樣去經歷的 所以這句話.
不妨在我們自己人生要這樣去操練.
接著就是43節.
說明者說 就大聲呼叫說 拉薩路出來.
這個就是一個信心和一個命令.
這個是神的時間 神的旨意.
好像將那個程式放在一個無人機.
放在那個機械人那裡 叫他來跳舞 春晚那樣.
但是這個是一個死了的人.
不是已經植入了晶片 是沒有晶片.
是一句說話 但是明明是一句希伯來文的說話.
是由耶穌無論那個聲音有多大的分貝.
都不會叫醒一個死人的 耶穌在這個時候.
就是神的時間 神的旨意.
記住仍然是那個大前提 要神得榮耀.
要神的兒子 因為這一句話.
使到周圍的人相信耶穌有這樣的復活的能力.
來得榮耀 所以這個時候.
只是一個人拉薩路復活 但是這個父善.
我們每一個人都相信了 耶穌就是復活.
耶穌就是生命 我們就帶著這個信念來生活.
來盼望 來傳福音 帶人去認識耶穌基督.
帶人得到永生 所以是很重要的.
所以當時只是為了一個人.
和當時現場的小部份的人.
但是這個故事的最重要就是神有足夠的能力.
讓所有相信祂的人 將來都可以得到永生.
所以我們其實是幫人進入永生的 很重要.
第44次 死人就出來手腳裹著布.
你想一下 裹著布 布裡面有什麼?.
布裡面就是那些香膏和那些抹藥.

$^{1241}$沉香 用很重的防腐劑 這些香料.
來把它包緊 接著用布來纏.
木乃伊大家見過了 纏得這麼緊.
紮到腫的樣子 是動都動不了的.
根本不想它會動的 但是拉撒路就這樣走出來了.
接著臉上還包著手巾 其實挺恐怖的.
雖然是白天 雖然有那麼多人 但是挺恐怖的.
這個木乃伊走出來 耶穌對他一番話.
解開叫他走 耶穌叫那些人幫他解開那些果絲布.
拉撒路就回家了 回家了 都難以想象.
就是一有了那個生命 拉撒路就再也不臭了.
不會有那些屍蟲一直走回去 一直拉著他.
應該不會的 整個生命立即有了生命力.
完全沒事了 就是說那些皮膚完全復原了.
因為稍後拉撒路還要坐直的.
所以這一篇經文就是伯大尼的《瑪利亞》.
他是非常之認識主耶穌的 很特別.
因為12章才記載瑪利亞來高沒耶穌的.
所以剛才作者約翰說的.
哪個瑪利亞呢?就是那個用坑高耶穌的瑪利亞.
其實那件事還沒有發生的 是發生在拉撒路之後的.
拉撒路復活之後的 所以瑪利亞就學到東西了.
學到功課了 她知道耶穌既然現在能夠叫她弟弟復活.
就是說耶穌能夠允許她將來死了.
就是她自己死了 她自己都會復活這樣的意思.
所以瑪利亞其實沒有人叫她的 她就很快.
在耶穌一次的坐直 去幫她洗腳 用坑高抹一抹.
就好像預示了將來我未必能夠做這個動作.
去告滿你的身體 但現在我就在這個飲宴的場合.
本來只需要用水去洗腳就可以了.
現在我就用我最好的去相信你將來死了必然會復活.
所以瑪利亞的信心又增長了 成長了.
今天我們可能不講得太長了.
因為11章還有一點點經文 12章一段.
不如我們留到下一次再講這一段.
「瑪利亞用坑高抹著」.
今天不好意思 只講了一個人.
就是這個發達人的瑪利亞.
希望我們都用一點時間 保羅逐個.
遲一點就會少一點了 不會太長的了.

$^{1281}$保羅所描述的這些民安的朋友.
他們的身份 他們的國籍 其實是多元的.
變成我們能夠一起去看一下在一個基督的教會裡面.
多元是不是一個問題呢?.
不同的階層是不是一個問題呢?.
是不是我們刻意要將不同的階層分開去呢?.
等等 這些在當時的制度.
男女是不能夠平等的.
奴隸永遠是一個代替品來的.
是不是我們就要這樣斷定了.
要不就你死我活.
不應該再存在這樣的不公義的制度.
我們要推翻它.
是不是要立即就要推倒重來呢?.
立即就要按個鍵就要發生呢?.
因為你有你的制度.
我有我們新的教會 一幫屬神的兒女.
就是說教會再不是地上的普通的一幫人.
教會就是神的家 神的兒女聚會.
是有新的那個家的法則.
就是聖經的教導真理 我們一起禱告.
天父說我們今天用了相當長的時間.
真的很奇妙 這個家庭是主耶穌的所愛的.
他們不是什麼大人物 可能有些錢都不給.
但是他們就是這麼單純的愛主 接待主.
聆聽主的道 還有可能在其它的生活裡面.
和耶穌基督都有一定的關聯.
有可能耶穌最後的受難周.
回到伯大利就是住在他們家裡.
多謝主讓我們看到.
你透過和瑪利亞 瑪達的對話.
讓我們認定你就是復活的主.
你就是那一位的復活 你就是那一位真理的真神.
我們凡是活著相信的人 是永遠不死的.
幫助我們能夠在多困難 我們都認定.
我們是屬主 我們有聖靈 我們可以在主裡面.
去將人帶進永恆的救恩裡面 願神幫助我們.
也都讓我們知道 你已經是復活 在天上的主.
就是說 你隨時會再回來去接教會.
去接每一個聖徒 以致我們每天的生活.

$^{1321}$其實都應該有一個信念 就是等候主的再臨.
求你幫助 奉主耶穌基督的名 阿們.
\newpage



\section{羅馬書}
\label{sec:Ex_WASW50ZA}
\textbf{2026-04-10 周五晚查經班《羅馬書(四十一):分題:來自保羅的問安(四)》 講員:曾裔貴牧師}
\newline
\newline
連結: \href{https://youtube.com/watch?v=Ex-WASW50ZA}{\texttt{https://youtube.com/watch?v=Ex-WASW50ZA}} ~~~~ 語音日期: 2026-04-11
\newline
\newline
\hyperref[sec:HW8W_v51HY8]{\small{< < < PREV SERMON < < <}}
~
\hyperref[sec:index_chronic]{\small{[返順時目]}}
~
\hyperref[sec:index_scriptual]{\small{[返順卷目]}}
~
\hyperref[sec:WmefrJ6J_Lg]{\small{> > > NEXT SERMON > > >}}
\newline
\newline
$^{1}$各位聽眾晚安.
好 和大家有一個簡單的禱告.
再一次我們多謝主讓我們在復活節.
在今年我們都有不同的教會.
有不同的領袖 在特別的節期.
我們都能夠可以來到去親近你.
還有很近距離的去認識十字架的耶穌基督.
和我們的生命關係.
同時我們又能夠在復活節.
有特別的聖餐的聚會.
可以和頂智滿一同的領袖聖餐.
和我們在下午都有報道會.
都見到神的福音的大能.
盼望我們在過了復活節之後的平常的日子.
我們都能夠想起主豐厚的大愛.
我們都能夠繼續有機會去傳福音.
來到去把握 求你帶領我們今晚.
我們還有好幾個人物要繼續來到討論的時候.
求主給我們都是能夠用心.
雖然那些名字和我們是相當陌生的.
但是都給我們能夠從其中透過保羅的教導.
給我們了解當時初期教會他們的情況.
他們的生活的形態,聚會的情況等等.
願神帶領我們 我們這樣仰望禱告.
奉耶穌基督的名 阿們.
我們又可以 雖然停了都有兩個星期.
因為有復活節 還有對上一次都有不方便.
所以我們就要恢復.
對上一次我們講過.
因為提到在《羅馬書》提到.
有一位應該很普通的人物.
就是保羅對他的文安.
這一位為主多受勞苦的瑪利亞.
為了你們的緣故多受勞苦.
就是第六節的 所以趁機會和大家揭一揭.
在聖經裡面 特別是新約聖經裡面.
關乎到 特別是在福音書裡面提到.
關乎到瑪利亞的 其實有總共六個.
我們一路和大家這樣整理了.
來到今天最後這個 當然還有一個.

$^{41}$就是這個格羅伯的妻子瑪利亞.
我們上一次講過拉薩路的一家人.
講過馬大和瑪利亞.
我們還有一段聖經 其實都算是相當重要.
就是12章的《約翰福音》.
這段聖經就再一次讓我們看到.
伯大尼這個家庭和耶穌的關係.
裡面日於言表的那份親密 對耶穌的愛.
這三姐弟對主耶穌的那份的擺上.
是衝破當時那個社會的一些文化倫理.
但是讓我們看到 主耶穌對他們所做的.
是帶著一份很深的欣賞和接納的.
反而就凸顯了有些人的虛假.
而且是門徒猶大.
所以作者約翰他去回憶這些事件的時候.
似乎他加上了很多他個人的事後的觀察.
和神學的反省.
所以《約翰福音》是一本多重神學反省的書.
就不只是記載一些事件的過程.
12章的第一節開始就是講伯大尼的瑪利亞.
「逾越節前六日 耶穌來到伯大尼.
就是他叫拉撒路從死裡復活之處.
有人在那裡給耶穌預備筵席.
馬大侍候 拉撒路也在那與耶穌坐席的人中」.
似乎就是為自己死了復活的弟弟來大排筵席.
雖然沒有刻意說是他們這個家庭來擺筵席.
但是在那個描述裡面都給我們暗示.
那個目的就是為拉撒路的復活.
因為下面有很多的猶太人.
因為聽到拉撒路從死裡復活就來看他.
有沒有參加筵席你就不知道.
有沒有這麼大的空間去容納這麼多人都不知道.
但是就真是為這件事.
有很多猶太人在周圍來看.
因為真是一個很大的神蹟.
一個從死在墳墓裡自己走出來.
還是穆賴伊這樣跳出來.
就是綁著手腳的.
其實是一個很大的一個事件.
因為當時如果沒有直播.

$^{81}$沒有這些傳播的媒體.
這樣來串流的平台等等.
但是這件事已經是對當時的耶路撒冷相當轟動.
而且還要是預節的前六天.
接著就告訴我們那個重點的主角出現.
也是這張聖經上半部份提到.
好像作者告訴我們.
真正愛主的有三種的形態.
我們可以看看哪三類的人.
三種的愛主的形態.
第三節 瑪利亞就拿著一斤.
極貴的珍拿達香膏抹耶穌的腳.
又用自己頭髮去擦.
屋裡就滿了高的香氣.
所以這裡就不要跟《路加夫音》的第七章.
那個叫做有罪的女人.
走到法利賽人西門的家裡.
不請自來的來去告末耶穌.
兩個人兩件事其實是不同的.
所以那裡還在說有罪的女人是沒有名沒有姓的.
好像不是很想去將她標籤.
但是這裡很明顯是有名有姓.
這個伯大尼的瑪利亞.
其實可能是在她家裡.
瑪利亞這次不是坐在那裡聽耶穌講道.
而是她做了一個很奇怪的舉動.
就是用香膏來抹耶穌的腳.
其實很多當時的一些倫理習俗.
就是重要的貴賓有些奴隸洗腳.
這個是一個很普遍的習俗.
如果有一些重要的賢侄.
那個奴隸 那個家的奴隸.
就要幫每一個來去參與這個飲宴.
或者這些賢侄來去洗腳的.
但是為什麼這次又沒有呢?.
是不是好像第十三章的有人福音.
耶穌在這個禹節的晚上.
沒有人幫耶穌洗腳 而門徒又不肯呢?.
這裡告訴我們 這裡就是這麼說的.
她就是這樣用香膏來抹耶穌的腳.

$^{121}$特別要形容那個當時的氣氛.
屋裡面就滿了高的香氣.
就是說眾人皆知 眾人皆聞到.
就是說有了感官上的滿足.
不只是耶穌一個人去享用.
而是整個樓房都能夠知道瑪利亞所做.
所以就引發了一些人有些話.
有一個門徒就是那將要賣耶穌的戛略人猶大.
所以這些肯定是事後的記錄.
路吉 對不起 約翰福音作者.
這一位的使徒約翰 他是大概.
教父告訴我們大概是主後的85至90年左右.
應該不會去到太後的.
就是約翰來寫這一本書.
所以過了一段相當長的時間.
當約翰要來記載耶穌快要離世之前.
一件很關鍵的事件.
就是這個用香膏抹耶穌.
到底是不是因為有罪來抹耶穌呢?.
我們等一下看一下有沒有分別呢?.
那就很不同了.
接著第五節 這個猶大就說.
這香膏為什麼不賣三十兩銀子.
周濟窮人呢?.
他說者說 並不是掛念窮人.
乃因他是個賊 又帶著錢囊.
尚取其中所存的.
耶穌說 猶他吧.
他是為我安葬之日存留的.
因為常有窮人和你們同在.
只是你們不常有我.
約翰當時在他見到的現場.
他參與在那個現場的十二使徒.
聽到猶大其中這個常常都是管錢的門徒.
提出了一個質疑.
對耶穌為什麼要用這麼貴重的香膏呢?.
還有他很快就知道那個價值.
就是說猶大衛什麼會這麼通.
那些婦女香膏的行情呢?.
就是你買一瓶香水.

$^{161}$你去百貨公司 通常地下的.
是吸引一些女士 很方便買到香水.
但是你都不知道買什麼味道 什麼用途有沒有呢?.
我們男士是完全不知道的.
每樣味道差不多的味道.
但是女士就不同了 分到不同的味道.
但是為什麼猶大又會這樣?.
當然這個是題外話.
不過這裡就帶出了一個很奇怪的想法.
似乎是相當合理的.
因為他提出的那個質詢.
這些錢可以去周濟窮人.
30兩很多 真的很多的去周濟窮人.
為什麼要來到好像有一種的浪費呢?.
這樣對於主耶穌的只是抹腳.
一個一次性的用途.
但是耶穌就說.
猶他吧 他是為我安葬之日存留的.
因為常有窮人和你們同在.
只是你們不常有我.
這樣一個的描述.
我們可以先讀下面.
因為整個故事都有整體的討論的.
第九節 有許多猶太人知道耶穌在哪裡.
就來了 不但是為耶穌的緣故.
也是為看祂從死裡所復活的拉撒路.
果然作者真的告訴我們.
很多猶太人在這個時候.
真是很想了解一下.
這個死人復活是不是真確的呢?.
當然見到真確之後又會有什麼行動呢?.
對他們的內心有什麼的催促呢?.
這個才是問題所在.
因為你看到那些事情.
而你自己沒有作出相應的一個神學的判斷.
或者一個道德的抉擇.
所以所做的事情都跟你無關.
就是你看到一個人在你旁邊復活.
你可能都是無動於衷的.
這就是告訴我們當時的群眾就是這樣的.

$^{201}$接著第十節突然間出現這句話.
但祭司長雙爾連拉撒路也要殺了.
因為有好些猶太人為拉撒路的緣故.
回去信了耶穌.
這當然是第一站他們真是相信這個神蹟.
又相信耶穌有一個這麼大的能力.
就是說祂非比尋常的身份.
這是不是已經導致了他們的恩信稱義呢?.
因為他們認同有一個事實.
耶穌真是行了神蹟.
叫了這個拉撒路復活.
不是一種魔術.
但是看到他們自己要認信耶穌.
會不會有一個過程呢?.
因為現在看到耶穌是大能.
但是再過多六天.
耶穌在十字架上的時候.
不僅沒有能力彰顯.
甚至是一個無助者的白白犧牲.
而看不到有任何的能力等等.
所以這個就是猶太人.
如果不是用信心繼續踏上去.
跟隨,相信耶穌.
可能那個衝擊使得這個初步的認信.
都是會失落的.
所以很難怪很多人聽完報道會.
聽完很好的見證.
當下還沒有離開會場之前.
可能他們是信得很好.
很願意舉手決志.
甚至乎都表示認同.
但是當回到家裡.
回到世界 回到生活.
回到他自己的困難的時候.
如果他沒有相應的.
就著這位耶穌既然能夠幫助哪些人.
但是是不是真的幫助到我呢?.
如果他那一下決志再去跟隨.
那個信心就堅固了.
那就會更加去認定.

$^{241}$他是決心去跟隨耶穌.
因為由那個認同.
由那個相信到他自己個人的確認.
其實可能是有一個過程的.
有些人當然很快.
有些人可能慢慢就會流失了.
所以這裡再一次反映給我們知道.
就有這樣的情況.
所以你會看到第十二節.
就是如字的第五天.
耶穌就騎驢進城了.
就是這個終之主日.
耶穌第二天有許多上來過節的人.
聽見耶穌將要到耶路撒冷.
就拿著鐘樹枝出去迎接他.
接著還要喊叫和珊娜.
奉主名來的以色列王.
是應當稱頌的.
這段經文其實撒迦利亞書.
就是說那個謙和的君王騎著驢.
就是驢軀來進城的.
撒迦利亞書的九章九節.
就是這段經文來的.
所以第十四節.
耶穌得了一個驢驢仔.
就騎上如經上所記的說.
「錫安的民不要懼怕.
你的王騎著驢驢來了」.
第十六節.
「這些事門徒豈先不明白.
等到耶穌得了榮耀以後.
才想起這話是指著他寫的.
並且眾人果然向他這樣行了」.
所以約翰福音是事後的一個記錄.
是有神學的取向的.
是作者刻意地記載一些的事件.
然後他要解釋這些事件.
還有他要引用舊約的經文.
還要講回當時當下的現場.
包括作者自己都不明白那些事情.

$^{281}$覺得只是巧合.
只是一個大家都接受到的.
一個很美麗的場景.
因為拉薩路被拯救,復活.
家主擺筵席,有香膏的陶抹.
接著猶大只是一個插曲.
他提出一個善意.
為什麼不周濟窮人.
大家都說「是喎 是喎 是喎」.
相信連約翰都說「是喎 是喎」.
因為約翰是事後自己來贖罪,補救.
因為他還沒有看到耶穌得到榮耀之後.
包括聖靈的工作,聖靈的降臨.
聖靈的啟迪,讓道使徒們.
他們真是有一個很大的.
對整個主耶穌犧牲的十大事件.
他們有足夠的新心靈的力量去承載.
去消化,去整理.
所以約翰福音是一個事後的保屬.
真是將他整套神學告訴我們.
所以約翰福音憤外好看,憤外深奧.
就是因為他將他的神學告訴我們.
就是這位主是怎樣為我們展現.
他是神兒子的身份和能力.
所以這裡是告訴我們.
一切的事情原來是得到榮耀之後.
是主耶穌得到榮耀之後.
所以當下的現場的這個現席.
只不過是一個現場的現席.
猶大講得很合理.
這個門徒真是非常節儉,非常慳家.
非常愛窮人.
就是說跟隨耶穌的目的就是要愛窮人.
跟隨耶穌的目的就是不要浪費.
馬上衝口而出的對馬尼亞的一個.
未到去譴責 但是都是一種暗示.
為什麼你一個小女子會這樣浪費這麼多的錢.
如果你早點給我 我是管錢,管家來的.
那我去變賣了它 那就可以多些人受惠.
所以其它的門徒沒有一個出聲.

$^{321}$不過耶穌為他平反.
到底耶穌講的這句話.
有多少的門徒聽得到,聽得明白.
或者聽得明白之後又願意心悅誠服呢?.
因為不容易的 耶穌為什麼你會這樣講的?.
為什麼你會這樣來叛責被猶大的?.
猶大這個很有point的 為什麼你會這樣呢?.
為什麼要這樣來做一個對立呢?.
兩者是沒有衝突的.
你可以用其它的東西代替的.
用清水就可以了 解決了.
為什麼要用那個香膏呢?.
還要是好像是有牌子的珍拿達香膏.
一斤的貴重的珍拿達香膏.
所以沒有人知道背後馬尼亞的內心.
那個心跡 來到它想表達什麼.
所以耶穌這樣講的時候.
就真的相當重要和寶貴.
因為猶大如果它不是這樣講.
這件事就是這樣過了.
或者在猶大的內心.
充滿了一種覺得浪費,可惜.
但是當它這樣講出來.
耶穌把握這個現場的.
為馬尼亞來平反的一個機會.
不是要盲撐這個女孩子.
不是要來為她的浪費去辯護.
原來耶穌講了一個相當重要的事實.
這個告訴我們不是周濟窮人是錯 是不對.
以後因為約翰福音之後.
所有人都不可以去周濟窮人了.
不要浪費了 不是這樣的.
什麼都要做在教會,做在耶穌身上.
不是這樣的.
不過是那次的事件.
因為是一個關鍵的.
為耶穌安葬的一個這樣的舉措.
所以這裡是我們難以想像.
耶穌這樣來稱讚馬尼亞.
他有一種屬靈的眼光.

$^{361}$知道我的主愛我們的這位老師.
這位拉比.
這位現在在這一刻還是人盡皆愛的老師.
因為又能夠行神蹟.
又到處去坐直.
又不斷地引發很多人對他的議論.
很多人因為拉薩路又信耶穌.
甚至乎外邦人那些希臘人都要來找耶穌.
因為已經成為了當時耶路撒冷的一個很重要的焦點人物.
大家都不斷將那個好像壓住在耶穌身上.
祂就是我們的拯救者.
祂就是我們成為了我們的屬靈的領袖.
去推翻羅馬帝國.
讓我們不要再做他們的奴隸.
做他們那些仍然是受壓制的國家的人.
這就告訴我們.
不說猶自可 耶穌說了是一個關鍵的.
因為祂為我的安葬而存留.
安葬而存留 大家一聽就明白.
因為死了就必須盡快要來去包紮.
用香料包好 這是一個極重要的死亡的禮儀.
所以大家都明白.
但是對一個未死而這樣去進行的包紮.
進行的高抹是匪夷所思.
從來沒有這樣的人會這樣做.
就算是這個亞利馬泰的約瑟.
不過他自己也是買定很多.
所以買定棺材.
猶太人就買定那個墳墓.
一個用石造的墳墓.
當然是預備好多的香料.
那些果絲布等等.
但是都是要等死了才能做到的.
但是現在瑪利亞這麼出人意表的.
完全沒有其它的門徒能夠留意到.
他的內心原來是這樣為主.
就是說他很清楚耶穌曾經說過.
他真是很快就會受害.
就是第三天就會復活.
接著就會升天的了.

$^{401}$就是說真是越來越近.
耶穌基督要在這個節期.
被祭司長,文士所棄絕.
然後就被殺害等等的.
所以這個就是關鍵的把握的時機.
所以關鍵的把握時機.
我想我們人真是需要在信仰上要關鍵的把握.
我想現在在這個社會.
特別香港人這麼忙碌,這麼辛苦.
那個心靈有時候這麼困倦.
大家看一下過去那五天的長假就知道了.
差不多有兩,三成的晚上的晚飯.
是七點多已經沒有人坐了.
那你可以想象得到已經成為了一個常態了.
那你可以想象得到.
如果航空公司會說我限時出一些免費的飛.
大家去搶呢 相信那個把握時機.
很多人用很多部手機叫上姨媽,姑姐.
一起在那個時間一起按上網.
盡快,鬥快去搶 這個就是掌握那個時機.
是我們現代人其實是相當明白,相當了解的.
什麼的時機,什麼時候會派小學名校的入學申請等等.
我想我們是很清楚的.
寧願早一天去排隊 我們都要拿到頭籌.
以致可以有機會去接見 諸如此類.
我們很緊張 連那些派的福袋.
派的那些長洲的搶包山的一些包.
以前是真的你排隊 我才派給你.
盂蘭盛會派那些平安米.
又是要老人家去排隊 這樣來大家搶的.
早期的免費報紙 又是每天早上那些老人家.
又是這樣排隊 就是希望搶到.
那你可以想象得到 這個就是那個契機.
那個時機的時間.
所以原來在福音書裡面.
有一個這樣的女流之輩.
她不是什麼跟隨耶穌的使徒.
也都不是有什麼神學的裝備.
對神學有很深的認識.
很多的辯論等等.

$^{441}$但是她就是很清楚,認順.
耶穌自己親口說過 祂會在耶路撒冷受害.
被祭司長這樣去陷害,誣告.
然後會在彼拉多的手下這樣受難.
你可以想象得到 瑪麗亞記在心裡.
現在把握到一個機會 就是拉薩路重死復活.
所以似乎整個的復活 不僅僅是為瑪麗亞,馬大.
原來是為後來信的人.
所以耶穌為什麼要故意.
來到耶路撒冷停留三天.
才去伯大利.
就讓我們看到神的時間.
所以現在拉薩路的復活.
又襯托耶穌很快就要引來整個耶路撒冷的群眾.
必須要內心對發生的事情.
有一個道德的反思.
怎樣去理解耶穌和你個人的生命關係.
或者和這個民族的生命的關係呢?.
不是 你見到這些異象 見到你只是認同.
只是見到而已 沒有任何道德的對你的判斷.
就不能夠幫助你的生命的轉變.
所以這裡就讓我們看到.
瑪麗亞有一個生命的轉變.
深奧到一個地步 屬靈到一個地步.
是沒有人當場可以來去欣賞.
看到這個舉措原來是這麼屬靈.
原來是這麼偉大.
原來是抓緊著耶穌去告立這個死亡的那個告立.
所以完全沒有任何的門徒.
可以來到去有這種這麼有智慧的眼光.
所以是不是叫做屬靈人參透萬事呢?.
所謂參透萬事.
不是說他未卜先知 不是說他去樣樣東西.
他就樣樣皆能 相信不是這樣.
相信是對一個屬靈的眼光判斷.
是有一種神秘的 可能是心腸.
可能是智慧 可能是人生的閱歷的總和.
很多這些事情這樣來判斷.
譬如我剛剛上個星期日.
我在教會講復活節的崇拜.

$^{481}$講到我自己剛剛出來傳道的時候.
跟我很熟的團契的姐妹 她的妹妹.
那個時候她已經是中五了.
也都小時候一起回教會.
後來就慢慢的掉了 很少回來.
開始讀不成書 有時在家裡又跟姐姐的關係不好.
後來就知道她帶了一個男孩 有個男孩拍拖.
後來姐姐叫她帶回教會 給我們看一下.
我們一看就基本上都不用什麼屬靈判斷了.
都不需要有什麼屬靈人參透萬事.
因為她整個樣子是瘦猛猛的.
不是她的身體有什麼疾病 而是她是一個吸毒者.
已經一眼就看到是有不良嗜好的.
我們已經很快的跟這個妹妹說.
我們說 當然不只是我們 有傳道人.
叫她快點跟他分手 就是不聽.
跟著就更加不回教會 跟家人的關係更加差.
後來就跟了一個男孩結婚 應該是同居 不是結婚.
後來輾轉接著就生了一個兒子 叫做阿謙.
我上個星期就說過 其實我們每個人一看.
就已經知道你們是不會有前面的將來.
因為那個男孩根本上在那段時間就說.
我會戒毒 我是不會再這樣做了.
就真的找工作去做 但是慢慢的.
當結婚之後 開始又再出現那些問題.
最終是這位姐妹跟這個丈夫離婚.
因為她真的是遊手好閒 又繼續吸毒.
就令到她太痛苦了 所以你會看到.
我們說的這個是基本的判斷.
瑪利亞有一種屬靈的智慧 參透一件事.
就是主耶穌說過 祂會在耶路撒冷受害.
所以祂把握這次祂弟弟的筵席.
祂把握 是不是祂珍藏的一個嫁妝呢?.
有解經家就說 這些可能是猶太人的女孩子.
那些大姐的嫁妝來的 極其珍貴.
不只是價值 真是一個全家之寶.
一個這麼重要的 好像類似一些金的戒指.
金的手錶 爸爸的手錶 真是全家之寶.
可想而知 瑪利亞真的在找這個機會.
她不是想找其它的代替.

$^{521}$而是用她最珍貴的香膏就告滿耶穌基督.
只是幫她洗腳 一般人就覺得是奴隸做.
用水 不需要這麼浪費 猶太人突然間跳出來.
就不是她的愛心上腦爆棚 而是她內心.
約翰 因為對她的人生已經是足夠的概觀定論之後.
約翰在這一刻寫 就一點都不留情.
就將 其實不是約翰當時在現場揭發她的動機.
而是約翰自己經過了一段長時間.
這個本來是自己的童工的使徒.
自己最後就是這樣 因為出賣耶穌.
最終都是要自殺身亡.
所以對猶太人的描述 約翰這樣講.
《他說者話》第六節 是約翰神學的論述.
《他說者話》 並不是掛念窮人.
乃因他是個賊 又帶著錢囊.
尚取其中所存的 當時你又不講.
所以人就是沒有那個高度 沒有那個屬靈的成長.
還沒有去到那個階段 就是說暗示.
原來我們基督徒是不斷可以長大的.
我們的屬靈 不斷我們在神國裡面參與.
熱心服侍 我們的屬靈的眼光是會長大的.
如果你現在寫約翰福音 還要為猶太人去平反.
那就是說你根本就沒有那種屬靈的眼光了.
現在這樣講 就不是對那個死人不敬.
而是他用神學的論述 就說這個猶太人.
原來當時 即是原來當時 約翰自己都不知道.
原來當時他的動機是這樣的.
是他從來不是想著掛念窮人.
掛念窮人可能只是大衛這樣做而已.
他真正的動機是因為他帶著錢囊.
因為他是個賊來的 他會那個錢囊有多少.
我就拿多少 我拿剩一些 大家都不為意.
不會突然間拿光 這樣的來到中飽私囊的猶太.
所以約翰現在過了這麼長的時間.
為這個猶太人做一個好像無致命的描述.
約翰福音充斥對他的責備 對猶太人的責備.
第六章已經有了 耶穌從起頭就知道.
對不起 十三章 從起頭就知道.
誰人要出賣他 約翰就是不斷做這些神學反省的.
所以我仍然說 瑪利亞是一種極高的屬靈的智慧.

$^{561}$其實是不是因為她的內心實在是非常清.
對耶穌的愛是相當的單純.
這個真是需要我們一生要追求.
對這一位的救主 有不同的人會信耶穌.
有不同的入門的 年輕的小朋友什麼都不懂.
上教會很好玩 因為兒童崇拜很好玩.
又生日會 又唱歌 完全沒有壓力的.
又見到自己的好朋友 青年人上教會.
如果你是基督教學校 上團計就認識女孩子.
或者有時候會見到老師 好感又會好一點.
又會有一班人的聚會 本來這些都是健康的.
但總是會有這個動機 到你可能就是.
很想追那個女孩子 她不聽你的話 不理你.
可能你就不上團計了 成年人 家長.
因為香港是有計分制 基督教學校.
大家都覺得始終是好的 所以你就算不信耶穌.
你都很想你的子女進基督教學校 怎麼可以進呢?.
因為要有分 所以先報名 接著自己上教會.
接著就洗禮 洗禮就有10分了.
認識傳道人 牧師 寫推薦信.
又可能會再多一些分 你可以想像得到.
比如我們去錦瑟幼稚園 伯大利.
伯大利沒有什麼人需要寫推薦信的.
反而叫我們寫另一間學校的推薦信.
孫德的徐澤林 我任內 包括張牧的任內.
都寫了很多的 你就可以想像得到.
本來這些不是錯的 但是那個入門.
你接觸到福音 認識了耶穌基督.
所以很重要的就是 基督教最重要的是耶穌基督.
最重要的不是基督教裡面的副產品.
有些人很喜歡回屯契的 他不去崇拜都回屯契的.
他搬了很遠 他不去教會崇拜 他都回來屯契.
本來也不是錯的 也好 但是他就是在那個屯契.
到後來他連他自己本來應該回去崇拜的都不回了.
都不再崇拜了 不再敬拜神 但是他喜歡屯契.
因為那幫人認識了幾十年.
那就慢慢 基督教最重要的原來是屯契.
基督教最重要的就是可以彼此相愛.
彼此相愛就是小撮那幫自己談得來的人.
所以其實基督教的重點是基督.

$^{601}$基督才是整個基督教的核心.
所以馬寧他在這一刻知道耶穌基督.
是他生命裡面最重要 不怕別人嘲笑他.
不怕一個浪費 也不怕做在主耶穌身上.
會得到什麼樣的回報 因為最大的.
最單純的愛做在主耶穌身上.
我想大家都去到一定年紀.
我們應該開始有很多的東西.
我們是要給予別人的 因為沒什麼用了.
開始用不了那麼多了 但是我們的內心最深處.
是不是一份單純對耶穌的愛呢?.
因為子女就開始會遠離我們.
有他們自己 應該是要這樣.
要有他們自己的獨立生活.
我們現在唯一的 都是可能最基本的.
我們已經很滿足 但是最深處.
真是唯有基督是我們的生命 是我們的一切.
這個是我們一生要追求.
而且永遠是會屬於我們自己.
所以真是求神幫助我們 我們看到馬寧.
所以這段聖經很多人講的.
這段聖經也相當多人熟悉的.
這段聖經我們就這樣隨便揭一揭.
已經看到有三類人 猶太人 看到拉撒路復活.
知道拉撒路復活 刻意整班人上來.
跟隨耶穌去看一下 哇 果然耶穌你真是厲害.
是我們猶太人的領袖 我們就要跟隨你了.
所以第二天當耶穌要準備.
其實是應驗撒迦利亞書騎驢進城.
所以那些人 原來是那些猶太人.
是他們自發的 不是那些請回來的臨時演員.
是他們很自發的 要去拿那些棕樹枝.
來鋪在地上給耶穌過 應驗了撒迦利亞書.
稍微揭一揭 撒迦利亞書的九章.
第九節 錫安的民應當大大喜樂.
耶路撒冷的民應當歡呼.
看啊 你的王來到你這裡 他是公義的.
並且施行拯救 謙謙和和的騎著驢.
就是騎著驢的驅子 我必除滅以法年的戰車.
特別是強調那個威武 類似驕傲.

$^{641}$和耶路撒冷的戰馬 爭戰的弓也必除滅.
他必向列國講和平.
他的權柄必從這海管到那海.
從大河管到地極.
接著下面我們就不詳細再讀了.
都有一點點的關聯的.
告訴我們 原來主耶穌是以祂的犧牲和平.
是用祂的謙和來得勝.
這些任何一個猶太人都不會明白.
或者今天我們未信耶穌的人都不會明白.
明明是一個這麼軟弱的死囚.
說自己講到自己是偉大的救主.
人類的救主 世界的真神.
講到這麼偉大.
但是聖經就是這樣這麼平鋪直敘的預言給我們知道.
撒迦利亞書就是主前大概四百多年.
就已經告訴我們.
我們的主耶穌騎驢進城.
就是謙謙和和的用驢驅進城.
是你的王 所以是和平之君.
所以基督教是在講和平 用愛 用和平去統治.
所以它要除滅以法言的戰車.
耶路撒冷的戰馬.
所以不是靠武力.
所以回來的時候.
這些群眾 原來他們的愛主.
就是將那些鐘樹枝擺在路上.
一個很盛大的儀式.
歡迎我們的君王來降臨.
這些其實是很政治敏感的.
因為羅馬皇帝要去巡行.
巡視一些地方 他所攻略的城池.
那些皇帝去到的時候.
就是這樣的一個歡迎的儀式.
雖然不是用橄欖的樹枝.
可能他們用香花來讓那些君王降臨.
但是現在是耶穌要這樣來進城.
所以你看到耶穌是用驢驅.
是不會有殺傷力 是謙和的.
不是穿上盔甲.

$^{681}$不是穿上皇族的漂亮的紙袍.
而是很普通的一個主耶穌.
平日的主耶穌.
但是對一大群的群眾.
但是他們的愛是不能夠持久的.
他們的愛不是改變他們的生命.
他們是見到有好處.
見到耶穌現在有這麼大的受歡迎.
所以就愛耶穌 所以就跟隨耶穌.
差不多是沒有成本.
到要付代價 要有成本.
他們就會寫耶穌而去.
他們就會在十字架下面.
聲討耶穌的群眾.
所以第一類的人特別多.
所以回來教會看熱鬧.
哇 這裡有一個很特別的崇拜.
很特別的子爺的聖餐.
又有聖誕節,壽苦節.
又有平安夜.
特別多人忽然之間有那個靈性.
很想滿足一下靈性.
不是不好的 但是要轉化.
很多人拋過橄欖枝的人.
棕樹枝的群眾.
又是攻擊耶穌的一幫群眾.
所以耶穌什麼時候一轉變了呢?.
應該說在你的內心 耶穌轉變了.
就會要你付代價的時候.
就是你要對你的神學有一個反思.
對你的生命有一個反思.
這樣來作出一個回應.
所以群眾的回應很自然.
我已經沒有跟耶穌得到利益.
我就只能夠就著我自己本身的.
難道我會歸附祂來抗衡羅馬帝國嗎?.
難道我要跟阿彼拉多.
跟當時的宗教領袖為敵嗎?.
我當然要站在這一幫人當中.
我怎麼會站在耶穌這一邊呢?.

$^{721}$這是第一類的作者告訴我們當時的那個場景.
第二類就是原來已經是耶穌呼召的猶大了.
原來祂的愛主 原來祂的內心.
一直都被人感覺到祂是很掛念窮人.
因為祂是負責派錢的那個.
一直都是跟隨大衛.
而且更恐怖的就是12個的使徒.
都沒有一個能夠所謂的明察秋毫.
為眼看到猶大的內心的那個叛變.
內心裡面的對主的不信.
和祂自己那個身份 那個的使徒.
祂都覺得是要有一個回報.
所以祂面對重大.
因為祂是更加給群眾更加了解主耶穌的為人.
更加日夕和耶穌和11個的同胞,同伴.
來跟隨耶穌的.
應該給其他群眾對耶穌有更深的了解,認識的.
但是奈何祂內心的那個賊性.
內心裡面的貪婪.
都挺多對祂的描述的 約翰福音.
我們之前好像有講過的.
第六章就有提到猶大.
第六章的70節 耶穌說.
「我不是揀選了你們12個門徒嗎?.
但你們中間有一個是魔鬼.
耶穌說:是指著加略人西門的兒子猶大說的.
他本是12個門徒裡的一個.
後來要賣耶穌的」.
我相信第71節 不是耶穌當下馬上揭穿.
如果揭穿就不會有這麼多後來的.
還有一年多 因為這個時候是五餅二魚.
還有一年半的時間.
所以第71節似乎是約翰幫耶穌做一個詮釋.
所以約翰真是現在這一刻就看得很清楚.
所以第二類的愛主.
有一大幫的門徒跟隨耶穌都沒有出聲.
但是有其中一個猶大就出聲.
對於這個事件 對於一個小女子.
用這麼重價值的香膏來高抹耶穌.
看不到是為耶穌安葬 看不到背後的心跡.

$^{761}$也不去提問一下 這樣做也挺特別的.
有沒有特別的意思呢?都沒有了.
都是馬上衝口而出 突然之間.
非常義正嚴詞的,嚴肅的 為耶穌這個團隊辯護.
實際是為他自己內心的賊性.
所以這是第二類的愛主.
不僅自己得不到生命的成長.
而是原來裡面雖然跟隨耶穌.
所以有一段聖經說 本是十二使徒的一個.
後來要出賣耶穌 本是 就是說耶穌真是選擇他的.
明知他是有這樣的內心都選擇他.
但是真是他自己選擇了這條黑暗的路.
所以第二類的愛主.
原來是沒有自己不斷的生命的扭轉.
生命的更新 所以弟兄姊妹.
信主原來生命的更新就更加難了.
真是要自己有一個很大的反思.
真是不斷要在神的裡面.
透過不同的人事物給自己有很多的反省.
如果不是 有時就是越來越藏起自己.
那個內心就會落在一種越來越的狹隘.
將別人的說話都看為是攻擊.
沒有了一個反思的能力.
因為耶穌是一而再的來勸導他.
他都聽不入耳.
因為他內心根本不肯面對自己的貪婪.
就帶來一個這樣的情況 其實也不是很難.
你真是覺得穩食艱難 跟隨耶穌都是不行的.
那你就直接跟耶穌說 哎呀 還是不要了 我還是回去.
因為始終我要養妻活兒.
我都要做回我自己的本行.
或者去在世界裡面有一些更加好的謀生的工作.
因為跟隨耶穌 找朝紅得晚 怎麼辦呢?.
家裡人又有媚艷 可以跟耶穌講.
將你內心 但是不是 以這個職位成為.
可以更加去為自己謀取更多的金錢.
而耶穌不是助長他 縱容他 而是給他機會.
這個微妙之處 都是一種藝術.
這個就是耶穌跟猶大之間的三年多的相處.
但是第二類的門徒 第二類的愛主.

$^{801}$就已經沒有群眾那些這麼路人甲的了.
就是一見到開心就很歡喜.
一見到不開心就沒有興趣.
就有很多這樣的情況.
就是如果平時的聚會 平時的講完.
平時的那個setting.
大家都不會感覺有什麼特別的興趣.
只是慣性的聚會這樣.
但是如果一個慣性的聚會.
我們的心都有準備好.
那個很慣性的聚會.
我們自己都會有很大的得著.
我們自己都會有一個很大的滿足.
因為我們每一次都不要管是不是特別.
我們自己都為自己預備一個特別的空間.
讓耶穌來進入.
讓我們帶著一個敬拜朝聖的心去敬拜神.
這樣每一個習慣 習以為常都可以是特別.
相反 多特別 死人復活.
都過了一個禮拜之後.
這幫群眾曾經 可惜當時沒有相機.
拍了這幫人的相片 馬上直播給他們看.
一個禮拜前你們還是這樣來到開心的和珊娜.
奉主明來的 是應當稱頌的 還是引聖經的.
你就可以想象得到 這個就是第一類的愛主.
但是猶大是第二類的愛主.
猶大就是為自己的內心沒有真正的跟隨耶穌.
其實它從來沒有愛主.
從來沒有在基督的救恩裡面有份.
它在基督教的事業上有份 有它名字的州制.
有它很多的足跡 去過不同的地方.
來派錢給窮人有東西吃.
但是它只是在基督教的事業上 曾經很熱心.
從來都沒有愛過主 這個就是猶大第二類.
所以其它的門徒是一樣的.
當時連約翰都是一樣的 但是他們會成長.
會悔改 會不斷的有聖靈工作.
所以約翰現在在晚年的時候寫的.
這個約翰福音是充滿神學的反省.
那你可以想像得到 原來第二類的門徒.

$^{841}$其實已經很好了 很願意跟隨耶穌的了.
但是都是要不斷我們內心.
一定要問清楚我們的動機.
所以整個獻香膏是動機的問題.
不是價值的問題.
所以耶穌都稱讚那個兩個小錢的窮寡婦.
路加福音23章還是21章.
那個窮寡婦奉獻了兩個小錢.
那些門徒就看不起了 但是耶穌叫那些門徒來.
你看看這個婦人 她將她僅有的.
完全的奉獻給神.
所以來得到主耶穌的讚賞.
你可以想像得到 原來價值背後的那個動機.
是最重要的 所以我們一生.
神都會不斷的考驗我們內心的動機.
所以來到我們這個階段 子女都長大了.
其實我們更加要追求最單純的動機.
真是那個事奉是越來越很純全的.
是要土主的喜悅 真是已經可以來到去.
在神的角度裡面 很想是有份.
能夠明智被記在生命冊.
就是最大的一個福氣 所以為什麼這麼深刻呢?.
就是有一位很出色的.
就是福音派很重要的神學家.
這一位 我突然間這一刻忘記了名字.
他有一天去英國 去到Westminster的教會.
去找到鍾馬田牧師 一位鍾馬田醫生.
Dr. Martyn Loyal-Jones.
很出名的60年代的一位神的僕人.
在英國的教會復興.
他的教會在倫敦相當重要的一個教會.
每一個星期日的崇拜 應該是晚上還是什麼.
都是講羅馬書 又是這樣verse by verse的.
講了12年 講羅馬書 每次講一節.
真是激勵了很多很多英國的信徒.
去到晚年的時候 這一位神的僕人.
年輕的一位神學家 去探望鍾馬田牧師.
問了他一個問題.
就是說現在你曾經在英語世界這麼出名.
帶來這麼大的復興 幫助了這麼多的人.

$^{881}$他的著作等等 現在你只能夠好像在一個陡塞.
什麼都做不了 你會怎麼樣呢?.
現在好像類似都被人忘記了.
鍾馬田牧師說 因為他去探望他的時候.
鍾馬田牧師說 我已經為你有很多的禱告.
其實不認識的 因為知道這一位年輕的神學家.
又是被神大用 在Trinity的一位神學家.
很重要的神學家 最終還是忘記了他的名字.
今晚給大家看一下他的名字 我很敬重.
我買了很多他的書去看 相當有insight的.
這一位神學家 鍾馬田牧師說.
我已經為你有很多的禱告 為你的服侍等等.
這個神學家就很詫異 因為自己是後輩.
根本都不認識鍾馬田牧師.
為什麼他會認識我 為我祈禱呢?.
為我的事奉祈禱呢? 真是很屬靈的長輩.
有一種這樣的眼光 知道 可能聽過的講道等等.
知道這一位年輕人被神用 就為他祈禱.
因為他在英國 這一位神學家在美國.
他說很難忘 鍾馬田牧師說了一句話.
因為是他問他 你會不會覺得有一些的失落.
你曾經那麼風光 那麼被神用.
現在你什麼都做不了了 好像一個很普通的人.
慢慢慢慢 一路一路走下坡了.
鍾馬田牧師引述了一句聖經.
到現在我還是很深刻 很深刻.
路加芬第十章引述那句話.
就是耶穌差遣70個人出去傳道.
兩個兩個回來 就興高采烈.
哇 耶穌啊 因為你的名字連鬼都服了我們.
他們很厲害 狐假虎威 很興奮.
一對一對回來 70個 就是有35對.
在那裡大家 你沒有我這麼厲害 我更厲害.
我更厲害 大家都很熱鬧.
耶穌講了一句話.
不要因為鬼服了你們的歡喜.
要為你們的名字記錄在天上歡喜.
鍾馬田牧師回應這位年輕的神學家的這句話.
說令到他相當深刻.
就是說 我是很開心 我現在是被人淡忘了.

$^{921}$我的事奉已經完結在地上了.
我等候的就是我要去天上見主.
弟兄姊妹 是不是我們的渴慕呢?.
我們的牧師的身份越來越感到尊貴.
真是幫很多人進入永恆.
那種感覺是越來越好像強烈.
就是永恆的真實 我們好像摸得到.
我們在那個聚會的儀式裡面.
不斷去將這份感覺 自己的感動.
去傳遞給那些哀傷的家庭.
讓他們都希望感覺到是真實的.
愛我們的家人只是暫時的離開.
在天下是會見的 是會永恆的重聚的.
越來越感受到這個事奉帶給我自己的一份尊貴.
一份福氣 所以這句話成為了我自己.
因為這位牧者的這樣講解.
成為了對他就是一生的影響.
對我第三者看到這個描述.
也是令我成為了我很重要的座右銘.
就是我現在事奉的動力.
行不行不要緊 盡了力 不要因為鬼服了你.
這次服了你 下次可能就不行了.
下次可能很艱難了 但是不要緊.
盡了力 做了就行了 問心無愧.
但是更寶貴更重要的是我們事奉的人.
將來在天下見主 所以第三類愛主的.
其實就是瑪利亞 第三類不是做很多的事情.
表面上都沒有什麼人知道.
原來他這麼深的肅靈的眼光.
對耶穌的那份擺上 但是耶穌看到了.
耶穌看到了 當然不是每一個人.
耶穌都馬上告訴我們.
這個人其實內心是這麼高尚,這麼肅靈.
很多人都未必看到的 但是我們要默默去做.
我們要繼續不斷 這個階段就是應該.
那個動機越來越純全.
單純為主,愛主,為教會 我們去擺上.
所以我們求神幫助.
用了一個小時講了瑪利亞這一段.
不好意思 都是有感動要講完的.

$^{961}$就是格羅伯的妻子 就只有幾段.
主要是《約翰福音》這一段.
所以我們可以再看這裡 19章.
19章 這個就是耶穌被釘的時候.
我們可能不要看太多了.
19章的第23節開始.
丁丁既然將耶穌釘在十字架上.
就拿他的衣服分為四份.
每丁一份 又拿他的女衣.
這件女衣原來沒有縫兒 是上下一片織成的.
他們就彼此說 我們不要撕開.
只要捏扣看誰得著 只要應驗經上的話說.
他們分了我的外衣 為我的女衣捏扣.
這是詩篇的22篇 18節的經文.
就是那個受苦的僕人.
第25節 丁丁果然作了這事.
丁丁又是完全不為意 原來自己在應驗詩篇.
就是在那個現場 當下的智慧.
一件到底的衣服 你撕爛它 分就沒用了.
不如大家捏扣 抽籤.
誰得到就一整件衣服.
一個死囚的衣服 是不是當時的物質這麼窮呢?.
一個死囚穿著的衣服 都覺得是貴重.
是不是因為他的質料呢?.
是不是因為他的身份 耶穌穿過.
所以真的要穿 等於世界盃的冠軍.
如果在朗拿道穿過的那件衣服 給了我.
如果我是喜歡巴西的 那就不同了.
那件衣服 我不喜歡穿.
這個在朗拿道穿過 世界盃拿冠軍的.
那就完全不同了 可能就是這樣也說不定.
但是原來應驗聖經 主耶穌死的時候.
真是赤裸裸 連祂僅有的一件身體的一件.
衣服本身就是遮蓋羞恥.
就是說耶穌基督 我們的救主.
祂為我們受盡羞恥凌辱.
所以我們的罪 現在是在永恆裡面.
已經再沒有那個追討 因為耶穌已經擔待了.
所以你可以想象得到 主耶穌的徹底的痛苦.
連這些都不會隱藏.

$^{1001}$第25節 這裡才寶貴.
站在耶穌十字架旁邊的.
與他母親的姐妹.
並格羅巴的妻子瑪利亞.
和莫大拉的瑪利亞.
這裡總共有三個瑪利亞.
耶穌的媽媽瑪利亞 格羅巴的妻子瑪利亞.
格羅巴 有些人就說是路加24章.
去驗馬五師的那個格羅巴 就是他太太.
有可能是 但未必是 也說不定的.
耶穌見母親和他所愛的門徒站著.
接著就托付給約翰.
但是這裡就已經告訴我們有四個的姐妹.
在現場 而且那個現場是相當之恐怖.
血腥殘忍的現場 但是作者約翰.
因為他也是在現場 陪伴著這四個的姐妹.
就是說 姐妹那個的耐力.
那種的內心的純潔.
為了一個真理.
他們那個感情完全擺上去.
有時候好像都忘記了那個案發現場的危機.
危險 好像都置危險不顧.
其實你的生命危險 因為你是和耶穌有關.
但是因為是女流之輩.
那些丁丁都沒有對他有什麼下手.
可能當是群眾.
因為最主要的那些門徒都走了.
所以就沒有刻意去針對這幾個姐妹.
但是從開始 耶穌出來傳道.
一直就見證著 似乎這四個姐妹.
有些是現在才出名的.
就是這個格羅巴的妻子瑪利亞.
這樣來和耶穌的媽媽.
和莫大拉的瑪利亞 耶穌媽媽的姐妹.
一起這樣去陪伴耶穌.
所以單是這樣.
就是耶穌原來是默默的忠誠的支持者,跟隨者.
又是勇敢的見證人.
所以女士在最需要的時候總是會出現.
往往真的是的 姐妹,媽媽的角色.

$^{1041}$在地球上真的是很重要的.
很多的男士,很多的男孩子成長.
都一定有媽媽的協助.
往往爸爸是缺席的 香港都不算嚴重.
美國就相當嚴重了.
因為那些罪案是非常之嚴重.
尤其是黑人的社會,社區.
布魯克林等等 很多的青少年的黑人.
就已經要麼坐牢 要麼就槍戰死了.
所以很多時候都是fatherless的家庭.
所以都是由媽媽養大的那些子女,那些男孩子.
所以對媽媽的感情非常深厚.
但是慢慢如果讀不了書 跳不出那個貧窮.
那個魔咒 又是被毒品,被很多的幫派.
這樣來去操控 又是走上爸爸的路.
又是要麼坐牢 在監獄裡面.
要麼很快就死等等 所以這些問題.
美國其實是不斷出現 是fatherless.
但是媽媽的責任就很大.
所以如果媽媽是敬虔 真是帶來很大的影響.
很多出名的傳道人 都是他的媽媽是很敬虔.
所以這些都是我們值得去思想的.
那是不是真是耶穌的親屬呢?.
那又不敢講 有一些解經家敢講.
所以不能夠百分之一百講是.
所以今晚就只能夠去到這裡.
我們下一次就繼續下去.
希望我能不能夠下一次講完.
因為都有很多張PowerPoint都未必夠.
希望我們盡快完成 我們就做一個總結.
看看什麼時候我們就開始出埃及記了.
因為一直去準備的時候.
都發現有很多東西值得講的.
可能有些大家覺得很悶的 我們可能快一點.
點出重點就算了 今晚我們就停在這裡.
天父啊 我們看到瑪利亞這一份愛主的心.
她能夠在逾越節受難的時候的前六天.
大家還是沉醉在耶穌能夠呼風喚雨行神蹟的時候.
她已經有這一份眼光 不知道是不是有一份內心的痛.
她要為主耶穌準備 好像準備後事那樣.

$^{1081}$就是要告末耶穌.
當然這些香膏的香氣不能夠去到耶穌死的時候.
但是那個背後的屬靈的含義.
絕對不是一個物質的那個象徵.
這份屬靈的含義 就是瑪利亞非常清楚.
主耶穌來這個世界就是要拯救罪人.
多謝主讓我們看到你為瑪利亞所講的.
為你的安葬所存留 多謝主讓我們在地上.
我現在或者我們一班弟兄姊妹.
我們都是每年每月不同的節期,不同的日子.
或者很尋常的平日 我們都是這樣跟隨.
我們都是這樣在地上生活.
但是如果我們有瑪利亞的心.
我們的人生就會很不同.
我們能夠在平常的日子.
我們都會知道主耶穌的同在.
就是給我們最大的安慰,幫助和力量.
以致我們能夠無論尋常與不尋常的日子.
我們都是這樣始終如一的單純的愛你,跟隨你.
為你的國度我們能夠擺上的.
以致當我們走完了人生的路.
我們是帶著一份平靜,安穩的心.
等待那個永恆的黎明.
好像宗馬田牧師所引用的經文.
主耶穌所講 不要因為鬼.
不要因為我們的成功而我們沾沾自喜.
而是因為我們的名字能夠在恩典裡面.
被記錄在生命冊而來得歡喜.
這個是非常大的 真是永恆的智慧.
是來自我們的主對當時的70個門徒講.
不斷有所得著 我們都能夠有啟發,有反思.
重點是我們能夠不斷透過神的話語去反思.
使我們的生命更加的趨向為主而活.
我們這樣感恩禱告 奉耶穌基督的名 阿們.
\newpage



\section{羅馬書}
\label{sec:WmefrJ6J_Lg}
\textbf{2026-05-08 周五晚查經班《羅馬書(四十二):總論:福音帶來與神和好,所以教會必須向萬民宣講福音.》 講員:曾裔貴牧師}
\newline
\newline
連結: \href{https://youtube.com/watch?v=WmefrJ6J-Lg}{\texttt{https://youtube.com/watch?v=WmefrJ6J-Lg}} ~~~~ 語音日期: 2026-05-10
\newline
\newline
\hyperref[sec:Ex_WASW50ZA]{\small{< < < PREV SERMON < < <}}
~
\hyperref[sec:index_chronic]{\small{[返順時目]}}
~
\hyperref[sec:index_scriptual]{\small{[返順卷目]}}
~
\hyperref[sec:GROoST1kr4I]{\small{> > > NEXT SERMON > > >}}
\newline
\newline
$^{1}$我們一起同心禱告.
因天賦給我們又可以在今晚.
來到去有最後一講的《羅馬書》.
都給我們有42講 能夠同鼎姐妹一同.
從第一章一直到第16章.
真是看到使徒保羅是有系統的.
這樣去將救恩的原委.
為什麼需要福音 講得很淋漓盡致.
也讓我們能夠看到我們成為了基督徒之後.
是有一個很榮耀的身份.
聖靈的內著和神兒女.
神的後製這個身份是何等的尊貴.
所以在地上的人如果不認識耶穌.
人生無論是多麼在地上有成就.
當然我們不是抹殺這些成就,這些努力.
但是在永恆裡面是一無是處.
是非常之可惜 所以現在那個宣講的責任.
就是落在教會和每一個時代的弟兄姊妹的身上.
求神幫助我們 因為福音不能夠自己去傳開.
必須是屬主的兒女有意識地.
甚至乎是要有那個承擔.
來將福音傳到最需要的人 求神幫助.
今天晚上我們有這個總論.
願神讓我們在大概一個多小時的時間.
都可以有一個溫習 願神帶領.
下個星期當我們開新的《出埃及記》的時候.
又讓我們能夠從40章的經文.
和在整個摩西五經的脈絡裡面.
我們能夠梳理到摩西或者五經的作者.
呈現給我們認識的那個出埃及的時代.
神怎樣去揀選你的選民 猶太人.
以色列人成為國度 願神帶領我們.
我們今晚仰望等候在你的面前.
奉耶穌基督的名 阿們.
首先多謝各位參加一段長時間的羅馬書研讀.
抱歉今晚有一點點喉嚨不舒服.
所以可能聲線不能夠這麼大.
希望都能夠 因為我預備了很多張PowerPoint.
我希望都能夠要快一點的去.
這個題目都差不多是必須要這個總題.

$^{41}$真的告訴我們福音是來自與神的和好.
來自耶穌基督 神所差遣.
從死復活的耶穌基督.
我們就能夠可以透過主耶穌.
就得著這個永恆生命.
這個永恆生命就是神的義.
與神和好的這個新生命.
所以既然這個這麼重要的信息.
保羅在一個整個大的脈絡裡面.
絕對不是將它分了很少的段落.
而是一個很清楚的一個大的脈絡.
那個總題就是教會.
特別以羅馬修信的教會為那個首先的單位.
必須要向萬民宣講福音.
因為萬民都需要救恩.
就是每一個民族都需要神的福音.
所以我們今天就在這個總題裡面.
等一下跟大家再做一個總結.
我們一直在過去用了好幾堂的時間.
去講到16章1到第16節.
是很特別的所有保羅書信13卷的保羅書信.
都沒有這麼長的文案.
在這個文案裡面我們梳理了.
雖然那些人物我們大部份都不認識.
但是我們又梳理到在古代羅馬的時候.
保羅的時代有些名字.
根據一些解經家他們的研究.
有些是貴族的名字.
有些是奴隸 明顯是奴隸的名字.
更加特別的就是有男的 不出奇.
但是特別的就是有女的.
而且不止一個在保羅的文案裡面.
要不就是他們可以在教會裡面.
有相當顯赫的地位 例如妃婢.
另外有其它的 譬如瑪麗娜.
我們之前都講過了.
所以等一下我們看這些PowerPoint的時候.
我們就知道 所以如果我們一個宏觀的看.
保羅這樣大量的文案 接近26位.
因為有一個是夫婦這樣文案.

$^{81}$就讓我們看到保羅似乎在羅馬的教會.
認識相當多的弟兄姊妹.
而且這些不是跟他經常一起同宮的.
宣教旅程的那些同宮.
而是在羅馬教會安頓.
在意大利教會裡面來服侍的.
這樣就帶出了保羅其實.
雖然沒有去到羅馬建立教會.
但是他都認識來自不同的地方.
最後落戶在意大利羅馬的教會.
這樣來服侍 所以這樣的文案.
以及準備他自己很快要去到羅馬.
跟他們有商交.
將一些肅靈的恩賜分給他們的時候.
接著保羅就說要蒙他們的頌恆.
這個16章後面的經文.
接著他就會去意大利之後.
就要去到西班牙那裡去傳道.
所以你可以想象得到.
在歐洲的意大利都是中部.
一直下去到南面 再南面 歐洲最南面.
就是西班牙 葡萄牙的地方.
你可以想象得到.
保羅進入東歐 保加利亞.
就是菲律賓那一帶.
馬其頓 就是現在的希臘等等的地方.
接著本來就要上北一點的.
高加索的那一帶 加拿大等等的地方.
但是神就帶領他一直由東歐.
接著就過去 最終去到哥林多.
接著就再回去以弗所.
就是過去土耳其那一邊.
所以當時他的足跡都是在東歐.
和中歐的地方.
但是後來輾轉 當然北歐就好像沒怎麼提他.
可能那個時候北歐都還沒有太開發.
所以最後保羅在他的心願裡面.
就是想去西班牙宣教.
就是說在他能夠用腳.
當然也是坐船去環島的地中海一帶的地方.

$^{121}$都是保羅向這些外邦人傳道的足跡.
所以他能夠認識這些弟兄姊妹.
更加讓我們看到保羅的胸襟.
保羅的神學觀.
保羅他貴為一個這麼重要.
他相當的時間是一個拉比.
拉比就是說那個嚴謹的訓練.
嚴謹訓練自己嚴謹.
有時候就慢慢越來越變成很原教旨.
當自己要求自己很高的時候.
那個人就不知不覺.
好像生人勿近那樣 不苟言笑.
還有什麼都很快就上綱上線.
對不對真理 有沒有違反神的旨意.
很容易不單止是對自己.
尤其是更加容易是對其他人.
所以保羅能夠這樣去對女性的民安.
對奴隸的民安 對外邦人的民安.
但是這樣已經讓我們看到.
這位使徒他的生命.
被主耶穌的救恩改變到一個地步.
就是裡裡外外都是被神的愛充滿.
既保有自己對真理的認識.
但同時間對其他人.
對一些事情的那個新的在救恩裡面的眼光.
這樣來去開闊就真是不簡單.
簡直是世外高人.
就是能夠去到一個地步融會貫通.
救恩如果我們不能夠融會貫通.
就很大的問題.
有時候對自己過度的苛刻.
有時候對別人過度的要求.
都是帶來有時候拆毀性的破壞.
我所指的是關係性的破壞.
所以我們看到保羅就是在這些民安裡面.
這麼簡單的一個問候.
都引發出我們可以有這麼多的學習.
當然我們一而再講就是.
在26個民安裡面17節都是要再讀多一次.
你可以知道那個修辭那個經文.

$^{161}$帶給我們的就是民安裡面都要留心.
如果關乎到真理,關乎到錯謀,關乎到異端.
這個就跟民安本來沒有抵觸.
但是就要警告就要去留心.
17節16章的17節.
「弟兄們!那些離間你們.
叫你們跌倒背負所學之道的人.
我勸你們要留意躲避他們.
因為這樣的人不服事我們的主基督.
只服事自己的肚腹.
用花言巧語誘惑那些老實人的心.
你們的信服已經傳於眾人.
所以我為你們歡喜.
但我願意你們在善上聰明.
在惡上愚拙.
賜平安的神快要將撒旦踐踏在你們腳下.
願我主耶穌基督的恩.
常和你們同在」.
那你就可以想像得到.
這一段聖經是相當的嚴肅的警告.
要羅馬教會當然可能是帶領的一些領袖.
就要留心就是哪些人會帶來有離間的.
就是那個心靈或者是一些真理上面.
以致叫那些弟兄姊妹原本是很好的.
因為保羅都稱讚他們.
他們是非常之信服的.
老實的人 而且是眾人皆知的.
所以保羅都為他們去感恩.
為他們歡喜.
希望他們繼續可以在善上聰明.
在惡上愚拙.
當然這樣的修辭就是告訴我們.
要向善不要向惡.
要學習一些好的榜樣.
一些好的習慣.
一些帶來我們好像耶穌這樣的身份.
這樣的氣質的基督徒.
一些不好的東西.
例如是是非非.
例如中傷,挑撥,離間.

$^{201}$或者是講一些錯誤的道理.
或者在教會裡面散播一些謠言.
這些就是在惡事上面要愚拙.
相信我們不是突然間蠢了.
不懂得說.
因為我們的心.
亞國書告訴我們.
我們的心又可以出.
揭開給大家看.
所以我們要留意.
在我們的心靈裡面.
其實要常常地去小心.
第三章很有趣的描述的形容詞.
它說 這是第五節.
它說 舌頭在白體裡也是最小.
卻能夠說大話.
看下最小的火能夠點著最大的樹林.
舌頭就是火.
在我們白體中.
舌頭是個罪惡的世界.
能夠污穢全身.
也能夠把生命的輪子點起來.
就是說 端扶你怎麼用我們的心靈.
用我們的舌頭.
當然不是舌頭.
是我們整個思考的機制.
我們的思考.
我們內心裡面的動機.
我們對人的態度.
我們很多裡面的運作.
如果我們裡面是奸險的小人.
我們說的話肯定不會是道德仁義的.
如果我們裡面是一個很坦蕩蕩的基督徒.
有真理 有聖潔 有仁義.
我們在我們跟人的相處裡面.
我們就能夠帶出主耶穌的那份恩情.
所以作者雅各也告訴我們.
原來我們的舌頭是可以駕馭的.
如果我們不知道.
那就會帶來非常大的毀滅.

$^{241}$並且是從地獄裡面點著.
各類的走獸,飛禽,昆蟲,水族.
本來都可以制服.
也已經被人制服.
就是很形象化的描述.
但是唯獨舌頭沒有人能夠制服.
也在地球上沒有人會覺得舌頭是最大的惡獸.
我們中國人也有一些周處除三害.
最大的禍害是他自己管人.
我們都知道.
所以聖經真的很奇妙.
告訴我們原來每個人的心靈.
在主耶穌的救恩裡面.
我們還需要不斷地學習.
學習向善.
所以這裡真是呼應保羅所說.
羅馬教會就是要在神的恩典裡面.
既然已經成為了基督徒.
罪惡的泉源已經被聖靈盡著光照.
以及出死入生.
以及有了聖靈在我們裡面.
已經來印證我們是神的兒女.
所以我們就要在這個清潔的泉源裡面.
不斷地讓自己的生命成為聖潔.
所以他說唯獨舌頭沒有人能夠制服.
是不止息的惡物.
滿了害死人的毒氣.
就是說用一些舌頭的.
無論是說謊.
惡毒的言語.
有轟.
或者是有陰謀.
陰險.
或者是一些處心積慮的.
要對教會帶來的.
或者對世界帶來的破壞.
這裡說我們用舌頭頌讚那為主為父.
又用舌頭咒助那照著神形象被造的人.
你想一下將兩樣東西.
去到極致.

$^{281}$去到最極致.
可以想像得到.
我們這一頭去敬拜神.
那一頭就開始咒助其他人.
所以聖經告訴我們.
絕對是隨時都可以發生的.
但是亞伯告訴我們.
我的弟兄們 這是不應當的.
泉源從一個眼裡能發出甜苦兩樣的水.
我的弟兄們 無花果樹能生橄欖.
葡萄樹能結無花果.
鹹水裡也不能發出甜水來.
所以就是告訴我們.
我們必須要留意我們的泉源.
所以保羅就告訴我們.
要羅馬的信徒.
一方面要躲避這樣的人.
叫我們跌倒的人.
還有背負所學之道的那些行徑.
所以如果是小圈子這樣去講是非.
聽是非 盡量可面則面.
除非真的需要去處理.
因為在教會裡面 譬如我們作為領袖.
我們需要有時處理不同的糾紛,誤會等等.
但是這裡告訴我們.
我們真的必須要學習 要留意躲避.
而且我們真的要在善事上聰明.
也就是說 其實就是心靈裡面的那個洩頭.
或者說整個心靈的那個泉源.
我們要留意 我們要在主的恩典裡面.
去取的是生命的活水.
對人必須要有一份的寬恕,憐憫,愛等等.
但是在惡上要愚處.
所以不是蠢的意思叫做愚處.
是我們真的知道.
這些事情是會影響我和神的關係.
和人的關係的.
我刻意去學習 不敢再做了.
這樣就會變成我們一方面不是一個審判官.
我們都是在神恩典裡面.

$^{321}$蒙搖樹的神的審判的天秤.
神的義我們是虧缺的.
所以我們就能夠在神裡面去學習行善.
所以勉於行善 在善事上.
這個是一生我們在基督徒裡面都要追求的.
因為洗了禮 不再回主的學.
其實就沒有什麼機會很有系統的將神的道.
去教訓的時候.
如果你又不上網聽道 又不上網去學查經.
又不回教會的主一學.
那你靠一個星期的 只有半個小時的聖經教導.
有時候那些主題都未必講得那麼深入.
所以往往弟兄姊妹就會停滯在.
表面的那個屬靈的生命裡面.
就不能夠有認真的對自己的心靈的光照和學習.
所以我們求主讓我們能夠真的看到.
所以你想像一下 在第16章這麼強烈的對比下.
保羅就是很想讓羅曼的弟兄姊妹去學習.
因為他列舉的那26個人名全部是美好的榜樣.
全部在教會裡面有好的明星 有好的服侍.
有不同程度的參與在神國的福音工作.
所以都是不一而足的.
不同的恩賜 有些做得很多 很明顯的.
好像菲比 菲比她就是一個堅格理教會的女執事.
現在她要去到羅馬了.
就要求羅馬的弟兄姊妹要在主裡面愛心的接待.
不只是去會所吃一頓飯 而是要合乎聖徒的體統.
因為他是那個資助人 保羅是他受恩的人.
那個恩就是資助保羅傳福音的菲比.
所以在當時羅馬有這樣的資助人的制度底下.
保羅很公開的讚賞這個菲比.
可想而知弟兄姊妹我們就是要在人生裡面.
必須找到你自己在福音裡面的位置.
神怎樣在福音裡面揀選你 使用你.
你清楚自己的位分 你自己的位置.
專心的投入下去 成為你生命的事業.
就是福音的工程.
我們之前講過了一些 我們不再回頭了.
這裡特別提到安多尼古和尤尼亞.
這樣的字根裡面 都是讓我們看到.

$^{361}$中文就告訴我們 在使徒當中是有明望.
可以是兩個解釋 中文都已經浮現了.
就是說在使徒的行列之中.
對夫婦都有一定的認識.
對他們都有相當的了解.
還有都覺得他們是非常之愛主.
但是更多的解經家說 可能他們是使徒.
他們本身就是使徒.
當然了 在耶穌的年代.
十二使徒就是一個專稱 專人來的.
因為是主耶穌親自去選召的.
但是當保羅蒙召之後.
一方面他不是要替補馬提亞等等.
而是保羅是一個特別的選召.
接著可能就是這一幫的使徒之後.
使徒這一個的位份 名字就開始比較.
不是說濫了 而是那個行列一直增加.
不過就一定和十二使徒是不相同的.
因為十二使徒是跟隨著耶穌出入.
這是路加醫生告訴我們的這樣的分別.
但是使徒在保羅的年代.
或者保羅的書信裡面.
就會有一個比較一般性的統稱.
會不會就是在跟隨主的門徒當中.
會再有一些的弟兄姊妹.
他們那個蒙召.
當然絕對不是做十二使徒的身份的見證人的使徒.
但是他們在教會裡面的那個被關注.
那個服侍就好像都成為了一個使徒那樣.
所以這裡也有另外一個解釋.
就是有解經家認為.
他們本身這對夫婦就是使徒 就是這個意思.
不過保羅強調的就是.
他們是比保羅更加早在主裡面.
保羅也提到了他自己.
這是哥林多前書15章.
跟大家很快的揭一揭.
給大家有一點印象.
保羅他自己對自己那麼後的蒙召.
那種的情況 15章.

$^{401}$臨前的15章.
第一節開始.
弟兄們 我如今把先前所傳給你們的福音.
告訴你們知道.
這福音你們也領受了.
又靠著站立得住.
並且你們若不是徒然相信.
能以持守我所傳給你們的.
就必因這福音得救.
福音就是主耶穌的死而復活.
所以第三節.
我當日所領受又傳給你們的.
第一就是基督照聖經所說.
為我們的罪死了.
而且埋葬了又照聖經所說.
第三天復活了.
並且顯給基法漢.
然後顯給十二使徒漢.
是特別的一個行列.
連保羅自己也認同這十二使徒是特別的身份.
是獨一無二的了 不會再有的了.
但是也不等於保羅.
就好像成為了另類的人那樣.
就不是的.
大家當時的教會對保羅.
不會覺得他是低於十二使徒的.
但是也不會覺得他是特別的一個人那樣.
然後第六節.
後來一時顯給五百多弟兄漢.
其中一大半到如今還在.
卻也有已經垂了丁.
以後顯給雅各漢.
再顯給眾使徒漢.
留意第八節.
沒了也顯給我漢.
我如同未到產期而生的人一般.
就是吳祝月這樣來描述自己.
就是說其實本來可有可無的了.
本來是難產的了.
本來根本就不會出現在這個世界的了.

$^{441}$但是神的憐憫.
我原是使徒中最少的.
不配稱為使徒.
因為我從前迫不得神的教會.
然而我今日成了何等人.
是蒙神的恩才成的.
並且祂所賜我的恩.
不是徒然的.
我比眾使徒格外勞苦.
這原不是我.
乃是神的恩與我同在.
不愧是我是眾使徒.
我們如此存.
你們也如此信了.
這麼長的十一節的聖經.
只是想印證.
耶穌真是復活了.
真是向過很多活生生的人.
來顯現,證明這一位的耶穌.
降世成為人的耶穌.
現在在天上.
而且是透過復活,升天.
用眼都看得到的一個.
這樣的肉身升天回到天堂的這位救主.
現在他們就是為了這個真確的事件.
來作見證.
保羅在這樣描述的時候.
一路這樣數12使徒.
接著有500多個弟兄.
接著一般性的教會.
大家都見得到.
可能就是耶路撒冷的.
當時的福音還沒有傳到周圍都是.
接著就是拿國.
也就是說可能就是耶穌肉身的弟弟.
接著有其他的就是基發等等.
接著保羅最後才說他自己.
就是很謙卑地這樣來講述.
還有他給使徒們格外的勞苦.
就好像他要去補課.

$^{481}$他過去迫不及待教會.
他要做多一些的救恩,救人的靈魂.
讓更多的外邦的人.
外邦的教會可以信主.
整個腦子裡面都是猶太教的律法.
很多的律法的要求等等的.
一個這樣的律法的法利賽人.
這樣來做一個向外邦人傳道.
你可以想像到本身這個身份已經是錯配的了.
但是神的恩典可以令到保羅.
那個生命的反轉.
那個對真理的融會貫通.
將他整個生命的任督二脈都打通了.
成為游刃有餘的一個出類拔萃的基督徒.
傳道者 這個就是保羅了.
所以他說他是好像還沒有滿月出生的.
很「恩」 不夠營養這樣來長大的.
但是原來恩典可以使我們後天努力的.
恩典使我們後來一路一路成長 一路一路進步的.
所以弟兄姊妹不要小看自己年紀大了.
沒有記性了 沒什麼學歷了.
或者覺得自己什麼用都沒有了.
身體有很多毛病了 千萬不要這樣看.
使徒保羅 他是有很多學問.
但是如果不是經過聖靈的手術工作.
他的學問是害人的.
是會使很多人不僅不認識基督.
反而是更加懼怕這一位的神.
現在他將自己放在神的救恩的那個恩典裡面.
說神可以用他過去的經驗,他的學歷.
來轉化 讓外邦人清清楚楚的明白.
救恩是屬於萬民的.
這樣就可以想象到每一個人.
必須要將自己放在神的手中.
不要有一種妄自菲薄.
過去原來都可以被神用的一個方法.
所以我們如果的原生家庭.
或者我們自己以往的成長背景.
這些一定是有它的限制的.
但是在神裡面就沒有限制了.

$^{521}$我們就要讓我們自己的背景.
都來去交給神.
所以這段經文都很寶貴的.
所以提到使徒.
保羅就這樣來描述他自己.
所以有可能去到保羅的年間.
這個使徒的這個稱呼或者這個身份.
可以再簡單一點也說不定的.
所以這個尤尼亞希利文的原文是陰性.
所以似乎是女性的使徒.
很特別 所以這些字根等等.
都是解經家他們去考究的.
但是保羅說是他的親屬 一同坐牢.
而且在使徒當中有名望.
教保羅更加早在基督裡面.
保羅對他們尊敬有加.
我們都要學習.
不過我們這一輩就因為不是後現代的.
千禧年或者現在的Z世代.
當然沒有貶義的.
不是說年輕就一定是對人沒有禮貌.
或者那個表達上.
我想我們那一輩一定是會對長輩,老師.
我們會有基本的禮貌,修養等等.
我想大部份香港人.
我們這一輩60後,50後.
都會有這樣的尊師重道.
因為都挺重的.
但是原來這些的美德.
你看一下保羅這樣的自律行間的描述.
你就可以想像到保羅對這些人.
所以我們不知道他是不是使徒.
但是對這些在教會裡面的.
服侍的安多尼古尤尼亞.
似乎這對夫婦是很敬重.
很欣賞他們在主裡面比保羅更早.
保羅是很後期信主.
一信就一個突變.
起初這個少年人去抓基督徒.
看到司提反殉道在偷笑.

$^{561}$看著那些劊子手的衣服.
就是說已經告訴我們.
在《四道行主》第七章.
路加醫生就已經將保羅他的成長背景.
很快地用一兩節聖經給我們知道.
就是說這個人一定是凶神惡煞.
一定是滿肚饑鱗.
就是猶太教的饑鱗.
滿有希臘的學識.
對很多希臘的學問是很順暢,很熟悉.
但是那個人的心地舊約的背景.
來搞到很混亂.
他還以為自己是很事奉神的.
在當時更多的猶太教裡面.
都是會誤解了耶穌和福音.
當時都沒有所謂的有系統的福音的講解.
福音書都沒有.
但是神就透過一個這麼特別的恩典.
去拯救保羅.
使得他真的成為了一個.
好像是當時神大用他的一個器皿.
所以原來每一個人 特別年輕人.
如果有些特別的經歷能夠信主.
那個生命的蛻變是難以估計的那種力量.
但是如果沒有被主耶穌的愛征服.
就可以是一個很大的破壞者.
破壞人的童心 破壞教會的合一.
破壞很多東西 所以我們那個腦袋.
真的要不斷更新.
我們的神學 不是神學需要不斷更新.
是我們自己以為的神學要更新.
我其實很慶幸.
因為剛剛星期三在洪水橋準備祈禱會.
好像星期三我在大法那裡吃飯.
在大法看到一對弟兄姊妹.
認識了幾十年的弟兄姊妹.
他們上的是我舞堂的那個教會.
現在在洪水橋有另外一間叫天守圍福音堂.
就是我的成長背景的傳道人的教會.
所以其實相當熟的.

$^{601}$但是那個背景的情況就是非常.
因為基要派非常嚴謹地走聖經.
不知不覺對其他的弟兄姊妹.
基本上就是很抽離 沒有大家的來往等等.
有時候就是跳不出我們自己已經有一套的神學觀.
去框住了 所以保羅現在能夠在教會裡面.
他現在是以神的使徒的身份.
去對這些弟兄姊妹的問安.
這位使徒這麼多經歷 這麼多的服侍.
這麼被神大用 但是對這些長輩.
其實可能很普通的人而已.
他能夠很虛心地多謝這些人.
這對夫婦在使徒當中是有名望的.
他還教了保羅在基督裡面重要祖.
這樣的描述就可以想象到.
相信我想我們在我們自己成長的環境裡面.
我們一點都不難對我們的老師.
對傳道者那個從心裡面的尊敬.
我想我們一點都不難的.
如果難的話 那就可能很後現代的了.
或者你真的很自我中心的了.
因為如果你看到服侍主的人.
你都不能夠欣賞和尊敬.
其實就相當大的問題.
因為人家這樣擺上.
你連這樣擺上都沒有.
你都這樣對人家有一種藐視.
那你想一下就差很遠.
那你可以想象到我們要學習保羅.
保羅貴為一個這麼大的使徒.
你想一下他寫了16章的經文.
多少的真理在羅馬書裡面有系統的講述.
但是當講到一些人的時候.
很普通的一對夫婦都來對人家的讚賞,尊敬.
這些就是畢生我們的學習.
都應該就是要這樣.
又向我在主裡面所親愛的暗白尼文安.
這個親密的稱呼在羅馬書16章.
只是用在幾個 以巴盧是大古 這些是奴隸.
暗白尼是拉丁文.

$^{641}$羅馬帝國常見於奴隸或者是被釋放的奴隸.
在主裡面蒙恩愛的弟兄.
在基督裡面不再分自主為奴或男或女.
很強的保羅的教導.
在救恩裡面每個人在主耶穌裡面都是神的兒女.
當然在服事上有不同的崗位.
彼得明顯就是耶路撒冷的使徒.
保羅是愛邦人的使徒.
大家的範疇不同.
有些人比較明顯的恩賜等等.
但是保羅在那個生命裡面.
大家同樣蒙神的恩典的.
這個生命的改變上面是一樣的.
所以菲比是呂執事.
阿彌陀佛家裡面的人.
這個可能是希臘家族有關的.
暗白尼就代表著奴隸的階層.
保羅都這樣來說.
這個暗白尼向我在主裡面所親愛的暗白尼文安.
對奴隸都有這樣的尊重.
可能已經成為了自由人了.
因為他是奴隸.
但是如果他的主人是一個很大的家族.
那個奴隸都有相當的自由度和權力的.
但是這個是羅馬帝國當時的治理底下.
為了維持那個國家.
所以他很多的人差不多七成就是奴隸.
所以不是我們想像當中.
每天晚上要住在柴房.
沒得吃 不准睡.
當然要工作.
但是那個身份不是我們今天所看的.
這個就讓我們看到保羅在治理行間的文安.
都對一個奴隸有一個這樣的讚賞.
在主裡面我所親愛的暗白尼.
這裡說又向我們在基督裡的童宮以巴盧.
和我所親愛的士大古文安.
又向在基督裡面經過考驗的亞比利安.
三國的人.
你可以想像到保羅不斷地去描述這些女性.

$^{681}$我們在基督裡面的童宮.
伯吉拉夫婦 菲比 以巴盧.
這樣的稱呼.
這裡告訴我們 這些就是多元背景之下的縮影.
士大古 亞比利 有不同的階層.
很正常 在羅馬的教會.
所以我們感恩 在錦繡堂.
由我們信主一直在洗禮.
我們會有同非裔 印裔.
遲些可能我們泰裔也有些人.
還有本地的 這樣一起洗禮.
就已經見證著我們是多元的教會.
我們都是能夠彼此相交.
彼此相愛 是一件非常美麗的事.
我們在聯合崇拜 又會有菲律賓的姐妹.
印尼的姐妹 來到無論獻唱.
或者用他們的舞蹈來呈獻美給神.
所以保羅現在在這樣的民安裡面.
雖然這些名字已經過了二千年.
對我們來說毫無意義.
但是當我們明白了這些的字根.
和當時的文化背景.
解經家如果他們沒有理解錯.
就已經告訴我們知道.
保羅是不是刻意的將這些不同行列的人.
擺在一起 襯托出一幅美麗的圖畫.
就是在羅馬教會裡面已經有不同的階層.
一起在文化的分歧裡面去融合.
所以我們都是不斷的彼此提醒.
你1至6或者1至5 現在上午班.
1至5可能你周圍穿梭去開會 國際的會議.
因為現在飛機的方便.
可能你星期一飛到美國.
接著就在那裡開會.
可能星期三你已經去了歐洲.
接著再回來 可能我們是這樣的人.
也都更多的就是我們每天早上就要逼968.
每天很辛苦的去到中環的飯盒.
接著我們很繁忙的工作.
甚至有些是基層的工作.

$^{721}$每天做一些很沉悶的工作.
甚至我們是打掃一些公廁的.
如果都是弟兄姊妹 我們要學習 我們要記得.
我們來到星期天 一般性的說.
因為大家有不同的 有些輪班有些什麼.
按一般性 大家去到星期天.
都可以將他的工作的服裝脫下.
星期大家 當如我們大家都穿教會的T恤.
大家都穿那個T恤.
那份外的那個T恤就來的緊要了.
就是在基督裡面 我們是合一.
在基督裡面 我們在神的眼中.
我們都是神的兒女.
A是跨國企業的CEO.
或者他經常要去到世界各地.
B天天就是這樣活天的辦公室.
早上7點就要出門口.
968 9點就要去到上環或者中環.
甚至灣仔這樣才算.
黃昏6點又要坐車回來.
天天如是 又辛苦又悶 有時候要加班.
但是C原來是很基層的.
原來在我們附近的麥當勞是做收拾東西,擦桌子的.
去到星期天來到教會.
大家換了教會的那個T恤.
大家是神的兒女.
這個是不容易去學的.
是不容易放下自己的.
高的如果學得不好會小看.
低的就會有自卑感.
其實只不過是那個工作和工作的收入待遇.
那個背景,那個情況帶來這些分歧而已.
這些分歧本身是沒有需要去打破的.
但是那個生命就是要打破了.
就是我們來到神的家裡面.
為什麼我們真的要彼此在主裡面來相交相愛呢?.
大家要很坦誠的分享.
大家要衷心的看到對方是神的兒女.
是神救恩的一個產物.
所以你看保羅這樣來對這些人的名字.

$^{761}$好像只是講一句而已 只是提一提他們的名字.
但是都已經在那個日語言表裡面.
保羅對這些人 他們的階層.
他不會因為他的階層特別要哀愁奉承.
或者對一些基層的,沒有學問的,沒有社會地位的.
沒有教會勢力的 就有很多的看不起.
不會的 如果有的話.
根本就沒有資格做教會的領袖,做使徒.
因為救恩就是帶給我們在基督裡面.
是要合而為一.
在保羅後期的書信是大量這些經文的.
所以這些是非常嚴肅的真理來的.
如果我們一生不去學習.
我們這件事是將來神一定審判的.
我們可以不同大家看事物的角度.
大家來自不同的背景,思考的模式都不同的.
但是一定要看對方是神的兒女.
除非他是異端.
就算他是異端你都要帶他回到真理.
就不要當人家是仇人.
比如看這兩段 已忽所書的第二章.
第二章第十一節.
已忽所書的第二章第十一節.
所以你們應當紀念你們從前.
按照肉身來講是外邦人.
就是已忽所教會裡面他們拜的那個大安娜.
那個生育的神,有雨水的神那個大安娜.
在已忽所那個地方還有一個遺址.
有一條柱就是那個大安娜的廟.
那些大的石是供奉的.
你可以想像當時已忽所是一個相當拜偶像鼎盛的地方.
但是很奇妙在已忽所有一個很大的廣場.
保羅就在那裡傳道.
接著就說是稱為沒有受過國禮.
這個名字源自那些貧人守在玉山上.
稱為受過國禮之人.
那就是猶太人所起的.
就是說你們未受過國禮.
就是說你們是非我們猶太人的族類.
那就是說你們不是亞伯拉罕的子孫.

$^{801}$那你們就跟我們的神是沒有份沒有關.
所以第十二節.
那時你們與基督無關在以色列國民以外.
在所應許的諸約上.
就是亞伯拉罕的約,大衛的約,所羅門的約.
撒母耳那個應許的約等等.
這個就是撒母耳上的應該是十八章.
我等一下找一下 我忘記了那裡.
在所應許的諸約上是局外人.
並且活在世上是沒有子望,沒有神.
就是你們是外邦人的時候.
就是未在救恩裡面的時候.
你們的子望只是今生今世.
你們有很多的學問.
有當時的亞里士多德,柏拉圖.
很多這些哲學家的 蘇格拉底.
所以保羅都認識很多的.
但是原來在猶太人的眼中.
外邦人是與神,與造物主無分無關的.
但是留意十三節 你們從前遠離神的人.
如今卻在基督耶穌裡靠著他的血.
已經得親近了.
不是因為猶太人現在豁達.
現在大開門戶讓外邦人可以信.
不關猶太人的事 是基督耶穌的救恩.
所以猶太人自己都相信.
是基督耶穌來到這個世界.
因為耶穌基督的緣故.
是靠著他的十價的流血,他的犧牲.
就現在本來遠離 現在是得親近.
就是說原來我們在真理,在主耶穌裡面.
我們是可以隨時去到神那裡.
去到永恆的主那裡.
我們本來是遠離的 現在是得到親近.
第十四節 為什麼呢?.
因他使我們和睦 或者說他本身就是和睦.
將兩下 這裡是上下文很清楚的.
就是猶太人和外邦人將兩下合而為一.
拆毀了中間隔斷的牆.
而且以自己的身體廢掉冤仇.

$^{841}$什麼是隔斷的牆呢?什麼是那些冤仇呢?.
就很清楚了 就是原來十價的律法.
就是那記載律法上的規條.
說要將兩下 兩下就是這兩班的人.
這兩類的人 藉著自己就是耶穌.
造成一個失人 如此就成就了這個救恩.
帶來的人類的和睦.
記載十字架上滅了冤仇.
便藉著十字架使兩下歸為一體.
與神和好了 就是這個和平的福音.
所以 希望神都幫助我們在地上 在教會.
我們成為一個和睦之子.
我們成為一個能夠促進那個合一的人.
能夠讓我們的笑容 我們內心的親切.
我們的心底裡面的那個渴慕.
帶來的就是神家的合一.
在愛裡面同歸於一.
這個其實是很好的一件事.
我們如果整個的心有一個同歸於一的心.
我們就充滿了一份平安 盼望 笑容.
就是這樣可以讓人得到這個和平.
所以 如果是使人為敵的人 就有禍了.
所以你可以看到那個登山補份就是這樣.
散播謠言的人就嚴重了.
因為帶來的就是分裂 帶來的就是猜疑.
帶來的就是分化 帶來的就是那個仇恨.
所以我們記得千萬不要這樣.
因為神必定審判的 這些罪.
阿彌陀佛一家的民安.
就是說很清楚了 阿彌陀佛已經是死了.
他的家人繼承了這個信仰.
就是已經是安息在主裡面的一位弟兄了.
所以保羅又再一次說.
向我的親戚希羅天民安.
向拿其訴家在主裡面的人民安.
是保羅的親戚和拿其訴家裡.
似乎解經家告訴我們是尼祿王身邊的.
一個重要的秘書.
所以當時的王族都有信主的.
那你就想像得到 福音可以令到那些奴隸信主.

$^{881}$在王族裡面都有些是信主的.
譬如在《士道行傳》的17章都有提到.
有些尊貴的希臘的婦女信主的.
17章後面.
這裡有幾段 譬如可以說17章.
看後面那裡.
32節這一段.
眾人聽見從死裡面復活.
就有譏笑他 又有人說.
我們不再聽你講這個了.
於是保羅從他們當中出去.
但有幾個人貼近他 信了主.
其中有阿略巴古的官刁尼修.
阿略巴古就是那個地方.
另一個婦人名叫大瑪麗.
還有別人一同順從.
這個就是在當時雅典的地方.
一個相當繁盛的地方.
阿略巴古一個聚集.
來討論學問哲學的地方.
在巴特勞神殿下面的地方.
巴特勞神殿就是現在在希臘的市中心.
很容易走上去的.
相當宏偉壯觀的一個.
到今天 又經過地震.
還有一個這麼大的規模的巴特勞神殿.
是非常壯觀的.
如果只是去希臘 一定要去這個巴特勞神殿.
很壯觀.
你可以想像得到那個規模有多大.
我們今天去都看到是非常壯觀.
何況當時在保羅的年代.
是敬奉這個巴特勞的天神.
你就可以想像得到.
希臘人他們是真的非常迷信.
但是又是非常充滿哲學的哲理等等.
但是保羅對這個拿奇數.
都有一個很大的稱讚.
就是這樣了 接著有些女性.
傑出的女性.

$^{921}$土菲娜氏,土庫薩氏,比瑟氏.
我們能夠知道的就是這麼多了.
只有一個名字.
保羅就是這樣來對他們的稱讚.
我們完這一段之後.
十六章我們可以當是看了一點.
最後的那個問安 第21節.
與我同宮的提摩泰.
和我的親屬路求耶孫.
索西巴德問你們安.
當時就是在哥林多寫這封羅馬書.
還有人幫保羅代筆寫信的德彪.
在主裡面問他們安.
接待我 又接待全教會的該由.
問你們安.
城內哥林多城的管銀庫的爾拉都.
和兄弟闊土.
都是有相當的地方官,政府官員信主的.
問你們安.
所以基督教因為真的講求真理.
還有真是一個傳播和平福音.
所以真是有不同的階層的人信主.
保羅的年間羅馬是一個相當大的地方.
哥林多已經是相當富庶的地方.
所以有這些有學問,有知識.
有一定權力的官員信主.
一點都不奇怪.
所以某程度我們就不要對立.
我們也要為一些高官信主來禱告.
因為他們有時候在那個位置.
有一些相當難的決定.
或者有相當的掣肘.
但是他們如果能夠有智慧.
去在他們的崗位上面行使那個權力.
其實帶來的也是一個相當大的祝福.
制定的政策的背後.
如果是從基督教的那個精髓,精神去設想.
又會帶來祝福會很多人.
所以千萬不要對權力一概否定.
當然用野心去追求權力是聖經所禁止.

$^{961}$但是權力本身是很忠誠的.
我們坐得高,能夠決策得好.
有足夠的社會資源.
其實是可以帶來福音更好的那個祝福的.
比如我不知道譚志源弟兄.
他現在是在馬會裡面的福利事業的總經理.
相當的行使的無論是錢.
或者他每年制定的不同的計劃.
如果他有基督的恩典,有那個真理.
有福音在他內心.
他帶來的很多事情.
就未必這麼快就說進馬會,馬會做.
未必需要這樣的 他是相當好的弟兄.
他是原浸成長的 鄉郊元朗的兒子.
讀的書,考AO等等.
所以有些原浸現在的理事,執事.
都和他一起玩泥沙大.
我們對有權力的人不要這麼快否定.
反而我們要為這樣的弟兄姊妹禱告.
當然我們如果是用野心去追求權力.
聖經說我們就不好 所以我們求神幫助.
這裡說城內的管銀庫的以拉都.
和兄弟闊土都問你們安.
所以保羅其實他沒有抗拒財富.
沒有抗拒那些權力的.
不過有一段經文大家都可能聽過.
提摩太前書的六章十節.
貪財是萬惡之根.
有人貪戀錢財就離氣爭道.
而且用很多仇富來將自己刺透.
其實是貪財 財富本身是很中性的.
所以保羅就一點都不是偏激的.
他是很正面去看事物的.
所以原來我們都看到這些經文.
是當時的社會是有這些官員的.
可能包括我們今晚要為現在信主.
表示自己是信主的應該是康樂.
不是 那個是什麼?我突然間忘記名字了.
經常都見報的體育局局長.
另外那個局長也是信主的.

$^{1001}$所以我們都要為不同的弟兄姊妹在公營的事業.
或者在職場上坐在高位要為他們祈禱的.
第25節 唯有神能夠照我所傳的福音.
和所說的耶穌基督.
並照永古隱藏不言的奧秘.
堅固你們的心 用這樣的祝福.
不是輕輕的 神祝福你.
是神要保羅所說的.
保羅告訴別人的萬民需要福音.
因為萬民都在神的咒詛底下.
因為虧欠神的義.
現在因為基督耶穌.
神的義就加給有罪的罪人.
就可以成為恩信稱義了.
得到的是神的義.
所以保羅這樣來祝願他們 堅固他們.
保羅所說的福音就是耶穌基督.
照著永古隱藏不言的奧秘.
本來是沒有人知道的.
用神祂自己的榮耀智慧隱藏的.
但是現在保羅的年間.
神的兒子出現之後.
這個奧秘就揭開了.
所以我們今天得到的心靈的堅固.
就是我們認識了神這個奧秘.
就是福音 就是耶穌基督.
所以耶穌基督就能夠堅固我們的心靈了.
接著告訴我們這個奧秘如今顯明出來.
而且按照永生神的命.
直眾先知的書 指示萬國的民.
使他們信服真道.
所以你看到就是最後的結論.
都是萬民需要聽到福音.
因為萬民都是在罪惡裡面.
所以保羅就說按照永生神的命令.
這個奧秘現在已經很清楚的啟示出來了.
明白福音的人就還沒有得到這個奧秘.
所以我們就要將這個奧秘有系統的講明白.
讓還沒有聽到福音的人聽到.
所以是信服真道 萬國的民信服真道.

$^{1041}$願榮耀因耶穌基督 歸於獨一全智的神.
榮耀歸給這位榮耀的神.
這裡特別講到獨一全智.
他的榮耀就是.
因為保羅的羅馬書都是整個的論述.
其中如果再精簡就是這個福音.
所以他發明 他成就.
他現在叫我們去傳揚.
這個是獨一全智的神.
要得到一切的榮耀.
所以我們傳福音是將人帶進神的榮耀裡面.
我們去傳福音 保羅第十章羅馬書.
報喜信傳福音的人.
他們的腳中何等的佳美.
我們不是阿Q精神.
而是我們真的帶著寶貴的福音來跟別人講.
我們是帶著佳美的腳中.
全的人本身已經得到那個福氣.
所以最後還有25分鐘.
跟大家很簡單的講一講.
整個羅馬書最重要的寫作目的.
很有系統的來展述福音的真理.
要揭的經文 因為大家可能已經有點遺忘了.
盼望大家有時間就重複重複的.
從第1到第16章 不厭其煩的自己讀通它 讀明白它.
一章的11節 羅馬書.
因為我切切想見你們.
要把一些屬靈的恩賜分給你們.
使你們可以堅固.
這樣我在你們中間 因你與我彼此的信心.
可以得到安慰.
弟兄們 我不是願意你們不知道.
我屢次定義去到你們那裡.
就是要在你們中間得些果子.
如同在其餘的愛邦人中一樣.
只是到如今仍有阻隔.
這個就是那個很清楚的寫作的動機.
就是要讓他們都能夠分享.
神在保羅生命裡面的那個恩賜.
那個恩典 這樣就一定是講述.

$^{1081}$關乎到一個整傳的福音的信息.
15章我們之前有提到的23,24節.
23節開始 但如今在這裡再沒有可傳的地方.
而且這好幾年 我切心想望到西班牙去.
可以到你們那裡 盼望從你們那裡經過.
得見你們 先與你們彼此的交往.
心裡面稍微滿足 然後蒙你們送幸.
這個就是很清楚的.
他很想有系統地將這個福音.
來去展述這個真理.
這個就是保羅告訴我們 他要來講述的.
很重要就是促進猶太人與外邦人的信徒的合一.
最重要的就是應該是10章這段經文.
第10章 12節.
猶太人與希臘人都沒有分別.
因為眾人同有一位主 他也厚待一切求告他的人.
因為凡求告主名就必得救.
是在講猶太人,外邦人.
現在必須要在主的救恩裡面.
來去認定這位真神已經賜下唯一的福音.
就是因順稱義的耶穌的犧牲,十價的愛.
所以我們就是要將這個合一去傳遞.
人類只是分了屬於猶太教,屬於律法.
和沒有律法兩者 都是活在一個虛無裡面.
在神的救恩以外.
現在我們要帶人去求告主.
就是一個很大的生命改變.
我之前都用文字跟大家分享過.
那位王永康先生真是很奇妙的.
他居然主動透過車房的弟兄找我.
一頓茶跟他一起去求告.
叫他自己晚上去祈禱.
其實輾轉都是大約四個星期.
他後來又回過一兩次北宣.
接著自己差不多很痛苦要離樓的時候.
叫太太去北宣找牧者在他床邊講福音,洗禮.
安息主懷 我去到靈堂.
其實不認識那些喪家 後來自我介紹.
他太太突然間想起就是曾牧師.
就跟我進去後堂 還有他兩個子女.

$^{1121}$就講述爸爸,丈夫自從那晚之後.
晚晚祈禱 還有開始回一兩次北宣的聚會.
接著真是很清楚 堅持叫家人叫北宣的牧者來.
因為太太不認識我 又沒有我的電話.
所以找不到我 所以北宣的牧者又去跟他洗禮.
真是很感恩 求告主明就必得救.
神給他很多的恩典.
其實他很年輕的時候 只顧著做生意.
看那些照片播放出來在靈堂.
他應該是一些馬主.
原來在身邊有很多做生意的朋友都勸他信主.
至少我車房的弟兄.
因為他就是幫這個車房的弟兄去修車.
所以很多年的朋友 但是一直都不肯信.
所以那次他自己突然間說.
記起我 要找我喝茶 要為他祈禱.
就是一個這樣的很小的動作.
就救了他自己的生命 所以都祈禱.
希望他的家人都能夠信主.
所以這裡告訴我們 猶太人,外邦人.
未在救恩裡面 全部都需要求告主明才得救.
這裡已經是所有的了 無論你是貴族.
無論你是平民,是奴隸,無論你是猶太的拉比.
或者你是外邦的,拜賈神的,祭司等等.
所有的人都需要.
萬民都需要求告主明就必得救了.
第三段我們可以來到講.
可以表達的就是這個 剛才我們讀了 不再讀了.
就是他要去羅馬宣教 接著再路經.
接著就去到西班牙 梵蒂岡大家看到.
在羅馬的境內 天主教.
希望裡面的人都能夠不要被一些教廷的教義.
死守一些表面的做法.
而忘記了心靈的那個割禮,歸順主.
第四段都重要.
他是要糾正錯誤的理解律法.
所以在《自律行間》他都有這樣的一種教導.
三章第二十節都重要了.
所以凡有血氣.
沒有一個人因為行律法能夠在神的面前稱義.

$^{1161}$對猶太人來說是一個相當大的當頭棒喝.
他自己就是一個律法的專家.
沒有一個人一輩子守十戒就可以稱義.
守足十戒都不可以稱義 也都不可能守足.
因為我們的心靈都是敗壞.
因為律法那個真正的意思.
只不過是叫你知道什麼是罪.
定義那個罪給你知道.
但是你自己不可以去解決自己的罪念,罪惡等等.
所以凡有血氣.
沒有一個因行律法能夠在神的面前稱義.
是相當清楚,斬釘截鐵的.
譬如第六章.
這些經文就浮在整個大的論述裡面.
第六章這樣怎麼說呢?.
我們是不是可以刑在罪中 叫恩典顯多芒.
這些就是信了主一個錯誤扭曲了福音.
白白恩典的定義.
是不是可以隨隨便便這樣做基督徒呢?.
我們可以肆無忌憚這樣繼續.
未信耶穌之前肆無忌憚是與生俱來的天性.
但是信了主有了聖靈之後.
我們就不能夠也都不會再肆無忌憚了.
這裡說 斷乎不可 我們在罪上死了的人.
這個就是信了主之後的生命改變.
就是說我們已經有一個生命.
所以用的說法就是在罪上死了.
豈可刑在罪中活著呢?.
豈不知我們者受洗歸入基督耶穌的人.
是受洗或者受斬 歸入他的死.
所以保羅不僅要糾正 守律法也是不行的.
不是律法本身有問題 是人守不了律法.
但是相對於恩典 我們又不可以因為信了主.
就隨隨便便這樣 我們就不用守律法了.
我們可以肆無忌憚這樣犯罪了 也是不行的.
所以很清楚的將這四個重要的目的.
在《羅馬書》裡面來交代.
就是說四個很重要的核心的中心思想.
當然有整個的《羅馬書》充滿了那個神學.
差不多是它的原創.

$^{1201}$保羅很清楚這個因信而得的義.
一個justification 就是被宣判無罪的意思.
透過信心被宣判無罪.
所以不是靠恩典的人.
我們就不可能可以靠自己來稱義.
所以是靠恩典 by faith 透過信.
所以恩典本身不是一個行為.
也是一個信心的領受的禮物.
這個就是那個神學的重點了.
神的義在第一到第四章就相當重的講神的義.
這個就不再詳細的複述了.
譬如一章十七節就很清楚的講了.
我也是稍微讀一讀 一章十七節.
因為神的義正在這福音上顯明出來.
這義是本於信 已至於信.
如經上所記 義人必因信得心.
清清楚楚的告訴我們.
福音就是神的義 神自己本質的義.
你可以說都是救恩帶來的那個義.
我們本來是不義的.
我們本來就是在神眼中無義可言的人.
現在得到的是神的義 而且只是需要憑信心.
接上這一位生命的主宰.
我們就得到這個神的義了.
賜給一切相信的人.
所以這裡就是普世的人都需要福音.
都沒有分別 規缺了神的榮耀.
普世的人的罪 那個需要就是救恩.
這個就是羅馬的鬥獸場.
一個當時對基督徒一個極盡侮辱的迫害的地方.
現在那個原址還在.
但是救恩繼續傳遞.
尼祿王都過去了 很多那些暴君都過去了.
但是救恩仍然在 這個這樣的原址就印證著.
裡面充滿了宣教士,基督徒的血.
但是成為了福音的種子,萌芽的地方.
所以我們求神幫助我們能夠看到.
讓更多人能夠看到.
其實一生人如果沒有主耶穌.
其實是規缺神的榮耀.

$^{1241}$一生人都不能夠活在神的那個恩典和愛.
和那個Blessing的底下.
無論有些人很成功,很順利,很有錢.
生活如意 但是在神的啟示,保羅的教導.
告訴我們 規缺了最重要的生而為人.
那個最大的虧欠就是神的榮耀.
所以有時候我們看那些可憐 就是基層可憐.
有缺憾可憐 一生要為自己的身體掙扎.
求存的可憐 很多的缺憾可憐.
家庭的悲哀可憐.
一出生就可能是無父無母的孤兒可憐.
所以我們人類都分了很多不同的等級的可憐.
但是我們很少人會感受到.
規缺神榮耀的人是最大的可憐.
如果他在有生之年有很多機會去信主.
而拒絕,而硬心 而以為是根本不稀罕.
這樣的人就真真正正在永恆裡面的可憐.
因為真是離開物質的世界.
是真是有永恆 是真是有靈界的存在.
所以這裡就是告訴我們.
我們無論去到什麼樣的人生階段.
我們都必須要熱誠地參與救恩的工作.
不同的範疇去參與救恩.
出錢,出力,祈禱等等.
這個都是我們需要去做的.
遮住了 在基督裡面的新生命.
聖靈的 先拿走這個.
對不起 回頭.
這裡也是一個很大的神學課題.
基督裡面的新生命.
用了相當多的經文.
第六章到第八章都是不斷描述.
罪仍然未完全的滅絕.
但恩典其妙的地方就是.
罪 我們這個人沒有改變.
但是新的力量就是聖靈進入我們內心.
可以慢慢讓我們成聖.
慢慢讓我們越來越像神的兒女.
這個就是那個生命力.
在基督裡面的新身份.

$^{1281}$和聖靈的生命在我們裡面.
所以六章的十八節.
不再作罪的奴僕 作義的奴僕.
就是說一生就是追求義.
神的義加了給我們.
但是我們都要在生命裡面追求神的義.
聖靈就進入我們內心.
見證我們就是神的兒女.
十一章 這裡大家都可以回頭看一看.
這個關乎到簡選 關乎到神的主權.
十一章.
我且說 神棄絕了他的百姓們 斷乎沒有.
因為我都是以色列人.
亞伯拉罕的後裔 屬別雅文茲派.
神並沒有棄絕他如先所知道的百姓.
你們豈不曉得經上論道以利亞.
是怎麼說的呢?.
他在神面前怎樣控告以色列人呢?.
主啊 他們殺了你的先知 拆了你的祭壇.
這是聶王紀上的十九章.
神的回答是什麼呢?.
我為自己留下七千人.
是未曾向巴黎屈膝.
所以第五節 如今也是這樣.
照著簡選的恩典 還有所留的餘數.
既是出於恩典 就不是在乎行為.
不然呢 恩典就不是恩典.
這個上下文是在說猶太人的.
保羅很明顯是在說他自己都是蒙恩典.
神的簡選 預定他信主的.
所以他這麼多的敵黨神.
都沒有令到他在救恩的以外好像永遠被棄絕.
所以保羅就從這個概念.
他知道神真的有他的主權.
有他自己的簡選.
不過照著簡選的恩典.
所以這個我們都一起去學習.
應該這麼說的.
沒有人知道誰被神簡選的.
因為唯有神才知道.

$^{1321}$但是當我們向他傳福音.
他生命有改變.
這樣的人其實就是冥冥中神簡選了.
似乎好像到底是雞先還是蛋先呢?.
因還是果呢?.
我們不需要想得這麼複雜的.
因為沒有人會知道自己是被神簡選.
更加不會有人在那裡去嬉皮笑臉.
或者在那裡去謾罵神.
是啊 因為我不信耶穌 因為祂沒有簡我.
那我就不信了.
很多人會這樣說的.
太輕忽 以為永生是可有可無的東西.
但是如果我們明白救恩這麼重要和寶貴.
能夠信到主 決志改變生命.
這個其實就是真是神的簡選.
所以我們人都不知道到底是因還是果.
但是無論如何.
我們從我們的角度 地上的人.
基督徒我們只知道信主 求告主.
他就可以有永生.
所以我們要帶領萬民求告神.
得到神的救恩.
至於簡選那些.
就不是我們有智慧的去分辨.
所以我們就千萬不要扮演神.
對了 該死了 你沒有被神簡選了.
所以我不向你傳福音了 那就不行了.
所以這裡告訴我們 信仰的群體.
這裡說的是倫理的相交 相愛.
真的信仰就是應該有信徒的生活.
這些都不侵然了 因為都已經說過很多次了.
所以最重要的核心.
全人類都需要救恩.
神藉著耶穌基督的十價復活.
因信稱罪人為義.
並藉聖靈使人成聖.
最終萬民要歸榮 要吸神.
很簡單的一句 做了《羅馬書》的總結.
都感恩 多謝各位陪我42堂.

$^{1361}$都是自己去研讀.
將自己讀完的跟大家分享.
所以盼望大家都對這一本書.
首先不要陌生.
不要只是滿足於背一兩節金句.
而是有系統地 很快的逐章去研讀.
大概你能夠有個脈絡.
然後別人問你隨便可以揭曉的.
這個是那個基本要求.
這個是很大的工程來的 因為16章.
讀完都差不多三四個小時 一次讀完.
所以你想一下還要放很多的時間.
放很多的心機去研讀.
希望大家有一種拋磚引玉.
我們下一卷就會看《出埃及記》40章.
其實分段相當簡單 比《羅馬書》沒那麼複雜.
舊約不是我們想象當中那麼難讀的.
當然很多裡面的東西.
我們可能到時用圖像化.
也都相當容易理解的.
我們一起有個簡單的感恩禱告.
多謝主容許我們來到去.
很膚淺的將《羅馬書》去研讀.
其實在歷史歷代 《羅馬書》的注釋書.
在圖書館 教會的圖書館.
或者是神學院的圖書館.
是最多人寫的.
連加爾文都寫過《羅馬書》的注釋.
已經是500多年前的牧者.
歷史歷代很多的研究.
到今天都繼續還有不同的牧者.
神學院的教授去研讀《羅馬書》.
這本書差不多在基督教的圈子.
讓我們很有系統地對這個一信清二的福音.
來認識 天父啊 感動我們.
一起向萬民傳福音.
用盡我們人生的智慧,財富,資源.
來去投身在福音的行列.
讓更多人得到主的福音.
求你賜給我們不同的弟兄姊妹.

$^{1401}$領受不同的群體.
有些人可能是向SEN.
有些人可能是向一些貧窮的人.
失學的 或者是一些可能孤兒等等.
求主讓我們領受我們的異象.
以致我們大家在不同的領域為主去傳揚福音.
因為人生真是很苦短的.
幫助我們不要浪費在一些只是吃喝玩樂.
反而就是能夠有衣有食.
我們就當知足這樣來服侍.
\newpage



\section{出埃及記}
\label{sec:GROoST1kr4I}
\textbf{2026-05-15 周五晚查經班《出埃及記研讀 進入應許之地(一)》 講員:曾裔貴牧師}
\newline
\newline
連結: \href{https://youtube.com/watch?v=GROoST1kr4I}{\texttt{https://youtube.com/watch?v=GROoST1kr4I}} ~~~~ 語音日期: 2026-05-18
\newline
\newline
\hyperref[sec:WmefrJ6J_Lg]{\small{< < < PREV SERMON < < <}}
~
\hyperref[sec:index_chronic]{\small{[返順時目]}}
~
\hyperref[sec:index_scriptual]{\small{[返順卷目]}}
~
\hyperref[sec:kZF3v6CUrpM]{\small{> > > NEXT SERMON > > >}}
\newline
\newline
$^{1}$好了 各位姐妹晚安.
能夠和大家開一本新的書卷.
我們一起禱告.
親愛的天父.
晚上我們能夠將一個新的時段.
盼望帶來一個新的學習,新的發現.
特別《春艾讀記》裡面有一些資料.
是相當沉悶的.
求主讓我們能夠從一個宏觀去看這本書.
因為它是佈局在整個摩西五經.
一個新的時代.
願神幫助我們.
也讓我們透過看《春艾讀記》.
不單純是字面裡面的上文下理我們要明白.
但是在整個五經的佈局裡面.
都讓我們能夠掌握,認識,求主帶領.
也讓我們可以有更多弟兄姊妹.
都能夠透過這本書得到那個鼓勵,幫助.
我們仰望神.
我們這樣禱告奉耶穌基督的名.
阿們.
也感恩能夠和大家有機會去看不同的書卷.
選取一本相對比較簡單的經文,經卷.
需要用的時間去準備又不算很多.
因為它的分段都相當清晰.
主要有三個大段落.
不過我們看羅馬書.
因為是一些書信的格式.
是一些神學的宣告.
和我們在舊約一些的敘事文體.
是有些不相同的.
這裡更加是一個歷史的敘事文體.
如果再簡單,再清楚一點.
就是很清楚告訴我們.
神曾經在歷史裡面.
他進入人類的歷史.
藉著摩西,藉著以色列這個民族.
來進行那個拯救.
而這個拯救就是他的計劃.
所以在神學家的描述裡面.

$^{41}$這個叫做歷史的敘事文體.
所以我們大概有一個這樣的觀念.
我們這樣看的時候.
我們就能夠從裡面的佈局.
從經文的發展.
我們就能夠看到神的救贖.
所以我們就從這個角度去看《出埃及記》.
其實很好看的.
還有你會看到神的手不斷出現.
當然他也透過摩西.
透過他被他揀選的人來成為代言人.
這個你可以告訴我們.
就是在其後的舊約.
特別是《先知書》.
神就正正式式有一些很清楚.
神揀選的代言人.
所以這些就是先知的角色.
所以這個《出埃及記》.
其實整個的主題是很簡單的.
因為它是剛剛坐落在以色列人.
透過阿富汗.
接著就是以色 雅各.
雅各的兒子 第四代了.
第四代就是約瑟.
當然約瑟還有11個兄弟.
但是以約瑟為重要的單位.
約瑟就成為了埃及的宰相.
拯救了那個時代.
因為大饑荒而來的拯救.
所以他因為約瑟的緣故.
這個民族 以色列人或者希伯來人.
就可以在埃及很安然的居住.
安然居住又怎麼樣呢?.
這個是不是一個結局呢?.
當然一定不是.
所以就是出埃及 就是離開埃及.
這個Exodus就是希伯來文.
或者叫做離開.
所以如果原文的《巴西五經》.
每一個一卷書的開頭就是頭一句.

$^{81}$所以離開就是出埃及記.
所以也是對的.
用中文的我們的和合本.
出埃及記是一個敘事的記錄.
所以這裡告訴我們.
就是從埃及離開進入迦南應許之地.
大體上這一段書.
你可以有很簡單的三個大的段落.
一到第十八章就是提到神的救贖.
當然就是藉著摩西.
這裡就有好幾章.
一到第四章就是揀選摩西.
接著就開始整理好摩西的心情.
他願意成為代言人.
又敢於回到埃及.
和當時的法老來討價還價.
來帶領自己的族民離開埃及.
去到應許之地.
接著應該就是法老的抵抗.
出爾反爾 不肯接受.
就帶來神要用那個十災來對付埃及.
輾轉以色列人就要過紅海.
接著曠野的情況.
所以這裡第十八章.
就是他們去到西乃山腳的時候.
所以第二個大段落.
又是很清楚的.
就是神和以色列人.
以十災為最重要的標準.
那個綱領來立約.
就是說以色列人不再是一個奴隸的國度.
接著就成為神兒女的國度.
一直去到二十四章.
二十五到四十章就有兩個段落.
應該是三十二章做那個分界.
就是摩西上西乃山.
其實很簡單的.
二十五到四十章.
我們用同在.
什麼叫同在呢?.

$^{121}$就是神既然選了希伯來這個民族.
神就應許在他們日加拿.
以及將來都與這個民族同在.
神的同在.
所以神的同在.
他講了就可以了.
但不是的.
因為神要住在百姓當中.
就需要有一個回望.
需要有一些保護.
不是保護神.
是保護以色列人.
因為神的同在.
就帶來他們能夠親近這位聖潔的神.
所以就需要做一個回望.
所以由二十五章到四十章.
是在講神頒佈了那些規條.
到底那個回望要怎樣做.
摩西就領受了那些在山上的設計圖之後.
下山之後就突然間三十二章.
就有一個金流瀑的事件.
就擾攘了一段時間.
後來重新再講.
這次就不是神講的那個設計圖.
而是摩西找了兩個人.
給撒列,阿荷利,阿伯兩個設計師.
可以做很精細的手工的.
這兩個設計師.
來帶領著回眾做好那個回望.
當然要收集很多的物資.
金銀銅或者一些的懷墨.
那些的卷,那些的布,海狗皮等等.
還有一些木樁等等.
這樣來收集就做好了回望.
所以整個《出埃及記》.
就是大概上是這三個大卷.
其實你已經掌握到整本書的脈絡和氣氛.
就是說神救了這個民族.
神要這個民族從此有一個新的身份定位.
就是神的選民.

$^{161}$就在眾多當時的加拿大拜偶像的民族當中.
他們不是特別的標籤.
不是特別的聰明.
更加不是特別的敬虔.
都是同一般的人一樣.
不過只因為有他們的祖宗亞伯拉罕.
曾經遇見神 神帶他離開這個加拿大的梧椅.
就是有法拉底河 底加里斯河的那個地方.
輾轉就去到以色列的地方來去寄居.
所以那個約就是從亞伯拉罕開始的.
所以現在你可以想像得到.
《出埃及記》就等於是神要找數.
神的時間到 神的那個計劃的時間到.
神就真是以整個的以色列民族.
來去拯救他們 離開了這個維羅斯加.
就進入加拿大.
所以他們不是特別給加拿人特別的蒙恩.
特別的蒙愛 特別的與別不同.
而是神定義透過亞伯拉罕的子孫.
來去施恩給萬族萬民.
所以這個也是一個很重要的.
在舊約的重要的信息.
所以很難怪就是加拿大書第三章.
應該是十九節 就說我們信了主耶穌.
在耶穌裡面 我們都是亞伯拉罕的子孫.
很特別的 其實同外邦人沒有需要這樣說.
為什麼要這樣說呢? 當然就是因為加拿大書.
有些猶太人還是很死硬.
就是說你們不行過禮 不做一些事情.
就不像我們這樣蒙恩的.
我們是亞伯拉罕的後裔等等.
保羅說其實不是 恩訊就可以得到恩典了.
不過我們恩訊.
我們就是屬靈的亞伯拉罕的子孫.
這樣的意思 第三章就是加拿大書.
可以和大家講一講這段經文.
因為我們都是會看亞伯拉罕.
一直都是出埃及 一定是從很遠古的立約.
來到出埃及的時候.
就兌現 就是說神的時間到.

$^{201}$他就進入這個以色列人的那個歷史.
所以不是以色列人的發跡史.
是神的救贖的歷史.
所以大家一定要搞清楚我們看聖經.
就不是以色列人的一些他們的過去.
例如美國史 例如歐洲史.
例如我們的香港歷史.
不是 是神救贖的計劃的那個歷史.
第三章 看看哪裡 我們不要看那麼多.
可能第26節 加拿大書的三章第26節.
他說 所以 當然這是回應上文講的.
就不是我們行律法就可以得到.
律法是我們奎摩的師父來的.
引我們去到基督那裡使我們能夠恩信稱義.
第25節 這個恩信得救的你既然來到.
我們從此就不再是師父的手下.
就是說我們不需要再守律法了.
律法是引起我們一個的認識.
對自己的身心靈的認識.
對自己能夠看到自己的一個罪性的認識.
看到自己不能夠滿足律法的要求的那個認識.
以致使我們去呼求耶穌的那個救恩.
耶穌的拯救 所以第26節.
就明顯是一個結論.
他說 所以 就是說如果我們認識自己的時候.
我們就會走到主那裡 以信來得到祂的拯救.
所以你們因信基督耶穌 都是臣的兒子.
你們受贈不入基督的 都是披戴基督了.
並不分猶太人 希臘尼人 自主 圍盧 或者男或者女.
因為你們在基督耶穌裡 都成為一了.
這個差不多是一個非常不能再改變的一個永恆的真理.
任何人在基督裡面是新造的人.
不是他自己一個人 是所有信耶穌的人.
都在神國度 神眼中 神的救贖裡面.
我們就是完全成為一.
就是說在主的眼中.
再沒有分我們地上的這些身份.
民族的階層 或者是性別等等.
而是我們大家在神的眼中.
我們有不同的恩賜 我們有不同的出身.

$^{241}$但是我們在救恩裡面 我們都是神兒女這個身份.
所以都是神的兒子 這個身份.
當然所以是雙重性 因為我們還活在這個地球.
還在我們那個各樣的限制底下.
但是我們在神眼中 所以我們在教會.
一個有形體的教會 就是教堂.
我們去到禮拜堂 神就是要我們不斷去訓練自己.
既是六天的一個工作的人.
一個家庭的人 一個在社會上的公民.
但是一進入到教堂 其實這個是教會.
就是在神眼中的那個普世的那個宿命身體.
我們就是神的兒子.
所以同時具備這兩個身份.
所以為什麼往往很多的抗爭.
就是我們很多時候沒有刻意去學習.
看自己是神兒子 看其他人都是神兒子.
這個課題出了一些亂.
所以我們很多時候就會有比較.
就會有很多很多不同的問題出現.
但是其實如果我們這樣明白的時候.
我們就不需要看自己過大或者過小.
但是其實講到這裡已經足夠了.
因為加泰的教會的人.
他們只需要這些分別.
就是說在救恩裡面不需要再有分別了.
怎樣可以得到救恩呢?不是去換取.
不是靠自己 而是在救恩裡面.
我們就立即是神的兒子了.
我們大家都是平等的.
在神眼中的那個領袖的身份.
不是那個恩賜 不是那個社會地位.
不是那個我們與生俱來的那個身份.
我經常說我是出生在務工區長大.
家裡是有兩個媽媽的.
那個成長是非常惡劣的.
周邊的環境惡劣.
長大已經有附近都有不同的黑社會的字頭的.
那些朋友當然不是好朋友.
都是我哥那些隔壁的鄰居.
都知道他是黑社會的.

$^{281}$很多有紋身 又打架.
在這樣的環境長大.
這樣的環境再加上自己家裡的環境都是很複雜.
但是無論如何在救恩裡面.
整個曾蕊桂這個人就是神的兒子了.
所以你越看自己這個身份.
那個寶貴 那個平安 那個盼望就會越來越大了.
還有要這樣去看其他人的時候.
你就會學會尊重了.
不同的弟兄姊妹有他們的那個神愛他們.
神擺在你面前和你一起成長的這樣的身份.
這樣的學習.
這樣看的時候我們就會很立體看到實的榮耀.
不然的話 每次都看自己 那就自卑了.
或者自大了 或者我們稍微厲害個人一點.
我們就會有很多的不可一世 或者很自我.
這些全部都是和救恩是會有一個極黨的.
但是必須告訴我們 其實已經講完就夠了.
但是29節就很特別了.
他說 你們既然現在屬乎基督.
就是亞伯拉罕的後裔.
是照著應許承受產業.
這裡就已經給了我們兩個很重要的向導了.
其實根本不需要講我們是亞伯拉罕的後裔.
實在不是 我們不是以色列人.
我們不是希伯來人.
但是保羅告訴加拿大教會的人知道.
我們原來每一個 無論什麼民族都好.
原來在救恩裡面 神預備的耶穌基督.
就使得我們是神的兒子 這個我們明白.
但是原來 當我們既然現在屬基督.
我們就是亞伯拉罕的後裔.
就是說 亞伯拉罕那個真蒙應許的兒子.
不是以色列 原來是耶穌基督.
就是說 神對亞伯拉罕的那個應許.
從來沒有停止 甚至不是只是給猶太人.
這樣就很特別 我們基督徒就比較少這樣去想.
自己是屬靈的 永恆國度裡面的.
那個亞伯拉罕的子孫.
所以原來我們既是神的兒女.

$^{321}$我們原來又是亞伯拉罕的子孫.
亞伯拉罕的子孫 其中一樣東西就很重要了.
也都是《春拉及記》的主題.
我們就是要進入神應許之地.
所以我們今天不是地上的應許.
是將來永恆的屬靈產業.
所以我們應該很期待將來我們是有產業的.
所以保羅也是這樣說的.
他說 我們信了主耶穌.
《聶不鎖宣》第一章是這樣說的.
我們將來是有屬天的基業的.
當然他沒有說很詳細.
到底可不可以再具體一些.
是樓宇,樓房,官面,還是王國.
他也不知道 或者說不出這麼詳細.
《聶不鎖宣》的第一章.
由第十三節開始 一章的十三節 《聶不鎖宣》.
你們既聽見真理的道.
就是信耶穌 報道會.
或者有人和我們個別談道 我們信主.
就是那叫你們得救的福音.
所以真理的道 就是福音 就是耶穌.
也信了基督.
所以保羅是很清楚的.
無論用什麼入手都好.
無論談論創造也好 罪惡也好.
或者這個世界是很虛空也好.
最終是要帶領那個決志者.
他知道自己需要迎順基督 向主悔改.
順了基督 既然順他.
就受了所應許的聖靈為印記.
又要讓他知道 他一打開他的心靈相信.
這個相信帶來一個非常巨大的改變.
就是聖靈會進入他的內心.
所以他裡面有一個心 和有一個聖靈.
很奇妙地 聖靈就會永遠住在他裡面.
一直來孕育他的成長.
這個就是聖靈的印記.
第十四節.
「這聖靈是我們得基業的群聚.

$^{361}$直等到神之民被贖.
史丹利榮耀 得著稱讚」.
原來我們有聖靈作為一個訂金.
作為一個群聚.
作為一個我們去買東西的時候.
一個很重要的一個大的訂金.
就是告訴我們.
這一位買家一定會兌現他的承諾.
就是他一定會到時完成那個交易.
現在我們有什麼交易呢?.
就是我們是得到基業的一個群聚.
聖靈就住在我們裡面.
就是說神在這個地球上救贖了我們.
祂一定會帶我們進入那個永恆.
得著將來永恆的基業.
所以基督徒要期望的.
是我們得著這個永恆將來的基業.
而且最重要的不是我們有多大的儲蓄.
而是要使到神的榮耀.
得著稱讚.
在永恆裡面.
還是主永遠得著那個榮耀,得著地上救贖的人.
對他的歌頌,對他的敬拜,對他的誇耀.
所以這樣就給大家有一個概念.
原來舊約講的應許.
新約就不是再是物質的以色列地.
不是現在的以色列的土地.
而是我們將來是有一個永恆的基業.
是基督徒.
而且聖靈是一直告訴我們.
那個就是神已經附上了的病根.
神的三位一體的神.
三一的真神的聖靈.
其實是神來的.
但是就住在我們每一個信耶穌的人的心裡.
無論你是一個不識字的婆婆.
或者一個很厲害,很有學問的教授,博士.
或者你是一個很擁有權力的一個君王等等.
一個跨國企業的CEO也好.
同樣都是得到這位聖靈.

$^{401}$是我們將來得到基業的病根.
保羅這樣理解.
不是貶低聖靈.
而是神自己的救贖的計劃.
使我們能夠有一個安心.
有一個渴慕.
我們將來就是追求這個永恆的基業.
所以我們進入春海球記.
我們就很不同了.
因為我們是新約的人.
我們現在有機會看回舊約.
當時代摩西的時代.
他們是沒有明天發生.
他們知道的.
他們沒有水晶球的.
我們就有了.
我們已經回去看回.
所以我們某程度上挺有優勢的.
去看回他們.
所以間古之間.
我們現在保羅告訴我們.
我們是同樣的救贖的子民.
這就意味著埃及在屬靈的層面.
就是我們舊有的過去的罪性.
我們過去的罪惡的場景.
就是可能我們的生活裡面.
有不同形式的罪.
有些你的地方尊貴一點的.
乾淨一點的.
文明一點的.
可能你的罪就沒有那麼明顯.
比如你在一些黑社會打滾.
晚上去一些夜場.
那些明顯就是每晚都爆瓶.
每晚都爆粗打架.
那些就比較明顯一點.
或者是販毒.
或者是逼良為娼.
比如潑流街.
在那裡生活的人就會藍白一點.

$^{441}$但是不代表.
只有那些人才是遺老之家.
其實我們未信耶穌就是遺老之家的人.
我們是被救贖出來的人.
所以《出埃及記》正正就是在說神的救贖.
所以去到十二章.
有羔羊的寶血.
大家在身上有很多這一類的經文.
我們這一位無殘疾的基督.
已經被殺獻祭了.
這個就已經告訴我們.
我們其實都是需要那個救贖.
所以原來新約其實已經在反映.
《出埃及》有很多很多這些經文.
新約的作者尤其是保羅.
差不多很多的概念.
都是以《出埃及》的那些場景.
來作為他寫的那些教義.
你可以說是靈感.
可能更加是神聖靈給他的感動.
就不單純是他自己.
因為對舊約滾瓜爛熟.
所以他就隨便可以駕馭那些經文.
其實就未必是.
反而就是固然應該這麼說.
舊約固然他真的很熟.
但是當他要寫到這些事情的時候.
他有新的靈感.
我相信不是他自己特別爆發的那些靈感.
而是聖靈在他裡面.
輔助著他的內心本身有的知識.
來到他好像共野一路那樣.
實際怎麼運作呢?.
要回去問保羅才知道.
譬如在《哥羅格前書》第十章.
就說過鴻海 很特別的.
所以你看看.
原來在《出埃及記》.
是一個很重要的一個母體.
因為舊約往後的經文.

$^{481}$就是《先知時代》.
到連耶穌的十字架.
他是無塵世的羔羊.
到使徒們.
都是以出埃及這些舊屬作為物本的.
《哥羅格前書》的第十章.
所以其實舊約都不是不重要的.
因為如果沒有舊約.
我們就不能夠明白新約很多的經文.
第十章的第一節.
你看到他怎麼用屬靈的那種詩意.
去說鴻海.
第十章明明是在說.
他很清楚的告訴我們.
是在說我們不可以.
因為當時哥林多教會.
一些頂子們以前是拜開偶像.
其實個個都拜的.
因為城邦是不能不拜的.
個個都要拜.
那些市面上賣的玉全部都拜偶像.
所以除非你不生存.
但是第十章就說.
弟兄們 我不願意你們不曉得.
我們的祖宗.
就是以色列和出埃及.
十二章,十三章那裡.
我們的祖宗從前都在雲下.
都從海中經過.
都在雲裡,海裡受洗或受咒.
歸了摩西.
這樣就很靈異了.
到底保羅想暗示什麼呢?.
並且似乎經過一個.
其實那個是審判來的.
是那些埃及的追兵全死了.
覆滅他們 把那些水合起來.
將所有的追兵.
記得是埃及那些勇猛的軍隊.
不是平民百姓.

$^{521}$軍隊是善戰的.
就是說 你想一下那些海軍陸戰隊.
你想一下那些最精銳的部隊.
是要在海邊抱著木頭.
24個小時要磨練你的意志.
多疲倦 那些海浪怎麼衝擊.
你都不放棄.
一放棄 你就要被淘汰了.
那些真的是特種部隊.
所以埃及的追兵.
都不是隨便的善男信女.
但是那個所謂被雲裡,海裡受洗.
其實是一個審判.
是一個神的憤怒的審判.
要救的是以色列人.
平安地過了那個乾地.
接著那些海水就蓋在一起.
就將那些追兵 還騎著那些戰車.
全部就殺了.
所以這裡說歸了摩西.
就是屬靈的含義.
摩西就等同於那個救主那樣.
留意就是三節了.
並且都吃了一樣的零食.
也都喝了一樣的零水.
一直都在說一種屬靈的暗示.
就是說他們吃的馬蘭.
明顯是吃馬蘭 不是吃零食.
喝的不是零水.
是那個馬拉的苦水.
第16章 《出海特記》.
其實以色列人就發怨言.
那些苦水苦的怎麼辦?.
馬拉就是苦的意思.
那麼摩西就向神祈禱.
因為百姓在那裡苦苦哀求.
神就指示摩西.
讓他看到有一棵樹.
他就斬了那棵樹.
扔在那些馬拉的水裡.

$^{561}$那些泉水裡就變甜了.
甜就是辣鵝米.
這裡告訴我們.
所害的是出於隨著他們的靈門石.
那門石就是基督.
非常靈異 阿保羅.
很誇張的 那個盆石.
盆石就是硬的 是固定的.
為什麼你搬著盆石跟著那些百姓呢?.
其實是那些百姓一路走.
準備進入這個西奈半島.
準備進入迦南.
就告訴我們.
原來一直追隨著他們.
有得吃 有得喝.
是這位眾物主 是一位好像是能夠.
不是一個不動的偶像.
是一個一直追隨著他們的基督.
所以相當靈異 阿保羅.
很誇張 將一些本來就是歷史的一些事情.
本來不是開心的事情.
是百姓不斷發怨言.
我們沒東西吃了.
我們在埃及都有黃瓜,韭菜.
現在沒得吃了.
其實一直都是發怨言的.
其實所以保羅 你不要開玩笑.
保羅其實就是借這段經文.
來去責備他們那個問題.
所以你看到第四節.
但他們中間多半是神不喜歡的人.
所以再曠野倒閉.
就是說保羅用四節聖經.
就將整個出埃及的事件.
來做一個濃縮的一個靈異解釋.
就是說整體來說.
40年的出埃及的那個歷史.
就是因為不斷的發怨言.
因為不斷的不信神.
因為仍然懷戀著舊有的那個.

$^{601}$埃及的那些所謂美好.
就是說情願做奴隸.
現在你都不知道帶我們去哪裡.
或者說這一位神那麼多的掣肘.
那麼多的繁文縟節.
我們不喜歡 我們不喜歡摩西里.
所以就帶來就有這樣的倒閉.
是神所不喜歡的人.
第六節那個dare for就出現在這裡.
保羅講了一大段的出埃及的歷史.
好像幫我們做了一個功課一樣.
給我們大概知道出埃及的那些事情.
當然橫誇的還有萌素記.
就不在我們這一次的言經稍微的範圍.
但是始終都會涉略到一點點的.
接著第六節.
重點就在這裡.
這些事都是我們的尷尬.
叫我們不要貪戀惡事.
像貪門就是以色列人那樣貪戀的.
也不要拜偶像.
像貪門有人拜的.
如經上所記 百姓坐下吃炕 起來挖沙.
果然就是出埃及的三十二章的拜金牛毒的事件.
所以那個同在二十五到四十章的建造回望.
中間那個分水嶺就是三十二章.
這裡說就是這樣逼阿倫來拜偶像.
第八節 我們也不要行奸淫.
像貪門有人行的.
一天就倒閉了二萬三千人.
這個就是那個菲尼蝦的事件.
就是那個十平的事件.
就是經文述記第二十章的經文.
告訴我們知道的.
接著也不要試探主.
像貪門有人試探的.
就被蛇鎖滅.
那個是同蛇的事件.
所以又是跟文述記有關.
你們也不要發怨言.

$^{641}$像貪門有人發怨言的.
就被滅名的鎖滅.
貪門遭遇這些事.
都要作為尷尬.
並且寫在經上.
正是警戒我們這末世的人.
所以自己以為站得穩的.
站得穩就是說.
自己比那些以色列人更加的.
我們不會這樣.
但實際上我們的行徑都是這麼軟弱.
我們有沒有一種警醒,謹慎的心呢?.
那就會跌倒了.
所以你們所遇見的試探.
就是說如果我們很謹慎地看自己.
我們就會在神那裡有那個力量.
因為我們不斷去依靠神.
在我們的生活裡.
不敢自誇.
也不敢去做一些神喜悅的事.
所以你們所遇見的試探.
無非是人所能領受.
神是信實的.
就是說既然他救贖我們.
他一定會讓我們在我們人生不同的際遇.
我們都能夠經過.
國書也提到.
我們千萬不要以為神會被惡試探.
因為他不會試探人.
也不會被惡試探.
因為各樣的美善的恩賜.
都是從這位重光之父而來的.
所以這件事講了這麼多的警告.
其實再一次想突顯.
神其實是信實的主.
他應許了我們.
他一定會成就.
所以果然就是出埃及.
很多的百姓基本上大部份都死了.
只剩下約書馬和加林.

$^{681}$可以進入加林.
所以我們記得這裡.
十三節其實是相當重要的.
這是一個整個段落的警告之後的安慰.
受試探的時候.
總要給你們開一條出路.
叫你們能忍受得住.
接著就講實際的警告,血面和牧羊.
叫哥林多教會千萬不要拜牛雀.
就好像他們這樣拜的.
還有在聖餐裡面千萬不要混亂.
所以我們隨便已經拿到出埃及的這些精髓.
保羅都會這樣去借一些歷史的事件.
作為牧羊的時候.
對哥林多教會可以順手拿來.
將它轉化成為一個屬靈的教導,血面.
甚至將過河海都成為一個屬靈的新的體驗.
就是說原來我們都一樣在主的奉庇底下.
我們已經脫離了那個奴隸.
罪的奴隸的身份,罪的世界.
我們已經在主的救贖裡面.
我們要準備進入永恆的天家.
不過原來都危機重重的.
我們又會感覺到屬靈的枯竭,淡而無味.
我們經常都覺得自己乾旱.
我們就要找一些吃和喝.
所以我們就慢慢用舊有的價值觀.
追求物質的世界,我們的享樂.
追求我們人生的.
我們要滿足我們自己很多的內心裡面的慾望.
很多的慾念.
我們很多的價值觀.
時時都喜歡爭鬥.
時時都喜歡嫉妒.
時時都喜歡來跟別人比較.
或者不抵得別人的屬靈過我們.
我們會用很多這些東西來包裝自己.
所以原來以色列人都是吃過這些零食.
喝過這些零水 也有盆石隨從他們.
所以我們盼望我們真是每一次來敬拜神.

$^{721}$我們要記得我們真是這個身份.
不是地上的身份不真實.
而是的而且確是很真實的.
地上的身份.
譬如我這幾天喉嚨不舒服.
說話都要遷就著.
都晚餐吃那些喉糖等等.
要喝琵琶糕這樣.
就是很真實的.
你有一點點不舒服.
你都不能夠這麼自然的說話.
說久了就會咳嗽.
但是我們更加要認定.
我們有屬靈永恆的身份.
所以如果我們的價值觀只是停留在地上.
而不是有衣有食就當知足.
以全心全力去服侍神.
我們就一定會有一個限制.
不是打折扣 再嚴重一點.
如果我們的自我再膨脹一點.
這個靈盤石隨從我們的這個靈盤石基督.
有時都會醒我們而去.
其實是我們寫主而去.
就不是主醒我們而去.
你想想那個回望一直都在以色列當中.
那些人一個一個的倒閉.
有金牛毒事件過了.
過了這一站.
接著還有銅蛇的事件.
所以其實外敵不是那些阿摩利人.
不是那些格加薩人.
不是那些巨人.
那些巨人反而不會傷害到他們.
反而就是我們自己貪戀.
世俗就會出現這些問題.
所以新約很好幫我們做了一些屬靈的解釋.
原來當我們稍後進入那個經文的世界.
在神學的敘事,救贖的歷史.
我們透過神給我們的視覺.
我們進入那個歷史裡面.

$^{761}$我們看的那些只不過是實實在在.
神揀選摩西 進入回到埃及.
來和埃及的法老王對質的時候.
這些全部是很表面的經文.
但是原來新約的作者.
小勞保羅他有這樣的屬靈的眼光和慧憤.
借出埃及來責備哥林多教會.
因為哥林多教會其實信得很好.
但是問題就是他們有以色列人的影子.
其實每個人都有這個影子的.
信了主耶穌仍然是一個離開了物質的埃及.
但是心靈的埃及就沒有離開.
所以他們就會有這種不斷的懷戀.
求神幫助我們.
真是我們要越來越所謂定性.
信主越來應該是越單純.
對主越來越有一個熱切的愛慕.
對於未信的人的那個靈魂.
我們真是有一個心 有一個焦急.
這樣來到很想看神給我們的恩賜.
盡量的運用.
好 我們很快的 隨便的選幾段新約給大家看.
其實你就會看到原來出埃及是這麼精彩的一件事.
說出埃及的說明是歷史的神學.
舊熟的歷史的神學.
你不能夠不看《創世記》最後的一段經文.
就是一個很重要的 差不多是一個誓言.
是約瑟要那個時代的兄弟.
那些長老要嚴嚴起一個毒誓.
就是我的骸骨不可以落葬在埃及.
其實當時的宰相雖然有新的皇兄氣.
已經不認識了 但是約瑟在最風光的時候.
那個埃及王法老最尊敬他的時候.
他就110歲 安然的離世.
但是居然嚴嚴的警告自己的兄弟.
不可以將他葬在埃及地.
因為當時他們就住在埃及地.
你就可以想象到 這個就是一個伏筆.
這個就是出埃及記的那個重要的開始.
跟大家看一段經文.

$^{801}$接著我們就再看回剛才那段.
這個應該是 如果我沒有記錯.
應該是約瑟拿記的 我看看是不是十二章.
因為我沒有記下來.
可能我不敢講 我有沒有記錯.
約瑟拿記的.
如果我找不到 可能我先回去再找給大家.
可能這裡沒有了 十二章.
那看看是不是 我很快看一下.
對不起 應該是二十四章才對.
三十二節.
或者三十一節.
約瑟拿在世和約瑟拿死後.
當然後記 因為不是他寫的.
那些知道耶和華為以色列人所行諸事的長老.
還在的時候 以色列人事奉耶和華.
第三十二節 大家可以記住.
以色列人從埃及所帶來約瑟的骸骨.
裝埋在士劍 就是在那國從前.
用一百塊銀子向士劍的父親哈末的子孫.
所買的那塊地雷.
如果我沒有記錯 就是三十八章還是三十九章.
所買的那塊地雷.
這就作了約瑟子孫的產業 就是士劍.
接著亞倫的兒子以尼亞薩都死了.
就把它葬在他兒子菲林哈以法連山地所得的小山上.
約瑟拿記的最後的一段.
都要確認他們進了迦南之後.
住了一段長時間.
真正帶領以色列人進迦南的.
就是約書亞就是摩西的幫手.
就是第五代 真的進了迦南之後.
真的將約瑟的骸骨.
葬在這個士劍那國應許的那個地方.
預備好的那個墳地.
就意味著原來亞伯拉罕那個遺願.
對亞伯拉罕和他的子孫的那個祝福.
那個的應許 一直這樣對峙.
所以以色列人對這個應許.
從來不敢掉以輕心.

$^{841}$所以有信心的希伯來人.
就會承受這個應許.
所以亞伯拉罕以薩那國.
所以不是耶穌在那國承受那個應許.
所以耶穌是不信的.
就算他是以薩的兒子.
他都沒有福氣承受應許.
因為他自己就是一個不信的人.
根本不是稀罕神的應許.
也都不認為這些應許有什麼重要.
但是那國就很看重.
就是說他雖然他自己一生出來.
他不是命不好.
是他裡面一樣那個罪性很強.
非常狡猾 但是神都可以收服他.
讓他這個人能夠具備那個屬靈生命的改變.
來承受那個應許.
所以是亞伯拉罕以薩那國的神.
所以又是可圈可點的神.
讓我們看到有信心的人.
就會接上神的應許.
所以如果我們每一個弟兄姊妹.
真的看到我們無恩信主.
原來是這麼美麗,這麼寶貴.
我們的生命就不斷改變.
真的越來越靚仔靚女.
應該說那個品性裡面的生命流露出來.
真的有一份的愛.
對人有一份的誠信.
對人有一份的接納.
對人會有一份慈悲.
因為神的聖靈會改變我們的生命.
所以你看到兩兄弟 孖仔一起長大.
我們廣東話說這麼高這麼大.
但是以數和雅各兩個是完全詭異的人生.
一個是得不到任何的祝福.
因為他根本就沒有那個信心.
他是一個不信的人.
但是雅各由他自己貪戀.
到慢慢他自己整個生命.

$^{881}$就能夠有這樣的福氣.
承受神的應許.
所以藥室的骸骨是相當重要的一個瀑布.
就是出埃及記 他們就帶著.
所以出埃及記應該是第十二章.
他們就帶著這些骸骨.
很快的記一下 出埃及記是第十二章.
對不起 應該是第十三章.
因為也是要再記這些經文.
第十七節.
「白老容百姓去的時候.
非利士地的道路須近.
神卻不領他們從那裡走.
因為神說恐怕百姓遇見打仗後悔.
就回埃及去.
所以神領百姓繞道而行.
走紅海曠野的路」.
就是走陀螺路的 因為一定要經過紅海.
「以色列人出埃及地.
都帶著兵器上去」第十九節.
這些經文都是要大家記住.
偶然間這樣記載.
似乎不用這麼誇張.
你現在趕著逃亡.
你還要拿著那個老祖宗的骸骨.
應該放在金塔 而不是用棺木埋在地底.
一直都是用金塔裝著那些骸骨.
這樣等待出埃及.
那些骸骨放在那裡.
就好像刺眼刺鼻.
以色列人什麼時候可以出埃及?.
可能大部分人都忘記了.
因為越來越艱難.
新的法老出現 新王出現.
就越來越壓迫那些百姓.
還說什麼出埃及.
現在連命都沒有了.
油焦滾焦不得了.
連命都沒有了.
每天都做苦工.

$^{921}$已經沒有人記得那些什麼棺木 什麼應許.
大概如果我們能夠所謂穿越時空.
回去出埃及的時代.
我們就可能看到很多的仇人慘墓.
大家都不知道怎麼辦.
哪個公民領袖可以帶領我們.
可以和我們一起對抗這些法老.
我們那個Good old days 弱勢的年代.
弱勢做宰相的時候.
我們那個社會地位多高.
讓我們住在那個歌山地.
又可以牧羊的 很豐富的.
現在就慘了.
他們都渴望著一位的救主出現.
他們都渴望著.
是不是真的他們祖宗的神.
來介入他們的痛苦呢?.
所以這些都是我們信耶穌.
個別丙姐妹信耶穌的時候的那些心路歷程.
所以有時候我們真的浸禮.
雖然你就算沒有丙姐妹你認識的來浸.
當然大家都忙碌.
但是甚少人是不忙碌.
他都不去浸禮的聚會.
其實你應該爭著去做一些的崗位.
即便是鋪鋪布 鋪鋪毛巾.
其實重要的是.
你看到有新信主的弟兄姊妹.
他們講見證.
神怎樣在他們還沒有得救之前.
預備他們的心靈.
其實偶然間你會聽到有些人落在一種很大的困苦裡面.
其實就是出埃及記的那個生命.
但是不要忘記.
原來真正的那個最重要的神的應許.
就是那副不是那個骷髏骨頭.
不是雕本.
是那個刺眼,刺鼻的那個金塔.
就是弱身的骸骨.
大家保住它 供奉著它.

$^{961}$沒有人敢來撞它.
大家知道這是一個局勢.
但是這裡原來就是神應許的一個人記.
擺了給你們.
時候還沒有到 等吧.
所以以色列人 人始終是人.
等下又是一代 等下又是一代.
越來越艱難就發怨言.
就向天發怨言.
這個壞天 這位神都不是幫我們.
就是類似這樣.
但是神真的對人類的愛就是.
仍然是這樣人大量集體的來控訴.
但是神都能夠將它可以轉化成為.
聽到是百姓的哀號.
所以十三章的十九節.
摩西把弱身的骸骨一同帶去.
因為弱身曾叫以色列人嚴嚴的起誓.
對他們說 神必眷顧你們.
你們要把我的骸骨從這裡一同帶上去.
這裡就是出埃及的歷史的起點.
也是創世記的伏筆.
所以摩西五經是一個連續劇.
是有五卷 所以叫做五經.
我們要回到這個讀世.
今晚最主要是開始要講這裡.
所以讀書讀記必須要認識弱身的棺木.
第五十章第十五節.
五十章是創世記 對不起.
五十章是創世記第十五節.
弱身的哥哥們見父親死了就說.
或者弱身懷恨我們.
照著我們從前代他一切的惡.
足足的報復我們.
他們就打發人去見弱身說.
你父親未死以先吩咐說.
你們要對弱身這樣說.
從前你哥哥們惡待你.
求你饒恕他們的過犯和罪惡.
如今求你饒恕你父親神之僕人的過犯.

$^{1001}$丹門對弱身說者說.
弱身就哭了.
因為太低估了這個屬靈的弱身.
他能夠解夢 能夠有屬靈的智慧.
所以千萬不要看屬靈人.
用自己的標準去看屬靈人.
這是很重要的.
因為人家屬靈的人.
看的事物看的事情比我們更加闊.
比我們更加有實的心腸.
所以我們不明白可以問.
千萬不要以為弱身說.
我爸爸死了就會報復了.
地上的人因為真的是這樣.
父母離世家就會散.
只顧及自己的家庭.
自己成家立室的家庭.
各護各家.
所以就一定會有這樣的想法.
但是我們要看.
不要用地上的價值觀去看.
這些屬靈的人.
所以弱身就哭了.
第十八節.
他的哥哥們又來俯伏在他面前說.
我們是你的僕人.
弱身對他們說.
不要害怕.
我豈能代替神呢.
所以屬靈人看事物是不同的.
他明白了神的計劃.
在他自己弱身的生命裡面.
是很慘的.
被人賣豬仔賣去做奴隸.
然後要坐牢.
然後幾乎死.
但是突然之間一夜就成為.
由一個監犯變成法老的宰相.
戴的是擁有十權的戒指.
運籌維握.

$^{1041}$而且是為了要他.
不是讓他做大官.
是為了他要拯救其他人.
所以原來神那個裝備.
是帶來弱身.
真是很慘.
一個人真是叫天不應的痛苦.
其實那二十年都不是說笑的.
有二十年這樣的磨練.
所以如果我們對一個已經磨練了二十年.
在神裡面這樣來將他訓練成材的人.
我們用一個這麼人性的水平看人家.
你想想我們那個屬靈眼光是多麼的弱.
多麼的弱雞.
人家屬靈人不是這樣看事物的.
弱身是說你們不要怕.
我怎麼能夠代替神呢?.
從前你們的意思是要害我.
但是神的意思就是說.
祂在經過了這麼多年的生命的成長.
四年在祂內心的工作.
祂現在真是一個屬靈人看透萬事.
祂知道原來祂是神手中的一個拯救者.
祂是神在那個時代的一個很重要的一個人.
所以弱身已經變成屬靈人了.
所以祂說的話看到是神的心意.
既然現在是神的心意.
那就是說那些害祂的兄弟.
真是原原本本是因為嫉妒,是因為激氣.
是因為分家產,是因為弱身很囂張.
最小的時候因為別男未出世.
所以拉子,拉心肝.
那些兄弟就會這樣看了.
但是弱身當時也是這樣看.
當時也是很激氣,很不明白.
但是經過了生命的歷練.
所以我們人生信耶穌最寶貴的.
不只是洗完禮之後就再見了.
跟教會拜拜.
那是很可惜的.

$^{1081}$信了主耶穌之後.
才開始神在我們生命裡面的不斷的裝備.
不斷給我們人生不同的際遇,不同的磨練.
讓我們成為屬靈人.
有屬靈的眼光的人.
不是變成不可一世,目中無人的屬靈人.
硬邦邦,望人非常嚴肅,為真理.
其實不是 應該是越來越隨和.
越來越有主的愛,有主的寬容.
但是對一些事情.
我們也會有我們的看法,價值觀等等.
你可以想像得到.
這句話真的非常有份量.
就是一個真的很悲慘的人生的弱失.
跟爸爸,跟父家已經是完全的割裂了.
差不多是二十年.
來到他自己成為一個被擄到很遠方的一個人.
就是說他的身世真的是非常可憐的.
但是他看到神的計劃.
要保存許多人的性命.
成就今天的光景.
所以現在你們不要害怕.
我便養活你們.
和你們的婦人,孩子.
於是若失用親愛的話.
安慰他們.
所以真的是屬靈人.
說出的是真的非常寶貴的生命的說話.
所以我們真的要去追求.
我們真的要好好讀經.
神將來用我們每一個弟兄姊妹.
因為經驗過不同的人生歷練.
的確是一個有經驗的人.
是一個過來人說.
就不是變成一種頭腦的知識.
因為用生命經驗去講出來.
我們的信心,我們的愛心,我們的盼望.
都應該是有體驗的.
所以盼望.
如果我們去到這樣的年紀.

$^{1121}$都不爭去事奉.
所謂爭不是用我們的血氣.
就是說沒有人做我做.
沒有人做我去做.
我很羨慕服侍神.
我很羨慕去幫弟兄姊妹.
這樣的羨慕.
今天有多少弟兄姊妹.
有時候我都.
不要緊.
我們都被人說我們是一個退休的樂園.
中產的退休的樂園.
如果退休在你的生命裡面.
這個只是一個樂園.
而我們信了主.
不好好地用我們的人生去服侍主.
我們就損失了我們自己生命的成長.
很多崗位不夠人.
我們是不是闊佬難理.
看不到就不知道.
因為我們很想去參與那些事奉.
配合牧者不同的安排.
我們去做呢?.
所以其實講的是我們那個心態怎樣轉變的.
是不是我們今天我們的教會.
還沒有牧養教導弟兄姊妹.
深化到一個地步.
我們渴望成為一個淑女人被主去用呢?.
所以若是講一句話.
就已經給他內心的那種感染力.
令到那些兄弟立刻無地自容.
但是是一份愛的重新的接納.
因為若是不是去碰他們.
不是現在來跟他們翻舊帳.
若是告訴他們知道.
這個是神在我生命裡面的工作.
就是說幾十年過去了.
若是已經跟他的兄弟.
原來距離很大了.
若是已經成為淑女人了.

$^{1161}$所以弟兄姊妹.
你跟你一起順主使禮的弟兄姊妹.
如果我們是有追求.
至少我們來查經班已經是追求了.
一個半小時,一個星期.
你已經損失了一次的劇集了.
這個時間是做了什麼劇集?我不知道.
你想一下你損失了一個一晚的聚會.
但是如果我們有追求.
其實你不要說笑.
時間是一個考驗.
過了一段長時間之後.
如果你默默地在事奉裡面.
用心去學習,用心去做.
跟你一起使禮的弟兄姊妹.
我不是要大家比較.
而是有一個不能回頭的落差.
一個人就越來越熟了.
一個人就慢慢越來越熟悉了.
慢慢越來越回到眷戀的舊友的埃及了.
所以若是這一番說話.
令到他的兄弟真的是非常之欽佩.
非常之敬服.
所以屬靈人帶來的是一個黏合的合一.
不是一個分裂.
屬靈人如果個個都不敢靠近你.
那就是你的問題了.
因為屬靈的人應該就是令到其他人.
更加有一份平安.
更加有一份恩典和祝福的.
接著就是二十二節.
若失和他父親的眷屬都住在埃及.
若失沒了一百一十歲.
若失得見爾法年第三代的子孫.
馬來西的孫 馬吉的兒子.
也仰在若失的膝上.
若失對他弟兄們說.
我要死了 但神必定看顧你們.
領你們從這地上去.
這個就是約書馬記的經文.

$^{1201}$道他起誓所應許.
給亞伯拉罕以撒那國之地.
就是迦南的應許之地.
第二十五節.
若失叫以色列的子孫起誓.
說神必定看顧你們.
你們要把我的骸骨從這裡搬上去.
所以不是那條詩.
根本那條詩無論你怎樣用那些埃及最好的防護.
它都會變骨的 沒有的了.
多大的防護都沒有用.
當然可能這個埃及的木乃伊可能會好一些.
例如留蘭的那些古詩.
中國裡面真是有一個完整的詩體.
但是無論如何.
若失要的不是那個詩體.
要的是他的骸骨就已經夠收貨了.
你拿我的骸骨就行了.
就是不會強人所難.
因為你怎樣防護呢?.
防多少年呢?沒有需要.
所以我初信的時候都很害怕.
火花行不行呢?.
後來越來越對聖經認識.
根本不在乎你這個人的保存得好和不好.
因為將來就會有一個榮耀的身體.
是屬於我 也是屬於你自己的.
所以是不是火化,土葬或者是殺落海.
都完全沒有問題.
重點真是我們要保留的是那個對永恆的應許.
古代的人若失都已經很清楚.
對永恆的主的信心,信靠和盼望.
原來人家的骸骨都可以說話.
人家的骸骨都可以有一種屬靈的影響力.
原來表面就是一個獨臣.
但實際上我和我將來的子子孫孫.
如果都是有信心的亞伯拉罕的子孫.
都是對永恆的這個不如比利爾迦納地.
很實在的物質的地圖.
因為神的計劃不是一跳就跳到新約.

$^{1241}$都是一個歷史的進程.
這樣去讓他們得到一個真是應許之地.
但是真正踏進應許之地的約書馬紀.
距離有40年有多.
所以這個都是告訴我們神的計劃和時間.
但是告訴我們有一個人只要他有信心.
就可以孕育信心的後代.
亞伯拉罕然後就將那個應許.
等於一個家族的寶貴的傳承.
就是這一份信心對神的信心.
轉嫁到以撒.
以撒就去到雅各.
雅各都離開了.
就由約瑟幫他送終.
約瑟都死了.
所以很完整的創世記.
真的告訴我們神對遠古的先祖的應許.
透過約瑟的棺材.
約瑟的骸骨.
來作為這一本書的應許的標誌.
一個圖騰.
以色列人什麼都沒有.
但是有骸骨.
但是很多時候已經忘記了.
但是感恩就是他們都不敢將這個骸骨.
來輕輕的處置.
所以一個教會.
到底我們對於不同的前輩留給我們的.
到底是一個只是硬梆梆的硬件的教堂.
建築物.
還是一個很重要的信心的遺產.
我們能夠孕育有一代,二代.
甘孫我想已經是第三,第四代了.
我們的信心是怎樣呢?.
我們是不是真的和前人那個應許的信心.
來到一路銜接呢?.
所以一定不是多人有這樣的信心的.
慢慢他們就會不知不覺就會在那個.
所謂現實的生活這麼殘酷,這麼現實.
就慢慢忘記了這些應許.

$^{1281}$所以我們要怎樣讓這些應許.
一路一路來到讓我們的下一代.
都能夠繼續這樣傳承呢?.
這些是當政者必須要做的功課.
必須要讓頂子妹很透徹地得到這個模樣.
所以很奇妙地.
創世記以一個棺木來收這本書的結束.
很噁心的結束.
但是原來是這麼重要的一脈相承的應許.
所以我們要很快地揭一揭.
創世記的第十二章.
第十二章第一節.
耶和華對阿伯藍說.
你要離開本地本族父家.
和我所要指示你的地去.
就是將來以色列人要進的迦勒地.
我必叫你成為大國.
我必賜福給你.
叫你的名為大.
你也要叫別人得福.
為你祝福的我必賜福於他.
就助你的我必就助他.
地上的萬族都要因你得福.
阿伯藍就照著神的顯現.
神的應許就開展他的人生.
就是離開當時的宗教文化的發源地.
這個吾爾就一路輾轉.
就是75歲的時候就出哈蘭.
和他的太太撒萊和他的侄羅德.
和他的爸爸塔拉.
一起離開準備進入迦勒地.
所以第五節就說.
他們已經將一切財物,人口.
都帶往迦勒地去.
他們就到了迦勒地.
聖經就是很簡單的.
但是這個過程.
你想想沿著右法拉底河.
一路沿著河這樣走.
右法拉底河現在最近說越來越乾固.

$^{1321}$那個水位越來越低.
可能到2030年會完全乾了.
就似乎好像《應驗聖經》說的寓言.
將來大河都會乾坡.
人們就會在那裡經過.
現在應該是伊拉克一帶的地方.
但是無論如何.
這個就是神對阿伯藍在吾爾的時候.
向他顯現的那個應許.
所以這個祝福一路到今天.
我們承受了.
因為我們是藉著基督.
我們都是亞伯拉罕的子孫.
我們都得到這個永恆的福祿.
我們都能夠有這個這樣的萬俗的福.
所以如果看不到我們屬靈的身份.
我們是神的兒女.
我們是相信主耶穌.
相信將來的天堂.
我們將這個福音告訴別人.
如果那些人還是不明白.
攻擊就坐.
他們就會無神的就坐.
當然就是他們成為了將來在永恆裡面的.
那個痛苦的火壺.
用了這個時間就做了這個引言.
講出埃及記.
還有一點都講了.
所以死亡沒有抵消到神的這個應許.
大家明白了.
我們要揭第十五章.
這個就是神在迦蘭的時候.
就向亞伯蘭有另外一個.
再一個 double confirm.
再給他兼顧這個約.
更加具體的約.
不是籠統.
第十五章標題都說神和亞伯蘭立約.
這些告訴我們.
這是以後耶和華在異象中有話對亞伯蘭說.

$^{1361}$你不要懼怕.
因為當時有五王的戰爭.
四王和五王的戰爭.
有羅德和亞伯蘭的產業.
被那些人擄掠了.
所以有這個戰爭在.
出現了一個很大的懼怕.
這是第十一節.
四王就把索多瑪和摩拉所有的財物.
並一切的糧食都擄掠去了.
又把亞伯蘭的侄兒羅德.
因為羅德就住在索多瑪.
就和羅德的財物擄掠去了.
當時羅德正住在索多瑪.
所以有這個情況.
後來當然亞伯蘭和他三百多個的家丁.
就去打這場仗.
最終打贏了.
三百一十八個的壯丁.
所以第十五章就是.
當時亞伯蘭應該都思前想後有一定的懼怕.
神跟他說你不需要懼怕.
我是你的盾牌.
必定大大的賞賜你.
亞伯蘭當然就是說.
他連羅德都不再和他一起.
羅德本來可以繼承都不行.
不如就是我的僕人.
因為我現在都沒得生.
我的太太薩萊都沒得生.
所以就有這個狀況.
這是一個很現實.
亞伯蘭都不是不相信.
但實際上他都沒有後裔.
何來有將來的後裔可以承繼產業呢?.
然後長話短說.
第十二節.
日頭正落的時候.
亞伯蘭沉沉地睡了.
忽然有驚人的大黑暗落在他身上.

$^{1401}$耶和華對亞伯蘭說.
你要的確知道.
其實沉沉地睡了.
是報夢還是嚇醒他呢?.
你要的確知道.
你的後裔必寄居別人的地.
相當嚴肅的氣氛.
和亞伯蘭這樣說.
又服侍那地的人.
這就是出埃及的時候.
他們落在這樣的光景.
但是現在亞伯蘭在迦林地.
其實只是85歲.
還有一百多年他才死的.
然後還有四百年他們才出埃及.
現在這樣說的時候.
就是說很長的時間都還沒兌現.
但是你要留意.
又服侍那地的人.
那地的人當然你都不知道在哪裡.
但其實就是埃及.
那地的人要輔待他們四百年.
並且他們所要服侍的那國.
我要懲罰.
後來他們必帶著許多財物.
就是第十二章說的.
以色列人就擄掠了埃及人許多財物.
就拿走.
因為要用許多的材料.
將來在路上要建造他們的國度.
他們的回望等等.
從哪裡出來.
第十五節.
但是要在大壽宿.
平平安安的歸到你列祖那裡.
被人埋葬.
葬在迦林地.
到了第四代.
他們必回到此地.
到了第四代.

$^{1441}$亞伯拉罕數了第一代.
以撒那國約瑟.
可能如果不計算.
亞伯拉罕就是摩西.
第五代就是第四代.
但是要在大壽宿.
接著第四代.
他們必回到此地.
不過留意這裡.
因為阿摩利人的罪孽.
還會有滿盈.
就是說原來那個時機未到.
就不是說神不立刻按了按鈕.
而是按了按鈕.
都未必會有一個好的結果.
包括百姓神的選民.
希伯來人未必會明白珍惜.
未必會感恩.
還有迦林地的原居民.
原來在神的眼中.
都還給他們有一個悔改的時間,機會.
阿摩利人的罪孽還未滿盈.
就是說神容讓他們去按他們的自由意志.
去信主,去犯罪,去拜偶像.
去抵擋真神,做物主.
但是希伯來人仍然是一群很渺小的遊牧民族.
但是與別不同的.
他們是有信心的希伯來人.
因為有亞伯拉罕在.
因為有亞伯拉罕的家族的傳承.
口述的歷史.
來告訴後代.
我們是希伯來人的神的子民.
但是迦林的原居民.
無論是阿摩利人.
那就是說似乎是統稱了.
就是說迦林地的那些七族的人.
他們的罪惡未滿盈.
所以神的時間就未到.
所以求主給我們都能夠看到.

$^{1481}$這裡有神的主權.
救贖的計劃的一種很美麗的計劃.
真的要等到第四代才會施行的拯救.
不是神不能夠拯救.
所以我們就要去看到.
這個就是那個應許的很重要的兌現.
接著告訴我們第十七節.
日落天黑.
不料有無煙的籠.
並燒著的火把.
很特別的 其實是什麼來的呢?.
是不是當時的迦南人.
或者當時的這些民族.
行一個納藥的記號呢?.
無煙的籠 就是有一個火籠.
是無煙的 有燒著的火把.
從那些玉塊中經過.
玉塊就是神叫阿巴拿.
之前上人說的.
就要把對碟放好.
就是一隻羊 tong 了它.
一tong 一分為二左右.
但是要疊整齊.
似乎是一個納藥的儀式.
似乎是那個納藥者要走的那條血路.
所以表面上好像很恐怖.
但其實是當時代的一個納藥的記號.
神都用這個記號.
當日耶和華與阿巴拿納藥.
我以切割你的後裔.
從埃及河直到伯拉大河之地.
就是基尼人 基尼洗人 格摩尼人.
黑人 比利洗人 尼弗欽人.
阿摩尼人 迦南人 格加薩人.
耶布斯人之地.
所以其實在一個昂爬虎穴裡面.
阿摩尼人的罪未滿刑.
也就是說 耶布斯人的罪都未滿刑.
也就是說 迦南地的人.
其實可能只是一個統稱的代表.

$^{1521}$也就是說 整個迦南地的人.
一直拜斗像 積蓄神的憤怒.
都還未到滿刑.
未到一個徹底的拯救的開始.
但是神已經和伯蘭有這個眼光.
當時他只有85歲.
就有這個眼光和你立了一個這樣的永約.
你的子孫將來就要寄居在這個地方.
有多大呢?.
從埃及河直到伯拉大河.
就是由伯拉大河之地.
你想一下 從埃及河到伯拉大河.
非常大的地方.
你可以想像到整個的希伯來人.
迦南地就是屬於他們的.
就是這七族的人居住的地方.
這個就是神和伯蘭的立約.
所以這個是很重要的.
出埃及的那個開始.
這些我們說了.
所以你會看到 出埃及我們看的經文.
你就會看到這個神學.
和一個修辭的敘事文體的一個連繫.
寄居和輔代就對應出埃及的第一章的時候.
因為神和伯蘭說.
你的子孫要寄居那個地方.
開始的時候相安無事.
和平共處 而且有很高的社會地位.
好像有一個最惠國待遇.
給你們這些希伯來人那樣.
因為你們因為約瑟的緣故.
你們蒙恩 但是後來好景不常.
突然間一個新王出現.
他也是埃及王.
但是可能就是其他的外族.
打贏了舊有的埃及 征服了.
又是叫做埃及 又是叫做法老.
但是就不認識那些歷史了.
或者根本不介意否定那些舊有的歷史.
就是說希伯來人如果只是寄望於法老.

$^{1561}$就慘了 因為原來好景可以不常的.
原來很多東西可以突然間改變的.
但是如果我們的信心是放在那副骸骨.
是放在伯拉罕的那個應許.
以色列就不需要懼怕了.
所以總有些屬靈人是等待那個應許的.
所以你看到寄居和輔待.
第一章就已經出現了.
因為神說了.
那地的人會輔待你們的子孫400年.
那麼神就會懲罰那個國家.
這個就是15章那裡就對應出埃及的神童.
就是透過摩西神童法老的那個對質.
就是十災 然後就會帶著財物出來.
所以對應的就是他們離開埃及滅長子.
就是反而埃及王請求以色列人盡快走.
和給那些財物他們.
所以是應驗了《創世記》的15章的經文.
所以我們由下個星期開始.
我們一起去看第一章的《出埃及記》.
這些我們剛才講了.
大概的情況就是這樣.
所以信心的種子就是亞伯蘭相信耶和華.
15章的第五節亞伯蘭信耶和華.
耶和華就以此為他的意.
所以以色列人出埃及.
其實就不是以色列人的建國史.
是神的透過亞伯蘭.
對亞伯蘭的應許的一個神的救贖計劃的一個成就.
一個的歷史.
所以舊約就稱為救贖的歷史.
所以不是一個偶然的.
是神已經應許了.
在耿古就已經立定了的旨意.
所以這個是一個定旨.
預言的一個必然的發生.
今天我們就講到這裡.
和大家一起做一個祈禱.
幾乎我們感恩 因為我們這樣揭舊約的時候.
12章的創世記.

$^{1601}$15章的神和亞伯蘭的立約.
而且這個是一個永久.
因為是一個用這樣的儀式.
就是有無煙的爐和燒著的火把.
在那些玉塊當中經過.
以當時代的一個立約的雙方的儀式.
就是說這個是永不更改的.
更加我們今天相信.
在主耶穌的寶血救贖永恆的約.
將來我們就在永恆裡面.
會來見到神 得到永恆的產業.
所以求神幫助我們.
當我們看似乎是舊約的一些歷史記載.
有些是特別建造回望的那麼巨細無遺的.
那些針線 那些的尺寸等等.
讓我們都能夠可以明白裡面背後的含義.
更加我們看到這裡原來講的是應許.
講的是一個民族得到應許的寶貴.
他們慢慢從奴隸成為神愛的.
或者無神愛的子民.
所以神幫助我們感宣家.
幫助我們香港的教會.
幫助我們在這個時代的教會.
無論世代如何變遷.
讓我們看到我們是從舊屬進入永恆的安息.
進入將來永恆的天間.
所以今生真的讓我們有衣有食就知足.
我們的心是放在永恆 放在舊恩裡面.
能夠與神同工才是最寶貴,最重要.
在不同的地方有不同的弟兄姊妹.
有時有不同的艱難.
所以盼望這一段的聖經.
今晚的應許,安慰,立約.
成為我們大家心靈的盼望.
如果有弟兄姊妹是落在屬靈的幽谷.
好像沒有盼望的.
讓我們看到摩西,看到這個骸骨.
看到亞特蘭的應許.
就讓我們有這個心靈的燃點的盼望.
我們這樣禱告奉耶穌基督的名.

$^{1641}$阿們.
\newpage



\section{出埃及記}
\label{sec:kZF3v6CUrpM}
\textbf{2026-05-22 周五晚查經班《出埃及記研讀(二)應許在黑暗中增長》 講員:曾裔貴牧師}
\newline
\newline
連結: \href{https://youtube.com/watch?v=kZF3v6CUrpM}{\texttt{https://youtube.com/watch?v=kZF3v6CUrpM}} ~~~~ 語音日期: 2026-05-23
\newline
\newline
\hyperref[sec:GROoST1kr4I]{\small{< < < PREV SERMON < < <}}
~
\hyperref[sec:index_chronic]{\small{[返順時目]}}
~
\hyperref[sec:index_scriptual]{\small{[返順卷目]}}
~
\hyperref[sec:4cUSeyCs18Q]{\small{> > > NEXT SERMON > > >}}
\newline
\newline
$^{1}$好 弟子們晚安.
我們開始第二堂《出埃及記》的查經.
先跟大家同心圖告.
多謝天父 我們又可以在晚上放下我們的工作.
也有一個短短的假期.
我們可以多在屬靈的方面去追求.
去親近主 去認識你.
還有當我們去研讀舊約經卷的時候.
讓我們能夠看到的 就是神理的救贖的計劃.
是怎樣在遠古的時候來去成就.
我們上一堂都已經看到 在《創世記》15章.
神理不單止是對亞伯蘭有一個祝福.
允許祂能夠生一個兒子.
但是更加讓我們看到的就是.
將來亞伯拉罕的子孫就是要承受神的國度.
而且這個過程是鉅細無遺地在《出埃及記》.
讓我們看到這個歷史的發展.
求主幫助我們 因為我們在新約.
在我們每天的生活.
我們未必能夠有一幅很清楚的圖畫.
我們看到前面的路.
但是當我們去回望以色列人他們過紅海.
他們在西乃山成為神的國度的時候.
還有他們準備進入迦南的過程.
好像我們在新約都有看到.
在哥林多前書我們有引述的.
就是很多的百姓都是在整個過程是沒有信心.
是神所不起悅 因為他們就是這樣.
只是憑著那個眼見 不是憑信心.
這樣來得到應許之地.
所以我們看到這個都是一個很重要的尷尬.
絕對不是舊約.
就是我們根本沒有需要再深入的研究.
求神幫助我們能夠看到背後一些的事情.
其實和我們今天成為基督徒息息相關的時候.
我們就更加要謹慎 願主帶領我們.
今晚我們開始第二課.
盼望神都給我們有一個很宏觀的看到.
那個發展的那個敘事.
以致我們看到作者那個寫作的一個企劃.

$^{41}$我們能夠在其中找到一些很重要的端倪.
成為我們自己的生命的學習和力量.
我們仰望禱告奉耶穌基督的名 阿們.
好 我們第二課很快和大家溫習.
相信大家都應該有印象.
《創世記》15章的那個應許.
和《創世記》的第50章.
表面上就是虐失這個這麼重要的.
當時真是因為他而帶來整個的希伯來的民族.
因為他而蒙福.
但是最終他還是要離開世界110歲.
但是告訴我們他有一個要求以色列人.
他們的子孫要有一個誓約.
就是不可以幫他埋葬在埃及.
意味著他都要出埃及.
不過是用他的棺木 可能不方便的.
就是用骨灰罌來裝著.
這樣來進入這個迦南.
這裡也讓我們看到真正的應許.
憑信心連死了都可以說.
這個希伯來書挺重的.
對於大概有差不多41個信心偉人的一個論述.
所以虐失的棺木和神對亞伯拉罕的應許.
400年之後 他們的子孫.
就是虐失的子孫就會進入應許之地.
這個應許是以色列人一直都承載著.
一直他們都在等待出埃及的.
這就是上一次我們說了.
迦南人的罪還沒有滿刑.
所以那個時間一直要等待 一直要拖.
相信頭那段時間都是非常的太平.
非常的安居樂業 非常的慢慢的生兒養女.
就已經在那裡 在埃及慢慢的融化.
融入了那個地方.
也都沒有這麼強烈的應許.
那個要呼求神快點兌現祖宗亞伯拉罕的那個應許.
但是我們看到神自己去工作.
就給我們看到 這裡首先有一個串連.
帶我們回到埃及的那幫人.
就是現在在埃及的新興的王之前.

$^{81}$就是這70人來到孕育.
相信那個時候都已經有二三百萬人了.
有400年的時間 都一點都不出奇.
因為據聞當時只是出埃及的男丁都60萬.
所以大概有婦女 有小朋友有200萬人.
都不是一件稀奇的事.
我們很快看看 出埃及記的第一章.
一到第五節.
以色列的眾子 各大家眷 和雅各一同來到埃及.
他們的名字記在下面.
有劉汴 西冕 利美 尤戴.
這些是利亞所生的 最大的那個太太.
以撒加 西布倫 卞瓦敏 丹 拿忽哈利 迦德 瓦切.
凡從雅各而生的共有70人.
約瑟已經在埃及 所以沒有記在約瑟的名字.
約瑟和他的弟兄 並那一代的人都死了.
因為已經是在這個時候.
是在說出埃及的時間的前後的光景.
所以用一個70人的家譜.
來串連創世記400年前的應許.
當然 亞伯蘭的時候就不是400年.
因為已經到了第四代 就是約瑟和他的子孫.
所以這一番話 其實應該第50章.
就有講到他們落埃及的時候 那個人數的.
很快和大家來到去稍微揭一揭.
應該就是46章.
先看看記載的家譜.
這裡就有一個很詳盡的家譜.
當然我們不需要讀完這麼多.
第46章的第一節 46章 創世記.
就是將兩段經文串連了.
就是說400年神都沒有忘記.
這一幫這麼小的民族.
46章就是以色列帶著一切所有的.
以色列就是那國改了名字之後的名字.
起身來到別氏巴.
別氏巴就是他們竹壇的地方.
就獻祭給他父親以撒的神.
夜間神在異象中對以色列說.
雅各 雅各 他說我在這裡.

$^{121}$這個都很特別的 神在異象中呼叫以色列.
就叫回他舊未改變的時候的名字.
不是語名這麼簡單.
是他未真正認識神的時候的那個雅各.
就是他那個好像未重生之前的那個本性的舊人.
神紀念他這個人 神現在看到他是.
以色列這個新的名字.
以色列就是神的王子這樣的意思.
我在這裡第三節 神說我是神.
就是你父親的神 你下埃及去不要害怕.
因為我必使你在那裡成為大族.
我要和你同下埃及去 也必定帶你上來.
約瑟必給你誦宗 雅各就從別氏巴起行.
就是臨落埃及的原因就是.
因為約瑟就派那些隨從.
相當大的一個歡迎的團隊.
就是有了他們要用的那些車.
來載爸爸 載兄弟.
載家中的一些財物 衣服等等.
其實其他的都不用拿了 因為埃及什麼都有的.
所以這次這個歡迎團隊很大很豪大.
但是雅各要下去之前.
和一個已經20年失散了的兒子來相認.
雅各做了一件很轟動的事情.
就是帶著所有的家眷去求問神.
去別氏巴來獻祭.
不單只是一個感恩 而是求問.
我應不應該下去呢?用什麼形態下去呢?.
你可以想象得到一個老人家.
其實他都已經是瘸腿的了.
一個兒子 最愛的兒子 失散了20年.
已經是差不多等於死了.
但是現在約瑟派他的隨從.
和他的兄弟回來.
說約瑟現在做了埃及的宰相.
這樣對於雅各來講.
是一件可以即時心臟病入ICU的.
這樣的一個消息來的.
但是他可以很平靜.
他要來到別氏巴求問神.

$^{161}$在這個重大的決定上面.
一個都挺重要的分水嶺.
到底神的心意是怎麼樣呢?.
這個是已經蒙恩 生命得到很大的改變.
和神的關係是非常密切的.
雅各他有這樣的一個看似很簡單的動作.
實質上是非常之敬虔的.
我不知道我們一些舉家移民.
突然之間要所有的家人.
連同爸爸媽媽都要一起移民的.
有沒有去禮拜堂來感受禮堂.
安靜一會兒 祈禱.
然後見到勵善牧師.
或者陳牧師 謝牧師來到去講.
我們要離開了 你為我們祈禱.
看看神怎麼帶領我們整個家族.
去到加拿大或者英國.
我暫時就好像沒有見到 沒有聽過.
原來很敬虔的 原來臨走之前.
回禮拜堂 回別氏巴去求問神.
神向雅各顯然 叫他那個還沒有改變之前.
在比魯伊勒神幫他改名的32章那裡.
這樣來去稱呼他.
就是說 雅各你今天能夠這樣的.
屬靈生命的成長 都是我的改變的能力.
都是我的聖靈在你的內心.
不斷做這些聖靈的手術.
你能夠有一個很自我中心.
好什麼都抓住 什麼都為自己著想的雅各.
變成現在你擴大到.
其實都不用問阿貴 你的兒子失散了20年.
已經死了 現在在埃及做宰相.
那些歡迎的隊伍 那些隨從.
埃及那些僕人已經來到你家門口了.
基本上不用想了 所以不是去還是不去的問題.
而是整件事他要在神的那個裡面再去消化.
其實這個真的很敬虔 真的挺重要的.
所以我們的心底裡面有沒有禮拜堂呢?.
有沒有屬靈的長者給我們大家一起.
不一定要求問的 可能是分享.

$^{201}$可能是讓一些屬靈的前輩知道你的事情.
為你感恩 為你祈禱 為你新的去向.
都有一個祝福 原來是很屬靈的 真的很重要的.
雅各本來什麼都不用想的了 他很有計謀的.
想事情很快的 轉數很快的.
他哥哥都不夠他鬥的 這樣的一個雅各.
他現在在這個時候帶著一切所有.
起身去到別墅 跟著聖經就真的很鉅細無遺了.
既然有了神所謂開綠燈 這個開綠燈更加不是說.
好像根本上就是他的心已經想去的了.
不過神你告訴我 我就更加有印證了.
其實就不是 他真的將他的內心帶到神那裡.
他現在要去埃及 神給他的應許就是.
不僅你可以去埃及 我跟你一起去埃及.
神可以無處不在 但是對一個現在要舉家.
整個家族要離開這個迦南應許之地.
大家不要忘記 雅各已經在應許之地.
是老祖宗留給我們的 是神應許我們的.
現在是因為有天災,有饑荒 所以我才要下去.
所以我這樣下去 和我的爸爸以撒.
和我的爺爺亞伯蘭 他們都去過埃及.
不過是犯罪 是軟弱的心靈.
現在雅各去埃及就不是軟弱 是有神的默認.
還有那個計劃就是你們整個民族.
70人就要寄居到下面 然後要等400年才可以走.
而且是以一個奴隸脫離那個困所.
這樣離開 昂然的進入迦南 就是回到迦南.
所以其實下去要回來是很容易的.
那些饑荒解決了之後.
都可以由蘭塞那個很好的地方舉家回來迦南.
但是他們清楚一件事就是 神的時間不容許.
他們就一直在埃及生兒養女.
所以這個都是很戲劇性的.
就是怎樣可以知道不能夠再回來這個迦南地呢?.
迦南明明是神應許的 所以下去之後.
他們就差不多永久了 400年就安居在埃及了.
簡單讀一讀那些人名.
第五節 雅各就從別氏巴起行.
以色列的兒子們使他們的父親雅各.
和他們的妻子兒女都坐在法老為雅各送來的車上.

$^{241}$沒錯了 他們又帶著在迦南地所得的生充貨財.
來到埃及 雅各和他的一切子孫都一同來了.
雅各把他的兒子,孫子,女兒,孫女.
並他的子子孫孫一同帶到埃及.
一個都不剩 決心按照神的指引.
舉家離開 好像不眷戀迦南地.
其實又是很難 已經說明.
迦南就是你們最終的那個應許之地了.
已經是第三代亞伯拉罕爾薩雅各.
第三代已經在這裡住了.
包括第四代的那12個兒子.
這樣迦南就是最終的應該是終點了.
但原來不是 還有一個旅途就是要到埃及.
雅各很聽話 全家70人.
老幼所有的家眷全部一起到埃及.
我們有時候很少留意這些經文.
你看到一個真是信心的人.
他的前半生 先當一半.
他的前半生是不斷為自己去抓住神的祝福.
但現在他是完全開放自己的心靈.
讓神的祝福一直伴隨他.
你想想多麼美麗的人生 去到晚年的時候.
他去埃及的時候應該就是130歲了.
他147歲死的 還有17年.
可以過一個埃及的太平盛世的生活.
雅各130歲 他的屬靈生命.
去到一種不再為自己求 充滿的就是基督.
就是那首聖詩 孫迅博士那首聖詩.
Himself 前要的是祝福 今要主自己.
那首的詩歌 雅各就是這樣了.
神要我現在下去 不是去visit一下.
好吧 去加拿大 因為我姐姐的女兒快要結婚.
很想我媽媽搭飛機去多倫多.
她說不行啊 我90歲去不了了 這麼累.
頭等都頂不住了 15個小時.
你想一下 如果她真的決心.
不是啊 孫女結婚 當然她會回來.
但是如果過去走一走也好.
這些都是一個visit 不會長期的居留.
但是雅各清楚神的旨意.

$^{281}$就是現在你放下加拿大 將來都會是你們的.
你想一下 用錢買的那些墳地和那些地圖.
現在要拱手讓給加拿大人.
還有舉家下去 70人下去.
都不知道法老和他們的埃及人會不會對他們不好呢.
這些全部未知數 所以有時候舊約.
其實有很多空間給我們基督徒去想的.
給我們去代日的那個光景裡面.
你看到屬靈人他們的行事真的很不簡單.
真是去求問神 所以你可以想象到.
他現在以一個很屬靈的長者的一個位份.
帶著全家下去埃及 不需要懼怕.
因為是你 永遠是你的.
應許之地已經給了希伯來人 就是希伯來人了.
不需要爭 爭都沒有用 因為神的時間還沒有到.
這些的豁達 這些的信心 這些的屬靈的眼光.
今天我們有多少人 隨著我們信主來.
我們屬靈的眼光能夠一直這樣成長.
有一種這樣的屬靈的高瞻遠矚.
如果沒有 其實就是一個教會的悲哀.
大家都是這樣聚會.
如果有些教會是很多敬虔的人出現.
而有些教會是沒有.
那只是因為我們根本從來就沒有遇見過主.
沒有認真的追求神的國度.
我們只是掛著一個名字.
開開心心做一個禮拜日基督徒.
我們以為就已經很足夠了 但其實不是的.
真正的追求敬虔 其實是有很多磨難的.
至少有一個交戰 嘩 迦南是我的.
我現在下去打一張卡 上來就OK了.
你有沒有可能叫我全家舉家下去埃及呢?.
沒有可能的 但是這一種屬靈的信心,盼望.
和謙卑的讓神的旨意去成就.
就真的不是那麼容易.
有時候好像解不明白那樣.
但是我們讀完這些經文.
其實我們就應該要有我們自己的代入.
屬靈的生命成長 我們都應該有這樣的眼光.
這個名字就記在下面了 來到埃及第八節.

$^{321}$來到埃及的以色列人名字記在下面.
雅各和他的兒孫 雅各的長子劉弁.
然後這次是一個很詳細的家譜.
所以連劉弁有多少的兒子.
利美或者等等都鉅細無遺地記載.
詳細我就不講那麼多了.
接著就是去到第26節.
乃與雅各同到埃及的.
除了他兒婦之外 凡從他所生的共有66人.
還有約瑟在埃及所生的兩個兒子.
雅各家來到埃及的共有70人.
當然也算約瑟在內 總共4個.
就變成70個了 你可以想像得到.
第一章出埃及的 剛才我們讀的.
是一個很簡單的家譜.
就是這12兄弟孕育出的.
有些是在迦南已經出生的.
有些是去到埃及之後再出生的.
就已經將整件事.
作者很巧妙地用文學的修辭.
就來將這個400年之後的事情.
連繫起來 好像這個家譜.
一個很簡單的家譜12個人.
就來連繫了整個的故事.
就是說我們有很多的思想.
雅各 當然還有一些記載就是.
一下去之後 雅各就以整個家族的代表.
去見這個法老王 告訴他們.
告訴法老知道那個來龍去脈神.
不如我讀給大家聽 可能大家又是不算很熟悉的.
47章.
第七節 47章的第七節.
若失領他父親雅各進到法老面前.
雅各就給法老祝福 法老問雅各說.
你平生的年日是多少呢?雅各對法老說.
我寄居在世的年日是130歲.
我平生的年日又少又苦.
不及我列祖在世寄居的年日.
因為他有相當長的年輕的時候.
跟他哥哥有很多的角力.

$^{361}$跟神的那個應許又有很多的角力.
或者他自己的人就是滿腦子都是那些計謀.
就是一個心計很重的一個雅各.
所以他信了主 用我們新約說.
他認識了神 認識了自己之後.
他的生命有改變之後.
他發現原來一切都是在神手中.
原來一切都是有安排.
原來一切都是恩典的時候.
他看他過去那種競爭.
真是又短又苦.
所以我們盼望.
我們能夠早日有我們自己的發現.
不再是我乃是基督.
我們這樣能夠人生就過得輕散.
很多事情絕對不是躺平 兩樣東西.
躺平是不思進取是不對的.
我們應該是很有那個渴慕神的心.
甚至年紀越來越大.
渴慕天家 渴慕見主.
這樣的心應該與日俱增.
但是仍然我們有一口氣.
都為福音義不容辭.
是用不同的方式.
年輕的用多一些力氣.
晚年了用多一些的心智.
和我們的禱告.
開始為不同的後輩去祈禱.
譬如有意識地我這個禮拜.
是為兒童牧區去祈禱.
下個禮拜為青年的牧者.
青年的士工去祈禱.
教會需要這樣的智慧人.
教會需要這樣的屬靈人.
不需要你要走到哪裡.
而是你的家裡的那個房間.
你的桌子就已經是一個聖殿的祭壇.
很重要的就是我們的心靈.
就是去到神的面前去禱告.
甚至神給你有感動.

$^{401}$哪一個是可造之材的.
你為他祈禱.
你為他來交託給神.
要神去用他 去磨練他等等.
這些都是很重要的.
屬靈的長輩要做的事情.
屬靈的長輩不一定是很大年紀的.
很大年紀應該就是屬靈的長輩.
兩樣東西在教會不應該是沒有的.
沒有就是大件事了.
沒有就是一個悲哀.
我們應該有這樣的眼光去看不同的事物.
我們看到年輕人不應該不斷說.
一孩不如一孩 其實不應該的.
應該就是為我們的年輕人禱告.
將他們帶到神的面前.
讓他們能夠經歷神使他們的生命成長.
所以雅各就是這樣來對自己的認識.
雅各又給法老祝福 就從法老面前出去了.
約瑟遵著法老的命.
把埃及國最好的地 就是南塞境內的地.
給他父親和弟兄居住 作為產業.
約瑟用糧食奉養他父親和他弟兄.
並他父親全家的眷屬.
到時照國家的人口奉養他們.
一個養活全個家族.
因為他已經貴為全國的宰相.
神給他有極大的智慧 運籌為額.
就是神的時間到了 神就大用約瑟.
但是經歷過被兄弟的背叛.
被那些人賣豬仔去了埃及.
經歷了波提佛家裡有試探.
要被成為階下囚.
接著在監獄裡面幫兩位埃及法老的王.
兩個重要的大臣.
一個酒精 一個是管理飯食的膳長.
也是這樣得不到報答.
最終在監獄裡面一直都是等候.
好像無了期的絕望.
這個光景有20年的時間.

$^{441}$就可以想象到20年對約瑟來說.
原來是一個黑暗的一個試煉.
神是萬有的主.
祂是不會試探人 也不會被惡試探.
祂是萬福的恩元.
所以每一個忠心在神的國度裡面的弟兄姊妹.
最終都會 當然我們不應該叫出人頭地.
儘管用約瑟這樣說.
所以屬靈的眼光.
當然有些人可以很成功被神大用.
但是有些人可能繼續一生都在一些疾病.
在一些困苦裡面這樣來渡過.
但是不要緊的.
神會讓那些事情得到最終的超說.
所以約瑟原來是一個無辜的人.
約瑟原來是可以有這麼大的智慧.
來管理整個國家的人民的安危.
因為有七年的封面.
之後有七年的大的饑荒.
所以約瑟就做了這樣的安排.
運籌為額的這樣來管理.
使得整個國度都來避免這一次天災的毀滅.
所以約瑟就是這樣.
這個當然是另外一個故事.
約瑟就是這樣被神來揀選.
所以我想也是的.
被神揀選使用的僕人.
一定是不可能沒有一些困難.
其實作為一個公眾人物.
任何一個公眾人物無論那個幅度有多大.
你只是一個普通的一個小小的公眾人物.
你坐在一個櫃檯裡面.
那些人來問問題都會被人罵.
更何況你管理一家公司.
一個教會或者一個國度.
所以一定是有很多的困難的.
所以困難就不要緊的.
原來可以面對困難的那個人才更加重要.
所以神是先訓練人有能力面對困難.
這樣的成長就是千金不換了.

$^{481}$所以基督徒不要怕困難.
如果神在我們生命裡面是有一些的考驗.
我們就要學習去承擔.
所以當我們有擔當的時候.
神又會將更大的祝福,崗位來給我們.
所以這裡就讓我們看到.
雅各就是這樣來到去埃及.
後面那段都很寶貴.
而且是很值得我們去讀一讀的.
第27節.
以色列人住在埃及的哥山地.
就是蘭塞那個地方.
他們在那裡置了產業.
並且生育甚多.
意味著沒有那麼快返上迦南.
甚至乎暫時不返迦南.
放棄這個應許之地.
相當大膽的.
所以如果不是有確實神的應許.
神的心意 他們是不敢這樣做的.
但是現在我們就看到 雅各住在埃及17年.
雅各平生的年日是147歲.
接著第29節.
以色列的死期臨近了.
他就叫了他兒子約瑟來說.
我約在你眼前蒙恩.
請你把手放在我大腿底下.
用慈愛和誠實待我.
請你不要將我撞在埃及.
我與我祖父同垂的時候.
你要將我帶出埃及.
撞在他們所撞的地方.
約瑟說 我必遵著你的命而行.
雅各說 你要向我起誓.
約瑟就向他起了誓.
於是以色列在床頭上敬拜神.
這個就是雅各的一生.
最後的17年.
他130歲的時候就下埃及.
在埃及住最好的地方.

$^{521}$有最好的那些兒孫.
伴隨著他17年的安享晚年在埃及.
所以雅各或者以色列在埃及.
一直這樣來到等候.
不過是他自己一個人可以回到加拿大.
就是他的骸骨上加拿大.
所以又是一個很浩大的葬禮.
都差不多是僅次於法老的國葬.
約瑟幫他的父親.
棺木運上加拿大以撒和亞伯蘭的墓地.
就是麥比拉田間的洞穴.
這樣去葬他.
就是說他們都是等候應許之地的.
所以約瑟亦都難免他將來離世.
他都是要他的後代要炎炎的起逝.
似乎當時約瑟死都已經是有一些改變.
不容易再有這樣的單獨的國葬.
要來到去上加拿大.
所以就出現了我們看的就是.
《出埃及記》的這一段經文.
所以這些串聯都是很重要的.
作者都是很有文學的修養.
將兩樣東西成為一個橋樑.
就是這個家譜.
一到第五節的家譜.
來告訴我們亞國家已經在埃及.
一路等候 一路繁衍 一路生養眾多.
生養眾多必然會帶來一個威脅.
埃及是否容許這樣的不同國籍的人入籍埃及呢?.
還是這些仍然是一些異鄉的人會有異見呢?.
那就麻煩了.
所以我們就要看到.
果然這些本來是美好的過去的恩典和祝福.
成為新王 新的轉變帶來的威脅.
帶來的窒度等等.
所以我們看到這個情況.
作者都用幾巧妙的說話告訴我們.
當埃及王 當然是新王.
可能是內部的叛亂.
又或者可能根本是富不過第三,第四代的埃及王.

$^{561}$根本就不再紀念約瑟的功績.
或者刻意淡化這個功績.
就帶來你的多就變成威脅了.
你的強大就成為了我們埃及人的眼紅了.
所以你想一下 如果作者用這麼重的字眼.
第六節開始.
約瑟和他的弟兄那一代的人都死了.
不過第七節 以色列人生養眾多.
並且繁茂 極其的強盛.
滿了那個地.
你想一下這些肯定是用四個這樣的形容詞.
是這麼正面講以色列人的繁衍.
48年可以是很大的倍數了.
70人可以變成200萬人.
所以這裡有一個很大的張力.
就是視乎那個領導階層.
怎麼看這個生養眾多,繁茂,強大.
以及遮滿了整個地呢?.
當然是誇張了 埃及都相當大的.
但是如果用這樣的張力呢.
就告訴我們知道 這裡就是對應了.
首先就是如果在希伯來人的眼中.
這個就是神給我們的祝福了.
十二章,十五章 不再詳細讀了.
三段都是應許亞伯蘭.
你的後裔好像海邊的沙,天上的星那麼多.
我要為你和你的後裔去立約.
賜福他們去迦南地.
所以你可以看到48年之後.
在整個的埃及不僅沒有萎縮.
而是這個民族按照著神的應許和祝福.
一直這樣去成長 多到一個地步.
根本他們已經看到神的豐富的供應.
以及對這個民族一直的掌管.
所以我們都是要不斷回望.
因為接下來我們有45週年的年報和刊物.
都很想去回顧神過去在錦仙家的發展.
無論去到孫德堂,洪恩堂,嘉欣堂.
我們一直的發展.
所以如果我們能夠從這個角度去看.

$^{601}$你就看到現在在埃及有200萬人.
是一個很大很明顯的祝福.
我們應該常常對於神在過去的那個應許.
現在看回那個祝福.
我們應該帶著一份的感恩和謙卑.
來看到都不是哪個人,哪個牧者的策略.
對應成功 而是真的我們要看到是神.
在不同的弟兄姊妹裡面的工作,那個感動.
以致我們願意這樣一起擺上.
所以現在就是你看到這個祝福.
但是聖經都告訴我們.
既然阿摩利人的罪孽滿盈.
15章我們上個禮拜讀的.
就是成為了埃及人對以色列希伯來人的恐懼.
因為阿摩利人的罪孽滿盈.
加拿大人的罪孽滿盈.
神要懲罰加拿大人.
所以原來就激動了埃及在那個政治的體系裡面.
出現了一個所謂的新王.
這個新 歷史就不是太多著墨.
因為古代的時間也都沒有那麼仔細記載.
這些內部的衝突.
特別如果是自己的內亂,自己的奪權.
有時候都報喜不報憂.
未必會寫得那麼詳細.
所以在這樣的情況下.
這個不認識約瑟的新王.
我就找了一些資料.
一般的解經家其實有兩個解釋的.
我覺得這個解釋應該相對比較近伴一點.
這個出埃及的歷史當時的背景.
大概是公元前的12世紀.
他不僅在帝國裡面的權力去到巔峰.
而且他更加理解了出埃及的罪事的張力的關鍵.
就是一個很重要的歷史參照的人.
這個叫拉美西斯兒世.
是他的兄弟感受到希伯來人是一個威脅.
所以他就要用一些計謀來針對以色列人.
第八節開始.
由不認識約瑟的新王起來治理埃及.

$^{641}$對他的百姓說 看啊 這以色列民.
聖經很特別的 這個這麼重要的人.
帶來一個這麼大的敵黨神,對抗神.
甚至要將整個以色列族毀滅.
用一種淘汰近乎這麼殘忍的來去殺嬰.
讓他慢慢的武藍丁被同化.
都是一些種族滅絕來的.
但是聖經就好像沒什麼興趣.
記載這個君王拉美西斯.
歷史就說他挺厲害的 相當傲慢.
自稱是神 當然他的國力也很強.
但是聖經作者就沒有興趣記載他的人名.
反而那兩個無名的小卒.
就是那兩個修身婆.
這個斯弗拉和普亞兩個是有名有姓的.
所以聖經又告訴我們.
原來往往我們不要小看小人物.
在聖經裡面 經常是那些不起眼的小人物.
是被神用的 比如在《路加福音》記載的.
那些牧羊人按照著他們上班的當更的更次.
來到在野外看守羊群 又凍.
不是美麗的 聖誕卡的那個氣氛.
其實很凍 又天黑.
突然間有些大光 天使顯現 嚇死的.
就是蒙蒙鬆鬆的 真是嚇死的.
就是你回家 如果門口的那個看更.
睡了覺兩三點 你進去狂歡.
看完演唱會12點多回來.
那個看更睡了 你叫他一叫 他整個彈起的.
就是很害怕的 所以這些牧羊人.
其實他們是很低下階層來的.
但是神向他們顯現.
不向這個希律王顯現的.
那些牧羊人是去到伯利恆.
看到那個嬰孩,包著布,臥在馬草.
這樣的氣氛的.
你見到有些很無名小卒在路加坊經常都有的.
例如那個寡婦阿娜 她已經83歲了.
在聖殿來敬拜神.
有那個老人家西敏 很多這些市井的人.

$^{681}$幫耶穌揹著自駕的那個.
加里尼人那個叫做西門 還是那個魯夫.
那個名字我現在忘記了.
很多這些無名小卒的 比如在舊約也有的.
以利沙那個僕人叫基哈西.
記載在這個阿曼的家裡.
對不起 那個賴曼將軍的家裡.
有一個侍女 小的侍女.
是曾經以色列被人擄了去這個亞蘭國.
在這個阿曼將軍的下面做一個小侍女.
因為她一句話 你這麼大麻瘋.
你都搞不定 又沒有神醫又不行.
醫生又醫不了 唉 簡單啦.
你來到我們以色列 我們有先知.
你來找以利沙 必定會好的.
一個小侍女一句話 又改變了這個賴曼的一生.
成為了敬虔的基督徒 賴曼.
你可以想象得到 神用小的人物.
所以這兩個修身婆 斯弗拉和普亞.
但是對當時的這個埃及王 隻字不提.
他叫聖明神 他的姓名.
不是作者摩西不知道.
摩西是那個時代的人.
其實是非常顯赫的君王.
能夠夠膽來統治埃及.
又能夠覺得希伯來的人是危險.
要做一個滅絕 他們的這個殺機.
其實是相當大的勇氣才能做到的.
第十節 賴曼 我們不如用巧計待他們.
恐怕他們多起來 日後如果遇到什麼爭戰的事.
就聯合我們的仇敵 攻擊我們.
離開這地去了 我想現代都是.
不同國家都有不同的移民政策.
你要能夠入籍 甚至乎.
你是不是真的可以入籍之後.
就可以參選成為國會的一些議員等等呢?.
很多的東西 或者你的人數多.
會不會沖淡了將來我們要投票的時候.
一些重大的決策呢?.
任何一個地方的國家都會想的.

$^{721}$所以你看到現在英國也好.
歐洲不同的國家也好.
都有一些保護自己的政策.
都對外面的移民實施有很多很多的界限的原因.
都因為真的是這樣的 也都絕對是火口飛的.
不過用一些這樣的手段.
它的內心就是基於一個嫉妒.
以及基於一個恐慌的懼怕.
所以第十一節 於是埃及人派督工的轄制他們.
加重擔苦害他們.
他們為法老建造兩座積火城.
就是北東和南塞 只是越發苦害他們.
他們越發多起來 越發蔓延.
埃及人就因以色列人受罰.
埃及人嚴嚴的使得以色列人作弓.
使他們因作苦弓 覺得命苦.
無論是和泥 是造磚 是作田間.
各樣的工 在一切的工上都嚴嚴的待他們.
開始不再是友善了.
不再是大家和睦共處了 是因為這個威脅.
但是如果我們是新約的基督徒.
看這個歷史的時候 是神的時間到.
希伯來人可以準備回到迦南地 應許之地.
永久居住在應許之地.
但是神可以透過這樣的一個局勢的轉變.
這個局勢絕對不是神去感動這個埃及王.
讓他突然間變成一個惡心.
完全是出於他的本心 出於他為自己的國家去打算.
出於他為了自己的民族的利益來去針對.
是為針對而針對這樣來對這個希伯來人.
實施一些這樣的苦待 用盡各種的手段.
就是要讓那些藍丁做到有日無夜.
做到筋疲力盡 要致服他們.
你們可以在這裡住 但是你們不會有我們的福利.
你不會有我們埃及人那個原先的享受.
和各方面的福利.
所以慢慢你不喜歡就麻煩你離開.
你就不要想著 因為你所謂的老祖宗虐室.
你們可以得到一個最惠國的待遇.
你們不會得到最好的福祉 那個福利等等.

$^{761}$這些其實真的我們隨時都看到的 在不同的國家.
香港至少七年 一下來之後 無論你多艱難都好.
你都要七年才能夠所謂的永久居留.
才可以申請公屋 申請綜援等等.
這些都某程度為了某個地方的本地人.
又不能夠說絕對錯 絕對是苛政的.
最多你說能夠NGO 我們可以發揮一些的聯盟幫助.
因為你來到這個地方.
表面上你應該是有足夠的那個金錢.
才能夠去到一個新的地方的.
沒有想過你這個移民的團聚.
是需要政府馬上接濟你的.
所以就算去加拿大都是移民團聚.
都看你有沒有那個足夠的錢等等的.
當然投資移民那些零計.
技術移民那些就是根本很想你去.
那就是都是借助你那個技術.
但是一般的你不是在這裡有交稅.
不是在這裡有任何的貢獻.
更加不是在這裡出世等等.
任何國家都要加一點點的限制.
是無可厚非的 不過這裡就好像要刺激他們.
讓他們自己甘心樂意這樣去離開埃及.
或者甘心樂意成為一個不敢反抗的.
二級的二等的公民.
這樣讓自己知道你們是二等公民.
你們不是埃及人.
所以這個就是法老想用的政策.
對於本地人來說是全力配合的.
因為一講到那個資源的稀少和爭奪.
這個就真是任何本土主義.
任何的民族派都是受落的.
所以現在歐洲的極左 英國的極左.
其實保護主義帶來的就是那種心態.
見到人家多了 勤奮了 賺到錢了.
那就看不過眼了.
因為自己的生活的節奏被打亂了.
那你就會有兩種的制度.
又肯存錢 又肯買房子 又肯勤奮工作.
又叫子女無論如何都要出人頭地.

$^{801}$你要讀書 你都要比別人更加勤力.
因為有這個危機感的新移民.
或者二等公民.
所以對於埃及人來說.
他們覺得這樣的苦袋.
就是說肯給你工作.
你屬於我們這個地方的.
不過你就要加倍的努力.
當然這是一個剝削.
這是一個有計劃的一種對這個民族的一個針對.
這樣的一種仇視.
源於一個懼怕和內心的疾度.
但是聖經告訴我們.
以色列人繼續在神應許亞伯拉罕的福上面.
越發的爭奪 越發的一路一路的發敗.
當然不斷的受壓.
一般的以色列人始終所謂的佛都有火.
如果神又不介入.
因為他們看不到的就是神的介入.
看到的就是法老越來越.
將那條綁著他們的繩子越拉越緊.
在全方位上面對以色列人.
實施這些嚴苛的新的政策.
就是奴役他們.
接著第15節 第二個計謀.
有希伯來的兩個修身婆.
一名斯法拉 一名普雅.
埃及王對他們說.
你們為希伯來婦人修身.
看她們臨盆的時候.
如果是男孩就把她殺了.
若是女孩就留她存活.
很少要男孩不要女孩.
很少的 因為男孩才是自己的繁衍.
那你就可以想象得到這個計謀.
其實真是人都看得到 但是就是要這樣做.
做到出了面 就是說原來我們用那種壓迫.
已經壓到最盡了.
不單止這個壓迫不影響了他們.
一個一個 你要吃樹皮了.

$^{841}$當然是回大陸了 香港都沒有東西吃了.
三年零八個月 又這麼多的困難.
不如慢慢自己甘心情願離開.
又或者無論如何.
都做一個絕對順服的一個二等公民.
但是兩者都不是的時候.
埃及人都不知道怎麼辦.
終於都想到一個獨計了.
就是由源頭開始.
似乎又是暗中了 不敢太貿然.
因為那些壯丁還在做那些奴役.
還在做那些幫法老來建那些公司.
那些積課城 還有做那些磚.
所以如果很明目立法.
就群情洶湧去反對了.
所以似乎由那個源頭開始.
就是說收買兩個收生婆.
叫她 當然不是她們兩個做.
因為這麼多婦女 她們又生養這麼眾多.
一定是這兩個好像是一個領班.
類似贊沃醫院的總護士長.
安排那些護士看到希伯來的婦人生兒子.
就用各種的辦法.
那個媽媽都痛得不得了的時候.
就馬上接了她 幫她剪臍帶的時候.
就弄死她 說她難產.
這樣都真是很殘忍 相當之惡毒.
因為任何一個小生命 無論什麼國籍.
一生出來 其實樣子是相當漂亮的.
差不多近乎天使那樣的嬰兒.
沒有人忍心對一個嬰兒.
差不多是人類當中最軟弱.
最無助的一個生命.
你殺一個男的大人 你都要跟他打過.
或者你要暗算他 用飛刀 用槍.
但是一個嬰兒 根本上是他什麼都不懂.
什麼都沒有辦法去抵抗的.
那個媽媽又生了.
又是在那個劇痛裡面都起不了身的.
所以你一抱了她 就抱了她了.

$^{881}$一交換了就交換她了.
所以這個都是一個挺惡毒的一個計謀.
所以如果這兩個收生婆.
這麼忍心做這些事情.
那個報應都相當大的.
因為真是生命來的.
不過第17節就告訴我們.
在這樣的時機 出現在這樣的時代.
都好像真的偏偏選中你 你沒得選擇的.
問題就是 這個就變成不是一個生存這麼簡單.
而是一個關乎信仰的表態,抉擇.
但是收生婆敬畏神.
不照埃及王的吩咐行.
竟存留男孩的性命.
所以《一節聖經》就將這兩個無名小卒的英雄.
斯弗拉和普雅這兩個希伯來的收生婆敬畏神.
在那個時代 那些男丁就出去做苦工.
那些女的就可能做其它的家務.
但是這兩個是做收生的.
就是產科來的 產科的護士.
可想而知在那個風浪的尖端的口.
面對一個很大的挑戰,威脅.
問題就是現在她內心的信仰.
能不能夠讓她站穩在神的旨意,計劃裡面.
神現在定義時間到了.
以色列人因為生養眾多 神要兌現他的應許.
以色列人要離開圍爐之家的埃及.
當年不是圍爐的 當年是和睦共處的.
但是現在以色列人要成為奴隸.
因為是外在的力量逼使他們做奴隸.
但是無論如何 神的計劃要他們安然的離開.
上去這個迦南地.
所以整個責任落在這兩個女士身上.
很大的誘惑 做還是不做.
做就要面對那個良心很嚴重的譴責.
就是你殺這些鷹派 就是你殺不下手的.
所以我們看羽裳片的那些傳說等等.
就是狸貓換太子這些事情.
真是諸狗族的 但是不忍心這個太子.
和這個皇后雖然失寵.

$^{921}$但是真正的太子 那些大臣.
真是不會埋於那個良知 都要真是這樣做.
保存那個太子的生命 將來成長.
就幫媽媽和那個國家來恢復那個帝位.
所以人生如果我們真是要去到這些抉擇.
所以這兩個女士成為了以色列人的國母.
你想一下比摩西還要來得厲害.
沒有她們兩個都不會有摩西 摩西都快要淹死了.
本來摩西因為三個月大已經不可以再收藏了.
就要送走她 因為似乎家家戶戶找藍丁了.
現在就是暗中在產房看到有藍的就要來殺.
那個時候也沒有那些超聲波.
去照誰是男還是女等等.
古代就沒有這些所謂 不知道希伯來人會不會.
肚子圓一點尖一點 這個女的 我不知道.
但是無論如何 一生出來就要抉擇.
是不是要將她拿去送給那些埃及的士兵.
這個是希伯來的男嬰要拿去殺的.
要拿去焚化的 很嚴重的這些罪惡.
所以埃及王將來要面對神的刑罰都相當大的.
因為由他想出來 首先就是將那些爸爸.
父輩的藍丁 壯丁全部奴役 奴役到死.
奴役到手無寸鐵 根本沒有力量反抗.
也都不敢反抗 但是越來越精壯.
沒有什麼傷損 仍然一樣生養眾多.
可想而知 他們都看到.
這個其實是神的祝福的一個印證.
但是埃及人 他們怎麼能夠從這個角度.
去看神的祝福呢?只會看到.
這個威脅是越來越白熱化 越來越嚴重.
所以當修生婆不肯的時候.
第18節 埃及王召了修生婆來說.
你們為什麼作者是全流男孩的姓名呢?.
相信不是一兩個的 是當時很多.
修生婆怎麼去發動那些一種這樣的抗爭.
一種這樣的偽像 這個這麼重要的理念.
好像不抵抗運動那樣 去貿然的抗議.
管理層對於這些新生命的屠殺.
是不能夠接受 是必然要反對.
所以他們的解釋 第19節.

$^{961}$似乎都是暗中進行 因為如果是明顯的.
根本就不需要理會 直接就抓去殺.
就是說暗地裡 又不敢讓那些壯丁群起洶湧.
只怪真是不幸 阿B就夭折了.
一出生就夭折了 那就沒辦法了.
又要等下一次了.
所以這個似乎是暗中低調的做法.
但是當這樣的時候 修生婆怎麼樣呢?.
修生婆對法老說 因為希伯來婦人.
與埃及的婦人不同 希伯來婦人本是建築.
修生婆還沒有到 她們已經生產了.
很特別的一個解釋 其實都是因為她敬畏神.
也是因為這樣 神給她們有一點點的智慧.
來勉強的說了 希伯來的婦人.
很自然的 個個都盛產 個個生完已經是男孩子了.
已經是好像公告於天下.
不能夠再採取一些暗中的 將她抱走.
將她去屠殺 所以可能也都有一些.
類似政府的生死註冊處.
要登記孩子的出世紙.
這個生命從那一刻開始 在香港政府就有記錄了.
就是如果她無故消失 就要去人口登記處.
要解釋為什麼她什麼 所以到死的時候.
我們就要去拿死亡證 就要註銷她這個人.
在香港的那個戶籍 例如曾銳貴下個月死了.
就要去生死註冊處 要拿回那個死亡證.
所以其實從出世到死 政府都有你的記錄.
你現在去到這個歲數 現在在這個歲數裡面.
還有多少人在香港出世的 到現在.
比如60歲了 60歲的 1960幾年出世的.
現在有多少人呢 全部有記錄的.
有了所有的姓名的 所以呢.
這裡就是他們已經沒得再隱藏了 不可以殺了.
所以就做了人口登記 所以差不多就是這樣.
應該再要加重那個刑罰 所以後面就這樣說.
第20節 未講之前 神就祝福修生婆.
神厚待修生婆 以色列人多起來 極其的強盛.
修生婆因為敬畏神 神便叫他們成立家室.
他們自己都有自己的美滿家庭.
修生婆 也都是似乎一直他們出埃及.

$^{1001}$這兩個無名的英雄的故事都會不斷流傳在.
因為那些希伯來的婦人 10月懷胎.
就算入院 可能在家裡 那個時候沒有醫院.
都要可能幾天臨盆要生.
那修生婆都會看一下 生了沒有.
那個陣痛去到多少度.
有沒有那些胎水 這樣就會看.
那個過程都會對話 這兩個都是希伯來的婦人.
都會用希伯來話對話 就是說那些生孩子的媽媽.
都會知道自己的孩子 如果是男孩子.
都可能會死的 所以這兩個無名英雄的故事.
應該就是已經在以色列人當中流傳.
成為了民族的英雄 他們的勇敢.
他們的勇氣可以救回那一代人的南丁.
可以成為了 可以出生 可以成長.
可以將來出埃及等等.
所以你可以想像到 第二章就記載摩西.
那個被拯救等等 你就可以想像得到.
當時救嬰的行動都是一直在暗中進行.
我們都要有一些思想 有一些想像力.
這兩位老人家慢慢一直隨著那個政策.
越來越嚴峻 他們要抵擋的那個黑暗的勢力.
就已經都不是他們兩個人能夠可以左右了.
但是神就厚待這兩個修身婆.
他們有美滿的家庭 但是聖經就沒有記載.
他們的丈夫 或者他們的兒女等等.
不過在當時的以色列人當中.
我們今天就沒有什麼感覺.
但是對於當時的希伯來 特別那些婦女.
對這兩位的阿瑪 真是愛護有加.
沒有將自己的兒子送給法老.
來將他弄死 你想一下是多大的一個恩情.
救一個生命 所以這些在神永恆的角度裡面.
將來我們都會見到這兩位的姐妹.
在永恆裡面 她真是救了很多很多的人.
當然22節 法老吩咐他的眾民說.
以色列人所生的男孩 你們都要雕在河裡.
一切的女孩 你們要存留她的性命.
就是說由暗中,低調,由醫學的角度.
一出生就夭折 沒辦法 救不了.

$^{1041}$那個臍帶綁著頸 或者有其它的不知道什麼病.
那個嬰兒就呼吸不了 一拍他的時候不懂得哭.
各種的醫學理由 總之就是淡化.
男丁為什麼一出生就會死呢?.
為什麼這麼多男丁會夭折呢?為什麼呢?.
或者根本消失了呢?都沒辦法.
但是後來就不再低調地,隱瞞地做了.
而是吩咐埃及人要這樣.
就是說挨家挨戶去找以色列的嬰兒.
男的嬰兒就要將他弄死了.
你可以想象得到無名的英雄的偉大.
所以神在歷世歷代都興起無名的英雄的.
所以當然我們求神給我們.
不是我們自己自我膨脹.
總之我們在神給我們的崗位.
我們有我們自己要站著的地方.
我們就為主去站著.
我在預備,準備的時候.
我突然間想起一個很重要的人跟大家分享.
這些我們都講了 最後我想跟大家分享一個人.
都跟這兩個修生婆是完全一樣的處境.
這裡告訴我們法老的恐怕.
就是那些以色列人越來越多.
越來越怕他們背叛 因為是一個威脅.
你想一下如果我們中國人在英國.
如果越來越多你的票數是左右你的決策.
你都要對中國人過年都要穿上唐服.
學一下寫書法 叫你的子女學一下普通話.
都因為這個選票的人口相當大.
比如印度裔等等 其實很實在的.
但是現在它怕的更加嚴重.
就是怕它聯合那些外族來侵吞.
是不是因為它以前都是外族進入了埃及.
去打敗了當地的埃及王.
而它自己成為新的法老呢?.
這些都有這個可能性的 但是無論如何.
人的恐怕和修生婆的敬畏神.
大家都是fear 但是對神的敬畏.
就使得修生婆得到自己百姓的愛戴.
和最終都能夠成就神的心意.

$^{1081}$我想起有一個人 不知道大家有沒有聽過.
這個是當時在45年 應該是44年.
當然應該是40年了 當時納粹德國.
這個叫做參元遷謀 他是一個東正教徒.
相當有學問的一個領事館的外交的領事.
就駐守在立陶宛 雖然立陶宛是很小的地方.
當時40年的時候 納粹德國開始和日本.
和意大利有俗稱聯盟.
所以日本就在東亞的亞洲來打發動戰爭.
德國就在歐洲發動戰爭 要大家一路這樣打下去.
就大家共同併吞世界.
所以當時它就生在這樣的一個大時代.
當時德國已經入侵波蘭的了.
大量的波蘭的猶太人就走了.
當時仍然還是中立的立陶宛.
就要急需一個過境的簽證.
就是離開立陶宛之後.
就是希望經過蘇聯,日本再轉去美國.
因為美國沒有仗打的.
接著就去中國的上海.
因為中國還是盟軍的參戰的.
就是站在美國那邊的.
所以當時都有.
因為上海當時是需要簽證就可以進的.
但是要先通過蘇聯的國境等等.
所以如果沒有簽證.
你沒有一些領事的簽證.
蘇聯就不會開放給你經過.
所以就困在立陶宛等待納粹來進攻.
於是在40年的時候.
因為這個參議員遷謀.
有很多的猶太人去到日本的領事館.
因為其它的領事館都關閉.
都不再發簽證了.
因為不想受威脅.
所以日本當時還沒有.
因為日本和納粹德國是聯盟.
所以它就打電報回去日本的領事館.
就說現在有這樣的情況.
猶太人現在基於那個人道主義.

$^{1121}$是不是應該發簽證給他們去經過蘇聯.
由蘇聯再去其它的地方.
例如去日本 因為你有日本簽證.
因為去上海不需要簽證.
但是你要通過蘇聯的國境.
三次的急電.
當時是用電報 會有WhatsApp等等.
但是日本的領事館基於不想得罪納粹德國.
所以就禁止參議員遷謀發這些簽證.
那時候已經有千多個納粹的猶太人.
去日本的領事館每天去求可以發簽證.
但是它一直在等待總領事館的批文.
外交部的批文 但是堅決.
如果你做的話 你就是得罪我們日本的國策.
所以就禁止它做 還要它立即關閉領事館.
要它盡快離開立陶宛.
所以基於這樣的原因 它就有很大的爭執.
它不知道怎麼辦 基於它是日本的領事.
它當然一定要服從.
但是它看到的就是有千多個.
在它的領事館外面的猶太人.
所以它都好像這兩個修生婆一樣.
當時它得到它太太的無限的支持.
所以一個言內助真的很重要.
給它最重要的就是信心 內心那種的安慰.
太太用耶利米蘇來安慰丈夫.
它不管那麼多了 被解除它的職務的威脅.
那個危險 它就發簽證.
但是因為那個時間非常緊迫.
它是必須盡快關閉領事館.
要它立即離開立陶宛.
所以它連續十天不斷地簽發.
去到一個地步 根本上手已經軟了.
後來它急於要走 它都不管那麼多.
就這樣簽個名 蓋了日本領事的人.
就坐火車的時候 一大疊就扔出去.
你撿到 你自己填回你的名字 你就過了.
最終當然了 因為領事館知道它的情況.
就送了它去蘇聯 在蘇聯囚禁了一段時間.
後來和平了 大概一個五幾年.

$^{1161}$它就送回日本 回到日本之後.
果然外交部低調地暗中解除它一切的職務.
也都不對它所做過的事有任何的表揚.
所以它一下子失業 它後來在五幾年的時候.
因為日本都很艱難 因為戰後的重建.
很多日本人都很慘 自殺等等.
它不能夠做外交家 它只能夠為了戶口.
就在自己的家鄉裡面 做那些售賣燈泡.
做那些雜工等等 後來後來輾轉.
才能夠找到一些翻譯的工作去做.
就這樣過了它的人生 不過.
過了差不多30年之後 這位猶太人.
耶福斯雅尼希尼 就是1985年的時候.
大概那個時間 就是這個少年人.
當時拿到他的簽證 可以和家人離開了.
後來輾轉回到以色列讀書勤奮 做了外交官.
就成為了以色列的外交官 但是他.
當然就是很想和日本的外交部去找回.
參元遷墓現在在哪裡 我們很想來多謝他.
但是外交部完全沒有任何的回應.
也都拒絕透露參元遷墓現在在哪裡.
是不是已經離開世界.
這個猶太人他很有智慧 他很有那份感恩的心.
就是他申請調任以色列駐東京的大使館的領事.
當他去到東京的時候 他就民間自己去搜尋.
自己來找這個參元遷墓 所以就是.
在1985年的時候 終於找到參元遷墓.
就有一個真是世紀的擁抱 彼此的流淚.
因為參元都沒有想到 原來過了30年了.
他都沒有想到自己的簽證可以救到差不多接近.
有6000人的猶太人 就是他們的後代.
差不多了 所以這樣的情況就可以想象得到.
但是他在1986年 參元就因為年紀太大了.
就離開世界 後來以色列就將一個最高的榮譽.
是給外國人最高的一個榮譽 就給了這個參元遷墓.
現在以色列有一個紀念碑 是紀念這個參元遷墓的.
這位東正教的基督徒.
這位當時要面對一個這麼大的一個將領.
他基於的就是他內心的良心 對神的敬畏.
太太對他的支持 他做了這個真是很英勇的決定.

$^{1201}$所以這個是亞洲版的《舒特拉名單》.
這樣的一個故事 突然間想起.
很近似這兩個修生婆 都很令人感動的.
在立陶宛有他的紀念碑.
紀念這位當年在立陶宛做過的這個義舉.
救了這麼多的猶太人.
所以猶太人對日本人都有這個好感.
所以日本早年的首相 我忘記了是哪一個首相.
去到以色列的時候 都在這個銅像的面前.
有紀念 又感受到我們日本人的前輩.
有這麼大的義舉等等 這些都是題外話.
不過真是我們人生 我們要看看埃及.
這個民族在那個困難裡面.
神用的是兩個小人物.
後來神用這個被拯救的摩西.
因為摩西不是嬰兒的時候想去淹死他.
而是他長大了之後都不能夠再隱藏了.
本來當時埃及人周家周戶去抓嬰兒.
藍丁要來將他整死.
但是仍然我們看到有不同的人.
為著神的工作 來做這樣的準備.
所以今晚我們都求主讓我們能夠透過.
我們先講第一章這個引言.
作者都用了很長的篇幅告訴我們.
這個民族的過去和現在的光景.
現在的光景表面上很慘 寄人籬下.
還有受壓迫 但原來從來都不是寄人籬下.
是神叫那個落埃及.
是要他們全家落埃及 是一個後代.
是一個等待 所以如果我們都能夠.
從一個宏觀去看回我們自己的人生.
或者我們家族的人生.
我們就不會有時慌失失.
因為以色列人第二章開始.
就不斷的去呼天搶地.
很多的哀聲去到神那裡.
所以如果我們能夠用宏觀去看.
我們當然容易 因為我們看舊約的歷史.
但如果我們在那個時代.
我們一樣的 白天做到死 沒力氣了.

$^{1241}$回到家裡你都是這樣繼續痛苦.
去到安息日 不僅是沒得安息.
更加可能是又怕你結集太多人.
所以可能那種哀怯哀求.
都是一種可能是整個民族是一種無望.
是一種找不到前路 找不到明天.
但原來這個找不到明天.
才是神整個救贖計劃的其中一點.
第三章就有一個拯救者回來了.
就是由這個曠野回到埃及的摩西.
所以神的時間到了.
我們就會看到神那個改變的能力.
可以一下子扭轉那個局勢.
願神幫助我們 我們都生活在現在這個世代.
我們求主讓我們能夠個人 整個教會.
整個香港的教會都能夠找到我們的定位.
我們在神的面前 我們怎樣去做.
譬如現在有很多新的香港人來了.
我們應該用一種擁抱的心 不要分你我.
福音是沒有分你我的.
不分猶太人 希利利人.
不分圍爐的自主的 不分男女的.
我們應該擁抱 不應該看他們是新移民.
是沒有教養的 不應該這樣看.
我們應該使他們因為福音.
他們慢慢成為有教養 成為有禮貌.
等於我們上一代的人一樣是沒有禮貌.
從大陸走難下來.
所以我們應該看到我們自己時代的定位.
既然現在有一大幫人離開了香港.
去到歐洲 他們有他們神的帶領.
我們既然在這裡的時候.
我們應該就要擁抱這個時代的人.
應該同一個命運一起在這個空間裡面.
一起去尋索福音 那個生命改變.
我們一起禱告 天父啊.
我們看到原來小人物在你的手中.
在你的眼裡從來都不是小.
很多時候是我們對於權力 對於我們自己的驕傲.
很多時候我們會有一種看不起其他人.

$^{1281}$或者我們有一種很容易的自大,自義的心.
以致我們很多時候會對比我們厲害的人.
就會去膜拜,巴結.
很多時候我們都會是這樣.
有時候甚至在教會.
我們都會看到有時候那些權力的爭奪.
我們真的很盼望神保守 叫我們不要這樣.
反而我們盡心 在每一個時代.
神給我們的崗位 我們盡忠.
有些話我們真的要講 來維繫那個真理的.
給我們帶著一個謙卑的心.
不是一個自大的心來去陳述.
也都有很多的事情 我們需要用勇氣.
好像兩個修生婆 好像這位日本的參議員千茂.
在他的人生 因為這個這樣的動作.
使得他差不多下半生都是很老倒.
得不到日本人的愛戴.
如果不是這個猶太人.
千里迢迢的來尋找他.
化身做一個駐東京的領事.
來暗中的去民間去報館去尋找.
真的永遠沒有人知道有這些這麼重要的異舉.
多謝主讓我們看到你的救贖計劃.
在以色列人當中 願神幫助我們.
我們在21世紀 我們仍然在香港這個大都會.
現在大灣區很多 7千多萬人口.
整個廣東省已經根本都做不完.
更何況有其他的省市 所以我們真的很盼望.
我們不要拒絕對內地的那個厭惡.
反而我們應該對不同的國家的人.
很多都沒有福音的.
盼望我們能夠 神給我們有不同的機會.
不同的辦法去將很多內地的同胞.
帶進救恩裡面 現在差不多很多的大學.
那個內地生比本地生還要多 求主幫助.
知道未必容易 但是希望無論是政策上面.
或者教會能夠結合更多的力量.
例如我們盼望神興起很多基督徒的講師,教授.
他們能夠在福音的事公上面.
一方面他們的生命見證 還有他們的組織.

$^{1321}$能夠帶給教會 可以有不同的方便.
可以帶領很多的內地的大學生來信主.
願神在這些事公上面給我們有智慧的心靈.
給我們有一份的思考 怎麼做是更好.
我們仰望著 多謝你讓我們在今晚.
一起來聆聽你自己的說話.
我們獻上感恩 奉耶穌基督的名 阿們.
\newpage



\section{出埃及記}
\label{sec:4cUSeyCs18Q}
\textbf{2026-05-29 周五晚查經班《出埃及記研讀(三)神所用的小人物:收生婆,母親,公主如何瓦解暴政》 講員:曾裔貴牧師}
\newline
\newline
連結: \href{https://youtube.com/watch?v=4cUSeyCs18Q}{\texttt{https://youtube.com/watch?v=4cUSeyCs18Q}} ~~~~ 語音日期: 2026-05-30
\newline
\newline
\hyperref[sec:kZF3v6CUrpM]{\small{< < < PREV SERMON < < <}}
~
\hyperref[sec:index_chronic]{\small{[返順時目]}}
~
\hyperref[sec:index_scriptual]{\small{[返順卷目]}}
~
\hyperref[sec:HMEU7PxsF0Y]{\small{> > > NEXT SERMON > > >}}
\newline
\newline
$^{1}$晚上可以一齊去研讀《出埃及記》.
和大家一齊同心禱告.
多謝天父我們在晚上再一次打開聖經.
我們很盼望我們研讀的《出埃及記》.
我們能夠從波西作者的敘述.
一直由《創世記》的應許.
到現在準備要《出埃及記》的時候.
讓我們看到神是如何來預備你自己的子民.
他們的心路歷程.
或者他們在極大的壓迫之下.
他們可能真的會垂頭喪氣.
因為那個壓迫的逼迫是非常的嚴重.
多謝主往往一個的事工能夠成就.
是因為有神興起的工人.
是因為有一個離開了埃及.
回到埃及的摩西阿魯等等.
我們看到神的準備.
所以我們都盼望透過在這幾章.
對於摩西那個人生的準備.
來到他有深入的認識.
以致當我們來到我們自己的時代.
我們自己要接上福音的棚的時候.
都能夠有所的頓悟 有所的明白.
有所的反思.
我們的人生是如何來交託給主.
摩西是在八十歲再一次蒙召.
當然我們絕對不是說要等到八十歲.
這個不是一個必然的真理.
而是在舊約的時代 神自己的作為.
所以我們盼望神給我們.
都有一份的整全的對真理的認識的心.
以致我們能夠在神面前都有這樣的追求.
我們在今晚等候在你的面前.
我們這樣祈禱 奉耶穌基督的名.
阿們.
我們進入第三講.
今晚用的題目.
其實大家第一章第二章.
上個禮拜大家已經掌握到了.
那兩個修身婆.

$^{41}$斯瑪拉和這個.
突然間又忘記了她的名字.
那個生命的流露.
和普雅 對不起.
我們就可以看到.
來到第二章仍然都有這樣的小人物.
而且不止一個.
甚至歸為法老的女兒.
這個埃及公主.
到摩西自己的媽媽約基別.
和她的姐姐米利奄.
可以想象到神出人意表的作為.
所以我們今晚就用這個題目.
神所容的小人物.
是一個大時代.
因為當時約瑟已經離開世界很長時間了.
有認識約瑟的身亡.
那個重點是在說.
對於以色列人.
不再是一份對他們懷有感恩的一個民族.
而是一個有阻礙埃及自己的發展.
甚至乎會成為一個敵人.
如果將來有戰爭.
可能這個民族會聯絡其他的敵人來攻打.
所以變成人越多就成為一個遺患.
所以這個在政治上都是很合理的.
在歷史上我們都有一些文獻告訴我們.
的確是有這樣的問題發生過.
我們上個星期都講過.
當時的那個王.
我們大概都猜到他那個.
因為他自己那個無論是勇力.
對於自己治理的埃及.
又開始有一種做神的運動.
就是自稱為神這樣敬拜.
所以很自然地對於希伯來人.
就有這種敵視和有一種威脅.
我們上個星期看到的就是修身堂.
今天我們可以來看第一到第十節.
其實主要是想描述.

$^{81}$波西在這樣的大時代.
當時就不只是一個嬰兒被殺或者被拯救.
作者就是透過一個敘事告訴我們.
是一個本來是男孩子.
希伯來的男孩子 嬰孩就一定會死.
死是似乎唯一的那個路.
一出生你隱藏不到 被搜尋到.
你就要死了 如果不是.
可能都會受到牽連的那些父母親.
但是我們就看到神是怎樣來到.
透過頭十節告訴我們 預備了這個嬰兒.
第十一節就是一個很重要的轉捩點.
因為第十一節他已經是長大成人了.
而且還有兩個階段 就是嬰兒的餵哺階段.
就由自己的媽媽這個弱基別來養大.
接著到了一個階段 似乎開始懂事了 懂得走路了.
就送回到法老女兒的王宮裡面.
在法老的王宮裡面來被養育.
那個養育不再是物質這麼簡單了.
而是供書教學 所以這個《善行傳》第七章.
當施提法來描述的時候.
就告訴我們摩西學盡了埃及人一切的學問.
說話能力都是很有條理 很有智慧的.
這樣的一個摩西 所以我們就看這一段聖經.
由一到第十節 第二章.
有一個尼美加的人娶了一個尼美女子為妻.
作者很快就告訴我們 尼美這個支派.
似乎將來都有相當重要的實用的.
當然現在純粹是告訴我們摩西是哪一個支派.
就是亞國哪一個兒子的那個族人的兒女.
這樣來描述 但是將來我們去到尼美記.
應該就去到文授記更加具體.
就告訴我們尼美支派的重要性.
有一個尼美加的人娶了一個尼美女子為妻.
女人懷孕 這個女人就是約基別.
她的爸爸是暗烈.
我反過來告訴大家.
生一個兒子 見他最尾就藏了他三個月.
後來不能再藏 因為要長大.
可能會做家族戶去搜尋.

$^{121}$就娶了一個蒲草霜 抹上石漆和石油.
這樣的描述和希伯來文的原文.
和那個方舟的描述是完全一樣的.
而且用的那個字 這個蒲草霜.
就好像一個舊增版的方舟那樣的.
所以都很特別的.
抹上了石漆和石油 將孩子放在裡頭.
打箱子 擱在河邊 當然是尼羅河.
那些羅迪當中 孩子的姐姐遠遠站不住.
都相當懂事的了 米玲林.
都不是跟摩西那麼大 可能一歲兩歲.
你不可能可以這樣遠遠站不住.
可能都是十幾歲的姐姐.
在家裡就已經幫忙做家務.
就可以知道 要知道她就是她的弟弟.
摩西究竟會怎麼樣呢?到底會不會被人抓了.
有那些士兵來巡查發現 因為孩子始終會哭.
或者一直衝到尼羅河.
一直流到一些地方就淹死了.
各樣的情況都會發生的.
但是聖經就告訴我們.
突然之間有這樣的戲劇性的一個轉變.
就是貴為最重要的法老的女兒.
法老的女兒來到河邊洗澡.
她的侍女們在河邊行走.
就是說有侍從 有不同的宮女到處看著.
就是說不可以有其他的人靠近.
她看見雙子在爐的中就打發一個婢女拿來.
她打開雙子看見那孩子 孩子哭了.
摩西都挺醒目的.
馬上大哭的吸引了法老的女兒的注意.
接著就說 她就可憐他說.
這是希伯來人的一個孩子.
這是相當重要的一個描述.
法老的女兒 我們不需要猜.
以為她是誰的棄嬰.
而是相當清楚.
就是說法老的女兒知道爸爸的禁令.
要殺所有希伯來的婦人所生的男嬰.
所以法老的女兒現在看到這個哭的嬰兒.

$^{161}$又認得出她是希伯來人的男嬰.
所以其實在這樣的情況下.
是沒得逃的 就是一定會死的.
你沒辦法想象這樣的戲劇發生.
就是第七節.
孩子的姐姐對法老的女兒說.
就是說 她遠遠看到之後.
因為她是女性 一個小女孩.
不會引起宮女的猜疑.
或者怕著來到去暗殺等等.
不是在擾亂.
所以沒有了一種的懼怕.
米利亞就走到她的身邊.
就對法老的女兒說.
我去在希伯來的婦人中.
叫一個奶媽來 為你奶姐的孩子.
可以不可以 相當戲劇性的.
我們拍電影都不可能這樣.
通常我們都會盡量收起來.
因為那個是一個嚴重的違反.
現在國家法老下的命令.
所以這樣的情況下 我們很難想象.
一個法老的女兒.
看到這個嬰兒而有一個惻隱的心.
可憐她 這個希伯來的男嬰很可憐.
她都還沒有去想怎麼決定.
可憐不可憐 父王的禁令難違.
但是當米利亞真的鼓起勇氣.
跟法老的女兒這樣說的時候.
都想象得到米利亞她的人生將來去到.
《出海道記》第15章.
米利亞就是女先知.
來去帶領那些婦女擊鼓跳舞.
因為過紅海之後 米利亞就這樣來去做戰神.
帶領所有的希伯來的婦人.
都認為是相當有獨當一面的性格.
可能平常做事情的那種領導力.
這個米利亞 所以在這裡都看到一點.
一個小朋友 當然十多歲了 我估計.
就能夠夠膽這樣直去跟這個法老的女兒.

$^{201}$有這樣的一個建議.
所以是相當勇氣才行的.
你想一下有時候那些姐妹 謙卑是好的.
但是有時候叫她做事 叫很多次她都不敢.
因為性格就是太內揚 太柔弱.
什麼呢?就是不用我出面就行了 我會做的.
但是本來這個性格是好的.
但往往就是真的要她出來做一個小小的領導的時候.
有時候就是那種的遲疑 那種懼怕.
但是不如這個天生的性格.
所以你看到米利亞在這個時候把握這個機會.
就不是她的媽媽出面 而是米利亞出面.
第八節 法老的女兒說.
可以 同類就是叫孩子的母親來.
法老的女兒對她說 你把這孩子抱去.
為我奶她 我必給你功格.
那就是相當的戲劇性.
那就是說 叫這個媽媽 就是這個若姬別.
就要回家去養育這個嬰兒.
就不是抓了這個若姬別 入宮成為奶媽.
就在法老的女兒的皇宮成為另外一個宮女 婢女.
這樣來侍候 她的工作就是.
因為她剛剛生完 她有奶的.
就可以來用人奶餵自己的嬰兒.
而不是叫她回家.
因為下面第九節就說.
法老的女兒 對不起 第九節中間那裡.
婦人就抱了孩子去奶她 孩子漸長.
婦人把她帶到法老的女兒那裡.
就作了她的兒子 她就是法老的女兒.
她給孩子起名叫摩西 意思說.
因我把她從水裡拉上來.
摩西就是被拯救 在水裡被拯救的意思.
Moses這個字.
我們休息一下這一段聖經.
我們要學的功課就是Prophetess神的護理.
我們看到有很多東西表面上是非常巧合的.
法老的女兒去河邊 可能這個都是慣常.
但是摩西已經不能夠隱藏在自己家裡了.
就是說周圍那些item的士兵.

$^{241}$是會逐家逐戶 挨家挨戶來去搜查.
有沒有這樣的嬰兒 很容易認的.
認嬰兒就一播出影片 就看到是男的還是女的.
男的就要放進去尼羅河浸死.
這個是第一章說了的.
所以我們又看到好幾個巧合.
在尼羅河本來是男嬰的葬身之地.
這個葡萄箱就成為了摩西的拯救.
使得他能夠從水裡面拯救出來.
就好像和那個羅亞的洪水一樣.
真的被水拯救出來.
當然這個是彼得前書三章十九節那段經文.
藉著水的救 這是彼得說的那八個人.
這就告訴我們 好幾個巧合.
我們要學習一個神學.
就是在《舊約》裡面.
我們看到作者的鋪排 那個敘事.
很多東西表面上是很巧合.
但是當我們現在來看的時候.
我們要用一個神聖的護理.
就是說神有祂的預知 有祂的安排.
有祂在命運當中 祂的主宰.
我們要學習就是原來連法老的女兒.
她那個惻隱的心都是神的主宰.
都是神的保守.
使得那件事好像成為巧合.
摩西被發現 摩西又在哭.
法老的女兒那一刻就不是因為那個皇命.
來怕將她交給自己的一些禁卒等等.
而是反而就是叫一個希伯來的婦人.
來去養育這個嬰兒.
養好之後 懂事了 懂得走就走了.
戒賴了 似乎真的要戒賴才可以入宮.
似乎都是有這樣的思考,盤算等等.
所以這個就變成一個很寶貴的一個劇本.
給我們有很多的空間去思想.
為什麼法老的女兒都能夠有這樣的憐憫的心呢?.
是不是就是說原來法老的命令.
不是每個人都願意去遵守 原來埃及.
不是每個人都是壞人 不是每個人都認同.

$^{281}$這些希伯來人是我們的阻礙 是我們的敵人.
就是說 原來神都可以在我們的行為.
不可能改變的 因為第四章打後.
其實嚴格來說應該從第七章開始.
一直到第十二章 都是十個災的描述.
就是法老不斷將自己的心變肝癌.
那些災禍一放鬆 他就馬上將心變成肝癌.
所以你可以想象得到.
代表著那個敵對是沒有改變.
一直都是 只會更加更加的敵黨神.
但是在這樣的敵黨裡面.
我們又看到法老的女兒是有神的預備.
所以這個就叫做神聖的confidence.
一個狐狸神管理著.
一個人有那個良知 有那個憐憫的惻隱的心.
這個如果在我們讀神學的裡面.
這個就叫做普通的狐狸.
或者說普通的一個恩典.
其實任何的人 包括一些很壞的惡人.
就算是黑手黨那些勢不斬板的.
他們自己對自己的兒女.
對自己的兄弟都很有那個愛心.
這個是與生俱來 我們每個人都有的.
但是當我們要學習聖經講的那個特殊的恩典.
是有聖靈 有神 對神認識的 神的愛.
我們才會擁有的.
所以我們就不是一般性的恩典.
所以這個神聖的狐狸.
我們就要在這裡.
如果我們這樣來用這個神學.
就是神學家做了這樣的詮釋.
我們這樣看這些經文 這些歷史.
就很有意思了 神有祂的準備.
一個本來 其實每個男孩子都是要在尼羅河死的.
但是神可以救起了一個.
而這個就是將來的以色列的拯救者.
這個人將來帶領自己的族民過埃及和過紅海.
來進入那個應許之地的摩西.
所以這個神聖的狐狸都很重要.
我們要留心去看 神冥冥中有掌管.

$^{321}$譬如我們都很熟悉的就是.
當我們看耶穌就快要進入那個受苦期.
最後晚餐 最後那個聖周的時候.
那隻驢子 耶穌都知道哪一隻驢子.
就叫那些門徒去到某處地方.
就拉那隻驢子來就行了.
這些都是神的護理.
很多東西好像表面上真的很巧合.
但是原來是在神的掌管裡面的.
所以我們就看到 法老的女兒居然扮演著一個.
在一個敵人的陣營裡面的 不是女嬰外合.
而是神可以在法老的女兒的內心.
種下這樣的一個感動,一個憐憫.
使她成為神可以用的 這裡都有提到.
就是神手中可以顛覆法老政權的關鍵的那個棋子.
所以如果這樣看人生.
這樣看我們自己有時候每天的一些遭遇.
一些事情 或者我們回望我們過去.
曾經經歷過,走過的路.
我們就應該要有這個肅靈的觸覺.
去看神在我們生命.
每一個弟兄姊妹生命裡面的那個護理.
那就會很不同了.
因為我們知道一定每一樣東西.
有神自己的許可 才會納到我們身上.
一方面我們就不會那麼激氣.
我們也都寬容 我們也都能夠面對的時候.
多一份謙卑去依靠神.
這個就是聖經說的神學家.
所給我們看到的這個護理.
表面上很簡單 但實際上其實.
這種巧合一點都不容易發生.
既然神要興起 神的時間到.
我們第一課已經說了 神跟亞伯蘭說了.
你要的確知道 你們的後裔要寄居那個地方.
那個地方的人就要來逼迫你們了.
第四代你們要回來.
在迦南地得到這個應許的地圖.
為什麼要這麼久呢?.
因為那個迦南地的人.

$^{361}$以這個阿摩利人做一個單位.
其實整個迦南地 他們的罪惡還沒有滿盈.
就是說現在 神的時間到了.
神的救贖的時間到了.
神的時間既然到了.
他用什麼都可以成為他手中的一個器皿.
所以法老的女兒都可以是神手中的器皿.
我都很深刻 我初信的時候.
聽一些長輩說 那個時候是挺嚴重的.
就是迫不及待的基督徒在內地.
但是居然在唐山的大地震七幾年.
我們信主肯定是很後期的 我自己.
聽這些長輩引述內地的一些基督徒說.
在唐山的大地震之前.
有一些監獄是困住了基督徒的 內地的.
不知道是什麼原因.
要將一些基督徒連夜遷去第二個城市的監獄.
這樣就避開那一劫.
是聽內地的一些信徒這樣說出來的.
你可以想像得到 原來神都可以管理很多很多的事情.
絕對不是因為有什麼自然的現象.
因為當時根本沒有人知道那個大地震是這麼嚴重的.
你就可以想像得到神的護理.
我們要記得基督徒永遠要知道神是會包底的.
神真的會使一些事情萬事互相效力的.
母親和姐姐的智慧在絕境裡面創造這個生機.
的確是要有她的勇氣 有媽媽 真是盤死的了.
因為那個是男嬰 男嬰不是有兩萬元獎金.
現在香港就生兒子生女兒都有了.
但是那個時候生兒子就要抓那個兒子浸死的.
生女兒就沒得獎 其實都應該慢慢就同化了.
但是他們真是有一個很大的勇敢.
所以基督徒都要訓練我們的勇敢.
我們要為神的真理 記得勇敢.
和我們亂來這樣無事的這樣罵人是兩回事.
這個智慧是他們能夠去到那一刻.
好像一張王牌那樣拿出來用.
因為關乎到她的弟弟 她的兒子.
這個性命的安危 他們肯定不知道.
自己的兒子將來是神用的 神揀選的.

$^{401}$第三章就已經知道了 但是他們是沒有水晶球的.
是我們今天 我們後世當然知道了.
我們看著事情發生 但是米利安和她的媽媽.
是完全不知道自己原來已經被神揀選了的家庭.
可以將來做這麼大的事工.
成為在舊約世界僅次於亞伯拉罕的摩西.
多麼重要的一個人 而且能夠上到西奈山上面.
面對面見到神 摩西就是這樣的一位.
世上少有的摩西為人謙和 勝過世上的眾人.
這個民訴記 我忘記了是第12章的一段經文.
就告訴我們神對摩西的描述 那個讚許.
這個就是告訴我們 媽媽和姐姐.
完全不知道自己已經被神選上的.
只是為了那個時代保存自己的男嬰的生命.
其實即是說有其它的男嬰都已經死了.
摩西如果她是公元前2700年 大概是這樣.
你可以想像到 在她的人口登記.
和她那年出生的那些很多已經死了.
所以一點都不戲劇性.
其實可能真的只有她一個被拯救出來的.
當然她的哥哥亞倫大概不止一歲還是兩歲.
還是三歲 好像是保存了.
聖經沒有告訴我們怎樣保存.
只是記載給我們知道 摩西的保存.
因為整個第12章 其實整個《春海舊記》.
那個主角都是摩西 所以聖經只是好像.
給我們一個示範 摩西就是這樣被神奇蹟地拯救.
其它的男嬰有沒有這樣拯救呢?.
應該是有的 但是用什麼方法?.
是不是又是法老的屢言 聖經完全沒有交代.
但是又給我們看到那個情況.
有可能就是之前沒有那個革命.
之前是低調地叫希伯來人的修身婆.
見到男嬰就掐死她 擇地說她是男產.
但是那個時候還是一個不是公開的一個命令.
可能是暗中的 所以過了那個歲數.
逃避了那個災難 大了 大定了.
法老都不敢太過的明目張膽.
所以摩西的哥哥亞倫就得以避開這一劫.
這個是我們可以去參賞的.

$^{441}$因為既然現在是有對男嬰有這樣的禁令.
那就是說一視同仁 所有的希伯來的男嬰.
都是必須面對這個滅絕的災難的.
但是摩西就是其中一個神將她拯救.
一個蒲草商 公主 剛才我們提過了.
公主的聯盟 敵對的陣營裡面.
她現在是一個有良知的那個殘忍的心的公主.
法老的女兒 所以原來我們每一個人.
在神的計劃裡面 我們都要扮演著一個角色.
神安排我們人生總是有那些關鍵的拼圖.
所以我們真的學習怎樣去將自己放在神的手中.
我想學習馴服 學習我們.
而且要用我們的腦 用我們的分析力.
用我們當前的環境去看我們下一步應該怎樣走.
我們應該怎樣做 但是我想總的來說.
我們的心是要為神的國度.
要為神的榮耀來看自己的人生.
這樣就會成為了一個在神不知不覺.
真的要很多年後的回望才看到那個拼圖的.
大家找過這些拼圖 沒了一塊就不完整了.
但是如果那塊都能夠真的放上去.
那拼圖就拼好了.
所以我們人生都是學拼圖的其中一粒.
神會將我們拼好.
可能大家都熟悉這些經文.
《彼得前書》的二章真的是講這個拼圖的.
《彼得前書》是這樣講的.
彼得的名字就是Petra 就是小的石子.
第二章的第一節 它說所以.
這個是《彼得前書》.
你們既然間除去一切的惡毒 詭詐.
並假善疾度和一切毀謗的話.
就要愛慕那純正的靈奶.
像財生的嬰孩愛慕那一樣.
叫你們因此漸長以致得救.
這樣做一個對比.
就是說純正的靈奶就沒有雜質.
但是原來影響我們吸收的.
就是我們裡面那個胃口.
如果我們裡面沒有胃口.

$^{481}$那些純正的靈奶就進不了我們的心靈.
所以原來真是有一個相輔相成的.
如果我們經常都想那些傷筋害理的事.
惡毒的言語,詭詐.
全部是和那個思想言語有關的.
假的善 即是說很虛假的 假到連自己都騙了.
疾度 又是內心裡面的一個扭爛.
看不順眼的別人.
和一切毀謗的說話 就是慣性了.
見高就拜 見低就踩.
總之就是對人的一個毀謗.
那就不能夠愛慕.
比如我昨天 星期四我們有周間的福音團體.
一般我們完了就會和老人家過去對面.
吃飯,聊天,喝茶 老人家很開心的.
昨天可能真的被迫38度去拿個車載老人家.
因為他們走很遠 走不了的.
載他們過去 其實很近的 教會對面.
但是一上車 就好像桑拿房一樣.
38度 然後開車 慢慢把老人家載上車.
開盡冷氣 過到去 一個燈位就過到去.
進到去 馬上開始涼了.
坐下之後 在酒樓 咦 有點不妥.
為什麼開始涼的手不停的冒汗.
那個人很累 還有完全沒有胃口.
可能真的是開始感冒了.
什麼胃口都不想吃 什麼都不想吃.
只吃了半個叉燒包就走了 要早走一點.
原來我們基本上.
只是我們自己本身的身體那個脾胃.
如果有一點點的不舒服.
我們都會抗拒那個食慾.
所以彼得告訴我們 原來真的不是開玩笑的.
原來很多人能夠成長 不愛慕38度.
其實不是什麼 不是整個度不好.
是人內心裡面有很多不完全的雜質.
不肯去改變 上文已經說了.
這個度就是永與常存的這個福音.
就是凡有血氣都過去 大主的度是永存的.
所得到的福音就是這個度.

$^{521}$所以其實神已經在我們生命裡面做了.
救恩最難的那個部分.
所以我們要栽種 原來就是第一節.
我們要意識到我們自己是一個這樣的人.
然後讓我們神的度.
在我們的生命慢慢地改變我們自己.
不過留意下面更加重要 更加精彩.
就是那個拼圖 第三節.
「你們如果嘗過主恩的滋味」.
大家都知道這是詩篇 應該是三十二篇.
「就必如此」 就是上文所說.
「愛慕凌乃」 好像一個財生的 屬靈的嬰兒.
大量地吸神的度 很想去認識神.
第四節 所以有時候很不開心的.
因為有時候見到頂尖的姐妹.
我絕對沒有任何的貶義.
但是你看她 因為看一次就可能誤會.
看兩次可能又是你主觀.
看三次 不是了 一年 兩年 三年.
她都是這樣的表現.
其實就是剛剛第一,第二節那個反思.
那個反照 就是她根本愛慕神的凌乃.
來到神的家 來到神的殿 來到禮拜堂.
就是要看慕神 敬拜神.
祂的真理在我們心靈裡面去養育我們.
來到的時候 根本不是太在意你說什麼.
根本不會有那個心靈的動盪.
已經看她是一次,兩次 不是一個誤會.
是她一個的形態來的.
來到教會 拿起手機.
還不怕被人看到她在看其他的東西.
在看其他的東西 在說道理.
你想一下 這樣的弟兄姊妹.
其實她自己在浪費自己屬靈的生命的成長.
下面告訴我們 如果有主的真理在我們內心.
這裡說 主乃活成Living Stone.
這個開始是一個比喻.
固然是被人所棄.
就是主耶穌在十字架 在各個灘.
來犧牲為我們而死.

$^{561}$也是應驗了以塞翁亞書的28章.
祂是堅固的石頭 是那個防角的石頭.
卻是被神所揀選 所寶貴.
第五節 留意第五節.
你們來到主面前.
就是我們這班信主的弟兄姊妹.
也就像活石頭.
接著就被建造成為靈宮.
作聖潔的祭司.
藉著耶穌基督奉獻神所閱立的靈祭.
原來這個靈宮這個Spiritual House.
在永恆裡面的將來的工程的圓滿.
原來就是我們每一個弟兄姊妹.
我們就是其中的一塊活石.
我們都是在主的裡面.
我們也就像活石.
主先就是那個真正的Living Stone.
但是我們都是有這個生命力的一個.
其實應該就是說 你可以說是關鍵的拼圖.
也是彼得所說的活石.
我們在這個活石的時候.
我們就是將來被建造成為靈宮的那個材料.
所以我們今天我們就要記住.
我們在這個靈宮.
講明這個是神聖的House.
一定是關乎獻祭.
所以有我們聖結的祭司.
我們是藉著耶穌基督奉獻給神.
神閱立的屬靈的祭.
相信就是我們的敬拜.
對眾物主的認識.
對祂的榮耀的尊崇.
所以主禱文真的是很重要的.
將那個次序.
最重要的次序就是.
願人都尊你的名.
眾物主的聖名為聖.
你的光度降臨.
你的旨意全部和我們人類無關的.
第四個才是我們日常的飲食.

$^{601}$赦免我們的罪.
接著就是保守我們不要遇見魔鬼,蝕客.
你可以想像得到.
我們現在在這個靈宮裡面.
原來我們有沒有意識到.
我們是其中一塊的活石呢?.
我們有這種作用.
來成為靈宮的建造材料.
在這個建造過程.
一個活石給了我們血道.
我們要靠著主耶穌.
不斷奉獻神閱立的靈陣.
例如將人帶進那個敬拜裡面.
將人不認識救恩的.
慢慢使他們明白救恩.
一起同心事奉.
在主裡面相愛.
幫助有需要的人.
甚至建立神給我們一些人生.
神要我們建立的事業.
就是為神的工作去做.
神的過渡的葡萄園.
是有很多的工作的.
不同的工種的.
我們有沒有找到我們自己的人生的方向.
那個目標呢?.
所以這裡就是告訴我們.
原來我們是一個活石.
原來我們是一個拼圖.
原來神會將我們擺在.
祂的永恆的那個計劃裡面的.
所以我們記得.
原來我們都是一個奉獻者.
我們來奉獻給神.
反正我們不再看了.
因為否則會拆得太遠了.
記得我們是一個拼圖.
告訴我們波西從嬰孩到成為拯救者.
頭十節就告訴我們.
這個不是一個孤立的溫馨的童話故事.

$^{641}$而是作者精心埋下的一個顛覆性的伏筆.
因為在第十一節.
摩西就立即意識到自己的人生.
那個巨大的落差.
他的身份令他感受到.
從來都沒有的那種困惑.
因為唯有他在埃及的王宮裡面.
成為一個享盡榮華.
和學盡埃及的文化.
我們揭一揭他的《士徒行傳》第七章.
看看新約的士徒們.
他們怎樣引述舊約的歷史.
第七章的第一節.
司提反當眾的申訴.
大祭司就說.
這些事果然有蚊.
因為司提反被人抓去長老那裡.
因為他在民間傳道.
這是第八節.
司提反滿德恩惠能力.
在民間行了大騎士和神蹟.
有一些人就對他開始有不滿了.
第十節.
司提反是以智慧和聖靈說話.
眾人抵擋不住.
就是說又有point又有真理.
那個argument又不是為反而反.
更加說話具有那種智慧.
那種謙卑.
那種不是咄咄逼人.
而是很有那種對別人的體諒包容.
又將那些事實陳明出來.
真是見到神藉著他使很多人能夠認識耶穌基督.
接著就擺出人來說.
這個人是旁讀窩西和神的說話.
接著又找一些人暗中買通一些人.
聳到很多百姓長老和民事.
就忽然來捉拿他.
將他帶到公會那裡設下假見證.
這個人說話不住的糟踐聖所和律法.

$^{681}$我們曾聽見他說.
這那撒旦人耶穌要毀壞此地.
這些相當嚴重的指控.
也要改變我們摩西所交託我們的律法規條.
你想一想在那個時代保羅都還沒有出現.
那使徒們怎樣去拆解這些這麼大的.
在猶太教裡面的一個很嚴重的一個分水嶺.
到底現在這些基督徒跟隨耶穌的.
是不是要將整個猶太教拆毀呢?.
似乎是這樣 表面的證供是成立的.
所以必然有這樣的衝突是會發生的.
不過首當其衝的這麼白熱化的.
就是斯大林這位的執事.
所以第15節 在公會裡坐著的人都定睛看他.
見他的面貌 保障天使的面貌.
後來開始說話和最後臨死之前.
都是這樣被人見到他.
作者都很有興趣描述這位衛主在眾人.
就是那時大家暴亂的暴徒圍繞的地方.
斯大林還在氣定神閒地站起來說道.
其實當時你應該很害怕 或者很生氣.
或者你是很情緒化.
很被那個氣氛激動.
你都不懂怎麼說話.
但是你看到這個就是一個很平靜的見證人.
見到他好像天使的面貌.
天使的面貌 你可以想象得到.
這樣的描述,這樣的形容.
但是第七章 斯大林其實是將.
我們都很熟悉的《創世記》.
來跟我們很簡單的濃縮了 都說了一次.
所以雖然有天使的面貌.
但是仍然是很有邏輯的 很有意思的.
還有說話是很合乎舊約的.
因為現在那個指控就是.
你和你們的跟隨者是要糟踐摩西的律法.
甚至說叫人不要再跟摩西的規條去做人.
整個猶太教就是這樣就會崩潰了.
誰可以承擔這個猶太教崩潰的一個這樣的罪名呢?.
斯大林 你想一下 在這樣的情況下.

$^{721}$他怎麼可能來作見證呢?.
大家想一下 現在我們不是很多公開的辯論.
但是你試一下在無論是一些社群的媒體.
例如Facebook或者小紅書,IG.
你稍微說一些不中聽的東西.
那些評論就對你那種人身攻擊.
那種完全沒有理智的謾罵起底.
是非常恐怖的網絡欺凌.
但是你想一下 當時斯大林面對的.
是百姓,長老,文士和假見證的人.
來對他的攻擊 而且那個攻擊.
是直接最致命的針對你這幫猶太人.
現在你相信了那個邪教的耶穌.
現在你叫我們猶太教不要守摩西的律法.
其實相當嚴重的 第七章.
大祭司就說 這些事果然有妄.
先是反說 諸位父兄請聽.
當日我們的祖宗亞伯拉罕在米索波大米.
就是當時的那個吾爾還未住哈蘭的時候.
榮耀的神向他顯現 對他說.
你要離開本地和親族.
往我所要指示你的地方去.
他就離開迦勒底人之地 住在哈蘭.
他父親死了以後 神從哪裡搬到你們現在.
就是以色列 以色列地現在所住之地.
在這個地方 神並沒有給他產業.
他就是亞伯拉罕 連蠟燭之地也沒有給他.
但應許要將這地賜給他和他的後裔為業.
就是《創世記》十五章 我們第一堂所說的.
第四代他們就可以回來了.
那時他還會有兒子 第六節 神說.
他的後裔必寄居外邦.
那裡的人要叫他們作奴僕 輔待他們四百年.
這個就是十五章的經文 《創世記》.
神又說 使他們作奴僕的那國 就是埃及.
我要懲罰 以後他們要出來 在這地方事奉我.
接著第八節 神又賜他割荔的藥.
於是亞伯拉罕失了耳實.
割荔就是《創世記》十七章.
所以神和亞伯拉罕納藥的記號.

$^{761}$記得不要搞錯 是割荔.
神和以色列人納藥的記號.
是安息人 所以不要搞錯了.
第八日 甘他行了割荔.
以殺身那國 哪國生十二位先祖.
很快用一兩句就將一些無關重要.
這次的argument的事情就省略.
斯蒂法要將那個重點帶進.
怎麼從舊約帶進去耶穌那裡.
因為重點就是 如果我們跟隨耶穌.
原來有些事情真的改變了.
不是再守一些教條了.
不是再像猶太教那樣死守教條就可以了.
因為真正的應許就是耶穌這位神的言辭了.
所以我們都留意一下斯蒂法.
這些重要的領袖是怎麼來到去.
其實當時根本沒有機會給他組織安靜.
我問一下AI應該怎麼組織.
完全沒有空間可以這樣做的.
就是斯蒂法平時都已經是不斷對舊約相當的純熟.
也都是有聖靈的感動 聖靈的同在.
有天使的那個平安 那個面孔.
這樣來到去 近乎急才.
來到去馬上要講一篇道.
你可以想像得到 聖靈的感動可以帶給一個人.
超出他自己所能夠的那種力量.
第九節 先祖日渡約瑟 把他埋到埃及去.
神卿與他同在 救他脫離一切苦難.
由使他在埃及王法老面前得恩典 有智慧.
法老就派他作埃及國的宰相 兼管全家.
這些經文全部是我們很熟悉的創世記的經文.
第十一節 後來埃及和迦南全地遭遇饑荒.
大受艱難 我們的祖宗就絕了糧.
哪個聽見在埃及有糧 就打發我們的祖宗.
初次往那裡去 第二次約瑟與弟兄們相認.
他的親族也被法老知道了.
第十四節 約瑟就打發弟兄請父親那國.
和全家七十五個人都來.
這裡有少許春油 差那五個.
這些是聖經的一些難題.

$^{801}$到底是怎樣計算的呢?.
是否計算在埃及生的約瑟自己兩個兒子呢?.
這些我們可以看看的.
第十五節 於是那國下了埃及.
後來他和我們的祖宗都死在那裡.
又被帶到士劍 葬於亞伯拉罕在士劍用銀紙.
從哈末子孫買來的墳墓裡.
就是說他的爸爸這個那國.
他下了埃及之後有十七年.
那裡就有他十二個兒子供養他.
讓他有一個安享晚年.
這個是雅各漂泊的人生.
終於他還是要在埃及那裡死.
不過他就首先離開埃及.
上回去葬在士劍.
葬在亞伯拉罕和他爸爸以撒的墳地.
又是等候那個應許的來到.
第十七節 及至神應許亞伯拉罕的日期將到.
以色列民在埃及興盛眾多.
直到有不曉得若生的新王興起.
就是我們今天所讀的.
他用詭計待我們的宗族.
首先就是逼他做苦工.
慢慢想淘汰那些男嬰.
就是越來越不能生 每次生都是生女.
這樣都不足為奇.
某些時代 女多於男的都是很合理的.
因為沒得選擇的 你一生下來.
不是說你用什麼方法.
什麼時候 初一,十五就會生女多一點的.
沒人知道的 這個是神秘的.
你可以用超聲波知道是男還是女.
但是要生男還是女就沒得選擇的.
不是到我們去選擇的.
但是法老說 不是啊 我是皇帝啊.
我是神啊 我可以選擇啊.
不過是用一種非法的手段.
用一種強迫的手段.
令到這個民族可能一代,兩代已經沒有人.
來成為它的威脅了 全部是女的.

$^{841}$就全部嫁給埃及人 使它完全的同化.
這些是種族的滅絕.
所以就是詭計 苦害我們的祖宗.
叫他們丟棄嬰孩 使嬰孩不能存活.
第二十節 那時 摩西生下來盡美非凡.
在他父親家裡撫養了三個月.
他被丟棄的時候.
法老的女兒拾了去養為自己的兒子.
第二十節 就是這裡了.
摩西學了埃及人一切的學問.
說話行事都有才能.
這個就是聖經給我們知道的.
新約的作者都給我們看到的這樣的情況.
或者我們可以看希伯來書11章的經文.
都是很短的.
可能要看後一點.
或者第21節 11章的第21節.
就稱為信心篇的不同時代的跟隨者.
他們在那個時代是怎樣用信心跟隨神.
所以某程度上不應該叫信心偉人.
是在那個時代是用什麼形式去跟隨神.
有些人 但是以你這樣 就要在獅子坑.
不是每個人都要在獅子坑的.
有些人可能一生都很順利的.
第21節.
那國恩爵信臨死的時候.
給約瑟的兩個兒子各自祝福.
扶著像頭敬拜神.
約瑟恩爵信臨終的時候.
提到以色列族將來要出埃及.
並為自己的骸骨留下遺命.
第一章我們已經說了.
《措死記》第50章就是那個棺材.
擺在埃及 就成為那個結束.
其實就是非常有詩意地告訴我們.
這個就是有信心的人的等候應許.
所以記得 我們無論好死不好死.
所謂好死就是 突然間心臟去打 馬上死.
或者搞你幾年 那些癌症痛得不得了.
又要用化療 又要電療等等 又要做手術.

$^{881}$煎熬得不得了.
這樣的人都大有人在 基督徒.
無論如何 我們要學聖經告訴我們.
這些人都是憑著信心離開世界.
所以約瑟當然是好 因為他真是壽終正寢.
他的兒孫都一直供奉他 直到他離世.
但是真正的信心就是.
他為自己的骸骨留遺命.
你們將來出埃及 一定要帶著我的骸骨上去.
不要趕忙離開埃及 因為要逃跑.
你要帶著我的骸骨 所以約瑟的骸骨.
應該放在一個很重要的 應該放在回望也不奇怪.
或者放在柯西的房間 因為隨時一走就拿了.
等於總統的核武器的行李箱.
有專人和總統一起 去哪裡都要帶著那個箱子.
我們一定要記住.
我們總有一天會離開我們這個物質的世代.
但是神的國度還沒有來.
我們都是憑著應許這樣離開.
所以我們要記住 我們要學習.
盡我們的人生來為主而活.
完成了過去了 我們不需要留戀.
也不需要懼怕 也不需要覺得很唏噓.
因為我們是需要被人忘記.
因為神才得到最大的榮耀.
所以留遺命的那個骸骨就象徵他們的信心.
原來希伯來孫都是這樣看的.
他們都是這樣看50章《創世記》那個棺材的.
原來真的是這樣 不過林軒第23節.
摩西森下來 他的父母見他是個俊美的孩子.
就恩著信 說他的父母 說到這個約經別.
應該我剛才記得了 就是暗蘭 他的爸爸.
原來他們是恩著信.
所以沒有將他交出去給那些兵丁.
將他藏了三個月 並不怕亡命 果然是.
就是說逐家逐戶都會來搜查.
並不怕亡命 但是又可以講一個故事.
在二戰的時候 應該叫日治 香港.
我以前補課的一位德高望重的前輩.
那個時候他剛剛出來傳道.

$^{921}$你想一下 1939年一直去到1945年才投降.
3008個月 香港淪陷 在深水Po 那些.
因為那個時候很艱難 那個時候那些房子.
全部是那些截前舊房子 是木屋來的 木樓.
有些樓梯的時候 他們會拿木頭來燒 燒火.
因為很艱難 日軍統治.
日軍逐家逐戶去搜有沒有那些革命分子.
這位傳道人 當時也是帶領著我們.
在軍塘的長者 當時他們是一個很年輕的階段.
可能是30多歲 出來傳道 在深水Po 一帶.
有一次日軍真的去抓人 逐家逐戶去抓人.
他藏起了一些姐妹在一些房間裡面.
因為那些日軍就是要抓那些長者.
很奇妙的 他沒得躲了.
逐家逐戶去搜 開門 打開那些房間 那些箱.
居然其中一個日軍 他親眼真是經歷了.
他親自經歷這件事 他說那個日軍.
當然空神落實 真是拿著那些長矛的刀.
那些槍來抓人 但是打開了那個.
因為他們藏起了他 打開那個房間等等.
居然打開之後立刻離開 就沒有再看.
接著就走了 那個姐妹就倖免於難.
我忘記了是師母太太 還是其他的姐妹.
你可以想象得到 原來神有時候都可以.
這樣來得這麼戲劇性.
讓這些事情會在我們面前發生.
所以不到你不用信心 去到這些非常時期.
去倚靠神 因為你已經沒得再藏起來了.
所以殷卓相將他藏了三個月 並不怕和平.
到沒得收藏了 真是會大聲喊 會被揭發.
可能連父母都會受到威脅 所以第24節.
摩西殷卓相長大了之後.
就不肯稱為法老女兒之子.
他寧可和神的百姓同受苦害.
也不願暫時享受罪終之樂.
他捍衛基督受的靈欲.
當然他不知道什麼是為基督受的靈欲.
因為當時他不認識基督.
希伯來書的作者告訴我們.
他也是和我們一樣 這樣信靠神的人.

$^{961}$他願意為基督受的靈欲.
比埃及的財物更加寶貴.
因他賞望所要得的賞賜.
他因著信離開埃及 不怕黃奴.
因為他含深忍耐 如圖看見.
那不能看見的主.
這裡我們先不再說了.
因為涉及到這麼長的經文.
所以今天我們重點要看的.
其實就是第二章的第十一節.
剛才引述的新約的經文.
或者說新約的士徒.
他們是看舊約.
就是看到和我們今天說的信心的倚靠.
完全是同出一裂的.
所以都是同樣有聖靈的感動.
這些新約的作者告訴我們.
第十一節就是很平鋪直敘地.
記載摩西長大之後.
我們可以有很多的空間.
摩西當他讓法老的女兒進入皇宮.
來撫養 不再是食物的層面.
開始要豐富他作為將來繼承人的那個分娩.
真的不簡單 法老都破格.
佔領希伯來的摩西.
因為是法老的女兒的兒子.
法老知道他是希伯來人.
但是都容許他在皇宮出入.
收養他 法老都願意佔領.
或者從一個角度就是.
好 我用我們最好的埃及的文化.
埃及的宗教 埃及的學術去分屠他.
洗他的腦 慢慢他是希伯來人.
但是思考是埃及人的思想.
就是說希伯來人是沒有用的.
是霸佔我們的地圖的.
是會成為我們的敵人的.
如果由摩西講出來很重要.
因為他成為一個傀儡 他成為一個漢奸.
這樣再和希伯來人這樣去說.

$^{1001}$去鎮壓希伯來人 就非常有效了.
或者勸導希伯來人要歸順埃及 又可以.
所以很合理的 法老都容許摩西.
在皇宮裡面學宮廷的禮儀.
不只是學禮儀的.
是真的關乎到宗教,學問等等的.
所以為什麼達爾尼和三個朋友.
薩德拉米薩,阿曼尼哥.
為什麼他們一定不可以學習宮廷的禮儀.
因為是在講和他們的宗教都掛鉤的.
因為一定是要逼他們接受他們埃及的宗教的.
所以摩西他一直學習 告訴我們.
這裡沒有記載 他全都學了 學盡了.
埃及人一切的文化,說話,行事都是有智慧的.
這裡只是說他已經長大了.
但是原來他的內心.
都沒有被埃及的文化所熏陶,所洗腦,所灌輸.
或者那個厲害 因為他已經是將來的繼承人了.
這樣的一個身份 你可以想像得到.
那個代價是非常大的.
我如果不要這個身份 我看自己是希伯來人.
我回去和他們一起 你是一個奴隸.
還是一個王子呢?.
是王子和奴隸的那個對比.
所以頂錢的人 原來神用的人.
他是清楚自己為什麼要這樣走人生的路.
所以這個世界有一個行業.
當然因為我自己也是牧者 特別要敬重.
因為這個世界有一個行業.
這個世界是不需要的 不稀罕這樣的人.
但是這個世界有一個行業 就是有些人傻乎乎的.
他說我自己要終身服侍神 跟隨神.
我要入神學院 我放下我的工作.
那些父母簡直把兒子當成瘋了.
入神學院讀四年 什麼都沒有做 加學費.
然後就是這樣來到將來在教會裡面服侍.
這個世界有一個這樣的行業.
但是世人就覺得你有什麼生產力呢?.
你有什麼用呢?.
但是在神的眼中 原來摩西就是代表著.

$^{1041}$一個清楚自己的身份那個呼召.
所以清楚呼召和我們要認識自己的夢召.
決心走上這條事奉的路.
這樣服侍主的人 頂錢的人一定要敬重的.
因為神的僕人就是這樣要來將他自己的一切放下.
為了神的旨意 所以我們要求主讓我們能夠看到.
十一節這裡很重要 因為他出去.
對不起我有點乾 我先喝口水.
好 繼續 大家有看到這個圖.
《Good Will Hunting》 有沒有看過這套劇集.
相當好看的 說到這個Robin Williams.
羅賓威廉斯 右邊那個 就是一個心理學家.
左邊這個就是麥迪文 拍這套電影真人真事.
在波士頓有一個年輕人.
原來自小被父親嚴重虐待.
所以他長大之後非常反悔.
覺得這個世界欠他很多.
他就是一個很聰明很厲害的數學天才.
就變成一個很兩極化的人 在他的內心.
所以人格分裂 他年輕的時候.
就跟那些豬朋狗友青春到底.
但是那些豬朋狗友本來就是學識很低.
低下階層 做一些很粗重的勞動的工作.
但是他都是做這些工作.
他就進了麻省理工大學裡面做清潔工.
在大學裡面有些很難的數學的問題.
每年那些教授都會寫出來.
如果能夠解答到.
就會得到差不多是數學界的諾貝爾獎.
這個故事的主角麥迪文.
就是這樣用了一段很長的時間.
就解釋了這個數學的難題.
簡直是越來越想像.
因為連他教授自己都不懂得解釋.
就是一個這樣的年輕人.
但是他是一個很不羈,很放縱.
和很憤世嫉俗的一個年輕人.
所以就處在一個這樣的狀況.
那個教授發現過去這麼多年都沒有一個數學天才.
現在在他面前發現了.

$^{1081}$所以那個教授就很想去培養這個麥迪文.
這個Hunting 他的名字就叫做Hunting.
就希望幫助他 慢慢帶他出來.
他自己那個困惑的人生.
因為他一輩子都是被爸爸的後父.
用鐵棒來打.
小時候有虐待等等.
所以他長大之後就到處去撩交打.
所以有很多案例的.
但是就遇到這個心理學教授.
就是慢慢讓他認識自己是誰.
這套戲很好看的.
就是他終於後來接受了這個教授的心靈輔導.
慢慢他願意開放自己.
慢慢他願意接納自己的不完美.
慢慢願意勇於面對自己的目前的追求等等.
就是他認識了一個大學生.
後來就展開了愛情等等.
整個故事就是這樣.
對不起 劇透了.
其實就是這個 發現自己的身份.
我想我們每個人都要發現自己的身份.
姜濤就更加 我們有兩個身份.
我們原生家庭 到我們成為了主的兒女.
我們要保衛自己的身份.
我相信這裡多少人都是這樣被感動的.
當真是知道自己原來是被神愛.
無條件被愛.
這樣的人在世界找到一份這麼真的愛.
他的生命就會有一個扭轉.
生命開始會有一個很大的動力.
想變好.
所以很渴望聽到這個錄音的朋友們.
因為我肯定不會知道你們的背景,你們的成長.
但是如果你的成長是很艱難的.
摩西都艱難 他同歲的可能已經淹死了.
不知有多少.
其實很坎坷,很悲哀的那個悲哀時代.
但是神就要在那個悲哀時代裡面.
就是神的時間到了 就要興起摩西.

$^{1121}$所以表面上不是他刀槍不入.
而是神種種的巧合就令到他能夠生存.
但是現在說的是長大.
十一節是說他長大 那個衝擊還更加大.
不是每次都和你為敵,虐待你.
如果你是從小到大都是奴隸.
沒辦法 你都是這樣沒辦法.
都是要找磚做磚 都是一個這樣的狀況.
但是不是 唯獨摩西在皇宮裡面.
來到他被教養 成為將來的繼任的皇子.
這樣的衝擊就更加大.
所以Who am I 如果是我們.
不斷有意識地去問自己.
這個是很重要,很寶貴的.
我相信弟兄姊妹對於我們自己的出生.
我們自己的成長 這個應該不難去處理.
但是再問深一層就是.
神給我們過去有這樣的成長.
有什麼作用呢?.
有什麼可以令到我今天可以更加被神用呢?.
這個我們應該再深一步去問自己這個課題.
所以她就真是表達出她是希伯來人.
她做的比較魯莽的決定.
她的第十一節 後來摩西長大.
她出去到她弟兄那裡 看他們的床單.
見一個埃及人打希伯來人的一個弟兄.
就把埃及人打死了 藏在沙土裡.
第二天她出去 見有兩個希伯來人爭鬥.
就對那希夫人的說.
你為什麼打你同族的人呢?.
那人說誰納你作我們的首領和審判官呢?.
這句話都挺刺激到摩西.
她真是有一個很單純的心.
很想幫自己的民族.
她似乎已經有一定的權了.
因為成為了埃及的王子.
也都有一定的影響力.
但是現在這樣的一個.
這麼普通的一個遭遇.
就激發她殺了這個埃及人.

$^{1161}$但是兩個希伯來人的爭鬥.
對她的說話的一個這樣的言語的侮辱.
就令到摩西非常氣餒.
感覺很難受 因為她的心真是很想.
她不是想做領袖.
是很想站在希伯來人一邊.
來認同他們的仇恙兵.
認同他們現在的苦楚.
可能大家從長計議.
怎樣能夠叫法老明白.
真是諒解希伯來人從來沒有一個異心.
從來沒有一個要反對,要背叛埃及的.
可能她就扮演著這個橋樑的角色.
就來到兩邊可以做一個和平的使者.
但是這樣的狀況.
就已經令到她非常氣餒和懼怕.
摩西便懼怕說 這是必是被人知道了.
法老聽見這事就想殺摩西.
但摩西躲避法老 毫無畏懼地居住.
這個就是摩西首先自己的出埃及.
這是一個身份的危機.
摩西在這一刻長大.
就是說希伯來人媽媽的教育有多深影響呢?.
還有她一路成長 她有多少機會.
以前有一套電視 王逸華做的.
說她是一個很像流氓的.
張靜怡和吳靜松是哪個主角呢?.
一些貴族住在半山區 她就是一個流氓.
輾轉她就進了一間名校 寄宿學校.
開始他們學校的生活 後來長大等等.
大概摩西就好像一個貴族的生活.
但她內心卻是在希伯來人的苦況裡面.
小時候看到 有時如果把心一橫.
見一次當然會很觸動 但當習慣了自己享受的.
你就會慢慢忘記那些窮苦人家的痛苦.
因為那個對比 如果不再在你面前出現.
慢慢就會消失 你就會在你自己的圈子生活.
因為你的圈子出入的就是馬會.
出入的就是銅鑼灣的名店.
出入的就是一些很漂亮的地方.

$^{1201}$去就是直去坐頭等的 你不用擠巴士.
不用去坐普通的機位.
你已經在另一個行業裡面 根本永遠在你的腦海裡面.
看不到有窮這件事的.
但摩西長大之後能夠這樣.
這有十幾年 聖經是沒有興趣告訴我們那些細節.
司提反就說 摩西學盡了埃及人一切的文化.
說話,言語都有智慧 但是沒有告訴我們.
到底他有多少受希伯來人的文化影響.
而導致他有這樣的心理不平衡.
看到埃及人打希伯來人.
就令到他有一個很憤恨的心 要殺了這個埃及人.
看到兩個自己的同胞在爭鬥.
為什麼又要這樣去勸架呢?.
這裡有很多空間給我們去想的.
一個人的內心 到底是從多少歲開始.
就奠定他自己有希伯來的血統.
他無論如何 他留住的就是希伯來人的血.
到底這個這麼渺小的基因.
是不是會影響 能不能影響他一直的成長呢?.
還是後天的那個培養 後天給他的那些資源.
那些享受 更加會淡化了他作為希伯來人呢?.
所以摩西的烏日奶是來得比我們更加的重大.
我們都很難說這麼大的落差.
我們被人執了去養 認識自己的原生.
很少這樣的 通常電影才有.
但是摩西就是這樣的一個狀況.
但是神就在他的生命裡面.
讓他裡面居然有一些是什麼都改變不了的.
那個基因留著的希伯來人 純正的.
亞伯拉团,以撒,雅各,約瑟留著的.
那個對神的信靠的 那一個這樣的熱血.
那個的信心 所以十一節是一個很重要的.
一個轉捩點的分水嶺.
作者巧妙地告訴我們 頭十節就說這個寶寶.
在眾多的幫助者的幫助下.
他才可以養得大.
但是到他自己要決定自己的人生.
其實那個誘惑真的相當大.
一邊是埃及王子 一切的名話.

$^{1241}$一邊是希伯來人 跟隨神的這幫.
現在是奴隸 不是貴為一些很尊貴的民族.
但是所有都是埃及人的奴隸,奴役的.
所以一個這麼大的落差.
就是波西應該怎麼去選呢?.
所以今晚就講到這裡.
希望大家如果對自己的屬靈的身份.
自己成為基督徒產生一個共鳴感的.
我們都能夠彼此勸勉.
或者你找一些屬靈的牧者去分享.
使你能夠走出你自己的困局.
能夠抓緊神在你心靈裡面的帶領.
來去服侍神 不要停留在.
不斷去施花輪福 這樣去陷入一個.
自己內心的膠戰裡面.
本來預備了另外一個故事.
就應該不夠時間去講了.
鄭丹瑞 不夠時間去看了.
其中他人生的一個開竅,一個轉捩點.
就是他曾經主持一個晚間的深宵的節目.
當然是收音機的節目 香港電台.
後來有聽眾寫了一封信告訴他.
你不喜歡做就不要做.
我半夜需要有音樂陪伴我做凡伴的工作.
這一封信就令到鄭丹瑞從此在播音界.
就是一個很深刻的當頭棒喝.
那封投訴信一直放在他桌面.
成為他自己的警聽.
原來有人關心我的工作.
因為他真的很狼狽 狼狽得很厲害.
因為當時他本來在做日間的節目.
黃金時間的 突然間又沒有任何通知之下.
將他調去做深宵做一些四五點鐘的.
真是哪有人在聽他寫的 做那些節目.
他就很激氣 播那些歌 九紅搭八的.
真是九紅搭八的 他想著根本不會有人聽的.
所以就是這麼狼狽 但是那封投訴信.
就成為了他一個很大的.
就是他自己的一個 Who am I?.
就是他終於發現自己的生命原來需要這樣改變.

$^{1281}$他就以這個成為他自己的警聽.
我們一起禱告 多謝主.
我們今晚都開始講摩西了.
他的人生有這麼多的坎坷.
當然當時的時代是極艱難.
因為埃及王法老是要滅絕所有的藍丁.
多謝主將摩西拯救 預備他的童年.
也都將他擺在 很奇妙地擺在埃及王的王宮.
來得到教育 而且這些埃及的教育.
沒有令到他變了他的內心.
仍然都是亞伯拉罕這個信心的行列.
成為了神的用人.
天父 這個時代需要被你揀選.
忠於真理 忠於福音的工人.
向夢主興起更多 尤其是年輕的一代.
雖然很艱難 現在都在不斷談論.
關於樓價 關於工作.
關於AI取代很多青年人的工作等等.
我們知道每個時代都有那個艱難.
所以我們更加要為下一代的年輕的.
無論是牧者 或者年輕的信徒去禱告.
求主興起他們的心靈.
來真的看到神 你那個永恆的護理.
是在每一個人的心上 在你的計劃裡面.
所以我們盼望 我們都能夠擺在.
神你的永恆計劃裡面.
成為一個關鍵的一個聘託 如主幫助.
我們這樣感恩禱告 奉耶穌基督的名 阿們.
\newpage



\section{出埃及記}
\label{sec:HMEU7PxsF0Y}
\textbf{2026-06-05 周五晚查經班《出埃及記研讀(四)神國領袖的再造工程》 講員:曾裔貴牧師}
\newline
\newline
連結: \href{https://youtube.com/watch?v=HMEU7PxsF0Y}{\texttt{https://youtube.com/watch?v=HMEU7PxsF0Y}} ~~~~ 語音日期: 2026-06-07
\newline
\newline
\hyperref[sec:4cUSeyCs18Q]{\small{< < < PREV SERMON < < <}}
~
\hyperref[sec:index_chronic]{\small{[返順時目]}}
~
\hyperref[sec:index_scriptual]{\small{[返順卷目]}}
~
\hyperref[sec:cWahrBO910c]{\small{> > > NEXT SERMON > > >}}
\newline
\newline
$^{1}$好 靚姐妹我們同心禱告.
再一次我們感恩.
讓我們又可以過了一個新的禮拜.
我們又能夠來到晚上的茶經班.
來繼續去研讀《出埃及記》的經文記載.
還有作者想要帶給我們看到的.
神的救贖計劃,那個企劃,那個安排.
我們看到每一件事 好像那個時間.
都是有一個準備.
我們真是很看到神自己對你的民族.
對你自己揀選的以色列人.
其實應該說對亞伯拉罕.
神你在遠古之前的應許從來都沒有落空.
就是這樣按照著時間來去兌現.
我們看到神你將摩西.
本來一個在當時是真的會死於非命的嬰孩.
你將祂拯救 而且不單只是一種簡單的拯救.
而是你要將祂的心靈都來去改換一生.
祂能夠成為神手中的器皿.
這個才反而是我們需要每一個時代.
我們有心服侍主的弟兄姊妹.
需要來去學習 求主幫助.
我們將我們今天晚上的時間仰望在你的手中.
我們這樣禱告奉耶穌基督的名 阿們.
很感恩 我們很快就來到第四講了.
今晚我們的主題 一個分題.
就是神國的領袖.
神在他生命裡面的一個再造的工程.
講到明是工程真是有step by step.
真是有他人生的步伐.
他的升高是因為他首先要被降下.
這個《雅各書》一章十二還是第….
我不記得第幾節 卑微的弟兄升高.
就該起落 富足的降卑也該如此.
這裡就讓我們看到 今天我們就透過這個工程.
讓我們去思想一下 事奉者的生命.
是怎樣經歷神的拆毀和重建.
對當事人來講 從來都不是一件容易的事.
我們又看看那些的際遇 那些當時的環境.
有沒有什麼環境的因素.

$^{41}$有沒有什麼他自己的一個生命需要去學習.
需要去扭轉.
我們上個星期講到二章第十節.
就是和十一節是一個很重要的橋樑.
用一節聖經告訴我們.
波西在一個非常受到保護,受到栽培的環境之下.
成為一個出類拔萃的一個當時的埃及王子.
有些歷史的記載.
就是說他根本就可能在這個埃及王死.
他就可以順利的因為他的養母.
這個埃及的公主.
他可以來到成為那個王子 接著繼位的.
你可以想像得到一個這麼大的國度.
對於約瑟 他也尚且是一個宰相.
來到去做一些管理 帶領那個國家.
和當時的世代不要因為大饑荒而來到去死亡慘重.
這個是神的計劃 但是更加寶貴的.
原來神有更加深遠的計劃.
第四代的子孫 亞伯拉罕的第四代的子孫.
就會拯救以色列人 離開這個圍爐之間.
所以十一節我們可以從這一段再讀起.
所以這一大段聖經 其實是去到第四章的經文.
所以應該有四個部份 就是摩西他的生命的成長.
或者說神將他塑造 由埃及的王子.
變成為神的僕人 也可以這麼說.
變成為以色列的一位僕人領袖.
我們讀一讀經文 第十一節 二章的十一節.
後來摩西長大 他出去到他弟兄那裡.
我忘記加那個麥克風.
看他們的重擔 見一個埃及人.
見沒有人 就把埃及人打死了.
藏在沙土裡 第二天他出去.
見有兩個希伯來人爭鬥 就對那希婦人的說.
你為什麼打你同族的人呢?.
那人說 誰讓你作我們的首領和審判官呢?.
難道你要殺我 像殺那埃及人嗎?.
摩西便懼怕 說 這是必是被人知道了.
法老聽見這事 就想殺摩西.
但摩西躲避法老 逃往米甸地居住.
可以說是他自己第一個 其實是逃難.

$^{81}$不是昂然的 挺胸昂首的來到出埃及.
不是浩浩蕩蕩帶著神的恩典.
帶著神的那個權柄 這樣離開埃及.
現在是一個逃犯 一個需要亡命的逃犯.
可以想象得到那種落魄 那種懼怕 徬徨.
第十六節 有一天 他在井旁坐下.
米甸的祭司有七個女兒 他們來打水.
打滿了槽 要喝父親的群藥.
有牧羊的人來 把他們趕走了.
摩西卻起來幫助他們 又喝了他們的群藥.
他們來到父親留意那裡.
他說 今天你們為何來得這麼快呢?.
就是說過去都是有些當地的土匪.
或者那些流氓 霸佔了那個水源.
不排隊 所以欺負那些老弱婦女等等.
摩西就幫他們出頭做駕輛.
第十九節 他們說 有一個埃及人.
救我們脫離牧羊人的手.
並且為我們打水 喝了群藥.
記住他們所認得出的 摩西是一個埃及人.
也就是說整個人都改變了.
可能外觀 可能談吐.
可能在宮廷裡面的成長氛圍.
就變成不是一個落魄的奴隸的希伯來人.
當然摩西未必用希伯來話跟他們說.
第二十節 他對女兒們說.
那個人在哪裡 你們為什麼撇下他呢?.
你們去請他來吃飯 摩西甘心和那人同住.
那人把他的女兒西波拉給摩西為妻.
西波拉生了一個兒子 摩西給他起名叫甲醇.
意思說 因我在外邦作了寄居的.
第一個兒子就是甲醇 就是在曠野那裡生的.
第23節 過了多年.
其實總共是40年 摩西在曠野的生活.
因為稍後去到第三章.
摩西已經是80歲的老人家了.
所以這裡簡單介紹完摩西的出埃及記之後.
他個人出埃及之後 落戶在米甸的祭司的家裡.
甘心樂意成為一個女婿.
好像被嫁入這個家庭.

$^{121}$對摩西來說 是一個怎樣的心情呢?.
固然就是一個牧羊人.
由一個在埃及的皇宮長大 他一定不需要怎麼做事.
反而就是接觸不同的高級的學問.
不同的演說的技巧 管理的那些能力.
不斷的來到他 透過不同的經驗.
不同的事情 來到他展現.
是一個出類拔萃的有為青年.
也都是慢慢的 人家都覺得他已經洗了底了.
雖然他仍然是希伯來人.
但是那個沒得洗的 與生俱來的血統.
但是他整個人的舉止談吐.
都已經是埃及化的了.
所以這七個女士 認出的他是埃及人.
不過《留意聖經》有這些伏筆都很重要的.
第二章的23節到25節.
相信比摩西的描述更加重要.
整個《出埃及記》某程度上的鑰節.
就是這一段經文.
就是神要紀念他自己的子民.
23節 過了多年 埃及王死了.
以色列人因作苦工 就嘆息哀求.
他們的哀聲達於神.
神聽見他們的哀聲.
就紀念他與亞伯拉罕,以撒,雅各所納的約.
神看顧以色列人 也知道他們的苦情.
有好幾個動詞.
神聽見 神看顧 神知道.
好像知識,情感以至我們人類.
可以具備的整個的身心靈的一個人.
我們的知識 我們看到 我們看見.
我們明白 這樣的全人整個生命的關懷.
譬如你去問人一些事情 他根本不想理你的.
你不斷說話 他聽到 聽到他的心完全無心裝載.
完全那個眼神不屑地看著你.
根本裡面在想 你知道他在想其它事情.
那你怎麼說 你怎麼流淚呢 他完全無動於衷的.
那就不是聽見看見和知道了.
所以現在說的是 受苦的哀聲達於神.
但是他們很無助 不知道怎麼辦.

$^{161}$仍然都是以自己是希伯來人 怎麼可以解脫呢?.
是不是集體的逃難 自己就可以逃走呢?.
似乎曾經是一個上賓的一個地方的人.
現在成為了被拘禁 當然不是有監獄.
被嚴謹的看守的一些平民百姓.
不過是要做苦工的平民百姓 一定有得給你吃.
不然你哪有力氣去做事呢?.
但是你已經失去了你們民族的自由.
那個自信 還有你們民族自己那種的發展.
所以這個哀聲絕對不是開玩笑的.
我想我們有沒有見過一些大災難那種哀號呢?.
我們沒有機會 比如我們在汶川大地震.
我們沒有機會去當地現場去做一個支援.
去拯救 或者一些特大的災難.
不知道大家在你們成長的經驗有沒有經歷過呢?.
我小時候住在九龍塘的小西湖村.
試過因為上游 所謂上游就是猛鬼湖的九龍水塘.
現在都叫做蓄水湖 因為那時候有一些大的工程.
所以試過一場特大的暴雨.
所有的山泥隨著那些水都湧下來.
因為我們一條大的坑渠 一路出到去.
應該是九龍仔,石劍尾,大坑東.
就是大坑東球場那裡 一個大的坑.
當然是暗渠來的 一路出到去.
但是在我們村裡 一夜之間.
因為那些雨是非常非常的浩大的.
大到一個地步是我們難以想象 那時候我很小.
大概六,七歲 突然間半夜.
家裡很吵 吵得很厲害.
原來周圍的鄰居走上來我們家避難.
因為那些水已經高過一間木屋的天台.
就是那個頂 你可以想象有差不多10呎高的水.
所以我們家是在半山的 當然是走上石級的半山.
不是大家想象的半山區.
我們一起來看下去 那時候已經很多人了.
我看出窗外面 那些水好像紅海那樣.
是滿到很恐怖的.
所有低於我們的東西全部都沖了.
接著過了 應該第二天早上 後來都很晚了.
真的沒有雨了 那些村民陸續回到自己家裡.

$^{201}$什麼都沒有了 接著我們小朋友周圍去玩.
整個坑 坑有我們那時候小時候那麼高.
差不多有五呎高 闊度是兩張桌子的一個闊度.
大概我想有八呎闊 八呎闊一直延綿很長的.
從這裡出到差不多村口 你可以想象.
整個坑全部變成了泥沙 完全封滿了.
所以那些水是從表面沖下來的 是非常恐怖的.
很深的感受到那些居民那種哀號.
因為什麼都沒有了 一夜之間沒有了.
但是各位你想一下 這裡是在說持續很多年的迫害.
那個家的苦工 所以就使他們的哀號來得很實在.
所以你想一下 如果你在現在中國有30位的專家.
就去了剛果準備提供一些支援.
都是非常受保護的 可能在一些技術層面.
防護的層面 都未必會直接參與前線的救災.
伊波拉的病毒 你可以想象得到.
如果你是那個醫護人員 在那個醫院.
你看到一個一個倒下 一個一個流血.
一個一個失救而死 那種哀號大家會感受到.
所以這裡不是因為摩西來哀號.
摩西已經落戶了一個米甸的牧羊人的家.
來去做女婿 她已經好像淡泊名利那樣.
她已經好像是無可奈何地.
她自己打了敗仗 要離開埃及王對她的追殺.
因為似乎好像她殺了這個埃及人.
使得埃及王或者埃及王周邊的那些大臣.
有一個警覺 原來這個人是養不熟的.
原來這個人的內心是無間道的.
原來她的內心深處還是流著希伯來人的血的.
那怎麼辦?她現在是我們最內圍栽培的一個人.
居然間是這樣的一個反轉.
所以一定是追殺是很合理的.
不會再說這麼多年的恩情.
因為涉及到國家安危的了.
所以這裡我們可以看到第一個段落.
神怎麼準備 但是神怎麼準備摩西之前.
神給我們看到的就是.
或者說作者按照神給他的感動.
寫下了這一段是很重要的.
我想一切的事情帶來的轉機.

$^{241}$都是因為人看到神有工作.
人看到神從來沒有放棄我們人類.
所以雖然這裡不是很多人看的金句.
因為19章神同以色列立約那段就是金句.
但是其實這一段應該是更加重要.
《出埃及記》那個原因.
為什麼會有《出埃及記》的原因.
就是因為神聽見,神看顧,神知道.
是整個生命的投入.
這個就是神自己的作為.
好了 我們看看這個工程是怎麼來到.
拆毀和重建.
摩西我們剛才看了 靠一個血氣.
來到一往的情深,一番的熱血.
做了一個好像很魯莽的舉動.
就是將一個埃及人可能在城牢之下.
來到去打死了,來到去藏起了,掩飾著.
但是原來這件事是有其他人看到.
甚至乎是希伯來人看到.
而來到去有一點點的不解.
為什麼你如果是希伯來人.
為什麼你的生活會這麼養尊處優.
為什麼你可以來看我們怎麼落魄.
為什麼你看完之後你又回到你的皇宮.
為什麼你不是和我們一起呢?.
所以無論對摩西而言.
對希伯來人而言都是一個矛盾.
但是摩西很難去解釋他內心的深處.
他那份的情,他那份對自己民族的愛.
對自己他認識多少呢?.
他對耶和華,對亞伯拉罕.
這些口傳的歷史傳到他的心靈裡面又有多少呢?.
這些我們可能要後面再逐樣逐樣.
透過經文給我們看到.
但是無論如何 神現在容讓這件事情的發生.
一方面 摩西的為人真的有一份這樣.
出於愛的坦誠的衝動.
總好過一些嘴巴說到真的要付代價.
一隻手指都不肯動 明哲保身的人.
某個程度很難去透過那個行動.

$^{281}$因為神根本看不到你的行動.
怎麼知道你真的可以將來去委以重任呢?.
當然神絕對不是說讚賞一些衝動,魯莽.
但是有時那個衝動的舉手 走出去接受呼召.
很想成為傳道人 不同的神學院.
總是有大大小小的獻心會.
或者我們去夏令營 最後那天的一個獻心會.
總是有衝動 當然絕對不讚賞說是不是衝動.
但是有時那一刻的感動未必是一個衝動.
有時那一刻出於一份很坦誠的愛.
雖然那個行動是做錯了.
但是真的代表著原來摩西的內心.
神對他的驗證完全是沒有看錯的.
所以仍然回到我們的心靈.
所以我們其實真的不要怪說.
教會都不給我機會 教會都不給我事奉空間.
教會都不重用我 教會看不起我.
我們都沒得事奉 那我就不事奉了.
似乎那個邏輯好像很對的.
其實在教會其實不是這樣的.
應該是剛剛相反.
應該是一個默默的,安靜的聚會,禱告.
神會在聚會透過不同的人事物.
其實你是要全程的聽,全程的看.
和全程的知道 你慢慢會有感動的.
慢慢神會感動你 你會為著某些事情.
按你的感動來禱告 還有你會安靜的去留意.
你有什麼可以在這個教會可以被神去用.
不是因為你是很聰明,很想大展拳腳.
而是你真的很想在這個地方.
令到這裡的教會得到一些的幫助.
你也有你自己的恩賜 在神的感動裡面.
時間也吻合了 你又願意這樣走出來.
往往這樣安靜的人 是會被教會的牧者發現的.
會被教會的牧者邀請的 經過禱告.
覺得這樣的弟兄姊妹真的可以來服侍.
而不是大大聲,巴巴閉閉.
你給我做什麼 你給我做什麼 你不給我做.
那我就不做了 我就來盤劈了.
其實不是這樣的 我們真的要求主幫助我們.

$^{321}$每個人都是能夠在主的國度裡面.
有一點點的服侍 其實是莫大的恩典.
有沒有得做都不要緊 而重要的就是.
我們必須在那個地方.
神讓我們有神的心腸.
所以什麼是神的心腸呢?.
你想一下 整個民族現在在這樣的情況下.
他們能夠做的 只是不斷的哀號.
他們不會說 其實我們真的不知道.
有多少人一直等候著.
亞伯拉罕的應許第四代滿足了 事後滿了.
四百年之後就可以離開埃及了.
有多少人將創世記的第十五章.
神和亞伯拉罕用一個這麼嚴肅的立約.
來這樣去應許了 現在兌現了.
有多少的以色列人為了這個應許而禱告呢?.
而不灰心 而不會是一種不斷的哀號.
咒罵來這樣去嘆息呢?.
相信是很少這樣的人.
當然我們不知道當時的口傳歷史.
去到他們當中是會怎麼樣的 但是無論如何.
弟兄姊妹 其實我們不需要憂愁有沒有事奉.
我們準備好我們自己.
神自然會感動我們去事奉.
怎麼能夠有這一份的感動呢?.
其實就是23到25節.
我們應該看到一個地方的困苦.
看到一個地方的心靈的哀號.
譬如有很多人他們得不到真正的心靈平安.
他們可能有表面的平安.
例如他衣食無憂 他就表面平安.
例如他身體很靚仔 很好 血管又很靚.
那就很平安 但是這些是隨時都會失去的.
如果有一些突然間的打擊.
譬如有一個很大的危機在香港.
有很多人失業 那個哀號就很明顯了.
那個困苦就很 大家能夠看到了.
但是心靈的困苦流離.
好像沒有牧人的羊.
耶穌在《馬太福音》第九章.

$^{361}$又叫門徒這樣看的.
所以真正服侍主的弟兄姊妹.
首先是知道看顧,看見.
還有心靈裡面也都聽見等等.
這樣的弟兄姊妹.
她就開始會在神面前禱告.
求神讓她可以更好的等待.
裝備來被神去用.
所以弟兄姊妹 我們回到來.
我真的經常教弟兄姊妹.
回到來先安靜敬拜神.
回到來你坐下 暫時不要跟別人聊天.
當然禮貌打個招呼 點個點頭是沒有問題.
但是安靜自己 今天你是帶著整個的心靈.
其實你不帶都要帶的.
因為你整個人是一起這樣走的.
你不會說 現在我留下心在家裡.
我留下什麼在家裡 那是不可能的.
但是往往就是我們真的留下我們的心靈.
心思,意念在家裡.
或者一坐下 我們已經習慣.
快點想一下其它的東西.
那就很沒有意識地去聚會敬拜.
這樣的人其實是不可能得到任何的祝福.
和神給的Blessing.
因為對神毫無興趣 毫無心 這樣是不好的.
所以準備好自己 留心聽 留心看.
留心的感動 這個教會有什麼重擔呢?.
有什麼嘆息,哀求呢?.
要聽到人心靈的深處的哀求.
那神就會用你了 因為你已經有神的心了.
所以這件事已經是很重要的一個鑰節.
我們繼續下去.
首先就是神要讓他真正看到自己的失敗.
剛才我們已經讀了 不再讀了.
11至15節就是他的失敗 讓他看到.
一定是要被神再做的那個影子.
就是不是一個培訓班 而是一個生命的破碎.
是實實在在的生命的經歷.
摩西有他的感動 有他的熱血.

$^{401}$有他對自己同胞的愛.
但是這個不是他上幾個課程培訓領袖.
他其實本身就是領袖.
他在埃及已經是成為一個獨當一面的領袖了.
《史圖行傳》七章我們上個星期讀過了.
他是學盡一切埃及的文化,說話,行事.
都是有才能的 這樣的新約的作者斯提反.
對摩西的形容 就是說.
他們猶太人的那個歷史的典故.
是很清楚的告訴我們 摩西他未被神選召.
被神破碎之前 他真是有他自己.
可能就是他有知識而帶來的傲氣.
他自己的那個角色,那個氣質.
真是一個王子的氣質.
你可以想像得到 可能他的出入.
都是有一些隨從 只是可能他叫了隨從.
不要跟那麼貼 但是都可能有隨從.
可能他自己根本就是他的衣著,談吐等等.
都是和以色列人格格不入的.
但是這樣的胸襠.
怎麼能夠感受到以色列人的哀聲呢?.
所以神感受到以色列人.
但是未被神改變,破碎的摩西.
似乎感受不到百姓,同胞的那個哀聲.
將來他成為一個領袖.
當他要獨當一面的時候.
他不需要再經常上西瀨山.
去感受一下那些靈氣 感受一下神的哀聲.
感受一下神的那個痛心,痛苦.
因為他在山下帶著這些同胞.
他已經每天深深的體會他們的環境.
他們的那種小信.
還有他們那種因為有時候對神的不尊一.
而帶來的周而復始的問題.
所以摩西已經是代表著.
在山下有神的那個內心的知道,內心的看顧.
和內心的聽見等等的了.
所以其實今天我們教會.
任何教會需要成熟的弟兄姊妹.
蓮芝固然一定是很重要.

$^{441}$蓮芝必須是每個禮拜.
我們都要打開我們的新心靈去聽,去看.
去心靈裡面動容.
這樣我們就越來越有這些神的感動了.
所以絕對不是空降是錯的.
空降如果不明白很多的氛圍.
其實往往就是有時候眼高手低.
又會帶來一些衝撞的.
所以這次摩西就是一個空降的領袖.
希望能夠用他的影響力.
用他的身份,用他的知識來帶領.
但是一出師就失利了.
就已經因為這個希伯來人的一番單打的說話.
就令到他要急速的逃避埃及王對他的追殺.
生命的破碎有時候是很難說的.
神在不同的人都是度身訂造.
我們回望.
我想我們牧者肯定是首當其衝的一個示範單位.
我們經過不同的人生的階段,學習.
或者作為傳道之後.
有時候要面對不同的人對事物的研判的看法.
總會有很多的張力.
你怎麼去處理這些人際關係.
這些就會令到我們越來越磨練.
只能夠去到一個有時候有理都說不清的情況.
只能夠默默的等待.
甚至去到一個階段.
今生都沒有一個合理的辯解.
其實就不需要辯解了.
問心無愧 對準神就算了.
當然這個不是一個擋箭牌.
我們自己亂來都可以 絕對不是這個意思.
而是你已經盡了力 你已經問心無愧.
按照公義去執行一些事情.
每個人都有不同的看法.
每個人都有自己的執著.
這些不能怪對錯.
但總是需要有一個決定的時候.
認不認同都是一個決定.
做決定的人盡了自己的力就算了.

$^{481}$未必是百分之一百的準確.
因為如果是百分之一百的準確.
那個人就變成神人了.
一定是有瑕疵 一定是決定之後.
原來有一些事情是做得未必最好的.
都是這樣的 但是必須是經過不同的事件.
不斷去學習謙卑 吸取教訓.
下一次會有更加的成熟.
因為如果每一件事都是神的旨意走出來.
那我也想 但是絕對沒有這個可能.
因為人總是受很多的時空.
受很多的輿論 受很多的環境的牽制的.
不是你想做就做得到的.
如果想做就做得到 連特朗普都不行.
何況我們只是一個小小的教會.
大家真的是來自不同的背景的人.
大家都有不同的持份 更加沒有可能的.
邪教尚且有一個他們自己的核心團體.
何況我們是一個多元的教會.
就更加有不同的意見.
所以在這個過程 其實就差不多等於.
這些已經是很現成的場景.
讓作為帶領的人必須是要學習謙卑.
學習聆聽 學習怎樣在不同的處境.
來做一些決定.
有時候每一個決定的背後都很艱難.
因為實在不能夠左右逢源.
順得多情 失數意等等.
這些是司空見慣的.
任何一個組織 任何一個人有超過三個人的地方.
你就需要有一個這樣的處境的考慮.
和一定有那個協商帶來的認同與否定的.
所以這些是必然會出現.
所以在這樣的狀況底下.
帶領者就必須不是用自己的血氣.
所以神拆毀的似乎是這個天然的那個人.
摩西他的雄心是好的 是值得欣賞的.
但是這種這樣的火熱 如果不經過奧煉.
不經過一些的制止.
就只會帶來 會令到被他帶領的人.

$^{521}$就會有一個困難 或者是非常之難受.
所以這裡就是給我們看到.
摩西他在這個小小的段落.
他就立即已經失敗 失敗到他要逃亡.
失敗到他40年都沒有再敢回到埃及.
似乎那個王朝都結束了.
但是好像對他下達的這個追殺令.
都仍然還不生效.
當然似乎埃及都給他一個空間.
不會追究到底 追殺到底 去到天涯海角.
而是有這個命令存在.
就是等於他不可以再回去埃及了.
因為還有這個禁令.
但是就給他一個類似.
好像他就是這樣來到去保外就醫那樣.
來到外面 他不再成為一個威脅.
所以他安危 他現在怎麼樣.
都不再是對埃及構成威脅.
所以一方面當然埃及王以為.
他已經是成為一個沒有任何殺傷力的平民百姓.
過著一些淡薄名利的生活.
不再和希伯來這個民族.
就是他們的苦工來到有任何的聯繫 掛鉤.
所以解除這個威脅了.
但是原來另一個角度是神有意識地.
讓摩西過一個淡如清水的牧羊人的生活.
很難做一個徹底的一個對號入座的比較.
很難說我們能夠真的有這樣的人.
剛剛一出來有這個雄心壯志.
就立即被一些事件潑冷水打沉.
接著就是沉寂幾十年.
在香港的社會根本不容許你有40年的沉寂.
80歲才出來傳道 你做什麼呢?.
沒有可能的 我相信未必是那個時間的長短.
而是那個人對自己真的沒有了那個雄心壯志.
而來到好像只是保留著自己是一個.
內心知道的希伯來人就算了.
過不過希伯來人的那個風格的生活.
不知道摩西40年有沒有祈禱呢?.
有沒有自己希伯來人對神的守節呢?.

$^{561}$這些都不是很清楚 對不起.
因為那些節日是由摩西之後才開始的.
希伯來人之前只是受割禮的.
所以摩西是沒有幫他們兩個兒子.
夾信來行割禮的.
所以第四章神就追殺這個摩西的.
所以似乎他已經對於自己希伯來人的.
那些規矩就沒有了.
因為亞伯拉罕的子孫第八天就要行割禮的.
但是摩西就沒有幫這個夾信去做.
所以第一個階段就是一個生命的拆毀.
沒有拆毀那個信仰 是拆毀那個事奉的心智.
但是神的工作是很奇妙的.
不會使那個心智完全磨滅的.
只會使那個人在那個事奉上面.
重新有一個新的看法.
不再用自己的血氣 反而懂得去依靠神.
默默的帶著一份感恩,依靠的心去服侍.
這樣我不知道每個人要用多少年.
去經歷這樣的拆毀和重建再做.
我們應該在這個階段會多一些的經文.
多一些的可能背誦 多一些的消化裝備.
你能夠被神用 或者用的過程感覺不是很舒服.
或者別人都覺得對你不是很舒服.
因為你很多的血氣 很多的表面的氣焰等等.
但是一個慢慢真的明白 認識自己多了的人.
就會真的開始有一種從內心的謙卑.
從內心對自己的本性的不濟.
來去認罪和感恩 神仍然給自己有這一份的心.
所以內心裡面的那個感動 事奉的心智.
這些神不僅沒有磨滅 神是不斷.
好像不是溫水煮蛙 是慢火來燃點.
慢慢只會越來越有那個心智和那個感動.
所以剛剛出來的傳道 剛剛神話畢業.
上個星期一 剛剛過去星期一.
自費傳道的畢業禮 我都有被邀請參加.
其實很感動的 因為你看到這麼多的神學生.
畢業去到不同的工場 默默為每一位祈禱.
當然我不是認識他們 為一些我認識的祈禱.
都希望他們能夠帶著這一份感動.

$^{601}$這一份心智去工場.
但是也都真的要為他們這一份很單純的熱血.
來去祈禱 因為肯定不同的教會.
不是那些人不好 而是一定會有不同的挫折.
有不同的背景 不同的文化的衝擊.
就差不多是冷水 冷水潑熱血.
雖然是很無情 但是其實往往是事實.
原來這些冷水是令傳道人成長的必須的愉悅.
所以通常去到一些晚年的牧者.
那個內心 其實不是說你潑冷水我們不會介意.
而是給我們內心那份仍然能夠事奉的感動.
其實多冷的冷水都不介意.
不是我們沒有了人家.
不是我們沒有了自己的尊嚴 不是這個意思.
而是相比於神的愛 對微信的朋友.
對弟兄姊妹希望他們能夠成長.
那些一兩句的冷言冷語.
或者那些冷水 或者那些誤會又算得了什麼呢?.
能夠仍然有一個崗位去服侍.
是一件很大的恩典.
剛剛我寫了45週年那份特刊的那篇文章.
都真的有感動葡萄園的那個比喻.
《馬太福音》第20章 真的大家一定要細味.
我時時忙忙就會看看這段經文.
給自己有一個反思 我是二十多歲信主.
第二年就已經是蒙照 接著裝備.
接著經過好幾個人生不同的階段等等.
現在能夠仍然在葡萄園.
我內心真的不斷去很感動地.
你給我一錢銀紙 我是很感動的.
我能夠現在繼續這樣勞苦受熱.
勞苦受熱是那些很早進入葡萄園的人.
發脾氣,發脾氣的描述.
我們整天勞苦受熱.
那個A君晚上六點進去.
你都給他一錢銀紙.
主耶穌講這個比喻 那種全神.
那種心靈的慨嘆.
為什麼你這麼久在服侍的人生.
你還這樣看不明白神的愛,神的心意呢?.

$^{641}$其實是很大的一件事.
一個事奉主這麼多年的人.
他埋怨的就是我勞苦受熱.
我不公平 對我不公平.
你講這些話 好像你完全不認識主.
對那個主人是很心痛的.
所以我給這個比喻感動得很厲害.
我能夠為神的國度有一點點的勞苦.
有一點點的受熱.
我很快就得到獎賞 一錢銀紙.
主人一定會說 曾牧師,玉貴.
你這麼多年為我的國度勞苦受熱.
家主一定在他的內心 難道他沒有眼睛看的?.
僥倖心態的人拿到一錢銀紙.
當然僥倖就帶來莫大的高興.
但是一個從七點鐘就進去葡萄園的.
默默的去付出 不問收穫 不計較的辛苦.
來為了主 為了福音.
所以我們一點都不需要.
去有這種錯誤扭曲的心態.
如果兄弟姐妹聽到這樣的訊息.
很偶然的你能夠開到這段經文聽到.
如果你真的要開始讀回《馬太福音》的20章.
那段經文 耶穌很語重心長的.
祂真的很想勸勉跟隨祂,服侍祂的人.
其實那個默契已經很清楚了.
我們在這個葡萄園裡面.
我們看到越來越多人進來.
五點多還有人進來做.
六點鐘還有人進來做.
其實你要讚嘆這個家主那種內心的愛.
對人的那種的愛惜.
我們作為工人的 開心都來不及.
都不知道多興奮.
看到有不同的後輩能夠加入服侍主的行列.
所以去神學院的畢業禮.
是很興奮的 對一些老的牧者來說.
我也算是的 是很興奮的.
看到不同的年輕人可以踏上工場.
那種興奮就是一個快要收工的葡萄園的工人.

$^{681}$那個莫名的興奮.
所以弟兄姊妹拆毀自己的血氣.
放在自己在神的國度裡面.
時時從一份恩典,不配這樣來服侍.
那一錢銀子咬下去都是有血有汗.
是很值得感動的.
因為你那一錢銀子是由七點做到收工.
那一錢銀子不是僥倖.
是神看到你是由始至終.
忠心勞苦的服侍.
是很寶貴的將來的一個稱讚.
所以我們牧者等待的就是將來那個稱讚.
你在你的人生 在我安排你的崗位.
你已經盡心竭力.
這個就是最大的一個祝福和獎賞.
所以仍然真的勸勉大家.
如果我們完全對服侍主沒有興趣.
一談論到主 你就要轉台.
其實你要問一下你自己是否得救.
你是否善無永恆的天家.
將來在天堂永恆裡面和造物主.
和愛我們的耶穌基督那個永恆的相聚.
所以我們真的作為一個好像老戰友那樣.
事奉了一段長時間.
我們現在是下山的時候.
真是這些經文是越來越感動的.
就是這些經文的那個核心.
使我們事奉完全不需要再問那些得與失的了.
就是很衷心地去服侍.
很興奮地每天能夠參與服侍.
弟兄姊妹 第一個階段就是摩西它要被拆毀.
它的血氣.
所以什麼是我們大家不同的盲點.
我們的心靈的弱項.
總是有我們一些東西.
神要先來幫我們做一個更新的.
這個更新就好像葡萄樹和枝子那樣.
因為葡萄樹和枝子大家都可能識背的了.
那個比喻 我是真葡萄樹.
我父是栽培的人.

$^{721}$凡屬我不潔果子的他就剪去.
潔果子的他就修理乾淨.
所以要潔果子多呢 原來都需要修剪的.
如果不屬主的就修剪 丟下去燒.
但是屬於主的 修剪就是使得潔果子更多.
所以原來拆毀是為著更好的服侍.
血氣用聖靈 用神的愛去代替個人的血氣.
也是會帶來一個更大的祝福.
所以我引述的就是《約翰福音》第15章.
「尼了我 你們就不能作甚麼」.
這個就是葡萄樹的那個比喻.
真葡萄樹的比喻.
第二個段落 第二個工程的建造.
相信都是為摩西來準備.
這個剛才我們已經講了.
我們就不再詳細講這一段.
非常重要 神聽見 神紀念 神看顧.
所以我們記得要為不同的牧者.
他有這樣的新心靈.
來去禱告 必須是真的有神這個心腸.
才是可以說真正的牧者的情懷.
其實就是神的牧養情懷 神的時間到了.
摩西現在在曠野來到去落泊.
但是神很清楚 祂聽到紀念和看顧.
做物主其實根本不需要用眼去看.
不需要去留意 不需要去聽.
但是近乎祂是質疑去聽以色列人的哀嘆的聲音.
他們的哀聲達到神那裡.
舊約還有哪一段有這些哀聲呢?.
約拿書就有了.
我們繼續下去 第二個大段落.
就是摩西的呼召 異象對摩西的啟示.
我們就要從這裡開始 我們還有時間的.
我們可以看看第三章這個呼召.
不過留意是 我先讀了 第三章.
摩西牧羊他岳府 米顛濟斯 葉蒂羅的羊群.
一日領羊群往野外去 到了神的山 就是荷列山.
大家要留意西奈山和荷列山.
因為往後摩西都會經常在荷列山.
當然會在西奈山 因為他已經離開了埃及.

$^{761}$進入了加拿大邊境.
第二節 耶和華的使者從荊棘裡火焰中.
荊棘被火燒著 卻沒有燒毀.
這是一個很普通的現象.
說不上什麼特別的異象.
因為在曠野 那些荊棘會突然間燒著.
可能有人放火 更加多的就是.
那些無論是那些日照的折射等等.
都會很自然的發生 因為如果極度的乾旱.
是會有這樣的事情發生.
但是如果聖經說這個是異象.
我們就要留意 為什麼它是異象呢?.
摩西都是在40年 在曠野牧羊.
就是說每天它的主要的工場.
就是帶著那些羊群去野外 就是新鮮的草原.
來去安堂 吃草 休息 喝水.
一直這樣來 經過這樣的野外的生活.
不可以圈禁的 羊是需要到處走的等等.
所以它就可以這樣來 一路一路長大.
但是問題就在這裡了 第三節.
摩西說 我要過去看這大異象.
這荊棘為何沒有燒壞.
是啦 焚燒 但是完全是完好的.
那就是一個很奇妙的異象了.
一個自然界 如果燒就是燒焦.
燒完就沒有了 整片的曠野的草地.
或者那些樹 那些灌木就會燒壞了.
那些灰燼就可能熄滅 因為沒有東西再燒.
這個是很自然的一個現象.
但是不自然的地方就是.
不斷地燒都燒不壞的 不斷地燒.
我們都見過的 那些假的就會.
其實那些都不是真燒的 當然不會爛.
但是真是燒 而是燒不壞.
那就是一個很奇妙的事.
似乎是摩西從來都沒有見過的.
對他來說這個是第一個大異象.
因為稍後是由他 那支棍.
後來就不需要那支棍了.
就是連他自己說一句話 就馬上產生神跡.

$^{801}$在其他人的眼中就是大的異象神跡了.
但是現在摩西他是首先進入一個心靈的與神.
同在的一個敬拜.
所以蒙召必須是你遇見過神.
神向你有一些的感動,託付.
這個就是尋求異象的開始.
不是自己去製造一些異象.
可以是從一些現象 現就是現在的現.
如果異就是異端的異.
可以從現象變成異象.
但是如果一個人說他自己被神的呼召.
而他沒有任何異象 那也都是需要再等待.
因為是一定不耐煩的.
因為他只是一腔的熱誠.
這樣很快就會受到很多挫折 他就會離開的了.
第四節 因為要看這個異象的時候.
耶和華神見他過去要汗.
就從荊棘裡呼叫說 摩西 摩西.
他說 我在這裡.
舊約也有很多這樣的對自己的稱呼.
亞伯拉罕也試過了 摩西也試過了.
不同的神的僕人 先知都有這樣的情況.
撒母耳也是這樣 起初他不知道.
老的祭司阿爾尼.
因為撒母耳還是小朋友 住在會幕裡面.
神就向撒母耳顯現 讓他知道.
撒母耳就第一時間衝去找祭司 吵醒他.
祭司就說 你回去睡吧 我沒有叫你.
第二次又是找以利.
以利說 你回去睡吧 你不要搞事.
當然這個是我家的 就叫他回去.
其實如果是這樣的時候 以利應該很心安.
因為一個小朋友兩次來找你.
你知道他平時都不是惡作劇.
其實很悲哀就是神不再向你顯現 向你感動.
是向一個小朋友來感動.
作為那個過去接收神心意旨意的以利.
大祭司已經沒有神的默示.
所以第三次撒母耳再來.
其實以利心裡有數 心中都很難過.

$^{841}$因為神已經pass了.
好像我們玩那些啤牌.
pass就是下一個 以利已經被神pass了.
所以下一次你再聽到這樣的呼叫.
你就說 僕人敬聽 即我在這裡這樣的意思.
所以似乎都是一個先知初次遇見神.
會有這樣的自我的一個意識.
就是我在這裡.
這裡就很特別 很重要.
神要一個僕人首先認識甚麼是聖潔.
原來聖潔 固然我們去到聖殿 去到致聖所.
這個是張倫摩西為神建造的回望.
有那個聖所 就是Holy of Holies 致聖所.
就是比聖所還要深入的.
有一些黑暗的海狗皮遮住 不准進去的.
要讓他首先認識聖潔的定義.
甚麼是聖潔的定義呢?.
留意這裡告訴我們 第五節.
神說不要近前來 當把你腳上的鞋脫下來.
因為你所站之地是聖地.
所以這裡就已經很後現代 讓我們看到.
摩西根本連一個回望都沒有建造.
神他都沒有認識.
但是他現在首次接觸神聖的事物.
他現在就好像去到最神聖的致聖所.
其實是在燒一些雜草叢生的礦物.
在燒垃圾 燒一些樹木.
就是現在有一點點的山火 受控的山火.
其實說不上甚麼神聖的聖地.
但是神的同在就是聖地.
這些概念我們新約都未必很明白.
但是對摩西在這個時候已經明白了.
所以記得新約就更加高級.
不是你回到禮拜堂才是神聖.
其實任何的地方都是神聖的.
所以我們有神的聖靈.
我們有一個敬拜神的心.
任何一個地方都是神同在.
固然一個地方 一個建築物.
是我們一班人去朝見神.

$^{881}$那裡有主的名下的地方.
是奉主名聚會的地方.
所以我們不要有另外一種極端.
就是每個地方都是神聖的.
那我用不著回到你的感受堂.
用不著回到你的洪恩堂.
用不著去你的嘉欣堂.
我在我家就是神聖 現在神和我一起.
這樣講得那過得而得來.
還是偏激 還是極之極端.
這個只是其中一個的教導.
如果是這樣的時候 神就不會這麼麻煩.
叫摩西建造回眸.
《出海及記》第14章開始.
就是大家不想讀的經文.
就是開始建造回眸.
開始離開這個西賴之後.
就要開始建造回眸.
為什麼要建造一個這麼實脈的神聖的聖所呢?.
就是要讓所有的人類學習.
在一個地方大家一起聚集 學習敬拜.
原來那個心靈的成長.
是需要有一個這樣的群體的敬拜的.
所以個人自己的敬拜 你不要騙我.
不出一個月 它就不會再回教會了.
不可能這樣的.
所以神知道人類是一個群體的動物.
所以我們是必須有一個敬拜固定的地方.
那裡是彼此守望,彼此禱告,彼此相愛.
彼此合一的地方.
所以從來不是「我在天上的父」.
是主耶穌說「我們在天上的父」.
是一個群體的敬拜.
所以這裡就告訴我們.
首先祂要認識聖潔.
因為祂是代表著神回去拯救祂的百姓.
所以祂首先自己要認識自己在神的眼中.
其實不是祂的腳骯髒 其實祂現在在聖地.
所以以脫鞋做一個實物的教材.
讓祂看到現在你就是脫了鞋站在聖地.

$^{921}$不是突然間有一張地毯.
不是突然間好像變化了的一個金的階磚.
完全不是那樣東西 都是那個礦物.
但是祂站在那裡.
因為整個的氛圍 神現在是臨近祂.
向祂顯現.
所以每個弟兄姊妹都要有你自己個人的夢照.
個人的感動.
可能你的夢照未必真是做一個傳道人.
但是你的夢照,你的感動.
神在你的內心的同在.
是給你有不同的那個向度.
你可以是突然間 我認識一位姐妹.
她去完菲律賓跟我們一起去垃圾村.
她回來之後她很有感動.
她找到自己未來的服事的方向.
很想去幫助更多的菲律賓的小朋友.
這些就是神給不同的弟兄姊妹的感動.
感動她可以關心自己的未信的家人.
感動她可能在她的公司做一些橋樑.
帶人信主 所以不一定只是感動.
是要做傳道.
那個傳道可以是一個比較闊的層次去想的.
感動你可以是繼續默默的準備自己的主一學.
去教導,去跟別人分享你對聖經的那個認識.
去扶植下一代的年輕人.
這些都是感動來的 這些都是呼召來的.
現在這個呼召對於摩西來說就是第六節.
我是你父親的神 亞伯拉罕的神.
以撒的神 雅各的神.
摩西蒙上臉 因為怕看神.
就是說整個氛圍讓摩西害怕到一個地步.
因為神已經是用一種讓人不會被擊殺.
讓人又不會在超自然裡面被融化.
而能夠跟它這樣這麼鄰近.
是摩西一個莫大的福氣.
現在是懼怕的 因為大家沒有任何關係.
沒有那個磨合的在事奉裡面對神的深長的理解.
所以去到晚年的摩西 其實就不再懼怕神的威脅.
不會再怕見神 因為它已經慢慢慢慢.

$^{961}$來到去擁有神給它的內心那份的愛.
和那個生命的成長 所以這個都是我們大家要學習.
我們不是敬而遠之的對神的那種懼怕.
反而就是越事奉.
我們越跟神有一份很深的默契.
這個就是恩典帶來的祝福.
所以其實事奉是會給我們有生命很大的成長.
如果不事奉 其實你只不過是將你自己的問題.
沒有被人知道的意思.
一個事奉者就一定是要面對不同的人.
你就會不斷要裝備 要改變你自己.
但是我們這個改變 正正就是一個生命的成長.
所以有幾樣東西的.
首先就是它看到一個現象的異象 這個第一樣.
第二 它必須是聽到神的那個呼聲.
往往那個呼聲就是第七節.
我的百姓在埃及所受的困苦 我實在看見.
他們因受督公轄制所發的哀聲.
我也聽見了 和剛才二章二十三節是完全一樣的.
我原知道他們的痛苦.
我下來是要救他們脫離埃及人的手.
這個就是呼召了 你接不接受 你有沒有這個感動.
你是否願意代替神這樣來看見,聽見和去感受呢?.
這樣就是呼召的最大的原動力了.
就是你要代表著我 你沒有我的能力的.
你只是給我分受給你的能力.
但是人是不會接受你的.
但是你就要將這個感動給他們看到.
其實在你的生命裡面 越來越愛這個民族.
所以 鼎姐妹 如果你想做傳道.
第七節就是很重要的那個感動了.
你有沒有看到,有沒有聽到.
有沒有完全內心知道.
就是有意識地被這個困難來感動.
而你願意投入去 希望神的恩典改變那些問題.
甚至最終改變不了.
你已經做了神要你做的事情 就已經足夠了.
因為有些地方 馬太福音第十一章.
有些事情在將來 耶穌說 哥拉遜.
如果你們能夠悔改.

$^{1001}$你們就不用將來去到沉淪了.
如果你們在你們當中走的神跡.
如果當時走在索多瑪 俄摩拉.
索多瑪 俄摩拉的人趕著去悔改.
那就是沒有了 但是哥拉遜是有很多的神跡.
但是都是終不悔改.
所以傳道人不一定要成功的.
不一定有很大的成就的.
但是一定要有這個就是.
和神一樣 聽見,看見,感受,原原本本這樣知道.
然後有一個行動就是要下來.
拯救這些百姓脫離埃及人的手.
帶領他們去到那個美好,寬闊,流淚與密之地.
所以一個傳道人往往都是一個屬靈的嚮導.
一個帶領的人 或者一個導遊.
導遊就是你對那個目的地非常熟悉.
你帶領那個群體 沿途你要很熟悉.
這裡就是哪個地方 我們還沒有到的.
還有五個小時車程的.
這裡大家要停一停 要看了.
這裡有一段歷史要告訴你.
一個導遊 傳道人都是一個導遊.
是要帶領那些群眾去到一個更美好的應許之地.
所以是應該 牧者的出現是會帶給一個群體.
或多或少那個靈性是有一個成長.
更好就是屬靈的生命有一個長大.
可以結果子 心靈裡面有聖靈的果子.
人愛 喜樂 和平 忍耐 恩慈 良善.
信實 溫柔 節制 這樣的一個新心靈的成長.
所以那個導遊應該就是對那個周邊的環境很熟悉.
那個目的地他又非常熟悉 去過很多次.
這個角度拍照不漂亮 告訴你.
這個角度拍照就很漂亮了.
因為通常去到就瘋狂地就馬上拍了.
原來那個導遊說 不是啊 這邊才是最漂亮的.
這些就是經驗了 這些就是那個的形式了.
所以任何一個教會的牧者都是希望帶領那個群體.
去進入那個美好 流淚 與物之地.
所以拍《望弟兄姊妹》來一同去學習.
今天開始喉嚨都有點刷聲了 因為沒有停過.

$^{1041}$所以可能都要留到下一次再說.
本來想說這個例子 都不說了.
這個都是相當感人的一個見證.
這個弟兄是主信得很好的.
他的人生都是非常傳奇和坎坷的.
我們在齊同心土報.
你是呼召人進入葡萄園的主.
你說莊稼多 我們要求莊稼的主.
來收取莊稼 求你讓我們聽到,看見.
和感受到未信的朋友.
他們未必是身心靈.
讓我們隨時都看到他們的落魄,他們的痛苦.
可能在心靈裡面 他們有很多的困難.
可能家庭裡面 他們有很多說不出的困難.
在他們的內心 有時未必從他們的面容.
可以看到一些的憂傷,一些的困苦.
但是天父給我們有屬靈的眼光.
給我們有神的心腸.
就是真的看到以色列人的哀求,歎息.
而且是哀聲答於你的.
以致我們真的能夠帶著有心,有血,有肉的.
一份的愛 來幫助有需要的人.
給我們每一次的教書,我們每一次的聚會.
我們有每一次的領唱,事奉.
我們都是用心的求主幫助.
我們帶著今晚這個學習.
去禱告求主在這個訊息聽到的人的內心.
雖然他們在不同的地方.
甚至在歐洲的弟兄姊妹.
能夠有感動,有主的心腸.
去看自己現在的處境 一起去學習.
我們這樣感恩禱告 奉耶穌基督的名 阿們.
\newpage



\section{出埃及記}
\label{sec:cWahrBO910c}
\textbf{2026-06-26 周五晚查經班《出埃及記研讀(五)五個「我不能」與一個「我是」》 講員:曾裔貴牧師}
\newline
\newline
連結: \href{https://youtube.com/watch?v=cWahrBO910c}{\texttt{https://youtube.com/watch?v=cWahrBO910c}} ~~~~ 語音日期: 2026-06-27
\newline
\newline
\hyperref[sec:HMEU7PxsF0Y]{\small{< < < PREV SERMON < < <}}
~
\hyperref[sec:index_chronic]{\small{[返順時目]}}
~
\hyperref[sec:index_scriptual]{\small{[返順卷目]}}
~
\hyperref[sec:dj6fIzaDTMQ]{\small{> > > NEXT SERMON > > >}}
\newline
\newline
$^{1}$同心禱告 再一次感恩天父又招聚我們.
雖然停了兩個星期 又假期 又一些其它的原因.
但是多謝你讓我們再開始回到第五講的《出埃及記》.
有兩章半的經文是講到摩西的蒙召.
整個的心路歷程 整個人生的跌宕起伏.
我們真的看到神的記載的經文.
這一個敘事是鉅細無遺.
這樣將摩西的人生讓我們有非常立體的認識.
祂不是一個完全人.
祂更加不是一個我們眼中沒有任何過錯的人.
但是聖經都好像沒有刻意去掩飾.
而是將一個真正有性情的人擺在我們的面前.
多謝主 看到的是祂內心那份的感動.
那份對同胞的愛.
那份很想改變世界的一個心願,一個夢想.
天父啊 我們今天在我們急速的21世紀.
有時我們對於夢想.
有時我們都好像連望都不敢去望.
幫助我們能夠真是明白.
在你的永恆的國度裡面 我們是有份的.
在你的葡萄園裡面 我們都是有我們被召的崗位.
來做一點點的工作的 求你施恩帶領我們.
也都讓透過摩西的夢召.
都給我們真是能夠想一想神理的作為.
往往這些都是屬靈的原則.
是真是放諸於四海都是合用的.
求主帶領我們今晚的查經班.
我們仰望禱告奉耶穌基督的名 阿們.
好 我們來到第五講.
跟大家重溫一點.
今天的分題其實你都大概看到神的作為.
在第二章 其實應該由第十六節開始.
其實嚴格來說是十一節 一路去到四章.
所以是相當長的篇章 在佔據出埃及記.
雖然叫做四十章 但是對一個人的.
這麼細緻的刻畫他的人生.
你可以想像得到這幾章的重要性.
或者對摩西自己領受這個呼召的重要性.
往往這個起點 或者一個這樣的記號.
成為他日後面對很大的困難.

$^{41}$一些的障礙,攻擊.
都成為他一個很重要的後盾.
不是他自己一種的好高騖遠.
一種崇尚要領導,要操權.
而是往往回到那個起點.
是神揀選他來成為一個合用的器皿.
所以這裡就給我們在這幾章裡面.
看到一個大的原則 就是摩西的推搪.
用「我是不能夠的」 他用五個理由.
但是這裡很寶貴 就是神一直都說.
他是那一位「我是 I am that I am」.
的那一位的造物主 就是自由永有的意思.
我們開始進入經文裡面.
就給我們看到一個循環.
來將摩西的人生去展述出來.
我們上個星期 就是上一次.
我們都看到神怎樣將這個宮人.
從他自己很大的心願.
因為他已經學盡了埃及人的文化.
又自己雖然是已經生在皇宮裡面的王子.
有記載的經卷 就是一些的典籍說.
他可能會繼任這個法老的兒子.
就是說將來他就是法老王.
但是他這個內心.
擁留著的希伯來人的身份.
那個內心的民族意識.
相信都是由小朋友.
由他的姐姐米利安 由他的媽媽約基別.
或者他都有很多機會去看這一幫人.
所受到的不人道的虐待.
要做一些苦工 就是沒有了民族的尊嚴等等.
所以激發他很想去幫這個民族.
去爭取 例如一些談判 一些的公允.
但是他的計劃都沒有真正實踐.
就已經有一些這樣很簡單的意外.
錯手 也不可以說是激進.
而是在一種很心急的情況下.
就打死了一個埃及人.
這個是上一次我們討論過.
所以原來那個失敗 那個離開.

$^{81}$那個逃跑 去到離開埃及.
去到曠野的米甸 來到很落魄的一個人生.
寄人籬下的做一個入贅的一個男士.
這樣就希望平平淡淡.
沒有什麼人認識 這樣去過他自己的.
很安穩的牧羊的人生 或者餘生.
但是我們又看到 原來這個陳夕琦.
大概40年了 因為她是80歲的時候.
神就正式向她顯現 等了40年.
等到摩西都已經沒有任何的憧憬,夢想.
想幹一番事業的抱負 都慢慢的消磨.
都沒有那個意志 家人,太太,西伯拉,外父.
一個米甸的祭司葉蒂羅 都不會覺得.
她是可以做一番大事業.
真是以一個很重要的領袖 回去埃及.
還要帶領著200萬的勞工,希伯來人.
整個民族徹底的離開一個國家.
去建立另外一個國家.
難以想象是匪夷所思 是不是?.
如果有地震 你帶2000人離開那個災區.
都已經很偉大 或者非洲經常都有部落戰爭.
如果你是一個大的工廠的經理.
你要帶那些人 因為現在要撤僑了.
你要帶他們穿越那些戰區.
去到另外一個算是支持中國的.
讓你通過 然後去到碼頭.
有解放軍來到 有船去撤僑.
只是你要通過那個地方 那個戰區 都已經不容易了.
何況你帶著的是200萬人 整個國家的人.
一次過離開埃及.
你可以想象得到 從來古往今來都沒有試過.
整個國家去移民走 是難以想象的.
所以都不會在摩西的法蒙 她的老海.
日間 晚上都不會這樣想事情的.
但是原來這40年 就是神的時間和計劃.
一直這樣準備 所以你看聖經你就會看到.
神一個這麼大的工程.
原來默默的在預備一個人.
如果這個部份 舊約就相當充沛的.
其實新約都有 主耶穌都預備12個使徒.

$^{121}$都有很多的教導 差不多一半時間.
一年半都是教導 對不起 應該這麼說.
頭一年半 耶穌就是公開的傳道.
去累積到一定程度的反對.
很多人對耶穌的誤會 那個困惑.
那耶穌就在頭一年半就開始低調的.
弟兄姊妹如果還沒有關掉那個喇叭.
麻煩大家自己靜音.
那就不會影響其他弟兄姊妹.
你會看到 後期有人就說.
哥薩尼亞,菲律賓那裡是一個稍微的分水嶺.
耶穌就不再太公開的傳道.
不是說祂不去 而是不再是大規模.
有很多群眾這樣的 而之後後一年半.
主力訓練門徒 預備他們將來.
主耶穌升天之後 這幫使徒.
就會成為耶路撒冷的獨當一面的領袖.
所以你會看到 無論舊約 無論新約.
主耶穌訓練門徒 或者神預備他的工人.
都是會有明確的訓練過程的.
例如以賽亞書的頭六章 就是以賽亞的夢召.
耶利米書的頭六章 當然第一,二章比較明顯.
就是耶利米的夢召.
其實很多的先知都是這樣的.
例如撒母耳記上的頭三章.
又是說撒母耳由他的媽媽用信心去祈禱.
得到一個兒子 奉獻給神.
奉獻給一個混亂的時思時代.
成為神的一個很重要的人.
去扭轉整個國家的命運.
頭三章的撒母耳記 又是說撒母耳的成長.
甚至在聖殿裡面 在以利的面前.
就是一個不好的老師底下來栽培成長的.
你想想你將一個小朋友放在兒童主義學.
放他在錦宣或者在紅恩.
那個主要的老師又不教聖經.
說那些無無畏畏的東西.
有時候說一下風水 說一下命運 說一下命理.
你作為爸爸聽到就慘了 為什麼那個主要的老師會這樣.
當然這個是開玩笑的 其實沒有的.

$^{161}$接著他就帶那些小朋友走去做很多不好的事情.
接著他又躲在門口在麥當勞.
接著在外面吸煙 說這個是你的主要的老師.
如果你是那個爸爸 你是那個媽媽 你也會害怕.
但是聖經說 撒母耳是在以利的面前.
來去事奉耶和華 在頭三章講了兩次.
所以舊約相當明顯的 頭那一段的記載.
是慰藥這個先知 或者慰藥這個神的工人.
他是將來怎麼去面對那個時代的改變.
所以是需要那個工人自己先有那個生命改變.
裡面非常的硬朗 硬朗到好像to be water那樣.
那就可以真是無論別人的攻擊.
無論別人的誇獎 你都來到去用一個很平常的心.
看到一切的榮耀都是應該歸給神.
一切的攻擊 問心無愧就算了.
不需要介意別人怎麼樣的說法.
不是不介意 而是真是要有反思.
是不斷反思 所以這裡就是神預備的.
工人的那個formula 或者說一個的pattern.
給大家來看到 所以如果我們有心服侍的.
大家看看這個powerpoint 都相當重要.
一定是要先看到自己的軟弱.
看到自己的無助無有.
因為如果什麼都行 就不需要依靠神.
如果真正你覺得需要的時候.
那就變成你真是不斷去依靠神的一個機會.
和一個這樣的人生 接著就是那個預備期了.
神就為他預備在曠野的一個人生.
很奇妙的 很自然的 投奔了別人一個家庭.
那個岳父又看中他 而那個太太西波拉.
又很願意嫁給摩西 摩西已經.
我想男士就不是叫做洗盡原話的.
她已經沒有那些在埃及的傲氣.
穿那些王子的漂亮衣服 走到去是走爛的.
是急速的走爛 不是現在乾隆皇下江南.
或者一些很貴族的王子去渡假 不是的.
是真的走爛去到曠野 所以什麼都不計較了.
現在人家收留你 甚至願意將女兒許配給你.
你真是覺得是一個很大的祝福.
安安靜靜的就生了兩個兒子了.

$^{201}$接著原來我們上一次說的就是那個呼召.
很特別的 一般人看到的就是那些荊棘燒著.
燒不著 燒不壞 燒不壞 很尋常的.
曠野你一向都有山火的 有山火有什麼稀奇呢?.
摩西也不是又有山火 很稀奇的.
不可能的 你天天看到 基本上經常都有山火.
你就天天都稀奇 這樣不行的.
我《駐大公報》每年的清明 秋天就會多一些.
因為秋天沒有雨水 秋天的重陽.
一定有山火的 因為那些人免得麻煩.
你想一下要帶個鋤頭上山 鋤開一些草.
沒有人這麼傻的 拿個打火機就簡單了.
燒了那些草 有時候就是他們燒不完那些火種.
那些人走了之後 他們自己會再燒.
所以一般都會有山火的.
基本上你住在新界區 錦田一帶.
你就會知道重陽一定有山火的.
有很多時候是人為的 除非現在這樣的雨季.
突然間有山火 那就可能是天然的了.
有時候太陽的折射 令到太高溫 那就不奇怪.
但是如果是重陽等等 一定是有大部分人為的了.
所以你可以想象到 摩西這次看到這個異象.
它看起來是一個很大的異象 也就是說很特殊.
特殊到要引起摩西的關注.
他要過去看 為什麼會有山火.
而這些山火 居然燒那些荊棘.
那些荊棘是燒不壞的 荊棘根本是野草.
燒不燒得壞有什麼所謂的.
又不是一個很大很漂亮的葡萄園 是不是?.
現在不是一些很重要的花卉 去有大的問題.
例如如果你在法國 有一些很漂亮的葡萄園.
那個收成的季節 如果有大火是不得了的.
說馬上找一些人去搶購那些葡萄.
大家都會明白這些了 在一個曠野燒野草.
很小的事 但是居然燒野草而引發摩西高度的關注.
而他走過去 走過去之後聽到聲音.
所以這些是很獨特的 一般的現象.
成為神的異象 明顯就是一個很明顯的枯燥.
在這個獨特的人身上.
所以不需要去製造 人為的去複製.

$^{241}$去重複它 因為每個人神的呼召的臨到.
都有不同的方式的 因為呼召的是神.
我們不是現在聽完今晚的講道.
我們明天就走去行山 就看看周圍有沒有野火.
接著看到就走過去 等一下有沒有聲音.
那你就不會有的 因為神是用不同的方法.
在那個人 祂呼召的人的生命裡面.
那個channel好像是校對了那個臺.
可以接收到那個channel 接收到那個呼召的.
所以重要的是 我們準備好我們的生命.
去接受神的呼召 其實摩西根本沒有準備.
他沒有準備原來就是神的預備.
我們很難想像的 摩西現在不是.
我現在先這麼有空 牧羊嘛 我就讀一讀.
那些《孫子兵法》 將來怎麼去行軍調將.
將來怎麼去管理這麼多的百姓.
有沒有千夫長 百夫長 我都有了.
當年和我一起的 現在都80歲的.
那幫喜伯來的領袖 將來就可以給我用了.
相信摩西什麼都沒有 肯定不會有任何的連繫.
也都不敢回埃及了.
所以根本是一個心如止水的一個老人家.
但這個就是摩西 但是呼召就臨到.
今晚我們看的就是在最後這個部份.
就是神的那個裝備 預備他真是身心靈.
都能夠具備成為神的功能的那個想法.
上一次本來有一個很寶貴的一個見證.
長話短說 很感動的.
這位國立台灣大學的一個環境管理的博士生.
非常之振奮的一個心路歷程.
因為他其實是一個肯尼亞.
一個相當窮的非洲部落的一個小朋友.
窮到吃那些地瓜.
有時候一餐都未必能夠一天有一餐吃.
那種窮法是我們難以想象的.
不會有鞋穿 那些地下有時候會踩到東西.
又會有寄生蟲 所以他說他小時候.
晚上爸爸會幫他用一些針挑那些寄生蟲.
挑的時候流血,痛 他說這樣都不是最痛苦.
最痛苦就是要在傷口上灑點鹽 因為要消毒.

$^{281}$所以那些灑鹽的時候那種痛是難以想象的.
但是有一個人改變了Levin的一生.
就是這位叫做曾古德 大家可能不認識的了.
是在肯尼亞一個研究大猩猩的一位女士.
終生整個人生就是研究大猩猩的生態.
和他們的習性 拿了博士學位.
拿了很多學位 一生在肯尼亞做這個研究.
真是每天埋身去看,去觀察那些大猩猩的習性等等.
所以是一個人類學家裡面相當重要的一個博士.
就是去到肯尼亞 因為不懂得講當地話.
剛好這位Levin是相當勤力讀書,相當努力去成功的一個學生.
所以在小學,中學都已經是全校最厲害的.
學校安排了他和這位曾古德做一個翻譯.
所以就這樣的緣故讓他認識了曾古德.
曾古德這位女士就看到這個黑人的少年人.
這麼有上進心就鼓勵他讀書.
還有如果你真的經濟有困難 你來找我.
我可以幫你拿獎學金來讀書.
所以輾轉這位黑人少年就完成了高中.
後來就在升大學的時候.
曾古德這位女士就鼓勵他去台灣進修.
所以輾轉他就去到台灣之後讀大學.
接著就進了台灣大學 就是台灣最高的學府.
就完成了他的博士學位 由一句中文都不懂.
到他能夠講流利的國語和台灣的土話.
你可以想像到他的成長的故事.
就是受著曾古德這位女士,這位教授的影響.
帶來他整個人生就是這樣.
他應該是今年的台灣大學學生代表的支持人.
所以你今晚上網去聽.
很感動人的十幾分鐘的一個故事.
一個很勵志的故事.
所以原來我們去認識自己是誰.
還有有人能夠在你的生命裡面.
去加一點點的色彩.
帶給你有一個生命的盼望都是很重要的.
所以今天我們要灌輸給其他人的.
應該是盼望,夢想.
應該是對人的生命的建立.
例如扶他在最需要的時候.

$^{321}$扶他一把讓他去展現他自己的夢想.
所以這位Levin他會回去肯尼亞.
在那裡去幫助更多當地的窮苦的小朋友.
有機會像他這樣透過學業.
透過一路的成長去改變自己的命運.
和改變國家的命運.
所以偶然間上一次我看到這個YouTube的短片.
都很感動我自己.
所以你看到神下來拯救這位摩西.
所以其實是神給這個異象.
異象背後就是一個理想,一個夢想.
這個摩西.
所以我們今天在教會我們來查經班.
我們其實是要裝備下一代的信徒.
讓他們身懷那個理想.
胸懷一個普世的價值觀來服侍主.
所以如果我們有足夠的人生資源和經歷際遇.
而我們沒有去運用,沒有去讓神.
在我們的生命裡面去運用.
原因只是一個.
他還沒有尋找到他的那個夢想.
他還沒有看到那個異象.
可能是我們現在仍然在位的.
沒有將那個異象給到弟兄姊妹.
弟兄姊妹只會在自己有限的知識的領域.
追尋他認為的理想,他認為的異象.
例如反正生活都不太擔心了.
現在時間也多了 可以多點去旅行.
多點去遊歷 多點去做自己喜歡的事.
也不想太多煩惱的事.
不想太多人事的紛爭.
這些全部是很正確的思想.
但是有時忘記了就是.
想的仍然是在地上的那個短短的人生.
沒有為神的國度去思想.
沒有在主統文的六個向度.
有頭三個耶穌教我們的人生.
其實第一件事你要想的就是神的名得到榮耀.
神的國度要降臨在地上.
神的旨意要暢通在地上.

$^{361}$接著才是我們地上的日用飲食.
還有我們去面對仇敵的試探.
我們能夠有力可以去面對.
我們就能夠赦免人的罪.
因為我們自己都是蒙赦免.
人際關係可以搞好.
所以如果我們弟兄姊妹.
只是在自己的人際關係上面.
來去不斷的縮窄 不要再這麼複雜.
去看這麼多事情 做這麼多事情.
做回自己的家庭 做回自己.
自己開心享受就夠了.
有空就做一下這個 做一下那個.
已經夠開心了.
但是我們如果在我們這麼多的教會.
都只是地上的這個小圈.
而沒有了神國度的那個大圈.
那就是一件很嚴重的問題.
就是說在一個時代已經失落了夢想.
失落了意象.
所以今天透過摩西.
我們應該看到就是.
我們要將真正的意象永恆的夢想.
我們挑戰弟兄姊妹 讓他們去看到.
所以這個黑人的年輕人的故事.
是相當勵志 相當的正面.
他很快看完自己讀完書 自己很有成就.
他想到的就是.
他很快就組織一個非洲的.
十個非洲的留學生.
一起關心自己國家的那個環保的問題.
經濟怎麼樣能夠轉型 各方面.
那你可以想象得到.
真的會將自己的學有所成.
看到是一個理想,一個責任,一個意象.
所以今天我們看這一大段經文.
看看神怎麼裝備這個摩西.
使它成為一個終極的神的工人.
好了 整個佈局就是這麼多了.
我是不能夠講了五次.

$^{401}$但是最後就是應該是三章六節.
沒錯 或者後面都有了.
我們等一下讀聖經的時候再一起看.
那你就可以想象得到.
神這一句話帶來的那個震撼的影響.
就是對摩西,對往後我們每一個弟兄姊妹.
三章的十四節又有了.
我們正式下去了.
我們在經文裡面由三章十一節讀起.
如果長一點 第七節就是神給他的異象挑戰.
第七節 這個是我們上一次未詳細看得仔細的.
第七節開始.
耶和華說 我的百姓在埃及所受的困苦.
我實在看見了.
他們因受督公的轄制所發的哀聲.
我也聽見了 我原知道他們的痛苦.
我下來是要救他們脫離埃及人的手.
令他們出了那地.
到美好,寬闊,流淚與密之地.
就是到迦南人,黑人,阿摩利人.
阿摩利人大家有印象.
因為我們一開始第一章已經講過.
第一章已經講過 整個出埃及記.
其實就是回應神應許了亞伯拉罕的那個應許.
就是將來第四代 你的子孫.
就是希伯來人會受壓迫.
這樣來到後來就要離開.
去上來這個應許的迦南地.
為什麼要到第四代呢?.
因為神跟亞伯拉罕說了.
當時他還沒有改名 十五章的錯世紀.
因為阿摩利人的罪孽還沒有滿盈.
罪惡還沒有滿盈 所以要去到第四代.
那就是說什麼意思呢?.
一代一代的恩典都是這樣來給阿摩利人.
給迦南族的原居民 他們繼續的拜偶像.
很多的邪風異族,很多的迷信,很多的拜鬼.
很多的離開了真神 一代一代給恩典.
但是到第四代已經是累積到.
阿摩利人的罪孽滿盈了.

$^{441}$神就興起摩西 那個時間到了.
就回去拯救他的子民.
另一個角度就是審判迦南人.
所以這裡告訴我們 到迦南人,黑人,.
阿摩利人,比利西人,希美人,耶布斯人之地.
現在以色列人的哀聲達到我耳中.
我也看見埃及人怎樣欺壓他們.
原來神一直都知道的 一直都看到的.
還不是只能夠忍耐這些的痛苦,這些的不公.
這些的哀嘆聲,這些的哭泣.
眾日主好像真是 我們很難想像.
我們有時候忍不了 有時候就會衝動.
忍不了就會衝動 有時候做了再說.
但是神的時間又不會令到他有一種衝動.
但是那個忍耐又是真的很痛苦的.
因為看著自己的百姓受苦.
你可以想像得到 這個就是神的奇妙的.
那種永恆的平衡 就是神屬性的那種平衡.
既是忍無可忍 但是又能夠在那個時間的配合.
他說 使得你可以將我的百姓以色列人.
從埃及領出來 開始來了.
摩西感受到那個埋身的挑戰了.
意象就不是給你憧憬一下 很羨慕.
大家茶餘飯後聊聊天.
現在是針對摩西里.
整個曠野就只有你在那裡跟我說話.
你現在就是 你已經是俯伏在那個曠野.
其實沒有一張紅地毯在那裡 仍然是曠野.
仍然是一個泥地來的 仍然都是那些荊棘.
仍然是荊棘燒著的 但是現在神就要它跪下.
因為你所站的是聖地.
現在摩西是俯伏在這個曠野.
就不是真的有一張紅地毯.
但是現在很清楚 我要你帶領你的子民.
離開埃及 從埃及領出來.
這個意象 那個現象變成一個具體的意象.
具體一個的呼召 就很清晰了 對摩西來說.
對每一個的傳道人來說 都是要很清晰的.
所以全時間呼召的傳道人.
教會如果你說 不是的 教會現在不用的了.

$^{481}$都沒有傳道人 不如不需要傳道人這個職位.
現在說到都很容易 YouTube這麼多的影片.
高銘謙的 郭文慈的 李思敬的.
很多很敬虔的牧者的講到.
真是你聽一百堂都聽不完這麼多.
就是說你聽完十堂 一百堂.
大概有個模式 其實很多評論圖都可以講到的.
講得比很多傳道人還要好.
不過 仍然是這個問題困擾著福音派的牧者.
福音派的教會 為什麼教會必須是專職.
全時間呼召的牧者 傳道去帶領教會呢?.
為什麼不找一個CEO 學管理的去管理教會呢?.
完全沒有問題的 可以的.
很多教會都有行政主任.
譬如我們教會都有行政主任.
為什麼都要牧者 傳道 全時間的.
呼召的 來帶領教會呢?.
相信那個問題就是.
因為這個是神的教會 神的工作.
這個不是一個機構 這個不是一個公司.
雖然我們現在的註冊 根據政府的安排.
我們都是要以公司註冊處.
這樣註冊一個非牟利的群體.
那個是地上的一個 基於那個法律上的安排.
那個的管理 但是政府完全不會干預.
你的屬靈帶領 所以你們自己內部.
有一個屬靈的組織 是完全任由你們的.
所以屬靈的組織 就是執事會.
是看我們的體制是不是跟《聖經》的長老制.
如果是長老制 就是牧者 傳道人.
是當然的帶領者 帶領教會 全年.
或者一個大的方向的發展 是屬靈的方向.
所以這裡是一個很重要的分水嶺.
第十一節 摩西就清楚自己將來被呼召.
成為那個領導人 非從小可.
絕對不是開玩笑 要做一個帶領者.
他就因為聽得太明白 太清楚了.
是你將我的百姓 以色列人 從埃及領出來.
這個第十節 故此 故此就是說.
我聽到 我見到 我感受到 那個水深火熱.

$^{521}$我聽到 見到 我已經感同身受的痛.
已經是在說四代了 現在第十節就是故此.
Therefore 我現在就要呼召你 打發你.
我讓你看到那個現象.
我現在讓你聽到很清楚的聲音的呼召.
你現在就要回應這個calling.
所以我要打發你去見法老.
使得你可以將我的百姓 以色列人.
從埃及領出來 第十一節 第一個的不能.
我是誰呢?我是什麼人呢?.
竟能去見法老 將以色列人從埃及領出來呢?.
清清楚楚的對話 我是誰呢?.
是那個重點 就是他已經再不是以前的摩西.
以前的王子 以前的那個洪森.
壯志的年輕的40歲的摩西 仍然是摩西.
可能他的名字 他自己都忘記了.
是人將你從水裡面拉上來的 叫摩西.
他都是在水裡面被拯救出來的.
他自己都是一個被人拯救的.
一個本來在那個蒲草箱.
在尼羅河浸死的一個棄嬰.
但是現在我幫不了什麼忙了 我是誰呢?.
我現在苦心自問 我只不過是一個牧羊的老人家.
我有的是 等下那些像很重要的.
那些像就是神給他一個證據 陪他一起.
因為去到第十章打後 連那些像都不用了.
因為他已經有足夠的信心和那個屬靈的慰藩.
那些像都不需要了 他的說話就已經是能力了.
因為他扔掉那些像 那些像就變成蛇.
那你可以想象得到 原來人生都是循序漸進.
所以原來有板有眼的那個學習都是重要的.
所以我們不要天馬行空 不用的了 不要讀神學.
浪費時間了 就去傳就行了 有神同在的了.
摩西都要慢慢 由拿起那些像.
到慢慢那些像都不需要了.
已經感覺到自己真的有聖靈同在.
有神的那個慰藩了.
所以第一句 一般人都應該這樣問自己.
不知道大家回教會有沒有想過自己的照明呢?.
你說我都不是傳道.

$^{561}$那你有沒有想過自己的信徒的身份呢?.
一路的裝備呢?有沒有一個想法?.
有沒有一個曾經內心的一個 可能是很羨慕.
可能是一個野心 可能是一個你真的不值得.
看到很多事情 你很想去有一點點改變.
但是你又不知道從何入手.
總是這些無論是怎麼樣的.
有些可能是經過包裝.
或者沒有經過包裝的那個動機.
但是都是會從這個起點的.
所以很多時候就是我們很想去做事的時候.
就不是太認識自己是多少的斤兩.
或者是多少的屬靈的生命的那個狀況.
因為我都不會這樣盤點自己的.
總是覺得自己做到事情.
自己很想去幫一幫忙.
如果越是一信主 就已經這麼強烈的.
通常就是這個人真是可造之材了.
因為他這個心是非常單純.
因為他什麼都不懂的裡面.
他就已經有一個很內心的深處的感動.
很想被神用 就是說他的信主信得很好.
他真是有一個好像彼得田書所說的.
出死入生的經歷.
他就是一個很大的某程度上是一個劇變的信主.
這樣的弟兄姊妹很快就想到很想去服侍神.
很想被神用 這樣的時候.
如果能夠在這個階段有好的栽培.
有好的建立 讓他慢慢對教會.
對事奉 對很多的神學有足夠的裝備.
那就會好了 所以往往.
如果這樣的心態的弟兄姊妹.
尤其是這麼原始的那個內心感動 非常好.
當然很少人一開始就相當具野心.
相當具一個機心 我都扮一個裝備.
扮一個感動 其實我想混入教會.
在教會裡面我將來要坐高位的.
很少人會這樣的 如果真的有這樣的人.
他一路就會很順暢.
因為他內心的動機是非常非常的包裝.

$^{601}$非常危險的 但是這樣的人.
當他去到高位的時候就會出很多問題.
因為他根本從來就不是靠神在他生命裡面的工作.
所以摩西在這個時候 他講這個「我是誰」.
就是他曾經那個單純 那個的抱負.
徹底的被他自己的衝動去幻滅了.
所以他已經不敢再想自己是誰.
在這個時候的80歲的摩西問神這個問題.
就是一個充滿了謙卑,罪疚,不配.
這樣來去反問神 我是誰你知不知道?.
其實他肯定知道神知道他是誰.
因為連他自己絕對不是沒有了那個記憶.
相信他能夠很偶然間 他會真是午夜的夢回.
回想起年輕的年代 那個時候.
曾經去跟希伯來人認同.
後來錯手的打死了埃及人 後來就要這樣離開.
將自己的身份都埋沒了.
但是偶然間聖靈都會給他有一些的感動.
他會想起自己的那個抱負.
但是不敢再想下去了 夢也醒了.
因為知道自己的根梁 因為沒有超自然的介入.
他是何方神聖呢?沒有超自然的介入.
怎麼可以回去埃及呢?不是路遠.
不是不夠盤川 是根本沒有這樣的能力.
去說服自己 說服自己也容易.
說服家人 說服岳父 說服完了.
怎麼回去說服那些領袖呢?.
怎麼回去跟那些希伯來的領袖去說.
我是神差遣來的.
如果根本沒有神差遣他 他怎麼能夠回去呢?.
所以在這個時候 就是神的時間到差遣他.
所以摩西問這個「我是誰」這個推搪.
其實很坦誠的將他的內心 剖白在神的面前.
這樣的弟兄姊妹 如果已經經歷了這樣的.
枯燥 而又有一點點意志低沉.
又重新有人再問他.
咦 我看你不像是一個這樣的人生.
你應該還有意象 有夢想.
所以弟兄姊妹 我們作為一個成熟的.
在教會內的 我經常鼓勵弟兄姊妹.

$^{641}$你認領一下一些小朋友 一些年輕人.
這個認領你可以看一下一些程序表.
或者你甚至認識某一兩個 你多一點為他祈禱.
你求問神 他有什麼生命的突破改變.
或者為他的家庭祈禱.
我們很欠缺這些的教導.
我們很欠缺這些這樣的恩錫的弟兄姊妹.
我們就自己回教會 自己聚會.
聚會完之後 自己就去大家樂.
或者可以的就去會所 我們每個星期如是.
連我們自己都沒有夢想.
我們怎麼能夠去收神的夢想給下一代呢?.
如果我們有夢想的 未必是我們身體力行.
而是我們為下一代祈禱.
默默為每一個弟兄姊妹的年輕人祈禱.
為他的夢想祈禱.
這樣的時候 這個就是一個很重要的傳承.
如果我們很多人有這樣的眼光.
有這樣的眼界祈禱.
整個世界會不同的 教會也會不同的.
有些太過個人的我不能夠講了.
因為我不能夠代表我的兒子去講.
太多人知道了 那就不是那麼好.
因為他現在轉了去恩福.
他有很大的生命的突破.
不想講了 因為我們現在敢宣的年輕工作.
不是做得太好 不是童工不好.
而是很多的限制我們做得不好.
所以我們要祈禱.
所以我們要看到這些經文.
神現在在一個很落魄的摩西的心靈裡面.
收的是神那個很大的夢想給摩西.
摩西其實這句話不是代表著他有一個意興闌珊.
而是我真的不認識自己.
我真的知道我是一個怎樣的人.
我怎麼能夠做任何事情呢?.
叫我說服我的太太 西伯拉也不行.
我怎麼說服希伯來的領袖.
告訴他們神的時間.
神那個呼召的人就是我摩西呢?.

$^{681}$所以沒有可能 但是跟著神怎麼回答呢?.
第十二節 神說:我必與你同在.
你將百姓從埃及領出來之後.
你們必在這山上事奉我.
這就是我打發你去的證據.
開始已經這麼清楚的對號入座.
逐點的去拆解摩西內心深處的擔心.
那個懼怕 那個自己的無助.
摩西聽到的是神斬釘截鐵的.
我與你同在.
在很後期 民訴記 因為太遠了.
摩西聽到神說 因為在山上.
應該是出埃及的24章 對不起.
在山上 摩西上山來領受法版十誡.
神就突然間跟他說 山下已經出事了.
因為沒有你在那裡 那些百姓蠢蠢欲動.
不知道怎麼辦 因為怕那些埃及的追兵會來曠野.
在這個西賴曠野來攻擊我們.
我們還不走 還在那裡等運動.
等摩西40天都不回來 這樣他們害怕.
所以很情有可原的 行軍調將.
沒理由那個領袖走去告假.
上山上面 你上山度假還是上山什麼.
你說上山去見神 40天都不回來.
所以那些人就說 不行.
我們要做一個神來幫助我們 來領導我們.
就輾轉做了一個金牛毒出來.
不過神就這樣說 你回去.
你跟那些百姓進加拿大 我不跟你同去了.
就是這一句 摩西就害怕了.
如果你不跟我同去.
根本這個民族沒有可能有他們前面的路.
有他們的前途的 所以摩西就懇求神.
一定要跟這個民族同在 一起進去加拿大.
免得那些外面的加拿大人就會嘲笑我們.
這個應該是24章的經文 我們遲些揭露.
我們先不看這裡 第二個了.
第十三節 摩西對神說.
我到以色列那裡對他們說.
你們祖宗的神 打發我到你們這裡來.

$^{721}$因為上文已經說了 第三章的第六節.
我是你父親的神 摩西的爸爸.
這個叫做阿約基別的丈夫 就是暗蘭.
接著再上就是雅各那一代.
亞伯拉罕的神 以色列的神 雅各的神.
或者約瑟的神 所以很清楚是我們祖宗的神.
因為摩西都是希伯來人 他認祖歸宗摩西.
所以現在摩西就似乎好像過了第一關.
內心裡面有神同在這個應許.
原來摩西是代表神回去 不是一個人回去.
這個很重要 所以傳道者最重要的就是.
有神同在的伺候 有神同在的人生.
才是真正的伺候 讓人認識的是神.
所以摩西明白了第一點之後.
接著就開始有一個很具體的技術性的問題.
很抽象的 我回去 我告訴別人.
是祖宗的神差遣我來.
他們如果問他叫什麼名字 那就很奇怪.
其實摩西應該知道耶和華.
因為當時摩西還沒有寫這封書信的時間.
這個敘事的 摩西不知不知道呢?.
不奇怪 就是說在這個時候.
眾物主還沒有很透徹地將他的永恆屬性.
給亞伯拉罕,給以撒他們知道得這麼詳細.
但是現在似乎就不再是口對口的傳說.
口的口述 而是透過文字記錄.
摩西將來有機會去記下來 去認識.
對神的這個永恆的屬性 就會有更深的認識.
那叫什麼名字呢?我要對他們說什麼呢?.
哇 摩西很仔細 因為你代表的是誰呢?.
現在是你個人的一個心願 你不抵得 不甘心.
現在你在曠野 有一定的經濟實力.
有一定的人脈 有一定的理想和後盾.
那你就回去拯救我們了 到底你是.
有什麼權利來給你呢?.
這個就是權柄的問題了.
所以原來我們有這個名字 真的相當重要.
神對摩西說 我是自有永有的.
I am that I am 我就是那一位「我是」.
我就是那一位永遠存在者.

$^{761}$或者用我們的神學 我就是那一位自有永有.
就是說 什麼都沒有之前 我已經在.
我已經是那個being了 我已經是那個存在者了.
去到永恆 我都是那個存在者.
你可以想像得到 這個權柄真的相當之大.
所以基督徒,傳道人.
必須是領受這個呼召 是全時間終身的呼召.
是真的清清楚楚 神將我們分別出來.
做他的見證人 所以那個四年三年的裝備.
如果天主教是七年的裝備 都是讓我們對.
神是哪一位 有一個全面的裝備認識.
所以讀神學相當重要的 不是象牙塔.
而是造物主到底是哪一位.
因為我們現在讀的是2000年的教會歷史.
古代的猶太人的那些文獻.
還有造物主那個救贖的歷史計劃.
我們都要認識的 所以有系統神學.
有教會的歷史 又有其它的一些不同的神學.
我們都需要了解的 還有對人性要有了解的.
對於救恩我們又要全面的掌握的.
所以包括當然聖經的神學更加重要.
66卷書我們大概都要掌握那個脈絡.
所以有四年的全時間裝備是相當重要.
同一班 大家都是有那個蒙罩.
一起受裝備是很重要的那個過程.
有一些神學教育家 給我們他們的教導.
給我們的指引 都是相當重要.
而是在這個過程裡面 是對自己是誰.
還有我們為什麼要做傳道這個事件.
所謂何事 我們需要有一個理解.
我們需要內心裡面是有一個的確據.
所以這個是神的名字 這個就是第二方面.
所以我講到這裡 自有永有的神.
昨天我在福音團契和一些老人家.
因為星期四多數都是一些婦女和長者.
回到福音團契 我們昨天說就是日用的飲食.
今天就刺激我們 天天刺激我們.
就是造物主關心我們的這麼世俗.
這麼小的事 飲食他都關心.
原來主禱文既是這麼偉大,這麼浩瀚的神.

$^{801}$造物主的榮耀,名字,他的國度,他的旨意.
但是第四樣東西 主耶穌立即關心門徒.
或者往後的信徒,歷世歷代的信徒.
我們那個生活的每天的需要.
他都來叫我們禱告 去倚靠神.
這個禱告背後不是一種的懶惰.
不是一種的貪心 不是一種的揮霍.
而是一份感恩的心 一份有能力工作的人生.
在這個工作的過程 我們可以學習到刻苦.
學習到忍耐 學習到慷慨.
學習到原來恩典令到我們可以豐衣足食.
可以有生活的所需.
我們可以有餘去幫助有需要的人.
所以這裡講到的就是神對我們每個人都有了解.
不知道為什麼講著講著就講到去旅行.
就是說我們人生原來有時候可以去一下旅行.
我就跟他們講了 去旅行第一樣東西.
去到酒店 你們會做什麼的?.
他們就說check in 當然check in.
check in之後就上酒店 上房.
有些老人家說上房就敲敲門 我來了.
我差不多開玩笑的跟他們拍桌子.
我說我們這一位做物主跟我們同在.
他叫做耶和華 我們是基督徒.
我們給了錢去酒店.
除非你說怕人家裡面收拾房間的客人嚇到.
我來了 從來不會你去check in.
有一個沒收拾好房間的給你房間.
是從來不會的 除非那些很殘舊的酒店.
一般是擺好了 所以敲門代表什麼?.
代表裡面有一個更大的我們的神.
我差不多開玩笑的罵他們 不可能這樣的.
連我們自己有這麼偉大的真神.
我們自己這樣去小看.
我們要學世界的人這樣迷信.
在美房就不行了 很猛的.
沒有這樣的東西的 我們基督徒.
我們不要這樣的觀念的.
神現在告訴摩西 你要證據嗎?.
其實他們要證據嗎?.

$^{841}$你說什麼他們都不會相信你的.
因為他們根本沒有呼召 他們是希伯來人.
如果他們有呼召 他們就不需要再在工作裡面.
不斷的叫聲嘆氣 痛苦埋怨.
救贖者為什麼還不來 還不拯救我們.
現在我們很慘 就是因為他們沒有那個意象.
沒有那個夢想 所以已經是好像一個鹹魚.
日日營營役役的生活.
所以一放假就不要管那麼多 馬上飛去日本.
要放鬆了 放鬆回來又很繃緊.
其實就是沒有那個夢想 當然大家要聽清楚.
絕對不是去日本就是沒有夢想 不是這個意思.
而是已經變成一般的香港人這樣.
沒有神的那種生活.
或者已經沒有一個理想禱告的一個日常的狀況.
如果我們是有理想的 我們不會看那個是一個必須.
反過來說是一個休息是沒有問題.
但是不會將我的夢想,我的理想沖淡了.
因為很多時候就是被世俗的那個觀念.
沖淡了我們自己的夢想.
所以如果基督徒 很多基督徒是沒有了夢想.
那你怎麼影響那個世界呢?.
摩西神要的 我要證據 他們想我去展示手.
讓他們看到我是誰派來的 其實沒有用的.
因為他們只是開心一會兒.
因為當你一去治理 他們就馬上出各樣的怨言.
很快十六章已經有怨言.
就是沒有嗎哪錫就已經發脾氣了.
所以不是他們需要誰叫你.
而是摩西你需要知道我是自有永有的.
所以弟兄姊妹 這句話等於主耶穌.
對我們的日用飲食 今天就賜給我們.
你想一下 如果我們真是一個大人.
我們有一個家庭 我們要有一個負擔.
其實我們真是很感恩 能夠養妻活兒.
能夠感恩 能夠看到.
真是越來越活在一種感恩的心靈.
賺到錢 有積蓄 有富裕 能夠去慷慨.
所以主統文是教我們要學習感恩.
學習知足 學習慷慨.

$^{881}$學習用愛去幫助有需要的人.
因為一切都是神給我們的.
這個就是資本主義 這個就是新教的倫理.
基督徒是拼命的勤奮工作.
好高騖遠 不貪心.
基督徒是會努力的儲蓄 不會亂花錢.
基督徒是努力的奉獻給神.
這個就是新教的倫理 所以很重要的.
所以這一句話可以應用到不同的範疇.
所以如果真的是弟兄姊妹.
原來你自己也有這種不知覺的觀念.
不需要 你不要敲門 敲門代表什麼?.
我們去到 我插了時卡.
道門開了 我插了時卡 我就進去.
我就擺好我的行李 我就去吃飯 我就去玩.
天公地道 耶和華與我同在.
我為什麼要敲門?對不起 我麻煩了你們.
你們麻煩走開 不要搞我.
沒有這樣的事情 不需要.
我們是基督徒 自有永有的主.
就在我們的生命裡面.
所以弟兄姊妹 你真的要將這些觀念.
教其他的未信的朋友.
但是如果信了主都還有這些觀念.
其實很可憐 我們的信仰原來是這樣的.
我們的信仰原來是這麼多秘技的.
我們的信仰原來還是這麼迷信的.
我們還是很怕 我們不敢一個人睡酒店房.
其實不應該是這樣想的 應該我們要認識.
所以摩西神是給他有這個.
其實不是錦囊 是給他知道.
不僅僅我認識你是誰.
現在我將我的榮耀 我就是永生的耶和華.
與你同在 我將這個榮耀授權給你.
所以摩西其實他不是一個公民領袖.
他不是一個非常愛國的一個民族英雄.
他是一個很謙虛的 拿著一支牧羊杖的.
所以那些希伯來的領袖.
其實不知道是好笑還是很生氣.
喂 神啊 你找一個這樣的人來.

$^{921}$你應該像Rambo那樣來 現在要打仗.
你找一個老人家來拿著一支牧羊杖.
這個就是神的工人領袖.
所以大家不要看外表.
大家要看的是那個牧者傳道.
內心裡面對神的忠貞 這個是很重要的.
所以摩西 這五樣東西.
是令到摩西更加確認他自己的身份.
還有一個這麼獨特的時代 神造就這個工人.
所以一個人被神來改變.
可以帶來整個時代的轉變.
所以奉獻弟兄姊妹 你什麼都不會做不要緊.
你快點去在教會認領 認領可以是用心靈的.
不要藏著它 我要認領你 嚇死人了 不行的.
以為你拐彎抹角 你真的祈禱 神給你.
可能你真的見到某個小朋友.
我們有位姊妹在紅欣.
有時候上去三樓去洗手間.
見到一個胖胖的小朋友.
他們一家人最近來教會 自己上網找來教會.
很有趣的 現在他的太太 就是媽媽.
我們叫她盟女 二十多歲而已.
她現在開始上初訓班 她的小朋友很有趣的.
上去姊妹見到她就很開心.
其實我們要認領下一代 為他們祈禱.
有時候走過去鼓勵一下他.
有意識地給他一些禮物.
甚至乎是書籤,金句.
有意識地去栽培我們的下一代.
灌輸給他的是夢想.
給他能夠我們扶他一把.
現在神正在建立摩西.
希望摩西能夠將來 其實很快的.
都不用三個月 摩西就回到埃及了.
八十歲的老人家.
真是笑死人的 好像做電影那樣.
那個人什麼英雄領袖 原來是一個老人家.
拿著一張牧羊像的 嚇死那些領袖.
但是不要管 神的功能.
就是有他內心裡面的那團火 那個照明.

$^{961}$所以我是自有永有.
希望大家自己都真是對號入座.
在你自己的一些.
我們解決不了,衝破不了的一些我們的限制.
我們認定神是自有永有.
又說 你要對以色列人這樣說 那自有的打發.
我到你們這裡來 神又對摩西說.
你要對以色列人這樣說.
也就是說 你們祖宗的神.
就是亞伯拉罕的神 以撒的神 雅各的神.
打發我到你們這裡來.
也就是說 是我的名 直到永遠.
這也是我的紀念 直到萬代.
你去焦聚以色列的長老 對他們說.
也就是說 你們祖宗的神.
就是亞伯拉罕的神 以撒的神 雅各的神.
向我顯現說 我實在眷顧了你們.
我也看見埃及人怎樣待你們.
我也說 要將你們從埃及的困苦中領出來.
黃迦南人,黑人,阿摩利人,比利西人.
希美人,耶布斯人的地去.
就是到流奶與麥芝地.
他們必聽你的話.
你和以色列的長老要去見埃及王.
對他說 耶和華希伯來人的神.
遇見了我們 現在求你.
就是法老 用我們往曠野去走三天的路程.
為要祭祀耶和華 我們的神.
我知道須用大能的手.
埃及王也不容你們去.
我必伸手在埃及中間.
施行我一切的歧視 攻擊那地.
然後他才容你們去.
我必叫你們在埃及人眼前蒙恩.
你們去的時候 就不至於空手而去.
但各婦女必向他的淪陷.
並居住在他家裡的那些埃及的女士.
要金器,銀器和衣裳.
好給你們的兒女穿戴.
這樣你們就把埃及人的財物奪去了.

$^{1001}$哇 只是問了一個問題.
誰打發我回來.
這個似乎是神的身份.
神的授權 那個權柄的問題.
那個授權的問題.
但是現在呢 神回答了是相當具體清晰.
連那個路線圖都給了摩西.
你回去之後 你只是用我們比較俗的.
你響我的鐸 那些領袖就會聽了.
姑且信 因為都沒有辦法了.
那都當是吧 那就帶我們離開.
然後呢 法老就好像很順利.
一定會讓他們離開三天的路程.
去到一個再遠一點的地方.
其實就回到加拿大.
如果在《申命記》11章說.
由摩崖平原去加拿大.
其實只需要11天而已.
但是他們就是因為不信的緣故.
不肯進加拿大 就消磨了38年有多.
因為一年的裝備在西瀨曠野建造會幕.
用了一年的時間 足足一年.
第二年的正月初一.
這個就是《李美記經文》的開始.
聖殿那個會幕就立好了.
開始有些憲制的條例頒佈等等.
所以大概消磨了一年多.
接著就是神的審判 接著就是那些以色列人.
那一代的以色列人都要死光 才可以進加拿大.
就有38年多一點 就是這樣來到整個曠野.
其實曠野不是很大.
但是不斷的來到去兜圈.
所以這麼詳細的來到去幫助摩西去解釋.
不僅僅是將神的名字說出來.
而是有那個很清晰的藍圖給摩西.
所以摩西在這個時候有神的呼召.
就有神的工作救贖的計劃.
這樣來到去跟法老有這樣的周旋.
這個就是第二個摩西問的.
他們問你叫什麼名字.

$^{1041}$就說那個授權的問題 那個權柄的問題.
就說不是那個人要回去穿一套規格很實際的.
戴一個很漂亮的頭盔 有一枝長矛.
這樣回去打仗 做將軍 不用的.
原來一個牧羊人就是一個牧人領袖.
所以一個今天的傳道者是一個牧人的領袖.
自己先在中頂展會面前學做一個牧人.
才能夠做一個帶領的人.
所以是一個牧人的領袖.
所以我們不是用一些出類拔萃 不是錯的.
如果只是有出類拔萃.
而沒有了一個神面前的裝備.
對神對自己的認識都是不足夠的.
所以盼望神都幫助我們.
記住第二個神將他永恆的榮耀的身份.
他是自由榮耀者.
所以盼望頂展會 我們今天.
我們在21世紀 我們在神的面前.
我們不是天不怕地不怕 而是我們有神.
我們就不應該再有一些不對的錯誤的觀念.
神裝備我們 我們去依靠主.
這樣去生活 做神喜悅的事.
今天也不可以說太長了.
突然間我的喉嚨又開始乾了 不舒服.
我們先說第二部分吧 一起禱告.
因為當我們讀這些舊日的經文.
有時候我們都好像忘記了.
我們今天所信的是自有永有者.
是我們這一位造物的主宰.
是差遣你的愛子來這個世界的耶和華真神.
整個可見的宇宙和屬靈界的宇宙.
都是造物主你所創造.
是因為你的榮耀.
是因為你的命令才能夠存在.
好像保羅在《羅馬書11章》所說的.
一切都是歸於你 屬於你.
還有真是要尊榮你的 願神幫助我們.
我們這個小小的茶經團體.
能夠建立弟兄姊妹.
心靈能夠有一個很宏大的理想,夢想.

$^{1081}$為神而活 為神去得到一代的年輕人.
歸向神 為他們的下一代來尋求.
這個事奉神的心 盼望主.
興起在我們錦宣家,洪恩家,嘉欣堂的嘉欣家.
一起或者宣德的 都能夠有新一代的年輕人.
可以被主呼召,興起.
來成為下一代的牧者,下一代的見證人.
願神來幫助 我們仰望你.
這樣感恩禱告 奉耶穌基督的名 阿們.
\newpage



\section{出埃及記}
\label{sec:dj6fIzaDTMQ}
\textbf{2026-07-17 周五晚查經班《出埃及記研讀(六)舊約先知的蒙召》 講員:曾裔貴牧師}
\newline
\newline
連結: \href{https://youtube.com/watch?v=dj6fIzaDTMQ}{\texttt{https://youtube.com/watch?v=dj6fIzaDTMQ}} ~~~~ 語音日期: 2026-07-20
\newline
\newline
\hyperref[sec:cWahrBO910c]{\small{< < < PREV SERMON < < <}}
~
\hyperref[sec:index_chronic]{\small{[返順時目]}}
~
\hyperref[sec:index_scriptual]{\small{[返順卷目]}}
~
\hyperref[sec:FKTWi97mjmY]{\small{> > > NEXT SERMON > > >}}
\newline
\newline
$^{1}$很感恩再一次晚上和大家有機會去分享.
神的話語去查考《出埃及記》.
和大家先有一個同心禱告.
天父我們在晚上感恩 雖然我們停了兩次.
但我們又可以恢復在晚上的查經講道.
能夠真是將神的話語一起去重讀.
有系統的去研讀 願神透過《出埃及記》.
已經是超過四千年的作品.
讓我們能夠掌握摩西的夢召.
摩西的人生在他最後的四十年.
能夠被神的使用.
求主讓我們找到一些重要的邏輯.
或者一些重要的神的作為.
特別一些可能是可以複製的.
例如去到後期 舊約的先知都有不同的夢召.
都是聖經刻意要記載.
所以我們都盼望藉著摩西這個這麼重大.
這麼重要的人生的日子.
來到去在曠野夢召.
都使我們明白.
原來我們的主這一位自由永有的真神.
是呼召人進入你的工場.
人生這樣來服侍的.
願神透過摩西的記載.
讓我們能夠明白.
我們仰望禱告奉耶穌基督的名 阿們.
我們之前都有機會一直看到摩西.
說他如果對上兩次.
就應該是在荊棘叢當中.
其實一個很尋常 沒什麼特別的曠野.
那個不尋常的地方.
就是他見到平常本來燒著.
會燒壞的那些荊棘.
這次就是燒不壞.
他好奇心驅使.
他就靠近去看.
突然間聽到聲音.
所以呼召是很玲瓏浮凸的.
他真的能夠聽到.
造物主跟他說話.

$^{41}$告訴他摩西.
你現在見到的是這一位主宰.
你所站的就是聖地.
所以不是那個硬件本身.
因為仍然都是曠野.
不是他雲遊障礙.
反而是神的降臨.
神的同在.
摩西就成為了污穢.
他要除下了一個象徵.
他現在是不潔.
他要面對聖潔的神.
所以神也會讓他.
不會因為見到神而面對生命的威脅.
所以在這樣的夢召裡面.
他就是一直聽到神跟他的對話.
這裡有大概五部曲.
就是他不斷的去推搪.
上一次我們說到.
應該是三章十四節.
我先重新開那個檔案.
應該是提到.
摩西就說誰邀請我回去.
就是說他某程度都很細心.
會想到無端端自己走回去埃及.
就會帶來一定會有這些重要的問題.
所以他現在都為自己去做一個鋪排.
就是說你不要搞錯我不適合做那個代言人.
更加不適合做一個全民族的領袖.
帶領我自己的族民去離開埃及.
因為離開埃及是一件很大的事.
因為要對抗的是整個埃及的國度.
埃及王法老.
所以摩西其實他是過來人.
他是在王宮裡面因為各樣的原因.
就由那個公主將他養大.
成為一個養子.
所以他所認識的埃及的王國.
是一個非同小可.
非常強大的帝國.

$^{81}$所以他問這些問題.
是更加反映他本身就是一個有智慧的人.
不是衝動不是做了才算.
我們有大把的資源.
相信他是非常認識自己現在的光景.
以及又認識埃及裡面以色列國民的素質等等.
所以摩西問這些問題是非常恰當.
我想一個蒙召的人都應該會問這些問題.
應該首先要對自己有一個了解.
然後我們就要看到神到底在我們的生命裡面.
祂有什麼的印證呢.
有什麼的給我們的應許呢.
所以稍後我們又有機會.
去看看另外一個舊約的先知的蒙召.
或者一個當時的事司一個領袖.
三章的十四節我們上一次講到這裡.
神就跟摩西說.
他是誰呢.
叫什麼名字呢.
神就提到他自己的名字.
我是自有永有的.
一個非常足夠身份的認證.
自有永有的主.
打發摩西.
所以埃及無論什麼的神都不夠耶和華大.
而且是這位造物主.
差遣摩西回去.
所以我們要記得我們信的神.
是萬君之耶和華.
是自有永有的主.
所以我們就不需要懼怕.
上一次講完這裡之後.
第二天就有收到姐妹問我一些問題.
因為有個鄰居非常之有心.
送了一條開光了的鍊.
真是在廟宇裡面求回來的.
送給他的女兒.
是一份很深的一份大家鄰居的愛.
完全沒有惡意的.
但是姐妹就很擔心.

$^{121}$怎樣去不要得失.
但是又要婉轉.
所以不知道是不是剛剛聽完這裡.
我鼓勵她很正面的講.
但是講得來又不會令到對方很難堪等等.
其實很重要的.
我想我們在面對的人生.
我們都真是需要在不同的人生認定.
造物主在我們生命裡面的帶領.
既然是帶領.
有時那種奇妙.
有時真是我們現在越來越大個.
我們不是很信.
一祈禱有時真是會即刻應驗.
昨晚我出街吃飯.
那條街出了名都不用想有車位.
又不是有什麼特別重要的趕時間.
真是內心裡面祈禱.
如果可以給我有個車位.
一轉入條街道.
剛剛在我的車頭隔離.
就有個車準備上車開車走了.
其實在我前面有一架車剛剛過了準備離開的車.
即是說它已經離開了.
它不可以再退後.
反過來那個位置就給了我.
真是說時遲那時快.
兩秒鐘一轉入那個路口.
就有那個車位.
我又不是迷信.
但是我就感覺到最近有這樣的思想.
原來祈禱真是要單純.
好像我們有一個很單純的信心.
去倚靠神.
所以我們記住.
當然我們不要像求王大仙那樣.
即是說你應許我.
你成就我就會怎樣怎樣還神.
我們不需要這樣的.
眾物主是自有永有的真神.

$^{161}$即是說祂說有就有.
祂命立就立.
祂是將無有創造成萬有的.
本來就沒有這個宇宙.
本來就沒有我們人類的.
但是祂要這個世界這個宇宙出現.
祂要我們人類得到救恩.
所以我們真是在這個宇宙當中.
我們人類是唯一認識祂.
和可以親近祂.
當然就是後來需要透過耶穌的救恩.
我們才能夠和祂恢復神人的關係.
所以我們的信仰就不應該是.
好像一種說是這麼說.
怎會這麼深的經歷.
有時都是心理作用.
我們有時變成我們的信仰.
我們對這位眾物主.
有時我們根本就是不信.
這其實就不是正確的神學.
正確的神學我們既不迷信.
我們既不是追求每天都要有神蹟.
但是我們要記住.
我們是與這位眾物主同行.
同在一起的生活.
所以有時我們真是憑著一個簡單的信心.
來去依靠神.
有時會有很奇妙的事情出現.
有時我們預備講道.
突然間會有些靈感湧進來.
有時其實我上次講過一次.
我剛好看到一個廣告.
就激發我有那個靈感.
一直都想那個例子想不到.
後來輾轉那個廣告就引發我去搜尋一些事件的背後.
就有一個靈感去預備那個例子.
至少我講道的時候.
我能夠串通到我自己講道的那個經文.
所以那些靈感有時是突然間出現的.
與那件事本身無關的.

$^{201}$透過剛剛看開電視一個廣告.
就讓我能夠知道.
所以原來眾物主真是往往.
在我們意想不到的時候.
會讓我們看到他的作為.
所以摩西問這個問題.
他們問我你叫什麼名字.
這是一個很寶貴很重要的印證.
當然稍後還有三個神蹟.
摩西可以帶著回埃及的.
本來那支杖是他的.
他是一個牧羊人.
他都八十歲了.
所以那支杖本來是非常實用的.
但是那支杖變成一個可以行神蹟的.
他一扔那支杖那支杖就變蛇了.
他捉住那條尾巴就立即變回那支杖.
這是第一個神蹟.
第二個神蹟.
他的手放在衣袋立即變大麻蜂.
生大麻蜂.
然後再放回去.
然後就沒事了.
又好像正常那樣很光亮的雪白.
而第三個神蹟就是.
你在尼羅河拿些水那些是河水.
你一倒在那些乾地.
那些河水就變血了.
所以摩西帶著這三個印證.
三個神蹟.
回去埃及告訴他的國民.
到底誰差遣我來.
這些就是證據了.
這些是印證摩西這個領袖的身份.
所以他問這個問題.
其實都是我們自己要問.
所以如果今天.
我上一次都講過了.
如果你去酒店.
每次都不好意思.

$^{241}$我現在來住在你的酒店.
有怪莫怪.
你麻煩借尾.
如果你這樣.
我會很生氣的.
如果我們的弟兄姊妹.
如果到今天還停留在這些迷信.
其實我們所信的神.
比那些鬼還差.
我們要為那些鬼.
你不要這樣.
其實不需要.
我們有主.
就已經有神的同在.
我們真的不需要懼怕.
我們的主.
他是耶和華.
是自有永有的神.
所以這些金句.
不妨寫在自己的紀念冊.
寫在自己的簿子.
久不久坐飛機的時候.
很害怕的時候拿出來.
我的主是自有永有.
還用怕嗎?.
我們要記住.
第二個摩西問神的問題.
接著說.
這是神說.
又說.
你要對以色列人這樣說.
這是第十四節的下半節.
那自由的打發我到你們這裡來.
神又對摩西說.
你要對以色列人這樣說.
耶和華你們祖宗的神.
就是亞伯拉罕的神.
以色列的神.
雅各的神.
打發我到你們這裡來.

$^{281}$耶和華是我的名.
直到永遠.
這也是我的紀念.
直到萬代.
自有永有.
就是這位耶和華.
更加重要的是.
不是摩西是一個過客.
不是摩西是一個外人.
摩西是一路承接著.
亞伯拉罕的應許.
神對亞伯拉罕.
在《創世記》十五章.
說到的第四代.
亞伯拉罕的子孫.
就要興起.
就要回歸應許之地.
現在摩西.
就是接續那個應許的領袖.
所以神去到那個時候.
夠了.
就差遣摩西去成為領袖.
所以你可以說摩西.
對這件事有一個更大的保證.
因為他自己曾經有一段長時間.
是生活在埃及的文化.
埃及的養育底下.
他有時候真的會那個認受性.
不夠一個一直都在苦難裡面的同胞.
那個身份.
但是原來有不一定的.
摩西他現在.
雖然他有一段這麼長的時間.
離開了自己的同胞.
沒有受到那些做苦工的辛苦.
但是不代表他在皇宮裡面.
不是神的那個人育.
所以原來神的工作是很奇妙的.
連一個領袖都可以在埃及的皇宮裡面.
去被孕育.

$^{321}$埃及的文化.
不可以改變不了他的生命.
他裡面的DNA.
他裡面的初心.
他裡面的小時候.
可能是母親.
可能是他的姐姐米利安.
灌輸給他的希伯來的身份.
所以原來小時候所得到的生命教育.
原來可以是這麼重要的.
所以原來什麼裝備都不重要.
重要的是他真的保有了這個根.
屬靈的根.
使得摩西更加知道.
不過現在再經由神告訴他.
你不是一個離開了.
好像抽空了一個世代一個時空.
你從來都是亞伯拉罕以撒雅各的族裔.
那個英雄的兌現.
你就是我所準備.
為以色列所預備的那個救主.
那個拯救者.
就是摩西了.
所以第二個問題.
問你叫什麼呢.
摩西這樣問.
其實神給他有很充分的一個的應許.
印證.
就是神將他自己.
和他自己在亞伯拉罕時代.
應許了要做的.
現在兌現.
讓摩西有更加充分的歷史的脈絡.
原來神一直沒有放棄我們這個民族.
時間和那些苦楚.
原來只不過是讓這個民族更加成熟的一個種子.
或者是那些經歷等等.
所以摩西就是這樣來得到充分的歷史觀.
和神對他的那個啟示.
比一般人沒有知道那麼多.

$^{361}$一般人都已經忘記了自己的祖宗.
我們的祖宗怎樣好.
但是摩西所認識的祖宗.
原來神一路這樣來等待.
一路一代一代這樣過去.
以撒,雅各,約瑟.
現在到了摩西了.
到了有四十年他都要逃難.
似乎我們見不到有任何的救恩.
任何的出路.
但是現在居然一個寂寂無聞.
已經心如止水的老人家.
現在就好像見到一個很大的宏偉的屬靈的藍圖.
摩西可能年輕的時候有些內心的一個感動.
很想幫一下自己的國民.
但是現在摩西就真的得到一個非常清晰的神的救恩計劃.
就是說摩西我就是你們祖宗的神了.
我當年的應許我從來都沒有忘懷到.
一定會兌現我對這個民族.
阿巴拉罕你們的子孫那個的應許的.
所以這裡就是給摩西有一個很大的感恩的印證.
所以事奉的人是會比一般弟兄姊妹知道神的心意更多的.
事奉的人參與事奉的人或多或少你會知道更多的.
真的我們作為牧者參與在教會的弟兄姊妹的事奉,牧養.
有時比起一般的弟兄姊妹我們知道更加多的事情.
所以會使得我們更加有時看東西全面.
有時我們會更加客觀.
有時我們會有一種屬靈的眼光.
其實這些都是慢慢累積的日積月累的.
所以真是一個超過60歲以上的牧者傳道.
應該是更加有屬靈的眼光經驗.
看人看事物那個研判的客觀程度.
應該比一個剛剛出來事奉的.
或者從來都不參與事奉的一般弟兄姊妹.
肯定那個屬靈的閱歷屬靈的眼光肯定是會多.
摩西知道整個圖畫神的救贖計劃.
一直都說我是亞伯拉罕以撒雅各約瑟的神.
即是摩西從來沒有斷對這個民族的應許.
所以那些苦難其實是必須的元素.
不是神割席.

$^{401}$不是神離開了這個民族.
不是神的憤怒隨便你了不理你了.
其實就不是.
所以這就是第二個神對摩西的啟示.
第十六節我們看下去.
你去焦聚以色列的長老對他們說.
耶和華你們祖宗的神.
第二次說就是亞伯拉罕的神以撒的神雅各的神.
向我顯現說我實在眷顧了你們.
我也看見埃及人怎樣待你們.
我也說要將你們從埃及的困苦中領出來.
黃加林人黑人阿摩尼人比利西人.
希美人耶布斯人的地去.
就是到流奶與蜜之地.
第十八節他們必聽你的話.
你和以色列的長老要去見埃及王.
對他說耶和華希伯來人的神預見了我們.
現在求你用我們往曠野去走三天的路程.
為要祭祀耶和華我們的神.
又是一個敬拜.
離開了埃及離開了法老的奴役.
其實是要去到進入那個應許的神的聖地.
去事奉祭祀榮耀我們這一位的耶和華神.
所以救贖和進入那個應許對神的敬拜.
是一個套餐來的.
如果只有救贖沒有了那個生命的敬拜.
成為神國的子民的那個生命的人生的與神同在的閱歷.
那這個就是一半的救恩.
所以我們經常說.
我不知道是不是我們感受或者是雄恩.
雄恩就少了一點.
感受的獨有就是.
當然我不能夠武斷.
洗了禮之後很多人不再回主一河.
甚至連崇拜都不回.
當然他不回我們見不到他.
就不代表他沒有回崇拜.
我們希望他回其他教會的崇拜.
那是沒有問題的因為按大家的需要地區等等.
但是擔心的就是.

$^{441}$信了主很熱切洗了禮很熱心.
但是都不夠一年已經消失.
已經不再回教會.
這樣真的大有人在的.
這樣的情況不健康.
因為救贖之後是進入與神的那個身份的那個相交.
所以這裡說走三天的路程.
為要祭祀耶和華我們的神.
就是說是很清楚的.
現在埃及.
對不起.
現在以色列人還在埃及做奴隸.
摩西已經知道這個計劃了.
是救贖之後.
摩西不是說行了.
拯救了你們.
我現在回去做牧羊了.
不是.
是要帶領他們走三天的路程.
可能其實就是入迦南了.
就回去應許之地.
成為神國度的子民去敬拜神.
去事奉祭祀我們這一位造物主.
所以信耶穌其實就是得到拯救.
接著有教會的生活.
在教會裡面與不同的弟兄姊妹一起敬拜神.
第十九節.
「我知道須用大能的手」.
所以往後的第五章到十二章.
有一個題目是很清楚的.
就是神的大能拯救.
透過摩西.
透過十災.
來彰顯神的大能.
羅馬書第九章也告訴我們.
透過十災.
是要使法老的心剛硬.
是要使到神的主權是不會改變的.
神要因待誰就因待誰.
所以原來在第五到第十二章.

$^{481}$這一大段的十災事件.
原來在新約的作者特別是保羅.
他會看為是一個頗重要的.
神的主權的工作.
就是將法老作為一個示範單位.
所以我們去到那些經文再看.
因為在《出埃及記》其中.
起初就是法老的心剛硬.
因為這些這麼小的法術.
我都懂.
我的術士都懂得去扮.
去整出來.
第一個就是水變血.
那些術士都懂.
但是整整下.
那些術士就有些不懂了.
所以法老起初是他自己心硬的.
後來去到中間的時候.
是神使到法老的心硬.
那就很奇怪了.
明明他都很想悔改.
為什麼你又要令到他的心硬呢?.
所以去到那些經文我們又再去研讀.
因為有這樣的經文.
連和合本都有這樣講的.
是神使到法老的心硬.
所以這裡告訴我們.
第十九節這裡說.
埃及王一定不會容許你們去.
我必伸手在埃及中間施行我一切的騎士.
攻擊那地.
然後他才容你們去.
我必叫你們在埃及人眼前蒙恩.
你們去的時候.
就不至於空手而去.
但各婦女必向他的鄰舍.
並居住在他家裡的女人.
要金器銀器和衣裳.
好給你們的兒女穿戴.
這樣你們就把埃及人的財物奪去了.

$^{521}$都講了往後的情況.
所以摩西其實是多有那個power.
有多有那個地位.
去和法老王直接的周旋.
你想想是這樣平起平坐.
這樣來對法老.
所以背後的權柄.
一定是非常之大.
是令到摩西有這份的自信心.
有這份的勇氣.
所以麻煩我們是牧者.
我們是肅靈的領袖.
麻煩展現我們那個自信心.
在神的眼中.
我們是有這一種的自信.
不是一種驕傲.
是我們因為有神.
有聖經的真理.
我們是真的能夠.
好像保羅那樣他曾經可以去到面見阿基帕王.
都不懼怕.
其實他是一個囚犯.
阿基帕王應該是《少林寺傳》的二十六章.
大張聲勢地放了那些皇宮的宮廷的setting.
來見這個囚犯.
現在可以傳召這個囚犯來.
保羅連一點的羞恥感都沒有.
很有勇氣地來跟這個阿基帕王講.
跟他講福音.
所以令到阿基帕王都覺得.
他自己都幾乎想信.
所以這裡就告訴我們.
摩西展現的是能夠跟法老王.
來對抗對質的勇氣.
那個身份和信心.
所以弟兄姊妹我們與神同行.
我們就會不斷有一種對屬靈永恆的信心.
我們有這一種自信.
所以我們不要有一種的懦弱.
除非我們自己都不信.

$^{561}$我們自己都不信我們自己所說所讀的聖經.
其實就很可惜.
所以如果我們基督徒信主內.
展現不到一種自信心.
屬靈的眼光.
屬靈的平安.
屬天的福樂.
其實就很可憐.
因為地上的人可能比我們更多.
比我們更加有那些資源.
但是可能那些人他們沒有的.
就是我們永恆的歸宿.
對屬天的平安.
那個把持.
那個擁有.
那個確實的證據.
所以求主幫助我們.
第四章.
開始摩西又說第三樣東西.
摩西回答說.
他們必不信我.
也不聽我的話.
必說耶和華並沒有向你顯現.
耶和華對摩西說.
就是我所說的.
他們必不信我.
神就一次過將三個神蹟.
給摩西.
讓她真的有一個這樣的經歷.
所以每一個蒙召的弟兄姊妹.
真的會有一個很重要的見證.
很重要神在你生命裡面的印證.
允許使得你.
因為神蹟.
其實是為了什麼呢?.
讓你真的去到一些困難的時候.
又或者你面對一些逆境的時候.
面對一些逼迫的時候.
你就能夠再一次.
你不是思前想後.

$^{601}$你再一次回想神當初對你的呼召的印證.
所以很重要的.
如果我們沒有這些清清楚楚的印證.
我們就會有時面對很大的困難.
我們會灰心.
我們會後退.
我們就會有一種好像要退場那樣.
三個的神蹟.
第四章.
摩西回答說.
對不起 第二節.
耶和華對摩西說.
你手裡是什麼?.
她說是杖.
耶和華說.
雕在地上.
她一雕下去.
就變作蛇.
摩西便跑開.
耶和華對摩西說.
伸出手來.
拿住她的尾巴.
她必在你手中.
形變為杖.
如此.
好叫他們信耶和華.
他們祖宗的神.
又再一次說.
不厭其煩地說.
說到摩西做夢都記得.
我的祖宗就是亞伯拉罕爾薩.
雅各.
約瑟.
其實真的將屬靈的長輩.
拉回到摩西的心靈裡面.
其實都不要說這些不重要.
真的有時看看前輩的信心.
一些很出名的傳道人.
那些見證.
那些自傳.

$^{641}$其實很重要的.
因為.
真是神怎樣在他們的生命裡面帶領他們.
同樣的.
神都會在我們的時代.
我們的遭遇裡面.
帶領我們.
所以每一個傳道者.
都應該有自己的自傳.
有自己的見證.
有自己神怎樣帶領我們.
走這條事奉的路.
所以.
所以做傳道的好處就是.
你一定要有自傳.
因為太多經歷了.
將來你一定要想告訴別人.
神怎樣帶領你.
任何的牧者.
臨離世前寫的自傳.
都值得看的.
因為.
其實.
應該很簡單的.
沒有任何的驕傲.
只有恩典.
沒有任何人可以誇耀.
只有神的憐憫.
可以在這個服侍的工場裡面.
能夠事奉神.
而且是充滿了感恩.
所以.
每一本的自傳都是這麼真實.
反映神的愛和憐憫.
在一個牧者的生命裡面.
所以.
盼望大家盡量找不同的牧者的自傳來看.
真的可以得到很多幫助的.
剛才講的就是以撒瓦閣的神.
向你顯現.

$^{681}$第六節.
耶和華又對他說.
把手放在懷裡.
他就把手放在懷裡.
及至抽出來.
不料手掌了大麻風.
有說那樣白.
耶和華說.
再把手放在懷裡.
他就再把手放在懷裡.
及至從懷裡抽出來.
不料手已經復原.
與周身的肉一樣.
又說.
倘或他們不聽你的話.
也不信頭一個神跡.
他們必信第二個神跡.
這兩個神跡若都不信.
也不聽你的話.
你就從河裡取些水.
倒在旱地上.
你從河裡取的水.
必在旱地上變作血.
記住.
三個神跡.
都是為著印證摩西的身份.
而摩西的身份.
其實就是為神代言.
所以摩西的說話才是真正的重點.
不是神跡.
所以我們不要追求神跡.
而忘記了神的道.
所以靈因派.
上一次廖炳棠牧師的講道非常之好.
如果大家有興趣想聽.
其實很重要的教導.
靈因派的重點.
那個迷思.
或者他們的盲點.
就是為著要追求神跡.

$^{721}$而做很多很多的事.
所以崇尚的.
就是要超自然.
要很多很多很活潑的動作.
為著要達到那個氣氛.
他們認為的超自然的神跡等等.
這樣就變成有時為做而做.
為了要得到.
有時甚至乎不顧一切.
就變成一個很大的盲點.
三個的神跡.
一次過給摩西.
說他們不信摩西.
那就真的好像無話可說.
這個呼召來到一個地步.
是清楚到一個地步.
現在好像神要求摩西一樣.
摩西其實不應該再有任何的推搪.
摩西應該很清楚.
這個不是做夢.
不是自己現在突然間很想要發達.
我要回埃及.
不是.
他現在真的清楚.
自己仍然是一個牧羊人.
很樂樸的牧羊人.
拿著的真是那支杖.
但是他現在能夠可以行駛.
神要他行的神跡.
那支杖他什麼時候認為要扔.
作為一個見證一個神跡.
那支杖就會變成一條蛇.
後來勁到一個地步.
去到不知道第八災還是怎樣.
那支杖都不用了.
新手就已經可以了.
你可以想像得到.
原來慢慢慢慢都不需要那些物件.
你和神之間的那個關係.
就已經是一個力量了.

$^{761}$所以摩西現在就拿著這三個神跡.
第二個神跡就是那隻手放在衣袋.
衣袋就沒事的.
不關衣袋發霉有菌.
不是.
而是一個神跡.
比魔術師還厲害.
魔術師都要用東西遮住他不用.
手插進衣袋.
再拿出來就山大麻風了.
就是說如果誰不相信.
他把他的手插進衣袋等等.
任何人都會怕都會相信.
特別古代真的沒那麼多醫生.
沒那麼多的科學.
大麻風差不多是致命的.
所以摩西就能夠這樣.
你想一下.
能夠好像徘徊在污穢.
因為在利未記.
如果你有大麻風.
就要是污穢.
就要獨居在營外.
就是隔離營.
就是要去方艙.
就是古代的方艙.
原來已經有了摩西已經認識了.
但是摩西就徘徊在污穢和聖潔.
手一伸出來.
立即就聖潔.
就是說原來神可以透過摩西.
可以叫人污穢.
叫人聖潔.
所以真的.
神的僕人就有這種身份衛分.
當然不是我們行神跡.
新約已經不需要這樣了.
而是我們真的能夠用神的話語.
去勸勉人離開罪惡.
他們可以過聖潔的生活.

$^{801}$或者有些人不聽我們講的.
他自己要沉溺的.
他自己就在污穢裡面打滾.
最近我不知道有沒有講過.
最近都是因為有一些姐妹的街坊.
哀教會.
就是哀我去探望她的兒子.
幫她的兒子寫求情信.
因為是帶一些毒品.
她買一隻整隻的燒鵝來.
說你千萬不要搞這些東西.
我們作為牧者我們很願意.
後來輾轉.
上個禮拜都去了碧屋去探望這個小朋友.
16歲.
但是又是誤入歧途.
但是原來最近的報導.
原來誤入歧途的那些毒販.
去利用這些小朋友的賺快錢.
首先就跟他做朋友.
給他很多好處.
跟他又喝又吃.
接著要他做這些勾當.
就是帶一些毒品.
不用怕.
你這麼小年紀沒事的.
他們的大佬就說.
我們會找最好的律師保釋你的.
就算告到最多幾年就沒事了.
還有幾萬元的安家費.
這樣騙這些小朋友.
那次我去都很深感.
很深的內心的感觸.
因為進去那個探監的地方.
有10個窗口的一次.
全部那些.
因為那個是青少年的監獄.
全部的小朋友個個的樣子.
還白淨過我.
斯文過我.

$^{841}$那你可以想像得到.
那個問題多麼嚴重.
還有他那些探病的女朋友.
個個12,13歲的女孩子.
紋身的.
還是很幼稚很幼嫩的樣子.
Baby face的.
好像跟世界是不相同的.
這樣的環境原來每天在香港發生著.
所以現在.
我又因為搜查一些法律的那些關乎的那些東西.
說現在法庭有一個指引.
只是因為這樣的情況.
就不可以輕判.
因為那些毒販就是利用.
一般人的那個同情心.
這麼小就不應該去太重的刑罰.
其實所以法庭就定了一個量刑的指引.
不能夠輕判.
所以似乎未來這個小朋友.
我寫了求情信.
都不能夠輕判.
因為那個量都不少.
所以其實很可惜的.
所以記住.
摩西的手伸下去.
就變成大麻風.
再伸下去.
又會成為了.
聖潔乾淨.
所以其實我們作為牧者.
真是幫人進入那個真理.
他認識神.
他倚靠神.
他的生命會有改變.
他會得到那個聖潔.
他慢慢他們的生命生活都會有改善.
但是如果拒絕神的話語.
其實就有時候很可惜.
有時候我們見證著有些人墮落.

$^{881}$就會有那個生命的改變.
譬如我都講過了.
我在另外崇拜講道都講過.
我以前的舞堂舞會.
有一個女孩子.
是我們團契的妹妹.
那個時候她讀中四中五.
就開始無心向學.
曾經回來教會的.
中四到已經不想讀書.
不聽家人的話 經常出去玩.
後來就帶了一個男朋友.
那個男朋友我們一見到.
我們都帶他一段時間.
帶他一定的年紀.
我們就勸這個小姐妹.
她是信主的 但是貪玩.
怎樣都不肯聽我們講.
其實一看就知道她是一個.
吸毒的年輕人.
即是臉都瘦削得很厲害.
沒有什麼眼神等等.
但是都不理 後來真的嫁了這個.
跟這個男孩結了婚.
後來就真的發現.
她嫁了之後.
她真的不知道什麼原因.
可能失業等等.
就再一次吃白粉.
所以這個姐妹就很不開心.
也都很後悔.
但是輾轉都沒辦法了.
因為都嫁了.
後來還生了個兒子.
那個兒子 後來她.
都有帶她的兒子回來教會.
包括她自己都有回來教會.
但是丈夫仍然沉淪在毒海.
後來輾轉都沒辦法了.
都跟這個丈夫離婚.

$^{921}$離了婚.
後來她的兒子.
不知道四歲還是五歲左右.
就進了瑪嘉烈.
她的姐姐叫我快點去探望他.
還有另外的傳道人.
因為她的兒子患了血癌.
在瑪嘉烈離世.
離世之前是順了主.
很清楚順了主的小朋友.
四五歲左右.
很乖的.
這個姐妹你想一下.
這麼人生的坎坷.
真是很慘.
她真的很後悔.
當年為什麼不聽姐姐.
不聽教會的弟兄姊妹說.
嫁了這個男朋友.
於是過了一段很長時間.
她當然當時都離了婚.
由她自己一個人照顧這個兒子.
後來就離世了.
你可以想像到一個姐妹.
看著自己的兒子血癌離世.
是非常痛心的事.
過了很長的時間.
我想有二十年.
我都不是這幾年.
已經再久一點.
在元朗見到這個姐妹.
嫁了一位很好的弟兄.
在教會有新的生活開始.
不是在我的母會.
在元朗另外一間教會.
我不說名字.
是非常之好.
她整個人的生命改變.
有神給她一個新的婚姻.
新的開始.

$^{961}$她很感恩.
那個弟兄很敦厚.
很好的一位弟兄.
雖然我不認識那個弟兄.
但見到他們兩個.
是很好的一對弟兄姊妹.
去服侍.
去教會等等.
人生是有不同的階段.
我們作為牧者.
我們有時會看到這些人生不同的事情.
因為看著這個姐妹.
她是在教會.
小時候長大.
中四就開始學壞.
但原來我們牧者有時真的.
擠一隻手下去就大麻風.
擠一隻手出來又會變成聖潔.
所以我們都徘徊在.
見證著不同的人的生命改變.
我想做人的生命.
有時那種痛心.
困難.
眼淚.
禱告.
但有時又會有感恩.
又會為神的工作來讚歎.
所以我們記得.
摩西都有三個神蹟印證他.
所以我們作為牧者.
我們要信我們自己的道.
我們自己的真理.
我們去幫人.
用心去幫.
神總會有他的時間.
尤其是我們認為看不到果效.
我們都不要追求那個外表的果效.
我們忠心做好自己的工作.
這是第二個神蹟.
第三個就是水變血.

$^{1001}$尼羅河就是埃及人看為是.
明脈之河.
生命之河.
很proud of他們自己的尼羅河.
如果我們有一些很出色的名山大村.
等等.
我們去誇耀都無可厚非.
所以埃及人很誇耀他們自己這條生命的河.
孕育他們這個民族的生命.
但是在摩西的眼中.
我可以令到這條河的河水變血.
什麼生命都沒有.
所以原來我們牧者都是有這種.
將神告訴人.
他不信.
神就會令到他的生命.
有時遭遇不同的打擊.
但是我們將人帶進去那個永恆.
我們會見證著有些人信到主.
而離世.
他會有很大的平安.
和心靈的福樂.
這就是第三個.
摩西他們不會相信我的神.
神就說我給三個印證.
三個像錦囊一樣.
你帶著回埃及吧.
還有什麼可以講呢.
還有什麼可以推呢.
真的火滾了.
所以神要追殺摩西.
第四個.
第十節.
摩西對耶和華說.
主啊 我數日不是能言的人.
就是從你對僕人說話以後.
向他顯現.
不知道Encounter這個對話.
持續了多久.
是半天還是三個小時.

$^{1041}$不知道.
但是摩西就說.
由你向我對話到現在.
我都不是馬上變成能言善辯了.
固然是.
也是這樣.
我本是絕口不舌的.
耶和華對他說.
誰造人的口呢.
誰使到人口啞耳聾.
目明眼瞎呢.
豈不是我耶和華嗎.
現在去吧.
我必賜你口才.
指教你所當說的話.
這也是一個非常重要的恩典的印證.
不是一次過全給你.
好像你立即按個鍵.
你就變成能言善辯.
其實往往都是不斷經過了不同的成長.
不同的學習.
不斷的在勤奮.
在失敗裡面.
慢慢練就一份經驗的累積.
來到成長的.
所以參與事奉或者服侍主的弟兄姊妹.
都應該是要這樣來到體會.
一切都是從神那裡而來.
所以何時神要將那個恩賜給你.
適合你這樣去用呢.
那就會給的了.
所以我們的責任就是.
應該做一個海綿.
這樣掏空自己.
讓神來到去幫助我們.
這個掏空不是坐在那裡等.
而是我們都有我們自己的準備.
有我們自己的訓練.
有我們自己的學習.
尤其是虛心的學習.

$^{1081}$這樣來到去成長.
我有個同學.
他說他在他一段時間有空.
神樂園的同學.
他說他看了五百篇傳道人的講章.
當然是學習人家的思路.
怎樣掌握那個起承轉合.
但他看了五百篇.
就是很想看看這些傳道人.
怎樣去預備講章的時候.
怎樣去拿捏那些經文.
這些訓練學習是很重要的.
就看人家怎樣預備.
自己可以怎樣能夠從其中.
參詳到一些智慧.
或者一些巧妙的技巧.
這些都是很重要.
所以這裡也都不等同於.
我們不相信神是可以.
一直之間改變我們的生命.
但是神往往喜歡從一個自然的演進.
你慢慢能夠可以一路一路這樣去掌握.
越來越建立你在事奉裡面的那個信心.
那個自信心.
那個力量等等.
如果你說往往都要一個神童12歲站在港台.
真的能夠可以叫每一個人都心悅誠服.
其實這樣都還是一個不靚的圖畫.
有什麼理由一個12歲的小朋友.
可以叫下面五百的人.
來向你佩服呢.
有什麼意思呢.
應該就是說一個不斷事奉.
不斷面對不同的困難.
很多的限制.
但是仍然很勉勵地來服侍.
開始多白頭髮多皺紋.
這樣會不會令到其他的人.
更加的信服神的恩典.
在一些軟弱的人身上的作為嗎.

$^{1121}$如果一個人他能夠隨隨便便站起來.
就可以講到滔滔紅遍的.
不用怎麼預備的.
他可以隨時講幾個國家的語言的.
但是對一些軟弱的弟兄姊妹的回眾.
其實幫助不是很大.
引起的共鳴感動不算很大.
但是如果我們看著.
這樣的牧者一路這樣來成長.
一路這樣在很多的困難.
很多的掣肘.
很多的限制底下.
都仍然表達出這一份.
對主的甘心樂意的奉獻.
那個的心靈.
這樣帶來的就是.
可以感動很多人.
所以摩西他自己.
開始他看到自己的.
絕口笨舌.
對了 因為八十歲了.
因為四十年牧羊了.
見羊多過見人了.
在曠野了.
了無成就了.
不知道他有沒有一些童年的弟兄姊妹.
在埃及現在在奴隸的行列.
有沒有再通信呢.
等等我們沒有任何的資料.
但是摩西現在有的.
就是他絕口笨舌.
是一個理由.
更加當然是一個推搪.
因為已經第四個對話了.
這個推搪就針對到他自己的自身了.
真的盤點一下.
不要說笑了.
你給我幾個神蹟.
你這樣演一下.
我帶領的是200萬人.

$^{1161}$面對的是法老.
這樣的level的對手.
怎麼講得到呢.
我怎麼可能有這樣的能力呢.
相信摩西仍然掌握他自己的文化.
那個學養.
那些就不會消失的.
因為他真的學到成長了.
但是不能夠再很純熟的運用.
這也是事實的.
但是無論如何.
神就說一個挑戰了.
誰做人的呢.
神就講得很詳細.
誰使人口啞.
耳聾.
目明.
眼瞎.
全部相關性的.
同一個人的能夠講話的震動.
講到話.
有邏輯.
有思想.
有智慧的講話.
不是就這樣啞啞啞啞.
個個都講到.
但是講出來是有point.
有智慧.
有說服力.
有power.
去議價的.
這個不只是表面的絕口笨舌.
而是真的他能夠成為一個屬靈的領袖.
獨當一面的領袖.
這個摩西當然.
做夢都不會再這樣做夢了.
因為已經心如止水了.
但是現在神就要讓他有這樣的夢.
神就要讓他.
你擔心口才.

$^{1201}$那些不用擔心了.
將來我給你有更高的智慧.
更高的那個屬靈的領導力.
更高的那種高瞻遠矚.
甚至乎你成為那個民族的一個.
差不多國父的一個人物.
甚至乎在尼泊爾山.
與神同在這樣離世的.
這樣的一個摩西.
你可以想像得到.
我們永遠都不會知道.
神為我們預備將來人生的歸宿.
但是有一樣東西肯定的就是.
事奉神的人一定會經歷許多的艱難.
許多的難處.
因為這些艱難的難處.
使得我們更成熟.
使得我們更加具備屬靈的領袖的自信心.
屬靈的領導力.
以及那種真正能夠走完人生的路.
充滿的是一種屬靈的智慧.
而不是一種的委屈.
一種的不是內疚.
或者一種的對世界對人的那種憤怒.
其實最可惜的.
有時見到那些老人家.
最可惜就是老人家已經一來時日無多.
二來行動又這麼慢.
但是他最敏捷的.
居然是他腦裡面的那些仇恨.
那些執著是最可惜的.
所以一個屬靈的事奉神的人.
你會看到一個真是給神用的人.
他充滿的是感恩,平安.
不再有執著,不再有憤怒.
不再對人心懷仇恨.
是一種不解除的仇恨.
這個就真是根本沒有資格進入天堂.
天堂就沒有這些東西.
天堂是完全沒有仇恨的.

$^{1241}$所以如果我們屬主的兒女.
弟兄姊妹真的.
千萬不要有仇恨.
不要有那種驕傲.
不要有一種自私.
不要有一種你自己的那種執著.
我們真是要這樣去帶著這樣的平安去見主.
所以摩西又再一次無話可說.
神說我做這些一切的恩賜.
我現在要令到你站在一個領袖的角色.
你的絕口笨舌.
不是突然間好像三個月學會外語那樣.
而是他能夠回到其實是他的身份.
與神不斷接觸的身份.
讓他有自信心.
讓他有代表神的那種.
不可以說是威嚴.
那種屬靈的眼光.
屬靈的可能是磁場.
可能是屬靈的膽量.
因為對抗魔,對抗法老就要膽量.
對抗那些攻擊他的以色列領袖.
因為遲些在《民數記》第15章就有了.
可拉的後裔攻擊摩西.
摩西有的就是屬靈的內心的忍耐等等.
所以這裡就再一次告訴他.
整個人的結構,身份,人生的成長.
神都已經幫你計劃.
所以傳道者因為要參與事奉.
就一定會經歷了這些一般人沒有那麼多經歷的甜酸苦辣.
這樣的高低起伏和悲歡離合.
我們經常見面突然間又要去探病,垂危.
上星期我又漏了.
又是講泊車.
博愛就出名沒得泊車.
每次去都要排兩個位.
很麻煩.
現在的位越來越少.
就算你告訴他我急症都不行.
我又祈禱.

$^{1281}$因為我的童宮的姑仔的老爺就危殆.
因為真的醫生說沒辦法了.
因為肺癌已經擴散了.
年紀很大.
他又是沒有生存的意志等等.
因為突然間聽到童宮說可能今晚已經不行了.
下班馬上開車去.
又祈禱.
剛好只剩下一個位.
去到馬上拍到.
馬上衝上去.
祈完禱之後大概過了一個多小時.
他們家人都在.
那時候都未需要.
不試探病時間.
我都可以進去.
告訴他們是牧師.
我們去祈禱.
祈完禱大概一個多小時.
老人家離世.
都安詳.
但我就不在了.
很奇妙.
我從來未去過博愛.
可以馬上去到馬上有位.
但那次都不理那麼多.
拍下去才算.
果然馬上真的有位.
你會見到神有時有這些奇妙的安排.
但摩西他說你不要找我.
我是絕口笨舌.
的確我們應該有兩層的心靈.
我們記得我們不要有一種驕傲自大.
其實事奉多久你就知道.
你就會慢慢神就將你這些驕傲自大.
磨到平了.
不會再有了.
但有些人沒有.
但仍然沒有那個自信.
沒有那個事奉的自信心.

$^{1321}$其實這樣是因為我們對自己的身份.
未能夠充分的了解.
如果我們能夠了解.
其實是神揀選我們.
我們代表神去講神的道.
所以那個自信不是因為我膨脹.
我自我的膨脹.
或者我因為是牧師不是這樣.
而是神已經印證了我們被他選照.
這樣的自信.
所以這裡就是摩西自己都要學.
所以原來神都很顧念他的僕人.
那種的自卑.
所以都是由自卑一直去成長.
所以雅各書一章.
卑微的弟兄升高就該喜樂.
一章的八節.
富足的慢慢的降卑也該喜樂.
因為富足的人都會過去.
好像那些卑微的人都會過去.
太陽出來熱風刮起美容就會消沒.
富足人都會是這樣的.
所以如果富足在他人生能夠降卑.
看到一切都是神的祝福,恩典.
就會很不同了.
所以摩西都真的要面對自己.
最終其實是自己的心態.
自己願不願意接受這個差遣.
其實可以不接受的.
神揀過另外一個.
但是似乎在舊約的時候.
神的工作是讓一個不心甘情願的人.
慢慢都看到神的印證.
是千真萬確,是玲瓏浮凸.
所以我們在現代的蒙召傳道人.
我想我們都需要有這些印證.
其實是給我們整個人.
面對很大的事奉逆境的時候.
仍然不要忘記.
當初真的神感動,印證.

$^{1361}$走上這條路不是為了一種虛榮.
就會令到那個牧者在困難的裡面.
他又能夠有那個信心,那個力量.
最後講完第五個了.
對不起我先喝一口水.
有點乾 因為講得太多.
好 麻煩你.
第十二節 現在去吧.
我必賜你口才.
只教你所當說的話.
神真的要發怒了.
第十三節 摩西說.
主啊 你願意打發誰就打發誰去吧.
第十四節.
所以第五個推唐.
神就真的要來攻擊他了.
耶和華向摩西發怒說.
不是有你的哥哥利未人阿倫嗎?.
我知道他是能言的.
現在他出來迎接你.
這裡好像有一個聖經問題.
到底阿倫是一直住在埃及.
他們能夠大家有相交,見面.
還是阿倫一直跟隨摩西在米顛礦業生活呢?.
聖經沒有詳細講阿倫住在哪裡.
但是他現在出來.
似乎神又感動他找自己的弟弟.
似乎不是第一次的相認.
所以又很特別的.
聖經又沒有詳細解釋.
阿倫到底他在哪裡長大呢?.
是一直陪伴著自己的埃及同胞.
以色列人這樣長大.
還是他住在礦業呢?.
這裡是沒有提到的.
然後就說.
他出來迎接你.
一見你心裡就歡喜.
你要將當說的話傳給他.
我又要賜你和他口才.

$^{1401}$又要指教你們所當行的事.
他要替你對百姓說話.
你要以他當作口.
他要以你當作神.
你手裡要拿者杖.
好行神蹟.
第五個的印證.
第五個推搪的神秘的.
就是有一個現成的親哥哥阿倫.
將來就是他一個小小的阻攔.
就是做金牛毒的事件.
出埃及記三十二章的時候.
聽從百姓的呼喊.
那個納迪摩西因為上山.
在西瀨山上面同神同在四十日.
做約版的時候.
百姓就在山下來攻擊阿倫.
要他快些帶我們回去埃及.
因為帶我們離開埃及的那個摩西.
記住那個摩西已經不知在哪裡.
是否見難處就自己先走了.
他們已經忘記了.
摩西其實是上了西瀨山.
所以阿倫就是幫他們做了金牛毒.
因為阿倫不具備那個領袖的尊榮和力量.
和那個領袖的身份角色.
所以阿倫只能夠成為一個輔助的領袖.
摩西不應該再能夠推搪.
也可以這樣說.
很快摩西就會見到阿倫的出現.
印證神對他的應許.
暫時你還未能夠具備.
獨當一面的能言善辯.
未有這樣的智慧.
未有這樣的膽量.
我給你一個陪伴.
給你一個左右手.
就是你的哥哥阿倫.
他會代替你去向百姓說話.
後來摩西不用了.

$^{1441}$摩西後來已經絕對自己能夠獨當一面了.
所以又是一個漸進式的成長.
對於一個事奉的人.
我想我們今天講到這裡.
因為喉嚨都有點乾.
和有點不舒服.
我們暫時去到四章的十七節.
下一次我們再看神要追殺摩西.
因為他沒有幫他的兩個兒子行割禮這件事.
我們先停在這裡.
一起同心禱告.
多謝天父帶領我們今天晚上.
其實我們用了比較少的時間.
去講摩西的夢召.
這五個推搪.
帶來給我們看到.
在舊約裡面.
其實都是成為一個模式.
後來的先知夢召.
都似乎是這樣.
對自己對神的認識有欠缺.
但是神又有那種憐憫.
使這些夢召的先知.
無論耶利米提到自己是年幼.
自己是膽怯等等.
譬如基奠.
又是這樣的去求印證等等.
或者是以賽亞.
或者是撒母耳等等.
都是這樣的.
好像模仿摩西夢召的模式.
所以我們看到.
神透過這麼詳盡的呼召過程.
都讓我們後世參與事奉的弟兄姊妹.
都必須要認識.
我們所跟隨的主.
你是自有永有的神.
是耶和華.
是使人有智慧有口才的神.
所以求主幫助我們.

$^{1481}$讓我們在人生都能夠依靠聖靈.
因為聖靈就像保羅在哥林多前書所講.
十二章所講.
都是這位聖靈隨自己的意思.
將恩賜來分給不同的人.
所以求主幫助.
我們感恩讓我們比較詳細.
去將摩西的夢召.
去彼此的分析去分享.
我們多謝主.
願你感動弟兄姊妹.
參與在無論是一些教會有的崗位.
或者在外面不同的地方的事奉崗位.
大家都能夠忠心的去服侍主.
願你得到一切的榮耀.
這樣感恩禱告.
奉耶穌基督的名.
阿們.
\newpage



\section{出埃及記}
\label{sec:FKTWi97mjmY}
\textbf{2026-07-31 周五晚查經班《出埃及記研讀(七)舊約先知的蒙召(二)》 講員:曾裔貴牧師}
\newline
\newline
連結: \href{https://youtube.com/watch?v=FKTWi97mjmY}{\texttt{https://youtube.com/watch?v=FKTWi97mjmY}} ~~~~ 語音日期: 2026-08-03
\newline
\newline
\hyperref[sec:dj6fIzaDTMQ]{\small{< < < PREV SERMON < < <}}
~
\hyperref[sec:index_chronic]{\small{[返順時目]}}
~
\hyperref[sec:index_scriptual]{\small{[返順卷目]}}
~
\hyperref[sec:IlLnn7WOgiE]{\small{> > > NEXT SERMON > > >}}
\newline
\newline
$^{1}$各位姐妹我們一起同心禱告.
天父我們雖然停了兩次.
但是我們今晚又能夠再一次.
將《出埃及記》,摩西的夢召.
和他回埃及來帶領自己的族民.
來建立神的國度.
我們很盼望透過這些我們可能都熟悉的經文.
能夠深入去摩西的內心.
在神的救贖計劃.
在神對伯拉罕的應許.
我們一直這樣去研讀的時候.
也都給我們看到有一個模式.
就是在舊約的先知夢召.
都好像和他們的身份要有一個確認.
有一個很重要的勾劃.
使我們能夠知道原來摩西.
是一個多重要的劃時代的人物.
求神帶領我們.
我們都在每一個時代.
神都是呼召你自己的工人.
願神來去施恩.
繼續這樣來去帶領.
我們有大的,有小的事奉.
或者有些應該是比較明顯的.
有顯露的一些崗位.
也都有一些可以低調地在背後.
來去無論禱告.
或者可能是預備一些材料等等.
都同樣是一個在神眼中尊貴的服侍.
願神都帶領著我們大家.
因為當我們進入到出埃及記.
當他們成為了神的國度.
離開了埃及之後.
他們就有大量的工作要來參與.
無論是建造,回望,預備的材料.
或者他們要撰寫他們自己的律法書.
這些全部都是在抗疫的時候準備的.
願神都給我們看到.
神是一位有計劃.
是有真是很奇妙帶領的神.

$^{41}$今天晚上我們再一次聚集在你的面前.
這一輪都是非常之不穩定的天氣.
特別我們都仍然紀念到.
在無論是委內瑞拉.
其實已經沒有人再關心裡面的居民.
因為經過一個大地震.
超過5000人以上離世.
現在在日本,在紅本縣.
甚至現在可能福岡都可能有問題.
願神都紀念每一個地方的人.
他們的困苦,他們向你的禱告.
特別日本人他們很難去信主的.
盼望神真的來讓一些宣教士.
盡快趕到去災區來拯救幫助裡面的人.
也都讓日本的居民.
他們曉得向你禱告,願主施恩.
我們這樣仰望奉耶穌基督的名.
阿們.
好 我們沒有辦法.
上一個禮拜因為有些不方便.
所以停了一次不是很好意思.
我們應該對上一次講到.
摩西有三個印證.
就是神給他一個很清楚的confirmation.
就是他手拿著的那支杖可以變蛇.
另外就是他的手山大麻風.
這些是不是一種掩眼法呢?.
我們去到《出埃及記》的第七章.
那些術士都有這樣做到的.
我們可以遲些再看.
當然水變血都要讓他們看到.
生命之河 埃及人的生命之河.
神都可以使到他變成為血.
就是說什麼都沒有之餘.
簡直是一個讓他們看到.
如果不是依靠神 什麼都沒有了.
我們繼續下去.
重要的題目就是本來上個禮拜就已經要講的.
我只是比較快選取了.
機電作為那個對比.

$^{81}$其實我上一次都提過.
其實無論耶利米 以塞瓦.
甚至乎薩姆爾.
特別是在他開始蒙召的時候.
都是跟他們那個身份的那個.
對自己的認識.
有一個很重要的基礎.
所以這裡就讓我們一起思想一下.
有一個很共同的模式.
甚至乎有些解經家認為.
摩西就是一個劃時代的開始.
他的呼召就是印證著日後在不同的時代.
當然就是在講王國的時代.
有先知的興起.
有君王大衛王的出現.
還有當時的領袖.
就算是時事記.
頭幾個時事記都比較.
比如這個俄台涅.
包括甚至乎這個參孫.
都是跟他們那個身份的認證有關的.
所以多重要是這個我們要留心.
所以我們就從這個角度.
我們去看一看.
參考一下摩西和基典的那個蒙召.
那個模式都很相似的.
第一時間好像是君權神授那樣.
人就受不了了.
因為太大的一個召命.
所以一定是第一時間推遲的.
這個推遲是差不多每一個先知每一個時代.
都是聽到而沒有立刻反應.
都是一個恩物相承的.
所以我們就在其中都能夠找到一些人性的端倪.
人的推遲.
神怎樣來回答幫助他.
和神怎樣用他自己的屬性.
就是他的名字.
或者他行一些神蹟.
所以我們就看一看.

$^{121}$當然最後就是那個印證的方法.
摩西就有三個的印證.
基典大家都應該相當熟的了.
有些很戲劇性的羊毛的乾和濕.
我們可以看一看這個時事記.
應該就是吳張在十五二之前.
對不起 手太乾了.
那些聖經因為比較薄.
對不起 應該第六章才對.
時事記第六章.
我們從這四個的類別去看一看.
基典怎樣能夠和摩西的夢召是相似的.
我們看看.
以色列人又行耶和華眼中看為惡的事.
耶和華就把他們交在米顛人手裡七年.
米顛人壓制以色列人.
以色列人因為米顛人就在山中挖洞.
挖掘挖洞建造營寨.
就是被逼上山.
好像以前中國的古代.
那些少數民族被強逼上山上面.
山上面就是說那個一點都不是富庶.
第三節 以色列人每逢撒種之後.
米顛人 阿瑪利人和東方人都上來攻打他們.
對著他們安營 毀壞土產.
直到加薩沒有給以色列人留下食物.
牛羊 驢也沒有留下.
就是說非常之困難.
那個壓迫好像去到胸口那樣.
第五節 因為那些人帶著生畜 象棚來.
象蝗蟲那樣多.
人和駱駝無數.
都進入國內毀壞全地.
第六節 以色列人因為米顛人的緣故.
極其的窮乏.
就呼求耶和華.
這個呼求就是一脈相承.
就是我們在《出埃及記》第三章尾.
以色列人因受苦工而哀求耶和華.
所以這個哀求都是一個很重要的指標.

$^{161}$所以人的困難去向神呼求.
這個都是一個很重要的動力.
當然在現代.
特別我們現在去到一個無神的後現代.
我們多困難都寧願自毀.
都不寧願去依靠神.
去向造物主呼求.
所以這個呼求都需要有對神的認識.
如果根本看這個世界是沒有東西的.
物質世界以外是沒有東西的.
是虛無的.
這樣他就不會呼求.
所以原來猶太人他們多困難.
始終有一個很強烈的信仰在他們的內心.
所以特別香港總有人聽過耶穌.
總是看他去到什麼困難的時候.
聖靈有工作.
所以都奉獻我們每一位.
每一次我們的分享,講道,查經.
甚至唱詩歌.
我們都要帶著信心.
總是有些種子會落到他的心田.
他什麼時候發芽.
什麼時候向神呼求.
這個都不是我們能夠掌握.
但是聖靈就可以掌握.
所以現在米甸人這樣的壓迫.
以色列人就集體去呼求耶和華.
就是說爛船都有三根釘.
他們的信仰去到這麼凹凸不平.
拜偶像,拜迦南的神.
但是去到這樣的壓迫的時候.
他們窮極就呼天了.
他們沒有東西可以依靠.
他們又真的來將這一位祖宗的神.
再一次來去呼求.
所以當他去呼求的時候.
第八節.
耶和華就差遣先知到以色列人那裡.
對他們說.

$^{201}$耶和華以色列的神如此說.
我曾領你們從埃及上來.
出了遺牢之家.
救你們脫離埃及人的手.
並脫離一切欺壓你們之人的手.
把他們從你們面前趕出.
將他們的地賜給你們.
留意第九節這麼詳細.
第八第九節不是說機電.
是說神差遣先知.
在這個時候撒母耳還沒有出現.
甚至還沒有出生.
到底這個先知是誰呢?.
其實在第二章也有.
有神人來勸這個俄陀列.
到底這些人是很超然的什麼人呢?.
為什麼沒有先知時代開始.
因為撒母耳是第一個正式的先知時代.
甚至撒母耳是先知學校的校長.
因為他有神學院.
他有那些先知門徒.
就是神學院的前身.
但是現在在事先的時候.
是說那些民族的英雄.
或者一些審判官.
但是為什麼現在會有先知的出現呢?.
這些就是神秘了.
我們有時候都沒有留意.
原來在那個時候.
就算沒有先知的職系.
但是神都可以個別感動一些人.
但是這些是無名氏.
是無名先知.
但是他的說話就是代表著神.
其實某程度上是給他們一個警告.
譴責他們.
他們不聽神的話.
但是有什麼的.
第十一節.
耶和華的使者.

$^{241}$這裡到底是不是上文的先知呢?.
還是有另外一個人呢?.
這個我們又不是太知道了.
耶和華的使者.
還是天使呢?.
去了俄弗拉這個地方.
坐在亞比爾藉族的約阿斯的橡樹下面.
約阿斯的兒子基奠正在酒鑿那裡打麥子.
說要防備米顛人.
在酒鑿裡面打麥子.
當然不是釀酒.
因為吃都不能吃.
所以都挺艱難的.
通常你釀酒.
就是說你本身大豐收.
你才會這樣去做那個飲品.
一般都不會.
相信是那個基本的收成的麥子.
而且是要防備.
不是在地方外面大庭廣眾.
因為那些麥子是要那些坑.
那些坑坯要有風吹散它.
因為古代沒有那麼多機器.
都是人手去打那些和等等.
所以這裡說要防備米顛人.
第十二節.
耶和華的使者向基奠顯現對他說.
大能的勇士.
耶和華與你同在.
基奠說.
主啊.
耶和華若與我們同在.
我們何至遭遇這一切事呢.
我們的列祖不是向我們說.
耶和華領我們從埃及上來嗎.
他知道有列祖的口傳的歷史.
猶太人雖然沒有摩西五經.
但是.
對不起.
已經有了.

$^{281}$在《士師記》的時候.
已經有摩西五經.
有律法書的了.
有十誡的了.
也有一些在百姓當中的一些教導律法的人.
因為利未人就是教導律法的.
所以.
他們知道這些歷史.
這些背景.
他們現在的確已經在迦南地.
不過呢.
迦南地不是他們自己來當家作主.
而是周邊有很多的外族.
來攻擊以色列人的.
當然不斷地說那個原因.
就是他們好像一個螺旋一樣.
不斷地墮落向下.
所以每一次他們去到最極致的時候.
就向神呼求了.
一呼求.
神就有工作了.
就是讓他們有一個悔改的轉機.
這個悔改往往在《士師記》就是興起拯救者.
所以這是一個螺旋式的拯救的故事的記載.
一個敘事的文體.
這裡就是基奠這一位拯救者.
這裡告訴我們.
基奠和這位使者對話.
他真的很清楚.
本來我們是應該榮耀的.
有神同在.
代表我們應該是豐衣足食.
五穀豐收.
那些老人家可以坐在那些無花果樹下.
北京客等等.
《怡唱日書》的那些光景.
是這樣描述的.
現在就不會了.
晚上不會坐在那些大樹下乘涼.
但是古代或者在農村.

$^{321}$你能夠安居樂業.
坐在樹下.
那時候沒有那麼多娛樂.
也不會有YouTube.
也不會有上網.
你能夠坐在樹蔭下面.
有些人講故事.
大家很開心地聚集聊天.
笑完一晚就回家睡覺.
這些就是舊約時代那種和平的光景.
所以尤其是沒有威嚇.
不需要懼怕.
夜不抵護這樣的光景.
就是代表著那個民族是非常有神的恩典和同在.
所以這裡你可以看到他們已經有一個很大的戒備心.
戒備心就是說有一個很大的恐懼.
因為外族不斷地侵害.
這個侵害似乎不是持續一兩次.
是持續以年來計算的.
所以和當下的那個光景有些不同.
所以第十四節.
去從米顛人手裡拯救以色列人.
不是我差遣你去的魔.
第十五節基甸說主啊我又何能拯救以色列人呢.
那些措辭你就會感受到.
和摩西的那個模式是很有關聯的.
摩西見到神 懼怕 脫鞋 接著就開始對話.
開始就感懷身世.
他現在自己的光景.
怎麼能夠有那個力量或者那個號召力.
去拯救自己的百姓在這些外族人的入侵呢.
所以第一時間就是那個人的推辭.
這個就是基甸所說的.
基甸說第十五節.
主啊我又何能拯救以色列人呢.
我家在馬來西亞支派中是最貧窮的.
我在我父家是最微小的.
他就是馬來西亞支派.
所以我們就記住.
士師記告訴我們.

$^{361}$任何一個支派都有可能是史詩.
不只是猶大支派.
不只是以法連支派.
或者弱勢的支派等等.
其實馬來西亞和以法連就是弱勢的分支.
是他們兩個的兒子.
但是這裡就讓我們見到.
現在基甸對自己的身世是非常的了解.
其實某程度上是一個托詞的推搪.
因為神現在欽點他.
你仗著我跟你同在.
這份能力十四節所說的.
你將以色列人從米顛人的手中拯救出來.
第十六節我們又見到這些對話.
好像神一直像剝洋蔥一樣.
慢慢去讓被呼召的人.
來慢慢去重拾信心.
對神的倚靠.
所以兩樣東西其實都沒有衝突.
對自己越來越認識.
對神越來越信靠.
這個其實不是50/50.
是百分之一百和百分之一百.
充分對自己的認識.
但又充分對神有一個倚靠.
相信一個真是事奉神的人.
就會經歷這種百分之一百和百分之一百.
又不是叫交戰.
而是一種融合.
挺奇妙的.
不事奉就沒有這些困難.
不用推遲.
因為根本你都不稀罕事奉.
何來推遲呢.
你又不需要經歷神的同在帶給你的力量.
所以大家自己去想.
不參與事奉.
一直從來都沒有蒙召.
從來都沒有感動.
這個才是問題.

$^{401}$這裡讓我們看到神如何再次安慰他.
第十六節.
耶和華對他說.
我與你同在.
你就必擊打米顛人.
如擊打一人一樣.
基澱說.
我若在你眼前蒙恩.
求你給我一個證據.
使我知道與我說話的就是主.
求你不要離開這裡.
等我歸回.
將禮物帶來.
供在你面前.
即是獻供物.
主說我必等你回來.
似乎是耶和華的使者.
這裡有少許隱含著.
就是天使.
同在帶來的就是給基澱.
好像有少許心裡面的一個感動.
雖然一直都自己說不行.
但是既然這樣的時候.
你給我一個證據.
讓我知道我真正與神的使者對話.
其實這裡已經不是一個個人的主觀的幻覺.
第十九節.
基澱去預備了一隻山羊羔.
用一衣發細麵.
即是麵粉.
不是那些公仔麵.
是很幼細的.
磨到最幼的那些麵粉.
大家應該磨過那些麵粉.
作了無窖餅.
將肉放在糠內.
因為山羊羔.
即是山羊羔的那些肉.
把湯盛在壺中.
帶到橡樹下.

$^{441}$獻在使者面前.
神的使者吩咐基澱說.
將肉和無窖餅放在這盤石上.
把湯倒出來.
他就這樣行了.
耶和華的使者伸出手內的橡.
橡頭挨了肉和無窖餅.
就有火從盤石中出來.
好像自己會燃燒一樣.
神可以用各種的形式.
去彰顯祂的同在.
祂的榮耀.
這裡是說那個呼召.
神對那個時代的一個特殊的呼召.
這些是特殊性的呼召.
這個時候.
橡頭挨了之後火從.
接著就燒盡了肉.
和無窖餅.
耶和華的使者就不見了.
基澱看到是耶和華的使者.
他就害怕了.
哀哉主耶和華.
我不好了.
因為我這個「獨」字.
孤獨的「獨」.
那個「獨」音.
我獨面看見耶和華的使者.
即是兜口兜面.
面對面的意思.
23節.
耶和華對祂說.
說明剛才已經升上去了.
不在了.
祂現在突然間這樣祈禱.
這樣去懼怕.
神就馬上回答他.
你放心不要懼怕.
你必不致死.
第24節.

$^{481}$於是基澱在那裡為耶和華竹鳥一座壇.
起名叫耶和華沙龍.
沙龍就是Salem.
S-A-L-E-M.
希伯來文.
即是平安的意思.
賜平安的意思.
所以耶和華沙龍.
就是耶和華賜平安的意思.
所以幫小朋友改名叫Salem也挺好的.
沙龍.
第10.
接著.
祭壇在阿比爾祭族的俄弗拉直到如今.
就是舍斯基的時代.
這裡我們稍微停一停.
人的推遲.
因為實在太大了.
要做一個拯救者.
即是說.
現在在馬來西亞派當中.
找到這個基地.
他說自己在馬來西亞派是最微小的一個族.
在這個族裡面.
他自己也是最微小的.
這樣來自稱.
這樣來貶低自己.
但到底在舍斯的時代.
神如何去揀選這些舍斯呢.
是沒有一個絕對的公式的.
有些是很厲害的俄陀列.
因為家底幾豐厚的.
因為他的太太的父親.
突然間講到口唇邊.
大家稍等.
這個就是俄陀列.
這裡對他的形容.
搞錯了.
對的沒錯.
我一直都想說.

$^{521}$就是加納的女婿.
這個俄陀列.
所以.
每個的時詩.
我們是估計不到神在哪一個民族.
即是哪一個支派.
來去興起.
和那些的時詩.
有什麼個人之處.
我們都不能夠來去做一個.
好像一個公式.
這樣來去等待的.
神要興起的.
我們有時覺得匪夷所思的.
這個是時詩給我們的.
不是叫隨機性.
而是神如果要用的人.
有時他有多少的限制.
神都可以遮蓋他的限制.
而他能夠成為一個時代的一個事司.
一個領袖.
我們回來看看基典這裡的時候.
一直都是講那個呼召.
神的回答.
這裡其中強調的就是.
神的同在.
神的英文.
加力量給他.
就是匠著神的那個能力.
這裡是他自己在上面說的.
你匠著我給你的能力.
你就能夠從那些人的手裡面.
拯救神的子民.
以色列人.
接著第25節.
當那耶和華吩咐基典說.
你取你父親的牛來.
就是那七歲的第二隻牛.
並拆毀你父親為巴力所築的壇.
似乎當時大家家家戶戶都要拜巴力.

$^{561}$可能是一個護身符.
可能是一個避免米顛人對他們的攻擊.
所以原來基典的爸爸.
都似乎是一個族長.
帶頭去做這些拜巴力的事.
所以現在神就要求基典做一個名正言順的宣言.
就是神的拯救臨到.
不需要再屈就於這些偶像.
巴力這些事情上面.
所以你留意到巴力.
一直是以色列人的.
好像一條刺一樣.
很多代都是這樣來攻擊以色列人.
一直有時會好些.
有時以色列人會強些.
但最終都是不徹底的.
一直到我們知道的列王記尚.
去到以利亞的時候.
真是在這個加密山上面.
和這些巴力的先知.
那些假先知來對抗的.
現在基典已經要面對他們拜巴力.
拜巴力也是以色列人的一個死穴.
因為是需要有雨水.
又需要有豐收.
豐收不單止是牛羊那些田野.
而是生育.
對一個古代的民族來說.
生兒養女.
特別是生兒子.
還有就是一定要有雨水.
所以巴力就是種這兩樣東西.
生育和豐收的雨水.
所以你看到以利亞就笑那些巴力先知.
在那裡求雨.
又用刀插自己.
以利亞說不用了.
挖個洞.
叫那些火焚下來.
燒光那些水.

$^{601}$叫那些人在基純河拿些水來.
倒滿那些挖了的溝.
那些好像地道那些.
那些不是石油不是火水.
是真的水.
在基純河那裡拿的.
諷刺巴力的先知.
不單止不可以求他們的神去降雨.
反而以利亞就能夠了.
這個聶王紀上的十七章的經文.
這裡告訴我們.
要做這件事其實可大可小.
因為是在說那個支派.
馬來西亞支派的亞比爾族.
這樣來到一個神壇.
一個在廟裡面的.
一個不是私人的.
自己個人的供奉.
這個是帶來可能整個族裡的人.
都會攻擊基典的.
所以神這個吩咐.
好像是否要給他膽量.
要做一件這麼匪夷所思的事呢.
還是要很撇脫的跟爸爸拜偶像.
這個錯誤的勾當.
一刀兩斷 兩者可能都有的.
等下我們看回摩西你就看到了.
摩西神要追殺他.
因為他兩個兒子還沒有行割禮.
這裡就讓我們看到第26節.
在這盆石上整整齊齊的.
為耶和華你的神築一座壇.
將第二隻牛獻為梵祭.
用你所撼下的木偶作柴.
整個祭壇都要撥爛了它.
基典就從他僕人中挑了十個人.
照著耶和華吩咐他的行了.
他因怕父家和本城的人.
不敢再白奏行者事.
就在夜間行了.

$^{641}$即是說基典很清楚自己這一次這樣做.
即是說他踏前了一步.
願意去試一下.
是不是神真的要呼召我做先知.
應該說做事師.
來拯救我的民族.
脫離米顛人的奴役.
所以一樣也是這樣.
那個原因也是要百姓歸向神.
所以任何一個時代的傳道人.
牧者,主學老師,導師.
都是要帶領人去到神的恩典裡面.
所以我們自己是要先經歷恩典.
先經歷神自己的呼召帶來的洗禮.
所以是蠻近的.
沒有說時間.
如果夠膽就日間做.
光明磊落.
但是也真的不容易.
所以夜晚做.
你不要以為基典是很害怕,很軟弱,很懼怕.
其實也是一個智慧.
即是低調地做了.
但是始終也是要公諸於世.
始終也是要面對那些族民會追究這件事.
所以後著你會見到.
是他的父親反而給他有信心.
所以也不是這麼簡單.
所以任何的呼召都是心靈的忐忑.
所以我們真的要尊重終身蒙召的神的僕人.
他們真的走過一條決心走上服事主.
將人帶去神恩典的路.
我們就要留心.
人家真的在一個很大的考驗底下.
決心踏上這條信心之路.
所以你看到基典,強如摩西.
或者是耶利米,以賽亞等等.
他們可以繼續在他們自己的安書裡面.
沒有所謂的.
但是往往那個真的不能磨滅的心靈的呼召.

$^{681}$當真是確認之後.
就毅然要走上這條路.
所以大家看看有沒有機會.
真的看到如果我們自己教會能夠出到神學生.
真是決心讀神學.
其實我們要暗中代討.
在明處多一些鼓勵,拍拍肩膀等等.
給他有一些信心的支持等等.
這些是很重要的.
所以這裡的時候.
你想想要得罪爸爸固然已經不是簡單了.
再要得罪整個自己的支派的人.
更加是不容易.
還有原來當時也挺分化的.
所以基典日後要.
事先就是要使整個以色列再一次在一個螺旋的下滑的情況下.
慢慢地回到神那裡.
所以你要將一群人的心慢慢地扭轉.
其實一點都不容易.
所以這裡神就給他好像小試牛刀那樣.
其實這個事件也有很大的衝擊.
第28節.
「果然,城裡的人清早起來,見巴力的壇拆毀,壇旁的木偶撼下,.
第二隻牛獻在新竹的壇上,(新竹的就是對耶和華的敬拜的壇).
彼此說:「這是是誰作的?」.
他們訪查之後就說:.
「這是約阿斯的兒子基典作的.」.
城裡的人對約阿斯說:.
「將你兒子交出來,好致死他.」.
即是已經拜巴力了,但是不單止沒有覺得慚愧.
反而還要對付這個基典.
「因為他拆毀了巴力的壇,撼下了壇旁的木偶.」.
約阿斯回答:.
「站著攻擊他的眾人說:.
「你們是為巴力爭論嗎?你們要救他嗎?誰為他爭論?.
趁早將誰致死?巴力如果是神,有人拆毀他的壇,.
讓他為自己爭論吧.」.
即是說叫巴力來顯靈等等.
所以當日人稱基典為耶路巴力.
意思說:他拆毀巴力的壇,讓巴力與他爭論.

$^{721}$耶路似乎是爭論的意思.
反而幫基典改了一個外邦假神的名字.
他與巴力爭論,這挺特別的.
後來這個耶路巴力又變成一個花名.
所以是基典一次的呼召.
再加上他要走上信心之路.
每個信心之路都不同,有大有小.
例如摩西真的要走回埃及.
跟埃及的領袖說:我是希伯來人.
我現在不只是迎祖歸宗.
而是神叫我回來帶領你們離開埃及.
與法老有一個周旋,議價.
基典的夢召就是在這樣的憲制上.
首先他自己要踏出唯有敬拜耶和華神.
自己的信心,奉獻要很清楚.
當時他們拜巴力已經很普遍.
也深植在以色列人的心裡.
因為真的沒辦法,神那裡又不顯靈.
給我們有足夠的供應.
唯有因為生活艱難,只能夠來去敬拜巴力.
即是心底裡面可能都拜神的.
但實際的行動就拜巴力.
希望我們頂錢會不會說:近來真的不夠用.
現在金多寶二千萬,都不要管了.
那些人自己走進去投注站.
買電腦票,買多少錢,戴包圍.
我想我們不需要這樣依靠這些.
但是有時候很難說的.
有時候去到一個窮途末路.
心靈沒有神的那種依靠.
有時候有些人就會沒什麼心情.
有時候喜歡一些迷信的東西.
譬如一些體相,術數等等這些東西.
其實有時候正正反映到我們內心.
神的影響力已經很減低.
所以都求神幫助.
如果我們一直內心已經沒有神對我們的生命的影響.
即是我們裡面如果有一個度數的能量.
其實去到谷底真的更加要回教會.
正正因為這樣更加要回教會.

$^{761}$更加要再一次被神的愛充滿自己的心靈.
這個真的要求神幫助.
這裡告訴我們第三十三節.
那時米顛人,阿瑪力人和東方人都聚集過河.
在耶斯列平原安營.
一個很大的平地.
耶和華的靈降在基甸身上.
舊約就有這些特殊性的.
因為它已經有了這個身份.
有了神已經確認了它.
就是這個時代的馬來西亞派的士師.
也是當時的以色列人的一個士師,一個領袖.
所以神的靈就很自然降在這個基甸的身上.
所以這個某程度上都是一個很重要的指標.
他就吹角阿比爾藉族都聚集跟隨他.
他打發人走遍馬來西地.
馬來西人也聚集跟隨他.
又打發人去見阿切,西布倫,拿忽他利.
他們也都出來與他們會合.
所以大家其實要留意一下.
如果你熟悉河東和河西.
他們有六個支派的屠城.
和他們那個十二支派是怎樣分佈呢?.
其實如果我們稍微有一點點脈絡.
我們就會大概知道.
因為一定是跟馬來西他們那邊的人.
相近就會有一個支援.
但是以法連就不支援.
第36節基甸對神說.
你若果照著所說的話.
藉我手拯救以色列人.
我就把一團羊毛放在禾場上.
若單是羊毛上有露水.
別的地方都是乾的.
我就知道你必照著所說的話.
藉我手拯救以色列人.
次日清晨.
基甸起來.
見果然是這樣.
將羊毛擠一擠.

$^{801}$從羊毛中拿出滿盆的露水.
基甸又對神說.
求你不要向我發怒.
我再說這一次.
讓我將羊毛再試一次.
但願羊毛是乾的.
別的地方都有露水.
這夜神也如此一行.
獨羊毛上是乾的 別的地方都有露水.
這樣來反轉測試.
無話可說.
是對還是不對.
聖經的作者沒有給我們一個評價.
好像照著一路經文的發展.
基甸似乎給我們感覺到.
有一點柔弱的信心.
但仍然被神用.
在這樣的非常時期.
神仍然用基甸.
所以不是要有全足的信心.
神才會用.
不是要等我們自己完全裝備好.
神才會用我們.
這裡讓我們見到.
有時我們在按內心的感動.
合適聖經的真理.
我們踏上其實已經是回應神的呼召.
所以很難說我見到那個終點.
我才開始起行.
永遠這樣的人.
其實就停滯在他自己的安書裡.
不能踏上那個呼召.
所以我們是走前一步.
又見到神的那個供應.
和那個開路.
更重要就是.
因為蒙召一定是說那個印證.
證明我們真是神帶我們去做一些.
祂要我們去做的工作.
所以這裡就讓我們見到.

$^{841}$所以我就不再揭下面了.
下面就是在說.
他確認自己的身份之後.
就正正式式以全個民族的.
事司領袖的身份.
發動和米顛人的戰爭.
那裡就不需要再詳細說了.
反而就是這裡印證著.
摩西的蒙召.
原來在神要做工作的時期.
他就會興起不同的他的工人.
一定會有他的訓練.
有他的加力.
有他的印證.
這樣在那個過程裡面.
慢慢那個被選召呼召的人.
慢慢會由那個信心的增加.
能力的增加.
恩賜的提升.
這樣來慢慢去到更加具備.
人都真是印證.
這個真是神的僕人.
所以我們又透過這些.
又給我們自己有一個這樣的信心.
如果我們真是有一個渴慕服侍神的心.
史圖保羅在提摩太后書.
人如果有一個服侍的心.
是罕無越立的.
所以如果我們有這個事奉的心智.
神是會不斷的來更新我們的心態.
心靈.
他會加力給我們.
尤其是他要我們做的特定的工作.
他會給我們有力量.
我們回過頭來看摩西這裡.
我們剛才提到的就是.
多重要耶和華的名字.
耶和華的同在.
是最重要的印證.
所以我相信事奉神的人.

$^{881}$一生要不斷感受到神同在.
對神的屬性的認識.
那個依靠.
就不再是慢慢將神.
我們好像不需要.
就依靠自己的恩賜.
自己的人脈.
自己的經驗.
本來這些全部不應該抵觸的.
但是有時候就是不知不覺.
那個恥笑比掌.
神的同在.
有時候我們未必這麼著緊.
其實不是的.
那個才是真正最重要的力量.
當然不是說我們什麼都不用管.
我們就站上台就可以了.
就不是的.
我們的準備.
我們的勤奮.
都一定是.
相輔相成一定要具備的.
所以第五章回來出埃及記.
我們就正式看回這一段.
我們上一次都有機會看一點點的.
第五章.
後來摩西亞倫去對法老說.
耶和華以色列的神這樣說.
容我的百姓去.
在曠野向我守節.
法老說.
耶和華是誰.
使我聽他的話.
容以色列人去呢.
我不認識耶和華.
也不容以色列人去.
他們說.
希伯來人的神預見了我們.
求你容我們往曠野去.
走三天的路程.

$^{921}$祭祀耶和華我們的神.
免得他用瘟疫.
東冰攻擊我們.
埃及王對他們說.
就是法老.
摩西亞倫.
你們為什麼叫百姓抗攻呢.
你們去擔你們的擔子吧.
第四章的結尾.
以色列人那些領袖.
就非常之感恩.
因為出現了摩西.
當然還有亞倫.
所以他們就.
看到那些神蹟.
這個是剛才我們還沒讀的四章三十節.
亞倫將耶和華對摩西所說的一切話.
述說了一遍.
因為他是代言.
摩西一直都覺得自己是.
口才應對不好.
所以接著就叫摩西在他們面前.
在百姓的面前.
行了那些神蹟.
就是象變蛇.
手長大麻風和復原.
接著拿一點水上來.
水一倒在乾地.
就變成血.
三個神蹟去印證.
所以第三十一節.
四章三十一節.
百姓就信了.
以色列人聽見耶和華眷顧他們.
眷顧和監察都是同一個字根.
所以你想想眷顧,顧念和監察.
你可以想像得到.
他們用這些字眼.
就是說神看到實我們.
百姓他們沒有經歷過蒙召的考驗.

$^{961}$摩西都要和神爭論.
摩西都要有一個交戰的過程.
但是百姓看到外在神蹟的印證.
不是自己親身的個人成長的經歷.
所以是很虛的.
所以往往我們回來.
如果沒有參與任何事奉的弟兄姊妹.
就會很舒服,很享受.
你不需要付什麼代價.
意味著那些道也不會很深刻.
停留在你的心靈裡面.
如果你不回去之後.
認真的去揣摩,去經驗,去實踐.
但是講道的人或者事奉的人.
他能夠講得出.
是因為他過去曾經有一些的際遇,經歷.
他也都依靠神.
所以兩者是不同的.
所以他們拜服摩西是一個暫時性的.
因為現在是有一個很大的一個承諾.
我們帶你們回去.
立刻行三個神蹟.
就是說他們期望有更多的神蹟.
有更好的日子過了.
不用做奴隸了.
但是他們不知道就是有很多艱難要面對的.
是摩西來將這個身份再一次重新賦予給以色列人.
他們就想著是摩西自己去經歷.
摩西幫他們搞定了.
我們就堂而皇之的進去.
離開埃及.
就回到迦南地神應許的土地.
他們想著就是這樣進去的.
所以你看到一個信耶穌洗禮.
如果沒有經過什麼困難.
他的信仰根本沒有經過考驗.
所以往往事奉的人不斷看到不同的難處.
不同的困境.
他在那些困境上面克服.
他就會很深刻神的同在.

$^{1001}$神的供應和恩典.
所以由第五章開始.
其實很快.
當摩西一回來跟他們說.
法老不准.
還更加給更多的困難那些百姓的時候.
百姓立即就反轉頭攻擊摩西了.
說摩西你這樣搞到我們更慘.
中雞毛鴨血.
我們以前都叫做有覺好訓有得食.
因為都需要有力去做那些泥磚.
但現在還加重了那個苦力.
就是你們兩個搞出來的.
你們兩個兄弟.
所以你會看到沒有自己去面對那些困難.
因為神給的恩典絕對不是無痛的直接享受.
是在困難裡面來不斷地成長的.
所以小行傳應該十四章.
我們進入神的國.
必須經歷許多艱難.
是跟當時在似乎在勞斯德那裡受逼迫的信徒說的.
所以大家都很明白的.
所以我們在這個世界.
我們信了主.
我們以為一定是一帆風順.
有困難來.
有這樣的阻滯.
我們就開始懷疑我們的信仰.
其實根本不應該這樣的.
我漏了一段就是第十九節.
四章的十九節.
我忘記了講這個西波拉那個事件.
十八節開始.
摩西回到他岳父當時的祭司葉蒂羅那裡.
就是一些民間的祭司.
對他說求你容我回去.
見我在埃及的弟兄.
看他們還在不在.
葉蒂羅對摩西說.
你可以平安回去.

$^{1041}$第十九節.
耶和華在米甸對摩西說.
你要回埃及去.
因為尋殺你命的人都死了.
因為他殺了一個埃及人.
第二十節.
摩西就帶著妻子和兩個兒子.
叫他們騎上路回埃及地去.
摩西手裡拿著神的像.
耶和華對摩西說.
你回到埃及的時候.
要留意將我指示你的一切騎士.
行在法老面前.
但我要使他的心剛硬.
這裡在未來的這幾章不斷出現法老的心剛硬.
有時是神令到法老的心剛硬.
有時是法老自己.
見到那些災禍減退.
他自己就心硬.
所以不全部是神使到法老的心剛硬.
嚴謹地說.
他本來從來都是剛硬的.
他必不容百姓去.
你要對法老說.
耶和華這樣說.
以色列是我的兒子.
我的長子.
我對你說過.
容我的兒子去.
好事奉我.
你還是不肯容他去.
看啊.
我要殺你的長子.
都相當清楚這個呼召的內容.
他要做的.
他就要在有生之年.
帶領他的族民.
不單止離開埃及.
要成為神的兒子.
要去到神應許的地方.

$^{1081}$去敬拜事奉神.
徹徹底底.
和埃及無分無關.
要進去迦南.
所以一個這麼大的民族.
摩西要帶領這一幫人.
其實一點都不容易.
摩西其實也很清楚.
百姓的心靈.
他們的無論是思前想後的記憶.
雖然是奴隸.
但是都會有一個穩定的.
在埃及不會說不知道怎麼辦.
但是有一段時間.
離開了埃及這個保障之後.
在曠野.
如果沒有適當時間的物質供應.
百姓就會起哄.
就會反對摩西.
這些情況果然就.
去到《民數記》13章就已經出現了.
本來就是開始進去窺探.
迦南地帶了那些葡萄回來.
要有兩個人抬的.
但是10個探子.
除了迦納和約書亞.
10個探子都是當時的領袖.
那些支派的領袖.
當然是年輕的20歲以外.
應該是40歲以外.
回來報惡訊.
說裡面的那些阿納族人.
好像巨人一樣.
我們好像是蚱蜢一樣.
就是說我們不行的.
我們沒理由這麼沒有保障的進去.
所以這樣的訊息.
其實這個訊息本身不是一個惡訊.
但是百姓聽到.
就有一個很大的.

$^{1121}$對他們來說很難逾越的一個關.
所以就好像兩邊不討好.
唯一的就是遷怒於摩西.
搞成這樣.
其實已經過了紅海.
其實埃及的那些追兵已經淹死了.
是神蹟來的.
那些海水合起來了.
但是他們去到曠野.
真的沒事情做 太空閒了.
每天就在討論我們前面應該怎樣.
我們之前以往.
我們在埃及的那些日子是怎樣的呢?.
現在又炎熱.
又不知道什麼時候能進到迦南.
那怎樣去解決我們那些民生的需要呢?.
所以摩西其實自己一直越來越清楚.
這一份的事奉.
是一個吃力不討好的工作.
但是寶貴的地方就是.
每一個清楚走上這條事奉路的神的工人.
不同時代都是的.
正正是因為曾經推遲過.
神有他的確認.
神有他的恩典.
神有他的栽培.
所以往往都會走到.
一路事奉裡面.
在困難裡面成長.
所以每一個傳道人.
都是在困難裡面成為一個成長的牧者.
任何一個都是的.
當然我們都要很明白.
有弟兄姊妹有牧者.
是中途會離任.
所謂離任就是不再做傳道.
做回自己的一些他認為可以事奉的.
我起初有時候很抗拒.
這些我認為是便捷的.
但是事奉久了.

$^{1161}$真的體諒和明白.
走不到未必是我們要怪責.
反而我們要體諒.
他能夠在另外的場景.
不同的場合.
我認識的都有好幾位是做了保險.
在他的保險界裡面繼續為主作見證.
這些有時候真的要諒解.
有些就找我聊天.
有時有空就聊 有時就沒有.
但是真的要明白.
就算在神學院一起畢業.
有些都是不再擔任牧者的崗位.
起初我真的抗拒.
為什麼不走到尾 為什麼不依靠神.
但是真的每個人有不同的心靈.
所以其實任何一個夢召.
神從來沒有騙我們的.
其實我們是在困難裡面.
其實神帶我們進入困難.
在困難裡面去成長.
所以其實在困難裡面.
才會使我們更加認識神.
認識自己和認識這個呼召的寶貴.
如果不夢召 不事奉.
或者我們從來都沒有想過要事奉神.
我們就沒有這些困難.
因為我們自己所謂免疫了這些困難.
我們就少了在困難裡面經歷神.
而我們的生命成長的這些寶貴的福氣.
所以其實沒有一個傳道人想拿苦來受的.
是沒有的 肯定不會有的.
但是在困難裡面.
其實根本那個工作的性質已經預備了.
因為你做帶領人的工作一定是有這些問題.
所以其實摩西就是典型的一個.
好像它都不是神學院不要了.
是一個很大的夢召.
神已經預備了它.
曠野成為一個很好的訓練場.

$^{1201}$現在好像它所有的東西.
那些雄心壯志都已經完全磨滅.
現在有的那團火.
就是神燒不壞的.
它內心的對自己的族民的感動.
越來越感動的.
現在這一刻就很激氣的.
因為百姓見好就好了 摩西好了.
一有那些困難就馬上遷怒摩西的.
但是我們看看這裡本來講得很好的.
本來講得很清楚你的人生.
就是要帶領你的族民對抗這個法老.
帶他們去到一個地方敬拜神的.
但是24節就有一段很特別的經文.
都是聖經難提.
摩西在路上住宿的地方.
就是正在回埃及的.
耶和華遇見他 想要殺他.
西波拉就拿一塊火石.
割下他兒子的羊皮.
就是割泥.
雕在摩西腳前.
說你真是我的血囊了.
我們不是很明白什麼叫血囊.
這樣耶和華才放了他 他就是摩西.
西波拉說你因割泥就是血囊了.
到底是什麼意思呢?.
是不是摩西其實都知道自己是希伯來人.
因為神與亞伯拉罕立約.
世世代代他們的子孫都要南丁行割泥.
神與亞伯拉罕立約就是要行割泥.
神與摩西的立約就是守安息日.
所以大家不要搞錯了.
所以摩西可能還沒有幫他兩個兒子行割泥.
照經文的表面就是這個意思.
所以摩西幾乎變成一個血囊.
可能他害死兩個兒子.
其實有可能就是神既然要摩西回去.
都讓摩西意識到你要名正言順.
你的家庭都是跟隨神的 按神的旨意.

$^{1241}$如果所有的以色列人他們做奴隸都好.
都要很敬畏神 無論如何.
我們是亞伯拉罕的子孫.
我們就算在埃及我們都要幫我們的兒子行割泥.
所以摩西如果沒有做到.
可能就是很生氣 離開了埃及.
在曠野沒有人理會.
我們是不是希伯來人都無所謂.
所以他已經錯過了幫他兩個兒子行割泥.
但是神就好像對這件事其實不是耿耿於懷.
而是幫助他將來回去不會被那些叔伯兄弟抽籤他.
去說你自己都其身不正了 摩西.
我們個個都行割泥 你兩個兒子就不行割泥.
大家記得如果在十六章的《小徒行傳》.
保羅要帶提摩泰去傳道.
因為知道他爸爸是愛邦人 希臘人.
他媽媽是猶太人 就幫提摩泰行割泥.
但是如果對提多 保羅就沒有幫他行割泥.
你可以想像得到 因為要將他帶在以色列人當中.
提摩泰他小時候 爸爸就一定是說了算.
媽媽就沒有權力了.
所以提摩泰就沒有意識到自己需要行割泥.
但是他現在出來事奉.
他主要的工場會見到的就是猶太人.
所以這個半堂翻提摩泰.
保羅幫他行割泥.
就是說要讓他們領袖那些信猶太教的.
摩西的律法的 他們是行到足夠.
這樣的一種人情世故.
你可以想像到保羅都會對這件事.
對神的道的尊重 因為他們是猶太人.
他們的氛圍都是這樣.
所以這裡摩西反而自己.
因為已經40年住在米甸曠野.
已經沒有了那些意識了.
所以神又提醒他.
不是那個割泥有什麼特別.
而是因為你服侍的群體有這樣的背景.
所以摩西你都要有一種敏銳的心靈.
來去看回這件事.

$^{1281}$所以有時候這些就涉及到.
我們對那個群體的了解,尊重.
我們就能夠不要碰釘.
這裡都不只是單純是禮儀這麼簡單.
而是摩西都要真的以身作則.
可能他自己就行了割泥.
他兩個兒子就沒有行割泥.
他就去在神的山遇見摩西.
耶和華對亞倫說.
這個我們上一次都講過了.
因為摩西自己有一種這樣的自我的貶低.
但是將來其實他不用了.
亞倫反而是左手左腳.
所以原來都有一個過程.
慢慢越來越那個信心.
那個能夠慢慢因為神的恩典而獨當一面.
這些真是每一個傳道人的生命的寫照.
我當然不能夠說我自己獨當一面.
不過神真的讓我自己慢慢慢慢這樣摸索.
有很多的前輩.
無論偷師,模仿或者直接的對話.
我都講過很多次.
張牧在樓上他在我隔離房.
經常叫我過去跟他聊天.
談很多教會核心的一些問題.
受益良多 而且看到他.
因為很近 看到他怎樣去想一個議案.
怎樣部署一些事工.
他是怎樣去做 怎樣跟別人去談.
自己耳喧目染地看到很多.
很實在看到神的牧者.
他是怎樣來策劃.
怎樣跟不同的人去交往.
去討論那些事情.
那些領袖等等.
很多有這樣的機會.
神給我的一個這樣的恩典.
所以真是我的一個啟蒙的老師,恩師.
所以我自己現在是一個雄恩.
都真是這樣來秉承著很多的教誨各方面.

$^{1321}$有些真是受益良多的.
所以無論現在是人手的問題或者怎樣.
其實原來你回想三十多年.
原來慢慢真是好像摩西那樣.
有時需要別人幫你.
但慢慢摩西不用了.
連那支杖都不用經常變了.
連不需要伸杖舉手了.
根本他自己就是那個力量的來源.
因為跟神已經非常靠近了.
所以我想我們事奉應該都是這樣.
應該都是越來越在神的恩典裡面.
有一種靠神而能夠獨當一面的事奉.
獨當一面不是一種驕傲.
是真是慢慢倚靠神.
越來越有一種神給我們的那個慰憤帶來的信心.
不如就停在這裡.
我們講到第五章後面那裡開始.
正式就是跟八老的對抗和周旋.
後面那裡已經差不多是處理了的問題.
我們就停在這裡一起祈禱.
再一次我們將感恩歸給神.
和這位造物的主配受我們的讚美.
我們雖然在講一些已經發生了三千多年的.
摩西的夢召.
摩西的對抗法老.
和這個民族到今天他們仍然是神的選民.
主耶穌已經在猶太人當中降生.
來到他生在以色列.
來到他是瑪利亞的兒子.
已經成為我們的救贖主.
和已經升天坐在父神寶座右邊的主.
但當我們再看回舊約的時候.
其實都有很多信心功課我們可以學習.
但是摩西我們就一直可以有機會.
看他怎樣面對困難.
但是他始終不放棄這個身份.
神的僕人的身份.
帶領他們的民族成就一個時代的見證.
更加他有機會直接去領受神的約版.

$^{1361}$十誡和十誡所衍生的那些條例.
願你幫助我們.
在這個時代興起更多神的工人.
明天開始又是每年的華人教會的培靈會.
盼望神使用好幾位的牧者.
大大使用他們將神的話語去講出來.
去祝福弟兄姊妹.
使到不同時代不同年齡的弟兄姊妹.
都得到神的恩典.
阿們.
\newpage

\allsectionsfont{\centering}

\setlength\parindent{0pt}
\setlength{\columnsep}{1.25em}
\setlength{\parfillskip}{0pt}
\setlength{\tabcolsep}{1em}
\raggedbottom

\pagenumbering{gobble}


\newfontfamily\leftfont[Path=../fonts/fell_french_canon/, Ligatures=TeX, ItalicFont=IMFeFCit29C.otf, BoldFont=AveriaLibre-Bold.ttf]{IMFeFCrm29C.otf}
\newfontfamily\leftcitationfont[Path=../fonts/frankruehl/]{FrankRuehlCLM-Medium.ttf}
\newfontfamily\centerfont[Path=../fonts/garamond/, Ligatures=TeX, ItalicFont=EBGaramond-SemiBoldItalic.ttf]{EBGaramond-SemiBold.ttf}
\newfontfamily\rightfont[Path=../fonts/averia/, Ligatures=TeX, ItalicFont=AveriaLibre-RegularItalic.ttf, BoldFont=AveriaLibre-Bold.ttf, BoldItalicFont=AveriaLibre-BoldItalic.ttf]{AveriaLibre-Light.ttf}
\newfontfamily\rightcitationfont[Path=../fonts/rashi/]{Mekorot-Rashi.ttf}
\definecolor{hcolor}{HTML}{D3230C}
\definecolor{rcolor}{HTML}{D36F0C}
\newcommand{\chfont}[1]{\centerfont{\huge\textcolor{hcolor}{#1}}}
\newcommand{\leftcitation}[1]{\leftcitationfont{\Large\textcolor{hcolor}{#1}}}
\newcommand{\rightcitation}[1]{\rightcitationfont{\normalsize\textcolor{rcolor}{#1}}}
\newfontfamily\flowerfont[Path=../fonts/fell_flowers/]{IMFeFlow2.otf}

\begin{sloppypar}

\chapter*{\chfont{編按結語}}

\columnratio{0.5,0.5}\begin{paracol}{2}

\fontsize{11}{13}\leftfont \Large \leftcitation{א} \leftfont 余少好文．宏志博覽群書而不忘．善存藏經籍文獻備後時之用。\leftcitation{ב} \leftfont 歸主年時．受友所薦．聞道網海．\switchcolumn\fontsize{11}{13}\rightfont \Large \leftcitation{ח} \rightfont 有見粵道之危．國之封講道千言亦將就至．急之何則為？\leftcitation{ט} \rightfont 嘗聞猶太者之傳承．在其力守口述之

\end{paracol}


\columnratio{0.32,0.32,0.32}\begin{paracol}{3}

\fontsize{11}{13}\leftfont \Large 尤以吳約翰遜者 \switchcolumn[2]\fontsize{11}{13}\rightfont \Large 統．以煉千載不

\end{paracol}

\columnratio{0.32,0.32,0.32}
\begin{paracol}{3}\fontsize{11}{13}\leftfont \Large 為重．其載上之粵語講道緩緩入耳．收之藏其音頻．善妥整存．反復而嚼．受益無窮。\leftcitation{ג} \leftfont 我城我國既限．歷一四一九之不測．肺疫延年．信徒靈長屢受圍創．神州燈臺數盡指日可待．粵道之求與日俱增。\leftcitation{ד} \leftfont 觀乎社、經、法、媒、言、信、網之地．愈趨受鋤．自翔不果．授受壓力．粵道聖言亦愈漸艱難。\leftcitation{ה} \leftfont 況崇基例乎．學苑講道屢逆權勢者．其言末強受壓．舊章盡刪以存其身。講道釋數失傳．徒嘆奈何。

\switchcolumn

\fontsize{11}{13}\centerfont 
\begin{tikzpicture}
    \node (0,0) [xshift=-0.10cm, yshift=-1.0cm, opacity=0.10]{\includegraphics[width=0.30\textwidth]{../ot_frontcover.png}} ;
    \node (0,0) [xshift=+0.20cm, yshift=+2.0cm, opacity=0.10]{\includegraphics[width=0.20\textwidth]{../christ_on_cross.png}} ;
\end{tikzpicture}
\Large 

\leftcitation{ס} \centerfont 詩百又廿七載：
\leftcitation{ע} \centerfont 非耶和華建屋宇.則匠人之經營徒.
\leftcitation{פ} \centerfont 非耶和華衛城邑.則守者之儆醒徒.
\leftcitation{צ} \centerfont 余獻是卷予華人社區．願為福音流通之器．願獻斯微材為祭榮耀上帝．
\leftcitation{ק} \centerfont 阿門

\switchcolumn

\fontsize{11}{13}\rightfont \Large 滅．時越次聖殿期及當今。\leftcitation{י} \rightfont 猶太者力廣納之．筆錄以卷軸．便以傳、閱、頌、攜、守、鎖、抄、譯、釋、編，得書塔木德、密示拿等經傳．家喻戶曉．傳流若芳。\leftcitation{כ} \rightfont 猶太者文以載道．傳其口述．今我輩粵道之傳應當作如是．遂力行粵音識辨之法．載言載道．以盡忠傳粵道以待興。\leftcitation{ל} \rightfont 蒙下賜恩惠．無畏海量字音文書．既馭上帝之道．今廣及粵語講道．重駛編程之技．匯導粵音遂字稿．重塑講道現場．以傚猶太卷軸之舉便以傳流。\leftcitation{מ} \rightfont 是卷乃粵音口述傳之屬．莫通華文白話之語．

\end{paracol}

\columnratio{0.5,0.5}
\begin{paracol}{2}\fontsize{11}{13}\leftfont \Large \leftcitation{ו} \leftfont 斯殺一違儆百逆．既禁壓之．我輩聞風無奈．在所難免。\leftcitation{ז} \leftfont 另有異人例乎．以版權之名．脅網絡頻道之舉．同授礙予粵道之存流。

\switchcolumn

\fontsize{11}{13}\rightfont \Large 惟待後繼來者之傚．以譯釋傳之於神州華文地。\leftcitation{נ} \rightfont 今能排程驅馭圖靈以編彙文檔，其碼長共數千千亦無逢大礙．全蒙上帝保守。

\end{paracol}



\columnratio{1}\begin{paracol}{1}

\fontsize{11}{13}\rightfont \Large
~~~~~~~~~~~~~~~~~~~~~~~~~~~~~~~~~~~~~~~~~~~~~~~~~~~~~~~~~~~~~~~~~~~~~~~~~~~~~~~\leftcitation{ר} \rightfont 二零二三年二月一日

~~~~~~~~~~~~~~~~~~~~~~~~~~~~~~~~~~~~~~~~~~~~~~~~~~~~~~~~~~~~~~~~~~~~~~~~~~~~~~~\leftcitation{ש} \rightfont 米迦勒

~~~~~~~~~~~~~~~~~~~~~~~~~~~~~~~~~~~~~~~~~~~~~~~~~~~~~~~~~~~~~~~~~~~~~~~~~~~~~~~\leftcitation{ת} \rightfont 書於香港

\end{paracol}

\end{sloppypar}
\end{document}
